\documentclass{article}
\usepackage[utf8]{inputenc}
\usepackage{booktabs}
\usepackage{amsmath}
\usepackage{natbib}

\title{My Document}
\author{Your Name}
\date{\today}

\begin{document}

\maketitle

\section{Introduction}
This is a sample document demonstrating various \LaTeX{} features including mathematical equations, tables, and formatting. You can cite papers using the citation search feature, and they will be automatically added to your \texttt{cite.bib} file.

\section{Mathematical Expressions}
Here are some example equations:

\subsection{Quadratic Formula}
The quadratic formula is given by:
\[x = \frac{-b \pm \sqrt{b^2 - 4ac}}{2a}\]

\subsection{Matrix Example}
A sample matrix:
\[
\begin{pmatrix}
a_{11} & a_{12} & a_{13} \\
a_{21} & a_{22} & a_{23} \\
a_{31} & a_{32} & a_{33}
\end{pmatrix}
\]

\section{Tables}
Below is a sample table showing data:

\begin{table}[h]
\centering
\begin{tabular}{@{}lrrr@{}}
\toprule
\textbf{Category} & \textbf{Value 1} & \textbf{Value 2} & \textbf{Total} \\
\midrule
Item A & 10.5 & 20.3 & 30.8 \\
Item B & 15.2 & 25.1 & 40.3 \\
Item C & 22.1 & 12.7 & 34.8 \\
\bottomrule
\end{tabular}
\caption{Sample Data Table}
\label{tab:sample}
\end{table}

\section{Lists}
\subsection{Bullet Points}
\begin{itemize}
  \item First item
  \item Second item
  \item Third item with sub-items:
    \begin{itemize}
      \item Sub-item 1
      \item Sub-item 2
    \end{itemize}
\end{itemize}

\subsection{Numbered List}
\begin{enumerate}
  \item First step
  \item Second step
  \item Third step
\end{enumerate}

\section{Citations and References}
When you use the citation search feature, references will be automatically managed. For example, you might cite a paper like this: \cite{example2024}. All citations will be automatically added to your \texttt{cite.bib} file.

\section{Conclusion}
This template demonstrates basic \LaTeX{} features including mathematical formulas, tables, and lists. The bibliography system is now set up to work seamlessly with Bibby AI's citation management. Feel free to modify it for your needs!

\bibliographystyle{plain}
\bibliography{cite}

\end{document}\documentclass{article}
\usepackage[utf8]{inputenc}
\usepackage{booktabs}
\usepackage{amsmath}
\usepackage{natbib}

\title{My Document}
\author{Your Name}
\date{\today}

\begin{document}

\maketitle

\section{Introduction}
This is a sample document demonstrating various \LaTeX{} features including mathematical equations, tables, and formatting. You can cite papers using the citation search feature, and they will be automatically added to your \texttt{cite.bib} file.

\section{Mathematical Expressions}
Here are some example equations:

\subsection{Quadratic Formula}
The quadratic formula is given by:
\[x = \frac{-b \pm \sqrt{b^2 - 4ac}}{2a}\]

\subsection{Matrix Example}
A sample matrix:
\[
\begin{pmatrix}
a_{11} & a_{12} & a_{13} \\
a_{21} & a_{22} & a_{23} \\
a_{31} & a_{32} & a_{33}
\end{pmatrix}
\]

\section{Tables}
Below is a sample table showing data:

\begin{table}[h]
\centering
\begin{tabular}{@{}lrrr@{}}
\toprule
\textbf{Category} & \textbf{Value 1} & \textbf{Value 2} & \textbf{Total} \\
\midrule
Item A & 10.5 & 20.3 & 30.8 \\
Item B & 15.2 & 25.1 & 40.3 \\
Item C & 22.1 & 12.7 & 34.8 \\
\bottomrule
\end{tabular}
\caption{Sample Data Table}
\label{tab:sample}
\end{table}

\section{Lists}
\subsection{Bullet Points}
\begin{itemize}
  \item First item
  \item Second item
  \item Third item with sub-items:
    \begin{itemize}
      \item Sub-item 1
      \item Sub-item 2
    \end{itemize}
\end{itemize}

\subsection{Numbered List}
\begin{enumerate}
  \item First step
  \item Second step
  \item Third step
\end{enumerate}

\section{Citations and References}
When you use the citation search feature, references will be automatically managed. For example, you might cite a paper like this: \cite{example2024}. All citations will be automatically added to your \texttt{cite.bib} file.

\section{Conclusion}
This template demonstrates basic \LaTeX{} features including mathematical formulas, tables, and lists. The bibliography system is now set up to work seamlessly with Bibby AI's citation management. Feel free to modify it for your needs!

\bibliographystyle{plain}
\bibliography{cite}

\end{document}\documentclass{article}
\usepackage[utf8]{inputenc}
\usepackage{booktabs}
\usepackage{amsmath}
\usepackage{natbib}

\title{My Document}
\author{Your Name}
\date{\today}

\begin{document}

\maketitle

\section{Introduction}
This is a sample document demonstrating various \LaTeX{} features including mathematical equations, tables, and formatting. You can cite papers using the citation search feature, and they will be automatically added to your \texttt{cite.bib} file.

\section{Mathematical Expressions}
Here are some example equations:

\subsection{Quadratic Formula}
The quadratic formula is given by:
\[x = \frac{-b \pm \sqrt{b^2 - 4ac}}{2a}\]

\subsection{Matrix Example}
A sample matrix:
\[
\begin{pmatrix}
a_{11} & a_{12} & a_{13} \\
a_{21} & a_{22} & a_{23} \\
a_{31} & a_{32} & a_{33}
\end{pmatrix}
\]

\section{Tables}
Below is a sample table showing data:

\begin{table}[h]
\centering
\begin{tabular}{@{}lrrr@{}}
\toprule
\textbf{Category} & \textbf{Value 1} & \textbf{Value 2} & \textbf{Total} \\
\midrule
Item A & 10.5 & 20.3 & 30.8 \\
Item B & 15.2 & 25.1 & 40.3 \\
Item C & 22.1 & 12.7 & 34.8 \\
\bottomrule
\end{tabular}
\caption{Sample Data Table}
\label{tab:sample}
\end{table}

\section{Lists}
\subsection{Bullet Points}
\begin{itemize}
  \item First item
  \item Second item
  \item Third item with sub-items:
    \begin{itemize}
      \item Sub-item 1
      \item Sub-item 2
    \end{itemize}
\end{itemize}

\subsection{Numbered List}
\begin{enumerate}
  \item First step
  \item Second step
  \item Third step
\end{enumerate}

\section{Citations and References}
When you use the citation search feature, references will be automatically managed. For example, you might cite a paper like this: \cite{example2024}. All citations will be automatically added to your \texttt{cite.bib} file.

\section{Conclusion}
This template demonstrates basic \LaTeX{} features including mathematical formulas, tables, and lists. The bibliography system is now set up to work seamlessly with Bibby AI's citation management. Feel free to modify it for your needs!

\bibliographystyle{plain}
\bibliography{cite}

\end{document}\documentclass{article}
\usepackage[utf8]{inputenc}
\usepackage{booktabs}
\usepackage{amsmath}
\usepackage{natbib}

\title{My Document}
\author{Your Name}
\date{\today}

\begin{document}

\maketitle

\section{Introduction}
This is a sample document demonstrating various \LaTeX{} features including mathematical equations, tables, and formatting. You can cite papers using the citation search feature, and they will be automatically added to your \texttt{cite.bib} file.

\section{Mathematical Expressions}
Here are some example equations:

\subsection{Quadratic Formula}
The quadratic formula is given by:
\[x = \frac{-b \pm \sqrt{b^2 - 4ac}}{2a}\]

\subsection{Matrix Example}
A sample matrix:
\[
\begin{pmatrix}
a_{11} & a_{12} & a_{13} \\
a_{21} & a_{22} & a_{23} \\
a_{31} & a_{32} & a_{33}
\end{pmatrix}
\]

\section{Tables}
Below is a sample table showing data:

\begin{table}[h]
\centering
\begin{tabular}{@{}lrrr@{}}
\toprule
\textbf{Category} & \textbf{Value 1} & \textbf{Value 2} & \textbf{Total} \\
\midrule
Item A & 10.5 & 20.3 & 30.8 \\
Item B & 15.2 & 25.1 & 40.3 \\
Item C & 22.1 & 12.7 & 34.8 \\
\bottomrule
\end{tabular}
\caption{Sample Data Table}
\label{tab:sample}
\end{table}

\section{Lists}
\subsection{Bullet Points}
\begin{itemize}
  \item First item
  \item Second item
  \item Third item with sub-items:
    \begin{itemize}
      \item Sub-item 1
      \item Sub-item 2
    \end{itemize}
\end{itemize}

\subsection{Numbered List}
\begin{enumerate}
  \item First step
  \item Second step
  \item Third step
\end{enumerate}

\section{Citations and References}
When you use the citation search feature, references will be automatically managed. For example, you might cite a paper like this: \cite{example2024}. All citations will be automatically added to your \texttt{cite.bib} file.

\section{Conclusion}
This template demonstrates basic \LaTeX{} features including mathematical formulas, tables, and lists. The bibliography system is now set up to work seamlessly with Bibby AI's citation management. Feel free to modify it for your needs!

\bibliographystyle{plain}
\bibliography{cite}

\end{document}\documentclass{article}
\usepackage[utf8]{inputenc}
\usepackage{booktabs}
\usepackage{amsmath}
\usepackage{natbib}

\title{My Document}
\author{Your Name}
\date{\today}

\begin{document}

\maketitle

\section{Introduction}
This is a sample document demonstrating various \LaTeX{} features including mathematical equations, tables, and formatting. You can cite papers using the citation search feature, and they will be automatically added to your \texttt{cite.bib} file.

\section{Mathematical Expressions}
Here are some example equations:

\subsection{Quadratic Formula}
The quadratic formula is given by:
\[x = \frac{-b \pm \sqrt{b^2 - 4ac}}{2a}\]

\subsection{Matrix Example}
A sample matrix:
\[
\begin{pmatrix}
a_{11} & a_{12} & a_{13} \\
a_{21} & a_{22} & a_{23} \\
a_{31} & a_{32} & a_{33}
\end{pmatrix}
\]

\section{Tables}
Below is a sample table showing data:

\begin{table}[h]
\centering
\begin{tabular}{@{}lrrr@{}}
\toprule
\textbf{Category} & \textbf{Value 1} & \textbf{Value 2} & \textbf{Total} \\
\midrule
Item A & 10.5 & 20.3 & 30.8 \\
Item B & 15.2 & 25.1 & 40.3 \\
Item C & 22.1 & 12.7 & 34.8 \\
\bottomrule
\end{tabular}
\caption{Sample Data Table}
\label{tab:sample}
\end{table}

\section{Lists}
\subsection{Bullet Points}
\begin{itemize}
  \item First item
  \item Second item
  \item Third item with sub-items:
    \begin{itemize}
      \item Sub-item 1
      \item Sub-item 2
    \end{itemize}
\end{itemize}

\subsection{Numbered List}
\begin{enumerate}
  \item First step
  \item Second step
  \item Third step
\end{enumerate}

\section{Citations and References}
When you use the citation search feature, references will be automatically managed. For example, you might cite a paper like this: \cite{example2024}. All citations will be automatically added to your \texttt{cite.bib} file.

\section{Conclusion}
This template demonstrates basic \LaTeX{} features including mathematical formulas, tables, and lists. The bibliography system is now set up to work seamlessly with Bibby AI's citation management. Feel free to modify it for your needs!

\bibliographystyle{plain}
\bibliography{cite}

\end{document}\documentclass{article}
\usepackage[utf8]{inputenc}
\usepackage{booktabs}
\usepackage{amsmath}
\usepackage{natbib}

\title{My Document}
\author{Your Name}
\date{\today}

\begin{document}

\maketitle

\section{Introduction}
This is a sample document demonstrating various \LaTeX{} features including mathematical equations, tables, and formatting. You can cite papers using the citation search feature, and they will be automatically added to your \texttt{cite.bib} file.

\section{Mathematical Expressions}
Here are some example equations:

\subsection{Quadratic Formula}
The quadratic formula is given by:
\[x = \frac{-b \pm \sqrt{b^2 - 4ac}}{2a}\]

\subsection{Matrix Example}
A sample matrix:
\[
\begin{pmatrix}
a_{11} & a_{12} & a_{13} \\
a_{21} & a_{22} & a_{23} \\
a_{31} & a_{32} & a_{33}
\end{pmatrix}
\]

\section{Tables}
Below is a sample table showing data:

\begin{table}[h]
\centering
\begin{tabular}{@{}lrrr@{}}
\toprule
\textbf{Category} & \textbf{Value 1} & \textbf{Value 2} & \textbf{Total} \\
\midrule
Item A & 10.5 & 20.3 & 30.8 \\
Item B & 15.2 & 25.1 & 40.3 \\
Item C & 22.1 & 12.7 & 34.8 \\
\bottomrule
\end{tabular}
\caption{Sample Data Table}
\label{tab:sample}
\end{table}

\section{Lists}
\subsection{Bullet Points}
\begin{itemize}
  \item First item
  \item Second item
  \item Third item with sub-items:
    \begin{itemize}
      \item Sub-item 1
      \item Sub-item 2
    \end{itemize}
\end{itemize}

\subsection{Numbered List}
\begin{enumerate}
  \item First step
  \item Second step
  \item Third step
\end{enumerate}

\section{Citations and References}
When you use the citation search feature, references will be automatically managed. For example, you might cite a paper like this: \cite{example2024}. All citations will be automatically added to your \texttt{cite.bib} file.

\section{Conclusion}
This template demonstrates basic \LaTeX{} features including mathematical formulas, tables, and lists. The bibliography system is now set up to work seamlessly with Bibby AI's citation management. Feel free to modify it for your needs!

\bibliographystyle{plain}
\bibliography{cite}

\end{document}\documentclass{article}
\usepackage[utf8]{inputenc}
\usepackage{booktabs}
\usepackage{amsmath}
\usepackage{natbib}

\title{My Document}
\author{Your Name}
\date{\today}

\begin{document}

\maketitle

\section{Introduction}
This is a sample document demonstrating various \LaTeX{} features including mathematical equations, tables, and formatting. You can cite papers using the citation search feature, and they will be automatically added to your \texttt{cite.bib} file.

\section{Mathematical Expressions}
Here are some example equations:

\subsection{Quadratic Formula}
The quadratic formula is given by:
\[x = \frac{-b \pm \sqrt{b^2 - 4ac}}{2a}\]

\subsection{Matrix Example}
A sample matrix:
\[
\begin{pmatrix}
a_{11} & a_{12} & a_{13} \\
a_{21} & a_{22} & a_{23} \\
a_{31} & a_{32} & a_{33}
\end{pmatrix}
\]

\section{Tables}
Below is a sample table showing data:

\begin{table}[h]
\centering
\begin{tabular}{@{}lrrr@{}}
\toprule
\textbf{Category} & \textbf{Value 1} & \textbf{Value 2} & \textbf{Total} \\
\midrule
Item A & 10.5 & 20.3 & 30.8 \\
Item B & 15.2 & 25.1 & 40.3 \\
Item C & 22.1 & 12.7 & 34.8 \\
\bottomrule
\end{tabular}
\caption{Sample Data Table}
\label{tab:sample}
\end{table}

\section{Lists}
\subsection{Bullet Points}
\begin{itemize}
  \item First item
  \item Second item
  \item Third item with sub-items:
    \begin{itemize}
      \item Sub-item 1
      \item Sub-item 2
    \end{itemize}
\end{itemize}

\subsection{Numbered List}
\begin{enumerate}
  \item First step
  \item Second step
  \item Third step
\end{enumerate}

\section{Citations and References}
When you use the citation search feature, references will be automatically managed. For example, you might cite a paper like this: \cite{example2024}. All citations will be automatically added to your \texttt{cite.bib} file.

\section{Conclusion}
This template demonstrates basic \LaTeX{} features including mathematical formulas, tables, and lists. The bibliography system is now set up to work seamlessly with Bibby AI's citation management. Feel free to modify it for your needs!

\bibliographystyle{plain}
\bibliography{cite}

\end{document}\documentclass{article}
\usepackage[utf8]{inputenc}
\usepackage{booktabs}
\usepackage{amsmath}
\usepackage{natbib}

\title{My Document}
\author{Your Name}
\date{\today}

\begin{document}

\maketitle

\section{Introduction}
This is a sample document demonstrating various \LaTeX{} features including mathematical equations, tables, and formatting. You can cite papers using the citation search feature, and they will be automatically added to your \texttt{cite.bib} file.

\section{Mathematical Expressions}
Here are some example equations:

\subsection{Quadratic Formula}
The quadratic formula is given by:
\[x = \frac{-b \pm \sqrt{b^2 - 4ac}}{2a}\]

\subsection{Matrix Example}
A sample matrix:
\[
\begin{pmatrix}
a_{11} & a_{12} & a_{13} \\
a_{21} & a_{22} & a_{23} \\
a_{31} & a_{32} & a_{33}
\end{pmatrix}
\]

\section{Tables}
Below is a sample table showing data:

\begin{table}[h]
\centering
\begin{tabular}{@{}lrrr@{}}
\toprule
\textbf{Category} & \textbf{Value 1} & \textbf{Value 2} & \textbf{Total} \\
\midrule
Item A & 10.5 & 20.3 & 30.8 \\
Item B & 15.2 & 25.1 & 40.3 \\
Item C & 22.1 & 12.7 & 34.8 \\
\bottomrule
\end{tabular}
\caption{Sample Data Table}
\label{tab:sample}
\end{table}

\section{Lists}
\subsection{Bullet Points}
\begin{itemize}
  \item First item
  \item Second item
  \item Third item with sub-items:
    \begin{itemize}
      \item Sub-item 1
      \item Sub-item 2
    \end{itemize}
\end{itemize}

\subsection{Numbered List}
\begin{enumerate}
  \item First step
  \item Second step
  \item Third step
\end{enumerate}

\section{Citations and References}
When you use the citation search feature, references will be automatically managed. For example, you might cite a paper like this: \cite{example2024}. All citations will be automatically added to your \texttt{cite.bib} file.

\section{Conclusion}
This template demonstrates basic \LaTeX{} features including mathematical formulas, tables, and lists. The bibliography system is now set up to work seamlessly with Bibby AI's citation management. Feel free to modify it for your needs!

\bibliographystyle{plain}
\bibliography{cite}

\end{document}\documentclass{article}
\usepackage[utf8]{inputenc}
\usepackage{booktabs}
\usepackage{amsmath}
\usepackage{natbib}

\title{My Document}
\author{Your Name}
\date{\today}

\begin{document}

\maketitle

\section{Introduction}
This is a sample document demonstrating various \LaTeX{} features including mathematical equations, tables, and formatting. You can cite papers using the citation search feature, and they will be automatically added to your \texttt{cite.bib} file.

\section{Mathematical Expressions}
Here are some example equations:

\subsection{Quadratic Formula}
The quadratic formula is given by:
\[x = \frac{-b \pm \sqrt{b^2 - 4ac}}{2a}\]

\subsection{Matrix Example}
A sample matrix:
\[
\begin{pmatrix}
a_{11} & a_{12} & a_{13} \\
a_{21} & a_{22} & a_{23} \\
a_{31} & a_{32} & a_{33}
\end{pmatrix}
\]

\section{Tables}
Below is a sample table showing data:

\begin{table}[h]
\centering
\begin{tabular}{@{}lrrr@{}}
\toprule
\textbf{Category} & \textbf{Value 1} & \textbf{Value 2} & \textbf{Total} \\
\midrule
Item A & 10.5 & 20.3 & 30.8 \\
Item B & 15.2 & 25.1 & 40.3 \\
Item C & 22.1 & 12.7 & 34.8 \\
\bottomrule
\end{tabular}
\caption{Sample Data Table}
\label{tab:sample}
\end{table}

\section{Lists}
\subsection{Bullet Points}
\begin{itemize}
  \item First item
  \item Second item
  \item Third item with sub-items:
    \begin{itemize}
      \item Sub-item 1
      \item Sub-item 2
    \end{itemize}
\end{itemize}

\subsection{Numbered List}
\begin{enumerate}
  \item First step
  \item Second step
  \item Third step
\end{enumerate}

\section{Citations and References}
When you use the citation search feature, references will be automatically managed. For example, you might cite a paper like this: \cite{example2024}. All citations will be automatically added to your \texttt{cite.bib} file.

\section{Conclusion}
This template demonstrates basic \LaTeX{} features including mathematical formulas, tables, and lists. The bibliography system is now set up to work seamlessly with Bibby AI's citation management. Feel free to modify it for your needs!

\bibliographystyle{plain}
\bibliography{cite}

\end{document}\documentclass{article}
\usepackage[utf8]{inputenc}
\usepackage{booktabs}
\usepackage{amsmath}
\usepackage{natbib}

\title{My Document}
\author{Your Name}
\date{\today}

\begin{document}

\maketitle

\section{Introduction}
This is a sample document demonstrating various \LaTeX{} features including mathematical equations, tables, and formatting. You can cite papers using the citation search feature, and they will be automatically added to your \texttt{cite.bib} file.

\section{Mathematical Expressions}
Here are some example equations:

\subsection{Quadratic Formula}
The quadratic formula is given by:
\[x = \frac{-b \pm \sqrt{b^2 - 4ac}}{2a}\]

\subsection{Matrix Example}
A sample matrix:
\[
\begin{pmatrix}
a_{11} & a_{12} & a_{13} \\
a_{21} & a_{22} & a_{23} \\
a_{31} & a_{32} & a_{33}
\end{pmatrix}
\]

\section{Tables}
Below is a sample table showing data:

\begin{table}[h]
\centering
\begin{tabular}{@{}lrrr@{}}
\toprule
\textbf{Category} & \textbf{Value 1} & \textbf{Value 2} & \textbf{Total} \\
\midrule
Item A & 10.5 & 20.3 & 30.8 \\
Item B & 15.2 & 25.1 & 40.3 \\
Item C & 22.1 & 12.7 & 34.8 \\
\bottomrule
\end{tabular}
\caption{Sample Data Table}
\label{tab:sample}
\end{table}

\section{Lists}
\subsection{Bullet Points}
\begin{itemize}
  \item First item
  \item Second item
  \item Third item with sub-items:
    \begin{itemize}
      \item Sub-item 1
      \item Sub-item 2
    \end{itemize}
\end{itemize}

\subsection{Numbered List}
\begin{enumerate}
  \item First step
  \item Second step
  \item Third step
\end{enumerate}

\section{Citations and References}
When you use the citation search feature, references will be automatically managed. For example, you might cite a paper like this: \cite{example2024}. All citations will be automatically added to your \texttt{cite.bib} file.

\section{Conclusion}
This template demonstrates basic \LaTeX{} features including mathematical formulas, tables, and lists. The bibliography system is now set up to work seamlessly with Bibby AI's citation management. Feel free to modify it for your needs!

\bibliographystyle{plain}
\bibliography{cite}

\end{document}\documentclass{article}
\usepackage[utf8]{inputenc}
\usepackage{booktabs}
\usepackage{amsmath}
\usepackage{natbib}

\title{My Document}
\author{Your Name}
\date{\today}

\begin{document}

\maketitle

\section{Introduction}
This is a sample document demonstrating various \LaTeX{} features including mathematical equations, tables, and formatting. You can cite papers using the citation search feature, and they will be automatically added to your \texttt{cite.bib} file.

\section{Mathematical Expressions}
Here are some example equations:

\subsection{Quadratic Formula}
The quadratic formula is given by:
\[x = \frac{-b \pm \sqrt{b^2 - 4ac}}{2a}\]

\subsection{Matrix Example}
A sample matrix:
\[
\begin{pmatrix}
a_{11} & a_{12} & a_{13} \\
a_{21} & a_{22} & a_{23} \\
a_{31} & a_{32} & a_{33}
\end{pmatrix}
\]

\section{Tables}
Below is a sample table showing data:

\begin{table}[h]
\centering
\begin{tabular}{@{}lrrr@{}}
\toprule
\textbf{Category} & \textbf{Value 1} & \textbf{Value 2} & \textbf{Total} \\
\midrule
Item A & 10.5 & 20.3 & 30.8 \\
Item B & 15.2 & 25.1 & 40.3 \\
Item C & 22.1 & 12.7 & 34.8 \\
\bottomrule
\end{tabular}
\caption{Sample Data Table}
\label{tab:sample}
\end{table}

\section{Lists}
\subsection{Bullet Points}
\begin{itemize}
  \item First item
  \item Second item
  \item Third item with sub-items:
    \begin{itemize}
      \item Sub-item 1
      \item Sub-item 2
    \end{itemize}
\end{itemize}

\subsection{Numbered List}
\begin{enumerate}
  \item First step
  \item Second step
  \item Third step
\end{enumerate}

\section{Citations and References}
When you use the citation search feature, references will be automatically managed. For example, you might cite a paper like this: \cite{example2024}. All citations will be automatically added to your \texttt{cite.bib} file.

\section{Conclusion}
This template demonstrates basic \LaTeX{} features including mathematical formulas, tables, and lists. The bibliography system is now set up to work seamlessly with Bibby AI's citation management. Feel free to modify it for your needs!

\bibliographystyle{plain}
\bibliography{cite}

\end{document}\documentclass{article}
\usepackage[utf8]{inputenc}
\usepackage{booktabs}
\usepackage{amsmath}
\usepackage{natbib}

\title{My Document}
\author{Your Name}
\date{\today}

\begin{document}

\maketitle

\section{Introduction}
This is a sample document demonstrating various \LaTeX{} features including mathematical equations, tables, and formatting. You can cite papers using the citation search feature, and they will be automatically added to your \texttt{cite.bib} file.

\section{Mathematical Expressions}
Here are some example equations:

\subsection{Quadratic Formula}
The quadratic formula is given by:
\[x = \frac{-b \pm \sqrt{b^2 - 4ac}}{2a}\]

\subsection{Matrix Example}
A sample matrix:
\[
\begin{pmatrix}
a_{11} & a_{12} & a_{13} \\
a_{21} & a_{22} & a_{23} \\
a_{31} & a_{32} & a_{33}
\end{pmatrix}
\]

\section{Tables}
Below is a sample table showing data:

\begin{table}[h]
\centering
\begin{tabular}{@{}lrrr@{}}
\toprule
\textbf{Category} & \textbf{Value 1} & \textbf{Value 2} & \textbf{Total} \\
\midrule
Item A & 10.5 & 20.3 & 30.8 \\
Item B & 15.2 & 25.1 & 40.3 \\
Item C & 22.1 & 12.7 & 34.8 \\
\bottomrule
\end{tabular}
\caption{Sample Data Table}
\label{tab:sample}
\end{table}

\section{Lists}
\subsection{Bullet Points}
\begin{itemize}
  \item First item
  \item Second item
  \item Third item with sub-items:
    \begin{itemize}
      \item Sub-item 1
      \item Sub-item 2
    \end{itemize}
\end{itemize}

\subsection{Numbered List}
\begin{enumerate}
  \item First step
  \item Second step
  \item Third step
\end{enumerate}

\section{Citations and References}
When you use the citation search feature, references will be automatically managed. For example, you might cite a paper like this: \cite{example2024}. All citations will be automatically added to your \texttt{cite.bib} file.

\section{Conclusion}
This template demonstrates basic \LaTeX{} features including mathematical formulas, tables, and lists. The bibliography system is now set up to work seamlessly with Bibby AI's citation management. Feel free to modify it for your needs!

\bibliographystyle{plain}
\bibliography{cite}

\end{document}\documentclass{article}
\usepackage[utf8]{inputenc}
\usepackage{booktabs}
\usepackage{amsmath}
\usepackage{natbib}

\title{My Document}
\author{Your Name}
\date{\today}

\begin{document}

\maketitle

\section{Introduction}
This is a sample document demonstrating various \LaTeX{} features including mathematical equations, tables, and formatting. You can cite papers using the citation search feature, and they will be automatically added to your \texttt{cite.bib} file.

\section{Mathematical Expressions}
Here are some example equations:

\subsection{Quadratic Formula}
The quadratic formula is given by:
\[x = \frac{-b \pm \sqrt{b^2 - 4ac}}{2a}\]

\subsection{Matrix Example}
A sample matrix:
\[
\begin{pmatrix}
a_{11} & a_{12} & a_{13} \\
a_{21} & a_{22} & a_{23} \\
a_{31} & a_{32} & a_{33}
\end{pmatrix}
\]

\section{Tables}
Below is a sample table showing data:

\begin{table}[h]
\centering
\begin{tabular}{@{}lrrr@{}}
\toprule
\textbf{Category} & \textbf{Value 1} & \textbf{Value 2} & \textbf{Total} \\
\midrule
Item A & 10.5 & 20.3 & 30.8 \\
Item B & 15.2 & 25.1 & 40.3 \\
Item C & 22.1 & 12.7 & 34.8 \\
\bottomrule
\end{tabular}
\caption{Sample Data Table}
\label{tab:sample}
\end{table}

\section{Lists}
\subsection{Bullet Points}
\begin{itemize}
  \item First item
  \item Second item
  \item Third item with sub-items:
    \begin{itemize}
      \item Sub-item 1
      \item Sub-item 2
    \end{itemize}
\end{itemize}

\subsection{Numbered List}
\begin{enumerate}
  \item First step
  \item Second step
  \item Third step
\end{enumerate}

\section{Citations and References}
When you use the citation search feature, references will be automatically managed. For example, you might cite a paper like this: \cite{example2024}. All citations will be automatically added to your \texttt{cite.bib} file.

\section{Conclusion}
This template demonstrates basic \LaTeX{} features including mathematical formulas, tables, and lists. The bibliography system is now set up to work seamlessly with Bibby AI's citation management. Feel free to modify it for your needs!

\bibliographystyle{plain}
\bibliography{cite}

\end{document}\documentclass{article}
\usepackage[utf8]{inputenc}
\usepackage{booktabs}
\usepackage{amsmath}
\usepackage{natbib}

\title{My Document}
\author{Your Name}
\date{\today}

\begin{document}

\maketitle

\section{Introduction}
This is a sample document demonstrating various \LaTeX{} features including mathematical equations, tables, and formatting. You can cite papers using the citation search feature, and they will be automatically added to your \texttt{cite.bib} file.

\section{Mathematical Expressions}
Here are some example equations:

\subsection{Quadratic Formula}
The quadratic formula is given by:
\[x = \frac{-b \pm \sqrt{b^2 - 4ac}}{2a}\]

\subsection{Matrix Example}
A sample matrix:
\[
\begin{pmatrix}
a_{11} & a_{12} & a_{13} \\
a_{21} & a_{22} & a_{23} \\
a_{31} & a_{32} & a_{33}
\end{pmatrix}
\]

\section{Tables}
Below is a sample table showing data:

\begin{table}[h]
\centering
\begin{tabular}{@{}lrrr@{}}
\toprule
\textbf{Category} & \textbf{Value 1} & \textbf{Value 2} & \textbf{Total} \\
\midrule
Item A & 10.5 & 20.3 & 30.8 \\
Item B & 15.2 & 25.1 & 40.3 \\
Item C & 22.1 & 12.7 & 34.8 \\
\bottomrule
\end{tabular}
\caption{Sample Data Table}
\label{tab:sample}
\end{table}

\section{Lists}
\subsection{Bullet Points}
\begin{itemize}
  \item First item
  \item Second item
  \item Third item with sub-items:
    \begin{itemize}
      \item Sub-item 1
      \item Sub-item 2
    \end{itemize}
\end{itemize}

\subsection{Numbered List}
\begin{enumerate}
  \item First step
  \item Second step
  \item Third step
\end{enumerate}

\section{Citations and References}
When you use the citation search feature, references will be automatically managed. For example, you might cite a paper like this: \cite{example2024}. All citations will be automatically added to your \texttt{cite.bib} file.

\section{Conclusion}
This template demonstrates basic \LaTeX{} features including mathematical formulas, tables, and lists. The bibliography system is now set up to work seamlessly with Bibby AI's citation management. Feel free to modify it for your needs!

\bibliographystyle{plain}
\bibliography{cite}

\end{document}\documentclass{article}
\usepackage[utf8]{inputenc}
\usepackage{booktabs}
\usepackage{amsmath}
\usepackage{natbib}

\title{My Document}
\author{Your Name}
\date{\today}

\begin{document}

\maketitle

\section{Introduction}
This is a sample document demonstrating various \LaTeX{} features including mathematical equations, tables, and formatting. You can cite papers using the citation search feature, and they will be automatically added to your \texttt{cite.bib} file.

\section{Mathematical Expressions}
Here are some example equations:

\subsection{Quadratic Formula}
The quadratic formula is given by:
\[x = \frac{-b \pm \sqrt{b^2 - 4ac}}{2a}\]

\subsection{Matrix Example}
A sample matrix:
\[
\begin{pmatrix}
a_{11} & a_{12} & a_{13} \\
a_{21} & a_{22} & a_{23} \\
a_{31} & a_{32} & a_{33}
\end{pmatrix}
\]

\section{Tables}
Below is a sample table showing data:

\begin{table}[h]
\centering
\begin{tabular}{@{}lrrr@{}}
\toprule
\textbf{Category} & \textbf{Value 1} & \textbf{Value 2} & \textbf{Total} \\
\midrule
Item A & 10.5 & 20.3 & 30.8 \\
Item B & 15.2 & 25.1 & 40.3 \\
Item C & 22.1 & 12.7 & 34.8 \\
\bottomrule
\end{tabular}
\caption{Sample Data Table}
\label{tab:sample}
\end{table}

\section{Lists}
\subsection{Bullet Points}
\begin{itemize}
  \item First item
  \item Second item
  \item Third item with sub-items:
    \begin{itemize}
      \item Sub-item 1
      \item Sub-item 2
    \end{itemize}
\end{itemize}

\subsection{Numbered List}
\begin{enumerate}
  \item First step
  \item Second step
  \item Third step
\end{enumerate}

\section{Citations and References}
When you use the citation search feature, references will be automatically managed. For example, you might cite a paper like this: \cite{example2024}. All citations will be automatically added to your \texttt{cite.bib} file.

\section{Conclusion}
This template demonstrates basic \LaTeX{} features including mathematical formulas, tables, and lists. The bibliography system is now set up to work seamlessly with Bibby AI's citation management. Feel free to modify it for your needs!

\bibliographystyle{plain}
\bibliography{cite}

\end{document}\documentclass{article}
\usepackage[utf8]{inputenc}
\usepackage{booktabs}
\usepackage{amsmath}
\usepackage{natbib}

\title{My Document}
\author{Your Name}
\date{\today}

\begin{document}

\maketitle

\section{Introduction}
This is a sample document demonstrating various \LaTeX{} features including mathematical equations, tables, and formatting. You can cite papers using the citation search feature, and they will be automatically added to your \texttt{cite.bib} file.

\section{Mathematical Expressions}
Here are some example equations:

\subsection{Quadratic Formula}
The quadratic formula is given by:
\[x = \frac{-b \pm \sqrt{b^2 - 4ac}}{2a}\]

\subsection{Matrix Example}
A sample matrix:
\[
\begin{pmatrix}
a_{11} & a_{12} & a_{13} \\
a_{21} & a_{22} & a_{23} \\
a_{31} & a_{32} & a_{33}
\end{pmatrix}
\]

\section{Tables}
Below is a sample table showing data:

\begin{table}[h]
\centering
\begin{tabular}{@{}lrrr@{}}
\toprule
\textbf{Category} & \textbf{Value 1} & \textbf{Value 2} & \textbf{Total} \\
\midrule
Item A & 10.5 & 20.3 & 30.8 \\
Item B & 15.2 & 25.1 & 40.3 \\
Item C & 22.1 & 12.7 & 34.8 \\
\bottomrule
\end{tabular}
\caption{Sample Data Table}
\label{tab:sample}
\end{table}

\section{Lists}
\subsection{Bullet Points}
\begin{itemize}
  \item First item
  \item Second item
  \item Third item with sub-items:
    \begin{itemize}
      \item Sub-item 1
      \item Sub-item 2
    \end{itemize}
\end{itemize}

\subsection{Numbered List}
\begin{enumerate}
  \item First step
  \item Second step
  \item Third step
\end{enumerate}

\section{Citations and References}
When you use the citation search feature, references will be automatically managed. For example, you might cite a paper like this: \cite{example2024}. All citations will be automatically added to your \texttt{cite.bib} file.

\section{Conclusion}
This template demonstrates basic \LaTeX{} features including mathematical formulas, tables, and lists. The bibliography system is now set up to work seamlessly with Bibby AI's citation management. Feel free to modify it for your needs!

\bibliographystyle{plain}
\bibliography{cite}

\end{document}\documentclass{article}
\usepackage[utf8]{inputenc}
\usepackage{booktabs}
\usepackage{amsmath}
\usepackage{natbib}

\title{My Document}
\author{Your Name}
\date{\today}

\begin{document}

\maketitle

\section{Introduction}
This is a sample document demonstrating various \LaTeX{} features including mathematical equations, tables, and formatting. You can cite papers using the citation search feature, and they will be automatically added to your \texttt{cite.bib} file.

\section{Mathematical Expressions}
Here are some example equations:

\subsection{Quadratic Formula}
The quadratic formula is given by:
\[x = \frac{-b \pm \sqrt{b^2 - 4ac}}{2a}\]

\subsection{Matrix Example}
A sample matrix:
\[
\begin{pmatrix}
a_{11} & a_{12} & a_{13} \\
a_{21} & a_{22} & a_{23} \\
a_{31} & a_{32} & a_{33}
\end{pmatrix}
\]

\section{Tables}
Below is a sample table showing data:

\begin{table}[h]
\centering
\begin{tabular}{@{}lrrr@{}}
\toprule
\textbf{Category} & \textbf{Value 1} & \textbf{Value 2} & \textbf{Total} \\
\midrule
Item A & 10.5 & 20.3 & 30.8 \\
Item B & 15.2 & 25.1 & 40.3 \\
Item C & 22.1 & 12.7 & 34.8 \\
\bottomrule
\end{tabular}
\caption{Sample Data Table}
\label{tab:sample}
\end{table}

\section{Lists}
\subsection{Bullet Points}
\begin{itemize}
  \item First item
  \item Second item
  \item Third item with sub-items:
    \begin{itemize}
      \item Sub-item 1
      \item Sub-item 2
    \end{itemize}
\end{itemize}

\subsection{Numbered List}
\begin{enumerate}
  \item First step
  \item Second step
  \item Third step
\end{enumerate}

\section{Citations and References}
When you use the citation search feature, references will be automatically managed. For example, you might cite a paper like this: \cite{example2024}. All citations will be automatically added to your \texttt{cite.bib} file.

\section{Conclusion}
This template demonstrates basic \LaTeX{} features including mathematical formulas, tables, and lists. The bibliography system is now set up to work seamlessly with Bibby AI's citation management. Feel free to modify it for your needs!

\bibliographystyle{plain}
\bibliography{cite}

\end{document}\documentclass{article}
\usepackage[utf8]{inputenc}
\usepackage{booktabs}
\usepackage{amsmath}
\usepackage{natbib}

\title{My Document}
\author{Your Name}
\date{\today}

\begin{document}

\maketitle

\section{Introduction}
This is a sample document demonstrating various \LaTeX{} features including mathematical equations, tables, and formatting. You can cite papers using the citation search feature, and they will be automatically added to your \texttt{cite.bib} file.

\section{Mathematical Expressions}
Here are some example equations:

\subsection{Quadratic Formula}
The quadratic formula is given by:
\[x = \frac{-b \pm \sqrt{b^2 - 4ac}}{2a}\]

\subsection{Matrix Example}
A sample matrix:
\[
\begin{pmatrix}
a_{11} & a_{12} & a_{13} \\
a_{21} & a_{22} & a_{23} \\
a_{31} & a_{32} & a_{33}
\end{pmatrix}
\]

\section{Tables}
Below is a sample table showing data:

\begin{table}[h]
\centering
\begin{tabular}{@{}lrrr@{}}
\toprule
\textbf{Category} & \textbf{Value 1} & \textbf{Value 2} & \textbf{Total} \\
\midrule
Item A & 10.5 & 20.3 & 30.8 \\
Item B & 15.2 & 25.1 & 40.3 \\
Item C & 22.1 & 12.7 & 34.8 \\
\bottomrule
\end{tabular}
\caption{Sample Data Table}
\label{tab:sample}
\end{table}

\section{Lists}
\subsection{Bullet Points}
\begin{itemize}
  \item First item
  \item Second item
  \item Third item with sub-items:
    \begin{itemize}
      \item Sub-item 1
      \item Sub-item 2
    \end{itemize}
\end{itemize}

\subsection{Numbered List}
\begin{enumerate}
  \item First step
  \item Second step
  \item Third step
\end{enumerate}

\section{Citations and References}
When you use the citation search feature, references will be automatically managed. For example, you might cite a paper like this: \cite{example2024}. All citations will be automatically added to your \texttt{cite.bib} file.

\section{Conclusion}
This template demonstrates basic \LaTeX{} features including mathematical formulas, tables, and lists. The bibliography system is now set up to work seamlessly with Bibby AI's citation management. Feel free to modify it for your needs!

\bibliographystyle{plain}
\bibliography{cite}

\end{document}\documentclass{article}
\usepackage[utf8]{inputenc}
\usepackage{booktabs}
\usepackage{amsmath}
\usepackage{natbib}

\title{My Document}
\author{Your Name}
\date{\today}

\begin{document}

\maketitle

\section{Introduction}
This is a sample document demonstrating various \LaTeX{} features including mathematical equations, tables, and formatting. You can cite papers using the citation search feature, and they will be automatically added to your \texttt{cite.bib} file.

\section{Mathematical Expressions}
Here are some example equations:

\subsection{Quadratic Formula}
The quadratic formula is given by:
\[x = \frac{-b \pm \sqrt{b^2 - 4ac}}{2a}\]

\subsection{Matrix Example}
A sample matrix:
\[
\begin{pmatrix}
a_{11} & a_{12} & a_{13} \\
a_{21} & a_{22} & a_{23} \\
a_{31} & a_{32} & a_{33}
\end{pmatrix}
\]

\section{Tables}
Below is a sample table showing data:

\begin{table}[h]
\centering
\begin{tabular}{@{}lrrr@{}}
\toprule
\textbf{Category} & \textbf{Value 1} & \textbf{Value 2} & \textbf{Total} \\
\midrule
Item A & 10.5 & 20.3 & 30.8 \\
Item B & 15.2 & 25.1 & 40.3 \\
Item C & 22.1 & 12.7 & 34.8 \\
\bottomrule
\end{tabular}
\caption{Sample Data Table}
\label{tab:sample}
\end{table}

\section{Lists}
\subsection{Bullet Points}
\begin{itemize}
  \item First item
  \item Second item
  \item Third item with sub-items:
    \begin{itemize}
      \item Sub-item 1
      \item Sub-item 2
    \end{itemize}
\end{itemize}

\subsection{Numbered List}
\begin{enumerate}
  \item First step
  \item Second step
  \item Third step
\end{enumerate}

\section{Citations and References}
When you use the citation search feature, references will be automatically managed. For example, you might cite a paper like this: \cite{example2024}. All citations will be automatically added to your \texttt{cite.bib} file.

\section{Conclusion}
This template demonstrates basic \LaTeX{} features including mathematical formulas, tables, and lists. The bibliography system is now set up to work seamlessly with Bibby AI's citation management. Feel free to modify it for your needs!

\bibliographystyle{plain}
\bibliography{cite}

\end{document}\documentclass{article}
\usepackage[utf8]{inputenc}
\usepackage{booktabs}
\usepackage{amsmath}
\usepackage{natbib}

\title{My Document}
\author{Your Name}
\date{\today}

\begin{document}

\maketitle

\section{Introduction}
This is a sample document demonstrating various \LaTeX{} features including mathematical equations, tables, and formatting. You can cite papers using the citation search feature, and they will be automatically added to your \texttt{cite.bib} file.

\section{Mathematical Expressions}
Here are some example equations:

\subsection{Quadratic Formula}
The quadratic formula is given by:
\[x = \frac{-b \pm \sqrt{b^2 - 4ac}}{2a}\]

\subsection{Matrix Example}
A sample matrix:
\[
\begin{pmatrix}
a_{11} & a_{12} & a_{13} \\
a_{21} & a_{22} & a_{23} \\
a_{31} & a_{32} & a_{33}
\end{pmatrix}
\]

\section{Tables}
Below is a sample table showing data:

\begin{table}[h]
\centering
\begin{tabular}{@{}lrrr@{}}
\toprule
\textbf{Category} & \textbf{Value 1} & \textbf{Value 2} & \textbf{Total} \\
\midrule
Item A & 10.5 & 20.3 & 30.8 \\
Item B & 15.2 & 25.1 & 40.3 \\
Item C & 22.1 & 12.7 & 34.8 \\
\bottomrule
\end{tabular}
\caption{Sample Data Table}
\label{tab:sample}
\end{table}

\section{Lists}
\subsection{Bullet Points}
\begin{itemize}
  \item First item
  \item Second item
  \item Third item with sub-items:
    \begin{itemize}
      \item Sub-item 1
      \item Sub-item 2
    \end{itemize}
\end{itemize}

\subsection{Numbered List}
\begin{enumerate}
  \item First step
  \item Second step
  \item Third step
\end{enumerate}

\section{Citations and References}
When you use the citation search feature, references will be automatically managed. For example, you might cite a paper like this: \cite{example2024}. All citations will be automatically added to your \texttt{cite.bib} file.

\section{Conclusion}
This template demonstrates basic \LaTeX{} features including mathematical formulas, tables, and lists. The bibliography system is now set up to work seamlessly with Bibby AI's citation management. Feel free to modify it for your needs!

\bibliographystyle{plain}
\bibliography{cite}

\end{document}\documentclass{article}
\usepackage[utf8]{inputenc}
\usepackage{booktabs}
\usepackage{amsmath}
\usepackage{natbib}

\title{My Document}
\author{Your Name}
\date{\today}

\begin{document}

\maketitle

\section{Introduction}
This is a sample document demonstrating various \LaTeX{} features including mathematical equations, tables, and formatting. You can cite papers using the citation search feature, and they will be automatically added to your \texttt{cite.bib} file.

\section{Mathematical Expressions}
Here are some example equations:

\subsection{Quadratic Formula}
The quadratic formula is given by:
\[x = \frac{-b \pm \sqrt{b^2 - 4ac}}{2a}\]

\subsection{Matrix Example}
A sample matrix:
\[
\begin{pmatrix}
a_{11} & a_{12} & a_{13} \\
a_{21} & a_{22} & a_{23} \\
a_{31} & a_{32} & a_{33}
\end{pmatrix}
\]

\section{Tables}
Below is a sample table showing data:

\begin{table}[h]
\centering
\begin{tabular}{@{}lrrr@{}}
\toprule
\textbf{Category} & \textbf{Value 1} & \textbf{Value 2} & \textbf{Total} \\
\midrule
Item A & 10.5 & 20.3 & 30.8 \\
Item B & 15.2 & 25.1 & 40.3 \\
Item C & 22.1 & 12.7 & 34.8 \\
\bottomrule
\end{tabular}
\caption{Sample Data Table}
\label{tab:sample}
\end{table}

\section{Lists}
\subsection{Bullet Points}
\begin{itemize}
  \item First item
  \item Second item
  \item Third item with sub-items:
    \begin{itemize}
      \item Sub-item 1
      \item Sub-item 2
    \end{itemize}
\end{itemize}

\subsection{Numbered List}
\begin{enumerate}
  \item First step
  \item Second step
  \item Third step
\end{enumerate}

\section{Citations and References}
When you use the citation search feature, references will be automatically managed. For example, you might cite a paper like this: \cite{example2024}. All citations will be automatically added to your \texttt{cite.bib} file.

\section{Conclusion}
This template demonstrates basic \LaTeX{} features including mathematical formulas, tables, and lists. The bibliography system is now set up to work seamlessly with Bibby AI's citation management. Feel free to modify it for your needs!

\bibliographystyle{plain}
\bibliography{cite}

\end{document}\documentclass{article}
\usepackage[utf8]{inputenc}
\usepackage{booktabs}
\usepackage{amsmath}
\usepackage{natbib}

\title{My Document}
\author{Your Name}
\date{\today}

\begin{document}

\maketitle

\section{Introduction}
This is a sample document demonstrating various \LaTeX{} features including mathematical equations, tables, and formatting. You can cite papers using the citation search feature, and they will be automatically added to your \texttt{cite.bib} file.

\section{Mathematical Expressions}
Here are some example equations:

\subsection{Quadratic Formula}
The quadratic formula is given by:
\[x = \frac{-b \pm \sqrt{b^2 - 4ac}}{2a}\]

\subsection{Matrix Example}
A sample matrix:
\[
\begin{pmatrix}
a_{11} & a_{12} & a_{13} \\
a_{21} & a_{22} & a_{23} \\
a_{31} & a_{32} & a_{33}
\end{pmatrix}
\]

\section{Tables}
Below is a sample table showing data:

\begin{table}[h]
\centering
\begin{tabular}{@{}lrrr@{}}
\toprule
\textbf{Category} & \textbf{Value 1} & \textbf{Value 2} & \textbf{Total} \\
\midrule
Item A & 10.5 & 20.3 & 30.8 \\
Item B & 15.2 & 25.1 & 40.3 \\
Item C & 22.1 & 12.7 & 34.8 \\
\bottomrule
\end{tabular}
\caption{Sample Data Table}
\label{tab:sample}
\end{table}

\section{Lists}
\subsection{Bullet Points}
\begin{itemize}
  \item First item
  \item Second item
  \item Third item with sub-items:
    \begin{itemize}
      \item Sub-item 1
      \item Sub-item 2
    \end{itemize}
\end{itemize}

\subsection{Numbered List}
\begin{enumerate}
  \item First step
  \item Second step
  \item Third step
\end{enumerate}

\section{Citations and References}
When you use the citation search feature, references will be automatically managed. For example, you might cite a paper like this: \cite{example2024}. All citations will be automatically added to your \texttt{cite.bib} file.

\section{Conclusion}
This template demonstrates basic \LaTeX{} features including mathematical formulas, tables, and lists. The bibliography system is now set up to work seamlessly with Bibby AI's citation management. Feel free to modify it for your needs!

\bibliographystyle{plain}
\bibliography{cite}

\end{document}\documentclass{article}
\usepackage[utf8]{inputenc}
\usepackage{booktabs}
\usepackage{amsmath}
\usepackage{natbib}

\title{My Document}
\author{Your Name}
\date{\today}

\begin{document}

\maketitle

\section{Introduction}
This is a sample document demonstrating various \LaTeX{} features including mathematical equations, tables, and formatting. You can cite papers using the citation search feature, and they will be automatically added to your \texttt{cite.bib} file.

\section{Mathematical Expressions}
Here are some example equations:

\subsection{Quadratic Formula}
The quadratic formula is given by:
\[x = \frac{-b \pm \sqrt{b^2 - 4ac}}{2a}\]

\subsection{Matrix Example}
A sample matrix:
\[
\begin{pmatrix}
a_{11} & a_{12} & a_{13} \\
a_{21} & a_{22} & a_{23} \\
a_{31} & a_{32} & a_{33}
\end{pmatrix}
\]

\section{Tables}
Below is a sample table showing data:

\begin{table}[h]
\centering
\begin{tabular}{@{}lrrr@{}}
\toprule
\textbf{Category} & \textbf{Value 1} & \textbf{Value 2} & \textbf{Total} \\
\midrule
Item A & 10.5 & 20.3 & 30.8 \\
Item B & 15.2 & 25.1 & 40.3 \\
Item C & 22.1 & 12.7 & 34.8 \\
\bottomrule
\end{tabular}
\caption{Sample Data Table}
\label{tab:sample}
\end{table}

\section{Lists}
\subsection{Bullet Points}
\begin{itemize}
  \item First item
  \item Second item
  \item Third item with sub-items:
    \begin{itemize}
      \item Sub-item 1
      \item Sub-item 2
    \end{itemize}
\end{itemize}

\subsection{Numbered List}
\begin{enumerate}
  \item First step
  \item Second step
  \item Third step
\end{enumerate}

\section{Citations and References}
When you use the citation search feature, references will be automatically managed. For example, you might cite a paper like this: \cite{example2024}. All citations will be automatically added to your \texttt{cite.bib} file.

\section{Conclusion}
This template demonstrates basic \LaTeX{} features including mathematical formulas, tables, and lists. The bibliography system is now set up to work seamlessly with Bibby AI's citation management. Feel free to modify it for your needs!

\bibliographystyle{plain}
\bibliography{cite}

\end{document}\documentclass{article}
\usepackage[utf8]{inputenc}
\usepackage{booktabs}
\usepackage{amsmath}
\usepackage{natbib}

\title{My Document}
\author{Your Name}
\date{\today}

\begin{document}

\maketitle

\section{Introduction}
This is a sample document demonstrating various \LaTeX{} features including mathematical equations, tables, and formatting. You can cite papers using the citation search feature, and they will be automatically added to your \texttt{cite.bib} file.

\section{Mathematical Expressions}
Here are some example equations:

\subsection{Quadratic Formula}
The quadratic formula is given by:
\[x = \frac{-b \pm \sqrt{b^2 - 4ac}}{2a}\]

\subsection{Matrix Example}
A sample matrix:
\[
\begin{pmatrix}
a_{11} & a_{12} & a_{13} \\
a_{21} & a_{22} & a_{23} \\
a_{31} & a_{32} & a_{33}
\end{pmatrix}
\]

\section{Tables}
Below is a sample table showing data:

\begin{table}[h]
\centering
\begin{tabular}{@{}lrrr@{}}
\toprule
\textbf{Category} & \textbf{Value 1} & \textbf{Value 2} & \textbf{Total} \\
\midrule
Item A & 10.5 & 20.3 & 30.8 \\
Item B & 15.2 & 25.1 & 40.3 \\
Item C & 22.1 & 12.7 & 34.8 \\
\bottomrule
\end{tabular}
\caption{Sample Data Table}
\label{tab:sample}
\end{table}

\section{Lists}
\subsection{Bullet Points}
\begin{itemize}
  \item First item
  \item Second item
  \item Third item with sub-items:
    \begin{itemize}
      \item Sub-item 1
      \item Sub-item 2
    \end{itemize}
\end{itemize}

\subsection{Numbered List}
\begin{enumerate}
  \item First step
  \item Second step
  \item Third step
\end{enumerate}

\section{Citations and References}
When you use the citation search feature, references will be automatically managed. For example, you might cite a paper like this: \cite{example2024}. All citations will be automatically added to your \texttt{cite.bib} file.

\section{Conclusion}
This template demonstrates basic \LaTeX{} features including mathematical formulas, tables, and lists. The bibliography system is now set up to work seamlessly with Bibby AI's citation management. Feel free to modify it for your needs!

\bibliographystyle{plain}
\bibliography{cite}

\end{document}\documentclass{article}
\usepackage[utf8]{inputenc}
\usepackage{booktabs}
\usepackage{amsmath}
\usepackage{natbib}

\title{My Document}
\author{Your Name}
\date{\today}

\begin{document}

\maketitle

\section{Introduction}
This is a sample document demonstrating various \LaTeX{} features including mathematical equations, tables, and formatting. You can cite papers using the citation search feature, and they will be automatically added to your \texttt{cite.bib} file.

\section{Mathematical Expressions}
Here are some example equations:

\subsection{Quadratic Formula}
The quadratic formula is given by:
\[x = \frac{-b \pm \sqrt{b^2 - 4ac}}{2a}\]

\subsection{Matrix Example}
A sample matrix:
\[
\begin{pmatrix}
a_{11} & a_{12} & a_{13} \\
a_{21} & a_{22} & a_{23} \\
a_{31} & a_{32} & a_{33}
\end{pmatrix}
\]

\section{Tables}
Below is a sample table showing data:

\begin{table}[h]
\centering
\begin{tabular}{@{}lrrr@{}}
\toprule
\textbf{Category} & \textbf{Value 1} & \textbf{Value 2} & \textbf{Total} \\
\midrule
Item A & 10.5 & 20.3 & 30.8 \\
Item B & 15.2 & 25.1 & 40.3 \\
Item C & 22.1 & 12.7 & 34.8 \\
\bottomrule
\end{tabular}
\caption{Sample Data Table}
\label{tab:sample}
\end{table}

\section{Lists}
\subsection{Bullet Points}
\begin{itemize}
  \item First item
  \item Second item
  \item Third item with sub-items:
    \begin{itemize}
      \item Sub-item 1
      \item Sub-item 2
    \end{itemize}
\end{itemize}

\subsection{Numbered List}
\begin{enumerate}
  \item First step
  \item Second step
  \item Third step
\end{enumerate}

\section{Citations and References}
When you use the citation search feature, references will be automatically managed. For example, you might cite a paper like this: \cite{example2024}. All citations will be automatically added to your \texttt{cite.bib} file.

\section{Conclusion}
This template demonstrates basic \LaTeX{} features including mathematical formulas, tables, and lists. The bibliography system is now set up to work seamlessly with Bibby AI's citation management. Feel free to modify it for your needs!

\bibliographystyle{plain}
\bibliography{cite}

\end{document}\documentclass{article}
\usepackage[utf8]{inputenc}
\usepackage{booktabs}
\usepackage{amsmath}
\usepackage{natbib}

\title{My Document}
\author{Your Name}
\date{\today}

\begin{document}

\maketitle

\section{Introduction}
This is a sample document demonstrating various \LaTeX{} features including mathematical equations, tables, and formatting. You can cite papers using the citation search feature, and they will be automatically added to your \texttt{cite.bib} file.

\section{Mathematical Expressions}
Here are some example equations:

\subsection{Quadratic Formula}
The quadratic formula is given by:
\[x = \frac{-b \pm \sqrt{b^2 - 4ac}}{2a}\]

\subsection{Matrix Example}
A sample matrix:
\[
\begin{pmatrix}
a_{11} & a_{12} & a_{13} \\
a_{21} & a_{22} & a_{23} \\
a_{31} & a_{32} & a_{33}
\end{pmatrix}
\]

\section{Tables}
Below is a sample table showing data:

\begin{table}[h]
\centering
\begin{tabular}{@{}lrrr@{}}
\toprule
\textbf{Category} & \textbf{Value 1} & \textbf{Value 2} & \textbf{Total} \\
\midrule
Item A & 10.5 & 20.3 & 30.8 \\
Item B & 15.2 & 25.1 & 40.3 \\
Item C & 22.1 & 12.7 & 34.8 \\
\bottomrule
\end{tabular}
\caption{Sample Data Table}
\label{tab:sample}
\end{table}

\section{Lists}
\subsection{Bullet Points}
\begin{itemize}
  \item First item
  \item Second item
  \item Third item with sub-items:
    \begin{itemize}
      \item Sub-item 1
      \item Sub-item 2
    \end{itemize}
\end{itemize}

\subsection{Numbered List}
\begin{enumerate}
  \item First step
  \item Second step
  \item Third step
\end{enumerate}

\section{Citations and References}
When you use the citation search feature, references will be automatically managed. For example, you might cite a paper like this: \cite{example2024}. All citations will be automatically added to your \texttt{cite.bib} file.

\section{Conclusion}
This template demonstrates basic \LaTeX{} features including mathematical formulas, tables, and lists. The bibliography system is now set up to work seamlessly with Bibby AI's citation management. Feel free to modify it for your needs!

\bibliographystyle{plain}
\bibliography{cite}

\end{document}\documentclass{article}
\usepackage[utf8]{inputenc}
\usepackage{booktabs}
\usepackage{amsmath}
\usepackage{natbib}

\title{My Document}
\author{Your Name}
\date{\today}

\begin{document}

\maketitle

\section{Introduction}
This is a sample document demonstrating various \LaTeX{} features including mathematical equations, tables, and formatting. You can cite papers using the citation search feature, and they will be automatically added to your \texttt{cite.bib} file.

\section{Mathematical Expressions}
Here are some example equations:

\subsection{Quadratic Formula}
The quadratic formula is given by:
\[x = \frac{-b \pm \sqrt{b^2 - 4ac}}{2a}\]

\subsection{Matrix Example}
A sample matrix:
\[
\begin{pmatrix}
a_{11} & a_{12} & a_{13} \\
a_{21} & a_{22} & a_{23} \\
a_{31} & a_{32} & a_{33}
\end{pmatrix}
\]

\section{Tables}
Below is a sample table showing data:

\begin{table}[h]
\centering
\begin{tabular}{@{}lrrr@{}}
\toprule
\textbf{Category} & \textbf{Value 1} & \textbf{Value 2} & \textbf{Total} \\
\midrule
Item A & 10.5 & 20.3 & 30.8 \\
Item B & 15.2 & 25.1 & 40.3 \\
Item C & 22.1 & 12.7 & 34.8 \\
\bottomrule
\end{tabular}
\caption{Sample Data Table}
\label{tab:sample}
\end{table}

\section{Lists}
\subsection{Bullet Points}
\begin{itemize}
  \item First item
  \item Second item
  \item Third item with sub-items:
    \begin{itemize}
      \item Sub-item 1
      \item Sub-item 2
    \end{itemize}
\end{itemize}

\subsection{Numbered List}
\begin{enumerate}
  \item First step
  \item Second step
  \item Third step
\end{enumerate}

\section{Citations and References}
When you use the citation search feature, references will be automatically managed. For example, you might cite a paper like this: \cite{example2024}. All citations will be automatically added to your \texttt{cite.bib} file.

\section{Conclusion}
This template demonstrates basic \LaTeX{} features including mathematical formulas, tables, and lists. The bibliography system is now set up to work seamlessly with Bibby AI's citation management. Feel free to modify it for your needs!

\bibliographystyle{plain}
\bibliography{cite}

\end{document}\documentclass{article}
\usepackage[utf8]{inputenc}
\usepackage{booktabs}
\usepackage{amsmath}
\usepackage{natbib}

\title{My Document}
\author{Your Name}
\date{\today}

\begin{document}

\maketitle

\section{Introduction}
This is a sample document demonstrating various \LaTeX{} features including mathematical equations, tables, and formatting. You can cite papers using the citation search feature, and they will be automatically added to your \texttt{cite.bib} file.

\section{Mathematical Expressions}
Here are some example equations:

\subsection{Quadratic Formula}
The quadratic formula is given by:
\[x = \frac{-b \pm \sqrt{b^2 - 4ac}}{2a}\]

\subsection{Matrix Example}
A sample matrix:
\[
\begin{pmatrix}
a_{11} & a_{12} & a_{13} \\
a_{21} & a_{22} & a_{23} \\
a_{31} & a_{32} & a_{33}
\end{pmatrix}
\]

\section{Tables}
Below is a sample table showing data:

\begin{table}[h]
\centering
\begin{tabular}{@{}lrrr@{}}
\toprule
\textbf{Category} & \textbf{Value 1} & \textbf{Value 2} & \textbf{Total} \\
\midrule
Item A & 10.5 & 20.3 & 30.8 \\
Item B & 15.2 & 25.1 & 40.3 \\
Item C & 22.1 & 12.7 & 34.8 \\
\bottomrule
\end{tabular}
\caption{Sample Data Table}
\label{tab:sample}
\end{table}

\section{Lists}
\subsection{Bullet Points}
\begin{itemize}
  \item First item
  \item Second item
  \item Third item with sub-items:
    \begin{itemize}
      \item Sub-item 1
      \item Sub-item 2
    \end{itemize}
\end{itemize}

\subsection{Numbered List}
\begin{enumerate}
  \item First step
  \item Second step
  \item Third step
\end{enumerate}

\section{Citations and References}
When you use the citation search feature, references will be automatically managed. For example, you might cite a paper like this: \cite{example2024}. All citations will be automatically added to your \texttt{cite.bib} file.

\section{Conclusion}
This template demonstrates basic \LaTeX{} features including mathematical formulas, tables, and lists. The bibliography system is now set up to work seamlessly with Bibby AI's citation management. Feel free to modify it for your needs!

\bibliographystyle{plain}
\bibliography{cite}

\end{document}\documentclass{article}
\usepackage[utf8]{inputenc}
\usepackage{booktabs}
\usepackage{amsmath}
\usepackage{natbib}

\title{My Document}
\author{Your Name}
\date{\today}

\begin{document}

\maketitle

\section{Introduction}
This is a sample document demonstrating various \LaTeX{} features including mathematical equations, tables, and formatting. You can cite papers using the citation search feature, and they will be automatically added to your \texttt{cite.bib} file.

\section{Mathematical Expressions}
Here are some example equations:

\subsection{Quadratic Formula}
The quadratic formula is given by:
\[x = \frac{-b \pm \sqrt{b^2 - 4ac}}{2a}\]

\subsection{Matrix Example}
A sample matrix:
\[
\begin{pmatrix}
a_{11} & a_{12} & a_{13} \\
a_{21} & a_{22} & a_{23} \\
a_{31} & a_{32} & a_{33}
\end{pmatrix}
\]

\section{Tables}
Below is a sample table showing data:

\begin{table}[h]
\centering
\begin{tabular}{@{}lrrr@{}}
\toprule
\textbf{Category} & \textbf{Value 1} & \textbf{Value 2} & \textbf{Total} \\
\midrule
Item A & 10.5 & 20.3 & 30.8 \\
Item B & 15.2 & 25.1 & 40.3 \\
Item C & 22.1 & 12.7 & 34.8 \\
\bottomrule
\end{tabular}
\caption{Sample Data Table}
\label{tab:sample}
\end{table}

\section{Lists}
\subsection{Bullet Points}
\begin{itemize}
  \item First item
  \item Second item
  \item Third item with sub-items:
    \begin{itemize}
      \item Sub-item 1
      \item Sub-item 2
    \end{itemize}
\end{itemize}

\subsection{Numbered List}
\begin{enumerate}
  \item First step
  \item Second step
  \item Third step
\end{enumerate}

\section{Citations and References}
When you use the citation search feature, references will be automatically managed. For example, you might cite a paper like this: \cite{example2024}. All citations will be automatically added to your \texttt{cite.bib} file.

\section{Conclusion}
This template demonstrates basic \LaTeX{} features including mathematical formulas, tables, and lists. The bibliography system is now set up to work seamlessly with Bibby AI's citation management. Feel free to modify it for your needs!

\bibliographystyle{plain}
\bibliography{cite}

\end{document}\documentclass{article}
\usepackage[utf8]{inputenc}
\usepackage{booktabs}
\usepackage{amsmath}
\usepackage{natbib}

\title{My Document}
\author{Your Name}
\date{\today}

\begin{document}

\maketitle

\section{Introduction}
This is a sample document demonstrating various \LaTeX{} features including mathematical equations, tables, and formatting. You can cite papers using the citation search feature, and they will be automatically added to your \texttt{cite.bib} file.

\section{Mathematical Expressions}
Here are some example equations:

\subsection{Quadratic Formula}
The quadratic formula is given by:
\[x = \frac{-b \pm \sqrt{b^2 - 4ac}}{2a}\]

\subsection{Matrix Example}
A sample matrix:
\[
\begin{pmatrix}
a_{11} & a_{12} & a_{13} \\
a_{21} & a_{22} & a_{23} \\
a_{31} & a_{32} & a_{33}
\end{pmatrix}
\]

\section{Tables}
Below is a sample table showing data:

\begin{table}[h]
\centering
\begin{tabular}{@{}lrrr@{}}
\toprule
\textbf{Category} & \textbf{Value 1} & \textbf{Value 2} & \textbf{Total} \\
\midrule
Item A & 10.5 & 20.3 & 30.8 \\
Item B & 15.2 & 25.1 & 40.3 \\
Item C & 22.1 & 12.7 & 34.8 \\
\bottomrule
\end{tabular}
\caption{Sample Data Table}
\label{tab:sample}
\end{table}

\section{Lists}
\subsection{Bullet Points}
\begin{itemize}
  \item First item
  \item Second item
  \item Third item with sub-items:
    \begin{itemize}
      \item Sub-item 1
      \item Sub-item 2
    \end{itemize}
\end{itemize}

\subsection{Numbered List}
\begin{enumerate}
  \item First step
  \item Second step
  \item Third step
\end{enumerate}

\section{Citations and References}
When you use the citation search feature, references will be automatically managed. For example, you might cite a paper like this: \cite{example2024}. All citations will be automatically added to your \texttt{cite.bib} file.

\section{Conclusion}
This template demonstrates basic \LaTeX{} features including mathematical formulas, tables, and lists. The bibliography system is now set up to work seamlessly with Bibby AI's citation management. Feel free to modify it for your needs!

\bibliographystyle{plain}
\bibliography{cite}

\end{document}\documentclass{article}
\usepackage[utf8]{inputenc}
\usepackage{booktabs}
\usepackage{amsmath}
\usepackage{natbib}

\title{My Document}
\author{Your Name}
\date{\today}

\begin{document}

\maketitle

\section{Introduction}
This is a sample document demonstrating various \LaTeX{} features including mathematical equations, tables, and formatting. You can cite papers using the citation search feature, and they will be automatically added to your \texttt{cite.bib} file.

\section{Mathematical Expressions}
Here are some example equations:

\subsection{Quadratic Formula}
The quadratic formula is given by:
\[x = \frac{-b \pm \sqrt{b^2 - 4ac}}{2a}\]

\subsection{Matrix Example}
A sample matrix:
\[
\begin{pmatrix}
a_{11} & a_{12} & a_{13} \\
a_{21} & a_{22} & a_{23} \\
a_{31} & a_{32} & a_{33}
\end{pmatrix}
\]

\section{Tables}
Below is a sample table showing data:

\begin{table}[h]
\centering
\begin{tabular}{@{}lrrr@{}}
\toprule
\textbf{Category} & \textbf{Value 1} & \textbf{Value 2} & \textbf{Total} \\
\midrule
Item A & 10.5 & 20.3 & 30.8 \\
Item B & 15.2 & 25.1 & 40.3 \\
Item C & 22.1 & 12.7 & 34.8 \\
\bottomrule
\end{tabular}
\caption{Sample Data Table}
\label{tab:sample}
\end{table}

\section{Lists}
\subsection{Bullet Points}
\begin{itemize}
  \item First item
  \item Second item
  \item Third item with sub-items:
    \begin{itemize}
      \item Sub-item 1
      \item Sub-item 2
    \end{itemize}
\end{itemize}

\subsection{Numbered List}
\begin{enumerate}
  \item First step
  \item Second step
  \item Third step
\end{enumerate}

\section{Citations and References}
When you use the citation search feature, references will be automatically managed. For example, you might cite a paper like this: \cite{example2024}. All citations will be automatically added to your \texttt{cite.bib} file.

\section{Conclusion}
This template demonstrates basic \LaTeX{} features including mathematical formulas, tables, and lists. The bibliography system is now set up to work seamlessly with Bibby AI's citation management. Feel free to modify it for your needs!

\bibliographystyle{plain}
\bibliography{cite}

\end{document}\documentclass{article}
\usepackage[utf8]{inputenc}
\usepackage{booktabs}
\usepackage{amsmath}
\usepackage{natbib}

\title{My Document}
\author{Your Name}
\date{\today}

\begin{document}

\maketitle

\section{Introduction}
This is a sample document demonstrating various \LaTeX{} features including mathematical equations, tables, and formatting. You can cite papers using the citation search feature, and they will be automatically added to your \texttt{cite.bib} file.

\section{Mathematical Expressions}
Here are some example equations:

\subsection{Quadratic Formula}
The quadratic formula is given by:
\[x = \frac{-b \pm \sqrt{b^2 - 4ac}}{2a}\]

\subsection{Matrix Example}
A sample matrix:
\[
\begin{pmatrix}
a_{11} & a_{12} & a_{13} \\
a_{21} & a_{22} & a_{23} \\
a_{31} & a_{32} & a_{33}
\end{pmatrix}
\]

\section{Tables}
Below is a sample table showing data:

\begin{table}[h]
\centering
\begin{tabular}{@{}lrrr@{}}
\toprule
\textbf{Category} & \textbf{Value 1} & \textbf{Value 2} & \textbf{Total} \\
\midrule
Item A & 10.5 & 20.3 & 30.8 \\
Item B & 15.2 & 25.1 & 40.3 \\
Item C & 22.1 & 12.7 & 34.8 \\
\bottomrule
\end{tabular}
\caption{Sample Data Table}
\label{tab:sample}
\end{table}

\section{Lists}
\subsection{Bullet Points}
\begin{itemize}
  \item First item
  \item Second item
  \item Third item with sub-items:
    \begin{itemize}
      \item Sub-item 1
      \item Sub-item 2
    \end{itemize}
\end{itemize}

\subsection{Numbered List}
\begin{enumerate}
  \item First step
  \item Second step
  \item Third step
\end{enumerate}

\section{Citations and References}
When you use the citation search feature, references will be automatically managed. For example, you might cite a paper like this: \cite{example2024}. All citations will be automatically added to your \texttt{cite.bib} file.

\section{Conclusion}
This template demonstrates basic \LaTeX{} features including mathematical formulas, tables, and lists. The bibliography system is now set up to work seamlessly with Bibby AI's citation management. Feel free to modify it for your needs!

\bibliographystyle{plain}
\bibliography{cite}

\end{document}\documentclass{article}
\usepackage[utf8]{inputenc}
\usepackage{booktabs}
\usepackage{amsmath}
\usepackage{natbib}

\title{My Document}
\author{Your Name}
\date{\today}

\begin{document}

\maketitle

\section{Introduction}
This is a sample document demonstrating various \LaTeX{} features including mathematical equations, tables, and formatting. You can cite papers using the citation search feature, and they will be automatically added to your \texttt{cite.bib} file.

\section{Mathematical Expressions}
Here are some example equations:

\subsection{Quadratic Formula}
The quadratic formula is given by:
\[x = \frac{-b \pm \sqrt{b^2 - 4ac}}{2a}\]

\subsection{Matrix Example}
A sample matrix:
\[
\begin{pmatrix}
a_{11} & a_{12} & a_{13} \\
a_{21} & a_{22} & a_{23} \\
a_{31} & a_{32} & a_{33}
\end{pmatrix}
\]

\section{Tables}
Below is a sample table showing data:

\begin{table}[h]
\centering
\begin{tabular}{@{}lrrr@{}}
\toprule
\textbf{Category} & \textbf{Value 1} & \textbf{Value 2} & \textbf{Total} \\
\midrule
Item A & 10.5 & 20.3 & 30.8 \\
Item B & 15.2 & 25.1 & 40.3 \\
Item C & 22.1 & 12.7 & 34.8 \\
\bottomrule
\end{tabular}
\caption{Sample Data Table}
\label{tab:sample}
\end{table}

\section{Lists}
\subsection{Bullet Points}
\begin{itemize}
  \item First item
  \item Second item
  \item Third item with sub-items:
    \begin{itemize}
      \item Sub-item 1
      \item Sub-item 2
    \end{itemize}
\end{itemize}

\subsection{Numbered List}
\begin{enumerate}
  \item First step
  \item Second step
  \item Third step
\end{enumerate}

\section{Citations and References}
When you use the citation search feature, references will be automatically managed. For example, you might cite a paper like this: \cite{example2024}. All citations will be automatically added to your \texttt{cite.bib} file.

\section{Conclusion}
This template demonstrates basic \LaTeX{} features including mathematical formulas, tables, and lists. The bibliography system is now set up to work seamlessly with Bibby AI's citation management. Feel free to modify it for your needs!

\bibliographystyle{plain}
\bibliography{cite}

\end{document}\documentclass{article}
\usepackage[utf8]{inputenc}
\usepackage{booktabs}
\usepackage{amsmath}
\usepackage{natbib}

\title{My Document}
\author{Your Name}
\date{\today}

\begin{document}

\maketitle

\section{Introduction}
This is a sample document demonstrating various \LaTeX{} features including mathematical equations, tables, and formatting. You can cite papers using the citation search feature, and they will be automatically added to your \texttt{cite.bib} file.

\section{Mathematical Expressions}
Here are some example equations:

\subsection{Quadratic Formula}
The quadratic formula is given by:
\[x = \frac{-b \pm \sqrt{b^2 - 4ac}}{2a}\]

\subsection{Matrix Example}
A sample matrix:
\[
\begin{pmatrix}
a_{11} & a_{12} & a_{13} \\
a_{21} & a_{22} & a_{23} \\
a_{31} & a_{32} & a_{33}
\end{pmatrix}
\]

\section{Tables}
Below is a sample table showing data:

\begin{table}[h]
\centering
\begin{tabular}{@{}lrrr@{}}
\toprule
\textbf{Category} & \textbf{Value 1} & \textbf{Value 2} & \textbf{Total} \\
\midrule
Item A & 10.5 & 20.3 & 30.8 \\
Item B & 15.2 & 25.1 & 40.3 \\
Item C & 22.1 & 12.7 & 34.8 \\
\bottomrule
\end{tabular}
\caption{Sample Data Table}
\label{tab:sample}
\end{table}

\section{Lists}
\subsection{Bullet Points}
\begin{itemize}
  \item First item
  \item Second item
  \item Third item with sub-items:
    \begin{itemize}
      \item Sub-item 1
      \item Sub-item 2
    \end{itemize}
\end{itemize}

\subsection{Numbered List}
\begin{enumerate}
  \item First step
  \item Second step
  \item Third step
\end{enumerate}

\section{Citations and References}
When you use the citation search feature, references will be automatically managed. For example, you might cite a paper like this: \cite{example2024}. All citations will be automatically added to your \texttt{cite.bib} file.

\section{Conclusion}
This template demonstrates basic \LaTeX{} features including mathematical formulas, tables, and lists. The bibliography system is now set up to work seamlessly with Bibby AI's citation management. Feel free to modify it for your needs!

\bibliographystyle{plain}
\bibliography{cite}

\end{document}\documentclass{article}
\usepackage[utf8]{inputenc}
\usepackage{booktabs}
\usepackage{amsmath}
\usepackage{natbib}

\title{My Document}
\author{Your Name}
\date{\today}

\begin{document}

\maketitle

\section{Introduction}
This is a sample document demonstrating various \LaTeX{} features including mathematical equations, tables, and formatting. You can cite papers using the citation search feature, and they will be automatically added to your \texttt{cite.bib} file.

\section{Mathematical Expressions}
Here are some example equations:

\subsection{Quadratic Formula}
The quadratic formula is given by:
\[x = \frac{-b \pm \sqrt{b^2 - 4ac}}{2a}\]

\subsection{Matrix Example}
A sample matrix:
\[
\begin{pmatrix}
a_{11} & a_{12} & a_{13} \\
a_{21} & a_{22} & a_{23} \\
a_{31} & a_{32} & a_{33}
\end{pmatrix}
\]

\section{Tables}
Below is a sample table showing data:

\begin{table}[h]
\centering
\begin{tabular}{@{}lrrr@{}}
\toprule
\textbf{Category} & \textbf{Value 1} & \textbf{Value 2} & \textbf{Total} \\
\midrule
Item A & 10.5 & 20.3 & 30.8 \\
Item B & 15.2 & 25.1 & 40.3 \\
Item C & 22.1 & 12.7 & 34.8 \\
\bottomrule
\end{tabular}
\caption{Sample Data Table}
\label{tab:sample}
\end{table}

\section{Lists}
\subsection{Bullet Points}
\begin{itemize}
  \item First item
  \item Second item
  \item Third item with sub-items:
    \begin{itemize}
      \item Sub-item 1
      \item Sub-item 2
    \end{itemize}
\end{itemize}

\subsection{Numbered List}
\begin{enumerate}
  \item First step
  \item Second step
  \item Third step
\end{enumerate}

\section{Citations and References}
When you use the citation search feature, references will be automatically managed. For example, you might cite a paper like this: \cite{example2024}. All citations will be automatically added to your \texttt{cite.bib} file.

\section{Conclusion}
This template demonstrates basic \LaTeX{} features including mathematical formulas, tables, and lists. The bibliography system is now set up to work seamlessly with Bibby AI's citation management. Feel free to modify it for your needs!

\bibliographystyle{plain}
\bibliography{cite}

\end{document}\documentclass{article}
\usepackage[utf8]{inputenc}
\usepackage{booktabs}
\usepackage{amsmath}
\usepackage{natbib}

\title{My Document}
\author{Your Name}
\date{\today}

\begin{document}

\maketitle

\section{Introduction}
This is a sample document demonstrating various \LaTeX{} features including mathematical equations, tables, and formatting. You can cite papers using the citation search feature, and they will be automatically added to your \texttt{cite.bib} file.

\section{Mathematical Expressions}
Here are some example equations:

\subsection{Quadratic Formula}
The quadratic formula is given by:
\[x = \frac{-b \pm \sqrt{b^2 - 4ac}}{2a}\]

\subsection{Matrix Example}
A sample matrix:
\[
\begin{pmatrix}
a_{11} & a_{12} & a_{13} \\
a_{21} & a_{22} & a_{23} \\
a_{31} & a_{32} & a_{33}
\end{pmatrix}
\]

\section{Tables}
Below is a sample table showing data:

\begin{table}[h]
\centering
\begin{tabular}{@{}lrrr@{}}
\toprule
\textbf{Category} & \textbf{Value 1} & \textbf{Value 2} & \textbf{Total} \\
\midrule
Item A & 10.5 & 20.3 & 30.8 \\
Item B & 15.2 & 25.1 & 40.3 \\
Item C & 22.1 & 12.7 & 34.8 \\
\bottomrule
\end{tabular}
\caption{Sample Data Table}
\label{tab:sample}
\end{table}

\section{Lists}
\subsection{Bullet Points}
\begin{itemize}
  \item First item
  \item Second item
  \item Third item with sub-items:
    \begin{itemize}
      \item Sub-item 1
      \item Sub-item 2
    \end{itemize}
\end{itemize}

\subsection{Numbered List}
\begin{enumerate}
  \item First step
  \item Second step
  \item Third step
\end{enumerate}

\section{Citations and References}
When you use the citation search feature, references will be automatically managed. For example, you might cite a paper like this: \cite{example2024}. All citations will be automatically added to your \texttt{cite.bib} file.

\section{Conclusion}
This template demonstrates basic \LaTeX{} features including mathematical formulas, tables, and lists. The bibliography system is now set up to work seamlessly with Bibby AI's citation management. Feel free to modify it for your needs!

\bibliographystyle{plain}
\bibliography{cite}

\end{document}\documentclass{article}
\usepackage[utf8]{inputenc}
\usepackage{booktabs}
\usepackage{amsmath}
\usepackage{natbib}

\title{My Document}
\author{Your Name}
\date{\today}

\begin{document}

\maketitle

\section{Introduction}
This is a sample document demonstrating various \LaTeX{} features including mathematical equations, tables, and formatting. You can cite papers using the citation search feature, and they will be automatically added to your \texttt{cite.bib} file.

\section{Mathematical Expressions}
Here are some example equations:

\subsection{Quadratic Formula}
The quadratic formula is given by:
\[x = \frac{-b \pm \sqrt{b^2 - 4ac}}{2a}\]

\subsection{Matrix Example}
A sample matrix:
\[
\begin{pmatrix}
a_{11} & a_{12} & a_{13} \\
a_{21} & a_{22} & a_{23} \\
a_{31} & a_{32} & a_{33}
\end{pmatrix}
\]

\section{Tables}
Below is a sample table showing data:

\begin{table}[h]
\centering
\begin{tabular}{@{}lrrr@{}}
\toprule
\textbf{Category} & \textbf{Value 1} & \textbf{Value 2} & \textbf{Total} \\
\midrule
Item A & 10.5 & 20.3 & 30.8 \\
Item B & 15.2 & 25.1 & 40.3 \\
Item C & 22.1 & 12.7 & 34.8 \\
\bottomrule
\end{tabular}
\caption{Sample Data Table}
\label{tab:sample}
\end{table}

\section{Lists}
\subsection{Bullet Points}
\begin{itemize}
  \item First item
  \item Second item
  \item Third item with sub-items:
    \begin{itemize}
      \item Sub-item 1
      \item Sub-item 2
    \end{itemize}
\end{itemize}

\subsection{Numbered List}
\begin{enumerate}
  \item First step
  \item Second step
  \item Third step
\end{enumerate}

\section{Citations and References}
When you use the citation search feature, references will be automatically managed. For example, you might cite a paper like this: \cite{example2024}. All citations will be automatically added to your \texttt{cite.bib} file.

\section{Conclusion}
This template demonstrates basic \LaTeX{} features including mathematical formulas, tables, and lists. The bibliography system is now set up to work seamlessly with Bibby AI's citation management. Feel free to modify it for your needs!

\bibliographystyle{plain}
\bibliography{cite}

\end{document}\documentclass{article}
\usepackage[utf8]{inputenc}
\usepackage{booktabs}
\usepackage{amsmath}
\usepackage{natbib}

\title{My Document}
\author{Your Name}
\date{\today}

\begin{document}

\maketitle

\section{Introduction}
This is a sample document demonstrating various \LaTeX{} features including mathematical equations, tables, and formatting. You can cite papers using the citation search feature, and they will be automatically added to your \texttt{cite.bib} file.

\section{Mathematical Expressions}
Here are some example equations:

\subsection{Quadratic Formula}
The quadratic formula is given by:
\[x = \frac{-b \pm \sqrt{b^2 - 4ac}}{2a}\]

\subsection{Matrix Example}
A sample matrix:
\[
\begin{pmatrix}
a_{11} & a_{12} & a_{13} \\
a_{21} & a_{22} & a_{23} \\
a_{31} & a_{32} & a_{33}
\end{pmatrix}
\]

\section{Tables}
Below is a sample table showing data:

\begin{table}[h]
\centering
\begin{tabular}{@{}lrrr@{}}
\toprule
\textbf{Category} & \textbf{Value 1} & \textbf{Value 2} & \textbf{Total} \\
\midrule
Item A & 10.5 & 20.3 & 30.8 \\
Item B & 15.2 & 25.1 & 40.3 \\
Item C & 22.1 & 12.7 & 34.8 \\
\bottomrule
\end{tabular}
\caption{Sample Data Table}
\label{tab:sample}
\end{table}

\section{Lists}
\subsection{Bullet Points}
\begin{itemize}
  \item First item
  \item Second item
  \item Third item with sub-items:
    \begin{itemize}
      \item Sub-item 1
      \item Sub-item 2
    \end{itemize}
\end{itemize}

\subsection{Numbered List}
\begin{enumerate}
  \item First step
  \item Second step
  \item Third step
\end{enumerate}

\section{Citations and References}
When you use the citation search feature, references will be automatically managed. For example, you might cite a paper like this: \cite{example2024}. All citations will be automatically added to your \texttt{cite.bib} file.

\section{Conclusion}
This template demonstrates basic \LaTeX{} features including mathematical formulas, tables, and lists. The bibliography system is now set up to work seamlessly with Bibby AI's citation management. Feel free to modify it for your needs!

\bibliographystyle{plain}
\bibliography{cite}

\end{document}\documentclass{article}
\usepackage[utf8]{inputenc}
\usepackage{booktabs}
\usepackage{amsmath}
\usepackage{natbib}

\title{My Document}
\author{Your Name}
\date{\today}

\begin{document}

\maketitle

\section{Introduction}
This is a sample document demonstrating various \LaTeX{} features including mathematical equations, tables, and formatting. You can cite papers using the citation search feature, and they will be automatically added to your \texttt{cite.bib} file.

\section{Mathematical Expressions}
Here are some example equations:

\subsection{Quadratic Formula}
The quadratic formula is given by:
\[x = \frac{-b \pm \sqrt{b^2 - 4ac}}{2a}\]

\subsection{Matrix Example}
A sample matrix:
\[
\begin{pmatrix}
a_{11} & a_{12} & a_{13} \\
a_{21} & a_{22} & a_{23} \\
a_{31} & a_{32} & a_{33}
\end{pmatrix}
\]

\section{Tables}
Below is a sample table showing data:

\begin{table}[h]
\centering
\begin{tabular}{@{}lrrr@{}}
\toprule
\textbf{Category} & \textbf{Value 1} & \textbf{Value 2} & \textbf{Total} \\
\midrule
Item A & 10.5 & 20.3 & 30.8 \\
Item B & 15.2 & 25.1 & 40.3 \\
Item C & 22.1 & 12.7 & 34.8 \\
\bottomrule
\end{tabular}
\caption{Sample Data Table}
\label{tab:sample}
\end{table}

\section{Lists}
\subsection{Bullet Points}
\begin{itemize}
  \item First item
  \item Second item
  \item Third item with sub-items:
    \begin{itemize}
      \item Sub-item 1
      \item Sub-item 2
    \end{itemize}
\end{itemize}

\subsection{Numbered List}
\begin{enumerate}
  \item First step
  \item Second step
  \item Third step
\end{enumerate}

\section{Citations and References}
When you use the citation search feature, references will be automatically managed. For example, you might cite a paper like this: \cite{example2024}. All citations will be automatically added to your \texttt{cite.bib} file.

\section{Conclusion}
This template demonstrates basic \LaTeX{} features including mathematical formulas, tables, and lists. The bibliography system is now set up to work seamlessly with Bibby AI's citation management. Feel free to modify it for your needs!

\bibliographystyle{plain}
\bibliography{cite}

\end{document}\documentclass{article}
\usepackage[utf8]{inputenc}
\usepackage{booktabs}
\usepackage{amsmath}
\usepackage{natbib}

\title{My Document}
\author{Your Name}
\date{\today}

\begin{document}

\maketitle

\section{Introduction}
This is a sample document demonstrating various \LaTeX{} features including mathematical equations, tables, and formatting. You can cite papers using the citation search feature, and they will be automatically added to your \texttt{cite.bib} file.

\section{Mathematical Expressions}
Here are some example equations:

\subsection{Quadratic Formula}
The quadratic formula is given by:
\[x = \frac{-b \pm \sqrt{b^2 - 4ac}}{2a}\]

\subsection{Matrix Example}
A sample matrix:
\[
\begin{pmatrix}
a_{11} & a_{12} & a_{13} \\
a_{21} & a_{22} & a_{23} \\
a_{31} & a_{32} & a_{33}
\end{pmatrix}
\]

\section{Tables}
Below is a sample table showing data:

\begin{table}[h]
\centering
\begin{tabular}{@{}lrrr@{}}
\toprule
\textbf{Category} & \textbf{Value 1} & \textbf{Value 2} & \textbf{Total} \\
\midrule
Item A & 10.5 & 20.3 & 30.8 \\
Item B & 15.2 & 25.1 & 40.3 \\
Item C & 22.1 & 12.7 & 34.8 \\
\bottomrule
\end{tabular}
\caption{Sample Data Table}
\label{tab:sample}
\end{table}

\section{Lists}
\subsection{Bullet Points}
\begin{itemize}
  \item First item
  \item Second item
  \item Third item with sub-items:
    \begin{itemize}
      \item Sub-item 1
      \item Sub-item 2
    \end{itemize}
\end{itemize}

\subsection{Numbered List}
\begin{enumerate}
  \item First step
  \item Second step
  \item Third step
\end{enumerate}

\section{Citations and References}
When you use the citation search feature, references will be automatically managed. For example, you might cite a paper like this: \cite{example2024}. All citations will be automatically added to your \texttt{cite.bib} file.

\section{Conclusion}
This template demonstrates basic \LaTeX{} features including mathematical formulas, tables, and lists. The bibliography system is now set up to work seamlessly with Bibby AI's citation management. Feel free to modify it for your needs!

\bibliographystyle{plain}
\bibliography{cite}

\end{document}\documentclass{article}
\usepackage[utf8]{inputenc}
\usepackage{booktabs}
\usepackage{amsmath}
\usepackage{natbib}

\title{My Document}
\author{Your Name}
\date{\today}

\begin{document}

\maketitle

\section{Introduction}
This is a sample document demonstrating various \LaTeX{} features including mathematical equations, tables, and formatting. You can cite papers using the citation search feature, and they will be automatically added to your \texttt{cite.bib} file.

\section{Mathematical Expressions}
Here are some example equations:

\subsection{Quadratic Formula}
The quadratic formula is given by:
\[x = \frac{-b \pm \sqrt{b^2 - 4ac}}{2a}\]

\subsection{Matrix Example}
A sample matrix:
\[
\begin{pmatrix}
a_{11} & a_{12} & a_{13} \\
a_{21} & a_{22} & a_{23} \\
a_{31} & a_{32} & a_{33}
\end{pmatrix}
\]

\section{Tables}
Below is a sample table showing data:

\begin{table}[h]
\centering
\begin{tabular}{@{}lrrr@{}}
\toprule
\textbf{Category} & \textbf{Value 1} & \textbf{Value 2} & \textbf{Total} \\
\midrule
Item A & 10.5 & 20.3 & 30.8 \\
Item B & 15.2 & 25.1 & 40.3 \\
Item C & 22.1 & 12.7 & 34.8 \\
\bottomrule
\end{tabular}
\caption{Sample Data Table}
\label{tab:sample}
\end{table}

\section{Lists}
\subsection{Bullet Points}
\begin{itemize}
  \item First item
  \item Second item
  \item Third item with sub-items:
    \begin{itemize}
      \item Sub-item 1
      \item Sub-item 2
    \end{itemize}
\end{itemize}

\subsection{Numbered List}
\begin{enumerate}
  \item First step
  \item Second step
  \item Third step
\end{enumerate}

\section{Citations and References}
When you use the citation search feature, references will be automatically managed. For example, you might cite a paper like this: \cite{example2024}. All citations will be automatically added to your \texttt{cite.bib} file.

\section{Conclusion}
This template demonstrates basic \LaTeX{} features including mathematical formulas, tables, and lists. The bibliography system is now set up to work seamlessly with Bibby AI's citation management. Feel free to modify it for your needs!

\bibliographystyle{plain}
\bibliography{cite}

\end{document}\documentclass{article}
\usepackage[utf8]{inputenc}
\usepackage{booktabs}
\usepackage{amsmath}
\usepackage{natbib}

\title{My Document}
\author{Your Name}
\date{\today}

\begin{document}

\maketitle

\section{Introduction}
This is a sample document demonstrating various \LaTeX{} features including mathematical equations, tables, and formatting. You can cite papers using the citation search feature, and they will be automatically added to your \texttt{cite.bib} file.

\section{Mathematical Expressions}
Here are some example equations:

\subsection{Quadratic Formula}
The quadratic formula is given by:
\[x = \frac{-b \pm \sqrt{b^2 - 4ac}}{2a}\]

\subsection{Matrix Example}
A sample matrix:
\[
\begin{pmatrix}
a_{11} & a_{12} & a_{13} \\
a_{21} & a_{22} & a_{23} \\
a_{31} & a_{32} & a_{33}
\end{pmatrix}
\]

\section{Tables}
Below is a sample table showing data:

\begin{table}[h]
\centering
\begin{tabular}{@{}lrrr@{}}
\toprule
\textbf{Category} & \textbf{Value 1} & \textbf{Value 2} & \textbf{Total} \\
\midrule
Item A & 10.5 & 20.3 & 30.8 \\
Item B & 15.2 & 25.1 & 40.3 \\
Item C & 22.1 & 12.7 & 34.8 \\
\bottomrule
\end{tabular}
\caption{Sample Data Table}
\label{tab:sample}
\end{table}

\section{Lists}
\subsection{Bullet Points}
\begin{itemize}
  \item First item
  \item Second item
  \item Third item with sub-items:
    \begin{itemize}
      \item Sub-item 1
      \item Sub-item 2
    \end{itemize}
\end{itemize}

\subsection{Numbered List}
\begin{enumerate}
  \item First step
  \item Second step
  \item Third step
\end{enumerate}

\section{Citations and References}
When you use the citation search feature, references will be automatically managed. For example, you might cite a paper like this: \cite{example2024}. All citations will be automatically added to your \texttt{cite.bib} file.

\section{Conclusion}
This template demonstrates basic \LaTeX{} features including mathematical formulas, tables, and lists. The bibliography system is now set up to work seamlessly with Bibby AI's citation management. Feel free to modify it for your needs!

\bibliographystyle{plain}
\bibliography{cite}

\end{document}\documentclass{article}
\usepackage[utf8]{inputenc}
\usepackage{booktabs}
\usepackage{amsmath}
\usepackage{natbib}

\title{My Document}
\author{Your Name}
\date{\today}

\begin{document}

\maketitle

\section{Introduction}
This is a sample document demonstrating various \LaTeX{} features including mathematical equations, tables, and formatting. You can cite papers using the citation search feature, and they will be automatically added to your \texttt{cite.bib} file.

\section{Mathematical Expressions}
Here are some example equations:

\subsection{Quadratic Formula}
The quadratic formula is given by:
\[x = \frac{-b \pm \sqrt{b^2 - 4ac}}{2a}\]

\subsection{Matrix Example}
A sample matrix:
\[
\begin{pmatrix}
a_{11} & a_{12} & a_{13} \\
a_{21} & a_{22} & a_{23} \\
a_{31} & a_{32} & a_{33}
\end{pmatrix}
\]

\section{Tables}
Below is a sample table showing data:

\begin{table}[h]
\centering
\begin{tabular}{@{}lrrr@{}}
\toprule
\textbf{Category} & \textbf{Value 1} & \textbf{Value 2} & \textbf{Total} \\
\midrule
Item A & 10.5 & 20.3 & 30.8 \\
Item B & 15.2 & 25.1 & 40.3 \\
Item C & 22.1 & 12.7 & 34.8 \\
\bottomrule
\end{tabular}
\caption{Sample Data Table}
\label{tab:sample}
\end{table}

\section{Lists}
\subsection{Bullet Points}
\begin{itemize}
  \item First item
  \item Second item
  \item Third item with sub-items:
    \begin{itemize}
      \item Sub-item 1
      \item Sub-item 2
    \end{itemize}
\end{itemize}

\subsection{Numbered List}
\begin{enumerate}
  \item First step
  \item Second step
  \item Third step
\end{enumerate}

\section{Citations and References}
When you use the citation search feature, references will be automatically managed. For example, you might cite a paper like this: \cite{example2024}. All citations will be automatically added to your \texttt{cite.bib} file.

\section{Conclusion}
This template demonstrates basic \LaTeX{} features including mathematical formulas, tables, and lists. The bibliography system is now set up to work seamlessly with Bibby AI's citation management. Feel free to modify it for your needs!

\bibliographystyle{plain}
\bibliography{cite}

\end{document}\documentclass{article}
\usepackage[utf8]{inputenc}
\usepackage{booktabs}
\usepackage{amsmath}
\usepackage{natbib}

\title{My Document}
\author{Your Name}
\date{\today}

\begin{document}

\maketitle

\section{Introduction}
This is a sample document demonstrating various \LaTeX{} features including mathematical equations, tables, and formatting. You can cite papers using the citation search feature, and they will be automatically added to your \texttt{cite.bib} file.

\section{Mathematical Expressions}
Here are some example equations:

\subsection{Quadratic Formula}
The quadratic formula is given by:
\[x = \frac{-b \pm \sqrt{b^2 - 4ac}}{2a}\]

\subsection{Matrix Example}
A sample matrix:
\[
\begin{pmatrix}
a_{11} & a_{12} & a_{13} \\
a_{21} & a_{22} & a_{23} \\
a_{31} & a_{32} & a_{33}
\end{pmatrix}
\]

\section{Tables}
Below is a sample table showing data:

\begin{table}[h]
\centering
\begin{tabular}{@{}lrrr@{}}
\toprule
\textbf{Category} & \textbf{Value 1} & \textbf{Value 2} & \textbf{Total} \\
\midrule
Item A & 10.5 & 20.3 & 30.8 \\
Item B & 15.2 & 25.1 & 40.3 \\
Item C & 22.1 & 12.7 & 34.8 \\
\bottomrule
\end{tabular}
\caption{Sample Data Table}
\label{tab:sample}
\end{table}

\section{Lists}
\subsection{Bullet Points}
\begin{itemize}
  \item First item
  \item Second item
  \item Third item with sub-items:
    \begin{itemize}
      \item Sub-item 1
      \item Sub-item 2
    \end{itemize}
\end{itemize}

\subsection{Numbered List}
\begin{enumerate}
  \item First step
  \item Second step
  \item Third step
\end{enumerate}

\section{Citations and References}
When you use the citation search feature, references will be automatically managed. For example, you might cite a paper like this: \cite{example2024}. All citations will be automatically added to your \texttt{cite.bib} file.

\section{Conclusion}
This template demonstrates basic \LaTeX{} features including mathematical formulas, tables, and lists. The bibliography system is now set up to work seamlessly with Bibby AI's citation management. Feel free to modify it for your needs!

\bibliographystyle{plain}
\bibliography{cite}

\end{document}\documentclass{article}
\usepackage[utf8]{inputenc}
\usepackage{booktabs}
\usepackage{amsmath}
\usepackage{natbib}

\title{My Document}
\author{Your Name}
\date{\today}

\begin{document}

\maketitle

\section{Introduction}
This is a sample document demonstrating various \LaTeX{} features including mathematical equations, tables, and formatting. You can cite papers using the citation search feature, and they will be automatically added to your \texttt{cite.bib} file.

\section{Mathematical Expressions}
Here are some example equations:

\subsection{Quadratic Formula}
The quadratic formula is given by:
\[x = \frac{-b \pm \sqrt{b^2 - 4ac}}{2a}\]

\subsection{Matrix Example}
A sample matrix:
\[
\begin{pmatrix}
a_{11} & a_{12} & a_{13} \\
a_{21} & a_{22} & a_{23} \\
a_{31} & a_{32} & a_{33}
\end{pmatrix}
\]

\section{Tables}
Below is a sample table showing data:

\begin{table}[h]
\centering
\begin{tabular}{@{}lrrr@{}}
\toprule
\textbf{Category} & \textbf{Value 1} & \textbf{Value 2} & \textbf{Total} \\
\midrule
Item A & 10.5 & 20.3 & 30.8 \\
Item B & 15.2 & 25.1 & 40.3 \\
Item C & 22.1 & 12.7 & 34.8 \\
\bottomrule
\end{tabular}
\caption{Sample Data Table}
\label{tab:sample}
\end{table}

\section{Lists}
\subsection{Bullet Points}
\begin{itemize}
  \item First item
  \item Second item
  \item Third item with sub-items:
    \begin{itemize}
      \item Sub-item 1
      \item Sub-item 2
    \end{itemize}
\end{itemize}

\subsection{Numbered List}
\begin{enumerate}
  \item First step
  \item Second step
  \item Third step
\end{enumerate}

\section{Citations and References}
When you use the citation search feature, references will be automatically managed. For example, you might cite a paper like this: \cite{example2024}. All citations will be automatically added to your \texttt{cite.bib} file.

\section{Conclusion}
This template demonstrates basic \LaTeX{} features including mathematical formulas, tables, and lists. The bibliography system is now set up to work seamlessly with Bibby AI's citation management. Feel free to modify it for your needs!

\bibliographystyle{plain}
\bibliography{cite}

\end{document}\documentclass{article}
\usepackage[utf8]{inputenc}
\usepackage{booktabs}
\usepackage{amsmath}
\usepackage{natbib}

\title{My Document}
\author{Your Name}
\date{\today}

\begin{document}

\maketitle

\section{Introduction}
This is a sample document demonstrating various \LaTeX{} features including mathematical equations, tables, and formatting. You can cite papers using the citation search feature, and they will be automatically added to your \texttt{cite.bib} file.

\section{Mathematical Expressions}
Here are some example equations:

\subsection{Quadratic Formula}
The quadratic formula is given by:
\[x = \frac{-b \pm \sqrt{b^2 - 4ac}}{2a}\]

\subsection{Matrix Example}
A sample matrix:
\[
\begin{pmatrix}
a_{11} & a_{12} & a_{13} \\
a_{21} & a_{22} & a_{23} \\
a_{31} & a_{32} & a_{33}
\end{pmatrix}
\]

\section{Tables}
Below is a sample table showing data:

\begin{table}[h]
\centering
\begin{tabular}{@{}lrrr@{}}
\toprule
\textbf{Category} & \textbf{Value 1} & \textbf{Value 2} & \textbf{Total} \\
\midrule
Item A & 10.5 & 20.3 & 30.8 \\
Item B & 15.2 & 25.1 & 40.3 \\
Item C & 22.1 & 12.7 & 34.8 \\
\bottomrule
\end{tabular}
\caption{Sample Data Table}
\label{tab:sample}
\end{table}

\section{Lists}
\subsection{Bullet Points}
\begin{itemize}
  \item First item
  \item Second item
  \item Third item with sub-items:
    \begin{itemize}
      \item Sub-item 1
      \item Sub-item 2
    \end{itemize}
\end{itemize}

\subsection{Numbered List}
\begin{enumerate}
  \item First step
  \item Second step
  \item Third step
\end{enumerate}

\section{Citations and References}
When you use the citation search feature, references will be automatically managed. For example, you might cite a paper like this: \cite{example2024}. All citations will be automatically added to your \texttt{cite.bib} file.

\section{Conclusion}
This template demonstrates basic \LaTeX{} features including mathematical formulas, tables, and lists. The bibliography system is now set up to work seamlessly with Bibby AI's citation management. Feel free to modify it for your needs!

\bibliographystyle{plain}
\bibliography{cite}

\end{document}\documentclass{article}
\usepackage[utf8]{inputenc}
\usepackage{booktabs}
\usepackage{amsmath}
\usepackage{natbib}

\title{My Document}
\author{Your Name}
\date{\today}

\begin{document}

\maketitle

\section{Introduction}
This is a sample document demonstrating various \LaTeX{} features including mathematical equations, tables, and formatting. You can cite papers using the citation search feature, and they will be automatically added to your \texttt{cite.bib} file.

\section{Mathematical Expressions}
Here are some example equations:

\subsection{Quadratic Formula}
The quadratic formula is given by:
\[x = \frac{-b \pm \sqrt{b^2 - 4ac}}{2a}\]

\subsection{Matrix Example}
A sample matrix:
\[
\begin{pmatrix}
a_{11} & a_{12} & a_{13} \\
a_{21} & a_{22} & a_{23} \\
a_{31} & a_{32} & a_{33}
\end{pmatrix}
\]

\section{Tables}
Below is a sample table showing data:

\begin{table}[h]
\centering
\begin{tabular}{@{}lrrr@{}}
\toprule
\textbf{Category} & \textbf{Value 1} & \textbf{Value 2} & \textbf{Total} \\
\midrule
Item A & 10.5 & 20.3 & 30.8 \\
Item B & 15.2 & 25.1 & 40.3 \\
Item C & 22.1 & 12.7 & 34.8 \\
\bottomrule
\end{tabular}
\caption{Sample Data Table}
\label{tab:sample}
\end{table}

\section{Lists}
\subsection{Bullet Points}
\begin{itemize}
  \item First item
  \item Second item
  \item Third item with sub-items:
    \begin{itemize}
      \item Sub-item 1
      \item Sub-item 2
    \end{itemize}
\end{itemize}

\subsection{Numbered List}
\begin{enumerate}
  \item First step
  \item Second step
  \item Third step
\end{enumerate}

\section{Citations and References}
When you use the citation search feature, references will be automatically managed. For example, you might cite a paper like this: \cite{example2024}. All citations will be automatically added to your \texttt{cite.bib} file.

\section{Conclusion}
This template demonstrates basic \LaTeX{} features including mathematical formulas, tables, and lists. The bibliography system is now set up to work seamlessly with Bibby AI's citation management. Feel free to modify it for your needs!

\bibliographystyle{plain}
\bibliography{cite}

\end{document}\documentclass{article}
\usepackage[utf8]{inputenc}
\usepackage{booktabs}
\usepackage{amsmath}
\usepackage{natbib}

\title{My Document}
\author{Your Name}
\date{\today}

\begin{document}

\maketitle

\section{Introduction}
This is a sample document demonstrating various \LaTeX{} features including mathematical equations, tables, and formatting. You can cite papers using the citation search feature, and they will be automatically added to your \texttt{cite.bib} file.

\section{Mathematical Expressions}
Here are some example equations:

\subsection{Quadratic Formula}
The quadratic formula is given by:
\[x = \frac{-b \pm \sqrt{b^2 - 4ac}}{2a}\]

\subsection{Matrix Example}
A sample matrix:
\[
\begin{pmatrix}
a_{11} & a_{12} & a_{13} \\
a_{21} & a_{22} & a_{23} \\
a_{31} & a_{32} & a_{33}
\end{pmatrix}
\]

\section{Tables}
Below is a sample table showing data:

\begin{table}[h]
\centering
\begin{tabular}{@{}lrrr@{}}
\toprule
\textbf{Category} & \textbf{Value 1} & \textbf{Value 2} & \textbf{Total} \\
\midrule
Item A & 10.5 & 20.3 & 30.8 \\
Item B & 15.2 & 25.1 & 40.3 \\
Item C & 22.1 & 12.7 & 34.8 \\
\bottomrule
\end{tabular}
\caption{Sample Data Table}
\label{tab:sample}
\end{table}

\section{Lists}
\subsection{Bullet Points}
\begin{itemize}
  \item First item
  \item Second item
  \item Third item with sub-items:
    \begin{itemize}
      \item Sub-item 1
      \item Sub-item 2
    \end{itemize}
\end{itemize}

\subsection{Numbered List}
\begin{enumerate}
  \item First step
  \item Second step
  \item Third step
\end{enumerate}

\section{Citations and References}
When you use the citation search feature, references will be automatically managed. For example, you might cite a paper like this: \cite{example2024}. All citations will be automatically added to your \texttt{cite.bib} file.

\section{Conclusion}
This template demonstrates basic \LaTeX{} features including mathematical formulas, tables, and lists. The bibliography system is now set up to work seamlessly with Bibby AI's citation management. Feel free to modify it for your needs!

\bibliographystyle{plain}
\bibliography{cite}

\end{document}\documentclass{article}
\usepackage[utf8]{inputenc}
\usepackage{booktabs}
\usepackage{amsmath}
\usepackage{natbib}

\title{My Document}
\author{Your Name}
\date{\today}

\begin{document}

\maketitle

\section{Introduction}
This is a sample document demonstrating various \LaTeX{} features including mathematical equations, tables, and formatting. You can cite papers using the citation search feature, and they will be automatically added to your \texttt{cite.bib} file.

\section{Mathematical Expressions}
Here are some example equations:

\subsection{Quadratic Formula}
The quadratic formula is given by:
\[x = \frac{-b \pm \sqrt{b^2 - 4ac}}{2a}\]

\subsection{Matrix Example}
A sample matrix:
\[
\begin{pmatrix}
a_{11} & a_{12} & a_{13} \\
a_{21} & a_{22} & a_{23} \\
a_{31} & a_{32} & a_{33}
\end{pmatrix}
\]

\section{Tables}
Below is a sample table showing data:

\begin{table}[h]
\centering
\begin{tabular}{@{}lrrr@{}}
\toprule
\textbf{Category} & \textbf{Value 1} & \textbf{Value 2} & \textbf{Total} \\
\midrule
Item A & 10.5 & 20.3 & 30.8 \\
Item B & 15.2 & 25.1 & 40.3 \\
Item C & 22.1 & 12.7 & 34.8 \\
\bottomrule
\end{tabular}
\caption{Sample Data Table}
\label{tab:sample}
\end{table}

\section{Lists}
\subsection{Bullet Points}
\begin{itemize}
  \item First item
  \item Second item
  \item Third item with sub-items:
    \begin{itemize}
      \item Sub-item 1
      \item Sub-item 2
    \end{itemize}
\end{itemize}

\subsection{Numbered List}
\begin{enumerate}
  \item First step
  \item Second step
  \item Third step
\end{enumerate}

\section{Citations and References}
When you use the citation search feature, references will be automatically managed. For example, you might cite a paper like this: \cite{example2024}. All citations will be automatically added to your \texttt{cite.bib} file.

\section{Conclusion}
This template demonstrates basic \LaTeX{} features including mathematical formulas, tables, and lists. The bibliography system is now set up to work seamlessly with Bibby AI's citation management. Feel free to modify it for your needs!

\bibliographystyle{plain}
\bibliography{cite}

\end{document}\documentclass{article}
\usepackage[utf8]{inputenc}
\usepackage{booktabs}
\usepackage{amsmath}
\usepackage{natbib}

\title{My Document}
\author{Your Name}
\date{\today}

\begin{document}

\maketitle

\section{Introduction}
This is a sample document demonstrating various \LaTeX{} features including mathematical equations, tables, and formatting. You can cite papers using the citation search feature, and they will be automatically added to your \texttt{cite.bib} file.

\section{Mathematical Expressions}
Here are some example equations:

\subsection{Quadratic Formula}
The quadratic formula is given by:
\[x = \frac{-b \pm \sqrt{b^2 - 4ac}}{2a}\]

\subsection{Matrix Example}
A sample matrix:
\[
\begin{pmatrix}
a_{11} & a_{12} & a_{13} \\
a_{21} & a_{22} & a_{23} \\
a_{31} & a_{32} & a_{33}
\end{pmatrix}
\]

\section{Tables}
Below is a sample table showing data:

\begin{table}[h]
\centering
\begin{tabular}{@{}lrrr@{}}
\toprule
\textbf{Category} & \textbf{Value 1} & \textbf{Value 2} & \textbf{Total} \\
\midrule
Item A & 10.5 & 20.3 & 30.8 \\
Item B & 15.2 & 25.1 & 40.3 \\
Item C & 22.1 & 12.7 & 34.8 \\
\bottomrule
\end{tabular}
\caption{Sample Data Table}
\label{tab:sample}
\end{table}

\section{Lists}
\subsection{Bullet Points}
\begin{itemize}
  \item First item
  \item Second item
  \item Third item with sub-items:
    \begin{itemize}
      \item Sub-item 1
      \item Sub-item 2
    \end{itemize}
\end{itemize}

\subsection{Numbered List}
\begin{enumerate}
  \item First step
  \item Second step
  \item Third step
\end{enumerate}

\section{Citations and References}
When you use the citation search feature, references will be automatically managed. For example, you might cite a paper like this: \cite{example2024}. All citations will be automatically added to your \texttt{cite.bib} file.

\section{Conclusion}
This template demonstrates basic \LaTeX{} features including mathematical formulas, tables, and lists. The bibliography system is now set up to work seamlessly with Bibby AI's citation management. Feel free to modify it for your needs!

\bibliographystyle{plain}
\bibliography{cite}

\end{document}\documentclass{article}
\usepackage[utf8]{inputenc}
\usepackage{booktabs}
\usepackage{amsmath}
\usepackage{natbib}

\title{My Document}
\author{Your Name}
\date{\today}

\begin{document}

\maketitle

\section{Introduction}
This is a sample document demonstrating various \LaTeX{} features including mathematical equations, tables, and formatting. You can cite papers using the citation search feature, and they will be automatically added to your \texttt{cite.bib} file.

\section{Mathematical Expressions}
Here are some example equations:

\subsection{Quadratic Formula}
The quadratic formula is given by:
\[x = \frac{-b \pm \sqrt{b^2 - 4ac}}{2a}\]

\subsection{Matrix Example}
A sample matrix:
\[
\begin{pmatrix}
a_{11} & a_{12} & a_{13} \\
a_{21} & a_{22} & a_{23} \\
a_{31} & a_{32} & a_{33}
\end{pmatrix}
\]

\section{Tables}
Below is a sample table showing data:

\begin{table}[h]
\centering
\begin{tabular}{@{}lrrr@{}}
\toprule
\textbf{Category} & \textbf{Value 1} & \textbf{Value 2} & \textbf{Total} \\
\midrule
Item A & 10.5 & 20.3 & 30.8 \\
Item B & 15.2 & 25.1 & 40.3 \\
Item C & 22.1 & 12.7 & 34.8 \\
\bottomrule
\end{tabular}
\caption{Sample Data Table}
\label{tab:sample}
\end{table}

\section{Lists}
\subsection{Bullet Points}
\begin{itemize}
  \item First item
  \item Second item
  \item Third item with sub-items:
    \begin{itemize}
      \item Sub-item 1
      \item Sub-item 2
    \end{itemize}
\end{itemize}

\subsection{Numbered List}
\begin{enumerate}
  \item First step
  \item Second step
  \item Third step
\end{enumerate}

\section{Citations and References}
When you use the citation search feature, references will be automatically managed. For example, you might cite a paper like this: \cite{example2024}. All citations will be automatically added to your \texttt{cite.bib} file.

\section{Conclusion}
This template demonstrates basic \LaTeX{} features including mathematical formulas, tables, and lists. The bibliography system is now set up to work seamlessly with Bibby AI's citation management. Feel free to modify it for your needs!

\bibliographystyle{plain}
\bibliography{cite}

\end{document}\documentclass{article}
\usepackage[utf8]{inputenc}
\usepackage{booktabs}
\usepackage{amsmath}
\usepackage{natbib}

\title{My Document}
\author{Your Name}
\date{\today}

\begin{document}

\maketitle

\section{Introduction}
This is a sample document demonstrating various \LaTeX{} features including mathematical equations, tables, and formatting. You can cite papers using the citation search feature, and they will be automatically added to your \texttt{cite.bib} file.

\section{Mathematical Expressions}
Here are some example equations:

\subsection{Quadratic Formula}
The quadratic formula is given by:
\[x = \frac{-b \pm \sqrt{b^2 - 4ac}}{2a}\]

\subsection{Matrix Example}
A sample matrix:
\[
\begin{pmatrix}
a_{11} & a_{12} & a_{13} \\
a_{21} & a_{22} & a_{23} \\
a_{31} & a_{32} & a_{33}
\end{pmatrix}
\]

\section{Tables}
Below is a sample table showing data:

\begin{table}[h]
\centering
\begin{tabular}{@{}lrrr@{}}
\toprule
\textbf{Category} & \textbf{Value 1} & \textbf{Value 2} & \textbf{Total} \\
\midrule
Item A & 10.5 & 20.3 & 30.8 \\
Item B & 15.2 & 25.1 & 40.3 \\
Item C & 22.1 & 12.7 & 34.8 \\
\bottomrule
\end{tabular}
\caption{Sample Data Table}
\label{tab:sample}
\end{table}

\section{Lists}
\subsection{Bullet Points}
\begin{itemize}
  \item First item
  \item Second item
  \item Third item with sub-items:
    \begin{itemize}
      \item Sub-item 1
      \item Sub-item 2
    \end{itemize}
\end{itemize}

\subsection{Numbered List}
\begin{enumerate}
  \item First step
  \item Second step
  \item Third step
\end{enumerate}

\section{Citations and References}
When you use the citation search feature, references will be automatically managed. For example, you might cite a paper like this: \cite{example2024}. All citations will be automatically added to your \texttt{cite.bib} file.

\section{Conclusion}
This template demonstrates basic \LaTeX{} features including mathematical formulas, tables, and lists. The bibliography system is now set up to work seamlessly with Bibby AI's citation management. Feel free to modify it for your needs!

\bibliographystyle{plain}
\bibliography{cite}

\end{document}\documentclass{article}
\usepackage[utf8]{inputenc}
\usepackage{booktabs}
\usepackage{amsmath}
\usepackage{natbib}

\title{My Document}
\author{Your Name}
\date{\today}

\begin{document}

\maketitle

\section{Introduction}
This is a sample document demonstrating various \LaTeX{} features including mathematical equations, tables, and formatting. You can cite papers using the citation search feature, and they will be automatically added to your \texttt{cite.bib} file.

\section{Mathematical Expressions}
Here are some example equations:

\subsection{Quadratic Formula}
The quadratic formula is given by:
\[x = \frac{-b \pm \sqrt{b^2 - 4ac}}{2a}\]

\subsection{Matrix Example}
A sample matrix:
\[
\begin{pmatrix}
a_{11} & a_{12} & a_{13} \\
a_{21} & a_{22} & a_{23} \\
a_{31} & a_{32} & a_{33}
\end{pmatrix}
\]

\section{Tables}
Below is a sample table showing data:

\begin{table}[h]
\centering
\begin{tabular}{@{}lrrr@{}}
\toprule
\textbf{Category} & \textbf{Value 1} & \textbf{Value 2} & \textbf{Total} \\
\midrule
Item A & 10.5 & 20.3 & 30.8 \\
Item B & 15.2 & 25.1 & 40.3 \\
Item C & 22.1 & 12.7 & 34.8 \\
\bottomrule
\end{tabular}
\caption{Sample Data Table}
\label{tab:sample}
\end{table}

\section{Lists}
\subsection{Bullet Points}
\begin{itemize}
  \item First item
  \item Second item
  \item Third item with sub-items:
    \begin{itemize}
      \item Sub-item 1
      \item Sub-item 2
    \end{itemize}
\end{itemize}

\subsection{Numbered List}
\begin{enumerate}
  \item First step
  \item Second step
  \item Third step
\end{enumerate}

\section{Citations and References}
When you use the citation search feature, references will be automatically managed. For example, you might cite a paper like this: \cite{example2024}. All citations will be automatically added to your \texttt{cite.bib} file.

\section{Conclusion}
This template demonstrates basic \LaTeX{} features including mathematical formulas, tables, and lists. The bibliography system is now set up to work seamlessly with Bibby AI's citation management. Feel free to modify it for your needs!

\bibliographystyle{plain}
\bibliography{cite}

\end{document}\documentclass{article}
\usepackage[utf8]{inputenc}
\usepackage{booktabs}
\usepackage{amsmath}
\usepackage{natbib}

\title{My Document}
\author{Your Name}
\date{\today}

\begin{document}

\maketitle

\section{Introduction}
This is a sample document demonstrating various \LaTeX{} features including mathematical equations, tables, and formatting. You can cite papers using the citation search feature, and they will be automatically added to your \texttt{cite.bib} file.

\section{Mathematical Expressions}
Here are some example equations:

\subsection{Quadratic Formula}
The quadratic formula is given by:
\[x = \frac{-b \pm \sqrt{b^2 - 4ac}}{2a}\]

\subsection{Matrix Example}
A sample matrix:
\[
\begin{pmatrix}
a_{11} & a_{12} & a_{13} \\
a_{21} & a_{22} & a_{23} \\
a_{31} & a_{32} & a_{33}
\end{pmatrix}
\]

\section{Tables}
Below is a sample table showing data:

\begin{table}[h]
\centering
\begin{tabular}{@{}lrrr@{}}
\toprule
\textbf{Category} & \textbf{Value 1} & \textbf{Value 2} & \textbf{Total} \\
\midrule
Item A & 10.5 & 20.3 & 30.8 \\
Item B & 15.2 & 25.1 & 40.3 \\
Item C & 22.1 & 12.7 & 34.8 \\
\bottomrule
\end{tabular}
\caption{Sample Data Table}
\label{tab:sample}
\end{table}

\section{Lists}
\subsection{Bullet Points}
\begin{itemize}
  \item First item
  \item Second item
  \item Third item with sub-items:
    \begin{itemize}
      \item Sub-item 1
      \item Sub-item 2
    \end{itemize}
\end{itemize}

\subsection{Numbered List}
\begin{enumerate}
  \item First step
  \item Second step
  \item Third step
\end{enumerate}

\section{Citations and References}
When you use the citation search feature, references will be automatically managed. For example, you might cite a paper like this: \cite{example2024}. All citations will be automatically added to your \texttt{cite.bib} file.

\section{Conclusion}
This template demonstrates basic \LaTeX{} features including mathematical formulas, tables, and lists. The bibliography system is now set up to work seamlessly with Bibby AI's citation management. Feel free to modify it for your needs!

\bibliographystyle{plain}
\bibliography{cite}

\end{document}\documentclass{article}
\usepackage[utf8]{inputenc}
\usepackage{booktabs}
\usepackage{amsmath}
\usepackage{natbib}

\title{My Document}
\author{Your Name}
\date{\today}

\begin{document}

\maketitle

\section{Introduction}
This is a sample document demonstrating various \LaTeX{} features including mathematical equations, tables, and formatting. You can cite papers using the citation search feature, and they will be automatically added to your \texttt{cite.bib} file.

\section{Mathematical Expressions}
Here are some example equations:

\subsection{Quadratic Formula}
The quadratic formula is given by:
\[x = \frac{-b \pm \sqrt{b^2 - 4ac}}{2a}\]

\subsection{Matrix Example}
A sample matrix:
\[
\begin{pmatrix}
a_{11} & a_{12} & a_{13} \\
a_{21} & a_{22} & a_{23} \\
a_{31} & a_{32} & a_{33}
\end{pmatrix}
\]

\section{Tables}
Below is a sample table showing data:

\begin{table}[h]
\centering
\begin{tabular}{@{}lrrr@{}}
\toprule
\textbf{Category} & \textbf{Value 1} & \textbf{Value 2} & \textbf{Total} \\
\midrule
Item A & 10.5 & 20.3 & 30.8 \\
Item B & 15.2 & 25.1 & 40.3 \\
Item C & 22.1 & 12.7 & 34.8 \\
\bottomrule
\end{tabular}
\caption{Sample Data Table}
\label{tab:sample}
\end{table}

\section{Lists}
\subsection{Bullet Points}
\begin{itemize}
  \item First item
  \item Second item
  \item Third item with sub-items:
    \begin{itemize}
      \item Sub-item 1
      \item Sub-item 2
    \end{itemize}
\end{itemize}

\subsection{Numbered List}
\begin{enumerate}
  \item First step
  \item Second step
  \item Third step
\end{enumerate}

\section{Citations and References}
When you use the citation search feature, references will be automatically managed. For example, you might cite a paper like this: \cite{example2024}. All citations will be automatically added to your \texttt{cite.bib} file.

\section{Conclusion}
This template demonstrates basic \LaTeX{} features including mathematical formulas, tables, and lists. The bibliography system is now set up to work seamlessly with Bibby AI's citation management. Feel free to modify it for your needs!

\bibliographystyle{plain}
\bibliography{cite}

\end{document}\documentclass{article}
\usepackage[utf8]{inputenc}
\usepackage{booktabs}
\usepackage{amsmath}
\usepackage{natbib}

\title{My Document}
\author{Your Name}
\date{\today}

\begin{document}

\maketitle

\section{Introduction}
This is a sample document demonstrating various \LaTeX{} features including mathematical equations, tables, and formatting. You can cite papers using the citation search feature, and they will be automatically added to your \texttt{cite.bib} file.

\section{Mathematical Expressions}
Here are some example equations:

\subsection{Quadratic Formula}
The quadratic formula is given by:
\[x = \frac{-b \pm \sqrt{b^2 - 4ac}}{2a}\]

\subsection{Matrix Example}
A sample matrix:
\[
\begin{pmatrix}
a_{11} & a_{12} & a_{13} \\
a_{21} & a_{22} & a_{23} \\
a_{31} & a_{32} & a_{33}
\end{pmatrix}
\]

\section{Tables}
Below is a sample table showing data:

\begin{table}[h]
\centering
\begin{tabular}{@{}lrrr@{}}
\toprule
\textbf{Category} & \textbf{Value 1} & \textbf{Value 2} & \textbf{Total} \\
\midrule
Item A & 10.5 & 20.3 & 30.8 \\
Item B & 15.2 & 25.1 & 40.3 \\
Item C & 22.1 & 12.7 & 34.8 \\
\bottomrule
\end{tabular}
\caption{Sample Data Table}
\label{tab:sample}
\end{table}

\section{Lists}
\subsection{Bullet Points}
\begin{itemize}
  \item First item
  \item Second item
  \item Third item with sub-items:
    \begin{itemize}
      \item Sub-item 1
      \item Sub-item 2
    \end{itemize}
\end{itemize}

\subsection{Numbered List}
\begin{enumerate}
  \item First step
  \item Second step
  \item Third step
\end{enumerate}

\section{Citations and References}
When you use the citation search feature, references will be automatically managed. For example, you might cite a paper like this: \cite{example2024}. All citations will be automatically added to your \texttt{cite.bib} file.

\section{Conclusion}
This template demonstrates basic \LaTeX{} features including mathematical formulas, tables, and lists. The bibliography system is now set up to work seamlessly with Bibby AI's citation management. Feel free to modify it for your needs!

\bibliographystyle{plain}
\bibliography{cite}

\end{document}\documentclass{article}
\usepackage[utf8]{inputenc}
\usepackage{booktabs}
\usepackage{amsmath}
\usepackage{natbib}

\title{My Document}
\author{Your Name}
\date{\today}

\begin{document}

\maketitle

\section{Introduction}
This is a sample document demonstrating various \LaTeX{} features including mathematical equations, tables, and formatting. You can cite papers using the citation search feature, and they will be automatically added to your \texttt{cite.bib} file.

\section{Mathematical Expressions}
Here are some example equations:

\subsection{Quadratic Formula}
The quadratic formula is given by:
\[x = \frac{-b \pm \sqrt{b^2 - 4ac}}{2a}\]

\subsection{Matrix Example}
A sample matrix:
\[
\begin{pmatrix}
a_{11} & a_{12} & a_{13} \\
a_{21} & a_{22} & a_{23} \\
a_{31} & a_{32} & a_{33}
\end{pmatrix}
\]

\section{Tables}
Below is a sample table showing data:

\begin{table}[h]
\centering
\begin{tabular}{@{}lrrr@{}}
\toprule
\textbf{Category} & \textbf{Value 1} & \textbf{Value 2} & \textbf{Total} \\
\midrule
Item A & 10.5 & 20.3 & 30.8 \\
Item B & 15.2 & 25.1 & 40.3 \\
Item C & 22.1 & 12.7 & 34.8 \\
\bottomrule
\end{tabular}
\caption{Sample Data Table}
\label{tab:sample}
\end{table}

\section{Lists}
\subsection{Bullet Points}
\begin{itemize}
  \item First item
  \item Second item
  \item Third item with sub-items:
    \begin{itemize}
      \item Sub-item 1
      \item Sub-item 2
    \end{itemize}
\end{itemize}

\subsection{Numbered List}
\begin{enumerate}
  \item First step
  \item Second step
  \item Third step
\end{enumerate}

\section{Citations and References}
When you use the citation search feature, references will be automatically managed. For example, you might cite a paper like this: \cite{example2024}. All citations will be automatically added to your \texttt{cite.bib} file.

\section{Conclusion}
This template demonstrates basic \LaTeX{} features including mathematical formulas, tables, and lists. The bibliography system is now set up to work seamlessly with Bibby AI's citation management. Feel free to modify it for your needs!

\bibliographystyle{plain}
\bibliography{cite}

\end{document}\documentclass{article}
\usepackage[utf8]{inputenc}
\usepackage{booktabs}
\usepackage{amsmath}
\usepackage{natbib}

\title{My Document}
\author{Your Name}
\date{\today}

\begin{document}

\maketitle

\section{Introduction}
This is a sample document demonstrating various \LaTeX{} features including mathematical equations, tables, and formatting. You can cite papers using the citation search feature, and they will be automatically added to your \texttt{cite.bib} file.

\section{Mathematical Expressions}
Here are some example equations:

\subsection{Quadratic Formula}
The quadratic formula is given by:
\[x = \frac{-b \pm \sqrt{b^2 - 4ac}}{2a}\]

\subsection{Matrix Example}
A sample matrix:
\[
\begin{pmatrix}
a_{11} & a_{12} & a_{13} \\
a_{21} & a_{22} & a_{23} \\
a_{31} & a_{32} & a_{33}
\end{pmatrix}
\]

\section{Tables}
Below is a sample table showing data:

\begin{table}[h]
\centering
\begin{tabular}{@{}lrrr@{}}
\toprule
\textbf{Category} & \textbf{Value 1} & \textbf{Value 2} & \textbf{Total} \\
\midrule
Item A & 10.5 & 20.3 & 30.8 \\
Item B & 15.2 & 25.1 & 40.3 \\
Item C & 22.1 & 12.7 & 34.8 \\
\bottomrule
\end{tabular}
\caption{Sample Data Table}
\label{tab:sample}
\end{table}

\section{Lists}
\subsection{Bullet Points}
\begin{itemize}
  \item First item
  \item Second item
  \item Third item with sub-items:
    \begin{itemize}
      \item Sub-item 1
      \item Sub-item 2
    \end{itemize}
\end{itemize}

\subsection{Numbered List}
\begin{enumerate}
  \item First step
  \item Second step
  \item Third step
\end{enumerate}

\section{Citations and References}
When you use the citation search feature, references will be automatically managed. For example, you might cite a paper like this: \cite{example2024}. All citations will be automatically added to your \texttt{cite.bib} file.

\section{Conclusion}
This template demonstrates basic \LaTeX{} features including mathematical formulas, tables, and lists. The bibliography system is now set up to work seamlessly with Bibby AI's citation management. Feel free to modify it for your needs!

\bibliographystyle{plain}
\bibliography{cite}

\end{document}\documentclass{article}
\usepackage[utf8]{inputenc}
\usepackage{booktabs}
\usepackage{amsmath}
\usepackage{natbib}

\title{My Document}
\author{Your Name}
\date{\today}

\begin{document}

\maketitle

\section{Introduction}
This is a sample document demonstrating various \LaTeX{} features including mathematical equations, tables, and formatting. You can cite papers using the citation search feature, and they will be automatically added to your \texttt{cite.bib} file.

\section{Mathematical Expressions}
Here are some example equations:

\subsection{Quadratic Formula}
The quadratic formula is given by:
\[x = \frac{-b \pm \sqrt{b^2 - 4ac}}{2a}\]

\subsection{Matrix Example}
A sample matrix:
\[
\begin{pmatrix}
a_{11} & a_{12} & a_{13} \\
a_{21} & a_{22} & a_{23} \\
a_{31} & a_{32} & a_{33}
\end{pmatrix}
\]

\section{Tables}
Below is a sample table showing data:

\begin{table}[h]
\centering
\begin{tabular}{@{}lrrr@{}}
\toprule
\textbf{Category} & \textbf{Value 1} & \textbf{Value 2} & \textbf{Total} \\
\midrule
Item A & 10.5 & 20.3 & 30.8 \\
Item B & 15.2 & 25.1 & 40.3 \\
Item C & 22.1 & 12.7 & 34.8 \\
\bottomrule
\end{tabular}
\caption{Sample Data Table}
\label{tab:sample}
\end{table}

\section{Lists}
\subsection{Bullet Points}
\begin{itemize}
  \item First item
  \item Second item
  \item Third item with sub-items:
    \begin{itemize}
      \item Sub-item 1
      \item Sub-item 2
    \end{itemize}
\end{itemize}

\subsection{Numbered List}
\begin{enumerate}
  \item First step
  \item Second step
  \item Third step
\end{enumerate}

\section{Citations and References}
When you use the citation search feature, references will be automatically managed. For example, you might cite a paper like this: \cite{example2024}. All citations will be automatically added to your \texttt{cite.bib} file.

\section{Conclusion}
This template demonstrates basic \LaTeX{} features including mathematical formulas, tables, and lists. The bibliography system is now set up to work seamlessly with Bibby AI's citation management. Feel free to modify it for your needs!

\bibliographystyle{plain}
\bibliography{cite}

\end{document}\documentclass{article}
\usepackage[utf8]{inputenc}
\usepackage{booktabs}
\usepackage{amsmath}
\usepackage{natbib}

\title{My Document}
\author{Your Name}
\date{\today}

\begin{document}

\maketitle

\section{Introduction}
This is a sample document demonstrating various \LaTeX{} features including mathematical equations, tables, and formatting. You can cite papers using the citation search feature, and they will be automatically added to your \texttt{cite.bib} file.

\section{Mathematical Expressions}
Here are some example equations:

\subsection{Quadratic Formula}
The quadratic formula is given by:
\[x = \frac{-b \pm \sqrt{b^2 - 4ac}}{2a}\]

\subsection{Matrix Example}
A sample matrix:
\[
\begin{pmatrix}
a_{11} & a_{12} & a_{13} \\
a_{21} & a_{22} & a_{23} \\
a_{31} & a_{32} & a_{33}
\end{pmatrix}
\]

\section{Tables}
Below is a sample table showing data:

\begin{table}[h]
\centering
\begin{tabular}{@{}lrrr@{}}
\toprule
\textbf{Category} & \textbf{Value 1} & \textbf{Value 2} & \textbf{Total} \\
\midrule
Item A & 10.5 & 20.3 & 30.8 \\
Item B & 15.2 & 25.1 & 40.3 \\
Item C & 22.1 & 12.7 & 34.8 \\
\bottomrule
\end{tabular}
\caption{Sample Data Table}
\label{tab:sample}
\end{table}

\section{Lists}
\subsection{Bullet Points}
\begin{itemize}
  \item First item
  \item Second item
  \item Third item with sub-items:
    \begin{itemize}
      \item Sub-item 1
      \item Sub-item 2
    \end{itemize}
\end{itemize}

\subsection{Numbered List}
\begin{enumerate}
  \item First step
  \item Second step
  \item Third step
\end{enumerate}

\section{Citations and References}
When you use the citation search feature, references will be automatically managed. For example, you might cite a paper like this: \cite{example2024}. All citations will be automatically added to your \texttt{cite.bib} file.

\section{Conclusion}
This template demonstrates basic \LaTeX{} features including mathematical formulas, tables, and lists. The bibliography system is now set up to work seamlessly with Bibby AI's citation management. Feel free to modify it for your needs!

\bibliographystyle{plain}
\bibliography{cite}

\end{document}\documentclass{article}
\usepackage[utf8]{inputenc}
\usepackage{booktabs}
\usepackage{amsmath}
\usepackage{natbib}

\title{My Document}
\author{Your Name}
\date{\today}

\begin{document}

\maketitle

\section{Introduction}
This is a sample document demonstrating various \LaTeX{} features including mathematical equations, tables, and formatting. You can cite papers using the citation search feature, and they will be automatically added to your \texttt{cite.bib} file.

\section{Mathematical Expressions}
Here are some example equations:

\subsection{Quadratic Formula}
The quadratic formula is given by:
\[x = \frac{-b \pm \sqrt{b^2 - 4ac}}{2a}\]

\subsection{Matrix Example}
A sample matrix:
\[
\begin{pmatrix}
a_{11} & a_{12} & a_{13} \\
a_{21} & a_{22} & a_{23} \\
a_{31} & a_{32} & a_{33}
\end{pmatrix}
\]

\section{Tables}
Below is a sample table showing data:

\begin{table}[h]
\centering
\begin{tabular}{@{}lrrr@{}}
\toprule
\textbf{Category} & \textbf{Value 1} & \textbf{Value 2} & \textbf{Total} \\
\midrule
Item A & 10.5 & 20.3 & 30.8 \\
Item B & 15.2 & 25.1 & 40.3 \\
Item C & 22.1 & 12.7 & 34.8 \\
\bottomrule
\end{tabular}
\caption{Sample Data Table}
\label{tab:sample}
\end{table}

\section{Lists}
\subsection{Bullet Points}
\begin{itemize}
  \item First item
  \item Second item
  \item Third item with sub-items:
    \begin{itemize}
      \item Sub-item 1
      \item Sub-item 2
    \end{itemize}
\end{itemize}

\subsection{Numbered List}
\begin{enumerate}
  \item First step
  \item Second step
  \item Third step
\end{enumerate}

\section{Citations and References}
When you use the citation search feature, references will be automatically managed. For example, you might cite a paper like this: \cite{example2024}. All citations will be automatically added to your \texttt{cite.bib} file.

\section{Conclusion}
This template demonstrates basic \LaTeX{} features including mathematical formulas, tables, and lists. The bibliography system is now set up to work seamlessly with Bibby AI's citation management. Feel free to modify it for your needs!

\bibliographystyle{plain}
\bibliography{cite}

\end{document}\documentclass{article}
\usepackage[utf8]{inputenc}
\usepackage{booktabs}
\usepackage{amsmath}
\usepackage{natbib}

\title{My Document}
\author{Your Name}
\date{\today}

\begin{document}

\maketitle

\section{Introduction}
This is a sample document demonstrating various \LaTeX{} features including mathematical equations, tables, and formatting. You can cite papers using the citation search feature, and they will be automatically added to your \texttt{cite.bib} file.

\section{Mathematical Expressions}
Here are some example equations:

\subsection{Quadratic Formula}
The quadratic formula is given by:
\[x = \frac{-b \pm \sqrt{b^2 - 4ac}}{2a}\]

\subsection{Matrix Example}
A sample matrix:
\[
\begin{pmatrix}
a_{11} & a_{12} & a_{13} \\
a_{21} & a_{22} & a_{23} \\
a_{31} & a_{32} & a_{33}
\end{pmatrix}
\]

\section{Tables}
Below is a sample table showing data:

\begin{table}[h]
\centering
\begin{tabular}{@{}lrrr@{}}
\toprule
\textbf{Category} & \textbf{Value 1} & \textbf{Value 2} & \textbf{Total} \\
\midrule
Item A & 10.5 & 20.3 & 30.8 \\
Item B & 15.2 & 25.1 & 40.3 \\
Item C & 22.1 & 12.7 & 34.8 \\
\bottomrule
\end{tabular}
\caption{Sample Data Table}
\label{tab:sample}
\end{table}

\section{Lists}
\subsection{Bullet Points}
\begin{itemize}
  \item First item
  \item Second item
  \item Third item with sub-items:
    \begin{itemize}
      \item Sub-item 1
      \item Sub-item 2
    \end{itemize}
\end{itemize}

\subsection{Numbered List}
\begin{enumerate}
  \item First step
  \item Second step
  \item Third step
\end{enumerate}

\section{Citations and References}
When you use the citation search feature, references will be automatically managed. For example, you might cite a paper like this: \cite{example2024}. All citations will be automatically added to your \texttt{cite.bib} file.

\section{Conclusion}
This template demonstrates basic \LaTeX{} features including mathematical formulas, tables, and lists. The bibliography system is now set up to work seamlessly with Bibby AI's citation management. Feel free to modify it for your needs!

\bibliographystyle{plain}
\bibliography{cite}

\end{document}\documentclass{article}
\usepackage[utf8]{inputenc}
\usepackage{booktabs}
\usepackage{amsmath}
\usepackage{natbib}

\title{My Document}
\author{Your Name}
\date{\today}

\begin{document}

\maketitle

\section{Introduction}
This is a sample document demonstrating various \LaTeX{} features including mathematical equations, tables, and formatting. You can cite papers using the citation search feature, and they will be automatically added to your \texttt{cite.bib} file.

\section{Mathematical Expressions}
Here are some example equations:

\subsection{Quadratic Formula}
The quadratic formula is given by:
\[x = \frac{-b \pm \sqrt{b^2 - 4ac}}{2a}\]

\subsection{Matrix Example}
A sample matrix:
\[
\begin{pmatrix}
a_{11} & a_{12} & a_{13} \\
a_{21} & a_{22} & a_{23} \\
a_{31} & a_{32} & a_{33}
\end{pmatrix}
\]

\section{Tables}
Below is a sample table showing data:

\begin{table}[h]
\centering
\begin{tabular}{@{}lrrr@{}}
\toprule
\textbf{Category} & \textbf{Value 1} & \textbf{Value 2} & \textbf{Total} \\
\midrule
Item A & 10.5 & 20.3 & 30.8 \\
Item B & 15.2 & 25.1 & 40.3 \\
Item C & 22.1 & 12.7 & 34.8 \\
\bottomrule
\end{tabular}
\caption{Sample Data Table}
\label{tab:sample}
\end{table}

\section{Lists}
\subsection{Bullet Points}
\begin{itemize}
  \item First item
  \item Second item
  \item Third item with sub-items:
    \begin{itemize}
      \item Sub-item 1
      \item Sub-item 2
    \end{itemize}
\end{itemize}

\subsection{Numbered List}
\begin{enumerate}
  \item First step
  \item Second step
  \item Third step
\end{enumerate}

\section{Citations and References}
When you use the citation search feature, references will be automatically managed. For example, you might cite a paper like this: \cite{example2024}. All citations will be automatically added to your \texttt{cite.bib} file.

\section{Conclusion}
This template demonstrates basic \LaTeX{} features including mathematical formulas, tables, and lists. The bibliography system is now set up to work seamlessly with Bibby AI's citation management. Feel free to modify it for your needs!

\bibliographystyle{plain}
\bibliography{cite}

\end{document}\documentclass{article}
\usepackage[utf8]{inputenc}
\usepackage{booktabs}
\usepackage{amsmath}
\usepackage{natbib}

\title{My Document}
\author{Your Name}
\date{\today}

\begin{document}

\maketitle

\section{Introduction}
This is a sample document demonstrating various \LaTeX{} features including mathematical equations, tables, and formatting. You can cite papers using the citation search feature, and they will be automatically added to your \texttt{cite.bib} file.

\section{Mathematical Expressions}
Here are some example equations:

\subsection{Quadratic Formula}
The quadratic formula is given by:
\[x = \frac{-b \pm \sqrt{b^2 - 4ac}}{2a}\]

\subsection{Matrix Example}
A sample matrix:
\[
\begin{pmatrix}
a_{11} & a_{12} & a_{13} \\
a_{21} & a_{22} & a_{23} \\
a_{31} & a_{32} & a_{33}
\end{pmatrix}
\]

\section{Tables}
Below is a sample table showing data:

\begin{table}[h]
\centering
\begin{tabular}{@{}lrrr@{}}
\toprule
\textbf{Category} & \textbf{Value 1} & \textbf{Value 2} & \textbf{Total} \\
\midrule
Item A & 10.5 & 20.3 & 30.8 \\
Item B & 15.2 & 25.1 & 40.3 \\
Item C & 22.1 & 12.7 & 34.8 \\
\bottomrule
\end{tabular}
\caption{Sample Data Table}
\label{tab:sample}
\end{table}

\section{Lists}
\subsection{Bullet Points}
\begin{itemize}
  \item First item
  \item Second item
  \item Third item with sub-items:
    \begin{itemize}
      \item Sub-item 1
      \item Sub-item 2
    \end{itemize}
\end{itemize}

\subsection{Numbered List}
\begin{enumerate}
  \item First step
  \item Second step
  \item Third step
\end{enumerate}

\section{Citations and References}
When you use the citation search feature, references will be automatically managed. For example, you might cite a paper like this: \cite{example2024}. All citations will be automatically added to your \texttt{cite.bib} file.

\section{Conclusion}
This template demonstrates basic \LaTeX{} features including mathematical formulas, tables, and lists. The bibliography system is now set up to work seamlessly with Bibby AI's citation management. Feel free to modify it for your needs!

\bibliographystyle{plain}
\bibliography{cite}

\end{document}\documentclass{article}
\usepackage[utf8]{inputenc}
\usepackage{booktabs}
\usepackage{amsmath}
\usepackage{natbib}

\title{My Document}
\author{Your Name}
\date{\today}

\begin{document}

\maketitle

\section{Introduction}
This is a sample document demonstrating various \LaTeX{} features including mathematical equations, tables, and formatting. You can cite papers using the citation search feature, and they will be automatically added to your \texttt{cite.bib} file.

\section{Mathematical Expressions}
Here are some example equations:

\subsection{Quadratic Formula}
The quadratic formula is given by:
\[x = \frac{-b \pm \sqrt{b^2 - 4ac}}{2a}\]

\subsection{Matrix Example}
A sample matrix:
\[
\begin{pmatrix}
a_{11} & a_{12} & a_{13} \\
a_{21} & a_{22} & a_{23} \\
a_{31} & a_{32} & a_{33}
\end{pmatrix}
\]

\section{Tables}
Below is a sample table showing data:

\begin{table}[h]
\centering
\begin{tabular}{@{}lrrr@{}}
\toprule
\textbf{Category} & \textbf{Value 1} & \textbf{Value 2} & \textbf{Total} \\
\midrule
Item A & 10.5 & 20.3 & 30.8 \\
Item B & 15.2 & 25.1 & 40.3 \\
Item C & 22.1 & 12.7 & 34.8 \\
\bottomrule
\end{tabular}
\caption{Sample Data Table}
\label{tab:sample}
\end{table}

\section{Lists}
\subsection{Bullet Points}
\begin{itemize}
  \item First item
  \item Second item
  \item Third item with sub-items:
    \begin{itemize}
      \item Sub-item 1
      \item Sub-item 2
    \end{itemize}
\end{itemize}

\subsection{Numbered List}
\begin{enumerate}
  \item First step
  \item Second step
  \item Third step
\end{enumerate}

\section{Citations and References}
When you use the citation search feature, references will be automatically managed. For example, you might cite a paper like this: \cite{example2024}. All citations will be automatically added to your \texttt{cite.bib} file.

\section{Conclusion}
This template demonstrates basic \LaTeX{} features including mathematical formulas, tables, and lists. The bibliography system is now set up to work seamlessly with Bibby AI's citation management. Feel free to modify it for your needs!

\bibliographystyle{plain}
\bibliography{cite}

\end{document}\documentclass{article}
\usepackage[utf8]{inputenc}
\usepackage{booktabs}
\usepackage{amsmath}
\usepackage{natbib}

\title{My Document}
\author{Your Name}
\date{\today}

\begin{document}

\maketitle

\section{Introduction}
This is a sample document demonstrating various \LaTeX{} features including mathematical equations, tables, and formatting. You can cite papers using the citation search feature, and they will be automatically added to your \texttt{cite.bib} file.

\section{Mathematical Expressions}
Here are some example equations:

\subsection{Quadratic Formula}
The quadratic formula is given by:
\[x = \frac{-b \pm \sqrt{b^2 - 4ac}}{2a}\]

\subsection{Matrix Example}
A sample matrix:
\[
\begin{pmatrix}
a_{11} & a_{12} & a_{13} \\
a_{21} & a_{22} & a_{23} \\
a_{31} & a_{32} & a_{33}
\end{pmatrix}
\]

\section{Tables}
Below is a sample table showing data:

\begin{table}[h]
\centering
\begin{tabular}{@{}lrrr@{}}
\toprule
\textbf{Category} & \textbf{Value 1} & \textbf{Value 2} & \textbf{Total} \\
\midrule
Item A & 10.5 & 20.3 & 30.8 \\
Item B & 15.2 & 25.1 & 40.3 \\
Item C & 22.1 & 12.7 & 34.8 \\
\bottomrule
\end{tabular}
\caption{Sample Data Table}
\label{tab:sample}
\end{table}

\section{Lists}
\subsection{Bullet Points}
\begin{itemize}
  \item First item
  \item Second item
  \item Third item with sub-items:
    \begin{itemize}
      \item Sub-item 1
      \item Sub-item 2
    \end{itemize}
\end{itemize}

\subsection{Numbered List}
\begin{enumerate}
  \item First step
  \item Second step
  \item Third step
\end{enumerate}

\section{Citations and References}
When you use the citation search feature, references will be automatically managed. For example, you might cite a paper like this: \cite{example2024}. All citations will be automatically added to your \texttt{cite.bib} file.

\section{Conclusion}
This template demonstrates basic \LaTeX{} features including mathematical formulas, tables, and lists. The bibliography system is now set up to work seamlessly with Bibby AI's citation management. Feel free to modify it for your needs!

\bibliographystyle{plain}
\bibliography{cite}

\end{document}\documentclass{article}
\usepackage[utf8]{inputenc}
\usepackage{booktabs}
\usepackage{amsmath}
\usepackage{natbib}

\title{My Document}
\author{Your Name}
\date{\today}

\begin{document}

\maketitle

\section{Introduction}
This is a sample document demonstrating various \LaTeX{} features including mathematical equations, tables, and formatting. You can cite papers using the citation search feature, and they will be automatically added to your \texttt{cite.bib} file.

\section{Mathematical Expressions}
Here are some example equations:

\subsection{Quadratic Formula}
The quadratic formula is given by:
\[x = \frac{-b \pm \sqrt{b^2 - 4ac}}{2a}\]

\subsection{Matrix Example}
A sample matrix:
\[
\begin{pmatrix}
a_{11} & a_{12} & a_{13} \\
a_{21} & a_{22} & a_{23} \\
a_{31} & a_{32} & a_{33}
\end{pmatrix}
\]

\section{Tables}
Below is a sample table showing data:

\begin{table}[h]
\centering
\begin{tabular}{@{}lrrr@{}}
\toprule
\textbf{Category} & \textbf{Value 1} & \textbf{Value 2} & \textbf{Total} \\
\midrule
Item A & 10.5 & 20.3 & 30.8 \\
Item B & 15.2 & 25.1 & 40.3 \\
Item C & 22.1 & 12.7 & 34.8 \\
\bottomrule
\end{tabular}
\caption{Sample Data Table}
\label{tab:sample}
\end{table}

\section{Lists}
\subsection{Bullet Points}
\begin{itemize}
  \item First item
  \item Second item
  \item Third item with sub-items:
    \begin{itemize}
      \item Sub-item 1
      \item Sub-item 2
    \end{itemize}
\end{itemize}

\subsection{Numbered List}
\begin{enumerate}
  \item First step
  \item Second step
  \item Third step
\end{enumerate}

\section{Citations and References}
When you use the citation search feature, references will be automatically managed. For example, you might cite a paper like this: \cite{example2024}. All citations will be automatically added to your \texttt{cite.bib} file.

\section{Conclusion}
This template demonstrates basic \LaTeX{} features including mathematical formulas, tables, and lists. The bibliography system is now set up to work seamlessly with Bibby AI's citation management. Feel free to modify it for your needs!

\bibliographystyle{plain}
\bibliography{cite}

\end{document}\documentclass{article}
\usepackage[utf8]{inputenc}
\usepackage{booktabs}
\usepackage{amsmath}
\usepackage{natbib}

\title{My Document}
\author{Your Name}
\date{\today}

\begin{document}

\maketitle

\section{Introduction}
This is a sample document demonstrating various \LaTeX{} features including mathematical equations, tables, and formatting. You can cite papers using the citation search feature, and they will be automatically added to your \texttt{cite.bib} file.

\section{Mathematical Expressions}
Here are some example equations:

\subsection{Quadratic Formula}
The quadratic formula is given by:
\[x = \frac{-b \pm \sqrt{b^2 - 4ac}}{2a}\]

\subsection{Matrix Example}
A sample matrix:
\[
\begin{pmatrix}
a_{11} & a_{12} & a_{13} \\
a_{21} & a_{22} & a_{23} \\
a_{31} & a_{32} & a_{33}
\end{pmatrix}
\]

\section{Tables}
Below is a sample table showing data:

\begin{table}[h]
\centering
\begin{tabular}{@{}lrrr@{}}
\toprule
\textbf{Category} & \textbf{Value 1} & \textbf{Value 2} & \textbf{Total} \\
\midrule
Item A & 10.5 & 20.3 & 30.8 \\
Item B & 15.2 & 25.1 & 40.3 \\
Item C & 22.1 & 12.7 & 34.8 \\
\bottomrule
\end{tabular}
\caption{Sample Data Table}
\label{tab:sample}
\end{table}

\section{Lists}
\subsection{Bullet Points}
\begin{itemize}
  \item First item
  \item Second item
  \item Third item with sub-items:
    \begin{itemize}
      \item Sub-item 1
      \item Sub-item 2
    \end{itemize}
\end{itemize}

\subsection{Numbered List}
\begin{enumerate}
  \item First step
  \item Second step
  \item Third step
\end{enumerate}

\section{Citations and References}
When you use the citation search feature, references will be automatically managed. For example, you might cite a paper like this: \cite{example2024}. All citations will be automatically added to your \texttt{cite.bib} file.

\section{Conclusion}
This template demonstrates basic \LaTeX{} features including mathematical formulas, tables, and lists. The bibliography system is now set up to work seamlessly with Bibby AI's citation management. Feel free to modify it for your needs!

\bibliographystyle{plain}
\bibliography{cite}

\end{document}\documentclass{article}
\usepackage[utf8]{inputenc}
\usepackage{booktabs}
\usepackage{amsmath}
\usepackage{natbib}

\title{My Document}
\author{Your Name}
\date{\today}

\begin{document}

\maketitle

\section{Introduction}
This is a sample document demonstrating various \LaTeX{} features including mathematical equations, tables, and formatting. You can cite papers using the citation search feature, and they will be automatically added to your \texttt{cite.bib} file.

\section{Mathematical Expressions}
Here are some example equations:

\subsection{Quadratic Formula}
The quadratic formula is given by:
\[x = \frac{-b \pm \sqrt{b^2 - 4ac}}{2a}\]

\subsection{Matrix Example}
A sample matrix:
\[
\begin{pmatrix}
a_{11} & a_{12} & a_{13} \\
a_{21} & a_{22} & a_{23} \\
a_{31} & a_{32} & a_{33}
\end{pmatrix}
\]

\section{Tables}
Below is a sample table showing data:

\begin{table}[h]
\centering
\begin{tabular}{@{}lrrr@{}}
\toprule
\textbf{Category} & \textbf{Value 1} & \textbf{Value 2} & \textbf{Total} \\
\midrule
Item A & 10.5 & 20.3 & 30.8 \\
Item B & 15.2 & 25.1 & 40.3 \\
Item C & 22.1 & 12.7 & 34.8 \\
\bottomrule
\end{tabular}
\caption{Sample Data Table}
\label{tab:sample}
\end{table}

\section{Lists}
\subsection{Bullet Points}
\begin{itemize}
  \item First item
  \item Second item
  \item Third item with sub-items:
    \begin{itemize}
      \item Sub-item 1
      \item Sub-item 2
    \end{itemize}
\end{itemize}

\subsection{Numbered List}
\begin{enumerate}
  \item First step
  \item Second step
  \item Third step
\end{enumerate}

\section{Citations and References}
When you use the citation search feature, references will be automatically managed. For example, you might cite a paper like this: \cite{example2024}. All citations will be automatically added to your \texttt{cite.bib} file.

\section{Conclusion}
This template demonstrates basic \LaTeX{} features including mathematical formulas, tables, and lists. The bibliography system is now set up to work seamlessly with Bibby AI's citation management. Feel free to modify it for your needs!

\bibliographystyle{plain}
\bibliography{cite}

\end{document}\documentclass{article}
\usepackage[utf8]{inputenc}
\usepackage{booktabs}
\usepackage{amsmath}
\usepackage{natbib}

\title{My Document}
\author{Your Name}
\date{\today}

\begin{document}

\maketitle

\section{Introduction}
This is a sample document demonstrating various \LaTeX{} features including mathematical equations, tables, and formatting. You can cite papers using the citation search feature, and they will be automatically added to your \texttt{cite.bib} file.

\section{Mathematical Expressions}
Here are some example equations:

\subsection{Quadratic Formula}
The quadratic formula is given by:
\[x = \frac{-b \pm \sqrt{b^2 - 4ac}}{2a}\]

\subsection{Matrix Example}
A sample matrix:
\[
\begin{pmatrix}
a_{11} & a_{12} & a_{13} \\
a_{21} & a_{22} & a_{23} \\
a_{31} & a_{32} & a_{33}
\end{pmatrix}
\]

\section{Tables}
Below is a sample table showing data:

\begin{table}[h]
\centering
\begin{tabular}{@{}lrrr@{}}
\toprule
\textbf{Category} & \textbf{Value 1} & \textbf{Value 2} & \textbf{Total} \\
\midrule
Item A & 10.5 & 20.3 & 30.8 \\
Item B & 15.2 & 25.1 & 40.3 \\
Item C & 22.1 & 12.7 & 34.8 \\
\bottomrule
\end{tabular}
\caption{Sample Data Table}
\label{tab:sample}
\end{table}

\section{Lists}
\subsection{Bullet Points}
\begin{itemize}
  \item First item
  \item Second item
  \item Third item with sub-items:
    \begin{itemize}
      \item Sub-item 1
      \item Sub-item 2
    \end{itemize}
\end{itemize}

\subsection{Numbered List}
\begin{enumerate}
  \item First step
  \item Second step
  \item Third step
\end{enumerate}

\section{Citations and References}
When you use the citation search feature, references will be automatically managed. For example, you might cite a paper like this: \cite{example2024}. All citations will be automatically added to your \texttt{cite.bib} file.

\section{Conclusion}
This template demonstrates basic \LaTeX{} features including mathematical formulas, tables, and lists. The bibliography system is now set up to work seamlessly with Bibby AI's citation management. Feel free to modify it for your needs!

\bibliographystyle{plain}
\bibliography{cite}

\end{document}\documentclass{article}
\usepackage[utf8]{inputenc}
\usepackage{booktabs}
\usepackage{amsmath}
\usepackage{natbib}

\title{My Document}
\author{Your Name}
\date{\today}

\begin{document}

\maketitle

\section{Introduction}
This is a sample document demonstrating various \LaTeX{} features including mathematical equations, tables, and formatting. You can cite papers using the citation search feature, and they will be automatically added to your \texttt{cite.bib} file.

\section{Mathematical Expressions}
Here are some example equations:

\subsection{Quadratic Formula}
The quadratic formula is given by:
\[x = \frac{-b \pm \sqrt{b^2 - 4ac}}{2a}\]

\subsection{Matrix Example}
A sample matrix:
\[
\begin{pmatrix}
a_{11} & a_{12} & a_{13} \\
a_{21} & a_{22} & a_{23} \\
a_{31} & a_{32} & a_{33}
\end{pmatrix}
\]

\section{Tables}
Below is a sample table showing data:

\begin{table}[h]
\centering
\begin{tabular}{@{}lrrr@{}}
\toprule
\textbf{Category} & \textbf{Value 1} & \textbf{Value 2} & \textbf{Total} \\
\midrule
Item A & 10.5 & 20.3 & 30.8 \\
Item B & 15.2 & 25.1 & 40.3 \\
Item C & 22.1 & 12.7 & 34.8 \\
\bottomrule
\end{tabular}
\caption{Sample Data Table}
\label{tab:sample}
\end{table}

\section{Lists}
\subsection{Bullet Points}
\begin{itemize}
  \item First item
  \item Second item
  \item Third item with sub-items:
    \begin{itemize}
      \item Sub-item 1
      \item Sub-item 2
    \end{itemize}
\end{itemize}

\subsection{Numbered List}
\begin{enumerate}
  \item First step
  \item Second step
  \item Third step
\end{enumerate}

\section{Citations and References}
When you use the citation search feature, references will be automatically managed. For example, you might cite a paper like this: \cite{example2024}. All citations will be automatically added to your \texttt{cite.bib} file.

\section{Conclusion}
This template demonstrates basic \LaTeX{} features including mathematical formulas, tables, and lists. The bibliography system is now set up to work seamlessly with Bibby AI's citation management. Feel free to modify it for your needs!

\bibliographystyle{plain}
\bibliography{cite}

\end{document}\documentclass{article}
\usepackage[utf8]{inputenc}
\usepackage{booktabs}
\usepackage{amsmath}
\usepackage{natbib}

\title{My Document}
\author{Your Name}
\date{\today}

\begin{document}

\maketitle

\section{Introduction}
This is a sample document demonstrating various \LaTeX{} features including mathematical equations, tables, and formatting. You can cite papers using the citation search feature, and they will be automatically added to your \texttt{cite.bib} file.

\section{Mathematical Expressions}
Here are some example equations:

\subsection{Quadratic Formula}
The quadratic formula is given by:
\[x = \frac{-b \pm \sqrt{b^2 - 4ac}}{2a}\]

\subsection{Matrix Example}
A sample matrix:
\[
\begin{pmatrix}
a_{11} & a_{12} & a_{13} \\
a_{21} & a_{22} & a_{23} \\
a_{31} & a_{32} & a_{33}
\end{pmatrix}
\]

\section{Tables}
Below is a sample table showing data:

\begin{table}[h]
\centering
\begin{tabular}{@{}lrrr@{}}
\toprule
\textbf{Category} & \textbf{Value 1} & \textbf{Value 2} & \textbf{Total} \\
\midrule
Item A & 10.5 & 20.3 & 30.8 \\
Item B & 15.2 & 25.1 & 40.3 \\
Item C & 22.1 & 12.7 & 34.8 \\
\bottomrule
\end{tabular}
\caption{Sample Data Table}
\label{tab:sample}
\end{table}

\section{Lists}
\subsection{Bullet Points}
\begin{itemize}
  \item First item
  \item Second item
  \item Third item with sub-items:
    \begin{itemize}
      \item Sub-item 1
      \item Sub-item 2
    \end{itemize}
\end{itemize}

\subsection{Numbered List}
\begin{enumerate}
  \item First step
  \item Second step
  \item Third step
\end{enumerate}

\section{Citations and References}
When you use the citation search feature, references will be automatically managed. For example, you might cite a paper like this: \cite{example2024}. All citations will be automatically added to your \texttt{cite.bib} file.

\section{Conclusion}
This template demonstrates basic \LaTeX{} features including mathematical formulas, tables, and lists. The bibliography system is now set up to work seamlessly with Bibby AI's citation management. Feel free to modify it for your needs!

\bibliographystyle{plain}
\bibliography{cite}

\end{document}\documentclass{article}
\usepackage[utf8]{inputenc}
\usepackage{booktabs}
\usepackage{amsmath}
\usepackage{natbib}

\title{My Document}
\author{Your Name}
\date{\today}

\begin{document}

\maketitle

\section{Introduction}
This is a sample document demonstrating various \LaTeX{} features including mathematical equations, tables, and formatting. You can cite papers using the citation search feature, and they will be automatically added to your \texttt{cite.bib} file.

\section{Mathematical Expressions}
Here are some example equations:

\subsection{Quadratic Formula}
The quadratic formula is given by:
\[x = \frac{-b \pm \sqrt{b^2 - 4ac}}{2a}\]

\subsection{Matrix Example}
A sample matrix:
\[
\begin{pmatrix}
a_{11} & a_{12} & a_{13} \\
a_{21} & a_{22} & a_{23} \\
a_{31} & a_{32} & a_{33}
\end{pmatrix}
\]

\section{Tables}
Below is a sample table showing data:

\begin{table}[h]
\centering
\begin{tabular}{@{}lrrr@{}}
\toprule
\textbf{Category} & \textbf{Value 1} & \textbf{Value 2} & \textbf{Total} \\
\midrule
Item A & 10.5 & 20.3 & 30.8 \\
Item B & 15.2 & 25.1 & 40.3 \\
Item C & 22.1 & 12.7 & 34.8 \\
\bottomrule
\end{tabular}
\caption{Sample Data Table}
\label{tab:sample}
\end{table}

\section{Lists}
\subsection{Bullet Points}
\begin{itemize}
  \item First item
  \item Second item
  \item Third item with sub-items:
    \begin{itemize}
      \item Sub-item 1
      \item Sub-item 2
    \end{itemize}
\end{itemize}

\subsection{Numbered List}
\begin{enumerate}
  \item First step
  \item Second step
  \item Third step
\end{enumerate}

\section{Citations and References}
When you use the citation search feature, references will be automatically managed. For example, you might cite a paper like this: \cite{example2024}. All citations will be automatically added to your \texttt{cite.bib} file.

\section{Conclusion}
This template demonstrates basic \LaTeX{} features including mathematical formulas, tables, and lists. The bibliography system is now set up to work seamlessly with Bibby AI's citation management. Feel free to modify it for your needs!

\bibliographystyle{plain}
\bibliography{cite}

\end{document}\documentclass{article}
\usepackage[utf8]{inputenc}
\usepackage{booktabs}
\usepackage{amsmath}
\usepackage{natbib}

\title{My Document}
\author{Your Name}
\date{\today}

\begin{document}

\maketitle

\section{Introduction}
This is a sample document demonstrating various \LaTeX{} features including mathematical equations, tables, and formatting. You can cite papers using the citation search feature, and they will be automatically added to your \texttt{cite.bib} file.

\section{Mathematical Expressions}
Here are some example equations:

\subsection{Quadratic Formula}
The quadratic formula is given by:
\[x = \frac{-b \pm \sqrt{b^2 - 4ac}}{2a}\]

\subsection{Matrix Example}
A sample matrix:
\[
\begin{pmatrix}
a_{11} & a_{12} & a_{13} \\
a_{21} & a_{22} & a_{23} \\
a_{31} & a_{32} & a_{33}
\end{pmatrix}
\]

\section{Tables}
Below is a sample table showing data:

\begin{table}[h]
\centering
\begin{tabular}{@{}lrrr@{}}
\toprule
\textbf{Category} & \textbf{Value 1} & \textbf{Value 2} & \textbf{Total} \\
\midrule
Item A & 10.5 & 20.3 & 30.8 \\
Item B & 15.2 & 25.1 & 40.3 \\
Item C & 22.1 & 12.7 & 34.8 \\
\bottomrule
\end{tabular}
\caption{Sample Data Table}
\label{tab:sample}
\end{table}

\section{Lists}
\subsection{Bullet Points}
\begin{itemize}
  \item First item
  \item Second item
  \item Third item with sub-items:
    \begin{itemize}
      \item Sub-item 1
      \item Sub-item 2
    \end{itemize}
\end{itemize}

\subsection{Numbered List}
\begin{enumerate}
  \item First step
  \item Second step
  \item Third step
\end{enumerate}

\section{Citations and References}
When you use the citation search feature, references will be automatically managed. For example, you might cite a paper like this: \cite{example2024}. All citations will be automatically added to your \texttt{cite.bib} file.

\section{Conclusion}
This template demonstrates basic \LaTeX{} features including mathematical formulas, tables, and lists. The bibliography system is now set up to work seamlessly with Bibby AI's citation management. Feel free to modify it for your needs!

\bibliographystyle{plain}
\bibliography{cite}

\end{document}\documentclass{article}
\usepackage[utf8]{inputenc}
\usepackage{booktabs}
\usepackage{amsmath}
\usepackage{natbib}

\title{My Document}
\author{Your Name}
\date{\today}

\begin{document}

\maketitle

\section{Introduction}
This is a sample document demonstrating various \LaTeX{} features including mathematical equations, tables, and formatting. You can cite papers using the citation search feature, and they will be automatically added to your \texttt{cite.bib} file.

\section{Mathematical Expressions}
Here are some example equations:

\subsection{Quadratic Formula}
The quadratic formula is given by:
\[x = \frac{-b \pm \sqrt{b^2 - 4ac}}{2a}\]

\subsection{Matrix Example}
A sample matrix:
\[
\begin{pmatrix}
a_{11} & a_{12} & a_{13} \\
a_{21} & a_{22} & a_{23} \\
a_{31} & a_{32} & a_{33}
\end{pmatrix}
\]

\section{Tables}
Below is a sample table showing data:

\begin{table}[h]
\centering
\begin{tabular}{@{}lrrr@{}}
\toprule
\textbf{Category} & \textbf{Value 1} & \textbf{Value 2} & \textbf{Total} \\
\midrule
Item A & 10.5 & 20.3 & 30.8 \\
Item B & 15.2 & 25.1 & 40.3 \\
Item C & 22.1 & 12.7 & 34.8 \\
\bottomrule
\end{tabular}
\caption{Sample Data Table}
\label{tab:sample}
\end{table}

\section{Lists}
\subsection{Bullet Points}
\begin{itemize}
  \item First item
  \item Second item
  \item Third item with sub-items:
    \begin{itemize}
      \item Sub-item 1
      \item Sub-item 2
    \end{itemize}
\end{itemize}

\subsection{Numbered List}
\begin{enumerate}
  \item First step
  \item Second step
  \item Third step
\end{enumerate}

\section{Citations and References}
When you use the citation search feature, references will be automatically managed. For example, you might cite a paper like this: \cite{example2024}. All citations will be automatically added to your \texttt{cite.bib} file.

\section{Conclusion}
This template demonstrates basic \LaTeX{} features including mathematical formulas, tables, and lists. The bibliography system is now set up to work seamlessly with Bibby AI's citation management. Feel free to modify it for your needs!

\bibliographystyle{plain}
\bibliography{cite}

\end{document}\documentclass{article}
\usepackage[utf8]{inputenc}
\usepackage{booktabs}
\usepackage{amsmath}
\usepackage{natbib}

\title{My Document}
\author{Your Name}
\date{\today}

\begin{document}

\maketitle

\section{Introduction}
This is a sample document demonstrating various \LaTeX{} features including mathematical equations, tables, and formatting. You can cite papers using the citation search feature, and they will be automatically added to your \texttt{cite.bib} file.

\section{Mathematical Expressions}
Here are some example equations:

\subsection{Quadratic Formula}
The quadratic formula is given by:
\[x = \frac{-b \pm \sqrt{b^2 - 4ac}}{2a}\]

\subsection{Matrix Example}
A sample matrix:
\[
\begin{pmatrix}
a_{11} & a_{12} & a_{13} \\
a_{21} & a_{22} & a_{23} \\
a_{31} & a_{32} & a_{33}
\end{pmatrix}
\]

\section{Tables}
Below is a sample table showing data:

\begin{table}[h]
\centering
\begin{tabular}{@{}lrrr@{}}
\toprule
\textbf{Category} & \textbf{Value 1} & \textbf{Value 2} & \textbf{Total} \\
\midrule
Item A & 10.5 & 20.3 & 30.8 \\
Item B & 15.2 & 25.1 & 40.3 \\
Item C & 22.1 & 12.7 & 34.8 \\
\bottomrule
\end{tabular}
\caption{Sample Data Table}
\label{tab:sample}
\end{table}

\section{Lists}
\subsection{Bullet Points}
\begin{itemize}
  \item First item
  \item Second item
  \item Third item with sub-items:
    \begin{itemize}
      \item Sub-item 1
      \item Sub-item 2
    \end{itemize}
\end{itemize}

\subsection{Numbered List}
\begin{enumerate}
  \item First step
  \item Second step
  \item Third step
\end{enumerate}

\section{Citations and References}
When you use the citation search feature, references will be automatically managed. For example, you might cite a paper like this: \cite{example2024}. All citations will be automatically added to your \texttt{cite.bib} file.

\section{Conclusion}
This template demonstrates basic \LaTeX{} features including mathematical formulas, tables, and lists. The bibliography system is now set up to work seamlessly with Bibby AI's citation management. Feel free to modify it for your needs!

\bibliographystyle{plain}
\bibliography{cite}

\end{document}\documentclass{article}
\usepackage[utf8]{inputenc}
\usepackage{booktabs}
\usepackage{amsmath}
\usepackage{natbib}

\title{My Document}
\author{Your Name}
\date{\today}

\begin{document}

\maketitle

\section{Introduction}
This is a sample document demonstrating various \LaTeX{} features including mathematical equations, tables, and formatting. You can cite papers using the citation search feature, and they will be automatically added to your \texttt{cite.bib} file.

\section{Mathematical Expressions}
Here are some example equations:

\subsection{Quadratic Formula}
The quadratic formula is given by:
\[x = \frac{-b \pm \sqrt{b^2 - 4ac}}{2a}\]

\subsection{Matrix Example}
A sample matrix:
\[
\begin{pmatrix}
a_{11} & a_{12} & a_{13} \\
a_{21} & a_{22} & a_{23} \\
a_{31} & a_{32} & a_{33}
\end{pmatrix}
\]

\section{Tables}
Below is a sample table showing data:

\begin{table}[h]
\centering
\begin{tabular}{@{}lrrr@{}}
\toprule
\textbf{Category} & \textbf{Value 1} & \textbf{Value 2} & \textbf{Total} \\
\midrule
Item A & 10.5 & 20.3 & 30.8 \\
Item B & 15.2 & 25.1 & 40.3 \\
Item C & 22.1 & 12.7 & 34.8 \\
\bottomrule
\end{tabular}
\caption{Sample Data Table}
\label{tab:sample}
\end{table}

\section{Lists}
\subsection{Bullet Points}
\begin{itemize}
  \item First item
  \item Second item
  \item Third item with sub-items:
    \begin{itemize}
      \item Sub-item 1
      \item Sub-item 2
    \end{itemize}
\end{itemize}

\subsection{Numbered List}
\begin{enumerate}
  \item First step
  \item Second step
  \item Third step
\end{enumerate}

\section{Citations and References}
When you use the citation search feature, references will be automatically managed. For example, you might cite a paper like this: \cite{example2024}. All citations will be automatically added to your \texttt{cite.bib} file.

\section{Conclusion}
This template demonstrates basic \LaTeX{} features including mathematical formulas, tables, and lists. The bibliography system is now set up to work seamlessly with Bibby AI's citation management. Feel free to modify it for your needs!

\bibliographystyle{plain}
\bibliography{cite}

\end{document}\documentclass{article}
\usepackage[utf8]{inputenc}
\usepackage{booktabs}
\usepackage{amsmath}
\usepackage{natbib}

\title{My Document}
\author{Your Name}
\date{\today}

\begin{document}

\maketitle

\section{Introduction}
This is a sample document demonstrating various \LaTeX{} features including mathematical equations, tables, and formatting. You can cite papers using the citation search feature, and they will be automatically added to your \texttt{cite.bib} file.

\section{Mathematical Expressions}
Here are some example equations:

\subsection{Quadratic Formula}
The quadratic formula is given by:
\[x = \frac{-b \pm \sqrt{b^2 - 4ac}}{2a}\]

\subsection{Matrix Example}
A sample matrix:
\[
\begin{pmatrix}
a_{11} & a_{12} & a_{13} \\
a_{21} & a_{22} & a_{23} \\
a_{31} & a_{32} & a_{33}
\end{pmatrix}
\]

\section{Tables}
Below is a sample table showing data:

\begin{table}[h]
\centering
\begin{tabular}{@{}lrrr@{}}
\toprule
\textbf{Category} & \textbf{Value 1} & \textbf{Value 2} & \textbf{Total} \\
\midrule
Item A & 10.5 & 20.3 & 30.8 \\
Item B & 15.2 & 25.1 & 40.3 \\
Item C & 22.1 & 12.7 & 34.8 \\
\bottomrule
\end{tabular}
\caption{Sample Data Table}
\label{tab:sample}
\end{table}

\section{Lists}
\subsection{Bullet Points}
\begin{itemize}
  \item First item
  \item Second item
  \item Third item with sub-items:
    \begin{itemize}
      \item Sub-item 1
      \item Sub-item 2
    \end{itemize}
\end{itemize}

\subsection{Numbered List}
\begin{enumerate}
  \item First step
  \item Second step
  \item Third step
\end{enumerate}

\section{Citations and References}
When you use the citation search feature, references will be automatically managed. For example, you might cite a paper like this: \cite{example2024}. All citations will be automatically added to your \texttt{cite.bib} file.

\section{Conclusion}
This template demonstrates basic \LaTeX{} features including mathematical formulas, tables, and lists. The bibliography system is now set up to work seamlessly with Bibby AI's citation management. Feel free to modify it for your needs!

\bibliographystyle{plain}
\bibliography{cite}

\end{document}\documentclass{article}
\usepackage[utf8]{inputenc}
\usepackage{booktabs}
\usepackage{amsmath}
\usepackage{natbib}

\title{My Document}
\author{Your Name}
\date{\today}

\begin{document}

\maketitle

\section{Introduction}
This is a sample document demonstrating various \LaTeX{} features including mathematical equations, tables, and formatting. You can cite papers using the citation search feature, and they will be automatically added to your \texttt{cite.bib} file.

\section{Mathematical Expressions}
Here are some example equations:

\subsection{Quadratic Formula}
The quadratic formula is given by:
\[x = \frac{-b \pm \sqrt{b^2 - 4ac}}{2a}\]

\subsection{Matrix Example}
A sample matrix:
\[
\begin{pmatrix}
a_{11} & a_{12} & a_{13} \\
a_{21} & a_{22} & a_{23} \\
a_{31} & a_{32} & a_{33}
\end{pmatrix}
\]

\section{Tables}
Below is a sample table showing data:

\begin{table}[h]
\centering
\begin{tabular}{@{}lrrr@{}}
\toprule
\textbf{Category} & \textbf{Value 1} & \textbf{Value 2} & \textbf{Total} \\
\midrule
Item A & 10.5 & 20.3 & 30.8 \\
Item B & 15.2 & 25.1 & 40.3 \\
Item C & 22.1 & 12.7 & 34.8 \\
\bottomrule
\end{tabular}
\caption{Sample Data Table}
\label{tab:sample}
\end{table}

\section{Lists}
\subsection{Bullet Points}
\begin{itemize}
  \item First item
  \item Second item
  \item Third item with sub-items:
    \begin{itemize}
      \item Sub-item 1
      \item Sub-item 2
    \end{itemize}
\end{itemize}

\subsection{Numbered List}
\begin{enumerate}
  \item First step
  \item Second step
  \item Third step
\end{enumerate}

\section{Citations and References}
When you use the citation search feature, references will be automatically managed. For example, you might cite a paper like this: \cite{example2024}. All citations will be automatically added to your \texttt{cite.bib} file.

\section{Conclusion}
This template demonstrates basic \LaTeX{} features including mathematical formulas, tables, and lists. The bibliography system is now set up to work seamlessly with Bibby AI's citation management. Feel free to modify it for your needs!

\bibliographystyle{plain}
\bibliography{cite}

\end{document}\documentclass{article}
\usepackage[utf8]{inputenc}
\usepackage{booktabs}
\usepackage{amsmath}
\usepackage{natbib}

\title{My Document}
\author{Your Name}
\date{\today}

\begin{document}

\maketitle

\section{Introduction}
This is a sample document demonstrating various \LaTeX{} features including mathematical equations, tables, and formatting. You can cite papers using the citation search feature, and they will be automatically added to your \texttt{cite.bib} file.

\section{Mathematical Expressions}
Here are some example equations:

\subsection{Quadratic Formula}
The quadratic formula is given by:
\[x = \frac{-b \pm \sqrt{b^2 - 4ac}}{2a}\]

\subsection{Matrix Example}
A sample matrix:
\[
\begin{pmatrix}
a_{11} & a_{12} & a_{13} \\
a_{21} & a_{22} & a_{23} \\
a_{31} & a_{32} & a_{33}
\end{pmatrix}
\]

\section{Tables}
Below is a sample table showing data:

\begin{table}[h]
\centering
\begin{tabular}{@{}lrrr@{}}
\toprule
\textbf{Category} & \textbf{Value 1} & \textbf{Value 2} & \textbf{Total} \\
\midrule
Item A & 10.5 & 20.3 & 30.8 \\
Item B & 15.2 & 25.1 & 40.3 \\
Item C & 22.1 & 12.7 & 34.8 \\
\bottomrule
\end{tabular}
\caption{Sample Data Table}
\label{tab:sample}
\end{table}

\section{Lists}
\subsection{Bullet Points}
\begin{itemize}
  \item First item
  \item Second item
  \item Third item with sub-items:
    \begin{itemize}
      \item Sub-item 1
      \item Sub-item 2
    \end{itemize}
\end{itemize}

\subsection{Numbered List}
\begin{enumerate}
  \item First step
  \item Second step
  \item Third step
\end{enumerate}

\section{Citations and References}
When you use the citation search feature, references will be automatically managed. For example, you might cite a paper like this: \cite{example2024}. All citations will be automatically added to your \texttt{cite.bib} file.

\section{Conclusion}
This template demonstrates basic \LaTeX{} features including mathematical formulas, tables, and lists. The bibliography system is now set up to work seamlessly with Bibby AI's citation management. Feel free to modify it for your needs!

\bibliographystyle{plain}
\bibliography{cite}

\end{document}\documentclass{article}
\usepackage[utf8]{inputenc}
\usepackage{booktabs}
\usepackage{amsmath}
\usepackage{natbib}

\title{My Document}
\author{Your Name}
\date{\today}

\begin{document}

\maketitle

\section{Introduction}
This is a sample document demonstrating various \LaTeX{} features including mathematical equations, tables, and formatting. You can cite papers using the citation search feature, and they will be automatically added to your \texttt{cite.bib} file.

\section{Mathematical Expressions}
Here are some example equations:

\subsection{Quadratic Formula}
The quadratic formula is given by:
\[x = \frac{-b \pm \sqrt{b^2 - 4ac}}{2a}\]

\subsection{Matrix Example}
A sample matrix:
\[
\begin{pmatrix}
a_{11} & a_{12} & a_{13} \\
a_{21} & a_{22} & a_{23} \\
a_{31} & a_{32} & a_{33}
\end{pmatrix}
\]

\section{Tables}
Below is a sample table showing data:

\begin{table}[h]
\centering
\begin{tabular}{@{}lrrr@{}}
\toprule
\textbf{Category} & \textbf{Value 1} & \textbf{Value 2} & \textbf{Total} \\
\midrule
Item A & 10.5 & 20.3 & 30.8 \\
Item B & 15.2 & 25.1 & 40.3 \\
Item C & 22.1 & 12.7 & 34.8 \\
\bottomrule
\end{tabular}
\caption{Sample Data Table}
\label{tab:sample}
\end{table}

\section{Lists}
\subsection{Bullet Points}
\begin{itemize}
  \item First item
  \item Second item
  \item Third item with sub-items:
    \begin{itemize}
      \item Sub-item 1
      \item Sub-item 2
    \end{itemize}
\end{itemize}

\subsection{Numbered List}
\begin{enumerate}
  \item First step
  \item Second step
  \item Third step
\end{enumerate}

\section{Citations and References}
When you use the citation search feature, references will be automatically managed. For example, you might cite a paper like this: \cite{example2024}. All citations will be automatically added to your \texttt{cite.bib} file.

\section{Conclusion}
This template demonstrates basic \LaTeX{} features including mathematical formulas, tables, and lists. The bibliography system is now set up to work seamlessly with Bibby AI's citation management. Feel free to modify it for your needs!

\bibliographystyle{plain}
\bibliography{cite}

\end{document}\documentclass{article}
\usepackage[utf8]{inputenc}
\usepackage{booktabs}
\usepackage{amsmath}
\usepackage{natbib}

\title{My Document}
\author{Your Name}
\date{\today}

\begin{document}

\maketitle

\section{Introduction}
This is a sample document demonstrating various \LaTeX{} features including mathematical equations, tables, and formatting. You can cite papers using the citation search feature, and they will be automatically added to your \texttt{cite.bib} file.

\section{Mathematical Expressions}
Here are some example equations:

\subsection{Quadratic Formula}
The quadratic formula is given by:
\[x = \frac{-b \pm \sqrt{b^2 - 4ac}}{2a}\]

\subsection{Matrix Example}
A sample matrix:
\[
\begin{pmatrix}
a_{11} & a_{12} & a_{13} \\
a_{21} & a_{22} & a_{23} \\
a_{31} & a_{32} & a_{33}
\end{pmatrix}
\]

\section{Tables}
Below is a sample table showing data:

\begin{table}[h]
\centering
\begin{tabular}{@{}lrrr@{}}
\toprule
\textbf{Category} & \textbf{Value 1} & \textbf{Value 2} & \textbf{Total} \\
\midrule
Item A & 10.5 & 20.3 & 30.8 \\
Item B & 15.2 & 25.1 & 40.3 \\
Item C & 22.1 & 12.7 & 34.8 \\
\bottomrule
\end{tabular}
\caption{Sample Data Table}
\label{tab:sample}
\end{table}

\section{Lists}
\subsection{Bullet Points}
\begin{itemize}
  \item First item
  \item Second item
  \item Third item with sub-items:
    \begin{itemize}
      \item Sub-item 1
      \item Sub-item 2
    \end{itemize}
\end{itemize}

\subsection{Numbered List}
\begin{enumerate}
  \item First step
  \item Second step
  \item Third step
\end{enumerate}

\section{Citations and References}
When you use the citation search feature, references will be automatically managed. For example, you might cite a paper like this: \cite{example2024}. All citations will be automatically added to your \texttt{cite.bib} file.

\section{Conclusion}
This template demonstrates basic \LaTeX{} features including mathematical formulas, tables, and lists. The bibliography system is now set up to work seamlessly with Bibby AI's citation management. Feel free to modify it for your needs!

\bibliographystyle{plain}
\bibliography{cite}

\end{document}\documentclass{article}
\usepackage[utf8]{inputenc}
\usepackage{booktabs}
\usepackage{amsmath}
\usepackage{natbib}

\title{My Document}
\author{Your Name}
\date{\today}

\begin{document}

\maketitle

\section{Introduction}
This is a sample document demonstrating various \LaTeX{} features including mathematical equations, tables, and formatting. You can cite papers using the citation search feature, and they will be automatically added to your \texttt{cite.bib} file.

\section{Mathematical Expressions}
Here are some example equations:

\subsection{Quadratic Formula}
The quadratic formula is given by:
\[x = \frac{-b \pm \sqrt{b^2 - 4ac}}{2a}\]

\subsection{Matrix Example}
A sample matrix:
\[
\begin{pmatrix}
a_{11} & a_{12} & a_{13} \\
a_{21} & a_{22} & a_{23} \\
a_{31} & a_{32} & a_{33}
\end{pmatrix}
\]

\section{Tables}
Below is a sample table showing data:

\begin{table}[h]
\centering
\begin{tabular}{@{}lrrr@{}}
\toprule
\textbf{Category} & \textbf{Value 1} & \textbf{Value 2} & \textbf{Total} \\
\midrule
Item A & 10.5 & 20.3 & 30.8 \\
Item B & 15.2 & 25.1 & 40.3 \\
Item C & 22.1 & 12.7 & 34.8 \\
\bottomrule
\end{tabular}
\caption{Sample Data Table}
\label{tab:sample}
\end{table}

\section{Lists}
\subsection{Bullet Points}
\begin{itemize}
  \item First item
  \item Second item
  \item Third item with sub-items:
    \begin{itemize}
      \item Sub-item 1
      \item Sub-item 2
    \end{itemize}
\end{itemize}

\subsection{Numbered List}
\begin{enumerate}
  \item First step
  \item Second step
  \item Third step
\end{enumerate}

\section{Citations and References}
When you use the citation search feature, references will be automatically managed. For example, you might cite a paper like this: \cite{example2024}. All citations will be automatically added to your \texttt{cite.bib} file.

\section{Conclusion}
This template demonstrates basic \LaTeX{} features including mathematical formulas, tables, and lists. The bibliography system is now set up to work seamlessly with Bibby AI's citation management. Feel free to modify it for your needs!

\bibliographystyle{plain}
\bibliography{cite}

\end{document}\documentclass{article}
\usepackage[utf8]{inputenc}
\usepackage{booktabs}
\usepackage{amsmath}
\usepackage{natbib}

\title{My Document}
\author{Your Name}
\date{\today}

\begin{document}

\maketitle

\section{Introduction}
This is a sample document demonstrating various \LaTeX{} features including mathematical equations, tables, and formatting. You can cite papers using the citation search feature, and they will be automatically added to your \texttt{cite.bib} file.

\section{Mathematical Expressions}
Here are some example equations:

\subsection{Quadratic Formula}
The quadratic formula is given by:
\[x = \frac{-b \pm \sqrt{b^2 - 4ac}}{2a}\]

\subsection{Matrix Example}
A sample matrix:
\[
\begin{pmatrix}
a_{11} & a_{12} & a_{13} \\
a_{21} & a_{22} & a_{23} \\
a_{31} & a_{32} & a_{33}
\end{pmatrix}
\]

\section{Tables}
Below is a sample table showing data:

\begin{table}[h]
\centering
\begin{tabular}{@{}lrrr@{}}
\toprule
\textbf{Category} & \textbf{Value 1} & \textbf{Value 2} & \textbf{Total} \\
\midrule
Item A & 10.5 & 20.3 & 30.8 \\
Item B & 15.2 & 25.1 & 40.3 \\
Item C & 22.1 & 12.7 & 34.8 \\
\bottomrule
\end{tabular}
\caption{Sample Data Table}
\label{tab:sample}
\end{table}

\section{Lists}
\subsection{Bullet Points}
\begin{itemize}
  \item First item
  \item Second item
  \item Third item with sub-items:
    \begin{itemize}
      \item Sub-item 1
      \item Sub-item 2
    \end{itemize}
\end{itemize}

\subsection{Numbered List}
\begin{enumerate}
  \item First step
  \item Second step
  \item Third step
\end{enumerate}

\section{Citations and References}
When you use the citation search feature, references will be automatically managed. For example, you might cite a paper like this: \cite{example2024}. All citations will be automatically added to your \texttt{cite.bib} file.

\section{Conclusion}
This template demonstrates basic \LaTeX{} features including mathematical formulas, tables, and lists. The bibliography system is now set up to work seamlessly with Bibby AI's citation management. Feel free to modify it for your needs!

\bibliographystyle{plain}
\bibliography{cite}

\end{document}\documentclass{article}
\usepackage[utf8]{inputenc}
\usepackage{booktabs}
\usepackage{amsmath}
\usepackage{natbib}

\title{My Document}
\author{Your Name}
\date{\today}

\begin{document}

\maketitle

\section{Introduction}
This is a sample document demonstrating various \LaTeX{} features including mathematical equations, tables, and formatting. You can cite papers using the citation search feature, and they will be automatically added to your \texttt{cite.bib} file.

\section{Mathematical Expressions}
Here are some example equations:

\subsection{Quadratic Formula}
The quadratic formula is given by:
\[x = \frac{-b \pm \sqrt{b^2 - 4ac}}{2a}\]

\subsection{Matrix Example}
A sample matrix:
\[
\begin{pmatrix}
a_{11} & a_{12} & a_{13} \\
a_{21} & a_{22} & a_{23} \\
a_{31} & a_{32} & a_{33}
\end{pmatrix}
\]

\section{Tables}
Below is a sample table showing data:

\begin{table}[h]
\centering
\begin{tabular}{@{}lrrr@{}}
\toprule
\textbf{Category} & \textbf{Value 1} & \textbf{Value 2} & \textbf{Total} \\
\midrule
Item A & 10.5 & 20.3 & 30.8 \\
Item B & 15.2 & 25.1 & 40.3 \\
Item C & 22.1 & 12.7 & 34.8 \\
\bottomrule
\end{tabular}
\caption{Sample Data Table}
\label{tab:sample}
\end{table}

\section{Lists}
\subsection{Bullet Points}
\begin{itemize}
  \item First item
  \item Second item
  \item Third item with sub-items:
    \begin{itemize}
      \item Sub-item 1
      \item Sub-item 2
    \end{itemize}
\end{itemize}

\subsection{Numbered List}
\begin{enumerate}
  \item First step
  \item Second step
  \item Third step
\end{enumerate}

\section{Citations and References}
When you use the citation search feature, references will be automatically managed. For example, you might cite a paper like this: \cite{example2024}. All citations will be automatically added to your \texttt{cite.bib} file.

\section{Conclusion}
This template demonstrates basic \LaTeX{} features including mathematical formulas, tables, and lists. The bibliography system is now set up to work seamlessly with Bibby AI's citation management. Feel free to modify it for your needs!

\bibliographystyle{plain}
\bibliography{cite}

\end{document}\documentclass{article}
\usepackage[utf8]{inputenc}
\usepackage{booktabs}
\usepackage{amsmath}
\usepackage{natbib}

\title{My Document}
\author{Your Name}
\date{\today}

\begin{document}

\maketitle

\section{Introduction}
This is a sample document demonstrating various \LaTeX{} features including mathematical equations, tables, and formatting. You can cite papers using the citation search feature, and they will be automatically added to your \texttt{cite.bib} file.

\section{Mathematical Expressions}
Here are some example equations:

\subsection{Quadratic Formula}
The quadratic formula is given by:
\[x = \frac{-b \pm \sqrt{b^2 - 4ac}}{2a}\]

\subsection{Matrix Example}
A sample matrix:
\[
\begin{pmatrix}
a_{11} & a_{12} & a_{13} \\
a_{21} & a_{22} & a_{23} \\
a_{31} & a_{32} & a_{33}
\end{pmatrix}
\]

\section{Tables}
Below is a sample table showing data:

\begin{table}[h]
\centering
\begin{tabular}{@{}lrrr@{}}
\toprule
\textbf{Category} & \textbf{Value 1} & \textbf{Value 2} & \textbf{Total} \\
\midrule
Item A & 10.5 & 20.3 & 30.8 \\
Item B & 15.2 & 25.1 & 40.3 \\
Item C & 22.1 & 12.7 & 34.8 \\
\bottomrule
\end{tabular}
\caption{Sample Data Table}
\label{tab:sample}
\end{table}

\section{Lists}
\subsection{Bullet Points}
\begin{itemize}
  \item First item
  \item Second item
  \item Third item with sub-items:
    \begin{itemize}
      \item Sub-item 1
      \item Sub-item 2
    \end{itemize}
\end{itemize}

\subsection{Numbered List}
\begin{enumerate}
  \item First step
  \item Second step
  \item Third step
\end{enumerate}

\section{Citations and References}
When you use the citation search feature, references will be automatically managed. For example, you might cite a paper like this: \cite{example2024}. All citations will be automatically added to your \texttt{cite.bib} file.

\section{Conclusion}
This template demonstrates basic \LaTeX{} features including mathematical formulas, tables, and lists. The bibliography system is now set up to work seamlessly with Bibby AI's citation management. Feel free to modify it for your needs!

\bibliographystyle{plain}
\bibliography{cite}

\end{document}\documentclass{article}
\usepackage[utf8]{inputenc}
\usepackage{booktabs}
\usepackage{amsmath}
\usepackage{natbib}

\title{My Document}
\author{Your Name}
\date{\today}

\begin{document}

\maketitle

\section{Introduction}
This is a sample document demonstrating various \LaTeX{} features including mathematical equations, tables, and formatting. You can cite papers using the citation search feature, and they will be automatically added to your \texttt{cite.bib} file.

\section{Mathematical Expressions}
Here are some example equations:

\subsection{Quadratic Formula}
The quadratic formula is given by:
\[x = \frac{-b \pm \sqrt{b^2 - 4ac}}{2a}\]

\subsection{Matrix Example}
A sample matrix:
\[
\begin{pmatrix}
a_{11} & a_{12} & a_{13} \\
a_{21} & a_{22} & a_{23} \\
a_{31} & a_{32} & a_{33}
\end{pmatrix}
\]

\section{Tables}
Below is a sample table showing data:

\begin{table}[h]
\centering
\begin{tabular}{@{}lrrr@{}}
\toprule
\textbf{Category} & \textbf{Value 1} & \textbf{Value 2} & \textbf{Total} \\
\midrule
Item A & 10.5 & 20.3 & 30.8 \\
Item B & 15.2 & 25.1 & 40.3 \\
Item C & 22.1 & 12.7 & 34.8 \\
\bottomrule
\end{tabular}
\caption{Sample Data Table}
\label{tab:sample}
\end{table}

\section{Lists}
\subsection{Bullet Points}
\begin{itemize}
  \item First item
  \item Second item
  \item Third item with sub-items:
    \begin{itemize}
      \item Sub-item 1
      \item Sub-item 2
    \end{itemize}
\end{itemize}

\subsection{Numbered List}
\begin{enumerate}
  \item First step
  \item Second step
  \item Third step
\end{enumerate}

\section{Citations and References}
When you use the citation search feature, references will be automatically managed. For example, you might cite a paper like this: \cite{example2024}. All citations will be automatically added to your \texttt{cite.bib} file.

\section{Conclusion}
This template demonstrates basic \LaTeX{} features including mathematical formulas, tables, and lists. The bibliography system is now set up to work seamlessly with Bibby AI's citation management. Feel free to modify it for your needs!

\bibliographystyle{plain}
\bibliography{cite}

\end{document}\documentclass{article}
\usepackage[utf8]{inputenc}
\usepackage{booktabs}
\usepackage{amsmath}
\usepackage{natbib}

\title{My Document}
\author{Your Name}
\date{\today}

\begin{document}

\maketitle

\section{Introduction}
This is a sample document demonstrating various \LaTeX{} features including mathematical equations, tables, and formatting. You can cite papers using the citation search feature, and they will be automatically added to your \texttt{cite.bib} file.

\section{Mathematical Expressions}
Here are some example equations:

\subsection{Quadratic Formula}
The quadratic formula is given by:
\[x = \frac{-b \pm \sqrt{b^2 - 4ac}}{2a}\]

\subsection{Matrix Example}
A sample matrix:
\[
\begin{pmatrix}
a_{11} & a_{12} & a_{13} \\
a_{21} & a_{22} & a_{23} \\
a_{31} & a_{32} & a_{33}
\end{pmatrix}
\]

\section{Tables}
Below is a sample table showing data:

\begin{table}[h]
\centering
\begin{tabular}{@{}lrrr@{}}
\toprule
\textbf{Category} & \textbf{Value 1} & \textbf{Value 2} & \textbf{Total} \\
\midrule
Item A & 10.5 & 20.3 & 30.8 \\
Item B & 15.2 & 25.1 & 40.3 \\
Item C & 22.1 & 12.7 & 34.8 \\
\bottomrule
\end{tabular}
\caption{Sample Data Table}
\label{tab:sample}
\end{table}

\section{Lists}
\subsection{Bullet Points}
\begin{itemize}
  \item First item
  \item Second item
  \item Third item with sub-items:
    \begin{itemize}
      \item Sub-item 1
      \item Sub-item 2
    \end{itemize}
\end{itemize}

\subsection{Numbered List}
\begin{enumerate}
  \item First step
  \item Second step
  \item Third step
\end{enumerate}

\section{Citations and References}
When you use the citation search feature, references will be automatically managed. For example, you might cite a paper like this: \cite{example2024}. All citations will be automatically added to your \texttt{cite.bib} file.

\section{Conclusion}
This template demonstrates basic \LaTeX{} features including mathematical formulas, tables, and lists. The bibliography system is now set up to work seamlessly with Bibby AI's citation management. Feel free to modify it for your needs!

\bibliographystyle{plain}
\bibliography{cite}

\end{document}\documentclass{article}
\usepackage[utf8]{inputenc}
\usepackage{booktabs}
\usepackage{amsmath}
\usepackage{natbib}

\title{My Document}
\author{Your Name}
\date{\today}

\begin{document}

\maketitle

\section{Introduction}
This is a sample document demonstrating various \LaTeX{} features including mathematical equations, tables, and formatting. You can cite papers using the citation search feature, and they will be automatically added to your \texttt{cite.bib} file.

\section{Mathematical Expressions}
Here are some example equations:

\subsection{Quadratic Formula}
The quadratic formula is given by:
\[x = \frac{-b \pm \sqrt{b^2 - 4ac}}{2a}\]

\subsection{Matrix Example}
A sample matrix:
\[
\begin{pmatrix}
a_{11} & a_{12} & a_{13} \\
a_{21} & a_{22} & a_{23} \\
a_{31} & a_{32} & a_{33}
\end{pmatrix}
\]

\section{Tables}
Below is a sample table showing data:

\begin{table}[h]
\centering
\begin{tabular}{@{}lrrr@{}}
\toprule
\textbf{Category} & \textbf{Value 1} & \textbf{Value 2} & \textbf{Total} \\
\midrule
Item A & 10.5 & 20.3 & 30.8 \\
Item B & 15.2 & 25.1 & 40.3 \\
Item C & 22.1 & 12.7 & 34.8 \\
\bottomrule
\end{tabular}
\caption{Sample Data Table}
\label{tab:sample}
\end{table}

\section{Lists}
\subsection{Bullet Points}
\begin{itemize}
  \item First item
  \item Second item
  \item Third item with sub-items:
    \begin{itemize}
      \item Sub-item 1
      \item Sub-item 2
    \end{itemize}
\end{itemize}

\subsection{Numbered List}
\begin{enumerate}
  \item First step
  \item Second step
  \item Third step
\end{enumerate}

\section{Citations and References}
When you use the citation search feature, references will be automatically managed. For example, you might cite a paper like this: \cite{example2024}. All citations will be automatically added to your \texttt{cite.bib} file.

\section{Conclusion}
This template demonstrates basic \LaTeX{} features including mathematical formulas, tables, and lists. The bibliography system is now set up to work seamlessly with Bibby AI's citation management. Feel free to modify it for your needs!

\bibliographystyle{plain}
\bibliography{cite}

\end{document}\documentclass{article}
\usepackage[utf8]{inputenc}
\usepackage{booktabs}
\usepackage{amsmath}
\usepackage{natbib}

\title{My Document}
\author{Your Name}
\date{\today}

\begin{document}

\maketitle

\section{Introduction}
This is a sample document demonstrating various \LaTeX{} features including mathematical equations, tables, and formatting. You can cite papers using the citation search feature, and they will be automatically added to your \texttt{cite.bib} file.

\section{Mathematical Expressions}
Here are some example equations:

\subsection{Quadratic Formula}
The quadratic formula is given by:
\[x = \frac{-b \pm \sqrt{b^2 - 4ac}}{2a}\]

\subsection{Matrix Example}
A sample matrix:
\[
\begin{pmatrix}
a_{11} & a_{12} & a_{13} \\
a_{21} & a_{22} & a_{23} \\
a_{31} & a_{32} & a_{33}
\end{pmatrix}
\]

\section{Tables}
Below is a sample table showing data:

\begin{table}[h]
\centering
\begin{tabular}{@{}lrrr@{}}
\toprule
\textbf{Category} & \textbf{Value 1} & \textbf{Value 2} & \textbf{Total} \\
\midrule
Item A & 10.5 & 20.3 & 30.8 \\
Item B & 15.2 & 25.1 & 40.3 \\
Item C & 22.1 & 12.7 & 34.8 \\
\bottomrule
\end{tabular}
\caption{Sample Data Table}
\label{tab:sample}
\end{table}

\section{Lists}
\subsection{Bullet Points}
\begin{itemize}
  \item First item
  \item Second item
  \item Third item with sub-items:
    \begin{itemize}
      \item Sub-item 1
      \item Sub-item 2
    \end{itemize}
\end{itemize}

\subsection{Numbered List}
\begin{enumerate}
  \item First step
  \item Second step
  \item Third step
\end{enumerate}

\section{Citations and References}
When you use the citation search feature, references will be automatically managed. For example, you might cite a paper like this: \cite{example2024}. All citations will be automatically added to your \texttt{cite.bib} file.

\section{Conclusion}
This template demonstrates basic \LaTeX{} features including mathematical formulas, tables, and lists. The bibliography system is now set up to work seamlessly with Bibby AI's citation management. Feel free to modify it for your needs!

\bibliographystyle{plain}
\bibliography{cite}

\end{document}\documentclass{article}
\usepackage[utf8]{inputenc}
\usepackage{booktabs}
\usepackage{amsmath}
\usepackage{natbib}

\title{My Document}
\author{Your Name}
\date{\today}

\begin{document}

\maketitle

\section{Introduction}
This is a sample document demonstrating various \LaTeX{} features including mathematical equations, tables, and formatting. You can cite papers using the citation search feature, and they will be automatically added to your \texttt{cite.bib} file.

\section{Mathematical Expressions}
Here are some example equations:

\subsection{Quadratic Formula}
The quadratic formula is given by:
\[x = \frac{-b \pm \sqrt{b^2 - 4ac}}{2a}\]

\subsection{Matrix Example}
A sample matrix:
\[
\begin{pmatrix}
a_{11} & a_{12} & a_{13} \\
a_{21} & a_{22} & a_{23} \\
a_{31} & a_{32} & a_{33}
\end{pmatrix}
\]

\section{Tables}
Below is a sample table showing data:

\begin{table}[h]
\centering
\begin{tabular}{@{}lrrr@{}}
\toprule
\textbf{Category} & \textbf{Value 1} & \textbf{Value 2} & \textbf{Total} \\
\midrule
Item A & 10.5 & 20.3 & 30.8 \\
Item B & 15.2 & 25.1 & 40.3 \\
Item C & 22.1 & 12.7 & 34.8 \\
\bottomrule
\end{tabular}
\caption{Sample Data Table}
\label{tab:sample}
\end{table}

\section{Lists}
\subsection{Bullet Points}
\begin{itemize}
  \item First item
  \item Second item
  \item Third item with sub-items:
    \begin{itemize}
      \item Sub-item 1
      \item Sub-item 2
    \end{itemize}
\end{itemize}

\subsection{Numbered List}
\begin{enumerate}
  \item First step
  \item Second step
  \item Third step
\end{enumerate}

\section{Citations and References}
When you use the citation search feature, references will be automatically managed. For example, you might cite a paper like this: \cite{example2024}. All citations will be automatically added to your \texttt{cite.bib} file.

\section{Conclusion}
This template demonstrates basic \LaTeX{} features including mathematical formulas, tables, and lists. The bibliography system is now set up to work seamlessly with Bibby AI's citation management. Feel free to modify it for your needs!

\bibliographystyle{plain}
\bibliography{cite}

\end{document}\documentclass{article}
\usepackage[utf8]{inputenc}
\usepackage{booktabs}
\usepackage{amsmath}
\usepackage{natbib}

\title{My Document}
\author{Your Name}
\date{\today}

\begin{document}

\maketitle

\section{Introduction}
This is a sample document demonstrating various \LaTeX{} features including mathematical equations, tables, and formatting. You can cite papers using the citation search feature, and they will be automatically added to your \texttt{cite.bib} file.

\section{Mathematical Expressions}
Here are some example equations:

\subsection{Quadratic Formula}
The quadratic formula is given by:
\[x = \frac{-b \pm \sqrt{b^2 - 4ac}}{2a}\]

\subsection{Matrix Example}
A sample matrix:
\[
\begin{pmatrix}
a_{11} & a_{12} & a_{13} \\
a_{21} & a_{22} & a_{23} \\
a_{31} & a_{32} & a_{33}
\end{pmatrix}
\]

\section{Tables}
Below is a sample table showing data:

\begin{table}[h]
\centering
\begin{tabular}{@{}lrrr@{}}
\toprule
\textbf{Category} & \textbf{Value 1} & \textbf{Value 2} & \textbf{Total} \\
\midrule
Item A & 10.5 & 20.3 & 30.8 \\
Item B & 15.2 & 25.1 & 40.3 \\
Item C & 22.1 & 12.7 & 34.8 \\
\bottomrule
\end{tabular}
\caption{Sample Data Table}
\label{tab:sample}
\end{table}

\section{Lists}
\subsection{Bullet Points}
\begin{itemize}
  \item First item
  \item Second item
  \item Third item with sub-items:
    \begin{itemize}
      \item Sub-item 1
      \item Sub-item 2
    \end{itemize}
\end{itemize}

\subsection{Numbered List}
\begin{enumerate}
  \item First step
  \item Second step
  \item Third step
\end{enumerate}

\section{Citations and References}
When you use the citation search feature, references will be automatically managed. For example, you might cite a paper like this: \cite{example2024}. All citations will be automatically added to your \texttt{cite.bib} file.

\section{Conclusion}
This template demonstrates basic \LaTeX{} features including mathematical formulas, tables, and lists. The bibliography system is now set up to work seamlessly with Bibby AI's citation management. Feel free to modify it for your needs!

\bibliographystyle{plain}
\bibliography{cite}

\end{document}\documentclass{article}
\usepackage[utf8]{inputenc}
\usepackage{booktabs}
\usepackage{amsmath}
\usepackage{natbib}

\title{My Document}
\author{Your Name}
\date{\today}

\begin{document}

\maketitle

\section{Introduction}
This is a sample document demonstrating various \LaTeX{} features including mathematical equations, tables, and formatting. You can cite papers using the citation search feature, and they will be automatically added to your \texttt{cite.bib} file.

\section{Mathematical Expressions}
Here are some example equations:

\subsection{Quadratic Formula}
The quadratic formula is given by:
\[x = \frac{-b \pm \sqrt{b^2 - 4ac}}{2a}\]

\subsection{Matrix Example}
A sample matrix:
\[
\begin{pmatrix}
a_{11} & a_{12} & a_{13} \\
a_{21} & a_{22} & a_{23} \\
a_{31} & a_{32} & a_{33}
\end{pmatrix}
\]

\section{Tables}
Below is a sample table showing data:

\begin{table}[h]
\centering
\begin{tabular}{@{}lrrr@{}}
\toprule
\textbf{Category} & \textbf{Value 1} & \textbf{Value 2} & \textbf{Total} \\
\midrule
Item A & 10.5 & 20.3 & 30.8 \\
Item B & 15.2 & 25.1 & 40.3 \\
Item C & 22.1 & 12.7 & 34.8 \\
\bottomrule
\end{tabular}
\caption{Sample Data Table}
\label{tab:sample}
\end{table}

\section{Lists}
\subsection{Bullet Points}
\begin{itemize}
  \item First item
  \item Second item
  \item Third item with sub-items:
    \begin{itemize}
      \item Sub-item 1
      \item Sub-item 2
    \end{itemize}
\end{itemize}

\subsection{Numbered List}
\begin{enumerate}
  \item First step
  \item Second step
  \item Third step
\end{enumerate}

\section{Citations and References}
When you use the citation search feature, references will be automatically managed. For example, you might cite a paper like this: \cite{example2024}. All citations will be automatically added to your \texttt{cite.bib} file.

\section{Conclusion}
This template demonstrates basic \LaTeX{} features including mathematical formulas, tables, and lists. The bibliography system is now set up to work seamlessly with Bibby AI's citation management. Feel free to modify it for your needs!

\bibliographystyle{plain}
\bibliography{cite}

\end{document}\documentclass{article}
\usepackage[utf8]{inputenc}
\usepackage{booktabs}
\usepackage{amsmath}
\usepackage{natbib}

\title{My Document}
\author{Your Name}
\date{\today}

\begin{document}

\maketitle

\section{Introduction}
This is a sample document demonstrating various \LaTeX{} features including mathematical equations, tables, and formatting. You can cite papers using the citation search feature, and they will be automatically added to your \texttt{cite.bib} file.

\section{Mathematical Expressions}
Here are some example equations:

\subsection{Quadratic Formula}
The quadratic formula is given by:
\[x = \frac{-b \pm \sqrt{b^2 - 4ac}}{2a}\]

\subsection{Matrix Example}
A sample matrix:
\[
\begin{pmatrix}
a_{11} & a_{12} & a_{13} \\
a_{21} & a_{22} & a_{23} \\
a_{31} & a_{32} & a_{33}
\end{pmatrix}
\]

\section{Tables}
Below is a sample table showing data:

\begin{table}[h]
\centering
\begin{tabular}{@{}lrrr@{}}
\toprule
\textbf{Category} & \textbf{Value 1} & \textbf{Value 2} & \textbf{Total} \\
\midrule
Item A & 10.5 & 20.3 & 30.8 \\
Item B & 15.2 & 25.1 & 40.3 \\
Item C & 22.1 & 12.7 & 34.8 \\
\bottomrule
\end{tabular}
\caption{Sample Data Table}
\label{tab:sample}
\end{table}

\section{Lists}
\subsection{Bullet Points}
\begin{itemize}
  \item First item
  \item Second item
  \item Third item with sub-items:
    \begin{itemize}
      \item Sub-item 1
      \item Sub-item 2
    \end{itemize}
\end{itemize}

\subsection{Numbered List}
\begin{enumerate}
  \item First step
  \item Second step
  \item Third step
\end{enumerate}

\section{Citations and References}
When you use the citation search feature, references will be automatically managed. For example, you might cite a paper like this: \cite{example2024}. All citations will be automatically added to your \texttt{cite.bib} file.

\section{Conclusion}
This template demonstrates basic \LaTeX{} features including mathematical formulas, tables, and lists. The bibliography system is now set up to work seamlessly with Bibby AI's citation management. Feel free to modify it for your needs!

\bibliographystyle{plain}
\bibliography{cite}

\end{document}\documentclass{article}
\usepackage[utf8]{inputenc}
\usepackage{booktabs}
\usepackage{amsmath}
\usepackage{natbib}

\title{My Document}
\author{Your Name}
\date{\today}

\begin{document}

\maketitle

\section{Introduction}
This is a sample document demonstrating various \LaTeX{} features including mathematical equations, tables, and formatting. You can cite papers using the citation search feature, and they will be automatically added to your \texttt{cite.bib} file.

\section{Mathematical Expressions}
Here are some example equations:

\subsection{Quadratic Formula}
The quadratic formula is given by:
\[x = \frac{-b \pm \sqrt{b^2 - 4ac}}{2a}\]

\subsection{Matrix Example}
A sample matrix:
\[
\begin{pmatrix}
a_{11} & a_{12} & a_{13} \\
a_{21} & a_{22} & a_{23} \\
a_{31} & a_{32} & a_{33}
\end{pmatrix}
\]

\section{Tables}
Below is a sample table showing data:

\begin{table}[h]
\centering
\begin{tabular}{@{}lrrr@{}}
\toprule
\textbf{Category} & \textbf{Value 1} & \textbf{Value 2} & \textbf{Total} \\
\midrule
Item A & 10.5 & 20.3 & 30.8 \\
Item B & 15.2 & 25.1 & 40.3 \\
Item C & 22.1 & 12.7 & 34.8 \\
\bottomrule
\end{tabular}
\caption{Sample Data Table}
\label{tab:sample}
\end{table}

\section{Lists}
\subsection{Bullet Points}
\begin{itemize}
  \item First item
  \item Second item
  \item Third item with sub-items:
    \begin{itemize}
      \item Sub-item 1
      \item Sub-item 2
    \end{itemize}
\end{itemize}

\subsection{Numbered List}
\begin{enumerate}
  \item First step
  \item Second step
  \item Third step
\end{enumerate}

\section{Citations and References}
When you use the citation search feature, references will be automatically managed. For example, you might cite a paper like this: \cite{example2024}. All citations will be automatically added to your \texttt{cite.bib} file.

\section{Conclusion}
This template demonstrates basic \LaTeX{} features including mathematical formulas, tables, and lists. The bibliography system is now set up to work seamlessly with Bibby AI's citation management. Feel free to modify it for your needs!

\bibliographystyle{plain}
\bibliography{cite}

\end{document}\documentclass{article}
\usepackage[utf8]{inputenc}
\usepackage{booktabs}
\usepackage{amsmath}
\usepackage{natbib}

\title{My Document}
\author{Your Name}
\date{\today}

\begin{document}

\maketitle

\section{Introduction}
This is a sample document demonstrating various \LaTeX{} features including mathematical equations, tables, and formatting. You can cite papers using the citation search feature, and they will be automatically added to your \texttt{cite.bib} file.

\section{Mathematical Expressions}
Here are some example equations:

\subsection{Quadratic Formula}
The quadratic formula is given by:
\[x = \frac{-b \pm \sqrt{b^2 - 4ac}}{2a}\]

\subsection{Matrix Example}
A sample matrix:
\[
\begin{pmatrix}
a_{11} & a_{12} & a_{13} \\
a_{21} & a_{22} & a_{23} \\
a_{31} & a_{32} & a_{33}
\end{pmatrix}
\]

\section{Tables}
Below is a sample table showing data:

\begin{table}[h]
\centering
\begin{tabular}{@{}lrrr@{}}
\toprule
\textbf{Category} & \textbf{Value 1} & \textbf{Value 2} & \textbf{Total} \\
\midrule
Item A & 10.5 & 20.3 & 30.8 \\
Item B & 15.2 & 25.1 & 40.3 \\
Item C & 22.1 & 12.7 & 34.8 \\
\bottomrule
\end{tabular}
\caption{Sample Data Table}
\label{tab:sample}
\end{table}

\section{Lists}
\subsection{Bullet Points}
\begin{itemize}
  \item First item
  \item Second item
  \item Third item with sub-items:
    \begin{itemize}
      \item Sub-item 1
      \item Sub-item 2
    \end{itemize}
\end{itemize}

\subsection{Numbered List}
\begin{enumerate}
  \item First step
  \item Second step
  \item Third step
\end{enumerate}

\section{Citations and References}
When you use the citation search feature, references will be automatically managed. For example, you might cite a paper like this: \cite{example2024}. All citations will be automatically added to your \texttt{cite.bib} file.

\section{Conclusion}
This template demonstrates basic \LaTeX{} features including mathematical formulas, tables, and lists. The bibliography system is now set up to work seamlessly with Bibby AI's citation management. Feel free to modify it for your needs!

\bibliographystyle{plain}
\bibliography{cite}

\end{document}\documentclass{article}
\usepackage[utf8]{inputenc}
\usepackage{booktabs}
\usepackage{amsmath}
\usepackage{natbib}

\title{My Document}
\author{Your Name}
\date{\today}

\begin{document}

\maketitle

\section{Introduction}
This is a sample document demonstrating various \LaTeX{} features including mathematical equations, tables, and formatting. You can cite papers using the citation search feature, and they will be automatically added to your \texttt{cite.bib} file.

\section{Mathematical Expressions}
Here are some example equations:

\subsection{Quadratic Formula}
The quadratic formula is given by:
\[x = \frac{-b \pm \sqrt{b^2 - 4ac}}{2a}\]

\subsection{Matrix Example}
A sample matrix:
\[
\begin{pmatrix}
a_{11} & a_{12} & a_{13} \\
a_{21} & a_{22} & a_{23} \\
a_{31} & a_{32} & a_{33}
\end{pmatrix}
\]

\section{Tables}
Below is a sample table showing data:

\begin{table}[h]
\centering
\begin{tabular}{@{}lrrr@{}}
\toprule
\textbf{Category} & \textbf{Value 1} & \textbf{Value 2} & \textbf{Total} \\
\midrule
Item A & 10.5 & 20.3 & 30.8 \\
Item B & 15.2 & 25.1 & 40.3 \\
Item C & 22.1 & 12.7 & 34.8 \\
\bottomrule
\end{tabular}
\caption{Sample Data Table}
\label{tab:sample}
\end{table}

\section{Lists}
\subsection{Bullet Points}
\begin{itemize}
  \item First item
  \item Second item
  \item Third item with sub-items:
    \begin{itemize}
      \item Sub-item 1
      \item Sub-item 2
    \end{itemize}
\end{itemize}

\subsection{Numbered List}
\begin{enumerate}
  \item First step
  \item Second step
  \item Third step
\end{enumerate}

\section{Citations and References}
When you use the citation search feature, references will be automatically managed. For example, you might cite a paper like this: \cite{example2024}. All citations will be automatically added to your \texttt{cite.bib} file.

\section{Conclusion}
This template demonstrates basic \LaTeX{} features including mathematical formulas, tables, and lists. The bibliography system is now set up to work seamlessly with Bibby AI's citation management. Feel free to modify it for your needs!

\bibliographystyle{plain}
\bibliography{cite}

\end{document}\documentclass{article}
\usepackage[utf8]{inputenc}
\usepackage{booktabs}
\usepackage{amsmath}
\usepackage{natbib}

\title{My Document}
\author{Your Name}
\date{\today}

\begin{document}

\maketitle

\section{Introduction}
This is a sample document demonstrating various \LaTeX{} features including mathematical equations, tables, and formatting. You can cite papers using the citation search feature, and they will be automatically added to your \texttt{cite.bib} file.

\section{Mathematical Expressions}
Here are some example equations:

\subsection{Quadratic Formula}
The quadratic formula is given by:
\[x = \frac{-b \pm \sqrt{b^2 - 4ac}}{2a}\]

\subsection{Matrix Example}
A sample matrix:
\[
\begin{pmatrix}
a_{11} & a_{12} & a_{13} \\
a_{21} & a_{22} & a_{23} \\
a_{31} & a_{32} & a_{33}
\end{pmatrix}
\]

\section{Tables}
Below is a sample table showing data:

\begin{table}[h]
\centering
\begin{tabular}{@{}lrrr@{}}
\toprule
\textbf{Category} & \textbf{Value 1} & \textbf{Value 2} & \textbf{Total} \\
\midrule
Item A & 10.5 & 20.3 & 30.8 \\
Item B & 15.2 & 25.1 & 40.3 \\
Item C & 22.1 & 12.7 & 34.8 \\
\bottomrule
\end{tabular}
\caption{Sample Data Table}
\label{tab:sample}
\end{table}

\section{Lists}
\subsection{Bullet Points}
\begin{itemize}
  \item First item
  \item Second item
  \item Third item with sub-items:
    \begin{itemize}
      \item Sub-item 1
      \item Sub-item 2
    \end{itemize}
\end{itemize}

\subsection{Numbered List}
\begin{enumerate}
  \item First step
  \item Second step
  \item Third step
\end{enumerate}

\section{Citations and References}
When you use the citation search feature, references will be automatically managed. For example, you might cite a paper like this: \cite{example2024}. All citations will be automatically added to your \texttt{cite.bib} file.

\section{Conclusion}
This template demonstrates basic \LaTeX{} features including mathematical formulas, tables, and lists. The bibliography system is now set up to work seamlessly with Bibby AI's citation management. Feel free to modify it for your needs!

\bibliographystyle{plain}
\bibliography{cite}

\end{document}\documentclass{article}
\usepackage[utf8]{inputenc}
\usepackage{booktabs}
\usepackage{amsmath}
\usepackage{natbib}

\title{My Document}
\author{Your Name}
\date{\today}

\begin{document}

\maketitle

\section{Introduction}
This is a sample document demonstrating various \LaTeX{} features including mathematical equations, tables, and formatting. You can cite papers using the citation search feature, and they will be automatically added to your \texttt{cite.bib} file.

\section{Mathematical Expressions}
Here are some example equations:

\subsection{Quadratic Formula}
The quadratic formula is given by:
\[x = \frac{-b \pm \sqrt{b^2 - 4ac}}{2a}\]

\subsection{Matrix Example}
A sample matrix:
\[
\begin{pmatrix}
a_{11} & a_{12} & a_{13} \\
a_{21} & a_{22} & a_{23} \\
a_{31} & a_{32} & a_{33}
\end{pmatrix}
\]

\section{Tables}
Below is a sample table showing data:

\begin{table}[h]
\centering
\begin{tabular}{@{}lrrr@{}}
\toprule
\textbf{Category} & \textbf{Value 1} & \textbf{Value 2} & \textbf{Total} \\
\midrule
Item A & 10.5 & 20.3 & 30.8 \\
Item B & 15.2 & 25.1 & 40.3 \\
Item C & 22.1 & 12.7 & 34.8 \\
\bottomrule
\end{tabular}
\caption{Sample Data Table}
\label{tab:sample}
\end{table}

\section{Lists}
\subsection{Bullet Points}
\begin{itemize}
  \item First item
  \item Second item
  \item Third item with sub-items:
    \begin{itemize}
      \item Sub-item 1
      \item Sub-item 2
    \end{itemize}
\end{itemize}

\subsection{Numbered List}
\begin{enumerate}
  \item First step
  \item Second step
  \item Third step
\end{enumerate}

\section{Citations and References}
When you use the citation search feature, references will be automatically managed. For example, you might cite a paper like this: \cite{example2024}. All citations will be automatically added to your \texttt{cite.bib} file.

\section{Conclusion}
This template demonstrates basic \LaTeX{} features including mathematical formulas, tables, and lists. The bibliography system is now set up to work seamlessly with Bibby AI's citation management. Feel free to modify it for your needs!

\bibliographystyle{plain}
\bibliography{cite}

\end{document}\documentclass{article}
\usepackage[utf8]{inputenc}
\usepackage{booktabs}
\usepackage{amsmath}
\usepackage{natbib}

\title{My Document}
\author{Your Name}
\date{\today}

\begin{document}

\maketitle

\section{Introduction}
This is a sample document demonstrating various \LaTeX{} features including mathematical equations, tables, and formatting. You can cite papers using the citation search feature, and they will be automatically added to your \texttt{cite.bib} file.

\section{Mathematical Expressions}
Here are some example equations:

\subsection{Quadratic Formula}
The quadratic formula is given by:
\[x = \frac{-b \pm \sqrt{b^2 - 4ac}}{2a}\]

\subsection{Matrix Example}
A sample matrix:
\[
\begin{pmatrix}
a_{11} & a_{12} & a_{13} \\
a_{21} & a_{22} & a_{23} \\
a_{31} & a_{32} & a_{33}
\end{pmatrix}
\]

\section{Tables}
Below is a sample table showing data:

\begin{table}[h]
\centering
\begin{tabular}{@{}lrrr@{}}
\toprule
\textbf{Category} & \textbf{Value 1} & \textbf{Value 2} & \textbf{Total} \\
\midrule
Item A & 10.5 & 20.3 & 30.8 \\
Item B & 15.2 & 25.1 & 40.3 \\
Item C & 22.1 & 12.7 & 34.8 \\
\bottomrule
\end{tabular}
\caption{Sample Data Table}
\label{tab:sample}
\end{table}

\section{Lists}
\subsection{Bullet Points}
\begin{itemize}
  \item First item
  \item Second item
  \item Third item with sub-items:
    \begin{itemize}
      \item Sub-item 1
      \item Sub-item 2
    \end{itemize}
\end{itemize}

\subsection{Numbered List}
\begin{enumerate}
  \item First step
  \item Second step
  \item Third step
\end{enumerate}

\section{Citations and References}
When you use the citation search feature, references will be automatically managed. For example, you might cite a paper like this: \cite{example2024}. All citations will be automatically added to your \texttt{cite.bib} file.

\section{Conclusion}
This template demonstrates basic \LaTeX{} features including mathematical formulas, tables, and lists. The bibliography system is now set up to work seamlessly with Bibby AI's citation management. Feel free to modify it for your needs!

\bibliographystyle{plain}
\bibliography{cite}

\end{document}\documentclass{article}
\usepackage[utf8]{inputenc}
\usepackage{booktabs}
\usepackage{amsmath}
\usepackage{natbib}

\title{My Document}
\author{Your Name}
\date{\today}

\begin{document}

\maketitle

\section{Introduction}
This is a sample document demonstrating various \LaTeX{} features including mathematical equations, tables, and formatting. You can cite papers using the citation search feature, and they will be automatically added to your \texttt{cite.bib} file.

\section{Mathematical Expressions}
Here are some example equations:

\subsection{Quadratic Formula}
The quadratic formula is given by:
\[x = \frac{-b \pm \sqrt{b^2 - 4ac}}{2a}\]

\subsection{Matrix Example}
A sample matrix:
\[
\begin{pmatrix}
a_{11} & a_{12} & a_{13} \\
a_{21} & a_{22} & a_{23} \\
a_{31} & a_{32} & a_{33}
\end{pmatrix}
\]

\section{Tables}
Below is a sample table showing data:

\begin{table}[h]
\centering
\begin{tabular}{@{}lrrr@{}}
\toprule
\textbf{Category} & \textbf{Value 1} & \textbf{Value 2} & \textbf{Total} \\
\midrule
Item A & 10.5 & 20.3 & 30.8 \\
Item B & 15.2 & 25.1 & 40.3 \\
Item C & 22.1 & 12.7 & 34.8 \\
\bottomrule
\end{tabular}
\caption{Sample Data Table}
\label{tab:sample}
\end{table}

\section{Lists}
\subsection{Bullet Points}
\begin{itemize}
  \item First item
  \item Second item
  \item Third item with sub-items:
    \begin{itemize}
      \item Sub-item 1
      \item Sub-item 2
    \end{itemize}
\end{itemize}

\subsection{Numbered List}
\begin{enumerate}
  \item First step
  \item Second step
  \item Third step
\end{enumerate}

\section{Citations and References}
When you use the citation search feature, references will be automatically managed. For example, you might cite a paper like this: \cite{example2024}. All citations will be automatically added to your \texttt{cite.bib} file.

\section{Conclusion}
This template demonstrates basic \LaTeX{} features including mathematical formulas, tables, and lists. The bibliography system is now set up to work seamlessly with Bibby AI's citation management. Feel free to modify it for your needs!

\bibliographystyle{plain}
\bibliography{cite}

\end{document}\documentclass{article}
\usepackage[utf8]{inputenc}
\usepackage{booktabs}
\usepackage{amsmath}
\usepackage{natbib}

\title{My Document}
\author{Your Name}
\date{\today}

\begin{document}

\maketitle

\section{Introduction}
This is a sample document demonstrating various \LaTeX{} features including mathematical equations, tables, and formatting. You can cite papers using the citation search feature, and they will be automatically added to your \texttt{cite.bib} file.

\section{Mathematical Expressions}
Here are some example equations:

\subsection{Quadratic Formula}
The quadratic formula is given by:
\[x = \frac{-b \pm \sqrt{b^2 - 4ac}}{2a}\]

\subsection{Matrix Example}
A sample matrix:
\[
\begin{pmatrix}
a_{11} & a_{12} & a_{13} \\
a_{21} & a_{22} & a_{23} \\
a_{31} & a_{32} & a_{33}
\end{pmatrix}
\]

\section{Tables}
Below is a sample table showing data:

\begin{table}[h]
\centering
\begin{tabular}{@{}lrrr@{}}
\toprule
\textbf{Category} & \textbf{Value 1} & \textbf{Value 2} & \textbf{Total} \\
\midrule
Item A & 10.5 & 20.3 & 30.8 \\
Item B & 15.2 & 25.1 & 40.3 \\
Item C & 22.1 & 12.7 & 34.8 \\
\bottomrule
\end{tabular}
\caption{Sample Data Table}
\label{tab:sample}
\end{table}

\section{Lists}
\subsection{Bullet Points}
\begin{itemize}
  \item First item
  \item Second item
  \item Third item with sub-items:
    \begin{itemize}
      \item Sub-item 1
      \item Sub-item 2
    \end{itemize}
\end{itemize}

\subsection{Numbered List}
\begin{enumerate}
  \item First step
  \item Second step
  \item Third step
\end{enumerate}

\section{Citations and References}
When you use the citation search feature, references will be automatically managed. For example, you might cite a paper like this: \cite{example2024}. All citations will be automatically added to your \texttt{cite.bib} file.

\section{Conclusion}
This template demonstrates basic \LaTeX{} features including mathematical formulas, tables, and lists. The bibliography system is now set up to work seamlessly with Bibby AI's citation management. Feel free to modify it for your needs!

\bibliographystyle{plain}
\bibliography{cite}

\end{document}\documentclass{article}
\usepackage[utf8]{inputenc}
\usepackage{booktabs}
\usepackage{amsmath}
\usepackage{natbib}

\title{My Document}
\author{Your Name}
\date{\today}

\begin{document}

\maketitle

\section{Introduction}
This is a sample document demonstrating various \LaTeX{} features including mathematical equations, tables, and formatting. You can cite papers using the citation search feature, and they will be automatically added to your \texttt{cite.bib} file.

\section{Mathematical Expressions}
Here are some example equations:

\subsection{Quadratic Formula}
The quadratic formula is given by:
\[x = \frac{-b \pm \sqrt{b^2 - 4ac}}{2a}\]

\subsection{Matrix Example}
A sample matrix:
\[
\begin{pmatrix}
a_{11} & a_{12} & a_{13} \\
a_{21} & a_{22} & a_{23} \\
a_{31} & a_{32} & a_{33}
\end{pmatrix}
\]

\section{Tables}
Below is a sample table showing data:

\begin{table}[h]
\centering
\begin{tabular}{@{}lrrr@{}}
\toprule
\textbf{Category} & \textbf{Value 1} & \textbf{Value 2} & \textbf{Total} \\
\midrule
Item A & 10.5 & 20.3 & 30.8 \\
Item B & 15.2 & 25.1 & 40.3 \\
Item C & 22.1 & 12.7 & 34.8 \\
\bottomrule
\end{tabular}
\caption{Sample Data Table}
\label{tab:sample}
\end{table}

\section{Lists}
\subsection{Bullet Points}
\begin{itemize}
  \item First item
  \item Second item
  \item Third item with sub-items:
    \begin{itemize}
      \item Sub-item 1
      \item Sub-item 2
    \end{itemize}
\end{itemize}

\subsection{Numbered List}
\begin{enumerate}
  \item First step
  \item Second step
  \item Third step
\end{enumerate}

\section{Citations and References}
When you use the citation search feature, references will be automatically managed. For example, you might cite a paper like this: \cite{example2024}. All citations will be automatically added to your \texttt{cite.bib} file.

\section{Conclusion}
This template demonstrates basic \LaTeX{} features including mathematical formulas, tables, and lists. The bibliography system is now set up to work seamlessly with Bibby AI's citation management. Feel free to modify it for your needs!

\bibliographystyle{plain}
\bibliography{cite}

\end{document}\documentclass{article}
\usepackage[utf8]{inputenc}
\usepackage{booktabs}
\usepackage{amsmath}
\usepackage{natbib}

\title{My Document}
\author{Your Name}
\date{\today}

\begin{document}

\maketitle

\section{Introduction}
This is a sample document demonstrating various \LaTeX{} features including mathematical equations, tables, and formatting. You can cite papers using the citation search feature, and they will be automatically added to your \texttt{cite.bib} file.

\section{Mathematical Expressions}
Here are some example equations:

\subsection{Quadratic Formula}
The quadratic formula is given by:
\[x = \frac{-b \pm \sqrt{b^2 - 4ac}}{2a}\]

\subsection{Matrix Example}
A sample matrix:
\[
\begin{pmatrix}
a_{11} & a_{12} & a_{13} \\
a_{21} & a_{22} & a_{23} \\
a_{31} & a_{32} & a_{33}
\end{pmatrix}
\]

\section{Tables}
Below is a sample table showing data:

\begin{table}[h]
\centering
\begin{tabular}{@{}lrrr@{}}
\toprule
\textbf{Category} & \textbf{Value 1} & \textbf{Value 2} & \textbf{Total} \\
\midrule
Item A & 10.5 & 20.3 & 30.8 \\
Item B & 15.2 & 25.1 & 40.3 \\
Item C & 22.1 & 12.7 & 34.8 \\
\bottomrule
\end{tabular}
\caption{Sample Data Table}
\label{tab:sample}
\end{table}

\section{Lists}
\subsection{Bullet Points}
\begin{itemize}
  \item First item
  \item Second item
  \item Third item with sub-items:
    \begin{itemize}
      \item Sub-item 1
      \item Sub-item 2
    \end{itemize}
\end{itemize}

\subsection{Numbered List}
\begin{enumerate}
  \item First step
  \item Second step
  \item Third step
\end{enumerate}

\section{Citations and References}
When you use the citation search feature, references will be automatically managed. For example, you might cite a paper like this: \cite{example2024}. All citations will be automatically added to your \texttt{cite.bib} file.

\section{Conclusion}
This template demonstrates basic \LaTeX{} features including mathematical formulas, tables, and lists. The bibliography system is now set up to work seamlessly with Bibby AI's citation management. Feel free to modify it for your needs!

\bibliographystyle{plain}
\bibliography{cite}

\end{document}\documentclass{article}
\usepackage[utf8]{inputenc}
\usepackage{booktabs}
\usepackage{amsmath}
\usepackage{natbib}

\title{My Document}
\author{Your Name}
\date{\today}

\begin{document}

\maketitle

\section{Introduction}
This is a sample document demonstrating various \LaTeX{} features including mathematical equations, tables, and formatting. You can cite papers using the citation search feature, and they will be automatically added to your \texttt{cite.bib} file.

\section{Mathematical Expressions}
Here are some example equations:

\subsection{Quadratic Formula}
The quadratic formula is given by:
\[x = \frac{-b \pm \sqrt{b^2 - 4ac}}{2a}\]

\subsection{Matrix Example}
A sample matrix:
\[
\begin{pmatrix}
a_{11} & a_{12} & a_{13} \\
a_{21} & a_{22} & a_{23} \\
a_{31} & a_{32} & a_{33}
\end{pmatrix}
\]

\section{Tables}
Below is a sample table showing data:

\begin{table}[h]
\centering
\begin{tabular}{@{}lrrr@{}}
\toprule
\textbf{Category} & \textbf{Value 1} & \textbf{Value 2} & \textbf{Total} \\
\midrule
Item A & 10.5 & 20.3 & 30.8 \\
Item B & 15.2 & 25.1 & 40.3 \\
Item C & 22.1 & 12.7 & 34.8 \\
\bottomrule
\end{tabular}
\caption{Sample Data Table}
\label{tab:sample}
\end{table}

\section{Lists}
\subsection{Bullet Points}
\begin{itemize}
  \item First item
  \item Second item
  \item Third item with sub-items:
    \begin{itemize}
      \item Sub-item 1
      \item Sub-item 2
    \end{itemize}
\end{itemize}

\subsection{Numbered List}
\begin{enumerate}
  \item First step
  \item Second step
  \item Third step
\end{enumerate}

\section{Citations and References}
When you use the citation search feature, references will be automatically managed. For example, you might cite a paper like this: \cite{example2024}. All citations will be automatically added to your \texttt{cite.bib} file.

\section{Conclusion}
This template demonstrates basic \LaTeX{} features including mathematical formulas, tables, and lists. The bibliography system is now set up to work seamlessly with Bibby AI's citation management. Feel free to modify it for your needs!

\bibliographystyle{plain}
\bibliography{cite}

\end{document}\documentclass{article}
\usepackage[utf8]{inputenc}
\usepackage{booktabs}
\usepackage{amsmath}
\usepackage{natbib}

\title{My Document}
\author{Your Name}
\date{\today}

\begin{document}

\maketitle

\section{Introduction}
This is a sample document demonstrating various \LaTeX{} features including mathematical equations, tables, and formatting. You can cite papers using the citation search feature, and they will be automatically added to your \texttt{cite.bib} file.

\section{Mathematical Expressions}
Here are some example equations:

\subsection{Quadratic Formula}
The quadratic formula is given by:
\[x = \frac{-b \pm \sqrt{b^2 - 4ac}}{2a}\]

\subsection{Matrix Example}
A sample matrix:
\[
\begin{pmatrix}
a_{11} & a_{12} & a_{13} \\
a_{21} & a_{22} & a_{23} \\
a_{31} & a_{32} & a_{33}
\end{pmatrix}
\]

\section{Tables}
Below is a sample table showing data:

\begin{table}[h]
\centering
\begin{tabular}{@{}lrrr@{}}
\toprule
\textbf{Category} & \textbf{Value 1} & \textbf{Value 2} & \textbf{Total} \\
\midrule
Item A & 10.5 & 20.3 & 30.8 \\
Item B & 15.2 & 25.1 & 40.3 \\
Item C & 22.1 & 12.7 & 34.8 \\
\bottomrule
\end{tabular}
\caption{Sample Data Table}
\label{tab:sample}
\end{table}

\section{Lists}
\subsection{Bullet Points}
\begin{itemize}
  \item First item
  \item Second item
  \item Third item with sub-items:
    \begin{itemize}
      \item Sub-item 1
      \item Sub-item 2
    \end{itemize}
\end{itemize}

\subsection{Numbered List}
\begin{enumerate}
  \item First step
  \item Second step
  \item Third step
\end{enumerate}

\section{Citations and References}
When you use the citation search feature, references will be automatically managed. For example, you might cite a paper like this: \cite{example2024}. All citations will be automatically added to your \texttt{cite.bib} file.

\section{Conclusion}
This template demonstrates basic \LaTeX{} features including mathematical formulas, tables, and lists. The bibliography system is now set up to work seamlessly with Bibby AI's citation management. Feel free to modify it for your needs!

\bibliographystyle{plain}
\bibliography{cite}

\end{document}\documentclass{article}
\usepackage[utf8]{inputenc}
\usepackage{booktabs}
\usepackage{amsmath}
\usepackage{natbib}

\title{My Document}
\author{Your Name}
\date{\today}

\begin{document}

\maketitle

\section{Introduction}
This is a sample document demonstrating various \LaTeX{} features including mathematical equations, tables, and formatting. You can cite papers using the citation search feature, and they will be automatically added to your \texttt{cite.bib} file.

\section{Mathematical Expressions}
Here are some example equations:

\subsection{Quadratic Formula}
The quadratic formula is given by:
\[x = \frac{-b \pm \sqrt{b^2 - 4ac}}{2a}\]

\subsection{Matrix Example}
A sample matrix:
\[
\begin{pmatrix}
a_{11} & a_{12} & a_{13} \\
a_{21} & a_{22} & a_{23} \\
a_{31} & a_{32} & a_{33}
\end{pmatrix}
\]

\section{Tables}
Below is a sample table showing data:

\begin{table}[h]
\centering
\begin{tabular}{@{}lrrr@{}}
\toprule
\textbf{Category} & \textbf{Value 1} & \textbf{Value 2} & \textbf{Total} \\
\midrule
Item A & 10.5 & 20.3 & 30.8 \\
Item B & 15.2 & 25.1 & 40.3 \\
Item C & 22.1 & 12.7 & 34.8 \\
\bottomrule
\end{tabular}
\caption{Sample Data Table}
\label{tab:sample}
\end{table}

\section{Lists}
\subsection{Bullet Points}
\begin{itemize}
  \item First item
  \item Second item
  \item Third item with sub-items:
    \begin{itemize}
      \item Sub-item 1
      \item Sub-item 2
    \end{itemize}
\end{itemize}

\subsection{Numbered List}
\begin{enumerate}
  \item First step
  \item Second step
  \item Third step
\end{enumerate}

\section{Citations and References}
When you use the citation search feature, references will be automatically managed. For example, you might cite a paper like this: \cite{example2024}. All citations will be automatically added to your \texttt{cite.bib} file.

\section{Conclusion}
This template demonstrates basic \LaTeX{} features including mathematical formulas, tables, and lists. The bibliography system is now set up to work seamlessly with Bibby AI's citation management. Feel free to modify it for your needs!

\bibliographystyle{plain}
\bibliography{cite}

\end{document}\documentclass{article}
\usepackage[utf8]{inputenc}
\usepackage{booktabs}
\usepackage{amsmath}
\usepackage{natbib}

\title{My Document}
\author{Your Name}
\date{\today}

\begin{document}

\maketitle

\section{Introduction}
This is a sample document demonstrating various \LaTeX{} features including mathematical equations, tables, and formatting. You can cite papers using the citation search feature, and they will be automatically added to your \texttt{cite.bib} file.

\section{Mathematical Expressions}
Here are some example equations:

\subsection{Quadratic Formula}
The quadratic formula is given by:
\[x = \frac{-b \pm \sqrt{b^2 - 4ac}}{2a}\]

\subsection{Matrix Example}
A sample matrix:
\[
\begin{pmatrix}
a_{11} & a_{12} & a_{13} \\
a_{21} & a_{22} & a_{23} \\
a_{31} & a_{32} & a_{33}
\end{pmatrix}
\]

\section{Tables}
Below is a sample table showing data:

\begin{table}[h]
\centering
\begin{tabular}{@{}lrrr@{}}
\toprule
\textbf{Category} & \textbf{Value 1} & \textbf{Value 2} & \textbf{Total} \\
\midrule
Item A & 10.5 & 20.3 & 30.8 \\
Item B & 15.2 & 25.1 & 40.3 \\
Item C & 22.1 & 12.7 & 34.8 \\
\bottomrule
\end{tabular}
\caption{Sample Data Table}
\label{tab:sample}
\end{table}

\section{Lists}
\subsection{Bullet Points}
\begin{itemize}
  \item First item
  \item Second item
  \item Third item with sub-items:
    \begin{itemize}
      \item Sub-item 1
      \item Sub-item 2
    \end{itemize}
\end{itemize}

\subsection{Numbered List}
\begin{enumerate}
  \item First step
  \item Second step
  \item Third step
\end{enumerate}

\section{Citations and References}
When you use the citation search feature, references will be automatically managed. For example, you might cite a paper like this: \cite{example2024}. All citations will be automatically added to your \texttt{cite.bib} file.

\section{Conclusion}
This template demonstrates basic \LaTeX{} features including mathematical formulas, tables, and lists. The bibliography system is now set up to work seamlessly with Bibby AI's citation management. Feel free to modify it for your needs!

\bibliographystyle{plain}
\bibliography{cite}

\end{document}\documentclass{article}
\usepackage[utf8]{inputenc}
\usepackage{booktabs}
\usepackage{amsmath}
\usepackage{natbib}

\title{My Document}
\author{Your Name}
\date{\today}

\begin{document}

\maketitle

\section{Introduction}
This is a sample document demonstrating various \LaTeX{} features including mathematical equations, tables, and formatting. You can cite papers using the citation search feature, and they will be automatically added to your \texttt{cite.bib} file.

\section{Mathematical Expressions}
Here are some example equations:

\subsection{Quadratic Formula}
The quadratic formula is given by:
\[x = \frac{-b \pm \sqrt{b^2 - 4ac}}{2a}\]

\subsection{Matrix Example}
A sample matrix:
\[
\begin{pmatrix}
a_{11} & a_{12} & a_{13} \\
a_{21} & a_{22} & a_{23} \\
a_{31} & a_{32} & a_{33}
\end{pmatrix}
\]

\section{Tables}
Below is a sample table showing data:

\begin{table}[h]
\centering
\begin{tabular}{@{}lrrr@{}}
\toprule
\textbf{Category} & \textbf{Value 1} & \textbf{Value 2} & \textbf{Total} \\
\midrule
Item A & 10.5 & 20.3 & 30.8 \\
Item B & 15.2 & 25.1 & 40.3 \\
Item C & 22.1 & 12.7 & 34.8 \\
\bottomrule
\end{tabular}
\caption{Sample Data Table}
\label{tab:sample}
\end{table}

\section{Lists}
\subsection{Bullet Points}
\begin{itemize}
  \item First item
  \item Second item
  \item Third item with sub-items:
    \begin{itemize}
      \item Sub-item 1
      \item Sub-item 2
    \end{itemize}
\end{itemize}

\subsection{Numbered List}
\begin{enumerate}
  \item First step
  \item Second step
  \item Third step
\end{enumerate}

\section{Citations and References}
When you use the citation search feature, references will be automatically managed. For example, you might cite a paper like this: \cite{example2024}. All citations will be automatically added to your \texttt{cite.bib} file.

\section{Conclusion}
This template demonstrates basic \LaTeX{} features including mathematical formulas, tables, and lists. The bibliography system is now set up to work seamlessly with Bibby AI's citation management. Feel free to modify it for your needs!

\bibliographystyle{plain}
\bibliography{cite}

\end{document}\documentclass{article}
\usepackage[utf8]{inputenc}
\usepackage{booktabs}
\usepackage{amsmath}
\usepackage{natbib}

\title{My Document}
\author{Your Name}
\date{\today}

\begin{document}

\maketitle

\section{Introduction}
This is a sample document demonstrating various \LaTeX{} features including mathematical equations, tables, and formatting. You can cite papers using the citation search feature, and they will be automatically added to your \texttt{cite.bib} file.

\section{Mathematical Expressions}
Here are some example equations:

\subsection{Quadratic Formula}
The quadratic formula is given by:
\[x = \frac{-b \pm \sqrt{b^2 - 4ac}}{2a}\]

\subsection{Matrix Example}
A sample matrix:
\[
\begin{pmatrix}
a_{11} & a_{12} & a_{13} \\
a_{21} & a_{22} & a_{23} \\
a_{31} & a_{32} & a_{33}
\end{pmatrix}
\]

\section{Tables}
Below is a sample table showing data:

\begin{table}[h]
\centering
\begin{tabular}{@{}lrrr@{}}
\toprule
\textbf{Category} & \textbf{Value 1} & \textbf{Value 2} & \textbf{Total} \\
\midrule
Item A & 10.5 & 20.3 & 30.8 \\
Item B & 15.2 & 25.1 & 40.3 \\
Item C & 22.1 & 12.7 & 34.8 \\
\bottomrule
\end{tabular}
\caption{Sample Data Table}
\label{tab:sample}
\end{table}

\section{Lists}
\subsection{Bullet Points}
\begin{itemize}
  \item First item
  \item Second item
  \item Third item with sub-items:
    \begin{itemize}
      \item Sub-item 1
      \item Sub-item 2
    \end{itemize}
\end{itemize}

\subsection{Numbered List}
\begin{enumerate}
  \item First step
  \item Second step
  \item Third step
\end{enumerate}

\section{Citations and References}
When you use the citation search feature, references will be automatically managed. For example, you might cite a paper like this: \cite{example2024}. All citations will be automatically added to your \texttt{cite.bib} file.

\section{Conclusion}
This template demonstrates basic \LaTeX{} features including mathematical formulas, tables, and lists. The bibliography system is now set up to work seamlessly with Bibby AI's citation management. Feel free to modify it for your needs!

\bibliographystyle{plain}
\bibliography{cite}

\end{document}\documentclass{article}
\usepackage[utf8]{inputenc}
\usepackage{booktabs}
\usepackage{amsmath}
\usepackage{natbib}

\title{My Document}
\author{Your Name}
\date{\today}

\begin{document}

\maketitle

\section{Introduction}
This is a sample document demonstrating various \LaTeX{} features including mathematical equations, tables, and formatting. You can cite papers using the citation search feature, and they will be automatically added to your \texttt{cite.bib} file.

\section{Mathematical Expressions}
Here are some example equations:

\subsection{Quadratic Formula}
The quadratic formula is given by:
\[x = \frac{-b \pm \sqrt{b^2 - 4ac}}{2a}\]

\subsection{Matrix Example}
A sample matrix:
\[
\begin{pmatrix}
a_{11} & a_{12} & a_{13} \\
a_{21} & a_{22} & a_{23} \\
a_{31} & a_{32} & a_{33}
\end{pmatrix}
\]

\section{Tables}
Below is a sample table showing data:

\begin{table}[h]
\centering
\begin{tabular}{@{}lrrr@{}}
\toprule
\textbf{Category} & \textbf{Value 1} & \textbf{Value 2} & \textbf{Total} \\
\midrule
Item A & 10.5 & 20.3 & 30.8 \\
Item B & 15.2 & 25.1 & 40.3 \\
Item C & 22.1 & 12.7 & 34.8 \\
\bottomrule
\end{tabular}
\caption{Sample Data Table}
\label{tab:sample}
\end{table}

\section{Lists}
\subsection{Bullet Points}
\begin{itemize}
  \item First item
  \item Second item
  \item Third item with sub-items:
    \begin{itemize}
      \item Sub-item 1
      \item Sub-item 2
    \end{itemize}
\end{itemize}

\subsection{Numbered List}
\begin{enumerate}
  \item First step
  \item Second step
  \item Third step
\end{enumerate}

\section{Citations and References}
When you use the citation search feature, references will be automatically managed. For example, you might cite a paper like this: \cite{example2024}. All citations will be automatically added to your \texttt{cite.bib} file.

\section{Conclusion}
This template demonstrates basic \LaTeX{} features including mathematical formulas, tables, and lists. The bibliography system is now set up to work seamlessly with Bibby AI's citation management. Feel free to modify it for your needs!

\bibliographystyle{plain}
\bibliography{cite}

\end{document}\documentclass{article}
\usepackage[utf8]{inputenc}
\usepackage{booktabs}
\usepackage{amsmath}
\usepackage{natbib}

\title{My Document}
\author{Your Name}
\date{\today}

\begin{document}

\maketitle

\section{Introduction}
This is a sample document demonstrating various \LaTeX{} features including mathematical equations, tables, and formatting. You can cite papers using the citation search feature, and they will be automatically added to your \texttt{cite.bib} file.

\section{Mathematical Expressions}
Here are some example equations:

\subsection{Quadratic Formula}
The quadratic formula is given by:
\[x = \frac{-b \pm \sqrt{b^2 - 4ac}}{2a}\]

\subsection{Matrix Example}
A sample matrix:
\[
\begin{pmatrix}
a_{11} & a_{12} & a_{13} \\
a_{21} & a_{22} & a_{23} \\
a_{31} & a_{32} & a_{33}
\end{pmatrix}
\]

\section{Tables}
Below is a sample table showing data:

\begin{table}[h]
\centering
\begin{tabular}{@{}lrrr@{}}
\toprule
\textbf{Category} & \textbf{Value 1} & \textbf{Value 2} & \textbf{Total} \\
\midrule
Item A & 10.5 & 20.3 & 30.8 \\
Item B & 15.2 & 25.1 & 40.3 \\
Item C & 22.1 & 12.7 & 34.8 \\
\bottomrule
\end{tabular}
\caption{Sample Data Table}
\label{tab:sample}
\end{table}

\section{Lists}
\subsection{Bullet Points}
\begin{itemize}
  \item First item
  \item Second item
  \item Third item with sub-items:
    \begin{itemize}
      \item Sub-item 1
      \item Sub-item 2
    \end{itemize}
\end{itemize}

\subsection{Numbered List}
\begin{enumerate}
  \item First step
  \item Second step
  \item Third step
\end{enumerate}

\section{Citations and References}
When you use the citation search feature, references will be automatically managed. For example, you might cite a paper like this: \cite{example2024}. All citations will be automatically added to your \texttt{cite.bib} file.

\section{Conclusion}
This template demonstrates basic \LaTeX{} features including mathematical formulas, tables, and lists. The bibliography system is now set up to work seamlessly with Bibby AI's citation management. Feel free to modify it for your needs!

\bibliographystyle{plain}
\bibliography{cite}

\end{document}\documentclass{article}
\usepackage[utf8]{inputenc}
\usepackage{booktabs}
\usepackage{amsmath}
\usepackage{natbib}

\title{My Document}
\author{Your Name}
\date{\today}

\begin{document}

\maketitle

\section{Introduction}
This is a sample document demonstrating various \LaTeX{} features including mathematical equations, tables, and formatting. You can cite papers using the citation search feature, and they will be automatically added to your \texttt{cite.bib} file.

\section{Mathematical Expressions}
Here are some example equations:

\subsection{Quadratic Formula}
The quadratic formula is given by:
\[x = \frac{-b \pm \sqrt{b^2 - 4ac}}{2a}\]

\subsection{Matrix Example}
A sample matrix:
\[
\begin{pmatrix}
a_{11} & a_{12} & a_{13} \\
a_{21} & a_{22} & a_{23} \\
a_{31} & a_{32} & a_{33}
\end{pmatrix}
\]

\section{Tables}
Below is a sample table showing data:

\begin{table}[h]
\centering
\begin{tabular}{@{}lrrr@{}}
\toprule
\textbf{Category} & \textbf{Value 1} & \textbf{Value 2} & \textbf{Total} \\
\midrule
Item A & 10.5 & 20.3 & 30.8 \\
Item B & 15.2 & 25.1 & 40.3 \\
Item C & 22.1 & 12.7 & 34.8 \\
\bottomrule
\end{tabular}
\caption{Sample Data Table}
\label{tab:sample}
\end{table}

\section{Lists}
\subsection{Bullet Points}
\begin{itemize}
  \item First item
  \item Second item
  \item Third item with sub-items:
    \begin{itemize}
      \item Sub-item 1
      \item Sub-item 2
    \end{itemize}
\end{itemize}

\subsection{Numbered List}
\begin{enumerate}
  \item First step
  \item Second step
  \item Third step
\end{enumerate}

\section{Citations and References}
When you use the citation search feature, references will be automatically managed. For example, you might cite a paper like this: \cite{example2024}. All citations will be automatically added to your \texttt{cite.bib} file.

\section{Conclusion}
This template demonstrates basic \LaTeX{} features including mathematical formulas, tables, and lists. The bibliography system is now set up to work seamlessly with Bibby AI's citation management. Feel free to modify it for your needs!

\bibliographystyle{plain}
\bibliography{cite}

\end{document}\documentclass{article}
\usepackage[utf8]{inputenc}
\usepackage{booktabs}
\usepackage{amsmath}
\usepackage{natbib}

\title{My Document}
\author{Your Name}
\date{\today}

\begin{document}

\maketitle

\section{Introduction}
This is a sample document demonstrating various \LaTeX{} features including mathematical equations, tables, and formatting. You can cite papers using the citation search feature, and they will be automatically added to your \texttt{cite.bib} file.

\section{Mathematical Expressions}
Here are some example equations:

\subsection{Quadratic Formula}
The quadratic formula is given by:
\[x = \frac{-b \pm \sqrt{b^2 - 4ac}}{2a}\]

\subsection{Matrix Example}
A sample matrix:
\[
\begin{pmatrix}
a_{11} & a_{12} & a_{13} \\
a_{21} & a_{22} & a_{23} \\
a_{31} & a_{32} & a_{33}
\end{pmatrix}
\]

\section{Tables}
Below is a sample table showing data:

\begin{table}[h]
\centering
\begin{tabular}{@{}lrrr@{}}
\toprule
\textbf{Category} & \textbf{Value 1} & \textbf{Value 2} & \textbf{Total} \\
\midrule
Item A & 10.5 & 20.3 & 30.8 \\
Item B & 15.2 & 25.1 & 40.3 \\
Item C & 22.1 & 12.7 & 34.8 \\
\bottomrule
\end{tabular}
\caption{Sample Data Table}
\label{tab:sample}
\end{table}

\section{Lists}
\subsection{Bullet Points}
\begin{itemize}
  \item First item
  \item Second item
  \item Third item with sub-items:
    \begin{itemize}
      \item Sub-item 1
      \item Sub-item 2
    \end{itemize}
\end{itemize}

\subsection{Numbered List}
\begin{enumerate}
  \item First step
  \item Second step
  \item Third step
\end{enumerate}

\section{Citations and References}
When you use the citation search feature, references will be automatically managed. For example, you might cite a paper like this: \cite{example2024}. All citations will be automatically added to your \texttt{cite.bib} file.

\section{Conclusion}
This template demonstrates basic \LaTeX{} features including mathematical formulas, tables, and lists. The bibliography system is now set up to work seamlessly with Bibby AI's citation management. Feel free to modify it for your needs!

\bibliographystyle{plain}
\bibliography{cite}

\end{document}\documentclass{article}
\usepackage[utf8]{inputenc}
\usepackage{booktabs}
\usepackage{amsmath}
\usepackage{natbib}

\title{My Document}
\author{Your Name}
\date{\today}

\begin{document}

\maketitle

\section{Introduction}
This is a sample document demonstrating various \LaTeX{} features including mathematical equations, tables, and formatting. You can cite papers using the citation search feature, and they will be automatically added to your \texttt{cite.bib} file.

\section{Mathematical Expressions}
Here are some example equations:

\subsection{Quadratic Formula}
The quadratic formula is given by:
\[x = \frac{-b \pm \sqrt{b^2 - 4ac}}{2a}\]

\subsection{Matrix Example}
A sample matrix:
\[
\begin{pmatrix}
a_{11} & a_{12} & a_{13} \\
a_{21} & a_{22} & a_{23} \\
a_{31} & a_{32} & a_{33}
\end{pmatrix}
\]

\section{Tables}
Below is a sample table showing data:

\begin{table}[h]
\centering
\begin{tabular}{@{}lrrr@{}}
\toprule
\textbf{Category} & \textbf{Value 1} & \textbf{Value 2} & \textbf{Total} \\
\midrule
Item A & 10.5 & 20.3 & 30.8 \\
Item B & 15.2 & 25.1 & 40.3 \\
Item C & 22.1 & 12.7 & 34.8 \\
\bottomrule
\end{tabular}
\caption{Sample Data Table}
\label{tab:sample}
\end{table}

\section{Lists}
\subsection{Bullet Points}
\begin{itemize}
  \item First item
  \item Second item
  \item Third item with sub-items:
    \begin{itemize}
      \item Sub-item 1
      \item Sub-item 2
    \end{itemize}
\end{itemize}

\subsection{Numbered List}
\begin{enumerate}
  \item First step
  \item Second step
  \item Third step
\end{enumerate}

\section{Citations and References}
When you use the citation search feature, references will be automatically managed. For example, you might cite a paper like this: \cite{example2024}. All citations will be automatically added to your \texttt{cite.bib} file.

\section{Conclusion}
This template demonstrates basic \LaTeX{} features including mathematical formulas, tables, and lists. The bibliography system is now set up to work seamlessly with Bibby AI's citation management. Feel free to modify it for your needs!

\bibliographystyle{plain}
\bibliography{cite}

\end{document}\documentclass{article}
\usepackage[utf8]{inputenc}
\usepackage{booktabs}
\usepackage{amsmath}
\usepackage{natbib}

\title{My Document}
\author{Your Name}
\date{\today}

\begin{document}

\maketitle

\section{Introduction}
This is a sample document demonstrating various \LaTeX{} features including mathematical equations, tables, and formatting. You can cite papers using the citation search feature, and they will be automatically added to your \texttt{cite.bib} file.

\section{Mathematical Expressions}
Here are some example equations:

\subsection{Quadratic Formula}
The quadratic formula is given by:
\[x = \frac{-b \pm \sqrt{b^2 - 4ac}}{2a}\]

\subsection{Matrix Example}
A sample matrix:
\[
\begin{pmatrix}
a_{11} & a_{12} & a_{13} \\
a_{21} & a_{22} & a_{23} \\
a_{31} & a_{32} & a_{33}
\end{pmatrix}
\]

\section{Tables}
Below is a sample table showing data:

\begin{table}[h]
\centering
\begin{tabular}{@{}lrrr@{}}
\toprule
\textbf{Category} & \textbf{Value 1} & \textbf{Value 2} & \textbf{Total} \\
\midrule
Item A & 10.5 & 20.3 & 30.8 \\
Item B & 15.2 & 25.1 & 40.3 \\
Item C & 22.1 & 12.7 & 34.8 \\
\bottomrule
\end{tabular}
\caption{Sample Data Table}
\label{tab:sample}
\end{table}

\section{Lists}
\subsection{Bullet Points}
\begin{itemize}
  \item First item
  \item Second item
  \item Third item with sub-items:
    \begin{itemize}
      \item Sub-item 1
      \item Sub-item 2
    \end{itemize}
\end{itemize}

\subsection{Numbered List}
\begin{enumerate}
  \item First step
  \item Second step
  \item Third step
\end{enumerate}

\section{Citations and References}
When you use the citation search feature, references will be automatically managed. For example, you might cite a paper like this: \cite{example2024}. All citations will be automatically added to your \texttt{cite.bib} file.

\section{Conclusion}
This template demonstrates basic \LaTeX{} features including mathematical formulas, tables, and lists. The bibliography system is now set up to work seamlessly with Bibby AI's citation management. Feel free to modify it for your needs!

\bibliographystyle{plain}
\bibliography{cite}

\end{document}\documentclass{article}
\usepackage[utf8]{inputenc}
\usepackage{booktabs}
\usepackage{amsmath}
\usepackage{natbib}

\title{My Document}
\author{Your Name}
\date{\today}

\begin{document}

\maketitle

\section{Introduction}
This is a sample document demonstrating various \LaTeX{} features including mathematical equations, tables, and formatting. You can cite papers using the citation search feature, and they will be automatically added to your \texttt{cite.bib} file.

\section{Mathematical Expressions}
Here are some example equations:

\subsection{Quadratic Formula}
The quadratic formula is given by:
\[x = \frac{-b \pm \sqrt{b^2 - 4ac}}{2a}\]

\subsection{Matrix Example}
A sample matrix:
\[
\begin{pmatrix}
a_{11} & a_{12} & a_{13} \\
a_{21} & a_{22} & a_{23} \\
a_{31} & a_{32} & a_{33}
\end{pmatrix}
\]

\section{Tables}
Below is a sample table showing data:

\begin{table}[h]
\centering
\begin{tabular}{@{}lrrr@{}}
\toprule
\textbf{Category} & \textbf{Value 1} & \textbf{Value 2} & \textbf{Total} \\
\midrule
Item A & 10.5 & 20.3 & 30.8 \\
Item B & 15.2 & 25.1 & 40.3 \\
Item C & 22.1 & 12.7 & 34.8 \\
\bottomrule
\end{tabular}
\caption{Sample Data Table}
\label{tab:sample}
\end{table}

\section{Lists}
\subsection{Bullet Points}
\begin{itemize}
  \item First item
  \item Second item
  \item Third item with sub-items:
    \begin{itemize}
      \item Sub-item 1
      \item Sub-item 2
    \end{itemize}
\end{itemize}

\subsection{Numbered List}
\begin{enumerate}
  \item First step
  \item Second step
  \item Third step
\end{enumerate}

\section{Citations and References}
When you use the citation search feature, references will be automatically managed. For example, you might cite a paper like this: \cite{example2024}. All citations will be automatically added to your \texttt{cite.bib} file.

\section{Conclusion}
This template demonstrates basic \LaTeX{} features including mathematical formulas, tables, and lists. The bibliography system is now set up to work seamlessly with Bibby AI's citation management. Feel free to modify it for your needs!

\bibliographystyle{plain}
\bibliography{cite}

\end{document}\documentclass{article}
\usepackage[utf8]{inputenc}
\usepackage{booktabs}
\usepackage{amsmath}
\usepackage{natbib}

\title{My Document}
\author{Your Name}
\date{\today}

\begin{document}

\maketitle

\section{Introduction}
This is a sample document demonstrating various \LaTeX{} features including mathematical equations, tables, and formatting. You can cite papers using the citation search feature, and they will be automatically added to your \texttt{cite.bib} file.

\section{Mathematical Expressions}
Here are some example equations:

\subsection{Quadratic Formula}
The quadratic formula is given by:
\[x = \frac{-b \pm \sqrt{b^2 - 4ac}}{2a}\]

\subsection{Matrix Example}
A sample matrix:
\[
\begin{pmatrix}
a_{11} & a_{12} & a_{13} \\
a_{21} & a_{22} & a_{23} \\
a_{31} & a_{32} & a_{33}
\end{pmatrix}
\]

\section{Tables}
Below is a sample table showing data:

\begin{table}[h]
\centering
\begin{tabular}{@{}lrrr@{}}
\toprule
\textbf{Category} & \textbf{Value 1} & \textbf{Value 2} & \textbf{Total} \\
\midrule
Item A & 10.5 & 20.3 & 30.8 \\
Item B & 15.2 & 25.1 & 40.3 \\
Item C & 22.1 & 12.7 & 34.8 \\
\bottomrule
\end{tabular}
\caption{Sample Data Table}
\label{tab:sample}
\end{table}

\section{Lists}
\subsection{Bullet Points}
\begin{itemize}
  \item First item
  \item Second item
  \item Third item with sub-items:
    \begin{itemize}
      \item Sub-item 1
      \item Sub-item 2
    \end{itemize}
\end{itemize}

\subsection{Numbered List}
\begin{enumerate}
  \item First step
  \item Second step
  \item Third step
\end{enumerate}

\section{Citations and References}
When you use the citation search feature, references will be automatically managed. For example, you might cite a paper like this: \cite{example2024}. All citations will be automatically added to your \texttt{cite.bib} file.

\section{Conclusion}
This template demonstrates basic \LaTeX{} features including mathematical formulas, tables, and lists. The bibliography system is now set up to work seamlessly with Bibby AI's citation management. Feel free to modify it for your needs!

\bibliographystyle{plain}
\bibliography{cite}

\end{document}\documentclass{article}
\usepackage[utf8]{inputenc}
\usepackage{booktabs}
\usepackage{amsmath}
\usepackage{natbib}

\title{My Document}
\author{Your Name}
\date{\today}

\begin{document}

\maketitle

\section{Introduction}
This is a sample document demonstrating various \LaTeX{} features including mathematical equations, tables, and formatting. You can cite papers using the citation search feature, and they will be automatically added to your \texttt{cite.bib} file.

\section{Mathematical Expressions}
Here are some example equations:

\subsection{Quadratic Formula}
The quadratic formula is given by:
\[x = \frac{-b \pm \sqrt{b^2 - 4ac}}{2a}\]

\subsection{Matrix Example}
A sample matrix:
\[
\begin{pmatrix}
a_{11} & a_{12} & a_{13} \\
a_{21} & a_{22} & a_{23} \\
a_{31} & a_{32} & a_{33}
\end{pmatrix}
\]

\section{Tables}
Below is a sample table showing data:

\begin{table}[h]
\centering
\begin{tabular}{@{}lrrr@{}}
\toprule
\textbf{Category} & \textbf{Value 1} & \textbf{Value 2} & \textbf{Total} \\
\midrule
Item A & 10.5 & 20.3 & 30.8 \\
Item B & 15.2 & 25.1 & 40.3 \\
Item C & 22.1 & 12.7 & 34.8 \\
\bottomrule
\end{tabular}
\caption{Sample Data Table}
\label{tab:sample}
\end{table}

\section{Lists}
\subsection{Bullet Points}
\begin{itemize}
  \item First item
  \item Second item
  \item Third item with sub-items:
    \begin{itemize}
      \item Sub-item 1
      \item Sub-item 2
    \end{itemize}
\end{itemize}

\subsection{Numbered List}
\begin{enumerate}
  \item First step
  \item Second step
  \item Third step
\end{enumerate}

\section{Citations and References}
When you use the citation search feature, references will be automatically managed. For example, you might cite a paper like this: \cite{example2024}. All citations will be automatically added to your \texttt{cite.bib} file.

\section{Conclusion}
This template demonstrates basic \LaTeX{} features including mathematical formulas, tables, and lists. The bibliography system is now set up to work seamlessly with Bibby AI's citation management. Feel free to modify it for your needs!

\bibliographystyle{plain}
\bibliography{cite}

\end{document}\documentclass{article}
\usepackage[utf8]{inputenc}
\usepackage{booktabs}
\usepackage{amsmath}
\usepackage{natbib}

\title{My Document}
\author{Your Name}
\date{\today}

\begin{document}

\maketitle

\section{Introduction}
This is a sample document demonstrating various \LaTeX{} features including mathematical equations, tables, and formatting. You can cite papers using the citation search feature, and they will be automatically added to your \texttt{cite.bib} file.

\section{Mathematical Expressions}
Here are some example equations:

\subsection{Quadratic Formula}
The quadratic formula is given by:
\[x = \frac{-b \pm \sqrt{b^2 - 4ac}}{2a}\]

\subsection{Matrix Example}
A sample matrix:
\[
\begin{pmatrix}
a_{11} & a_{12} & a_{13} \\
a_{21} & a_{22} & a_{23} \\
a_{31} & a_{32} & a_{33}
\end{pmatrix}
\]

\section{Tables}
Below is a sample table showing data:

\begin{table}[h]
\centering
\begin{tabular}{@{}lrrr@{}}
\toprule
\textbf{Category} & \textbf{Value 1} & \textbf{Value 2} & \textbf{Total} \\
\midrule
Item A & 10.5 & 20.3 & 30.8 \\
Item B & 15.2 & 25.1 & 40.3 \\
Item C & 22.1 & 12.7 & 34.8 \\
\bottomrule
\end{tabular}
\caption{Sample Data Table}
\label{tab:sample}
\end{table}

\section{Lists}
\subsection{Bullet Points}
\begin{itemize}
  \item First item
  \item Second item
  \item Third item with sub-items:
    \begin{itemize}
      \item Sub-item 1
      \item Sub-item 2
    \end{itemize}
\end{itemize}

\subsection{Numbered List}
\begin{enumerate}
  \item First step
  \item Second step
  \item Third step
\end{enumerate}

\section{Citations and References}
When you use the citation search feature, references will be automatically managed. For example, you might cite a paper like this: \cite{example2024}. All citations will be automatically added to your \texttt{cite.bib} file.

\section{Conclusion}
This template demonstrates basic \LaTeX{} features including mathematical formulas, tables, and lists. The bibliography system is now set up to work seamlessly with Bibby AI's citation management. Feel free to modify it for your needs!

\bibliographystyle{plain}
\bibliography{cite}

\end{document}\documentclass{article}
\usepackage[utf8]{inputenc}
\usepackage{booktabs}
\usepackage{amsmath}
\usepackage{natbib}

\title{My Document}
\author{Your Name}
\date{\today}

\begin{document}

\maketitle

\section{Introduction}
This is a sample document demonstrating various \LaTeX{} features including mathematical equations, tables, and formatting. You can cite papers using the citation search feature, and they will be automatically added to your \texttt{cite.bib} file.

\section{Mathematical Expressions}
Here are some example equations:

\subsection{Quadratic Formula}
The quadratic formula is given by:
\[x = \frac{-b \pm \sqrt{b^2 - 4ac}}{2a}\]

\subsection{Matrix Example}
A sample matrix:
\[
\begin{pmatrix}
a_{11} & a_{12} & a_{13} \\
a_{21} & a_{22} & a_{23} \\
a_{31} & a_{32} & a_{33}
\end{pmatrix}
\]

\section{Tables}
Below is a sample table showing data:

\begin{table}[h]
\centering
\begin{tabular}{@{}lrrr@{}}
\toprule
\textbf{Category} & \textbf{Value 1} & \textbf{Value 2} & \textbf{Total} \\
\midrule
Item A & 10.5 & 20.3 & 30.8 \\
Item B & 15.2 & 25.1 & 40.3 \\
Item C & 22.1 & 12.7 & 34.8 \\
\bottomrule
\end{tabular}
\caption{Sample Data Table}
\label{tab:sample}
\end{table}

\section{Lists}
\subsection{Bullet Points}
\begin{itemize}
  \item First item
  \item Second item
  \item Third item with sub-items:
    \begin{itemize}
      \item Sub-item 1
      \item Sub-item 2
    \end{itemize}
\end{itemize}

\subsection{Numbered List}
\begin{enumerate}
  \item First step
  \item Second step
  \item Third step
\end{enumerate}

\section{Citations and References}
When you use the citation search feature, references will be automatically managed. For example, you might cite a paper like this: \cite{example2024}. All citations will be automatically added to your \texttt{cite.bib} file.

\section{Conclusion}
This template demonstrates basic \LaTeX{} features including mathematical formulas, tables, and lists. The bibliography system is now set up to work seamlessly with Bibby AI's citation management. Feel free to modify it for your needs!

\bibliographystyle{plain}
\bibliography{cite}

\end{document}\documentclass{article}
\usepackage[utf8]{inputenc}
\usepackage{booktabs}
\usepackage{amsmath}
\usepackage{natbib}

\title{My Document}
\author{Your Name}
\date{\today}

\begin{document}

\maketitle

\section{Introduction}
This is a sample document demonstrating various \LaTeX{} features including mathematical equations, tables, and formatting. You can cite papers using the citation search feature, and they will be automatically added to your \texttt{cite.bib} file.

\section{Mathematical Expressions}
Here are some example equations:

\subsection{Quadratic Formula}
The quadratic formula is given by:
\[x = \frac{-b \pm \sqrt{b^2 - 4ac}}{2a}\]

\subsection{Matrix Example}
A sample matrix:
\[
\begin{pmatrix}
a_{11} & a_{12} & a_{13} \\
a_{21} & a_{22} & a_{23} \\
a_{31} & a_{32} & a_{33}
\end{pmatrix}
\]

\section{Tables}
Below is a sample table showing data:

\begin{table}[h]
\centering
\begin{tabular}{@{}lrrr@{}}
\toprule
\textbf{Category} & \textbf{Value 1} & \textbf{Value 2} & \textbf{Total} \\
\midrule
Item A & 10.5 & 20.3 & 30.8 \\
Item B & 15.2 & 25.1 & 40.3 \\
Item C & 22.1 & 12.7 & 34.8 \\
\bottomrule
\end{tabular}
\caption{Sample Data Table}
\label{tab:sample}
\end{table}

\section{Lists}
\subsection{Bullet Points}
\begin{itemize}
  \item First item
  \item Second item
  \item Third item with sub-items:
    \begin{itemize}
      \item Sub-item 1
      \item Sub-item 2
    \end{itemize}
\end{itemize}

\subsection{Numbered List}
\begin{enumerate}
  \item First step
  \item Second step
  \item Third step
\end{enumerate}

\section{Citations and References}
When you use the citation search feature, references will be automatically managed. For example, you might cite a paper like this: \cite{example2024}. All citations will be automatically added to your \texttt{cite.bib} file.

\section{Conclusion}
This template demonstrates basic \LaTeX{} features including mathematical formulas, tables, and lists. The bibliography system is now set up to work seamlessly with Bibby AI's citation management. Feel free to modify it for your needs!

\bibliographystyle{plain}
\bibliography{cite}

\end{document}\documentclass{article}
\usepackage[utf8]{inputenc}
\usepackage{booktabs}
\usepackage{amsmath}
\usepackage{natbib}

\title{My Document}
\author{Your Name}
\date{\today}

\begin{document}

\maketitle

\section{Introduction}
This is a sample document demonstrating various \LaTeX{} features including mathematical equations, tables, and formatting. You can cite papers using the citation search feature, and they will be automatically added to your \texttt{cite.bib} file.

\section{Mathematical Expressions}
Here are some example equations:

\subsection{Quadratic Formula}
The quadratic formula is given by:
\[x = \frac{-b \pm \sqrt{b^2 - 4ac}}{2a}\]

\subsection{Matrix Example}
A sample matrix:
\[
\begin{pmatrix}
a_{11} & a_{12} & a_{13} \\
a_{21} & a_{22} & a_{23} \\
a_{31} & a_{32} & a_{33}
\end{pmatrix}
\]

\section{Tables}
Below is a sample table showing data:

\begin{table}[h]
\centering
\begin{tabular}{@{}lrrr@{}}
\toprule
\textbf{Category} & \textbf{Value 1} & \textbf{Value 2} & \textbf{Total} \\
\midrule
Item A & 10.5 & 20.3 & 30.8 \\
Item B & 15.2 & 25.1 & 40.3 \\
Item C & 22.1 & 12.7 & 34.8 \\
\bottomrule
\end{tabular}
\caption{Sample Data Table}
\label{tab:sample}
\end{table}

\section{Lists}
\subsection{Bullet Points}
\begin{itemize}
  \item First item
  \item Second item
  \item Third item with sub-items:
    \begin{itemize}
      \item Sub-item 1
      \item Sub-item 2
    \end{itemize}
\end{itemize}

\subsection{Numbered List}
\begin{enumerate}
  \item First step
  \item Second step
  \item Third step
\end{enumerate}

\section{Citations and References}
When you use the citation search feature, references will be automatically managed. For example, you might cite a paper like this: \cite{example2024}. All citations will be automatically added to your \texttt{cite.bib} file.

\section{Conclusion}
This template demonstrates basic \LaTeX{} features including mathematical formulas, tables, and lists. The bibliography system is now set up to work seamlessly with Bibby AI's citation management. Feel free to modify it for your needs!

\bibliographystyle{plain}
\bibliography{cite}

\end{document}\documentclass{article}
\usepackage[utf8]{inputenc}
\usepackage{booktabs}
\usepackage{amsmath}
\usepackage{natbib}

\title{My Document}
\author{Your Name}
\date{\today}

\begin{document}

\maketitle

\section{Introduction}
This is a sample document demonstrating various \LaTeX{} features including mathematical equations, tables, and formatting. You can cite papers using the citation search feature, and they will be automatically added to your \texttt{cite.bib} file.

\section{Mathematical Expressions}
Here are some example equations:

\subsection{Quadratic Formula}
The quadratic formula is given by:
\[x = \frac{-b \pm \sqrt{b^2 - 4ac}}{2a}\]

\subsection{Matrix Example}
A sample matrix:
\[
\begin{pmatrix}
a_{11} & a_{12} & a_{13} \\
a_{21} & a_{22} & a_{23} \\
a_{31} & a_{32} & a_{33}
\end{pmatrix}
\]

\section{Tables}
Below is a sample table showing data:

\begin{table}[h]
\centering
\begin{tabular}{@{}lrrr@{}}
\toprule
\textbf{Category} & \textbf{Value 1} & \textbf{Value 2} & \textbf{Total} \\
\midrule
Item A & 10.5 & 20.3 & 30.8 \\
Item B & 15.2 & 25.1 & 40.3 \\
Item C & 22.1 & 12.7 & 34.8 \\
\bottomrule
\end{tabular}
\caption{Sample Data Table}
\label{tab:sample}
\end{table}

\section{Lists}
\subsection{Bullet Points}
\begin{itemize}
  \item First item
  \item Second item
  \item Third item with sub-items:
    \begin{itemize}
      \item Sub-item 1
      \item Sub-item 2
    \end{itemize}
\end{itemize}

\subsection{Numbered List}
\begin{enumerate}
  \item First step
  \item Second step
  \item Third step
\end{enumerate}

\section{Citations and References}
When you use the citation search feature, references will be automatically managed. For example, you might cite a paper like this: \cite{example2024}. All citations will be automatically added to your \texttt{cite.bib} file.

\section{Conclusion}
This template demonstrates basic \LaTeX{} features including mathematical formulas, tables, and lists. The bibliography system is now set up to work seamlessly with Bibby AI's citation management. Feel free to modify it for your needs!

\bibliographystyle{plain}
\bibliography{cite}

\end{document}\documentclass{article}
\usepackage[utf8]{inputenc}
\usepackage{booktabs}
\usepackage{amsmath}
\usepackage{natbib}

\title{My Document}
\author{Your Name}
\date{\today}

\begin{document}

\maketitle

\section{Introduction}
This is a sample document demonstrating various \LaTeX{} features including mathematical equations, tables, and formatting. You can cite papers using the citation search feature, and they will be automatically added to your \texttt{cite.bib} file.

\section{Mathematical Expressions}
Here are some example equations:

\subsection{Quadratic Formula}
The quadratic formula is given by:
\[x = \frac{-b \pm \sqrt{b^2 - 4ac}}{2a}\]

\subsection{Matrix Example}
A sample matrix:
\[
\begin{pmatrix}
a_{11} & a_{12} & a_{13} \\
a_{21} & a_{22} & a_{23} \\
a_{31} & a_{32} & a_{33}
\end{pmatrix}
\]

\section{Tables}
Below is a sample table showing data:

\begin{table}[h]
\centering
\begin{tabular}{@{}lrrr@{}}
\toprule
\textbf{Category} & \textbf{Value 1} & \textbf{Value 2} & \textbf{Total} \\
\midrule
Item A & 10.5 & 20.3 & 30.8 \\
Item B & 15.2 & 25.1 & 40.3 \\
Item C & 22.1 & 12.7 & 34.8 \\
\bottomrule
\end{tabular}
\caption{Sample Data Table}
\label{tab:sample}
\end{table}

\section{Lists}
\subsection{Bullet Points}
\begin{itemize}
  \item First item
  \item Second item
  \item Third item with sub-items:
    \begin{itemize}
      \item Sub-item 1
      \item Sub-item 2
    \end{itemize}
\end{itemize}

\subsection{Numbered List}
\begin{enumerate}
  \item First step
  \item Second step
  \item Third step
\end{enumerate}

\section{Citations and References}
When you use the citation search feature, references will be automatically managed. For example, you might cite a paper like this: \cite{example2024}. All citations will be automatically added to your \texttt{cite.bib} file.

\section{Conclusion}
This template demonstrates basic \LaTeX{} features including mathematical formulas, tables, and lists. The bibliography system is now set up to work seamlessly with Bibby AI's citation management. Feel free to modify it for your needs!

\bibliographystyle{plain}
\bibliography{cite}

\end{document}\documentclass{article}
\usepackage[utf8]{inputenc}
\usepackage{booktabs}
\usepackage{amsmath}
\usepackage{natbib}

\title{My Document}
\author{Your Name}
\date{\today}

\begin{document}

\maketitle

\section{Introduction}
This is a sample document demonstrating various \LaTeX{} features including mathematical equations, tables, and formatting. You can cite papers using the citation search feature, and they will be automatically added to your \texttt{cite.bib} file.

\section{Mathematical Expressions}
Here are some example equations:

\subsection{Quadratic Formula}
The quadratic formula is given by:
\[x = \frac{-b \pm \sqrt{b^2 - 4ac}}{2a}\]

\subsection{Matrix Example}
A sample matrix:
\[
\begin{pmatrix}
a_{11} & a_{12} & a_{13} \\
a_{21} & a_{22} & a_{23} \\
a_{31} & a_{32} & a_{33}
\end{pmatrix}
\]

\section{Tables}
Below is a sample table showing data:

\begin{table}[h]
\centering
\begin{tabular}{@{}lrrr@{}}
\toprule
\textbf{Category} & \textbf{Value 1} & \textbf{Value 2} & \textbf{Total} \\
\midrule
Item A & 10.5 & 20.3 & 30.8 \\
Item B & 15.2 & 25.1 & 40.3 \\
Item C & 22.1 & 12.7 & 34.8 \\
\bottomrule
\end{tabular}
\caption{Sample Data Table}
\label{tab:sample}
\end{table}

\section{Lists}
\subsection{Bullet Points}
\begin{itemize}
  \item First item
  \item Second item
  \item Third item with sub-items:
    \begin{itemize}
      \item Sub-item 1
      \item Sub-item 2
    \end{itemize}
\end{itemize}

\subsection{Numbered List}
\begin{enumerate}
  \item First step
  \item Second step
  \item Third step
\end{enumerate}

\section{Citations and References}
When you use the citation search feature, references will be automatically managed. For example, you might cite a paper like this: \cite{example2024}. All citations will be automatically added to your \texttt{cite.bib} file.

\section{Conclusion}
This template demonstrates basic \LaTeX{} features including mathematical formulas, tables, and lists. The bibliography system is now set up to work seamlessly with Bibby AI's citation management. Feel free to modify it for your needs!

\bibliographystyle{plain}
\bibliography{cite}

\end{document}\documentclass{article}
\usepackage[utf8]{inputenc}
\usepackage{booktabs}
\usepackage{amsmath}
\usepackage{natbib}

\title{My Document}
\author{Your Name}
\date{\today}

\begin{document}

\maketitle

\section{Introduction}
This is a sample document demonstrating various \LaTeX{} features including mathematical equations, tables, and formatting. You can cite papers using the citation search feature, and they will be automatically added to your \texttt{cite.bib} file.

\section{Mathematical Expressions}
Here are some example equations:

\subsection{Quadratic Formula}
The quadratic formula is given by:
\[x = \frac{-b \pm \sqrt{b^2 - 4ac}}{2a}\]

\subsection{Matrix Example}
A sample matrix:
\[
\begin{pmatrix}
a_{11} & a_{12} & a_{13} \\
a_{21} & a_{22} & a_{23} \\
a_{31} & a_{32} & a_{33}
\end{pmatrix}
\]

\section{Tables}
Below is a sample table showing data:

\begin{table}[h]
\centering
\begin{tabular}{@{}lrrr@{}}
\toprule
\textbf{Category} & \textbf{Value 1} & \textbf{Value 2} & \textbf{Total} \\
\midrule
Item A & 10.5 & 20.3 & 30.8 \\
Item B & 15.2 & 25.1 & 40.3 \\
Item C & 22.1 & 12.7 & 34.8 \\
\bottomrule
\end{tabular}
\caption{Sample Data Table}
\label{tab:sample}
\end{table}

\section{Lists}
\subsection{Bullet Points}
\begin{itemize}
  \item First item
  \item Second item
  \item Third item with sub-items:
    \begin{itemize}
      \item Sub-item 1
      \item Sub-item 2
    \end{itemize}
\end{itemize}

\subsection{Numbered List}
\begin{enumerate}
  \item First step
  \item Second step
  \item Third step
\end{enumerate}

\section{Citations and References}
When you use the citation search feature, references will be automatically managed. For example, you might cite a paper like this: \cite{example2024}. All citations will be automatically added to your \texttt{cite.bib} file.

\section{Conclusion}
This template demonstrates basic \LaTeX{} features including mathematical formulas, tables, and lists. The bibliography system is now set up to work seamlessly with Bibby AI's citation management. Feel free to modify it for your needs!

\bibliographystyle{plain}
\bibliography{cite}

\end{document}\documentclass{article}
\usepackage[utf8]{inputenc}
\usepackage{booktabs}
\usepackage{amsmath}
\usepackage{natbib}

\title{My Document}
\author{Your Name}
\date{\today}

\begin{document}

\maketitle

\section{Introduction}
This is a sample document demonstrating various \LaTeX{} features including mathematical equations, tables, and formatting. You can cite papers using the citation search feature, and they will be automatically added to your \texttt{cite.bib} file.

\section{Mathematical Expressions}
Here are some example equations:

\subsection{Quadratic Formula}
The quadratic formula is given by:
\[x = \frac{-b \pm \sqrt{b^2 - 4ac}}{2a}\]

\subsection{Matrix Example}
A sample matrix:
\[
\begin{pmatrix}
a_{11} & a_{12} & a_{13} \\
a_{21} & a_{22} & a_{23} \\
a_{31} & a_{32} & a_{33}
\end{pmatrix}
\]

\section{Tables}
Below is a sample table showing data:

\begin{table}[h]
\centering
\begin{tabular}{@{}lrrr@{}}
\toprule
\textbf{Category} & \textbf{Value 1} & \textbf{Value 2} & \textbf{Total} \\
\midrule
Item A & 10.5 & 20.3 & 30.8 \\
Item B & 15.2 & 25.1 & 40.3 \\
Item C & 22.1 & 12.7 & 34.8 \\
\bottomrule
\end{tabular}
\caption{Sample Data Table}
\label{tab:sample}
\end{table}

\section{Lists}
\subsection{Bullet Points}
\begin{itemize}
  \item First item
  \item Second item
  \item Third item with sub-items:
    \begin{itemize}
      \item Sub-item 1
      \item Sub-item 2
    \end{itemize}
\end{itemize}

\subsection{Numbered List}
\begin{enumerate}
  \item First step
  \item Second step
  \item Third step
\end{enumerate}

\section{Citations and References}
When you use the citation search feature, references will be automatically managed. For example, you might cite a paper like this: \cite{example2024}. All citations will be automatically added to your \texttt{cite.bib} file.

\section{Conclusion}
This template demonstrates basic \LaTeX{} features including mathematical formulas, tables, and lists. The bibliography system is now set up to work seamlessly with Bibby AI's citation management. Feel free to modify it for your needs!

\bibliographystyle{plain}
\bibliography{cite}

\end{document}\documentclass{article}
\usepackage[utf8]{inputenc}
\usepackage{booktabs}
\usepackage{amsmath}
\usepackage{natbib}

\title{My Document}
\author{Your Name}
\date{\today}

\begin{document}

\maketitle

\section{Introduction}
This is a sample document demonstrating various \LaTeX{} features including mathematical equations, tables, and formatting. You can cite papers using the citation search feature, and they will be automatically added to your \texttt{cite.bib} file.

\section{Mathematical Expressions}
Here are some example equations:

\subsection{Quadratic Formula}
The quadratic formula is given by:
\[x = \frac{-b \pm \sqrt{b^2 - 4ac}}{2a}\]

\subsection{Matrix Example}
A sample matrix:
\[
\begin{pmatrix}
a_{11} & a_{12} & a_{13} \\
a_{21} & a_{22} & a_{23} \\
a_{31} & a_{32} & a_{33}
\end{pmatrix}
\]

\section{Tables}
Below is a sample table showing data:

\begin{table}[h]
\centering
\begin{tabular}{@{}lrrr@{}}
\toprule
\textbf{Category} & \textbf{Value 1} & \textbf{Value 2} & \textbf{Total} \\
\midrule
Item A & 10.5 & 20.3 & 30.8 \\
Item B & 15.2 & 25.1 & 40.3 \\
Item C & 22.1 & 12.7 & 34.8 \\
\bottomrule
\end{tabular}
\caption{Sample Data Table}
\label{tab:sample}
\end{table}

\section{Lists}
\subsection{Bullet Points}
\begin{itemize}
  \item First item
  \item Second item
  \item Third item with sub-items:
    \begin{itemize}
      \item Sub-item 1
      \item Sub-item 2
    \end{itemize}
\end{itemize}

\subsection{Numbered List}
\begin{enumerate}
  \item First step
  \item Second step
  \item Third step
\end{enumerate}

\section{Citations and References}
When you use the citation search feature, references will be automatically managed. For example, you might cite a paper like this: \cite{example2024}. All citations will be automatically added to your \texttt{cite.bib} file.

\section{Conclusion}
This template demonstrates basic \LaTeX{} features including mathematical formulas, tables, and lists. The bibliography system is now set up to work seamlessly with Bibby AI's citation management. Feel free to modify it for your needs!

\bibliographystyle{plain}
\bibliography{cite}

\end{document}\documentclass{article}
\usepackage[utf8]{inputenc}
\usepackage{booktabs}
\usepackage{amsmath}
\usepackage{natbib}

\title{My Document}
\author{Your Name}
\date{\today}

\begin{document}

\maketitle

\section{Introduction}
This is a sample document demonstrating various \LaTeX{} features including mathematical equations, tables, and formatting. You can cite papers using the citation search feature, and they will be automatically added to your \texttt{cite.bib} file.

\section{Mathematical Expressions}
Here are some example equations:

\subsection{Quadratic Formula}
The quadratic formula is given by:
\[x = \frac{-b \pm \sqrt{b^2 - 4ac}}{2a}\]

\subsection{Matrix Example}
A sample matrix:
\[
\begin{pmatrix}
a_{11} & a_{12} & a_{13} \\
a_{21} & a_{22} & a_{23} \\
a_{31} & a_{32} & a_{33}
\end{pmatrix}
\]

\section{Tables}
Below is a sample table showing data:

\begin{table}[h]
\centering
\begin{tabular}{@{}lrrr@{}}
\toprule
\textbf{Category} & \textbf{Value 1} & \textbf{Value 2} & \textbf{Total} \\
\midrule
Item A & 10.5 & 20.3 & 30.8 \\
Item B & 15.2 & 25.1 & 40.3 \\
Item C & 22.1 & 12.7 & 34.8 \\
\bottomrule
\end{tabular}
\caption{Sample Data Table}
\label{tab:sample}
\end{table}

\section{Lists}
\subsection{Bullet Points}
\begin{itemize}
  \item First item
  \item Second item
  \item Third item with sub-items:
    \begin{itemize}
      \item Sub-item 1
      \item Sub-item 2
    \end{itemize}
\end{itemize}

\subsection{Numbered List}
\begin{enumerate}
  \item First step
  \item Second step
  \item Third step
\end{enumerate}

\section{Citations and References}
When you use the citation search feature, references will be automatically managed. For example, you might cite a paper like this: \cite{example2024}. All citations will be automatically added to your \texttt{cite.bib} file.

\section{Conclusion}
This template demonstrates basic \LaTeX{} features including mathematical formulas, tables, and lists. The bibliography system is now set up to work seamlessly with Bibby AI's citation management. Feel free to modify it for your needs!

\bibliographystyle{plain}
\bibliography{cite}

\end{document}\documentclass{article}
\usepackage[utf8]{inputenc}
\usepackage{booktabs}
\usepackage{amsmath}
\usepackage{natbib}

\title{My Document}
\author{Your Name}
\date{\today}

\begin{document}

\maketitle

\section{Introduction}
This is a sample document demonstrating various \LaTeX{} features including mathematical equations, tables, and formatting. You can cite papers using the citation search feature, and they will be automatically added to your \texttt{cite.bib} file.

\section{Mathematical Expressions}
Here are some example equations:

\subsection{Quadratic Formula}
The quadratic formula is given by:
\[x = \frac{-b \pm \sqrt{b^2 - 4ac}}{2a}\]

\subsection{Matrix Example}
A sample matrix:
\[
\begin{pmatrix}
a_{11} & a_{12} & a_{13} \\
a_{21} & a_{22} & a_{23} \\
a_{31} & a_{32} & a_{33}
\end{pmatrix}
\]

\section{Tables}
Below is a sample table showing data:

\begin{table}[h]
\centering
\begin{tabular}{@{}lrrr@{}}
\toprule
\textbf{Category} & \textbf{Value 1} & \textbf{Value 2} & \textbf{Total} \\
\midrule
Item A & 10.5 & 20.3 & 30.8 \\
Item B & 15.2 & 25.1 & 40.3 \\
Item C & 22.1 & 12.7 & 34.8 \\
\bottomrule
\end{tabular}
\caption{Sample Data Table}
\label{tab:sample}
\end{table}

\section{Lists}
\subsection{Bullet Points}
\begin{itemize}
  \item First item
  \item Second item
  \item Third item with sub-items:
    \begin{itemize}
      \item Sub-item 1
      \item Sub-item 2
    \end{itemize}
\end{itemize}

\subsection{Numbered List}
\begin{enumerate}
  \item First step
  \item Second step
  \item Third step
\end{enumerate}

\section{Citations and References}
When you use the citation search feature, references will be automatically managed. For example, you might cite a paper like this: \cite{example2024}. All citations will be automatically added to your \texttt{cite.bib} file.

\section{Conclusion}
This template demonstrates basic \LaTeX{} features including mathematical formulas, tables, and lists. The bibliography system is now set up to work seamlessly with Bibby AI's citation management. Feel free to modify it for your needs!

\bibliographystyle{plain}
\bibliography{cite}

\end{document}\documentclass{article}
\usepackage[utf8]{inputenc}
\usepackage{booktabs}
\usepackage{amsmath}
\usepackage{natbib}

\title{My Document}
\author{Your Name}
\date{\today}

\begin{document}

\maketitle

\section{Introduction}
This is a sample document demonstrating various \LaTeX{} features including mathematical equations, tables, and formatting. You can cite papers using the citation search feature, and they will be automatically added to your \texttt{cite.bib} file.

\section{Mathematical Expressions}
Here are some example equations:

\subsection{Quadratic Formula}
The quadratic formula is given by:
\[x = \frac{-b \pm \sqrt{b^2 - 4ac}}{2a}\]

\subsection{Matrix Example}
A sample matrix:
\[
\begin{pmatrix}
a_{11} & a_{12} & a_{13} \\
a_{21} & a_{22} & a_{23} \\
a_{31} & a_{32} & a_{33}
\end{pmatrix}
\]

\section{Tables}
Below is a sample table showing data:

\begin{table}[h]
\centering
\begin{tabular}{@{}lrrr@{}}
\toprule
\textbf{Category} & \textbf{Value 1} & \textbf{Value 2} & \textbf{Total} \\
\midrule
Item A & 10.5 & 20.3 & 30.8 \\
Item B & 15.2 & 25.1 & 40.3 \\
Item C & 22.1 & 12.7 & 34.8 \\
\bottomrule
\end{tabular}
\caption{Sample Data Table}
\label{tab:sample}
\end{table}

\section{Lists}
\subsection{Bullet Points}
\begin{itemize}
  \item First item
  \item Second item
  \item Third item with sub-items:
    \begin{itemize}
      \item Sub-item 1
      \item Sub-item 2
    \end{itemize}
\end{itemize}

\subsection{Numbered List}
\begin{enumerate}
  \item First step
  \item Second step
  \item Third step
\end{enumerate}

\section{Citations and References}
When you use the citation search feature, references will be automatically managed. For example, you might cite a paper like this: \cite{example2024}. All citations will be automatically added to your \texttt{cite.bib} file.

\section{Conclusion}
This template demonstrates basic \LaTeX{} features including mathematical formulas, tables, and lists. The bibliography system is now set up to work seamlessly with Bibby AI's citation management. Feel free to modify it for your needs!

\bibliographystyle{plain}
\bibliography{cite}

\end{document}\documentclass{article}
\usepackage[utf8]{inputenc}
\usepackage{booktabs}
\usepackage{amsmath}
\usepackage{natbib}

\title{My Document}
\author{Your Name}
\date{\today}

\begin{document}

\maketitle

\section{Introduction}
This is a sample document demonstrating various \LaTeX{} features including mathematical equations, tables, and formatting. You can cite papers using the citation search feature, and they will be automatically added to your \texttt{cite.bib} file.

\section{Mathematical Expressions}
Here are some example equations:

\subsection{Quadratic Formula}
The quadratic formula is given by:
\[x = \frac{-b \pm \sqrt{b^2 - 4ac}}{2a}\]

\subsection{Matrix Example}
A sample matrix:
\[
\begin{pmatrix}
a_{11} & a_{12} & a_{13} \\
a_{21} & a_{22} & a_{23} \\
a_{31} & a_{32} & a_{33}
\end{pmatrix}
\]

\section{Tables}
Below is a sample table showing data:

\begin{table}[h]
\centering
\begin{tabular}{@{}lrrr@{}}
\toprule
\textbf{Category} & \textbf{Value 1} & \textbf{Value 2} & \textbf{Total} \\
\midrule
Item A & 10.5 & 20.3 & 30.8 \\
Item B & 15.2 & 25.1 & 40.3 \\
Item C & 22.1 & 12.7 & 34.8 \\
\bottomrule
\end{tabular}
\caption{Sample Data Table}
\label{tab:sample}
\end{table}

\section{Lists}
\subsection{Bullet Points}
\begin{itemize}
  \item First item
  \item Second item
  \item Third item with sub-items:
    \begin{itemize}
      \item Sub-item 1
      \item Sub-item 2
    \end{itemize}
\end{itemize}

\subsection{Numbered List}
\begin{enumerate}
  \item First step
  \item Second step
  \item Third step
\end{enumerate}

\section{Citations and References}
When you use the citation search feature, references will be automatically managed. For example, you might cite a paper like this: \cite{example2024}. All citations will be automatically added to your \texttt{cite.bib} file.

\section{Conclusion}
This template demonstrates basic \LaTeX{} features including mathematical formulas, tables, and lists. The bibliography system is now set up to work seamlessly with Bibby AI's citation management. Feel free to modify it for your needs!

\bibliographystyle{plain}
\bibliography{cite}

\end{document}\documentclass{article}
\usepackage[utf8]{inputenc}
\usepackage{booktabs}
\usepackage{amsmath}
\usepackage{natbib}

\title{My Document}
\author{Your Name}
\date{\today}

\begin{document}

\maketitle

\section{Introduction}
This is a sample document demonstrating various \LaTeX{} features including mathematical equations, tables, and formatting. You can cite papers using the citation search feature, and they will be automatically added to your \texttt{cite.bib} file.

\section{Mathematical Expressions}
Here are some example equations:

\subsection{Quadratic Formula}
The quadratic formula is given by:
\[x = \frac{-b \pm \sqrt{b^2 - 4ac}}{2a}\]

\subsection{Matrix Example}
A sample matrix:
\[
\begin{pmatrix}
a_{11} & a_{12} & a_{13} \\
a_{21} & a_{22} & a_{23} \\
a_{31} & a_{32} & a_{33}
\end{pmatrix}
\]

\section{Tables}
Below is a sample table showing data:

\begin{table}[h]
\centering
\begin{tabular}{@{}lrrr@{}}
\toprule
\textbf{Category} & \textbf{Value 1} & \textbf{Value 2} & \textbf{Total} \\
\midrule
Item A & 10.5 & 20.3 & 30.8 \\
Item B & 15.2 & 25.1 & 40.3 \\
Item C & 22.1 & 12.7 & 34.8 \\
\bottomrule
\end{tabular}
\caption{Sample Data Table}
\label{tab:sample}
\end{table}

\section{Lists}
\subsection{Bullet Points}
\begin{itemize}
  \item First item
  \item Second item
  \item Third item with sub-items:
    \begin{itemize}
      \item Sub-item 1
      \item Sub-item 2
    \end{itemize}
\end{itemize}

\subsection{Numbered List}
\begin{enumerate}
  \item First step
  \item Second step
  \item Third step
\end{enumerate}

\section{Citations and References}
When you use the citation search feature, references will be automatically managed. For example, you might cite a paper like this: \cite{example2024}. All citations will be automatically added to your \texttt{cite.bib} file.

\section{Conclusion}
This template demonstrates basic \LaTeX{} features including mathematical formulas, tables, and lists. The bibliography system is now set up to work seamlessly with Bibby AI's citation management. Feel free to modify it for your needs!

\bibliographystyle{plain}
\bibliography{cite}

\end{document}\documentclass{article}
\usepackage[utf8]{inputenc}
\usepackage{booktabs}
\usepackage{amsmath}
\usepackage{natbib}

\title{My Document}
\author{Your Name}
\date{\today}

\begin{document}

\maketitle

\section{Introduction}
This is a sample document demonstrating various \LaTeX{} features including mathematical equations, tables, and formatting. You can cite papers using the citation search feature, and they will be automatically added to your \texttt{cite.bib} file.

\section{Mathematical Expressions}
Here are some example equations:

\subsection{Quadratic Formula}
The quadratic formula is given by:
\[x = \frac{-b \pm \sqrt{b^2 - 4ac}}{2a}\]

\subsection{Matrix Example}
A sample matrix:
\[
\begin{pmatrix}
a_{11} & a_{12} & a_{13} \\
a_{21} & a_{22} & a_{23} \\
a_{31} & a_{32} & a_{33}
\end{pmatrix}
\]

\section{Tables}
Below is a sample table showing data:

\begin{table}[h]
\centering
\begin{tabular}{@{}lrrr@{}}
\toprule
\textbf{Category} & \textbf{Value 1} & \textbf{Value 2} & \textbf{Total} \\
\midrule
Item A & 10.5 & 20.3 & 30.8 \\
Item B & 15.2 & 25.1 & 40.3 \\
Item C & 22.1 & 12.7 & 34.8 \\
\bottomrule
\end{tabular}
\caption{Sample Data Table}
\label{tab:sample}
\end{table}

\section{Lists}
\subsection{Bullet Points}
\begin{itemize}
  \item First item
  \item Second item
  \item Third item with sub-items:
    \begin{itemize}
      \item Sub-item 1
      \item Sub-item 2
    \end{itemize}
\end{itemize}

\subsection{Numbered List}
\begin{enumerate}
  \item First step
  \item Second step
  \item Third step
\end{enumerate}

\section{Citations and References}
When you use the citation search feature, references will be automatically managed. For example, you might cite a paper like this: \cite{example2024}. All citations will be automatically added to your \texttt{cite.bib} file.

\section{Conclusion}
This template demonstrates basic \LaTeX{} features including mathematical formulas, tables, and lists. The bibliography system is now set up to work seamlessly with Bibby AI's citation management. Feel free to modify it for your needs!

\bibliographystyle{plain}
\bibliography{cite}

\end{document}\documentclass{article}
\usepackage[utf8]{inputenc}
\usepackage{booktabs}
\usepackage{amsmath}
\usepackage{natbib}

\title{My Document}
\author{Your Name}
\date{\today}

\begin{document}

\maketitle

\section{Introduction}
This is a sample document demonstrating various \LaTeX{} features including mathematical equations, tables, and formatting. You can cite papers using the citation search feature, and they will be automatically added to your \texttt{cite.bib} file.

\section{Mathematical Expressions}
Here are some example equations:

\subsection{Quadratic Formula}
The quadratic formula is given by:
\[x = \frac{-b \pm \sqrt{b^2 - 4ac}}{2a}\]

\subsection{Matrix Example}
A sample matrix:
\[
\begin{pmatrix}
a_{11} & a_{12} & a_{13} \\
a_{21} & a_{22} & a_{23} \\
a_{31} & a_{32} & a_{33}
\end{pmatrix}
\]

\section{Tables}
Below is a sample table showing data:

\begin{table}[h]
\centering
\begin{tabular}{@{}lrrr@{}}
\toprule
\textbf{Category} & \textbf{Value 1} & \textbf{Value 2} & \textbf{Total} \\
\midrule
Item A & 10.5 & 20.3 & 30.8 \\
Item B & 15.2 & 25.1 & 40.3 \\
Item C & 22.1 & 12.7 & 34.8 \\
\bottomrule
\end{tabular}
\caption{Sample Data Table}
\label{tab:sample}
\end{table}

\section{Lists}
\subsection{Bullet Points}
\begin{itemize}
  \item First item
  \item Second item
  \item Third item with sub-items:
    \begin{itemize}
      \item Sub-item 1
      \item Sub-item 2
    \end{itemize}
\end{itemize}

\subsection{Numbered List}
\begin{enumerate}
  \item First step
  \item Second step
  \item Third step
\end{enumerate}

\section{Citations and References}
When you use the citation search feature, references will be automatically managed. For example, you might cite a paper like this: \cite{example2024}. All citations will be automatically added to your \texttt{cite.bib} file.

\section{Conclusion}
This template demonstrates basic \LaTeX{} features including mathematical formulas, tables, and lists. The bibliography system is now set up to work seamlessly with Bibby AI's citation management. Feel free to modify it for your needs!

\bibliographystyle{plain}
\bibliography{cite}

\end{document}\documentclass{article}
\usepackage[utf8]{inputenc}
\usepackage{booktabs}
\usepackage{amsmath}
\usepackage{natbib}

\title{My Document}
\author{Your Name}
\date{\today}

\begin{document}

\maketitle

\section{Introduction}
This is a sample document demonstrating various \LaTeX{} features including mathematical equations, tables, and formatting. You can cite papers using the citation search feature, and they will be automatically added to your \texttt{cite.bib} file.

\section{Mathematical Expressions}
Here are some example equations:

\subsection{Quadratic Formula}
The quadratic formula is given by:
\[x = \frac{-b \pm \sqrt{b^2 - 4ac}}{2a}\]

\subsection{Matrix Example}
A sample matrix:
\[
\begin{pmatrix}
a_{11} & a_{12} & a_{13} \\
a_{21} & a_{22} & a_{23} \\
a_{31} & a_{32} & a_{33}
\end{pmatrix}
\]

\section{Tables}
Below is a sample table showing data:

\begin{table}[h]
\centering
\begin{tabular}{@{}lrrr@{}}
\toprule
\textbf{Category} & \textbf{Value 1} & \textbf{Value 2} & \textbf{Total} \\
\midrule
Item A & 10.5 & 20.3 & 30.8 \\
Item B & 15.2 & 25.1 & 40.3 \\
Item C & 22.1 & 12.7 & 34.8 \\
\bottomrule
\end{tabular}
\caption{Sample Data Table}
\label{tab:sample}
\end{table}

\section{Lists}
\subsection{Bullet Points}
\begin{itemize}
  \item First item
  \item Second item
  \item Third item with sub-items:
    \begin{itemize}
      \item Sub-item 1
      \item Sub-item 2
    \end{itemize}
\end{itemize}

\subsection{Numbered List}
\begin{enumerate}
  \item First step
  \item Second step
  \item Third step
\end{enumerate}

\section{Citations and References}
When you use the citation search feature, references will be automatically managed. For example, you might cite a paper like this: \cite{example2024}. All citations will be automatically added to your \texttt{cite.bib} file.

\section{Conclusion}
This template demonstrates basic \LaTeX{} features including mathematical formulas, tables, and lists. The bibliography system is now set up to work seamlessly with Bibby AI's citation management. Feel free to modify it for your needs!

\bibliographystyle{plain}
\bibliography{cite}

\end{document}\documentclass{article}
\usepackage[utf8]{inputenc}
\usepackage{booktabs}
\usepackage{amsmath}
\usepackage{natbib}

\title{My Document}
\author{Your Name}
\date{\today}

\begin{document}

\maketitle

\section{Introduction}
This is a sample document demonstrating various \LaTeX{} features including mathematical equations, tables, and formatting. You can cite papers using the citation search feature, and they will be automatically added to your \texttt{cite.bib} file.

\section{Mathematical Expressions}
Here are some example equations:

\subsection{Quadratic Formula}
The quadratic formula is given by:
\[x = \frac{-b \pm \sqrt{b^2 - 4ac}}{2a}\]

\subsection{Matrix Example}
A sample matrix:
\[
\begin{pmatrix}
a_{11} & a_{12} & a_{13} \\
a_{21} & a_{22} & a_{23} \\
a_{31} & a_{32} & a_{33}
\end{pmatrix}
\]

\section{Tables}
Below is a sample table showing data:

\begin{table}[h]
\centering
\begin{tabular}{@{}lrrr@{}}
\toprule
\textbf{Category} & \textbf{Value 1} & \textbf{Value 2} & \textbf{Total} \\
\midrule
Item A & 10.5 & 20.3 & 30.8 \\
Item B & 15.2 & 25.1 & 40.3 \\
Item C & 22.1 & 12.7 & 34.8 \\
\bottomrule
\end{tabular}
\caption{Sample Data Table}
\label{tab:sample}
\end{table}

\section{Lists}
\subsection{Bullet Points}
\begin{itemize}
  \item First item
  \item Second item
  \item Third item with sub-items:
    \begin{itemize}
      \item Sub-item 1
      \item Sub-item 2
    \end{itemize}
\end{itemize}

\subsection{Numbered List}
\begin{enumerate}
  \item First step
  \item Second step
  \item Third step
\end{enumerate}

\section{Citations and References}
When you use the citation search feature, references will be automatically managed. For example, you might cite a paper like this: \cite{example2024}. All citations will be automatically added to your \texttt{cite.bib} file.

\section{Conclusion}
This template demonstrates basic \LaTeX{} features including mathematical formulas, tables, and lists. The bibliography system is now set up to work seamlessly with Bibby AI's citation management. Feel free to modify it for your needs!

\bibliographystyle{plain}
\bibliography{cite}

\end{document}\documentclass{article}
\usepackage[utf8]{inputenc}
\usepackage{booktabs}
\usepackage{amsmath}
\usepackage{natbib}

\title{My Document}
\author{Your Name}
\date{\today}

\begin{document}

\maketitle

\section{Introduction}
This is a sample document demonstrating various \LaTeX{} features including mathematical equations, tables, and formatting. You can cite papers using the citation search feature, and they will be automatically added to your \texttt{cite.bib} file.

\section{Mathematical Expressions}
Here are some example equations:

\subsection{Quadratic Formula}
The quadratic formula is given by:
\[x = \frac{-b \pm \sqrt{b^2 - 4ac}}{2a}\]

\subsection{Matrix Example}
A sample matrix:
\[
\begin{pmatrix}
a_{11} & a_{12} & a_{13} \\
a_{21} & a_{22} & a_{23} \\
a_{31} & a_{32} & a_{33}
\end{pmatrix}
\]

\section{Tables}
Below is a sample table showing data:

\begin{table}[h]
\centering
\begin{tabular}{@{}lrrr@{}}
\toprule
\textbf{Category} & \textbf{Value 1} & \textbf{Value 2} & \textbf{Total} \\
\midrule
Item A & 10.5 & 20.3 & 30.8 \\
Item B & 15.2 & 25.1 & 40.3 \\
Item C & 22.1 & 12.7 & 34.8 \\
\bottomrule
\end{tabular}
\caption{Sample Data Table}
\label{tab:sample}
\end{table}

\section{Lists}
\subsection{Bullet Points}
\begin{itemize}
  \item First item
  \item Second item
  \item Third item with sub-items:
    \begin{itemize}
      \item Sub-item 1
      \item Sub-item 2
    \end{itemize}
\end{itemize}

\subsection{Numbered List}
\begin{enumerate}
  \item First step
  \item Second step
  \item Third step
\end{enumerate}

\section{Citations and References}
When you use the citation search feature, references will be automatically managed. For example, you might cite a paper like this: \cite{example2024}. All citations will be automatically added to your \texttt{cite.bib} file.

\section{Conclusion}
This template demonstrates basic \LaTeX{} features including mathematical formulas, tables, and lists. The bibliography system is now set up to work seamlessly with Bibby AI's citation management. Feel free to modify it for your needs!

\bibliographystyle{plain}
\bibliography{cite}

\end{document}\documentclass{article}
\usepackage[utf8]{inputenc}
\usepackage{booktabs}
\usepackage{amsmath}
\usepackage{natbib}

\title{My Document}
\author{Your Name}
\date{\today}

\begin{document}

\maketitle

\section{Introduction}
This is a sample document demonstrating various \LaTeX{} features including mathematical equations, tables, and formatting. You can cite papers using the citation search feature, and they will be automatically added to your \texttt{cite.bib} file.

\section{Mathematical Expressions}
Here are some example equations:

\subsection{Quadratic Formula}
The quadratic formula is given by:
\[x = \frac{-b \pm \sqrt{b^2 - 4ac}}{2a}\]

\subsection{Matrix Example}
A sample matrix:
\[
\begin{pmatrix}
a_{11} & a_{12} & a_{13} \\
a_{21} & a_{22} & a_{23} \\
a_{31} & a_{32} & a_{33}
\end{pmatrix}
\]

\section{Tables}
Below is a sample table showing data:

\begin{table}[h]
\centering
\begin{tabular}{@{}lrrr@{}}
\toprule
\textbf{Category} & \textbf{Value 1} & \textbf{Value 2} & \textbf{Total} \\
\midrule
Item A & 10.5 & 20.3 & 30.8 \\
Item B & 15.2 & 25.1 & 40.3 \\
Item C & 22.1 & 12.7 & 34.8 \\
\bottomrule
\end{tabular}
\caption{Sample Data Table}
\label{tab:sample}
\end{table}

\section{Lists}
\subsection{Bullet Points}
\begin{itemize}
  \item First item
  \item Second item
  \item Third item with sub-items:
    \begin{itemize}
      \item Sub-item 1
      \item Sub-item 2
    \end{itemize}
\end{itemize}

\subsection{Numbered List}
\begin{enumerate}
  \item First step
  \item Second step
  \item Third step
\end{enumerate}

\section{Citations and References}
When you use the citation search feature, references will be automatically managed. For example, you might cite a paper like this: \cite{example2024}. All citations will be automatically added to your \texttt{cite.bib} file.

\section{Conclusion}
This template demonstrates basic \LaTeX{} features including mathematical formulas, tables, and lists. The bibliography system is now set up to work seamlessly with Bibby AI's citation management. Feel free to modify it for your needs!

\bibliographystyle{plain}
\bibliography{cite}

\end{document}\documentclass{article}
\usepackage[utf8]{inputenc}
\usepackage{booktabs}
\usepackage{amsmath}
\usepackage{natbib}

\title{My Document}
\author{Your Name}
\date{\today}

\begin{document}

\maketitle

\section{Introduction}
This is a sample document demonstrating various \LaTeX{} features including mathematical equations, tables, and formatting. You can cite papers using the citation search feature, and they will be automatically added to your \texttt{cite.bib} file.

\section{Mathematical Expressions}
Here are some example equations:

\subsection{Quadratic Formula}
The quadratic formula is given by:
\[x = \frac{-b \pm \sqrt{b^2 - 4ac}}{2a}\]

\subsection{Matrix Example}
A sample matrix:
\[
\begin{pmatrix}
a_{11} & a_{12} & a_{13} \\
a_{21} & a_{22} & a_{23} \\
a_{31} & a_{32} & a_{33}
\end{pmatrix}
\]

\section{Tables}
Below is a sample table showing data:

\begin{table}[h]
\centering
\begin{tabular}{@{}lrrr@{}}
\toprule
\textbf{Category} & \textbf{Value 1} & \textbf{Value 2} & \textbf{Total} \\
\midrule
Item A & 10.5 & 20.3 & 30.8 \\
Item B & 15.2 & 25.1 & 40.3 \\
Item C & 22.1 & 12.7 & 34.8 \\
\bottomrule
\end{tabular}
\caption{Sample Data Table}
\label{tab:sample}
\end{table}

\section{Lists}
\subsection{Bullet Points}
\begin{itemize}
  \item First item
  \item Second item
  \item Third item with sub-items:
    \begin{itemize}
      \item Sub-item 1
      \item Sub-item 2
    \end{itemize}
\end{itemize}

\subsection{Numbered List}
\begin{enumerate}
  \item First step
  \item Second step
  \item Third step
\end{enumerate}

\section{Citations and References}
When you use the citation search feature, references will be automatically managed. For example, you might cite a paper like this: \cite{example2024}. All citations will be automatically added to your \texttt{cite.bib} file.

\section{Conclusion}
This template demonstrates basic \LaTeX{} features including mathematical formulas, tables, and lists. The bibliography system is now set up to work seamlessly with Bibby AI's citation management. Feel free to modify it for your needs!

\bibliographystyle{plain}
\bibliography{cite}

\end{document}\documentclass{article}
\usepackage[utf8]{inputenc}
\usepackage{booktabs}
\usepackage{amsmath}
\usepackage{natbib}

\title{My Document}
\author{Your Name}
\date{\today}

\begin{document}

\maketitle

\section{Introduction}
This is a sample document demonstrating various \LaTeX{} features including mathematical equations, tables, and formatting. You can cite papers using the citation search feature, and they will be automatically added to your \texttt{cite.bib} file.

\section{Mathematical Expressions}
Here are some example equations:

\subsection{Quadratic Formula}
The quadratic formula is given by:
\[x = \frac{-b \pm \sqrt{b^2 - 4ac}}{2a}\]

\subsection{Matrix Example}
A sample matrix:
\[
\begin{pmatrix}
a_{11} & a_{12} & a_{13} \\
a_{21} & a_{22} & a_{23} \\
a_{31} & a_{32} & a_{33}
\end{pmatrix}
\]

\section{Tables}
Below is a sample table showing data:

\begin{table}[h]
\centering
\begin{tabular}{@{}lrrr@{}}
\toprule
\textbf{Category} & \textbf{Value 1} & \textbf{Value 2} & \textbf{Total} \\
\midrule
Item A & 10.5 & 20.3 & 30.8 \\
Item B & 15.2 & 25.1 & 40.3 \\
Item C & 22.1 & 12.7 & 34.8 \\
\bottomrule
\end{tabular}
\caption{Sample Data Table}
\label{tab:sample}
\end{table}

\section{Lists}
\subsection{Bullet Points}
\begin{itemize}
  \item First item
  \item Second item
  \item Third item with sub-items:
    \begin{itemize}
      \item Sub-item 1
      \item Sub-item 2
    \end{itemize}
\end{itemize}

\subsection{Numbered List}
\begin{enumerate}
  \item First step
  \item Second step
  \item Third step
\end{enumerate}

\section{Citations and References}
When you use the citation search feature, references will be automatically managed. For example, you might cite a paper like this: \cite{example2024}. All citations will be automatically added to your \texttt{cite.bib} file.

\section{Conclusion}
This template demonstrates basic \LaTeX{} features including mathematical formulas, tables, and lists. The bibliography system is now set up to work seamlessly with Bibby AI's citation management. Feel free to modify it for your needs!

\bibliographystyle{plain}
\bibliography{cite}

\end{document}\documentclass{article}
\usepackage[utf8]{inputenc}
\usepackage{booktabs}
\usepackage{amsmath}
\usepackage{natbib}

\title{My Document}
\author{Your Name}
\date{\today}

\begin{document}

\maketitle

\section{Introduction}
This is a sample document demonstrating various \LaTeX{} features including mathematical equations, tables, and formatting. You can cite papers using the citation search feature, and they will be automatically added to your \texttt{cite.bib} file.

\section{Mathematical Expressions}
Here are some example equations:

\subsection{Quadratic Formula}
The quadratic formula is given by:
\[x = \frac{-b \pm \sqrt{b^2 - 4ac}}{2a}\]

\subsection{Matrix Example}
A sample matrix:
\[
\begin{pmatrix}
a_{11} & a_{12} & a_{13} \\
a_{21} & a_{22} & a_{23} \\
a_{31} & a_{32} & a_{33}
\end{pmatrix}
\]

\section{Tables}
Below is a sample table showing data:

\begin{table}[h]
\centering
\begin{tabular}{@{}lrrr@{}}
\toprule
\textbf{Category} & \textbf{Value 1} & \textbf{Value 2} & \textbf{Total} \\
\midrule
Item A & 10.5 & 20.3 & 30.8 \\
Item B & 15.2 & 25.1 & 40.3 \\
Item C & 22.1 & 12.7 & 34.8 \\
\bottomrule
\end{tabular}
\caption{Sample Data Table}
\label{tab:sample}
\end{table}

\section{Lists}
\subsection{Bullet Points}
\begin{itemize}
  \item First item
  \item Second item
  \item Third item with sub-items:
    \begin{itemize}
      \item Sub-item 1
      \item Sub-item 2
    \end{itemize}
\end{itemize}

\subsection{Numbered List}
\begin{enumerate}
  \item First step
  \item Second step
  \item Third step
\end{enumerate}

\section{Citations and References}
When you use the citation search feature, references will be automatically managed. For example, you might cite a paper like this: \cite{example2024}. All citations will be automatically added to your \texttt{cite.bib} file.

\section{Conclusion}
This template demonstrates basic \LaTeX{} features including mathematical formulas, tables, and lists. The bibliography system is now set up to work seamlessly with Bibby AI's citation management. Feel free to modify it for your needs!

\bibliographystyle{plain}
\bibliography{cite}

\end{document}\documentclass{article}
\usepackage[utf8]{inputenc}
\usepackage{booktabs}
\usepackage{amsmath}
\usepackage{natbib}

\title{My Document}
\author{Your Name}
\date{\today}

\begin{document}

\maketitle

\section{Introduction}
This is a sample document demonstrating various \LaTeX{} features including mathematical equations, tables, and formatting. You can cite papers using the citation search feature, and they will be automatically added to your \texttt{cite.bib} file.

\section{Mathematical Expressions}
Here are some example equations:

\subsection{Quadratic Formula}
The quadratic formula is given by:
\[x = \frac{-b \pm \sqrt{b^2 - 4ac}}{2a}\]

\subsection{Matrix Example}
A sample matrix:
\[
\begin{pmatrix}
a_{11} & a_{12} & a_{13} \\
a_{21} & a_{22} & a_{23} \\
a_{31} & a_{32} & a_{33}
\end{pmatrix}
\]

\section{Tables}
Below is a sample table showing data:

\begin{table}[h]
\centering
\begin{tabular}{@{}lrrr@{}}
\toprule
\textbf{Category} & \textbf{Value 1} & \textbf{Value 2} & \textbf{Total} \\
\midrule
Item A & 10.5 & 20.3 & 30.8 \\
Item B & 15.2 & 25.1 & 40.3 \\
Item C & 22.1 & 12.7 & 34.8 \\
\bottomrule
\end{tabular}
\caption{Sample Data Table}
\label{tab:sample}
\end{table}

\section{Lists}
\subsection{Bullet Points}
\begin{itemize}
  \item First item
  \item Second item
  \item Third item with sub-items:
    \begin{itemize}
      \item Sub-item 1
      \item Sub-item 2
    \end{itemize}
\end{itemize}

\subsection{Numbered List}
\begin{enumerate}
  \item First step
  \item Second step
  \item Third step
\end{enumerate}

\section{Citations and References}
When you use the citation search feature, references will be automatically managed. For example, you might cite a paper like this: \cite{example2024}. All citations will be automatically added to your \texttt{cite.bib} file.

\section{Conclusion}
This template demonstrates basic \LaTeX{} features including mathematical formulas, tables, and lists. The bibliography system is now set up to work seamlessly with Bibby AI's citation management. Feel free to modify it for your needs!

\bibliographystyle{plain}
\bibliography{cite}

\end{document}\documentclass{article}
\usepackage[utf8]{inputenc}
\usepackage{booktabs}
\usepackage{amsmath}
\usepackage{natbib}

\title{My Document}
\author{Your Name}
\date{\today}

\begin{document}

\maketitle

\section{Introduction}
This is a sample document demonstrating various \LaTeX{} features including mathematical equations, tables, and formatting. You can cite papers using the citation search feature, and they will be automatically added to your \texttt{cite.bib} file.

\section{Mathematical Expressions}
Here are some example equations:

\subsection{Quadratic Formula}
The quadratic formula is given by:
\[x = \frac{-b \pm \sqrt{b^2 - 4ac}}{2a}\]

\subsection{Matrix Example}
A sample matrix:
\[
\begin{pmatrix}
a_{11} & a_{12} & a_{13} \\
a_{21} & a_{22} & a_{23} \\
a_{31} & a_{32} & a_{33}
\end{pmatrix}
\]

\section{Tables}
Below is a sample table showing data:

\begin{table}[h]
\centering
\begin{tabular}{@{}lrrr@{}}
\toprule
\textbf{Category} & \textbf{Value 1} & \textbf{Value 2} & \textbf{Total} \\
\midrule
Item A & 10.5 & 20.3 & 30.8 \\
Item B & 15.2 & 25.1 & 40.3 \\
Item C & 22.1 & 12.7 & 34.8 \\
\bottomrule
\end{tabular}
\caption{Sample Data Table}
\label{tab:sample}
\end{table}

\section{Lists}
\subsection{Bullet Points}
\begin{itemize}
  \item First item
  \item Second item
  \item Third item with sub-items:
    \begin{itemize}
      \item Sub-item 1
      \item Sub-item 2
    \end{itemize}
\end{itemize}

\subsection{Numbered List}
\begin{enumerate}
  \item First step
  \item Second step
  \item Third step
\end{enumerate}

\section{Citations and References}
When you use the citation search feature, references will be automatically managed. For example, you might cite a paper like this: \cite{example2024}. All citations will be automatically added to your \texttt{cite.bib} file.

\section{Conclusion}
This template demonstrates basic \LaTeX{} features including mathematical formulas, tables, and lists. The bibliography system is now set up to work seamlessly with Bibby AI's citation management. Feel free to modify it for your needs!

\bibliographystyle{plain}
\bibliography{cite}

\end{document}\documentclass{article}
\usepackage[utf8]{inputenc}
\usepackage{booktabs}
\usepackage{amsmath}
\usepackage{natbib}

\title{My Document}
\author{Your Name}
\date{\today}

\begin{document}

\maketitle

\section{Introduction}
This is a sample document demonstrating various \LaTeX{} features including mathematical equations, tables, and formatting. You can cite papers using the citation search feature, and they will be automatically added to your \texttt{cite.bib} file.

\section{Mathematical Expressions}
Here are some example equations:

\subsection{Quadratic Formula}
The quadratic formula is given by:
\[x = \frac{-b \pm \sqrt{b^2 - 4ac}}{2a}\]

\subsection{Matrix Example}
A sample matrix:
\[
\begin{pmatrix}
a_{11} & a_{12} & a_{13} \\
a_{21} & a_{22} & a_{23} \\
a_{31} & a_{32} & a_{33}
\end{pmatrix}
\]

\section{Tables}
Below is a sample table showing data:

\begin{table}[h]
\centering
\begin{tabular}{@{}lrrr@{}}
\toprule
\textbf{Category} & \textbf{Value 1} & \textbf{Value 2} & \textbf{Total} \\
\midrule
Item A & 10.5 & 20.3 & 30.8 \\
Item B & 15.2 & 25.1 & 40.3 \\
Item C & 22.1 & 12.7 & 34.8 \\
\bottomrule
\end{tabular}
\caption{Sample Data Table}
\label{tab:sample}
\end{table}

\section{Lists}
\subsection{Bullet Points}
\begin{itemize}
  \item First item
  \item Second item
  \item Third item with sub-items:
    \begin{itemize}
      \item Sub-item 1
      \item Sub-item 2
    \end{itemize}
\end{itemize}

\subsection{Numbered List}
\begin{enumerate}
  \item First step
  \item Second step
  \item Third step
\end{enumerate}

\section{Citations and References}
When you use the citation search feature, references will be automatically managed. For example, you might cite a paper like this: \cite{example2024}. All citations will be automatically added to your \texttt{cite.bib} file.

\section{Conclusion}
This template demonstrates basic \LaTeX{} features including mathematical formulas, tables, and lists. The bibliography system is now set up to work seamlessly with Bibby AI's citation management. Feel free to modify it for your needs!

\bibliographystyle{plain}
\bibliography{cite}

\end{document}\documentclass{article}
\usepackage[utf8]{inputenc}
\usepackage{booktabs}
\usepackage{amsmath}
\usepackage{natbib}

\title{My Document}
\author{Your Name}
\date{\today}

\begin{document}

\maketitle

\section{Introduction}
This is a sample document demonstrating various \LaTeX{} features including mathematical equations, tables, and formatting. You can cite papers using the citation search feature, and they will be automatically added to your \texttt{cite.bib} file.

\section{Mathematical Expressions}
Here are some example equations:

\subsection{Quadratic Formula}
The quadratic formula is given by:
\[x = \frac{-b \pm \sqrt{b^2 - 4ac}}{2a}\]

\subsection{Matrix Example}
A sample matrix:
\[
\begin{pmatrix}
a_{11} & a_{12} & a_{13} \\
a_{21} & a_{22} & a_{23} \\
a_{31} & a_{32} & a_{33}
\end{pmatrix}
\]

\section{Tables}
Below is a sample table showing data:

\begin{table}[h]
\centering
\begin{tabular}{@{}lrrr@{}}
\toprule
\textbf{Category} & \textbf{Value 1} & \textbf{Value 2} & \textbf{Total} \\
\midrule
Item A & 10.5 & 20.3 & 30.8 \\
Item B & 15.2 & 25.1 & 40.3 \\
Item C & 22.1 & 12.7 & 34.8 \\
\bottomrule
\end{tabular}
\caption{Sample Data Table}
\label{tab:sample}
\end{table}

\section{Lists}
\subsection{Bullet Points}
\begin{itemize}
  \item First item
  \item Second item
  \item Third item with sub-items:
    \begin{itemize}
      \item Sub-item 1
      \item Sub-item 2
    \end{itemize}
\end{itemize}

\subsection{Numbered List}
\begin{enumerate}
  \item First step
  \item Second step
  \item Third step
\end{enumerate}

\section{Citations and References}
When you use the citation search feature, references will be automatically managed. For example, you might cite a paper like this: \cite{example2024}. All citations will be automatically added to your \texttt{cite.bib} file.

\section{Conclusion}
This template demonstrates basic \LaTeX{} features including mathematical formulas, tables, and lists. The bibliography system is now set up to work seamlessly with Bibby AI's citation management. Feel free to modify it for your needs!

\bibliographystyle{plain}
\bibliography{cite}

\end{document}\documentclass{article}
\usepackage[utf8]{inputenc}
\usepackage{booktabs}
\usepackage{amsmath}
\usepackage{natbib}

\title{My Document}
\author{Your Name}
\date{\today}

\begin{document}

\maketitle

\section{Introduction}
This is a sample document demonstrating various \LaTeX{} features including mathematical equations, tables, and formatting. You can cite papers using the citation search feature, and they will be automatically added to your \texttt{cite.bib} file.

\section{Mathematical Expressions}
Here are some example equations:

\subsection{Quadratic Formula}
The quadratic formula is given by:
\[x = \frac{-b \pm \sqrt{b^2 - 4ac}}{2a}\]

\subsection{Matrix Example}
A sample matrix:
\[
\begin{pmatrix}
a_{11} & a_{12} & a_{13} \\
a_{21} & a_{22} & a_{23} \\
a_{31} & a_{32} & a_{33}
\end{pmatrix}
\]

\section{Tables}
Below is a sample table showing data:

\begin{table}[h]
\centering
\begin{tabular}{@{}lrrr@{}}
\toprule
\textbf{Category} & \textbf{Value 1} & \textbf{Value 2} & \textbf{Total} \\
\midrule
Item A & 10.5 & 20.3 & 30.8 \\
Item B & 15.2 & 25.1 & 40.3 \\
Item C & 22.1 & 12.7 & 34.8 \\
\bottomrule
\end{tabular}
\caption{Sample Data Table}
\label{tab:sample}
\end{table}

\section{Lists}
\subsection{Bullet Points}
\begin{itemize}
  \item First item
  \item Second item
  \item Third item with sub-items:
    \begin{itemize}
      \item Sub-item 1
      \item Sub-item 2
    \end{itemize}
\end{itemize}

\subsection{Numbered List}
\begin{enumerate}
  \item First step
  \item Second step
  \item Third step
\end{enumerate}

\section{Citations and References}
When you use the citation search feature, references will be automatically managed. For example, you might cite a paper like this: \cite{example2024}. All citations will be automatically added to your \texttt{cite.bib} file.

\section{Conclusion}
This template demonstrates basic \LaTeX{} features including mathematical formulas, tables, and lists. The bibliography system is now set up to work seamlessly with Bibby AI's citation management. Feel free to modify it for your needs!

\bibliographystyle{plain}
\bibliography{cite}

\end{document}\documentclass{article}
\usepackage[utf8]{inputenc}
\usepackage{booktabs}
\usepackage{amsmath}
\usepackage{natbib}

\title{My Document}
\author{Your Name}
\date{\today}

\begin{document}

\maketitle

\section{Introduction}
This is a sample document demonstrating various \LaTeX{} features including mathematical equations, tables, and formatting. You can cite papers using the citation search feature, and they will be automatically added to your \texttt{cite.bib} file.

\section{Mathematical Expressions}
Here are some example equations:

\subsection{Quadratic Formula}
The quadratic formula is given by:
\[x = \frac{-b \pm \sqrt{b^2 - 4ac}}{2a}\]

\subsection{Matrix Example}
A sample matrix:
\[
\begin{pmatrix}
a_{11} & a_{12} & a_{13} \\
a_{21} & a_{22} & a_{23} \\
a_{31} & a_{32} & a_{33}
\end{pmatrix}
\]

\section{Tables}
Below is a sample table showing data:

\begin{table}[h]
\centering
\begin{tabular}{@{}lrrr@{}}
\toprule
\textbf{Category} & \textbf{Value 1} & \textbf{Value 2} & \textbf{Total} \\
\midrule
Item A & 10.5 & 20.3 & 30.8 \\
Item B & 15.2 & 25.1 & 40.3 \\
Item C & 22.1 & 12.7 & 34.8 \\
\bottomrule
\end{tabular}
\caption{Sample Data Table}
\label{tab:sample}
\end{table}

\section{Lists}
\subsection{Bullet Points}
\begin{itemize}
  \item First item
  \item Second item
  \item Third item with sub-items:
    \begin{itemize}
      \item Sub-item 1
      \item Sub-item 2
    \end{itemize}
\end{itemize}

\subsection{Numbered List}
\begin{enumerate}
  \item First step
  \item Second step
  \item Third step
\end{enumerate}

\section{Citations and References}
When you use the citation search feature, references will be automatically managed. For example, you might cite a paper like this: \cite{example2024}. All citations will be automatically added to your \texttt{cite.bib} file.

\section{Conclusion}
This template demonstrates basic \LaTeX{} features including mathematical formulas, tables, and lists. The bibliography system is now set up to work seamlessly with Bibby AI's citation management. Feel free to modify it for your needs!

\bibliographystyle{plain}
\bibliography{cite}

\end{document}\documentclass{article}
\usepackage[utf8]{inputenc}
\usepackage{booktabs}
\usepackage{amsmath}
\usepackage{natbib}

\title{My Document}
\author{Your Name}
\date{\today}

\begin{document}

\maketitle

\section{Introduction}
This is a sample document demonstrating various \LaTeX{} features including mathematical equations, tables, and formatting. You can cite papers using the citation search feature, and they will be automatically added to your \texttt{cite.bib} file.

\section{Mathematical Expressions}
Here are some example equations:

\subsection{Quadratic Formula}
The quadratic formula is given by:
\[x = \frac{-b \pm \sqrt{b^2 - 4ac}}{2a}\]

\subsection{Matrix Example}
A sample matrix:
\[
\begin{pmatrix}
a_{11} & a_{12} & a_{13} \\
a_{21} & a_{22} & a_{23} \\
a_{31} & a_{32} & a_{33}
\end{pmatrix}
\]

\section{Tables}
Below is a sample table showing data:

\begin{table}[h]
\centering
\begin{tabular}{@{}lrrr@{}}
\toprule
\textbf{Category} & \textbf{Value 1} & \textbf{Value 2} & \textbf{Total} \\
\midrule
Item A & 10.5 & 20.3 & 30.8 \\
Item B & 15.2 & 25.1 & 40.3 \\
Item C & 22.1 & 12.7 & 34.8 \\
\bottomrule
\end{tabular}
\caption{Sample Data Table}
\label{tab:sample}
\end{table}

\section{Lists}
\subsection{Bullet Points}
\begin{itemize}
  \item First item
  \item Second item
  \item Third item with sub-items:
    \begin{itemize}
      \item Sub-item 1
      \item Sub-item 2
    \end{itemize}
\end{itemize}

\subsection{Numbered List}
\begin{enumerate}
  \item First step
  \item Second step
  \item Third step
\end{enumerate}

\section{Citations and References}
When you use the citation search feature, references will be automatically managed. For example, you might cite a paper like this: \cite{example2024}. All citations will be automatically added to your \texttt{cite.bib} file.

\section{Conclusion}
This template demonstrates basic \LaTeX{} features including mathematical formulas, tables, and lists. The bibliography system is now set up to work seamlessly with Bibby AI's citation management. Feel free to modify it for your needs!

\bibliographystyle{plain}
\bibliography{cite}

\end{document}\documentclass{article}
\usepackage[utf8]{inputenc}
\usepackage{booktabs}
\usepackage{amsmath}
\usepackage{natbib}

\title{My Document}
\author{Your Name}
\date{\today}

\begin{document}

\maketitle

\section{Introduction}
This is a sample document demonstrating various \LaTeX{} features including mathematical equations, tables, and formatting. You can cite papers using the citation search feature, and they will be automatically added to your \texttt{cite.bib} file.

\section{Mathematical Expressions}
Here are some example equations:

\subsection{Quadratic Formula}
The quadratic formula is given by:
\[x = \frac{-b \pm \sqrt{b^2 - 4ac}}{2a}\]

\subsection{Matrix Example}
A sample matrix:
\[
\begin{pmatrix}
a_{11} & a_{12} & a_{13} \\
a_{21} & a_{22} & a_{23} \\
a_{31} & a_{32} & a_{33}
\end{pmatrix}
\]

\section{Tables}
Below is a sample table showing data:

\begin{table}[h]
\centering
\begin{tabular}{@{}lrrr@{}}
\toprule
\textbf{Category} & \textbf{Value 1} & \textbf{Value 2} & \textbf{Total} \\
\midrule
Item A & 10.5 & 20.3 & 30.8 \\
Item B & 15.2 & 25.1 & 40.3 \\
Item C & 22.1 & 12.7 & 34.8 \\
\bottomrule
\end{tabular}
\caption{Sample Data Table}
\label{tab:sample}
\end{table}

\section{Lists}
\subsection{Bullet Points}
\begin{itemize}
  \item First item
  \item Second item
  \item Third item with sub-items:
    \begin{itemize}
      \item Sub-item 1
      \item Sub-item 2
    \end{itemize}
\end{itemize}

\subsection{Numbered List}
\begin{enumerate}
  \item First step
  \item Second step
  \item Third step
\end{enumerate}

\section{Citations and References}
When you use the citation search feature, references will be automatically managed. For example, you might cite a paper like this: \cite{example2024}. All citations will be automatically added to your \texttt{cite.bib} file.

\section{Conclusion}
This template demonstrates basic \LaTeX{} features including mathematical formulas, tables, and lists. The bibliography system is now set up to work seamlessly with Bibby AI's citation management. Feel free to modify it for your needs!

\bibliographystyle{plain}
\bibliography{cite}

\end{document}\documentclass{article}
\usepackage[utf8]{inputenc}
\usepackage{booktabs}
\usepackage{amsmath}
\usepackage{natbib}

\title{My Document}
\author{Your Name}
\date{\today}

\begin{document}

\maketitle

\section{Introduction}
This is a sample document demonstrating various \LaTeX{} features including mathematical equations, tables, and formatting. You can cite papers using the citation search feature, and they will be automatically added to your \texttt{cite.bib} file.

\section{Mathematical Expressions}
Here are some example equations:

\subsection{Quadratic Formula}
The quadratic formula is given by:
\[x = \frac{-b \pm \sqrt{b^2 - 4ac}}{2a}\]

\subsection{Matrix Example}
A sample matrix:
\[
\begin{pmatrix}
a_{11} & a_{12} & a_{13} \\
a_{21} & a_{22} & a_{23} \\
a_{31} & a_{32} & a_{33}
\end{pmatrix}
\]

\section{Tables}
Below is a sample table showing data:

\begin{table}[h]
\centering
\begin{tabular}{@{}lrrr@{}}
\toprule
\textbf{Category} & \textbf{Value 1} & \textbf{Value 2} & \textbf{Total} \\
\midrule
Item A & 10.5 & 20.3 & 30.8 \\
Item B & 15.2 & 25.1 & 40.3 \\
Item C & 22.1 & 12.7 & 34.8 \\
\bottomrule
\end{tabular}
\caption{Sample Data Table}
\label{tab:sample}
\end{table}

\section{Lists}
\subsection{Bullet Points}
\begin{itemize}
  \item First item
  \item Second item
  \item Third item with sub-items:
    \begin{itemize}
      \item Sub-item 1
      \item Sub-item 2
    \end{itemize}
\end{itemize}

\subsection{Numbered List}
\begin{enumerate}
  \item First step
  \item Second step
  \item Third step
\end{enumerate}

\section{Citations and References}
When you use the citation search feature, references will be automatically managed. For example, you might cite a paper like this: \cite{example2024}. All citations will be automatically added to your \texttt{cite.bib} file.

\section{Conclusion}
This template demonstrates basic \LaTeX{} features including mathematical formulas, tables, and lists. The bibliography system is now set up to work seamlessly with Bibby AI's citation management. Feel free to modify it for your needs!

\bibliographystyle{plain}
\bibliography{cite}

\end{document}\documentclass{article}
\usepackage[utf8]{inputenc}
\usepackage{booktabs}
\usepackage{amsmath}
\usepackage{natbib}

\title{My Document}
\author{Your Name}
\date{\today}

\begin{document}

\maketitle

\section{Introduction}
This is a sample document demonstrating various \LaTeX{} features including mathematical equations, tables, and formatting. You can cite papers using the citation search feature, and they will be automatically added to your \texttt{cite.bib} file.

\section{Mathematical Expressions}
Here are some example equations:

\subsection{Quadratic Formula}
The quadratic formula is given by:
\[x = \frac{-b \pm \sqrt{b^2 - 4ac}}{2a}\]

\subsection{Matrix Example}
A sample matrix:
\[
\begin{pmatrix}
a_{11} & a_{12} & a_{13} \\
a_{21} & a_{22} & a_{23} \\
a_{31} & a_{32} & a_{33}
\end{pmatrix}
\]

\section{Tables}
Below is a sample table showing data:

\begin{table}[h]
\centering
\begin{tabular}{@{}lrrr@{}}
\toprule
\textbf{Category} & \textbf{Value 1} & \textbf{Value 2} & \textbf{Total} \\
\midrule
Item A & 10.5 & 20.3 & 30.8 \\
Item B & 15.2 & 25.1 & 40.3 \\
Item C & 22.1 & 12.7 & 34.8 \\
\bottomrule
\end{tabular}
\caption{Sample Data Table}
\label{tab:sample}
\end{table}

\section{Lists}
\subsection{Bullet Points}
\begin{itemize}
  \item First item
  \item Second item
  \item Third item with sub-items:
    \begin{itemize}
      \item Sub-item 1
      \item Sub-item 2
    \end{itemize}
\end{itemize}

\subsection{Numbered List}
\begin{enumerate}
  \item First step
  \item Second step
  \item Third step
\end{enumerate}

\section{Citations and References}
When you use the citation search feature, references will be automatically managed. For example, you might cite a paper like this: \cite{example2024}. All citations will be automatically added to your \texttt{cite.bib} file.

\section{Conclusion}
This template demonstrates basic \LaTeX{} features including mathematical formulas, tables, and lists. The bibliography system is now set up to work seamlessly with Bibby AI's citation management. Feel free to modify it for your needs!

\bibliographystyle{plain}
\bibliography{cite}

\end{document}\documentclass{article}
\usepackage[utf8]{inputenc}
\usepackage{booktabs}
\usepackage{amsmath}
\usepackage{natbib}

\title{My Document}
\author{Your Name}
\date{\today}

\begin{document}

\maketitle

\section{Introduction}
This is a sample document demonstrating various \LaTeX{} features including mathematical equations, tables, and formatting. You can cite papers using the citation search feature, and they will be automatically added to your \texttt{cite.bib} file.

\section{Mathematical Expressions}
Here are some example equations:

\subsection{Quadratic Formula}
The quadratic formula is given by:
\[x = \frac{-b \pm \sqrt{b^2 - 4ac}}{2a}\]

\subsection{Matrix Example}
A sample matrix:
\[
\begin{pmatrix}
a_{11} & a_{12} & a_{13} \\
a_{21} & a_{22} & a_{23} \\
a_{31} & a_{32} & a_{33}
\end{pmatrix}
\]

\section{Tables}
Below is a sample table showing data:

\begin{table}[h]
\centering
\begin{tabular}{@{}lrrr@{}}
\toprule
\textbf{Category} & \textbf{Value 1} & \textbf{Value 2} & \textbf{Total} \\
\midrule
Item A & 10.5 & 20.3 & 30.8 \\
Item B & 15.2 & 25.1 & 40.3 \\
Item C & 22.1 & 12.7 & 34.8 \\
\bottomrule
\end{tabular}
\caption{Sample Data Table}
\label{tab:sample}
\end{table}

\section{Lists}
\subsection{Bullet Points}
\begin{itemize}
  \item First item
  \item Second item
  \item Third item with sub-items:
    \begin{itemize}
      \item Sub-item 1
      \item Sub-item 2
    \end{itemize}
\end{itemize}

\subsection{Numbered List}
\begin{enumerate}
  \item First step
  \item Second step
  \item Third step
\end{enumerate}

\section{Citations and References}
When you use the citation search feature, references will be automatically managed. For example, you might cite a paper like this: \cite{example2024}. All citations will be automatically added to your \texttt{cite.bib} file.

\section{Conclusion}
This template demonstrates basic \LaTeX{} features including mathematical formulas, tables, and lists. The bibliography system is now set up to work seamlessly with Bibby AI's citation management. Feel free to modify it for your needs!

\bibliographystyle{plain}
\bibliography{cite}

\end{document}\documentclass{article}
\usepackage[utf8]{inputenc}
\usepackage{booktabs}
\usepackage{amsmath}
\usepackage{natbib}

\title{My Document}
\author{Your Name}
\date{\today}

\begin{document}

\maketitle

\section{Introduction}
This is a sample document demonstrating various \LaTeX{} features including mathematical equations, tables, and formatting. You can cite papers using the citation search feature, and they will be automatically added to your \texttt{cite.bib} file.

\section{Mathematical Expressions}
Here are some example equations:

\subsection{Quadratic Formula}
The quadratic formula is given by:
\[x = \frac{-b \pm \sqrt{b^2 - 4ac}}{2a}\]

\subsection{Matrix Example}
A sample matrix:
\[
\begin{pmatrix}
a_{11} & a_{12} & a_{13} \\
a_{21} & a_{22} & a_{23} \\
a_{31} & a_{32} & a_{33}
\end{pmatrix}
\]

\section{Tables}
Below is a sample table showing data:

\begin{table}[h]
\centering
\begin{tabular}{@{}lrrr@{}}
\toprule
\textbf{Category} & \textbf{Value 1} & \textbf{Value 2} & \textbf{Total} \\
\midrule
Item A & 10.5 & 20.3 & 30.8 \\
Item B & 15.2 & 25.1 & 40.3 \\
Item C & 22.1 & 12.7 & 34.8 \\
\bottomrule
\end{tabular}
\caption{Sample Data Table}
\label{tab:sample}
\end{table}

\section{Lists}
\subsection{Bullet Points}
\begin{itemize}
  \item First item
  \item Second item
  \item Third item with sub-items:
    \begin{itemize}
      \item Sub-item 1
      \item Sub-item 2
    \end{itemize}
\end{itemize}

\subsection{Numbered List}
\begin{enumerate}
  \item First step
  \item Second step
  \item Third step
\end{enumerate}

\section{Citations and References}
When you use the citation search feature, references will be automatically managed. For example, you might cite a paper like this: \cite{example2024}. All citations will be automatically added to your \texttt{cite.bib} file.

\section{Conclusion}
This template demonstrates basic \LaTeX{} features including mathematical formulas, tables, and lists. The bibliography system is now set up to work seamlessly with Bibby AI's citation management. Feel free to modify it for your needs!

\bibliographystyle{plain}
\bibliography{cite}

\end{document}\documentclass{article}
\usepackage[utf8]{inputenc}
\usepackage{booktabs}
\usepackage{amsmath}
\usepackage{natbib}

\title{My Document}
\author{Your Name}
\date{\today}

\begin{document}

\maketitle

\section{Introduction}
This is a sample document demonstrating various \LaTeX{} features including mathematical equations, tables, and formatting. You can cite papers using the citation search feature, and they will be automatically added to your \texttt{cite.bib} file.

\section{Mathematical Expressions}
Here are some example equations:

\subsection{Quadratic Formula}
The quadratic formula is given by:
\[x = \frac{-b \pm \sqrt{b^2 - 4ac}}{2a}\]

\subsection{Matrix Example}
A sample matrix:
\[
\begin{pmatrix}
a_{11} & a_{12} & a_{13} \\
a_{21} & a_{22} & a_{23} \\
a_{31} & a_{32} & a_{33}
\end{pmatrix}
\]

\section{Tables}
Below is a sample table showing data:

\begin{table}[h]
\centering
\begin{tabular}{@{}lrrr@{}}
\toprule
\textbf{Category} & \textbf{Value 1} & \textbf{Value 2} & \textbf{Total} \\
\midrule
Item A & 10.5 & 20.3 & 30.8 \\
Item B & 15.2 & 25.1 & 40.3 \\
Item C & 22.1 & 12.7 & 34.8 \\
\bottomrule
\end{tabular}
\caption{Sample Data Table}
\label{tab:sample}
\end{table}

\section{Lists}
\subsection{Bullet Points}
\begin{itemize}
  \item First item
  \item Second item
  \item Third item with sub-items:
    \begin{itemize}
      \item Sub-item 1
      \item Sub-item 2
    \end{itemize}
\end{itemize}

\subsection{Numbered List}
\begin{enumerate}
  \item First step
  \item Second step
  \item Third step
\end{enumerate}

\section{Citations and References}
When you use the citation search feature, references will be automatically managed. For example, you might cite a paper like this: \cite{example2024}. All citations will be automatically added to your \texttt{cite.bib} file.

\section{Conclusion}
This template demonstrates basic \LaTeX{} features including mathematical formulas, tables, and lists. The bibliography system is now set up to work seamlessly with Bibby AI's citation management. Feel free to modify it for your needs!

\bibliographystyle{plain}
\bibliography{cite}

\end{document}\documentclass{article}
\usepackage[utf8]{inputenc}
\usepackage{booktabs}
\usepackage{amsmath}
\usepackage{natbib}

\title{My Document}
\author{Your Name}
\date{\today}

\begin{document}

\maketitle

\section{Introduction}
This is a sample document demonstrating various \LaTeX{} features including mathematical equations, tables, and formatting. You can cite papers using the citation search feature, and they will be automatically added to your \texttt{cite.bib} file.

\section{Mathematical Expressions}
Here are some example equations:

\subsection{Quadratic Formula}
The quadratic formula is given by:
\[x = \frac{-b \pm \sqrt{b^2 - 4ac}}{2a}\]

\subsection{Matrix Example}
A sample matrix:
\[
\begin{pmatrix}
a_{11} & a_{12} & a_{13} \\
a_{21} & a_{22} & a_{23} \\
a_{31} & a_{32} & a_{33}
\end{pmatrix}
\]

\section{Tables}
Below is a sample table showing data:

\begin{table}[h]
\centering
\begin{tabular}{@{}lrrr@{}}
\toprule
\textbf{Category} & \textbf{Value 1} & \textbf{Value 2} & \textbf{Total} \\
\midrule
Item A & 10.5 & 20.3 & 30.8 \\
Item B & 15.2 & 25.1 & 40.3 \\
Item C & 22.1 & 12.7 & 34.8 \\
\bottomrule
\end{tabular}
\caption{Sample Data Table}
\label{tab:sample}
\end{table}

\section{Lists}
\subsection{Bullet Points}
\begin{itemize}
  \item First item
  \item Second item
  \item Third item with sub-items:
    \begin{itemize}
      \item Sub-item 1
      \item Sub-item 2
    \end{itemize}
\end{itemize}

\subsection{Numbered List}
\begin{enumerate}
  \item First step
  \item Second step
  \item Third step
\end{enumerate}

\section{Citations and References}
When you use the citation search feature, references will be automatically managed. For example, you might cite a paper like this: \cite{example2024}. All citations will be automatically added to your \texttt{cite.bib} file.

\section{Conclusion}
This template demonstrates basic \LaTeX{} features including mathematical formulas, tables, and lists. The bibliography system is now set up to work seamlessly with Bibby AI's citation management. Feel free to modify it for your needs!

\bibliographystyle{plain}
\bibliography{cite}

\end{document}\documentclass{article}
\usepackage[utf8]{inputenc}
\usepackage{booktabs}
\usepackage{amsmath}
\usepackage{natbib}

\title{My Document}
\author{Your Name}
\date{\today}

\begin{document}

\maketitle

\section{Introduction}
This is a sample document demonstrating various \LaTeX{} features including mathematical equations, tables, and formatting. You can cite papers using the citation search feature, and they will be automatically added to your \texttt{cite.bib} file.

\section{Mathematical Expressions}
Here are some example equations:

\subsection{Quadratic Formula}
The quadratic formula is given by:
\[x = \frac{-b \pm \sqrt{b^2 - 4ac}}{2a}\]

\subsection{Matrix Example}
A sample matrix:
\[
\begin{pmatrix}
a_{11} & a_{12} & a_{13} \\
a_{21} & a_{22} & a_{23} \\
a_{31} & a_{32} & a_{33}
\end{pmatrix}
\]

\section{Tables}
Below is a sample table showing data:

\begin{table}[h]
\centering
\begin{tabular}{@{}lrrr@{}}
\toprule
\textbf{Category} & \textbf{Value 1} & \textbf{Value 2} & \textbf{Total} \\
\midrule
Item A & 10.5 & 20.3 & 30.8 \\
Item B & 15.2 & 25.1 & 40.3 \\
Item C & 22.1 & 12.7 & 34.8 \\
\bottomrule
\end{tabular}
\caption{Sample Data Table}
\label{tab:sample}
\end{table}

\section{Lists}
\subsection{Bullet Points}
\begin{itemize}
  \item First item
  \item Second item
  \item Third item with sub-items:
    \begin{itemize}
      \item Sub-item 1
      \item Sub-item 2
    \end{itemize}
\end{itemize}

\subsection{Numbered List}
\begin{enumerate}
  \item First step
  \item Second step
  \item Third step
\end{enumerate}

\section{Citations and References}
When you use the citation search feature, references will be automatically managed. For example, you might cite a paper like this: \cite{example2024}. All citations will be automatically added to your \texttt{cite.bib} file.

\section{Conclusion}
This template demonstrates basic \LaTeX{} features including mathematical formulas, tables, and lists. The bibliography system is now set up to work seamlessly with Bibby AI's citation management. Feel free to modify it for your needs!

\bibliographystyle{plain}
\bibliography{cite}

\end{document}\documentclass{article}
\usepackage[utf8]{inputenc}
\usepackage{booktabs}
\usepackage{amsmath}
\usepackage{natbib}

\title{My Document}
\author{Your Name}
\date{\today}

\begin{document}

\maketitle

\section{Introduction}
This is a sample document demonstrating various \LaTeX{} features including mathematical equations, tables, and formatting. You can cite papers using the citation search feature, and they will be automatically added to your \texttt{cite.bib} file.

\section{Mathematical Expressions}
Here are some example equations:

\subsection{Quadratic Formula}
The quadratic formula is given by:
\[x = \frac{-b \pm \sqrt{b^2 - 4ac}}{2a}\]

\subsection{Matrix Example}
A sample matrix:
\[
\begin{pmatrix}
a_{11} & a_{12} & a_{13} \\
a_{21} & a_{22} & a_{23} \\
a_{31} & a_{32} & a_{33}
\end{pmatrix}
\]

\section{Tables}
Below is a sample table showing data:

\begin{table}[h]
\centering
\begin{tabular}{@{}lrrr@{}}
\toprule
\textbf{Category} & \textbf{Value 1} & \textbf{Value 2} & \textbf{Total} \\
\midrule
Item A & 10.5 & 20.3 & 30.8 \\
Item B & 15.2 & 25.1 & 40.3 \\
Item C & 22.1 & 12.7 & 34.8 \\
\bottomrule
\end{tabular}
\caption{Sample Data Table}
\label{tab:sample}
\end{table}

\section{Lists}
\subsection{Bullet Points}
\begin{itemize}
  \item First item
  \item Second item
  \item Third item with sub-items:
    \begin{itemize}
      \item Sub-item 1
      \item Sub-item 2
    \end{itemize}
\end{itemize}

\subsection{Numbered List}
\begin{enumerate}
  \item First step
  \item Second step
  \item Third step
\end{enumerate}

\section{Citations and References}
When you use the citation search feature, references will be automatically managed. For example, you might cite a paper like this: \cite{example2024}. All citations will be automatically added to your \texttt{cite.bib} file.

\section{Conclusion}
This template demonstrates basic \LaTeX{} features including mathematical formulas, tables, and lists. The bibliography system is now set up to work seamlessly with Bibby AI's citation management. Feel free to modify it for your needs!

\bibliographystyle{plain}
\bibliography{cite}

\end{document}\documentclass{article}
\usepackage[utf8]{inputenc}
\usepackage{booktabs}
\usepackage{amsmath}
\usepackage{natbib}

\title{My Document}
\author{Your Name}
\date{\today}

\begin{document}

\maketitle

\section{Introduction}
This is a sample document demonstrating various \LaTeX{} features including mathematical equations, tables, and formatting. You can cite papers using the citation search feature, and they will be automatically added to your \texttt{cite.bib} file.

\section{Mathematical Expressions}
Here are some example equations:

\subsection{Quadratic Formula}
The quadratic formula is given by:
\[x = \frac{-b \pm \sqrt{b^2 - 4ac}}{2a}\]

\subsection{Matrix Example}
A sample matrix:
\[
\begin{pmatrix}
a_{11} & a_{12} & a_{13} \\
a_{21} & a_{22} & a_{23} \\
a_{31} & a_{32} & a_{33}
\end{pmatrix}
\]

\section{Tables}
Below is a sample table showing data:

\begin{table}[h]
\centering
\begin{tabular}{@{}lrrr@{}}
\toprule
\textbf{Category} & \textbf{Value 1} & \textbf{Value 2} & \textbf{Total} \\
\midrule
Item A & 10.5 & 20.3 & 30.8 \\
Item B & 15.2 & 25.1 & 40.3 \\
Item C & 22.1 & 12.7 & 34.8 \\
\bottomrule
\end{tabular}
\caption{Sample Data Table}
\label{tab:sample}
\end{table}

\section{Lists}
\subsection{Bullet Points}
\begin{itemize}
  \item First item
  \item Second item
  \item Third item with sub-items:
    \begin{itemize}
      \item Sub-item 1
      \item Sub-item 2
    \end{itemize}
\end{itemize}

\subsection{Numbered List}
\begin{enumerate}
  \item First step
  \item Second step
  \item Third step
\end{enumerate}

\section{Citations and References}
When you use the citation search feature, references will be automatically managed. For example, you might cite a paper like this: \cite{example2024}. All citations will be automatically added to your \texttt{cite.bib} file.

\section{Conclusion}
This template demonstrates basic \LaTeX{} features including mathematical formulas, tables, and lists. The bibliography system is now set up to work seamlessly with Bibby AI's citation management. Feel free to modify it for your needs!

\bibliographystyle{plain}
\bibliography{cite}

\end{document}\documentclass{article}
\usepackage[utf8]{inputenc}
\usepackage{booktabs}
\usepackage{amsmath}
\usepackage{natbib}

\title{My Document}
\author{Your Name}
\date{\today}

\begin{document}

\maketitle

\section{Introduction}
This is a sample document demonstrating various \LaTeX{} features including mathematical equations, tables, and formatting. You can cite papers using the citation search feature, and they will be automatically added to your \texttt{cite.bib} file.

\section{Mathematical Expressions}
Here are some example equations:

\subsection{Quadratic Formula}
The quadratic formula is given by:
\[x = \frac{-b \pm \sqrt{b^2 - 4ac}}{2a}\]

\subsection{Matrix Example}
A sample matrix:
\[
\begin{pmatrix}
a_{11} & a_{12} & a_{13} \\
a_{21} & a_{22} & a_{23} \\
a_{31} & a_{32} & a_{33}
\end{pmatrix}
\]

\section{Tables}
Below is a sample table showing data:

\begin{table}[h]
\centering
\begin{tabular}{@{}lrrr@{}}
\toprule
\textbf{Category} & \textbf{Value 1} & \textbf{Value 2} & \textbf{Total} \\
\midrule
Item A & 10.5 & 20.3 & 30.8 \\
Item B & 15.2 & 25.1 & 40.3 \\
Item C & 22.1 & 12.7 & 34.8 \\
\bottomrule
\end{tabular}
\caption{Sample Data Table}
\label{tab:sample}
\end{table}

\section{Lists}
\subsection{Bullet Points}
\begin{itemize}
  \item First item
  \item Second item
  \item Third item with sub-items:
    \begin{itemize}
      \item Sub-item 1
      \item Sub-item 2
    \end{itemize}
\end{itemize}

\subsection{Numbered List}
\begin{enumerate}
  \item First step
  \item Second step
  \item Third step
\end{enumerate}

\section{Citations and References}
When you use the citation search feature, references will be automatically managed. For example, you might cite a paper like this: \cite{example2024}. All citations will be automatically added to your \texttt{cite.bib} file.

\section{Conclusion}
This template demonstrates basic \LaTeX{} features including mathematical formulas, tables, and lists. The bibliography system is now set up to work seamlessly with Bibby AI's citation management. Feel free to modify it for your needs!

\bibliographystyle{plain}
\bibliography{cite}

\end{document}\documentclass{article}
\usepackage[utf8]{inputenc}
\usepackage{booktabs}
\usepackage{amsmath}
\usepackage{natbib}

\title{My Document}
\author{Your Name}
\date{\today}

\begin{document}

\maketitle

\section{Introduction}
This is a sample document demonstrating various \LaTeX{} features including mathematical equations, tables, and formatting. You can cite papers using the citation search feature, and they will be automatically added to your \texttt{cite.bib} file.

\section{Mathematical Expressions}
Here are some example equations:

\subsection{Quadratic Formula}
The quadratic formula is given by:
\[x = \frac{-b \pm \sqrt{b^2 - 4ac}}{2a}\]

\subsection{Matrix Example}
A sample matrix:
\[
\begin{pmatrix}
a_{11} & a_{12} & a_{13} \\
a_{21} & a_{22} & a_{23} \\
a_{31} & a_{32} & a_{33}
\end{pmatrix}
\]

\section{Tables}
Below is a sample table showing data:

\begin{table}[h]
\centering
\begin{tabular}{@{}lrrr@{}}
\toprule
\textbf{Category} & \textbf{Value 1} & \textbf{Value 2} & \textbf{Total} \\
\midrule
Item A & 10.5 & 20.3 & 30.8 \\
Item B & 15.2 & 25.1 & 40.3 \\
Item C & 22.1 & 12.7 & 34.8 \\
\bottomrule
\end{tabular}
\caption{Sample Data Table}
\label{tab:sample}
\end{table}

\section{Lists}
\subsection{Bullet Points}
\begin{itemize}
  \item First item
  \item Second item
  \item Third item with sub-items:
    \begin{itemize}
      \item Sub-item 1
      \item Sub-item 2
    \end{itemize}
\end{itemize}

\subsection{Numbered List}
\begin{enumerate}
  \item First step
  \item Second step
  \item Third step
\end{enumerate}

\section{Citations and References}
When you use the citation search feature, references will be automatically managed. For example, you might cite a paper like this: \cite{example2024}. All citations will be automatically added to your \texttt{cite.bib} file.

\section{Conclusion}
This template demonstrates basic \LaTeX{} features including mathematical formulas, tables, and lists. The bibliography system is now set up to work seamlessly with Bibby AI's citation management. Feel free to modify it for your needs!

\bibliographystyle{plain}
\bibliography{cite}

\end{document}\documentclass{article}
\usepackage[utf8]{inputenc}
\usepackage{booktabs}
\usepackage{amsmath}
\usepackage{natbib}

\title{My Document}
\author{Your Name}
\date{\today}

\begin{document}

\maketitle

\section{Introduction}
This is a sample document demonstrating various \LaTeX{} features including mathematical equations, tables, and formatting. You can cite papers using the citation search feature, and they will be automatically added to your \texttt{cite.bib} file.

\section{Mathematical Expressions}
Here are some example equations:

\subsection{Quadratic Formula}
The quadratic formula is given by:
\[x = \frac{-b \pm \sqrt{b^2 - 4ac}}{2a}\]

\subsection{Matrix Example}
A sample matrix:
\[
\begin{pmatrix}
a_{11} & a_{12} & a_{13} \\
a_{21} & a_{22} & a_{23} \\
a_{31} & a_{32} & a_{33}
\end{pmatrix}
\]

\section{Tables}
Below is a sample table showing data:

\begin{table}[h]
\centering
\begin{tabular}{@{}lrrr@{}}
\toprule
\textbf{Category} & \textbf{Value 1} & \textbf{Value 2} & \textbf{Total} \\
\midrule
Item A & 10.5 & 20.3 & 30.8 \\
Item B & 15.2 & 25.1 & 40.3 \\
Item C & 22.1 & 12.7 & 34.8 \\
\bottomrule
\end{tabular}
\caption{Sample Data Table}
\label{tab:sample}
\end{table}

\section{Lists}
\subsection{Bullet Points}
\begin{itemize}
  \item First item
  \item Second item
  \item Third item with sub-items:
    \begin{itemize}
      \item Sub-item 1
      \item Sub-item 2
    \end{itemize}
\end{itemize}

\subsection{Numbered List}
\begin{enumerate}
  \item First step
  \item Second step
  \item Third step
\end{enumerate}

\section{Citations and References}
When you use the citation search feature, references will be automatically managed. For example, you might cite a paper like this: \cite{example2024}. All citations will be automatically added to your \texttt{cite.bib} file.

\section{Conclusion}
This template demonstrates basic \LaTeX{} features including mathematical formulas, tables, and lists. The bibliography system is now set up to work seamlessly with Bibby AI's citation management. Feel free to modify it for your needs!

\bibliographystyle{plain}
\bibliography{cite}

\end{document}\documentclass{article}
\usepackage[utf8]{inputenc}
\usepackage{booktabs}
\usepackage{amsmath}
\usepackage{natbib}

\title{My Document}
\author{Your Name}
\date{\today}

\begin{document}

\maketitle

\section{Introduction}
This is a sample document demonstrating various \LaTeX{} features including mathematical equations, tables, and formatting. You can cite papers using the citation search feature, and they will be automatically added to your \texttt{cite.bib} file.

\section{Mathematical Expressions}
Here are some example equations:

\subsection{Quadratic Formula}
The quadratic formula is given by:
\[x = \frac{-b \pm \sqrt{b^2 - 4ac}}{2a}\]

\subsection{Matrix Example}
A sample matrix:
\[
\begin{pmatrix}
a_{11} & a_{12} & a_{13} \\
a_{21} & a_{22} & a_{23} \\
a_{31} & a_{32} & a_{33}
\end{pmatrix}
\]

\section{Tables}
Below is a sample table showing data:

\begin{table}[h]
\centering
\begin{tabular}{@{}lrrr@{}}
\toprule
\textbf{Category} & \textbf{Value 1} & \textbf{Value 2} & \textbf{Total} \\
\midrule
Item A & 10.5 & 20.3 & 30.8 \\
Item B & 15.2 & 25.1 & 40.3 \\
Item C & 22.1 & 12.7 & 34.8 \\
\bottomrule
\end{tabular}
\caption{Sample Data Table}
\label{tab:sample}
\end{table}

\section{Lists}
\subsection{Bullet Points}
\begin{itemize}
  \item First item
  \item Second item
  \item Third item with sub-items:
    \begin{itemize}
      \item Sub-item 1
      \item Sub-item 2
    \end{itemize}
\end{itemize}

\subsection{Numbered List}
\begin{enumerate}
  \item First step
  \item Second step
  \item Third step
\end{enumerate}

\section{Citations and References}
When you use the citation search feature, references will be automatically managed. For example, you might cite a paper like this: \cite{example2024}. All citations will be automatically added to your \texttt{cite.bib} file.

\section{Conclusion}
This template demonstrates basic \LaTeX{} features including mathematical formulas, tables, and lists. The bibliography system is now set up to work seamlessly with Bibby AI's citation management. Feel free to modify it for your needs!

\bibliographystyle{plain}
\bibliography{cite}

\end{document}\documentclass{article}
\usepackage[utf8]{inputenc}
\usepackage{booktabs}
\usepackage{amsmath}
\usepackage{natbib}

\title{My Document}
\author{Your Name}
\date{\today}

\begin{document}

\maketitle

\section{Introduction}
This is a sample document demonstrating various \LaTeX{} features including mathematical equations, tables, and formatting. You can cite papers using the citation search feature, and they will be automatically added to your \texttt{cite.bib} file.

\section{Mathematical Expressions}
Here are some example equations:

\subsection{Quadratic Formula}
The quadratic formula is given by:
\[x = \frac{-b \pm \sqrt{b^2 - 4ac}}{2a}\]

\subsection{Matrix Example}
A sample matrix:
\[
\begin{pmatrix}
a_{11} & a_{12} & a_{13} \\
a_{21} & a_{22} & a_{23} \\
a_{31} & a_{32} & a_{33}
\end{pmatrix}
\]

\section{Tables}
Below is a sample table showing data:

\begin{table}[h]
\centering
\begin{tabular}{@{}lrrr@{}}
\toprule
\textbf{Category} & \textbf{Value 1} & \textbf{Value 2} & \textbf{Total} \\
\midrule
Item A & 10.5 & 20.3 & 30.8 \\
Item B & 15.2 & 25.1 & 40.3 \\
Item C & 22.1 & 12.7 & 34.8 \\
\bottomrule
\end{tabular}
\caption{Sample Data Table}
\label{tab:sample}
\end{table}

\section{Lists}
\subsection{Bullet Points}
\begin{itemize}
  \item First item
  \item Second item
  \item Third item with sub-items:
    \begin{itemize}
      \item Sub-item 1
      \item Sub-item 2
    \end{itemize}
\end{itemize}

\subsection{Numbered List}
\begin{enumerate}
  \item First step
  \item Second step
  \item Third step
\end{enumerate}

\section{Citations and References}
When you use the citation search feature, references will be automatically managed. For example, you might cite a paper like this: \cite{example2024}. All citations will be automatically added to your \texttt{cite.bib} file.

\section{Conclusion}
This template demonstrates basic \LaTeX{} features including mathematical formulas, tables, and lists. The bibliography system is now set up to work seamlessly with Bibby AI's citation management. Feel free to modify it for your needs!

\bibliographystyle{plain}
\bibliography{cite}

\end{document}\documentclass{article}
\usepackage[utf8]{inputenc}
\usepackage{booktabs}
\usepackage{amsmath}
\usepackage{natbib}

\title{My Document}
\author{Your Name}
\date{\today}

\begin{document}

\maketitle

\section{Introduction}
This is a sample document demonstrating various \LaTeX{} features including mathematical equations, tables, and formatting. You can cite papers using the citation search feature, and they will be automatically added to your \texttt{cite.bib} file.

\section{Mathematical Expressions}
Here are some example equations:

\subsection{Quadratic Formula}
The quadratic formula is given by:
\[x = \frac{-b \pm \sqrt{b^2 - 4ac}}{2a}\]

\subsection{Matrix Example}
A sample matrix:
\[
\begin{pmatrix}
a_{11} & a_{12} & a_{13} \\
a_{21} & a_{22} & a_{23} \\
a_{31} & a_{32} & a_{33}
\end{pmatrix}
\]

\section{Tables}
Below is a sample table showing data:

\begin{table}[h]
\centering
\begin{tabular}{@{}lrrr@{}}
\toprule
\textbf{Category} & \textbf{Value 1} & \textbf{Value 2} & \textbf{Total} \\
\midrule
Item A & 10.5 & 20.3 & 30.8 \\
Item B & 15.2 & 25.1 & 40.3 \\
Item C & 22.1 & 12.7 & 34.8 \\
\bottomrule
\end{tabular}
\caption{Sample Data Table}
\label{tab:sample}
\end{table}

\section{Lists}
\subsection{Bullet Points}
\begin{itemize}
  \item First item
  \item Second item
  \item Third item with sub-items:
    \begin{itemize}
      \item Sub-item 1
      \item Sub-item 2
    \end{itemize}
\end{itemize}

\subsection{Numbered List}
\begin{enumerate}
  \item First step
  \item Second step
  \item Third step
\end{enumerate}

\section{Citations and References}
When you use the citation search feature, references will be automatically managed. For example, you might cite a paper like this: \cite{example2024}. All citations will be automatically added to your \texttt{cite.bib} file.

\section{Conclusion}
This template demonstrates basic \LaTeX{} features including mathematical formulas, tables, and lists. The bibliography system is now set up to work seamlessly with Bibby AI's citation management. Feel free to modify it for your needs!

\bibliographystyle{plain}
\bibliography{cite}

\end{document}\documentclass{article}
\usepackage[utf8]{inputenc}
\usepackage{booktabs}
\usepackage{amsmath}
\usepackage{natbib}

\title{My Document}
\author{Your Name}
\date{\today}

\begin{document}

\maketitle

\section{Introduction}
This is a sample document demonstrating various \LaTeX{} features including mathematical equations, tables, and formatting. You can cite papers using the citation search feature, and they will be automatically added to your \texttt{cite.bib} file.

\section{Mathematical Expressions}
Here are some example equations:

\subsection{Quadratic Formula}
The quadratic formula is given by:
\[x = \frac{-b \pm \sqrt{b^2 - 4ac}}{2a}\]

\subsection{Matrix Example}
A sample matrix:
\[
\begin{pmatrix}
a_{11} & a_{12} & a_{13} \\
a_{21} & a_{22} & a_{23} \\
a_{31} & a_{32} & a_{33}
\end{pmatrix}
\]

\section{Tables}
Below is a sample table showing data:

\begin{table}[h]
\centering
\begin{tabular}{@{}lrrr@{}}
\toprule
\textbf{Category} & \textbf{Value 1} & \textbf{Value 2} & \textbf{Total} \\
\midrule
Item A & 10.5 & 20.3 & 30.8 \\
Item B & 15.2 & 25.1 & 40.3 \\
Item C & 22.1 & 12.7 & 34.8 \\
\bottomrule
\end{tabular}
\caption{Sample Data Table}
\label{tab:sample}
\end{table}

\section{Lists}
\subsection{Bullet Points}
\begin{itemize}
  \item First item
  \item Second item
  \item Third item with sub-items:
    \begin{itemize}
      \item Sub-item 1
      \item Sub-item 2
    \end{itemize}
\end{itemize}

\subsection{Numbered List}
\begin{enumerate}
  \item First step
  \item Second step
  \item Third step
\end{enumerate}

\section{Citations and References}
When you use the citation search feature, references will be automatically managed. For example, you might cite a paper like this: \cite{example2024}. All citations will be automatically added to your \texttt{cite.bib} file.

\section{Conclusion}
This template demonstrates basic \LaTeX{} features including mathematical formulas, tables, and lists. The bibliography system is now set up to work seamlessly with Bibby AI's citation management. Feel free to modify it for your needs!

\bibliographystyle{plain}
\bibliography{cite}

\end{document}\documentclass{article}
\usepackage[utf8]{inputenc}
\usepackage{booktabs}
\usepackage{amsmath}
\usepackage{natbib}

\title{My Document}
\author{Your Name}
\date{\today}

\begin{document}

\maketitle

\section{Introduction}
This is a sample document demonstrating various \LaTeX{} features including mathematical equations, tables, and formatting. You can cite papers using the citation search feature, and they will be automatically added to your \texttt{cite.bib} file.

\section{Mathematical Expressions}
Here are some example equations:

\subsection{Quadratic Formula}
The quadratic formula is given by:
\[x = \frac{-b \pm \sqrt{b^2 - 4ac}}{2a}\]

\subsection{Matrix Example}
A sample matrix:
\[
\begin{pmatrix}
a_{11} & a_{12} & a_{13} \\
a_{21} & a_{22} & a_{23} \\
a_{31} & a_{32} & a_{33}
\end{pmatrix}
\]

\section{Tables}
Below is a sample table showing data:

\begin{table}[h]
\centering
\begin{tabular}{@{}lrrr@{}}
\toprule
\textbf{Category} & \textbf{Value 1} & \textbf{Value 2} & \textbf{Total} \\
\midrule
Item A & 10.5 & 20.3 & 30.8 \\
Item B & 15.2 & 25.1 & 40.3 \\
Item C & 22.1 & 12.7 & 34.8 \\
\bottomrule
\end{tabular}
\caption{Sample Data Table}
\label{tab:sample}
\end{table}

\section{Lists}
\subsection{Bullet Points}
\begin{itemize}
  \item First item
  \item Second item
  \item Third item with sub-items:
    \begin{itemize}
      \item Sub-item 1
      \item Sub-item 2
    \end{itemize}
\end{itemize}

\subsection{Numbered List}
\begin{enumerate}
  \item First step
  \item Second step
  \item Third step
\end{enumerate}

\section{Citations and References}
When you use the citation search feature, references will be automatically managed. For example, you might cite a paper like this: \cite{example2024}. All citations will be automatically added to your \texttt{cite.bib} file.

\section{Conclusion}
This template demonstrates basic \LaTeX{} features including mathematical formulas, tables, and lists. The bibliography system is now set up to work seamlessly with Bibby AI's citation management. Feel free to modify it for your needs!

\bibliographystyle{plain}
\bibliography{cite}

\end{document}\documentclass{article}
\usepackage[utf8]{inputenc}
\usepackage{booktabs}
\usepackage{amsmath}
\usepackage{natbib}

\title{My Document}
\author{Your Name}
\date{\today}

\begin{document}

\maketitle

\section{Introduction}
This is a sample document demonstrating various \LaTeX{} features including mathematical equations, tables, and formatting. You can cite papers using the citation search feature, and they will be automatically added to your \texttt{cite.bib} file.

\section{Mathematical Expressions}
Here are some example equations:

\subsection{Quadratic Formula}
The quadratic formula is given by:
\[x = \frac{-b \pm \sqrt{b^2 - 4ac}}{2a}\]

\subsection{Matrix Example}
A sample matrix:
\[
\begin{pmatrix}
a_{11} & a_{12} & a_{13} \\
a_{21} & a_{22} & a_{23} \\
a_{31} & a_{32} & a_{33}
\end{pmatrix}
\]

\section{Tables}
Below is a sample table showing data:

\begin{table}[h]
\centering
\begin{tabular}{@{}lrrr@{}}
\toprule
\textbf{Category} & \textbf{Value 1} & \textbf{Value 2} & \textbf{Total} \\
\midrule
Item A & 10.5 & 20.3 & 30.8 \\
Item B & 15.2 & 25.1 & 40.3 \\
Item C & 22.1 & 12.7 & 34.8 \\
\bottomrule
\end{tabular}
\caption{Sample Data Table}
\label{tab:sample}
\end{table}

\section{Lists}
\subsection{Bullet Points}
\begin{itemize}
  \item First item
  \item Second item
  \item Third item with sub-items:
    \begin{itemize}
      \item Sub-item 1
      \item Sub-item 2
    \end{itemize}
\end{itemize}

\subsection{Numbered List}
\begin{enumerate}
  \item First step
  \item Second step
  \item Third step
\end{enumerate}

\section{Citations and References}
When you use the citation search feature, references will be automatically managed. For example, you might cite a paper like this: \cite{example2024}. All citations will be automatically added to your \texttt{cite.bib} file.

\section{Conclusion}
This template demonstrates basic \LaTeX{} features including mathematical formulas, tables, and lists. The bibliography system is now set up to work seamlessly with Bibby AI's citation management. Feel free to modify it for your needs!

\bibliographystyle{plain}
\bibliography{cite}

\end{document}\documentclass{article}
\usepackage[utf8]{inputenc}
\usepackage{booktabs}
\usepackage{amsmath}
\usepackage{natbib}

\title{My Document}
\author{Your Name}
\date{\today}

\begin{document}

\maketitle

\section{Introduction}
This is a sample document demonstrating various \LaTeX{} features including mathematical equations, tables, and formatting. You can cite papers using the citation search feature, and they will be automatically added to your \texttt{cite.bib} file.

\section{Mathematical Expressions}
Here are some example equations:

\subsection{Quadratic Formula}
The quadratic formula is given by:
\[x = \frac{-b \pm \sqrt{b^2 - 4ac}}{2a}\]

\subsection{Matrix Example}
A sample matrix:
\[
\begin{pmatrix}
a_{11} & a_{12} & a_{13} \\
a_{21} & a_{22} & a_{23} \\
a_{31} & a_{32} & a_{33}
\end{pmatrix}
\]

\section{Tables}
Below is a sample table showing data:

\begin{table}[h]
\centering
\begin{tabular}{@{}lrrr@{}}
\toprule
\textbf{Category} & \textbf{Value 1} & \textbf{Value 2} & \textbf{Total} \\
\midrule
Item A & 10.5 & 20.3 & 30.8 \\
Item B & 15.2 & 25.1 & 40.3 \\
Item C & 22.1 & 12.7 & 34.8 \\
\bottomrule
\end{tabular}
\caption{Sample Data Table}
\label{tab:sample}
\end{table}

\section{Lists}
\subsection{Bullet Points}
\begin{itemize}
  \item First item
  \item Second item
  \item Third item with sub-items:
    \begin{itemize}
      \item Sub-item 1
      \item Sub-item 2
    \end{itemize}
\end{itemize}

\subsection{Numbered List}
\begin{enumerate}
  \item First step
  \item Second step
  \item Third step
\end{enumerate}

\section{Citations and References}
When you use the citation search feature, references will be automatically managed. For example, you might cite a paper like this: \cite{example2024}. All citations will be automatically added to your \texttt{cite.bib} file.

\section{Conclusion}
This template demonstrates basic \LaTeX{} features including mathematical formulas, tables, and lists. The bibliography system is now set up to work seamlessly with Bibby AI's citation management. Feel free to modify it for your needs!

\bibliographystyle{plain}
\bibliography{cite}

\end{document}\documentclass{article}
\usepackage[utf8]{inputenc}
\usepackage{booktabs}
\usepackage{amsmath}
\usepackage{natbib}

\title{My Document}
\author{Your Name}
\date{\today}

\begin{document}

\maketitle

\section{Introduction}
This is a sample document demonstrating various \LaTeX{} features including mathematical equations, tables, and formatting. You can cite papers using the citation search feature, and they will be automatically added to your \texttt{cite.bib} file.

\section{Mathematical Expressions}
Here are some example equations:

\subsection{Quadratic Formula}
The quadratic formula is given by:
\[x = \frac{-b \pm \sqrt{b^2 - 4ac}}{2a}\]

\subsection{Matrix Example}
A sample matrix:
\[
\begin{pmatrix}
a_{11} & a_{12} & a_{13} \\
a_{21} & a_{22} & a_{23} \\
a_{31} & a_{32} & a_{33}
\end{pmatrix}
\]

\section{Tables}
Below is a sample table showing data:

\begin{table}[h]
\centering
\begin{tabular}{@{}lrrr@{}}
\toprule
\textbf{Category} & \textbf{Value 1} & \textbf{Value 2} & \textbf{Total} \\
\midrule
Item A & 10.5 & 20.3 & 30.8 \\
Item B & 15.2 & 25.1 & 40.3 \\
Item C & 22.1 & 12.7 & 34.8 \\
\bottomrule
\end{tabular}
\caption{Sample Data Table}
\label{tab:sample}
\end{table}

\section{Lists}
\subsection{Bullet Points}
\begin{itemize}
  \item First item
  \item Second item
  \item Third item with sub-items:
    \begin{itemize}
      \item Sub-item 1
      \item Sub-item 2
    \end{itemize}
\end{itemize}

\subsection{Numbered List}
\begin{enumerate}
  \item First step
  \item Second step
  \item Third step
\end{enumerate}

\section{Citations and References}
When you use the citation search feature, references will be automatically managed. For example, you might cite a paper like this: \cite{example2024}. All citations will be automatically added to your \texttt{cite.bib} file.

\section{Conclusion}
This template demonstrates basic \LaTeX{} features including mathematical formulas, tables, and lists. The bibliography system is now set up to work seamlessly with Bibby AI's citation management. Feel free to modify it for your needs!

\bibliographystyle{plain}
\bibliography{cite}

\end{document}\documentclass{article}
\usepackage[utf8]{inputenc}
\usepackage{booktabs}
\usepackage{amsmath}
\usepackage{natbib}

\title{My Document}
\author{Your Name}
\date{\today}

\begin{document}

\maketitle

\section{Introduction}
This is a sample document demonstrating various \LaTeX{} features including mathematical equations, tables, and formatting. You can cite papers using the citation search feature, and they will be automatically added to your \texttt{cite.bib} file.

\section{Mathematical Expressions}
Here are some example equations:

\subsection{Quadratic Formula}
The quadratic formula is given by:
\[x = \frac{-b \pm \sqrt{b^2 - 4ac}}{2a}\]

\subsection{Matrix Example}
A sample matrix:
\[
\begin{pmatrix}
a_{11} & a_{12} & a_{13} \\
a_{21} & a_{22} & a_{23} \\
a_{31} & a_{32} & a_{33}
\end{pmatrix}
\]

\section{Tables}
Below is a sample table showing data:

\begin{table}[h]
\centering
\begin{tabular}{@{}lrrr@{}}
\toprule
\textbf{Category} & \textbf{Value 1} & \textbf{Value 2} & \textbf{Total} \\
\midrule
Item A & 10.5 & 20.3 & 30.8 \\
Item B & 15.2 & 25.1 & 40.3 \\
Item C & 22.1 & 12.7 & 34.8 \\
\bottomrule
\end{tabular}
\caption{Sample Data Table}
\label{tab:sample}
\end{table}

\section{Lists}
\subsection{Bullet Points}
\begin{itemize}
  \item First item
  \item Second item
  \item Third item with sub-items:
    \begin{itemize}
      \item Sub-item 1
      \item Sub-item 2
    \end{itemize}
\end{itemize}

\subsection{Numbered List}
\begin{enumerate}
  \item First step
  \item Second step
  \item Third step
\end{enumerate}

\section{Citations and References}
When you use the citation search feature, references will be automatically managed. For example, you might cite a paper like this: \cite{example2024}. All citations will be automatically added to your \texttt{cite.bib} file.

\section{Conclusion}
This template demonstrates basic \LaTeX{} features including mathematical formulas, tables, and lists. The bibliography system is now set up to work seamlessly with Bibby AI's citation management. Feel free to modify it for your needs!

\bibliographystyle{plain}
\bibliography{cite}

\end{document}\documentclass{article}
\usepackage[utf8]{inputenc}
\usepackage{booktabs}
\usepackage{amsmath}
\usepackage{natbib}

\title{My Document}
\author{Your Name}
\date{\today}

\begin{document}

\maketitle

\section{Introduction}
This is a sample document demonstrating various \LaTeX{} features including mathematical equations, tables, and formatting. You can cite papers using the citation search feature, and they will be automatically added to your \texttt{cite.bib} file.

\section{Mathematical Expressions}
Here are some example equations:

\subsection{Quadratic Formula}
The quadratic formula is given by:
\[x = \frac{-b \pm \sqrt{b^2 - 4ac}}{2a}\]

\subsection{Matrix Example}
A sample matrix:
\[
\begin{pmatrix}
a_{11} & a_{12} & a_{13} \\
a_{21} & a_{22} & a_{23} \\
a_{31} & a_{32} & a_{33}
\end{pmatrix}
\]

\section{Tables}
Below is a sample table showing data:

\begin{table}[h]
\centering
\begin{tabular}{@{}lrrr@{}}
\toprule
\textbf{Category} & \textbf{Value 1} & \textbf{Value 2} & \textbf{Total} \\
\midrule
Item A & 10.5 & 20.3 & 30.8 \\
Item B & 15.2 & 25.1 & 40.3 \\
Item C & 22.1 & 12.7 & 34.8 \\
\bottomrule
\end{tabular}
\caption{Sample Data Table}
\label{tab:sample}
\end{table}

\section{Lists}
\subsection{Bullet Points}
\begin{itemize}
  \item First item
  \item Second item
  \item Third item with sub-items:
    \begin{itemize}
      \item Sub-item 1
      \item Sub-item 2
    \end{itemize}
\end{itemize}

\subsection{Numbered List}
\begin{enumerate}
  \item First step
  \item Second step
  \item Third step
\end{enumerate}

\section{Citations and References}
When you use the citation search feature, references will be automatically managed. For example, you might cite a paper like this: \cite{example2024}. All citations will be automatically added to your \texttt{cite.bib} file.

\section{Conclusion}
This template demonstrates basic \LaTeX{} features including mathematical formulas, tables, and lists. The bibliography system is now set up to work seamlessly with Bibby AI's citation management. Feel free to modify it for your needs!

\bibliographystyle{plain}
\bibliography{cite}

\end{document}\documentclass{article}
\usepackage[utf8]{inputenc}
\usepackage{booktabs}
\usepackage{amsmath}
\usepackage{natbib}

\title{My Document}
\author{Your Name}
\date{\today}

\begin{document}

\maketitle

\section{Introduction}
This is a sample document demonstrating various \LaTeX{} features including mathematical equations, tables, and formatting. You can cite papers using the citation search feature, and they will be automatically added to your \texttt{cite.bib} file.

\section{Mathematical Expressions}
Here are some example equations:

\subsection{Quadratic Formula}
The quadratic formula is given by:
\[x = \frac{-b \pm \sqrt{b^2 - 4ac}}{2a}\]

\subsection{Matrix Example}
A sample matrix:
\[
\begin{pmatrix}
a_{11} & a_{12} & a_{13} \\
a_{21} & a_{22} & a_{23} \\
a_{31} & a_{32} & a_{33}
\end{pmatrix}
\]

\section{Tables}
Below is a sample table showing data:

\begin{table}[h]
\centering
\begin{tabular}{@{}lrrr@{}}
\toprule
\textbf{Category} & \textbf{Value 1} & \textbf{Value 2} & \textbf{Total} \\
\midrule
Item A & 10.5 & 20.3 & 30.8 \\
Item B & 15.2 & 25.1 & 40.3 \\
Item C & 22.1 & 12.7 & 34.8 \\
\bottomrule
\end{tabular}
\caption{Sample Data Table}
\label{tab:sample}
\end{table}

\section{Lists}
\subsection{Bullet Points}
\begin{itemize}
  \item First item
  \item Second item
  \item Third item with sub-items:
    \begin{itemize}
      \item Sub-item 1
      \item Sub-item 2
    \end{itemize}
\end{itemize}

\subsection{Numbered List}
\begin{enumerate}
  \item First step
  \item Second step
  \item Third step
\end{enumerate}

\section{Citations and References}
When you use the citation search feature, references will be automatically managed. For example, you might cite a paper like this: \cite{example2024}. All citations will be automatically added to your \texttt{cite.bib} file.

\section{Conclusion}
This template demonstrates basic \LaTeX{} features including mathematical formulas, tables, and lists. The bibliography system is now set up to work seamlessly with Bibby AI's citation management. Feel free to modify it for your needs!

\bibliographystyle{plain}
\bibliography{cite}

\end{document}\documentclass{article}
\usepackage[utf8]{inputenc}
\usepackage{booktabs}
\usepackage{amsmath}
\usepackage{natbib}

\title{My Document}
\author{Your Name}
\date{\today}

\begin{document}

\maketitle

\section{Introduction}
This is a sample document demonstrating various \LaTeX{} features including mathematical equations, tables, and formatting. You can cite papers using the citation search feature, and they will be automatically added to your \texttt{cite.bib} file.

\section{Mathematical Expressions}
Here are some example equations:

\subsection{Quadratic Formula}
The quadratic formula is given by:
\[x = \frac{-b \pm \sqrt{b^2 - 4ac}}{2a}\]

\subsection{Matrix Example}
A sample matrix:
\[
\begin{pmatrix}
a_{11} & a_{12} & a_{13} \\
a_{21} & a_{22} & a_{23} \\
a_{31} & a_{32} & a_{33}
\end{pmatrix}
\]

\section{Tables}
Below is a sample table showing data:

\begin{table}[h]
\centering
\begin{tabular}{@{}lrrr@{}}
\toprule
\textbf{Category} & \textbf{Value 1} & \textbf{Value 2} & \textbf{Total} \\
\midrule
Item A & 10.5 & 20.3 & 30.8 \\
Item B & 15.2 & 25.1 & 40.3 \\
Item C & 22.1 & 12.7 & 34.8 \\
\bottomrule
\end{tabular}
\caption{Sample Data Table}
\label{tab:sample}
\end{table}

\section{Lists}
\subsection{Bullet Points}
\begin{itemize}
  \item First item
  \item Second item
  \item Third item with sub-items:
    \begin{itemize}
      \item Sub-item 1
      \item Sub-item 2
    \end{itemize}
\end{itemize}

\subsection{Numbered List}
\begin{enumerate}
  \item First step
  \item Second step
  \item Third step
\end{enumerate}

\section{Citations and References}
When you use the citation search feature, references will be automatically managed. For example, you might cite a paper like this: \cite{example2024}. All citations will be automatically added to your \texttt{cite.bib} file.

\section{Conclusion}
This template demonstrates basic \LaTeX{} features including mathematical formulas, tables, and lists. The bibliography system is now set up to work seamlessly with Bibby AI's citation management. Feel free to modify it for your needs!

\bibliographystyle{plain}
\bibliography{cite}

\end{document}\documentclass{article}
\usepackage[utf8]{inputenc}
\usepackage{booktabs}
\usepackage{amsmath}
\usepackage{natbib}

\title{My Document}
\author{Your Name}
\date{\today}

\begin{document}

\maketitle

\section{Introduction}
This is a sample document demonstrating various \LaTeX{} features including mathematical equations, tables, and formatting. You can cite papers using the citation search feature, and they will be automatically added to your \texttt{cite.bib} file.

\section{Mathematical Expressions}
Here are some example equations:

\subsection{Quadratic Formula}
The quadratic formula is given by:
\[x = \frac{-b \pm \sqrt{b^2 - 4ac}}{2a}\]

\subsection{Matrix Example}
A sample matrix:
\[
\begin{pmatrix}
a_{11} & a_{12} & a_{13} \\
a_{21} & a_{22} & a_{23} \\
a_{31} & a_{32} & a_{33}
\end{pmatrix}
\]

\section{Tables}
Below is a sample table showing data:

\begin{table}[h]
\centering
\begin{tabular}{@{}lrrr@{}}
\toprule
\textbf{Category} & \textbf{Value 1} & \textbf{Value 2} & \textbf{Total} \\
\midrule
Item A & 10.5 & 20.3 & 30.8 \\
Item B & 15.2 & 25.1 & 40.3 \\
Item C & 22.1 & 12.7 & 34.8 \\
\bottomrule
\end{tabular}
\caption{Sample Data Table}
\label{tab:sample}
\end{table}

\section{Lists}
\subsection{Bullet Points}
\begin{itemize}
  \item First item
  \item Second item
  \item Third item with sub-items:
    \begin{itemize}
      \item Sub-item 1
      \item Sub-item 2
    \end{itemize}
\end{itemize}

\subsection{Numbered List}
\begin{enumerate}
  \item First step
  \item Second step
  \item Third step
\end{enumerate}

\section{Citations and References}
When you use the citation search feature, references will be automatically managed. For example, you might cite a paper like this: \cite{example2024}. All citations will be automatically added to your \texttt{cite.bib} file.

\section{Conclusion}
This template demonstrates basic \LaTeX{} features including mathematical formulas, tables, and lists. The bibliography system is now set up to work seamlessly with Bibby AI's citation management. Feel free to modify it for your needs!

\bibliographystyle{plain}
\bibliography{cite}

\end{document}\documentclass{article}
\usepackage[utf8]{inputenc}
\usepackage{booktabs}
\usepackage{amsmath}
\usepackage{natbib}

\title{My Document}
\author{Your Name}
\date{\today}

\begin{document}

\maketitle

\section{Introduction}
This is a sample document demonstrating various \LaTeX{} features including mathematical equations, tables, and formatting. You can cite papers using the citation search feature, and they will be automatically added to your \texttt{cite.bib} file.

\section{Mathematical Expressions}
Here are some example equations:

\subsection{Quadratic Formula}
The quadratic formula is given by:
\[x = \frac{-b \pm \sqrt{b^2 - 4ac}}{2a}\]

\subsection{Matrix Example}
A sample matrix:
\[
\begin{pmatrix}
a_{11} & a_{12} & a_{13} \\
a_{21} & a_{22} & a_{23} \\
a_{31} & a_{32} & a_{33}
\end{pmatrix}
\]

\section{Tables}
Below is a sample table showing data:

\begin{table}[h]
\centering
\begin{tabular}{@{}lrrr@{}}
\toprule
\textbf{Category} & \textbf{Value 1} & \textbf{Value 2} & \textbf{Total} \\
\midrule
Item A & 10.5 & 20.3 & 30.8 \\
Item B & 15.2 & 25.1 & 40.3 \\
Item C & 22.1 & 12.7 & 34.8 \\
\bottomrule
\end{tabular}
\caption{Sample Data Table}
\label{tab:sample}
\end{table}

\section{Lists}
\subsection{Bullet Points}
\begin{itemize}
  \item First item
  \item Second item
  \item Third item with sub-items:
    \begin{itemize}
      \item Sub-item 1
      \item Sub-item 2
    \end{itemize}
\end{itemize}

\subsection{Numbered List}
\begin{enumerate}
  \item First step
  \item Second step
  \item Third step
\end{enumerate}

\section{Citations and References}
When you use the citation search feature, references will be automatically managed. For example, you might cite a paper like this: \cite{example2024}. All citations will be automatically added to your \texttt{cite.bib} file.

\section{Conclusion}
This template demonstrates basic \LaTeX{} features including mathematical formulas, tables, and lists. The bibliography system is now set up to work seamlessly with Bibby AI's citation management. Feel free to modify it for your needs!

\bibliographystyle{plain}
\bibliography{cite}

\end{document}\documentclass{article}
\usepackage[utf8]{inputenc}
\usepackage{booktabs}
\usepackage{amsmath}
\usepackage{natbib}

\title{My Document}
\author{Your Name}
\date{\today}

\begin{document}

\maketitle

\section{Introduction}
This is a sample document demonstrating various \LaTeX{} features including mathematical equations, tables, and formatting. You can cite papers using the citation search feature, and they will be automatically added to your \texttt{cite.bib} file.

\section{Mathematical Expressions}
Here are some example equations:

\subsection{Quadratic Formula}
The quadratic formula is given by:
\[x = \frac{-b \pm \sqrt{b^2 - 4ac}}{2a}\]

\subsection{Matrix Example}
A sample matrix:
\[
\begin{pmatrix}
a_{11} & a_{12} & a_{13} \\
a_{21} & a_{22} & a_{23} \\
a_{31} & a_{32} & a_{33}
\end{pmatrix}
\]

\section{Tables}
Below is a sample table showing data:

\begin{table}[h]
\centering
\begin{tabular}{@{}lrrr@{}}
\toprule
\textbf{Category} & \textbf{Value 1} & \textbf{Value 2} & \textbf{Total} \\
\midrule
Item A & 10.5 & 20.3 & 30.8 \\
Item B & 15.2 & 25.1 & 40.3 \\
Item C & 22.1 & 12.7 & 34.8 \\
\bottomrule
\end{tabular}
\caption{Sample Data Table}
\label{tab:sample}
\end{table}

\section{Lists}
\subsection{Bullet Points}
\begin{itemize}
  \item First item
  \item Second item
  \item Third item with sub-items:
    \begin{itemize}
      \item Sub-item 1
      \item Sub-item 2
    \end{itemize}
\end{itemize}

\subsection{Numbered List}
\begin{enumerate}
  \item First step
  \item Second step
  \item Third step
\end{enumerate}

\section{Citations and References}
When you use the citation search feature, references will be automatically managed. For example, you might cite a paper like this: \cite{example2024}. All citations will be automatically added to your \texttt{cite.bib} file.

\section{Conclusion}
This template demonstrates basic \LaTeX{} features including mathematical formulas, tables, and lists. The bibliography system is now set up to work seamlessly with Bibby AI's citation management. Feel free to modify it for your needs!

\bibliographystyle{plain}
\bibliography{cite}

\end{document}\documentclass{article}
\usepackage[utf8]{inputenc}
\usepackage{booktabs}
\usepackage{amsmath}
\usepackage{natbib}

\title{My Document}
\author{Your Name}
\date{\today}

\begin{document}

\maketitle

\section{Introduction}
This is a sample document demonstrating various \LaTeX{} features including mathematical equations, tables, and formatting. You can cite papers using the citation search feature, and they will be automatically added to your \texttt{cite.bib} file.

\section{Mathematical Expressions}
Here are some example equations:

\subsection{Quadratic Formula}
The quadratic formula is given by:
\[x = \frac{-b \pm \sqrt{b^2 - 4ac}}{2a}\]

\subsection{Matrix Example}
A sample matrix:
\[
\begin{pmatrix}
a_{11} & a_{12} & a_{13} \\
a_{21} & a_{22} & a_{23} \\
a_{31} & a_{32} & a_{33}
\end{pmatrix}
\]

\section{Tables}
Below is a sample table showing data:

\begin{table}[h]
\centering
\begin{tabular}{@{}lrrr@{}}
\toprule
\textbf{Category} & \textbf{Value 1} & \textbf{Value 2} & \textbf{Total} \\
\midrule
Item A & 10.5 & 20.3 & 30.8 \\
Item B & 15.2 & 25.1 & 40.3 \\
Item C & 22.1 & 12.7 & 34.8 \\
\bottomrule
\end{tabular}
\caption{Sample Data Table}
\label{tab:sample}
\end{table}

\section{Lists}
\subsection{Bullet Points}
\begin{itemize}
  \item First item
  \item Second item
  \item Third item with sub-items:
    \begin{itemize}
      \item Sub-item 1
      \item Sub-item 2
    \end{itemize}
\end{itemize}

\subsection{Numbered List}
\begin{enumerate}
  \item First step
  \item Second step
  \item Third step
\end{enumerate}

\section{Citations and References}
When you use the citation search feature, references will be automatically managed. For example, you might cite a paper like this: \cite{example2024}. All citations will be automatically added to your \texttt{cite.bib} file.

\section{Conclusion}
This template demonstrates basic \LaTeX{} features including mathematical formulas, tables, and lists. The bibliography system is now set up to work seamlessly with Bibby AI's citation management. Feel free to modify it for your needs!

\bibliographystyle{plain}
\bibliography{cite}

\end{document}\documentclass{article}
\usepackage[utf8]{inputenc}
\usepackage{booktabs}
\usepackage{amsmath}
\usepackage{natbib}

\title{My Document}
\author{Your Name}
\date{\today}

\begin{document}

\maketitle

\section{Introduction}
This is a sample document demonstrating various \LaTeX{} features including mathematical equations, tables, and formatting. You can cite papers using the citation search feature, and they will be automatically added to your \texttt{cite.bib} file.

\section{Mathematical Expressions}
Here are some example equations:

\subsection{Quadratic Formula}
The quadratic formula is given by:
\[x = \frac{-b \pm \sqrt{b^2 - 4ac}}{2a}\]

\subsection{Matrix Example}
A sample matrix:
\[
\begin{pmatrix}
a_{11} & a_{12} & a_{13} \\
a_{21} & a_{22} & a_{23} \\
a_{31} & a_{32} & a_{33}
\end{pmatrix}
\]

\section{Tables}
Below is a sample table showing data:

\begin{table}[h]
\centering
\begin{tabular}{@{}lrrr@{}}
\toprule
\textbf{Category} & \textbf{Value 1} & \textbf{Value 2} & \textbf{Total} \\
\midrule
Item A & 10.5 & 20.3 & 30.8 \\
Item B & 15.2 & 25.1 & 40.3 \\
Item C & 22.1 & 12.7 & 34.8 \\
\bottomrule
\end{tabular}
\caption{Sample Data Table}
\label{tab:sample}
\end{table}

\section{Lists}
\subsection{Bullet Points}
\begin{itemize}
  \item First item
  \item Second item
  \item Third item with sub-items:
    \begin{itemize}
      \item Sub-item 1
      \item Sub-item 2
    \end{itemize}
\end{itemize}

\subsection{Numbered List}
\begin{enumerate}
  \item First step
  \item Second step
  \item Third step
\end{enumerate}

\section{Citations and References}
When you use the citation search feature, references will be automatically managed. For example, you might cite a paper like this: \cite{example2024}. All citations will be automatically added to your \texttt{cite.bib} file.

\section{Conclusion}
This template demonstrates basic \LaTeX{} features including mathematical formulas, tables, and lists. The bibliography system is now set up to work seamlessly with Bibby AI's citation management. Feel free to modify it for your needs!

\bibliographystyle{plain}
\bibliography{cite}

\end{document}\documentclass{article}
\usepackage[utf8]{inputenc}
\usepackage{booktabs}
\usepackage{amsmath}
\usepackage{natbib}

\title{My Document}
\author{Your Name}
\date{\today}

\begin{document}

\maketitle

\section{Introduction}
This is a sample document demonstrating various \LaTeX{} features including mathematical equations, tables, and formatting. You can cite papers using the citation search feature, and they will be automatically added to your \texttt{cite.bib} file.

\section{Mathematical Expressions}
Here are some example equations:

\subsection{Quadratic Formula}
The quadratic formula is given by:
\[x = \frac{-b \pm \sqrt{b^2 - 4ac}}{2a}\]

\subsection{Matrix Example}
A sample matrix:
\[
\begin{pmatrix}
a_{11} & a_{12} & a_{13} \\
a_{21} & a_{22} & a_{23} \\
a_{31} & a_{32} & a_{33}
\end{pmatrix}
\]

\section{Tables}
Below is a sample table showing data:

\begin{table}[h]
\centering
\begin{tabular}{@{}lrrr@{}}
\toprule
\textbf{Category} & \textbf{Value 1} & \textbf{Value 2} & \textbf{Total} \\
\midrule
Item A & 10.5 & 20.3 & 30.8 \\
Item B & 15.2 & 25.1 & 40.3 \\
Item C & 22.1 & 12.7 & 34.8 \\
\bottomrule
\end{tabular}
\caption{Sample Data Table}
\label{tab:sample}
\end{table}

\section{Lists}
\subsection{Bullet Points}
\begin{itemize}
  \item First item
  \item Second item
  \item Third item with sub-items:
    \begin{itemize}
      \item Sub-item 1
      \item Sub-item 2
    \end{itemize}
\end{itemize}

\subsection{Numbered List}
\begin{enumerate}
  \item First step
  \item Second step
  \item Third step
\end{enumerate}

\section{Citations and References}
When you use the citation search feature, references will be automatically managed. For example, you might cite a paper like this: \cite{example2024}. All citations will be automatically added to your \texttt{cite.bib} file.

\section{Conclusion}
This template demonstrates basic \LaTeX{} features including mathematical formulas, tables, and lists. The bibliography system is now set up to work seamlessly with Bibby AI's citation management. Feel free to modify it for your needs!

\bibliographystyle{plain}
\bibliography{cite}

\end{document}\documentclass{article}
\usepackage[utf8]{inputenc}
\usepackage{booktabs}
\usepackage{amsmath}
\usepackage{natbib}

\title{My Document}
\author{Your Name}
\date{\today}

\begin{document}

\maketitle

\section{Introduction}
This is a sample document demonstrating various \LaTeX{} features including mathematical equations, tables, and formatting. You can cite papers using the citation search feature, and they will be automatically added to your \texttt{cite.bib} file.

\section{Mathematical Expressions}
Here are some example equations:

\subsection{Quadratic Formula}
The quadratic formula is given by:
\[x = \frac{-b \pm \sqrt{b^2 - 4ac}}{2a}\]

\subsection{Matrix Example}
A sample matrix:
\[
\begin{pmatrix}
a_{11} & a_{12} & a_{13} \\
a_{21} & a_{22} & a_{23} \\
a_{31} & a_{32} & a_{33}
\end{pmatrix}
\]

\section{Tables}
Below is a sample table showing data:

\begin{table}[h]
\centering
\begin{tabular}{@{}lrrr@{}}
\toprule
\textbf{Category} & \textbf{Value 1} & \textbf{Value 2} & \textbf{Total} \\
\midrule
Item A & 10.5 & 20.3 & 30.8 \\
Item B & 15.2 & 25.1 & 40.3 \\
Item C & 22.1 & 12.7 & 34.8 \\
\bottomrule
\end{tabular}
\caption{Sample Data Table}
\label{tab:sample}
\end{table}

\section{Lists}
\subsection{Bullet Points}
\begin{itemize}
  \item First item
  \item Second item
  \item Third item with sub-items:
    \begin{itemize}
      \item Sub-item 1
      \item Sub-item 2
    \end{itemize}
\end{itemize}

\subsection{Numbered List}
\begin{enumerate}
  \item First step
  \item Second step
  \item Third step
\end{enumerate}

\section{Citations and References}
When you use the citation search feature, references will be automatically managed. For example, you might cite a paper like this: \cite{example2024}. All citations will be automatically added to your \texttt{cite.bib} file.

\section{Conclusion}
This template demonstrates basic \LaTeX{} features including mathematical formulas, tables, and lists. The bibliography system is now set up to work seamlessly with Bibby AI's citation management. Feel free to modify it for your needs!

\bibliographystyle{plain}
\bibliography{cite}

\end{document}\documentclass{article}
\usepackage[utf8]{inputenc}
\usepackage{booktabs}
\usepackage{amsmath}
\usepackage{natbib}

\title{My Document}
\author{Your Name}
\date{\today}

\begin{document}

\maketitle

\section{Introduction}
This is a sample document demonstrating various \LaTeX{} features including mathematical equations, tables, and formatting. You can cite papers using the citation search feature, and they will be automatically added to your \texttt{cite.bib} file.

\section{Mathematical Expressions}
Here are some example equations:

\subsection{Quadratic Formula}
The quadratic formula is given by:
\[x = \frac{-b \pm \sqrt{b^2 - 4ac}}{2a}\]

\subsection{Matrix Example}
A sample matrix:
\[
\begin{pmatrix}
a_{11} & a_{12} & a_{13} \\
a_{21} & a_{22} & a_{23} \\
a_{31} & a_{32} & a_{33}
\end{pmatrix}
\]

\section{Tables}
Below is a sample table showing data:

\begin{table}[h]
\centering
\begin{tabular}{@{}lrrr@{}}
\toprule
\textbf{Category} & \textbf{Value 1} & \textbf{Value 2} & \textbf{Total} \\
\midrule
Item A & 10.5 & 20.3 & 30.8 \\
Item B & 15.2 & 25.1 & 40.3 \\
Item C & 22.1 & 12.7 & 34.8 \\
\bottomrule
\end{tabular}
\caption{Sample Data Table}
\label{tab:sample}
\end{table}

\section{Lists}
\subsection{Bullet Points}
\begin{itemize}
  \item First item
  \item Second item
  \item Third item with sub-items:
    \begin{itemize}
      \item Sub-item 1
      \item Sub-item 2
    \end{itemize}
\end{itemize}

\subsection{Numbered List}
\begin{enumerate}
  \item First step
  \item Second step
  \item Third step
\end{enumerate}

\section{Citations and References}
When you use the citation search feature, references will be automatically managed. For example, you might cite a paper like this: \cite{example2024}. All citations will be automatically added to your \texttt{cite.bib} file.

\section{Conclusion}
This template demonstrates basic \LaTeX{} features including mathematical formulas, tables, and lists. The bibliography system is now set up to work seamlessly with Bibby AI's citation management. Feel free to modify it for your needs!

\bibliographystyle{plain}
\bibliography{cite}

\end{document}\documentclass{article}
\usepackage[utf8]{inputenc}
\usepackage{booktabs}
\usepackage{amsmath}
\usepackage{natbib}

\title{My Document}
\author{Your Name}
\date{\today}

\begin{document}

\maketitle

\section{Introduction}
This is a sample document demonstrating various \LaTeX{} features including mathematical equations, tables, and formatting. You can cite papers using the citation search feature, and they will be automatically added to your \texttt{cite.bib} file.

\section{Mathematical Expressions}
Here are some example equations:

\subsection{Quadratic Formula}
The quadratic formula is given by:
\[x = \frac{-b \pm \sqrt{b^2 - 4ac}}{2a}\]

\subsection{Matrix Example}
A sample matrix:
\[
\begin{pmatrix}
a_{11} & a_{12} & a_{13} \\
a_{21} & a_{22} & a_{23} \\
a_{31} & a_{32} & a_{33}
\end{pmatrix}
\]

\section{Tables}
Below is a sample table showing data:

\begin{table}[h]
\centering
\begin{tabular}{@{}lrrr@{}}
\toprule
\textbf{Category} & \textbf{Value 1} & \textbf{Value 2} & \textbf{Total} \\
\midrule
Item A & 10.5 & 20.3 & 30.8 \\
Item B & 15.2 & 25.1 & 40.3 \\
Item C & 22.1 & 12.7 & 34.8 \\
\bottomrule
\end{tabular}
\caption{Sample Data Table}
\label{tab:sample}
\end{table}

\section{Lists}
\subsection{Bullet Points}
\begin{itemize}
  \item First item
  \item Second item
  \item Third item with sub-items:
    \begin{itemize}
      \item Sub-item 1
      \item Sub-item 2
    \end{itemize}
\end{itemize}

\subsection{Numbered List}
\begin{enumerate}
  \item First step
  \item Second step
  \item Third step
\end{enumerate}

\section{Citations and References}
When you use the citation search feature, references will be automatically managed. For example, you might cite a paper like this: \cite{example2024}. All citations will be automatically added to your \texttt{cite.bib} file.

\section{Conclusion}
This template demonstrates basic \LaTeX{} features including mathematical formulas, tables, and lists. The bibliography system is now set up to work seamlessly with Bibby AI's citation management. Feel free to modify it for your needs!

\bibliographystyle{plain}
\bibliography{cite}

\end{document}\documentclass{article}
\usepackage[utf8]{inputenc}
\usepackage{booktabs}
\usepackage{amsmath}
\usepackage{natbib}

\title{My Document}
\author{Your Name}
\date{\today}

\begin{document}

\maketitle

\section{Introduction}
This is a sample document demonstrating various \LaTeX{} features including mathematical equations, tables, and formatting. You can cite papers using the citation search feature, and they will be automatically added to your \texttt{cite.bib} file.

\section{Mathematical Expressions}
Here are some example equations:

\subsection{Quadratic Formula}
The quadratic formula is given by:
\[x = \frac{-b \pm \sqrt{b^2 - 4ac}}{2a}\]

\subsection{Matrix Example}
A sample matrix:
\[
\begin{pmatrix}
a_{11} & a_{12} & a_{13} \\
a_{21} & a_{22} & a_{23} \\
a_{31} & a_{32} & a_{33}
\end{pmatrix}
\]

\section{Tables}
Below is a sample table showing data:

\begin{table}[h]
\centering
\begin{tabular}{@{}lrrr@{}}
\toprule
\textbf{Category} & \textbf{Value 1} & \textbf{Value 2} & \textbf{Total} \\
\midrule
Item A & 10.5 & 20.3 & 30.8 \\
Item B & 15.2 & 25.1 & 40.3 \\
Item C & 22.1 & 12.7 & 34.8 \\
\bottomrule
\end{tabular}
\caption{Sample Data Table}
\label{tab:sample}
\end{table}

\section{Lists}
\subsection{Bullet Points}
\begin{itemize}
  \item First item
  \item Second item
  \item Third item with sub-items:
    \begin{itemize}
      \item Sub-item 1
      \item Sub-item 2
    \end{itemize}
\end{itemize}

\subsection{Numbered List}
\begin{enumerate}
  \item First step
  \item Second step
  \item Third step
\end{enumerate}

\section{Citations and References}
When you use the citation search feature, references will be automatically managed. For example, you might cite a paper like this: \cite{example2024}. All citations will be automatically added to your \texttt{cite.bib} file.

\section{Conclusion}
This template demonstrates basic \LaTeX{} features including mathematical formulas, tables, and lists. The bibliography system is now set up to work seamlessly with Bibby AI's citation management. Feel free to modify it for your needs!

\bibliographystyle{plain}
\bibliography{cite}

\end{document}\documentclass{article}
\usepackage[utf8]{inputenc}
\usepackage{booktabs}
\usepackage{amsmath}
\usepackage{natbib}

\title{My Document}
\author{Your Name}
\date{\today}

\begin{document}

\maketitle

\section{Introduction}
This is a sample document demonstrating various \LaTeX{} features including mathematical equations, tables, and formatting. You can cite papers using the citation search feature, and they will be automatically added to your \texttt{cite.bib} file.

\section{Mathematical Expressions}
Here are some example equations:

\subsection{Quadratic Formula}
The quadratic formula is given by:
\[x = \frac{-b \pm \sqrt{b^2 - 4ac}}{2a}\]

\subsection{Matrix Example}
A sample matrix:
\[
\begin{pmatrix}
a_{11} & a_{12} & a_{13} \\
a_{21} & a_{22} & a_{23} \\
a_{31} & a_{32} & a_{33}
\end{pmatrix}
\]

\section{Tables}
Below is a sample table showing data:

\begin{table}[h]
\centering
\begin{tabular}{@{}lrrr@{}}
\toprule
\textbf{Category} & \textbf{Value 1} & \textbf{Value 2} & \textbf{Total} \\
\midrule
Item A & 10.5 & 20.3 & 30.8 \\
Item B & 15.2 & 25.1 & 40.3 \\
Item C & 22.1 & 12.7 & 34.8 \\
\bottomrule
\end{tabular}
\caption{Sample Data Table}
\label{tab:sample}
\end{table}

\section{Lists}
\subsection{Bullet Points}
\begin{itemize}
  \item First item
  \item Second item
  \item Third item with sub-items:
    \begin{itemize}
      \item Sub-item 1
      \item Sub-item 2
    \end{itemize}
\end{itemize}

\subsection{Numbered List}
\begin{enumerate}
  \item First step
  \item Second step
  \item Third step
\end{enumerate}

\section{Citations and References}
When you use the citation search feature, references will be automatically managed. For example, you might cite a paper like this: \cite{example2024}. All citations will be automatically added to your \texttt{cite.bib} file.

\section{Conclusion}
This template demonstrates basic \LaTeX{} features including mathematical formulas, tables, and lists. The bibliography system is now set up to work seamlessly with Bibby AI's citation management. Feel free to modify it for your needs!

\bibliographystyle{plain}
\bibliography{cite}

\end{document}\documentclass{article}
\usepackage[utf8]{inputenc}
\usepackage{booktabs}
\usepackage{amsmath}
\usepackage{natbib}

\title{My Document}
\author{Your Name}
\date{\today}

\begin{document}

\maketitle

\section{Introduction}
This is a sample document demonstrating various \LaTeX{} features including mathematical equations, tables, and formatting. You can cite papers using the citation search feature, and they will be automatically added to your \texttt{cite.bib} file.

\section{Mathematical Expressions}
Here are some example equations:

\subsection{Quadratic Formula}
The quadratic formula is given by:
\[x = \frac{-b \pm \sqrt{b^2 - 4ac}}{2a}\]

\subsection{Matrix Example}
A sample matrix:
\[
\begin{pmatrix}
a_{11} & a_{12} & a_{13} \\
a_{21} & a_{22} & a_{23} \\
a_{31} & a_{32} & a_{33}
\end{pmatrix}
\]

\section{Tables}
Below is a sample table showing data:

\begin{table}[h]
\centering
\begin{tabular}{@{}lrrr@{}}
\toprule
\textbf{Category} & \textbf{Value 1} & \textbf{Value 2} & \textbf{Total} \\
\midrule
Item A & 10.5 & 20.3 & 30.8 \\
Item B & 15.2 & 25.1 & 40.3 \\
Item C & 22.1 & 12.7 & 34.8 \\
\bottomrule
\end{tabular}
\caption{Sample Data Table}
\label{tab:sample}
\end{table}

\section{Lists}
\subsection{Bullet Points}
\begin{itemize}
  \item First item
  \item Second item
  \item Third item with sub-items:
    \begin{itemize}
      \item Sub-item 1
      \item Sub-item 2
    \end{itemize}
\end{itemize}

\subsection{Numbered List}
\begin{enumerate}
  \item First step
  \item Second step
  \item Third step
\end{enumerate}

\section{Citations and References}
When you use the citation search feature, references will be automatically managed. For example, you might cite a paper like this: \cite{example2024}. All citations will be automatically added to your \texttt{cite.bib} file.

\section{Conclusion}
This template demonstrates basic \LaTeX{} features including mathematical formulas, tables, and lists. The bibliography system is now set up to work seamlessly with Bibby AI's citation management. Feel free to modify it for your needs!

\bibliographystyle{plain}
\bibliography{cite}

\end{document}\documentclass{article}
\usepackage[utf8]{inputenc}
\usepackage{booktabs}
\usepackage{amsmath}
\usepackage{natbib}

\title{My Document}
\author{Your Name}
\date{\today}

\begin{document}

\maketitle

\section{Introduction}
This is a sample document demonstrating various \LaTeX{} features including mathematical equations, tables, and formatting. You can cite papers using the citation search feature, and they will be automatically added to your \texttt{cite.bib} file.

\section{Mathematical Expressions}
Here are some example equations:

\subsection{Quadratic Formula}
The quadratic formula is given by:
\[x = \frac{-b \pm \sqrt{b^2 - 4ac}}{2a}\]

\subsection{Matrix Example}
A sample matrix:
\[
\begin{pmatrix}
a_{11} & a_{12} & a_{13} \\
a_{21} & a_{22} & a_{23} \\
a_{31} & a_{32} & a_{33}
\end{pmatrix}
\]

\section{Tables}
Below is a sample table showing data:

\begin{table}[h]
\centering
\begin{tabular}{@{}lrrr@{}}
\toprule
\textbf{Category} & \textbf{Value 1} & \textbf{Value 2} & \textbf{Total} \\
\midrule
Item A & 10.5 & 20.3 & 30.8 \\
Item B & 15.2 & 25.1 & 40.3 \\
Item C & 22.1 & 12.7 & 34.8 \\
\bottomrule
\end{tabular}
\caption{Sample Data Table}
\label{tab:sample}
\end{table}

\section{Lists}
\subsection{Bullet Points}
\begin{itemize}
  \item First item
  \item Second item
  \item Third item with sub-items:
    \begin{itemize}
      \item Sub-item 1
      \item Sub-item 2
    \end{itemize}
\end{itemize}

\subsection{Numbered List}
\begin{enumerate}
  \item First step
  \item Second step
  \item Third step
\end{enumerate}

\section{Citations and References}
When you use the citation search feature, references will be automatically managed. For example, you might cite a paper like this: \cite{example2024}. All citations will be automatically added to your \texttt{cite.bib} file.

\section{Conclusion}
This template demonstrates basic \LaTeX{} features including mathematical formulas, tables, and lists. The bibliography system is now set up to work seamlessly with Bibby AI's citation management. Feel free to modify it for your needs!

\bibliographystyle{plain}
\bibliography{cite}

\end{document}\documentclass{article}
\usepackage[utf8]{inputenc}
\usepackage{booktabs}
\usepackage{amsmath}
\usepackage{natbib}

\title{My Document}
\author{Your Name}
\date{\today}

\begin{document}

\maketitle

\section{Introduction}
This is a sample document demonstrating various \LaTeX{} features including mathematical equations, tables, and formatting. You can cite papers using the citation search feature, and they will be automatically added to your \texttt{cite.bib} file.

\section{Mathematical Expressions}
Here are some example equations:

\subsection{Quadratic Formula}
The quadratic formula is given by:
\[x = \frac{-b \pm \sqrt{b^2 - 4ac}}{2a}\]

\subsection{Matrix Example}
A sample matrix:
\[
\begin{pmatrix}
a_{11} & a_{12} & a_{13} \\
a_{21} & a_{22} & a_{23} \\
a_{31} & a_{32} & a_{33}
\end{pmatrix}
\]

\section{Tables}
Below is a sample table showing data:

\begin{table}[h]
\centering
\begin{tabular}{@{}lrrr@{}}
\toprule
\textbf{Category} & \textbf{Value 1} & \textbf{Value 2} & \textbf{Total} \\
\midrule
Item A & 10.5 & 20.3 & 30.8 \\
Item B & 15.2 & 25.1 & 40.3 \\
Item C & 22.1 & 12.7 & 34.8 \\
\bottomrule
\end{tabular}
\caption{Sample Data Table}
\label{tab:sample}
\end{table}

\section{Lists}
\subsection{Bullet Points}
\begin{itemize}
  \item First item
  \item Second item
  \item Third item with sub-items:
    \begin{itemize}
      \item Sub-item 1
      \item Sub-item 2
    \end{itemize}
\end{itemize}

\subsection{Numbered List}
\begin{enumerate}
  \item First step
  \item Second step
  \item Third step
\end{enumerate}

\section{Citations and References}
When you use the citation search feature, references will be automatically managed. For example, you might cite a paper like this: \cite{example2024}. All citations will be automatically added to your \texttt{cite.bib} file.

\section{Conclusion}
This template demonstrates basic \LaTeX{} features including mathematical formulas, tables, and lists. The bibliography system is now set up to work seamlessly with Bibby AI's citation management. Feel free to modify it for your needs!

\bibliographystyle{plain}
\bibliography{cite}

\end{document}\documentclass{article}
\usepackage[utf8]{inputenc}
\usepackage{booktabs}
\usepackage{amsmath}
\usepackage{natbib}

\title{My Document}
\author{Your Name}
\date{\today}

\begin{document}

\maketitle

\section{Introduction}
This is a sample document demonstrating various \LaTeX{} features including mathematical equations, tables, and formatting. You can cite papers using the citation search feature, and they will be automatically added to your \texttt{cite.bib} file.

\section{Mathematical Expressions}
Here are some example equations:

\subsection{Quadratic Formula}
The quadratic formula is given by:
\[x = \frac{-b \pm \sqrt{b^2 - 4ac}}{2a}\]

\subsection{Matrix Example}
A sample matrix:
\[
\begin{pmatrix}
a_{11} & a_{12} & a_{13} \\
a_{21} & a_{22} & a_{23} \\
a_{31} & a_{32} & a_{33}
\end{pmatrix}
\]

\section{Tables}
Below is a sample table showing data:

\begin{table}[h]
\centering
\begin{tabular}{@{}lrrr@{}}
\toprule
\textbf{Category} & \textbf{Value 1} & \textbf{Value 2} & \textbf{Total} \\
\midrule
Item A & 10.5 & 20.3 & 30.8 \\
Item B & 15.2 & 25.1 & 40.3 \\
Item C & 22.1 & 12.7 & 34.8 \\
\bottomrule
\end{tabular}
\caption{Sample Data Table}
\label{tab:sample}
\end{table}

\section{Lists}
\subsection{Bullet Points}
\begin{itemize}
  \item First item
  \item Second item
  \item Third item with sub-items:
    \begin{itemize}
      \item Sub-item 1
      \item Sub-item 2
    \end{itemize}
\end{itemize}

\subsection{Numbered List}
\begin{enumerate}
  \item First step
  \item Second step
  \item Third step
\end{enumerate}

\section{Citations and References}
When you use the citation search feature, references will be automatically managed. For example, you might cite a paper like this: \cite{example2024}. All citations will be automatically added to your \texttt{cite.bib} file.

\section{Conclusion}
This template demonstrates basic \LaTeX{} features including mathematical formulas, tables, and lists. The bibliography system is now set up to work seamlessly with Bibby AI's citation management. Feel free to modify it for your needs!

\bibliographystyle{plain}
\bibliography{cite}

\end{document}\documentclass{article}
\usepackage[utf8]{inputenc}
\usepackage{booktabs}
\usepackage{amsmath}
\usepackage{natbib}

\title{My Document}
\author{Your Name}
\date{\today}

\begin{document}

\maketitle

\section{Introduction}
This is a sample document demonstrating various \LaTeX{} features including mathematical equations, tables, and formatting. You can cite papers using the citation search feature, and they will be automatically added to your \texttt{cite.bib} file.

\section{Mathematical Expressions}
Here are some example equations:

\subsection{Quadratic Formula}
The quadratic formula is given by:
\[x = \frac{-b \pm \sqrt{b^2 - 4ac}}{2a}\]

\subsection{Matrix Example}
A sample matrix:
\[
\begin{pmatrix}
a_{11} & a_{12} & a_{13} \\
a_{21} & a_{22} & a_{23} \\
a_{31} & a_{32} & a_{33}
\end{pmatrix}
\]

\section{Tables}
Below is a sample table showing data:

\begin{table}[h]
\centering
\begin{tabular}{@{}lrrr@{}}
\toprule
\textbf{Category} & \textbf{Value 1} & \textbf{Value 2} & \textbf{Total} \\
\midrule
Item A & 10.5 & 20.3 & 30.8 \\
Item B & 15.2 & 25.1 & 40.3 \\
Item C & 22.1 & 12.7 & 34.8 \\
\bottomrule
\end{tabular}
\caption{Sample Data Table}
\label{tab:sample}
\end{table}

\section{Lists}
\subsection{Bullet Points}
\begin{itemize}
  \item First item
  \item Second item
  \item Third item with sub-items:
    \begin{itemize}
      \item Sub-item 1
      \item Sub-item 2
    \end{itemize}
\end{itemize}

\subsection{Numbered List}
\begin{enumerate}
  \item First step
  \item Second step
  \item Third step
\end{enumerate}

\section{Citations and References}
When you use the citation search feature, references will be automatically managed. For example, you might cite a paper like this: \cite{example2024}. All citations will be automatically added to your \texttt{cite.bib} file.

\section{Conclusion}
This template demonstrates basic \LaTeX{} features including mathematical formulas, tables, and lists. The bibliography system is now set up to work seamlessly with Bibby AI's citation management. Feel free to modify it for your needs!

\bibliographystyle{plain}
\bibliography{cite}

\end{document}\documentclass{article}
\usepackage[utf8]{inputenc}
\usepackage{booktabs}
\usepackage{amsmath}
\usepackage{natbib}

\title{My Document}
\author{Your Name}
\date{\today}

\begin{document}

\maketitle

\section{Introduction}
This is a sample document demonstrating various \LaTeX{} features including mathematical equations, tables, and formatting. You can cite papers using the citation search feature, and they will be automatically added to your \texttt{cite.bib} file.

\section{Mathematical Expressions}
Here are some example equations:

\subsection{Quadratic Formula}
The quadratic formula is given by:
\[x = \frac{-b \pm \sqrt{b^2 - 4ac}}{2a}\]

\subsection{Matrix Example}
A sample matrix:
\[
\begin{pmatrix}
a_{11} & a_{12} & a_{13} \\
a_{21} & a_{22} & a_{23} \\
a_{31} & a_{32} & a_{33}
\end{pmatrix}
\]

\section{Tables}
Below is a sample table showing data:

\begin{table}[h]
\centering
\begin{tabular}{@{}lrrr@{}}
\toprule
\textbf{Category} & \textbf{Value 1} & \textbf{Value 2} & \textbf{Total} \\
\midrule
Item A & 10.5 & 20.3 & 30.8 \\
Item B & 15.2 & 25.1 & 40.3 \\
Item C & 22.1 & 12.7 & 34.8 \\
\bottomrule
\end{tabular}
\caption{Sample Data Table}
\label{tab:sample}
\end{table}

\section{Lists}
\subsection{Bullet Points}
\begin{itemize}
  \item First item
  \item Second item
  \item Third item with sub-items:
    \begin{itemize}
      \item Sub-item 1
      \item Sub-item 2
    \end{itemize}
\end{itemize}

\subsection{Numbered List}
\begin{enumerate}
  \item First step
  \item Second step
  \item Third step
\end{enumerate}

\section{Citations and References}
When you use the citation search feature, references will be automatically managed. For example, you might cite a paper like this: \cite{example2024}. All citations will be automatically added to your \texttt{cite.bib} file.

\section{Conclusion}
This template demonstrates basic \LaTeX{} features including mathematical formulas, tables, and lists. The bibliography system is now set up to work seamlessly with Bibby AI's citation management. Feel free to modify it for your needs!

\bibliographystyle{plain}
\bibliography{cite}

\end{document}\documentclass{article}
\usepackage[utf8]{inputenc}
\usepackage{booktabs}
\usepackage{amsmath}
\usepackage{natbib}

\title{My Document}
\author{Your Name}
\date{\today}

\begin{document}

\maketitle

\section{Introduction}
This is a sample document demonstrating various \LaTeX{} features including mathematical equations, tables, and formatting. You can cite papers using the citation search feature, and they will be automatically added to your \texttt{cite.bib} file.

\section{Mathematical Expressions}
Here are some example equations:

\subsection{Quadratic Formula}
The quadratic formula is given by:
\[x = \frac{-b \pm \sqrt{b^2 - 4ac}}{2a}\]

\subsection{Matrix Example}
A sample matrix:
\[
\begin{pmatrix}
a_{11} & a_{12} & a_{13} \\
a_{21} & a_{22} & a_{23} \\
a_{31} & a_{32} & a_{33}
\end{pmatrix}
\]

\section{Tables}
Below is a sample table showing data:

\begin{table}[h]
\centering
\begin{tabular}{@{}lrrr@{}}
\toprule
\textbf{Category} & \textbf{Value 1} & \textbf{Value 2} & \textbf{Total} \\
\midrule
Item A & 10.5 & 20.3 & 30.8 \\
Item B & 15.2 & 25.1 & 40.3 \\
Item C & 22.1 & 12.7 & 34.8 \\
\bottomrule
\end{tabular}
\caption{Sample Data Table}
\label{tab:sample}
\end{table}

\section{Lists}
\subsection{Bullet Points}
\begin{itemize}
  \item First item
  \item Second item
  \item Third item with sub-items:
    \begin{itemize}
      \item Sub-item 1
      \item Sub-item 2
    \end{itemize}
\end{itemize}

\subsection{Numbered List}
\begin{enumerate}
  \item First step
  \item Second step
  \item Third step
\end{enumerate}

\section{Citations and References}
When you use the citation search feature, references will be automatically managed. For example, you might cite a paper like this: \cite{example2024}. All citations will be automatically added to your \texttt{cite.bib} file.

\section{Conclusion}
This template demonstrates basic \LaTeX{} features including mathematical formulas, tables, and lists. The bibliography system is now set up to work seamlessly with Bibby AI's citation management. Feel free to modify it for your needs!

\bibliographystyle{plain}
\bibliography{cite}

\end{document}\documentclass{article}
\usepackage[utf8]{inputenc}
\usepackage{booktabs}
\usepackage{amsmath}
\usepackage{natbib}

\title{My Document}
\author{Your Name}
\date{\today}

\begin{document}

\maketitle

\section{Introduction}
This is a sample document demonstrating various \LaTeX{} features including mathematical equations, tables, and formatting. You can cite papers using the citation search feature, and they will be automatically added to your \texttt{cite.bib} file.

\section{Mathematical Expressions}
Here are some example equations:

\subsection{Quadratic Formula}
The quadratic formula is given by:
\[x = \frac{-b \pm \sqrt{b^2 - 4ac}}{2a}\]

\subsection{Matrix Example}
A sample matrix:
\[
\begin{pmatrix}
a_{11} & a_{12} & a_{13} \\
a_{21} & a_{22} & a_{23} \\
a_{31} & a_{32} & a_{33}
\end{pmatrix}
\]

\section{Tables}
Below is a sample table showing data:

\begin{table}[h]
\centering
\begin{tabular}{@{}lrrr@{}}
\toprule
\textbf{Category} & \textbf{Value 1} & \textbf{Value 2} & \textbf{Total} \\
\midrule
Item A & 10.5 & 20.3 & 30.8 \\
Item B & 15.2 & 25.1 & 40.3 \\
Item C & 22.1 & 12.7 & 34.8 \\
\bottomrule
\end{tabular}
\caption{Sample Data Table}
\label{tab:sample}
\end{table}

\section{Lists}
\subsection{Bullet Points}
\begin{itemize}
  \item First item
  \item Second item
  \item Third item with sub-items:
    \begin{itemize}
      \item Sub-item 1
      \item Sub-item 2
    \end{itemize}
\end{itemize}

\subsection{Numbered List}
\begin{enumerate}
  \item First step
  \item Second step
  \item Third step
\end{enumerate}

\section{Citations and References}
When you use the citation search feature, references will be automatically managed. For example, you might cite a paper like this: \cite{example2024}. All citations will be automatically added to your \texttt{cite.bib} file.

\section{Conclusion}
This template demonstrates basic \LaTeX{} features including mathematical formulas, tables, and lists. The bibliography system is now set up to work seamlessly with Bibby AI's citation management. Feel free to modify it for your needs!

\bibliographystyle{plain}
\bibliography{cite}

\end{document}\documentclass{article}
\usepackage[utf8]{inputenc}
\usepackage{booktabs}
\usepackage{amsmath}
\usepackage{natbib}

\title{My Document}
\author{Your Name}
\date{\today}

\begin{document}

\maketitle

\section{Introduction}
This is a sample document demonstrating various \LaTeX{} features including mathematical equations, tables, and formatting. You can cite papers using the citation search feature, and they will be automatically added to your \texttt{cite.bib} file.

\section{Mathematical Expressions}
Here are some example equations:

\subsection{Quadratic Formula}
The quadratic formula is given by:
\[x = \frac{-b \pm \sqrt{b^2 - 4ac}}{2a}\]

\subsection{Matrix Example}
A sample matrix:
\[
\begin{pmatrix}
a_{11} & a_{12} & a_{13} \\
a_{21} & a_{22} & a_{23} \\
a_{31} & a_{32} & a_{33}
\end{pmatrix}
\]

\section{Tables}
Below is a sample table showing data:

\begin{table}[h]
\centering
\begin{tabular}{@{}lrrr@{}}
\toprule
\textbf{Category} & \textbf{Value 1} & \textbf{Value 2} & \textbf{Total} \\
\midrule
Item A & 10.5 & 20.3 & 30.8 \\
Item B & 15.2 & 25.1 & 40.3 \\
Item C & 22.1 & 12.7 & 34.8 \\
\bottomrule
\end{tabular}
\caption{Sample Data Table}
\label{tab:sample}
\end{table}

\section{Lists}
\subsection{Bullet Points}
\begin{itemize}
  \item First item
  \item Second item
  \item Third item with sub-items:
    \begin{itemize}
      \item Sub-item 1
      \item Sub-item 2
    \end{itemize}
\end{itemize}

\subsection{Numbered List}
\begin{enumerate}
  \item First step
  \item Second step
  \item Third step
\end{enumerate}

\section{Citations and References}
When you use the citation search feature, references will be automatically managed. For example, you might cite a paper like this: \cite{example2024}. All citations will be automatically added to your \texttt{cite.bib} file.

\section{Conclusion}
This template demonstrates basic \LaTeX{} features including mathematical formulas, tables, and lists. The bibliography system is now set up to work seamlessly with Bibby AI's citation management. Feel free to modify it for your needs!

\bibliographystyle{plain}
\bibliography{cite}

\end{document}\documentclass{article}
\usepackage[utf8]{inputenc}
\usepackage{booktabs}
\usepackage{amsmath}
\usepackage{natbib}

\title{My Document}
\author{Your Name}
\date{\today}

\begin{document}

\maketitle

\section{Introduction}
This is a sample document demonstrating various \LaTeX{} features including mathematical equations, tables, and formatting. You can cite papers using the citation search feature, and they will be automatically added to your \texttt{cite.bib} file.

\section{Mathematical Expressions}
Here are some example equations:

\subsection{Quadratic Formula}
The quadratic formula is given by:
\[x = \frac{-b \pm \sqrt{b^2 - 4ac}}{2a}\]

\subsection{Matrix Example}
A sample matrix:
\[
\begin{pmatrix}
a_{11} & a_{12} & a_{13} \\
a_{21} & a_{22} & a_{23} \\
a_{31} & a_{32} & a_{33}
\end{pmatrix}
\]

\section{Tables}
Below is a sample table showing data:

\begin{table}[h]
\centering
\begin{tabular}{@{}lrrr@{}}
\toprule
\textbf{Category} & \textbf{Value 1} & \textbf{Value 2} & \textbf{Total} \\
\midrule
Item A & 10.5 & 20.3 & 30.8 \\
Item B & 15.2 & 25.1 & 40.3 \\
Item C & 22.1 & 12.7 & 34.8 \\
\bottomrule
\end{tabular}
\caption{Sample Data Table}
\label{tab:sample}
\end{table}

\section{Lists}
\subsection{Bullet Points}
\begin{itemize}
  \item First item
  \item Second item
  \item Third item with sub-items:
    \begin{itemize}
      \item Sub-item 1
      \item Sub-item 2
    \end{itemize}
\end{itemize}

\subsection{Numbered List}
\begin{enumerate}
  \item First step
  \item Second step
  \item Third step
\end{enumerate}

\section{Citations and References}
When you use the citation search feature, references will be automatically managed. For example, you might cite a paper like this: \cite{example2024}. All citations will be automatically added to your \texttt{cite.bib} file.

\section{Conclusion}
This template demonstrates basic \LaTeX{} features including mathematical formulas, tables, and lists. The bibliography system is now set up to work seamlessly with Bibby AI's citation management. Feel free to modify it for your needs!

\bibliographystyle{plain}
\bibliography{cite}

\end{document}\documentclass{article}
\usepackage[utf8]{inputenc}
\usepackage{booktabs}
\usepackage{amsmath}
\usepackage{natbib}

\title{My Document}
\author{Your Name}
\date{\today}

\begin{document}

\maketitle

\section{Introduction}
This is a sample document demonstrating various \LaTeX{} features including mathematical equations, tables, and formatting. You can cite papers using the citation search feature, and they will be automatically added to your \texttt{cite.bib} file.

\section{Mathematical Expressions}
Here are some example equations:

\subsection{Quadratic Formula}
The quadratic formula is given by:
\[x = \frac{-b \pm \sqrt{b^2 - 4ac}}{2a}\]

\subsection{Matrix Example}
A sample matrix:
\[
\begin{pmatrix}
a_{11} & a_{12} & a_{13} \\
a_{21} & a_{22} & a_{23} \\
a_{31} & a_{32} & a_{33}
\end{pmatrix}
\]

\section{Tables}
Below is a sample table showing data:

\begin{table}[h]
\centering
\begin{tabular}{@{}lrrr@{}}
\toprule
\textbf{Category} & \textbf{Value 1} & \textbf{Value 2} & \textbf{Total} \\
\midrule
Item A & 10.5 & 20.3 & 30.8 \\
Item B & 15.2 & 25.1 & 40.3 \\
Item C & 22.1 & 12.7 & 34.8 \\
\bottomrule
\end{tabular}
\caption{Sample Data Table}
\label{tab:sample}
\end{table}

\section{Lists}
\subsection{Bullet Points}
\begin{itemize}
  \item First item
  \item Second item
  \item Third item with sub-items:
    \begin{itemize}
      \item Sub-item 1
      \item Sub-item 2
    \end{itemize}
\end{itemize}

\subsection{Numbered List}
\begin{enumerate}
  \item First step
  \item Second step
  \item Third step
\end{enumerate}

\section{Citations and References}
When you use the citation search feature, references will be automatically managed. For example, you might cite a paper like this: \cite{example2024}. All citations will be automatically added to your \texttt{cite.bib} file.

\section{Conclusion}
This template demonstrates basic \LaTeX{} features including mathematical formulas, tables, and lists. The bibliography system is now set up to work seamlessly with Bibby AI's citation management. Feel free to modify it for your needs!

\bibliographystyle{plain}
\bibliography{cite}

\end{document}\documentclass{article}
\usepackage[utf8]{inputenc}
\usepackage{booktabs}
\usepackage{amsmath}
\usepackage{natbib}

\title{My Document}
\author{Your Name}
\date{\today}

\begin{document}

\maketitle

\section{Introduction}
This is a sample document demonstrating various \LaTeX{} features including mathematical equations, tables, and formatting. You can cite papers using the citation search feature, and they will be automatically added to your \texttt{cite.bib} file.

\section{Mathematical Expressions}
Here are some example equations:

\subsection{Quadratic Formula}
The quadratic formula is given by:
\[x = \frac{-b \pm \sqrt{b^2 - 4ac}}{2a}\]

\subsection{Matrix Example}
A sample matrix:
\[
\begin{pmatrix}
a_{11} & a_{12} & a_{13} \\
a_{21} & a_{22} & a_{23} \\
a_{31} & a_{32} & a_{33}
\end{pmatrix}
\]

\section{Tables}
Below is a sample table showing data:

\begin{table}[h]
\centering
\begin{tabular}{@{}lrrr@{}}
\toprule
\textbf{Category} & \textbf{Value 1} & \textbf{Value 2} & \textbf{Total} \\
\midrule
Item A & 10.5 & 20.3 & 30.8 \\
Item B & 15.2 & 25.1 & 40.3 \\
Item C & 22.1 & 12.7 & 34.8 \\
\bottomrule
\end{tabular}
\caption{Sample Data Table}
\label{tab:sample}
\end{table}

\section{Lists}
\subsection{Bullet Points}
\begin{itemize}
  \item First item
  \item Second item
  \item Third item with sub-items:
    \begin{itemize}
      \item Sub-item 1
      \item Sub-item 2
    \end{itemize}
\end{itemize}

\subsection{Numbered List}
\begin{enumerate}
  \item First step
  \item Second step
  \item Third step
\end{enumerate}

\section{Citations and References}
When you use the citation search feature, references will be automatically managed. For example, you might cite a paper like this: \cite{example2024}. All citations will be automatically added to your \texttt{cite.bib} file.

\section{Conclusion}
This template demonstrates basic \LaTeX{} features including mathematical formulas, tables, and lists. The bibliography system is now set up to work seamlessly with Bibby AI's citation management. Feel free to modify it for your needs!

\bibliographystyle{plain}
\bibliography{cite}

\end{document}\documentclass{article}
\usepackage[utf8]{inputenc}
\usepackage{booktabs}
\usepackage{amsmath}
\usepackage{natbib}

\title{My Document}
\author{Your Name}
\date{\today}

\begin{document}

\maketitle

\section{Introduction}
This is a sample document demonstrating various \LaTeX{} features including mathematical equations, tables, and formatting. You can cite papers using the citation search feature, and they will be automatically added to your \texttt{cite.bib} file.

\section{Mathematical Expressions}
Here are some example equations:

\subsection{Quadratic Formula}
The quadratic formula is given by:
\[x = \frac{-b \pm \sqrt{b^2 - 4ac}}{2a}\]

\subsection{Matrix Example}
A sample matrix:
\[
\begin{pmatrix}
a_{11} & a_{12} & a_{13} \\
a_{21} & a_{22} & a_{23} \\
a_{31} & a_{32} & a_{33}
\end{pmatrix}
\]

\section{Tables}
Below is a sample table showing data:

\begin{table}[h]
\centering
\begin{tabular}{@{}lrrr@{}}
\toprule
\textbf{Category} & \textbf{Value 1} & \textbf{Value 2} & \textbf{Total} \\
\midrule
Item A & 10.5 & 20.3 & 30.8 \\
Item B & 15.2 & 25.1 & 40.3 \\
Item C & 22.1 & 12.7 & 34.8 \\
\bottomrule
\end{tabular}
\caption{Sample Data Table}
\label{tab:sample}
\end{table}

\section{Lists}
\subsection{Bullet Points}
\begin{itemize}
  \item First item
  \item Second item
  \item Third item with sub-items:
    \begin{itemize}
      \item Sub-item 1
      \item Sub-item 2
    \end{itemize}
\end{itemize}

\subsection{Numbered List}
\begin{enumerate}
  \item First step
  \item Second step
  \item Third step
\end{enumerate}

\section{Citations and References}
When you use the citation search feature, references will be automatically managed. For example, you might cite a paper like this: \cite{example2024}. All citations will be automatically added to your \texttt{cite.bib} file.

\section{Conclusion}
This template demonstrates basic \LaTeX{} features including mathematical formulas, tables, and lists. The bibliography system is now set up to work seamlessly with Bibby AI's citation management. Feel free to modify it for your needs!

\bibliographystyle{plain}
\bibliography{cite}

\end{document}\documentclass{article}
\usepackage[utf8]{inputenc}
\usepackage{booktabs}
\usepackage{amsmath}
\usepackage{natbib}

\title{My Document}
\author{Your Name}
\date{\today}

\begin{document}

\maketitle

\section{Introduction}
This is a sample document demonstrating various \LaTeX{} features including mathematical equations, tables, and formatting. You can cite papers using the citation search feature, and they will be automatically added to your \texttt{cite.bib} file.

\section{Mathematical Expressions}
Here are some example equations:

\subsection{Quadratic Formula}
The quadratic formula is given by:
\[x = \frac{-b \pm \sqrt{b^2 - 4ac}}{2a}\]

\subsection{Matrix Example}
A sample matrix:
\[
\begin{pmatrix}
a_{11} & a_{12} & a_{13} \\
a_{21} & a_{22} & a_{23} \\
a_{31} & a_{32} & a_{33}
\end{pmatrix}
\]

\section{Tables}
Below is a sample table showing data:

\begin{table}[h]
\centering
\begin{tabular}{@{}lrrr@{}}
\toprule
\textbf{Category} & \textbf{Value 1} & \textbf{Value 2} & \textbf{Total} \\
\midrule
Item A & 10.5 & 20.3 & 30.8 \\
Item B & 15.2 & 25.1 & 40.3 \\
Item C & 22.1 & 12.7 & 34.8 \\
\bottomrule
\end{tabular}
\caption{Sample Data Table}
\label{tab:sample}
\end{table}

\section{Lists}
\subsection{Bullet Points}
\begin{itemize}
  \item First item
  \item Second item
  \item Third item with sub-items:
    \begin{itemize}
      \item Sub-item 1
      \item Sub-item 2
    \end{itemize}
\end{itemize}

\subsection{Numbered List}
\begin{enumerate}
  \item First step
  \item Second step
  \item Third step
\end{enumerate}

\section{Citations and References}
When you use the citation search feature, references will be automatically managed. For example, you might cite a paper like this: \cite{example2024}. All citations will be automatically added to your \texttt{cite.bib} file.

\section{Conclusion}
This template demonstrates basic \LaTeX{} features including mathematical formulas, tables, and lists. The bibliography system is now set up to work seamlessly with Bibby AI's citation management. Feel free to modify it for your needs!

\bibliographystyle{plain}
\bibliography{cite}

\end{document}\documentclass{article}
\usepackage[utf8]{inputenc}
\usepackage{booktabs}
\usepackage{amsmath}
\usepackage{natbib}

\title{My Document}
\author{Your Name}
\date{\today}

\begin{document}

\maketitle

\section{Introduction}
This is a sample document demonstrating various \LaTeX{} features including mathematical equations, tables, and formatting. You can cite papers using the citation search feature, and they will be automatically added to your \texttt{cite.bib} file.

\section{Mathematical Expressions}
Here are some example equations:

\subsection{Quadratic Formula}
The quadratic formula is given by:
\[x = \frac{-b \pm \sqrt{b^2 - 4ac}}{2a}\]

\subsection{Matrix Example}
A sample matrix:
\[
\begin{pmatrix}
a_{11} & a_{12} & a_{13} \\
a_{21} & a_{22} & a_{23} \\
a_{31} & a_{32} & a_{33}
\end{pmatrix}
\]

\section{Tables}
Below is a sample table showing data:

\begin{table}[h]
\centering
\begin{tabular}{@{}lrrr@{}}
\toprule
\textbf{Category} & \textbf{Value 1} & \textbf{Value 2} & \textbf{Total} \\
\midrule
Item A & 10.5 & 20.3 & 30.8 \\
Item B & 15.2 & 25.1 & 40.3 \\
Item C & 22.1 & 12.7 & 34.8 \\
\bottomrule
\end{tabular}
\caption{Sample Data Table}
\label{tab:sample}
\end{table}

\section{Lists}
\subsection{Bullet Points}
\begin{itemize}
  \item First item
  \item Second item
  \item Third item with sub-items:
    \begin{itemize}
      \item Sub-item 1
      \item Sub-item 2
    \end{itemize}
\end{itemize}

\subsection{Numbered List}
\begin{enumerate}
  \item First step
  \item Second step
  \item Third step
\end{enumerate}

\section{Citations and References}
When you use the citation search feature, references will be automatically managed. For example, you might cite a paper like this: \cite{example2024}. All citations will be automatically added to your \texttt{cite.bib} file.

\section{Conclusion}
This template demonstrates basic \LaTeX{} features including mathematical formulas, tables, and lists. The bibliography system is now set up to work seamlessly with Bibby AI's citation management. Feel free to modify it for your needs!

\bibliographystyle{plain}
\bibliography{cite}

\end{document}\documentclass{article}
\usepackage[utf8]{inputenc}
\usepackage{booktabs}
\usepackage{amsmath}
\usepackage{natbib}

\title{My Document}
\author{Your Name}
\date{\today}

\begin{document}

\maketitle

\section{Introduction}
This is a sample document demonstrating various \LaTeX{} features including mathematical equations, tables, and formatting. You can cite papers using the citation search feature, and they will be automatically added to your \texttt{cite.bib} file.

\section{Mathematical Expressions}
Here are some example equations:

\subsection{Quadratic Formula}
The quadratic formula is given by:
\[x = \frac{-b \pm \sqrt{b^2 - 4ac}}{2a}\]

\subsection{Matrix Example}
A sample matrix:
\[
\begin{pmatrix}
a_{11} & a_{12} & a_{13} \\
a_{21} & a_{22} & a_{23} \\
a_{31} & a_{32} & a_{33}
\end{pmatrix}
\]

\section{Tables}
Below is a sample table showing data:

\begin{table}[h]
\centering
\begin{tabular}{@{}lrrr@{}}
\toprule
\textbf{Category} & \textbf{Value 1} & \textbf{Value 2} & \textbf{Total} \\
\midrule
Item A & 10.5 & 20.3 & 30.8 \\
Item B & 15.2 & 25.1 & 40.3 \\
Item C & 22.1 & 12.7 & 34.8 \\
\bottomrule
\end{tabular}
\caption{Sample Data Table}
\label{tab:sample}
\end{table}

\section{Lists}
\subsection{Bullet Points}
\begin{itemize}
  \item First item
  \item Second item
  \item Third item with sub-items:
    \begin{itemize}
      \item Sub-item 1
      \item Sub-item 2
    \end{itemize}
\end{itemize}

\subsection{Numbered List}
\begin{enumerate}
  \item First step
  \item Second step
  \item Third step
\end{enumerate}

\section{Citations and References}
When you use the citation search feature, references will be automatically managed. For example, you might cite a paper like this: \cite{example2024}. All citations will be automatically added to your \texttt{cite.bib} file.

\section{Conclusion}
This template demonstrates basic \LaTeX{} features including mathematical formulas, tables, and lists. The bibliography system is now set up to work seamlessly with Bibby AI's citation management. Feel free to modify it for your needs!

\bibliographystyle{plain}
\bibliography{cite}

\end{document}\documentclass{article}
\usepackage[utf8]{inputenc}
\usepackage{booktabs}
\usepackage{amsmath}
\usepackage{natbib}

\title{My Document}
\author{Your Name}
\date{\today}

\begin{document}

\maketitle

\section{Introduction}
This is a sample document demonstrating various \LaTeX{} features including mathematical equations, tables, and formatting. You can cite papers using the citation search feature, and they will be automatically added to your \texttt{cite.bib} file.

\section{Mathematical Expressions}
Here are some example equations:

\subsection{Quadratic Formula}
The quadratic formula is given by:
\[x = \frac{-b \pm \sqrt{b^2 - 4ac}}{2a}\]

\subsection{Matrix Example}
A sample matrix:
\[
\begin{pmatrix}
a_{11} & a_{12} & a_{13} \\
a_{21} & a_{22} & a_{23} \\
a_{31} & a_{32} & a_{33}
\end{pmatrix}
\]

\section{Tables}
Below is a sample table showing data:

\begin{table}[h]
\centering
\begin{tabular}{@{}lrrr@{}}
\toprule
\textbf{Category} & \textbf{Value 1} & \textbf{Value 2} & \textbf{Total} \\
\midrule
Item A & 10.5 & 20.3 & 30.8 \\
Item B & 15.2 & 25.1 & 40.3 \\
Item C & 22.1 & 12.7 & 34.8 \\
\bottomrule
\end{tabular}
\caption{Sample Data Table}
\label{tab:sample}
\end{table}

\section{Lists}
\subsection{Bullet Points}
\begin{itemize}
  \item First item
  \item Second item
  \item Third item with sub-items:
    \begin{itemize}
      \item Sub-item 1
      \item Sub-item 2
    \end{itemize}
\end{itemize}

\subsection{Numbered List}
\begin{enumerate}
  \item First step
  \item Second step
  \item Third step
\end{enumerate}

\section{Citations and References}
When you use the citation search feature, references will be automatically managed. For example, you might cite a paper like this: \cite{example2024}. All citations will be automatically added to your \texttt{cite.bib} file.

\section{Conclusion}
This template demonstrates basic \LaTeX{} features including mathematical formulas, tables, and lists. The bibliography system is now set up to work seamlessly with Bibby AI's citation management. Feel free to modify it for your needs!

\bibliographystyle{plain}
\bibliography{cite}

\end{document}\documentclass{article}
\usepackage[utf8]{inputenc}
\usepackage{booktabs}
\usepackage{amsmath}
\usepackage{natbib}

\title{My Document}
\author{Your Name}
\date{\today}

\begin{document}

\maketitle

\section{Introduction}
This is a sample document demonstrating various \LaTeX{} features including mathematical equations, tables, and formatting. You can cite papers using the citation search feature, and they will be automatically added to your \texttt{cite.bib} file.

\section{Mathematical Expressions}
Here are some example equations:

\subsection{Quadratic Formula}
The quadratic formula is given by:
\[x = \frac{-b \pm \sqrt{b^2 - 4ac}}{2a}\]

\subsection{Matrix Example}
A sample matrix:
\[
\begin{pmatrix}
a_{11} & a_{12} & a_{13} \\
a_{21} & a_{22} & a_{23} \\
a_{31} & a_{32} & a_{33}
\end{pmatrix}
\]

\section{Tables}
Below is a sample table showing data:

\begin{table}[h]
\centering
\begin{tabular}{@{}lrrr@{}}
\toprule
\textbf{Category} & \textbf{Value 1} & \textbf{Value 2} & \textbf{Total} \\
\midrule
Item A & 10.5 & 20.3 & 30.8 \\
Item B & 15.2 & 25.1 & 40.3 \\
Item C & 22.1 & 12.7 & 34.8 \\
\bottomrule
\end{tabular}
\caption{Sample Data Table}
\label{tab:sample}
\end{table}

\section{Lists}
\subsection{Bullet Points}
\begin{itemize}
  \item First item
  \item Second item
  \item Third item with sub-items:
    \begin{itemize}
      \item Sub-item 1
      \item Sub-item 2
    \end{itemize}
\end{itemize}

\subsection{Numbered List}
\begin{enumerate}
  \item First step
  \item Second step
  \item Third step
\end{enumerate}

\section{Citations and References}
When you use the citation search feature, references will be automatically managed. For example, you might cite a paper like this: \cite{example2024}. All citations will be automatically added to your \texttt{cite.bib} file.

\section{Conclusion}
This template demonstrates basic \LaTeX{} features including mathematical formulas, tables, and lists. The bibliography system is now set up to work seamlessly with Bibby AI's citation management. Feel free to modify it for your needs!

\bibliographystyle{plain}
\bibliography{cite}

\end{document}\documentclass{article}
\usepackage[utf8]{inputenc}
\usepackage{booktabs}
\usepackage{amsmath}
\usepackage{natbib}

\title{My Document}
\author{Your Name}
\date{\today}

\begin{document}

\maketitle

\section{Introduction}
This is a sample document demonstrating various \LaTeX{} features including mathematical equations, tables, and formatting. You can cite papers using the citation search feature, and they will be automatically added to your \texttt{cite.bib} file.

\section{Mathematical Expressions}
Here are some example equations:

\subsection{Quadratic Formula}
The quadratic formula is given by:
\[x = \frac{-b \pm \sqrt{b^2 - 4ac}}{2a}\]

\subsection{Matrix Example}
A sample matrix:
\[
\begin{pmatrix}
a_{11} & a_{12} & a_{13} \\
a_{21} & a_{22} & a_{23} \\
a_{31} & a_{32} & a_{33}
\end{pmatrix}
\]

\section{Tables}
Below is a sample table showing data:

\begin{table}[h]
\centering
\begin{tabular}{@{}lrrr@{}}
\toprule
\textbf{Category} & \textbf{Value 1} & \textbf{Value 2} & \textbf{Total} \\
\midrule
Item A & 10.5 & 20.3 & 30.8 \\
Item B & 15.2 & 25.1 & 40.3 \\
Item C & 22.1 & 12.7 & 34.8 \\
\bottomrule
\end{tabular}
\caption{Sample Data Table}
\label{tab:sample}
\end{table}

\section{Lists}
\subsection{Bullet Points}
\begin{itemize}
  \item First item
  \item Second item
  \item Third item with sub-items:
    \begin{itemize}
      \item Sub-item 1
      \item Sub-item 2
    \end{itemize}
\end{itemize}

\subsection{Numbered List}
\begin{enumerate}
  \item First step
  \item Second step
  \item Third step
\end{enumerate}

\section{Citations and References}
When you use the citation search feature, references will be automatically managed. For example, you might cite a paper like this: \cite{example2024}. All citations will be automatically added to your \texttt{cite.bib} file.

\section{Conclusion}
This template demonstrates basic \LaTeX{} features including mathematical formulas, tables, and lists. The bibliography system is now set up to work seamlessly with Bibby AI's citation management. Feel free to modify it for your needs!

\bibliographystyle{plain}
\bibliography{cite}

\end{document}\documentclass{article}
\usepackage[utf8]{inputenc}
\usepackage{booktabs}
\usepackage{amsmath}
\usepackage{natbib}

\title{My Document}
\author{Your Name}
\date{\today}

\begin{document}

\maketitle

\section{Introduction}
This is a sample document demonstrating various \LaTeX{} features including mathematical equations, tables, and formatting. You can cite papers using the citation search feature, and they will be automatically added to your \texttt{cite.bib} file.

\section{Mathematical Expressions}
Here are some example equations:

\subsection{Quadratic Formula}
The quadratic formula is given by:
\[x = \frac{-b \pm \sqrt{b^2 - 4ac}}{2a}\]

\subsection{Matrix Example}
A sample matrix:
\[
\begin{pmatrix}
a_{11} & a_{12} & a_{13} \\
a_{21} & a_{22} & a_{23} \\
a_{31} & a_{32} & a_{33}
\end{pmatrix}
\]

\section{Tables}
Below is a sample table showing data:

\begin{table}[h]
\centering
\begin{tabular}{@{}lrrr@{}}
\toprule
\textbf{Category} & \textbf{Value 1} & \textbf{Value 2} & \textbf{Total} \\
\midrule
Item A & 10.5 & 20.3 & 30.8 \\
Item B & 15.2 & 25.1 & 40.3 \\
Item C & 22.1 & 12.7 & 34.8 \\
\bottomrule
\end{tabular}
\caption{Sample Data Table}
\label{tab:sample}
\end{table}

\section{Lists}
\subsection{Bullet Points}
\begin{itemize}
  \item First item
  \item Second item
  \item Third item with sub-items:
    \begin{itemize}
      \item Sub-item 1
      \item Sub-item 2
    \end{itemize}
\end{itemize}

\subsection{Numbered List}
\begin{enumerate}
  \item First step
  \item Second step
  \item Third step
\end{enumerate}

\section{Citations and References}
When you use the citation search feature, references will be automatically managed. For example, you might cite a paper like this: \cite{example2024}. All citations will be automatically added to your \texttt{cite.bib} file.

\section{Conclusion}
This template demonstrates basic \LaTeX{} features including mathematical formulas, tables, and lists. The bibliography system is now set up to work seamlessly with Bibby AI's citation management. Feel free to modify it for your needs!

\bibliographystyle{plain}
\bibliography{cite}

\end{document}\documentclass{article}
\usepackage[utf8]{inputenc}
\usepackage{booktabs}
\usepackage{amsmath}
\usepackage{natbib}

\title{My Document}
\author{Your Name}
\date{\today}

\begin{document}

\maketitle

\section{Introduction}
This is a sample document demonstrating various \LaTeX{} features including mathematical equations, tables, and formatting. You can cite papers using the citation search feature, and they will be automatically added to your \texttt{cite.bib} file.

\section{Mathematical Expressions}
Here are some example equations:

\subsection{Quadratic Formula}
The quadratic formula is given by:
\[x = \frac{-b \pm \sqrt{b^2 - 4ac}}{2a}\]

\subsection{Matrix Example}
A sample matrix:
\[
\begin{pmatrix}
a_{11} & a_{12} & a_{13} \\
a_{21} & a_{22} & a_{23} \\
a_{31} & a_{32} & a_{33}
\end{pmatrix}
\]

\section{Tables}
Below is a sample table showing data:

\begin{table}[h]
\centering
\begin{tabular}{@{}lrrr@{}}
\toprule
\textbf{Category} & \textbf{Value 1} & \textbf{Value 2} & \textbf{Total} \\
\midrule
Item A & 10.5 & 20.3 & 30.8 \\
Item B & 15.2 & 25.1 & 40.3 \\
Item C & 22.1 & 12.7 & 34.8 \\
\bottomrule
\end{tabular}
\caption{Sample Data Table}
\label{tab:sample}
\end{table}

\section{Lists}
\subsection{Bullet Points}
\begin{itemize}
  \item First item
  \item Second item
  \item Third item with sub-items:
    \begin{itemize}
      \item Sub-item 1
      \item Sub-item 2
    \end{itemize}
\end{itemize}

\subsection{Numbered List}
\begin{enumerate}
  \item First step
  \item Second step
  \item Third step
\end{enumerate}

\section{Citations and References}
When you use the citation search feature, references will be automatically managed. For example, you might cite a paper like this: \cite{example2024}. All citations will be automatically added to your \texttt{cite.bib} file.

\section{Conclusion}
This template demonstrates basic \LaTeX{} features including mathematical formulas, tables, and lists. The bibliography system is now set up to work seamlessly with Bibby AI's citation management. Feel free to modify it for your needs!

\bibliographystyle{plain}
\bibliography{cite}

\end{document}\documentclass{article}
\usepackage[utf8]{inputenc}
\usepackage{booktabs}
\usepackage{amsmath}
\usepackage{natbib}

\title{My Document}
\author{Your Name}
\date{\today}

\begin{document}

\maketitle

\section{Introduction}
This is a sample document demonstrating various \LaTeX{} features including mathematical equations, tables, and formatting. You can cite papers using the citation search feature, and they will be automatically added to your \texttt{cite.bib} file.

\section{Mathematical Expressions}
Here are some example equations:

\subsection{Quadratic Formula}
The quadratic formula is given by:
\[x = \frac{-b \pm \sqrt{b^2 - 4ac}}{2a}\]

\subsection{Matrix Example}
A sample matrix:
\[
\begin{pmatrix}
a_{11} & a_{12} & a_{13} \\
a_{21} & a_{22} & a_{23} \\
a_{31} & a_{32} & a_{33}
\end{pmatrix}
\]

\section{Tables}
Below is a sample table showing data:

\begin{table}[h]
\centering
\begin{tabular}{@{}lrrr@{}}
\toprule
\textbf{Category} & \textbf{Value 1} & \textbf{Value 2} & \textbf{Total} \\
\midrule
Item A & 10.5 & 20.3 & 30.8 \\
Item B & 15.2 & 25.1 & 40.3 \\
Item C & 22.1 & 12.7 & 34.8 \\
\bottomrule
\end{tabular}
\caption{Sample Data Table}
\label{tab:sample}
\end{table}

\section{Lists}
\subsection{Bullet Points}
\begin{itemize}
  \item First item
  \item Second item
  \item Third item with sub-items:
    \begin{itemize}
      \item Sub-item 1
      \item Sub-item 2
    \end{itemize}
\end{itemize}

\subsection{Numbered List}
\begin{enumerate}
  \item First step
  \item Second step
  \item Third step
\end{enumerate}

\section{Citations and References}
When you use the citation search feature, references will be automatically managed. For example, you might cite a paper like this: \cite{example2024}. All citations will be automatically added to your \texttt{cite.bib} file.

\section{Conclusion}
This template demonstrates basic \LaTeX{} features including mathematical formulas, tables, and lists. The bibliography system is now set up to work seamlessly with Bibby AI's citation management. Feel free to modify it for your needs!

\bibliographystyle{plain}
\bibliography{cite}

\end{document}\documentclass{article}
\usepackage[utf8]{inputenc}
\usepackage{booktabs}
\usepackage{amsmath}
\usepackage{natbib}

\title{My Document}
\author{Your Name}
\date{\today}

\begin{document}

\maketitle

\section{Introduction}
This is a sample document demonstrating various \LaTeX{} features including mathematical equations, tables, and formatting. You can cite papers using the citation search feature, and they will be automatically added to your \texttt{cite.bib} file.

\section{Mathematical Expressions}
Here are some example equations:

\subsection{Quadratic Formula}
The quadratic formula is given by:
\[x = \frac{-b \pm \sqrt{b^2 - 4ac}}{2a}\]

\subsection{Matrix Example}
A sample matrix:
\[
\begin{pmatrix}
a_{11} & a_{12} & a_{13} \\
a_{21} & a_{22} & a_{23} \\
a_{31} & a_{32} & a_{33}
\end{pmatrix}
\]

\section{Tables}
Below is a sample table showing data:

\begin{table}[h]
\centering
\begin{tabular}{@{}lrrr@{}}
\toprule
\textbf{Category} & \textbf{Value 1} & \textbf{Value 2} & \textbf{Total} \\
\midrule
Item A & 10.5 & 20.3 & 30.8 \\
Item B & 15.2 & 25.1 & 40.3 \\
Item C & 22.1 & 12.7 & 34.8 \\
\bottomrule
\end{tabular}
\caption{Sample Data Table}
\label{tab:sample}
\end{table}

\section{Lists}
\subsection{Bullet Points}
\begin{itemize}
  \item First item
  \item Second item
  \item Third item with sub-items:
    \begin{itemize}
      \item Sub-item 1
      \item Sub-item 2
    \end{itemize}
\end{itemize}

\subsection{Numbered List}
\begin{enumerate}
  \item First step
  \item Second step
  \item Third step
\end{enumerate}

\section{Citations and References}
When you use the citation search feature, references will be automatically managed. For example, you might cite a paper like this: \cite{example2024}. All citations will be automatically added to your \texttt{cite.bib} file.

\section{Conclusion}
This template demonstrates basic \LaTeX{} features including mathematical formulas, tables, and lists. The bibliography system is now set up to work seamlessly with Bibby AI's citation management. Feel free to modify it for your needs!

\bibliographystyle{plain}
\bibliography{cite}

\end{document}\documentclass{article}
\usepackage[utf8]{inputenc}
\usepackage{booktabs}
\usepackage{amsmath}
\usepackage{natbib}

\title{My Document}
\author{Your Name}
\date{\today}

\begin{document}

\maketitle

\section{Introduction}
This is a sample document demonstrating various \LaTeX{} features including mathematical equations, tables, and formatting. You can cite papers using the citation search feature, and they will be automatically added to your \texttt{cite.bib} file.

\section{Mathematical Expressions}
Here are some example equations:

\subsection{Quadratic Formula}
The quadratic formula is given by:
\[x = \frac{-b \pm \sqrt{b^2 - 4ac}}{2a}\]

\subsection{Matrix Example}
A sample matrix:
\[
\begin{pmatrix}
a_{11} & a_{12} & a_{13} \\
a_{21} & a_{22} & a_{23} \\
a_{31} & a_{32} & a_{33}
\end{pmatrix}
\]

\section{Tables}
Below is a sample table showing data:

\begin{table}[h]
\centering
\begin{tabular}{@{}lrrr@{}}
\toprule
\textbf{Category} & \textbf{Value 1} & \textbf{Value 2} & \textbf{Total} \\
\midrule
Item A & 10.5 & 20.3 & 30.8 \\
Item B & 15.2 & 25.1 & 40.3 \\
Item C & 22.1 & 12.7 & 34.8 \\
\bottomrule
\end{tabular}
\caption{Sample Data Table}
\label{tab:sample}
\end{table}

\section{Lists}
\subsection{Bullet Points}
\begin{itemize}
  \item First item
  \item Second item
  \item Third item with sub-items:
    \begin{itemize}
      \item Sub-item 1
      \item Sub-item 2
    \end{itemize}
\end{itemize}

\subsection{Numbered List}
\begin{enumerate}
  \item First step
  \item Second step
  \item Third step
\end{enumerate}

\section{Citations and References}
When you use the citation search feature, references will be automatically managed. For example, you might cite a paper like this: \cite{example2024}. All citations will be automatically added to your \texttt{cite.bib} file.

\section{Conclusion}
This template demonstrates basic \LaTeX{} features including mathematical formulas, tables, and lists. The bibliography system is now set up to work seamlessly with Bibby AI's citation management. Feel free to modify it for your needs!

\bibliographystyle{plain}
\bibliography{cite}

\end{document}\documentclass{article}
\usepackage[utf8]{inputenc}
\usepackage{booktabs}
\usepackage{amsmath}
\usepackage{natbib}

\title{My Document}
\author{Your Name}
\date{\today}

\begin{document}

\maketitle

\section{Introduction}
This is a sample document demonstrating various \LaTeX{} features including mathematical equations, tables, and formatting. You can cite papers using the citation search feature, and they will be automatically added to your \texttt{cite.bib} file.

\section{Mathematical Expressions}
Here are some example equations:

\subsection{Quadratic Formula}
The quadratic formula is given by:
\[x = \frac{-b \pm \sqrt{b^2 - 4ac}}{2a}\]

\subsection{Matrix Example}
A sample matrix:
\[
\begin{pmatrix}
a_{11} & a_{12} & a_{13} \\
a_{21} & a_{22} & a_{23} \\
a_{31} & a_{32} & a_{33}
\end{pmatrix}
\]

\section{Tables}
Below is a sample table showing data:

\begin{table}[h]
\centering
\begin{tabular}{@{}lrrr@{}}
\toprule
\textbf{Category} & \textbf{Value 1} & \textbf{Value 2} & \textbf{Total} \\
\midrule
Item A & 10.5 & 20.3 & 30.8 \\
Item B & 15.2 & 25.1 & 40.3 \\
Item C & 22.1 & 12.7 & 34.8 \\
\bottomrule
\end{tabular}
\caption{Sample Data Table}
\label{tab:sample}
\end{table}

\section{Lists}
\subsection{Bullet Points}
\begin{itemize}
  \item First item
  \item Second item
  \item Third item with sub-items:
    \begin{itemize}
      \item Sub-item 1
      \item Sub-item 2
    \end{itemize}
\end{itemize}

\subsection{Numbered List}
\begin{enumerate}
  \item First step
  \item Second step
  \item Third step
\end{enumerate}

\section{Citations and References}
When you use the citation search feature, references will be automatically managed. For example, you might cite a paper like this: \cite{example2024}. All citations will be automatically added to your \texttt{cite.bib} file.

\section{Conclusion}
This template demonstrates basic \LaTeX{} features including mathematical formulas, tables, and lists. The bibliography system is now set up to work seamlessly with Bibby AI's citation management. Feel free to modify it for your needs!

\bibliographystyle{plain}
\bibliography{cite}

\end{document}\documentclass{article}
\usepackage[utf8]{inputenc}
\usepackage{booktabs}
\usepackage{amsmath}
\usepackage{natbib}

\title{My Document}
\author{Your Name}
\date{\today}

\begin{document}

\maketitle

\section{Introduction}
This is a sample document demonstrating various \LaTeX{} features including mathematical equations, tables, and formatting. You can cite papers using the citation search feature, and they will be automatically added to your \texttt{cite.bib} file.

\section{Mathematical Expressions}
Here are some example equations:

\subsection{Quadratic Formula}
The quadratic formula is given by:
\[x = \frac{-b \pm \sqrt{b^2 - 4ac}}{2a}\]

\subsection{Matrix Example}
A sample matrix:
\[
\begin{pmatrix}
a_{11} & a_{12} & a_{13} \\
a_{21} & a_{22} & a_{23} \\
a_{31} & a_{32} & a_{33}
\end{pmatrix}
\]

\section{Tables}
Below is a sample table showing data:

\begin{table}[h]
\centering
\begin{tabular}{@{}lrrr@{}}
\toprule
\textbf{Category} & \textbf{Value 1} & \textbf{Value 2} & \textbf{Total} \\
\midrule
Item A & 10.5 & 20.3 & 30.8 \\
Item B & 15.2 & 25.1 & 40.3 \\
Item C & 22.1 & 12.7 & 34.8 \\
\bottomrule
\end{tabular}
\caption{Sample Data Table}
\label{tab:sample}
\end{table}

\section{Lists}
\subsection{Bullet Points}
\begin{itemize}
  \item First item
  \item Second item
  \item Third item with sub-items:
    \begin{itemize}
      \item Sub-item 1
      \item Sub-item 2
    \end{itemize}
\end{itemize}

\subsection{Numbered List}
\begin{enumerate}
  \item First step
  \item Second step
  \item Third step
\end{enumerate}

\section{Citations and References}
When you use the citation search feature, references will be automatically managed. For example, you might cite a paper like this: \cite{example2024}. All citations will be automatically added to your \texttt{cite.bib} file.

\section{Conclusion}
This template demonstrates basic \LaTeX{} features including mathematical formulas, tables, and lists. The bibliography system is now set up to work seamlessly with Bibby AI's citation management. Feel free to modify it for your needs!

\bibliographystyle{plain}
\bibliography{cite}

\end{document}\documentclass{article}
\usepackage[utf8]{inputenc}
\usepackage{booktabs}
\usepackage{amsmath}
\usepackage{natbib}

\title{My Document}
\author{Your Name}
\date{\today}

\begin{document}

\maketitle

\section{Introduction}
This is a sample document demonstrating various \LaTeX{} features including mathematical equations, tables, and formatting. You can cite papers using the citation search feature, and they will be automatically added to your \texttt{cite.bib} file.

\section{Mathematical Expressions}
Here are some example equations:

\subsection{Quadratic Formula}
The quadratic formula is given by:
\[x = \frac{-b \pm \sqrt{b^2 - 4ac}}{2a}\]

\subsection{Matrix Example}
A sample matrix:
\[
\begin{pmatrix}
a_{11} & a_{12} & a_{13} \\
a_{21} & a_{22} & a_{23} \\
a_{31} & a_{32} & a_{33}
\end{pmatrix}
\]

\section{Tables}
Below is a sample table showing data:

\begin{table}[h]
\centering
\begin{tabular}{@{}lrrr@{}}
\toprule
\textbf{Category} & \textbf{Value 1} & \textbf{Value 2} & \textbf{Total} \\
\midrule
Item A & 10.5 & 20.3 & 30.8 \\
Item B & 15.2 & 25.1 & 40.3 \\
Item C & 22.1 & 12.7 & 34.8 \\
\bottomrule
\end{tabular}
\caption{Sample Data Table}
\label{tab:sample}
\end{table}

\section{Lists}
\subsection{Bullet Points}
\begin{itemize}
  \item First item
  \item Second item
  \item Third item with sub-items:
    \begin{itemize}
      \item Sub-item 1
      \item Sub-item 2
    \end{itemize}
\end{itemize}

\subsection{Numbered List}
\begin{enumerate}
  \item First step
  \item Second step
  \item Third step
\end{enumerate}

\section{Citations and References}
When you use the citation search feature, references will be automatically managed. For example, you might cite a paper like this: \cite{example2024}. All citations will be automatically added to your \texttt{cite.bib} file.

\section{Conclusion}
This template demonstrates basic \LaTeX{} features including mathematical formulas, tables, and lists. The bibliography system is now set up to work seamlessly with Bibby AI's citation management. Feel free to modify it for your needs!

\bibliographystyle{plain}
\bibliography{cite}

\end{document}\documentclass{article}
\usepackage[utf8]{inputenc}
\usepackage{booktabs}
\usepackage{amsmath}
\usepackage{natbib}

\title{My Document}
\author{Your Name}
\date{\today}

\begin{document}

\maketitle

\section{Introduction}
This is a sample document demonstrating various \LaTeX{} features including mathematical equations, tables, and formatting. You can cite papers using the citation search feature, and they will be automatically added to your \texttt{cite.bib} file.

\section{Mathematical Expressions}
Here are some example equations:

\subsection{Quadratic Formula}
The quadratic formula is given by:
\[x = \frac{-b \pm \sqrt{b^2 - 4ac}}{2a}\]

\subsection{Matrix Example}
A sample matrix:
\[
\begin{pmatrix}
a_{11} & a_{12} & a_{13} \\
a_{21} & a_{22} & a_{23} \\
a_{31} & a_{32} & a_{33}
\end{pmatrix}
\]

\section{Tables}
Below is a sample table showing data:

\begin{table}[h]
\centering
\begin{tabular}{@{}lrrr@{}}
\toprule
\textbf{Category} & \textbf{Value 1} & \textbf{Value 2} & \textbf{Total} \\
\midrule
Item A & 10.5 & 20.3 & 30.8 \\
Item B & 15.2 & 25.1 & 40.3 \\
Item C & 22.1 & 12.7 & 34.8 \\
\bottomrule
\end{tabular}
\caption{Sample Data Table}
\label{tab:sample}
\end{table}

\section{Lists}
\subsection{Bullet Points}
\begin{itemize}
  \item First item
  \item Second item
  \item Third item with sub-items:
    \begin{itemize}
      \item Sub-item 1
      \item Sub-item 2
    \end{itemize}
\end{itemize}

\subsection{Numbered List}
\begin{enumerate}
  \item First step
  \item Second step
  \item Third step
\end{enumerate}

\section{Citations and References}
When you use the citation search feature, references will be automatically managed. For example, you might cite a paper like this: \cite{example2024}. All citations will be automatically added to your \texttt{cite.bib} file.

\section{Conclusion}
This template demonstrates basic \LaTeX{} features including mathematical formulas, tables, and lists. The bibliography system is now set up to work seamlessly with Bibby AI's citation management. Feel free to modify it for your needs!

\bibliographystyle{plain}
\bibliography{cite}

\end{document}\documentclass{article}
\usepackage[utf8]{inputenc}
\usepackage{booktabs}
\usepackage{amsmath}
\usepackage{natbib}

\title{My Document}
\author{Your Name}
\date{\today}

\begin{document}

\maketitle

\section{Introduction}
This is a sample document demonstrating various \LaTeX{} features including mathematical equations, tables, and formatting. You can cite papers using the citation search feature, and they will be automatically added to your \texttt{cite.bib} file.

\section{Mathematical Expressions}
Here are some example equations:

\subsection{Quadratic Formula}
The quadratic formula is given by:
\[x = \frac{-b \pm \sqrt{b^2 - 4ac}}{2a}\]

\subsection{Matrix Example}
A sample matrix:
\[
\begin{pmatrix}
a_{11} & a_{12} & a_{13} \\
a_{21} & a_{22} & a_{23} \\
a_{31} & a_{32} & a_{33}
\end{pmatrix}
\]

\section{Tables}
Below is a sample table showing data:

\begin{table}[h]
\centering
\begin{tabular}{@{}lrrr@{}}
\toprule
\textbf{Category} & \textbf{Value 1} & \textbf{Value 2} & \textbf{Total} \\
\midrule
Item A & 10.5 & 20.3 & 30.8 \\
Item B & 15.2 & 25.1 & 40.3 \\
Item C & 22.1 & 12.7 & 34.8 \\
\bottomrule
\end{tabular}
\caption{Sample Data Table}
\label{tab:sample}
\end{table}

\section{Lists}
\subsection{Bullet Points}
\begin{itemize}
  \item First item
  \item Second item
  \item Third item with sub-items:
    \begin{itemize}
      \item Sub-item 1
      \item Sub-item 2
    \end{itemize}
\end{itemize}

\subsection{Numbered List}
\begin{enumerate}
  \item First step
  \item Second step
  \item Third step
\end{enumerate}

\section{Citations and References}
When you use the citation search feature, references will be automatically managed. For example, you might cite a paper like this: \cite{example2024}. All citations will be automatically added to your \texttt{cite.bib} file.

\section{Conclusion}
This template demonstrates basic \LaTeX{} features including mathematical formulas, tables, and lists. The bibliography system is now set up to work seamlessly with Bibby AI's citation management. Feel free to modify it for your needs!

\bibliographystyle{plain}
\bibliography{cite}

\end{document}\documentclass{article}
\usepackage[utf8]{inputenc}
\usepackage{booktabs}
\usepackage{amsmath}
\usepackage{natbib}

\title{My Document}
\author{Your Name}
\date{\today}

\begin{document}

\maketitle

\section{Introduction}
This is a sample document demonstrating various \LaTeX{} features including mathematical equations, tables, and formatting. You can cite papers using the citation search feature, and they will be automatically added to your \texttt{cite.bib} file.

\section{Mathematical Expressions}
Here are some example equations:

\subsection{Quadratic Formula}
The quadratic formula is given by:
\[x = \frac{-b \pm \sqrt{b^2 - 4ac}}{2a}\]

\subsection{Matrix Example}
A sample matrix:
\[
\begin{pmatrix}
a_{11} & a_{12} & a_{13} \\
a_{21} & a_{22} & a_{23} \\
a_{31} & a_{32} & a_{33}
\end{pmatrix}
\]

\section{Tables}
Below is a sample table showing data:

\begin{table}[h]
\centering
\begin{tabular}{@{}lrrr@{}}
\toprule
\textbf{Category} & \textbf{Value 1} & \textbf{Value 2} & \textbf{Total} \\
\midrule
Item A & 10.5 & 20.3 & 30.8 \\
Item B & 15.2 & 25.1 & 40.3 \\
Item C & 22.1 & 12.7 & 34.8 \\
\bottomrule
\end{tabular}
\caption{Sample Data Table}
\label{tab:sample}
\end{table}

\section{Lists}
\subsection{Bullet Points}
\begin{itemize}
  \item First item
  \item Second item
  \item Third item with sub-items:
    \begin{itemize}
      \item Sub-item 1
      \item Sub-item 2
    \end{itemize}
\end{itemize}

\subsection{Numbered List}
\begin{enumerate}
  \item First step
  \item Second step
  \item Third step
\end{enumerate}

\section{Citations and References}
When you use the citation search feature, references will be automatically managed. For example, you might cite a paper like this: \cite{example2024}. All citations will be automatically added to your \texttt{cite.bib} file.

\section{Conclusion}
This template demonstrates basic \LaTeX{} features including mathematical formulas, tables, and lists. The bibliography system is now set up to work seamlessly with Bibby AI's citation management. Feel free to modify it for your needs!

\bibliographystyle{plain}
\bibliography{cite}

\end{document}\documentclass{article}
\usepackage[utf8]{inputenc}
\usepackage{booktabs}
\usepackage{amsmath}
\usepackage{natbib}

\title{My Document}
\author{Your Name}
\date{\today}

\begin{document}

\maketitle

\section{Introduction}
This is a sample document demonstrating various \LaTeX{} features including mathematical equations, tables, and formatting. You can cite papers using the citation search feature, and they will be automatically added to your \texttt{cite.bib} file.

\section{Mathematical Expressions}
Here are some example equations:

\subsection{Quadratic Formula}
The quadratic formula is given by:
\[x = \frac{-b \pm \sqrt{b^2 - 4ac}}{2a}\]

\subsection{Matrix Example}
A sample matrix:
\[
\begin{pmatrix}
a_{11} & a_{12} & a_{13} \\
a_{21} & a_{22} & a_{23} \\
a_{31} & a_{32} & a_{33}
\end{pmatrix}
\]

\section{Tables}
Below is a sample table showing data:

\begin{table}[h]
\centering
\begin{tabular}{@{}lrrr@{}}
\toprule
\textbf{Category} & \textbf{Value 1} & \textbf{Value 2} & \textbf{Total} \\
\midrule
Item A & 10.5 & 20.3 & 30.8 \\
Item B & 15.2 & 25.1 & 40.3 \\
Item C & 22.1 & 12.7 & 34.8 \\
\bottomrule
\end{tabular}
\caption{Sample Data Table}
\label{tab:sample}
\end{table}

\section{Lists}
\subsection{Bullet Points}
\begin{itemize}
  \item First item
  \item Second item
  \item Third item with sub-items:
    \begin{itemize}
      \item Sub-item 1
      \item Sub-item 2
    \end{itemize}
\end{itemize}

\subsection{Numbered List}
\begin{enumerate}
  \item First step
  \item Second step
  \item Third step
\end{enumerate}

\section{Citations and References}
When you use the citation search feature, references will be automatically managed. For example, you might cite a paper like this: \cite{example2024}. All citations will be automatically added to your \texttt{cite.bib} file.

\section{Conclusion}
This template demonstrates basic \LaTeX{} features including mathematical formulas, tables, and lists. The bibliography system is now set up to work seamlessly with Bibby AI's citation management. Feel free to modify it for your needs!

\bibliographystyle{plain}
\bibliography{cite}

\end{document}\documentclass{article}
\usepackage[utf8]{inputenc}
\usepackage{booktabs}
\usepackage{amsmath}
\usepackage{natbib}

\title{My Document}
\author{Your Name}
\date{\today}

\begin{document}

\maketitle

\section{Introduction}
This is a sample document demonstrating various \LaTeX{} features including mathematical equations, tables, and formatting. You can cite papers using the citation search feature, and they will be automatically added to your \texttt{cite.bib} file.

\section{Mathematical Expressions}
Here are some example equations:

\subsection{Quadratic Formula}
The quadratic formula is given by:
\[x = \frac{-b \pm \sqrt{b^2 - 4ac}}{2a}\]

\subsection{Matrix Example}
A sample matrix:
\[
\begin{pmatrix}
a_{11} & a_{12} & a_{13} \\
a_{21} & a_{22} & a_{23} \\
a_{31} & a_{32} & a_{33}
\end{pmatrix}
\]

\section{Tables}
Below is a sample table showing data:

\begin{table}[h]
\centering
\begin{tabular}{@{}lrrr@{}}
\toprule
\textbf{Category} & \textbf{Value 1} & \textbf{Value 2} & \textbf{Total} \\
\midrule
Item A & 10.5 & 20.3 & 30.8 \\
Item B & 15.2 & 25.1 & 40.3 \\
Item C & 22.1 & 12.7 & 34.8 \\
\bottomrule
\end{tabular}
\caption{Sample Data Table}
\label{tab:sample}
\end{table}

\section{Lists}
\subsection{Bullet Points}
\begin{itemize}
  \item First item
  \item Second item
  \item Third item with sub-items:
    \begin{itemize}
      \item Sub-item 1
      \item Sub-item 2
    \end{itemize}
\end{itemize}

\subsection{Numbered List}
\begin{enumerate}
  \item First step
  \item Second step
  \item Third step
\end{enumerate}

\section{Citations and References}
When you use the citation search feature, references will be automatically managed. For example, you might cite a paper like this: \cite{example2024}. All citations will be automatically added to your \texttt{cite.bib} file.

\section{Conclusion}
This template demonstrates basic \LaTeX{} features including mathematical formulas, tables, and lists. The bibliography system is now set up to work seamlessly with Bibby AI's citation management. Feel free to modify it for your needs!

\bibliographystyle{plain}
\bibliography{cite}

\end{document}\documentclass{article}
\usepackage[utf8]{inputenc}
\usepackage{booktabs}
\usepackage{amsmath}
\usepackage{natbib}

\title{My Document}
\author{Your Name}
\date{\today}

\begin{document}

\maketitle

\section{Introduction}
This is a sample document demonstrating various \LaTeX{} features including mathematical equations, tables, and formatting. You can cite papers using the citation search feature, and they will be automatically added to your \texttt{cite.bib} file.

\section{Mathematical Expressions}
Here are some example equations:

\subsection{Quadratic Formula}
The quadratic formula is given by:
\[x = \frac{-b \pm \sqrt{b^2 - 4ac}}{2a}\]

\subsection{Matrix Example}
A sample matrix:
\[
\begin{pmatrix}
a_{11} & a_{12} & a_{13} \\
a_{21} & a_{22} & a_{23} \\
a_{31} & a_{32} & a_{33}
\end{pmatrix}
\]

\section{Tables}
Below is a sample table showing data:

\begin{table}[h]
\centering
\begin{tabular}{@{}lrrr@{}}
\toprule
\textbf{Category} & \textbf{Value 1} & \textbf{Value 2} & \textbf{Total} \\
\midrule
Item A & 10.5 & 20.3 & 30.8 \\
Item B & 15.2 & 25.1 & 40.3 \\
Item C & 22.1 & 12.7 & 34.8 \\
\bottomrule
\end{tabular}
\caption{Sample Data Table}
\label{tab:sample}
\end{table}

\section{Lists}
\subsection{Bullet Points}
\begin{itemize}
  \item First item
  \item Second item
  \item Third item with sub-items:
    \begin{itemize}
      \item Sub-item 1
      \item Sub-item 2
    \end{itemize}
\end{itemize}

\subsection{Numbered List}
\begin{enumerate}
  \item First step
  \item Second step
  \item Third step
\end{enumerate}

\section{Citations and References}
When you use the citation search feature, references will be automatically managed. For example, you might cite a paper like this: \cite{example2024}. All citations will be automatically added to your \texttt{cite.bib} file.

\section{Conclusion}
This template demonstrates basic \LaTeX{} features including mathematical formulas, tables, and lists. The bibliography system is now set up to work seamlessly with Bibby AI's citation management. Feel free to modify it for your needs!

\bibliographystyle{plain}
\bibliography{cite}

\end{document}\documentclass{article}
\usepackage[utf8]{inputenc}
\usepackage{booktabs}
\usepackage{amsmath}
\usepackage{natbib}

\title{My Document}
\author{Your Name}
\date{\today}

\begin{document}

\maketitle

\section{Introduction}
This is a sample document demonstrating various \LaTeX{} features including mathematical equations, tables, and formatting. You can cite papers using the citation search feature, and they will be automatically added to your \texttt{cite.bib} file.

\section{Mathematical Expressions}
Here are some example equations:

\subsection{Quadratic Formula}
The quadratic formula is given by:
\[x = \frac{-b \pm \sqrt{b^2 - 4ac}}{2a}\]

\subsection{Matrix Example}
A sample matrix:
\[
\begin{pmatrix}
a_{11} & a_{12} & a_{13} \\
a_{21} & a_{22} & a_{23} \\
a_{31} & a_{32} & a_{33}
\end{pmatrix}
\]

\section{Tables}
Below is a sample table showing data:

\begin{table}[h]
\centering
\begin{tabular}{@{}lrrr@{}}
\toprule
\textbf{Category} & \textbf{Value 1} & \textbf{Value 2} & \textbf{Total} \\
\midrule
Item A & 10.5 & 20.3 & 30.8 \\
Item B & 15.2 & 25.1 & 40.3 \\
Item C & 22.1 & 12.7 & 34.8 \\
\bottomrule
\end{tabular}
\caption{Sample Data Table}
\label{tab:sample}
\end{table}

\section{Lists}
\subsection{Bullet Points}
\begin{itemize}
  \item First item
  \item Second item
  \item Third item with sub-items:
    \begin{itemize}
      \item Sub-item 1
      \item Sub-item 2
    \end{itemize}
\end{itemize}

\subsection{Numbered List}
\begin{enumerate}
  \item First step
  \item Second step
  \item Third step
\end{enumerate}

\section{Citations and References}
When you use the citation search feature, references will be automatically managed. For example, you might cite a paper like this: \cite{example2024}. All citations will be automatically added to your \texttt{cite.bib} file.

\section{Conclusion}
This template demonstrates basic \LaTeX{} features including mathematical formulas, tables, and lists. The bibliography system is now set up to work seamlessly with Bibby AI's citation management. Feel free to modify it for your needs!

\bibliographystyle{plain}
\bibliography{cite}

\end{document}\documentclass{article}
\usepackage[utf8]{inputenc}
\usepackage{booktabs}
\usepackage{amsmath}
\usepackage{natbib}

\title{My Document}
\author{Your Name}
\date{\today}

\begin{document}

\maketitle

\section{Introduction}
This is a sample document demonstrating various \LaTeX{} features including mathematical equations, tables, and formatting. You can cite papers using the citation search feature, and they will be automatically added to your \texttt{cite.bib} file.

\section{Mathematical Expressions}
Here are some example equations:

\subsection{Quadratic Formula}
The quadratic formula is given by:
\[x = \frac{-b \pm \sqrt{b^2 - 4ac}}{2a}\]

\subsection{Matrix Example}
A sample matrix:
\[
\begin{pmatrix}
a_{11} & a_{12} & a_{13} \\
a_{21} & a_{22} & a_{23} \\
a_{31} & a_{32} & a_{33}
\end{pmatrix}
\]

\section{Tables}
Below is a sample table showing data:

\begin{table}[h]
\centering
\begin{tabular}{@{}lrrr@{}}
\toprule
\textbf{Category} & \textbf{Value 1} & \textbf{Value 2} & \textbf{Total} \\
\midrule
Item A & 10.5 & 20.3 & 30.8 \\
Item B & 15.2 & 25.1 & 40.3 \\
Item C & 22.1 & 12.7 & 34.8 \\
\bottomrule
\end{tabular}
\caption{Sample Data Table}
\label{tab:sample}
\end{table}

\section{Lists}
\subsection{Bullet Points}
\begin{itemize}
  \item First item
  \item Second item
  \item Third item with sub-items:
    \begin{itemize}
      \item Sub-item 1
      \item Sub-item 2
    \end{itemize}
\end{itemize}

\subsection{Numbered List}
\begin{enumerate}
  \item First step
  \item Second step
  \item Third step
\end{enumerate}

\section{Citations and References}
When you use the citation search feature, references will be automatically managed. For example, you might cite a paper like this: \cite{example2024}. All citations will be automatically added to your \texttt{cite.bib} file.

\section{Conclusion}
This template demonstrates basic \LaTeX{} features including mathematical formulas, tables, and lists. The bibliography system is now set up to work seamlessly with Bibby AI's citation management. Feel free to modify it for your needs!

\bibliographystyle{plain}
\bibliography{cite}

\end{document}\documentclass{article}
\usepackage[utf8]{inputenc}
\usepackage{booktabs}
\usepackage{amsmath}
\usepackage{natbib}

\title{My Document}
\author{Your Name}
\date{\today}

\begin{document}

\maketitle

\section{Introduction}
This is a sample document demonstrating various \LaTeX{} features including mathematical equations, tables, and formatting. You can cite papers using the citation search feature, and they will be automatically added to your \texttt{cite.bib} file.

\section{Mathematical Expressions}
Here are some example equations:

\subsection{Quadratic Formula}
The quadratic formula is given by:
\[x = \frac{-b \pm \sqrt{b^2 - 4ac}}{2a}\]

\subsection{Matrix Example}
A sample matrix:
\[
\begin{pmatrix}
a_{11} & a_{12} & a_{13} \\
a_{21} & a_{22} & a_{23} \\
a_{31} & a_{32} & a_{33}
\end{pmatrix}
\]

\section{Tables}
Below is a sample table showing data:

\begin{table}[h]
\centering
\begin{tabular}{@{}lrrr@{}}
\toprule
\textbf{Category} & \textbf{Value 1} & \textbf{Value 2} & \textbf{Total} \\
\midrule
Item A & 10.5 & 20.3 & 30.8 \\
Item B & 15.2 & 25.1 & 40.3 \\
Item C & 22.1 & 12.7 & 34.8 \\
\bottomrule
\end{tabular}
\caption{Sample Data Table}
\label{tab:sample}
\end{table}

\section{Lists}
\subsection{Bullet Points}
\begin{itemize}
  \item First item
  \item Second item
  \item Third item with sub-items:
    \begin{itemize}
      \item Sub-item 1
      \item Sub-item 2
    \end{itemize}
\end{itemize}

\subsection{Numbered List}
\begin{enumerate}
  \item First step
  \item Second step
  \item Third step
\end{enumerate}

\section{Citations and References}
When you use the citation search feature, references will be automatically managed. For example, you might cite a paper like this: \cite{example2024}. All citations will be automatically added to your \texttt{cite.bib} file.

\section{Conclusion}
This template demonstrates basic \LaTeX{} features including mathematical formulas, tables, and lists. The bibliography system is now set up to work seamlessly with Bibby AI's citation management. Feel free to modify it for your needs!

\bibliographystyle{plain}
\bibliography{cite}

\end{document}\documentclass{article}
\usepackage[utf8]{inputenc}
\usepackage{booktabs}
\usepackage{amsmath}
\usepackage{natbib}

\title{My Document}
\author{Your Name}
\date{\today}

\begin{document}

\maketitle

\section{Introduction}
This is a sample document demonstrating various \LaTeX{} features including mathematical equations, tables, and formatting. You can cite papers using the citation search feature, and they will be automatically added to your \texttt{cite.bib} file.

\section{Mathematical Expressions}
Here are some example equations:

\subsection{Quadratic Formula}
The quadratic formula is given by:
\[x = \frac{-b \pm \sqrt{b^2 - 4ac}}{2a}\]

\subsection{Matrix Example}
A sample matrix:
\[
\begin{pmatrix}
a_{11} & a_{12} & a_{13} \\
a_{21} & a_{22} & a_{23} \\
a_{31} & a_{32} & a_{33}
\end{pmatrix}
\]

\section{Tables}
Below is a sample table showing data:

\begin{table}[h]
\centering
\begin{tabular}{@{}lrrr@{}}
\toprule
\textbf{Category} & \textbf{Value 1} & \textbf{Value 2} & \textbf{Total} \\
\midrule
Item A & 10.5 & 20.3 & 30.8 \\
Item B & 15.2 & 25.1 & 40.3 \\
Item C & 22.1 & 12.7 & 34.8 \\
\bottomrule
\end{tabular}
\caption{Sample Data Table}
\label{tab:sample}
\end{table}

\section{Lists}
\subsection{Bullet Points}
\begin{itemize}
  \item First item
  \item Second item
  \item Third item with sub-items:
    \begin{itemize}
      \item Sub-item 1
      \item Sub-item 2
    \end{itemize}
\end{itemize}

\subsection{Numbered List}
\begin{enumerate}
  \item First step
  \item Second step
  \item Third step
\end{enumerate}

\section{Citations and References}
When you use the citation search feature, references will be automatically managed. For example, you might cite a paper like this: \cite{example2024}. All citations will be automatically added to your \texttt{cite.bib} file.

\section{Conclusion}
This template demonstrates basic \LaTeX{} features including mathematical formulas, tables, and lists. The bibliography system is now set up to work seamlessly with Bibby AI's citation management. Feel free to modify it for your needs!

\bibliographystyle{plain}
\bibliography{cite}

\end{document}\documentclass{article}
\usepackage[utf8]{inputenc}
\usepackage{booktabs}
\usepackage{amsmath}
\usepackage{natbib}

\title{My Document}
\author{Your Name}
\date{\today}

\begin{document}

\maketitle

\section{Introduction}
This is a sample document demonstrating various \LaTeX{} features including mathematical equations, tables, and formatting. You can cite papers using the citation search feature, and they will be automatically added to your \texttt{cite.bib} file.

\section{Mathematical Expressions}
Here are some example equations:

\subsection{Quadratic Formula}
The quadratic formula is given by:
\[x = \frac{-b \pm \sqrt{b^2 - 4ac}}{2a}\]

\subsection{Matrix Example}
A sample matrix:
\[
\begin{pmatrix}
a_{11} & a_{12} & a_{13} \\
a_{21} & a_{22} & a_{23} \\
a_{31} & a_{32} & a_{33}
\end{pmatrix}
\]

\section{Tables}
Below is a sample table showing data:

\begin{table}[h]
\centering
\begin{tabular}{@{}lrrr@{}}
\toprule
\textbf{Category} & \textbf{Value 1} & \textbf{Value 2} & \textbf{Total} \\
\midrule
Item A & 10.5 & 20.3 & 30.8 \\
Item B & 15.2 & 25.1 & 40.3 \\
Item C & 22.1 & 12.7 & 34.8 \\
\bottomrule
\end{tabular}
\caption{Sample Data Table}
\label{tab:sample}
\end{table}

\section{Lists}
\subsection{Bullet Points}
\begin{itemize}
  \item First item
  \item Second item
  \item Third item with sub-items:
    \begin{itemize}
      \item Sub-item 1
      \item Sub-item 2
    \end{itemize}
\end{itemize}

\subsection{Numbered List}
\begin{enumerate}
  \item First step
  \item Second step
  \item Third step
\end{enumerate}

\section{Citations and References}
When you use the citation search feature, references will be automatically managed. For example, you might cite a paper like this: \cite{example2024}. All citations will be automatically added to your \texttt{cite.bib} file.

\section{Conclusion}
This template demonstrates basic \LaTeX{} features including mathematical formulas, tables, and lists. The bibliography system is now set up to work seamlessly with Bibby AI's citation management. Feel free to modify it for your needs!

\bibliographystyle{plain}
\bibliography{cite}

\end{document}\documentclass{article}
\usepackage[utf8]{inputenc}
\usepackage{booktabs}
\usepackage{amsmath}
\usepackage{natbib}

\title{My Document}
\author{Your Name}
\date{\today}

\begin{document}

\maketitle

\section{Introduction}
This is a sample document demonstrating various \LaTeX{} features including mathematical equations, tables, and formatting. You can cite papers using the citation search feature, and they will be automatically added to your \texttt{cite.bib} file.

\section{Mathematical Expressions}
Here are some example equations:

\subsection{Quadratic Formula}
The quadratic formula is given by:
\[x = \frac{-b \pm \sqrt{b^2 - 4ac}}{2a}\]

\subsection{Matrix Example}
A sample matrix:
\[
\begin{pmatrix}
a_{11} & a_{12} & a_{13} \\
a_{21} & a_{22} & a_{23} \\
a_{31} & a_{32} & a_{33}
\end{pmatrix}
\]

\section{Tables}
Below is a sample table showing data:

\begin{table}[h]
\centering
\begin{tabular}{@{}lrrr@{}}
\toprule
\textbf{Category} & \textbf{Value 1} & \textbf{Value 2} & \textbf{Total} \\
\midrule
Item A & 10.5 & 20.3 & 30.8 \\
Item B & 15.2 & 25.1 & 40.3 \\
Item C & 22.1 & 12.7 & 34.8 \\
\bottomrule
\end{tabular}
\caption{Sample Data Table}
\label{tab:sample}
\end{table}

\section{Lists}
\subsection{Bullet Points}
\begin{itemize}
  \item First item
  \item Second item
  \item Third item with sub-items:
    \begin{itemize}
      \item Sub-item 1
      \item Sub-item 2
    \end{itemize}
\end{itemize}

\subsection{Numbered List}
\begin{enumerate}
  \item First step
  \item Second step
  \item Third step
\end{enumerate}

\section{Citations and References}
When you use the citation search feature, references will be automatically managed. For example, you might cite a paper like this: \cite{example2024}. All citations will be automatically added to your \texttt{cite.bib} file.

\section{Conclusion}
This template demonstrates basic \LaTeX{} features including mathematical formulas, tables, and lists. The bibliography system is now set up to work seamlessly with Bibby AI's citation management. Feel free to modify it for your needs!

\bibliographystyle{plain}
\bibliography{cite}

\end{document}\documentclass{article}
\usepackage[utf8]{inputenc}
\usepackage{booktabs}
\usepackage{amsmath}
\usepackage{natbib}

\title{My Document}
\author{Your Name}
\date{\today}

\begin{document}

\maketitle

\section{Introduction}
This is a sample document demonstrating various \LaTeX{} features including mathematical equations, tables, and formatting. You can cite papers using the citation search feature, and they will be automatically added to your \texttt{cite.bib} file.

\section{Mathematical Expressions}
Here are some example equations:

\subsection{Quadratic Formula}
The quadratic formula is given by:
\[x = \frac{-b \pm \sqrt{b^2 - 4ac}}{2a}\]

\subsection{Matrix Example}
A sample matrix:
\[
\begin{pmatrix}
a_{11} & a_{12} & a_{13} \\
a_{21} & a_{22} & a_{23} \\
a_{31} & a_{32} & a_{33}
\end{pmatrix}
\]

\section{Tables}
Below is a sample table showing data:

\begin{table}[h]
\centering
\begin{tabular}{@{}lrrr@{}}
\toprule
\textbf{Category} & \textbf{Value 1} & \textbf{Value 2} & \textbf{Total} \\
\midrule
Item A & 10.5 & 20.3 & 30.8 \\
Item B & 15.2 & 25.1 & 40.3 \\
Item C & 22.1 & 12.7 & 34.8 \\
\bottomrule
\end{tabular}
\caption{Sample Data Table}
\label{tab:sample}
\end{table}

\section{Lists}
\subsection{Bullet Points}
\begin{itemize}
  \item First item
  \item Second item
  \item Third item with sub-items:
    \begin{itemize}
      \item Sub-item 1
      \item Sub-item 2
    \end{itemize}
\end{itemize}

\subsection{Numbered List}
\begin{enumerate}
  \item First step
  \item Second step
  \item Third step
\end{enumerate}

\section{Citations and References}
When you use the citation search feature, references will be automatically managed. For example, you might cite a paper like this: \cite{example2024}. All citations will be automatically added to your \texttt{cite.bib} file.

\section{Conclusion}
This template demonstrates basic \LaTeX{} features including mathematical formulas, tables, and lists. The bibliography system is now set up to work seamlessly with Bibby AI's citation management. Feel free to modify it for your needs!

\bibliographystyle{plain}
\bibliography{cite}

\end{document}\documentclass{article}
\usepackage[utf8]{inputenc}
\usepackage{booktabs}
\usepackage{amsmath}
\usepackage{natbib}

\title{My Document}
\author{Your Name}
\date{\today}

\begin{document}

\maketitle

\section{Introduction}
This is a sample document demonstrating various \LaTeX{} features including mathematical equations, tables, and formatting. You can cite papers using the citation search feature, and they will be automatically added to your \texttt{cite.bib} file.

\section{Mathematical Expressions}
Here are some example equations:

\subsection{Quadratic Formula}
The quadratic formula is given by:
\[x = \frac{-b \pm \sqrt{b^2 - 4ac}}{2a}\]

\subsection{Matrix Example}
A sample matrix:
\[
\begin{pmatrix}
a_{11} & a_{12} & a_{13} \\
a_{21} & a_{22} & a_{23} \\
a_{31} & a_{32} & a_{33}
\end{pmatrix}
\]

\section{Tables}
Below is a sample table showing data:

\begin{table}[h]
\centering
\begin{tabular}{@{}lrrr@{}}
\toprule
\textbf{Category} & \textbf{Value 1} & \textbf{Value 2} & \textbf{Total} \\
\midrule
Item A & 10.5 & 20.3 & 30.8 \\
Item B & 15.2 & 25.1 & 40.3 \\
Item C & 22.1 & 12.7 & 34.8 \\
\bottomrule
\end{tabular}
\caption{Sample Data Table}
\label{tab:sample}
\end{table}

\section{Lists}
\subsection{Bullet Points}
\begin{itemize}
  \item First item
  \item Second item
  \item Third item with sub-items:
    \begin{itemize}
      \item Sub-item 1
      \item Sub-item 2
    \end{itemize}
\end{itemize}

\subsection{Numbered List}
\begin{enumerate}
  \item First step
  \item Second step
  \item Third step
\end{enumerate}

\section{Citations and References}
When you use the citation search feature, references will be automatically managed. For example, you might cite a paper like this: \cite{example2024}. All citations will be automatically added to your \texttt{cite.bib} file.

\section{Conclusion}
This template demonstrates basic \LaTeX{} features including mathematical formulas, tables, and lists. The bibliography system is now set up to work seamlessly with Bibby AI's citation management. Feel free to modify it for your needs!

\bibliographystyle{plain}
\bibliography{cite}

\end{document}\documentclass{article}
\usepackage[utf8]{inputenc}
\usepackage{booktabs}
\usepackage{amsmath}
\usepackage{natbib}

\title{My Document}
\author{Your Name}
\date{\today}

\begin{document}

\maketitle

\section{Introduction}
This is a sample document demonstrating various \LaTeX{} features including mathematical equations, tables, and formatting. You can cite papers using the citation search feature, and they will be automatically added to your \texttt{cite.bib} file.

\section{Mathematical Expressions}
Here are some example equations:

\subsection{Quadratic Formula}
The quadratic formula is given by:
\[x = \frac{-b \pm \sqrt{b^2 - 4ac}}{2a}\]

\subsection{Matrix Example}
A sample matrix:
\[
\begin{pmatrix}
a_{11} & a_{12} & a_{13} \\
a_{21} & a_{22} & a_{23} \\
a_{31} & a_{32} & a_{33}
\end{pmatrix}
\]

\section{Tables}
Below is a sample table showing data:

\begin{table}[h]
\centering
\begin{tabular}{@{}lrrr@{}}
\toprule
\textbf{Category} & \textbf{Value 1} & \textbf{Value 2} & \textbf{Total} \\
\midrule
Item A & 10.5 & 20.3 & 30.8 \\
Item B & 15.2 & 25.1 & 40.3 \\
Item C & 22.1 & 12.7 & 34.8 \\
\bottomrule
\end{tabular}
\caption{Sample Data Table}
\label{tab:sample}
\end{table}

\section{Lists}
\subsection{Bullet Points}
\begin{itemize}
  \item First item
  \item Second item
  \item Third item with sub-items:
    \begin{itemize}
      \item Sub-item 1
      \item Sub-item 2
    \end{itemize}
\end{itemize}

\subsection{Numbered List}
\begin{enumerate}
  \item First step
  \item Second step
  \item Third step
\end{enumerate}

\section{Citations and References}
When you use the citation search feature, references will be automatically managed. For example, you might cite a paper like this: \cite{example2024}. All citations will be automatically added to your \texttt{cite.bib} file.

\section{Conclusion}
This template demonstrates basic \LaTeX{} features including mathematical formulas, tables, and lists. The bibliography system is now set up to work seamlessly with Bibby AI's citation management. Feel free to modify it for your needs!

\bibliographystyle{plain}
\bibliography{cite}

\end{document}\documentclass{article}
\usepackage[utf8]{inputenc}
\usepackage{booktabs}
\usepackage{amsmath}
\usepackage{natbib}

\title{My Document}
\author{Your Name}
\date{\today}

\begin{document}

\maketitle

\section{Introduction}
This is a sample document demonstrating various \LaTeX{} features including mathematical equations, tables, and formatting. You can cite papers using the citation search feature, and they will be automatically added to your \texttt{cite.bib} file.

\section{Mathematical Expressions}
Here are some example equations:

\subsection{Quadratic Formula}
The quadratic formula is given by:
\[x = \frac{-b \pm \sqrt{b^2 - 4ac}}{2a}\]

\subsection{Matrix Example}
A sample matrix:
\[
\begin{pmatrix}
a_{11} & a_{12} & a_{13} \\
a_{21} & a_{22} & a_{23} \\
a_{31} & a_{32} & a_{33}
\end{pmatrix}
\]

\section{Tables}
Below is a sample table showing data:

\begin{table}[h]
\centering
\begin{tabular}{@{}lrrr@{}}
\toprule
\textbf{Category} & \textbf{Value 1} & \textbf{Value 2} & \textbf{Total} \\
\midrule
Item A & 10.5 & 20.3 & 30.8 \\
Item B & 15.2 & 25.1 & 40.3 \\
Item C & 22.1 & 12.7 & 34.8 \\
\bottomrule
\end{tabular}
\caption{Sample Data Table}
\label{tab:sample}
\end{table}

\section{Lists}
\subsection{Bullet Points}
\begin{itemize}
  \item First item
  \item Second item
  \item Third item with sub-items:
    \begin{itemize}
      \item Sub-item 1
      \item Sub-item 2
    \end{itemize}
\end{itemize}

\subsection{Numbered List}
\begin{enumerate}
  \item First step
  \item Second step
  \item Third step
\end{enumerate}

\section{Citations and References}
When you use the citation search feature, references will be automatically managed. For example, you might cite a paper like this: \cite{example2024}. All citations will be automatically added to your \texttt{cite.bib} file.

\section{Conclusion}
This template demonstrates basic \LaTeX{} features including mathematical formulas, tables, and lists. The bibliography system is now set up to work seamlessly with Bibby AI's citation management. Feel free to modify it for your needs!

\bibliographystyle{plain}
\bibliography{cite}

\end{document}\documentclass{article}
\usepackage[utf8]{inputenc}
\usepackage{booktabs}
\usepackage{amsmath}
\usepackage{natbib}

\title{My Document}
\author{Your Name}
\date{\today}

\begin{document}

\maketitle

\section{Introduction}
This is a sample document demonstrating various \LaTeX{} features including mathematical equations, tables, and formatting. You can cite papers using the citation search feature, and they will be automatically added to your \texttt{cite.bib} file.

\section{Mathematical Expressions}
Here are some example equations:

\subsection{Quadratic Formula}
The quadratic formula is given by:
\[x = \frac{-b \pm \sqrt{b^2 - 4ac}}{2a}\]

\subsection{Matrix Example}
A sample matrix:
\[
\begin{pmatrix}
a_{11} & a_{12} & a_{13} \\
a_{21} & a_{22} & a_{23} \\
a_{31} & a_{32} & a_{33}
\end{pmatrix}
\]

\section{Tables}
Below is a sample table showing data:

\begin{table}[h]
\centering
\begin{tabular}{@{}lrrr@{}}
\toprule
\textbf{Category} & \textbf{Value 1} & \textbf{Value 2} & \textbf{Total} \\
\midrule
Item A & 10.5 & 20.3 & 30.8 \\
Item B & 15.2 & 25.1 & 40.3 \\
Item C & 22.1 & 12.7 & 34.8 \\
\bottomrule
\end{tabular}
\caption{Sample Data Table}
\label{tab:sample}
\end{table}

\section{Lists}
\subsection{Bullet Points}
\begin{itemize}
  \item First item
  \item Second item
  \item Third item with sub-items:
    \begin{itemize}
      \item Sub-item 1
      \item Sub-item 2
    \end{itemize}
\end{itemize}

\subsection{Numbered List}
\begin{enumerate}
  \item First step
  \item Second step
  \item Third step
\end{enumerate}

\section{Citations and References}
When you use the citation search feature, references will be automatically managed. For example, you might cite a paper like this: \cite{example2024}. All citations will be automatically added to your \texttt{cite.bib} file.

\section{Conclusion}
This template demonstrates basic \LaTeX{} features including mathematical formulas, tables, and lists. The bibliography system is now set up to work seamlessly with Bibby AI's citation management. Feel free to modify it for your needs!

\bibliographystyle{plain}
\bibliography{cite}

\end{document}\documentclass{article}
\usepackage[utf8]{inputenc}
\usepackage{booktabs}
\usepackage{amsmath}
\usepackage{natbib}

\title{My Document}
\author{Your Name}
\date{\today}

\begin{document}

\maketitle

\section{Introduction}
This is a sample document demonstrating various \LaTeX{} features including mathematical equations, tables, and formatting. You can cite papers using the citation search feature, and they will be automatically added to your \texttt{cite.bib} file.

\section{Mathematical Expressions}
Here are some example equations:

\subsection{Quadratic Formula}
The quadratic formula is given by:
\[x = \frac{-b \pm \sqrt{b^2 - 4ac}}{2a}\]

\subsection{Matrix Example}
A sample matrix:
\[
\begin{pmatrix}
a_{11} & a_{12} & a_{13} \\
a_{21} & a_{22} & a_{23} \\
a_{31} & a_{32} & a_{33}
\end{pmatrix}
\]

\section{Tables}
Below is a sample table showing data:

\begin{table}[h]
\centering
\begin{tabular}{@{}lrrr@{}}
\toprule
\textbf{Category} & \textbf{Value 1} & \textbf{Value 2} & \textbf{Total} \\
\midrule
Item A & 10.5 & 20.3 & 30.8 \\
Item B & 15.2 & 25.1 & 40.3 \\
Item C & 22.1 & 12.7 & 34.8 \\
\bottomrule
\end{tabular}
\caption{Sample Data Table}
\label{tab:sample}
\end{table}

\section{Lists}
\subsection{Bullet Points}
\begin{itemize}
  \item First item
  \item Second item
  \item Third item with sub-items:
    \begin{itemize}
      \item Sub-item 1
      \item Sub-item 2
    \end{itemize}
\end{itemize}

\subsection{Numbered List}
\begin{enumerate}
  \item First step
  \item Second step
  \item Third step
\end{enumerate}

\section{Citations and References}
When you use the citation search feature, references will be automatically managed. For example, you might cite a paper like this: \cite{example2024}. All citations will be automatically added to your \texttt{cite.bib} file.

\section{Conclusion}
This template demonstrates basic \LaTeX{} features including mathematical formulas, tables, and lists. The bibliography system is now set up to work seamlessly with Bibby AI's citation management. Feel free to modify it for your needs!

\bibliographystyle{plain}
\bibliography{cite}

\end{document}\documentclass{article}
\usepackage[utf8]{inputenc}
\usepackage{booktabs}
\usepackage{amsmath}
\usepackage{natbib}

\title{My Document}
\author{Your Name}
\date{\today}

\begin{document}

\maketitle

\section{Introduction}
This is a sample document demonstrating various \LaTeX{} features including mathematical equations, tables, and formatting. You can cite papers using the citation search feature, and they will be automatically added to your \texttt{cite.bib} file.

\section{Mathematical Expressions}
Here are some example equations:

\subsection{Quadratic Formula}
The quadratic formula is given by:
\[x = \frac{-b \pm \sqrt{b^2 - 4ac}}{2a}\]

\subsection{Matrix Example}
A sample matrix:
\[
\begin{pmatrix}
a_{11} & a_{12} & a_{13} \\
a_{21} & a_{22} & a_{23} \\
a_{31} & a_{32} & a_{33}
\end{pmatrix}
\]

\section{Tables}
Below is a sample table showing data:

\begin{table}[h]
\centering
\begin{tabular}{@{}lrrr@{}}
\toprule
\textbf{Category} & \textbf{Value 1} & \textbf{Value 2} & \textbf{Total} \\
\midrule
Item A & 10.5 & 20.3 & 30.8 \\
Item B & 15.2 & 25.1 & 40.3 \\
Item C & 22.1 & 12.7 & 34.8 \\
\bottomrule
\end{tabular}
\caption{Sample Data Table}
\label{tab:sample}
\end{table}

\section{Lists}
\subsection{Bullet Points}
\begin{itemize}
  \item First item
  \item Second item
  \item Third item with sub-items:
    \begin{itemize}
      \item Sub-item 1
      \item Sub-item 2
    \end{itemize}
\end{itemize}

\subsection{Numbered List}
\begin{enumerate}
  \item First step
  \item Second step
  \item Third step
\end{enumerate}

\section{Citations and References}
When you use the citation search feature, references will be automatically managed. For example, you might cite a paper like this: \cite{example2024}. All citations will be automatically added to your \texttt{cite.bib} file.

\section{Conclusion}
This template demonstrates basic \LaTeX{} features including mathematical formulas, tables, and lists. The bibliography system is now set up to work seamlessly with Bibby AI's citation management. Feel free to modify it for your needs!

\bibliographystyle{plain}
\bibliography{cite}

\end{document}\documentclass{article}
\usepackage[utf8]{inputenc}
\usepackage{booktabs}
\usepackage{amsmath}
\usepackage{natbib}

\title{My Document}
\author{Your Name}
\date{\today}

\begin{document}

\maketitle

\section{Introduction}
This is a sample document demonstrating various \LaTeX{} features including mathematical equations, tables, and formatting. You can cite papers using the citation search feature, and they will be automatically added to your \texttt{cite.bib} file.

\section{Mathematical Expressions}
Here are some example equations:

\subsection{Quadratic Formula}
The quadratic formula is given by:
\[x = \frac{-b \pm \sqrt{b^2 - 4ac}}{2a}\]

\subsection{Matrix Example}
A sample matrix:
\[
\begin{pmatrix}
a_{11} & a_{12} & a_{13} \\
a_{21} & a_{22} & a_{23} \\
a_{31} & a_{32} & a_{33}
\end{pmatrix}
\]

\section{Tables}
Below is a sample table showing data:

\begin{table}[h]
\centering
\begin{tabular}{@{}lrrr@{}}
\toprule
\textbf{Category} & \textbf{Value 1} & \textbf{Value 2} & \textbf{Total} \\
\midrule
Item A & 10.5 & 20.3 & 30.8 \\
Item B & 15.2 & 25.1 & 40.3 \\
Item C & 22.1 & 12.7 & 34.8 \\
\bottomrule
\end{tabular}
\caption{Sample Data Table}
\label{tab:sample}
\end{table}

\section{Lists}
\subsection{Bullet Points}
\begin{itemize}
  \item First item
  \item Second item
  \item Third item with sub-items:
    \begin{itemize}
      \item Sub-item 1
      \item Sub-item 2
    \end{itemize}
\end{itemize}

\subsection{Numbered List}
\begin{enumerate}
  \item First step
  \item Second step
  \item Third step
\end{enumerate}

\section{Citations and References}
When you use the citation search feature, references will be automatically managed. For example, you might cite a paper like this: \cite{example2024}. All citations will be automatically added to your \texttt{cite.bib} file.

\section{Conclusion}
This template demonstrates basic \LaTeX{} features including mathematical formulas, tables, and lists. The bibliography system is now set up to work seamlessly with Bibby AI's citation management. Feel free to modify it for your needs!

\bibliographystyle{plain}
\bibliography{cite}

\end{document}\documentclass{article}
\usepackage[utf8]{inputenc}
\usepackage{booktabs}
\usepackage{amsmath}
\usepackage{natbib}

\title{My Document}
\author{Your Name}
\date{\today}

\begin{document}

\maketitle

\section{Introduction}
This is a sample document demonstrating various \LaTeX{} features including mathematical equations, tables, and formatting. You can cite papers using the citation search feature, and they will be automatically added to your \texttt{cite.bib} file.

\section{Mathematical Expressions}
Here are some example equations:

\subsection{Quadratic Formula}
The quadratic formula is given by:
\[x = \frac{-b \pm \sqrt{b^2 - 4ac}}{2a}\]

\subsection{Matrix Example}
A sample matrix:
\[
\begin{pmatrix}
a_{11} & a_{12} & a_{13} \\
a_{21} & a_{22} & a_{23} \\
a_{31} & a_{32} & a_{33}
\end{pmatrix}
\]

\section{Tables}
Below is a sample table showing data:

\begin{table}[h]
\centering
\begin{tabular}{@{}lrrr@{}}
\toprule
\textbf{Category} & \textbf{Value 1} & \textbf{Value 2} & \textbf{Total} \\
\midrule
Item A & 10.5 & 20.3 & 30.8 \\
Item B & 15.2 & 25.1 & 40.3 \\
Item C & 22.1 & 12.7 & 34.8 \\
\bottomrule
\end{tabular}
\caption{Sample Data Table}
\label{tab:sample}
\end{table}

\section{Lists}
\subsection{Bullet Points}
\begin{itemize}
  \item First item
  \item Second item
  \item Third item with sub-items:
    \begin{itemize}
      \item Sub-item 1
      \item Sub-item 2
    \end{itemize}
\end{itemize}

\subsection{Numbered List}
\begin{enumerate}
  \item First step
  \item Second step
  \item Third step
\end{enumerate}

\section{Citations and References}
When you use the citation search feature, references will be automatically managed. For example, you might cite a paper like this: \cite{example2024}. All citations will be automatically added to your \texttt{cite.bib} file.

\section{Conclusion}
This template demonstrates basic \LaTeX{} features including mathematical formulas, tables, and lists. The bibliography system is now set up to work seamlessly with Bibby AI's citation management. Feel free to modify it for your needs!

\bibliographystyle{plain}
\bibliography{cite}

\end{document}\documentclass{article}
\usepackage[utf8]{inputenc}
\usepackage{booktabs}
\usepackage{amsmath}
\usepackage{natbib}

\title{My Document}
\author{Your Name}
\date{\today}

\begin{document}

\maketitle

\section{Introduction}
This is a sample document demonstrating various \LaTeX{} features including mathematical equations, tables, and formatting. You can cite papers using the citation search feature, and they will be automatically added to your \texttt{cite.bib} file.

\section{Mathematical Expressions}
Here are some example equations:

\subsection{Quadratic Formula}
The quadratic formula is given by:
\[x = \frac{-b \pm \sqrt{b^2 - 4ac}}{2a}\]

\subsection{Matrix Example}
A sample matrix:
\[
\begin{pmatrix}
a_{11} & a_{12} & a_{13} \\
a_{21} & a_{22} & a_{23} \\
a_{31} & a_{32} & a_{33}
\end{pmatrix}
\]

\section{Tables}
Below is a sample table showing data:

\begin{table}[h]
\centering
\begin{tabular}{@{}lrrr@{}}
\toprule
\textbf{Category} & \textbf{Value 1} & \textbf{Value 2} & \textbf{Total} \\
\midrule
Item A & 10.5 & 20.3 & 30.8 \\
Item B & 15.2 & 25.1 & 40.3 \\
Item C & 22.1 & 12.7 & 34.8 \\
\bottomrule
\end{tabular}
\caption{Sample Data Table}
\label{tab:sample}
\end{table}

\section{Lists}
\subsection{Bullet Points}
\begin{itemize}
  \item First item
  \item Second item
  \item Third item with sub-items:
    \begin{itemize}
      \item Sub-item 1
      \item Sub-item 2
    \end{itemize}
\end{itemize}

\subsection{Numbered List}
\begin{enumerate}
  \item First step
  \item Second step
  \item Third step
\end{enumerate}

\section{Citations and References}
When you use the citation search feature, references will be automatically managed. For example, you might cite a paper like this: \cite{example2024}. All citations will be automatically added to your \texttt{cite.bib} file.

\section{Conclusion}
This template demonstrates basic \LaTeX{} features including mathematical formulas, tables, and lists. The bibliography system is now set up to work seamlessly with Bibby AI's citation management. Feel free to modify it for your needs!

\bibliographystyle{plain}
\bibliography{cite}

\end{document}\documentclass{article}
\usepackage[utf8]{inputenc}
\usepackage{booktabs}
\usepackage{amsmath}
\usepackage{natbib}

\title{My Document}
\author{Your Name}
\date{\today}

\begin{document}

\maketitle

\section{Introduction}
This is a sample document demonstrating various \LaTeX{} features including mathematical equations, tables, and formatting. You can cite papers using the citation search feature, and they will be automatically added to your \texttt{cite.bib} file.

\section{Mathematical Expressions}
Here are some example equations:

\subsection{Quadratic Formula}
The quadratic formula is given by:
\[x = \frac{-b \pm \sqrt{b^2 - 4ac}}{2a}\]

\subsection{Matrix Example}
A sample matrix:
\[
\begin{pmatrix}
a_{11} & a_{12} & a_{13} \\
a_{21} & a_{22} & a_{23} \\
a_{31} & a_{32} & a_{33}
\end{pmatrix}
\]

\section{Tables}
Below is a sample table showing data:

\begin{table}[h]
\centering
\begin{tabular}{@{}lrrr@{}}
\toprule
\textbf{Category} & \textbf{Value 1} & \textbf{Value 2} & \textbf{Total} \\
\midrule
Item A & 10.5 & 20.3 & 30.8 \\
Item B & 15.2 & 25.1 & 40.3 \\
Item C & 22.1 & 12.7 & 34.8 \\
\bottomrule
\end{tabular}
\caption{Sample Data Table}
\label{tab:sample}
\end{table}

\section{Lists}
\subsection{Bullet Points}
\begin{itemize}
  \item First item
  \item Second item
  \item Third item with sub-items:
    \begin{itemize}
      \item Sub-item 1
      \item Sub-item 2
    \end{itemize}
\end{itemize}

\subsection{Numbered List}
\begin{enumerate}
  \item First step
  \item Second step
  \item Third step
\end{enumerate}

\section{Citations and References}
When you use the citation search feature, references will be automatically managed. For example, you might cite a paper like this: \cite{example2024}. All citations will be automatically added to your \texttt{cite.bib} file.

\section{Conclusion}
This template demonstrates basic \LaTeX{} features including mathematical formulas, tables, and lists. The bibliography system is now set up to work seamlessly with Bibby AI's citation management. Feel free to modify it for your needs!

\bibliographystyle{plain}
\bibliography{cite}

\end{document}\documentclass{article}
\usepackage[utf8]{inputenc}
\usepackage{booktabs}
\usepackage{amsmath}
\usepackage{natbib}

\title{My Document}
\author{Your Name}
\date{\today}

\begin{document}

\maketitle

\section{Introduction}
This is a sample document demonstrating various \LaTeX{} features including mathematical equations, tables, and formatting. You can cite papers using the citation search feature, and they will be automatically added to your \texttt{cite.bib} file.

\section{Mathematical Expressions}
Here are some example equations:

\subsection{Quadratic Formula}
The quadratic formula is given by:
\[x = \frac{-b \pm \sqrt{b^2 - 4ac}}{2a}\]

\subsection{Matrix Example}
A sample matrix:
\[
\begin{pmatrix}
a_{11} & a_{12} & a_{13} \\
a_{21} & a_{22} & a_{23} \\
a_{31} & a_{32} & a_{33}
\end{pmatrix}
\]

\section{Tables}
Below is a sample table showing data:

\begin{table}[h]
\centering
\begin{tabular}{@{}lrrr@{}}
\toprule
\textbf{Category} & \textbf{Value 1} & \textbf{Value 2} & \textbf{Total} \\
\midrule
Item A & 10.5 & 20.3 & 30.8 \\
Item B & 15.2 & 25.1 & 40.3 \\
Item C & 22.1 & 12.7 & 34.8 \\
\bottomrule
\end{tabular}
\caption{Sample Data Table}
\label{tab:sample}
\end{table}

\section{Lists}
\subsection{Bullet Points}
\begin{itemize}
  \item First item
  \item Second item
  \item Third item with sub-items:
    \begin{itemize}
      \item Sub-item 1
      \item Sub-item 2
    \end{itemize}
\end{itemize}

\subsection{Numbered List}
\begin{enumerate}
  \item First step
  \item Second step
  \item Third step
\end{enumerate}

\section{Citations and References}
When you use the citation search feature, references will be automatically managed. For example, you might cite a paper like this: \cite{example2024}. All citations will be automatically added to your \texttt{cite.bib} file.

\section{Conclusion}
This template demonstrates basic \LaTeX{} features including mathematical formulas, tables, and lists. The bibliography system is now set up to work seamlessly with Bibby AI's citation management. Feel free to modify it for your needs!

\bibliographystyle{plain}
\bibliography{cite}

\end{document}\documentclass{article}
\usepackage[utf8]{inputenc}
\usepackage{booktabs}
\usepackage{amsmath}
\usepackage{natbib}

\title{My Document}
\author{Your Name}
\date{\today}

\begin{document}

\maketitle

\section{Introduction}
This is a sample document demonstrating various \LaTeX{} features including mathematical equations, tables, and formatting. You can cite papers using the citation search feature, and they will be automatically added to your \texttt{cite.bib} file.

\section{Mathematical Expressions}
Here are some example equations:

\subsection{Quadratic Formula}
The quadratic formula is given by:
\[x = \frac{-b \pm \sqrt{b^2 - 4ac}}{2a}\]

\subsection{Matrix Example}
A sample matrix:
\[
\begin{pmatrix}
a_{11} & a_{12} & a_{13} \\
a_{21} & a_{22} & a_{23} \\
a_{31} & a_{32} & a_{33}
\end{pmatrix}
\]

\section{Tables}
Below is a sample table showing data:

\begin{table}[h]
\centering
\begin{tabular}{@{}lrrr@{}}
\toprule
\textbf{Category} & \textbf{Value 1} & \textbf{Value 2} & \textbf{Total} \\
\midrule
Item A & 10.5 & 20.3 & 30.8 \\
Item B & 15.2 & 25.1 & 40.3 \\
Item C & 22.1 & 12.7 & 34.8 \\
\bottomrule
\end{tabular}
\caption{Sample Data Table}
\label{tab:sample}
\end{table}

\section{Lists}
\subsection{Bullet Points}
\begin{itemize}
  \item First item
  \item Second item
  \item Third item with sub-items:
    \begin{itemize}
      \item Sub-item 1
      \item Sub-item 2
    \end{itemize}
\end{itemize}

\subsection{Numbered List}
\begin{enumerate}
  \item First step
  \item Second step
  \item Third step
\end{enumerate}

\section{Citations and References}
When you use the citation search feature, references will be automatically managed. For example, you might cite a paper like this: \cite{example2024}. All citations will be automatically added to your \texttt{cite.bib} file.

\section{Conclusion}
This template demonstrates basic \LaTeX{} features including mathematical formulas, tables, and lists. The bibliography system is now set up to work seamlessly with Bibby AI's citation management. Feel free to modify it for your needs!

\bibliographystyle{plain}
\bibliography{cite}

\end{document}\documentclass{article}
\usepackage[utf8]{inputenc}
\usepackage{booktabs}
\usepackage{amsmath}
\usepackage{natbib}

\title{My Document}
\author{Your Name}
\date{\today}

\begin{document}

\maketitle

\section{Introduction}
This is a sample document demonstrating various \LaTeX{} features including mathematical equations, tables, and formatting. You can cite papers using the citation search feature, and they will be automatically added to your \texttt{cite.bib} file.

\section{Mathematical Expressions}
Here are some example equations:

\subsection{Quadratic Formula}
The quadratic formula is given by:
\[x = \frac{-b \pm \sqrt{b^2 - 4ac}}{2a}\]

\subsection{Matrix Example}
A sample matrix:
\[
\begin{pmatrix}
a_{11} & a_{12} & a_{13} \\
a_{21} & a_{22} & a_{23} \\
a_{31} & a_{32} & a_{33}
\end{pmatrix}
\]

\section{Tables}
Below is a sample table showing data:

\begin{table}[h]
\centering
\begin{tabular}{@{}lrrr@{}}
\toprule
\textbf{Category} & \textbf{Value 1} & \textbf{Value 2} & \textbf{Total} \\
\midrule
Item A & 10.5 & 20.3 & 30.8 \\
Item B & 15.2 & 25.1 & 40.3 \\
Item C & 22.1 & 12.7 & 34.8 \\
\bottomrule
\end{tabular}
\caption{Sample Data Table}
\label{tab:sample}
\end{table}

\section{Lists}
\subsection{Bullet Points}
\begin{itemize}
  \item First item
  \item Second item
  \item Third item with sub-items:
    \begin{itemize}
      \item Sub-item 1
      \item Sub-item 2
    \end{itemize}
\end{itemize}

\subsection{Numbered List}
\begin{enumerate}
  \item First step
  \item Second step
  \item Third step
\end{enumerate}

\section{Citations and References}
When you use the citation search feature, references will be automatically managed. For example, you might cite a paper like this: \cite{example2024}. All citations will be automatically added to your \texttt{cite.bib} file.

\section{Conclusion}
This template demonstrates basic \LaTeX{} features including mathematical formulas, tables, and lists. The bibliography system is now set up to work seamlessly with Bibby AI's citation management. Feel free to modify it for your needs!

\bibliographystyle{plain}
\bibliography{cite}

\end{document}\documentclass{article}
\usepackage[utf8]{inputenc}
\usepackage{booktabs}
\usepackage{amsmath}
\usepackage{natbib}

\title{My Document}
\author{Your Name}
\date{\today}

\begin{document}

\maketitle

\section{Introduction}
This is a sample document demonstrating various \LaTeX{} features including mathematical equations, tables, and formatting. You can cite papers using the citation search feature, and they will be automatically added to your \texttt{cite.bib} file.

\section{Mathematical Expressions}
Here are some example equations:

\subsection{Quadratic Formula}
The quadratic formula is given by:
\[x = \frac{-b \pm \sqrt{b^2 - 4ac}}{2a}\]

\subsection{Matrix Example}
A sample matrix:
\[
\begin{pmatrix}
a_{11} & a_{12} & a_{13} \\
a_{21} & a_{22} & a_{23} \\
a_{31} & a_{32} & a_{33}
\end{pmatrix}
\]

\section{Tables}
Below is a sample table showing data:

\begin{table}[h]
\centering
\begin{tabular}{@{}lrrr@{}}
\toprule
\textbf{Category} & \textbf{Value 1} & \textbf{Value 2} & \textbf{Total} \\
\midrule
Item A & 10.5 & 20.3 & 30.8 \\
Item B & 15.2 & 25.1 & 40.3 \\
Item C & 22.1 & 12.7 & 34.8 \\
\bottomrule
\end{tabular}
\caption{Sample Data Table}
\label{tab:sample}
\end{table}

\section{Lists}
\subsection{Bullet Points}
\begin{itemize}
  \item First item
  \item Second item
  \item Third item with sub-items:
    \begin{itemize}
      \item Sub-item 1
      \item Sub-item 2
    \end{itemize}
\end{itemize}

\subsection{Numbered List}
\begin{enumerate}
  \item First step
  \item Second step
  \item Third step
\end{enumerate}

\section{Citations and References}
When you use the citation search feature, references will be automatically managed. For example, you might cite a paper like this: \cite{example2024}. All citations will be automatically added to your \texttt{cite.bib} file.

\section{Conclusion}
This template demonstrates basic \LaTeX{} features including mathematical formulas, tables, and lists. The bibliography system is now set up to work seamlessly with Bibby AI's citation management. Feel free to modify it for your needs!

\bibliographystyle{plain}
\bibliography{cite}

\end{document}\documentclass{article}
\usepackage[utf8]{inputenc}
\usepackage{booktabs}
\usepackage{amsmath}
\usepackage{natbib}

\title{My Document}
\author{Your Name}
\date{\today}

\begin{document}

\maketitle

\section{Introduction}
This is a sample document demonstrating various \LaTeX{} features including mathematical equations, tables, and formatting. You can cite papers using the citation search feature, and they will be automatically added to your \texttt{cite.bib} file.

\section{Mathematical Expressions}
Here are some example equations:

\subsection{Quadratic Formula}
The quadratic formula is given by:
\[x = \frac{-b \pm \sqrt{b^2 - 4ac}}{2a}\]

\subsection{Matrix Example}
A sample matrix:
\[
\begin{pmatrix}
a_{11} & a_{12} & a_{13} \\
a_{21} & a_{22} & a_{23} \\
a_{31} & a_{32} & a_{33}
\end{pmatrix}
\]

\section{Tables}
Below is a sample table showing data:

\begin{table}[h]
\centering
\begin{tabular}{@{}lrrr@{}}
\toprule
\textbf{Category} & \textbf{Value 1} & \textbf{Value 2} & \textbf{Total} \\
\midrule
Item A & 10.5 & 20.3 & 30.8 \\
Item B & 15.2 & 25.1 & 40.3 \\
Item C & 22.1 & 12.7 & 34.8 \\
\bottomrule
\end{tabular}
\caption{Sample Data Table}
\label{tab:sample}
\end{table}

\section{Lists}
\subsection{Bullet Points}
\begin{itemize}
  \item First item
  \item Second item
  \item Third item with sub-items:
    \begin{itemize}
      \item Sub-item 1
      \item Sub-item 2
    \end{itemize}
\end{itemize}

\subsection{Numbered List}
\begin{enumerate}
  \item First step
  \item Second step
  \item Third step
\end{enumerate}

\section{Citations and References}
When you use the citation search feature, references will be automatically managed. For example, you might cite a paper like this: \cite{example2024}. All citations will be automatically added to your \texttt{cite.bib} file.

\section{Conclusion}
This template demonstrates basic \LaTeX{} features including mathematical formulas, tables, and lists. The bibliography system is now set up to work seamlessly with Bibby AI's citation management. Feel free to modify it for your needs!

\bibliographystyle{plain}
\bibliography{cite}

\end{document}\documentclass{article}
\usepackage[utf8]{inputenc}
\usepackage{booktabs}
\usepackage{amsmath}
\usepackage{natbib}

\title{My Document}
\author{Your Name}
\date{\today}

\begin{document}

\maketitle

\section{Introduction}
This is a sample document demonstrating various \LaTeX{} features including mathematical equations, tables, and formatting. You can cite papers using the citation search feature, and they will be automatically added to your \texttt{cite.bib} file.

\section{Mathematical Expressions}
Here are some example equations:

\subsection{Quadratic Formula}
The quadratic formula is given by:
\[x = \frac{-b \pm \sqrt{b^2 - 4ac}}{2a}\]

\subsection{Matrix Example}
A sample matrix:
\[
\begin{pmatrix}
a_{11} & a_{12} & a_{13} \\
a_{21} & a_{22} & a_{23} \\
a_{31} & a_{32} & a_{33}
\end{pmatrix}
\]

\section{Tables}
Below is a sample table showing data:

\begin{table}[h]
\centering
\begin{tabular}{@{}lrrr@{}}
\toprule
\textbf{Category} & \textbf{Value 1} & \textbf{Value 2} & \textbf{Total} \\
\midrule
Item A & 10.5 & 20.3 & 30.8 \\
Item B & 15.2 & 25.1 & 40.3 \\
Item C & 22.1 & 12.7 & 34.8 \\
\bottomrule
\end{tabular}
\caption{Sample Data Table}
\label{tab:sample}
\end{table}

\section{Lists}
\subsection{Bullet Points}
\begin{itemize}
  \item First item
  \item Second item
  \item Third item with sub-items:
    \begin{itemize}
      \item Sub-item 1
      \item Sub-item 2
    \end{itemize}
\end{itemize}

\subsection{Numbered List}
\begin{enumerate}
  \item First step
  \item Second step
  \item Third step
\end{enumerate}

\section{Citations and References}
When you use the citation search feature, references will be automatically managed. For example, you might cite a paper like this: \cite{example2024}. All citations will be automatically added to your \texttt{cite.bib} file.

\section{Conclusion}
This template demonstrates basic \LaTeX{} features including mathematical formulas, tables, and lists. The bibliography system is now set up to work seamlessly with Bibby AI's citation management. Feel free to modify it for your needs!

\bibliographystyle{plain}
\bibliography{cite}

\end{document}\documentclass{article}
\usepackage[utf8]{inputenc}
\usepackage{booktabs}
\usepackage{amsmath}
\usepackage{natbib}

\title{My Document}
\author{Your Name}
\date{\today}

\begin{document}

\maketitle

\section{Introduction}
This is a sample document demonstrating various \LaTeX{} features including mathematical equations, tables, and formatting. You can cite papers using the citation search feature, and they will be automatically added to your \texttt{cite.bib} file.

\section{Mathematical Expressions}
Here are some example equations:

\subsection{Quadratic Formula}
The quadratic formula is given by:
\[x = \frac{-b \pm \sqrt{b^2 - 4ac}}{2a}\]

\subsection{Matrix Example}
A sample matrix:
\[
\begin{pmatrix}
a_{11} & a_{12} & a_{13} \\
a_{21} & a_{22} & a_{23} \\
a_{31} & a_{32} & a_{33}
\end{pmatrix}
\]

\section{Tables}
Below is a sample table showing data:

\begin{table}[h]
\centering
\begin{tabular}{@{}lrrr@{}}
\toprule
\textbf{Category} & \textbf{Value 1} & \textbf{Value 2} & \textbf{Total} \\
\midrule
Item A & 10.5 & 20.3 & 30.8 \\
Item B & 15.2 & 25.1 & 40.3 \\
Item C & 22.1 & 12.7 & 34.8 \\
\bottomrule
\end{tabular}
\caption{Sample Data Table}
\label{tab:sample}
\end{table}

\section{Lists}
\subsection{Bullet Points}
\begin{itemize}
  \item First item
  \item Second item
  \item Third item with sub-items:
    \begin{itemize}
      \item Sub-item 1
      \item Sub-item 2
    \end{itemize}
\end{itemize}

\subsection{Numbered List}
\begin{enumerate}
  \item First step
  \item Second step
  \item Third step
\end{enumerate}

\section{Citations and References}
When you use the citation search feature, references will be automatically managed. For example, you might cite a paper like this: \cite{example2024}. All citations will be automatically added to your \texttt{cite.bib} file.

\section{Conclusion}
This template demonstrates basic \LaTeX{} features including mathematical formulas, tables, and lists. The bibliography system is now set up to work seamlessly with Bibby AI's citation management. Feel free to modify it for your needs!

\bibliographystyle{plain}
\bibliography{cite}

\end{document}\documentclass{article}
\usepackage[utf8]{inputenc}
\usepackage{booktabs}
\usepackage{amsmath}
\usepackage{natbib}

\title{My Document}
\author{Your Name}
\date{\today}

\begin{document}

\maketitle

\section{Introduction}
This is a sample document demonstrating various \LaTeX{} features including mathematical equations, tables, and formatting. You can cite papers using the citation search feature, and they will be automatically added to your \texttt{cite.bib} file.

\section{Mathematical Expressions}
Here are some example equations:

\subsection{Quadratic Formula}
The quadratic formula is given by:
\[x = \frac{-b \pm \sqrt{b^2 - 4ac}}{2a}\]

\subsection{Matrix Example}
A sample matrix:
\[
\begin{pmatrix}
a_{11} & a_{12} & a_{13} \\
a_{21} & a_{22} & a_{23} \\
a_{31} & a_{32} & a_{33}
\end{pmatrix}
\]

\section{Tables}
Below is a sample table showing data:

\begin{table}[h]
\centering
\begin{tabular}{@{}lrrr@{}}
\toprule
\textbf{Category} & \textbf{Value 1} & \textbf{Value 2} & \textbf{Total} \\
\midrule
Item A & 10.5 & 20.3 & 30.8 \\
Item B & 15.2 & 25.1 & 40.3 \\
Item C & 22.1 & 12.7 & 34.8 \\
\bottomrule
\end{tabular}
\caption{Sample Data Table}
\label{tab:sample}
\end{table}

\section{Lists}
\subsection{Bullet Points}
\begin{itemize}
  \item First item
  \item Second item
  \item Third item with sub-items:
    \begin{itemize}
      \item Sub-item 1
      \item Sub-item 2
    \end{itemize}
\end{itemize}

\subsection{Numbered List}
\begin{enumerate}
  \item First step
  \item Second step
  \item Third step
\end{enumerate}

\section{Citations and References}
When you use the citation search feature, references will be automatically managed. For example, you might cite a paper like this: \cite{example2024}. All citations will be automatically added to your \texttt{cite.bib} file.

\section{Conclusion}
This template demonstrates basic \LaTeX{} features including mathematical formulas, tables, and lists. The bibliography system is now set up to work seamlessly with Bibby AI's citation management. Feel free to modify it for your needs!

\bibliographystyle{plain}
\bibliography{cite}

\end{document}\documentclass{article}
\usepackage[utf8]{inputenc}
\usepackage{booktabs}
\usepackage{amsmath}
\usepackage{natbib}

\title{My Document}
\author{Your Name}
\date{\today}

\begin{document}

\maketitle

\section{Introduction}
This is a sample document demonstrating various \LaTeX{} features including mathematical equations, tables, and formatting. You can cite papers using the citation search feature, and they will be automatically added to your \texttt{cite.bib} file.

\section{Mathematical Expressions}
Here are some example equations:

\subsection{Quadratic Formula}
The quadratic formula is given by:
\[x = \frac{-b \pm \sqrt{b^2 - 4ac}}{2a}\]

\subsection{Matrix Example}
A sample matrix:
\[
\begin{pmatrix}
a_{11} & a_{12} & a_{13} \\
a_{21} & a_{22} & a_{23} \\
a_{31} & a_{32} & a_{33}
\end{pmatrix}
\]

\section{Tables}
Below is a sample table showing data:

\begin{table}[h]
\centering
\begin{tabular}{@{}lrrr@{}}
\toprule
\textbf{Category} & \textbf{Value 1} & \textbf{Value 2} & \textbf{Total} \\
\midrule
Item A & 10.5 & 20.3 & 30.8 \\
Item B & 15.2 & 25.1 & 40.3 \\
Item C & 22.1 & 12.7 & 34.8 \\
\bottomrule
\end{tabular}
\caption{Sample Data Table}
\label{tab:sample}
\end{table}

\section{Lists}
\subsection{Bullet Points}
\begin{itemize}
  \item First item
  \item Second item
  \item Third item with sub-items:
    \begin{itemize}
      \item Sub-item 1
      \item Sub-item 2
    \end{itemize}
\end{itemize}

\subsection{Numbered List}
\begin{enumerate}
  \item First step
  \item Second step
  \item Third step
\end{enumerate}

\section{Citations and References}
When you use the citation search feature, references will be automatically managed. For example, you might cite a paper like this: \cite{example2024}. All citations will be automatically added to your \texttt{cite.bib} file.

\section{Conclusion}
This template demonstrates basic \LaTeX{} features including mathematical formulas, tables, and lists. The bibliography system is now set up to work seamlessly with Bibby AI's citation management. Feel free to modify it for your needs!

\bibliographystyle{plain}
\bibliography{cite}

\end{document}\documentclass{article}
\usepackage[utf8]{inputenc}
\usepackage{booktabs}
\usepackage{amsmath}
\usepackage{natbib}

\title{My Document}
\author{Your Name}
\date{\today}

\begin{document}

\maketitle

\section{Introduction}
This is a sample document demonstrating various \LaTeX{} features including mathematical equations, tables, and formatting. You can cite papers using the citation search feature, and they will be automatically added to your \texttt{cite.bib} file.

\section{Mathematical Expressions}
Here are some example equations:

\subsection{Quadratic Formula}
The quadratic formula is given by:
\[x = \frac{-b \pm \sqrt{b^2 - 4ac}}{2a}\]

\subsection{Matrix Example}
A sample matrix:
\[
\begin{pmatrix}
a_{11} & a_{12} & a_{13} \\
a_{21} & a_{22} & a_{23} \\
a_{31} & a_{32} & a_{33}
\end{pmatrix}
\]

\section{Tables}
Below is a sample table showing data:

\begin{table}[h]
\centering
\begin{tabular}{@{}lrrr@{}}
\toprule
\textbf{Category} & \textbf{Value 1} & \textbf{Value 2} & \textbf{Total} \\
\midrule
Item A & 10.5 & 20.3 & 30.8 \\
Item B & 15.2 & 25.1 & 40.3 \\
Item C & 22.1 & 12.7 & 34.8 \\
\bottomrule
\end{tabular}
\caption{Sample Data Table}
\label{tab:sample}
\end{table}

\section{Lists}
\subsection{Bullet Points}
\begin{itemize}
  \item First item
  \item Second item
  \item Third item with sub-items:
    \begin{itemize}
      \item Sub-item 1
      \item Sub-item 2
    \end{itemize}
\end{itemize}

\subsection{Numbered List}
\begin{enumerate}
  \item First step
  \item Second step
  \item Third step
\end{enumerate}

\section{Citations and References}
When you use the citation search feature, references will be automatically managed. For example, you might cite a paper like this: \cite{example2024}. All citations will be automatically added to your \texttt{cite.bib} file.

\section{Conclusion}
This template demonstrates basic \LaTeX{} features including mathematical formulas, tables, and lists. The bibliography system is now set up to work seamlessly with Bibby AI's citation management. Feel free to modify it for your needs!

\bibliographystyle{plain}
\bibliography{cite}

\end{document}\documentclass{article}
\usepackage[utf8]{inputenc}
\usepackage{booktabs}
\usepackage{amsmath}
\usepackage{natbib}

\title{My Document}
\author{Your Name}
\date{\today}

\begin{document}

\maketitle

\section{Introduction}
This is a sample document demonstrating various \LaTeX{} features including mathematical equations, tables, and formatting. You can cite papers using the citation search feature, and they will be automatically added to your \texttt{cite.bib} file.

\section{Mathematical Expressions}
Here are some example equations:

\subsection{Quadratic Formula}
The quadratic formula is given by:
\[x = \frac{-b \pm \sqrt{b^2 - 4ac}}{2a}\]

\subsection{Matrix Example}
A sample matrix:
\[
\begin{pmatrix}
a_{11} & a_{12} & a_{13} \\
a_{21} & a_{22} & a_{23} \\
a_{31} & a_{32} & a_{33}
\end{pmatrix}
\]

\section{Tables}
Below is a sample table showing data:

\begin{table}[h]
\centering
\begin{tabular}{@{}lrrr@{}}
\toprule
\textbf{Category} & \textbf{Value 1} & \textbf{Value 2} & \textbf{Total} \\
\midrule
Item A & 10.5 & 20.3 & 30.8 \\
Item B & 15.2 & 25.1 & 40.3 \\
Item C & 22.1 & 12.7 & 34.8 \\
\bottomrule
\end{tabular}
\caption{Sample Data Table}
\label{tab:sample}
\end{table}

\section{Lists}
\subsection{Bullet Points}
\begin{itemize}
  \item First item
  \item Second item
  \item Third item with sub-items:
    \begin{itemize}
      \item Sub-item 1
      \item Sub-item 2
    \end{itemize}
\end{itemize}

\subsection{Numbered List}
\begin{enumerate}
  \item First step
  \item Second step
  \item Third step
\end{enumerate}

\section{Citations and References}
When you use the citation search feature, references will be automatically managed. For example, you might cite a paper like this: \cite{example2024}. All citations will be automatically added to your \texttt{cite.bib} file.

\section{Conclusion}
This template demonstrates basic \LaTeX{} features including mathematical formulas, tables, and lists. The bibliography system is now set up to work seamlessly with Bibby AI's citation management. Feel free to modify it for your needs!

\bibliographystyle{plain}
\bibliography{cite}

\end{document}\documentclass{article}
\usepackage[utf8]{inputenc}
\usepackage{booktabs}
\usepackage{amsmath}
\usepackage{natbib}

\title{My Document}
\author{Your Name}
\date{\today}

\begin{document}

\maketitle

\section{Introduction}
This is a sample document demonstrating various \LaTeX{} features including mathematical equations, tables, and formatting. You can cite papers using the citation search feature, and they will be automatically added to your \texttt{cite.bib} file.

\section{Mathematical Expressions}
Here are some example equations:

\subsection{Quadratic Formula}
The quadratic formula is given by:
\[x = \frac{-b \pm \sqrt{b^2 - 4ac}}{2a}\]

\subsection{Matrix Example}
A sample matrix:
\[
\begin{pmatrix}
a_{11} & a_{12} & a_{13} \\
a_{21} & a_{22} & a_{23} \\
a_{31} & a_{32} & a_{33}
\end{pmatrix}
\]

\section{Tables}
Below is a sample table showing data:

\begin{table}[h]
\centering
\begin{tabular}{@{}lrrr@{}}
\toprule
\textbf{Category} & \textbf{Value 1} & \textbf{Value 2} & \textbf{Total} \\
\midrule
Item A & 10.5 & 20.3 & 30.8 \\
Item B & 15.2 & 25.1 & 40.3 \\
Item C & 22.1 & 12.7 & 34.8 \\
\bottomrule
\end{tabular}
\caption{Sample Data Table}
\label{tab:sample}
\end{table}

\section{Lists}
\subsection{Bullet Points}
\begin{itemize}
  \item First item
  \item Second item
  \item Third item with sub-items:
    \begin{itemize}
      \item Sub-item 1
      \item Sub-item 2
    \end{itemize}
\end{itemize}

\subsection{Numbered List}
\begin{enumerate}
  \item First step
  \item Second step
  \item Third step
\end{enumerate}

\section{Citations and References}
When you use the citation search feature, references will be automatically managed. For example, you might cite a paper like this: \cite{example2024}. All citations will be automatically added to your \texttt{cite.bib} file.

\section{Conclusion}
This template demonstrates basic \LaTeX{} features including mathematical formulas, tables, and lists. The bibliography system is now set up to work seamlessly with Bibby AI's citation management. Feel free to modify it for your needs!

\bibliographystyle{plain}
\bibliography{cite}

\end{document}\documentclass{article}
\usepackage[utf8]{inputenc}
\usepackage{booktabs}
\usepackage{amsmath}
\usepackage{natbib}

\title{My Document}
\author{Your Name}
\date{\today}

\begin{document}

\maketitle

\section{Introduction}
This is a sample document demonstrating various \LaTeX{} features including mathematical equations, tables, and formatting. You can cite papers using the citation search feature, and they will be automatically added to your \texttt{cite.bib} file.

\section{Mathematical Expressions}
Here are some example equations:

\subsection{Quadratic Formula}
The quadratic formula is given by:
\[x = \frac{-b \pm \sqrt{b^2 - 4ac}}{2a}\]

\subsection{Matrix Example}
A sample matrix:
\[
\begin{pmatrix}
a_{11} & a_{12} & a_{13} \\
a_{21} & a_{22} & a_{23} \\
a_{31} & a_{32} & a_{33}
\end{pmatrix}
\]

\section{Tables}
Below is a sample table showing data:

\begin{table}[h]
\centering
\begin{tabular}{@{}lrrr@{}}
\toprule
\textbf{Category} & \textbf{Value 1} & \textbf{Value 2} & \textbf{Total} \\
\midrule
Item A & 10.5 & 20.3 & 30.8 \\
Item B & 15.2 & 25.1 & 40.3 \\
Item C & 22.1 & 12.7 & 34.8 \\
\bottomrule
\end{tabular}
\caption{Sample Data Table}
\label{tab:sample}
\end{table}

\section{Lists}
\subsection{Bullet Points}
\begin{itemize}
  \item First item
  \item Second item
  \item Third item with sub-items:
    \begin{itemize}
      \item Sub-item 1
      \item Sub-item 2
    \end{itemize}
\end{itemize}

\subsection{Numbered List}
\begin{enumerate}
  \item First step
  \item Second step
  \item Third step
\end{enumerate}

\section{Citations and References}
When you use the citation search feature, references will be automatically managed. For example, you might cite a paper like this: \cite{example2024}. All citations will be automatically added to your \texttt{cite.bib} file.

\section{Conclusion}
This template demonstrates basic \LaTeX{} features including mathematical formulas, tables, and lists. The bibliography system is now set up to work seamlessly with Bibby AI's citation management. Feel free to modify it for your needs!

\bibliographystyle{plain}
\bibliography{cite}

\end{document}\documentclass{article}
\usepackage[utf8]{inputenc}
\usepackage{booktabs}
\usepackage{amsmath}
\usepackage{natbib}

\title{My Document}
\author{Your Name}
\date{\today}

\begin{document}

\maketitle

\section{Introduction}
This is a sample document demonstrating various \LaTeX{} features including mathematical equations, tables, and formatting. You can cite papers using the citation search feature, and they will be automatically added to your \texttt{cite.bib} file.

\section{Mathematical Expressions}
Here are some example equations:

\subsection{Quadratic Formula}
The quadratic formula is given by:
\[x = \frac{-b \pm \sqrt{b^2 - 4ac}}{2a}\]

\subsection{Matrix Example}
A sample matrix:
\[
\begin{pmatrix}
a_{11} & a_{12} & a_{13} \\
a_{21} & a_{22} & a_{23} \\
a_{31} & a_{32} & a_{33}
\end{pmatrix}
\]

\section{Tables}
Below is a sample table showing data:

\begin{table}[h]
\centering
\begin{tabular}{@{}lrrr@{}}
\toprule
\textbf{Category} & \textbf{Value 1} & \textbf{Value 2} & \textbf{Total} \\
\midrule
Item A & 10.5 & 20.3 & 30.8 \\
Item B & 15.2 & 25.1 & 40.3 \\
Item C & 22.1 & 12.7 & 34.8 \\
\bottomrule
\end{tabular}
\caption{Sample Data Table}
\label{tab:sample}
\end{table}

\section{Lists}
\subsection{Bullet Points}
\begin{itemize}
  \item First item
  \item Second item
  \item Third item with sub-items:
    \begin{itemize}
      \item Sub-item 1
      \item Sub-item 2
    \end{itemize}
\end{itemize}

\subsection{Numbered List}
\begin{enumerate}
  \item First step
  \item Second step
  \item Third step
\end{enumerate}

\section{Citations and References}
When you use the citation search feature, references will be automatically managed. For example, you might cite a paper like this: \cite{example2024}. All citations will be automatically added to your \texttt{cite.bib} file.

\section{Conclusion}
This template demonstrates basic \LaTeX{} features including mathematical formulas, tables, and lists. The bibliography system is now set up to work seamlessly with Bibby AI's citation management. Feel free to modify it for your needs!

\bibliographystyle{plain}
\bibliography{cite}

\end{document}\documentclass{article}
\usepackage[utf8]{inputenc}
\usepackage{booktabs}
\usepackage{amsmath}
\usepackage{natbib}

\title{My Document}
\author{Your Name}
\date{\today}

\begin{document}

\maketitle

\section{Introduction}
This is a sample document demonstrating various \LaTeX{} features including mathematical equations, tables, and formatting. You can cite papers using the citation search feature, and they will be automatically added to your \texttt{cite.bib} file.

\section{Mathematical Expressions}
Here are some example equations:

\subsection{Quadratic Formula}
The quadratic formula is given by:
\[x = \frac{-b \pm \sqrt{b^2 - 4ac}}{2a}\]

\subsection{Matrix Example}
A sample matrix:
\[
\begin{pmatrix}
a_{11} & a_{12} & a_{13} \\
a_{21} & a_{22} & a_{23} \\
a_{31} & a_{32} & a_{33}
\end{pmatrix}
\]

\section{Tables}
Below is a sample table showing data:

\begin{table}[h]
\centering
\begin{tabular}{@{}lrrr@{}}
\toprule
\textbf{Category} & \textbf{Value 1} & \textbf{Value 2} & \textbf{Total} \\
\midrule
Item A & 10.5 & 20.3 & 30.8 \\
Item B & 15.2 & 25.1 & 40.3 \\
Item C & 22.1 & 12.7 & 34.8 \\
\bottomrule
\end{tabular}
\caption{Sample Data Table}
\label{tab:sample}
\end{table}

\section{Lists}
\subsection{Bullet Points}
\begin{itemize}
  \item First item
  \item Second item
  \item Third item with sub-items:
    \begin{itemize}
      \item Sub-item 1
      \item Sub-item 2
    \end{itemize}
\end{itemize}

\subsection{Numbered List}
\begin{enumerate}
  \item First step
  \item Second step
  \item Third step
\end{enumerate}

\section{Citations and References}
When you use the citation search feature, references will be automatically managed. For example, you might cite a paper like this: \cite{example2024}. All citations will be automatically added to your \texttt{cite.bib} file.

\section{Conclusion}
This template demonstrates basic \LaTeX{} features including mathematical formulas, tables, and lists. The bibliography system is now set up to work seamlessly with Bibby AI's citation management. Feel free to modify it for your needs!

\bibliographystyle{plain}
\bibliography{cite}

\end{document}\documentclass{article}
\usepackage[utf8]{inputenc}
\usepackage{booktabs}
\usepackage{amsmath}
\usepackage{natbib}

\title{My Document}
\author{Your Name}
\date{\today}

\begin{document}

\maketitle

\section{Introduction}
This is a sample document demonstrating various \LaTeX{} features including mathematical equations, tables, and formatting. You can cite papers using the citation search feature, and they will be automatically added to your \texttt{cite.bib} file.

\section{Mathematical Expressions}
Here are some example equations:

\subsection{Quadratic Formula}
The quadratic formula is given by:
\[x = \frac{-b \pm \sqrt{b^2 - 4ac}}{2a}\]

\subsection{Matrix Example}
A sample matrix:
\[
\begin{pmatrix}
a_{11} & a_{12} & a_{13} \\
a_{21} & a_{22} & a_{23} \\
a_{31} & a_{32} & a_{33}
\end{pmatrix}
\]

\section{Tables}
Below is a sample table showing data:

\begin{table}[h]
\centering
\begin{tabular}{@{}lrrr@{}}
\toprule
\textbf{Category} & \textbf{Value 1} & \textbf{Value 2} & \textbf{Total} \\
\midrule
Item A & 10.5 & 20.3 & 30.8 \\
Item B & 15.2 & 25.1 & 40.3 \\
Item C & 22.1 & 12.7 & 34.8 \\
\bottomrule
\end{tabular}
\caption{Sample Data Table}
\label{tab:sample}
\end{table}

\section{Lists}
\subsection{Bullet Points}
\begin{itemize}
  \item First item
  \item Second item
  \item Third item with sub-items:
    \begin{itemize}
      \item Sub-item 1
      \item Sub-item 2
    \end{itemize}
\end{itemize}

\subsection{Numbered List}
\begin{enumerate}
  \item First step
  \item Second step
  \item Third step
\end{enumerate}

\section{Citations and References}
When you use the citation search feature, references will be automatically managed. For example, you might cite a paper like this: \cite{example2024}. All citations will be automatically added to your \texttt{cite.bib} file.

\section{Conclusion}
This template demonstrates basic \LaTeX{} features including mathematical formulas, tables, and lists. The bibliography system is now set up to work seamlessly with Bibby AI's citation management. Feel free to modify it for your needs!

\bibliographystyle{plain}
\bibliography{cite}

\end{document}\documentclass{article}
\usepackage[utf8]{inputenc}
\usepackage{booktabs}
\usepackage{amsmath}
\usepackage{natbib}

\title{My Document}
\author{Your Name}
\date{\today}

\begin{document}

\maketitle

\section{Introduction}
This is a sample document demonstrating various \LaTeX{} features including mathematical equations, tables, and formatting. You can cite papers using the citation search feature, and they will be automatically added to your \texttt{cite.bib} file.

\section{Mathematical Expressions}
Here are some example equations:

\subsection{Quadratic Formula}
The quadratic formula is given by:
\[x = \frac{-b \pm \sqrt{b^2 - 4ac}}{2a}\]

\subsection{Matrix Example}
A sample matrix:
\[
\begin{pmatrix}
a_{11} & a_{12} & a_{13} \\
a_{21} & a_{22} & a_{23} \\
a_{31} & a_{32} & a_{33}
\end{pmatrix}
\]

\section{Tables}
Below is a sample table showing data:

\begin{table}[h]
\centering
\begin{tabular}{@{}lrrr@{}}
\toprule
\textbf{Category} & \textbf{Value 1} & \textbf{Value 2} & \textbf{Total} \\
\midrule
Item A & 10.5 & 20.3 & 30.8 \\
Item B & 15.2 & 25.1 & 40.3 \\
Item C & 22.1 & 12.7 & 34.8 \\
\bottomrule
\end{tabular}
\caption{Sample Data Table}
\label{tab:sample}
\end{table}

\section{Lists}
\subsection{Bullet Points}
\begin{itemize}
  \item First item
  \item Second item
  \item Third item with sub-items:
    \begin{itemize}
      \item Sub-item 1
      \item Sub-item 2
    \end{itemize}
\end{itemize}

\subsection{Numbered List}
\begin{enumerate}
  \item First step
  \item Second step
  \item Third step
\end{enumerate}

\section{Citations and References}
When you use the citation search feature, references will be automatically managed. For example, you might cite a paper like this: \cite{example2024}. All citations will be automatically added to your \texttt{cite.bib} file.

\section{Conclusion}
This template demonstrates basic \LaTeX{} features including mathematical formulas, tables, and lists. The bibliography system is now set up to work seamlessly with Bibby AI's citation management. Feel free to modify it for your needs!

\bibliographystyle{plain}
\bibliography{cite}

\end{document}\documentclass{article}
\usepackage[utf8]{inputenc}
\usepackage{booktabs}
\usepackage{amsmath}
\usepackage{natbib}

\title{My Document}
\author{Your Name}
\date{\today}

\begin{document}

\maketitle

\section{Introduction}
This is a sample document demonstrating various \LaTeX{} features including mathematical equations, tables, and formatting. You can cite papers using the citation search feature, and they will be automatically added to your \texttt{cite.bib} file.

\section{Mathematical Expressions}
Here are some example equations:

\subsection{Quadratic Formula}
The quadratic formula is given by:
\[x = \frac{-b \pm \sqrt{b^2 - 4ac}}{2a}\]

\subsection{Matrix Example}
A sample matrix:
\[
\begin{pmatrix}
a_{11} & a_{12} & a_{13} \\
a_{21} & a_{22} & a_{23} \\
a_{31} & a_{32} & a_{33}
\end{pmatrix}
\]

\section{Tables}
Below is a sample table showing data:

\begin{table}[h]
\centering
\begin{tabular}{@{}lrrr@{}}
\toprule
\textbf{Category} & \textbf{Value 1} & \textbf{Value 2} & \textbf{Total} \\
\midrule
Item A & 10.5 & 20.3 & 30.8 \\
Item B & 15.2 & 25.1 & 40.3 \\
Item C & 22.1 & 12.7 & 34.8 \\
\bottomrule
\end{tabular}
\caption{Sample Data Table}
\label{tab:sample}
\end{table}

\section{Lists}
\subsection{Bullet Points}
\begin{itemize}
  \item First item
  \item Second item
  \item Third item with sub-items:
    \begin{itemize}
      \item Sub-item 1
      \item Sub-item 2
    \end{itemize}
\end{itemize}

\subsection{Numbered List}
\begin{enumerate}
  \item First step
  \item Second step
  \item Third step
\end{enumerate}

\section{Citations and References}
When you use the citation search feature, references will be automatically managed. For example, you might cite a paper like this: \cite{example2024}. All citations will be automatically added to your \texttt{cite.bib} file.

\section{Conclusion}
This template demonstrates basic \LaTeX{} features including mathematical formulas, tables, and lists. The bibliography system is now set up to work seamlessly with Bibby AI's citation management. Feel free to modify it for your needs!

\bibliographystyle{plain}
\bibliography{cite}

\end{document}\documentclass{article}
\usepackage[utf8]{inputenc}
\usepackage{booktabs}
\usepackage{amsmath}
\usepackage{natbib}

\title{My Document}
\author{Your Name}
\date{\today}

\begin{document}

\maketitle

\section{Introduction}
This is a sample document demonstrating various \LaTeX{} features including mathematical equations, tables, and formatting. You can cite papers using the citation search feature, and they will be automatically added to your \texttt{cite.bib} file.

\section{Mathematical Expressions}
Here are some example equations:

\subsection{Quadratic Formula}
The quadratic formula is given by:
\[x = \frac{-b \pm \sqrt{b^2 - 4ac}}{2a}\]

\subsection{Matrix Example}
A sample matrix:
\[
\begin{pmatrix}
a_{11} & a_{12} & a_{13} \\
a_{21} & a_{22} & a_{23} \\
a_{31} & a_{32} & a_{33}
\end{pmatrix}
\]

\section{Tables}
Below is a sample table showing data:

\begin{table}[h]
\centering
\begin{tabular}{@{}lrrr@{}}
\toprule
\textbf{Category} & \textbf{Value 1} & \textbf{Value 2} & \textbf{Total} \\
\midrule
Item A & 10.5 & 20.3 & 30.8 \\
Item B & 15.2 & 25.1 & 40.3 \\
Item C & 22.1 & 12.7 & 34.8 \\
\bottomrule
\end{tabular}
\caption{Sample Data Table}
\label{tab:sample}
\end{table}

\section{Lists}
\subsection{Bullet Points}
\begin{itemize}
  \item First item
  \item Second item
  \item Third item with sub-items:
    \begin{itemize}
      \item Sub-item 1
      \item Sub-item 2
    \end{itemize}
\end{itemize}

\subsection{Numbered List}
\begin{enumerate}
  \item First step
  \item Second step
  \item Third step
\end{enumerate}

\section{Citations and References}
When you use the citation search feature, references will be automatically managed. For example, you might cite a paper like this: \cite{example2024}. All citations will be automatically added to your \texttt{cite.bib} file.

\section{Conclusion}
This template demonstrates basic \LaTeX{} features including mathematical formulas, tables, and lists. The bibliography system is now set up to work seamlessly with Bibby AI's citation management. Feel free to modify it for your needs!

\bibliographystyle{plain}
\bibliography{cite}

\end{document}\documentclass{article}
\usepackage[utf8]{inputenc}
\usepackage{booktabs}
\usepackage{amsmath}
\usepackage{natbib}

\title{My Document}
\author{Your Name}
\date{\today}

\begin{document}

\maketitle

\section{Introduction}
This is a sample document demonstrating various \LaTeX{} features including mathematical equations, tables, and formatting. You can cite papers using the citation search feature, and they will be automatically added to your \texttt{cite.bib} file.

\section{Mathematical Expressions}
Here are some example equations:

\subsection{Quadratic Formula}
The quadratic formula is given by:
\[x = \frac{-b \pm \sqrt{b^2 - 4ac}}{2a}\]

\subsection{Matrix Example}
A sample matrix:
\[
\begin{pmatrix}
a_{11} & a_{12} & a_{13} \\
a_{21} & a_{22} & a_{23} \\
a_{31} & a_{32} & a_{33}
\end{pmatrix}
\]

\section{Tables}
Below is a sample table showing data:

\begin{table}[h]
\centering
\begin{tabular}{@{}lrrr@{}}
\toprule
\textbf{Category} & \textbf{Value 1} & \textbf{Value 2} & \textbf{Total} \\
\midrule
Item A & 10.5 & 20.3 & 30.8 \\
Item B & 15.2 & 25.1 & 40.3 \\
Item C & 22.1 & 12.7 & 34.8 \\
\bottomrule
\end{tabular}
\caption{Sample Data Table}
\label{tab:sample}
\end{table}

\section{Lists}
\subsection{Bullet Points}
\begin{itemize}
  \item First item
  \item Second item
  \item Third item with sub-items:
    \begin{itemize}
      \item Sub-item 1
      \item Sub-item 2
    \end{itemize}
\end{itemize}

\subsection{Numbered List}
\begin{enumerate}
  \item First step
  \item Second step
  \item Third step
\end{enumerate}

\section{Citations and References}
When you use the citation search feature, references will be automatically managed. For example, you might cite a paper like this: \cite{example2024}. All citations will be automatically added to your \texttt{cite.bib} file.

\section{Conclusion}
This template demonstrates basic \LaTeX{} features including mathematical formulas, tables, and lists. The bibliography system is now set up to work seamlessly with Bibby AI's citation management. Feel free to modify it for your needs!

\bibliographystyle{plain}
\bibliography{cite}

\end{document}\documentclass{article}
\usepackage[utf8]{inputenc}
\usepackage{booktabs}
\usepackage{amsmath}
\usepackage{natbib}

\title{My Document}
\author{Your Name}
\date{\today}

\begin{document}

\maketitle

\section{Introduction}
This is a sample document demonstrating various \LaTeX{} features including mathematical equations, tables, and formatting. You can cite papers using the citation search feature, and they will be automatically added to your \texttt{cite.bib} file.

\section{Mathematical Expressions}
Here are some example equations:

\subsection{Quadratic Formula}
The quadratic formula is given by:
\[x = \frac{-b \pm \sqrt{b^2 - 4ac}}{2a}\]

\subsection{Matrix Example}
A sample matrix:
\[
\begin{pmatrix}
a_{11} & a_{12} & a_{13} \\
a_{21} & a_{22} & a_{23} \\
a_{31} & a_{32} & a_{33}
\end{pmatrix}
\]

\section{Tables}
Below is a sample table showing data:

\begin{table}[h]
\centering
\begin{tabular}{@{}lrrr@{}}
\toprule
\textbf{Category} & \textbf{Value 1} & \textbf{Value 2} & \textbf{Total} \\
\midrule
Item A & 10.5 & 20.3 & 30.8 \\
Item B & 15.2 & 25.1 & 40.3 \\
Item C & 22.1 & 12.7 & 34.8 \\
\bottomrule
\end{tabular}
\caption{Sample Data Table}
\label{tab:sample}
\end{table}

\section{Lists}
\subsection{Bullet Points}
\begin{itemize}
  \item First item
  \item Second item
  \item Third item with sub-items:
    \begin{itemize}
      \item Sub-item 1
      \item Sub-item 2
    \end{itemize}
\end{itemize}

\subsection{Numbered List}
\begin{enumerate}
  \item First step
  \item Second step
  \item Third step
\end{enumerate}

\section{Citations and References}
When you use the citation search feature, references will be automatically managed. For example, you might cite a paper like this: \cite{example2024}. All citations will be automatically added to your \texttt{cite.bib} file.

\section{Conclusion}
This template demonstrates basic \LaTeX{} features including mathematical formulas, tables, and lists. The bibliography system is now set up to work seamlessly with Bibby AI's citation management. Feel free to modify it for your needs!

\bibliographystyle{plain}
\bibliography{cite}

\end{document}\documentclass{article}
\usepackage[utf8]{inputenc}
\usepackage{booktabs}
\usepackage{amsmath}
\usepackage{natbib}

\title{My Document}
\author{Your Name}
\date{\today}

\begin{document}

\maketitle

\section{Introduction}
This is a sample document demonstrating various \LaTeX{} features including mathematical equations, tables, and formatting. You can cite papers using the citation search feature, and they will be automatically added to your \texttt{cite.bib} file.

\section{Mathematical Expressions}
Here are some example equations:

\subsection{Quadratic Formula}
The quadratic formula is given by:
\[x = \frac{-b \pm \sqrt{b^2 - 4ac}}{2a}\]

\subsection{Matrix Example}
A sample matrix:
\[
\begin{pmatrix}
a_{11} & a_{12} & a_{13} \\
a_{21} & a_{22} & a_{23} \\
a_{31} & a_{32} & a_{33}
\end{pmatrix}
\]

\section{Tables}
Below is a sample table showing data:

\begin{table}[h]
\centering
\begin{tabular}{@{}lrrr@{}}
\toprule
\textbf{Category} & \textbf{Value 1} & \textbf{Value 2} & \textbf{Total} \\
\midrule
Item A & 10.5 & 20.3 & 30.8 \\
Item B & 15.2 & 25.1 & 40.3 \\
Item C & 22.1 & 12.7 & 34.8 \\
\bottomrule
\end{tabular}
\caption{Sample Data Table}
\label{tab:sample}
\end{table}

\section{Lists}
\subsection{Bullet Points}
\begin{itemize}
  \item First item
  \item Second item
  \item Third item with sub-items:
    \begin{itemize}
      \item Sub-item 1
      \item Sub-item 2
    \end{itemize}
\end{itemize}

\subsection{Numbered List}
\begin{enumerate}
  \item First step
  \item Second step
  \item Third step
\end{enumerate}

\section{Citations and References}
When you use the citation search feature, references will be automatically managed. For example, you might cite a paper like this: \cite{example2024}. All citations will be automatically added to your \texttt{cite.bib} file.

\section{Conclusion}
This template demonstrates basic \LaTeX{} features including mathematical formulas, tables, and lists. The bibliography system is now set up to work seamlessly with Bibby AI's citation management. Feel free to modify it for your needs!

\bibliographystyle{plain}
\bibliography{cite}

\end{document}\documentclass{article}
\usepackage[utf8]{inputenc}
\usepackage{booktabs}
\usepackage{amsmath}
\usepackage{natbib}

\title{My Document}
\author{Your Name}
\date{\today}

\begin{document}

\maketitle

\section{Introduction}
This is a sample document demonstrating various \LaTeX{} features including mathematical equations, tables, and formatting. You can cite papers using the citation search feature, and they will be automatically added to your \texttt{cite.bib} file.

\section{Mathematical Expressions}
Here are some example equations:

\subsection{Quadratic Formula}
The quadratic formula is given by:
\[x = \frac{-b \pm \sqrt{b^2 - 4ac}}{2a}\]

\subsection{Matrix Example}
A sample matrix:
\[
\begin{pmatrix}
a_{11} & a_{12} & a_{13} \\
a_{21} & a_{22} & a_{23} \\
a_{31} & a_{32} & a_{33}
\end{pmatrix}
\]

\section{Tables}
Below is a sample table showing data:

\begin{table}[h]
\centering
\begin{tabular}{@{}lrrr@{}}
\toprule
\textbf{Category} & \textbf{Value 1} & \textbf{Value 2} & \textbf{Total} \\
\midrule
Item A & 10.5 & 20.3 & 30.8 \\
Item B & 15.2 & 25.1 & 40.3 \\
Item C & 22.1 & 12.7 & 34.8 \\
\bottomrule
\end{tabular}
\caption{Sample Data Table}
\label{tab:sample}
\end{table}

\section{Lists}
\subsection{Bullet Points}
\begin{itemize}
  \item First item
  \item Second item
  \item Third item with sub-items:
    \begin{itemize}
      \item Sub-item 1
      \item Sub-item 2
    \end{itemize}
\end{itemize}

\subsection{Numbered List}
\begin{enumerate}
  \item First step
  \item Second step
  \item Third step
\end{enumerate}

\section{Citations and References}
When you use the citation search feature, references will be automatically managed. For example, you might cite a paper like this: \cite{example2024}. All citations will be automatically added to your \texttt{cite.bib} file.

\section{Conclusion}
This template demonstrates basic \LaTeX{} features including mathematical formulas, tables, and lists. The bibliography system is now set up to work seamlessly with Bibby AI's citation management. Feel free to modify it for your needs!

\bibliographystyle{plain}
\bibliography{cite}

\end{document}\documentclass{article}
\usepackage[utf8]{inputenc}
\usepackage{booktabs}
\usepackage{amsmath}
\usepackage{natbib}

\title{My Document}
\author{Your Name}
\date{\today}

\begin{document}

\maketitle

\section{Introduction}
This is a sample document demonstrating various \LaTeX{} features including mathematical equations, tables, and formatting. You can cite papers using the citation search feature, and they will be automatically added to your \texttt{cite.bib} file.

\section{Mathematical Expressions}
Here are some example equations:

\subsection{Quadratic Formula}
The quadratic formula is given by:
\[x = \frac{-b \pm \sqrt{b^2 - 4ac}}{2a}\]

\subsection{Matrix Example}
A sample matrix:
\[
\begin{pmatrix}
a_{11} & a_{12} & a_{13} \\
a_{21} & a_{22} & a_{23} \\
a_{31} & a_{32} & a_{33}
\end{pmatrix}
\]

\section{Tables}
Below is a sample table showing data:

\begin{table}[h]
\centering
\begin{tabular}{@{}lrrr@{}}
\toprule
\textbf{Category} & \textbf{Value 1} & \textbf{Value 2} & \textbf{Total} \\
\midrule
Item A & 10.5 & 20.3 & 30.8 \\
Item B & 15.2 & 25.1 & 40.3 \\
Item C & 22.1 & 12.7 & 34.8 \\
\bottomrule
\end{tabular}
\caption{Sample Data Table}
\label{tab:sample}
\end{table}

\section{Lists}
\subsection{Bullet Points}
\begin{itemize}
  \item First item
  \item Second item
  \item Third item with sub-items:
    \begin{itemize}
      \item Sub-item 1
      \item Sub-item 2
    \end{itemize}
\end{itemize}

\subsection{Numbered List}
\begin{enumerate}
  \item First step
  \item Second step
  \item Third step
\end{enumerate}

\section{Citations and References}
When you use the citation search feature, references will be automatically managed. For example, you might cite a paper like this: \cite{example2024}. All citations will be automatically added to your \texttt{cite.bib} file.

\section{Conclusion}
This template demonstrates basic \LaTeX{} features including mathematical formulas, tables, and lists. The bibliography system is now set up to work seamlessly with Bibby AI's citation management. Feel free to modify it for your needs!

\bibliographystyle{plain}
\bibliography{cite}

\end{document}\documentclass{article}
\usepackage[utf8]{inputenc}
\usepackage{booktabs}
\usepackage{amsmath}
\usepackage{natbib}

\title{My Document}
\author{Your Name}
\date{\today}

\begin{document}

\maketitle

\section{Introduction}
This is a sample document demonstrating various \LaTeX{} features including mathematical equations, tables, and formatting. You can cite papers using the citation search feature, and they will be automatically added to your \texttt{cite.bib} file.

\section{Mathematical Expressions}
Here are some example equations:

\subsection{Quadratic Formula}
The quadratic formula is given by:
\[x = \frac{-b \pm \sqrt{b^2 - 4ac}}{2a}\]

\subsection{Matrix Example}
A sample matrix:
\[
\begin{pmatrix}
a_{11} & a_{12} & a_{13} \\
a_{21} & a_{22} & a_{23} \\
a_{31} & a_{32} & a_{33}
\end{pmatrix}
\]

\section{Tables}
Below is a sample table showing data:

\begin{table}[h]
\centering
\begin{tabular}{@{}lrrr@{}}
\toprule
\textbf{Category} & \textbf{Value 1} & \textbf{Value 2} & \textbf{Total} \\
\midrule
Item A & 10.5 & 20.3 & 30.8 \\
Item B & 15.2 & 25.1 & 40.3 \\
Item C & 22.1 & 12.7 & 34.8 \\
\bottomrule
\end{tabular}
\caption{Sample Data Table}
\label{tab:sample}
\end{table}

\section{Lists}
\subsection{Bullet Points}
\begin{itemize}
  \item First item
  \item Second item
  \item Third item with sub-items:
    \begin{itemize}
      \item Sub-item 1
      \item Sub-item 2
    \end{itemize}
\end{itemize}

\subsection{Numbered List}
\begin{enumerate}
  \item First step
  \item Second step
  \item Third step
\end{enumerate}

\section{Citations and References}
When you use the citation search feature, references will be automatically managed. For example, you might cite a paper like this: \cite{example2024}. All citations will be automatically added to your \texttt{cite.bib} file.

\section{Conclusion}
This template demonstrates basic \LaTeX{} features including mathematical formulas, tables, and lists. The bibliography system is now set up to work seamlessly with Bibby AI's citation management. Feel free to modify it for your needs!

\bibliographystyle{plain}
\bibliography{cite}

\end{document}\documentclass{article}
\usepackage[utf8]{inputenc}
\usepackage{booktabs}
\usepackage{amsmath}
\usepackage{natbib}

\title{My Document}
\author{Your Name}
\date{\today}

\begin{document}

\maketitle

\section{Introduction}
This is a sample document demonstrating various \LaTeX{} features including mathematical equations, tables, and formatting. You can cite papers using the citation search feature, and they will be automatically added to your \texttt{cite.bib} file.

\section{Mathematical Expressions}
Here are some example equations:

\subsection{Quadratic Formula}
The quadratic formula is given by:
\[x = \frac{-b \pm \sqrt{b^2 - 4ac}}{2a}\]

\subsection{Matrix Example}
A sample matrix:
\[
\begin{pmatrix}
a_{11} & a_{12} & a_{13} \\
a_{21} & a_{22} & a_{23} \\
a_{31} & a_{32} & a_{33}
\end{pmatrix}
\]

\section{Tables}
Below is a sample table showing data:

\begin{table}[h]
\centering
\begin{tabular}{@{}lrrr@{}}
\toprule
\textbf{Category} & \textbf{Value 1} & \textbf{Value 2} & \textbf{Total} \\
\midrule
Item A & 10.5 & 20.3 & 30.8 \\
Item B & 15.2 & 25.1 & 40.3 \\
Item C & 22.1 & 12.7 & 34.8 \\
\bottomrule
\end{tabular}
\caption{Sample Data Table}
\label{tab:sample}
\end{table}

\section{Lists}
\subsection{Bullet Points}
\begin{itemize}
  \item First item
  \item Second item
  \item Third item with sub-items:
    \begin{itemize}
      \item Sub-item 1
      \item Sub-item 2
    \end{itemize}
\end{itemize}

\subsection{Numbered List}
\begin{enumerate}
  \item First step
  \item Second step
  \item Third step
\end{enumerate}

\section{Citations and References}
When you use the citation search feature, references will be automatically managed. For example, you might cite a paper like this: \cite{example2024}. All citations will be automatically added to your \texttt{cite.bib} file.

\section{Conclusion}
This template demonstrates basic \LaTeX{} features including mathematical formulas, tables, and lists. The bibliography system is now set up to work seamlessly with Bibby AI's citation management. Feel free to modify it for your needs!

\bibliographystyle{plain}
\bibliography{cite}

\end{document}\documentclass{article}
\usepackage[utf8]{inputenc}
\usepackage{booktabs}
\usepackage{amsmath}
\usepackage{natbib}

\title{My Document}
\author{Your Name}
\date{\today}

\begin{document}

\maketitle

\section{Introduction}
This is a sample document demonstrating various \LaTeX{} features including mathematical equations, tables, and formatting. You can cite papers using the citation search feature, and they will be automatically added to your \texttt{cite.bib} file.

\section{Mathematical Expressions}
Here are some example equations:

\subsection{Quadratic Formula}
The quadratic formula is given by:
\[x = \frac{-b \pm \sqrt{b^2 - 4ac}}{2a}\]

\subsection{Matrix Example}
A sample matrix:
\[
\begin{pmatrix}
a_{11} & a_{12} & a_{13} \\
a_{21} & a_{22} & a_{23} \\
a_{31} & a_{32} & a_{33}
\end{pmatrix}
\]

\section{Tables}
Below is a sample table showing data:

\begin{table}[h]
\centering
\begin{tabular}{@{}lrrr@{}}
\toprule
\textbf{Category} & \textbf{Value 1} & \textbf{Value 2} & \textbf{Total} \\
\midrule
Item A & 10.5 & 20.3 & 30.8 \\
Item B & 15.2 & 25.1 & 40.3 \\
Item C & 22.1 & 12.7 & 34.8 \\
\bottomrule
\end{tabular}
\caption{Sample Data Table}
\label{tab:sample}
\end{table}

\section{Lists}
\subsection{Bullet Points}
\begin{itemize}
  \item First item
  \item Second item
  \item Third item with sub-items:
    \begin{itemize}
      \item Sub-item 1
      \item Sub-item 2
    \end{itemize}
\end{itemize}

\subsection{Numbered List}
\begin{enumerate}
  \item First step
  \item Second step
  \item Third step
\end{enumerate}

\section{Citations and References}
When you use the citation search feature, references will be automatically managed. For example, you might cite a paper like this: \cite{example2024}. All citations will be automatically added to your \texttt{cite.bib} file.

\section{Conclusion}
This template demonstrates basic \LaTeX{} features including mathematical formulas, tables, and lists. The bibliography system is now set up to work seamlessly with Bibby AI's citation management. Feel free to modify it for your needs!

\bibliographystyle{plain}
\bibliography{cite}

\end{document}\documentclass{article}
\usepackage[utf8]{inputenc}
\usepackage{booktabs}
\usepackage{amsmath}
\usepackage{natbib}

\title{My Document}
\author{Your Name}
\date{\today}

\begin{document}

\maketitle

\section{Introduction}
This is a sample document demonstrating various \LaTeX{} features including mathematical equations, tables, and formatting. You can cite papers using the citation search feature, and they will be automatically added to your \texttt{cite.bib} file.

\section{Mathematical Expressions}
Here are some example equations:

\subsection{Quadratic Formula}
The quadratic formula is given by:
\[x = \frac{-b \pm \sqrt{b^2 - 4ac}}{2a}\]

\subsection{Matrix Example}
A sample matrix:
\[
\begin{pmatrix}
a_{11} & a_{12} & a_{13} \\
a_{21} & a_{22} & a_{23} \\
a_{31} & a_{32} & a_{33}
\end{pmatrix}
\]

\section{Tables}
Below is a sample table showing data:

\begin{table}[h]
\centering
\begin{tabular}{@{}lrrr@{}}
\toprule
\textbf{Category} & \textbf{Value 1} & \textbf{Value 2} & \textbf{Total} \\
\midrule
Item A & 10.5 & 20.3 & 30.8 \\
Item B & 15.2 & 25.1 & 40.3 \\
Item C & 22.1 & 12.7 & 34.8 \\
\bottomrule
\end{tabular}
\caption{Sample Data Table}
\label{tab:sample}
\end{table}

\section{Lists}
\subsection{Bullet Points}
\begin{itemize}
  \item First item
  \item Second item
  \item Third item with sub-items:
    \begin{itemize}
      \item Sub-item 1
      \item Sub-item 2
    \end{itemize}
\end{itemize}

\subsection{Numbered List}
\begin{enumerate}
  \item First step
  \item Second step
  \item Third step
\end{enumerate}

\section{Citations and References}
When you use the citation search feature, references will be automatically managed. For example, you might cite a paper like this: \cite{example2024}. All citations will be automatically added to your \texttt{cite.bib} file.

\section{Conclusion}
This template demonstrates basic \LaTeX{} features including mathematical formulas, tables, and lists. The bibliography system is now set up to work seamlessly with Bibby AI's citation management. Feel free to modify it for your needs!

\bibliographystyle{plain}
\bibliography{cite}

\end{document}\documentclass{article}
\usepackage[utf8]{inputenc}
\usepackage{booktabs}
\usepackage{amsmath}
\usepackage{natbib}

\title{My Document}
\author{Your Name}
\date{\today}

\begin{document}

\maketitle

\section{Introduction}
This is a sample document demonstrating various \LaTeX{} features including mathematical equations, tables, and formatting. You can cite papers using the citation search feature, and they will be automatically added to your \texttt{cite.bib} file.

\section{Mathematical Expressions}
Here are some example equations:

\subsection{Quadratic Formula}
The quadratic formula is given by:
\[x = \frac{-b \pm \sqrt{b^2 - 4ac}}{2a}\]

\subsection{Matrix Example}
A sample matrix:
\[
\begin{pmatrix}
a_{11} & a_{12} & a_{13} \\
a_{21} & a_{22} & a_{23} \\
a_{31} & a_{32} & a_{33}
\end{pmatrix}
\]

\section{Tables}
Below is a sample table showing data:

\begin{table}[h]
\centering
\begin{tabular}{@{}lrrr@{}}
\toprule
\textbf{Category} & \textbf{Value 1} & \textbf{Value 2} & \textbf{Total} \\
\midrule
Item A & 10.5 & 20.3 & 30.8 \\
Item B & 15.2 & 25.1 & 40.3 \\
Item C & 22.1 & 12.7 & 34.8 \\
\bottomrule
\end{tabular}
\caption{Sample Data Table}
\label{tab:sample}
\end{table}

\section{Lists}
\subsection{Bullet Points}
\begin{itemize}
  \item First item
  \item Second item
  \item Third item with sub-items:
    \begin{itemize}
      \item Sub-item 1
      \item Sub-item 2
    \end{itemize}
\end{itemize}

\subsection{Numbered List}
\begin{enumerate}
  \item First step
  \item Second step
  \item Third step
\end{enumerate}

\section{Citations and References}
When you use the citation search feature, references will be automatically managed. For example, you might cite a paper like this: \cite{example2024}. All citations will be automatically added to your \texttt{cite.bib} file.

\section{Conclusion}
This template demonstrates basic \LaTeX{} features including mathematical formulas, tables, and lists. The bibliography system is now set up to work seamlessly with Bibby AI's citation management. Feel free to modify it for your needs!

\bibliographystyle{plain}
\bibliography{cite}

\end{document}\documentclass{article}
\usepackage[utf8]{inputenc}
\usepackage{booktabs}
\usepackage{amsmath}
\usepackage{natbib}

\title{My Document}
\author{Your Name}
\date{\today}

\begin{document}

\maketitle

\section{Introduction}
This is a sample document demonstrating various \LaTeX{} features including mathematical equations, tables, and formatting. You can cite papers using the citation search feature, and they will be automatically added to your \texttt{cite.bib} file.

\section{Mathematical Expressions}
Here are some example equations:

\subsection{Quadratic Formula}
The quadratic formula is given by:
\[x = \frac{-b \pm \sqrt{b^2 - 4ac}}{2a}\]

\subsection{Matrix Example}
A sample matrix:
\[
\begin{pmatrix}
a_{11} & a_{12} & a_{13} \\
a_{21} & a_{22} & a_{23} \\
a_{31} & a_{32} & a_{33}
\end{pmatrix}
\]

\section{Tables}
Below is a sample table showing data:

\begin{table}[h]
\centering
\begin{tabular}{@{}lrrr@{}}
\toprule
\textbf{Category} & \textbf{Value 1} & \textbf{Value 2} & \textbf{Total} \\
\midrule
Item A & 10.5 & 20.3 & 30.8 \\
Item B & 15.2 & 25.1 & 40.3 \\
Item C & 22.1 & 12.7 & 34.8 \\
\bottomrule
\end{tabular}
\caption{Sample Data Table}
\label{tab:sample}
\end{table}

\section{Lists}
\subsection{Bullet Points}
\begin{itemize}
  \item First item
  \item Second item
  \item Third item with sub-items:
    \begin{itemize}
      \item Sub-item 1
      \item Sub-item 2
    \end{itemize}
\end{itemize}

\subsection{Numbered List}
\begin{enumerate}
  \item First step
  \item Second step
  \item Third step
\end{enumerate}

\section{Citations and References}
When you use the citation search feature, references will be automatically managed. For example, you might cite a paper like this: \cite{example2024}. All citations will be automatically added to your \texttt{cite.bib} file.

\section{Conclusion}
This template demonstrates basic \LaTeX{} features including mathematical formulas, tables, and lists. The bibliography system is now set up to work seamlessly with Bibby AI's citation management. Feel free to modify it for your needs!

\bibliographystyle{plain}
\bibliography{cite}

\end{document}\documentclass{article}
\usepackage[utf8]{inputenc}
\usepackage{booktabs}
\usepackage{amsmath}
\usepackage{natbib}

\title{My Document}
\author{Your Name}
\date{\today}

\begin{document}

\maketitle

\section{Introduction}
This is a sample document demonstrating various \LaTeX{} features including mathematical equations, tables, and formatting. You can cite papers using the citation search feature, and they will be automatically added to your \texttt{cite.bib} file.

\section{Mathematical Expressions}
Here are some example equations:

\subsection{Quadratic Formula}
The quadratic formula is given by:
\[x = \frac{-b \pm \sqrt{b^2 - 4ac}}{2a}\]

\subsection{Matrix Example}
A sample matrix:
\[
\begin{pmatrix}
a_{11} & a_{12} & a_{13} \\
a_{21} & a_{22} & a_{23} \\
a_{31} & a_{32} & a_{33}
\end{pmatrix}
\]

\section{Tables}
Below is a sample table showing data:

\begin{table}[h]
\centering
\begin{tabular}{@{}lrrr@{}}
\toprule
\textbf{Category} & \textbf{Value 1} & \textbf{Value 2} & \textbf{Total} \\
\midrule
Item A & 10.5 & 20.3 & 30.8 \\
Item B & 15.2 & 25.1 & 40.3 \\
Item C & 22.1 & 12.7 & 34.8 \\
\bottomrule
\end{tabular}
\caption{Sample Data Table}
\label{tab:sample}
\end{table}

\section{Lists}
\subsection{Bullet Points}
\begin{itemize}
  \item First item
  \item Second item
  \item Third item with sub-items:
    \begin{itemize}
      \item Sub-item 1
      \item Sub-item 2
    \end{itemize}
\end{itemize}

\subsection{Numbered List}
\begin{enumerate}
  \item First step
  \item Second step
  \item Third step
\end{enumerate}

\section{Citations and References}
When you use the citation search feature, references will be automatically managed. For example, you might cite a paper like this: \cite{example2024}. All citations will be automatically added to your \texttt{cite.bib} file.

\section{Conclusion}
This template demonstrates basic \LaTeX{} features including mathematical formulas, tables, and lists. The bibliography system is now set up to work seamlessly with Bibby AI's citation management. Feel free to modify it for your needs!

\bibliographystyle{plain}
\bibliography{cite}

\end{document}\documentclass{article}
\usepackage[utf8]{inputenc}
\usepackage{booktabs}
\usepackage{amsmath}
\usepackage{natbib}

\title{My Document}
\author{Your Name}
\date{\today}

\begin{document}

\maketitle

\section{Introduction}
This is a sample document demonstrating various \LaTeX{} features including mathematical equations, tables, and formatting. You can cite papers using the citation search feature, and they will be automatically added to your \texttt{cite.bib} file.

\section{Mathematical Expressions}
Here are some example equations:

\subsection{Quadratic Formula}
The quadratic formula is given by:
\[x = \frac{-b \pm \sqrt{b^2 - 4ac}}{2a}\]

\subsection{Matrix Example}
A sample matrix:
\[
\begin{pmatrix}
a_{11} & a_{12} & a_{13} \\
a_{21} & a_{22} & a_{23} \\
a_{31} & a_{32} & a_{33}
\end{pmatrix}
\]

\section{Tables}
Below is a sample table showing data:

\begin{table}[h]
\centering
\begin{tabular}{@{}lrrr@{}}
\toprule
\textbf{Category} & \textbf{Value 1} & \textbf{Value 2} & \textbf{Total} \\
\midrule
Item A & 10.5 & 20.3 & 30.8 \\
Item B & 15.2 & 25.1 & 40.3 \\
Item C & 22.1 & 12.7 & 34.8 \\
\bottomrule
\end{tabular}
\caption{Sample Data Table}
\label{tab:sample}
\end{table}

\section{Lists}
\subsection{Bullet Points}
\begin{itemize}
  \item First item
  \item Second item
  \item Third item with sub-items:
    \begin{itemize}
      \item Sub-item 1
      \item Sub-item 2
    \end{itemize}
\end{itemize}

\subsection{Numbered List}
\begin{enumerate}
  \item First step
  \item Second step
  \item Third step
\end{enumerate}

\section{Citations and References}
When you use the citation search feature, references will be automatically managed. For example, you might cite a paper like this: \cite{example2024}. All citations will be automatically added to your \texttt{cite.bib} file.

\section{Conclusion}
This template demonstrates basic \LaTeX{} features including mathematical formulas, tables, and lists. The bibliography system is now set up to work seamlessly with Bibby AI's citation management. Feel free to modify it for your needs!

\bibliographystyle{plain}
\bibliography{cite}

\end{document}\documentclass{article}
\usepackage[utf8]{inputenc}
\usepackage{booktabs}
\usepackage{amsmath}
\usepackage{natbib}

\title{My Document}
\author{Your Name}
\date{\today}

\begin{document}

\maketitle

\section{Introduction}
This is a sample document demonstrating various \LaTeX{} features including mathematical equations, tables, and formatting. You can cite papers using the citation search feature, and they will be automatically added to your \texttt{cite.bib} file.

\section{Mathematical Expressions}
Here are some example equations:

\subsection{Quadratic Formula}
The quadratic formula is given by:
\[x = \frac{-b \pm \sqrt{b^2 - 4ac}}{2a}\]

\subsection{Matrix Example}
A sample matrix:
\[
\begin{pmatrix}
a_{11} & a_{12} & a_{13} \\
a_{21} & a_{22} & a_{23} \\
a_{31} & a_{32} & a_{33}
\end{pmatrix}
\]

\section{Tables}
Below is a sample table showing data:

\begin{table}[h]
\centering
\begin{tabular}{@{}lrrr@{}}
\toprule
\textbf{Category} & \textbf{Value 1} & \textbf{Value 2} & \textbf{Total} \\
\midrule
Item A & 10.5 & 20.3 & 30.8 \\
Item B & 15.2 & 25.1 & 40.3 \\
Item C & 22.1 & 12.7 & 34.8 \\
\bottomrule
\end{tabular}
\caption{Sample Data Table}
\label{tab:sample}
\end{table}

\section{Lists}
\subsection{Bullet Points}
\begin{itemize}
  \item First item
  \item Second item
  \item Third item with sub-items:
    \begin{itemize}
      \item Sub-item 1
      \item Sub-item 2
    \end{itemize}
\end{itemize}

\subsection{Numbered List}
\begin{enumerate}
  \item First step
  \item Second step
  \item Third step
\end{enumerate}

\section{Citations and References}
When you use the citation search feature, references will be automatically managed. For example, you might cite a paper like this: \cite{example2024}. All citations will be automatically added to your \texttt{cite.bib} file.

\section{Conclusion}
This template demonstrates basic \LaTeX{} features including mathematical formulas, tables, and lists. The bibliography system is now set up to work seamlessly with Bibby AI's citation management. Feel free to modify it for your needs!

\bibliographystyle{plain}
\bibliography{cite}

\end{document}\documentclass{article}
\usepackage[utf8]{inputenc}
\usepackage{booktabs}
\usepackage{amsmath}
\usepackage{natbib}

\title{My Document}
\author{Your Name}
\date{\today}

\begin{document}

\maketitle

\section{Introduction}
This is a sample document demonstrating various \LaTeX{} features including mathematical equations, tables, and formatting. You can cite papers using the citation search feature, and they will be automatically added to your \texttt{cite.bib} file.

\section{Mathematical Expressions}
Here are some example equations:

\subsection{Quadratic Formula}
The quadratic formula is given by:
\[x = \frac{-b \pm \sqrt{b^2 - 4ac}}{2a}\]

\subsection{Matrix Example}
A sample matrix:
\[
\begin{pmatrix}
a_{11} & a_{12} & a_{13} \\
a_{21} & a_{22} & a_{23} \\
a_{31} & a_{32} & a_{33}
\end{pmatrix}
\]

\section{Tables}
Below is a sample table showing data:

\begin{table}[h]
\centering
\begin{tabular}{@{}lrrr@{}}
\toprule
\textbf{Category} & \textbf{Value 1} & \textbf{Value 2} & \textbf{Total} \\
\midrule
Item A & 10.5 & 20.3 & 30.8 \\
Item B & 15.2 & 25.1 & 40.3 \\
Item C & 22.1 & 12.7 & 34.8 \\
\bottomrule
\end{tabular}
\caption{Sample Data Table}
\label{tab:sample}
\end{table}

\section{Lists}
\subsection{Bullet Points}
\begin{itemize}
  \item First item
  \item Second item
  \item Third item with sub-items:
    \begin{itemize}
      \item Sub-item 1
      \item Sub-item 2
    \end{itemize}
\end{itemize}

\subsection{Numbered List}
\begin{enumerate}
  \item First step
  \item Second step
  \item Third step
\end{enumerate}

\section{Citations and References}
When you use the citation search feature, references will be automatically managed. For example, you might cite a paper like this: \cite{example2024}. All citations will be automatically added to your \texttt{cite.bib} file.

\section{Conclusion}
This template demonstrates basic \LaTeX{} features including mathematical formulas, tables, and lists. The bibliography system is now set up to work seamlessly with Bibby AI's citation management. Feel free to modify it for your needs!

\bibliographystyle{plain}
\bibliography{cite}

\end{document}\documentclass{article}
\usepackage[utf8]{inputenc}
\usepackage{booktabs}
\usepackage{amsmath}
\usepackage{natbib}

\title{My Document}
\author{Your Name}
\date{\today}

\begin{document}

\maketitle

\section{Introduction}
This is a sample document demonstrating various \LaTeX{} features including mathematical equations, tables, and formatting. You can cite papers using the citation search feature, and they will be automatically added to your \texttt{cite.bib} file.

\section{Mathematical Expressions}
Here are some example equations:

\subsection{Quadratic Formula}
The quadratic formula is given by:
\[x = \frac{-b \pm \sqrt{b^2 - 4ac}}{2a}\]

\subsection{Matrix Example}
A sample matrix:
\[
\begin{pmatrix}
a_{11} & a_{12} & a_{13} \\
a_{21} & a_{22} & a_{23} \\
a_{31} & a_{32} & a_{33}
\end{pmatrix}
\]

\section{Tables}
Below is a sample table showing data:

\begin{table}[h]
\centering
\begin{tabular}{@{}lrrr@{}}
\toprule
\textbf{Category} & \textbf{Value 1} & \textbf{Value 2} & \textbf{Total} \\
\midrule
Item A & 10.5 & 20.3 & 30.8 \\
Item B & 15.2 & 25.1 & 40.3 \\
Item C & 22.1 & 12.7 & 34.8 \\
\bottomrule
\end{tabular}
\caption{Sample Data Table}
\label{tab:sample}
\end{table}

\section{Lists}
\subsection{Bullet Points}
\begin{itemize}
  \item First item
  \item Second item
  \item Third item with sub-items:
    \begin{itemize}
      \item Sub-item 1
      \item Sub-item 2
    \end{itemize}
\end{itemize}

\subsection{Numbered List}
\begin{enumerate}
  \item First step
  \item Second step
  \item Third step
\end{enumerate}

\section{Citations and References}
When you use the citation search feature, references will be automatically managed. For example, you might cite a paper like this: \cite{example2024}. All citations will be automatically added to your \texttt{cite.bib} file.

\section{Conclusion}
This template demonstrates basic \LaTeX{} features including mathematical formulas, tables, and lists. The bibliography system is now set up to work seamlessly with Bibby AI's citation management. Feel free to modify it for your needs!

\bibliographystyle{plain}
\bibliography{cite}

\end{document}\documentclass{article}
\usepackage[utf8]{inputenc}
\usepackage{booktabs}
\usepackage{amsmath}
\usepackage{natbib}

\title{My Document}
\author{Your Name}
\date{\today}

\begin{document}

\maketitle

\section{Introduction}
This is a sample document demonstrating various \LaTeX{} features including mathematical equations, tables, and formatting. You can cite papers using the citation search feature, and they will be automatically added to your \texttt{cite.bib} file.

\section{Mathematical Expressions}
Here are some example equations:

\subsection{Quadratic Formula}
The quadratic formula is given by:
\[x = \frac{-b \pm \sqrt{b^2 - 4ac}}{2a}\]

\subsection{Matrix Example}
A sample matrix:
\[
\begin{pmatrix}
a_{11} & a_{12} & a_{13} \\
a_{21} & a_{22} & a_{23} \\
a_{31} & a_{32} & a_{33}
\end{pmatrix}
\]

\section{Tables}
Below is a sample table showing data:

\begin{table}[h]
\centering
\begin{tabular}{@{}lrrr@{}}
\toprule
\textbf{Category} & \textbf{Value 1} & \textbf{Value 2} & \textbf{Total} \\
\midrule
Item A & 10.5 & 20.3 & 30.8 \\
Item B & 15.2 & 25.1 & 40.3 \\
Item C & 22.1 & 12.7 & 34.8 \\
\bottomrule
\end{tabular}
\caption{Sample Data Table}
\label{tab:sample}
\end{table}

\section{Lists}
\subsection{Bullet Points}
\begin{itemize}
  \item First item
  \item Second item
  \item Third item with sub-items:
    \begin{itemize}
      \item Sub-item 1
      \item Sub-item 2
    \end{itemize}
\end{itemize}

\subsection{Numbered List}
\begin{enumerate}
  \item First step
  \item Second step
  \item Third step
\end{enumerate}

\section{Citations and References}
When you use the citation search feature, references will be automatically managed. For example, you might cite a paper like this: \cite{example2024}. All citations will be automatically added to your \texttt{cite.bib} file.

\section{Conclusion}
This template demonstrates basic \LaTeX{} features including mathematical formulas, tables, and lists. The bibliography system is now set up to work seamlessly with Bibby AI's citation management. Feel free to modify it for your needs!

\bibliographystyle{plain}
\bibliography{cite}

\end{document}\documentclass{article}
\usepackage[utf8]{inputenc}
\usepackage{booktabs}
\usepackage{amsmath}
\usepackage{natbib}

\title{My Document}
\author{Your Name}
\date{\today}

\begin{document}

\maketitle

\section{Introduction}
This is a sample document demonstrating various \LaTeX{} features including mathematical equations, tables, and formatting. You can cite papers using the citation search feature, and they will be automatically added to your \texttt{cite.bib} file.

\section{Mathematical Expressions}
Here are some example equations:

\subsection{Quadratic Formula}
The quadratic formula is given by:
\[x = \frac{-b \pm \sqrt{b^2 - 4ac}}{2a}\]

\subsection{Matrix Example}
A sample matrix:
\[
\begin{pmatrix}
a_{11} & a_{12} & a_{13} \\
a_{21} & a_{22} & a_{23} \\
a_{31} & a_{32} & a_{33}
\end{pmatrix}
\]

\section{Tables}
Below is a sample table showing data:

\begin{table}[h]
\centering
\begin{tabular}{@{}lrrr@{}}
\toprule
\textbf{Category} & \textbf{Value 1} & \textbf{Value 2} & \textbf{Total} \\
\midrule
Item A & 10.5 & 20.3 & 30.8 \\
Item B & 15.2 & 25.1 & 40.3 \\
Item C & 22.1 & 12.7 & 34.8 \\
\bottomrule
\end{tabular}
\caption{Sample Data Table}
\label{tab:sample}
\end{table}

\section{Lists}
\subsection{Bullet Points}
\begin{itemize}
  \item First item
  \item Second item
  \item Third item with sub-items:
    \begin{itemize}
      \item Sub-item 1
      \item Sub-item 2
    \end{itemize}
\end{itemize}

\subsection{Numbered List}
\begin{enumerate}
  \item First step
  \item Second step
  \item Third step
\end{enumerate}

\section{Citations and References}
When you use the citation search feature, references will be automatically managed. For example, you might cite a paper like this: \cite{example2024}. All citations will be automatically added to your \texttt{cite.bib} file.

\section{Conclusion}
This template demonstrates basic \LaTeX{} features including mathematical formulas, tables, and lists. The bibliography system is now set up to work seamlessly with Bibby AI's citation management. Feel free to modify it for your needs!

\bibliographystyle{plain}
\bibliography{cite}

\end{document}\documentclass{article}
\usepackage[utf8]{inputenc}
\usepackage{booktabs}
\usepackage{amsmath}
\usepackage{natbib}

\title{My Document}
\author{Your Name}
\date{\today}

\begin{document}

\maketitle

\section{Introduction}
This is a sample document demonstrating various \LaTeX{} features including mathematical equations, tables, and formatting. You can cite papers using the citation search feature, and they will be automatically added to your \texttt{cite.bib} file.

\section{Mathematical Expressions}
Here are some example equations:

\subsection{Quadratic Formula}
The quadratic formula is given by:
\[x = \frac{-b \pm \sqrt{b^2 - 4ac}}{2a}\]

\subsection{Matrix Example}
A sample matrix:
\[
\begin{pmatrix}
a_{11} & a_{12} & a_{13} \\
a_{21} & a_{22} & a_{23} \\
a_{31} & a_{32} & a_{33}
\end{pmatrix}
\]

\section{Tables}
Below is a sample table showing data:

\begin{table}[h]
\centering
\begin{tabular}{@{}lrrr@{}}
\toprule
\textbf{Category} & \textbf{Value 1} & \textbf{Value 2} & \textbf{Total} \\
\midrule
Item A & 10.5 & 20.3 & 30.8 \\
Item B & 15.2 & 25.1 & 40.3 \\
Item C & 22.1 & 12.7 & 34.8 \\
\bottomrule
\end{tabular}
\caption{Sample Data Table}
\label{tab:sample}
\end{table}

\section{Lists}
\subsection{Bullet Points}
\begin{itemize}
  \item First item
  \item Second item
  \item Third item with sub-items:
    \begin{itemize}
      \item Sub-item 1
      \item Sub-item 2
    \end{itemize}
\end{itemize}

\subsection{Numbered List}
\begin{enumerate}
  \item First step
  \item Second step
  \item Third step
\end{enumerate}

\section{Citations and References}
When you use the citation search feature, references will be automatically managed. For example, you might cite a paper like this: \cite{example2024}. All citations will be automatically added to your \texttt{cite.bib} file.

\section{Conclusion}
This template demonstrates basic \LaTeX{} features including mathematical formulas, tables, and lists. The bibliography system is now set up to work seamlessly with Bibby AI's citation management. Feel free to modify it for your needs!

\bibliographystyle{plain}
\bibliography{cite}

\end{document}\documentclass{article}
\usepackage[utf8]{inputenc}
\usepackage{booktabs}
\usepackage{amsmath}
\usepackage{natbib}

\title{My Document}
\author{Your Name}
\date{\today}

\begin{document}

\maketitle

\section{Introduction}
This is a sample document demonstrating various \LaTeX{} features including mathematical equations, tables, and formatting. You can cite papers using the citation search feature, and they will be automatically added to your \texttt{cite.bib} file.

\section{Mathematical Expressions}
Here are some example equations:

\subsection{Quadratic Formula}
The quadratic formula is given by:
\[x = \frac{-b \pm \sqrt{b^2 - 4ac}}{2a}\]

\subsection{Matrix Example}
A sample matrix:
\[
\begin{pmatrix}
a_{11} & a_{12} & a_{13} \\
a_{21} & a_{22} & a_{23} \\
a_{31} & a_{32} & a_{33}
\end{pmatrix}
\]

\section{Tables}
Below is a sample table showing data:

\begin{table}[h]
\centering
\begin{tabular}{@{}lrrr@{}}
\toprule
\textbf{Category} & \textbf{Value 1} & \textbf{Value 2} & \textbf{Total} \\
\midrule
Item A & 10.5 & 20.3 & 30.8 \\
Item B & 15.2 & 25.1 & 40.3 \\
Item C & 22.1 & 12.7 & 34.8 \\
\bottomrule
\end{tabular}
\caption{Sample Data Table}
\label{tab:sample}
\end{table}

\section{Lists}
\subsection{Bullet Points}
\begin{itemize}
  \item First item
  \item Second item
  \item Third item with sub-items:
    \begin{itemize}
      \item Sub-item 1
      \item Sub-item 2
    \end{itemize}
\end{itemize}

\subsection{Numbered List}
\begin{enumerate}
  \item First step
  \item Second step
  \item Third step
\end{enumerate}

\section{Citations and References}
When you use the citation search feature, references will be automatically managed. For example, you might cite a paper like this: \cite{example2024}. All citations will be automatically added to your \texttt{cite.bib} file.

\section{Conclusion}
This template demonstrates basic \LaTeX{} features including mathematical formulas, tables, and lists. The bibliography system is now set up to work seamlessly with Bibby AI's citation management. Feel free to modify it for your needs!

\bibliographystyle{plain}
\bibliography{cite}

\end{document}\documentclass{article}
\usepackage[utf8]{inputenc}
\usepackage{booktabs}
\usepackage{amsmath}
\usepackage{natbib}

\title{My Document}
\author{Your Name}
\date{\today}

\begin{document}

\maketitle

\section{Introduction}
This is a sample document demonstrating various \LaTeX{} features including mathematical equations, tables, and formatting. You can cite papers using the citation search feature, and they will be automatically added to your \texttt{cite.bib} file.

\section{Mathematical Expressions}
Here are some example equations:

\subsection{Quadratic Formula}
The quadratic formula is given by:
\[x = \frac{-b \pm \sqrt{b^2 - 4ac}}{2a}\]

\subsection{Matrix Example}
A sample matrix:
\[
\begin{pmatrix}
a_{11} & a_{12} & a_{13} \\
a_{21} & a_{22} & a_{23} \\
a_{31} & a_{32} & a_{33}
\end{pmatrix}
\]

\section{Tables}
Below is a sample table showing data:

\begin{table}[h]
\centering
\begin{tabular}{@{}lrrr@{}}
\toprule
\textbf{Category} & \textbf{Value 1} & \textbf{Value 2} & \textbf{Total} \\
\midrule
Item A & 10.5 & 20.3 & 30.8 \\
Item B & 15.2 & 25.1 & 40.3 \\
Item C & 22.1 & 12.7 & 34.8 \\
\bottomrule
\end{tabular}
\caption{Sample Data Table}
\label{tab:sample}
\end{table}

\section{Lists}
\subsection{Bullet Points}
\begin{itemize}
  \item First item
  \item Second item
  \item Third item with sub-items:
    \begin{itemize}
      \item Sub-item 1
      \item Sub-item 2
    \end{itemize}
\end{itemize}

\subsection{Numbered List}
\begin{enumerate}
  \item First step
  \item Second step
  \item Third step
\end{enumerate}

\section{Citations and References}
When you use the citation search feature, references will be automatically managed. For example, you might cite a paper like this: \cite{example2024}. All citations will be automatically added to your \texttt{cite.bib} file.

\section{Conclusion}
This template demonstrates basic \LaTeX{} features including mathematical formulas, tables, and lists. The bibliography system is now set up to work seamlessly with Bibby AI's citation management. Feel free to modify it for your needs!

\bibliographystyle{plain}
\bibliography{cite}

\end{document}\documentclass{article}
\usepackage[utf8]{inputenc}
\usepackage{booktabs}
\usepackage{amsmath}
\usepackage{natbib}

\title{My Document}
\author{Your Name}
\date{\today}

\begin{document}

\maketitle

\section{Introduction}
This is a sample document demonstrating various \LaTeX{} features including mathematical equations, tables, and formatting. You can cite papers using the citation search feature, and they will be automatically added to your \texttt{cite.bib} file.

\section{Mathematical Expressions}
Here are some example equations:

\subsection{Quadratic Formula}
The quadratic formula is given by:
\[x = \frac{-b \pm \sqrt{b^2 - 4ac}}{2a}\]

\subsection{Matrix Example}
A sample matrix:
\[
\begin{pmatrix}
a_{11} & a_{12} & a_{13} \\
a_{21} & a_{22} & a_{23} \\
a_{31} & a_{32} & a_{33}
\end{pmatrix}
\]

\section{Tables}
Below is a sample table showing data:

\begin{table}[h]
\centering
\begin{tabular}{@{}lrrr@{}}
\toprule
\textbf{Category} & \textbf{Value 1} & \textbf{Value 2} & \textbf{Total} \\
\midrule
Item A & 10.5 & 20.3 & 30.8 \\
Item B & 15.2 & 25.1 & 40.3 \\
Item C & 22.1 & 12.7 & 34.8 \\
\bottomrule
\end{tabular}
\caption{Sample Data Table}
\label{tab:sample}
\end{table}

\section{Lists}
\subsection{Bullet Points}
\begin{itemize}
  \item First item
  \item Second item
  \item Third item with sub-items:
    \begin{itemize}
      \item Sub-item 1
      \item Sub-item 2
    \end{itemize}
\end{itemize}

\subsection{Numbered List}
\begin{enumerate}
  \item First step
  \item Second step
  \item Third step
\end{enumerate}

\section{Citations and References}
When you use the citation search feature, references will be automatically managed. For example, you might cite a paper like this: \cite{example2024}. All citations will be automatically added to your \texttt{cite.bib} file.

\section{Conclusion}
This template demonstrates basic \LaTeX{} features including mathematical formulas, tables, and lists. The bibliography system is now set up to work seamlessly with Bibby AI's citation management. Feel free to modify it for your needs!

\bibliographystyle{plain}
\bibliography{cite}

\end{document}\documentclass{article}
\usepackage[utf8]{inputenc}
\usepackage{booktabs}
\usepackage{amsmath}
\usepackage{natbib}

\title{My Document}
\author{Your Name}
\date{\today}

\begin{document}

\maketitle

\section{Introduction}
This is a sample document demonstrating various \LaTeX{} features including mathematical equations, tables, and formatting. You can cite papers using the citation search feature, and they will be automatically added to your \texttt{cite.bib} file.

\section{Mathematical Expressions}
Here are some example equations:

\subsection{Quadratic Formula}
The quadratic formula is given by:
\[x = \frac{-b \pm \sqrt{b^2 - 4ac}}{2a}\]

\subsection{Matrix Example}
A sample matrix:
\[
\begin{pmatrix}
a_{11} & a_{12} & a_{13} \\
a_{21} & a_{22} & a_{23} \\
a_{31} & a_{32} & a_{33}
\end{pmatrix}
\]

\section{Tables}
Below is a sample table showing data:

\begin{table}[h]
\centering
\begin{tabular}{@{}lrrr@{}}
\toprule
\textbf{Category} & \textbf{Value 1} & \textbf{Value 2} & \textbf{Total} \\
\midrule
Item A & 10.5 & 20.3 & 30.8 \\
Item B & 15.2 & 25.1 & 40.3 \\
Item C & 22.1 & 12.7 & 34.8 \\
\bottomrule
\end{tabular}
\caption{Sample Data Table}
\label{tab:sample}
\end{table}

\section{Lists}
\subsection{Bullet Points}
\begin{itemize}
  \item First item
  \item Second item
  \item Third item with sub-items:
    \begin{itemize}
      \item Sub-item 1
      \item Sub-item 2
    \end{itemize}
\end{itemize}

\subsection{Numbered List}
\begin{enumerate}
  \item First step
  \item Second step
  \item Third step
\end{enumerate}

\section{Citations and References}
When you use the citation search feature, references will be automatically managed. For example, you might cite a paper like this: \cite{example2024}. All citations will be automatically added to your \texttt{cite.bib} file.

\section{Conclusion}
This template demonstrates basic \LaTeX{} features including mathematical formulas, tables, and lists. The bibliography system is now set up to work seamlessly with Bibby AI's citation management. Feel free to modify it for your needs!

\bibliographystyle{plain}
\bibliography{cite}

\end{document}\documentclass{article}
\usepackage[utf8]{inputenc}
\usepackage{booktabs}
\usepackage{amsmath}
\usepackage{natbib}

\title{My Document}
\author{Your Name}
\date{\today}

\begin{document}

\maketitle

\section{Introduction}
This is a sample document demonstrating various \LaTeX{} features including mathematical equations, tables, and formatting. You can cite papers using the citation search feature, and they will be automatically added to your \texttt{cite.bib} file.

\section{Mathematical Expressions}
Here are some example equations:

\subsection{Quadratic Formula}
The quadratic formula is given by:
\[x = \frac{-b \pm \sqrt{b^2 - 4ac}}{2a}\]

\subsection{Matrix Example}
A sample matrix:
\[
\begin{pmatrix}
a_{11} & a_{12} & a_{13} \\
a_{21} & a_{22} & a_{23} \\
a_{31} & a_{32} & a_{33}
\end{pmatrix}
\]

\section{Tables}
Below is a sample table showing data:

\begin{table}[h]
\centering
\begin{tabular}{@{}lrrr@{}}
\toprule
\textbf{Category} & \textbf{Value 1} & \textbf{Value 2} & \textbf{Total} \\
\midrule
Item A & 10.5 & 20.3 & 30.8 \\
Item B & 15.2 & 25.1 & 40.3 \\
Item C & 22.1 & 12.7 & 34.8 \\
\bottomrule
\end{tabular}
\caption{Sample Data Table}
\label{tab:sample}
\end{table}

\section{Lists}
\subsection{Bullet Points}
\begin{itemize}
  \item First item
  \item Second item
  \item Third item with sub-items:
    \begin{itemize}
      \item Sub-item 1
      \item Sub-item 2
    \end{itemize}
\end{itemize}

\subsection{Numbered List}
\begin{enumerate}
  \item First step
  \item Second step
  \item Third step
\end{enumerate}

\section{Citations and References}
When you use the citation search feature, references will be automatically managed. For example, you might cite a paper like this: \cite{example2024}. All citations will be automatically added to your \texttt{cite.bib} file.

\section{Conclusion}
This template demonstrates basic \LaTeX{} features including mathematical formulas, tables, and lists. The bibliography system is now set up to work seamlessly with Bibby AI's citation management. Feel free to modify it for your needs!

\bibliographystyle{plain}
\bibliography{cite}

\end{document}\documentclass{article}
\usepackage[utf8]{inputenc}
\usepackage{booktabs}
\usepackage{amsmath}
\usepackage{natbib}

\title{My Document}
\author{Your Name}
\date{\today}

\begin{document}

\maketitle

\section{Introduction}
This is a sample document demonstrating various \LaTeX{} features including mathematical equations, tables, and formatting. You can cite papers using the citation search feature, and they will be automatically added to your \texttt{cite.bib} file.

\section{Mathematical Expressions}
Here are some example equations:

\subsection{Quadratic Formula}
The quadratic formula is given by:
\[x = \frac{-b \pm \sqrt{b^2 - 4ac}}{2a}\]

\subsection{Matrix Example}
A sample matrix:
\[
\begin{pmatrix}
a_{11} & a_{12} & a_{13} \\
a_{21} & a_{22} & a_{23} \\
a_{31} & a_{32} & a_{33}
\end{pmatrix}
\]

\section{Tables}
Below is a sample table showing data:

\begin{table}[h]
\centering
\begin{tabular}{@{}lrrr@{}}
\toprule
\textbf{Category} & \textbf{Value 1} & \textbf{Value 2} & \textbf{Total} \\
\midrule
Item A & 10.5 & 20.3 & 30.8 \\
Item B & 15.2 & 25.1 & 40.3 \\
Item C & 22.1 & 12.7 & 34.8 \\
\bottomrule
\end{tabular}
\caption{Sample Data Table}
\label{tab:sample}
\end{table}

\section{Lists}
\subsection{Bullet Points}
\begin{itemize}
  \item First item
  \item Second item
  \item Third item with sub-items:
    \begin{itemize}
      \item Sub-item 1
      \item Sub-item 2
    \end{itemize}
\end{itemize}

\subsection{Numbered List}
\begin{enumerate}
  \item First step
  \item Second step
  \item Third step
\end{enumerate}

\section{Citations and References}
When you use the citation search feature, references will be automatically managed. For example, you might cite a paper like this: \cite{example2024}. All citations will be automatically added to your \texttt{cite.bib} file.

\section{Conclusion}
This template demonstrates basic \LaTeX{} features including mathematical formulas, tables, and lists. The bibliography system is now set up to work seamlessly with Bibby AI's citation management. Feel free to modify it for your needs!

\bibliographystyle{plain}
\bibliography{cite}

\end{document}\documentclass{article}
\usepackage[utf8]{inputenc}
\usepackage{booktabs}
\usepackage{amsmath}
\usepackage{natbib}

\title{My Document}
\author{Your Name}
\date{\today}

\begin{document}

\maketitle

\section{Introduction}
This is a sample document demonstrating various \LaTeX{} features including mathematical equations, tables, and formatting. You can cite papers using the citation search feature, and they will be automatically added to your \texttt{cite.bib} file.

\section{Mathematical Expressions}
Here are some example equations:

\subsection{Quadratic Formula}
The quadratic formula is given by:
\[x = \frac{-b \pm \sqrt{b^2 - 4ac}}{2a}\]

\subsection{Matrix Example}
A sample matrix:
\[
\begin{pmatrix}
a_{11} & a_{12} & a_{13} \\
a_{21} & a_{22} & a_{23} \\
a_{31} & a_{32} & a_{33}
\end{pmatrix}
\]

\section{Tables}
Below is a sample table showing data:

\begin{table}[h]
\centering
\begin{tabular}{@{}lrrr@{}}
\toprule
\textbf{Category} & \textbf{Value 1} & \textbf{Value 2} & \textbf{Total} \\
\midrule
Item A & 10.5 & 20.3 & 30.8 \\
Item B & 15.2 & 25.1 & 40.3 \\
Item C & 22.1 & 12.7 & 34.8 \\
\bottomrule
\end{tabular}
\caption{Sample Data Table}
\label{tab:sample}
\end{table}

\section{Lists}
\subsection{Bullet Points}
\begin{itemize}
  \item First item
  \item Second item
  \item Third item with sub-items:
    \begin{itemize}
      \item Sub-item 1
      \item Sub-item 2
    \end{itemize}
\end{itemize}

\subsection{Numbered List}
\begin{enumerate}
  \item First step
  \item Second step
  \item Third step
\end{enumerate}

\section{Citations and References}
When you use the citation search feature, references will be automatically managed. For example, you might cite a paper like this: \cite{example2024}. All citations will be automatically added to your \texttt{cite.bib} file.

\section{Conclusion}
This template demonstrates basic \LaTeX{} features including mathematical formulas, tables, and lists. The bibliography system is now set up to work seamlessly with Bibby AI's citation management. Feel free to modify it for your needs!

\bibliographystyle{plain}
\bibliography{cite}

\end{document}\documentclass{article}
\usepackage[utf8]{inputenc}
\usepackage{booktabs}
\usepackage{amsmath}
\usepackage{natbib}

\title{My Document}
\author{Your Name}
\date{\today}

\begin{document}

\maketitle

\section{Introduction}
This is a sample document demonstrating various \LaTeX{} features including mathematical equations, tables, and formatting. You can cite papers using the citation search feature, and they will be automatically added to your \texttt{cite.bib} file.

\section{Mathematical Expressions}
Here are some example equations:

\subsection{Quadratic Formula}
The quadratic formula is given by:
\[x = \frac{-b \pm \sqrt{b^2 - 4ac}}{2a}\]

\subsection{Matrix Example}
A sample matrix:
\[
\begin{pmatrix}
a_{11} & a_{12} & a_{13} \\
a_{21} & a_{22} & a_{23} \\
a_{31} & a_{32} & a_{33}
\end{pmatrix}
\]

\section{Tables}
Below is a sample table showing data:

\begin{table}[h]
\centering
\begin{tabular}{@{}lrrr@{}}
\toprule
\textbf{Category} & \textbf{Value 1} & \textbf{Value 2} & \textbf{Total} \\
\midrule
Item A & 10.5 & 20.3 & 30.8 \\
Item B & 15.2 & 25.1 & 40.3 \\
Item C & 22.1 & 12.7 & 34.8 \\
\bottomrule
\end{tabular}
\caption{Sample Data Table}
\label{tab:sample}
\end{table}

\section{Lists}
\subsection{Bullet Points}
\begin{itemize}
  \item First item
  \item Second item
  \item Third item with sub-items:
    \begin{itemize}
      \item Sub-item 1
      \item Sub-item 2
    \end{itemize}
\end{itemize}

\subsection{Numbered List}
\begin{enumerate}
  \item First step
  \item Second step
  \item Third step
\end{enumerate}

\section{Citations and References}
When you use the citation search feature, references will be automatically managed. For example, you might cite a paper like this: \cite{example2024}. All citations will be automatically added to your \texttt{cite.bib} file.

\section{Conclusion}
This template demonstrates basic \LaTeX{} features including mathematical formulas, tables, and lists. The bibliography system is now set up to work seamlessly with Bibby AI's citation management. Feel free to modify it for your needs!

\bibliographystyle{plain}
\bibliography{cite}

\end{document}\documentclass{article}
\usepackage[utf8]{inputenc}
\usepackage{booktabs}
\usepackage{amsmath}
\usepackage{natbib}

\title{My Document}
\author{Your Name}
\date{\today}

\begin{document}

\maketitle

\section{Introduction}
This is a sample document demonstrating various \LaTeX{} features including mathematical equations, tables, and formatting. You can cite papers using the citation search feature, and they will be automatically added to your \texttt{cite.bib} file.

\section{Mathematical Expressions}
Here are some example equations:

\subsection{Quadratic Formula}
The quadratic formula is given by:
\[x = \frac{-b \pm \sqrt{b^2 - 4ac}}{2a}\]

\subsection{Matrix Example}
A sample matrix:
\[
\begin{pmatrix}
a_{11} & a_{12} & a_{13} \\
a_{21} & a_{22} & a_{23} \\
a_{31} & a_{32} & a_{33}
\end{pmatrix}
\]

\section{Tables}
Below is a sample table showing data:

\begin{table}[h]
\centering
\begin{tabular}{@{}lrrr@{}}
\toprule
\textbf{Category} & \textbf{Value 1} & \textbf{Value 2} & \textbf{Total} \\
\midrule
Item A & 10.5 & 20.3 & 30.8 \\
Item B & 15.2 & 25.1 & 40.3 \\
Item C & 22.1 & 12.7 & 34.8 \\
\bottomrule
\end{tabular}
\caption{Sample Data Table}
\label{tab:sample}
\end{table}

\section{Lists}
\subsection{Bullet Points}
\begin{itemize}
  \item First item
  \item Second item
  \item Third item with sub-items:
    \begin{itemize}
      \item Sub-item 1
      \item Sub-item 2
    \end{itemize}
\end{itemize}

\subsection{Numbered List}
\begin{enumerate}
  \item First step
  \item Second step
  \item Third step
\end{enumerate}

\section{Citations and References}
When you use the citation search feature, references will be automatically managed. For example, you might cite a paper like this: \cite{example2024}. All citations will be automatically added to your \texttt{cite.bib} file.

\section{Conclusion}
This template demonstrates basic \LaTeX{} features including mathematical formulas, tables, and lists. The bibliography system is now set up to work seamlessly with Bibby AI's citation management. Feel free to modify it for your needs!

\bibliographystyle{plain}
\bibliography{cite}

\end{document}\documentclass{article}
\usepackage[utf8]{inputenc}
\usepackage{booktabs}
\usepackage{amsmath}
\usepackage{natbib}

\title{My Document}
\author{Your Name}
\date{\today}

\begin{document}

\maketitle

\section{Introduction}
This is a sample document demonstrating various \LaTeX{} features including mathematical equations, tables, and formatting. You can cite papers using the citation search feature, and they will be automatically added to your \texttt{cite.bib} file.

\section{Mathematical Expressions}
Here are some example equations:

\subsection{Quadratic Formula}
The quadratic formula is given by:
\[x = \frac{-b \pm \sqrt{b^2 - 4ac}}{2a}\]

\subsection{Matrix Example}
A sample matrix:
\[
\begin{pmatrix}
a_{11} & a_{12} & a_{13} \\
a_{21} & a_{22} & a_{23} \\
a_{31} & a_{32} & a_{33}
\end{pmatrix}
\]

\section{Tables}
Below is a sample table showing data:

\begin{table}[h]
\centering
\begin{tabular}{@{}lrrr@{}}
\toprule
\textbf{Category} & \textbf{Value 1} & \textbf{Value 2} & \textbf{Total} \\
\midrule
Item A & 10.5 & 20.3 & 30.8 \\
Item B & 15.2 & 25.1 & 40.3 \\
Item C & 22.1 & 12.7 & 34.8 \\
\bottomrule
\end{tabular}
\caption{Sample Data Table}
\label{tab:sample}
\end{table}

\section{Lists}
\subsection{Bullet Points}
\begin{itemize}
  \item First item
  \item Second item
  \item Third item with sub-items:
    \begin{itemize}
      \item Sub-item 1
      \item Sub-item 2
    \end{itemize}
\end{itemize}

\subsection{Numbered List}
\begin{enumerate}
  \item First step
  \item Second step
  \item Third step
\end{enumerate}

\section{Citations and References}
When you use the citation search feature, references will be automatically managed. For example, you might cite a paper like this: \cite{example2024}. All citations will be automatically added to your \texttt{cite.bib} file.

\section{Conclusion}
This template demonstrates basic \LaTeX{} features including mathematical formulas, tables, and lists. The bibliography system is now set up to work seamlessly with Bibby AI's citation management. Feel free to modify it for your needs!

\bibliographystyle{plain}
\bibliography{cite}

\end{document}\documentclass{article}
\usepackage[utf8]{inputenc}
\usepackage{booktabs}
\usepackage{amsmath}
\usepackage{natbib}

\title{My Document}
\author{Your Name}
\date{\today}

\begin{document}

\maketitle

\section{Introduction}
This is a sample document demonstrating various \LaTeX{} features including mathematical equations, tables, and formatting. You can cite papers using the citation search feature, and they will be automatically added to your \texttt{cite.bib} file.

\section{Mathematical Expressions}
Here are some example equations:

\subsection{Quadratic Formula}
The quadratic formula is given by:
\[x = \frac{-b \pm \sqrt{b^2 - 4ac}}{2a}\]

\subsection{Matrix Example}
A sample matrix:
\[
\begin{pmatrix}
a_{11} & a_{12} & a_{13} \\
a_{21} & a_{22} & a_{23} \\
a_{31} & a_{32} & a_{33}
\end{pmatrix}
\]

\section{Tables}
Below is a sample table showing data:

\begin{table}[h]
\centering
\begin{tabular}{@{}lrrr@{}}
\toprule
\textbf{Category} & \textbf{Value 1} & \textbf{Value 2} & \textbf{Total} \\
\midrule
Item A & 10.5 & 20.3 & 30.8 \\
Item B & 15.2 & 25.1 & 40.3 \\
Item C & 22.1 & 12.7 & 34.8 \\
\bottomrule
\end{tabular}
\caption{Sample Data Table}
\label{tab:sample}
\end{table}

\section{Lists}
\subsection{Bullet Points}
\begin{itemize}
  \item First item
  \item Second item
  \item Third item with sub-items:
    \begin{itemize}
      \item Sub-item 1
      \item Sub-item 2
    \end{itemize}
\end{itemize}

\subsection{Numbered List}
\begin{enumerate}
  \item First step
  \item Second step
  \item Third step
\end{enumerate}

\section{Citations and References}
When you use the citation search feature, references will be automatically managed. For example, you might cite a paper like this: \cite{example2024}. All citations will be automatically added to your \texttt{cite.bib} file.

\section{Conclusion}
This template demonstrates basic \LaTeX{} features including mathematical formulas, tables, and lists. The bibliography system is now set up to work seamlessly with Bibby AI's citation management. Feel free to modify it for your needs!

\bibliographystyle{plain}
\bibliography{cite}

\end{document}\documentclass{article}
\usepackage[utf8]{inputenc}
\usepackage{booktabs}
\usepackage{amsmath}
\usepackage{natbib}

\title{My Document}
\author{Your Name}
\date{\today}

\begin{document}

\maketitle

\section{Introduction}
This is a sample document demonstrating various \LaTeX{} features including mathematical equations, tables, and formatting. You can cite papers using the citation search feature, and they will be automatically added to your \texttt{cite.bib} file.

\section{Mathematical Expressions}
Here are some example equations:

\subsection{Quadratic Formula}
The quadratic formula is given by:
\[x = \frac{-b \pm \sqrt{b^2 - 4ac}}{2a}\]

\subsection{Matrix Example}
A sample matrix:
\[
\begin{pmatrix}
a_{11} & a_{12} & a_{13} \\
a_{21} & a_{22} & a_{23} \\
a_{31} & a_{32} & a_{33}
\end{pmatrix}
\]

\section{Tables}
Below is a sample table showing data:

\begin{table}[h]
\centering
\begin{tabular}{@{}lrrr@{}}
\toprule
\textbf{Category} & \textbf{Value 1} & \textbf{Value 2} & \textbf{Total} \\
\midrule
Item A & 10.5 & 20.3 & 30.8 \\
Item B & 15.2 & 25.1 & 40.3 \\
Item C & 22.1 & 12.7 & 34.8 \\
\bottomrule
\end{tabular}
\caption{Sample Data Table}
\label{tab:sample}
\end{table}

\section{Lists}
\subsection{Bullet Points}
\begin{itemize}
  \item First item
  \item Second item
  \item Third item with sub-items:
    \begin{itemize}
      \item Sub-item 1
      \item Sub-item 2
    \end{itemize}
\end{itemize}

\subsection{Numbered List}
\begin{enumerate}
  \item First step
  \item Second step
  \item Third step
\end{enumerate}

\section{Citations and References}
When you use the citation search feature, references will be automatically managed. For example, you might cite a paper like this: \cite{example2024}. All citations will be automatically added to your \texttt{cite.bib} file.

\section{Conclusion}
This template demonstrates basic \LaTeX{} features including mathematical formulas, tables, and lists. The bibliography system is now set up to work seamlessly with Bibby AI's citation management. Feel free to modify it for your needs!

\bibliographystyle{plain}
\bibliography{cite}

\end{document}\documentclass{article}
\usepackage[utf8]{inputenc}
\usepackage{booktabs}
\usepackage{amsmath}
\usepackage{natbib}

\title{My Document}
\author{Your Name}
\date{\today}

\begin{document}

\maketitle

\section{Introduction}
This is a sample document demonstrating various \LaTeX{} features including mathematical equations, tables, and formatting. You can cite papers using the citation search feature, and they will be automatically added to your \texttt{cite.bib} file.

\section{Mathematical Expressions}
Here are some example equations:

\subsection{Quadratic Formula}
The quadratic formula is given by:
\[x = \frac{-b \pm \sqrt{b^2 - 4ac}}{2a}\]

\subsection{Matrix Example}
A sample matrix:
\[
\begin{pmatrix}
a_{11} & a_{12} & a_{13} \\
a_{21} & a_{22} & a_{23} \\
a_{31} & a_{32} & a_{33}
\end{pmatrix}
\]

\section{Tables}
Below is a sample table showing data:

\begin{table}[h]
\centering
\begin{tabular}{@{}lrrr@{}}
\toprule
\textbf{Category} & \textbf{Value 1} & \textbf{Value 2} & \textbf{Total} \\
\midrule
Item A & 10.5 & 20.3 & 30.8 \\
Item B & 15.2 & 25.1 & 40.3 \\
Item C & 22.1 & 12.7 & 34.8 \\
\bottomrule
\end{tabular}
\caption{Sample Data Table}
\label{tab:sample}
\end{table}

\section{Lists}
\subsection{Bullet Points}
\begin{itemize}
  \item First item
  \item Second item
  \item Third item with sub-items:
    \begin{itemize}
      \item Sub-item 1
      \item Sub-item 2
    \end{itemize}
\end{itemize}

\subsection{Numbered List}
\begin{enumerate}
  \item First step
  \item Second step
  \item Third step
\end{enumerate}

\section{Citations and References}
When you use the citation search feature, references will be automatically managed. For example, you might cite a paper like this: \cite{example2024}. All citations will be automatically added to your \texttt{cite.bib} file.

\section{Conclusion}
This template demonstrates basic \LaTeX{} features including mathematical formulas, tables, and lists. The bibliography system is now set up to work seamlessly with Bibby AI's citation management. Feel free to modify it for your needs!

\bibliographystyle{plain}
\bibliography{cite}

\end{document}\documentclass{article}
\usepackage[utf8]{inputenc}
\usepackage{booktabs}
\usepackage{amsmath}
\usepackage{natbib}

\title{My Document}
\author{Your Name}
\date{\today}

\begin{document}

\maketitle

\section{Introduction}
This is a sample document demonstrating various \LaTeX{} features including mathematical equations, tables, and formatting. You can cite papers using the citation search feature, and they will be automatically added to your \texttt{cite.bib} file.

\section{Mathematical Expressions}
Here are some example equations:

\subsection{Quadratic Formula}
The quadratic formula is given by:
\[x = \frac{-b \pm \sqrt{b^2 - 4ac}}{2a}\]

\subsection{Matrix Example}
A sample matrix:
\[
\begin{pmatrix}
a_{11} & a_{12} & a_{13} \\
a_{21} & a_{22} & a_{23} \\
a_{31} & a_{32} & a_{33}
\end{pmatrix}
\]

\section{Tables}
Below is a sample table showing data:

\begin{table}[h]
\centering
\begin{tabular}{@{}lrrr@{}}
\toprule
\textbf{Category} & \textbf{Value 1} & \textbf{Value 2} & \textbf{Total} \\
\midrule
Item A & 10.5 & 20.3 & 30.8 \\
Item B & 15.2 & 25.1 & 40.3 \\
Item C & 22.1 & 12.7 & 34.8 \\
\bottomrule
\end{tabular}
\caption{Sample Data Table}
\label{tab:sample}
\end{table}

\section{Lists}
\subsection{Bullet Points}
\begin{itemize}
  \item First item
  \item Second item
  \item Third item with sub-items:
    \begin{itemize}
      \item Sub-item 1
      \item Sub-item 2
    \end{itemize}
\end{itemize}

\subsection{Numbered List}
\begin{enumerate}
  \item First step
  \item Second step
  \item Third step
\end{enumerate}

\section{Citations and References}
When you use the citation search feature, references will be automatically managed. For example, you might cite a paper like this: \cite{example2024}. All citations will be automatically added to your \texttt{cite.bib} file.

\section{Conclusion}
This template demonstrates basic \LaTeX{} features including mathematical formulas, tables, and lists. The bibliography system is now set up to work seamlessly with Bibby AI's citation management. Feel free to modify it for your needs!

\bibliographystyle{plain}
\bibliography{cite}

\end{document}\documentclass{article}
\usepackage[utf8]{inputenc}
\usepackage{booktabs}
\usepackage{amsmath}
\usepackage{natbib}

\title{My Document}
\author{Your Name}
\date{\today}

\begin{document}

\maketitle

\section{Introduction}
This is a sample document demonstrating various \LaTeX{} features including mathematical equations, tables, and formatting. You can cite papers using the citation search feature, and they will be automatically added to your \texttt{cite.bib} file.

\section{Mathematical Expressions}
Here are some example equations:

\subsection{Quadratic Formula}
The quadratic formula is given by:
\[x = \frac{-b \pm \sqrt{b^2 - 4ac}}{2a}\]

\subsection{Matrix Example}
A sample matrix:
\[
\begin{pmatrix}
a_{11} & a_{12} & a_{13} \\
a_{21} & a_{22} & a_{23} \\
a_{31} & a_{32} & a_{33}
\end{pmatrix}
\]

\section{Tables}
Below is a sample table showing data:

\begin{table}[h]
\centering
\begin{tabular}{@{}lrrr@{}}
\toprule
\textbf{Category} & \textbf{Value 1} & \textbf{Value 2} & \textbf{Total} \\
\midrule
Item A & 10.5 & 20.3 & 30.8 \\
Item B & 15.2 & 25.1 & 40.3 \\
Item C & 22.1 & 12.7 & 34.8 \\
\bottomrule
\end{tabular}
\caption{Sample Data Table}
\label{tab:sample}
\end{table}

\section{Lists}
\subsection{Bullet Points}
\begin{itemize}
  \item First item
  \item Second item
  \item Third item with sub-items:
    \begin{itemize}
      \item Sub-item 1
      \item Sub-item 2
    \end{itemize}
\end{itemize}

\subsection{Numbered List}
\begin{enumerate}
  \item First step
  \item Second step
  \item Third step
\end{enumerate}

\section{Citations and References}
When you use the citation search feature, references will be automatically managed. For example, you might cite a paper like this: \cite{example2024}. All citations will be automatically added to your \texttt{cite.bib} file.

\section{Conclusion}
This template demonstrates basic \LaTeX{} features including mathematical formulas, tables, and lists. The bibliography system is now set up to work seamlessly with Bibby AI's citation management. Feel free to modify it for your needs!

\bibliographystyle{plain}
\bibliography{cite}

\end{document}\documentclass{article}
\usepackage[utf8]{inputenc}
\usepackage{booktabs}
\usepackage{amsmath}
\usepackage{natbib}

\title{My Document}
\author{Your Name}
\date{\today}

\begin{document}

\maketitle

\section{Introduction}
This is a sample document demonstrating various \LaTeX{} features including mathematical equations, tables, and formatting. You can cite papers using the citation search feature, and they will be automatically added to your \texttt{cite.bib} file.

\section{Mathematical Expressions}
Here are some example equations:

\subsection{Quadratic Formula}
The quadratic formula is given by:
\[x = \frac{-b \pm \sqrt{b^2 - 4ac}}{2a}\]

\subsection{Matrix Example}
A sample matrix:
\[
\begin{pmatrix}
a_{11} & a_{12} & a_{13} \\
a_{21} & a_{22} & a_{23} \\
a_{31} & a_{32} & a_{33}
\end{pmatrix}
\]

\section{Tables}
Below is a sample table showing data:

\begin{table}[h]
\centering
\begin{tabular}{@{}lrrr@{}}
\toprule
\textbf{Category} & \textbf{Value 1} & \textbf{Value 2} & \textbf{Total} \\
\midrule
Item A & 10.5 & 20.3 & 30.8 \\
Item B & 15.2 & 25.1 & 40.3 \\
Item C & 22.1 & 12.7 & 34.8 \\
\bottomrule
\end{tabular}
\caption{Sample Data Table}
\label{tab:sample}
\end{table}

\section{Lists}
\subsection{Bullet Points}
\begin{itemize}
  \item First item
  \item Second item
  \item Third item with sub-items:
    \begin{itemize}
      \item Sub-item 1
      \item Sub-item 2
    \end{itemize}
\end{itemize}

\subsection{Numbered List}
\begin{enumerate}
  \item First step
  \item Second step
  \item Third step
\end{enumerate}

\section{Citations and References}
When you use the citation search feature, references will be automatically managed. For example, you might cite a paper like this: \cite{example2024}. All citations will be automatically added to your \texttt{cite.bib} file.

\section{Conclusion}
This template demonstrates basic \LaTeX{} features including mathematical formulas, tables, and lists. The bibliography system is now set up to work seamlessly with Bibby AI's citation management. Feel free to modify it for your needs!

\bibliographystyle{plain}
\bibliography{cite}

\end{document}\documentclass{article}
\usepackage[utf8]{inputenc}
\usepackage{booktabs}
\usepackage{amsmath}
\usepackage{natbib}

\title{My Document}
\author{Your Name}
\date{\today}

\begin{document}

\maketitle

\section{Introduction}
This is a sample document demonstrating various \LaTeX{} features including mathematical equations, tables, and formatting. You can cite papers using the citation search feature, and they will be automatically added to your \texttt{cite.bib} file.

\section{Mathematical Expressions}
Here are some example equations:

\subsection{Quadratic Formula}
The quadratic formula is given by:
\[x = \frac{-b \pm \sqrt{b^2 - 4ac}}{2a}\]

\subsection{Matrix Example}
A sample matrix:
\[
\begin{pmatrix}
a_{11} & a_{12} & a_{13} \\
a_{21} & a_{22} & a_{23} \\
a_{31} & a_{32} & a_{33}
\end{pmatrix}
\]

\section{Tables}
Below is a sample table showing data:

\begin{table}[h]
\centering
\begin{tabular}{@{}lrrr@{}}
\toprule
\textbf{Category} & \textbf{Value 1} & \textbf{Value 2} & \textbf{Total} \\
\midrule
Item A & 10.5 & 20.3 & 30.8 \\
Item B & 15.2 & 25.1 & 40.3 \\
Item C & 22.1 & 12.7 & 34.8 \\
\bottomrule
\end{tabular}
\caption{Sample Data Table}
\label{tab:sample}
\end{table}

\section{Lists}
\subsection{Bullet Points}
\begin{itemize}
  \item First item
  \item Second item
  \item Third item with sub-items:
    \begin{itemize}
      \item Sub-item 1
      \item Sub-item 2
    \end{itemize}
\end{itemize}

\subsection{Numbered List}
\begin{enumerate}
  \item First step
  \item Second step
  \item Third step
\end{enumerate}

\section{Citations and References}
When you use the citation search feature, references will be automatically managed. For example, you might cite a paper like this: \cite{example2024}. All citations will be automatically added to your \texttt{cite.bib} file.

\section{Conclusion}
This template demonstrates basic \LaTeX{} features including mathematical formulas, tables, and lists. The bibliography system is now set up to work seamlessly with Bibby AI's citation management. Feel free to modify it for your needs!

\bibliographystyle{plain}
\bibliography{cite}

\end{document}\documentclass{article}
\usepackage[utf8]{inputenc}
\usepackage{booktabs}
\usepackage{amsmath}
\usepackage{natbib}

\title{My Document}
\author{Your Name}
\date{\today}

\begin{document}

\maketitle

\section{Introduction}
This is a sample document demonstrating various \LaTeX{} features including mathematical equations, tables, and formatting. You can cite papers using the citation search feature, and they will be automatically added to your \texttt{cite.bib} file.

\section{Mathematical Expressions}
Here are some example equations:

\subsection{Quadratic Formula}
The quadratic formula is given by:
\[x = \frac{-b \pm \sqrt{b^2 - 4ac}}{2a}\]

\subsection{Matrix Example}
A sample matrix:
\[
\begin{pmatrix}
a_{11} & a_{12} & a_{13} \\
a_{21} & a_{22} & a_{23} \\
a_{31} & a_{32} & a_{33}
\end{pmatrix}
\]

\section{Tables}
Below is a sample table showing data:

\begin{table}[h]
\centering
\begin{tabular}{@{}lrrr@{}}
\toprule
\textbf{Category} & \textbf{Value 1} & \textbf{Value 2} & \textbf{Total} \\
\midrule
Item A & 10.5 & 20.3 & 30.8 \\
Item B & 15.2 & 25.1 & 40.3 \\
Item C & 22.1 & 12.7 & 34.8 \\
\bottomrule
\end{tabular}
\caption{Sample Data Table}
\label{tab:sample}
\end{table}

\section{Lists}
\subsection{Bullet Points}
\begin{itemize}
  \item First item
  \item Second item
  \item Third item with sub-items:
    \begin{itemize}
      \item Sub-item 1
      \item Sub-item 2
    \end{itemize}
\end{itemize}

\subsection{Numbered List}
\begin{enumerate}
  \item First step
  \item Second step
  \item Third step
\end{enumerate}

\section{Citations and References}
When you use the citation search feature, references will be automatically managed. For example, you might cite a paper like this: \cite{example2024}. All citations will be automatically added to your \texttt{cite.bib} file.

\section{Conclusion}
This template demonstrates basic \LaTeX{} features including mathematical formulas, tables, and lists. The bibliography system is now set up to work seamlessly with Bibby AI's citation management. Feel free to modify it for your needs!

\bibliographystyle{plain}
\bibliography{cite}

\end{document}\documentclass{article}
\usepackage[utf8]{inputenc}
\usepackage{booktabs}
\usepackage{amsmath}
\usepackage{natbib}

\title{My Document}
\author{Your Name}
\date{\today}

\begin{document}

\maketitle

\section{Introduction}
This is a sample document demonstrating various \LaTeX{} features including mathematical equations, tables, and formatting. You can cite papers using the citation search feature, and they will be automatically added to your \texttt{cite.bib} file.

\section{Mathematical Expressions}
Here are some example equations:

\subsection{Quadratic Formula}
The quadratic formula is given by:
\[x = \frac{-b \pm \sqrt{b^2 - 4ac}}{2a}\]

\subsection{Matrix Example}
A sample matrix:
\[
\begin{pmatrix}
a_{11} & a_{12} & a_{13} \\
a_{21} & a_{22} & a_{23} \\
a_{31} & a_{32} & a_{33}
\end{pmatrix}
\]

\section{Tables}
Below is a sample table showing data:

\begin{table}[h]
\centering
\begin{tabular}{@{}lrrr@{}}
\toprule
\textbf{Category} & \textbf{Value 1} & \textbf{Value 2} & \textbf{Total} \\
\midrule
Item A & 10.5 & 20.3 & 30.8 \\
Item B & 15.2 & 25.1 & 40.3 \\
Item C & 22.1 & 12.7 & 34.8 \\
\bottomrule
\end{tabular}
\caption{Sample Data Table}
\label{tab:sample}
\end{table}

\section{Lists}
\subsection{Bullet Points}
\begin{itemize}
  \item First item
  \item Second item
  \item Third item with sub-items:
    \begin{itemize}
      \item Sub-item 1
      \item Sub-item 2
    \end{itemize}
\end{itemize}

\subsection{Numbered List}
\begin{enumerate}
  \item First step
  \item Second step
  \item Third step
\end{enumerate}

\section{Citations and References}
When you use the citation search feature, references will be automatically managed. For example, you might cite a paper like this: \cite{example2024}. All citations will be automatically added to your \texttt{cite.bib} file.

\section{Conclusion}
This template demonstrates basic \LaTeX{} features including mathematical formulas, tables, and lists. The bibliography system is now set up to work seamlessly with Bibby AI's citation management. Feel free to modify it for your needs!

\bibliographystyle{plain}
\bibliography{cite}

\end{document}\documentclass{article}
\usepackage[utf8]{inputenc}
\usepackage{booktabs}
\usepackage{amsmath}
\usepackage{natbib}

\title{My Document}
\author{Your Name}
\date{\today}

\begin{document}

\maketitle

\section{Introduction}
This is a sample document demonstrating various \LaTeX{} features including mathematical equations, tables, and formatting. You can cite papers using the citation search feature, and they will be automatically added to your \texttt{cite.bib} file.

\section{Mathematical Expressions}
Here are some example equations:

\subsection{Quadratic Formula}
The quadratic formula is given by:
\[x = \frac{-b \pm \sqrt{b^2 - 4ac}}{2a}\]

\subsection{Matrix Example}
A sample matrix:
\[
\begin{pmatrix}
a_{11} & a_{12} & a_{13} \\
a_{21} & a_{22} & a_{23} \\
a_{31} & a_{32} & a_{33}
\end{pmatrix}
\]

\section{Tables}
Below is a sample table showing data:

\begin{table}[h]
\centering
\begin{tabular}{@{}lrrr@{}}
\toprule
\textbf{Category} & \textbf{Value 1} & \textbf{Value 2} & \textbf{Total} \\
\midrule
Item A & 10.5 & 20.3 & 30.8 \\
Item B & 15.2 & 25.1 & 40.3 \\
Item C & 22.1 & 12.7 & 34.8 \\
\bottomrule
\end{tabular}
\caption{Sample Data Table}
\label{tab:sample}
\end{table}

\section{Lists}
\subsection{Bullet Points}
\begin{itemize}
  \item First item
  \item Second item
  \item Third item with sub-items:
    \begin{itemize}
      \item Sub-item 1
      \item Sub-item 2
    \end{itemize}
\end{itemize}

\subsection{Numbered List}
\begin{enumerate}
  \item First step
  \item Second step
  \item Third step
\end{enumerate}

\section{Citations and References}
When you use the citation search feature, references will be automatically managed. For example, you might cite a paper like this: \cite{example2024}. All citations will be automatically added to your \texttt{cite.bib} file.

\section{Conclusion}
This template demonstrates basic \LaTeX{} features including mathematical formulas, tables, and lists. The bibliography system is now set up to work seamlessly with Bibby AI's citation management. Feel free to modify it for your needs!

\bibliographystyle{plain}
\bibliography{cite}

\end{document}\documentclass{article}
\usepackage[utf8]{inputenc}
\usepackage{booktabs}
\usepackage{amsmath}
\usepackage{natbib}

\title{My Document}
\author{Your Name}
\date{\today}

\begin{document}

\maketitle

\section{Introduction}
This is a sample document demonstrating various \LaTeX{} features including mathematical equations, tables, and formatting. You can cite papers using the citation search feature, and they will be automatically added to your \texttt{cite.bib} file.

\section{Mathematical Expressions}
Here are some example equations:

\subsection{Quadratic Formula}
The quadratic formula is given by:
\[x = \frac{-b \pm \sqrt{b^2 - 4ac}}{2a}\]

\subsection{Matrix Example}
A sample matrix:
\[
\begin{pmatrix}
a_{11} & a_{12} & a_{13} \\
a_{21} & a_{22} & a_{23} \\
a_{31} & a_{32} & a_{33}
\end{pmatrix}
\]

\section{Tables}
Below is a sample table showing data:

\begin{table}[h]
\centering
\begin{tabular}{@{}lrrr@{}}
\toprule
\textbf{Category} & \textbf{Value 1} & \textbf{Value 2} & \textbf{Total} \\
\midrule
Item A & 10.5 & 20.3 & 30.8 \\
Item B & 15.2 & 25.1 & 40.3 \\
Item C & 22.1 & 12.7 & 34.8 \\
\bottomrule
\end{tabular}
\caption{Sample Data Table}
\label{tab:sample}
\end{table}

\section{Lists}
\subsection{Bullet Points}
\begin{itemize}
  \item First item
  \item Second item
  \item Third item with sub-items:
    \begin{itemize}
      \item Sub-item 1
      \item Sub-item 2
    \end{itemize}
\end{itemize}

\subsection{Numbered List}
\begin{enumerate}
  \item First step
  \item Second step
  \item Third step
\end{enumerate}

\section{Citations and References}
When you use the citation search feature, references will be automatically managed. For example, you might cite a paper like this: \cite{example2024}. All citations will be automatically added to your \texttt{cite.bib} file.

\section{Conclusion}
This template demonstrates basic \LaTeX{} features including mathematical formulas, tables, and lists. The bibliography system is now set up to work seamlessly with Bibby AI's citation management. Feel free to modify it for your needs!

\bibliographystyle{plain}
\bibliography{cite}

\end{document}\documentclass{article}
\usepackage[utf8]{inputenc}
\usepackage{booktabs}
\usepackage{amsmath}
\usepackage{natbib}

\title{My Document}
\author{Your Name}
\date{\today}

\begin{document}

\maketitle

\section{Introduction}
This is a sample document demonstrating various \LaTeX{} features including mathematical equations, tables, and formatting. You can cite papers using the citation search feature, and they will be automatically added to your \texttt{cite.bib} file.

\section{Mathematical Expressions}
Here are some example equations:

\subsection{Quadratic Formula}
The quadratic formula is given by:
\[x = \frac{-b \pm \sqrt{b^2 - 4ac}}{2a}\]

\subsection{Matrix Example}
A sample matrix:
\[
\begin{pmatrix}
a_{11} & a_{12} & a_{13} \\
a_{21} & a_{22} & a_{23} \\
a_{31} & a_{32} & a_{33}
\end{pmatrix}
\]

\section{Tables}
Below is a sample table showing data:

\begin{table}[h]
\centering
\begin{tabular}{@{}lrrr@{}}
\toprule
\textbf{Category} & \textbf{Value 1} & \textbf{Value 2} & \textbf{Total} \\
\midrule
Item A & 10.5 & 20.3 & 30.8 \\
Item B & 15.2 & 25.1 & 40.3 \\
Item C & 22.1 & 12.7 & 34.8 \\
\bottomrule
\end{tabular}
\caption{Sample Data Table}
\label{tab:sample}
\end{table}

\section{Lists}
\subsection{Bullet Points}
\begin{itemize}
  \item First item
  \item Second item
  \item Third item with sub-items:
    \begin{itemize}
      \item Sub-item 1
      \item Sub-item 2
    \end{itemize}
\end{itemize}

\subsection{Numbered List}
\begin{enumerate}
  \item First step
  \item Second step
  \item Third step
\end{enumerate}

\section{Citations and References}
When you use the citation search feature, references will be automatically managed. For example, you might cite a paper like this: \cite{example2024}. All citations will be automatically added to your \texttt{cite.bib} file.

\section{Conclusion}
This template demonstrates basic \LaTeX{} features including mathematical formulas, tables, and lists. The bibliography system is now set up to work seamlessly with Bibby AI's citation management. Feel free to modify it for your needs!

\bibliographystyle{plain}
\bibliography{cite}

\end{document}\documentclass{article}
\usepackage[utf8]{inputenc}
\usepackage{booktabs}
\usepackage{amsmath}
\usepackage{natbib}

\title{My Document}
\author{Your Name}
\date{\today}

\begin{document}

\maketitle

\section{Introduction}
This is a sample document demonstrating various \LaTeX{} features including mathematical equations, tables, and formatting. You can cite papers using the citation search feature, and they will be automatically added to your \texttt{cite.bib} file.

\section{Mathematical Expressions}
Here are some example equations:

\subsection{Quadratic Formula}
The quadratic formula is given by:
\[x = \frac{-b \pm \sqrt{b^2 - 4ac}}{2a}\]

\subsection{Matrix Example}
A sample matrix:
\[
\begin{pmatrix}
a_{11} & a_{12} & a_{13} \\
a_{21} & a_{22} & a_{23} \\
a_{31} & a_{32} & a_{33}
\end{pmatrix}
\]

\section{Tables}
Below is a sample table showing data:

\begin{table}[h]
\centering
\begin{tabular}{@{}lrrr@{}}
\toprule
\textbf{Category} & \textbf{Value 1} & \textbf{Value 2} & \textbf{Total} \\
\midrule
Item A & 10.5 & 20.3 & 30.8 \\
Item B & 15.2 & 25.1 & 40.3 \\
Item C & 22.1 & 12.7 & 34.8 \\
\bottomrule
\end{tabular}
\caption{Sample Data Table}
\label{tab:sample}
\end{table}

\section{Lists}
\subsection{Bullet Points}
\begin{itemize}
  \item First item
  \item Second item
  \item Third item with sub-items:
    \begin{itemize}
      \item Sub-item 1
      \item Sub-item 2
    \end{itemize}
\end{itemize}

\subsection{Numbered List}
\begin{enumerate}
  \item First step
  \item Second step
  \item Third step
\end{enumerate}

\section{Citations and References}
When you use the citation search feature, references will be automatically managed. For example, you might cite a paper like this: \cite{example2024}. All citations will be automatically added to your \texttt{cite.bib} file.

\section{Conclusion}
This template demonstrates basic \LaTeX{} features including mathematical formulas, tables, and lists. The bibliography system is now set up to work seamlessly with Bibby AI's citation management. Feel free to modify it for your needs!

\bibliographystyle{plain}
\bibliography{cite}

\end{document}\documentclass{article}
\usepackage[utf8]{inputenc}
\usepackage{booktabs}
\usepackage{amsmath}
\usepackage{natbib}

\title{My Document}
\author{Your Name}
\date{\today}

\begin{document}

\maketitle

\section{Introduction}
This is a sample document demonstrating various \LaTeX{} features including mathematical equations, tables, and formatting. You can cite papers using the citation search feature, and they will be automatically added to your \texttt{cite.bib} file.

\section{Mathematical Expressions}
Here are some example equations:

\subsection{Quadratic Formula}
The quadratic formula is given by:
\[x = \frac{-b \pm \sqrt{b^2 - 4ac}}{2a}\]

\subsection{Matrix Example}
A sample matrix:
\[
\begin{pmatrix}
a_{11} & a_{12} & a_{13} \\
a_{21} & a_{22} & a_{23} \\
a_{31} & a_{32} & a_{33}
\end{pmatrix}
\]

\section{Tables}
Below is a sample table showing data:

\begin{table}[h]
\centering
\begin{tabular}{@{}lrrr@{}}
\toprule
\textbf{Category} & \textbf{Value 1} & \textbf{Value 2} & \textbf{Total} \\
\midrule
Item A & 10.5 & 20.3 & 30.8 \\
Item B & 15.2 & 25.1 & 40.3 \\
Item C & 22.1 & 12.7 & 34.8 \\
\bottomrule
\end{tabular}
\caption{Sample Data Table}
\label{tab:sample}
\end{table}

\section{Lists}
\subsection{Bullet Points}
\begin{itemize}
  \item First item
  \item Second item
  \item Third item with sub-items:
    \begin{itemize}
      \item Sub-item 1
      \item Sub-item 2
    \end{itemize}
\end{itemize}

\subsection{Numbered List}
\begin{enumerate}
  \item First step
  \item Second step
  \item Third step
\end{enumerate}

\section{Citations and References}
When you use the citation search feature, references will be automatically managed. For example, you might cite a paper like this: \cite{example2024}. All citations will be automatically added to your \texttt{cite.bib} file.

\section{Conclusion}
This template demonstrates basic \LaTeX{} features including mathematical formulas, tables, and lists. The bibliography system is now set up to work seamlessly with Bibby AI's citation management. Feel free to modify it for your needs!

\bibliographystyle{plain}
\bibliography{cite}

\end{document}\documentclass{article}
\usepackage[utf8]{inputenc}
\usepackage{booktabs}
\usepackage{amsmath}
\usepackage{natbib}

\title{My Document}
\author{Your Name}
\date{\today}

\begin{document}

\maketitle

\section{Introduction}
This is a sample document demonstrating various \LaTeX{} features including mathematical equations, tables, and formatting. You can cite papers using the citation search feature, and they will be automatically added to your \texttt{cite.bib} file.

\section{Mathematical Expressions}
Here are some example equations:

\subsection{Quadratic Formula}
The quadratic formula is given by:
\[x = \frac{-b \pm \sqrt{b^2 - 4ac}}{2a}\]

\subsection{Matrix Example}
A sample matrix:
\[
\begin{pmatrix}
a_{11} & a_{12} & a_{13} \\
a_{21} & a_{22} & a_{23} \\
a_{31} & a_{32} & a_{33}
\end{pmatrix}
\]

\section{Tables}
Below is a sample table showing data:

\begin{table}[h]
\centering
\begin{tabular}{@{}lrrr@{}}
\toprule
\textbf{Category} & \textbf{Value 1} & \textbf{Value 2} & \textbf{Total} \\
\midrule
Item A & 10.5 & 20.3 & 30.8 \\
Item B & 15.2 & 25.1 & 40.3 \\
Item C & 22.1 & 12.7 & 34.8 \\
\bottomrule
\end{tabular}
\caption{Sample Data Table}
\label{tab:sample}
\end{table}

\section{Lists}
\subsection{Bullet Points}
\begin{itemize}
  \item First item
  \item Second item
  \item Third item with sub-items:
    \begin{itemize}
      \item Sub-item 1
      \item Sub-item 2
    \end{itemize}
\end{itemize}

\subsection{Numbered List}
\begin{enumerate}
  \item First step
  \item Second step
  \item Third step
\end{enumerate}

\section{Citations and References}
When you use the citation search feature, references will be automatically managed. For example, you might cite a paper like this: \cite{example2024}. All citations will be automatically added to your \texttt{cite.bib} file.

\section{Conclusion}
This template demonstrates basic \LaTeX{} features including mathematical formulas, tables, and lists. The bibliography system is now set up to work seamlessly with Bibby AI's citation management. Feel free to modify it for your needs!

\bibliographystyle{plain}
\bibliography{cite}

\end{document}\documentclass{article}
\usepackage[utf8]{inputenc}
\usepackage{booktabs}
\usepackage{amsmath}
\usepackage{natbib}

\title{My Document}
\author{Your Name}
\date{\today}

\begin{document}

\maketitle

\section{Introduction}
This is a sample document demonstrating various \LaTeX{} features including mathematical equations, tables, and formatting. You can cite papers using the citation search feature, and they will be automatically added to your \texttt{cite.bib} file.

\section{Mathematical Expressions}
Here are some example equations:

\subsection{Quadratic Formula}
The quadratic formula is given by:
\[x = \frac{-b \pm \sqrt{b^2 - 4ac}}{2a}\]

\subsection{Matrix Example}
A sample matrix:
\[
\begin{pmatrix}
a_{11} & a_{12} & a_{13} \\
a_{21} & a_{22} & a_{23} \\
a_{31} & a_{32} & a_{33}
\end{pmatrix}
\]

\section{Tables}
Below is a sample table showing data:

\begin{table}[h]
\centering
\begin{tabular}{@{}lrrr@{}}
\toprule
\textbf{Category} & \textbf{Value 1} & \textbf{Value 2} & \textbf{Total} \\
\midrule
Item A & 10.5 & 20.3 & 30.8 \\
Item B & 15.2 & 25.1 & 40.3 \\
Item C & 22.1 & 12.7 & 34.8 \\
\bottomrule
\end{tabular}
\caption{Sample Data Table}
\label{tab:sample}
\end{table}

\section{Lists}
\subsection{Bullet Points}
\begin{itemize}
  \item First item
  \item Second item
  \item Third item with sub-items:
    \begin{itemize}
      \item Sub-item 1
      \item Sub-item 2
    \end{itemize}
\end{itemize}

\subsection{Numbered List}
\begin{enumerate}
  \item First step
  \item Second step
  \item Third step
\end{enumerate}

\section{Citations and References}
When you use the citation search feature, references will be automatically managed. For example, you might cite a paper like this: \cite{example2024}. All citations will be automatically added to your \texttt{cite.bib} file.

\section{Conclusion}
This template demonstrates basic \LaTeX{} features including mathematical formulas, tables, and lists. The bibliography system is now set up to work seamlessly with Bibby AI's citation management. Feel free to modify it for your needs!

\bibliographystyle{plain}
\bibliography{cite}

\end{document}\documentclass{article}
\usepackage[utf8]{inputenc}
\usepackage{booktabs}
\usepackage{amsmath}
\usepackage{natbib}

\title{My Document}
\author{Your Name}
\date{\today}

\begin{document}

\maketitle

\section{Introduction}
This is a sample document demonstrating various \LaTeX{} features including mathematical equations, tables, and formatting. You can cite papers using the citation search feature, and they will be automatically added to your \texttt{cite.bib} file.

\section{Mathematical Expressions}
Here are some example equations:

\subsection{Quadratic Formula}
The quadratic formula is given by:
\[x = \frac{-b \pm \sqrt{b^2 - 4ac}}{2a}\]

\subsection{Matrix Example}
A sample matrix:
\[
\begin{pmatrix}
a_{11} & a_{12} & a_{13} \\
a_{21} & a_{22} & a_{23} \\
a_{31} & a_{32} & a_{33}
\end{pmatrix}
\]

\section{Tables}
Below is a sample table showing data:

\begin{table}[h]
\centering
\begin{tabular}{@{}lrrr@{}}
\toprule
\textbf{Category} & \textbf{Value 1} & \textbf{Value 2} & \textbf{Total} \\
\midrule
Item A & 10.5 & 20.3 & 30.8 \\
Item B & 15.2 & 25.1 & 40.3 \\
Item C & 22.1 & 12.7 & 34.8 \\
\bottomrule
\end{tabular}
\caption{Sample Data Table}
\label{tab:sample}
\end{table}

\section{Lists}
\subsection{Bullet Points}
\begin{itemize}
  \item First item
  \item Second item
  \item Third item with sub-items:
    \begin{itemize}
      \item Sub-item 1
      \item Sub-item 2
    \end{itemize}
\end{itemize}

\subsection{Numbered List}
\begin{enumerate}
  \item First step
  \item Second step
  \item Third step
\end{enumerate}

\section{Citations and References}
When you use the citation search feature, references will be automatically managed. For example, you might cite a paper like this: \cite{example2024}. All citations will be automatically added to your \texttt{cite.bib} file.

\section{Conclusion}
This template demonstrates basic \LaTeX{} features including mathematical formulas, tables, and lists. The bibliography system is now set up to work seamlessly with Bibby AI's citation management. Feel free to modify it for your needs!

\bibliographystyle{plain}
\bibliography{cite}

\end{document}\documentclass{article}
\usepackage[utf8]{inputenc}
\usepackage{booktabs}
\usepackage{amsmath}
\usepackage{natbib}

\title{My Document}
\author{Your Name}
\date{\today}

\begin{document}

\maketitle

\section{Introduction}
This is a sample document demonstrating various \LaTeX{} features including mathematical equations, tables, and formatting. You can cite papers using the citation search feature, and they will be automatically added to your \texttt{cite.bib} file.

\section{Mathematical Expressions}
Here are some example equations:

\subsection{Quadratic Formula}
The quadratic formula is given by:
\[x = \frac{-b \pm \sqrt{b^2 - 4ac}}{2a}\]

\subsection{Matrix Example}
A sample matrix:
\[
\begin{pmatrix}
a_{11} & a_{12} & a_{13} \\
a_{21} & a_{22} & a_{23} \\
a_{31} & a_{32} & a_{33}
\end{pmatrix}
\]

\section{Tables}
Below is a sample table showing data:

\begin{table}[h]
\centering
\begin{tabular}{@{}lrrr@{}}
\toprule
\textbf{Category} & \textbf{Value 1} & \textbf{Value 2} & \textbf{Total} \\
\midrule
Item A & 10.5 & 20.3 & 30.8 \\
Item B & 15.2 & 25.1 & 40.3 \\
Item C & 22.1 & 12.7 & 34.8 \\
\bottomrule
\end{tabular}
\caption{Sample Data Table}
\label{tab:sample}
\end{table}

\section{Lists}
\subsection{Bullet Points}
\begin{itemize}
  \item First item
  \item Second item
  \item Third item with sub-items:
    \begin{itemize}
      \item Sub-item 1
      \item Sub-item 2
    \end{itemize}
\end{itemize}

\subsection{Numbered List}
\begin{enumerate}
  \item First step
  \item Second step
  \item Third step
\end{enumerate}

\section{Citations and References}
When you use the citation search feature, references will be automatically managed. For example, you might cite a paper like this: \cite{example2024}. All citations will be automatically added to your \texttt{cite.bib} file.

\section{Conclusion}
This template demonstrates basic \LaTeX{} features including mathematical formulas, tables, and lists. The bibliography system is now set up to work seamlessly with Bibby AI's citation management. Feel free to modify it for your needs!

\bibliographystyle{plain}
\bibliography{cite}

\end{document}\documentclass{article}
\usepackage[utf8]{inputenc}
\usepackage{booktabs}
\usepackage{amsmath}
\usepackage{natbib}

\title{My Document}
\author{Your Name}
\date{\today}

\begin{document}

\maketitle

\section{Introduction}
This is a sample document demonstrating various \LaTeX{} features including mathematical equations, tables, and formatting. You can cite papers using the citation search feature, and they will be automatically added to your \texttt{cite.bib} file.

\section{Mathematical Expressions}
Here are some example equations:

\subsection{Quadratic Formula}
The quadratic formula is given by:
\[x = \frac{-b \pm \sqrt{b^2 - 4ac}}{2a}\]

\subsection{Matrix Example}
A sample matrix:
\[
\begin{pmatrix}
a_{11} & a_{12} & a_{13} \\
a_{21} & a_{22} & a_{23} \\
a_{31} & a_{32} & a_{33}
\end{pmatrix}
\]

\section{Tables}
Below is a sample table showing data:

\begin{table}[h]
\centering
\begin{tabular}{@{}lrrr@{}}
\toprule
\textbf{Category} & \textbf{Value 1} & \textbf{Value 2} & \textbf{Total} \\
\midrule
Item A & 10.5 & 20.3 & 30.8 \\
Item B & 15.2 & 25.1 & 40.3 \\
Item C & 22.1 & 12.7 & 34.8 \\
\bottomrule
\end{tabular}
\caption{Sample Data Table}
\label{tab:sample}
\end{table}

\section{Lists}
\subsection{Bullet Points}
\begin{itemize}
  \item First item
  \item Second item
  \item Third item with sub-items:
    \begin{itemize}
      \item Sub-item 1
      \item Sub-item 2
    \end{itemize}
\end{itemize}

\subsection{Numbered List}
\begin{enumerate}
  \item First step
  \item Second step
  \item Third step
\end{enumerate}

\section{Citations and References}
When you use the citation search feature, references will be automatically managed. For example, you might cite a paper like this: \cite{example2024}. All citations will be automatically added to your \texttt{cite.bib} file.

\section{Conclusion}
This template demonstrates basic \LaTeX{} features including mathematical formulas, tables, and lists. The bibliography system is now set up to work seamlessly with Bibby AI's citation management. Feel free to modify it for your needs!

\bibliographystyle{plain}
\bibliography{cite}

\end{document}\documentclass{article}
\usepackage[utf8]{inputenc}
\usepackage{booktabs}
\usepackage{amsmath}
\usepackage{natbib}

\title{My Document}
\author{Your Name}
\date{\today}

\begin{document}

\maketitle

\section{Introduction}
This is a sample document demonstrating various \LaTeX{} features including mathematical equations, tables, and formatting. You can cite papers using the citation search feature, and they will be automatically added to your \texttt{cite.bib} file.

\section{Mathematical Expressions}
Here are some example equations:

\subsection{Quadratic Formula}
The quadratic formula is given by:
\[x = \frac{-b \pm \sqrt{b^2 - 4ac}}{2a}\]

\subsection{Matrix Example}
A sample matrix:
\[
\begin{pmatrix}
a_{11} & a_{12} & a_{13} \\
a_{21} & a_{22} & a_{23} \\
a_{31} & a_{32} & a_{33}
\end{pmatrix}
\]

\section{Tables}
Below is a sample table showing data:

\begin{table}[h]
\centering
\begin{tabular}{@{}lrrr@{}}
\toprule
\textbf{Category} & \textbf{Value 1} & \textbf{Value 2} & \textbf{Total} \\
\midrule
Item A & 10.5 & 20.3 & 30.8 \\
Item B & 15.2 & 25.1 & 40.3 \\
Item C & 22.1 & 12.7 & 34.8 \\
\bottomrule
\end{tabular}
\caption{Sample Data Table}
\label{tab:sample}
\end{table}

\section{Lists}
\subsection{Bullet Points}
\begin{itemize}
  \item First item
  \item Second item
  \item Third item with sub-items:
    \begin{itemize}
      \item Sub-item 1
      \item Sub-item 2
    \end{itemize}
\end{itemize}

\subsection{Numbered List}
\begin{enumerate}
  \item First step
  \item Second step
  \item Third step
\end{enumerate}

\section{Citations and References}
When you use the citation search feature, references will be automatically managed. For example, you might cite a paper like this: \cite{example2024}. All citations will be automatically added to your \texttt{cite.bib} file.

\section{Conclusion}
This template demonstrates basic \LaTeX{} features including mathematical formulas, tables, and lists. The bibliography system is now set up to work seamlessly with Bibby AI's citation management. Feel free to modify it for your needs!

\bibliographystyle{plain}
\bibliography{cite}

\end{document}\documentclass{article}
\usepackage[utf8]{inputenc}
\usepackage{booktabs}
\usepackage{amsmath}
\usepackage{natbib}

\title{My Document}
\author{Your Name}
\date{\today}

\begin{document}

\maketitle

\section{Introduction}
This is a sample document demonstrating various \LaTeX{} features including mathematical equations, tables, and formatting. You can cite papers using the citation search feature, and they will be automatically added to your \texttt{cite.bib} file.

\section{Mathematical Expressions}
Here are some example equations:

\subsection{Quadratic Formula}
The quadratic formula is given by:
\[x = \frac{-b \pm \sqrt{b^2 - 4ac}}{2a}\]

\subsection{Matrix Example}
A sample matrix:
\[
\begin{pmatrix}
a_{11} & a_{12} & a_{13} \\
a_{21} & a_{22} & a_{23} \\
a_{31} & a_{32} & a_{33}
\end{pmatrix}
\]

\section{Tables}
Below is a sample table showing data:

\begin{table}[h]
\centering
\begin{tabular}{@{}lrrr@{}}
\toprule
\textbf{Category} & \textbf{Value 1} & \textbf{Value 2} & \textbf{Total} \\
\midrule
Item A & 10.5 & 20.3 & 30.8 \\
Item B & 15.2 & 25.1 & 40.3 \\
Item C & 22.1 & 12.7 & 34.8 \\
\bottomrule
\end{tabular}
\caption{Sample Data Table}
\label{tab:sample}
\end{table}

\section{Lists}
\subsection{Bullet Points}
\begin{itemize}
  \item First item
  \item Second item
  \item Third item with sub-items:
    \begin{itemize}
      \item Sub-item 1
      \item Sub-item 2
    \end{itemize}
\end{itemize}

\subsection{Numbered List}
\begin{enumerate}
  \item First step
  \item Second step
  \item Third step
\end{enumerate}

\section{Citations and References}
When you use the citation search feature, references will be automatically managed. For example, you might cite a paper like this: \cite{example2024}. All citations will be automatically added to your \texttt{cite.bib} file.

\section{Conclusion}
This template demonstrates basic \LaTeX{} features including mathematical formulas, tables, and lists. The bibliography system is now set up to work seamlessly with Bibby AI's citation management. Feel free to modify it for your needs!

\bibliographystyle{plain}
\bibliography{cite}

\end{document}\documentclass{article}
\usepackage[utf8]{inputenc}
\usepackage{booktabs}
\usepackage{amsmath}
\usepackage{natbib}

\title{My Document}
\author{Your Name}
\date{\today}

\begin{document}

\maketitle

\section{Introduction}
This is a sample document demonstrating various \LaTeX{} features including mathematical equations, tables, and formatting. You can cite papers using the citation search feature, and they will be automatically added to your \texttt{cite.bib} file.

\section{Mathematical Expressions}
Here are some example equations:

\subsection{Quadratic Formula}
The quadratic formula is given by:
\[x = \frac{-b \pm \sqrt{b^2 - 4ac}}{2a}\]

\subsection{Matrix Example}
A sample matrix:
\[
\begin{pmatrix}
a_{11} & a_{12} & a_{13} \\
a_{21} & a_{22} & a_{23} \\
a_{31} & a_{32} & a_{33}
\end{pmatrix}
\]

\section{Tables}
Below is a sample table showing data:

\begin{table}[h]
\centering
\begin{tabular}{@{}lrrr@{}}
\toprule
\textbf{Category} & \textbf{Value 1} & \textbf{Value 2} & \textbf{Total} \\
\midrule
Item A & 10.5 & 20.3 & 30.8 \\
Item B & 15.2 & 25.1 & 40.3 \\
Item C & 22.1 & 12.7 & 34.8 \\
\bottomrule
\end{tabular}
\caption{Sample Data Table}
\label{tab:sample}
\end{table}

\section{Lists}
\subsection{Bullet Points}
\begin{itemize}
  \item First item
  \item Second item
  \item Third item with sub-items:
    \begin{itemize}
      \item Sub-item 1
      \item Sub-item 2
    \end{itemize}
\end{itemize}

\subsection{Numbered List}
\begin{enumerate}
  \item First step
  \item Second step
  \item Third step
\end{enumerate}

\section{Citations and References}
When you use the citation search feature, references will be automatically managed. For example, you might cite a paper like this: \cite{example2024}. All citations will be automatically added to your \texttt{cite.bib} file.

\section{Conclusion}
This template demonstrates basic \LaTeX{} features including mathematical formulas, tables, and lists. The bibliography system is now set up to work seamlessly with Bibby AI's citation management. Feel free to modify it for your needs!

\bibliographystyle{plain}
\bibliography{cite}

\end{document}\documentclass{article}
\usepackage[utf8]{inputenc}
\usepackage{booktabs}
\usepackage{amsmath}
\usepackage{natbib}

\title{My Document}
\author{Your Name}
\date{\today}

\begin{document}

\maketitle

\section{Introduction}
This is a sample document demonstrating various \LaTeX{} features including mathematical equations, tables, and formatting. You can cite papers using the citation search feature, and they will be automatically added to your \texttt{cite.bib} file.

\section{Mathematical Expressions}
Here are some example equations:

\subsection{Quadratic Formula}
The quadratic formula is given by:
\[x = \frac{-b \pm \sqrt{b^2 - 4ac}}{2a}\]

\subsection{Matrix Example}
A sample matrix:
\[
\begin{pmatrix}
a_{11} & a_{12} & a_{13} \\
a_{21} & a_{22} & a_{23} \\
a_{31} & a_{32} & a_{33}
\end{pmatrix}
\]

\section{Tables}
Below is a sample table showing data:

\begin{table}[h]
\centering
\begin{tabular}{@{}lrrr@{}}
\toprule
\textbf{Category} & \textbf{Value 1} & \textbf{Value 2} & \textbf{Total} \\
\midrule
Item A & 10.5 & 20.3 & 30.8 \\
Item B & 15.2 & 25.1 & 40.3 \\
Item C & 22.1 & 12.7 & 34.8 \\
\bottomrule
\end{tabular}
\caption{Sample Data Table}
\label{tab:sample}
\end{table}

\section{Lists}
\subsection{Bullet Points}
\begin{itemize}
  \item First item
  \item Second item
  \item Third item with sub-items:
    \begin{itemize}
      \item Sub-item 1
      \item Sub-item 2
    \end{itemize}
\end{itemize}

\subsection{Numbered List}
\begin{enumerate}
  \item First step
  \item Second step
  \item Third step
\end{enumerate}

\section{Citations and References}
When you use the citation search feature, references will be automatically managed. For example, you might cite a paper like this: \cite{example2024}. All citations will be automatically added to your \texttt{cite.bib} file.

\section{Conclusion}
This template demonstrates basic \LaTeX{} features including mathematical formulas, tables, and lists. The bibliography system is now set up to work seamlessly with Bibby AI's citation management. Feel free to modify it for your needs!

\bibliographystyle{plain}
\bibliography{cite}

\end{document}\documentclass{article}
\usepackage[utf8]{inputenc}
\usepackage{booktabs}
\usepackage{amsmath}
\usepackage{natbib}

\title{My Document}
\author{Your Name}
\date{\today}

\begin{document}

\maketitle

\section{Introduction}
This is a sample document demonstrating various \LaTeX{} features including mathematical equations, tables, and formatting. You can cite papers using the citation search feature, and they will be automatically added to your \texttt{cite.bib} file.

\section{Mathematical Expressions}
Here are some example equations:

\subsection{Quadratic Formula}
The quadratic formula is given by:
\[x = \frac{-b \pm \sqrt{b^2 - 4ac}}{2a}\]

\subsection{Matrix Example}
A sample matrix:
\[
\begin{pmatrix}
a_{11} & a_{12} & a_{13} \\
a_{21} & a_{22} & a_{23} \\
a_{31} & a_{32} & a_{33}
\end{pmatrix}
\]

\section{Tables}
Below is a sample table showing data:

\begin{table}[h]
\centering
\begin{tabular}{@{}lrrr@{}}
\toprule
\textbf{Category} & \textbf{Value 1} & \textbf{Value 2} & \textbf{Total} \\
\midrule
Item A & 10.5 & 20.3 & 30.8 \\
Item B & 15.2 & 25.1 & 40.3 \\
Item C & 22.1 & 12.7 & 34.8 \\
\bottomrule
\end{tabular}
\caption{Sample Data Table}
\label{tab:sample}
\end{table}

\section{Lists}
\subsection{Bullet Points}
\begin{itemize}
  \item First item
  \item Second item
  \item Third item with sub-items:
    \begin{itemize}
      \item Sub-item 1
      \item Sub-item 2
    \end{itemize}
\end{itemize}

\subsection{Numbered List}
\begin{enumerate}
  \item First step
  \item Second step
  \item Third step
\end{enumerate}

\section{Citations and References}
When you use the citation search feature, references will be automatically managed. For example, you might cite a paper like this: \cite{example2024}. All citations will be automatically added to your \texttt{cite.bib} file.

\section{Conclusion}
This template demonstrates basic \LaTeX{} features including mathematical formulas, tables, and lists. The bibliography system is now set up to work seamlessly with Bibby AI's citation management. Feel free to modify it for your needs!

\bibliographystyle{plain}
\bibliography{cite}

\end{document}\documentclass{article}
\usepackage[utf8]{inputenc}
\usepackage{booktabs}
\usepackage{amsmath}
\usepackage{natbib}

\title{My Document}
\author{Your Name}
\date{\today}

\begin{document}

\maketitle

\section{Introduction}
This is a sample document demonstrating various \LaTeX{} features including mathematical equations, tables, and formatting. You can cite papers using the citation search feature, and they will be automatically added to your \texttt{cite.bib} file.

\section{Mathematical Expressions}
Here are some example equations:

\subsection{Quadratic Formula}
The quadratic formula is given by:
\[x = \frac{-b \pm \sqrt{b^2 - 4ac}}{2a}\]

\subsection{Matrix Example}
A sample matrix:
\[
\begin{pmatrix}
a_{11} & a_{12} & a_{13} \\
a_{21} & a_{22} & a_{23} \\
a_{31} & a_{32} & a_{33}
\end{pmatrix}
\]

\section{Tables}
Below is a sample table showing data:

\begin{table}[h]
\centering
\begin{tabular}{@{}lrrr@{}}
\toprule
\textbf{Category} & \textbf{Value 1} & \textbf{Value 2} & \textbf{Total} \\
\midrule
Item A & 10.5 & 20.3 & 30.8 \\
Item B & 15.2 & 25.1 & 40.3 \\
Item C & 22.1 & 12.7 & 34.8 \\
\bottomrule
\end{tabular}
\caption{Sample Data Table}
\label{tab:sample}
\end{table}

\section{Lists}
\subsection{Bullet Points}
\begin{itemize}
  \item First item
  \item Second item
  \item Third item with sub-items:
    \begin{itemize}
      \item Sub-item 1
      \item Sub-item 2
    \end{itemize}
\end{itemize}

\subsection{Numbered List}
\begin{enumerate}
  \item First step
  \item Second step
  \item Third step
\end{enumerate}

\section{Citations and References}
When you use the citation search feature, references will be automatically managed. For example, you might cite a paper like this: \cite{example2024}. All citations will be automatically added to your \texttt{cite.bib} file.

\section{Conclusion}
This template demonstrates basic \LaTeX{} features including mathematical formulas, tables, and lists. The bibliography system is now set up to work seamlessly with Bibby AI's citation management. Feel free to modify it for your needs!

\bibliographystyle{plain}
\bibliography{cite}

\end{document}\documentclass{article}
\usepackage[utf8]{inputenc}
\usepackage{booktabs}
\usepackage{amsmath}
\usepackage{natbib}

\title{My Document}
\author{Your Name}
\date{\today}

\begin{document}

\maketitle

\section{Introduction}
This is a sample document demonstrating various \LaTeX{} features including mathematical equations, tables, and formatting. You can cite papers using the citation search feature, and they will be automatically added to your \texttt{cite.bib} file.

\section{Mathematical Expressions}
Here are some example equations:

\subsection{Quadratic Formula}
The quadratic formula is given by:
\[x = \frac{-b \pm \sqrt{b^2 - 4ac}}{2a}\]

\subsection{Matrix Example}
A sample matrix:
\[
\begin{pmatrix}
a_{11} & a_{12} & a_{13} \\
a_{21} & a_{22} & a_{23} \\
a_{31} & a_{32} & a_{33}
\end{pmatrix}
\]

\section{Tables}
Below is a sample table showing data:

\begin{table}[h]
\centering
\begin{tabular}{@{}lrrr@{}}
\toprule
\textbf{Category} & \textbf{Value 1} & \textbf{Value 2} & \textbf{Total} \\
\midrule
Item A & 10.5 & 20.3 & 30.8 \\
Item B & 15.2 & 25.1 & 40.3 \\
Item C & 22.1 & 12.7 & 34.8 \\
\bottomrule
\end{tabular}
\caption{Sample Data Table}
\label{tab:sample}
\end{table}

\section{Lists}
\subsection{Bullet Points}
\begin{itemize}
  \item First item
  \item Second item
  \item Third item with sub-items:
    \begin{itemize}
      \item Sub-item 1
      \item Sub-item 2
    \end{itemize}
\end{itemize}

\subsection{Numbered List}
\begin{enumerate}
  \item First step
  \item Second step
  \item Third step
\end{enumerate}

\section{Citations and References}
When you use the citation search feature, references will be automatically managed. For example, you might cite a paper like this: \cite{example2024}. All citations will be automatically added to your \texttt{cite.bib} file.

\section{Conclusion}
This template demonstrates basic \LaTeX{} features including mathematical formulas, tables, and lists. The bibliography system is now set up to work seamlessly with Bibby AI's citation management. Feel free to modify it for your needs!

\bibliographystyle{plain}
\bibliography{cite}

\end{document}\documentclass{article}
\usepackage[utf8]{inputenc}
\usepackage{booktabs}
\usepackage{amsmath}
\usepackage{natbib}

\title{My Document}
\author{Your Name}
\date{\today}

\begin{document}

\maketitle

\section{Introduction}
This is a sample document demonstrating various \LaTeX{} features including mathematical equations, tables, and formatting. You can cite papers using the citation search feature, and they will be automatically added to your \texttt{cite.bib} file.

\section{Mathematical Expressions}
Here are some example equations:

\subsection{Quadratic Formula}
The quadratic formula is given by:
\[x = \frac{-b \pm \sqrt{b^2 - 4ac}}{2a}\]

\subsection{Matrix Example}
A sample matrix:
\[
\begin{pmatrix}
a_{11} & a_{12} & a_{13} \\
a_{21} & a_{22} & a_{23} \\
a_{31} & a_{32} & a_{33}
\end{pmatrix}
\]

\section{Tables}
Below is a sample table showing data:

\begin{table}[h]
\centering
\begin{tabular}{@{}lrrr@{}}
\toprule
\textbf{Category} & \textbf{Value 1} & \textbf{Value 2} & \textbf{Total} \\
\midrule
Item A & 10.5 & 20.3 & 30.8 \\
Item B & 15.2 & 25.1 & 40.3 \\
Item C & 22.1 & 12.7 & 34.8 \\
\bottomrule
\end{tabular}
\caption{Sample Data Table}
\label{tab:sample}
\end{table}

\section{Lists}
\subsection{Bullet Points}
\begin{itemize}
  \item First item
  \item Second item
  \item Third item with sub-items:
    \begin{itemize}
      \item Sub-item 1
      \item Sub-item 2
    \end{itemize}
\end{itemize}

\subsection{Numbered List}
\begin{enumerate}
  \item First step
  \item Second step
  \item Third step
\end{enumerate}

\section{Citations and References}
When you use the citation search feature, references will be automatically managed. For example, you might cite a paper like this: \cite{example2024}. All citations will be automatically added to your \texttt{cite.bib} file.

\section{Conclusion}
This template demonstrates basic \LaTeX{} features including mathematical formulas, tables, and lists. The bibliography system is now set up to work seamlessly with Bibby AI's citation management. Feel free to modify it for your needs!

\bibliographystyle{plain}
\bibliography{cite}

\end{document}\documentclass{article}
\usepackage[utf8]{inputenc}
\usepackage{booktabs}
\usepackage{amsmath}
\usepackage{natbib}

\title{My Document}
\author{Your Name}
\date{\today}

\begin{document}

\maketitle

\section{Introduction}
This is a sample document demonstrating various \LaTeX{} features including mathematical equations, tables, and formatting. You can cite papers using the citation search feature, and they will be automatically added to your \texttt{cite.bib} file.

\section{Mathematical Expressions}
Here are some example equations:

\subsection{Quadratic Formula}
The quadratic formula is given by:
\[x = \frac{-b \pm \sqrt{b^2 - 4ac}}{2a}\]

\subsection{Matrix Example}
A sample matrix:
\[
\begin{pmatrix}
a_{11} & a_{12} & a_{13} \\
a_{21} & a_{22} & a_{23} \\
a_{31} & a_{32} & a_{33}
\end{pmatrix}
\]

\section{Tables}
Below is a sample table showing data:

\begin{table}[h]
\centering
\begin{tabular}{@{}lrrr@{}}
\toprule
\textbf{Category} & \textbf{Value 1} & \textbf{Value 2} & \textbf{Total} \\
\midrule
Item A & 10.5 & 20.3 & 30.8 \\
Item B & 15.2 & 25.1 & 40.3 \\
Item C & 22.1 & 12.7 & 34.8 \\
\bottomrule
\end{tabular}
\caption{Sample Data Table}
\label{tab:sample}
\end{table}

\section{Lists}
\subsection{Bullet Points}
\begin{itemize}
  \item First item
  \item Second item
  \item Third item with sub-items:
    \begin{itemize}
      \item Sub-item 1
      \item Sub-item 2
    \end{itemize}
\end{itemize}

\subsection{Numbered List}
\begin{enumerate}
  \item First step
  \item Second step
  \item Third step
\end{enumerate}

\section{Citations and References}
When you use the citation search feature, references will be automatically managed. For example, you might cite a paper like this: \cite{example2024}. All citations will be automatically added to your \texttt{cite.bib} file.

\section{Conclusion}
This template demonstrates basic \LaTeX{} features including mathematical formulas, tables, and lists. The bibliography system is now set up to work seamlessly with Bibby AI's citation management. Feel free to modify it for your needs!

\bibliographystyle{plain}
\bibliography{cite}

\end{document}\documentclass{article}
\usepackage[utf8]{inputenc}
\usepackage{booktabs}
\usepackage{amsmath}
\usepackage{natbib}

\title{My Document}
\author{Your Name}
\date{\today}

\begin{document}

\maketitle

\section{Introduction}
This is a sample document demonstrating various \LaTeX{} features including mathematical equations, tables, and formatting. You can cite papers using the citation search feature, and they will be automatically added to your \texttt{cite.bib} file.

\section{Mathematical Expressions}
Here are some example equations:

\subsection{Quadratic Formula}
The quadratic formula is given by:
\[x = \frac{-b \pm \sqrt{b^2 - 4ac}}{2a}\]

\subsection{Matrix Example}
A sample matrix:
\[
\begin{pmatrix}
a_{11} & a_{12} & a_{13} \\
a_{21} & a_{22} & a_{23} \\
a_{31} & a_{32} & a_{33}
\end{pmatrix}
\]

\section{Tables}
Below is a sample table showing data:

\begin{table}[h]
\centering
\begin{tabular}{@{}lrrr@{}}
\toprule
\textbf{Category} & \textbf{Value 1} & \textbf{Value 2} & \textbf{Total} \\
\midrule
Item A & 10.5 & 20.3 & 30.8 \\
Item B & 15.2 & 25.1 & 40.3 \\
Item C & 22.1 & 12.7 & 34.8 \\
\bottomrule
\end{tabular}
\caption{Sample Data Table}
\label{tab:sample}
\end{table}

\section{Lists}
\subsection{Bullet Points}
\begin{itemize}
  \item First item
  \item Second item
  \item Third item with sub-items:
    \begin{itemize}
      \item Sub-item 1
      \item Sub-item 2
    \end{itemize}
\end{itemize}

\subsection{Numbered List}
\begin{enumerate}
  \item First step
  \item Second step
  \item Third step
\end{enumerate}

\section{Citations and References}
When you use the citation search feature, references will be automatically managed. For example, you might cite a paper like this: \cite{example2024}. All citations will be automatically added to your \texttt{cite.bib} file.

\section{Conclusion}
This template demonstrates basic \LaTeX{} features including mathematical formulas, tables, and lists. The bibliography system is now set up to work seamlessly with Bibby AI's citation management. Feel free to modify it for your needs!

\bibliographystyle{plain}
\bibliography{cite}

\end{document}\documentclass{article}
\usepackage[utf8]{inputenc}
\usepackage{booktabs}
\usepackage{amsmath}
\usepackage{natbib}

\title{My Document}
\author{Your Name}
\date{\today}

\begin{document}

\maketitle

\section{Introduction}
This is a sample document demonstrating various \LaTeX{} features including mathematical equations, tables, and formatting. You can cite papers using the citation search feature, and they will be automatically added to your \texttt{cite.bib} file.

\section{Mathematical Expressions}
Here are some example equations:

\subsection{Quadratic Formula}
The quadratic formula is given by:
\[x = \frac{-b \pm \sqrt{b^2 - 4ac}}{2a}\]

\subsection{Matrix Example}
A sample matrix:
\[
\begin{pmatrix}
a_{11} & a_{12} & a_{13} \\
a_{21} & a_{22} & a_{23} \\
a_{31} & a_{32} & a_{33}
\end{pmatrix}
\]

\section{Tables}
Below is a sample table showing data:

\begin{table}[h]
\centering
\begin{tabular}{@{}lrrr@{}}
\toprule
\textbf{Category} & \textbf{Value 1} & \textbf{Value 2} & \textbf{Total} \\
\midrule
Item A & 10.5 & 20.3 & 30.8 \\
Item B & 15.2 & 25.1 & 40.3 \\
Item C & 22.1 & 12.7 & 34.8 \\
\bottomrule
\end{tabular}
\caption{Sample Data Table}
\label{tab:sample}
\end{table}

\section{Lists}
\subsection{Bullet Points}
\begin{itemize}
  \item First item
  \item Second item
  \item Third item with sub-items:
    \begin{itemize}
      \item Sub-item 1
      \item Sub-item 2
    \end{itemize}
\end{itemize}

\subsection{Numbered List}
\begin{enumerate}
  \item First step
  \item Second step
  \item Third step
\end{enumerate}

\section{Citations and References}
When you use the citation search feature, references will be automatically managed. For example, you might cite a paper like this: \cite{example2024}. All citations will be automatically added to your \texttt{cite.bib} file.

\section{Conclusion}
This template demonstrates basic \LaTeX{} features including mathematical formulas, tables, and lists. The bibliography system is now set up to work seamlessly with Bibby AI's citation management. Feel free to modify it for your needs!

\bibliographystyle{plain}
\bibliography{cite}

\end{document}\documentclass{article}
\usepackage[utf8]{inputenc}
\usepackage{booktabs}
\usepackage{amsmath}
\usepackage{natbib}

\title{My Document}
\author{Your Name}
\date{\today}

\begin{document}

\maketitle

\section{Introduction}
This is a sample document demonstrating various \LaTeX{} features including mathematical equations, tables, and formatting. You can cite papers using the citation search feature, and they will be automatically added to your \texttt{cite.bib} file.

\section{Mathematical Expressions}
Here are some example equations:

\subsection{Quadratic Formula}
The quadratic formula is given by:
\[x = \frac{-b \pm \sqrt{b^2 - 4ac}}{2a}\]

\subsection{Matrix Example}
A sample matrix:
\[
\begin{pmatrix}
a_{11} & a_{12} & a_{13} \\
a_{21} & a_{22} & a_{23} \\
a_{31} & a_{32} & a_{33}
\end{pmatrix}
\]

\section{Tables}
Below is a sample table showing data:

\begin{table}[h]
\centering
\begin{tabular}{@{}lrrr@{}}
\toprule
\textbf{Category} & \textbf{Value 1} & \textbf{Value 2} & \textbf{Total} \\
\midrule
Item A & 10.5 & 20.3 & 30.8 \\
Item B & 15.2 & 25.1 & 40.3 \\
Item C & 22.1 & 12.7 & 34.8 \\
\bottomrule
\end{tabular}
\caption{Sample Data Table}
\label{tab:sample}
\end{table}

\section{Lists}
\subsection{Bullet Points}
\begin{itemize}
  \item First item
  \item Second item
  \item Third item with sub-items:
    \begin{itemize}
      \item Sub-item 1
      \item Sub-item 2
    \end{itemize}
\end{itemize}

\subsection{Numbered List}
\begin{enumerate}
  \item First step
  \item Second step
  \item Third step
\end{enumerate}

\section{Citations and References}
When you use the citation search feature, references will be automatically managed. For example, you might cite a paper like this: \cite{example2024}. All citations will be automatically added to your \texttt{cite.bib} file.

\section{Conclusion}
This template demonstrates basic \LaTeX{} features including mathematical formulas, tables, and lists. The bibliography system is now set up to work seamlessly with Bibby AI's citation management. Feel free to modify it for your needs!

\bibliographystyle{plain}
\bibliography{cite}

\end{document}\documentclass{article}
\usepackage[utf8]{inputenc}
\usepackage{booktabs}
\usepackage{amsmath}
\usepackage{natbib}

\title{My Document}
\author{Your Name}
\date{\today}

\begin{document}

\maketitle

\section{Introduction}
This is a sample document demonstrating various \LaTeX{} features including mathematical equations, tables, and formatting. You can cite papers using the citation search feature, and they will be automatically added to your \texttt{cite.bib} file.

\section{Mathematical Expressions}
Here are some example equations:

\subsection{Quadratic Formula}
The quadratic formula is given by:
\[x = \frac{-b \pm \sqrt{b^2 - 4ac}}{2a}\]

\subsection{Matrix Example}
A sample matrix:
\[
\begin{pmatrix}
a_{11} & a_{12} & a_{13} \\
a_{21} & a_{22} & a_{23} \\
a_{31} & a_{32} & a_{33}
\end{pmatrix}
\]

\section{Tables}
Below is a sample table showing data:

\begin{table}[h]
\centering
\begin{tabular}{@{}lrrr@{}}
\toprule
\textbf{Category} & \textbf{Value 1} & \textbf{Value 2} & \textbf{Total} \\
\midrule
Item A & 10.5 & 20.3 & 30.8 \\
Item B & 15.2 & 25.1 & 40.3 \\
Item C & 22.1 & 12.7 & 34.8 \\
\bottomrule
\end{tabular}
\caption{Sample Data Table}
\label{tab:sample}
\end{table}

\section{Lists}
\subsection{Bullet Points}
\begin{itemize}
  \item First item
  \item Second item
  \item Third item with sub-items:
    \begin{itemize}
      \item Sub-item 1
      \item Sub-item 2
    \end{itemize}
\end{itemize}

\subsection{Numbered List}
\begin{enumerate}
  \item First step
  \item Second step
  \item Third step
\end{enumerate}

\section{Citations and References}
When you use the citation search feature, references will be automatically managed. For example, you might cite a paper like this: \cite{example2024}. All citations will be automatically added to your \texttt{cite.bib} file.

\section{Conclusion}
This template demonstrates basic \LaTeX{} features including mathematical formulas, tables, and lists. The bibliography system is now set up to work seamlessly with Bibby AI's citation management. Feel free to modify it for your needs!

\bibliographystyle{plain}
\bibliography{cite}

\end{document}\documentclass{article}
\usepackage[utf8]{inputenc}
\usepackage{booktabs}
\usepackage{amsmath}
\usepackage{natbib}

\title{My Document}
\author{Your Name}
\date{\today}

\begin{document}

\maketitle

\section{Introduction}
This is a sample document demonstrating various \LaTeX{} features including mathematical equations, tables, and formatting. You can cite papers using the citation search feature, and they will be automatically added to your \texttt{cite.bib} file.

\section{Mathematical Expressions}
Here are some example equations:

\subsection{Quadratic Formula}
The quadratic formula is given by:
\[x = \frac{-b \pm \sqrt{b^2 - 4ac}}{2a}\]

\subsection{Matrix Example}
A sample matrix:
\[
\begin{pmatrix}
a_{11} & a_{12} & a_{13} \\
a_{21} & a_{22} & a_{23} \\
a_{31} & a_{32} & a_{33}
\end{pmatrix}
\]

\section{Tables}
Below is a sample table showing data:

\begin{table}[h]
\centering
\begin{tabular}{@{}lrrr@{}}
\toprule
\textbf{Category} & \textbf{Value 1} & \textbf{Value 2} & \textbf{Total} \\
\midrule
Item A & 10.5 & 20.3 & 30.8 \\
Item B & 15.2 & 25.1 & 40.3 \\
Item C & 22.1 & 12.7 & 34.8 \\
\bottomrule
\end{tabular}
\caption{Sample Data Table}
\label{tab:sample}
\end{table}

\section{Lists}
\subsection{Bullet Points}
\begin{itemize}
  \item First item
  \item Second item
  \item Third item with sub-items:
    \begin{itemize}
      \item Sub-item 1
      \item Sub-item 2
    \end{itemize}
\end{itemize}

\subsection{Numbered List}
\begin{enumerate}
  \item First step
  \item Second step
  \item Third step
\end{enumerate}

\section{Citations and References}
When you use the citation search feature, references will be automatically managed. For example, you might cite a paper like this: \cite{example2024}. All citations will be automatically added to your \texttt{cite.bib} file.

\section{Conclusion}
This template demonstrates basic \LaTeX{} features including mathematical formulas, tables, and lists. The bibliography system is now set up to work seamlessly with Bibby AI's citation management. Feel free to modify it for your needs!

\bibliographystyle{plain}
\bibliography{cite}

\end{document}\documentclass{article}
\usepackage[utf8]{inputenc}
\usepackage{booktabs}
\usepackage{amsmath}
\usepackage{natbib}

\title{My Document}
\author{Your Name}
\date{\today}

\begin{document}

\maketitle

\section{Introduction}
This is a sample document demonstrating various \LaTeX{} features including mathematical equations, tables, and formatting. You can cite papers using the citation search feature, and they will be automatically added to your \texttt{cite.bib} file.

\section{Mathematical Expressions}
Here are some example equations:

\subsection{Quadratic Formula}
The quadratic formula is given by:
\[x = \frac{-b \pm \sqrt{b^2 - 4ac}}{2a}\]

\subsection{Matrix Example}
A sample matrix:
\[
\begin{pmatrix}
a_{11} & a_{12} & a_{13} \\
a_{21} & a_{22} & a_{23} \\
a_{31} & a_{32} & a_{33}
\end{pmatrix}
\]

\section{Tables}
Below is a sample table showing data:

\begin{table}[h]
\centering
\begin{tabular}{@{}lrrr@{}}
\toprule
\textbf{Category} & \textbf{Value 1} & \textbf{Value 2} & \textbf{Total} \\
\midrule
Item A & 10.5 & 20.3 & 30.8 \\
Item B & 15.2 & 25.1 & 40.3 \\
Item C & 22.1 & 12.7 & 34.8 \\
\bottomrule
\end{tabular}
\caption{Sample Data Table}
\label{tab:sample}
\end{table}

\section{Lists}
\subsection{Bullet Points}
\begin{itemize}
  \item First item
  \item Second item
  \item Third item with sub-items:
    \begin{itemize}
      \item Sub-item 1
      \item Sub-item 2
    \end{itemize}
\end{itemize}

\subsection{Numbered List}
\begin{enumerate}
  \item First step
  \item Second step
  \item Third step
\end{enumerate}

\section{Citations and References}
When you use the citation search feature, references will be automatically managed. For example, you might cite a paper like this: \cite{example2024}. All citations will be automatically added to your \texttt{cite.bib} file.

\section{Conclusion}
This template demonstrates basic \LaTeX{} features including mathematical formulas, tables, and lists. The bibliography system is now set up to work seamlessly with Bibby AI's citation management. Feel free to modify it for your needs!

\bibliographystyle{plain}
\bibliography{cite}

\end{document}\documentclass{article}
\usepackage[utf8]{inputenc}
\usepackage{booktabs}
\usepackage{amsmath}
\usepackage{natbib}

\title{My Document}
\author{Your Name}
\date{\today}

\begin{document}

\maketitle

\section{Introduction}
This is a sample document demonstrating various \LaTeX{} features including mathematical equations, tables, and formatting. You can cite papers using the citation search feature, and they will be automatically added to your \texttt{cite.bib} file.

\section{Mathematical Expressions}
Here are some example equations:

\subsection{Quadratic Formula}
The quadratic formula is given by:
\[x = \frac{-b \pm \sqrt{b^2 - 4ac}}{2a}\]

\subsection{Matrix Example}
A sample matrix:
\[
\begin{pmatrix}
a_{11} & a_{12} & a_{13} \\
a_{21} & a_{22} & a_{23} \\
a_{31} & a_{32} & a_{33}
\end{pmatrix}
\]

\section{Tables}
Below is a sample table showing data:

\begin{table}[h]
\centering
\begin{tabular}{@{}lrrr@{}}
\toprule
\textbf{Category} & \textbf{Value 1} & \textbf{Value 2} & \textbf{Total} \\
\midrule
Item A & 10.5 & 20.3 & 30.8 \\
Item B & 15.2 & 25.1 & 40.3 \\
Item C & 22.1 & 12.7 & 34.8 \\
\bottomrule
\end{tabular}
\caption{Sample Data Table}
\label{tab:sample}
\end{table}

\section{Lists}
\subsection{Bullet Points}
\begin{itemize}
  \item First item
  \item Second item
  \item Third item with sub-items:
    \begin{itemize}
      \item Sub-item 1
      \item Sub-item 2
    \end{itemize}
\end{itemize}

\subsection{Numbered List}
\begin{enumerate}
  \item First step
  \item Second step
  \item Third step
\end{enumerate}

\section{Citations and References}
When you use the citation search feature, references will be automatically managed. For example, you might cite a paper like this: \cite{example2024}. All citations will be automatically added to your \texttt{cite.bib} file.

\section{Conclusion}
This template demonstrates basic \LaTeX{} features including mathematical formulas, tables, and lists. The bibliography system is now set up to work seamlessly with Bibby AI's citation management. Feel free to modify it for your needs!

\bibliographystyle{plain}
\bibliography{cite}

\end{document}\documentclass{article}
\usepackage[utf8]{inputenc}
\usepackage{booktabs}
\usepackage{amsmath}
\usepackage{natbib}

\title{My Document}
\author{Your Name}
\date{\today}

\begin{document}

\maketitle

\section{Introduction}
This is a sample document demonstrating various \LaTeX{} features including mathematical equations, tables, and formatting. You can cite papers using the citation search feature, and they will be automatically added to your \texttt{cite.bib} file.

\section{Mathematical Expressions}
Here are some example equations:

\subsection{Quadratic Formula}
The quadratic formula is given by:
\[x = \frac{-b \pm \sqrt{b^2 - 4ac}}{2a}\]

\subsection{Matrix Example}
A sample matrix:
\[
\begin{pmatrix}
a_{11} & a_{12} & a_{13} \\
a_{21} & a_{22} & a_{23} \\
a_{31} & a_{32} & a_{33}
\end{pmatrix}
\]

\section{Tables}
Below is a sample table showing data:

\begin{table}[h]
\centering
\begin{tabular}{@{}lrrr@{}}
\toprule
\textbf{Category} & \textbf{Value 1} & \textbf{Value 2} & \textbf{Total} \\
\midrule
Item A & 10.5 & 20.3 & 30.8 \\
Item B & 15.2 & 25.1 & 40.3 \\
Item C & 22.1 & 12.7 & 34.8 \\
\bottomrule
\end{tabular}
\caption{Sample Data Table}
\label{tab:sample}
\end{table}

\section{Lists}
\subsection{Bullet Points}
\begin{itemize}
  \item First item
  \item Second item
  \item Third item with sub-items:
    \begin{itemize}
      \item Sub-item 1
      \item Sub-item 2
    \end{itemize}
\end{itemize}

\subsection{Numbered List}
\begin{enumerate}
  \item First step
  \item Second step
  \item Third step
\end{enumerate}

\section{Citations and References}
When you use the citation search feature, references will be automatically managed. For example, you might cite a paper like this: \cite{example2024}. All citations will be automatically added to your \texttt{cite.bib} file.

\section{Conclusion}
This template demonstrates basic \LaTeX{} features including mathematical formulas, tables, and lists. The bibliography system is now set up to work seamlessly with Bibby AI's citation management. Feel free to modify it for your needs!

\bibliographystyle{plain}
\bibliography{cite}

\end{document}\documentclass{article}
\usepackage[utf8]{inputenc}
\usepackage{booktabs}
\usepackage{amsmath}
\usepackage{natbib}

\title{My Document}
\author{Your Name}
\date{\today}

\begin{document}

\maketitle

\section{Introduction}
This is a sample document demonstrating various \LaTeX{} features including mathematical equations, tables, and formatting. You can cite papers using the citation search feature, and they will be automatically added to your \texttt{cite.bib} file.

\section{Mathematical Expressions}
Here are some example equations:

\subsection{Quadratic Formula}
The quadratic formula is given by:
\[x = \frac{-b \pm \sqrt{b^2 - 4ac}}{2a}\]

\subsection{Matrix Example}
A sample matrix:
\[
\begin{pmatrix}
a_{11} & a_{12} & a_{13} \\
a_{21} & a_{22} & a_{23} \\
a_{31} & a_{32} & a_{33}
\end{pmatrix}
\]

\section{Tables}
Below is a sample table showing data:

\begin{table}[h]
\centering
\begin{tabular}{@{}lrrr@{}}
\toprule
\textbf{Category} & \textbf{Value 1} & \textbf{Value 2} & \textbf{Total} \\
\midrule
Item A & 10.5 & 20.3 & 30.8 \\
Item B & 15.2 & 25.1 & 40.3 \\
Item C & 22.1 & 12.7 & 34.8 \\
\bottomrule
\end{tabular}
\caption{Sample Data Table}
\label{tab:sample}
\end{table}

\section{Lists}
\subsection{Bullet Points}
\begin{itemize}
  \item First item
  \item Second item
  \item Third item with sub-items:
    \begin{itemize}
      \item Sub-item 1
      \item Sub-item 2
    \end{itemize}
\end{itemize}

\subsection{Numbered List}
\begin{enumerate}
  \item First step
  \item Second step
  \item Third step
\end{enumerate}

\section{Citations and References}
When you use the citation search feature, references will be automatically managed. For example, you might cite a paper like this: \cite{example2024}. All citations will be automatically added to your \texttt{cite.bib} file.

\section{Conclusion}
This template demonstrates basic \LaTeX{} features including mathematical formulas, tables, and lists. The bibliography system is now set up to work seamlessly with Bibby AI's citation management. Feel free to modify it for your needs!

\bibliographystyle{plain}
\bibliography{cite}

\end{document}\documentclass{article}
\usepackage[utf8]{inputenc}
\usepackage{booktabs}
\usepackage{amsmath}
\usepackage{natbib}

\title{My Document}
\author{Your Name}
\date{\today}

\begin{document}

\maketitle

\section{Introduction}
This is a sample document demonstrating various \LaTeX{} features including mathematical equations, tables, and formatting. You can cite papers using the citation search feature, and they will be automatically added to your \texttt{cite.bib} file.

\section{Mathematical Expressions}
Here are some example equations:

\subsection{Quadratic Formula}
The quadratic formula is given by:
\[x = \frac{-b \pm \sqrt{b^2 - 4ac}}{2a}\]

\subsection{Matrix Example}
A sample matrix:
\[
\begin{pmatrix}
a_{11} & a_{12} & a_{13} \\
a_{21} & a_{22} & a_{23} \\
a_{31} & a_{32} & a_{33}
\end{pmatrix}
\]

\section{Tables}
Below is a sample table showing data:

\begin{table}[h]
\centering
\begin{tabular}{@{}lrrr@{}}
\toprule
\textbf{Category} & \textbf{Value 1} & \textbf{Value 2} & \textbf{Total} \\
\midrule
Item A & 10.5 & 20.3 & 30.8 \\
Item B & 15.2 & 25.1 & 40.3 \\
Item C & 22.1 & 12.7 & 34.8 \\
\bottomrule
\end{tabular}
\caption{Sample Data Table}
\label{tab:sample}
\end{table}

\section{Lists}
\subsection{Bullet Points}
\begin{itemize}
  \item First item
  \item Second item
  \item Third item with sub-items:
    \begin{itemize}
      \item Sub-item 1
      \item Sub-item 2
    \end{itemize}
\end{itemize}

\subsection{Numbered List}
\begin{enumerate}
  \item First step
  \item Second step
  \item Third step
\end{enumerate}

\section{Citations and References}
When you use the citation search feature, references will be automatically managed. For example, you might cite a paper like this: \cite{example2024}. All citations will be automatically added to your \texttt{cite.bib} file.

\section{Conclusion}
This template demonstrates basic \LaTeX{} features including mathematical formulas, tables, and lists. The bibliography system is now set up to work seamlessly with Bibby AI's citation management. Feel free to modify it for your needs!

\bibliographystyle{plain}
\bibliography{cite}

\end{document}\documentclass{article}
\usepackage[utf8]{inputenc}
\usepackage{booktabs}
\usepackage{amsmath}
\usepackage{natbib}

\title{My Document}
\author{Your Name}
\date{\today}

\begin{document}

\maketitle

\section{Introduction}
This is a sample document demonstrating various \LaTeX{} features including mathematical equations, tables, and formatting. You can cite papers using the citation search feature, and they will be automatically added to your \texttt{cite.bib} file.

\section{Mathematical Expressions}
Here are some example equations:

\subsection{Quadratic Formula}
The quadratic formula is given by:
\[x = \frac{-b \pm \sqrt{b^2 - 4ac}}{2a}\]

\subsection{Matrix Example}
A sample matrix:
\[
\begin{pmatrix}
a_{11} & a_{12} & a_{13} \\
a_{21} & a_{22} & a_{23} \\
a_{31} & a_{32} & a_{33}
\end{pmatrix}
\]

\section{Tables}
Below is a sample table showing data:

\begin{table}[h]
\centering
\begin{tabular}{@{}lrrr@{}}
\toprule
\textbf{Category} & \textbf{Value 1} & \textbf{Value 2} & \textbf{Total} \\
\midrule
Item A & 10.5 & 20.3 & 30.8 \\
Item B & 15.2 & 25.1 & 40.3 \\
Item C & 22.1 & 12.7 & 34.8 \\
\bottomrule
\end{tabular}
\caption{Sample Data Table}
\label{tab:sample}
\end{table}

\section{Lists}
\subsection{Bullet Points}
\begin{itemize}
  \item First item
  \item Second item
  \item Third item with sub-items:
    \begin{itemize}
      \item Sub-item 1
      \item Sub-item 2
    \end{itemize}
\end{itemize}

\subsection{Numbered List}
\begin{enumerate}
  \item First step
  \item Second step
  \item Third step
\end{enumerate}

\section{Citations and References}
When you use the citation search feature, references will be automatically managed. For example, you might cite a paper like this: \cite{example2024}. All citations will be automatically added to your \texttt{cite.bib} file.

\section{Conclusion}
This template demonstrates basic \LaTeX{} features including mathematical formulas, tables, and lists. The bibliography system is now set up to work seamlessly with Bibby AI's citation management. Feel free to modify it for your needs!

\bibliographystyle{plain}
\bibliography{cite}

\end{document}\documentclass{article}
\usepackage[utf8]{inputenc}
\usepackage{booktabs}
\usepackage{amsmath}
\usepackage{natbib}

\title{My Document}
\author{Your Name}
\date{\today}

\begin{document}

\maketitle

\section{Introduction}
This is a sample document demonstrating various \LaTeX{} features including mathematical equations, tables, and formatting. You can cite papers using the citation search feature, and they will be automatically added to your \texttt{cite.bib} file.

\section{Mathematical Expressions}
Here are some example equations:

\subsection{Quadratic Formula}
The quadratic formula is given by:
\[x = \frac{-b \pm \sqrt{b^2 - 4ac}}{2a}\]

\subsection{Matrix Example}
A sample matrix:
\[
\begin{pmatrix}
a_{11} & a_{12} & a_{13} \\
a_{21} & a_{22} & a_{23} \\
a_{31} & a_{32} & a_{33}
\end{pmatrix}
\]

\section{Tables}
Below is a sample table showing data:

\begin{table}[h]
\centering
\begin{tabular}{@{}lrrr@{}}
\toprule
\textbf{Category} & \textbf{Value 1} & \textbf{Value 2} & \textbf{Total} \\
\midrule
Item A & 10.5 & 20.3 & 30.8 \\
Item B & 15.2 & 25.1 & 40.3 \\
Item C & 22.1 & 12.7 & 34.8 \\
\bottomrule
\end{tabular}
\caption{Sample Data Table}
\label{tab:sample}
\end{table}

\section{Lists}
\subsection{Bullet Points}
\begin{itemize}
  \item First item
  \item Second item
  \item Third item with sub-items:
    \begin{itemize}
      \item Sub-item 1
      \item Sub-item 2
    \end{itemize}
\end{itemize}

\subsection{Numbered List}
\begin{enumerate}
  \item First step
  \item Second step
  \item Third step
\end{enumerate}

\section{Citations and References}
When you use the citation search feature, references will be automatically managed. For example, you might cite a paper like this: \cite{example2024}. All citations will be automatically added to your \texttt{cite.bib} file.

\section{Conclusion}
This template demonstrates basic \LaTeX{} features including mathematical formulas, tables, and lists. The bibliography system is now set up to work seamlessly with Bibby AI's citation management. Feel free to modify it for your needs!

\bibliographystyle{plain}
\bibliography{cite}

\end{document}\documentclass{article}
\usepackage[utf8]{inputenc}
\usepackage{booktabs}
\usepackage{amsmath}
\usepackage{natbib}

\title{My Document}
\author{Your Name}
\date{\today}

\begin{document}

\maketitle

\section{Introduction}
This is a sample document demonstrating various \LaTeX{} features including mathematical equations, tables, and formatting. You can cite papers using the citation search feature, and they will be automatically added to your \texttt{cite.bib} file.

\section{Mathematical Expressions}
Here are some example equations:

\subsection{Quadratic Formula}
The quadratic formula is given by:
\[x = \frac{-b \pm \sqrt{b^2 - 4ac}}{2a}\]

\subsection{Matrix Example}
A sample matrix:
\[
\begin{pmatrix}
a_{11} & a_{12} & a_{13} \\
a_{21} & a_{22} & a_{23} \\
a_{31} & a_{32} & a_{33}
\end{pmatrix}
\]

\section{Tables}
Below is a sample table showing data:

\begin{table}[h]
\centering
\begin{tabular}{@{}lrrr@{}}
\toprule
\textbf{Category} & \textbf{Value 1} & \textbf{Value 2} & \textbf{Total} \\
\midrule
Item A & 10.5 & 20.3 & 30.8 \\
Item B & 15.2 & 25.1 & 40.3 \\
Item C & 22.1 & 12.7 & 34.8 \\
\bottomrule
\end{tabular}
\caption{Sample Data Table}
\label{tab:sample}
\end{table}

\section{Lists}
\subsection{Bullet Points}
\begin{itemize}
  \item First item
  \item Second item
  \item Third item with sub-items:
    \begin{itemize}
      \item Sub-item 1
      \item Sub-item 2
    \end{itemize}
\end{itemize}

\subsection{Numbered List}
\begin{enumerate}
  \item First step
  \item Second step
  \item Third step
\end{enumerate}

\section{Citations and References}
When you use the citation search feature, references will be automatically managed. For example, you might cite a paper like this: \cite{example2024}. All citations will be automatically added to your \texttt{cite.bib} file.

\section{Conclusion}
This template demonstrates basic \LaTeX{} features including mathematical formulas, tables, and lists. The bibliography system is now set up to work seamlessly with Bibby AI's citation management. Feel free to modify it for your needs!

\bibliographystyle{plain}
\bibliography{cite}

\end{document}\documentclass{article}
\usepackage[utf8]{inputenc}
\usepackage{booktabs}
\usepackage{amsmath}
\usepackage{natbib}

\title{My Document}
\author{Your Name}
\date{\today}

\begin{document}

\maketitle

\section{Introduction}
This is a sample document demonstrating various \LaTeX{} features including mathematical equations, tables, and formatting. You can cite papers using the citation search feature, and they will be automatically added to your \texttt{cite.bib} file.

\section{Mathematical Expressions}
Here are some example equations:

\subsection{Quadratic Formula}
The quadratic formula is given by:
\[x = \frac{-b \pm \sqrt{b^2 - 4ac}}{2a}\]

\subsection{Matrix Example}
A sample matrix:
\[
\begin{pmatrix}
a_{11} & a_{12} & a_{13} \\
a_{21} & a_{22} & a_{23} \\
a_{31} & a_{32} & a_{33}
\end{pmatrix}
\]

\section{Tables}
Below is a sample table showing data:

\begin{table}[h]
\centering
\begin{tabular}{@{}lrrr@{}}
\toprule
\textbf{Category} & \textbf{Value 1} & \textbf{Value 2} & \textbf{Total} \\
\midrule
Item A & 10.5 & 20.3 & 30.8 \\
Item B & 15.2 & 25.1 & 40.3 \\
Item C & 22.1 & 12.7 & 34.8 \\
\bottomrule
\end{tabular}
\caption{Sample Data Table}
\label{tab:sample}
\end{table}

\section{Lists}
\subsection{Bullet Points}
\begin{itemize}
  \item First item
  \item Second item
  \item Third item with sub-items:
    \begin{itemize}
      \item Sub-item 1
      \item Sub-item 2
    \end{itemize}
\end{itemize}

\subsection{Numbered List}
\begin{enumerate}
  \item First step
  \item Second step
  \item Third step
\end{enumerate}

\section{Citations and References}
When you use the citation search feature, references will be automatically managed. For example, you might cite a paper like this: \cite{example2024}. All citations will be automatically added to your \texttt{cite.bib} file.

\section{Conclusion}
This template demonstrates basic \LaTeX{} features including mathematical formulas, tables, and lists. The bibliography system is now set up to work seamlessly with Bibby AI's citation management. Feel free to modify it for your needs!

\bibliographystyle{plain}
\bibliography{cite}

\end{document}\documentclass{article}
\usepackage[utf8]{inputenc}
\usepackage{booktabs}
\usepackage{amsmath}
\usepackage{natbib}

\title{My Document}
\author{Your Name}
\date{\today}

\begin{document}

\maketitle

\section{Introduction}
This is a sample document demonstrating various \LaTeX{} features including mathematical equations, tables, and formatting. You can cite papers using the citation search feature, and they will be automatically added to your \texttt{cite.bib} file.

\section{Mathematical Expressions}
Here are some example equations:

\subsection{Quadratic Formula}
The quadratic formula is given by:
\[x = \frac{-b \pm \sqrt{b^2 - 4ac}}{2a}\]

\subsection{Matrix Example}
A sample matrix:
\[
\begin{pmatrix}
a_{11} & a_{12} & a_{13} \\
a_{21} & a_{22} & a_{23} \\
a_{31} & a_{32} & a_{33}
\end{pmatrix}
\]

\section{Tables}
Below is a sample table showing data:

\begin{table}[h]
\centering
\begin{tabular}{@{}lrrr@{}}
\toprule
\textbf{Category} & \textbf{Value 1} & \textbf{Value 2} & \textbf{Total} \\
\midrule
Item A & 10.5 & 20.3 & 30.8 \\
Item B & 15.2 & 25.1 & 40.3 \\
Item C & 22.1 & 12.7 & 34.8 \\
\bottomrule
\end{tabular}
\caption{Sample Data Table}
\label{tab:sample}
\end{table}

\section{Lists}
\subsection{Bullet Points}
\begin{itemize}
  \item First item
  \item Second item
  \item Third item with sub-items:
    \begin{itemize}
      \item Sub-item 1
      \item Sub-item 2
    \end{itemize}
\end{itemize}

\subsection{Numbered List}
\begin{enumerate}
  \item First step
  \item Second step
  \item Third step
\end{enumerate}

\section{Citations and References}
When you use the citation search feature, references will be automatically managed. For example, you might cite a paper like this: \cite{example2024}. All citations will be automatically added to your \texttt{cite.bib} file.

\section{Conclusion}
This template demonstrates basic \LaTeX{} features including mathematical formulas, tables, and lists. The bibliography system is now set up to work seamlessly with Bibby AI's citation management. Feel free to modify it for your needs!

\bibliographystyle{plain}
\bibliography{cite}

\end{document}\documentclass{article}
\usepackage[utf8]{inputenc}
\usepackage{booktabs}
\usepackage{amsmath}
\usepackage{natbib}

\title{My Document}
\author{Your Name}
\date{\today}

\begin{document}

\maketitle

\section{Introduction}
This is a sample document demonstrating various \LaTeX{} features including mathematical equations, tables, and formatting. You can cite papers using the citation search feature, and they will be automatically added to your \texttt{cite.bib} file.

\section{Mathematical Expressions}
Here are some example equations:

\subsection{Quadratic Formula}
The quadratic formula is given by:
\[x = \frac{-b \pm \sqrt{b^2 - 4ac}}{2a}\]

\subsection{Matrix Example}
A sample matrix:
\[
\begin{pmatrix}
a_{11} & a_{12} & a_{13} \\
a_{21} & a_{22} & a_{23} \\
a_{31} & a_{32} & a_{33}
\end{pmatrix}
\]

\section{Tables}
Below is a sample table showing data:

\begin{table}[h]
\centering
\begin{tabular}{@{}lrrr@{}}
\toprule
\textbf{Category} & \textbf{Value 1} & \textbf{Value 2} & \textbf{Total} \\
\midrule
Item A & 10.5 & 20.3 & 30.8 \\
Item B & 15.2 & 25.1 & 40.3 \\
Item C & 22.1 & 12.7 & 34.8 \\
\bottomrule
\end{tabular}
\caption{Sample Data Table}
\label{tab:sample}
\end{table}

\section{Lists}
\subsection{Bullet Points}
\begin{itemize}
  \item First item
  \item Second item
  \item Third item with sub-items:
    \begin{itemize}
      \item Sub-item 1
      \item Sub-item 2
    \end{itemize}
\end{itemize}

\subsection{Numbered List}
\begin{enumerate}
  \item First step
  \item Second step
  \item Third step
\end{enumerate}

\section{Citations and References}
When you use the citation search feature, references will be automatically managed. For example, you might cite a paper like this: \cite{example2024}. All citations will be automatically added to your \texttt{cite.bib} file.

\section{Conclusion}
This template demonstrates basic \LaTeX{} features including mathematical formulas, tables, and lists. The bibliography system is now set up to work seamlessly with Bibby AI's citation management. Feel free to modify it for your needs!

\bibliographystyle{plain}
\bibliography{cite}

\end{document}\documentclass{article}
\usepackage[utf8]{inputenc}
\usepackage{booktabs}
\usepackage{amsmath}
\usepackage{natbib}

\title{My Document}
\author{Your Name}
\date{\today}

\begin{document}

\maketitle

\section{Introduction}
This is a sample document demonstrating various \LaTeX{} features including mathematical equations, tables, and formatting. You can cite papers using the citation search feature, and they will be automatically added to your \texttt{cite.bib} file.

\section{Mathematical Expressions}
Here are some example equations:

\subsection{Quadratic Formula}
The quadratic formula is given by:
\[x = \frac{-b \pm \sqrt{b^2 - 4ac}}{2a}\]

\subsection{Matrix Example}
A sample matrix:
\[
\begin{pmatrix}
a_{11} & a_{12} & a_{13} \\
a_{21} & a_{22} & a_{23} \\
a_{31} & a_{32} & a_{33}
\end{pmatrix}
\]

\section{Tables}
Below is a sample table showing data:

\begin{table}[h]
\centering
\begin{tabular}{@{}lrrr@{}}
\toprule
\textbf{Category} & \textbf{Value 1} & \textbf{Value 2} & \textbf{Total} \\
\midrule
Item A & 10.5 & 20.3 & 30.8 \\
Item B & 15.2 & 25.1 & 40.3 \\
Item C & 22.1 & 12.7 & 34.8 \\
\bottomrule
\end{tabular}
\caption{Sample Data Table}
\label{tab:sample}
\end{table}

\section{Lists}
\subsection{Bullet Points}
\begin{itemize}
  \item First item
  \item Second item
  \item Third item with sub-items:
    \begin{itemize}
      \item Sub-item 1
      \item Sub-item 2
    \end{itemize}
\end{itemize}

\subsection{Numbered List}
\begin{enumerate}
  \item First step
  \item Second step
  \item Third step
\end{enumerate}

\section{Citations and References}
When you use the citation search feature, references will be automatically managed. For example, you might cite a paper like this: \cite{example2024}. All citations will be automatically added to your \texttt{cite.bib} file.

\section{Conclusion}
This template demonstrates basic \LaTeX{} features including mathematical formulas, tables, and lists. The bibliography system is now set up to work seamlessly with Bibby AI's citation management. Feel free to modify it for your needs!

\bibliographystyle{plain}
\bibliography{cite}

\end{document}\documentclass{article}
\usepackage[utf8]{inputenc}
\usepackage{booktabs}
\usepackage{amsmath}
\usepackage{natbib}

\title{My Document}
\author{Your Name}
\date{\today}

\begin{document}

\maketitle

\section{Introduction}
This is a sample document demonstrating various \LaTeX{} features including mathematical equations, tables, and formatting. You can cite papers using the citation search feature, and they will be automatically added to your \texttt{cite.bib} file.

\section{Mathematical Expressions}
Here are some example equations:

\subsection{Quadratic Formula}
The quadratic formula is given by:
\[x = \frac{-b \pm \sqrt{b^2 - 4ac}}{2a}\]

\subsection{Matrix Example}
A sample matrix:
\[
\begin{pmatrix}
a_{11} & a_{12} & a_{13} \\
a_{21} & a_{22} & a_{23} \\
a_{31} & a_{32} & a_{33}
\end{pmatrix}
\]

\section{Tables}
Below is a sample table showing data:

\begin{table}[h]
\centering
\begin{tabular}{@{}lrrr@{}}
\toprule
\textbf{Category} & \textbf{Value 1} & \textbf{Value 2} & \textbf{Total} \\
\midrule
Item A & 10.5 & 20.3 & 30.8 \\
Item B & 15.2 & 25.1 & 40.3 \\
Item C & 22.1 & 12.7 & 34.8 \\
\bottomrule
\end{tabular}
\caption{Sample Data Table}
\label{tab:sample}
\end{table}

\section{Lists}
\subsection{Bullet Points}
\begin{itemize}
  \item First item
  \item Second item
  \item Third item with sub-items:
    \begin{itemize}
      \item Sub-item 1
      \item Sub-item 2
    \end{itemize}
\end{itemize}

\subsection{Numbered List}
\begin{enumerate}
  \item First step
  \item Second step
  \item Third step
\end{enumerate}

\section{Citations and References}
When you use the citation search feature, references will be automatically managed. For example, you might cite a paper like this: \cite{example2024}. All citations will be automatically added to your \texttt{cite.bib} file.

\section{Conclusion}
This template demonstrates basic \LaTeX{} features including mathematical formulas, tables, and lists. The bibliography system is now set up to work seamlessly with Bibby AI's citation management. Feel free to modify it for your needs!

\bibliographystyle{plain}
\bibliography{cite}

\end{document}\documentclass{article}
\usepackage[utf8]{inputenc}
\usepackage{booktabs}
\usepackage{amsmath}
\usepackage{natbib}

\title{My Document}
\author{Your Name}
\date{\today}

\begin{document}

\maketitle

\section{Introduction}
This is a sample document demonstrating various \LaTeX{} features including mathematical equations, tables, and formatting. You can cite papers using the citation search feature, and they will be automatically added to your \texttt{cite.bib} file.

\section{Mathematical Expressions}
Here are some example equations:

\subsection{Quadratic Formula}
The quadratic formula is given by:
\[x = \frac{-b \pm \sqrt{b^2 - 4ac}}{2a}\]

\subsection{Matrix Example}
A sample matrix:
\[
\begin{pmatrix}
a_{11} & a_{12} & a_{13} \\
a_{21} & a_{22} & a_{23} \\
a_{31} & a_{32} & a_{33}
\end{pmatrix}
\]

\section{Tables}
Below is a sample table showing data:

\begin{table}[h]
\centering
\begin{tabular}{@{}lrrr@{}}
\toprule
\textbf{Category} & \textbf{Value 1} & \textbf{Value 2} & \textbf{Total} \\
\midrule
Item A & 10.5 & 20.3 & 30.8 \\
Item B & 15.2 & 25.1 & 40.3 \\
Item C & 22.1 & 12.7 & 34.8 \\
\bottomrule
\end{tabular}
\caption{Sample Data Table}
\label{tab:sample}
\end{table}

\section{Lists}
\subsection{Bullet Points}
\begin{itemize}
  \item First item
  \item Second item
  \item Third item with sub-items:
    \begin{itemize}
      \item Sub-item 1
      \item Sub-item 2
    \end{itemize}
\end{itemize}

\subsection{Numbered List}
\begin{enumerate}
  \item First step
  \item Second step
  \item Third step
\end{enumerate}

\section{Citations and References}
When you use the citation search feature, references will be automatically managed. For example, you might cite a paper like this: \cite{example2024}. All citations will be automatically added to your \texttt{cite.bib} file.

\section{Conclusion}
This template demonstrates basic \LaTeX{} features including mathematical formulas, tables, and lists. The bibliography system is now set up to work seamlessly with Bibby AI's citation management. Feel free to modify it for your needs!

\bibliographystyle{plain}
\bibliography{cite}

\end{document}\documentclass{article}
\usepackage[utf8]{inputenc}
\usepackage{booktabs}
\usepackage{amsmath}
\usepackage{natbib}

\title{My Document}
\author{Your Name}
\date{\today}

\begin{document}

\maketitle

\section{Introduction}
This is a sample document demonstrating various \LaTeX{} features including mathematical equations, tables, and formatting. You can cite papers using the citation search feature, and they will be automatically added to your \texttt{cite.bib} file.

\section{Mathematical Expressions}
Here are some example equations:

\subsection{Quadratic Formula}
The quadratic formula is given by:
\[x = \frac{-b \pm \sqrt{b^2 - 4ac}}{2a}\]

\subsection{Matrix Example}
A sample matrix:
\[
\begin{pmatrix}
a_{11} & a_{12} & a_{13} \\
a_{21} & a_{22} & a_{23} \\
a_{31} & a_{32} & a_{33}
\end{pmatrix}
\]

\section{Tables}
Below is a sample table showing data:

\begin{table}[h]
\centering
\begin{tabular}{@{}lrrr@{}}
\toprule
\textbf{Category} & \textbf{Value 1} & \textbf{Value 2} & \textbf{Total} \\
\midrule
Item A & 10.5 & 20.3 & 30.8 \\
Item B & 15.2 & 25.1 & 40.3 \\
Item C & 22.1 & 12.7 & 34.8 \\
\bottomrule
\end{tabular}
\caption{Sample Data Table}
\label{tab:sample}
\end{table}

\section{Lists}
\subsection{Bullet Points}
\begin{itemize}
  \item First item
  \item Second item
  \item Third item with sub-items:
    \begin{itemize}
      \item Sub-item 1
      \item Sub-item 2
    \end{itemize}
\end{itemize}

\subsection{Numbered List}
\begin{enumerate}
  \item First step
  \item Second step
  \item Third step
\end{enumerate}

\section{Citations and References}
When you use the citation search feature, references will be automatically managed. For example, you might cite a paper like this: \cite{example2024}. All citations will be automatically added to your \texttt{cite.bib} file.

\section{Conclusion}
This template demonstrates basic \LaTeX{} features including mathematical formulas, tables, and lists. The bibliography system is now set up to work seamlessly with Bibby AI's citation management. Feel free to modify it for your needs!

\bibliographystyle{plain}
\bibliography{cite}

\end{document}\documentclass{article}
\usepackage[utf8]{inputenc}
\usepackage{booktabs}
\usepackage{amsmath}
\usepackage{natbib}

\title{My Document}
\author{Your Name}
\date{\today}

\begin{document}

\maketitle

\section{Introduction}
This is a sample document demonstrating various \LaTeX{} features including mathematical equations, tables, and formatting. You can cite papers using the citation search feature, and they will be automatically added to your \texttt{cite.bib} file.

\section{Mathematical Expressions}
Here are some example equations:

\subsection{Quadratic Formula}
The quadratic formula is given by:
\[x = \frac{-b \pm \sqrt{b^2 - 4ac}}{2a}\]

\subsection{Matrix Example}
A sample matrix:
\[
\begin{pmatrix}
a_{11} & a_{12} & a_{13} \\
a_{21} & a_{22} & a_{23} \\
a_{31} & a_{32} & a_{33}
\end{pmatrix}
\]

\section{Tables}
Below is a sample table showing data:

\begin{table}[h]
\centering
\begin{tabular}{@{}lrrr@{}}
\toprule
\textbf{Category} & \textbf{Value 1} & \textbf{Value 2} & \textbf{Total} \\
\midrule
Item A & 10.5 & 20.3 & 30.8 \\
Item B & 15.2 & 25.1 & 40.3 \\
Item C & 22.1 & 12.7 & 34.8 \\
\bottomrule
\end{tabular}
\caption{Sample Data Table}
\label{tab:sample}
\end{table}

\section{Lists}
\subsection{Bullet Points}
\begin{itemize}
  \item First item
  \item Second item
  \item Third item with sub-items:
    \begin{itemize}
      \item Sub-item 1
      \item Sub-item 2
    \end{itemize}
\end{itemize}

\subsection{Numbered List}
\begin{enumerate}
  \item First step
  \item Second step
  \item Third step
\end{enumerate}

\section{Citations and References}
When you use the citation search feature, references will be automatically managed. For example, you might cite a paper like this: \cite{example2024}. All citations will be automatically added to your \texttt{cite.bib} file.

\section{Conclusion}
This template demonstrates basic \LaTeX{} features including mathematical formulas, tables, and lists. The bibliography system is now set up to work seamlessly with Bibby AI's citation management. Feel free to modify it for your needs!

\bibliographystyle{plain}
\bibliography{cite}

\end{document}\documentclass{article}
\usepackage[utf8]{inputenc}
\usepackage{booktabs}
\usepackage{amsmath}
\usepackage{natbib}

\title{My Document}
\author{Your Name}
\date{\today}

\begin{document}

\maketitle

\section{Introduction}
This is a sample document demonstrating various \LaTeX{} features including mathematical equations, tables, and formatting. You can cite papers using the citation search feature, and they will be automatically added to your \texttt{cite.bib} file.

\section{Mathematical Expressions}
Here are some example equations:

\subsection{Quadratic Formula}
The quadratic formula is given by:
\[x = \frac{-b \pm \sqrt{b^2 - 4ac}}{2a}\]

\subsection{Matrix Example}
A sample matrix:
\[
\begin{pmatrix}
a_{11} & a_{12} & a_{13} \\
a_{21} & a_{22} & a_{23} \\
a_{31} & a_{32} & a_{33}
\end{pmatrix}
\]

\section{Tables}
Below is a sample table showing data:

\begin{table}[h]
\centering
\begin{tabular}{@{}lrrr@{}}
\toprule
\textbf{Category} & \textbf{Value 1} & \textbf{Value 2} & \textbf{Total} \\
\midrule
Item A & 10.5 & 20.3 & 30.8 \\
Item B & 15.2 & 25.1 & 40.3 \\
Item C & 22.1 & 12.7 & 34.8 \\
\bottomrule
\end{tabular}
\caption{Sample Data Table}
\label{tab:sample}
\end{table}

\section{Lists}
\subsection{Bullet Points}
\begin{itemize}
  \item First item
  \item Second item
  \item Third item with sub-items:
    \begin{itemize}
      \item Sub-item 1
      \item Sub-item 2
    \end{itemize}
\end{itemize}

\subsection{Numbered List}
\begin{enumerate}
  \item First step
  \item Second step
  \item Third step
\end{enumerate}

\section{Citations and References}
When you use the citation search feature, references will be automatically managed. For example, you might cite a paper like this: \cite{example2024}. All citations will be automatically added to your \texttt{cite.bib} file.

\section{Conclusion}
This template demonstrates basic \LaTeX{} features including mathematical formulas, tables, and lists. The bibliography system is now set up to work seamlessly with Bibby AI's citation management. Feel free to modify it for your needs!

\bibliographystyle{plain}
\bibliography{cite}

\end{document}\documentclass{article}
\usepackage[utf8]{inputenc}
\usepackage{booktabs}
\usepackage{amsmath}
\usepackage{natbib}

\title{My Document}
\author{Your Name}
\date{\today}

\begin{document}

\maketitle

\section{Introduction}
This is a sample document demonstrating various \LaTeX{} features including mathematical equations, tables, and formatting. You can cite papers using the citation search feature, and they will be automatically added to your \texttt{cite.bib} file.

\section{Mathematical Expressions}
Here are some example equations:

\subsection{Quadratic Formula}
The quadratic formula is given by:
\[x = \frac{-b \pm \sqrt{b^2 - 4ac}}{2a}\]

\subsection{Matrix Example}
A sample matrix:
\[
\begin{pmatrix}
a_{11} & a_{12} & a_{13} \\
a_{21} & a_{22} & a_{23} \\
a_{31} & a_{32} & a_{33}
\end{pmatrix}
\]

\section{Tables}
Below is a sample table showing data:

\begin{table}[h]
\centering
\begin{tabular}{@{}lrrr@{}}
\toprule
\textbf{Category} & \textbf{Value 1} & \textbf{Value 2} & \textbf{Total} \\
\midrule
Item A & 10.5 & 20.3 & 30.8 \\
Item B & 15.2 & 25.1 & 40.3 \\
Item C & 22.1 & 12.7 & 34.8 \\
\bottomrule
\end{tabular}
\caption{Sample Data Table}
\label{tab:sample}
\end{table}

\section{Lists}
\subsection{Bullet Points}
\begin{itemize}
  \item First item
  \item Second item
  \item Third item with sub-items:
    \begin{itemize}
      \item Sub-item 1
      \item Sub-item 2
    \end{itemize}
\end{itemize}

\subsection{Numbered List}
\begin{enumerate}
  \item First step
  \item Second step
  \item Third step
\end{enumerate}

\section{Citations and References}
When you use the citation search feature, references will be automatically managed. For example, you might cite a paper like this: \cite{example2024}. All citations will be automatically added to your \texttt{cite.bib} file.

\section{Conclusion}
This template demonstrates basic \LaTeX{} features including mathematical formulas, tables, and lists. The bibliography system is now set up to work seamlessly with Bibby AI's citation management. Feel free to modify it for your needs!

\bibliographystyle{plain}
\bibliography{cite}

\end{document}\documentclass{article}
\usepackage[utf8]{inputenc}
\usepackage{booktabs}
\usepackage{amsmath}
\usepackage{natbib}

\title{My Document}
\author{Your Name}
\date{\today}

\begin{document}

\maketitle

\section{Introduction}
This is a sample document demonstrating various \LaTeX{} features including mathematical equations, tables, and formatting. You can cite papers using the citation search feature, and they will be automatically added to your \texttt{cite.bib} file.

\section{Mathematical Expressions}
Here are some example equations:

\subsection{Quadratic Formula}
The quadratic formula is given by:
\[x = \frac{-b \pm \sqrt{b^2 - 4ac}}{2a}\]

\subsection{Matrix Example}
A sample matrix:
\[
\begin{pmatrix}
a_{11} & a_{12} & a_{13} \\
a_{21} & a_{22} & a_{23} \\
a_{31} & a_{32} & a_{33}
\end{pmatrix}
\]

\section{Tables}
Below is a sample table showing data:

\begin{table}[h]
\centering
\begin{tabular}{@{}lrrr@{}}
\toprule
\textbf{Category} & \textbf{Value 1} & \textbf{Value 2} & \textbf{Total} \\
\midrule
Item A & 10.5 & 20.3 & 30.8 \\
Item B & 15.2 & 25.1 & 40.3 \\
Item C & 22.1 & 12.7 & 34.8 \\
\bottomrule
\end{tabular}
\caption{Sample Data Table}
\label{tab:sample}
\end{table}

\section{Lists}
\subsection{Bullet Points}
\begin{itemize}
  \item First item
  \item Second item
  \item Third item with sub-items:
    \begin{itemize}
      \item Sub-item 1
      \item Sub-item 2
    \end{itemize}
\end{itemize}

\subsection{Numbered List}
\begin{enumerate}
  \item First step
  \item Second step
  \item Third step
\end{enumerate}

\section{Citations and References}
When you use the citation search feature, references will be automatically managed. For example, you might cite a paper like this: \cite{example2024}. All citations will be automatically added to your \texttt{cite.bib} file.

\section{Conclusion}
This template demonstrates basic \LaTeX{} features including mathematical formulas, tables, and lists. The bibliography system is now set up to work seamlessly with Bibby AI's citation management. Feel free to modify it for your needs!

\bibliographystyle{plain}
\bibliography{cite}

\end{document}\documentclass{article}
\usepackage[utf8]{inputenc}
\usepackage{booktabs}
\usepackage{amsmath}
\usepackage{natbib}

\title{My Document}
\author{Your Name}
\date{\today}

\begin{document}

\maketitle

\section{Introduction}
This is a sample document demonstrating various \LaTeX{} features including mathematical equations, tables, and formatting. You can cite papers using the citation search feature, and they will be automatically added to your \texttt{cite.bib} file.

\section{Mathematical Expressions}
Here are some example equations:

\subsection{Quadratic Formula}
The quadratic formula is given by:
\[x = \frac{-b \pm \sqrt{b^2 - 4ac}}{2a}\]

\subsection{Matrix Example}
A sample matrix:
\[
\begin{pmatrix}
a_{11} & a_{12} & a_{13} \\
a_{21} & a_{22} & a_{23} \\
a_{31} & a_{32} & a_{33}
\end{pmatrix}
\]

\section{Tables}
Below is a sample table showing data:

\begin{table}[h]
\centering
\begin{tabular}{@{}lrrr@{}}
\toprule
\textbf{Category} & \textbf{Value 1} & \textbf{Value 2} & \textbf{Total} \\
\midrule
Item A & 10.5 & 20.3 & 30.8 \\
Item B & 15.2 & 25.1 & 40.3 \\
Item C & 22.1 & 12.7 & 34.8 \\
\bottomrule
\end{tabular}
\caption{Sample Data Table}
\label{tab:sample}
\end{table}

\section{Lists}
\subsection{Bullet Points}
\begin{itemize}
  \item First item
  \item Second item
  \item Third item with sub-items:
    \begin{itemize}
      \item Sub-item 1
      \item Sub-item 2
    \end{itemize}
\end{itemize}

\subsection{Numbered List}
\begin{enumerate}
  \item First step
  \item Second step
  \item Third step
\end{enumerate}

\section{Citations and References}
When you use the citation search feature, references will be automatically managed. For example, you might cite a paper like this: \cite{example2024}. All citations will be automatically added to your \texttt{cite.bib} file.

\section{Conclusion}
This template demonstrates basic \LaTeX{} features including mathematical formulas, tables, and lists. The bibliography system is now set up to work seamlessly with Bibby AI's citation management. Feel free to modify it for your needs!

\bibliographystyle{plain}
\bibliography{cite}

\end{document}\documentclass{article}
\usepackage[utf8]{inputenc}
\usepackage{booktabs}
\usepackage{amsmath}
\usepackage{natbib}

\title{My Document}
\author{Your Name}
\date{\today}

\begin{document}

\maketitle

\section{Introduction}
This is a sample document demonstrating various \LaTeX{} features including mathematical equations, tables, and formatting. You can cite papers using the citation search feature, and they will be automatically added to your \texttt{cite.bib} file.

\section{Mathematical Expressions}
Here are some example equations:

\subsection{Quadratic Formula}
The quadratic formula is given by:
\[x = \frac{-b \pm \sqrt{b^2 - 4ac}}{2a}\]

\subsection{Matrix Example}
A sample matrix:
\[
\begin{pmatrix}
a_{11} & a_{12} & a_{13} \\
a_{21} & a_{22} & a_{23} \\
a_{31} & a_{32} & a_{33}
\end{pmatrix}
\]

\section{Tables}
Below is a sample table showing data:

\begin{table}[h]
\centering
\begin{tabular}{@{}lrrr@{}}
\toprule
\textbf{Category} & \textbf{Value 1} & \textbf{Value 2} & \textbf{Total} \\
\midrule
Item A & 10.5 & 20.3 & 30.8 \\
Item B & 15.2 & 25.1 & 40.3 \\
Item C & 22.1 & 12.7 & 34.8 \\
\bottomrule
\end{tabular}
\caption{Sample Data Table}
\label{tab:sample}
\end{table}

\section{Lists}
\subsection{Bullet Points}
\begin{itemize}
  \item First item
  \item Second item
  \item Third item with sub-items:
    \begin{itemize}
      \item Sub-item 1
      \item Sub-item 2
    \end{itemize}
\end{itemize}

\subsection{Numbered List}
\begin{enumerate}
  \item First step
  \item Second step
  \item Third step
\end{enumerate}

\section{Citations and References}
When you use the citation search feature, references will be automatically managed. For example, you might cite a paper like this: \cite{example2024}. All citations will be automatically added to your \texttt{cite.bib} file.

\section{Conclusion}
This template demonstrates basic \LaTeX{} features including mathematical formulas, tables, and lists. The bibliography system is now set up to work seamlessly with Bibby AI's citation management. Feel free to modify it for your needs!

\bibliographystyle{plain}
\bibliography{cite}

\end{document}\documentclass{article}
\usepackage[utf8]{inputenc}
\usepackage{booktabs}
\usepackage{amsmath}
\usepackage{natbib}

\title{My Document}
\author{Your Name}
\date{\today}

\begin{document}

\maketitle

\section{Introduction}
This is a sample document demonstrating various \LaTeX{} features including mathematical equations, tables, and formatting. You can cite papers using the citation search feature, and they will be automatically added to your \texttt{cite.bib} file.

\section{Mathematical Expressions}
Here are some example equations:

\subsection{Quadratic Formula}
The quadratic formula is given by:
\[x = \frac{-b \pm \sqrt{b^2 - 4ac}}{2a}\]

\subsection{Matrix Example}
A sample matrix:
\[
\begin{pmatrix}
a_{11} & a_{12} & a_{13} \\
a_{21} & a_{22} & a_{23} \\
a_{31} & a_{32} & a_{33}
\end{pmatrix}
\]

\section{Tables}
Below is a sample table showing data:

\begin{table}[h]
\centering
\begin{tabular}{@{}lrrr@{}}
\toprule
\textbf{Category} & \textbf{Value 1} & \textbf{Value 2} & \textbf{Total} \\
\midrule
Item A & 10.5 & 20.3 & 30.8 \\
Item B & 15.2 & 25.1 & 40.3 \\
Item C & 22.1 & 12.7 & 34.8 \\
\bottomrule
\end{tabular}
\caption{Sample Data Table}
\label{tab:sample}
\end{table}

\section{Lists}
\subsection{Bullet Points}
\begin{itemize}
  \item First item
  \item Second item
  \item Third item with sub-items:
    \begin{itemize}
      \item Sub-item 1
      \item Sub-item 2
    \end{itemize}
\end{itemize}

\subsection{Numbered List}
\begin{enumerate}
  \item First step
  \item Second step
  \item Third step
\end{enumerate}

\section{Citations and References}
When you use the citation search feature, references will be automatically managed. For example, you might cite a paper like this: \cite{example2024}. All citations will be automatically added to your \texttt{cite.bib} file.

\section{Conclusion}
This template demonstrates basic \LaTeX{} features including mathematical formulas, tables, and lists. The bibliography system is now set up to work seamlessly with Bibby AI's citation management. Feel free to modify it for your needs!

\bibliographystyle{plain}
\bibliography{cite}

\end{document}\documentclass{article}
\usepackage[utf8]{inputenc}
\usepackage{booktabs}
\usepackage{amsmath}
\usepackage{natbib}

\title{My Document}
\author{Your Name}
\date{\today}

\begin{document}

\maketitle

\section{Introduction}
This is a sample document demonstrating various \LaTeX{} features including mathematical equations, tables, and formatting. You can cite papers using the citation search feature, and they will be automatically added to your \texttt{cite.bib} file.

\section{Mathematical Expressions}
Here are some example equations:

\subsection{Quadratic Formula}
The quadratic formula is given by:
\[x = \frac{-b \pm \sqrt{b^2 - 4ac}}{2a}\]

\subsection{Matrix Example}
A sample matrix:
\[
\begin{pmatrix}
a_{11} & a_{12} & a_{13} \\
a_{21} & a_{22} & a_{23} \\
a_{31} & a_{32} & a_{33}
\end{pmatrix}
\]

\section{Tables}
Below is a sample table showing data:

\begin{table}[h]
\centering
\begin{tabular}{@{}lrrr@{}}
\toprule
\textbf{Category} & \textbf{Value 1} & \textbf{Value 2} & \textbf{Total} \\
\midrule
Item A & 10.5 & 20.3 & 30.8 \\
Item B & 15.2 & 25.1 & 40.3 \\
Item C & 22.1 & 12.7 & 34.8 \\
\bottomrule
\end{tabular}
\caption{Sample Data Table}
\label{tab:sample}
\end{table}

\section{Lists}
\subsection{Bullet Points}
\begin{itemize}
  \item First item
  \item Second item
  \item Third item with sub-items:
    \begin{itemize}
      \item Sub-item 1
      \item Sub-item 2
    \end{itemize}
\end{itemize}

\subsection{Numbered List}
\begin{enumerate}
  \item First step
  \item Second step
  \item Third step
\end{enumerate}

\section{Citations and References}
When you use the citation search feature, references will be automatically managed. For example, you might cite a paper like this: \cite{example2024}. All citations will be automatically added to your \texttt{cite.bib} file.

\section{Conclusion}
This template demonstrates basic \LaTeX{} features including mathematical formulas, tables, and lists. The bibliography system is now set up to work seamlessly with Bibby AI's citation management. Feel free to modify it for your needs!

\bibliographystyle{plain}
\bibliography{cite}

\end{document}\documentclass{article}
\usepackage[utf8]{inputenc}
\usepackage{booktabs}
\usepackage{amsmath}
\usepackage{natbib}

\title{My Document}
\author{Your Name}
\date{\today}

\begin{document}

\maketitle

\section{Introduction}
This is a sample document demonstrating various \LaTeX{} features including mathematical equations, tables, and formatting. You can cite papers using the citation search feature, and they will be automatically added to your \texttt{cite.bib} file.

\section{Mathematical Expressions}
Here are some example equations:

\subsection{Quadratic Formula}
The quadratic formula is given by:
\[x = \frac{-b \pm \sqrt{b^2 - 4ac}}{2a}\]

\subsection{Matrix Example}
A sample matrix:
\[
\begin{pmatrix}
a_{11} & a_{12} & a_{13} \\
a_{21} & a_{22} & a_{23} \\
a_{31} & a_{32} & a_{33}
\end{pmatrix}
\]

\section{Tables}
Below is a sample table showing data:

\begin{table}[h]
\centering
\begin{tabular}{@{}lrrr@{}}
\toprule
\textbf{Category} & \textbf{Value 1} & \textbf{Value 2} & \textbf{Total} \\
\midrule
Item A & 10.5 & 20.3 & 30.8 \\
Item B & 15.2 & 25.1 & 40.3 \\
Item C & 22.1 & 12.7 & 34.8 \\
\bottomrule
\end{tabular}
\caption{Sample Data Table}
\label{tab:sample}
\end{table}

\section{Lists}
\subsection{Bullet Points}
\begin{itemize}
  \item First item
  \item Second item
  \item Third item with sub-items:
    \begin{itemize}
      \item Sub-item 1
      \item Sub-item 2
    \end{itemize}
\end{itemize}

\subsection{Numbered List}
\begin{enumerate}
  \item First step
  \item Second step
  \item Third step
\end{enumerate}

\section{Citations and References}
When you use the citation search feature, references will be automatically managed. For example, you might cite a paper like this: \cite{example2024}. All citations will be automatically added to your \texttt{cite.bib} file.

\section{Conclusion}
This template demonstrates basic \LaTeX{} features including mathematical formulas, tables, and lists. The bibliography system is now set up to work seamlessly with Bibby AI's citation management. Feel free to modify it for your needs!

\bibliographystyle{plain}
\bibliography{cite}

\end{document}\documentclass{article}
\usepackage[utf8]{inputenc}
\usepackage{booktabs}
\usepackage{amsmath}
\usepackage{natbib}

\title{My Document}
\author{Your Name}
\date{\today}

\begin{document}

\maketitle

\section{Introduction}
This is a sample document demonstrating various \LaTeX{} features including mathematical equations, tables, and formatting. You can cite papers using the citation search feature, and they will be automatically added to your \texttt{cite.bib} file.

\section{Mathematical Expressions}
Here are some example equations:

\subsection{Quadratic Formula}
The quadratic formula is given by:
\[x = \frac{-b \pm \sqrt{b^2 - 4ac}}{2a}\]

\subsection{Matrix Example}
A sample matrix:
\[
\begin{pmatrix}
a_{11} & a_{12} & a_{13} \\
a_{21} & a_{22} & a_{23} \\
a_{31} & a_{32} & a_{33}
\end{pmatrix}
\]

\section{Tables}
Below is a sample table showing data:

\begin{table}[h]
\centering
\begin{tabular}{@{}lrrr@{}}
\toprule
\textbf{Category} & \textbf{Value 1} & \textbf{Value 2} & \textbf{Total} \\
\midrule
Item A & 10.5 & 20.3 & 30.8 \\
Item B & 15.2 & 25.1 & 40.3 \\
Item C & 22.1 & 12.7 & 34.8 \\
\bottomrule
\end{tabular}
\caption{Sample Data Table}
\label{tab:sample}
\end{table}

\section{Lists}
\subsection{Bullet Points}
\begin{itemize}
  \item First item
  \item Second item
  \item Third item with sub-items:
    \begin{itemize}
      \item Sub-item 1
      \item Sub-item 2
    \end{itemize}
\end{itemize}

\subsection{Numbered List}
\begin{enumerate}
  \item First step
  \item Second step
  \item Third step
\end{enumerate}

\section{Citations and References}
When you use the citation search feature, references will be automatically managed. For example, you might cite a paper like this: \cite{example2024}. All citations will be automatically added to your \texttt{cite.bib} file.

\section{Conclusion}
This template demonstrates basic \LaTeX{} features including mathematical formulas, tables, and lists. The bibliography system is now set up to work seamlessly with Bibby AI's citation management. Feel free to modify it for your needs!

\bibliographystyle{plain}
\bibliography{cite}

\end{document}\documentclass{article}
\usepackage[utf8]{inputenc}
\usepackage{booktabs}
\usepackage{amsmath}
\usepackage{natbib}

\title{My Document}
\author{Your Name}
\date{\today}

\begin{document}

\maketitle

\section{Introduction}
This is a sample document demonstrating various \LaTeX{} features including mathematical equations, tables, and formatting. You can cite papers using the citation search feature, and they will be automatically added to your \texttt{cite.bib} file.

\section{Mathematical Expressions}
Here are some example equations:

\subsection{Quadratic Formula}
The quadratic formula is given by:
\[x = \frac{-b \pm \sqrt{b^2 - 4ac}}{2a}\]

\subsection{Matrix Example}
A sample matrix:
\[
\begin{pmatrix}
a_{11} & a_{12} & a_{13} \\
a_{21} & a_{22} & a_{23} \\
a_{31} & a_{32} & a_{33}
\end{pmatrix}
\]

\section{Tables}
Below is a sample table showing data:

\begin{table}[h]
\centering
\begin{tabular}{@{}lrrr@{}}
\toprule
\textbf{Category} & \textbf{Value 1} & \textbf{Value 2} & \textbf{Total} \\
\midrule
Item A & 10.5 & 20.3 & 30.8 \\
Item B & 15.2 & 25.1 & 40.3 \\
Item C & 22.1 & 12.7 & 34.8 \\
\bottomrule
\end{tabular}
\caption{Sample Data Table}
\label{tab:sample}
\end{table}

\section{Lists}
\subsection{Bullet Points}
\begin{itemize}
  \item First item
  \item Second item
  \item Third item with sub-items:
    \begin{itemize}
      \item Sub-item 1
      \item Sub-item 2
    \end{itemize}
\end{itemize}

\subsection{Numbered List}
\begin{enumerate}
  \item First step
  \item Second step
  \item Third step
\end{enumerate}

\section{Citations and References}
When you use the citation search feature, references will be automatically managed. For example, you might cite a paper like this: \cite{example2024}. All citations will be automatically added to your \texttt{cite.bib} file.

\section{Conclusion}
This template demonstrates basic \LaTeX{} features including mathematical formulas, tables, and lists. The bibliography system is now set up to work seamlessly with Bibby AI's citation management. Feel free to modify it for your needs!

\bibliographystyle{plain}
\bibliography{cite}

\end{document}\documentclass{article}
\usepackage[utf8]{inputenc}
\usepackage{booktabs}
\usepackage{amsmath}
\usepackage{natbib}

\title{My Document}
\author{Your Name}
\date{\today}

\begin{document}

\maketitle

\section{Introduction}
This is a sample document demonstrating various \LaTeX{} features including mathematical equations, tables, and formatting. You can cite papers using the citation search feature, and they will be automatically added to your \texttt{cite.bib} file.

\section{Mathematical Expressions}
Here are some example equations:

\subsection{Quadratic Formula}
The quadratic formula is given by:
\[x = \frac{-b \pm \sqrt{b^2 - 4ac}}{2a}\]

\subsection{Matrix Example}
A sample matrix:
\[
\begin{pmatrix}
a_{11} & a_{12} & a_{13} \\
a_{21} & a_{22} & a_{23} \\
a_{31} & a_{32} & a_{33}
\end{pmatrix}
\]

\section{Tables}
Below is a sample table showing data:

\begin{table}[h]
\centering
\begin{tabular}{@{}lrrr@{}}
\toprule
\textbf{Category} & \textbf{Value 1} & \textbf{Value 2} & \textbf{Total} \\
\midrule
Item A & 10.5 & 20.3 & 30.8 \\
Item B & 15.2 & 25.1 & 40.3 \\
Item C & 22.1 & 12.7 & 34.8 \\
\bottomrule
\end{tabular}
\caption{Sample Data Table}
\label{tab:sample}
\end{table}

\section{Lists}
\subsection{Bullet Points}
\begin{itemize}
  \item First item
  \item Second item
  \item Third item with sub-items:
    \begin{itemize}
      \item Sub-item 1
      \item Sub-item 2
    \end{itemize}
\end{itemize}

\subsection{Numbered List}
\begin{enumerate}
  \item First step
  \item Second step
  \item Third step
\end{enumerate}

\section{Citations and References}
When you use the citation search feature, references will be automatically managed. For example, you might cite a paper like this: \cite{example2024}. All citations will be automatically added to your \texttt{cite.bib} file.

\section{Conclusion}
This template demonstrates basic \LaTeX{} features including mathematical formulas, tables, and lists. The bibliography system is now set up to work seamlessly with Bibby AI's citation management. Feel free to modify it for your needs!

\bibliographystyle{plain}
\bibliography{cite}

\end{document}\documentclass{article}
\usepackage[utf8]{inputenc}
\usepackage{booktabs}
\usepackage{amsmath}
\usepackage{natbib}

\title{My Document}
\author{Your Name}
\date{\today}

\begin{document}

\maketitle

\section{Introduction}
This is a sample document demonstrating various \LaTeX{} features including mathematical equations, tables, and formatting. You can cite papers using the citation search feature, and they will be automatically added to your \texttt{cite.bib} file.

\section{Mathematical Expressions}
Here are some example equations:

\subsection{Quadratic Formula}
The quadratic formula is given by:
\[x = \frac{-b \pm \sqrt{b^2 - 4ac}}{2a}\]

\subsection{Matrix Example}
A sample matrix:
\[
\begin{pmatrix}
a_{11} & a_{12} & a_{13} \\
a_{21} & a_{22} & a_{23} \\
a_{31} & a_{32} & a_{33}
\end{pmatrix}
\]

\section{Tables}
Below is a sample table showing data:

\begin{table}[h]
\centering
\begin{tabular}{@{}lrrr@{}}
\toprule
\textbf{Category} & \textbf{Value 1} & \textbf{Value 2} & \textbf{Total} \\
\midrule
Item A & 10.5 & 20.3 & 30.8 \\
Item B & 15.2 & 25.1 & 40.3 \\
Item C & 22.1 & 12.7 & 34.8 \\
\bottomrule
\end{tabular}
\caption{Sample Data Table}
\label{tab:sample}
\end{table}

\section{Lists}
\subsection{Bullet Points}
\begin{itemize}
  \item First item
  \item Second item
  \item Third item with sub-items:
    \begin{itemize}
      \item Sub-item 1
      \item Sub-item 2
    \end{itemize}
\end{itemize}

\subsection{Numbered List}
\begin{enumerate}
  \item First step
  \item Second step
  \item Third step
\end{enumerate}

\section{Citations and References}
When you use the citation search feature, references will be automatically managed. For example, you might cite a paper like this: \cite{example2024}. All citations will be automatically added to your \texttt{cite.bib} file.

\section{Conclusion}
This template demonstrates basic \LaTeX{} features including mathematical formulas, tables, and lists. The bibliography system is now set up to work seamlessly with Bibby AI's citation management. Feel free to modify it for your needs!

\bibliographystyle{plain}
\bibliography{cite}

\end{document}\documentclass{article}
\usepackage[utf8]{inputenc}
\usepackage{booktabs}
\usepackage{amsmath}
\usepackage{natbib}

\title{My Document}
\author{Your Name}
\date{\today}

\begin{document}

\maketitle

\section{Introduction}
This is a sample document demonstrating various \LaTeX{} features including mathematical equations, tables, and formatting. You can cite papers using the citation search feature, and they will be automatically added to your \texttt{cite.bib} file.

\section{Mathematical Expressions}
Here are some example equations:

\subsection{Quadratic Formula}
The quadratic formula is given by:
\[x = \frac{-b \pm \sqrt{b^2 - 4ac}}{2a}\]

\subsection{Matrix Example}
A sample matrix:
\[
\begin{pmatrix}
a_{11} & a_{12} & a_{13} \\
a_{21} & a_{22} & a_{23} \\
a_{31} & a_{32} & a_{33}
\end{pmatrix}
\]

\section{Tables}
Below is a sample table showing data:

\begin{table}[h]
\centering
\begin{tabular}{@{}lrrr@{}}
\toprule
\textbf{Category} & \textbf{Value 1} & \textbf{Value 2} & \textbf{Total} \\
\midrule
Item A & 10.5 & 20.3 & 30.8 \\
Item B & 15.2 & 25.1 & 40.3 \\
Item C & 22.1 & 12.7 & 34.8 \\
\bottomrule
\end{tabular}
\caption{Sample Data Table}
\label{tab:sample}
\end{table}

\section{Lists}
\subsection{Bullet Points}
\begin{itemize}
  \item First item
  \item Second item
  \item Third item with sub-items:
    \begin{itemize}
      \item Sub-item 1
      \item Sub-item 2
    \end{itemize}
\end{itemize}

\subsection{Numbered List}
\begin{enumerate}
  \item First step
  \item Second step
  \item Third step
\end{enumerate}

\section{Citations and References}
When you use the citation search feature, references will be automatically managed. For example, you might cite a paper like this: \cite{example2024}. All citations will be automatically added to your \texttt{cite.bib} file.

\section{Conclusion}
This template demonstrates basic \LaTeX{} features including mathematical formulas, tables, and lists. The bibliography system is now set up to work seamlessly with Bibby AI's citation management. Feel free to modify it for your needs!

\bibliographystyle{plain}
\bibliography{cite}

\end{document}\documentclass{article}
\usepackage[utf8]{inputenc}
\usepackage{booktabs}
\usepackage{amsmath}
\usepackage{natbib}

\title{My Document}
\author{Your Name}
\date{\today}

\begin{document}

\maketitle

\section{Introduction}
This is a sample document demonstrating various \LaTeX{} features including mathematical equations, tables, and formatting. You can cite papers using the citation search feature, and they will be automatically added to your \texttt{cite.bib} file.

\section{Mathematical Expressions}
Here are some example equations:

\subsection{Quadratic Formula}
The quadratic formula is given by:
\[x = \frac{-b \pm \sqrt{b^2 - 4ac}}{2a}\]

\subsection{Matrix Example}
A sample matrix:
\[
\begin{pmatrix}
a_{11} & a_{12} & a_{13} \\
a_{21} & a_{22} & a_{23} \\
a_{31} & a_{32} & a_{33}
\end{pmatrix}
\]

\section{Tables}
Below is a sample table showing data:

\begin{table}[h]
\centering
\begin{tabular}{@{}lrrr@{}}
\toprule
\textbf{Category} & \textbf{Value 1} & \textbf{Value 2} & \textbf{Total} \\
\midrule
Item A & 10.5 & 20.3 & 30.8 \\
Item B & 15.2 & 25.1 & 40.3 \\
Item C & 22.1 & 12.7 & 34.8 \\
\bottomrule
\end{tabular}
\caption{Sample Data Table}
\label{tab:sample}
\end{table}

\section{Lists}
\subsection{Bullet Points}
\begin{itemize}
  \item First item
  \item Second item
  \item Third item with sub-items:
    \begin{itemize}
      \item Sub-item 1
      \item Sub-item 2
    \end{itemize}
\end{itemize}

\subsection{Numbered List}
\begin{enumerate}
  \item First step
  \item Second step
  \item Third step
\end{enumerate}

\section{Citations and References}
When you use the citation search feature, references will be automatically managed. For example, you might cite a paper like this: \cite{example2024}. All citations will be automatically added to your \texttt{cite.bib} file.

\section{Conclusion}
This template demonstrates basic \LaTeX{} features including mathematical formulas, tables, and lists. The bibliography system is now set up to work seamlessly with Bibby AI's citation management. Feel free to modify it for your needs!

\bibliographystyle{plain}
\bibliography{cite}

\end{document}\documentclass{article}
\usepackage[utf8]{inputenc}
\usepackage{booktabs}
\usepackage{amsmath}
\usepackage{natbib}

\title{My Document}
\author{Your Name}
\date{\today}

\begin{document}

\maketitle

\section{Introduction}
This is a sample document demonstrating various \LaTeX{} features including mathematical equations, tables, and formatting. You can cite papers using the citation search feature, and they will be automatically added to your \texttt{cite.bib} file.

\section{Mathematical Expressions}
Here are some example equations:

\subsection{Quadratic Formula}
The quadratic formula is given by:
\[x = \frac{-b \pm \sqrt{b^2 - 4ac}}{2a}\]

\subsection{Matrix Example}
A sample matrix:
\[
\begin{pmatrix}
a_{11} & a_{12} & a_{13} \\
a_{21} & a_{22} & a_{23} \\
a_{31} & a_{32} & a_{33}
\end{pmatrix}
\]

\section{Tables}
Below is a sample table showing data:

\begin{table}[h]
\centering
\begin{tabular}{@{}lrrr@{}}
\toprule
\textbf{Category} & \textbf{Value 1} & \textbf{Value 2} & \textbf{Total} \\
\midrule
Item A & 10.5 & 20.3 & 30.8 \\
Item B & 15.2 & 25.1 & 40.3 \\
Item C & 22.1 & 12.7 & 34.8 \\
\bottomrule
\end{tabular}
\caption{Sample Data Table}
\label{tab:sample}
\end{table}

\section{Lists}
\subsection{Bullet Points}
\begin{itemize}
  \item First item
  \item Second item
  \item Third item with sub-items:
    \begin{itemize}
      \item Sub-item 1
      \item Sub-item 2
    \end{itemize}
\end{itemize}

\subsection{Numbered List}
\begin{enumerate}
  \item First step
  \item Second step
  \item Third step
\end{enumerate}

\section{Citations and References}
When you use the citation search feature, references will be automatically managed. For example, you might cite a paper like this: \cite{example2024}. All citations will be automatically added to your \texttt{cite.bib} file.

\section{Conclusion}
This template demonstrates basic \LaTeX{} features including mathematical formulas, tables, and lists. The bibliography system is now set up to work seamlessly with Bibby AI's citation management. Feel free to modify it for your needs!

\bibliographystyle{plain}
\bibliography{cite}

\end{document}\documentclass{article}
\usepackage[utf8]{inputenc}
\usepackage{booktabs}
\usepackage{amsmath}
\usepackage{natbib}

\title{My Document}
\author{Your Name}
\date{\today}

\begin{document}

\maketitle

\section{Introduction}
This is a sample document demonstrating various \LaTeX{} features including mathematical equations, tables, and formatting. You can cite papers using the citation search feature, and they will be automatically added to your \texttt{cite.bib} file.

\section{Mathematical Expressions}
Here are some example equations:

\subsection{Quadratic Formula}
The quadratic formula is given by:
\[x = \frac{-b \pm \sqrt{b^2 - 4ac}}{2a}\]

\subsection{Matrix Example}
A sample matrix:
\[
\begin{pmatrix}
a_{11} & a_{12} & a_{13} \\
a_{21} & a_{22} & a_{23} \\
a_{31} & a_{32} & a_{33}
\end{pmatrix}
\]

\section{Tables}
Below is a sample table showing data:

\begin{table}[h]
\centering
\begin{tabular}{@{}lrrr@{}}
\toprule
\textbf{Category} & \textbf{Value 1} & \textbf{Value 2} & \textbf{Total} \\
\midrule
Item A & 10.5 & 20.3 & 30.8 \\
Item B & 15.2 & 25.1 & 40.3 \\
Item C & 22.1 & 12.7 & 34.8 \\
\bottomrule
\end{tabular}
\caption{Sample Data Table}
\label{tab:sample}
\end{table}

\section{Lists}
\subsection{Bullet Points}
\begin{itemize}
  \item First item
  \item Second item
  \item Third item with sub-items:
    \begin{itemize}
      \item Sub-item 1
      \item Sub-item 2
    \end{itemize}
\end{itemize}

\subsection{Numbered List}
\begin{enumerate}
  \item First step
  \item Second step
  \item Third step
\end{enumerate}

\section{Citations and References}
When you use the citation search feature, references will be automatically managed. For example, you might cite a paper like this: \cite{example2024}. All citations will be automatically added to your \texttt{cite.bib} file.

\section{Conclusion}
This template demonstrates basic \LaTeX{} features including mathematical formulas, tables, and lists. The bibliography system is now set up to work seamlessly with Bibby AI's citation management. Feel free to modify it for your needs!

\bibliographystyle{plain}
\bibliography{cite}

\end{document}\documentclass{article}
\usepackage[utf8]{inputenc}
\usepackage{booktabs}
\usepackage{amsmath}
\usepackage{natbib}

\title{My Document}
\author{Your Name}
\date{\today}

\begin{document}

\maketitle

\section{Introduction}
This is a sample document demonstrating various \LaTeX{} features including mathematical equations, tables, and formatting. You can cite papers using the citation search feature, and they will be automatically added to your \texttt{cite.bib} file.

\section{Mathematical Expressions}
Here are some example equations:

\subsection{Quadratic Formula}
The quadratic formula is given by:
\[x = \frac{-b \pm \sqrt{b^2 - 4ac}}{2a}\]

\subsection{Matrix Example}
A sample matrix:
\[
\begin{pmatrix}
a_{11} & a_{12} & a_{13} \\
a_{21} & a_{22} & a_{23} \\
a_{31} & a_{32} & a_{33}
\end{pmatrix}
\]

\section{Tables}
Below is a sample table showing data:

\begin{table}[h]
\centering
\begin{tabular}{@{}lrrr@{}}
\toprule
\textbf{Category} & \textbf{Value 1} & \textbf{Value 2} & \textbf{Total} \\
\midrule
Item A & 10.5 & 20.3 & 30.8 \\
Item B & 15.2 & 25.1 & 40.3 \\
Item C & 22.1 & 12.7 & 34.8 \\
\bottomrule
\end{tabular}
\caption{Sample Data Table}
\label{tab:sample}
\end{table}

\section{Lists}
\subsection{Bullet Points}
\begin{itemize}
  \item First item
  \item Second item
  \item Third item with sub-items:
    \begin{itemize}
      \item Sub-item 1
      \item Sub-item 2
    \end{itemize}
\end{itemize}

\subsection{Numbered List}
\begin{enumerate}
  \item First step
  \item Second step
  \item Third step
\end{enumerate}

\section{Citations and References}
When you use the citation search feature, references will be automatically managed. For example, you might cite a paper like this: \cite{example2024}. All citations will be automatically added to your \texttt{cite.bib} file.

\section{Conclusion}
This template demonstrates basic \LaTeX{} features including mathematical formulas, tables, and lists. The bibliography system is now set up to work seamlessly with Bibby AI's citation management. Feel free to modify it for your needs!

\bibliographystyle{plain}
\bibliography{cite}

\end{document}\documentclass{article}
\usepackage[utf8]{inputenc}
\usepackage{booktabs}
\usepackage{amsmath}
\usepackage{natbib}

\title{My Document}
\author{Your Name}
\date{\today}

\begin{document}

\maketitle

\section{Introduction}
This is a sample document demonstrating various \LaTeX{} features including mathematical equations, tables, and formatting. You can cite papers using the citation search feature, and they will be automatically added to your \texttt{cite.bib} file.

\section{Mathematical Expressions}
Here are some example equations:

\subsection{Quadratic Formula}
The quadratic formula is given by:
\[x = \frac{-b \pm \sqrt{b^2 - 4ac}}{2a}\]

\subsection{Matrix Example}
A sample matrix:
\[
\begin{pmatrix}
a_{11} & a_{12} & a_{13} \\
a_{21} & a_{22} & a_{23} \\
a_{31} & a_{32} & a_{33}
\end{pmatrix}
\]

\section{Tables}
Below is a sample table showing data:

\begin{table}[h]
\centering
\begin{tabular}{@{}lrrr@{}}
\toprule
\textbf{Category} & \textbf{Value 1} & \textbf{Value 2} & \textbf{Total} \\
\midrule
Item A & 10.5 & 20.3 & 30.8 \\
Item B & 15.2 & 25.1 & 40.3 \\
Item C & 22.1 & 12.7 & 34.8 \\
\bottomrule
\end{tabular}
\caption{Sample Data Table}
\label{tab:sample}
\end{table}

\section{Lists}
\subsection{Bullet Points}
\begin{itemize}
  \item First item
  \item Second item
  \item Third item with sub-items:
    \begin{itemize}
      \item Sub-item 1
      \item Sub-item 2
    \end{itemize}
\end{itemize}

\subsection{Numbered List}
\begin{enumerate}
  \item First step
  \item Second step
  \item Third step
\end{enumerate}

\section{Citations and References}
When you use the citation search feature, references will be automatically managed. For example, you might cite a paper like this: \cite{example2024}. All citations will be automatically added to your \texttt{cite.bib} file.

\section{Conclusion}
This template demonstrates basic \LaTeX{} features including mathematical formulas, tables, and lists. The bibliography system is now set up to work seamlessly with Bibby AI's citation management. Feel free to modify it for your needs!

\bibliographystyle{plain}
\bibliography{cite}

\end{document}\documentclass{article}
\usepackage[utf8]{inputenc}
\usepackage{booktabs}
\usepackage{amsmath}
\usepackage{natbib}

\title{My Document}
\author{Your Name}
\date{\today}

\begin{document}

\maketitle

\section{Introduction}
This is a sample document demonstrating various \LaTeX{} features including mathematical equations, tables, and formatting. You can cite papers using the citation search feature, and they will be automatically added to your \texttt{cite.bib} file.

\section{Mathematical Expressions}
Here are some example equations:

\subsection{Quadratic Formula}
The quadratic formula is given by:
\[x = \frac{-b \pm \sqrt{b^2 - 4ac}}{2a}\]

\subsection{Matrix Example}
A sample matrix:
\[
\begin{pmatrix}
a_{11} & a_{12} & a_{13} \\
a_{21} & a_{22} & a_{23} \\
a_{31} & a_{32} & a_{33}
\end{pmatrix}
\]

\section{Tables}
Below is a sample table showing data:

\begin{table}[h]
\centering
\begin{tabular}{@{}lrrr@{}}
\toprule
\textbf{Category} & \textbf{Value 1} & \textbf{Value 2} & \textbf{Total} \\
\midrule
Item A & 10.5 & 20.3 & 30.8 \\
Item B & 15.2 & 25.1 & 40.3 \\
Item C & 22.1 & 12.7 & 34.8 \\
\bottomrule
\end{tabular}
\caption{Sample Data Table}
\label{tab:sample}
\end{table}

\section{Lists}
\subsection{Bullet Points}
\begin{itemize}
  \item First item
  \item Second item
  \item Third item with sub-items:
    \begin{itemize}
      \item Sub-item 1
      \item Sub-item 2
    \end{itemize}
\end{itemize}

\subsection{Numbered List}
\begin{enumerate}
  \item First step
  \item Second step
  \item Third step
\end{enumerate}

\section{Citations and References}
When you use the citation search feature, references will be automatically managed. For example, you might cite a paper like this: \cite{example2024}. All citations will be automatically added to your \texttt{cite.bib} file.

\section{Conclusion}
This template demonstrates basic \LaTeX{} features including mathematical formulas, tables, and lists. The bibliography system is now set up to work seamlessly with Bibby AI's citation management. Feel free to modify it for your needs!

\bibliographystyle{plain}
\bibliography{cite}

\end{document}\documentclass{article}
\usepackage[utf8]{inputenc}
\usepackage{booktabs}
\usepackage{amsmath}
\usepackage{natbib}

\title{My Document}
\author{Your Name}
\date{\today}

\begin{document}

\maketitle

\section{Introduction}
This is a sample document demonstrating various \LaTeX{} features including mathematical equations, tables, and formatting. You can cite papers using the citation search feature, and they will be automatically added to your \texttt{cite.bib} file.

\section{Mathematical Expressions}
Here are some example equations:

\subsection{Quadratic Formula}
The quadratic formula is given by:
\[x = \frac{-b \pm \sqrt{b^2 - 4ac}}{2a}\]

\subsection{Matrix Example}
A sample matrix:
\[
\begin{pmatrix}
a_{11} & a_{12} & a_{13} \\
a_{21} & a_{22} & a_{23} \\
a_{31} & a_{32} & a_{33}
\end{pmatrix}
\]

\section{Tables}
Below is a sample table showing data:

\begin{table}[h]
\centering
\begin{tabular}{@{}lrrr@{}}
\toprule
\textbf{Category} & \textbf{Value 1} & \textbf{Value 2} & \textbf{Total} \\
\midrule
Item A & 10.5 & 20.3 & 30.8 \\
Item B & 15.2 & 25.1 & 40.3 \\
Item C & 22.1 & 12.7 & 34.8 \\
\bottomrule
\end{tabular}
\caption{Sample Data Table}
\label{tab:sample}
\end{table}

\section{Lists}
\subsection{Bullet Points}
\begin{itemize}
  \item First item
  \item Second item
  \item Third item with sub-items:
    \begin{itemize}
      \item Sub-item 1
      \item Sub-item 2
    \end{itemize}
\end{itemize}

\subsection{Numbered List}
\begin{enumerate}
  \item First step
  \item Second step
  \item Third step
\end{enumerate}

\section{Citations and References}
When you use the citation search feature, references will be automatically managed. For example, you might cite a paper like this: \cite{example2024}. All citations will be automatically added to your \texttt{cite.bib} file.

\section{Conclusion}
This template demonstrates basic \LaTeX{} features including mathematical formulas, tables, and lists. The bibliography system is now set up to work seamlessly with Bibby AI's citation management. Feel free to modify it for your needs!

\bibliographystyle{plain}
\bibliography{cite}

\end{document}\documentclass{article}
\usepackage[utf8]{inputenc}
\usepackage{booktabs}
\usepackage{amsmath}
\usepackage{natbib}

\title{My Document}
\author{Your Name}
\date{\today}

\begin{document}

\maketitle

\section{Introduction}
This is a sample document demonstrating various \LaTeX{} features including mathematical equations, tables, and formatting. You can cite papers using the citation search feature, and they will be automatically added to your \texttt{cite.bib} file.

\section{Mathematical Expressions}
Here are some example equations:

\subsection{Quadratic Formula}
The quadratic formula is given by:
\[x = \frac{-b \pm \sqrt{b^2 - 4ac}}{2a}\]

\subsection{Matrix Example}
A sample matrix:
\[
\begin{pmatrix}
a_{11} & a_{12} & a_{13} \\
a_{21} & a_{22} & a_{23} \\
a_{31} & a_{32} & a_{33}
\end{pmatrix}
\]

\section{Tables}
Below is a sample table showing data:

\begin{table}[h]
\centering
\begin{tabular}{@{}lrrr@{}}
\toprule
\textbf{Category} & \textbf{Value 1} & \textbf{Value 2} & \textbf{Total} \\
\midrule
Item A & 10.5 & 20.3 & 30.8 \\
Item B & 15.2 & 25.1 & 40.3 \\
Item C & 22.1 & 12.7 & 34.8 \\
\bottomrule
\end{tabular}
\caption{Sample Data Table}
\label{tab:sample}
\end{table}

\section{Lists}
\subsection{Bullet Points}
\begin{itemize}
  \item First item
  \item Second item
  \item Third item with sub-items:
    \begin{itemize}
      \item Sub-item 1
      \item Sub-item 2
    \end{itemize}
\end{itemize}

\subsection{Numbered List}
\begin{enumerate}
  \item First step
  \item Second step
  \item Third step
\end{enumerate}

\section{Citations and References}
When you use the citation search feature, references will be automatically managed. For example, you might cite a paper like this: \cite{example2024}. All citations will be automatically added to your \texttt{cite.bib} file.

\section{Conclusion}
This template demonstrates basic \LaTeX{} features including mathematical formulas, tables, and lists. The bibliography system is now set up to work seamlessly with Bibby AI's citation management. Feel free to modify it for your needs!

\bibliographystyle{plain}
\bibliography{cite}

\end{document}\documentclass{article}
\usepackage[utf8]{inputenc}
\usepackage{booktabs}
\usepackage{amsmath}
\usepackage{natbib}

\title{My Document}
\author{Your Name}
\date{\today}

\begin{document}

\maketitle

\section{Introduction}
This is a sample document demonstrating various \LaTeX{} features including mathematical equations, tables, and formatting. You can cite papers using the citation search feature, and they will be automatically added to your \texttt{cite.bib} file.

\section{Mathematical Expressions}
Here are some example equations:

\subsection{Quadratic Formula}
The quadratic formula is given by:
\[x = \frac{-b \pm \sqrt{b^2 - 4ac}}{2a}\]

\subsection{Matrix Example}
A sample matrix:
\[
\begin{pmatrix}
a_{11} & a_{12} & a_{13} \\
a_{21} & a_{22} & a_{23} \\
a_{31} & a_{32} & a_{33}
\end{pmatrix}
\]

\section{Tables}
Below is a sample table showing data:

\begin{table}[h]
\centering
\begin{tabular}{@{}lrrr@{}}
\toprule
\textbf{Category} & \textbf{Value 1} & \textbf{Value 2} & \textbf{Total} \\
\midrule
Item A & 10.5 & 20.3 & 30.8 \\
Item B & 15.2 & 25.1 & 40.3 \\
Item C & 22.1 & 12.7 & 34.8 \\
\bottomrule
\end{tabular}
\caption{Sample Data Table}
\label{tab:sample}
\end{table}

\section{Lists}
\subsection{Bullet Points}
\begin{itemize}
  \item First item
  \item Second item
  \item Third item with sub-items:
    \begin{itemize}
      \item Sub-item 1
      \item Sub-item 2
    \end{itemize}
\end{itemize}

\subsection{Numbered List}
\begin{enumerate}
  \item First step
  \item Second step
  \item Third step
\end{enumerate}

\section{Citations and References}
When you use the citation search feature, references will be automatically managed. For example, you might cite a paper like this: \cite{example2024}. All citations will be automatically added to your \texttt{cite.bib} file.

\section{Conclusion}
This template demonstrates basic \LaTeX{} features including mathematical formulas, tables, and lists. The bibliography system is now set up to work seamlessly with Bibby AI's citation management. Feel free to modify it for your needs!

\bibliographystyle{plain}
\bibliography{cite}

\end{document}\documentclass{article}
\usepackage[utf8]{inputenc}
\usepackage{booktabs}
\usepackage{amsmath}
\usepackage{natbib}

\title{My Document}
\author{Your Name}
\date{\today}

\begin{document}

\maketitle

\section{Introduction}
This is a sample document demonstrating various \LaTeX{} features including mathematical equations, tables, and formatting. You can cite papers using the citation search feature, and they will be automatically added to your \texttt{cite.bib} file.

\section{Mathematical Expressions}
Here are some example equations:

\subsection{Quadratic Formula}
The quadratic formula is given by:
\[x = \frac{-b \pm \sqrt{b^2 - 4ac}}{2a}\]

\subsection{Matrix Example}
A sample matrix:
\[
\begin{pmatrix}
a_{11} & a_{12} & a_{13} \\
a_{21} & a_{22} & a_{23} \\
a_{31} & a_{32} & a_{33}
\end{pmatrix}
\]

\section{Tables}
Below is a sample table showing data:

\begin{table}[h]
\centering
\begin{tabular}{@{}lrrr@{}}
\toprule
\textbf{Category} & \textbf{Value 1} & \textbf{Value 2} & \textbf{Total} \\
\midrule
Item A & 10.5 & 20.3 & 30.8 \\
Item B & 15.2 & 25.1 & 40.3 \\
Item C & 22.1 & 12.7 & 34.8 \\
\bottomrule
\end{tabular}
\caption{Sample Data Table}
\label{tab:sample}
\end{table}

\section{Lists}
\subsection{Bullet Points}
\begin{itemize}
  \item First item
  \item Second item
  \item Third item with sub-items:
    \begin{itemize}
      \item Sub-item 1
      \item Sub-item 2
    \end{itemize}
\end{itemize}

\subsection{Numbered List}
\begin{enumerate}
  \item First step
  \item Second step
  \item Third step
\end{enumerate}

\section{Citations and References}
When you use the citation search feature, references will be automatically managed. For example, you might cite a paper like this: \cite{example2024}. All citations will be automatically added to your \texttt{cite.bib} file.

\section{Conclusion}
This template demonstrates basic \LaTeX{} features including mathematical formulas, tables, and lists. The bibliography system is now set up to work seamlessly with Bibby AI's citation management. Feel free to modify it for your needs!

\bibliographystyle{plain}
\bibliography{cite}

\end{document}\documentclass{article}
\usepackage[utf8]{inputenc}
\usepackage{booktabs}
\usepackage{amsmath}
\usepackage{natbib}

\title{My Document}
\author{Your Name}
\date{\today}

\begin{document}

\maketitle

\section{Introduction}
This is a sample document demonstrating various \LaTeX{} features including mathematical equations, tables, and formatting. You can cite papers using the citation search feature, and they will be automatically added to your \texttt{cite.bib} file.

\section{Mathematical Expressions}
Here are some example equations:

\subsection{Quadratic Formula}
The quadratic formula is given by:
\[x = \frac{-b \pm \sqrt{b^2 - 4ac}}{2a}\]

\subsection{Matrix Example}
A sample matrix:
\[
\begin{pmatrix}
a_{11} & a_{12} & a_{13} \\
a_{21} & a_{22} & a_{23} \\
a_{31} & a_{32} & a_{33}
\end{pmatrix}
\]

\section{Tables}
Below is a sample table showing data:

\begin{table}[h]
\centering
\begin{tabular}{@{}lrrr@{}}
\toprule
\textbf{Category} & \textbf{Value 1} & \textbf{Value 2} & \textbf{Total} \\
\midrule
Item A & 10.5 & 20.3 & 30.8 \\
Item B & 15.2 & 25.1 & 40.3 \\
Item C & 22.1 & 12.7 & 34.8 \\
\bottomrule
\end{tabular}
\caption{Sample Data Table}
\label{tab:sample}
\end{table}

\section{Lists}
\subsection{Bullet Points}
\begin{itemize}
  \item First item
  \item Second item
  \item Third item with sub-items:
    \begin{itemize}
      \item Sub-item 1
      \item Sub-item 2
    \end{itemize}
\end{itemize}

\subsection{Numbered List}
\begin{enumerate}
  \item First step
  \item Second step
  \item Third step
\end{enumerate}

\section{Citations and References}
When you use the citation search feature, references will be automatically managed. For example, you might cite a paper like this: \cite{example2024}. All citations will be automatically added to your \texttt{cite.bib} file.

\section{Conclusion}
This template demonstrates basic \LaTeX{} features including mathematical formulas, tables, and lists. The bibliography system is now set up to work seamlessly with Bibby AI's citation management. Feel free to modify it for your needs!

\bibliographystyle{plain}
\bibliography{cite}

\end{document}\documentclass{article}
\usepackage[utf8]{inputenc}
\usepackage{booktabs}
\usepackage{amsmath}
\usepackage{natbib}

\title{My Document}
\author{Your Name}
\date{\today}

\begin{document}

\maketitle

\section{Introduction}
This is a sample document demonstrating various \LaTeX{} features including mathematical equations, tables, and formatting. You can cite papers using the citation search feature, and they will be automatically added to your \texttt{cite.bib} file.

\section{Mathematical Expressions}
Here are some example equations:

\subsection{Quadratic Formula}
The quadratic formula is given by:
\[x = \frac{-b \pm \sqrt{b^2 - 4ac}}{2a}\]

\subsection{Matrix Example}
A sample matrix:
\[
\begin{pmatrix}
a_{11} & a_{12} & a_{13} \\
a_{21} & a_{22} & a_{23} \\
a_{31} & a_{32} & a_{33}
\end{pmatrix}
\]

\section{Tables}
Below is a sample table showing data:

\begin{table}[h]
\centering
\begin{tabular}{@{}lrrr@{}}
\toprule
\textbf{Category} & \textbf{Value 1} & \textbf{Value 2} & \textbf{Total} \\
\midrule
Item A & 10.5 & 20.3 & 30.8 \\
Item B & 15.2 & 25.1 & 40.3 \\
Item C & 22.1 & 12.7 & 34.8 \\
\bottomrule
\end{tabular}
\caption{Sample Data Table}
\label{tab:sample}
\end{table}

\section{Lists}
\subsection{Bullet Points}
\begin{itemize}
  \item First item
  \item Second item
  \item Third item with sub-items:
    \begin{itemize}
      \item Sub-item 1
      \item Sub-item 2
    \end{itemize}
\end{itemize}

\subsection{Numbered List}
\begin{enumerate}
  \item First step
  \item Second step
  \item Third step
\end{enumerate}

\section{Citations and References}
When you use the citation search feature, references will be automatically managed. For example, you might cite a paper like this: \cite{example2024}. All citations will be automatically added to your \texttt{cite.bib} file.

\section{Conclusion}
This template demonstrates basic \LaTeX{} features including mathematical formulas, tables, and lists. The bibliography system is now set up to work seamlessly with Bibby AI's citation management. Feel free to modify it for your needs!

\bibliographystyle{plain}
\bibliography{cite}

\end{document}\documentclass{article}
\usepackage[utf8]{inputenc}
\usepackage{booktabs}
\usepackage{amsmath}
\usepackage{natbib}

\title{My Document}
\author{Your Name}
\date{\today}

\begin{document}

\maketitle

\section{Introduction}
This is a sample document demonstrating various \LaTeX{} features including mathematical equations, tables, and formatting. You can cite papers using the citation search feature, and they will be automatically added to your \texttt{cite.bib} file.

\section{Mathematical Expressions}
Here are some example equations:

\subsection{Quadratic Formula}
The quadratic formula is given by:
\[x = \frac{-b \pm \sqrt{b^2 - 4ac}}{2a}\]

\subsection{Matrix Example}
A sample matrix:
\[
\begin{pmatrix}
a_{11} & a_{12} & a_{13} \\
a_{21} & a_{22} & a_{23} \\
a_{31} & a_{32} & a_{33}
\end{pmatrix}
\]

\section{Tables}
Below is a sample table showing data:

\begin{table}[h]
\centering
\begin{tabular}{@{}lrrr@{}}
\toprule
\textbf{Category} & \textbf{Value 1} & \textbf{Value 2} & \textbf{Total} \\
\midrule
Item A & 10.5 & 20.3 & 30.8 \\
Item B & 15.2 & 25.1 & 40.3 \\
Item C & 22.1 & 12.7 & 34.8 \\
\bottomrule
\end{tabular}
\caption{Sample Data Table}
\label{tab:sample}
\end{table}

\section{Lists}
\subsection{Bullet Points}
\begin{itemize}
  \item First item
  \item Second item
  \item Third item with sub-items:
    \begin{itemize}
      \item Sub-item 1
      \item Sub-item 2
    \end{itemize}
\end{itemize}

\subsection{Numbered List}
\begin{enumerate}
  \item First step
  \item Second step
  \item Third step
\end{enumerate}

\section{Citations and References}
When you use the citation search feature, references will be automatically managed. For example, you might cite a paper like this: \cite{example2024}. All citations will be automatically added to your \texttt{cite.bib} file.

\section{Conclusion}
This template demonstrates basic \LaTeX{} features including mathematical formulas, tables, and lists. The bibliography system is now set up to work seamlessly with Bibby AI's citation management. Feel free to modify it for your needs!

\bibliographystyle{plain}
\bibliography{cite}

\end{document}\documentclass{article}
\usepackage[utf8]{inputenc}
\usepackage{booktabs}
\usepackage{amsmath}
\usepackage{natbib}

\title{My Document}
\author{Your Name}
\date{\today}

\begin{document}

\maketitle

\section{Introduction}
This is a sample document demonstrating various \LaTeX{} features including mathematical equations, tables, and formatting. You can cite papers using the citation search feature, and they will be automatically added to your \texttt{cite.bib} file.

\section{Mathematical Expressions}
Here are some example equations:

\subsection{Quadratic Formula}
The quadratic formula is given by:
\[x = \frac{-b \pm \sqrt{b^2 - 4ac}}{2a}\]

\subsection{Matrix Example}
A sample matrix:
\[
\begin{pmatrix}
a_{11} & a_{12} & a_{13} \\
a_{21} & a_{22} & a_{23} \\
a_{31} & a_{32} & a_{33}
\end{pmatrix}
\]

\section{Tables}
Below is a sample table showing data:

\begin{table}[h]
\centering
\begin{tabular}{@{}lrrr@{}}
\toprule
\textbf{Category} & \textbf{Value 1} & \textbf{Value 2} & \textbf{Total} \\
\midrule
Item A & 10.5 & 20.3 & 30.8 \\
Item B & 15.2 & 25.1 & 40.3 \\
Item C & 22.1 & 12.7 & 34.8 \\
\bottomrule
\end{tabular}
\caption{Sample Data Table}
\label{tab:sample}
\end{table}

\section{Lists}
\subsection{Bullet Points}
\begin{itemize}
  \item First item
  \item Second item
  \item Third item with sub-items:
    \begin{itemize}
      \item Sub-item 1
      \item Sub-item 2
    \end{itemize}
\end{itemize}

\subsection{Numbered List}
\begin{enumerate}
  \item First step
  \item Second step
  \item Third step
\end{enumerate}

\section{Citations and References}
When you use the citation search feature, references will be automatically managed. For example, you might cite a paper like this: \cite{example2024}. All citations will be automatically added to your \texttt{cite.bib} file.

\section{Conclusion}
This template demonstrates basic \LaTeX{} features including mathematical formulas, tables, and lists. The bibliography system is now set up to work seamlessly with Bibby AI's citation management. Feel free to modify it for your needs!

\bibliographystyle{plain}
\bibliography{cite}

\end{document}\documentclass{article}
\usepackage[utf8]{inputenc}
\usepackage{booktabs}
\usepackage{amsmath}
\usepackage{natbib}

\title{My Document}
\author{Your Name}
\date{\today}

\begin{document}

\maketitle

\section{Introduction}
This is a sample document demonstrating various \LaTeX{} features including mathematical equations, tables, and formatting. You can cite papers using the citation search feature, and they will be automatically added to your \texttt{cite.bib} file.

\section{Mathematical Expressions}
Here are some example equations:

\subsection{Quadratic Formula}
The quadratic formula is given by:
\[x = \frac{-b \pm \sqrt{b^2 - 4ac}}{2a}\]

\subsection{Matrix Example}
A sample matrix:
\[
\begin{pmatrix}
a_{11} & a_{12} & a_{13} \\
a_{21} & a_{22} & a_{23} \\
a_{31} & a_{32} & a_{33}
\end{pmatrix}
\]

\section{Tables}
Below is a sample table showing data:

\begin{table}[h]
\centering
\begin{tabular}{@{}lrrr@{}}
\toprule
\textbf{Category} & \textbf{Value 1} & \textbf{Value 2} & \textbf{Total} \\
\midrule
Item A & 10.5 & 20.3 & 30.8 \\
Item B & 15.2 & 25.1 & 40.3 \\
Item C & 22.1 & 12.7 & 34.8 \\
\bottomrule
\end{tabular}
\caption{Sample Data Table}
\label{tab:sample}
\end{table}

\section{Lists}
\subsection{Bullet Points}
\begin{itemize}
  \item First item
  \item Second item
  \item Third item with sub-items:
    \begin{itemize}
      \item Sub-item 1
      \item Sub-item 2
    \end{itemize}
\end{itemize}

\subsection{Numbered List}
\begin{enumerate}
  \item First step
  \item Second step
  \item Third step
\end{enumerate}

\section{Citations and References}
When you use the citation search feature, references will be automatically managed. For example, you might cite a paper like this: \cite{example2024}. All citations will be automatically added to your \texttt{cite.bib} file.

\section{Conclusion}
This template demonstrates basic \LaTeX{} features including mathematical formulas, tables, and lists. The bibliography system is now set up to work seamlessly with Bibby AI's citation management. Feel free to modify it for your needs!

\bibliographystyle{plain}
\bibliography{cite}

\end{document}\documentclass{article}
\usepackage[utf8]{inputenc}
\usepackage{booktabs}
\usepackage{amsmath}
\usepackage{natbib}

\title{My Document}
\author{Your Name}
\date{\today}

\begin{document}

\maketitle

\section{Introduction}
This is a sample document demonstrating various \LaTeX{} features including mathematical equations, tables, and formatting. You can cite papers using the citation search feature, and they will be automatically added to your \texttt{cite.bib} file.

\section{Mathematical Expressions}
Here are some example equations:

\subsection{Quadratic Formula}
The quadratic formula is given by:
\[x = \frac{-b \pm \sqrt{b^2 - 4ac}}{2a}\]

\subsection{Matrix Example}
A sample matrix:
\[
\begin{pmatrix}
a_{11} & a_{12} & a_{13} \\
a_{21} & a_{22} & a_{23} \\
a_{31} & a_{32} & a_{33}
\end{pmatrix}
\]

\section{Tables}
Below is a sample table showing data:

\begin{table}[h]
\centering
\begin{tabular}{@{}lrrr@{}}
\toprule
\textbf{Category} & \textbf{Value 1} & \textbf{Value 2} & \textbf{Total} \\
\midrule
Item A & 10.5 & 20.3 & 30.8 \\
Item B & 15.2 & 25.1 & 40.3 \\
Item C & 22.1 & 12.7 & 34.8 \\
\bottomrule
\end{tabular}
\caption{Sample Data Table}
\label{tab:sample}
\end{table}

\section{Lists}
\subsection{Bullet Points}
\begin{itemize}
  \item First item
  \item Second item
  \item Third item with sub-items:
    \begin{itemize}
      \item Sub-item 1
      \item Sub-item 2
    \end{itemize}
\end{itemize}

\subsection{Numbered List}
\begin{enumerate}
  \item First step
  \item Second step
  \item Third step
\end{enumerate}

\section{Citations and References}
When you use the citation search feature, references will be automatically managed. For example, you might cite a paper like this: \cite{example2024}. All citations will be automatically added to your \texttt{cite.bib} file.

\section{Conclusion}
This template demonstrates basic \LaTeX{} features including mathematical formulas, tables, and lists. The bibliography system is now set up to work seamlessly with Bibby AI's citation management. Feel free to modify it for your needs!

\bibliographystyle{plain}
\bibliography{cite}

\end{document}\documentclass{article}
\usepackage[utf8]{inputenc}
\usepackage{booktabs}
\usepackage{amsmath}
\usepackage{natbib}

\title{My Document}
\author{Your Name}
\date{\today}

\begin{document}

\maketitle

\section{Introduction}
This is a sample document demonstrating various \LaTeX{} features including mathematical equations, tables, and formatting. You can cite papers using the citation search feature, and they will be automatically added to your \texttt{cite.bib} file.

\section{Mathematical Expressions}
Here are some example equations:

\subsection{Quadratic Formula}
The quadratic formula is given by:
\[x = \frac{-b \pm \sqrt{b^2 - 4ac}}{2a}\]

\subsection{Matrix Example}
A sample matrix:
\[
\begin{pmatrix}
a_{11} & a_{12} & a_{13} \\
a_{21} & a_{22} & a_{23} \\
a_{31} & a_{32} & a_{33}
\end{pmatrix}
\]

\section{Tables}
Below is a sample table showing data:

\begin{table}[h]
\centering
\begin{tabular}{@{}lrrr@{}}
\toprule
\textbf{Category} & \textbf{Value 1} & \textbf{Value 2} & \textbf{Total} \\
\midrule
Item A & 10.5 & 20.3 & 30.8 \\
Item B & 15.2 & 25.1 & 40.3 \\
Item C & 22.1 & 12.7 & 34.8 \\
\bottomrule
\end{tabular}
\caption{Sample Data Table}
\label{tab:sample}
\end{table}

\section{Lists}
\subsection{Bullet Points}
\begin{itemize}
  \item First item
  \item Second item
  \item Third item with sub-items:
    \begin{itemize}
      \item Sub-item 1
      \item Sub-item 2
    \end{itemize}
\end{itemize}

\subsection{Numbered List}
\begin{enumerate}
  \item First step
  \item Second step
  \item Third step
\end{enumerate}

\section{Citations and References}
When you use the citation search feature, references will be automatically managed. For example, you might cite a paper like this: \cite{example2024}. All citations will be automatically added to your \texttt{cite.bib} file.

\section{Conclusion}
This template demonstrates basic \LaTeX{} features including mathematical formulas, tables, and lists. The bibliography system is now set up to work seamlessly with Bibby AI's citation management. Feel free to modify it for your needs!

\bibliographystyle{plain}
\bibliography{cite}

\end{document}\documentclass{article}
\usepackage[utf8]{inputenc}
\usepackage{booktabs}
\usepackage{amsmath}
\usepackage{natbib}

\title{My Document}
\author{Your Name}
\date{\today}

\begin{document}

\maketitle

\section{Introduction}
This is a sample document demonstrating various \LaTeX{} features including mathematical equations, tables, and formatting. You can cite papers using the citation search feature, and they will be automatically added to your \texttt{cite.bib} file.

\section{Mathematical Expressions}
Here are some example equations:

\subsection{Quadratic Formula}
The quadratic formula is given by:
\[x = \frac{-b \pm \sqrt{b^2 - 4ac}}{2a}\]

\subsection{Matrix Example}
A sample matrix:
\[
\begin{pmatrix}
a_{11} & a_{12} & a_{13} \\
a_{21} & a_{22} & a_{23} \\
a_{31} & a_{32} & a_{33}
\end{pmatrix}
\]

\section{Tables}
Below is a sample table showing data:

\begin{table}[h]
\centering
\begin{tabular}{@{}lrrr@{}}
\toprule
\textbf{Category} & \textbf{Value 1} & \textbf{Value 2} & \textbf{Total} \\
\midrule
Item A & 10.5 & 20.3 & 30.8 \\
Item B & 15.2 & 25.1 & 40.3 \\
Item C & 22.1 & 12.7 & 34.8 \\
\bottomrule
\end{tabular}
\caption{Sample Data Table}
\label{tab:sample}
\end{table}

\section{Lists}
\subsection{Bullet Points}
\begin{itemize}
  \item First item
  \item Second item
  \item Third item with sub-items:
    \begin{itemize}
      \item Sub-item 1
      \item Sub-item 2
    \end{itemize}
\end{itemize}

\subsection{Numbered List}
\begin{enumerate}
  \item First step
  \item Second step
  \item Third step
\end{enumerate}

\section{Citations and References}
When you use the citation search feature, references will be automatically managed. For example, you might cite a paper like this: \cite{example2024}. All citations will be automatically added to your \texttt{cite.bib} file.

\section{Conclusion}
This template demonstrates basic \LaTeX{} features including mathematical formulas, tables, and lists. The bibliography system is now set up to work seamlessly with Bibby AI's citation management. Feel free to modify it for your needs!

\bibliographystyle{plain}
\bibliography{cite}

\end{document}\documentclass{article}
\usepackage[utf8]{inputenc}
\usepackage{booktabs}
\usepackage{amsmath}
\usepackage{natbib}

\title{My Document}
\author{Your Name}
\date{\today}

\begin{document}

\maketitle

\section{Introduction}
This is a sample document demonstrating various \LaTeX{} features including mathematical equations, tables, and formatting. You can cite papers using the citation search feature, and they will be automatically added to your \texttt{cite.bib} file.

\section{Mathematical Expressions}
Here are some example equations:

\subsection{Quadratic Formula}
The quadratic formula is given by:
\[x = \frac{-b \pm \sqrt{b^2 - 4ac}}{2a}\]

\subsection{Matrix Example}
A sample matrix:
\[
\begin{pmatrix}
a_{11} & a_{12} & a_{13} \\
a_{21} & a_{22} & a_{23} \\
a_{31} & a_{32} & a_{33}
\end{pmatrix}
\]

\section{Tables}
Below is a sample table showing data:

\begin{table}[h]
\centering
\begin{tabular}{@{}lrrr@{}}
\toprule
\textbf{Category} & \textbf{Value 1} & \textbf{Value 2} & \textbf{Total} \\
\midrule
Item A & 10.5 & 20.3 & 30.8 \\
Item B & 15.2 & 25.1 & 40.3 \\
Item C & 22.1 & 12.7 & 34.8 \\
\bottomrule
\end{tabular}
\caption{Sample Data Table}
\label{tab:sample}
\end{table}

\section{Lists}
\subsection{Bullet Points}
\begin{itemize}
  \item First item
  \item Second item
  \item Third item with sub-items:
    \begin{itemize}
      \item Sub-item 1
      \item Sub-item 2
    \end{itemize}
\end{itemize}

\subsection{Numbered List}
\begin{enumerate}
  \item First step
  \item Second step
  \item Third step
\end{enumerate}

\section{Citations and References}
When you use the citation search feature, references will be automatically managed. For example, you might cite a paper like this: \cite{example2024}. All citations will be automatically added to your \texttt{cite.bib} file.

\section{Conclusion}
This template demonstrates basic \LaTeX{} features including mathematical formulas, tables, and lists. The bibliography system is now set up to work seamlessly with Bibby AI's citation management. Feel free to modify it for your needs!

\bibliographystyle{plain}
\bibliography{cite}

\end{document}\documentclass{article}
\usepackage[utf8]{inputenc}
\usepackage{booktabs}
\usepackage{amsmath}
\usepackage{natbib}

\title{My Document}
\author{Your Name}
\date{\today}

\begin{document}

\maketitle

\section{Introduction}
This is a sample document demonstrating various \LaTeX{} features including mathematical equations, tables, and formatting. You can cite papers using the citation search feature, and they will be automatically added to your \texttt{cite.bib} file.

\section{Mathematical Expressions}
Here are some example equations:

\subsection{Quadratic Formula}
The quadratic formula is given by:
\[x = \frac{-b \pm \sqrt{b^2 - 4ac}}{2a}\]

\subsection{Matrix Example}
A sample matrix:
\[
\begin{pmatrix}
a_{11} & a_{12} & a_{13} \\
a_{21} & a_{22} & a_{23} \\
a_{31} & a_{32} & a_{33}
\end{pmatrix}
\]

\section{Tables}
Below is a sample table showing data:

\begin{table}[h]
\centering
\begin{tabular}{@{}lrrr@{}}
\toprule
\textbf{Category} & \textbf{Value 1} & \textbf{Value 2} & \textbf{Total} \\
\midrule
Item A & 10.5 & 20.3 & 30.8 \\
Item B & 15.2 & 25.1 & 40.3 \\
Item C & 22.1 & 12.7 & 34.8 \\
\bottomrule
\end{tabular}
\caption{Sample Data Table}
\label{tab:sample}
\end{table}

\section{Lists}
\subsection{Bullet Points}
\begin{itemize}
  \item First item
  \item Second item
  \item Third item with sub-items:
    \begin{itemize}
      \item Sub-item 1
      \item Sub-item 2
    \end{itemize}
\end{itemize}

\subsection{Numbered List}
\begin{enumerate}
  \item First step
  \item Second step
  \item Third step
\end{enumerate}

\section{Citations and References}
When you use the citation search feature, references will be automatically managed. For example, you might cite a paper like this: \cite{example2024}. All citations will be automatically added to your \texttt{cite.bib} file.

\section{Conclusion}
This template demonstrates basic \LaTeX{} features including mathematical formulas, tables, and lists. The bibliography system is now set up to work seamlessly with Bibby AI's citation management. Feel free to modify it for your needs!

\bibliographystyle{plain}
\bibliography{cite}

\end{document}\documentclass{article}
\usepackage[utf8]{inputenc}
\usepackage{booktabs}
\usepackage{amsmath}
\usepackage{natbib}

\title{My Document}
\author{Your Name}
\date{\today}

\begin{document}

\maketitle

\section{Introduction}
This is a sample document demonstrating various \LaTeX{} features including mathematical equations, tables, and formatting. You can cite papers using the citation search feature, and they will be automatically added to your \texttt{cite.bib} file.

\section{Mathematical Expressions}
Here are some example equations:

\subsection{Quadratic Formula}
The quadratic formula is given by:
\[x = \frac{-b \pm \sqrt{b^2 - 4ac}}{2a}\]

\subsection{Matrix Example}
A sample matrix:
\[
\begin{pmatrix}
a_{11} & a_{12} & a_{13} \\
a_{21} & a_{22} & a_{23} \\
a_{31} & a_{32} & a_{33}
\end{pmatrix}
\]

\section{Tables}
Below is a sample table showing data:

\begin{table}[h]
\centering
\begin{tabular}{@{}lrrr@{}}
\toprule
\textbf{Category} & \textbf{Value 1} & \textbf{Value 2} & \textbf{Total} \\
\midrule
Item A & 10.5 & 20.3 & 30.8 \\
Item B & 15.2 & 25.1 & 40.3 \\
Item C & 22.1 & 12.7 & 34.8 \\
\bottomrule
\end{tabular}
\caption{Sample Data Table}
\label{tab:sample}
\end{table}

\section{Lists}
\subsection{Bullet Points}
\begin{itemize}
  \item First item
  \item Second item
  \item Third item with sub-items:
    \begin{itemize}
      \item Sub-item 1
      \item Sub-item 2
    \end{itemize}
\end{itemize}

\subsection{Numbered List}
\begin{enumerate}
  \item First step
  \item Second step
  \item Third step
\end{enumerate}

\section{Citations and References}
When you use the citation search feature, references will be automatically managed. For example, you might cite a paper like this: \cite{example2024}. All citations will be automatically added to your \texttt{cite.bib} file.

\section{Conclusion}
This template demonstrates basic \LaTeX{} features including mathematical formulas, tables, and lists. The bibliography system is now set up to work seamlessly with Bibby AI's citation management. Feel free to modify it for your needs!

\bibliographystyle{plain}
\bibliography{cite}

\end{document}\documentclass{article}
\usepackage[utf8]{inputenc}
\usepackage{booktabs}
\usepackage{amsmath}
\usepackage{natbib}

\title{My Document}
\author{Your Name}
\date{\today}

\begin{document}

\maketitle

\section{Introduction}
This is a sample document demonstrating various \LaTeX{} features including mathematical equations, tables, and formatting. You can cite papers using the citation search feature, and they will be automatically added to your \texttt{cite.bib} file.

\section{Mathematical Expressions}
Here are some example equations:

\subsection{Quadratic Formula}
The quadratic formula is given by:
\[x = \frac{-b \pm \sqrt{b^2 - 4ac}}{2a}\]

\subsection{Matrix Example}
A sample matrix:
\[
\begin{pmatrix}
a_{11} & a_{12} & a_{13} \\
a_{21} & a_{22} & a_{23} \\
a_{31} & a_{32} & a_{33}
\end{pmatrix}
\]

\section{Tables}
Below is a sample table showing data:

\begin{table}[h]
\centering
\begin{tabular}{@{}lrrr@{}}
\toprule
\textbf{Category} & \textbf{Value 1} & \textbf{Value 2} & \textbf{Total} \\
\midrule
Item A & 10.5 & 20.3 & 30.8 \\
Item B & 15.2 & 25.1 & 40.3 \\
Item C & 22.1 & 12.7 & 34.8 \\
\bottomrule
\end{tabular}
\caption{Sample Data Table}
\label{tab:sample}
\end{table}

\section{Lists}
\subsection{Bullet Points}
\begin{itemize}
  \item First item
  \item Second item
  \item Third item with sub-items:
    \begin{itemize}
      \item Sub-item 1
      \item Sub-item 2
    \end{itemize}
\end{itemize}

\subsection{Numbered List}
\begin{enumerate}
  \item First step
  \item Second step
  \item Third step
\end{enumerate}

\section{Citations and References}
When you use the citation search feature, references will be automatically managed. For example, you might cite a paper like this: \cite{example2024}. All citations will be automatically added to your \texttt{cite.bib} file.

\section{Conclusion}
This template demonstrates basic \LaTeX{} features including mathematical formulas, tables, and lists. The bibliography system is now set up to work seamlessly with Bibby AI's citation management. Feel free to modify it for your needs!

\bibliographystyle{plain}
\bibliography{cite}

\end{document}\documentclass{article}
\usepackage[utf8]{inputenc}
\usepackage{booktabs}
\usepackage{amsmath}
\usepackage{natbib}

\title{My Document}
\author{Your Name}
\date{\today}

\begin{document}

\maketitle

\section{Introduction}
This is a sample document demonstrating various \LaTeX{} features including mathematical equations, tables, and formatting. You can cite papers using the citation search feature, and they will be automatically added to your \texttt{cite.bib} file.

\section{Mathematical Expressions}
Here are some example equations:

\subsection{Quadratic Formula}
The quadratic formula is given by:
\[x = \frac{-b \pm \sqrt{b^2 - 4ac}}{2a}\]

\subsection{Matrix Example}
A sample matrix:
\[
\begin{pmatrix}
a_{11} & a_{12} & a_{13} \\
a_{21} & a_{22} & a_{23} \\
a_{31} & a_{32} & a_{33}
\end{pmatrix}
\]

\section{Tables}
Below is a sample table showing data:

\begin{table}[h]
\centering
\begin{tabular}{@{}lrrr@{}}
\toprule
\textbf{Category} & \textbf{Value 1} & \textbf{Value 2} & \textbf{Total} \\
\midrule
Item A & 10.5 & 20.3 & 30.8 \\
Item B & 15.2 & 25.1 & 40.3 \\
Item C & 22.1 & 12.7 & 34.8 \\
\bottomrule
\end{tabular}
\caption{Sample Data Table}
\label{tab:sample}
\end{table}

\section{Lists}
\subsection{Bullet Points}
\begin{itemize}
  \item First item
  \item Second item
  \item Third item with sub-items:
    \begin{itemize}
      \item Sub-item 1
      \item Sub-item 2
    \end{itemize}
\end{itemize}

\subsection{Numbered List}
\begin{enumerate}
  \item First step
  \item Second step
  \item Third step
\end{enumerate}

\section{Citations and References}
When you use the citation search feature, references will be automatically managed. For example, you might cite a paper like this: \cite{example2024}. All citations will be automatically added to your \texttt{cite.bib} file.

\section{Conclusion}
This template demonstrates basic \LaTeX{} features including mathematical formulas, tables, and lists. The bibliography system is now set up to work seamlessly with Bibby AI's citation management. Feel free to modify it for your needs!

\bibliographystyle{plain}
\bibliography{cite}

\end{document}\documentclass{article}
\usepackage[utf8]{inputenc}
\usepackage{booktabs}
\usepackage{amsmath}
\usepackage{natbib}

\title{My Document}
\author{Your Name}
\date{\today}

\begin{document}

\maketitle

\section{Introduction}
This is a sample document demonstrating various \LaTeX{} features including mathematical equations, tables, and formatting. You can cite papers using the citation search feature, and they will be automatically added to your \texttt{cite.bib} file.

\section{Mathematical Expressions}
Here are some example equations:

\subsection{Quadratic Formula}
The quadratic formula is given by:
\[x = \frac{-b \pm \sqrt{b^2 - 4ac}}{2a}\]

\subsection{Matrix Example}
A sample matrix:
\[
\begin{pmatrix}
a_{11} & a_{12} & a_{13} \\
a_{21} & a_{22} & a_{23} \\
a_{31} & a_{32} & a_{33}
\end{pmatrix}
\]

\section{Tables}
Below is a sample table showing data:

\begin{table}[h]
\centering
\begin{tabular}{@{}lrrr@{}}
\toprule
\textbf{Category} & \textbf{Value 1} & \textbf{Value 2} & \textbf{Total} \\
\midrule
Item A & 10.5 & 20.3 & 30.8 \\
Item B & 15.2 & 25.1 & 40.3 \\
Item C & 22.1 & 12.7 & 34.8 \\
\bottomrule
\end{tabular}
\caption{Sample Data Table}
\label{tab:sample}
\end{table}

\section{Lists}
\subsection{Bullet Points}
\begin{itemize}
  \item First item
  \item Second item
  \item Third item with sub-items:
    \begin{itemize}
      \item Sub-item 1
      \item Sub-item 2
    \end{itemize}
\end{itemize}

\subsection{Numbered List}
\begin{enumerate}
  \item First step
  \item Second step
  \item Third step
\end{enumerate}

\section{Citations and References}
When you use the citation search feature, references will be automatically managed. For example, you might cite a paper like this: \cite{example2024}. All citations will be automatically added to your \texttt{cite.bib} file.

\section{Conclusion}
This template demonstrates basic \LaTeX{} features including mathematical formulas, tables, and lists. The bibliography system is now set up to work seamlessly with Bibby AI's citation management. Feel free to modify it for your needs!

\bibliographystyle{plain}
\bibliography{cite}

\end{document}\documentclass{article}
\usepackage[utf8]{inputenc}
\usepackage{booktabs}
\usepackage{amsmath}
\usepackage{natbib}

\title{My Document}
\author{Your Name}
\date{\today}

\begin{document}

\maketitle

\section{Introduction}
This is a sample document demonstrating various \LaTeX{} features including mathematical equations, tables, and formatting. You can cite papers using the citation search feature, and they will be automatically added to your \texttt{cite.bib} file.

\section{Mathematical Expressions}
Here are some example equations:

\subsection{Quadratic Formula}
The quadratic formula is given by:
\[x = \frac{-b \pm \sqrt{b^2 - 4ac}}{2a}\]

\subsection{Matrix Example}
A sample matrix:
\[
\begin{pmatrix}
a_{11} & a_{12} & a_{13} \\
a_{21} & a_{22} & a_{23} \\
a_{31} & a_{32} & a_{33}
\end{pmatrix}
\]

\section{Tables}
Below is a sample table showing data:

\begin{table}[h]
\centering
\begin{tabular}{@{}lrrr@{}}
\toprule
\textbf{Category} & \textbf{Value 1} & \textbf{Value 2} & \textbf{Total} \\
\midrule
Item A & 10.5 & 20.3 & 30.8 \\
Item B & 15.2 & 25.1 & 40.3 \\
Item C & 22.1 & 12.7 & 34.8 \\
\bottomrule
\end{tabular}
\caption{Sample Data Table}
\label{tab:sample}
\end{table}

\section{Lists}
\subsection{Bullet Points}
\begin{itemize}
  \item First item
  \item Second item
  \item Third item with sub-items:
    \begin{itemize}
      \item Sub-item 1
      \item Sub-item 2
    \end{itemize}
\end{itemize}

\subsection{Numbered List}
\begin{enumerate}
  \item First step
  \item Second step
  \item Third step
\end{enumerate}

\section{Citations and References}
When you use the citation search feature, references will be automatically managed. For example, you might cite a paper like this: \cite{example2024}. All citations will be automatically added to your \texttt{cite.bib} file.

\section{Conclusion}
This template demonstrates basic \LaTeX{} features including mathematical formulas, tables, and lists. The bibliography system is now set up to work seamlessly with Bibby AI's citation management. Feel free to modify it for your needs!

\bibliographystyle{plain}
\bibliography{cite}

\end{document}\documentclass{article}
\usepackage[utf8]{inputenc}
\usepackage{booktabs}
\usepackage{amsmath}
\usepackage{natbib}

\title{My Document}
\author{Your Name}
\date{\today}

\begin{document}

\maketitle

\section{Introduction}
This is a sample document demonstrating various \LaTeX{} features including mathematical equations, tables, and formatting. You can cite papers using the citation search feature, and they will be automatically added to your \texttt{cite.bib} file.

\section{Mathematical Expressions}
Here are some example equations:

\subsection{Quadratic Formula}
The quadratic formula is given by:
\[x = \frac{-b \pm \sqrt{b^2 - 4ac}}{2a}\]

\subsection{Matrix Example}
A sample matrix:
\[
\begin{pmatrix}
a_{11} & a_{12} & a_{13} \\
a_{21} & a_{22} & a_{23} \\
a_{31} & a_{32} & a_{33}
\end{pmatrix}
\]

\section{Tables}
Below is a sample table showing data:

\begin{table}[h]
\centering
\begin{tabular}{@{}lrrr@{}}
\toprule
\textbf{Category} & \textbf{Value 1} & \textbf{Value 2} & \textbf{Total} \\
\midrule
Item A & 10.5 & 20.3 & 30.8 \\
Item B & 15.2 & 25.1 & 40.3 \\
Item C & 22.1 & 12.7 & 34.8 \\
\bottomrule
\end{tabular}
\caption{Sample Data Table}
\label{tab:sample}
\end{table}

\section{Lists}
\subsection{Bullet Points}
\begin{itemize}
  \item First item
  \item Second item
  \item Third item with sub-items:
    \begin{itemize}
      \item Sub-item 1
      \item Sub-item 2
    \end{itemize}
\end{itemize}

\subsection{Numbered List}
\begin{enumerate}
  \item First step
  \item Second step
  \item Third step
\end{enumerate}

\section{Citations and References}
When you use the citation search feature, references will be automatically managed. For example, you might cite a paper like this: \cite{example2024}. All citations will be automatically added to your \texttt{cite.bib} file.

\section{Conclusion}
This template demonstrates basic \LaTeX{} features including mathematical formulas, tables, and lists. The bibliography system is now set up to work seamlessly with Bibby AI's citation management. Feel free to modify it for your needs!

\bibliographystyle{plain}
\bibliography{cite}

\end{document}\documentclass{article}
\usepackage[utf8]{inputenc}
\usepackage{booktabs}
\usepackage{amsmath}
\usepackage{natbib}

\title{My Document}
\author{Your Name}
\date{\today}

\begin{document}

\maketitle

\section{Introduction}
This is a sample document demonstrating various \LaTeX{} features including mathematical equations, tables, and formatting. You can cite papers using the citation search feature, and they will be automatically added to your \texttt{cite.bib} file.

\section{Mathematical Expressions}
Here are some example equations:

\subsection{Quadratic Formula}
The quadratic formula is given by:
\[x = \frac{-b \pm \sqrt{b^2 - 4ac}}{2a}\]

\subsection{Matrix Example}
A sample matrix:
\[
\begin{pmatrix}
a_{11} & a_{12} & a_{13} \\
a_{21} & a_{22} & a_{23} \\
a_{31} & a_{32} & a_{33}
\end{pmatrix}
\]

\section{Tables}
Below is a sample table showing data:

\begin{table}[h]
\centering
\begin{tabular}{@{}lrrr@{}}
\toprule
\textbf{Category} & \textbf{Value 1} & \textbf{Value 2} & \textbf{Total} \\
\midrule
Item A & 10.5 & 20.3 & 30.8 \\
Item B & 15.2 & 25.1 & 40.3 \\
Item C & 22.1 & 12.7 & 34.8 \\
\bottomrule
\end{tabular}
\caption{Sample Data Table}
\label{tab:sample}
\end{table}

\section{Lists}
\subsection{Bullet Points}
\begin{itemize}
  \item First item
  \item Second item
  \item Third item with sub-items:
    \begin{itemize}
      \item Sub-item 1
      \item Sub-item 2
    \end{itemize}
\end{itemize}

\subsection{Numbered List}
\begin{enumerate}
  \item First step
  \item Second step
  \item Third step
\end{enumerate}

\section{Citations and References}
When you use the citation search feature, references will be automatically managed. For example, you might cite a paper like this: \cite{example2024}. All citations will be automatically added to your \texttt{cite.bib} file.

\section{Conclusion}
This template demonstrates basic \LaTeX{} features including mathematical formulas, tables, and lists. The bibliography system is now set up to work seamlessly with Bibby AI's citation management. Feel free to modify it for your needs!

\bibliographystyle{plain}
\bibliography{cite}

\end{document}\documentclass{article}
\usepackage[utf8]{inputenc}
\usepackage{booktabs}
\usepackage{amsmath}
\usepackage{natbib}

\title{My Document}
\author{Your Name}
\date{\today}

\begin{document}

\maketitle

\section{Introduction}
This is a sample document demonstrating various \LaTeX{} features including mathematical equations, tables, and formatting. You can cite papers using the citation search feature, and they will be automatically added to your \texttt{cite.bib} file.

\section{Mathematical Expressions}
Here are some example equations:

\subsection{Quadratic Formula}
The quadratic formula is given by:
\[x = \frac{-b \pm \sqrt{b^2 - 4ac}}{2a}\]

\subsection{Matrix Example}
A sample matrix:
\[
\begin{pmatrix}
a_{11} & a_{12} & a_{13} \\
a_{21} & a_{22} & a_{23} \\
a_{31} & a_{32} & a_{33}
\end{pmatrix}
\]

\section{Tables}
Below is a sample table showing data:

\begin{table}[h]
\centering
\begin{tabular}{@{}lrrr@{}}
\toprule
\textbf{Category} & \textbf{Value 1} & \textbf{Value 2} & \textbf{Total} \\
\midrule
Item A & 10.5 & 20.3 & 30.8 \\
Item B & 15.2 & 25.1 & 40.3 \\
Item C & 22.1 & 12.7 & 34.8 \\
\bottomrule
\end{tabular}
\caption{Sample Data Table}
\label{tab:sample}
\end{table}

\section{Lists}
\subsection{Bullet Points}
\begin{itemize}
  \item First item
  \item Second item
  \item Third item with sub-items:
    \begin{itemize}
      \item Sub-item 1
      \item Sub-item 2
    \end{itemize}
\end{itemize}

\subsection{Numbered List}
\begin{enumerate}
  \item First step
  \item Second step
  \item Third step
\end{enumerate}

\section{Citations and References}
When you use the citation search feature, references will be automatically managed. For example, you might cite a paper like this: \cite{example2024}. All citations will be automatically added to your \texttt{cite.bib} file.

\section{Conclusion}
This template demonstrates basic \LaTeX{} features including mathematical formulas, tables, and lists. The bibliography system is now set up to work seamlessly with Bibby AI's citation management. Feel free to modify it for your needs!

\bibliographystyle{plain}
\bibliography{cite}

\end{document}\documentclass{article}
\usepackage[utf8]{inputenc}
\usepackage{booktabs}
\usepackage{amsmath}
\usepackage{natbib}

\title{My Document}
\author{Your Name}
\date{\today}

\begin{document}

\maketitle

\section{Introduction}
This is a sample document demonstrating various \LaTeX{} features including mathematical equations, tables, and formatting. You can cite papers using the citation search feature, and they will be automatically added to your \texttt{cite.bib} file.

\section{Mathematical Expressions}
Here are some example equations:

\subsection{Quadratic Formula}
The quadratic formula is given by:
\[x = \frac{-b \pm \sqrt{b^2 - 4ac}}{2a}\]

\subsection{Matrix Example}
A sample matrix:
\[
\begin{pmatrix}
a_{11} & a_{12} & a_{13} \\
a_{21} & a_{22} & a_{23} \\
a_{31} & a_{32} & a_{33}
\end{pmatrix}
\]

\section{Tables}
Below is a sample table showing data:

\begin{table}[h]
\centering
\begin{tabular}{@{}lrrr@{}}
\toprule
\textbf{Category} & \textbf{Value 1} & \textbf{Value 2} & \textbf{Total} \\
\midrule
Item A & 10.5 & 20.3 & 30.8 \\
Item B & 15.2 & 25.1 & 40.3 \\
Item C & 22.1 & 12.7 & 34.8 \\
\bottomrule
\end{tabular}
\caption{Sample Data Table}
\label{tab:sample}
\end{table}

\section{Lists}
\subsection{Bullet Points}
\begin{itemize}
  \item First item
  \item Second item
  \item Third item with sub-items:
    \begin{itemize}
      \item Sub-item 1
      \item Sub-item 2
    \end{itemize}
\end{itemize}

\subsection{Numbered List}
\begin{enumerate}
  \item First step
  \item Second step
  \item Third step
\end{enumerate}

\section{Citations and References}
When you use the citation search feature, references will be automatically managed. For example, you might cite a paper like this: \cite{example2024}. All citations will be automatically added to your \texttt{cite.bib} file.

\section{Conclusion}
This template demonstrates basic \LaTeX{} features including mathematical formulas, tables, and lists. The bibliography system is now set up to work seamlessly with Bibby AI's citation management. Feel free to modify it for your needs!

\bibliographystyle{plain}
\bibliography{cite}

\end{document}\documentclass{article}
\usepackage[utf8]{inputenc}
\usepackage{booktabs}
\usepackage{amsmath}
\usepackage{natbib}

\title{My Document}
\author{Your Name}
\date{\today}

\begin{document}

\maketitle

\section{Introduction}
This is a sample document demonstrating various \LaTeX{} features including mathematical equations, tables, and formatting. You can cite papers using the citation search feature, and they will be automatically added to your \texttt{cite.bib} file.

\section{Mathematical Expressions}
Here are some example equations:

\subsection{Quadratic Formula}
The quadratic formula is given by:
\[x = \frac{-b \pm \sqrt{b^2 - 4ac}}{2a}\]

\subsection{Matrix Example}
A sample matrix:
\[
\begin{pmatrix}
a_{11} & a_{12} & a_{13} \\
a_{21} & a_{22} & a_{23} \\
a_{31} & a_{32} & a_{33}
\end{pmatrix}
\]

\section{Tables}
Below is a sample table showing data:

\begin{table}[h]
\centering
\begin{tabular}{@{}lrrr@{}}
\toprule
\textbf{Category} & \textbf{Value 1} & \textbf{Value 2} & \textbf{Total} \\
\midrule
Item A & 10.5 & 20.3 & 30.8 \\
Item B & 15.2 & 25.1 & 40.3 \\
Item C & 22.1 & 12.7 & 34.8 \\
\bottomrule
\end{tabular}
\caption{Sample Data Table}
\label{tab:sample}
\end{table}

\section{Lists}
\subsection{Bullet Points}
\begin{itemize}
  \item First item
  \item Second item
  \item Third item with sub-items:
    \begin{itemize}
      \item Sub-item 1
      \item Sub-item 2
    \end{itemize}
\end{itemize}

\subsection{Numbered List}
\begin{enumerate}
  \item First step
  \item Second step
  \item Third step
\end{enumerate}

\section{Citations and References}
When you use the citation search feature, references will be automatically managed. For example, you might cite a paper like this: \cite{example2024}. All citations will be automatically added to your \texttt{cite.bib} file.

\section{Conclusion}
This template demonstrates basic \LaTeX{} features including mathematical formulas, tables, and lists. The bibliography system is now set up to work seamlessly with Bibby AI's citation management. Feel free to modify it for your needs!

\bibliographystyle{plain}
\bibliography{cite}

\end{document}\documentclass{article}
\usepackage[utf8]{inputenc}
\usepackage{booktabs}
\usepackage{amsmath}
\usepackage{natbib}

\title{My Document}
\author{Your Name}
\date{\today}

\begin{document}

\maketitle

\section{Introduction}
This is a sample document demonstrating various \LaTeX{} features including mathematical equations, tables, and formatting. You can cite papers using the citation search feature, and they will be automatically added to your \texttt{cite.bib} file.

\section{Mathematical Expressions}
Here are some example equations:

\subsection{Quadratic Formula}
The quadratic formula is given by:
\[x = \frac{-b \pm \sqrt{b^2 - 4ac}}{2a}\]

\subsection{Matrix Example}
A sample matrix:
\[
\begin{pmatrix}
a_{11} & a_{12} & a_{13} \\
a_{21} & a_{22} & a_{23} \\
a_{31} & a_{32} & a_{33}
\end{pmatrix}
\]

\section{Tables}
Below is a sample table showing data:

\begin{table}[h]
\centering
\begin{tabular}{@{}lrrr@{}}
\toprule
\textbf{Category} & \textbf{Value 1} & \textbf{Value 2} & \textbf{Total} \\
\midrule
Item A & 10.5 & 20.3 & 30.8 \\
Item B & 15.2 & 25.1 & 40.3 \\
Item C & 22.1 & 12.7 & 34.8 \\
\bottomrule
\end{tabular}
\caption{Sample Data Table}
\label{tab:sample}
\end{table}

\section{Lists}
\subsection{Bullet Points}
\begin{itemize}
  \item First item
  \item Second item
  \item Third item with sub-items:
    \begin{itemize}
      \item Sub-item 1
      \item Sub-item 2
    \end{itemize}
\end{itemize}

\subsection{Numbered List}
\begin{enumerate}
  \item First step
  \item Second step
  \item Third step
\end{enumerate}

\section{Citations and References}
When you use the citation search feature, references will be automatically managed. For example, you might cite a paper like this: \cite{example2024}. All citations will be automatically added to your \texttt{cite.bib} file.

\section{Conclusion}
This template demonstrates basic \LaTeX{} features including mathematical formulas, tables, and lists. The bibliography system is now set up to work seamlessly with Bibby AI's citation management. Feel free to modify it for your needs!

\bibliographystyle{plain}
\bibliography{cite}

\end{document}\documentclass{article}
\usepackage[utf8]{inputenc}
\usepackage{booktabs}
\usepackage{amsmath}
\usepackage{natbib}

\title{My Document}
\author{Your Name}
\date{\today}

\begin{document}

\maketitle

\section{Introduction}
This is a sample document demonstrating various \LaTeX{} features including mathematical equations, tables, and formatting. You can cite papers using the citation search feature, and they will be automatically added to your \texttt{cite.bib} file.

\section{Mathematical Expressions}
Here are some example equations:

\subsection{Quadratic Formula}
The quadratic formula is given by:
\[x = \frac{-b \pm \sqrt{b^2 - 4ac}}{2a}\]

\subsection{Matrix Example}
A sample matrix:
\[
\begin{pmatrix}
a_{11} & a_{12} & a_{13} \\
a_{21} & a_{22} & a_{23} \\
a_{31} & a_{32} & a_{33}
\end{pmatrix}
\]

\section{Tables}
Below is a sample table showing data:

\begin{table}[h]
\centering
\begin{tabular}{@{}lrrr@{}}
\toprule
\textbf{Category} & \textbf{Value 1} & \textbf{Value 2} & \textbf{Total} \\
\midrule
Item A & 10.5 & 20.3 & 30.8 \\
Item B & 15.2 & 25.1 & 40.3 \\
Item C & 22.1 & 12.7 & 34.8 \\
\bottomrule
\end{tabular}
\caption{Sample Data Table}
\label{tab:sample}
\end{table}

\section{Lists}
\subsection{Bullet Points}
\begin{itemize}
  \item First item
  \item Second item
  \item Third item with sub-items:
    \begin{itemize}
      \item Sub-item 1
      \item Sub-item 2
    \end{itemize}
\end{itemize}

\subsection{Numbered List}
\begin{enumerate}
  \item First step
  \item Second step
  \item Third step
\end{enumerate}

\section{Citations and References}
When you use the citation search feature, references will be automatically managed. For example, you might cite a paper like this: \cite{example2024}. All citations will be automatically added to your \texttt{cite.bib} file.

\section{Conclusion}
This template demonstrates basic \LaTeX{} features including mathematical formulas, tables, and lists. The bibliography system is now set up to work seamlessly with Bibby AI's citation management. Feel free to modify it for your needs!

\bibliographystyle{plain}
\bibliography{cite}

\end{document}\documentclass{article}
\usepackage[utf8]{inputenc}
\usepackage{booktabs}
\usepackage{amsmath}
\usepackage{natbib}

\title{My Document}
\author{Your Name}
\date{\today}

\begin{document}

\maketitle

\section{Introduction}
This is a sample document demonstrating various \LaTeX{} features including mathematical equations, tables, and formatting. You can cite papers using the citation search feature, and they will be automatically added to your \texttt{cite.bib} file.

\section{Mathematical Expressions}
Here are some example equations:

\subsection{Quadratic Formula}
The quadratic formula is given by:
\[x = \frac{-b \pm \sqrt{b^2 - 4ac}}{2a}\]

\subsection{Matrix Example}
A sample matrix:
\[
\begin{pmatrix}
a_{11} & a_{12} & a_{13} \\
a_{21} & a_{22} & a_{23} \\
a_{31} & a_{32} & a_{33}
\end{pmatrix}
\]

\section{Tables}
Below is a sample table showing data:

\begin{table}[h]
\centering
\begin{tabular}{@{}lrrr@{}}
\toprule
\textbf{Category} & \textbf{Value 1} & \textbf{Value 2} & \textbf{Total} \\
\midrule
Item A & 10.5 & 20.3 & 30.8 \\
Item B & 15.2 & 25.1 & 40.3 \\
Item C & 22.1 & 12.7 & 34.8 \\
\bottomrule
\end{tabular}
\caption{Sample Data Table}
\label{tab:sample}
\end{table}

\section{Lists}
\subsection{Bullet Points}
\begin{itemize}
  \item First item
  \item Second item
  \item Third item with sub-items:
    \begin{itemize}
      \item Sub-item 1
      \item Sub-item 2
    \end{itemize}
\end{itemize}

\subsection{Numbered List}
\begin{enumerate}
  \item First step
  \item Second step
  \item Third step
\end{enumerate}

\section{Citations and References}
When you use the citation search feature, references will be automatically managed. For example, you might cite a paper like this: \cite{example2024}. All citations will be automatically added to your \texttt{cite.bib} file.

\section{Conclusion}
This template demonstrates basic \LaTeX{} features including mathematical formulas, tables, and lists. The bibliography system is now set up to work seamlessly with Bibby AI's citation management. Feel free to modify it for your needs!

\bibliographystyle{plain}
\bibliography{cite}

\end{document}\documentclass{article}
\usepackage[utf8]{inputenc}
\usepackage{booktabs}
\usepackage{amsmath}
\usepackage{natbib}

\title{My Document}
\author{Your Name}
\date{\today}

\begin{document}

\maketitle

\section{Introduction}
This is a sample document demonstrating various \LaTeX{} features including mathematical equations, tables, and formatting. You can cite papers using the citation search feature, and they will be automatically added to your \texttt{cite.bib} file.

\section{Mathematical Expressions}
Here are some example equations:

\subsection{Quadratic Formula}
The quadratic formula is given by:
\[x = \frac{-b \pm \sqrt{b^2 - 4ac}}{2a}\]

\subsection{Matrix Example}
A sample matrix:
\[
\begin{pmatrix}
a_{11} & a_{12} & a_{13} \\
a_{21} & a_{22} & a_{23} \\
a_{31} & a_{32} & a_{33}
\end{pmatrix}
\]

\section{Tables}
Below is a sample table showing data:

\begin{table}[h]
\centering
\begin{tabular}{@{}lrrr@{}}
\toprule
\textbf{Category} & \textbf{Value 1} & \textbf{Value 2} & \textbf{Total} \\
\midrule
Item A & 10.5 & 20.3 & 30.8 \\
Item B & 15.2 & 25.1 & 40.3 \\
Item C & 22.1 & 12.7 & 34.8 \\
\bottomrule
\end{tabular}
\caption{Sample Data Table}
\label{tab:sample}
\end{table}

\section{Lists}
\subsection{Bullet Points}
\begin{itemize}
  \item First item
  \item Second item
  \item Third item with sub-items:
    \begin{itemize}
      \item Sub-item 1
      \item Sub-item 2
    \end{itemize}
\end{itemize}

\subsection{Numbered List}
\begin{enumerate}
  \item First step
  \item Second step
  \item Third step
\end{enumerate}

\section{Citations and References}
When you use the citation search feature, references will be automatically managed. For example, you might cite a paper like this: \cite{example2024}. All citations will be automatically added to your \texttt{cite.bib} file.

\section{Conclusion}
This template demonstrates basic \LaTeX{} features including mathematical formulas, tables, and lists. The bibliography system is now set up to work seamlessly with Bibby AI's citation management. Feel free to modify it for your needs!

\bibliographystyle{plain}
\bibliography{cite}

\end{document}\documentclass{article}
\usepackage[utf8]{inputenc}
\usepackage{booktabs}
\usepackage{amsmath}
\usepackage{natbib}

\title{My Document}
\author{Your Name}
\date{\today}

\begin{document}

\maketitle

\section{Introduction}
This is a sample document demonstrating various \LaTeX{} features including mathematical equations, tables, and formatting. You can cite papers using the citation search feature, and they will be automatically added to your \texttt{cite.bib} file.

\section{Mathematical Expressions}
Here are some example equations:

\subsection{Quadratic Formula}
The quadratic formula is given by:
\[x = \frac{-b \pm \sqrt{b^2 - 4ac}}{2a}\]

\subsection{Matrix Example}
A sample matrix:
\[
\begin{pmatrix}
a_{11} & a_{12} & a_{13} \\
a_{21} & a_{22} & a_{23} \\
a_{31} & a_{32} & a_{33}
\end{pmatrix}
\]

\section{Tables}
Below is a sample table showing data:

\begin{table}[h]
\centering
\begin{tabular}{@{}lrrr@{}}
\toprule
\textbf{Category} & \textbf{Value 1} & \textbf{Value 2} & \textbf{Total} \\
\midrule
Item A & 10.5 & 20.3 & 30.8 \\
Item B & 15.2 & 25.1 & 40.3 \\
Item C & 22.1 & 12.7 & 34.8 \\
\bottomrule
\end{tabular}
\caption{Sample Data Table}
\label{tab:sample}
\end{table}

\section{Lists}
\subsection{Bullet Points}
\begin{itemize}
  \item First item
  \item Second item
  \item Third item with sub-items:
    \begin{itemize}
      \item Sub-item 1
      \item Sub-item 2
    \end{itemize}
\end{itemize}

\subsection{Numbered List}
\begin{enumerate}
  \item First step
  \item Second step
  \item Third step
\end{enumerate}

\section{Citations and References}
When you use the citation search feature, references will be automatically managed. For example, you might cite a paper like this: \cite{example2024}. All citations will be automatically added to your \texttt{cite.bib} file.

\section{Conclusion}
This template demonstrates basic \LaTeX{} features including mathematical formulas, tables, and lists. The bibliography system is now set up to work seamlessly with Bibby AI's citation management. Feel free to modify it for your needs!

\bibliographystyle{plain}
\bibliography{cite}

\end{document}\documentclass{article}
\usepackage[utf8]{inputenc}
\usepackage{booktabs}
\usepackage{amsmath}
\usepackage{natbib}

\title{My Document}
\author{Your Name}
\date{\today}

\begin{document}

\maketitle

\section{Introduction}
This is a sample document demonstrating various \LaTeX{} features including mathematical equations, tables, and formatting. You can cite papers using the citation search feature, and they will be automatically added to your \texttt{cite.bib} file.

\section{Mathematical Expressions}
Here are some example equations:

\subsection{Quadratic Formula}
The quadratic formula is given by:
\[x = \frac{-b \pm \sqrt{b^2 - 4ac}}{2a}\]

\subsection{Matrix Example}
A sample matrix:
\[
\begin{pmatrix}
a_{11} & a_{12} & a_{13} \\
a_{21} & a_{22} & a_{23} \\
a_{31} & a_{32} & a_{33}
\end{pmatrix}
\]

\section{Tables}
Below is a sample table showing data:

\begin{table}[h]
\centering
\begin{tabular}{@{}lrrr@{}}
\toprule
\textbf{Category} & \textbf{Value 1} & \textbf{Value 2} & \textbf{Total} \\
\midrule
Item A & 10.5 & 20.3 & 30.8 \\
Item B & 15.2 & 25.1 & 40.3 \\
Item C & 22.1 & 12.7 & 34.8 \\
\bottomrule
\end{tabular}
\caption{Sample Data Table}
\label{tab:sample}
\end{table}

\section{Lists}
\subsection{Bullet Points}
\begin{itemize}
  \item First item
  \item Second item
  \item Third item with sub-items:
    \begin{itemize}
      \item Sub-item 1
      \item Sub-item 2
    \end{itemize}
\end{itemize}

\subsection{Numbered List}
\begin{enumerate}
  \item First step
  \item Second step
  \item Third step
\end{enumerate}

\section{Citations and References}
When you use the citation search feature, references will be automatically managed. For example, you might cite a paper like this: \cite{example2024}. All citations will be automatically added to your \texttt{cite.bib} file.

\section{Conclusion}
This template demonstrates basic \LaTeX{} features including mathematical formulas, tables, and lists. The bibliography system is now set up to work seamlessly with Bibby AI's citation management. Feel free to modify it for your needs!

\bibliographystyle{plain}
\bibliography{cite}

\end{document}\documentclass{article}
\usepackage[utf8]{inputenc}
\usepackage{booktabs}
\usepackage{amsmath}
\usepackage{natbib}

\title{My Document}
\author{Your Name}
\date{\today}

\begin{document}

\maketitle

\section{Introduction}
This is a sample document demonstrating various \LaTeX{} features including mathematical equations, tables, and formatting. You can cite papers using the citation search feature, and they will be automatically added to your \texttt{cite.bib} file.

\section{Mathematical Expressions}
Here are some example equations:

\subsection{Quadratic Formula}
The quadratic formula is given by:
\[x = \frac{-b \pm \sqrt{b^2 - 4ac}}{2a}\]

\subsection{Matrix Example}
A sample matrix:
\[
\begin{pmatrix}
a_{11} & a_{12} & a_{13} \\
a_{21} & a_{22} & a_{23} \\
a_{31} & a_{32} & a_{33}
\end{pmatrix}
\]

\section{Tables}
Below is a sample table showing data:

\begin{table}[h]
\centering
\begin{tabular}{@{}lrrr@{}}
\toprule
\textbf{Category} & \textbf{Value 1} & \textbf{Value 2} & \textbf{Total} \\
\midrule
Item A & 10.5 & 20.3 & 30.8 \\
Item B & 15.2 & 25.1 & 40.3 \\
Item C & 22.1 & 12.7 & 34.8 \\
\bottomrule
\end{tabular}
\caption{Sample Data Table}
\label{tab:sample}
\end{table}

\section{Lists}
\subsection{Bullet Points}
\begin{itemize}
  \item First item
  \item Second item
  \item Third item with sub-items:
    \begin{itemize}
      \item Sub-item 1
      \item Sub-item 2
    \end{itemize}
\end{itemize}

\subsection{Numbered List}
\begin{enumerate}
  \item First step
  \item Second step
  \item Third step
\end{enumerate}

\section{Citations and References}
When you use the citation search feature, references will be automatically managed. For example, you might cite a paper like this: \cite{example2024}. All citations will be automatically added to your \texttt{cite.bib} file.

\section{Conclusion}
This template demonstrates basic \LaTeX{} features including mathematical formulas, tables, and lists. The bibliography system is now set up to work seamlessly with Bibby AI's citation management. Feel free to modify it for your needs!

\bibliographystyle{plain}
\bibliography{cite}

\end{document}\documentclass{article}
\usepackage[utf8]{inputenc}
\usepackage{booktabs}
\usepackage{amsmath}
\usepackage{natbib}

\title{My Document}
\author{Your Name}
\date{\today}

\begin{document}

\maketitle

\section{Introduction}
This is a sample document demonstrating various \LaTeX{} features including mathematical equations, tables, and formatting. You can cite papers using the citation search feature, and they will be automatically added to your \texttt{cite.bib} file.

\section{Mathematical Expressions}
Here are some example equations:

\subsection{Quadratic Formula}
The quadratic formula is given by:
\[x = \frac{-b \pm \sqrt{b^2 - 4ac}}{2a}\]

\subsection{Matrix Example}
A sample matrix:
\[
\begin{pmatrix}
a_{11} & a_{12} & a_{13} \\
a_{21} & a_{22} & a_{23} \\
a_{31} & a_{32} & a_{33}
\end{pmatrix}
\]

\section{Tables}
Below is a sample table showing data:

\begin{table}[h]
\centering
\begin{tabular}{@{}lrrr@{}}
\toprule
\textbf{Category} & \textbf{Value 1} & \textbf{Value 2} & \textbf{Total} \\
\midrule
Item A & 10.5 & 20.3 & 30.8 \\
Item B & 15.2 & 25.1 & 40.3 \\
Item C & 22.1 & 12.7 & 34.8 \\
\bottomrule
\end{tabular}
\caption{Sample Data Table}
\label{tab:sample}
\end{table}

\section{Lists}
\subsection{Bullet Points}
\begin{itemize}
  \item First item
  \item Second item
  \item Third item with sub-items:
    \begin{itemize}
      \item Sub-item 1
      \item Sub-item 2
    \end{itemize}
\end{itemize}

\subsection{Numbered List}
\begin{enumerate}
  \item First step
  \item Second step
  \item Third step
\end{enumerate}

\section{Citations and References}
When you use the citation search feature, references will be automatically managed. For example, you might cite a paper like this: \cite{example2024}. All citations will be automatically added to your \texttt{cite.bib} file.

\section{Conclusion}
This template demonstrates basic \LaTeX{} features including mathematical formulas, tables, and lists. The bibliography system is now set up to work seamlessly with Bibby AI's citation management. Feel free to modify it for your needs!

\bibliographystyle{plain}
\bibliography{cite}

\end{document}\documentclass{article}
\usepackage[utf8]{inputenc}
\usepackage{booktabs}
\usepackage{amsmath}
\usepackage{natbib}

\title{My Document}
\author{Your Name}
\date{\today}

\begin{document}

\maketitle

\section{Introduction}
This is a sample document demonstrating various \LaTeX{} features including mathematical equations, tables, and formatting. You can cite papers using the citation search feature, and they will be automatically added to your \texttt{cite.bib} file.

\section{Mathematical Expressions}
Here are some example equations:

\subsection{Quadratic Formula}
The quadratic formula is given by:
\[x = \frac{-b \pm \sqrt{b^2 - 4ac}}{2a}\]

\subsection{Matrix Example}
A sample matrix:
\[
\begin{pmatrix}
a_{11} & a_{12} & a_{13} \\
a_{21} & a_{22} & a_{23} \\
a_{31} & a_{32} & a_{33}
\end{pmatrix}
\]

\section{Tables}
Below is a sample table showing data:

\begin{table}[h]
\centering
\begin{tabular}{@{}lrrr@{}}
\toprule
\textbf{Category} & \textbf{Value 1} & \textbf{Value 2} & \textbf{Total} \\
\midrule
Item A & 10.5 & 20.3 & 30.8 \\
Item B & 15.2 & 25.1 & 40.3 \\
Item C & 22.1 & 12.7 & 34.8 \\
\bottomrule
\end{tabular}
\caption{Sample Data Table}
\label{tab:sample}
\end{table}

\section{Lists}
\subsection{Bullet Points}
\begin{itemize}
  \item First item
  \item Second item
  \item Third item with sub-items:
    \begin{itemize}
      \item Sub-item 1
      \item Sub-item 2
    \end{itemize}
\end{itemize}

\subsection{Numbered List}
\begin{enumerate}
  \item First step
  \item Second step
  \item Third step
\end{enumerate}

\section{Citations and References}
When you use the citation search feature, references will be automatically managed. For example, you might cite a paper like this: \cite{example2024}. All citations will be automatically added to your \texttt{cite.bib} file.

\section{Conclusion}
This template demonstrates basic \LaTeX{} features including mathematical formulas, tables, and lists. The bibliography system is now set up to work seamlessly with Bibby AI's citation management. Feel free to modify it for your needs!

\bibliographystyle{plain}
\bibliography{cite}

\end{document}\documentclass{article}
\usepackage[utf8]{inputenc}
\usepackage{booktabs}
\usepackage{amsmath}
\usepackage{natbib}

\title{My Document}
\author{Your Name}
\date{\today}

\begin{document}

\maketitle

\section{Introduction}
This is a sample document demonstrating various \LaTeX{} features including mathematical equations, tables, and formatting. You can cite papers using the citation search feature, and they will be automatically added to your \texttt{cite.bib} file.

\section{Mathematical Expressions}
Here are some example equations:

\subsection{Quadratic Formula}
The quadratic formula is given by:
\[x = \frac{-b \pm \sqrt{b^2 - 4ac}}{2a}\]

\subsection{Matrix Example}
A sample matrix:
\[
\begin{pmatrix}
a_{11} & a_{12} & a_{13} \\
a_{21} & a_{22} & a_{23} \\
a_{31} & a_{32} & a_{33}
\end{pmatrix}
\]

\section{Tables}
Below is a sample table showing data:

\begin{table}[h]
\centering
\begin{tabular}{@{}lrrr@{}}
\toprule
\textbf{Category} & \textbf{Value 1} & \textbf{Value 2} & \textbf{Total} \\
\midrule
Item A & 10.5 & 20.3 & 30.8 \\
Item B & 15.2 & 25.1 & 40.3 \\
Item C & 22.1 & 12.7 & 34.8 \\
\bottomrule
\end{tabular}
\caption{Sample Data Table}
\label{tab:sample}
\end{table}

\section{Lists}
\subsection{Bullet Points}
\begin{itemize}
  \item First item
  \item Second item
  \item Third item with sub-items:
    \begin{itemize}
      \item Sub-item 1
      \item Sub-item 2
    \end{itemize}
\end{itemize}

\subsection{Numbered List}
\begin{enumerate}
  \item First step
  \item Second step
  \item Third step
\end{enumerate}

\section{Citations and References}
When you use the citation search feature, references will be automatically managed. For example, you might cite a paper like this: \cite{example2024}. All citations will be automatically added to your \texttt{cite.bib} file.

\section{Conclusion}
This template demonstrates basic \LaTeX{} features including mathematical formulas, tables, and lists. The bibliography system is now set up to work seamlessly with Bibby AI's citation management. Feel free to modify it for your needs!

\bibliographystyle{plain}
\bibliography{cite}

\end{document}\documentclass{article}
\usepackage[utf8]{inputenc}
\usepackage{booktabs}
\usepackage{amsmath}
\usepackage{natbib}

\title{My Document}
\author{Your Name}
\date{\today}

\begin{document}

\maketitle

\section{Introduction}
This is a sample document demonstrating various \LaTeX{} features including mathematical equations, tables, and formatting. You can cite papers using the citation search feature, and they will be automatically added to your \texttt{cite.bib} file.

\section{Mathematical Expressions}
Here are some example equations:

\subsection{Quadratic Formula}
The quadratic formula is given by:
\[x = \frac{-b \pm \sqrt{b^2 - 4ac}}{2a}\]

\subsection{Matrix Example}
A sample matrix:
\[
\begin{pmatrix}
a_{11} & a_{12} & a_{13} \\
a_{21} & a_{22} & a_{23} \\
a_{31} & a_{32} & a_{33}
\end{pmatrix}
\]

\section{Tables}
Below is a sample table showing data:

\begin{table}[h]
\centering
\begin{tabular}{@{}lrrr@{}}
\toprule
\textbf{Category} & \textbf{Value 1} & \textbf{Value 2} & \textbf{Total} \\
\midrule
Item A & 10.5 & 20.3 & 30.8 \\
Item B & 15.2 & 25.1 & 40.3 \\
Item C & 22.1 & 12.7 & 34.8 \\
\bottomrule
\end{tabular}
\caption{Sample Data Table}
\label{tab:sample}
\end{table}

\section{Lists}
\subsection{Bullet Points}
\begin{itemize}
  \item First item
  \item Second item
  \item Third item with sub-items:
    \begin{itemize}
      \item Sub-item 1
      \item Sub-item 2
    \end{itemize}
\end{itemize}

\subsection{Numbered List}
\begin{enumerate}
  \item First step
  \item Second step
  \item Third step
\end{enumerate}

\section{Citations and References}
When you use the citation search feature, references will be automatically managed. For example, you might cite a paper like this: \cite{example2024}. All citations will be automatically added to your \texttt{cite.bib} file.

\section{Conclusion}
This template demonstrates basic \LaTeX{} features including mathematical formulas, tables, and lists. The bibliography system is now set up to work seamlessly with Bibby AI's citation management. Feel free to modify it for your needs!

\bibliographystyle{plain}
\bibliography{cite}

\end{document}\documentclass{article}
\usepackage[utf8]{inputenc}
\usepackage{booktabs}
\usepackage{amsmath}
\usepackage{natbib}

\title{My Document}
\author{Your Name}
\date{\today}

\begin{document}

\maketitle

\section{Introduction}
This is a sample document demonstrating various \LaTeX{} features including mathematical equations, tables, and formatting. You can cite papers using the citation search feature, and they will be automatically added to your \texttt{cite.bib} file.

\section{Mathematical Expressions}
Here are some example equations:

\subsection{Quadratic Formula}
The quadratic formula is given by:
\[x = \frac{-b \pm \sqrt{b^2 - 4ac}}{2a}\]

\subsection{Matrix Example}
A sample matrix:
\[
\begin{pmatrix}
a_{11} & a_{12} & a_{13} \\
a_{21} & a_{22} & a_{23} \\
a_{31} & a_{32} & a_{33}
\end{pmatrix}
\]

\section{Tables}
Below is a sample table showing data:

\begin{table}[h]
\centering
\begin{tabular}{@{}lrrr@{}}
\toprule
\textbf{Category} & \textbf{Value 1} & \textbf{Value 2} & \textbf{Total} \\
\midrule
Item A & 10.5 & 20.3 & 30.8 \\
Item B & 15.2 & 25.1 & 40.3 \\
Item C & 22.1 & 12.7 & 34.8 \\
\bottomrule
\end{tabular}
\caption{Sample Data Table}
\label{tab:sample}
\end{table}

\section{Lists}
\subsection{Bullet Points}
\begin{itemize}
  \item First item
  \item Second item
  \item Third item with sub-items:
    \begin{itemize}
      \item Sub-item 1
      \item Sub-item 2
    \end{itemize}
\end{itemize}

\subsection{Numbered List}
\begin{enumerate}
  \item First step
  \item Second step
  \item Third step
\end{enumerate}

\section{Citations and References}
When you use the citation search feature, references will be automatically managed. For example, you might cite a paper like this: \cite{example2024}. All citations will be automatically added to your \texttt{cite.bib} file.

\section{Conclusion}
This template demonstrates basic \LaTeX{} features including mathematical formulas, tables, and lists. The bibliography system is now set up to work seamlessly with Bibby AI's citation management. Feel free to modify it for your needs!

\bibliographystyle{plain}
\bibliography{cite}

\end{document}\documentclass{article}
\usepackage[utf8]{inputenc}
\usepackage{booktabs}
\usepackage{amsmath}
\usepackage{natbib}

\title{My Document}
\author{Your Name}
\date{\today}

\begin{document}

\maketitle

\section{Introduction}
This is a sample document demonstrating various \LaTeX{} features including mathematical equations, tables, and formatting. You can cite papers using the citation search feature, and they will be automatically added to your \texttt{cite.bib} file.

\section{Mathematical Expressions}
Here are some example equations:

\subsection{Quadratic Formula}
The quadratic formula is given by:
\[x = \frac{-b \pm \sqrt{b^2 - 4ac}}{2a}\]

\subsection{Matrix Example}
A sample matrix:
\[
\begin{pmatrix}
a_{11} & a_{12} & a_{13} \\
a_{21} & a_{22} & a_{23} \\
a_{31} & a_{32} & a_{33}
\end{pmatrix}
\]

\section{Tables}
Below is a sample table showing data:

\begin{table}[h]
\centering
\begin{tabular}{@{}lrrr@{}}
\toprule
\textbf{Category} & \textbf{Value 1} & \textbf{Value 2} & \textbf{Total} \\
\midrule
Item A & 10.5 & 20.3 & 30.8 \\
Item B & 15.2 & 25.1 & 40.3 \\
Item C & 22.1 & 12.7 & 34.8 \\
\bottomrule
\end{tabular}
\caption{Sample Data Table}
\label{tab:sample}
\end{table}

\section{Lists}
\subsection{Bullet Points}
\begin{itemize}
  \item First item
  \item Second item
  \item Third item with sub-items:
    \begin{itemize}
      \item Sub-item 1
      \item Sub-item 2
    \end{itemize}
\end{itemize}

\subsection{Numbered List}
\begin{enumerate}
  \item First step
  \item Second step
  \item Third step
\end{enumerate}

\section{Citations and References}
When you use the citation search feature, references will be automatically managed. For example, you might cite a paper like this: \cite{example2024}. All citations will be automatically added to your \texttt{cite.bib} file.

\section{Conclusion}
This template demonstrates basic \LaTeX{} features including mathematical formulas, tables, and lists. The bibliography system is now set up to work seamlessly with Bibby AI's citation management. Feel free to modify it for your needs!

\bibliographystyle{plain}
\bibliography{cite}

\end{document}\documentclass{article}
\usepackage[utf8]{inputenc}
\usepackage{booktabs}
\usepackage{amsmath}
\usepackage{natbib}

\title{My Document}
\author{Your Name}
\date{\today}

\begin{document}

\maketitle

\section{Introduction}
This is a sample document demonstrating various \LaTeX{} features including mathematical equations, tables, and formatting. You can cite papers using the citation search feature, and they will be automatically added to your \texttt{cite.bib} file.

\section{Mathematical Expressions}
Here are some example equations:

\subsection{Quadratic Formula}
The quadratic formula is given by:
\[x = \frac{-b \pm \sqrt{b^2 - 4ac}}{2a}\]

\subsection{Matrix Example}
A sample matrix:
\[
\begin{pmatrix}
a_{11} & a_{12} & a_{13} \\
a_{21} & a_{22} & a_{23} \\
a_{31} & a_{32} & a_{33}
\end{pmatrix}
\]

\section{Tables}
Below is a sample table showing data:

\begin{table}[h]
\centering
\begin{tabular}{@{}lrrr@{}}
\toprule
\textbf{Category} & \textbf{Value 1} & \textbf{Value 2} & \textbf{Total} \\
\midrule
Item A & 10.5 & 20.3 & 30.8 \\
Item B & 15.2 & 25.1 & 40.3 \\
Item C & 22.1 & 12.7 & 34.8 \\
\bottomrule
\end{tabular}
\caption{Sample Data Table}
\label{tab:sample}
\end{table}

\section{Lists}
\subsection{Bullet Points}
\begin{itemize}
  \item First item
  \item Second item
  \item Third item with sub-items:
    \begin{itemize}
      \item Sub-item 1
      \item Sub-item 2
    \end{itemize}
\end{itemize}

\subsection{Numbered List}
\begin{enumerate}
  \item First step
  \item Second step
  \item Third step
\end{enumerate}

\section{Citations and References}
When you use the citation search feature, references will be automatically managed. For example, you might cite a paper like this: \cite{example2024}. All citations will be automatically added to your \texttt{cite.bib} file.

\section{Conclusion}
This template demonstrates basic \LaTeX{} features including mathematical formulas, tables, and lists. The bibliography system is now set up to work seamlessly with Bibby AI's citation management. Feel free to modify it for your needs!

\bibliographystyle{plain}
\bibliography{cite}

\end{document}\documentclass{article}
\usepackage[utf8]{inputenc}
\usepackage{booktabs}
\usepackage{amsmath}
\usepackage{natbib}

\title{My Document}
\author{Your Name}
\date{\today}

\begin{document}

\maketitle

\section{Introduction}
This is a sample document demonstrating various \LaTeX{} features including mathematical equations, tables, and formatting. You can cite papers using the citation search feature, and they will be automatically added to your \texttt{cite.bib} file.

\section{Mathematical Expressions}
Here are some example equations:

\subsection{Quadratic Formula}
The quadratic formula is given by:
\[x = \frac{-b \pm \sqrt{b^2 - 4ac}}{2a}\]

\subsection{Matrix Example}
A sample matrix:
\[
\begin{pmatrix}
a_{11} & a_{12} & a_{13} \\
a_{21} & a_{22} & a_{23} \\
a_{31} & a_{32} & a_{33}
\end{pmatrix}
\]

\section{Tables}
Below is a sample table showing data:

\begin{table}[h]
\centering
\begin{tabular}{@{}lrrr@{}}
\toprule
\textbf{Category} & \textbf{Value 1} & \textbf{Value 2} & \textbf{Total} \\
\midrule
Item A & 10.5 & 20.3 & 30.8 \\
Item B & 15.2 & 25.1 & 40.3 \\
Item C & 22.1 & 12.7 & 34.8 \\
\bottomrule
\end{tabular}
\caption{Sample Data Table}
\label{tab:sample}
\end{table}

\section{Lists}
\subsection{Bullet Points}
\begin{itemize}
  \item First item
  \item Second item
  \item Third item with sub-items:
    \begin{itemize}
      \item Sub-item 1
      \item Sub-item 2
    \end{itemize}
\end{itemize}

\subsection{Numbered List}
\begin{enumerate}
  \item First step
  \item Second step
  \item Third step
\end{enumerate}

\section{Citations and References}
When you use the citation search feature, references will be automatically managed. For example, you might cite a paper like this: \cite{example2024}. All citations will be automatically added to your \texttt{cite.bib} file.

\section{Conclusion}
This template demonstrates basic \LaTeX{} features including mathematical formulas, tables, and lists. The bibliography system is now set up to work seamlessly with Bibby AI's citation management. Feel free to modify it for your needs!

\bibliographystyle{plain}
\bibliography{cite}

\end{document}\documentclass{article}
\usepackage[utf8]{inputenc}
\usepackage{booktabs}
\usepackage{amsmath}
\usepackage{natbib}

\title{My Document}
\author{Your Name}
\date{\today}

\begin{document}

\maketitle

\section{Introduction}
This is a sample document demonstrating various \LaTeX{} features including mathematical equations, tables, and formatting. You can cite papers using the citation search feature, and they will be automatically added to your \texttt{cite.bib} file.

\section{Mathematical Expressions}
Here are some example equations:

\subsection{Quadratic Formula}
The quadratic formula is given by:
\[x = \frac{-b \pm \sqrt{b^2 - 4ac}}{2a}\]

\subsection{Matrix Example}
A sample matrix:
\[
\begin{pmatrix}
a_{11} & a_{12} & a_{13} \\
a_{21} & a_{22} & a_{23} \\
a_{31} & a_{32} & a_{33}
\end{pmatrix}
\]

\section{Tables}
Below is a sample table showing data:

\begin{table}[h]
\centering
\begin{tabular}{@{}lrrr@{}}
\toprule
\textbf{Category} & \textbf{Value 1} & \textbf{Value 2} & \textbf{Total} \\
\midrule
Item A & 10.5 & 20.3 & 30.8 \\
Item B & 15.2 & 25.1 & 40.3 \\
Item C & 22.1 & 12.7 & 34.8 \\
\bottomrule
\end{tabular}
\caption{Sample Data Table}
\label{tab:sample}
\end{table}

\section{Lists}
\subsection{Bullet Points}
\begin{itemize}
  \item First item
  \item Second item
  \item Third item with sub-items:
    \begin{itemize}
      \item Sub-item 1
      \item Sub-item 2
    \end{itemize}
\end{itemize}

\subsection{Numbered List}
\begin{enumerate}
  \item First step
  \item Second step
  \item Third step
\end{enumerate}

\section{Citations and References}
When you use the citation search feature, references will be automatically managed. For example, you might cite a paper like this: \cite{example2024}. All citations will be automatically added to your \texttt{cite.bib} file.

\section{Conclusion}
This template demonstrates basic \LaTeX{} features including mathematical formulas, tables, and lists. The bibliography system is now set up to work seamlessly with Bibby AI's citation management. Feel free to modify it for your needs!

\bibliographystyle{plain}
\bibliography{cite}

\end{document}\documentclass{article}
\usepackage[utf8]{inputenc}
\usepackage{booktabs}
\usepackage{amsmath}
\usepackage{natbib}

\title{My Document}
\author{Your Name}
\date{\today}

\begin{document}

\maketitle

\section{Introduction}
This is a sample document demonstrating various \LaTeX{} features including mathematical equations, tables, and formatting. You can cite papers using the citation search feature, and they will be automatically added to your \texttt{cite.bib} file.

\section{Mathematical Expressions}
Here are some example equations:

\subsection{Quadratic Formula}
The quadratic formula is given by:
\[x = \frac{-b \pm \sqrt{b^2 - 4ac}}{2a}\]

\subsection{Matrix Example}
A sample matrix:
\[
\begin{pmatrix}
a_{11} & a_{12} & a_{13} \\
a_{21} & a_{22} & a_{23} \\
a_{31} & a_{32} & a_{33}
\end{pmatrix}
\]

\section{Tables}
Below is a sample table showing data:

\begin{table}[h]
\centering
\begin{tabular}{@{}lrrr@{}}
\toprule
\textbf{Category} & \textbf{Value 1} & \textbf{Value 2} & \textbf{Total} \\
\midrule
Item A & 10.5 & 20.3 & 30.8 \\
Item B & 15.2 & 25.1 & 40.3 \\
Item C & 22.1 & 12.7 & 34.8 \\
\bottomrule
\end{tabular}
\caption{Sample Data Table}
\label{tab:sample}
\end{table}

\section{Lists}
\subsection{Bullet Points}
\begin{itemize}
  \item First item
  \item Second item
  \item Third item with sub-items:
    \begin{itemize}
      \item Sub-item 1
      \item Sub-item 2
    \end{itemize}
\end{itemize}

\subsection{Numbered List}
\begin{enumerate}
  \item First step
  \item Second step
  \item Third step
\end{enumerate}

\section{Citations and References}
When you use the citation search feature, references will be automatically managed. For example, you might cite a paper like this: \cite{example2024}. All citations will be automatically added to your \texttt{cite.bib} file.

\section{Conclusion}
This template demonstrates basic \LaTeX{} features including mathematical formulas, tables, and lists. The bibliography system is now set up to work seamlessly with Bibby AI's citation management. Feel free to modify it for your needs!

\bibliographystyle{plain}
\bibliography{cite}

\end{document}\documentclass{article}
\usepackage[utf8]{inputenc}
\usepackage{booktabs}
\usepackage{amsmath}
\usepackage{natbib}

\title{My Document}
\author{Your Name}
\date{\today}

\begin{document}

\maketitle

\section{Introduction}
This is a sample document demonstrating various \LaTeX{} features including mathematical equations, tables, and formatting. You can cite papers using the citation search feature, and they will be automatically added to your \texttt{cite.bib} file.

\section{Mathematical Expressions}
Here are some example equations:

\subsection{Quadratic Formula}
The quadratic formula is given by:
\[x = \frac{-b \pm \sqrt{b^2 - 4ac}}{2a}\]

\subsection{Matrix Example}
A sample matrix:
\[
\begin{pmatrix}
a_{11} & a_{12} & a_{13} \\
a_{21} & a_{22} & a_{23} \\
a_{31} & a_{32} & a_{33}
\end{pmatrix}
\]

\section{Tables}
Below is a sample table showing data:

\begin{table}[h]
\centering
\begin{tabular}{@{}lrrr@{}}
\toprule
\textbf{Category} & \textbf{Value 1} & \textbf{Value 2} & \textbf{Total} \\
\midrule
Item A & 10.5 & 20.3 & 30.8 \\
Item B & 15.2 & 25.1 & 40.3 \\
Item C & 22.1 & 12.7 & 34.8 \\
\bottomrule
\end{tabular}
\caption{Sample Data Table}
\label{tab:sample}
\end{table}

\section{Lists}
\subsection{Bullet Points}
\begin{itemize}
  \item First item
  \item Second item
  \item Third item with sub-items:
    \begin{itemize}
      \item Sub-item 1
      \item Sub-item 2
    \end{itemize}
\end{itemize}

\subsection{Numbered List}
\begin{enumerate}
  \item First step
  \item Second step
  \item Third step
\end{enumerate}

\section{Citations and References}
When you use the citation search feature, references will be automatically managed. For example, you might cite a paper like this: \cite{example2024}. All citations will be automatically added to your \texttt{cite.bib} file.

\section{Conclusion}
This template demonstrates basic \LaTeX{} features including mathematical formulas, tables, and lists. The bibliography system is now set up to work seamlessly with Bibby AI's citation management. Feel free to modify it for your needs!

\bibliographystyle{plain}
\bibliography{cite}

\end{document}\documentclass{article}
\usepackage[utf8]{inputenc}
\usepackage{booktabs}
\usepackage{amsmath}
\usepackage{natbib}

\title{My Document}
\author{Your Name}
\date{\today}

\begin{document}

\maketitle

\section{Introduction}
This is a sample document demonstrating various \LaTeX{} features including mathematical equations, tables, and formatting. You can cite papers using the citation search feature, and they will be automatically added to your \texttt{cite.bib} file.

\section{Mathematical Expressions}
Here are some example equations:

\subsection{Quadratic Formula}
The quadratic formula is given by:
\[x = \frac{-b \pm \sqrt{b^2 - 4ac}}{2a}\]

\subsection{Matrix Example}
A sample matrix:
\[
\begin{pmatrix}
a_{11} & a_{12} & a_{13} \\
a_{21} & a_{22} & a_{23} \\
a_{31} & a_{32} & a_{33}
\end{pmatrix}
\]

\section{Tables}
Below is a sample table showing data:

\begin{table}[h]
\centering
\begin{tabular}{@{}lrrr@{}}
\toprule
\textbf{Category} & \textbf{Value 1} & \textbf{Value 2} & \textbf{Total} \\
\midrule
Item A & 10.5 & 20.3 & 30.8 \\
Item B & 15.2 & 25.1 & 40.3 \\
Item C & 22.1 & 12.7 & 34.8 \\
\bottomrule
\end{tabular}
\caption{Sample Data Table}
\label{tab:sample}
\end{table}

\section{Lists}
\subsection{Bullet Points}
\begin{itemize}
  \item First item
  \item Second item
  \item Third item with sub-items:
    \begin{itemize}
      \item Sub-item 1
      \item Sub-item 2
    \end{itemize}
\end{itemize}

\subsection{Numbered List}
\begin{enumerate}
  \item First step
  \item Second step
  \item Third step
\end{enumerate}

\section{Citations and References}
When you use the citation search feature, references will be automatically managed. For example, you might cite a paper like this: \cite{example2024}. All citations will be automatically added to your \texttt{cite.bib} file.

\section{Conclusion}
This template demonstrates basic \LaTeX{} features including mathematical formulas, tables, and lists. The bibliography system is now set up to work seamlessly with Bibby AI's citation management. Feel free to modify it for your needs!

\bibliographystyle{plain}
\bibliography{cite}

\end{document}\documentclass{article}
\usepackage[utf8]{inputenc}
\usepackage{booktabs}
\usepackage{amsmath}
\usepackage{natbib}

\title{My Document}
\author{Your Name}
\date{\today}

\begin{document}

\maketitle

\section{Introduction}
This is a sample document demonstrating various \LaTeX{} features including mathematical equations, tables, and formatting. You can cite papers using the citation search feature, and they will be automatically added to your \texttt{cite.bib} file.

\section{Mathematical Expressions}
Here are some example equations:

\subsection{Quadratic Formula}
The quadratic formula is given by:
\[x = \frac{-b \pm \sqrt{b^2 - 4ac}}{2a}\]

\subsection{Matrix Example}
A sample matrix:
\[
\begin{pmatrix}
a_{11} & a_{12} & a_{13} \\
a_{21} & a_{22} & a_{23} \\
a_{31} & a_{32} & a_{33}
\end{pmatrix}
\]

\section{Tables}
Below is a sample table showing data:

\begin{table}[h]
\centering
\begin{tabular}{@{}lrrr@{}}
\toprule
\textbf{Category} & \textbf{Value 1} & \textbf{Value 2} & \textbf{Total} \\
\midrule
Item A & 10.5 & 20.3 & 30.8 \\
Item B & 15.2 & 25.1 & 40.3 \\
Item C & 22.1 & 12.7 & 34.8 \\
\bottomrule
\end{tabular}
\caption{Sample Data Table}
\label{tab:sample}
\end{table}

\section{Lists}
\subsection{Bullet Points}
\begin{itemize}
  \item First item
  \item Second item
  \item Third item with sub-items:
    \begin{itemize}
      \item Sub-item 1
      \item Sub-item 2
    \end{itemize}
\end{itemize}

\subsection{Numbered List}
\begin{enumerate}
  \item First step
  \item Second step
  \item Third step
\end{enumerate}

\section{Citations and References}
When you use the citation search feature, references will be automatically managed. For example, you might cite a paper like this: \cite{example2024}. All citations will be automatically added to your \texttt{cite.bib} file.

\section{Conclusion}
This template demonstrates basic \LaTeX{} features including mathematical formulas, tables, and lists. The bibliography system is now set up to work seamlessly with Bibby AI's citation management. Feel free to modify it for your needs!

\bibliographystyle{plain}
\bibliography{cite}

\end{document}\documentclass{article}
\usepackage[utf8]{inputenc}
\usepackage{booktabs}
\usepackage{amsmath}
\usepackage{natbib}

\title{My Document}
\author{Your Name}
\date{\today}

\begin{document}

\maketitle

\section{Introduction}
This is a sample document demonstrating various \LaTeX{} features including mathematical equations, tables, and formatting. You can cite papers using the citation search feature, and they will be automatically added to your \texttt{cite.bib} file.

\section{Mathematical Expressions}
Here are some example equations:

\subsection{Quadratic Formula}
The quadratic formula is given by:
\[x = \frac{-b \pm \sqrt{b^2 - 4ac}}{2a}\]

\subsection{Matrix Example}
A sample matrix:
\[
\begin{pmatrix}
a_{11} & a_{12} & a_{13} \\
a_{21} & a_{22} & a_{23} \\
a_{31} & a_{32} & a_{33}
\end{pmatrix}
\]

\section{Tables}
Below is a sample table showing data:

\begin{table}[h]
\centering
\begin{tabular}{@{}lrrr@{}}
\toprule
\textbf{Category} & \textbf{Value 1} & \textbf{Value 2} & \textbf{Total} \\
\midrule
Item A & 10.5 & 20.3 & 30.8 \\
Item B & 15.2 & 25.1 & 40.3 \\
Item C & 22.1 & 12.7 & 34.8 \\
\bottomrule
\end{tabular}
\caption{Sample Data Table}
\label{tab:sample}
\end{table}

\section{Lists}
\subsection{Bullet Points}
\begin{itemize}
  \item First item
  \item Second item
  \item Third item with sub-items:
    \begin{itemize}
      \item Sub-item 1
      \item Sub-item 2
    \end{itemize}
\end{itemize}

\subsection{Numbered List}
\begin{enumerate}
  \item First step
  \item Second step
  \item Third step
\end{enumerate}

\section{Citations and References}
When you use the citation search feature, references will be automatically managed. For example, you might cite a paper like this: \cite{example2024}. All citations will be automatically added to your \texttt{cite.bib} file.

\section{Conclusion}
This template demonstrates basic \LaTeX{} features including mathematical formulas, tables, and lists. The bibliography system is now set up to work seamlessly with Bibby AI's citation management. Feel free to modify it for your needs!

\bibliographystyle{plain}
\bibliography{cite}

\end{document}\documentclass{article}
\usepackage[utf8]{inputenc}
\usepackage{booktabs}
\usepackage{amsmath}
\usepackage{natbib}

\title{My Document}
\author{Your Name}
\date{\today}

\begin{document}

\maketitle

\section{Introduction}
This is a sample document demonstrating various \LaTeX{} features including mathematical equations, tables, and formatting. You can cite papers using the citation search feature, and they will be automatically added to your \texttt{cite.bib} file.

\section{Mathematical Expressions}
Here are some example equations:

\subsection{Quadratic Formula}
The quadratic formula is given by:
\[x = \frac{-b \pm \sqrt{b^2 - 4ac}}{2a}\]

\subsection{Matrix Example}
A sample matrix:
\[
\begin{pmatrix}
a_{11} & a_{12} & a_{13} \\
a_{21} & a_{22} & a_{23} \\
a_{31} & a_{32} & a_{33}
\end{pmatrix}
\]

\section{Tables}
Below is a sample table showing data:

\begin{table}[h]
\centering
\begin{tabular}{@{}lrrr@{}}
\toprule
\textbf{Category} & \textbf{Value 1} & \textbf{Value 2} & \textbf{Total} \\
\midrule
Item A & 10.5 & 20.3 & 30.8 \\
Item B & 15.2 & 25.1 & 40.3 \\
Item C & 22.1 & 12.7 & 34.8 \\
\bottomrule
\end{tabular}
\caption{Sample Data Table}
\label{tab:sample}
\end{table}

\section{Lists}
\subsection{Bullet Points}
\begin{itemize}
  \item First item
  \item Second item
  \item Third item with sub-items:
    \begin{itemize}
      \item Sub-item 1
      \item Sub-item 2
    \end{itemize}
\end{itemize}

\subsection{Numbered List}
\begin{enumerate}
  \item First step
  \item Second step
  \item Third step
\end{enumerate}

\section{Citations and References}
When you use the citation search feature, references will be automatically managed. For example, you might cite a paper like this: \cite{example2024}. All citations will be automatically added to your \texttt{cite.bib} file.

\section{Conclusion}
This template demonstrates basic \LaTeX{} features including mathematical formulas, tables, and lists. The bibliography system is now set up to work seamlessly with Bibby AI's citation management. Feel free to modify it for your needs!

\bibliographystyle{plain}
\bibliography{cite}

\end{document}\documentclass{article}
\usepackage[utf8]{inputenc}
\usepackage{booktabs}
\usepackage{amsmath}
\usepackage{natbib}

\title{My Document}
\author{Your Name}
\date{\today}

\begin{document}

\maketitle

\section{Introduction}
This is a sample document demonstrating various \LaTeX{} features including mathematical equations, tables, and formatting. You can cite papers using the citation search feature, and they will be automatically added to your \texttt{cite.bib} file.

\section{Mathematical Expressions}
Here are some example equations:

\subsection{Quadratic Formula}
The quadratic formula is given by:
\[x = \frac{-b \pm \sqrt{b^2 - 4ac}}{2a}\]

\subsection{Matrix Example}
A sample matrix:
\[
\begin{pmatrix}
a_{11} & a_{12} & a_{13} \\
a_{21} & a_{22} & a_{23} \\
a_{31} & a_{32} & a_{33}
\end{pmatrix}
\]

\section{Tables}
Below is a sample table showing data:

\begin{table}[h]
\centering
\begin{tabular}{@{}lrrr@{}}
\toprule
\textbf{Category} & \textbf{Value 1} & \textbf{Value 2} & \textbf{Total} \\
\midrule
Item A & 10.5 & 20.3 & 30.8 \\
Item B & 15.2 & 25.1 & 40.3 \\
Item C & 22.1 & 12.7 & 34.8 \\
\bottomrule
\end{tabular}
\caption{Sample Data Table}
\label{tab:sample}
\end{table}

\section{Lists}
\subsection{Bullet Points}
\begin{itemize}
  \item First item
  \item Second item
  \item Third item with sub-items:
    \begin{itemize}
      \item Sub-item 1
      \item Sub-item 2
    \end{itemize}
\end{itemize}

\subsection{Numbered List}
\begin{enumerate}
  \item First step
  \item Second step
  \item Third step
\end{enumerate}

\section{Citations and References}
When you use the citation search feature, references will be automatically managed. For example, you might cite a paper like this: \cite{example2024}. All citations will be automatically added to your \texttt{cite.bib} file.

\section{Conclusion}
This template demonstrates basic \LaTeX{} features including mathematical formulas, tables, and lists. The bibliography system is now set up to work seamlessly with Bibby AI's citation management. Feel free to modify it for your needs!

\bibliographystyle{plain}
\bibliography{cite}

\end{document}\documentclass{article}
\usepackage[utf8]{inputenc}
\usepackage{booktabs}
\usepackage{amsmath}
\usepackage{natbib}

\title{My Document}
\author{Your Name}
\date{\today}

\begin{document}

\maketitle

\section{Introduction}
This is a sample document demonstrating various \LaTeX{} features including mathematical equations, tables, and formatting. You can cite papers using the citation search feature, and they will be automatically added to your \texttt{cite.bib} file.

\section{Mathematical Expressions}
Here are some example equations:

\subsection{Quadratic Formula}
The quadratic formula is given by:
\[x = \frac{-b \pm \sqrt{b^2 - 4ac}}{2a}\]

\subsection{Matrix Example}
A sample matrix:
\[
\begin{pmatrix}
a_{11} & a_{12} & a_{13} \\
a_{21} & a_{22} & a_{23} \\
a_{31} & a_{32} & a_{33}
\end{pmatrix}
\]

\section{Tables}
Below is a sample table showing data:

\begin{table}[h]
\centering
\begin{tabular}{@{}lrrr@{}}
\toprule
\textbf{Category} & \textbf{Value 1} & \textbf{Value 2} & \textbf{Total} \\
\midrule
Item A & 10.5 & 20.3 & 30.8 \\
Item B & 15.2 & 25.1 & 40.3 \\
Item C & 22.1 & 12.7 & 34.8 \\
\bottomrule
\end{tabular}
\caption{Sample Data Table}
\label{tab:sample}
\end{table}

\section{Lists}
\subsection{Bullet Points}
\begin{itemize}
  \item First item
  \item Second item
  \item Third item with sub-items:
    \begin{itemize}
      \item Sub-item 1
      \item Sub-item 2
    \end{itemize}
\end{itemize}

\subsection{Numbered List}
\begin{enumerate}
  \item First step
  \item Second step
  \item Third step
\end{enumerate}

\section{Citations and References}
When you use the citation search feature, references will be automatically managed. For example, you might cite a paper like this: \cite{example2024}. All citations will be automatically added to your \texttt{cite.bib} file.

\section{Conclusion}
This template demonstrates basic \LaTeX{} features including mathematical formulas, tables, and lists. The bibliography system is now set up to work seamlessly with Bibby AI's citation management. Feel free to modify it for your needs!

\bibliographystyle{plain}
\bibliography{cite}

\end{document}\documentclass{article}
\usepackage[utf8]{inputenc}
\usepackage{booktabs}
\usepackage{amsmath}
\usepackage{natbib}

\title{My Document}
\author{Your Name}
\date{\today}

\begin{document}

\maketitle

\section{Introduction}
This is a sample document demonstrating various \LaTeX{} features including mathematical equations, tables, and formatting. You can cite papers using the citation search feature, and they will be automatically added to your \texttt{cite.bib} file.

\section{Mathematical Expressions}
Here are some example equations:

\subsection{Quadratic Formula}
The quadratic formula is given by:
\[x = \frac{-b \pm \sqrt{b^2 - 4ac}}{2a}\]

\subsection{Matrix Example}
A sample matrix:
\[
\begin{pmatrix}
a_{11} & a_{12} & a_{13} \\
a_{21} & a_{22} & a_{23} \\
a_{31} & a_{32} & a_{33}
\end{pmatrix}
\]

\section{Tables}
Below is a sample table showing data:

\begin{table}[h]
\centering
\begin{tabular}{@{}lrrr@{}}
\toprule
\textbf{Category} & \textbf{Value 1} & \textbf{Value 2} & \textbf{Total} \\
\midrule
Item A & 10.5 & 20.3 & 30.8 \\
Item B & 15.2 & 25.1 & 40.3 \\
Item C & 22.1 & 12.7 & 34.8 \\
\bottomrule
\end{tabular}
\caption{Sample Data Table}
\label{tab:sample}
\end{table}

\section{Lists}
\subsection{Bullet Points}
\begin{itemize}
  \item First item
  \item Second item
  \item Third item with sub-items:
    \begin{itemize}
      \item Sub-item 1
      \item Sub-item 2
    \end{itemize}
\end{itemize}

\subsection{Numbered List}
\begin{enumerate}
  \item First step
  \item Second step
  \item Third step
\end{enumerate}

\section{Citations and References}
When you use the citation search feature, references will be automatically managed. For example, you might cite a paper like this: \cite{example2024}. All citations will be automatically added to your \texttt{cite.bib} file.

\section{Conclusion}
This template demonstrates basic \LaTeX{} features including mathematical formulas, tables, and lists. The bibliography system is now set up to work seamlessly with Bibby AI's citation management. Feel free to modify it for your needs!

\bibliographystyle{plain}
\bibliography{cite}

\end{document}\documentclass{article}
\usepackage[utf8]{inputenc}
\usepackage{booktabs}
\usepackage{amsmath}
\usepackage{natbib}

\title{My Document}
\author{Your Name}
\date{\today}

\begin{document}

\maketitle

\section{Introduction}
This is a sample document demonstrating various \LaTeX{} features including mathematical equations, tables, and formatting. You can cite papers using the citation search feature, and they will be automatically added to your \texttt{cite.bib} file.

\section{Mathematical Expressions}
Here are some example equations:

\subsection{Quadratic Formula}
The quadratic formula is given by:
\[x = \frac{-b \pm \sqrt{b^2 - 4ac}}{2a}\]

\subsection{Matrix Example}
A sample matrix:
\[
\begin{pmatrix}
a_{11} & a_{12} & a_{13} \\
a_{21} & a_{22} & a_{23} \\
a_{31} & a_{32} & a_{33}
\end{pmatrix}
\]

\section{Tables}
Below is a sample table showing data:

\begin{table}[h]
\centering
\begin{tabular}{@{}lrrr@{}}
\toprule
\textbf{Category} & \textbf{Value 1} & \textbf{Value 2} & \textbf{Total} \\
\midrule
Item A & 10.5 & 20.3 & 30.8 \\
Item B & 15.2 & 25.1 & 40.3 \\
Item C & 22.1 & 12.7 & 34.8 \\
\bottomrule
\end{tabular}
\caption{Sample Data Table}
\label{tab:sample}
\end{table}

\section{Lists}
\subsection{Bullet Points}
\begin{itemize}
  \item First item
  \item Second item
  \item Third item with sub-items:
    \begin{itemize}
      \item Sub-item 1
      \item Sub-item 2
    \end{itemize}
\end{itemize}

\subsection{Numbered List}
\begin{enumerate}
  \item First step
  \item Second step
  \item Third step
\end{enumerate}

\section{Citations and References}
When you use the citation search feature, references will be automatically managed. For example, you might cite a paper like this: \cite{example2024}. All citations will be automatically added to your \texttt{cite.bib} file.

\section{Conclusion}
This template demonstrates basic \LaTeX{} features including mathematical formulas, tables, and lists. The bibliography system is now set up to work seamlessly with Bibby AI's citation management. Feel free to modify it for your needs!

\bibliographystyle{plain}
\bibliography{cite}

\end{document}\documentclass{article}
\usepackage[utf8]{inputenc}
\usepackage{booktabs}
\usepackage{amsmath}
\usepackage{natbib}

\title{My Document}
\author{Your Name}
\date{\today}

\begin{document}

\maketitle

\section{Introduction}
This is a sample document demonstrating various \LaTeX{} features including mathematical equations, tables, and formatting. You can cite papers using the citation search feature, and they will be automatically added to your \texttt{cite.bib} file.

\section{Mathematical Expressions}
Here are some example equations:

\subsection{Quadratic Formula}
The quadratic formula is given by:
\[x = \frac{-b \pm \sqrt{b^2 - 4ac}}{2a}\]

\subsection{Matrix Example}
A sample matrix:
\[
\begin{pmatrix}
a_{11} & a_{12} & a_{13} \\
a_{21} & a_{22} & a_{23} \\
a_{31} & a_{32} & a_{33}
\end{pmatrix}
\]

\section{Tables}
Below is a sample table showing data:

\begin{table}[h]
\centering
\begin{tabular}{@{}lrrr@{}}
\toprule
\textbf{Category} & \textbf{Value 1} & \textbf{Value 2} & \textbf{Total} \\
\midrule
Item A & 10.5 & 20.3 & 30.8 \\
Item B & 15.2 & 25.1 & 40.3 \\
Item C & 22.1 & 12.7 & 34.8 \\
\bottomrule
\end{tabular}
\caption{Sample Data Table}
\label{tab:sample}
\end{table}

\section{Lists}
\subsection{Bullet Points}
\begin{itemize}
  \item First item
  \item Second item
  \item Third item with sub-items:
    \begin{itemize}
      \item Sub-item 1
      \item Sub-item 2
    \end{itemize}
\end{itemize}

\subsection{Numbered List}
\begin{enumerate}
  \item First step
  \item Second step
  \item Third step
\end{enumerate}

\section{Citations and References}
When you use the citation search feature, references will be automatically managed. For example, you might cite a paper like this: \cite{example2024}. All citations will be automatically added to your \texttt{cite.bib} file.

\section{Conclusion}
This template demonstrates basic \LaTeX{} features including mathematical formulas, tables, and lists. The bibliography system is now set up to work seamlessly with Bibby AI's citation management. Feel free to modify it for your needs!

\bibliographystyle{plain}
\bibliography{cite}

\end{document}\documentclass{article}
\usepackage[utf8]{inputenc}
\usepackage{booktabs}
\usepackage{amsmath}
\usepackage{natbib}

\title{My Document}
\author{Your Name}
\date{\today}

\begin{document}

\maketitle

\section{Introduction}
This is a sample document demonstrating various \LaTeX{} features including mathematical equations, tables, and formatting. You can cite papers using the citation search feature, and they will be automatically added to your \texttt{cite.bib} file.

\section{Mathematical Expressions}
Here are some example equations:

\subsection{Quadratic Formula}
The quadratic formula is given by:
\[x = \frac{-b \pm \sqrt{b^2 - 4ac}}{2a}\]

\subsection{Matrix Example}
A sample matrix:
\[
\begin{pmatrix}
a_{11} & a_{12} & a_{13} \\
a_{21} & a_{22} & a_{23} \\
a_{31} & a_{32} & a_{33}
\end{pmatrix}
\]

\section{Tables}
Below is a sample table showing data:

\begin{table}[h]
\centering
\begin{tabular}{@{}lrrr@{}}
\toprule
\textbf{Category} & \textbf{Value 1} & \textbf{Value 2} & \textbf{Total} \\
\midrule
Item A & 10.5 & 20.3 & 30.8 \\
Item B & 15.2 & 25.1 & 40.3 \\
Item C & 22.1 & 12.7 & 34.8 \\
\bottomrule
\end{tabular}
\caption{Sample Data Table}
\label{tab:sample}
\end{table}

\section{Lists}
\subsection{Bullet Points}
\begin{itemize}
  \item First item
  \item Second item
  \item Third item with sub-items:
    \begin{itemize}
      \item Sub-item 1
      \item Sub-item 2
    \end{itemize}
\end{itemize}

\subsection{Numbered List}
\begin{enumerate}
  \item First step
  \item Second step
  \item Third step
\end{enumerate}

\section{Citations and References}
When you use the citation search feature, references will be automatically managed. For example, you might cite a paper like this: \cite{example2024}. All citations will be automatically added to your \texttt{cite.bib} file.

\section{Conclusion}
This template demonstrates basic \LaTeX{} features including mathematical formulas, tables, and lists. The bibliography system is now set up to work seamlessly with Bibby AI's citation management. Feel free to modify it for your needs!

\bibliographystyle{plain}
\bibliography{cite}

\end{document}\documentclass{article}
\usepackage[utf8]{inputenc}
\usepackage{booktabs}
\usepackage{amsmath}
\usepackage{natbib}

\title{My Document}
\author{Your Name}
\date{\today}

\begin{document}

\maketitle

\section{Introduction}
This is a sample document demonstrating various \LaTeX{} features including mathematical equations, tables, and formatting. You can cite papers using the citation search feature, and they will be automatically added to your \texttt{cite.bib} file.

\section{Mathematical Expressions}
Here are some example equations:

\subsection{Quadratic Formula}
The quadratic formula is given by:
\[x = \frac{-b \pm \sqrt{b^2 - 4ac}}{2a}\]

\subsection{Matrix Example}
A sample matrix:
\[
\begin{pmatrix}
a_{11} & a_{12} & a_{13} \\
a_{21} & a_{22} & a_{23} \\
a_{31} & a_{32} & a_{33}
\end{pmatrix}
\]

\section{Tables}
Below is a sample table showing data:

\begin{table}[h]
\centering
\begin{tabular}{@{}lrrr@{}}
\toprule
\textbf{Category} & \textbf{Value 1} & \textbf{Value 2} & \textbf{Total} \\
\midrule
Item A & 10.5 & 20.3 & 30.8 \\
Item B & 15.2 & 25.1 & 40.3 \\
Item C & 22.1 & 12.7 & 34.8 \\
\bottomrule
\end{tabular}
\caption{Sample Data Table}
\label{tab:sample}
\end{table}

\section{Lists}
\subsection{Bullet Points}
\begin{itemize}
  \item First item
  \item Second item
  \item Third item with sub-items:
    \begin{itemize}
      \item Sub-item 1
      \item Sub-item 2
    \end{itemize}
\end{itemize}

\subsection{Numbered List}
\begin{enumerate}
  \item First step
  \item Second step
  \item Third step
\end{enumerate}

\section{Citations and References}
When you use the citation search feature, references will be automatically managed. For example, you might cite a paper like this: \cite{example2024}. All citations will be automatically added to your \texttt{cite.bib} file.

\section{Conclusion}
This template demonstrates basic \LaTeX{} features including mathematical formulas, tables, and lists. The bibliography system is now set up to work seamlessly with Bibby AI's citation management. Feel free to modify it for your needs!

\bibliographystyle{plain}
\bibliography{cite}

\end{document}\documentclass{article}
\usepackage[utf8]{inputenc}
\usepackage{booktabs}
\usepackage{amsmath}
\usepackage{natbib}

\title{My Document}
\author{Your Name}
\date{\today}

\begin{document}

\maketitle

\section{Introduction}
This is a sample document demonstrating various \LaTeX{} features including mathematical equations, tables, and formatting. You can cite papers using the citation search feature, and they will be automatically added to your \texttt{cite.bib} file.

\section{Mathematical Expressions}
Here are some example equations:

\subsection{Quadratic Formula}
The quadratic formula is given by:
\[x = \frac{-b \pm \sqrt{b^2 - 4ac}}{2a}\]

\subsection{Matrix Example}
A sample matrix:
\[
\begin{pmatrix}
a_{11} & a_{12} & a_{13} \\
a_{21} & a_{22} & a_{23} \\
a_{31} & a_{32} & a_{33}
\end{pmatrix}
\]

\section{Tables}
Below is a sample table showing data:

\begin{table}[h]
\centering
\begin{tabular}{@{}lrrr@{}}
\toprule
\textbf{Category} & \textbf{Value 1} & \textbf{Value 2} & \textbf{Total} \\
\midrule
Item A & 10.5 & 20.3 & 30.8 \\
Item B & 15.2 & 25.1 & 40.3 \\
Item C & 22.1 & 12.7 & 34.8 \\
\bottomrule
\end{tabular}
\caption{Sample Data Table}
\label{tab:sample}
\end{table}

\section{Lists}
\subsection{Bullet Points}
\begin{itemize}
  \item First item
  \item Second item
  \item Third item with sub-items:
    \begin{itemize}
      \item Sub-item 1
      \item Sub-item 2
    \end{itemize}
\end{itemize}

\subsection{Numbered List}
\begin{enumerate}
  \item First step
  \item Second step
  \item Third step
\end{enumerate}

\section{Citations and References}
When you use the citation search feature, references will be automatically managed. For example, you might cite a paper like this: \cite{example2024}. All citations will be automatically added to your \texttt{cite.bib} file.

\section{Conclusion}
This template demonstrates basic \LaTeX{} features including mathematical formulas, tables, and lists. The bibliography system is now set up to work seamlessly with Bibby AI's citation management. Feel free to modify it for your needs!

\bibliographystyle{plain}
\bibliography{cite}

\end{document}\documentclass{article}
\usepackage[utf8]{inputenc}
\usepackage{booktabs}
\usepackage{amsmath}
\usepackage{natbib}

\title{My Document}
\author{Your Name}
\date{\today}

\begin{document}

\maketitle

\section{Introduction}
This is a sample document demonstrating various \LaTeX{} features including mathematical equations, tables, and formatting. You can cite papers using the citation search feature, and they will be automatically added to your \texttt{cite.bib} file.

\section{Mathematical Expressions}
Here are some example equations:

\subsection{Quadratic Formula}
The quadratic formula is given by:
\[x = \frac{-b \pm \sqrt{b^2 - 4ac}}{2a}\]

\subsection{Matrix Example}
A sample matrix:
\[
\begin{pmatrix}
a_{11} & a_{12} & a_{13} \\
a_{21} & a_{22} & a_{23} \\
a_{31} & a_{32} & a_{33}
\end{pmatrix}
\]

\section{Tables}
Below is a sample table showing data:

\begin{table}[h]
\centering
\begin{tabular}{@{}lrrr@{}}
\toprule
\textbf{Category} & \textbf{Value 1} & \textbf{Value 2} & \textbf{Total} \\
\midrule
Item A & 10.5 & 20.3 & 30.8 \\
Item B & 15.2 & 25.1 & 40.3 \\
Item C & 22.1 & 12.7 & 34.8 \\
\bottomrule
\end{tabular}
\caption{Sample Data Table}
\label{tab:sample}
\end{table}

\section{Lists}
\subsection{Bullet Points}
\begin{itemize}
  \item First item
  \item Second item
  \item Third item with sub-items:
    \begin{itemize}
      \item Sub-item 1
      \item Sub-item 2
    \end{itemize}
\end{itemize}

\subsection{Numbered List}
\begin{enumerate}
  \item First step
  \item Second step
  \item Third step
\end{enumerate}

\section{Citations and References}
When you use the citation search feature, references will be automatically managed. For example, you might cite a paper like this: \cite{example2024}. All citations will be automatically added to your \texttt{cite.bib} file.

\section{Conclusion}
This template demonstrates basic \LaTeX{} features including mathematical formulas, tables, and lists. The bibliography system is now set up to work seamlessly with Bibby AI's citation management. Feel free to modify it for your needs!

\bibliographystyle{plain}
\bibliography{cite}

\end{document}\documentclass{article}
\usepackage[utf8]{inputenc}
\usepackage{booktabs}
\usepackage{amsmath}
\usepackage{natbib}

\title{My Document}
\author{Your Name}
\date{\today}

\begin{document}

\maketitle

\section{Introduction}
This is a sample document demonstrating various \LaTeX{} features including mathematical equations, tables, and formatting. You can cite papers using the citation search feature, and they will be automatically added to your \texttt{cite.bib} file.

\section{Mathematical Expressions}
Here are some example equations:

\subsection{Quadratic Formula}
The quadratic formula is given by:
\[x = \frac{-b \pm \sqrt{b^2 - 4ac}}{2a}\]

\subsection{Matrix Example}
A sample matrix:
\[
\begin{pmatrix}
a_{11} & a_{12} & a_{13} \\
a_{21} & a_{22} & a_{23} \\
a_{31} & a_{32} & a_{33}
\end{pmatrix}
\]

\section{Tables}
Below is a sample table showing data:

\begin{table}[h]
\centering
\begin{tabular}{@{}lrrr@{}}
\toprule
\textbf{Category} & \textbf{Value 1} & \textbf{Value 2} & \textbf{Total} \\
\midrule
Item A & 10.5 & 20.3 & 30.8 \\
Item B & 15.2 & 25.1 & 40.3 \\
Item C & 22.1 & 12.7 & 34.8 \\
\bottomrule
\end{tabular}
\caption{Sample Data Table}
\label{tab:sample}
\end{table}

\section{Lists}
\subsection{Bullet Points}
\begin{itemize}
  \item First item
  \item Second item
  \item Third item with sub-items:
    \begin{itemize}
      \item Sub-item 1
      \item Sub-item 2
    \end{itemize}
\end{itemize}

\subsection{Numbered List}
\begin{enumerate}
  \item First step
  \item Second step
  \item Third step
\end{enumerate}

\section{Citations and References}
When you use the citation search feature, references will be automatically managed. For example, you might cite a paper like this: \cite{example2024}. All citations will be automatically added to your \texttt{cite.bib} file.

\section{Conclusion}
This template demonstrates basic \LaTeX{} features including mathematical formulas, tables, and lists. The bibliography system is now set up to work seamlessly with Bibby AI's citation management. Feel free to modify it for your needs!

\bibliographystyle{plain}
\bibliography{cite}

\end{document}\documentclass{article}
\usepackage[utf8]{inputenc}
\usepackage{booktabs}
\usepackage{amsmath}
\usepackage{natbib}

\title{My Document}
\author{Your Name}
\date{\today}

\begin{document}

\maketitle

\section{Introduction}
This is a sample document demonstrating various \LaTeX{} features including mathematical equations, tables, and formatting. You can cite papers using the citation search feature, and they will be automatically added to your \texttt{cite.bib} file.

\section{Mathematical Expressions}
Here are some example equations:

\subsection{Quadratic Formula}
The quadratic formula is given by:
\[x = \frac{-b \pm \sqrt{b^2 - 4ac}}{2a}\]

\subsection{Matrix Example}
A sample matrix:
\[
\begin{pmatrix}
a_{11} & a_{12} & a_{13} \\
a_{21} & a_{22} & a_{23} \\
a_{31} & a_{32} & a_{33}
\end{pmatrix}
\]

\section{Tables}
Below is a sample table showing data:

\begin{table}[h]
\centering
\begin{tabular}{@{}lrrr@{}}
\toprule
\textbf{Category} & \textbf{Value 1} & \textbf{Value 2} & \textbf{Total} \\
\midrule
Item A & 10.5 & 20.3 & 30.8 \\
Item B & 15.2 & 25.1 & 40.3 \\
Item C & 22.1 & 12.7 & 34.8 \\
\bottomrule
\end{tabular}
\caption{Sample Data Table}
\label{tab:sample}
\end{table}

\section{Lists}
\subsection{Bullet Points}
\begin{itemize}
  \item First item
  \item Second item
  \item Third item with sub-items:
    \begin{itemize}
      \item Sub-item 1
      \item Sub-item 2
    \end{itemize}
\end{itemize}

\subsection{Numbered List}
\begin{enumerate}
  \item First step
  \item Second step
  \item Third step
\end{enumerate}

\section{Citations and References}
When you use the citation search feature, references will be automatically managed. For example, you might cite a paper like this: \cite{example2024}. All citations will be automatically added to your \texttt{cite.bib} file.

\section{Conclusion}
This template demonstrates basic \LaTeX{} features including mathematical formulas, tables, and lists. The bibliography system is now set up to work seamlessly with Bibby AI's citation management. Feel free to modify it for your needs!

\bibliographystyle{plain}
\bibliography{cite}

\end{document}\documentclass{article}
\usepackage[utf8]{inputenc}
\usepackage{booktabs}
\usepackage{amsmath}
\usepackage{natbib}

\title{My Document}
\author{Your Name}
\date{\today}

\begin{document}

\maketitle

\section{Introduction}
This is a sample document demonstrating various \LaTeX{} features including mathematical equations, tables, and formatting. You can cite papers using the citation search feature, and they will be automatically added to your \texttt{cite.bib} file.

\section{Mathematical Expressions}
Here are some example equations:

\subsection{Quadratic Formula}
The quadratic formula is given by:
\[x = \frac{-b \pm \sqrt{b^2 - 4ac}}{2a}\]

\subsection{Matrix Example}
A sample matrix:
\[
\begin{pmatrix}
a_{11} & a_{12} & a_{13} \\
a_{21} & a_{22} & a_{23} \\
a_{31} & a_{32} & a_{33}
\end{pmatrix}
\]

\section{Tables}
Below is a sample table showing data:

\begin{table}[h]
\centering
\begin{tabular}{@{}lrrr@{}}
\toprule
\textbf{Category} & \textbf{Value 1} & \textbf{Value 2} & \textbf{Total} \\
\midrule
Item A & 10.5 & 20.3 & 30.8 \\
Item B & 15.2 & 25.1 & 40.3 \\
Item C & 22.1 & 12.7 & 34.8 \\
\bottomrule
\end{tabular}
\caption{Sample Data Table}
\label{tab:sample}
\end{table}

\section{Lists}
\subsection{Bullet Points}
\begin{itemize}
  \item First item
  \item Second item
  \item Third item with sub-items:
    \begin{itemize}
      \item Sub-item 1
      \item Sub-item 2
    \end{itemize}
\end{itemize}

\subsection{Numbered List}
\begin{enumerate}
  \item First step
  \item Second step
  \item Third step
\end{enumerate}

\section{Citations and References}
When you use the citation search feature, references will be automatically managed. For example, you might cite a paper like this: \cite{example2024}. All citations will be automatically added to your \texttt{cite.bib} file.

\section{Conclusion}
This template demonstrates basic \LaTeX{} features including mathematical formulas, tables, and lists. The bibliography system is now set up to work seamlessly with Bibby AI's citation management. Feel free to modify it for your needs!

\bibliographystyle{plain}
\bibliography{cite}

\end{document}\documentclass{article}
\usepackage[utf8]{inputenc}
\usepackage{booktabs}
\usepackage{amsmath}
\usepackage{natbib}

\title{My Document}
\author{Your Name}
\date{\today}

\begin{document}

\maketitle

\section{Introduction}
This is a sample document demonstrating various \LaTeX{} features including mathematical equations, tables, and formatting. You can cite papers using the citation search feature, and they will be automatically added to your \texttt{cite.bib} file.

\section{Mathematical Expressions}
Here are some example equations:

\subsection{Quadratic Formula}
The quadratic formula is given by:
\[x = \frac{-b \pm \sqrt{b^2 - 4ac}}{2a}\]

\subsection{Matrix Example}
A sample matrix:
\[
\begin{pmatrix}
a_{11} & a_{12} & a_{13} \\
a_{21} & a_{22} & a_{23} \\
a_{31} & a_{32} & a_{33}
\end{pmatrix}
\]

\section{Tables}
Below is a sample table showing data:

\begin{table}[h]
\centering
\begin{tabular}{@{}lrrr@{}}
\toprule
\textbf{Category} & \textbf{Value 1} & \textbf{Value 2} & \textbf{Total} \\
\midrule
Item A & 10.5 & 20.3 & 30.8 \\
Item B & 15.2 & 25.1 & 40.3 \\
Item C & 22.1 & 12.7 & 34.8 \\
\bottomrule
\end{tabular}
\caption{Sample Data Table}
\label{tab:sample}
\end{table}

\section{Lists}
\subsection{Bullet Points}
\begin{itemize}
  \item First item
  \item Second item
  \item Third item with sub-items:
    \begin{itemize}
      \item Sub-item 1
      \item Sub-item 2
    \end{itemize}
\end{itemize}

\subsection{Numbered List}
\begin{enumerate}
  \item First step
  \item Second step
  \item Third step
\end{enumerate}

\section{Citations and References}
When you use the citation search feature, references will be automatically managed. For example, you might cite a paper like this: \cite{example2024}. All citations will be automatically added to your \texttt{cite.bib} file.

\section{Conclusion}
This template demonstrates basic \LaTeX{} features including mathematical formulas, tables, and lists. The bibliography system is now set up to work seamlessly with Bibby AI's citation management. Feel free to modify it for your needs!

\bibliographystyle{plain}
\bibliography{cite}

\end{document}\documentclass{article}
\usepackage[utf8]{inputenc}
\usepackage{booktabs}
\usepackage{amsmath}
\usepackage{natbib}

\title{My Document}
\author{Your Name}
\date{\today}

\begin{document}

\maketitle

\section{Introduction}
This is a sample document demonstrating various \LaTeX{} features including mathematical equations, tables, and formatting. You can cite papers using the citation search feature, and they will be automatically added to your \texttt{cite.bib} file.

\section{Mathematical Expressions}
Here are some example equations:

\subsection{Quadratic Formula}
The quadratic formula is given by:
\[x = \frac{-b \pm \sqrt{b^2 - 4ac}}{2a}\]

\subsection{Matrix Example}
A sample matrix:
\[
\begin{pmatrix}
a_{11} & a_{12} & a_{13} \\
a_{21} & a_{22} & a_{23} \\
a_{31} & a_{32} & a_{33}
\end{pmatrix}
\]

\section{Tables}
Below is a sample table showing data:

\begin{table}[h]
\centering
\begin{tabular}{@{}lrrr@{}}
\toprule
\textbf{Category} & \textbf{Value 1} & \textbf{Value 2} & \textbf{Total} \\
\midrule
Item A & 10.5 & 20.3 & 30.8 \\
Item B & 15.2 & 25.1 & 40.3 \\
Item C & 22.1 & 12.7 & 34.8 \\
\bottomrule
\end{tabular}
\caption{Sample Data Table}
\label{tab:sample}
\end{table}

\section{Lists}
\subsection{Bullet Points}
\begin{itemize}
  \item First item
  \item Second item
  \item Third item with sub-items:
    \begin{itemize}
      \item Sub-item 1
      \item Sub-item 2
    \end{itemize}
\end{itemize}

\subsection{Numbered List}
\begin{enumerate}
  \item First step
  \item Second step
  \item Third step
\end{enumerate}

\section{Citations and References}
When you use the citation search feature, references will be automatically managed. For example, you might cite a paper like this: \cite{example2024}. All citations will be automatically added to your \texttt{cite.bib} file.

\section{Conclusion}
This template demonstrates basic \LaTeX{} features including mathematical formulas, tables, and lists. The bibliography system is now set up to work seamlessly with Bibby AI's citation management. Feel free to modify it for your needs!

\bibliographystyle{plain}
\bibliography{cite}

\end{document}\documentclass{article}
\usepackage[utf8]{inputenc}
\usepackage{booktabs}
\usepackage{amsmath}
\usepackage{natbib}

\title{My Document}
\author{Your Name}
\date{\today}

\begin{document}

\maketitle

\section{Introduction}
This is a sample document demonstrating various \LaTeX{} features including mathematical equations, tables, and formatting. You can cite papers using the citation search feature, and they will be automatically added to your \texttt{cite.bib} file.

\section{Mathematical Expressions}
Here are some example equations:

\subsection{Quadratic Formula}
The quadratic formula is given by:
\[x = \frac{-b \pm \sqrt{b^2 - 4ac}}{2a}\]

\subsection{Matrix Example}
A sample matrix:
\[
\begin{pmatrix}
a_{11} & a_{12} & a_{13} \\
a_{21} & a_{22} & a_{23} \\
a_{31} & a_{32} & a_{33}
\end{pmatrix}
\]

\section{Tables}
Below is a sample table showing data:

\begin{table}[h]
\centering
\begin{tabular}{@{}lrrr@{}}
\toprule
\textbf{Category} & \textbf{Value 1} & \textbf{Value 2} & \textbf{Total} \\
\midrule
Item A & 10.5 & 20.3 & 30.8 \\
Item B & 15.2 & 25.1 & 40.3 \\
Item C & 22.1 & 12.7 & 34.8 \\
\bottomrule
\end{tabular}
\caption{Sample Data Table}
\label{tab:sample}
\end{table}

\section{Lists}
\subsection{Bullet Points}
\begin{itemize}
  \item First item
  \item Second item
  \item Third item with sub-items:
    \begin{itemize}
      \item Sub-item 1
      \item Sub-item 2
    \end{itemize}
\end{itemize}

\subsection{Numbered List}
\begin{enumerate}
  \item First step
  \item Second step
  \item Third step
\end{enumerate}

\section{Citations and References}
When you use the citation search feature, references will be automatically managed. For example, you might cite a paper like this: \cite{example2024}. All citations will be automatically added to your \texttt{cite.bib} file.

\section{Conclusion}
This template demonstrates basic \LaTeX{} features including mathematical formulas, tables, and lists. The bibliography system is now set up to work seamlessly with Bibby AI's citation management. Feel free to modify it for your needs!

\bibliographystyle{plain}
\bibliography{cite}

\end{document}\documentclass{article}
\usepackage[utf8]{inputenc}
\usepackage{booktabs}
\usepackage{amsmath}
\usepackage{natbib}

\title{My Document}
\author{Your Name}
\date{\today}

\begin{document}

\maketitle

\section{Introduction}
This is a sample document demonstrating various \LaTeX{} features including mathematical equations, tables, and formatting. You can cite papers using the citation search feature, and they will be automatically added to your \texttt{cite.bib} file.

\section{Mathematical Expressions}
Here are some example equations:

\subsection{Quadratic Formula}
The quadratic formula is given by:
\[x = \frac{-b \pm \sqrt{b^2 - 4ac}}{2a}\]

\subsection{Matrix Example}
A sample matrix:
\[
\begin{pmatrix}
a_{11} & a_{12} & a_{13} \\
a_{21} & a_{22} & a_{23} \\
a_{31} & a_{32} & a_{33}
\end{pmatrix}
\]

\section{Tables}
Below is a sample table showing data:

\begin{table}[h]
\centering
\begin{tabular}{@{}lrrr@{}}
\toprule
\textbf{Category} & \textbf{Value 1} & \textbf{Value 2} & \textbf{Total} \\
\midrule
Item A & 10.5 & 20.3 & 30.8 \\
Item B & 15.2 & 25.1 & 40.3 \\
Item C & 22.1 & 12.7 & 34.8 \\
\bottomrule
\end{tabular}
\caption{Sample Data Table}
\label{tab:sample}
\end{table}

\section{Lists}
\subsection{Bullet Points}
\begin{itemize}
  \item First item
  \item Second item
  \item Third item with sub-items:
    \begin{itemize}
      \item Sub-item 1
      \item Sub-item 2
    \end{itemize}
\end{itemize}

\subsection{Numbered List}
\begin{enumerate}
  \item First step
  \item Second step
  \item Third step
\end{enumerate}

\section{Citations and References}
When you use the citation search feature, references will be automatically managed. For example, you might cite a paper like this: \cite{example2024}. All citations will be automatically added to your \texttt{cite.bib} file.

\section{Conclusion}
This template demonstrates basic \LaTeX{} features including mathematical formulas, tables, and lists. The bibliography system is now set up to work seamlessly with Bibby AI's citation management. Feel free to modify it for your needs!

\bibliographystyle{plain}
\bibliography{cite}

\end{document}\documentclass{article}
\usepackage[utf8]{inputenc}
\usepackage{booktabs}
\usepackage{amsmath}
\usepackage{natbib}

\title{My Document}
\author{Your Name}
\date{\today}

\begin{document}

\maketitle

\section{Introduction}
This is a sample document demonstrating various \LaTeX{} features including mathematical equations, tables, and formatting. You can cite papers using the citation search feature, and they will be automatically added to your \texttt{cite.bib} file.

\section{Mathematical Expressions}
Here are some example equations:

\subsection{Quadratic Formula}
The quadratic formula is given by:
\[x = \frac{-b \pm \sqrt{b^2 - 4ac}}{2a}\]

\subsection{Matrix Example}
A sample matrix:
\[
\begin{pmatrix}
a_{11} & a_{12} & a_{13} \\
a_{21} & a_{22} & a_{23} \\
a_{31} & a_{32} & a_{33}
\end{pmatrix}
\]

\section{Tables}
Below is a sample table showing data:

\begin{table}[h]
\centering
\begin{tabular}{@{}lrrr@{}}
\toprule
\textbf{Category} & \textbf{Value 1} & \textbf{Value 2} & \textbf{Total} \\
\midrule
Item A & 10.5 & 20.3 & 30.8 \\
Item B & 15.2 & 25.1 & 40.3 \\
Item C & 22.1 & 12.7 & 34.8 \\
\bottomrule
\end{tabular}
\caption{Sample Data Table}
\label{tab:sample}
\end{table}

\section{Lists}
\subsection{Bullet Points}
\begin{itemize}
  \item First item
  \item Second item
  \item Third item with sub-items:
    \begin{itemize}
      \item Sub-item 1
      \item Sub-item 2
    \end{itemize}
\end{itemize}

\subsection{Numbered List}
\begin{enumerate}
  \item First step
  \item Second step
  \item Third step
\end{enumerate}

\section{Citations and References}
When you use the citation search feature, references will be automatically managed. For example, you might cite a paper like this: \cite{example2024}. All citations will be automatically added to your \texttt{cite.bib} file.

\section{Conclusion}
This template demonstrates basic \LaTeX{} features including mathematical formulas, tables, and lists. The bibliography system is now set up to work seamlessly with Bibby AI's citation management. Feel free to modify it for your needs!

\bibliographystyle{plain}
\bibliography{cite}

\end{document}\documentclass{article}
\usepackage[utf8]{inputenc}
\usepackage{booktabs}
\usepackage{amsmath}
\usepackage{natbib}

\title{My Document}
\author{Your Name}
\date{\today}

\begin{document}

\maketitle

\section{Introduction}
This is a sample document demonstrating various \LaTeX{} features including mathematical equations, tables, and formatting. You can cite papers using the citation search feature, and they will be automatically added to your \texttt{cite.bib} file.

\section{Mathematical Expressions}
Here are some example equations:

\subsection{Quadratic Formula}
The quadratic formula is given by:
\[x = \frac{-b \pm \sqrt{b^2 - 4ac}}{2a}\]

\subsection{Matrix Example}
A sample matrix:
\[
\begin{pmatrix}
a_{11} & a_{12} & a_{13} \\
a_{21} & a_{22} & a_{23} \\
a_{31} & a_{32} & a_{33}
\end{pmatrix}
\]

\section{Tables}
Below is a sample table showing data:

\begin{table}[h]
\centering
\begin{tabular}{@{}lrrr@{}}
\toprule
\textbf{Category} & \textbf{Value 1} & \textbf{Value 2} & \textbf{Total} \\
\midrule
Item A & 10.5 & 20.3 & 30.8 \\
Item B & 15.2 & 25.1 & 40.3 \\
Item C & 22.1 & 12.7 & 34.8 \\
\bottomrule
\end{tabular}
\caption{Sample Data Table}
\label{tab:sample}
\end{table}

\section{Lists}
\subsection{Bullet Points}
\begin{itemize}
  \item First item
  \item Second item
  \item Third item with sub-items:
    \begin{itemize}
      \item Sub-item 1
      \item Sub-item 2
    \end{itemize}
\end{itemize}

\subsection{Numbered List}
\begin{enumerate}
  \item First step
  \item Second step
  \item Third step
\end{enumerate}

\section{Citations and References}
When you use the citation search feature, references will be automatically managed. For example, you might cite a paper like this: \cite{example2024}. All citations will be automatically added to your \texttt{cite.bib} file.

\section{Conclusion}
This template demonstrates basic \LaTeX{} features including mathematical formulas, tables, and lists. The bibliography system is now set up to work seamlessly with Bibby AI's citation management. Feel free to modify it for your needs!

\bibliographystyle{plain}
\bibliography{cite}

\end{document}\documentclass{article}
\usepackage[utf8]{inputenc}
\usepackage{booktabs}
\usepackage{amsmath}
\usepackage{natbib}

\title{My Document}
\author{Your Name}
\date{\today}

\begin{document}

\maketitle

\section{Introduction}
This is a sample document demonstrating various \LaTeX{} features including mathematical equations, tables, and formatting. You can cite papers using the citation search feature, and they will be automatically added to your \texttt{cite.bib} file.

\section{Mathematical Expressions}
Here are some example equations:

\subsection{Quadratic Formula}
The quadratic formula is given by:
\[x = \frac{-b \pm \sqrt{b^2 - 4ac}}{2a}\]

\subsection{Matrix Example}
A sample matrix:
\[
\begin{pmatrix}
a_{11} & a_{12} & a_{13} \\
a_{21} & a_{22} & a_{23} \\
a_{31} & a_{32} & a_{33}
\end{pmatrix}
\]

\section{Tables}
Below is a sample table showing data:

\begin{table}[h]
\centering
\begin{tabular}{@{}lrrr@{}}
\toprule
\textbf{Category} & \textbf{Value 1} & \textbf{Value 2} & \textbf{Total} \\
\midrule
Item A & 10.5 & 20.3 & 30.8 \\
Item B & 15.2 & 25.1 & 40.3 \\
Item C & 22.1 & 12.7 & 34.8 \\
\bottomrule
\end{tabular}
\caption{Sample Data Table}
\label{tab:sample}
\end{table}

\section{Lists}
\subsection{Bullet Points}
\begin{itemize}
  \item First item
  \item Second item
  \item Third item with sub-items:
    \begin{itemize}
      \item Sub-item 1
      \item Sub-item 2
    \end{itemize}
\end{itemize}

\subsection{Numbered List}
\begin{enumerate}
  \item First step
  \item Second step
  \item Third step
\end{enumerate}

\section{Citations and References}
When you use the citation search feature, references will be automatically managed. For example, you might cite a paper like this: \cite{example2024}. All citations will be automatically added to your \texttt{cite.bib} file.

\section{Conclusion}
This template demonstrates basic \LaTeX{} features including mathematical formulas, tables, and lists. The bibliography system is now set up to work seamlessly with Bibby AI's citation management. Feel free to modify it for your needs!

\bibliographystyle{plain}
\bibliography{cite}

\end{document}\documentclass{article}
\usepackage[utf8]{inputenc}
\usepackage{booktabs}
\usepackage{amsmath}
\usepackage{natbib}

\title{My Document}
\author{Your Name}
\date{\today}

\begin{document}

\maketitle

\section{Introduction}
This is a sample document demonstrating various \LaTeX{} features including mathematical equations, tables, and formatting. You can cite papers using the citation search feature, and they will be automatically added to your \texttt{cite.bib} file.

\section{Mathematical Expressions}
Here are some example equations:

\subsection{Quadratic Formula}
The quadratic formula is given by:
\[x = \frac{-b \pm \sqrt{b^2 - 4ac}}{2a}\]

\subsection{Matrix Example}
A sample matrix:
\[
\begin{pmatrix}
a_{11} & a_{12} & a_{13} \\
a_{21} & a_{22} & a_{23} \\
a_{31} & a_{32} & a_{33}
\end{pmatrix}
\]

\section{Tables}
Below is a sample table showing data:

\begin{table}[h]
\centering
\begin{tabular}{@{}lrrr@{}}
\toprule
\textbf{Category} & \textbf{Value 1} & \textbf{Value 2} & \textbf{Total} \\
\midrule
Item A & 10.5 & 20.3 & 30.8 \\
Item B & 15.2 & 25.1 & 40.3 \\
Item C & 22.1 & 12.7 & 34.8 \\
\bottomrule
\end{tabular}
\caption{Sample Data Table}
\label{tab:sample}
\end{table}

\section{Lists}
\subsection{Bullet Points}
\begin{itemize}
  \item First item
  \item Second item
  \item Third item with sub-items:
    \begin{itemize}
      \item Sub-item 1
      \item Sub-item 2
    \end{itemize}
\end{itemize}

\subsection{Numbered List}
\begin{enumerate}
  \item First step
  \item Second step
  \item Third step
\end{enumerate}

\section{Citations and References}
When you use the citation search feature, references will be automatically managed. For example, you might cite a paper like this: \cite{example2024}. All citations will be automatically added to your \texttt{cite.bib} file.

\section{Conclusion}
This template demonstrates basic \LaTeX{} features including mathematical formulas, tables, and lists. The bibliography system is now set up to work seamlessly with Bibby AI's citation management. Feel free to modify it for your needs!

\bibliographystyle{plain}
\bibliography{cite}

\end{document}\documentclass{article}
\usepackage[utf8]{inputenc}
\usepackage{booktabs}
\usepackage{amsmath}
\usepackage{natbib}

\title{My Document}
\author{Your Name}
\date{\today}

\begin{document}

\maketitle

\section{Introduction}
This is a sample document demonstrating various \LaTeX{} features including mathematical equations, tables, and formatting. You can cite papers using the citation search feature, and they will be automatically added to your \texttt{cite.bib} file.

\section{Mathematical Expressions}
Here are some example equations:

\subsection{Quadratic Formula}
The quadratic formula is given by:
\[x = \frac{-b \pm \sqrt{b^2 - 4ac}}{2a}\]

\subsection{Matrix Example}
A sample matrix:
\[
\begin{pmatrix}
a_{11} & a_{12} & a_{13} \\
a_{21} & a_{22} & a_{23} \\
a_{31} & a_{32} & a_{33}
\end{pmatrix}
\]

\section{Tables}
Below is a sample table showing data:

\begin{table}[h]
\centering
\begin{tabular}{@{}lrrr@{}}
\toprule
\textbf{Category} & \textbf{Value 1} & \textbf{Value 2} & \textbf{Total} \\
\midrule
Item A & 10.5 & 20.3 & 30.8 \\
Item B & 15.2 & 25.1 & 40.3 \\
Item C & 22.1 & 12.7 & 34.8 \\
\bottomrule
\end{tabular}
\caption{Sample Data Table}
\label{tab:sample}
\end{table}

\section{Lists}
\subsection{Bullet Points}
\begin{itemize}
  \item First item
  \item Second item
  \item Third item with sub-items:
    \begin{itemize}
      \item Sub-item 1
      \item Sub-item 2
    \end{itemize}
\end{itemize}

\subsection{Numbered List}
\begin{enumerate}
  \item First step
  \item Second step
  \item Third step
\end{enumerate}

\section{Citations and References}
When you use the citation search feature, references will be automatically managed. For example, you might cite a paper like this: \cite{example2024}. All citations will be automatically added to your \texttt{cite.bib} file.

\section{Conclusion}
This template demonstrates basic \LaTeX{} features including mathematical formulas, tables, and lists. The bibliography system is now set up to work seamlessly with Bibby AI's citation management. Feel free to modify it for your needs!

\bibliographystyle{plain}
\bibliography{cite}

\end{document}\documentclass{article}
\usepackage[utf8]{inputenc}
\usepackage{booktabs}
\usepackage{amsmath}
\usepackage{natbib}

\title{My Document}
\author{Your Name}
\date{\today}

\begin{document}

\maketitle

\section{Introduction}
This is a sample document demonstrating various \LaTeX{} features including mathematical equations, tables, and formatting. You can cite papers using the citation search feature, and they will be automatically added to your \texttt{cite.bib} file.

\section{Mathematical Expressions}
Here are some example equations:

\subsection{Quadratic Formula}
The quadratic formula is given by:
\[x = \frac{-b \pm \sqrt{b^2 - 4ac}}{2a}\]

\subsection{Matrix Example}
A sample matrix:
\[
\begin{pmatrix}
a_{11} & a_{12} & a_{13} \\
a_{21} & a_{22} & a_{23} \\
a_{31} & a_{32} & a_{33}
\end{pmatrix}
\]

\section{Tables}
Below is a sample table showing data:

\begin{table}[h]
\centering
\begin{tabular}{@{}lrrr@{}}
\toprule
\textbf{Category} & \textbf{Value 1} & \textbf{Value 2} & \textbf{Total} \\
\midrule
Item A & 10.5 & 20.3 & 30.8 \\
Item B & 15.2 & 25.1 & 40.3 \\
Item C & 22.1 & 12.7 & 34.8 \\
\bottomrule
\end{tabular}
\caption{Sample Data Table}
\label{tab:sample}
\end{table}

\section{Lists}
\subsection{Bullet Points}
\begin{itemize}
  \item First item
  \item Second item
  \item Third item with sub-items:
    \begin{itemize}
      \item Sub-item 1
      \item Sub-item 2
    \end{itemize}
\end{itemize}

\subsection{Numbered List}
\begin{enumerate}
  \item First step
  \item Second step
  \item Third step
\end{enumerate}

\section{Citations and References}
When you use the citation search feature, references will be automatically managed. For example, you might cite a paper like this: \cite{example2024}. All citations will be automatically added to your \texttt{cite.bib} file.

\section{Conclusion}
This template demonstrates basic \LaTeX{} features including mathematical formulas, tables, and lists. The bibliography system is now set up to work seamlessly with Bibby AI's citation management. Feel free to modify it for your needs!

\bibliographystyle{plain}
\bibliography{cite}

\end{document}\documentclass{article}
\usepackage[utf8]{inputenc}
\usepackage{booktabs}
\usepackage{amsmath}
\usepackage{natbib}

\title{My Document}
\author{Your Name}
\date{\today}

\begin{document}

\maketitle

\section{Introduction}
This is a sample document demonstrating various \LaTeX{} features including mathematical equations, tables, and formatting. You can cite papers using the citation search feature, and they will be automatically added to your \texttt{cite.bib} file.

\section{Mathematical Expressions}
Here are some example equations:

\subsection{Quadratic Formula}
The quadratic formula is given by:
\[x = \frac{-b \pm \sqrt{b^2 - 4ac}}{2a}\]

\subsection{Matrix Example}
A sample matrix:
\[
\begin{pmatrix}
a_{11} & a_{12} & a_{13} \\
a_{21} & a_{22} & a_{23} \\
a_{31} & a_{32} & a_{33}
\end{pmatrix}
\]

\section{Tables}
Below is a sample table showing data:

\begin{table}[h]
\centering
\begin{tabular}{@{}lrrr@{}}
\toprule
\textbf{Category} & \textbf{Value 1} & \textbf{Value 2} & \textbf{Total} \\
\midrule
Item A & 10.5 & 20.3 & 30.8 \\
Item B & 15.2 & 25.1 & 40.3 \\
Item C & 22.1 & 12.7 & 34.8 \\
\bottomrule
\end{tabular}
\caption{Sample Data Table}
\label{tab:sample}
\end{table}

\section{Lists}
\subsection{Bullet Points}
\begin{itemize}
  \item First item
  \item Second item
  \item Third item with sub-items:
    \begin{itemize}
      \item Sub-item 1
      \item Sub-item 2
    \end{itemize}
\end{itemize}

\subsection{Numbered List}
\begin{enumerate}
  \item First step
  \item Second step
  \item Third step
\end{enumerate}

\section{Citations and References}
When you use the citation search feature, references will be automatically managed. For example, you might cite a paper like this: \cite{example2024}. All citations will be automatically added to your \texttt{cite.bib} file.

\section{Conclusion}
This template demonstrates basic \LaTeX{} features including mathematical formulas, tables, and lists. The bibliography system is now set up to work seamlessly with Bibby AI's citation management. Feel free to modify it for your needs!

\bibliographystyle{plain}
\bibliography{cite}

\end{document}\documentclass{article}
\usepackage[utf8]{inputenc}
\usepackage{booktabs}
\usepackage{amsmath}
\usepackage{natbib}

\title{My Document}
\author{Your Name}
\date{\today}

\begin{document}

\maketitle

\section{Introduction}
This is a sample document demonstrating various \LaTeX{} features including mathematical equations, tables, and formatting. You can cite papers using the citation search feature, and they will be automatically added to your \texttt{cite.bib} file.

\section{Mathematical Expressions}
Here are some example equations:

\subsection{Quadratic Formula}
The quadratic formula is given by:
\[x = \frac{-b \pm \sqrt{b^2 - 4ac}}{2a}\]

\subsection{Matrix Example}
A sample matrix:
\[
\begin{pmatrix}
a_{11} & a_{12} & a_{13} \\
a_{21} & a_{22} & a_{23} \\
a_{31} & a_{32} & a_{33}
\end{pmatrix}
\]

\section{Tables}
Below is a sample table showing data:

\begin{table}[h]
\centering
\begin{tabular}{@{}lrrr@{}}
\toprule
\textbf{Category} & \textbf{Value 1} & \textbf{Value 2} & \textbf{Total} \\
\midrule
Item A & 10.5 & 20.3 & 30.8 \\
Item B & 15.2 & 25.1 & 40.3 \\
Item C & 22.1 & 12.7 & 34.8 \\
\bottomrule
\end{tabular}
\caption{Sample Data Table}
\label{tab:sample}
\end{table}

\section{Lists}
\subsection{Bullet Points}
\begin{itemize}
  \item First item
  \item Second item
  \item Third item with sub-items:
    \begin{itemize}
      \item Sub-item 1
      \item Sub-item 2
    \end{itemize}
\end{itemize}

\subsection{Numbered List}
\begin{enumerate}
  \item First step
  \item Second step
  \item Third step
\end{enumerate}

\section{Citations and References}
When you use the citation search feature, references will be automatically managed. For example, you might cite a paper like this: \cite{example2024}. All citations will be automatically added to your \texttt{cite.bib} file.

\section{Conclusion}
This template demonstrates basic \LaTeX{} features including mathematical formulas, tables, and lists. The bibliography system is now set up to work seamlessly with Bibby AI's citation management. Feel free to modify it for your needs!

\bibliographystyle{plain}
\bibliography{cite}

\end{document}\documentclass{article}
\usepackage[utf8]{inputenc}
\usepackage{booktabs}
\usepackage{amsmath}
\usepackage{natbib}

\title{My Document}
\author{Your Name}
\date{\today}

\begin{document}

\maketitle

\section{Introduction}
This is a sample document demonstrating various \LaTeX{} features including mathematical equations, tables, and formatting. You can cite papers using the citation search feature, and they will be automatically added to your \texttt{cite.bib} file.

\section{Mathematical Expressions}
Here are some example equations:

\subsection{Quadratic Formula}
The quadratic formula is given by:
\[x = \frac{-b \pm \sqrt{b^2 - 4ac}}{2a}\]

\subsection{Matrix Example}
A sample matrix:
\[
\begin{pmatrix}
a_{11} & a_{12} & a_{13} \\
a_{21} & a_{22} & a_{23} \\
a_{31} & a_{32} & a_{33}
\end{pmatrix}
\]

\section{Tables}
Below is a sample table showing data:

\begin{table}[h]
\centering
\begin{tabular}{@{}lrrr@{}}
\toprule
\textbf{Category} & \textbf{Value 1} & \textbf{Value 2} & \textbf{Total} \\
\midrule
Item A & 10.5 & 20.3 & 30.8 \\
Item B & 15.2 & 25.1 & 40.3 \\
Item C & 22.1 & 12.7 & 34.8 \\
\bottomrule
\end{tabular}
\caption{Sample Data Table}
\label{tab:sample}
\end{table}

\section{Lists}
\subsection{Bullet Points}
\begin{itemize}
  \item First item
  \item Second item
  \item Third item with sub-items:
    \begin{itemize}
      \item Sub-item 1
      \item Sub-item 2
    \end{itemize}
\end{itemize}

\subsection{Numbered List}
\begin{enumerate}
  \item First step
  \item Second step
  \item Third step
\end{enumerate}

\section{Citations and References}
When you use the citation search feature, references will be automatically managed. For example, you might cite a paper like this: \cite{example2024}. All citations will be automatically added to your \texttt{cite.bib} file.

\section{Conclusion}
This template demonstrates basic \LaTeX{} features including mathematical formulas, tables, and lists. The bibliography system is now set up to work seamlessly with Bibby AI's citation management. Feel free to modify it for your needs!

\bibliographystyle{plain}
\bibliography{cite}

\end{document}\documentclass{article}
\usepackage[utf8]{inputenc}
\usepackage{booktabs}
\usepackage{amsmath}
\usepackage{natbib}

\title{My Document}
\author{Your Name}
\date{\today}

\begin{document}

\maketitle

\section{Introduction}
This is a sample document demonstrating various \LaTeX{} features including mathematical equations, tables, and formatting. You can cite papers using the citation search feature, and they will be automatically added to your \texttt{cite.bib} file.

\section{Mathematical Expressions}
Here are some example equations:

\subsection{Quadratic Formula}
The quadratic formula is given by:
\[x = \frac{-b \pm \sqrt{b^2 - 4ac}}{2a}\]

\subsection{Matrix Example}
A sample matrix:
\[
\begin{pmatrix}
a_{11} & a_{12} & a_{13} \\
a_{21} & a_{22} & a_{23} \\
a_{31} & a_{32} & a_{33}
\end{pmatrix}
\]

\section{Tables}
Below is a sample table showing data:

\begin{table}[h]
\centering
\begin{tabular}{@{}lrrr@{}}
\toprule
\textbf{Category} & \textbf{Value 1} & \textbf{Value 2} & \textbf{Total} \\
\midrule
Item A & 10.5 & 20.3 & 30.8 \\
Item B & 15.2 & 25.1 & 40.3 \\
Item C & 22.1 & 12.7 & 34.8 \\
\bottomrule
\end{tabular}
\caption{Sample Data Table}
\label{tab:sample}
\end{table}

\section{Lists}
\subsection{Bullet Points}
\begin{itemize}
  \item First item
  \item Second item
  \item Third item with sub-items:
    \begin{itemize}
      \item Sub-item 1
      \item Sub-item 2
    \end{itemize}
\end{itemize}

\subsection{Numbered List}
\begin{enumerate}
  \item First step
  \item Second step
  \item Third step
\end{enumerate}

\section{Citations and References}
When you use the citation search feature, references will be automatically managed. For example, you might cite a paper like this: \cite{example2024}. All citations will be automatically added to your \texttt{cite.bib} file.

\section{Conclusion}
This template demonstrates basic \LaTeX{} features including mathematical formulas, tables, and lists. The bibliography system is now set up to work seamlessly with Bibby AI's citation management. Feel free to modify it for your needs!

\bibliographystyle{plain}
\bibliography{cite}

\end{document}\documentclass{article}
\usepackage[utf8]{inputenc}
\usepackage{booktabs}
\usepackage{amsmath}
\usepackage{natbib}

\title{My Document}
\author{Your Name}
\date{\today}

\begin{document}

\maketitle

\section{Introduction}
This is a sample document demonstrating various \LaTeX{} features including mathematical equations, tables, and formatting. You can cite papers using the citation search feature, and they will be automatically added to your \texttt{cite.bib} file.

\section{Mathematical Expressions}
Here are some example equations:

\subsection{Quadratic Formula}
The quadratic formula is given by:
\[x = \frac{-b \pm \sqrt{b^2 - 4ac}}{2a}\]

\subsection{Matrix Example}
A sample matrix:
\[
\begin{pmatrix}
a_{11} & a_{12} & a_{13} \\
a_{21} & a_{22} & a_{23} \\
a_{31} & a_{32} & a_{33}
\end{pmatrix}
\]

\section{Tables}
Below is a sample table showing data:

\begin{table}[h]
\centering
\begin{tabular}{@{}lrrr@{}}
\toprule
\textbf{Category} & \textbf{Value 1} & \textbf{Value 2} & \textbf{Total} \\
\midrule
Item A & 10.5 & 20.3 & 30.8 \\
Item B & 15.2 & 25.1 & 40.3 \\
Item C & 22.1 & 12.7 & 34.8 \\
\bottomrule
\end{tabular}
\caption{Sample Data Table}
\label{tab:sample}
\end{table}

\section{Lists}
\subsection{Bullet Points}
\begin{itemize}
  \item First item
  \item Second item
  \item Third item with sub-items:
    \begin{itemize}
      \item Sub-item 1
      \item Sub-item 2
    \end{itemize}
\end{itemize}

\subsection{Numbered List}
\begin{enumerate}
  \item First step
  \item Second step
  \item Third step
\end{enumerate}

\section{Citations and References}
When you use the citation search feature, references will be automatically managed. For example, you might cite a paper like this: \cite{example2024}. All citations will be automatically added to your \texttt{cite.bib} file.

\section{Conclusion}
This template demonstrates basic \LaTeX{} features including mathematical formulas, tables, and lists. The bibliography system is now set up to work seamlessly with Bibby AI's citation management. Feel free to modify it for your needs!

\bibliographystyle{plain}
\bibliography{cite}

\end{document}\documentclass{article}
\usepackage[utf8]{inputenc}
\usepackage{booktabs}
\usepackage{amsmath}
\usepackage{natbib}

\title{My Document}
\author{Your Name}
\date{\today}

\begin{document}

\maketitle

\section{Introduction}
This is a sample document demonstrating various \LaTeX{} features including mathematical equations, tables, and formatting. You can cite papers using the citation search feature, and they will be automatically added to your \texttt{cite.bib} file.

\section{Mathematical Expressions}
Here are some example equations:

\subsection{Quadratic Formula}
The quadratic formula is given by:
\[x = \frac{-b \pm \sqrt{b^2 - 4ac}}{2a}\]

\subsection{Matrix Example}
A sample matrix:
\[
\begin{pmatrix}
a_{11} & a_{12} & a_{13} \\
a_{21} & a_{22} & a_{23} \\
a_{31} & a_{32} & a_{33}
\end{pmatrix}
\]

\section{Tables}
Below is a sample table showing data:

\begin{table}[h]
\centering
\begin{tabular}{@{}lrrr@{}}
\toprule
\textbf{Category} & \textbf{Value 1} & \textbf{Value 2} & \textbf{Total} \\
\midrule
Item A & 10.5 & 20.3 & 30.8 \\
Item B & 15.2 & 25.1 & 40.3 \\
Item C & 22.1 & 12.7 & 34.8 \\
\bottomrule
\end{tabular}
\caption{Sample Data Table}
\label{tab:sample}
\end{table}

\section{Lists}
\subsection{Bullet Points}
\begin{itemize}
  \item First item
  \item Second item
  \item Third item with sub-items:
    \begin{itemize}
      \item Sub-item 1
      \item Sub-item 2
    \end{itemize}
\end{itemize}

\subsection{Numbered List}
\begin{enumerate}
  \item First step
  \item Second step
  \item Third step
\end{enumerate}

\section{Citations and References}
When you use the citation search feature, references will be automatically managed. For example, you might cite a paper like this: \cite{example2024}. All citations will be automatically added to your \texttt{cite.bib} file.

\section{Conclusion}
This template demonstrates basic \LaTeX{} features including mathematical formulas, tables, and lists. The bibliography system is now set up to work seamlessly with Bibby AI's citation management. Feel free to modify it for your needs!

\bibliographystyle{plain}
\bibliography{cite}

\end{document}\documentclass{article}
\usepackage[utf8]{inputenc}
\usepackage{booktabs}
\usepackage{amsmath}
\usepackage{natbib}

\title{My Document}
\author{Your Name}
\date{\today}

\begin{document}

\maketitle

\section{Introduction}
This is a sample document demonstrating various \LaTeX{} features including mathematical equations, tables, and formatting. You can cite papers using the citation search feature, and they will be automatically added to your \texttt{cite.bib} file.

\section{Mathematical Expressions}
Here are some example equations:

\subsection{Quadratic Formula}
The quadratic formula is given by:
\[x = \frac{-b \pm \sqrt{b^2 - 4ac}}{2a}\]

\subsection{Matrix Example}
A sample matrix:
\[
\begin{pmatrix}
a_{11} & a_{12} & a_{13} \\
a_{21} & a_{22} & a_{23} \\
a_{31} & a_{32} & a_{33}
\end{pmatrix}
\]

\section{Tables}
Below is a sample table showing data:

\begin{table}[h]
\centering
\begin{tabular}{@{}lrrr@{}}
\toprule
\textbf{Category} & \textbf{Value 1} & \textbf{Value 2} & \textbf{Total} \\
\midrule
Item A & 10.5 & 20.3 & 30.8 \\
Item B & 15.2 & 25.1 & 40.3 \\
Item C & 22.1 & 12.7 & 34.8 \\
\bottomrule
\end{tabular}
\caption{Sample Data Table}
\label{tab:sample}
\end{table}

\section{Lists}
\subsection{Bullet Points}
\begin{itemize}
  \item First item
  \item Second item
  \item Third item with sub-items:
    \begin{itemize}
      \item Sub-item 1
      \item Sub-item 2
    \end{itemize}
\end{itemize}

\subsection{Numbered List}
\begin{enumerate}
  \item First step
  \item Second step
  \item Third step
\end{enumerate}

\section{Citations and References}
When you use the citation search feature, references will be automatically managed. For example, you might cite a paper like this: \cite{example2024}. All citations will be automatically added to your \texttt{cite.bib} file.

\section{Conclusion}
This template demonstrates basic \LaTeX{} features including mathematical formulas, tables, and lists. The bibliography system is now set up to work seamlessly with Bibby AI's citation management. Feel free to modify it for your needs!

\bibliographystyle{plain}
\bibliography{cite}

\end{document}\documentclass{article}
\usepackage[utf8]{inputenc}
\usepackage{booktabs}
\usepackage{amsmath}
\usepackage{natbib}

\title{My Document}
\author{Your Name}
\date{\today}

\begin{document}

\maketitle

\section{Introduction}
This is a sample document demonstrating various \LaTeX{} features including mathematical equations, tables, and formatting. You can cite papers using the citation search feature, and they will be automatically added to your \texttt{cite.bib} file.

\section{Mathematical Expressions}
Here are some example equations:

\subsection{Quadratic Formula}
The quadratic formula is given by:
\[x = \frac{-b \pm \sqrt{b^2 - 4ac}}{2a}\]

\subsection{Matrix Example}
A sample matrix:
\[
\begin{pmatrix}
a_{11} & a_{12} & a_{13} \\
a_{21} & a_{22} & a_{23} \\
a_{31} & a_{32} & a_{33}
\end{pmatrix}
\]

\section{Tables}
Below is a sample table showing data:

\begin{table}[h]
\centering
\begin{tabular}{@{}lrrr@{}}
\toprule
\textbf{Category} & \textbf{Value 1} & \textbf{Value 2} & \textbf{Total} \\
\midrule
Item A & 10.5 & 20.3 & 30.8 \\
Item B & 15.2 & 25.1 & 40.3 \\
Item C & 22.1 & 12.7 & 34.8 \\
\bottomrule
\end{tabular}
\caption{Sample Data Table}
\label{tab:sample}
\end{table}

\section{Lists}
\subsection{Bullet Points}
\begin{itemize}
  \item First item
  \item Second item
  \item Third item with sub-items:
    \begin{itemize}
      \item Sub-item 1
      \item Sub-item 2
    \end{itemize}
\end{itemize}

\subsection{Numbered List}
\begin{enumerate}
  \item First step
  \item Second step
  \item Third step
\end{enumerate}

\section{Citations and References}
When you use the citation search feature, references will be automatically managed. For example, you might cite a paper like this: \cite{example2024}. All citations will be automatically added to your \texttt{cite.bib} file.

\section{Conclusion}
This template demonstrates basic \LaTeX{} features including mathematical formulas, tables, and lists. The bibliography system is now set up to work seamlessly with Bibby AI's citation management. Feel free to modify it for your needs!

\bibliographystyle{plain}
\bibliography{cite}

\end{document}\documentclass{article}
\usepackage[utf8]{inputenc}
\usepackage{booktabs}
\usepackage{amsmath}
\usepackage{natbib}

\title{My Document}
\author{Your Name}
\date{\today}

\begin{document}

\maketitle

\section{Introduction}
This is a sample document demonstrating various \LaTeX{} features including mathematical equations, tables, and formatting. You can cite papers using the citation search feature, and they will be automatically added to your \texttt{cite.bib} file.

\section{Mathematical Expressions}
Here are some example equations:

\subsection{Quadratic Formula}
The quadratic formula is given by:
\[x = \frac{-b \pm \sqrt{b^2 - 4ac}}{2a}\]

\subsection{Matrix Example}
A sample matrix:
\[
\begin{pmatrix}
a_{11} & a_{12} & a_{13} \\
a_{21} & a_{22} & a_{23} \\
a_{31} & a_{32} & a_{33}
\end{pmatrix}
\]

\section{Tables}
Below is a sample table showing data:

\begin{table}[h]
\centering
\begin{tabular}{@{}lrrr@{}}
\toprule
\textbf{Category} & \textbf{Value 1} & \textbf{Value 2} & \textbf{Total} \\
\midrule
Item A & 10.5 & 20.3 & 30.8 \\
Item B & 15.2 & 25.1 & 40.3 \\
Item C & 22.1 & 12.7 & 34.8 \\
\bottomrule
\end{tabular}
\caption{Sample Data Table}
\label{tab:sample}
\end{table}

\section{Lists}
\subsection{Bullet Points}
\begin{itemize}
  \item First item
  \item Second item
  \item Third item with sub-items:
    \begin{itemize}
      \item Sub-item 1
      \item Sub-item 2
    \end{itemize}
\end{itemize}

\subsection{Numbered List}
\begin{enumerate}
  \item First step
  \item Second step
  \item Third step
\end{enumerate}

\section{Citations and References}
When you use the citation search feature, references will be automatically managed. For example, you might cite a paper like this: \cite{example2024}. All citations will be automatically added to your \texttt{cite.bib} file.

\section{Conclusion}
This template demonstrates basic \LaTeX{} features including mathematical formulas, tables, and lists. The bibliography system is now set up to work seamlessly with Bibby AI's citation management. Feel free to modify it for your needs!

\bibliographystyle{plain}
\bibliography{cite}

\end{document}\documentclass{article}
\usepackage[utf8]{inputenc}
\usepackage{booktabs}
\usepackage{amsmath}
\usepackage{natbib}

\title{My Document}
\author{Your Name}
\date{\today}

\begin{document}

\maketitle

\section{Introduction}
This is a sample document demonstrating various \LaTeX{} features including mathematical equations, tables, and formatting. You can cite papers using the citation search feature, and they will be automatically added to your \texttt{cite.bib} file.

\section{Mathematical Expressions}
Here are some example equations:

\subsection{Quadratic Formula}
The quadratic formula is given by:
\[x = \frac{-b \pm \sqrt{b^2 - 4ac}}{2a}\]

\subsection{Matrix Example}
A sample matrix:
\[
\begin{pmatrix}
a_{11} & a_{12} & a_{13} \\
a_{21} & a_{22} & a_{23} \\
a_{31} & a_{32} & a_{33}
\end{pmatrix}
\]

\section{Tables}
Below is a sample table showing data:

\begin{table}[h]
\centering
\begin{tabular}{@{}lrrr@{}}
\toprule
\textbf{Category} & \textbf{Value 1} & \textbf{Value 2} & \textbf{Total} \\
\midrule
Item A & 10.5 & 20.3 & 30.8 \\
Item B & 15.2 & 25.1 & 40.3 \\
Item C & 22.1 & 12.7 & 34.8 \\
\bottomrule
\end{tabular}
\caption{Sample Data Table}
\label{tab:sample}
\end{table}

\section{Lists}
\subsection{Bullet Points}
\begin{itemize}
  \item First item
  \item Second item
  \item Third item with sub-items:
    \begin{itemize}
      \item Sub-item 1
      \item Sub-item 2
    \end{itemize}
\end{itemize}

\subsection{Numbered List}
\begin{enumerate}
  \item First step
  \item Second step
  \item Third step
\end{enumerate}

\section{Citations and References}
When you use the citation search feature, references will be automatically managed. For example, you might cite a paper like this: \cite{example2024}. All citations will be automatically added to your \texttt{cite.bib} file.

\section{Conclusion}
This template demonstrates basic \LaTeX{} features including mathematical formulas, tables, and lists. The bibliography system is now set up to work seamlessly with Bibby AI's citation management. Feel free to modify it for your needs!

\bibliographystyle{plain}
\bibliography{cite}

\end{document}\documentclass{article}
\usepackage[utf8]{inputenc}
\usepackage{booktabs}
\usepackage{amsmath}
\usepackage{natbib}

\title{My Document}
\author{Your Name}
\date{\today}

\begin{document}

\maketitle

\section{Introduction}
This is a sample document demonstrating various \LaTeX{} features including mathematical equations, tables, and formatting. You can cite papers using the citation search feature, and they will be automatically added to your \texttt{cite.bib} file.

\section{Mathematical Expressions}
Here are some example equations:

\subsection{Quadratic Formula}
The quadratic formula is given by:
\[x = \frac{-b \pm \sqrt{b^2 - 4ac}}{2a}\]

\subsection{Matrix Example}
A sample matrix:
\[
\begin{pmatrix}
a_{11} & a_{12} & a_{13} \\
a_{21} & a_{22} & a_{23} \\
a_{31} & a_{32} & a_{33}
\end{pmatrix}
\]

\section{Tables}
Below is a sample table showing data:

\begin{table}[h]
\centering
\begin{tabular}{@{}lrrr@{}}
\toprule
\textbf{Category} & \textbf{Value 1} & \textbf{Value 2} & \textbf{Total} \\
\midrule
Item A & 10.5 & 20.3 & 30.8 \\
Item B & 15.2 & 25.1 & 40.3 \\
Item C & 22.1 & 12.7 & 34.8 \\
\bottomrule
\end{tabular}
\caption{Sample Data Table}
\label{tab:sample}
\end{table}

\section{Lists}
\subsection{Bullet Points}
\begin{itemize}
  \item First item
  \item Second item
  \item Third item with sub-items:
    \begin{itemize}
      \item Sub-item 1
      \item Sub-item 2
    \end{itemize}
\end{itemize}

\subsection{Numbered List}
\begin{enumerate}
  \item First step
  \item Second step
  \item Third step
\end{enumerate}

\section{Citations and References}
When you use the citation search feature, references will be automatically managed. For example, you might cite a paper like this: \cite{example2024}. All citations will be automatically added to your \texttt{cite.bib} file.

\section{Conclusion}
This template demonstrates basic \LaTeX{} features including mathematical formulas, tables, and lists. The bibliography system is now set up to work seamlessly with Bibby AI's citation management. Feel free to modify it for your needs!

\bibliographystyle{plain}
\bibliography{cite}

\end{document}\documentclass{article}
\usepackage[utf8]{inputenc}
\usepackage{booktabs}
\usepackage{amsmath}
\usepackage{natbib}

\title{My Document}
\author{Your Name}
\date{\today}

\begin{document}

\maketitle

\section{Introduction}
This is a sample document demonstrating various \LaTeX{} features including mathematical equations, tables, and formatting. You can cite papers using the citation search feature, and they will be automatically added to your \texttt{cite.bib} file.

\section{Mathematical Expressions}
Here are some example equations:

\subsection{Quadratic Formula}
The quadratic formula is given by:
\[x = \frac{-b \pm \sqrt{b^2 - 4ac}}{2a}\]

\subsection{Matrix Example}
A sample matrix:
\[
\begin{pmatrix}
a_{11} & a_{12} & a_{13} \\
a_{21} & a_{22} & a_{23} \\
a_{31} & a_{32} & a_{33}
\end{pmatrix}
\]

\section{Tables}
Below is a sample table showing data:

\begin{table}[h]
\centering
\begin{tabular}{@{}lrrr@{}}
\toprule
\textbf{Category} & \textbf{Value 1} & \textbf{Value 2} & \textbf{Total} \\
\midrule
Item A & 10.5 & 20.3 & 30.8 \\
Item B & 15.2 & 25.1 & 40.3 \\
Item C & 22.1 & 12.7 & 34.8 \\
\bottomrule
\end{tabular}
\caption{Sample Data Table}
\label{tab:sample}
\end{table}

\section{Lists}
\subsection{Bullet Points}
\begin{itemize}
  \item First item
  \item Second item
  \item Third item with sub-items:
    \begin{itemize}
      \item Sub-item 1
      \item Sub-item 2
    \end{itemize}
\end{itemize}

\subsection{Numbered List}
\begin{enumerate}
  \item First step
  \item Second step
  \item Third step
\end{enumerate}

\section{Citations and References}
When you use the citation search feature, references will be automatically managed. For example, you might cite a paper like this: \cite{example2024}. All citations will be automatically added to your \texttt{cite.bib} file.

\section{Conclusion}
This template demonstrates basic \LaTeX{} features including mathematical formulas, tables, and lists. The bibliography system is now set up to work seamlessly with Bibby AI's citation management. Feel free to modify it for your needs!

\bibliographystyle{plain}
\bibliography{cite}

\end{document}\documentclass{article}
\usepackage[utf8]{inputenc}
\usepackage{booktabs}
\usepackage{amsmath}
\usepackage{natbib}

\title{My Document}
\author{Your Name}
\date{\today}

\begin{document}

\maketitle

\section{Introduction}
This is a sample document demonstrating various \LaTeX{} features including mathematical equations, tables, and formatting. You can cite papers using the citation search feature, and they will be automatically added to your \texttt{cite.bib} file.

\section{Mathematical Expressions}
Here are some example equations:

\subsection{Quadratic Formula}
The quadratic formula is given by:
\[x = \frac{-b \pm \sqrt{b^2 - 4ac}}{2a}\]

\subsection{Matrix Example}
A sample matrix:
\[
\begin{pmatrix}
a_{11} & a_{12} & a_{13} \\
a_{21} & a_{22} & a_{23} \\
a_{31} & a_{32} & a_{33}
\end{pmatrix}
\]

\section{Tables}
Below is a sample table showing data:

\begin{table}[h]
\centering
\begin{tabular}{@{}lrrr@{}}
\toprule
\textbf{Category} & \textbf{Value 1} & \textbf{Value 2} & \textbf{Total} \\
\midrule
Item A & 10.5 & 20.3 & 30.8 \\
Item B & 15.2 & 25.1 & 40.3 \\
Item C & 22.1 & 12.7 & 34.8 \\
\bottomrule
\end{tabular}
\caption{Sample Data Table}
\label{tab:sample}
\end{table}

\section{Lists}
\subsection{Bullet Points}
\begin{itemize}
  \item First item
  \item Second item
  \item Third item with sub-items:
    \begin{itemize}
      \item Sub-item 1
      \item Sub-item 2
    \end{itemize}
\end{itemize}

\subsection{Numbered List}
\begin{enumerate}
  \item First step
  \item Second step
  \item Third step
\end{enumerate}

\section{Citations and References}
When you use the citation search feature, references will be automatically managed. For example, you might cite a paper like this: \cite{example2024}. All citations will be automatically added to your \texttt{cite.bib} file.

\section{Conclusion}
This template demonstrates basic \LaTeX{} features including mathematical formulas, tables, and lists. The bibliography system is now set up to work seamlessly with Bibby AI's citation management. Feel free to modify it for your needs!

\bibliographystyle{plain}
\bibliography{cite}

\end{document}\documentclass{article}
\usepackage[utf8]{inputenc}
\usepackage{booktabs}
\usepackage{amsmath}
\usepackage{natbib}

\title{My Document}
\author{Your Name}
\date{\today}

\begin{document}

\maketitle

\section{Introduction}
This is a sample document demonstrating various \LaTeX{} features including mathematical equations, tables, and formatting. You can cite papers using the citation search feature, and they will be automatically added to your \texttt{cite.bib} file.

\section{Mathematical Expressions}
Here are some example equations:

\subsection{Quadratic Formula}
The quadratic formula is given by:
\[x = \frac{-b \pm \sqrt{b^2 - 4ac}}{2a}\]

\subsection{Matrix Example}
A sample matrix:
\[
\begin{pmatrix}
a_{11} & a_{12} & a_{13} \\
a_{21} & a_{22} & a_{23} \\
a_{31} & a_{32} & a_{33}
\end{pmatrix}
\]

\section{Tables}
Below is a sample table showing data:

\begin{table}[h]
\centering
\begin{tabular}{@{}lrrr@{}}
\toprule
\textbf{Category} & \textbf{Value 1} & \textbf{Value 2} & \textbf{Total} \\
\midrule
Item A & 10.5 & 20.3 & 30.8 \\
Item B & 15.2 & 25.1 & 40.3 \\
Item C & 22.1 & 12.7 & 34.8 \\
\bottomrule
\end{tabular}
\caption{Sample Data Table}
\label{tab:sample}
\end{table}

\section{Lists}
\subsection{Bullet Points}
\begin{itemize}
  \item First item
  \item Second item
  \item Third item with sub-items:
    \begin{itemize}
      \item Sub-item 1
      \item Sub-item 2
    \end{itemize}
\end{itemize}

\subsection{Numbered List}
\begin{enumerate}
  \item First step
  \item Second step
  \item Third step
\end{enumerate}

\section{Citations and References}
When you use the citation search feature, references will be automatically managed. For example, you might cite a paper like this: \cite{example2024}. All citations will be automatically added to your \texttt{cite.bib} file.

\section{Conclusion}
This template demonstrates basic \LaTeX{} features including mathematical formulas, tables, and lists. The bibliography system is now set up to work seamlessly with Bibby AI's citation management. Feel free to modify it for your needs!

\bibliographystyle{plain}
\bibliography{cite}

\end{document}\documentclass{article}
\usepackage[utf8]{inputenc}
\usepackage{booktabs}
\usepackage{amsmath}
\usepackage{natbib}

\title{My Document}
\author{Your Name}
\date{\today}

\begin{document}

\maketitle

\section{Introduction}
This is a sample document demonstrating various \LaTeX{} features including mathematical equations, tables, and formatting. You can cite papers using the citation search feature, and they will be automatically added to your \texttt{cite.bib} file.

\section{Mathematical Expressions}
Here are some example equations:

\subsection{Quadratic Formula}
The quadratic formula is given by:
\[x = \frac{-b \pm \sqrt{b^2 - 4ac}}{2a}\]

\subsection{Matrix Example}
A sample matrix:
\[
\begin{pmatrix}
a_{11} & a_{12} & a_{13} \\
a_{21} & a_{22} & a_{23} \\
a_{31} & a_{32} & a_{33}
\end{pmatrix}
\]

\section{Tables}
Below is a sample table showing data:

\begin{table}[h]
\centering
\begin{tabular}{@{}lrrr@{}}
\toprule
\textbf{Category} & \textbf{Value 1} & \textbf{Value 2} & \textbf{Total} \\
\midrule
Item A & 10.5 & 20.3 & 30.8 \\
Item B & 15.2 & 25.1 & 40.3 \\
Item C & 22.1 & 12.7 & 34.8 \\
\bottomrule
\end{tabular}
\caption{Sample Data Table}
\label{tab:sample}
\end{table}

\section{Lists}
\subsection{Bullet Points}
\begin{itemize}
  \item First item
  \item Second item
  \item Third item with sub-items:
    \begin{itemize}
      \item Sub-item 1
      \item Sub-item 2
    \end{itemize}
\end{itemize}

\subsection{Numbered List}
\begin{enumerate}
  \item First step
  \item Second step
  \item Third step
\end{enumerate}

\section{Citations and References}
When you use the citation search feature, references will be automatically managed. For example, you might cite a paper like this: \cite{example2024}. All citations will be automatically added to your \texttt{cite.bib} file.

\section{Conclusion}
This template demonstrates basic \LaTeX{} features including mathematical formulas, tables, and lists. The bibliography system is now set up to work seamlessly with Bibby AI's citation management. Feel free to modify it for your needs!

\bibliographystyle{plain}
\bibliography{cite}

\end{document}\documentclass{article}
\usepackage[utf8]{inputenc}
\usepackage{booktabs}
\usepackage{amsmath}
\usepackage{natbib}

\title{My Document}
\author{Your Name}
\date{\today}

\begin{document}

\maketitle

\section{Introduction}
This is a sample document demonstrating various \LaTeX{} features including mathematical equations, tables, and formatting. You can cite papers using the citation search feature, and they will be automatically added to your \texttt{cite.bib} file.

\section{Mathematical Expressions}
Here are some example equations:

\subsection{Quadratic Formula}
The quadratic formula is given by:
\[x = \frac{-b \pm \sqrt{b^2 - 4ac}}{2a}\]

\subsection{Matrix Example}
A sample matrix:
\[
\begin{pmatrix}
a_{11} & a_{12} & a_{13} \\
a_{21} & a_{22} & a_{23} \\
a_{31} & a_{32} & a_{33}
\end{pmatrix}
\]

\section{Tables}
Below is a sample table showing data:

\begin{table}[h]
\centering
\begin{tabular}{@{}lrrr@{}}
\toprule
\textbf{Category} & \textbf{Value 1} & \textbf{Value 2} & \textbf{Total} \\
\midrule
Item A & 10.5 & 20.3 & 30.8 \\
Item B & 15.2 & 25.1 & 40.3 \\
Item C & 22.1 & 12.7 & 34.8 \\
\bottomrule
\end{tabular}
\caption{Sample Data Table}
\label{tab:sample}
\end{table}

\section{Lists}
\subsection{Bullet Points}
\begin{itemize}
  \item First item
  \item Second item
  \item Third item with sub-items:
    \begin{itemize}
      \item Sub-item 1
      \item Sub-item 2
    \end{itemize}
\end{itemize}

\subsection{Numbered List}
\begin{enumerate}
  \item First step
  \item Second step
  \item Third step
\end{enumerate}

\section{Citations and References}
When you use the citation search feature, references will be automatically managed. For example, you might cite a paper like this: \cite{example2024}. All citations will be automatically added to your \texttt{cite.bib} file.

\section{Conclusion}
This template demonstrates basic \LaTeX{} features including mathematical formulas, tables, and lists. The bibliography system is now set up to work seamlessly with Bibby AI's citation management. Feel free to modify it for your needs!

\bibliographystyle{plain}
\bibliography{cite}

\end{document}\documentclass{article}
\usepackage[utf8]{inputenc}
\usepackage{booktabs}
\usepackage{amsmath}
\usepackage{natbib}

\title{My Document}
\author{Your Name}
\date{\today}

\begin{document}

\maketitle

\section{Introduction}
This is a sample document demonstrating various \LaTeX{} features including mathematical equations, tables, and formatting. You can cite papers using the citation search feature, and they will be automatically added to your \texttt{cite.bib} file.

\section{Mathematical Expressions}
Here are some example equations:

\subsection{Quadratic Formula}
The quadratic formula is given by:
\[x = \frac{-b \pm \sqrt{b^2 - 4ac}}{2a}\]

\subsection{Matrix Example}
A sample matrix:
\[
\begin{pmatrix}
a_{11} & a_{12} & a_{13} \\
a_{21} & a_{22} & a_{23} \\
a_{31} & a_{32} & a_{33}
\end{pmatrix}
\]

\section{Tables}
Below is a sample table showing data:

\begin{table}[h]
\centering
\begin{tabular}{@{}lrrr@{}}
\toprule
\textbf{Category} & \textbf{Value 1} & \textbf{Value 2} & \textbf{Total} \\
\midrule
Item A & 10.5 & 20.3 & 30.8 \\
Item B & 15.2 & 25.1 & 40.3 \\
Item C & 22.1 & 12.7 & 34.8 \\
\bottomrule
\end{tabular}
\caption{Sample Data Table}
\label{tab:sample}
\end{table}

\section{Lists}
\subsection{Bullet Points}
\begin{itemize}
  \item First item
  \item Second item
  \item Third item with sub-items:
    \begin{itemize}
      \item Sub-item 1
      \item Sub-item 2
    \end{itemize}
\end{itemize}

\subsection{Numbered List}
\begin{enumerate}
  \item First step
  \item Second step
  \item Third step
\end{enumerate}

\section{Citations and References}
When you use the citation search feature, references will be automatically managed. For example, you might cite a paper like this: \cite{example2024}. All citations will be automatically added to your \texttt{cite.bib} file.

\section{Conclusion}
This template demonstrates basic \LaTeX{} features including mathematical formulas, tables, and lists. The bibliography system is now set up to work seamlessly with Bibby AI's citation management. Feel free to modify it for your needs!

\bibliographystyle{plain}
\bibliography{cite}

\end{document}\documentclass{article}
\usepackage[utf8]{inputenc}
\usepackage{booktabs}
\usepackage{amsmath}
\usepackage{natbib}

\title{My Document}
\author{Your Name}
\date{\today}

\begin{document}

\maketitle

\section{Introduction}
This is a sample document demonstrating various \LaTeX{} features including mathematical equations, tables, and formatting. You can cite papers using the citation search feature, and they will be automatically added to your \texttt{cite.bib} file.

\section{Mathematical Expressions}
Here are some example equations:

\subsection{Quadratic Formula}
The quadratic formula is given by:
\[x = \frac{-b \pm \sqrt{b^2 - 4ac}}{2a}\]

\subsection{Matrix Example}
A sample matrix:
\[
\begin{pmatrix}
a_{11} & a_{12} & a_{13} \\
a_{21} & a_{22} & a_{23} \\
a_{31} & a_{32} & a_{33}
\end{pmatrix}
\]

\section{Tables}
Below is a sample table showing data:

\begin{table}[h]
\centering
\begin{tabular}{@{}lrrr@{}}
\toprule
\textbf{Category} & \textbf{Value 1} & \textbf{Value 2} & \textbf{Total} \\
\midrule
Item A & 10.5 & 20.3 & 30.8 \\
Item B & 15.2 & 25.1 & 40.3 \\
Item C & 22.1 & 12.7 & 34.8 \\
\bottomrule
\end{tabular}
\caption{Sample Data Table}
\label{tab:sample}
\end{table}

\section{Lists}
\subsection{Bullet Points}
\begin{itemize}
  \item First item
  \item Second item
  \item Third item with sub-items:
    \begin{itemize}
      \item Sub-item 1
      \item Sub-item 2
    \end{itemize}
\end{itemize}

\subsection{Numbered List}
\begin{enumerate}
  \item First step
  \item Second step
  \item Third step
\end{enumerate}

\section{Citations and References}
When you use the citation search feature, references will be automatically managed. For example, you might cite a paper like this: \cite{example2024}. All citations will be automatically added to your \texttt{cite.bib} file.

\section{Conclusion}
This template demonstrates basic \LaTeX{} features including mathematical formulas, tables, and lists. The bibliography system is now set up to work seamlessly with Bibby AI's citation management. Feel free to modify it for your needs!

\bibliographystyle{plain}
\bibliography{cite}

\end{document}\documentclass{article}
\usepackage[utf8]{inputenc}
\usepackage{booktabs}
\usepackage{amsmath}
\usepackage{natbib}

\title{My Document}
\author{Your Name}
\date{\today}

\begin{document}

\maketitle

\section{Introduction}
This is a sample document demonstrating various \LaTeX{} features including mathematical equations, tables, and formatting. You can cite papers using the citation search feature, and they will be automatically added to your \texttt{cite.bib} file.

\section{Mathematical Expressions}
Here are some example equations:

\subsection{Quadratic Formula}
The quadratic formula is given by:
\[x = \frac{-b \pm \sqrt{b^2 - 4ac}}{2a}\]

\subsection{Matrix Example}
A sample matrix:
\[
\begin{pmatrix}
a_{11} & a_{12} & a_{13} \\
a_{21} & a_{22} & a_{23} \\
a_{31} & a_{32} & a_{33}
\end{pmatrix}
\]

\section{Tables}
Below is a sample table showing data:

\begin{table}[h]
\centering
\begin{tabular}{@{}lrrr@{}}
\toprule
\textbf{Category} & \textbf{Value 1} & \textbf{Value 2} & \textbf{Total} \\
\midrule
Item A & 10.5 & 20.3 & 30.8 \\
Item B & 15.2 & 25.1 & 40.3 \\
Item C & 22.1 & 12.7 & 34.8 \\
\bottomrule
\end{tabular}
\caption{Sample Data Table}
\label{tab:sample}
\end{table}

\section{Lists}
\subsection{Bullet Points}
\begin{itemize}
  \item First item
  \item Second item
  \item Third item with sub-items:
    \begin{itemize}
      \item Sub-item 1
      \item Sub-item 2
    \end{itemize}
\end{itemize}

\subsection{Numbered List}
\begin{enumerate}
  \item First step
  \item Second step
  \item Third step
\end{enumerate}

\section{Citations and References}
When you use the citation search feature, references will be automatically managed. For example, you might cite a paper like this: \cite{example2024}. All citations will be automatically added to your \texttt{cite.bib} file.

\section{Conclusion}
This template demonstrates basic \LaTeX{} features including mathematical formulas, tables, and lists. The bibliography system is now set up to work seamlessly with Bibby AI's citation management. Feel free to modify it for your needs!

\bibliographystyle{plain}
\bibliography{cite}

\end{document}\documentclass{article}
\usepackage[utf8]{inputenc}
\usepackage{booktabs}
\usepackage{amsmath}
\usepackage{natbib}

\title{My Document}
\author{Your Name}
\date{\today}

\begin{document}

\maketitle

\section{Introduction}
This is a sample document demonstrating various \LaTeX{} features including mathematical equations, tables, and formatting. You can cite papers using the citation search feature, and they will be automatically added to your \texttt{cite.bib} file.

\section{Mathematical Expressions}
Here are some example equations:

\subsection{Quadratic Formula}
The quadratic formula is given by:
\[x = \frac{-b \pm \sqrt{b^2 - 4ac}}{2a}\]

\subsection{Matrix Example}
A sample matrix:
\[
\begin{pmatrix}
a_{11} & a_{12} & a_{13} \\
a_{21} & a_{22} & a_{23} \\
a_{31} & a_{32} & a_{33}
\end{pmatrix}
\]

\section{Tables}
Below is a sample table showing data:

\begin{table}[h]
\centering
\begin{tabular}{@{}lrrr@{}}
\toprule
\textbf{Category} & \textbf{Value 1} & \textbf{Value 2} & \textbf{Total} \\
\midrule
Item A & 10.5 & 20.3 & 30.8 \\
Item B & 15.2 & 25.1 & 40.3 \\
Item C & 22.1 & 12.7 & 34.8 \\
\bottomrule
\end{tabular}
\caption{Sample Data Table}
\label{tab:sample}
\end{table}

\section{Lists}
\subsection{Bullet Points}
\begin{itemize}
  \item First item
  \item Second item
  \item Third item with sub-items:
    \begin{itemize}
      \item Sub-item 1
      \item Sub-item 2
    \end{itemize}
\end{itemize}

\subsection{Numbered List}
\begin{enumerate}
  \item First step
  \item Second step
  \item Third step
\end{enumerate}

\section{Citations and References}
When you use the citation search feature, references will be automatically managed. For example, you might cite a paper like this: \cite{example2024}. All citations will be automatically added to your \texttt{cite.bib} file.

\section{Conclusion}
This template demonstrates basic \LaTeX{} features including mathematical formulas, tables, and lists. The bibliography system is now set up to work seamlessly with Bibby AI's citation management. Feel free to modify it for your needs!

\bibliographystyle{plain}
\bibliography{cite}

\end{document}\documentclass{article}
\usepackage[utf8]{inputenc}
\usepackage{booktabs}
\usepackage{amsmath}
\usepackage{natbib}

\title{My Document}
\author{Your Name}
\date{\today}

\begin{document}

\maketitle

\section{Introduction}
This is a sample document demonstrating various \LaTeX{} features including mathematical equations, tables, and formatting. You can cite papers using the citation search feature, and they will be automatically added to your \texttt{cite.bib} file.

\section{Mathematical Expressions}
Here are some example equations:

\subsection{Quadratic Formula}
The quadratic formula is given by:
\[x = \frac{-b \pm \sqrt{b^2 - 4ac}}{2a}\]

\subsection{Matrix Example}
A sample matrix:
\[
\begin{pmatrix}
a_{11} & a_{12} & a_{13} \\
a_{21} & a_{22} & a_{23} \\
a_{31} & a_{32} & a_{33}
\end{pmatrix}
\]

\section{Tables}
Below is a sample table showing data:

\begin{table}[h]
\centering
\begin{tabular}{@{}lrrr@{}}
\toprule
\textbf{Category} & \textbf{Value 1} & \textbf{Value 2} & \textbf{Total} \\
\midrule
Item A & 10.5 & 20.3 & 30.8 \\
Item B & 15.2 & 25.1 & 40.3 \\
Item C & 22.1 & 12.7 & 34.8 \\
\bottomrule
\end{tabular}
\caption{Sample Data Table}
\label{tab:sample}
\end{table}

\section{Lists}
\subsection{Bullet Points}
\begin{itemize}
  \item First item
  \item Second item
  \item Third item with sub-items:
    \begin{itemize}
      \item Sub-item 1
      \item Sub-item 2
    \end{itemize}
\end{itemize}

\subsection{Numbered List}
\begin{enumerate}
  \item First step
  \item Second step
  \item Third step
\end{enumerate}

\section{Citations and References}
When you use the citation search feature, references will be automatically managed. For example, you might cite a paper like this: \cite{example2024}. All citations will be automatically added to your \texttt{cite.bib} file.

\section{Conclusion}
This template demonstrates basic \LaTeX{} features including mathematical formulas, tables, and lists. The bibliography system is now set up to work seamlessly with Bibby AI's citation management. Feel free to modify it for your needs!

\bibliographystyle{plain}
\bibliography{cite}

\end{document}\documentclass{article}
\usepackage[utf8]{inputenc}
\usepackage{booktabs}
\usepackage{amsmath}
\usepackage{natbib}

\title{My Document}
\author{Your Name}
\date{\today}

\begin{document}

\maketitle

\section{Introduction}
This is a sample document demonstrating various \LaTeX{} features including mathematical equations, tables, and formatting. You can cite papers using the citation search feature, and they will be automatically added to your \texttt{cite.bib} file.

\section{Mathematical Expressions}
Here are some example equations:

\subsection{Quadratic Formula}
The quadratic formula is given by:
\[x = \frac{-b \pm \sqrt{b^2 - 4ac}}{2a}\]

\subsection{Matrix Example}
A sample matrix:
\[
\begin{pmatrix}
a_{11} & a_{12} & a_{13} \\
a_{21} & a_{22} & a_{23} \\
a_{31} & a_{32} & a_{33}
\end{pmatrix}
\]

\section{Tables}
Below is a sample table showing data:

\begin{table}[h]
\centering
\begin{tabular}{@{}lrrr@{}}
\toprule
\textbf{Category} & \textbf{Value 1} & \textbf{Value 2} & \textbf{Total} \\
\midrule
Item A & 10.5 & 20.3 & 30.8 \\
Item B & 15.2 & 25.1 & 40.3 \\
Item C & 22.1 & 12.7 & 34.8 \\
\bottomrule
\end{tabular}
\caption{Sample Data Table}
\label{tab:sample}
\end{table}

\section{Lists}
\subsection{Bullet Points}
\begin{itemize}
  \item First item
  \item Second item
  \item Third item with sub-items:
    \begin{itemize}
      \item Sub-item 1
      \item Sub-item 2
    \end{itemize}
\end{itemize}

\subsection{Numbered List}
\begin{enumerate}
  \item First step
  \item Second step
  \item Third step
\end{enumerate}

\section{Citations and References}
When you use the citation search feature, references will be automatically managed. For example, you might cite a paper like this: \cite{example2024}. All citations will be automatically added to your \texttt{cite.bib} file.

\section{Conclusion}
This template demonstrates basic \LaTeX{} features including mathematical formulas, tables, and lists. The bibliography system is now set up to work seamlessly with Bibby AI's citation management. Feel free to modify it for your needs!

\bibliographystyle{plain}
\bibliography{cite}

\end{document}\documentclass{article}
\usepackage[utf8]{inputenc}
\usepackage{booktabs}
\usepackage{amsmath}
\usepackage{natbib}

\title{My Document}
\author{Your Name}
\date{\today}

\begin{document}

\maketitle

\section{Introduction}
This is a sample document demonstrating various \LaTeX{} features including mathematical equations, tables, and formatting. You can cite papers using the citation search feature, and they will be automatically added to your \texttt{cite.bib} file.

\section{Mathematical Expressions}
Here are some example equations:

\subsection{Quadratic Formula}
The quadratic formula is given by:
\[x = \frac{-b \pm \sqrt{b^2 - 4ac}}{2a}\]

\subsection{Matrix Example}
A sample matrix:
\[
\begin{pmatrix}
a_{11} & a_{12} & a_{13} \\
a_{21} & a_{22} & a_{23} \\
a_{31} & a_{32} & a_{33}
\end{pmatrix}
\]

\section{Tables}
Below is a sample table showing data:

\begin{table}[h]
\centering
\begin{tabular}{@{}lrrr@{}}
\toprule
\textbf{Category} & \textbf{Value 1} & \textbf{Value 2} & \textbf{Total} \\
\midrule
Item A & 10.5 & 20.3 & 30.8 \\
Item B & 15.2 & 25.1 & 40.3 \\
Item C & 22.1 & 12.7 & 34.8 \\
\bottomrule
\end{tabular}
\caption{Sample Data Table}
\label{tab:sample}
\end{table}

\section{Lists}
\subsection{Bullet Points}
\begin{itemize}
  \item First item
  \item Second item
  \item Third item with sub-items:
    \begin{itemize}
      \item Sub-item 1
      \item Sub-item 2
    \end{itemize}
\end{itemize}

\subsection{Numbered List}
\begin{enumerate}
  \item First step
  \item Second step
  \item Third step
\end{enumerate}

\section{Citations and References}
When you use the citation search feature, references will be automatically managed. For example, you might cite a paper like this: \cite{example2024}. All citations will be automatically added to your \texttt{cite.bib} file.

\section{Conclusion}
This template demonstrates basic \LaTeX{} features including mathematical formulas, tables, and lists. The bibliography system is now set up to work seamlessly with Bibby AI's citation management. Feel free to modify it for your needs!

\bibliographystyle{plain}
\bibliography{cite}

\end{document}\documentclass{article}
\usepackage[utf8]{inputenc}
\usepackage{booktabs}
\usepackage{amsmath}
\usepackage{natbib}

\title{My Document}
\author{Your Name}
\date{\today}

\begin{document}

\maketitle

\section{Introduction}
This is a sample document demonstrating various \LaTeX{} features including mathematical equations, tables, and formatting. You can cite papers using the citation search feature, and they will be automatically added to your \texttt{cite.bib} file.

\section{Mathematical Expressions}
Here are some example equations:

\subsection{Quadratic Formula}
The quadratic formula is given by:
\[x = \frac{-b \pm \sqrt{b^2 - 4ac}}{2a}\]

\subsection{Matrix Example}
A sample matrix:
\[
\begin{pmatrix}
a_{11} & a_{12} & a_{13} \\
a_{21} & a_{22} & a_{23} \\
a_{31} & a_{32} & a_{33}
\end{pmatrix}
\]

\section{Tables}
Below is a sample table showing data:

\begin{table}[h]
\centering
\begin{tabular}{@{}lrrr@{}}
\toprule
\textbf{Category} & \textbf{Value 1} & \textbf{Value 2} & \textbf{Total} \\
\midrule
Item A & 10.5 & 20.3 & 30.8 \\
Item B & 15.2 & 25.1 & 40.3 \\
Item C & 22.1 & 12.7 & 34.8 \\
\bottomrule
\end{tabular}
\caption{Sample Data Table}
\label{tab:sample}
\end{table}

\section{Lists}
\subsection{Bullet Points}
\begin{itemize}
  \item First item
  \item Second item
  \item Third item with sub-items:
    \begin{itemize}
      \item Sub-item 1
      \item Sub-item 2
    \end{itemize}
\end{itemize}

\subsection{Numbered List}
\begin{enumerate}
  \item First step
  \item Second step
  \item Third step
\end{enumerate}

\section{Citations and References}
When you use the citation search feature, references will be automatically managed. For example, you might cite a paper like this: \cite{example2024}. All citations will be automatically added to your \texttt{cite.bib} file.

\section{Conclusion}
This template demonstrates basic \LaTeX{} features including mathematical formulas, tables, and lists. The bibliography system is now set up to work seamlessly with Bibby AI's citation management. Feel free to modify it for your needs!

\bibliographystyle{plain}
\bibliography{cite}

\end{document}\documentclass{article}
\usepackage[utf8]{inputenc}
\usepackage{booktabs}
\usepackage{amsmath}
\usepackage{natbib}

\title{My Document}
\author{Your Name}
\date{\today}

\begin{document}

\maketitle

\section{Introduction}
This is a sample document demonstrating various \LaTeX{} features including mathematical equations, tables, and formatting. You can cite papers using the citation search feature, and they will be automatically added to your \texttt{cite.bib} file.

\section{Mathematical Expressions}
Here are some example equations:

\subsection{Quadratic Formula}
The quadratic formula is given by:
\[x = \frac{-b \pm \sqrt{b^2 - 4ac}}{2a}\]

\subsection{Matrix Example}
A sample matrix:
\[
\begin{pmatrix}
a_{11} & a_{12} & a_{13} \\
a_{21} & a_{22} & a_{23} \\
a_{31} & a_{32} & a_{33}
\end{pmatrix}
\]

\section{Tables}
Below is a sample table showing data:

\begin{table}[h]
\centering
\begin{tabular}{@{}lrrr@{}}
\toprule
\textbf{Category} & \textbf{Value 1} & \textbf{Value 2} & \textbf{Total} \\
\midrule
Item A & 10.5 & 20.3 & 30.8 \\
Item B & 15.2 & 25.1 & 40.3 \\
Item C & 22.1 & 12.7 & 34.8 \\
\bottomrule
\end{tabular}
\caption{Sample Data Table}
\label{tab:sample}
\end{table}

\section{Lists}
\subsection{Bullet Points}
\begin{itemize}
  \item First item
  \item Second item
  \item Third item with sub-items:
    \begin{itemize}
      \item Sub-item 1
      \item Sub-item 2
    \end{itemize}
\end{itemize}

\subsection{Numbered List}
\begin{enumerate}
  \item First step
  \item Second step
  \item Third step
\end{enumerate}

\section{Citations and References}
When you use the citation search feature, references will be automatically managed. For example, you might cite a paper like this: \cite{example2024}. All citations will be automatically added to your \texttt{cite.bib} file.

\section{Conclusion}
This template demonstrates basic \LaTeX{} features including mathematical formulas, tables, and lists. The bibliography system is now set up to work seamlessly with Bibby AI's citation management. Feel free to modify it for your needs!

\bibliographystyle{plain}
\bibliography{cite}

\end{document}\documentclass{article}
\usepackage[utf8]{inputenc}
\usepackage{booktabs}
\usepackage{amsmath}
\usepackage{natbib}

\title{My Document}
\author{Your Name}
\date{\today}

\begin{document}

\maketitle

\section{Introduction}
This is a sample document demonstrating various \LaTeX{} features including mathematical equations, tables, and formatting. You can cite papers using the citation search feature, and they will be automatically added to your \texttt{cite.bib} file.

\section{Mathematical Expressions}
Here are some example equations:

\subsection{Quadratic Formula}
The quadratic formula is given by:
\[x = \frac{-b \pm \sqrt{b^2 - 4ac}}{2a}\]

\subsection{Matrix Example}
A sample matrix:
\[
\begin{pmatrix}
a_{11} & a_{12} & a_{13} \\
a_{21} & a_{22} & a_{23} \\
a_{31} & a_{32} & a_{33}
\end{pmatrix}
\]

\section{Tables}
Below is a sample table showing data:

\begin{table}[h]
\centering
\begin{tabular}{@{}lrrr@{}}
\toprule
\textbf{Category} & \textbf{Value 1} & \textbf{Value 2} & \textbf{Total} \\
\midrule
Item A & 10.5 & 20.3 & 30.8 \\
Item B & 15.2 & 25.1 & 40.3 \\
Item C & 22.1 & 12.7 & 34.8 \\
\bottomrule
\end{tabular}
\caption{Sample Data Table}
\label{tab:sample}
\end{table}

\section{Lists}
\subsection{Bullet Points}
\begin{itemize}
  \item First item
  \item Second item
  \item Third item with sub-items:
    \begin{itemize}
      \item Sub-item 1
      \item Sub-item 2
    \end{itemize}
\end{itemize}

\subsection{Numbered List}
\begin{enumerate}
  \item First step
  \item Second step
  \item Third step
\end{enumerate}

\section{Citations and References}
When you use the citation search feature, references will be automatically managed. For example, you might cite a paper like this: \cite{example2024}. All citations will be automatically added to your \texttt{cite.bib} file.

\section{Conclusion}
This template demonstrates basic \LaTeX{} features including mathematical formulas, tables, and lists. The bibliography system is now set up to work seamlessly with Bibby AI's citation management. Feel free to modify it for your needs!

\bibliographystyle{plain}
\bibliography{cite}

\end{document}\documentclass{article}
\usepackage[utf8]{inputenc}
\usepackage{booktabs}
\usepackage{amsmath}
\usepackage{natbib}

\title{My Document}
\author{Your Name}
\date{\today}

\begin{document}

\maketitle

\section{Introduction}
This is a sample document demonstrating various \LaTeX{} features including mathematical equations, tables, and formatting. You can cite papers using the citation search feature, and they will be automatically added to your \texttt{cite.bib} file.

\section{Mathematical Expressions}
Here are some example equations:

\subsection{Quadratic Formula}
The quadratic formula is given by:
\[x = \frac{-b \pm \sqrt{b^2 - 4ac}}{2a}\]

\subsection{Matrix Example}
A sample matrix:
\[
\begin{pmatrix}
a_{11} & a_{12} & a_{13} \\
a_{21} & a_{22} & a_{23} \\
a_{31} & a_{32} & a_{33}
\end{pmatrix}
\]

\section{Tables}
Below is a sample table showing data:

\begin{table}[h]
\centering
\begin{tabular}{@{}lrrr@{}}
\toprule
\textbf{Category} & \textbf{Value 1} & \textbf{Value 2} & \textbf{Total} \\
\midrule
Item A & 10.5 & 20.3 & 30.8 \\
Item B & 15.2 & 25.1 & 40.3 \\
Item C & 22.1 & 12.7 & 34.8 \\
\bottomrule
\end{tabular}
\caption{Sample Data Table}
\label{tab:sample}
\end{table}

\section{Lists}
\subsection{Bullet Points}
\begin{itemize}
  \item First item
  \item Second item
  \item Third item with sub-items:
    \begin{itemize}
      \item Sub-item 1
      \item Sub-item 2
    \end{itemize}
\end{itemize}

\subsection{Numbered List}
\begin{enumerate}
  \item First step
  \item Second step
  \item Third step
\end{enumerate}

\section{Citations and References}
When you use the citation search feature, references will be automatically managed. For example, you might cite a paper like this: \cite{example2024}. All citations will be automatically added to your \texttt{cite.bib} file.

\section{Conclusion}
This template demonstrates basic \LaTeX{} features including mathematical formulas, tables, and lists. The bibliography system is now set up to work seamlessly with Bibby AI's citation management. Feel free to modify it for your needs!

\bibliographystyle{plain}
\bibliography{cite}

\end{document}\documentclass{article}
\usepackage[utf8]{inputenc}
\usepackage{booktabs}
\usepackage{amsmath}
\usepackage{natbib}

\title{My Document}
\author{Your Name}
\date{\today}

\begin{document}

\maketitle

\section{Introduction}
This is a sample document demonstrating various \LaTeX{} features including mathematical equations, tables, and formatting. You can cite papers using the citation search feature, and they will be automatically added to your \texttt{cite.bib} file.

\section{Mathematical Expressions}
Here are some example equations:

\subsection{Quadratic Formula}
The quadratic formula is given by:
\[x = \frac{-b \pm \sqrt{b^2 - 4ac}}{2a}\]

\subsection{Matrix Example}
A sample matrix:
\[
\begin{pmatrix}
a_{11} & a_{12} & a_{13} \\
a_{21} & a_{22} & a_{23} \\
a_{31} & a_{32} & a_{33}
\end{pmatrix}
\]

\section{Tables}
Below is a sample table showing data:

\begin{table}[h]
\centering
\begin{tabular}{@{}lrrr@{}}
\toprule
\textbf{Category} & \textbf{Value 1} & \textbf{Value 2} & \textbf{Total} \\
\midrule
Item A & 10.5 & 20.3 & 30.8 \\
Item B & 15.2 & 25.1 & 40.3 \\
Item C & 22.1 & 12.7 & 34.8 \\
\bottomrule
\end{tabular}
\caption{Sample Data Table}
\label{tab:sample}
\end{table}

\section{Lists}
\subsection{Bullet Points}
\begin{itemize}
  \item First item
  \item Second item
  \item Third item with sub-items:
    \begin{itemize}
      \item Sub-item 1
      \item Sub-item 2
    \end{itemize}
\end{itemize}

\subsection{Numbered List}
\begin{enumerate}
  \item First step
  \item Second step
  \item Third step
\end{enumerate}

\section{Citations and References}
When you use the citation search feature, references will be automatically managed. For example, you might cite a paper like this: \cite{example2024}. All citations will be automatically added to your \texttt{cite.bib} file.

\section{Conclusion}
This template demonstrates basic \LaTeX{} features including mathematical formulas, tables, and lists. The bibliography system is now set up to work seamlessly with Bibby AI's citation management. Feel free to modify it for your needs!

\bibliographystyle{plain}
\bibliography{cite}

\end{document}\documentclass{article}
\usepackage[utf8]{inputenc}
\usepackage{booktabs}
\usepackage{amsmath}
\usepackage{natbib}

\title{My Document}
\author{Your Name}
\date{\today}

\begin{document}

\maketitle

\section{Introduction}
This is a sample document demonstrating various \LaTeX{} features including mathematical equations, tables, and formatting. You can cite papers using the citation search feature, and they will be automatically added to your \texttt{cite.bib} file.

\section{Mathematical Expressions}
Here are some example equations:

\subsection{Quadratic Formula}
The quadratic formula is given by:
\[x = \frac{-b \pm \sqrt{b^2 - 4ac}}{2a}\]

\subsection{Matrix Example}
A sample matrix:
\[
\begin{pmatrix}
a_{11} & a_{12} & a_{13} \\
a_{21} & a_{22} & a_{23} \\
a_{31} & a_{32} & a_{33}
\end{pmatrix}
\]

\section{Tables}
Below is a sample table showing data:

\begin{table}[h]
\centering
\begin{tabular}{@{}lrrr@{}}
\toprule
\textbf{Category} & \textbf{Value 1} & \textbf{Value 2} & \textbf{Total} \\
\midrule
Item A & 10.5 & 20.3 & 30.8 \\
Item B & 15.2 & 25.1 & 40.3 \\
Item C & 22.1 & 12.7 & 34.8 \\
\bottomrule
\end{tabular}
\caption{Sample Data Table}
\label{tab:sample}
\end{table}

\section{Lists}
\subsection{Bullet Points}
\begin{itemize}
  \item First item
  \item Second item
  \item Third item with sub-items:
    \begin{itemize}
      \item Sub-item 1
      \item Sub-item 2
    \end{itemize}
\end{itemize}

\subsection{Numbered List}
\begin{enumerate}
  \item First step
  \item Second step
  \item Third step
\end{enumerate}

\section{Citations and References}
When you use the citation search feature, references will be automatically managed. For example, you might cite a paper like this: \cite{example2024}. All citations will be automatically added to your \texttt{cite.bib} file.

\section{Conclusion}
This template demonstrates basic \LaTeX{} features including mathematical formulas, tables, and lists. The bibliography system is now set up to work seamlessly with Bibby AI's citation management. Feel free to modify it for your needs!

\bibliographystyle{plain}
\bibliography{cite}

\end{document}\documentclass{article}
\usepackage[utf8]{inputenc}
\usepackage{booktabs}
\usepackage{amsmath}
\usepackage{natbib}

\title{My Document}
\author{Your Name}
\date{\today}

\begin{document}

\maketitle

\section{Introduction}
This is a sample document demonstrating various \LaTeX{} features including mathematical equations, tables, and formatting. You can cite papers using the citation search feature, and they will be automatically added to your \texttt{cite.bib} file.

\section{Mathematical Expressions}
Here are some example equations:

\subsection{Quadratic Formula}
The quadratic formula is given by:
\[x = \frac{-b \pm \sqrt{b^2 - 4ac}}{2a}\]

\subsection{Matrix Example}
A sample matrix:
\[
\begin{pmatrix}
a_{11} & a_{12} & a_{13} \\
a_{21} & a_{22} & a_{23} \\
a_{31} & a_{32} & a_{33}
\end{pmatrix}
\]

\section{Tables}
Below is a sample table showing data:

\begin{table}[h]
\centering
\begin{tabular}{@{}lrrr@{}}
\toprule
\textbf{Category} & \textbf{Value 1} & \textbf{Value 2} & \textbf{Total} \\
\midrule
Item A & 10.5 & 20.3 & 30.8 \\
Item B & 15.2 & 25.1 & 40.3 \\
Item C & 22.1 & 12.7 & 34.8 \\
\bottomrule
\end{tabular}
\caption{Sample Data Table}
\label{tab:sample}
\end{table}

\section{Lists}
\subsection{Bullet Points}
\begin{itemize}
  \item First item
  \item Second item
  \item Third item with sub-items:
    \begin{itemize}
      \item Sub-item 1
      \item Sub-item 2
    \end{itemize}
\end{itemize}

\subsection{Numbered List}
\begin{enumerate}
  \item First step
  \item Second step
  \item Third step
\end{enumerate}

\section{Citations and References}
When you use the citation search feature, references will be automatically managed. For example, you might cite a paper like this: \cite{example2024}. All citations will be automatically added to your \texttt{cite.bib} file.

\section{Conclusion}
This template demonstrates basic \LaTeX{} features including mathematical formulas, tables, and lists. The bibliography system is now set up to work seamlessly with Bibby AI's citation management. Feel free to modify it for your needs!

\bibliographystyle{plain}
\bibliography{cite}

\end{document}\documentclass{article}
\usepackage[utf8]{inputenc}
\usepackage{booktabs}
\usepackage{amsmath}
\usepackage{natbib}

\title{My Document}
\author{Your Name}
\date{\today}

\begin{document}

\maketitle

\section{Introduction}
This is a sample document demonstrating various \LaTeX{} features including mathematical equations, tables, and formatting. You can cite papers using the citation search feature, and they will be automatically added to your \texttt{cite.bib} file.

\section{Mathematical Expressions}
Here are some example equations:

\subsection{Quadratic Formula}
The quadratic formula is given by:
\[x = \frac{-b \pm \sqrt{b^2 - 4ac}}{2a}\]

\subsection{Matrix Example}
A sample matrix:
\[
\begin{pmatrix}
a_{11} & a_{12} & a_{13} \\
a_{21} & a_{22} & a_{23} \\
a_{31} & a_{32} & a_{33}
\end{pmatrix}
\]

\section{Tables}
Below is a sample table showing data:

\begin{table}[h]
\centering
\begin{tabular}{@{}lrrr@{}}
\toprule
\textbf{Category} & \textbf{Value 1} & \textbf{Value 2} & \textbf{Total} \\
\midrule
Item A & 10.5 & 20.3 & 30.8 \\
Item B & 15.2 & 25.1 & 40.3 \\
Item C & 22.1 & 12.7 & 34.8 \\
\bottomrule
\end{tabular}
\caption{Sample Data Table}
\label{tab:sample}
\end{table}

\section{Lists}
\subsection{Bullet Points}
\begin{itemize}
  \item First item
  \item Second item
  \item Third item with sub-items:
    \begin{itemize}
      \item Sub-item 1
      \item Sub-item 2
    \end{itemize}
\end{itemize}

\subsection{Numbered List}
\begin{enumerate}
  \item First step
  \item Second step
  \item Third step
\end{enumerate}

\section{Citations and References}
When you use the citation search feature, references will be automatically managed. For example, you might cite a paper like this: \cite{example2024}. All citations will be automatically added to your \texttt{cite.bib} file.

\section{Conclusion}
This template demonstrates basic \LaTeX{} features including mathematical formulas, tables, and lists. The bibliography system is now set up to work seamlessly with Bibby AI's citation management. Feel free to modify it for your needs!

\bibliographystyle{plain}
\bibliography{cite}

\end{document}\documentclass{article}
\usepackage[utf8]{inputenc}
\usepackage{booktabs}
\usepackage{amsmath}
\usepackage{natbib}

\title{My Document}
\author{Your Name}
\date{\today}

\begin{document}

\maketitle

\section{Introduction}
This is a sample document demonstrating various \LaTeX{} features including mathematical equations, tables, and formatting. You can cite papers using the citation search feature, and they will be automatically added to your \texttt{cite.bib} file.

\section{Mathematical Expressions}
Here are some example equations:

\subsection{Quadratic Formula}
The quadratic formula is given by:
\[x = \frac{-b \pm \sqrt{b^2 - 4ac}}{2a}\]

\subsection{Matrix Example}
A sample matrix:
\[
\begin{pmatrix}
a_{11} & a_{12} & a_{13} \\
a_{21} & a_{22} & a_{23} \\
a_{31} & a_{32} & a_{33}
\end{pmatrix}
\]

\section{Tables}
Below is a sample table showing data:

\begin{table}[h]
\centering
\begin{tabular}{@{}lrrr@{}}
\toprule
\textbf{Category} & \textbf{Value 1} & \textbf{Value 2} & \textbf{Total} \\
\midrule
Item A & 10.5 & 20.3 & 30.8 \\
Item B & 15.2 & 25.1 & 40.3 \\
Item C & 22.1 & 12.7 & 34.8 \\
\bottomrule
\end{tabular}
\caption{Sample Data Table}
\label{tab:sample}
\end{table}

\section{Lists}
\subsection{Bullet Points}
\begin{itemize}
  \item First item
  \item Second item
  \item Third item with sub-items:
    \begin{itemize}
      \item Sub-item 1
      \item Sub-item 2
    \end{itemize}
\end{itemize}

\subsection{Numbered List}
\begin{enumerate}
  \item First step
  \item Second step
  \item Third step
\end{enumerate}

\section{Citations and References}
When you use the citation search feature, references will be automatically managed. For example, you might cite a paper like this: \cite{example2024}. All citations will be automatically added to your \texttt{cite.bib} file.

\section{Conclusion}
This template demonstrates basic \LaTeX{} features including mathematical formulas, tables, and lists. The bibliography system is now set up to work seamlessly with Bibby AI's citation management. Feel free to modify it for your needs!

\bibliographystyle{plain}
\bibliography{cite}

\end{document}\documentclass{article}
\usepackage[utf8]{inputenc}
\usepackage{booktabs}
\usepackage{amsmath}
\usepackage{natbib}

\title{My Document}
\author{Your Name}
\date{\today}

\begin{document}

\maketitle

\section{Introduction}
This is a sample document demonstrating various \LaTeX{} features including mathematical equations, tables, and formatting. You can cite papers using the citation search feature, and they will be automatically added to your \texttt{cite.bib} file.

\section{Mathematical Expressions}
Here are some example equations:

\subsection{Quadratic Formula}
The quadratic formula is given by:
\[x = \frac{-b \pm \sqrt{b^2 - 4ac}}{2a}\]

\subsection{Matrix Example}
A sample matrix:
\[
\begin{pmatrix}
a_{11} & a_{12} & a_{13} \\
a_{21} & a_{22} & a_{23} \\
a_{31} & a_{32} & a_{33}
\end{pmatrix}
\]

\section{Tables}
Below is a sample table showing data:

\begin{table}[h]
\centering
\begin{tabular}{@{}lrrr@{}}
\toprule
\textbf{Category} & \textbf{Value 1} & \textbf{Value 2} & \textbf{Total} \\
\midrule
Item A & 10.5 & 20.3 & 30.8 \\
Item B & 15.2 & 25.1 & 40.3 \\
Item C & 22.1 & 12.7 & 34.8 \\
\bottomrule
\end{tabular}
\caption{Sample Data Table}
\label{tab:sample}
\end{table}

\section{Lists}
\subsection{Bullet Points}
\begin{itemize}
  \item First item
  \item Second item
  \item Third item with sub-items:
    \begin{itemize}
      \item Sub-item 1
      \item Sub-item 2
    \end{itemize}
\end{itemize}

\subsection{Numbered List}
\begin{enumerate}
  \item First step
  \item Second step
  \item Third step
\end{enumerate}

\section{Citations and References}
When you use the citation search feature, references will be automatically managed. For example, you might cite a paper like this: \cite{example2024}. All citations will be automatically added to your \texttt{cite.bib} file.

\section{Conclusion}
This template demonstrates basic \LaTeX{} features including mathematical formulas, tables, and lists. The bibliography system is now set up to work seamlessly with Bibby AI's citation management. Feel free to modify it for your needs!

\bibliographystyle{plain}
\bibliography{cite}

\end{document}\documentclass{article}
\usepackage[utf8]{inputenc}
\usepackage{booktabs}
\usepackage{amsmath}
\usepackage{natbib}

\title{My Document}
\author{Your Name}
\date{\today}

\begin{document}

\maketitle

\section{Introduction}
This is a sample document demonstrating various \LaTeX{} features including mathematical equations, tables, and formatting. You can cite papers using the citation search feature, and they will be automatically added to your \texttt{cite.bib} file.

\section{Mathematical Expressions}
Here are some example equations:

\subsection{Quadratic Formula}
The quadratic formula is given by:
\[x = \frac{-b \pm \sqrt{b^2 - 4ac}}{2a}\]

\subsection{Matrix Example}
A sample matrix:
\[
\begin{pmatrix}
a_{11} & a_{12} & a_{13} \\
a_{21} & a_{22} & a_{23} \\
a_{31} & a_{32} & a_{33}
\end{pmatrix}
\]

\section{Tables}
Below is a sample table showing data:

\begin{table}[h]
\centering
\begin{tabular}{@{}lrrr@{}}
\toprule
\textbf{Category} & \textbf{Value 1} & \textbf{Value 2} & \textbf{Total} \\
\midrule
Item A & 10.5 & 20.3 & 30.8 \\
Item B & 15.2 & 25.1 & 40.3 \\
Item C & 22.1 & 12.7 & 34.8 \\
\bottomrule
\end{tabular}
\caption{Sample Data Table}
\label{tab:sample}
\end{table}

\section{Lists}
\subsection{Bullet Points}
\begin{itemize}
  \item First item
  \item Second item
  \item Third item with sub-items:
    \begin{itemize}
      \item Sub-item 1
      \item Sub-item 2
    \end{itemize}
\end{itemize}

\subsection{Numbered List}
\begin{enumerate}
  \item First step
  \item Second step
  \item Third step
\end{enumerate}

\section{Citations and References}
When you use the citation search feature, references will be automatically managed. For example, you might cite a paper like this: \cite{example2024}. All citations will be automatically added to your \texttt{cite.bib} file.

\section{Conclusion}
This template demonstrates basic \LaTeX{} features including mathematical formulas, tables, and lists. The bibliography system is now set up to work seamlessly with Bibby AI's citation management. Feel free to modify it for your needs!

\bibliographystyle{plain}
\bibliography{cite}

\end{document}\documentclass{article}
\usepackage[utf8]{inputenc}
\usepackage{booktabs}
\usepackage{amsmath}
\usepackage{natbib}

\title{My Document}
\author{Your Name}
\date{\today}

\begin{document}

\maketitle

\section{Introduction}
This is a sample document demonstrating various \LaTeX{} features including mathematical equations, tables, and formatting. You can cite papers using the citation search feature, and they will be automatically added to your \texttt{cite.bib} file.

\section{Mathematical Expressions}
Here are some example equations:

\subsection{Quadratic Formula}
The quadratic formula is given by:
\[x = \frac{-b \pm \sqrt{b^2 - 4ac}}{2a}\]

\subsection{Matrix Example}
A sample matrix:
\[
\begin{pmatrix}
a_{11} & a_{12} & a_{13} \\
a_{21} & a_{22} & a_{23} \\
a_{31} & a_{32} & a_{33}
\end{pmatrix}
\]

\section{Tables}
Below is a sample table showing data:

\begin{table}[h]
\centering
\begin{tabular}{@{}lrrr@{}}
\toprule
\textbf{Category} & \textbf{Value 1} & \textbf{Value 2} & \textbf{Total} \\
\midrule
Item A & 10.5 & 20.3 & 30.8 \\
Item B & 15.2 & 25.1 & 40.3 \\
Item C & 22.1 & 12.7 & 34.8 \\
\bottomrule
\end{tabular}
\caption{Sample Data Table}
\label{tab:sample}
\end{table}

\section{Lists}
\subsection{Bullet Points}
\begin{itemize}
  \item First item
  \item Second item
  \item Third item with sub-items:
    \begin{itemize}
      \item Sub-item 1
      \item Sub-item 2
    \end{itemize}
\end{itemize}

\subsection{Numbered List}
\begin{enumerate}
  \item First step
  \item Second step
  \item Third step
\end{enumerate}

\section{Citations and References}
When you use the citation search feature, references will be automatically managed. For example, you might cite a paper like this: \cite{example2024}. All citations will be automatically added to your \texttt{cite.bib} file.

\section{Conclusion}
This template demonstrates basic \LaTeX{} features including mathematical formulas, tables, and lists. The bibliography system is now set up to work seamlessly with Bibby AI's citation management. Feel free to modify it for your needs!

\bibliographystyle{plain}
\bibliography{cite}

\end{document}\documentclass{article}
\usepackage[utf8]{inputenc}
\usepackage{booktabs}
\usepackage{amsmath}
\usepackage{natbib}

\title{My Document}
\author{Your Name}
\date{\today}

\begin{document}

\maketitle

\section{Introduction}
This is a sample document demonstrating various \LaTeX{} features including mathematical equations, tables, and formatting. You can cite papers using the citation search feature, and they will be automatically added to your \texttt{cite.bib} file.

\section{Mathematical Expressions}
Here are some example equations:

\subsection{Quadratic Formula}
The quadratic formula is given by:
\[x = \frac{-b \pm \sqrt{b^2 - 4ac}}{2a}\]

\subsection{Matrix Example}
A sample matrix:
\[
\begin{pmatrix}
a_{11} & a_{12} & a_{13} \\
a_{21} & a_{22} & a_{23} \\
a_{31} & a_{32} & a_{33}
\end{pmatrix}
\]

\section{Tables}
Below is a sample table showing data:

\begin{table}[h]
\centering
\begin{tabular}{@{}lrrr@{}}
\toprule
\textbf{Category} & \textbf{Value 1} & \textbf{Value 2} & \textbf{Total} \\
\midrule
Item A & 10.5 & 20.3 & 30.8 \\
Item B & 15.2 & 25.1 & 40.3 \\
Item C & 22.1 & 12.7 & 34.8 \\
\bottomrule
\end{tabular}
\caption{Sample Data Table}
\label{tab:sample}
\end{table}

\section{Lists}
\subsection{Bullet Points}
\begin{itemize}
  \item First item
  \item Second item
  \item Third item with sub-items:
    \begin{itemize}
      \item Sub-item 1
      \item Sub-item 2
    \end{itemize}
\end{itemize}

\subsection{Numbered List}
\begin{enumerate}
  \item First step
  \item Second step
  \item Third step
\end{enumerate}

\section{Citations and References}
When you use the citation search feature, references will be automatically managed. For example, you might cite a paper like this: \cite{example2024}. All citations will be automatically added to your \texttt{cite.bib} file.

\section{Conclusion}
This template demonstrates basic \LaTeX{} features including mathematical formulas, tables, and lists. The bibliography system is now set up to work seamlessly with Bibby AI's citation management. Feel free to modify it for your needs!

\bibliographystyle{plain}
\bibliography{cite}

\end{document}\documentclass{article}
\usepackage[utf8]{inputenc}
\usepackage{booktabs}
\usepackage{amsmath}
\usepackage{natbib}

\title{My Document}
\author{Your Name}
\date{\today}

\begin{document}

\maketitle

\section{Introduction}
This is a sample document demonstrating various \LaTeX{} features including mathematical equations, tables, and formatting. You can cite papers using the citation search feature, and they will be automatically added to your \texttt{cite.bib} file.

\section{Mathematical Expressions}
Here are some example equations:

\subsection{Quadratic Formula}
The quadratic formula is given by:
\[x = \frac{-b \pm \sqrt{b^2 - 4ac}}{2a}\]

\subsection{Matrix Example}
A sample matrix:
\[
\begin{pmatrix}
a_{11} & a_{12} & a_{13} \\
a_{21} & a_{22} & a_{23} \\
a_{31} & a_{32} & a_{33}
\end{pmatrix}
\]

\section{Tables}
Below is a sample table showing data:

\begin{table}[h]
\centering
\begin{tabular}{@{}lrrr@{}}
\toprule
\textbf{Category} & \textbf{Value 1} & \textbf{Value 2} & \textbf{Total} \\
\midrule
Item A & 10.5 & 20.3 & 30.8 \\
Item B & 15.2 & 25.1 & 40.3 \\
Item C & 22.1 & 12.7 & 34.8 \\
\bottomrule
\end{tabular}
\caption{Sample Data Table}
\label{tab:sample}
\end{table}

\section{Lists}
\subsection{Bullet Points}
\begin{itemize}
  \item First item
  \item Second item
  \item Third item with sub-items:
    \begin{itemize}
      \item Sub-item 1
      \item Sub-item 2
    \end{itemize}
\end{itemize}

\subsection{Numbered List}
\begin{enumerate}
  \item First step
  \item Second step
  \item Third step
\end{enumerate}

\section{Citations and References}
When you use the citation search feature, references will be automatically managed. For example, you might cite a paper like this: \cite{example2024}. All citations will be automatically added to your \texttt{cite.bib} file.

\section{Conclusion}
This template demonstrates basic \LaTeX{} features including mathematical formulas, tables, and lists. The bibliography system is now set up to work seamlessly with Bibby AI's citation management. Feel free to modify it for your needs!

\bibliographystyle{plain}
\bibliography{cite}

\end{document}\documentclass{article}
\usepackage[utf8]{inputenc}
\usepackage{booktabs}
\usepackage{amsmath}
\usepackage{natbib}

\title{My Document}
\author{Your Name}
\date{\today}

\begin{document}

\maketitle

\section{Introduction}
This is a sample document demonstrating various \LaTeX{} features including mathematical equations, tables, and formatting. You can cite papers using the citation search feature, and they will be automatically added to your \texttt{cite.bib} file.

\section{Mathematical Expressions}
Here are some example equations:

\subsection{Quadratic Formula}
The quadratic formula is given by:
\[x = \frac{-b \pm \sqrt{b^2 - 4ac}}{2a}\]

\subsection{Matrix Example}
A sample matrix:
\[
\begin{pmatrix}
a_{11} & a_{12} & a_{13} \\
a_{21} & a_{22} & a_{23} \\
a_{31} & a_{32} & a_{33}
\end{pmatrix}
\]

\section{Tables}
Below is a sample table showing data:

\begin{table}[h]
\centering
\begin{tabular}{@{}lrrr@{}}
\toprule
\textbf{Category} & \textbf{Value 1} & \textbf{Value 2} & \textbf{Total} \\
\midrule
Item A & 10.5 & 20.3 & 30.8 \\
Item B & 15.2 & 25.1 & 40.3 \\
Item C & 22.1 & 12.7 & 34.8 \\
\bottomrule
\end{tabular}
\caption{Sample Data Table}
\label{tab:sample}
\end{table}

\section{Lists}
\subsection{Bullet Points}
\begin{itemize}
  \item First item
  \item Second item
  \item Third item with sub-items:
    \begin{itemize}
      \item Sub-item 1
      \item Sub-item 2
    \end{itemize}
\end{itemize}

\subsection{Numbered List}
\begin{enumerate}
  \item First step
  \item Second step
  \item Third step
\end{enumerate}

\section{Citations and References}
When you use the citation search feature, references will be automatically managed. For example, you might cite a paper like this: \cite{example2024}. All citations will be automatically added to your \texttt{cite.bib} file.

\section{Conclusion}
This template demonstrates basic \LaTeX{} features including mathematical formulas, tables, and lists. The bibliography system is now set up to work seamlessly with Bibby AI's citation management. Feel free to modify it for your needs!

\bibliographystyle{plain}
\bibliography{cite}

\end{document}\documentclass{article}
\usepackage[utf8]{inputenc}
\usepackage{booktabs}
\usepackage{amsmath}
\usepackage{natbib}

\title{My Document}
\author{Your Name}
\date{\today}

\begin{document}

\maketitle

\section{Introduction}
This is a sample document demonstrating various \LaTeX{} features including mathematical equations, tables, and formatting. You can cite papers using the citation search feature, and they will be automatically added to your \texttt{cite.bib} file.

\section{Mathematical Expressions}
Here are some example equations:

\subsection{Quadratic Formula}
The quadratic formula is given by:
\[x = \frac{-b \pm \sqrt{b^2 - 4ac}}{2a}\]

\subsection{Matrix Example}
A sample matrix:
\[
\begin{pmatrix}
a_{11} & a_{12} & a_{13} \\
a_{21} & a_{22} & a_{23} \\
a_{31} & a_{32} & a_{33}
\end{pmatrix}
\]

\section{Tables}
Below is a sample table showing data:

\begin{table}[h]
\centering
\begin{tabular}{@{}lrrr@{}}
\toprule
\textbf{Category} & \textbf{Value 1} & \textbf{Value 2} & \textbf{Total} \\
\midrule
Item A & 10.5 & 20.3 & 30.8 \\
Item B & 15.2 & 25.1 & 40.3 \\
Item C & 22.1 & 12.7 & 34.8 \\
\bottomrule
\end{tabular}
\caption{Sample Data Table}
\label{tab:sample}
\end{table}

\section{Lists}
\subsection{Bullet Points}
\begin{itemize}
  \item First item
  \item Second item
  \item Third item with sub-items:
    \begin{itemize}
      \item Sub-item 1
      \item Sub-item 2
    \end{itemize}
\end{itemize}

\subsection{Numbered List}
\begin{enumerate}
  \item First step
  \item Second step
  \item Third step
\end{enumerate}

\section{Citations and References}
When you use the citation search feature, references will be automatically managed. For example, you might cite a paper like this: \cite{example2024}. All citations will be automatically added to your \texttt{cite.bib} file.

\section{Conclusion}
This template demonstrates basic \LaTeX{} features including mathematical formulas, tables, and lists. The bibliography system is now set up to work seamlessly with Bibby AI's citation management. Feel free to modify it for your needs!

\bibliographystyle{plain}
\bibliography{cite}

\end{document}\documentclass{article}
\usepackage[utf8]{inputenc}
\usepackage{booktabs}
\usepackage{amsmath}
\usepackage{natbib}

\title{My Document}
\author{Your Name}
\date{\today}

\begin{document}

\maketitle

\section{Introduction}
This is a sample document demonstrating various \LaTeX{} features including mathematical equations, tables, and formatting. You can cite papers using the citation search feature, and they will be automatically added to your \texttt{cite.bib} file.

\section{Mathematical Expressions}
Here are some example equations:

\subsection{Quadratic Formula}
The quadratic formula is given by:
\[x = \frac{-b \pm \sqrt{b^2 - 4ac}}{2a}\]

\subsection{Matrix Example}
A sample matrix:
\[
\begin{pmatrix}
a_{11} & a_{12} & a_{13} \\
a_{21} & a_{22} & a_{23} \\
a_{31} & a_{32} & a_{33}
\end{pmatrix}
\]

\section{Tables}
Below is a sample table showing data:

\begin{table}[h]
\centering
\begin{tabular}{@{}lrrr@{}}
\toprule
\textbf{Category} & \textbf{Value 1} & \textbf{Value 2} & \textbf{Total} \\
\midrule
Item A & 10.5 & 20.3 & 30.8 \\
Item B & 15.2 & 25.1 & 40.3 \\
Item C & 22.1 & 12.7 & 34.8 \\
\bottomrule
\end{tabular}
\caption{Sample Data Table}
\label{tab:sample}
\end{table}

\section{Lists}
\subsection{Bullet Points}
\begin{itemize}
  \item First item
  \item Second item
  \item Third item with sub-items:
    \begin{itemize}
      \item Sub-item 1
      \item Sub-item 2
    \end{itemize}
\end{itemize}

\subsection{Numbered List}
\begin{enumerate}
  \item First step
  \item Second step
  \item Third step
\end{enumerate}

\section{Citations and References}
When you use the citation search feature, references will be automatically managed. For example, you might cite a paper like this: \cite{example2024}. All citations will be automatically added to your \texttt{cite.bib} file.

\section{Conclusion}
This template demonstrates basic \LaTeX{} features including mathematical formulas, tables, and lists. The bibliography system is now set up to work seamlessly with Bibby AI's citation management. Feel free to modify it for your needs!

\bibliographystyle{plain}
\bibliography{cite}

\end{document}\documentclass{article}
\usepackage[utf8]{inputenc}
\usepackage{booktabs}
\usepackage{amsmath}
\usepackage{natbib}

\title{My Document}
\author{Your Name}
\date{\today}

\begin{document}

\maketitle

\section{Introduction}
This is a sample document demonstrating various \LaTeX{} features including mathematical equations, tables, and formatting. You can cite papers using the citation search feature, and they will be automatically added to your \texttt{cite.bib} file.

\section{Mathematical Expressions}
Here are some example equations:

\subsection{Quadratic Formula}
The quadratic formula is given by:
\[x = \frac{-b \pm \sqrt{b^2 - 4ac}}{2a}\]

\subsection{Matrix Example}
A sample matrix:
\[
\begin{pmatrix}
a_{11} & a_{12} & a_{13} \\
a_{21} & a_{22} & a_{23} \\
a_{31} & a_{32} & a_{33}
\end{pmatrix}
\]

\section{Tables}
Below is a sample table showing data:

\begin{table}[h]
\centering
\begin{tabular}{@{}lrrr@{}}
\toprule
\textbf{Category} & \textbf{Value 1} & \textbf{Value 2} & \textbf{Total} \\
\midrule
Item A & 10.5 & 20.3 & 30.8 \\
Item B & 15.2 & 25.1 & 40.3 \\
Item C & 22.1 & 12.7 & 34.8 \\
\bottomrule
\end{tabular}
\caption{Sample Data Table}
\label{tab:sample}
\end{table}

\section{Lists}
\subsection{Bullet Points}
\begin{itemize}
  \item First item
  \item Second item
  \item Third item with sub-items:
    \begin{itemize}
      \item Sub-item 1
      \item Sub-item 2
    \end{itemize}
\end{itemize}

\subsection{Numbered List}
\begin{enumerate}
  \item First step
  \item Second step
  \item Third step
\end{enumerate}

\section{Citations and References}
When you use the citation search feature, references will be automatically managed. For example, you might cite a paper like this: \cite{example2024}. All citations will be automatically added to your \texttt{cite.bib} file.

\section{Conclusion}
This template demonstrates basic \LaTeX{} features including mathematical formulas, tables, and lists. The bibliography system is now set up to work seamlessly with Bibby AI's citation management. Feel free to modify it for your needs!

\bibliographystyle{plain}
\bibliography{cite}

\end{document}\documentclass{article}
\usepackage[utf8]{inputenc}
\usepackage{booktabs}
\usepackage{amsmath}
\usepackage{natbib}

\title{My Document}
\author{Your Name}
\date{\today}

\begin{document}

\maketitle

\section{Introduction}
This is a sample document demonstrating various \LaTeX{} features including mathematical equations, tables, and formatting. You can cite papers using the citation search feature, and they will be automatically added to your \texttt{cite.bib} file.

\section{Mathematical Expressions}
Here are some example equations:

\subsection{Quadratic Formula}
The quadratic formula is given by:
\[x = \frac{-b \pm \sqrt{b^2 - 4ac}}{2a}\]

\subsection{Matrix Example}
A sample matrix:
\[
\begin{pmatrix}
a_{11} & a_{12} & a_{13} \\
a_{21} & a_{22} & a_{23} \\
a_{31} & a_{32} & a_{33}
\end{pmatrix}
\]

\section{Tables}
Below is a sample table showing data:

\begin{table}[h]
\centering
\begin{tabular}{@{}lrrr@{}}
\toprule
\textbf{Category} & \textbf{Value 1} & \textbf{Value 2} & \textbf{Total} \\
\midrule
Item A & 10.5 & 20.3 & 30.8 \\
Item B & 15.2 & 25.1 & 40.3 \\
Item C & 22.1 & 12.7 & 34.8 \\
\bottomrule
\end{tabular}
\caption{Sample Data Table}
\label{tab:sample}
\end{table}

\section{Lists}
\subsection{Bullet Points}
\begin{itemize}
  \item First item
  \item Second item
  \item Third item with sub-items:
    \begin{itemize}
      \item Sub-item 1
      \item Sub-item 2
    \end{itemize}
\end{itemize}

\subsection{Numbered List}
\begin{enumerate}
  \item First step
  \item Second step
  \item Third step
\end{enumerate}

\section{Citations and References}
When you use the citation search feature, references will be automatically managed. For example, you might cite a paper like this: \cite{example2024}. All citations will be automatically added to your \texttt{cite.bib} file.

\section{Conclusion}
This template demonstrates basic \LaTeX{} features including mathematical formulas, tables, and lists. The bibliography system is now set up to work seamlessly with Bibby AI's citation management. Feel free to modify it for your needs!

\bibliographystyle{plain}
\bibliography{cite}

\end{document}\documentclass{article}
\usepackage[utf8]{inputenc}
\usepackage{booktabs}
\usepackage{amsmath}
\usepackage{natbib}

\title{My Document}
\author{Your Name}
\date{\today}

\begin{document}

\maketitle

\section{Introduction}
This is a sample document demonstrating various \LaTeX{} features including mathematical equations, tables, and formatting. You can cite papers using the citation search feature, and they will be automatically added to your \texttt{cite.bib} file.

\section{Mathematical Expressions}
Here are some example equations:

\subsection{Quadratic Formula}
The quadratic formula is given by:
\[x = \frac{-b \pm \sqrt{b^2 - 4ac}}{2a}\]

\subsection{Matrix Example}
A sample matrix:
\[
\begin{pmatrix}
a_{11} & a_{12} & a_{13} \\
a_{21} & a_{22} & a_{23} \\
a_{31} & a_{32} & a_{33}
\end{pmatrix}
\]

\section{Tables}
Below is a sample table showing data:

\begin{table}[h]
\centering
\begin{tabular}{@{}lrrr@{}}
\toprule
\textbf{Category} & \textbf{Value 1} & \textbf{Value 2} & \textbf{Total} \\
\midrule
Item A & 10.5 & 20.3 & 30.8 \\
Item B & 15.2 & 25.1 & 40.3 \\
Item C & 22.1 & 12.7 & 34.8 \\
\bottomrule
\end{tabular}
\caption{Sample Data Table}
\label{tab:sample}
\end{table}

\section{Lists}
\subsection{Bullet Points}
\begin{itemize}
  \item First item
  \item Second item
  \item Third item with sub-items:
    \begin{itemize}
      \item Sub-item 1
      \item Sub-item 2
    \end{itemize}
\end{itemize}

\subsection{Numbered List}
\begin{enumerate}
  \item First step
  \item Second step
  \item Third step
\end{enumerate}

\section{Citations and References}
When you use the citation search feature, references will be automatically managed. For example, you might cite a paper like this: \cite{example2024}. All citations will be automatically added to your \texttt{cite.bib} file.

\section{Conclusion}
This template demonstrates basic \LaTeX{} features including mathematical formulas, tables, and lists. The bibliography system is now set up to work seamlessly with Bibby AI's citation management. Feel free to modify it for your needs!

\bibliographystyle{plain}
\bibliography{cite}

\end{document}\documentclass{article}
\usepackage[utf8]{inputenc}
\usepackage{booktabs}
\usepackage{amsmath}
\usepackage{natbib}

\title{My Document}
\author{Your Name}
\date{\today}

\begin{document}

\maketitle

\section{Introduction}
This is a sample document demonstrating various \LaTeX{} features including mathematical equations, tables, and formatting. You can cite papers using the citation search feature, and they will be automatically added to your \texttt{cite.bib} file.

\section{Mathematical Expressions}
Here are some example equations:

\subsection{Quadratic Formula}
The quadratic formula is given by:
\[x = \frac{-b \pm \sqrt{b^2 - 4ac}}{2a}\]

\subsection{Matrix Example}
A sample matrix:
\[
\begin{pmatrix}
a_{11} & a_{12} & a_{13} \\
a_{21} & a_{22} & a_{23} \\
a_{31} & a_{32} & a_{33}
\end{pmatrix}
\]

\section{Tables}
Below is a sample table showing data:

\begin{table}[h]
\centering
\begin{tabular}{@{}lrrr@{}}
\toprule
\textbf{Category} & \textbf{Value 1} & \textbf{Value 2} & \textbf{Total} \\
\midrule
Item A & 10.5 & 20.3 & 30.8 \\
Item B & 15.2 & 25.1 & 40.3 \\
Item C & 22.1 & 12.7 & 34.8 \\
\bottomrule
\end{tabular}
\caption{Sample Data Table}
\label{tab:sample}
\end{table}

\section{Lists}
\subsection{Bullet Points}
\begin{itemize}
  \item First item
  \item Second item
  \item Third item with sub-items:
    \begin{itemize}
      \item Sub-item 1
      \item Sub-item 2
    \end{itemize}
\end{itemize}

\subsection{Numbered List}
\begin{enumerate}
  \item First step
  \item Second step
  \item Third step
\end{enumerate}

\section{Citations and References}
When you use the citation search feature, references will be automatically managed. For example, you might cite a paper like this: \cite{example2024}. All citations will be automatically added to your \texttt{cite.bib} file.

\section{Conclusion}
This template demonstrates basic \LaTeX{} features including mathematical formulas, tables, and lists. The bibliography system is now set up to work seamlessly with Bibby AI's citation management. Feel free to modify it for your needs!

\bibliographystyle{plain}
\bibliography{cite}

\end{document}\documentclass{article}
\usepackage[utf8]{inputenc}
\usepackage{booktabs}
\usepackage{amsmath}
\usepackage{natbib}

\title{My Document}
\author{Your Name}
\date{\today}

\begin{document}

\maketitle

\section{Introduction}
This is a sample document demonstrating various \LaTeX{} features including mathematical equations, tables, and formatting. You can cite papers using the citation search feature, and they will be automatically added to your \texttt{cite.bib} file.

\section{Mathematical Expressions}
Here are some example equations:

\subsection{Quadratic Formula}
The quadratic formula is given by:
\[x = \frac{-b \pm \sqrt{b^2 - 4ac}}{2a}\]

\subsection{Matrix Example}
A sample matrix:
\[
\begin{pmatrix}
a_{11} & a_{12} & a_{13} \\
a_{21} & a_{22} & a_{23} \\
a_{31} & a_{32} & a_{33}
\end{pmatrix}
\]

\section{Tables}
Below is a sample table showing data:

\begin{table}[h]
\centering
\begin{tabular}{@{}lrrr@{}}
\toprule
\textbf{Category} & \textbf{Value 1} & \textbf{Value 2} & \textbf{Total} \\
\midrule
Item A & 10.5 & 20.3 & 30.8 \\
Item B & 15.2 & 25.1 & 40.3 \\
Item C & 22.1 & 12.7 & 34.8 \\
\bottomrule
\end{tabular}
\caption{Sample Data Table}
\label{tab:sample}
\end{table}

\section{Lists}
\subsection{Bullet Points}
\begin{itemize}
  \item First item
  \item Second item
  \item Third item with sub-items:
    \begin{itemize}
      \item Sub-item 1
      \item Sub-item 2
    \end{itemize}
\end{itemize}

\subsection{Numbered List}
\begin{enumerate}
  \item First step
  \item Second step
  \item Third step
\end{enumerate}

\section{Citations and References}
When you use the citation search feature, references will be automatically managed. For example, you might cite a paper like this: \cite{example2024}. All citations will be automatically added to your \texttt{cite.bib} file.

\section{Conclusion}
This template demonstrates basic \LaTeX{} features including mathematical formulas, tables, and lists. The bibliography system is now set up to work seamlessly with Bibby AI's citation management. Feel free to modify it for your needs!

\bibliographystyle{plain}
\bibliography{cite}

\end{document}\documentclass{article}
\usepackage[utf8]{inputenc}
\usepackage{booktabs}
\usepackage{amsmath}
\usepackage{natbib}

\title{My Document}
\author{Your Name}
\date{\today}

\begin{document}

\maketitle

\section{Introduction}
This is a sample document demonstrating various \LaTeX{} features including mathematical equations, tables, and formatting. You can cite papers using the citation search feature, and they will be automatically added to your \texttt{cite.bib} file.

\section{Mathematical Expressions}
Here are some example equations:

\subsection{Quadratic Formula}
The quadratic formula is given by:
\[x = \frac{-b \pm \sqrt{b^2 - 4ac}}{2a}\]

\subsection{Matrix Example}
A sample matrix:
\[
\begin{pmatrix}
a_{11} & a_{12} & a_{13} \\
a_{21} & a_{22} & a_{23} \\
a_{31} & a_{32} & a_{33}
\end{pmatrix}
\]

\section{Tables}
Below is a sample table showing data:

\begin{table}[h]
\centering
\begin{tabular}{@{}lrrr@{}}
\toprule
\textbf{Category} & \textbf{Value 1} & \textbf{Value 2} & \textbf{Total} \\
\midrule
Item A & 10.5 & 20.3 & 30.8 \\
Item B & 15.2 & 25.1 & 40.3 \\
Item C & 22.1 & 12.7 & 34.8 \\
\bottomrule
\end{tabular}
\caption{Sample Data Table}
\label{tab:sample}
\end{table}

\section{Lists}
\subsection{Bullet Points}
\begin{itemize}
  \item First item
  \item Second item
  \item Third item with sub-items:
    \begin{itemize}
      \item Sub-item 1
      \item Sub-item 2
    \end{itemize}
\end{itemize}

\subsection{Numbered List}
\begin{enumerate}
  \item First step
  \item Second step
  \item Third step
\end{enumerate}

\section{Citations and References}
When you use the citation search feature, references will be automatically managed. For example, you might cite a paper like this: \cite{example2024}. All citations will be automatically added to your \texttt{cite.bib} file.

\section{Conclusion}
This template demonstrates basic \LaTeX{} features including mathematical formulas, tables, and lists. The bibliography system is now set up to work seamlessly with Bibby AI's citation management. Feel free to modify it for your needs!

\bibliographystyle{plain}
\bibliography{cite}

\end{document}\documentclass{article}
\usepackage[utf8]{inputenc}
\usepackage{booktabs}
\usepackage{amsmath}
\usepackage{natbib}

\title{My Document}
\author{Your Name}
\date{\today}

\begin{document}

\maketitle

\section{Introduction}
This is a sample document demonstrating various \LaTeX{} features including mathematical equations, tables, and formatting. You can cite papers using the citation search feature, and they will be automatically added to your \texttt{cite.bib} file.

\section{Mathematical Expressions}
Here are some example equations:

\subsection{Quadratic Formula}
The quadratic formula is given by:
\[x = \frac{-b \pm \sqrt{b^2 - 4ac}}{2a}\]

\subsection{Matrix Example}
A sample matrix:
\[
\begin{pmatrix}
a_{11} & a_{12} & a_{13} \\
a_{21} & a_{22} & a_{23} \\
a_{31} & a_{32} & a_{33}
\end{pmatrix}
\]

\section{Tables}
Below is a sample table showing data:

\begin{table}[h]
\centering
\begin{tabular}{@{}lrrr@{}}
\toprule
\textbf{Category} & \textbf{Value 1} & \textbf{Value 2} & \textbf{Total} \\
\midrule
Item A & 10.5 & 20.3 & 30.8 \\
Item B & 15.2 & 25.1 & 40.3 \\
Item C & 22.1 & 12.7 & 34.8 \\
\bottomrule
\end{tabular}
\caption{Sample Data Table}
\label{tab:sample}
\end{table}

\section{Lists}
\subsection{Bullet Points}
\begin{itemize}
  \item First item
  \item Second item
  \item Third item with sub-items:
    \begin{itemize}
      \item Sub-item 1
      \item Sub-item 2
    \end{itemize}
\end{itemize}

\subsection{Numbered List}
\begin{enumerate}
  \item First step
  \item Second step
  \item Third step
\end{enumerate}

\section{Citations and References}
When you use the citation search feature, references will be automatically managed. For example, you might cite a paper like this: \cite{example2024}. All citations will be automatically added to your \texttt{cite.bib} file.

\section{Conclusion}
This template demonstrates basic \LaTeX{} features including mathematical formulas, tables, and lists. The bibliography system is now set up to work seamlessly with Bibby AI's citation management. Feel free to modify it for your needs!

\bibliographystyle{plain}
\bibliography{cite}

\end{document}\documentclass{article}
\usepackage[utf8]{inputenc}
\usepackage{booktabs}
\usepackage{amsmath}
\usepackage{natbib}

\title{My Document}
\author{Your Name}
\date{\today}

\begin{document}

\maketitle

\section{Introduction}
This is a sample document demonstrating various \LaTeX{} features including mathematical equations, tables, and formatting. You can cite papers using the citation search feature, and they will be automatically added to your \texttt{cite.bib} file.

\section{Mathematical Expressions}
Here are some example equations:

\subsection{Quadratic Formula}
The quadratic formula is given by:
\[x = \frac{-b \pm \sqrt{b^2 - 4ac}}{2a}\]

\subsection{Matrix Example}
A sample matrix:
\[
\begin{pmatrix}
a_{11} & a_{12} & a_{13} \\
a_{21} & a_{22} & a_{23} \\
a_{31} & a_{32} & a_{33}
\end{pmatrix}
\]

\section{Tables}
Below is a sample table showing data:

\begin{table}[h]
\centering
\begin{tabular}{@{}lrrr@{}}
\toprule
\textbf{Category} & \textbf{Value 1} & \textbf{Value 2} & \textbf{Total} \\
\midrule
Item A & 10.5 & 20.3 & 30.8 \\
Item B & 15.2 & 25.1 & 40.3 \\
Item C & 22.1 & 12.7 & 34.8 \\
\bottomrule
\end{tabular}
\caption{Sample Data Table}
\label{tab:sample}
\end{table}

\section{Lists}
\subsection{Bullet Points}
\begin{itemize}
  \item First item
  \item Second item
  \item Third item with sub-items:
    \begin{itemize}
      \item Sub-item 1
      \item Sub-item 2
    \end{itemize}
\end{itemize}

\subsection{Numbered List}
\begin{enumerate}
  \item First step
  \item Second step
  \item Third step
\end{enumerate}

\section{Citations and References}
When you use the citation search feature, references will be automatically managed. For example, you might cite a paper like this: \cite{example2024}. All citations will be automatically added to your \texttt{cite.bib} file.

\section{Conclusion}
This template demonstrates basic \LaTeX{} features including mathematical formulas, tables, and lists. The bibliography system is now set up to work seamlessly with Bibby AI's citation management. Feel free to modify it for your needs!

\bibliographystyle{plain}
\bibliography{cite}

\end{document}\documentclass{article}
\usepackage[utf8]{inputenc}
\usepackage{booktabs}
\usepackage{amsmath}
\usepackage{natbib}

\title{My Document}
\author{Your Name}
\date{\today}

\begin{document}

\maketitle

\section{Introduction}
This is a sample document demonstrating various \LaTeX{} features including mathematical equations, tables, and formatting. You can cite papers using the citation search feature, and they will be automatically added to your \texttt{cite.bib} file.

\section{Mathematical Expressions}
Here are some example equations:

\subsection{Quadratic Formula}
The quadratic formula is given by:
\[x = \frac{-b \pm \sqrt{b^2 - 4ac}}{2a}\]

\subsection{Matrix Example}
A sample matrix:
\[
\begin{pmatrix}
a_{11} & a_{12} & a_{13} \\
a_{21} & a_{22} & a_{23} \\
a_{31} & a_{32} & a_{33}
\end{pmatrix}
\]

\section{Tables}
Below is a sample table showing data:

\begin{table}[h]
\centering
\begin{tabular}{@{}lrrr@{}}
\toprule
\textbf{Category} & \textbf{Value 1} & \textbf{Value 2} & \textbf{Total} \\
\midrule
Item A & 10.5 & 20.3 & 30.8 \\
Item B & 15.2 & 25.1 & 40.3 \\
Item C & 22.1 & 12.7 & 34.8 \\
\bottomrule
\end{tabular}
\caption{Sample Data Table}
\label{tab:sample}
\end{table}

\section{Lists}
\subsection{Bullet Points}
\begin{itemize}
  \item First item
  \item Second item
  \item Third item with sub-items:
    \begin{itemize}
      \item Sub-item 1
      \item Sub-item 2
    \end{itemize}
\end{itemize}

\subsection{Numbered List}
\begin{enumerate}
  \item First step
  \item Second step
  \item Third step
\end{enumerate}

\section{Citations and References}
When you use the citation search feature, references will be automatically managed. For example, you might cite a paper like this: \cite{example2024}. All citations will be automatically added to your \texttt{cite.bib} file.

\section{Conclusion}
This template demonstrates basic \LaTeX{} features including mathematical formulas, tables, and lists. The bibliography system is now set up to work seamlessly with Bibby AI's citation management. Feel free to modify it for your needs!

\bibliographystyle{plain}
\bibliography{cite}

\end{document}\documentclass{article}
\usepackage[utf8]{inputenc}
\usepackage{booktabs}
\usepackage{amsmath}
\usepackage{natbib}

\title{My Document}
\author{Your Name}
\date{\today}

\begin{document}

\maketitle

\section{Introduction}
This is a sample document demonstrating various \LaTeX{} features including mathematical equations, tables, and formatting. You can cite papers using the citation search feature, and they will be automatically added to your \texttt{cite.bib} file.

\section{Mathematical Expressions}
Here are some example equations:

\subsection{Quadratic Formula}
The quadratic formula is given by:
\[x = \frac{-b \pm \sqrt{b^2 - 4ac}}{2a}\]

\subsection{Matrix Example}
A sample matrix:
\[
\begin{pmatrix}
a_{11} & a_{12} & a_{13} \\
a_{21} & a_{22} & a_{23} \\
a_{31} & a_{32} & a_{33}
\end{pmatrix}
\]

\section{Tables}
Below is a sample table showing data:

\begin{table}[h]
\centering
\begin{tabular}{@{}lrrr@{}}
\toprule
\textbf{Category} & \textbf{Value 1} & \textbf{Value 2} & \textbf{Total} \\
\midrule
Item A & 10.5 & 20.3 & 30.8 \\
Item B & 15.2 & 25.1 & 40.3 \\
Item C & 22.1 & 12.7 & 34.8 \\
\bottomrule
\end{tabular}
\caption{Sample Data Table}
\label{tab:sample}
\end{table}

\section{Lists}
\subsection{Bullet Points}
\begin{itemize}
  \item First item
  \item Second item
  \item Third item with sub-items:
    \begin{itemize}
      \item Sub-item 1
      \item Sub-item 2
    \end{itemize}
\end{itemize}

\subsection{Numbered List}
\begin{enumerate}
  \item First step
  \item Second step
  \item Third step
\end{enumerate}

\section{Citations and References}
When you use the citation search feature, references will be automatically managed. For example, you might cite a paper like this: \cite{example2024}. All citations will be automatically added to your \texttt{cite.bib} file.

\section{Conclusion}
This template demonstrates basic \LaTeX{} features including mathematical formulas, tables, and lists. The bibliography system is now set up to work seamlessly with Bibby AI's citation management. Feel free to modify it for your needs!

\bibliographystyle{plain}
\bibliography{cite}

\end{document}\documentclass{article}
\usepackage[utf8]{inputenc}
\usepackage{booktabs}
\usepackage{amsmath}
\usepackage{natbib}

\title{My Document}
\author{Your Name}
\date{\today}

\begin{document}

\maketitle

\section{Introduction}
This is a sample document demonstrating various \LaTeX{} features including mathematical equations, tables, and formatting. You can cite papers using the citation search feature, and they will be automatically added to your \texttt{cite.bib} file.

\section{Mathematical Expressions}
Here are some example equations:

\subsection{Quadratic Formula}
The quadratic formula is given by:
\[x = \frac{-b \pm \sqrt{b^2 - 4ac}}{2a}\]

\subsection{Matrix Example}
A sample matrix:
\[
\begin{pmatrix}
a_{11} & a_{12} & a_{13} \\
a_{21} & a_{22} & a_{23} \\
a_{31} & a_{32} & a_{33}
\end{pmatrix}
\]

\section{Tables}
Below is a sample table showing data:

\begin{table}[h]
\centering
\begin{tabular}{@{}lrrr@{}}
\toprule
\textbf{Category} & \textbf{Value 1} & \textbf{Value 2} & \textbf{Total} \\
\midrule
Item A & 10.5 & 20.3 & 30.8 \\
Item B & 15.2 & 25.1 & 40.3 \\
Item C & 22.1 & 12.7 & 34.8 \\
\bottomrule
\end{tabular}
\caption{Sample Data Table}
\label{tab:sample}
\end{table}

\section{Lists}
\subsection{Bullet Points}
\begin{itemize}
  \item First item
  \item Second item
  \item Third item with sub-items:
    \begin{itemize}
      \item Sub-item 1
      \item Sub-item 2
    \end{itemize}
\end{itemize}

\subsection{Numbered List}
\begin{enumerate}
  \item First step
  \item Second step
  \item Third step
\end{enumerate}

\section{Citations and References}
When you use the citation search feature, references will be automatically managed. For example, you might cite a paper like this: \cite{example2024}. All citations will be automatically added to your \texttt{cite.bib} file.

\section{Conclusion}
This template demonstrates basic \LaTeX{} features including mathematical formulas, tables, and lists. The bibliography system is now set up to work seamlessly with Bibby AI's citation management. Feel free to modify it for your needs!

\bibliographystyle{plain}
\bibliography{cite}

\end{document}\documentclass{article}
\usepackage[utf8]{inputenc}
\usepackage{booktabs}
\usepackage{amsmath}
\usepackage{natbib}

\title{My Document}
\author{Your Name}
\date{\today}

\begin{document}

\maketitle

\section{Introduction}
This is a sample document demonstrating various \LaTeX{} features including mathematical equations, tables, and formatting. You can cite papers using the citation search feature, and they will be automatically added to your \texttt{cite.bib} file.

\section{Mathematical Expressions}
Here are some example equations:

\subsection{Quadratic Formula}
The quadratic formula is given by:
\[x = \frac{-b \pm \sqrt{b^2 - 4ac}}{2a}\]

\subsection{Matrix Example}
A sample matrix:
\[
\begin{pmatrix}
a_{11} & a_{12} & a_{13} \\
a_{21} & a_{22} & a_{23} \\
a_{31} & a_{32} & a_{33}
\end{pmatrix}
\]

\section{Tables}
Below is a sample table showing data:

\begin{table}[h]
\centering
\begin{tabular}{@{}lrrr@{}}
\toprule
\textbf{Category} & \textbf{Value 1} & \textbf{Value 2} & \textbf{Total} \\
\midrule
Item A & 10.5 & 20.3 & 30.8 \\
Item B & 15.2 & 25.1 & 40.3 \\
Item C & 22.1 & 12.7 & 34.8 \\
\bottomrule
\end{tabular}
\caption{Sample Data Table}
\label{tab:sample}
\end{table}

\section{Lists}
\subsection{Bullet Points}
\begin{itemize}
  \item First item
  \item Second item
  \item Third item with sub-items:
    \begin{itemize}
      \item Sub-item 1
      \item Sub-item 2
    \end{itemize}
\end{itemize}

\subsection{Numbered List}
\begin{enumerate}
  \item First step
  \item Second step
  \item Third step
\end{enumerate}

\section{Citations and References}
When you use the citation search feature, references will be automatically managed. For example, you might cite a paper like this: \cite{example2024}. All citations will be automatically added to your \texttt{cite.bib} file.

\section{Conclusion}
This template demonstrates basic \LaTeX{} features including mathematical formulas, tables, and lists. The bibliography system is now set up to work seamlessly with Bibby AI's citation management. Feel free to modify it for your needs!

\bibliographystyle{plain}
\bibliography{cite}

\end{document}\documentclass{article}
\usepackage[utf8]{inputenc}
\usepackage{booktabs}
\usepackage{amsmath}
\usepackage{natbib}

\title{My Document}
\author{Your Name}
\date{\today}

\begin{document}

\maketitle

\section{Introduction}
This is a sample document demonstrating various \LaTeX{} features including mathematical equations, tables, and formatting. You can cite papers using the citation search feature, and they will be automatically added to your \texttt{cite.bib} file.

\section{Mathematical Expressions}
Here are some example equations:

\subsection{Quadratic Formula}
The quadratic formula is given by:
\[x = \frac{-b \pm \sqrt{b^2 - 4ac}}{2a}\]

\subsection{Matrix Example}
A sample matrix:
\[
\begin{pmatrix}
a_{11} & a_{12} & a_{13} \\
a_{21} & a_{22} & a_{23} \\
a_{31} & a_{32} & a_{33}
\end{pmatrix}
\]

\section{Tables}
Below is a sample table showing data:

\begin{table}[h]
\centering
\begin{tabular}{@{}lrrr@{}}
\toprule
\textbf{Category} & \textbf{Value 1} & \textbf{Value 2} & \textbf{Total} \\
\midrule
Item A & 10.5 & 20.3 & 30.8 \\
Item B & 15.2 & 25.1 & 40.3 \\
Item C & 22.1 & 12.7 & 34.8 \\
\bottomrule
\end{tabular}
\caption{Sample Data Table}
\label{tab:sample}
\end{table}

\section{Lists}
\subsection{Bullet Points}
\begin{itemize}
  \item First item
  \item Second item
  \item Third item with sub-items:
    \begin{itemize}
      \item Sub-item 1
      \item Sub-item 2
    \end{itemize}
\end{itemize}

\subsection{Numbered List}
\begin{enumerate}
  \item First step
  \item Second step
  \item Third step
\end{enumerate}

\section{Citations and References}
When you use the citation search feature, references will be automatically managed. For example, you might cite a paper like this: \cite{example2024}. All citations will be automatically added to your \texttt{cite.bib} file.

\section{Conclusion}
This template demonstrates basic \LaTeX{} features including mathematical formulas, tables, and lists. The bibliography system is now set up to work seamlessly with Bibby AI's citation management. Feel free to modify it for your needs!

\bibliographystyle{plain}
\bibliography{cite}

\end{document}\documentclass{article}
\usepackage[utf8]{inputenc}
\usepackage{booktabs}
\usepackage{amsmath}
\usepackage{natbib}

\title{My Document}
\author{Your Name}
\date{\today}

\begin{document}

\maketitle

\section{Introduction}
This is a sample document demonstrating various \LaTeX{} features including mathematical equations, tables, and formatting. You can cite papers using the citation search feature, and they will be automatically added to your \texttt{cite.bib} file.

\section{Mathematical Expressions}
Here are some example equations:

\subsection{Quadratic Formula}
The quadratic formula is given by:
\[x = \frac{-b \pm \sqrt{b^2 - 4ac}}{2a}\]

\subsection{Matrix Example}
A sample matrix:
\[
\begin{pmatrix}
a_{11} & a_{12} & a_{13} \\
a_{21} & a_{22} & a_{23} \\
a_{31} & a_{32} & a_{33}
\end{pmatrix}
\]

\section{Tables}
Below is a sample table showing data:

\begin{table}[h]
\centering
\begin{tabular}{@{}lrrr@{}}
\toprule
\textbf{Category} & \textbf{Value 1} & \textbf{Value 2} & \textbf{Total} \\
\midrule
Item A & 10.5 & 20.3 & 30.8 \\
Item B & 15.2 & 25.1 & 40.3 \\
Item C & 22.1 & 12.7 & 34.8 \\
\bottomrule
\end{tabular}
\caption{Sample Data Table}
\label{tab:sample}
\end{table}

\section{Lists}
\subsection{Bullet Points}
\begin{itemize}
  \item First item
  \item Second item
  \item Third item with sub-items:
    \begin{itemize}
      \item Sub-item 1
      \item Sub-item 2
    \end{itemize}
\end{itemize}

\subsection{Numbered List}
\begin{enumerate}
  \item First step
  \item Second step
  \item Third step
\end{enumerate}

\section{Citations and References}
When you use the citation search feature, references will be automatically managed. For example, you might cite a paper like this: \cite{example2024}. All citations will be automatically added to your \texttt{cite.bib} file.

\section{Conclusion}
This template demonstrates basic \LaTeX{} features including mathematical formulas, tables, and lists. The bibliography system is now set up to work seamlessly with Bibby AI's citation management. Feel free to modify it for your needs!

\bibliographystyle{plain}
\bibliography{cite}

\end{document}\documentclass{article}
\usepackage[utf8]{inputenc}
\usepackage{booktabs}
\usepackage{amsmath}
\usepackage{natbib}

\title{My Document}
\author{Your Name}
\date{\today}

\begin{document}

\maketitle

\section{Introduction}
This is a sample document demonstrating various \LaTeX{} features including mathematical equations, tables, and formatting. You can cite papers using the citation search feature, and they will be automatically added to your \texttt{cite.bib} file.

\section{Mathematical Expressions}
Here are some example equations:

\subsection{Quadratic Formula}
The quadratic formula is given by:
\[x = \frac{-b \pm \sqrt{b^2 - 4ac}}{2a}\]

\subsection{Matrix Example}
A sample matrix:
\[
\begin{pmatrix}
a_{11} & a_{12} & a_{13} \\
a_{21} & a_{22} & a_{23} \\
a_{31} & a_{32} & a_{33}
\end{pmatrix}
\]

\section{Tables}
Below is a sample table showing data:

\begin{table}[h]
\centering
\begin{tabular}{@{}lrrr@{}}
\toprule
\textbf{Category} & \textbf{Value 1} & \textbf{Value 2} & \textbf{Total} \\
\midrule
Item A & 10.5 & 20.3 & 30.8 \\
Item B & 15.2 & 25.1 & 40.3 \\
Item C & 22.1 & 12.7 & 34.8 \\
\bottomrule
\end{tabular}
\caption{Sample Data Table}
\label{tab:sample}
\end{table}

\section{Lists}
\subsection{Bullet Points}
\begin{itemize}
  \item First item
  \item Second item
  \item Third item with sub-items:
    \begin{itemize}
      \item Sub-item 1
      \item Sub-item 2
    \end{itemize}
\end{itemize}

\subsection{Numbered List}
\begin{enumerate}
  \item First step
  \item Second step
  \item Third step
\end{enumerate}

\section{Citations and References}
When you use the citation search feature, references will be automatically managed. For example, you might cite a paper like this: \cite{example2024}. All citations will be automatically added to your \texttt{cite.bib} file.

\section{Conclusion}
This template demonstrates basic \LaTeX{} features including mathematical formulas, tables, and lists. The bibliography system is now set up to work seamlessly with Bibby AI's citation management. Feel free to modify it for your needs!

\bibliographystyle{plain}
\bibliography{cite}

\end{document}\documentclass{article}
\usepackage[utf8]{inputenc}
\usepackage{booktabs}
\usepackage{amsmath}
\usepackage{natbib}

\title{My Document}
\author{Your Name}
\date{\today}

\begin{document}

\maketitle

\section{Introduction}
This is a sample document demonstrating various \LaTeX{} features including mathematical equations, tables, and formatting. You can cite papers using the citation search feature, and they will be automatically added to your \texttt{cite.bib} file.

\section{Mathematical Expressions}
Here are some example equations:

\subsection{Quadratic Formula}
The quadratic formula is given by:
\[x = \frac{-b \pm \sqrt{b^2 - 4ac}}{2a}\]

\subsection{Matrix Example}
A sample matrix:
\[
\begin{pmatrix}
a_{11} & a_{12} & a_{13} \\
a_{21} & a_{22} & a_{23} \\
a_{31} & a_{32} & a_{33}
\end{pmatrix}
\]

\section{Tables}
Below is a sample table showing data:

\begin{table}[h]
\centering
\begin{tabular}{@{}lrrr@{}}
\toprule
\textbf{Category} & \textbf{Value 1} & \textbf{Value 2} & \textbf{Total} \\
\midrule
Item A & 10.5 & 20.3 & 30.8 \\
Item B & 15.2 & 25.1 & 40.3 \\
Item C & 22.1 & 12.7 & 34.8 \\
\bottomrule
\end{tabular}
\caption{Sample Data Table}
\label{tab:sample}
\end{table}

\section{Lists}
\subsection{Bullet Points}
\begin{itemize}
  \item First item
  \item Second item
  \item Third item with sub-items:
    \begin{itemize}
      \item Sub-item 1
      \item Sub-item 2
    \end{itemize}
\end{itemize}

\subsection{Numbered List}
\begin{enumerate}
  \item First step
  \item Second step
  \item Third step
\end{enumerate}

\section{Citations and References}
When you use the citation search feature, references will be automatically managed. For example, you might cite a paper like this: \cite{example2024}. All citations will be automatically added to your \texttt{cite.bib} file.

\section{Conclusion}
This template demonstrates basic \LaTeX{} features including mathematical formulas, tables, and lists. The bibliography system is now set up to work seamlessly with Bibby AI's citation management. Feel free to modify it for your needs!

\bibliographystyle{plain}
\bibliography{cite}

\end{document}\documentclass{article}
\usepackage[utf8]{inputenc}
\usepackage{booktabs}
\usepackage{amsmath}
\usepackage{natbib}

\title{My Document}
\author{Your Name}
\date{\today}

\begin{document}

\maketitle

\section{Introduction}
This is a sample document demonstrating various \LaTeX{} features including mathematical equations, tables, and formatting. You can cite papers using the citation search feature, and they will be automatically added to your \texttt{cite.bib} file.

\section{Mathematical Expressions}
Here are some example equations:

\subsection{Quadratic Formula}
The quadratic formula is given by:
\[x = \frac{-b \pm \sqrt{b^2 - 4ac}}{2a}\]

\subsection{Matrix Example}
A sample matrix:
\[
\begin{pmatrix}
a_{11} & a_{12} & a_{13} \\
a_{21} & a_{22} & a_{23} \\
a_{31} & a_{32} & a_{33}
\end{pmatrix}
\]

\section{Tables}
Below is a sample table showing data:

\begin{table}[h]
\centering
\begin{tabular}{@{}lrrr@{}}
\toprule
\textbf{Category} & \textbf{Value 1} & \textbf{Value 2} & \textbf{Total} \\
\midrule
Item A & 10.5 & 20.3 & 30.8 \\
Item B & 15.2 & 25.1 & 40.3 \\
Item C & 22.1 & 12.7 & 34.8 \\
\bottomrule
\end{tabular}
\caption{Sample Data Table}
\label{tab:sample}
\end{table}

\section{Lists}
\subsection{Bullet Points}
\begin{itemize}
  \item First item
  \item Second item
  \item Third item with sub-items:
    \begin{itemize}
      \item Sub-item 1
      \item Sub-item 2
    \end{itemize}
\end{itemize}

\subsection{Numbered List}
\begin{enumerate}
  \item First step
  \item Second step
  \item Third step
\end{enumerate}

\section{Citations and References}
When you use the citation search feature, references will be automatically managed. For example, you might cite a paper like this: \cite{example2024}. All citations will be automatically added to your \texttt{cite.bib} file.

\section{Conclusion}
This template demonstrates basic \LaTeX{} features including mathematical formulas, tables, and lists. The bibliography system is now set up to work seamlessly with Bibby AI's citation management. Feel free to modify it for your needs!

\bibliographystyle{plain}
\bibliography{cite}

\end{document}\documentclass{article}
\usepackage[utf8]{inputenc}
\usepackage{booktabs}
\usepackage{amsmath}
\usepackage{natbib}

\title{My Document}
\author{Your Name}
\date{\today}

\begin{document}

\maketitle

\section{Introduction}
This is a sample document demonstrating various \LaTeX{} features including mathematical equations, tables, and formatting. You can cite papers using the citation search feature, and they will be automatically added to your \texttt{cite.bib} file.

\section{Mathematical Expressions}
Here are some example equations:

\subsection{Quadratic Formula}
The quadratic formula is given by:
\[x = \frac{-b \pm \sqrt{b^2 - 4ac}}{2a}\]

\subsection{Matrix Example}
A sample matrix:
\[
\begin{pmatrix}
a_{11} & a_{12} & a_{13} \\
a_{21} & a_{22} & a_{23} \\
a_{31} & a_{32} & a_{33}
\end{pmatrix}
\]

\section{Tables}
Below is a sample table showing data:

\begin{table}[h]
\centering
\begin{tabular}{@{}lrrr@{}}
\toprule
\textbf{Category} & \textbf{Value 1} & \textbf{Value 2} & \textbf{Total} \\
\midrule
Item A & 10.5 & 20.3 & 30.8 \\
Item B & 15.2 & 25.1 & 40.3 \\
Item C & 22.1 & 12.7 & 34.8 \\
\bottomrule
\end{tabular}
\caption{Sample Data Table}
\label{tab:sample}
\end{table}

\section{Lists}
\subsection{Bullet Points}
\begin{itemize}
  \item First item
  \item Second item
  \item Third item with sub-items:
    \begin{itemize}
      \item Sub-item 1
      \item Sub-item 2
    \end{itemize}
\end{itemize}

\subsection{Numbered List}
\begin{enumerate}
  \item First step
  \item Second step
  \item Third step
\end{enumerate}

\section{Citations and References}
When you use the citation search feature, references will be automatically managed. For example, you might cite a paper like this: \cite{example2024}. All citations will be automatically added to your \texttt{cite.bib} file.

\section{Conclusion}
This template demonstrates basic \LaTeX{} features including mathematical formulas, tables, and lists. The bibliography system is now set up to work seamlessly with Bibby AI's citation management. Feel free to modify it for your needs!

\bibliographystyle{plain}
\bibliography{cite}

\end{document}\documentclass{article}
\usepackage[utf8]{inputenc}
\usepackage{booktabs}
\usepackage{amsmath}
\usepackage{natbib}

\title{My Document}
\author{Your Name}
\date{\today}

\begin{document}

\maketitle

\section{Introduction}
This is a sample document demonstrating various \LaTeX{} features including mathematical equations, tables, and formatting. You can cite papers using the citation search feature, and they will be automatically added to your \texttt{cite.bib} file.

\section{Mathematical Expressions}
Here are some example equations:

\subsection{Quadratic Formula}
The quadratic formula is given by:
\[x = \frac{-b \pm \sqrt{b^2 - 4ac}}{2a}\]

\subsection{Matrix Example}
A sample matrix:
\[
\begin{pmatrix}
a_{11} & a_{12} & a_{13} \\
a_{21} & a_{22} & a_{23} \\
a_{31} & a_{32} & a_{33}
\end{pmatrix}
\]

\section{Tables}
Below is a sample table showing data:

\begin{table}[h]
\centering
\begin{tabular}{@{}lrrr@{}}
\toprule
\textbf{Category} & \textbf{Value 1} & \textbf{Value 2} & \textbf{Total} \\
\midrule
Item A & 10.5 & 20.3 & 30.8 \\
Item B & 15.2 & 25.1 & 40.3 \\
Item C & 22.1 & 12.7 & 34.8 \\
\bottomrule
\end{tabular}
\caption{Sample Data Table}
\label{tab:sample}
\end{table}

\section{Lists}
\subsection{Bullet Points}
\begin{itemize}
  \item First item
  \item Second item
  \item Third item with sub-items:
    \begin{itemize}
      \item Sub-item 1
      \item Sub-item 2
    \end{itemize}
\end{itemize}

\subsection{Numbered List}
\begin{enumerate}
  \item First step
  \item Second step
  \item Third step
\end{enumerate}

\section{Citations and References}
When you use the citation search feature, references will be automatically managed. For example, you might cite a paper like this: \cite{example2024}. All citations will be automatically added to your \texttt{cite.bib} file.

\section{Conclusion}
This template demonstrates basic \LaTeX{} features including mathematical formulas, tables, and lists. The bibliography system is now set up to work seamlessly with Bibby AI's citation management. Feel free to modify it for your needs!

\bibliographystyle{plain}
\bibliography{cite}

\end{document}\documentclass{article}
\usepackage[utf8]{inputenc}
\usepackage{booktabs}
\usepackage{amsmath}
\usepackage{natbib}

\title{My Document}
\author{Your Name}
\date{\today}

\begin{document}

\maketitle

\section{Introduction}
This is a sample document demonstrating various \LaTeX{} features including mathematical equations, tables, and formatting. You can cite papers using the citation search feature, and they will be automatically added to your \texttt{cite.bib} file.

\section{Mathematical Expressions}
Here are some example equations:

\subsection{Quadratic Formula}
The quadratic formula is given by:
\[x = \frac{-b \pm \sqrt{b^2 - 4ac}}{2a}\]

\subsection{Matrix Example}
A sample matrix:
\[
\begin{pmatrix}
a_{11} & a_{12} & a_{13} \\
a_{21} & a_{22} & a_{23} \\
a_{31} & a_{32} & a_{33}
\end{pmatrix}
\]

\section{Tables}
Below is a sample table showing data:

\begin{table}[h]
\centering
\begin{tabular}{@{}lrrr@{}}
\toprule
\textbf{Category} & \textbf{Value 1} & \textbf{Value 2} & \textbf{Total} \\
\midrule
Item A & 10.5 & 20.3 & 30.8 \\
Item B & 15.2 & 25.1 & 40.3 \\
Item C & 22.1 & 12.7 & 34.8 \\
\bottomrule
\end{tabular}
\caption{Sample Data Table}
\label{tab:sample}
\end{table}

\section{Lists}
\subsection{Bullet Points}
\begin{itemize}
  \item First item
  \item Second item
  \item Third item with sub-items:
    \begin{itemize}
      \item Sub-item 1
      \item Sub-item 2
    \end{itemize}
\end{itemize}

\subsection{Numbered List}
\begin{enumerate}
  \item First step
  \item Second step
  \item Third step
\end{enumerate}

\section{Citations and References}
When you use the citation search feature, references will be automatically managed. For example, you might cite a paper like this: \cite{example2024}. All citations will be automatically added to your \texttt{cite.bib} file.

\section{Conclusion}
This template demonstrates basic \LaTeX{} features including mathematical formulas, tables, and lists. The bibliography system is now set up to work seamlessly with Bibby AI's citation management. Feel free to modify it for your needs!

\bibliographystyle{plain}
\bibliography{cite}

\end{document}\documentclass{article}
\usepackage[utf8]{inputenc}
\usepackage{booktabs}
\usepackage{amsmath}
\usepackage{natbib}

\title{My Document}
\author{Your Name}
\date{\today}

\begin{document}

\maketitle

\section{Introduction}
This is a sample document demonstrating various \LaTeX{} features including mathematical equations, tables, and formatting. You can cite papers using the citation search feature, and they will be automatically added to your \texttt{cite.bib} file.

\section{Mathematical Expressions}
Here are some example equations:

\subsection{Quadratic Formula}
The quadratic formula is given by:
\[x = \frac{-b \pm \sqrt{b^2 - 4ac}}{2a}\]

\subsection{Matrix Example}
A sample matrix:
\[
\begin{pmatrix}
a_{11} & a_{12} & a_{13} \\
a_{21} & a_{22} & a_{23} \\
a_{31} & a_{32} & a_{33}
\end{pmatrix}
\]

\section{Tables}
Below is a sample table showing data:

\begin{table}[h]
\centering
\begin{tabular}{@{}lrrr@{}}
\toprule
\textbf{Category} & \textbf{Value 1} & \textbf{Value 2} & \textbf{Total} \\
\midrule
Item A & 10.5 & 20.3 & 30.8 \\
Item B & 15.2 & 25.1 & 40.3 \\
Item C & 22.1 & 12.7 & 34.8 \\
\bottomrule
\end{tabular}
\caption{Sample Data Table}
\label{tab:sample}
\end{table}

\section{Lists}
\subsection{Bullet Points}
\begin{itemize}
  \item First item
  \item Second item
  \item Third item with sub-items:
    \begin{itemize}
      \item Sub-item 1
      \item Sub-item 2
    \end{itemize}
\end{itemize}

\subsection{Numbered List}
\begin{enumerate}
  \item First step
  \item Second step
  \item Third step
\end{enumerate}

\section{Citations and References}
When you use the citation search feature, references will be automatically managed. For example, you might cite a paper like this: \cite{example2024}. All citations will be automatically added to your \texttt{cite.bib} file.

\section{Conclusion}
This template demonstrates basic \LaTeX{} features including mathematical formulas, tables, and lists. The bibliography system is now set up to work seamlessly with Bibby AI's citation management. Feel free to modify it for your needs!

\bibliographystyle{plain}
\bibliography{cite}

\end{document}\documentclass{article}
\usepackage[utf8]{inputenc}
\usepackage{booktabs}
\usepackage{amsmath}
\usepackage{natbib}

\title{My Document}
\author{Your Name}
\date{\today}

\begin{document}

\maketitle

\section{Introduction}
This is a sample document demonstrating various \LaTeX{} features including mathematical equations, tables, and formatting. You can cite papers using the citation search feature, and they will be automatically added to your \texttt{cite.bib} file.

\section{Mathematical Expressions}
Here are some example equations:

\subsection{Quadratic Formula}
The quadratic formula is given by:
\[x = \frac{-b \pm \sqrt{b^2 - 4ac}}{2a}\]

\subsection{Matrix Example}
A sample matrix:
\[
\begin{pmatrix}
a_{11} & a_{12} & a_{13} \\
a_{21} & a_{22} & a_{23} \\
a_{31} & a_{32} & a_{33}
\end{pmatrix}
\]

\section{Tables}
Below is a sample table showing data:

\begin{table}[h]
\centering
\begin{tabular}{@{}lrrr@{}}
\toprule
\textbf{Category} & \textbf{Value 1} & \textbf{Value 2} & \textbf{Total} \\
\midrule
Item A & 10.5 & 20.3 & 30.8 \\
Item B & 15.2 & 25.1 & 40.3 \\
Item C & 22.1 & 12.7 & 34.8 \\
\bottomrule
\end{tabular}
\caption{Sample Data Table}
\label{tab:sample}
\end{table}

\section{Lists}
\subsection{Bullet Points}
\begin{itemize}
  \item First item
  \item Second item
  \item Third item with sub-items:
    \begin{itemize}
      \item Sub-item 1
      \item Sub-item 2
    \end{itemize}
\end{itemize}

\subsection{Numbered List}
\begin{enumerate}
  \item First step
  \item Second step
  \item Third step
\end{enumerate}

\section{Citations and References}
When you use the citation search feature, references will be automatically managed. For example, you might cite a paper like this: \cite{example2024}. All citations will be automatically added to your \texttt{cite.bib} file.

\section{Conclusion}
This template demonstrates basic \LaTeX{} features including mathematical formulas, tables, and lists. The bibliography system is now set up to work seamlessly with Bibby AI's citation management. Feel free to modify it for your needs!

\bibliographystyle{plain}
\bibliography{cite}

\end{document}\documentclass{article}
\usepackage[utf8]{inputenc}
\usepackage{booktabs}
\usepackage{amsmath}
\usepackage{natbib}

\title{My Document}
\author{Your Name}
\date{\today}

\begin{document}

\maketitle

\section{Introduction}
This is a sample document demonstrating various \LaTeX{} features including mathematical equations, tables, and formatting. You can cite papers using the citation search feature, and they will be automatically added to your \texttt{cite.bib} file.

\section{Mathematical Expressions}
Here are some example equations:

\subsection{Quadratic Formula}
The quadratic formula is given by:
\[x = \frac{-b \pm \sqrt{b^2 - 4ac}}{2a}\]

\subsection{Matrix Example}
A sample matrix:
\[
\begin{pmatrix}
a_{11} & a_{12} & a_{13} \\
a_{21} & a_{22} & a_{23} \\
a_{31} & a_{32} & a_{33}
\end{pmatrix}
\]

\section{Tables}
Below is a sample table showing data:

\begin{table}[h]
\centering
\begin{tabular}{@{}lrrr@{}}
\toprule
\textbf{Category} & \textbf{Value 1} & \textbf{Value 2} & \textbf{Total} \\
\midrule
Item A & 10.5 & 20.3 & 30.8 \\
Item B & 15.2 & 25.1 & 40.3 \\
Item C & 22.1 & 12.7 & 34.8 \\
\bottomrule
\end{tabular}
\caption{Sample Data Table}
\label{tab:sample}
\end{table}

\section{Lists}
\subsection{Bullet Points}
\begin{itemize}
  \item First item
  \item Second item
  \item Third item with sub-items:
    \begin{itemize}
      \item Sub-item 1
      \item Sub-item 2
    \end{itemize}
\end{itemize}

\subsection{Numbered List}
\begin{enumerate}
  \item First step
  \item Second step
  \item Third step
\end{enumerate}

\section{Citations and References}
When you use the citation search feature, references will be automatically managed. For example, you might cite a paper like this: \cite{example2024}. All citations will be automatically added to your \texttt{cite.bib} file.

\section{Conclusion}
This template demonstrates basic \LaTeX{} features including mathematical formulas, tables, and lists. The bibliography system is now set up to work seamlessly with Bibby AI's citation management. Feel free to modify it for your needs!

\bibliographystyle{plain}
\bibliography{cite}

\end{document}\documentclass{article}
\usepackage[utf8]{inputenc}
\usepackage{booktabs}
\usepackage{amsmath}
\usepackage{natbib}

\title{My Document}
\author{Your Name}
\date{\today}

\begin{document}

\maketitle

\section{Introduction}
This is a sample document demonstrating various \LaTeX{} features including mathematical equations, tables, and formatting. You can cite papers using the citation search feature, and they will be automatically added to your \texttt{cite.bib} file.

\section{Mathematical Expressions}
Here are some example equations:

\subsection{Quadratic Formula}
The quadratic formula is given by:
\[x = \frac{-b \pm \sqrt{b^2 - 4ac}}{2a}\]

\subsection{Matrix Example}
A sample matrix:
\[
\begin{pmatrix}
a_{11} & a_{12} & a_{13} \\
a_{21} & a_{22} & a_{23} \\
a_{31} & a_{32} & a_{33}
\end{pmatrix}
\]

\section{Tables}
Below is a sample table showing data:

\begin{table}[h]
\centering
\begin{tabular}{@{}lrrr@{}}
\toprule
\textbf{Category} & \textbf{Value 1} & \textbf{Value 2} & \textbf{Total} \\
\midrule
Item A & 10.5 & 20.3 & 30.8 \\
Item B & 15.2 & 25.1 & 40.3 \\
Item C & 22.1 & 12.7 & 34.8 \\
\bottomrule
\end{tabular}
\caption{Sample Data Table}
\label{tab:sample}
\end{table}

\section{Lists}
\subsection{Bullet Points}
\begin{itemize}
  \item First item
  \item Second item
  \item Third item with sub-items:
    \begin{itemize}
      \item Sub-item 1
      \item Sub-item 2
    \end{itemize}
\end{itemize}

\subsection{Numbered List}
\begin{enumerate}
  \item First step
  \item Second step
  \item Third step
\end{enumerate}

\section{Citations and References}
When you use the citation search feature, references will be automatically managed. For example, you might cite a paper like this: \cite{example2024}. All citations will be automatically added to your \texttt{cite.bib} file.

\section{Conclusion}
This template demonstrates basic \LaTeX{} features including mathematical formulas, tables, and lists. The bibliography system is now set up to work seamlessly with Bibby AI's citation management. Feel free to modify it for your needs!

\bibliographystyle{plain}
\bibliography{cite}

\end{document}\documentclass{article}
\usepackage[utf8]{inputenc}
\usepackage{booktabs}
\usepackage{amsmath}
\usepackage{natbib}

\title{My Document}
\author{Your Name}
\date{\today}

\begin{document}

\maketitle

\section{Introduction}
This is a sample document demonstrating various \LaTeX{} features including mathematical equations, tables, and formatting. You can cite papers using the citation search feature, and they will be automatically added to your \texttt{cite.bib} file.

\section{Mathematical Expressions}
Here are some example equations:

\subsection{Quadratic Formula}
The quadratic formula is given by:
\[x = \frac{-b \pm \sqrt{b^2 - 4ac}}{2a}\]

\subsection{Matrix Example}
A sample matrix:
\[
\begin{pmatrix}
a_{11} & a_{12} & a_{13} \\
a_{21} & a_{22} & a_{23} \\
a_{31} & a_{32} & a_{33}
\end{pmatrix}
\]

\section{Tables}
Below is a sample table showing data:

\begin{table}[h]
\centering
\begin{tabular}{@{}lrrr@{}}
\toprule
\textbf{Category} & \textbf{Value 1} & \textbf{Value 2} & \textbf{Total} \\
\midrule
Item A & 10.5 & 20.3 & 30.8 \\
Item B & 15.2 & 25.1 & 40.3 \\
Item C & 22.1 & 12.7 & 34.8 \\
\bottomrule
\end{tabular}
\caption{Sample Data Table}
\label{tab:sample}
\end{table}

\section{Lists}
\subsection{Bullet Points}
\begin{itemize}
  \item First item
  \item Second item
  \item Third item with sub-items:
    \begin{itemize}
      \item Sub-item 1
      \item Sub-item 2
    \end{itemize}
\end{itemize}

\subsection{Numbered List}
\begin{enumerate}
  \item First step
  \item Second step
  \item Third step
\end{enumerate}

\section{Citations and References}
When you use the citation search feature, references will be automatically managed. For example, you might cite a paper like this: \cite{example2024}. All citations will be automatically added to your \texttt{cite.bib} file.

\section{Conclusion}
This template demonstrates basic \LaTeX{} features including mathematical formulas, tables, and lists. The bibliography system is now set up to work seamlessly with Bibby AI's citation management. Feel free to modify it for your needs!

\bibliographystyle{plain}
\bibliography{cite}

\end{document}\documentclass{article}
\usepackage[utf8]{inputenc}
\usepackage{booktabs}
\usepackage{amsmath}
\usepackage{natbib}

\title{My Document}
\author{Your Name}
\date{\today}

\begin{document}

\maketitle

\section{Introduction}
This is a sample document demonstrating various \LaTeX{} features including mathematical equations, tables, and formatting. You can cite papers using the citation search feature, and they will be automatically added to your \texttt{cite.bib} file.

\section{Mathematical Expressions}
Here are some example equations:

\subsection{Quadratic Formula}
The quadratic formula is given by:
\[x = \frac{-b \pm \sqrt{b^2 - 4ac}}{2a}\]

\subsection{Matrix Example}
A sample matrix:
\[
\begin{pmatrix}
a_{11} & a_{12} & a_{13} \\
a_{21} & a_{22} & a_{23} \\
a_{31} & a_{32} & a_{33}
\end{pmatrix}
\]

\section{Tables}
Below is a sample table showing data:

\begin{table}[h]
\centering
\begin{tabular}{@{}lrrr@{}}
\toprule
\textbf{Category} & \textbf{Value 1} & \textbf{Value 2} & \textbf{Total} \\
\midrule
Item A & 10.5 & 20.3 & 30.8 \\
Item B & 15.2 & 25.1 & 40.3 \\
Item C & 22.1 & 12.7 & 34.8 \\
\bottomrule
\end{tabular}
\caption{Sample Data Table}
\label{tab:sample}
\end{table}

\section{Lists}
\subsection{Bullet Points}
\begin{itemize}
  \item First item
  \item Second item
  \item Third item with sub-items:
    \begin{itemize}
      \item Sub-item 1
      \item Sub-item 2
    \end{itemize}
\end{itemize}

\subsection{Numbered List}
\begin{enumerate}
  \item First step
  \item Second step
  \item Third step
\end{enumerate}

\section{Citations and References}
When you use the citation search feature, references will be automatically managed. For example, you might cite a paper like this: \cite{example2024}. All citations will be automatically added to your \texttt{cite.bib} file.

\section{Conclusion}
This template demonstrates basic \LaTeX{} features including mathematical formulas, tables, and lists. The bibliography system is now set up to work seamlessly with Bibby AI's citation management. Feel free to modify it for your needs!

\bibliographystyle{plain}
\bibliography{cite}

\end{document}\documentclass{article}
\usepackage[utf8]{inputenc}
\usepackage{booktabs}
\usepackage{amsmath}
\usepackage{natbib}

\title{My Document}
\author{Your Name}
\date{\today}

\begin{document}

\maketitle

\section{Introduction}
This is a sample document demonstrating various \LaTeX{} features including mathematical equations, tables, and formatting. You can cite papers using the citation search feature, and they will be automatically added to your \texttt{cite.bib} file.

\section{Mathematical Expressions}
Here are some example equations:

\subsection{Quadratic Formula}
The quadratic formula is given by:
\[x = \frac{-b \pm \sqrt{b^2 - 4ac}}{2a}\]

\subsection{Matrix Example}
A sample matrix:
\[
\begin{pmatrix}
a_{11} & a_{12} & a_{13} \\
a_{21} & a_{22} & a_{23} \\
a_{31} & a_{32} & a_{33}
\end{pmatrix}
\]

\section{Tables}
Below is a sample table showing data:

\begin{table}[h]
\centering
\begin{tabular}{@{}lrrr@{}}
\toprule
\textbf{Category} & \textbf{Value 1} & \textbf{Value 2} & \textbf{Total} \\
\midrule
Item A & 10.5 & 20.3 & 30.8 \\
Item B & 15.2 & 25.1 & 40.3 \\
Item C & 22.1 & 12.7 & 34.8 \\
\bottomrule
\end{tabular}
\caption{Sample Data Table}
\label{tab:sample}
\end{table}

\section{Lists}
\subsection{Bullet Points}
\begin{itemize}
  \item First item
  \item Second item
  \item Third item with sub-items:
    \begin{itemize}
      \item Sub-item 1
      \item Sub-item 2
    \end{itemize}
\end{itemize}

\subsection{Numbered List}
\begin{enumerate}
  \item First step
  \item Second step
  \item Third step
\end{enumerate}

\section{Citations and References}
When you use the citation search feature, references will be automatically managed. For example, you might cite a paper like this: \cite{example2024}. All citations will be automatically added to your \texttt{cite.bib} file.

\section{Conclusion}
This template demonstrates basic \LaTeX{} features including mathematical formulas, tables, and lists. The bibliography system is now set up to work seamlessly with Bibby AI's citation management. Feel free to modify it for your needs!

\bibliographystyle{plain}
\bibliography{cite}

\end{document}\documentclass{article}
\usepackage[utf8]{inputenc}
\usepackage{booktabs}
\usepackage{amsmath}
\usepackage{natbib}

\title{My Document}
\author{Your Name}
\date{\today}

\begin{document}

\maketitle

\section{Introduction}
This is a sample document demonstrating various \LaTeX{} features including mathematical equations, tables, and formatting. You can cite papers using the citation search feature, and they will be automatically added to your \texttt{cite.bib} file.

\section{Mathematical Expressions}
Here are some example equations:

\subsection{Quadratic Formula}
The quadratic formula is given by:
\[x = \frac{-b \pm \sqrt{b^2 - 4ac}}{2a}\]

\subsection{Matrix Example}
A sample matrix:
\[
\begin{pmatrix}
a_{11} & a_{12} & a_{13} \\
a_{21} & a_{22} & a_{23} \\
a_{31} & a_{32} & a_{33}
\end{pmatrix}
\]

\section{Tables}
Below is a sample table showing data:

\begin{table}[h]
\centering
\begin{tabular}{@{}lrrr@{}}
\toprule
\textbf{Category} & \textbf{Value 1} & \textbf{Value 2} & \textbf{Total} \\
\midrule
Item A & 10.5 & 20.3 & 30.8 \\
Item B & 15.2 & 25.1 & 40.3 \\
Item C & 22.1 & 12.7 & 34.8 \\
\bottomrule
\end{tabular}
\caption{Sample Data Table}
\label{tab:sample}
\end{table}

\section{Lists}
\subsection{Bullet Points}
\begin{itemize}
  \item First item
  \item Second item
  \item Third item with sub-items:
    \begin{itemize}
      \item Sub-item 1
      \item Sub-item 2
    \end{itemize}
\end{itemize}

\subsection{Numbered List}
\begin{enumerate}
  \item First step
  \item Second step
  \item Third step
\end{enumerate}

\section{Citations and References}
When you use the citation search feature, references will be automatically managed. For example, you might cite a paper like this: \cite{example2024}. All citations will be automatically added to your \texttt{cite.bib} file.

\section{Conclusion}
This template demonstrates basic \LaTeX{} features including mathematical formulas, tables, and lists. The bibliography system is now set up to work seamlessly with Bibby AI's citation management. Feel free to modify it for your needs!

\bibliographystyle{plain}
\bibliography{cite}

\end{document}\documentclass{article}
\usepackage[utf8]{inputenc}
\usepackage{booktabs}
\usepackage{amsmath}
\usepackage{natbib}

\title{My Document}
\author{Your Name}
\date{\today}

\begin{document}

\maketitle

\section{Introduction}
This is a sample document demonstrating various \LaTeX{} features including mathematical equations, tables, and formatting. You can cite papers using the citation search feature, and they will be automatically added to your \texttt{cite.bib} file.

\section{Mathematical Expressions}
Here are some example equations:

\subsection{Quadratic Formula}
The quadratic formula is given by:
\[x = \frac{-b \pm \sqrt{b^2 - 4ac}}{2a}\]

\subsection{Matrix Example}
A sample matrix:
\[
\begin{pmatrix}
a_{11} & a_{12} & a_{13} \\
a_{21} & a_{22} & a_{23} \\
a_{31} & a_{32} & a_{33}
\end{pmatrix}
\]

\section{Tables}
Below is a sample table showing data:

\begin{table}[h]
\centering
\begin{tabular}{@{}lrrr@{}}
\toprule
\textbf{Category} & \textbf{Value 1} & \textbf{Value 2} & \textbf{Total} \\
\midrule
Item A & 10.5 & 20.3 & 30.8 \\
Item B & 15.2 & 25.1 & 40.3 \\
Item C & 22.1 & 12.7 & 34.8 \\
\bottomrule
\end{tabular}
\caption{Sample Data Table}
\label{tab:sample}
\end{table}

\section{Lists}
\subsection{Bullet Points}
\begin{itemize}
  \item First item
  \item Second item
  \item Third item with sub-items:
    \begin{itemize}
      \item Sub-item 1
      \item Sub-item 2
    \end{itemize}
\end{itemize}

\subsection{Numbered List}
\begin{enumerate}
  \item First step
  \item Second step
  \item Third step
\end{enumerate}

\section{Citations and References}
When you use the citation search feature, references will be automatically managed. For example, you might cite a paper like this: \cite{example2024}. All citations will be automatically added to your \texttt{cite.bib} file.

\section{Conclusion}
This template demonstrates basic \LaTeX{} features including mathematical formulas, tables, and lists. The bibliography system is now set up to work seamlessly with Bibby AI's citation management. Feel free to modify it for your needs!

\bibliographystyle{plain}
\bibliography{cite}

\end{document}\documentclass{article}
\usepackage[utf8]{inputenc}
\usepackage{booktabs}
\usepackage{amsmath}
\usepackage{natbib}

\title{My Document}
\author{Your Name}
\date{\today}

\begin{document}

\maketitle

\section{Introduction}
This is a sample document demonstrating various \LaTeX{} features including mathematical equations, tables, and formatting. You can cite papers using the citation search feature, and they will be automatically added to your \texttt{cite.bib} file.

\section{Mathematical Expressions}
Here are some example equations:

\subsection{Quadratic Formula}
The quadratic formula is given by:
\[x = \frac{-b \pm \sqrt{b^2 - 4ac}}{2a}\]

\subsection{Matrix Example}
A sample matrix:
\[
\begin{pmatrix}
a_{11} & a_{12} & a_{13} \\
a_{21} & a_{22} & a_{23} \\
a_{31} & a_{32} & a_{33}
\end{pmatrix}
\]

\section{Tables}
Below is a sample table showing data:

\begin{table}[h]
\centering
\begin{tabular}{@{}lrrr@{}}
\toprule
\textbf{Category} & \textbf{Value 1} & \textbf{Value 2} & \textbf{Total} \\
\midrule
Item A & 10.5 & 20.3 & 30.8 \\
Item B & 15.2 & 25.1 & 40.3 \\
Item C & 22.1 & 12.7 & 34.8 \\
\bottomrule
\end{tabular}
\caption{Sample Data Table}
\label{tab:sample}
\end{table}

\section{Lists}
\subsection{Bullet Points}
\begin{itemize}
  \item First item
  \item Second item
  \item Third item with sub-items:
    \begin{itemize}
      \item Sub-item 1
      \item Sub-item 2
    \end{itemize}
\end{itemize}

\subsection{Numbered List}
\begin{enumerate}
  \item First step
  \item Second step
  \item Third step
\end{enumerate}

\section{Citations and References}
When you use the citation search feature, references will be automatically managed. For example, you might cite a paper like this: \cite{example2024}. All citations will be automatically added to your \texttt{cite.bib} file.

\section{Conclusion}
This template demonstrates basic \LaTeX{} features including mathematical formulas, tables, and lists. The bibliography system is now set up to work seamlessly with Bibby AI's citation management. Feel free to modify it for your needs!

\bibliographystyle{plain}
\bibliography{cite}

\end{document}\documentclass{article}
\usepackage[utf8]{inputenc}
\usepackage{booktabs}
\usepackage{amsmath}
\usepackage{natbib}

\title{My Document}
\author{Your Name}
\date{\today}

\begin{document}

\maketitle

\section{Introduction}
This is a sample document demonstrating various \LaTeX{} features including mathematical equations, tables, and formatting. You can cite papers using the citation search feature, and they will be automatically added to your \texttt{cite.bib} file.

\section{Mathematical Expressions}
Here are some example equations:

\subsection{Quadratic Formula}
The quadratic formula is given by:
\[x = \frac{-b \pm \sqrt{b^2 - 4ac}}{2a}\]

\subsection{Matrix Example}
A sample matrix:
\[
\begin{pmatrix}
a_{11} & a_{12} & a_{13} \\
a_{21} & a_{22} & a_{23} \\
a_{31} & a_{32} & a_{33}
\end{pmatrix}
\]

\section{Tables}
Below is a sample table showing data:

\begin{table}[h]
\centering
\begin{tabular}{@{}lrrr@{}}
\toprule
\textbf{Category} & \textbf{Value 1} & \textbf{Value 2} & \textbf{Total} \\
\midrule
Item A & 10.5 & 20.3 & 30.8 \\
Item B & 15.2 & 25.1 & 40.3 \\
Item C & 22.1 & 12.7 & 34.8 \\
\bottomrule
\end{tabular}
\caption{Sample Data Table}
\label{tab:sample}
\end{table}

\section{Lists}
\subsection{Bullet Points}
\begin{itemize}
  \item First item
  \item Second item
  \item Third item with sub-items:
    \begin{itemize}
      \item Sub-item 1
      \item Sub-item 2
    \end{itemize}
\end{itemize}

\subsection{Numbered List}
\begin{enumerate}
  \item First step
  \item Second step
  \item Third step
\end{enumerate}

\section{Citations and References}
When you use the citation search feature, references will be automatically managed. For example, you might cite a paper like this: \cite{example2024}. All citations will be automatically added to your \texttt{cite.bib} file.

\section{Conclusion}
This template demonstrates basic \LaTeX{} features including mathematical formulas, tables, and lists. The bibliography system is now set up to work seamlessly with Bibby AI's citation management. Feel free to modify it for your needs!

\bibliographystyle{plain}
\bibliography{cite}

\end{document}\documentclass{article}
\usepackage[utf8]{inputenc}
\usepackage{booktabs}
\usepackage{amsmath}
\usepackage{natbib}

\title{My Document}
\author{Your Name}
\date{\today}

\begin{document}

\maketitle

\section{Introduction}
This is a sample document demonstrating various \LaTeX{} features including mathematical equations, tables, and formatting. You can cite papers using the citation search feature, and they will be automatically added to your \texttt{cite.bib} file.

\section{Mathematical Expressions}
Here are some example equations:

\subsection{Quadratic Formula}
The quadratic formula is given by:
\[x = \frac{-b \pm \sqrt{b^2 - 4ac}}{2a}\]

\subsection{Matrix Example}
A sample matrix:
\[
\begin{pmatrix}
a_{11} & a_{12} & a_{13} \\
a_{21} & a_{22} & a_{23} \\
a_{31} & a_{32} & a_{33}
\end{pmatrix}
\]

\section{Tables}
Below is a sample table showing data:

\begin{table}[h]
\centering
\begin{tabular}{@{}lrrr@{}}
\toprule
\textbf{Category} & \textbf{Value 1} & \textbf{Value 2} & \textbf{Total} \\
\midrule
Item A & 10.5 & 20.3 & 30.8 \\
Item B & 15.2 & 25.1 & 40.3 \\
Item C & 22.1 & 12.7 & 34.8 \\
\bottomrule
\end{tabular}
\caption{Sample Data Table}
\label{tab:sample}
\end{table}

\section{Lists}
\subsection{Bullet Points}
\begin{itemize}
  \item First item
  \item Second item
  \item Third item with sub-items:
    \begin{itemize}
      \item Sub-item 1
      \item Sub-item 2
    \end{itemize}
\end{itemize}

\subsection{Numbered List}
\begin{enumerate}
  \item First step
  \item Second step
  \item Third step
\end{enumerate}

\section{Citations and References}
When you use the citation search feature, references will be automatically managed. For example, you might cite a paper like this: \cite{example2024}. All citations will be automatically added to your \texttt{cite.bib} file.

\section{Conclusion}
This template demonstrates basic \LaTeX{} features including mathematical formulas, tables, and lists. The bibliography system is now set up to work seamlessly with Bibby AI's citation management. Feel free to modify it for your needs!

\bibliographystyle{plain}
\bibliography{cite}

\end{document}\documentclass{article}
\usepackage[utf8]{inputenc}
\usepackage{booktabs}
\usepackage{amsmath}
\usepackage{natbib}

\title{My Document}
\author{Your Name}
\date{\today}

\begin{document}

\maketitle

\section{Introduction}
This is a sample document demonstrating various \LaTeX{} features including mathematical equations, tables, and formatting. You can cite papers using the citation search feature, and they will be automatically added to your \texttt{cite.bib} file.

\section{Mathematical Expressions}
Here are some example equations:

\subsection{Quadratic Formula}
The quadratic formula is given by:
\[x = \frac{-b \pm \sqrt{b^2 - 4ac}}{2a}\]

\subsection{Matrix Example}
A sample matrix:
\[
\begin{pmatrix}
a_{11} & a_{12} & a_{13} \\
a_{21} & a_{22} & a_{23} \\
a_{31} & a_{32} & a_{33}
\end{pmatrix}
\]

\section{Tables}
Below is a sample table showing data:

\begin{table}[h]
\centering
\begin{tabular}{@{}lrrr@{}}
\toprule
\textbf{Category} & \textbf{Value 1} & \textbf{Value 2} & \textbf{Total} \\
\midrule
Item A & 10.5 & 20.3 & 30.8 \\
Item B & 15.2 & 25.1 & 40.3 \\
Item C & 22.1 & 12.7 & 34.8 \\
\bottomrule
\end{tabular}
\caption{Sample Data Table}
\label{tab:sample}
\end{table}

\section{Lists}
\subsection{Bullet Points}
\begin{itemize}
  \item First item
  \item Second item
  \item Third item with sub-items:
    \begin{itemize}
      \item Sub-item 1
      \item Sub-item 2
    \end{itemize}
\end{itemize}

\subsection{Numbered List}
\begin{enumerate}
  \item First step
  \item Second step
  \item Third step
\end{enumerate}

\section{Citations and References}
When you use the citation search feature, references will be automatically managed. For example, you might cite a paper like this: \cite{example2024}. All citations will be automatically added to your \texttt{cite.bib} file.

\section{Conclusion}
This template demonstrates basic \LaTeX{} features including mathematical formulas, tables, and lists. The bibliography system is now set up to work seamlessly with Bibby AI's citation management. Feel free to modify it for your needs!

\bibliographystyle{plain}
\bibliography{cite}

\end{document}\documentclass{article}
\usepackage[utf8]{inputenc}
\usepackage{booktabs}
\usepackage{amsmath}
\usepackage{natbib}

\title{My Document}
\author{Your Name}
\date{\today}

\begin{document}

\maketitle

\section{Introduction}
This is a sample document demonstrating various \LaTeX{} features including mathematical equations, tables, and formatting. You can cite papers using the citation search feature, and they will be automatically added to your \texttt{cite.bib} file.

\section{Mathematical Expressions}
Here are some example equations:

\subsection{Quadratic Formula}
The quadratic formula is given by:
\[x = \frac{-b \pm \sqrt{b^2 - 4ac}}{2a}\]

\subsection{Matrix Example}
A sample matrix:
\[
\begin{pmatrix}
a_{11} & a_{12} & a_{13} \\
a_{21} & a_{22} & a_{23} \\
a_{31} & a_{32} & a_{33}
\end{pmatrix}
\]

\section{Tables}
Below is a sample table showing data:

\begin{table}[h]
\centering
\begin{tabular}{@{}lrrr@{}}
\toprule
\textbf{Category} & \textbf{Value 1} & \textbf{Value 2} & \textbf{Total} \\
\midrule
Item A & 10.5 & 20.3 & 30.8 \\
Item B & 15.2 & 25.1 & 40.3 \\
Item C & 22.1 & 12.7 & 34.8 \\
\bottomrule
\end{tabular}
\caption{Sample Data Table}
\label{tab:sample}
\end{table}

\section{Lists}
\subsection{Bullet Points}
\begin{itemize}
  \item First item
  \item Second item
  \item Third item with sub-items:
    \begin{itemize}
      \item Sub-item 1
      \item Sub-item 2
    \end{itemize}
\end{itemize}

\subsection{Numbered List}
\begin{enumerate}
  \item First step
  \item Second step
  \item Third step
\end{enumerate}

\section{Citations and References}
When you use the citation search feature, references will be automatically managed. For example, you might cite a paper like this: \cite{example2024}. All citations will be automatically added to your \texttt{cite.bib} file.

\section{Conclusion}
This template demonstrates basic \LaTeX{} features including mathematical formulas, tables, and lists. The bibliography system is now set up to work seamlessly with Bibby AI's citation management. Feel free to modify it for your needs!

\bibliographystyle{plain}
\bibliography{cite}

\end{document}\documentclass{article}
\usepackage[utf8]{inputenc}
\usepackage{booktabs}
\usepackage{amsmath}
\usepackage{natbib}

\title{My Document}
\author{Your Name}
\date{\today}

\begin{document}

\maketitle

\section{Introduction}
This is a sample document demonstrating various \LaTeX{} features including mathematical equations, tables, and formatting. You can cite papers using the citation search feature, and they will be automatically added to your \texttt{cite.bib} file.

\section{Mathematical Expressions}
Here are some example equations:

\subsection{Quadratic Formula}
The quadratic formula is given by:
\[x = \frac{-b \pm \sqrt{b^2 - 4ac}}{2a}\]

\subsection{Matrix Example}
A sample matrix:
\[
\begin{pmatrix}
a_{11} & a_{12} & a_{13} \\
a_{21} & a_{22} & a_{23} \\
a_{31} & a_{32} & a_{33}
\end{pmatrix}
\]

\section{Tables}
Below is a sample table showing data:

\begin{table}[h]
\centering
\begin{tabular}{@{}lrrr@{}}
\toprule
\textbf{Category} & \textbf{Value 1} & \textbf{Value 2} & \textbf{Total} \\
\midrule
Item A & 10.5 & 20.3 & 30.8 \\
Item B & 15.2 & 25.1 & 40.3 \\
Item C & 22.1 & 12.7 & 34.8 \\
\bottomrule
\end{tabular}
\caption{Sample Data Table}
\label{tab:sample}
\end{table}

\section{Lists}
\subsection{Bullet Points}
\begin{itemize}
  \item First item
  \item Second item
  \item Third item with sub-items:
    \begin{itemize}
      \item Sub-item 1
      \item Sub-item 2
    \end{itemize}
\end{itemize}

\subsection{Numbered List}
\begin{enumerate}
  \item First step
  \item Second step
  \item Third step
\end{enumerate}

\section{Citations and References}
When you use the citation search feature, references will be automatically managed. For example, you might cite a paper like this: \cite{example2024}. All citations will be automatically added to your \texttt{cite.bib} file.

\section{Conclusion}
This template demonstrates basic \LaTeX{} features including mathematical formulas, tables, and lists. The bibliography system is now set up to work seamlessly with Bibby AI's citation management. Feel free to modify it for your needs!

\bibliographystyle{plain}
\bibliography{cite}

\end{document}\documentclass{article}
\usepackage[utf8]{inputenc}
\usepackage{booktabs}
\usepackage{amsmath}
\usepackage{natbib}

\title{My Document}
\author{Your Name}
\date{\today}

\begin{document}

\maketitle

\section{Introduction}
This is a sample document demonstrating various \LaTeX{} features including mathematical equations, tables, and formatting. You can cite papers using the citation search feature, and they will be automatically added to your \texttt{cite.bib} file.

\section{Mathematical Expressions}
Here are some example equations:

\subsection{Quadratic Formula}
The quadratic formula is given by:
\[x = \frac{-b \pm \sqrt{b^2 - 4ac}}{2a}\]

\subsection{Matrix Example}
A sample matrix:
\[
\begin{pmatrix}
a_{11} & a_{12} & a_{13} \\
a_{21} & a_{22} & a_{23} \\
a_{31} & a_{32} & a_{33}
\end{pmatrix}
\]

\section{Tables}
Below is a sample table showing data:

\begin{table}[h]
\centering
\begin{tabular}{@{}lrrr@{}}
\toprule
\textbf{Category} & \textbf{Value 1} & \textbf{Value 2} & \textbf{Total} \\
\midrule
Item A & 10.5 & 20.3 & 30.8 \\
Item B & 15.2 & 25.1 & 40.3 \\
Item C & 22.1 & 12.7 & 34.8 \\
\bottomrule
\end{tabular}
\caption{Sample Data Table}
\label{tab:sample}
\end{table}

\section{Lists}
\subsection{Bullet Points}
\begin{itemize}
  \item First item
  \item Second item
  \item Third item with sub-items:
    \begin{itemize}
      \item Sub-item 1
      \item Sub-item 2
    \end{itemize}
\end{itemize}

\subsection{Numbered List}
\begin{enumerate}
  \item First step
  \item Second step
  \item Third step
\end{enumerate}

\section{Citations and References}
When you use the citation search feature, references will be automatically managed. For example, you might cite a paper like this: \cite{example2024}. All citations will be automatically added to your \texttt{cite.bib} file.

\section{Conclusion}
This template demonstrates basic \LaTeX{} features including mathematical formulas, tables, and lists. The bibliography system is now set up to work seamlessly with Bibby AI's citation management. Feel free to modify it for your needs!

\bibliographystyle{plain}
\bibliography{cite}

\end{document}\documentclass{article}
\usepackage[utf8]{inputenc}
\usepackage{booktabs}
\usepackage{amsmath}
\usepackage{natbib}

\title{My Document}
\author{Your Name}
\date{\today}

\begin{document}

\maketitle

\section{Introduction}
This is a sample document demonstrating various \LaTeX{} features including mathematical equations, tables, and formatting. You can cite papers using the citation search feature, and they will be automatically added to your \texttt{cite.bib} file.

\section{Mathematical Expressions}
Here are some example equations:

\subsection{Quadratic Formula}
The quadratic formula is given by:
\[x = \frac{-b \pm \sqrt{b^2 - 4ac}}{2a}\]

\subsection{Matrix Example}
A sample matrix:
\[
\begin{pmatrix}
a_{11} & a_{12} & a_{13} \\
a_{21} & a_{22} & a_{23} \\
a_{31} & a_{32} & a_{33}
\end{pmatrix}
\]

\section{Tables}
Below is a sample table showing data:

\begin{table}[h]
\centering
\begin{tabular}{@{}lrrr@{}}
\toprule
\textbf{Category} & \textbf{Value 1} & \textbf{Value 2} & \textbf{Total} \\
\midrule
Item A & 10.5 & 20.3 & 30.8 \\
Item B & 15.2 & 25.1 & 40.3 \\
Item C & 22.1 & 12.7 & 34.8 \\
\bottomrule
\end{tabular}
\caption{Sample Data Table}
\label{tab:sample}
\end{table}

\section{Lists}
\subsection{Bullet Points}
\begin{itemize}
  \item First item
  \item Second item
  \item Third item with sub-items:
    \begin{itemize}
      \item Sub-item 1
      \item Sub-item 2
    \end{itemize}
\end{itemize}

\subsection{Numbered List}
\begin{enumerate}
  \item First step
  \item Second step
  \item Third step
\end{enumerate}

\section{Citations and References}
When you use the citation search feature, references will be automatically managed. For example, you might cite a paper like this: \cite{example2024}. All citations will be automatically added to your \texttt{cite.bib} file.

\section{Conclusion}
This template demonstrates basic \LaTeX{} features including mathematical formulas, tables, and lists. The bibliography system is now set up to work seamlessly with Bibby AI's citation management. Feel free to modify it for your needs!

\bibliographystyle{plain}
\bibliography{cite}

\end{document}\documentclass{article}
\usepackage[utf8]{inputenc}
\usepackage{booktabs}
\usepackage{amsmath}
\usepackage{natbib}

\title{My Document}
\author{Your Name}
\date{\today}

\begin{document}

\maketitle

\section{Introduction}
This is a sample document demonstrating various \LaTeX{} features including mathematical equations, tables, and formatting. You can cite papers using the citation search feature, and they will be automatically added to your \texttt{cite.bib} file.

\section{Mathematical Expressions}
Here are some example equations:

\subsection{Quadratic Formula}
The quadratic formula is given by:
\[x = \frac{-b \pm \sqrt{b^2 - 4ac}}{2a}\]

\subsection{Matrix Example}
A sample matrix:
\[
\begin{pmatrix}
a_{11} & a_{12} & a_{13} \\
a_{21} & a_{22} & a_{23} \\
a_{31} & a_{32} & a_{33}
\end{pmatrix}
\]

\section{Tables}
Below is a sample table showing data:

\begin{table}[h]
\centering
\begin{tabular}{@{}lrrr@{}}
\toprule
\textbf{Category} & \textbf{Value 1} & \textbf{Value 2} & \textbf{Total} \\
\midrule
Item A & 10.5 & 20.3 & 30.8 \\
Item B & 15.2 & 25.1 & 40.3 \\
Item C & 22.1 & 12.7 & 34.8 \\
\bottomrule
\end{tabular}
\caption{Sample Data Table}
\label{tab:sample}
\end{table}

\section{Lists}
\subsection{Bullet Points}
\begin{itemize}
  \item First item
  \item Second item
  \item Third item with sub-items:
    \begin{itemize}
      \item Sub-item 1
      \item Sub-item 2
    \end{itemize}
\end{itemize}

\subsection{Numbered List}
\begin{enumerate}
  \item First step
  \item Second step
  \item Third step
\end{enumerate}

\section{Citations and References}
When you use the citation search feature, references will be automatically managed. For example, you might cite a paper like this: \cite{example2024}. All citations will be automatically added to your \texttt{cite.bib} file.

\section{Conclusion}
This template demonstrates basic \LaTeX{} features including mathematical formulas, tables, and lists. The bibliography system is now set up to work seamlessly with Bibby AI's citation management. Feel free to modify it for your needs!

\bibliographystyle{plain}
\bibliography{cite}

\end{document}\documentclass{article}
\usepackage[utf8]{inputenc}
\usepackage{booktabs}
\usepackage{amsmath}
\usepackage{natbib}

\title{My Document}
\author{Your Name}
\date{\today}

\begin{document}

\maketitle

\section{Introduction}
This is a sample document demonstrating various \LaTeX{} features including mathematical equations, tables, and formatting. You can cite papers using the citation search feature, and they will be automatically added to your \texttt{cite.bib} file.

\section{Mathematical Expressions}
Here are some example equations:

\subsection{Quadratic Formula}
The quadratic formula is given by:
\[x = \frac{-b \pm \sqrt{b^2 - 4ac}}{2a}\]

\subsection{Matrix Example}
A sample matrix:
\[
\begin{pmatrix}
a_{11} & a_{12} & a_{13} \\
a_{21} & a_{22} & a_{23} \\
a_{31} & a_{32} & a_{33}
\end{pmatrix}
\]

\section{Tables}
Below is a sample table showing data:

\begin{table}[h]
\centering
\begin{tabular}{@{}lrrr@{}}
\toprule
\textbf{Category} & \textbf{Value 1} & \textbf{Value 2} & \textbf{Total} \\
\midrule
Item A & 10.5 & 20.3 & 30.8 \\
Item B & 15.2 & 25.1 & 40.3 \\
Item C & 22.1 & 12.7 & 34.8 \\
\bottomrule
\end{tabular}
\caption{Sample Data Table}
\label{tab:sample}
\end{table}

\section{Lists}
\subsection{Bullet Points}
\begin{itemize}
  \item First item
  \item Second item
  \item Third item with sub-items:
    \begin{itemize}
      \item Sub-item 1
      \item Sub-item 2
    \end{itemize}
\end{itemize}

\subsection{Numbered List}
\begin{enumerate}
  \item First step
  \item Second step
  \item Third step
\end{enumerate}

\section{Citations and References}
When you use the citation search feature, references will be automatically managed. For example, you might cite a paper like this: \cite{example2024}. All citations will be automatically added to your \texttt{cite.bib} file.

\section{Conclusion}
This template demonstrates basic \LaTeX{} features including mathematical formulas, tables, and lists. The bibliography system is now set up to work seamlessly with Bibby AI's citation management. Feel free to modify it for your needs!

\bibliographystyle{plain}
\bibliography{cite}

\end{document}\documentclass{article}
\usepackage[utf8]{inputenc}
\usepackage{booktabs}
\usepackage{amsmath}
\usepackage{natbib}

\title{My Document}
\author{Your Name}
\date{\today}

\begin{document}

\maketitle

\section{Introduction}
This is a sample document demonstrating various \LaTeX{} features including mathematical equations, tables, and formatting. You can cite papers using the citation search feature, and they will be automatically added to your \texttt{cite.bib} file.

\section{Mathematical Expressions}
Here are some example equations:

\subsection{Quadratic Formula}
The quadratic formula is given by:
\[x = \frac{-b \pm \sqrt{b^2 - 4ac}}{2a}\]

\subsection{Matrix Example}
A sample matrix:
\[
\begin{pmatrix}
a_{11} & a_{12} & a_{13} \\
a_{21} & a_{22} & a_{23} \\
a_{31} & a_{32} & a_{33}
\end{pmatrix}
\]

\section{Tables}
Below is a sample table showing data:

\begin{table}[h]
\centering
\begin{tabular}{@{}lrrr@{}}
\toprule
\textbf{Category} & \textbf{Value 1} & \textbf{Value 2} & \textbf{Total} \\
\midrule
Item A & 10.5 & 20.3 & 30.8 \\
Item B & 15.2 & 25.1 & 40.3 \\
Item C & 22.1 & 12.7 & 34.8 \\
\bottomrule
\end{tabular}
\caption{Sample Data Table}
\label{tab:sample}
\end{table}

\section{Lists}
\subsection{Bullet Points}
\begin{itemize}
  \item First item
  \item Second item
  \item Third item with sub-items:
    \begin{itemize}
      \item Sub-item 1
      \item Sub-item 2
    \end{itemize}
\end{itemize}

\subsection{Numbered List}
\begin{enumerate}
  \item First step
  \item Second step
  \item Third step
\end{enumerate}

\section{Citations and References}
When you use the citation search feature, references will be automatically managed. For example, you might cite a paper like this: \cite{example2024}. All citations will be automatically added to your \texttt{cite.bib} file.

\section{Conclusion}
This template demonstrates basic \LaTeX{} features including mathematical formulas, tables, and lists. The bibliography system is now set up to work seamlessly with Bibby AI's citation management. Feel free to modify it for your needs!

\bibliographystyle{plain}
\bibliography{cite}

\end{document}\documentclass{article}
\usepackage[utf8]{inputenc}
\usepackage{booktabs}
\usepackage{amsmath}
\usepackage{natbib}

\title{My Document}
\author{Your Name}
\date{\today}

\begin{document}

\maketitle

\section{Introduction}
This is a sample document demonstrating various \LaTeX{} features including mathematical equations, tables, and formatting. You can cite papers using the citation search feature, and they will be automatically added to your \texttt{cite.bib} file.

\section{Mathematical Expressions}
Here are some example equations:

\subsection{Quadratic Formula}
The quadratic formula is given by:
\[x = \frac{-b \pm \sqrt{b^2 - 4ac}}{2a}\]

\subsection{Matrix Example}
A sample matrix:
\[
\begin{pmatrix}
a_{11} & a_{12} & a_{13} \\
a_{21} & a_{22} & a_{23} \\
a_{31} & a_{32} & a_{33}
\end{pmatrix}
\]

\section{Tables}
Below is a sample table showing data:

\begin{table}[h]
\centering
\begin{tabular}{@{}lrrr@{}}
\toprule
\textbf{Category} & \textbf{Value 1} & \textbf{Value 2} & \textbf{Total} \\
\midrule
Item A & 10.5 & 20.3 & 30.8 \\
Item B & 15.2 & 25.1 & 40.3 \\
Item C & 22.1 & 12.7 & 34.8 \\
\bottomrule
\end{tabular}
\caption{Sample Data Table}
\label{tab:sample}
\end{table}

\section{Lists}
\subsection{Bullet Points}
\begin{itemize}
  \item First item
  \item Second item
  \item Third item with sub-items:
    \begin{itemize}
      \item Sub-item 1
      \item Sub-item 2
    \end{itemize}
\end{itemize}

\subsection{Numbered List}
\begin{enumerate}
  \item First step
  \item Second step
  \item Third step
\end{enumerate}

\section{Citations and References}
When you use the citation search feature, references will be automatically managed. For example, you might cite a paper like this: \cite{example2024}. All citations will be automatically added to your \texttt{cite.bib} file.

\section{Conclusion}
This template demonstrates basic \LaTeX{} features including mathematical formulas, tables, and lists. The bibliography system is now set up to work seamlessly with Bibby AI's citation management. Feel free to modify it for your needs!

\bibliographystyle{plain}
\bibliography{cite}

\end{document}\documentclass{article}
\usepackage[utf8]{inputenc}
\usepackage{booktabs}
\usepackage{amsmath}
\usepackage{natbib}

\title{My Document}
\author{Your Name}
\date{\today}

\begin{document}

\maketitle

\section{Introduction}
This is a sample document demonstrating various \LaTeX{} features including mathematical equations, tables, and formatting. You can cite papers using the citation search feature, and they will be automatically added to your \texttt{cite.bib} file.

\section{Mathematical Expressions}
Here are some example equations:

\subsection{Quadratic Formula}
The quadratic formula is given by:
\[x = \frac{-b \pm \sqrt{b^2 - 4ac}}{2a}\]

\subsection{Matrix Example}
A sample matrix:
\[
\begin{pmatrix}
a_{11} & a_{12} & a_{13} \\
a_{21} & a_{22} & a_{23} \\
a_{31} & a_{32} & a_{33}
\end{pmatrix}
\]

\section{Tables}
Below is a sample table showing data:

\begin{table}[h]
\centering
\begin{tabular}{@{}lrrr@{}}
\toprule
\textbf{Category} & \textbf{Value 1} & \textbf{Value 2} & \textbf{Total} \\
\midrule
Item A & 10.5 & 20.3 & 30.8 \\
Item B & 15.2 & 25.1 & 40.3 \\
Item C & 22.1 & 12.7 & 34.8 \\
\bottomrule
\end{tabular}
\caption{Sample Data Table}
\label{tab:sample}
\end{table}

\section{Lists}
\subsection{Bullet Points}
\begin{itemize}
  \item First item
  \item Second item
  \item Third item with sub-items:
    \begin{itemize}
      \item Sub-item 1
      \item Sub-item 2
    \end{itemize}
\end{itemize}

\subsection{Numbered List}
\begin{enumerate}
  \item First step
  \item Second step
  \item Third step
\end{enumerate}

\section{Citations and References}
When you use the citation search feature, references will be automatically managed. For example, you might cite a paper like this: \cite{example2024}. All citations will be automatically added to your \texttt{cite.bib} file.

\section{Conclusion}
This template demonstrates basic \LaTeX{} features including mathematical formulas, tables, and lists. The bibliography system is now set up to work seamlessly with Bibby AI's citation management. Feel free to modify it for your needs!

\bibliographystyle{plain}
\bibliography{cite}

\end{document}\documentclass{article}
\usepackage[utf8]{inputenc}
\usepackage{booktabs}
\usepackage{amsmath}
\usepackage{natbib}

\title{My Document}
\author{Your Name}
\date{\today}

\begin{document}

\maketitle

\section{Introduction}
This is a sample document demonstrating various \LaTeX{} features including mathematical equations, tables, and formatting. You can cite papers using the citation search feature, and they will be automatically added to your \texttt{cite.bib} file.

\section{Mathematical Expressions}
Here are some example equations:

\subsection{Quadratic Formula}
The quadratic formula is given by:
\[x = \frac{-b \pm \sqrt{b^2 - 4ac}}{2a}\]

\subsection{Matrix Example}
A sample matrix:
\[
\begin{pmatrix}
a_{11} & a_{12} & a_{13} \\
a_{21} & a_{22} & a_{23} \\
a_{31} & a_{32} & a_{33}
\end{pmatrix}
\]

\section{Tables}
Below is a sample table showing data:

\begin{table}[h]
\centering
\begin{tabular}{@{}lrrr@{}}
\toprule
\textbf{Category} & \textbf{Value 1} & \textbf{Value 2} & \textbf{Total} \\
\midrule
Item A & 10.5 & 20.3 & 30.8 \\
Item B & 15.2 & 25.1 & 40.3 \\
Item C & 22.1 & 12.7 & 34.8 \\
\bottomrule
\end{tabular}
\caption{Sample Data Table}
\label{tab:sample}
\end{table}

\section{Lists}
\subsection{Bullet Points}
\begin{itemize}
  \item First item
  \item Second item
  \item Third item with sub-items:
    \begin{itemize}
      \item Sub-item 1
      \item Sub-item 2
    \end{itemize}
\end{itemize}

\subsection{Numbered List}
\begin{enumerate}
  \item First step
  \item Second step
  \item Third step
\end{enumerate}

\section{Citations and References}
When you use the citation search feature, references will be automatically managed. For example, you might cite a paper like this: \cite{example2024}. All citations will be automatically added to your \texttt{cite.bib} file.

\section{Conclusion}
This template demonstrates basic \LaTeX{} features including mathematical formulas, tables, and lists. The bibliography system is now set up to work seamlessly with Bibby AI's citation management. Feel free to modify it for your needs!

\bibliographystyle{plain}
\bibliography{cite}

\end{document}\documentclass{article}
\usepackage[utf8]{inputenc}
\usepackage{booktabs}
\usepackage{amsmath}
\usepackage{natbib}

\title{My Document}
\author{Your Name}
\date{\today}

\begin{document}

\maketitle

\section{Introduction}
This is a sample document demonstrating various \LaTeX{} features including mathematical equations, tables, and formatting. You can cite papers using the citation search feature, and they will be automatically added to your \texttt{cite.bib} file.

\section{Mathematical Expressions}
Here are some example equations:

\subsection{Quadratic Formula}
The quadratic formula is given by:
\[x = \frac{-b \pm \sqrt{b^2 - 4ac}}{2a}\]

\subsection{Matrix Example}
A sample matrix:
\[
\begin{pmatrix}
a_{11} & a_{12} & a_{13} \\
a_{21} & a_{22} & a_{23} \\
a_{31} & a_{32} & a_{33}
\end{pmatrix}
\]

\section{Tables}
Below is a sample table showing data:

\begin{table}[h]
\centering
\begin{tabular}{@{}lrrr@{}}
\toprule
\textbf{Category} & \textbf{Value 1} & \textbf{Value 2} & \textbf{Total} \\
\midrule
Item A & 10.5 & 20.3 & 30.8 \\
Item B & 15.2 & 25.1 & 40.3 \\
Item C & 22.1 & 12.7 & 34.8 \\
\bottomrule
\end{tabular}
\caption{Sample Data Table}
\label{tab:sample}
\end{table}

\section{Lists}
\subsection{Bullet Points}
\begin{itemize}
  \item First item
  \item Second item
  \item Third item with sub-items:
    \begin{itemize}
      \item Sub-item 1
      \item Sub-item 2
    \end{itemize}
\end{itemize}

\subsection{Numbered List}
\begin{enumerate}
  \item First step
  \item Second step
  \item Third step
\end{enumerate}

\section{Citations and References}
When you use the citation search feature, references will be automatically managed. For example, you might cite a paper like this: \cite{example2024}. All citations will be automatically added to your \texttt{cite.bib} file.

\section{Conclusion}
This template demonstrates basic \LaTeX{} features including mathematical formulas, tables, and lists. The bibliography system is now set up to work seamlessly with Bibby AI's citation management. Feel free to modify it for your needs!

\bibliographystyle{plain}
\bibliography{cite}

\end{document}\documentclass{article}
\usepackage[utf8]{inputenc}
\usepackage{booktabs}
\usepackage{amsmath}
\usepackage{natbib}

\title{My Document}
\author{Your Name}
\date{\today}

\begin{document}

\maketitle

\section{Introduction}
This is a sample document demonstrating various \LaTeX{} features including mathematical equations, tables, and formatting. You can cite papers using the citation search feature, and they will be automatically added to your \texttt{cite.bib} file.

\section{Mathematical Expressions}
Here are some example equations:

\subsection{Quadratic Formula}
The quadratic formula is given by:
\[x = \frac{-b \pm \sqrt{b^2 - 4ac}}{2a}\]

\subsection{Matrix Example}
A sample matrix:
\[
\begin{pmatrix}
a_{11} & a_{12} & a_{13} \\
a_{21} & a_{22} & a_{23} \\
a_{31} & a_{32} & a_{33}
\end{pmatrix}
\]

\section{Tables}
Below is a sample table showing data:

\begin{table}[h]
\centering
\begin{tabular}{@{}lrrr@{}}
\toprule
\textbf{Category} & \textbf{Value 1} & \textbf{Value 2} & \textbf{Total} \\
\midrule
Item A & 10.5 & 20.3 & 30.8 \\
Item B & 15.2 & 25.1 & 40.3 \\
Item C & 22.1 & 12.7 & 34.8 \\
\bottomrule
\end{tabular}
\caption{Sample Data Table}
\label{tab:sample}
\end{table}

\section{Lists}
\subsection{Bullet Points}
\begin{itemize}
  \item First item
  \item Second item
  \item Third item with sub-items:
    \begin{itemize}
      \item Sub-item 1
      \item Sub-item 2
    \end{itemize}
\end{itemize}

\subsection{Numbered List}
\begin{enumerate}
  \item First step
  \item Second step
  \item Third step
\end{enumerate}

\section{Citations and References}
When you use the citation search feature, references will be automatically managed. For example, you might cite a paper like this: \cite{example2024}. All citations will be automatically added to your \texttt{cite.bib} file.

\section{Conclusion}
This template demonstrates basic \LaTeX{} features including mathematical formulas, tables, and lists. The bibliography system is now set up to work seamlessly with Bibby AI's citation management. Feel free to modify it for your needs!

\bibliographystyle{plain}
\bibliography{cite}

\end{document}\documentclass{article}
\usepackage[utf8]{inputenc}
\usepackage{booktabs}
\usepackage{amsmath}
\usepackage{natbib}

\title{My Document}
\author{Your Name}
\date{\today}

\begin{document}

\maketitle

\section{Introduction}
This is a sample document demonstrating various \LaTeX{} features including mathematical equations, tables, and formatting. You can cite papers using the citation search feature, and they will be automatically added to your \texttt{cite.bib} file.

\section{Mathematical Expressions}
Here are some example equations:

\subsection{Quadratic Formula}
The quadratic formula is given by:
\[x = \frac{-b \pm \sqrt{b^2 - 4ac}}{2a}\]

\subsection{Matrix Example}
A sample matrix:
\[
\begin{pmatrix}
a_{11} & a_{12} & a_{13} \\
a_{21} & a_{22} & a_{23} \\
a_{31} & a_{32} & a_{33}
\end{pmatrix}
\]

\section{Tables}
Below is a sample table showing data:

\begin{table}[h]
\centering
\begin{tabular}{@{}lrrr@{}}
\toprule
\textbf{Category} & \textbf{Value 1} & \textbf{Value 2} & \textbf{Total} \\
\midrule
Item A & 10.5 & 20.3 & 30.8 \\
Item B & 15.2 & 25.1 & 40.3 \\
Item C & 22.1 & 12.7 & 34.8 \\
\bottomrule
\end{tabular}
\caption{Sample Data Table}
\label{tab:sample}
\end{table}

\section{Lists}
\subsection{Bullet Points}
\begin{itemize}
  \item First item
  \item Second item
  \item Third item with sub-items:
    \begin{itemize}
      \item Sub-item 1
      \item Sub-item 2
    \end{itemize}
\end{itemize}

\subsection{Numbered List}
\begin{enumerate}
  \item First step
  \item Second step
  \item Third step
\end{enumerate}

\section{Citations and References}
When you use the citation search feature, references will be automatically managed. For example, you might cite a paper like this: \cite{example2024}. All citations will be automatically added to your \texttt{cite.bib} file.

\section{Conclusion}
This template demonstrates basic \LaTeX{} features including mathematical formulas, tables, and lists. The bibliography system is now set up to work seamlessly with Bibby AI's citation management. Feel free to modify it for your needs!

\bibliographystyle{plain}
\bibliography{cite}

\end{document}\documentclass{article}
\usepackage[utf8]{inputenc}
\usepackage{booktabs}
\usepackage{amsmath}
\usepackage{natbib}

\title{My Document}
\author{Your Name}
\date{\today}

\begin{document}

\maketitle

\section{Introduction}
This is a sample document demonstrating various \LaTeX{} features including mathematical equations, tables, and formatting. You can cite papers using the citation search feature, and they will be automatically added to your \texttt{cite.bib} file.

\section{Mathematical Expressions}
Here are some example equations:

\subsection{Quadratic Formula}
The quadratic formula is given by:
\[x = \frac{-b \pm \sqrt{b^2 - 4ac}}{2a}\]

\subsection{Matrix Example}
A sample matrix:
\[
\begin{pmatrix}
a_{11} & a_{12} & a_{13} \\
a_{21} & a_{22} & a_{23} \\
a_{31} & a_{32} & a_{33}
\end{pmatrix}
\]

\section{Tables}
Below is a sample table showing data:

\begin{table}[h]
\centering
\begin{tabular}{@{}lrrr@{}}
\toprule
\textbf{Category} & \textbf{Value 1} & \textbf{Value 2} & \textbf{Total} \\
\midrule
Item A & 10.5 & 20.3 & 30.8 \\
Item B & 15.2 & 25.1 & 40.3 \\
Item C & 22.1 & 12.7 & 34.8 \\
\bottomrule
\end{tabular}
\caption{Sample Data Table}
\label{tab:sample}
\end{table}

\section{Lists}
\subsection{Bullet Points}
\begin{itemize}
  \item First item
  \item Second item
  \item Third item with sub-items:
    \begin{itemize}
      \item Sub-item 1
      \item Sub-item 2
    \end{itemize}
\end{itemize}

\subsection{Numbered List}
\begin{enumerate}
  \item First step
  \item Second step
  \item Third step
\end{enumerate}

\section{Citations and References}
When you use the citation search feature, references will be automatically managed. For example, you might cite a paper like this: \cite{example2024}. All citations will be automatically added to your \texttt{cite.bib} file.

\section{Conclusion}
This template demonstrates basic \LaTeX{} features including mathematical formulas, tables, and lists. The bibliography system is now set up to work seamlessly with Bibby AI's citation management. Feel free to modify it for your needs!

\bibliographystyle{plain}
\bibliography{cite}

\end{document}\documentclass{article}
\usepackage[utf8]{inputenc}
\usepackage{booktabs}
\usepackage{amsmath}
\usepackage{natbib}

\title{My Document}
\author{Your Name}
\date{\today}

\begin{document}

\maketitle

\section{Introduction}
This is a sample document demonstrating various \LaTeX{} features including mathematical equations, tables, and formatting. You can cite papers using the citation search feature, and they will be automatically added to your \texttt{cite.bib} file.

\section{Mathematical Expressions}
Here are some example equations:

\subsection{Quadratic Formula}
The quadratic formula is given by:
\[x = \frac{-b \pm \sqrt{b^2 - 4ac}}{2a}\]

\subsection{Matrix Example}
A sample matrix:
\[
\begin{pmatrix}
a_{11} & a_{12} & a_{13} \\
a_{21} & a_{22} & a_{23} \\
a_{31} & a_{32} & a_{33}
\end{pmatrix}
\]

\section{Tables}
Below is a sample table showing data:

\begin{table}[h]
\centering
\begin{tabular}{@{}lrrr@{}}
\toprule
\textbf{Category} & \textbf{Value 1} & \textbf{Value 2} & \textbf{Total} \\
\midrule
Item A & 10.5 & 20.3 & 30.8 \\
Item B & 15.2 & 25.1 & 40.3 \\
Item C & 22.1 & 12.7 & 34.8 \\
\bottomrule
\end{tabular}
\caption{Sample Data Table}
\label{tab:sample}
\end{table}

\section{Lists}
\subsection{Bullet Points}
\begin{itemize}
  \item First item
  \item Second item
  \item Third item with sub-items:
    \begin{itemize}
      \item Sub-item 1
      \item Sub-item 2
    \end{itemize}
\end{itemize}

\subsection{Numbered List}
\begin{enumerate}
  \item First step
  \item Second step
  \item Third step
\end{enumerate}

\section{Citations and References}
When you use the citation search feature, references will be automatically managed. For example, you might cite a paper like this: \cite{example2024}. All citations will be automatically added to your \texttt{cite.bib} file.

\section{Conclusion}
This template demonstrates basic \LaTeX{} features including mathematical formulas, tables, and lists. The bibliography system is now set up to work seamlessly with Bibby AI's citation management. Feel free to modify it for your needs!

\bibliographystyle{plain}
\bibliography{cite}

\end{document}\documentclass{article}
\usepackage[utf8]{inputenc}
\usepackage{booktabs}
\usepackage{amsmath}
\usepackage{natbib}

\title{My Document}
\author{Your Name}
\date{\today}

\begin{document}

\maketitle

\section{Introduction}
This is a sample document demonstrating various \LaTeX{} features including mathematical equations, tables, and formatting. You can cite papers using the citation search feature, and they will be automatically added to your \texttt{cite.bib} file.

\section{Mathematical Expressions}
Here are some example equations:

\subsection{Quadratic Formula}
The quadratic formula is given by:
\[x = \frac{-b \pm \sqrt{b^2 - 4ac}}{2a}\]

\subsection{Matrix Example}
A sample matrix:
\[
\begin{pmatrix}
a_{11} & a_{12} & a_{13} \\
a_{21} & a_{22} & a_{23} \\
a_{31} & a_{32} & a_{33}
\end{pmatrix}
\]

\section{Tables}
Below is a sample table showing data:

\begin{table}[h]
\centering
\begin{tabular}{@{}lrrr@{}}
\toprule
\textbf{Category} & \textbf{Value 1} & \textbf{Value 2} & \textbf{Total} \\
\midrule
Item A & 10.5 & 20.3 & 30.8 \\
Item B & 15.2 & 25.1 & 40.3 \\
Item C & 22.1 & 12.7 & 34.8 \\
\bottomrule
\end{tabular}
\caption{Sample Data Table}
\label{tab:sample}
\end{table}

\section{Lists}
\subsection{Bullet Points}
\begin{itemize}
  \item First item
  \item Second item
  \item Third item with sub-items:
    \begin{itemize}
      \item Sub-item 1
      \item Sub-item 2
    \end{itemize}
\end{itemize}

\subsection{Numbered List}
\begin{enumerate}
  \item First step
  \item Second step
  \item Third step
\end{enumerate}

\section{Citations and References}
When you use the citation search feature, references will be automatically managed. For example, you might cite a paper like this: \cite{example2024}. All citations will be automatically added to your \texttt{cite.bib} file.

\section{Conclusion}
This template demonstrates basic \LaTeX{} features including mathematical formulas, tables, and lists. The bibliography system is now set up to work seamlessly with Bibby AI's citation management. Feel free to modify it for your needs!

\bibliographystyle{plain}
\bibliography{cite}

\end{document}\documentclass{article}
\usepackage[utf8]{inputenc}
\usepackage{booktabs}
\usepackage{amsmath}
\usepackage{natbib}

\title{My Document}
\author{Your Name}
\date{\today}

\begin{document}

\maketitle

\section{Introduction}
This is a sample document demonstrating various \LaTeX{} features including mathematical equations, tables, and formatting. You can cite papers using the citation search feature, and they will be automatically added to your \texttt{cite.bib} file.

\section{Mathematical Expressions}
Here are some example equations:

\subsection{Quadratic Formula}
The quadratic formula is given by:
\[x = \frac{-b \pm \sqrt{b^2 - 4ac}}{2a}\]

\subsection{Matrix Example}
A sample matrix:
\[
\begin{pmatrix}
a_{11} & a_{12} & a_{13} \\
a_{21} & a_{22} & a_{23} \\
a_{31} & a_{32} & a_{33}
\end{pmatrix}
\]

\section{Tables}
Below is a sample table showing data:

\begin{table}[h]
\centering
\begin{tabular}{@{}lrrr@{}}
\toprule
\textbf{Category} & \textbf{Value 1} & \textbf{Value 2} & \textbf{Total} \\
\midrule
Item A & 10.5 & 20.3 & 30.8 \\
Item B & 15.2 & 25.1 & 40.3 \\
Item C & 22.1 & 12.7 & 34.8 \\
\bottomrule
\end{tabular}
\caption{Sample Data Table}
\label{tab:sample}
\end{table}

\section{Lists}
\subsection{Bullet Points}
\begin{itemize}
  \item First item
  \item Second item
  \item Third item with sub-items:
    \begin{itemize}
      \item Sub-item 1
      \item Sub-item 2
    \end{itemize}
\end{itemize}

\subsection{Numbered List}
\begin{enumerate}
  \item First step
  \item Second step
  \item Third step
\end{enumerate}

\section{Citations and References}
When you use the citation search feature, references will be automatically managed. For example, you might cite a paper like this: \cite{example2024}. All citations will be automatically added to your \texttt{cite.bib} file.

\section{Conclusion}
This template demonstrates basic \LaTeX{} features including mathematical formulas, tables, and lists. The bibliography system is now set up to work seamlessly with Bibby AI's citation management. Feel free to modify it for your needs!

\bibliographystyle{plain}
\bibliography{cite}

\end{document}\documentclass{article}
\usepackage[utf8]{inputenc}
\usepackage{booktabs}
\usepackage{amsmath}
\usepackage{natbib}

\title{My Document}
\author{Your Name}
\date{\today}

\begin{document}

\maketitle

\section{Introduction}
This is a sample document demonstrating various \LaTeX{} features including mathematical equations, tables, and formatting. You can cite papers using the citation search feature, and they will be automatically added to your \texttt{cite.bib} file.

\section{Mathematical Expressions}
Here are some example equations:

\subsection{Quadratic Formula}
The quadratic formula is given by:
\[x = \frac{-b \pm \sqrt{b^2 - 4ac}}{2a}\]

\subsection{Matrix Example}
A sample matrix:
\[
\begin{pmatrix}
a_{11} & a_{12} & a_{13} \\
a_{21} & a_{22} & a_{23} \\
a_{31} & a_{32} & a_{33}
\end{pmatrix}
\]

\section{Tables}
Below is a sample table showing data:

\begin{table}[h]
\centering
\begin{tabular}{@{}lrrr@{}}
\toprule
\textbf{Category} & \textbf{Value 1} & \textbf{Value 2} & \textbf{Total} \\
\midrule
Item A & 10.5 & 20.3 & 30.8 \\
Item B & 15.2 & 25.1 & 40.3 \\
Item C & 22.1 & 12.7 & 34.8 \\
\bottomrule
\end{tabular}
\caption{Sample Data Table}
\label{tab:sample}
\end{table}

\section{Lists}
\subsection{Bullet Points}
\begin{itemize}
  \item First item
  \item Second item
  \item Third item with sub-items:
    \begin{itemize}
      \item Sub-item 1
      \item Sub-item 2
    \end{itemize}
\end{itemize}

\subsection{Numbered List}
\begin{enumerate}
  \item First step
  \item Second step
  \item Third step
\end{enumerate}

\section{Citations and References}
When you use the citation search feature, references will be automatically managed. For example, you might cite a paper like this: \cite{example2024}. All citations will be automatically added to your \texttt{cite.bib} file.

\section{Conclusion}
This template demonstrates basic \LaTeX{} features including mathematical formulas, tables, and lists. The bibliography system is now set up to work seamlessly with Bibby AI's citation management. Feel free to modify it for your needs!

\bibliographystyle{plain}
\bibliography{cite}

\end{document}\documentclass{article}
\usepackage[utf8]{inputenc}
\usepackage{booktabs}
\usepackage{amsmath}
\usepackage{natbib}

\title{My Document}
\author{Your Name}
\date{\today}

\begin{document}

\maketitle

\section{Introduction}
This is a sample document demonstrating various \LaTeX{} features including mathematical equations, tables, and formatting. You can cite papers using the citation search feature, and they will be automatically added to your \texttt{cite.bib} file.

\section{Mathematical Expressions}
Here are some example equations:

\subsection{Quadratic Formula}
The quadratic formula is given by:
\[x = \frac{-b \pm \sqrt{b^2 - 4ac}}{2a}\]

\subsection{Matrix Example}
A sample matrix:
\[
\begin{pmatrix}
a_{11} & a_{12} & a_{13} \\
a_{21} & a_{22} & a_{23} \\
a_{31} & a_{32} & a_{33}
\end{pmatrix}
\]

\section{Tables}
Below is a sample table showing data:

\begin{table}[h]
\centering
\begin{tabular}{@{}lrrr@{}}
\toprule
\textbf{Category} & \textbf{Value 1} & \textbf{Value 2} & \textbf{Total} \\
\midrule
Item A & 10.5 & 20.3 & 30.8 \\
Item B & 15.2 & 25.1 & 40.3 \\
Item C & 22.1 & 12.7 & 34.8 \\
\bottomrule
\end{tabular}
\caption{Sample Data Table}
\label{tab:sample}
\end{table}

\section{Lists}
\subsection{Bullet Points}
\begin{itemize}
  \item First item
  \item Second item
  \item Third item with sub-items:
    \begin{itemize}
      \item Sub-item 1
      \item Sub-item 2
    \end{itemize}
\end{itemize}

\subsection{Numbered List}
\begin{enumerate}
  \item First step
  \item Second step
  \item Third step
\end{enumerate}

\section{Citations and References}
When you use the citation search feature, references will be automatically managed. For example, you might cite a paper like this: \cite{example2024}. All citations will be automatically added to your \texttt{cite.bib} file.

\section{Conclusion}
This template demonstrates basic \LaTeX{} features including mathematical formulas, tables, and lists. The bibliography system is now set up to work seamlessly with Bibby AI's citation management. Feel free to modify it for your needs!

\bibliographystyle{plain}
\bibliography{cite}

\end{document}\documentclass{article}
\usepackage[utf8]{inputenc}
\usepackage{booktabs}
\usepackage{amsmath}
\usepackage{natbib}

\title{My Document}
\author{Your Name}
\date{\today}

\begin{document}

\maketitle

\section{Introduction}
This is a sample document demonstrating various \LaTeX{} features including mathematical equations, tables, and formatting. You can cite papers using the citation search feature, and they will be automatically added to your \texttt{cite.bib} file.

\section{Mathematical Expressions}
Here are some example equations:

\subsection{Quadratic Formula}
The quadratic formula is given by:
\[x = \frac{-b \pm \sqrt{b^2 - 4ac}}{2a}\]

\subsection{Matrix Example}
A sample matrix:
\[
\begin{pmatrix}
a_{11} & a_{12} & a_{13} \\
a_{21} & a_{22} & a_{23} \\
a_{31} & a_{32} & a_{33}
\end{pmatrix}
\]

\section{Tables}
Below is a sample table showing data:

\begin{table}[h]
\centering
\begin{tabular}{@{}lrrr@{}}
\toprule
\textbf{Category} & \textbf{Value 1} & \textbf{Value 2} & \textbf{Total} \\
\midrule
Item A & 10.5 & 20.3 & 30.8 \\
Item B & 15.2 & 25.1 & 40.3 \\
Item C & 22.1 & 12.7 & 34.8 \\
\bottomrule
\end{tabular}
\caption{Sample Data Table}
\label{tab:sample}
\end{table}

\section{Lists}
\subsection{Bullet Points}
\begin{itemize}
  \item First item
  \item Second item
  \item Third item with sub-items:
    \begin{itemize}
      \item Sub-item 1
      \item Sub-item 2
    \end{itemize}
\end{itemize}

\subsection{Numbered List}
\begin{enumerate}
  \item First step
  \item Second step
  \item Third step
\end{enumerate}

\section{Citations and References}
When you use the citation search feature, references will be automatically managed. For example, you might cite a paper like this: \cite{example2024}. All citations will be automatically added to your \texttt{cite.bib} file.

\section{Conclusion}
This template demonstrates basic \LaTeX{} features including mathematical formulas, tables, and lists. The bibliography system is now set up to work seamlessly with Bibby AI's citation management. Feel free to modify it for your needs!

\bibliographystyle{plain}
\bibliography{cite}

\end{document}\documentclass{article}
\usepackage[utf8]{inputenc}
\usepackage{booktabs}
\usepackage{amsmath}
\usepackage{natbib}

\title{My Document}
\author{Your Name}
\date{\today}

\begin{document}

\maketitle

\section{Introduction}
This is a sample document demonstrating various \LaTeX{} features including mathematical equations, tables, and formatting. You can cite papers using the citation search feature, and they will be automatically added to your \texttt{cite.bib} file.

\section{Mathematical Expressions}
Here are some example equations:

\subsection{Quadratic Formula}
The quadratic formula is given by:
\[x = \frac{-b \pm \sqrt{b^2 - 4ac}}{2a}\]

\subsection{Matrix Example}
A sample matrix:
\[
\begin{pmatrix}
a_{11} & a_{12} & a_{13} \\
a_{21} & a_{22} & a_{23} \\
a_{31} & a_{32} & a_{33}
\end{pmatrix}
\]

\section{Tables}
Below is a sample table showing data:

\begin{table}[h]
\centering
\begin{tabular}{@{}lrrr@{}}
\toprule
\textbf{Category} & \textbf{Value 1} & \textbf{Value 2} & \textbf{Total} \\
\midrule
Item A & 10.5 & 20.3 & 30.8 \\
Item B & 15.2 & 25.1 & 40.3 \\
Item C & 22.1 & 12.7 & 34.8 \\
\bottomrule
\end{tabular}
\caption{Sample Data Table}
\label{tab:sample}
\end{table}

\section{Lists}
\subsection{Bullet Points}
\begin{itemize}
  \item First item
  \item Second item
  \item Third item with sub-items:
    \begin{itemize}
      \item Sub-item 1
      \item Sub-item 2
    \end{itemize}
\end{itemize}

\subsection{Numbered List}
\begin{enumerate}
  \item First step
  \item Second step
  \item Third step
\end{enumerate}

\section{Citations and References}
When you use the citation search feature, references will be automatically managed. For example, you might cite a paper like this: \cite{example2024}. All citations will be automatically added to your \texttt{cite.bib} file.

\section{Conclusion}
This template demonstrates basic \LaTeX{} features including mathematical formulas, tables, and lists. The bibliography system is now set up to work seamlessly with Bibby AI's citation management. Feel free to modify it for your needs!

\bibliographystyle{plain}
\bibliography{cite}

\end{document}\documentclass{article}
\usepackage[utf8]{inputenc}
\usepackage{booktabs}
\usepackage{amsmath}
\usepackage{natbib}

\title{My Document}
\author{Your Name}
\date{\today}

\begin{document}

\maketitle

\section{Introduction}
This is a sample document demonstrating various \LaTeX{} features including mathematical equations, tables, and formatting. You can cite papers using the citation search feature, and they will be automatically added to your \texttt{cite.bib} file.

\section{Mathematical Expressions}
Here are some example equations:

\subsection{Quadratic Formula}
The quadratic formula is given by:
\[x = \frac{-b \pm \sqrt{b^2 - 4ac}}{2a}\]

\subsection{Matrix Example}
A sample matrix:
\[
\begin{pmatrix}
a_{11} & a_{12} & a_{13} \\
a_{21} & a_{22} & a_{23} \\
a_{31} & a_{32} & a_{33}
\end{pmatrix}
\]

\section{Tables}
Below is a sample table showing data:

\begin{table}[h]
\centering
\begin{tabular}{@{}lrrr@{}}
\toprule
\textbf{Category} & \textbf{Value 1} & \textbf{Value 2} & \textbf{Total} \\
\midrule
Item A & 10.5 & 20.3 & 30.8 \\
Item B & 15.2 & 25.1 & 40.3 \\
Item C & 22.1 & 12.7 & 34.8 \\
\bottomrule
\end{tabular}
\caption{Sample Data Table}
\label{tab:sample}
\end{table}

\section{Lists}
\subsection{Bullet Points}
\begin{itemize}
  \item First item
  \item Second item
  \item Third item with sub-items:
    \begin{itemize}
      \item Sub-item 1
      \item Sub-item 2
    \end{itemize}
\end{itemize}

\subsection{Numbered List}
\begin{enumerate}
  \item First step
  \item Second step
  \item Third step
\end{enumerate}

\section{Citations and References}
When you use the citation search feature, references will be automatically managed. For example, you might cite a paper like this: \cite{example2024}. All citations will be automatically added to your \texttt{cite.bib} file.

\section{Conclusion}
This template demonstrates basic \LaTeX{} features including mathematical formulas, tables, and lists. The bibliography system is now set up to work seamlessly with Bibby AI's citation management. Feel free to modify it for your needs!

\bibliographystyle{plain}
\bibliography{cite}

\end{document}\documentclass{article}
\usepackage[utf8]{inputenc}
\usepackage{booktabs}
\usepackage{amsmath}
\usepackage{natbib}

\title{My Document}
\author{Your Name}
\date{\today}

\begin{document}

\maketitle

\section{Introduction}
This is a sample document demonstrating various \LaTeX{} features including mathematical equations, tables, and formatting. You can cite papers using the citation search feature, and they will be automatically added to your \texttt{cite.bib} file.

\section{Mathematical Expressions}
Here are some example equations:

\subsection{Quadratic Formula}
The quadratic formula is given by:
\[x = \frac{-b \pm \sqrt{b^2 - 4ac}}{2a}\]

\subsection{Matrix Example}
A sample matrix:
\[
\begin{pmatrix}
a_{11} & a_{12} & a_{13} \\
a_{21} & a_{22} & a_{23} \\
a_{31} & a_{32} & a_{33}
\end{pmatrix}
\]

\section{Tables}
Below is a sample table showing data:

\begin{table}[h]
\centering
\begin{tabular}{@{}lrrr@{}}
\toprule
\textbf{Category} & \textbf{Value 1} & \textbf{Value 2} & \textbf{Total} \\
\midrule
Item A & 10.5 & 20.3 & 30.8 \\
Item B & 15.2 & 25.1 & 40.3 \\
Item C & 22.1 & 12.7 & 34.8 \\
\bottomrule
\end{tabular}
\caption{Sample Data Table}
\label{tab:sample}
\end{table}

\section{Lists}
\subsection{Bullet Points}
\begin{itemize}
  \item First item
  \item Second item
  \item Third item with sub-items:
    \begin{itemize}
      \item Sub-item 1
      \item Sub-item 2
    \end{itemize}
\end{itemize}

\subsection{Numbered List}
\begin{enumerate}
  \item First step
  \item Second step
  \item Third step
\end{enumerate}

\section{Citations and References}
When you use the citation search feature, references will be automatically managed. For example, you might cite a paper like this: \cite{example2024}. All citations will be automatically added to your \texttt{cite.bib} file.

\section{Conclusion}
This template demonstrates basic \LaTeX{} features including mathematical formulas, tables, and lists. The bibliography system is now set up to work seamlessly with Bibby AI's citation management. Feel free to modify it for your needs!

\bibliographystyle{plain}
\bibliography{cite}

\end{document}\documentclass{article}
\usepackage[utf8]{inputenc}
\usepackage{booktabs}
\usepackage{amsmath}
\usepackage{natbib}

\title{My Document}
\author{Your Name}
\date{\today}

\begin{document}

\maketitle

\section{Introduction}
This is a sample document demonstrating various \LaTeX{} features including mathematical equations, tables, and formatting. You can cite papers using the citation search feature, and they will be automatically added to your \texttt{cite.bib} file.

\section{Mathematical Expressions}
Here are some example equations:

\subsection{Quadratic Formula}
The quadratic formula is given by:
\[x = \frac{-b \pm \sqrt{b^2 - 4ac}}{2a}\]

\subsection{Matrix Example}
A sample matrix:
\[
\begin{pmatrix}
a_{11} & a_{12} & a_{13} \\
a_{21} & a_{22} & a_{23} \\
a_{31} & a_{32} & a_{33}
\end{pmatrix}
\]

\section{Tables}
Below is a sample table showing data:

\begin{table}[h]
\centering
\begin{tabular}{@{}lrrr@{}}
\toprule
\textbf{Category} & \textbf{Value 1} & \textbf{Value 2} & \textbf{Total} \\
\midrule
Item A & 10.5 & 20.3 & 30.8 \\
Item B & 15.2 & 25.1 & 40.3 \\
Item C & 22.1 & 12.7 & 34.8 \\
\bottomrule
\end{tabular}
\caption{Sample Data Table}
\label{tab:sample}
\end{table}

\section{Lists}
\subsection{Bullet Points}
\begin{itemize}
  \item First item
  \item Second item
  \item Third item with sub-items:
    \begin{itemize}
      \item Sub-item 1
      \item Sub-item 2
    \end{itemize}
\end{itemize}

\subsection{Numbered List}
\begin{enumerate}
  \item First step
  \item Second step
  \item Third step
\end{enumerate}

\section{Citations and References}
When you use the citation search feature, references will be automatically managed. For example, you might cite a paper like this: \cite{example2024}. All citations will be automatically added to your \texttt{cite.bib} file.

\section{Conclusion}
This template demonstrates basic \LaTeX{} features including mathematical formulas, tables, and lists. The bibliography system is now set up to work seamlessly with Bibby AI's citation management. Feel free to modify it for your needs!

\bibliographystyle{plain}
\bibliography{cite}

\end{document}\documentclass{article}
\usepackage[utf8]{inputenc}
\usepackage{booktabs}
\usepackage{amsmath}
\usepackage{natbib}

\title{My Document}
\author{Your Name}
\date{\today}

\begin{document}

\maketitle

\section{Introduction}
This is a sample document demonstrating various \LaTeX{} features including mathematical equations, tables, and formatting. You can cite papers using the citation search feature, and they will be automatically added to your \texttt{cite.bib} file.

\section{Mathematical Expressions}
Here are some example equations:

\subsection{Quadratic Formula}
The quadratic formula is given by:
\[x = \frac{-b \pm \sqrt{b^2 - 4ac}}{2a}\]

\subsection{Matrix Example}
A sample matrix:
\[
\begin{pmatrix}
a_{11} & a_{12} & a_{13} \\
a_{21} & a_{22} & a_{23} \\
a_{31} & a_{32} & a_{33}
\end{pmatrix}
\]

\section{Tables}
Below is a sample table showing data:

\begin{table}[h]
\centering
\begin{tabular}{@{}lrrr@{}}
\toprule
\textbf{Category} & \textbf{Value 1} & \textbf{Value 2} & \textbf{Total} \\
\midrule
Item A & 10.5 & 20.3 & 30.8 \\
Item B & 15.2 & 25.1 & 40.3 \\
Item C & 22.1 & 12.7 & 34.8 \\
\bottomrule
\end{tabular}
\caption{Sample Data Table}
\label{tab:sample}
\end{table}

\section{Lists}
\subsection{Bullet Points}
\begin{itemize}
  \item First item
  \item Second item
  \item Third item with sub-items:
    \begin{itemize}
      \item Sub-item 1
      \item Sub-item 2
    \end{itemize}
\end{itemize}

\subsection{Numbered List}
\begin{enumerate}
  \item First step
  \item Second step
  \item Third step
\end{enumerate}

\section{Citations and References}
When you use the citation search feature, references will be automatically managed. For example, you might cite a paper like this: \cite{example2024}. All citations will be automatically added to your \texttt{cite.bib} file.

\section{Conclusion}
This template demonstrates basic \LaTeX{} features including mathematical formulas, tables, and lists. The bibliography system is now set up to work seamlessly with Bibby AI's citation management. Feel free to modify it for your needs!

\bibliographystyle{plain}
\bibliography{cite}

\end{document}\documentclass{article}
\usepackage[utf8]{inputenc}
\usepackage{booktabs}
\usepackage{amsmath}
\usepackage{natbib}

\title{My Document}
\author{Your Name}
\date{\today}

\begin{document}

\maketitle

\section{Introduction}
This is a sample document demonstrating various \LaTeX{} features including mathematical equations, tables, and formatting. You can cite papers using the citation search feature, and they will be automatically added to your \texttt{cite.bib} file.

\section{Mathematical Expressions}
Here are some example equations:

\subsection{Quadratic Formula}
The quadratic formula is given by:
\[x = \frac{-b \pm \sqrt{b^2 - 4ac}}{2a}\]

\subsection{Matrix Example}
A sample matrix:
\[
\begin{pmatrix}
a_{11} & a_{12} & a_{13} \\
a_{21} & a_{22} & a_{23} \\
a_{31} & a_{32} & a_{33}
\end{pmatrix}
\]

\section{Tables}
Below is a sample table showing data:

\begin{table}[h]
\centering
\begin{tabular}{@{}lrrr@{}}
\toprule
\textbf{Category} & \textbf{Value 1} & \textbf{Value 2} & \textbf{Total} \\
\midrule
Item A & 10.5 & 20.3 & 30.8 \\
Item B & 15.2 & 25.1 & 40.3 \\
Item C & 22.1 & 12.7 & 34.8 \\
\bottomrule
\end{tabular}
\caption{Sample Data Table}
\label{tab:sample}
\end{table}

\section{Lists}
\subsection{Bullet Points}
\begin{itemize}
  \item First item
  \item Second item
  \item Third item with sub-items:
    \begin{itemize}
      \item Sub-item 1
      \item Sub-item 2
    \end{itemize}
\end{itemize}

\subsection{Numbered List}
\begin{enumerate}
  \item First step
  \item Second step
  \item Third step
\end{enumerate}

\section{Citations and References}
When you use the citation search feature, references will be automatically managed. For example, you might cite a paper like this: \cite{example2024}. All citations will be automatically added to your \texttt{cite.bib} file.

\section{Conclusion}
This template demonstrates basic \LaTeX{} features including mathematical formulas, tables, and lists. The bibliography system is now set up to work seamlessly with Bibby AI's citation management. Feel free to modify it for your needs!

\bibliographystyle{plain}
\bibliography{cite}

\end{document}\documentclass{article}
\usepackage[utf8]{inputenc}
\usepackage{booktabs}
\usepackage{amsmath}
\usepackage{natbib}

\title{My Document}
\author{Your Name}
\date{\today}

\begin{document}

\maketitle

\section{Introduction}
This is a sample document demonstrating various \LaTeX{} features including mathematical equations, tables, and formatting. You can cite papers using the citation search feature, and they will be automatically added to your \texttt{cite.bib} file.

\section{Mathematical Expressions}
Here are some example equations:

\subsection{Quadratic Formula}
The quadratic formula is given by:
\[x = \frac{-b \pm \sqrt{b^2 - 4ac}}{2a}\]

\subsection{Matrix Example}
A sample matrix:
\[
\begin{pmatrix}
a_{11} & a_{12} & a_{13} \\
a_{21} & a_{22} & a_{23} \\
a_{31} & a_{32} & a_{33}
\end{pmatrix}
\]

\section{Tables}
Below is a sample table showing data:

\begin{table}[h]
\centering
\begin{tabular}{@{}lrrr@{}}
\toprule
\textbf{Category} & \textbf{Value 1} & \textbf{Value 2} & \textbf{Total} \\
\midrule
Item A & 10.5 & 20.3 & 30.8 \\
Item B & 15.2 & 25.1 & 40.3 \\
Item C & 22.1 & 12.7 & 34.8 \\
\bottomrule
\end{tabular}
\caption{Sample Data Table}
\label{tab:sample}
\end{table}

\section{Lists}
\subsection{Bullet Points}
\begin{itemize}
  \item First item
  \item Second item
  \item Third item with sub-items:
    \begin{itemize}
      \item Sub-item 1
      \item Sub-item 2
    \end{itemize}
\end{itemize}

\subsection{Numbered List}
\begin{enumerate}
  \item First step
  \item Second step
  \item Third step
\end{enumerate}

\section{Citations and References}
When you use the citation search feature, references will be automatically managed. For example, you might cite a paper like this: \cite{example2024}. All citations will be automatically added to your \texttt{cite.bib} file.

\section{Conclusion}
This template demonstrates basic \LaTeX{} features including mathematical formulas, tables, and lists. The bibliography system is now set up to work seamlessly with Bibby AI's citation management. Feel free to modify it for your needs!

\bibliographystyle{plain}
\bibliography{cite}

\end{document}\documentclass{article}
\usepackage[utf8]{inputenc}
\usepackage{booktabs}
\usepackage{amsmath}
\usepackage{natbib}

\title{My Document}
\author{Your Name}
\date{\today}

\begin{document}

\maketitle

\section{Introduction}
This is a sample document demonstrating various \LaTeX{} features including mathematical equations, tables, and formatting. You can cite papers using the citation search feature, and they will be automatically added to your \texttt{cite.bib} file.

\section{Mathematical Expressions}
Here are some example equations:

\subsection{Quadratic Formula}
The quadratic formula is given by:
\[x = \frac{-b \pm \sqrt{b^2 - 4ac}}{2a}\]

\subsection{Matrix Example}
A sample matrix:
\[
\begin{pmatrix}
a_{11} & a_{12} & a_{13} \\
a_{21} & a_{22} & a_{23} \\
a_{31} & a_{32} & a_{33}
\end{pmatrix}
\]

\section{Tables}
Below is a sample table showing data:

\begin{table}[h]
\centering
\begin{tabular}{@{}lrrr@{}}
\toprule
\textbf{Category} & \textbf{Value 1} & \textbf{Value 2} & \textbf{Total} \\
\midrule
Item A & 10.5 & 20.3 & 30.8 \\
Item B & 15.2 & 25.1 & 40.3 \\
Item C & 22.1 & 12.7 & 34.8 \\
\bottomrule
\end{tabular}
\caption{Sample Data Table}
\label{tab:sample}
\end{table}

\section{Lists}
\subsection{Bullet Points}
\begin{itemize}
  \item First item
  \item Second item
  \item Third item with sub-items:
    \begin{itemize}
      \item Sub-item 1
      \item Sub-item 2
    \end{itemize}
\end{itemize}

\subsection{Numbered List}
\begin{enumerate}
  \item First step
  \item Second step
  \item Third step
\end{enumerate}

\section{Citations and References}
When you use the citation search feature, references will be automatically managed. For example, you might cite a paper like this: \cite{example2024}. All citations will be automatically added to your \texttt{cite.bib} file.

\section{Conclusion}
This template demonstrates basic \LaTeX{} features including mathematical formulas, tables, and lists. The bibliography system is now set up to work seamlessly with Bibby AI's citation management. Feel free to modify it for your needs!

\bibliographystyle{plain}
\bibliography{cite}

\end{document}\documentclass{article}
\usepackage[utf8]{inputenc}
\usepackage{booktabs}
\usepackage{amsmath}
\usepackage{natbib}

\title{My Document}
\author{Your Name}
\date{\today}

\begin{document}

\maketitle

\section{Introduction}
This is a sample document demonstrating various \LaTeX{} features including mathematical equations, tables, and formatting. You can cite papers using the citation search feature, and they will be automatically added to your \texttt{cite.bib} file.

\section{Mathematical Expressions}
Here are some example equations:

\subsection{Quadratic Formula}
The quadratic formula is given by:
\[x = \frac{-b \pm \sqrt{b^2 - 4ac}}{2a}\]

\subsection{Matrix Example}
A sample matrix:
\[
\begin{pmatrix}
a_{11} & a_{12} & a_{13} \\
a_{21} & a_{22} & a_{23} \\
a_{31} & a_{32} & a_{33}
\end{pmatrix}
\]

\section{Tables}
Below is a sample table showing data:

\begin{table}[h]
\centering
\begin{tabular}{@{}lrrr@{}}
\toprule
\textbf{Category} & \textbf{Value 1} & \textbf{Value 2} & \textbf{Total} \\
\midrule
Item A & 10.5 & 20.3 & 30.8 \\
Item B & 15.2 & 25.1 & 40.3 \\
Item C & 22.1 & 12.7 & 34.8 \\
\bottomrule
\end{tabular}
\caption{Sample Data Table}
\label{tab:sample}
\end{table}

\section{Lists}
\subsection{Bullet Points}
\begin{itemize}
  \item First item
  \item Second item
  \item Third item with sub-items:
    \begin{itemize}
      \item Sub-item 1
      \item Sub-item 2
    \end{itemize}
\end{itemize}

\subsection{Numbered List}
\begin{enumerate}
  \item First step
  \item Second step
  \item Third step
\end{enumerate}

\section{Citations and References}
When you use the citation search feature, references will be automatically managed. For example, you might cite a paper like this: \cite{example2024}. All citations will be automatically added to your \texttt{cite.bib} file.

\section{Conclusion}
This template demonstrates basic \LaTeX{} features including mathematical formulas, tables, and lists. The bibliography system is now set up to work seamlessly with Bibby AI's citation management. Feel free to modify it for your needs!

\bibliographystyle{plain}
\bibliography{cite}

\end{document}\documentclass{article}
\usepackage[utf8]{inputenc}
\usepackage{booktabs}
\usepackage{amsmath}
\usepackage{natbib}

\title{My Document}
\author{Your Name}
\date{\today}

\begin{document}

\maketitle

\section{Introduction}
This is a sample document demonstrating various \LaTeX{} features including mathematical equations, tables, and formatting. You can cite papers using the citation search feature, and they will be automatically added to your \texttt{cite.bib} file.

\section{Mathematical Expressions}
Here are some example equations:

\subsection{Quadratic Formula}
The quadratic formula is given by:
\[x = \frac{-b \pm \sqrt{b^2 - 4ac}}{2a}\]

\subsection{Matrix Example}
A sample matrix:
\[
\begin{pmatrix}
a_{11} & a_{12} & a_{13} \\
a_{21} & a_{22} & a_{23} \\
a_{31} & a_{32} & a_{33}
\end{pmatrix}
\]

\section{Tables}
Below is a sample table showing data:

\begin{table}[h]
\centering
\begin{tabular}{@{}lrrr@{}}
\toprule
\textbf{Category} & \textbf{Value 1} & \textbf{Value 2} & \textbf{Total} \\
\midrule
Item A & 10.5 & 20.3 & 30.8 \\
Item B & 15.2 & 25.1 & 40.3 \\
Item C & 22.1 & 12.7 & 34.8 \\
\bottomrule
\end{tabular}
\caption{Sample Data Table}
\label{tab:sample}
\end{table}

\section{Lists}
\subsection{Bullet Points}
\begin{itemize}
  \item First item
  \item Second item
  \item Third item with sub-items:
    \begin{itemize}
      \item Sub-item 1
      \item Sub-item 2
    \end{itemize}
\end{itemize}

\subsection{Numbered List}
\begin{enumerate}
  \item First step
  \item Second step
  \item Third step
\end{enumerate}

\section{Citations and References}
When you use the citation search feature, references will be automatically managed. For example, you might cite a paper like this: \cite{example2024}. All citations will be automatically added to your \texttt{cite.bib} file.

\section{Conclusion}
This template demonstrates basic \LaTeX{} features including mathematical formulas, tables, and lists. The bibliography system is now set up to work seamlessly with Bibby AI's citation management. Feel free to modify it for your needs!

\bibliographystyle{plain}
\bibliography{cite}

\end{document}\documentclass{article}
\usepackage[utf8]{inputenc}
\usepackage{booktabs}
\usepackage{amsmath}
\usepackage{natbib}

\title{My Document}
\author{Your Name}
\date{\today}

\begin{document}

\maketitle

\section{Introduction}
This is a sample document demonstrating various \LaTeX{} features including mathematical equations, tables, and formatting. You can cite papers using the citation search feature, and they will be automatically added to your \texttt{cite.bib} file.

\section{Mathematical Expressions}
Here are some example equations:

\subsection{Quadratic Formula}
The quadratic formula is given by:
\[x = \frac{-b \pm \sqrt{b^2 - 4ac}}{2a}\]

\subsection{Matrix Example}
A sample matrix:
\[
\begin{pmatrix}
a_{11} & a_{12} & a_{13} \\
a_{21} & a_{22} & a_{23} \\
a_{31} & a_{32} & a_{33}
\end{pmatrix}
\]

\section{Tables}
Below is a sample table showing data:

\begin{table}[h]
\centering
\begin{tabular}{@{}lrrr@{}}
\toprule
\textbf{Category} & \textbf{Value 1} & \textbf{Value 2} & \textbf{Total} \\
\midrule
Item A & 10.5 & 20.3 & 30.8 \\
Item B & 15.2 & 25.1 & 40.3 \\
Item C & 22.1 & 12.7 & 34.8 \\
\bottomrule
\end{tabular}
\caption{Sample Data Table}
\label{tab:sample}
\end{table}

\section{Lists}
\subsection{Bullet Points}
\begin{itemize}
  \item First item
  \item Second item
  \item Third item with sub-items:
    \begin{itemize}
      \item Sub-item 1
      \item Sub-item 2
    \end{itemize}
\end{itemize}

\subsection{Numbered List}
\begin{enumerate}
  \item First step
  \item Second step
  \item Third step
\end{enumerate}

\section{Citations and References}
When you use the citation search feature, references will be automatically managed. For example, you might cite a paper like this: \cite{example2024}. All citations will be automatically added to your \texttt{cite.bib} file.

\section{Conclusion}
This template demonstrates basic \LaTeX{} features including mathematical formulas, tables, and lists. The bibliography system is now set up to work seamlessly with Bibby AI's citation management. Feel free to modify it for your needs!

\bibliographystyle{plain}
\bibliography{cite}

\end{document}\documentclass{article}
\usepackage[utf8]{inputenc}
\usepackage{booktabs}
\usepackage{amsmath}
\usepackage{natbib}

\title{My Document}
\author{Your Name}
\date{\today}

\begin{document}

\maketitle

\section{Introduction}
This is a sample document demonstrating various \LaTeX{} features including mathematical equations, tables, and formatting. You can cite papers using the citation search feature, and they will be automatically added to your \texttt{cite.bib} file.

\section{Mathematical Expressions}
Here are some example equations:

\subsection{Quadratic Formula}
The quadratic formula is given by:
\[x = \frac{-b \pm \sqrt{b^2 - 4ac}}{2a}\]

\subsection{Matrix Example}
A sample matrix:
\[
\begin{pmatrix}
a_{11} & a_{12} & a_{13} \\
a_{21} & a_{22} & a_{23} \\
a_{31} & a_{32} & a_{33}
\end{pmatrix}
\]

\section{Tables}
Below is a sample table showing data:

\begin{table}[h]
\centering
\begin{tabular}{@{}lrrr@{}}
\toprule
\textbf{Category} & \textbf{Value 1} & \textbf{Value 2} & \textbf{Total} \\
\midrule
Item A & 10.5 & 20.3 & 30.8 \\
Item B & 15.2 & 25.1 & 40.3 \\
Item C & 22.1 & 12.7 & 34.8 \\
\bottomrule
\end{tabular}
\caption{Sample Data Table}
\label{tab:sample}
\end{table}

\section{Lists}
\subsection{Bullet Points}
\begin{itemize}
  \item First item
  \item Second item
  \item Third item with sub-items:
    \begin{itemize}
      \item Sub-item 1
      \item Sub-item 2
    \end{itemize}
\end{itemize}

\subsection{Numbered List}
\begin{enumerate}
  \item First step
  \item Second step
  \item Third step
\end{enumerate}

\section{Citations and References}
When you use the citation search feature, references will be automatically managed. For example, you might cite a paper like this: \cite{example2024}. All citations will be automatically added to your \texttt{cite.bib} file.

\section{Conclusion}
This template demonstrates basic \LaTeX{} features including mathematical formulas, tables, and lists. The bibliography system is now set up to work seamlessly with Bibby AI's citation management. Feel free to modify it for your needs!

\bibliographystyle{plain}
\bibliography{cite}

\end{document}\documentclass{article}
\usepackage[utf8]{inputenc}
\usepackage{booktabs}
\usepackage{amsmath}
\usepackage{natbib}

\title{My Document}
\author{Your Name}
\date{\today}

\begin{document}

\maketitle

\section{Introduction}
This is a sample document demonstrating various \LaTeX{} features including mathematical equations, tables, and formatting. You can cite papers using the citation search feature, and they will be automatically added to your \texttt{cite.bib} file.

\section{Mathematical Expressions}
Here are some example equations:

\subsection{Quadratic Formula}
The quadratic formula is given by:
\[x = \frac{-b \pm \sqrt{b^2 - 4ac}}{2a}\]

\subsection{Matrix Example}
A sample matrix:
\[
\begin{pmatrix}
a_{11} & a_{12} & a_{13} \\
a_{21} & a_{22} & a_{23} \\
a_{31} & a_{32} & a_{33}
\end{pmatrix}
\]

\section{Tables}
Below is a sample table showing data:

\begin{table}[h]
\centering
\begin{tabular}{@{}lrrr@{}}
\toprule
\textbf{Category} & \textbf{Value 1} & \textbf{Value 2} & \textbf{Total} \\
\midrule
Item A & 10.5 & 20.3 & 30.8 \\
Item B & 15.2 & 25.1 & 40.3 \\
Item C & 22.1 & 12.7 & 34.8 \\
\bottomrule
\end{tabular}
\caption{Sample Data Table}
\label{tab:sample}
\end{table}

\section{Lists}
\subsection{Bullet Points}
\begin{itemize}
  \item First item
  \item Second item
  \item Third item with sub-items:
    \begin{itemize}
      \item Sub-item 1
      \item Sub-item 2
    \end{itemize}
\end{itemize}

\subsection{Numbered List}
\begin{enumerate}
  \item First step
  \item Second step
  \item Third step
\end{enumerate}

\section{Citations and References}
When you use the citation search feature, references will be automatically managed. For example, you might cite a paper like this: \cite{example2024}. All citations will be automatically added to your \texttt{cite.bib} file.

\section{Conclusion}
This template demonstrates basic \LaTeX{} features including mathematical formulas, tables, and lists. The bibliography system is now set up to work seamlessly with Bibby AI's citation management. Feel free to modify it for your needs!

\bibliographystyle{plain}
\bibliography{cite}

\end{document}\documentclass{article}
\usepackage[utf8]{inputenc}
\usepackage{booktabs}
\usepackage{amsmath}
\usepackage{natbib}

\title{My Document}
\author{Your Name}
\date{\today}

\begin{document}

\maketitle

\section{Introduction}
This is a sample document demonstrating various \LaTeX{} features including mathematical equations, tables, and formatting. You can cite papers using the citation search feature, and they will be automatically added to your \texttt{cite.bib} file.

\section{Mathematical Expressions}
Here are some example equations:

\subsection{Quadratic Formula}
The quadratic formula is given by:
\[x = \frac{-b \pm \sqrt{b^2 - 4ac}}{2a}\]

\subsection{Matrix Example}
A sample matrix:
\[
\begin{pmatrix}
a_{11} & a_{12} & a_{13} \\
a_{21} & a_{22} & a_{23} \\
a_{31} & a_{32} & a_{33}
\end{pmatrix}
\]

\section{Tables}
Below is a sample table showing data:

\begin{table}[h]
\centering
\begin{tabular}{@{}lrrr@{}}
\toprule
\textbf{Category} & \textbf{Value 1} & \textbf{Value 2} & \textbf{Total} \\
\midrule
Item A & 10.5 & 20.3 & 30.8 \\
Item B & 15.2 & 25.1 & 40.3 \\
Item C & 22.1 & 12.7 & 34.8 \\
\bottomrule
\end{tabular}
\caption{Sample Data Table}
\label{tab:sample}
\end{table}

\section{Lists}
\subsection{Bullet Points}
\begin{itemize}
  \item First item
  \item Second item
  \item Third item with sub-items:
    \begin{itemize}
      \item Sub-item 1
      \item Sub-item 2
    \end{itemize}
\end{itemize}

\subsection{Numbered List}
\begin{enumerate}
  \item First step
  \item Second step
  \item Third step
\end{enumerate}

\section{Citations and References}
When you use the citation search feature, references will be automatically managed. For example, you might cite a paper like this: \cite{example2024}. All citations will be automatically added to your \texttt{cite.bib} file.

\section{Conclusion}
This template demonstrates basic \LaTeX{} features including mathematical formulas, tables, and lists. The bibliography system is now set up to work seamlessly with Bibby AI's citation management. Feel free to modify it for your needs!

\bibliographystyle{plain}
\bibliography{cite}

\end{document}\documentclass{article}
\usepackage[utf8]{inputenc}
\usepackage{booktabs}
\usepackage{amsmath}
\usepackage{natbib}

\title{My Document}
\author{Your Name}
\date{\today}

\begin{document}

\maketitle

\section{Introduction}
This is a sample document demonstrating various \LaTeX{} features including mathematical equations, tables, and formatting. You can cite papers using the citation search feature, and they will be automatically added to your \texttt{cite.bib} file.

\section{Mathematical Expressions}
Here are some example equations:

\subsection{Quadratic Formula}
The quadratic formula is given by:
\[x = \frac{-b \pm \sqrt{b^2 - 4ac}}{2a}\]

\subsection{Matrix Example}
A sample matrix:
\[
\begin{pmatrix}
a_{11} & a_{12} & a_{13} \\
a_{21} & a_{22} & a_{23} \\
a_{31} & a_{32} & a_{33}
\end{pmatrix}
\]

\section{Tables}
Below is a sample table showing data:

\begin{table}[h]
\centering
\begin{tabular}{@{}lrrr@{}}
\toprule
\textbf{Category} & \textbf{Value 1} & \textbf{Value 2} & \textbf{Total} \\
\midrule
Item A & 10.5 & 20.3 & 30.8 \\
Item B & 15.2 & 25.1 & 40.3 \\
Item C & 22.1 & 12.7 & 34.8 \\
\bottomrule
\end{tabular}
\caption{Sample Data Table}
\label{tab:sample}
\end{table}

\section{Lists}
\subsection{Bullet Points}
\begin{itemize}
  \item First item
  \item Second item
  \item Third item with sub-items:
    \begin{itemize}
      \item Sub-item 1
      \item Sub-item 2
    \end{itemize}
\end{itemize}

\subsection{Numbered List}
\begin{enumerate}
  \item First step
  \item Second step
  \item Third step
\end{enumerate}

\section{Citations and References}
When you use the citation search feature, references will be automatically managed. For example, you might cite a paper like this: \cite{example2024}. All citations will be automatically added to your \texttt{cite.bib} file.

\section{Conclusion}
This template demonstrates basic \LaTeX{} features including mathematical formulas, tables, and lists. The bibliography system is now set up to work seamlessly with Bibby AI's citation management. Feel free to modify it for your needs!

\bibliographystyle{plain}
\bibliography{cite}

\end{document}\documentclass{article}
\usepackage[utf8]{inputenc}
\usepackage{booktabs}
\usepackage{amsmath}
\usepackage{natbib}

\title{My Document}
\author{Your Name}
\date{\today}

\begin{document}

\maketitle

\section{Introduction}
This is a sample document demonstrating various \LaTeX{} features including mathematical equations, tables, and formatting. You can cite papers using the citation search feature, and they will be automatically added to your \texttt{cite.bib} file.

\section{Mathematical Expressions}
Here are some example equations:

\subsection{Quadratic Formula}
The quadratic formula is given by:
\[x = \frac{-b \pm \sqrt{b^2 - 4ac}}{2a}\]

\subsection{Matrix Example}
A sample matrix:
\[
\begin{pmatrix}
a_{11} & a_{12} & a_{13} \\
a_{21} & a_{22} & a_{23} \\
a_{31} & a_{32} & a_{33}
\end{pmatrix}
\]

\section{Tables}
Below is a sample table showing data:

\begin{table}[h]
\centering
\begin{tabular}{@{}lrrr@{}}
\toprule
\textbf{Category} & \textbf{Value 1} & \textbf{Value 2} & \textbf{Total} \\
\midrule
Item A & 10.5 & 20.3 & 30.8 \\
Item B & 15.2 & 25.1 & 40.3 \\
Item C & 22.1 & 12.7 & 34.8 \\
\bottomrule
\end{tabular}
\caption{Sample Data Table}
\label{tab:sample}
\end{table}

\section{Lists}
\subsection{Bullet Points}
\begin{itemize}
  \item First item
  \item Second item
  \item Third item with sub-items:
    \begin{itemize}
      \item Sub-item 1
      \item Sub-item 2
    \end{itemize}
\end{itemize}

\subsection{Numbered List}
\begin{enumerate}
  \item First step
  \item Second step
  \item Third step
\end{enumerate}

\section{Citations and References}
When you use the citation search feature, references will be automatically managed. For example, you might cite a paper like this: \cite{example2024}. All citations will be automatically added to your \texttt{cite.bib} file.

\section{Conclusion}
This template demonstrates basic \LaTeX{} features including mathematical formulas, tables, and lists. The bibliography system is now set up to work seamlessly with Bibby AI's citation management. Feel free to modify it for your needs!

\bibliographystyle{plain}
\bibliography{cite}

\end{document}\documentclass{article}
\usepackage[utf8]{inputenc}
\usepackage{booktabs}
\usepackage{amsmath}
\usepackage{natbib}

\title{My Document}
\author{Your Name}
\date{\today}

\begin{document}

\maketitle

\section{Introduction}
This is a sample document demonstrating various \LaTeX{} features including mathematical equations, tables, and formatting. You can cite papers using the citation search feature, and they will be automatically added to your \texttt{cite.bib} file.

\section{Mathematical Expressions}
Here are some example equations:

\subsection{Quadratic Formula}
The quadratic formula is given by:
\[x = \frac{-b \pm \sqrt{b^2 - 4ac}}{2a}\]

\subsection{Matrix Example}
A sample matrix:
\[
\begin{pmatrix}
a_{11} & a_{12} & a_{13} \\
a_{21} & a_{22} & a_{23} \\
a_{31} & a_{32} & a_{33}
\end{pmatrix}
\]

\section{Tables}
Below is a sample table showing data:

\begin{table}[h]
\centering
\begin{tabular}{@{}lrrr@{}}
\toprule
\textbf{Category} & \textbf{Value 1} & \textbf{Value 2} & \textbf{Total} \\
\midrule
Item A & 10.5 & 20.3 & 30.8 \\
Item B & 15.2 & 25.1 & 40.3 \\
Item C & 22.1 & 12.7 & 34.8 \\
\bottomrule
\end{tabular}
\caption{Sample Data Table}
\label{tab:sample}
\end{table}

\section{Lists}
\subsection{Bullet Points}
\begin{itemize}
  \item First item
  \item Second item
  \item Third item with sub-items:
    \begin{itemize}
      \item Sub-item 1
      \item Sub-item 2
    \end{itemize}
\end{itemize}

\subsection{Numbered List}
\begin{enumerate}
  \item First step
  \item Second step
  \item Third step
\end{enumerate}

\section{Citations and References}
When you use the citation search feature, references will be automatically managed. For example, you might cite a paper like this: \cite{example2024}. All citations will be automatically added to your \texttt{cite.bib} file.

\section{Conclusion}
This template demonstrates basic \LaTeX{} features including mathematical formulas, tables, and lists. The bibliography system is now set up to work seamlessly with Bibby AI's citation management. Feel free to modify it for your needs!

\bibliographystyle{plain}
\bibliography{cite}

\end{document}\documentclass{article}
\usepackage[utf8]{inputenc}
\usepackage{booktabs}
\usepackage{amsmath}
\usepackage{natbib}

\title{My Document}
\author{Your Name}
\date{\today}

\begin{document}

\maketitle

\section{Introduction}
This is a sample document demonstrating various \LaTeX{} features including mathematical equations, tables, and formatting. You can cite papers using the citation search feature, and they will be automatically added to your \texttt{cite.bib} file.

\section{Mathematical Expressions}
Here are some example equations:

\subsection{Quadratic Formula}
The quadratic formula is given by:
\[x = \frac{-b \pm \sqrt{b^2 - 4ac}}{2a}\]

\subsection{Matrix Example}
A sample matrix:
\[
\begin{pmatrix}
a_{11} & a_{12} & a_{13} \\
a_{21} & a_{22} & a_{23} \\
a_{31} & a_{32} & a_{33}
\end{pmatrix}
\]

\section{Tables}
Below is a sample table showing data:

\begin{table}[h]
\centering
\begin{tabular}{@{}lrrr@{}}
\toprule
\textbf{Category} & \textbf{Value 1} & \textbf{Value 2} & \textbf{Total} \\
\midrule
Item A & 10.5 & 20.3 & 30.8 \\
Item B & 15.2 & 25.1 & 40.3 \\
Item C & 22.1 & 12.7 & 34.8 \\
\bottomrule
\end{tabular}
\caption{Sample Data Table}
\label{tab:sample}
\end{table}

\section{Lists}
\subsection{Bullet Points}
\begin{itemize}
  \item First item
  \item Second item
  \item Third item with sub-items:
    \begin{itemize}
      \item Sub-item 1
      \item Sub-item 2
    \end{itemize}
\end{itemize}

\subsection{Numbered List}
\begin{enumerate}
  \item First step
  \item Second step
  \item Third step
\end{enumerate}

\section{Citations and References}
When you use the citation search feature, references will be automatically managed. For example, you might cite a paper like this: \cite{example2024}. All citations will be automatically added to your \texttt{cite.bib} file.

\section{Conclusion}
This template demonstrates basic \LaTeX{} features including mathematical formulas, tables, and lists. The bibliography system is now set up to work seamlessly with Bibby AI's citation management. Feel free to modify it for your needs!

\bibliographystyle{plain}
\bibliography{cite}

\end{document}\documentclass{article}
\usepackage[utf8]{inputenc}
\usepackage{booktabs}
\usepackage{amsmath}
\usepackage{natbib}

\title{My Document}
\author{Your Name}
\date{\today}

\begin{document}

\maketitle

\section{Introduction}
This is a sample document demonstrating various \LaTeX{} features including mathematical equations, tables, and formatting. You can cite papers using the citation search feature, and they will be automatically added to your \texttt{cite.bib} file.

\section{Mathematical Expressions}
Here are some example equations:

\subsection{Quadratic Formula}
The quadratic formula is given by:
\[x = \frac{-b \pm \sqrt{b^2 - 4ac}}{2a}\]

\subsection{Matrix Example}
A sample matrix:
\[
\begin{pmatrix}
a_{11} & a_{12} & a_{13} \\
a_{21} & a_{22} & a_{23} \\
a_{31} & a_{32} & a_{33}
\end{pmatrix}
\]

\section{Tables}
Below is a sample table showing data:

\begin{table}[h]
\centering
\begin{tabular}{@{}lrrr@{}}
\toprule
\textbf{Category} & \textbf{Value 1} & \textbf{Value 2} & \textbf{Total} \\
\midrule
Item A & 10.5 & 20.3 & 30.8 \\
Item B & 15.2 & 25.1 & 40.3 \\
Item C & 22.1 & 12.7 & 34.8 \\
\bottomrule
\end{tabular}
\caption{Sample Data Table}
\label{tab:sample}
\end{table}

\section{Lists}
\subsection{Bullet Points}
\begin{itemize}
  \item First item
  \item Second item
  \item Third item with sub-items:
    \begin{itemize}
      \item Sub-item 1
      \item Sub-item 2
    \end{itemize}
\end{itemize}

\subsection{Numbered List}
\begin{enumerate}
  \item First step
  \item Second step
  \item Third step
\end{enumerate}

\section{Citations and References}
When you use the citation search feature, references will be automatically managed. For example, you might cite a paper like this: \cite{example2024}. All citations will be automatically added to your \texttt{cite.bib} file.

\section{Conclusion}
This template demonstrates basic \LaTeX{} features including mathematical formulas, tables, and lists. The bibliography system is now set up to work seamlessly with Bibby AI's citation management. Feel free to modify it for your needs!

\bibliographystyle{plain}
\bibliography{cite}

\end{document}\documentclass{article}
\usepackage[utf8]{inputenc}
\usepackage{booktabs}
\usepackage{amsmath}
\usepackage{natbib}

\title{My Document}
\author{Your Name}
\date{\today}

\begin{document}

\maketitle

\section{Introduction}
This is a sample document demonstrating various \LaTeX{} features including mathematical equations, tables, and formatting. You can cite papers using the citation search feature, and they will be automatically added to your \texttt{cite.bib} file.

\section{Mathematical Expressions}
Here are some example equations:

\subsection{Quadratic Formula}
The quadratic formula is given by:
\[x = \frac{-b \pm \sqrt{b^2 - 4ac}}{2a}\]

\subsection{Matrix Example}
A sample matrix:
\[
\begin{pmatrix}
a_{11} & a_{12} & a_{13} \\
a_{21} & a_{22} & a_{23} \\
a_{31} & a_{32} & a_{33}
\end{pmatrix}
\]

\section{Tables}
Below is a sample table showing data:

\begin{table}[h]
\centering
\begin{tabular}{@{}lrrr@{}}
\toprule
\textbf{Category} & \textbf{Value 1} & \textbf{Value 2} & \textbf{Total} \\
\midrule
Item A & 10.5 & 20.3 & 30.8 \\
Item B & 15.2 & 25.1 & 40.3 \\
Item C & 22.1 & 12.7 & 34.8 \\
\bottomrule
\end{tabular}
\caption{Sample Data Table}
\label{tab:sample}
\end{table}

\section{Lists}
\subsection{Bullet Points}
\begin{itemize}
  \item First item
  \item Second item
  \item Third item with sub-items:
    \begin{itemize}
      \item Sub-item 1
      \item Sub-item 2
    \end{itemize}
\end{itemize}

\subsection{Numbered List}
\begin{enumerate}
  \item First step
  \item Second step
  \item Third step
\end{enumerate}

\section{Citations and References}
When you use the citation search feature, references will be automatically managed. For example, you might cite a paper like this: \cite{example2024}. All citations will be automatically added to your \texttt{cite.bib} file.

\section{Conclusion}
This template demonstrates basic \LaTeX{} features including mathematical formulas, tables, and lists. The bibliography system is now set up to work seamlessly with Bibby AI's citation management. Feel free to modify it for your needs!

\bibliographystyle{plain}
\bibliography{cite}

\end{document}\documentclass{article}
\usepackage[utf8]{inputenc}
\usepackage{booktabs}
\usepackage{amsmath}
\usepackage{natbib}

\title{My Document}
\author{Your Name}
\date{\today}

\begin{document}

\maketitle

\section{Introduction}
This is a sample document demonstrating various \LaTeX{} features including mathematical equations, tables, and formatting. You can cite papers using the citation search feature, and they will be automatically added to your \texttt{cite.bib} file.

\section{Mathematical Expressions}
Here are some example equations:

\subsection{Quadratic Formula}
The quadratic formula is given by:
\[x = \frac{-b \pm \sqrt{b^2 - 4ac}}{2a}\]

\subsection{Matrix Example}
A sample matrix:
\[
\begin{pmatrix}
a_{11} & a_{12} & a_{13} \\
a_{21} & a_{22} & a_{23} \\
a_{31} & a_{32} & a_{33}
\end{pmatrix}
\]

\section{Tables}
Below is a sample table showing data:

\begin{table}[h]
\centering
\begin{tabular}{@{}lrrr@{}}
\toprule
\textbf{Category} & \textbf{Value 1} & \textbf{Value 2} & \textbf{Total} \\
\midrule
Item A & 10.5 & 20.3 & 30.8 \\
Item B & 15.2 & 25.1 & 40.3 \\
Item C & 22.1 & 12.7 & 34.8 \\
\bottomrule
\end{tabular}
\caption{Sample Data Table}
\label{tab:sample}
\end{table}

\section{Lists}
\subsection{Bullet Points}
\begin{itemize}
  \item First item
  \item Second item
  \item Third item with sub-items:
    \begin{itemize}
      \item Sub-item 1
      \item Sub-item 2
    \end{itemize}
\end{itemize}

\subsection{Numbered List}
\begin{enumerate}
  \item First step
  \item Second step
  \item Third step
\end{enumerate}

\section{Citations and References}
When you use the citation search feature, references will be automatically managed. For example, you might cite a paper like this: \cite{example2024}. All citations will be automatically added to your \texttt{cite.bib} file.

\section{Conclusion}
This template demonstrates basic \LaTeX{} features including mathematical formulas, tables, and lists. The bibliography system is now set up to work seamlessly with Bibby AI's citation management. Feel free to modify it for your needs!

\bibliographystyle{plain}
\bibliography{cite}

\end{document}\documentclass{article}
\usepackage[utf8]{inputenc}
\usepackage{booktabs}
\usepackage{amsmath}
\usepackage{natbib}

\title{My Document}
\author{Your Name}
\date{\today}

\begin{document}

\maketitle

\section{Introduction}
This is a sample document demonstrating various \LaTeX{} features including mathematical equations, tables, and formatting. You can cite papers using the citation search feature, and they will be automatically added to your \texttt{cite.bib} file.

\section{Mathematical Expressions}
Here are some example equations:

\subsection{Quadratic Formula}
The quadratic formula is given by:
\[x = \frac{-b \pm \sqrt{b^2 - 4ac}}{2a}\]

\subsection{Matrix Example}
A sample matrix:
\[
\begin{pmatrix}
a_{11} & a_{12} & a_{13} \\
a_{21} & a_{22} & a_{23} \\
a_{31} & a_{32} & a_{33}
\end{pmatrix}
\]

\section{Tables}
Below is a sample table showing data:

\begin{table}[h]
\centering
\begin{tabular}{@{}lrrr@{}}
\toprule
\textbf{Category} & \textbf{Value 1} & \textbf{Value 2} & \textbf{Total} \\
\midrule
Item A & 10.5 & 20.3 & 30.8 \\
Item B & 15.2 & 25.1 & 40.3 \\
Item C & 22.1 & 12.7 & 34.8 \\
\bottomrule
\end{tabular}
\caption{Sample Data Table}
\label{tab:sample}
\end{table}

\section{Lists}
\subsection{Bullet Points}
\begin{itemize}
  \item First item
  \item Second item
  \item Third item with sub-items:
    \begin{itemize}
      \item Sub-item 1
      \item Sub-item 2
    \end{itemize}
\end{itemize}

\subsection{Numbered List}
\begin{enumerate}
  \item First step
  \item Second step
  \item Third step
\end{enumerate}

\section{Citations and References}
When you use the citation search feature, references will be automatically managed. For example, you might cite a paper like this: \cite{example2024}. All citations will be automatically added to your \texttt{cite.bib} file.

\section{Conclusion}
This template demonstrates basic \LaTeX{} features including mathematical formulas, tables, and lists. The bibliography system is now set up to work seamlessly with Bibby AI's citation management. Feel free to modify it for your needs!

\bibliographystyle{plain}
\bibliography{cite}

\end{document}\documentclass{article}
\usepackage[utf8]{inputenc}
\usepackage{booktabs}
\usepackage{amsmath}
\usepackage{natbib}

\title{My Document}
\author{Your Name}
\date{\today}

\begin{document}

\maketitle

\section{Introduction}
This is a sample document demonstrating various \LaTeX{} features including mathematical equations, tables, and formatting. You can cite papers using the citation search feature, and they will be automatically added to your \texttt{cite.bib} file.

\section{Mathematical Expressions}
Here are some example equations:

\subsection{Quadratic Formula}
The quadratic formula is given by:
\[x = \frac{-b \pm \sqrt{b^2 - 4ac}}{2a}\]

\subsection{Matrix Example}
A sample matrix:
\[
\begin{pmatrix}
a_{11} & a_{12} & a_{13} \\
a_{21} & a_{22} & a_{23} \\
a_{31} & a_{32} & a_{33}
\end{pmatrix}
\]

\section{Tables}
Below is a sample table showing data:

\begin{table}[h]
\centering
\begin{tabular}{@{}lrrr@{}}
\toprule
\textbf{Category} & \textbf{Value 1} & \textbf{Value 2} & \textbf{Total} \\
\midrule
Item A & 10.5 & 20.3 & 30.8 \\
Item B & 15.2 & 25.1 & 40.3 \\
Item C & 22.1 & 12.7 & 34.8 \\
\bottomrule
\end{tabular}
\caption{Sample Data Table}
\label{tab:sample}
\end{table}

\section{Lists}
\subsection{Bullet Points}
\begin{itemize}
  \item First item
  \item Second item
  \item Third item with sub-items:
    \begin{itemize}
      \item Sub-item 1
      \item Sub-item 2
    \end{itemize}
\end{itemize}

\subsection{Numbered List}
\begin{enumerate}
  \item First step
  \item Second step
  \item Third step
\end{enumerate}

\section{Citations and References}
When you use the citation search feature, references will be automatically managed. For example, you might cite a paper like this: \cite{example2024}. All citations will be automatically added to your \texttt{cite.bib} file.

\section{Conclusion}
This template demonstrates basic \LaTeX{} features including mathematical formulas, tables, and lists. The bibliography system is now set up to work seamlessly with Bibby AI's citation management. Feel free to modify it for your needs!

\bibliographystyle{plain}
\bibliography{cite}

\end{document}\documentclass{article}
\usepackage[utf8]{inputenc}
\usepackage{booktabs}
\usepackage{amsmath}
\usepackage{natbib}

\title{My Document}
\author{Your Name}
\date{\today}

\begin{document}

\maketitle

\section{Introduction}
This is a sample document demonstrating various \LaTeX{} features including mathematical equations, tables, and formatting. You can cite papers using the citation search feature, and they will be automatically added to your \texttt{cite.bib} file.

\section{Mathematical Expressions}
Here are some example equations:

\subsection{Quadratic Formula}
The quadratic formula is given by:
\[x = \frac{-b \pm \sqrt{b^2 - 4ac}}{2a}\]

\subsection{Matrix Example}
A sample matrix:
\[
\begin{pmatrix}
a_{11} & a_{12} & a_{13} \\
a_{21} & a_{22} & a_{23} \\
a_{31} & a_{32} & a_{33}
\end{pmatrix}
\]

\section{Tables}
Below is a sample table showing data:

\begin{table}[h]
\centering
\begin{tabular}{@{}lrrr@{}}
\toprule
\textbf{Category} & \textbf{Value 1} & \textbf{Value 2} & \textbf{Total} \\
\midrule
Item A & 10.5 & 20.3 & 30.8 \\
Item B & 15.2 & 25.1 & 40.3 \\
Item C & 22.1 & 12.7 & 34.8 \\
\bottomrule
\end{tabular}
\caption{Sample Data Table}
\label{tab:sample}
\end{table}

\section{Lists}
\subsection{Bullet Points}
\begin{itemize}
  \item First item
  \item Second item
  \item Third item with sub-items:
    \begin{itemize}
      \item Sub-item 1
      \item Sub-item 2
    \end{itemize}
\end{itemize}

\subsection{Numbered List}
\begin{enumerate}
  \item First step
  \item Second step
  \item Third step
\end{enumerate}

\section{Citations and References}
When you use the citation search feature, references will be automatically managed. For example, you might cite a paper like this: \cite{example2024}. All citations will be automatically added to your \texttt{cite.bib} file.

\section{Conclusion}
This template demonstrates basic \LaTeX{} features including mathematical formulas, tables, and lists. The bibliography system is now set up to work seamlessly with Bibby AI's citation management. Feel free to modify it for your needs!

\bibliographystyle{plain}
\bibliography{cite}

\end{document}\documentclass{article}
\usepackage[utf8]{inputenc}
\usepackage{booktabs}
\usepackage{amsmath}
\usepackage{natbib}

\title{My Document}
\author{Your Name}
\date{\today}

\begin{document}

\maketitle

\section{Introduction}
This is a sample document demonstrating various \LaTeX{} features including mathematical equations, tables, and formatting. You can cite papers using the citation search feature, and they will be automatically added to your \texttt{cite.bib} file.

\section{Mathematical Expressions}
Here are some example equations:

\subsection{Quadratic Formula}
The quadratic formula is given by:
\[x = \frac{-b \pm \sqrt{b^2 - 4ac}}{2a}\]

\subsection{Matrix Example}
A sample matrix:
\[
\begin{pmatrix}
a_{11} & a_{12} & a_{13} \\
a_{21} & a_{22} & a_{23} \\
a_{31} & a_{32} & a_{33}
\end{pmatrix}
\]

\section{Tables}
Below is a sample table showing data:

\begin{table}[h]
\centering
\begin{tabular}{@{}lrrr@{}}
\toprule
\textbf{Category} & \textbf{Value 1} & \textbf{Value 2} & \textbf{Total} \\
\midrule
Item A & 10.5 & 20.3 & 30.8 \\
Item B & 15.2 & 25.1 & 40.3 \\
Item C & 22.1 & 12.7 & 34.8 \\
\bottomrule
\end{tabular}
\caption{Sample Data Table}
\label{tab:sample}
\end{table}

\section{Lists}
\subsection{Bullet Points}
\begin{itemize}
  \item First item
  \item Second item
  \item Third item with sub-items:
    \begin{itemize}
      \item Sub-item 1
      \item Sub-item 2
    \end{itemize}
\end{itemize}

\subsection{Numbered List}
\begin{enumerate}
  \item First step
  \item Second step
  \item Third step
\end{enumerate}

\section{Citations and References}
When you use the citation search feature, references will be automatically managed. For example, you might cite a paper like this: \cite{example2024}. All citations will be automatically added to your \texttt{cite.bib} file.

\section{Conclusion}
This template demonstrates basic \LaTeX{} features including mathematical formulas, tables, and lists. The bibliography system is now set up to work seamlessly with Bibby AI's citation management. Feel free to modify it for your needs!

\bibliographystyle{plain}
\bibliography{cite}

\end{document}\documentclass{article}
\usepackage[utf8]{inputenc}
\usepackage{booktabs}
\usepackage{amsmath}
\usepackage{natbib}

\title{My Document}
\author{Your Name}
\date{\today}

\begin{document}

\maketitle

\section{Introduction}
This is a sample document demonstrating various \LaTeX{} features including mathematical equations, tables, and formatting. You can cite papers using the citation search feature, and they will be automatically added to your \texttt{cite.bib} file.

\section{Mathematical Expressions}
Here are some example equations:

\subsection{Quadratic Formula}
The quadratic formula is given by:
\[x = \frac{-b \pm \sqrt{b^2 - 4ac}}{2a}\]

\subsection{Matrix Example}
A sample matrix:
\[
\begin{pmatrix}
a_{11} & a_{12} & a_{13} \\
a_{21} & a_{22} & a_{23} \\
a_{31} & a_{32} & a_{33}
\end{pmatrix}
\]

\section{Tables}
Below is a sample table showing data:

\begin{table}[h]
\centering
\begin{tabular}{@{}lrrr@{}}
\toprule
\textbf{Category} & \textbf{Value 1} & \textbf{Value 2} & \textbf{Total} \\
\midrule
Item A & 10.5 & 20.3 & 30.8 \\
Item B & 15.2 & 25.1 & 40.3 \\
Item C & 22.1 & 12.7 & 34.8 \\
\bottomrule
\end{tabular}
\caption{Sample Data Table}
\label{tab:sample}
\end{table}

\section{Lists}
\subsection{Bullet Points}
\begin{itemize}
  \item First item
  \item Second item
  \item Third item with sub-items:
    \begin{itemize}
      \item Sub-item 1
      \item Sub-item 2
    \end{itemize}
\end{itemize}

\subsection{Numbered List}
\begin{enumerate}
  \item First step
  \item Second step
  \item Third step
\end{enumerate}

\section{Citations and References}
When you use the citation search feature, references will be automatically managed. For example, you might cite a paper like this: \cite{example2024}. All citations will be automatically added to your \texttt{cite.bib} file.

\section{Conclusion}
This template demonstrates basic \LaTeX{} features including mathematical formulas, tables, and lists. The bibliography system is now set up to work seamlessly with Bibby AI's citation management. Feel free to modify it for your needs!

\bibliographystyle{plain}
\bibliography{cite}

\end{document}\documentclass{article}
\usepackage[utf8]{inputenc}
\usepackage{booktabs}
\usepackage{amsmath}
\usepackage{natbib}

\title{My Document}
\author{Your Name}
\date{\today}

\begin{document}

\maketitle

\section{Introduction}
This is a sample document demonstrating various \LaTeX{} features including mathematical equations, tables, and formatting. You can cite papers using the citation search feature, and they will be automatically added to your \texttt{cite.bib} file.

\section{Mathematical Expressions}
Here are some example equations:

\subsection{Quadratic Formula}
The quadratic formula is given by:
\[x = \frac{-b \pm \sqrt{b^2 - 4ac}}{2a}\]

\subsection{Matrix Example}
A sample matrix:
\[
\begin{pmatrix}
a_{11} & a_{12} & a_{13} \\
a_{21} & a_{22} & a_{23} \\
a_{31} & a_{32} & a_{33}
\end{pmatrix}
\]

\section{Tables}
Below is a sample table showing data:

\begin{table}[h]
\centering
\begin{tabular}{@{}lrrr@{}}
\toprule
\textbf{Category} & \textbf{Value 1} & \textbf{Value 2} & \textbf{Total} \\
\midrule
Item A & 10.5 & 20.3 & 30.8 \\
Item B & 15.2 & 25.1 & 40.3 \\
Item C & 22.1 & 12.7 & 34.8 \\
\bottomrule
\end{tabular}
\caption{Sample Data Table}
\label{tab:sample}
\end{table}

\section{Lists}
\subsection{Bullet Points}
\begin{itemize}
  \item First item
  \item Second item
  \item Third item with sub-items:
    \begin{itemize}
      \item Sub-item 1
      \item Sub-item 2
    \end{itemize}
\end{itemize}

\subsection{Numbered List}
\begin{enumerate}
  \item First step
  \item Second step
  \item Third step
\end{enumerate}

\section{Citations and References}
When you use the citation search feature, references will be automatically managed. For example, you might cite a paper like this: \cite{example2024}. All citations will be automatically added to your \texttt{cite.bib} file.

\section{Conclusion}
This template demonstrates basic \LaTeX{} features including mathematical formulas, tables, and lists. The bibliography system is now set up to work seamlessly with Bibby AI's citation management. Feel free to modify it for your needs!

\bibliographystyle{plain}
\bibliography{cite}

\end{document}\documentclass{article}
\usepackage[utf8]{inputenc}
\usepackage{booktabs}
\usepackage{amsmath}
\usepackage{natbib}

\title{My Document}
\author{Your Name}
\date{\today}

\begin{document}

\maketitle

\section{Introduction}
This is a sample document demonstrating various \LaTeX{} features including mathematical equations, tables, and formatting. You can cite papers using the citation search feature, and they will be automatically added to your \texttt{cite.bib} file.

\section{Mathematical Expressions}
Here are some example equations:

\subsection{Quadratic Formula}
The quadratic formula is given by:
\[x = \frac{-b \pm \sqrt{b^2 - 4ac}}{2a}\]

\subsection{Matrix Example}
A sample matrix:
\[
\begin{pmatrix}
a_{11} & a_{12} & a_{13} \\
a_{21} & a_{22} & a_{23} \\
a_{31} & a_{32} & a_{33}
\end{pmatrix}
\]

\section{Tables}
Below is a sample table showing data:

\begin{table}[h]
\centering
\begin{tabular}{@{}lrrr@{}}
\toprule
\textbf{Category} & \textbf{Value 1} & \textbf{Value 2} & \textbf{Total} \\
\midrule
Item A & 10.5 & 20.3 & 30.8 \\
Item B & 15.2 & 25.1 & 40.3 \\
Item C & 22.1 & 12.7 & 34.8 \\
\bottomrule
\end{tabular}
\caption{Sample Data Table}
\label{tab:sample}
\end{table}

\section{Lists}
\subsection{Bullet Points}
\begin{itemize}
  \item First item
  \item Second item
  \item Third item with sub-items:
    \begin{itemize}
      \item Sub-item 1
      \item Sub-item 2
    \end{itemize}
\end{itemize}

\subsection{Numbered List}
\begin{enumerate}
  \item First step
  \item Second step
  \item Third step
\end{enumerate}

\section{Citations and References}
When you use the citation search feature, references will be automatically managed. For example, you might cite a paper like this: \cite{example2024}. All citations will be automatically added to your \texttt{cite.bib} file.

\section{Conclusion}
This template demonstrates basic \LaTeX{} features including mathematical formulas, tables, and lists. The bibliography system is now set up to work seamlessly with Bibby AI's citation management. Feel free to modify it for your needs!

\bibliographystyle{plain}
\bibliography{cite}

\end{document}\documentclass{article}
\usepackage[utf8]{inputenc}
\usepackage{booktabs}
\usepackage{amsmath}
\usepackage{natbib}

\title{My Document}
\author{Your Name}
\date{\today}

\begin{document}

\maketitle

\section{Introduction}
This is a sample document demonstrating various \LaTeX{} features including mathematical equations, tables, and formatting. You can cite papers using the citation search feature, and they will be automatically added to your \texttt{cite.bib} file.

\section{Mathematical Expressions}
Here are some example equations:

\subsection{Quadratic Formula}
The quadratic formula is given by:
\[x = \frac{-b \pm \sqrt{b^2 - 4ac}}{2a}\]

\subsection{Matrix Example}
A sample matrix:
\[
\begin{pmatrix}
a_{11} & a_{12} & a_{13} \\
a_{21} & a_{22} & a_{23} \\
a_{31} & a_{32} & a_{33}
\end{pmatrix}
\]

\section{Tables}
Below is a sample table showing data:

\begin{table}[h]
\centering
\begin{tabular}{@{}lrrr@{}}
\toprule
\textbf{Category} & \textbf{Value 1} & \textbf{Value 2} & \textbf{Total} \\
\midrule
Item A & 10.5 & 20.3 & 30.8 \\
Item B & 15.2 & 25.1 & 40.3 \\
Item C & 22.1 & 12.7 & 34.8 \\
\bottomrule
\end{tabular}
\caption{Sample Data Table}
\label{tab:sample}
\end{table}

\section{Lists}
\subsection{Bullet Points}
\begin{itemize}
  \item First item
  \item Second item
  \item Third item with sub-items:
    \begin{itemize}
      \item Sub-item 1
      \item Sub-item 2
    \end{itemize}
\end{itemize}

\subsection{Numbered List}
\begin{enumerate}
  \item First step
  \item Second step
  \item Third step
\end{enumerate}

\section{Citations and References}
When you use the citation search feature, references will be automatically managed. For example, you might cite a paper like this: \cite{example2024}. All citations will be automatically added to your \texttt{cite.bib} file.

\section{Conclusion}
This template demonstrates basic \LaTeX{} features including mathematical formulas, tables, and lists. The bibliography system is now set up to work seamlessly with Bibby AI's citation management. Feel free to modify it for your needs!

\bibliographystyle{plain}
\bibliography{cite}

\end{document}\documentclass{article}
\usepackage[utf8]{inputenc}
\usepackage{booktabs}
\usepackage{amsmath}
\usepackage{natbib}

\title{My Document}
\author{Your Name}
\date{\today}

\begin{document}

\maketitle

\section{Introduction}
This is a sample document demonstrating various \LaTeX{} features including mathematical equations, tables, and formatting. You can cite papers using the citation search feature, and they will be automatically added to your \texttt{cite.bib} file.

\section{Mathematical Expressions}
Here are some example equations:

\subsection{Quadratic Formula}
The quadratic formula is given by:
\[x = \frac{-b \pm \sqrt{b^2 - 4ac}}{2a}\]

\subsection{Matrix Example}
A sample matrix:
\[
\begin{pmatrix}
a_{11} & a_{12} & a_{13} \\
a_{21} & a_{22} & a_{23} \\
a_{31} & a_{32} & a_{33}
\end{pmatrix}
\]

\section{Tables}
Below is a sample table showing data:

\begin{table}[h]
\centering
\begin{tabular}{@{}lrrr@{}}
\toprule
\textbf{Category} & \textbf{Value 1} & \textbf{Value 2} & \textbf{Total} \\
\midrule
Item A & 10.5 & 20.3 & 30.8 \\
Item B & 15.2 & 25.1 & 40.3 \\
Item C & 22.1 & 12.7 & 34.8 \\
\bottomrule
\end{tabular}
\caption{Sample Data Table}
\label{tab:sample}
\end{table}

\section{Lists}
\subsection{Bullet Points}
\begin{itemize}
  \item First item
  \item Second item
  \item Third item with sub-items:
    \begin{itemize}
      \item Sub-item 1
      \item Sub-item 2
    \end{itemize}
\end{itemize}

\subsection{Numbered List}
\begin{enumerate}
  \item First step
  \item Second step
  \item Third step
\end{enumerate}

\section{Citations and References}
When you use the citation search feature, references will be automatically managed. For example, you might cite a paper like this: \cite{example2024}. All citations will be automatically added to your \texttt{cite.bib} file.

\section{Conclusion}
This template demonstrates basic \LaTeX{} features including mathematical formulas, tables, and lists. The bibliography system is now set up to work seamlessly with Bibby AI's citation management. Feel free to modify it for your needs!

\bibliographystyle{plain}
\bibliography{cite}

\end{document}\documentclass{article}
\usepackage[utf8]{inputenc}
\usepackage{booktabs}
\usepackage{amsmath}
\usepackage{natbib}

\title{My Document}
\author{Your Name}
\date{\today}

\begin{document}

\maketitle

\section{Introduction}
This is a sample document demonstrating various \LaTeX{} features including mathematical equations, tables, and formatting. You can cite papers using the citation search feature, and they will be automatically added to your \texttt{cite.bib} file.

\section{Mathematical Expressions}
Here are some example equations:

\subsection{Quadratic Formula}
The quadratic formula is given by:
\[x = \frac{-b \pm \sqrt{b^2 - 4ac}}{2a}\]

\subsection{Matrix Example}
A sample matrix:
\[
\begin{pmatrix}
a_{11} & a_{12} & a_{13} \\
a_{21} & a_{22} & a_{23} \\
a_{31} & a_{32} & a_{33}
\end{pmatrix}
\]

\section{Tables}
Below is a sample table showing data:

\begin{table}[h]
\centering
\begin{tabular}{@{}lrrr@{}}
\toprule
\textbf{Category} & \textbf{Value 1} & \textbf{Value 2} & \textbf{Total} \\
\midrule
Item A & 10.5 & 20.3 & 30.8 \\
Item B & 15.2 & 25.1 & 40.3 \\
Item C & 22.1 & 12.7 & 34.8 \\
\bottomrule
\end{tabular}
\caption{Sample Data Table}
\label{tab:sample}
\end{table}

\section{Lists}
\subsection{Bullet Points}
\begin{itemize}
  \item First item
  \item Second item
  \item Third item with sub-items:
    \begin{itemize}
      \item Sub-item 1
      \item Sub-item 2
    \end{itemize}
\end{itemize}

\subsection{Numbered List}
\begin{enumerate}
  \item First step
  \item Second step
  \item Third step
\end{enumerate}

\section{Citations and References}
When you use the citation search feature, references will be automatically managed. For example, you might cite a paper like this: \cite{example2024}. All citations will be automatically added to your \texttt{cite.bib} file.

\section{Conclusion}
This template demonstrates basic \LaTeX{} features including mathematical formulas, tables, and lists. The bibliography system is now set up to work seamlessly with Bibby AI's citation management. Feel free to modify it for your needs!

\bibliographystyle{plain}
\bibliography{cite}

\end{document}\documentclass{article}
\usepackage[utf8]{inputenc}
\usepackage{booktabs}
\usepackage{amsmath}
\usepackage{natbib}

\title{My Document}
\author{Your Name}
\date{\today}

\begin{document}

\maketitle

\section{Introduction}
This is a sample document demonstrating various \LaTeX{} features including mathematical equations, tables, and formatting. You can cite papers using the citation search feature, and they will be automatically added to your \texttt{cite.bib} file.

\section{Mathematical Expressions}
Here are some example equations:

\subsection{Quadratic Formula}
The quadratic formula is given by:
\[x = \frac{-b \pm \sqrt{b^2 - 4ac}}{2a}\]

\subsection{Matrix Example}
A sample matrix:
\[
\begin{pmatrix}
a_{11} & a_{12} & a_{13} \\
a_{21} & a_{22} & a_{23} \\
a_{31} & a_{32} & a_{33}
\end{pmatrix}
\]

\section{Tables}
Below is a sample table showing data:

\begin{table}[h]
\centering
\begin{tabular}{@{}lrrr@{}}
\toprule
\textbf{Category} & \textbf{Value 1} & \textbf{Value 2} & \textbf{Total} \\
\midrule
Item A & 10.5 & 20.3 & 30.8 \\
Item B & 15.2 & 25.1 & 40.3 \\
Item C & 22.1 & 12.7 & 34.8 \\
\bottomrule
\end{tabular}
\caption{Sample Data Table}
\label{tab:sample}
\end{table}

\section{Lists}
\subsection{Bullet Points}
\begin{itemize}
  \item First item
  \item Second item
  \item Third item with sub-items:
    \begin{itemize}
      \item Sub-item 1
      \item Sub-item 2
    \end{itemize}
\end{itemize}

\subsection{Numbered List}
\begin{enumerate}
  \item First step
  \item Second step
  \item Third step
\end{enumerate}

\section{Citations and References}
When you use the citation search feature, references will be automatically managed. For example, you might cite a paper like this: \cite{example2024}. All citations will be automatically added to your \texttt{cite.bib} file.

\section{Conclusion}
This template demonstrates basic \LaTeX{} features including mathematical formulas, tables, and lists. The bibliography system is now set up to work seamlessly with Bibby AI's citation management. Feel free to modify it for your needs!

\bibliographystyle{plain}
\bibliography{cite}

\end{document}\documentclass{article}
\usepackage[utf8]{inputenc}
\usepackage{booktabs}
\usepackage{amsmath}
\usepackage{natbib}

\title{My Document}
\author{Your Name}
\date{\today}

\begin{document}

\maketitle

\section{Introduction}
This is a sample document demonstrating various \LaTeX{} features including mathematical equations, tables, and formatting. You can cite papers using the citation search feature, and they will be automatically added to your \texttt{cite.bib} file.

\section{Mathematical Expressions}
Here are some example equations:

\subsection{Quadratic Formula}
The quadratic formula is given by:
\[x = \frac{-b \pm \sqrt{b^2 - 4ac}}{2a}\]

\subsection{Matrix Example}
A sample matrix:
\[
\begin{pmatrix}
a_{11} & a_{12} & a_{13} \\
a_{21} & a_{22} & a_{23} \\
a_{31} & a_{32} & a_{33}
\end{pmatrix}
\]

\section{Tables}
Below is a sample table showing data:

\begin{table}[h]
\centering
\begin{tabular}{@{}lrrr@{}}
\toprule
\textbf{Category} & \textbf{Value 1} & \textbf{Value 2} & \textbf{Total} \\
\midrule
Item A & 10.5 & 20.3 & 30.8 \\
Item B & 15.2 & 25.1 & 40.3 \\
Item C & 22.1 & 12.7 & 34.8 \\
\bottomrule
\end{tabular}
\caption{Sample Data Table}
\label{tab:sample}
\end{table}

\section{Lists}
\subsection{Bullet Points}
\begin{itemize}
  \item First item
  \item Second item
  \item Third item with sub-items:
    \begin{itemize}
      \item Sub-item 1
      \item Sub-item 2
    \end{itemize}
\end{itemize}

\subsection{Numbered List}
\begin{enumerate}
  \item First step
  \item Second step
  \item Third step
\end{enumerate}

\section{Citations and References}
When you use the citation search feature, references will be automatically managed. For example, you might cite a paper like this: \cite{example2024}. All citations will be automatically added to your \texttt{cite.bib} file.

\section{Conclusion}
This template demonstrates basic \LaTeX{} features including mathematical formulas, tables, and lists. The bibliography system is now set up to work seamlessly with Bibby AI's citation management. Feel free to modify it for your needs!

\bibliographystyle{plain}
\bibliography{cite}

\end{document}\documentclass{article}
\usepackage[utf8]{inputenc}
\usepackage{booktabs}
\usepackage{amsmath}
\usepackage{natbib}

\title{My Document}
\author{Your Name}
\date{\today}

\begin{document}

\maketitle

\section{Introduction}
This is a sample document demonstrating various \LaTeX{} features including mathematical equations, tables, and formatting. You can cite papers using the citation search feature, and they will be automatically added to your \texttt{cite.bib} file.

\section{Mathematical Expressions}
Here are some example equations:

\subsection{Quadratic Formula}
The quadratic formula is given by:
\[x = \frac{-b \pm \sqrt{b^2 - 4ac}}{2a}\]

\subsection{Matrix Example}
A sample matrix:
\[
\begin{pmatrix}
a_{11} & a_{12} & a_{13} \\
a_{21} & a_{22} & a_{23} \\
a_{31} & a_{32} & a_{33}
\end{pmatrix}
\]

\section{Tables}
Below is a sample table showing data:

\begin{table}[h]
\centering
\begin{tabular}{@{}lrrr@{}}
\toprule
\textbf{Category} & \textbf{Value 1} & \textbf{Value 2} & \textbf{Total} \\
\midrule
Item A & 10.5 & 20.3 & 30.8 \\
Item B & 15.2 & 25.1 & 40.3 \\
Item C & 22.1 & 12.7 & 34.8 \\
\bottomrule
\end{tabular}
\caption{Sample Data Table}
\label{tab:sample}
\end{table}

\section{Lists}
\subsection{Bullet Points}
\begin{itemize}
  \item First item
  \item Second item
  \item Third item with sub-items:
    \begin{itemize}
      \item Sub-item 1
      \item Sub-item 2
    \end{itemize}
\end{itemize}

\subsection{Numbered List}
\begin{enumerate}
  \item First step
  \item Second step
  \item Third step
\end{enumerate}

\section{Citations and References}
When you use the citation search feature, references will be automatically managed. For example, you might cite a paper like this: \cite{example2024}. All citations will be automatically added to your \texttt{cite.bib} file.

\section{Conclusion}
This template demonstrates basic \LaTeX{} features including mathematical formulas, tables, and lists. The bibliography system is now set up to work seamlessly with Bibby AI's citation management. Feel free to modify it for your needs!

\bibliographystyle{plain}
\bibliography{cite}

\end{document}\documentclass{article}
\usepackage[utf8]{inputenc}
\usepackage{booktabs}
\usepackage{amsmath}
\usepackage{natbib}

\title{My Document}
\author{Your Name}
\date{\today}

\begin{document}

\maketitle

\section{Introduction}
This is a sample document demonstrating various \LaTeX{} features including mathematical equations, tables, and formatting. You can cite papers using the citation search feature, and they will be automatically added to your \texttt{cite.bib} file.

\section{Mathematical Expressions}
Here are some example equations:

\subsection{Quadratic Formula}
The quadratic formula is given by:
\[x = \frac{-b \pm \sqrt{b^2 - 4ac}}{2a}\]

\subsection{Matrix Example}
A sample matrix:
\[
\begin{pmatrix}
a_{11} & a_{12} & a_{13} \\
a_{21} & a_{22} & a_{23} \\
a_{31} & a_{32} & a_{33}
\end{pmatrix}
\]

\section{Tables}
Below is a sample table showing data:

\begin{table}[h]
\centering
\begin{tabular}{@{}lrrr@{}}
\toprule
\textbf{Category} & \textbf{Value 1} & \textbf{Value 2} & \textbf{Total} \\
\midrule
Item A & 10.5 & 20.3 & 30.8 \\
Item B & 15.2 & 25.1 & 40.3 \\
Item C & 22.1 & 12.7 & 34.8 \\
\bottomrule
\end{tabular}
\caption{Sample Data Table}
\label{tab:sample}
\end{table}

\section{Lists}
\subsection{Bullet Points}
\begin{itemize}
  \item First item
  \item Second item
  \item Third item with sub-items:
    \begin{itemize}
      \item Sub-item 1
      \item Sub-item 2
    \end{itemize}
\end{itemize}

\subsection{Numbered List}
\begin{enumerate}
  \item First step
  \item Second step
  \item Third step
\end{enumerate}

\section{Citations and References}
When you use the citation search feature, references will be automatically managed. For example, you might cite a paper like this: \cite{example2024}. All citations will be automatically added to your \texttt{cite.bib} file.

\section{Conclusion}
This template demonstrates basic \LaTeX{} features including mathematical formulas, tables, and lists. The bibliography system is now set up to work seamlessly with Bibby AI's citation management. Feel free to modify it for your needs!

\bibliographystyle{plain}
\bibliography{cite}

\end{document}\documentclass{article}
\usepackage[utf8]{inputenc}
\usepackage{booktabs}
\usepackage{amsmath}
\usepackage{natbib}

\title{My Document}
\author{Your Name}
\date{\today}

\begin{document}

\maketitle

\section{Introduction}
This is a sample document demonstrating various \LaTeX{} features including mathematical equations, tables, and formatting. You can cite papers using the citation search feature, and they will be automatically added to your \texttt{cite.bib} file.

\section{Mathematical Expressions}
Here are some example equations:

\subsection{Quadratic Formula}
The quadratic formula is given by:
\[x = \frac{-b \pm \sqrt{b^2 - 4ac}}{2a}\]

\subsection{Matrix Example}
A sample matrix:
\[
\begin{pmatrix}
a_{11} & a_{12} & a_{13} \\
a_{21} & a_{22} & a_{23} \\
a_{31} & a_{32} & a_{33}
\end{pmatrix}
\]

\section{Tables}
Below is a sample table showing data:

\begin{table}[h]
\centering
\begin{tabular}{@{}lrrr@{}}
\toprule
\textbf{Category} & \textbf{Value 1} & \textbf{Value 2} & \textbf{Total} \\
\midrule
Item A & 10.5 & 20.3 & 30.8 \\
Item B & 15.2 & 25.1 & 40.3 \\
Item C & 22.1 & 12.7 & 34.8 \\
\bottomrule
\end{tabular}
\caption{Sample Data Table}
\label{tab:sample}
\end{table}

\section{Lists}
\subsection{Bullet Points}
\begin{itemize}
  \item First item
  \item Second item
  \item Third item with sub-items:
    \begin{itemize}
      \item Sub-item 1
      \item Sub-item 2
    \end{itemize}
\end{itemize}

\subsection{Numbered List}
\begin{enumerate}
  \item First step
  \item Second step
  \item Third step
\end{enumerate}

\section{Citations and References}
When you use the citation search feature, references will be automatically managed. For example, you might cite a paper like this: \cite{example2024}. All citations will be automatically added to your \texttt{cite.bib} file.

\section{Conclusion}
This template demonstrates basic \LaTeX{} features including mathematical formulas, tables, and lists. The bibliography system is now set up to work seamlessly with Bibby AI's citation management. Feel free to modify it for your needs!

\bibliographystyle{plain}
\bibliography{cite}

\end{document}\documentclass{article}
\usepackage[utf8]{inputenc}
\usepackage{booktabs}
\usepackage{amsmath}
\usepackage{natbib}

\title{My Document}
\author{Your Name}
\date{\today}

\begin{document}

\maketitle

\section{Introduction}
This is a sample document demonstrating various \LaTeX{} features including mathematical equations, tables, and formatting. You can cite papers using the citation search feature, and they will be automatically added to your \texttt{cite.bib} file.

\section{Mathematical Expressions}
Here are some example equations:

\subsection{Quadratic Formula}
The quadratic formula is given by:
\[x = \frac{-b \pm \sqrt{b^2 - 4ac}}{2a}\]

\subsection{Matrix Example}
A sample matrix:
\[
\begin{pmatrix}
a_{11} & a_{12} & a_{13} \\
a_{21} & a_{22} & a_{23} \\
a_{31} & a_{32} & a_{33}
\end{pmatrix}
\]

\section{Tables}
Below is a sample table showing data:

\begin{table}[h]
\centering
\begin{tabular}{@{}lrrr@{}}
\toprule
\textbf{Category} & \textbf{Value 1} & \textbf{Value 2} & \textbf{Total} \\
\midrule
Item A & 10.5 & 20.3 & 30.8 \\
Item B & 15.2 & 25.1 & 40.3 \\
Item C & 22.1 & 12.7 & 34.8 \\
\bottomrule
\end{tabular}
\caption{Sample Data Table}
\label{tab:sample}
\end{table}

\section{Lists}
\subsection{Bullet Points}
\begin{itemize}
  \item First item
  \item Second item
  \item Third item with sub-items:
    \begin{itemize}
      \item Sub-item 1
      \item Sub-item 2
    \end{itemize}
\end{itemize}

\subsection{Numbered List}
\begin{enumerate}
  \item First step
  \item Second step
  \item Third step
\end{enumerate}

\section{Citations and References}
When you use the citation search feature, references will be automatically managed. For example, you might cite a paper like this: \cite{example2024}. All citations will be automatically added to your \texttt{cite.bib} file.

\section{Conclusion}
This template demonstrates basic \LaTeX{} features including mathematical formulas, tables, and lists. The bibliography system is now set up to work seamlessly with Bibby AI's citation management. Feel free to modify it for your needs!

\bibliographystyle{plain}
\bibliography{cite}

\end{document}\documentclass{article}
\usepackage[utf8]{inputenc}
\usepackage{booktabs}
\usepackage{amsmath}
\usepackage{natbib}

\title{My Document}
\author{Your Name}
\date{\today}

\begin{document}

\maketitle

\section{Introduction}
This is a sample document demonstrating various \LaTeX{} features including mathematical equations, tables, and formatting. You can cite papers using the citation search feature, and they will be automatically added to your \texttt{cite.bib} file.

\section{Mathematical Expressions}
Here are some example equations:

\subsection{Quadratic Formula}
The quadratic formula is given by:
\[x = \frac{-b \pm \sqrt{b^2 - 4ac}}{2a}\]

\subsection{Matrix Example}
A sample matrix:
\[
\begin{pmatrix}
a_{11} & a_{12} & a_{13} \\
a_{21} & a_{22} & a_{23} \\
a_{31} & a_{32} & a_{33}
\end{pmatrix}
\]

\section{Tables}
Below is a sample table showing data:

\begin{table}[h]
\centering
\begin{tabular}{@{}lrrr@{}}
\toprule
\textbf{Category} & \textbf{Value 1} & \textbf{Value 2} & \textbf{Total} \\
\midrule
Item A & 10.5 & 20.3 & 30.8 \\
Item B & 15.2 & 25.1 & 40.3 \\
Item C & 22.1 & 12.7 & 34.8 \\
\bottomrule
\end{tabular}
\caption{Sample Data Table}
\label{tab:sample}
\end{table}

\section{Lists}
\subsection{Bullet Points}
\begin{itemize}
  \item First item
  \item Second item
  \item Third item with sub-items:
    \begin{itemize}
      \item Sub-item 1
      \item Sub-item 2
    \end{itemize}
\end{itemize}

\subsection{Numbered List}
\begin{enumerate}
  \item First step
  \item Second step
  \item Third step
\end{enumerate}

\section{Citations and References}
When you use the citation search feature, references will be automatically managed. For example, you might cite a paper like this: \cite{example2024}. All citations will be automatically added to your \texttt{cite.bib} file.

\section{Conclusion}
This template demonstrates basic \LaTeX{} features including mathematical formulas, tables, and lists. The bibliography system is now set up to work seamlessly with Bibby AI's citation management. Feel free to modify it for your needs!

\bibliographystyle{plain}
\bibliography{cite}

\end{document}\documentclass{article}
\usepackage[utf8]{inputenc}
\usepackage{booktabs}
\usepackage{amsmath}
\usepackage{natbib}

\title{My Document}
\author{Your Name}
\date{\today}

\begin{document}

\maketitle

\section{Introduction}
This is a sample document demonstrating various \LaTeX{} features including mathematical equations, tables, and formatting. You can cite papers using the citation search feature, and they will be automatically added to your \texttt{cite.bib} file.

\section{Mathematical Expressions}
Here are some example equations:

\subsection{Quadratic Formula}
The quadratic formula is given by:
\[x = \frac{-b \pm \sqrt{b^2 - 4ac}}{2a}\]

\subsection{Matrix Example}
A sample matrix:
\[
\begin{pmatrix}
a_{11} & a_{12} & a_{13} \\
a_{21} & a_{22} & a_{23} \\
a_{31} & a_{32} & a_{33}
\end{pmatrix}
\]

\section{Tables}
Below is a sample table showing data:

\begin{table}[h]
\centering
\begin{tabular}{@{}lrrr@{}}
\toprule
\textbf{Category} & \textbf{Value 1} & \textbf{Value 2} & \textbf{Total} \\
\midrule
Item A & 10.5 & 20.3 & 30.8 \\
Item B & 15.2 & 25.1 & 40.3 \\
Item C & 22.1 & 12.7 & 34.8 \\
\bottomrule
\end{tabular}
\caption{Sample Data Table}
\label{tab:sample}
\end{table}

\section{Lists}
\subsection{Bullet Points}
\begin{itemize}
  \item First item
  \item Second item
  \item Third item with sub-items:
    \begin{itemize}
      \item Sub-item 1
      \item Sub-item 2
    \end{itemize}
\end{itemize}

\subsection{Numbered List}
\begin{enumerate}
  \item First step
  \item Second step
  \item Third step
\end{enumerate}

\section{Citations and References}
When you use the citation search feature, references will be automatically managed. For example, you might cite a paper like this: \cite{example2024}. All citations will be automatically added to your \texttt{cite.bib} file.

\section{Conclusion}
This template demonstrates basic \LaTeX{} features including mathematical formulas, tables, and lists. The bibliography system is now set up to work seamlessly with Bibby AI's citation management. Feel free to modify it for your needs!

\bibliographystyle{plain}
\bibliography{cite}

\end{document}\documentclass{article}
\usepackage[utf8]{inputenc}
\usepackage{booktabs}
\usepackage{amsmath}
\usepackage{natbib}

\title{My Document}
\author{Your Name}
\date{\today}

\begin{document}

\maketitle

\section{Introduction}
This is a sample document demonstrating various \LaTeX{} features including mathematical equations, tables, and formatting. You can cite papers using the citation search feature, and they will be automatically added to your \texttt{cite.bib} file.

\section{Mathematical Expressions}
Here are some example equations:

\subsection{Quadratic Formula}
The quadratic formula is given by:
\[x = \frac{-b \pm \sqrt{b^2 - 4ac}}{2a}\]

\subsection{Matrix Example}
A sample matrix:
\[
\begin{pmatrix}
a_{11} & a_{12} & a_{13} \\
a_{21} & a_{22} & a_{23} \\
a_{31} & a_{32} & a_{33}
\end{pmatrix}
\]

\section{Tables}
Below is a sample table showing data:

\begin{table}[h]
\centering
\begin{tabular}{@{}lrrr@{}}
\toprule
\textbf{Category} & \textbf{Value 1} & \textbf{Value 2} & \textbf{Total} \\
\midrule
Item A & 10.5 & 20.3 & 30.8 \\
Item B & 15.2 & 25.1 & 40.3 \\
Item C & 22.1 & 12.7 & 34.8 \\
\bottomrule
\end{tabular}
\caption{Sample Data Table}
\label{tab:sample}
\end{table}

\section{Lists}
\subsection{Bullet Points}
\begin{itemize}
  \item First item
  \item Second item
  \item Third item with sub-items:
    \begin{itemize}
      \item Sub-item 1
      \item Sub-item 2
    \end{itemize}
\end{itemize}

\subsection{Numbered List}
\begin{enumerate}
  \item First step
  \item Second step
  \item Third step
\end{enumerate}

\section{Citations and References}
When you use the citation search feature, references will be automatically managed. For example, you might cite a paper like this: \cite{example2024}. All citations will be automatically added to your \texttt{cite.bib} file.

\section{Conclusion}
This template demonstrates basic \LaTeX{} features including mathematical formulas, tables, and lists. The bibliography system is now set up to work seamlessly with Bibby AI's citation management. Feel free to modify it for your needs!

\bibliographystyle{plain}
\bibliography{cite}

\end{document}\documentclass{article}
\usepackage[utf8]{inputenc}
\usepackage{booktabs}
\usepackage{amsmath}
\usepackage{natbib}

\title{My Document}
\author{Your Name}
\date{\today}

\begin{document}

\maketitle

\section{Introduction}
This is a sample document demonstrating various \LaTeX{} features including mathematical equations, tables, and formatting. You can cite papers using the citation search feature, and they will be automatically added to your \texttt{cite.bib} file.

\section{Mathematical Expressions}
Here are some example equations:

\subsection{Quadratic Formula}
The quadratic formula is given by:
\[x = \frac{-b \pm \sqrt{b^2 - 4ac}}{2a}\]

\subsection{Matrix Example}
A sample matrix:
\[
\begin{pmatrix}
a_{11} & a_{12} & a_{13} \\
a_{21} & a_{22} & a_{23} \\
a_{31} & a_{32} & a_{33}
\end{pmatrix}
\]

\section{Tables}
Below is a sample table showing data:

\begin{table}[h]
\centering
\begin{tabular}{@{}lrrr@{}}
\toprule
\textbf{Category} & \textbf{Value 1} & \textbf{Value 2} & \textbf{Total} \\
\midrule
Item A & 10.5 & 20.3 & 30.8 \\
Item B & 15.2 & 25.1 & 40.3 \\
Item C & 22.1 & 12.7 & 34.8 \\
\bottomrule
\end{tabular}
\caption{Sample Data Table}
\label{tab:sample}
\end{table}

\section{Lists}
\subsection{Bullet Points}
\begin{itemize}
  \item First item
  \item Second item
  \item Third item with sub-items:
    \begin{itemize}
      \item Sub-item 1
      \item Sub-item 2
    \end{itemize}
\end{itemize}

\subsection{Numbered List}
\begin{enumerate}
  \item First step
  \item Second step
  \item Third step
\end{enumerate}

\section{Citations and References}
When you use the citation search feature, references will be automatically managed. For example, you might cite a paper like this: \cite{example2024}. All citations will be automatically added to your \texttt{cite.bib} file.

\section{Conclusion}
This template demonstrates basic \LaTeX{} features including mathematical formulas, tables, and lists. The bibliography system is now set up to work seamlessly with Bibby AI's citation management. Feel free to modify it for your needs!

\bibliographystyle{plain}
\bibliography{cite}

\end{document}\documentclass{article}
\usepackage[utf8]{inputenc}
\usepackage{booktabs}
\usepackage{amsmath}
\usepackage{natbib}

\title{My Document}
\author{Your Name}
\date{\today}

\begin{document}

\maketitle

\section{Introduction}
This is a sample document demonstrating various \LaTeX{} features including mathematical equations, tables, and formatting. You can cite papers using the citation search feature, and they will be automatically added to your \texttt{cite.bib} file.

\section{Mathematical Expressions}
Here are some example equations:

\subsection{Quadratic Formula}
The quadratic formula is given by:
\[x = \frac{-b \pm \sqrt{b^2 - 4ac}}{2a}\]

\subsection{Matrix Example}
A sample matrix:
\[
\begin{pmatrix}
a_{11} & a_{12} & a_{13} \\
a_{21} & a_{22} & a_{23} \\
a_{31} & a_{32} & a_{33}
\end{pmatrix}
\]

\section{Tables}
Below is a sample table showing data:

\begin{table}[h]
\centering
\begin{tabular}{@{}lrrr@{}}
\toprule
\textbf{Category} & \textbf{Value 1} & \textbf{Value 2} & \textbf{Total} \\
\midrule
Item A & 10.5 & 20.3 & 30.8 \\
Item B & 15.2 & 25.1 & 40.3 \\
Item C & 22.1 & 12.7 & 34.8 \\
\bottomrule
\end{tabular}
\caption{Sample Data Table}
\label{tab:sample}
\end{table}

\section{Lists}
\subsection{Bullet Points}
\begin{itemize}
  \item First item
  \item Second item
  \item Third item with sub-items:
    \begin{itemize}
      \item Sub-item 1
      \item Sub-item 2
    \end{itemize}
\end{itemize}

\subsection{Numbered List}
\begin{enumerate}
  \item First step
  \item Second step
  \item Third step
\end{enumerate}

\section{Citations and References}
When you use the citation search feature, references will be automatically managed. For example, you might cite a paper like this: \cite{example2024}. All citations will be automatically added to your \texttt{cite.bib} file.

\section{Conclusion}
This template demonstrates basic \LaTeX{} features including mathematical formulas, tables, and lists. The bibliography system is now set up to work seamlessly with Bibby AI's citation management. Feel free to modify it for your needs!

\bibliographystyle{plain}
\bibliography{cite}

\end{document}\documentclass{article}
\usepackage[utf8]{inputenc}
\usepackage{booktabs}
\usepackage{amsmath}
\usepackage{natbib}

\title{My Document}
\author{Your Name}
\date{\today}

\begin{document}

\maketitle

\section{Introduction}
This is a sample document demonstrating various \LaTeX{} features including mathematical equations, tables, and formatting. You can cite papers using the citation search feature, and they will be automatically added to your \texttt{cite.bib} file.

\section{Mathematical Expressions}
Here are some example equations:

\subsection{Quadratic Formula}
The quadratic formula is given by:
\[x = \frac{-b \pm \sqrt{b^2 - 4ac}}{2a}\]

\subsection{Matrix Example}
A sample matrix:
\[
\begin{pmatrix}
a_{11} & a_{12} & a_{13} \\
a_{21} & a_{22} & a_{23} \\
a_{31} & a_{32} & a_{33}
\end{pmatrix}
\]

\section{Tables}
Below is a sample table showing data:

\begin{table}[h]
\centering
\begin{tabular}{@{}lrrr@{}}
\toprule
\textbf{Category} & \textbf{Value 1} & \textbf{Value 2} & \textbf{Total} \\
\midrule
Item A & 10.5 & 20.3 & 30.8 \\
Item B & 15.2 & 25.1 & 40.3 \\
Item C & 22.1 & 12.7 & 34.8 \\
\bottomrule
\end{tabular}
\caption{Sample Data Table}
\label{tab:sample}
\end{table}

\section{Lists}
\subsection{Bullet Points}
\begin{itemize}
  \item First item
  \item Second item
  \item Third item with sub-items:
    \begin{itemize}
      \item Sub-item 1
      \item Sub-item 2
    \end{itemize}
\end{itemize}

\subsection{Numbered List}
\begin{enumerate}
  \item First step
  \item Second step
  \item Third step
\end{enumerate}

\section{Citations and References}
When you use the citation search feature, references will be automatically managed. For example, you might cite a paper like this: \cite{example2024}. All citations will be automatically added to your \texttt{cite.bib} file.

\section{Conclusion}
This template demonstrates basic \LaTeX{} features including mathematical formulas, tables, and lists. The bibliography system is now set up to work seamlessly with Bibby AI's citation management. Feel free to modify it for your needs!

\bibliographystyle{plain}
\bibliography{cite}

\end{document}\documentclass{article}
\usepackage[utf8]{inputenc}
\usepackage{booktabs}
\usepackage{amsmath}
\usepackage{natbib}

\title{My Document}
\author{Your Name}
\date{\today}

\begin{document}

\maketitle

\section{Introduction}
This is a sample document demonstrating various \LaTeX{} features including mathematical equations, tables, and formatting. You can cite papers using the citation search feature, and they will be automatically added to your \texttt{cite.bib} file.

\section{Mathematical Expressions}
Here are some example equations:

\subsection{Quadratic Formula}
The quadratic formula is given by:
\[x = \frac{-b \pm \sqrt{b^2 - 4ac}}{2a}\]

\subsection{Matrix Example}
A sample matrix:
\[
\begin{pmatrix}
a_{11} & a_{12} & a_{13} \\
a_{21} & a_{22} & a_{23} \\
a_{31} & a_{32} & a_{33}
\end{pmatrix}
\]

\section{Tables}
Below is a sample table showing data:

\begin{table}[h]
\centering
\begin{tabular}{@{}lrrr@{}}
\toprule
\textbf{Category} & \textbf{Value 1} & \textbf{Value 2} & \textbf{Total} \\
\midrule
Item A & 10.5 & 20.3 & 30.8 \\
Item B & 15.2 & 25.1 & 40.3 \\
Item C & 22.1 & 12.7 & 34.8 \\
\bottomrule
\end{tabular}
\caption{Sample Data Table}
\label{tab:sample}
\end{table}

\section{Lists}
\subsection{Bullet Points}
\begin{itemize}
  \item First item
  \item Second item
  \item Third item with sub-items:
    \begin{itemize}
      \item Sub-item 1
      \item Sub-item 2
    \end{itemize}
\end{itemize}

\subsection{Numbered List}
\begin{enumerate}
  \item First step
  \item Second step
  \item Third step
\end{enumerate}

\section{Citations and References}
When you use the citation search feature, references will be automatically managed. For example, you might cite a paper like this: \cite{example2024}. All citations will be automatically added to your \texttt{cite.bib} file.

\section{Conclusion}
This template demonstrates basic \LaTeX{} features including mathematical formulas, tables, and lists. The bibliography system is now set up to work seamlessly with Bibby AI's citation management. Feel free to modify it for your needs!

\bibliographystyle{plain}
\bibliography{cite}

\end{document}\documentclass{article}
\usepackage[utf8]{inputenc}
\usepackage{booktabs}
\usepackage{amsmath}
\usepackage{natbib}

\title{My Document}
\author{Your Name}
\date{\today}

\begin{document}

\maketitle

\section{Introduction}
This is a sample document demonstrating various \LaTeX{} features including mathematical equations, tables, and formatting. You can cite papers using the citation search feature, and they will be automatically added to your \texttt{cite.bib} file.

\section{Mathematical Expressions}
Here are some example equations:

\subsection{Quadratic Formula}
The quadratic formula is given by:
\[x = \frac{-b \pm \sqrt{b^2 - 4ac}}{2a}\]

\subsection{Matrix Example}
A sample matrix:
\[
\begin{pmatrix}
a_{11} & a_{12} & a_{13} \\
a_{21} & a_{22} & a_{23} \\
a_{31} & a_{32} & a_{33}
\end{pmatrix}
\]

\section{Tables}
Below is a sample table showing data:

\begin{table}[h]
\centering
\begin{tabular}{@{}lrrr@{}}
\toprule
\textbf{Category} & \textbf{Value 1} & \textbf{Value 2} & \textbf{Total} \\
\midrule
Item A & 10.5 & 20.3 & 30.8 \\
Item B & 15.2 & 25.1 & 40.3 \\
Item C & 22.1 & 12.7 & 34.8 \\
\bottomrule
\end{tabular}
\caption{Sample Data Table}
\label{tab:sample}
\end{table}

\section{Lists}
\subsection{Bullet Points}
\begin{itemize}
  \item First item
  \item Second item
  \item Third item with sub-items:
    \begin{itemize}
      \item Sub-item 1
      \item Sub-item 2
    \end{itemize}
\end{itemize}

\subsection{Numbered List}
\begin{enumerate}
  \item First step
  \item Second step
  \item Third step
\end{enumerate}

\section{Citations and References}
When you use the citation search feature, references will be automatically managed. For example, you might cite a paper like this: \cite{example2024}. All citations will be automatically added to your \texttt{cite.bib} file.

\section{Conclusion}
This template demonstrates basic \LaTeX{} features including mathematical formulas, tables, and lists. The bibliography system is now set up to work seamlessly with Bibby AI's citation management. Feel free to modify it for your needs!

\bibliographystyle{plain}
\bibliography{cite}

\end{document}\documentclass{article}
\usepackage[utf8]{inputenc}
\usepackage{booktabs}
\usepackage{amsmath}
\usepackage{natbib}

\title{My Document}
\author{Your Name}
\date{\today}

\begin{document}

\maketitle

\section{Introduction}
This is a sample document demonstrating various \LaTeX{} features including mathematical equations, tables, and formatting. You can cite papers using the citation search feature, and they will be automatically added to your \texttt{cite.bib} file.

\section{Mathematical Expressions}
Here are some example equations:

\subsection{Quadratic Formula}
The quadratic formula is given by:
\[x = \frac{-b \pm \sqrt{b^2 - 4ac}}{2a}\]

\subsection{Matrix Example}
A sample matrix:
\[
\begin{pmatrix}
a_{11} & a_{12} & a_{13} \\
a_{21} & a_{22} & a_{23} \\
a_{31} & a_{32} & a_{33}
\end{pmatrix}
\]

\section{Tables}
Below is a sample table showing data:

\begin{table}[h]
\centering
\begin{tabular}{@{}lrrr@{}}
\toprule
\textbf{Category} & \textbf{Value 1} & \textbf{Value 2} & \textbf{Total} \\
\midrule
Item A & 10.5 & 20.3 & 30.8 \\
Item B & 15.2 & 25.1 & 40.3 \\
Item C & 22.1 & 12.7 & 34.8 \\
\bottomrule
\end{tabular}
\caption{Sample Data Table}
\label{tab:sample}
\end{table}

\section{Lists}
\subsection{Bullet Points}
\begin{itemize}
  \item First item
  \item Second item
  \item Third item with sub-items:
    \begin{itemize}
      \item Sub-item 1
      \item Sub-item 2
    \end{itemize}
\end{itemize}

\subsection{Numbered List}
\begin{enumerate}
  \item First step
  \item Second step
  \item Third step
\end{enumerate}

\section{Citations and References}
When you use the citation search feature, references will be automatically managed. For example, you might cite a paper like this: \cite{example2024}. All citations will be automatically added to your \texttt{cite.bib} file.

\section{Conclusion}
This template demonstrates basic \LaTeX{} features including mathematical formulas, tables, and lists. The bibliography system is now set up to work seamlessly with Bibby AI's citation management. Feel free to modify it for your needs!

\bibliographystyle{plain}
\bibliography{cite}

\end{document}\documentclass{article}
\usepackage[utf8]{inputenc}
\usepackage{booktabs}
\usepackage{amsmath}
\usepackage{natbib}

\title{My Document}
\author{Your Name}
\date{\today}

\begin{document}

\maketitle

\section{Introduction}
This is a sample document demonstrating various \LaTeX{} features including mathematical equations, tables, and formatting. You can cite papers using the citation search feature, and they will be automatically added to your \texttt{cite.bib} file.

\section{Mathematical Expressions}
Here are some example equations:

\subsection{Quadratic Formula}
The quadratic formula is given by:
\[x = \frac{-b \pm \sqrt{b^2 - 4ac}}{2a}\]

\subsection{Matrix Example}
A sample matrix:
\[
\begin{pmatrix}
a_{11} & a_{12} & a_{13} \\
a_{21} & a_{22} & a_{23} \\
a_{31} & a_{32} & a_{33}
\end{pmatrix}
\]

\section{Tables}
Below is a sample table showing data:

\begin{table}[h]
\centering
\begin{tabular}{@{}lrrr@{}}
\toprule
\textbf{Category} & \textbf{Value 1} & \textbf{Value 2} & \textbf{Total} \\
\midrule
Item A & 10.5 & 20.3 & 30.8 \\
Item B & 15.2 & 25.1 & 40.3 \\
Item C & 22.1 & 12.7 & 34.8 \\
\bottomrule
\end{tabular}
\caption{Sample Data Table}
\label{tab:sample}
\end{table}

\section{Lists}
\subsection{Bullet Points}
\begin{itemize}
  \item First item
  \item Second item
  \item Third item with sub-items:
    \begin{itemize}
      \item Sub-item 1
      \item Sub-item 2
    \end{itemize}
\end{itemize}

\subsection{Numbered List}
\begin{enumerate}
  \item First step
  \item Second step
  \item Third step
\end{enumerate}

\section{Citations and References}
When you use the citation search feature, references will be automatically managed. For example, you might cite a paper like this: \cite{example2024}. All citations will be automatically added to your \texttt{cite.bib} file.

\section{Conclusion}
This template demonstrates basic \LaTeX{} features including mathematical formulas, tables, and lists. The bibliography system is now set up to work seamlessly with Bibby AI's citation management. Feel free to modify it for your needs!

\bibliographystyle{plain}
\bibliography{cite}

\end{document}\documentclass{article}
\usepackage[utf8]{inputenc}
\usepackage{booktabs}
\usepackage{amsmath}
\usepackage{natbib}

\title{My Document}
\author{Your Name}
\date{\today}

\begin{document}

\maketitle

\section{Introduction}
This is a sample document demonstrating various \LaTeX{} features including mathematical equations, tables, and formatting. You can cite papers using the citation search feature, and they will be automatically added to your \texttt{cite.bib} file.

\section{Mathematical Expressions}
Here are some example equations:

\subsection{Quadratic Formula}
The quadratic formula is given by:
\[x = \frac{-b \pm \sqrt{b^2 - 4ac}}{2a}\]

\subsection{Matrix Example}
A sample matrix:
\[
\begin{pmatrix}
a_{11} & a_{12} & a_{13} \\
a_{21} & a_{22} & a_{23} \\
a_{31} & a_{32} & a_{33}
\end{pmatrix}
\]

\section{Tables}
Below is a sample table showing data:

\begin{table}[h]
\centering
\begin{tabular}{@{}lrrr@{}}
\toprule
\textbf{Category} & \textbf{Value 1} & \textbf{Value 2} & \textbf{Total} \\
\midrule
Item A & 10.5 & 20.3 & 30.8 \\
Item B & 15.2 & 25.1 & 40.3 \\
Item C & 22.1 & 12.7 & 34.8 \\
\bottomrule
\end{tabular}
\caption{Sample Data Table}
\label{tab:sample}
\end{table}

\section{Lists}
\subsection{Bullet Points}
\begin{itemize}
  \item First item
  \item Second item
  \item Third item with sub-items:
    \begin{itemize}
      \item Sub-item 1
      \item Sub-item 2
    \end{itemize}
\end{itemize}

\subsection{Numbered List}
\begin{enumerate}
  \item First step
  \item Second step
  \item Third step
\end{enumerate}

\section{Citations and References}
When you use the citation search feature, references will be automatically managed. For example, you might cite a paper like this: \cite{example2024}. All citations will be automatically added to your \texttt{cite.bib} file.

\section{Conclusion}
This template demonstrates basic \LaTeX{} features including mathematical formulas, tables, and lists. The bibliography system is now set up to work seamlessly with Bibby AI's citation management. Feel free to modify it for your needs!

\bibliographystyle{plain}
\bibliography{cite}

\end{document}\documentclass{article}
\usepackage[utf8]{inputenc}
\usepackage{booktabs}
\usepackage{amsmath}
\usepackage{natbib}

\title{My Document}
\author{Your Name}
\date{\today}

\begin{document}

\maketitle

\section{Introduction}
This is a sample document demonstrating various \LaTeX{} features including mathematical equations, tables, and formatting. You can cite papers using the citation search feature, and they will be automatically added to your \texttt{cite.bib} file.

\section{Mathematical Expressions}
Here are some example equations:

\subsection{Quadratic Formula}
The quadratic formula is given by:
\[x = \frac{-b \pm \sqrt{b^2 - 4ac}}{2a}\]

\subsection{Matrix Example}
A sample matrix:
\[
\begin{pmatrix}
a_{11} & a_{12} & a_{13} \\
a_{21} & a_{22} & a_{23} \\
a_{31} & a_{32} & a_{33}
\end{pmatrix}
\]

\section{Tables}
Below is a sample table showing data:

\begin{table}[h]
\centering
\begin{tabular}{@{}lrrr@{}}
\toprule
\textbf{Category} & \textbf{Value 1} & \textbf{Value 2} & \textbf{Total} \\
\midrule
Item A & 10.5 & 20.3 & 30.8 \\
Item B & 15.2 & 25.1 & 40.3 \\
Item C & 22.1 & 12.7 & 34.8 \\
\bottomrule
\end{tabular}
\caption{Sample Data Table}
\label{tab:sample}
\end{table}

\section{Lists}
\subsection{Bullet Points}
\begin{itemize}
  \item First item
  \item Second item
  \item Third item with sub-items:
    \begin{itemize}
      \item Sub-item 1
      \item Sub-item 2
    \end{itemize}
\end{itemize}

\subsection{Numbered List}
\begin{enumerate}
  \item First step
  \item Second step
  \item Third step
\end{enumerate}

\section{Citations and References}
When you use the citation search feature, references will be automatically managed. For example, you might cite a paper like this: \cite{example2024}. All citations will be automatically added to your \texttt{cite.bib} file.

\section{Conclusion}
This template demonstrates basic \LaTeX{} features including mathematical formulas, tables, and lists. The bibliography system is now set up to work seamlessly with Bibby AI's citation management. Feel free to modify it for your needs!

\bibliographystyle{plain}
\bibliography{cite}

\end{document}\documentclass{article}
\usepackage[utf8]{inputenc}
\usepackage{booktabs}
\usepackage{amsmath}
\usepackage{natbib}

\title{My Document}
\author{Your Name}
\date{\today}

\begin{document}

\maketitle

\section{Introduction}
This is a sample document demonstrating various \LaTeX{} features including mathematical equations, tables, and formatting. You can cite papers using the citation search feature, and they will be automatically added to your \texttt{cite.bib} file.

\section{Mathematical Expressions}
Here are some example equations:

\subsection{Quadratic Formula}
The quadratic formula is given by:
\[x = \frac{-b \pm \sqrt{b^2 - 4ac}}{2a}\]

\subsection{Matrix Example}
A sample matrix:
\[
\begin{pmatrix}
a_{11} & a_{12} & a_{13} \\
a_{21} & a_{22} & a_{23} \\
a_{31} & a_{32} & a_{33}
\end{pmatrix}
\]

\section{Tables}
Below is a sample table showing data:

\begin{table}[h]
\centering
\begin{tabular}{@{}lrrr@{}}
\toprule
\textbf{Category} & \textbf{Value 1} & \textbf{Value 2} & \textbf{Total} \\
\midrule
Item A & 10.5 & 20.3 & 30.8 \\
Item B & 15.2 & 25.1 & 40.3 \\
Item C & 22.1 & 12.7 & 34.8 \\
\bottomrule
\end{tabular}
\caption{Sample Data Table}
\label{tab:sample}
\end{table}

\section{Lists}
\subsection{Bullet Points}
\begin{itemize}
  \item First item
  \item Second item
  \item Third item with sub-items:
    \begin{itemize}
      \item Sub-item 1
      \item Sub-item 2
    \end{itemize}
\end{itemize}

\subsection{Numbered List}
\begin{enumerate}
  \item First step
  \item Second step
  \item Third step
\end{enumerate}

\section{Citations and References}
When you use the citation search feature, references will be automatically managed. For example, you might cite a paper like this: \cite{example2024}. All citations will be automatically added to your \texttt{cite.bib} file.

\section{Conclusion}
This template demonstrates basic \LaTeX{} features including mathematical formulas, tables, and lists. The bibliography system is now set up to work seamlessly with Bibby AI's citation management. Feel free to modify it for your needs!

\bibliographystyle{plain}
\bibliography{cite}

\end{document}\documentclass{article}
\usepackage[utf8]{inputenc}
\usepackage{booktabs}
\usepackage{amsmath}
\usepackage{natbib}

\title{My Document}
\author{Your Name}
\date{\today}

\begin{document}

\maketitle

\section{Introduction}
This is a sample document demonstrating various \LaTeX{} features including mathematical equations, tables, and formatting. You can cite papers using the citation search feature, and they will be automatically added to your \texttt{cite.bib} file.

\section{Mathematical Expressions}
Here are some example equations:

\subsection{Quadratic Formula}
The quadratic formula is given by:
\[x = \frac{-b \pm \sqrt{b^2 - 4ac}}{2a}\]

\subsection{Matrix Example}
A sample matrix:
\[
\begin{pmatrix}
a_{11} & a_{12} & a_{13} \\
a_{21} & a_{22} & a_{23} \\
a_{31} & a_{32} & a_{33}
\end{pmatrix}
\]

\section{Tables}
Below is a sample table showing data:

\begin{table}[h]
\centering
\begin{tabular}{@{}lrrr@{}}
\toprule
\textbf{Category} & \textbf{Value 1} & \textbf{Value 2} & \textbf{Total} \\
\midrule
Item A & 10.5 & 20.3 & 30.8 \\
Item B & 15.2 & 25.1 & 40.3 \\
Item C & 22.1 & 12.7 & 34.8 \\
\bottomrule
\end{tabular}
\caption{Sample Data Table}
\label{tab:sample}
\end{table}

\section{Lists}
\subsection{Bullet Points}
\begin{itemize}
  \item First item
  \item Second item
  \item Third item with sub-items:
    \begin{itemize}
      \item Sub-item 1
      \item Sub-item 2
    \end{itemize}
\end{itemize}

\subsection{Numbered List}
\begin{enumerate}
  \item First step
  \item Second step
  \item Third step
\end{enumerate}

\section{Citations and References}
When you use the citation search feature, references will be automatically managed. For example, you might cite a paper like this: \cite{example2024}. All citations will be automatically added to your \texttt{cite.bib} file.

\section{Conclusion}
This template demonstrates basic \LaTeX{} features including mathematical formulas, tables, and lists. The bibliography system is now set up to work seamlessly with Bibby AI's citation management. Feel free to modify it for your needs!

\bibliographystyle{plain}
\bibliography{cite}

\end{document}\documentclass{article}
\usepackage[utf8]{inputenc}
\usepackage{booktabs}
\usepackage{amsmath}
\usepackage{natbib}

\title{My Document}
\author{Your Name}
\date{\today}

\begin{document}

\maketitle

\section{Introduction}
This is a sample document demonstrating various \LaTeX{} features including mathematical equations, tables, and formatting. You can cite papers using the citation search feature, and they will be automatically added to your \texttt{cite.bib} file.

\section{Mathematical Expressions}
Here are some example equations:

\subsection{Quadratic Formula}
The quadratic formula is given by:
\[x = \frac{-b \pm \sqrt{b^2 - 4ac}}{2a}\]

\subsection{Matrix Example}
A sample matrix:
\[
\begin{pmatrix}
a_{11} & a_{12} & a_{13} \\
a_{21} & a_{22} & a_{23} \\
a_{31} & a_{32} & a_{33}
\end{pmatrix}
\]

\section{Tables}
Below is a sample table showing data:

\begin{table}[h]
\centering
\begin{tabular}{@{}lrrr@{}}
\toprule
\textbf{Category} & \textbf{Value 1} & \textbf{Value 2} & \textbf{Total} \\
\midrule
Item A & 10.5 & 20.3 & 30.8 \\
Item B & 15.2 & 25.1 & 40.3 \\
Item C & 22.1 & 12.7 & 34.8 \\
\bottomrule
\end{tabular}
\caption{Sample Data Table}
\label{tab:sample}
\end{table}

\section{Lists}
\subsection{Bullet Points}
\begin{itemize}
  \item First item
  \item Second item
  \item Third item with sub-items:
    \begin{itemize}
      \item Sub-item 1
      \item Sub-item 2
    \end{itemize}
\end{itemize}

\subsection{Numbered List}
\begin{enumerate}
  \item First step
  \item Second step
  \item Third step
\end{enumerate}

\section{Citations and References}
When you use the citation search feature, references will be automatically managed. For example, you might cite a paper like this: \cite{example2024}. All citations will be automatically added to your \texttt{cite.bib} file.

\section{Conclusion}
This template demonstrates basic \LaTeX{} features including mathematical formulas, tables, and lists. The bibliography system is now set up to work seamlessly with Bibby AI's citation management. Feel free to modify it for your needs!

\bibliographystyle{plain}
\bibliography{cite}

\end{document}\documentclass{article}
\usepackage[utf8]{inputenc}
\usepackage{booktabs}
\usepackage{amsmath}
\usepackage{natbib}

\title{My Document}
\author{Your Name}
\date{\today}

\begin{document}

\maketitle

\section{Introduction}
This is a sample document demonstrating various \LaTeX{} features including mathematical equations, tables, and formatting. You can cite papers using the citation search feature, and they will be automatically added to your \texttt{cite.bib} file.

\section{Mathematical Expressions}
Here are some example equations:

\subsection{Quadratic Formula}
The quadratic formula is given by:
\[x = \frac{-b \pm \sqrt{b^2 - 4ac}}{2a}\]

\subsection{Matrix Example}
A sample matrix:
\[
\begin{pmatrix}
a_{11} & a_{12} & a_{13} \\
a_{21} & a_{22} & a_{23} \\
a_{31} & a_{32} & a_{33}
\end{pmatrix}
\]

\section{Tables}
Below is a sample table showing data:

\begin{table}[h]
\centering
\begin{tabular}{@{}lrrr@{}}
\toprule
\textbf{Category} & \textbf{Value 1} & \textbf{Value 2} & \textbf{Total} \\
\midrule
Item A & 10.5 & 20.3 & 30.8 \\
Item B & 15.2 & 25.1 & 40.3 \\
Item C & 22.1 & 12.7 & 34.8 \\
\bottomrule
\end{tabular}
\caption{Sample Data Table}
\label{tab:sample}
\end{table}

\section{Lists}
\subsection{Bullet Points}
\begin{itemize}
  \item First item
  \item Second item
  \item Third item with sub-items:
    \begin{itemize}
      \item Sub-item 1
      \item Sub-item 2
    \end{itemize}
\end{itemize}

\subsection{Numbered List}
\begin{enumerate}
  \item First step
  \item Second step
  \item Third step
\end{enumerate}

\section{Citations and References}
When you use the citation search feature, references will be automatically managed. For example, you might cite a paper like this: \cite{example2024}. All citations will be automatically added to your \texttt{cite.bib} file.

\section{Conclusion}
This template demonstrates basic \LaTeX{} features including mathematical formulas, tables, and lists. The bibliography system is now set up to work seamlessly with Bibby AI's citation management. Feel free to modify it for your needs!

\bibliographystyle{plain}
\bibliography{cite}

\end{document}\documentclass{article}
\usepackage[utf8]{inputenc}
\usepackage{booktabs}
\usepackage{amsmath}
\usepackage{natbib}

\title{My Document}
\author{Your Name}
\date{\today}

\begin{document}

\maketitle

\section{Introduction}
This is a sample document demonstrating various \LaTeX{} features including mathematical equations, tables, and formatting. You can cite papers using the citation search feature, and they will be automatically added to your \texttt{cite.bib} file.

\section{Mathematical Expressions}
Here are some example equations:

\subsection{Quadratic Formula}
The quadratic formula is given by:
\[x = \frac{-b \pm \sqrt{b^2 - 4ac}}{2a}\]

\subsection{Matrix Example}
A sample matrix:
\[
\begin{pmatrix}
a_{11} & a_{12} & a_{13} \\
a_{21} & a_{22} & a_{23} \\
a_{31} & a_{32} & a_{33}
\end{pmatrix}
\]

\section{Tables}
Below is a sample table showing data:

\begin{table}[h]
\centering
\begin{tabular}{@{}lrrr@{}}
\toprule
\textbf{Category} & \textbf{Value 1} & \textbf{Value 2} & \textbf{Total} \\
\midrule
Item A & 10.5 & 20.3 & 30.8 \\
Item B & 15.2 & 25.1 & 40.3 \\
Item C & 22.1 & 12.7 & 34.8 \\
\bottomrule
\end{tabular}
\caption{Sample Data Table}
\label{tab:sample}
\end{table}

\section{Lists}
\subsection{Bullet Points}
\begin{itemize}
  \item First item
  \item Second item
  \item Third item with sub-items:
    \begin{itemize}
      \item Sub-item 1
      \item Sub-item 2
    \end{itemize}
\end{itemize}

\subsection{Numbered List}
\begin{enumerate}
  \item First step
  \item Second step
  \item Third step
\end{enumerate}

\section{Citations and References}
When you use the citation search feature, references will be automatically managed. For example, you might cite a paper like this: \cite{example2024}. All citations will be automatically added to your \texttt{cite.bib} file.

\section{Conclusion}
This template demonstrates basic \LaTeX{} features including mathematical formulas, tables, and lists. The bibliography system is now set up to work seamlessly with Bibby AI's citation management. Feel free to modify it for your needs!

\bibliographystyle{plain}
\bibliography{cite}

\end{document}\documentclass{article}
\usepackage[utf8]{inputenc}
\usepackage{booktabs}
\usepackage{amsmath}
\usepackage{natbib}

\title{My Document}
\author{Your Name}
\date{\today}

\begin{document}

\maketitle

\section{Introduction}
This is a sample document demonstrating various \LaTeX{} features including mathematical equations, tables, and formatting. You can cite papers using the citation search feature, and they will be automatically added to your \texttt{cite.bib} file.

\section{Mathematical Expressions}
Here are some example equations:

\subsection{Quadratic Formula}
The quadratic formula is given by:
\[x = \frac{-b \pm \sqrt{b^2 - 4ac}}{2a}\]

\subsection{Matrix Example}
A sample matrix:
\[
\begin{pmatrix}
a_{11} & a_{12} & a_{13} \\
a_{21} & a_{22} & a_{23} \\
a_{31} & a_{32} & a_{33}
\end{pmatrix}
\]

\section{Tables}
Below is a sample table showing data:

\begin{table}[h]
\centering
\begin{tabular}{@{}lrrr@{}}
\toprule
\textbf{Category} & \textbf{Value 1} & \textbf{Value 2} & \textbf{Total} \\
\midrule
Item A & 10.5 & 20.3 & 30.8 \\
Item B & 15.2 & 25.1 & 40.3 \\
Item C & 22.1 & 12.7 & 34.8 \\
\bottomrule
\end{tabular}
\caption{Sample Data Table}
\label{tab:sample}
\end{table}

\section{Lists}
\subsection{Bullet Points}
\begin{itemize}
  \item First item
  \item Second item
  \item Third item with sub-items:
    \begin{itemize}
      \item Sub-item 1
      \item Sub-item 2
    \end{itemize}
\end{itemize}

\subsection{Numbered List}
\begin{enumerate}
  \item First step
  \item Second step
  \item Third step
\end{enumerate}

\section{Citations and References}
When you use the citation search feature, references will be automatically managed. For example, you might cite a paper like this: \cite{example2024}. All citations will be automatically added to your \texttt{cite.bib} file.

\section{Conclusion}
This template demonstrates basic \LaTeX{} features including mathematical formulas, tables, and lists. The bibliography system is now set up to work seamlessly with Bibby AI's citation management. Feel free to modify it for your needs!

\bibliographystyle{plain}
\bibliography{cite}

\end{document}\documentclass{article}
\usepackage[utf8]{inputenc}
\usepackage{booktabs}
\usepackage{amsmath}
\usepackage{natbib}

\title{My Document}
\author{Your Name}
\date{\today}

\begin{document}

\maketitle

\section{Introduction}
This is a sample document demonstrating various \LaTeX{} features including mathematical equations, tables, and formatting. You can cite papers using the citation search feature, and they will be automatically added to your \texttt{cite.bib} file.

\section{Mathematical Expressions}
Here are some example equations:

\subsection{Quadratic Formula}
The quadratic formula is given by:
\[x = \frac{-b \pm \sqrt{b^2 - 4ac}}{2a}\]

\subsection{Matrix Example}
A sample matrix:
\[
\begin{pmatrix}
a_{11} & a_{12} & a_{13} \\
a_{21} & a_{22} & a_{23} \\
a_{31} & a_{32} & a_{33}
\end{pmatrix}
\]

\section{Tables}
Below is a sample table showing data:

\begin{table}[h]
\centering
\begin{tabular}{@{}lrrr@{}}
\toprule
\textbf{Category} & \textbf{Value 1} & \textbf{Value 2} & \textbf{Total} \\
\midrule
Item A & 10.5 & 20.3 & 30.8 \\
Item B & 15.2 & 25.1 & 40.3 \\
Item C & 22.1 & 12.7 & 34.8 \\
\bottomrule
\end{tabular}
\caption{Sample Data Table}
\label{tab:sample}
\end{table}

\section{Lists}
\subsection{Bullet Points}
\begin{itemize}
  \item First item
  \item Second item
  \item Third item with sub-items:
    \begin{itemize}
      \item Sub-item 1
      \item Sub-item 2
    \end{itemize}
\end{itemize}

\subsection{Numbered List}
\begin{enumerate}
  \item First step
  \item Second step
  \item Third step
\end{enumerate}

\section{Citations and References}
When you use the citation search feature, references will be automatically managed. For example, you might cite a paper like this: \cite{example2024}. All citations will be automatically added to your \texttt{cite.bib} file.

\section{Conclusion}
This template demonstrates basic \LaTeX{} features including mathematical formulas, tables, and lists. The bibliography system is now set up to work seamlessly with Bibby AI's citation management. Feel free to modify it for your needs!

\bibliographystyle{plain}
\bibliography{cite}

\end{document}\documentclass{article}
\usepackage[utf8]{inputenc}
\usepackage{booktabs}
\usepackage{amsmath}
\usepackage{natbib}

\title{My Document}
\author{Your Name}
\date{\today}

\begin{document}

\maketitle

\section{Introduction}
This is a sample document demonstrating various \LaTeX{} features including mathematical equations, tables, and formatting. You can cite papers using the citation search feature, and they will be automatically added to your \texttt{cite.bib} file.

\section{Mathematical Expressions}
Here are some example equations:

\subsection{Quadratic Formula}
The quadratic formula is given by:
\[x = \frac{-b \pm \sqrt{b^2 - 4ac}}{2a}\]

\subsection{Matrix Example}
A sample matrix:
\[
\begin{pmatrix}
a_{11} & a_{12} & a_{13} \\
a_{21} & a_{22} & a_{23} \\
a_{31} & a_{32} & a_{33}
\end{pmatrix}
\]

\section{Tables}
Below is a sample table showing data:

\begin{table}[h]
\centering
\begin{tabular}{@{}lrrr@{}}
\toprule
\textbf{Category} & \textbf{Value 1} & \textbf{Value 2} & \textbf{Total} \\
\midrule
Item A & 10.5 & 20.3 & 30.8 \\
Item B & 15.2 & 25.1 & 40.3 \\
Item C & 22.1 & 12.7 & 34.8 \\
\bottomrule
\end{tabular}
\caption{Sample Data Table}
\label{tab:sample}
\end{table}

\section{Lists}
\subsection{Bullet Points}
\begin{itemize}
  \item First item
  \item Second item
  \item Third item with sub-items:
    \begin{itemize}
      \item Sub-item 1
      \item Sub-item 2
    \end{itemize}
\end{itemize}

\subsection{Numbered List}
\begin{enumerate}
  \item First step
  \item Second step
  \item Third step
\end{enumerate}

\section{Citations and References}
When you use the citation search feature, references will be automatically managed. For example, you might cite a paper like this: \cite{example2024}. All citations will be automatically added to your \texttt{cite.bib} file.

\section{Conclusion}
This template demonstrates basic \LaTeX{} features including mathematical formulas, tables, and lists. The bibliography system is now set up to work seamlessly with Bibby AI's citation management. Feel free to modify it for your needs!

\bibliographystyle{plain}
\bibliography{cite}

\end{document}\documentclass{article}
\usepackage[utf8]{inputenc}
\usepackage{booktabs}
\usepackage{amsmath}
\usepackage{natbib}

\title{My Document}
\author{Your Name}
\date{\today}

\begin{document}

\maketitle

\section{Introduction}
This is a sample document demonstrating various \LaTeX{} features including mathematical equations, tables, and formatting. You can cite papers using the citation search feature, and they will be automatically added to your \texttt{cite.bib} file.

\section{Mathematical Expressions}
Here are some example equations:

\subsection{Quadratic Formula}
The quadratic formula is given by:
\[x = \frac{-b \pm \sqrt{b^2 - 4ac}}{2a}\]

\subsection{Matrix Example}
A sample matrix:
\[
\begin{pmatrix}
a_{11} & a_{12} & a_{13} \\
a_{21} & a_{22} & a_{23} \\
a_{31} & a_{32} & a_{33}
\end{pmatrix}
\]

\section{Tables}
Below is a sample table showing data:

\begin{table}[h]
\centering
\begin{tabular}{@{}lrrr@{}}
\toprule
\textbf{Category} & \textbf{Value 1} & \textbf{Value 2} & \textbf{Total} \\
\midrule
Item A & 10.5 & 20.3 & 30.8 \\
Item B & 15.2 & 25.1 & 40.3 \\
Item C & 22.1 & 12.7 & 34.8 \\
\bottomrule
\end{tabular}
\caption{Sample Data Table}
\label{tab:sample}
\end{table}

\section{Lists}
\subsection{Bullet Points}
\begin{itemize}
  \item First item
  \item Second item
  \item Third item with sub-items:
    \begin{itemize}
      \item Sub-item 1
      \item Sub-item 2
    \end{itemize}
\end{itemize}

\subsection{Numbered List}
\begin{enumerate}
  \item First step
  \item Second step
  \item Third step
\end{enumerate}

\section{Citations and References}
When you use the citation search feature, references will be automatically managed. For example, you might cite a paper like this: \cite{example2024}. All citations will be automatically added to your \texttt{cite.bib} file.

\section{Conclusion}
This template demonstrates basic \LaTeX{} features including mathematical formulas, tables, and lists. The bibliography system is now set up to work seamlessly with Bibby AI's citation management. Feel free to modify it for your needs!

\bibliographystyle{plain}
\bibliography{cite}

\end{document}\documentclass{article}
\usepackage[utf8]{inputenc}
\usepackage{booktabs}
\usepackage{amsmath}
\usepackage{natbib}

\title{My Document}
\author{Your Name}
\date{\today}

\begin{document}

\maketitle

\section{Introduction}
This is a sample document demonstrating various \LaTeX{} features including mathematical equations, tables, and formatting. You can cite papers using the citation search feature, and they will be automatically added to your \texttt{cite.bib} file.

\section{Mathematical Expressions}
Here are some example equations:

\subsection{Quadratic Formula}
The quadratic formula is given by:
\[x = \frac{-b \pm \sqrt{b^2 - 4ac}}{2a}\]

\subsection{Matrix Example}
A sample matrix:
\[
\begin{pmatrix}
a_{11} & a_{12} & a_{13} \\
a_{21} & a_{22} & a_{23} \\
a_{31} & a_{32} & a_{33}
\end{pmatrix}
\]

\section{Tables}
Below is a sample table showing data:

\begin{table}[h]
\centering
\begin{tabular}{@{}lrrr@{}}
\toprule
\textbf{Category} & \textbf{Value 1} & \textbf{Value 2} & \textbf{Total} \\
\midrule
Item A & 10.5 & 20.3 & 30.8 \\
Item B & 15.2 & 25.1 & 40.3 \\
Item C & 22.1 & 12.7 & 34.8 \\
\bottomrule
\end{tabular}
\caption{Sample Data Table}
\label{tab:sample}
\end{table}

\section{Lists}
\subsection{Bullet Points}
\begin{itemize}
  \item First item
  \item Second item
  \item Third item with sub-items:
    \begin{itemize}
      \item Sub-item 1
      \item Sub-item 2
    \end{itemize}
\end{itemize}

\subsection{Numbered List}
\begin{enumerate}
  \item First step
  \item Second step
  \item Third step
\end{enumerate}

\section{Citations and References}
When you use the citation search feature, references will be automatically managed. For example, you might cite a paper like this: \cite{example2024}. All citations will be automatically added to your \texttt{cite.bib} file.

\section{Conclusion}
This template demonstrates basic \LaTeX{} features including mathematical formulas, tables, and lists. The bibliography system is now set up to work seamlessly with Bibby AI's citation management. Feel free to modify it for your needs!

\bibliographystyle{plain}
\bibliography{cite}

\end{document}\documentclass{article}
\usepackage[utf8]{inputenc}
\usepackage{booktabs}
\usepackage{amsmath}
\usepackage{natbib}

\title{My Document}
\author{Your Name}
\date{\today}

\begin{document}

\maketitle

\section{Introduction}
This is a sample document demonstrating various \LaTeX{} features including mathematical equations, tables, and formatting. You can cite papers using the citation search feature, and they will be automatically added to your \texttt{cite.bib} file.

\section{Mathematical Expressions}
Here are some example equations:

\subsection{Quadratic Formula}
The quadratic formula is given by:
\[x = \frac{-b \pm \sqrt{b^2 - 4ac}}{2a}\]

\subsection{Matrix Example}
A sample matrix:
\[
\begin{pmatrix}
a_{11} & a_{12} & a_{13} \\
a_{21} & a_{22} & a_{23} \\
a_{31} & a_{32} & a_{33}
\end{pmatrix}
\]

\section{Tables}
Below is a sample table showing data:

\begin{table}[h]
\centering
\begin{tabular}{@{}lrrr@{}}
\toprule
\textbf{Category} & \textbf{Value 1} & \textbf{Value 2} & \textbf{Total} \\
\midrule
Item A & 10.5 & 20.3 & 30.8 \\
Item B & 15.2 & 25.1 & 40.3 \\
Item C & 22.1 & 12.7 & 34.8 \\
\bottomrule
\end{tabular}
\caption{Sample Data Table}
\label{tab:sample}
\end{table}

\section{Lists}
\subsection{Bullet Points}
\begin{itemize}
  \item First item
  \item Second item
  \item Third item with sub-items:
    \begin{itemize}
      \item Sub-item 1
      \item Sub-item 2
    \end{itemize}
\end{itemize}

\subsection{Numbered List}
\begin{enumerate}
  \item First step
  \item Second step
  \item Third step
\end{enumerate}

\section{Citations and References}
When you use the citation search feature, references will be automatically managed. For example, you might cite a paper like this: \cite{example2024}. All citations will be automatically added to your \texttt{cite.bib} file.

\section{Conclusion}
This template demonstrates basic \LaTeX{} features including mathematical formulas, tables, and lists. The bibliography system is now set up to work seamlessly with Bibby AI's citation management. Feel free to modify it for your needs!

\bibliographystyle{plain}
\bibliography{cite}

\end{document}\documentclass{article}
\usepackage[utf8]{inputenc}
\usepackage{booktabs}
\usepackage{amsmath}
\usepackage{natbib}

\title{My Document}
\author{Your Name}
\date{\today}

\begin{document}

\maketitle

\section{Introduction}
This is a sample document demonstrating various \LaTeX{} features including mathematical equations, tables, and formatting. You can cite papers using the citation search feature, and they will be automatically added to your \texttt{cite.bib} file.

\section{Mathematical Expressions}
Here are some example equations:

\subsection{Quadratic Formula}
The quadratic formula is given by:
\[x = \frac{-b \pm \sqrt{b^2 - 4ac}}{2a}\]

\subsection{Matrix Example}
A sample matrix:
\[
\begin{pmatrix}
a_{11} & a_{12} & a_{13} \\
a_{21} & a_{22} & a_{23} \\
a_{31} & a_{32} & a_{33}
\end{pmatrix}
\]

\section{Tables}
Below is a sample table showing data:

\begin{table}[h]
\centering
\begin{tabular}{@{}lrrr@{}}
\toprule
\textbf{Category} & \textbf{Value 1} & \textbf{Value 2} & \textbf{Total} \\
\midrule
Item A & 10.5 & 20.3 & 30.8 \\
Item B & 15.2 & 25.1 & 40.3 \\
Item C & 22.1 & 12.7 & 34.8 \\
\bottomrule
\end{tabular}
\caption{Sample Data Table}
\label{tab:sample}
\end{table}

\section{Lists}
\subsection{Bullet Points}
\begin{itemize}
  \item First item
  \item Second item
  \item Third item with sub-items:
    \begin{itemize}
      \item Sub-item 1
      \item Sub-item 2
    \end{itemize}
\end{itemize}

\subsection{Numbered List}
\begin{enumerate}
  \item First step
  \item Second step
  \item Third step
\end{enumerate}

\section{Citations and References}
When you use the citation search feature, references will be automatically managed. For example, you might cite a paper like this: \cite{example2024}. All citations will be automatically added to your \texttt{cite.bib} file.

\section{Conclusion}
This template demonstrates basic \LaTeX{} features including mathematical formulas, tables, and lists. The bibliography system is now set up to work seamlessly with Bibby AI's citation management. Feel free to modify it for your needs!

\bibliographystyle{plain}
\bibliography{cite}

\end{document}\documentclass{article}
\usepackage[utf8]{inputenc}
\usepackage{booktabs}
\usepackage{amsmath}
\usepackage{natbib}

\title{My Document}
\author{Your Name}
\date{\today}

\begin{document}

\maketitle

\section{Introduction}
This is a sample document demonstrating various \LaTeX{} features including mathematical equations, tables, and formatting. You can cite papers using the citation search feature, and they will be automatically added to your \texttt{cite.bib} file.

\section{Mathematical Expressions}
Here are some example equations:

\subsection{Quadratic Formula}
The quadratic formula is given by:
\[x = \frac{-b \pm \sqrt{b^2 - 4ac}}{2a}\]

\subsection{Matrix Example}
A sample matrix:
\[
\begin{pmatrix}
a_{11} & a_{12} & a_{13} \\
a_{21} & a_{22} & a_{23} \\
a_{31} & a_{32} & a_{33}
\end{pmatrix}
\]

\section{Tables}
Below is a sample table showing data:

\begin{table}[h]
\centering
\begin{tabular}{@{}lrrr@{}}
\toprule
\textbf{Category} & \textbf{Value 1} & \textbf{Value 2} & \textbf{Total} \\
\midrule
Item A & 10.5 & 20.3 & 30.8 \\
Item B & 15.2 & 25.1 & 40.3 \\
Item C & 22.1 & 12.7 & 34.8 \\
\bottomrule
\end{tabular}
\caption{Sample Data Table}
\label{tab:sample}
\end{table}

\section{Lists}
\subsection{Bullet Points}
\begin{itemize}
  \item First item
  \item Second item
  \item Third item with sub-items:
    \begin{itemize}
      \item Sub-item 1
      \item Sub-item 2
    \end{itemize}
\end{itemize}

\subsection{Numbered List}
\begin{enumerate}
  \item First step
  \item Second step
  \item Third step
\end{enumerate}

\section{Citations and References}
When you use the citation search feature, references will be automatically managed. For example, you might cite a paper like this: \cite{example2024}. All citations will be automatically added to your \texttt{cite.bib} file.

\section{Conclusion}
This template demonstrates basic \LaTeX{} features including mathematical formulas, tables, and lists. The bibliography system is now set up to work seamlessly with Bibby AI's citation management. Feel free to modify it for your needs!

\bibliographystyle{plain}
\bibliography{cite}

\end{document}\documentclass{article}
\usepackage[utf8]{inputenc}
\usepackage{booktabs}
\usepackage{amsmath}
\usepackage{natbib}

\title{My Document}
\author{Your Name}
\date{\today}

\begin{document}

\maketitle

\section{Introduction}
This is a sample document demonstrating various \LaTeX{} features including mathematical equations, tables, and formatting. You can cite papers using the citation search feature, and they will be automatically added to your \texttt{cite.bib} file.

\section{Mathematical Expressions}
Here are some example equations:

\subsection{Quadratic Formula}
The quadratic formula is given by:
\[x = \frac{-b \pm \sqrt{b^2 - 4ac}}{2a}\]

\subsection{Matrix Example}
A sample matrix:
\[
\begin{pmatrix}
a_{11} & a_{12} & a_{13} \\
a_{21} & a_{22} & a_{23} \\
a_{31} & a_{32} & a_{33}
\end{pmatrix}
\]

\section{Tables}
Below is a sample table showing data:

\begin{table}[h]
\centering
\begin{tabular}{@{}lrrr@{}}
\toprule
\textbf{Category} & \textbf{Value 1} & \textbf{Value 2} & \textbf{Total} \\
\midrule
Item A & 10.5 & 20.3 & 30.8 \\
Item B & 15.2 & 25.1 & 40.3 \\
Item C & 22.1 & 12.7 & 34.8 \\
\bottomrule
\end{tabular}
\caption{Sample Data Table}
\label{tab:sample}
\end{table}

\section{Lists}
\subsection{Bullet Points}
\begin{itemize}
  \item First item
  \item Second item
  \item Third item with sub-items:
    \begin{itemize}
      \item Sub-item 1
      \item Sub-item 2
    \end{itemize}
\end{itemize}

\subsection{Numbered List}
\begin{enumerate}
  \item First step
  \item Second step
  \item Third step
\end{enumerate}

\section{Citations and References}
When you use the citation search feature, references will be automatically managed. For example, you might cite a paper like this: \cite{example2024}. All citations will be automatically added to your \texttt{cite.bib} file.

\section{Conclusion}
This template demonstrates basic \LaTeX{} features including mathematical formulas, tables, and lists. The bibliography system is now set up to work seamlessly with Bibby AI's citation management. Feel free to modify it for your needs!

\bibliographystyle{plain}
\bibliography{cite}

\end{document}\documentclass{article}
\usepackage[utf8]{inputenc}
\usepackage{booktabs}
\usepackage{amsmath}
\usepackage{natbib}

\title{My Document}
\author{Your Name}
\date{\today}

\begin{document}

\maketitle

\section{Introduction}
This is a sample document demonstrating various \LaTeX{} features including mathematical equations, tables, and formatting. You can cite papers using the citation search feature, and they will be automatically added to your \texttt{cite.bib} file.

\section{Mathematical Expressions}
Here are some example equations:

\subsection{Quadratic Formula}
The quadratic formula is given by:
\[x = \frac{-b \pm \sqrt{b^2 - 4ac}}{2a}\]

\subsection{Matrix Example}
A sample matrix:
\[
\begin{pmatrix}
a_{11} & a_{12} & a_{13} \\
a_{21} & a_{22} & a_{23} \\
a_{31} & a_{32} & a_{33}
\end{pmatrix}
\]

\section{Tables}
Below is a sample table showing data:

\begin{table}[h]
\centering
\begin{tabular}{@{}lrrr@{}}
\toprule
\textbf{Category} & \textbf{Value 1} & \textbf{Value 2} & \textbf{Total} \\
\midrule
Item A & 10.5 & 20.3 & 30.8 \\
Item B & 15.2 & 25.1 & 40.3 \\
Item C & 22.1 & 12.7 & 34.8 \\
\bottomrule
\end{tabular}
\caption{Sample Data Table}
\label{tab:sample}
\end{table}

\section{Lists}
\subsection{Bullet Points}
\begin{itemize}
  \item First item
  \item Second item
  \item Third item with sub-items:
    \begin{itemize}
      \item Sub-item 1
      \item Sub-item 2
    \end{itemize}
\end{itemize}

\subsection{Numbered List}
\begin{enumerate}
  \item First step
  \item Second step
  \item Third step
\end{enumerate}

\section{Citations and References}
When you use the citation search feature, references will be automatically managed. For example, you might cite a paper like this: \cite{example2024}. All citations will be automatically added to your \texttt{cite.bib} file.

\section{Conclusion}
This template demonstrates basic \LaTeX{} features including mathematical formulas, tables, and lists. The bibliography system is now set up to work seamlessly with Bibby AI's citation management. Feel free to modify it for your needs!

\bibliographystyle{plain}
\bibliography{cite}

\end{document}\documentclass{article}
\usepackage[utf8]{inputenc}
\usepackage{booktabs}
\usepackage{amsmath}
\usepackage{natbib}

\title{My Document}
\author{Your Name}
\date{\today}

\begin{document}

\maketitle

\section{Introduction}
This is a sample document demonstrating various \LaTeX{} features including mathematical equations, tables, and formatting. You can cite papers using the citation search feature, and they will be automatically added to your \texttt{cite.bib} file.

\section{Mathematical Expressions}
Here are some example equations:

\subsection{Quadratic Formula}
The quadratic formula is given by:
\[x = \frac{-b \pm \sqrt{b^2 - 4ac}}{2a}\]

\subsection{Matrix Example}
A sample matrix:
\[
\begin{pmatrix}
a_{11} & a_{12} & a_{13} \\
a_{21} & a_{22} & a_{23} \\
a_{31} & a_{32} & a_{33}
\end{pmatrix}
\]

\section{Tables}
Below is a sample table showing data:

\begin{table}[h]
\centering
\begin{tabular}{@{}lrrr@{}}
\toprule
\textbf{Category} & \textbf{Value 1} & \textbf{Value 2} & \textbf{Total} \\
\midrule
Item A & 10.5 & 20.3 & 30.8 \\
Item B & 15.2 & 25.1 & 40.3 \\
Item C & 22.1 & 12.7 & 34.8 \\
\bottomrule
\end{tabular}
\caption{Sample Data Table}
\label{tab:sample}
\end{table}

\section{Lists}
\subsection{Bullet Points}
\begin{itemize}
  \item First item
  \item Second item
  \item Third item with sub-items:
    \begin{itemize}
      \item Sub-item 1
      \item Sub-item 2
    \end{itemize}
\end{itemize}

\subsection{Numbered List}
\begin{enumerate}
  \item First step
  \item Second step
  \item Third step
\end{enumerate}

\section{Citations and References}
When you use the citation search feature, references will be automatically managed. For example, you might cite a paper like this: \cite{example2024}. All citations will be automatically added to your \texttt{cite.bib} file.

\section{Conclusion}
This template demonstrates basic \LaTeX{} features including mathematical formulas, tables, and lists. The bibliography system is now set up to work seamlessly with Bibby AI's citation management. Feel free to modify it for your needs!

\bibliographystyle{plain}
\bibliography{cite}

\end{document}\documentclass{article}
\usepackage[utf8]{inputenc}
\usepackage{booktabs}
\usepackage{amsmath}
\usepackage{natbib}

\title{My Document}
\author{Your Name}
\date{\today}

\begin{document}

\maketitle

\section{Introduction}
This is a sample document demonstrating various \LaTeX{} features including mathematical equations, tables, and formatting. You can cite papers using the citation search feature, and they will be automatically added to your \texttt{cite.bib} file.

\section{Mathematical Expressions}
Here are some example equations:

\subsection{Quadratic Formula}
The quadratic formula is given by:
\[x = \frac{-b \pm \sqrt{b^2 - 4ac}}{2a}\]

\subsection{Matrix Example}
A sample matrix:
\[
\begin{pmatrix}
a_{11} & a_{12} & a_{13} \\
a_{21} & a_{22} & a_{23} \\
a_{31} & a_{32} & a_{33}
\end{pmatrix}
\]

\section{Tables}
Below is a sample table showing data:

\begin{table}[h]
\centering
\begin{tabular}{@{}lrrr@{}}
\toprule
\textbf{Category} & \textbf{Value 1} & \textbf{Value 2} & \textbf{Total} \\
\midrule
Item A & 10.5 & 20.3 & 30.8 \\
Item B & 15.2 & 25.1 & 40.3 \\
Item C & 22.1 & 12.7 & 34.8 \\
\bottomrule
\end{tabular}
\caption{Sample Data Table}
\label{tab:sample}
\end{table}

\section{Lists}
\subsection{Bullet Points}
\begin{itemize}
  \item First item
  \item Second item
  \item Third item with sub-items:
    \begin{itemize}
      \item Sub-item 1
      \item Sub-item 2
    \end{itemize}
\end{itemize}

\subsection{Numbered List}
\begin{enumerate}
  \item First step
  \item Second step
  \item Third step
\end{enumerate}

\section{Citations and References}
When you use the citation search feature, references will be automatically managed. For example, you might cite a paper like this: \cite{example2024}. All citations will be automatically added to your \texttt{cite.bib} file.

\section{Conclusion}
This template demonstrates basic \LaTeX{} features including mathematical formulas, tables, and lists. The bibliography system is now set up to work seamlessly with Bibby AI's citation management. Feel free to modify it for your needs!

\bibliographystyle{plain}
\bibliography{cite}

\end{document}\documentclass{article}
\usepackage[utf8]{inputenc}
\usepackage{booktabs}
\usepackage{amsmath}
\usepackage{natbib}

\title{My Document}
\author{Your Name}
\date{\today}

\begin{document}

\maketitle

\section{Introduction}
This is a sample document demonstrating various \LaTeX{} features including mathematical equations, tables, and formatting. You can cite papers using the citation search feature, and they will be automatically added to your \texttt{cite.bib} file.

\section{Mathematical Expressions}
Here are some example equations:

\subsection{Quadratic Formula}
The quadratic formula is given by:
\[x = \frac{-b \pm \sqrt{b^2 - 4ac}}{2a}\]

\subsection{Matrix Example}
A sample matrix:
\[
\begin{pmatrix}
a_{11} & a_{12} & a_{13} \\
a_{21} & a_{22} & a_{23} \\
a_{31} & a_{32} & a_{33}
\end{pmatrix}
\]

\section{Tables}
Below is a sample table showing data:

\begin{table}[h]
\centering
\begin{tabular}{@{}lrrr@{}}
\toprule
\textbf{Category} & \textbf{Value 1} & \textbf{Value 2} & \textbf{Total} \\
\midrule
Item A & 10.5 & 20.3 & 30.8 \\
Item B & 15.2 & 25.1 & 40.3 \\
Item C & 22.1 & 12.7 & 34.8 \\
\bottomrule
\end{tabular}
\caption{Sample Data Table}
\label{tab:sample}
\end{table}

\section{Lists}
\subsection{Bullet Points}
\begin{itemize}
  \item First item
  \item Second item
  \item Third item with sub-items:
    \begin{itemize}
      \item Sub-item 1
      \item Sub-item 2
    \end{itemize}
\end{itemize}

\subsection{Numbered List}
\begin{enumerate}
  \item First step
  \item Second step
  \item Third step
\end{enumerate}

\section{Citations and References}
When you use the citation search feature, references will be automatically managed. For example, you might cite a paper like this: \cite{example2024}. All citations will be automatically added to your \texttt{cite.bib} file.

\section{Conclusion}
This template demonstrates basic \LaTeX{} features including mathematical formulas, tables, and lists. The bibliography system is now set up to work seamlessly with Bibby AI's citation management. Feel free to modify it for your needs!

\bibliographystyle{plain}
\bibliography{cite}

\end{document}\documentclass{article}
\usepackage[utf8]{inputenc}
\usepackage{booktabs}
\usepackage{amsmath}
\usepackage{natbib}

\title{My Document}
\author{Your Name}
\date{\today}

\begin{document}

\maketitle

\section{Introduction}
This is a sample document demonstrating various \LaTeX{} features including mathematical equations, tables, and formatting. You can cite papers using the citation search feature, and they will be automatically added to your \texttt{cite.bib} file.

\section{Mathematical Expressions}
Here are some example equations:

\subsection{Quadratic Formula}
The quadratic formula is given by:
\[x = \frac{-b \pm \sqrt{b^2 - 4ac}}{2a}\]

\subsection{Matrix Example}
A sample matrix:
\[
\begin{pmatrix}
a_{11} & a_{12} & a_{13} \\
a_{21} & a_{22} & a_{23} \\
a_{31} & a_{32} & a_{33}
\end{pmatrix}
\]

\section{Tables}
Below is a sample table showing data:

\begin{table}[h]
\centering
\begin{tabular}{@{}lrrr@{}}
\toprule
\textbf{Category} & \textbf{Value 1} & \textbf{Value 2} & \textbf{Total} \\
\midrule
Item A & 10.5 & 20.3 & 30.8 \\
Item B & 15.2 & 25.1 & 40.3 \\
Item C & 22.1 & 12.7 & 34.8 \\
\bottomrule
\end{tabular}
\caption{Sample Data Table}
\label{tab:sample}
\end{table}

\section{Lists}
\subsection{Bullet Points}
\begin{itemize}
  \item First item
  \item Second item
  \item Third item with sub-items:
    \begin{itemize}
      \item Sub-item 1
      \item Sub-item 2
    \end{itemize}
\end{itemize}

\subsection{Numbered List}
\begin{enumerate}
  \item First step
  \item Second step
  \item Third step
\end{enumerate}

\section{Citations and References}
When you use the citation search feature, references will be automatically managed. For example, you might cite a paper like this: \cite{example2024}. All citations will be automatically added to your \texttt{cite.bib} file.

\section{Conclusion}
This template demonstrates basic \LaTeX{} features including mathematical formulas, tables, and lists. The bibliography system is now set up to work seamlessly with Bibby AI's citation management. Feel free to modify it for your needs!

\bibliographystyle{plain}
\bibliography{cite}

\end{document}\documentclass{article}
\usepackage[utf8]{inputenc}
\usepackage{booktabs}
\usepackage{amsmath}
\usepackage{natbib}

\title{My Document}
\author{Your Name}
\date{\today}

\begin{document}

\maketitle

\section{Introduction}
This is a sample document demonstrating various \LaTeX{} features including mathematical equations, tables, and formatting. You can cite papers using the citation search feature, and they will be automatically added to your \texttt{cite.bib} file.

\section{Mathematical Expressions}
Here are some example equations:

\subsection{Quadratic Formula}
The quadratic formula is given by:
\[x = \frac{-b \pm \sqrt{b^2 - 4ac}}{2a}\]

\subsection{Matrix Example}
A sample matrix:
\[
\begin{pmatrix}
a_{11} & a_{12} & a_{13} \\
a_{21} & a_{22} & a_{23} \\
a_{31} & a_{32} & a_{33}
\end{pmatrix}
\]

\section{Tables}
Below is a sample table showing data:

\begin{table}[h]
\centering
\begin{tabular}{@{}lrrr@{}}
\toprule
\textbf{Category} & \textbf{Value 1} & \textbf{Value 2} & \textbf{Total} \\
\midrule
Item A & 10.5 & 20.3 & 30.8 \\
Item B & 15.2 & 25.1 & 40.3 \\
Item C & 22.1 & 12.7 & 34.8 \\
\bottomrule
\end{tabular}
\caption{Sample Data Table}
\label{tab:sample}
\end{table}

\section{Lists}
\subsection{Bullet Points}
\begin{itemize}
  \item First item
  \item Second item
  \item Third item with sub-items:
    \begin{itemize}
      \item Sub-item 1
      \item Sub-item 2
    \end{itemize}
\end{itemize}

\subsection{Numbered List}
\begin{enumerate}
  \item First step
  \item Second step
  \item Third step
\end{enumerate}

\section{Citations and References}
When you use the citation search feature, references will be automatically managed. For example, you might cite a paper like this: \cite{example2024}. All citations will be automatically added to your \texttt{cite.bib} file.

\section{Conclusion}
This template demonstrates basic \LaTeX{} features including mathematical formulas, tables, and lists. The bibliography system is now set up to work seamlessly with Bibby AI's citation management. Feel free to modify it for your needs!

\bibliographystyle{plain}
\bibliography{cite}

\end{document}\documentclass{article}
\usepackage[utf8]{inputenc}
\usepackage{booktabs}
\usepackage{amsmath}
\usepackage{natbib}

\title{My Document}
\author{Your Name}
\date{\today}

\begin{document}

\maketitle

\section{Introduction}
This is a sample document demonstrating various \LaTeX{} features including mathematical equations, tables, and formatting. You can cite papers using the citation search feature, and they will be automatically added to your \texttt{cite.bib} file.

\section{Mathematical Expressions}
Here are some example equations:

\subsection{Quadratic Formula}
The quadratic formula is given by:
\[x = \frac{-b \pm \sqrt{b^2 - 4ac}}{2a}\]

\subsection{Matrix Example}
A sample matrix:
\[
\begin{pmatrix}
a_{11} & a_{12} & a_{13} \\
a_{21} & a_{22} & a_{23} \\
a_{31} & a_{32} & a_{33}
\end{pmatrix}
\]

\section{Tables}
Below is a sample table showing data:

\begin{table}[h]
\centering
\begin{tabular}{@{}lrrr@{}}
\toprule
\textbf{Category} & \textbf{Value 1} & \textbf{Value 2} & \textbf{Total} \\
\midrule
Item A & 10.5 & 20.3 & 30.8 \\
Item B & 15.2 & 25.1 & 40.3 \\
Item C & 22.1 & 12.7 & 34.8 \\
\bottomrule
\end{tabular}
\caption{Sample Data Table}
\label{tab:sample}
\end{table}

\section{Lists}
\subsection{Bullet Points}
\begin{itemize}
  \item First item
  \item Second item
  \item Third item with sub-items:
    \begin{itemize}
      \item Sub-item 1
      \item Sub-item 2
    \end{itemize}
\end{itemize}

\subsection{Numbered List}
\begin{enumerate}
  \item First step
  \item Second step
  \item Third step
\end{enumerate}

\section{Citations and References}
When you use the citation search feature, references will be automatically managed. For example, you might cite a paper like this: \cite{example2024}. All citations will be automatically added to your \texttt{cite.bib} file.

\section{Conclusion}
This template demonstrates basic \LaTeX{} features including mathematical formulas, tables, and lists. The bibliography system is now set up to work seamlessly with Bibby AI's citation management. Feel free to modify it for your needs!

\bibliographystyle{plain}
\bibliography{cite}

\end{document}\documentclass{article}
\usepackage[utf8]{inputenc}
\usepackage{booktabs}
\usepackage{amsmath}
\usepackage{natbib}

\title{My Document}
\author{Your Name}
\date{\today}

\begin{document}

\maketitle

\section{Introduction}
This is a sample document demonstrating various \LaTeX{} features including mathematical equations, tables, and formatting. You can cite papers using the citation search feature, and they will be automatically added to your \texttt{cite.bib} file.

\section{Mathematical Expressions}
Here are some example equations:

\subsection{Quadratic Formula}
The quadratic formula is given by:
\[x = \frac{-b \pm \sqrt{b^2 - 4ac}}{2a}\]

\subsection{Matrix Example}
A sample matrix:
\[
\begin{pmatrix}
a_{11} & a_{12} & a_{13} \\
a_{21} & a_{22} & a_{23} \\
a_{31} & a_{32} & a_{33}
\end{pmatrix}
\]

\section{Tables}
Below is a sample table showing data:

\begin{table}[h]
\centering
\begin{tabular}{@{}lrrr@{}}
\toprule
\textbf{Category} & \textbf{Value 1} & \textbf{Value 2} & \textbf{Total} \\
\midrule
Item A & 10.5 & 20.3 & 30.8 \\
Item B & 15.2 & 25.1 & 40.3 \\
Item C & 22.1 & 12.7 & 34.8 \\
\bottomrule
\end{tabular}
\caption{Sample Data Table}
\label{tab:sample}
\end{table}

\section{Lists}
\subsection{Bullet Points}
\begin{itemize}
  \item First item
  \item Second item
  \item Third item with sub-items:
    \begin{itemize}
      \item Sub-item 1
      \item Sub-item 2
    \end{itemize}
\end{itemize}

\subsection{Numbered List}
\begin{enumerate}
  \item First step
  \item Second step
  \item Third step
\end{enumerate}

\section{Citations and References}
When you use the citation search feature, references will be automatically managed. For example, you might cite a paper like this: \cite{example2024}. All citations will be automatically added to your \texttt{cite.bib} file.

\section{Conclusion}
This template demonstrates basic \LaTeX{} features including mathematical formulas, tables, and lists. The bibliography system is now set up to work seamlessly with Bibby AI's citation management. Feel free to modify it for your needs!

\bibliographystyle{plain}
\bibliography{cite}

\end{document}\documentclass{article}
\usepackage[utf8]{inputenc}
\usepackage{booktabs}
\usepackage{amsmath}
\usepackage{natbib}

\title{My Document}
\author{Your Name}
\date{\today}

\begin{document}

\maketitle

\section{Introduction}
This is a sample document demonstrating various \LaTeX{} features including mathematical equations, tables, and formatting. You can cite papers using the citation search feature, and they will be automatically added to your \texttt{cite.bib} file.

\section{Mathematical Expressions}
Here are some example equations:

\subsection{Quadratic Formula}
The quadratic formula is given by:
\[x = \frac{-b \pm \sqrt{b^2 - 4ac}}{2a}\]

\subsection{Matrix Example}
A sample matrix:
\[
\begin{pmatrix}
a_{11} & a_{12} & a_{13} \\
a_{21} & a_{22} & a_{23} \\
a_{31} & a_{32} & a_{33}
\end{pmatrix}
\]

\section{Tables}
Below is a sample table showing data:

\begin{table}[h]
\centering
\begin{tabular}{@{}lrrr@{}}
\toprule
\textbf{Category} & \textbf{Value 1} & \textbf{Value 2} & \textbf{Total} \\
\midrule
Item A & 10.5 & 20.3 & 30.8 \\
Item B & 15.2 & 25.1 & 40.3 \\
Item C & 22.1 & 12.7 & 34.8 \\
\bottomrule
\end{tabular}
\caption{Sample Data Table}
\label{tab:sample}
\end{table}

\section{Lists}
\subsection{Bullet Points}
\begin{itemize}
  \item First item
  \item Second item
  \item Third item with sub-items:
    \begin{itemize}
      \item Sub-item 1
      \item Sub-item 2
    \end{itemize}
\end{itemize}

\subsection{Numbered List}
\begin{enumerate}
  \item First step
  \item Second step
  \item Third step
\end{enumerate}

\section{Citations and References}
When you use the citation search feature, references will be automatically managed. For example, you might cite a paper like this: \cite{example2024}. All citations will be automatically added to your \texttt{cite.bib} file.

\section{Conclusion}
This template demonstrates basic \LaTeX{} features including mathematical formulas, tables, and lists. The bibliography system is now set up to work seamlessly with Bibby AI's citation management. Feel free to modify it for your needs!

\bibliographystyle{plain}
\bibliography{cite}

\end{document}\documentclass{article}
\usepackage[utf8]{inputenc}
\usepackage{booktabs}
\usepackage{amsmath}
\usepackage{natbib}

\title{My Document}
\author{Your Name}
\date{\today}

\begin{document}

\maketitle

\section{Introduction}
This is a sample document demonstrating various \LaTeX{} features including mathematical equations, tables, and formatting. You can cite papers using the citation search feature, and they will be automatically added to your \texttt{cite.bib} file.

\section{Mathematical Expressions}
Here are some example equations:

\subsection{Quadratic Formula}
The quadratic formula is given by:
\[x = \frac{-b \pm \sqrt{b^2 - 4ac}}{2a}\]

\subsection{Matrix Example}
A sample matrix:
\[
\begin{pmatrix}
a_{11} & a_{12} & a_{13} \\
a_{21} & a_{22} & a_{23} \\
a_{31} & a_{32} & a_{33}
\end{pmatrix}
\]

\section{Tables}
Below is a sample table showing data:

\begin{table}[h]
\centering
\begin{tabular}{@{}lrrr@{}}
\toprule
\textbf{Category} & \textbf{Value 1} & \textbf{Value 2} & \textbf{Total} \\
\midrule
Item A & 10.5 & 20.3 & 30.8 \\
Item B & 15.2 & 25.1 & 40.3 \\
Item C & 22.1 & 12.7 & 34.8 \\
\bottomrule
\end{tabular}
\caption{Sample Data Table}
\label{tab:sample}
\end{table}

\section{Lists}
\subsection{Bullet Points}
\begin{itemize}
  \item First item
  \item Second item
  \item Third item with sub-items:
    \begin{itemize}
      \item Sub-item 1
      \item Sub-item 2
    \end{itemize}
\end{itemize}

\subsection{Numbered List}
\begin{enumerate}
  \item First step
  \item Second step
  \item Third step
\end{enumerate}

\section{Citations and References}
When you use the citation search feature, references will be automatically managed. For example, you might cite a paper like this: \cite{example2024}. All citations will be automatically added to your \texttt{cite.bib} file.

\section{Conclusion}
This template demonstrates basic \LaTeX{} features including mathematical formulas, tables, and lists. The bibliography system is now set up to work seamlessly with Bibby AI's citation management. Feel free to modify it for your needs!

\bibliographystyle{plain}
\bibliography{cite}

\end{document}\documentclass{article}
\usepackage[utf8]{inputenc}
\usepackage{booktabs}
\usepackage{amsmath}
\usepackage{natbib}

\title{My Document}
\author{Your Name}
\date{\today}

\begin{document}

\maketitle

\section{Introduction}
This is a sample document demonstrating various \LaTeX{} features including mathematical equations, tables, and formatting. You can cite papers using the citation search feature, and they will be automatically added to your \texttt{cite.bib} file.

\section{Mathematical Expressions}
Here are some example equations:

\subsection{Quadratic Formula}
The quadratic formula is given by:
\[x = \frac{-b \pm \sqrt{b^2 - 4ac}}{2a}\]

\subsection{Matrix Example}
A sample matrix:
\[
\begin{pmatrix}
a_{11} & a_{12} & a_{13} \\
a_{21} & a_{22} & a_{23} \\
a_{31} & a_{32} & a_{33}
\end{pmatrix}
\]

\section{Tables}
Below is a sample table showing data:

\begin{table}[h]
\centering
\begin{tabular}{@{}lrrr@{}}
\toprule
\textbf{Category} & \textbf{Value 1} & \textbf{Value 2} & \textbf{Total} \\
\midrule
Item A & 10.5 & 20.3 & 30.8 \\
Item B & 15.2 & 25.1 & 40.3 \\
Item C & 22.1 & 12.7 & 34.8 \\
\bottomrule
\end{tabular}
\caption{Sample Data Table}
\label{tab:sample}
\end{table}

\section{Lists}
\subsection{Bullet Points}
\begin{itemize}
  \item First item
  \item Second item
  \item Third item with sub-items:
    \begin{itemize}
      \item Sub-item 1
      \item Sub-item 2
    \end{itemize}
\end{itemize}

\subsection{Numbered List}
\begin{enumerate}
  \item First step
  \item Second step
  \item Third step
\end{enumerate}

\section{Citations and References}
When you use the citation search feature, references will be automatically managed. For example, you might cite a paper like this: \cite{example2024}. All citations will be automatically added to your \texttt{cite.bib} file.

\section{Conclusion}
This template demonstrates basic \LaTeX{} features including mathematical formulas, tables, and lists. The bibliography system is now set up to work seamlessly with Bibby AI's citation management. Feel free to modify it for your needs!

\bibliographystyle{plain}
\bibliography{cite}

\end{document}\documentclass{article}
\usepackage[utf8]{inputenc}
\usepackage{booktabs}
\usepackage{amsmath}
\usepackage{natbib}

\title{My Document}
\author{Your Name}
\date{\today}

\begin{document}

\maketitle

\section{Introduction}
This is a sample document demonstrating various \LaTeX{} features including mathematical equations, tables, and formatting. You can cite papers using the citation search feature, and they will be automatically added to your \texttt{cite.bib} file.

\section{Mathematical Expressions}
Here are some example equations:

\subsection{Quadratic Formula}
The quadratic formula is given by:
\[x = \frac{-b \pm \sqrt{b^2 - 4ac}}{2a}\]

\subsection{Matrix Example}
A sample matrix:
\[
\begin{pmatrix}
a_{11} & a_{12} & a_{13} \\
a_{21} & a_{22} & a_{23} \\
a_{31} & a_{32} & a_{33}
\end{pmatrix}
\]

\section{Tables}
Below is a sample table showing data:

\begin{table}[h]
\centering
\begin{tabular}{@{}lrrr@{}}
\toprule
\textbf{Category} & \textbf{Value 1} & \textbf{Value 2} & \textbf{Total} \\
\midrule
Item A & 10.5 & 20.3 & 30.8 \\
Item B & 15.2 & 25.1 & 40.3 \\
Item C & 22.1 & 12.7 & 34.8 \\
\bottomrule
\end{tabular}
\caption{Sample Data Table}
\label{tab:sample}
\end{table}

\section{Lists}
\subsection{Bullet Points}
\begin{itemize}
  \item First item
  \item Second item
  \item Third item with sub-items:
    \begin{itemize}
      \item Sub-item 1
      \item Sub-item 2
    \end{itemize}
\end{itemize}

\subsection{Numbered List}
\begin{enumerate}
  \item First step
  \item Second step
  \item Third step
\end{enumerate}

\section{Citations and References}
When you use the citation search feature, references will be automatically managed. For example, you might cite a paper like this: \cite{example2024}. All citations will be automatically added to your \texttt{cite.bib} file.

\section{Conclusion}
This template demonstrates basic \LaTeX{} features including mathematical formulas, tables, and lists. The bibliography system is now set up to work seamlessly with Bibby AI's citation management. Feel free to modify it for your needs!

\bibliographystyle{plain}
\bibliography{cite}

\end{document}\documentclass{article}
\usepackage[utf8]{inputenc}
\usepackage{booktabs}
\usepackage{amsmath}
\usepackage{natbib}

\title{My Document}
\author{Your Name}
\date{\today}

\begin{document}

\maketitle

\section{Introduction}
This is a sample document demonstrating various \LaTeX{} features including mathematical equations, tables, and formatting. You can cite papers using the citation search feature, and they will be automatically added to your \texttt{cite.bib} file.

\section{Mathematical Expressions}
Here are some example equations:

\subsection{Quadratic Formula}
The quadratic formula is given by:
\[x = \frac{-b \pm \sqrt{b^2 - 4ac}}{2a}\]

\subsection{Matrix Example}
A sample matrix:
\[
\begin{pmatrix}
a_{11} & a_{12} & a_{13} \\
a_{21} & a_{22} & a_{23} \\
a_{31} & a_{32} & a_{33}
\end{pmatrix}
\]

\section{Tables}
Below is a sample table showing data:

\begin{table}[h]
\centering
\begin{tabular}{@{}lrrr@{}}
\toprule
\textbf{Category} & \textbf{Value 1} & \textbf{Value 2} & \textbf{Total} \\
\midrule
Item A & 10.5 & 20.3 & 30.8 \\
Item B & 15.2 & 25.1 & 40.3 \\
Item C & 22.1 & 12.7 & 34.8 \\
\bottomrule
\end{tabular}
\caption{Sample Data Table}
\label{tab:sample}
\end{table}

\section{Lists}
\subsection{Bullet Points}
\begin{itemize}
  \item First item
  \item Second item
  \item Third item with sub-items:
    \begin{itemize}
      \item Sub-item 1
      \item Sub-item 2
    \end{itemize}
\end{itemize}

\subsection{Numbered List}
\begin{enumerate}
  \item First step
  \item Second step
  \item Third step
\end{enumerate}

\section{Citations and References}
When you use the citation search feature, references will be automatically managed. For example, you might cite a paper like this: \cite{example2024}. All citations will be automatically added to your \texttt{cite.bib} file.

\section{Conclusion}
This template demonstrates basic \LaTeX{} features including mathematical formulas, tables, and lists. The bibliography system is now set up to work seamlessly with Bibby AI's citation management. Feel free to modify it for your needs!

\bibliographystyle{plain}
\bibliography{cite}

\end{document}\documentclass{article}
\usepackage[utf8]{inputenc}
\usepackage{booktabs}
\usepackage{amsmath}
\usepackage{natbib}

\title{My Document}
\author{Your Name}
\date{\today}

\begin{document}

\maketitle

\section{Introduction}
This is a sample document demonstrating various \LaTeX{} features including mathematical equations, tables, and formatting. You can cite papers using the citation search feature, and they will be automatically added to your \texttt{cite.bib} file.

\section{Mathematical Expressions}
Here are some example equations:

\subsection{Quadratic Formula}
The quadratic formula is given by:
\[x = \frac{-b \pm \sqrt{b^2 - 4ac}}{2a}\]

\subsection{Matrix Example}
A sample matrix:
\[
\begin{pmatrix}
a_{11} & a_{12} & a_{13} \\
a_{21} & a_{22} & a_{23} \\
a_{31} & a_{32} & a_{33}
\end{pmatrix}
\]

\section{Tables}
Below is a sample table showing data:

\begin{table}[h]
\centering
\begin{tabular}{@{}lrrr@{}}
\toprule
\textbf{Category} & \textbf{Value 1} & \textbf{Value 2} & \textbf{Total} \\
\midrule
Item A & 10.5 & 20.3 & 30.8 \\
Item B & 15.2 & 25.1 & 40.3 \\
Item C & 22.1 & 12.7 & 34.8 \\
\bottomrule
\end{tabular}
\caption{Sample Data Table}
\label{tab:sample}
\end{table}

\section{Lists}
\subsection{Bullet Points}
\begin{itemize}
  \item First item
  \item Second item
  \item Third item with sub-items:
    \begin{itemize}
      \item Sub-item 1
      \item Sub-item 2
    \end{itemize}
\end{itemize}

\subsection{Numbered List}
\begin{enumerate}
  \item First step
  \item Second step
  \item Third step
\end{enumerate}

\section{Citations and References}
When you use the citation search feature, references will be automatically managed. For example, you might cite a paper like this: \cite{example2024}. All citations will be automatically added to your \texttt{cite.bib} file.

\section{Conclusion}
This template demonstrates basic \LaTeX{} features including mathematical formulas, tables, and lists. The bibliography system is now set up to work seamlessly with Bibby AI's citation management. Feel free to modify it for your needs!

\bibliographystyle{plain}
\bibliography{cite}

\end{document}\documentclass{article}
\usepackage[utf8]{inputenc}
\usepackage{booktabs}
\usepackage{amsmath}
\usepackage{natbib}

\title{My Document}
\author{Your Name}
\date{\today}

\begin{document}

\maketitle

\section{Introduction}
This is a sample document demonstrating various \LaTeX{} features including mathematical equations, tables, and formatting. You can cite papers using the citation search feature, and they will be automatically added to your \texttt{cite.bib} file.

\section{Mathematical Expressions}
Here are some example equations:

\subsection{Quadratic Formula}
The quadratic formula is given by:
\[x = \frac{-b \pm \sqrt{b^2 - 4ac}}{2a}\]

\subsection{Matrix Example}
A sample matrix:
\[
\begin{pmatrix}
a_{11} & a_{12} & a_{13} \\
a_{21} & a_{22} & a_{23} \\
a_{31} & a_{32} & a_{33}
\end{pmatrix}
\]

\section{Tables}
Below is a sample table showing data:

\begin{table}[h]
\centering
\begin{tabular}{@{}lrrr@{}}
\toprule
\textbf{Category} & \textbf{Value 1} & \textbf{Value 2} & \textbf{Total} \\
\midrule
Item A & 10.5 & 20.3 & 30.8 \\
Item B & 15.2 & 25.1 & 40.3 \\
Item C & 22.1 & 12.7 & 34.8 \\
\bottomrule
\end{tabular}
\caption{Sample Data Table}
\label{tab:sample}
\end{table}

\section{Lists}
\subsection{Bullet Points}
\begin{itemize}
  \item First item
  \item Second item
  \item Third item with sub-items:
    \begin{itemize}
      \item Sub-item 1
      \item Sub-item 2
    \end{itemize}
\end{itemize}

\subsection{Numbered List}
\begin{enumerate}
  \item First step
  \item Second step
  \item Third step
\end{enumerate}

\section{Citations and References}
When you use the citation search feature, references will be automatically managed. For example, you might cite a paper like this: \cite{example2024}. All citations will be automatically added to your \texttt{cite.bib} file.

\section{Conclusion}
This template demonstrates basic \LaTeX{} features including mathematical formulas, tables, and lists. The bibliography system is now set up to work seamlessly with Bibby AI's citation management. Feel free to modify it for your needs!

\bibliographystyle{plain}
\bibliography{cite}

\end{document}\documentclass{article}
\usepackage[utf8]{inputenc}
\usepackage{booktabs}
\usepackage{amsmath}
\usepackage{natbib}

\title{My Document}
\author{Your Name}
\date{\today}

\begin{document}

\maketitle

\section{Introduction}
This is a sample document demonstrating various \LaTeX{} features including mathematical equations, tables, and formatting. You can cite papers using the citation search feature, and they will be automatically added to your \texttt{cite.bib} file.

\section{Mathematical Expressions}
Here are some example equations:

\subsection{Quadratic Formula}
The quadratic formula is given by:
\[x = \frac{-b \pm \sqrt{b^2 - 4ac}}{2a}\]

\subsection{Matrix Example}
A sample matrix:
\[
\begin{pmatrix}
a_{11} & a_{12} & a_{13} \\
a_{21} & a_{22} & a_{23} \\
a_{31} & a_{32} & a_{33}
\end{pmatrix}
\]

\section{Tables}
Below is a sample table showing data:

\begin{table}[h]
\centering
\begin{tabular}{@{}lrrr@{}}
\toprule
\textbf{Category} & \textbf{Value 1} & \textbf{Value 2} & \textbf{Total} \\
\midrule
Item A & 10.5 & 20.3 & 30.8 \\
Item B & 15.2 & 25.1 & 40.3 \\
Item C & 22.1 & 12.7 & 34.8 \\
\bottomrule
\end{tabular}
\caption{Sample Data Table}
\label{tab:sample}
\end{table}

\section{Lists}
\subsection{Bullet Points}
\begin{itemize}
  \item First item
  \item Second item
  \item Third item with sub-items:
    \begin{itemize}
      \item Sub-item 1
      \item Sub-item 2
    \end{itemize}
\end{itemize}

\subsection{Numbered List}
\begin{enumerate}
  \item First step
  \item Second step
  \item Third step
\end{enumerate}

\section{Citations and References}
When you use the citation search feature, references will be automatically managed. For example, you might cite a paper like this: \cite{example2024}. All citations will be automatically added to your \texttt{cite.bib} file.

\section{Conclusion}
This template demonstrates basic \LaTeX{} features including mathematical formulas, tables, and lists. The bibliography system is now set up to work seamlessly with Bibby AI's citation management. Feel free to modify it for your needs!

\bibliographystyle{plain}
\bibliography{cite}

\end{document}\documentclass{article}
\usepackage[utf8]{inputenc}
\usepackage{booktabs}
\usepackage{amsmath}
\usepackage{natbib}

\title{My Document}
\author{Your Name}
\date{\today}

\begin{document}

\maketitle

\section{Introduction}
This is a sample document demonstrating various \LaTeX{} features including mathematical equations, tables, and formatting. You can cite papers using the citation search feature, and they will be automatically added to your \texttt{cite.bib} file.

\section{Mathematical Expressions}
Here are some example equations:

\subsection{Quadratic Formula}
The quadratic formula is given by:
\[x = \frac{-b \pm \sqrt{b^2 - 4ac}}{2a}\]

\subsection{Matrix Example}
A sample matrix:
\[
\begin{pmatrix}
a_{11} & a_{12} & a_{13} \\
a_{21} & a_{22} & a_{23} \\
a_{31} & a_{32} & a_{33}
\end{pmatrix}
\]

\section{Tables}
Below is a sample table showing data:

\begin{table}[h]
\centering
\begin{tabular}{@{}lrrr@{}}
\toprule
\textbf{Category} & \textbf{Value 1} & \textbf{Value 2} & \textbf{Total} \\
\midrule
Item A & 10.5 & 20.3 & 30.8 \\
Item B & 15.2 & 25.1 & 40.3 \\
Item C & 22.1 & 12.7 & 34.8 \\
\bottomrule
\end{tabular}
\caption{Sample Data Table}
\label{tab:sample}
\end{table}

\section{Lists}
\subsection{Bullet Points}
\begin{itemize}
  \item First item
  \item Second item
  \item Third item with sub-items:
    \begin{itemize}
      \item Sub-item 1
      \item Sub-item 2
    \end{itemize}
\end{itemize}

\subsection{Numbered List}
\begin{enumerate}
  \item First step
  \item Second step
  \item Third step
\end{enumerate}

\section{Citations and References}
When you use the citation search feature, references will be automatically managed. For example, you might cite a paper like this: \cite{example2024}. All citations will be automatically added to your \texttt{cite.bib} file.

\section{Conclusion}
This template demonstrates basic \LaTeX{} features including mathematical formulas, tables, and lists. The bibliography system is now set up to work seamlessly with Bibby AI's citation management. Feel free to modify it for your needs!

\bibliographystyle{plain}
\bibliography{cite}

\end{document}\documentclass{article}
\usepackage[utf8]{inputenc}
\usepackage{booktabs}
\usepackage{amsmath}
\usepackage{natbib}

\title{My Document}
\author{Your Name}
\date{\today}

\begin{document}

\maketitle

\section{Introduction}
This is a sample document demonstrating various \LaTeX{} features including mathematical equations, tables, and formatting. You can cite papers using the citation search feature, and they will be automatically added to your \texttt{cite.bib} file.

\section{Mathematical Expressions}
Here are some example equations:

\subsection{Quadratic Formula}
The quadratic formula is given by:
\[x = \frac{-b \pm \sqrt{b^2 - 4ac}}{2a}\]

\subsection{Matrix Example}
A sample matrix:
\[
\begin{pmatrix}
a_{11} & a_{12} & a_{13} \\
a_{21} & a_{22} & a_{23} \\
a_{31} & a_{32} & a_{33}
\end{pmatrix}
\]

\section{Tables}
Below is a sample table showing data:

\begin{table}[h]
\centering
\begin{tabular}{@{}lrrr@{}}
\toprule
\textbf{Category} & \textbf{Value 1} & \textbf{Value 2} & \textbf{Total} \\
\midrule
Item A & 10.5 & 20.3 & 30.8 \\
Item B & 15.2 & 25.1 & 40.3 \\
Item C & 22.1 & 12.7 & 34.8 \\
\bottomrule
\end{tabular}
\caption{Sample Data Table}
\label{tab:sample}
\end{table}

\section{Lists}
\subsection{Bullet Points}
\begin{itemize}
  \item First item
  \item Second item
  \item Third item with sub-items:
    \begin{itemize}
      \item Sub-item 1
      \item Sub-item 2
    \end{itemize}
\end{itemize}

\subsection{Numbered List}
\begin{enumerate}
  \item First step
  \item Second step
  \item Third step
\end{enumerate}

\section{Citations and References}
When you use the citation search feature, references will be automatically managed. For example, you might cite a paper like this: \cite{example2024}. All citations will be automatically added to your \texttt{cite.bib} file.

\section{Conclusion}
This template demonstrates basic \LaTeX{} features including mathematical formulas, tables, and lists. The bibliography system is now set up to work seamlessly with Bibby AI's citation management. Feel free to modify it for your needs!

\bibliographystyle{plain}
\bibliography{cite}

\end{document}\documentclass{article}
\usepackage[utf8]{inputenc}
\usepackage{booktabs}
\usepackage{amsmath}
\usepackage{natbib}

\title{My Document}
\author{Your Name}
\date{\today}

\begin{document}

\maketitle

\section{Introduction}
This is a sample document demonstrating various \LaTeX{} features including mathematical equations, tables, and formatting. You can cite papers using the citation search feature, and they will be automatically added to your \texttt{cite.bib} file.

\section{Mathematical Expressions}
Here are some example equations:

\subsection{Quadratic Formula}
The quadratic formula is given by:
\[x = \frac{-b \pm \sqrt{b^2 - 4ac}}{2a}\]

\subsection{Matrix Example}
A sample matrix:
\[
\begin{pmatrix}
a_{11} & a_{12} & a_{13} \\
a_{21} & a_{22} & a_{23} \\
a_{31} & a_{32} & a_{33}
\end{pmatrix}
\]

\section{Tables}
Below is a sample table showing data:

\begin{table}[h]
\centering
\begin{tabular}{@{}lrrr@{}}
\toprule
\textbf{Category} & \textbf{Value 1} & \textbf{Value 2} & \textbf{Total} \\
\midrule
Item A & 10.5 & 20.3 & 30.8 \\
Item B & 15.2 & 25.1 & 40.3 \\
Item C & 22.1 & 12.7 & 34.8 \\
\bottomrule
\end{tabular}
\caption{Sample Data Table}
\label{tab:sample}
\end{table}

\section{Lists}
\subsection{Bullet Points}
\begin{itemize}
  \item First item
  \item Second item
  \item Third item with sub-items:
    \begin{itemize}
      \item Sub-item 1
      \item Sub-item 2
    \end{itemize}
\end{itemize}

\subsection{Numbered List}
\begin{enumerate}
  \item First step
  \item Second step
  \item Third step
\end{enumerate}

\section{Citations and References}
When you use the citation search feature, references will be automatically managed. For example, you might cite a paper like this: \cite{example2024}. All citations will be automatically added to your \texttt{cite.bib} file.

\section{Conclusion}
This template demonstrates basic \LaTeX{} features including mathematical formulas, tables, and lists. The bibliography system is now set up to work seamlessly with Bibby AI's citation management. Feel free to modify it for your needs!

\bibliographystyle{plain}
\bibliography{cite}

\end{document}\documentclass{article}
\usepackage[utf8]{inputenc}
\usepackage{booktabs}
\usepackage{amsmath}
\usepackage{natbib}

\title{My Document}
\author{Your Name}
\date{\today}

\begin{document}

\maketitle

\section{Introduction}
This is a sample document demonstrating various \LaTeX{} features including mathematical equations, tables, and formatting. You can cite papers using the citation search feature, and they will be automatically added to your \texttt{cite.bib} file.

\section{Mathematical Expressions}
Here are some example equations:

\subsection{Quadratic Formula}
The quadratic formula is given by:
\[x = \frac{-b \pm \sqrt{b^2 - 4ac}}{2a}\]

\subsection{Matrix Example}
A sample matrix:
\[
\begin{pmatrix}
a_{11} & a_{12} & a_{13} \\
a_{21} & a_{22} & a_{23} \\
a_{31} & a_{32} & a_{33}
\end{pmatrix}
\]

\section{Tables}
Below is a sample table showing data:

\begin{table}[h]
\centering
\begin{tabular}{@{}lrrr@{}}
\toprule
\textbf{Category} & \textbf{Value 1} & \textbf{Value 2} & \textbf{Total} \\
\midrule
Item A & 10.5 & 20.3 & 30.8 \\
Item B & 15.2 & 25.1 & 40.3 \\
Item C & 22.1 & 12.7 & 34.8 \\
\bottomrule
\end{tabular}
\caption{Sample Data Table}
\label{tab:sample}
\end{table}

\section{Lists}
\subsection{Bullet Points}
\begin{itemize}
  \item First item
  \item Second item
  \item Third item with sub-items:
    \begin{itemize}
      \item Sub-item 1
      \item Sub-item 2
    \end{itemize}
\end{itemize}

\subsection{Numbered List}
\begin{enumerate}
  \item First step
  \item Second step
  \item Third step
\end{enumerate}

\section{Citations and References}
When you use the citation search feature, references will be automatically managed. For example, you might cite a paper like this: \cite{example2024}. All citations will be automatically added to your \texttt{cite.bib} file.

\section{Conclusion}
This template demonstrates basic \LaTeX{} features including mathematical formulas, tables, and lists. The bibliography system is now set up to work seamlessly with Bibby AI's citation management. Feel free to modify it for your needs!

\bibliographystyle{plain}
\bibliography{cite}

\end{document}\documentclass{article}
\usepackage[utf8]{inputenc}
\usepackage{booktabs}
\usepackage{amsmath}
\usepackage{natbib}

\title{My Document}
\author{Your Name}
\date{\today}

\begin{document}

\maketitle

\section{Introduction}
This is a sample document demonstrating various \LaTeX{} features including mathematical equations, tables, and formatting. You can cite papers using the citation search feature, and they will be automatically added to your \texttt{cite.bib} file.

\section{Mathematical Expressions}
Here are some example equations:

\subsection{Quadratic Formula}
The quadratic formula is given by:
\[x = \frac{-b \pm \sqrt{b^2 - 4ac}}{2a}\]

\subsection{Matrix Example}
A sample matrix:
\[
\begin{pmatrix}
a_{11} & a_{12} & a_{13} \\
a_{21} & a_{22} & a_{23} \\
a_{31} & a_{32} & a_{33}
\end{pmatrix}
\]

\section{Tables}
Below is a sample table showing data:

\begin{table}[h]
\centering
\begin{tabular}{@{}lrrr@{}}
\toprule
\textbf{Category} & \textbf{Value 1} & \textbf{Value 2} & \textbf{Total} \\
\midrule
Item A & 10.5 & 20.3 & 30.8 \\
Item B & 15.2 & 25.1 & 40.3 \\
Item C & 22.1 & 12.7 & 34.8 \\
\bottomrule
\end{tabular}
\caption{Sample Data Table}
\label{tab:sample}
\end{table}

\section{Lists}
\subsection{Bullet Points}
\begin{itemize}
  \item First item
  \item Second item
  \item Third item with sub-items:
    \begin{itemize}
      \item Sub-item 1
      \item Sub-item 2
    \end{itemize}
\end{itemize}

\subsection{Numbered List}
\begin{enumerate}
  \item First step
  \item Second step
  \item Third step
\end{enumerate}

\section{Citations and References}
When you use the citation search feature, references will be automatically managed. For example, you might cite a paper like this: \cite{example2024}. All citations will be automatically added to your \texttt{cite.bib} file.

\section{Conclusion}
This template demonstrates basic \LaTeX{} features including mathematical formulas, tables, and lists. The bibliography system is now set up to work seamlessly with Bibby AI's citation management. Feel free to modify it for your needs!

\bibliographystyle{plain}
\bibliography{cite}

\end{document}\documentclass{article}
\usepackage[utf8]{inputenc}
\usepackage{booktabs}
\usepackage{amsmath}
\usepackage{natbib}

\title{My Document}
\author{Your Name}
\date{\today}

\begin{document}

\maketitle

\section{Introduction}
This is a sample document demonstrating various \LaTeX{} features including mathematical equations, tables, and formatting. You can cite papers using the citation search feature, and they will be automatically added to your \texttt{cite.bib} file.

\section{Mathematical Expressions}
Here are some example equations:

\subsection{Quadratic Formula}
The quadratic formula is given by:
\[x = \frac{-b \pm \sqrt{b^2 - 4ac}}{2a}\]

\subsection{Matrix Example}
A sample matrix:
\[
\begin{pmatrix}
a_{11} & a_{12} & a_{13} \\
a_{21} & a_{22} & a_{23} \\
a_{31} & a_{32} & a_{33}
\end{pmatrix}
\]

\section{Tables}
Below is a sample table showing data:

\begin{table}[h]
\centering
\begin{tabular}{@{}lrrr@{}}
\toprule
\textbf{Category} & \textbf{Value 1} & \textbf{Value 2} & \textbf{Total} \\
\midrule
Item A & 10.5 & 20.3 & 30.8 \\
Item B & 15.2 & 25.1 & 40.3 \\
Item C & 22.1 & 12.7 & 34.8 \\
\bottomrule
\end{tabular}
\caption{Sample Data Table}
\label{tab:sample}
\end{table}

\section{Lists}
\subsection{Bullet Points}
\begin{itemize}
  \item First item
  \item Second item
  \item Third item with sub-items:
    \begin{itemize}
      \item Sub-item 1
      \item Sub-item 2
    \end{itemize}
\end{itemize}

\subsection{Numbered List}
\begin{enumerate}
  \item First step
  \item Second step
  \item Third step
\end{enumerate}

\section{Citations and References}
When you use the citation search feature, references will be automatically managed. For example, you might cite a paper like this: \cite{example2024}. All citations will be automatically added to your \texttt{cite.bib} file.

\section{Conclusion}
This template demonstrates basic \LaTeX{} features including mathematical formulas, tables, and lists. The bibliography system is now set up to work seamlessly with Bibby AI's citation management. Feel free to modify it for your needs!

\bibliographystyle{plain}
\bibliography{cite}

\end{document}\documentclass{article}
\usepackage[utf8]{inputenc}
\usepackage{booktabs}
\usepackage{amsmath}
\usepackage{natbib}

\title{My Document}
\author{Your Name}
\date{\today}

\begin{document}

\maketitle

\section{Introduction}
This is a sample document demonstrating various \LaTeX{} features including mathematical equations, tables, and formatting. You can cite papers using the citation search feature, and they will be automatically added to your \texttt{cite.bib} file.

\section{Mathematical Expressions}
Here are some example equations:

\subsection{Quadratic Formula}
The quadratic formula is given by:
\[x = \frac{-b \pm \sqrt{b^2 - 4ac}}{2a}\]

\subsection{Matrix Example}
A sample matrix:
\[
\begin{pmatrix}
a_{11} & a_{12} & a_{13} \\
a_{21} & a_{22} & a_{23} \\
a_{31} & a_{32} & a_{33}
\end{pmatrix}
\]

\section{Tables}
Below is a sample table showing data:

\begin{table}[h]
\centering
\begin{tabular}{@{}lrrr@{}}
\toprule
\textbf{Category} & \textbf{Value 1} & \textbf{Value 2} & \textbf{Total} \\
\midrule
Item A & 10.5 & 20.3 & 30.8 \\
Item B & 15.2 & 25.1 & 40.3 \\
Item C & 22.1 & 12.7 & 34.8 \\
\bottomrule
\end{tabular}
\caption{Sample Data Table}
\label{tab:sample}
\end{table}

\section{Lists}
\subsection{Bullet Points}
\begin{itemize}
  \item First item
  \item Second item
  \item Third item with sub-items:
    \begin{itemize}
      \item Sub-item 1
      \item Sub-item 2
    \end{itemize}
\end{itemize}

\subsection{Numbered List}
\begin{enumerate}
  \item First step
  \item Second step
  \item Third step
\end{enumerate}

\section{Citations and References}
When you use the citation search feature, references will be automatically managed. For example, you might cite a paper like this: \cite{example2024}. All citations will be automatically added to your \texttt{cite.bib} file.

\section{Conclusion}
This template demonstrates basic \LaTeX{} features including mathematical formulas, tables, and lists. The bibliography system is now set up to work seamlessly with Bibby AI's citation management. Feel free to modify it for your needs!

\bibliographystyle{plain}
\bibliography{cite}

\end{document}\documentclass{article}
\usepackage[utf8]{inputenc}
\usepackage{booktabs}
\usepackage{amsmath}
\usepackage{natbib}

\title{My Document}
\author{Your Name}
\date{\today}

\begin{document}

\maketitle

\section{Introduction}
This is a sample document demonstrating various \LaTeX{} features including mathematical equations, tables, and formatting. You can cite papers using the citation search feature, and they will be automatically added to your \texttt{cite.bib} file.

\section{Mathematical Expressions}
Here are some example equations:

\subsection{Quadratic Formula}
The quadratic formula is given by:
\[x = \frac{-b \pm \sqrt{b^2 - 4ac}}{2a}\]

\subsection{Matrix Example}
A sample matrix:
\[
\begin{pmatrix}
a_{11} & a_{12} & a_{13} \\
a_{21} & a_{22} & a_{23} \\
a_{31} & a_{32} & a_{33}
\end{pmatrix}
\]

\section{Tables}
Below is a sample table showing data:

\begin{table}[h]
\centering
\begin{tabular}{@{}lrrr@{}}
\toprule
\textbf{Category} & \textbf{Value 1} & \textbf{Value 2} & \textbf{Total} \\
\midrule
Item A & 10.5 & 20.3 & 30.8 \\
Item B & 15.2 & 25.1 & 40.3 \\
Item C & 22.1 & 12.7 & 34.8 \\
\bottomrule
\end{tabular}
\caption{Sample Data Table}
\label{tab:sample}
\end{table}

\section{Lists}
\subsection{Bullet Points}
\begin{itemize}
  \item First item
  \item Second item
  \item Third item with sub-items:
    \begin{itemize}
      \item Sub-item 1
      \item Sub-item 2
    \end{itemize}
\end{itemize}

\subsection{Numbered List}
\begin{enumerate}
  \item First step
  \item Second step
  \item Third step
\end{enumerate}

\section{Citations and References}
When you use the citation search feature, references will be automatically managed. For example, you might cite a paper like this: \cite{example2024}. All citations will be automatically added to your \texttt{cite.bib} file.

\section{Conclusion}
This template demonstrates basic \LaTeX{} features including mathematical formulas, tables, and lists. The bibliography system is now set up to work seamlessly with Bibby AI's citation management. Feel free to modify it for your needs!

\bibliographystyle{plain}
\bibliography{cite}

\end{document}\documentclass{article}
\usepackage[utf8]{inputenc}
\usepackage{booktabs}
\usepackage{amsmath}
\usepackage{natbib}

\title{My Document}
\author{Your Name}
\date{\today}

\begin{document}

\maketitle

\section{Introduction}
This is a sample document demonstrating various \LaTeX{} features including mathematical equations, tables, and formatting. You can cite papers using the citation search feature, and they will be automatically added to your \texttt{cite.bib} file.

\section{Mathematical Expressions}
Here are some example equations:

\subsection{Quadratic Formula}
The quadratic formula is given by:
\[x = \frac{-b \pm \sqrt{b^2 - 4ac}}{2a}\]

\subsection{Matrix Example}
A sample matrix:
\[
\begin{pmatrix}
a_{11} & a_{12} & a_{13} \\
a_{21} & a_{22} & a_{23} \\
a_{31} & a_{32} & a_{33}
\end{pmatrix}
\]

\section{Tables}
Below is a sample table showing data:

\begin{table}[h]
\centering
\begin{tabular}{@{}lrrr@{}}
\toprule
\textbf{Category} & \textbf{Value 1} & \textbf{Value 2} & \textbf{Total} \\
\midrule
Item A & 10.5 & 20.3 & 30.8 \\
Item B & 15.2 & 25.1 & 40.3 \\
Item C & 22.1 & 12.7 & 34.8 \\
\bottomrule
\end{tabular}
\caption{Sample Data Table}
\label{tab:sample}
\end{table}

\section{Lists}
\subsection{Bullet Points}
\begin{itemize}
  \item First item
  \item Second item
  \item Third item with sub-items:
    \begin{itemize}
      \item Sub-item 1
      \item Sub-item 2
    \end{itemize}
\end{itemize}

\subsection{Numbered List}
\begin{enumerate}
  \item First step
  \item Second step
  \item Third step
\end{enumerate}

\section{Citations and References}
When you use the citation search feature, references will be automatically managed. For example, you might cite a paper like this: \cite{example2024}. All citations will be automatically added to your \texttt{cite.bib} file.

\section{Conclusion}
This template demonstrates basic \LaTeX{} features including mathematical formulas, tables, and lists. The bibliography system is now set up to work seamlessly with Bibby AI's citation management. Feel free to modify it for your needs!

\bibliographystyle{plain}
\bibliography{cite}

\end{document}\documentclass{article}
\usepackage[utf8]{inputenc}
\usepackage{booktabs}
\usepackage{amsmath}
\usepackage{natbib}

\title{My Document}
\author{Your Name}
\date{\today}

\begin{document}

\maketitle

\section{Introduction}
This is a sample document demonstrating various \LaTeX{} features including mathematical equations, tables, and formatting. You can cite papers using the citation search feature, and they will be automatically added to your \texttt{cite.bib} file.

\section{Mathematical Expressions}
Here are some example equations:

\subsection{Quadratic Formula}
The quadratic formula is given by:
\[x = \frac{-b \pm \sqrt{b^2 - 4ac}}{2a}\]

\subsection{Matrix Example}
A sample matrix:
\[
\begin{pmatrix}
a_{11} & a_{12} & a_{13} \\
a_{21} & a_{22} & a_{23} \\
a_{31} & a_{32} & a_{33}
\end{pmatrix}
\]

\section{Tables}
Below is a sample table showing data:

\begin{table}[h]
\centering
\begin{tabular}{@{}lrrr@{}}
\toprule
\textbf{Category} & \textbf{Value 1} & \textbf{Value 2} & \textbf{Total} \\
\midrule
Item A & 10.5 & 20.3 & 30.8 \\
Item B & 15.2 & 25.1 & 40.3 \\
Item C & 22.1 & 12.7 & 34.8 \\
\bottomrule
\end{tabular}
\caption{Sample Data Table}
\label{tab:sample}
\end{table}

\section{Lists}
\subsection{Bullet Points}
\begin{itemize}
  \item First item
  \item Second item
  \item Third item with sub-items:
    \begin{itemize}
      \item Sub-item 1
      \item Sub-item 2
    \end{itemize}
\end{itemize}

\subsection{Numbered List}
\begin{enumerate}
  \item First step
  \item Second step
  \item Third step
\end{enumerate}

\section{Citations and References}
When you use the citation search feature, references will be automatically managed. For example, you might cite a paper like this: \cite{example2024}. All citations will be automatically added to your \texttt{cite.bib} file.

\section{Conclusion}
This template demonstrates basic \LaTeX{} features including mathematical formulas, tables, and lists. The bibliography system is now set up to work seamlessly with Bibby AI's citation management. Feel free to modify it for your needs!

\bibliographystyle{plain}
\bibliography{cite}

\end{document}\documentclass{article}
\usepackage[utf8]{inputenc}
\usepackage{booktabs}
\usepackage{amsmath}
\usepackage{natbib}

\title{My Document}
\author{Your Name}
\date{\today}

\begin{document}

\maketitle

\section{Introduction}
This is a sample document demonstrating various \LaTeX{} features including mathematical equations, tables, and formatting. You can cite papers using the citation search feature, and they will be automatically added to your \texttt{cite.bib} file.

\section{Mathematical Expressions}
Here are some example equations:

\subsection{Quadratic Formula}
The quadratic formula is given by:
\[x = \frac{-b \pm \sqrt{b^2 - 4ac}}{2a}\]

\subsection{Matrix Example}
A sample matrix:
\[
\begin{pmatrix}
a_{11} & a_{12} & a_{13} \\
a_{21} & a_{22} & a_{23} \\
a_{31} & a_{32} & a_{33}
\end{pmatrix}
\]

\section{Tables}
Below is a sample table showing data:

\begin{table}[h]
\centering
\begin{tabular}{@{}lrrr@{}}
\toprule
\textbf{Category} & \textbf{Value 1} & \textbf{Value 2} & \textbf{Total} \\
\midrule
Item A & 10.5 & 20.3 & 30.8 \\
Item B & 15.2 & 25.1 & 40.3 \\
Item C & 22.1 & 12.7 & 34.8 \\
\bottomrule
\end{tabular}
\caption{Sample Data Table}
\label{tab:sample}
\end{table}

\section{Lists}
\subsection{Bullet Points}
\begin{itemize}
  \item First item
  \item Second item
  \item Third item with sub-items:
    \begin{itemize}
      \item Sub-item 1
      \item Sub-item 2
    \end{itemize}
\end{itemize}

\subsection{Numbered List}
\begin{enumerate}
  \item First step
  \item Second step
  \item Third step
\end{enumerate}

\section{Citations and References}
When you use the citation search feature, references will be automatically managed. For example, you might cite a paper like this: \cite{example2024}. All citations will be automatically added to your \texttt{cite.bib} file.

\section{Conclusion}
This template demonstrates basic \LaTeX{} features including mathematical formulas, tables, and lists. The bibliography system is now set up to work seamlessly with Bibby AI's citation management. Feel free to modify it for your needs!

\bibliographystyle{plain}
\bibliography{cite}

\end{document}\documentclass{article}
\usepackage[utf8]{inputenc}
\usepackage{booktabs}
\usepackage{amsmath}
\usepackage{natbib}

\title{My Document}
\author{Your Name}
\date{\today}

\begin{document}

\maketitle

\section{Introduction}
This is a sample document demonstrating various \LaTeX{} features including mathematical equations, tables, and formatting. You can cite papers using the citation search feature, and they will be automatically added to your \texttt{cite.bib} file.

\section{Mathematical Expressions}
Here are some example equations:

\subsection{Quadratic Formula}
The quadratic formula is given by:
\[x = \frac{-b \pm \sqrt{b^2 - 4ac}}{2a}\]

\subsection{Matrix Example}
A sample matrix:
\[
\begin{pmatrix}
a_{11} & a_{12} & a_{13} \\
a_{21} & a_{22} & a_{23} \\
a_{31} & a_{32} & a_{33}
\end{pmatrix}
\]

\section{Tables}
Below is a sample table showing data:

\begin{table}[h]
\centering
\begin{tabular}{@{}lrrr@{}}
\toprule
\textbf{Category} & \textbf{Value 1} & \textbf{Value 2} & \textbf{Total} \\
\midrule
Item A & 10.5 & 20.3 & 30.8 \\
Item B & 15.2 & 25.1 & 40.3 \\
Item C & 22.1 & 12.7 & 34.8 \\
\bottomrule
\end{tabular}
\caption{Sample Data Table}
\label{tab:sample}
\end{table}

\section{Lists}
\subsection{Bullet Points}
\begin{itemize}
  \item First item
  \item Second item
  \item Third item with sub-items:
    \begin{itemize}
      \item Sub-item 1
      \item Sub-item 2
    \end{itemize}
\end{itemize}

\subsection{Numbered List}
\begin{enumerate}
  \item First step
  \item Second step
  \item Third step
\end{enumerate}

\section{Citations and References}
When you use the citation search feature, references will be automatically managed. For example, you might cite a paper like this: \cite{example2024}. All citations will be automatically added to your \texttt{cite.bib} file.

\section{Conclusion}
This template demonstrates basic \LaTeX{} features including mathematical formulas, tables, and lists. The bibliography system is now set up to work seamlessly with Bibby AI's citation management. Feel free to modify it for your needs!

\bibliographystyle{plain}
\bibliography{cite}

\end{document}\documentclass{article}
\usepackage[utf8]{inputenc}
\usepackage{booktabs}
\usepackage{amsmath}
\usepackage{natbib}

\title{My Document}
\author{Your Name}
\date{\today}

\begin{document}

\maketitle

\section{Introduction}
This is a sample document demonstrating various \LaTeX{} features including mathematical equations, tables, and formatting. You can cite papers using the citation search feature, and they will be automatically added to your \texttt{cite.bib} file.

\section{Mathematical Expressions}
Here are some example equations:

\subsection{Quadratic Formula}
The quadratic formula is given by:
\[x = \frac{-b \pm \sqrt{b^2 - 4ac}}{2a}\]

\subsection{Matrix Example}
A sample matrix:
\[
\begin{pmatrix}
a_{11} & a_{12} & a_{13} \\
a_{21} & a_{22} & a_{23} \\
a_{31} & a_{32} & a_{33}
\end{pmatrix}
\]

\section{Tables}
Below is a sample table showing data:

\begin{table}[h]
\centering
\begin{tabular}{@{}lrrr@{}}
\toprule
\textbf{Category} & \textbf{Value 1} & \textbf{Value 2} & \textbf{Total} \\
\midrule
Item A & 10.5 & 20.3 & 30.8 \\
Item B & 15.2 & 25.1 & 40.3 \\
Item C & 22.1 & 12.7 & 34.8 \\
\bottomrule
\end{tabular}
\caption{Sample Data Table}
\label{tab:sample}
\end{table}

\section{Lists}
\subsection{Bullet Points}
\begin{itemize}
  \item First item
  \item Second item
  \item Third item with sub-items:
    \begin{itemize}
      \item Sub-item 1
      \item Sub-item 2
    \end{itemize}
\end{itemize}

\subsection{Numbered List}
\begin{enumerate}
  \item First step
  \item Second step
  \item Third step
\end{enumerate}

\section{Citations and References}
When you use the citation search feature, references will be automatically managed. For example, you might cite a paper like this: \cite{example2024}. All citations will be automatically added to your \texttt{cite.bib} file.

\section{Conclusion}
This template demonstrates basic \LaTeX{} features including mathematical formulas, tables, and lists. The bibliography system is now set up to work seamlessly with Bibby AI's citation management. Feel free to modify it for your needs!

\bibliographystyle{plain}
\bibliography{cite}

\end{document}\documentclass{article}
\usepackage[utf8]{inputenc}
\usepackage{booktabs}
\usepackage{amsmath}
\usepackage{natbib}

\title{My Document}
\author{Your Name}
\date{\today}

\begin{document}

\maketitle

\section{Introduction}
This is a sample document demonstrating various \LaTeX{} features including mathematical equations, tables, and formatting. You can cite papers using the citation search feature, and they will be automatically added to your \texttt{cite.bib} file.

\section{Mathematical Expressions}
Here are some example equations:

\subsection{Quadratic Formula}
The quadratic formula is given by:
\[x = \frac{-b \pm \sqrt{b^2 - 4ac}}{2a}\]

\subsection{Matrix Example}
A sample matrix:
\[
\begin{pmatrix}
a_{11} & a_{12} & a_{13} \\
a_{21} & a_{22} & a_{23} \\
a_{31} & a_{32} & a_{33}
\end{pmatrix}
\]

\section{Tables}
Below is a sample table showing data:

\begin{table}[h]
\centering
\begin{tabular}{@{}lrrr@{}}
\toprule
\textbf{Category} & \textbf{Value 1} & \textbf{Value 2} & \textbf{Total} \\
\midrule
Item A & 10.5 & 20.3 & 30.8 \\
Item B & 15.2 & 25.1 & 40.3 \\
Item C & 22.1 & 12.7 & 34.8 \\
\bottomrule
\end{tabular}
\caption{Sample Data Table}
\label{tab:sample}
\end{table}

\section{Lists}
\subsection{Bullet Points}
\begin{itemize}
  \item First item
  \item Second item
  \item Third item with sub-items:
    \begin{itemize}
      \item Sub-item 1
      \item Sub-item 2
    \end{itemize}
\end{itemize}

\subsection{Numbered List}
\begin{enumerate}
  \item First step
  \item Second step
  \item Third step
\end{enumerate}

\section{Citations and References}
When you use the citation search feature, references will be automatically managed. For example, you might cite a paper like this: \cite{example2024}. All citations will be automatically added to your \texttt{cite.bib} file.

\section{Conclusion}
This template demonstrates basic \LaTeX{} features including mathematical formulas, tables, and lists. The bibliography system is now set up to work seamlessly with Bibby AI's citation management. Feel free to modify it for your needs!

\bibliographystyle{plain}
\bibliography{cite}

\end{document}\documentclass{article}
\usepackage[utf8]{inputenc}
\usepackage{booktabs}
\usepackage{amsmath}
\usepackage{natbib}

\title{My Document}
\author{Your Name}
\date{\today}

\begin{document}

\maketitle

\section{Introduction}
This is a sample document demonstrating various \LaTeX{} features including mathematical equations, tables, and formatting. You can cite papers using the citation search feature, and they will be automatically added to your \texttt{cite.bib} file.

\section{Mathematical Expressions}
Here are some example equations:

\subsection{Quadratic Formula}
The quadratic formula is given by:
\[x = \frac{-b \pm \sqrt{b^2 - 4ac}}{2a}\]

\subsection{Matrix Example}
A sample matrix:
\[
\begin{pmatrix}
a_{11} & a_{12} & a_{13} \\
a_{21} & a_{22} & a_{23} \\
a_{31} & a_{32} & a_{33}
\end{pmatrix}
\]

\section{Tables}
Below is a sample table showing data:

\begin{table}[h]
\centering
\begin{tabular}{@{}lrrr@{}}
\toprule
\textbf{Category} & \textbf{Value 1} & \textbf{Value 2} & \textbf{Total} \\
\midrule
Item A & 10.5 & 20.3 & 30.8 \\
Item B & 15.2 & 25.1 & 40.3 \\
Item C & 22.1 & 12.7 & 34.8 \\
\bottomrule
\end{tabular}
\caption{Sample Data Table}
\label{tab:sample}
\end{table}

\section{Lists}
\subsection{Bullet Points}
\begin{itemize}
  \item First item
  \item Second item
  \item Third item with sub-items:
    \begin{itemize}
      \item Sub-item 1
      \item Sub-item 2
    \end{itemize}
\end{itemize}

\subsection{Numbered List}
\begin{enumerate}
  \item First step
  \item Second step
  \item Third step
\end{enumerate}

\section{Citations and References}
When you use the citation search feature, references will be automatically managed. For example, you might cite a paper like this: \cite{example2024}. All citations will be automatically added to your \texttt{cite.bib} file.

\section{Conclusion}
This template demonstrates basic \LaTeX{} features including mathematical formulas, tables, and lists. The bibliography system is now set up to work seamlessly with Bibby AI's citation management. Feel free to modify it for your needs!

\bibliographystyle{plain}
\bibliography{cite}

\end{document}\documentclass{article}
\usepackage[utf8]{inputenc}
\usepackage{booktabs}
\usepackage{amsmath}
\usepackage{natbib}

\title{My Document}
\author{Your Name}
\date{\today}

\begin{document}

\maketitle

\section{Introduction}
This is a sample document demonstrating various \LaTeX{} features including mathematical equations, tables, and formatting. You can cite papers using the citation search feature, and they will be automatically added to your \texttt{cite.bib} file.

\section{Mathematical Expressions}
Here are some example equations:

\subsection{Quadratic Formula}
The quadratic formula is given by:
\[x = \frac{-b \pm \sqrt{b^2 - 4ac}}{2a}\]

\subsection{Matrix Example}
A sample matrix:
\[
\begin{pmatrix}
a_{11} & a_{12} & a_{13} \\
a_{21} & a_{22} & a_{23} \\
a_{31} & a_{32} & a_{33}
\end{pmatrix}
\]

\section{Tables}
Below is a sample table showing data:

\begin{table}[h]
\centering
\begin{tabular}{@{}lrrr@{}}
\toprule
\textbf{Category} & \textbf{Value 1} & \textbf{Value 2} & \textbf{Total} \\
\midrule
Item A & 10.5 & 20.3 & 30.8 \\
Item B & 15.2 & 25.1 & 40.3 \\
Item C & 22.1 & 12.7 & 34.8 \\
\bottomrule
\end{tabular}
\caption{Sample Data Table}
\label{tab:sample}
\end{table}

\section{Lists}
\subsection{Bullet Points}
\begin{itemize}
  \item First item
  \item Second item
  \item Third item with sub-items:
    \begin{itemize}
      \item Sub-item 1
      \item Sub-item 2
    \end{itemize}
\end{itemize}

\subsection{Numbered List}
\begin{enumerate}
  \item First step
  \item Second step
  \item Third step
\end{enumerate}

\section{Citations and References}
When you use the citation search feature, references will be automatically managed. For example, you might cite a paper like this: \cite{example2024}. All citations will be automatically added to your \texttt{cite.bib} file.

\section{Conclusion}
This template demonstrates basic \LaTeX{} features including mathematical formulas, tables, and lists. The bibliography system is now set up to work seamlessly with Bibby AI's citation management. Feel free to modify it for your needs!

\bibliographystyle{plain}
\bibliography{cite}

\end{document}\documentclass{article}
\usepackage[utf8]{inputenc}
\usepackage{booktabs}
\usepackage{amsmath}
\usepackage{natbib}

\title{My Document}
\author{Your Name}
\date{\today}

\begin{document}

\maketitle

\section{Introduction}
This is a sample document demonstrating various \LaTeX{} features including mathematical equations, tables, and formatting. You can cite papers using the citation search feature, and they will be automatically added to your \texttt{cite.bib} file.

\section{Mathematical Expressions}
Here are some example equations:

\subsection{Quadratic Formula}
The quadratic formula is given by:
\[x = \frac{-b \pm \sqrt{b^2 - 4ac}}{2a}\]

\subsection{Matrix Example}
A sample matrix:
\[
\begin{pmatrix}
a_{11} & a_{12} & a_{13} \\
a_{21} & a_{22} & a_{23} \\
a_{31} & a_{32} & a_{33}
\end{pmatrix}
\]

\section{Tables}
Below is a sample table showing data:

\begin{table}[h]
\centering
\begin{tabular}{@{}lrrr@{}}
\toprule
\textbf{Category} & \textbf{Value 1} & \textbf{Value 2} & \textbf{Total} \\
\midrule
Item A & 10.5 & 20.3 & 30.8 \\
Item B & 15.2 & 25.1 & 40.3 \\
Item C & 22.1 & 12.7 & 34.8 \\
\bottomrule
\end{tabular}
\caption{Sample Data Table}
\label{tab:sample}
\end{table}

\section{Lists}
\subsection{Bullet Points}
\begin{itemize}
  \item First item
  \item Second item
  \item Third item with sub-items:
    \begin{itemize}
      \item Sub-item 1
      \item Sub-item 2
    \end{itemize}
\end{itemize}

\subsection{Numbered List}
\begin{enumerate}
  \item First step
  \item Second step
  \item Third step
\end{enumerate}

\section{Citations and References}
When you use the citation search feature, references will be automatically managed. For example, you might cite a paper like this: \cite{example2024}. All citations will be automatically added to your \texttt{cite.bib} file.

\section{Conclusion}
This template demonstrates basic \LaTeX{} features including mathematical formulas, tables, and lists. The bibliography system is now set up to work seamlessly with Bibby AI's citation management. Feel free to modify it for your needs!

\bibliographystyle{plain}
\bibliography{cite}

\end{document}\documentclass{article}
\usepackage[utf8]{inputenc}
\usepackage{booktabs}
\usepackage{amsmath}
\usepackage{natbib}

\title{My Document}
\author{Your Name}
\date{\today}

\begin{document}

\maketitle

\section{Introduction}
This is a sample document demonstrating various \LaTeX{} features including mathematical equations, tables, and formatting. You can cite papers using the citation search feature, and they will be automatically added to your \texttt{cite.bib} file.

\section{Mathematical Expressions}
Here are some example equations:

\subsection{Quadratic Formula}
The quadratic formula is given by:
\[x = \frac{-b \pm \sqrt{b^2 - 4ac}}{2a}\]

\subsection{Matrix Example}
A sample matrix:
\[
\begin{pmatrix}
a_{11} & a_{12} & a_{13} \\
a_{21} & a_{22} & a_{23} \\
a_{31} & a_{32} & a_{33}
\end{pmatrix}
\]

\section{Tables}
Below is a sample table showing data:

\begin{table}[h]
\centering
\begin{tabular}{@{}lrrr@{}}
\toprule
\textbf{Category} & \textbf{Value 1} & \textbf{Value 2} & \textbf{Total} \\
\midrule
Item A & 10.5 & 20.3 & 30.8 \\
Item B & 15.2 & 25.1 & 40.3 \\
Item C & 22.1 & 12.7 & 34.8 \\
\bottomrule
\end{tabular}
\caption{Sample Data Table}
\label{tab:sample}
\end{table}

\section{Lists}
\subsection{Bullet Points}
\begin{itemize}
  \item First item
  \item Second item
  \item Third item with sub-items:
    \begin{itemize}
      \item Sub-item 1
      \item Sub-item 2
    \end{itemize}
\end{itemize}

\subsection{Numbered List}
\begin{enumerate}
  \item First step
  \item Second step
  \item Third step
\end{enumerate}

\section{Citations and References}
When you use the citation search feature, references will be automatically managed. For example, you might cite a paper like this: \cite{example2024}. All citations will be automatically added to your \texttt{cite.bib} file.

\section{Conclusion}
This template demonstrates basic \LaTeX{} features including mathematical formulas, tables, and lists. The bibliography system is now set up to work seamlessly with Bibby AI's citation management. Feel free to modify it for your needs!

\bibliographystyle{plain}
\bibliography{cite}

\end{document}\documentclass{article}
\usepackage[utf8]{inputenc}
\usepackage{booktabs}
\usepackage{amsmath}
\usepackage{natbib}

\title{My Document}
\author{Your Name}
\date{\today}

\begin{document}

\maketitle

\section{Introduction}
This is a sample document demonstrating various \LaTeX{} features including mathematical equations, tables, and formatting. You can cite papers using the citation search feature, and they will be automatically added to your \texttt{cite.bib} file.

\section{Mathematical Expressions}
Here are some example equations:

\subsection{Quadratic Formula}
The quadratic formula is given by:
\[x = \frac{-b \pm \sqrt{b^2 - 4ac}}{2a}\]

\subsection{Matrix Example}
A sample matrix:
\[
\begin{pmatrix}
a_{11} & a_{12} & a_{13} \\
a_{21} & a_{22} & a_{23} \\
a_{31} & a_{32} & a_{33}
\end{pmatrix}
\]

\section{Tables}
Below is a sample table showing data:

\begin{table}[h]
\centering
\begin{tabular}{@{}lrrr@{}}
\toprule
\textbf{Category} & \textbf{Value 1} & \textbf{Value 2} & \textbf{Total} \\
\midrule
Item A & 10.5 & 20.3 & 30.8 \\
Item B & 15.2 & 25.1 & 40.3 \\
Item C & 22.1 & 12.7 & 34.8 \\
\bottomrule
\end{tabular}
\caption{Sample Data Table}
\label{tab:sample}
\end{table}

\section{Lists}
\subsection{Bullet Points}
\begin{itemize}
  \item First item
  \item Second item
  \item Third item with sub-items:
    \begin{itemize}
      \item Sub-item 1
      \item Sub-item 2
    \end{itemize}
\end{itemize}

\subsection{Numbered List}
\begin{enumerate}
  \item First step
  \item Second step
  \item Third step
\end{enumerate}

\section{Citations and References}
When you use the citation search feature, references will be automatically managed. For example, you might cite a paper like this: \cite{example2024}. All citations will be automatically added to your \texttt{cite.bib} file.

\section{Conclusion}
This template demonstrates basic \LaTeX{} features including mathematical formulas, tables, and lists. The bibliography system is now set up to work seamlessly with Bibby AI's citation management. Feel free to modify it for your needs!

\bibliographystyle{plain}
\bibliography{cite}

\end{document}\documentclass{article}
\usepackage[utf8]{inputenc}
\usepackage{booktabs}
\usepackage{amsmath}
\usepackage{natbib}

\title{My Document}
\author{Your Name}
\date{\today}

\begin{document}

\maketitle

\section{Introduction}
This is a sample document demonstrating various \LaTeX{} features including mathematical equations, tables, and formatting. You can cite papers using the citation search feature, and they will be automatically added to your \texttt{cite.bib} file.

\section{Mathematical Expressions}
Here are some example equations:

\subsection{Quadratic Formula}
The quadratic formula is given by:
\[x = \frac{-b \pm \sqrt{b^2 - 4ac}}{2a}\]

\subsection{Matrix Example}
A sample matrix:
\[
\begin{pmatrix}
a_{11} & a_{12} & a_{13} \\
a_{21} & a_{22} & a_{23} \\
a_{31} & a_{32} & a_{33}
\end{pmatrix}
\]

\section{Tables}
Below is a sample table showing data:

\begin{table}[h]
\centering
\begin{tabular}{@{}lrrr@{}}
\toprule
\textbf{Category} & \textbf{Value 1} & \textbf{Value 2} & \textbf{Total} \\
\midrule
Item A & 10.5 & 20.3 & 30.8 \\
Item B & 15.2 & 25.1 & 40.3 \\
Item C & 22.1 & 12.7 & 34.8 \\
\bottomrule
\end{tabular}
\caption{Sample Data Table}
\label{tab:sample}
\end{table}

\section{Lists}
\subsection{Bullet Points}
\begin{itemize}
  \item First item
  \item Second item
  \item Third item with sub-items:
    \begin{itemize}
      \item Sub-item 1
      \item Sub-item 2
    \end{itemize}
\end{itemize}

\subsection{Numbered List}
\begin{enumerate}
  \item First step
  \item Second step
  \item Third step
\end{enumerate}

\section{Citations and References}
When you use the citation search feature, references will be automatically managed. For example, you might cite a paper like this: \cite{example2024}. All citations will be automatically added to your \texttt{cite.bib} file.

\section{Conclusion}
This template demonstrates basic \LaTeX{} features including mathematical formulas, tables, and lists. The bibliography system is now set up to work seamlessly with Bibby AI's citation management. Feel free to modify it for your needs!

\bibliographystyle{plain}
\bibliography{cite}

\end{document}\documentclass{article}
\usepackage[utf8]{inputenc}
\usepackage{booktabs}
\usepackage{amsmath}
\usepackage{natbib}

\title{My Document}
\author{Your Name}
\date{\today}

\begin{document}

\maketitle

\section{Introduction}
This is a sample document demonstrating various \LaTeX{} features including mathematical equations, tables, and formatting. You can cite papers using the citation search feature, and they will be automatically added to your \texttt{cite.bib} file.

\section{Mathematical Expressions}
Here are some example equations:

\subsection{Quadratic Formula}
The quadratic formula is given by:
\[x = \frac{-b \pm \sqrt{b^2 - 4ac}}{2a}\]

\subsection{Matrix Example}
A sample matrix:
\[
\begin{pmatrix}
a_{11} & a_{12} & a_{13} \\
a_{21} & a_{22} & a_{23} \\
a_{31} & a_{32} & a_{33}
\end{pmatrix}
\]

\section{Tables}
Below is a sample table showing data:

\begin{table}[h]
\centering
\begin{tabular}{@{}lrrr@{}}
\toprule
\textbf{Category} & \textbf{Value 1} & \textbf{Value 2} & \textbf{Total} \\
\midrule
Item A & 10.5 & 20.3 & 30.8 \\
Item B & 15.2 & 25.1 & 40.3 \\
Item C & 22.1 & 12.7 & 34.8 \\
\bottomrule
\end{tabular}
\caption{Sample Data Table}
\label{tab:sample}
\end{table}

\section{Lists}
\subsection{Bullet Points}
\begin{itemize}
  \item First item
  \item Second item
  \item Third item with sub-items:
    \begin{itemize}
      \item Sub-item 1
      \item Sub-item 2
    \end{itemize}
\end{itemize}

\subsection{Numbered List}
\begin{enumerate}
  \item First step
  \item Second step
  \item Third step
\end{enumerate}

\section{Citations and References}
When you use the citation search feature, references will be automatically managed. For example, you might cite a paper like this: \cite{example2024}. All citations will be automatically added to your \texttt{cite.bib} file.

\section{Conclusion}
This template demonstrates basic \LaTeX{} features including mathematical formulas, tables, and lists. The bibliography system is now set up to work seamlessly with Bibby AI's citation management. Feel free to modify it for your needs!

\bibliographystyle{plain}
\bibliography{cite}

\end{document}\documentclass{article}
\usepackage[utf8]{inputenc}
\usepackage{booktabs}
\usepackage{amsmath}
\usepackage{natbib}

\title{My Document}
\author{Your Name}
\date{\today}

\begin{document}

\maketitle

\section{Introduction}
This is a sample document demonstrating various \LaTeX{} features including mathematical equations, tables, and formatting. You can cite papers using the citation search feature, and they will be automatically added to your \texttt{cite.bib} file.

\section{Mathematical Expressions}
Here are some example equations:

\subsection{Quadratic Formula}
The quadratic formula is given by:
\[x = \frac{-b \pm \sqrt{b^2 - 4ac}}{2a}\]

\subsection{Matrix Example}
A sample matrix:
\[
\begin{pmatrix}
a_{11} & a_{12} & a_{13} \\
a_{21} & a_{22} & a_{23} \\
a_{31} & a_{32} & a_{33}
\end{pmatrix}
\]

\section{Tables}
Below is a sample table showing data:

\begin{table}[h]
\centering
\begin{tabular}{@{}lrrr@{}}
\toprule
\textbf{Category} & \textbf{Value 1} & \textbf{Value 2} & \textbf{Total} \\
\midrule
Item A & 10.5 & 20.3 & 30.8 \\
Item B & 15.2 & 25.1 & 40.3 \\
Item C & 22.1 & 12.7 & 34.8 \\
\bottomrule
\end{tabular}
\caption{Sample Data Table}
\label{tab:sample}
\end{table}

\section{Lists}
\subsection{Bullet Points}
\begin{itemize}
  \item First item
  \item Second item
  \item Third item with sub-items:
    \begin{itemize}
      \item Sub-item 1
      \item Sub-item 2
    \end{itemize}
\end{itemize}

\subsection{Numbered List}
\begin{enumerate}
  \item First step
  \item Second step
  \item Third step
\end{enumerate}

\section{Citations and References}
When you use the citation search feature, references will be automatically managed. For example, you might cite a paper like this: \cite{example2024}. All citations will be automatically added to your \texttt{cite.bib} file.

\section{Conclusion}
This template demonstrates basic \LaTeX{} features including mathematical formulas, tables, and lists. The bibliography system is now set up to work seamlessly with Bibby AI's citation management. Feel free to modify it for your needs!

\bibliographystyle{plain}
\bibliography{cite}

\end{document}\documentclass{article}
\usepackage[utf8]{inputenc}
\usepackage{booktabs}
\usepackage{amsmath}
\usepackage{natbib}

\title{My Document}
\author{Your Name}
\date{\today}

\begin{document}

\maketitle

\section{Introduction}
This is a sample document demonstrating various \LaTeX{} features including mathematical equations, tables, and formatting. You can cite papers using the citation search feature, and they will be automatically added to your \texttt{cite.bib} file.

\section{Mathematical Expressions}
Here are some example equations:

\subsection{Quadratic Formula}
The quadratic formula is given by:
\[x = \frac{-b \pm \sqrt{b^2 - 4ac}}{2a}\]

\subsection{Matrix Example}
A sample matrix:
\[
\begin{pmatrix}
a_{11} & a_{12} & a_{13} \\
a_{21} & a_{22} & a_{23} \\
a_{31} & a_{32} & a_{33}
\end{pmatrix}
\]

\section{Tables}
Below is a sample table showing data:

\begin{table}[h]
\centering
\begin{tabular}{@{}lrrr@{}}
\toprule
\textbf{Category} & \textbf{Value 1} & \textbf{Value 2} & \textbf{Total} \\
\midrule
Item A & 10.5 & 20.3 & 30.8 \\
Item B & 15.2 & 25.1 & 40.3 \\
Item C & 22.1 & 12.7 & 34.8 \\
\bottomrule
\end{tabular}
\caption{Sample Data Table}
\label{tab:sample}
\end{table}

\section{Lists}
\subsection{Bullet Points}
\begin{itemize}
  \item First item
  \item Second item
  \item Third item with sub-items:
    \begin{itemize}
      \item Sub-item 1
      \item Sub-item 2
    \end{itemize}
\end{itemize}

\subsection{Numbered List}
\begin{enumerate}
  \item First step
  \item Second step
  \item Third step
\end{enumerate}

\section{Citations and References}
When you use the citation search feature, references will be automatically managed. For example, you might cite a paper like this: \cite{example2024}. All citations will be automatically added to your \texttt{cite.bib} file.

\section{Conclusion}
This template demonstrates basic \LaTeX{} features including mathematical formulas, tables, and lists. The bibliography system is now set up to work seamlessly with Bibby AI's citation management. Feel free to modify it for your needs!

\bibliographystyle{plain}
\bibliography{cite}

\end{document}\documentclass{article}
\usepackage[utf8]{inputenc}
\usepackage{booktabs}
\usepackage{amsmath}
\usepackage{natbib}

\title{My Document}
\author{Your Name}
\date{\today}

\begin{document}

\maketitle

\section{Introduction}
This is a sample document demonstrating various \LaTeX{} features including mathematical equations, tables, and formatting. You can cite papers using the citation search feature, and they will be automatically added to your \texttt{cite.bib} file.

\section{Mathematical Expressions}
Here are some example equations:

\subsection{Quadratic Formula}
The quadratic formula is given by:
\[x = \frac{-b \pm \sqrt{b^2 - 4ac}}{2a}\]

\subsection{Matrix Example}
A sample matrix:
\[
\begin{pmatrix}
a_{11} & a_{12} & a_{13} \\
a_{21} & a_{22} & a_{23} \\
a_{31} & a_{32} & a_{33}
\end{pmatrix}
\]

\section{Tables}
Below is a sample table showing data:

\begin{table}[h]
\centering
\begin{tabular}{@{}lrrr@{}}
\toprule
\textbf{Category} & \textbf{Value 1} & \textbf{Value 2} & \textbf{Total} \\
\midrule
Item A & 10.5 & 20.3 & 30.8 \\
Item B & 15.2 & 25.1 & 40.3 \\
Item C & 22.1 & 12.7 & 34.8 \\
\bottomrule
\end{tabular}
\caption{Sample Data Table}
\label{tab:sample}
\end{table}

\section{Lists}
\subsection{Bullet Points}
\begin{itemize}
  \item First item
  \item Second item
  \item Third item with sub-items:
    \begin{itemize}
      \item Sub-item 1
      \item Sub-item 2
    \end{itemize}
\end{itemize}

\subsection{Numbered List}
\begin{enumerate}
  \item First step
  \item Second step
  \item Third step
\end{enumerate}

\section{Citations and References}
When you use the citation search feature, references will be automatically managed. For example, you might cite a paper like this: \cite{example2024}. All citations will be automatically added to your \texttt{cite.bib} file.

\section{Conclusion}
This template demonstrates basic \LaTeX{} features including mathematical formulas, tables, and lists. The bibliography system is now set up to work seamlessly with Bibby AI's citation management. Feel free to modify it for your needs!

\bibliographystyle{plain}
\bibliography{cite}

\end{document}\documentclass{article}
\usepackage[utf8]{inputenc}
\usepackage{booktabs}
\usepackage{amsmath}
\usepackage{natbib}

\title{My Document}
\author{Your Name}
\date{\today}

\begin{document}

\maketitle

\section{Introduction}
This is a sample document demonstrating various \LaTeX{} features including mathematical equations, tables, and formatting. You can cite papers using the citation search feature, and they will be automatically added to your \texttt{cite.bib} file.

\section{Mathematical Expressions}
Here are some example equations:

\subsection{Quadratic Formula}
The quadratic formula is given by:
\[x = \frac{-b \pm \sqrt{b^2 - 4ac}}{2a}\]

\subsection{Matrix Example}
A sample matrix:
\[
\begin{pmatrix}
a_{11} & a_{12} & a_{13} \\
a_{21} & a_{22} & a_{23} \\
a_{31} & a_{32} & a_{33}
\end{pmatrix}
\]

\section{Tables}
Below is a sample table showing data:

\begin{table}[h]
\centering
\begin{tabular}{@{}lrrr@{}}
\toprule
\textbf{Category} & \textbf{Value 1} & \textbf{Value 2} & \textbf{Total} \\
\midrule
Item A & 10.5 & 20.3 & 30.8 \\
Item B & 15.2 & 25.1 & 40.3 \\
Item C & 22.1 & 12.7 & 34.8 \\
\bottomrule
\end{tabular}
\caption{Sample Data Table}
\label{tab:sample}
\end{table}

\section{Lists}
\subsection{Bullet Points}
\begin{itemize}
  \item First item
  \item Second item
  \item Third item with sub-items:
    \begin{itemize}
      \item Sub-item 1
      \item Sub-item 2
    \end{itemize}
\end{itemize}

\subsection{Numbered List}
\begin{enumerate}
  \item First step
  \item Second step
  \item Third step
\end{enumerate}

\section{Citations and References}
When you use the citation search feature, references will be automatically managed. For example, you might cite a paper like this: \cite{example2024}. All citations will be automatically added to your \texttt{cite.bib} file.

\section{Conclusion}
This template demonstrates basic \LaTeX{} features including mathematical formulas, tables, and lists. The bibliography system is now set up to work seamlessly with Bibby AI's citation management. Feel free to modify it for your needs!

\bibliographystyle{plain}
\bibliography{cite}

\end{document}\documentclass{article}
\usepackage[utf8]{inputenc}
\usepackage{booktabs}
\usepackage{amsmath}
\usepackage{natbib}

\title{My Document}
\author{Your Name}
\date{\today}

\begin{document}

\maketitle

\section{Introduction}
This is a sample document demonstrating various \LaTeX{} features including mathematical equations, tables, and formatting. You can cite papers using the citation search feature, and they will be automatically added to your \texttt{cite.bib} file.

\section{Mathematical Expressions}
Here are some example equations:

\subsection{Quadratic Formula}
The quadratic formula is given by:
\[x = \frac{-b \pm \sqrt{b^2 - 4ac}}{2a}\]

\subsection{Matrix Example}
A sample matrix:
\[
\begin{pmatrix}
a_{11} & a_{12} & a_{13} \\
a_{21} & a_{22} & a_{23} \\
a_{31} & a_{32} & a_{33}
\end{pmatrix}
\]

\section{Tables}
Below is a sample table showing data:

\begin{table}[h]
\centering
\begin{tabular}{@{}lrrr@{}}
\toprule
\textbf{Category} & \textbf{Value 1} & \textbf{Value 2} & \textbf{Total} \\
\midrule
Item A & 10.5 & 20.3 & 30.8 \\
Item B & 15.2 & 25.1 & 40.3 \\
Item C & 22.1 & 12.7 & 34.8 \\
\bottomrule
\end{tabular}
\caption{Sample Data Table}
\label{tab:sample}
\end{table}

\section{Lists}
\subsection{Bullet Points}
\begin{itemize}
  \item First item
  \item Second item
  \item Third item with sub-items:
    \begin{itemize}
      \item Sub-item 1
      \item Sub-item 2
    \end{itemize}
\end{itemize}

\subsection{Numbered List}
\begin{enumerate}
  \item First step
  \item Second step
  \item Third step
\end{enumerate}

\section{Citations and References}
When you use the citation search feature, references will be automatically managed. For example, you might cite a paper like this: \cite{example2024}. All citations will be automatically added to your \texttt{cite.bib} file.

\section{Conclusion}
This template demonstrates basic \LaTeX{} features including mathematical formulas, tables, and lists. The bibliography system is now set up to work seamlessly with Bibby AI's citation management. Feel free to modify it for your needs!

\bibliographystyle{plain}
\bibliography{cite}

\end{document}\documentclass{article}
\usepackage[utf8]{inputenc}
\usepackage{booktabs}
\usepackage{amsmath}
\usepackage{natbib}

\title{My Document}
\author{Your Name}
\date{\today}

\begin{document}

\maketitle

\section{Introduction}
This is a sample document demonstrating various \LaTeX{} features including mathematical equations, tables, and formatting. You can cite papers using the citation search feature, and they will be automatically added to your \texttt{cite.bib} file.

\section{Mathematical Expressions}
Here are some example equations:

\subsection{Quadratic Formula}
The quadratic formula is given by:
\[x = \frac{-b \pm \sqrt{b^2 - 4ac}}{2a}\]

\subsection{Matrix Example}
A sample matrix:
\[
\begin{pmatrix}
a_{11} & a_{12} & a_{13} \\
a_{21} & a_{22} & a_{23} \\
a_{31} & a_{32} & a_{33}
\end{pmatrix}
\]

\section{Tables}
Below is a sample table showing data:

\begin{table}[h]
\centering
\begin{tabular}{@{}lrrr@{}}
\toprule
\textbf{Category} & \textbf{Value 1} & \textbf{Value 2} & \textbf{Total} \\
\midrule
Item A & 10.5 & 20.3 & 30.8 \\
Item B & 15.2 & 25.1 & 40.3 \\
Item C & 22.1 & 12.7 & 34.8 \\
\bottomrule
\end{tabular}
\caption{Sample Data Table}
\label{tab:sample}
\end{table}

\section{Lists}
\subsection{Bullet Points}
\begin{itemize}
  \item First item
  \item Second item
  \item Third item with sub-items:
    \begin{itemize}
      \item Sub-item 1
      \item Sub-item 2
    \end{itemize}
\end{itemize}

\subsection{Numbered List}
\begin{enumerate}
  \item First step
  \item Second step
  \item Third step
\end{enumerate}

\section{Citations and References}
When you use the citation search feature, references will be automatically managed. For example, you might cite a paper like this: \cite{example2024}. All citations will be automatically added to your \texttt{cite.bib} file.

\section{Conclusion}
This template demonstrates basic \LaTeX{} features including mathematical formulas, tables, and lists. The bibliography system is now set up to work seamlessly with Bibby AI's citation management. Feel free to modify it for your needs!

\bibliographystyle{plain}
\bibliography{cite}

\end{document}\documentclass{article}
\usepackage[utf8]{inputenc}
\usepackage{booktabs}
\usepackage{amsmath}
\usepackage{natbib}

\title{My Document}
\author{Your Name}
\date{\today}

\begin{document}

\maketitle

\section{Introduction}
This is a sample document demonstrating various \LaTeX{} features including mathematical equations, tables, and formatting. You can cite papers using the citation search feature, and they will be automatically added to your \texttt{cite.bib} file.

\section{Mathematical Expressions}
Here are some example equations:

\subsection{Quadratic Formula}
The quadratic formula is given by:
\[x = \frac{-b \pm \sqrt{b^2 - 4ac}}{2a}\]

\subsection{Matrix Example}
A sample matrix:
\[
\begin{pmatrix}
a_{11} & a_{12} & a_{13} \\
a_{21} & a_{22} & a_{23} \\
a_{31} & a_{32} & a_{33}
\end{pmatrix}
\]

\section{Tables}
Below is a sample table showing data:

\begin{table}[h]
\centering
\begin{tabular}{@{}lrrr@{}}
\toprule
\textbf{Category} & \textbf{Value 1} & \textbf{Value 2} & \textbf{Total} \\
\midrule
Item A & 10.5 & 20.3 & 30.8 \\
Item B & 15.2 & 25.1 & 40.3 \\
Item C & 22.1 & 12.7 & 34.8 \\
\bottomrule
\end{tabular}
\caption{Sample Data Table}
\label{tab:sample}
\end{table}

\section{Lists}
\subsection{Bullet Points}
\begin{itemize}
  \item First item
  \item Second item
  \item Third item with sub-items:
    \begin{itemize}
      \item Sub-item 1
      \item Sub-item 2
    \end{itemize}
\end{itemize}

\subsection{Numbered List}
\begin{enumerate}
  \item First step
  \item Second step
  \item Third step
\end{enumerate}

\section{Citations and References}
When you use the citation search feature, references will be automatically managed. For example, you might cite a paper like this: \cite{example2024}. All citations will be automatically added to your \texttt{cite.bib} file.

\section{Conclusion}
This template demonstrates basic \LaTeX{} features including mathematical formulas, tables, and lists. The bibliography system is now set up to work seamlessly with Bibby AI's citation management. Feel free to modify it for your needs!

\bibliographystyle{plain}
\bibliography{cite}

\end{document}\documentclass{article}
\usepackage[utf8]{inputenc}
\usepackage{booktabs}
\usepackage{amsmath}
\usepackage{natbib}

\title{My Document}
\author{Your Name}
\date{\today}

\begin{document}

\maketitle

\section{Introduction}
This is a sample document demonstrating various \LaTeX{} features including mathematical equations, tables, and formatting. You can cite papers using the citation search feature, and they will be automatically added to your \texttt{cite.bib} file.

\section{Mathematical Expressions}
Here are some example equations:

\subsection{Quadratic Formula}
The quadratic formula is given by:
\[x = \frac{-b \pm \sqrt{b^2 - 4ac}}{2a}\]

\subsection{Matrix Example}
A sample matrix:
\[
\begin{pmatrix}
a_{11} & a_{12} & a_{13} \\
a_{21} & a_{22} & a_{23} \\
a_{31} & a_{32} & a_{33}
\end{pmatrix}
\]

\section{Tables}
Below is a sample table showing data:

\begin{table}[h]
\centering
\begin{tabular}{@{}lrrr@{}}
\toprule
\textbf{Category} & \textbf{Value 1} & \textbf{Value 2} & \textbf{Total} \\
\midrule
Item A & 10.5 & 20.3 & 30.8 \\
Item B & 15.2 & 25.1 & 40.3 \\
Item C & 22.1 & 12.7 & 34.8 \\
\bottomrule
\end{tabular}
\caption{Sample Data Table}
\label{tab:sample}
\end{table}

\section{Lists}
\subsection{Bullet Points}
\begin{itemize}
  \item First item
  \item Second item
  \item Third item with sub-items:
    \begin{itemize}
      \item Sub-item 1
      \item Sub-item 2
    \end{itemize}
\end{itemize}

\subsection{Numbered List}
\begin{enumerate}
  \item First step
  \item Second step
  \item Third step
\end{enumerate}

\section{Citations and References}
When you use the citation search feature, references will be automatically managed. For example, you might cite a paper like this: \cite{example2024}. All citations will be automatically added to your \texttt{cite.bib} file.

\section{Conclusion}
This template demonstrates basic \LaTeX{} features including mathematical formulas, tables, and lists. The bibliography system is now set up to work seamlessly with Bibby AI's citation management. Feel free to modify it for your needs!

\bibliographystyle{plain}
\bibliography{cite}

\end{document}\documentclass{article}
\usepackage[utf8]{inputenc}
\usepackage{booktabs}
\usepackage{amsmath}
\usepackage{natbib}

\title{My Document}
\author{Your Name}
\date{\today}

\begin{document}

\maketitle

\section{Introduction}
This is a sample document demonstrating various \LaTeX{} features including mathematical equations, tables, and formatting. You can cite papers using the citation search feature, and they will be automatically added to your \texttt{cite.bib} file.

\section{Mathematical Expressions}
Here are some example equations:

\subsection{Quadratic Formula}
The quadratic formula is given by:
\[x = \frac{-b \pm \sqrt{b^2 - 4ac}}{2a}\]

\subsection{Matrix Example}
A sample matrix:
\[
\begin{pmatrix}
a_{11} & a_{12} & a_{13} \\
a_{21} & a_{22} & a_{23} \\
a_{31} & a_{32} & a_{33}
\end{pmatrix}
\]

\section{Tables}
Below is a sample table showing data:

\begin{table}[h]
\centering
\begin{tabular}{@{}lrrr@{}}
\toprule
\textbf{Category} & \textbf{Value 1} & \textbf{Value 2} & \textbf{Total} \\
\midrule
Item A & 10.5 & 20.3 & 30.8 \\
Item B & 15.2 & 25.1 & 40.3 \\
Item C & 22.1 & 12.7 & 34.8 \\
\bottomrule
\end{tabular}
\caption{Sample Data Table}
\label{tab:sample}
\end{table}

\section{Lists}
\subsection{Bullet Points}
\begin{itemize}
  \item First item
  \item Second item
  \item Third item with sub-items:
    \begin{itemize}
      \item Sub-item 1
      \item Sub-item 2
    \end{itemize}
\end{itemize}

\subsection{Numbered List}
\begin{enumerate}
  \item First step
  \item Second step
  \item Third step
\end{enumerate}

\section{Citations and References}
When you use the citation search feature, references will be automatically managed. For example, you might cite a paper like this: \cite{example2024}. All citations will be automatically added to your \texttt{cite.bib} file.

\section{Conclusion}
This template demonstrates basic \LaTeX{} features including mathematical formulas, tables, and lists. The bibliography system is now set up to work seamlessly with Bibby AI's citation management. Feel free to modify it for your needs!

\bibliographystyle{plain}
\bibliography{cite}

\end{document}\documentclass{article}
\usepackage[utf8]{inputenc}
\usepackage{booktabs}
\usepackage{amsmath}
\usepackage{natbib}

\title{My Document}
\author{Your Name}
\date{\today}

\begin{document}

\maketitle

\section{Introduction}
This is a sample document demonstrating various \LaTeX{} features including mathematical equations, tables, and formatting. You can cite papers using the citation search feature, and they will be automatically added to your \texttt{cite.bib} file.

\section{Mathematical Expressions}
Here are some example equations:

\subsection{Quadratic Formula}
The quadratic formula is given by:
\[x = \frac{-b \pm \sqrt{b^2 - 4ac}}{2a}\]

\subsection{Matrix Example}
A sample matrix:
\[
\begin{pmatrix}
a_{11} & a_{12} & a_{13} \\
a_{21} & a_{22} & a_{23} \\
a_{31} & a_{32} & a_{33}
\end{pmatrix}
\]

\section{Tables}
Below is a sample table showing data:

\begin{table}[h]
\centering
\begin{tabular}{@{}lrrr@{}}
\toprule
\textbf{Category} & \textbf{Value 1} & \textbf{Value 2} & \textbf{Total} \\
\midrule
Item A & 10.5 & 20.3 & 30.8 \\
Item B & 15.2 & 25.1 & 40.3 \\
Item C & 22.1 & 12.7 & 34.8 \\
\bottomrule
\end{tabular}
\caption{Sample Data Table}
\label{tab:sample}
\end{table}

\section{Lists}
\subsection{Bullet Points}
\begin{itemize}
  \item First item
  \item Second item
  \item Third item with sub-items:
    \begin{itemize}
      \item Sub-item 1
      \item Sub-item 2
    \end{itemize}
\end{itemize}

\subsection{Numbered List}
\begin{enumerate}
  \item First step
  \item Second step
  \item Third step
\end{enumerate}

\section{Citations and References}
When you use the citation search feature, references will be automatically managed. For example, you might cite a paper like this: \cite{example2024}. All citations will be automatically added to your \texttt{cite.bib} file.

\section{Conclusion}
This template demonstrates basic \LaTeX{} features including mathematical formulas, tables, and lists. The bibliography system is now set up to work seamlessly with Bibby AI's citation management. Feel free to modify it for your needs!

\bibliographystyle{plain}
\bibliography{cite}

\end{document}\documentclass{article}
\usepackage[utf8]{inputenc}
\usepackage{booktabs}
\usepackage{amsmath}
\usepackage{natbib}

\title{My Document}
\author{Your Name}
\date{\today}

\begin{document}

\maketitle

\section{Introduction}
This is a sample document demonstrating various \LaTeX{} features including mathematical equations, tables, and formatting. You can cite papers using the citation search feature, and they will be automatically added to your \texttt{cite.bib} file.

\section{Mathematical Expressions}
Here are some example equations:

\subsection{Quadratic Formula}
The quadratic formula is given by:
\[x = \frac{-b \pm \sqrt{b^2 - 4ac}}{2a}\]

\subsection{Matrix Example}
A sample matrix:
\[
\begin{pmatrix}
a_{11} & a_{12} & a_{13} \\
a_{21} & a_{22} & a_{23} \\
a_{31} & a_{32} & a_{33}
\end{pmatrix}
\]

\section{Tables}
Below is a sample table showing data:

\begin{table}[h]
\centering
\begin{tabular}{@{}lrrr@{}}
\toprule
\textbf{Category} & \textbf{Value 1} & \textbf{Value 2} & \textbf{Total} \\
\midrule
Item A & 10.5 & 20.3 & 30.8 \\
Item B & 15.2 & 25.1 & 40.3 \\
Item C & 22.1 & 12.7 & 34.8 \\
\bottomrule
\end{tabular}
\caption{Sample Data Table}
\label{tab:sample}
\end{table}

\section{Lists}
\subsection{Bullet Points}
\begin{itemize}
  \item First item
  \item Second item
  \item Third item with sub-items:
    \begin{itemize}
      \item Sub-item 1
      \item Sub-item 2
    \end{itemize}
\end{itemize}

\subsection{Numbered List}
\begin{enumerate}
  \item First step
  \item Second step
  \item Third step
\end{enumerate}

\section{Citations and References}
When you use the citation search feature, references will be automatically managed. For example, you might cite a paper like this: \cite{example2024}. All citations will be automatically added to your \texttt{cite.bib} file.

\section{Conclusion}
This template demonstrates basic \LaTeX{} features including mathematical formulas, tables, and lists. The bibliography system is now set up to work seamlessly with Bibby AI's citation management. Feel free to modify it for your needs!

\bibliographystyle{plain}
\bibliography{cite}

\end{document}\documentclass{article}
\usepackage[utf8]{inputenc}
\usepackage{booktabs}
\usepackage{amsmath}
\usepackage{natbib}

\title{My Document}
\author{Your Name}
\date{\today}

\begin{document}

\maketitle

\section{Introduction}
This is a sample document demonstrating various \LaTeX{} features including mathematical equations, tables, and formatting. You can cite papers using the citation search feature, and they will be automatically added to your \texttt{cite.bib} file.

\section{Mathematical Expressions}
Here are some example equations:

\subsection{Quadratic Formula}
The quadratic formula is given by:
\[x = \frac{-b \pm \sqrt{b^2 - 4ac}}{2a}\]

\subsection{Matrix Example}
A sample matrix:
\[
\begin{pmatrix}
a_{11} & a_{12} & a_{13} \\
a_{21} & a_{22} & a_{23} \\
a_{31} & a_{32} & a_{33}
\end{pmatrix}
\]

\section{Tables}
Below is a sample table showing data:

\begin{table}[h]
\centering
\begin{tabular}{@{}lrrr@{}}
\toprule
\textbf{Category} & \textbf{Value 1} & \textbf{Value 2} & \textbf{Total} \\
\midrule
Item A & 10.5 & 20.3 & 30.8 \\
Item B & 15.2 & 25.1 & 40.3 \\
Item C & 22.1 & 12.7 & 34.8 \\
\bottomrule
\end{tabular}
\caption{Sample Data Table}
\label{tab:sample}
\end{table}

\section{Lists}
\subsection{Bullet Points}
\begin{itemize}
  \item First item
  \item Second item
  \item Third item with sub-items:
    \begin{itemize}
      \item Sub-item 1
      \item Sub-item 2
    \end{itemize}
\end{itemize}

\subsection{Numbered List}
\begin{enumerate}
  \item First step
  \item Second step
  \item Third step
\end{enumerate}

\section{Citations and References}
When you use the citation search feature, references will be automatically managed. For example, you might cite a paper like this: \cite{example2024}. All citations will be automatically added to your \texttt{cite.bib} file.

\section{Conclusion}
This template demonstrates basic \LaTeX{} features including mathematical formulas, tables, and lists. The bibliography system is now set up to work seamlessly with Bibby AI's citation management. Feel free to modify it for your needs!

\bibliographystyle{plain}
\bibliography{cite}

\end{document}\documentclass{article}
\usepackage[utf8]{inputenc}
\usepackage{booktabs}
\usepackage{amsmath}
\usepackage{natbib}

\title{My Document}
\author{Your Name}
\date{\today}

\begin{document}

\maketitle

\section{Introduction}
This is a sample document demonstrating various \LaTeX{} features including mathematical equations, tables, and formatting. You can cite papers using the citation search feature, and they will be automatically added to your \texttt{cite.bib} file.

\section{Mathematical Expressions}
Here are some example equations:

\subsection{Quadratic Formula}
The quadratic formula is given by:
\[x = \frac{-b \pm \sqrt{b^2 - 4ac}}{2a}\]

\subsection{Matrix Example}
A sample matrix:
\[
\begin{pmatrix}
a_{11} & a_{12} & a_{13} \\
a_{21} & a_{22} & a_{23} \\
a_{31} & a_{32} & a_{33}
\end{pmatrix}
\]

\section{Tables}
Below is a sample table showing data:

\begin{table}[h]
\centering
\begin{tabular}{@{}lrrr@{}}
\toprule
\textbf{Category} & \textbf{Value 1} & \textbf{Value 2} & \textbf{Total} \\
\midrule
Item A & 10.5 & 20.3 & 30.8 \\
Item B & 15.2 & 25.1 & 40.3 \\
Item C & 22.1 & 12.7 & 34.8 \\
\bottomrule
\end{tabular}
\caption{Sample Data Table}
\label{tab:sample}
\end{table}

\section{Lists}
\subsection{Bullet Points}
\begin{itemize}
  \item First item
  \item Second item
  \item Third item with sub-items:
    \begin{itemize}
      \item Sub-item 1
      \item Sub-item 2
    \end{itemize}
\end{itemize}

\subsection{Numbered List}
\begin{enumerate}
  \item First step
  \item Second step
  \item Third step
\end{enumerate}

\section{Citations and References}
When you use the citation search feature, references will be automatically managed. For example, you might cite a paper like this: \cite{example2024}. All citations will be automatically added to your \texttt{cite.bib} file.

\section{Conclusion}
This template demonstrates basic \LaTeX{} features including mathematical formulas, tables, and lists. The bibliography system is now set up to work seamlessly with Bibby AI's citation management. Feel free to modify it for your needs!

\bibliographystyle{plain}
\bibliography{cite}

\end{document}\documentclass{article}
\usepackage[utf8]{inputenc}
\usepackage{booktabs}
\usepackage{amsmath}
\usepackage{natbib}

\title{My Document}
\author{Your Name}
\date{\today}

\begin{document}

\maketitle

\section{Introduction}
This is a sample document demonstrating various \LaTeX{} features including mathematical equations, tables, and formatting. You can cite papers using the citation search feature, and they will be automatically added to your \texttt{cite.bib} file.

\section{Mathematical Expressions}
Here are some example equations:

\subsection{Quadratic Formula}
The quadratic formula is given by:
\[x = \frac{-b \pm \sqrt{b^2 - 4ac}}{2a}\]

\subsection{Matrix Example}
A sample matrix:
\[
\begin{pmatrix}
a_{11} & a_{12} & a_{13} \\
a_{21} & a_{22} & a_{23} \\
a_{31} & a_{32} & a_{33}
\end{pmatrix}
\]

\section{Tables}
Below is a sample table showing data:

\begin{table}[h]
\centering
\begin{tabular}{@{}lrrr@{}}
\toprule
\textbf{Category} & \textbf{Value 1} & \textbf{Value 2} & \textbf{Total} \\
\midrule
Item A & 10.5 & 20.3 & 30.8 \\
Item B & 15.2 & 25.1 & 40.3 \\
Item C & 22.1 & 12.7 & 34.8 \\
\bottomrule
\end{tabular}
\caption{Sample Data Table}
\label{tab:sample}
\end{table}

\section{Lists}
\subsection{Bullet Points}
\begin{itemize}
  \item First item
  \item Second item
  \item Third item with sub-items:
    \begin{itemize}
      \item Sub-item 1
      \item Sub-item 2
    \end{itemize}
\end{itemize}

\subsection{Numbered List}
\begin{enumerate}
  \item First step
  \item Second step
  \item Third step
\end{enumerate}

\section{Citations and References}
When you use the citation search feature, references will be automatically managed. For example, you might cite a paper like this: \cite{example2024}. All citations will be automatically added to your \texttt{cite.bib} file.

\section{Conclusion}
This template demonstrates basic \LaTeX{} features including mathematical formulas, tables, and lists. The bibliography system is now set up to work seamlessly with Bibby AI's citation management. Feel free to modify it for your needs!

\bibliographystyle{plain}
\bibliography{cite}

\end{document}\documentclass{article}
\usepackage[utf8]{inputenc}
\usepackage{booktabs}
\usepackage{amsmath}
\usepackage{natbib}

\title{My Document}
\author{Your Name}
\date{\today}

\begin{document}

\maketitle

\section{Introduction}
This is a sample document demonstrating various \LaTeX{} features including mathematical equations, tables, and formatting. You can cite papers using the citation search feature, and they will be automatically added to your \texttt{cite.bib} file.

\section{Mathematical Expressions}
Here are some example equations:

\subsection{Quadratic Formula}
The quadratic formula is given by:
\[x = \frac{-b \pm \sqrt{b^2 - 4ac}}{2a}\]

\subsection{Matrix Example}
A sample matrix:
\[
\begin{pmatrix}
a_{11} & a_{12} & a_{13} \\
a_{21} & a_{22} & a_{23} \\
a_{31} & a_{32} & a_{33}
\end{pmatrix}
\]

\section{Tables}
Below is a sample table showing data:

\begin{table}[h]
\centering
\begin{tabular}{@{}lrrr@{}}
\toprule
\textbf{Category} & \textbf{Value 1} & \textbf{Value 2} & \textbf{Total} \\
\midrule
Item A & 10.5 & 20.3 & 30.8 \\
Item B & 15.2 & 25.1 & 40.3 \\
Item C & 22.1 & 12.7 & 34.8 \\
\bottomrule
\end{tabular}
\caption{Sample Data Table}
\label{tab:sample}
\end{table}

\section{Lists}
\subsection{Bullet Points}
\begin{itemize}
  \item First item
  \item Second item
  \item Third item with sub-items:
    \begin{itemize}
      \item Sub-item 1
      \item Sub-item 2
    \end{itemize}
\end{itemize}

\subsection{Numbered List}
\begin{enumerate}
  \item First step
  \item Second step
  \item Third step
\end{enumerate}

\section{Citations and References}
When you use the citation search feature, references will be automatically managed. For example, you might cite a paper like this: \cite{example2024}. All citations will be automatically added to your \texttt{cite.bib} file.

\section{Conclusion}
This template demonstrates basic \LaTeX{} features including mathematical formulas, tables, and lists. The bibliography system is now set up to work seamlessly with Bibby AI's citation management. Feel free to modify it for your needs!

\bibliographystyle{plain}
\bibliography{cite}

\end{document}\documentclass{article}
\usepackage[utf8]{inputenc}
\usepackage{booktabs}
\usepackage{amsmath}
\usepackage{natbib}

\title{My Document}
\author{Your Name}
\date{\today}

\begin{document}

\maketitle

\section{Introduction}
This is a sample document demonstrating various \LaTeX{} features including mathematical equations, tables, and formatting. You can cite papers using the citation search feature, and they will be automatically added to your \texttt{cite.bib} file.

\section{Mathematical Expressions}
Here are some example equations:

\subsection{Quadratic Formula}
The quadratic formula is given by:
\[x = \frac{-b \pm \sqrt{b^2 - 4ac}}{2a}\]

\subsection{Matrix Example}
A sample matrix:
\[
\begin{pmatrix}
a_{11} & a_{12} & a_{13} \\
a_{21} & a_{22} & a_{23} \\
a_{31} & a_{32} & a_{33}
\end{pmatrix}
\]

\section{Tables}
Below is a sample table showing data:

\begin{table}[h]
\centering
\begin{tabular}{@{}lrrr@{}}
\toprule
\textbf{Category} & \textbf{Value 1} & \textbf{Value 2} & \textbf{Total} \\
\midrule
Item A & 10.5 & 20.3 & 30.8 \\
Item B & 15.2 & 25.1 & 40.3 \\
Item C & 22.1 & 12.7 & 34.8 \\
\bottomrule
\end{tabular}
\caption{Sample Data Table}
\label{tab:sample}
\end{table}

\section{Lists}
\subsection{Bullet Points}
\begin{itemize}
  \item First item
  \item Second item
  \item Third item with sub-items:
    \begin{itemize}
      \item Sub-item 1
      \item Sub-item 2
    \end{itemize}
\end{itemize}

\subsection{Numbered List}
\begin{enumerate}
  \item First step
  \item Second step
  \item Third step
\end{enumerate}

\section{Citations and References}
When you use the citation search feature, references will be automatically managed. For example, you might cite a paper like this: \cite{example2024}. All citations will be automatically added to your \texttt{cite.bib} file.

\section{Conclusion}
This template demonstrates basic \LaTeX{} features including mathematical formulas, tables, and lists. The bibliography system is now set up to work seamlessly with Bibby AI's citation management. Feel free to modify it for your needs!

\bibliographystyle{plain}
\bibliography{cite}

\end{document}\documentclass{article}
\usepackage[utf8]{inputenc}
\usepackage{booktabs}
\usepackage{amsmath}
\usepackage{natbib}

\title{My Document}
\author{Your Name}
\date{\today}

\begin{document}

\maketitle

\section{Introduction}
This is a sample document demonstrating various \LaTeX{} features including mathematical equations, tables, and formatting. You can cite papers using the citation search feature, and they will be automatically added to your \texttt{cite.bib} file.

\section{Mathematical Expressions}
Here are some example equations:

\subsection{Quadratic Formula}
The quadratic formula is given by:
\[x = \frac{-b \pm \sqrt{b^2 - 4ac}}{2a}\]

\subsection{Matrix Example}
A sample matrix:
\[
\begin{pmatrix}
a_{11} & a_{12} & a_{13} \\
a_{21} & a_{22} & a_{23} \\
a_{31} & a_{32} & a_{33}
\end{pmatrix}
\]

\section{Tables}
Below is a sample table showing data:

\begin{table}[h]
\centering
\begin{tabular}{@{}lrrr@{}}
\toprule
\textbf{Category} & \textbf{Value 1} & \textbf{Value 2} & \textbf{Total} \\
\midrule
Item A & 10.5 & 20.3 & 30.8 \\
Item B & 15.2 & 25.1 & 40.3 \\
Item C & 22.1 & 12.7 & 34.8 \\
\bottomrule
\end{tabular}
\caption{Sample Data Table}
\label{tab:sample}
\end{table}

\section{Lists}
\subsection{Bullet Points}
\begin{itemize}
  \item First item
  \item Second item
  \item Third item with sub-items:
    \begin{itemize}
      \item Sub-item 1
      \item Sub-item 2
    \end{itemize}
\end{itemize}

\subsection{Numbered List}
\begin{enumerate}
  \item First step
  \item Second step
  \item Third step
\end{enumerate}

\section{Citations and References}
When you use the citation search feature, references will be automatically managed. For example, you might cite a paper like this: \cite{example2024}. All citations will be automatically added to your \texttt{cite.bib} file.

\section{Conclusion}
This template demonstrates basic \LaTeX{} features including mathematical formulas, tables, and lists. The bibliography system is now set up to work seamlessly with Bibby AI's citation management. Feel free to modify it for your needs!

\bibliographystyle{plain}
\bibliography{cite}

\end{document}\documentclass{article}
\usepackage[utf8]{inputenc}
\usepackage{booktabs}
\usepackage{amsmath}
\usepackage{natbib}

\title{My Document}
\author{Your Name}
\date{\today}

\begin{document}

\maketitle

\section{Introduction}
This is a sample document demonstrating various \LaTeX{} features including mathematical equations, tables, and formatting. You can cite papers using the citation search feature, and they will be automatically added to your \texttt{cite.bib} file.

\section{Mathematical Expressions}
Here are some example equations:

\subsection{Quadratic Formula}
The quadratic formula is given by:
\[x = \frac{-b \pm \sqrt{b^2 - 4ac}}{2a}\]

\subsection{Matrix Example}
A sample matrix:
\[
\begin{pmatrix}
a_{11} & a_{12} & a_{13} \\
a_{21} & a_{22} & a_{23} \\
a_{31} & a_{32} & a_{33}
\end{pmatrix}
\]

\section{Tables}
Below is a sample table showing data:

\begin{table}[h]
\centering
\begin{tabular}{@{}lrrr@{}}
\toprule
\textbf{Category} & \textbf{Value 1} & \textbf{Value 2} & \textbf{Total} \\
\midrule
Item A & 10.5 & 20.3 & 30.8 \\
Item B & 15.2 & 25.1 & 40.3 \\
Item C & 22.1 & 12.7 & 34.8 \\
\bottomrule
\end{tabular}
\caption{Sample Data Table}
\label{tab:sample}
\end{table}

\section{Lists}
\subsection{Bullet Points}
\begin{itemize}
  \item First item
  \item Second item
  \item Third item with sub-items:
    \begin{itemize}
      \item Sub-item 1
      \item Sub-item 2
    \end{itemize}
\end{itemize}

\subsection{Numbered List}
\begin{enumerate}
  \item First step
  \item Second step
  \item Third step
\end{enumerate}

\section{Citations and References}
When you use the citation search feature, references will be automatically managed. For example, you might cite a paper like this: \cite{example2024}. All citations will be automatically added to your \texttt{cite.bib} file.

\section{Conclusion}
This template demonstrates basic \LaTeX{} features including mathematical formulas, tables, and lists. The bibliography system is now set up to work seamlessly with Bibby AI's citation management. Feel free to modify it for your needs!

\bibliographystyle{plain}
\bibliography{cite}

\end{document}\documentclass{article}
\usepackage[utf8]{inputenc}
\usepackage{booktabs}
\usepackage{amsmath}
\usepackage{natbib}

\title{My Document}
\author{Your Name}
\date{\today}

\begin{document}

\maketitle

\section{Introduction}
This is a sample document demonstrating various \LaTeX{} features including mathematical equations, tables, and formatting. You can cite papers using the citation search feature, and they will be automatically added to your \texttt{cite.bib} file.

\section{Mathematical Expressions}
Here are some example equations:

\subsection{Quadratic Formula}
The quadratic formula is given by:
\[x = \frac{-b \pm \sqrt{b^2 - 4ac}}{2a}\]

\subsection{Matrix Example}
A sample matrix:
\[
\begin{pmatrix}
a_{11} & a_{12} & a_{13} \\
a_{21} & a_{22} & a_{23} \\
a_{31} & a_{32} & a_{33}
\end{pmatrix}
\]

\section{Tables}
Below is a sample table showing data:

\begin{table}[h]
\centering
\begin{tabular}{@{}lrrr@{}}
\toprule
\textbf{Category} & \textbf{Value 1} & \textbf{Value 2} & \textbf{Total} \\
\midrule
Item A & 10.5 & 20.3 & 30.8 \\
Item B & 15.2 & 25.1 & 40.3 \\
Item C & 22.1 & 12.7 & 34.8 \\
\bottomrule
\end{tabular}
\caption{Sample Data Table}
\label{tab:sample}
\end{table}

\section{Lists}
\subsection{Bullet Points}
\begin{itemize}
  \item First item
  \item Second item
  \item Third item with sub-items:
    \begin{itemize}
      \item Sub-item 1
      \item Sub-item 2
    \end{itemize}
\end{itemize}

\subsection{Numbered List}
\begin{enumerate}
  \item First step
  \item Second step
  \item Third step
\end{enumerate}

\section{Citations and References}
When you use the citation search feature, references will be automatically managed. For example, you might cite a paper like this: \cite{example2024}. All citations will be automatically added to your \texttt{cite.bib} file.

\section{Conclusion}
This template demonstrates basic \LaTeX{} features including mathematical formulas, tables, and lists. The bibliography system is now set up to work seamlessly with Bibby AI's citation management. Feel free to modify it for your needs!

\bibliographystyle{plain}
\bibliography{cite}

\end{document}\documentclass{article}
\usepackage[utf8]{inputenc}
\usepackage{booktabs}
\usepackage{amsmath}
\usepackage{natbib}

\title{My Document}
\author{Your Name}
\date{\today}

\begin{document}

\maketitle

\section{Introduction}
This is a sample document demonstrating various \LaTeX{} features including mathematical equations, tables, and formatting. You can cite papers using the citation search feature, and they will be automatically added to your \texttt{cite.bib} file.

\section{Mathematical Expressions}
Here are some example equations:

\subsection{Quadratic Formula}
The quadratic formula is given by:
\[x = \frac{-b \pm \sqrt{b^2 - 4ac}}{2a}\]

\subsection{Matrix Example}
A sample matrix:
\[
\begin{pmatrix}
a_{11} & a_{12} & a_{13} \\
a_{21} & a_{22} & a_{23} \\
a_{31} & a_{32} & a_{33}
\end{pmatrix}
\]

\section{Tables}
Below is a sample table showing data:

\begin{table}[h]
\centering
\begin{tabular}{@{}lrrr@{}}
\toprule
\textbf{Category} & \textbf{Value 1} & \textbf{Value 2} & \textbf{Total} \\
\midrule
Item A & 10.5 & 20.3 & 30.8 \\
Item B & 15.2 & 25.1 & 40.3 \\
Item C & 22.1 & 12.7 & 34.8 \\
\bottomrule
\end{tabular}
\caption{Sample Data Table}
\label{tab:sample}
\end{table}

\section{Lists}
\subsection{Bullet Points}
\begin{itemize}
  \item First item
  \item Second item
  \item Third item with sub-items:
    \begin{itemize}
      \item Sub-item 1
      \item Sub-item 2
    \end{itemize}
\end{itemize}

\subsection{Numbered List}
\begin{enumerate}
  \item First step
  \item Second step
  \item Third step
\end{enumerate}

\section{Citations and References}
When you use the citation search feature, references will be automatically managed. For example, you might cite a paper like this: \cite{example2024}. All citations will be automatically added to your \texttt{cite.bib} file.

\section{Conclusion}
This template demonstrates basic \LaTeX{} features including mathematical formulas, tables, and lists. The bibliography system is now set up to work seamlessly with Bibby AI's citation management. Feel free to modify it for your needs!

\bibliographystyle{plain}
\bibliography{cite}

\end{document}\documentclass{article}
\usepackage[utf8]{inputenc}
\usepackage{booktabs}
\usepackage{amsmath}
\usepackage{natbib}

\title{My Document}
\author{Your Name}
\date{\today}

\begin{document}

\maketitle

\section{Introduction}
This is a sample document demonstrating various \LaTeX{} features including mathematical equations, tables, and formatting. You can cite papers using the citation search feature, and they will be automatically added to your \texttt{cite.bib} file.

\section{Mathematical Expressions}
Here are some example equations:

\subsection{Quadratic Formula}
The quadratic formula is given by:
\[x = \frac{-b \pm \sqrt{b^2 - 4ac}}{2a}\]

\subsection{Matrix Example}
A sample matrix:
\[
\begin{pmatrix}
a_{11} & a_{12} & a_{13} \\
a_{21} & a_{22} & a_{23} \\
a_{31} & a_{32} & a_{33}
\end{pmatrix}
\]

\section{Tables}
Below is a sample table showing data:

\begin{table}[h]
\centering
\begin{tabular}{@{}lrrr@{}}
\toprule
\textbf{Category} & \textbf{Value 1} & \textbf{Value 2} & \textbf{Total} \\
\midrule
Item A & 10.5 & 20.3 & 30.8 \\
Item B & 15.2 & 25.1 & 40.3 \\
Item C & 22.1 & 12.7 & 34.8 \\
\bottomrule
\end{tabular}
\caption{Sample Data Table}
\label{tab:sample}
\end{table}

\section{Lists}
\subsection{Bullet Points}
\begin{itemize}
  \item First item
  \item Second item
  \item Third item with sub-items:
    \begin{itemize}
      \item Sub-item 1
      \item Sub-item 2
    \end{itemize}
\end{itemize}

\subsection{Numbered List}
\begin{enumerate}
  \item First step
  \item Second step
  \item Third step
\end{enumerate}

\section{Citations and References}
When you use the citation search feature, references will be automatically managed. For example, you might cite a paper like this: \cite{example2024}. All citations will be automatically added to your \texttt{cite.bib} file.

\section{Conclusion}
This template demonstrates basic \LaTeX{} features including mathematical formulas, tables, and lists. The bibliography system is now set up to work seamlessly with Bibby AI's citation management. Feel free to modify it for your needs!

\bibliographystyle{plain}
\bibliography{cite}

\end{document}\documentclass{article}
\usepackage[utf8]{inputenc}
\usepackage{booktabs}
\usepackage{amsmath}
\usepackage{natbib}

\title{My Document}
\author{Your Name}
\date{\today}

\begin{document}

\maketitle

\section{Introduction}
This is a sample document demonstrating various \LaTeX{} features including mathematical equations, tables, and formatting. You can cite papers using the citation search feature, and they will be automatically added to your \texttt{cite.bib} file.

\section{Mathematical Expressions}
Here are some example equations:

\subsection{Quadratic Formula}
The quadratic formula is given by:
\[x = \frac{-b \pm \sqrt{b^2 - 4ac}}{2a}\]

\subsection{Matrix Example}
A sample matrix:
\[
\begin{pmatrix}
a_{11} & a_{12} & a_{13} \\
a_{21} & a_{22} & a_{23} \\
a_{31} & a_{32} & a_{33}
\end{pmatrix}
\]

\section{Tables}
Below is a sample table showing data:

\begin{table}[h]
\centering
\begin{tabular}{@{}lrrr@{}}
\toprule
\textbf{Category} & \textbf{Value 1} & \textbf{Value 2} & \textbf{Total} \\
\midrule
Item A & 10.5 & 20.3 & 30.8 \\
Item B & 15.2 & 25.1 & 40.3 \\
Item C & 22.1 & 12.7 & 34.8 \\
\bottomrule
\end{tabular}
\caption{Sample Data Table}
\label{tab:sample}
\end{table}

\section{Lists}
\subsection{Bullet Points}
\begin{itemize}
  \item First item
  \item Second item
  \item Third item with sub-items:
    \begin{itemize}
      \item Sub-item 1
      \item Sub-item 2
    \end{itemize}
\end{itemize}

\subsection{Numbered List}
\begin{enumerate}
  \item First step
  \item Second step
  \item Third step
\end{enumerate}

\section{Citations and References}
When you use the citation search feature, references will be automatically managed. For example, you might cite a paper like this: \cite{example2024}. All citations will be automatically added to your \texttt{cite.bib} file.

\section{Conclusion}
This template demonstrates basic \LaTeX{} features including mathematical formulas, tables, and lists. The bibliography system is now set up to work seamlessly with Bibby AI's citation management. Feel free to modify it for your needs!

\bibliographystyle{plain}
\bibliography{cite}

\end{document}\documentclass{article}
\usepackage[utf8]{inputenc}
\usepackage{booktabs}
\usepackage{amsmath}
\usepackage{natbib}

\title{My Document}
\author{Your Name}
\date{\today}

\begin{document}

\maketitle

\section{Introduction}
This is a sample document demonstrating various \LaTeX{} features including mathematical equations, tables, and formatting. You can cite papers using the citation search feature, and they will be automatically added to your \texttt{cite.bib} file.

\section{Mathematical Expressions}
Here are some example equations:

\subsection{Quadratic Formula}
The quadratic formula is given by:
\[x = \frac{-b \pm \sqrt{b^2 - 4ac}}{2a}\]

\subsection{Matrix Example}
A sample matrix:
\[
\begin{pmatrix}
a_{11} & a_{12} & a_{13} \\
a_{21} & a_{22} & a_{23} \\
a_{31} & a_{32} & a_{33}
\end{pmatrix}
\]

\section{Tables}
Below is a sample table showing data:

\begin{table}[h]
\centering
\begin{tabular}{@{}lrrr@{}}
\toprule
\textbf{Category} & \textbf{Value 1} & \textbf{Value 2} & \textbf{Total} \\
\midrule
Item A & 10.5 & 20.3 & 30.8 \\
Item B & 15.2 & 25.1 & 40.3 \\
Item C & 22.1 & 12.7 & 34.8 \\
\bottomrule
\end{tabular}
\caption{Sample Data Table}
\label{tab:sample}
\end{table}

\section{Lists}
\subsection{Bullet Points}
\begin{itemize}
  \item First item
  \item Second item
  \item Third item with sub-items:
    \begin{itemize}
      \item Sub-item 1
      \item Sub-item 2
    \end{itemize}
\end{itemize}

\subsection{Numbered List}
\begin{enumerate}
  \item First step
  \item Second step
  \item Third step
\end{enumerate}

\section{Citations and References}
When you use the citation search feature, references will be automatically managed. For example, you might cite a paper like this: \cite{example2024}. All citations will be automatically added to your \texttt{cite.bib} file.

\section{Conclusion}
This template demonstrates basic \LaTeX{} features including mathematical formulas, tables, and lists. The bibliography system is now set up to work seamlessly with Bibby AI's citation management. Feel free to modify it for your needs!

\bibliographystyle{plain}
\bibliography{cite}

\end{document}\documentclass{article}
\usepackage[utf8]{inputenc}
\usepackage{booktabs}
\usepackage{amsmath}
\usepackage{natbib}

\title{My Document}
\author{Your Name}
\date{\today}

\begin{document}

\maketitle

\section{Introduction}
This is a sample document demonstrating various \LaTeX{} features including mathematical equations, tables, and formatting. You can cite papers using the citation search feature, and they will be automatically added to your \texttt{cite.bib} file.

\section{Mathematical Expressions}
Here are some example equations:

\subsection{Quadratic Formula}
The quadratic formula is given by:
\[x = \frac{-b \pm \sqrt{b^2 - 4ac}}{2a}\]

\subsection{Matrix Example}
A sample matrix:
\[
\begin{pmatrix}
a_{11} & a_{12} & a_{13} \\
a_{21} & a_{22} & a_{23} \\
a_{31} & a_{32} & a_{33}
\end{pmatrix}
\]

\section{Tables}
Below is a sample table showing data:

\begin{table}[h]
\centering
\begin{tabular}{@{}lrrr@{}}
\toprule
\textbf{Category} & \textbf{Value 1} & \textbf{Value 2} & \textbf{Total} \\
\midrule
Item A & 10.5 & 20.3 & 30.8 \\
Item B & 15.2 & 25.1 & 40.3 \\
Item C & 22.1 & 12.7 & 34.8 \\
\bottomrule
\end{tabular}
\caption{Sample Data Table}
\label{tab:sample}
\end{table}

\section{Lists}
\subsection{Bullet Points}
\begin{itemize}
  \item First item
  \item Second item
  \item Third item with sub-items:
    \begin{itemize}
      \item Sub-item 1
      \item Sub-item 2
    \end{itemize}
\end{itemize}

\subsection{Numbered List}
\begin{enumerate}
  \item First step
  \item Second step
  \item Third step
\end{enumerate}

\section{Citations and References}
When you use the citation search feature, references will be automatically managed. For example, you might cite a paper like this: \cite{example2024}. All citations will be automatically added to your \texttt{cite.bib} file.

\section{Conclusion}
This template demonstrates basic \LaTeX{} features including mathematical formulas, tables, and lists. The bibliography system is now set up to work seamlessly with Bibby AI's citation management. Feel free to modify it for your needs!

\bibliographystyle{plain}
\bibliography{cite}

\end{document}\documentclass{article}
\usepackage[utf8]{inputenc}
\usepackage{booktabs}
\usepackage{amsmath}
\usepackage{natbib}

\title{My Document}
\author{Your Name}
\date{\today}

\begin{document}

\maketitle

\section{Introduction}
This is a sample document demonstrating various \LaTeX{} features including mathematical equations, tables, and formatting. You can cite papers using the citation search feature, and they will be automatically added to your \texttt{cite.bib} file.

\section{Mathematical Expressions}
Here are some example equations:

\subsection{Quadratic Formula}
The quadratic formula is given by:
\[x = \frac{-b \pm \sqrt{b^2 - 4ac}}{2a}\]

\subsection{Matrix Example}
A sample matrix:
\[
\begin{pmatrix}
a_{11} & a_{12} & a_{13} \\
a_{21} & a_{22} & a_{23} \\
a_{31} & a_{32} & a_{33}
\end{pmatrix}
\]

\section{Tables}
Below is a sample table showing data:

\begin{table}[h]
\centering
\begin{tabular}{@{}lrrr@{}}
\toprule
\textbf{Category} & \textbf{Value 1} & \textbf{Value 2} & \textbf{Total} \\
\midrule
Item A & 10.5 & 20.3 & 30.8 \\
Item B & 15.2 & 25.1 & 40.3 \\
Item C & 22.1 & 12.7 & 34.8 \\
\bottomrule
\end{tabular}
\caption{Sample Data Table}
\label{tab:sample}
\end{table}

\section{Lists}
\subsection{Bullet Points}
\begin{itemize}
  \item First item
  \item Second item
  \item Third item with sub-items:
    \begin{itemize}
      \item Sub-item 1
      \item Sub-item 2
    \end{itemize}
\end{itemize}

\subsection{Numbered List}
\begin{enumerate}
  \item First step
  \item Second step
  \item Third step
\end{enumerate}

\section{Citations and References}
When you use the citation search feature, references will be automatically managed. For example, you might cite a paper like this: \cite{example2024}. All citations will be automatically added to your \texttt{cite.bib} file.

\section{Conclusion}
This template demonstrates basic \LaTeX{} features including mathematical formulas, tables, and lists. The bibliography system is now set up to work seamlessly with Bibby AI's citation management. Feel free to modify it for your needs!

\bibliographystyle{plain}
\bibliography{cite}

\end{document}\documentclass{article}
\usepackage[utf8]{inputenc}
\usepackage{booktabs}
\usepackage{amsmath}
\usepackage{natbib}

\title{My Document}
\author{Your Name}
\date{\today}

\begin{document}

\maketitle

\section{Introduction}
This is a sample document demonstrating various \LaTeX{} features including mathematical equations, tables, and formatting. You can cite papers using the citation search feature, and they will be automatically added to your \texttt{cite.bib} file.

\section{Mathematical Expressions}
Here are some example equations:

\subsection{Quadratic Formula}
The quadratic formula is given by:
\[x = \frac{-b \pm \sqrt{b^2 - 4ac}}{2a}\]

\subsection{Matrix Example}
A sample matrix:
\[
\begin{pmatrix}
a_{11} & a_{12} & a_{13} \\
a_{21} & a_{22} & a_{23} \\
a_{31} & a_{32} & a_{33}
\end{pmatrix}
\]

\section{Tables}
Below is a sample table showing data:

\begin{table}[h]
\centering
\begin{tabular}{@{}lrrr@{}}
\toprule
\textbf{Category} & \textbf{Value 1} & \textbf{Value 2} & \textbf{Total} \\
\midrule
Item A & 10.5 & 20.3 & 30.8 \\
Item B & 15.2 & 25.1 & 40.3 \\
Item C & 22.1 & 12.7 & 34.8 \\
\bottomrule
\end{tabular}
\caption{Sample Data Table}
\label{tab:sample}
\end{table}

\section{Lists}
\subsection{Bullet Points}
\begin{itemize}
  \item First item
  \item Second item
  \item Third item with sub-items:
    \begin{itemize}
      \item Sub-item 1
      \item Sub-item 2
    \end{itemize}
\end{itemize}

\subsection{Numbered List}
\begin{enumerate}
  \item First step
  \item Second step
  \item Third step
\end{enumerate}

\section{Citations and References}
When you use the citation search feature, references will be automatically managed. For example, you might cite a paper like this: \cite{example2024}. All citations will be automatically added to your \texttt{cite.bib} file.

\section{Conclusion}
This template demonstrates basic \LaTeX{} features including mathematical formulas, tables, and lists. The bibliography system is now set up to work seamlessly with Bibby AI's citation management. Feel free to modify it for your needs!

\bibliographystyle{plain}
\bibliography{cite}

\end{document}\documentclass{article}
\usepackage[utf8]{inputenc}
\usepackage{booktabs}
\usepackage{amsmath}
\usepackage{natbib}

\title{My Document}
\author{Your Name}
\date{\today}

\begin{document}

\maketitle

\section{Introduction}
This is a sample document demonstrating various \LaTeX{} features including mathematical equations, tables, and formatting. You can cite papers using the citation search feature, and they will be automatically added to your \texttt{cite.bib} file.

\section{Mathematical Expressions}
Here are some example equations:

\subsection{Quadratic Formula}
The quadratic formula is given by:
\[x = \frac{-b \pm \sqrt{b^2 - 4ac}}{2a}\]

\subsection{Matrix Example}
A sample matrix:
\[
\begin{pmatrix}
a_{11} & a_{12} & a_{13} \\
a_{21} & a_{22} & a_{23} \\
a_{31} & a_{32} & a_{33}
\end{pmatrix}
\]

\section{Tables}
Below is a sample table showing data:

\begin{table}[h]
\centering
\begin{tabular}{@{}lrrr@{}}
\toprule
\textbf{Category} & \textbf{Value 1} & \textbf{Value 2} & \textbf{Total} \\
\midrule
Item A & 10.5 & 20.3 & 30.8 \\
Item B & 15.2 & 25.1 & 40.3 \\
Item C & 22.1 & 12.7 & 34.8 \\
\bottomrule
\end{tabular}
\caption{Sample Data Table}
\label{tab:sample}
\end{table}

\section{Lists}
\subsection{Bullet Points}
\begin{itemize}
  \item First item
  \item Second item
  \item Third item with sub-items:
    \begin{itemize}
      \item Sub-item 1
      \item Sub-item 2
    \end{itemize}
\end{itemize}

\subsection{Numbered List}
\begin{enumerate}
  \item First step
  \item Second step
  \item Third step
\end{enumerate}

\section{Citations and References}
When you use the citation search feature, references will be automatically managed. For example, you might cite a paper like this: \cite{example2024}. All citations will be automatically added to your \texttt{cite.bib} file.

\section{Conclusion}
This template demonstrates basic \LaTeX{} features including mathematical formulas, tables, and lists. The bibliography system is now set up to work seamlessly with Bibby AI's citation management. Feel free to modify it for your needs!

\bibliographystyle{plain}
\bibliography{cite}

\end{document}\documentclass{article}
\usepackage[utf8]{inputenc}
\usepackage{booktabs}
\usepackage{amsmath}
\usepackage{natbib}

\title{My Document}
\author{Your Name}
\date{\today}

\begin{document}

\maketitle

\section{Introduction}
This is a sample document demonstrating various \LaTeX{} features including mathematical equations, tables, and formatting. You can cite papers using the citation search feature, and they will be automatically added to your \texttt{cite.bib} file.

\section{Mathematical Expressions}
Here are some example equations:

\subsection{Quadratic Formula}
The quadratic formula is given by:
\[x = \frac{-b \pm \sqrt{b^2 - 4ac}}{2a}\]

\subsection{Matrix Example}
A sample matrix:
\[
\begin{pmatrix}
a_{11} & a_{12} & a_{13} \\
a_{21} & a_{22} & a_{23} \\
a_{31} & a_{32} & a_{33}
\end{pmatrix}
\]

\section{Tables}
Below is a sample table showing data:

\begin{table}[h]
\centering
\begin{tabular}{@{}lrrr@{}}
\toprule
\textbf{Category} & \textbf{Value 1} & \textbf{Value 2} & \textbf{Total} \\
\midrule
Item A & 10.5 & 20.3 & 30.8 \\
Item B & 15.2 & 25.1 & 40.3 \\
Item C & 22.1 & 12.7 & 34.8 \\
\bottomrule
\end{tabular}
\caption{Sample Data Table}
\label{tab:sample}
\end{table}

\section{Lists}
\subsection{Bullet Points}
\begin{itemize}
  \item First item
  \item Second item
  \item Third item with sub-items:
    \begin{itemize}
      \item Sub-item 1
      \item Sub-item 2
    \end{itemize}
\end{itemize}

\subsection{Numbered List}
\begin{enumerate}
  \item First step
  \item Second step
  \item Third step
\end{enumerate}

\section{Citations and References}
When you use the citation search feature, references will be automatically managed. For example, you might cite a paper like this: \cite{example2024}. All citations will be automatically added to your \texttt{cite.bib} file.

\section{Conclusion}
This template demonstrates basic \LaTeX{} features including mathematical formulas, tables, and lists. The bibliography system is now set up to work seamlessly with Bibby AI's citation management. Feel free to modify it for your needs!

\bibliographystyle{plain}
\bibliography{cite}

\end{document}\documentclass{article}
\usepackage[utf8]{inputenc}
\usepackage{booktabs}
\usepackage{amsmath}
\usepackage{natbib}

\title{My Document}
\author{Your Name}
\date{\today}

\begin{document}

\maketitle

\section{Introduction}
This is a sample document demonstrating various \LaTeX{} features including mathematical equations, tables, and formatting. You can cite papers using the citation search feature, and they will be automatically added to your \texttt{cite.bib} file.

\section{Mathematical Expressions}
Here are some example equations:

\subsection{Quadratic Formula}
The quadratic formula is given by:
\[x = \frac{-b \pm \sqrt{b^2 - 4ac}}{2a}\]

\subsection{Matrix Example}
A sample matrix:
\[
\begin{pmatrix}
a_{11} & a_{12} & a_{13} \\
a_{21} & a_{22} & a_{23} \\
a_{31} & a_{32} & a_{33}
\end{pmatrix}
\]

\section{Tables}
Below is a sample table showing data:

\begin{table}[h]
\centering
\begin{tabular}{@{}lrrr@{}}
\toprule
\textbf{Category} & \textbf{Value 1} & \textbf{Value 2} & \textbf{Total} \\
\midrule
Item A & 10.5 & 20.3 & 30.8 \\
Item B & 15.2 & 25.1 & 40.3 \\
Item C & 22.1 & 12.7 & 34.8 \\
\bottomrule
\end{tabular}
\caption{Sample Data Table}
\label{tab:sample}
\end{table}

\section{Lists}
\subsection{Bullet Points}
\begin{itemize}
  \item First item
  \item Second item
  \item Third item with sub-items:
    \begin{itemize}
      \item Sub-item 1
      \item Sub-item 2
    \end{itemize}
\end{itemize}

\subsection{Numbered List}
\begin{enumerate}
  \item First step
  \item Second step
  \item Third step
\end{enumerate}

\section{Citations and References}
When you use the citation search feature, references will be automatically managed. For example, you might cite a paper like this: \cite{example2024}. All citations will be automatically added to your \texttt{cite.bib} file.

\section{Conclusion}
This template demonstrates basic \LaTeX{} features including mathematical formulas, tables, and lists. The bibliography system is now set up to work seamlessly with Bibby AI's citation management. Feel free to modify it for your needs!

\bibliographystyle{plain}
\bibliography{cite}

\end{document}\documentclass{article}
\usepackage[utf8]{inputenc}
\usepackage{booktabs}
\usepackage{amsmath}
\usepackage{natbib}

\title{My Document}
\author{Your Name}
\date{\today}

\begin{document}

\maketitle

\section{Introduction}
This is a sample document demonstrating various \LaTeX{} features including mathematical equations, tables, and formatting. You can cite papers using the citation search feature, and they will be automatically added to your \texttt{cite.bib} file.

\section{Mathematical Expressions}
Here are some example equations:

\subsection{Quadratic Formula}
The quadratic formula is given by:
\[x = \frac{-b \pm \sqrt{b^2 - 4ac}}{2a}\]

\subsection{Matrix Example}
A sample matrix:
\[
\begin{pmatrix}
a_{11} & a_{12} & a_{13} \\
a_{21} & a_{22} & a_{23} \\
a_{31} & a_{32} & a_{33}
\end{pmatrix}
\]

\section{Tables}
Below is a sample table showing data:

\begin{table}[h]
\centering
\begin{tabular}{@{}lrrr@{}}
\toprule
\textbf{Category} & \textbf{Value 1} & \textbf{Value 2} & \textbf{Total} \\
\midrule
Item A & 10.5 & 20.3 & 30.8 \\
Item B & 15.2 & 25.1 & 40.3 \\
Item C & 22.1 & 12.7 & 34.8 \\
\bottomrule
\end{tabular}
\caption{Sample Data Table}
\label{tab:sample}
\end{table}

\section{Lists}
\subsection{Bullet Points}
\begin{itemize}
  \item First item
  \item Second item
  \item Third item with sub-items:
    \begin{itemize}
      \item Sub-item 1
      \item Sub-item 2
    \end{itemize}
\end{itemize}

\subsection{Numbered List}
\begin{enumerate}
  \item First step
  \item Second step
  \item Third step
\end{enumerate}

\section{Citations and References}
When you use the citation search feature, references will be automatically managed. For example, you might cite a paper like this: \cite{example2024}. All citations will be automatically added to your \texttt{cite.bib} file.

\section{Conclusion}
This template demonstrates basic \LaTeX{} features including mathematical formulas, tables, and lists. The bibliography system is now set up to work seamlessly with Bibby AI's citation management. Feel free to modify it for your needs!

\bibliographystyle{plain}
\bibliography{cite}

\end{document}\documentclass{article}
\usepackage[utf8]{inputenc}
\usepackage{booktabs}
\usepackage{amsmath}
\usepackage{natbib}

\title{My Document}
\author{Your Name}
\date{\today}

\begin{document}

\maketitle

\section{Introduction}
This is a sample document demonstrating various \LaTeX{} features including mathematical equations, tables, and formatting. You can cite papers using the citation search feature, and they will be automatically added to your \texttt{cite.bib} file.

\section{Mathematical Expressions}
Here are some example equations:

\subsection{Quadratic Formula}
The quadratic formula is given by:
\[x = \frac{-b \pm \sqrt{b^2 - 4ac}}{2a}\]

\subsection{Matrix Example}
A sample matrix:
\[
\begin{pmatrix}
a_{11} & a_{12} & a_{13} \\
a_{21} & a_{22} & a_{23} \\
a_{31} & a_{32} & a_{33}
\end{pmatrix}
\]

\section{Tables}
Below is a sample table showing data:

\begin{table}[h]
\centering
\begin{tabular}{@{}lrrr@{}}
\toprule
\textbf{Category} & \textbf{Value 1} & \textbf{Value 2} & \textbf{Total} \\
\midrule
Item A & 10.5 & 20.3 & 30.8 \\
Item B & 15.2 & 25.1 & 40.3 \\
Item C & 22.1 & 12.7 & 34.8 \\
\bottomrule
\end{tabular}
\caption{Sample Data Table}
\label{tab:sample}
\end{table}

\section{Lists}
\subsection{Bullet Points}
\begin{itemize}
  \item First item
  \item Second item
  \item Third item with sub-items:
    \begin{itemize}
      \item Sub-item 1
      \item Sub-item 2
    \end{itemize}
\end{itemize}

\subsection{Numbered List}
\begin{enumerate}
  \item First step
  \item Second step
  \item Third step
\end{enumerate}

\section{Citations and References}
When you use the citation search feature, references will be automatically managed. For example, you might cite a paper like this: \cite{example2024}. All citations will be automatically added to your \texttt{cite.bib} file.

\section{Conclusion}
This template demonstrates basic \LaTeX{} features including mathematical formulas, tables, and lists. The bibliography system is now set up to work seamlessly with Bibby AI's citation management. Feel free to modify it for your needs!

\bibliographystyle{plain}
\bibliography{cite}

\end{document}\documentclass{article}
\usepackage[utf8]{inputenc}
\usepackage{booktabs}
\usepackage{amsmath}
\usepackage{natbib}

\title{My Document}
\author{Your Name}
\date{\today}

\begin{document}

\maketitle

\section{Introduction}
This is a sample document demonstrating various \LaTeX{} features including mathematical equations, tables, and formatting. You can cite papers using the citation search feature, and they will be automatically added to your \texttt{cite.bib} file.

\section{Mathematical Expressions}
Here are some example equations:

\subsection{Quadratic Formula}
The quadratic formula is given by:
\[x = \frac{-b \pm \sqrt{b^2 - 4ac}}{2a}\]

\subsection{Matrix Example}
A sample matrix:
\[
\begin{pmatrix}
a_{11} & a_{12} & a_{13} \\
a_{21} & a_{22} & a_{23} \\
a_{31} & a_{32} & a_{33}
\end{pmatrix}
\]

\section{Tables}
Below is a sample table showing data:

\begin{table}[h]
\centering
\begin{tabular}{@{}lrrr@{}}
\toprule
\textbf{Category} & \textbf{Value 1} & \textbf{Value 2} & \textbf{Total} \\
\midrule
Item A & 10.5 & 20.3 & 30.8 \\
Item B & 15.2 & 25.1 & 40.3 \\
Item C & 22.1 & 12.7 & 34.8 \\
\bottomrule
\end{tabular}
\caption{Sample Data Table}
\label{tab:sample}
\end{table}

\section{Lists}
\subsection{Bullet Points}
\begin{itemize}
  \item First item
  \item Second item
  \item Third item with sub-items:
    \begin{itemize}
      \item Sub-item 1
      \item Sub-item 2
    \end{itemize}
\end{itemize}

\subsection{Numbered List}
\begin{enumerate}
  \item First step
  \item Second step
  \item Third step
\end{enumerate}

\section{Citations and References}
When you use the citation search feature, references will be automatically managed. For example, you might cite a paper like this: \cite{example2024}. All citations will be automatically added to your \texttt{cite.bib} file.

\section{Conclusion}
This template demonstrates basic \LaTeX{} features including mathematical formulas, tables, and lists. The bibliography system is now set up to work seamlessly with Bibby AI's citation management. Feel free to modify it for your needs!

\bibliographystyle{plain}
\bibliography{cite}

\end{document}\documentclass{article}
\usepackage[utf8]{inputenc}
\usepackage{booktabs}
\usepackage{amsmath}
\usepackage{natbib}

\title{My Document}
\author{Your Name}
\date{\today}

\begin{document}

\maketitle

\section{Introduction}
This is a sample document demonstrating various \LaTeX{} features including mathematical equations, tables, and formatting. You can cite papers using the citation search feature, and they will be automatically added to your \texttt{cite.bib} file.

\section{Mathematical Expressions}
Here are some example equations:

\subsection{Quadratic Formula}
The quadratic formula is given by:
\[x = \frac{-b \pm \sqrt{b^2 - 4ac}}{2a}\]

\subsection{Matrix Example}
A sample matrix:
\[
\begin{pmatrix}
a_{11} & a_{12} & a_{13} \\
a_{21} & a_{22} & a_{23} \\
a_{31} & a_{32} & a_{33}
\end{pmatrix}
\]

\section{Tables}
Below is a sample table showing data:

\begin{table}[h]
\centering
\begin{tabular}{@{}lrrr@{}}
\toprule
\textbf{Category} & \textbf{Value 1} & \textbf{Value 2} & \textbf{Total} \\
\midrule
Item A & 10.5 & 20.3 & 30.8 \\
Item B & 15.2 & 25.1 & 40.3 \\
Item C & 22.1 & 12.7 & 34.8 \\
\bottomrule
\end{tabular}
\caption{Sample Data Table}
\label{tab:sample}
\end{table}

\section{Lists}
\subsection{Bullet Points}
\begin{itemize}
  \item First item
  \item Second item
  \item Third item with sub-items:
    \begin{itemize}
      \item Sub-item 1
      \item Sub-item 2
    \end{itemize}
\end{itemize}

\subsection{Numbered List}
\begin{enumerate}
  \item First step
  \item Second step
  \item Third step
\end{enumerate}

\section{Citations and References}
When you use the citation search feature, references will be automatically managed. For example, you might cite a paper like this: \cite{example2024}. All citations will be automatically added to your \texttt{cite.bib} file.

\section{Conclusion}
This template demonstrates basic \LaTeX{} features including mathematical formulas, tables, and lists. The bibliography system is now set up to work seamlessly with Bibby AI's citation management. Feel free to modify it for your needs!

\bibliographystyle{plain}
\bibliography{cite}

\end{document}\documentclass{article}
\usepackage[utf8]{inputenc}
\usepackage{booktabs}
\usepackage{amsmath}
\usepackage{natbib}

\title{My Document}
\author{Your Name}
\date{\today}

\begin{document}

\maketitle

\section{Introduction}
This is a sample document demonstrating various \LaTeX{} features including mathematical equations, tables, and formatting. You can cite papers using the citation search feature, and they will be automatically added to your \texttt{cite.bib} file.

\section{Mathematical Expressions}
Here are some example equations:

\subsection{Quadratic Formula}
The quadratic formula is given by:
\[x = \frac{-b \pm \sqrt{b^2 - 4ac}}{2a}\]

\subsection{Matrix Example}
A sample matrix:
\[
\begin{pmatrix}
a_{11} & a_{12} & a_{13} \\
a_{21} & a_{22} & a_{23} \\
a_{31} & a_{32} & a_{33}
\end{pmatrix}
\]

\section{Tables}
Below is a sample table showing data:

\begin{table}[h]
\centering
\begin{tabular}{@{}lrrr@{}}
\toprule
\textbf{Category} & \textbf{Value 1} & \textbf{Value 2} & \textbf{Total} \\
\midrule
Item A & 10.5 & 20.3 & 30.8 \\
Item B & 15.2 & 25.1 & 40.3 \\
Item C & 22.1 & 12.7 & 34.8 \\
\bottomrule
\end{tabular}
\caption{Sample Data Table}
\label{tab:sample}
\end{table}

\section{Lists}
\subsection{Bullet Points}
\begin{itemize}
  \item First item
  \item Second item
  \item Third item with sub-items:
    \begin{itemize}
      \item Sub-item 1
      \item Sub-item 2
    \end{itemize}
\end{itemize}

\subsection{Numbered List}
\begin{enumerate}
  \item First step
  \item Second step
  \item Third step
\end{enumerate}

\section{Citations and References}
When you use the citation search feature, references will be automatically managed. For example, you might cite a paper like this: \cite{example2024}. All citations will be automatically added to your \texttt{cite.bib} file.

\section{Conclusion}
This template demonstrates basic \LaTeX{} features including mathematical formulas, tables, and lists. The bibliography system is now set up to work seamlessly with Bibby AI's citation management. Feel free to modify it for your needs!

\bibliographystyle{plain}
\bibliography{cite}

\end{document}\documentclass{article}
\usepackage[utf8]{inputenc}
\usepackage{booktabs}
\usepackage{amsmath}
\usepackage{natbib}

\title{My Document}
\author{Your Name}
\date{\today}

\begin{document}

\maketitle

\section{Introduction}
This is a sample document demonstrating various \LaTeX{} features including mathematical equations, tables, and formatting. You can cite papers using the citation search feature, and they will be automatically added to your \texttt{cite.bib} file.

\section{Mathematical Expressions}
Here are some example equations:

\subsection{Quadratic Formula}
The quadratic formula is given by:
\[x = \frac{-b \pm \sqrt{b^2 - 4ac}}{2a}\]

\subsection{Matrix Example}
A sample matrix:
\[
\begin{pmatrix}
a_{11} & a_{12} & a_{13} \\
a_{21} & a_{22} & a_{23} \\
a_{31} & a_{32} & a_{33}
\end{pmatrix}
\]

\section{Tables}
Below is a sample table showing data:

\begin{table}[h]
\centering
\begin{tabular}{@{}lrrr@{}}
\toprule
\textbf{Category} & \textbf{Value 1} & \textbf{Value 2} & \textbf{Total} \\
\midrule
Item A & 10.5 & 20.3 & 30.8 \\
Item B & 15.2 & 25.1 & 40.3 \\
Item C & 22.1 & 12.7 & 34.8 \\
\bottomrule
\end{tabular}
\caption{Sample Data Table}
\label{tab:sample}
\end{table}

\section{Lists}
\subsection{Bullet Points}
\begin{itemize}
  \item First item
  \item Second item
  \item Third item with sub-items:
    \begin{itemize}
      \item Sub-item 1
      \item Sub-item 2
    \end{itemize}
\end{itemize}

\subsection{Numbered List}
\begin{enumerate}
  \item First step
  \item Second step
  \item Third step
\end{enumerate}

\section{Citations and References}
When you use the citation search feature, references will be automatically managed. For example, you might cite a paper like this: \cite{example2024}. All citations will be automatically added to your \texttt{cite.bib} file.

\section{Conclusion}
This template demonstrates basic \LaTeX{} features including mathematical formulas, tables, and lists. The bibliography system is now set up to work seamlessly with Bibby AI's citation management. Feel free to modify it for your needs!

\bibliographystyle{plain}
\bibliography{cite}

\end{document}\documentclass{article}
\usepackage[utf8]{inputenc}
\usepackage{booktabs}
\usepackage{amsmath}
\usepackage{natbib}

\title{My Document}
\author{Your Name}
\date{\today}

\begin{document}

\maketitle

\section{Introduction}
This is a sample document demonstrating various \LaTeX{} features including mathematical equations, tables, and formatting. You can cite papers using the citation search feature, and they will be automatically added to your \texttt{cite.bib} file.

\section{Mathematical Expressions}
Here are some example equations:

\subsection{Quadratic Formula}
The quadratic formula is given by:
\[x = \frac{-b \pm \sqrt{b^2 - 4ac}}{2a}\]

\subsection{Matrix Example}
A sample matrix:
\[
\begin{pmatrix}
a_{11} & a_{12} & a_{13} \\
a_{21} & a_{22} & a_{23} \\
a_{31} & a_{32} & a_{33}
\end{pmatrix}
\]

\section{Tables}
Below is a sample table showing data:

\begin{table}[h]
\centering
\begin{tabular}{@{}lrrr@{}}
\toprule
\textbf{Category} & \textbf{Value 1} & \textbf{Value 2} & \textbf{Total} \\
\midrule
Item A & 10.5 & 20.3 & 30.8 \\
Item B & 15.2 & 25.1 & 40.3 \\
Item C & 22.1 & 12.7 & 34.8 \\
\bottomrule
\end{tabular}
\caption{Sample Data Table}
\label{tab:sample}
\end{table}

\section{Lists}
\subsection{Bullet Points}
\begin{itemize}
  \item First item
  \item Second item
  \item Third item with sub-items:
    \begin{itemize}
      \item Sub-item 1
      \item Sub-item 2
    \end{itemize}
\end{itemize}

\subsection{Numbered List}
\begin{enumerate}
  \item First step
  \item Second step
  \item Third step
\end{enumerate}

\section{Citations and References}
When you use the citation search feature, references will be automatically managed. For example, you might cite a paper like this: \cite{example2024}. All citations will be automatically added to your \texttt{cite.bib} file.

\section{Conclusion}
This template demonstrates basic \LaTeX{} features including mathematical formulas, tables, and lists. The bibliography system is now set up to work seamlessly with Bibby AI's citation management. Feel free to modify it for your needs!

\bibliographystyle{plain}
\bibliography{cite}

\end{document}\documentclass{article}
\usepackage[utf8]{inputenc}
\usepackage{booktabs}
\usepackage{amsmath}
\usepackage{natbib}

\title{My Document}
\author{Your Name}
\date{\today}

\begin{document}

\maketitle

\section{Introduction}
This is a sample document demonstrating various \LaTeX{} features including mathematical equations, tables, and formatting. You can cite papers using the citation search feature, and they will be automatically added to your \texttt{cite.bib} file.

\section{Mathematical Expressions}
Here are some example equations:

\subsection{Quadratic Formula}
The quadratic formula is given by:
\[x = \frac{-b \pm \sqrt{b^2 - 4ac}}{2a}\]

\subsection{Matrix Example}
A sample matrix:
\[
\begin{pmatrix}
a_{11} & a_{12} & a_{13} \\
a_{21} & a_{22} & a_{23} \\
a_{31} & a_{32} & a_{33}
\end{pmatrix}
\]

\section{Tables}
Below is a sample table showing data:

\begin{table}[h]
\centering
\begin{tabular}{@{}lrrr@{}}
\toprule
\textbf{Category} & \textbf{Value 1} & \textbf{Value 2} & \textbf{Total} \\
\midrule
Item A & 10.5 & 20.3 & 30.8 \\
Item B & 15.2 & 25.1 & 40.3 \\
Item C & 22.1 & 12.7 & 34.8 \\
\bottomrule
\end{tabular}
\caption{Sample Data Table}
\label{tab:sample}
\end{table}

\section{Lists}
\subsection{Bullet Points}
\begin{itemize}
  \item First item
  \item Second item
  \item Third item with sub-items:
    \begin{itemize}
      \item Sub-item 1
      \item Sub-item 2
    \end{itemize}
\end{itemize}

\subsection{Numbered List}
\begin{enumerate}
  \item First step
  \item Second step
  \item Third step
\end{enumerate}

\section{Citations and References}
When you use the citation search feature, references will be automatically managed. For example, you might cite a paper like this: \cite{example2024}. All citations will be automatically added to your \texttt{cite.bib} file.

\section{Conclusion}
This template demonstrates basic \LaTeX{} features including mathematical formulas, tables, and lists. The bibliography system is now set up to work seamlessly with Bibby AI's citation management. Feel free to modify it for your needs!

\bibliographystyle{plain}
\bibliography{cite}

\end{document}\documentclass{article}
\usepackage[utf8]{inputenc}
\usepackage{booktabs}
\usepackage{amsmath}
\usepackage{natbib}

\title{My Document}
\author{Your Name}
\date{\today}

\begin{document}

\maketitle

\section{Introduction}
This is a sample document demonstrating various \LaTeX{} features including mathematical equations, tables, and formatting. You can cite papers using the citation search feature, and they will be automatically added to your \texttt{cite.bib} file.

\section{Mathematical Expressions}
Here are some example equations:

\subsection{Quadratic Formula}
The quadratic formula is given by:
\[x = \frac{-b \pm \sqrt{b^2 - 4ac}}{2a}\]

\subsection{Matrix Example}
A sample matrix:
\[
\begin{pmatrix}
a_{11} & a_{12} & a_{13} \\
a_{21} & a_{22} & a_{23} \\
a_{31} & a_{32} & a_{33}
\end{pmatrix}
\]

\section{Tables}
Below is a sample table showing data:

\begin{table}[h]
\centering
\begin{tabular}{@{}lrrr@{}}
\toprule
\textbf{Category} & \textbf{Value 1} & \textbf{Value 2} & \textbf{Total} \\
\midrule
Item A & 10.5 & 20.3 & 30.8 \\
Item B & 15.2 & 25.1 & 40.3 \\
Item C & 22.1 & 12.7 & 34.8 \\
\bottomrule
\end{tabular}
\caption{Sample Data Table}
\label{tab:sample}
\end{table}

\section{Lists}
\subsection{Bullet Points}
\begin{itemize}
  \item First item
  \item Second item
  \item Third item with sub-items:
    \begin{itemize}
      \item Sub-item 1
      \item Sub-item 2
    \end{itemize}
\end{itemize}

\subsection{Numbered List}
\begin{enumerate}
  \item First step
  \item Second step
  \item Third step
\end{enumerate}

\section{Citations and References}
When you use the citation search feature, references will be automatically managed. For example, you might cite a paper like this: \cite{example2024}. All citations will be automatically added to your \texttt{cite.bib} file.

\section{Conclusion}
This template demonstrates basic \LaTeX{} features including mathematical formulas, tables, and lists. The bibliography system is now set up to work seamlessly with Bibby AI's citation management. Feel free to modify it for your needs!

\bibliographystyle{plain}
\bibliography{cite}

\end{document}\documentclass{article}
\usepackage[utf8]{inputenc}
\usepackage{booktabs}
\usepackage{amsmath}
\usepackage{natbib}

\title{My Document}
\author{Your Name}
\date{\today}

\begin{document}

\maketitle

\section{Introduction}
This is a sample document demonstrating various \LaTeX{} features including mathematical equations, tables, and formatting. You can cite papers using the citation search feature, and they will be automatically added to your \texttt{cite.bib} file.

\section{Mathematical Expressions}
Here are some example equations:

\subsection{Quadratic Formula}
The quadratic formula is given by:
\[x = \frac{-b \pm \sqrt{b^2 - 4ac}}{2a}\]

\subsection{Matrix Example}
A sample matrix:
\[
\begin{pmatrix}
a_{11} & a_{12} & a_{13} \\
a_{21} & a_{22} & a_{23} \\
a_{31} & a_{32} & a_{33}
\end{pmatrix}
\]

\section{Tables}
Below is a sample table showing data:

\begin{table}[h]
\centering
\begin{tabular}{@{}lrrr@{}}
\toprule
\textbf{Category} & \textbf{Value 1} & \textbf{Value 2} & \textbf{Total} \\
\midrule
Item A & 10.5 & 20.3 & 30.8 \\
Item B & 15.2 & 25.1 & 40.3 \\
Item C & 22.1 & 12.7 & 34.8 \\
\bottomrule
\end{tabular}
\caption{Sample Data Table}
\label{tab:sample}
\end{table}

\section{Lists}
\subsection{Bullet Points}
\begin{itemize}
  \item First item
  \item Second item
  \item Third item with sub-items:
    \begin{itemize}
      \item Sub-item 1
      \item Sub-item 2
    \end{itemize}
\end{itemize}

\subsection{Numbered List}
\begin{enumerate}
  \item First step
  \item Second step
  \item Third step
\end{enumerate}

\section{Citations and References}
When you use the citation search feature, references will be automatically managed. For example, you might cite a paper like this: \cite{example2024}. All citations will be automatically added to your \texttt{cite.bib} file.

\section{Conclusion}
This template demonstrates basic \LaTeX{} features including mathematical formulas, tables, and lists. The bibliography system is now set up to work seamlessly with Bibby AI's citation management. Feel free to modify it for your needs!

\bibliographystyle{plain}
\bibliography{cite}

\end{document}\documentclass{article}
\usepackage[utf8]{inputenc}
\usepackage{booktabs}
\usepackage{amsmath}
\usepackage{natbib}

\title{My Document}
\author{Your Name}
\date{\today}

\begin{document}

\maketitle

\section{Introduction}
This is a sample document demonstrating various \LaTeX{} features including mathematical equations, tables, and formatting. You can cite papers using the citation search feature, and they will be automatically added to your \texttt{cite.bib} file.

\section{Mathematical Expressions}
Here are some example equations:

\subsection{Quadratic Formula}
The quadratic formula is given by:
\[x = \frac{-b \pm \sqrt{b^2 - 4ac}}{2a}\]

\subsection{Matrix Example}
A sample matrix:
\[
\begin{pmatrix}
a_{11} & a_{12} & a_{13} \\
a_{21} & a_{22} & a_{23} \\
a_{31} & a_{32} & a_{33}
\end{pmatrix}
\]

\section{Tables}
Below is a sample table showing data:

\begin{table}[h]
\centering
\begin{tabular}{@{}lrrr@{}}
\toprule
\textbf{Category} & \textbf{Value 1} & \textbf{Value 2} & \textbf{Total} \\
\midrule
Item A & 10.5 & 20.3 & 30.8 \\
Item B & 15.2 & 25.1 & 40.3 \\
Item C & 22.1 & 12.7 & 34.8 \\
\bottomrule
\end{tabular}
\caption{Sample Data Table}
\label{tab:sample}
\end{table}

\section{Lists}
\subsection{Bullet Points}
\begin{itemize}
  \item First item
  \item Second item
  \item Third item with sub-items:
    \begin{itemize}
      \item Sub-item 1
      \item Sub-item 2
    \end{itemize}
\end{itemize}

\subsection{Numbered List}
\begin{enumerate}
  \item First step
  \item Second step
  \item Third step
\end{enumerate}

\section{Citations and References}
When you use the citation search feature, references will be automatically managed. For example, you might cite a paper like this: \cite{example2024}. All citations will be automatically added to your \texttt{cite.bib} file.

\section{Conclusion}
This template demonstrates basic \LaTeX{} features including mathematical formulas, tables, and lists. The bibliography system is now set up to work seamlessly with Bibby AI's citation management. Feel free to modify it for your needs!

\bibliographystyle{plain}
\bibliography{cite}

\end{document}\documentclass{article}
\usepackage[utf8]{inputenc}
\usepackage{booktabs}
\usepackage{amsmath}
\usepackage{natbib}

\title{My Document}
\author{Your Name}
\date{\today}

\begin{document}

\maketitle

\section{Introduction}
This is a sample document demonstrating various \LaTeX{} features including mathematical equations, tables, and formatting. You can cite papers using the citation search feature, and they will be automatically added to your \texttt{cite.bib} file.

\section{Mathematical Expressions}
Here are some example equations:

\subsection{Quadratic Formula}
The quadratic formula is given by:
\[x = \frac{-b \pm \sqrt{b^2 - 4ac}}{2a}\]

\subsection{Matrix Example}
A sample matrix:
\[
\begin{pmatrix}
a_{11} & a_{12} & a_{13} \\
a_{21} & a_{22} & a_{23} \\
a_{31} & a_{32} & a_{33}
\end{pmatrix}
\]

\section{Tables}
Below is a sample table showing data:

\begin{table}[h]
\centering
\begin{tabular}{@{}lrrr@{}}
\toprule
\textbf{Category} & \textbf{Value 1} & \textbf{Value 2} & \textbf{Total} \\
\midrule
Item A & 10.5 & 20.3 & 30.8 \\
Item B & 15.2 & 25.1 & 40.3 \\
Item C & 22.1 & 12.7 & 34.8 \\
\bottomrule
\end{tabular}
\caption{Sample Data Table}
\label{tab:sample}
\end{table}

\section{Lists}
\subsection{Bullet Points}
\begin{itemize}
  \item First item
  \item Second item
  \item Third item with sub-items:
    \begin{itemize}
      \item Sub-item 1
      \item Sub-item 2
    \end{itemize}
\end{itemize}

\subsection{Numbered List}
\begin{enumerate}
  \item First step
  \item Second step
  \item Third step
\end{enumerate}

\section{Citations and References}
When you use the citation search feature, references will be automatically managed. For example, you might cite a paper like this: \cite{example2024}. All citations will be automatically added to your \texttt{cite.bib} file.

\section{Conclusion}
This template demonstrates basic \LaTeX{} features including mathematical formulas, tables, and lists. The bibliography system is now set up to work seamlessly with Bibby AI's citation management. Feel free to modify it for your needs!

\bibliographystyle{plain}
\bibliography{cite}

\end{document}\documentclass{article}
\usepackage[utf8]{inputenc}
\usepackage{booktabs}
\usepackage{amsmath}
\usepackage{natbib}

\title{My Document}
\author{Your Name}
\date{\today}

\begin{document}

\maketitle

\section{Introduction}
This is a sample document demonstrating various \LaTeX{} features including mathematical equations, tables, and formatting. You can cite papers using the citation search feature, and they will be automatically added to your \texttt{cite.bib} file.

\section{Mathematical Expressions}
Here are some example equations:

\subsection{Quadratic Formula}
The quadratic formula is given by:
\[x = \frac{-b \pm \sqrt{b^2 - 4ac}}{2a}\]

\subsection{Matrix Example}
A sample matrix:
\[
\begin{pmatrix}
a_{11} & a_{12} & a_{13} \\
a_{21} & a_{22} & a_{23} \\
a_{31} & a_{32} & a_{33}
\end{pmatrix}
\]

\section{Tables}
Below is a sample table showing data:

\begin{table}[h]
\centering
\begin{tabular}{@{}lrrr@{}}
\toprule
\textbf{Category} & \textbf{Value 1} & \textbf{Value 2} & \textbf{Total} \\
\midrule
Item A & 10.5 & 20.3 & 30.8 \\
Item B & 15.2 & 25.1 & 40.3 \\
Item C & 22.1 & 12.7 & 34.8 \\
\bottomrule
\end{tabular}
\caption{Sample Data Table}
\label{tab:sample}
\end{table}

\section{Lists}
\subsection{Bullet Points}
\begin{itemize}
  \item First item
  \item Second item
  \item Third item with sub-items:
    \begin{itemize}
      \item Sub-item 1
      \item Sub-item 2
    \end{itemize}
\end{itemize}

\subsection{Numbered List}
\begin{enumerate}
  \item First step
  \item Second step
  \item Third step
\end{enumerate}

\section{Citations and References}
When you use the citation search feature, references will be automatically managed. For example, you might cite a paper like this: \cite{example2024}. All citations will be automatically added to your \texttt{cite.bib} file.

\section{Conclusion}
This template demonstrates basic \LaTeX{} features including mathematical formulas, tables, and lists. The bibliography system is now set up to work seamlessly with Bibby AI's citation management. Feel free to modify it for your needs!

\bibliographystyle{plain}
\bibliography{cite}

\end{document}\documentclass{article}
\usepackage[utf8]{inputenc}
\usepackage{booktabs}
\usepackage{amsmath}
\usepackage{natbib}

\title{My Document}
\author{Your Name}
\date{\today}

\begin{document}

\maketitle

\section{Introduction}
This is a sample document demonstrating various \LaTeX{} features including mathematical equations, tables, and formatting. You can cite papers using the citation search feature, and they will be automatically added to your \texttt{cite.bib} file.

\section{Mathematical Expressions}
Here are some example equations:

\subsection{Quadratic Formula}
The quadratic formula is given by:
\[x = \frac{-b \pm \sqrt{b^2 - 4ac}}{2a}\]

\subsection{Matrix Example}
A sample matrix:
\[
\begin{pmatrix}
a_{11} & a_{12} & a_{13} \\
a_{21} & a_{22} & a_{23} \\
a_{31} & a_{32} & a_{33}
\end{pmatrix}
\]

\section{Tables}
Below is a sample table showing data:

\begin{table}[h]
\centering
\begin{tabular}{@{}lrrr@{}}
\toprule
\textbf{Category} & \textbf{Value 1} & \textbf{Value 2} & \textbf{Total} \\
\midrule
Item A & 10.5 & 20.3 & 30.8 \\
Item B & 15.2 & 25.1 & 40.3 \\
Item C & 22.1 & 12.7 & 34.8 \\
\bottomrule
\end{tabular}
\caption{Sample Data Table}
\label{tab:sample}
\end{table}

\section{Lists}
\subsection{Bullet Points}
\begin{itemize}
  \item First item
  \item Second item
  \item Third item with sub-items:
    \begin{itemize}
      \item Sub-item 1
      \item Sub-item 2
    \end{itemize}
\end{itemize}

\subsection{Numbered List}
\begin{enumerate}
  \item First step
  \item Second step
  \item Third step
\end{enumerate}

\section{Citations and References}
When you use the citation search feature, references will be automatically managed. For example, you might cite a paper like this: \cite{example2024}. All citations will be automatically added to your \texttt{cite.bib} file.

\section{Conclusion}
This template demonstrates basic \LaTeX{} features including mathematical formulas, tables, and lists. The bibliography system is now set up to work seamlessly with Bibby AI's citation management. Feel free to modify it for your needs!

\bibliographystyle{plain}
\bibliography{cite}

\end{document}\documentclass{article}
\usepackage[utf8]{inputenc}
\usepackage{booktabs}
\usepackage{amsmath}
\usepackage{natbib}

\title{My Document}
\author{Your Name}
\date{\today}

\begin{document}

\maketitle

\section{Introduction}
This is a sample document demonstrating various \LaTeX{} features including mathematical equations, tables, and formatting. You can cite papers using the citation search feature, and they will be automatically added to your \texttt{cite.bib} file.

\section{Mathematical Expressions}
Here are some example equations:

\subsection{Quadratic Formula}
The quadratic formula is given by:
\[x = \frac{-b \pm \sqrt{b^2 - 4ac}}{2a}\]

\subsection{Matrix Example}
A sample matrix:
\[
\begin{pmatrix}
a_{11} & a_{12} & a_{13} \\
a_{21} & a_{22} & a_{23} \\
a_{31} & a_{32} & a_{33}
\end{pmatrix}
\]

\section{Tables}
Below is a sample table showing data:

\begin{table}[h]
\centering
\begin{tabular}{@{}lrrr@{}}
\toprule
\textbf{Category} & \textbf{Value 1} & \textbf{Value 2} & \textbf{Total} \\
\midrule
Item A & 10.5 & 20.3 & 30.8 \\
Item B & 15.2 & 25.1 & 40.3 \\
Item C & 22.1 & 12.7 & 34.8 \\
\bottomrule
\end{tabular}
\caption{Sample Data Table}
\label{tab:sample}
\end{table}

\section{Lists}
\subsection{Bullet Points}
\begin{itemize}
  \item First item
  \item Second item
  \item Third item with sub-items:
    \begin{itemize}
      \item Sub-item 1
      \item Sub-item 2
    \end{itemize}
\end{itemize}

\subsection{Numbered List}
\begin{enumerate}
  \item First step
  \item Second step
  \item Third step
\end{enumerate}

\section{Citations and References}
When you use the citation search feature, references will be automatically managed. For example, you might cite a paper like this: \cite{example2024}. All citations will be automatically added to your \texttt{cite.bib} file.

\section{Conclusion}
This template demonstrates basic \LaTeX{} features including mathematical formulas, tables, and lists. The bibliography system is now set up to work seamlessly with Bibby AI's citation management. Feel free to modify it for your needs!

\bibliographystyle{plain}
\bibliography{cite}

\end{document}\documentclass{article}
\usepackage[utf8]{inputenc}
\usepackage{booktabs}
\usepackage{amsmath}
\usepackage{natbib}

\title{My Document}
\author{Your Name}
\date{\today}

\begin{document}

\maketitle

\section{Introduction}
This is a sample document demonstrating various \LaTeX{} features including mathematical equations, tables, and formatting. You can cite papers using the citation search feature, and they will be automatically added to your \texttt{cite.bib} file.

\section{Mathematical Expressions}
Here are some example equations:

\subsection{Quadratic Formula}
The quadratic formula is given by:
\[x = \frac{-b \pm \sqrt{b^2 - 4ac}}{2a}\]

\subsection{Matrix Example}
A sample matrix:
\[
\begin{pmatrix}
a_{11} & a_{12} & a_{13} \\
a_{21} & a_{22} & a_{23} \\
a_{31} & a_{32} & a_{33}
\end{pmatrix}
\]

\section{Tables}
Below is a sample table showing data:

\begin{table}[h]
\centering
\begin{tabular}{@{}lrrr@{}}
\toprule
\textbf{Category} & \textbf{Value 1} & \textbf{Value 2} & \textbf{Total} \\
\midrule
Item A & 10.5 & 20.3 & 30.8 \\
Item B & 15.2 & 25.1 & 40.3 \\
Item C & 22.1 & 12.7 & 34.8 \\
\bottomrule
\end{tabular}
\caption{Sample Data Table}
\label{tab:sample}
\end{table}

\section{Lists}
\subsection{Bullet Points}
\begin{itemize}
  \item First item
  \item Second item
  \item Third item with sub-items:
    \begin{itemize}
      \item Sub-item 1
      \item Sub-item 2
    \end{itemize}
\end{itemize}

\subsection{Numbered List}
\begin{enumerate}
  \item First step
  \item Second step
  \item Third step
\end{enumerate}

\section{Citations and References}
When you use the citation search feature, references will be automatically managed. For example, you might cite a paper like this: \cite{example2024}. All citations will be automatically added to your \texttt{cite.bib} file.

\section{Conclusion}
This template demonstrates basic \LaTeX{} features including mathematical formulas, tables, and lists. The bibliography system is now set up to work seamlessly with Bibby AI's citation management. Feel free to modify it for your needs!

\bibliographystyle{plain}
\bibliography{cite}

\end{document}\documentclass{article}
\usepackage[utf8]{inputenc}
\usepackage{booktabs}
\usepackage{amsmath}
\usepackage{natbib}

\title{My Document}
\author{Your Name}
\date{\today}

\begin{document}

\maketitle

\section{Introduction}
This is a sample document demonstrating various \LaTeX{} features including mathematical equations, tables, and formatting. You can cite papers using the citation search feature, and they will be automatically added to your \texttt{cite.bib} file.

\section{Mathematical Expressions}
Here are some example equations:

\subsection{Quadratic Formula}
The quadratic formula is given by:
\[x = \frac{-b \pm \sqrt{b^2 - 4ac}}{2a}\]

\subsection{Matrix Example}
A sample matrix:
\[
\begin{pmatrix}
a_{11} & a_{12} & a_{13} \\
a_{21} & a_{22} & a_{23} \\
a_{31} & a_{32} & a_{33}
\end{pmatrix}
\]

\section{Tables}
Below is a sample table showing data:

\begin{table}[h]
\centering
\begin{tabular}{@{}lrrr@{}}
\toprule
\textbf{Category} & \textbf{Value 1} & \textbf{Value 2} & \textbf{Total} \\
\midrule
Item A & 10.5 & 20.3 & 30.8 \\
Item B & 15.2 & 25.1 & 40.3 \\
Item C & 22.1 & 12.7 & 34.8 \\
\bottomrule
\end{tabular}
\caption{Sample Data Table}
\label{tab:sample}
\end{table}

\section{Lists}
\subsection{Bullet Points}
\begin{itemize}
  \item First item
  \item Second item
  \item Third item with sub-items:
    \begin{itemize}
      \item Sub-item 1
      \item Sub-item 2
    \end{itemize}
\end{itemize}

\subsection{Numbered List}
\begin{enumerate}
  \item First step
  \item Second step
  \item Third step
\end{enumerate}

\section{Citations and References}
When you use the citation search feature, references will be automatically managed. For example, you might cite a paper like this: \cite{example2024}. All citations will be automatically added to your \texttt{cite.bib} file.

\section{Conclusion}
This template demonstrates basic \LaTeX{} features including mathematical formulas, tables, and lists. The bibliography system is now set up to work seamlessly with Bibby AI's citation management. Feel free to modify it for your needs!

\bibliographystyle{plain}
\bibliography{cite}

\end{document}\documentclass{article}
\usepackage[utf8]{inputenc}
\usepackage{booktabs}
\usepackage{amsmath}
\usepackage{natbib}

\title{My Document}
\author{Your Name}
\date{\today}

\begin{document}

\maketitle

\section{Introduction}
This is a sample document demonstrating various \LaTeX{} features including mathematical equations, tables, and formatting. You can cite papers using the citation search feature, and they will be automatically added to your \texttt{cite.bib} file.

\section{Mathematical Expressions}
Here are some example equations:

\subsection{Quadratic Formula}
The quadratic formula is given by:
\[x = \frac{-b \pm \sqrt{b^2 - 4ac}}{2a}\]

\subsection{Matrix Example}
A sample matrix:
\[
\begin{pmatrix}
a_{11} & a_{12} & a_{13} \\
a_{21} & a_{22} & a_{23} \\
a_{31} & a_{32} & a_{33}
\end{pmatrix}
\]

\section{Tables}
Below is a sample table showing data:

\begin{table}[h]
\centering
\begin{tabular}{@{}lrrr@{}}
\toprule
\textbf{Category} & \textbf{Value 1} & \textbf{Value 2} & \textbf{Total} \\
\midrule
Item A & 10.5 & 20.3 & 30.8 \\
Item B & 15.2 & 25.1 & 40.3 \\
Item C & 22.1 & 12.7 & 34.8 \\
\bottomrule
\end{tabular}
\caption{Sample Data Table}
\label{tab:sample}
\end{table}

\section{Lists}
\subsection{Bullet Points}
\begin{itemize}
  \item First item
  \item Second item
  \item Third item with sub-items:
    \begin{itemize}
      \item Sub-item 1
      \item Sub-item 2
    \end{itemize}
\end{itemize}

\subsection{Numbered List}
\begin{enumerate}
  \item First step
  \item Second step
  \item Third step
\end{enumerate}

\section{Citations and References}
When you use the citation search feature, references will be automatically managed. For example, you might cite a paper like this: \cite{example2024}. All citations will be automatically added to your \texttt{cite.bib} file.

\section{Conclusion}
This template demonstrates basic \LaTeX{} features including mathematical formulas, tables, and lists. The bibliography system is now set up to work seamlessly with Bibby AI's citation management. Feel free to modify it for your needs!

\bibliographystyle{plain}
\bibliography{cite}

\end{document}\documentclass{article}
\usepackage[utf8]{inputenc}
\usepackage{booktabs}
\usepackage{amsmath}
\usepackage{natbib}

\title{My Document}
\author{Your Name}
\date{\today}

\begin{document}

\maketitle

\section{Introduction}
This is a sample document demonstrating various \LaTeX{} features including mathematical equations, tables, and formatting. You can cite papers using the citation search feature, and they will be automatically added to your \texttt{cite.bib} file.

\section{Mathematical Expressions}
Here are some example equations:

\subsection{Quadratic Formula}
The quadratic formula is given by:
\[x = \frac{-b \pm \sqrt{b^2 - 4ac}}{2a}\]

\subsection{Matrix Example}
A sample matrix:
\[
\begin{pmatrix}
a_{11} & a_{12} & a_{13} \\
a_{21} & a_{22} & a_{23} \\
a_{31} & a_{32} & a_{33}
\end{pmatrix}
\]

\section{Tables}
Below is a sample table showing data:

\begin{table}[h]
\centering
\begin{tabular}{@{}lrrr@{}}
\toprule
\textbf{Category} & \textbf{Value 1} & \textbf{Value 2} & \textbf{Total} \\
\midrule
Item A & 10.5 & 20.3 & 30.8 \\
Item B & 15.2 & 25.1 & 40.3 \\
Item C & 22.1 & 12.7 & 34.8 \\
\bottomrule
\end{tabular}
\caption{Sample Data Table}
\label{tab:sample}
\end{table}

\section{Lists}
\subsection{Bullet Points}
\begin{itemize}
  \item First item
  \item Second item
  \item Third item with sub-items:
    \begin{itemize}
      \item Sub-item 1
      \item Sub-item 2
    \end{itemize}
\end{itemize}

\subsection{Numbered List}
\begin{enumerate}
  \item First step
  \item Second step
  \item Third step
\end{enumerate}

\section{Citations and References}
When you use the citation search feature, references will be automatically managed. For example, you might cite a paper like this: \cite{example2024}. All citations will be automatically added to your \texttt{cite.bib} file.

\section{Conclusion}
This template demonstrates basic \LaTeX{} features including mathematical formulas, tables, and lists. The bibliography system is now set up to work seamlessly with Bibby AI's citation management. Feel free to modify it for your needs!

\bibliographystyle{plain}
\bibliography{cite}

\end{document}\documentclass{article}
\usepackage[utf8]{inputenc}
\usepackage{booktabs}
\usepackage{amsmath}
\usepackage{natbib}

\title{My Document}
\author{Your Name}
\date{\today}

\begin{document}

\maketitle

\section{Introduction}
This is a sample document demonstrating various \LaTeX{} features including mathematical equations, tables, and formatting. You can cite papers using the citation search feature, and they will be automatically added to your \texttt{cite.bib} file.

\section{Mathematical Expressions}
Here are some example equations:

\subsection{Quadratic Formula}
The quadratic formula is given by:
\[x = \frac{-b \pm \sqrt{b^2 - 4ac}}{2a}\]

\subsection{Matrix Example}
A sample matrix:
\[
\begin{pmatrix}
a_{11} & a_{12} & a_{13} \\
a_{21} & a_{22} & a_{23} \\
a_{31} & a_{32} & a_{33}
\end{pmatrix}
\]

\section{Tables}
Below is a sample table showing data:

\begin{table}[h]
\centering
\begin{tabular}{@{}lrrr@{}}
\toprule
\textbf{Category} & \textbf{Value 1} & \textbf{Value 2} & \textbf{Total} \\
\midrule
Item A & 10.5 & 20.3 & 30.8 \\
Item B & 15.2 & 25.1 & 40.3 \\
Item C & 22.1 & 12.7 & 34.8 \\
\bottomrule
\end{tabular}
\caption{Sample Data Table}
\label{tab:sample}
\end{table}

\section{Lists}
\subsection{Bullet Points}
\begin{itemize}
  \item First item
  \item Second item
  \item Third item with sub-items:
    \begin{itemize}
      \item Sub-item 1
      \item Sub-item 2
    \end{itemize}
\end{itemize}

\subsection{Numbered List}
\begin{enumerate}
  \item First step
  \item Second step
  \item Third step
\end{enumerate}

\section{Citations and References}
When you use the citation search feature, references will be automatically managed. For example, you might cite a paper like this: \cite{example2024}. All citations will be automatically added to your \texttt{cite.bib} file.

\section{Conclusion}
This template demonstrates basic \LaTeX{} features including mathematical formulas, tables, and lists. The bibliography system is now set up to work seamlessly with Bibby AI's citation management. Feel free to modify it for your needs!

\bibliographystyle{plain}
\bibliography{cite}

\end{document}\documentclass{article}
\usepackage[utf8]{inputenc}
\usepackage{booktabs}
\usepackage{amsmath}
\usepackage{natbib}

\title{My Document}
\author{Your Name}
\date{\today}

\begin{document}

\maketitle

\section{Introduction}
This is a sample document demonstrating various \LaTeX{} features including mathematical equations, tables, and formatting. You can cite papers using the citation search feature, and they will be automatically added to your \texttt{cite.bib} file.

\section{Mathematical Expressions}
Here are some example equations:

\subsection{Quadratic Formula}
The quadratic formula is given by:
\[x = \frac{-b \pm \sqrt{b^2 - 4ac}}{2a}\]

\subsection{Matrix Example}
A sample matrix:
\[
\begin{pmatrix}
a_{11} & a_{12} & a_{13} \\
a_{21} & a_{22} & a_{23} \\
a_{31} & a_{32} & a_{33}
\end{pmatrix}
\]

\section{Tables}
Below is a sample table showing data:

\begin{table}[h]
\centering
\begin{tabular}{@{}lrrr@{}}
\toprule
\textbf{Category} & \textbf{Value 1} & \textbf{Value 2} & \textbf{Total} \\
\midrule
Item A & 10.5 & 20.3 & 30.8 \\
Item B & 15.2 & 25.1 & 40.3 \\
Item C & 22.1 & 12.7 & 34.8 \\
\bottomrule
\end{tabular}
\caption{Sample Data Table}
\label{tab:sample}
\end{table}

\section{Lists}
\subsection{Bullet Points}
\begin{itemize}
  \item First item
  \item Second item
  \item Third item with sub-items:
    \begin{itemize}
      \item Sub-item 1
      \item Sub-item 2
    \end{itemize}
\end{itemize}

\subsection{Numbered List}
\begin{enumerate}
  \item First step
  \item Second step
  \item Third step
\end{enumerate}

\section{Citations and References}
When you use the citation search feature, references will be automatically managed. For example, you might cite a paper like this: \cite{example2024}. All citations will be automatically added to your \texttt{cite.bib} file.

\section{Conclusion}
This template demonstrates basic \LaTeX{} features including mathematical formulas, tables, and lists. The bibliography system is now set up to work seamlessly with Bibby AI's citation management. Feel free to modify it for your needs!

\bibliographystyle{plain}
\bibliography{cite}

\end{document}\documentclass{article}
\usepackage[utf8]{inputenc}
\usepackage{booktabs}
\usepackage{amsmath}
\usepackage{natbib}

\title{My Document}
\author{Your Name}
\date{\today}

\begin{document}

\maketitle

\section{Introduction}
This is a sample document demonstrating various \LaTeX{} features including mathematical equations, tables, and formatting. You can cite papers using the citation search feature, and they will be automatically added to your \texttt{cite.bib} file.

\section{Mathematical Expressions}
Here are some example equations:

\subsection{Quadratic Formula}
The quadratic formula is given by:
\[x = \frac{-b \pm \sqrt{b^2 - 4ac}}{2a}\]

\subsection{Matrix Example}
A sample matrix:
\[
\begin{pmatrix}
a_{11} & a_{12} & a_{13} \\
a_{21} & a_{22} & a_{23} \\
a_{31} & a_{32} & a_{33}
\end{pmatrix}
\]

\section{Tables}
Below is a sample table showing data:

\begin{table}[h]
\centering
\begin{tabular}{@{}lrrr@{}}
\toprule
\textbf{Category} & \textbf{Value 1} & \textbf{Value 2} & \textbf{Total} \\
\midrule
Item A & 10.5 & 20.3 & 30.8 \\
Item B & 15.2 & 25.1 & 40.3 \\
Item C & 22.1 & 12.7 & 34.8 \\
\bottomrule
\end{tabular}
\caption{Sample Data Table}
\label{tab:sample}
\end{table}

\section{Lists}
\subsection{Bullet Points}
\begin{itemize}
  \item First item
  \item Second item
  \item Third item with sub-items:
    \begin{itemize}
      \item Sub-item 1
      \item Sub-item 2
    \end{itemize}
\end{itemize}

\subsection{Numbered List}
\begin{enumerate}
  \item First step
  \item Second step
  \item Third step
\end{enumerate}

\section{Citations and References}
When you use the citation search feature, references will be automatically managed. For example, you might cite a paper like this: \cite{example2024}. All citations will be automatically added to your \texttt{cite.bib} file.

\section{Conclusion}
This template demonstrates basic \LaTeX{} features including mathematical formulas, tables, and lists. The bibliography system is now set up to work seamlessly with Bibby AI's citation management. Feel free to modify it for your needs!

\bibliographystyle{plain}
\bibliography{cite}

\end{document}\documentclass{article}
\usepackage[utf8]{inputenc}
\usepackage{booktabs}
\usepackage{amsmath}
\usepackage{natbib}

\title{My Document}
\author{Your Name}
\date{\today}

\begin{document}

\maketitle

\section{Introduction}
This is a sample document demonstrating various \LaTeX{} features including mathematical equations, tables, and formatting. You can cite papers using the citation search feature, and they will be automatically added to your \texttt{cite.bib} file.

\section{Mathematical Expressions}
Here are some example equations:

\subsection{Quadratic Formula}
The quadratic formula is given by:
\[x = \frac{-b \pm \sqrt{b^2 - 4ac}}{2a}\]

\subsection{Matrix Example}
A sample matrix:
\[
\begin{pmatrix}
a_{11} & a_{12} & a_{13} \\
a_{21} & a_{22} & a_{23} \\
a_{31} & a_{32} & a_{33}
\end{pmatrix}
\]

\section{Tables}
Below is a sample table showing data:

\begin{table}[h]
\centering
\begin{tabular}{@{}lrrr@{}}
\toprule
\textbf{Category} & \textbf{Value 1} & \textbf{Value 2} & \textbf{Total} \\
\midrule
Item A & 10.5 & 20.3 & 30.8 \\
Item B & 15.2 & 25.1 & 40.3 \\
Item C & 22.1 & 12.7 & 34.8 \\
\bottomrule
\end{tabular}
\caption{Sample Data Table}
\label{tab:sample}
\end{table}

\section{Lists}
\subsection{Bullet Points}
\begin{itemize}
  \item First item
  \item Second item
  \item Third item with sub-items:
    \begin{itemize}
      \item Sub-item 1
      \item Sub-item 2
    \end{itemize}
\end{itemize}

\subsection{Numbered List}
\begin{enumerate}
  \item First step
  \item Second step
  \item Third step
\end{enumerate}

\section{Citations and References}
When you use the citation search feature, references will be automatically managed. For example, you might cite a paper like this: \cite{example2024}. All citations will be automatically added to your \texttt{cite.bib} file.

\section{Conclusion}
This template demonstrates basic \LaTeX{} features including mathematical formulas, tables, and lists. The bibliography system is now set up to work seamlessly with Bibby AI's citation management. Feel free to modify it for your needs!

\bibliographystyle{plain}
\bibliography{cite}

\end{document}\documentclass{article}
\usepackage[utf8]{inputenc}
\usepackage{booktabs}
\usepackage{amsmath}
\usepackage{natbib}

\title{My Document}
\author{Your Name}
\date{\today}

\begin{document}

\maketitle

\section{Introduction}
This is a sample document demonstrating various \LaTeX{} features including mathematical equations, tables, and formatting. You can cite papers using the citation search feature, and they will be automatically added to your \texttt{cite.bib} file.

\section{Mathematical Expressions}
Here are some example equations:

\subsection{Quadratic Formula}
The quadratic formula is given by:
\[x = \frac{-b \pm \sqrt{b^2 - 4ac}}{2a}\]

\subsection{Matrix Example}
A sample matrix:
\[
\begin{pmatrix}
a_{11} & a_{12} & a_{13} \\
a_{21} & a_{22} & a_{23} \\
a_{31} & a_{32} & a_{33}
\end{pmatrix}
\]

\section{Tables}
Below is a sample table showing data:

\begin{table}[h]
\centering
\begin{tabular}{@{}lrrr@{}}
\toprule
\textbf{Category} & \textbf{Value 1} & \textbf{Value 2} & \textbf{Total} \\
\midrule
Item A & 10.5 & 20.3 & 30.8 \\
Item B & 15.2 & 25.1 & 40.3 \\
Item C & 22.1 & 12.7 & 34.8 \\
\bottomrule
\end{tabular}
\caption{Sample Data Table}
\label{tab:sample}
\end{table}

\section{Lists}
\subsection{Bullet Points}
\begin{itemize}
  \item First item
  \item Second item
  \item Third item with sub-items:
    \begin{itemize}
      \item Sub-item 1
      \item Sub-item 2
    \end{itemize}
\end{itemize}

\subsection{Numbered List}
\begin{enumerate}
  \item First step
  \item Second step
  \item Third step
\end{enumerate}

\section{Citations and References}
When you use the citation search feature, references will be automatically managed. For example, you might cite a paper like this: \cite{example2024}. All citations will be automatically added to your \texttt{cite.bib} file.

\section{Conclusion}
This template demonstrates basic \LaTeX{} features including mathematical formulas, tables, and lists. The bibliography system is now set up to work seamlessly with Bibby AI's citation management. Feel free to modify it for your needs!

\bibliographystyle{plain}
\bibliography{cite}

\end{document}\documentclass{article}
\usepackage[utf8]{inputenc}
\usepackage{booktabs}
\usepackage{amsmath}
\usepackage{natbib}

\title{My Document}
\author{Your Name}
\date{\today}

\begin{document}

\maketitle

\section{Introduction}
This is a sample document demonstrating various \LaTeX{} features including mathematical equations, tables, and formatting. You can cite papers using the citation search feature, and they will be automatically added to your \texttt{cite.bib} file.

\section{Mathematical Expressions}
Here are some example equations:

\subsection{Quadratic Formula}
The quadratic formula is given by:
\[x = \frac{-b \pm \sqrt{b^2 - 4ac}}{2a}\]

\subsection{Matrix Example}
A sample matrix:
\[
\begin{pmatrix}
a_{11} & a_{12} & a_{13} \\
a_{21} & a_{22} & a_{23} \\
a_{31} & a_{32} & a_{33}
\end{pmatrix}
\]

\section{Tables}
Below is a sample table showing data:

\begin{table}[h]
\centering
\begin{tabular}{@{}lrrr@{}}
\toprule
\textbf{Category} & \textbf{Value 1} & \textbf{Value 2} & \textbf{Total} \\
\midrule
Item A & 10.5 & 20.3 & 30.8 \\
Item B & 15.2 & 25.1 & 40.3 \\
Item C & 22.1 & 12.7 & 34.8 \\
\bottomrule
\end{tabular}
\caption{Sample Data Table}
\label{tab:sample}
\end{table}

\section{Lists}
\subsection{Bullet Points}
\begin{itemize}
  \item First item
  \item Second item
  \item Third item with sub-items:
    \begin{itemize}
      \item Sub-item 1
      \item Sub-item 2
    \end{itemize}
\end{itemize}

\subsection{Numbered List}
\begin{enumerate}
  \item First step
  \item Second step
  \item Third step
\end{enumerate}

\section{Citations and References}
When you use the citation search feature, references will be automatically managed. For example, you might cite a paper like this: \cite{example2024}. All citations will be automatically added to your \texttt{cite.bib} file.

\section{Conclusion}
This template demonstrates basic \LaTeX{} features including mathematical formulas, tables, and lists. The bibliography system is now set up to work seamlessly with Bibby AI's citation management. Feel free to modify it for your needs!

\bibliographystyle{plain}
\bibliography{cite}

\end{document}\documentclass{article}
\usepackage[utf8]{inputenc}
\usepackage{booktabs}
\usepackage{amsmath}
\usepackage{natbib}

\title{My Document}
\author{Your Name}
\date{\today}

\begin{document}

\maketitle

\section{Introduction}
This is a sample document demonstrating various \LaTeX{} features including mathematical equations, tables, and formatting. You can cite papers using the citation search feature, and they will be automatically added to your \texttt{cite.bib} file.

\section{Mathematical Expressions}
Here are some example equations:

\subsection{Quadratic Formula}
The quadratic formula is given by:
\[x = \frac{-b \pm \sqrt{b^2 - 4ac}}{2a}\]

\subsection{Matrix Example}
A sample matrix:
\[
\begin{pmatrix}
a_{11} & a_{12} & a_{13} \\
a_{21} & a_{22} & a_{23} \\
a_{31} & a_{32} & a_{33}
\end{pmatrix}
\]

\section{Tables}
Below is a sample table showing data:

\begin{table}[h]
\centering
\begin{tabular}{@{}lrrr@{}}
\toprule
\textbf{Category} & \textbf{Value 1} & \textbf{Value 2} & \textbf{Total} \\
\midrule
Item A & 10.5 & 20.3 & 30.8 \\
Item B & 15.2 & 25.1 & 40.3 \\
Item C & 22.1 & 12.7 & 34.8 \\
\bottomrule
\end{tabular}
\caption{Sample Data Table}
\label{tab:sample}
\end{table}

\section{Lists}
\subsection{Bullet Points}
\begin{itemize}
  \item First item
  \item Second item
  \item Third item with sub-items:
    \begin{itemize}
      \item Sub-item 1
      \item Sub-item 2
    \end{itemize}
\end{itemize}

\subsection{Numbered List}
\begin{enumerate}
  \item First step
  \item Second step
  \item Third step
\end{enumerate}

\section{Citations and References}
When you use the citation search feature, references will be automatically managed. For example, you might cite a paper like this: \cite{example2024}. All citations will be automatically added to your \texttt{cite.bib} file.

\section{Conclusion}
This template demonstrates basic \LaTeX{} features including mathematical formulas, tables, and lists. The bibliography system is now set up to work seamlessly with Bibby AI's citation management. Feel free to modify it for your needs!

\bibliographystyle{plain}
\bibliography{cite}

\end{document}\documentclass{article}
\usepackage[utf8]{inputenc}
\usepackage{booktabs}
\usepackage{amsmath}
\usepackage{natbib}

\title{My Document}
\author{Your Name}
\date{\today}

\begin{document}

\maketitle

\section{Introduction}
This is a sample document demonstrating various \LaTeX{} features including mathematical equations, tables, and formatting. You can cite papers using the citation search feature, and they will be automatically added to your \texttt{cite.bib} file.

\section{Mathematical Expressions}
Here are some example equations:

\subsection{Quadratic Formula}
The quadratic formula is given by:
\[x = \frac{-b \pm \sqrt{b^2 - 4ac}}{2a}\]

\subsection{Matrix Example}
A sample matrix:
\[
\begin{pmatrix}
a_{11} & a_{12} & a_{13} \\
a_{21} & a_{22} & a_{23} \\
a_{31} & a_{32} & a_{33}
\end{pmatrix}
\]

\section{Tables}
Below is a sample table showing data:

\begin{table}[h]
\centering
\begin{tabular}{@{}lrrr@{}}
\toprule
\textbf{Category} & \textbf{Value 1} & \textbf{Value 2} & \textbf{Total} \\
\midrule
Item A & 10.5 & 20.3 & 30.8 \\
Item B & 15.2 & 25.1 & 40.3 \\
Item C & 22.1 & 12.7 & 34.8 \\
\bottomrule
\end{tabular}
\caption{Sample Data Table}
\label{tab:sample}
\end{table}

\section{Lists}
\subsection{Bullet Points}
\begin{itemize}
  \item First item
  \item Second item
  \item Third item with sub-items:
    \begin{itemize}
      \item Sub-item 1
      \item Sub-item 2
    \end{itemize}
\end{itemize}

\subsection{Numbered List}
\begin{enumerate}
  \item First step
  \item Second step
  \item Third step
\end{enumerate}

\section{Citations and References}
When you use the citation search feature, references will be automatically managed. For example, you might cite a paper like this: \cite{example2024}. All citations will be automatically added to your \texttt{cite.bib} file.

\section{Conclusion}
This template demonstrates basic \LaTeX{} features including mathematical formulas, tables, and lists. The bibliography system is now set up to work seamlessly with Bibby AI's citation management. Feel free to modify it for your needs!

\bibliographystyle{plain}
\bibliography{cite}

\end{document}\documentclass{article}
\usepackage[utf8]{inputenc}
\usepackage{booktabs}
\usepackage{amsmath}
\usepackage{natbib}

\title{My Document}
\author{Your Name}
\date{\today}

\begin{document}

\maketitle

\section{Introduction}
This is a sample document demonstrating various \LaTeX{} features including mathematical equations, tables, and formatting. You can cite papers using the citation search feature, and they will be automatically added to your \texttt{cite.bib} file.

\section{Mathematical Expressions}
Here are some example equations:

\subsection{Quadratic Formula}
The quadratic formula is given by:
\[x = \frac{-b \pm \sqrt{b^2 - 4ac}}{2a}\]

\subsection{Matrix Example}
A sample matrix:
\[
\begin{pmatrix}
a_{11} & a_{12} & a_{13} \\
a_{21} & a_{22} & a_{23} \\
a_{31} & a_{32} & a_{33}
\end{pmatrix}
\]

\section{Tables}
Below is a sample table showing data:

\begin{table}[h]
\centering
\begin{tabular}{@{}lrrr@{}}
\toprule
\textbf{Category} & \textbf{Value 1} & \textbf{Value 2} & \textbf{Total} \\
\midrule
Item A & 10.5 & 20.3 & 30.8 \\
Item B & 15.2 & 25.1 & 40.3 \\
Item C & 22.1 & 12.7 & 34.8 \\
\bottomrule
\end{tabular}
\caption{Sample Data Table}
\label{tab:sample}
\end{table}

\section{Lists}
\subsection{Bullet Points}
\begin{itemize}
  \item First item
  \item Second item
  \item Third item with sub-items:
    \begin{itemize}
      \item Sub-item 1
      \item Sub-item 2
    \end{itemize}
\end{itemize}

\subsection{Numbered List}
\begin{enumerate}
  \item First step
  \item Second step
  \item Third step
\end{enumerate}

\section{Citations and References}
When you use the citation search feature, references will be automatically managed. For example, you might cite a paper like this: \cite{example2024}. All citations will be automatically added to your \texttt{cite.bib} file.

\section{Conclusion}
This template demonstrates basic \LaTeX{} features including mathematical formulas, tables, and lists. The bibliography system is now set up to work seamlessly with Bibby AI's citation management. Feel free to modify it for your needs!

\bibliographystyle{plain}
\bibliography{cite}

\end{document}\documentclass{article}
\usepackage[utf8]{inputenc}
\usepackage{booktabs}
\usepackage{amsmath}
\usepackage{natbib}

\title{My Document}
\author{Your Name}
\date{\today}

\begin{document}

\maketitle

\section{Introduction}
This is a sample document demonstrating various \LaTeX{} features including mathematical equations, tables, and formatting. You can cite papers using the citation search feature, and they will be automatically added to your \texttt{cite.bib} file.

\section{Mathematical Expressions}
Here are some example equations:

\subsection{Quadratic Formula}
The quadratic formula is given by:
\[x = \frac{-b \pm \sqrt{b^2 - 4ac}}{2a}\]

\subsection{Matrix Example}
A sample matrix:
\[
\begin{pmatrix}
a_{11} & a_{12} & a_{13} \\
a_{21} & a_{22} & a_{23} \\
a_{31} & a_{32} & a_{33}
\end{pmatrix}
\]

\section{Tables}
Below is a sample table showing data:

\begin{table}[h]
\centering
\begin{tabular}{@{}lrrr@{}}
\toprule
\textbf{Category} & \textbf{Value 1} & \textbf{Value 2} & \textbf{Total} \\
\midrule
Item A & 10.5 & 20.3 & 30.8 \\
Item B & 15.2 & 25.1 & 40.3 \\
Item C & 22.1 & 12.7 & 34.8 \\
\bottomrule
\end{tabular}
\caption{Sample Data Table}
\label{tab:sample}
\end{table}

\section{Lists}
\subsection{Bullet Points}
\begin{itemize}
  \item First item
  \item Second item
  \item Third item with sub-items:
    \begin{itemize}
      \item Sub-item 1
      \item Sub-item 2
    \end{itemize}
\end{itemize}

\subsection{Numbered List}
\begin{enumerate}
  \item First step
  \item Second step
  \item Third step
\end{enumerate}

\section{Citations and References}
When you use the citation search feature, references will be automatically managed. For example, you might cite a paper like this: \cite{example2024}. All citations will be automatically added to your \texttt{cite.bib} file.

\section{Conclusion}
This template demonstrates basic \LaTeX{} features including mathematical formulas, tables, and lists. The bibliography system is now set up to work seamlessly with Bibby AI's citation management. Feel free to modify it for your needs!

\bibliographystyle{plain}
\bibliography{cite}

\end{document}\documentclass{article}
\usepackage[utf8]{inputenc}
\usepackage{booktabs}
\usepackage{amsmath}
\usepackage{natbib}

\title{My Document}
\author{Your Name}
\date{\today}

\begin{document}

\maketitle

\section{Introduction}
This is a sample document demonstrating various \LaTeX{} features including mathematical equations, tables, and formatting. You can cite papers using the citation search feature, and they will be automatically added to your \texttt{cite.bib} file.

\section{Mathematical Expressions}
Here are some example equations:

\subsection{Quadratic Formula}
The quadratic formula is given by:
\[x = \frac{-b \pm \sqrt{b^2 - 4ac}}{2a}\]

\subsection{Matrix Example}
A sample matrix:
\[
\begin{pmatrix}
a_{11} & a_{12} & a_{13} \\
a_{21} & a_{22} & a_{23} \\
a_{31} & a_{32} & a_{33}
\end{pmatrix}
\]

\section{Tables}
Below is a sample table showing data:

\begin{table}[h]
\centering
\begin{tabular}{@{}lrrr@{}}
\toprule
\textbf{Category} & \textbf{Value 1} & \textbf{Value 2} & \textbf{Total} \\
\midrule
Item A & 10.5 & 20.3 & 30.8 \\
Item B & 15.2 & 25.1 & 40.3 \\
Item C & 22.1 & 12.7 & 34.8 \\
\bottomrule
\end{tabular}
\caption{Sample Data Table}
\label{tab:sample}
\end{table}

\section{Lists}
\subsection{Bullet Points}
\begin{itemize}
  \item First item
  \item Second item
  \item Third item with sub-items:
    \begin{itemize}
      \item Sub-item 1
      \item Sub-item 2
    \end{itemize}
\end{itemize}

\subsection{Numbered List}
\begin{enumerate}
  \item First step
  \item Second step
  \item Third step
\end{enumerate}

\section{Citations and References}
When you use the citation search feature, references will be automatically managed. For example, you might cite a paper like this: \cite{example2024}. All citations will be automatically added to your \texttt{cite.bib} file.

\section{Conclusion}
This template demonstrates basic \LaTeX{} features including mathematical formulas, tables, and lists. The bibliography system is now set up to work seamlessly with Bibby AI's citation management. Feel free to modify it for your needs!

\bibliographystyle{plain}
\bibliography{cite}

\end{document}\documentclass{article}
\usepackage[utf8]{inputenc}
\usepackage{booktabs}
\usepackage{amsmath}
\usepackage{natbib}

\title{My Document}
\author{Your Name}
\date{\today}

\begin{document}

\maketitle

\section{Introduction}
This is a sample document demonstrating various \LaTeX{} features including mathematical equations, tables, and formatting. You can cite papers using the citation search feature, and they will be automatically added to your \texttt{cite.bib} file.

\section{Mathematical Expressions}
Here are some example equations:

\subsection{Quadratic Formula}
The quadratic formula is given by:
\[x = \frac{-b \pm \sqrt{b^2 - 4ac}}{2a}\]

\subsection{Matrix Example}
A sample matrix:
\[
\begin{pmatrix}
a_{11} & a_{12} & a_{13} \\
a_{21} & a_{22} & a_{23} \\
a_{31} & a_{32} & a_{33}
\end{pmatrix}
\]

\section{Tables}
Below is a sample table showing data:

\begin{table}[h]
\centering
\begin{tabular}{@{}lrrr@{}}
\toprule
\textbf{Category} & \textbf{Value 1} & \textbf{Value 2} & \textbf{Total} \\
\midrule
Item A & 10.5 & 20.3 & 30.8 \\
Item B & 15.2 & 25.1 & 40.3 \\
Item C & 22.1 & 12.7 & 34.8 \\
\bottomrule
\end{tabular}
\caption{Sample Data Table}
\label{tab:sample}
\end{table}

\section{Lists}
\subsection{Bullet Points}
\begin{itemize}
  \item First item
  \item Second item
  \item Third item with sub-items:
    \begin{itemize}
      \item Sub-item 1
      \item Sub-item 2
    \end{itemize}
\end{itemize}

\subsection{Numbered List}
\begin{enumerate}
  \item First step
  \item Second step
  \item Third step
\end{enumerate}

\section{Citations and References}
When you use the citation search feature, references will be automatically managed. For example, you might cite a paper like this: \cite{example2024}. All citations will be automatically added to your \texttt{cite.bib} file.

\section{Conclusion}
This template demonstrates basic \LaTeX{} features including mathematical formulas, tables, and lists. The bibliography system is now set up to work seamlessly with Bibby AI's citation management. Feel free to modify it for your needs!

\bibliographystyle{plain}
\bibliography{cite}

\end{document}\documentclass{article}
\usepackage[utf8]{inputenc}
\usepackage{booktabs}
\usepackage{amsmath}
\usepackage{natbib}

\title{My Document}
\author{Your Name}
\date{\today}

\begin{document}

\maketitle

\section{Introduction}
This is a sample document demonstrating various \LaTeX{} features including mathematical equations, tables, and formatting. You can cite papers using the citation search feature, and they will be automatically added to your \texttt{cite.bib} file.

\section{Mathematical Expressions}
Here are some example equations:

\subsection{Quadratic Formula}
The quadratic formula is given by:
\[x = \frac{-b \pm \sqrt{b^2 - 4ac}}{2a}\]

\subsection{Matrix Example}
A sample matrix:
\[
\begin{pmatrix}
a_{11} & a_{12} & a_{13} \\
a_{21} & a_{22} & a_{23} \\
a_{31} & a_{32} & a_{33}
\end{pmatrix}
\]

\section{Tables}
Below is a sample table showing data:

\begin{table}[h]
\centering
\begin{tabular}{@{}lrrr@{}}
\toprule
\textbf{Category} & \textbf{Value 1} & \textbf{Value 2} & \textbf{Total} \\
\midrule
Item A & 10.5 & 20.3 & 30.8 \\
Item B & 15.2 & 25.1 & 40.3 \\
Item C & 22.1 & 12.7 & 34.8 \\
\bottomrule
\end{tabular}
\caption{Sample Data Table}
\label{tab:sample}
\end{table}

\section{Lists}
\subsection{Bullet Points}
\begin{itemize}
  \item First item
  \item Second item
  \item Third item with sub-items:
    \begin{itemize}
      \item Sub-item 1
      \item Sub-item 2
    \end{itemize}
\end{itemize}

\subsection{Numbered List}
\begin{enumerate}
  \item First step
  \item Second step
  \item Third step
\end{enumerate}

\section{Citations and References}
When you use the citation search feature, references will be automatically managed. For example, you might cite a paper like this: \cite{example2024}. All citations will be automatically added to your \texttt{cite.bib} file.

\section{Conclusion}
This template demonstrates basic \LaTeX{} features including mathematical formulas, tables, and lists. The bibliography system is now set up to work seamlessly with Bibby AI's citation management. Feel free to modify it for your needs!

\bibliographystyle{plain}
\bibliography{cite}

\end{document}\documentclass{article}
\usepackage[utf8]{inputenc}
\usepackage{booktabs}
\usepackage{amsmath}
\usepackage{natbib}

\title{My Document}
\author{Your Name}
\date{\today}

\begin{document}

\maketitle

\section{Introduction}
This is a sample document demonstrating various \LaTeX{} features including mathematical equations, tables, and formatting. You can cite papers using the citation search feature, and they will be automatically added to your \texttt{cite.bib} file.

\section{Mathematical Expressions}
Here are some example equations:

\subsection{Quadratic Formula}
The quadratic formula is given by:
\[x = \frac{-b \pm \sqrt{b^2 - 4ac}}{2a}\]

\subsection{Matrix Example}
A sample matrix:
\[
\begin{pmatrix}
a_{11} & a_{12} & a_{13} \\
a_{21} & a_{22} & a_{23} \\
a_{31} & a_{32} & a_{33}
\end{pmatrix}
\]

\section{Tables}
Below is a sample table showing data:

\begin{table}[h]
\centering
\begin{tabular}{@{}lrrr@{}}
\toprule
\textbf{Category} & \textbf{Value 1} & \textbf{Value 2} & \textbf{Total} \\
\midrule
Item A & 10.5 & 20.3 & 30.8 \\
Item B & 15.2 & 25.1 & 40.3 \\
Item C & 22.1 & 12.7 & 34.8 \\
\bottomrule
\end{tabular}
\caption{Sample Data Table}
\label{tab:sample}
\end{table}

\section{Lists}
\subsection{Bullet Points}
\begin{itemize}
  \item First item
  \item Second item
  \item Third item with sub-items:
    \begin{itemize}
      \item Sub-item 1
      \item Sub-item 2
    \end{itemize}
\end{itemize}

\subsection{Numbered List}
\begin{enumerate}
  \item First step
  \item Second step
  \item Third step
\end{enumerate}

\section{Citations and References}
When you use the citation search feature, references will be automatically managed. For example, you might cite a paper like this: \cite{example2024}. All citations will be automatically added to your \texttt{cite.bib} file.

\section{Conclusion}
This template demonstrates basic \LaTeX{} features including mathematical formulas, tables, and lists. The bibliography system is now set up to work seamlessly with Bibby AI's citation management. Feel free to modify it for your needs!

\bibliographystyle{plain}
\bibliography{cite}

\end{document}\documentclass{article}
\usepackage[utf8]{inputenc}
\usepackage{booktabs}
\usepackage{amsmath}
\usepackage{natbib}

\title{My Document}
\author{Your Name}
\date{\today}

\begin{document}

\maketitle

\section{Introduction}
This is a sample document demonstrating various \LaTeX{} features including mathematical equations, tables, and formatting. You can cite papers using the citation search feature, and they will be automatically added to your \texttt{cite.bib} file.

\section{Mathematical Expressions}
Here are some example equations:

\subsection{Quadratic Formula}
The quadratic formula is given by:
\[x = \frac{-b \pm \sqrt{b^2 - 4ac}}{2a}\]

\subsection{Matrix Example}
A sample matrix:
\[
\begin{pmatrix}
a_{11} & a_{12} & a_{13} \\
a_{21} & a_{22} & a_{23} \\
a_{31} & a_{32} & a_{33}
\end{pmatrix}
\]

\section{Tables}
Below is a sample table showing data:

\begin{table}[h]
\centering
\begin{tabular}{@{}lrrr@{}}
\toprule
\textbf{Category} & \textbf{Value 1} & \textbf{Value 2} & \textbf{Total} \\
\midrule
Item A & 10.5 & 20.3 & 30.8 \\
Item B & 15.2 & 25.1 & 40.3 \\
Item C & 22.1 & 12.7 & 34.8 \\
\bottomrule
\end{tabular}
\caption{Sample Data Table}
\label{tab:sample}
\end{table}

\section{Lists}
\subsection{Bullet Points}
\begin{itemize}
  \item First item
  \item Second item
  \item Third item with sub-items:
    \begin{itemize}
      \item Sub-item 1
      \item Sub-item 2
    \end{itemize}
\end{itemize}

\subsection{Numbered List}
\begin{enumerate}
  \item First step
  \item Second step
  \item Third step
\end{enumerate}

\section{Citations and References}
When you use the citation search feature, references will be automatically managed. For example, you might cite a paper like this: \cite{example2024}. All citations will be automatically added to your \texttt{cite.bib} file.

\section{Conclusion}
This template demonstrates basic \LaTeX{} features including mathematical formulas, tables, and lists. The bibliography system is now set up to work seamlessly with Bibby AI's citation management. Feel free to modify it for your needs!

\bibliographystyle{plain}
\bibliography{cite}

\end{document}\documentclass{article}
\usepackage[utf8]{inputenc}
\usepackage{booktabs}
\usepackage{amsmath}
\usepackage{natbib}

\title{My Document}
\author{Your Name}
\date{\today}

\begin{document}

\maketitle

\section{Introduction}
This is a sample document demonstrating various \LaTeX{} features including mathematical equations, tables, and formatting. You can cite papers using the citation search feature, and they will be automatically added to your \texttt{cite.bib} file.

\section{Mathematical Expressions}
Here are some example equations:

\subsection{Quadratic Formula}
The quadratic formula is given by:
\[x = \frac{-b \pm \sqrt{b^2 - 4ac}}{2a}\]

\subsection{Matrix Example}
A sample matrix:
\[
\begin{pmatrix}
a_{11} & a_{12} & a_{13} \\
a_{21} & a_{22} & a_{23} \\
a_{31} & a_{32} & a_{33}
\end{pmatrix}
\]

\section{Tables}
Below is a sample table showing data:

\begin{table}[h]
\centering
\begin{tabular}{@{}lrrr@{}}
\toprule
\textbf{Category} & \textbf{Value 1} & \textbf{Value 2} & \textbf{Total} \\
\midrule
Item A & 10.5 & 20.3 & 30.8 \\
Item B & 15.2 & 25.1 & 40.3 \\
Item C & 22.1 & 12.7 & 34.8 \\
\bottomrule
\end{tabular}
\caption{Sample Data Table}
\label{tab:sample}
\end{table}

\section{Lists}
\subsection{Bullet Points}
\begin{itemize}
  \item First item
  \item Second item
  \item Third item with sub-items:
    \begin{itemize}
      \item Sub-item 1
      \item Sub-item 2
    \end{itemize}
\end{itemize}

\subsection{Numbered List}
\begin{enumerate}
  \item First step
  \item Second step
  \item Third step
\end{enumerate}

\section{Citations and References}
When you use the citation search feature, references will be automatically managed. For example, you might cite a paper like this: \cite{example2024}. All citations will be automatically added to your \texttt{cite.bib} file.

\section{Conclusion}
This template demonstrates basic \LaTeX{} features including mathematical formulas, tables, and lists. The bibliography system is now set up to work seamlessly with Bibby AI's citation management. Feel free to modify it for your needs!

\bibliographystyle{plain}
\bibliography{cite}

\end{document}\documentclass{article}
\usepackage[utf8]{inputenc}
\usepackage{booktabs}
\usepackage{amsmath}
\usepackage{natbib}

\title{My Document}
\author{Your Name}
\date{\today}

\begin{document}

\maketitle

\section{Introduction}
This is a sample document demonstrating various \LaTeX{} features including mathematical equations, tables, and formatting. You can cite papers using the citation search feature, and they will be automatically added to your \texttt{cite.bib} file.

\section{Mathematical Expressions}
Here are some example equations:

\subsection{Quadratic Formula}
The quadratic formula is given by:
\[x = \frac{-b \pm \sqrt{b^2 - 4ac}}{2a}\]

\subsection{Matrix Example}
A sample matrix:
\[
\begin{pmatrix}
a_{11} & a_{12} & a_{13} \\
a_{21} & a_{22} & a_{23} \\
a_{31} & a_{32} & a_{33}
\end{pmatrix}
\]

\section{Tables}
Below is a sample table showing data:

\begin{table}[h]
\centering
\begin{tabular}{@{}lrrr@{}}
\toprule
\textbf{Category} & \textbf{Value 1} & \textbf{Value 2} & \textbf{Total} \\
\midrule
Item A & 10.5 & 20.3 & 30.8 \\
Item B & 15.2 & 25.1 & 40.3 \\
Item C & 22.1 & 12.7 & 34.8 \\
\bottomrule
\end{tabular}
\caption{Sample Data Table}
\label{tab:sample}
\end{table}

\section{Lists}
\subsection{Bullet Points}
\begin{itemize}
  \item First item
  \item Second item
  \item Third item with sub-items:
    \begin{itemize}
      \item Sub-item 1
      \item Sub-item 2
    \end{itemize}
\end{itemize}

\subsection{Numbered List}
\begin{enumerate}
  \item First step
  \item Second step
  \item Third step
\end{enumerate}

\section{Citations and References}
When you use the citation search feature, references will be automatically managed. For example, you might cite a paper like this: \cite{example2024}. All citations will be automatically added to your \texttt{cite.bib} file.

\section{Conclusion}
This template demonstrates basic \LaTeX{} features including mathematical formulas, tables, and lists. The bibliography system is now set up to work seamlessly with Bibby AI's citation management. Feel free to modify it for your needs!

\bibliographystyle{plain}
\bibliography{cite}

\end{document}\documentclass{article}
\usepackage[utf8]{inputenc}
\usepackage{booktabs}
\usepackage{amsmath}
\usepackage{natbib}

\title{My Document}
\author{Your Name}
\date{\today}

\begin{document}

\maketitle

\section{Introduction}
This is a sample document demonstrating various \LaTeX{} features including mathematical equations, tables, and formatting. You can cite papers using the citation search feature, and they will be automatically added to your \texttt{cite.bib} file.

\section{Mathematical Expressions}
Here are some example equations:

\subsection{Quadratic Formula}
The quadratic formula is given by:
\[x = \frac{-b \pm \sqrt{b^2 - 4ac}}{2a}\]

\subsection{Matrix Example}
A sample matrix:
\[
\begin{pmatrix}
a_{11} & a_{12} & a_{13} \\
a_{21} & a_{22} & a_{23} \\
a_{31} & a_{32} & a_{33}
\end{pmatrix}
\]

\section{Tables}
Below is a sample table showing data:

\begin{table}[h]
\centering
\begin{tabular}{@{}lrrr@{}}
\toprule
\textbf{Category} & \textbf{Value 1} & \textbf{Value 2} & \textbf{Total} \\
\midrule
Item A & 10.5 & 20.3 & 30.8 \\
Item B & 15.2 & 25.1 & 40.3 \\
Item C & 22.1 & 12.7 & 34.8 \\
\bottomrule
\end{tabular}
\caption{Sample Data Table}
\label{tab:sample}
\end{table}

\section{Lists}
\subsection{Bullet Points}
\begin{itemize}
  \item First item
  \item Second item
  \item Third item with sub-items:
    \begin{itemize}
      \item Sub-item 1
      \item Sub-item 2
    \end{itemize}
\end{itemize}

\subsection{Numbered List}
\begin{enumerate}
  \item First step
  \item Second step
  \item Third step
\end{enumerate}

\section{Citations and References}
When you use the citation search feature, references will be automatically managed. For example, you might cite a paper like this: \cite{example2024}. All citations will be automatically added to your \texttt{cite.bib} file.

\section{Conclusion}
This template demonstrates basic \LaTeX{} features including mathematical formulas, tables, and lists. The bibliography system is now set up to work seamlessly with Bibby AI's citation management. Feel free to modify it for your needs!

\bibliographystyle{plain}
\bibliography{cite}

\end{document}\documentclass{article}
\usepackage[utf8]{inputenc}
\usepackage{booktabs}
\usepackage{amsmath}
\usepackage{natbib}

\title{My Document}
\author{Your Name}
\date{\today}

\begin{document}

\maketitle

\section{Introduction}
This is a sample document demonstrating various \LaTeX{} features including mathematical equations, tables, and formatting. You can cite papers using the citation search feature, and they will be automatically added to your \texttt{cite.bib} file.

\section{Mathematical Expressions}
Here are some example equations:

\subsection{Quadratic Formula}
The quadratic formula is given by:
\[x = \frac{-b \pm \sqrt{b^2 - 4ac}}{2a}\]

\subsection{Matrix Example}
A sample matrix:
\[
\begin{pmatrix}
a_{11} & a_{12} & a_{13} \\
a_{21} & a_{22} & a_{23} \\
a_{31} & a_{32} & a_{33}
\end{pmatrix}
\]

\section{Tables}
Below is a sample table showing data:

\begin{table}[h]
\centering
\begin{tabular}{@{}lrrr@{}}
\toprule
\textbf{Category} & \textbf{Value 1} & \textbf{Value 2} & \textbf{Total} \\
\midrule
Item A & 10.5 & 20.3 & 30.8 \\
Item B & 15.2 & 25.1 & 40.3 \\
Item C & 22.1 & 12.7 & 34.8 \\
\bottomrule
\end{tabular}
\caption{Sample Data Table}
\label{tab:sample}
\end{table}

\section{Lists}
\subsection{Bullet Points}
\begin{itemize}
  \item First item
  \item Second item
  \item Third item with sub-items:
    \begin{itemize}
      \item Sub-item 1
      \item Sub-item 2
    \end{itemize}
\end{itemize}

\subsection{Numbered List}
\begin{enumerate}
  \item First step
  \item Second step
  \item Third step
\end{enumerate}

\section{Citations and References}
When you use the citation search feature, references will be automatically managed. For example, you might cite a paper like this: \cite{example2024}. All citations will be automatically added to your \texttt{cite.bib} file.

\section{Conclusion}
This template demonstrates basic \LaTeX{} features including mathematical formulas, tables, and lists. The bibliography system is now set up to work seamlessly with Bibby AI's citation management. Feel free to modify it for your needs!

\bibliographystyle{plain}
\bibliography{cite}

\end{document}\documentclass{article}
\usepackage[utf8]{inputenc}
\usepackage{booktabs}
\usepackage{amsmath}
\usepackage{natbib}

\title{My Document}
\author{Your Name}
\date{\today}

\begin{document}

\maketitle

\section{Introduction}
This is a sample document demonstrating various \LaTeX{} features including mathematical equations, tables, and formatting. You can cite papers using the citation search feature, and they will be automatically added to your \texttt{cite.bib} file.

\section{Mathematical Expressions}
Here are some example equations:

\subsection{Quadratic Formula}
The quadratic formula is given by:
\[x = \frac{-b \pm \sqrt{b^2 - 4ac}}{2a}\]

\subsection{Matrix Example}
A sample matrix:
\[
\begin{pmatrix}
a_{11} & a_{12} & a_{13} \\
a_{21} & a_{22} & a_{23} \\
a_{31} & a_{32} & a_{33}
\end{pmatrix}
\]

\section{Tables}
Below is a sample table showing data:

\begin{table}[h]
\centering
\begin{tabular}{@{}lrrr@{}}
\toprule
\textbf{Category} & \textbf{Value 1} & \textbf{Value 2} & \textbf{Total} \\
\midrule
Item A & 10.5 & 20.3 & 30.8 \\
Item B & 15.2 & 25.1 & 40.3 \\
Item C & 22.1 & 12.7 & 34.8 \\
\bottomrule
\end{tabular}
\caption{Sample Data Table}
\label{tab:sample}
\end{table}

\section{Lists}
\subsection{Bullet Points}
\begin{itemize}
  \item First item
  \item Second item
  \item Third item with sub-items:
    \begin{itemize}
      \item Sub-item 1
      \item Sub-item 2
    \end{itemize}
\end{itemize}

\subsection{Numbered List}
\begin{enumerate}
  \item First step
  \item Second step
  \item Third step
\end{enumerate}

\section{Citations and References}
When you use the citation search feature, references will be automatically managed. For example, you might cite a paper like this: \cite{example2024}. All citations will be automatically added to your \texttt{cite.bib} file.

\section{Conclusion}
This template demonstrates basic \LaTeX{} features including mathematical formulas, tables, and lists. The bibliography system is now set up to work seamlessly with Bibby AI's citation management. Feel free to modify it for your needs!

\bibliographystyle{plain}
\bibliography{cite}

\end{document}\documentclass{article}
\usepackage[utf8]{inputenc}
\usepackage{booktabs}
\usepackage{amsmath}
\usepackage{natbib}

\title{My Document}
\author{Your Name}
\date{\today}

\begin{document}

\maketitle

\section{Introduction}
This is a sample document demonstrating various \LaTeX{} features including mathematical equations, tables, and formatting. You can cite papers using the citation search feature, and they will be automatically added to your \texttt{cite.bib} file.

\section{Mathematical Expressions}
Here are some example equations:

\subsection{Quadratic Formula}
The quadratic formula is given by:
\[x = \frac{-b \pm \sqrt{b^2 - 4ac}}{2a}\]

\subsection{Matrix Example}
A sample matrix:
\[
\begin{pmatrix}
a_{11} & a_{12} & a_{13} \\
a_{21} & a_{22} & a_{23} \\
a_{31} & a_{32} & a_{33}
\end{pmatrix}
\]

\section{Tables}
Below is a sample table showing data:

\begin{table}[h]
\centering
\begin{tabular}{@{}lrrr@{}}
\toprule
\textbf{Category} & \textbf{Value 1} & \textbf{Value 2} & \textbf{Total} \\
\midrule
Item A & 10.5 & 20.3 & 30.8 \\
Item B & 15.2 & 25.1 & 40.3 \\
Item C & 22.1 & 12.7 & 34.8 \\
\bottomrule
\end{tabular}
\caption{Sample Data Table}
\label{tab:sample}
\end{table}

\section{Lists}
\subsection{Bullet Points}
\begin{itemize}
  \item First item
  \item Second item
  \item Third item with sub-items:
    \begin{itemize}
      \item Sub-item 1
      \item Sub-item 2
    \end{itemize}
\end{itemize}

\subsection{Numbered List}
\begin{enumerate}
  \item First step
  \item Second step
  \item Third step
\end{enumerate}

\section{Citations and References}
When you use the citation search feature, references will be automatically managed. For example, you might cite a paper like this: \cite{example2024}. All citations will be automatically added to your \texttt{cite.bib} file.

\section{Conclusion}
This template demonstrates basic \LaTeX{} features including mathematical formulas, tables, and lists. The bibliography system is now set up to work seamlessly with Bibby AI's citation management. Feel free to modify it for your needs!

\bibliographystyle{plain}
\bibliography{cite}

\end{document}\documentclass{article}
\usepackage[utf8]{inputenc}
\usepackage{booktabs}
\usepackage{amsmath}
\usepackage{natbib}

\title{My Document}
\author{Your Name}
\date{\today}

\begin{document}

\maketitle

\section{Introduction}
This is a sample document demonstrating various \LaTeX{} features including mathematical equations, tables, and formatting. You can cite papers using the citation search feature, and they will be automatically added to your \texttt{cite.bib} file.

\section{Mathematical Expressions}
Here are some example equations:

\subsection{Quadratic Formula}
The quadratic formula is given by:
\[x = \frac{-b \pm \sqrt{b^2 - 4ac}}{2a}\]

\subsection{Matrix Example}
A sample matrix:
\[
\begin{pmatrix}
a_{11} & a_{12} & a_{13} \\
a_{21} & a_{22} & a_{23} \\
a_{31} & a_{32} & a_{33}
\end{pmatrix}
\]

\section{Tables}
Below is a sample table showing data:

\begin{table}[h]
\centering
\begin{tabular}{@{}lrrr@{}}
\toprule
\textbf{Category} & \textbf{Value 1} & \textbf{Value 2} & \textbf{Total} \\
\midrule
Item A & 10.5 & 20.3 & 30.8 \\
Item B & 15.2 & 25.1 & 40.3 \\
Item C & 22.1 & 12.7 & 34.8 \\
\bottomrule
\end{tabular}
\caption{Sample Data Table}
\label{tab:sample}
\end{table}

\section{Lists}
\subsection{Bullet Points}
\begin{itemize}
  \item First item
  \item Second item
  \item Third item with sub-items:
    \begin{itemize}
      \item Sub-item 1
      \item Sub-item 2
    \end{itemize}
\end{itemize}

\subsection{Numbered List}
\begin{enumerate}
  \item First step
  \item Second step
  \item Third step
\end{enumerate}

\section{Citations and References}
When you use the citation search feature, references will be automatically managed. For example, you might cite a paper like this: \cite{example2024}. All citations will be automatically added to your \texttt{cite.bib} file.

\section{Conclusion}
This template demonstrates basic \LaTeX{} features including mathematical formulas, tables, and lists. The bibliography system is now set up to work seamlessly with Bibby AI's citation management. Feel free to modify it for your needs!

\bibliographystyle{plain}
\bibliography{cite}

\end{document}\documentclass{article}
\usepackage[utf8]{inputenc}
\usepackage{booktabs}
\usepackage{amsmath}
\usepackage{natbib}

\title{My Document}
\author{Your Name}
\date{\today}

\begin{document}

\maketitle

\section{Introduction}
This is a sample document demonstrating various \LaTeX{} features including mathematical equations, tables, and formatting. You can cite papers using the citation search feature, and they will be automatically added to your \texttt{cite.bib} file.

\section{Mathematical Expressions}
Here are some example equations:

\subsection{Quadratic Formula}
The quadratic formula is given by:
\[x = \frac{-b \pm \sqrt{b^2 - 4ac}}{2a}\]

\subsection{Matrix Example}
A sample matrix:
\[
\begin{pmatrix}
a_{11} & a_{12} & a_{13} \\
a_{21} & a_{22} & a_{23} \\
a_{31} & a_{32} & a_{33}
\end{pmatrix}
\]

\section{Tables}
Below is a sample table showing data:

\begin{table}[h]
\centering
\begin{tabular}{@{}lrrr@{}}
\toprule
\textbf{Category} & \textbf{Value 1} & \textbf{Value 2} & \textbf{Total} \\
\midrule
Item A & 10.5 & 20.3 & 30.8 \\
Item B & 15.2 & 25.1 & 40.3 \\
Item C & 22.1 & 12.7 & 34.8 \\
\bottomrule
\end{tabular}
\caption{Sample Data Table}
\label{tab:sample}
\end{table}

\section{Lists}
\subsection{Bullet Points}
\begin{itemize}
  \item First item
  \item Second item
  \item Third item with sub-items:
    \begin{itemize}
      \item Sub-item 1
      \item Sub-item 2
    \end{itemize}
\end{itemize}

\subsection{Numbered List}
\begin{enumerate}
  \item First step
  \item Second step
  \item Third step
\end{enumerate}

\section{Citations and References}
When you use the citation search feature, references will be automatically managed. For example, you might cite a paper like this: \cite{example2024}. All citations will be automatically added to your \texttt{cite.bib} file.

\section{Conclusion}
This template demonstrates basic \LaTeX{} features including mathematical formulas, tables, and lists. The bibliography system is now set up to work seamlessly with Bibby AI's citation management. Feel free to modify it for your needs!

\bibliographystyle{plain}
\bibliography{cite}

\end{document}\documentclass{article}
\usepackage[utf8]{inputenc}
\usepackage{booktabs}
\usepackage{amsmath}
\usepackage{natbib}

\title{My Document}
\author{Your Name}
\date{\today}

\begin{document}

\maketitle

\section{Introduction}
This is a sample document demonstrating various \LaTeX{} features including mathematical equations, tables, and formatting. You can cite papers using the citation search feature, and they will be automatically added to your \texttt{cite.bib} file.

\section{Mathematical Expressions}
Here are some example equations:

\subsection{Quadratic Formula}
The quadratic formula is given by:
\[x = \frac{-b \pm \sqrt{b^2 - 4ac}}{2a}\]

\subsection{Matrix Example}
A sample matrix:
\[
\begin{pmatrix}
a_{11} & a_{12} & a_{13} \\
a_{21} & a_{22} & a_{23} \\
a_{31} & a_{32} & a_{33}
\end{pmatrix}
\]

\section{Tables}
Below is a sample table showing data:

\begin{table}[h]
\centering
\begin{tabular}{@{}lrrr@{}}
\toprule
\textbf{Category} & \textbf{Value 1} & \textbf{Value 2} & \textbf{Total} \\
\midrule
Item A & 10.5 & 20.3 & 30.8 \\
Item B & 15.2 & 25.1 & 40.3 \\
Item C & 22.1 & 12.7 & 34.8 \\
\bottomrule
\end{tabular}
\caption{Sample Data Table}
\label{tab:sample}
\end{table}

\section{Lists}
\subsection{Bullet Points}
\begin{itemize}
  \item First item
  \item Second item
  \item Third item with sub-items:
    \begin{itemize}
      \item Sub-item 1
      \item Sub-item 2
    \end{itemize}
\end{itemize}

\subsection{Numbered List}
\begin{enumerate}
  \item First step
  \item Second step
  \item Third step
\end{enumerate}

\section{Citations and References}
When you use the citation search feature, references will be automatically managed. For example, you might cite a paper like this: \cite{example2024}. All citations will be automatically added to your \texttt{cite.bib} file.

\section{Conclusion}
This template demonstrates basic \LaTeX{} features including mathematical formulas, tables, and lists. The bibliography system is now set up to work seamlessly with Bibby AI's citation management. Feel free to modify it for your needs!

\bibliographystyle{plain}
\bibliography{cite}

\end{document}\documentclass{article}
\usepackage[utf8]{inputenc}
\usepackage{booktabs}
\usepackage{amsmath}
\usepackage{natbib}

\title{My Document}
\author{Your Name}
\date{\today}

\begin{document}

\maketitle

\section{Introduction}
This is a sample document demonstrating various \LaTeX{} features including mathematical equations, tables, and formatting. You can cite papers using the citation search feature, and they will be automatically added to your \texttt{cite.bib} file.

\section{Mathematical Expressions}
Here are some example equations:

\subsection{Quadratic Formula}
The quadratic formula is given by:
\[x = \frac{-b \pm \sqrt{b^2 - 4ac}}{2a}\]

\subsection{Matrix Example}
A sample matrix:
\[
\begin{pmatrix}
a_{11} & a_{12} & a_{13} \\
a_{21} & a_{22} & a_{23} \\
a_{31} & a_{32} & a_{33}
\end{pmatrix}
\]

\section{Tables}
Below is a sample table showing data:

\begin{table}[h]
\centering
\begin{tabular}{@{}lrrr@{}}
\toprule
\textbf{Category} & \textbf{Value 1} & \textbf{Value 2} & \textbf{Total} \\
\midrule
Item A & 10.5 & 20.3 & 30.8 \\
Item B & 15.2 & 25.1 & 40.3 \\
Item C & 22.1 & 12.7 & 34.8 \\
\bottomrule
\end{tabular}
\caption{Sample Data Table}
\label{tab:sample}
\end{table}

\section{Lists}
\subsection{Bullet Points}
\begin{itemize}
  \item First item
  \item Second item
  \item Third item with sub-items:
    \begin{itemize}
      \item Sub-item 1
      \item Sub-item 2
    \end{itemize}
\end{itemize}

\subsection{Numbered List}
\begin{enumerate}
  \item First step
  \item Second step
  \item Third step
\end{enumerate}

\section{Citations and References}
When you use the citation search feature, references will be automatically managed. For example, you might cite a paper like this: \cite{example2024}. All citations will be automatically added to your \texttt{cite.bib} file.

\section{Conclusion}
This template demonstrates basic \LaTeX{} features including mathematical formulas, tables, and lists. The bibliography system is now set up to work seamlessly with Bibby AI's citation management. Feel free to modify it for your needs!

\bibliographystyle{plain}
\bibliography{cite}

\end{document}\documentclass{article}
\usepackage[utf8]{inputenc}
\usepackage{booktabs}
\usepackage{amsmath}
\usepackage{natbib}

\title{My Document}
\author{Your Name}
\date{\today}

\begin{document}

\maketitle

\section{Introduction}
This is a sample document demonstrating various \LaTeX{} features including mathematical equations, tables, and formatting. You can cite papers using the citation search feature, and they will be automatically added to your \texttt{cite.bib} file.

\section{Mathematical Expressions}
Here are some example equations:

\subsection{Quadratic Formula}
The quadratic formula is given by:
\[x = \frac{-b \pm \sqrt{b^2 - 4ac}}{2a}\]

\subsection{Matrix Example}
A sample matrix:
\[
\begin{pmatrix}
a_{11} & a_{12} & a_{13} \\
a_{21} & a_{22} & a_{23} \\
a_{31} & a_{32} & a_{33}
\end{pmatrix}
\]

\section{Tables}
Below is a sample table showing data:

\begin{table}[h]
\centering
\begin{tabular}{@{}lrrr@{}}
\toprule
\textbf{Category} & \textbf{Value 1} & \textbf{Value 2} & \textbf{Total} \\
\midrule
Item A & 10.5 & 20.3 & 30.8 \\
Item B & 15.2 & 25.1 & 40.3 \\
Item C & 22.1 & 12.7 & 34.8 \\
\bottomrule
\end{tabular}
\caption{Sample Data Table}
\label{tab:sample}
\end{table}

\section{Lists}
\subsection{Bullet Points}
\begin{itemize}
  \item First item
  \item Second item
  \item Third item with sub-items:
    \begin{itemize}
      \item Sub-item 1
      \item Sub-item 2
    \end{itemize}
\end{itemize}

\subsection{Numbered List}
\begin{enumerate}
  \item First step
  \item Second step
  \item Third step
\end{enumerate}

\section{Citations and References}
When you use the citation search feature, references will be automatically managed. For example, you might cite a paper like this: \cite{example2024}. All citations will be automatically added to your \texttt{cite.bib} file.

\section{Conclusion}
This template demonstrates basic \LaTeX{} features including mathematical formulas, tables, and lists. The bibliography system is now set up to work seamlessly with Bibby AI's citation management. Feel free to modify it for your needs!

\bibliographystyle{plain}
\bibliography{cite}

\end{document}\documentclass{article}
\usepackage[utf8]{inputenc}
\usepackage{booktabs}
\usepackage{amsmath}
\usepackage{natbib}

\title{My Document}
\author{Your Name}
\date{\today}

\begin{document}

\maketitle

\section{Introduction}
This is a sample document demonstrating various \LaTeX{} features including mathematical equations, tables, and formatting. You can cite papers using the citation search feature, and they will be automatically added to your \texttt{cite.bib} file.

\section{Mathematical Expressions}
Here are some example equations:

\subsection{Quadratic Formula}
The quadratic formula is given by:
\[x = \frac{-b \pm \sqrt{b^2 - 4ac}}{2a}\]

\subsection{Matrix Example}
A sample matrix:
\[
\begin{pmatrix}
a_{11} & a_{12} & a_{13} \\
a_{21} & a_{22} & a_{23} \\
a_{31} & a_{32} & a_{33}
\end{pmatrix}
\]

\section{Tables}
Below is a sample table showing data:

\begin{table}[h]
\centering
\begin{tabular}{@{}lrrr@{}}
\toprule
\textbf{Category} & \textbf{Value 1} & \textbf{Value 2} & \textbf{Total} \\
\midrule
Item A & 10.5 & 20.3 & 30.8 \\
Item B & 15.2 & 25.1 & 40.3 \\
Item C & 22.1 & 12.7 & 34.8 \\
\bottomrule
\end{tabular}
\caption{Sample Data Table}
\label{tab:sample}
\end{table}

\section{Lists}
\subsection{Bullet Points}
\begin{itemize}
  \item First item
  \item Second item
  \item Third item with sub-items:
    \begin{itemize}
      \item Sub-item 1
      \item Sub-item 2
    \end{itemize}
\end{itemize}

\subsection{Numbered List}
\begin{enumerate}
  \item First step
  \item Second step
  \item Third step
\end{enumerate}

\section{Citations and References}
When you use the citation search feature, references will be automatically managed. For example, you might cite a paper like this: \cite{example2024}. All citations will be automatically added to your \texttt{cite.bib} file.

\section{Conclusion}
This template demonstrates basic \LaTeX{} features including mathematical formulas, tables, and lists. The bibliography system is now set up to work seamlessly with Bibby AI's citation management. Feel free to modify it for your needs!

\bibliographystyle{plain}
\bibliography{cite}

\end{document}\documentclass{article}
\usepackage[utf8]{inputenc}
\usepackage{booktabs}
\usepackage{amsmath}
\usepackage{natbib}

\title{My Document}
\author{Your Name}
\date{\today}

\begin{document}

\maketitle

\section{Introduction}
This is a sample document demonstrating various \LaTeX{} features including mathematical equations, tables, and formatting. You can cite papers using the citation search feature, and they will be automatically added to your \texttt{cite.bib} file.

\section{Mathematical Expressions}
Here are some example equations:

\subsection{Quadratic Formula}
The quadratic formula is given by:
\[x = \frac{-b \pm \sqrt{b^2 - 4ac}}{2a}\]

\subsection{Matrix Example}
A sample matrix:
\[
\begin{pmatrix}
a_{11} & a_{12} & a_{13} \\
a_{21} & a_{22} & a_{23} \\
a_{31} & a_{32} & a_{33}
\end{pmatrix}
\]

\section{Tables}
Below is a sample table showing data:

\begin{table}[h]
\centering
\begin{tabular}{@{}lrrr@{}}
\toprule
\textbf{Category} & \textbf{Value 1} & \textbf{Value 2} & \textbf{Total} \\
\midrule
Item A & 10.5 & 20.3 & 30.8 \\
Item B & 15.2 & 25.1 & 40.3 \\
Item C & 22.1 & 12.7 & 34.8 \\
\bottomrule
\end{tabular}
\caption{Sample Data Table}
\label{tab:sample}
\end{table}

\section{Lists}
\subsection{Bullet Points}
\begin{itemize}
  \item First item
  \item Second item
  \item Third item with sub-items:
    \begin{itemize}
      \item Sub-item 1
      \item Sub-item 2
    \end{itemize}
\end{itemize}

\subsection{Numbered List}
\begin{enumerate}
  \item First step
  \item Second step
  \item Third step
\end{enumerate}

\section{Citations and References}
When you use the citation search feature, references will be automatically managed. For example, you might cite a paper like this: \cite{example2024}. All citations will be automatically added to your \texttt{cite.bib} file.

\section{Conclusion}
This template demonstrates basic \LaTeX{} features including mathematical formulas, tables, and lists. The bibliography system is now set up to work seamlessly with Bibby AI's citation management. Feel free to modify it for your needs!

\bibliographystyle{plain}
\bibliography{cite}

\end{document}\documentclass{article}
\usepackage[utf8]{inputenc}
\usepackage{booktabs}
\usepackage{amsmath}
\usepackage{natbib}

\title{My Document}
\author{Your Name}
\date{\today}

\begin{document}

\maketitle

\section{Introduction}
This is a sample document demonstrating various \LaTeX{} features including mathematical equations, tables, and formatting. You can cite papers using the citation search feature, and they will be automatically added to your \texttt{cite.bib} file.

\section{Mathematical Expressions}
Here are some example equations:

\subsection{Quadratic Formula}
The quadratic formula is given by:
\[x = \frac{-b \pm \sqrt{b^2 - 4ac}}{2a}\]

\subsection{Matrix Example}
A sample matrix:
\[
\begin{pmatrix}
a_{11} & a_{12} & a_{13} \\
a_{21} & a_{22} & a_{23} \\
a_{31} & a_{32} & a_{33}
\end{pmatrix}
\]

\section{Tables}
Below is a sample table showing data:

\begin{table}[h]
\centering
\begin{tabular}{@{}lrrr@{}}
\toprule
\textbf{Category} & \textbf{Value 1} & \textbf{Value 2} & \textbf{Total} \\
\midrule
Item A & 10.5 & 20.3 & 30.8 \\
Item B & 15.2 & 25.1 & 40.3 \\
Item C & 22.1 & 12.7 & 34.8 \\
\bottomrule
\end{tabular}
\caption{Sample Data Table}
\label{tab:sample}
\end{table}

\section{Lists}
\subsection{Bullet Points}
\begin{itemize}
  \item First item
  \item Second item
  \item Third item with sub-items:
    \begin{itemize}
      \item Sub-item 1
      \item Sub-item 2
    \end{itemize}
\end{itemize}

\subsection{Numbered List}
\begin{enumerate}
  \item First step
  \item Second step
  \item Third step
\end{enumerate}

\section{Citations and References}
When you use the citation search feature, references will be automatically managed. For example, you might cite a paper like this: \cite{example2024}. All citations will be automatically added to your \texttt{cite.bib} file.

\section{Conclusion}
This template demonstrates basic \LaTeX{} features including mathematical formulas, tables, and lists. The bibliography system is now set up to work seamlessly with Bibby AI's citation management. Feel free to modify it for your needs!

\bibliographystyle{plain}
\bibliography{cite}

\end{document}\documentclass{article}
\usepackage[utf8]{inputenc}
\usepackage{booktabs}
\usepackage{amsmath}
\usepackage{natbib}

\title{My Document}
\author{Your Name}
\date{\today}

\begin{document}

\maketitle

\section{Introduction}
This is a sample document demonstrating various \LaTeX{} features including mathematical equations, tables, and formatting. You can cite papers using the citation search feature, and they will be automatically added to your \texttt{cite.bib} file.

\section{Mathematical Expressions}
Here are some example equations:

\subsection{Quadratic Formula}
The quadratic formula is given by:
\[x = \frac{-b \pm \sqrt{b^2 - 4ac}}{2a}\]

\subsection{Matrix Example}
A sample matrix:
\[
\begin{pmatrix}
a_{11} & a_{12} & a_{13} \\
a_{21} & a_{22} & a_{23} \\
a_{31} & a_{32} & a_{33}
\end{pmatrix}
\]

\section{Tables}
Below is a sample table showing data:

\begin{table}[h]
\centering
\begin{tabular}{@{}lrrr@{}}
\toprule
\textbf{Category} & \textbf{Value 1} & \textbf{Value 2} & \textbf{Total} \\
\midrule
Item A & 10.5 & 20.3 & 30.8 \\
Item B & 15.2 & 25.1 & 40.3 \\
Item C & 22.1 & 12.7 & 34.8 \\
\bottomrule
\end{tabular}
\caption{Sample Data Table}
\label{tab:sample}
\end{table}

\section{Lists}
\subsection{Bullet Points}
\begin{itemize}
  \item First item
  \item Second item
  \item Third item with sub-items:
    \begin{itemize}
      \item Sub-item 1
      \item Sub-item 2
    \end{itemize}
\end{itemize}

\subsection{Numbered List}
\begin{enumerate}
  \item First step
  \item Second step
  \item Third step
\end{enumerate}

\section{Citations and References}
When you use the citation search feature, references will be automatically managed. For example, you might cite a paper like this: \cite{example2024}. All citations will be automatically added to your \texttt{cite.bib} file.

\section{Conclusion}
This template demonstrates basic \LaTeX{} features including mathematical formulas, tables, and lists. The bibliography system is now set up to work seamlessly with Bibby AI's citation management. Feel free to modify it for your needs!

\bibliographystyle{plain}
\bibliography{cite}

\end{document}\documentclass{article}
\usepackage[utf8]{inputenc}
\usepackage{booktabs}
\usepackage{amsmath}
\usepackage{natbib}

\title{My Document}
\author{Your Name}
\date{\today}

\begin{document}

\maketitle

\section{Introduction}
This is a sample document demonstrating various \LaTeX{} features including mathematical equations, tables, and formatting. You can cite papers using the citation search feature, and they will be automatically added to your \texttt{cite.bib} file.

\section{Mathematical Expressions}
Here are some example equations:

\subsection{Quadratic Formula}
The quadratic formula is given by:
\[x = \frac{-b \pm \sqrt{b^2 - 4ac}}{2a}\]

\subsection{Matrix Example}
A sample matrix:
\[
\begin{pmatrix}
a_{11} & a_{12} & a_{13} \\
a_{21} & a_{22} & a_{23} \\
a_{31} & a_{32} & a_{33}
\end{pmatrix}
\]

\section{Tables}
Below is a sample table showing data:

\begin{table}[h]
\centering
\begin{tabular}{@{}lrrr@{}}
\toprule
\textbf{Category} & \textbf{Value 1} & \textbf{Value 2} & \textbf{Total} \\
\midrule
Item A & 10.5 & 20.3 & 30.8 \\
Item B & 15.2 & 25.1 & 40.3 \\
Item C & 22.1 & 12.7 & 34.8 \\
\bottomrule
\end{tabular}
\caption{Sample Data Table}
\label{tab:sample}
\end{table}

\section{Lists}
\subsection{Bullet Points}
\begin{itemize}
  \item First item
  \item Second item
  \item Third item with sub-items:
    \begin{itemize}
      \item Sub-item 1
      \item Sub-item 2
    \end{itemize}
\end{itemize}

\subsection{Numbered List}
\begin{enumerate}
  \item First step
  \item Second step
  \item Third step
\end{enumerate}

\section{Citations and References}
When you use the citation search feature, references will be automatically managed. For example, you might cite a paper like this: \cite{example2024}. All citations will be automatically added to your \texttt{cite.bib} file.

\section{Conclusion}
This template demonstrates basic \LaTeX{} features including mathematical formulas, tables, and lists. The bibliography system is now set up to work seamlessly with Bibby AI's citation management. Feel free to modify it for your needs!

\bibliographystyle{plain}
\bibliography{cite}

\end{document}\documentclass{article}
\usepackage[utf8]{inputenc}
\usepackage{booktabs}
\usepackage{amsmath}
\usepackage{natbib}

\title{My Document}
\author{Your Name}
\date{\today}

\begin{document}

\maketitle

\section{Introduction}
This is a sample document demonstrating various \LaTeX{} features including mathematical equations, tables, and formatting. You can cite papers using the citation search feature, and they will be automatically added to your \texttt{cite.bib} file.

\section{Mathematical Expressions}
Here are some example equations:

\subsection{Quadratic Formula}
The quadratic formula is given by:
\[x = \frac{-b \pm \sqrt{b^2 - 4ac}}{2a}\]

\subsection{Matrix Example}
A sample matrix:
\[
\begin{pmatrix}
a_{11} & a_{12} & a_{13} \\
a_{21} & a_{22} & a_{23} \\
a_{31} & a_{32} & a_{33}
\end{pmatrix}
\]

\section{Tables}
Below is a sample table showing data:

\begin{table}[h]
\centering
\begin{tabular}{@{}lrrr@{}}
\toprule
\textbf{Category} & \textbf{Value 1} & \textbf{Value 2} & \textbf{Total} \\
\midrule
Item A & 10.5 & 20.3 & 30.8 \\
Item B & 15.2 & 25.1 & 40.3 \\
Item C & 22.1 & 12.7 & 34.8 \\
\bottomrule
\end{tabular}
\caption{Sample Data Table}
\label{tab:sample}
\end{table}

\section{Lists}
\subsection{Bullet Points}
\begin{itemize}
  \item First item
  \item Second item
  \item Third item with sub-items:
    \begin{itemize}
      \item Sub-item 1
      \item Sub-item 2
    \end{itemize}
\end{itemize}

\subsection{Numbered List}
\begin{enumerate}
  \item First step
  \item Second step
  \item Third step
\end{enumerate}

\section{Citations and References}
When you use the citation search feature, references will be automatically managed. For example, you might cite a paper like this: \cite{example2024}. All citations will be automatically added to your \texttt{cite.bib} file.

\section{Conclusion}
This template demonstrates basic \LaTeX{} features including mathematical formulas, tables, and lists. The bibliography system is now set up to work seamlessly with Bibby AI's citation management. Feel free to modify it for your needs!

\bibliographystyle{plain}
\bibliography{cite}

\end{document}\documentclass{article}
\usepackage[utf8]{inputenc}
\usepackage{booktabs}
\usepackage{amsmath}
\usepackage{natbib}

\title{My Document}
\author{Your Name}
\date{\today}

\begin{document}

\maketitle

\section{Introduction}
This is a sample document demonstrating various \LaTeX{} features including mathematical equations, tables, and formatting. You can cite papers using the citation search feature, and they will be automatically added to your \texttt{cite.bib} file.

\section{Mathematical Expressions}
Here are some example equations:

\subsection{Quadratic Formula}
The quadratic formula is given by:
\[x = \frac{-b \pm \sqrt{b^2 - 4ac}}{2a}\]

\subsection{Matrix Example}
A sample matrix:
\[
\begin{pmatrix}
a_{11} & a_{12} & a_{13} \\
a_{21} & a_{22} & a_{23} \\
a_{31} & a_{32} & a_{33}
\end{pmatrix}
\]

\section{Tables}
Below is a sample table showing data:

\begin{table}[h]
\centering
\begin{tabular}{@{}lrrr@{}}
\toprule
\textbf{Category} & \textbf{Value 1} & \textbf{Value 2} & \textbf{Total} \\
\midrule
Item A & 10.5 & 20.3 & 30.8 \\
Item B & 15.2 & 25.1 & 40.3 \\
Item C & 22.1 & 12.7 & 34.8 \\
\bottomrule
\end{tabular}
\caption{Sample Data Table}
\label{tab:sample}
\end{table}

\section{Lists}
\subsection{Bullet Points}
\begin{itemize}
  \item First item
  \item Second item
  \item Third item with sub-items:
    \begin{itemize}
      \item Sub-item 1
      \item Sub-item 2
    \end{itemize}
\end{itemize}

\subsection{Numbered List}
\begin{enumerate}
  \item First step
  \item Second step
  \item Third step
\end{enumerate}

\section{Citations and References}
When you use the citation search feature, references will be automatically managed. For example, you might cite a paper like this: \cite{example2024}. All citations will be automatically added to your \texttt{cite.bib} file.

\section{Conclusion}
This template demonstrates basic \LaTeX{} features including mathematical formulas, tables, and lists. The bibliography system is now set up to work seamlessly with Bibby AI's citation management. Feel free to modify it for your needs!

\bibliographystyle{plain}
\bibliography{cite}

\end{document}\documentclass{article}
\usepackage[utf8]{inputenc}
\usepackage{booktabs}
\usepackage{amsmath}
\usepackage{natbib}

\title{My Document}
\author{Your Name}
\date{\today}

\begin{document}

\maketitle

\section{Introduction}
This is a sample document demonstrating various \LaTeX{} features including mathematical equations, tables, and formatting. You can cite papers using the citation search feature, and they will be automatically added to your \texttt{cite.bib} file.

\section{Mathematical Expressions}
Here are some example equations:

\subsection{Quadratic Formula}
The quadratic formula is given by:
\[x = \frac{-b \pm \sqrt{b^2 - 4ac}}{2a}\]

\subsection{Matrix Example}
A sample matrix:
\[
\begin{pmatrix}
a_{11} & a_{12} & a_{13} \\
a_{21} & a_{22} & a_{23} \\
a_{31} & a_{32} & a_{33}
\end{pmatrix}
\]

\section{Tables}
Below is a sample table showing data:

\begin{table}[h]
\centering
\begin{tabular}{@{}lrrr@{}}
\toprule
\textbf{Category} & \textbf{Value 1} & \textbf{Value 2} & \textbf{Total} \\
\midrule
Item A & 10.5 & 20.3 & 30.8 \\
Item B & 15.2 & 25.1 & 40.3 \\
Item C & 22.1 & 12.7 & 34.8 \\
\bottomrule
\end{tabular}
\caption{Sample Data Table}
\label{tab:sample}
\end{table}

\section{Lists}
\subsection{Bullet Points}
\begin{itemize}
  \item First item
  \item Second item
  \item Third item with sub-items:
    \begin{itemize}
      \item Sub-item 1
      \item Sub-item 2
    \end{itemize}
\end{itemize}

\subsection{Numbered List}
\begin{enumerate}
  \item First step
  \item Second step
  \item Third step
\end{enumerate}

\section{Citations and References}
When you use the citation search feature, references will be automatically managed. For example, you might cite a paper like this: \cite{example2024}. All citations will be automatically added to your \texttt{cite.bib} file.

\section{Conclusion}
This template demonstrates basic \LaTeX{} features including mathematical formulas, tables, and lists. The bibliography system is now set up to work seamlessly with Bibby AI's citation management. Feel free to modify it for your needs!

\bibliographystyle{plain}
\bibliography{cite}

\end{document}\documentclass{article}
\usepackage[utf8]{inputenc}
\usepackage{booktabs}
\usepackage{amsmath}
\usepackage{natbib}

\title{My Document}
\author{Your Name}
\date{\today}

\begin{document}

\maketitle

\section{Introduction}
This is a sample document demonstrating various \LaTeX{} features including mathematical equations, tables, and formatting. You can cite papers using the citation search feature, and they will be automatically added to your \texttt{cite.bib} file.

\section{Mathematical Expressions}
Here are some example equations:

\subsection{Quadratic Formula}
The quadratic formula is given by:
\[x = \frac{-b \pm \sqrt{b^2 - 4ac}}{2a}\]

\subsection{Matrix Example}
A sample matrix:
\[
\begin{pmatrix}
a_{11} & a_{12} & a_{13} \\
a_{21} & a_{22} & a_{23} \\
a_{31} & a_{32} & a_{33}
\end{pmatrix}
\]

\section{Tables}
Below is a sample table showing data:

\begin{table}[h]
\centering
\begin{tabular}{@{}lrrr@{}}
\toprule
\textbf{Category} & \textbf{Value 1} & \textbf{Value 2} & \textbf{Total} \\
\midrule
Item A & 10.5 & 20.3 & 30.8 \\
Item B & 15.2 & 25.1 & 40.3 \\
Item C & 22.1 & 12.7 & 34.8 \\
\bottomrule
\end{tabular}
\caption{Sample Data Table}
\label{tab:sample}
\end{table}

\section{Lists}
\subsection{Bullet Points}
\begin{itemize}
  \item First item
  \item Second item
  \item Third item with sub-items:
    \begin{itemize}
      \item Sub-item 1
      \item Sub-item 2
    \end{itemize}
\end{itemize}

\subsection{Numbered List}
\begin{enumerate}
  \item First step
  \item Second step
  \item Third step
\end{enumerate}

\section{Citations and References}
When you use the citation search feature, references will be automatically managed. For example, you might cite a paper like this: \cite{example2024}. All citations will be automatically added to your \texttt{cite.bib} file.

\section{Conclusion}
This template demonstrates basic \LaTeX{} features including mathematical formulas, tables, and lists. The bibliography system is now set up to work seamlessly with Bibby AI's citation management. Feel free to modify it for your needs!

\bibliographystyle{plain}
\bibliography{cite}

\end{document}\documentclass{article}
\usepackage[utf8]{inputenc}
\usepackage{booktabs}
\usepackage{amsmath}
\usepackage{natbib}

\title{My Document}
\author{Your Name}
\date{\today}

\begin{document}

\maketitle

\section{Introduction}
This is a sample document demonstrating various \LaTeX{} features including mathematical equations, tables, and formatting. You can cite papers using the citation search feature, and they will be automatically added to your \texttt{cite.bib} file.

\section{Mathematical Expressions}
Here are some example equations:

\subsection{Quadratic Formula}
The quadratic formula is given by:
\[x = \frac{-b \pm \sqrt{b^2 - 4ac}}{2a}\]

\subsection{Matrix Example}
A sample matrix:
\[
\begin{pmatrix}
a_{11} & a_{12} & a_{13} \\
a_{21} & a_{22} & a_{23} \\
a_{31} & a_{32} & a_{33}
\end{pmatrix}
\]

\section{Tables}
Below is a sample table showing data:

\begin{table}[h]
\centering
\begin{tabular}{@{}lrrr@{}}
\toprule
\textbf{Category} & \textbf{Value 1} & \textbf{Value 2} & \textbf{Total} \\
\midrule
Item A & 10.5 & 20.3 & 30.8 \\
Item B & 15.2 & 25.1 & 40.3 \\
Item C & 22.1 & 12.7 & 34.8 \\
\bottomrule
\end{tabular}
\caption{Sample Data Table}
\label{tab:sample}
\end{table}

\section{Lists}
\subsection{Bullet Points}
\begin{itemize}
  \item First item
  \item Second item
  \item Third item with sub-items:
    \begin{itemize}
      \item Sub-item 1
      \item Sub-item 2
    \end{itemize}
\end{itemize}

\subsection{Numbered List}
\begin{enumerate}
  \item First step
  \item Second step
  \item Third step
\end{enumerate}

\section{Citations and References}
When you use the citation search feature, references will be automatically managed. For example, you might cite a paper like this: \cite{example2024}. All citations will be automatically added to your \texttt{cite.bib} file.

\section{Conclusion}
This template demonstrates basic \LaTeX{} features including mathematical formulas, tables, and lists. The bibliography system is now set up to work seamlessly with Bibby AI's citation management. Feel free to modify it for your needs!

\bibliographystyle{plain}
\bibliography{cite}

\end{document}\documentclass{article}
\usepackage[utf8]{inputenc}
\usepackage{booktabs}
\usepackage{amsmath}
\usepackage{natbib}

\title{My Document}
\author{Your Name}
\date{\today}

\begin{document}

\maketitle

\section{Introduction}
This is a sample document demonstrating various \LaTeX{} features including mathematical equations, tables, and formatting. You can cite papers using the citation search feature, and they will be automatically added to your \texttt{cite.bib} file.

\section{Mathematical Expressions}
Here are some example equations:

\subsection{Quadratic Formula}
The quadratic formula is given by:
\[x = \frac{-b \pm \sqrt{b^2 - 4ac}}{2a}\]

\subsection{Matrix Example}
A sample matrix:
\[
\begin{pmatrix}
a_{11} & a_{12} & a_{13} \\
a_{21} & a_{22} & a_{23} \\
a_{31} & a_{32} & a_{33}
\end{pmatrix}
\]

\section{Tables}
Below is a sample table showing data:

\begin{table}[h]
\centering
\begin{tabular}{@{}lrrr@{}}
\toprule
\textbf{Category} & \textbf{Value 1} & \textbf{Value 2} & \textbf{Total} \\
\midrule
Item A & 10.5 & 20.3 & 30.8 \\
Item B & 15.2 & 25.1 & 40.3 \\
Item C & 22.1 & 12.7 & 34.8 \\
\bottomrule
\end{tabular}
\caption{Sample Data Table}
\label{tab:sample}
\end{table}

\section{Lists}
\subsection{Bullet Points}
\begin{itemize}
  \item First item
  \item Second item
  \item Third item with sub-items:
    \begin{itemize}
      \item Sub-item 1
      \item Sub-item 2
    \end{itemize}
\end{itemize}

\subsection{Numbered List}
\begin{enumerate}
  \item First step
  \item Second step
  \item Third step
\end{enumerate}

\section{Citations and References}
When you use the citation search feature, references will be automatically managed. For example, you might cite a paper like this: \cite{example2024}. All citations will be automatically added to your \texttt{cite.bib} file.

\section{Conclusion}
This template demonstrates basic \LaTeX{} features including mathematical formulas, tables, and lists. The bibliography system is now set up to work seamlessly with Bibby AI's citation management. Feel free to modify it for your needs!

\bibliographystyle{plain}
\bibliography{cite}

\end{document}\documentclass{article}
\usepackage[utf8]{inputenc}
\usepackage{booktabs}
\usepackage{amsmath}
\usepackage{natbib}

\title{My Document}
\author{Your Name}
\date{\today}

\begin{document}

\maketitle

\section{Introduction}
This is a sample document demonstrating various \LaTeX{} features including mathematical equations, tables, and formatting. You can cite papers using the citation search feature, and they will be automatically added to your \texttt{cite.bib} file.

\section{Mathematical Expressions}
Here are some example equations:

\subsection{Quadratic Formula}
The quadratic formula is given by:
\[x = \frac{-b \pm \sqrt{b^2 - 4ac}}{2a}\]

\subsection{Matrix Example}
A sample matrix:
\[
\begin{pmatrix}
a_{11} & a_{12} & a_{13} \\
a_{21} & a_{22} & a_{23} \\
a_{31} & a_{32} & a_{33}
\end{pmatrix}
\]

\section{Tables}
Below is a sample table showing data:

\begin{table}[h]
\centering
\begin{tabular}{@{}lrrr@{}}
\toprule
\textbf{Category} & \textbf{Value 1} & \textbf{Value 2} & \textbf{Total} \\
\midrule
Item A & 10.5 & 20.3 & 30.8 \\
Item B & 15.2 & 25.1 & 40.3 \\
Item C & 22.1 & 12.7 & 34.8 \\
\bottomrule
\end{tabular}
\caption{Sample Data Table}
\label{tab:sample}
\end{table}

\section{Lists}
\subsection{Bullet Points}
\begin{itemize}
  \item First item
  \item Second item
  \item Third item with sub-items:
    \begin{itemize}
      \item Sub-item 1
      \item Sub-item 2
    \end{itemize}
\end{itemize}

\subsection{Numbered List}
\begin{enumerate}
  \item First step
  \item Second step
  \item Third step
\end{enumerate}

\section{Citations and References}
When you use the citation search feature, references will be automatically managed. For example, you might cite a paper like this: \cite{example2024}. All citations will be automatically added to your \texttt{cite.bib} file.

\section{Conclusion}
This template demonstrates basic \LaTeX{} features including mathematical formulas, tables, and lists. The bibliography system is now set up to work seamlessly with Bibby AI's citation management. Feel free to modify it for your needs!

\bibliographystyle{plain}
\bibliography{cite}

\end{document}\documentclass{article}
\usepackage[utf8]{inputenc}
\usepackage{booktabs}
\usepackage{amsmath}
\usepackage{natbib}

\title{My Document}
\author{Your Name}
\date{\today}

\begin{document}

\maketitle

\section{Introduction}
This is a sample document demonstrating various \LaTeX{} features including mathematical equations, tables, and formatting. You can cite papers using the citation search feature, and they will be automatically added to your \texttt{cite.bib} file.

\section{Mathematical Expressions}
Here are some example equations:

\subsection{Quadratic Formula}
The quadratic formula is given by:
\[x = \frac{-b \pm \sqrt{b^2 - 4ac}}{2a}\]

\subsection{Matrix Example}
A sample matrix:
\[
\begin{pmatrix}
a_{11} & a_{12} & a_{13} \\
a_{21} & a_{22} & a_{23} \\
a_{31} & a_{32} & a_{33}
\end{pmatrix}
\]

\section{Tables}
Below is a sample table showing data:

\begin{table}[h]
\centering
\begin{tabular}{@{}lrrr@{}}
\toprule
\textbf{Category} & \textbf{Value 1} & \textbf{Value 2} & \textbf{Total} \\
\midrule
Item A & 10.5 & 20.3 & 30.8 \\
Item B & 15.2 & 25.1 & 40.3 \\
Item C & 22.1 & 12.7 & 34.8 \\
\bottomrule
\end{tabular}
\caption{Sample Data Table}
\label{tab:sample}
\end{table}

\section{Lists}
\subsection{Bullet Points}
\begin{itemize}
  \item First item
  \item Second item
  \item Third item with sub-items:
    \begin{itemize}
      \item Sub-item 1
      \item Sub-item 2
    \end{itemize}
\end{itemize}

\subsection{Numbered List}
\begin{enumerate}
  \item First step
  \item Second step
  \item Third step
\end{enumerate}

\section{Citations and References}
When you use the citation search feature, references will be automatically managed. For example, you might cite a paper like this: \cite{example2024}. All citations will be automatically added to your \texttt{cite.bib} file.

\section{Conclusion}
This template demonstrates basic \LaTeX{} features including mathematical formulas, tables, and lists. The bibliography system is now set up to work seamlessly with Bibby AI's citation management. Feel free to modify it for your needs!

\bibliographystyle{plain}
\bibliography{cite}

\end{document}\documentclass{article}
\usepackage[utf8]{inputenc}
\usepackage{booktabs}
\usepackage{amsmath}
\usepackage{natbib}

\title{My Document}
\author{Your Name}
\date{\today}

\begin{document}

\maketitle

\section{Introduction}
This is a sample document demonstrating various \LaTeX{} features including mathematical equations, tables, and formatting. You can cite papers using the citation search feature, and they will be automatically added to your \texttt{cite.bib} file.

\section{Mathematical Expressions}
Here are some example equations:

\subsection{Quadratic Formula}
The quadratic formula is given by:
\[x = \frac{-b \pm \sqrt{b^2 - 4ac}}{2a}\]

\subsection{Matrix Example}
A sample matrix:
\[
\begin{pmatrix}
a_{11} & a_{12} & a_{13} \\
a_{21} & a_{22} & a_{23} \\
a_{31} & a_{32} & a_{33}
\end{pmatrix}
\]

\section{Tables}
Below is a sample table showing data:

\begin{table}[h]
\centering
\begin{tabular}{@{}lrrr@{}}
\toprule
\textbf{Category} & \textbf{Value 1} & \textbf{Value 2} & \textbf{Total} \\
\midrule
Item A & 10.5 & 20.3 & 30.8 \\
Item B & 15.2 & 25.1 & 40.3 \\
Item C & 22.1 & 12.7 & 34.8 \\
\bottomrule
\end{tabular}
\caption{Sample Data Table}
\label{tab:sample}
\end{table}

\section{Lists}
\subsection{Bullet Points}
\begin{itemize}
  \item First item
  \item Second item
  \item Third item with sub-items:
    \begin{itemize}
      \item Sub-item 1
      \item Sub-item 2
    \end{itemize}
\end{itemize}

\subsection{Numbered List}
\begin{enumerate}
  \item First step
  \item Second step
  \item Third step
\end{enumerate}

\section{Citations and References}
When you use the citation search feature, references will be automatically managed. For example, you might cite a paper like this: \cite{example2024}. All citations will be automatically added to your \texttt{cite.bib} file.

\section{Conclusion}
This template demonstrates basic \LaTeX{} features including mathematical formulas, tables, and lists. The bibliography system is now set up to work seamlessly with Bibby AI's citation management. Feel free to modify it for your needs!

\bibliographystyle{plain}
\bibliography{cite}

\end{document}\documentclass{article}
\usepackage[utf8]{inputenc}
\usepackage{booktabs}
\usepackage{amsmath}
\usepackage{natbib}

\title{My Document}
\author{Your Name}
\date{\today}

\begin{document}

\maketitle

\section{Introduction}
This is a sample document demonstrating various \LaTeX{} features including mathematical equations, tables, and formatting. You can cite papers using the citation search feature, and they will be automatically added to your \texttt{cite.bib} file.

\section{Mathematical Expressions}
Here are some example equations:

\subsection{Quadratic Formula}
The quadratic formula is given by:
\[x = \frac{-b \pm \sqrt{b^2 - 4ac}}{2a}\]

\subsection{Matrix Example}
A sample matrix:
\[
\begin{pmatrix}
a_{11} & a_{12} & a_{13} \\
a_{21} & a_{22} & a_{23} \\
a_{31} & a_{32} & a_{33}
\end{pmatrix}
\]

\section{Tables}
Below is a sample table showing data:

\begin{table}[h]
\centering
\begin{tabular}{@{}lrrr@{}}
\toprule
\textbf{Category} & \textbf{Value 1} & \textbf{Value 2} & \textbf{Total} \\
\midrule
Item A & 10.5 & 20.3 & 30.8 \\
Item B & 15.2 & 25.1 & 40.3 \\
Item C & 22.1 & 12.7 & 34.8 \\
\bottomrule
\end{tabular}
\caption{Sample Data Table}
\label{tab:sample}
\end{table}

\section{Lists}
\subsection{Bullet Points}
\begin{itemize}
  \item First item
  \item Second item
  \item Third item with sub-items:
    \begin{itemize}
      \item Sub-item 1
      \item Sub-item 2
    \end{itemize}
\end{itemize}

\subsection{Numbered List}
\begin{enumerate}
  \item First step
  \item Second step
  \item Third step
\end{enumerate}

\section{Citations and References}
When you use the citation search feature, references will be automatically managed. For example, you might cite a paper like this: \cite{example2024}. All citations will be automatically added to your \texttt{cite.bib} file.

\section{Conclusion}
This template demonstrates basic \LaTeX{} features including mathematical formulas, tables, and lists. The bibliography system is now set up to work seamlessly with Bibby AI's citation management. Feel free to modify it for your needs!

\bibliographystyle{plain}
\bibliography{cite}

\end{document}\documentclass{article}
\usepackage[utf8]{inputenc}
\usepackage{booktabs}
\usepackage{amsmath}
\usepackage{natbib}

\title{My Document}
\author{Your Name}
\date{\today}

\begin{document}

\maketitle

\section{Introduction}
This is a sample document demonstrating various \LaTeX{} features including mathematical equations, tables, and formatting. You can cite papers using the citation search feature, and they will be automatically added to your \texttt{cite.bib} file.

\section{Mathematical Expressions}
Here are some example equations:

\subsection{Quadratic Formula}
The quadratic formula is given by:
\[x = \frac{-b \pm \sqrt{b^2 - 4ac}}{2a}\]

\subsection{Matrix Example}
A sample matrix:
\[
\begin{pmatrix}
a_{11} & a_{12} & a_{13} \\
a_{21} & a_{22} & a_{23} \\
a_{31} & a_{32} & a_{33}
\end{pmatrix}
\]

\section{Tables}
Below is a sample table showing data:

\begin{table}[h]
\centering
\begin{tabular}{@{}lrrr@{}}
\toprule
\textbf{Category} & \textbf{Value 1} & \textbf{Value 2} & \textbf{Total} \\
\midrule
Item A & 10.5 & 20.3 & 30.8 \\
Item B & 15.2 & 25.1 & 40.3 \\
Item C & 22.1 & 12.7 & 34.8 \\
\bottomrule
\end{tabular}
\caption{Sample Data Table}
\label{tab:sample}
\end{table}

\section{Lists}
\subsection{Bullet Points}
\begin{itemize}
  \item First item
  \item Second item
  \item Third item with sub-items:
    \begin{itemize}
      \item Sub-item 1
      \item Sub-item 2
    \end{itemize}
\end{itemize}

\subsection{Numbered List}
\begin{enumerate}
  \item First step
  \item Second step
  \item Third step
\end{enumerate}

\section{Citations and References}
When you use the citation search feature, references will be automatically managed. For example, you might cite a paper like this: \cite{example2024}. All citations will be automatically added to your \texttt{cite.bib} file.

\section{Conclusion}
This template demonstrates basic \LaTeX{} features including mathematical formulas, tables, and lists. The bibliography system is now set up to work seamlessly with Bibby AI's citation management. Feel free to modify it for your needs!

\bibliographystyle{plain}
\bibliography{cite}

\end{document}\documentclass{article}
\usepackage[utf8]{inputenc}
\usepackage{booktabs}
\usepackage{amsmath}
\usepackage{natbib}

\title{My Document}
\author{Your Name}
\date{\today}

\begin{document}

\maketitle

\section{Introduction}
This is a sample document demonstrating various \LaTeX{} features including mathematical equations, tables, and formatting. You can cite papers using the citation search feature, and they will be automatically added to your \texttt{cite.bib} file.

\section{Mathematical Expressions}
Here are some example equations:

\subsection{Quadratic Formula}
The quadratic formula is given by:
\[x = \frac{-b \pm \sqrt{b^2 - 4ac}}{2a}\]

\subsection{Matrix Example}
A sample matrix:
\[
\begin{pmatrix}
a_{11} & a_{12} & a_{13} \\
a_{21} & a_{22} & a_{23} \\
a_{31} & a_{32} & a_{33}
\end{pmatrix}
\]

\section{Tables}
Below is a sample table showing data:

\begin{table}[h]
\centering
\begin{tabular}{@{}lrrr@{}}
\toprule
\textbf{Category} & \textbf{Value 1} & \textbf{Value 2} & \textbf{Total} \\
\midrule
Item A & 10.5 & 20.3 & 30.8 \\
Item B & 15.2 & 25.1 & 40.3 \\
Item C & 22.1 & 12.7 & 34.8 \\
\bottomrule
\end{tabular}
\caption{Sample Data Table}
\label{tab:sample}
\end{table}

\section{Lists}
\subsection{Bullet Points}
\begin{itemize}
  \item First item
  \item Second item
  \item Third item with sub-items:
    \begin{itemize}
      \item Sub-item 1
      \item Sub-item 2
    \end{itemize}
\end{itemize}

\subsection{Numbered List}
\begin{enumerate}
  \item First step
  \item Second step
  \item Third step
\end{enumerate}

\section{Citations and References}
When you use the citation search feature, references will be automatically managed. For example, you might cite a paper like this: \cite{example2024}. All citations will be automatically added to your \texttt{cite.bib} file.

\section{Conclusion}
This template demonstrates basic \LaTeX{} features including mathematical formulas, tables, and lists. The bibliography system is now set up to work seamlessly with Bibby AI's citation management. Feel free to modify it for your needs!

\bibliographystyle{plain}
\bibliography{cite}

\end{document}\documentclass{article}
\usepackage[utf8]{inputenc}
\usepackage{booktabs}
\usepackage{amsmath}
\usepackage{natbib}

\title{My Document}
\author{Your Name}
\date{\today}

\begin{document}

\maketitle

\section{Introduction}
This is a sample document demonstrating various \LaTeX{} features including mathematical equations, tables, and formatting. You can cite papers using the citation search feature, and they will be automatically added to your \texttt{cite.bib} file.

\section{Mathematical Expressions}
Here are some example equations:

\subsection{Quadratic Formula}
The quadratic formula is given by:
\[x = \frac{-b \pm \sqrt{b^2 - 4ac}}{2a}\]

\subsection{Matrix Example}
A sample matrix:
\[
\begin{pmatrix}
a_{11} & a_{12} & a_{13} \\
a_{21} & a_{22} & a_{23} \\
a_{31} & a_{32} & a_{33}
\end{pmatrix}
\]

\section{Tables}
Below is a sample table showing data:

\begin{table}[h]
\centering
\begin{tabular}{@{}lrrr@{}}
\toprule
\textbf{Category} & \textbf{Value 1} & \textbf{Value 2} & \textbf{Total} \\
\midrule
Item A & 10.5 & 20.3 & 30.8 \\
Item B & 15.2 & 25.1 & 40.3 \\
Item C & 22.1 & 12.7 & 34.8 \\
\bottomrule
\end{tabular}
\caption{Sample Data Table}
\label{tab:sample}
\end{table}

\section{Lists}
\subsection{Bullet Points}
\begin{itemize}
  \item First item
  \item Second item
  \item Third item with sub-items:
    \begin{itemize}
      \item Sub-item 1
      \item Sub-item 2
    \end{itemize}
\end{itemize}

\subsection{Numbered List}
\begin{enumerate}
  \item First step
  \item Second step
  \item Third step
\end{enumerate}

\section{Citations and References}
When you use the citation search feature, references will be automatically managed. For example, you might cite a paper like this: \cite{example2024}. All citations will be automatically added to your \texttt{cite.bib} file.

\section{Conclusion}
This template demonstrates basic \LaTeX{} features including mathematical formulas, tables, and lists. The bibliography system is now set up to work seamlessly with Bibby AI's citation management. Feel free to modify it for your needs!

\bibliographystyle{plain}
\bibliography{cite}

\end{document}\documentclass{article}
\usepackage[utf8]{inputenc}
\usepackage{booktabs}
\usepackage{amsmath}
\usepackage{natbib}

\title{My Document}
\author{Your Name}
\date{\today}

\begin{document}

\maketitle

\section{Introduction}
This is a sample document demonstrating various \LaTeX{} features including mathematical equations, tables, and formatting. You can cite papers using the citation search feature, and they will be automatically added to your \texttt{cite.bib} file.

\section{Mathematical Expressions}
Here are some example equations:

\subsection{Quadratic Formula}
The quadratic formula is given by:
\[x = \frac{-b \pm \sqrt{b^2 - 4ac}}{2a}\]

\subsection{Matrix Example}
A sample matrix:
\[
\begin{pmatrix}
a_{11} & a_{12} & a_{13} \\
a_{21} & a_{22} & a_{23} \\
a_{31} & a_{32} & a_{33}
\end{pmatrix}
\]

\section{Tables}
Below is a sample table showing data:

\begin{table}[h]
\centering
\begin{tabular}{@{}lrrr@{}}
\toprule
\textbf{Category} & \textbf{Value 1} & \textbf{Value 2} & \textbf{Total} \\
\midrule
Item A & 10.5 & 20.3 & 30.8 \\
Item B & 15.2 & 25.1 & 40.3 \\
Item C & 22.1 & 12.7 & 34.8 \\
\bottomrule
\end{tabular}
\caption{Sample Data Table}
\label{tab:sample}
\end{table}

\section{Lists}
\subsection{Bullet Points}
\begin{itemize}
  \item First item
  \item Second item
  \item Third item with sub-items:
    \begin{itemize}
      \item Sub-item 1
      \item Sub-item 2
    \end{itemize}
\end{itemize}

\subsection{Numbered List}
\begin{enumerate}
  \item First step
  \item Second step
  \item Third step
\end{enumerate}

\section{Citations and References}
When you use the citation search feature, references will be automatically managed. For example, you might cite a paper like this: \cite{example2024}. All citations will be automatically added to your \texttt{cite.bib} file.

\section{Conclusion}
This template demonstrates basic \LaTeX{} features including mathematical formulas, tables, and lists. The bibliography system is now set up to work seamlessly with Bibby AI's citation management. Feel free to modify it for your needs!

\bibliographystyle{plain}
\bibliography{cite}

\end{document}\documentclass{article}
\usepackage[utf8]{inputenc}
\usepackage{booktabs}
\usepackage{amsmath}
\usepackage{natbib}

\title{My Document}
\author{Your Name}
\date{\today}

\begin{document}

\maketitle

\section{Introduction}
This is a sample document demonstrating various \LaTeX{} features including mathematical equations, tables, and formatting. You can cite papers using the citation search feature, and they will be automatically added to your \texttt{cite.bib} file.

\section{Mathematical Expressions}
Here are some example equations:

\subsection{Quadratic Formula}
The quadratic formula is given by:
\[x = \frac{-b \pm \sqrt{b^2 - 4ac}}{2a}\]

\subsection{Matrix Example}
A sample matrix:
\[
\begin{pmatrix}
a_{11} & a_{12} & a_{13} \\
a_{21} & a_{22} & a_{23} \\
a_{31} & a_{32} & a_{33}
\end{pmatrix}
\]

\section{Tables}
Below is a sample table showing data:

\begin{table}[h]
\centering
\begin{tabular}{@{}lrrr@{}}
\toprule
\textbf{Category} & \textbf{Value 1} & \textbf{Value 2} & \textbf{Total} \\
\midrule
Item A & 10.5 & 20.3 & 30.8 \\
Item B & 15.2 & 25.1 & 40.3 \\
Item C & 22.1 & 12.7 & 34.8 \\
\bottomrule
\end{tabular}
\caption{Sample Data Table}
\label{tab:sample}
\end{table}

\section{Lists}
\subsection{Bullet Points}
\begin{itemize}
  \item First item
  \item Second item
  \item Third item with sub-items:
    \begin{itemize}
      \item Sub-item 1
      \item Sub-item 2
    \end{itemize}
\end{itemize}

\subsection{Numbered List}
\begin{enumerate}
  \item First step
  \item Second step
  \item Third step
\end{enumerate}

\section{Citations and References}
When you use the citation search feature, references will be automatically managed. For example, you might cite a paper like this: \cite{example2024}. All citations will be automatically added to your \texttt{cite.bib} file.

\section{Conclusion}
This template demonstrates basic \LaTeX{} features including mathematical formulas, tables, and lists. The bibliography system is now set up to work seamlessly with Bibby AI's citation management. Feel free to modify it for your needs!

\bibliographystyle{plain}
\bibliography{cite}

\end{document}\documentclass{article}
\usepackage[utf8]{inputenc}
\usepackage{booktabs}
\usepackage{amsmath}
\usepackage{natbib}

\title{My Document}
\author{Your Name}
\date{\today}

\begin{document}

\maketitle

\section{Introduction}
This is a sample document demonstrating various \LaTeX{} features including mathematical equations, tables, and formatting. You can cite papers using the citation search feature, and they will be automatically added to your \texttt{cite.bib} file.

\section{Mathematical Expressions}
Here are some example equations:

\subsection{Quadratic Formula}
The quadratic formula is given by:
\[x = \frac{-b \pm \sqrt{b^2 - 4ac}}{2a}\]

\subsection{Matrix Example}
A sample matrix:
\[
\begin{pmatrix}
a_{11} & a_{12} & a_{13} \\
a_{21} & a_{22} & a_{23} \\
a_{31} & a_{32} & a_{33}
\end{pmatrix}
\]

\section{Tables}
Below is a sample table showing data:

\begin{table}[h]
\centering
\begin{tabular}{@{}lrrr@{}}
\toprule
\textbf{Category} & \textbf{Value 1} & \textbf{Value 2} & \textbf{Total} \\
\midrule
Item A & 10.5 & 20.3 & 30.8 \\
Item B & 15.2 & 25.1 & 40.3 \\
Item C & 22.1 & 12.7 & 34.8 \\
\bottomrule
\end{tabular}
\caption{Sample Data Table}
\label{tab:sample}
\end{table}

\section{Lists}
\subsection{Bullet Points}
\begin{itemize}
  \item First item
  \item Second item
  \item Third item with sub-items:
    \begin{itemize}
      \item Sub-item 1
      \item Sub-item 2
    \end{itemize}
\end{itemize}

\subsection{Numbered List}
\begin{enumerate}
  \item First step
  \item Second step
  \item Third step
\end{enumerate}

\section{Citations and References}
When you use the citation search feature, references will be automatically managed. For example, you might cite a paper like this: \cite{example2024}. All citations will be automatically added to your \texttt{cite.bib} file.

\section{Conclusion}
This template demonstrates basic \LaTeX{} features including mathematical formulas, tables, and lists. The bibliography system is now set up to work seamlessly with Bibby AI's citation management. Feel free to modify it for your needs!

\bibliographystyle{plain}
\bibliography{cite}

\end{document}\documentclass{article}
\usepackage[utf8]{inputenc}
\usepackage{booktabs}
\usepackage{amsmath}
\usepackage{natbib}

\title{My Document}
\author{Your Name}
\date{\today}

\begin{document}

\maketitle

\section{Introduction}
This is a sample document demonstrating various \LaTeX{} features including mathematical equations, tables, and formatting. You can cite papers using the citation search feature, and they will be automatically added to your \texttt{cite.bib} file.

\section{Mathematical Expressions}
Here are some example equations:

\subsection{Quadratic Formula}
The quadratic formula is given by:
\[x = \frac{-b \pm \sqrt{b^2 - 4ac}}{2a}\]

\subsection{Matrix Example}
A sample matrix:
\[
\begin{pmatrix}
a_{11} & a_{12} & a_{13} \\
a_{21} & a_{22} & a_{23} \\
a_{31} & a_{32} & a_{33}
\end{pmatrix}
\]

\section{Tables}
Below is a sample table showing data:

\begin{table}[h]
\centering
\begin{tabular}{@{}lrrr@{}}
\toprule
\textbf{Category} & \textbf{Value 1} & \textbf{Value 2} & \textbf{Total} \\
\midrule
Item A & 10.5 & 20.3 & 30.8 \\
Item B & 15.2 & 25.1 & 40.3 \\
Item C & 22.1 & 12.7 & 34.8 \\
\bottomrule
\end{tabular}
\caption{Sample Data Table}
\label{tab:sample}
\end{table}

\section{Lists}
\subsection{Bullet Points}
\begin{itemize}
  \item First item
  \item Second item
  \item Third item with sub-items:
    \begin{itemize}
      \item Sub-item 1
      \item Sub-item 2
    \end{itemize}
\end{itemize}

\subsection{Numbered List}
\begin{enumerate}
  \item First step
  \item Second step
  \item Third step
\end{enumerate}

\section{Citations and References}
When you use the citation search feature, references will be automatically managed. For example, you might cite a paper like this: \cite{example2024}. All citations will be automatically added to your \texttt{cite.bib} file.

\section{Conclusion}
This template demonstrates basic \LaTeX{} features including mathematical formulas, tables, and lists. The bibliography system is now set up to work seamlessly with Bibby AI's citation management. Feel free to modify it for your needs!

\bibliographystyle{plain}
\bibliography{cite}

\end{document}\documentclass{article}
\usepackage[utf8]{inputenc}
\usepackage{booktabs}
\usepackage{amsmath}
\usepackage{natbib}

\title{My Document}
\author{Your Name}
\date{\today}

\begin{document}

\maketitle

\section{Introduction}
This is a sample document demonstrating various \LaTeX{} features including mathematical equations, tables, and formatting. You can cite papers using the citation search feature, and they will be automatically added to your \texttt{cite.bib} file.

\section{Mathematical Expressions}
Here are some example equations:

\subsection{Quadratic Formula}
The quadratic formula is given by:
\[x = \frac{-b \pm \sqrt{b^2 - 4ac}}{2a}\]

\subsection{Matrix Example}
A sample matrix:
\[
\begin{pmatrix}
a_{11} & a_{12} & a_{13} \\
a_{21} & a_{22} & a_{23} \\
a_{31} & a_{32} & a_{33}
\end{pmatrix}
\]

\section{Tables}
Below is a sample table showing data:

\begin{table}[h]
\centering
\begin{tabular}{@{}lrrr@{}}
\toprule
\textbf{Category} & \textbf{Value 1} & \textbf{Value 2} & \textbf{Total} \\
\midrule
Item A & 10.5 & 20.3 & 30.8 \\
Item B & 15.2 & 25.1 & 40.3 \\
Item C & 22.1 & 12.7 & 34.8 \\
\bottomrule
\end{tabular}
\caption{Sample Data Table}
\label{tab:sample}
\end{table}

\section{Lists}
\subsection{Bullet Points}
\begin{itemize}
  \item First item
  \item Second item
  \item Third item with sub-items:
    \begin{itemize}
      \item Sub-item 1
      \item Sub-item 2
    \end{itemize}
\end{itemize}

\subsection{Numbered List}
\begin{enumerate}
  \item First step
  \item Second step
  \item Third step
\end{enumerate}

\section{Citations and References}
When you use the citation search feature, references will be automatically managed. For example, you might cite a paper like this: \cite{example2024}. All citations will be automatically added to your \texttt{cite.bib} file.

\section{Conclusion}
This template demonstrates basic \LaTeX{} features including mathematical formulas, tables, and lists. The bibliography system is now set up to work seamlessly with Bibby AI's citation management. Feel free to modify it for your needs!

\bibliographystyle{plain}
\bibliography{cite}

\end{document}\documentclass{article}
\usepackage[utf8]{inputenc}
\usepackage{booktabs}
\usepackage{amsmath}
\usepackage{natbib}

\title{My Document}
\author{Your Name}
\date{\today}

\begin{document}

\maketitle

\section{Introduction}
This is a sample document demonstrating various \LaTeX{} features including mathematical equations, tables, and formatting. You can cite papers using the citation search feature, and they will be automatically added to your \texttt{cite.bib} file.

\section{Mathematical Expressions}
Here are some example equations:

\subsection{Quadratic Formula}
The quadratic formula is given by:
\[x = \frac{-b \pm \sqrt{b^2 - 4ac}}{2a}\]

\subsection{Matrix Example}
A sample matrix:
\[
\begin{pmatrix}
a_{11} & a_{12} & a_{13} \\
a_{21} & a_{22} & a_{23} \\
a_{31} & a_{32} & a_{33}
\end{pmatrix}
\]

\section{Tables}
Below is a sample table showing data:

\begin{table}[h]
\centering
\begin{tabular}{@{}lrrr@{}}
\toprule
\textbf{Category} & \textbf{Value 1} & \textbf{Value 2} & \textbf{Total} \\
\midrule
Item A & 10.5 & 20.3 & 30.8 \\
Item B & 15.2 & 25.1 & 40.3 \\
Item C & 22.1 & 12.7 & 34.8 \\
\bottomrule
\end{tabular}
\caption{Sample Data Table}
\label{tab:sample}
\end{table}

\section{Lists}
\subsection{Bullet Points}
\begin{itemize}
  \item First item
  \item Second item
  \item Third item with sub-items:
    \begin{itemize}
      \item Sub-item 1
      \item Sub-item 2
    \end{itemize}
\end{itemize}

\subsection{Numbered List}
\begin{enumerate}
  \item First step
  \item Second step
  \item Third step
\end{enumerate}

\section{Citations and References}
When you use the citation search feature, references will be automatically managed. For example, you might cite a paper like this: \cite{example2024}. All citations will be automatically added to your \texttt{cite.bib} file.

\section{Conclusion}
This template demonstrates basic \LaTeX{} features including mathematical formulas, tables, and lists. The bibliography system is now set up to work seamlessly with Bibby AI's citation management. Feel free to modify it for your needs!

\bibliographystyle{plain}
\bibliography{cite}

\end{document}\documentclass{article}
\usepackage[utf8]{inputenc}
\usepackage{booktabs}
\usepackage{amsmath}
\usepackage{natbib}

\title{My Document}
\author{Your Name}
\date{\today}

\begin{document}

\maketitle

\section{Introduction}
This is a sample document demonstrating various \LaTeX{} features including mathematical equations, tables, and formatting. You can cite papers using the citation search feature, and they will be automatically added to your \texttt{cite.bib} file.

\section{Mathematical Expressions}
Here are some example equations:

\subsection{Quadratic Formula}
The quadratic formula is given by:
\[x = \frac{-b \pm \sqrt{b^2 - 4ac}}{2a}\]

\subsection{Matrix Example}
A sample matrix:
\[
\begin{pmatrix}
a_{11} & a_{12} & a_{13} \\
a_{21} & a_{22} & a_{23} \\
a_{31} & a_{32} & a_{33}
\end{pmatrix}
\]

\section{Tables}
Below is a sample table showing data:

\begin{table}[h]
\centering
\begin{tabular}{@{}lrrr@{}}
\toprule
\textbf{Category} & \textbf{Value 1} & \textbf{Value 2} & \textbf{Total} \\
\midrule
Item A & 10.5 & 20.3 & 30.8 \\
Item B & 15.2 & 25.1 & 40.3 \\
Item C & 22.1 & 12.7 & 34.8 \\
\bottomrule
\end{tabular}
\caption{Sample Data Table}
\label{tab:sample}
\end{table}

\section{Lists}
\subsection{Bullet Points}
\begin{itemize}
  \item First item
  \item Second item
  \item Third item with sub-items:
    \begin{itemize}
      \item Sub-item 1
      \item Sub-item 2
    \end{itemize}
\end{itemize}

\subsection{Numbered List}
\begin{enumerate}
  \item First step
  \item Second step
  \item Third step
\end{enumerate}

\section{Citations and References}
When you use the citation search feature, references will be automatically managed. For example, you might cite a paper like this: \cite{example2024}. All citations will be automatically added to your \texttt{cite.bib} file.

\section{Conclusion}
This template demonstrates basic \LaTeX{} features including mathematical formulas, tables, and lists. The bibliography system is now set up to work seamlessly with Bibby AI's citation management. Feel free to modify it for your needs!

\bibliographystyle{plain}
\bibliography{cite}

\end{document}\documentclass{article}
\usepackage[utf8]{inputenc}
\usepackage{booktabs}
\usepackage{amsmath}
\usepackage{natbib}

\title{My Document}
\author{Your Name}
\date{\today}

\begin{document}

\maketitle

\section{Introduction}
This is a sample document demonstrating various \LaTeX{} features including mathematical equations, tables, and formatting. You can cite papers using the citation search feature, and they will be automatically added to your \texttt{cite.bib} file.

\section{Mathematical Expressions}
Here are some example equations:

\subsection{Quadratic Formula}
The quadratic formula is given by:
\[x = \frac{-b \pm \sqrt{b^2 - 4ac}}{2a}\]

\subsection{Matrix Example}
A sample matrix:
\[
\begin{pmatrix}
a_{11} & a_{12} & a_{13} \\
a_{21} & a_{22} & a_{23} \\
a_{31} & a_{32} & a_{33}
\end{pmatrix}
\]

\section{Tables}
Below is a sample table showing data:

\begin{table}[h]
\centering
\begin{tabular}{@{}lrrr@{}}
\toprule
\textbf{Category} & \textbf{Value 1} & \textbf{Value 2} & \textbf{Total} \\
\midrule
Item A & 10.5 & 20.3 & 30.8 \\
Item B & 15.2 & 25.1 & 40.3 \\
Item C & 22.1 & 12.7 & 34.8 \\
\bottomrule
\end{tabular}
\caption{Sample Data Table}
\label{tab:sample}
\end{table}

\section{Lists}
\subsection{Bullet Points}
\begin{itemize}
  \item First item
  \item Second item
  \item Third item with sub-items:
    \begin{itemize}
      \item Sub-item 1
      \item Sub-item 2
    \end{itemize}
\end{itemize}

\subsection{Numbered List}
\begin{enumerate}
  \item First step
  \item Second step
  \item Third step
\end{enumerate}

\section{Citations and References}
When you use the citation search feature, references will be automatically managed. For example, you might cite a paper like this: \cite{example2024}. All citations will be automatically added to your \texttt{cite.bib} file.

\section{Conclusion}
This template demonstrates basic \LaTeX{} features including mathematical formulas, tables, and lists. The bibliography system is now set up to work seamlessly with Bibby AI's citation management. Feel free to modify it for your needs!

\bibliographystyle{plain}
\bibliography{cite}

\end{document}\documentclass{article}
\usepackage[utf8]{inputenc}
\usepackage{booktabs}
\usepackage{amsmath}
\usepackage{natbib}

\title{My Document}
\author{Your Name}
\date{\today}

\begin{document}

\maketitle

\section{Introduction}
This is a sample document demonstrating various \LaTeX{} features including mathematical equations, tables, and formatting. You can cite papers using the citation search feature, and they will be automatically added to your \texttt{cite.bib} file.

\section{Mathematical Expressions}
Here are some example equations:

\subsection{Quadratic Formula}
The quadratic formula is given by:
\[x = \frac{-b \pm \sqrt{b^2 - 4ac}}{2a}\]

\subsection{Matrix Example}
A sample matrix:
\[
\begin{pmatrix}
a_{11} & a_{12} & a_{13} \\
a_{21} & a_{22} & a_{23} \\
a_{31} & a_{32} & a_{33}
\end{pmatrix}
\]

\section{Tables}
Below is a sample table showing data:

\begin{table}[h]
\centering
\begin{tabular}{@{}lrrr@{}}
\toprule
\textbf{Category} & \textbf{Value 1} & \textbf{Value 2} & \textbf{Total} \\
\midrule
Item A & 10.5 & 20.3 & 30.8 \\
Item B & 15.2 & 25.1 & 40.3 \\
Item C & 22.1 & 12.7 & 34.8 \\
\bottomrule
\end{tabular}
\caption{Sample Data Table}
\label{tab:sample}
\end{table}

\section{Lists}
\subsection{Bullet Points}
\begin{itemize}
  \item First item
  \item Second item
  \item Third item with sub-items:
    \begin{itemize}
      \item Sub-item 1
      \item Sub-item 2
    \end{itemize}
\end{itemize}

\subsection{Numbered List}
\begin{enumerate}
  \item First step
  \item Second step
  \item Third step
\end{enumerate}

\section{Citations and References}
When you use the citation search feature, references will be automatically managed. For example, you might cite a paper like this: \cite{example2024}. All citations will be automatically added to your \texttt{cite.bib} file.

\section{Conclusion}
This template demonstrates basic \LaTeX{} features including mathematical formulas, tables, and lists. The bibliography system is now set up to work seamlessly with Bibby AI's citation management. Feel free to modify it for your needs!

\bibliographystyle{plain}
\bibliography{cite}

\end{document}\documentclass{article}
\usepackage[utf8]{inputenc}
\usepackage{booktabs}
\usepackage{amsmath}
\usepackage{natbib}

\title{My Document}
\author{Your Name}
\date{\today}

\begin{document}

\maketitle

\section{Introduction}
This is a sample document demonstrating various \LaTeX{} features including mathematical equations, tables, and formatting. You can cite papers using the citation search feature, and they will be automatically added to your \texttt{cite.bib} file.

\section{Mathematical Expressions}
Here are some example equations:

\subsection{Quadratic Formula}
The quadratic formula is given by:
\[x = \frac{-b \pm \sqrt{b^2 - 4ac}}{2a}\]

\subsection{Matrix Example}
A sample matrix:
\[
\begin{pmatrix}
a_{11} & a_{12} & a_{13} \\
a_{21} & a_{22} & a_{23} \\
a_{31} & a_{32} & a_{33}
\end{pmatrix}
\]

\section{Tables}
Below is a sample table showing data:

\begin{table}[h]
\centering
\begin{tabular}{@{}lrrr@{}}
\toprule
\textbf{Category} & \textbf{Value 1} & \textbf{Value 2} & \textbf{Total} \\
\midrule
Item A & 10.5 & 20.3 & 30.8 \\
Item B & 15.2 & 25.1 & 40.3 \\
Item C & 22.1 & 12.7 & 34.8 \\
\bottomrule
\end{tabular}
\caption{Sample Data Table}
\label{tab:sample}
\end{table}

\section{Lists}
\subsection{Bullet Points}
\begin{itemize}
  \item First item
  \item Second item
  \item Third item with sub-items:
    \begin{itemize}
      \item Sub-item 1
      \item Sub-item 2
    \end{itemize}
\end{itemize}

\subsection{Numbered List}
\begin{enumerate}
  \item First step
  \item Second step
  \item Third step
\end{enumerate}

\section{Citations and References}
When you use the citation search feature, references will be automatically managed. For example, you might cite a paper like this: \cite{example2024}. All citations will be automatically added to your \texttt{cite.bib} file.

\section{Conclusion}
This template demonstrates basic \LaTeX{} features including mathematical formulas, tables, and lists. The bibliography system is now set up to work seamlessly with Bibby AI's citation management. Feel free to modify it for your needs!

\bibliographystyle{plain}
\bibliography{cite}

\end{document}\documentclass{article}
\usepackage[utf8]{inputenc}
\usepackage{booktabs}
\usepackage{amsmath}
\usepackage{natbib}

\title{My Document}
\author{Your Name}
\date{\today}

\begin{document}

\maketitle

\section{Introduction}
This is a sample document demonstrating various \LaTeX{} features including mathematical equations, tables, and formatting. You can cite papers using the citation search feature, and they will be automatically added to your \texttt{cite.bib} file.

\section{Mathematical Expressions}
Here are some example equations:

\subsection{Quadratic Formula}
The quadratic formula is given by:
\[x = \frac{-b \pm \sqrt{b^2 - 4ac}}{2a}\]

\subsection{Matrix Example}
A sample matrix:
\[
\begin{pmatrix}
a_{11} & a_{12} & a_{13} \\
a_{21} & a_{22} & a_{23} \\
a_{31} & a_{32} & a_{33}
\end{pmatrix}
\]

\section{Tables}
Below is a sample table showing data:

\begin{table}[h]
\centering
\begin{tabular}{@{}lrrr@{}}
\toprule
\textbf{Category} & \textbf{Value 1} & \textbf{Value 2} & \textbf{Total} \\
\midrule
Item A & 10.5 & 20.3 & 30.8 \\
Item B & 15.2 & 25.1 & 40.3 \\
Item C & 22.1 & 12.7 & 34.8 \\
\bottomrule
\end{tabular}
\caption{Sample Data Table}
\label{tab:sample}
\end{table}

\section{Lists}
\subsection{Bullet Points}
\begin{itemize}
  \item First item
  \item Second item
  \item Third item with sub-items:
    \begin{itemize}
      \item Sub-item 1
      \item Sub-item 2
    \end{itemize}
\end{itemize}

\subsection{Numbered List}
\begin{enumerate}
  \item First step
  \item Second step
  \item Third step
\end{enumerate}

\section{Citations and References}
When you use the citation search feature, references will be automatically managed. For example, you might cite a paper like this: \cite{example2024}. All citations will be automatically added to your \texttt{cite.bib} file.

\section{Conclusion}
This template demonstrates basic \LaTeX{} features including mathematical formulas, tables, and lists. The bibliography system is now set up to work seamlessly with Bibby AI's citation management. Feel free to modify it for your needs!

\bibliographystyle{plain}
\bibliography{cite}

\end{document}\documentclass{article}
\usepackage[utf8]{inputenc}
\usepackage{booktabs}
\usepackage{amsmath}
\usepackage{natbib}

\title{My Document}
\author{Your Name}
\date{\today}

\begin{document}

\maketitle

\section{Introduction}
This is a sample document demonstrating various \LaTeX{} features including mathematical equations, tables, and formatting. You can cite papers using the citation search feature, and they will be automatically added to your \texttt{cite.bib} file.

\section{Mathematical Expressions}
Here are some example equations:

\subsection{Quadratic Formula}
The quadratic formula is given by:
\[x = \frac{-b \pm \sqrt{b^2 - 4ac}}{2a}\]

\subsection{Matrix Example}
A sample matrix:
\[
\begin{pmatrix}
a_{11} & a_{12} & a_{13} \\
a_{21} & a_{22} & a_{23} \\
a_{31} & a_{32} & a_{33}
\end{pmatrix}
\]

\section{Tables}
Below is a sample table showing data:

\begin{table}[h]
\centering
\begin{tabular}{@{}lrrr@{}}
\toprule
\textbf{Category} & \textbf{Value 1} & \textbf{Value 2} & \textbf{Total} \\
\midrule
Item A & 10.5 & 20.3 & 30.8 \\
Item B & 15.2 & 25.1 & 40.3 \\
Item C & 22.1 & 12.7 & 34.8 \\
\bottomrule
\end{tabular}
\caption{Sample Data Table}
\label{tab:sample}
\end{table}

\section{Lists}
\subsection{Bullet Points}
\begin{itemize}
  \item First item
  \item Second item
  \item Third item with sub-items:
    \begin{itemize}
      \item Sub-item 1
      \item Sub-item 2
    \end{itemize}
\end{itemize}

\subsection{Numbered List}
\begin{enumerate}
  \item First step
  \item Second step
  \item Third step
\end{enumerate}

\section{Citations and References}
When you use the citation search feature, references will be automatically managed. For example, you might cite a paper like this: \cite{example2024}. All citations will be automatically added to your \texttt{cite.bib} file.

\section{Conclusion}
This template demonstrates basic \LaTeX{} features including mathematical formulas, tables, and lists. The bibliography system is now set up to work seamlessly with Bibby AI's citation management. Feel free to modify it for your needs!

\bibliographystyle{plain}
\bibliography{cite}

\end{document}\documentclass{article}
\usepackage[utf8]{inputenc}
\usepackage{booktabs}
\usepackage{amsmath}
\usepackage{natbib}

\title{My Document}
\author{Your Name}
\date{\today}

\begin{document}

\maketitle

\section{Introduction}
This is a sample document demonstrating various \LaTeX{} features including mathematical equations, tables, and formatting. You can cite papers using the citation search feature, and they will be automatically added to your \texttt{cite.bib} file.

\section{Mathematical Expressions}
Here are some example equations:

\subsection{Quadratic Formula}
The quadratic formula is given by:
\[x = \frac{-b \pm \sqrt{b^2 - 4ac}}{2a}\]

\subsection{Matrix Example}
A sample matrix:
\[
\begin{pmatrix}
a_{11} & a_{12} & a_{13} \\
a_{21} & a_{22} & a_{23} \\
a_{31} & a_{32} & a_{33}
\end{pmatrix}
\]

\section{Tables}
Below is a sample table showing data:

\begin{table}[h]
\centering
\begin{tabular}{@{}lrrr@{}}
\toprule
\textbf{Category} & \textbf{Value 1} & \textbf{Value 2} & \textbf{Total} \\
\midrule
Item A & 10.5 & 20.3 & 30.8 \\
Item B & 15.2 & 25.1 & 40.3 \\
Item C & 22.1 & 12.7 & 34.8 \\
\bottomrule
\end{tabular}
\caption{Sample Data Table}
\label{tab:sample}
\end{table}

\section{Lists}
\subsection{Bullet Points}
\begin{itemize}
  \item First item
  \item Second item
  \item Third item with sub-items:
    \begin{itemize}
      \item Sub-item 1
      \item Sub-item 2
    \end{itemize}
\end{itemize}

\subsection{Numbered List}
\begin{enumerate}
  \item First step
  \item Second step
  \item Third step
\end{enumerate}

\section{Citations and References}
When you use the citation search feature, references will be automatically managed. For example, you might cite a paper like this: \cite{example2024}. All citations will be automatically added to your \texttt{cite.bib} file.

\section{Conclusion}
This template demonstrates basic \LaTeX{} features including mathematical formulas, tables, and lists. The bibliography system is now set up to work seamlessly with Bibby AI's citation management. Feel free to modify it for your needs!

\bibliographystyle{plain}
\bibliography{cite}

\end{document}\documentclass{article}
\usepackage[utf8]{inputenc}
\usepackage{booktabs}
\usepackage{amsmath}
\usepackage{natbib}

\title{My Document}
\author{Your Name}
\date{\today}

\begin{document}

\maketitle

\section{Introduction}
This is a sample document demonstrating various \LaTeX{} features including mathematical equations, tables, and formatting. You can cite papers using the citation search feature, and they will be automatically added to your \texttt{cite.bib} file.

\section{Mathematical Expressions}
Here are some example equations:

\subsection{Quadratic Formula}
The quadratic formula is given by:
\[x = \frac{-b \pm \sqrt{b^2 - 4ac}}{2a}\]

\subsection{Matrix Example}
A sample matrix:
\[
\begin{pmatrix}
a_{11} & a_{12} & a_{13} \\
a_{21} & a_{22} & a_{23} \\
a_{31} & a_{32} & a_{33}
\end{pmatrix}
\]

\section{Tables}
Below is a sample table showing data:

\begin{table}[h]
\centering
\begin{tabular}{@{}lrrr@{}}
\toprule
\textbf{Category} & \textbf{Value 1} & \textbf{Value 2} & \textbf{Total} \\
\midrule
Item A & 10.5 & 20.3 & 30.8 \\
Item B & 15.2 & 25.1 & 40.3 \\
Item C & 22.1 & 12.7 & 34.8 \\
\bottomrule
\end{tabular}
\caption{Sample Data Table}
\label{tab:sample}
\end{table}

\section{Lists}
\subsection{Bullet Points}
\begin{itemize}
  \item First item
  \item Second item
  \item Third item with sub-items:
    \begin{itemize}
      \item Sub-item 1
      \item Sub-item 2
    \end{itemize}
\end{itemize}

\subsection{Numbered List}
\begin{enumerate}
  \item First step
  \item Second step
  \item Third step
\end{enumerate}

\section{Citations and References}
When you use the citation search feature, references will be automatically managed. For example, you might cite a paper like this: \cite{example2024}. All citations will be automatically added to your \texttt{cite.bib} file.

\section{Conclusion}
This template demonstrates basic \LaTeX{} features including mathematical formulas, tables, and lists. The bibliography system is now set up to work seamlessly with Bibby AI's citation management. Feel free to modify it for your needs!

\bibliographystyle{plain}
\bibliography{cite}

\end{document}\documentclass{article}
\usepackage[utf8]{inputenc}
\usepackage{booktabs}
\usepackage{amsmath}
\usepackage{natbib}

\title{My Document}
\author{Your Name}
\date{\today}

\begin{document}

\maketitle

\section{Introduction}
This is a sample document demonstrating various \LaTeX{} features including mathematical equations, tables, and formatting. You can cite papers using the citation search feature, and they will be automatically added to your \texttt{cite.bib} file.

\section{Mathematical Expressions}
Here are some example equations:

\subsection{Quadratic Formula}
The quadratic formula is given by:
\[x = \frac{-b \pm \sqrt{b^2 - 4ac}}{2a}\]

\subsection{Matrix Example}
A sample matrix:
\[
\begin{pmatrix}
a_{11} & a_{12} & a_{13} \\
a_{21} & a_{22} & a_{23} \\
a_{31} & a_{32} & a_{33}
\end{pmatrix}
\]

\section{Tables}
Below is a sample table showing data:

\begin{table}[h]
\centering
\begin{tabular}{@{}lrrr@{}}
\toprule
\textbf{Category} & \textbf{Value 1} & \textbf{Value 2} & \textbf{Total} \\
\midrule
Item A & 10.5 & 20.3 & 30.8 \\
Item B & 15.2 & 25.1 & 40.3 \\
Item C & 22.1 & 12.7 & 34.8 \\
\bottomrule
\end{tabular}
\caption{Sample Data Table}
\label{tab:sample}
\end{table}

\section{Lists}
\subsection{Bullet Points}
\begin{itemize}
  \item First item
  \item Second item
  \item Third item with sub-items:
    \begin{itemize}
      \item Sub-item 1
      \item Sub-item 2
    \end{itemize}
\end{itemize}

\subsection{Numbered List}
\begin{enumerate}
  \item First step
  \item Second step
  \item Third step
\end{enumerate}

\section{Citations and References}
When you use the citation search feature, references will be automatically managed. For example, you might cite a paper like this: \cite{example2024}. All citations will be automatically added to your \texttt{cite.bib} file.

\section{Conclusion}
This template demonstrates basic \LaTeX{} features including mathematical formulas, tables, and lists. The bibliography system is now set up to work seamlessly with Bibby AI's citation management. Feel free to modify it for your needs!

\bibliographystyle{plain}
\bibliography{cite}

\end{document}\documentclass{article}
\usepackage[utf8]{inputenc}
\usepackage{booktabs}
\usepackage{amsmath}
\usepackage{natbib}

\title{My Document}
\author{Your Name}
\date{\today}

\begin{document}

\maketitle

\section{Introduction}
This is a sample document demonstrating various \LaTeX{} features including mathematical equations, tables, and formatting. You can cite papers using the citation search feature, and they will be automatically added to your \texttt{cite.bib} file.

\section{Mathematical Expressions}
Here are some example equations:

\subsection{Quadratic Formula}
The quadratic formula is given by:
\[x = \frac{-b \pm \sqrt{b^2 - 4ac}}{2a}\]

\subsection{Matrix Example}
A sample matrix:
\[
\begin{pmatrix}
a_{11} & a_{12} & a_{13} \\
a_{21} & a_{22} & a_{23} \\
a_{31} & a_{32} & a_{33}
\end{pmatrix}
\]

\section{Tables}
Below is a sample table showing data:

\begin{table}[h]
\centering
\begin{tabular}{@{}lrrr@{}}
\toprule
\textbf{Category} & \textbf{Value 1} & \textbf{Value 2} & \textbf{Total} \\
\midrule
Item A & 10.5 & 20.3 & 30.8 \\
Item B & 15.2 & 25.1 & 40.3 \\
Item C & 22.1 & 12.7 & 34.8 \\
\bottomrule
\end{tabular}
\caption{Sample Data Table}
\label{tab:sample}
\end{table}

\section{Lists}
\subsection{Bullet Points}
\begin{itemize}
  \item First item
  \item Second item
  \item Third item with sub-items:
    \begin{itemize}
      \item Sub-item 1
      \item Sub-item 2
    \end{itemize}
\end{itemize}

\subsection{Numbered List}
\begin{enumerate}
  \item First step
  \item Second step
  \item Third step
\end{enumerate}

\section{Citations and References}
When you use the citation search feature, references will be automatically managed. For example, you might cite a paper like this: \cite{example2024}. All citations will be automatically added to your \texttt{cite.bib} file.

\section{Conclusion}
This template demonstrates basic \LaTeX{} features including mathematical formulas, tables, and lists. The bibliography system is now set up to work seamlessly with Bibby AI's citation management. Feel free to modify it for your needs!

\bibliographystyle{plain}
\bibliography{cite}

\end{document}\documentclass{article}
\usepackage[utf8]{inputenc}
\usepackage{booktabs}
\usepackage{amsmath}
\usepackage{natbib}

\title{My Document}
\author{Your Name}
\date{\today}

\begin{document}

\maketitle

\section{Introduction}
This is a sample document demonstrating various \LaTeX{} features including mathematical equations, tables, and formatting. You can cite papers using the citation search feature, and they will be automatically added to your \texttt{cite.bib} file.

\section{Mathematical Expressions}
Here are some example equations:

\subsection{Quadratic Formula}
The quadratic formula is given by:
\[x = \frac{-b \pm \sqrt{b^2 - 4ac}}{2a}\]

\subsection{Matrix Example}
A sample matrix:
\[
\begin{pmatrix}
a_{11} & a_{12} & a_{13} \\
a_{21} & a_{22} & a_{23} \\
a_{31} & a_{32} & a_{33}
\end{pmatrix}
\]

\section{Tables}
Below is a sample table showing data:

\begin{table}[h]
\centering
\begin{tabular}{@{}lrrr@{}}
\toprule
\textbf{Category} & \textbf{Value 1} & \textbf{Value 2} & \textbf{Total} \\
\midrule
Item A & 10.5 & 20.3 & 30.8 \\
Item B & 15.2 & 25.1 & 40.3 \\
Item C & 22.1 & 12.7 & 34.8 \\
\bottomrule
\end{tabular}
\caption{Sample Data Table}
\label{tab:sample}
\end{table}

\section{Lists}
\subsection{Bullet Points}
\begin{itemize}
  \item First item
  \item Second item
  \item Third item with sub-items:
    \begin{itemize}
      \item Sub-item 1
      \item Sub-item 2
    \end{itemize}
\end{itemize}

\subsection{Numbered List}
\begin{enumerate}
  \item First step
  \item Second step
  \item Third step
\end{enumerate}

\section{Citations and References}
When you use the citation search feature, references will be automatically managed. For example, you might cite a paper like this: \cite{example2024}. All citations will be automatically added to your \texttt{cite.bib} file.

\section{Conclusion}
This template demonstrates basic \LaTeX{} features including mathematical formulas, tables, and lists. The bibliography system is now set up to work seamlessly with Bibby AI's citation management. Feel free to modify it for your needs!

\bibliographystyle{plain}
\bibliography{cite}

\end{document}\documentclass{article}
\usepackage[utf8]{inputenc}
\usepackage{booktabs}
\usepackage{amsmath}
\usepackage{natbib}

\title{My Document}
\author{Your Name}
\date{\today}

\begin{document}

\maketitle

\section{Introduction}
This is a sample document demonstrating various \LaTeX{} features including mathematical equations, tables, and formatting. You can cite papers using the citation search feature, and they will be automatically added to your \texttt{cite.bib} file.

\section{Mathematical Expressions}
Here are some example equations:

\subsection{Quadratic Formula}
The quadratic formula is given by:
\[x = \frac{-b \pm \sqrt{b^2 - 4ac}}{2a}\]

\subsection{Matrix Example}
A sample matrix:
\[
\begin{pmatrix}
a_{11} & a_{12} & a_{13} \\
a_{21} & a_{22} & a_{23} \\
a_{31} & a_{32} & a_{33}
\end{pmatrix}
\]

\section{Tables}
Below is a sample table showing data:

\begin{table}[h]
\centering
\begin{tabular}{@{}lrrr@{}}
\toprule
\textbf{Category} & \textbf{Value 1} & \textbf{Value 2} & \textbf{Total} \\
\midrule
Item A & 10.5 & 20.3 & 30.8 \\
Item B & 15.2 & 25.1 & 40.3 \\
Item C & 22.1 & 12.7 & 34.8 \\
\bottomrule
\end{tabular}
\caption{Sample Data Table}
\label{tab:sample}
\end{table}

\section{Lists}
\subsection{Bullet Points}
\begin{itemize}
  \item First item
  \item Second item
  \item Third item with sub-items:
    \begin{itemize}
      \item Sub-item 1
      \item Sub-item 2
    \end{itemize}
\end{itemize}

\subsection{Numbered List}
\begin{enumerate}
  \item First step
  \item Second step
  \item Third step
\end{enumerate}

\section{Citations and References}
When you use the citation search feature, references will be automatically managed. For example, you might cite a paper like this: \cite{example2024}. All citations will be automatically added to your \texttt{cite.bib} file.

\section{Conclusion}
This template demonstrates basic \LaTeX{} features including mathematical formulas, tables, and lists. The bibliography system is now set up to work seamlessly with Bibby AI's citation management. Feel free to modify it for your needs!

\bibliographystyle{plain}
\bibliography{cite}

\end{document}\documentclass{article}
\usepackage[utf8]{inputenc}
\usepackage{booktabs}
\usepackage{amsmath}
\usepackage{natbib}

\title{My Document}
\author{Your Name}
\date{\today}

\begin{document}

\maketitle

\section{Introduction}
This is a sample document demonstrating various \LaTeX{} features including mathematical equations, tables, and formatting. You can cite papers using the citation search feature, and they will be automatically added to your \texttt{cite.bib} file.

\section{Mathematical Expressions}
Here are some example equations:

\subsection{Quadratic Formula}
The quadratic formula is given by:
\[x = \frac{-b \pm \sqrt{b^2 - 4ac}}{2a}\]

\subsection{Matrix Example}
A sample matrix:
\[
\begin{pmatrix}
a_{11} & a_{12} & a_{13} \\
a_{21} & a_{22} & a_{23} \\
a_{31} & a_{32} & a_{33}
\end{pmatrix}
\]

\section{Tables}
Below is a sample table showing data:

\begin{table}[h]
\centering
\begin{tabular}{@{}lrrr@{}}
\toprule
\textbf{Category} & \textbf{Value 1} & \textbf{Value 2} & \textbf{Total} \\
\midrule
Item A & 10.5 & 20.3 & 30.8 \\
Item B & 15.2 & 25.1 & 40.3 \\
Item C & 22.1 & 12.7 & 34.8 \\
\bottomrule
\end{tabular}
\caption{Sample Data Table}
\label{tab:sample}
\end{table}

\section{Lists}
\subsection{Bullet Points}
\begin{itemize}
  \item First item
  \item Second item
  \item Third item with sub-items:
    \begin{itemize}
      \item Sub-item 1
      \item Sub-item 2
    \end{itemize}
\end{itemize}

\subsection{Numbered List}
\begin{enumerate}
  \item First step
  \item Second step
  \item Third step
\end{enumerate}

\section{Citations and References}
When you use the citation search feature, references will be automatically managed. For example, you might cite a paper like this: \cite{example2024}. All citations will be automatically added to your \texttt{cite.bib} file.

\section{Conclusion}
This template demonstrates basic \LaTeX{} features including mathematical formulas, tables, and lists. The bibliography system is now set up to work seamlessly with Bibby AI's citation management. Feel free to modify it for your needs!

\bibliographystyle{plain}
\bibliography{cite}

\end{document}\documentclass{article}
\usepackage[utf8]{inputenc}
\usepackage{booktabs}
\usepackage{amsmath}
\usepackage{natbib}

\title{My Document}
\author{Your Name}
\date{\today}

\begin{document}

\maketitle

\section{Introduction}
This is a sample document demonstrating various \LaTeX{} features including mathematical equations, tables, and formatting. You can cite papers using the citation search feature, and they will be automatically added to your \texttt{cite.bib} file.

\section{Mathematical Expressions}
Here are some example equations:

\subsection{Quadratic Formula}
The quadratic formula is given by:
\[x = \frac{-b \pm \sqrt{b^2 - 4ac}}{2a}\]

\subsection{Matrix Example}
A sample matrix:
\[
\begin{pmatrix}
a_{11} & a_{12} & a_{13} \\
a_{21} & a_{22} & a_{23} \\
a_{31} & a_{32} & a_{33}
\end{pmatrix}
\]

\section{Tables}
Below is a sample table showing data:

\begin{table}[h]
\centering
\begin{tabular}{@{}lrrr@{}}
\toprule
\textbf{Category} & \textbf{Value 1} & \textbf{Value 2} & \textbf{Total} \\
\midrule
Item A & 10.5 & 20.3 & 30.8 \\
Item B & 15.2 & 25.1 & 40.3 \\
Item C & 22.1 & 12.7 & 34.8 \\
\bottomrule
\end{tabular}
\caption{Sample Data Table}
\label{tab:sample}
\end{table}

\section{Lists}
\subsection{Bullet Points}
\begin{itemize}
  \item First item
  \item Second item
  \item Third item with sub-items:
    \begin{itemize}
      \item Sub-item 1
      \item Sub-item 2
    \end{itemize}
\end{itemize}

\subsection{Numbered List}
\begin{enumerate}
  \item First step
  \item Second step
  \item Third step
\end{enumerate}

\section{Citations and References}
When you use the citation search feature, references will be automatically managed. For example, you might cite a paper like this: \cite{example2024}. All citations will be automatically added to your \texttt{cite.bib} file.

\section{Conclusion}
This template demonstrates basic \LaTeX{} features including mathematical formulas, tables, and lists. The bibliography system is now set up to work seamlessly with Bibby AI's citation management. Feel free to modify it for your needs!

\bibliographystyle{plain}
\bibliography{cite}

\end{document}\documentclass{article}
\usepackage[utf8]{inputenc}
\usepackage{booktabs}
\usepackage{amsmath}
\usepackage{natbib}

\title{My Document}
\author{Your Name}
\date{\today}

\begin{document}

\maketitle

\section{Introduction}
This is a sample document demonstrating various \LaTeX{} features including mathematical equations, tables, and formatting. You can cite papers using the citation search feature, and they will be automatically added to your \texttt{cite.bib} file.

\section{Mathematical Expressions}
Here are some example equations:

\subsection{Quadratic Formula}
The quadratic formula is given by:
\[x = \frac{-b \pm \sqrt{b^2 - 4ac}}{2a}\]

\subsection{Matrix Example}
A sample matrix:
\[
\begin{pmatrix}
a_{11} & a_{12} & a_{13} \\
a_{21} & a_{22} & a_{23} \\
a_{31} & a_{32} & a_{33}
\end{pmatrix}
\]

\section{Tables}
Below is a sample table showing data:

\begin{table}[h]
\centering
\begin{tabular}{@{}lrrr@{}}
\toprule
\textbf{Category} & \textbf{Value 1} & \textbf{Value 2} & \textbf{Total} \\
\midrule
Item A & 10.5 & 20.3 & 30.8 \\
Item B & 15.2 & 25.1 & 40.3 \\
Item C & 22.1 & 12.7 & 34.8 \\
\bottomrule
\end{tabular}
\caption{Sample Data Table}
\label{tab:sample}
\end{table}

\section{Lists}
\subsection{Bullet Points}
\begin{itemize}
  \item First item
  \item Second item
  \item Third item with sub-items:
    \begin{itemize}
      \item Sub-item 1
      \item Sub-item 2
    \end{itemize}
\end{itemize}

\subsection{Numbered List}
\begin{enumerate}
  \item First step
  \item Second step
  \item Third step
\end{enumerate}

\section{Citations and References}
When you use the citation search feature, references will be automatically managed. For example, you might cite a paper like this: \cite{example2024}. All citations will be automatically added to your \texttt{cite.bib} file.

\section{Conclusion}
This template demonstrates basic \LaTeX{} features including mathematical formulas, tables, and lists. The bibliography system is now set up to work seamlessly with Bibby AI's citation management. Feel free to modify it for your needs!

\bibliographystyle{plain}
\bibliography{cite}

\end{document}\documentclass{article}
\usepackage[utf8]{inputenc}
\usepackage{booktabs}
\usepackage{amsmath}
\usepackage{natbib}

\title{My Document}
\author{Your Name}
\date{\today}

\begin{document}

\maketitle

\section{Introduction}
This is a sample document demonstrating various \LaTeX{} features including mathematical equations, tables, and formatting. You can cite papers using the citation search feature, and they will be automatically added to your \texttt{cite.bib} file.

\section{Mathematical Expressions}
Here are some example equations:

\subsection{Quadratic Formula}
The quadratic formula is given by:
\[x = \frac{-b \pm \sqrt{b^2 - 4ac}}{2a}\]

\subsection{Matrix Example}
A sample matrix:
\[
\begin{pmatrix}
a_{11} & a_{12} & a_{13} \\
a_{21} & a_{22} & a_{23} \\
a_{31} & a_{32} & a_{33}
\end{pmatrix}
\]

\section{Tables}
Below is a sample table showing data:

\begin{table}[h]
\centering
\begin{tabular}{@{}lrrr@{}}
\toprule
\textbf{Category} & \textbf{Value 1} & \textbf{Value 2} & \textbf{Total} \\
\midrule
Item A & 10.5 & 20.3 & 30.8 \\
Item B & 15.2 & 25.1 & 40.3 \\
Item C & 22.1 & 12.7 & 34.8 \\
\bottomrule
\end{tabular}
\caption{Sample Data Table}
\label{tab:sample}
\end{table}

\section{Lists}
\subsection{Bullet Points}
\begin{itemize}
  \item First item
  \item Second item
  \item Third item with sub-items:
    \begin{itemize}
      \item Sub-item 1
      \item Sub-item 2
    \end{itemize}
\end{itemize}

\subsection{Numbered List}
\begin{enumerate}
  \item First step
  \item Second step
  \item Third step
\end{enumerate}

\section{Citations and References}
When you use the citation search feature, references will be automatically managed. For example, you might cite a paper like this: \cite{example2024}. All citations will be automatically added to your \texttt{cite.bib} file.

\section{Conclusion}
This template demonstrates basic \LaTeX{} features including mathematical formulas, tables, and lists. The bibliography system is now set up to work seamlessly with Bibby AI's citation management. Feel free to modify it for your needs!

\bibliographystyle{plain}
\bibliography{cite}

\end{document}\documentclass{article}
\usepackage[utf8]{inputenc}
\usepackage{booktabs}
\usepackage{amsmath}
\usepackage{natbib}

\title{My Document}
\author{Your Name}
\date{\today}

\begin{document}

\maketitle

\section{Introduction}
This is a sample document demonstrating various \LaTeX{} features including mathematical equations, tables, and formatting. You can cite papers using the citation search feature, and they will be automatically added to your \texttt{cite.bib} file.

\section{Mathematical Expressions}
Here are some example equations:

\subsection{Quadratic Formula}
The quadratic formula is given by:
\[x = \frac{-b \pm \sqrt{b^2 - 4ac}}{2a}\]

\subsection{Matrix Example}
A sample matrix:
\[
\begin{pmatrix}
a_{11} & a_{12} & a_{13} \\
a_{21} & a_{22} & a_{23} \\
a_{31} & a_{32} & a_{33}
\end{pmatrix}
\]

\section{Tables}
Below is a sample table showing data:

\begin{table}[h]
\centering
\begin{tabular}{@{}lrrr@{}}
\toprule
\textbf{Category} & \textbf{Value 1} & \textbf{Value 2} & \textbf{Total} \\
\midrule
Item A & 10.5 & 20.3 & 30.8 \\
Item B & 15.2 & 25.1 & 40.3 \\
Item C & 22.1 & 12.7 & 34.8 \\
\bottomrule
\end{tabular}
\caption{Sample Data Table}
\label{tab:sample}
\end{table}

\section{Lists}
\subsection{Bullet Points}
\begin{itemize}
  \item First item
  \item Second item
  \item Third item with sub-items:
    \begin{itemize}
      \item Sub-item 1
      \item Sub-item 2
    \end{itemize}
\end{itemize}

\subsection{Numbered List}
\begin{enumerate}
  \item First step
  \item Second step
  \item Third step
\end{enumerate}

\section{Citations and References}
When you use the citation search feature, references will be automatically managed. For example, you might cite a paper like this: \cite{example2024}. All citations will be automatically added to your \texttt{cite.bib} file.

\section{Conclusion}
This template demonstrates basic \LaTeX{} features including mathematical formulas, tables, and lists. The bibliography system is now set up to work seamlessly with Bibby AI's citation management. Feel free to modify it for your needs!

\bibliographystyle{plain}
\bibliography{cite}

\end{document}\documentclass{article}
\usepackage[utf8]{inputenc}
\usepackage{booktabs}
\usepackage{amsmath}
\usepackage{natbib}

\title{My Document}
\author{Your Name}
\date{\today}

\begin{document}

\maketitle

\section{Introduction}
This is a sample document demonstrating various \LaTeX{} features including mathematical equations, tables, and formatting. You can cite papers using the citation search feature, and they will be automatically added to your \texttt{cite.bib} file.

\section{Mathematical Expressions}
Here are some example equations:

\subsection{Quadratic Formula}
The quadratic formula is given by:
\[x = \frac{-b \pm \sqrt{b^2 - 4ac}}{2a}\]

\subsection{Matrix Example}
A sample matrix:
\[
\begin{pmatrix}
a_{11} & a_{12} & a_{13} \\
a_{21} & a_{22} & a_{23} \\
a_{31} & a_{32} & a_{33}
\end{pmatrix}
\]

\section{Tables}
Below is a sample table showing data:

\begin{table}[h]
\centering
\begin{tabular}{@{}lrrr@{}}
\toprule
\textbf{Category} & \textbf{Value 1} & \textbf{Value 2} & \textbf{Total} \\
\midrule
Item A & 10.5 & 20.3 & 30.8 \\
Item B & 15.2 & 25.1 & 40.3 \\
Item C & 22.1 & 12.7 & 34.8 \\
\bottomrule
\end{tabular}
\caption{Sample Data Table}
\label{tab:sample}
\end{table}

\section{Lists}
\subsection{Bullet Points}
\begin{itemize}
  \item First item
  \item Second item
  \item Third item with sub-items:
    \begin{itemize}
      \item Sub-item 1
      \item Sub-item 2
    \end{itemize}
\end{itemize}

\subsection{Numbered List}
\begin{enumerate}
  \item First step
  \item Second step
  \item Third step
\end{enumerate}

\section{Citations and References}
When you use the citation search feature, references will be automatically managed. For example, you might cite a paper like this: \cite{example2024}. All citations will be automatically added to your \texttt{cite.bib} file.

\section{Conclusion}
This template demonstrates basic \LaTeX{} features including mathematical formulas, tables, and lists. The bibliography system is now set up to work seamlessly with Bibby AI's citation management. Feel free to modify it for your needs!

\bibliographystyle{plain}
\bibliography{cite}

\end{document}\documentclass{article}
\usepackage[utf8]{inputenc}
\usepackage{booktabs}
\usepackage{amsmath}
\usepackage{natbib}

\title{My Document}
\author{Your Name}
\date{\today}

\begin{document}

\maketitle

\section{Introduction}
This is a sample document demonstrating various \LaTeX{} features including mathematical equations, tables, and formatting. You can cite papers using the citation search feature, and they will be automatically added to your \texttt{cite.bib} file.

\section{Mathematical Expressions}
Here are some example equations:

\subsection{Quadratic Formula}
The quadratic formula is given by:
\[x = \frac{-b \pm \sqrt{b^2 - 4ac}}{2a}\]

\subsection{Matrix Example}
A sample matrix:
\[
\begin{pmatrix}
a_{11} & a_{12} & a_{13} \\
a_{21} & a_{22} & a_{23} \\
a_{31} & a_{32} & a_{33}
\end{pmatrix}
\]

\section{Tables}
Below is a sample table showing data:

\begin{table}[h]
\centering
\begin{tabular}{@{}lrrr@{}}
\toprule
\textbf{Category} & \textbf{Value 1} & \textbf{Value 2} & \textbf{Total} \\
\midrule
Item A & 10.5 & 20.3 & 30.8 \\
Item B & 15.2 & 25.1 & 40.3 \\
Item C & 22.1 & 12.7 & 34.8 \\
\bottomrule
\end{tabular}
\caption{Sample Data Table}
\label{tab:sample}
\end{table}

\section{Lists}
\subsection{Bullet Points}
\begin{itemize}
  \item First item
  \item Second item
  \item Third item with sub-items:
    \begin{itemize}
      \item Sub-item 1
      \item Sub-item 2
    \end{itemize}
\end{itemize}

\subsection{Numbered List}
\begin{enumerate}
  \item First step
  \item Second step
  \item Third step
\end{enumerate}

\section{Citations and References}
When you use the citation search feature, references will be automatically managed. For example, you might cite a paper like this: \cite{example2024}. All citations will be automatically added to your \texttt{cite.bib} file.

\section{Conclusion}
This template demonstrates basic \LaTeX{} features including mathematical formulas, tables, and lists. The bibliography system is now set up to work seamlessly with Bibby AI's citation management. Feel free to modify it for your needs!

\bibliographystyle{plain}
\bibliography{cite}

\end{document}\documentclass{article}
\usepackage[utf8]{inputenc}
\usepackage{booktabs}
\usepackage{amsmath}
\usepackage{natbib}

\title{My Document}
\author{Your Name}
\date{\today}

\begin{document}

\maketitle

\section{Introduction}
This is a sample document demonstrating various \LaTeX{} features including mathematical equations, tables, and formatting. You can cite papers using the citation search feature, and they will be automatically added to your \texttt{cite.bib} file.

\section{Mathematical Expressions}
Here are some example equations:

\subsection{Quadratic Formula}
The quadratic formula is given by:
\[x = \frac{-b \pm \sqrt{b^2 - 4ac}}{2a}\]

\subsection{Matrix Example}
A sample matrix:
\[
\begin{pmatrix}
a_{11} & a_{12} & a_{13} \\
a_{21} & a_{22} & a_{23} \\
a_{31} & a_{32} & a_{33}
\end{pmatrix}
\]

\section{Tables}
Below is a sample table showing data:

\begin{table}[h]
\centering
\begin{tabular}{@{}lrrr@{}}
\toprule
\textbf{Category} & \textbf{Value 1} & \textbf{Value 2} & \textbf{Total} \\
\midrule
Item A & 10.5 & 20.3 & 30.8 \\
Item B & 15.2 & 25.1 & 40.3 \\
Item C & 22.1 & 12.7 & 34.8 \\
\bottomrule
\end{tabular}
\caption{Sample Data Table}
\label{tab:sample}
\end{table}

\section{Lists}
\subsection{Bullet Points}
\begin{itemize}
  \item First item
  \item Second item
  \item Third item with sub-items:
    \begin{itemize}
      \item Sub-item 1
      \item Sub-item 2
    \end{itemize}
\end{itemize}

\subsection{Numbered List}
\begin{enumerate}
  \item First step
  \item Second step
  \item Third step
\end{enumerate}

\section{Citations and References}
When you use the citation search feature, references will be automatically managed. For example, you might cite a paper like this: \cite{example2024}. All citations will be automatically added to your \texttt{cite.bib} file.

\section{Conclusion}
This template demonstrates basic \LaTeX{} features including mathematical formulas, tables, and lists. The bibliography system is now set up to work seamlessly with Bibby AI's citation management. Feel free to modify it for your needs!

\bibliographystyle{plain}
\bibliography{cite}

\end{document}\documentclass{article}
\usepackage[utf8]{inputenc}
\usepackage{booktabs}
\usepackage{amsmath}
\usepackage{natbib}

\title{My Document}
\author{Your Name}
\date{\today}

\begin{document}

\maketitle

\section{Introduction}
This is a sample document demonstrating various \LaTeX{} features including mathematical equations, tables, and formatting. You can cite papers using the citation search feature, and they will be automatically added to your \texttt{cite.bib} file.

\section{Mathematical Expressions}
Here are some example equations:

\subsection{Quadratic Formula}
The quadratic formula is given by:
\[x = \frac{-b \pm \sqrt{b^2 - 4ac}}{2a}\]

\subsection{Matrix Example}
A sample matrix:
\[
\begin{pmatrix}
a_{11} & a_{12} & a_{13} \\
a_{21} & a_{22} & a_{23} \\
a_{31} & a_{32} & a_{33}
\end{pmatrix}
\]

\section{Tables}
Below is a sample table showing data:

\begin{table}[h]
\centering
\begin{tabular}{@{}lrrr@{}}
\toprule
\textbf{Category} & \textbf{Value 1} & \textbf{Value 2} & \textbf{Total} \\
\midrule
Item A & 10.5 & 20.3 & 30.8 \\
Item B & 15.2 & 25.1 & 40.3 \\
Item C & 22.1 & 12.7 & 34.8 \\
\bottomrule
\end{tabular}
\caption{Sample Data Table}
\label{tab:sample}
\end{table}

\section{Lists}
\subsection{Bullet Points}
\begin{itemize}
  \item First item
  \item Second item
  \item Third item with sub-items:
    \begin{itemize}
      \item Sub-item 1
      \item Sub-item 2
    \end{itemize}
\end{itemize}

\subsection{Numbered List}
\begin{enumerate}
  \item First step
  \item Second step
  \item Third step
\end{enumerate}

\section{Citations and References}
When you use the citation search feature, references will be automatically managed. For example, you might cite a paper like this: \cite{example2024}. All citations will be automatically added to your \texttt{cite.bib} file.

\section{Conclusion}
This template demonstrates basic \LaTeX{} features including mathematical formulas, tables, and lists. The bibliography system is now set up to work seamlessly with Bibby AI's citation management. Feel free to modify it for your needs!

\bibliographystyle{plain}
\bibliography{cite}

\end{document}\documentclass{article}
\usepackage[utf8]{inputenc}
\usepackage{booktabs}
\usepackage{amsmath}
\usepackage{natbib}

\title{My Document}
\author{Your Name}
\date{\today}

\begin{document}

\maketitle

\section{Introduction}
This is a sample document demonstrating various \LaTeX{} features including mathematical equations, tables, and formatting. You can cite papers using the citation search feature, and they will be automatically added to your \texttt{cite.bib} file.

\section{Mathematical Expressions}
Here are some example equations:

\subsection{Quadratic Formula}
The quadratic formula is given by:
\[x = \frac{-b \pm \sqrt{b^2 - 4ac}}{2a}\]

\subsection{Matrix Example}
A sample matrix:
\[
\begin{pmatrix}
a_{11} & a_{12} & a_{13} \\
a_{21} & a_{22} & a_{23} \\
a_{31} & a_{32} & a_{33}
\end{pmatrix}
\]

\section{Tables}
Below is a sample table showing data:

\begin{table}[h]
\centering
\begin{tabular}{@{}lrrr@{}}
\toprule
\textbf{Category} & \textbf{Value 1} & \textbf{Value 2} & \textbf{Total} \\
\midrule
Item A & 10.5 & 20.3 & 30.8 \\
Item B & 15.2 & 25.1 & 40.3 \\
Item C & 22.1 & 12.7 & 34.8 \\
\bottomrule
\end{tabular}
\caption{Sample Data Table}
\label{tab:sample}
\end{table}

\section{Lists}
\subsection{Bullet Points}
\begin{itemize}
  \item First item
  \item Second item
  \item Third item with sub-items:
    \begin{itemize}
      \item Sub-item 1
      \item Sub-item 2
    \end{itemize}
\end{itemize}

\subsection{Numbered List}
\begin{enumerate}
  \item First step
  \item Second step
  \item Third step
\end{enumerate}

\section{Citations and References}
When you use the citation search feature, references will be automatically managed. For example, you might cite a paper like this: \cite{example2024}. All citations will be automatically added to your \texttt{cite.bib} file.

\section{Conclusion}
This template demonstrates basic \LaTeX{} features including mathematical formulas, tables, and lists. The bibliography system is now set up to work seamlessly with Bibby AI's citation management. Feel free to modify it for your needs!

\bibliographystyle{plain}
\bibliography{cite}

\end{document}\documentclass{article}
\usepackage[utf8]{inputenc}
\usepackage{booktabs}
\usepackage{amsmath}
\usepackage{natbib}

\title{My Document}
\author{Your Name}
\date{\today}

\begin{document}

\maketitle

\section{Introduction}
This is a sample document demonstrating various \LaTeX{} features including mathematical equations, tables, and formatting. You can cite papers using the citation search feature, and they will be automatically added to your \texttt{cite.bib} file.

\section{Mathematical Expressions}
Here are some example equations:

\subsection{Quadratic Formula}
The quadratic formula is given by:
\[x = \frac{-b \pm \sqrt{b^2 - 4ac}}{2a}\]

\subsection{Matrix Example}
A sample matrix:
\[
\begin{pmatrix}
a_{11} & a_{12} & a_{13} \\
a_{21} & a_{22} & a_{23} \\
a_{31} & a_{32} & a_{33}
\end{pmatrix}
\]

\section{Tables}
Below is a sample table showing data:

\begin{table}[h]
\centering
\begin{tabular}{@{}lrrr@{}}
\toprule
\textbf{Category} & \textbf{Value 1} & \textbf{Value 2} & \textbf{Total} \\
\midrule
Item A & 10.5 & 20.3 & 30.8 \\
Item B & 15.2 & 25.1 & 40.3 \\
Item C & 22.1 & 12.7 & 34.8 \\
\bottomrule
\end{tabular}
\caption{Sample Data Table}
\label{tab:sample}
\end{table}

\section{Lists}
\subsection{Bullet Points}
\begin{itemize}
  \item First item
  \item Second item
  \item Third item with sub-items:
    \begin{itemize}
      \item Sub-item 1
      \item Sub-item 2
    \end{itemize}
\end{itemize}

\subsection{Numbered List}
\begin{enumerate}
  \item First step
  \item Second step
  \item Third step
\end{enumerate}

\section{Citations and References}
When you use the citation search feature, references will be automatically managed. For example, you might cite a paper like this: \cite{example2024}. All citations will be automatically added to your \texttt{cite.bib} file.

\section{Conclusion}
This template demonstrates basic \LaTeX{} features including mathematical formulas, tables, and lists. The bibliography system is now set up to work seamlessly with Bibby AI's citation management. Feel free to modify it for your needs!

\bibliographystyle{plain}
\bibliography{cite}

\end{document}\documentclass{article}
\usepackage[utf8]{inputenc}
\usepackage{booktabs}
\usepackage{amsmath}
\usepackage{natbib}

\title{My Document}
\author{Your Name}
\date{\today}

\begin{document}

\maketitle

\section{Introduction}
This is a sample document demonstrating various \LaTeX{} features including mathematical equations, tables, and formatting. You can cite papers using the citation search feature, and they will be automatically added to your \texttt{cite.bib} file.

\section{Mathematical Expressions}
Here are some example equations:

\subsection{Quadratic Formula}
The quadratic formula is given by:
\[x = \frac{-b \pm \sqrt{b^2 - 4ac}}{2a}\]

\subsection{Matrix Example}
A sample matrix:
\[
\begin{pmatrix}
a_{11} & a_{12} & a_{13} \\
a_{21} & a_{22} & a_{23} \\
a_{31} & a_{32} & a_{33}
\end{pmatrix}
\]

\section{Tables}
Below is a sample table showing data:

\begin{table}[h]
\centering
\begin{tabular}{@{}lrrr@{}}
\toprule
\textbf{Category} & \textbf{Value 1} & \textbf{Value 2} & \textbf{Total} \\
\midrule
Item A & 10.5 & 20.3 & 30.8 \\
Item B & 15.2 & 25.1 & 40.3 \\
Item C & 22.1 & 12.7 & 34.8 \\
\bottomrule
\end{tabular}
\caption{Sample Data Table}
\label{tab:sample}
\end{table}

\section{Lists}
\subsection{Bullet Points}
\begin{itemize}
  \item First item
  \item Second item
  \item Third item with sub-items:
    \begin{itemize}
      \item Sub-item 1
      \item Sub-item 2
    \end{itemize}
\end{itemize}

\subsection{Numbered List}
\begin{enumerate}
  \item First step
  \item Second step
  \item Third step
\end{enumerate}

\section{Citations and References}
When you use the citation search feature, references will be automatically managed. For example, you might cite a paper like this: \cite{example2024}. All citations will be automatically added to your \texttt{cite.bib} file.

\section{Conclusion}
This template demonstrates basic \LaTeX{} features including mathematical formulas, tables, and lists. The bibliography system is now set up to work seamlessly with Bibby AI's citation management. Feel free to modify it for your needs!

\bibliographystyle{plain}
\bibliography{cite}

\end{document}\documentclass{article}
\usepackage[utf8]{inputenc}
\usepackage{booktabs}
\usepackage{amsmath}
\usepackage{natbib}

\title{My Document}
\author{Your Name}
\date{\today}

\begin{document}

\maketitle

\section{Introduction}
This is a sample document demonstrating various \LaTeX{} features including mathematical equations, tables, and formatting. You can cite papers using the citation search feature, and they will be automatically added to your \texttt{cite.bib} file.

\section{Mathematical Expressions}
Here are some example equations:

\subsection{Quadratic Formula}
The quadratic formula is given by:
\[x = \frac{-b \pm \sqrt{b^2 - 4ac}}{2a}\]

\subsection{Matrix Example}
A sample matrix:
\[
\begin{pmatrix}
a_{11} & a_{12} & a_{13} \\
a_{21} & a_{22} & a_{23} \\
a_{31} & a_{32} & a_{33}
\end{pmatrix}
\]

\section{Tables}
Below is a sample table showing data:

\begin{table}[h]
\centering
\begin{tabular}{@{}lrrr@{}}
\toprule
\textbf{Category} & \textbf{Value 1} & \textbf{Value 2} & \textbf{Total} \\
\midrule
Item A & 10.5 & 20.3 & 30.8 \\
Item B & 15.2 & 25.1 & 40.3 \\
Item C & 22.1 & 12.7 & 34.8 \\
\bottomrule
\end{tabular}
\caption{Sample Data Table}
\label{tab:sample}
\end{table}

\section{Lists}
\subsection{Bullet Points}
\begin{itemize}
  \item First item
  \item Second item
  \item Third item with sub-items:
    \begin{itemize}
      \item Sub-item 1
      \item Sub-item 2
    \end{itemize}
\end{itemize}

\subsection{Numbered List}
\begin{enumerate}
  \item First step
  \item Second step
  \item Third step
\end{enumerate}

\section{Citations and References}
When you use the citation search feature, references will be automatically managed. For example, you might cite a paper like this: \cite{example2024}. All citations will be automatically added to your \texttt{cite.bib} file.

\section{Conclusion}
This template demonstrates basic \LaTeX{} features including mathematical formulas, tables, and lists. The bibliography system is now set up to work seamlessly with Bibby AI's citation management. Feel free to modify it for your needs!

\bibliographystyle{plain}
\bibliography{cite}

\end{document}\documentclass{article}
\usepackage[utf8]{inputenc}
\usepackage{booktabs}
\usepackage{amsmath}
\usepackage{natbib}

\title{My Document}
\author{Your Name}
\date{\today}

\begin{document}

\maketitle

\section{Introduction}
This is a sample document demonstrating various \LaTeX{} features including mathematical equations, tables, and formatting. You can cite papers using the citation search feature, and they will be automatically added to your \texttt{cite.bib} file.

\section{Mathematical Expressions}
Here are some example equations:

\subsection{Quadratic Formula}
The quadratic formula is given by:
\[x = \frac{-b \pm \sqrt{b^2 - 4ac}}{2a}\]

\subsection{Matrix Example}
A sample matrix:
\[
\begin{pmatrix}
a_{11} & a_{12} & a_{13} \\
a_{21} & a_{22} & a_{23} \\
a_{31} & a_{32} & a_{33}
\end{pmatrix}
\]

\section{Tables}
Below is a sample table showing data:

\begin{table}[h]
\centering
\begin{tabular}{@{}lrrr@{}}
\toprule
\textbf{Category} & \textbf{Value 1} & \textbf{Value 2} & \textbf{Total} \\
\midrule
Item A & 10.5 & 20.3 & 30.8 \\
Item B & 15.2 & 25.1 & 40.3 \\
Item C & 22.1 & 12.7 & 34.8 \\
\bottomrule
\end{tabular}
\caption{Sample Data Table}
\label{tab:sample}
\end{table}

\section{Lists}
\subsection{Bullet Points}
\begin{itemize}
  \item First item
  \item Second item
  \item Third item with sub-items:
    \begin{itemize}
      \item Sub-item 1
      \item Sub-item 2
    \end{itemize}
\end{itemize}

\subsection{Numbered List}
\begin{enumerate}
  \item First step
  \item Second step
  \item Third step
\end{enumerate}

\section{Citations and References}
When you use the citation search feature, references will be automatically managed. For example, you might cite a paper like this: \cite{example2024}. All citations will be automatically added to your \texttt{cite.bib} file.

\section{Conclusion}
This template demonstrates basic \LaTeX{} features including mathematical formulas, tables, and lists. The bibliography system is now set up to work seamlessly with Bibby AI's citation management. Feel free to modify it for your needs!

\bibliographystyle{plain}
\bibliography{cite}

\end{document}\documentclass{article}
\usepackage[utf8]{inputenc}
\usepackage{booktabs}
\usepackage{amsmath}
\usepackage{natbib}

\title{My Document}
\author{Your Name}
\date{\today}

\begin{document}

\maketitle

\section{Introduction}
This is a sample document demonstrating various \LaTeX{} features including mathematical equations, tables, and formatting. You can cite papers using the citation search feature, and they will be automatically added to your \texttt{cite.bib} file.

\section{Mathematical Expressions}
Here are some example equations:

\subsection{Quadratic Formula}
The quadratic formula is given by:
\[x = \frac{-b \pm \sqrt{b^2 - 4ac}}{2a}\]

\subsection{Matrix Example}
A sample matrix:
\[
\begin{pmatrix}
a_{11} & a_{12} & a_{13} \\
a_{21} & a_{22} & a_{23} \\
a_{31} & a_{32} & a_{33}
\end{pmatrix}
\]

\section{Tables}
Below is a sample table showing data:

\begin{table}[h]
\centering
\begin{tabular}{@{}lrrr@{}}
\toprule
\textbf{Category} & \textbf{Value 1} & \textbf{Value 2} & \textbf{Total} \\
\midrule
Item A & 10.5 & 20.3 & 30.8 \\
Item B & 15.2 & 25.1 & 40.3 \\
Item C & 22.1 & 12.7 & 34.8 \\
\bottomrule
\end{tabular}
\caption{Sample Data Table}
\label{tab:sample}
\end{table}

\section{Lists}
\subsection{Bullet Points}
\begin{itemize}
  \item First item
  \item Second item
  \item Third item with sub-items:
    \begin{itemize}
      \item Sub-item 1
      \item Sub-item 2
    \end{itemize}
\end{itemize}

\subsection{Numbered List}
\begin{enumerate}
  \item First step
  \item Second step
  \item Third step
\end{enumerate}

\section{Citations and References}
When you use the citation search feature, references will be automatically managed. For example, you might cite a paper like this: \cite{example2024}. All citations will be automatically added to your \texttt{cite.bib} file.

\section{Conclusion}
This template demonstrates basic \LaTeX{} features including mathematical formulas, tables, and lists. The bibliography system is now set up to work seamlessly with Bibby AI's citation management. Feel free to modify it for your needs!

\bibliographystyle{plain}
\bibliography{cite}

\end{document}\documentclass{article}
\usepackage[utf8]{inputenc}
\usepackage{booktabs}
\usepackage{amsmath}
\usepackage{natbib}

\title{My Document}
\author{Your Name}
\date{\today}

\begin{document}

\maketitle

\section{Introduction}
This is a sample document demonstrating various \LaTeX{} features including mathematical equations, tables, and formatting. You can cite papers using the citation search feature, and they will be automatically added to your \texttt{cite.bib} file.

\section{Mathematical Expressions}
Here are some example equations:

\subsection{Quadratic Formula}
The quadratic formula is given by:
\[x = \frac{-b \pm \sqrt{b^2 - 4ac}}{2a}\]

\subsection{Matrix Example}
A sample matrix:
\[
\begin{pmatrix}
a_{11} & a_{12} & a_{13} \\
a_{21} & a_{22} & a_{23} \\
a_{31} & a_{32} & a_{33}
\end{pmatrix}
\]

\section{Tables}
Below is a sample table showing data:

\begin{table}[h]
\centering
\begin{tabular}{@{}lrrr@{}}
\toprule
\textbf{Category} & \textbf{Value 1} & \textbf{Value 2} & \textbf{Total} \\
\midrule
Item A & 10.5 & 20.3 & 30.8 \\
Item B & 15.2 & 25.1 & 40.3 \\
Item C & 22.1 & 12.7 & 34.8 \\
\bottomrule
\end{tabular}
\caption{Sample Data Table}
\label{tab:sample}
\end{table}

\section{Lists}
\subsection{Bullet Points}
\begin{itemize}
  \item First item
  \item Second item
  \item Third item with sub-items:
    \begin{itemize}
      \item Sub-item 1
      \item Sub-item 2
    \end{itemize}
\end{itemize}

\subsection{Numbered List}
\begin{enumerate}
  \item First step
  \item Second step
  \item Third step
\end{enumerate}

\section{Citations and References}
When you use the citation search feature, references will be automatically managed. For example, you might cite a paper like this: \cite{example2024}. All citations will be automatically added to your \texttt{cite.bib} file.

\section{Conclusion}
This template demonstrates basic \LaTeX{} features including mathematical formulas, tables, and lists. The bibliography system is now set up to work seamlessly with Bibby AI's citation management. Feel free to modify it for your needs!

\bibliographystyle{plain}
\bibliography{cite}

\end{document}\documentclass{article}
\usepackage[utf8]{inputenc}
\usepackage{booktabs}
\usepackage{amsmath}
\usepackage{natbib}

\title{My Document}
\author{Your Name}
\date{\today}

\begin{document}

\maketitle

\section{Introduction}
This is a sample document demonstrating various \LaTeX{} features including mathematical equations, tables, and formatting. You can cite papers using the citation search feature, and they will be automatically added to your \texttt{cite.bib} file.

\section{Mathematical Expressions}
Here are some example equations:

\subsection{Quadratic Formula}
The quadratic formula is given by:
\[x = \frac{-b \pm \sqrt{b^2 - 4ac}}{2a}\]

\subsection{Matrix Example}
A sample matrix:
\[
\begin{pmatrix}
a_{11} & a_{12} & a_{13} \\
a_{21} & a_{22} & a_{23} \\
a_{31} & a_{32} & a_{33}
\end{pmatrix}
\]

\section{Tables}
Below is a sample table showing data:

\begin{table}[h]
\centering
\begin{tabular}{@{}lrrr@{}}
\toprule
\textbf{Category} & \textbf{Value 1} & \textbf{Value 2} & \textbf{Total} \\
\midrule
Item A & 10.5 & 20.3 & 30.8 \\
Item B & 15.2 & 25.1 & 40.3 \\
Item C & 22.1 & 12.7 & 34.8 \\
\bottomrule
\end{tabular}
\caption{Sample Data Table}
\label{tab:sample}
\end{table}

\section{Lists}
\subsection{Bullet Points}
\begin{itemize}
  \item First item
  \item Second item
  \item Third item with sub-items:
    \begin{itemize}
      \item Sub-item 1
      \item Sub-item 2
    \end{itemize}
\end{itemize}

\subsection{Numbered List}
\begin{enumerate}
  \item First step
  \item Second step
  \item Third step
\end{enumerate}

\section{Citations and References}
When you use the citation search feature, references will be automatically managed. For example, you might cite a paper like this: \cite{example2024}. All citations will be automatically added to your \texttt{cite.bib} file.

\section{Conclusion}
This template demonstrates basic \LaTeX{} features including mathematical formulas, tables, and lists. The bibliography system is now set up to work seamlessly with Bibby AI's citation management. Feel free to modify it for your needs!

\bibliographystyle{plain}
\bibliography{cite}

\end{document}\documentclass{article}
\usepackage[utf8]{inputenc}
\usepackage{booktabs}
\usepackage{amsmath}
\usepackage{natbib}

\title{My Document}
\author{Your Name}
\date{\today}

\begin{document}

\maketitle

\section{Introduction}
This is a sample document demonstrating various \LaTeX{} features including mathematical equations, tables, and formatting. You can cite papers using the citation search feature, and they will be automatically added to your \texttt{cite.bib} file.

\section{Mathematical Expressions}
Here are some example equations:

\subsection{Quadratic Formula}
The quadratic formula is given by:
\[x = \frac{-b \pm \sqrt{b^2 - 4ac}}{2a}\]

\subsection{Matrix Example}
A sample matrix:
\[
\begin{pmatrix}
a_{11} & a_{12} & a_{13} \\
a_{21} & a_{22} & a_{23} \\
a_{31} & a_{32} & a_{33}
\end{pmatrix}
\]

\section{Tables}
Below is a sample table showing data:

\begin{table}[h]
\centering
\begin{tabular}{@{}lrrr@{}}
\toprule
\textbf{Category} & \textbf{Value 1} & \textbf{Value 2} & \textbf{Total} \\
\midrule
Item A & 10.5 & 20.3 & 30.8 \\
Item B & 15.2 & 25.1 & 40.3 \\
Item C & 22.1 & 12.7 & 34.8 \\
\bottomrule
\end{tabular}
\caption{Sample Data Table}
\label{tab:sample}
\end{table}

\section{Lists}
\subsection{Bullet Points}
\begin{itemize}
  \item First item
  \item Second item
  \item Third item with sub-items:
    \begin{itemize}
      \item Sub-item 1
      \item Sub-item 2
    \end{itemize}
\end{itemize}

\subsection{Numbered List}
\begin{enumerate}
  \item First step
  \item Second step
  \item Third step
\end{enumerate}

\section{Citations and References}
When you use the citation search feature, references will be automatically managed. For example, you might cite a paper like this: \cite{example2024}. All citations will be automatically added to your \texttt{cite.bib} file.

\section{Conclusion}
This template demonstrates basic \LaTeX{} features including mathematical formulas, tables, and lists. The bibliography system is now set up to work seamlessly with Bibby AI's citation management. Feel free to modify it for your needs!

\bibliographystyle{plain}
\bibliography{cite}

\end{document}\documentclass{article}
\usepackage[utf8]{inputenc}
\usepackage{booktabs}
\usepackage{amsmath}
\usepackage{natbib}

\title{My Document}
\author{Your Name}
\date{\today}

\begin{document}

\maketitle

\section{Introduction}
This is a sample document demonstrating various \LaTeX{} features including mathematical equations, tables, and formatting. You can cite papers using the citation search feature, and they will be automatically added to your \texttt{cite.bib} file.

\section{Mathematical Expressions}
Here are some example equations:

\subsection{Quadratic Formula}
The quadratic formula is given by:
\[x = \frac{-b \pm \sqrt{b^2 - 4ac}}{2a}\]

\subsection{Matrix Example}
A sample matrix:
\[
\begin{pmatrix}
a_{11} & a_{12} & a_{13} \\
a_{21} & a_{22} & a_{23} \\
a_{31} & a_{32} & a_{33}
\end{pmatrix}
\]

\section{Tables}
Below is a sample table showing data:

\begin{table}[h]
\centering
\begin{tabular}{@{}lrrr@{}}
\toprule
\textbf{Category} & \textbf{Value 1} & \textbf{Value 2} & \textbf{Total} \\
\midrule
Item A & 10.5 & 20.3 & 30.8 \\
Item B & 15.2 & 25.1 & 40.3 \\
Item C & 22.1 & 12.7 & 34.8 \\
\bottomrule
\end{tabular}
\caption{Sample Data Table}
\label{tab:sample}
\end{table}

\section{Lists}
\subsection{Bullet Points}
\begin{itemize}
  \item First item
  \item Second item
  \item Third item with sub-items:
    \begin{itemize}
      \item Sub-item 1
      \item Sub-item 2
    \end{itemize}
\end{itemize}

\subsection{Numbered List}
\begin{enumerate}
  \item First step
  \item Second step
  \item Third step
\end{enumerate}

\section{Citations and References}
When you use the citation search feature, references will be automatically managed. For example, you might cite a paper like this: \cite{example2024}. All citations will be automatically added to your \texttt{cite.bib} file.

\section{Conclusion}
This template demonstrates basic \LaTeX{} features including mathematical formulas, tables, and lists. The bibliography system is now set up to work seamlessly with Bibby AI's citation management. Feel free to modify it for your needs!

\bibliographystyle{plain}
\bibliography{cite}

\end{document}\documentclass{article}
\usepackage[utf8]{inputenc}
\usepackage{booktabs}
\usepackage{amsmath}
\usepackage{natbib}

\title{My Document}
\author{Your Name}
\date{\today}

\begin{document}

\maketitle

\section{Introduction}
This is a sample document demonstrating various \LaTeX{} features including mathematical equations, tables, and formatting. You can cite papers using the citation search feature, and they will be automatically added to your \texttt{cite.bib} file.

\section{Mathematical Expressions}
Here are some example equations:

\subsection{Quadratic Formula}
The quadratic formula is given by:
\[x = \frac{-b \pm \sqrt{b^2 - 4ac}}{2a}\]

\subsection{Matrix Example}
A sample matrix:
\[
\begin{pmatrix}
a_{11} & a_{12} & a_{13} \\
a_{21} & a_{22} & a_{23} \\
a_{31} & a_{32} & a_{33}
\end{pmatrix}
\]

\section{Tables}
Below is a sample table showing data:

\begin{table}[h]
\centering
\begin{tabular}{@{}lrrr@{}}
\toprule
\textbf{Category} & \textbf{Value 1} & \textbf{Value 2} & \textbf{Total} \\
\midrule
Item A & 10.5 & 20.3 & 30.8 \\
Item B & 15.2 & 25.1 & 40.3 \\
Item C & 22.1 & 12.7 & 34.8 \\
\bottomrule
\end{tabular}
\caption{Sample Data Table}
\label{tab:sample}
\end{table}

\section{Lists}
\subsection{Bullet Points}
\begin{itemize}
  \item First item
  \item Second item
  \item Third item with sub-items:
    \begin{itemize}
      \item Sub-item 1
      \item Sub-item 2
    \end{itemize}
\end{itemize}

\subsection{Numbered List}
\begin{enumerate}
  \item First step
  \item Second step
  \item Third step
\end{enumerate}

\section{Citations and References}
When you use the citation search feature, references will be automatically managed. For example, you might cite a paper like this: \cite{example2024}. All citations will be automatically added to your \texttt{cite.bib} file.

\section{Conclusion}
This template demonstrates basic \LaTeX{} features including mathematical formulas, tables, and lists. The bibliography system is now set up to work seamlessly with Bibby AI's citation management. Feel free to modify it for your needs!

\bibliographystyle{plain}
\bibliography{cite}

\end{document}\documentclass{article}
\usepackage[utf8]{inputenc}
\usepackage{booktabs}
\usepackage{amsmath}
\usepackage{natbib}

\title{My Document}
\author{Your Name}
\date{\today}

\begin{document}

\maketitle

\section{Introduction}
This is a sample document demonstrating various \LaTeX{} features including mathematical equations, tables, and formatting. You can cite papers using the citation search feature, and they will be automatically added to your \texttt{cite.bib} file.

\section{Mathematical Expressions}
Here are some example equations:

\subsection{Quadratic Formula}
The quadratic formula is given by:
\[x = \frac{-b \pm \sqrt{b^2 - 4ac}}{2a}\]

\subsection{Matrix Example}
A sample matrix:
\[
\begin{pmatrix}
a_{11} & a_{12} & a_{13} \\
a_{21} & a_{22} & a_{23} \\
a_{31} & a_{32} & a_{33}
\end{pmatrix}
\]

\section{Tables}
Below is a sample table showing data:

\begin{table}[h]
\centering
\begin{tabular}{@{}lrrr@{}}
\toprule
\textbf{Category} & \textbf{Value 1} & \textbf{Value 2} & \textbf{Total} \\
\midrule
Item A & 10.5 & 20.3 & 30.8 \\
Item B & 15.2 & 25.1 & 40.3 \\
Item C & 22.1 & 12.7 & 34.8 \\
\bottomrule
\end{tabular}
\caption{Sample Data Table}
\label{tab:sample}
\end{table}

\section{Lists}
\subsection{Bullet Points}
\begin{itemize}
  \item First item
  \item Second item
  \item Third item with sub-items:
    \begin{itemize}
      \item Sub-item 1
      \item Sub-item 2
    \end{itemize}
\end{itemize}

\subsection{Numbered List}
\begin{enumerate}
  \item First step
  \item Second step
  \item Third step
\end{enumerate}

\section{Citations and References}
When you use the citation search feature, references will be automatically managed. For example, you might cite a paper like this: \cite{example2024}. All citations will be automatically added to your \texttt{cite.bib} file.

\section{Conclusion}
This template demonstrates basic \LaTeX{} features including mathematical formulas, tables, and lists. The bibliography system is now set up to work seamlessly with Bibby AI's citation management. Feel free to modify it for your needs!

\bibliographystyle{plain}
\bibliography{cite}

\end{document}\documentclass{article}
\usepackage[utf8]{inputenc}
\usepackage{booktabs}
\usepackage{amsmath}
\usepackage{natbib}

\title{My Document}
\author{Your Name}
\date{\today}

\begin{document}

\maketitle

\section{Introduction}
This is a sample document demonstrating various \LaTeX{} features including mathematical equations, tables, and formatting. You can cite papers using the citation search feature, and they will be automatically added to your \texttt{cite.bib} file.

\section{Mathematical Expressions}
Here are some example equations:

\subsection{Quadratic Formula}
The quadratic formula is given by:
\[x = \frac{-b \pm \sqrt{b^2 - 4ac}}{2a}\]

\subsection{Matrix Example}
A sample matrix:
\[
\begin{pmatrix}
a_{11} & a_{12} & a_{13} \\
a_{21} & a_{22} & a_{23} \\
a_{31} & a_{32} & a_{33}
\end{pmatrix}
\]

\section{Tables}
Below is a sample table showing data:

\begin{table}[h]
\centering
\begin{tabular}{@{}lrrr@{}}
\toprule
\textbf{Category} & \textbf{Value 1} & \textbf{Value 2} & \textbf{Total} \\
\midrule
Item A & 10.5 & 20.3 & 30.8 \\
Item B & 15.2 & 25.1 & 40.3 \\
Item C & 22.1 & 12.7 & 34.8 \\
\bottomrule
\end{tabular}
\caption{Sample Data Table}
\label{tab:sample}
\end{table}

\section{Lists}
\subsection{Bullet Points}
\begin{itemize}
  \item First item
  \item Second item
  \item Third item with sub-items:
    \begin{itemize}
      \item Sub-item 1
      \item Sub-item 2
    \end{itemize}
\end{itemize}

\subsection{Numbered List}
\begin{enumerate}
  \item First step
  \item Second step
  \item Third step
\end{enumerate}

\section{Citations and References}
When you use the citation search feature, references will be automatically managed. For example, you might cite a paper like this: \cite{example2024}. All citations will be automatically added to your \texttt{cite.bib} file.

\section{Conclusion}
This template demonstrates basic \LaTeX{} features including mathematical formulas, tables, and lists. The bibliography system is now set up to work seamlessly with Bibby AI's citation management. Feel free to modify it for your needs!

\bibliographystyle{plain}
\bibliography{cite}

\end{document}\documentclass{article}
\usepackage[utf8]{inputenc}
\usepackage{booktabs}
\usepackage{amsmath}
\usepackage{natbib}

\title{My Document}
\author{Your Name}
\date{\today}

\begin{document}

\maketitle

\section{Introduction}
This is a sample document demonstrating various \LaTeX{} features including mathematical equations, tables, and formatting. You can cite papers using the citation search feature, and they will be automatically added to your \texttt{cite.bib} file.

\section{Mathematical Expressions}
Here are some example equations:

\subsection{Quadratic Formula}
The quadratic formula is given by:
\[x = \frac{-b \pm \sqrt{b^2 - 4ac}}{2a}\]

\subsection{Matrix Example}
A sample matrix:
\[
\begin{pmatrix}
a_{11} & a_{12} & a_{13} \\
a_{21} & a_{22} & a_{23} \\
a_{31} & a_{32} & a_{33}
\end{pmatrix}
\]

\section{Tables}
Below is a sample table showing data:

\begin{table}[h]
\centering
\begin{tabular}{@{}lrrr@{}}
\toprule
\textbf{Category} & \textbf{Value 1} & \textbf{Value 2} & \textbf{Total} \\
\midrule
Item A & 10.5 & 20.3 & 30.8 \\
Item B & 15.2 & 25.1 & 40.3 \\
Item C & 22.1 & 12.7 & 34.8 \\
\bottomrule
\end{tabular}
\caption{Sample Data Table}
\label{tab:sample}
\end{table}

\section{Lists}
\subsection{Bullet Points}
\begin{itemize}
  \item First item
  \item Second item
  \item Third item with sub-items:
    \begin{itemize}
      \item Sub-item 1
      \item Sub-item 2
    \end{itemize}
\end{itemize}

\subsection{Numbered List}
\begin{enumerate}
  \item First step
  \item Second step
  \item Third step
\end{enumerate}

\section{Citations and References}
When you use the citation search feature, references will be automatically managed. For example, you might cite a paper like this: \cite{example2024}. All citations will be automatically added to your \texttt{cite.bib} file.

\section{Conclusion}
This template demonstrates basic \LaTeX{} features including mathematical formulas, tables, and lists. The bibliography system is now set up to work seamlessly with Bibby AI's citation management. Feel free to modify it for your needs!

\bibliographystyle{plain}
\bibliography{cite}

\end{document}\documentclass{article}
\usepackage[utf8]{inputenc}
\usepackage{booktabs}
\usepackage{amsmath}
\usepackage{natbib}

\title{My Document}
\author{Your Name}
\date{\today}

\begin{document}

\maketitle

\section{Introduction}
This is a sample document demonstrating various \LaTeX{} features including mathematical equations, tables, and formatting. You can cite papers using the citation search feature, and they will be automatically added to your \texttt{cite.bib} file.

\section{Mathematical Expressions}
Here are some example equations:

\subsection{Quadratic Formula}
The quadratic formula is given by:
\[x = \frac{-b \pm \sqrt{b^2 - 4ac}}{2a}\]

\subsection{Matrix Example}
A sample matrix:
\[
\begin{pmatrix}
a_{11} & a_{12} & a_{13} \\
a_{21} & a_{22} & a_{23} \\
a_{31} & a_{32} & a_{33}
\end{pmatrix}
\]

\section{Tables}
Below is a sample table showing data:

\begin{table}[h]
\centering
\begin{tabular}{@{}lrrr@{}}
\toprule
\textbf{Category} & \textbf{Value 1} & \textbf{Value 2} & \textbf{Total} \\
\midrule
Item A & 10.5 & 20.3 & 30.8 \\
Item B & 15.2 & 25.1 & 40.3 \\
Item C & 22.1 & 12.7 & 34.8 \\
\bottomrule
\end{tabular}
\caption{Sample Data Table}
\label{tab:sample}
\end{table}

\section{Lists}
\subsection{Bullet Points}
\begin{itemize}
  \item First item
  \item Second item
  \item Third item with sub-items:
    \begin{itemize}
      \item Sub-item 1
      \item Sub-item 2
    \end{itemize}
\end{itemize}

\subsection{Numbered List}
\begin{enumerate}
  \item First step
  \item Second step
  \item Third step
\end{enumerate}

\section{Citations and References}
When you use the citation search feature, references will be automatically managed. For example, you might cite a paper like this: \cite{example2024}. All citations will be automatically added to your \texttt{cite.bib} file.

\section{Conclusion}
This template demonstrates basic \LaTeX{} features including mathematical formulas, tables, and lists. The bibliography system is now set up to work seamlessly with Bibby AI's citation management. Feel free to modify it for your needs!

\bibliographystyle{plain}
\bibliography{cite}

\end{document}\documentclass{article}
\usepackage[utf8]{inputenc}
\usepackage{booktabs}
\usepackage{amsmath}
\usepackage{natbib}

\title{My Document}
\author{Your Name}
\date{\today}

\begin{document}

\maketitle

\section{Introduction}
This is a sample document demonstrating various \LaTeX{} features including mathematical equations, tables, and formatting. You can cite papers using the citation search feature, and they will be automatically added to your \texttt{cite.bib} file.

\section{Mathematical Expressions}
Here are some example equations:

\subsection{Quadratic Formula}
The quadratic formula is given by:
\[x = \frac{-b \pm \sqrt{b^2 - 4ac}}{2a}\]

\subsection{Matrix Example}
A sample matrix:
\[
\begin{pmatrix}
a_{11} & a_{12} & a_{13} \\
a_{21} & a_{22} & a_{23} \\
a_{31} & a_{32} & a_{33}
\end{pmatrix}
\]

\section{Tables}
Below is a sample table showing data:

\begin{table}[h]
\centering
\begin{tabular}{@{}lrrr@{}}
\toprule
\textbf{Category} & \textbf{Value 1} & \textbf{Value 2} & \textbf{Total} \\
\midrule
Item A & 10.5 & 20.3 & 30.8 \\
Item B & 15.2 & 25.1 & 40.3 \\
Item C & 22.1 & 12.7 & 34.8 \\
\bottomrule
\end{tabular}
\caption{Sample Data Table}
\label{tab:sample}
\end{table}

\section{Lists}
\subsection{Bullet Points}
\begin{itemize}
  \item First item
  \item Second item
  \item Third item with sub-items:
    \begin{itemize}
      \item Sub-item 1
      \item Sub-item 2
    \end{itemize}
\end{itemize}

\subsection{Numbered List}
\begin{enumerate}
  \item First step
  \item Second step
  \item Third step
\end{enumerate}

\section{Citations and References}
When you use the citation search feature, references will be automatically managed. For example, you might cite a paper like this: \cite{example2024}. All citations will be automatically added to your \texttt{cite.bib} file.

\section{Conclusion}
This template demonstrates basic \LaTeX{} features including mathematical formulas, tables, and lists. The bibliography system is now set up to work seamlessly with Bibby AI's citation management. Feel free to modify it for your needs!

\bibliographystyle{plain}
\bibliography{cite}

\end{document}\documentclass{article}
\usepackage[utf8]{inputenc}
\usepackage{booktabs}
\usepackage{amsmath}
\usepackage{natbib}

\title{My Document}
\author{Your Name}
\date{\today}

\begin{document}

\maketitle

\section{Introduction}
This is a sample document demonstrating various \LaTeX{} features including mathematical equations, tables, and formatting. You can cite papers using the citation search feature, and they will be automatically added to your \texttt{cite.bib} file.

\section{Mathematical Expressions}
Here are some example equations:

\subsection{Quadratic Formula}
The quadratic formula is given by:
\[x = \frac{-b \pm \sqrt{b^2 - 4ac}}{2a}\]

\subsection{Matrix Example}
A sample matrix:
\[
\begin{pmatrix}
a_{11} & a_{12} & a_{13} \\
a_{21} & a_{22} & a_{23} \\
a_{31} & a_{32} & a_{33}
\end{pmatrix}
\]

\section{Tables}
Below is a sample table showing data:

\begin{table}[h]
\centering
\begin{tabular}{@{}lrrr@{}}
\toprule
\textbf{Category} & \textbf{Value 1} & \textbf{Value 2} & \textbf{Total} \\
\midrule
Item A & 10.5 & 20.3 & 30.8 \\
Item B & 15.2 & 25.1 & 40.3 \\
Item C & 22.1 & 12.7 & 34.8 \\
\bottomrule
\end{tabular}
\caption{Sample Data Table}
\label{tab:sample}
\end{table}

\section{Lists}
\subsection{Bullet Points}
\begin{itemize}
  \item First item
  \item Second item
  \item Third item with sub-items:
    \begin{itemize}
      \item Sub-item 1
      \item Sub-item 2
    \end{itemize}
\end{itemize}

\subsection{Numbered List}
\begin{enumerate}
  \item First step
  \item Second step
  \item Third step
\end{enumerate}

\section{Citations and References}
When you use the citation search feature, references will be automatically managed. For example, you might cite a paper like this: \cite{example2024}. All citations will be automatically added to your \texttt{cite.bib} file.

\section{Conclusion}
This template demonstrates basic \LaTeX{} features including mathematical formulas, tables, and lists. The bibliography system is now set up to work seamlessly with Bibby AI's citation management. Feel free to modify it for your needs!

\bibliographystyle{plain}
\bibliography{cite}

\end{document}\documentclass{article}
\usepackage[utf8]{inputenc}
\usepackage{booktabs}
\usepackage{amsmath}
\usepackage{natbib}

\title{My Document}
\author{Your Name}
\date{\today}

\begin{document}

\maketitle

\section{Introduction}
This is a sample document demonstrating various \LaTeX{} features including mathematical equations, tables, and formatting. You can cite papers using the citation search feature, and they will be automatically added to your \texttt{cite.bib} file.

\section{Mathematical Expressions}
Here are some example equations:

\subsection{Quadratic Formula}
The quadratic formula is given by:
\[x = \frac{-b \pm \sqrt{b^2 - 4ac}}{2a}\]

\subsection{Matrix Example}
A sample matrix:
\[
\begin{pmatrix}
a_{11} & a_{12} & a_{13} \\
a_{21} & a_{22} & a_{23} \\
a_{31} & a_{32} & a_{33}
\end{pmatrix}
\]

\section{Tables}
Below is a sample table showing data:

\begin{table}[h]
\centering
\begin{tabular}{@{}lrrr@{}}
\toprule
\textbf{Category} & \textbf{Value 1} & \textbf{Value 2} & \textbf{Total} \\
\midrule
Item A & 10.5 & 20.3 & 30.8 \\
Item B & 15.2 & 25.1 & 40.3 \\
Item C & 22.1 & 12.7 & 34.8 \\
\bottomrule
\end{tabular}
\caption{Sample Data Table}
\label{tab:sample}
\end{table}

\section{Lists}
\subsection{Bullet Points}
\begin{itemize}
  \item First item
  \item Second item
  \item Third item with sub-items:
    \begin{itemize}
      \item Sub-item 1
      \item Sub-item 2
    \end{itemize}
\end{itemize}

\subsection{Numbered List}
\begin{enumerate}
  \item First step
  \item Second step
  \item Third step
\end{enumerate}

\section{Citations and References}
When you use the citation search feature, references will be automatically managed. For example, you might cite a paper like this: \cite{example2024}. All citations will be automatically added to your \texttt{cite.bib} file.

\section{Conclusion}
This template demonstrates basic \LaTeX{} features including mathematical formulas, tables, and lists. The bibliography system is now set up to work seamlessly with Bibby AI's citation management. Feel free to modify it for your needs!

\bibliographystyle{plain}
\bibliography{cite}

\end{document}\documentclass{article}
\usepackage[utf8]{inputenc}
\usepackage{booktabs}
\usepackage{amsmath}
\usepackage{natbib}

\title{My Document}
\author{Your Name}
\date{\today}

\begin{document}

\maketitle

\section{Introduction}
This is a sample document demonstrating various \LaTeX{} features including mathematical equations, tables, and formatting. You can cite papers using the citation search feature, and they will be automatically added to your \texttt{cite.bib} file.

\section{Mathematical Expressions}
Here are some example equations:

\subsection{Quadratic Formula}
The quadratic formula is given by:
\[x = \frac{-b \pm \sqrt{b^2 - 4ac}}{2a}\]

\subsection{Matrix Example}
A sample matrix:
\[
\begin{pmatrix}
a_{11} & a_{12} & a_{13} \\
a_{21} & a_{22} & a_{23} \\
a_{31} & a_{32} & a_{33}
\end{pmatrix}
\]

\section{Tables}
Below is a sample table showing data:

\begin{table}[h]
\centering
\begin{tabular}{@{}lrrr@{}}
\toprule
\textbf{Category} & \textbf{Value 1} & \textbf{Value 2} & \textbf{Total} \\
\midrule
Item A & 10.5 & 20.3 & 30.8 \\
Item B & 15.2 & 25.1 & 40.3 \\
Item C & 22.1 & 12.7 & 34.8 \\
\bottomrule
\end{tabular}
\caption{Sample Data Table}
\label{tab:sample}
\end{table}

\section{Lists}
\subsection{Bullet Points}
\begin{itemize}
  \item First item
  \item Second item
  \item Third item with sub-items:
    \begin{itemize}
      \item Sub-item 1
      \item Sub-item 2
    \end{itemize}
\end{itemize}

\subsection{Numbered List}
\begin{enumerate}
  \item First step
  \item Second step
  \item Third step
\end{enumerate}

\section{Citations and References}
When you use the citation search feature, references will be automatically managed. For example, you might cite a paper like this: \cite{example2024}. All citations will be automatically added to your \texttt{cite.bib} file.

\section{Conclusion}
This template demonstrates basic \LaTeX{} features including mathematical formulas, tables, and lists. The bibliography system is now set up to work seamlessly with Bibby AI's citation management. Feel free to modify it for your needs!

\bibliographystyle{plain}
\bibliography{cite}

\end{document}\documentclass{article}
\usepackage[utf8]{inputenc}
\usepackage{booktabs}
\usepackage{amsmath}
\usepackage{natbib}

\title{My Document}
\author{Your Name}
\date{\today}

\begin{document}

\maketitle

\section{Introduction}
This is a sample document demonstrating various \LaTeX{} features including mathematical equations, tables, and formatting. You can cite papers using the citation search feature, and they will be automatically added to your \texttt{cite.bib} file.

\section{Mathematical Expressions}
Here are some example equations:

\subsection{Quadratic Formula}
The quadratic formula is given by:
\[x = \frac{-b \pm \sqrt{b^2 - 4ac}}{2a}\]

\subsection{Matrix Example}
A sample matrix:
\[
\begin{pmatrix}
a_{11} & a_{12} & a_{13} \\
a_{21} & a_{22} & a_{23} \\
a_{31} & a_{32} & a_{33}
\end{pmatrix}
\]

\section{Tables}
Below is a sample table showing data:

\begin{table}[h]
\centering
\begin{tabular}{@{}lrrr@{}}
\toprule
\textbf{Category} & \textbf{Value 1} & \textbf{Value 2} & \textbf{Total} \\
\midrule
Item A & 10.5 & 20.3 & 30.8 \\
Item B & 15.2 & 25.1 & 40.3 \\
Item C & 22.1 & 12.7 & 34.8 \\
\bottomrule
\end{tabular}
\caption{Sample Data Table}
\label{tab:sample}
\end{table}

\section{Lists}
\subsection{Bullet Points}
\begin{itemize}
  \item First item
  \item Second item
  \item Third item with sub-items:
    \begin{itemize}
      \item Sub-item 1
      \item Sub-item 2
    \end{itemize}
\end{itemize}

\subsection{Numbered List}
\begin{enumerate}
  \item First step
  \item Second step
  \item Third step
\end{enumerate}

\section{Citations and References}
When you use the citation search feature, references will be automatically managed. For example, you might cite a paper like this: \cite{example2024}. All citations will be automatically added to your \texttt{cite.bib} file.

\section{Conclusion}
This template demonstrates basic \LaTeX{} features including mathematical formulas, tables, and lists. The bibliography system is now set up to work seamlessly with Bibby AI's citation management. Feel free to modify it for your needs!

\bibliographystyle{plain}
\bibliography{cite}

\end{document}\documentclass{article}
\usepackage[utf8]{inputenc}
\usepackage{booktabs}
\usepackage{amsmath}
\usepackage{natbib}

\title{My Document}
\author{Your Name}
\date{\today}

\begin{document}

\maketitle

\section{Introduction}
This is a sample document demonstrating various \LaTeX{} features including mathematical equations, tables, and formatting. You can cite papers using the citation search feature, and they will be automatically added to your \texttt{cite.bib} file.

\section{Mathematical Expressions}
Here are some example equations:

\subsection{Quadratic Formula}
The quadratic formula is given by:
\[x = \frac{-b \pm \sqrt{b^2 - 4ac}}{2a}\]

\subsection{Matrix Example}
A sample matrix:
\[
\begin{pmatrix}
a_{11} & a_{12} & a_{13} \\
a_{21} & a_{22} & a_{23} \\
a_{31} & a_{32} & a_{33}
\end{pmatrix}
\]

\section{Tables}
Below is a sample table showing data:

\begin{table}[h]
\centering
\begin{tabular}{@{}lrrr@{}}
\toprule
\textbf{Category} & \textbf{Value 1} & \textbf{Value 2} & \textbf{Total} \\
\midrule
Item A & 10.5 & 20.3 & 30.8 \\
Item B & 15.2 & 25.1 & 40.3 \\
Item C & 22.1 & 12.7 & 34.8 \\
\bottomrule
\end{tabular}
\caption{Sample Data Table}
\label{tab:sample}
\end{table}

\section{Lists}
\subsection{Bullet Points}
\begin{itemize}
  \item First item
  \item Second item
  \item Third item with sub-items:
    \begin{itemize}
      \item Sub-item 1
      \item Sub-item 2
    \end{itemize}
\end{itemize}

\subsection{Numbered List}
\begin{enumerate}
  \item First step
  \item Second step
  \item Third step
\end{enumerate}

\section{Citations and References}
When you use the citation search feature, references will be automatically managed. For example, you might cite a paper like this: \cite{example2024}. All citations will be automatically added to your \texttt{cite.bib} file.

\section{Conclusion}
This template demonstrates basic \LaTeX{} features including mathematical formulas, tables, and lists. The bibliography system is now set up to work seamlessly with Bibby AI's citation management. Feel free to modify it for your needs!

\bibliographystyle{plain}
\bibliography{cite}

\end{document}\documentclass{article}
\usepackage[utf8]{inputenc}
\usepackage{booktabs}
\usepackage{amsmath}
\usepackage{natbib}

\title{My Document}
\author{Your Name}
\date{\today}

\begin{document}

\maketitle

\section{Introduction}
This is a sample document demonstrating various \LaTeX{} features including mathematical equations, tables, and formatting. You can cite papers using the citation search feature, and they will be automatically added to your \texttt{cite.bib} file.

\section{Mathematical Expressions}
Here are some example equations:

\subsection{Quadratic Formula}
The quadratic formula is given by:
\[x = \frac{-b \pm \sqrt{b^2 - 4ac}}{2a}\]

\subsection{Matrix Example}
A sample matrix:
\[
\begin{pmatrix}
a_{11} & a_{12} & a_{13} \\
a_{21} & a_{22} & a_{23} \\
a_{31} & a_{32} & a_{33}
\end{pmatrix}
\]

\section{Tables}
Below is a sample table showing data:

\begin{table}[h]
\centering
\begin{tabular}{@{}lrrr@{}}
\toprule
\textbf{Category} & \textbf{Value 1} & \textbf{Value 2} & \textbf{Total} \\
\midrule
Item A & 10.5 & 20.3 & 30.8 \\
Item B & 15.2 & 25.1 & 40.3 \\
Item C & 22.1 & 12.7 & 34.8 \\
\bottomrule
\end{tabular}
\caption{Sample Data Table}
\label{tab:sample}
\end{table}

\section{Lists}
\subsection{Bullet Points}
\begin{itemize}
  \item First item
  \item Second item
  \item Third item with sub-items:
    \begin{itemize}
      \item Sub-item 1
      \item Sub-item 2
    \end{itemize}
\end{itemize}

\subsection{Numbered List}
\begin{enumerate}
  \item First step
  \item Second step
  \item Third step
\end{enumerate}

\section{Citations and References}
When you use the citation search feature, references will be automatically managed. For example, you might cite a paper like this: \cite{example2024}. All citations will be automatically added to your \texttt{cite.bib} file.

\section{Conclusion}
This template demonstrates basic \LaTeX{} features including mathematical formulas, tables, and lists. The bibliography system is now set up to work seamlessly with Bibby AI's citation management. Feel free to modify it for your needs!

\bibliographystyle{plain}
\bibliography{cite}

\end{document}\documentclass{article}
\usepackage[utf8]{inputenc}
\usepackage{booktabs}
\usepackage{amsmath}
\usepackage{natbib}

\title{My Document}
\author{Your Name}
\date{\today}

\begin{document}

\maketitle

\section{Introduction}
This is a sample document demonstrating various \LaTeX{} features including mathematical equations, tables, and formatting. You can cite papers using the citation search feature, and they will be automatically added to your \texttt{cite.bib} file.

\section{Mathematical Expressions}
Here are some example equations:

\subsection{Quadratic Formula}
The quadratic formula is given by:
\[x = \frac{-b \pm \sqrt{b^2 - 4ac}}{2a}\]

\subsection{Matrix Example}
A sample matrix:
\[
\begin{pmatrix}
a_{11} & a_{12} & a_{13} \\
a_{21} & a_{22} & a_{23} \\
a_{31} & a_{32} & a_{33}
\end{pmatrix}
\]

\section{Tables}
Below is a sample table showing data:

\begin{table}[h]
\centering
\begin{tabular}{@{}lrrr@{}}
\toprule
\textbf{Category} & \textbf{Value 1} & \textbf{Value 2} & \textbf{Total} \\
\midrule
Item A & 10.5 & 20.3 & 30.8 \\
Item B & 15.2 & 25.1 & 40.3 \\
Item C & 22.1 & 12.7 & 34.8 \\
\bottomrule
\end{tabular}
\caption{Sample Data Table}
\label{tab:sample}
\end{table}

\section{Lists}
\subsection{Bullet Points}
\begin{itemize}
  \item First item
  \item Second item
  \item Third item with sub-items:
    \begin{itemize}
      \item Sub-item 1
      \item Sub-item 2
    \end{itemize}
\end{itemize}

\subsection{Numbered List}
\begin{enumerate}
  \item First step
  \item Second step
  \item Third step
\end{enumerate}

\section{Citations and References}
When you use the citation search feature, references will be automatically managed. For example, you might cite a paper like this: \cite{example2024}. All citations will be automatically added to your \texttt{cite.bib} file.

\section{Conclusion}
This template demonstrates basic \LaTeX{} features including mathematical formulas, tables, and lists. The bibliography system is now set up to work seamlessly with Bibby AI's citation management. Feel free to modify it for your needs!

\bibliographystyle{plain}
\bibliography{cite}

\end{document}\documentclass{article}
\usepackage[utf8]{inputenc}
\usepackage{booktabs}
\usepackage{amsmath}
\usepackage{natbib}

\title{My Document}
\author{Your Name}
\date{\today}

\begin{document}

\maketitle

\section{Introduction}
This is a sample document demonstrating various \LaTeX{} features including mathematical equations, tables, and formatting. You can cite papers using the citation search feature, and they will be automatically added to your \texttt{cite.bib} file.

\section{Mathematical Expressions}
Here are some example equations:

\subsection{Quadratic Formula}
The quadratic formula is given by:
\[x = \frac{-b \pm \sqrt{b^2 - 4ac}}{2a}\]

\subsection{Matrix Example}
A sample matrix:
\[
\begin{pmatrix}
a_{11} & a_{12} & a_{13} \\
a_{21} & a_{22} & a_{23} \\
a_{31} & a_{32} & a_{33}
\end{pmatrix}
\]

\section{Tables}
Below is a sample table showing data:

\begin{table}[h]
\centering
\begin{tabular}{@{}lrrr@{}}
\toprule
\textbf{Category} & \textbf{Value 1} & \textbf{Value 2} & \textbf{Total} \\
\midrule
Item A & 10.5 & 20.3 & 30.8 \\
Item B & 15.2 & 25.1 & 40.3 \\
Item C & 22.1 & 12.7 & 34.8 \\
\bottomrule
\end{tabular}
\caption{Sample Data Table}
\label{tab:sample}
\end{table}

\section{Lists}
\subsection{Bullet Points}
\begin{itemize}
  \item First item
  \item Second item
  \item Third item with sub-items:
    \begin{itemize}
      \item Sub-item 1
      \item Sub-item 2
    \end{itemize}
\end{itemize}

\subsection{Numbered List}
\begin{enumerate}
  \item First step
  \item Second step
  \item Third step
\end{enumerate}

\section{Citations and References}
When you use the citation search feature, references will be automatically managed. For example, you might cite a paper like this: \cite{example2024}. All citations will be automatically added to your \texttt{cite.bib} file.

\section{Conclusion}
This template demonstrates basic \LaTeX{} features including mathematical formulas, tables, and lists. The bibliography system is now set up to work seamlessly with Bibby AI's citation management. Feel free to modify it for your needs!

\bibliographystyle{plain}
\bibliography{cite}

\end{document}\documentclass{article}
\usepackage[utf8]{inputenc}
\usepackage{booktabs}
\usepackage{amsmath}
\usepackage{natbib}

\title{My Document}
\author{Your Name}
\date{\today}

\begin{document}

\maketitle

\section{Introduction}
This is a sample document demonstrating various \LaTeX{} features including mathematical equations, tables, and formatting. You can cite papers using the citation search feature, and they will be automatically added to your \texttt{cite.bib} file.

\section{Mathematical Expressions}
Here are some example equations:

\subsection{Quadratic Formula}
The quadratic formula is given by:
\[x = \frac{-b \pm \sqrt{b^2 - 4ac}}{2a}\]

\subsection{Matrix Example}
A sample matrix:
\[
\begin{pmatrix}
a_{11} & a_{12} & a_{13} \\
a_{21} & a_{22} & a_{23} \\
a_{31} & a_{32} & a_{33}
\end{pmatrix}
\]

\section{Tables}
Below is a sample table showing data:

\begin{table}[h]
\centering
\begin{tabular}{@{}lrrr@{}}
\toprule
\textbf{Category} & \textbf{Value 1} & \textbf{Value 2} & \textbf{Total} \\
\midrule
Item A & 10.5 & 20.3 & 30.8 \\
Item B & 15.2 & 25.1 & 40.3 \\
Item C & 22.1 & 12.7 & 34.8 \\
\bottomrule
\end{tabular}
\caption{Sample Data Table}
\label{tab:sample}
\end{table}

\section{Lists}
\subsection{Bullet Points}
\begin{itemize}
  \item First item
  \item Second item
  \item Third item with sub-items:
    \begin{itemize}
      \item Sub-item 1
      \item Sub-item 2
    \end{itemize}
\end{itemize}

\subsection{Numbered List}
\begin{enumerate}
  \item First step
  \item Second step
  \item Third step
\end{enumerate}

\section{Citations and References}
When you use the citation search feature, references will be automatically managed. For example, you might cite a paper like this: \cite{example2024}. All citations will be automatically added to your \texttt{cite.bib} file.

\section{Conclusion}
This template demonstrates basic \LaTeX{} features including mathematical formulas, tables, and lists. The bibliography system is now set up to work seamlessly with Bibby AI's citation management. Feel free to modify it for your needs!

\bibliographystyle{plain}
\bibliography{cite}

\end{document}\documentclass{article}
\usepackage[utf8]{inputenc}
\usepackage{booktabs}
\usepackage{amsmath}
\usepackage{natbib}

\title{My Document}
\author{Your Name}
\date{\today}

\begin{document}

\maketitle

\section{Introduction}
This is a sample document demonstrating various \LaTeX{} features including mathematical equations, tables, and formatting. You can cite papers using the citation search feature, and they will be automatically added to your \texttt{cite.bib} file.

\section{Mathematical Expressions}
Here are some example equations:

\subsection{Quadratic Formula}
The quadratic formula is given by:
\[x = \frac{-b \pm \sqrt{b^2 - 4ac}}{2a}\]

\subsection{Matrix Example}
A sample matrix:
\[
\begin{pmatrix}
a_{11} & a_{12} & a_{13} \\
a_{21} & a_{22} & a_{23} \\
a_{31} & a_{32} & a_{33}
\end{pmatrix}
\]

\section{Tables}
Below is a sample table showing data:

\begin{table}[h]
\centering
\begin{tabular}{@{}lrrr@{}}
\toprule
\textbf{Category} & \textbf{Value 1} & \textbf{Value 2} & \textbf{Total} \\
\midrule
Item A & 10.5 & 20.3 & 30.8 \\
Item B & 15.2 & 25.1 & 40.3 \\
Item C & 22.1 & 12.7 & 34.8 \\
\bottomrule
\end{tabular}
\caption{Sample Data Table}
\label{tab:sample}
\end{table}

\section{Lists}
\subsection{Bullet Points}
\begin{itemize}
  \item First item
  \item Second item
  \item Third item with sub-items:
    \begin{itemize}
      \item Sub-item 1
      \item Sub-item 2
    \end{itemize}
\end{itemize}

\subsection{Numbered List}
\begin{enumerate}
  \item First step
  \item Second step
  \item Third step
\end{enumerate}

\section{Citations and References}
When you use the citation search feature, references will be automatically managed. For example, you might cite a paper like this: \cite{example2024}. All citations will be automatically added to your \texttt{cite.bib} file.

\section{Conclusion}
This template demonstrates basic \LaTeX{} features including mathematical formulas, tables, and lists. The bibliography system is now set up to work seamlessly with Bibby AI's citation management. Feel free to modify it for your needs!

\bibliographystyle{plain}
\bibliography{cite}

\end{document}\documentclass{article}
\usepackage[utf8]{inputenc}
\usepackage{booktabs}
\usepackage{amsmath}
\usepackage{natbib}

\title{My Document}
\author{Your Name}
\date{\today}

\begin{document}

\maketitle

\section{Introduction}
This is a sample document demonstrating various \LaTeX{} features including mathematical equations, tables, and formatting. You can cite papers using the citation search feature, and they will be automatically added to your \texttt{cite.bib} file.

\section{Mathematical Expressions}
Here are some example equations:

\subsection{Quadratic Formula}
The quadratic formula is given by:
\[x = \frac{-b \pm \sqrt{b^2 - 4ac}}{2a}\]

\subsection{Matrix Example}
A sample matrix:
\[
\begin{pmatrix}
a_{11} & a_{12} & a_{13} \\
a_{21} & a_{22} & a_{23} \\
a_{31} & a_{32} & a_{33}
\end{pmatrix}
\]

\section{Tables}
Below is a sample table showing data:

\begin{table}[h]
\centering
\begin{tabular}{@{}lrrr@{}}
\toprule
\textbf{Category} & \textbf{Value 1} & \textbf{Value 2} & \textbf{Total} \\
\midrule
Item A & 10.5 & 20.3 & 30.8 \\
Item B & 15.2 & 25.1 & 40.3 \\
Item C & 22.1 & 12.7 & 34.8 \\
\bottomrule
\end{tabular}
\caption{Sample Data Table}
\label{tab:sample}
\end{table}

\section{Lists}
\subsection{Bullet Points}
\begin{itemize}
  \item First item
  \item Second item
  \item Third item with sub-items:
    \begin{itemize}
      \item Sub-item 1
      \item Sub-item 2
    \end{itemize}
\end{itemize}

\subsection{Numbered List}
\begin{enumerate}
  \item First step
  \item Second step
  \item Third step
\end{enumerate}

\section{Citations and References}
When you use the citation search feature, references will be automatically managed. For example, you might cite a paper like this: \cite{example2024}. All citations will be automatically added to your \texttt{cite.bib} file.

\section{Conclusion}
This template demonstrates basic \LaTeX{} features including mathematical formulas, tables, and lists. The bibliography system is now set up to work seamlessly with Bibby AI's citation management. Feel free to modify it for your needs!

\bibliographystyle{plain}
\bibliography{cite}

\end{document}\documentclass{article}
\usepackage[utf8]{inputenc}
\usepackage{booktabs}
\usepackage{amsmath}
\usepackage{natbib}

\title{My Document}
\author{Your Name}
\date{\today}

\begin{document}

\maketitle

\section{Introduction}
This is a sample document demonstrating various \LaTeX{} features including mathematical equations, tables, and formatting. You can cite papers using the citation search feature, and they will be automatically added to your \texttt{cite.bib} file.

\section{Mathematical Expressions}
Here are some example equations:

\subsection{Quadratic Formula}
The quadratic formula is given by:
\[x = \frac{-b \pm \sqrt{b^2 - 4ac}}{2a}\]

\subsection{Matrix Example}
A sample matrix:
\[
\begin{pmatrix}
a_{11} & a_{12} & a_{13} \\
a_{21} & a_{22} & a_{23} \\
a_{31} & a_{32} & a_{33}
\end{pmatrix}
\]

\section{Tables}
Below is a sample table showing data:

\begin{table}[h]
\centering
\begin{tabular}{@{}lrrr@{}}
\toprule
\textbf{Category} & \textbf{Value 1} & \textbf{Value 2} & \textbf{Total} \\
\midrule
Item A & 10.5 & 20.3 & 30.8 \\
Item B & 15.2 & 25.1 & 40.3 \\
Item C & 22.1 & 12.7 & 34.8 \\
\bottomrule
\end{tabular}
\caption{Sample Data Table}
\label{tab:sample}
\end{table}

\section{Lists}
\subsection{Bullet Points}
\begin{itemize}
  \item First item
  \item Second item
  \item Third item with sub-items:
    \begin{itemize}
      \item Sub-item 1
      \item Sub-item 2
    \end{itemize}
\end{itemize}

\subsection{Numbered List}
\begin{enumerate}
  \item First step
  \item Second step
  \item Third step
\end{enumerate}

\section{Citations and References}
When you use the citation search feature, references will be automatically managed. For example, you might cite a paper like this: \cite{example2024}. All citations will be automatically added to your \texttt{cite.bib} file.

\section{Conclusion}
This template demonstrates basic \LaTeX{} features including mathematical formulas, tables, and lists. The bibliography system is now set up to work seamlessly with Bibby AI's citation management. Feel free to modify it for your needs!

\bibliographystyle{plain}
\bibliography{cite}

\end{document}\documentclass{article}
\usepackage[utf8]{inputenc}
\usepackage{booktabs}
\usepackage{amsmath}
\usepackage{natbib}

\title{My Document}
\author{Your Name}
\date{\today}

\begin{document}

\maketitle

\section{Introduction}
This is a sample document demonstrating various \LaTeX{} features including mathematical equations, tables, and formatting. You can cite papers using the citation search feature, and they will be automatically added to your \texttt{cite.bib} file.

\section{Mathematical Expressions}
Here are some example equations:

\subsection{Quadratic Formula}
The quadratic formula is given by:
\[x = \frac{-b \pm \sqrt{b^2 - 4ac}}{2a}\]

\subsection{Matrix Example}
A sample matrix:
\[
\begin{pmatrix}
a_{11} & a_{12} & a_{13} \\
a_{21} & a_{22} & a_{23} \\
a_{31} & a_{32} & a_{33}
\end{pmatrix}
\]

\section{Tables}
Below is a sample table showing data:

\begin{table}[h]
\centering
\begin{tabular}{@{}lrrr@{}}
\toprule
\textbf{Category} & \textbf{Value 1} & \textbf{Value 2} & \textbf{Total} \\
\midrule
Item A & 10.5 & 20.3 & 30.8 \\
Item B & 15.2 & 25.1 & 40.3 \\
Item C & 22.1 & 12.7 & 34.8 \\
\bottomrule
\end{tabular}
\caption{Sample Data Table}
\label{tab:sample}
\end{table}

\section{Lists}
\subsection{Bullet Points}
\begin{itemize}
  \item First item
  \item Second item
  \item Third item with sub-items:
    \begin{itemize}
      \item Sub-item 1
      \item Sub-item 2
    \end{itemize}
\end{itemize}

\subsection{Numbered List}
\begin{enumerate}
  \item First step
  \item Second step
  \item Third step
\end{enumerate}

\section{Citations and References}
When you use the citation search feature, references will be automatically managed. For example, you might cite a paper like this: \cite{example2024}. All citations will be automatically added to your \texttt{cite.bib} file.

\section{Conclusion}
This template demonstrates basic \LaTeX{} features including mathematical formulas, tables, and lists. The bibliography system is now set up to work seamlessly with Bibby AI's citation management. Feel free to modify it for your needs!

\bibliographystyle{plain}
\bibliography{cite}

\end{document}\documentclass{article}
\usepackage[utf8]{inputenc}
\usepackage{booktabs}
\usepackage{amsmath}
\usepackage{natbib}

\title{My Document}
\author{Your Name}
\date{\today}

\begin{document}

\maketitle

\section{Introduction}
This is a sample document demonstrating various \LaTeX{} features including mathematical equations, tables, and formatting. You can cite papers using the citation search feature, and they will be automatically added to your \texttt{cite.bib} file.

\section{Mathematical Expressions}
Here are some example equations:

\subsection{Quadratic Formula}
The quadratic formula is given by:
\[x = \frac{-b \pm \sqrt{b^2 - 4ac}}{2a}\]

\subsection{Matrix Example}
A sample matrix:
\[
\begin{pmatrix}
a_{11} & a_{12} & a_{13} \\
a_{21} & a_{22} & a_{23} \\
a_{31} & a_{32} & a_{33}
\end{pmatrix}
\]

\section{Tables}
Below is a sample table showing data:

\begin{table}[h]
\centering
\begin{tabular}{@{}lrrr@{}}
\toprule
\textbf{Category} & \textbf{Value 1} & \textbf{Value 2} & \textbf{Total} \\
\midrule
Item A & 10.5 & 20.3 & 30.8 \\
Item B & 15.2 & 25.1 & 40.3 \\
Item C & 22.1 & 12.7 & 34.8 \\
\bottomrule
\end{tabular}
\caption{Sample Data Table}
\label{tab:sample}
\end{table}

\section{Lists}
\subsection{Bullet Points}
\begin{itemize}
  \item First item
  \item Second item
  \item Third item with sub-items:
    \begin{itemize}
      \item Sub-item 1
      \item Sub-item 2
    \end{itemize}
\end{itemize}

\subsection{Numbered List}
\begin{enumerate}
  \item First step
  \item Second step
  \item Third step
\end{enumerate}

\section{Citations and References}
When you use the citation search feature, references will be automatically managed. For example, you might cite a paper like this: \cite{example2024}. All citations will be automatically added to your \texttt{cite.bib} file.

\section{Conclusion}
This template demonstrates basic \LaTeX{} features including mathematical formulas, tables, and lists. The bibliography system is now set up to work seamlessly with Bibby AI's citation management. Feel free to modify it for your needs!

\bibliographystyle{plain}
\bibliography{cite}

\end{document}\documentclass{article}
\usepackage[utf8]{inputenc}
\usepackage{booktabs}
\usepackage{amsmath}
\usepackage{natbib}

\title{My Document}
\author{Your Name}
\date{\today}

\begin{document}

\maketitle

\section{Introduction}
This is a sample document demonstrating various \LaTeX{} features including mathematical equations, tables, and formatting. You can cite papers using the citation search feature, and they will be automatically added to your \texttt{cite.bib} file.

\section{Mathematical Expressions}
Here are some example equations:

\subsection{Quadratic Formula}
The quadratic formula is given by:
\[x = \frac{-b \pm \sqrt{b^2 - 4ac}}{2a}\]

\subsection{Matrix Example}
A sample matrix:
\[
\begin{pmatrix}
a_{11} & a_{12} & a_{13} \\
a_{21} & a_{22} & a_{23} \\
a_{31} & a_{32} & a_{33}
\end{pmatrix}
\]

\section{Tables}
Below is a sample table showing data:

\begin{table}[h]
\centering
\begin{tabular}{@{}lrrr@{}}
\toprule
\textbf{Category} & \textbf{Value 1} & \textbf{Value 2} & \textbf{Total} \\
\midrule
Item A & 10.5 & 20.3 & 30.8 \\
Item B & 15.2 & 25.1 & 40.3 \\
Item C & 22.1 & 12.7 & 34.8 \\
\bottomrule
\end{tabular}
\caption{Sample Data Table}
\label{tab:sample}
\end{table}

\section{Lists}
\subsection{Bullet Points}
\begin{itemize}
  \item First item
  \item Second item
  \item Third item with sub-items:
    \begin{itemize}
      \item Sub-item 1
      \item Sub-item 2
    \end{itemize}
\end{itemize}

\subsection{Numbered List}
\begin{enumerate}
  \item First step
  \item Second step
  \item Third step
\end{enumerate}

\section{Citations and References}
When you use the citation search feature, references will be automatically managed. For example, you might cite a paper like this: \cite{example2024}. All citations will be automatically added to your \texttt{cite.bib} file.

\section{Conclusion}
This template demonstrates basic \LaTeX{} features including mathematical formulas, tables, and lists. The bibliography system is now set up to work seamlessly with Bibby AI's citation management. Feel free to modify it for your needs!

\bibliographystyle{plain}
\bibliography{cite}

\end{document}\documentclass{article}
\usepackage[utf8]{inputenc}
\usepackage{booktabs}
\usepackage{amsmath}
\usepackage{natbib}

\title{My Document}
\author{Your Name}
\date{\today}

\begin{document}

\maketitle

\section{Introduction}
This is a sample document demonstrating various \LaTeX{} features including mathematical equations, tables, and formatting. You can cite papers using the citation search feature, and they will be automatically added to your \texttt{cite.bib} file.

\section{Mathematical Expressions}
Here are some example equations:

\subsection{Quadratic Formula}
The quadratic formula is given by:
\[x = \frac{-b \pm \sqrt{b^2 - 4ac}}{2a}\]

\subsection{Matrix Example}
A sample matrix:
\[
\begin{pmatrix}
a_{11} & a_{12} & a_{13} \\
a_{21} & a_{22} & a_{23} \\
a_{31} & a_{32} & a_{33}
\end{pmatrix}
\]

\section{Tables}
Below is a sample table showing data:

\begin{table}[h]
\centering
\begin{tabular}{@{}lrrr@{}}
\toprule
\textbf{Category} & \textbf{Value 1} & \textbf{Value 2} & \textbf{Total} \\
\midrule
Item A & 10.5 & 20.3 & 30.8 \\
Item B & 15.2 & 25.1 & 40.3 \\
Item C & 22.1 & 12.7 & 34.8 \\
\bottomrule
\end{tabular}
\caption{Sample Data Table}
\label{tab:sample}
\end{table}

\section{Lists}
\subsection{Bullet Points}
\begin{itemize}
  \item First item
  \item Second item
  \item Third item with sub-items:
    \begin{itemize}
      \item Sub-item 1
      \item Sub-item 2
    \end{itemize}
\end{itemize}

\subsection{Numbered List}
\begin{enumerate}
  \item First step
  \item Second step
  \item Third step
\end{enumerate}

\section{Citations and References}
When you use the citation search feature, references will be automatically managed. For example, you might cite a paper like this: \cite{example2024}. All citations will be automatically added to your \texttt{cite.bib} file.

\section{Conclusion}
This template demonstrates basic \LaTeX{} features including mathematical formulas, tables, and lists. The bibliography system is now set up to work seamlessly with Bibby AI's citation management. Feel free to modify it for your needs!

\bibliographystyle{plain}
\bibliography{cite}

\end{document}\documentclass{article}
\usepackage[utf8]{inputenc}
\usepackage{booktabs}
\usepackage{amsmath}
\usepackage{natbib}

\title{My Document}
\author{Your Name}
\date{\today}

\begin{document}

\maketitle

\section{Introduction}
This is a sample document demonstrating various \LaTeX{} features including mathematical equations, tables, and formatting. You can cite papers using the citation search feature, and they will be automatically added to your \texttt{cite.bib} file.

\section{Mathematical Expressions}
Here are some example equations:

\subsection{Quadratic Formula}
The quadratic formula is given by:
\[x = \frac{-b \pm \sqrt{b^2 - 4ac}}{2a}\]

\subsection{Matrix Example}
A sample matrix:
\[
\begin{pmatrix}
a_{11} & a_{12} & a_{13} \\
a_{21} & a_{22} & a_{23} \\
a_{31} & a_{32} & a_{33}
\end{pmatrix}
\]

\section{Tables}
Below is a sample table showing data:

\begin{table}[h]
\centering
\begin{tabular}{@{}lrrr@{}}
\toprule
\textbf{Category} & \textbf{Value 1} & \textbf{Value 2} & \textbf{Total} \\
\midrule
Item A & 10.5 & 20.3 & 30.8 \\
Item B & 15.2 & 25.1 & 40.3 \\
Item C & 22.1 & 12.7 & 34.8 \\
\bottomrule
\end{tabular}
\caption{Sample Data Table}
\label{tab:sample}
\end{table}

\section{Lists}
\subsection{Bullet Points}
\begin{itemize}
  \item First item
  \item Second item
  \item Third item with sub-items:
    \begin{itemize}
      \item Sub-item 1
      \item Sub-item 2
    \end{itemize}
\end{itemize}

\subsection{Numbered List}
\begin{enumerate}
  \item First step
  \item Second step
  \item Third step
\end{enumerate}

\section{Citations and References}
When you use the citation search feature, references will be automatically managed. For example, you might cite a paper like this: \cite{example2024}. All citations will be automatically added to your \texttt{cite.bib} file.

\section{Conclusion}
This template demonstrates basic \LaTeX{} features including mathematical formulas, tables, and lists. The bibliography system is now set up to work seamlessly with Bibby AI's citation management. Feel free to modify it for your needs!

\bibliographystyle{plain}
\bibliography{cite}

\end{document}\documentclass{article}
\usepackage[utf8]{inputenc}
\usepackage{booktabs}
\usepackage{amsmath}
\usepackage{natbib}

\title{My Document}
\author{Your Name}
\date{\today}

\begin{document}

\maketitle

\section{Introduction}
This is a sample document demonstrating various \LaTeX{} features including mathematical equations, tables, and formatting. You can cite papers using the citation search feature, and they will be automatically added to your \texttt{cite.bib} file.

\section{Mathematical Expressions}
Here are some example equations:

\subsection{Quadratic Formula}
The quadratic formula is given by:
\[x = \frac{-b \pm \sqrt{b^2 - 4ac}}{2a}\]

\subsection{Matrix Example}
A sample matrix:
\[
\begin{pmatrix}
a_{11} & a_{12} & a_{13} \\
a_{21} & a_{22} & a_{23} \\
a_{31} & a_{32} & a_{33}
\end{pmatrix}
\]

\section{Tables}
Below is a sample table showing data:

\begin{table}[h]
\centering
\begin{tabular}{@{}lrrr@{}}
\toprule
\textbf{Category} & \textbf{Value 1} & \textbf{Value 2} & \textbf{Total} \\
\midrule
Item A & 10.5 & 20.3 & 30.8 \\
Item B & 15.2 & 25.1 & 40.3 \\
Item C & 22.1 & 12.7 & 34.8 \\
\bottomrule
\end{tabular}
\caption{Sample Data Table}
\label{tab:sample}
\end{table}

\section{Lists}
\subsection{Bullet Points}
\begin{itemize}
  \item First item
  \item Second item
  \item Third item with sub-items:
    \begin{itemize}
      \item Sub-item 1
      \item Sub-item 2
    \end{itemize}
\end{itemize}

\subsection{Numbered List}
\begin{enumerate}
  \item First step
  \item Second step
  \item Third step
\end{enumerate}

\section{Citations and References}
When you use the citation search feature, references will be automatically managed. For example, you might cite a paper like this: \cite{example2024}. All citations will be automatically added to your \texttt{cite.bib} file.

\section{Conclusion}
This template demonstrates basic \LaTeX{} features including mathematical formulas, tables, and lists. The bibliography system is now set up to work seamlessly with Bibby AI's citation management. Feel free to modify it for your needs!

\bibliographystyle{plain}
\bibliography{cite}

\end{document}\documentclass{article}
\usepackage[utf8]{inputenc}
\usepackage{booktabs}
\usepackage{amsmath}
\usepackage{natbib}

\title{My Document}
\author{Your Name}
\date{\today}

\begin{document}

\maketitle

\section{Introduction}
This is a sample document demonstrating various \LaTeX{} features including mathematical equations, tables, and formatting. You can cite papers using the citation search feature, and they will be automatically added to your \texttt{cite.bib} file.

\section{Mathematical Expressions}
Here are some example equations:

\subsection{Quadratic Formula}
The quadratic formula is given by:
\[x = \frac{-b \pm \sqrt{b^2 - 4ac}}{2a}\]

\subsection{Matrix Example}
A sample matrix:
\[
\begin{pmatrix}
a_{11} & a_{12} & a_{13} \\
a_{21} & a_{22} & a_{23} \\
a_{31} & a_{32} & a_{33}
\end{pmatrix}
\]

\section{Tables}
Below is a sample table showing data:

\begin{table}[h]
\centering
\begin{tabular}{@{}lrrr@{}}
\toprule
\textbf{Category} & \textbf{Value 1} & \textbf{Value 2} & \textbf{Total} \\
\midrule
Item A & 10.5 & 20.3 & 30.8 \\
Item B & 15.2 & 25.1 & 40.3 \\
Item C & 22.1 & 12.7 & 34.8 \\
\bottomrule
\end{tabular}
\caption{Sample Data Table}
\label{tab:sample}
\end{table}

\section{Lists}
\subsection{Bullet Points}
\begin{itemize}
  \item First item
  \item Second item
  \item Third item with sub-items:
    \begin{itemize}
      \item Sub-item 1
      \item Sub-item 2
    \end{itemize}
\end{itemize}

\subsection{Numbered List}
\begin{enumerate}
  \item First step
  \item Second step
  \item Third step
\end{enumerate}

\section{Citations and References}
When you use the citation search feature, references will be automatically managed. For example, you might cite a paper like this: \cite{example2024}. All citations will be automatically added to your \texttt{cite.bib} file.

\section{Conclusion}
This template demonstrates basic \LaTeX{} features including mathematical formulas, tables, and lists. The bibliography system is now set up to work seamlessly with Bibby AI's citation management. Feel free to modify it for your needs!

\bibliographystyle{plain}
\bibliography{cite}

\end{document}\documentclass{article}
\usepackage[utf8]{inputenc}
\usepackage{booktabs}
\usepackage{amsmath}
\usepackage{natbib}

\title{My Document}
\author{Your Name}
\date{\today}

\begin{document}

\maketitle

\section{Introduction}
This is a sample document demonstrating various \LaTeX{} features including mathematical equations, tables, and formatting. You can cite papers using the citation search feature, and they will be automatically added to your \texttt{cite.bib} file.

\section{Mathematical Expressions}
Here are some example equations:

\subsection{Quadratic Formula}
The quadratic formula is given by:
\[x = \frac{-b \pm \sqrt{b^2 - 4ac}}{2a}\]

\subsection{Matrix Example}
A sample matrix:
\[
\begin{pmatrix}
a_{11} & a_{12} & a_{13} \\
a_{21} & a_{22} & a_{23} \\
a_{31} & a_{32} & a_{33}
\end{pmatrix}
\]

\section{Tables}
Below is a sample table showing data:

\begin{table}[h]
\centering
\begin{tabular}{@{}lrrr@{}}
\toprule
\textbf{Category} & \textbf{Value 1} & \textbf{Value 2} & \textbf{Total} \\
\midrule
Item A & 10.5 & 20.3 & 30.8 \\
Item B & 15.2 & 25.1 & 40.3 \\
Item C & 22.1 & 12.7 & 34.8 \\
\bottomrule
\end{tabular}
\caption{Sample Data Table}
\label{tab:sample}
\end{table}

\section{Lists}
\subsection{Bullet Points}
\begin{itemize}
  \item First item
  \item Second item
  \item Third item with sub-items:
    \begin{itemize}
      \item Sub-item 1
      \item Sub-item 2
    \end{itemize}
\end{itemize}

\subsection{Numbered List}
\begin{enumerate}
  \item First step
  \item Second step
  \item Third step
\end{enumerate}

\section{Citations and References}
When you use the citation search feature, references will be automatically managed. For example, you might cite a paper like this: \cite{example2024}. All citations will be automatically added to your \texttt{cite.bib} file.

\section{Conclusion}
This template demonstrates basic \LaTeX{} features including mathematical formulas, tables, and lists. The bibliography system is now set up to work seamlessly with Bibby AI's citation management. Feel free to modify it for your needs!

\bibliographystyle{plain}
\bibliography{cite}

\end{document}\documentclass{article}
\usepackage[utf8]{inputenc}
\usepackage{booktabs}
\usepackage{amsmath}
\usepackage{natbib}

\title{My Document}
\author{Your Name}
\date{\today}

\begin{document}

\maketitle

\section{Introduction}
This is a sample document demonstrating various \LaTeX{} features including mathematical equations, tables, and formatting. You can cite papers using the citation search feature, and they will be automatically added to your \texttt{cite.bib} file.

\section{Mathematical Expressions}
Here are some example equations:

\subsection{Quadratic Formula}
The quadratic formula is given by:
\[x = \frac{-b \pm \sqrt{b^2 - 4ac}}{2a}\]

\subsection{Matrix Example}
A sample matrix:
\[
\begin{pmatrix}
a_{11} & a_{12} & a_{13} \\
a_{21} & a_{22} & a_{23} \\
a_{31} & a_{32} & a_{33}
\end{pmatrix}
\]

\section{Tables}
Below is a sample table showing data:

\begin{table}[h]
\centering
\begin{tabular}{@{}lrrr@{}}
\toprule
\textbf{Category} & \textbf{Value 1} & \textbf{Value 2} & \textbf{Total} \\
\midrule
Item A & 10.5 & 20.3 & 30.8 \\
Item B & 15.2 & 25.1 & 40.3 \\
Item C & 22.1 & 12.7 & 34.8 \\
\bottomrule
\end{tabular}
\caption{Sample Data Table}
\label{tab:sample}
\end{table}

\section{Lists}
\subsection{Bullet Points}
\begin{itemize}
  \item First item
  \item Second item
  \item Third item with sub-items:
    \begin{itemize}
      \item Sub-item 1
      \item Sub-item 2
    \end{itemize}
\end{itemize}

\subsection{Numbered List}
\begin{enumerate}
  \item First step
  \item Second step
  \item Third step
\end{enumerate}

\section{Citations and References}
When you use the citation search feature, references will be automatically managed. For example, you might cite a paper like this: \cite{example2024}. All citations will be automatically added to your \texttt{cite.bib} file.

\section{Conclusion}
This template demonstrates basic \LaTeX{} features including mathematical formulas, tables, and lists. The bibliography system is now set up to work seamlessly with Bibby AI's citation management. Feel free to modify it for your needs!

\bibliographystyle{plain}
\bibliography{cite}

\end{document}\documentclass{article}
\usepackage[utf8]{inputenc}
\usepackage{booktabs}
\usepackage{amsmath}
\usepackage{natbib}

\title{My Document}
\author{Your Name}
\date{\today}

\begin{document}

\maketitle

\section{Introduction}
This is a sample document demonstrating various \LaTeX{} features including mathematical equations, tables, and formatting. You can cite papers using the citation search feature, and they will be automatically added to your \texttt{cite.bib} file.

\section{Mathematical Expressions}
Here are some example equations:

\subsection{Quadratic Formula}
The quadratic formula is given by:
\[x = \frac{-b \pm \sqrt{b^2 - 4ac}}{2a}\]

\subsection{Matrix Example}
A sample matrix:
\[
\begin{pmatrix}
a_{11} & a_{12} & a_{13} \\
a_{21} & a_{22} & a_{23} \\
a_{31} & a_{32} & a_{33}
\end{pmatrix}
\]

\section{Tables}
Below is a sample table showing data:

\begin{table}[h]
\centering
\begin{tabular}{@{}lrrr@{}}
\toprule
\textbf{Category} & \textbf{Value 1} & \textbf{Value 2} & \textbf{Total} \\
\midrule
Item A & 10.5 & 20.3 & 30.8 \\
Item B & 15.2 & 25.1 & 40.3 \\
Item C & 22.1 & 12.7 & 34.8 \\
\bottomrule
\end{tabular}
\caption{Sample Data Table}
\label{tab:sample}
\end{table}

\section{Lists}
\subsection{Bullet Points}
\begin{itemize}
  \item First item
  \item Second item
  \item Third item with sub-items:
    \begin{itemize}
      \item Sub-item 1
      \item Sub-item 2
    \end{itemize}
\end{itemize}

\subsection{Numbered List}
\begin{enumerate}
  \item First step
  \item Second step
  \item Third step
\end{enumerate}

\section{Citations and References}
When you use the citation search feature, references will be automatically managed. For example, you might cite a paper like this: \cite{example2024}. All citations will be automatically added to your \texttt{cite.bib} file.

\section{Conclusion}
This template demonstrates basic \LaTeX{} features including mathematical formulas, tables, and lists. The bibliography system is now set up to work seamlessly with Bibby AI's citation management. Feel free to modify it for your needs!

\bibliographystyle{plain}
\bibliography{cite}

\end{document}\documentclass{article}
\usepackage[utf8]{inputenc}
\usepackage{booktabs}
\usepackage{amsmath}
\usepackage{natbib}

\title{My Document}
\author{Your Name}
\date{\today}

\begin{document}

\maketitle

\section{Introduction}
This is a sample document demonstrating various \LaTeX{} features including mathematical equations, tables, and formatting. You can cite papers using the citation search feature, and they will be automatically added to your \texttt{cite.bib} file.

\section{Mathematical Expressions}
Here are some example equations:

\subsection{Quadratic Formula}
The quadratic formula is given by:
\[x = \frac{-b \pm \sqrt{b^2 - 4ac}}{2a}\]

\subsection{Matrix Example}
A sample matrix:
\[
\begin{pmatrix}
a_{11} & a_{12} & a_{13} \\
a_{21} & a_{22} & a_{23} \\
a_{31} & a_{32} & a_{33}
\end{pmatrix}
\]

\section{Tables}
Below is a sample table showing data:

\begin{table}[h]
\centering
\begin{tabular}{@{}lrrr@{}}
\toprule
\textbf{Category} & \textbf{Value 1} & \textbf{Value 2} & \textbf{Total} \\
\midrule
Item A & 10.5 & 20.3 & 30.8 \\
Item B & 15.2 & 25.1 & 40.3 \\
Item C & 22.1 & 12.7 & 34.8 \\
\bottomrule
\end{tabular}
\caption{Sample Data Table}
\label{tab:sample}
\end{table}

\section{Lists}
\subsection{Bullet Points}
\begin{itemize}
  \item First item
  \item Second item
  \item Third item with sub-items:
    \begin{itemize}
      \item Sub-item 1
      \item Sub-item 2
    \end{itemize}
\end{itemize}

\subsection{Numbered List}
\begin{enumerate}
  \item First step
  \item Second step
  \item Third step
\end{enumerate}

\section{Citations and References}
When you use the citation search feature, references will be automatically managed. For example, you might cite a paper like this: \cite{example2024}. All citations will be automatically added to your \texttt{cite.bib} file.

\section{Conclusion}
This template demonstrates basic \LaTeX{} features including mathematical formulas, tables, and lists. The bibliography system is now set up to work seamlessly with Bibby AI's citation management. Feel free to modify it for your needs!

\bibliographystyle{plain}
\bibliography{cite}

\end{document}\documentclass{article}
\usepackage[utf8]{inputenc}
\usepackage{booktabs}
\usepackage{amsmath}
\usepackage{natbib}

\title{My Document}
\author{Your Name}
\date{\today}

\begin{document}

\maketitle

\section{Introduction}
This is a sample document demonstrating various \LaTeX{} features including mathematical equations, tables, and formatting. You can cite papers using the citation search feature, and they will be automatically added to your \texttt{cite.bib} file.

\section{Mathematical Expressions}
Here are some example equations:

\subsection{Quadratic Formula}
The quadratic formula is given by:
\[x = \frac{-b \pm \sqrt{b^2 - 4ac}}{2a}\]

\subsection{Matrix Example}
A sample matrix:
\[
\begin{pmatrix}
a_{11} & a_{12} & a_{13} \\
a_{21} & a_{22} & a_{23} \\
a_{31} & a_{32} & a_{33}
\end{pmatrix}
\]

\section{Tables}
Below is a sample table showing data:

\begin{table}[h]
\centering
\begin{tabular}{@{}lrrr@{}}
\toprule
\textbf{Category} & \textbf{Value 1} & \textbf{Value 2} & \textbf{Total} \\
\midrule
Item A & 10.5 & 20.3 & 30.8 \\
Item B & 15.2 & 25.1 & 40.3 \\
Item C & 22.1 & 12.7 & 34.8 \\
\bottomrule
\end{tabular}
\caption{Sample Data Table}
\label{tab:sample}
\end{table}

\section{Lists}
\subsection{Bullet Points}
\begin{itemize}
  \item First item
  \item Second item
  \item Third item with sub-items:
    \begin{itemize}
      \item Sub-item 1
      \item Sub-item 2
    \end{itemize}
\end{itemize}

\subsection{Numbered List}
\begin{enumerate}
  \item First step
  \item Second step
  \item Third step
\end{enumerate}

\section{Citations and References}
When you use the citation search feature, references will be automatically managed. For example, you might cite a paper like this: \cite{example2024}. All citations will be automatically added to your \texttt{cite.bib} file.

\section{Conclusion}
This template demonstrates basic \LaTeX{} features including mathematical formulas, tables, and lists. The bibliography system is now set up to work seamlessly with Bibby AI's citation management. Feel free to modify it for your needs!

\bibliographystyle{plain}
\bibliography{cite}

\end{document}\documentclass{article}
\usepackage[utf8]{inputenc}
\usepackage{booktabs}
\usepackage{amsmath}
\usepackage{natbib}

\title{My Document}
\author{Your Name}
\date{\today}

\begin{document}

\maketitle

\section{Introduction}
This is a sample document demonstrating various \LaTeX{} features including mathematical equations, tables, and formatting. You can cite papers using the citation search feature, and they will be automatically added to your \texttt{cite.bib} file.

\section{Mathematical Expressions}
Here are some example equations:

\subsection{Quadratic Formula}
The quadratic formula is given by:
\[x = \frac{-b \pm \sqrt{b^2 - 4ac}}{2a}\]

\subsection{Matrix Example}
A sample matrix:
\[
\begin{pmatrix}
a_{11} & a_{12} & a_{13} \\
a_{21} & a_{22} & a_{23} \\
a_{31} & a_{32} & a_{33}
\end{pmatrix}
\]

\section{Tables}
Below is a sample table showing data:

\begin{table}[h]
\centering
\begin{tabular}{@{}lrrr@{}}
\toprule
\textbf{Category} & \textbf{Value 1} & \textbf{Value 2} & \textbf{Total} \\
\midrule
Item A & 10.5 & 20.3 & 30.8 \\
Item B & 15.2 & 25.1 & 40.3 \\
Item C & 22.1 & 12.7 & 34.8 \\
\bottomrule
\end{tabular}
\caption{Sample Data Table}
\label{tab:sample}
\end{table}

\section{Lists}
\subsection{Bullet Points}
\begin{itemize}
  \item First item
  \item Second item
  \item Third item with sub-items:
    \begin{itemize}
      \item Sub-item 1
      \item Sub-item 2
    \end{itemize}
\end{itemize}

\subsection{Numbered List}
\begin{enumerate}
  \item First step
  \item Second step
  \item Third step
\end{enumerate}

\section{Citations and References}
When you use the citation search feature, references will be automatically managed. For example, you might cite a paper like this: \cite{example2024}. All citations will be automatically added to your \texttt{cite.bib} file.

\section{Conclusion}
This template demonstrates basic \LaTeX{} features including mathematical formulas, tables, and lists. The bibliography system is now set up to work seamlessly with Bibby AI's citation management. Feel free to modify it for your needs!

\bibliographystyle{plain}
\bibliography{cite}

\end{document}\documentclass{article}
\usepackage[utf8]{inputenc}
\usepackage{booktabs}
\usepackage{amsmath}
\usepackage{natbib}

\title{My Document}
\author{Your Name}
\date{\today}

\begin{document}

\maketitle

\section{Introduction}
This is a sample document demonstrating various \LaTeX{} features including mathematical equations, tables, and formatting. You can cite papers using the citation search feature, and they will be automatically added to your \texttt{cite.bib} file.

\section{Mathematical Expressions}
Here are some example equations:

\subsection{Quadratic Formula}
The quadratic formula is given by:
\[x = \frac{-b \pm \sqrt{b^2 - 4ac}}{2a}\]

\subsection{Matrix Example}
A sample matrix:
\[
\begin{pmatrix}
a_{11} & a_{12} & a_{13} \\
a_{21} & a_{22} & a_{23} \\
a_{31} & a_{32} & a_{33}
\end{pmatrix}
\]

\section{Tables}
Below is a sample table showing data:

\begin{table}[h]
\centering
\begin{tabular}{@{}lrrr@{}}
\toprule
\textbf{Category} & \textbf{Value 1} & \textbf{Value 2} & \textbf{Total} \\
\midrule
Item A & 10.5 & 20.3 & 30.8 \\
Item B & 15.2 & 25.1 & 40.3 \\
Item C & 22.1 & 12.7 & 34.8 \\
\bottomrule
\end{tabular}
\caption{Sample Data Table}
\label{tab:sample}
\end{table}

\section{Lists}
\subsection{Bullet Points}
\begin{itemize}
  \item First item
  \item Second item
  \item Third item with sub-items:
    \begin{itemize}
      \item Sub-item 1
      \item Sub-item 2
    \end{itemize}
\end{itemize}

\subsection{Numbered List}
\begin{enumerate}
  \item First step
  \item Second step
  \item Third step
\end{enumerate}

\section{Citations and References}
When you use the citation search feature, references will be automatically managed. For example, you might cite a paper like this: \cite{example2024}. All citations will be automatically added to your \texttt{cite.bib} file.

\section{Conclusion}
This template demonstrates basic \LaTeX{} features including mathematical formulas, tables, and lists. The bibliography system is now set up to work seamlessly with Bibby AI's citation management. Feel free to modify it for your needs!

\bibliographystyle{plain}
\bibliography{cite}

\end{document}\documentclass{article}
\usepackage[utf8]{inputenc}
\usepackage{booktabs}
\usepackage{amsmath}
\usepackage{natbib}

\title{My Document}
\author{Your Name}
\date{\today}

\begin{document}

\maketitle

\section{Introduction}
This is a sample document demonstrating various \LaTeX{} features including mathematical equations, tables, and formatting. You can cite papers using the citation search feature, and they will be automatically added to your \texttt{cite.bib} file.

\section{Mathematical Expressions}
Here are some example equations:

\subsection{Quadratic Formula}
The quadratic formula is given by:
\[x = \frac{-b \pm \sqrt{b^2 - 4ac}}{2a}\]

\subsection{Matrix Example}
A sample matrix:
\[
\begin{pmatrix}
a_{11} & a_{12} & a_{13} \\
a_{21} & a_{22} & a_{23} \\
a_{31} & a_{32} & a_{33}
\end{pmatrix}
\]

\section{Tables}
Below is a sample table showing data:

\begin{table}[h]
\centering
\begin{tabular}{@{}lrrr@{}}
\toprule
\textbf{Category} & \textbf{Value 1} & \textbf{Value 2} & \textbf{Total} \\
\midrule
Item A & 10.5 & 20.3 & 30.8 \\
Item B & 15.2 & 25.1 & 40.3 \\
Item C & 22.1 & 12.7 & 34.8 \\
\bottomrule
\end{tabular}
\caption{Sample Data Table}
\label{tab:sample}
\end{table}

\section{Lists}
\subsection{Bullet Points}
\begin{itemize}
  \item First item
  \item Second item
  \item Third item with sub-items:
    \begin{itemize}
      \item Sub-item 1
      \item Sub-item 2
    \end{itemize}
\end{itemize}

\subsection{Numbered List}
\begin{enumerate}
  \item First step
  \item Second step
  \item Third step
\end{enumerate}

\section{Citations and References}
When you use the citation search feature, references will be automatically managed. For example, you might cite a paper like this: \cite{example2024}. All citations will be automatically added to your \texttt{cite.bib} file.

\section{Conclusion}
This template demonstrates basic \LaTeX{} features including mathematical formulas, tables, and lists. The bibliography system is now set up to work seamlessly with Bibby AI's citation management. Feel free to modify it for your needs!

\bibliographystyle{plain}
\bibliography{cite}

\end{document}\documentclass{article}
\usepackage[utf8]{inputenc}
\usepackage{booktabs}
\usepackage{amsmath}
\usepackage{natbib}

\title{My Document}
\author{Your Name}
\date{\today}

\begin{document}

\maketitle

\section{Introduction}
This is a sample document demonstrating various \LaTeX{} features including mathematical equations, tables, and formatting. You can cite papers using the citation search feature, and they will be automatically added to your \texttt{cite.bib} file.

\section{Mathematical Expressions}
Here are some example equations:

\subsection{Quadratic Formula}
The quadratic formula is given by:
\[x = \frac{-b \pm \sqrt{b^2 - 4ac}}{2a}\]

\subsection{Matrix Example}
A sample matrix:
\[
\begin{pmatrix}
a_{11} & a_{12} & a_{13} \\
a_{21} & a_{22} & a_{23} \\
a_{31} & a_{32} & a_{33}
\end{pmatrix}
\]

\section{Tables}
Below is a sample table showing data:

\begin{table}[h]
\centering
\begin{tabular}{@{}lrrr@{}}
\toprule
\textbf{Category} & \textbf{Value 1} & \textbf{Value 2} & \textbf{Total} \\
\midrule
Item A & 10.5 & 20.3 & 30.8 \\
Item B & 15.2 & 25.1 & 40.3 \\
Item C & 22.1 & 12.7 & 34.8 \\
\bottomrule
\end{tabular}
\caption{Sample Data Table}
\label{tab:sample}
\end{table}

\section{Lists}
\subsection{Bullet Points}
\begin{itemize}
  \item First item
  \item Second item
  \item Third item with sub-items:
    \begin{itemize}
      \item Sub-item 1
      \item Sub-item 2
    \end{itemize}
\end{itemize}

\subsection{Numbered List}
\begin{enumerate}
  \item First step
  \item Second step
  \item Third step
\end{enumerate}

\section{Citations and References}
When you use the citation search feature, references will be automatically managed. For example, you might cite a paper like this: \cite{example2024}. All citations will be automatically added to your \texttt{cite.bib} file.

\section{Conclusion}
This template demonstrates basic \LaTeX{} features including mathematical formulas, tables, and lists. The bibliography system is now set up to work seamlessly with Bibby AI's citation management. Feel free to modify it for your needs!

\bibliographystyle{plain}
\bibliography{cite}

\end{document}\documentclass{article}
\usepackage[utf8]{inputenc}
\usepackage{booktabs}
\usepackage{amsmath}
\usepackage{natbib}

\title{My Document}
\author{Your Name}
\date{\today}

\begin{document}

\maketitle

\section{Introduction}
This is a sample document demonstrating various \LaTeX{} features including mathematical equations, tables, and formatting. You can cite papers using the citation search feature, and they will be automatically added to your \texttt{cite.bib} file.

\section{Mathematical Expressions}
Here are some example equations:

\subsection{Quadratic Formula}
The quadratic formula is given by:
\[x = \frac{-b \pm \sqrt{b^2 - 4ac}}{2a}\]

\subsection{Matrix Example}
A sample matrix:
\[
\begin{pmatrix}
a_{11} & a_{12} & a_{13} \\
a_{21} & a_{22} & a_{23} \\
a_{31} & a_{32} & a_{33}
\end{pmatrix}
\]

\section{Tables}
Below is a sample table showing data:

\begin{table}[h]
\centering
\begin{tabular}{@{}lrrr@{}}
\toprule
\textbf{Category} & \textbf{Value 1} & \textbf{Value 2} & \textbf{Total} \\
\midrule
Item A & 10.5 & 20.3 & 30.8 \\
Item B & 15.2 & 25.1 & 40.3 \\
Item C & 22.1 & 12.7 & 34.8 \\
\bottomrule
\end{tabular}
\caption{Sample Data Table}
\label{tab:sample}
\end{table}

\section{Lists}
\subsection{Bullet Points}
\begin{itemize}
  \item First item
  \item Second item
  \item Third item with sub-items:
    \begin{itemize}
      \item Sub-item 1
      \item Sub-item 2
    \end{itemize}
\end{itemize}

\subsection{Numbered List}
\begin{enumerate}
  \item First step
  \item Second step
  \item Third step
\end{enumerate}

\section{Citations and References}
When you use the citation search feature, references will be automatically managed. For example, you might cite a paper like this: \cite{example2024}. All citations will be automatically added to your \texttt{cite.bib} file.

\section{Conclusion}
This template demonstrates basic \LaTeX{} features including mathematical formulas, tables, and lists. The bibliography system is now set up to work seamlessly with Bibby AI's citation management. Feel free to modify it for your needs!

\bibliographystyle{plain}
\bibliography{cite}

\end{document}\documentclass{article}
\usepackage[utf8]{inputenc}
\usepackage{booktabs}
\usepackage{amsmath}
\usepackage{natbib}

\title{My Document}
\author{Your Name}
\date{\today}

\begin{document}

\maketitle

\section{Introduction}
This is a sample document demonstrating various \LaTeX{} features including mathematical equations, tables, and formatting. You can cite papers using the citation search feature, and they will be automatically added to your \texttt{cite.bib} file.

\section{Mathematical Expressions}
Here are some example equations:

\subsection{Quadratic Formula}
The quadratic formula is given by:
\[x = \frac{-b \pm \sqrt{b^2 - 4ac}}{2a}\]

\subsection{Matrix Example}
A sample matrix:
\[
\begin{pmatrix}
a_{11} & a_{12} & a_{13} \\
a_{21} & a_{22} & a_{23} \\
a_{31} & a_{32} & a_{33}
\end{pmatrix}
\]

\section{Tables}
Below is a sample table showing data:

\begin{table}[h]
\centering
\begin{tabular}{@{}lrrr@{}}
\toprule
\textbf{Category} & \textbf{Value 1} & \textbf{Value 2} & \textbf{Total} \\
\midrule
Item A & 10.5 & 20.3 & 30.8 \\
Item B & 15.2 & 25.1 & 40.3 \\
Item C & 22.1 & 12.7 & 34.8 \\
\bottomrule
\end{tabular}
\caption{Sample Data Table}
\label{tab:sample}
\end{table}

\section{Lists}
\subsection{Bullet Points}
\begin{itemize}
  \item First item
  \item Second item
  \item Third item with sub-items:
    \begin{itemize}
      \item Sub-item 1
      \item Sub-item 2
    \end{itemize}
\end{itemize}

\subsection{Numbered List}
\begin{enumerate}
  \item First step
  \item Second step
  \item Third step
\end{enumerate}

\section{Citations and References}
When you use the citation search feature, references will be automatically managed. For example, you might cite a paper like this: \cite{example2024}. All citations will be automatically added to your \texttt{cite.bib} file.

\section{Conclusion}
This template demonstrates basic \LaTeX{} features including mathematical formulas, tables, and lists. The bibliography system is now set up to work seamlessly with Bibby AI's citation management. Feel free to modify it for your needs!

\bibliographystyle{plain}
\bibliography{cite}

\end{document}\documentclass{article}
\usepackage[utf8]{inputenc}
\usepackage{booktabs}
\usepackage{amsmath}
\usepackage{natbib}

\title{My Document}
\author{Your Name}
\date{\today}

\begin{document}

\maketitle

\section{Introduction}
This is a sample document demonstrating various \LaTeX{} features including mathematical equations, tables, and formatting. You can cite papers using the citation search feature, and they will be automatically added to your \texttt{cite.bib} file.

\section{Mathematical Expressions}
Here are some example equations:

\subsection{Quadratic Formula}
The quadratic formula is given by:
\[x = \frac{-b \pm \sqrt{b^2 - 4ac}}{2a}\]

\subsection{Matrix Example}
A sample matrix:
\[
\begin{pmatrix}
a_{11} & a_{12} & a_{13} \\
a_{21} & a_{22} & a_{23} \\
a_{31} & a_{32} & a_{33}
\end{pmatrix}
\]

\section{Tables}
Below is a sample table showing data:

\begin{table}[h]
\centering
\begin{tabular}{@{}lrrr@{}}
\toprule
\textbf{Category} & \textbf{Value 1} & \textbf{Value 2} & \textbf{Total} \\
\midrule
Item A & 10.5 & 20.3 & 30.8 \\
Item B & 15.2 & 25.1 & 40.3 \\
Item C & 22.1 & 12.7 & 34.8 \\
\bottomrule
\end{tabular}
\caption{Sample Data Table}
\label{tab:sample}
\end{table}

\section{Lists}
\subsection{Bullet Points}
\begin{itemize}
  \item First item
  \item Second item
  \item Third item with sub-items:
    \begin{itemize}
      \item Sub-item 1
      \item Sub-item 2
    \end{itemize}
\end{itemize}

\subsection{Numbered List}
\begin{enumerate}
  \item First step
  \item Second step
  \item Third step
\end{enumerate}

\section{Citations and References}
When you use the citation search feature, references will be automatically managed. For example, you might cite a paper like this: \cite{example2024}. All citations will be automatically added to your \texttt{cite.bib} file.

\section{Conclusion}
This template demonstrates basic \LaTeX{} features including mathematical formulas, tables, and lists. The bibliography system is now set up to work seamlessly with Bibby AI's citation management. Feel free to modify it for your needs!

\bibliographystyle{plain}
\bibliography{cite}

\end{document}\documentclass{article}
\usepackage[utf8]{inputenc}
\usepackage{booktabs}
\usepackage{amsmath}
\usepackage{natbib}

\title{My Document}
\author{Your Name}
\date{\today}

\begin{document}

\maketitle

\section{Introduction}
This is a sample document demonstrating various \LaTeX{} features including mathematical equations, tables, and formatting. You can cite papers using the citation search feature, and they will be automatically added to your \texttt{cite.bib} file.

\section{Mathematical Expressions}
Here are some example equations:

\subsection{Quadratic Formula}
The quadratic formula is given by:
\[x = \frac{-b \pm \sqrt{b^2 - 4ac}}{2a}\]

\subsection{Matrix Example}
A sample matrix:
\[
\begin{pmatrix}
a_{11} & a_{12} & a_{13} \\
a_{21} & a_{22} & a_{23} \\
a_{31} & a_{32} & a_{33}
\end{pmatrix}
\]

\section{Tables}
Below is a sample table showing data:

\begin{table}[h]
\centering
\begin{tabular}{@{}lrrr@{}}
\toprule
\textbf{Category} & \textbf{Value 1} & \textbf{Value 2} & \textbf{Total} \\
\midrule
Item A & 10.5 & 20.3 & 30.8 \\
Item B & 15.2 & 25.1 & 40.3 \\
Item C & 22.1 & 12.7 & 34.8 \\
\bottomrule
\end{tabular}
\caption{Sample Data Table}
\label{tab:sample}
\end{table}

\section{Lists}
\subsection{Bullet Points}
\begin{itemize}
  \item First item
  \item Second item
  \item Third item with sub-items:
    \begin{itemize}
      \item Sub-item 1
      \item Sub-item 2
    \end{itemize}
\end{itemize}

\subsection{Numbered List}
\begin{enumerate}
  \item First step
  \item Second step
  \item Third step
\end{enumerate}

\section{Citations and References}
When you use the citation search feature, references will be automatically managed. For example, you might cite a paper like this: \cite{example2024}. All citations will be automatically added to your \texttt{cite.bib} file.

\section{Conclusion}
This template demonstrates basic \LaTeX{} features including mathematical formulas, tables, and lists. The bibliography system is now set up to work seamlessly with Bibby AI's citation management. Feel free to modify it for your needs!

\bibliographystyle{plain}
\bibliography{cite}

\end{document}\documentclass{article}
\usepackage[utf8]{inputenc}
\usepackage{booktabs}
\usepackage{amsmath}
\usepackage{natbib}

\title{My Document}
\author{Your Name}
\date{\today}

\begin{document}

\maketitle

\section{Introduction}
This is a sample document demonstrating various \LaTeX{} features including mathematical equations, tables, and formatting. You can cite papers using the citation search feature, and they will be automatically added to your \texttt{cite.bib} file.

\section{Mathematical Expressions}
Here are some example equations:

\subsection{Quadratic Formula}
The quadratic formula is given by:
\[x = \frac{-b \pm \sqrt{b^2 - 4ac}}{2a}\]

\subsection{Matrix Example}
A sample matrix:
\[
\begin{pmatrix}
a_{11} & a_{12} & a_{13} \\
a_{21} & a_{22} & a_{23} \\
a_{31} & a_{32} & a_{33}
\end{pmatrix}
\]

\section{Tables}
Below is a sample table showing data:

\begin{table}[h]
\centering
\begin{tabular}{@{}lrrr@{}}
\toprule
\textbf{Category} & \textbf{Value 1} & \textbf{Value 2} & \textbf{Total} \\
\midrule
Item A & 10.5 & 20.3 & 30.8 \\
Item B & 15.2 & 25.1 & 40.3 \\
Item C & 22.1 & 12.7 & 34.8 \\
\bottomrule
\end{tabular}
\caption{Sample Data Table}
\label{tab:sample}
\end{table}

\section{Lists}
\subsection{Bullet Points}
\begin{itemize}
  \item First item
  \item Second item
  \item Third item with sub-items:
    \begin{itemize}
      \item Sub-item 1
      \item Sub-item 2
    \end{itemize}
\end{itemize}

\subsection{Numbered List}
\begin{enumerate}
  \item First step
  \item Second step
  \item Third step
\end{enumerate}

\section{Citations and References}
When you use the citation search feature, references will be automatically managed. For example, you might cite a paper like this: \cite{example2024}. All citations will be automatically added to your \texttt{cite.bib} file.

\section{Conclusion}
This template demonstrates basic \LaTeX{} features including mathematical formulas, tables, and lists. The bibliography system is now set up to work seamlessly with Bibby AI's citation management. Feel free to modify it for your needs!

\bibliographystyle{plain}
\bibliography{cite}

\end{document}\documentclass{article}
\usepackage[utf8]{inputenc}
\usepackage{booktabs}
\usepackage{amsmath}
\usepackage{natbib}

\title{My Document}
\author{Your Name}
\date{\today}

\begin{document}

\maketitle

\section{Introduction}
This is a sample document demonstrating various \LaTeX{} features including mathematical equations, tables, and formatting. You can cite papers using the citation search feature, and they will be automatically added to your \texttt{cite.bib} file.

\section{Mathematical Expressions}
Here are some example equations:

\subsection{Quadratic Formula}
The quadratic formula is given by:
\[x = \frac{-b \pm \sqrt{b^2 - 4ac}}{2a}\]

\subsection{Matrix Example}
A sample matrix:
\[
\begin{pmatrix}
a_{11} & a_{12} & a_{13} \\
a_{21} & a_{22} & a_{23} \\
a_{31} & a_{32} & a_{33}
\end{pmatrix}
\]

\section{Tables}
Below is a sample table showing data:

\begin{table}[h]
\centering
\begin{tabular}{@{}lrrr@{}}
\toprule
\textbf{Category} & \textbf{Value 1} & \textbf{Value 2} & \textbf{Total} \\
\midrule
Item A & 10.5 & 20.3 & 30.8 \\
Item B & 15.2 & 25.1 & 40.3 \\
Item C & 22.1 & 12.7 & 34.8 \\
\bottomrule
\end{tabular}
\caption{Sample Data Table}
\label{tab:sample}
\end{table}

\section{Lists}
\subsection{Bullet Points}
\begin{itemize}
  \item First item
  \item Second item
  \item Third item with sub-items:
    \begin{itemize}
      \item Sub-item 1
      \item Sub-item 2
    \end{itemize}
\end{itemize}

\subsection{Numbered List}
\begin{enumerate}
  \item First step
  \item Second step
  \item Third step
\end{enumerate}

\section{Citations and References}
When you use the citation search feature, references will be automatically managed. For example, you might cite a paper like this: \cite{example2024}. All citations will be automatically added to your \texttt{cite.bib} file.

\section{Conclusion}
This template demonstrates basic \LaTeX{} features including mathematical formulas, tables, and lists. The bibliography system is now set up to work seamlessly with Bibby AI's citation management. Feel free to modify it for your needs!

\bibliographystyle{plain}
\bibliography{cite}

\end{document}\documentclass{article}
\usepackage[utf8]{inputenc}
\usepackage{booktabs}
\usepackage{amsmath}
\usepackage{natbib}

\title{My Document}
\author{Your Name}
\date{\today}

\begin{document}

\maketitle

\section{Introduction}
This is a sample document demonstrating various \LaTeX{} features including mathematical equations, tables, and formatting. You can cite papers using the citation search feature, and they will be automatically added to your \texttt{cite.bib} file.

\section{Mathematical Expressions}
Here are some example equations:

\subsection{Quadratic Formula}
The quadratic formula is given by:
\[x = \frac{-b \pm \sqrt{b^2 - 4ac}}{2a}\]

\subsection{Matrix Example}
A sample matrix:
\[
\begin{pmatrix}
a_{11} & a_{12} & a_{13} \\
a_{21} & a_{22} & a_{23} \\
a_{31} & a_{32} & a_{33}
\end{pmatrix}
\]

\section{Tables}
Below is a sample table showing data:

\begin{table}[h]
\centering
\begin{tabular}{@{}lrrr@{}}
\toprule
\textbf{Category} & \textbf{Value 1} & \textbf{Value 2} & \textbf{Total} \\
\midrule
Item A & 10.5 & 20.3 & 30.8 \\
Item B & 15.2 & 25.1 & 40.3 \\
Item C & 22.1 & 12.7 & 34.8 \\
\bottomrule
\end{tabular}
\caption{Sample Data Table}
\label{tab:sample}
\end{table}

\section{Lists}
\subsection{Bullet Points}
\begin{itemize}
  \item First item
  \item Second item
  \item Third item with sub-items:
    \begin{itemize}
      \item Sub-item 1
      \item Sub-item 2
    \end{itemize}
\end{itemize}

\subsection{Numbered List}
\begin{enumerate}
  \item First step
  \item Second step
  \item Third step
\end{enumerate}

\section{Citations and References}
When you use the citation search feature, references will be automatically managed. For example, you might cite a paper like this: \cite{example2024}. All citations will be automatically added to your \texttt{cite.bib} file.

\section{Conclusion}
This template demonstrates basic \LaTeX{} features including mathematical formulas, tables, and lists. The bibliography system is now set up to work seamlessly with Bibby AI's citation management. Feel free to modify it for your needs!

\bibliographystyle{plain}
\bibliography{cite}

\end{document}\documentclass{article}
\usepackage[utf8]{inputenc}
\usepackage{booktabs}
\usepackage{amsmath}
\usepackage{natbib}

\title{My Document}
\author{Your Name}
\date{\today}

\begin{document}

\maketitle

\section{Introduction}
This is a sample document demonstrating various \LaTeX{} features including mathematical equations, tables, and formatting. You can cite papers using the citation search feature, and they will be automatically added to your \texttt{cite.bib} file.

\section{Mathematical Expressions}
Here are some example equations:

\subsection{Quadratic Formula}
The quadratic formula is given by:
\[x = \frac{-b \pm \sqrt{b^2 - 4ac}}{2a}\]

\subsection{Matrix Example}
A sample matrix:
\[
\begin{pmatrix}
a_{11} & a_{12} & a_{13} \\
a_{21} & a_{22} & a_{23} \\
a_{31} & a_{32} & a_{33}
\end{pmatrix}
\]

\section{Tables}
Below is a sample table showing data:

\begin{table}[h]
\centering
\begin{tabular}{@{}lrrr@{}}
\toprule
\textbf{Category} & \textbf{Value 1} & \textbf{Value 2} & \textbf{Total} \\
\midrule
Item A & 10.5 & 20.3 & 30.8 \\
Item B & 15.2 & 25.1 & 40.3 \\
Item C & 22.1 & 12.7 & 34.8 \\
\bottomrule
\end{tabular}
\caption{Sample Data Table}
\label{tab:sample}
\end{table}

\section{Lists}
\subsection{Bullet Points}
\begin{itemize}
  \item First item
  \item Second item
  \item Third item with sub-items:
    \begin{itemize}
      \item Sub-item 1
      \item Sub-item 2
    \end{itemize}
\end{itemize}

\subsection{Numbered List}
\begin{enumerate}
  \item First step
  \item Second step
  \item Third step
\end{enumerate}

\section{Citations and References}
When you use the citation search feature, references will be automatically managed. For example, you might cite a paper like this: \cite{example2024}. All citations will be automatically added to your \texttt{cite.bib} file.

\section{Conclusion}
This template demonstrates basic \LaTeX{} features including mathematical formulas, tables, and lists. The bibliography system is now set up to work seamlessly with Bibby AI's citation management. Feel free to modify it for your needs!

\bibliographystyle{plain}
\bibliography{cite}

\end{document}\documentclass{article}
\usepackage[utf8]{inputenc}
\usepackage{booktabs}
\usepackage{amsmath}
\usepackage{natbib}

\title{My Document}
\author{Your Name}
\date{\today}

\begin{document}

\maketitle

\section{Introduction}
This is a sample document demonstrating various \LaTeX{} features including mathematical equations, tables, and formatting. You can cite papers using the citation search feature, and they will be automatically added to your \texttt{cite.bib} file.

\section{Mathematical Expressions}
Here are some example equations:

\subsection{Quadratic Formula}
The quadratic formula is given by:
\[x = \frac{-b \pm \sqrt{b^2 - 4ac}}{2a}\]

\subsection{Matrix Example}
A sample matrix:
\[
\begin{pmatrix}
a_{11} & a_{12} & a_{13} \\
a_{21} & a_{22} & a_{23} \\
a_{31} & a_{32} & a_{33}
\end{pmatrix}
\]

\section{Tables}
Below is a sample table showing data:

\begin{table}[h]
\centering
\begin{tabular}{@{}lrrr@{}}
\toprule
\textbf{Category} & \textbf{Value 1} & \textbf{Value 2} & \textbf{Total} \\
\midrule
Item A & 10.5 & 20.3 & 30.8 \\
Item B & 15.2 & 25.1 & 40.3 \\
Item C & 22.1 & 12.7 & 34.8 \\
\bottomrule
\end{tabular}
\caption{Sample Data Table}
\label{tab:sample}
\end{table}

\section{Lists}
\subsection{Bullet Points}
\begin{itemize}
  \item First item
  \item Second item
  \item Third item with sub-items:
    \begin{itemize}
      \item Sub-item 1
      \item Sub-item 2
    \end{itemize}
\end{itemize}

\subsection{Numbered List}
\begin{enumerate}
  \item First step
  \item Second step
  \item Third step
\end{enumerate}

\section{Citations and References}
When you use the citation search feature, references will be automatically managed. For example, you might cite a paper like this: \cite{example2024}. All citations will be automatically added to your \texttt{cite.bib} file.

\section{Conclusion}
This template demonstrates basic \LaTeX{} features including mathematical formulas, tables, and lists. The bibliography system is now set up to work seamlessly with Bibby AI's citation management. Feel free to modify it for your needs!

\bibliographystyle{plain}
\bibliography{cite}

\end{document}\documentclass{article}
\usepackage[utf8]{inputenc}
\usepackage{booktabs}
\usepackage{amsmath}
\usepackage{natbib}

\title{My Document}
\author{Your Name}
\date{\today}

\begin{document}

\maketitle

\section{Introduction}
This is a sample document demonstrating various \LaTeX{} features including mathematical equations, tables, and formatting. You can cite papers using the citation search feature, and they will be automatically added to your \texttt{cite.bib} file.

\section{Mathematical Expressions}
Here are some example equations:

\subsection{Quadratic Formula}
The quadratic formula is given by:
\[x = \frac{-b \pm \sqrt{b^2 - 4ac}}{2a}\]

\subsection{Matrix Example}
A sample matrix:
\[
\begin{pmatrix}
a_{11} & a_{12} & a_{13} \\
a_{21} & a_{22} & a_{23} \\
a_{31} & a_{32} & a_{33}
\end{pmatrix}
\]

\section{Tables}
Below is a sample table showing data:

\begin{table}[h]
\centering
\begin{tabular}{@{}lrrr@{}}
\toprule
\textbf{Category} & \textbf{Value 1} & \textbf{Value 2} & \textbf{Total} \\
\midrule
Item A & 10.5 & 20.3 & 30.8 \\
Item B & 15.2 & 25.1 & 40.3 \\
Item C & 22.1 & 12.7 & 34.8 \\
\bottomrule
\end{tabular}
\caption{Sample Data Table}
\label{tab:sample}
\end{table}

\section{Lists}
\subsection{Bullet Points}
\begin{itemize}
  \item First item
  \item Second item
  \item Third item with sub-items:
    \begin{itemize}
      \item Sub-item 1
      \item Sub-item 2
    \end{itemize}
\end{itemize}

\subsection{Numbered List}
\begin{enumerate}
  \item First step
  \item Second step
  \item Third step
\end{enumerate}

\section{Citations and References}
When you use the citation search feature, references will be automatically managed. For example, you might cite a paper like this: \cite{example2024}. All citations will be automatically added to your \texttt{cite.bib} file.

\section{Conclusion}
This template demonstrates basic \LaTeX{} features including mathematical formulas, tables, and lists. The bibliography system is now set up to work seamlessly with Bibby AI's citation management. Feel free to modify it for your needs!

\bibliographystyle{plain}
\bibliography{cite}

\end{document}\documentclass{article}
\usepackage[utf8]{inputenc}
\usepackage{booktabs}
\usepackage{amsmath}
\usepackage{natbib}

\title{My Document}
\author{Your Name}
\date{\today}

\begin{document}

\maketitle

\section{Introduction}
This is a sample document demonstrating various \LaTeX{} features including mathematical equations, tables, and formatting. You can cite papers using the citation search feature, and they will be automatically added to your \texttt{cite.bib} file.

\section{Mathematical Expressions}
Here are some example equations:

\subsection{Quadratic Formula}
The quadratic formula is given by:
\[x = \frac{-b \pm \sqrt{b^2 - 4ac}}{2a}\]

\subsection{Matrix Example}
A sample matrix:
\[
\begin{pmatrix}
a_{11} & a_{12} & a_{13} \\
a_{21} & a_{22} & a_{23} \\
a_{31} & a_{32} & a_{33}
\end{pmatrix}
\]

\section{Tables}
Below is a sample table showing data:

\begin{table}[h]
\centering
\begin{tabular}{@{}lrrr@{}}
\toprule
\textbf{Category} & \textbf{Value 1} & \textbf{Value 2} & \textbf{Total} \\
\midrule
Item A & 10.5 & 20.3 & 30.8 \\
Item B & 15.2 & 25.1 & 40.3 \\
Item C & 22.1 & 12.7 & 34.8 \\
\bottomrule
\end{tabular}
\caption{Sample Data Table}
\label{tab:sample}
\end{table}

\section{Lists}
\subsection{Bullet Points}
\begin{itemize}
  \item First item
  \item Second item
  \item Third item with sub-items:
    \begin{itemize}
      \item Sub-item 1
      \item Sub-item 2
    \end{itemize}
\end{itemize}

\subsection{Numbered List}
\begin{enumerate}
  \item First step
  \item Second step
  \item Third step
\end{enumerate}

\section{Citations and References}
When you use the citation search feature, references will be automatically managed. For example, you might cite a paper like this: \cite{example2024}. All citations will be automatically added to your \texttt{cite.bib} file.

\section{Conclusion}
This template demonstrates basic \LaTeX{} features including mathematical formulas, tables, and lists. The bibliography system is now set up to work seamlessly with Bibby AI's citation management. Feel free to modify it for your needs!

\bibliographystyle{plain}
\bibliography{cite}

\end{document}\documentclass{article}
\usepackage[utf8]{inputenc}
\usepackage{booktabs}
\usepackage{amsmath}
\usepackage{natbib}

\title{My Document}
\author{Your Name}
\date{\today}

\begin{document}

\maketitle

\section{Introduction}
This is a sample document demonstrating various \LaTeX{} features including mathematical equations, tables, and formatting. You can cite papers using the citation search feature, and they will be automatically added to your \texttt{cite.bib} file.

\section{Mathematical Expressions}
Here are some example equations:

\subsection{Quadratic Formula}
The quadratic formula is given by:
\[x = \frac{-b \pm \sqrt{b^2 - 4ac}}{2a}\]

\subsection{Matrix Example}
A sample matrix:
\[
\begin{pmatrix}
a_{11} & a_{12} & a_{13} \\
a_{21} & a_{22} & a_{23} \\
a_{31} & a_{32} & a_{33}
\end{pmatrix}
\]

\section{Tables}
Below is a sample table showing data:

\begin{table}[h]
\centering
\begin{tabular}{@{}lrrr@{}}
\toprule
\textbf{Category} & \textbf{Value 1} & \textbf{Value 2} & \textbf{Total} \\
\midrule
Item A & 10.5 & 20.3 & 30.8 \\
Item B & 15.2 & 25.1 & 40.3 \\
Item C & 22.1 & 12.7 & 34.8 \\
\bottomrule
\end{tabular}
\caption{Sample Data Table}
\label{tab:sample}
\end{table}

\section{Lists}
\subsection{Bullet Points}
\begin{itemize}
  \item First item
  \item Second item
  \item Third item with sub-items:
    \begin{itemize}
      \item Sub-item 1
      \item Sub-item 2
    \end{itemize}
\end{itemize}

\subsection{Numbered List}
\begin{enumerate}
  \item First step
  \item Second step
  \item Third step
\end{enumerate}

\section{Citations and References}
When you use the citation search feature, references will be automatically managed. For example, you might cite a paper like this: \cite{example2024}. All citations will be automatically added to your \texttt{cite.bib} file.

\section{Conclusion}
This template demonstrates basic \LaTeX{} features including mathematical formulas, tables, and lists. The bibliography system is now set up to work seamlessly with Bibby AI's citation management. Feel free to modify it for your needs!

\bibliographystyle{plain}
\bibliography{cite}

\end{document}\documentclass{article}
\usepackage[utf8]{inputenc}
\usepackage{booktabs}
\usepackage{amsmath}
\usepackage{natbib}

\title{My Document}
\author{Your Name}
\date{\today}

\begin{document}

\maketitle

\section{Introduction}
This is a sample document demonstrating various \LaTeX{} features including mathematical equations, tables, and formatting. You can cite papers using the citation search feature, and they will be automatically added to your \texttt{cite.bib} file.

\section{Mathematical Expressions}
Here are some example equations:

\subsection{Quadratic Formula}
The quadratic formula is given by:
\[x = \frac{-b \pm \sqrt{b^2 - 4ac}}{2a}\]

\subsection{Matrix Example}
A sample matrix:
\[
\begin{pmatrix}
a_{11} & a_{12} & a_{13} \\
a_{21} & a_{22} & a_{23} \\
a_{31} & a_{32} & a_{33}
\end{pmatrix}
\]

\section{Tables}
Below is a sample table showing data:

\begin{table}[h]
\centering
\begin{tabular}{@{}lrrr@{}}
\toprule
\textbf{Category} & \textbf{Value 1} & \textbf{Value 2} & \textbf{Total} \\
\midrule
Item A & 10.5 & 20.3 & 30.8 \\
Item B & 15.2 & 25.1 & 40.3 \\
Item C & 22.1 & 12.7 & 34.8 \\
\bottomrule
\end{tabular}
\caption{Sample Data Table}
\label{tab:sample}
\end{table}

\section{Lists}
\subsection{Bullet Points}
\begin{itemize}
  \item First item
  \item Second item
  \item Third item with sub-items:
    \begin{itemize}
      \item Sub-item 1
      \item Sub-item 2
    \end{itemize}
\end{itemize}

\subsection{Numbered List}
\begin{enumerate}
  \item First step
  \item Second step
  \item Third step
\end{enumerate}

\section{Citations and References}
When you use the citation search feature, references will be automatically managed. For example, you might cite a paper like this: \cite{example2024}. All citations will be automatically added to your \texttt{cite.bib} file.

\section{Conclusion}
This template demonstrates basic \LaTeX{} features including mathematical formulas, tables, and lists. The bibliography system is now set up to work seamlessly with Bibby AI's citation management. Feel free to modify it for your needs!

\bibliographystyle{plain}
\bibliography{cite}

\end{document}\documentclass{article}
\usepackage[utf8]{inputenc}
\usepackage{booktabs}
\usepackage{amsmath}
\usepackage{natbib}

\title{My Document}
\author{Your Name}
\date{\today}

\begin{document}

\maketitle

\section{Introduction}
This is a sample document demonstrating various \LaTeX{} features including mathematical equations, tables, and formatting. You can cite papers using the citation search feature, and they will be automatically added to your \texttt{cite.bib} file.

\section{Mathematical Expressions}
Here are some example equations:

\subsection{Quadratic Formula}
The quadratic formula is given by:
\[x = \frac{-b \pm \sqrt{b^2 - 4ac}}{2a}\]

\subsection{Matrix Example}
A sample matrix:
\[
\begin{pmatrix}
a_{11} & a_{12} & a_{13} \\
a_{21} & a_{22} & a_{23} \\
a_{31} & a_{32} & a_{33}
\end{pmatrix}
\]

\section{Tables}
Below is a sample table showing data:

\begin{table}[h]
\centering
\begin{tabular}{@{}lrrr@{}}
\toprule
\textbf{Category} & \textbf{Value 1} & \textbf{Value 2} & \textbf{Total} \\
\midrule
Item A & 10.5 & 20.3 & 30.8 \\
Item B & 15.2 & 25.1 & 40.3 \\
Item C & 22.1 & 12.7 & 34.8 \\
\bottomrule
\end{tabular}
\caption{Sample Data Table}
\label{tab:sample}
\end{table}

\section{Lists}
\subsection{Bullet Points}
\begin{itemize}
  \item First item
  \item Second item
  \item Third item with sub-items:
    \begin{itemize}
      \item Sub-item 1
      \item Sub-item 2
    \end{itemize}
\end{itemize}

\subsection{Numbered List}
\begin{enumerate}
  \item First step
  \item Second step
  \item Third step
\end{enumerate}

\section{Citations and References}
When you use the citation search feature, references will be automatically managed. For example, you might cite a paper like this: \cite{example2024}. All citations will be automatically added to your \texttt{cite.bib} file.

\section{Conclusion}
This template demonstrates basic \LaTeX{} features including mathematical formulas, tables, and lists. The bibliography system is now set up to work seamlessly with Bibby AI's citation management. Feel free to modify it for your needs!

\bibliographystyle{plain}
\bibliography{cite}

\end{document}\documentclass{article}
\usepackage[utf8]{inputenc}
\usepackage{booktabs}
\usepackage{amsmath}
\usepackage{natbib}

\title{My Document}
\author{Your Name}
\date{\today}

\begin{document}

\maketitle

\section{Introduction}
This is a sample document demonstrating various \LaTeX{} features including mathematical equations, tables, and formatting. You can cite papers using the citation search feature, and they will be automatically added to your \texttt{cite.bib} file.

\section{Mathematical Expressions}
Here are some example equations:

\subsection{Quadratic Formula}
The quadratic formula is given by:
\[x = \frac{-b \pm \sqrt{b^2 - 4ac}}{2a}\]

\subsection{Matrix Example}
A sample matrix:
\[
\begin{pmatrix}
a_{11} & a_{12} & a_{13} \\
a_{21} & a_{22} & a_{23} \\
a_{31} & a_{32} & a_{33}
\end{pmatrix}
\]

\section{Tables}
Below is a sample table showing data:

\begin{table}[h]
\centering
\begin{tabular}{@{}lrrr@{}}
\toprule
\textbf{Category} & \textbf{Value 1} & \textbf{Value 2} & \textbf{Total} \\
\midrule
Item A & 10.5 & 20.3 & 30.8 \\
Item B & 15.2 & 25.1 & 40.3 \\
Item C & 22.1 & 12.7 & 34.8 \\
\bottomrule
\end{tabular}
\caption{Sample Data Table}
\label{tab:sample}
\end{table}

\section{Lists}
\subsection{Bullet Points}
\begin{itemize}
  \item First item
  \item Second item
  \item Third item with sub-items:
    \begin{itemize}
      \item Sub-item 1
      \item Sub-item 2
    \end{itemize}
\end{itemize}

\subsection{Numbered List}
\begin{enumerate}
  \item First step
  \item Second step
  \item Third step
\end{enumerate}

\section{Citations and References}
When you use the citation search feature, references will be automatically managed. For example, you might cite a paper like this: \cite{example2024}. All citations will be automatically added to your \texttt{cite.bib} file.

\section{Conclusion}
This template demonstrates basic \LaTeX{} features including mathematical formulas, tables, and lists. The bibliography system is now set up to work seamlessly with Bibby AI's citation management. Feel free to modify it for your needs!

\bibliographystyle{plain}
\bibliography{cite}

\end{document}\documentclass{article}
\usepackage[utf8]{inputenc}
\usepackage{booktabs}
\usepackage{amsmath}
\usepackage{natbib}

\title{My Document}
\author{Your Name}
\date{\today}

\begin{document}

\maketitle

\section{Introduction}
This is a sample document demonstrating various \LaTeX{} features including mathematical equations, tables, and formatting. You can cite papers using the citation search feature, and they will be automatically added to your \texttt{cite.bib} file.

\section{Mathematical Expressions}
Here are some example equations:

\subsection{Quadratic Formula}
The quadratic formula is given by:
\[x = \frac{-b \pm \sqrt{b^2 - 4ac}}{2a}\]

\subsection{Matrix Example}
A sample matrix:
\[
\begin{pmatrix}
a_{11} & a_{12} & a_{13} \\
a_{21} & a_{22} & a_{23} \\
a_{31} & a_{32} & a_{33}
\end{pmatrix}
\]

\section{Tables}
Below is a sample table showing data:

\begin{table}[h]
\centering
\begin{tabular}{@{}lrrr@{}}
\toprule
\textbf{Category} & \textbf{Value 1} & \textbf{Value 2} & \textbf{Total} \\
\midrule
Item A & 10.5 & 20.3 & 30.8 \\
Item B & 15.2 & 25.1 & 40.3 \\
Item C & 22.1 & 12.7 & 34.8 \\
\bottomrule
\end{tabular}
\caption{Sample Data Table}
\label{tab:sample}
\end{table}

\section{Lists}
\subsection{Bullet Points}
\begin{itemize}
  \item First item
  \item Second item
  \item Third item with sub-items:
    \begin{itemize}
      \item Sub-item 1
      \item Sub-item 2
    \end{itemize}
\end{itemize}

\subsection{Numbered List}
\begin{enumerate}
  \item First step
  \item Second step
  \item Third step
\end{enumerate}

\section{Citations and References}
When you use the citation search feature, references will be automatically managed. For example, you might cite a paper like this: \cite{example2024}. All citations will be automatically added to your \texttt{cite.bib} file.

\section{Conclusion}
This template demonstrates basic \LaTeX{} features including mathematical formulas, tables, and lists. The bibliography system is now set up to work seamlessly with Bibby AI's citation management. Feel free to modify it for your needs!

\bibliographystyle{plain}
\bibliography{cite}

\end{document}\documentclass{article}
\usepackage[utf8]{inputenc}
\usepackage{booktabs}
\usepackage{amsmath}
\usepackage{natbib}

\title{My Document}
\author{Your Name}
\date{\today}

\begin{document}

\maketitle

\section{Introduction}
This is a sample document demonstrating various \LaTeX{} features including mathematical equations, tables, and formatting. You can cite papers using the citation search feature, and they will be automatically added to your \texttt{cite.bib} file.

\section{Mathematical Expressions}
Here are some example equations:

\subsection{Quadratic Formula}
The quadratic formula is given by:
\[x = \frac{-b \pm \sqrt{b^2 - 4ac}}{2a}\]

\subsection{Matrix Example}
A sample matrix:
\[
\begin{pmatrix}
a_{11} & a_{12} & a_{13} \\
a_{21} & a_{22} & a_{23} \\
a_{31} & a_{32} & a_{33}
\end{pmatrix}
\]

\section{Tables}
Below is a sample table showing data:

\begin{table}[h]
\centering
\begin{tabular}{@{}lrrr@{}}
\toprule
\textbf{Category} & \textbf{Value 1} & \textbf{Value 2} & \textbf{Total} \\
\midrule
Item A & 10.5 & 20.3 & 30.8 \\
Item B & 15.2 & 25.1 & 40.3 \\
Item C & 22.1 & 12.7 & 34.8 \\
\bottomrule
\end{tabular}
\caption{Sample Data Table}
\label{tab:sample}
\end{table}

\section{Lists}
\subsection{Bullet Points}
\begin{itemize}
  \item First item
  \item Second item
  \item Third item with sub-items:
    \begin{itemize}
      \item Sub-item 1
      \item Sub-item 2
    \end{itemize}
\end{itemize}

\subsection{Numbered List}
\begin{enumerate}
  \item First step
  \item Second step
  \item Third step
\end{enumerate}

\section{Citations and References}
When you use the citation search feature, references will be automatically managed. For example, you might cite a paper like this: \cite{example2024}. All citations will be automatically added to your \texttt{cite.bib} file.

\section{Conclusion}
This template demonstrates basic \LaTeX{} features including mathematical formulas, tables, and lists. The bibliography system is now set up to work seamlessly with Bibby AI's citation management. Feel free to modify it for your needs!

\bibliographystyle{plain}
\bibliography{cite}

\end{document}\documentclass{article}
\usepackage[utf8]{inputenc}
\usepackage{booktabs}
\usepackage{amsmath}
\usepackage{natbib}

\title{My Document}
\author{Your Name}
\date{\today}

\begin{document}

\maketitle

\section{Introduction}
This is a sample document demonstrating various \LaTeX{} features including mathematical equations, tables, and formatting. You can cite papers using the citation search feature, and they will be automatically added to your \texttt{cite.bib} file.

\section{Mathematical Expressions}
Here are some example equations:

\subsection{Quadratic Formula}
The quadratic formula is given by:
\[x = \frac{-b \pm \sqrt{b^2 - 4ac}}{2a}\]

\subsection{Matrix Example}
A sample matrix:
\[
\begin{pmatrix}
a_{11} & a_{12} & a_{13} \\
a_{21} & a_{22} & a_{23} \\
a_{31} & a_{32} & a_{33}
\end{pmatrix}
\]

\section{Tables}
Below is a sample table showing data:

\begin{table}[h]
\centering
\begin{tabular}{@{}lrrr@{}}
\toprule
\textbf{Category} & \textbf{Value 1} & \textbf{Value 2} & \textbf{Total} \\
\midrule
Item A & 10.5 & 20.3 & 30.8 \\
Item B & 15.2 & 25.1 & 40.3 \\
Item C & 22.1 & 12.7 & 34.8 \\
\bottomrule
\end{tabular}
\caption{Sample Data Table}
\label{tab:sample}
\end{table}

\section{Lists}
\subsection{Bullet Points}
\begin{itemize}
  \item First item
  \item Second item
  \item Third item with sub-items:
    \begin{itemize}
      \item Sub-item 1
      \item Sub-item 2
    \end{itemize}
\end{itemize}

\subsection{Numbered List}
\begin{enumerate}
  \item First step
  \item Second step
  \item Third step
\end{enumerate}

\section{Citations and References}
When you use the citation search feature, references will be automatically managed. For example, you might cite a paper like this: \cite{example2024}. All citations will be automatically added to your \texttt{cite.bib} file.

\section{Conclusion}
This template demonstrates basic \LaTeX{} features including mathematical formulas, tables, and lists. The bibliography system is now set up to work seamlessly with Bibby AI's citation management. Feel free to modify it for your needs!

\bibliographystyle{plain}
\bibliography{cite}

\end{document}\documentclass{article}
\usepackage[utf8]{inputenc}
\usepackage{booktabs}
\usepackage{amsmath}
\usepackage{natbib}

\title{My Document}
\author{Your Name}
\date{\today}

\begin{document}

\maketitle

\section{Introduction}
This is a sample document demonstrating various \LaTeX{} features including mathematical equations, tables, and formatting. You can cite papers using the citation search feature, and they will be automatically added to your \texttt{cite.bib} file.

\section{Mathematical Expressions}
Here are some example equations:

\subsection{Quadratic Formula}
The quadratic formula is given by:
\[x = \frac{-b \pm \sqrt{b^2 - 4ac}}{2a}\]

\subsection{Matrix Example}
A sample matrix:
\[
\begin{pmatrix}
a_{11} & a_{12} & a_{13} \\
a_{21} & a_{22} & a_{23} \\
a_{31} & a_{32} & a_{33}
\end{pmatrix}
\]

\section{Tables}
Below is a sample table showing data:

\begin{table}[h]
\centering
\begin{tabular}{@{}lrrr@{}}
\toprule
\textbf{Category} & \textbf{Value 1} & \textbf{Value 2} & \textbf{Total} \\
\midrule
Item A & 10.5 & 20.3 & 30.8 \\
Item B & 15.2 & 25.1 & 40.3 \\
Item C & 22.1 & 12.7 & 34.8 \\
\bottomrule
\end{tabular}
\caption{Sample Data Table}
\label{tab:sample}
\end{table}

\section{Lists}
\subsection{Bullet Points}
\begin{itemize}
  \item First item
  \item Second item
  \item Third item with sub-items:
    \begin{itemize}
      \item Sub-item 1
      \item Sub-item 2
    \end{itemize}
\end{itemize}

\subsection{Numbered List}
\begin{enumerate}
  \item First step
  \item Second step
  \item Third step
\end{enumerate}

\section{Citations and References}
When you use the citation search feature, references will be automatically managed. For example, you might cite a paper like this: \cite{example2024}. All citations will be automatically added to your \texttt{cite.bib} file.

\section{Conclusion}
This template demonstrates basic \LaTeX{} features including mathematical formulas, tables, and lists. The bibliography system is now set up to work seamlessly with Bibby AI's citation management. Feel free to modify it for your needs!

\bibliographystyle{plain}
\bibliography{cite}

\end{document}\documentclass{article}
\usepackage[utf8]{inputenc}
\usepackage{booktabs}
\usepackage{amsmath}
\usepackage{natbib}

\title{My Document}
\author{Your Name}
\date{\today}

\begin{document}

\maketitle

\section{Introduction}
This is a sample document demonstrating various \LaTeX{} features including mathematical equations, tables, and formatting. You can cite papers using the citation search feature, and they will be automatically added to your \texttt{cite.bib} file.

\section{Mathematical Expressions}
Here are some example equations:

\subsection{Quadratic Formula}
The quadratic formula is given by:
\[x = \frac{-b \pm \sqrt{b^2 - 4ac}}{2a}\]

\subsection{Matrix Example}
A sample matrix:
\[
\begin{pmatrix}
a_{11} & a_{12} & a_{13} \\
a_{21} & a_{22} & a_{23} \\
a_{31} & a_{32} & a_{33}
\end{pmatrix}
\]

\section{Tables}
Below is a sample table showing data:

\begin{table}[h]
\centering
\begin{tabular}{@{}lrrr@{}}
\toprule
\textbf{Category} & \textbf{Value 1} & \textbf{Value 2} & \textbf{Total} \\
\midrule
Item A & 10.5 & 20.3 & 30.8 \\
Item B & 15.2 & 25.1 & 40.3 \\
Item C & 22.1 & 12.7 & 34.8 \\
\bottomrule
\end{tabular}
\caption{Sample Data Table}
\label{tab:sample}
\end{table}

\section{Lists}
\subsection{Bullet Points}
\begin{itemize}
  \item First item
  \item Second item
  \item Third item with sub-items:
    \begin{itemize}
      \item Sub-item 1
      \item Sub-item 2
    \end{itemize}
\end{itemize}

\subsection{Numbered List}
\begin{enumerate}
  \item First step
  \item Second step
  \item Third step
\end{enumerate}

\section{Citations and References}
When you use the citation search feature, references will be automatically managed. For example, you might cite a paper like this: \cite{example2024}. All citations will be automatically added to your \texttt{cite.bib} file.

\section{Conclusion}
This template demonstrates basic \LaTeX{} features including mathematical formulas, tables, and lists. The bibliography system is now set up to work seamlessly with Bibby AI's citation management. Feel free to modify it for your needs!

\bibliographystyle{plain}
\bibliography{cite}

\end{document}\documentclass{article}
\usepackage[utf8]{inputenc}
\usepackage{booktabs}
\usepackage{amsmath}
\usepackage{natbib}

\title{My Document}
\author{Your Name}
\date{\today}

\begin{document}

\maketitle

\section{Introduction}
This is a sample document demonstrating various \LaTeX{} features including mathematical equations, tables, and formatting. You can cite papers using the citation search feature, and they will be automatically added to your \texttt{cite.bib} file.

\section{Mathematical Expressions}
Here are some example equations:

\subsection{Quadratic Formula}
The quadratic formula is given by:
\[x = \frac{-b \pm \sqrt{b^2 - 4ac}}{2a}\]

\subsection{Matrix Example}
A sample matrix:
\[
\begin{pmatrix}
a_{11} & a_{12} & a_{13} \\
a_{21} & a_{22} & a_{23} \\
a_{31} & a_{32} & a_{33}
\end{pmatrix}
\]

\section{Tables}
Below is a sample table showing data:

\begin{table}[h]
\centering
\begin{tabular}{@{}lrrr@{}}
\toprule
\textbf{Category} & \textbf{Value 1} & \textbf{Value 2} & \textbf{Total} \\
\midrule
Item A & 10.5 & 20.3 & 30.8 \\
Item B & 15.2 & 25.1 & 40.3 \\
Item C & 22.1 & 12.7 & 34.8 \\
\bottomrule
\end{tabular}
\caption{Sample Data Table}
\label{tab:sample}
\end{table}

\section{Lists}
\subsection{Bullet Points}
\begin{itemize}
  \item First item
  \item Second item
  \item Third item with sub-items:
    \begin{itemize}
      \item Sub-item 1
      \item Sub-item 2
    \end{itemize}
\end{itemize}

\subsection{Numbered List}
\begin{enumerate}
  \item First step
  \item Second step
  \item Third step
\end{enumerate}

\section{Citations and References}
When you use the citation search feature, references will be automatically managed. For example, you might cite a paper like this: \cite{example2024}. All citations will be automatically added to your \texttt{cite.bib} file.

\section{Conclusion}
This template demonstrates basic \LaTeX{} features including mathematical formulas, tables, and lists. The bibliography system is now set up to work seamlessly with Bibby AI's citation management. Feel free to modify it for your needs!

\bibliographystyle{plain}
\bibliography{cite}

\end{document}\documentclass{article}
\usepackage[utf8]{inputenc}
\usepackage{booktabs}
\usepackage{amsmath}
\usepackage{natbib}

\title{My Document}
\author{Your Name}
\date{\today}

\begin{document}

\maketitle

\section{Introduction}
This is a sample document demonstrating various \LaTeX{} features including mathematical equations, tables, and formatting. You can cite papers using the citation search feature, and they will be automatically added to your \texttt{cite.bib} file.

\section{Mathematical Expressions}
Here are some example equations:

\subsection{Quadratic Formula}
The quadratic formula is given by:
\[x = \frac{-b \pm \sqrt{b^2 - 4ac}}{2a}\]

\subsection{Matrix Example}
A sample matrix:
\[
\begin{pmatrix}
a_{11} & a_{12} & a_{13} \\
a_{21} & a_{22} & a_{23} \\
a_{31} & a_{32} & a_{33}
\end{pmatrix}
\]

\section{Tables}
Below is a sample table showing data:

\begin{table}[h]
\centering
\begin{tabular}{@{}lrrr@{}}
\toprule
\textbf{Category} & \textbf{Value 1} & \textbf{Value 2} & \textbf{Total} \\
\midrule
Item A & 10.5 & 20.3 & 30.8 \\
Item B & 15.2 & 25.1 & 40.3 \\
Item C & 22.1 & 12.7 & 34.8 \\
\bottomrule
\end{tabular}
\caption{Sample Data Table}
\label{tab:sample}
\end{table}

\section{Lists}
\subsection{Bullet Points}
\begin{itemize}
  \item First item
  \item Second item
  \item Third item with sub-items:
    \begin{itemize}
      \item Sub-item 1
      \item Sub-item 2
    \end{itemize}
\end{itemize}

\subsection{Numbered List}
\begin{enumerate}
  \item First step
  \item Second step
  \item Third step
\end{enumerate}

\section{Citations and References}
When you use the citation search feature, references will be automatically managed. For example, you might cite a paper like this: \cite{example2024}. All citations will be automatically added to your \texttt{cite.bib} file.

\section{Conclusion}
This template demonstrates basic \LaTeX{} features including mathematical formulas, tables, and lists. The bibliography system is now set up to work seamlessly with Bibby AI's citation management. Feel free to modify it for your needs!

\bibliographystyle{plain}
\bibliography{cite}

\end{document}\documentclass{article}
\usepackage[utf8]{inputenc}
\usepackage{booktabs}
\usepackage{amsmath}
\usepackage{natbib}

\title{My Document}
\author{Your Name}
\date{\today}

\begin{document}

\maketitle

\section{Introduction}
This is a sample document demonstrating various \LaTeX{} features including mathematical equations, tables, and formatting. You can cite papers using the citation search feature, and they will be automatically added to your \texttt{cite.bib} file.

\section{Mathematical Expressions}
Here are some example equations:

\subsection{Quadratic Formula}
The quadratic formula is given by:
\[x = \frac{-b \pm \sqrt{b^2 - 4ac}}{2a}\]

\subsection{Matrix Example}
A sample matrix:
\[
\begin{pmatrix}
a_{11} & a_{12} & a_{13} \\
a_{21} & a_{22} & a_{23} \\
a_{31} & a_{32} & a_{33}
\end{pmatrix}
\]

\section{Tables}
Below is a sample table showing data:

\begin{table}[h]
\centering
\begin{tabular}{@{}lrrr@{}}
\toprule
\textbf{Category} & \textbf{Value 1} & \textbf{Value 2} & \textbf{Total} \\
\midrule
Item A & 10.5 & 20.3 & 30.8 \\
Item B & 15.2 & 25.1 & 40.3 \\
Item C & 22.1 & 12.7 & 34.8 \\
\bottomrule
\end{tabular}
\caption{Sample Data Table}
\label{tab:sample}
\end{table}

\section{Lists}
\subsection{Bullet Points}
\begin{itemize}
  \item First item
  \item Second item
  \item Third item with sub-items:
    \begin{itemize}
      \item Sub-item 1
      \item Sub-item 2
    \end{itemize}
\end{itemize}

\subsection{Numbered List}
\begin{enumerate}
  \item First step
  \item Second step
  \item Third step
\end{enumerate}

\section{Citations and References}
When you use the citation search feature, references will be automatically managed. For example, you might cite a paper like this: \cite{example2024}. All citations will be automatically added to your \texttt{cite.bib} file.

\section{Conclusion}
This template demonstrates basic \LaTeX{} features including mathematical formulas, tables, and lists. The bibliography system is now set up to work seamlessly with Bibby AI's citation management. Feel free to modify it for your needs!

\bibliographystyle{plain}
\bibliography{cite}

\end{document}\documentclass{article}
\usepackage[utf8]{inputenc}
\usepackage{booktabs}
\usepackage{amsmath}
\usepackage{natbib}

\title{My Document}
\author{Your Name}
\date{\today}

\begin{document}

\maketitle

\section{Introduction}
This is a sample document demonstrating various \LaTeX{} features including mathematical equations, tables, and formatting. You can cite papers using the citation search feature, and they will be automatically added to your \texttt{cite.bib} file.

\section{Mathematical Expressions}
Here are some example equations:

\subsection{Quadratic Formula}
The quadratic formula is given by:
\[x = \frac{-b \pm \sqrt{b^2 - 4ac}}{2a}\]

\subsection{Matrix Example}
A sample matrix:
\[
\begin{pmatrix}
a_{11} & a_{12} & a_{13} \\
a_{21} & a_{22} & a_{23} \\
a_{31} & a_{32} & a_{33}
\end{pmatrix}
\]

\section{Tables}
Below is a sample table showing data:

\begin{table}[h]
\centering
\begin{tabular}{@{}lrrr@{}}
\toprule
\textbf{Category} & \textbf{Value 1} & \textbf{Value 2} & \textbf{Total} \\
\midrule
Item A & 10.5 & 20.3 & 30.8 \\
Item B & 15.2 & 25.1 & 40.3 \\
Item C & 22.1 & 12.7 & 34.8 \\
\bottomrule
\end{tabular}
\caption{Sample Data Table}
\label{tab:sample}
\end{table}

\section{Lists}
\subsection{Bullet Points}
\begin{itemize}
  \item First item
  \item Second item
  \item Third item with sub-items:
    \begin{itemize}
      \item Sub-item 1
      \item Sub-item 2
    \end{itemize}
\end{itemize}

\subsection{Numbered List}
\begin{enumerate}
  \item First step
  \item Second step
  \item Third step
\end{enumerate}

\section{Citations and References}
When you use the citation search feature, references will be automatically managed. For example, you might cite a paper like this: \cite{example2024}. All citations will be automatically added to your \texttt{cite.bib} file.

\section{Conclusion}
This template demonstrates basic \LaTeX{} features including mathematical formulas, tables, and lists. The bibliography system is now set up to work seamlessly with Bibby AI's citation management. Feel free to modify it for your needs!

\bibliographystyle{plain}
\bibliography{cite}

\end{document}\documentclass{article}
\usepackage[utf8]{inputenc}
\usepackage{booktabs}
\usepackage{amsmath}
\usepackage{natbib}

\title{My Document}
\author{Your Name}
\date{\today}

\begin{document}

\maketitle

\section{Introduction}
This is a sample document demonstrating various \LaTeX{} features including mathematical equations, tables, and formatting. You can cite papers using the citation search feature, and they will be automatically added to your \texttt{cite.bib} file.

\section{Mathematical Expressions}
Here are some example equations:

\subsection{Quadratic Formula}
The quadratic formula is given by:
\[x = \frac{-b \pm \sqrt{b^2 - 4ac}}{2a}\]

\subsection{Matrix Example}
A sample matrix:
\[
\begin{pmatrix}
a_{11} & a_{12} & a_{13} \\
a_{21} & a_{22} & a_{23} \\
a_{31} & a_{32} & a_{33}
\end{pmatrix}
\]

\section{Tables}
Below is a sample table showing data:

\begin{table}[h]
\centering
\begin{tabular}{@{}lrrr@{}}
\toprule
\textbf{Category} & \textbf{Value 1} & \textbf{Value 2} & \textbf{Total} \\
\midrule
Item A & 10.5 & 20.3 & 30.8 \\
Item B & 15.2 & 25.1 & 40.3 \\
Item C & 22.1 & 12.7 & 34.8 \\
\bottomrule
\end{tabular}
\caption{Sample Data Table}
\label{tab:sample}
\end{table}

\section{Lists}
\subsection{Bullet Points}
\begin{itemize}
  \item First item
  \item Second item
  \item Third item with sub-items:
    \begin{itemize}
      \item Sub-item 1
      \item Sub-item 2
    \end{itemize}
\end{itemize}

\subsection{Numbered List}
\begin{enumerate}
  \item First step
  \item Second step
  \item Third step
\end{enumerate}

\section{Citations and References}
When you use the citation search feature, references will be automatically managed. For example, you might cite a paper like this: \cite{example2024}. All citations will be automatically added to your \texttt{cite.bib} file.

\section{Conclusion}
This template demonstrates basic \LaTeX{} features including mathematical formulas, tables, and lists. The bibliography system is now set up to work seamlessly with Bibby AI's citation management. Feel free to modify it for your needs!

\bibliographystyle{plain}
\bibliography{cite}

\end{document}\documentclass{article}
\usepackage[utf8]{inputenc}
\usepackage{booktabs}
\usepackage{amsmath}
\usepackage{natbib}

\title{My Document}
\author{Your Name}
\date{\today}

\begin{document}

\maketitle

\section{Introduction}
This is a sample document demonstrating various \LaTeX{} features including mathematical equations, tables, and formatting. You can cite papers using the citation search feature, and they will be automatically added to your \texttt{cite.bib} file.

\section{Mathematical Expressions}
Here are some example equations:

\subsection{Quadratic Formula}
The quadratic formula is given by:
\[x = \frac{-b \pm \sqrt{b^2 - 4ac}}{2a}\]

\subsection{Matrix Example}
A sample matrix:
\[
\begin{pmatrix}
a_{11} & a_{12} & a_{13} \\
a_{21} & a_{22} & a_{23} \\
a_{31} & a_{32} & a_{33}
\end{pmatrix}
\]

\section{Tables}
Below is a sample table showing data:

\begin{table}[h]
\centering
\begin{tabular}{@{}lrrr@{}}
\toprule
\textbf{Category} & \textbf{Value 1} & \textbf{Value 2} & \textbf{Total} \\
\midrule
Item A & 10.5 & 20.3 & 30.8 \\
Item B & 15.2 & 25.1 & 40.3 \\
Item C & 22.1 & 12.7 & 34.8 \\
\bottomrule
\end{tabular}
\caption{Sample Data Table}
\label{tab:sample}
\end{table}

\section{Lists}
\subsection{Bullet Points}
\begin{itemize}
  \item First item
  \item Second item
  \item Third item with sub-items:
    \begin{itemize}
      \item Sub-item 1
      \item Sub-item 2
    \end{itemize}
\end{itemize}

\subsection{Numbered List}
\begin{enumerate}
  \item First step
  \item Second step
  \item Third step
\end{enumerate}

\section{Citations and References}
When you use the citation search feature, references will be automatically managed. For example, you might cite a paper like this: \cite{example2024}. All citations will be automatically added to your \texttt{cite.bib} file.

\section{Conclusion}
This template demonstrates basic \LaTeX{} features including mathematical formulas, tables, and lists. The bibliography system is now set up to work seamlessly with Bibby AI's citation management. Feel free to modify it for your needs!

\bibliographystyle{plain}
\bibliography{cite}

\end{document}\documentclass{article}
\usepackage[utf8]{inputenc}
\usepackage{booktabs}
\usepackage{amsmath}
\usepackage{natbib}

\title{My Document}
\author{Your Name}
\date{\today}

\begin{document}

\maketitle

\section{Introduction}
This is a sample document demonstrating various \LaTeX{} features including mathematical equations, tables, and formatting. You can cite papers using the citation search feature, and they will be automatically added to your \texttt{cite.bib} file.

\section{Mathematical Expressions}
Here are some example equations:

\subsection{Quadratic Formula}
The quadratic formula is given by:
\[x = \frac{-b \pm \sqrt{b^2 - 4ac}}{2a}\]

\subsection{Matrix Example}
A sample matrix:
\[
\begin{pmatrix}
a_{11} & a_{12} & a_{13} \\
a_{21} & a_{22} & a_{23} \\
a_{31} & a_{32} & a_{33}
\end{pmatrix}
\]

\section{Tables}
Below is a sample table showing data:

\begin{table}[h]
\centering
\begin{tabular}{@{}lrrr@{}}
\toprule
\textbf{Category} & \textbf{Value 1} & \textbf{Value 2} & \textbf{Total} \\
\midrule
Item A & 10.5 & 20.3 & 30.8 \\
Item B & 15.2 & 25.1 & 40.3 \\
Item C & 22.1 & 12.7 & 34.8 \\
\bottomrule
\end{tabular}
\caption{Sample Data Table}
\label{tab:sample}
\end{table}

\section{Lists}
\subsection{Bullet Points}
\begin{itemize}
  \item First item
  \item Second item
  \item Third item with sub-items:
    \begin{itemize}
      \item Sub-item 1
      \item Sub-item 2
    \end{itemize}
\end{itemize}

\subsection{Numbered List}
\begin{enumerate}
  \item First step
  \item Second step
  \item Third step
\end{enumerate}

\section{Citations and References}
When you use the citation search feature, references will be automatically managed. For example, you might cite a paper like this: \cite{example2024}. All citations will be automatically added to your \texttt{cite.bib} file.

\section{Conclusion}
This template demonstrates basic \LaTeX{} features including mathematical formulas, tables, and lists. The bibliography system is now set up to work seamlessly with Bibby AI's citation management. Feel free to modify it for your needs!

\bibliographystyle{plain}
\bibliography{cite}

\end{document}\documentclass{article}
\usepackage[utf8]{inputenc}
\usepackage{booktabs}
\usepackage{amsmath}
\usepackage{natbib}

\title{My Document}
\author{Your Name}
\date{\today}

\begin{document}

\maketitle

\section{Introduction}
This is a sample document demonstrating various \LaTeX{} features including mathematical equations, tables, and formatting. You can cite papers using the citation search feature, and they will be automatically added to your \texttt{cite.bib} file.

\section{Mathematical Expressions}
Here are some example equations:

\subsection{Quadratic Formula}
The quadratic formula is given by:
\[x = \frac{-b \pm \sqrt{b^2 - 4ac}}{2a}\]

\subsection{Matrix Example}
A sample matrix:
\[
\begin{pmatrix}
a_{11} & a_{12} & a_{13} \\
a_{21} & a_{22} & a_{23} \\
a_{31} & a_{32} & a_{33}
\end{pmatrix}
\]

\section{Tables}
Below is a sample table showing data:

\begin{table}[h]
\centering
\begin{tabular}{@{}lrrr@{}}
\toprule
\textbf{Category} & \textbf{Value 1} & \textbf{Value 2} & \textbf{Total} \\
\midrule
Item A & 10.5 & 20.3 & 30.8 \\
Item B & 15.2 & 25.1 & 40.3 \\
Item C & 22.1 & 12.7 & 34.8 \\
\bottomrule
\end{tabular}
\caption{Sample Data Table}
\label{tab:sample}
\end{table}

\section{Lists}
\subsection{Bullet Points}
\begin{itemize}
  \item First item
  \item Second item
  \item Third item with sub-items:
    \begin{itemize}
      \item Sub-item 1
      \item Sub-item 2
    \end{itemize}
\end{itemize}

\subsection{Numbered List}
\begin{enumerate}
  \item First step
  \item Second step
  \item Third step
\end{enumerate}

\section{Citations and References}
When you use the citation search feature, references will be automatically managed. For example, you might cite a paper like this: \cite{example2024}. All citations will be automatically added to your \texttt{cite.bib} file.

\section{Conclusion}
This template demonstrates basic \LaTeX{} features including mathematical formulas, tables, and lists. The bibliography system is now set up to work seamlessly with Bibby AI's citation management. Feel free to modify it for your needs!

\bibliographystyle{plain}
\bibliography{cite}

\end{document}\documentclass{article}
\usepackage[utf8]{inputenc}
\usepackage{booktabs}
\usepackage{amsmath}
\usepackage{natbib}

\title{My Document}
\author{Your Name}
\date{\today}

\begin{document}

\maketitle

\section{Introduction}
This is a sample document demonstrating various \LaTeX{} features including mathematical equations, tables, and formatting. You can cite papers using the citation search feature, and they will be automatically added to your \texttt{cite.bib} file.

\section{Mathematical Expressions}
Here are some example equations:

\subsection{Quadratic Formula}
The quadratic formula is given by:
\[x = \frac{-b \pm \sqrt{b^2 - 4ac}}{2a}\]

\subsection{Matrix Example}
A sample matrix:
\[
\begin{pmatrix}
a_{11} & a_{12} & a_{13} \\
a_{21} & a_{22} & a_{23} \\
a_{31} & a_{32} & a_{33}
\end{pmatrix}
\]

\section{Tables}
Below is a sample table showing data:

\begin{table}[h]
\centering
\begin{tabular}{@{}lrrr@{}}
\toprule
\textbf{Category} & \textbf{Value 1} & \textbf{Value 2} & \textbf{Total} \\
\midrule
Item A & 10.5 & 20.3 & 30.8 \\
Item B & 15.2 & 25.1 & 40.3 \\
Item C & 22.1 & 12.7 & 34.8 \\
\bottomrule
\end{tabular}
\caption{Sample Data Table}
\label{tab:sample}
\end{table}

\section{Lists}
\subsection{Bullet Points}
\begin{itemize}
  \item First item
  \item Second item
  \item Third item with sub-items:
    \begin{itemize}
      \item Sub-item 1
      \item Sub-item 2
    \end{itemize}
\end{itemize}

\subsection{Numbered List}
\begin{enumerate}
  \item First step
  \item Second step
  \item Third step
\end{enumerate}

\section{Citations and References}
When you use the citation search feature, references will be automatically managed. For example, you might cite a paper like this: \cite{example2024}. All citations will be automatically added to your \texttt{cite.bib} file.

\section{Conclusion}
This template demonstrates basic \LaTeX{} features including mathematical formulas, tables, and lists. The bibliography system is now set up to work seamlessly with Bibby AI's citation management. Feel free to modify it for your needs!

\bibliographystyle{plain}
\bibliography{cite}

\end{document}\documentclass{article}
\usepackage[utf8]{inputenc}
\usepackage{booktabs}
\usepackage{amsmath}
\usepackage{natbib}

\title{My Document}
\author{Your Name}
\date{\today}

\begin{document}

\maketitle

\section{Introduction}
This is a sample document demonstrating various \LaTeX{} features including mathematical equations, tables, and formatting. You can cite papers using the citation search feature, and they will be automatically added to your \texttt{cite.bib} file.

\section{Mathematical Expressions}
Here are some example equations:

\subsection{Quadratic Formula}
The quadratic formula is given by:
\[x = \frac{-b \pm \sqrt{b^2 - 4ac}}{2a}\]

\subsection{Matrix Example}
A sample matrix:
\[
\begin{pmatrix}
a_{11} & a_{12} & a_{13} \\
a_{21} & a_{22} & a_{23} \\
a_{31} & a_{32} & a_{33}
\end{pmatrix}
\]

\section{Tables}
Below is a sample table showing data:

\begin{table}[h]
\centering
\begin{tabular}{@{}lrrr@{}}
\toprule
\textbf{Category} & \textbf{Value 1} & \textbf{Value 2} & \textbf{Total} \\
\midrule
Item A & 10.5 & 20.3 & 30.8 \\
Item B & 15.2 & 25.1 & 40.3 \\
Item C & 22.1 & 12.7 & 34.8 \\
\bottomrule
\end{tabular}
\caption{Sample Data Table}
\label{tab:sample}
\end{table}

\section{Lists}
\subsection{Bullet Points}
\begin{itemize}
  \item First item
  \item Second item
  \item Third item with sub-items:
    \begin{itemize}
      \item Sub-item 1
      \item Sub-item 2
    \end{itemize}
\end{itemize}

\subsection{Numbered List}
\begin{enumerate}
  \item First step
  \item Second step
  \item Third step
\end{enumerate}

\section{Citations and References}
When you use the citation search feature, references will be automatically managed. For example, you might cite a paper like this: \cite{example2024}. All citations will be automatically added to your \texttt{cite.bib} file.

\section{Conclusion}
This template demonstrates basic \LaTeX{} features including mathematical formulas, tables, and lists. The bibliography system is now set up to work seamlessly with Bibby AI's citation management. Feel free to modify it for your needs!

\bibliographystyle{plain}
\bibliography{cite}

\end{document}\documentclass{article}
\usepackage[utf8]{inputenc}
\usepackage{booktabs}
\usepackage{amsmath}
\usepackage{natbib}

\title{My Document}
\author{Your Name}
\date{\today}

\begin{document}

\maketitle

\section{Introduction}
This is a sample document demonstrating various \LaTeX{} features including mathematical equations, tables, and formatting. You can cite papers using the citation search feature, and they will be automatically added to your \texttt{cite.bib} file.

\section{Mathematical Expressions}
Here are some example equations:

\subsection{Quadratic Formula}
The quadratic formula is given by:
\[x = \frac{-b \pm \sqrt{b^2 - 4ac}}{2a}\]

\subsection{Matrix Example}
A sample matrix:
\[
\begin{pmatrix}
a_{11} & a_{12} & a_{13} \\
a_{21} & a_{22} & a_{23} \\
a_{31} & a_{32} & a_{33}
\end{pmatrix}
\]

\section{Tables}
Below is a sample table showing data:

\begin{table}[h]
\centering
\begin{tabular}{@{}lrrr@{}}
\toprule
\textbf{Category} & \textbf{Value 1} & \textbf{Value 2} & \textbf{Total} \\
\midrule
Item A & 10.5 & 20.3 & 30.8 \\
Item B & 15.2 & 25.1 & 40.3 \\
Item C & 22.1 & 12.7 & 34.8 \\
\bottomrule
\end{tabular}
\caption{Sample Data Table}
\label{tab:sample}
\end{table}

\section{Lists}
\subsection{Bullet Points}
\begin{itemize}
  \item First item
  \item Second item
  \item Third item with sub-items:
    \begin{itemize}
      \item Sub-item 1
      \item Sub-item 2
    \end{itemize}
\end{itemize}

\subsection{Numbered List}
\begin{enumerate}
  \item First step
  \item Second step
  \item Third step
\end{enumerate}

\section{Citations and References}
When you use the citation search feature, references will be automatically managed. For example, you might cite a paper like this: \cite{example2024}. All citations will be automatically added to your \texttt{cite.bib} file.

\section{Conclusion}
This template demonstrates basic \LaTeX{} features including mathematical formulas, tables, and lists. The bibliography system is now set up to work seamlessly with Bibby AI's citation management. Feel free to modify it for your needs!

\bibliographystyle{plain}
\bibliography{cite}

\end{document}\documentclass{article}
\usepackage[utf8]{inputenc}
\usepackage{booktabs}
\usepackage{amsmath}
\usepackage{natbib}

\title{My Document}
\author{Your Name}
\date{\today}

\begin{document}

\maketitle

\section{Introduction}
This is a sample document demonstrating various \LaTeX{} features including mathematical equations, tables, and formatting. You can cite papers using the citation search feature, and they will be automatically added to your \texttt{cite.bib} file.

\section{Mathematical Expressions}
Here are some example equations:

\subsection{Quadratic Formula}
The quadratic formula is given by:
\[x = \frac{-b \pm \sqrt{b^2 - 4ac}}{2a}\]

\subsection{Matrix Example}
A sample matrix:
\[
\begin{pmatrix}
a_{11} & a_{12} & a_{13} \\
a_{21} & a_{22} & a_{23} \\
a_{31} & a_{32} & a_{33}
\end{pmatrix}
\]

\section{Tables}
Below is a sample table showing data:

\begin{table}[h]
\centering
\begin{tabular}{@{}lrrr@{}}
\toprule
\textbf{Category} & \textbf{Value 1} & \textbf{Value 2} & \textbf{Total} \\
\midrule
Item A & 10.5 & 20.3 & 30.8 \\
Item B & 15.2 & 25.1 & 40.3 \\
Item C & 22.1 & 12.7 & 34.8 \\
\bottomrule
\end{tabular}
\caption{Sample Data Table}
\label{tab:sample}
\end{table}

\section{Lists}
\subsection{Bullet Points}
\begin{itemize}
  \item First item
  \item Second item
  \item Third item with sub-items:
    \begin{itemize}
      \item Sub-item 1
      \item Sub-item 2
    \end{itemize}
\end{itemize}

\subsection{Numbered List}
\begin{enumerate}
  \item First step
  \item Second step
  \item Third step
\end{enumerate}

\section{Citations and References}
When you use the citation search feature, references will be automatically managed. For example, you might cite a paper like this: \cite{example2024}. All citations will be automatically added to your \texttt{cite.bib} file.

\section{Conclusion}
This template demonstrates basic \LaTeX{} features including mathematical formulas, tables, and lists. The bibliography system is now set up to work seamlessly with Bibby AI's citation management. Feel free to modify it for your needs!

\bibliographystyle{plain}
\bibliography{cite}

\end{document}\documentclass{article}
\usepackage[utf8]{inputenc}
\usepackage{booktabs}
\usepackage{amsmath}
\usepackage{natbib}

\title{My Document}
\author{Your Name}
\date{\today}

\begin{document}

\maketitle

\section{Introduction}
This is a sample document demonstrating various \LaTeX{} features including mathematical equations, tables, and formatting. You can cite papers using the citation search feature, and they will be automatically added to your \texttt{cite.bib} file.

\section{Mathematical Expressions}
Here are some example equations:

\subsection{Quadratic Formula}
The quadratic formula is given by:
\[x = \frac{-b \pm \sqrt{b^2 - 4ac}}{2a}\]

\subsection{Matrix Example}
A sample matrix:
\[
\begin{pmatrix}
a_{11} & a_{12} & a_{13} \\
a_{21} & a_{22} & a_{23} \\
a_{31} & a_{32} & a_{33}
\end{pmatrix}
\]

\section{Tables}
Below is a sample table showing data:

\begin{table}[h]
\centering
\begin{tabular}{@{}lrrr@{}}
\toprule
\textbf{Category} & \textbf{Value 1} & \textbf{Value 2} & \textbf{Total} \\
\midrule
Item A & 10.5 & 20.3 & 30.8 \\
Item B & 15.2 & 25.1 & 40.3 \\
Item C & 22.1 & 12.7 & 34.8 \\
\bottomrule
\end{tabular}
\caption{Sample Data Table}
\label{tab:sample}
\end{table}

\section{Lists}
\subsection{Bullet Points}
\begin{itemize}
  \item First item
  \item Second item
  \item Third item with sub-items:
    \begin{itemize}
      \item Sub-item 1
      \item Sub-item 2
    \end{itemize}
\end{itemize}

\subsection{Numbered List}
\begin{enumerate}
  \item First step
  \item Second step
  \item Third step
\end{enumerate}

\section{Citations and References}
When you use the citation search feature, references will be automatically managed. For example, you might cite a paper like this: \cite{example2024}. All citations will be automatically added to your \texttt{cite.bib} file.

\section{Conclusion}
This template demonstrates basic \LaTeX{} features including mathematical formulas, tables, and lists. The bibliography system is now set up to work seamlessly with Bibby AI's citation management. Feel free to modify it for your needs!

\bibliographystyle{plain}
\bibliography{cite}

\end{document}\documentclass{article}
\usepackage[utf8]{inputenc}
\usepackage{booktabs}
\usepackage{amsmath}
\usepackage{natbib}

\title{My Document}
\author{Your Name}
\date{\today}

\begin{document}

\maketitle

\section{Introduction}
This is a sample document demonstrating various \LaTeX{} features including mathematical equations, tables, and formatting. You can cite papers using the citation search feature, and they will be automatically added to your \texttt{cite.bib} file.

\section{Mathematical Expressions}
Here are some example equations:

\subsection{Quadratic Formula}
The quadratic formula is given by:
\[x = \frac{-b \pm \sqrt{b^2 - 4ac}}{2a}\]

\subsection{Matrix Example}
A sample matrix:
\[
\begin{pmatrix}
a_{11} & a_{12} & a_{13} \\
a_{21} & a_{22} & a_{23} \\
a_{31} & a_{32} & a_{33}
\end{pmatrix}
\]

\section{Tables}
Below is a sample table showing data:

\begin{table}[h]
\centering
\begin{tabular}{@{}lrrr@{}}
\toprule
\textbf{Category} & \textbf{Value 1} & \textbf{Value 2} & \textbf{Total} \\
\midrule
Item A & 10.5 & 20.3 & 30.8 \\
Item B & 15.2 & 25.1 & 40.3 \\
Item C & 22.1 & 12.7 & 34.8 \\
\bottomrule
\end{tabular}
\caption{Sample Data Table}
\label{tab:sample}
\end{table}

\section{Lists}
\subsection{Bullet Points}
\begin{itemize}
  \item First item
  \item Second item
  \item Third item with sub-items:
    \begin{itemize}
      \item Sub-item 1
      \item Sub-item 2
    \end{itemize}
\end{itemize}

\subsection{Numbered List}
\begin{enumerate}
  \item First step
  \item Second step
  \item Third step
\end{enumerate}

\section{Citations and References}
When you use the citation search feature, references will be automatically managed. For example, you might cite a paper like this: \cite{example2024}. All citations will be automatically added to your \texttt{cite.bib} file.

\section{Conclusion}
This template demonstrates basic \LaTeX{} features including mathematical formulas, tables, and lists. The bibliography system is now set up to work seamlessly with Bibby AI's citation management. Feel free to modify it for your needs!

\bibliographystyle{plain}
\bibliography{cite}

\end{document}\documentclass{article}
\usepackage[utf8]{inputenc}
\usepackage{booktabs}
\usepackage{amsmath}
\usepackage{natbib}

\title{My Document}
\author{Your Name}
\date{\today}

\begin{document}

\maketitle

\section{Introduction}
This is a sample document demonstrating various \LaTeX{} features including mathematical equations, tables, and formatting. You can cite papers using the citation search feature, and they will be automatically added to your \texttt{cite.bib} file.

\section{Mathematical Expressions}
Here are some example equations:

\subsection{Quadratic Formula}
The quadratic formula is given by:
\[x = \frac{-b \pm \sqrt{b^2 - 4ac}}{2a}\]

\subsection{Matrix Example}
A sample matrix:
\[
\begin{pmatrix}
a_{11} & a_{12} & a_{13} \\
a_{21} & a_{22} & a_{23} \\
a_{31} & a_{32} & a_{33}
\end{pmatrix}
\]

\section{Tables}
Below is a sample table showing data:

\begin{table}[h]
\centering
\begin{tabular}{@{}lrrr@{}}
\toprule
\textbf{Category} & \textbf{Value 1} & \textbf{Value 2} & \textbf{Total} \\
\midrule
Item A & 10.5 & 20.3 & 30.8 \\
Item B & 15.2 & 25.1 & 40.3 \\
Item C & 22.1 & 12.7 & 34.8 \\
\bottomrule
\end{tabular}
\caption{Sample Data Table}
\label{tab:sample}
\end{table}

\section{Lists}
\subsection{Bullet Points}
\begin{itemize}
  \item First item
  \item Second item
  \item Third item with sub-items:
    \begin{itemize}
      \item Sub-item 1
      \item Sub-item 2
    \end{itemize}
\end{itemize}

\subsection{Numbered List}
\begin{enumerate}
  \item First step
  \item Second step
  \item Third step
\end{enumerate}

\section{Citations and References}
When you use the citation search feature, references will be automatically managed. For example, you might cite a paper like this: \cite{example2024}. All citations will be automatically added to your \texttt{cite.bib} file.

\section{Conclusion}
This template demonstrates basic \LaTeX{} features including mathematical formulas, tables, and lists. The bibliography system is now set up to work seamlessly with Bibby AI's citation management. Feel free to modify it for your needs!

\bibliographystyle{plain}
\bibliography{cite}

\end{document}\documentclass{article}
\usepackage[utf8]{inputenc}
\usepackage{booktabs}
\usepackage{amsmath}
\usepackage{natbib}

\title{My Document}
\author{Your Name}
\date{\today}

\begin{document}

\maketitle

\section{Introduction}
This is a sample document demonstrating various \LaTeX{} features including mathematical equations, tables, and formatting. You can cite papers using the citation search feature, and they will be automatically added to your \texttt{cite.bib} file.

\section{Mathematical Expressions}
Here are some example equations:

\subsection{Quadratic Formula}
The quadratic formula is given by:
\[x = \frac{-b \pm \sqrt{b^2 - 4ac}}{2a}\]

\subsection{Matrix Example}
A sample matrix:
\[
\begin{pmatrix}
a_{11} & a_{12} & a_{13} \\
a_{21} & a_{22} & a_{23} \\
a_{31} & a_{32} & a_{33}
\end{pmatrix}
\]

\section{Tables}
Below is a sample table showing data:

\begin{table}[h]
\centering
\begin{tabular}{@{}lrrr@{}}
\toprule
\textbf{Category} & \textbf{Value 1} & \textbf{Value 2} & \textbf{Total} \\
\midrule
Item A & 10.5 & 20.3 & 30.8 \\
Item B & 15.2 & 25.1 & 40.3 \\
Item C & 22.1 & 12.7 & 34.8 \\
\bottomrule
\end{tabular}
\caption{Sample Data Table}
\label{tab:sample}
\end{table}

\section{Lists}
\subsection{Bullet Points}
\begin{itemize}
  \item First item
  \item Second item
  \item Third item with sub-items:
    \begin{itemize}
      \item Sub-item 1
      \item Sub-item 2
    \end{itemize}
\end{itemize}

\subsection{Numbered List}
\begin{enumerate}
  \item First step
  \item Second step
  \item Third step
\end{enumerate}

\section{Citations and References}
When you use the citation search feature, references will be automatically managed. For example, you might cite a paper like this: \cite{example2024}. All citations will be automatically added to your \texttt{cite.bib} file.

\section{Conclusion}
This template demonstrates basic \LaTeX{} features including mathematical formulas, tables, and lists. The bibliography system is now set up to work seamlessly with Bibby AI's citation management. Feel free to modify it for your needs!

\bibliographystyle{plain}
\bibliography{cite}

\end{document}\documentclass{article}
\usepackage[utf8]{inputenc}
\usepackage{booktabs}
\usepackage{amsmath}
\usepackage{natbib}

\title{My Document}
\author{Your Name}
\date{\today}

\begin{document}

\maketitle

\section{Introduction}
This is a sample document demonstrating various \LaTeX{} features including mathematical equations, tables, and formatting. You can cite papers using the citation search feature, and they will be automatically added to your \texttt{cite.bib} file.

\section{Mathematical Expressions}
Here are some example equations:

\subsection{Quadratic Formula}
The quadratic formula is given by:
\[x = \frac{-b \pm \sqrt{b^2 - 4ac}}{2a}\]

\subsection{Matrix Example}
A sample matrix:
\[
\begin{pmatrix}
a_{11} & a_{12} & a_{13} \\
a_{21} & a_{22} & a_{23} \\
a_{31} & a_{32} & a_{33}
\end{pmatrix}
\]

\section{Tables}
Below is a sample table showing data:

\begin{table}[h]
\centering
\begin{tabular}{@{}lrrr@{}}
\toprule
\textbf{Category} & \textbf{Value 1} & \textbf{Value 2} & \textbf{Total} \\
\midrule
Item A & 10.5 & 20.3 & 30.8 \\
Item B & 15.2 & 25.1 & 40.3 \\
Item C & 22.1 & 12.7 & 34.8 \\
\bottomrule
\end{tabular}
\caption{Sample Data Table}
\label{tab:sample}
\end{table}

\section{Lists}
\subsection{Bullet Points}
\begin{itemize}
  \item First item
  \item Second item
  \item Third item with sub-items:
    \begin{itemize}
      \item Sub-item 1
      \item Sub-item 2
    \end{itemize}
\end{itemize}

\subsection{Numbered List}
\begin{enumerate}
  \item First step
  \item Second step
  \item Third step
\end{enumerate}

\section{Citations and References}
When you use the citation search feature, references will be automatically managed. For example, you might cite a paper like this: \cite{example2024}. All citations will be automatically added to your \texttt{cite.bib} file.

\section{Conclusion}
This template demonstrates basic \LaTeX{} features including mathematical formulas, tables, and lists. The bibliography system is now set up to work seamlessly with Bibby AI's citation management. Feel free to modify it for your needs!

\bibliographystyle{plain}
\bibliography{cite}

\end{document}\documentclass{article}
\usepackage[utf8]{inputenc}
\usepackage{booktabs}
\usepackage{amsmath}
\usepackage{natbib}

\title{My Document}
\author{Your Name}
\date{\today}

\begin{document}

\maketitle

\section{Introduction}
This is a sample document demonstrating various \LaTeX{} features including mathematical equations, tables, and formatting. You can cite papers using the citation search feature, and they will be automatically added to your \texttt{cite.bib} file.

\section{Mathematical Expressions}
Here are some example equations:

\subsection{Quadratic Formula}
The quadratic formula is given by:
\[x = \frac{-b \pm \sqrt{b^2 - 4ac}}{2a}\]

\subsection{Matrix Example}
A sample matrix:
\[
\begin{pmatrix}
a_{11} & a_{12} & a_{13} \\
a_{21} & a_{22} & a_{23} \\
a_{31} & a_{32} & a_{33}
\end{pmatrix}
\]

\section{Tables}
Below is a sample table showing data:

\begin{table}[h]
\centering
\begin{tabular}{@{}lrrr@{}}
\toprule
\textbf{Category} & \textbf{Value 1} & \textbf{Value 2} & \textbf{Total} \\
\midrule
Item A & 10.5 & 20.3 & 30.8 \\
Item B & 15.2 & 25.1 & 40.3 \\
Item C & 22.1 & 12.7 & 34.8 \\
\bottomrule
\end{tabular}
\caption{Sample Data Table}
\label{tab:sample}
\end{table}

\section{Lists}
\subsection{Bullet Points}
\begin{itemize}
  \item First item
  \item Second item
  \item Third item with sub-items:
    \begin{itemize}
      \item Sub-item 1
      \item Sub-item 2
    \end{itemize}
\end{itemize}

\subsection{Numbered List}
\begin{enumerate}
  \item First step
  \item Second step
  \item Third step
\end{enumerate}

\section{Citations and References}
When you use the citation search feature, references will be automatically managed. For example, you might cite a paper like this: \cite{example2024}. All citations will be automatically added to your \texttt{cite.bib} file.

\section{Conclusion}
This template demonstrates basic \LaTeX{} features including mathematical formulas, tables, and lists. The bibliography system is now set up to work seamlessly with Bibby AI's citation management. Feel free to modify it for your needs!

\bibliographystyle{plain}
\bibliography{cite}

\end{document}\documentclass{article}
\usepackage[utf8]{inputenc}
\usepackage{booktabs}
\usepackage{amsmath}
\usepackage{natbib}

\title{My Document}
\author{Your Name}
\date{\today}

\begin{document}

\maketitle

\section{Introduction}
This is a sample document demonstrating various \LaTeX{} features including mathematical equations, tables, and formatting. You can cite papers using the citation search feature, and they will be automatically added to your \texttt{cite.bib} file.

\section{Mathematical Expressions}
Here are some example equations:

\subsection{Quadratic Formula}
The quadratic formula is given by:
\[x = \frac{-b \pm \sqrt{b^2 - 4ac}}{2a}\]

\subsection{Matrix Example}
A sample matrix:
\[
\begin{pmatrix}
a_{11} & a_{12} & a_{13} \\
a_{21} & a_{22} & a_{23} \\
a_{31} & a_{32} & a_{33}
\end{pmatrix}
\]

\section{Tables}
Below is a sample table showing data:

\begin{table}[h]
\centering
\begin{tabular}{@{}lrrr@{}}
\toprule
\textbf{Category} & \textbf{Value 1} & \textbf{Value 2} & \textbf{Total} \\
\midrule
Item A & 10.5 & 20.3 & 30.8 \\
Item B & 15.2 & 25.1 & 40.3 \\
Item C & 22.1 & 12.7 & 34.8 \\
\bottomrule
\end{tabular}
\caption{Sample Data Table}
\label{tab:sample}
\end{table}

\section{Lists}
\subsection{Bullet Points}
\begin{itemize}
  \item First item
  \item Second item
  \item Third item with sub-items:
    \begin{itemize}
      \item Sub-item 1
      \item Sub-item 2
    \end{itemize}
\end{itemize}

\subsection{Numbered List}
\begin{enumerate}
  \item First step
  \item Second step
  \item Third step
\end{enumerate}

\section{Citations and References}
When you use the citation search feature, references will be automatically managed. For example, you might cite a paper like this: \cite{example2024}. All citations will be automatically added to your \texttt{cite.bib} file.

\section{Conclusion}
This template demonstrates basic \LaTeX{} features including mathematical formulas, tables, and lists. The bibliography system is now set up to work seamlessly with Bibby AI's citation management. Feel free to modify it for your needs!

\bibliographystyle{plain}
\bibliography{cite}

\end{document}\documentclass{article}
\usepackage[utf8]{inputenc}
\usepackage{booktabs}
\usepackage{amsmath}
\usepackage{natbib}

\title{My Document}
\author{Your Name}
\date{\today}

\begin{document}

\maketitle

\section{Introduction}
This is a sample document demonstrating various \LaTeX{} features including mathematical equations, tables, and formatting. You can cite papers using the citation search feature, and they will be automatically added to your \texttt{cite.bib} file.

\section{Mathematical Expressions}
Here are some example equations:

\subsection{Quadratic Formula}
The quadratic formula is given by:
\[x = \frac{-b \pm \sqrt{b^2 - 4ac}}{2a}\]

\subsection{Matrix Example}
A sample matrix:
\[
\begin{pmatrix}
a_{11} & a_{12} & a_{13} \\
a_{21} & a_{22} & a_{23} \\
a_{31} & a_{32} & a_{33}
\end{pmatrix}
\]

\section{Tables}
Below is a sample table showing data:

\begin{table}[h]
\centering
\begin{tabular}{@{}lrrr@{}}
\toprule
\textbf{Category} & \textbf{Value 1} & \textbf{Value 2} & \textbf{Total} \\
\midrule
Item A & 10.5 & 20.3 & 30.8 \\
Item B & 15.2 & 25.1 & 40.3 \\
Item C & 22.1 & 12.7 & 34.8 \\
\bottomrule
\end{tabular}
\caption{Sample Data Table}
\label{tab:sample}
\end{table}

\section{Lists}
\subsection{Bullet Points}
\begin{itemize}
  \item First item
  \item Second item
  \item Third item with sub-items:
    \begin{itemize}
      \item Sub-item 1
      \item Sub-item 2
    \end{itemize}
\end{itemize}

\subsection{Numbered List}
\begin{enumerate}
  \item First step
  \item Second step
  \item Third step
\end{enumerate}

\section{Citations and References}
When you use the citation search feature, references will be automatically managed. For example, you might cite a paper like this: \cite{example2024}. All citations will be automatically added to your \texttt{cite.bib} file.

\section{Conclusion}
This template demonstrates basic \LaTeX{} features including mathematical formulas, tables, and lists. The bibliography system is now set up to work seamlessly with Bibby AI's citation management. Feel free to modify it for your needs!

\bibliographystyle{plain}
\bibliography{cite}

\end{document}\documentclass{article}
\usepackage[utf8]{inputenc}
\usepackage{booktabs}
\usepackage{amsmath}
\usepackage{natbib}

\title{My Document}
\author{Your Name}
\date{\today}

\begin{document}

\maketitle

\section{Introduction}
This is a sample document demonstrating various \LaTeX{} features including mathematical equations, tables, and formatting. You can cite papers using the citation search feature, and they will be automatically added to your \texttt{cite.bib} file.

\section{Mathematical Expressions}
Here are some example equations:

\subsection{Quadratic Formula}
The quadratic formula is given by:
\[x = \frac{-b \pm \sqrt{b^2 - 4ac}}{2a}\]

\subsection{Matrix Example}
A sample matrix:
\[
\begin{pmatrix}
a_{11} & a_{12} & a_{13} \\
a_{21} & a_{22} & a_{23} \\
a_{31} & a_{32} & a_{33}
\end{pmatrix}
\]

\section{Tables}
Below is a sample table showing data:

\begin{table}[h]
\centering
\begin{tabular}{@{}lrrr@{}}
\toprule
\textbf{Category} & \textbf{Value 1} & \textbf{Value 2} & \textbf{Total} \\
\midrule
Item A & 10.5 & 20.3 & 30.8 \\
Item B & 15.2 & 25.1 & 40.3 \\
Item C & 22.1 & 12.7 & 34.8 \\
\bottomrule
\end{tabular}
\caption{Sample Data Table}
\label{tab:sample}
\end{table}

\section{Lists}
\subsection{Bullet Points}
\begin{itemize}
  \item First item
  \item Second item
  \item Third item with sub-items:
    \begin{itemize}
      \item Sub-item 1
      \item Sub-item 2
    \end{itemize}
\end{itemize}

\subsection{Numbered List}
\begin{enumerate}
  \item First step
  \item Second step
  \item Third step
\end{enumerate}

\section{Citations and References}
When you use the citation search feature, references will be automatically managed. For example, you might cite a paper like this: \cite{example2024}. All citations will be automatically added to your \texttt{cite.bib} file.

\section{Conclusion}
This template demonstrates basic \LaTeX{} features including mathematical formulas, tables, and lists. The bibliography system is now set up to work seamlessly with Bibby AI's citation management. Feel free to modify it for your needs!

\bibliographystyle{plain}
\bibliography{cite}

\end{document}\documentclass{article}
\usepackage[utf8]{inputenc}
\usepackage{booktabs}
\usepackage{amsmath}
\usepackage{natbib}

\title{My Document}
\author{Your Name}
\date{\today}

\begin{document}

\maketitle

\section{Introduction}
This is a sample document demonstrating various \LaTeX{} features including mathematical equations, tables, and formatting. You can cite papers using the citation search feature, and they will be automatically added to your \texttt{cite.bib} file.

\section{Mathematical Expressions}
Here are some example equations:

\subsection{Quadratic Formula}
The quadratic formula is given by:
\[x = \frac{-b \pm \sqrt{b^2 - 4ac}}{2a}\]

\subsection{Matrix Example}
A sample matrix:
\[
\begin{pmatrix}
a_{11} & a_{12} & a_{13} \\
a_{21} & a_{22} & a_{23} \\
a_{31} & a_{32} & a_{33}
\end{pmatrix}
\]

\section{Tables}
Below is a sample table showing data:

\begin{table}[h]
\centering
\begin{tabular}{@{}lrrr@{}}
\toprule
\textbf{Category} & \textbf{Value 1} & \textbf{Value 2} & \textbf{Total} \\
\midrule
Item A & 10.5 & 20.3 & 30.8 \\
Item B & 15.2 & 25.1 & 40.3 \\
Item C & 22.1 & 12.7 & 34.8 \\
\bottomrule
\end{tabular}
\caption{Sample Data Table}
\label{tab:sample}
\end{table}

\section{Lists}
\subsection{Bullet Points}
\begin{itemize}
  \item First item
  \item Second item
  \item Third item with sub-items:
    \begin{itemize}
      \item Sub-item 1
      \item Sub-item 2
    \end{itemize}
\end{itemize}

\subsection{Numbered List}
\begin{enumerate}
  \item First step
  \item Second step
  \item Third step
\end{enumerate}

\section{Citations and References}
When you use the citation search feature, references will be automatically managed. For example, you might cite a paper like this: \cite{example2024}. All citations will be automatically added to your \texttt{cite.bib} file.

\section{Conclusion}
This template demonstrates basic \LaTeX{} features including mathematical formulas, tables, and lists. The bibliography system is now set up to work seamlessly with Bibby AI's citation management. Feel free to modify it for your needs!

\bibliographystyle{plain}
\bibliography{cite}

\end{document}\documentclass{article}
\usepackage[utf8]{inputenc}
\usepackage{booktabs}
\usepackage{amsmath}
\usepackage{natbib}

\title{My Document}
\author{Your Name}
\date{\today}

\begin{document}

\maketitle

\section{Introduction}
This is a sample document demonstrating various \LaTeX{} features including mathematical equations, tables, and formatting. You can cite papers using the citation search feature, and they will be automatically added to your \texttt{cite.bib} file.

\section{Mathematical Expressions}
Here are some example equations:

\subsection{Quadratic Formula}
The quadratic formula is given by:
\[x = \frac{-b \pm \sqrt{b^2 - 4ac}}{2a}\]

\subsection{Matrix Example}
A sample matrix:
\[
\begin{pmatrix}
a_{11} & a_{12} & a_{13} \\
a_{21} & a_{22} & a_{23} \\
a_{31} & a_{32} & a_{33}
\end{pmatrix}
\]

\section{Tables}
Below is a sample table showing data:

\begin{table}[h]
\centering
\begin{tabular}{@{}lrrr@{}}
\toprule
\textbf{Category} & \textbf{Value 1} & \textbf{Value 2} & \textbf{Total} \\
\midrule
Item A & 10.5 & 20.3 & 30.8 \\
Item B & 15.2 & 25.1 & 40.3 \\
Item C & 22.1 & 12.7 & 34.8 \\
\bottomrule
\end{tabular}
\caption{Sample Data Table}
\label{tab:sample}
\end{table}

\section{Lists}
\subsection{Bullet Points}
\begin{itemize}
  \item First item
  \item Second item
  \item Third item with sub-items:
    \begin{itemize}
      \item Sub-item 1
      \item Sub-item 2
    \end{itemize}
\end{itemize}

\subsection{Numbered List}
\begin{enumerate}
  \item First step
  \item Second step
  \item Third step
\end{enumerate}

\section{Citations and References}
When you use the citation search feature, references will be automatically managed. For example, you might cite a paper like this: \cite{example2024}. All citations will be automatically added to your \texttt{cite.bib} file.

\section{Conclusion}
This template demonstrates basic \LaTeX{} features including mathematical formulas, tables, and lists. The bibliography system is now set up to work seamlessly with Bibby AI's citation management. Feel free to modify it for your needs!

\bibliographystyle{plain}
\bibliography{cite}

\end{document}\documentclass{article}
\usepackage[utf8]{inputenc}
\usepackage{booktabs}
\usepackage{amsmath}
\usepackage{natbib}

\title{My Document}
\author{Your Name}
\date{\today}

\begin{document}

\maketitle

\section{Introduction}
This is a sample document demonstrating various \LaTeX{} features including mathematical equations, tables, and formatting. You can cite papers using the citation search feature, and they will be automatically added to your \texttt{cite.bib} file.

\section{Mathematical Expressions}
Here are some example equations:

\subsection{Quadratic Formula}
The quadratic formula is given by:
\[x = \frac{-b \pm \sqrt{b^2 - 4ac}}{2a}\]

\subsection{Matrix Example}
A sample matrix:
\[
\begin{pmatrix}
a_{11} & a_{12} & a_{13} \\
a_{21} & a_{22} & a_{23} \\
a_{31} & a_{32} & a_{33}
\end{pmatrix}
\]

\section{Tables}
Below is a sample table showing data:

\begin{table}[h]
\centering
\begin{tabular}{@{}lrrr@{}}
\toprule
\textbf{Category} & \textbf{Value 1} & \textbf{Value 2} & \textbf{Total} \\
\midrule
Item A & 10.5 & 20.3 & 30.8 \\
Item B & 15.2 & 25.1 & 40.3 \\
Item C & 22.1 & 12.7 & 34.8 \\
\bottomrule
\end{tabular}
\caption{Sample Data Table}
\label{tab:sample}
\end{table}

\section{Lists}
\subsection{Bullet Points}
\begin{itemize}
  \item First item
  \item Second item
  \item Third item with sub-items:
    \begin{itemize}
      \item Sub-item 1
      \item Sub-item 2
    \end{itemize}
\end{itemize}

\subsection{Numbered List}
\begin{enumerate}
  \item First step
  \item Second step
  \item Third step
\end{enumerate}

\section{Citations and References}
When you use the citation search feature, references will be automatically managed. For example, you might cite a paper like this: \cite{example2024}. All citations will be automatically added to your \texttt{cite.bib} file.

\section{Conclusion}
This template demonstrates basic \LaTeX{} features including mathematical formulas, tables, and lists. The bibliography system is now set up to work seamlessly with Bibby AI's citation management. Feel free to modify it for your needs!

\bibliographystyle{plain}
\bibliography{cite}

\end{document}\documentclass{article}
\usepackage[utf8]{inputenc}
\usepackage{booktabs}
\usepackage{amsmath}
\usepackage{natbib}

\title{My Document}
\author{Your Name}
\date{\today}

\begin{document}

\maketitle

\section{Introduction}
This is a sample document demonstrating various \LaTeX{} features including mathematical equations, tables, and formatting. You can cite papers using the citation search feature, and they will be automatically added to your \texttt{cite.bib} file.

\section{Mathematical Expressions}
Here are some example equations:

\subsection{Quadratic Formula}
The quadratic formula is given by:
\[x = \frac{-b \pm \sqrt{b^2 - 4ac}}{2a}\]

\subsection{Matrix Example}
A sample matrix:
\[
\begin{pmatrix}
a_{11} & a_{12} & a_{13} \\
a_{21} & a_{22} & a_{23} \\
a_{31} & a_{32} & a_{33}
\end{pmatrix}
\]

\section{Tables}
Below is a sample table showing data:

\begin{table}[h]
\centering
\begin{tabular}{@{}lrrr@{}}
\toprule
\textbf{Category} & \textbf{Value 1} & \textbf{Value 2} & \textbf{Total} \\
\midrule
Item A & 10.5 & 20.3 & 30.8 \\
Item B & 15.2 & 25.1 & 40.3 \\
Item C & 22.1 & 12.7 & 34.8 \\
\bottomrule
\end{tabular}
\caption{Sample Data Table}
\label{tab:sample}
\end{table}

\section{Lists}
\subsection{Bullet Points}
\begin{itemize}
  \item First item
  \item Second item
  \item Third item with sub-items:
    \begin{itemize}
      \item Sub-item 1
      \item Sub-item 2
    \end{itemize}
\end{itemize}

\subsection{Numbered List}
\begin{enumerate}
  \item First step
  \item Second step
  \item Third step
\end{enumerate}

\section{Citations and References}
When you use the citation search feature, references will be automatically managed. For example, you might cite a paper like this: \cite{example2024}. All citations will be automatically added to your \texttt{cite.bib} file.

\section{Conclusion}
This template demonstrates basic \LaTeX{} features including mathematical formulas, tables, and lists. The bibliography system is now set up to work seamlessly with Bibby AI's citation management. Feel free to modify it for your needs!

\bibliographystyle{plain}
\bibliography{cite}

\end{document}\documentclass{article}
\usepackage[utf8]{inputenc}
\usepackage{booktabs}
\usepackage{amsmath}
\usepackage{natbib}

\title{My Document}
\author{Your Name}
\date{\today}

\begin{document}

\maketitle

\section{Introduction}
This is a sample document demonstrating various \LaTeX{} features including mathematical equations, tables, and formatting. You can cite papers using the citation search feature, and they will be automatically added to your \texttt{cite.bib} file.

\section{Mathematical Expressions}
Here are some example equations:

\subsection{Quadratic Formula}
The quadratic formula is given by:
\[x = \frac{-b \pm \sqrt{b^2 - 4ac}}{2a}\]

\subsection{Matrix Example}
A sample matrix:
\[
\begin{pmatrix}
a_{11} & a_{12} & a_{13} \\
a_{21} & a_{22} & a_{23} \\
a_{31} & a_{32} & a_{33}
\end{pmatrix}
\]

\section{Tables}
Below is a sample table showing data:

\begin{table}[h]
\centering
\begin{tabular}{@{}lrrr@{}}
\toprule
\textbf{Category} & \textbf{Value 1} & \textbf{Value 2} & \textbf{Total} \\
\midrule
Item A & 10.5 & 20.3 & 30.8 \\
Item B & 15.2 & 25.1 & 40.3 \\
Item C & 22.1 & 12.7 & 34.8 \\
\bottomrule
\end{tabular}
\caption{Sample Data Table}
\label{tab:sample}
\end{table}

\section{Lists}
\subsection{Bullet Points}
\begin{itemize}
  \item First item
  \item Second item
  \item Third item with sub-items:
    \begin{itemize}
      \item Sub-item 1
      \item Sub-item 2
    \end{itemize}
\end{itemize}

\subsection{Numbered List}
\begin{enumerate}
  \item First step
  \item Second step
  \item Third step
\end{enumerate}

\section{Citations and References}
When you use the citation search feature, references will be automatically managed. For example, you might cite a paper like this: \cite{example2024}. All citations will be automatically added to your \texttt{cite.bib} file.

\section{Conclusion}
This template demonstrates basic \LaTeX{} features including mathematical formulas, tables, and lists. The bibliography system is now set up to work seamlessly with Bibby AI's citation management. Feel free to modify it for your needs!

\bibliographystyle{plain}
\bibliography{cite}

\end{document}\documentclass{article}
\usepackage[utf8]{inputenc}
\usepackage{booktabs}
\usepackage{amsmath}
\usepackage{natbib}

\title{My Document}
\author{Your Name}
\date{\today}

\begin{document}

\maketitle

\section{Introduction}
This is a sample document demonstrating various \LaTeX{} features including mathematical equations, tables, and formatting. You can cite papers using the citation search feature, and they will be automatically added to your \texttt{cite.bib} file.

\section{Mathematical Expressions}
Here are some example equations:

\subsection{Quadratic Formula}
The quadratic formula is given by:
\[x = \frac{-b \pm \sqrt{b^2 - 4ac}}{2a}\]

\subsection{Matrix Example}
A sample matrix:
\[
\begin{pmatrix}
a_{11} & a_{12} & a_{13} \\
a_{21} & a_{22} & a_{23} \\
a_{31} & a_{32} & a_{33}
\end{pmatrix}
\]

\section{Tables}
Below is a sample table showing data:

\begin{table}[h]
\centering
\begin{tabular}{@{}lrrr@{}}
\toprule
\textbf{Category} & \textbf{Value 1} & \textbf{Value 2} & \textbf{Total} \\
\midrule
Item A & 10.5 & 20.3 & 30.8 \\
Item B & 15.2 & 25.1 & 40.3 \\
Item C & 22.1 & 12.7 & 34.8 \\
\bottomrule
\end{tabular}
\caption{Sample Data Table}
\label{tab:sample}
\end{table}

\section{Lists}
\subsection{Bullet Points}
\begin{itemize}
  \item First item
  \item Second item
  \item Third item with sub-items:
    \begin{itemize}
      \item Sub-item 1
      \item Sub-item 2
    \end{itemize}
\end{itemize}

\subsection{Numbered List}
\begin{enumerate}
  \item First step
  \item Second step
  \item Third step
\end{enumerate}

\section{Citations and References}
When you use the citation search feature, references will be automatically managed. For example, you might cite a paper like this: \cite{example2024}. All citations will be automatically added to your \texttt{cite.bib} file.

\section{Conclusion}
This template demonstrates basic \LaTeX{} features including mathematical formulas, tables, and lists. The bibliography system is now set up to work seamlessly with Bibby AI's citation management. Feel free to modify it for your needs!

\bibliographystyle{plain}
\bibliography{cite}

\end{document}\documentclass{article}
\usepackage[utf8]{inputenc}
\usepackage{booktabs}
\usepackage{amsmath}
\usepackage{natbib}

\title{My Document}
\author{Your Name}
\date{\today}

\begin{document}

\maketitle

\section{Introduction}
This is a sample document demonstrating various \LaTeX{} features including mathematical equations, tables, and formatting. You can cite papers using the citation search feature, and they will be automatically added to your \texttt{cite.bib} file.

\section{Mathematical Expressions}
Here are some example equations:

\subsection{Quadratic Formula}
The quadratic formula is given by:
\[x = \frac{-b \pm \sqrt{b^2 - 4ac}}{2a}\]

\subsection{Matrix Example}
A sample matrix:
\[
\begin{pmatrix}
a_{11} & a_{12} & a_{13} \\
a_{21} & a_{22} & a_{23} \\
a_{31} & a_{32} & a_{33}
\end{pmatrix}
\]

\section{Tables}
Below is a sample table showing data:

\begin{table}[h]
\centering
\begin{tabular}{@{}lrrr@{}}
\toprule
\textbf{Category} & \textbf{Value 1} & \textbf{Value 2} & \textbf{Total} \\
\midrule
Item A & 10.5 & 20.3 & 30.8 \\
Item B & 15.2 & 25.1 & 40.3 \\
Item C & 22.1 & 12.7 & 34.8 \\
\bottomrule
\end{tabular}
\caption{Sample Data Table}
\label{tab:sample}
\end{table}

\section{Lists}
\subsection{Bullet Points}
\begin{itemize}
  \item First item
  \item Second item
  \item Third item with sub-items:
    \begin{itemize}
      \item Sub-item 1
      \item Sub-item 2
    \end{itemize}
\end{itemize}

\subsection{Numbered List}
\begin{enumerate}
  \item First step
  \item Second step
  \item Third step
\end{enumerate}

\section{Citations and References}
When you use the citation search feature, references will be automatically managed. For example, you might cite a paper like this: \cite{example2024}. All citations will be automatically added to your \texttt{cite.bib} file.

\section{Conclusion}
This template demonstrates basic \LaTeX{} features including mathematical formulas, tables, and lists. The bibliography system is now set up to work seamlessly with Bibby AI's citation management. Feel free to modify it for your needs!

\bibliographystyle{plain}
\bibliography{cite}

\end{document}\documentclass{article}
\usepackage[utf8]{inputenc}
\usepackage{booktabs}
\usepackage{amsmath}
\usepackage{natbib}

\title{My Document}
\author{Your Name}
\date{\today}

\begin{document}

\maketitle

\section{Introduction}
This is a sample document demonstrating various \LaTeX{} features including mathematical equations, tables, and formatting. You can cite papers using the citation search feature, and they will be automatically added to your \texttt{cite.bib} file.

\section{Mathematical Expressions}
Here are some example equations:

\subsection{Quadratic Formula}
The quadratic formula is given by:
\[x = \frac{-b \pm \sqrt{b^2 - 4ac}}{2a}\]

\subsection{Matrix Example}
A sample matrix:
\[
\begin{pmatrix}
a_{11} & a_{12} & a_{13} \\
a_{21} & a_{22} & a_{23} \\
a_{31} & a_{32} & a_{33}
\end{pmatrix}
\]

\section{Tables}
Below is a sample table showing data:

\begin{table}[h]
\centering
\begin{tabular}{@{}lrrr@{}}
\toprule
\textbf{Category} & \textbf{Value 1} & \textbf{Value 2} & \textbf{Total} \\
\midrule
Item A & 10.5 & 20.3 & 30.8 \\
Item B & 15.2 & 25.1 & 40.3 \\
Item C & 22.1 & 12.7 & 34.8 \\
\bottomrule
\end{tabular}
\caption{Sample Data Table}
\label{tab:sample}
\end{table}

\section{Lists}
\subsection{Bullet Points}
\begin{itemize}
  \item First item
  \item Second item
  \item Third item with sub-items:
    \begin{itemize}
      \item Sub-item 1
      \item Sub-item 2
    \end{itemize}
\end{itemize}

\subsection{Numbered List}
\begin{enumerate}
  \item First step
  \item Second step
  \item Third step
\end{enumerate}

\section{Citations and References}
When you use the citation search feature, references will be automatically managed. For example, you might cite a paper like this: \cite{example2024}. All citations will be automatically added to your \texttt{cite.bib} file.

\section{Conclusion}
This template demonstrates basic \LaTeX{} features including mathematical formulas, tables, and lists. The bibliography system is now set up to work seamlessly with Bibby AI's citation management. Feel free to modify it for your needs!

\bibliographystyle{plain}
\bibliography{cite}

\end{document}\documentclass{article}
\usepackage[utf8]{inputenc}
\usepackage{booktabs}
\usepackage{amsmath}
\usepackage{natbib}

\title{My Document}
\author{Your Name}
\date{\today}

\begin{document}

\maketitle

\section{Introduction}
This is a sample document demonstrating various \LaTeX{} features including mathematical equations, tables, and formatting. You can cite papers using the citation search feature, and they will be automatically added to your \texttt{cite.bib} file.

\section{Mathematical Expressions}
Here are some example equations:

\subsection{Quadratic Formula}
The quadratic formula is given by:
\[x = \frac{-b \pm \sqrt{b^2 - 4ac}}{2a}\]

\subsection{Matrix Example}
A sample matrix:
\[
\begin{pmatrix}
a_{11} & a_{12} & a_{13} \\
a_{21} & a_{22} & a_{23} \\
a_{31} & a_{32} & a_{33}
\end{pmatrix}
\]

\section{Tables}
Below is a sample table showing data:

\begin{table}[h]
\centering
\begin{tabular}{@{}lrrr@{}}
\toprule
\textbf{Category} & \textbf{Value 1} & \textbf{Value 2} & \textbf{Total} \\
\midrule
Item A & 10.5 & 20.3 & 30.8 \\
Item B & 15.2 & 25.1 & 40.3 \\
Item C & 22.1 & 12.7 & 34.8 \\
\bottomrule
\end{tabular}
\caption{Sample Data Table}
\label{tab:sample}
\end{table}

\section{Lists}
\subsection{Bullet Points}
\begin{itemize}
  \item First item
  \item Second item
  \item Third item with sub-items:
    \begin{itemize}
      \item Sub-item 1
      \item Sub-item 2
    \end{itemize}
\end{itemize}

\subsection{Numbered List}
\begin{enumerate}
  \item First step
  \item Second step
  \item Third step
\end{enumerate}

\section{Citations and References}
When you use the citation search feature, references will be automatically managed. For example, you might cite a paper like this: \cite{example2024}. All citations will be automatically added to your \texttt{cite.bib} file.

\section{Conclusion}
This template demonstrates basic \LaTeX{} features including mathematical formulas, tables, and lists. The bibliography system is now set up to work seamlessly with Bibby AI's citation management. Feel free to modify it for your needs!

\bibliographystyle{plain}
\bibliography{cite}

\end{document}\documentclass{article}
\usepackage[utf8]{inputenc}
\usepackage{booktabs}
\usepackage{amsmath}
\usepackage{natbib}

\title{My Document}
\author{Your Name}
\date{\today}

\begin{document}

\maketitle

\section{Introduction}
This is a sample document demonstrating various \LaTeX{} features including mathematical equations, tables, and formatting. You can cite papers using the citation search feature, and they will be automatically added to your \texttt{cite.bib} file.

\section{Mathematical Expressions}
Here are some example equations:

\subsection{Quadratic Formula}
The quadratic formula is given by:
\[x = \frac{-b \pm \sqrt{b^2 - 4ac}}{2a}\]

\subsection{Matrix Example}
A sample matrix:
\[
\begin{pmatrix}
a_{11} & a_{12} & a_{13} \\
a_{21} & a_{22} & a_{23} \\
a_{31} & a_{32} & a_{33}
\end{pmatrix}
\]

\section{Tables}
Below is a sample table showing data:

\begin{table}[h]
\centering
\begin{tabular}{@{}lrrr@{}}
\toprule
\textbf{Category} & \textbf{Value 1} & \textbf{Value 2} & \textbf{Total} \\
\midrule
Item A & 10.5 & 20.3 & 30.8 \\
Item B & 15.2 & 25.1 & 40.3 \\
Item C & 22.1 & 12.7 & 34.8 \\
\bottomrule
\end{tabular}
\caption{Sample Data Table}
\label{tab:sample}
\end{table}

\section{Lists}
\subsection{Bullet Points}
\begin{itemize}
  \item First item
  \item Second item
  \item Third item with sub-items:
    \begin{itemize}
      \item Sub-item 1
      \item Sub-item 2
    \end{itemize}
\end{itemize}

\subsection{Numbered List}
\begin{enumerate}
  \item First step
  \item Second step
  \item Third step
\end{enumerate}

\section{Citations and References}
When you use the citation search feature, references will be automatically managed. For example, you might cite a paper like this: \cite{example2024}. All citations will be automatically added to your \texttt{cite.bib} file.

\section{Conclusion}
This template demonstrates basic \LaTeX{} features including mathematical formulas, tables, and lists. The bibliography system is now set up to work seamlessly with Bibby AI's citation management. Feel free to modify it for your needs!

\bibliographystyle{plain}
\bibliography{cite}

\end{document}\documentclass{article}
\usepackage[utf8]{inputenc}
\usepackage{booktabs}
\usepackage{amsmath}
\usepackage{natbib}

\title{My Document}
\author{Your Name}
\date{\today}

\begin{document}

\maketitle

\section{Introduction}
This is a sample document demonstrating various \LaTeX{} features including mathematical equations, tables, and formatting. You can cite papers using the citation search feature, and they will be automatically added to your \texttt{cite.bib} file.

\section{Mathematical Expressions}
Here are some example equations:

\subsection{Quadratic Formula}
The quadratic formula is given by:
\[x = \frac{-b \pm \sqrt{b^2 - 4ac}}{2a}\]

\subsection{Matrix Example}
A sample matrix:
\[
\begin{pmatrix}
a_{11} & a_{12} & a_{13} \\
a_{21} & a_{22} & a_{23} \\
a_{31} & a_{32} & a_{33}
\end{pmatrix}
\]

\section{Tables}
Below is a sample table showing data:

\begin{table}[h]
\centering
\begin{tabular}{@{}lrrr@{}}
\toprule
\textbf{Category} & \textbf{Value 1} & \textbf{Value 2} & \textbf{Total} \\
\midrule
Item A & 10.5 & 20.3 & 30.8 \\
Item B & 15.2 & 25.1 & 40.3 \\
Item C & 22.1 & 12.7 & 34.8 \\
\bottomrule
\end{tabular}
\caption{Sample Data Table}
\label{tab:sample}
\end{table}

\section{Lists}
\subsection{Bullet Points}
\begin{itemize}
  \item First item
  \item Second item
  \item Third item with sub-items:
    \begin{itemize}
      \item Sub-item 1
      \item Sub-item 2
    \end{itemize}
\end{itemize}

\subsection{Numbered List}
\begin{enumerate}
  \item First step
  \item Second step
  \item Third step
\end{enumerate}

\section{Citations and References}
When you use the citation search feature, references will be automatically managed. For example, you might cite a paper like this: \cite{example2024}. All citations will be automatically added to your \texttt{cite.bib} file.

\section{Conclusion}
This template demonstrates basic \LaTeX{} features including mathematical formulas, tables, and lists. The bibliography system is now set up to work seamlessly with Bibby AI's citation management. Feel free to modify it for your needs!

\bibliographystyle{plain}
\bibliography{cite}

\end{document}\documentclass{article}
\usepackage[utf8]{inputenc}
\usepackage{booktabs}
\usepackage{amsmath}
\usepackage{natbib}

\title{My Document}
\author{Your Name}
\date{\today}

\begin{document}

\maketitle

\section{Introduction}
This is a sample document demonstrating various \LaTeX{} features including mathematical equations, tables, and formatting. You can cite papers using the citation search feature, and they will be automatically added to your \texttt{cite.bib} file.

\section{Mathematical Expressions}
Here are some example equations:

\subsection{Quadratic Formula}
The quadratic formula is given by:
\[x = \frac{-b \pm \sqrt{b^2 - 4ac}}{2a}\]

\subsection{Matrix Example}
A sample matrix:
\[
\begin{pmatrix}
a_{11} & a_{12} & a_{13} \\
a_{21} & a_{22} & a_{23} \\
a_{31} & a_{32} & a_{33}
\end{pmatrix}
\]

\section{Tables}
Below is a sample table showing data:

\begin{table}[h]
\centering
\begin{tabular}{@{}lrrr@{}}
\toprule
\textbf{Category} & \textbf{Value 1} & \textbf{Value 2} & \textbf{Total} \\
\midrule
Item A & 10.5 & 20.3 & 30.8 \\
Item B & 15.2 & 25.1 & 40.3 \\
Item C & 22.1 & 12.7 & 34.8 \\
\bottomrule
\end{tabular}
\caption{Sample Data Table}
\label{tab:sample}
\end{table}

\section{Lists}
\subsection{Bullet Points}
\begin{itemize}
  \item First item
  \item Second item
  \item Third item with sub-items:
    \begin{itemize}
      \item Sub-item 1
      \item Sub-item 2
    \end{itemize}
\end{itemize}

\subsection{Numbered List}
\begin{enumerate}
  \item First step
  \item Second step
  \item Third step
\end{enumerate}

\section{Citations and References}
When you use the citation search feature, references will be automatically managed. For example, you might cite a paper like this: \cite{example2024}. All citations will be automatically added to your \texttt{cite.bib} file.

\section{Conclusion}
This template demonstrates basic \LaTeX{} features including mathematical formulas, tables, and lists. The bibliography system is now set up to work seamlessly with Bibby AI's citation management. Feel free to modify it for your needs!

\bibliographystyle{plain}
\bibliography{cite}

\end{document}\documentclass{article}
\usepackage[utf8]{inputenc}
\usepackage{booktabs}
\usepackage{amsmath}
\usepackage{natbib}

\title{My Document}
\author{Your Name}
\date{\today}

\begin{document}

\maketitle

\section{Introduction}
This is a sample document demonstrating various \LaTeX{} features including mathematical equations, tables, and formatting. You can cite papers using the citation search feature, and they will be automatically added to your \texttt{cite.bib} file.

\section{Mathematical Expressions}
Here are some example equations:

\subsection{Quadratic Formula}
The quadratic formula is given by:
\[x = \frac{-b \pm \sqrt{b^2 - 4ac}}{2a}\]

\subsection{Matrix Example}
A sample matrix:
\[
\begin{pmatrix}
a_{11} & a_{12} & a_{13} \\
a_{21} & a_{22} & a_{23} \\
a_{31} & a_{32} & a_{33}
\end{pmatrix}
\]

\section{Tables}
Below is a sample table showing data:

\begin{table}[h]
\centering
\begin{tabular}{@{}lrrr@{}}
\toprule
\textbf{Category} & \textbf{Value 1} & \textbf{Value 2} & \textbf{Total} \\
\midrule
Item A & 10.5 & 20.3 & 30.8 \\
Item B & 15.2 & 25.1 & 40.3 \\
Item C & 22.1 & 12.7 & 34.8 \\
\bottomrule
\end{tabular}
\caption{Sample Data Table}
\label{tab:sample}
\end{table}

\section{Lists}
\subsection{Bullet Points}
\begin{itemize}
  \item First item
  \item Second item
  \item Third item with sub-items:
    \begin{itemize}
      \item Sub-item 1
      \item Sub-item 2
    \end{itemize}
\end{itemize}

\subsection{Numbered List}
\begin{enumerate}
  \item First step
  \item Second step
  \item Third step
\end{enumerate}

\section{Citations and References}
When you use the citation search feature, references will be automatically managed. For example, you might cite a paper like this: \cite{example2024}. All citations will be automatically added to your \texttt{cite.bib} file.

\section{Conclusion}
This template demonstrates basic \LaTeX{} features including mathematical formulas, tables, and lists. The bibliography system is now set up to work seamlessly with Bibby AI's citation management. Feel free to modify it for your needs!

\bibliographystyle{plain}
\bibliography{cite}

\end{document}\documentclass{article}
\usepackage[utf8]{inputenc}
\usepackage{booktabs}
\usepackage{amsmath}
\usepackage{natbib}

\title{My Document}
\author{Your Name}
\date{\today}

\begin{document}

\maketitle

\section{Introduction}
This is a sample document demonstrating various \LaTeX{} features including mathematical equations, tables, and formatting. You can cite papers using the citation search feature, and they will be automatically added to your \texttt{cite.bib} file.

\section{Mathematical Expressions}
Here are some example equations:

\subsection{Quadratic Formula}
The quadratic formula is given by:
\[x = \frac{-b \pm \sqrt{b^2 - 4ac}}{2a}\]

\subsection{Matrix Example}
A sample matrix:
\[
\begin{pmatrix}
a_{11} & a_{12} & a_{13} \\
a_{21} & a_{22} & a_{23} \\
a_{31} & a_{32} & a_{33}
\end{pmatrix}
\]

\section{Tables}
Below is a sample table showing data:

\begin{table}[h]
\centering
\begin{tabular}{@{}lrrr@{}}
\toprule
\textbf{Category} & \textbf{Value 1} & \textbf{Value 2} & \textbf{Total} \\
\midrule
Item A & 10.5 & 20.3 & 30.8 \\
Item B & 15.2 & 25.1 & 40.3 \\
Item C & 22.1 & 12.7 & 34.8 \\
\bottomrule
\end{tabular}
\caption{Sample Data Table}
\label{tab:sample}
\end{table}

\section{Lists}
\subsection{Bullet Points}
\begin{itemize}
  \item First item
  \item Second item
  \item Third item with sub-items:
    \begin{itemize}
      \item Sub-item 1
      \item Sub-item 2
    \end{itemize}
\end{itemize}

\subsection{Numbered List}
\begin{enumerate}
  \item First step
  \item Second step
  \item Third step
\end{enumerate}

\section{Citations and References}
When you use the citation search feature, references will be automatically managed. For example, you might cite a paper like this: \cite{example2024}. All citations will be automatically added to your \texttt{cite.bib} file.

\section{Conclusion}
This template demonstrates basic \LaTeX{} features including mathematical formulas, tables, and lists. The bibliography system is now set up to work seamlessly with Bibby AI's citation management. Feel free to modify it for your needs!

\bibliographystyle{plain}
\bibliography{cite}

\end{document}\documentclass{article}
\usepackage[utf8]{inputenc}
\usepackage{booktabs}
\usepackage{amsmath}
\usepackage{natbib}

\title{My Document}
\author{Your Name}
\date{\today}

\begin{document}

\maketitle

\section{Introduction}
This is a sample document demonstrating various \LaTeX{} features including mathematical equations, tables, and formatting. You can cite papers using the citation search feature, and they will be automatically added to your \texttt{cite.bib} file.

\section{Mathematical Expressions}
Here are some example equations:

\subsection{Quadratic Formula}
The quadratic formula is given by:
\[x = \frac{-b \pm \sqrt{b^2 - 4ac}}{2a}\]

\subsection{Matrix Example}
A sample matrix:
\[
\begin{pmatrix}
a_{11} & a_{12} & a_{13} \\
a_{21} & a_{22} & a_{23} \\
a_{31} & a_{32} & a_{33}
\end{pmatrix}
\]

\section{Tables}
Below is a sample table showing data:

\begin{table}[h]
\centering
\begin{tabular}{@{}lrrr@{}}
\toprule
\textbf{Category} & \textbf{Value 1} & \textbf{Value 2} & \textbf{Total} \\
\midrule
Item A & 10.5 & 20.3 & 30.8 \\
Item B & 15.2 & 25.1 & 40.3 \\
Item C & 22.1 & 12.7 & 34.8 \\
\bottomrule
\end{tabular}
\caption{Sample Data Table}
\label{tab:sample}
\end{table}

\section{Lists}
\subsection{Bullet Points}
\begin{itemize}
  \item First item
  \item Second item
  \item Third item with sub-items:
    \begin{itemize}
      \item Sub-item 1
      \item Sub-item 2
    \end{itemize}
\end{itemize}

\subsection{Numbered List}
\begin{enumerate}
  \item First step
  \item Second step
  \item Third step
\end{enumerate}

\section{Citations and References}
When you use the citation search feature, references will be automatically managed. For example, you might cite a paper like this: \cite{example2024}. All citations will be automatically added to your \texttt{cite.bib} file.

\section{Conclusion}
This template demonstrates basic \LaTeX{} features including mathematical formulas, tables, and lists. The bibliography system is now set up to work seamlessly with Bibby AI's citation management. Feel free to modify it for your needs!

\bibliographystyle{plain}
\bibliography{cite}

\end{document}\documentclass{article}
\usepackage[utf8]{inputenc}
\usepackage{booktabs}
\usepackage{amsmath}
\usepackage{natbib}

\title{My Document}
\author{Your Name}
\date{\today}

\begin{document}

\maketitle

\section{Introduction}
This is a sample document demonstrating various \LaTeX{} features including mathematical equations, tables, and formatting. You can cite papers using the citation search feature, and they will be automatically added to your \texttt{cite.bib} file.

\section{Mathematical Expressions}
Here are some example equations:

\subsection{Quadratic Formula}
The quadratic formula is given by:
\[x = \frac{-b \pm \sqrt{b^2 - 4ac}}{2a}\]

\subsection{Matrix Example}
A sample matrix:
\[
\begin{pmatrix}
a_{11} & a_{12} & a_{13} \\
a_{21} & a_{22} & a_{23} \\
a_{31} & a_{32} & a_{33}
\end{pmatrix}
\]

\section{Tables}
Below is a sample table showing data:

\begin{table}[h]
\centering
\begin{tabular}{@{}lrrr@{}}
\toprule
\textbf{Category} & \textbf{Value 1} & \textbf{Value 2} & \textbf{Total} \\
\midrule
Item A & 10.5 & 20.3 & 30.8 \\
Item B & 15.2 & 25.1 & 40.3 \\
Item C & 22.1 & 12.7 & 34.8 \\
\bottomrule
\end{tabular}
\caption{Sample Data Table}
\label{tab:sample}
\end{table}

\section{Lists}
\subsection{Bullet Points}
\begin{itemize}
  \item First item
  \item Second item
  \item Third item with sub-items:
    \begin{itemize}
      \item Sub-item 1
      \item Sub-item 2
    \end{itemize}
\end{itemize}

\subsection{Numbered List}
\begin{enumerate}
  \item First step
  \item Second step
  \item Third step
\end{enumerate}

\section{Citations and References}
When you use the citation search feature, references will be automatically managed. For example, you might cite a paper like this: \cite{example2024}. All citations will be automatically added to your \texttt{cite.bib} file.

\section{Conclusion}
This template demonstrates basic \LaTeX{} features including mathematical formulas, tables, and lists. The bibliography system is now set up to work seamlessly with Bibby AI's citation management. Feel free to modify it for your needs!

\bibliographystyle{plain}
\bibliography{cite}

\end{document}\documentclass{article}
\usepackage[utf8]{inputenc}
\usepackage{booktabs}
\usepackage{amsmath}
\usepackage{natbib}

\title{My Document}
\author{Your Name}
\date{\today}

\begin{document}

\maketitle

\section{Introduction}
This is a sample document demonstrating various \LaTeX{} features including mathematical equations, tables, and formatting. You can cite papers using the citation search feature, and they will be automatically added to your \texttt{cite.bib} file.

\section{Mathematical Expressions}
Here are some example equations:

\subsection{Quadratic Formula}
The quadratic formula is given by:
\[x = \frac{-b \pm \sqrt{b^2 - 4ac}}{2a}\]

\subsection{Matrix Example}
A sample matrix:
\[
\begin{pmatrix}
a_{11} & a_{12} & a_{13} \\
a_{21} & a_{22} & a_{23} \\
a_{31} & a_{32} & a_{33}
\end{pmatrix}
\]

\section{Tables}
Below is a sample table showing data:

\begin{table}[h]
\centering
\begin{tabular}{@{}lrrr@{}}
\toprule
\textbf{Category} & \textbf{Value 1} & \textbf{Value 2} & \textbf{Total} \\
\midrule
Item A & 10.5 & 20.3 & 30.8 \\
Item B & 15.2 & 25.1 & 40.3 \\
Item C & 22.1 & 12.7 & 34.8 \\
\bottomrule
\end{tabular}
\caption{Sample Data Table}
\label{tab:sample}
\end{table}

\section{Lists}
\subsection{Bullet Points}
\begin{itemize}
  \item First item
  \item Second item
  \item Third item with sub-items:
    \begin{itemize}
      \item Sub-item 1
      \item Sub-item 2
    \end{itemize}
\end{itemize}

\subsection{Numbered List}
\begin{enumerate}
  \item First step
  \item Second step
  \item Third step
\end{enumerate}

\section{Citations and References}
When you use the citation search feature, references will be automatically managed. For example, you might cite a paper like this: \cite{example2024}. All citations will be automatically added to your \texttt{cite.bib} file.

\section{Conclusion}
This template demonstrates basic \LaTeX{} features including mathematical formulas, tables, and lists. The bibliography system is now set up to work seamlessly with Bibby AI's citation management. Feel free to modify it for your needs!

\bibliographystyle{plain}
\bibliography{cite}

\end{document}\documentclass{article}
\usepackage[utf8]{inputenc}
\usepackage{booktabs}
\usepackage{amsmath}
\usepackage{natbib}

\title{My Document}
\author{Your Name}
\date{\today}

\begin{document}

\maketitle

\section{Introduction}
This is a sample document demonstrating various \LaTeX{} features including mathematical equations, tables, and formatting. You can cite papers using the citation search feature, and they will be automatically added to your \texttt{cite.bib} file.

\section{Mathematical Expressions}
Here are some example equations:

\subsection{Quadratic Formula}
The quadratic formula is given by:
\[x = \frac{-b \pm \sqrt{b^2 - 4ac}}{2a}\]

\subsection{Matrix Example}
A sample matrix:
\[
\begin{pmatrix}
a_{11} & a_{12} & a_{13} \\
a_{21} & a_{22} & a_{23} \\
a_{31} & a_{32} & a_{33}
\end{pmatrix}
\]

\section{Tables}
Below is a sample table showing data:

\begin{table}[h]
\centering
\begin{tabular}{@{}lrrr@{}}
\toprule
\textbf{Category} & \textbf{Value 1} & \textbf{Value 2} & \textbf{Total} \\
\midrule
Item A & 10.5 & 20.3 & 30.8 \\
Item B & 15.2 & 25.1 & 40.3 \\
Item C & 22.1 & 12.7 & 34.8 \\
\bottomrule
\end{tabular}
\caption{Sample Data Table}
\label{tab:sample}
\end{table}

\section{Lists}
\subsection{Bullet Points}
\begin{itemize}
  \item First item
  \item Second item
  \item Third item with sub-items:
    \begin{itemize}
      \item Sub-item 1
      \item Sub-item 2
    \end{itemize}
\end{itemize}

\subsection{Numbered List}
\begin{enumerate}
  \item First step
  \item Second step
  \item Third step
\end{enumerate}

\section{Citations and References}
When you use the citation search feature, references will be automatically managed. For example, you might cite a paper like this: \cite{example2024}. All citations will be automatically added to your \texttt{cite.bib} file.

\section{Conclusion}
This template demonstrates basic \LaTeX{} features including mathematical formulas, tables, and lists. The bibliography system is now set up to work seamlessly with Bibby AI's citation management. Feel free to modify it for your needs!

\bibliographystyle{plain}
\bibliography{cite}

\end{document}\documentclass{article}
\usepackage[utf8]{inputenc}
\usepackage{booktabs}
\usepackage{amsmath}
\usepackage{natbib}

\title{My Document}
\author{Your Name}
\date{\today}

\begin{document}

\maketitle

\section{Introduction}
This is a sample document demonstrating various \LaTeX{} features including mathematical equations, tables, and formatting. You can cite papers using the citation search feature, and they will be automatically added to your \texttt{cite.bib} file.

\section{Mathematical Expressions}
Here are some example equations:

\subsection{Quadratic Formula}
The quadratic formula is given by:
\[x = \frac{-b \pm \sqrt{b^2 - 4ac}}{2a}\]

\subsection{Matrix Example}
A sample matrix:
\[
\begin{pmatrix}
a_{11} & a_{12} & a_{13} \\
a_{21} & a_{22} & a_{23} \\
a_{31} & a_{32} & a_{33}
\end{pmatrix}
\]

\section{Tables}
Below is a sample table showing data:

\begin{table}[h]
\centering
\begin{tabular}{@{}lrrr@{}}
\toprule
\textbf{Category} & \textbf{Value 1} & \textbf{Value 2} & \textbf{Total} \\
\midrule
Item A & 10.5 & 20.3 & 30.8 \\
Item B & 15.2 & 25.1 & 40.3 \\
Item C & 22.1 & 12.7 & 34.8 \\
\bottomrule
\end{tabular}
\caption{Sample Data Table}
\label{tab:sample}
\end{table}

\section{Lists}
\subsection{Bullet Points}
\begin{itemize}
  \item First item
  \item Second item
  \item Third item with sub-items:
    \begin{itemize}
      \item Sub-item 1
      \item Sub-item 2
    \end{itemize}
\end{itemize}

\subsection{Numbered List}
\begin{enumerate}
  \item First step
  \item Second step
  \item Third step
\end{enumerate}

\section{Citations and References}
When you use the citation search feature, references will be automatically managed. For example, you might cite a paper like this: \cite{example2024}. All citations will be automatically added to your \texttt{cite.bib} file.

\section{Conclusion}
This template demonstrates basic \LaTeX{} features including mathematical formulas, tables, and lists. The bibliography system is now set up to work seamlessly with Bibby AI's citation management. Feel free to modify it for your needs!

\bibliographystyle{plain}
\bibliography{cite}

\end{document}\documentclass{article}
\usepackage[utf8]{inputenc}
\usepackage{booktabs}
\usepackage{amsmath}
\usepackage{natbib}

\title{My Document}
\author{Your Name}
\date{\today}

\begin{document}

\maketitle

\section{Introduction}
This is a sample document demonstrating various \LaTeX{} features including mathematical equations, tables, and formatting. You can cite papers using the citation search feature, and they will be automatically added to your \texttt{cite.bib} file.

\section{Mathematical Expressions}
Here are some example equations:

\subsection{Quadratic Formula}
The quadratic formula is given by:
\[x = \frac{-b \pm \sqrt{b^2 - 4ac}}{2a}\]

\subsection{Matrix Example}
A sample matrix:
\[
\begin{pmatrix}
a_{11} & a_{12} & a_{13} \\
a_{21} & a_{22} & a_{23} \\
a_{31} & a_{32} & a_{33}
\end{pmatrix}
\]

\section{Tables}
Below is a sample table showing data:

\begin{table}[h]
\centering
\begin{tabular}{@{}lrrr@{}}
\toprule
\textbf{Category} & \textbf{Value 1} & \textbf{Value 2} & \textbf{Total} \\
\midrule
Item A & 10.5 & 20.3 & 30.8 \\
Item B & 15.2 & 25.1 & 40.3 \\
Item C & 22.1 & 12.7 & 34.8 \\
\bottomrule
\end{tabular}
\caption{Sample Data Table}
\label{tab:sample}
\end{table}

\section{Lists}
\subsection{Bullet Points}
\begin{itemize}
  \item First item
  \item Second item
  \item Third item with sub-items:
    \begin{itemize}
      \item Sub-item 1
      \item Sub-item 2
    \end{itemize}
\end{itemize}

\subsection{Numbered List}
\begin{enumerate}
  \item First step
  \item Second step
  \item Third step
\end{enumerate}

\section{Citations and References}
When you use the citation search feature, references will be automatically managed. For example, you might cite a paper like this: \cite{example2024}. All citations will be automatically added to your \texttt{cite.bib} file.

\section{Conclusion}
This template demonstrates basic \LaTeX{} features including mathematical formulas, tables, and lists. The bibliography system is now set up to work seamlessly with Bibby AI's citation management. Feel free to modify it for your needs!

\bibliographystyle{plain}
\bibliography{cite}

\end{document}\documentclass{article}
\usepackage[utf8]{inputenc}
\usepackage{booktabs}
\usepackage{amsmath}
\usepackage{natbib}

\title{My Document}
\author{Your Name}
\date{\today}

\begin{document}

\maketitle

\section{Introduction}
This is a sample document demonstrating various \LaTeX{} features including mathematical equations, tables, and formatting. You can cite papers using the citation search feature, and they will be automatically added to your \texttt{cite.bib} file.

\section{Mathematical Expressions}
Here are some example equations:

\subsection{Quadratic Formula}
The quadratic formula is given by:
\[x = \frac{-b \pm \sqrt{b^2 - 4ac}}{2a}\]

\subsection{Matrix Example}
A sample matrix:
\[
\begin{pmatrix}
a_{11} & a_{12} & a_{13} \\
a_{21} & a_{22} & a_{23} \\
a_{31} & a_{32} & a_{33}
\end{pmatrix}
\]

\section{Tables}
Below is a sample table showing data:

\begin{table}[h]
\centering
\begin{tabular}{@{}lrrr@{}}
\toprule
\textbf{Category} & \textbf{Value 1} & \textbf{Value 2} & \textbf{Total} \\
\midrule
Item A & 10.5 & 20.3 & 30.8 \\
Item B & 15.2 & 25.1 & 40.3 \\
Item C & 22.1 & 12.7 & 34.8 \\
\bottomrule
\end{tabular}
\caption{Sample Data Table}
\label{tab:sample}
\end{table}

\section{Lists}
\subsection{Bullet Points}
\begin{itemize}
  \item First item
  \item Second item
  \item Third item with sub-items:
    \begin{itemize}
      \item Sub-item 1
      \item Sub-item 2
    \end{itemize}
\end{itemize}

\subsection{Numbered List}
\begin{enumerate}
  \item First step
  \item Second step
  \item Third step
\end{enumerate}

\section{Citations and References}
When you use the citation search feature, references will be automatically managed. For example, you might cite a paper like this: \cite{example2024}. All citations will be automatically added to your \texttt{cite.bib} file.

\section{Conclusion}
This template demonstrates basic \LaTeX{} features including mathematical formulas, tables, and lists. The bibliography system is now set up to work seamlessly with Bibby AI's citation management. Feel free to modify it for your needs!

\bibliographystyle{plain}
\bibliography{cite}

\end{document}\documentclass{article}
\usepackage[utf8]{inputenc}
\usepackage{booktabs}
\usepackage{amsmath}
\usepackage{natbib}

\title{My Document}
\author{Your Name}
\date{\today}

\begin{document}

\maketitle

\section{Introduction}
This is a sample document demonstrating various \LaTeX{} features including mathematical equations, tables, and formatting. You can cite papers using the citation search feature, and they will be automatically added to your \texttt{cite.bib} file.

\section{Mathematical Expressions}
Here are some example equations:

\subsection{Quadratic Formula}
The quadratic formula is given by:
\[x = \frac{-b \pm \sqrt{b^2 - 4ac}}{2a}\]

\subsection{Matrix Example}
A sample matrix:
\[
\begin{pmatrix}
a_{11} & a_{12} & a_{13} \\
a_{21} & a_{22} & a_{23} \\
a_{31} & a_{32} & a_{33}
\end{pmatrix}
\]

\section{Tables}
Below is a sample table showing data:

\begin{table}[h]
\centering
\begin{tabular}{@{}lrrr@{}}
\toprule
\textbf{Category} & \textbf{Value 1} & \textbf{Value 2} & \textbf{Total} \\
\midrule
Item A & 10.5 & 20.3 & 30.8 \\
Item B & 15.2 & 25.1 & 40.3 \\
Item C & 22.1 & 12.7 & 34.8 \\
\bottomrule
\end{tabular}
\caption{Sample Data Table}
\label{tab:sample}
\end{table}

\section{Lists}
\subsection{Bullet Points}
\begin{itemize}
  \item First item
  \item Second item
  \item Third item with sub-items:
    \begin{itemize}
      \item Sub-item 1
      \item Sub-item 2
    \end{itemize}
\end{itemize}

\subsection{Numbered List}
\begin{enumerate}
  \item First step
  \item Second step
  \item Third step
\end{enumerate}

\section{Citations and References}
When you use the citation search feature, references will be automatically managed. For example, you might cite a paper like this: \cite{example2024}. All citations will be automatically added to your \texttt{cite.bib} file.

\section{Conclusion}
This template demonstrates basic \LaTeX{} features including mathematical formulas, tables, and lists. The bibliography system is now set up to work seamlessly with Bibby AI's citation management. Feel free to modify it for your needs!

\bibliographystyle{plain}
\bibliography{cite}

\end{document}\documentclass{article}
\usepackage[utf8]{inputenc}
\usepackage{booktabs}
\usepackage{amsmath}
\usepackage{natbib}

\title{My Document}
\author{Your Name}
\date{\today}

\begin{document}

\maketitle

\section{Introduction}
This is a sample document demonstrating various \LaTeX{} features including mathematical equations, tables, and formatting. You can cite papers using the citation search feature, and they will be automatically added to your \texttt{cite.bib} file.

\section{Mathematical Expressions}
Here are some example equations:

\subsection{Quadratic Formula}
The quadratic formula is given by:
\[x = \frac{-b \pm \sqrt{b^2 - 4ac}}{2a}\]

\subsection{Matrix Example}
A sample matrix:
\[
\begin{pmatrix}
a_{11} & a_{12} & a_{13} \\
a_{21} & a_{22} & a_{23} \\
a_{31} & a_{32} & a_{33}
\end{pmatrix}
\]

\section{Tables}
Below is a sample table showing data:

\begin{table}[h]
\centering
\begin{tabular}{@{}lrrr@{}}
\toprule
\textbf{Category} & \textbf{Value 1} & \textbf{Value 2} & \textbf{Total} \\
\midrule
Item A & 10.5 & 20.3 & 30.8 \\
Item B & 15.2 & 25.1 & 40.3 \\
Item C & 22.1 & 12.7 & 34.8 \\
\bottomrule
\end{tabular}
\caption{Sample Data Table}
\label{tab:sample}
\end{table}

\section{Lists}
\subsection{Bullet Points}
\begin{itemize}
  \item First item
  \item Second item
  \item Third item with sub-items:
    \begin{itemize}
      \item Sub-item 1
      \item Sub-item 2
    \end{itemize}
\end{itemize}

\subsection{Numbered List}
\begin{enumerate}
  \item First step
  \item Second step
  \item Third step
\end{enumerate}

\section{Citations and References}
When you use the citation search feature, references will be automatically managed. For example, you might cite a paper like this: \cite{example2024}. All citations will be automatically added to your \texttt{cite.bib} file.

\section{Conclusion}
This template demonstrates basic \LaTeX{} features including mathematical formulas, tables, and lists. The bibliography system is now set up to work seamlessly with Bibby AI's citation management. Feel free to modify it for your needs!

\bibliographystyle{plain}
\bibliography{cite}

\end{document}\documentclass{article}
\usepackage[utf8]{inputenc}
\usepackage{booktabs}
\usepackage{amsmath}
\usepackage{natbib}

\title{My Document}
\author{Your Name}
\date{\today}

\begin{document}

\maketitle

\section{Introduction}
This is a sample document demonstrating various \LaTeX{} features including mathematical equations, tables, and formatting. You can cite papers using the citation search feature, and they will be automatically added to your \texttt{cite.bib} file.

\section{Mathematical Expressions}
Here are some example equations:

\subsection{Quadratic Formula}
The quadratic formula is given by:
\[x = \frac{-b \pm \sqrt{b^2 - 4ac}}{2a}\]

\subsection{Matrix Example}
A sample matrix:
\[
\begin{pmatrix}
a_{11} & a_{12} & a_{13} \\
a_{21} & a_{22} & a_{23} \\
a_{31} & a_{32} & a_{33}
\end{pmatrix}
\]

\section{Tables}
Below is a sample table showing data:

\begin{table}[h]
\centering
\begin{tabular}{@{}lrrr@{}}
\toprule
\textbf{Category} & \textbf{Value 1} & \textbf{Value 2} & \textbf{Total} \\
\midrule
Item A & 10.5 & 20.3 & 30.8 \\
Item B & 15.2 & 25.1 & 40.3 \\
Item C & 22.1 & 12.7 & 34.8 \\
\bottomrule
\end{tabular}
\caption{Sample Data Table}
\label{tab:sample}
\end{table}

\section{Lists}
\subsection{Bullet Points}
\begin{itemize}
  \item First item
  \item Second item
  \item Third item with sub-items:
    \begin{itemize}
      \item Sub-item 1
      \item Sub-item 2
    \end{itemize}
\end{itemize}

\subsection{Numbered List}
\begin{enumerate}
  \item First step
  \item Second step
  \item Third step
\end{enumerate}

\section{Citations and References}
When you use the citation search feature, references will be automatically managed. For example, you might cite a paper like this: \cite{example2024}. All citations will be automatically added to your \texttt{cite.bib} file.

\section{Conclusion}
This template demonstrates basic \LaTeX{} features including mathematical formulas, tables, and lists. The bibliography system is now set up to work seamlessly with Bibby AI's citation management. Feel free to modify it for your needs!

\bibliographystyle{plain}
\bibliography{cite}

\end{document}\documentclass{article}
\usepackage[utf8]{inputenc}
\usepackage{booktabs}
\usepackage{amsmath}
\usepackage{natbib}

\title{My Document}
\author{Your Name}
\date{\today}

\begin{document}

\maketitle

\section{Introduction}
This is a sample document demonstrating various \LaTeX{} features including mathematical equations, tables, and formatting. You can cite papers using the citation search feature, and they will be automatically added to your \texttt{cite.bib} file.

\section{Mathematical Expressions}
Here are some example equations:

\subsection{Quadratic Formula}
The quadratic formula is given by:
\[x = \frac{-b \pm \sqrt{b^2 - 4ac}}{2a}\]

\subsection{Matrix Example}
A sample matrix:
\[
\begin{pmatrix}
a_{11} & a_{12} & a_{13} \\
a_{21} & a_{22} & a_{23} \\
a_{31} & a_{32} & a_{33}
\end{pmatrix}
\]

\section{Tables}
Below is a sample table showing data:

\begin{table}[h]
\centering
\begin{tabular}{@{}lrrr@{}}
\toprule
\textbf{Category} & \textbf{Value 1} & \textbf{Value 2} & \textbf{Total} \\
\midrule
Item A & 10.5 & 20.3 & 30.8 \\
Item B & 15.2 & 25.1 & 40.3 \\
Item C & 22.1 & 12.7 & 34.8 \\
\bottomrule
\end{tabular}
\caption{Sample Data Table}
\label{tab:sample}
\end{table}

\section{Lists}
\subsection{Bullet Points}
\begin{itemize}
  \item First item
  \item Second item
  \item Third item with sub-items:
    \begin{itemize}
      \item Sub-item 1
      \item Sub-item 2
    \end{itemize}
\end{itemize}

\subsection{Numbered List}
\begin{enumerate}
  \item First step
  \item Second step
  \item Third step
\end{enumerate}

\section{Citations and References}
When you use the citation search feature, references will be automatically managed. For example, you might cite a paper like this: \cite{example2024}. All citations will be automatically added to your \texttt{cite.bib} file.

\section{Conclusion}
This template demonstrates basic \LaTeX{} features including mathematical formulas, tables, and lists. The bibliography system is now set up to work seamlessly with Bibby AI's citation management. Feel free to modify it for your needs!

\bibliographystyle{plain}
\bibliography{cite}

\end{document}\documentclass{article}
\usepackage[utf8]{inputenc}
\usepackage{booktabs}
\usepackage{amsmath}
\usepackage{natbib}

\title{My Document}
\author{Your Name}
\date{\today}

\begin{document}

\maketitle

\section{Introduction}
This is a sample document demonstrating various \LaTeX{} features including mathematical equations, tables, and formatting. You can cite papers using the citation search feature, and they will be automatically added to your \texttt{cite.bib} file.

\section{Mathematical Expressions}
Here are some example equations:

\subsection{Quadratic Formula}
The quadratic formula is given by:
\[x = \frac{-b \pm \sqrt{b^2 - 4ac}}{2a}\]

\subsection{Matrix Example}
A sample matrix:
\[
\begin{pmatrix}
a_{11} & a_{12} & a_{13} \\
a_{21} & a_{22} & a_{23} \\
a_{31} & a_{32} & a_{33}
\end{pmatrix}
\]

\section{Tables}
Below is a sample table showing data:

\begin{table}[h]
\centering
\begin{tabular}{@{}lrrr@{}}
\toprule
\textbf{Category} & \textbf{Value 1} & \textbf{Value 2} & \textbf{Total} \\
\midrule
Item A & 10.5 & 20.3 & 30.8 \\
Item B & 15.2 & 25.1 & 40.3 \\
Item C & 22.1 & 12.7 & 34.8 \\
\bottomrule
\end{tabular}
\caption{Sample Data Table}
\label{tab:sample}
\end{table}

\section{Lists}
\subsection{Bullet Points}
\begin{itemize}
  \item First item
  \item Second item
  \item Third item with sub-items:
    \begin{itemize}
      \item Sub-item 1
      \item Sub-item 2
    \end{itemize}
\end{itemize}

\subsection{Numbered List}
\begin{enumerate}
  \item First step
  \item Second step
  \item Third step
\end{enumerate}

\section{Citations and References}
When you use the citation search feature, references will be automatically managed. For example, you might cite a paper like this: \cite{example2024}. All citations will be automatically added to your \texttt{cite.bib} file.

\section{Conclusion}
This template demonstrates basic \LaTeX{} features including mathematical formulas, tables, and lists. The bibliography system is now set up to work seamlessly with Bibby AI's citation management. Feel free to modify it for your needs!

\bibliographystyle{plain}
\bibliography{cite}

\end{document}\documentclass{article}
\usepackage[utf8]{inputenc}
\usepackage{booktabs}
\usepackage{amsmath}
\usepackage{natbib}

\title{My Document}
\author{Your Name}
\date{\today}

\begin{document}

\maketitle

\section{Introduction}
This is a sample document demonstrating various \LaTeX{} features including mathematical equations, tables, and formatting. You can cite papers using the citation search feature, and they will be automatically added to your \texttt{cite.bib} file.

\section{Mathematical Expressions}
Here are some example equations:

\subsection{Quadratic Formula}
The quadratic formula is given by:
\[x = \frac{-b \pm \sqrt{b^2 - 4ac}}{2a}\]

\subsection{Matrix Example}
A sample matrix:
\[
\begin{pmatrix}
a_{11} & a_{12} & a_{13} \\
a_{21} & a_{22} & a_{23} \\
a_{31} & a_{32} & a_{33}
\end{pmatrix}
\]

\section{Tables}
Below is a sample table showing data:

\begin{table}[h]
\centering
\begin{tabular}{@{}lrrr@{}}
\toprule
\textbf{Category} & \textbf{Value 1} & \textbf{Value 2} & \textbf{Total} \\
\midrule
Item A & 10.5 & 20.3 & 30.8 \\
Item B & 15.2 & 25.1 & 40.3 \\
Item C & 22.1 & 12.7 & 34.8 \\
\bottomrule
\end{tabular}
\caption{Sample Data Table}
\label{tab:sample}
\end{table}

\section{Lists}
\subsection{Bullet Points}
\begin{itemize}
  \item First item
  \item Second item
  \item Third item with sub-items:
    \begin{itemize}
      \item Sub-item 1
      \item Sub-item 2
    \end{itemize}
\end{itemize}

\subsection{Numbered List}
\begin{enumerate}
  \item First step
  \item Second step
  \item Third step
\end{enumerate}

\section{Citations and References}
When you use the citation search feature, references will be automatically managed. For example, you might cite a paper like this: \cite{example2024}. All citations will be automatically added to your \texttt{cite.bib} file.

\section{Conclusion}
This template demonstrates basic \LaTeX{} features including mathematical formulas, tables, and lists. The bibliography system is now set up to work seamlessly with Bibby AI's citation management. Feel free to modify it for your needs!

\bibliographystyle{plain}
\bibliography{cite}

\end{document}\documentclass{article}
\usepackage[utf8]{inputenc}
\usepackage{booktabs}
\usepackage{amsmath}
\usepackage{natbib}

\title{My Document}
\author{Your Name}
\date{\today}

\begin{document}

\maketitle

\section{Introduction}
This is a sample document demonstrating various \LaTeX{} features including mathematical equations, tables, and formatting. You can cite papers using the citation search feature, and they will be automatically added to your \texttt{cite.bib} file.

\section{Mathematical Expressions}
Here are some example equations:

\subsection{Quadratic Formula}
The quadratic formula is given by:
\[x = \frac{-b \pm \sqrt{b^2 - 4ac}}{2a}\]

\subsection{Matrix Example}
A sample matrix:
\[
\begin{pmatrix}
a_{11} & a_{12} & a_{13} \\
a_{21} & a_{22} & a_{23} \\
a_{31} & a_{32} & a_{33}
\end{pmatrix}
\]

\section{Tables}
Below is a sample table showing data:

\begin{table}[h]
\centering
\begin{tabular}{@{}lrrr@{}}
\toprule
\textbf{Category} & \textbf{Value 1} & \textbf{Value 2} & \textbf{Total} \\
\midrule
Item A & 10.5 & 20.3 & 30.8 \\
Item B & 15.2 & 25.1 & 40.3 \\
Item C & 22.1 & 12.7 & 34.8 \\
\bottomrule
\end{tabular}
\caption{Sample Data Table}
\label{tab:sample}
\end{table}

\section{Lists}
\subsection{Bullet Points}
\begin{itemize}
  \item First item
  \item Second item
  \item Third item with sub-items:
    \begin{itemize}
      \item Sub-item 1
      \item Sub-item 2
    \end{itemize}
\end{itemize}

\subsection{Numbered List}
\begin{enumerate}
  \item First step
  \item Second step
  \item Third step
\end{enumerate}

\section{Citations and References}
When you use the citation search feature, references will be automatically managed. For example, you might cite a paper like this: \cite{example2024}. All citations will be automatically added to your \texttt{cite.bib} file.

\section{Conclusion}
This template demonstrates basic \LaTeX{} features including mathematical formulas, tables, and lists. The bibliography system is now set up to work seamlessly with Bibby AI's citation management. Feel free to modify it for your needs!

\bibliographystyle{plain}
\bibliography{cite}

\end{document}\documentclass{article}
\usepackage[utf8]{inputenc}
\usepackage{booktabs}
\usepackage{amsmath}
\usepackage{natbib}

\title{My Document}
\author{Your Name}
\date{\today}

\begin{document}

\maketitle

\section{Introduction}
This is a sample document demonstrating various \LaTeX{} features including mathematical equations, tables, and formatting. You can cite papers using the citation search feature, and they will be automatically added to your \texttt{cite.bib} file.

\section{Mathematical Expressions}
Here are some example equations:

\subsection{Quadratic Formula}
The quadratic formula is given by:
\[x = \frac{-b \pm \sqrt{b^2 - 4ac}}{2a}\]

\subsection{Matrix Example}
A sample matrix:
\[
\begin{pmatrix}
a_{11} & a_{12} & a_{13} \\
a_{21} & a_{22} & a_{23} \\
a_{31} & a_{32} & a_{33}
\end{pmatrix}
\]

\section{Tables}
Below is a sample table showing data:

\begin{table}[h]
\centering
\begin{tabular}{@{}lrrr@{}}
\toprule
\textbf{Category} & \textbf{Value 1} & \textbf{Value 2} & \textbf{Total} \\
\midrule
Item A & 10.5 & 20.3 & 30.8 \\
Item B & 15.2 & 25.1 & 40.3 \\
Item C & 22.1 & 12.7 & 34.8 \\
\bottomrule
\end{tabular}
\caption{Sample Data Table}
\label{tab:sample}
\end{table}

\section{Lists}
\subsection{Bullet Points}
\begin{itemize}
  \item First item
  \item Second item
  \item Third item with sub-items:
    \begin{itemize}
      \item Sub-item 1
      \item Sub-item 2
    \end{itemize}
\end{itemize}

\subsection{Numbered List}
\begin{enumerate}
  \item First step
  \item Second step
  \item Third step
\end{enumerate}

\section{Citations and References}
When you use the citation search feature, references will be automatically managed. For example, you might cite a paper like this: \cite{example2024}. All citations will be automatically added to your \texttt{cite.bib} file.

\section{Conclusion}
This template demonstrates basic \LaTeX{} features including mathematical formulas, tables, and lists. The bibliography system is now set up to work seamlessly with Bibby AI's citation management. Feel free to modify it for your needs!

\bibliographystyle{plain}
\bibliography{cite}

\end{document}\documentclass{article}
\usepackage[utf8]{inputenc}
\usepackage{booktabs}
\usepackage{amsmath}
\usepackage{natbib}

\title{My Document}
\author{Your Name}
\date{\today}

\begin{document}

\maketitle

\section{Introduction}
This is a sample document demonstrating various \LaTeX{} features including mathematical equations, tables, and formatting. You can cite papers using the citation search feature, and they will be automatically added to your \texttt{cite.bib} file.

\section{Mathematical Expressions}
Here are some example equations:

\subsection{Quadratic Formula}
The quadratic formula is given by:
\[x = \frac{-b \pm \sqrt{b^2 - 4ac}}{2a}\]

\subsection{Matrix Example}
A sample matrix:
\[
\begin{pmatrix}
a_{11} & a_{12} & a_{13} \\
a_{21} & a_{22} & a_{23} \\
a_{31} & a_{32} & a_{33}
\end{pmatrix}
\]

\section{Tables}
Below is a sample table showing data:

\begin{table}[h]
\centering
\begin{tabular}{@{}lrrr@{}}
\toprule
\textbf{Category} & \textbf{Value 1} & \textbf{Value 2} & \textbf{Total} \\
\midrule
Item A & 10.5 & 20.3 & 30.8 \\
Item B & 15.2 & 25.1 & 40.3 \\
Item C & 22.1 & 12.7 & 34.8 \\
\bottomrule
\end{tabular}
\caption{Sample Data Table}
\label{tab:sample}
\end{table}

\section{Lists}
\subsection{Bullet Points}
\begin{itemize}
  \item First item
  \item Second item
  \item Third item with sub-items:
    \begin{itemize}
      \item Sub-item 1
      \item Sub-item 2
    \end{itemize}
\end{itemize}

\subsection{Numbered List}
\begin{enumerate}
  \item First step
  \item Second step
  \item Third step
\end{enumerate}

\section{Citations and References}
When you use the citation search feature, references will be automatically managed. For example, you might cite a paper like this: \cite{example2024}. All citations will be automatically added to your \texttt{cite.bib} file.

\section{Conclusion}
This template demonstrates basic \LaTeX{} features including mathematical formulas, tables, and lists. The bibliography system is now set up to work seamlessly with Bibby AI's citation management. Feel free to modify it for your needs!

\bibliographystyle{plain}
\bibliography{cite}

\end{document}\documentclass{article}
\usepackage[utf8]{inputenc}
\usepackage{booktabs}
\usepackage{amsmath}
\usepackage{natbib}

\title{My Document}
\author{Your Name}
\date{\today}

\begin{document}

\maketitle

\section{Introduction}
This is a sample document demonstrating various \LaTeX{} features including mathematical equations, tables, and formatting. You can cite papers using the citation search feature, and they will be automatically added to your \texttt{cite.bib} file.

\section{Mathematical Expressions}
Here are some example equations:

\subsection{Quadratic Formula}
The quadratic formula is given by:
\[x = \frac{-b \pm \sqrt{b^2 - 4ac}}{2a}\]

\subsection{Matrix Example}
A sample matrix:
\[
\begin{pmatrix}
a_{11} & a_{12} & a_{13} \\
a_{21} & a_{22} & a_{23} \\
a_{31} & a_{32} & a_{33}
\end{pmatrix}
\]

\section{Tables}
Below is a sample table showing data:

\begin{table}[h]
\centering
\begin{tabular}{@{}lrrr@{}}
\toprule
\textbf{Category} & \textbf{Value 1} & \textbf{Value 2} & \textbf{Total} \\
\midrule
Item A & 10.5 & 20.3 & 30.8 \\
Item B & 15.2 & 25.1 & 40.3 \\
Item C & 22.1 & 12.7 & 34.8 \\
\bottomrule
\end{tabular}
\caption{Sample Data Table}
\label{tab:sample}
\end{table}

\section{Lists}
\subsection{Bullet Points}
\begin{itemize}
  \item First item
  \item Second item
  \item Third item with sub-items:
    \begin{itemize}
      \item Sub-item 1
      \item Sub-item 2
    \end{itemize}
\end{itemize}

\subsection{Numbered List}
\begin{enumerate}
  \item First step
  \item Second step
  \item Third step
\end{enumerate}

\section{Citations and References}
When you use the citation search feature, references will be automatically managed. For example, you might cite a paper like this: \cite{example2024}. All citations will be automatically added to your \texttt{cite.bib} file.

\section{Conclusion}
This template demonstrates basic \LaTeX{} features including mathematical formulas, tables, and lists. The bibliography system is now set up to work seamlessly with Bibby AI's citation management. Feel free to modify it for your needs!

\bibliographystyle{plain}
\bibliography{cite}

\end{document}\documentclass{article}
\usepackage[utf8]{inputenc}
\usepackage{booktabs}
\usepackage{amsmath}
\usepackage{natbib}

\title{My Document}
\author{Your Name}
\date{\today}

\begin{document}

\maketitle

\section{Introduction}
This is a sample document demonstrating various \LaTeX{} features including mathematical equations, tables, and formatting. You can cite papers using the citation search feature, and they will be automatically added to your \texttt{cite.bib} file.

\section{Mathematical Expressions}
Here are some example equations:

\subsection{Quadratic Formula}
The quadratic formula is given by:
\[x = \frac{-b \pm \sqrt{b^2 - 4ac}}{2a}\]

\subsection{Matrix Example}
A sample matrix:
\[
\begin{pmatrix}
a_{11} & a_{12} & a_{13} \\
a_{21} & a_{22} & a_{23} \\
a_{31} & a_{32} & a_{33}
\end{pmatrix}
\]

\section{Tables}
Below is a sample table showing data:

\begin{table}[h]
\centering
\begin{tabular}{@{}lrrr@{}}
\toprule
\textbf{Category} & \textbf{Value 1} & \textbf{Value 2} & \textbf{Total} \\
\midrule
Item A & 10.5 & 20.3 & 30.8 \\
Item B & 15.2 & 25.1 & 40.3 \\
Item C & 22.1 & 12.7 & 34.8 \\
\bottomrule
\end{tabular}
\caption{Sample Data Table}
\label{tab:sample}
\end{table}

\section{Lists}
\subsection{Bullet Points}
\begin{itemize}
  \item First item
  \item Second item
  \item Third item with sub-items:
    \begin{itemize}
      \item Sub-item 1
      \item Sub-item 2
    \end{itemize}
\end{itemize}

\subsection{Numbered List}
\begin{enumerate}
  \item First step
  \item Second step
  \item Third step
\end{enumerate}

\section{Citations and References}
When you use the citation search feature, references will be automatically managed. For example, you might cite a paper like this: \cite{example2024}. All citations will be automatically added to your \texttt{cite.bib} file.

\section{Conclusion}
This template demonstrates basic \LaTeX{} features including mathematical formulas, tables, and lists. The bibliography system is now set up to work seamlessly with Bibby AI's citation management. Feel free to modify it for your needs!

\bibliographystyle{plain}
\bibliography{cite}

\end{document}\documentclass{article}
\usepackage[utf8]{inputenc}
\usepackage{booktabs}
\usepackage{amsmath}
\usepackage{natbib}

\title{My Document}
\author{Your Name}
\date{\today}

\begin{document}

\maketitle

\section{Introduction}
This is a sample document demonstrating various \LaTeX{} features including mathematical equations, tables, and formatting. You can cite papers using the citation search feature, and they will be automatically added to your \texttt{cite.bib} file.

\section{Mathematical Expressions}
Here are some example equations:

\subsection{Quadratic Formula}
The quadratic formula is given by:
\[x = \frac{-b \pm \sqrt{b^2 - 4ac}}{2a}\]

\subsection{Matrix Example}
A sample matrix:
\[
\begin{pmatrix}
a_{11} & a_{12} & a_{13} \\
a_{21} & a_{22} & a_{23} \\
a_{31} & a_{32} & a_{33}
\end{pmatrix}
\]

\section{Tables}
Below is a sample table showing data:

\begin{table}[h]
\centering
\begin{tabular}{@{}lrrr@{}}
\toprule
\textbf{Category} & \textbf{Value 1} & \textbf{Value 2} & \textbf{Total} \\
\midrule
Item A & 10.5 & 20.3 & 30.8 \\
Item B & 15.2 & 25.1 & 40.3 \\
Item C & 22.1 & 12.7 & 34.8 \\
\bottomrule
\end{tabular}
\caption{Sample Data Table}
\label{tab:sample}
\end{table}

\section{Lists}
\subsection{Bullet Points}
\begin{itemize}
  \item First item
  \item Second item
  \item Third item with sub-items:
    \begin{itemize}
      \item Sub-item 1
      \item Sub-item 2
    \end{itemize}
\end{itemize}

\subsection{Numbered List}
\begin{enumerate}
  \item First step
  \item Second step
  \item Third step
\end{enumerate}

\section{Citations and References}
When you use the citation search feature, references will be automatically managed. For example, you might cite a paper like this: \cite{example2024}. All citations will be automatically added to your \texttt{cite.bib} file.

\section{Conclusion}
This template demonstrates basic \LaTeX{} features including mathematical formulas, tables, and lists. The bibliography system is now set up to work seamlessly with Bibby AI's citation management. Feel free to modify it for your needs!

\bibliographystyle{plain}
\bibliography{cite}

\end{document}\documentclass{article}
\usepackage[utf8]{inputenc}
\usepackage{booktabs}
\usepackage{amsmath}
\usepackage{natbib}

\title{My Document}
\author{Your Name}
\date{\today}

\begin{document}

\maketitle

\section{Introduction}
This is a sample document demonstrating various \LaTeX{} features including mathematical equations, tables, and formatting. You can cite papers using the citation search feature, and they will be automatically added to your \texttt{cite.bib} file.

\section{Mathematical Expressions}
Here are some example equations:

\subsection{Quadratic Formula}
The quadratic formula is given by:
\[x = \frac{-b \pm \sqrt{b^2 - 4ac}}{2a}\]

\subsection{Matrix Example}
A sample matrix:
\[
\begin{pmatrix}
a_{11} & a_{12} & a_{13} \\
a_{21} & a_{22} & a_{23} \\
a_{31} & a_{32} & a_{33}
\end{pmatrix}
\]

\section{Tables}
Below is a sample table showing data:

\begin{table}[h]
\centering
\begin{tabular}{@{}lrrr@{}}
\toprule
\textbf{Category} & \textbf{Value 1} & \textbf{Value 2} & \textbf{Total} \\
\midrule
Item A & 10.5 & 20.3 & 30.8 \\
Item B & 15.2 & 25.1 & 40.3 \\
Item C & 22.1 & 12.7 & 34.8 \\
\bottomrule
\end{tabular}
\caption{Sample Data Table}
\label{tab:sample}
\end{table}

\section{Lists}
\subsection{Bullet Points}
\begin{itemize}
  \item First item
  \item Second item
  \item Third item with sub-items:
    \begin{itemize}
      \item Sub-item 1
      \item Sub-item 2
    \end{itemize}
\end{itemize}

\subsection{Numbered List}
\begin{enumerate}
  \item First step
  \item Second step
  \item Third step
\end{enumerate}

\section{Citations and References}
When you use the citation search feature, references will be automatically managed. For example, you might cite a paper like this: \cite{example2024}. All citations will be automatically added to your \texttt{cite.bib} file.

\section{Conclusion}
This template demonstrates basic \LaTeX{} features including mathematical formulas, tables, and lists. The bibliography system is now set up to work seamlessly with Bibby AI's citation management. Feel free to modify it for your needs!

\bibliographystyle{plain}
\bibliography{cite}

\end{document}\documentclass{article}
\usepackage[utf8]{inputenc}
\usepackage{booktabs}
\usepackage{amsmath}
\usepackage{natbib}

\title{My Document}
\author{Your Name}
\date{\today}

\begin{document}

\maketitle

\section{Introduction}
This is a sample document demonstrating various \LaTeX{} features including mathematical equations, tables, and formatting. You can cite papers using the citation search feature, and they will be automatically added to your \texttt{cite.bib} file.

\section{Mathematical Expressions}
Here are some example equations:

\subsection{Quadratic Formula}
The quadratic formula is given by:
\[x = \frac{-b \pm \sqrt{b^2 - 4ac}}{2a}\]

\subsection{Matrix Example}
A sample matrix:
\[
\begin{pmatrix}
a_{11} & a_{12} & a_{13} \\
a_{21} & a_{22} & a_{23} \\
a_{31} & a_{32} & a_{33}
\end{pmatrix}
\]

\section{Tables}
Below is a sample table showing data:

\begin{table}[h]
\centering
\begin{tabular}{@{}lrrr@{}}
\toprule
\textbf{Category} & \textbf{Value 1} & \textbf{Value 2} & \textbf{Total} \\
\midrule
Item A & 10.5 & 20.3 & 30.8 \\
Item B & 15.2 & 25.1 & 40.3 \\
Item C & 22.1 & 12.7 & 34.8 \\
\bottomrule
\end{tabular}
\caption{Sample Data Table}
\label{tab:sample}
\end{table}

\section{Lists}
\subsection{Bullet Points}
\begin{itemize}
  \item First item
  \item Second item
  \item Third item with sub-items:
    \begin{itemize}
      \item Sub-item 1
      \item Sub-item 2
    \end{itemize}
\end{itemize}

\subsection{Numbered List}
\begin{enumerate}
  \item First step
  \item Second step
  \item Third step
\end{enumerate}

\section{Citations and References}
When you use the citation search feature, references will be automatically managed. For example, you might cite a paper like this: \cite{example2024}. All citations will be automatically added to your \texttt{cite.bib} file.

\section{Conclusion}
This template demonstrates basic \LaTeX{} features including mathematical formulas, tables, and lists. The bibliography system is now set up to work seamlessly with Bibby AI's citation management. Feel free to modify it for your needs!

\bibliographystyle{plain}
\bibliography{cite}

\end{document}\documentclass{article}
\usepackage[utf8]{inputenc}
\usepackage{booktabs}
\usepackage{amsmath}
\usepackage{natbib}

\title{My Document}
\author{Your Name}
\date{\today}

\begin{document}

\maketitle

\section{Introduction}
This is a sample document demonstrating various \LaTeX{} features including mathematical equations, tables, and formatting. You can cite papers using the citation search feature, and they will be automatically added to your \texttt{cite.bib} file.

\section{Mathematical Expressions}
Here are some example equations:

\subsection{Quadratic Formula}
The quadratic formula is given by:
\[x = \frac{-b \pm \sqrt{b^2 - 4ac}}{2a}\]

\subsection{Matrix Example}
A sample matrix:
\[
\begin{pmatrix}
a_{11} & a_{12} & a_{13} \\
a_{21} & a_{22} & a_{23} \\
a_{31} & a_{32} & a_{33}
\end{pmatrix}
\]

\section{Tables}
Below is a sample table showing data:

\begin{table}[h]
\centering
\begin{tabular}{@{}lrrr@{}}
\toprule
\textbf{Category} & \textbf{Value 1} & \textbf{Value 2} & \textbf{Total} \\
\midrule
Item A & 10.5 & 20.3 & 30.8 \\
Item B & 15.2 & 25.1 & 40.3 \\
Item C & 22.1 & 12.7 & 34.8 \\
\bottomrule
\end{tabular}
\caption{Sample Data Table}
\label{tab:sample}
\end{table}

\section{Lists}
\subsection{Bullet Points}
\begin{itemize}
  \item First item
  \item Second item
  \item Third item with sub-items:
    \begin{itemize}
      \item Sub-item 1
      \item Sub-item 2
    \end{itemize}
\end{itemize}

\subsection{Numbered List}
\begin{enumerate}
  \item First step
  \item Second step
  \item Third step
\end{enumerate}

\section{Citations and References}
When you use the citation search feature, references will be automatically managed. For example, you might cite a paper like this: \cite{example2024}. All citations will be automatically added to your \texttt{cite.bib} file.

\section{Conclusion}
This template demonstrates basic \LaTeX{} features including mathematical formulas, tables, and lists. The bibliography system is now set up to work seamlessly with Bibby AI's citation management. Feel free to modify it for your needs!

\bibliographystyle{plain}
\bibliography{cite}

\end{document}\documentclass{article}
\usepackage[utf8]{inputenc}
\usepackage{booktabs}
\usepackage{amsmath}
\usepackage{natbib}

\title{My Document}
\author{Your Name}
\date{\today}

\begin{document}

\maketitle

\section{Introduction}
This is a sample document demonstrating various \LaTeX{} features including mathematical equations, tables, and formatting. You can cite papers using the citation search feature, and they will be automatically added to your \texttt{cite.bib} file.

\section{Mathematical Expressions}
Here are some example equations:

\subsection{Quadratic Formula}
The quadratic formula is given by:
\[x = \frac{-b \pm \sqrt{b^2 - 4ac}}{2a}\]

\subsection{Matrix Example}
A sample matrix:
\[
\begin{pmatrix}
a_{11} & a_{12} & a_{13} \\
a_{21} & a_{22} & a_{23} \\
a_{31} & a_{32} & a_{33}
\end{pmatrix}
\]

\section{Tables}
Below is a sample table showing data:

\begin{table}[h]
\centering
\begin{tabular}{@{}lrrr@{}}
\toprule
\textbf{Category} & \textbf{Value 1} & \textbf{Value 2} & \textbf{Total} \\
\midrule
Item A & 10.5 & 20.3 & 30.8 \\
Item B & 15.2 & 25.1 & 40.3 \\
Item C & 22.1 & 12.7 & 34.8 \\
\bottomrule
\end{tabular}
\caption{Sample Data Table}
\label{tab:sample}
\end{table}

\section{Lists}
\subsection{Bullet Points}
\begin{itemize}
  \item First item
  \item Second item
  \item Third item with sub-items:
    \begin{itemize}
      \item Sub-item 1
      \item Sub-item 2
    \end{itemize}
\end{itemize}

\subsection{Numbered List}
\begin{enumerate}
  \item First step
  \item Second step
  \item Third step
\end{enumerate}

\section{Citations and References}
When you use the citation search feature, references will be automatically managed. For example, you might cite a paper like this: \cite{example2024}. All citations will be automatically added to your \texttt{cite.bib} file.

\section{Conclusion}
This template demonstrates basic \LaTeX{} features including mathematical formulas, tables, and lists. The bibliography system is now set up to work seamlessly with Bibby AI's citation management. Feel free to modify it for your needs!

\bibliographystyle{plain}
\bibliography{cite}

\end{document}\documentclass{article}
\usepackage[utf8]{inputenc}
\usepackage{booktabs}
\usepackage{amsmath}
\usepackage{natbib}

\title{My Document}
\author{Your Name}
\date{\today}

\begin{document}

\maketitle

\section{Introduction}
This is a sample document demonstrating various \LaTeX{} features including mathematical equations, tables, and formatting. You can cite papers using the citation search feature, and they will be automatically added to your \texttt{cite.bib} file.

\section{Mathematical Expressions}
Here are some example equations:

\subsection{Quadratic Formula}
The quadratic formula is given by:
\[x = \frac{-b \pm \sqrt{b^2 - 4ac}}{2a}\]

\subsection{Matrix Example}
A sample matrix:
\[
\begin{pmatrix}
a_{11} & a_{12} & a_{13} \\
a_{21} & a_{22} & a_{23} \\
a_{31} & a_{32} & a_{33}
\end{pmatrix}
\]

\section{Tables}
Below is a sample table showing data:

\begin{table}[h]
\centering
\begin{tabular}{@{}lrrr@{}}
\toprule
\textbf{Category} & \textbf{Value 1} & \textbf{Value 2} & \textbf{Total} \\
\midrule
Item A & 10.5 & 20.3 & 30.8 \\
Item B & 15.2 & 25.1 & 40.3 \\
Item C & 22.1 & 12.7 & 34.8 \\
\bottomrule
\end{tabular}
\caption{Sample Data Table}
\label{tab:sample}
\end{table}

\section{Lists}
\subsection{Bullet Points}
\begin{itemize}
  \item First item
  \item Second item
  \item Third item with sub-items:
    \begin{itemize}
      \item Sub-item 1
      \item Sub-item 2
    \end{itemize}
\end{itemize}

\subsection{Numbered List}
\begin{enumerate}
  \item First step
  \item Second step
  \item Third step
\end{enumerate}

\section{Citations and References}
When you use the citation search feature, references will be automatically managed. For example, you might cite a paper like this: \cite{example2024}. All citations will be automatically added to your \texttt{cite.bib} file.

\section{Conclusion}
This template demonstrates basic \LaTeX{} features including mathematical formulas, tables, and lists. The bibliography system is now set up to work seamlessly with Bibby AI's citation management. Feel free to modify it for your needs!

\bibliographystyle{plain}
\bibliography{cite}

\end{document}\documentclass{article}
\usepackage[utf8]{inputenc}
\usepackage{booktabs}
\usepackage{amsmath}
\usepackage{natbib}

\title{My Document}
\author{Your Name}
\date{\today}

\begin{document}

\maketitle

\section{Introduction}
This is a sample document demonstrating various \LaTeX{} features including mathematical equations, tables, and formatting. You can cite papers using the citation search feature, and they will be automatically added to your \texttt{cite.bib} file.

\section{Mathematical Expressions}
Here are some example equations:

\subsection{Quadratic Formula}
The quadratic formula is given by:
\[x = \frac{-b \pm \sqrt{b^2 - 4ac}}{2a}\]

\subsection{Matrix Example}
A sample matrix:
\[
\begin{pmatrix}
a_{11} & a_{12} & a_{13} \\
a_{21} & a_{22} & a_{23} \\
a_{31} & a_{32} & a_{33}
\end{pmatrix}
\]

\section{Tables}
Below is a sample table showing data:

\begin{table}[h]
\centering
\begin{tabular}{@{}lrrr@{}}
\toprule
\textbf{Category} & \textbf{Value 1} & \textbf{Value 2} & \textbf{Total} \\
\midrule
Item A & 10.5 & 20.3 & 30.8 \\
Item B & 15.2 & 25.1 & 40.3 \\
Item C & 22.1 & 12.7 & 34.8 \\
\bottomrule
\end{tabular}
\caption{Sample Data Table}
\label{tab:sample}
\end{table}

\section{Lists}
\subsection{Bullet Points}
\begin{itemize}
  \item First item
  \item Second item
  \item Third item with sub-items:
    \begin{itemize}
      \item Sub-item 1
      \item Sub-item 2
    \end{itemize}
\end{itemize}

\subsection{Numbered List}
\begin{enumerate}
  \item First step
  \item Second step
  \item Third step
\end{enumerate}

\section{Citations and References}
When you use the citation search feature, references will be automatically managed. For example, you might cite a paper like this: \cite{example2024}. All citations will be automatically added to your \texttt{cite.bib} file.

\section{Conclusion}
This template demonstrates basic \LaTeX{} features including mathematical formulas, tables, and lists. The bibliography system is now set up to work seamlessly with Bibby AI's citation management. Feel free to modify it for your needs!

\bibliographystyle{plain}
\bibliography{cite}

\end{document}\documentclass{article}
\usepackage[utf8]{inputenc}
\usepackage{booktabs}
\usepackage{amsmath}
\usepackage{natbib}

\title{My Document}
\author{Your Name}
\date{\today}

\begin{document}

\maketitle

\section{Introduction}
This is a sample document demonstrating various \LaTeX{} features including mathematical equations, tables, and formatting. You can cite papers using the citation search feature, and they will be automatically added to your \texttt{cite.bib} file.

\section{Mathematical Expressions}
Here are some example equations:

\subsection{Quadratic Formula}
The quadratic formula is given by:
\[x = \frac{-b \pm \sqrt{b^2 - 4ac}}{2a}\]

\subsection{Matrix Example}
A sample matrix:
\[
\begin{pmatrix}
a_{11} & a_{12} & a_{13} \\
a_{21} & a_{22} & a_{23} \\
a_{31} & a_{32} & a_{33}
\end{pmatrix}
\]

\section{Tables}
Below is a sample table showing data:

\begin{table}[h]
\centering
\begin{tabular}{@{}lrrr@{}}
\toprule
\textbf{Category} & \textbf{Value 1} & \textbf{Value 2} & \textbf{Total} \\
\midrule
Item A & 10.5 & 20.3 & 30.8 \\
Item B & 15.2 & 25.1 & 40.3 \\
Item C & 22.1 & 12.7 & 34.8 \\
\bottomrule
\end{tabular}
\caption{Sample Data Table}
\label{tab:sample}
\end{table}

\section{Lists}
\subsection{Bullet Points}
\begin{itemize}
  \item First item
  \item Second item
  \item Third item with sub-items:
    \begin{itemize}
      \item Sub-item 1
      \item Sub-item 2
    \end{itemize}
\end{itemize}

\subsection{Numbered List}
\begin{enumerate}
  \item First step
  \item Second step
  \item Third step
\end{enumerate}

\section{Citations and References}
When you use the citation search feature, references will be automatically managed. For example, you might cite a paper like this: \cite{example2024}. All citations will be automatically added to your \texttt{cite.bib} file.

\section{Conclusion}
This template demonstrates basic \LaTeX{} features including mathematical formulas, tables, and lists. The bibliography system is now set up to work seamlessly with Bibby AI's citation management. Feel free to modify it for your needs!

\bibliographystyle{plain}
\bibliography{cite}

\end{document}\documentclass{article}
\usepackage[utf8]{inputenc}
\usepackage{booktabs}
\usepackage{amsmath}
\usepackage{natbib}

\title{My Document}
\author{Your Name}
\date{\today}

\begin{document}

\maketitle

\section{Introduction}
This is a sample document demonstrating various \LaTeX{} features including mathematical equations, tables, and formatting. You can cite papers using the citation search feature, and they will be automatically added to your \texttt{cite.bib} file.

\section{Mathematical Expressions}
Here are some example equations:

\subsection{Quadratic Formula}
The quadratic formula is given by:
\[x = \frac{-b \pm \sqrt{b^2 - 4ac}}{2a}\]

\subsection{Matrix Example}
A sample matrix:
\[
\begin{pmatrix}
a_{11} & a_{12} & a_{13} \\
a_{21} & a_{22} & a_{23} \\
a_{31} & a_{32} & a_{33}
\end{pmatrix}
\]

\section{Tables}
Below is a sample table showing data:

\begin{table}[h]
\centering
\begin{tabular}{@{}lrrr@{}}
\toprule
\textbf{Category} & \textbf{Value 1} & \textbf{Value 2} & \textbf{Total} \\
\midrule
Item A & 10.5 & 20.3 & 30.8 \\
Item B & 15.2 & 25.1 & 40.3 \\
Item C & 22.1 & 12.7 & 34.8 \\
\bottomrule
\end{tabular}
\caption{Sample Data Table}
\label{tab:sample}
\end{table}

\section{Lists}
\subsection{Bullet Points}
\begin{itemize}
  \item First item
  \item Second item
  \item Third item with sub-items:
    \begin{itemize}
      \item Sub-item 1
      \item Sub-item 2
    \end{itemize}
\end{itemize}

\subsection{Numbered List}
\begin{enumerate}
  \item First step
  \item Second step
  \item Third step
\end{enumerate}

\section{Citations and References}
When you use the citation search feature, references will be automatically managed. For example, you might cite a paper like this: \cite{example2024}. All citations will be automatically added to your \texttt{cite.bib} file.

\section{Conclusion}
This template demonstrates basic \LaTeX{} features including mathematical formulas, tables, and lists. The bibliography system is now set up to work seamlessly with Bibby AI's citation management. Feel free to modify it for your needs!

\bibliographystyle{plain}
\bibliography{cite}

\end{document}\documentclass{article}
\usepackage[utf8]{inputenc}
\usepackage{booktabs}
\usepackage{amsmath}
\usepackage{natbib}

\title{My Document}
\author{Your Name}
\date{\today}

\begin{document}

\maketitle

\section{Introduction}
This is a sample document demonstrating various \LaTeX{} features including mathematical equations, tables, and formatting. You can cite papers using the citation search feature, and they will be automatically added to your \texttt{cite.bib} file.

\section{Mathematical Expressions}
Here are some example equations:

\subsection{Quadratic Formula}
The quadratic formula is given by:
\[x = \frac{-b \pm \sqrt{b^2 - 4ac}}{2a}\]

\subsection{Matrix Example}
A sample matrix:
\[
\begin{pmatrix}
a_{11} & a_{12} & a_{13} \\
a_{21} & a_{22} & a_{23} \\
a_{31} & a_{32} & a_{33}
\end{pmatrix}
\]

\section{Tables}
Below is a sample table showing data:

\begin{table}[h]
\centering
\begin{tabular}{@{}lrrr@{}}
\toprule
\textbf{Category} & \textbf{Value 1} & \textbf{Value 2} & \textbf{Total} \\
\midrule
Item A & 10.5 & 20.3 & 30.8 \\
Item B & 15.2 & 25.1 & 40.3 \\
Item C & 22.1 & 12.7 & 34.8 \\
\bottomrule
\end{tabular}
\caption{Sample Data Table}
\label{tab:sample}
\end{table}

\section{Lists}
\subsection{Bullet Points}
\begin{itemize}
  \item First item
  \item Second item
  \item Third item with sub-items:
    \begin{itemize}
      \item Sub-item 1
      \item Sub-item 2
    \end{itemize}
\end{itemize}

\subsection{Numbered List}
\begin{enumerate}
  \item First step
  \item Second step
  \item Third step
\end{enumerate}

\section{Citations and References}
When you use the citation search feature, references will be automatically managed. For example, you might cite a paper like this: \cite{example2024}. All citations will be automatically added to your \texttt{cite.bib} file.

\section{Conclusion}
This template demonstrates basic \LaTeX{} features including mathematical formulas, tables, and lists. The bibliography system is now set up to work seamlessly with Bibby AI's citation management. Feel free to modify it for your needs!

\bibliographystyle{plain}
\bibliography{cite}

\end{document}\documentclass{article}
\usepackage[utf8]{inputenc}
\usepackage{booktabs}
\usepackage{amsmath}
\usepackage{natbib}

\title{My Document}
\author{Your Name}
\date{\today}

\begin{document}

\maketitle

\section{Introduction}
This is a sample document demonstrating various \LaTeX{} features including mathematical equations, tables, and formatting. You can cite papers using the citation search feature, and they will be automatically added to your \texttt{cite.bib} file.

\section{Mathematical Expressions}
Here are some example equations:

\subsection{Quadratic Formula}
The quadratic formula is given by:
\[x = \frac{-b \pm \sqrt{b^2 - 4ac}}{2a}\]

\subsection{Matrix Example}
A sample matrix:
\[
\begin{pmatrix}
a_{11} & a_{12} & a_{13} \\
a_{21} & a_{22} & a_{23} \\
a_{31} & a_{32} & a_{33}
\end{pmatrix}
\]

\section{Tables}
Below is a sample table showing data:

\begin{table}[h]
\centering
\begin{tabular}{@{}lrrr@{}}
\toprule
\textbf{Category} & \textbf{Value 1} & \textbf{Value 2} & \textbf{Total} \\
\midrule
Item A & 10.5 & 20.3 & 30.8 \\
Item B & 15.2 & 25.1 & 40.3 \\
Item C & 22.1 & 12.7 & 34.8 \\
\bottomrule
\end{tabular}
\caption{Sample Data Table}
\label{tab:sample}
\end{table}

\section{Lists}
\subsection{Bullet Points}
\begin{itemize}
  \item First item
  \item Second item
  \item Third item with sub-items:
    \begin{itemize}
      \item Sub-item 1
      \item Sub-item 2
    \end{itemize}
\end{itemize}

\subsection{Numbered List}
\begin{enumerate}
  \item First step
  \item Second step
  \item Third step
\end{enumerate}

\section{Citations and References}
When you use the citation search feature, references will be automatically managed. For example, you might cite a paper like this: \cite{example2024}. All citations will be automatically added to your \texttt{cite.bib} file.

\section{Conclusion}
This template demonstrates basic \LaTeX{} features including mathematical formulas, tables, and lists. The bibliography system is now set up to work seamlessly with Bibby AI's citation management. Feel free to modify it for your needs!

\bibliographystyle{plain}
\bibliography{cite}

\end{document}\documentclass{article}
\usepackage[utf8]{inputenc}
\usepackage{booktabs}
\usepackage{amsmath}
\usepackage{natbib}

\title{My Document}
\author{Your Name}
\date{\today}

\begin{document}

\maketitle

\section{Introduction}
This is a sample document demonstrating various \LaTeX{} features including mathematical equations, tables, and formatting. You can cite papers using the citation search feature, and they will be automatically added to your \texttt{cite.bib} file.

\section{Mathematical Expressions}
Here are some example equations:

\subsection{Quadratic Formula}
The quadratic formula is given by:
\[x = \frac{-b \pm \sqrt{b^2 - 4ac}}{2a}\]

\subsection{Matrix Example}
A sample matrix:
\[
\begin{pmatrix}
a_{11} & a_{12} & a_{13} \\
a_{21} & a_{22} & a_{23} \\
a_{31} & a_{32} & a_{33}
\end{pmatrix}
\]

\section{Tables}
Below is a sample table showing data:

\begin{table}[h]
\centering
\begin{tabular}{@{}lrrr@{}}
\toprule
\textbf{Category} & \textbf{Value 1} & \textbf{Value 2} & \textbf{Total} \\
\midrule
Item A & 10.5 & 20.3 & 30.8 \\
Item B & 15.2 & 25.1 & 40.3 \\
Item C & 22.1 & 12.7 & 34.8 \\
\bottomrule
\end{tabular}
\caption{Sample Data Table}
\label{tab:sample}
\end{table}

\section{Lists}
\subsection{Bullet Points}
\begin{itemize}
  \item First item
  \item Second item
  \item Third item with sub-items:
    \begin{itemize}
      \item Sub-item 1
      \item Sub-item 2
    \end{itemize}
\end{itemize}

\subsection{Numbered List}
\begin{enumerate}
  \item First step
  \item Second step
  \item Third step
\end{enumerate}

\section{Citations and References}
When you use the citation search feature, references will be automatically managed. For example, you might cite a paper like this: \cite{example2024}. All citations will be automatically added to your \texttt{cite.bib} file.

\section{Conclusion}
This template demonstrates basic \LaTeX{} features including mathematical formulas, tables, and lists. The bibliography system is now set up to work seamlessly with Bibby AI's citation management. Feel free to modify it for your needs!

\bibliographystyle{plain}
\bibliography{cite}

\end{document}\documentclass{article}
\usepackage[utf8]{inputenc}
\usepackage{booktabs}
\usepackage{amsmath}
\usepackage{natbib}

\title{My Document}
\author{Your Name}
\date{\today}

\begin{document}

\maketitle

\section{Introduction}
This is a sample document demonstrating various \LaTeX{} features including mathematical equations, tables, and formatting. You can cite papers using the citation search feature, and they will be automatically added to your \texttt{cite.bib} file.

\section{Mathematical Expressions}
Here are some example equations:

\subsection{Quadratic Formula}
The quadratic formula is given by:
\[x = \frac{-b \pm \sqrt{b^2 - 4ac}}{2a}\]

\subsection{Matrix Example}
A sample matrix:
\[
\begin{pmatrix}
a_{11} & a_{12} & a_{13} \\
a_{21} & a_{22} & a_{23} \\
a_{31} & a_{32} & a_{33}
\end{pmatrix}
\]

\section{Tables}
Below is a sample table showing data:

\begin{table}[h]
\centering
\begin{tabular}{@{}lrrr@{}}
\toprule
\textbf{Category} & \textbf{Value 1} & \textbf{Value 2} & \textbf{Total} \\
\midrule
Item A & 10.5 & 20.3 & 30.8 \\
Item B & 15.2 & 25.1 & 40.3 \\
Item C & 22.1 & 12.7 & 34.8 \\
\bottomrule
\end{tabular}
\caption{Sample Data Table}
\label{tab:sample}
\end{table}

\section{Lists}
\subsection{Bullet Points}
\begin{itemize}
  \item First item
  \item Second item
  \item Third item with sub-items:
    \begin{itemize}
      \item Sub-item 1
      \item Sub-item 2
    \end{itemize}
\end{itemize}

\subsection{Numbered List}
\begin{enumerate}
  \item First step
  \item Second step
  \item Third step
\end{enumerate}

\section{Citations and References}
When you use the citation search feature, references will be automatically managed. For example, you might cite a paper like this: \cite{example2024}. All citations will be automatically added to your \texttt{cite.bib} file.

\section{Conclusion}
This template demonstrates basic \LaTeX{} features including mathematical formulas, tables, and lists. The bibliography system is now set up to work seamlessly with Bibby AI's citation management. Feel free to modify it for your needs!

\bibliographystyle{plain}
\bibliography{cite}

\end{document}\documentclass{article}
\usepackage[utf8]{inputenc}
\usepackage{booktabs}
\usepackage{amsmath}
\usepackage{natbib}

\title{My Document}
\author{Your Name}
\date{\today}

\begin{document}

\maketitle

\section{Introduction}
This is a sample document demonstrating various \LaTeX{} features including mathematical equations, tables, and formatting. You can cite papers using the citation search feature, and they will be automatically added to your \texttt{cite.bib} file.

\section{Mathematical Expressions}
Here are some example equations:

\subsection{Quadratic Formula}
The quadratic formula is given by:
\[x = \frac{-b \pm \sqrt{b^2 - 4ac}}{2a}\]

\subsection{Matrix Example}
A sample matrix:
\[
\begin{pmatrix}
a_{11} & a_{12} & a_{13} \\
a_{21} & a_{22} & a_{23} \\
a_{31} & a_{32} & a_{33}
\end{pmatrix}
\]

\section{Tables}
Below is a sample table showing data:

\begin{table}[h]
\centering
\begin{tabular}{@{}lrrr@{}}
\toprule
\textbf{Category} & \textbf{Value 1} & \textbf{Value 2} & \textbf{Total} \\
\midrule
Item A & 10.5 & 20.3 & 30.8 \\
Item B & 15.2 & 25.1 & 40.3 \\
Item C & 22.1 & 12.7 & 34.8 \\
\bottomrule
\end{tabular}
\caption{Sample Data Table}
\label{tab:sample}
\end{table}

\section{Lists}
\subsection{Bullet Points}
\begin{itemize}
  \item First item
  \item Second item
  \item Third item with sub-items:
    \begin{itemize}
      \item Sub-item 1
      \item Sub-item 2
    \end{itemize}
\end{itemize}

\subsection{Numbered List}
\begin{enumerate}
  \item First step
  \item Second step
  \item Third step
\end{enumerate}

\section{Citations and References}
When you use the citation search feature, references will be automatically managed. For example, you might cite a paper like this: \cite{example2024}. All citations will be automatically added to your \texttt{cite.bib} file.

\section{Conclusion}
This template demonstrates basic \LaTeX{} features including mathematical formulas, tables, and lists. The bibliography system is now set up to work seamlessly with Bibby AI's citation management. Feel free to modify it for your needs!

\bibliographystyle{plain}
\bibliography{cite}

\end{document}\documentclass{article}
\usepackage[utf8]{inputenc}
\usepackage{booktabs}
\usepackage{amsmath}
\usepackage{natbib}

\title{My Document}
\author{Your Name}
\date{\today}

\begin{document}

\maketitle

\section{Introduction}
This is a sample document demonstrating various \LaTeX{} features including mathematical equations, tables, and formatting. You can cite papers using the citation search feature, and they will be automatically added to your \texttt{cite.bib} file.

\section{Mathematical Expressions}
Here are some example equations:

\subsection{Quadratic Formula}
The quadratic formula is given by:
\[x = \frac{-b \pm \sqrt{b^2 - 4ac}}{2a}\]

\subsection{Matrix Example}
A sample matrix:
\[
\begin{pmatrix}
a_{11} & a_{12} & a_{13} \\
a_{21} & a_{22} & a_{23} \\
a_{31} & a_{32} & a_{33}
\end{pmatrix}
\]

\section{Tables}
Below is a sample table showing data:

\begin{table}[h]
\centering
\begin{tabular}{@{}lrrr@{}}
\toprule
\textbf{Category} & \textbf{Value 1} & \textbf{Value 2} & \textbf{Total} \\
\midrule
Item A & 10.5 & 20.3 & 30.8 \\
Item B & 15.2 & 25.1 & 40.3 \\
Item C & 22.1 & 12.7 & 34.8 \\
\bottomrule
\end{tabular}
\caption{Sample Data Table}
\label{tab:sample}
\end{table}

\section{Lists}
\subsection{Bullet Points}
\begin{itemize}
  \item First item
  \item Second item
  \item Third item with sub-items:
    \begin{itemize}
      \item Sub-item 1
      \item Sub-item 2
    \end{itemize}
\end{itemize}

\subsection{Numbered List}
\begin{enumerate}
  \item First step
  \item Second step
  \item Third step
\end{enumerate}

\section{Citations and References}
When you use the citation search feature, references will be automatically managed. For example, you might cite a paper like this: \cite{example2024}. All citations will be automatically added to your \texttt{cite.bib} file.

\section{Conclusion}
This template demonstrates basic \LaTeX{} features including mathematical formulas, tables, and lists. The bibliography system is now set up to work seamlessly with Bibby AI's citation management. Feel free to modify it for your needs!

\bibliographystyle{plain}
\bibliography{cite}

\end{document}\documentclass{article}
\usepackage[utf8]{inputenc}
\usepackage{booktabs}
\usepackage{amsmath}
\usepackage{natbib}

\title{My Document}
\author{Your Name}
\date{\today}

\begin{document}

\maketitle

\section{Introduction}
This is a sample document demonstrating various \LaTeX{} features including mathematical equations, tables, and formatting. You can cite papers using the citation search feature, and they will be automatically added to your \texttt{cite.bib} file.

\section{Mathematical Expressions}
Here are some example equations:

\subsection{Quadratic Formula}
The quadratic formula is given by:
\[x = \frac{-b \pm \sqrt{b^2 - 4ac}}{2a}\]

\subsection{Matrix Example}
A sample matrix:
\[
\begin{pmatrix}
a_{11} & a_{12} & a_{13} \\
a_{21} & a_{22} & a_{23} \\
a_{31} & a_{32} & a_{33}
\end{pmatrix}
\]

\section{Tables}
Below is a sample table showing data:

\begin{table}[h]
\centering
\begin{tabular}{@{}lrrr@{}}
\toprule
\textbf{Category} & \textbf{Value 1} & \textbf{Value 2} & \textbf{Total} \\
\midrule
Item A & 10.5 & 20.3 & 30.8 \\
Item B & 15.2 & 25.1 & 40.3 \\
Item C & 22.1 & 12.7 & 34.8 \\
\bottomrule
\end{tabular}
\caption{Sample Data Table}
\label{tab:sample}
\end{table}

\section{Lists}
\subsection{Bullet Points}
\begin{itemize}
  \item First item
  \item Second item
  \item Third item with sub-items:
    \begin{itemize}
      \item Sub-item 1
      \item Sub-item 2
    \end{itemize}
\end{itemize}

\subsection{Numbered List}
\begin{enumerate}
  \item First step
  \item Second step
  \item Third step
\end{enumerate}

\section{Citations and References}
When you use the citation search feature, references will be automatically managed. For example, you might cite a paper like this: \cite{example2024}. All citations will be automatically added to your \texttt{cite.bib} file.

\section{Conclusion}
This template demonstrates basic \LaTeX{} features including mathematical formulas, tables, and lists. The bibliography system is now set up to work seamlessly with Bibby AI's citation management. Feel free to modify it for your needs!

\bibliographystyle{plain}
\bibliography{cite}

\end{document}\documentclass{article}
\usepackage[utf8]{inputenc}
\usepackage{booktabs}
\usepackage{amsmath}
\usepackage{natbib}

\title{My Document}
\author{Your Name}
\date{\today}

\begin{document}

\maketitle

\section{Introduction}
This is a sample document demonstrating various \LaTeX{} features including mathematical equations, tables, and formatting. You can cite papers using the citation search feature, and they will be automatically added to your \texttt{cite.bib} file.

\section{Mathematical Expressions}
Here are some example equations:

\subsection{Quadratic Formula}
The quadratic formula is given by:
\[x = \frac{-b \pm \sqrt{b^2 - 4ac}}{2a}\]

\subsection{Matrix Example}
A sample matrix:
\[
\begin{pmatrix}
a_{11} & a_{12} & a_{13} \\
a_{21} & a_{22} & a_{23} \\
a_{31} & a_{32} & a_{33}
\end{pmatrix}
\]

\section{Tables}
Below is a sample table showing data:

\begin{table}[h]
\centering
\begin{tabular}{@{}lrrr@{}}
\toprule
\textbf{Category} & \textbf{Value 1} & \textbf{Value 2} & \textbf{Total} \\
\midrule
Item A & 10.5 & 20.3 & 30.8 \\
Item B & 15.2 & 25.1 & 40.3 \\
Item C & 22.1 & 12.7 & 34.8 \\
\bottomrule
\end{tabular}
\caption{Sample Data Table}
\label{tab:sample}
\end{table}

\section{Lists}
\subsection{Bullet Points}
\begin{itemize}
  \item First item
  \item Second item
  \item Third item with sub-items:
    \begin{itemize}
      \item Sub-item 1
      \item Sub-item 2
    \end{itemize}
\end{itemize}

\subsection{Numbered List}
\begin{enumerate}
  \item First step
  \item Second step
  \item Third step
\end{enumerate}

\section{Citations and References}
When you use the citation search feature, references will be automatically managed. For example, you might cite a paper like this: \cite{example2024}. All citations will be automatically added to your \texttt{cite.bib} file.

\section{Conclusion}
This template demonstrates basic \LaTeX{} features including mathematical formulas, tables, and lists. The bibliography system is now set up to work seamlessly with Bibby AI's citation management. Feel free to modify it for your needs!

\bibliographystyle{plain}
\bibliography{cite}

\end{document}\documentclass{article}
\usepackage[utf8]{inputenc}
\usepackage{booktabs}
\usepackage{amsmath}
\usepackage{natbib}

\title{My Document}
\author{Your Name}
\date{\today}

\begin{document}

\maketitle

\section{Introduction}
This is a sample document demonstrating various \LaTeX{} features including mathematical equations, tables, and formatting. You can cite papers using the citation search feature, and they will be automatically added to your \texttt{cite.bib} file.

\section{Mathematical Expressions}
Here are some example equations:

\subsection{Quadratic Formula}
The quadratic formula is given by:
\[x = \frac{-b \pm \sqrt{b^2 - 4ac}}{2a}\]

\subsection{Matrix Example}
A sample matrix:
\[
\begin{pmatrix}
a_{11} & a_{12} & a_{13} \\
a_{21} & a_{22} & a_{23} \\
a_{31} & a_{32} & a_{33}
\end{pmatrix}
\]

\section{Tables}
Below is a sample table showing data:

\begin{table}[h]
\centering
\begin{tabular}{@{}lrrr@{}}
\toprule
\textbf{Category} & \textbf{Value 1} & \textbf{Value 2} & \textbf{Total} \\
\midrule
Item A & 10.5 & 20.3 & 30.8 \\
Item B & 15.2 & 25.1 & 40.3 \\
Item C & 22.1 & 12.7 & 34.8 \\
\bottomrule
\end{tabular}
\caption{Sample Data Table}
\label{tab:sample}
\end{table}

\section{Lists}
\subsection{Bullet Points}
\begin{itemize}
  \item First item
  \item Second item
  \item Third item with sub-items:
    \begin{itemize}
      \item Sub-item 1
      \item Sub-item 2
    \end{itemize}
\end{itemize}

\subsection{Numbered List}
\begin{enumerate}
  \item First step
  \item Second step
  \item Third step
\end{enumerate}

\section{Citations and References}
When you use the citation search feature, references will be automatically managed. For example, you might cite a paper like this: \cite{example2024}. All citations will be automatically added to your \texttt{cite.bib} file.

\section{Conclusion}
This template demonstrates basic \LaTeX{} features including mathematical formulas, tables, and lists. The bibliography system is now set up to work seamlessly with Bibby AI's citation management. Feel free to modify it for your needs!

\bibliographystyle{plain}
\bibliography{cite}

\end{document}\documentclass{article}
\usepackage[utf8]{inputenc}
\usepackage{booktabs}
\usepackage{amsmath}
\usepackage{natbib}

\title{My Document}
\author{Your Name}
\date{\today}

\begin{document}

\maketitle

\section{Introduction}
This is a sample document demonstrating various \LaTeX{} features including mathematical equations, tables, and formatting. You can cite papers using the citation search feature, and they will be automatically added to your \texttt{cite.bib} file.

\section{Mathematical Expressions}
Here are some example equations:

\subsection{Quadratic Formula}
The quadratic formula is given by:
\[x = \frac{-b \pm \sqrt{b^2 - 4ac}}{2a}\]

\subsection{Matrix Example}
A sample matrix:
\[
\begin{pmatrix}
a_{11} & a_{12} & a_{13} \\
a_{21} & a_{22} & a_{23} \\
a_{31} & a_{32} & a_{33}
\end{pmatrix}
\]

\section{Tables}
Below is a sample table showing data:

\begin{table}[h]
\centering
\begin{tabular}{@{}lrrr@{}}
\toprule
\textbf{Category} & \textbf{Value 1} & \textbf{Value 2} & \textbf{Total} \\
\midrule
Item A & 10.5 & 20.3 & 30.8 \\
Item B & 15.2 & 25.1 & 40.3 \\
Item C & 22.1 & 12.7 & 34.8 \\
\bottomrule
\end{tabular}
\caption{Sample Data Table}
\label{tab:sample}
\end{table}

\section{Lists}
\subsection{Bullet Points}
\begin{itemize}
  \item First item
  \item Second item
  \item Third item with sub-items:
    \begin{itemize}
      \item Sub-item 1
      \item Sub-item 2
    \end{itemize}
\end{itemize}

\subsection{Numbered List}
\begin{enumerate}
  \item First step
  \item Second step
  \item Third step
\end{enumerate}

\section{Citations and References}
When you use the citation search feature, references will be automatically managed. For example, you might cite a paper like this: \cite{example2024}. All citations will be automatically added to your \texttt{cite.bib} file.

\section{Conclusion}
This template demonstrates basic \LaTeX{} features including mathematical formulas, tables, and lists. The bibliography system is now set up to work seamlessly with Bibby AI's citation management. Feel free to modify it for your needs!

\bibliographystyle{plain}
\bibliography{cite}

\end{document}\documentclass{article}
\usepackage[utf8]{inputenc}
\usepackage{booktabs}
\usepackage{amsmath}
\usepackage{natbib}

\title{My Document}
\author{Your Name}
\date{\today}

\begin{document}

\maketitle

\section{Introduction}
This is a sample document demonstrating various \LaTeX{} features including mathematical equations, tables, and formatting. You can cite papers using the citation search feature, and they will be automatically added to your \texttt{cite.bib} file.

\section{Mathematical Expressions}
Here are some example equations:

\subsection{Quadratic Formula}
The quadratic formula is given by:
\[x = \frac{-b \pm \sqrt{b^2 - 4ac}}{2a}\]

\subsection{Matrix Example}
A sample matrix:
\[
\begin{pmatrix}
a_{11} & a_{12} & a_{13} \\
a_{21} & a_{22} & a_{23} \\
a_{31} & a_{32} & a_{33}
\end{pmatrix}
\]

\section{Tables}
Below is a sample table showing data:

\begin{table}[h]
\centering
\begin{tabular}{@{}lrrr@{}}
\toprule
\textbf{Category} & \textbf{Value 1} & \textbf{Value 2} & \textbf{Total} \\
\midrule
Item A & 10.5 & 20.3 & 30.8 \\
Item B & 15.2 & 25.1 & 40.3 \\
Item C & 22.1 & 12.7 & 34.8 \\
\bottomrule
\end{tabular}
\caption{Sample Data Table}
\label{tab:sample}
\end{table}

\section{Lists}
\subsection{Bullet Points}
\begin{itemize}
  \item First item
  \item Second item
  \item Third item with sub-items:
    \begin{itemize}
      \item Sub-item 1
      \item Sub-item 2
    \end{itemize}
\end{itemize}

\subsection{Numbered List}
\begin{enumerate}
  \item First step
  \item Second step
  \item Third step
\end{enumerate}

\section{Citations and References}
When you use the citation search feature, references will be automatically managed. For example, you might cite a paper like this: \cite{example2024}. All citations will be automatically added to your \texttt{cite.bib} file.

\section{Conclusion}
This template demonstrates basic \LaTeX{} features including mathematical formulas, tables, and lists. The bibliography system is now set up to work seamlessly with Bibby AI's citation management. Feel free to modify it for your needs!

\bibliographystyle{plain}
\bibliography{cite}

\end{document}\documentclass{article}
\usepackage[utf8]{inputenc}
\usepackage{booktabs}
\usepackage{amsmath}
\usepackage{natbib}

\title{My Document}
\author{Your Name}
\date{\today}

\begin{document}

\maketitle

\section{Introduction}
This is a sample document demonstrating various \LaTeX{} features including mathematical equations, tables, and formatting. You can cite papers using the citation search feature, and they will be automatically added to your \texttt{cite.bib} file.

\section{Mathematical Expressions}
Here are some example equations:

\subsection{Quadratic Formula}
The quadratic formula is given by:
\[x = \frac{-b \pm \sqrt{b^2 - 4ac}}{2a}\]

\subsection{Matrix Example}
A sample matrix:
\[
\begin{pmatrix}
a_{11} & a_{12} & a_{13} \\
a_{21} & a_{22} & a_{23} \\
a_{31} & a_{32} & a_{33}
\end{pmatrix}
\]

\section{Tables}
Below is a sample table showing data:

\begin{table}[h]
\centering
\begin{tabular}{@{}lrrr@{}}
\toprule
\textbf{Category} & \textbf{Value 1} & \textbf{Value 2} & \textbf{Total} \\
\midrule
Item A & 10.5 & 20.3 & 30.8 \\
Item B & 15.2 & 25.1 & 40.3 \\
Item C & 22.1 & 12.7 & 34.8 \\
\bottomrule
\end{tabular}
\caption{Sample Data Table}
\label{tab:sample}
\end{table}

\section{Lists}
\subsection{Bullet Points}
\begin{itemize}
  \item First item
  \item Second item
  \item Third item with sub-items:
    \begin{itemize}
      \item Sub-item 1
      \item Sub-item 2
    \end{itemize}
\end{itemize}

\subsection{Numbered List}
\begin{enumerate}
  \item First step
  \item Second step
  \item Third step
\end{enumerate}

\section{Citations and References}
When you use the citation search feature, references will be automatically managed. For example, you might cite a paper like this: \cite{example2024}. All citations will be automatically added to your \texttt{cite.bib} file.

\section{Conclusion}
This template demonstrates basic \LaTeX{} features including mathematical formulas, tables, and lists. The bibliography system is now set up to work seamlessly with Bibby AI's citation management. Feel free to modify it for your needs!

\bibliographystyle{plain}
\bibliography{cite}

\end{document}\documentclass{article}
\usepackage[utf8]{inputenc}
\usepackage{booktabs}
\usepackage{amsmath}
\usepackage{natbib}

\title{My Document}
\author{Your Name}
\date{\today}

\begin{document}

\maketitle

\section{Introduction}
This is a sample document demonstrating various \LaTeX{} features including mathematical equations, tables, and formatting. You can cite papers using the citation search feature, and they will be automatically added to your \texttt{cite.bib} file.

\section{Mathematical Expressions}
Here are some example equations:

\subsection{Quadratic Formula}
The quadratic formula is given by:
\[x = \frac{-b \pm \sqrt{b^2 - 4ac}}{2a}\]

\subsection{Matrix Example}
A sample matrix:
\[
\begin{pmatrix}
a_{11} & a_{12} & a_{13} \\
a_{21} & a_{22} & a_{23} \\
a_{31} & a_{32} & a_{33}
\end{pmatrix}
\]

\section{Tables}
Below is a sample table showing data:

\begin{table}[h]
\centering
\begin{tabular}{@{}lrrr@{}}
\toprule
\textbf{Category} & \textbf{Value 1} & \textbf{Value 2} & \textbf{Total} \\
\midrule
Item A & 10.5 & 20.3 & 30.8 \\
Item B & 15.2 & 25.1 & 40.3 \\
Item C & 22.1 & 12.7 & 34.8 \\
\bottomrule
\end{tabular}
\caption{Sample Data Table}
\label{tab:sample}
\end{table}

\section{Lists}
\subsection{Bullet Points}
\begin{itemize}
  \item First item
  \item Second item
  \item Third item with sub-items:
    \begin{itemize}
      \item Sub-item 1
      \item Sub-item 2
    \end{itemize}
\end{itemize}

\subsection{Numbered List}
\begin{enumerate}
  \item First step
  \item Second step
  \item Third step
\end{enumerate}

\section{Citations and References}
When you use the citation search feature, references will be automatically managed. For example, you might cite a paper like this: \cite{example2024}. All citations will be automatically added to your \texttt{cite.bib} file.

\section{Conclusion}
This template demonstrates basic \LaTeX{} features including mathematical formulas, tables, and lists. The bibliography system is now set up to work seamlessly with Bibby AI's citation management. Feel free to modify it for your needs!

\bibliographystyle{plain}
\bibliography{cite}

\end{document}\documentclass{article}
\usepackage[utf8]{inputenc}
\usepackage{booktabs}
\usepackage{amsmath}
\usepackage{natbib}

\title{My Document}
\author{Your Name}
\date{\today}

\begin{document}

\maketitle

\section{Introduction}
This is a sample document demonstrating various \LaTeX{} features including mathematical equations, tables, and formatting. You can cite papers using the citation search feature, and they will be automatically added to your \texttt{cite.bib} file.

\section{Mathematical Expressions}
Here are some example equations:

\subsection{Quadratic Formula}
The quadratic formula is given by:
\[x = \frac{-b \pm \sqrt{b^2 - 4ac}}{2a}\]

\subsection{Matrix Example}
A sample matrix:
\[
\begin{pmatrix}
a_{11} & a_{12} & a_{13} \\
a_{21} & a_{22} & a_{23} \\
a_{31} & a_{32} & a_{33}
\end{pmatrix}
\]

\section{Tables}
Below is a sample table showing data:

\begin{table}[h]
\centering
\begin{tabular}{@{}lrrr@{}}
\toprule
\textbf{Category} & \textbf{Value 1} & \textbf{Value 2} & \textbf{Total} \\
\midrule
Item A & 10.5 & 20.3 & 30.8 \\
Item B & 15.2 & 25.1 & 40.3 \\
Item C & 22.1 & 12.7 & 34.8 \\
\bottomrule
\end{tabular}
\caption{Sample Data Table}
\label{tab:sample}
\end{table}

\section{Lists}
\subsection{Bullet Points}
\begin{itemize}
  \item First item
  \item Second item
  \item Third item with sub-items:
    \begin{itemize}
      \item Sub-item 1
      \item Sub-item 2
    \end{itemize}
\end{itemize}

\subsection{Numbered List}
\begin{enumerate}
  \item First step
  \item Second step
  \item Third step
\end{enumerate}

\section{Citations and References}
When you use the citation search feature, references will be automatically managed. For example, you might cite a paper like this: \cite{example2024}. All citations will be automatically added to your \texttt{cite.bib} file.

\section{Conclusion}
This template demonstrates basic \LaTeX{} features including mathematical formulas, tables, and lists. The bibliography system is now set up to work seamlessly with Bibby AI's citation management. Feel free to modify it for your needs!

\bibliographystyle{plain}
\bibliography{cite}

\end{document}\documentclass{article}
\usepackage[utf8]{inputenc}
\usepackage{booktabs}
\usepackage{amsmath}
\usepackage{natbib}

\title{My Document}
\author{Your Name}
\date{\today}

\begin{document}

\maketitle

\section{Introduction}
This is a sample document demonstrating various \LaTeX{} features including mathematical equations, tables, and formatting. You can cite papers using the citation search feature, and they will be automatically added to your \texttt{cite.bib} file.

\section{Mathematical Expressions}
Here are some example equations:

\subsection{Quadratic Formula}
The quadratic formula is given by:
\[x = \frac{-b \pm \sqrt{b^2 - 4ac}}{2a}\]

\subsection{Matrix Example}
A sample matrix:
\[
\begin{pmatrix}
a_{11} & a_{12} & a_{13} \\
a_{21} & a_{22} & a_{23} \\
a_{31} & a_{32} & a_{33}
\end{pmatrix}
\]

\section{Tables}
Below is a sample table showing data:

\begin{table}[h]
\centering
\begin{tabular}{@{}lrrr@{}}
\toprule
\textbf{Category} & \textbf{Value 1} & \textbf{Value 2} & \textbf{Total} \\
\midrule
Item A & 10.5 & 20.3 & 30.8 \\
Item B & 15.2 & 25.1 & 40.3 \\
Item C & 22.1 & 12.7 & 34.8 \\
\bottomrule
\end{tabular}
\caption{Sample Data Table}
\label{tab:sample}
\end{table}

\section{Lists}
\subsection{Bullet Points}
\begin{itemize}
  \item First item
  \item Second item
  \item Third item with sub-items:
    \begin{itemize}
      \item Sub-item 1
      \item Sub-item 2
    \end{itemize}
\end{itemize}

\subsection{Numbered List}
\begin{enumerate}
  \item First step
  \item Second step
  \item Third step
\end{enumerate}

\section{Citations and References}
When you use the citation search feature, references will be automatically managed. For example, you might cite a paper like this: \cite{example2024}. All citations will be automatically added to your \texttt{cite.bib} file.

\section{Conclusion}
This template demonstrates basic \LaTeX{} features including mathematical formulas, tables, and lists. The bibliography system is now set up to work seamlessly with Bibby AI's citation management. Feel free to modify it for your needs!

\bibliographystyle{plain}
\bibliography{cite}

\end{document}\documentclass{article}
\usepackage[utf8]{inputenc}
\usepackage{booktabs}
\usepackage{amsmath}
\usepackage{natbib}

\title{My Document}
\author{Your Name}
\date{\today}

\begin{document}

\maketitle

\section{Introduction}
This is a sample document demonstrating various \LaTeX{} features including mathematical equations, tables, and formatting. You can cite papers using the citation search feature, and they will be automatically added to your \texttt{cite.bib} file.

\section{Mathematical Expressions}
Here are some example equations:

\subsection{Quadratic Formula}
The quadratic formula is given by:
\[x = \frac{-b \pm \sqrt{b^2 - 4ac}}{2a}\]

\subsection{Matrix Example}
A sample matrix:
\[
\begin{pmatrix}
a_{11} & a_{12} & a_{13} \\
a_{21} & a_{22} & a_{23} \\
a_{31} & a_{32} & a_{33}
\end{pmatrix}
\]

\section{Tables}
Below is a sample table showing data:

\begin{table}[h]
\centering
\begin{tabular}{@{}lrrr@{}}
\toprule
\textbf{Category} & \textbf{Value 1} & \textbf{Value 2} & \textbf{Total} \\
\midrule
Item A & 10.5 & 20.3 & 30.8 \\
Item B & 15.2 & 25.1 & 40.3 \\
Item C & 22.1 & 12.7 & 34.8 \\
\bottomrule
\end{tabular}
\caption{Sample Data Table}
\label{tab:sample}
\end{table}

\section{Lists}
\subsection{Bullet Points}
\begin{itemize}
  \item First item
  \item Second item
  \item Third item with sub-items:
    \begin{itemize}
      \item Sub-item 1
      \item Sub-item 2
    \end{itemize}
\end{itemize}

\subsection{Numbered List}
\begin{enumerate}
  \item First step
  \item Second step
  \item Third step
\end{enumerate}

\section{Citations and References}
When you use the citation search feature, references will be automatically managed. For example, you might cite a paper like this: \cite{example2024}. All citations will be automatically added to your \texttt{cite.bib} file.

\section{Conclusion}
This template demonstrates basic \LaTeX{} features including mathematical formulas, tables, and lists. The bibliography system is now set up to work seamlessly with Bibby AI's citation management. Feel free to modify it for your needs!

\bibliographystyle{plain}
\bibliography{cite}

\end{document}\documentclass{article}
\usepackage[utf8]{inputenc}
\usepackage{booktabs}
\usepackage{amsmath}
\usepackage{natbib}

\title{My Document}
\author{Your Name}
\date{\today}

\begin{document}

\maketitle

\section{Introduction}
This is a sample document demonstrating various \LaTeX{} features including mathematical equations, tables, and formatting. You can cite papers using the citation search feature, and they will be automatically added to your \texttt{cite.bib} file.

\section{Mathematical Expressions}
Here are some example equations:

\subsection{Quadratic Formula}
The quadratic formula is given by:
\[x = \frac{-b \pm \sqrt{b^2 - 4ac}}{2a}\]

\subsection{Matrix Example}
A sample matrix:
\[
\begin{pmatrix}
a_{11} & a_{12} & a_{13} \\
a_{21} & a_{22} & a_{23} \\
a_{31} & a_{32} & a_{33}
\end{pmatrix}
\]

\section{Tables}
Below is a sample table showing data:

\begin{table}[h]
\centering
\begin{tabular}{@{}lrrr@{}}
\toprule
\textbf{Category} & \textbf{Value 1} & \textbf{Value 2} & \textbf{Total} \\
\midrule
Item A & 10.5 & 20.3 & 30.8 \\
Item B & 15.2 & 25.1 & 40.3 \\
Item C & 22.1 & 12.7 & 34.8 \\
\bottomrule
\end{tabular}
\caption{Sample Data Table}
\label{tab:sample}
\end{table}

\section{Lists}
\subsection{Bullet Points}
\begin{itemize}
  \item First item
  \item Second item
  \item Third item with sub-items:
    \begin{itemize}
      \item Sub-item 1
      \item Sub-item 2
    \end{itemize}
\end{itemize}

\subsection{Numbered List}
\begin{enumerate}
  \item First step
  \item Second step
  \item Third step
\end{enumerate}

\section{Citations and References}
When you use the citation search feature, references will be automatically managed. For example, you might cite a paper like this: \cite{example2024}. All citations will be automatically added to your \texttt{cite.bib} file.

\section{Conclusion}
This template demonstrates basic \LaTeX{} features including mathematical formulas, tables, and lists. The bibliography system is now set up to work seamlessly with Bibby AI's citation management. Feel free to modify it for your needs!

\bibliographystyle{plain}
\bibliography{cite}

\end{document}\documentclass{article}
\usepackage[utf8]{inputenc}
\usepackage{booktabs}
\usepackage{amsmath}
\usepackage{natbib}

\title{My Document}
\author{Your Name}
\date{\today}

\begin{document}

\maketitle

\section{Introduction}
This is a sample document demonstrating various \LaTeX{} features including mathematical equations, tables, and formatting. You can cite papers using the citation search feature, and they will be automatically added to your \texttt{cite.bib} file.

\section{Mathematical Expressions}
Here are some example equations:

\subsection{Quadratic Formula}
The quadratic formula is given by:
\[x = \frac{-b \pm \sqrt{b^2 - 4ac}}{2a}\]

\subsection{Matrix Example}
A sample matrix:
\[
\begin{pmatrix}
a_{11} & a_{12} & a_{13} \\
a_{21} & a_{22} & a_{23} \\
a_{31} & a_{32} & a_{33}
\end{pmatrix}
\]

\section{Tables}
Below is a sample table showing data:

\begin{table}[h]
\centering
\begin{tabular}{@{}lrrr@{}}
\toprule
\textbf{Category} & \textbf{Value 1} & \textbf{Value 2} & \textbf{Total} \\
\midrule
Item A & 10.5 & 20.3 & 30.8 \\
Item B & 15.2 & 25.1 & 40.3 \\
Item C & 22.1 & 12.7 & 34.8 \\
\bottomrule
\end{tabular}
\caption{Sample Data Table}
\label{tab:sample}
\end{table}

\section{Lists}
\subsection{Bullet Points}
\begin{itemize}
  \item First item
  \item Second item
  \item Third item with sub-items:
    \begin{itemize}
      \item Sub-item 1
      \item Sub-item 2
    \end{itemize}
\end{itemize}

\subsection{Numbered List}
\begin{enumerate}
  \item First step
  \item Second step
  \item Third step
\end{enumerate}

\section{Citations and References}
When you use the citation search feature, references will be automatically managed. For example, you might cite a paper like this: \cite{example2024}. All citations will be automatically added to your \texttt{cite.bib} file.

\section{Conclusion}
This template demonstrates basic \LaTeX{} features including mathematical formulas, tables, and lists. The bibliography system is now set up to work seamlessly with Bibby AI's citation management. Feel free to modify it for your needs!

\bibliographystyle{plain}
\bibliography{cite}

\end{document}\documentclass{article}
\usepackage[utf8]{inputenc}
\usepackage{booktabs}
\usepackage{amsmath}
\usepackage{natbib}

\title{My Document}
\author{Your Name}
\date{\today}

\begin{document}

\maketitle

\section{Introduction}
This is a sample document demonstrating various \LaTeX{} features including mathematical equations, tables, and formatting. You can cite papers using the citation search feature, and they will be automatically added to your \texttt{cite.bib} file.

\section{Mathematical Expressions}
Here are some example equations:

\subsection{Quadratic Formula}
The quadratic formula is given by:
\[x = \frac{-b \pm \sqrt{b^2 - 4ac}}{2a}\]

\subsection{Matrix Example}
A sample matrix:
\[
\begin{pmatrix}
a_{11} & a_{12} & a_{13} \\
a_{21} & a_{22} & a_{23} \\
a_{31} & a_{32} & a_{33}
\end{pmatrix}
\]

\section{Tables}
Below is a sample table showing data:

\begin{table}[h]
\centering
\begin{tabular}{@{}lrrr@{}}
\toprule
\textbf{Category} & \textbf{Value 1} & \textbf{Value 2} & \textbf{Total} \\
\midrule
Item A & 10.5 & 20.3 & 30.8 \\
Item B & 15.2 & 25.1 & 40.3 \\
Item C & 22.1 & 12.7 & 34.8 \\
\bottomrule
\end{tabular}
\caption{Sample Data Table}
\label{tab:sample}
\end{table}

\section{Lists}
\subsection{Bullet Points}
\begin{itemize}
  \item First item
  \item Second item
  \item Third item with sub-items:
    \begin{itemize}
      \item Sub-item 1
      \item Sub-item 2
    \end{itemize}
\end{itemize}

\subsection{Numbered List}
\begin{enumerate}
  \item First step
  \item Second step
  \item Third step
\end{enumerate}

\section{Citations and References}
When you use the citation search feature, references will be automatically managed. For example, you might cite a paper like this: \cite{example2024}. All citations will be automatically added to your \texttt{cite.bib} file.

\section{Conclusion}
This template demonstrates basic \LaTeX{} features including mathematical formulas, tables, and lists. The bibliography system is now set up to work seamlessly with Bibby AI's citation management. Feel free to modify it for your needs!

\bibliographystyle{plain}
\bibliography{cite}

\end{document}\documentclass{article}
\usepackage[utf8]{inputenc}
\usepackage{booktabs}
\usepackage{amsmath}
\usepackage{natbib}

\title{My Document}
\author{Your Name}
\date{\today}

\begin{document}

\maketitle

\section{Introduction}
This is a sample document demonstrating various \LaTeX{} features including mathematical equations, tables, and formatting. You can cite papers using the citation search feature, and they will be automatically added to your \texttt{cite.bib} file.

\section{Mathematical Expressions}
Here are some example equations:

\subsection{Quadratic Formula}
The quadratic formula is given by:
\[x = \frac{-b \pm \sqrt{b^2 - 4ac}}{2a}\]

\subsection{Matrix Example}
A sample matrix:
\[
\begin{pmatrix}
a_{11} & a_{12} & a_{13} \\
a_{21} & a_{22} & a_{23} \\
a_{31} & a_{32} & a_{33}
\end{pmatrix}
\]

\section{Tables}
Below is a sample table showing data:

\begin{table}[h]
\centering
\begin{tabular}{@{}lrrr@{}}
\toprule
\textbf{Category} & \textbf{Value 1} & \textbf{Value 2} & \textbf{Total} \\
\midrule
Item A & 10.5 & 20.3 & 30.8 \\
Item B & 15.2 & 25.1 & 40.3 \\
Item C & 22.1 & 12.7 & 34.8 \\
\bottomrule
\end{tabular}
\caption{Sample Data Table}
\label{tab:sample}
\end{table}

\section{Lists}
\subsection{Bullet Points}
\begin{itemize}
  \item First item
  \item Second item
  \item Third item with sub-items:
    \begin{itemize}
      \item Sub-item 1
      \item Sub-item 2
    \end{itemize}
\end{itemize}

\subsection{Numbered List}
\begin{enumerate}
  \item First step
  \item Second step
  \item Third step
\end{enumerate}

\section{Citations and References}
When you use the citation search feature, references will be automatically managed. For example, you might cite a paper like this: \cite{example2024}. All citations will be automatically added to your \texttt{cite.bib} file.

\section{Conclusion}
This template demonstrates basic \LaTeX{} features including mathematical formulas, tables, and lists. The bibliography system is now set up to work seamlessly with Bibby AI's citation management. Feel free to modify it for your needs!

\bibliographystyle{plain}
\bibliography{cite}

\end{document}\documentclass{article}
\usepackage[utf8]{inputenc}
\usepackage{booktabs}
\usepackage{amsmath}
\usepackage{natbib}

\title{My Document}
\author{Your Name}
\date{\today}

\begin{document}

\maketitle

\section{Introduction}
This is a sample document demonstrating various \LaTeX{} features including mathematical equations, tables, and formatting. You can cite papers using the citation search feature, and they will be automatically added to your \texttt{cite.bib} file.

\section{Mathematical Expressions}
Here are some example equations:

\subsection{Quadratic Formula}
The quadratic formula is given by:
\[x = \frac{-b \pm \sqrt{b^2 - 4ac}}{2a}\]

\subsection{Matrix Example}
A sample matrix:
\[
\begin{pmatrix}
a_{11} & a_{12} & a_{13} \\
a_{21} & a_{22} & a_{23} \\
a_{31} & a_{32} & a_{33}
\end{pmatrix}
\]

\section{Tables}
Below is a sample table showing data:

\begin{table}[h]
\centering
\begin{tabular}{@{}lrrr@{}}
\toprule
\textbf{Category} & \textbf{Value 1} & \textbf{Value 2} & \textbf{Total} \\
\midrule
Item A & 10.5 & 20.3 & 30.8 \\
Item B & 15.2 & 25.1 & 40.3 \\
Item C & 22.1 & 12.7 & 34.8 \\
\bottomrule
\end{tabular}
\caption{Sample Data Table}
\label{tab:sample}
\end{table}

\section{Lists}
\subsection{Bullet Points}
\begin{itemize}
  \item First item
  \item Second item
  \item Third item with sub-items:
    \begin{itemize}
      \item Sub-item 1
      \item Sub-item 2
    \end{itemize}
\end{itemize}

\subsection{Numbered List}
\begin{enumerate}
  \item First step
  \item Second step
  \item Third step
\end{enumerate}

\section{Citations and References}
When you use the citation search feature, references will be automatically managed. For example, you might cite a paper like this: \cite{example2024}. All citations will be automatically added to your \texttt{cite.bib} file.

\section{Conclusion}
This template demonstrates basic \LaTeX{} features including mathematical formulas, tables, and lists. The bibliography system is now set up to work seamlessly with Bibby AI's citation management. Feel free to modify it for your needs!

\bibliographystyle{plain}
\bibliography{cite}

\end{document}\documentclass{article}
\usepackage[utf8]{inputenc}
\usepackage{booktabs}
\usepackage{amsmath}
\usepackage{natbib}

\title{My Document}
\author{Your Name}
\date{\today}

\begin{document}

\maketitle

\section{Introduction}
This is a sample document demonstrating various \LaTeX{} features including mathematical equations, tables, and formatting. You can cite papers using the citation search feature, and they will be automatically added to your \texttt{cite.bib} file.

\section{Mathematical Expressions}
Here are some example equations:

\subsection{Quadratic Formula}
The quadratic formula is given by:
\[x = \frac{-b \pm \sqrt{b^2 - 4ac}}{2a}\]

\subsection{Matrix Example}
A sample matrix:
\[
\begin{pmatrix}
a_{11} & a_{12} & a_{13} \\
a_{21} & a_{22} & a_{23} \\
a_{31} & a_{32} & a_{33}
\end{pmatrix}
\]

\section{Tables}
Below is a sample table showing data:

\begin{table}[h]
\centering
\begin{tabular}{@{}lrrr@{}}
\toprule
\textbf{Category} & \textbf{Value 1} & \textbf{Value 2} & \textbf{Total} \\
\midrule
Item A & 10.5 & 20.3 & 30.8 \\
Item B & 15.2 & 25.1 & 40.3 \\
Item C & 22.1 & 12.7 & 34.8 \\
\bottomrule
\end{tabular}
\caption{Sample Data Table}
\label{tab:sample}
\end{table}

\section{Lists}
\subsection{Bullet Points}
\begin{itemize}
  \item First item
  \item Second item
  \item Third item with sub-items:
    \begin{itemize}
      \item Sub-item 1
      \item Sub-item 2
    \end{itemize}
\end{itemize}

\subsection{Numbered List}
\begin{enumerate}
  \item First step
  \item Second step
  \item Third step
\end{enumerate}

\section{Citations and References}
When you use the citation search feature, references will be automatically managed. For example, you might cite a paper like this: \cite{example2024}. All citations will be automatically added to your \texttt{cite.bib} file.

\section{Conclusion}
This template demonstrates basic \LaTeX{} features including mathematical formulas, tables, and lists. The bibliography system is now set up to work seamlessly with Bibby AI's citation management. Feel free to modify it for your needs!

\bibliographystyle{plain}
\bibliography{cite}

\end{document}\documentclass{article}
\usepackage[utf8]{inputenc}
\usepackage{booktabs}
\usepackage{amsmath}
\usepackage{natbib}

\title{My Document}
\author{Your Name}
\date{\today}

\begin{document}

\maketitle

\section{Introduction}
This is a sample document demonstrating various \LaTeX{} features including mathematical equations, tables, and formatting. You can cite papers using the citation search feature, and they will be automatically added to your \texttt{cite.bib} file.

\section{Mathematical Expressions}
Here are some example equations:

\subsection{Quadratic Formula}
The quadratic formula is given by:
\[x = \frac{-b \pm \sqrt{b^2 - 4ac}}{2a}\]

\subsection{Matrix Example}
A sample matrix:
\[
\begin{pmatrix}
a_{11} & a_{12} & a_{13} \\
a_{21} & a_{22} & a_{23} \\
a_{31} & a_{32} & a_{33}
\end{pmatrix}
\]

\section{Tables}
Below is a sample table showing data:

\begin{table}[h]
\centering
\begin{tabular}{@{}lrrr@{}}
\toprule
\textbf{Category} & \textbf{Value 1} & \textbf{Value 2} & \textbf{Total} \\
\midrule
Item A & 10.5 & 20.3 & 30.8 \\
Item B & 15.2 & 25.1 & 40.3 \\
Item C & 22.1 & 12.7 & 34.8 \\
\bottomrule
\end{tabular}
\caption{Sample Data Table}
\label{tab:sample}
\end{table}

\section{Lists}
\subsection{Bullet Points}
\begin{itemize}
  \item First item
  \item Second item
  \item Third item with sub-items:
    \begin{itemize}
      \item Sub-item 1
      \item Sub-item 2
    \end{itemize}
\end{itemize}

\subsection{Numbered List}
\begin{enumerate}
  \item First step
  \item Second step
  \item Third step
\end{enumerate}

\section{Citations and References}
When you use the citation search feature, references will be automatically managed. For example, you might cite a paper like this: \cite{example2024}. All citations will be automatically added to your \texttt{cite.bib} file.

\section{Conclusion}
This template demonstrates basic \LaTeX{} features including mathematical formulas, tables, and lists. The bibliography system is now set up to work seamlessly with Bibby AI's citation management. Feel free to modify it for your needs!

\bibliographystyle{plain}
\bibliography{cite}

\end{document}\documentclass{article}
\usepackage[utf8]{inputenc}
\usepackage{booktabs}
\usepackage{amsmath}
\usepackage{natbib}

\title{My Document}
\author{Your Name}
\date{\today}

\begin{document}

\maketitle

\section{Introduction}
This is a sample document demonstrating various \LaTeX{} features including mathematical equations, tables, and formatting. You can cite papers using the citation search feature, and they will be automatically added to your \texttt{cite.bib} file.

\section{Mathematical Expressions}
Here are some example equations:

\subsection{Quadratic Formula}
The quadratic formula is given by:
\[x = \frac{-b \pm \sqrt{b^2 - 4ac}}{2a}\]

\subsection{Matrix Example}
A sample matrix:
\[
\begin{pmatrix}
a_{11} & a_{12} & a_{13} \\
a_{21} & a_{22} & a_{23} \\
a_{31} & a_{32} & a_{33}
\end{pmatrix}
\]

\section{Tables}
Below is a sample table showing data:

\begin{table}[h]
\centering
\begin{tabular}{@{}lrrr@{}}
\toprule
\textbf{Category} & \textbf{Value 1} & \textbf{Value 2} & \textbf{Total} \\
\midrule
Item A & 10.5 & 20.3 & 30.8 \\
Item B & 15.2 & 25.1 & 40.3 \\
Item C & 22.1 & 12.7 & 34.8 \\
\bottomrule
\end{tabular}
\caption{Sample Data Table}
\label{tab:sample}
\end{table}

\section{Lists}
\subsection{Bullet Points}
\begin{itemize}
  \item First item
  \item Second item
  \item Third item with sub-items:
    \begin{itemize}
      \item Sub-item 1
      \item Sub-item 2
    \end{itemize}
\end{itemize}

\subsection{Numbered List}
\begin{enumerate}
  \item First step
  \item Second step
  \item Third step
\end{enumerate}

\section{Citations and References}
When you use the citation search feature, references will be automatically managed. For example, you might cite a paper like this: \cite{example2024}. All citations will be automatically added to your \texttt{cite.bib} file.

\section{Conclusion}
This template demonstrates basic \LaTeX{} features including mathematical formulas, tables, and lists. The bibliography system is now set up to work seamlessly with Bibby AI's citation management. Feel free to modify it for your needs!

\bibliographystyle{plain}
\bibliography{cite}

\end{document}\documentclass{article}
\usepackage[utf8]{inputenc}
\usepackage{booktabs}
\usepackage{amsmath}
\usepackage{natbib}

\title{My Document}
\author{Your Name}
\date{\today}

\begin{document}

\maketitle

\section{Introduction}
This is a sample document demonstrating various \LaTeX{} features including mathematical equations, tables, and formatting. You can cite papers using the citation search feature, and they will be automatically added to your \texttt{cite.bib} file.

\section{Mathematical Expressions}
Here are some example equations:

\subsection{Quadratic Formula}
The quadratic formula is given by:
\[x = \frac{-b \pm \sqrt{b^2 - 4ac}}{2a}\]

\subsection{Matrix Example}
A sample matrix:
\[
\begin{pmatrix}
a_{11} & a_{12} & a_{13} \\
a_{21} & a_{22} & a_{23} \\
a_{31} & a_{32} & a_{33}
\end{pmatrix}
\]

\section{Tables}
Below is a sample table showing data:

\begin{table}[h]
\centering
\begin{tabular}{@{}lrrr@{}}
\toprule
\textbf{Category} & \textbf{Value 1} & \textbf{Value 2} & \textbf{Total} \\
\midrule
Item A & 10.5 & 20.3 & 30.8 \\
Item B & 15.2 & 25.1 & 40.3 \\
Item C & 22.1 & 12.7 & 34.8 \\
\bottomrule
\end{tabular}
\caption{Sample Data Table}
\label{tab:sample}
\end{table}

\section{Lists}
\subsection{Bullet Points}
\begin{itemize}
  \item First item
  \item Second item
  \item Third item with sub-items:
    \begin{itemize}
      \item Sub-item 1
      \item Sub-item 2
    \end{itemize}
\end{itemize}

\subsection{Numbered List}
\begin{enumerate}
  \item First step
  \item Second step
  \item Third step
\end{enumerate}

\section{Citations and References}
When you use the citation search feature, references will be automatically managed. For example, you might cite a paper like this: \cite{example2024}. All citations will be automatically added to your \texttt{cite.bib} file.

\section{Conclusion}
This template demonstrates basic \LaTeX{} features including mathematical formulas, tables, and lists. The bibliography system is now set up to work seamlessly with Bibby AI's citation management. Feel free to modify it for your needs!

\bibliographystyle{plain}
\bibliography{cite}

\end{document}\documentclass{article}
\usepackage[utf8]{inputenc}
\usepackage{booktabs}
\usepackage{amsmath}
\usepackage{natbib}

\title{My Document}
\author{Your Name}
\date{\today}

\begin{document}

\maketitle

\section{Introduction}
This is a sample document demonstrating various \LaTeX{} features including mathematical equations, tables, and formatting. You can cite papers using the citation search feature, and they will be automatically added to your \texttt{cite.bib} file.

\section{Mathematical Expressions}
Here are some example equations:

\subsection{Quadratic Formula}
The quadratic formula is given by:
\[x = \frac{-b \pm \sqrt{b^2 - 4ac}}{2a}\]

\subsection{Matrix Example}
A sample matrix:
\[
\begin{pmatrix}
a_{11} & a_{12} & a_{13} \\
a_{21} & a_{22} & a_{23} \\
a_{31} & a_{32} & a_{33}
\end{pmatrix}
\]

\section{Tables}
Below is a sample table showing data:

\begin{table}[h]
\centering
\begin{tabular}{@{}lrrr@{}}
\toprule
\textbf{Category} & \textbf{Value 1} & \textbf{Value 2} & \textbf{Total} \\
\midrule
Item A & 10.5 & 20.3 & 30.8 \\
Item B & 15.2 & 25.1 & 40.3 \\
Item C & 22.1 & 12.7 & 34.8 \\
\bottomrule
\end{tabular}
\caption{Sample Data Table}
\label{tab:sample}
\end{table}

\section{Lists}
\subsection{Bullet Points}
\begin{itemize}
  \item First item
  \item Second item
  \item Third item with sub-items:
    \begin{itemize}
      \item Sub-item 1
      \item Sub-item 2
    \end{itemize}
\end{itemize}

\subsection{Numbered List}
\begin{enumerate}
  \item First step
  \item Second step
  \item Third step
\end{enumerate}

\section{Citations and References}
When you use the citation search feature, references will be automatically managed. For example, you might cite a paper like this: \cite{example2024}. All citations will be automatically added to your \texttt{cite.bib} file.

\section{Conclusion}
This template demonstrates basic \LaTeX{} features including mathematical formulas, tables, and lists. The bibliography system is now set up to work seamlessly with Bibby AI's citation management. Feel free to modify it for your needs!

\bibliographystyle{plain}
\bibliography{cite}

\end{document}\documentclass{article}
\usepackage[utf8]{inputenc}
\usepackage{booktabs}
\usepackage{amsmath}
\usepackage{natbib}

\title{My Document}
\author{Your Name}
\date{\today}

\begin{document}

\maketitle

\section{Introduction}
This is a sample document demonstrating various \LaTeX{} features including mathematical equations, tables, and formatting. You can cite papers using the citation search feature, and they will be automatically added to your \texttt{cite.bib} file.

\section{Mathematical Expressions}
Here are some example equations:

\subsection{Quadratic Formula}
The quadratic formula is given by:
\[x = \frac{-b \pm \sqrt{b^2 - 4ac}}{2a}\]

\subsection{Matrix Example}
A sample matrix:
\[
\begin{pmatrix}
a_{11} & a_{12} & a_{13} \\
a_{21} & a_{22} & a_{23} \\
a_{31} & a_{32} & a_{33}
\end{pmatrix}
\]

\section{Tables}
Below is a sample table showing data:

\begin{table}[h]
\centering
\begin{tabular}{@{}lrrr@{}}
\toprule
\textbf{Category} & \textbf{Value 1} & \textbf{Value 2} & \textbf{Total} \\
\midrule
Item A & 10.5 & 20.3 & 30.8 \\
Item B & 15.2 & 25.1 & 40.3 \\
Item C & 22.1 & 12.7 & 34.8 \\
\bottomrule
\end{tabular}
\caption{Sample Data Table}
\label{tab:sample}
\end{table}

\section{Lists}
\subsection{Bullet Points}
\begin{itemize}
  \item First item
  \item Second item
  \item Third item with sub-items:
    \begin{itemize}
      \item Sub-item 1
      \item Sub-item 2
    \end{itemize}
\end{itemize}

\subsection{Numbered List}
\begin{enumerate}
  \item First step
  \item Second step
  \item Third step
\end{enumerate}

\section{Citations and References}
When you use the citation search feature, references will be automatically managed. For example, you might cite a paper like this: \cite{example2024}. All citations will be automatically added to your \texttt{cite.bib} file.

\section{Conclusion}
This template demonstrates basic \LaTeX{} features including mathematical formulas, tables, and lists. The bibliography system is now set up to work seamlessly with Bibby AI's citation management. Feel free to modify it for your needs!

\bibliographystyle{plain}
\bibliography{cite}

\end{document}\documentclass{article}
\usepackage[utf8]{inputenc}
\usepackage{booktabs}
\usepackage{amsmath}
\usepackage{natbib}

\title{My Document}
\author{Your Name}
\date{\today}

\begin{document}

\maketitle

\section{Introduction}
This is a sample document demonstrating various \LaTeX{} features including mathematical equations, tables, and formatting. You can cite papers using the citation search feature, and they will be automatically added to your \texttt{cite.bib} file.

\section{Mathematical Expressions}
Here are some example equations:

\subsection{Quadratic Formula}
The quadratic formula is given by:
\[x = \frac{-b \pm \sqrt{b^2 - 4ac}}{2a}\]

\subsection{Matrix Example}
A sample matrix:
\[
\begin{pmatrix}
a_{11} & a_{12} & a_{13} \\
a_{21} & a_{22} & a_{23} \\
a_{31} & a_{32} & a_{33}
\end{pmatrix}
\]

\section{Tables}
Below is a sample table showing data:

\begin{table}[h]
\centering
\begin{tabular}{@{}lrrr@{}}
\toprule
\textbf{Category} & \textbf{Value 1} & \textbf{Value 2} & \textbf{Total} \\
\midrule
Item A & 10.5 & 20.3 & 30.8 \\
Item B & 15.2 & 25.1 & 40.3 \\
Item C & 22.1 & 12.7 & 34.8 \\
\bottomrule
\end{tabular}
\caption{Sample Data Table}
\label{tab:sample}
\end{table}

\section{Lists}
\subsection{Bullet Points}
\begin{itemize}
  \item First item
  \item Second item
  \item Third item with sub-items:
    \begin{itemize}
      \item Sub-item 1
      \item Sub-item 2
    \end{itemize}
\end{itemize}

\subsection{Numbered List}
\begin{enumerate}
  \item First step
  \item Second step
  \item Third step
\end{enumerate}

\section{Citations and References}
When you use the citation search feature, references will be automatically managed. For example, you might cite a paper like this: \cite{example2024}. All citations will be automatically added to your \texttt{cite.bib} file.

\section{Conclusion}
This template demonstrates basic \LaTeX{} features including mathematical formulas, tables, and lists. The bibliography system is now set up to work seamlessly with Bibby AI's citation management. Feel free to modify it for your needs!

\bibliographystyle{plain}
\bibliography{cite}

\end{document}\documentclass{article}
\usepackage[utf8]{inputenc}
\usepackage{booktabs}
\usepackage{amsmath}
\usepackage{natbib}

\title{My Document}
\author{Your Name}
\date{\today}

\begin{document}

\maketitle

\section{Introduction}
This is a sample document demonstrating various \LaTeX{} features including mathematical equations, tables, and formatting. You can cite papers using the citation search feature, and they will be automatically added to your \texttt{cite.bib} file.

\section{Mathematical Expressions}
Here are some example equations:

\subsection{Quadratic Formula}
The quadratic formula is given by:
\[x = \frac{-b \pm \sqrt{b^2 - 4ac}}{2a}\]

\subsection{Matrix Example}
A sample matrix:
\[
\begin{pmatrix}
a_{11} & a_{12} & a_{13} \\
a_{21} & a_{22} & a_{23} \\
a_{31} & a_{32} & a_{33}
\end{pmatrix}
\]

\section{Tables}
Below is a sample table showing data:

\begin{table}[h]
\centering
\begin{tabular}{@{}lrrr@{}}
\toprule
\textbf{Category} & \textbf{Value 1} & \textbf{Value 2} & \textbf{Total} \\
\midrule
Item A & 10.5 & 20.3 & 30.8 \\
Item B & 15.2 & 25.1 & 40.3 \\
Item C & 22.1 & 12.7 & 34.8 \\
\bottomrule
\end{tabular}
\caption{Sample Data Table}
\label{tab:sample}
\end{table}

\section{Lists}
\subsection{Bullet Points}
\begin{itemize}
  \item First item
  \item Second item
  \item Third item with sub-items:
    \begin{itemize}
      \item Sub-item 1
      \item Sub-item 2
    \end{itemize}
\end{itemize}

\subsection{Numbered List}
\begin{enumerate}
  \item First step
  \item Second step
  \item Third step
\end{enumerate}

\section{Citations and References}
When you use the citation search feature, references will be automatically managed. For example, you might cite a paper like this: \cite{example2024}. All citations will be automatically added to your \texttt{cite.bib} file.

\section{Conclusion}
This template demonstrates basic \LaTeX{} features including mathematical formulas, tables, and lists. The bibliography system is now set up to work seamlessly with Bibby AI's citation management. Feel free to modify it for your needs!

\bibliographystyle{plain}
\bibliography{cite}

\end{document}\documentclass{article}
\usepackage[utf8]{inputenc}
\usepackage{booktabs}
\usepackage{amsmath}
\usepackage{natbib}

\title{My Document}
\author{Your Name}
\date{\today}

\begin{document}

\maketitle

\section{Introduction}
This is a sample document demonstrating various \LaTeX{} features including mathematical equations, tables, and formatting. You can cite papers using the citation search feature, and they will be automatically added to your \texttt{cite.bib} file.

\section{Mathematical Expressions}
Here are some example equations:

\subsection{Quadratic Formula}
The quadratic formula is given by:
\[x = \frac{-b \pm \sqrt{b^2 - 4ac}}{2a}\]

\subsection{Matrix Example}
A sample matrix:
\[
\begin{pmatrix}
a_{11} & a_{12} & a_{13} \\
a_{21} & a_{22} & a_{23} \\
a_{31} & a_{32} & a_{33}
\end{pmatrix}
\]

\section{Tables}
Below is a sample table showing data:

\begin{table}[h]
\centering
\begin{tabular}{@{}lrrr@{}}
\toprule
\textbf{Category} & \textbf{Value 1} & \textbf{Value 2} & \textbf{Total} \\
\midrule
Item A & 10.5 & 20.3 & 30.8 \\
Item B & 15.2 & 25.1 & 40.3 \\
Item C & 22.1 & 12.7 & 34.8 \\
\bottomrule
\end{tabular}
\caption{Sample Data Table}
\label{tab:sample}
\end{table}

\section{Lists}
\subsection{Bullet Points}
\begin{itemize}
  \item First item
  \item Second item
  \item Third item with sub-items:
    \begin{itemize}
      \item Sub-item 1
      \item Sub-item 2
    \end{itemize}
\end{itemize}

\subsection{Numbered List}
\begin{enumerate}
  \item First step
  \item Second step
  \item Third step
\end{enumerate}

\section{Citations and References}
When you use the citation search feature, references will be automatically managed. For example, you might cite a paper like this: \cite{example2024}. All citations will be automatically added to your \texttt{cite.bib} file.

\section{Conclusion}
This template demonstrates basic \LaTeX{} features including mathematical formulas, tables, and lists. The bibliography system is now set up to work seamlessly with Bibby AI's citation management. Feel free to modify it for your needs!

\bibliographystyle{plain}
\bibliography{cite}

\end{document}\documentclass{article}
\usepackage[utf8]{inputenc}
\usepackage{booktabs}
\usepackage{amsmath}
\usepackage{natbib}

\title{My Document}
\author{Your Name}
\date{\today}

\begin{document}

\maketitle

\section{Introduction}
This is a sample document demonstrating various \LaTeX{} features including mathematical equations, tables, and formatting. You can cite papers using the citation search feature, and they will be automatically added to your \texttt{cite.bib} file.

\section{Mathematical Expressions}
Here are some example equations:

\subsection{Quadratic Formula}
The quadratic formula is given by:
\[x = \frac{-b \pm \sqrt{b^2 - 4ac}}{2a}\]

\subsection{Matrix Example}
A sample matrix:
\[
\begin{pmatrix}
a_{11} & a_{12} & a_{13} \\
a_{21} & a_{22} & a_{23} \\
a_{31} & a_{32} & a_{33}
\end{pmatrix}
\]

\section{Tables}
Below is a sample table showing data:

\begin{table}[h]
\centering
\begin{tabular}{@{}lrrr@{}}
\toprule
\textbf{Category} & \textbf{Value 1} & \textbf{Value 2} & \textbf{Total} \\
\midrule
Item A & 10.5 & 20.3 & 30.8 \\
Item B & 15.2 & 25.1 & 40.3 \\
Item C & 22.1 & 12.7 & 34.8 \\
\bottomrule
\end{tabular}
\caption{Sample Data Table}
\label{tab:sample}
\end{table}

\section{Lists}
\subsection{Bullet Points}
\begin{itemize}
  \item First item
  \item Second item
  \item Third item with sub-items:
    \begin{itemize}
      \item Sub-item 1
      \item Sub-item 2
    \end{itemize}
\end{itemize}

\subsection{Numbered List}
\begin{enumerate}
  \item First step
  \item Second step
  \item Third step
\end{enumerate}

\section{Citations and References}
When you use the citation search feature, references will be automatically managed. For example, you might cite a paper like this: \cite{example2024}. All citations will be automatically added to your \texttt{cite.bib} file.

\section{Conclusion}
This template demonstrates basic \LaTeX{} features including mathematical formulas, tables, and lists. The bibliography system is now set up to work seamlessly with Bibby AI's citation management. Feel free to modify it for your needs!

\bibliographystyle{plain}
\bibliography{cite}

\end{document}\documentclass{article}
\usepackage[utf8]{inputenc}
\usepackage{booktabs}
\usepackage{amsmath}
\usepackage{natbib}

\title{My Document}
\author{Your Name}
\date{\today}

\begin{document}

\maketitle

\section{Introduction}
This is a sample document demonstrating various \LaTeX{} features including mathematical equations, tables, and formatting. You can cite papers using the citation search feature, and they will be automatically added to your \texttt{cite.bib} file.

\section{Mathematical Expressions}
Here are some example equations:

\subsection{Quadratic Formula}
The quadratic formula is given by:
\[x = \frac{-b \pm \sqrt{b^2 - 4ac}}{2a}\]

\subsection{Matrix Example}
A sample matrix:
\[
\begin{pmatrix}
a_{11} & a_{12} & a_{13} \\
a_{21} & a_{22} & a_{23} \\
a_{31} & a_{32} & a_{33}
\end{pmatrix}
\]

\section{Tables}
Below is a sample table showing data:

\begin{table}[h]
\centering
\begin{tabular}{@{}lrrr@{}}
\toprule
\textbf{Category} & \textbf{Value 1} & \textbf{Value 2} & \textbf{Total} \\
\midrule
Item A & 10.5 & 20.3 & 30.8 \\
Item B & 15.2 & 25.1 & 40.3 \\
Item C & 22.1 & 12.7 & 34.8 \\
\bottomrule
\end{tabular}
\caption{Sample Data Table}
\label{tab:sample}
\end{table}

\section{Lists}
\subsection{Bullet Points}
\begin{itemize}
  \item First item
  \item Second item
  \item Third item with sub-items:
    \begin{itemize}
      \item Sub-item 1
      \item Sub-item 2
    \end{itemize}
\end{itemize}

\subsection{Numbered List}
\begin{enumerate}
  \item First step
  \item Second step
  \item Third step
\end{enumerate}

\section{Citations and References}
When you use the citation search feature, references will be automatically managed. For example, you might cite a paper like this: \cite{example2024}. All citations will be automatically added to your \texttt{cite.bib} file.

\section{Conclusion}
This template demonstrates basic \LaTeX{} features including mathematical formulas, tables, and lists. The bibliography system is now set up to work seamlessly with Bibby AI's citation management. Feel free to modify it for your needs!

\bibliographystyle{plain}
\bibliography{cite}

\end{document}\documentclass{article}
\usepackage[utf8]{inputenc}
\usepackage{booktabs}
\usepackage{amsmath}
\usepackage{natbib}

\title{My Document}
\author{Your Name}
\date{\today}

\begin{document}

\maketitle

\section{Introduction}
This is a sample document demonstrating various \LaTeX{} features including mathematical equations, tables, and formatting. You can cite papers using the citation search feature, and they will be automatically added to your \texttt{cite.bib} file.

\section{Mathematical Expressions}
Here are some example equations:

\subsection{Quadratic Formula}
The quadratic formula is given by:
\[x = \frac{-b \pm \sqrt{b^2 - 4ac}}{2a}\]

\subsection{Matrix Example}
A sample matrix:
\[
\begin{pmatrix}
a_{11} & a_{12} & a_{13} \\
a_{21} & a_{22} & a_{23} \\
a_{31} & a_{32} & a_{33}
\end{pmatrix}
\]

\section{Tables}
Below is a sample table showing data:

\begin{table}[h]
\centering
\begin{tabular}{@{}lrrr@{}}
\toprule
\textbf{Category} & \textbf{Value 1} & \textbf{Value 2} & \textbf{Total} \\
\midrule
Item A & 10.5 & 20.3 & 30.8 \\
Item B & 15.2 & 25.1 & 40.3 \\
Item C & 22.1 & 12.7 & 34.8 \\
\bottomrule
\end{tabular}
\caption{Sample Data Table}
\label{tab:sample}
\end{table}

\section{Lists}
\subsection{Bullet Points}
\begin{itemize}
  \item First item
  \item Second item
  \item Third item with sub-items:
    \begin{itemize}
      \item Sub-item 1
      \item Sub-item 2
    \end{itemize}
\end{itemize}

\subsection{Numbered List}
\begin{enumerate}
  \item First step
  \item Second step
  \item Third step
\end{enumerate}

\section{Citations and References}
When you use the citation search feature, references will be automatically managed. For example, you might cite a paper like this: \cite{example2024}. All citations will be automatically added to your \texttt{cite.bib} file.

\section{Conclusion}
This template demonstrates basic \LaTeX{} features including mathematical formulas, tables, and lists. The bibliography system is now set up to work seamlessly with Bibby AI's citation management. Feel free to modify it for your needs!

\bibliographystyle{plain}
\bibliography{cite}

\end{document}\documentclass{article}
\usepackage[utf8]{inputenc}
\usepackage{booktabs}
\usepackage{amsmath}
\usepackage{natbib}

\title{My Document}
\author{Your Name}
\date{\today}

\begin{document}

\maketitle

\section{Introduction}
This is a sample document demonstrating various \LaTeX{} features including mathematical equations, tables, and formatting. You can cite papers using the citation search feature, and they will be automatically added to your \texttt{cite.bib} file.

\section{Mathematical Expressions}
Here are some example equations:

\subsection{Quadratic Formula}
The quadratic formula is given by:
\[x = \frac{-b \pm \sqrt{b^2 - 4ac}}{2a}\]

\subsection{Matrix Example}
A sample matrix:
\[
\begin{pmatrix}
a_{11} & a_{12} & a_{13} \\
a_{21} & a_{22} & a_{23} \\
a_{31} & a_{32} & a_{33}
\end{pmatrix}
\]

\section{Tables}
Below is a sample table showing data:

\begin{table}[h]
\centering
\begin{tabular}{@{}lrrr@{}}
\toprule
\textbf{Category} & \textbf{Value 1} & \textbf{Value 2} & \textbf{Total} \\
\midrule
Item A & 10.5 & 20.3 & 30.8 \\
Item B & 15.2 & 25.1 & 40.3 \\
Item C & 22.1 & 12.7 & 34.8 \\
\bottomrule
\end{tabular}
\caption{Sample Data Table}
\label{tab:sample}
\end{table}

\section{Lists}
\subsection{Bullet Points}
\begin{itemize}
  \item First item
  \item Second item
  \item Third item with sub-items:
    \begin{itemize}
      \item Sub-item 1
      \item Sub-item 2
    \end{itemize}
\end{itemize}

\subsection{Numbered List}
\begin{enumerate}
  \item First step
  \item Second step
  \item Third step
\end{enumerate}

\section{Citations and References}
When you use the citation search feature, references will be automatically managed. For example, you might cite a paper like this: \cite{example2024}. All citations will be automatically added to your \texttt{cite.bib} file.

\section{Conclusion}
This template demonstrates basic \LaTeX{} features including mathematical formulas, tables, and lists. The bibliography system is now set up to work seamlessly with Bibby AI's citation management. Feel free to modify it for your needs!

\bibliographystyle{plain}
\bibliography{cite}

\end{document}\documentclass{article}
\usepackage[utf8]{inputenc}
\usepackage{booktabs}
\usepackage{amsmath}
\usepackage{natbib}

\title{My Document}
\author{Your Name}
\date{\today}

\begin{document}

\maketitle

\section{Introduction}
This is a sample document demonstrating various \LaTeX{} features including mathematical equations, tables, and formatting. You can cite papers using the citation search feature, and they will be automatically added to your \texttt{cite.bib} file.

\section{Mathematical Expressions}
Here are some example equations:

\subsection{Quadratic Formula}
The quadratic formula is given by:
\[x = \frac{-b \pm \sqrt{b^2 - 4ac}}{2a}\]

\subsection{Matrix Example}
A sample matrix:
\[
\begin{pmatrix}
a_{11} & a_{12} & a_{13} \\
a_{21} & a_{22} & a_{23} \\
a_{31} & a_{32} & a_{33}
\end{pmatrix}
\]

\section{Tables}
Below is a sample table showing data:

\begin{table}[h]
\centering
\begin{tabular}{@{}lrrr@{}}
\toprule
\textbf{Category} & \textbf{Value 1} & \textbf{Value 2} & \textbf{Total} \\
\midrule
Item A & 10.5 & 20.3 & 30.8 \\
Item B & 15.2 & 25.1 & 40.3 \\
Item C & 22.1 & 12.7 & 34.8 \\
\bottomrule
\end{tabular}
\caption{Sample Data Table}
\label{tab:sample}
\end{table}

\section{Lists}
\subsection{Bullet Points}
\begin{itemize}
  \item First item
  \item Second item
  \item Third item with sub-items:
    \begin{itemize}
      \item Sub-item 1
      \item Sub-item 2
    \end{itemize}
\end{itemize}

\subsection{Numbered List}
\begin{enumerate}
  \item First step
  \item Second step
  \item Third step
\end{enumerate}

\section{Citations and References}
When you use the citation search feature, references will be automatically managed. For example, you might cite a paper like this: \cite{example2024}. All citations will be automatically added to your \texttt{cite.bib} file.

\section{Conclusion}
This template demonstrates basic \LaTeX{} features including mathematical formulas, tables, and lists. The bibliography system is now set up to work seamlessly with Bibby AI's citation management. Feel free to modify it for your needs!

\bibliographystyle{plain}
\bibliography{cite}

\end{document}\documentclass{article}
\usepackage[utf8]{inputenc}
\usepackage{booktabs}
\usepackage{amsmath}
\usepackage{natbib}

\title{My Document}
\author{Your Name}
\date{\today}

\begin{document}

\maketitle

\section{Introduction}
This is a sample document demonstrating various \LaTeX{} features including mathematical equations, tables, and formatting. You can cite papers using the citation search feature, and they will be automatically added to your \texttt{cite.bib} file.

\section{Mathematical Expressions}
Here are some example equations:

\subsection{Quadratic Formula}
The quadratic formula is given by:
\[x = \frac{-b \pm \sqrt{b^2 - 4ac}}{2a}\]

\subsection{Matrix Example}
A sample matrix:
\[
\begin{pmatrix}
a_{11} & a_{12} & a_{13} \\
a_{21} & a_{22} & a_{23} \\
a_{31} & a_{32} & a_{33}
\end{pmatrix}
\]

\section{Tables}
Below is a sample table showing data:

\begin{table}[h]
\centering
\begin{tabular}{@{}lrrr@{}}
\toprule
\textbf{Category} & \textbf{Value 1} & \textbf{Value 2} & \textbf{Total} \\
\midrule
Item A & 10.5 & 20.3 & 30.8 \\
Item B & 15.2 & 25.1 & 40.3 \\
Item C & 22.1 & 12.7 & 34.8 \\
\bottomrule
\end{tabular}
\caption{Sample Data Table}
\label{tab:sample}
\end{table}

\section{Lists}
\subsection{Bullet Points}
\begin{itemize}
  \item First item
  \item Second item
  \item Third item with sub-items:
    \begin{itemize}
      \item Sub-item 1
      \item Sub-item 2
    \end{itemize}
\end{itemize}

\subsection{Numbered List}
\begin{enumerate}
  \item First step
  \item Second step
  \item Third step
\end{enumerate}

\section{Citations and References}
When you use the citation search feature, references will be automatically managed. For example, you might cite a paper like this: \cite{example2024}. All citations will be automatically added to your \texttt{cite.bib} file.

\section{Conclusion}
This template demonstrates basic \LaTeX{} features including mathematical formulas, tables, and lists. The bibliography system is now set up to work seamlessly with Bibby AI's citation management. Feel free to modify it for your needs!

\bibliographystyle{plain}
\bibliography{cite}

\end{document}\documentclass{article}
\usepackage[utf8]{inputenc}
\usepackage{booktabs}
\usepackage{amsmath}
\usepackage{natbib}

\title{My Document}
\author{Your Name}
\date{\today}

\begin{document}

\maketitle

\section{Introduction}
This is a sample document demonstrating various \LaTeX{} features including mathematical equations, tables, and formatting. You can cite papers using the citation search feature, and they will be automatically added to your \texttt{cite.bib} file.

\section{Mathematical Expressions}
Here are some example equations:

\subsection{Quadratic Formula}
The quadratic formula is given by:
\[x = \frac{-b \pm \sqrt{b^2 - 4ac}}{2a}\]

\subsection{Matrix Example}
A sample matrix:
\[
\begin{pmatrix}
a_{11} & a_{12} & a_{13} \\
a_{21} & a_{22} & a_{23} \\
a_{31} & a_{32} & a_{33}
\end{pmatrix}
\]

\section{Tables}
Below is a sample table showing data:

\begin{table}[h]
\centering
\begin{tabular}{@{}lrrr@{}}
\toprule
\textbf{Category} & \textbf{Value 1} & \textbf{Value 2} & \textbf{Total} \\
\midrule
Item A & 10.5 & 20.3 & 30.8 \\
Item B & 15.2 & 25.1 & 40.3 \\
Item C & 22.1 & 12.7 & 34.8 \\
\bottomrule
\end{tabular}
\caption{Sample Data Table}
\label{tab:sample}
\end{table}

\section{Lists}
\subsection{Bullet Points}
\begin{itemize}
  \item First item
  \item Second item
  \item Third item with sub-items:
    \begin{itemize}
      \item Sub-item 1
      \item Sub-item 2
    \end{itemize}
\end{itemize}

\subsection{Numbered List}
\begin{enumerate}
  \item First step
  \item Second step
  \item Third step
\end{enumerate}

\section{Citations and References}
When you use the citation search feature, references will be automatically managed. For example, you might cite a paper like this: \cite{example2024}. All citations will be automatically added to your \texttt{cite.bib} file.

\section{Conclusion}
This template demonstrates basic \LaTeX{} features including mathematical formulas, tables, and lists. The bibliography system is now set up to work seamlessly with Bibby AI's citation management. Feel free to modify it for your needs!

\bibliographystyle{plain}
\bibliography{cite}

\end{document}\documentclass{article}
\usepackage[utf8]{inputenc}
\usepackage{booktabs}
\usepackage{amsmath}
\usepackage{natbib}

\title{My Document}
\author{Your Name}
\date{\today}

\begin{document}

\maketitle

\section{Introduction}
This is a sample document demonstrating various \LaTeX{} features including mathematical equations, tables, and formatting. You can cite papers using the citation search feature, and they will be automatically added to your \texttt{cite.bib} file.

\section{Mathematical Expressions}
Here are some example equations:

\subsection{Quadratic Formula}
The quadratic formula is given by:
\[x = \frac{-b \pm \sqrt{b^2 - 4ac}}{2a}\]

\subsection{Matrix Example}
A sample matrix:
\[
\begin{pmatrix}
a_{11} & a_{12} & a_{13} \\
a_{21} & a_{22} & a_{23} \\
a_{31} & a_{32} & a_{33}
\end{pmatrix}
\]

\section{Tables}
Below is a sample table showing data:

\begin{table}[h]
\centering
\begin{tabular}{@{}lrrr@{}}
\toprule
\textbf{Category} & \textbf{Value 1} & \textbf{Value 2} & \textbf{Total} \\
\midrule
Item A & 10.5 & 20.3 & 30.8 \\
Item B & 15.2 & 25.1 & 40.3 \\
Item C & 22.1 & 12.7 & 34.8 \\
\bottomrule
\end{tabular}
\caption{Sample Data Table}
\label{tab:sample}
\end{table}

\section{Lists}
\subsection{Bullet Points}
\begin{itemize}
  \item First item
  \item Second item
  \item Third item with sub-items:
    \begin{itemize}
      \item Sub-item 1
      \item Sub-item 2
    \end{itemize}
\end{itemize}

\subsection{Numbered List}
\begin{enumerate}
  \item First step
  \item Second step
  \item Third step
\end{enumerate}

\section{Citations and References}
When you use the citation search feature, references will be automatically managed. For example, you might cite a paper like this: \cite{example2024}. All citations will be automatically added to your \texttt{cite.bib} file.

\section{Conclusion}
This template demonstrates basic \LaTeX{} features including mathematical formulas, tables, and lists. The bibliography system is now set up to work seamlessly with Bibby AI's citation management. Feel free to modify it for your needs!

\bibliographystyle{plain}
\bibliography{cite}

\end{document}\documentclass{article}
\usepackage[utf8]{inputenc}
\usepackage{booktabs}
\usepackage{amsmath}
\usepackage{natbib}

\title{My Document}
\author{Your Name}
\date{\today}

\begin{document}

\maketitle

\section{Introduction}
This is a sample document demonstrating various \LaTeX{} features including mathematical equations, tables, and formatting. You can cite papers using the citation search feature, and they will be automatically added to your \texttt{cite.bib} file.

\section{Mathematical Expressions}
Here are some example equations:

\subsection{Quadratic Formula}
The quadratic formula is given by:
\[x = \frac{-b \pm \sqrt{b^2 - 4ac}}{2a}\]

\subsection{Matrix Example}
A sample matrix:
\[
\begin{pmatrix}
a_{11} & a_{12} & a_{13} \\
a_{21} & a_{22} & a_{23} \\
a_{31} & a_{32} & a_{33}
\end{pmatrix}
\]

\section{Tables}
Below is a sample table showing data:

\begin{table}[h]
\centering
\begin{tabular}{@{}lrrr@{}}
\toprule
\textbf{Category} & \textbf{Value 1} & \textbf{Value 2} & \textbf{Total} \\
\midrule
Item A & 10.5 & 20.3 & 30.8 \\
Item B & 15.2 & 25.1 & 40.3 \\
Item C & 22.1 & 12.7 & 34.8 \\
\bottomrule
\end{tabular}
\caption{Sample Data Table}
\label{tab:sample}
\end{table}

\section{Lists}
\subsection{Bullet Points}
\begin{itemize}
  \item First item
  \item Second item
  \item Third item with sub-items:
    \begin{itemize}
      \item Sub-item 1
      \item Sub-item 2
    \end{itemize}
\end{itemize}

\subsection{Numbered List}
\begin{enumerate}
  \item First step
  \item Second step
  \item Third step
\end{enumerate}

\section{Citations and References}
When you use the citation search feature, references will be automatically managed. For example, you might cite a paper like this: \cite{example2024}. All citations will be automatically added to your \texttt{cite.bib} file.

\section{Conclusion}
This template demonstrates basic \LaTeX{} features including mathematical formulas, tables, and lists. The bibliography system is now set up to work seamlessly with Bibby AI's citation management. Feel free to modify it for your needs!

\bibliographystyle{plain}
\bibliography{cite}

\end{document}\documentclass{article}
\usepackage[utf8]{inputenc}
\usepackage{booktabs}
\usepackage{amsmath}
\usepackage{natbib}

\title{My Document}
\author{Your Name}
\date{\today}

\begin{document}

\maketitle

\section{Introduction}
This is a sample document demonstrating various \LaTeX{} features including mathematical equations, tables, and formatting. You can cite papers using the citation search feature, and they will be automatically added to your \texttt{cite.bib} file.

\section{Mathematical Expressions}
Here are some example equations:

\subsection{Quadratic Formula}
The quadratic formula is given by:
\[x = \frac{-b \pm \sqrt{b^2 - 4ac}}{2a}\]

\subsection{Matrix Example}
A sample matrix:
\[
\begin{pmatrix}
a_{11} & a_{12} & a_{13} \\
a_{21} & a_{22} & a_{23} \\
a_{31} & a_{32} & a_{33}
\end{pmatrix}
\]

\section{Tables}
Below is a sample table showing data:

\begin{table}[h]
\centering
\begin{tabular}{@{}lrrr@{}}
\toprule
\textbf{Category} & \textbf{Value 1} & \textbf{Value 2} & \textbf{Total} \\
\midrule
Item A & 10.5 & 20.3 & 30.8 \\
Item B & 15.2 & 25.1 & 40.3 \\
Item C & 22.1 & 12.7 & 34.8 \\
\bottomrule
\end{tabular}
\caption{Sample Data Table}
\label{tab:sample}
\end{table}

\section{Lists}
\subsection{Bullet Points}
\begin{itemize}
  \item First item
  \item Second item
  \item Third item with sub-items:
    \begin{itemize}
      \item Sub-item 1
      \item Sub-item 2
    \end{itemize}
\end{itemize}

\subsection{Numbered List}
\begin{enumerate}
  \item First step
  \item Second step
  \item Third step
\end{enumerate}

\section{Citations and References}
When you use the citation search feature, references will be automatically managed. For example, you might cite a paper like this: \cite{example2024}. All citations will be automatically added to your \texttt{cite.bib} file.

\section{Conclusion}
This template demonstrates basic \LaTeX{} features including mathematical formulas, tables, and lists. The bibliography system is now set up to work seamlessly with Bibby AI's citation management. Feel free to modify it for your needs!

\bibliographystyle{plain}
\bibliography{cite}

\end{document}\documentclass{article}
\usepackage[utf8]{inputenc}
\usepackage{booktabs}
\usepackage{amsmath}
\usepackage{natbib}

\title{My Document}
\author{Your Name}
\date{\today}

\begin{document}

\maketitle

\section{Introduction}
This is a sample document demonstrating various \LaTeX{} features including mathematical equations, tables, and formatting. You can cite papers using the citation search feature, and they will be automatically added to your \texttt{cite.bib} file.

\section{Mathematical Expressions}
Here are some example equations:

\subsection{Quadratic Formula}
The quadratic formula is given by:
\[x = \frac{-b \pm \sqrt{b^2 - 4ac}}{2a}\]

\subsection{Matrix Example}
A sample matrix:
\[
\begin{pmatrix}
a_{11} & a_{12} & a_{13} \\
a_{21} & a_{22} & a_{23} \\
a_{31} & a_{32} & a_{33}
\end{pmatrix}
\]

\section{Tables}
Below is a sample table showing data:

\begin{table}[h]
\centering
\begin{tabular}{@{}lrrr@{}}
\toprule
\textbf{Category} & \textbf{Value 1} & \textbf{Value 2} & \textbf{Total} \\
\midrule
Item A & 10.5 & 20.3 & 30.8 \\
Item B & 15.2 & 25.1 & 40.3 \\
Item C & 22.1 & 12.7 & 34.8 \\
\bottomrule
\end{tabular}
\caption{Sample Data Table}
\label{tab:sample}
\end{table}

\section{Lists}
\subsection{Bullet Points}
\begin{itemize}
  \item First item
  \item Second item
  \item Third item with sub-items:
    \begin{itemize}
      \item Sub-item 1
      \item Sub-item 2
    \end{itemize}
\end{itemize}

\subsection{Numbered List}
\begin{enumerate}
  \item First step
  \item Second step
  \item Third step
\end{enumerate}

\section{Citations and References}
When you use the citation search feature, references will be automatically managed. For example, you might cite a paper like this: \cite{example2024}. All citations will be automatically added to your \texttt{cite.bib} file.

\section{Conclusion}
This template demonstrates basic \LaTeX{} features including mathematical formulas, tables, and lists. The bibliography system is now set up to work seamlessly with Bibby AI's citation management. Feel free to modify it for your needs!

\bibliographystyle{plain}
\bibliography{cite}

\end{document}\documentclass{article}
\usepackage[utf8]{inputenc}
\usepackage{booktabs}
\usepackage{amsmath}
\usepackage{natbib}

\title{My Document}
\author{Your Name}
\date{\today}

\begin{document}

\maketitle

\section{Introduction}
This is a sample document demonstrating various \LaTeX{} features including mathematical equations, tables, and formatting. You can cite papers using the citation search feature, and they will be automatically added to your \texttt{cite.bib} file.

\section{Mathematical Expressions}
Here are some example equations:

\subsection{Quadratic Formula}
The quadratic formula is given by:
\[x = \frac{-b \pm \sqrt{b^2 - 4ac}}{2a}\]

\subsection{Matrix Example}
A sample matrix:
\[
\begin{pmatrix}
a_{11} & a_{12} & a_{13} \\
a_{21} & a_{22} & a_{23} \\
a_{31} & a_{32} & a_{33}
\end{pmatrix}
\]

\section{Tables}
Below is a sample table showing data:

\begin{table}[h]
\centering
\begin{tabular}{@{}lrrr@{}}
\toprule
\textbf{Category} & \textbf{Value 1} & \textbf{Value 2} & \textbf{Total} \\
\midrule
Item A & 10.5 & 20.3 & 30.8 \\
Item B & 15.2 & 25.1 & 40.3 \\
Item C & 22.1 & 12.7 & 34.8 \\
\bottomrule
\end{tabular}
\caption{Sample Data Table}
\label{tab:sample}
\end{table}

\section{Lists}
\subsection{Bullet Points}
\begin{itemize}
  \item First item
  \item Second item
  \item Third item with sub-items:
    \begin{itemize}
      \item Sub-item 1
      \item Sub-item 2
    \end{itemize}
\end{itemize}

\subsection{Numbered List}
\begin{enumerate}
  \item First step
  \item Second step
  \item Third step
\end{enumerate}

\section{Citations and References}
When you use the citation search feature, references will be automatically managed. For example, you might cite a paper like this: \cite{example2024}. All citations will be automatically added to your \texttt{cite.bib} file.

\section{Conclusion}
This template demonstrates basic \LaTeX{} features including mathematical formulas, tables, and lists. The bibliography system is now set up to work seamlessly with Bibby AI's citation management. Feel free to modify it for your needs!

\bibliographystyle{plain}
\bibliography{cite}

\end{document}\documentclass{article}
\usepackage[utf8]{inputenc}
\usepackage{booktabs}
\usepackage{amsmath}
\usepackage{natbib}

\title{My Document}
\author{Your Name}
\date{\today}

\begin{document}

\maketitle

\section{Introduction}
This is a sample document demonstrating various \LaTeX{} features including mathematical equations, tables, and formatting. You can cite papers using the citation search feature, and they will be automatically added to your \texttt{cite.bib} file.

\section{Mathematical Expressions}
Here are some example equations:

\subsection{Quadratic Formula}
The quadratic formula is given by:
\[x = \frac{-b \pm \sqrt{b^2 - 4ac}}{2a}\]

\subsection{Matrix Example}
A sample matrix:
\[
\begin{pmatrix}
a_{11} & a_{12} & a_{13} \\
a_{21} & a_{22} & a_{23} \\
a_{31} & a_{32} & a_{33}
\end{pmatrix}
\]

\section{Tables}
Below is a sample table showing data:

\begin{table}[h]
\centering
\begin{tabular}{@{}lrrr@{}}
\toprule
\textbf{Category} & \textbf{Value 1} & \textbf{Value 2} & \textbf{Total} \\
\midrule
Item A & 10.5 & 20.3 & 30.8 \\
Item B & 15.2 & 25.1 & 40.3 \\
Item C & 22.1 & 12.7 & 34.8 \\
\bottomrule
\end{tabular}
\caption{Sample Data Table}
\label{tab:sample}
\end{table}

\section{Lists}
\subsection{Bullet Points}
\begin{itemize}
  \item First item
  \item Second item
  \item Third item with sub-items:
    \begin{itemize}
      \item Sub-item 1
      \item Sub-item 2
    \end{itemize}
\end{itemize}

\subsection{Numbered List}
\begin{enumerate}
  \item First step
  \item Second step
  \item Third step
\end{enumerate}

\section{Citations and References}
When you use the citation search feature, references will be automatically managed. For example, you might cite a paper like this: \cite{example2024}. All citations will be automatically added to your \texttt{cite.bib} file.

\section{Conclusion}
This template demonstrates basic \LaTeX{} features including mathematical formulas, tables, and lists. The bibliography system is now set up to work seamlessly with Bibby AI's citation management. Feel free to modify it for your needs!

\bibliographystyle{plain}
\bibliography{cite}

\end{document}\documentclass{article}
\usepackage[utf8]{inputenc}
\usepackage{booktabs}
\usepackage{amsmath}
\usepackage{natbib}

\title{My Document}
\author{Your Name}
\date{\today}

\begin{document}

\maketitle

\section{Introduction}
This is a sample document demonstrating various \LaTeX{} features including mathematical equations, tables, and formatting. You can cite papers using the citation search feature, and they will be automatically added to your \texttt{cite.bib} file.

\section{Mathematical Expressions}
Here are some example equations:

\subsection{Quadratic Formula}
The quadratic formula is given by:
\[x = \frac{-b \pm \sqrt{b^2 - 4ac}}{2a}\]

\subsection{Matrix Example}
A sample matrix:
\[
\begin{pmatrix}
a_{11} & a_{12} & a_{13} \\
a_{21} & a_{22} & a_{23} \\
a_{31} & a_{32} & a_{33}
\end{pmatrix}
\]

\section{Tables}
Below is a sample table showing data:

\begin{table}[h]
\centering
\begin{tabular}{@{}lrrr@{}}
\toprule
\textbf{Category} & \textbf{Value 1} & \textbf{Value 2} & \textbf{Total} \\
\midrule
Item A & 10.5 & 20.3 & 30.8 \\
Item B & 15.2 & 25.1 & 40.3 \\
Item C & 22.1 & 12.7 & 34.8 \\
\bottomrule
\end{tabular}
\caption{Sample Data Table}
\label{tab:sample}
\end{table}

\section{Lists}
\subsection{Bullet Points}
\begin{itemize}
  \item First item
  \item Second item
  \item Third item with sub-items:
    \begin{itemize}
      \item Sub-item 1
      \item Sub-item 2
    \end{itemize}
\end{itemize}

\subsection{Numbered List}
\begin{enumerate}
  \item First step
  \item Second step
  \item Third step
\end{enumerate}

\section{Citations and References}
When you use the citation search feature, references will be automatically managed. For example, you might cite a paper like this: \cite{example2024}. All citations will be automatically added to your \texttt{cite.bib} file.

\section{Conclusion}
This template demonstrates basic \LaTeX{} features including mathematical formulas, tables, and lists. The bibliography system is now set up to work seamlessly with Bibby AI's citation management. Feel free to modify it for your needs!

\bibliographystyle{plain}
\bibliography{cite}

\end{document}\documentclass{article}
\usepackage[utf8]{inputenc}
\usepackage{booktabs}
\usepackage{amsmath}
\usepackage{natbib}

\title{My Document}
\author{Your Name}
\date{\today}

\begin{document}

\maketitle

\section{Introduction}
This is a sample document demonstrating various \LaTeX{} features including mathematical equations, tables, and formatting. You can cite papers using the citation search feature, and they will be automatically added to your \texttt{cite.bib} file.

\section{Mathematical Expressions}
Here are some example equations:

\subsection{Quadratic Formula}
The quadratic formula is given by:
\[x = \frac{-b \pm \sqrt{b^2 - 4ac}}{2a}\]

\subsection{Matrix Example}
A sample matrix:
\[
\begin{pmatrix}
a_{11} & a_{12} & a_{13} \\
a_{21} & a_{22} & a_{23} \\
a_{31} & a_{32} & a_{33}
\end{pmatrix}
\]

\section{Tables}
Below is a sample table showing data:

\begin{table}[h]
\centering
\begin{tabular}{@{}lrrr@{}}
\toprule
\textbf{Category} & \textbf{Value 1} & \textbf{Value 2} & \textbf{Total} \\
\midrule
Item A & 10.5 & 20.3 & 30.8 \\
Item B & 15.2 & 25.1 & 40.3 \\
Item C & 22.1 & 12.7 & 34.8 \\
\bottomrule
\end{tabular}
\caption{Sample Data Table}
\label{tab:sample}
\end{table}

\section{Lists}
\subsection{Bullet Points}
\begin{itemize}
  \item First item
  \item Second item
  \item Third item with sub-items:
    \begin{itemize}
      \item Sub-item 1
      \item Sub-item 2
    \end{itemize}
\end{itemize}

\subsection{Numbered List}
\begin{enumerate}
  \item First step
  \item Second step
  \item Third step
\end{enumerate}

\section{Citations and References}
When you use the citation search feature, references will be automatically managed. For example, you might cite a paper like this: \cite{example2024}. All citations will be automatically added to your \texttt{cite.bib} file.

\section{Conclusion}
This template demonstrates basic \LaTeX{} features including mathematical formulas, tables, and lists. The bibliography system is now set up to work seamlessly with Bibby AI's citation management. Feel free to modify it for your needs!

\bibliographystyle{plain}
\bibliography{cite}

\end{document}\documentclass{article}
\usepackage[utf8]{inputenc}
\usepackage{booktabs}
\usepackage{amsmath}
\usepackage{natbib}

\title{My Document}
\author{Your Name}
\date{\today}

\begin{document}

\maketitle

\section{Introduction}
This is a sample document demonstrating various \LaTeX{} features including mathematical equations, tables, and formatting. You can cite papers using the citation search feature, and they will be automatically added to your \texttt{cite.bib} file.

\section{Mathematical Expressions}
Here are some example equations:

\subsection{Quadratic Formula}
The quadratic formula is given by:
\[x = \frac{-b \pm \sqrt{b^2 - 4ac}}{2a}\]

\subsection{Matrix Example}
A sample matrix:
\[
\begin{pmatrix}
a_{11} & a_{12} & a_{13} \\
a_{21} & a_{22} & a_{23} \\
a_{31} & a_{32} & a_{33}
\end{pmatrix}
\]

\section{Tables}
Below is a sample table showing data:

\begin{table}[h]
\centering
\begin{tabular}{@{}lrrr@{}}
\toprule
\textbf{Category} & \textbf{Value 1} & \textbf{Value 2} & \textbf{Total} \\
\midrule
Item A & 10.5 & 20.3 & 30.8 \\
Item B & 15.2 & 25.1 & 40.3 \\
Item C & 22.1 & 12.7 & 34.8 \\
\bottomrule
\end{tabular}
\caption{Sample Data Table}
\label{tab:sample}
\end{table}

\section{Lists}
\subsection{Bullet Points}
\begin{itemize}
  \item First item
  \item Second item
  \item Third item with sub-items:
    \begin{itemize}
      \item Sub-item 1
      \item Sub-item 2
    \end{itemize}
\end{itemize}

\subsection{Numbered List}
\begin{enumerate}
  \item First step
  \item Second step
  \item Third step
\end{enumerate}

\section{Citations and References}
When you use the citation search feature, references will be automatically managed. For example, you might cite a paper like this: \cite{example2024}. All citations will be automatically added to your \texttt{cite.bib} file.

\section{Conclusion}
This template demonstrates basic \LaTeX{} features including mathematical formulas, tables, and lists. The bibliography system is now set up to work seamlessly with Bibby AI's citation management. Feel free to modify it for your needs!

\bibliographystyle{plain}
\bibliography{cite}

\end{document}\documentclass{article}
\usepackage[utf8]{inputenc}
\usepackage{booktabs}
\usepackage{amsmath}
\usepackage{natbib}

\title{My Document}
\author{Your Name}
\date{\today}

\begin{document}

\maketitle

\section{Introduction}
This is a sample document demonstrating various \LaTeX{} features including mathematical equations, tables, and formatting. You can cite papers using the citation search feature, and they will be automatically added to your \texttt{cite.bib} file.

\section{Mathematical Expressions}
Here are some example equations:

\subsection{Quadratic Formula}
The quadratic formula is given by:
\[x = \frac{-b \pm \sqrt{b^2 - 4ac}}{2a}\]

\subsection{Matrix Example}
A sample matrix:
\[
\begin{pmatrix}
a_{11} & a_{12} & a_{13} \\
a_{21} & a_{22} & a_{23} \\
a_{31} & a_{32} & a_{33}
\end{pmatrix}
\]

\section{Tables}
Below is a sample table showing data:

\begin{table}[h]
\centering
\begin{tabular}{@{}lrrr@{}}
\toprule
\textbf{Category} & \textbf{Value 1} & \textbf{Value 2} & \textbf{Total} \\
\midrule
Item A & 10.5 & 20.3 & 30.8 \\
Item B & 15.2 & 25.1 & 40.3 \\
Item C & 22.1 & 12.7 & 34.8 \\
\bottomrule
\end{tabular}
\caption{Sample Data Table}
\label{tab:sample}
\end{table}

\section{Lists}
\subsection{Bullet Points}
\begin{itemize}
  \item First item
  \item Second item
  \item Third item with sub-items:
    \begin{itemize}
      \item Sub-item 1
      \item Sub-item 2
    \end{itemize}
\end{itemize}

\subsection{Numbered List}
\begin{enumerate}
  \item First step
  \item Second step
  \item Third step
\end{enumerate}

\section{Citations and References}
When you use the citation search feature, references will be automatically managed. For example, you might cite a paper like this: \cite{example2024}. All citations will be automatically added to your \texttt{cite.bib} file.

\section{Conclusion}
This template demonstrates basic \LaTeX{} features including mathematical formulas, tables, and lists. The bibliography system is now set up to work seamlessly with Bibby AI's citation management. Feel free to modify it for your needs!

\bibliographystyle{plain}
\bibliography{cite}

\end{document}\documentclass{article}
\usepackage[utf8]{inputenc}
\usepackage{booktabs}
\usepackage{amsmath}
\usepackage{natbib}

\title{My Document}
\author{Your Name}
\date{\today}

\begin{document}

\maketitle

\section{Introduction}
This is a sample document demonstrating various \LaTeX{} features including mathematical equations, tables, and formatting. You can cite papers using the citation search feature, and they will be automatically added to your \texttt{cite.bib} file.

\section{Mathematical Expressions}
Here are some example equations:

\subsection{Quadratic Formula}
The quadratic formula is given by:
\[x = \frac{-b \pm \sqrt{b^2 - 4ac}}{2a}\]

\subsection{Matrix Example}
A sample matrix:
\[
\begin{pmatrix}
a_{11} & a_{12} & a_{13} \\
a_{21} & a_{22} & a_{23} \\
a_{31} & a_{32} & a_{33}
\end{pmatrix}
\]

\section{Tables}
Below is a sample table showing data:

\begin{table}[h]
\centering
\begin{tabular}{@{}lrrr@{}}
\toprule
\textbf{Category} & \textbf{Value 1} & \textbf{Value 2} & \textbf{Total} \\
\midrule
Item A & 10.5 & 20.3 & 30.8 \\
Item B & 15.2 & 25.1 & 40.3 \\
Item C & 22.1 & 12.7 & 34.8 \\
\bottomrule
\end{tabular}
\caption{Sample Data Table}
\label{tab:sample}
\end{table}

\section{Lists}
\subsection{Bullet Points}
\begin{itemize}
  \item First item
  \item Second item
  \item Third item with sub-items:
    \begin{itemize}
      \item Sub-item 1
      \item Sub-item 2
    \end{itemize}
\end{itemize}

\subsection{Numbered List}
\begin{enumerate}
  \item First step
  \item Second step
  \item Third step
\end{enumerate}

\section{Citations and References}
When you use the citation search feature, references will be automatically managed. For example, you might cite a paper like this: \cite{example2024}. All citations will be automatically added to your \texttt{cite.bib} file.

\section{Conclusion}
This template demonstrates basic \LaTeX{} features including mathematical formulas, tables, and lists. The bibliography system is now set up to work seamlessly with Bibby AI's citation management. Feel free to modify it for your needs!

\bibliographystyle{plain}
\bibliography{cite}

\end{document}\documentclass{article}
\usepackage[utf8]{inputenc}
\usepackage{booktabs}
\usepackage{amsmath}
\usepackage{natbib}

\title{My Document}
\author{Your Name}
\date{\today}

\begin{document}

\maketitle

\section{Introduction}
This is a sample document demonstrating various \LaTeX{} features including mathematical equations, tables, and formatting. You can cite papers using the citation search feature, and they will be automatically added to your \texttt{cite.bib} file.

\section{Mathematical Expressions}
Here are some example equations:

\subsection{Quadratic Formula}
The quadratic formula is given by:
\[x = \frac{-b \pm \sqrt{b^2 - 4ac}}{2a}\]

\subsection{Matrix Example}
A sample matrix:
\[
\begin{pmatrix}
a_{11} & a_{12} & a_{13} \\
a_{21} & a_{22} & a_{23} \\
a_{31} & a_{32} & a_{33}
\end{pmatrix}
\]

\section{Tables}
Below is a sample table showing data:

\begin{table}[h]
\centering
\begin{tabular}{@{}lrrr@{}}
\toprule
\textbf{Category} & \textbf{Value 1} & \textbf{Value 2} & \textbf{Total} \\
\midrule
Item A & 10.5 & 20.3 & 30.8 \\
Item B & 15.2 & 25.1 & 40.3 \\
Item C & 22.1 & 12.7 & 34.8 \\
\bottomrule
\end{tabular}
\caption{Sample Data Table}
\label{tab:sample}
\end{table}

\section{Lists}
\subsection{Bullet Points}
\begin{itemize}
  \item First item
  \item Second item
  \item Third item with sub-items:
    \begin{itemize}
      \item Sub-item 1
      \item Sub-item 2
    \end{itemize}
\end{itemize}

\subsection{Numbered List}
\begin{enumerate}
  \item First step
  \item Second step
  \item Third step
\end{enumerate}

\section{Citations and References}
When you use the citation search feature, references will be automatically managed. For example, you might cite a paper like this: \cite{example2024}. All citations will be automatically added to your \texttt{cite.bib} file.

\section{Conclusion}
This template demonstrates basic \LaTeX{} features including mathematical formulas, tables, and lists. The bibliography system is now set up to work seamlessly with Bibby AI's citation management. Feel free to modify it for your needs!

\bibliographystyle{plain}
\bibliography{cite}

\end{document}\documentclass{article}
\usepackage[utf8]{inputenc}
\usepackage{booktabs}
\usepackage{amsmath}
\usepackage{natbib}

\title{My Document}
\author{Your Name}
\date{\today}

\begin{document}

\maketitle

\section{Introduction}
This is a sample document demonstrating various \LaTeX{} features including mathematical equations, tables, and formatting. You can cite papers using the citation search feature, and they will be automatically added to your \texttt{cite.bib} file.

\section{Mathematical Expressions}
Here are some example equations:

\subsection{Quadratic Formula}
The quadratic formula is given by:
\[x = \frac{-b \pm \sqrt{b^2 - 4ac}}{2a}\]

\subsection{Matrix Example}
A sample matrix:
\[
\begin{pmatrix}
a_{11} & a_{12} & a_{13} \\
a_{21} & a_{22} & a_{23} \\
a_{31} & a_{32} & a_{33}
\end{pmatrix}
\]

\section{Tables}
Below is a sample table showing data:

\begin{table}[h]
\centering
\begin{tabular}{@{}lrrr@{}}
\toprule
\textbf{Category} & \textbf{Value 1} & \textbf{Value 2} & \textbf{Total} \\
\midrule
Item A & 10.5 & 20.3 & 30.8 \\
Item B & 15.2 & 25.1 & 40.3 \\
Item C & 22.1 & 12.7 & 34.8 \\
\bottomrule
\end{tabular}
\caption{Sample Data Table}
\label{tab:sample}
\end{table}

\section{Lists}
\subsection{Bullet Points}
\begin{itemize}
  \item First item
  \item Second item
  \item Third item with sub-items:
    \begin{itemize}
      \item Sub-item 1
      \item Sub-item 2
    \end{itemize}
\end{itemize}

\subsection{Numbered List}
\begin{enumerate}
  \item First step
  \item Second step
  \item Third step
\end{enumerate}

\section{Citations and References}
When you use the citation search feature, references will be automatically managed. For example, you might cite a paper like this: \cite{example2024}. All citations will be automatically added to your \texttt{cite.bib} file.

\section{Conclusion}
This template demonstrates basic \LaTeX{} features including mathematical formulas, tables, and lists. The bibliography system is now set up to work seamlessly with Bibby AI's citation management. Feel free to modify it for your needs!

\bibliographystyle{plain}
\bibliography{cite}

\end{document}\documentclass{article}
\usepackage[utf8]{inputenc}
\usepackage{booktabs}
\usepackage{amsmath}
\usepackage{natbib}

\title{My Document}
\author{Your Name}
\date{\today}

\begin{document}

\maketitle

\section{Introduction}
This is a sample document demonstrating various \LaTeX{} features including mathematical equations, tables, and formatting. You can cite papers using the citation search feature, and they will be automatically added to your \texttt{cite.bib} file.

\section{Mathematical Expressions}
Here are some example equations:

\subsection{Quadratic Formula}
The quadratic formula is given by:
\[x = \frac{-b \pm \sqrt{b^2 - 4ac}}{2a}\]

\subsection{Matrix Example}
A sample matrix:
\[
\begin{pmatrix}
a_{11} & a_{12} & a_{13} \\
a_{21} & a_{22} & a_{23} \\
a_{31} & a_{32} & a_{33}
\end{pmatrix}
\]

\section{Tables}
Below is a sample table showing data:

\begin{table}[h]
\centering
\begin{tabular}{@{}lrrr@{}}
\toprule
\textbf{Category} & \textbf{Value 1} & \textbf{Value 2} & \textbf{Total} \\
\midrule
Item A & 10.5 & 20.3 & 30.8 \\
Item B & 15.2 & 25.1 & 40.3 \\
Item C & 22.1 & 12.7 & 34.8 \\
\bottomrule
\end{tabular}
\caption{Sample Data Table}
\label{tab:sample}
\end{table}

\section{Lists}
\subsection{Bullet Points}
\begin{itemize}
  \item First item
  \item Second item
  \item Third item with sub-items:
    \begin{itemize}
      \item Sub-item 1
      \item Sub-item 2
    \end{itemize}
\end{itemize}

\subsection{Numbered List}
\begin{enumerate}
  \item First step
  \item Second step
  \item Third step
\end{enumerate}

\section{Citations and References}
When you use the citation search feature, references will be automatically managed. For example, you might cite a paper like this: \cite{example2024}. All citations will be automatically added to your \texttt{cite.bib} file.

\section{Conclusion}
This template demonstrates basic \LaTeX{} features including mathematical formulas, tables, and lists. The bibliography system is now set up to work seamlessly with Bibby AI's citation management. Feel free to modify it for your needs!

\bibliographystyle{plain}
\bibliography{cite}

\end{document}\documentclass{article}
\usepackage[utf8]{inputenc}
\usepackage{booktabs}
\usepackage{amsmath}
\usepackage{natbib}

\title{My Document}
\author{Your Name}
\date{\today}

\begin{document}

\maketitle

\section{Introduction}
This is a sample document demonstrating various \LaTeX{} features including mathematical equations, tables, and formatting. You can cite papers using the citation search feature, and they will be automatically added to your \texttt{cite.bib} file.

\section{Mathematical Expressions}
Here are some example equations:

\subsection{Quadratic Formula}
The quadratic formula is given by:
\[x = \frac{-b \pm \sqrt{b^2 - 4ac}}{2a}\]

\subsection{Matrix Example}
A sample matrix:
\[
\begin{pmatrix}
a_{11} & a_{12} & a_{13} \\
a_{21} & a_{22} & a_{23} \\
a_{31} & a_{32} & a_{33}
\end{pmatrix}
\]

\section{Tables}
Below is a sample table showing data:

\begin{table}[h]
\centering
\begin{tabular}{@{}lrrr@{}}
\toprule
\textbf{Category} & \textbf{Value 1} & \textbf{Value 2} & \textbf{Total} \\
\midrule
Item A & 10.5 & 20.3 & 30.8 \\
Item B & 15.2 & 25.1 & 40.3 \\
Item C & 22.1 & 12.7 & 34.8 \\
\bottomrule
\end{tabular}
\caption{Sample Data Table}
\label{tab:sample}
\end{table}

\section{Lists}
\subsection{Bullet Points}
\begin{itemize}
  \item First item
  \item Second item
  \item Third item with sub-items:
    \begin{itemize}
      \item Sub-item 1
      \item Sub-item 2
    \end{itemize}
\end{itemize}

\subsection{Numbered List}
\begin{enumerate}
  \item First step
  \item Second step
  \item Third step
\end{enumerate}

\section{Citations and References}
When you use the citation search feature, references will be automatically managed. For example, you might cite a paper like this: \cite{example2024}. All citations will be automatically added to your \texttt{cite.bib} file.

\section{Conclusion}
This template demonstrates basic \LaTeX{} features including mathematical formulas, tables, and lists. The bibliography system is now set up to work seamlessly with Bibby AI's citation management. Feel free to modify it for your needs!

\bibliographystyle{plain}
\bibliography{cite}

\end{document}\documentclass{article}
\usepackage[utf8]{inputenc}
\usepackage{booktabs}
\usepackage{amsmath}
\usepackage{natbib}

\title{My Document}
\author{Your Name}
\date{\today}

\begin{document}

\maketitle

\section{Introduction}
This is a sample document demonstrating various \LaTeX{} features including mathematical equations, tables, and formatting. You can cite papers using the citation search feature, and they will be automatically added to your \texttt{cite.bib} file.

\section{Mathematical Expressions}
Here are some example equations:

\subsection{Quadratic Formula}
The quadratic formula is given by:
\[x = \frac{-b \pm \sqrt{b^2 - 4ac}}{2a}\]

\subsection{Matrix Example}
A sample matrix:
\[
\begin{pmatrix}
a_{11} & a_{12} & a_{13} \\
a_{21} & a_{22} & a_{23} \\
a_{31} & a_{32} & a_{33}
\end{pmatrix}
\]

\section{Tables}
Below is a sample table showing data:

\begin{table}[h]
\centering
\begin{tabular}{@{}lrrr@{}}
\toprule
\textbf{Category} & \textbf{Value 1} & \textbf{Value 2} & \textbf{Total} \\
\midrule
Item A & 10.5 & 20.3 & 30.8 \\
Item B & 15.2 & 25.1 & 40.3 \\
Item C & 22.1 & 12.7 & 34.8 \\
\bottomrule
\end{tabular}
\caption{Sample Data Table}
\label{tab:sample}
\end{table}

\section{Lists}
\subsection{Bullet Points}
\begin{itemize}
  \item First item
  \item Second item
  \item Third item with sub-items:
    \begin{itemize}
      \item Sub-item 1
      \item Sub-item 2
    \end{itemize}
\end{itemize}

\subsection{Numbered List}
\begin{enumerate}
  \item First step
  \item Second step
  \item Third step
\end{enumerate}

\section{Citations and References}
When you use the citation search feature, references will be automatically managed. For example, you might cite a paper like this: \cite{example2024}. All citations will be automatically added to your \texttt{cite.bib} file.

\section{Conclusion}
This template demonstrates basic \LaTeX{} features including mathematical formulas, tables, and lists. The bibliography system is now set up to work seamlessly with Bibby AI's citation management. Feel free to modify it for your needs!

\bibliographystyle{plain}
\bibliography{cite}

\end{document}\documentclass{article}
\usepackage[utf8]{inputenc}
\usepackage{booktabs}
\usepackage{amsmath}
\usepackage{natbib}

\title{My Document}
\author{Your Name}
\date{\today}

\begin{document}

\maketitle

\section{Introduction}
This is a sample document demonstrating various \LaTeX{} features including mathematical equations, tables, and formatting. You can cite papers using the citation search feature, and they will be automatically added to your \texttt{cite.bib} file.

\section{Mathematical Expressions}
Here are some example equations:

\subsection{Quadratic Formula}
The quadratic formula is given by:
\[x = \frac{-b \pm \sqrt{b^2 - 4ac}}{2a}\]

\subsection{Matrix Example}
A sample matrix:
\[
\begin{pmatrix}
a_{11} & a_{12} & a_{13} \\
a_{21} & a_{22} & a_{23} \\
a_{31} & a_{32} & a_{33}
\end{pmatrix}
\]

\section{Tables}
Below is a sample table showing data:

\begin{table}[h]
\centering
\begin{tabular}{@{}lrrr@{}}
\toprule
\textbf{Category} & \textbf{Value 1} & \textbf{Value 2} & \textbf{Total} \\
\midrule
Item A & 10.5 & 20.3 & 30.8 \\
Item B & 15.2 & 25.1 & 40.3 \\
Item C & 22.1 & 12.7 & 34.8 \\
\bottomrule
\end{tabular}
\caption{Sample Data Table}
\label{tab:sample}
\end{table}

\section{Lists}
\subsection{Bullet Points}
\begin{itemize}
  \item First item
  \item Second item
  \item Third item with sub-items:
    \begin{itemize}
      \item Sub-item 1
      \item Sub-item 2
    \end{itemize}
\end{itemize}

\subsection{Numbered List}
\begin{enumerate}
  \item First step
  \item Second step
  \item Third step
\end{enumerate}

\section{Citations and References}
When you use the citation search feature, references will be automatically managed. For example, you might cite a paper like this: \cite{example2024}. All citations will be automatically added to your \texttt{cite.bib} file.

\section{Conclusion}
This template demonstrates basic \LaTeX{} features including mathematical formulas, tables, and lists. The bibliography system is now set up to work seamlessly with Bibby AI's citation management. Feel free to modify it for your needs!

\bibliographystyle{plain}
\bibliography{cite}

\end{document}\documentclass{article}
\usepackage[utf8]{inputenc}
\usepackage{booktabs}
\usepackage{amsmath}
\usepackage{natbib}

\title{My Document}
\author{Your Name}
\date{\today}

\begin{document}

\maketitle

\section{Introduction}
This is a sample document demonstrating various \LaTeX{} features including mathematical equations, tables, and formatting. You can cite papers using the citation search feature, and they will be automatically added to your \texttt{cite.bib} file.

\section{Mathematical Expressions}
Here are some example equations:

\subsection{Quadratic Formula}
The quadratic formula is given by:
\[x = \frac{-b \pm \sqrt{b^2 - 4ac}}{2a}\]

\subsection{Matrix Example}
A sample matrix:
\[
\begin{pmatrix}
a_{11} & a_{12} & a_{13} \\
a_{21} & a_{22} & a_{23} \\
a_{31} & a_{32} & a_{33}
\end{pmatrix}
\]

\section{Tables}
Below is a sample table showing data:

\begin{table}[h]
\centering
\begin{tabular}{@{}lrrr@{}}
\toprule
\textbf{Category} & \textbf{Value 1} & \textbf{Value 2} & \textbf{Total} \\
\midrule
Item A & 10.5 & 20.3 & 30.8 \\
Item B & 15.2 & 25.1 & 40.3 \\
Item C & 22.1 & 12.7 & 34.8 \\
\bottomrule
\end{tabular}
\caption{Sample Data Table}
\label{tab:sample}
\end{table}

\section{Lists}
\subsection{Bullet Points}
\begin{itemize}
  \item First item
  \item Second item
  \item Third item with sub-items:
    \begin{itemize}
      \item Sub-item 1
      \item Sub-item 2
    \end{itemize}
\end{itemize}

\subsection{Numbered List}
\begin{enumerate}
  \item First step
  \item Second step
  \item Third step
\end{enumerate}

\section{Citations and References}
When you use the citation search feature, references will be automatically managed. For example, you might cite a paper like this: \cite{example2024}. All citations will be automatically added to your \texttt{cite.bib} file.

\section{Conclusion}
This template demonstrates basic \LaTeX{} features including mathematical formulas, tables, and lists. The bibliography system is now set up to work seamlessly with Bibby AI's citation management. Feel free to modify it for your needs!

\bibliographystyle{plain}
\bibliography{cite}

\end{document}\documentclass{article}
\usepackage[utf8]{inputenc}
\usepackage{booktabs}
\usepackage{amsmath}
\usepackage{natbib}

\title{My Document}
\author{Your Name}
\date{\today}

\begin{document}

\maketitle

\section{Introduction}
This is a sample document demonstrating various \LaTeX{} features including mathematical equations, tables, and formatting. You can cite papers using the citation search feature, and they will be automatically added to your \texttt{cite.bib} file.

\section{Mathematical Expressions}
Here are some example equations:

\subsection{Quadratic Formula}
The quadratic formula is given by:
\[x = \frac{-b \pm \sqrt{b^2 - 4ac}}{2a}\]

\subsection{Matrix Example}
A sample matrix:
\[
\begin{pmatrix}
a_{11} & a_{12} & a_{13} \\
a_{21} & a_{22} & a_{23} \\
a_{31} & a_{32} & a_{33}
\end{pmatrix}
\]

\section{Tables}
Below is a sample table showing data:

\begin{table}[h]
\centering
\begin{tabular}{@{}lrrr@{}}
\toprule
\textbf{Category} & \textbf{Value 1} & \textbf{Value 2} & \textbf{Total} \\
\midrule
Item A & 10.5 & 20.3 & 30.8 \\
Item B & 15.2 & 25.1 & 40.3 \\
Item C & 22.1 & 12.7 & 34.8 \\
\bottomrule
\end{tabular}
\caption{Sample Data Table}
\label{tab:sample}
\end{table}

\section{Lists}
\subsection{Bullet Points}
\begin{itemize}
  \item First item
  \item Second item
  \item Third item with sub-items:
    \begin{itemize}
      \item Sub-item 1
      \item Sub-item 2
    \end{itemize}
\end{itemize}

\subsection{Numbered List}
\begin{enumerate}
  \item First step
  \item Second step
  \item Third step
\end{enumerate}

\section{Citations and References}
When you use the citation search feature, references will be automatically managed. For example, you might cite a paper like this: \cite{example2024}. All citations will be automatically added to your \texttt{cite.bib} file.

\section{Conclusion}
This template demonstrates basic \LaTeX{} features including mathematical formulas, tables, and lists. The bibliography system is now set up to work seamlessly with Bibby AI's citation management. Feel free to modify it for your needs!

\bibliographystyle{plain}
\bibliography{cite}

\end{document}\documentclass{article}
\usepackage[utf8]{inputenc}
\usepackage{booktabs}
\usepackage{amsmath}
\usepackage{natbib}

\title{My Document}
\author{Your Name}
\date{\today}

\begin{document}

\maketitle

\section{Introduction}
This is a sample document demonstrating various \LaTeX{} features including mathematical equations, tables, and formatting. You can cite papers using the citation search feature, and they will be automatically added to your \texttt{cite.bib} file.

\section{Mathematical Expressions}
Here are some example equations:

\subsection{Quadratic Formula}
The quadratic formula is given by:
\[x = \frac{-b \pm \sqrt{b^2 - 4ac}}{2a}\]

\subsection{Matrix Example}
A sample matrix:
\[
\begin{pmatrix}
a_{11} & a_{12} & a_{13} \\
a_{21} & a_{22} & a_{23} \\
a_{31} & a_{32} & a_{33}
\end{pmatrix}
\]

\section{Tables}
Below is a sample table showing data:

\begin{table}[h]
\centering
\begin{tabular}{@{}lrrr@{}}
\toprule
\textbf{Category} & \textbf{Value 1} & \textbf{Value 2} & \textbf{Total} \\
\midrule
Item A & 10.5 & 20.3 & 30.8 \\
Item B & 15.2 & 25.1 & 40.3 \\
Item C & 22.1 & 12.7 & 34.8 \\
\bottomrule
\end{tabular}
\caption{Sample Data Table}
\label{tab:sample}
\end{table}

\section{Lists}
\subsection{Bullet Points}
\begin{itemize}
  \item First item
  \item Second item
  \item Third item with sub-items:
    \begin{itemize}
      \item Sub-item 1
      \item Sub-item 2
    \end{itemize}
\end{itemize}

\subsection{Numbered List}
\begin{enumerate}
  \item First step
  \item Second step
  \item Third step
\end{enumerate}

\section{Citations and References}
When you use the citation search feature, references will be automatically managed. For example, you might cite a paper like this: \cite{example2024}. All citations will be automatically added to your \texttt{cite.bib} file.

\section{Conclusion}
This template demonstrates basic \LaTeX{} features including mathematical formulas, tables, and lists. The bibliography system is now set up to work seamlessly with Bibby AI's citation management. Feel free to modify it for your needs!

\bibliographystyle{plain}
\bibliography{cite}

\end{document}\documentclass{article}
\usepackage[utf8]{inputenc}
\usepackage{booktabs}
\usepackage{amsmath}
\usepackage{natbib}

\title{My Document}
\author{Your Name}
\date{\today}

\begin{document}

\maketitle

\section{Introduction}
This is a sample document demonstrating various \LaTeX{} features including mathematical equations, tables, and formatting. You can cite papers using the citation search feature, and they will be automatically added to your \texttt{cite.bib} file.

\section{Mathematical Expressions}
Here are some example equations:

\subsection{Quadratic Formula}
The quadratic formula is given by:
\[x = \frac{-b \pm \sqrt{b^2 - 4ac}}{2a}\]

\subsection{Matrix Example}
A sample matrix:
\[
\begin{pmatrix}
a_{11} & a_{12} & a_{13} \\
a_{21} & a_{22} & a_{23} \\
a_{31} & a_{32} & a_{33}
\end{pmatrix}
\]

\section{Tables}
Below is a sample table showing data:

\begin{table}[h]
\centering
\begin{tabular}{@{}lrrr@{}}
\toprule
\textbf{Category} & \textbf{Value 1} & \textbf{Value 2} & \textbf{Total} \\
\midrule
Item A & 10.5 & 20.3 & 30.8 \\
Item B & 15.2 & 25.1 & 40.3 \\
Item C & 22.1 & 12.7 & 34.8 \\
\bottomrule
\end{tabular}
\caption{Sample Data Table}
\label{tab:sample}
\end{table}

\section{Lists}
\subsection{Bullet Points}
\begin{itemize}
  \item First item
  \item Second item
  \item Third item with sub-items:
    \begin{itemize}
      \item Sub-item 1
      \item Sub-item 2
    \end{itemize}
\end{itemize}

\subsection{Numbered List}
\begin{enumerate}
  \item First step
  \item Second step
  \item Third step
\end{enumerate}

\section{Citations and References}
When you use the citation search feature, references will be automatically managed. For example, you might cite a paper like this: \cite{example2024}. All citations will be automatically added to your \texttt{cite.bib} file.

\section{Conclusion}
This template demonstrates basic \LaTeX{} features including mathematical formulas, tables, and lists. The bibliography system is now set up to work seamlessly with Bibby AI's citation management. Feel free to modify it for your needs!

\bibliographystyle{plain}
\bibliography{cite}

\end{document}\documentclass{article}
\usepackage[utf8]{inputenc}
\usepackage{booktabs}
\usepackage{amsmath}
\usepackage{natbib}

\title{My Document}
\author{Your Name}
\date{\today}

\begin{document}

\maketitle

\section{Introduction}
This is a sample document demonstrating various \LaTeX{} features including mathematical equations, tables, and formatting. You can cite papers using the citation search feature, and they will be automatically added to your \texttt{cite.bib} file.

\section{Mathematical Expressions}
Here are some example equations:

\subsection{Quadratic Formula}
The quadratic formula is given by:
\[x = \frac{-b \pm \sqrt{b^2 - 4ac}}{2a}\]

\subsection{Matrix Example}
A sample matrix:
\[
\begin{pmatrix}
a_{11} & a_{12} & a_{13} \\
a_{21} & a_{22} & a_{23} \\
a_{31} & a_{32} & a_{33}
\end{pmatrix}
\]

\section{Tables}
Below is a sample table showing data:

\begin{table}[h]
\centering
\begin{tabular}{@{}lrrr@{}}
\toprule
\textbf{Category} & \textbf{Value 1} & \textbf{Value 2} & \textbf{Total} \\
\midrule
Item A & 10.5 & 20.3 & 30.8 \\
Item B & 15.2 & 25.1 & 40.3 \\
Item C & 22.1 & 12.7 & 34.8 \\
\bottomrule
\end{tabular}
\caption{Sample Data Table}
\label{tab:sample}
\end{table}

\section{Lists}
\subsection{Bullet Points}
\begin{itemize}
  \item First item
  \item Second item
  \item Third item with sub-items:
    \begin{itemize}
      \item Sub-item 1
      \item Sub-item 2
    \end{itemize}
\end{itemize}

\subsection{Numbered List}
\begin{enumerate}
  \item First step
  \item Second step
  \item Third step
\end{enumerate}

\section{Citations and References}
When you use the citation search feature, references will be automatically managed. For example, you might cite a paper like this: \cite{example2024}. All citations will be automatically added to your \texttt{cite.bib} file.

\section{Conclusion}
This template demonstrates basic \LaTeX{} features including mathematical formulas, tables, and lists. The bibliography system is now set up to work seamlessly with Bibby AI's citation management. Feel free to modify it for your needs!

\bibliographystyle{plain}
\bibliography{cite}

\end{document}\documentclass{article}
\usepackage[utf8]{inputenc}
\usepackage{booktabs}
\usepackage{amsmath}
\usepackage{natbib}

\title{My Document}
\author{Your Name}
\date{\today}

\begin{document}

\maketitle

\section{Introduction}
This is a sample document demonstrating various \LaTeX{} features including mathematical equations, tables, and formatting. You can cite papers using the citation search feature, and they will be automatically added to your \texttt{cite.bib} file.

\section{Mathematical Expressions}
Here are some example equations:

\subsection{Quadratic Formula}
The quadratic formula is given by:
\[x = \frac{-b \pm \sqrt{b^2 - 4ac}}{2a}\]

\subsection{Matrix Example}
A sample matrix:
\[
\begin{pmatrix}
a_{11} & a_{12} & a_{13} \\
a_{21} & a_{22} & a_{23} \\
a_{31} & a_{32} & a_{33}
\end{pmatrix}
\]

\section{Tables}
Below is a sample table showing data:

\begin{table}[h]
\centering
\begin{tabular}{@{}lrrr@{}}
\toprule
\textbf{Category} & \textbf{Value 1} & \textbf{Value 2} & \textbf{Total} \\
\midrule
Item A & 10.5 & 20.3 & 30.8 \\
Item B & 15.2 & 25.1 & 40.3 \\
Item C & 22.1 & 12.7 & 34.8 \\
\bottomrule
\end{tabular}
\caption{Sample Data Table}
\label{tab:sample}
\end{table}

\section{Lists}
\subsection{Bullet Points}
\begin{itemize}
  \item First item
  \item Second item
  \item Third item with sub-items:
    \begin{itemize}
      \item Sub-item 1
      \item Sub-item 2
    \end{itemize}
\end{itemize}

\subsection{Numbered List}
\begin{enumerate}
  \item First step
  \item Second step
  \item Third step
\end{enumerate}

\section{Citations and References}
When you use the citation search feature, references will be automatically managed. For example, you might cite a paper like this: \cite{example2024}. All citations will be automatically added to your \texttt{cite.bib} file.

\section{Conclusion}
This template demonstrates basic \LaTeX{} features including mathematical formulas, tables, and lists. The bibliography system is now set up to work seamlessly with Bibby AI's citation management. Feel free to modify it for your needs!

\bibliographystyle{plain}
\bibliography{cite}

\end{document}\documentclass{article}
\usepackage[utf8]{inputenc}
\usepackage{booktabs}
\usepackage{amsmath}
\usepackage{natbib}

\title{My Document}
\author{Your Name}
\date{\today}

\begin{document}

\maketitle

\section{Introduction}
This is a sample document demonstrating various \LaTeX{} features including mathematical equations, tables, and formatting. You can cite papers using the citation search feature, and they will be automatically added to your \texttt{cite.bib} file.

\section{Mathematical Expressions}
Here are some example equations:

\subsection{Quadratic Formula}
The quadratic formula is given by:
\[x = \frac{-b \pm \sqrt{b^2 - 4ac}}{2a}\]

\subsection{Matrix Example}
A sample matrix:
\[
\begin{pmatrix}
a_{11} & a_{12} & a_{13} \\
a_{21} & a_{22} & a_{23} \\
a_{31} & a_{32} & a_{33}
\end{pmatrix}
\]

\section{Tables}
Below is a sample table showing data:

\begin{table}[h]
\centering
\begin{tabular}{@{}lrrr@{}}
\toprule
\textbf{Category} & \textbf{Value 1} & \textbf{Value 2} & \textbf{Total} \\
\midrule
Item A & 10.5 & 20.3 & 30.8 \\
Item B & 15.2 & 25.1 & 40.3 \\
Item C & 22.1 & 12.7 & 34.8 \\
\bottomrule
\end{tabular}
\caption{Sample Data Table}
\label{tab:sample}
\end{table}

\section{Lists}
\subsection{Bullet Points}
\begin{itemize}
  \item First item
  \item Second item
  \item Third item with sub-items:
    \begin{itemize}
      \item Sub-item 1
      \item Sub-item 2
    \end{itemize}
\end{itemize}

\subsection{Numbered List}
\begin{enumerate}
  \item First step
  \item Second step
  \item Third step
\end{enumerate}

\section{Citations and References}
When you use the citation search feature, references will be automatically managed. For example, you might cite a paper like this: \cite{example2024}. All citations will be automatically added to your \texttt{cite.bib} file.

\section{Conclusion}
This template demonstrates basic \LaTeX{} features including mathematical formulas, tables, and lists. The bibliography system is now set up to work seamlessly with Bibby AI's citation management. Feel free to modify it for your needs!

\bibliographystyle{plain}
\bibliography{cite}

\end{document}\documentclass{article}
\usepackage[utf8]{inputenc}
\usepackage{booktabs}
\usepackage{amsmath}
\usepackage{natbib}

\title{My Document}
\author{Your Name}
\date{\today}

\begin{document}

\maketitle

\section{Introduction}
This is a sample document demonstrating various \LaTeX{} features including mathematical equations, tables, and formatting. You can cite papers using the citation search feature, and they will be automatically added to your \texttt{cite.bib} file.

\section{Mathematical Expressions}
Here are some example equations:

\subsection{Quadratic Formula}
The quadratic formula is given by:
\[x = \frac{-b \pm \sqrt{b^2 - 4ac}}{2a}\]

\subsection{Matrix Example}
A sample matrix:
\[
\begin{pmatrix}
a_{11} & a_{12} & a_{13} \\
a_{21} & a_{22} & a_{23} \\
a_{31} & a_{32} & a_{33}
\end{pmatrix}
\]

\section{Tables}
Below is a sample table showing data:

\begin{table}[h]
\centering
\begin{tabular}{@{}lrrr@{}}
\toprule
\textbf{Category} & \textbf{Value 1} & \textbf{Value 2} & \textbf{Total} \\
\midrule
Item A & 10.5 & 20.3 & 30.8 \\
Item B & 15.2 & 25.1 & 40.3 \\
Item C & 22.1 & 12.7 & 34.8 \\
\bottomrule
\end{tabular}
\caption{Sample Data Table}
\label{tab:sample}
\end{table}

\section{Lists}
\subsection{Bullet Points}
\begin{itemize}
  \item First item
  \item Second item
  \item Third item with sub-items:
    \begin{itemize}
      \item Sub-item 1
      \item Sub-item 2
    \end{itemize}
\end{itemize}

\subsection{Numbered List}
\begin{enumerate}
  \item First step
  \item Second step
  \item Third step
\end{enumerate}

\section{Citations and References}
When you use the citation search feature, references will be automatically managed. For example, you might cite a paper like this: \cite{example2024}. All citations will be automatically added to your \texttt{cite.bib} file.

\section{Conclusion}
This template demonstrates basic \LaTeX{} features including mathematical formulas, tables, and lists. The bibliography system is now set up to work seamlessly with Bibby AI's citation management. Feel free to modify it for your needs!

\bibliographystyle{plain}
\bibliography{cite}

\end{document}\documentclass{article}
\usepackage[utf8]{inputenc}
\usepackage{booktabs}
\usepackage{amsmath}
\usepackage{natbib}

\title{My Document}
\author{Your Name}
\date{\today}

\begin{document}

\maketitle

\section{Introduction}
This is a sample document demonstrating various \LaTeX{} features including mathematical equations, tables, and formatting. You can cite papers using the citation search feature, and they will be automatically added to your \texttt{cite.bib} file.

\section{Mathematical Expressions}
Here are some example equations:

\subsection{Quadratic Formula}
The quadratic formula is given by:
\[x = \frac{-b \pm \sqrt{b^2 - 4ac}}{2a}\]

\subsection{Matrix Example}
A sample matrix:
\[
\begin{pmatrix}
a_{11} & a_{12} & a_{13} \\
a_{21} & a_{22} & a_{23} \\
a_{31} & a_{32} & a_{33}
\end{pmatrix}
\]

\section{Tables}
Below is a sample table showing data:

\begin{table}[h]
\centering
\begin{tabular}{@{}lrrr@{}}
\toprule
\textbf{Category} & \textbf{Value 1} & \textbf{Value 2} & \textbf{Total} \\
\midrule
Item A & 10.5 & 20.3 & 30.8 \\
Item B & 15.2 & 25.1 & 40.3 \\
Item C & 22.1 & 12.7 & 34.8 \\
\bottomrule
\end{tabular}
\caption{Sample Data Table}
\label{tab:sample}
\end{table}

\section{Lists}
\subsection{Bullet Points}
\begin{itemize}
  \item First item
  \item Second item
  \item Third item with sub-items:
    \begin{itemize}
      \item Sub-item 1
      \item Sub-item 2
    \end{itemize}
\end{itemize}

\subsection{Numbered List}
\begin{enumerate}
  \item First step
  \item Second step
  \item Third step
\end{enumerate}

\section{Citations and References}
When you use the citation search feature, references will be automatically managed. For example, you might cite a paper like this: \cite{example2024}. All citations will be automatically added to your \texttt{cite.bib} file.

\section{Conclusion}
This template demonstrates basic \LaTeX{} features including mathematical formulas, tables, and lists. The bibliography system is now set up to work seamlessly with Bibby AI's citation management. Feel free to modify it for your needs!

\bibliographystyle{plain}
\bibliography{cite}

\end{document}\documentclass{article}
\usepackage[utf8]{inputenc}
\usepackage{booktabs}
\usepackage{amsmath}
\usepackage{natbib}

\title{My Document}
\author{Your Name}
\date{\today}

\begin{document}

\maketitle

\section{Introduction}
This is a sample document demonstrating various \LaTeX{} features including mathematical equations, tables, and formatting. You can cite papers using the citation search feature, and they will be automatically added to your \texttt{cite.bib} file.

\section{Mathematical Expressions}
Here are some example equations:

\subsection{Quadratic Formula}
The quadratic formula is given by:
\[x = \frac{-b \pm \sqrt{b^2 - 4ac}}{2a}\]

\subsection{Matrix Example}
A sample matrix:
\[
\begin{pmatrix}
a_{11} & a_{12} & a_{13} \\
a_{21} & a_{22} & a_{23} \\
a_{31} & a_{32} & a_{33}
\end{pmatrix}
\]

\section{Tables}
Below is a sample table showing data:

\begin{table}[h]
\centering
\begin{tabular}{@{}lrrr@{}}
\toprule
\textbf{Category} & \textbf{Value 1} & \textbf{Value 2} & \textbf{Total} \\
\midrule
Item A & 10.5 & 20.3 & 30.8 \\
Item B & 15.2 & 25.1 & 40.3 \\
Item C & 22.1 & 12.7 & 34.8 \\
\bottomrule
\end{tabular}
\caption{Sample Data Table}
\label{tab:sample}
\end{table}

\section{Lists}
\subsection{Bullet Points}
\begin{itemize}
  \item First item
  \item Second item
  \item Third item with sub-items:
    \begin{itemize}
      \item Sub-item 1
      \item Sub-item 2
    \end{itemize}
\end{itemize}

\subsection{Numbered List}
\begin{enumerate}
  \item First step
  \item Second step
  \item Third step
\end{enumerate}

\section{Citations and References}
When you use the citation search feature, references will be automatically managed. For example, you might cite a paper like this: \cite{example2024}. All citations will be automatically added to your \texttt{cite.bib} file.

\section{Conclusion}
This template demonstrates basic \LaTeX{} features including mathematical formulas, tables, and lists. The bibliography system is now set up to work seamlessly with Bibby AI's citation management. Feel free to modify it for your needs!

\bibliographystyle{plain}
\bibliography{cite}

\end{document}\documentclass{article}
\usepackage[utf8]{inputenc}
\usepackage{booktabs}
\usepackage{amsmath}
\usepackage{natbib}

\title{My Document}
\author{Your Name}
\date{\today}

\begin{document}

\maketitle

\section{Introduction}
This is a sample document demonstrating various \LaTeX{} features including mathematical equations, tables, and formatting. You can cite papers using the citation search feature, and they will be automatically added to your \texttt{cite.bib} file.

\section{Mathematical Expressions}
Here are some example equations:

\subsection{Quadratic Formula}
The quadratic formula is given by:
\[x = \frac{-b \pm \sqrt{b^2 - 4ac}}{2a}\]

\subsection{Matrix Example}
A sample matrix:
\[
\begin{pmatrix}
a_{11} & a_{12} & a_{13} \\
a_{21} & a_{22} & a_{23} \\
a_{31} & a_{32} & a_{33}
\end{pmatrix}
\]

\section{Tables}
Below is a sample table showing data:

\begin{table}[h]
\centering
\begin{tabular}{@{}lrrr@{}}
\toprule
\textbf{Category} & \textbf{Value 1} & \textbf{Value 2} & \textbf{Total} \\
\midrule
Item A & 10.5 & 20.3 & 30.8 \\
Item B & 15.2 & 25.1 & 40.3 \\
Item C & 22.1 & 12.7 & 34.8 \\
\bottomrule
\end{tabular}
\caption{Sample Data Table}
\label{tab:sample}
\end{table}

\section{Lists}
\subsection{Bullet Points}
\begin{itemize}
  \item First item
  \item Second item
  \item Third item with sub-items:
    \begin{itemize}
      \item Sub-item 1
      \item Sub-item 2
    \end{itemize}
\end{itemize}

\subsection{Numbered List}
\begin{enumerate}
  \item First step
  \item Second step
  \item Third step
\end{enumerate}

\section{Citations and References}
When you use the citation search feature, references will be automatically managed. For example, you might cite a paper like this: \cite{example2024}. All citations will be automatically added to your \texttt{cite.bib} file.

\section{Conclusion}
This template demonstrates basic \LaTeX{} features including mathematical formulas, tables, and lists. The bibliography system is now set up to work seamlessly with Bibby AI's citation management. Feel free to modify it for your needs!

\bibliographystyle{plain}
\bibliography{cite}

\end{document}\documentclass{article}
\usepackage[utf8]{inputenc}
\usepackage{booktabs}
\usepackage{amsmath}
\usepackage{natbib}

\title{My Document}
\author{Your Name}
\date{\today}

\begin{document}

\maketitle

\section{Introduction}
This is a sample document demonstrating various \LaTeX{} features including mathematical equations, tables, and formatting. You can cite papers using the citation search feature, and they will be automatically added to your \texttt{cite.bib} file.

\section{Mathematical Expressions}
Here are some example equations:

\subsection{Quadratic Formula}
The quadratic formula is given by:
\[x = \frac{-b \pm \sqrt{b^2 - 4ac}}{2a}\]

\subsection{Matrix Example}
A sample matrix:
\[
\begin{pmatrix}
a_{11} & a_{12} & a_{13} \\
a_{21} & a_{22} & a_{23} \\
a_{31} & a_{32} & a_{33}
\end{pmatrix}
\]

\section{Tables}
Below is a sample table showing data:

\begin{table}[h]
\centering
\begin{tabular}{@{}lrrr@{}}
\toprule
\textbf{Category} & \textbf{Value 1} & \textbf{Value 2} & \textbf{Total} \\
\midrule
Item A & 10.5 & 20.3 & 30.8 \\
Item B & 15.2 & 25.1 & 40.3 \\
Item C & 22.1 & 12.7 & 34.8 \\
\bottomrule
\end{tabular}
\caption{Sample Data Table}
\label{tab:sample}
\end{table}

\section{Lists}
\subsection{Bullet Points}
\begin{itemize}
  \item First item
  \item Second item
  \item Third item with sub-items:
    \begin{itemize}
      \item Sub-item 1
      \item Sub-item 2
    \end{itemize}
\end{itemize}

\subsection{Numbered List}
\begin{enumerate}
  \item First step
  \item Second step
  \item Third step
\end{enumerate}

\section{Citations and References}
When you use the citation search feature, references will be automatically managed. For example, you might cite a paper like this: \cite{example2024}. All citations will be automatically added to your \texttt{cite.bib} file.

\section{Conclusion}
This template demonstrates basic \LaTeX{} features including mathematical formulas, tables, and lists. The bibliography system is now set up to work seamlessly with Bibby AI's citation management. Feel free to modify it for your needs!

\bibliographystyle{plain}
\bibliography{cite}

\end{document}\documentclass{article}
\usepackage[utf8]{inputenc}
\usepackage{booktabs}
\usepackage{amsmath}
\usepackage{natbib}

\title{My Document}
\author{Your Name}
\date{\today}

\begin{document}

\maketitle

\section{Introduction}
This is a sample document demonstrating various \LaTeX{} features including mathematical equations, tables, and formatting. You can cite papers using the citation search feature, and they will be automatically added to your \texttt{cite.bib} file.

\section{Mathematical Expressions}
Here are some example equations:

\subsection{Quadratic Formula}
The quadratic formula is given by:
\[x = \frac{-b \pm \sqrt{b^2 - 4ac}}{2a}\]

\subsection{Matrix Example}
A sample matrix:
\[
\begin{pmatrix}
a_{11} & a_{12} & a_{13} \\
a_{21} & a_{22} & a_{23} \\
a_{31} & a_{32} & a_{33}
\end{pmatrix}
\]

\section{Tables}
Below is a sample table showing data:

\begin{table}[h]
\centering
\begin{tabular}{@{}lrrr@{}}
\toprule
\textbf{Category} & \textbf{Value 1} & \textbf{Value 2} & \textbf{Total} \\
\midrule
Item A & 10.5 & 20.3 & 30.8 \\
Item B & 15.2 & 25.1 & 40.3 \\
Item C & 22.1 & 12.7 & 34.8 \\
\bottomrule
\end{tabular}
\caption{Sample Data Table}
\label{tab:sample}
\end{table}

\section{Lists}
\subsection{Bullet Points}
\begin{itemize}
  \item First item
  \item Second item
  \item Third item with sub-items:
    \begin{itemize}
      \item Sub-item 1
      \item Sub-item 2
    \end{itemize}
\end{itemize}

\subsection{Numbered List}
\begin{enumerate}
  \item First step
  \item Second step
  \item Third step
\end{enumerate}

\section{Citations and References}
When you use the citation search feature, references will be automatically managed. For example, you might cite a paper like this: \cite{example2024}. All citations will be automatically added to your \texttt{cite.bib} file.

\section{Conclusion}
This template demonstrates basic \LaTeX{} features including mathematical formulas, tables, and lists. The bibliography system is now set up to work seamlessly with Bibby AI's citation management. Feel free to modify it for your needs!

\bibliographystyle{plain}
\bibliography{cite}

\end{document}\documentclass{article}
\usepackage[utf8]{inputenc}
\usepackage{booktabs}
\usepackage{amsmath}
\usepackage{natbib}

\title{My Document}
\author{Your Name}
\date{\today}

\begin{document}

\maketitle

\section{Introduction}
This is a sample document demonstrating various \LaTeX{} features including mathematical equations, tables, and formatting. You can cite papers using the citation search feature, and they will be automatically added to your \texttt{cite.bib} file.

\section{Mathematical Expressions}
Here are some example equations:

\subsection{Quadratic Formula}
The quadratic formula is given by:
\[x = \frac{-b \pm \sqrt{b^2 - 4ac}}{2a}\]

\subsection{Matrix Example}
A sample matrix:
\[
\begin{pmatrix}
a_{11} & a_{12} & a_{13} \\
a_{21} & a_{22} & a_{23} \\
a_{31} & a_{32} & a_{33}
\end{pmatrix}
\]

\section{Tables}
Below is a sample table showing data:

\begin{table}[h]
\centering
\begin{tabular}{@{}lrrr@{}}
\toprule
\textbf{Category} & \textbf{Value 1} & \textbf{Value 2} & \textbf{Total} \\
\midrule
Item A & 10.5 & 20.3 & 30.8 \\
Item B & 15.2 & 25.1 & 40.3 \\
Item C & 22.1 & 12.7 & 34.8 \\
\bottomrule
\end{tabular}
\caption{Sample Data Table}
\label{tab:sample}
\end{table}

\section{Lists}
\subsection{Bullet Points}
\begin{itemize}
  \item First item
  \item Second item
  \item Third item with sub-items:
    \begin{itemize}
      \item Sub-item 1
      \item Sub-item 2
    \end{itemize}
\end{itemize}

\subsection{Numbered List}
\begin{enumerate}
  \item First step
  \item Second step
  \item Third step
\end{enumerate}

\section{Citations and References}
When you use the citation search feature, references will be automatically managed. For example, you might cite a paper like this: \cite{example2024}. All citations will be automatically added to your \texttt{cite.bib} file.

\section{Conclusion}
This template demonstrates basic \LaTeX{} features including mathematical formulas, tables, and lists. The bibliography system is now set up to work seamlessly with Bibby AI's citation management. Feel free to modify it for your needs!

\bibliographystyle{plain}
\bibliography{cite}

\end{document}\documentclass{article}
\usepackage[utf8]{inputenc}
\usepackage{booktabs}
\usepackage{amsmath}
\usepackage{natbib}

\title{My Document}
\author{Your Name}
\date{\today}

\begin{document}

\maketitle

\section{Introduction}
This is a sample document demonstrating various \LaTeX{} features including mathematical equations, tables, and formatting. You can cite papers using the citation search feature, and they will be automatically added to your \texttt{cite.bib} file.

\section{Mathematical Expressions}
Here are some example equations:

\subsection{Quadratic Formula}
The quadratic formula is given by:
\[x = \frac{-b \pm \sqrt{b^2 - 4ac}}{2a}\]

\subsection{Matrix Example}
A sample matrix:
\[
\begin{pmatrix}
a_{11} & a_{12} & a_{13} \\
a_{21} & a_{22} & a_{23} \\
a_{31} & a_{32} & a_{33}
\end{pmatrix}
\]

\section{Tables}
Below is a sample table showing data:

\begin{table}[h]
\centering
\begin{tabular}{@{}lrrr@{}}
\toprule
\textbf{Category} & \textbf{Value 1} & \textbf{Value 2} & \textbf{Total} \\
\midrule
Item A & 10.5 & 20.3 & 30.8 \\
Item B & 15.2 & 25.1 & 40.3 \\
Item C & 22.1 & 12.7 & 34.8 \\
\bottomrule
\end{tabular}
\caption{Sample Data Table}
\label{tab:sample}
\end{table}

\section{Lists}
\subsection{Bullet Points}
\begin{itemize}
  \item First item
  \item Second item
  \item Third item with sub-items:
    \begin{itemize}
      \item Sub-item 1
      \item Sub-item 2
    \end{itemize}
\end{itemize}

\subsection{Numbered List}
\begin{enumerate}
  \item First step
  \item Second step
  \item Third step
\end{enumerate}

\section{Citations and References}
When you use the citation search feature, references will be automatically managed. For example, you might cite a paper like this: \cite{example2024}. All citations will be automatically added to your \texttt{cite.bib} file.

\section{Conclusion}
This template demonstrates basic \LaTeX{} features including mathematical formulas, tables, and lists. The bibliography system is now set up to work seamlessly with Bibby AI's citation management. Feel free to modify it for your needs!

\bibliographystyle{plain}
\bibliography{cite}

\end{document}\documentclass{article}
\usepackage[utf8]{inputenc}
\usepackage{booktabs}
\usepackage{amsmath}
\usepackage{natbib}

\title{My Document}
\author{Your Name}
\date{\today}

\begin{document}

\maketitle

\section{Introduction}
This is a sample document demonstrating various \LaTeX{} features including mathematical equations, tables, and formatting. You can cite papers using the citation search feature, and they will be automatically added to your \texttt{cite.bib} file.

\section{Mathematical Expressions}
Here are some example equations:

\subsection{Quadratic Formula}
The quadratic formula is given by:
\[x = \frac{-b \pm \sqrt{b^2 - 4ac}}{2a}\]

\subsection{Matrix Example}
A sample matrix:
\[
\begin{pmatrix}
a_{11} & a_{12} & a_{13} \\
a_{21} & a_{22} & a_{23} \\
a_{31} & a_{32} & a_{33}
\end{pmatrix}
\]

\section{Tables}
Below is a sample table showing data:

\begin{table}[h]
\centering
\begin{tabular}{@{}lrrr@{}}
\toprule
\textbf{Category} & \textbf{Value 1} & \textbf{Value 2} & \textbf{Total} \\
\midrule
Item A & 10.5 & 20.3 & 30.8 \\
Item B & 15.2 & 25.1 & 40.3 \\
Item C & 22.1 & 12.7 & 34.8 \\
\bottomrule
\end{tabular}
\caption{Sample Data Table}
\label{tab:sample}
\end{table}

\section{Lists}
\subsection{Bullet Points}
\begin{itemize}
  \item First item
  \item Second item
  \item Third item with sub-items:
    \begin{itemize}
      \item Sub-item 1
      \item Sub-item 2
    \end{itemize}
\end{itemize}

\subsection{Numbered List}
\begin{enumerate}
  \item First step
  \item Second step
  \item Third step
\end{enumerate}

\section{Citations and References}
When you use the citation search feature, references will be automatically managed. For example, you might cite a paper like this: \cite{example2024}. All citations will be automatically added to your \texttt{cite.bib} file.

\section{Conclusion}
This template demonstrates basic \LaTeX{} features including mathematical formulas, tables, and lists. The bibliography system is now set up to work seamlessly with Bibby AI's citation management. Feel free to modify it for your needs!

\bibliographystyle{plain}
\bibliography{cite}

\end{document}\documentclass{article}
\usepackage[utf8]{inputenc}
\usepackage{booktabs}
\usepackage{amsmath}
\usepackage{natbib}

\title{My Document}
\author{Your Name}
\date{\today}

\begin{document}

\maketitle

\section{Introduction}
This is a sample document demonstrating various \LaTeX{} features including mathematical equations, tables, and formatting. You can cite papers using the citation search feature, and they will be automatically added to your \texttt{cite.bib} file.

\section{Mathematical Expressions}
Here are some example equations:

\subsection{Quadratic Formula}
The quadratic formula is given by:
\[x = \frac{-b \pm \sqrt{b^2 - 4ac}}{2a}\]

\subsection{Matrix Example}
A sample matrix:
\[
\begin{pmatrix}
a_{11} & a_{12} & a_{13} \\
a_{21} & a_{22} & a_{23} \\
a_{31} & a_{32} & a_{33}
\end{pmatrix}
\]

\section{Tables}
Below is a sample table showing data:

\begin{table}[h]
\centering
\begin{tabular}{@{}lrrr@{}}
\toprule
\textbf{Category} & \textbf{Value 1} & \textbf{Value 2} & \textbf{Total} \\
\midrule
Item A & 10.5 & 20.3 & 30.8 \\
Item B & 15.2 & 25.1 & 40.3 \\
Item C & 22.1 & 12.7 & 34.8 \\
\bottomrule
\end{tabular}
\caption{Sample Data Table}
\label{tab:sample}
\end{table}

\section{Lists}
\subsection{Bullet Points}
\begin{itemize}
  \item First item
  \item Second item
  \item Third item with sub-items:
    \begin{itemize}
      \item Sub-item 1
      \item Sub-item 2
    \end{itemize}
\end{itemize}

\subsection{Numbered List}
\begin{enumerate}
  \item First step
  \item Second step
  \item Third step
\end{enumerate}

\section{Citations and References}
When you use the citation search feature, references will be automatically managed. For example, you might cite a paper like this: \cite{example2024}. All citations will be automatically added to your \texttt{cite.bib} file.

\section{Conclusion}
This template demonstrates basic \LaTeX{} features including mathematical formulas, tables, and lists. The bibliography system is now set up to work seamlessly with Bibby AI's citation management. Feel free to modify it for your needs!

\bibliographystyle{plain}
\bibliography{cite}

\end{document}\documentclass{article}
\usepackage[utf8]{inputenc}
\usepackage{booktabs}
\usepackage{amsmath}
\usepackage{natbib}

\title{My Document}
\author{Your Name}
\date{\today}

\begin{document}

\maketitle

\section{Introduction}
This is a sample document demonstrating various \LaTeX{} features including mathematical equations, tables, and formatting. You can cite papers using the citation search feature, and they will be automatically added to your \texttt{cite.bib} file.

\section{Mathematical Expressions}
Here are some example equations:

\subsection{Quadratic Formula}
The quadratic formula is given by:
\[x = \frac{-b \pm \sqrt{b^2 - 4ac}}{2a}\]

\subsection{Matrix Example}
A sample matrix:
\[
\begin{pmatrix}
a_{11} & a_{12} & a_{13} \\
a_{21} & a_{22} & a_{23} \\
a_{31} & a_{32} & a_{33}
\end{pmatrix}
\]

\section{Tables}
Below is a sample table showing data:

\begin{table}[h]
\centering
\begin{tabular}{@{}lrrr@{}}
\toprule
\textbf{Category} & \textbf{Value 1} & \textbf{Value 2} & \textbf{Total} \\
\midrule
Item A & 10.5 & 20.3 & 30.8 \\
Item B & 15.2 & 25.1 & 40.3 \\
Item C & 22.1 & 12.7 & 34.8 \\
\bottomrule
\end{tabular}
\caption{Sample Data Table}
\label{tab:sample}
\end{table}

\section{Lists}
\subsection{Bullet Points}
\begin{itemize}
  \item First item
  \item Second item
  \item Third item with sub-items:
    \begin{itemize}
      \item Sub-item 1
      \item Sub-item 2
    \end{itemize}
\end{itemize}

\subsection{Numbered List}
\begin{enumerate}
  \item First step
  \item Second step
  \item Third step
\end{enumerate}

\section{Citations and References}
When you use the citation search feature, references will be automatically managed. For example, you might cite a paper like this: \cite{example2024}. All citations will be automatically added to your \texttt{cite.bib} file.

\section{Conclusion}
This template demonstrates basic \LaTeX{} features including mathematical formulas, tables, and lists. The bibliography system is now set up to work seamlessly with Bibby AI's citation management. Feel free to modify it for your needs!

\bibliographystyle{plain}
\bibliography{cite}

\end{document}\documentclass{article}
\usepackage[utf8]{inputenc}
\usepackage{booktabs}
\usepackage{amsmath}
\usepackage{natbib}

\title{My Document}
\author{Your Name}
\date{\today}

\begin{document}

\maketitle

\section{Introduction}
This is a sample document demonstrating various \LaTeX{} features including mathematical equations, tables, and formatting. You can cite papers using the citation search feature, and they will be automatically added to your \texttt{cite.bib} file.

\section{Mathematical Expressions}
Here are some example equations:

\subsection{Quadratic Formula}
The quadratic formula is given by:
\[x = \frac{-b \pm \sqrt{b^2 - 4ac}}{2a}\]

\subsection{Matrix Example}
A sample matrix:
\[
\begin{pmatrix}
a_{11} & a_{12} & a_{13} \\
a_{21} & a_{22} & a_{23} \\
a_{31} & a_{32} & a_{33}
\end{pmatrix}
\]

\section{Tables}
Below is a sample table showing data:

\begin{table}[h]
\centering
\begin{tabular}{@{}lrrr@{}}
\toprule
\textbf{Category} & \textbf{Value 1} & \textbf{Value 2} & \textbf{Total} \\
\midrule
Item A & 10.5 & 20.3 & 30.8 \\
Item B & 15.2 & 25.1 & 40.3 \\
Item C & 22.1 & 12.7 & 34.8 \\
\bottomrule
\end{tabular}
\caption{Sample Data Table}
\label{tab:sample}
\end{table}

\section{Lists}
\subsection{Bullet Points}
\begin{itemize}
  \item First item
  \item Second item
  \item Third item with sub-items:
    \begin{itemize}
      \item Sub-item 1
      \item Sub-item 2
    \end{itemize}
\end{itemize}

\subsection{Numbered List}
\begin{enumerate}
  \item First step
  \item Second step
  \item Third step
\end{enumerate}

\section{Citations and References}
When you use the citation search feature, references will be automatically managed. For example, you might cite a paper like this: \cite{example2024}. All citations will be automatically added to your \texttt{cite.bib} file.

\section{Conclusion}
This template demonstrates basic \LaTeX{} features including mathematical formulas, tables, and lists. The bibliography system is now set up to work seamlessly with Bibby AI's citation management. Feel free to modify it for your needs!

\bibliographystyle{plain}
\bibliography{cite}

\end{document}\documentclass{article}
\usepackage[utf8]{inputenc}
\usepackage{booktabs}
\usepackage{amsmath}
\usepackage{natbib}

\title{My Document}
\author{Your Name}
\date{\today}

\begin{document}

\maketitle

\section{Introduction}
This is a sample document demonstrating various \LaTeX{} features including mathematical equations, tables, and formatting. You can cite papers using the citation search feature, and they will be automatically added to your \texttt{cite.bib} file.

\section{Mathematical Expressions}
Here are some example equations:

\subsection{Quadratic Formula}
The quadratic formula is given by:
\[x = \frac{-b \pm \sqrt{b^2 - 4ac}}{2a}\]

\subsection{Matrix Example}
A sample matrix:
\[
\begin{pmatrix}
a_{11} & a_{12} & a_{13} \\
a_{21} & a_{22} & a_{23} \\
a_{31} & a_{32} & a_{33}
\end{pmatrix}
\]

\section{Tables}
Below is a sample table showing data:

\begin{table}[h]
\centering
\begin{tabular}{@{}lrrr@{}}
\toprule
\textbf{Category} & \textbf{Value 1} & \textbf{Value 2} & \textbf{Total} \\
\midrule
Item A & 10.5 & 20.3 & 30.8 \\
Item B & 15.2 & 25.1 & 40.3 \\
Item C & 22.1 & 12.7 & 34.8 \\
\bottomrule
\end{tabular}
\caption{Sample Data Table}
\label{tab:sample}
\end{table}

\section{Lists}
\subsection{Bullet Points}
\begin{itemize}
  \item First item
  \item Second item
  \item Third item with sub-items:
    \begin{itemize}
      \item Sub-item 1
      \item Sub-item 2
    \end{itemize}
\end{itemize}

\subsection{Numbered List}
\begin{enumerate}
  \item First step
  \item Second step
  \item Third step
\end{enumerate}

\section{Citations and References}
When you use the citation search feature, references will be automatically managed. For example, you might cite a paper like this: \cite{example2024}. All citations will be automatically added to your \texttt{cite.bib} file.

\section{Conclusion}
This template demonstrates basic \LaTeX{} features including mathematical formulas, tables, and lists. The bibliography system is now set up to work seamlessly with Bibby AI's citation management. Feel free to modify it for your needs!

\bibliographystyle{plain}
\bibliography{cite}

\end{document}\documentclass{article}
\usepackage[utf8]{inputenc}
\usepackage{booktabs}
\usepackage{amsmath}
\usepackage{natbib}

\title{My Document}
\author{Your Name}
\date{\today}

\begin{document}

\maketitle

\section{Introduction}
This is a sample document demonstrating various \LaTeX{} features including mathematical equations, tables, and formatting. You can cite papers using the citation search feature, and they will be automatically added to your \texttt{cite.bib} file.

\section{Mathematical Expressions}
Here are some example equations:

\subsection{Quadratic Formula}
The quadratic formula is given by:
\[x = \frac{-b \pm \sqrt{b^2 - 4ac}}{2a}\]

\subsection{Matrix Example}
A sample matrix:
\[
\begin{pmatrix}
a_{11} & a_{12} & a_{13} \\
a_{21} & a_{22} & a_{23} \\
a_{31} & a_{32} & a_{33}
\end{pmatrix}
\]

\section{Tables}
Below is a sample table showing data:

\begin{table}[h]
\centering
\begin{tabular}{@{}lrrr@{}}
\toprule
\textbf{Category} & \textbf{Value 1} & \textbf{Value 2} & \textbf{Total} \\
\midrule
Item A & 10.5 & 20.3 & 30.8 \\
Item B & 15.2 & 25.1 & 40.3 \\
Item C & 22.1 & 12.7 & 34.8 \\
\bottomrule
\end{tabular}
\caption{Sample Data Table}
\label{tab:sample}
\end{table}

\section{Lists}
\subsection{Bullet Points}
\begin{itemize}
  \item First item
  \item Second item
  \item Third item with sub-items:
    \begin{itemize}
      \item Sub-item 1
      \item Sub-item 2
    \end{itemize}
\end{itemize}

\subsection{Numbered List}
\begin{enumerate}
  \item First step
  \item Second step
  \item Third step
\end{enumerate}

\section{Citations and References}
When you use the citation search feature, references will be automatically managed. For example, you might cite a paper like this: \cite{example2024}. All citations will be automatically added to your \texttt{cite.bib} file.

\section{Conclusion}
This template demonstrates basic \LaTeX{} features including mathematical formulas, tables, and lists. The bibliography system is now set up to work seamlessly with Bibby AI's citation management. Feel free to modify it for your needs!

\bibliographystyle{plain}
\bibliography{cite}

\end{document}\documentclass{article}
\usepackage[utf8]{inputenc}
\usepackage{booktabs}
\usepackage{amsmath}
\usepackage{natbib}

\title{My Document}
\author{Your Name}
\date{\today}

\begin{document}

\maketitle

\section{Introduction}
This is a sample document demonstrating various \LaTeX{} features including mathematical equations, tables, and formatting. You can cite papers using the citation search feature, and they will be automatically added to your \texttt{cite.bib} file.

\section{Mathematical Expressions}
Here are some example equations:

\subsection{Quadratic Formula}
The quadratic formula is given by:
\[x = \frac{-b \pm \sqrt{b^2 - 4ac}}{2a}\]

\subsection{Matrix Example}
A sample matrix:
\[
\begin{pmatrix}
a_{11} & a_{12} & a_{13} \\
a_{21} & a_{22} & a_{23} \\
a_{31} & a_{32} & a_{33}
\end{pmatrix}
\]

\section{Tables}
Below is a sample table showing data:

\begin{table}[h]
\centering
\begin{tabular}{@{}lrrr@{}}
\toprule
\textbf{Category} & \textbf{Value 1} & \textbf{Value 2} & \textbf{Total} \\
\midrule
Item A & 10.5 & 20.3 & 30.8 \\
Item B & 15.2 & 25.1 & 40.3 \\
Item C & 22.1 & 12.7 & 34.8 \\
\bottomrule
\end{tabular}
\caption{Sample Data Table}
\label{tab:sample}
\end{table}

\section{Lists}
\subsection{Bullet Points}
\begin{itemize}
  \item First item
  \item Second item
  \item Third item with sub-items:
    \begin{itemize}
      \item Sub-item 1
      \item Sub-item 2
    \end{itemize}
\end{itemize}

\subsection{Numbered List}
\begin{enumerate}
  \item First step
  \item Second step
  \item Third step
\end{enumerate}

\section{Citations and References}
When you use the citation search feature, references will be automatically managed. For example, you might cite a paper like this: \cite{example2024}. All citations will be automatically added to your \texttt{cite.bib} file.

\section{Conclusion}
This template demonstrates basic \LaTeX{} features including mathematical formulas, tables, and lists. The bibliography system is now set up to work seamlessly with Bibby AI's citation management. Feel free to modify it for your needs!

\bibliographystyle{plain}
\bibliography{cite}

\end{document}\documentclass{article}
\usepackage[utf8]{inputenc}
\usepackage{booktabs}
\usepackage{amsmath}
\usepackage{natbib}

\title{My Document}
\author{Your Name}
\date{\today}

\begin{document}

\maketitle

\section{Introduction}
This is a sample document demonstrating various \LaTeX{} features including mathematical equations, tables, and formatting. You can cite papers using the citation search feature, and they will be automatically added to your \texttt{cite.bib} file.

\section{Mathematical Expressions}
Here are some example equations:

\subsection{Quadratic Formula}
The quadratic formula is given by:
\[x = \frac{-b \pm \sqrt{b^2 - 4ac}}{2a}\]

\subsection{Matrix Example}
A sample matrix:
\[
\begin{pmatrix}
a_{11} & a_{12} & a_{13} \\
a_{21} & a_{22} & a_{23} \\
a_{31} & a_{32} & a_{33}
\end{pmatrix}
\]

\section{Tables}
Below is a sample table showing data:

\begin{table}[h]
\centering
\begin{tabular}{@{}lrrr@{}}
\toprule
\textbf{Category} & \textbf{Value 1} & \textbf{Value 2} & \textbf{Total} \\
\midrule
Item A & 10.5 & 20.3 & 30.8 \\
Item B & 15.2 & 25.1 & 40.3 \\
Item C & 22.1 & 12.7 & 34.8 \\
\bottomrule
\end{tabular}
\caption{Sample Data Table}
\label{tab:sample}
\end{table}

\section{Lists}
\subsection{Bullet Points}
\begin{itemize}
  \item First item
  \item Second item
  \item Third item with sub-items:
    \begin{itemize}
      \item Sub-item 1
      \item Sub-item 2
    \end{itemize}
\end{itemize}

\subsection{Numbered List}
\begin{enumerate}
  \item First step
  \item Second step
  \item Third step
\end{enumerate}

\section{Citations and References}
When you use the citation search feature, references will be automatically managed. For example, you might cite a paper like this: \cite{example2024}. All citations will be automatically added to your \texttt{cite.bib} file.

\section{Conclusion}
This template demonstrates basic \LaTeX{} features including mathematical formulas, tables, and lists. The bibliography system is now set up to work seamlessly with Bibby AI's citation management. Feel free to modify it for your needs!

\bibliographystyle{plain}
\bibliography{cite}

\end{document}\documentclass{article}
\usepackage[utf8]{inputenc}
\usepackage{booktabs}
\usepackage{amsmath}
\usepackage{natbib}

\title{My Document}
\author{Your Name}
\date{\today}

\begin{document}

\maketitle

\section{Introduction}
This is a sample document demonstrating various \LaTeX{} features including mathematical equations, tables, and formatting. You can cite papers using the citation search feature, and they will be automatically added to your \texttt{cite.bib} file.

\section{Mathematical Expressions}
Here are some example equations:

\subsection{Quadratic Formula}
The quadratic formula is given by:
\[x = \frac{-b \pm \sqrt{b^2 - 4ac}}{2a}\]

\subsection{Matrix Example}
A sample matrix:
\[
\begin{pmatrix}
a_{11} & a_{12} & a_{13} \\
a_{21} & a_{22} & a_{23} \\
a_{31} & a_{32} & a_{33}
\end{pmatrix}
\]

\section{Tables}
Below is a sample table showing data:

\begin{table}[h]
\centering
\begin{tabular}{@{}lrrr@{}}
\toprule
\textbf{Category} & \textbf{Value 1} & \textbf{Value 2} & \textbf{Total} \\
\midrule
Item A & 10.5 & 20.3 & 30.8 \\
Item B & 15.2 & 25.1 & 40.3 \\
Item C & 22.1 & 12.7 & 34.8 \\
\bottomrule
\end{tabular}
\caption{Sample Data Table}
\label{tab:sample}
\end{table}

\section{Lists}
\subsection{Bullet Points}
\begin{itemize}
  \item First item
  \item Second item
  \item Third item with sub-items:
    \begin{itemize}
      \item Sub-item 1
      \item Sub-item 2
    \end{itemize}
\end{itemize}

\subsection{Numbered List}
\begin{enumerate}
  \item First step
  \item Second step
  \item Third step
\end{enumerate}

\section{Citations and References}
When you use the citation search feature, references will be automatically managed. For example, you might cite a paper like this: \cite{example2024}. All citations will be automatically added to your \texttt{cite.bib} file.

\section{Conclusion}
This template demonstrates basic \LaTeX{} features including mathematical formulas, tables, and lists. The bibliography system is now set up to work seamlessly with Bibby AI's citation management. Feel free to modify it for your needs!

\bibliographystyle{plain}
\bibliography{cite}

\end{document}\documentclass{article}
\usepackage[utf8]{inputenc}
\usepackage{booktabs}
\usepackage{amsmath}
\usepackage{natbib}

\title{My Document}
\author{Your Name}
\date{\today}

\begin{document}

\maketitle

\section{Introduction}
This is a sample document demonstrating various \LaTeX{} features including mathematical equations, tables, and formatting. You can cite papers using the citation search feature, and they will be automatically added to your \texttt{cite.bib} file.

\section{Mathematical Expressions}
Here are some example equations:

\subsection{Quadratic Formula}
The quadratic formula is given by:
\[x = \frac{-b \pm \sqrt{b^2 - 4ac}}{2a}\]

\subsection{Matrix Example}
A sample matrix:
\[
\begin{pmatrix}
a_{11} & a_{12} & a_{13} \\
a_{21} & a_{22} & a_{23} \\
a_{31} & a_{32} & a_{33}
\end{pmatrix}
\]

\section{Tables}
Below is a sample table showing data:

\begin{table}[h]
\centering
\begin{tabular}{@{}lrrr@{}}
\toprule
\textbf{Category} & \textbf{Value 1} & \textbf{Value 2} & \textbf{Total} \\
\midrule
Item A & 10.5 & 20.3 & 30.8 \\
Item B & 15.2 & 25.1 & 40.3 \\
Item C & 22.1 & 12.7 & 34.8 \\
\bottomrule
\end{tabular}
\caption{Sample Data Table}
\label{tab:sample}
\end{table}

\section{Lists}
\subsection{Bullet Points}
\begin{itemize}
  \item First item
  \item Second item
  \item Third item with sub-items:
    \begin{itemize}
      \item Sub-item 1
      \item Sub-item 2
    \end{itemize}
\end{itemize}

\subsection{Numbered List}
\begin{enumerate}
  \item First step
  \item Second step
  \item Third step
\end{enumerate}

\section{Citations and References}
When you use the citation search feature, references will be automatically managed. For example, you might cite a paper like this: \cite{example2024}. All citations will be automatically added to your \texttt{cite.bib} file.

\section{Conclusion}
This template demonstrates basic \LaTeX{} features including mathematical formulas, tables, and lists. The bibliography system is now set up to work seamlessly with Bibby AI's citation management. Feel free to modify it for your needs!

\bibliographystyle{plain}
\bibliography{cite}

\end{document}\documentclass{article}
\usepackage[utf8]{inputenc}
\usepackage{booktabs}
\usepackage{amsmath}
\usepackage{natbib}

\title{My Document}
\author{Your Name}
\date{\today}

\begin{document}

\maketitle

\section{Introduction}
This is a sample document demonstrating various \LaTeX{} features including mathematical equations, tables, and formatting. You can cite papers using the citation search feature, and they will be automatically added to your \texttt{cite.bib} file.

\section{Mathematical Expressions}
Here are some example equations:

\subsection{Quadratic Formula}
The quadratic formula is given by:
\[x = \frac{-b \pm \sqrt{b^2 - 4ac}}{2a}\]

\subsection{Matrix Example}
A sample matrix:
\[
\begin{pmatrix}
a_{11} & a_{12} & a_{13} \\
a_{21} & a_{22} & a_{23} \\
a_{31} & a_{32} & a_{33}
\end{pmatrix}
\]

\section{Tables}
Below is a sample table showing data:

\begin{table}[h]
\centering
\begin{tabular}{@{}lrrr@{}}
\toprule
\textbf{Category} & \textbf{Value 1} & \textbf{Value 2} & \textbf{Total} \\
\midrule
Item A & 10.5 & 20.3 & 30.8 \\
Item B & 15.2 & 25.1 & 40.3 \\
Item C & 22.1 & 12.7 & 34.8 \\
\bottomrule
\end{tabular}
\caption{Sample Data Table}
\label{tab:sample}
\end{table}

\section{Lists}
\subsection{Bullet Points}
\begin{itemize}
  \item First item
  \item Second item
  \item Third item with sub-items:
    \begin{itemize}
      \item Sub-item 1
      \item Sub-item 2
    \end{itemize}
\end{itemize}

\subsection{Numbered List}
\begin{enumerate}
  \item First step
  \item Second step
  \item Third step
\end{enumerate}

\section{Citations and References}
When you use the citation search feature, references will be automatically managed. For example, you might cite a paper like this: \cite{example2024}. All citations will be automatically added to your \texttt{cite.bib} file.

\section{Conclusion}
This template demonstrates basic \LaTeX{} features including mathematical formulas, tables, and lists. The bibliography system is now set up to work seamlessly with Bibby AI's citation management. Feel free to modify it for your needs!

\bibliographystyle{plain}
\bibliography{cite}

\end{document}\documentclass{article}
\usepackage[utf8]{inputenc}
\usepackage{booktabs}
\usepackage{amsmath}
\usepackage{natbib}

\title{My Document}
\author{Your Name}
\date{\today}

\begin{document}

\maketitle

\section{Introduction}
This is a sample document demonstrating various \LaTeX{} features including mathematical equations, tables, and formatting. You can cite papers using the citation search feature, and they will be automatically added to your \texttt{cite.bib} file.

\section{Mathematical Expressions}
Here are some example equations:

\subsection{Quadratic Formula}
The quadratic formula is given by:
\[x = \frac{-b \pm \sqrt{b^2 - 4ac}}{2a}\]

\subsection{Matrix Example}
A sample matrix:
\[
\begin{pmatrix}
a_{11} & a_{12} & a_{13} \\
a_{21} & a_{22} & a_{23} \\
a_{31} & a_{32} & a_{33}
\end{pmatrix}
\]

\section{Tables}
Below is a sample table showing data:

\begin{table}[h]
\centering
\begin{tabular}{@{}lrrr@{}}
\toprule
\textbf{Category} & \textbf{Value 1} & \textbf{Value 2} & \textbf{Total} \\
\midrule
Item A & 10.5 & 20.3 & 30.8 \\
Item B & 15.2 & 25.1 & 40.3 \\
Item C & 22.1 & 12.7 & 34.8 \\
\bottomrule
\end{tabular}
\caption{Sample Data Table}
\label{tab:sample}
\end{table}

\section{Lists}
\subsection{Bullet Points}
\begin{itemize}
  \item First item
  \item Second item
  \item Third item with sub-items:
    \begin{itemize}
      \item Sub-item 1
      \item Sub-item 2
    \end{itemize}
\end{itemize}

\subsection{Numbered List}
\begin{enumerate}
  \item First step
  \item Second step
  \item Third step
\end{enumerate}

\section{Citations and References}
When you use the citation search feature, references will be automatically managed. For example, you might cite a paper like this: \cite{example2024}. All citations will be automatically added to your \texttt{cite.bib} file.

\section{Conclusion}
This template demonstrates basic \LaTeX{} features including mathematical formulas, tables, and lists. The bibliography system is now set up to work seamlessly with Bibby AI's citation management. Feel free to modify it for your needs!

\bibliographystyle{plain}
\bibliography{cite}

\end{document}\documentclass{article}
\usepackage[utf8]{inputenc}
\usepackage{booktabs}
\usepackage{amsmath}
\usepackage{natbib}

\title{My Document}
\author{Your Name}
\date{\today}

\begin{document}

\maketitle

\section{Introduction}
This is a sample document demonstrating various \LaTeX{} features including mathematical equations, tables, and formatting. You can cite papers using the citation search feature, and they will be automatically added to your \texttt{cite.bib} file.

\section{Mathematical Expressions}
Here are some example equations:

\subsection{Quadratic Formula}
The quadratic formula is given by:
\[x = \frac{-b \pm \sqrt{b^2 - 4ac}}{2a}\]

\subsection{Matrix Example}
A sample matrix:
\[
\begin{pmatrix}
a_{11} & a_{12} & a_{13} \\
a_{21} & a_{22} & a_{23} \\
a_{31} & a_{32} & a_{33}
\end{pmatrix}
\]

\section{Tables}
Below is a sample table showing data:

\begin{table}[h]
\centering
\begin{tabular}{@{}lrrr@{}}
\toprule
\textbf{Category} & \textbf{Value 1} & \textbf{Value 2} & \textbf{Total} \\
\midrule
Item A & 10.5 & 20.3 & 30.8 \\
Item B & 15.2 & 25.1 & 40.3 \\
Item C & 22.1 & 12.7 & 34.8 \\
\bottomrule
\end{tabular}
\caption{Sample Data Table}
\label{tab:sample}
\end{table}

\section{Lists}
\subsection{Bullet Points}
\begin{itemize}
  \item First item
  \item Second item
  \item Third item with sub-items:
    \begin{itemize}
      \item Sub-item 1
      \item Sub-item 2
    \end{itemize}
\end{itemize}

\subsection{Numbered List}
\begin{enumerate}
  \item First step
  \item Second step
  \item Third step
\end{enumerate}

\section{Citations and References}
When you use the citation search feature, references will be automatically managed. For example, you might cite a paper like this: \cite{example2024}. All citations will be automatically added to your \texttt{cite.bib} file.

\section{Conclusion}
This template demonstrates basic \LaTeX{} features including mathematical formulas, tables, and lists. The bibliography system is now set up to work seamlessly with Bibby AI's citation management. Feel free to modify it for your needs!

\bibliographystyle{plain}
\bibliography{cite}

\end{document}\documentclass{article}
\usepackage[utf8]{inputenc}
\usepackage{booktabs}
\usepackage{amsmath}
\usepackage{natbib}

\title{My Document}
\author{Your Name}
\date{\today}

\begin{document}

\maketitle

\section{Introduction}
This is a sample document demonstrating various \LaTeX{} features including mathematical equations, tables, and formatting. You can cite papers using the citation search feature, and they will be automatically added to your \texttt{cite.bib} file.

\section{Mathematical Expressions}
Here are some example equations:

\subsection{Quadratic Formula}
The quadratic formula is given by:
\[x = \frac{-b \pm \sqrt{b^2 - 4ac}}{2a}\]

\subsection{Matrix Example}
A sample matrix:
\[
\begin{pmatrix}
a_{11} & a_{12} & a_{13} \\
a_{21} & a_{22} & a_{23} \\
a_{31} & a_{32} & a_{33}
\end{pmatrix}
\]

\section{Tables}
Below is a sample table showing data:

\begin{table}[h]
\centering
\begin{tabular}{@{}lrrr@{}}
\toprule
\textbf{Category} & \textbf{Value 1} & \textbf{Value 2} & \textbf{Total} \\
\midrule
Item A & 10.5 & 20.3 & 30.8 \\
Item B & 15.2 & 25.1 & 40.3 \\
Item C & 22.1 & 12.7 & 34.8 \\
\bottomrule
\end{tabular}
\caption{Sample Data Table}
\label{tab:sample}
\end{table}

\section{Lists}
\subsection{Bullet Points}
\begin{itemize}
  \item First item
  \item Second item
  \item Third item with sub-items:
    \begin{itemize}
      \item Sub-item 1
      \item Sub-item 2
    \end{itemize}
\end{itemize}

\subsection{Numbered List}
\begin{enumerate}
  \item First step
  \item Second step
  \item Third step
\end{enumerate}

\section{Citations and References}
When you use the citation search feature, references will be automatically managed. For example, you might cite a paper like this: \cite{example2024}. All citations will be automatically added to your \texttt{cite.bib} file.

\section{Conclusion}
This template demonstrates basic \LaTeX{} features including mathematical formulas, tables, and lists. The bibliography system is now set up to work seamlessly with Bibby AI's citation management. Feel free to modify it for your needs!

\bibliographystyle{plain}
\bibliography{cite}

\end{document}\documentclass{article}
\usepackage[utf8]{inputenc}
\usepackage{booktabs}
\usepackage{amsmath}
\usepackage{natbib}

\title{My Document}
\author{Your Name}
\date{\today}

\begin{document}

\maketitle

\section{Introduction}
This is a sample document demonstrating various \LaTeX{} features including mathematical equations, tables, and formatting. You can cite papers using the citation search feature, and they will be automatically added to your \texttt{cite.bib} file.

\section{Mathematical Expressions}
Here are some example equations:

\subsection{Quadratic Formula}
The quadratic formula is given by:
\[x = \frac{-b \pm \sqrt{b^2 - 4ac}}{2a}\]

\subsection{Matrix Example}
A sample matrix:
\[
\begin{pmatrix}
a_{11} & a_{12} & a_{13} \\
a_{21} & a_{22} & a_{23} \\
a_{31} & a_{32} & a_{33}
\end{pmatrix}
\]

\section{Tables}
Below is a sample table showing data:

\begin{table}[h]
\centering
\begin{tabular}{@{}lrrr@{}}
\toprule
\textbf{Category} & \textbf{Value 1} & \textbf{Value 2} & \textbf{Total} \\
\midrule
Item A & 10.5 & 20.3 & 30.8 \\
Item B & 15.2 & 25.1 & 40.3 \\
Item C & 22.1 & 12.7 & 34.8 \\
\bottomrule
\end{tabular}
\caption{Sample Data Table}
\label{tab:sample}
\end{table}

\section{Lists}
\subsection{Bullet Points}
\begin{itemize}
  \item First item
  \item Second item
  \item Third item with sub-items:
    \begin{itemize}
      \item Sub-item 1
      \item Sub-item 2
    \end{itemize}
\end{itemize}

\subsection{Numbered List}
\begin{enumerate}
  \item First step
  \item Second step
  \item Third step
\end{enumerate}

\section{Citations and References}
When you use the citation search feature, references will be automatically managed. For example, you might cite a paper like this: \cite{example2024}. All citations will be automatically added to your \texttt{cite.bib} file.

\section{Conclusion}
This template demonstrates basic \LaTeX{} features including mathematical formulas, tables, and lists. The bibliography system is now set up to work seamlessly with Bibby AI's citation management. Feel free to modify it for your needs!

\bibliographystyle{plain}
\bibliography{cite}

\end{document}\documentclass{article}
\usepackage[utf8]{inputenc}
\usepackage{booktabs}
\usepackage{amsmath}
\usepackage{natbib}

\title{My Document}
\author{Your Name}
\date{\today}

\begin{document}

\maketitle

\section{Introduction}
This is a sample document demonstrating various \LaTeX{} features including mathematical equations, tables, and formatting. You can cite papers using the citation search feature, and they will be automatically added to your \texttt{cite.bib} file.

\section{Mathematical Expressions}
Here are some example equations:

\subsection{Quadratic Formula}
The quadratic formula is given by:
\[x = \frac{-b \pm \sqrt{b^2 - 4ac}}{2a}\]

\subsection{Matrix Example}
A sample matrix:
\[
\begin{pmatrix}
a_{11} & a_{12} & a_{13} \\
a_{21} & a_{22} & a_{23} \\
a_{31} & a_{32} & a_{33}
\end{pmatrix}
\]

\section{Tables}
Below is a sample table showing data:

\begin{table}[h]
\centering
\begin{tabular}{@{}lrrr@{}}
\toprule
\textbf{Category} & \textbf{Value 1} & \textbf{Value 2} & \textbf{Total} \\
\midrule
Item A & 10.5 & 20.3 & 30.8 \\
Item B & 15.2 & 25.1 & 40.3 \\
Item C & 22.1 & 12.7 & 34.8 \\
\bottomrule
\end{tabular}
\caption{Sample Data Table}
\label{tab:sample}
\end{table}

\section{Lists}
\subsection{Bullet Points}
\begin{itemize}
  \item First item
  \item Second item
  \item Third item with sub-items:
    \begin{itemize}
      \item Sub-item 1
      \item Sub-item 2
    \end{itemize}
\end{itemize}

\subsection{Numbered List}
\begin{enumerate}
  \item First step
  \item Second step
  \item Third step
\end{enumerate}

\section{Citations and References}
When you use the citation search feature, references will be automatically managed. For example, you might cite a paper like this: \cite{example2024}. All citations will be automatically added to your \texttt{cite.bib} file.

\section{Conclusion}
This template demonstrates basic \LaTeX{} features including mathematical formulas, tables, and lists. The bibliography system is now set up to work seamlessly with Bibby AI's citation management. Feel free to modify it for your needs!

\bibliographystyle{plain}
\bibliography{cite}

\end{document}\documentclass{article}
\usepackage[utf8]{inputenc}
\usepackage{booktabs}
\usepackage{amsmath}
\usepackage{natbib}

\title{My Document}
\author{Your Name}
\date{\today}

\begin{document}

\maketitle

\section{Introduction}
This is a sample document demonstrating various \LaTeX{} features including mathematical equations, tables, and formatting. You can cite papers using the citation search feature, and they will be automatically added to your \texttt{cite.bib} file.

\section{Mathematical Expressions}
Here are some example equations:

\subsection{Quadratic Formula}
The quadratic formula is given by:
\[x = \frac{-b \pm \sqrt{b^2 - 4ac}}{2a}\]

\subsection{Matrix Example}
A sample matrix:
\[
\begin{pmatrix}
a_{11} & a_{12} & a_{13} \\
a_{21} & a_{22} & a_{23} \\
a_{31} & a_{32} & a_{33}
\end{pmatrix}
\]

\section{Tables}
Below is a sample table showing data:

\begin{table}[h]
\centering
\begin{tabular}{@{}lrrr@{}}
\toprule
\textbf{Category} & \textbf{Value 1} & \textbf{Value 2} & \textbf{Total} \\
\midrule
Item A & 10.5 & 20.3 & 30.8 \\
Item B & 15.2 & 25.1 & 40.3 \\
Item C & 22.1 & 12.7 & 34.8 \\
\bottomrule
\end{tabular}
\caption{Sample Data Table}
\label{tab:sample}
\end{table}

\section{Lists}
\subsection{Bullet Points}
\begin{itemize}
  \item First item
  \item Second item
  \item Third item with sub-items:
    \begin{itemize}
      \item Sub-item 1
      \item Sub-item 2
    \end{itemize}
\end{itemize}

\subsection{Numbered List}
\begin{enumerate}
  \item First step
  \item Second step
  \item Third step
\end{enumerate}

\section{Citations and References}
When you use the citation search feature, references will be automatically managed. For example, you might cite a paper like this: \cite{example2024}. All citations will be automatically added to your \texttt{cite.bib} file.

\section{Conclusion}
This template demonstrates basic \LaTeX{} features including mathematical formulas, tables, and lists. The bibliography system is now set up to work seamlessly with Bibby AI's citation management. Feel free to modify it for your needs!

\bibliographystyle{plain}
\bibliography{cite}

\end{document}\documentclass{article}
\usepackage[utf8]{inputenc}
\usepackage{booktabs}
\usepackage{amsmath}
\usepackage{natbib}

\title{My Document}
\author{Your Name}
\date{\today}

\begin{document}

\maketitle

\section{Introduction}
This is a sample document demonstrating various \LaTeX{} features including mathematical equations, tables, and formatting. You can cite papers using the citation search feature, and they will be automatically added to your \texttt{cite.bib} file.

\section{Mathematical Expressions}
Here are some example equations:

\subsection{Quadratic Formula}
The quadratic formula is given by:
\[x = \frac{-b \pm \sqrt{b^2 - 4ac}}{2a}\]

\subsection{Matrix Example}
A sample matrix:
\[
\begin{pmatrix}
a_{11} & a_{12} & a_{13} \\
a_{21} & a_{22} & a_{23} \\
a_{31} & a_{32} & a_{33}
\end{pmatrix}
\]

\section{Tables}
Below is a sample table showing data:

\begin{table}[h]
\centering
\begin{tabular}{@{}lrrr@{}}
\toprule
\textbf{Category} & \textbf{Value 1} & \textbf{Value 2} & \textbf{Total} \\
\midrule
Item A & 10.5 & 20.3 & 30.8 \\
Item B & 15.2 & 25.1 & 40.3 \\
Item C & 22.1 & 12.7 & 34.8 \\
\bottomrule
\end{tabular}
\caption{Sample Data Table}
\label{tab:sample}
\end{table}

\section{Lists}
\subsection{Bullet Points}
\begin{itemize}
  \item First item
  \item Second item
  \item Third item with sub-items:
    \begin{itemize}
      \item Sub-item 1
      \item Sub-item 2
    \end{itemize}
\end{itemize}

\subsection{Numbered List}
\begin{enumerate}
  \item First step
  \item Second step
  \item Third step
\end{enumerate}

\section{Citations and References}
When you use the citation search feature, references will be automatically managed. For example, you might cite a paper like this: \cite{example2024}. All citations will be automatically added to your \texttt{cite.bib} file.

\section{Conclusion}
This template demonstrates basic \LaTeX{} features including mathematical formulas, tables, and lists. The bibliography system is now set up to work seamlessly with Bibby AI's citation management. Feel free to modify it for your needs!

\bibliographystyle{plain}
\bibliography{cite}

\end{document}\documentclass{article}
\usepackage[utf8]{inputenc}
\usepackage{booktabs}
\usepackage{amsmath}
\usepackage{natbib}

\title{My Document}
\author{Your Name}
\date{\today}

\begin{document}

\maketitle

\section{Introduction}
This is a sample document demonstrating various \LaTeX{} features including mathematical equations, tables, and formatting. You can cite papers using the citation search feature, and they will be automatically added to your \texttt{cite.bib} file.

\section{Mathematical Expressions}
Here are some example equations:

\subsection{Quadratic Formula}
The quadratic formula is given by:
\[x = \frac{-b \pm \sqrt{b^2 - 4ac}}{2a}\]

\subsection{Matrix Example}
A sample matrix:
\[
\begin{pmatrix}
a_{11} & a_{12} & a_{13} \\
a_{21} & a_{22} & a_{23} \\
a_{31} & a_{32} & a_{33}
\end{pmatrix}
\]

\section{Tables}
Below is a sample table showing data:

\begin{table}[h]
\centering
\begin{tabular}{@{}lrrr@{}}
\toprule
\textbf{Category} & \textbf{Value 1} & \textbf{Value 2} & \textbf{Total} \\
\midrule
Item A & 10.5 & 20.3 & 30.8 \\
Item B & 15.2 & 25.1 & 40.3 \\
Item C & 22.1 & 12.7 & 34.8 \\
\bottomrule
\end{tabular}
\caption{Sample Data Table}
\label{tab:sample}
\end{table}

\section{Lists}
\subsection{Bullet Points}
\begin{itemize}
  \item First item
  \item Second item
  \item Third item with sub-items:
    \begin{itemize}
      \item Sub-item 1
      \item Sub-item 2
    \end{itemize}
\end{itemize}

\subsection{Numbered List}
\begin{enumerate}
  \item First step
  \item Second step
  \item Third step
\end{enumerate}

\section{Citations and References}
When you use the citation search feature, references will be automatically managed. For example, you might cite a paper like this: \cite{example2024}. All citations will be automatically added to your \texttt{cite.bib} file.

\section{Conclusion}
This template demonstrates basic \LaTeX{} features including mathematical formulas, tables, and lists. The bibliography system is now set up to work seamlessly with Bibby AI's citation management. Feel free to modify it for your needs!

\bibliographystyle{plain}
\bibliography{cite}

\end{document}\documentclass{article}
\usepackage[utf8]{inputenc}
\usepackage{booktabs}
\usepackage{amsmath}
\usepackage{natbib}

\title{My Document}
\author{Your Name}
\date{\today}

\begin{document}

\maketitle

\section{Introduction}
This is a sample document demonstrating various \LaTeX{} features including mathematical equations, tables, and formatting. You can cite papers using the citation search feature, and they will be automatically added to your \texttt{cite.bib} file.

\section{Mathematical Expressions}
Here are some example equations:

\subsection{Quadratic Formula}
The quadratic formula is given by:
\[x = \frac{-b \pm \sqrt{b^2 - 4ac}}{2a}\]

\subsection{Matrix Example}
A sample matrix:
\[
\begin{pmatrix}
a_{11} & a_{12} & a_{13} \\
a_{21} & a_{22} & a_{23} \\
a_{31} & a_{32} & a_{33}
\end{pmatrix}
\]

\section{Tables}
Below is a sample table showing data:

\begin{table}[h]
\centering
\begin{tabular}{@{}lrrr@{}}
\toprule
\textbf{Category} & \textbf{Value 1} & \textbf{Value 2} & \textbf{Total} \\
\midrule
Item A & 10.5 & 20.3 & 30.8 \\
Item B & 15.2 & 25.1 & 40.3 \\
Item C & 22.1 & 12.7 & 34.8 \\
\bottomrule
\end{tabular}
\caption{Sample Data Table}
\label{tab:sample}
\end{table}

\section{Lists}
\subsection{Bullet Points}
\begin{itemize}
  \item First item
  \item Second item
  \item Third item with sub-items:
    \begin{itemize}
      \item Sub-item 1
      \item Sub-item 2
    \end{itemize}
\end{itemize}

\subsection{Numbered List}
\begin{enumerate}
  \item First step
  \item Second step
  \item Third step
\end{enumerate}

\section{Citations and References}
When you use the citation search feature, references will be automatically managed. For example, you might cite a paper like this: \cite{example2024}. All citations will be automatically added to your \texttt{cite.bib} file.

\section{Conclusion}
This template demonstrates basic \LaTeX{} features including mathematical formulas, tables, and lists. The bibliography system is now set up to work seamlessly with Bibby AI's citation management. Feel free to modify it for your needs!

\bibliographystyle{plain}
\bibliography{cite}

\end{document}\documentclass{article}
\usepackage[utf8]{inputenc}
\usepackage{booktabs}
\usepackage{amsmath}
\usepackage{natbib}

\title{My Document}
\author{Your Name}
\date{\today}

\begin{document}

\maketitle

\section{Introduction}
This is a sample document demonstrating various \LaTeX{} features including mathematical equations, tables, and formatting. You can cite papers using the citation search feature, and they will be automatically added to your \texttt{cite.bib} file.

\section{Mathematical Expressions}
Here are some example equations:

\subsection{Quadratic Formula}
The quadratic formula is given by:
\[x = \frac{-b \pm \sqrt{b^2 - 4ac}}{2a}\]

\subsection{Matrix Example}
A sample matrix:
\[
\begin{pmatrix}
a_{11} & a_{12} & a_{13} \\
a_{21} & a_{22} & a_{23} \\
a_{31} & a_{32} & a_{33}
\end{pmatrix}
\]

\section{Tables}
Below is a sample table showing data:

\begin{table}[h]
\centering
\begin{tabular}{@{}lrrr@{}}
\toprule
\textbf{Category} & \textbf{Value 1} & \textbf{Value 2} & \textbf{Total} \\
\midrule
Item A & 10.5 & 20.3 & 30.8 \\
Item B & 15.2 & 25.1 & 40.3 \\
Item C & 22.1 & 12.7 & 34.8 \\
\bottomrule
\end{tabular}
\caption{Sample Data Table}
\label{tab:sample}
\end{table}

\section{Lists}
\subsection{Bullet Points}
\begin{itemize}
  \item First item
  \item Second item
  \item Third item with sub-items:
    \begin{itemize}
      \item Sub-item 1
      \item Sub-item 2
    \end{itemize}
\end{itemize}

\subsection{Numbered List}
\begin{enumerate}
  \item First step
  \item Second step
  \item Third step
\end{enumerate}

\section{Citations and References}
When you use the citation search feature, references will be automatically managed. For example, you might cite a paper like this: \cite{example2024}. All citations will be automatically added to your \texttt{cite.bib} file.

\section{Conclusion}
This template demonstrates basic \LaTeX{} features including mathematical formulas, tables, and lists. The bibliography system is now set up to work seamlessly with Bibby AI's citation management. Feel free to modify it for your needs!

\bibliographystyle{plain}
\bibliography{cite}

\end{document}\documentclass{article}
\usepackage[utf8]{inputenc}
\usepackage{booktabs}
\usepackage{amsmath}
\usepackage{natbib}

\title{My Document}
\author{Your Name}
\date{\today}

\begin{document}

\maketitle

\section{Introduction}
This is a sample document demonstrating various \LaTeX{} features including mathematical equations, tables, and formatting. You can cite papers using the citation search feature, and they will be automatically added to your \texttt{cite.bib} file.

\section{Mathematical Expressions}
Here are some example equations:

\subsection{Quadratic Formula}
The quadratic formula is given by:
\[x = \frac{-b \pm \sqrt{b^2 - 4ac}}{2a}\]

\subsection{Matrix Example}
A sample matrix:
\[
\begin{pmatrix}
a_{11} & a_{12} & a_{13} \\
a_{21} & a_{22} & a_{23} \\
a_{31} & a_{32} & a_{33}
\end{pmatrix}
\]

\section{Tables}
Below is a sample table showing data:

\begin{table}[h]
\centering
\begin{tabular}{@{}lrrr@{}}
\toprule
\textbf{Category} & \textbf{Value 1} & \textbf{Value 2} & \textbf{Total} \\
\midrule
Item A & 10.5 & 20.3 & 30.8 \\
Item B & 15.2 & 25.1 & 40.3 \\
Item C & 22.1 & 12.7 & 34.8 \\
\bottomrule
\end{tabular}
\caption{Sample Data Table}
\label{tab:sample}
\end{table}

\section{Lists}
\subsection{Bullet Points}
\begin{itemize}
  \item First item
  \item Second item
  \item Third item with sub-items:
    \begin{itemize}
      \item Sub-item 1
      \item Sub-item 2
    \end{itemize}
\end{itemize}

\subsection{Numbered List}
\begin{enumerate}
  \item First step
  \item Second step
  \item Third step
\end{enumerate}

\section{Citations and References}
When you use the citation search feature, references will be automatically managed. For example, you might cite a paper like this: \cite{example2024}. All citations will be automatically added to your \texttt{cite.bib} file.

\section{Conclusion}
This template demonstrates basic \LaTeX{} features including mathematical formulas, tables, and lists. The bibliography system is now set up to work seamlessly with Bibby AI's citation management. Feel free to modify it for your needs!

\bibliographystyle{plain}
\bibliography{cite}

\end{document}\documentclass{article}
\usepackage[utf8]{inputenc}
\usepackage{booktabs}
\usepackage{amsmath}
\usepackage{natbib}

\title{My Document}
\author{Your Name}
\date{\today}

\begin{document}

\maketitle

\section{Introduction}
This is a sample document demonstrating various \LaTeX{} features including mathematical equations, tables, and formatting. You can cite papers using the citation search feature, and they will be automatically added to your \texttt{cite.bib} file.

\section{Mathematical Expressions}
Here are some example equations:

\subsection{Quadratic Formula}
The quadratic formula is given by:
\[x = \frac{-b \pm \sqrt{b^2 - 4ac}}{2a}\]

\subsection{Matrix Example}
A sample matrix:
\[
\begin{pmatrix}
a_{11} & a_{12} & a_{13} \\
a_{21} & a_{22} & a_{23} \\
a_{31} & a_{32} & a_{33}
\end{pmatrix}
\]

\section{Tables}
Below is a sample table showing data:

\begin{table}[h]
\centering
\begin{tabular}{@{}lrrr@{}}
\toprule
\textbf{Category} & \textbf{Value 1} & \textbf{Value 2} & \textbf{Total} \\
\midrule
Item A & 10.5 & 20.3 & 30.8 \\
Item B & 15.2 & 25.1 & 40.3 \\
Item C & 22.1 & 12.7 & 34.8 \\
\bottomrule
\end{tabular}
\caption{Sample Data Table}
\label{tab:sample}
\end{table}

\section{Lists}
\subsection{Bullet Points}
\begin{itemize}
  \item First item
  \item Second item
  \item Third item with sub-items:
    \begin{itemize}
      \item Sub-item 1
      \item Sub-item 2
    \end{itemize}
\end{itemize}

\subsection{Numbered List}
\begin{enumerate}
  \item First step
  \item Second step
  \item Third step
\end{enumerate}

\section{Citations and References}
When you use the citation search feature, references will be automatically managed. For example, you might cite a paper like this: \cite{example2024}. All citations will be automatically added to your \texttt{cite.bib} file.

\section{Conclusion}
This template demonstrates basic \LaTeX{} features including mathematical formulas, tables, and lists. The bibliography system is now set up to work seamlessly with Bibby AI's citation management. Feel free to modify it for your needs!

\bibliographystyle{plain}
\bibliography{cite}

\end{document}\documentclass{article}
\usepackage[utf8]{inputenc}
\usepackage{booktabs}
\usepackage{amsmath}
\usepackage{natbib}

\title{My Document}
\author{Your Name}
\date{\today}

\begin{document}

\maketitle

\section{Introduction}
This is a sample document demonstrating various \LaTeX{} features including mathematical equations, tables, and formatting. You can cite papers using the citation search feature, and they will be automatically added to your \texttt{cite.bib} file.

\section{Mathematical Expressions}
Here are some example equations:

\subsection{Quadratic Formula}
The quadratic formula is given by:
\[x = \frac{-b \pm \sqrt{b^2 - 4ac}}{2a}\]

\subsection{Matrix Example}
A sample matrix:
\[
\begin{pmatrix}
a_{11} & a_{12} & a_{13} \\
a_{21} & a_{22} & a_{23} \\
a_{31} & a_{32} & a_{33}
\end{pmatrix}
\]

\section{Tables}
Below is a sample table showing data:

\begin{table}[h]
\centering
\begin{tabular}{@{}lrrr@{}}
\toprule
\textbf{Category} & \textbf{Value 1} & \textbf{Value 2} & \textbf{Total} \\
\midrule
Item A & 10.5 & 20.3 & 30.8 \\
Item B & 15.2 & 25.1 & 40.3 \\
Item C & 22.1 & 12.7 & 34.8 \\
\bottomrule
\end{tabular}
\caption{Sample Data Table}
\label{tab:sample}
\end{table}

\section{Lists}
\subsection{Bullet Points}
\begin{itemize}
  \item First item
  \item Second item
  \item Third item with sub-items:
    \begin{itemize}
      \item Sub-item 1
      \item Sub-item 2
    \end{itemize}
\end{itemize}

\subsection{Numbered List}
\begin{enumerate}
  \item First step
  \item Second step
  \item Third step
\end{enumerate}

\section{Citations and References}
When you use the citation search feature, references will be automatically managed. For example, you might cite a paper like this: \cite{example2024}. All citations will be automatically added to your \texttt{cite.bib} file.

\section{Conclusion}
This template demonstrates basic \LaTeX{} features including mathematical formulas, tables, and lists. The bibliography system is now set up to work seamlessly with Bibby AI's citation management. Feel free to modify it for your needs!

\bibliographystyle{plain}
\bibliography{cite}

\end{document}\documentclass{article}
\usepackage[utf8]{inputenc}
\usepackage{booktabs}
\usepackage{amsmath}
\usepackage{natbib}

\title{My Document}
\author{Your Name}
\date{\today}

\begin{document}

\maketitle

\section{Introduction}
This is a sample document demonstrating various \LaTeX{} features including mathematical equations, tables, and formatting. You can cite papers using the citation search feature, and they will be automatically added to your \texttt{cite.bib} file.

\section{Mathematical Expressions}
Here are some example equations:

\subsection{Quadratic Formula}
The quadratic formula is given by:
\[x = \frac{-b \pm \sqrt{b^2 - 4ac}}{2a}\]

\subsection{Matrix Example}
A sample matrix:
\[
\begin{pmatrix}
a_{11} & a_{12} & a_{13} \\
a_{21} & a_{22} & a_{23} \\
a_{31} & a_{32} & a_{33}
\end{pmatrix}
\]

\section{Tables}
Below is a sample table showing data:

\begin{table}[h]
\centering
\begin{tabular}{@{}lrrr@{}}
\toprule
\textbf{Category} & \textbf{Value 1} & \textbf{Value 2} & \textbf{Total} \\
\midrule
Item A & 10.5 & 20.3 & 30.8 \\
Item B & 15.2 & 25.1 & 40.3 \\
Item C & 22.1 & 12.7 & 34.8 \\
\bottomrule
\end{tabular}
\caption{Sample Data Table}
\label{tab:sample}
\end{table}

\section{Lists}
\subsection{Bullet Points}
\begin{itemize}
  \item First item
  \item Second item
  \item Third item with sub-items:
    \begin{itemize}
      \item Sub-item 1
      \item Sub-item 2
    \end{itemize}
\end{itemize}

\subsection{Numbered List}
\begin{enumerate}
  \item First step
  \item Second step
  \item Third step
\end{enumerate}

\section{Citations and References}
When you use the citation search feature, references will be automatically managed. For example, you might cite a paper like this: \cite{example2024}. All citations will be automatically added to your \texttt{cite.bib} file.

\section{Conclusion}
This template demonstrates basic \LaTeX{} features including mathematical formulas, tables, and lists. The bibliography system is now set up to work seamlessly with Bibby AI's citation management. Feel free to modify it for your needs!

\bibliographystyle{plain}
\bibliography{cite}

\end{document}\documentclass{article}
\usepackage[utf8]{inputenc}
\usepackage{booktabs}
\usepackage{amsmath}
\usepackage{natbib}

\title{My Document}
\author{Your Name}
\date{\today}

\begin{document}

\maketitle

\section{Introduction}
This is a sample document demonstrating various \LaTeX{} features including mathematical equations, tables, and formatting. You can cite papers using the citation search feature, and they will be automatically added to your \texttt{cite.bib} file.

\section{Mathematical Expressions}
Here are some example equations:

\subsection{Quadratic Formula}
The quadratic formula is given by:
\[x = \frac{-b \pm \sqrt{b^2 - 4ac}}{2a}\]

\subsection{Matrix Example}
A sample matrix:
\[
\begin{pmatrix}
a_{11} & a_{12} & a_{13} \\
a_{21} & a_{22} & a_{23} \\
a_{31} & a_{32} & a_{33}
\end{pmatrix}
\]

\section{Tables}
Below is a sample table showing data:

\begin{table}[h]
\centering
\begin{tabular}{@{}lrrr@{}}
\toprule
\textbf{Category} & \textbf{Value 1} & \textbf{Value 2} & \textbf{Total} \\
\midrule
Item A & 10.5 & 20.3 & 30.8 \\
Item B & 15.2 & 25.1 & 40.3 \\
Item C & 22.1 & 12.7 & 34.8 \\
\bottomrule
\end{tabular}
\caption{Sample Data Table}
\label{tab:sample}
\end{table}

\section{Lists}
\subsection{Bullet Points}
\begin{itemize}
  \item First item
  \item Second item
  \item Third item with sub-items:
    \begin{itemize}
      \item Sub-item 1
      \item Sub-item 2
    \end{itemize}
\end{itemize}

\subsection{Numbered List}
\begin{enumerate}
  \item First step
  \item Second step
  \item Third step
\end{enumerate}

\section{Citations and References}
When you use the citation search feature, references will be automatically managed. For example, you might cite a paper like this: \cite{example2024}. All citations will be automatically added to your \texttt{cite.bib} file.

\section{Conclusion}
This template demonstrates basic \LaTeX{} features including mathematical formulas, tables, and lists. The bibliography system is now set up to work seamlessly with Bibby AI's citation management. Feel free to modify it for your needs!

\bibliographystyle{plain}
\bibliography{cite}

\end{document}\documentclass{article}
\usepackage[utf8]{inputenc}
\usepackage{booktabs}
\usepackage{amsmath}
\usepackage{natbib}

\title{My Document}
\author{Your Name}
\date{\today}

\begin{document}

\maketitle

\section{Introduction}
This is a sample document demonstrating various \LaTeX{} features including mathematical equations, tables, and formatting. You can cite papers using the citation search feature, and they will be automatically added to your \texttt{cite.bib} file.

\section{Mathematical Expressions}
Here are some example equations:

\subsection{Quadratic Formula}
The quadratic formula is given by:
\[x = \frac{-b \pm \sqrt{b^2 - 4ac}}{2a}\]

\subsection{Matrix Example}
A sample matrix:
\[
\begin{pmatrix}
a_{11} & a_{12} & a_{13} \\
a_{21} & a_{22} & a_{23} \\
a_{31} & a_{32} & a_{33}
\end{pmatrix}
\]

\section{Tables}
Below is a sample table showing data:

\begin{table}[h]
\centering
\begin{tabular}{@{}lrrr@{}}
\toprule
\textbf{Category} & \textbf{Value 1} & \textbf{Value 2} & \textbf{Total} \\
\midrule
Item A & 10.5 & 20.3 & 30.8 \\
Item B & 15.2 & 25.1 & 40.3 \\
Item C & 22.1 & 12.7 & 34.8 \\
\bottomrule
\end{tabular}
\caption{Sample Data Table}
\label{tab:sample}
\end{table}

\section{Lists}
\subsection{Bullet Points}
\begin{itemize}
  \item First item
  \item Second item
  \item Third item with sub-items:
    \begin{itemize}
      \item Sub-item 1
      \item Sub-item 2
    \end{itemize}
\end{itemize}

\subsection{Numbered List}
\begin{enumerate}
  \item First step
  \item Second step
  \item Third step
\end{enumerate}

\section{Citations and References}
When you use the citation search feature, references will be automatically managed. For example, you might cite a paper like this: \cite{example2024}. All citations will be automatically added to your \texttt{cite.bib} file.

\section{Conclusion}
This template demonstrates basic \LaTeX{} features including mathematical formulas, tables, and lists. The bibliography system is now set up to work seamlessly with Bibby AI's citation management. Feel free to modify it for your needs!

\bibliographystyle{plain}
\bibliography{cite}

\end{document}\documentclass{article}
\usepackage[utf8]{inputenc}
\usepackage{booktabs}
\usepackage{amsmath}
\usepackage{natbib}

\title{My Document}
\author{Your Name}
\date{\today}

\begin{document}

\maketitle

\section{Introduction}
This is a sample document demonstrating various \LaTeX{} features including mathematical equations, tables, and formatting. You can cite papers using the citation search feature, and they will be automatically added to your \texttt{cite.bib} file.

\section{Mathematical Expressions}
Here are some example equations:

\subsection{Quadratic Formula}
The quadratic formula is given by:
\[x = \frac{-b \pm \sqrt{b^2 - 4ac}}{2a}\]

\subsection{Matrix Example}
A sample matrix:
\[
\begin{pmatrix}
a_{11} & a_{12} & a_{13} \\
a_{21} & a_{22} & a_{23} \\
a_{31} & a_{32} & a_{33}
\end{pmatrix}
\]

\section{Tables}
Below is a sample table showing data:

\begin{table}[h]
\centering
\begin{tabular}{@{}lrrr@{}}
\toprule
\textbf{Category} & \textbf{Value 1} & \textbf{Value 2} & \textbf{Total} \\
\midrule
Item A & 10.5 & 20.3 & 30.8 \\
Item B & 15.2 & 25.1 & 40.3 \\
Item C & 22.1 & 12.7 & 34.8 \\
\bottomrule
\end{tabular}
\caption{Sample Data Table}
\label{tab:sample}
\end{table}

\section{Lists}
\subsection{Bullet Points}
\begin{itemize}
  \item First item
  \item Second item
  \item Third item with sub-items:
    \begin{itemize}
      \item Sub-item 1
      \item Sub-item 2
    \end{itemize}
\end{itemize}

\subsection{Numbered List}
\begin{enumerate}
  \item First step
  \item Second step
  \item Third step
\end{enumerate}

\section{Citations and References}
When you use the citation search feature, references will be automatically managed. For example, you might cite a paper like this: \cite{example2024}. All citations will be automatically added to your \texttt{cite.bib} file.

\section{Conclusion}
This template demonstrates basic \LaTeX{} features including mathematical formulas, tables, and lists. The bibliography system is now set up to work seamlessly with Bibby AI's citation management. Feel free to modify it for your needs!

\bibliographystyle{plain}
\bibliography{cite}

\end{document}\documentclass{article}
\usepackage[utf8]{inputenc}
\usepackage{booktabs}
\usepackage{amsmath}
\usepackage{natbib}

\title{My Document}
\author{Your Name}
\date{\today}

\begin{document}

\maketitle

\section{Introduction}
This is a sample document demonstrating various \LaTeX{} features including mathematical equations, tables, and formatting. You can cite papers using the citation search feature, and they will be automatically added to your \texttt{cite.bib} file.

\section{Mathematical Expressions}
Here are some example equations:

\subsection{Quadratic Formula}
The quadratic formula is given by:
\[x = \frac{-b \pm \sqrt{b^2 - 4ac}}{2a}\]

\subsection{Matrix Example}
A sample matrix:
\[
\begin{pmatrix}
a_{11} & a_{12} & a_{13} \\
a_{21} & a_{22} & a_{23} \\
a_{31} & a_{32} & a_{33}
\end{pmatrix}
\]

\section{Tables}
Below is a sample table showing data:

\begin{table}[h]
\centering
\begin{tabular}{@{}lrrr@{}}
\toprule
\textbf{Category} & \textbf{Value 1} & \textbf{Value 2} & \textbf{Total} \\
\midrule
Item A & 10.5 & 20.3 & 30.8 \\
Item B & 15.2 & 25.1 & 40.3 \\
Item C & 22.1 & 12.7 & 34.8 \\
\bottomrule
\end{tabular}
\caption{Sample Data Table}
\label{tab:sample}
\end{table}

\section{Lists}
\subsection{Bullet Points}
\begin{itemize}
  \item First item
  \item Second item
  \item Third item with sub-items:
    \begin{itemize}
      \item Sub-item 1
      \item Sub-item 2
    \end{itemize}
\end{itemize}

\subsection{Numbered List}
\begin{enumerate}
  \item First step
  \item Second step
  \item Third step
\end{enumerate}

\section{Citations and References}
When you use the citation search feature, references will be automatically managed. For example, you might cite a paper like this: \cite{example2024}. All citations will be automatically added to your \texttt{cite.bib} file.

\section{Conclusion}
This template demonstrates basic \LaTeX{} features including mathematical formulas, tables, and lists. The bibliography system is now set up to work seamlessly with Bibby AI's citation management. Feel free to modify it for your needs!

\bibliographystyle{plain}
\bibliography{cite}

\end{document}\documentclass{article}
\usepackage[utf8]{inputenc}
\usepackage{booktabs}
\usepackage{amsmath}
\usepackage{natbib}

\title{My Document}
\author{Your Name}
\date{\today}

\begin{document}

\maketitle

\section{Introduction}
This is a sample document demonstrating various \LaTeX{} features including mathematical equations, tables, and formatting. You can cite papers using the citation search feature, and they will be automatically added to your \texttt{cite.bib} file.

\section{Mathematical Expressions}
Here are some example equations:

\subsection{Quadratic Formula}
The quadratic formula is given by:
\[x = \frac{-b \pm \sqrt{b^2 - 4ac}}{2a}\]

\subsection{Matrix Example}
A sample matrix:
\[
\begin{pmatrix}
a_{11} & a_{12} & a_{13} \\
a_{21} & a_{22} & a_{23} \\
a_{31} & a_{32} & a_{33}
\end{pmatrix}
\]

\section{Tables}
Below is a sample table showing data:

\begin{table}[h]
\centering
\begin{tabular}{@{}lrrr@{}}
\toprule
\textbf{Category} & \textbf{Value 1} & \textbf{Value 2} & \textbf{Total} \\
\midrule
Item A & 10.5 & 20.3 & 30.8 \\
Item B & 15.2 & 25.1 & 40.3 \\
Item C & 22.1 & 12.7 & 34.8 \\
\bottomrule
\end{tabular}
\caption{Sample Data Table}
\label{tab:sample}
\end{table}

\section{Lists}
\subsection{Bullet Points}
\begin{itemize}
  \item First item
  \item Second item
  \item Third item with sub-items:
    \begin{itemize}
      \item Sub-item 1
      \item Sub-item 2
    \end{itemize}
\end{itemize}

\subsection{Numbered List}
\begin{enumerate}
  \item First step
  \item Second step
  \item Third step
\end{enumerate}

\section{Citations and References}
When you use the citation search feature, references will be automatically managed. For example, you might cite a paper like this: \cite{example2024}. All citations will be automatically added to your \texttt{cite.bib} file.

\section{Conclusion}
This template demonstrates basic \LaTeX{} features including mathematical formulas, tables, and lists. The bibliography system is now set up to work seamlessly with Bibby AI's citation management. Feel free to modify it for your needs!

\bibliographystyle{plain}
\bibliography{cite}

\end{document}\documentclass{article}
\usepackage[utf8]{inputenc}
\usepackage{booktabs}
\usepackage{amsmath}
\usepackage{natbib}

\title{My Document}
\author{Your Name}
\date{\today}

\begin{document}

\maketitle

\section{Introduction}
This is a sample document demonstrating various \LaTeX{} features including mathematical equations, tables, and formatting. You can cite papers using the citation search feature, and they will be automatically added to your \texttt{cite.bib} file.

\section{Mathematical Expressions}
Here are some example equations:

\subsection{Quadratic Formula}
The quadratic formula is given by:
\[x = \frac{-b \pm \sqrt{b^2 - 4ac}}{2a}\]

\subsection{Matrix Example}
A sample matrix:
\[
\begin{pmatrix}
a_{11} & a_{12} & a_{13} \\
a_{21} & a_{22} & a_{23} \\
a_{31} & a_{32} & a_{33}
\end{pmatrix}
\]

\section{Tables}
Below is a sample table showing data:

\begin{table}[h]
\centering
\begin{tabular}{@{}lrrr@{}}
\toprule
\textbf{Category} & \textbf{Value 1} & \textbf{Value 2} & \textbf{Total} \\
\midrule
Item A & 10.5 & 20.3 & 30.8 \\
Item B & 15.2 & 25.1 & 40.3 \\
Item C & 22.1 & 12.7 & 34.8 \\
\bottomrule
\end{tabular}
\caption{Sample Data Table}
\label{tab:sample}
\end{table}

\section{Lists}
\subsection{Bullet Points}
\begin{itemize}
  \item First item
  \item Second item
  \item Third item with sub-items:
    \begin{itemize}
      \item Sub-item 1
      \item Sub-item 2
    \end{itemize}
\end{itemize}

\subsection{Numbered List}
\begin{enumerate}
  \item First step
  \item Second step
  \item Third step
\end{enumerate}

\section{Citations and References}
When you use the citation search feature, references will be automatically managed. For example, you might cite a paper like this: \cite{example2024}. All citations will be automatically added to your \texttt{cite.bib} file.

\section{Conclusion}
This template demonstrates basic \LaTeX{} features including mathematical formulas, tables, and lists. The bibliography system is now set up to work seamlessly with Bibby AI's citation management. Feel free to modify it for your needs!

\bibliographystyle{plain}
\bibliography{cite}

\end{document}\documentclass{article}
\usepackage[utf8]{inputenc}
\usepackage{booktabs}
\usepackage{amsmath}
\usepackage{natbib}

\title{My Document}
\author{Your Name}
\date{\today}

\begin{document}

\maketitle

\section{Introduction}
This is a sample document demonstrating various \LaTeX{} features including mathematical equations, tables, and formatting. You can cite papers using the citation search feature, and they will be automatically added to your \texttt{cite.bib} file.

\section{Mathematical Expressions}
Here are some example equations:

\subsection{Quadratic Formula}
The quadratic formula is given by:
\[x = \frac{-b \pm \sqrt{b^2 - 4ac}}{2a}\]

\subsection{Matrix Example}
A sample matrix:
\[
\begin{pmatrix}
a_{11} & a_{12} & a_{13} \\
a_{21} & a_{22} & a_{23} \\
a_{31} & a_{32} & a_{33}
\end{pmatrix}
\]

\section{Tables}
Below is a sample table showing data:

\begin{table}[h]
\centering
\begin{tabular}{@{}lrrr@{}}
\toprule
\textbf{Category} & \textbf{Value 1} & \textbf{Value 2} & \textbf{Total} \\
\midrule
Item A & 10.5 & 20.3 & 30.8 \\
Item B & 15.2 & 25.1 & 40.3 \\
Item C & 22.1 & 12.7 & 34.8 \\
\bottomrule
\end{tabular}
\caption{Sample Data Table}
\label{tab:sample}
\end{table}

\section{Lists}
\subsection{Bullet Points}
\begin{itemize}
  \item First item
  \item Second item
  \item Third item with sub-items:
    \begin{itemize}
      \item Sub-item 1
      \item Sub-item 2
    \end{itemize}
\end{itemize}

\subsection{Numbered List}
\begin{enumerate}
  \item First step
  \item Second step
  \item Third step
\end{enumerate}

\section{Citations and References}
When you use the citation search feature, references will be automatically managed. For example, you might cite a paper like this: \cite{example2024}. All citations will be automatically added to your \texttt{cite.bib} file.

\section{Conclusion}
This template demonstrates basic \LaTeX{} features including mathematical formulas, tables, and lists. The bibliography system is now set up to work seamlessly with Bibby AI's citation management. Feel free to modify it for your needs!

\bibliographystyle{plain}
\bibliography{cite}

\end{document}\documentclass{article}
\usepackage[utf8]{inputenc}
\usepackage{booktabs}
\usepackage{amsmath}
\usepackage{natbib}

\title{My Document}
\author{Your Name}
\date{\today}

\begin{document}

\maketitle

\section{Introduction}
This is a sample document demonstrating various \LaTeX{} features including mathematical equations, tables, and formatting. You can cite papers using the citation search feature, and they will be automatically added to your \texttt{cite.bib} file.

\section{Mathematical Expressions}
Here are some example equations:

\subsection{Quadratic Formula}
The quadratic formula is given by:
\[x = \frac{-b \pm \sqrt{b^2 - 4ac}}{2a}\]

\subsection{Matrix Example}
A sample matrix:
\[
\begin{pmatrix}
a_{11} & a_{12} & a_{13} \\
a_{21} & a_{22} & a_{23} \\
a_{31} & a_{32} & a_{33}
\end{pmatrix}
\]

\section{Tables}
Below is a sample table showing data:

\begin{table}[h]
\centering
\begin{tabular}{@{}lrrr@{}}
\toprule
\textbf{Category} & \textbf{Value 1} & \textbf{Value 2} & \textbf{Total} \\
\midrule
Item A & 10.5 & 20.3 & 30.8 \\
Item B & 15.2 & 25.1 & 40.3 \\
Item C & 22.1 & 12.7 & 34.8 \\
\bottomrule
\end{tabular}
\caption{Sample Data Table}
\label{tab:sample}
\end{table}

\section{Lists}
\subsection{Bullet Points}
\begin{itemize}
  \item First item
  \item Second item
  \item Third item with sub-items:
    \begin{itemize}
      \item Sub-item 1
      \item Sub-item 2
    \end{itemize}
\end{itemize}

\subsection{Numbered List}
\begin{enumerate}
  \item First step
  \item Second step
  \item Third step
\end{enumerate}

\section{Citations and References}
When you use the citation search feature, references will be automatically managed. For example, you might cite a paper like this: \cite{example2024}. All citations will be automatically added to your \texttt{cite.bib} file.

\section{Conclusion}
This template demonstrates basic \LaTeX{} features including mathematical formulas, tables, and lists. The bibliography system is now set up to work seamlessly with Bibby AI's citation management. Feel free to modify it for your needs!

\bibliographystyle{plain}
\bibliography{cite}

\end{document}\documentclass{article}
\usepackage[utf8]{inputenc}
\usepackage{booktabs}
\usepackage{amsmath}
\usepackage{natbib}

\title{My Document}
\author{Your Name}
\date{\today}

\begin{document}

\maketitle

\section{Introduction}
This is a sample document demonstrating various \LaTeX{} features including mathematical equations, tables, and formatting. You can cite papers using the citation search feature, and they will be automatically added to your \texttt{cite.bib} file.

\section{Mathematical Expressions}
Here are some example equations:

\subsection{Quadratic Formula}
The quadratic formula is given by:
\[x = \frac{-b \pm \sqrt{b^2 - 4ac}}{2a}\]

\subsection{Matrix Example}
A sample matrix:
\[
\begin{pmatrix}
a_{11} & a_{12} & a_{13} \\
a_{21} & a_{22} & a_{23} \\
a_{31} & a_{32} & a_{33}
\end{pmatrix}
\]

\section{Tables}
Below is a sample table showing data:

\begin{table}[h]
\centering
\begin{tabular}{@{}lrrr@{}}
\toprule
\textbf{Category} & \textbf{Value 1} & \textbf{Value 2} & \textbf{Total} \\
\midrule
Item A & 10.5 & 20.3 & 30.8 \\
Item B & 15.2 & 25.1 & 40.3 \\
Item C & 22.1 & 12.7 & 34.8 \\
\bottomrule
\end{tabular}
\caption{Sample Data Table}
\label{tab:sample}
\end{table}

\section{Lists}
\subsection{Bullet Points}
\begin{itemize}
  \item First item
  \item Second item
  \item Third item with sub-items:
    \begin{itemize}
      \item Sub-item 1
      \item Sub-item 2
    \end{itemize}
\end{itemize}

\subsection{Numbered List}
\begin{enumerate}
  \item First step
  \item Second step
  \item Third step
\end{enumerate}

\section{Citations and References}
When you use the citation search feature, references will be automatically managed. For example, you might cite a paper like this: \cite{example2024}. All citations will be automatically added to your \texttt{cite.bib} file.

\section{Conclusion}
This template demonstrates basic \LaTeX{} features including mathematical formulas, tables, and lists. The bibliography system is now set up to work seamlessly with Bibby AI's citation management. Feel free to modify it for your needs!

\bibliographystyle{plain}
\bibliography{cite}

\end{document}\documentclass{article}
\usepackage[utf8]{inputenc}
\usepackage{booktabs}
\usepackage{amsmath}
\usepackage{natbib}

\title{My Document}
\author{Your Name}
\date{\today}

\begin{document}

\maketitle

\section{Introduction}
This is a sample document demonstrating various \LaTeX{} features including mathematical equations, tables, and formatting. You can cite papers using the citation search feature, and they will be automatically added to your \texttt{cite.bib} file.

\section{Mathematical Expressions}
Here are some example equations:

\subsection{Quadratic Formula}
The quadratic formula is given by:
\[x = \frac{-b \pm \sqrt{b^2 - 4ac}}{2a}\]

\subsection{Matrix Example}
A sample matrix:
\[
\begin{pmatrix}
a_{11} & a_{12} & a_{13} \\
a_{21} & a_{22} & a_{23} \\
a_{31} & a_{32} & a_{33}
\end{pmatrix}
\]

\section{Tables}
Below is a sample table showing data:

\begin{table}[h]
\centering
\begin{tabular}{@{}lrrr@{}}
\toprule
\textbf{Category} & \textbf{Value 1} & \textbf{Value 2} & \textbf{Total} \\
\midrule
Item A & 10.5 & 20.3 & 30.8 \\
Item B & 15.2 & 25.1 & 40.3 \\
Item C & 22.1 & 12.7 & 34.8 \\
\bottomrule
\end{tabular}
\caption{Sample Data Table}
\label{tab:sample}
\end{table}

\section{Lists}
\subsection{Bullet Points}
\begin{itemize}
  \item First item
  \item Second item
  \item Third item with sub-items:
    \begin{itemize}
      \item Sub-item 1
      \item Sub-item 2
    \end{itemize}
\end{itemize}

\subsection{Numbered List}
\begin{enumerate}
  \item First step
  \item Second step
  \item Third step
\end{enumerate}

\section{Citations and References}
When you use the citation search feature, references will be automatically managed. For example, you might cite a paper like this: \cite{example2024}. All citations will be automatically added to your \texttt{cite.bib} file.

\section{Conclusion}
This template demonstrates basic \LaTeX{} features including mathematical formulas, tables, and lists. The bibliography system is now set up to work seamlessly with Bibby AI's citation management. Feel free to modify it for your needs!

\bibliographystyle{plain}
\bibliography{cite}

\end{document}\documentclass{article}
\usepackage[utf8]{inputenc}
\usepackage{booktabs}
\usepackage{amsmath}
\usepackage{natbib}

\title{My Document}
\author{Your Name}
\date{\today}

\begin{document}

\maketitle

\section{Introduction}
This is a sample document demonstrating various \LaTeX{} features including mathematical equations, tables, and formatting. You can cite papers using the citation search feature, and they will be automatically added to your \texttt{cite.bib} file.

\section{Mathematical Expressions}
Here are some example equations:

\subsection{Quadratic Formula}
The quadratic formula is given by:
\[x = \frac{-b \pm \sqrt{b^2 - 4ac}}{2a}\]

\subsection{Matrix Example}
A sample matrix:
\[
\begin{pmatrix}
a_{11} & a_{12} & a_{13} \\
a_{21} & a_{22} & a_{23} \\
a_{31} & a_{32} & a_{33}
\end{pmatrix}
\]

\section{Tables}
Below is a sample table showing data:

\begin{table}[h]
\centering
\begin{tabular}{@{}lrrr@{}}
\toprule
\textbf{Category} & \textbf{Value 1} & \textbf{Value 2} & \textbf{Total} \\
\midrule
Item A & 10.5 & 20.3 & 30.8 \\
Item B & 15.2 & 25.1 & 40.3 \\
Item C & 22.1 & 12.7 & 34.8 \\
\bottomrule
\end{tabular}
\caption{Sample Data Table}
\label{tab:sample}
\end{table}

\section{Lists}
\subsection{Bullet Points}
\begin{itemize}
  \item First item
  \item Second item
  \item Third item with sub-items:
    \begin{itemize}
      \item Sub-item 1
      \item Sub-item 2
    \end{itemize}
\end{itemize}

\subsection{Numbered List}
\begin{enumerate}
  \item First step
  \item Second step
  \item Third step
\end{enumerate}

\section{Citations and References}
When you use the citation search feature, references will be automatically managed. For example, you might cite a paper like this: \cite{example2024}. All citations will be automatically added to your \texttt{cite.bib} file.

\section{Conclusion}
This template demonstrates basic \LaTeX{} features including mathematical formulas, tables, and lists. The bibliography system is now set up to work seamlessly with Bibby AI's citation management. Feel free to modify it for your needs!

\bibliographystyle{plain}
\bibliography{cite}

\end{document}\documentclass{article}
\usepackage[utf8]{inputenc}
\usepackage{booktabs}
\usepackage{amsmath}
\usepackage{natbib}

\title{My Document}
\author{Your Name}
\date{\today}

\begin{document}

\maketitle

\section{Introduction}
This is a sample document demonstrating various \LaTeX{} features including mathematical equations, tables, and formatting. You can cite papers using the citation search feature, and they will be automatically added to your \texttt{cite.bib} file.

\section{Mathematical Expressions}
Here are some example equations:

\subsection{Quadratic Formula}
The quadratic formula is given by:
\[x = \frac{-b \pm \sqrt{b^2 - 4ac}}{2a}\]

\subsection{Matrix Example}
A sample matrix:
\[
\begin{pmatrix}
a_{11} & a_{12} & a_{13} \\
a_{21} & a_{22} & a_{23} \\
a_{31} & a_{32} & a_{33}
\end{pmatrix}
\]

\section{Tables}
Below is a sample table showing data:

\begin{table}[h]
\centering
\begin{tabular}{@{}lrrr@{}}
\toprule
\textbf{Category} & \textbf{Value 1} & \textbf{Value 2} & \textbf{Total} \\
\midrule
Item A & 10.5 & 20.3 & 30.8 \\
Item B & 15.2 & 25.1 & 40.3 \\
Item C & 22.1 & 12.7 & 34.8 \\
\bottomrule
\end{tabular}
\caption{Sample Data Table}
\label{tab:sample}
\end{table}

\section{Lists}
\subsection{Bullet Points}
\begin{itemize}
  \item First item
  \item Second item
  \item Third item with sub-items:
    \begin{itemize}
      \item Sub-item 1
      \item Sub-item 2
    \end{itemize}
\end{itemize}

\subsection{Numbered List}
\begin{enumerate}
  \item First step
  \item Second step
  \item Third step
\end{enumerate}

\section{Citations and References}
When you use the citation search feature, references will be automatically managed. For example, you might cite a paper like this: \cite{example2024}. All citations will be automatically added to your \texttt{cite.bib} file.

\section{Conclusion}
This template demonstrates basic \LaTeX{} features including mathematical formulas, tables, and lists. The bibliography system is now set up to work seamlessly with Bibby AI's citation management. Feel free to modify it for your needs!

\bibliographystyle{plain}
\bibliography{cite}

\end{document}\documentclass{article}
\usepackage[utf8]{inputenc}
\usepackage{booktabs}
\usepackage{amsmath}
\usepackage{natbib}

\title{My Document}
\author{Your Name}
\date{\today}

\begin{document}

\maketitle

\section{Introduction}
This is a sample document demonstrating various \LaTeX{} features including mathematical equations, tables, and formatting. You can cite papers using the citation search feature, and they will be automatically added to your \texttt{cite.bib} file.

\section{Mathematical Expressions}
Here are some example equations:

\subsection{Quadratic Formula}
The quadratic formula is given by:
\[x = \frac{-b \pm \sqrt{b^2 - 4ac}}{2a}\]

\subsection{Matrix Example}
A sample matrix:
\[
\begin{pmatrix}
a_{11} & a_{12} & a_{13} \\
a_{21} & a_{22} & a_{23} \\
a_{31} & a_{32} & a_{33}
\end{pmatrix}
\]

\section{Tables}
Below is a sample table showing data:

\begin{table}[h]
\centering
\begin{tabular}{@{}lrrr@{}}
\toprule
\textbf{Category} & \textbf{Value 1} & \textbf{Value 2} & \textbf{Total} \\
\midrule
Item A & 10.5 & 20.3 & 30.8 \\
Item B & 15.2 & 25.1 & 40.3 \\
Item C & 22.1 & 12.7 & 34.8 \\
\bottomrule
\end{tabular}
\caption{Sample Data Table}
\label{tab:sample}
\end{table}

\section{Lists}
\subsection{Bullet Points}
\begin{itemize}
  \item First item
  \item Second item
  \item Third item with sub-items:
    \begin{itemize}
      \item Sub-item 1
      \item Sub-item 2
    \end{itemize}
\end{itemize}

\subsection{Numbered List}
\begin{enumerate}
  \item First step
  \item Second step
  \item Third step
\end{enumerate}

\section{Citations and References}
When you use the citation search feature, references will be automatically managed. For example, you might cite a paper like this: \cite{example2024}. All citations will be automatically added to your \texttt{cite.bib} file.

\section{Conclusion}
This template demonstrates basic \LaTeX{} features including mathematical formulas, tables, and lists. The bibliography system is now set up to work seamlessly with Bibby AI's citation management. Feel free to modify it for your needs!

\bibliographystyle{plain}
\bibliography{cite}

\end{document}\documentclass{article}
\usepackage[utf8]{inputenc}
\usepackage{booktabs}
\usepackage{amsmath}
\usepackage{natbib}

\title{My Document}
\author{Your Name}
\date{\today}

\begin{document}

\maketitle

\section{Introduction}
This is a sample document demonstrating various \LaTeX{} features including mathematical equations, tables, and formatting. You can cite papers using the citation search feature, and they will be automatically added to your \texttt{cite.bib} file.

\section{Mathematical Expressions}
Here are some example equations:

\subsection{Quadratic Formula}
The quadratic formula is given by:
\[x = \frac{-b \pm \sqrt{b^2 - 4ac}}{2a}\]

\subsection{Matrix Example}
A sample matrix:
\[
\begin{pmatrix}
a_{11} & a_{12} & a_{13} \\
a_{21} & a_{22} & a_{23} \\
a_{31} & a_{32} & a_{33}
\end{pmatrix}
\]

\section{Tables}
Below is a sample table showing data:

\begin{table}[h]
\centering
\begin{tabular}{@{}lrrr@{}}
\toprule
\textbf{Category} & \textbf{Value 1} & \textbf{Value 2} & \textbf{Total} \\
\midrule
Item A & 10.5 & 20.3 & 30.8 \\
Item B & 15.2 & 25.1 & 40.3 \\
Item C & 22.1 & 12.7 & 34.8 \\
\bottomrule
\end{tabular}
\caption{Sample Data Table}
\label{tab:sample}
\end{table}

\section{Lists}
\subsection{Bullet Points}
\begin{itemize}
  \item First item
  \item Second item
  \item Third item with sub-items:
    \begin{itemize}
      \item Sub-item 1
      \item Sub-item 2
    \end{itemize}
\end{itemize}

\subsection{Numbered List}
\begin{enumerate}
  \item First step
  \item Second step
  \item Third step
\end{enumerate}

\section{Citations and References}
When you use the citation search feature, references will be automatically managed. For example, you might cite a paper like this: \cite{example2024}. All citations will be automatically added to your \texttt{cite.bib} file.

\section{Conclusion}
This template demonstrates basic \LaTeX{} features including mathematical formulas, tables, and lists. The bibliography system is now set up to work seamlessly with Bibby AI's citation management. Feel free to modify it for your needs!

\bibliographystyle{plain}
\bibliography{cite}

\end{document}\documentclass{article}
\usepackage[utf8]{inputenc}
\usepackage{booktabs}
\usepackage{amsmath}
\usepackage{natbib}

\title{My Document}
\author{Your Name}
\date{\today}

\begin{document}

\maketitle

\section{Introduction}
This is a sample document demonstrating various \LaTeX{} features including mathematical equations, tables, and formatting. You can cite papers using the citation search feature, and they will be automatically added to your \texttt{cite.bib} file.

\section{Mathematical Expressions}
Here are some example equations:

\subsection{Quadratic Formula}
The quadratic formula is given by:
\[x = \frac{-b \pm \sqrt{b^2 - 4ac}}{2a}\]

\subsection{Matrix Example}
A sample matrix:
\[
\begin{pmatrix}
a_{11} & a_{12} & a_{13} \\
a_{21} & a_{22} & a_{23} \\
a_{31} & a_{32} & a_{33}
\end{pmatrix}
\]

\section{Tables}
Below is a sample table showing data:

\begin{table}[h]
\centering
\begin{tabular}{@{}lrrr@{}}
\toprule
\textbf{Category} & \textbf{Value 1} & \textbf{Value 2} & \textbf{Total} \\
\midrule
Item A & 10.5 & 20.3 & 30.8 \\
Item B & 15.2 & 25.1 & 40.3 \\
Item C & 22.1 & 12.7 & 34.8 \\
\bottomrule
\end{tabular}
\caption{Sample Data Table}
\label{tab:sample}
\end{table}

\section{Lists}
\subsection{Bullet Points}
\begin{itemize}
  \item First item
  \item Second item
  \item Third item with sub-items:
    \begin{itemize}
      \item Sub-item 1
      \item Sub-item 2
    \end{itemize}
\end{itemize}

\subsection{Numbered List}
\begin{enumerate}
  \item First step
  \item Second step
  \item Third step
\end{enumerate}

\section{Citations and References}
When you use the citation search feature, references will be automatically managed. For example, you might cite a paper like this: \cite{example2024}. All citations will be automatically added to your \texttt{cite.bib} file.

\section{Conclusion}
This template demonstrates basic \LaTeX{} features including mathematical formulas, tables, and lists. The bibliography system is now set up to work seamlessly with Bibby AI's citation management. Feel free to modify it for your needs!

\bibliographystyle{plain}
\bibliography{cite}

\end{document}\documentclass{article}
\usepackage[utf8]{inputenc}
\usepackage{booktabs}
\usepackage{amsmath}
\usepackage{natbib}

\title{My Document}
\author{Your Name}
\date{\today}

\begin{document}

\maketitle

\section{Introduction}
This is a sample document demonstrating various \LaTeX{} features including mathematical equations, tables, and formatting. You can cite papers using the citation search feature, and they will be automatically added to your \texttt{cite.bib} file.

\section{Mathematical Expressions}
Here are some example equations:

\subsection{Quadratic Formula}
The quadratic formula is given by:
\[x = \frac{-b \pm \sqrt{b^2 - 4ac}}{2a}\]

\subsection{Matrix Example}
A sample matrix:
\[
\begin{pmatrix}
a_{11} & a_{12} & a_{13} \\
a_{21} & a_{22} & a_{23} \\
a_{31} & a_{32} & a_{33}
\end{pmatrix}
\]

\section{Tables}
Below is a sample table showing data:

\begin{table}[h]
\centering
\begin{tabular}{@{}lrrr@{}}
\toprule
\textbf{Category} & \textbf{Value 1} & \textbf{Value 2} & \textbf{Total} \\
\midrule
Item A & 10.5 & 20.3 & 30.8 \\
Item B & 15.2 & 25.1 & 40.3 \\
Item C & 22.1 & 12.7 & 34.8 \\
\bottomrule
\end{tabular}
\caption{Sample Data Table}
\label{tab:sample}
\end{table}

\section{Lists}
\subsection{Bullet Points}
\begin{itemize}
  \item First item
  \item Second item
  \item Third item with sub-items:
    \begin{itemize}
      \item Sub-item 1
      \item Sub-item 2
    \end{itemize}
\end{itemize}

\subsection{Numbered List}
\begin{enumerate}
  \item First step
  \item Second step
  \item Third step
\end{enumerate}

\section{Citations and References}
When you use the citation search feature, references will be automatically managed. For example, you might cite a paper like this: \cite{example2024}. All citations will be automatically added to your \texttt{cite.bib} file.

\section{Conclusion}
This template demonstrates basic \LaTeX{} features including mathematical formulas, tables, and lists. The bibliography system is now set up to work seamlessly with Bibby AI's citation management. Feel free to modify it for your needs!

\bibliographystyle{plain}
\bibliography{cite}

\end{document}\documentclass{article}
\usepackage[utf8]{inputenc}
\usepackage{booktabs}
\usepackage{amsmath}
\usepackage{natbib}

\title{My Document}
\author{Your Name}
\date{\today}

\begin{document}

\maketitle

\section{Introduction}
This is a sample document demonstrating various \LaTeX{} features including mathematical equations, tables, and formatting. You can cite papers using the citation search feature, and they will be automatically added to your \texttt{cite.bib} file.

\section{Mathematical Expressions}
Here are some example equations:

\subsection{Quadratic Formula}
The quadratic formula is given by:
\[x = \frac{-b \pm \sqrt{b^2 - 4ac}}{2a}\]

\subsection{Matrix Example}
A sample matrix:
\[
\begin{pmatrix}
a_{11} & a_{12} & a_{13} \\
a_{21} & a_{22} & a_{23} \\
a_{31} & a_{32} & a_{33}
\end{pmatrix}
\]

\section{Tables}
Below is a sample table showing data:

\begin{table}[h]
\centering
\begin{tabular}{@{}lrrr@{}}
\toprule
\textbf{Category} & \textbf{Value 1} & \textbf{Value 2} & \textbf{Total} \\
\midrule
Item A & 10.5 & 20.3 & 30.8 \\
Item B & 15.2 & 25.1 & 40.3 \\
Item C & 22.1 & 12.7 & 34.8 \\
\bottomrule
\end{tabular}
\caption{Sample Data Table}
\label{tab:sample}
\end{table}

\section{Lists}
\subsection{Bullet Points}
\begin{itemize}
  \item First item
  \item Second item
  \item Third item with sub-items:
    \begin{itemize}
      \item Sub-item 1
      \item Sub-item 2
    \end{itemize}
\end{itemize}

\subsection{Numbered List}
\begin{enumerate}
  \item First step
  \item Second step
  \item Third step
\end{enumerate}

\section{Citations and References}
When you use the citation search feature, references will be automatically managed. For example, you might cite a paper like this: \cite{example2024}. All citations will be automatically added to your \texttt{cite.bib} file.

\section{Conclusion}
This template demonstrates basic \LaTeX{} features including mathematical formulas, tables, and lists. The bibliography system is now set up to work seamlessly with Bibby AI's citation management. Feel free to modify it for your needs!

\bibliographystyle{plain}
\bibliography{cite}

\end{document}\documentclass{article}
\usepackage[utf8]{inputenc}
\usepackage{booktabs}
\usepackage{amsmath}
\usepackage{natbib}

\title{My Document}
\author{Your Name}
\date{\today}

\begin{document}

\maketitle

\section{Introduction}
This is a sample document demonstrating various \LaTeX{} features including mathematical equations, tables, and formatting. You can cite papers using the citation search feature, and they will be automatically added to your \texttt{cite.bib} file.

\section{Mathematical Expressions}
Here are some example equations:

\subsection{Quadratic Formula}
The quadratic formula is given by:
\[x = \frac{-b \pm \sqrt{b^2 - 4ac}}{2a}\]

\subsection{Matrix Example}
A sample matrix:
\[
\begin{pmatrix}
a_{11} & a_{12} & a_{13} \\
a_{21} & a_{22} & a_{23} \\
a_{31} & a_{32} & a_{33}
\end{pmatrix}
\]

\section{Tables}
Below is a sample table showing data:

\begin{table}[h]
\centering
\begin{tabular}{@{}lrrr@{}}
\toprule
\textbf{Category} & \textbf{Value 1} & \textbf{Value 2} & \textbf{Total} \\
\midrule
Item A & 10.5 & 20.3 & 30.8 \\
Item B & 15.2 & 25.1 & 40.3 \\
Item C & 22.1 & 12.7 & 34.8 \\
\bottomrule
\end{tabular}
\caption{Sample Data Table}
\label{tab:sample}
\end{table}

\section{Lists}
\subsection{Bullet Points}
\begin{itemize}
  \item First item
  \item Second item
  \item Third item with sub-items:
    \begin{itemize}
      \item Sub-item 1
      \item Sub-item 2
    \end{itemize}
\end{itemize}

\subsection{Numbered List}
\begin{enumerate}
  \item First step
  \item Second step
  \item Third step
\end{enumerate}

\section{Citations and References}
When you use the citation search feature, references will be automatically managed. For example, you might cite a paper like this: \cite{example2024}. All citations will be automatically added to your \texttt{cite.bib} file.

\section{Conclusion}
This template demonstrates basic \LaTeX{} features including mathematical formulas, tables, and lists. The bibliography system is now set up to work seamlessly with Bibby AI's citation management. Feel free to modify it for your needs!

\bibliographystyle{plain}
\bibliography{cite}

\end{document}\documentclass{article}
\usepackage[utf8]{inputenc}
\usepackage{booktabs}
\usepackage{amsmath}
\usepackage{natbib}

\title{My Document}
\author{Your Name}
\date{\today}

\begin{document}

\maketitle

\section{Introduction}
This is a sample document demonstrating various \LaTeX{} features including mathematical equations, tables, and formatting. You can cite papers using the citation search feature, and they will be automatically added to your \texttt{cite.bib} file.

\section{Mathematical Expressions}
Here are some example equations:

\subsection{Quadratic Formula}
The quadratic formula is given by:
\[x = \frac{-b \pm \sqrt{b^2 - 4ac}}{2a}\]

\subsection{Matrix Example}
A sample matrix:
\[
\begin{pmatrix}
a_{11} & a_{12} & a_{13} \\
a_{21} & a_{22} & a_{23} \\
a_{31} & a_{32} & a_{33}
\end{pmatrix}
\]

\section{Tables}
Below is a sample table showing data:

\begin{table}[h]
\centering
\begin{tabular}{@{}lrrr@{}}
\toprule
\textbf{Category} & \textbf{Value 1} & \textbf{Value 2} & \textbf{Total} \\
\midrule
Item A & 10.5 & 20.3 & 30.8 \\
Item B & 15.2 & 25.1 & 40.3 \\
Item C & 22.1 & 12.7 & 34.8 \\
\bottomrule
\end{tabular}
\caption{Sample Data Table}
\label{tab:sample}
\end{table}

\section{Lists}
\subsection{Bullet Points}
\begin{itemize}
  \item First item
  \item Second item
  \item Third item with sub-items:
    \begin{itemize}
      \item Sub-item 1
      \item Sub-item 2
    \end{itemize}
\end{itemize}

\subsection{Numbered List}
\begin{enumerate}
  \item First step
  \item Second step
  \item Third step
\end{enumerate}

\section{Citations and References}
When you use the citation search feature, references will be automatically managed. For example, you might cite a paper like this: \cite{example2024}. All citations will be automatically added to your \texttt{cite.bib} file.

\section{Conclusion}
This template demonstrates basic \LaTeX{} features including mathematical formulas, tables, and lists. The bibliography system is now set up to work seamlessly with Bibby AI's citation management. Feel free to modify it for your needs!

\bibliographystyle{plain}
\bibliography{cite}

\end{document}\documentclass{article}
\usepackage[utf8]{inputenc}
\usepackage{booktabs}
\usepackage{amsmath}
\usepackage{natbib}

\title{My Document}
\author{Your Name}
\date{\today}

\begin{document}

\maketitle

\section{Introduction}
This is a sample document demonstrating various \LaTeX{} features including mathematical equations, tables, and formatting. You can cite papers using the citation search feature, and they will be automatically added to your \texttt{cite.bib} file.

\section{Mathematical Expressions}
Here are some example equations:

\subsection{Quadratic Formula}
The quadratic formula is given by:
\[x = \frac{-b \pm \sqrt{b^2 - 4ac}}{2a}\]

\subsection{Matrix Example}
A sample matrix:
\[
\begin{pmatrix}
a_{11} & a_{12} & a_{13} \\
a_{21} & a_{22} & a_{23} \\
a_{31} & a_{32} & a_{33}
\end{pmatrix}
\]

\section{Tables}
Below is a sample table showing data:

\begin{table}[h]
\centering
\begin{tabular}{@{}lrrr@{}}
\toprule
\textbf{Category} & \textbf{Value 1} & \textbf{Value 2} & \textbf{Total} \\
\midrule
Item A & 10.5 & 20.3 & 30.8 \\
Item B & 15.2 & 25.1 & 40.3 \\
Item C & 22.1 & 12.7 & 34.8 \\
\bottomrule
\end{tabular}
\caption{Sample Data Table}
\label{tab:sample}
\end{table}

\section{Lists}
\subsection{Bullet Points}
\begin{itemize}
  \item First item
  \item Second item
  \item Third item with sub-items:
    \begin{itemize}
      \item Sub-item 1
      \item Sub-item 2
    \end{itemize}
\end{itemize}

\subsection{Numbered List}
\begin{enumerate}
  \item First step
  \item Second step
  \item Third step
\end{enumerate}

\section{Citations and References}
When you use the citation search feature, references will be automatically managed. For example, you might cite a paper like this: \cite{example2024}. All citations will be automatically added to your \texttt{cite.bib} file.

\section{Conclusion}
This template demonstrates basic \LaTeX{} features including mathematical formulas, tables, and lists. The bibliography system is now set up to work seamlessly with Bibby AI's citation management. Feel free to modify it for your needs!

\bibliographystyle{plain}
\bibliography{cite}

\end{document}\documentclass{article}
\usepackage[utf8]{inputenc}
\usepackage{booktabs}
\usepackage{amsmath}
\usepackage{natbib}

\title{My Document}
\author{Your Name}
\date{\today}

\begin{document}

\maketitle

\section{Introduction}
This is a sample document demonstrating various \LaTeX{} features including mathematical equations, tables, and formatting. You can cite papers using the citation search feature, and they will be automatically added to your \texttt{cite.bib} file.

\section{Mathematical Expressions}
Here are some example equations:

\subsection{Quadratic Formula}
The quadratic formula is given by:
\[x = \frac{-b \pm \sqrt{b^2 - 4ac}}{2a}\]

\subsection{Matrix Example}
A sample matrix:
\[
\begin{pmatrix}
a_{11} & a_{12} & a_{13} \\
a_{21} & a_{22} & a_{23} \\
a_{31} & a_{32} & a_{33}
\end{pmatrix}
\]

\section{Tables}
Below is a sample table showing data:

\begin{table}[h]
\centering
\begin{tabular}{@{}lrrr@{}}
\toprule
\textbf{Category} & \textbf{Value 1} & \textbf{Value 2} & \textbf{Total} \\
\midrule
Item A & 10.5 & 20.3 & 30.8 \\
Item B & 15.2 & 25.1 & 40.3 \\
Item C & 22.1 & 12.7 & 34.8 \\
\bottomrule
\end{tabular}
\caption{Sample Data Table}
\label{tab:sample}
\end{table}

\section{Lists}
\subsection{Bullet Points}
\begin{itemize}
  \item First item
  \item Second item
  \item Third item with sub-items:
    \begin{itemize}
      \item Sub-item 1
      \item Sub-item 2
    \end{itemize}
\end{itemize}

\subsection{Numbered List}
\begin{enumerate}
  \item First step
  \item Second step
  \item Third step
\end{enumerate}

\section{Citations and References}
When you use the citation search feature, references will be automatically managed. For example, you might cite a paper like this: \cite{example2024}. All citations will be automatically added to your \texttt{cite.bib} file.

\section{Conclusion}
This template demonstrates basic \LaTeX{} features including mathematical formulas, tables, and lists. The bibliography system is now set up to work seamlessly with Bibby AI's citation management. Feel free to modify it for your needs!

\bibliographystyle{plain}
\bibliography{cite}

\end{document}\documentclass{article}
\usepackage[utf8]{inputenc}
\usepackage{booktabs}
\usepackage{amsmath}
\usepackage{natbib}

\title{My Document}
\author{Your Name}
\date{\today}

\begin{document}

\maketitle

\section{Introduction}
This is a sample document demonstrating various \LaTeX{} features including mathematical equations, tables, and formatting. You can cite papers using the citation search feature, and they will be automatically added to your \texttt{cite.bib} file.

\section{Mathematical Expressions}
Here are some example equations:

\subsection{Quadratic Formula}
The quadratic formula is given by:
\[x = \frac{-b \pm \sqrt{b^2 - 4ac}}{2a}\]

\subsection{Matrix Example}
A sample matrix:
\[
\begin{pmatrix}
a_{11} & a_{12} & a_{13} \\
a_{21} & a_{22} & a_{23} \\
a_{31} & a_{32} & a_{33}
\end{pmatrix}
\]

\section{Tables}
Below is a sample table showing data:

\begin{table}[h]
\centering
\begin{tabular}{@{}lrrr@{}}
\toprule
\textbf{Category} & \textbf{Value 1} & \textbf{Value 2} & \textbf{Total} \\
\midrule
Item A & 10.5 & 20.3 & 30.8 \\
Item B & 15.2 & 25.1 & 40.3 \\
Item C & 22.1 & 12.7 & 34.8 \\
\bottomrule
\end{tabular}
\caption{Sample Data Table}
\label{tab:sample}
\end{table}

\section{Lists}
\subsection{Bullet Points}
\begin{itemize}
  \item First item
  \item Second item
  \item Third item with sub-items:
    \begin{itemize}
      \item Sub-item 1
      \item Sub-item 2
    \end{itemize}
\end{itemize}

\subsection{Numbered List}
\begin{enumerate}
  \item First step
  \item Second step
  \item Third step
\end{enumerate}

\section{Citations and References}
When you use the citation search feature, references will be automatically managed. For example, you might cite a paper like this: \cite{example2024}. All citations will be automatically added to your \texttt{cite.bib} file.

\section{Conclusion}
This template demonstrates basic \LaTeX{} features including mathematical formulas, tables, and lists. The bibliography system is now set up to work seamlessly with Bibby AI's citation management. Feel free to modify it for your needs!

\bibliographystyle{plain}
\bibliography{cite}

\end{document}\documentclass{article}
\usepackage[utf8]{inputenc}
\usepackage{booktabs}
\usepackage{amsmath}
\usepackage{natbib}

\title{My Document}
\author{Your Name}
\date{\today}

\begin{document}

\maketitle

\section{Introduction}
This is a sample document demonstrating various \LaTeX{} features including mathematical equations, tables, and formatting. You can cite papers using the citation search feature, and they will be automatically added to your \texttt{cite.bib} file.

\section{Mathematical Expressions}
Here are some example equations:

\subsection{Quadratic Formula}
The quadratic formula is given by:
\[x = \frac{-b \pm \sqrt{b^2 - 4ac}}{2a}\]

\subsection{Matrix Example}
A sample matrix:
\[
\begin{pmatrix}
a_{11} & a_{12} & a_{13} \\
a_{21} & a_{22} & a_{23} \\
a_{31} & a_{32} & a_{33}
\end{pmatrix}
\]

\section{Tables}
Below is a sample table showing data:

\begin{table}[h]
\centering
\begin{tabular}{@{}lrrr@{}}
\toprule
\textbf{Category} & \textbf{Value 1} & \textbf{Value 2} & \textbf{Total} \\
\midrule
Item A & 10.5 & 20.3 & 30.8 \\
Item B & 15.2 & 25.1 & 40.3 \\
Item C & 22.1 & 12.7 & 34.8 \\
\bottomrule
\end{tabular}
\caption{Sample Data Table}
\label{tab:sample}
\end{table}

\section{Lists}
\subsection{Bullet Points}
\begin{itemize}
  \item First item
  \item Second item
  \item Third item with sub-items:
    \begin{itemize}
      \item Sub-item 1
      \item Sub-item 2
    \end{itemize}
\end{itemize}

\subsection{Numbered List}
\begin{enumerate}
  \item First step
  \item Second step
  \item Third step
\end{enumerate}

\section{Citations and References}
When you use the citation search feature, references will be automatically managed. For example, you might cite a paper like this: \cite{example2024}. All citations will be automatically added to your \texttt{cite.bib} file.

\section{Conclusion}
This template demonstrates basic \LaTeX{} features including mathematical formulas, tables, and lists. The bibliography system is now set up to work seamlessly with Bibby AI's citation management. Feel free to modify it for your needs!

\bibliographystyle{plain}
\bibliography{cite}

\end{document}\documentclass{article}
\usepackage[utf8]{inputenc}
\usepackage{booktabs}
\usepackage{amsmath}
\usepackage{natbib}

\title{My Document}
\author{Your Name}
\date{\today}

\begin{document}

\maketitle

\section{Introduction}
This is a sample document demonstrating various \LaTeX{} features including mathematical equations, tables, and formatting. You can cite papers using the citation search feature, and they will be automatically added to your \texttt{cite.bib} file.

\section{Mathematical Expressions}
Here are some example equations:

\subsection{Quadratic Formula}
The quadratic formula is given by:
\[x = \frac{-b \pm \sqrt{b^2 - 4ac}}{2a}\]

\subsection{Matrix Example}
A sample matrix:
\[
\begin{pmatrix}
a_{11} & a_{12} & a_{13} \\
a_{21} & a_{22} & a_{23} \\
a_{31} & a_{32} & a_{33}
\end{pmatrix}
\]

\section{Tables}
Below is a sample table showing data:

\begin{table}[h]
\centering
\begin{tabular}{@{}lrrr@{}}
\toprule
\textbf{Category} & \textbf{Value 1} & \textbf{Value 2} & \textbf{Total} \\
\midrule
Item A & 10.5 & 20.3 & 30.8 \\
Item B & 15.2 & 25.1 & 40.3 \\
Item C & 22.1 & 12.7 & 34.8 \\
\bottomrule
\end{tabular}
\caption{Sample Data Table}
\label{tab:sample}
\end{table}

\section{Lists}
\subsection{Bullet Points}
\begin{itemize}
  \item First item
  \item Second item
  \item Third item with sub-items:
    \begin{itemize}
      \item Sub-item 1
      \item Sub-item 2
    \end{itemize}
\end{itemize}

\subsection{Numbered List}
\begin{enumerate}
  \item First step
  \item Second step
  \item Third step
\end{enumerate}

\section{Citations and References}
When you use the citation search feature, references will be automatically managed. For example, you might cite a paper like this: \cite{example2024}. All citations will be automatically added to your \texttt{cite.bib} file.

\section{Conclusion}
This template demonstrates basic \LaTeX{} features including mathematical formulas, tables, and lists. The bibliography system is now set up to work seamlessly with Bibby AI's citation management. Feel free to modify it for your needs!

\bibliographystyle{plain}
\bibliography{cite}

\end{document}\documentclass{article}
\usepackage[utf8]{inputenc}
\usepackage{booktabs}
\usepackage{amsmath}
\usepackage{natbib}

\title{My Document}
\author{Your Name}
\date{\today}

\begin{document}

\maketitle

\section{Introduction}
This is a sample document demonstrating various \LaTeX{} features including mathematical equations, tables, and formatting. You can cite papers using the citation search feature, and they will be automatically added to your \texttt{cite.bib} file.

\section{Mathematical Expressions}
Here are some example equations:

\subsection{Quadratic Formula}
The quadratic formula is given by:
\[x = \frac{-b \pm \sqrt{b^2 - 4ac}}{2a}\]

\subsection{Matrix Example}
A sample matrix:
\[
\begin{pmatrix}
a_{11} & a_{12} & a_{13} \\
a_{21} & a_{22} & a_{23} \\
a_{31} & a_{32} & a_{33}
\end{pmatrix}
\]

\section{Tables}
Below is a sample table showing data:

\begin{table}[h]
\centering
\begin{tabular}{@{}lrrr@{}}
\toprule
\textbf{Category} & \textbf{Value 1} & \textbf{Value 2} & \textbf{Total} \\
\midrule
Item A & 10.5 & 20.3 & 30.8 \\
Item B & 15.2 & 25.1 & 40.3 \\
Item C & 22.1 & 12.7 & 34.8 \\
\bottomrule
\end{tabular}
\caption{Sample Data Table}
\label{tab:sample}
\end{table}

\section{Lists}
\subsection{Bullet Points}
\begin{itemize}
  \item First item
  \item Second item
  \item Third item with sub-items:
    \begin{itemize}
      \item Sub-item 1
      \item Sub-item 2
    \end{itemize}
\end{itemize}

\subsection{Numbered List}
\begin{enumerate}
  \item First step
  \item Second step
  \item Third step
\end{enumerate}

\section{Citations and References}
When you use the citation search feature, references will be automatically managed. For example, you might cite a paper like this: \cite{example2024}. All citations will be automatically added to your \texttt{cite.bib} file.

\section{Conclusion}
This template demonstrates basic \LaTeX{} features including mathematical formulas, tables, and lists. The bibliography system is now set up to work seamlessly with Bibby AI's citation management. Feel free to modify it for your needs!

\bibliographystyle{plain}
\bibliography{cite}

\end{document}\documentclass{article}
\usepackage[utf8]{inputenc}
\usepackage{booktabs}
\usepackage{amsmath}
\usepackage{natbib}

\title{My Document}
\author{Your Name}
\date{\today}

\begin{document}

\maketitle

\section{Introduction}
This is a sample document demonstrating various \LaTeX{} features including mathematical equations, tables, and formatting. You can cite papers using the citation search feature, and they will be automatically added to your \texttt{cite.bib} file.

\section{Mathematical Expressions}
Here are some example equations:

\subsection{Quadratic Formula}
The quadratic formula is given by:
\[x = \frac{-b \pm \sqrt{b^2 - 4ac}}{2a}\]

\subsection{Matrix Example}
A sample matrix:
\[
\begin{pmatrix}
a_{11} & a_{12} & a_{13} \\
a_{21} & a_{22} & a_{23} \\
a_{31} & a_{32} & a_{33}
\end{pmatrix}
\]

\section{Tables}
Below is a sample table showing data:

\begin{table}[h]
\centering
\begin{tabular}{@{}lrrr@{}}
\toprule
\textbf{Category} & \textbf{Value 1} & \textbf{Value 2} & \textbf{Total} \\
\midrule
Item A & 10.5 & 20.3 & 30.8 \\
Item B & 15.2 & 25.1 & 40.3 \\
Item C & 22.1 & 12.7 & 34.8 \\
\bottomrule
\end{tabular}
\caption{Sample Data Table}
\label{tab:sample}
\end{table}

\section{Lists}
\subsection{Bullet Points}
\begin{itemize}
  \item First item
  \item Second item
  \item Third item with sub-items:
    \begin{itemize}
      \item Sub-item 1
      \item Sub-item 2
    \end{itemize}
\end{itemize}

\subsection{Numbered List}
\begin{enumerate}
  \item First step
  \item Second step
  \item Third step
\end{enumerate}

\section{Citations and References}
When you use the citation search feature, references will be automatically managed. For example, you might cite a paper like this: \cite{example2024}. All citations will be automatically added to your \texttt{cite.bib} file.

\section{Conclusion}
This template demonstrates basic \LaTeX{} features including mathematical formulas, tables, and lists. The bibliography system is now set up to work seamlessly with Bibby AI's citation management. Feel free to modify it for your needs!

\bibliographystyle{plain}
\bibliography{cite}

\end{document}\documentclass{article}
\usepackage[utf8]{inputenc}
\usepackage{booktabs}
\usepackage{amsmath}
\usepackage{natbib}

\title{My Document}
\author{Your Name}
\date{\today}

\begin{document}

\maketitle

\section{Introduction}
This is a sample document demonstrating various \LaTeX{} features including mathematical equations, tables, and formatting. You can cite papers using the citation search feature, and they will be automatically added to your \texttt{cite.bib} file.

\section{Mathematical Expressions}
Here are some example equations:

\subsection{Quadratic Formula}
The quadratic formula is given by:
\[x = \frac{-b \pm \sqrt{b^2 - 4ac}}{2a}\]

\subsection{Matrix Example}
A sample matrix:
\[
\begin{pmatrix}
a_{11} & a_{12} & a_{13} \\
a_{21} & a_{22} & a_{23} \\
a_{31} & a_{32} & a_{33}
\end{pmatrix}
\]

\section{Tables}
Below is a sample table showing data:

\begin{table}[h]
\centering
\begin{tabular}{@{}lrrr@{}}
\toprule
\textbf{Category} & \textbf{Value 1} & \textbf{Value 2} & \textbf{Total} \\
\midrule
Item A & 10.5 & 20.3 & 30.8 \\
Item B & 15.2 & 25.1 & 40.3 \\
Item C & 22.1 & 12.7 & 34.8 \\
\bottomrule
\end{tabular}
\caption{Sample Data Table}
\label{tab:sample}
\end{table}

\section{Lists}
\subsection{Bullet Points}
\begin{itemize}
  \item First item
  \item Second item
  \item Third item with sub-items:
    \begin{itemize}
      \item Sub-item 1
      \item Sub-item 2
    \end{itemize}
\end{itemize}

\subsection{Numbered List}
\begin{enumerate}
  \item First step
  \item Second step
  \item Third step
\end{enumerate}

\section{Citations and References}
When you use the citation search feature, references will be automatically managed. For example, you might cite a paper like this: \cite{example2024}. All citations will be automatically added to your \texttt{cite.bib} file.

\section{Conclusion}
This template demonstrates basic \LaTeX{} features including mathematical formulas, tables, and lists. The bibliography system is now set up to work seamlessly with Bibby AI's citation management. Feel free to modify it for your needs!

\bibliographystyle{plain}
\bibliography{cite}

\end{document}\documentclass{article}
\usepackage[utf8]{inputenc}
\usepackage{booktabs}
\usepackage{amsmath}
\usepackage{natbib}

\title{My Document}
\author{Your Name}
\date{\today}

\begin{document}

\maketitle

\section{Introduction}
This is a sample document demonstrating various \LaTeX{} features including mathematical equations, tables, and formatting. You can cite papers using the citation search feature, and they will be automatically added to your \texttt{cite.bib} file.

\section{Mathematical Expressions}
Here are some example equations:

\subsection{Quadratic Formula}
The quadratic formula is given by:
\[x = \frac{-b \pm \sqrt{b^2 - 4ac}}{2a}\]

\subsection{Matrix Example}
A sample matrix:
\[
\begin{pmatrix}
a_{11} & a_{12} & a_{13} \\
a_{21} & a_{22} & a_{23} \\
a_{31} & a_{32} & a_{33}
\end{pmatrix}
\]

\section{Tables}
Below is a sample table showing data:

\begin{table}[h]
\centering
\begin{tabular}{@{}lrrr@{}}
\toprule
\textbf{Category} & \textbf{Value 1} & \textbf{Value 2} & \textbf{Total} \\
\midrule
Item A & 10.5 & 20.3 & 30.8 \\
Item B & 15.2 & 25.1 & 40.3 \\
Item C & 22.1 & 12.7 & 34.8 \\
\bottomrule
\end{tabular}
\caption{Sample Data Table}
\label{tab:sample}
\end{table}

\section{Lists}
\subsection{Bullet Points}
\begin{itemize}
  \item First item
  \item Second item
  \item Third item with sub-items:
    \begin{itemize}
      \item Sub-item 1
      \item Sub-item 2
    \end{itemize}
\end{itemize}

\subsection{Numbered List}
\begin{enumerate}
  \item First step
  \item Second step
  \item Third step
\end{enumerate}

\section{Citations and References}
When you use the citation search feature, references will be automatically managed. For example, you might cite a paper like this: \cite{example2024}. All citations will be automatically added to your \texttt{cite.bib} file.

\section{Conclusion}
This template demonstrates basic \LaTeX{} features including mathematical formulas, tables, and lists. The bibliography system is now set up to work seamlessly with Bibby AI's citation management. Feel free to modify it for your needs!

\bibliographystyle{plain}
\bibliography{cite}

\end{document}\documentclass{article}
\usepackage[utf8]{inputenc}
\usepackage{booktabs}
\usepackage{amsmath}
\usepackage{natbib}

\title{My Document}
\author{Your Name}
\date{\today}

\begin{document}

\maketitle

\section{Introduction}
This is a sample document demonstrating various \LaTeX{} features including mathematical equations, tables, and formatting. You can cite papers using the citation search feature, and they will be automatically added to your \texttt{cite.bib} file.

\section{Mathematical Expressions}
Here are some example equations:

\subsection{Quadratic Formula}
The quadratic formula is given by:
\[x = \frac{-b \pm \sqrt{b^2 - 4ac}}{2a}\]

\subsection{Matrix Example}
A sample matrix:
\[
\begin{pmatrix}
a_{11} & a_{12} & a_{13} \\
a_{21} & a_{22} & a_{23} \\
a_{31} & a_{32} & a_{33}
\end{pmatrix}
\]

\section{Tables}
Below is a sample table showing data:

\begin{table}[h]
\centering
\begin{tabular}{@{}lrrr@{}}
\toprule
\textbf{Category} & \textbf{Value 1} & \textbf{Value 2} & \textbf{Total} \\
\midrule
Item A & 10.5 & 20.3 & 30.8 \\
Item B & 15.2 & 25.1 & 40.3 \\
Item C & 22.1 & 12.7 & 34.8 \\
\bottomrule
\end{tabular}
\caption{Sample Data Table}
\label{tab:sample}
\end{table}

\section{Lists}
\subsection{Bullet Points}
\begin{itemize}
  \item First item
  \item Second item
  \item Third item with sub-items:
    \begin{itemize}
      \item Sub-item 1
      \item Sub-item 2
    \end{itemize}
\end{itemize}

\subsection{Numbered List}
\begin{enumerate}
  \item First step
  \item Second step
  \item Third step
\end{enumerate}

\section{Citations and References}
When you use the citation search feature, references will be automatically managed. For example, you might cite a paper like this: \cite{example2024}. All citations will be automatically added to your \texttt{cite.bib} file.

\section{Conclusion}
This template demonstrates basic \LaTeX{} features including mathematical formulas, tables, and lists. The bibliography system is now set up to work seamlessly with Bibby AI's citation management. Feel free to modify it for your needs!

\bibliographystyle{plain}
\bibliography{cite}

\end{document}\documentclass{article}
\usepackage[utf8]{inputenc}
\usepackage{booktabs}
\usepackage{amsmath}
\usepackage{natbib}

\title{My Document}
\author{Your Name}
\date{\today}

\begin{document}

\maketitle

\section{Introduction}
This is a sample document demonstrating various \LaTeX{} features including mathematical equations, tables, and formatting. You can cite papers using the citation search feature, and they will be automatically added to your \texttt{cite.bib} file.

\section{Mathematical Expressions}
Here are some example equations:

\subsection{Quadratic Formula}
The quadratic formula is given by:
\[x = \frac{-b \pm \sqrt{b^2 - 4ac}}{2a}\]

\subsection{Matrix Example}
A sample matrix:
\[
\begin{pmatrix}
a_{11} & a_{12} & a_{13} \\
a_{21} & a_{22} & a_{23} \\
a_{31} & a_{32} & a_{33}
\end{pmatrix}
\]

\section{Tables}
Below is a sample table showing data:

\begin{table}[h]
\centering
\begin{tabular}{@{}lrrr@{}}
\toprule
\textbf{Category} & \textbf{Value 1} & \textbf{Value 2} & \textbf{Total} \\
\midrule
Item A & 10.5 & 20.3 & 30.8 \\
Item B & 15.2 & 25.1 & 40.3 \\
Item C & 22.1 & 12.7 & 34.8 \\
\bottomrule
\end{tabular}
\caption{Sample Data Table}
\label{tab:sample}
\end{table}

\section{Lists}
\subsection{Bullet Points}
\begin{itemize}
  \item First item
  \item Second item
  \item Third item with sub-items:
    \begin{itemize}
      \item Sub-item 1
      \item Sub-item 2
    \end{itemize}
\end{itemize}

\subsection{Numbered List}
\begin{enumerate}
  \item First step
  \item Second step
  \item Third step
\end{enumerate}

\section{Citations and References}
When you use the citation search feature, references will be automatically managed. For example, you might cite a paper like this: \cite{example2024}. All citations will be automatically added to your \texttt{cite.bib} file.

\section{Conclusion}
This template demonstrates basic \LaTeX{} features including mathematical formulas, tables, and lists. The bibliography system is now set up to work seamlessly with Bibby AI's citation management. Feel free to modify it for your needs!

\bibliographystyle{plain}
\bibliography{cite}

\end{document}\documentclass{article}
\usepackage[utf8]{inputenc}
\usepackage{booktabs}
\usepackage{amsmath}
\usepackage{natbib}

\title{My Document}
\author{Your Name}
\date{\today}

\begin{document}

\maketitle

\section{Introduction}
This is a sample document demonstrating various \LaTeX{} features including mathematical equations, tables, and formatting. You can cite papers using the citation search feature, and they will be automatically added to your \texttt{cite.bib} file.

\section{Mathematical Expressions}
Here are some example equations:

\subsection{Quadratic Formula}
The quadratic formula is given by:
\[x = \frac{-b \pm \sqrt{b^2 - 4ac}}{2a}\]

\subsection{Matrix Example}
A sample matrix:
\[
\begin{pmatrix}
a_{11} & a_{12} & a_{13} \\
a_{21} & a_{22} & a_{23} \\
a_{31} & a_{32} & a_{33}
\end{pmatrix}
\]

\section{Tables}
Below is a sample table showing data:

\begin{table}[h]
\centering
\begin{tabular}{@{}lrrr@{}}
\toprule
\textbf{Category} & \textbf{Value 1} & \textbf{Value 2} & \textbf{Total} \\
\midrule
Item A & 10.5 & 20.3 & 30.8 \\
Item B & 15.2 & 25.1 & 40.3 \\
Item C & 22.1 & 12.7 & 34.8 \\
\bottomrule
\end{tabular}
\caption{Sample Data Table}
\label{tab:sample}
\end{table}

\section{Lists}
\subsection{Bullet Points}
\begin{itemize}
  \item First item
  \item Second item
  \item Third item with sub-items:
    \begin{itemize}
      \item Sub-item 1
      \item Sub-item 2
    \end{itemize}
\end{itemize}

\subsection{Numbered List}
\begin{enumerate}
  \item First step
  \item Second step
  \item Third step
\end{enumerate}

\section{Citations and References}
When you use the citation search feature, references will be automatically managed. For example, you might cite a paper like this: \cite{example2024}. All citations will be automatically added to your \texttt{cite.bib} file.

\section{Conclusion}
This template demonstrates basic \LaTeX{} features including mathematical formulas, tables, and lists. The bibliography system is now set up to work seamlessly with Bibby AI's citation management. Feel free to modify it for your needs!

\bibliographystyle{plain}
\bibliography{cite}

\end{document}\documentclass{article}
\usepackage[utf8]{inputenc}
\usepackage{booktabs}
\usepackage{amsmath}
\usepackage{natbib}

\title{My Document}
\author{Your Name}
\date{\today}

\begin{document}

\maketitle

\section{Introduction}
This is a sample document demonstrating various \LaTeX{} features including mathematical equations, tables, and formatting. You can cite papers using the citation search feature, and they will be automatically added to your \texttt{cite.bib} file.

\section{Mathematical Expressions}
Here are some example equations:

\subsection{Quadratic Formula}
The quadratic formula is given by:
\[x = \frac{-b \pm \sqrt{b^2 - 4ac}}{2a}\]

\subsection{Matrix Example}
A sample matrix:
\[
\begin{pmatrix}
a_{11} & a_{12} & a_{13} \\
a_{21} & a_{22} & a_{23} \\
a_{31} & a_{32} & a_{33}
\end{pmatrix}
\]

\section{Tables}
Below is a sample table showing data:

\begin{table}[h]
\centering
\begin{tabular}{@{}lrrr@{}}
\toprule
\textbf{Category} & \textbf{Value 1} & \textbf{Value 2} & \textbf{Total} \\
\midrule
Item A & 10.5 & 20.3 & 30.8 \\
Item B & 15.2 & 25.1 & 40.3 \\
Item C & 22.1 & 12.7 & 34.8 \\
\bottomrule
\end{tabular}
\caption{Sample Data Table}
\label{tab:sample}
\end{table}

\section{Lists}
\subsection{Bullet Points}
\begin{itemize}
  \item First item
  \item Second item
  \item Third item with sub-items:
    \begin{itemize}
      \item Sub-item 1
      \item Sub-item 2
    \end{itemize}
\end{itemize}

\subsection{Numbered List}
\begin{enumerate}
  \item First step
  \item Second step
  \item Third step
\end{enumerate}

\section{Citations and References}
When you use the citation search feature, references will be automatically managed. For example, you might cite a paper like this: \cite{example2024}. All citations will be automatically added to your \texttt{cite.bib} file.

\section{Conclusion}
This template demonstrates basic \LaTeX{} features including mathematical formulas, tables, and lists. The bibliography system is now set up to work seamlessly with Bibby AI's citation management. Feel free to modify it for your needs!

\bibliographystyle{plain}
\bibliography{cite}

\end{document}\documentclass{article}
\usepackage[utf8]{inputenc}
\usepackage{booktabs}
\usepackage{amsmath}
\usepackage{natbib}

\title{My Document}
\author{Your Name}
\date{\today}

\begin{document}

\maketitle

\section{Introduction}
This is a sample document demonstrating various \LaTeX{} features including mathematical equations, tables, and formatting. You can cite papers using the citation search feature, and they will be automatically added to your \texttt{cite.bib} file.

\section{Mathematical Expressions}
Here are some example equations:

\subsection{Quadratic Formula}
The quadratic formula is given by:
\[x = \frac{-b \pm \sqrt{b^2 - 4ac}}{2a}\]

\subsection{Matrix Example}
A sample matrix:
\[
\begin{pmatrix}
a_{11} & a_{12} & a_{13} \\
a_{21} & a_{22} & a_{23} \\
a_{31} & a_{32} & a_{33}
\end{pmatrix}
\]

\section{Tables}
Below is a sample table showing data:

\begin{table}[h]
\centering
\begin{tabular}{@{}lrrr@{}}
\toprule
\textbf{Category} & \textbf{Value 1} & \textbf{Value 2} & \textbf{Total} \\
\midrule
Item A & 10.5 & 20.3 & 30.8 \\
Item B & 15.2 & 25.1 & 40.3 \\
Item C & 22.1 & 12.7 & 34.8 \\
\bottomrule
\end{tabular}
\caption{Sample Data Table}
\label{tab:sample}
\end{table}

\section{Lists}
\subsection{Bullet Points}
\begin{itemize}
  \item First item
  \item Second item
  \item Third item with sub-items:
    \begin{itemize}
      \item Sub-item 1
      \item Sub-item 2
    \end{itemize}
\end{itemize}

\subsection{Numbered List}
\begin{enumerate}
  \item First step
  \item Second step
  \item Third step
\end{enumerate}

\section{Citations and References}
When you use the citation search feature, references will be automatically managed. For example, you might cite a paper like this: \cite{example2024}. All citations will be automatically added to your \texttt{cite.bib} file.

\section{Conclusion}
This template demonstrates basic \LaTeX{} features including mathematical formulas, tables, and lists. The bibliography system is now set up to work seamlessly with Bibby AI's citation management. Feel free to modify it for your needs!

\bibliographystyle{plain}
\bibliography{cite}

\end{document}\documentclass{article}
\usepackage[utf8]{inputenc}
\usepackage{booktabs}
\usepackage{amsmath}
\usepackage{natbib}

\title{My Document}
\author{Your Name}
\date{\today}

\begin{document}

\maketitle

\section{Introduction}
This is a sample document demonstrating various \LaTeX{} features including mathematical equations, tables, and formatting. You can cite papers using the citation search feature, and they will be automatically added to your \texttt{cite.bib} file.

\section{Mathematical Expressions}
Here are some example equations:

\subsection{Quadratic Formula}
The quadratic formula is given by:
\[x = \frac{-b \pm \sqrt{b^2 - 4ac}}{2a}\]

\subsection{Matrix Example}
A sample matrix:
\[
\begin{pmatrix}
a_{11} & a_{12} & a_{13} \\
a_{21} & a_{22} & a_{23} \\
a_{31} & a_{32} & a_{33}
\end{pmatrix}
\]

\section{Tables}
Below is a sample table showing data:

\begin{table}[h]
\centering
\begin{tabular}{@{}lrrr@{}}
\toprule
\textbf{Category} & \textbf{Value 1} & \textbf{Value 2} & \textbf{Total} \\
\midrule
Item A & 10.5 & 20.3 & 30.8 \\
Item B & 15.2 & 25.1 & 40.3 \\
Item C & 22.1 & 12.7 & 34.8 \\
\bottomrule
\end{tabular}
\caption{Sample Data Table}
\label{tab:sample}
\end{table}

\section{Lists}
\subsection{Bullet Points}
\begin{itemize}
  \item First item
  \item Second item
  \item Third item with sub-items:
    \begin{itemize}
      \item Sub-item 1
      \item Sub-item 2
    \end{itemize}
\end{itemize}

\subsection{Numbered List}
\begin{enumerate}
  \item First step
  \item Second step
  \item Third step
\end{enumerate}

\section{Citations and References}
When you use the citation search feature, references will be automatically managed. For example, you might cite a paper like this: \cite{example2024}. All citations will be automatically added to your \texttt{cite.bib} file.

\section{Conclusion}
This template demonstrates basic \LaTeX{} features including mathematical formulas, tables, and lists. The bibliography system is now set up to work seamlessly with Bibby AI's citation management. Feel free to modify it for your needs!

\bibliographystyle{plain}
\bibliography{cite}

\end{document}\documentclass{article}
\usepackage[utf8]{inputenc}
\usepackage{booktabs}
\usepackage{amsmath}
\usepackage{natbib}

\title{My Document}
\author{Your Name}
\date{\today}

\begin{document}

\maketitle

\section{Introduction}
This is a sample document demonstrating various \LaTeX{} features including mathematical equations, tables, and formatting. You can cite papers using the citation search feature, and they will be automatically added to your \texttt{cite.bib} file.

\section{Mathematical Expressions}
Here are some example equations:

\subsection{Quadratic Formula}
The quadratic formula is given by:
\[x = \frac{-b \pm \sqrt{b^2 - 4ac}}{2a}\]

\subsection{Matrix Example}
A sample matrix:
\[
\begin{pmatrix}
a_{11} & a_{12} & a_{13} \\
a_{21} & a_{22} & a_{23} \\
a_{31} & a_{32} & a_{33}
\end{pmatrix}
\]

\section{Tables}
Below is a sample table showing data:

\begin{table}[h]
\centering
\begin{tabular}{@{}lrrr@{}}
\toprule
\textbf{Category} & \textbf{Value 1} & \textbf{Value 2} & \textbf{Total} \\
\midrule
Item A & 10.5 & 20.3 & 30.8 \\
Item B & 15.2 & 25.1 & 40.3 \\
Item C & 22.1 & 12.7 & 34.8 \\
\bottomrule
\end{tabular}
\caption{Sample Data Table}
\label{tab:sample}
\end{table}

\section{Lists}
\subsection{Bullet Points}
\begin{itemize}
  \item First item
  \item Second item
  \item Third item with sub-items:
    \begin{itemize}
      \item Sub-item 1
      \item Sub-item 2
    \end{itemize}
\end{itemize}

\subsection{Numbered List}
\begin{enumerate}
  \item First step
  \item Second step
  \item Third step
\end{enumerate}

\section{Citations and References}
When you use the citation search feature, references will be automatically managed. For example, you might cite a paper like this: \cite{example2024}. All citations will be automatically added to your \texttt{cite.bib} file.

\section{Conclusion}
This template demonstrates basic \LaTeX{} features including mathematical formulas, tables, and lists. The bibliography system is now set up to work seamlessly with Bibby AI's citation management. Feel free to modify it for your needs!

\bibliographystyle{plain}
\bibliography{cite}

\end{document}\documentclass{article}
\usepackage[utf8]{inputenc}
\usepackage{booktabs}
\usepackage{amsmath}
\usepackage{natbib}

\title{My Document}
\author{Your Name}
\date{\today}

\begin{document}

\maketitle

\section{Introduction}
This is a sample document demonstrating various \LaTeX{} features including mathematical equations, tables, and formatting. You can cite papers using the citation search feature, and they will be automatically added to your \texttt{cite.bib} file.

\section{Mathematical Expressions}
Here are some example equations:

\subsection{Quadratic Formula}
The quadratic formula is given by:
\[x = \frac{-b \pm \sqrt{b^2 - 4ac}}{2a}\]

\subsection{Matrix Example}
A sample matrix:
\[
\begin{pmatrix}
a_{11} & a_{12} & a_{13} \\
a_{21} & a_{22} & a_{23} \\
a_{31} & a_{32} & a_{33}
\end{pmatrix}
\]

\section{Tables}
Below is a sample table showing data:

\begin{table}[h]
\centering
\begin{tabular}{@{}lrrr@{}}
\toprule
\textbf{Category} & \textbf{Value 1} & \textbf{Value 2} & \textbf{Total} \\
\midrule
Item A & 10.5 & 20.3 & 30.8 \\
Item B & 15.2 & 25.1 & 40.3 \\
Item C & 22.1 & 12.7 & 34.8 \\
\bottomrule
\end{tabular}
\caption{Sample Data Table}
\label{tab:sample}
\end{table}

\section{Lists}
\subsection{Bullet Points}
\begin{itemize}
  \item First item
  \item Second item
  \item Third item with sub-items:
    \begin{itemize}
      \item Sub-item 1
      \item Sub-item 2
    \end{itemize}
\end{itemize}

\subsection{Numbered List}
\begin{enumerate}
  \item First step
  \item Second step
  \item Third step
\end{enumerate}

\section{Citations and References}
When you use the citation search feature, references will be automatically managed. For example, you might cite a paper like this: \cite{example2024}. All citations will be automatically added to your \texttt{cite.bib} file.

\section{Conclusion}
This template demonstrates basic \LaTeX{} features including mathematical formulas, tables, and lists. The bibliography system is now set up to work seamlessly with Bibby AI's citation management. Feel free to modify it for your needs!

\bibliographystyle{plain}
\bibliography{cite}

\end{document}\documentclass{article}
\usepackage[utf8]{inputenc}
\usepackage{booktabs}
\usepackage{amsmath}
\usepackage{natbib}

\title{My Document}
\author{Your Name}
\date{\today}

\begin{document}

\maketitle

\section{Introduction}
This is a sample document demonstrating various \LaTeX{} features including mathematical equations, tables, and formatting. You can cite papers using the citation search feature, and they will be automatically added to your \texttt{cite.bib} file.

\section{Mathematical Expressions}
Here are some example equations:

\subsection{Quadratic Formula}
The quadratic formula is given by:
\[x = \frac{-b \pm \sqrt{b^2 - 4ac}}{2a}\]

\subsection{Matrix Example}
A sample matrix:
\[
\begin{pmatrix}
a_{11} & a_{12} & a_{13} \\
a_{21} & a_{22} & a_{23} \\
a_{31} & a_{32} & a_{33}
\end{pmatrix}
\]

\section{Tables}
Below is a sample table showing data:

\begin{table}[h]
\centering
\begin{tabular}{@{}lrrr@{}}
\toprule
\textbf{Category} & \textbf{Value 1} & \textbf{Value 2} & \textbf{Total} \\
\midrule
Item A & 10.5 & 20.3 & 30.8 \\
Item B & 15.2 & 25.1 & 40.3 \\
Item C & 22.1 & 12.7 & 34.8 \\
\bottomrule
\end{tabular}
\caption{Sample Data Table}
\label{tab:sample}
\end{table}

\section{Lists}
\subsection{Bullet Points}
\begin{itemize}
  \item First item
  \item Second item
  \item Third item with sub-items:
    \begin{itemize}
      \item Sub-item 1
      \item Sub-item 2
    \end{itemize}
\end{itemize}

\subsection{Numbered List}
\begin{enumerate}
  \item First step
  \item Second step
  \item Third step
\end{enumerate}

\section{Citations and References}
When you use the citation search feature, references will be automatically managed. For example, you might cite a paper like this: \cite{example2024}. All citations will be automatically added to your \texttt{cite.bib} file.

\section{Conclusion}
This template demonstrates basic \LaTeX{} features including mathematical formulas, tables, and lists. The bibliography system is now set up to work seamlessly with Bibby AI's citation management. Feel free to modify it for your needs!

\bibliographystyle{plain}
\bibliography{cite}

\end{document}\documentclass{article}
\usepackage[utf8]{inputenc}
\usepackage{booktabs}
\usepackage{amsmath}
\usepackage{natbib}

\title{My Document}
\author{Your Name}
\date{\today}

\begin{document}

\maketitle

\section{Introduction}
This is a sample document demonstrating various \LaTeX{} features including mathematical equations, tables, and formatting. You can cite papers using the citation search feature, and they will be automatically added to your \texttt{cite.bib} file.

\section{Mathematical Expressions}
Here are some example equations:

\subsection{Quadratic Formula}
The quadratic formula is given by:
\[x = \frac{-b \pm \sqrt{b^2 - 4ac}}{2a}\]

\subsection{Matrix Example}
A sample matrix:
\[
\begin{pmatrix}
a_{11} & a_{12} & a_{13} \\
a_{21} & a_{22} & a_{23} \\
a_{31} & a_{32} & a_{33}
\end{pmatrix}
\]

\section{Tables}
Below is a sample table showing data:

\begin{table}[h]
\centering
\begin{tabular}{@{}lrrr@{}}
\toprule
\textbf{Category} & \textbf{Value 1} & \textbf{Value 2} & \textbf{Total} \\
\midrule
Item A & 10.5 & 20.3 & 30.8 \\
Item B & 15.2 & 25.1 & 40.3 \\
Item C & 22.1 & 12.7 & 34.8 \\
\bottomrule
\end{tabular}
\caption{Sample Data Table}
\label{tab:sample}
\end{table}

\section{Lists}
\subsection{Bullet Points}
\begin{itemize}
  \item First item
  \item Second item
  \item Third item with sub-items:
    \begin{itemize}
      \item Sub-item 1
      \item Sub-item 2
    \end{itemize}
\end{itemize}

\subsection{Numbered List}
\begin{enumerate}
  \item First step
  \item Second step
  \item Third step
\end{enumerate}

\section{Citations and References}
When you use the citation search feature, references will be automatically managed. For example, you might cite a paper like this: \cite{example2024}. All citations will be automatically added to your \texttt{cite.bib} file.

\section{Conclusion}
This template demonstrates basic \LaTeX{} features including mathematical formulas, tables, and lists. The bibliography system is now set up to work seamlessly with Bibby AI's citation management. Feel free to modify it for your needs!

\bibliographystyle{plain}
\bibliography{cite}

\end{document}\documentclass{article}
\usepackage[utf8]{inputenc}
\usepackage{booktabs}
\usepackage{amsmath}
\usepackage{natbib}

\title{My Document}
\author{Your Name}
\date{\today}

\begin{document}

\maketitle

\section{Introduction}
This is a sample document demonstrating various \LaTeX{} features including mathematical equations, tables, and formatting. You can cite papers using the citation search feature, and they will be automatically added to your \texttt{cite.bib} file.

\section{Mathematical Expressions}
Here are some example equations:

\subsection{Quadratic Formula}
The quadratic formula is given by:
\[x = \frac{-b \pm \sqrt{b^2 - 4ac}}{2a}\]

\subsection{Matrix Example}
A sample matrix:
\[
\begin{pmatrix}
a_{11} & a_{12} & a_{13} \\
a_{21} & a_{22} & a_{23} \\
a_{31} & a_{32} & a_{33}
\end{pmatrix}
\]

\section{Tables}
Below is a sample table showing data:

\begin{table}[h]
\centering
\begin{tabular}{@{}lrrr@{}}
\toprule
\textbf{Category} & \textbf{Value 1} & \textbf{Value 2} & \textbf{Total} \\
\midrule
Item A & 10.5 & 20.3 & 30.8 \\
Item B & 15.2 & 25.1 & 40.3 \\
Item C & 22.1 & 12.7 & 34.8 \\
\bottomrule
\end{tabular}
\caption{Sample Data Table}
\label{tab:sample}
\end{table}

\section{Lists}
\subsection{Bullet Points}
\begin{itemize}
  \item First item
  \item Second item
  \item Third item with sub-items:
    \begin{itemize}
      \item Sub-item 1
      \item Sub-item 2
    \end{itemize}
\end{itemize}

\subsection{Numbered List}
\begin{enumerate}
  \item First step
  \item Second step
  \item Third step
\end{enumerate}

\section{Citations and References}
When you use the citation search feature, references will be automatically managed. For example, you might cite a paper like this: \cite{example2024}. All citations will be automatically added to your \texttt{cite.bib} file.

\section{Conclusion}
This template demonstrates basic \LaTeX{} features including mathematical formulas, tables, and lists. The bibliography system is now set up to work seamlessly with Bibby AI's citation management. Feel free to modify it for your needs!

\bibliographystyle{plain}
\bibliography{cite}

\end{document}\documentclass{article}
\usepackage[utf8]{inputenc}
\usepackage{booktabs}
\usepackage{amsmath}
\usepackage{natbib}

\title{My Document}
\author{Your Name}
\date{\today}

\begin{document}

\maketitle

\section{Introduction}
This is a sample document demonstrating various \LaTeX{} features including mathematical equations, tables, and formatting. You can cite papers using the citation search feature, and they will be automatically added to your \texttt{cite.bib} file.

\section{Mathematical Expressions}
Here are some example equations:

\subsection{Quadratic Formula}
The quadratic formula is given by:
\[x = \frac{-b \pm \sqrt{b^2 - 4ac}}{2a}\]

\subsection{Matrix Example}
A sample matrix:
\[
\begin{pmatrix}
a_{11} & a_{12} & a_{13} \\
a_{21} & a_{22} & a_{23} \\
a_{31} & a_{32} & a_{33}
\end{pmatrix}
\]

\section{Tables}
Below is a sample table showing data:

\begin{table}[h]
\centering
\begin{tabular}{@{}lrrr@{}}
\toprule
\textbf{Category} & \textbf{Value 1} & \textbf{Value 2} & \textbf{Total} \\
\midrule
Item A & 10.5 & 20.3 & 30.8 \\
Item B & 15.2 & 25.1 & 40.3 \\
Item C & 22.1 & 12.7 & 34.8 \\
\bottomrule
\end{tabular}
\caption{Sample Data Table}
\label{tab:sample}
\end{table}

\section{Lists}
\subsection{Bullet Points}
\begin{itemize}
  \item First item
  \item Second item
  \item Third item with sub-items:
    \begin{itemize}
      \item Sub-item 1
      \item Sub-item 2
    \end{itemize}
\end{itemize}

\subsection{Numbered List}
\begin{enumerate}
  \item First step
  \item Second step
  \item Third step
\end{enumerate}

\section{Citations and References}
When you use the citation search feature, references will be automatically managed. For example, you might cite a paper like this: \cite{example2024}. All citations will be automatically added to your \texttt{cite.bib} file.

\section{Conclusion}
This template demonstrates basic \LaTeX{} features including mathematical formulas, tables, and lists. The bibliography system is now set up to work seamlessly with Bibby AI's citation management. Feel free to modify it for your needs!

\bibliographystyle{plain}
\bibliography{cite}

\end{document}\documentclass{article}
\usepackage[utf8]{inputenc}
\usepackage{booktabs}
\usepackage{amsmath}
\usepackage{natbib}

\title{My Document}
\author{Your Name}
\date{\today}

\begin{document}

\maketitle

\section{Introduction}
This is a sample document demonstrating various \LaTeX{} features including mathematical equations, tables, and formatting. You can cite papers using the citation search feature, and they will be automatically added to your \texttt{cite.bib} file.

\section{Mathematical Expressions}
Here are some example equations:

\subsection{Quadratic Formula}
The quadratic formula is given by:
\[x = \frac{-b \pm \sqrt{b^2 - 4ac}}{2a}\]

\subsection{Matrix Example}
A sample matrix:
\[
\begin{pmatrix}
a_{11} & a_{12} & a_{13} \\
a_{21} & a_{22} & a_{23} \\
a_{31} & a_{32} & a_{33}
\end{pmatrix}
\]

\section{Tables}
Below is a sample table showing data:

\begin{table}[h]
\centering
\begin{tabular}{@{}lrrr@{}}
\toprule
\textbf{Category} & \textbf{Value 1} & \textbf{Value 2} & \textbf{Total} \\
\midrule
Item A & 10.5 & 20.3 & 30.8 \\
Item B & 15.2 & 25.1 & 40.3 \\
Item C & 22.1 & 12.7 & 34.8 \\
\bottomrule
\end{tabular}
\caption{Sample Data Table}
\label{tab:sample}
\end{table}

\section{Lists}
\subsection{Bullet Points}
\begin{itemize}
  \item First item
  \item Second item
  \item Third item with sub-items:
    \begin{itemize}
      \item Sub-item 1
      \item Sub-item 2
    \end{itemize}
\end{itemize}

\subsection{Numbered List}
\begin{enumerate}
  \item First step
  \item Second step
  \item Third step
\end{enumerate}

\section{Citations and References}
When you use the citation search feature, references will be automatically managed. For example, you might cite a paper like this: \cite{example2024}. All citations will be automatically added to your \texttt{cite.bib} file.

\section{Conclusion}
This template demonstrates basic \LaTeX{} features including mathematical formulas, tables, and lists. The bibliography system is now set up to work seamlessly with Bibby AI's citation management. Feel free to modify it for your needs!

\bibliographystyle{plain}
\bibliography{cite}

\end{document}\documentclass{article}
\usepackage[utf8]{inputenc}
\usepackage{booktabs}
\usepackage{amsmath}
\usepackage{natbib}

\title{My Document}
\author{Your Name}
\date{\today}

\begin{document}

\maketitle

\section{Introduction}
This is a sample document demonstrating various \LaTeX{} features including mathematical equations, tables, and formatting. You can cite papers using the citation search feature, and they will be automatically added to your \texttt{cite.bib} file.

\section{Mathematical Expressions}
Here are some example equations:

\subsection{Quadratic Formula}
The quadratic formula is given by:
\[x = \frac{-b \pm \sqrt{b^2 - 4ac}}{2a}\]

\subsection{Matrix Example}
A sample matrix:
\[
\begin{pmatrix}
a_{11} & a_{12} & a_{13} \\
a_{21} & a_{22} & a_{23} \\
a_{31} & a_{32} & a_{33}
\end{pmatrix}
\]

\section{Tables}
Below is a sample table showing data:

\begin{table}[h]
\centering
\begin{tabular}{@{}lrrr@{}}
\toprule
\textbf{Category} & \textbf{Value 1} & \textbf{Value 2} & \textbf{Total} \\
\midrule
Item A & 10.5 & 20.3 & 30.8 \\
Item B & 15.2 & 25.1 & 40.3 \\
Item C & 22.1 & 12.7 & 34.8 \\
\bottomrule
\end{tabular}
\caption{Sample Data Table}
\label{tab:sample}
\end{table}

\section{Lists}
\subsection{Bullet Points}
\begin{itemize}
  \item First item
  \item Second item
  \item Third item with sub-items:
    \begin{itemize}
      \item Sub-item 1
      \item Sub-item 2
    \end{itemize}
\end{itemize}

\subsection{Numbered List}
\begin{enumerate}
  \item First step
  \item Second step
  \item Third step
\end{enumerate}

\section{Citations and References}
When you use the citation search feature, references will be automatically managed. For example, you might cite a paper like this: \cite{example2024}. All citations will be automatically added to your \texttt{cite.bib} file.

\section{Conclusion}
This template demonstrates basic \LaTeX{} features including mathematical formulas, tables, and lists. The bibliography system is now set up to work seamlessly with Bibby AI's citation management. Feel free to modify it for your needs!

\bibliographystyle{plain}
\bibliography{cite}

\end{document}\documentclass{article}
\usepackage[utf8]{inputenc}
\usepackage{booktabs}
\usepackage{amsmath}
\usepackage{natbib}

\title{My Document}
\author{Your Name}
\date{\today}

\begin{document}

\maketitle

\section{Introduction}
This is a sample document demonstrating various \LaTeX{} features including mathematical equations, tables, and formatting. You can cite papers using the citation search feature, and they will be automatically added to your \texttt{cite.bib} file.

\section{Mathematical Expressions}
Here are some example equations:

\subsection{Quadratic Formula}
The quadratic formula is given by:
\[x = \frac{-b \pm \sqrt{b^2 - 4ac}}{2a}\]

\subsection{Matrix Example}
A sample matrix:
\[
\begin{pmatrix}
a_{11} & a_{12} & a_{13} \\
a_{21} & a_{22} & a_{23} \\
a_{31} & a_{32} & a_{33}
\end{pmatrix}
\]

\section{Tables}
Below is a sample table showing data:

\begin{table}[h]
\centering
\begin{tabular}{@{}lrrr@{}}
\toprule
\textbf{Category} & \textbf{Value 1} & \textbf{Value 2} & \textbf{Total} \\
\midrule
Item A & 10.5 & 20.3 & 30.8 \\
Item B & 15.2 & 25.1 & 40.3 \\
Item C & 22.1 & 12.7 & 34.8 \\
\bottomrule
\end{tabular}
\caption{Sample Data Table}
\label{tab:sample}
\end{table}

\section{Lists}
\subsection{Bullet Points}
\begin{itemize}
  \item First item
  \item Second item
  \item Third item with sub-items:
    \begin{itemize}
      \item Sub-item 1
      \item Sub-item 2
    \end{itemize}
\end{itemize}

\subsection{Numbered List}
\begin{enumerate}
  \item First step
  \item Second step
  \item Third step
\end{enumerate}

\section{Citations and References}
When you use the citation search feature, references will be automatically managed. For example, you might cite a paper like this: \cite{example2024}. All citations will be automatically added to your \texttt{cite.bib} file.

\section{Conclusion}
This template demonstrates basic \LaTeX{} features including mathematical formulas, tables, and lists. The bibliography system is now set up to work seamlessly with Bibby AI's citation management. Feel free to modify it for your needs!

\bibliographystyle{plain}
\bibliography{cite}

\end{document}\documentclass{article}
\usepackage[utf8]{inputenc}
\usepackage{booktabs}
\usepackage{amsmath}
\usepackage{natbib}

\title{My Document}
\author{Your Name}
\date{\today}

\begin{document}

\maketitle

\section{Introduction}
This is a sample document demonstrating various \LaTeX{} features including mathematical equations, tables, and formatting. You can cite papers using the citation search feature, and they will be automatically added to your \texttt{cite.bib} file.

\section{Mathematical Expressions}
Here are some example equations:

\subsection{Quadratic Formula}
The quadratic formula is given by:
\[x = \frac{-b \pm \sqrt{b^2 - 4ac}}{2a}\]

\subsection{Matrix Example}
A sample matrix:
\[
\begin{pmatrix}
a_{11} & a_{12} & a_{13} \\
a_{21} & a_{22} & a_{23} \\
a_{31} & a_{32} & a_{33}
\end{pmatrix}
\]

\section{Tables}
Below is a sample table showing data:

\begin{table}[h]
\centering
\begin{tabular}{@{}lrrr@{}}
\toprule
\textbf{Category} & \textbf{Value 1} & \textbf{Value 2} & \textbf{Total} \\
\midrule
Item A & 10.5 & 20.3 & 30.8 \\
Item B & 15.2 & 25.1 & 40.3 \\
Item C & 22.1 & 12.7 & 34.8 \\
\bottomrule
\end{tabular}
\caption{Sample Data Table}
\label{tab:sample}
\end{table}

\section{Lists}
\subsection{Bullet Points}
\begin{itemize}
  \item First item
  \item Second item
  \item Third item with sub-items:
    \begin{itemize}
      \item Sub-item 1
      \item Sub-item 2
    \end{itemize}
\end{itemize}

\subsection{Numbered List}
\begin{enumerate}
  \item First step
  \item Second step
  \item Third step
\end{enumerate}

\section{Citations and References}
When you use the citation search feature, references will be automatically managed. For example, you might cite a paper like this: \cite{example2024}. All citations will be automatically added to your \texttt{cite.bib} file.

\section{Conclusion}
This template demonstrates basic \LaTeX{} features including mathematical formulas, tables, and lists. The bibliography system is now set up to work seamlessly with Bibby AI's citation management. Feel free to modify it for your needs!

\bibliographystyle{plain}
\bibliography{cite}

\end{document}\documentclass{article}
\usepackage[utf8]{inputenc}
\usepackage{booktabs}
\usepackage{amsmath}
\usepackage{natbib}

\title{My Document}
\author{Your Name}
\date{\today}

\begin{document}

\maketitle

\section{Introduction}
This is a sample document demonstrating various \LaTeX{} features including mathematical equations, tables, and formatting. You can cite papers using the citation search feature, and they will be automatically added to your \texttt{cite.bib} file.

\section{Mathematical Expressions}
Here are some example equations:

\subsection{Quadratic Formula}
The quadratic formula is given by:
\[x = \frac{-b \pm \sqrt{b^2 - 4ac}}{2a}\]

\subsection{Matrix Example}
A sample matrix:
\[
\begin{pmatrix}
a_{11} & a_{12} & a_{13} \\
a_{21} & a_{22} & a_{23} \\
a_{31} & a_{32} & a_{33}
\end{pmatrix}
\]

\section{Tables}
Below is a sample table showing data:

\begin{table}[h]
\centering
\begin{tabular}{@{}lrrr@{}}
\toprule
\textbf{Category} & \textbf{Value 1} & \textbf{Value 2} & \textbf{Total} \\
\midrule
Item A & 10.5 & 20.3 & 30.8 \\
Item B & 15.2 & 25.1 & 40.3 \\
Item C & 22.1 & 12.7 & 34.8 \\
\bottomrule
\end{tabular}
\caption{Sample Data Table}
\label{tab:sample}
\end{table}

\section{Lists}
\subsection{Bullet Points}
\begin{itemize}
  \item First item
  \item Second item
  \item Third item with sub-items:
    \begin{itemize}
      \item Sub-item 1
      \item Sub-item 2
    \end{itemize}
\end{itemize}

\subsection{Numbered List}
\begin{enumerate}
  \item First step
  \item Second step
  \item Third step
\end{enumerate}

\section{Citations and References}
When you use the citation search feature, references will be automatically managed. For example, you might cite a paper like this: \cite{example2024}. All citations will be automatically added to your \texttt{cite.bib} file.

\section{Conclusion}
This template demonstrates basic \LaTeX{} features including mathematical formulas, tables, and lists. The bibliography system is now set up to work seamlessly with Bibby AI's citation management. Feel free to modify it for your needs!

\bibliographystyle{plain}
\bibliography{cite}

\end{document}\documentclass{article}
\usepackage[utf8]{inputenc}
\usepackage{booktabs}
\usepackage{amsmath}
\usepackage{natbib}

\title{My Document}
\author{Your Name}
\date{\today}

\begin{document}

\maketitle

\section{Introduction}
This is a sample document demonstrating various \LaTeX{} features including mathematical equations, tables, and formatting. You can cite papers using the citation search feature, and they will be automatically added to your \texttt{cite.bib} file.

\section{Mathematical Expressions}
Here are some example equations:

\subsection{Quadratic Formula}
The quadratic formula is given by:
\[x = \frac{-b \pm \sqrt{b^2 - 4ac}}{2a}\]

\subsection{Matrix Example}
A sample matrix:
\[
\begin{pmatrix}
a_{11} & a_{12} & a_{13} \\
a_{21} & a_{22} & a_{23} \\
a_{31} & a_{32} & a_{33}
\end{pmatrix}
\]

\section{Tables}
Below is a sample table showing data:

\begin{table}[h]
\centering
\begin{tabular}{@{}lrrr@{}}
\toprule
\textbf{Category} & \textbf{Value 1} & \textbf{Value 2} & \textbf{Total} \\
\midrule
Item A & 10.5 & 20.3 & 30.8 \\
Item B & 15.2 & 25.1 & 40.3 \\
Item C & 22.1 & 12.7 & 34.8 \\
\bottomrule
\end{tabular}
\caption{Sample Data Table}
\label{tab:sample}
\end{table}

\section{Lists}
\subsection{Bullet Points}
\begin{itemize}
  \item First item
  \item Second item
  \item Third item with sub-items:
    \begin{itemize}
      \item Sub-item 1
      \item Sub-item 2
    \end{itemize}
\end{itemize}

\subsection{Numbered List}
\begin{enumerate}
  \item First step
  \item Second step
  \item Third step
\end{enumerate}

\section{Citations and References}
When you use the citation search feature, references will be automatically managed. For example, you might cite a paper like this: \cite{example2024}. All citations will be automatically added to your \texttt{cite.bib} file.

\section{Conclusion}
This template demonstrates basic \LaTeX{} features including mathematical formulas, tables, and lists. The bibliography system is now set up to work seamlessly with Bibby AI's citation management. Feel free to modify it for your needs!

\bibliographystyle{plain}
\bibliography{cite}

\end{document}\documentclass{article}
\usepackage[utf8]{inputenc}
\usepackage{booktabs}
\usepackage{amsmath}
\usepackage{natbib}

\title{My Document}
\author{Your Name}
\date{\today}

\begin{document}

\maketitle

\section{Introduction}
This is a sample document demonstrating various \LaTeX{} features including mathematical equations, tables, and formatting. You can cite papers using the citation search feature, and they will be automatically added to your \texttt{cite.bib} file.

\section{Mathematical Expressions}
Here are some example equations:

\subsection{Quadratic Formula}
The quadratic formula is given by:
\[x = \frac{-b \pm \sqrt{b^2 - 4ac}}{2a}\]

\subsection{Matrix Example}
A sample matrix:
\[
\begin{pmatrix}
a_{11} & a_{12} & a_{13} \\
a_{21} & a_{22} & a_{23} \\
a_{31} & a_{32} & a_{33}
\end{pmatrix}
\]

\section{Tables}
Below is a sample table showing data:

\begin{table}[h]
\centering
\begin{tabular}{@{}lrrr@{}}
\toprule
\textbf{Category} & \textbf{Value 1} & \textbf{Value 2} & \textbf{Total} \\
\midrule
Item A & 10.5 & 20.3 & 30.8 \\
Item B & 15.2 & 25.1 & 40.3 \\
Item C & 22.1 & 12.7 & 34.8 \\
\bottomrule
\end{tabular}
\caption{Sample Data Table}
\label{tab:sample}
\end{table}

\section{Lists}
\subsection{Bullet Points}
\begin{itemize}
  \item First item
  \item Second item
  \item Third item with sub-items:
    \begin{itemize}
      \item Sub-item 1
      \item Sub-item 2
    \end{itemize}
\end{itemize}

\subsection{Numbered List}
\begin{enumerate}
  \item First step
  \item Second step
  \item Third step
\end{enumerate}

\section{Citations and References}
When you use the citation search feature, references will be automatically managed. For example, you might cite a paper like this: \cite{example2024}. All citations will be automatically added to your \texttt{cite.bib} file.

\section{Conclusion}
This template demonstrates basic \LaTeX{} features including mathematical formulas, tables, and lists. The bibliography system is now set up to work seamlessly with Bibby AI's citation management. Feel free to modify it for your needs!

\bibliographystyle{plain}
\bibliography{cite}

\end{document}\documentclass{article}
\usepackage[utf8]{inputenc}
\usepackage{booktabs}
\usepackage{amsmath}
\usepackage{natbib}

\title{My Document}
\author{Your Name}
\date{\today}

\begin{document}

\maketitle

\section{Introduction}
This is a sample document demonstrating various \LaTeX{} features including mathematical equations, tables, and formatting. You can cite papers using the citation search feature, and they will be automatically added to your \texttt{cite.bib} file.

\section{Mathematical Expressions}
Here are some example equations:

\subsection{Quadratic Formula}
The quadratic formula is given by:
\[x = \frac{-b \pm \sqrt{b^2 - 4ac}}{2a}\]

\subsection{Matrix Example}
A sample matrix:
\[
\begin{pmatrix}
a_{11} & a_{12} & a_{13} \\
a_{21} & a_{22} & a_{23} \\
a_{31} & a_{32} & a_{33}
\end{pmatrix}
\]

\section{Tables}
Below is a sample table showing data:

\begin{table}[h]
\centering
\begin{tabular}{@{}lrrr@{}}
\toprule
\textbf{Category} & \textbf{Value 1} & \textbf{Value 2} & \textbf{Total} \\
\midrule
Item A & 10.5 & 20.3 & 30.8 \\
Item B & 15.2 & 25.1 & 40.3 \\
Item C & 22.1 & 12.7 & 34.8 \\
\bottomrule
\end{tabular}
\caption{Sample Data Table}
\label{tab:sample}
\end{table}

\section{Lists}
\subsection{Bullet Points}
\begin{itemize}
  \item First item
  \item Second item
  \item Third item with sub-items:
    \begin{itemize}
      \item Sub-item 1
      \item Sub-item 2
    \end{itemize}
\end{itemize}

\subsection{Numbered List}
\begin{enumerate}
  \item First step
  \item Second step
  \item Third step
\end{enumerate}

\section{Citations and References}
When you use the citation search feature, references will be automatically managed. For example, you might cite a paper like this: \cite{example2024}. All citations will be automatically added to your \texttt{cite.bib} file.

\section{Conclusion}
This template demonstrates basic \LaTeX{} features including mathematical formulas, tables, and lists. The bibliography system is now set up to work seamlessly with Bibby AI's citation management. Feel free to modify it for your needs!

\bibliographystyle{plain}
\bibliography{cite}

\end{document}\documentclass{article}
\usepackage[utf8]{inputenc}
\usepackage{booktabs}
\usepackage{amsmath}
\usepackage{natbib}

\title{My Document}
\author{Your Name}
\date{\today}

\begin{document}

\maketitle

\section{Introduction}
This is a sample document demonstrating various \LaTeX{} features including mathematical equations, tables, and formatting. You can cite papers using the citation search feature, and they will be automatically added to your \texttt{cite.bib} file.

\section{Mathematical Expressions}
Here are some example equations:

\subsection{Quadratic Formula}
The quadratic formula is given by:
\[x = \frac{-b \pm \sqrt{b^2 - 4ac}}{2a}\]

\subsection{Matrix Example}
A sample matrix:
\[
\begin{pmatrix}
a_{11} & a_{12} & a_{13} \\
a_{21} & a_{22} & a_{23} \\
a_{31} & a_{32} & a_{33}
\end{pmatrix}
\]

\section{Tables}
Below is a sample table showing data:

\begin{table}[h]
\centering
\begin{tabular}{@{}lrrr@{}}
\toprule
\textbf{Category} & \textbf{Value 1} & \textbf{Value 2} & \textbf{Total} \\
\midrule
Item A & 10.5 & 20.3 & 30.8 \\
Item B & 15.2 & 25.1 & 40.3 \\
Item C & 22.1 & 12.7 & 34.8 \\
\bottomrule
\end{tabular}
\caption{Sample Data Table}
\label{tab:sample}
\end{table}

\section{Lists}
\subsection{Bullet Points}
\begin{itemize}
  \item First item
  \item Second item
  \item Third item with sub-items:
    \begin{itemize}
      \item Sub-item 1
      \item Sub-item 2
    \end{itemize}
\end{itemize}

\subsection{Numbered List}
\begin{enumerate}
  \item First step
  \item Second step
  \item Third step
\end{enumerate}

\section{Citations and References}
When you use the citation search feature, references will be automatically managed. For example, you might cite a paper like this: \cite{example2024}. All citations will be automatically added to your \texttt{cite.bib} file.

\section{Conclusion}
This template demonstrates basic \LaTeX{} features including mathematical formulas, tables, and lists. The bibliography system is now set up to work seamlessly with Bibby AI's citation management. Feel free to modify it for your needs!

\bibliographystyle{plain}
\bibliography{cite}

\end{document}\documentclass{article}
\usepackage[utf8]{inputenc}
\usepackage{booktabs}
\usepackage{amsmath}
\usepackage{natbib}

\title{My Document}
\author{Your Name}
\date{\today}

\begin{document}

\maketitle

\section{Introduction}
This is a sample document demonstrating various \LaTeX{} features including mathematical equations, tables, and formatting. You can cite papers using the citation search feature, and they will be automatically added to your \texttt{cite.bib} file.

\section{Mathematical Expressions}
Here are some example equations:

\subsection{Quadratic Formula}
The quadratic formula is given by:
\[x = \frac{-b \pm \sqrt{b^2 - 4ac}}{2a}\]

\subsection{Matrix Example}
A sample matrix:
\[
\begin{pmatrix}
a_{11} & a_{12} & a_{13} \\
a_{21} & a_{22} & a_{23} \\
a_{31} & a_{32} & a_{33}
\end{pmatrix}
\]

\section{Tables}
Below is a sample table showing data:

\begin{table}[h]
\centering
\begin{tabular}{@{}lrrr@{}}
\toprule
\textbf{Category} & \textbf{Value 1} & \textbf{Value 2} & \textbf{Total} \\
\midrule
Item A & 10.5 & 20.3 & 30.8 \\
Item B & 15.2 & 25.1 & 40.3 \\
Item C & 22.1 & 12.7 & 34.8 \\
\bottomrule
\end{tabular}
\caption{Sample Data Table}
\label{tab:sample}
\end{table}

\section{Lists}
\subsection{Bullet Points}
\begin{itemize}
  \item First item
  \item Second item
  \item Third item with sub-items:
    \begin{itemize}
      \item Sub-item 1
      \item Sub-item 2
    \end{itemize}
\end{itemize}

\subsection{Numbered List}
\begin{enumerate}
  \item First step
  \item Second step
  \item Third step
\end{enumerate}

\section{Citations and References}
When you use the citation search feature, references will be automatically managed. For example, you might cite a paper like this: \cite{example2024}. All citations will be automatically added to your \texttt{cite.bib} file.

\section{Conclusion}
This template demonstrates basic \LaTeX{} features including mathematical formulas, tables, and lists. The bibliography system is now set up to work seamlessly with Bibby AI's citation management. Feel free to modify it for your needs!

\bibliographystyle{plain}
\bibliography{cite}

\end{document}\documentclass{article}
\usepackage[utf8]{inputenc}
\usepackage{booktabs}
\usepackage{amsmath}
\usepackage{natbib}

\title{My Document}
\author{Your Name}
\date{\today}

\begin{document}

\maketitle

\section{Introduction}
This is a sample document demonstrating various \LaTeX{} features including mathematical equations, tables, and formatting. You can cite papers using the citation search feature, and they will be automatically added to your \texttt{cite.bib} file.

\section{Mathematical Expressions}
Here are some example equations:

\subsection{Quadratic Formula}
The quadratic formula is given by:
\[x = \frac{-b \pm \sqrt{b^2 - 4ac}}{2a}\]

\subsection{Matrix Example}
A sample matrix:
\[
\begin{pmatrix}
a_{11} & a_{12} & a_{13} \\
a_{21} & a_{22} & a_{23} \\
a_{31} & a_{32} & a_{33}
\end{pmatrix}
\]

\section{Tables}
Below is a sample table showing data:

\begin{table}[h]
\centering
\begin{tabular}{@{}lrrr@{}}
\toprule
\textbf{Category} & \textbf{Value 1} & \textbf{Value 2} & \textbf{Total} \\
\midrule
Item A & 10.5 & 20.3 & 30.8 \\
Item B & 15.2 & 25.1 & 40.3 \\
Item C & 22.1 & 12.7 & 34.8 \\
\bottomrule
\end{tabular}
\caption{Sample Data Table}
\label{tab:sample}
\end{table}

\section{Lists}
\subsection{Bullet Points}
\begin{itemize}
  \item First item
  \item Second item
  \item Third item with sub-items:
    \begin{itemize}
      \item Sub-item 1
      \item Sub-item 2
    \end{itemize}
\end{itemize}

\subsection{Numbered List}
\begin{enumerate}
  \item First step
  \item Second step
  \item Third step
\end{enumerate}

\section{Citations and References}
When you use the citation search feature, references will be automatically managed. For example, you might cite a paper like this: \cite{example2024}. All citations will be automatically added to your \texttt{cite.bib} file.

\section{Conclusion}
This template demonstrates basic \LaTeX{} features including mathematical formulas, tables, and lists. The bibliography system is now set up to work seamlessly with Bibby AI's citation management. Feel free to modify it for your needs!

\bibliographystyle{plain}
\bibliography{cite}

\end{document}\documentclass{article}
\usepackage[utf8]{inputenc}
\usepackage{booktabs}
\usepackage{amsmath}
\usepackage{natbib}

\title{My Document}
\author{Your Name}
\date{\today}

\begin{document}

\maketitle

\section{Introduction}
This is a sample document demonstrating various \LaTeX{} features including mathematical equations, tables, and formatting. You can cite papers using the citation search feature, and they will be automatically added to your \texttt{cite.bib} file.

\section{Mathematical Expressions}
Here are some example equations:

\subsection{Quadratic Formula}
The quadratic formula is given by:
\[x = \frac{-b \pm \sqrt{b^2 - 4ac}}{2a}\]

\subsection{Matrix Example}
A sample matrix:
\[
\begin{pmatrix}
a_{11} & a_{12} & a_{13} \\
a_{21} & a_{22} & a_{23} \\
a_{31} & a_{32} & a_{33}
\end{pmatrix}
\]

\section{Tables}
Below is a sample table showing data:

\begin{table}[h]
\centering
\begin{tabular}{@{}lrrr@{}}
\toprule
\textbf{Category} & \textbf{Value 1} & \textbf{Value 2} & \textbf{Total} \\
\midrule
Item A & 10.5 & 20.3 & 30.8 \\
Item B & 15.2 & 25.1 & 40.3 \\
Item C & 22.1 & 12.7 & 34.8 \\
\bottomrule
\end{tabular}
\caption{Sample Data Table}
\label{tab:sample}
\end{table}

\section{Lists}
\subsection{Bullet Points}
\begin{itemize}
  \item First item
  \item Second item
  \item Third item with sub-items:
    \begin{itemize}
      \item Sub-item 1
      \item Sub-item 2
    \end{itemize}
\end{itemize}

\subsection{Numbered List}
\begin{enumerate}
  \item First step
  \item Second step
  \item Third step
\end{enumerate}

\section{Citations and References}
When you use the citation search feature, references will be automatically managed. For example, you might cite a paper like this: \cite{example2024}. All citations will be automatically added to your \texttt{cite.bib} file.

\section{Conclusion}
This template demonstrates basic \LaTeX{} features including mathematical formulas, tables, and lists. The bibliography system is now set up to work seamlessly with Bibby AI's citation management. Feel free to modify it for your needs!

\bibliographystyle{plain}
\bibliography{cite}

\end{document}\documentclass{article}
\usepackage[utf8]{inputenc}
\usepackage{booktabs}
\usepackage{amsmath}
\usepackage{natbib}

\title{My Document}
\author{Your Name}
\date{\today}

\begin{document}

\maketitle

\section{Introduction}
This is a sample document demonstrating various \LaTeX{} features including mathematical equations, tables, and formatting. You can cite papers using the citation search feature, and they will be automatically added to your \texttt{cite.bib} file.

\section{Mathematical Expressions}
Here are some example equations:

\subsection{Quadratic Formula}
The quadratic formula is given by:
\[x = \frac{-b \pm \sqrt{b^2 - 4ac}}{2a}\]

\subsection{Matrix Example}
A sample matrix:
\[
\begin{pmatrix}
a_{11} & a_{12} & a_{13} \\
a_{21} & a_{22} & a_{23} \\
a_{31} & a_{32} & a_{33}
\end{pmatrix}
\]

\section{Tables}
Below is a sample table showing data:

\begin{table}[h]
\centering
\begin{tabular}{@{}lrrr@{}}
\toprule
\textbf{Category} & \textbf{Value 1} & \textbf{Value 2} & \textbf{Total} \\
\midrule
Item A & 10.5 & 20.3 & 30.8 \\
Item B & 15.2 & 25.1 & 40.3 \\
Item C & 22.1 & 12.7 & 34.8 \\
\bottomrule
\end{tabular}
\caption{Sample Data Table}
\label{tab:sample}
\end{table}

\section{Lists}
\subsection{Bullet Points}
\begin{itemize}
  \item First item
  \item Second item
  \item Third item with sub-items:
    \begin{itemize}
      \item Sub-item 1
      \item Sub-item 2
    \end{itemize}
\end{itemize}

\subsection{Numbered List}
\begin{enumerate}
  \item First step
  \item Second step
  \item Third step
\end{enumerate}

\section{Citations and References}
When you use the citation search feature, references will be automatically managed. For example, you might cite a paper like this: \cite{example2024}. All citations will be automatically added to your \texttt{cite.bib} file.

\section{Conclusion}
This template demonstrates basic \LaTeX{} features including mathematical formulas, tables, and lists. The bibliography system is now set up to work seamlessly with Bibby AI's citation management. Feel free to modify it for your needs!

\bibliographystyle{plain}
\bibliography{cite}

\end{document}\documentclass{article}
\usepackage[utf8]{inputenc}
\usepackage{booktabs}
\usepackage{amsmath}
\usepackage{natbib}

\title{My Document}
\author{Your Name}
\date{\today}

\begin{document}

\maketitle

\section{Introduction}
This is a sample document demonstrating various \LaTeX{} features including mathematical equations, tables, and formatting. You can cite papers using the citation search feature, and they will be automatically added to your \texttt{cite.bib} file.

\section{Mathematical Expressions}
Here are some example equations:

\subsection{Quadratic Formula}
The quadratic formula is given by:
\[x = \frac{-b \pm \sqrt{b^2 - 4ac}}{2a}\]

\subsection{Matrix Example}
A sample matrix:
\[
\begin{pmatrix}
a_{11} & a_{12} & a_{13} \\
a_{21} & a_{22} & a_{23} \\
a_{31} & a_{32} & a_{33}
\end{pmatrix}
\]

\section{Tables}
Below is a sample table showing data:

\begin{table}[h]
\centering
\begin{tabular}{@{}lrrr@{}}
\toprule
\textbf{Category} & \textbf{Value 1} & \textbf{Value 2} & \textbf{Total} \\
\midrule
Item A & 10.5 & 20.3 & 30.8 \\
Item B & 15.2 & 25.1 & 40.3 \\
Item C & 22.1 & 12.7 & 34.8 \\
\bottomrule
\end{tabular}
\caption{Sample Data Table}
\label{tab:sample}
\end{table}

\section{Lists}
\subsection{Bullet Points}
\begin{itemize}
  \item First item
  \item Second item
  \item Third item with sub-items:
    \begin{itemize}
      \item Sub-item 1
      \item Sub-item 2
    \end{itemize}
\end{itemize}

\subsection{Numbered List}
\begin{enumerate}
  \item First step
  \item Second step
  \item Third step
\end{enumerate}

\section{Citations and References}
When you use the citation search feature, references will be automatically managed. For example, you might cite a paper like this: \cite{example2024}. All citations will be automatically added to your \texttt{cite.bib} file.

\section{Conclusion}
This template demonstrates basic \LaTeX{} features including mathematical formulas, tables, and lists. The bibliography system is now set up to work seamlessly with Bibby AI's citation management. Feel free to modify it for your needs!

\bibliographystyle{plain}
\bibliography{cite}

\end{document}\documentclass{article}
\usepackage[utf8]{inputenc}
\usepackage{booktabs}
\usepackage{amsmath}
\usepackage{natbib}

\title{My Document}
\author{Your Name}
\date{\today}

\begin{document}

\maketitle

\section{Introduction}
This is a sample document demonstrating various \LaTeX{} features including mathematical equations, tables, and formatting. You can cite papers using the citation search feature, and they will be automatically added to your \texttt{cite.bib} file.

\section{Mathematical Expressions}
Here are some example equations:

\subsection{Quadratic Formula}
The quadratic formula is given by:
\[x = \frac{-b \pm \sqrt{b^2 - 4ac}}{2a}\]

\subsection{Matrix Example}
A sample matrix:
\[
\begin{pmatrix}
a_{11} & a_{12} & a_{13} \\
a_{21} & a_{22} & a_{23} \\
a_{31} & a_{32} & a_{33}
\end{pmatrix}
\]

\section{Tables}
Below is a sample table showing data:

\begin{table}[h]
\centering
\begin{tabular}{@{}lrrr@{}}
\toprule
\textbf{Category} & \textbf{Value 1} & \textbf{Value 2} & \textbf{Total} \\
\midrule
Item A & 10.5 & 20.3 & 30.8 \\
Item B & 15.2 & 25.1 & 40.3 \\
Item C & 22.1 & 12.7 & 34.8 \\
\bottomrule
\end{tabular}
\caption{Sample Data Table}
\label{tab:sample}
\end{table}

\section{Lists}
\subsection{Bullet Points}
\begin{itemize}
  \item First item
  \item Second item
  \item Third item with sub-items:
    \begin{itemize}
      \item Sub-item 1
      \item Sub-item 2
    \end{itemize}
\end{itemize}

\subsection{Numbered List}
\begin{enumerate}
  \item First step
  \item Second step
  \item Third step
\end{enumerate}

\section{Citations and References}
When you use the citation search feature, references will be automatically managed. For example, you might cite a paper like this: \cite{example2024}. All citations will be automatically added to your \texttt{cite.bib} file.

\section{Conclusion}
This template demonstrates basic \LaTeX{} features including mathematical formulas, tables, and lists. The bibliography system is now set up to work seamlessly with Bibby AI's citation management. Feel free to modify it for your needs!

\bibliographystyle{plain}
\bibliography{cite}

\end{document}\documentclass{article}
\usepackage[utf8]{inputenc}
\usepackage{booktabs}
\usepackage{amsmath}
\usepackage{natbib}

\title{My Document}
\author{Your Name}
\date{\today}

\begin{document}

\maketitle

\section{Introduction}
This is a sample document demonstrating various \LaTeX{} features including mathematical equations, tables, and formatting. You can cite papers using the citation search feature, and they will be automatically added to your \texttt{cite.bib} file.

\section{Mathematical Expressions}
Here are some example equations:

\subsection{Quadratic Formula}
The quadratic formula is given by:
\[x = \frac{-b \pm \sqrt{b^2 - 4ac}}{2a}\]

\subsection{Matrix Example}
A sample matrix:
\[
\begin{pmatrix}
a_{11} & a_{12} & a_{13} \\
a_{21} & a_{22} & a_{23} \\
a_{31} & a_{32} & a_{33}
\end{pmatrix}
\]

\section{Tables}
Below is a sample table showing data:

\begin{table}[h]
\centering
\begin{tabular}{@{}lrrr@{}}
\toprule
\textbf{Category} & \textbf{Value 1} & \textbf{Value 2} & \textbf{Total} \\
\midrule
Item A & 10.5 & 20.3 & 30.8 \\
Item B & 15.2 & 25.1 & 40.3 \\
Item C & 22.1 & 12.7 & 34.8 \\
\bottomrule
\end{tabular}
\caption{Sample Data Table}
\label{tab:sample}
\end{table}

\section{Lists}
\subsection{Bullet Points}
\begin{itemize}
  \item First item
  \item Second item
  \item Third item with sub-items:
    \begin{itemize}
      \item Sub-item 1
      \item Sub-item 2
    \end{itemize}
\end{itemize}

\subsection{Numbered List}
\begin{enumerate}
  \item First step
  \item Second step
  \item Third step
\end{enumerate}

\section{Citations and References}
When you use the citation search feature, references will be automatically managed. For example, you might cite a paper like this: \cite{example2024}. All citations will be automatically added to your \texttt{cite.bib} file.

\section{Conclusion}
This template demonstrates basic \LaTeX{} features including mathematical formulas, tables, and lists. The bibliography system is now set up to work seamlessly with Bibby AI's citation management. Feel free to modify it for your needs!

\bibliographystyle{plain}
\bibliography{cite}

\end{document}\documentclass{article}
\usepackage[utf8]{inputenc}
\usepackage{booktabs}
\usepackage{amsmath}
\usepackage{natbib}

\title{My Document}
\author{Your Name}
\date{\today}

\begin{document}

\maketitle

\section{Introduction}
This is a sample document demonstrating various \LaTeX{} features including mathematical equations, tables, and formatting. You can cite papers using the citation search feature, and they will be automatically added to your \texttt{cite.bib} file.

\section{Mathematical Expressions}
Here are some example equations:

\subsection{Quadratic Formula}
The quadratic formula is given by:
\[x = \frac{-b \pm \sqrt{b^2 - 4ac}}{2a}\]

\subsection{Matrix Example}
A sample matrix:
\[
\begin{pmatrix}
a_{11} & a_{12} & a_{13} \\
a_{21} & a_{22} & a_{23} \\
a_{31} & a_{32} & a_{33}
\end{pmatrix}
\]

\section{Tables}
Below is a sample table showing data:

\begin{table}[h]
\centering
\begin{tabular}{@{}lrrr@{}}
\toprule
\textbf{Category} & \textbf{Value 1} & \textbf{Value 2} & \textbf{Total} \\
\midrule
Item A & 10.5 & 20.3 & 30.8 \\
Item B & 15.2 & 25.1 & 40.3 \\
Item C & 22.1 & 12.7 & 34.8 \\
\bottomrule
\end{tabular}
\caption{Sample Data Table}
\label{tab:sample}
\end{table}

\section{Lists}
\subsection{Bullet Points}
\begin{itemize}
  \item First item
  \item Second item
  \item Third item with sub-items:
    \begin{itemize}
      \item Sub-item 1
      \item Sub-item 2
    \end{itemize}
\end{itemize}

\subsection{Numbered List}
\begin{enumerate}
  \item First step
  \item Second step
  \item Third step
\end{enumerate}

\section{Citations and References}
When you use the citation search feature, references will be automatically managed. For example, you might cite a paper like this: \cite{example2024}. All citations will be automatically added to your \texttt{cite.bib} file.

\section{Conclusion}
This template demonstrates basic \LaTeX{} features including mathematical formulas, tables, and lists. The bibliography system is now set up to work seamlessly with Bibby AI's citation management. Feel free to modify it for your needs!

\bibliographystyle{plain}
\bibliography{cite}

\end{document}\documentclass{article}
\usepackage[utf8]{inputenc}
\usepackage{booktabs}
\usepackage{amsmath}
\usepackage{natbib}

\title{My Document}
\author{Your Name}
\date{\today}

\begin{document}

\maketitle

\section{Introduction}
This is a sample document demonstrating various \LaTeX{} features including mathematical equations, tables, and formatting. You can cite papers using the citation search feature, and they will be automatically added to your \texttt{cite.bib} file.

\section{Mathematical Expressions}
Here are some example equations:

\subsection{Quadratic Formula}
The quadratic formula is given by:
\[x = \frac{-b \pm \sqrt{b^2 - 4ac}}{2a}\]

\subsection{Matrix Example}
A sample matrix:
\[
\begin{pmatrix}
a_{11} & a_{12} & a_{13} \\
a_{21} & a_{22} & a_{23} \\
a_{31} & a_{32} & a_{33}
\end{pmatrix}
\]

\section{Tables}
Below is a sample table showing data:

\begin{table}[h]
\centering
\begin{tabular}{@{}lrrr@{}}
\toprule
\textbf{Category} & \textbf{Value 1} & \textbf{Value 2} & \textbf{Total} \\
\midrule
Item A & 10.5 & 20.3 & 30.8 \\
Item B & 15.2 & 25.1 & 40.3 \\
Item C & 22.1 & 12.7 & 34.8 \\
\bottomrule
\end{tabular}
\caption{Sample Data Table}
\label{tab:sample}
\end{table}

\section{Lists}
\subsection{Bullet Points}
\begin{itemize}
  \item First item
  \item Second item
  \item Third item with sub-items:
    \begin{itemize}
      \item Sub-item 1
      \item Sub-item 2
    \end{itemize}
\end{itemize}

\subsection{Numbered List}
\begin{enumerate}
  \item First step
  \item Second step
  \item Third step
\end{enumerate}

\section{Citations and References}
When you use the citation search feature, references will be automatically managed. For example, you might cite a paper like this: \cite{example2024}. All citations will be automatically added to your \texttt{cite.bib} file.

\section{Conclusion}
This template demonstrates basic \LaTeX{} features including mathematical formulas, tables, and lists. The bibliography system is now set up to work seamlessly with Bibby AI's citation management. Feel free to modify it for your needs!

\bibliographystyle{plain}
\bibliography{cite}

\end{document}\documentclass{article}
\usepackage[utf8]{inputenc}
\usepackage{booktabs}
\usepackage{amsmath}
\usepackage{natbib}

\title{My Document}
\author{Your Name}
\date{\today}

\begin{document}

\maketitle

\section{Introduction}
This is a sample document demonstrating various \LaTeX{} features including mathematical equations, tables, and formatting. You can cite papers using the citation search feature, and they will be automatically added to your \texttt{cite.bib} file.

\section{Mathematical Expressions}
Here are some example equations:

\subsection{Quadratic Formula}
The quadratic formula is given by:
\[x = \frac{-b \pm \sqrt{b^2 - 4ac}}{2a}\]

\subsection{Matrix Example}
A sample matrix:
\[
\begin{pmatrix}
a_{11} & a_{12} & a_{13} \\
a_{21} & a_{22} & a_{23} \\
a_{31} & a_{32} & a_{33}
\end{pmatrix}
\]

\section{Tables}
Below is a sample table showing data:

\begin{table}[h]
\centering
\begin{tabular}{@{}lrrr@{}}
\toprule
\textbf{Category} & \textbf{Value 1} & \textbf{Value 2} & \textbf{Total} \\
\midrule
Item A & 10.5 & 20.3 & 30.8 \\
Item B & 15.2 & 25.1 & 40.3 \\
Item C & 22.1 & 12.7 & 34.8 \\
\bottomrule
\end{tabular}
\caption{Sample Data Table}
\label{tab:sample}
\end{table}

\section{Lists}
\subsection{Bullet Points}
\begin{itemize}
  \item First item
  \item Second item
  \item Third item with sub-items:
    \begin{itemize}
      \item Sub-item 1
      \item Sub-item 2
    \end{itemize}
\end{itemize}

\subsection{Numbered List}
\begin{enumerate}
  \item First step
  \item Second step
  \item Third step
\end{enumerate}

\section{Citations and References}
When you use the citation search feature, references will be automatically managed. For example, you might cite a paper like this: \cite{example2024}. All citations will be automatically added to your \texttt{cite.bib} file.

\section{Conclusion}
This template demonstrates basic \LaTeX{} features including mathematical formulas, tables, and lists. The bibliography system is now set up to work seamlessly with Bibby AI's citation management. Feel free to modify it for your needs!

\bibliographystyle{plain}
\bibliography{cite}

\end{document}\documentclass{article}
\usepackage[utf8]{inputenc}
\usepackage{booktabs}
\usepackage{amsmath}
\usepackage{natbib}

\title{My Document}
\author{Your Name}
\date{\today}

\begin{document}

\maketitle

\section{Introduction}
This is a sample document demonstrating various \LaTeX{} features including mathematical equations, tables, and formatting. You can cite papers using the citation search feature, and they will be automatically added to your \texttt{cite.bib} file.

\section{Mathematical Expressions}
Here are some example equations:

\subsection{Quadratic Formula}
The quadratic formula is given by:
\[x = \frac{-b \pm \sqrt{b^2 - 4ac}}{2a}\]

\subsection{Matrix Example}
A sample matrix:
\[
\begin{pmatrix}
a_{11} & a_{12} & a_{13} \\
a_{21} & a_{22} & a_{23} \\
a_{31} & a_{32} & a_{33}
\end{pmatrix}
\]

\section{Tables}
Below is a sample table showing data:

\begin{table}[h]
\centering
\begin{tabular}{@{}lrrr@{}}
\toprule
\textbf{Category} & \textbf{Value 1} & \textbf{Value 2} & \textbf{Total} \\
\midrule
Item A & 10.5 & 20.3 & 30.8 \\
Item B & 15.2 & 25.1 & 40.3 \\
Item C & 22.1 & 12.7 & 34.8 \\
\bottomrule
\end{tabular}
\caption{Sample Data Table}
\label{tab:sample}
\end{table}

\section{Lists}
\subsection{Bullet Points}
\begin{itemize}
  \item First item
  \item Second item
  \item Third item with sub-items:
    \begin{itemize}
      \item Sub-item 1
      \item Sub-item 2
    \end{itemize}
\end{itemize}

\subsection{Numbered List}
\begin{enumerate}
  \item First step
  \item Second step
  \item Third step
\end{enumerate}

\section{Citations and References}
When you use the citation search feature, references will be automatically managed. For example, you might cite a paper like this: \cite{example2024}. All citations will be automatically added to your \texttt{cite.bib} file.

\section{Conclusion}
This template demonstrates basic \LaTeX{} features including mathematical formulas, tables, and lists. The bibliography system is now set up to work seamlessly with Bibby AI's citation management. Feel free to modify it for your needs!

\bibliographystyle{plain}
\bibliography{cite}

\end{document}\documentclass{article}
\usepackage[utf8]{inputenc}
\usepackage{booktabs}
\usepackage{amsmath}
\usepackage{natbib}

\title{My Document}
\author{Your Name}
\date{\today}

\begin{document}

\maketitle

\section{Introduction}
This is a sample document demonstrating various \LaTeX{} features including mathematical equations, tables, and formatting. You can cite papers using the citation search feature, and they will be automatically added to your \texttt{cite.bib} file.

\section{Mathematical Expressions}
Here are some example equations:

\subsection{Quadratic Formula}
The quadratic formula is given by:
\[x = \frac{-b \pm \sqrt{b^2 - 4ac}}{2a}\]

\subsection{Matrix Example}
A sample matrix:
\[
\begin{pmatrix}
a_{11} & a_{12} & a_{13} \\
a_{21} & a_{22} & a_{23} \\
a_{31} & a_{32} & a_{33}
\end{pmatrix}
\]

\section{Tables}
Below is a sample table showing data:

\begin{table}[h]
\centering
\begin{tabular}{@{}lrrr@{}}
\toprule
\textbf{Category} & \textbf{Value 1} & \textbf{Value 2} & \textbf{Total} \\
\midrule
Item A & 10.5 & 20.3 & 30.8 \\
Item B & 15.2 & 25.1 & 40.3 \\
Item C & 22.1 & 12.7 & 34.8 \\
\bottomrule
\end{tabular}
\caption{Sample Data Table}
\label{tab:sample}
\end{table}

\section{Lists}
\subsection{Bullet Points}
\begin{itemize}
  \item First item
  \item Second item
  \item Third item with sub-items:
    \begin{itemize}
      \item Sub-item 1
      \item Sub-item 2
    \end{itemize}
\end{itemize}

\subsection{Numbered List}
\begin{enumerate}
  \item First step
  \item Second step
  \item Third step
\end{enumerate}

\section{Citations and References}
When you use the citation search feature, references will be automatically managed. For example, you might cite a paper like this: \cite{example2024}. All citations will be automatically added to your \texttt{cite.bib} file.

\section{Conclusion}
This template demonstrates basic \LaTeX{} features including mathematical formulas, tables, and lists. The bibliography system is now set up to work seamlessly with Bibby AI's citation management. Feel free to modify it for your needs!

\bibliographystyle{plain}
\bibliography{cite}

\end{document}\documentclass{article}
\usepackage[utf8]{inputenc}
\usepackage{booktabs}
\usepackage{amsmath}
\usepackage{natbib}

\title{My Document}
\author{Your Name}
\date{\today}

\begin{document}

\maketitle

\section{Introduction}
This is a sample document demonstrating various \LaTeX{} features including mathematical equations, tables, and formatting. You can cite papers using the citation search feature, and they will be automatically added to your \texttt{cite.bib} file.

\section{Mathematical Expressions}
Here are some example equations:

\subsection{Quadratic Formula}
The quadratic formula is given by:
\[x = \frac{-b \pm \sqrt{b^2 - 4ac}}{2a}\]

\subsection{Matrix Example}
A sample matrix:
\[
\begin{pmatrix}
a_{11} & a_{12} & a_{13} \\
a_{21} & a_{22} & a_{23} \\
a_{31} & a_{32} & a_{33}
\end{pmatrix}
\]

\section{Tables}
Below is a sample table showing data:

\begin{table}[h]
\centering
\begin{tabular}{@{}lrrr@{}}
\toprule
\textbf{Category} & \textbf{Value 1} & \textbf{Value 2} & \textbf{Total} \\
\midrule
Item A & 10.5 & 20.3 & 30.8 \\
Item B & 15.2 & 25.1 & 40.3 \\
Item C & 22.1 & 12.7 & 34.8 \\
\bottomrule
\end{tabular}
\caption{Sample Data Table}
\label{tab:sample}
\end{table}

\section{Lists}
\subsection{Bullet Points}
\begin{itemize}
  \item First item
  \item Second item
  \item Third item with sub-items:
    \begin{itemize}
      \item Sub-item 1
      \item Sub-item 2
    \end{itemize}
\end{itemize}

\subsection{Numbered List}
\begin{enumerate}
  \item First step
  \item Second step
  \item Third step
\end{enumerate}

\section{Citations and References}
When you use the citation search feature, references will be automatically managed. For example, you might cite a paper like this: \cite{example2024}. All citations will be automatically added to your \texttt{cite.bib} file.

\section{Conclusion}
This template demonstrates basic \LaTeX{} features including mathematical formulas, tables, and lists. The bibliography system is now set up to work seamlessly with Bibby AI's citation management. Feel free to modify it for your needs!

\bibliographystyle{plain}
\bibliography{cite}

\end{document}\documentclass{article}
\usepackage[utf8]{inputenc}
\usepackage{booktabs}
\usepackage{amsmath}
\usepackage{natbib}

\title{My Document}
\author{Your Name}
\date{\today}

\begin{document}

\maketitle

\section{Introduction}
This is a sample document demonstrating various \LaTeX{} features including mathematical equations, tables, and formatting. You can cite papers using the citation search feature, and they will be automatically added to your \texttt{cite.bib} file.

\section{Mathematical Expressions}
Here are some example equations:

\subsection{Quadratic Formula}
The quadratic formula is given by:
\[x = \frac{-b \pm \sqrt{b^2 - 4ac}}{2a}\]

\subsection{Matrix Example}
A sample matrix:
\[
\begin{pmatrix}
a_{11} & a_{12} & a_{13} \\
a_{21} & a_{22} & a_{23} \\
a_{31} & a_{32} & a_{33}
\end{pmatrix}
\]

\section{Tables}
Below is a sample table showing data:

\begin{table}[h]
\centering
\begin{tabular}{@{}lrrr@{}}
\toprule
\textbf{Category} & \textbf{Value 1} & \textbf{Value 2} & \textbf{Total} \\
\midrule
Item A & 10.5 & 20.3 & 30.8 \\
Item B & 15.2 & 25.1 & 40.3 \\
Item C & 22.1 & 12.7 & 34.8 \\
\bottomrule
\end{tabular}
\caption{Sample Data Table}
\label{tab:sample}
\end{table}

\section{Lists}
\subsection{Bullet Points}
\begin{itemize}
  \item First item
  \item Second item
  \item Third item with sub-items:
    \begin{itemize}
      \item Sub-item 1
      \item Sub-item 2
    \end{itemize}
\end{itemize}

\subsection{Numbered List}
\begin{enumerate}
  \item First step
  \item Second step
  \item Third step
\end{enumerate}

\section{Citations and References}
When you use the citation search feature, references will be automatically managed. For example, you might cite a paper like this: \cite{example2024}. All citations will be automatically added to your \texttt{cite.bib} file.

\section{Conclusion}
This template demonstrates basic \LaTeX{} features including mathematical formulas, tables, and lists. The bibliography system is now set up to work seamlessly with Bibby AI's citation management. Feel free to modify it for your needs!

\bibliographystyle{plain}
\bibliography{cite}

\end{document}\documentclass{article}
\usepackage[utf8]{inputenc}
\usepackage{booktabs}
\usepackage{amsmath}
\usepackage{natbib}

\title{My Document}
\author{Your Name}
\date{\today}

\begin{document}

\maketitle

\section{Introduction}
This is a sample document demonstrating various \LaTeX{} features including mathematical equations, tables, and formatting. You can cite papers using the citation search feature, and they will be automatically added to your \texttt{cite.bib} file.

\section{Mathematical Expressions}
Here are some example equations:

\subsection{Quadratic Formula}
The quadratic formula is given by:
\[x = \frac{-b \pm \sqrt{b^2 - 4ac}}{2a}\]

\subsection{Matrix Example}
A sample matrix:
\[
\begin{pmatrix}
a_{11} & a_{12} & a_{13} \\
a_{21} & a_{22} & a_{23} \\
a_{31} & a_{32} & a_{33}
\end{pmatrix}
\]

\section{Tables}
Below is a sample table showing data:

\begin{table}[h]
\centering
\begin{tabular}{@{}lrrr@{}}
\toprule
\textbf{Category} & \textbf{Value 1} & \textbf{Value 2} & \textbf{Total} \\
\midrule
Item A & 10.5 & 20.3 & 30.8 \\
Item B & 15.2 & 25.1 & 40.3 \\
Item C & 22.1 & 12.7 & 34.8 \\
\bottomrule
\end{tabular}
\caption{Sample Data Table}
\label{tab:sample}
\end{table}

\section{Lists}
\subsection{Bullet Points}
\begin{itemize}
  \item First item
  \item Second item
  \item Third item with sub-items:
    \begin{itemize}
      \item Sub-item 1
      \item Sub-item 2
    \end{itemize}
\end{itemize}

\subsection{Numbered List}
\begin{enumerate}
  \item First step
  \item Second step
  \item Third step
\end{enumerate}

\section{Citations and References}
When you use the citation search feature, references will be automatically managed. For example, you might cite a paper like this: \cite{example2024}. All citations will be automatically added to your \texttt{cite.bib} file.

\section{Conclusion}
This template demonstrates basic \LaTeX{} features including mathematical formulas, tables, and lists. The bibliography system is now set up to work seamlessly with Bibby AI's citation management. Feel free to modify it for your needs!

\bibliographystyle{plain}
\bibliography{cite}

\end{document}\documentclass{article}
\usepackage[utf8]{inputenc}
\usepackage{booktabs}
\usepackage{amsmath}
\usepackage{natbib}

\title{My Document}
\author{Your Name}
\date{\today}

\begin{document}

\maketitle

\section{Introduction}
This is a sample document demonstrating various \LaTeX{} features including mathematical equations, tables, and formatting. You can cite papers using the citation search feature, and they will be automatically added to your \texttt{cite.bib} file.

\section{Mathematical Expressions}
Here are some example equations:

\subsection{Quadratic Formula}
The quadratic formula is given by:
\[x = \frac{-b \pm \sqrt{b^2 - 4ac}}{2a}\]

\subsection{Matrix Example}
A sample matrix:
\[
\begin{pmatrix}
a_{11} & a_{12} & a_{13} \\
a_{21} & a_{22} & a_{23} \\
a_{31} & a_{32} & a_{33}
\end{pmatrix}
\]

\section{Tables}
Below is a sample table showing data:

\begin{table}[h]
\centering
\begin{tabular}{@{}lrrr@{}}
\toprule
\textbf{Category} & \textbf{Value 1} & \textbf{Value 2} & \textbf{Total} \\
\midrule
Item A & 10.5 & 20.3 & 30.8 \\
Item B & 15.2 & 25.1 & 40.3 \\
Item C & 22.1 & 12.7 & 34.8 \\
\bottomrule
\end{tabular}
\caption{Sample Data Table}
\label{tab:sample}
\end{table}

\section{Lists}
\subsection{Bullet Points}
\begin{itemize}
  \item First item
  \item Second item
  \item Third item with sub-items:
    \begin{itemize}
      \item Sub-item 1
      \item Sub-item 2
    \end{itemize}
\end{itemize}

\subsection{Numbered List}
\begin{enumerate}
  \item First step
  \item Second step
  \item Third step
\end{enumerate}

\section{Citations and References}
When you use the citation search feature, references will be automatically managed. For example, you might cite a paper like this: \cite{example2024}. All citations will be automatically added to your \texttt{cite.bib} file.

\section{Conclusion}
This template demonstrates basic \LaTeX{} features including mathematical formulas, tables, and lists. The bibliography system is now set up to work seamlessly with Bibby AI's citation management. Feel free to modify it for your needs!

\bibliographystyle{plain}
\bibliography{cite}

\end{document}\documentclass{article}
\usepackage[utf8]{inputenc}
\usepackage{booktabs}
\usepackage{amsmath}
\usepackage{natbib}

\title{My Document}
\author{Your Name}
\date{\today}

\begin{document}

\maketitle

\section{Introduction}
This is a sample document demonstrating various \LaTeX{} features including mathematical equations, tables, and formatting. You can cite papers using the citation search feature, and they will be automatically added to your \texttt{cite.bib} file.

\section{Mathematical Expressions}
Here are some example equations:

\subsection{Quadratic Formula}
The quadratic formula is given by:
\[x = \frac{-b \pm \sqrt{b^2 - 4ac}}{2a}\]

\subsection{Matrix Example}
A sample matrix:
\[
\begin{pmatrix}
a_{11} & a_{12} & a_{13} \\
a_{21} & a_{22} & a_{23} \\
a_{31} & a_{32} & a_{33}
\end{pmatrix}
\]

\section{Tables}
Below is a sample table showing data:

\begin{table}[h]
\centering
\begin{tabular}{@{}lrrr@{}}
\toprule
\textbf{Category} & \textbf{Value 1} & \textbf{Value 2} & \textbf{Total} \\
\midrule
Item A & 10.5 & 20.3 & 30.8 \\
Item B & 15.2 & 25.1 & 40.3 \\
Item C & 22.1 & 12.7 & 34.8 \\
\bottomrule
\end{tabular}
\caption{Sample Data Table}
\label{tab:sample}
\end{table}

\section{Lists}
\subsection{Bullet Points}
\begin{itemize}
  \item First item
  \item Second item
  \item Third item with sub-items:
    \begin{itemize}
      \item Sub-item 1
      \item Sub-item 2
    \end{itemize}
\end{itemize}

\subsection{Numbered List}
\begin{enumerate}
  \item First step
  \item Second step
  \item Third step
\end{enumerate}

\section{Citations and References}
When you use the citation search feature, references will be automatically managed. For example, you might cite a paper like this: \cite{example2024}. All citations will be automatically added to your \texttt{cite.bib} file.

\section{Conclusion}
This template demonstrates basic \LaTeX{} features including mathematical formulas, tables, and lists. The bibliography system is now set up to work seamlessly with Bibby AI's citation management. Feel free to modify it for your needs!

\bibliographystyle{plain}
\bibliography{cite}

\end{document}\documentclass{article}
\usepackage[utf8]{inputenc}
\usepackage{booktabs}
\usepackage{amsmath}
\usepackage{natbib}

\title{My Document}
\author{Your Name}
\date{\today}

\begin{document}

\maketitle

\section{Introduction}
This is a sample document demonstrating various \LaTeX{} features including mathematical equations, tables, and formatting. You can cite papers using the citation search feature, and they will be automatically added to your \texttt{cite.bib} file.

\section{Mathematical Expressions}
Here are some example equations:

\subsection{Quadratic Formula}
The quadratic formula is given by:
\[x = \frac{-b \pm \sqrt{b^2 - 4ac}}{2a}\]

\subsection{Matrix Example}
A sample matrix:
\[
\begin{pmatrix}
a_{11} & a_{12} & a_{13} \\
a_{21} & a_{22} & a_{23} \\
a_{31} & a_{32} & a_{33}
\end{pmatrix}
\]

\section{Tables}
Below is a sample table showing data:

\begin{table}[h]
\centering
\begin{tabular}{@{}lrrr@{}}
\toprule
\textbf{Category} & \textbf{Value 1} & \textbf{Value 2} & \textbf{Total} \\
\midrule
Item A & 10.5 & 20.3 & 30.8 \\
Item B & 15.2 & 25.1 & 40.3 \\
Item C & 22.1 & 12.7 & 34.8 \\
\bottomrule
\end{tabular}
\caption{Sample Data Table}
\label{tab:sample}
\end{table}

\section{Lists}
\subsection{Bullet Points}
\begin{itemize}
  \item First item
  \item Second item
  \item Third item with sub-items:
    \begin{itemize}
      \item Sub-item 1
      \item Sub-item 2
    \end{itemize}
\end{itemize}

\subsection{Numbered List}
\begin{enumerate}
  \item First step
  \item Second step
  \item Third step
\end{enumerate}

\section{Citations and References}
When you use the citation search feature, references will be automatically managed. For example, you might cite a paper like this: \cite{example2024}. All citations will be automatically added to your \texttt{cite.bib} file.

\section{Conclusion}
This template demonstrates basic \LaTeX{} features including mathematical formulas, tables, and lists. The bibliography system is now set up to work seamlessly with Bibby AI's citation management. Feel free to modify it for your needs!

\bibliographystyle{plain}
\bibliography{cite}

\end{document}\documentclass{article}
\usepackage[utf8]{inputenc}
\usepackage{booktabs}
\usepackage{amsmath}
\usepackage{natbib}

\title{My Document}
\author{Your Name}
\date{\today}

\begin{document}

\maketitle

\section{Introduction}
This is a sample document demonstrating various \LaTeX{} features including mathematical equations, tables, and formatting. You can cite papers using the citation search feature, and they will be automatically added to your \texttt{cite.bib} file.

\section{Mathematical Expressions}
Here are some example equations:

\subsection{Quadratic Formula}
The quadratic formula is given by:
\[x = \frac{-b \pm \sqrt{b^2 - 4ac}}{2a}\]

\subsection{Matrix Example}
A sample matrix:
\[
\begin{pmatrix}
a_{11} & a_{12} & a_{13} \\
a_{21} & a_{22} & a_{23} \\
a_{31} & a_{32} & a_{33}
\end{pmatrix}
\]

\section{Tables}
Below is a sample table showing data:

\begin{table}[h]
\centering
\begin{tabular}{@{}lrrr@{}}
\toprule
\textbf{Category} & \textbf{Value 1} & \textbf{Value 2} & \textbf{Total} \\
\midrule
Item A & 10.5 & 20.3 & 30.8 \\
Item B & 15.2 & 25.1 & 40.3 \\
Item C & 22.1 & 12.7 & 34.8 \\
\bottomrule
\end{tabular}
\caption{Sample Data Table}
\label{tab:sample}
\end{table}

\section{Lists}
\subsection{Bullet Points}
\begin{itemize}
  \item First item
  \item Second item
  \item Third item with sub-items:
    \begin{itemize}
      \item Sub-item 1
      \item Sub-item 2
    \end{itemize}
\end{itemize}

\subsection{Numbered List}
\begin{enumerate}
  \item First step
  \item Second step
  \item Third step
\end{enumerate}

\section{Citations and References}
When you use the citation search feature, references will be automatically managed. For example, you might cite a paper like this: \cite{example2024}. All citations will be automatically added to your \texttt{cite.bib} file.

\section{Conclusion}
This template demonstrates basic \LaTeX{} features including mathematical formulas, tables, and lists. The bibliography system is now set up to work seamlessly with Bibby AI's citation management. Feel free to modify it for your needs!

\bibliographystyle{plain}
\bibliography{cite}

\end{document}\documentclass{article}
\usepackage[utf8]{inputenc}
\usepackage{booktabs}
\usepackage{amsmath}
\usepackage{natbib}

\title{My Document}
\author{Your Name}
\date{\today}

\begin{document}

\maketitle

\section{Introduction}
This is a sample document demonstrating various \LaTeX{} features including mathematical equations, tables, and formatting. You can cite papers using the citation search feature, and they will be automatically added to your \texttt{cite.bib} file.

\section{Mathematical Expressions}
Here are some example equations:

\subsection{Quadratic Formula}
The quadratic formula is given by:
\[x = \frac{-b \pm \sqrt{b^2 - 4ac}}{2a}\]

\subsection{Matrix Example}
A sample matrix:
\[
\begin{pmatrix}
a_{11} & a_{12} & a_{13} \\
a_{21} & a_{22} & a_{23} \\
a_{31} & a_{32} & a_{33}
\end{pmatrix}
\]

\section{Tables}
Below is a sample table showing data:

\begin{table}[h]
\centering
\begin{tabular}{@{}lrrr@{}}
\toprule
\textbf{Category} & \textbf{Value 1} & \textbf{Value 2} & \textbf{Total} \\
\midrule
Item A & 10.5 & 20.3 & 30.8 \\
Item B & 15.2 & 25.1 & 40.3 \\
Item C & 22.1 & 12.7 & 34.8 \\
\bottomrule
\end{tabular}
\caption{Sample Data Table}
\label{tab:sample}
\end{table}

\section{Lists}
\subsection{Bullet Points}
\begin{itemize}
  \item First item
  \item Second item
  \item Third item with sub-items:
    \begin{itemize}
      \item Sub-item 1
      \item Sub-item 2
    \end{itemize}
\end{itemize}

\subsection{Numbered List}
\begin{enumerate}
  \item First step
  \item Second step
  \item Third step
\end{enumerate}

\section{Citations and References}
When you use the citation search feature, references will be automatically managed. For example, you might cite a paper like this: \cite{example2024}. All citations will be automatically added to your \texttt{cite.bib} file.

\section{Conclusion}
This template demonstrates basic \LaTeX{} features including mathematical formulas, tables, and lists. The bibliography system is now set up to work seamlessly with Bibby AI's citation management. Feel free to modify it for your needs!

\bibliographystyle{plain}
\bibliography{cite}

\end{document}\documentclass{article}
\usepackage[utf8]{inputenc}
\usepackage{booktabs}
\usepackage{amsmath}
\usepackage{natbib}

\title{My Document}
\author{Your Name}
\date{\today}

\begin{document}

\maketitle

\section{Introduction}
This is a sample document demonstrating various \LaTeX{} features including mathematical equations, tables, and formatting. You can cite papers using the citation search feature, and they will be automatically added to your \texttt{cite.bib} file.

\section{Mathematical Expressions}
Here are some example equations:

\subsection{Quadratic Formula}
The quadratic formula is given by:
\[x = \frac{-b \pm \sqrt{b^2 - 4ac}}{2a}\]

\subsection{Matrix Example}
A sample matrix:
\[
\begin{pmatrix}
a_{11} & a_{12} & a_{13} \\
a_{21} & a_{22} & a_{23} \\
a_{31} & a_{32} & a_{33}
\end{pmatrix}
\]

\section{Tables}
Below is a sample table showing data:

\begin{table}[h]
\centering
\begin{tabular}{@{}lrrr@{}}
\toprule
\textbf{Category} & \textbf{Value 1} & \textbf{Value 2} & \textbf{Total} \\
\midrule
Item A & 10.5 & 20.3 & 30.8 \\
Item B & 15.2 & 25.1 & 40.3 \\
Item C & 22.1 & 12.7 & 34.8 \\
\bottomrule
\end{tabular}
\caption{Sample Data Table}
\label{tab:sample}
\end{table}

\section{Lists}
\subsection{Bullet Points}
\begin{itemize}
  \item First item
  \item Second item
  \item Third item with sub-items:
    \begin{itemize}
      \item Sub-item 1
      \item Sub-item 2
    \end{itemize}
\end{itemize}

\subsection{Numbered List}
\begin{enumerate}
  \item First step
  \item Second step
  \item Third step
\end{enumerate}

\section{Citations and References}
When you use the citation search feature, references will be automatically managed. For example, you might cite a paper like this: \cite{example2024}. All citations will be automatically added to your \texttt{cite.bib} file.

\section{Conclusion}
This template demonstrates basic \LaTeX{} features including mathematical formulas, tables, and lists. The bibliography system is now set up to work seamlessly with Bibby AI's citation management. Feel free to modify it for your needs!

\bibliographystyle{plain}
\bibliography{cite}

\end{document}\documentclass{article}
\usepackage[utf8]{inputenc}
\usepackage{booktabs}
\usepackage{amsmath}
\usepackage{natbib}

\title{My Document}
\author{Your Name}
\date{\today}

\begin{document}

\maketitle

\section{Introduction}
This is a sample document demonstrating various \LaTeX{} features including mathematical equations, tables, and formatting. You can cite papers using the citation search feature, and they will be automatically added to your \texttt{cite.bib} file.

\section{Mathematical Expressions}
Here are some example equations:

\subsection{Quadratic Formula}
The quadratic formula is given by:
\[x = \frac{-b \pm \sqrt{b^2 - 4ac}}{2a}\]

\subsection{Matrix Example}
A sample matrix:
\[
\begin{pmatrix}
a_{11} & a_{12} & a_{13} \\
a_{21} & a_{22} & a_{23} \\
a_{31} & a_{32} & a_{33}
\end{pmatrix}
\]

\section{Tables}
Below is a sample table showing data:

\begin{table}[h]
\centering
\begin{tabular}{@{}lrrr@{}}
\toprule
\textbf{Category} & \textbf{Value 1} & \textbf{Value 2} & \textbf{Total} \\
\midrule
Item A & 10.5 & 20.3 & 30.8 \\
Item B & 15.2 & 25.1 & 40.3 \\
Item C & 22.1 & 12.7 & 34.8 \\
\bottomrule
\end{tabular}
\caption{Sample Data Table}
\label{tab:sample}
\end{table}

\section{Lists}
\subsection{Bullet Points}
\begin{itemize}
  \item First item
  \item Second item
  \item Third item with sub-items:
    \begin{itemize}
      \item Sub-item 1
      \item Sub-item 2
    \end{itemize}
\end{itemize}

\subsection{Numbered List}
\begin{enumerate}
  \item First step
  \item Second step
  \item Third step
\end{enumerate}

\section{Citations and References}
When you use the citation search feature, references will be automatically managed. For example, you might cite a paper like this: \cite{example2024}. All citations will be automatically added to your \texttt{cite.bib} file.

\section{Conclusion}
This template demonstrates basic \LaTeX{} features including mathematical formulas, tables, and lists. The bibliography system is now set up to work seamlessly with Bibby AI's citation management. Feel free to modify it for your needs!

\bibliographystyle{plain}
\bibliography{cite}

\end{document}\documentclass{article}
\usepackage[utf8]{inputenc}
\usepackage{booktabs}
\usepackage{amsmath}
\usepackage{natbib}

\title{My Document}
\author{Your Name}
\date{\today}

\begin{document}

\maketitle

\section{Introduction}
This is a sample document demonstrating various \LaTeX{} features including mathematical equations, tables, and formatting. You can cite papers using the citation search feature, and they will be automatically added to your \texttt{cite.bib} file.

\section{Mathematical Expressions}
Here are some example equations:

\subsection{Quadratic Formula}
The quadratic formula is given by:
\[x = \frac{-b \pm \sqrt{b^2 - 4ac}}{2a}\]

\subsection{Matrix Example}
A sample matrix:
\[
\begin{pmatrix}
a_{11} & a_{12} & a_{13} \\
a_{21} & a_{22} & a_{23} \\
a_{31} & a_{32} & a_{33}
\end{pmatrix}
\]

\section{Tables}
Below is a sample table showing data:

\begin{table}[h]
\centering
\begin{tabular}{@{}lrrr@{}}
\toprule
\textbf{Category} & \textbf{Value 1} & \textbf{Value 2} & \textbf{Total} \\
\midrule
Item A & 10.5 & 20.3 & 30.8 \\
Item B & 15.2 & 25.1 & 40.3 \\
Item C & 22.1 & 12.7 & 34.8 \\
\bottomrule
\end{tabular}
\caption{Sample Data Table}
\label{tab:sample}
\end{table}

\section{Lists}
\subsection{Bullet Points}
\begin{itemize}
  \item First item
  \item Second item
  \item Third item with sub-items:
    \begin{itemize}
      \item Sub-item 1
      \item Sub-item 2
    \end{itemize}
\end{itemize}

\subsection{Numbered List}
\begin{enumerate}
  \item First step
  \item Second step
  \item Third step
\end{enumerate}

\section{Citations and References}
When you use the citation search feature, references will be automatically managed. For example, you might cite a paper like this: \cite{example2024}. All citations will be automatically added to your \texttt{cite.bib} file.

\section{Conclusion}
This template demonstrates basic \LaTeX{} features including mathematical formulas, tables, and lists. The bibliography system is now set up to work seamlessly with Bibby AI's citation management. Feel free to modify it for your needs!

\bibliographystyle{plain}
\bibliography{cite}

\end{document}\documentclass{article}
\usepackage[utf8]{inputenc}
\usepackage{booktabs}
\usepackage{amsmath}
\usepackage{natbib}

\title{My Document}
\author{Your Name}
\date{\today}

\begin{document}

\maketitle

\section{Introduction}
This is a sample document demonstrating various \LaTeX{} features including mathematical equations, tables, and formatting. You can cite papers using the citation search feature, and they will be automatically added to your \texttt{cite.bib} file.

\section{Mathematical Expressions}
Here are some example equations:

\subsection{Quadratic Formula}
The quadratic formula is given by:
\[x = \frac{-b \pm \sqrt{b^2 - 4ac}}{2a}\]

\subsection{Matrix Example}
A sample matrix:
\[
\begin{pmatrix}
a_{11} & a_{12} & a_{13} \\
a_{21} & a_{22} & a_{23} \\
a_{31} & a_{32} & a_{33}
\end{pmatrix}
\]

\section{Tables}
Below is a sample table showing data:

\begin{table}[h]
\centering
\begin{tabular}{@{}lrrr@{}}
\toprule
\textbf{Category} & \textbf{Value 1} & \textbf{Value 2} & \textbf{Total} \\
\midrule
Item A & 10.5 & 20.3 & 30.8 \\
Item B & 15.2 & 25.1 & 40.3 \\
Item C & 22.1 & 12.7 & 34.8 \\
\bottomrule
\end{tabular}
\caption{Sample Data Table}
\label{tab:sample}
\end{table}

\section{Lists}
\subsection{Bullet Points}
\begin{itemize}
  \item First item
  \item Second item
  \item Third item with sub-items:
    \begin{itemize}
      \item Sub-item 1
      \item Sub-item 2
    \end{itemize}
\end{itemize}

\subsection{Numbered List}
\begin{enumerate}
  \item First step
  \item Second step
  \item Third step
\end{enumerate}

\section{Citations and References}
When you use the citation search feature, references will be automatically managed. For example, you might cite a paper like this: \cite{example2024}. All citations will be automatically added to your \texttt{cite.bib} file.

\section{Conclusion}
This template demonstrates basic \LaTeX{} features including mathematical formulas, tables, and lists. The bibliography system is now set up to work seamlessly with Bibby AI's citation management. Feel free to modify it for your needs!

\bibliographystyle{plain}
\bibliography{cite}

\end{document}\documentclass{article}
\usepackage[utf8]{inputenc}
\usepackage{booktabs}
\usepackage{amsmath}
\usepackage{natbib}

\title{My Document}
\author{Your Name}
\date{\today}

\begin{document}

\maketitle

\section{Introduction}
This is a sample document demonstrating various \LaTeX{} features including mathematical equations, tables, and formatting. You can cite papers using the citation search feature, and they will be automatically added to your \texttt{cite.bib} file.

\section{Mathematical Expressions}
Here are some example equations:

\subsection{Quadratic Formula}
The quadratic formula is given by:
\[x = \frac{-b \pm \sqrt{b^2 - 4ac}}{2a}\]

\subsection{Matrix Example}
A sample matrix:
\[
\begin{pmatrix}
a_{11} & a_{12} & a_{13} \\
a_{21} & a_{22} & a_{23} \\
a_{31} & a_{32} & a_{33}
\end{pmatrix}
\]

\section{Tables}
Below is a sample table showing data:

\begin{table}[h]
\centering
\begin{tabular}{@{}lrrr@{}}
\toprule
\textbf{Category} & \textbf{Value 1} & \textbf{Value 2} & \textbf{Total} \\
\midrule
Item A & 10.5 & 20.3 & 30.8 \\
Item B & 15.2 & 25.1 & 40.3 \\
Item C & 22.1 & 12.7 & 34.8 \\
\bottomrule
\end{tabular}
\caption{Sample Data Table}
\label{tab:sample}
\end{table}

\section{Lists}
\subsection{Bullet Points}
\begin{itemize}
  \item First item
  \item Second item
  \item Third item with sub-items:
    \begin{itemize}
      \item Sub-item 1
      \item Sub-item 2
    \end{itemize}
\end{itemize}

\subsection{Numbered List}
\begin{enumerate}
  \item First step
  \item Second step
  \item Third step
\end{enumerate}

\section{Citations and References}
When you use the citation search feature, references will be automatically managed. For example, you might cite a paper like this: \cite{example2024}. All citations will be automatically added to your \texttt{cite.bib} file.

\section{Conclusion}
This template demonstrates basic \LaTeX{} features including mathematical formulas, tables, and lists. The bibliography system is now set up to work seamlessly with Bibby AI's citation management. Feel free to modify it for your needs!

\bibliographystyle{plain}
\bibliography{cite}

\end{document}\documentclass{article}
\usepackage[utf8]{inputenc}
\usepackage{booktabs}
\usepackage{amsmath}
\usepackage{natbib}

\title{My Document}
\author{Your Name}
\date{\today}

\begin{document}

\maketitle

\section{Introduction}
This is a sample document demonstrating various \LaTeX{} features including mathematical equations, tables, and formatting. You can cite papers using the citation search feature, and they will be automatically added to your \texttt{cite.bib} file.

\section{Mathematical Expressions}
Here are some example equations:

\subsection{Quadratic Formula}
The quadratic formula is given by:
\[x = \frac{-b \pm \sqrt{b^2 - 4ac}}{2a}\]

\subsection{Matrix Example}
A sample matrix:
\[
\begin{pmatrix}
a_{11} & a_{12} & a_{13} \\
a_{21} & a_{22} & a_{23} \\
a_{31} & a_{32} & a_{33}
\end{pmatrix}
\]

\section{Tables}
Below is a sample table showing data:

\begin{table}[h]
\centering
\begin{tabular}{@{}lrrr@{}}
\toprule
\textbf{Category} & \textbf{Value 1} & \textbf{Value 2} & \textbf{Total} \\
\midrule
Item A & 10.5 & 20.3 & 30.8 \\
Item B & 15.2 & 25.1 & 40.3 \\
Item C & 22.1 & 12.7 & 34.8 \\
\bottomrule
\end{tabular}
\caption{Sample Data Table}
\label{tab:sample}
\end{table}

\section{Lists}
\subsection{Bullet Points}
\begin{itemize}
  \item First item
  \item Second item
  \item Third item with sub-items:
    \begin{itemize}
      \item Sub-item 1
      \item Sub-item 2
    \end{itemize}
\end{itemize}

\subsection{Numbered List}
\begin{enumerate}
  \item First step
  \item Second step
  \item Third step
\end{enumerate}

\section{Citations and References}
When you use the citation search feature, references will be automatically managed. For example, you might cite a paper like this: \cite{example2024}. All citations will be automatically added to your \texttt{cite.bib} file.

\section{Conclusion}
This template demonstrates basic \LaTeX{} features including mathematical formulas, tables, and lists. The bibliography system is now set up to work seamlessly with Bibby AI's citation management. Feel free to modify it for your needs!

\bibliographystyle{plain}
\bibliography{cite}

\end{document}\documentclass{article}
\usepackage[utf8]{inputenc}
\usepackage{booktabs}
\usepackage{amsmath}
\usepackage{natbib}

\title{My Document}
\author{Your Name}
\date{\today}

\begin{document}

\maketitle

\section{Introduction}
This is a sample document demonstrating various \LaTeX{} features including mathematical equations, tables, and formatting. You can cite papers using the citation search feature, and they will be automatically added to your \texttt{cite.bib} file.

\section{Mathematical Expressions}
Here are some example equations:

\subsection{Quadratic Formula}
The quadratic formula is given by:
\[x = \frac{-b \pm \sqrt{b^2 - 4ac}}{2a}\]

\subsection{Matrix Example}
A sample matrix:
\[
\begin{pmatrix}
a_{11} & a_{12} & a_{13} \\
a_{21} & a_{22} & a_{23} \\
a_{31} & a_{32} & a_{33}
\end{pmatrix}
\]

\section{Tables}
Below is a sample table showing data:

\begin{table}[h]
\centering
\begin{tabular}{@{}lrrr@{}}
\toprule
\textbf{Category} & \textbf{Value 1} & \textbf{Value 2} & \textbf{Total} \\
\midrule
Item A & 10.5 & 20.3 & 30.8 \\
Item B & 15.2 & 25.1 & 40.3 \\
Item C & 22.1 & 12.7 & 34.8 \\
\bottomrule
\end{tabular}
\caption{Sample Data Table}
\label{tab:sample}
\end{table}

\section{Lists}
\subsection{Bullet Points}
\begin{itemize}
  \item First item
  \item Second item
  \item Third item with sub-items:
    \begin{itemize}
      \item Sub-item 1
      \item Sub-item 2
    \end{itemize}
\end{itemize}

\subsection{Numbered List}
\begin{enumerate}
  \item First step
  \item Second step
  \item Third step
\end{enumerate}

\section{Citations and References}
When you use the citation search feature, references will be automatically managed. For example, you might cite a paper like this: \cite{example2024}. All citations will be automatically added to your \texttt{cite.bib} file.

\section{Conclusion}
This template demonstrates basic \LaTeX{} features including mathematical formulas, tables, and lists. The bibliography system is now set up to work seamlessly with Bibby AI's citation management. Feel free to modify it for your needs!

\bibliographystyle{plain}
\bibliography{cite}

\end{document}\documentclass{article}
\usepackage[utf8]{inputenc}
\usepackage{booktabs}
\usepackage{amsmath}
\usepackage{natbib}

\title{My Document}
\author{Your Name}
\date{\today}

\begin{document}

\maketitle

\section{Introduction}
This is a sample document demonstrating various \LaTeX{} features including mathematical equations, tables, and formatting. You can cite papers using the citation search feature, and they will be automatically added to your \texttt{cite.bib} file.

\section{Mathematical Expressions}
Here are some example equations:

\subsection{Quadratic Formula}
The quadratic formula is given by:
\[x = \frac{-b \pm \sqrt{b^2 - 4ac}}{2a}\]

\subsection{Matrix Example}
A sample matrix:
\[
\begin{pmatrix}
a_{11} & a_{12} & a_{13} \\
a_{21} & a_{22} & a_{23} \\
a_{31} & a_{32} & a_{33}
\end{pmatrix}
\]

\section{Tables}
Below is a sample table showing data:

\begin{table}[h]
\centering
\begin{tabular}{@{}lrrr@{}}
\toprule
\textbf{Category} & \textbf{Value 1} & \textbf{Value 2} & \textbf{Total} \\
\midrule
Item A & 10.5 & 20.3 & 30.8 \\
Item B & 15.2 & 25.1 & 40.3 \\
Item C & 22.1 & 12.7 & 34.8 \\
\bottomrule
\end{tabular}
\caption{Sample Data Table}
\label{tab:sample}
\end{table}

\section{Lists}
\subsection{Bullet Points}
\begin{itemize}
  \item First item
  \item Second item
  \item Third item with sub-items:
    \begin{itemize}
      \item Sub-item 1
      \item Sub-item 2
    \end{itemize}
\end{itemize}

\subsection{Numbered List}
\begin{enumerate}
  \item First step
  \item Second step
  \item Third step
\end{enumerate}

\section{Citations and References}
When you use the citation search feature, references will be automatically managed. For example, you might cite a paper like this: \cite{example2024}. All citations will be automatically added to your \texttt{cite.bib} file.

\section{Conclusion}
This template demonstrates basic \LaTeX{} features including mathematical formulas, tables, and lists. The bibliography system is now set up to work seamlessly with Bibby AI's citation management. Feel free to modify it for your needs!

\bibliographystyle{plain}
\bibliography{cite}

\end{document}\documentclass{article}
\usepackage[utf8]{inputenc}
\usepackage{booktabs}
\usepackage{amsmath}
\usepackage{natbib}

\title{My Document}
\author{Your Name}
\date{\today}

\begin{document}

\maketitle

\section{Introduction}
This is a sample document demonstrating various \LaTeX{} features including mathematical equations, tables, and formatting. You can cite papers using the citation search feature, and they will be automatically added to your \texttt{cite.bib} file.

\section{Mathematical Expressions}
Here are some example equations:

\subsection{Quadratic Formula}
The quadratic formula is given by:
\[x = \frac{-b \pm \sqrt{b^2 - 4ac}}{2a}\]

\subsection{Matrix Example}
A sample matrix:
\[
\begin{pmatrix}
a_{11} & a_{12} & a_{13} \\
a_{21} & a_{22} & a_{23} \\
a_{31} & a_{32} & a_{33}
\end{pmatrix}
\]

\section{Tables}
Below is a sample table showing data:

\begin{table}[h]
\centering
\begin{tabular}{@{}lrrr@{}}
\toprule
\textbf{Category} & \textbf{Value 1} & \textbf{Value 2} & \textbf{Total} \\
\midrule
Item A & 10.5 & 20.3 & 30.8 \\
Item B & 15.2 & 25.1 & 40.3 \\
Item C & 22.1 & 12.7 & 34.8 \\
\bottomrule
\end{tabular}
\caption{Sample Data Table}
\label{tab:sample}
\end{table}

\section{Lists}
\subsection{Bullet Points}
\begin{itemize}
  \item First item
  \item Second item
  \item Third item with sub-items:
    \begin{itemize}
      \item Sub-item 1
      \item Sub-item 2
    \end{itemize}
\end{itemize}

\subsection{Numbered List}
\begin{enumerate}
  \item First step
  \item Second step
  \item Third step
\end{enumerate}

\section{Citations and References}
When you use the citation search feature, references will be automatically managed. For example, you might cite a paper like this: \cite{example2024}. All citations will be automatically added to your \texttt{cite.bib} file.

\section{Conclusion}
This template demonstrates basic \LaTeX{} features including mathematical formulas, tables, and lists. The bibliography system is now set up to work seamlessly with Bibby AI's citation management. Feel free to modify it for your needs!

\bibliographystyle{plain}
\bibliography{cite}

\end{document}\documentclass{article}
\usepackage[utf8]{inputenc}
\usepackage{booktabs}
\usepackage{amsmath}
\usepackage{natbib}

\title{My Document}
\author{Your Name}
\date{\today}

\begin{document}

\maketitle

\section{Introduction}
This is a sample document demonstrating various \LaTeX{} features including mathematical equations, tables, and formatting. You can cite papers using the citation search feature, and they will be automatically added to your \texttt{cite.bib} file.

\section{Mathematical Expressions}
Here are some example equations:

\subsection{Quadratic Formula}
The quadratic formula is given by:
\[x = \frac{-b \pm \sqrt{b^2 - 4ac}}{2a}\]

\subsection{Matrix Example}
A sample matrix:
\[
\begin{pmatrix}
a_{11} & a_{12} & a_{13} \\
a_{21} & a_{22} & a_{23} \\
a_{31} & a_{32} & a_{33}
\end{pmatrix}
\]

\section{Tables}
Below is a sample table showing data:

\begin{table}[h]
\centering
\begin{tabular}{@{}lrrr@{}}
\toprule
\textbf{Category} & \textbf{Value 1} & \textbf{Value 2} & \textbf{Total} \\
\midrule
Item A & 10.5 & 20.3 & 30.8 \\
Item B & 15.2 & 25.1 & 40.3 \\
Item C & 22.1 & 12.7 & 34.8 \\
\bottomrule
\end{tabular}
\caption{Sample Data Table}
\label{tab:sample}
\end{table}

\section{Lists}
\subsection{Bullet Points}
\begin{itemize}
  \item First item
  \item Second item
  \item Third item with sub-items:
    \begin{itemize}
      \item Sub-item 1
      \item Sub-item 2
    \end{itemize}
\end{itemize}

\subsection{Numbered List}
\begin{enumerate}
  \item First step
  \item Second step
  \item Third step
\end{enumerate}

\section{Citations and References}
When you use the citation search feature, references will be automatically managed. For example, you might cite a paper like this: \cite{example2024}. All citations will be automatically added to your \texttt{cite.bib} file.

\section{Conclusion}
This template demonstrates basic \LaTeX{} features including mathematical formulas, tables, and lists. The bibliography system is now set up to work seamlessly with Bibby AI's citation management. Feel free to modify it for your needs!

\bibliographystyle{plain}
\bibliography{cite}

\end{document}\documentclass{article}
\usepackage[utf8]{inputenc}
\usepackage{booktabs}
\usepackage{amsmath}
\usepackage{natbib}

\title{My Document}
\author{Your Name}
\date{\today}

\begin{document}

\maketitle

\section{Introduction}
This is a sample document demonstrating various \LaTeX{} features including mathematical equations, tables, and formatting. You can cite papers using the citation search feature, and they will be automatically added to your \texttt{cite.bib} file.

\section{Mathematical Expressions}
Here are some example equations:

\subsection{Quadratic Formula}
The quadratic formula is given by:
\[x = \frac{-b \pm \sqrt{b^2 - 4ac}}{2a}\]

\subsection{Matrix Example}
A sample matrix:
\[
\begin{pmatrix}
a_{11} & a_{12} & a_{13} \\
a_{21} & a_{22} & a_{23} \\
a_{31} & a_{32} & a_{33}
\end{pmatrix}
\]

\section{Tables}
Below is a sample table showing data:

\begin{table}[h]
\centering
\begin{tabular}{@{}lrrr@{}}
\toprule
\textbf{Category} & \textbf{Value 1} & \textbf{Value 2} & \textbf{Total} \\
\midrule
Item A & 10.5 & 20.3 & 30.8 \\
Item B & 15.2 & 25.1 & 40.3 \\
Item C & 22.1 & 12.7 & 34.8 \\
\bottomrule
\end{tabular}
\caption{Sample Data Table}
\label{tab:sample}
\end{table}

\section{Lists}
\subsection{Bullet Points}
\begin{itemize}
  \item First item
  \item Second item
  \item Third item with sub-items:
    \begin{itemize}
      \item Sub-item 1
      \item Sub-item 2
    \end{itemize}
\end{itemize}

\subsection{Numbered List}
\begin{enumerate}
  \item First step
  \item Second step
  \item Third step
\end{enumerate}

\section{Citations and References}
When you use the citation search feature, references will be automatically managed. For example, you might cite a paper like this: \cite{example2024}. All citations will be automatically added to your \texttt{cite.bib} file.

\section{Conclusion}
This template demonstrates basic \LaTeX{} features including mathematical formulas, tables, and lists. The bibliography system is now set up to work seamlessly with Bibby AI's citation management. Feel free to modify it for your needs!

\bibliographystyle{plain}
\bibliography{cite}

\end{document}\documentclass{article}
\usepackage[utf8]{inputenc}
\usepackage{booktabs}
\usepackage{amsmath}
\usepackage{natbib}

\title{My Document}
\author{Your Name}
\date{\today}

\begin{document}

\maketitle

\section{Introduction}
This is a sample document demonstrating various \LaTeX{} features including mathematical equations, tables, and formatting. You can cite papers using the citation search feature, and they will be automatically added to your \texttt{cite.bib} file.

\section{Mathematical Expressions}
Here are some example equations:

\subsection{Quadratic Formula}
The quadratic formula is given by:
\[x = \frac{-b \pm \sqrt{b^2 - 4ac}}{2a}\]

\subsection{Matrix Example}
A sample matrix:
\[
\begin{pmatrix}
a_{11} & a_{12} & a_{13} \\
a_{21} & a_{22} & a_{23} \\
a_{31} & a_{32} & a_{33}
\end{pmatrix}
\]

\section{Tables}
Below is a sample table showing data:

\begin{table}[h]
\centering
\begin{tabular}{@{}lrrr@{}}
\toprule
\textbf{Category} & \textbf{Value 1} & \textbf{Value 2} & \textbf{Total} \\
\midrule
Item A & 10.5 & 20.3 & 30.8 \\
Item B & 15.2 & 25.1 & 40.3 \\
Item C & 22.1 & 12.7 & 34.8 \\
\bottomrule
\end{tabular}
\caption{Sample Data Table}
\label{tab:sample}
\end{table}

\section{Lists}
\subsection{Bullet Points}
\begin{itemize}
  \item First item
  \item Second item
  \item Third item with sub-items:
    \begin{itemize}
      \item Sub-item 1
      \item Sub-item 2
    \end{itemize}
\end{itemize}

\subsection{Numbered List}
\begin{enumerate}
  \item First step
  \item Second step
  \item Third step
\end{enumerate}

\section{Citations and References}
When you use the citation search feature, references will be automatically managed. For example, you might cite a paper like this: \cite{example2024}. All citations will be automatically added to your \texttt{cite.bib} file.

\section{Conclusion}
This template demonstrates basic \LaTeX{} features including mathematical formulas, tables, and lists. The bibliography system is now set up to work seamlessly with Bibby AI's citation management. Feel free to modify it for your needs!

\bibliographystyle{plain}
\bibliography{cite}

\end{document}\documentclass{article}
\usepackage[utf8]{inputenc}
\usepackage{booktabs}
\usepackage{amsmath}
\usepackage{natbib}

\title{My Document}
\author{Your Name}
\date{\today}

\begin{document}

\maketitle

\section{Introduction}
This is a sample document demonstrating various \LaTeX{} features including mathematical equations, tables, and formatting. You can cite papers using the citation search feature, and they will be automatically added to your \texttt{cite.bib} file.

\section{Mathematical Expressions}
Here are some example equations:

\subsection{Quadratic Formula}
The quadratic formula is given by:
\[x = \frac{-b \pm \sqrt{b^2 - 4ac}}{2a}\]

\subsection{Matrix Example}
A sample matrix:
\[
\begin{pmatrix}
a_{11} & a_{12} & a_{13} \\
a_{21} & a_{22} & a_{23} \\
a_{31} & a_{32} & a_{33}
\end{pmatrix}
\]

\section{Tables}
Below is a sample table showing data:

\begin{table}[h]
\centering
\begin{tabular}{@{}lrrr@{}}
\toprule
\textbf{Category} & \textbf{Value 1} & \textbf{Value 2} & \textbf{Total} \\
\midrule
Item A & 10.5 & 20.3 & 30.8 \\
Item B & 15.2 & 25.1 & 40.3 \\
Item C & 22.1 & 12.7 & 34.8 \\
\bottomrule
\end{tabular}
\caption{Sample Data Table}
\label{tab:sample}
\end{table}

\section{Lists}
\subsection{Bullet Points}
\begin{itemize}
  \item First item
  \item Second item
  \item Third item with sub-items:
    \begin{itemize}
      \item Sub-item 1
      \item Sub-item 2
    \end{itemize}
\end{itemize}

\subsection{Numbered List}
\begin{enumerate}
  \item First step
  \item Second step
  \item Third step
\end{enumerate}

\section{Citations and References}
When you use the citation search feature, references will be automatically managed. For example, you might cite a paper like this: \cite{example2024}. All citations will be automatically added to your \texttt{cite.bib} file.

\section{Conclusion}
This template demonstrates basic \LaTeX{} features including mathematical formulas, tables, and lists. The bibliography system is now set up to work seamlessly with Bibby AI's citation management. Feel free to modify it for your needs!

\bibliographystyle{plain}
\bibliography{cite}

\end{document}\documentclass{article}
\usepackage[utf8]{inputenc}
\usepackage{booktabs}
\usepackage{amsmath}
\usepackage{natbib}

\title{My Document}
\author{Your Name}
\date{\today}

\begin{document}

\maketitle

\section{Introduction}
This is a sample document demonstrating various \LaTeX{} features including mathematical equations, tables, and formatting. You can cite papers using the citation search feature, and they will be automatically added to your \texttt{cite.bib} file.

\section{Mathematical Expressions}
Here are some example equations:

\subsection{Quadratic Formula}
The quadratic formula is given by:
\[x = \frac{-b \pm \sqrt{b^2 - 4ac}}{2a}\]

\subsection{Matrix Example}
A sample matrix:
\[
\begin{pmatrix}
a_{11} & a_{12} & a_{13} \\
a_{21} & a_{22} & a_{23} \\
a_{31} & a_{32} & a_{33}
\end{pmatrix}
\]

\section{Tables}
Below is a sample table showing data:

\begin{table}[h]
\centering
\begin{tabular}{@{}lrrr@{}}
\toprule
\textbf{Category} & \textbf{Value 1} & \textbf{Value 2} & \textbf{Total} \\
\midrule
Item A & 10.5 & 20.3 & 30.8 \\
Item B & 15.2 & 25.1 & 40.3 \\
Item C & 22.1 & 12.7 & 34.8 \\
\bottomrule
\end{tabular}
\caption{Sample Data Table}
\label{tab:sample}
\end{table}

\section{Lists}
\subsection{Bullet Points}
\begin{itemize}
  \item First item
  \item Second item
  \item Third item with sub-items:
    \begin{itemize}
      \item Sub-item 1
      \item Sub-item 2
    \end{itemize}
\end{itemize}

\subsection{Numbered List}
\begin{enumerate}
  \item First step
  \item Second step
  \item Third step
\end{enumerate}

\section{Citations and References}
When you use the citation search feature, references will be automatically managed. For example, you might cite a paper like this: \cite{example2024}. All citations will be automatically added to your \texttt{cite.bib} file.

\section{Conclusion}
This template demonstrates basic \LaTeX{} features including mathematical formulas, tables, and lists. The bibliography system is now set up to work seamlessly with Bibby AI's citation management. Feel free to modify it for your needs!

\bibliographystyle{plain}
\bibliography{cite}

\end{document}\documentclass{article}
\usepackage[utf8]{inputenc}
\usepackage{booktabs}
\usepackage{amsmath}
\usepackage{natbib}

\title{My Document}
\author{Your Name}
\date{\today}

\begin{document}

\maketitle

\section{Introduction}
This is a sample document demonstrating various \LaTeX{} features including mathematical equations, tables, and formatting. You can cite papers using the citation search feature, and they will be automatically added to your \texttt{cite.bib} file.

\section{Mathematical Expressions}
Here are some example equations:

\subsection{Quadratic Formula}
The quadratic formula is given by:
\[x = \frac{-b \pm \sqrt{b^2 - 4ac}}{2a}\]

\subsection{Matrix Example}
A sample matrix:
\[
\begin{pmatrix}
a_{11} & a_{12} & a_{13} \\
a_{21} & a_{22} & a_{23} \\
a_{31} & a_{32} & a_{33}
\end{pmatrix}
\]

\section{Tables}
Below is a sample table showing data:

\begin{table}[h]
\centering
\begin{tabular}{@{}lrrr@{}}
\toprule
\textbf{Category} & \textbf{Value 1} & \textbf{Value 2} & \textbf{Total} \\
\midrule
Item A & 10.5 & 20.3 & 30.8 \\
Item B & 15.2 & 25.1 & 40.3 \\
Item C & 22.1 & 12.7 & 34.8 \\
\bottomrule
\end{tabular}
\caption{Sample Data Table}
\label{tab:sample}
\end{table}

\section{Lists}
\subsection{Bullet Points}
\begin{itemize}
  \item First item
  \item Second item
  \item Third item with sub-items:
    \begin{itemize}
      \item Sub-item 1
      \item Sub-item 2
    \end{itemize}
\end{itemize}

\subsection{Numbered List}
\begin{enumerate}
  \item First step
  \item Second step
  \item Third step
\end{enumerate}

\section{Citations and References}
When you use the citation search feature, references will be automatically managed. For example, you might cite a paper like this: \cite{example2024}. All citations will be automatically added to your \texttt{cite.bib} file.

\section{Conclusion}
This template demonstrates basic \LaTeX{} features including mathematical formulas, tables, and lists. The bibliography system is now set up to work seamlessly with Bibby AI's citation management. Feel free to modify it for your needs!

\bibliographystyle{plain}
\bibliography{cite}

\end{document}\documentclass{article}
\usepackage[utf8]{inputenc}
\usepackage{booktabs}
\usepackage{amsmath}
\usepackage{natbib}

\title{My Document}
\author{Your Name}
\date{\today}

\begin{document}

\maketitle

\section{Introduction}
This is a sample document demonstrating various \LaTeX{} features including mathematical equations, tables, and formatting. You can cite papers using the citation search feature, and they will be automatically added to your \texttt{cite.bib} file.

\section{Mathematical Expressions}
Here are some example equations:

\subsection{Quadratic Formula}
The quadratic formula is given by:
\[x = \frac{-b \pm \sqrt{b^2 - 4ac}}{2a}\]

\subsection{Matrix Example}
A sample matrix:
\[
\begin{pmatrix}
a_{11} & a_{12} & a_{13} \\
a_{21} & a_{22} & a_{23} \\
a_{31} & a_{32} & a_{33}
\end{pmatrix}
\]

\section{Tables}
Below is a sample table showing data:

\begin{table}[h]
\centering
\begin{tabular}{@{}lrrr@{}}
\toprule
\textbf{Category} & \textbf{Value 1} & \textbf{Value 2} & \textbf{Total} \\
\midrule
Item A & 10.5 & 20.3 & 30.8 \\
Item B & 15.2 & 25.1 & 40.3 \\
Item C & 22.1 & 12.7 & 34.8 \\
\bottomrule
\end{tabular}
\caption{Sample Data Table}
\label{tab:sample}
\end{table}

\section{Lists}
\subsection{Bullet Points}
\begin{itemize}
  \item First item
  \item Second item
  \item Third item with sub-items:
    \begin{itemize}
      \item Sub-item 1
      \item Sub-item 2
    \end{itemize}
\end{itemize}

\subsection{Numbered List}
\begin{enumerate}
  \item First step
  \item Second step
  \item Third step
\end{enumerate}

\section{Citations and References}
When you use the citation search feature, references will be automatically managed. For example, you might cite a paper like this: \cite{example2024}. All citations will be automatically added to your \texttt{cite.bib} file.

\section{Conclusion}
This template demonstrates basic \LaTeX{} features including mathematical formulas, tables, and lists. The bibliography system is now set up to work seamlessly with Bibby AI's citation management. Feel free to modify it for your needs!

\bibliographystyle{plain}
\bibliography{cite}

\end{document}\documentclass{article}
\usepackage[utf8]{inputenc}
\usepackage{booktabs}
\usepackage{amsmath}
\usepackage{natbib}

\title{My Document}
\author{Your Name}
\date{\today}

\begin{document}

\maketitle

\section{Introduction}
This is a sample document demonstrating various \LaTeX{} features including mathematical equations, tables, and formatting. You can cite papers using the citation search feature, and they will be automatically added to your \texttt{cite.bib} file.

\section{Mathematical Expressions}
Here are some example equations:

\subsection{Quadratic Formula}
The quadratic formula is given by:
\[x = \frac{-b \pm \sqrt{b^2 - 4ac}}{2a}\]

\subsection{Matrix Example}
A sample matrix:
\[
\begin{pmatrix}
a_{11} & a_{12} & a_{13} \\
a_{21} & a_{22} & a_{23} \\
a_{31} & a_{32} & a_{33}
\end{pmatrix}
\]

\section{Tables}
Below is a sample table showing data:

\begin{table}[h]
\centering
\begin{tabular}{@{}lrrr@{}}
\toprule
\textbf{Category} & \textbf{Value 1} & \textbf{Value 2} & \textbf{Total} \\
\midrule
Item A & 10.5 & 20.3 & 30.8 \\
Item B & 15.2 & 25.1 & 40.3 \\
Item C & 22.1 & 12.7 & 34.8 \\
\bottomrule
\end{tabular}
\caption{Sample Data Table}
\label{tab:sample}
\end{table}

\section{Lists}
\subsection{Bullet Points}
\begin{itemize}
  \item First item
  \item Second item
  \item Third item with sub-items:
    \begin{itemize}
      \item Sub-item 1
      \item Sub-item 2
    \end{itemize}
\end{itemize}

\subsection{Numbered List}
\begin{enumerate}
  \item First step
  \item Second step
  \item Third step
\end{enumerate}

\section{Citations and References}
When you use the citation search feature, references will be automatically managed. For example, you might cite a paper like this: \cite{example2024}. All citations will be automatically added to your \texttt{cite.bib} file.

\section{Conclusion}
This template demonstrates basic \LaTeX{} features including mathematical formulas, tables, and lists. The bibliography system is now set up to work seamlessly with Bibby AI's citation management. Feel free to modify it for your needs!

\bibliographystyle{plain}
\bibliography{cite}

\end{document}\documentclass{article}
\usepackage[utf8]{inputenc}
\usepackage{booktabs}
\usepackage{amsmath}
\usepackage{natbib}

\title{My Document}
\author{Your Name}
\date{\today}

\begin{document}

\maketitle

\section{Introduction}
This is a sample document demonstrating various \LaTeX{} features including mathematical equations, tables, and formatting. You can cite papers using the citation search feature, and they will be automatically added to your \texttt{cite.bib} file.

\section{Mathematical Expressions}
Here are some example equations:

\subsection{Quadratic Formula}
The quadratic formula is given by:
\[x = \frac{-b \pm \sqrt{b^2 - 4ac}}{2a}\]

\subsection{Matrix Example}
A sample matrix:
\[
\begin{pmatrix}
a_{11} & a_{12} & a_{13} \\
a_{21} & a_{22} & a_{23} \\
a_{31} & a_{32} & a_{33}
\end{pmatrix}
\]

\section{Tables}
Below is a sample table showing data:

\begin{table}[h]
\centering
\begin{tabular}{@{}lrrr@{}}
\toprule
\textbf{Category} & \textbf{Value 1} & \textbf{Value 2} & \textbf{Total} \\
\midrule
Item A & 10.5 & 20.3 & 30.8 \\
Item B & 15.2 & 25.1 & 40.3 \\
Item C & 22.1 & 12.7 & 34.8 \\
\bottomrule
\end{tabular}
\caption{Sample Data Table}
\label{tab:sample}
\end{table}

\section{Lists}
\subsection{Bullet Points}
\begin{itemize}
  \item First item
  \item Second item
  \item Third item with sub-items:
    \begin{itemize}
      \item Sub-item 1
      \item Sub-item 2
    \end{itemize}
\end{itemize}

\subsection{Numbered List}
\begin{enumerate}
  \item First step
  \item Second step
  \item Third step
\end{enumerate}

\section{Citations and References}
When you use the citation search feature, references will be automatically managed. For example, you might cite a paper like this: \cite{example2024}. All citations will be automatically added to your \texttt{cite.bib} file.

\section{Conclusion}
This template demonstrates basic \LaTeX{} features including mathematical formulas, tables, and lists. The bibliography system is now set up to work seamlessly with Bibby AI's citation management. Feel free to modify it for your needs!

\bibliographystyle{plain}
\bibliography{cite}

\end{document}\documentclass{article}
\usepackage[utf8]{inputenc}
\usepackage{booktabs}
\usepackage{amsmath}
\usepackage{natbib}

\title{My Document}
\author{Your Name}
\date{\today}

\begin{document}

\maketitle

\section{Introduction}
This is a sample document demonstrating various \LaTeX{} features including mathematical equations, tables, and formatting. You can cite papers using the citation search feature, and they will be automatically added to your \texttt{cite.bib} file.

\section{Mathematical Expressions}
Here are some example equations:

\subsection{Quadratic Formula}
The quadratic formula is given by:
\[x = \frac{-b \pm \sqrt{b^2 - 4ac}}{2a}\]

\subsection{Matrix Example}
A sample matrix:
\[
\begin{pmatrix}
a_{11} & a_{12} & a_{13} \\
a_{21} & a_{22} & a_{23} \\
a_{31} & a_{32} & a_{33}
\end{pmatrix}
\]

\section{Tables}
Below is a sample table showing data:

\begin{table}[h]
\centering
\begin{tabular}{@{}lrrr@{}}
\toprule
\textbf{Category} & \textbf{Value 1} & \textbf{Value 2} & \textbf{Total} \\
\midrule
Item A & 10.5 & 20.3 & 30.8 \\
Item B & 15.2 & 25.1 & 40.3 \\
Item C & 22.1 & 12.7 & 34.8 \\
\bottomrule
\end{tabular}
\caption{Sample Data Table}
\label{tab:sample}
\end{table}

\section{Lists}
\subsection{Bullet Points}
\begin{itemize}
  \item First item
  \item Second item
  \item Third item with sub-items:
    \begin{itemize}
      \item Sub-item 1
      \item Sub-item 2
    \end{itemize}
\end{itemize}

\subsection{Numbered List}
\begin{enumerate}
  \item First step
  \item Second step
  \item Third step
\end{enumerate}

\section{Citations and References}
When you use the citation search feature, references will be automatically managed. For example, you might cite a paper like this: \cite{example2024}. All citations will be automatically added to your \texttt{cite.bib} file.

\section{Conclusion}
This template demonstrates basic \LaTeX{} features including mathematical formulas, tables, and lists. The bibliography system is now set up to work seamlessly with Bibby AI's citation management. Feel free to modify it for your needs!

\bibliographystyle{plain}
\bibliography{cite}

\end{document}\documentclass{article}
\usepackage[utf8]{inputenc}
\usepackage{booktabs}
\usepackage{amsmath}
\usepackage{natbib}

\title{My Document}
\author{Your Name}
\date{\today}

\begin{document}

\maketitle

\section{Introduction}
This is a sample document demonstrating various \LaTeX{} features including mathematical equations, tables, and formatting. You can cite papers using the citation search feature, and they will be automatically added to your \texttt{cite.bib} file.

\section{Mathematical Expressions}
Here are some example equations:

\subsection{Quadratic Formula}
The quadratic formula is given by:
\[x = \frac{-b \pm \sqrt{b^2 - 4ac}}{2a}\]

\subsection{Matrix Example}
A sample matrix:
\[
\begin{pmatrix}
a_{11} & a_{12} & a_{13} \\
a_{21} & a_{22} & a_{23} \\
a_{31} & a_{32} & a_{33}
\end{pmatrix}
\]

\section{Tables}
Below is a sample table showing data:

\begin{table}[h]
\centering
\begin{tabular}{@{}lrrr@{}}
\toprule
\textbf{Category} & \textbf{Value 1} & \textbf{Value 2} & \textbf{Total} \\
\midrule
Item A & 10.5 & 20.3 & 30.8 \\
Item B & 15.2 & 25.1 & 40.3 \\
Item C & 22.1 & 12.7 & 34.8 \\
\bottomrule
\end{tabular}
\caption{Sample Data Table}
\label{tab:sample}
\end{table}

\section{Lists}
\subsection{Bullet Points}
\begin{itemize}
  \item First item
  \item Second item
  \item Third item with sub-items:
    \begin{itemize}
      \item Sub-item 1
      \item Sub-item 2
    \end{itemize}
\end{itemize}

\subsection{Numbered List}
\begin{enumerate}
  \item First step
  \item Second step
  \item Third step
\end{enumerate}

\section{Citations and References}
When you use the citation search feature, references will be automatically managed. For example, you might cite a paper like this: \cite{example2024}. All citations will be automatically added to your \texttt{cite.bib} file.

\section{Conclusion}
This template demonstrates basic \LaTeX{} features including mathematical formulas, tables, and lists. The bibliography system is now set up to work seamlessly with Bibby AI's citation management. Feel free to modify it for your needs!

\bibliographystyle{plain}
\bibliography{cite}

\end{document}\documentclass{article}
\usepackage[utf8]{inputenc}
\usepackage{booktabs}
\usepackage{amsmath}
\usepackage{natbib}

\title{My Document}
\author{Your Name}
\date{\today}

\begin{document}

\maketitle

\section{Introduction}
This is a sample document demonstrating various \LaTeX{} features including mathematical equations, tables, and formatting. You can cite papers using the citation search feature, and they will be automatically added to your \texttt{cite.bib} file.

\section{Mathematical Expressions}
Here are some example equations:

\subsection{Quadratic Formula}
The quadratic formula is given by:
\[x = \frac{-b \pm \sqrt{b^2 - 4ac}}{2a}\]

\subsection{Matrix Example}
A sample matrix:
\[
\begin{pmatrix}
a_{11} & a_{12} & a_{13} \\
a_{21} & a_{22} & a_{23} \\
a_{31} & a_{32} & a_{33}
\end{pmatrix}
\]

\section{Tables}
Below is a sample table showing data:

\begin{table}[h]
\centering
\begin{tabular}{@{}lrrr@{}}
\toprule
\textbf{Category} & \textbf{Value 1} & \textbf{Value 2} & \textbf{Total} \\
\midrule
Item A & 10.5 & 20.3 & 30.8 \\
Item B & 15.2 & 25.1 & 40.3 \\
Item C & 22.1 & 12.7 & 34.8 \\
\bottomrule
\end{tabular}
\caption{Sample Data Table}
\label{tab:sample}
\end{table}

\section{Lists}
\subsection{Bullet Points}
\begin{itemize}
  \item First item
  \item Second item
  \item Third item with sub-items:
    \begin{itemize}
      \item Sub-item 1
      \item Sub-item 2
    \end{itemize}
\end{itemize}

\subsection{Numbered List}
\begin{enumerate}
  \item First step
  \item Second step
  \item Third step
\end{enumerate}

\section{Citations and References}
When you use the citation search feature, references will be automatically managed. For example, you might cite a paper like this: \cite{example2024}. All citations will be automatically added to your \texttt{cite.bib} file.

\section{Conclusion}
This template demonstrates basic \LaTeX{} features including mathematical formulas, tables, and lists. The bibliography system is now set up to work seamlessly with Bibby AI's citation management. Feel free to modify it for your needs!

\bibliographystyle{plain}
\bibliography{cite}

\end{document}\documentclass{article}
\usepackage[utf8]{inputenc}
\usepackage{booktabs}
\usepackage{amsmath}
\usepackage{natbib}

\title{My Document}
\author{Your Name}
\date{\today}

\begin{document}

\maketitle

\section{Introduction}
This is a sample document demonstrating various \LaTeX{} features including mathematical equations, tables, and formatting. You can cite papers using the citation search feature, and they will be automatically added to your \texttt{cite.bib} file.

\section{Mathematical Expressions}
Here are some example equations:

\subsection{Quadratic Formula}
The quadratic formula is given by:
\[x = \frac{-b \pm \sqrt{b^2 - 4ac}}{2a}\]

\subsection{Matrix Example}
A sample matrix:
\[
\begin{pmatrix}
a_{11} & a_{12} & a_{13} \\
a_{21} & a_{22} & a_{23} \\
a_{31} & a_{32} & a_{33}
\end{pmatrix}
\]

\section{Tables}
Below is a sample table showing data:

\begin{table}[h]
\centering
\begin{tabular}{@{}lrrr@{}}
\toprule
\textbf{Category} & \textbf{Value 1} & \textbf{Value 2} & \textbf{Total} \\
\midrule
Item A & 10.5 & 20.3 & 30.8 \\
Item B & 15.2 & 25.1 & 40.3 \\
Item C & 22.1 & 12.7 & 34.8 \\
\bottomrule
\end{tabular}
\caption{Sample Data Table}
\label{tab:sample}
\end{table}

\section{Lists}
\subsection{Bullet Points}
\begin{itemize}
  \item First item
  \item Second item
  \item Third item with sub-items:
    \begin{itemize}
      \item Sub-item 1
      \item Sub-item 2
    \end{itemize}
\end{itemize}

\subsection{Numbered List}
\begin{enumerate}
  \item First step
  \item Second step
  \item Third step
\end{enumerate}

\section{Citations and References}
When you use the citation search feature, references will be automatically managed. For example, you might cite a paper like this: \cite{example2024}. All citations will be automatically added to your \texttt{cite.bib} file.

\section{Conclusion}
This template demonstrates basic \LaTeX{} features including mathematical formulas, tables, and lists. The bibliography system is now set up to work seamlessly with Bibby AI's citation management. Feel free to modify it for your needs!

\bibliographystyle{plain}
\bibliography{cite}

\end{document}\documentclass{article}
\usepackage[utf8]{inputenc}
\usepackage{booktabs}
\usepackage{amsmath}
\usepackage{natbib}

\title{My Document}
\author{Your Name}
\date{\today}

\begin{document}

\maketitle

\section{Introduction}
This is a sample document demonstrating various \LaTeX{} features including mathematical equations, tables, and formatting. You can cite papers using the citation search feature, and they will be automatically added to your \texttt{cite.bib} file.

\section{Mathematical Expressions}
Here are some example equations:

\subsection{Quadratic Formula}
The quadratic formula is given by:
\[x = \frac{-b \pm \sqrt{b^2 - 4ac}}{2a}\]

\subsection{Matrix Example}
A sample matrix:
\[
\begin{pmatrix}
a_{11} & a_{12} & a_{13} \\
a_{21} & a_{22} & a_{23} \\
a_{31} & a_{32} & a_{33}
\end{pmatrix}
\]

\section{Tables}
Below is a sample table showing data:

\begin{table}[h]
\centering
\begin{tabular}{@{}lrrr@{}}
\toprule
\textbf{Category} & \textbf{Value 1} & \textbf{Value 2} & \textbf{Total} \\
\midrule
Item A & 10.5 & 20.3 & 30.8 \\
Item B & 15.2 & 25.1 & 40.3 \\
Item C & 22.1 & 12.7 & 34.8 \\
\bottomrule
\end{tabular}
\caption{Sample Data Table}
\label{tab:sample}
\end{table}

\section{Lists}
\subsection{Bullet Points}
\begin{itemize}
  \item First item
  \item Second item
  \item Third item with sub-items:
    \begin{itemize}
      \item Sub-item 1
      \item Sub-item 2
    \end{itemize}
\end{itemize}

\subsection{Numbered List}
\begin{enumerate}
  \item First step
  \item Second step
  \item Third step
\end{enumerate}

\section{Citations and References}
When you use the citation search feature, references will be automatically managed. For example, you might cite a paper like this: \cite{example2024}. All citations will be automatically added to your \texttt{cite.bib} file.

\section{Conclusion}
This template demonstrates basic \LaTeX{} features including mathematical formulas, tables, and lists. The bibliography system is now set up to work seamlessly with Bibby AI's citation management. Feel free to modify it for your needs!

\bibliographystyle{plain}
\bibliography{cite}

\end{document}\documentclass{article}
\usepackage[utf8]{inputenc}
\usepackage{booktabs}
\usepackage{amsmath}
\usepackage{natbib}

\title{My Document}
\author{Your Name}
\date{\today}

\begin{document}

\maketitle

\section{Introduction}
This is a sample document demonstrating various \LaTeX{} features including mathematical equations, tables, and formatting. You can cite papers using the citation search feature, and they will be automatically added to your \texttt{cite.bib} file.

\section{Mathematical Expressions}
Here are some example equations:

\subsection{Quadratic Formula}
The quadratic formula is given by:
\[x = \frac{-b \pm \sqrt{b^2 - 4ac}}{2a}\]

\subsection{Matrix Example}
A sample matrix:
\[
\begin{pmatrix}
a_{11} & a_{12} & a_{13} \\
a_{21} & a_{22} & a_{23} \\
a_{31} & a_{32} & a_{33}
\end{pmatrix}
\]

\section{Tables}
Below is a sample table showing data:

\begin{table}[h]
\centering
\begin{tabular}{@{}lrrr@{}}
\toprule
\textbf{Category} & \textbf{Value 1} & \textbf{Value 2} & \textbf{Total} \\
\midrule
Item A & 10.5 & 20.3 & 30.8 \\
Item B & 15.2 & 25.1 & 40.3 \\
Item C & 22.1 & 12.7 & 34.8 \\
\bottomrule
\end{tabular}
\caption{Sample Data Table}
\label{tab:sample}
\end{table}

\section{Lists}
\subsection{Bullet Points}
\begin{itemize}
  \item First item
  \item Second item
  \item Third item with sub-items:
    \begin{itemize}
      \item Sub-item 1
      \item Sub-item 2
    \end{itemize}
\end{itemize}

\subsection{Numbered List}
\begin{enumerate}
  \item First step
  \item Second step
  \item Third step
\end{enumerate}

\section{Citations and References}
When you use the citation search feature, references will be automatically managed. For example, you might cite a paper like this: \cite{example2024}. All citations will be automatically added to your \texttt{cite.bib} file.

\section{Conclusion}
This template demonstrates basic \LaTeX{} features including mathematical formulas, tables, and lists. The bibliography system is now set up to work seamlessly with Bibby AI's citation management. Feel free to modify it for your needs!

\bibliographystyle{plain}
\bibliography{cite}

\end{document}\documentclass{article}
\usepackage[utf8]{inputenc}
\usepackage{booktabs}
\usepackage{amsmath}
\usepackage{natbib}

\title{My Document}
\author{Your Name}
\date{\today}

\begin{document}

\maketitle

\section{Introduction}
This is a sample document demonstrating various \LaTeX{} features including mathematical equations, tables, and formatting. You can cite papers using the citation search feature, and they will be automatically added to your \texttt{cite.bib} file.

\section{Mathematical Expressions}
Here are some example equations:

\subsection{Quadratic Formula}
The quadratic formula is given by:
\[x = \frac{-b \pm \sqrt{b^2 - 4ac}}{2a}\]

\subsection{Matrix Example}
A sample matrix:
\[
\begin{pmatrix}
a_{11} & a_{12} & a_{13} \\
a_{21} & a_{22} & a_{23} \\
a_{31} & a_{32} & a_{33}
\end{pmatrix}
\]

\section{Tables}
Below is a sample table showing data:

\begin{table}[h]
\centering
\begin{tabular}{@{}lrrr@{}}
\toprule
\textbf{Category} & \textbf{Value 1} & \textbf{Value 2} & \textbf{Total} \\
\midrule
Item A & 10.5 & 20.3 & 30.8 \\
Item B & 15.2 & 25.1 & 40.3 \\
Item C & 22.1 & 12.7 & 34.8 \\
\bottomrule
\end{tabular}
\caption{Sample Data Table}
\label{tab:sample}
\end{table}

\section{Lists}
\subsection{Bullet Points}
\begin{itemize}
  \item First item
  \item Second item
  \item Third item with sub-items:
    \begin{itemize}
      \item Sub-item 1
      \item Sub-item 2
    \end{itemize}
\end{itemize}

\subsection{Numbered List}
\begin{enumerate}
  \item First step
  \item Second step
  \item Third step
\end{enumerate}

\section{Citations and References}
When you use the citation search feature, references will be automatically managed. For example, you might cite a paper like this: \cite{example2024}. All citations will be automatically added to your \texttt{cite.bib} file.

\section{Conclusion}
This template demonstrates basic \LaTeX{} features including mathematical formulas, tables, and lists. The bibliography system is now set up to work seamlessly with Bibby AI's citation management. Feel free to modify it for your needs!

\bibliographystyle{plain}
\bibliography{cite}

\end{document}\documentclass{article}
\usepackage[utf8]{inputenc}
\usepackage{booktabs}
\usepackage{amsmath}
\usepackage{natbib}

\title{My Document}
\author{Your Name}
\date{\today}

\begin{document}

\maketitle

\section{Introduction}
This is a sample document demonstrating various \LaTeX{} features including mathematical equations, tables, and formatting. You can cite papers using the citation search feature, and they will be automatically added to your \texttt{cite.bib} file.

\section{Mathematical Expressions}
Here are some example equations:

\subsection{Quadratic Formula}
The quadratic formula is given by:
\[x = \frac{-b \pm \sqrt{b^2 - 4ac}}{2a}\]

\subsection{Matrix Example}
A sample matrix:
\[
\begin{pmatrix}
a_{11} & a_{12} & a_{13} \\
a_{21} & a_{22} & a_{23} \\
a_{31} & a_{32} & a_{33}
\end{pmatrix}
\]

\section{Tables}
Below is a sample table showing data:

\begin{table}[h]
\centering
\begin{tabular}{@{}lrrr@{}}
\toprule
\textbf{Category} & \textbf{Value 1} & \textbf{Value 2} & \textbf{Total} \\
\midrule
Item A & 10.5 & 20.3 & 30.8 \\
Item B & 15.2 & 25.1 & 40.3 \\
Item C & 22.1 & 12.7 & 34.8 \\
\bottomrule
\end{tabular}
\caption{Sample Data Table}
\label{tab:sample}
\end{table}

\section{Lists}
\subsection{Bullet Points}
\begin{itemize}
  \item First item
  \item Second item
  \item Third item with sub-items:
    \begin{itemize}
      \item Sub-item 1
      \item Sub-item 2
    \end{itemize}
\end{itemize}

\subsection{Numbered List}
\begin{enumerate}
  \item First step
  \item Second step
  \item Third step
\end{enumerate}

\section{Citations and References}
When you use the citation search feature, references will be automatically managed. For example, you might cite a paper like this: \cite{example2024}. All citations will be automatically added to your \texttt{cite.bib} file.

\section{Conclusion}
This template demonstrates basic \LaTeX{} features including mathematical formulas, tables, and lists. The bibliography system is now set up to work seamlessly with Bibby AI's citation management. Feel free to modify it for your needs!

\bibliographystyle{plain}
\bibliography{cite}

\end{document}\documentclass{article}
\usepackage[utf8]{inputenc}
\usepackage{booktabs}
\usepackage{amsmath}
\usepackage{natbib}

\title{My Document}
\author{Your Name}
\date{\today}

\begin{document}

\maketitle

\section{Introduction}
This is a sample document demonstrating various \LaTeX{} features including mathematical equations, tables, and formatting. You can cite papers using the citation search feature, and they will be automatically added to your \texttt{cite.bib} file.

\section{Mathematical Expressions}
Here are some example equations:

\subsection{Quadratic Formula}
The quadratic formula is given by:
\[x = \frac{-b \pm \sqrt{b^2 - 4ac}}{2a}\]

\subsection{Matrix Example}
A sample matrix:
\[
\begin{pmatrix}
a_{11} & a_{12} & a_{13} \\
a_{21} & a_{22} & a_{23} \\
a_{31} & a_{32} & a_{33}
\end{pmatrix}
\]

\section{Tables}
Below is a sample table showing data:

\begin{table}[h]
\centering
\begin{tabular}{@{}lrrr@{}}
\toprule
\textbf{Category} & \textbf{Value 1} & \textbf{Value 2} & \textbf{Total} \\
\midrule
Item A & 10.5 & 20.3 & 30.8 \\
Item B & 15.2 & 25.1 & 40.3 \\
Item C & 22.1 & 12.7 & 34.8 \\
\bottomrule
\end{tabular}
\caption{Sample Data Table}
\label{tab:sample}
\end{table}

\section{Lists}
\subsection{Bullet Points}
\begin{itemize}
  \item First item
  \item Second item
  \item Third item with sub-items:
    \begin{itemize}
      \item Sub-item 1
      \item Sub-item 2
    \end{itemize}
\end{itemize}

\subsection{Numbered List}
\begin{enumerate}
  \item First step
  \item Second step
  \item Third step
\end{enumerate}

\section{Citations and References}
When you use the citation search feature, references will be automatically managed. For example, you might cite a paper like this: \cite{example2024}. All citations will be automatically added to your \texttt{cite.bib} file.

\section{Conclusion}
This template demonstrates basic \LaTeX{} features including mathematical formulas, tables, and lists. The bibliography system is now set up to work seamlessly with Bibby AI's citation management. Feel free to modify it for your needs!

\bibliographystyle{plain}
\bibliography{cite}

\end{document}\documentclass{article}
\usepackage[utf8]{inputenc}
\usepackage{booktabs}
\usepackage{amsmath}
\usepackage{natbib}

\title{My Document}
\author{Your Name}
\date{\today}

\begin{document}

\maketitle

\section{Introduction}
This is a sample document demonstrating various \LaTeX{} features including mathematical equations, tables, and formatting. You can cite papers using the citation search feature, and they will be automatically added to your \texttt{cite.bib} file.

\section{Mathematical Expressions}
Here are some example equations:

\subsection{Quadratic Formula}
The quadratic formula is given by:
\[x = \frac{-b \pm \sqrt{b^2 - 4ac}}{2a}\]

\subsection{Matrix Example}
A sample matrix:
\[
\begin{pmatrix}
a_{11} & a_{12} & a_{13} \\
a_{21} & a_{22} & a_{23} \\
a_{31} & a_{32} & a_{33}
\end{pmatrix}
\]

\section{Tables}
Below is a sample table showing data:

\begin{table}[h]
\centering
\begin{tabular}{@{}lrrr@{}}
\toprule
\textbf{Category} & \textbf{Value 1} & \textbf{Value 2} & \textbf{Total} \\
\midrule
Item A & 10.5 & 20.3 & 30.8 \\
Item B & 15.2 & 25.1 & 40.3 \\
Item C & 22.1 & 12.7 & 34.8 \\
\bottomrule
\end{tabular}
\caption{Sample Data Table}
\label{tab:sample}
\end{table}

\section{Lists}
\subsection{Bullet Points}
\begin{itemize}
  \item First item
  \item Second item
  \item Third item with sub-items:
    \begin{itemize}
      \item Sub-item 1
      \item Sub-item 2
    \end{itemize}
\end{itemize}

\subsection{Numbered List}
\begin{enumerate}
  \item First step
  \item Second step
  \item Third step
\end{enumerate}

\section{Citations and References}
When you use the citation search feature, references will be automatically managed. For example, you might cite a paper like this: \cite{example2024}. All citations will be automatically added to your \texttt{cite.bib} file.

\section{Conclusion}
This template demonstrates basic \LaTeX{} features including mathematical formulas, tables, and lists. The bibliography system is now set up to work seamlessly with Bibby AI's citation management. Feel free to modify it for your needs!

\bibliographystyle{plain}
\bibliography{cite}

\end{document}\documentclass{article}
\usepackage[utf8]{inputenc}
\usepackage{booktabs}
\usepackage{amsmath}
\usepackage{natbib}

\title{My Document}
\author{Your Name}
\date{\today}

\begin{document}

\maketitle

\section{Introduction}
This is a sample document demonstrating various \LaTeX{} features including mathematical equations, tables, and formatting. You can cite papers using the citation search feature, and they will be automatically added to your \texttt{cite.bib} file.

\section{Mathematical Expressions}
Here are some example equations:

\subsection{Quadratic Formula}
The quadratic formula is given by:
\[x = \frac{-b \pm \sqrt{b^2 - 4ac}}{2a}\]

\subsection{Matrix Example}
A sample matrix:
\[
\begin{pmatrix}
a_{11} & a_{12} & a_{13} \\
a_{21} & a_{22} & a_{23} \\
a_{31} & a_{32} & a_{33}
\end{pmatrix}
\]

\section{Tables}
Below is a sample table showing data:

\begin{table}[h]
\centering
\begin{tabular}{@{}lrrr@{}}
\toprule
\textbf{Category} & \textbf{Value 1} & \textbf{Value 2} & \textbf{Total} \\
\midrule
Item A & 10.5 & 20.3 & 30.8 \\
Item B & 15.2 & 25.1 & 40.3 \\
Item C & 22.1 & 12.7 & 34.8 \\
\bottomrule
\end{tabular}
\caption{Sample Data Table}
\label{tab:sample}
\end{table}

\section{Lists}
\subsection{Bullet Points}
\begin{itemize}
  \item First item
  \item Second item
  \item Third item with sub-items:
    \begin{itemize}
      \item Sub-item 1
      \item Sub-item 2
    \end{itemize}
\end{itemize}

\subsection{Numbered List}
\begin{enumerate}
  \item First step
  \item Second step
  \item Third step
\end{enumerate}

\section{Citations and References}
When you use the citation search feature, references will be automatically managed. For example, you might cite a paper like this: \cite{example2024}. All citations will be automatically added to your \texttt{cite.bib} file.

\section{Conclusion}
This template demonstrates basic \LaTeX{} features including mathematical formulas, tables, and lists. The bibliography system is now set up to work seamlessly with Bibby AI's citation management. Feel free to modify it for your needs!

\bibliographystyle{plain}
\bibliography{cite}

\end{document}\documentclass{article}
\usepackage[utf8]{inputenc}
\usepackage{booktabs}
\usepackage{amsmath}
\usepackage{natbib}

\title{My Document}
\author{Your Name}
\date{\today}

\begin{document}

\maketitle

\section{Introduction}
This is a sample document demonstrating various \LaTeX{} features including mathematical equations, tables, and formatting. You can cite papers using the citation search feature, and they will be automatically added to your \texttt{cite.bib} file.

\section{Mathematical Expressions}
Here are some example equations:

\subsection{Quadratic Formula}
The quadratic formula is given by:
\[x = \frac{-b \pm \sqrt{b^2 - 4ac}}{2a}\]

\subsection{Matrix Example}
A sample matrix:
\[
\begin{pmatrix}
a_{11} & a_{12} & a_{13} \\
a_{21} & a_{22} & a_{23} \\
a_{31} & a_{32} & a_{33}
\end{pmatrix}
\]

\section{Tables}
Below is a sample table showing data:

\begin{table}[h]
\centering
\begin{tabular}{@{}lrrr@{}}
\toprule
\textbf{Category} & \textbf{Value 1} & \textbf{Value 2} & \textbf{Total} \\
\midrule
Item A & 10.5 & 20.3 & 30.8 \\
Item B & 15.2 & 25.1 & 40.3 \\
Item C & 22.1 & 12.7 & 34.8 \\
\bottomrule
\end{tabular}
\caption{Sample Data Table}
\label{tab:sample}
\end{table}

\section{Lists}
\subsection{Bullet Points}
\begin{itemize}
  \item First item
  \item Second item
  \item Third item with sub-items:
    \begin{itemize}
      \item Sub-item 1
      \item Sub-item 2
    \end{itemize}
\end{itemize}

\subsection{Numbered List}
\begin{enumerate}
  \item First step
  \item Second step
  \item Third step
\end{enumerate}

\section{Citations and References}
When you use the citation search feature, references will be automatically managed. For example, you might cite a paper like this: \cite{example2024}. All citations will be automatically added to your \texttt{cite.bib} file.

\section{Conclusion}
This template demonstrates basic \LaTeX{} features including mathematical formulas, tables, and lists. The bibliography system is now set up to work seamlessly with Bibby AI's citation management. Feel free to modify it for your needs!

\bibliographystyle{plain}
\bibliography{cite}

\end{document}\documentclass{article}
\usepackage[utf8]{inputenc}
\usepackage{booktabs}
\usepackage{amsmath}
\usepackage{natbib}

\title{My Document}
\author{Your Name}
\date{\today}

\begin{document}

\maketitle

\section{Introduction}
This is a sample document demonstrating various \LaTeX{} features including mathematical equations, tables, and formatting. You can cite papers using the citation search feature, and they will be automatically added to your \texttt{cite.bib} file.

\section{Mathematical Expressions}
Here are some example equations:

\subsection{Quadratic Formula}
The quadratic formula is given by:
\[x = \frac{-b \pm \sqrt{b^2 - 4ac}}{2a}\]

\subsection{Matrix Example}
A sample matrix:
\[
\begin{pmatrix}
a_{11} & a_{12} & a_{13} \\
a_{21} & a_{22} & a_{23} \\
a_{31} & a_{32} & a_{33}
\end{pmatrix}
\]

\section{Tables}
Below is a sample table showing data:

\begin{table}[h]
\centering
\begin{tabular}{@{}lrrr@{}}
\toprule
\textbf{Category} & \textbf{Value 1} & \textbf{Value 2} & \textbf{Total} \\
\midrule
Item A & 10.5 & 20.3 & 30.8 \\
Item B & 15.2 & 25.1 & 40.3 \\
Item C & 22.1 & 12.7 & 34.8 \\
\bottomrule
\end{tabular}
\caption{Sample Data Table}
\label{tab:sample}
\end{table}

\section{Lists}
\subsection{Bullet Points}
\begin{itemize}
  \item First item
  \item Second item
  \item Third item with sub-items:
    \begin{itemize}
      \item Sub-item 1
      \item Sub-item 2
    \end{itemize}
\end{itemize}

\subsection{Numbered List}
\begin{enumerate}
  \item First step
  \item Second step
  \item Third step
\end{enumerate}

\section{Citations and References}
When you use the citation search feature, references will be automatically managed. For example, you might cite a paper like this: \cite{example2024}. All citations will be automatically added to your \texttt{cite.bib} file.

\section{Conclusion}
This template demonstrates basic \LaTeX{} features including mathematical formulas, tables, and lists. The bibliography system is now set up to work seamlessly with Bibby AI's citation management. Feel free to modify it for your needs!

\bibliographystyle{plain}
\bibliography{cite}

\end{document}\documentclass{article}
\usepackage[utf8]{inputenc}
\usepackage{booktabs}
\usepackage{amsmath}
\usepackage{natbib}

\title{My Document}
\author{Your Name}
\date{\today}

\begin{document}

\maketitle

\section{Introduction}
This is a sample document demonstrating various \LaTeX{} features including mathematical equations, tables, and formatting. You can cite papers using the citation search feature, and they will be automatically added to your \texttt{cite.bib} file.

\section{Mathematical Expressions}
Here are some example equations:

\subsection{Quadratic Formula}
The quadratic formula is given by:
\[x = \frac{-b \pm \sqrt{b^2 - 4ac}}{2a}\]

\subsection{Matrix Example}
A sample matrix:
\[
\begin{pmatrix}
a_{11} & a_{12} & a_{13} \\
a_{21} & a_{22} & a_{23} \\
a_{31} & a_{32} & a_{33}
\end{pmatrix}
\]

\section{Tables}
Below is a sample table showing data:

\begin{table}[h]
\centering
\begin{tabular}{@{}lrrr@{}}
\toprule
\textbf{Category} & \textbf{Value 1} & \textbf{Value 2} & \textbf{Total} \\
\midrule
Item A & 10.5 & 20.3 & 30.8 \\
Item B & 15.2 & 25.1 & 40.3 \\
Item C & 22.1 & 12.7 & 34.8 \\
\bottomrule
\end{tabular}
\caption{Sample Data Table}
\label{tab:sample}
\end{table}

\section{Lists}
\subsection{Bullet Points}
\begin{itemize}
  \item First item
  \item Second item
  \item Third item with sub-items:
    \begin{itemize}
      \item Sub-item 1
      \item Sub-item 2
    \end{itemize}
\end{itemize}

\subsection{Numbered List}
\begin{enumerate}
  \item First step
  \item Second step
  \item Third step
\end{enumerate}

\section{Citations and References}
When you use the citation search feature, references will be automatically managed. For example, you might cite a paper like this: \cite{example2024}. All citations will be automatically added to your \texttt{cite.bib} file.

\section{Conclusion}
This template demonstrates basic \LaTeX{} features including mathematical formulas, tables, and lists. The bibliography system is now set up to work seamlessly with Bibby AI's citation management. Feel free to modify it for your needs!

\bibliographystyle{plain}
\bibliography{cite}

\end{document}\documentclass{article}
\usepackage[utf8]{inputenc}
\usepackage{booktabs}
\usepackage{amsmath}
\usepackage{natbib}

\title{My Document}
\author{Your Name}
\date{\today}

\begin{document}

\maketitle

\section{Introduction}
This is a sample document demonstrating various \LaTeX{} features including mathematical equations, tables, and formatting. You can cite papers using the citation search feature, and they will be automatically added to your \texttt{cite.bib} file.

\section{Mathematical Expressions}
Here are some example equations:

\subsection{Quadratic Formula}
The quadratic formula is given by:
\[x = \frac{-b \pm \sqrt{b^2 - 4ac}}{2a}\]

\subsection{Matrix Example}
A sample matrix:
\[
\begin{pmatrix}
a_{11} & a_{12} & a_{13} \\
a_{21} & a_{22} & a_{23} \\
a_{31} & a_{32} & a_{33}
\end{pmatrix}
\]

\section{Tables}
Below is a sample table showing data:

\begin{table}[h]
\centering
\begin{tabular}{@{}lrrr@{}}
\toprule
\textbf{Category} & \textbf{Value 1} & \textbf{Value 2} & \textbf{Total} \\
\midrule
Item A & 10.5 & 20.3 & 30.8 \\
Item B & 15.2 & 25.1 & 40.3 \\
Item C & 22.1 & 12.7 & 34.8 \\
\bottomrule
\end{tabular}
\caption{Sample Data Table}
\label{tab:sample}
\end{table}

\section{Lists}
\subsection{Bullet Points}
\begin{itemize}
  \item First item
  \item Second item
  \item Third item with sub-items:
    \begin{itemize}
      \item Sub-item 1
      \item Sub-item 2
    \end{itemize}
\end{itemize}

\subsection{Numbered List}
\begin{enumerate}
  \item First step
  \item Second step
  \item Third step
\end{enumerate}

\section{Citations and References}
When you use the citation search feature, references will be automatically managed. For example, you might cite a paper like this: \cite{example2024}. All citations will be automatically added to your \texttt{cite.bib} file.

\section{Conclusion}
This template demonstrates basic \LaTeX{} features including mathematical formulas, tables, and lists. The bibliography system is now set up to work seamlessly with Bibby AI's citation management. Feel free to modify it for your needs!

\bibliographystyle{plain}
\bibliography{cite}

\end{document}\documentclass{article}
\usepackage[utf8]{inputenc}
\usepackage{booktabs}
\usepackage{amsmath}
\usepackage{natbib}

\title{My Document}
\author{Your Name}
\date{\today}

\begin{document}

\maketitle

\section{Introduction}
This is a sample document demonstrating various \LaTeX{} features including mathematical equations, tables, and formatting. You can cite papers using the citation search feature, and they will be automatically added to your \texttt{cite.bib} file.

\section{Mathematical Expressions}
Here are some example equations:

\subsection{Quadratic Formula}
The quadratic formula is given by:
\[x = \frac{-b \pm \sqrt{b^2 - 4ac}}{2a}\]

\subsection{Matrix Example}
A sample matrix:
\[
\begin{pmatrix}
a_{11} & a_{12} & a_{13} \\
a_{21} & a_{22} & a_{23} \\
a_{31} & a_{32} & a_{33}
\end{pmatrix}
\]

\section{Tables}
Below is a sample table showing data:

\begin{table}[h]
\centering
\begin{tabular}{@{}lrrr@{}}
\toprule
\textbf{Category} & \textbf{Value 1} & \textbf{Value 2} & \textbf{Total} \\
\midrule
Item A & 10.5 & 20.3 & 30.8 \\
Item B & 15.2 & 25.1 & 40.3 \\
Item C & 22.1 & 12.7 & 34.8 \\
\bottomrule
\end{tabular}
\caption{Sample Data Table}
\label{tab:sample}
\end{table}

\section{Lists}
\subsection{Bullet Points}
\begin{itemize}
  \item First item
  \item Second item
  \item Third item with sub-items:
    \begin{itemize}
      \item Sub-item 1
      \item Sub-item 2
    \end{itemize}
\end{itemize}

\subsection{Numbered List}
\begin{enumerate}
  \item First step
  \item Second step
  \item Third step
\end{enumerate}

\section{Citations and References}
When you use the citation search feature, references will be automatically managed. For example, you might cite a paper like this: \cite{example2024}. All citations will be automatically added to your \texttt{cite.bib} file.

\section{Conclusion}
This template demonstrates basic \LaTeX{} features including mathematical formulas, tables, and lists. The bibliography system is now set up to work seamlessly with Bibby AI's citation management. Feel free to modify it for your needs!

\bibliographystyle{plain}
\bibliography{cite}

\end{document}\documentclass{article}
\usepackage[utf8]{inputenc}
\usepackage{booktabs}
\usepackage{amsmath}
\usepackage{natbib}

\title{My Document}
\author{Your Name}
\date{\today}

\begin{document}

\maketitle

\section{Introduction}
This is a sample document demonstrating various \LaTeX{} features including mathematical equations, tables, and formatting. You can cite papers using the citation search feature, and they will be automatically added to your \texttt{cite.bib} file.

\section{Mathematical Expressions}
Here are some example equations:

\subsection{Quadratic Formula}
The quadratic formula is given by:
\[x = \frac{-b \pm \sqrt{b^2 - 4ac}}{2a}\]

\subsection{Matrix Example}
A sample matrix:
\[
\begin{pmatrix}
a_{11} & a_{12} & a_{13} \\
a_{21} & a_{22} & a_{23} \\
a_{31} & a_{32} & a_{33}
\end{pmatrix}
\]

\section{Tables}
Below is a sample table showing data:

\begin{table}[h]
\centering
\begin{tabular}{@{}lrrr@{}}
\toprule
\textbf{Category} & \textbf{Value 1} & \textbf{Value 2} & \textbf{Total} \\
\midrule
Item A & 10.5 & 20.3 & 30.8 \\
Item B & 15.2 & 25.1 & 40.3 \\
Item C & 22.1 & 12.7 & 34.8 \\
\bottomrule
\end{tabular}
\caption{Sample Data Table}
\label{tab:sample}
\end{table}

\section{Lists}
\subsection{Bullet Points}
\begin{itemize}
  \item First item
  \item Second item
  \item Third item with sub-items:
    \begin{itemize}
      \item Sub-item 1
      \item Sub-item 2
    \end{itemize}
\end{itemize}

\subsection{Numbered List}
\begin{enumerate}
  \item First step
  \item Second step
  \item Third step
\end{enumerate}

\section{Citations and References}
When you use the citation search feature, references will be automatically managed. For example, you might cite a paper like this: \cite{example2024}. All citations will be automatically added to your \texttt{cite.bib} file.

\section{Conclusion}
This template demonstrates basic \LaTeX{} features including mathematical formulas, tables, and lists. The bibliography system is now set up to work seamlessly with Bibby AI's citation management. Feel free to modify it for your needs!

\bibliographystyle{plain}
\bibliography{cite}

\end{document}\documentclass{article}
\usepackage[utf8]{inputenc}
\usepackage{booktabs}
\usepackage{amsmath}
\usepackage{natbib}

\title{My Document}
\author{Your Name}
\date{\today}

\begin{document}

\maketitle

\section{Introduction}
This is a sample document demonstrating various \LaTeX{} features including mathematical equations, tables, and formatting. You can cite papers using the citation search feature, and they will be automatically added to your \texttt{cite.bib} file.

\section{Mathematical Expressions}
Here are some example equations:

\subsection{Quadratic Formula}
The quadratic formula is given by:
\[x = \frac{-b \pm \sqrt{b^2 - 4ac}}{2a}\]

\subsection{Matrix Example}
A sample matrix:
\[
\begin{pmatrix}
a_{11} & a_{12} & a_{13} \\
a_{21} & a_{22} & a_{23} \\
a_{31} & a_{32} & a_{33}
\end{pmatrix}
\]

\section{Tables}
Below is a sample table showing data:

\begin{table}[h]
\centering
\begin{tabular}{@{}lrrr@{}}
\toprule
\textbf{Category} & \textbf{Value 1} & \textbf{Value 2} & \textbf{Total} \\
\midrule
Item A & 10.5 & 20.3 & 30.8 \\
Item B & 15.2 & 25.1 & 40.3 \\
Item C & 22.1 & 12.7 & 34.8 \\
\bottomrule
\end{tabular}
\caption{Sample Data Table}
\label{tab:sample}
\end{table}

\section{Lists}
\subsection{Bullet Points}
\begin{itemize}
  \item First item
  \item Second item
  \item Third item with sub-items:
    \begin{itemize}
      \item Sub-item 1
      \item Sub-item 2
    \end{itemize}
\end{itemize}

\subsection{Numbered List}
\begin{enumerate}
  \item First step
  \item Second step
  \item Third step
\end{enumerate}

\section{Citations and References}
When you use the citation search feature, references will be automatically managed. For example, you might cite a paper like this: \cite{example2024}. All citations will be automatically added to your \texttt{cite.bib} file.

\section{Conclusion}
This template demonstrates basic \LaTeX{} features including mathematical formulas, tables, and lists. The bibliography system is now set up to work seamlessly with Bibby AI's citation management. Feel free to modify it for your needs!

\bibliographystyle{plain}
\bibliography{cite}

\end{document}\documentclass{article}
\usepackage[utf8]{inputenc}
\usepackage{booktabs}
\usepackage{amsmath}
\usepackage{natbib}

\title{My Document}
\author{Your Name}
\date{\today}

\begin{document}

\maketitle

\section{Introduction}
This is a sample document demonstrating various \LaTeX{} features including mathematical equations, tables, and formatting. You can cite papers using the citation search feature, and they will be automatically added to your \texttt{cite.bib} file.

\section{Mathematical Expressions}
Here are some example equations:

\subsection{Quadratic Formula}
The quadratic formula is given by:
\[x = \frac{-b \pm \sqrt{b^2 - 4ac}}{2a}\]

\subsection{Matrix Example}
A sample matrix:
\[
\begin{pmatrix}
a_{11} & a_{12} & a_{13} \\
a_{21} & a_{22} & a_{23} \\
a_{31} & a_{32} & a_{33}
\end{pmatrix}
\]

\section{Tables}
Below is a sample table showing data:

\begin{table}[h]
\centering
\begin{tabular}{@{}lrrr@{}}
\toprule
\textbf{Category} & \textbf{Value 1} & \textbf{Value 2} & \textbf{Total} \\
\midrule
Item A & 10.5 & 20.3 & 30.8 \\
Item B & 15.2 & 25.1 & 40.3 \\
Item C & 22.1 & 12.7 & 34.8 \\
\bottomrule
\end{tabular}
\caption{Sample Data Table}
\label{tab:sample}
\end{table}

\section{Lists}
\subsection{Bullet Points}
\begin{itemize}
  \item First item
  \item Second item
  \item Third item with sub-items:
    \begin{itemize}
      \item Sub-item 1
      \item Sub-item 2
    \end{itemize}
\end{itemize}

\subsection{Numbered List}
\begin{enumerate}
  \item First step
  \item Second step
  \item Third step
\end{enumerate}

\section{Citations and References}
When you use the citation search feature, references will be automatically managed. For example, you might cite a paper like this: \cite{example2024}. All citations will be automatically added to your \texttt{cite.bib} file.

\section{Conclusion}
This template demonstrates basic \LaTeX{} features including mathematical formulas, tables, and lists. The bibliography system is now set up to work seamlessly with Bibby AI's citation management. Feel free to modify it for your needs!

\bibliographystyle{plain}
\bibliography{cite}

\end{document}\documentclass{article}
\usepackage[utf8]{inputenc}
\usepackage{booktabs}
\usepackage{amsmath}
\usepackage{natbib}

\title{My Document}
\author{Your Name}
\date{\today}

\begin{document}

\maketitle

\section{Introduction}
This is a sample document demonstrating various \LaTeX{} features including mathematical equations, tables, and formatting. You can cite papers using the citation search feature, and they will be automatically added to your \texttt{cite.bib} file.

\section{Mathematical Expressions}
Here are some example equations:

\subsection{Quadratic Formula}
The quadratic formula is given by:
\[x = \frac{-b \pm \sqrt{b^2 - 4ac}}{2a}\]

\subsection{Matrix Example}
A sample matrix:
\[
\begin{pmatrix}
a_{11} & a_{12} & a_{13} \\
a_{21} & a_{22} & a_{23} \\
a_{31} & a_{32} & a_{33}
\end{pmatrix}
\]

\section{Tables}
Below is a sample table showing data:

\begin{table}[h]
\centering
\begin{tabular}{@{}lrrr@{}}
\toprule
\textbf{Category} & \textbf{Value 1} & \textbf{Value 2} & \textbf{Total} \\
\midrule
Item A & 10.5 & 20.3 & 30.8 \\
Item B & 15.2 & 25.1 & 40.3 \\
Item C & 22.1 & 12.7 & 34.8 \\
\bottomrule
\end{tabular}
\caption{Sample Data Table}
\label{tab:sample}
\end{table}

\section{Lists}
\subsection{Bullet Points}
\begin{itemize}
  \item First item
  \item Second item
  \item Third item with sub-items:
    \begin{itemize}
      \item Sub-item 1
      \item Sub-item 2
    \end{itemize}
\end{itemize}

\subsection{Numbered List}
\begin{enumerate}
  \item First step
  \item Second step
  \item Third step
\end{enumerate}

\section{Citations and References}
When you use the citation search feature, references will be automatically managed. For example, you might cite a paper like this: \cite{example2024}. All citations will be automatically added to your \texttt{cite.bib} file.

\section{Conclusion}
This template demonstrates basic \LaTeX{} features including mathematical formulas, tables, and lists. The bibliography system is now set up to work seamlessly with Bibby AI's citation management. Feel free to modify it for your needs!

\bibliographystyle{plain}
\bibliography{cite}

\end{document}\documentclass{article}
\usepackage[utf8]{inputenc}
\usepackage{booktabs}
\usepackage{amsmath}
\usepackage{natbib}

\title{My Document}
\author{Your Name}
\date{\today}

\begin{document}

\maketitle

\section{Introduction}
This is a sample document demonstrating various \LaTeX{} features including mathematical equations, tables, and formatting. You can cite papers using the citation search feature, and they will be automatically added to your \texttt{cite.bib} file.

\section{Mathematical Expressions}
Here are some example equations:

\subsection{Quadratic Formula}
The quadratic formula is given by:
\[x = \frac{-b \pm \sqrt{b^2 - 4ac}}{2a}\]

\subsection{Matrix Example}
A sample matrix:
\[
\begin{pmatrix}
a_{11} & a_{12} & a_{13} \\
a_{21} & a_{22} & a_{23} \\
a_{31} & a_{32} & a_{33}
\end{pmatrix}
\]

\section{Tables}
Below is a sample table showing data:

\begin{table}[h]
\centering
\begin{tabular}{@{}lrrr@{}}
\toprule
\textbf{Category} & \textbf{Value 1} & \textbf{Value 2} & \textbf{Total} \\
\midrule
Item A & 10.5 & 20.3 & 30.8 \\
Item B & 15.2 & 25.1 & 40.3 \\
Item C & 22.1 & 12.7 & 34.8 \\
\bottomrule
\end{tabular}
\caption{Sample Data Table}
\label{tab:sample}
\end{table}

\section{Lists}
\subsection{Bullet Points}
\begin{itemize}
  \item First item
  \item Second item
  \item Third item with sub-items:
    \begin{itemize}
      \item Sub-item 1
      \item Sub-item 2
    \end{itemize}
\end{itemize}

\subsection{Numbered List}
\begin{enumerate}
  \item First step
  \item Second step
  \item Third step
\end{enumerate}

\section{Citations and References}
When you use the citation search feature, references will be automatically managed. For example, you might cite a paper like this: \cite{example2024}. All citations will be automatically added to your \texttt{cite.bib} file.

\section{Conclusion}
This template demonstrates basic \LaTeX{} features including mathematical formulas, tables, and lists. The bibliography system is now set up to work seamlessly with Bibby AI's citation management. Feel free to modify it for your needs!

\bibliographystyle{plain}
\bibliography{cite}

\end{document}\documentclass{article}
\usepackage[utf8]{inputenc}
\usepackage{booktabs}
\usepackage{amsmath}
\usepackage{natbib}

\title{My Document}
\author{Your Name}
\date{\today}

\begin{document}

\maketitle

\section{Introduction}
This is a sample document demonstrating various \LaTeX{} features including mathematical equations, tables, and formatting. You can cite papers using the citation search feature, and they will be automatically added to your \texttt{cite.bib} file.

\section{Mathematical Expressions}
Here are some example equations:

\subsection{Quadratic Formula}
The quadratic formula is given by:
\[x = \frac{-b \pm \sqrt{b^2 - 4ac}}{2a}\]

\subsection{Matrix Example}
A sample matrix:
\[
\begin{pmatrix}
a_{11} & a_{12} & a_{13} \\
a_{21} & a_{22} & a_{23} \\
a_{31} & a_{32} & a_{33}
\end{pmatrix}
\]

\section{Tables}
Below is a sample table showing data:

\begin{table}[h]
\centering
\begin{tabular}{@{}lrrr@{}}
\toprule
\textbf{Category} & \textbf{Value 1} & \textbf{Value 2} & \textbf{Total} \\
\midrule
Item A & 10.5 & 20.3 & 30.8 \\
Item B & 15.2 & 25.1 & 40.3 \\
Item C & 22.1 & 12.7 & 34.8 \\
\bottomrule
\end{tabular}
\caption{Sample Data Table}
\label{tab:sample}
\end{table}

\section{Lists}
\subsection{Bullet Points}
\begin{itemize}
  \item First item
  \item Second item
  \item Third item with sub-items:
    \begin{itemize}
      \item Sub-item 1
      \item Sub-item 2
    \end{itemize}
\end{itemize}

\subsection{Numbered List}
\begin{enumerate}
  \item First step
  \item Second step
  \item Third step
\end{enumerate}

\section{Citations and References}
When you use the citation search feature, references will be automatically managed. For example, you might cite a paper like this: \cite{example2024}. All citations will be automatically added to your \texttt{cite.bib} file.

\section{Conclusion}
This template demonstrates basic \LaTeX{} features including mathematical formulas, tables, and lists. The bibliography system is now set up to work seamlessly with Bibby AI's citation management. Feel free to modify it for your needs!

\bibliographystyle{plain}
\bibliography{cite}

\end{document}\documentclass{article}
\usepackage[utf8]{inputenc}
\usepackage{booktabs}
\usepackage{amsmath}
\usepackage{natbib}

\title{My Document}
\author{Your Name}
\date{\today}

\begin{document}

\maketitle

\section{Introduction}
This is a sample document demonstrating various \LaTeX{} features including mathematical equations, tables, and formatting. You can cite papers using the citation search feature, and they will be automatically added to your \texttt{cite.bib} file.

\section{Mathematical Expressions}
Here are some example equations:

\subsection{Quadratic Formula}
The quadratic formula is given by:
\[x = \frac{-b \pm \sqrt{b^2 - 4ac}}{2a}\]

\subsection{Matrix Example}
A sample matrix:
\[
\begin{pmatrix}
a_{11} & a_{12} & a_{13} \\
a_{21} & a_{22} & a_{23} \\
a_{31} & a_{32} & a_{33}
\end{pmatrix}
\]

\section{Tables}
Below is a sample table showing data:

\begin{table}[h]
\centering
\begin{tabular}{@{}lrrr@{}}
\toprule
\textbf{Category} & \textbf{Value 1} & \textbf{Value 2} & \textbf{Total} \\
\midrule
Item A & 10.5 & 20.3 & 30.8 \\
Item B & 15.2 & 25.1 & 40.3 \\
Item C & 22.1 & 12.7 & 34.8 \\
\bottomrule
\end{tabular}
\caption{Sample Data Table}
\label{tab:sample}
\end{table}

\section{Lists}
\subsection{Bullet Points}
\begin{itemize}
  \item First item
  \item Second item
  \item Third item with sub-items:
    \begin{itemize}
      \item Sub-item 1
      \item Sub-item 2
    \end{itemize}
\end{itemize}

\subsection{Numbered List}
\begin{enumerate}
  \item First step
  \item Second step
  \item Third step
\end{enumerate}

\section{Citations and References}
When you use the citation search feature, references will be automatically managed. For example, you might cite a paper like this: \cite{example2024}. All citations will be automatically added to your \texttt{cite.bib} file.

\section{Conclusion}
This template demonstrates basic \LaTeX{} features including mathematical formulas, tables, and lists. The bibliography system is now set up to work seamlessly with Bibby AI's citation management. Feel free to modify it for your needs!

\bibliographystyle{plain}
\bibliography{cite}

\end{document}\documentclass{article}
\usepackage[utf8]{inputenc}
\usepackage{booktabs}
\usepackage{amsmath}
\usepackage{natbib}

\title{My Document}
\author{Your Name}
\date{\today}

\begin{document}

\maketitle

\section{Introduction}
This is a sample document demonstrating various \LaTeX{} features including mathematical equations, tables, and formatting. You can cite papers using the citation search feature, and they will be automatically added to your \texttt{cite.bib} file.

\section{Mathematical Expressions}
Here are some example equations:

\subsection{Quadratic Formula}
The quadratic formula is given by:
\[x = \frac{-b \pm \sqrt{b^2 - 4ac}}{2a}\]

\subsection{Matrix Example}
A sample matrix:
\[
\begin{pmatrix}
a_{11} & a_{12} & a_{13} \\
a_{21} & a_{22} & a_{23} \\
a_{31} & a_{32} & a_{33}
\end{pmatrix}
\]

\section{Tables}
Below is a sample table showing data:

\begin{table}[h]
\centering
\begin{tabular}{@{}lrrr@{}}
\toprule
\textbf{Category} & \textbf{Value 1} & \textbf{Value 2} & \textbf{Total} \\
\midrule
Item A & 10.5 & 20.3 & 30.8 \\
Item B & 15.2 & 25.1 & 40.3 \\
Item C & 22.1 & 12.7 & 34.8 \\
\bottomrule
\end{tabular}
\caption{Sample Data Table}
\label{tab:sample}
\end{table}

\section{Lists}
\subsection{Bullet Points}
\begin{itemize}
  \item First item
  \item Second item
  \item Third item with sub-items:
    \begin{itemize}
      \item Sub-item 1
      \item Sub-item 2
    \end{itemize}
\end{itemize}

\subsection{Numbered List}
\begin{enumerate}
  \item First step
  \item Second step
  \item Third step
\end{enumerate}

\section{Citations and References}
When you use the citation search feature, references will be automatically managed. For example, you might cite a paper like this: \cite{example2024}. All citations will be automatically added to your \texttt{cite.bib} file.

\section{Conclusion}
This template demonstrates basic \LaTeX{} features including mathematical formulas, tables, and lists. The bibliography system is now set up to work seamlessly with Bibby AI's citation management. Feel free to modify it for your needs!

\bibliographystyle{plain}
\bibliography{cite}

\end{document}\documentclass{article}
\usepackage[utf8]{inputenc}
\usepackage{booktabs}
\usepackage{amsmath}
\usepackage{natbib}

\title{My Document}
\author{Your Name}
\date{\today}

\begin{document}

\maketitle

\section{Introduction}
This is a sample document demonstrating various \LaTeX{} features including mathematical equations, tables, and formatting. You can cite papers using the citation search feature, and they will be automatically added to your \texttt{cite.bib} file.

\section{Mathematical Expressions}
Here are some example equations:

\subsection{Quadratic Formula}
The quadratic formula is given by:
\[x = \frac{-b \pm \sqrt{b^2 - 4ac}}{2a}\]

\subsection{Matrix Example}
A sample matrix:
\[
\begin{pmatrix}
a_{11} & a_{12} & a_{13} \\
a_{21} & a_{22} & a_{23} \\
a_{31} & a_{32} & a_{33}
\end{pmatrix}
\]

\section{Tables}
Below is a sample table showing data:

\begin{table}[h]
\centering
\begin{tabular}{@{}lrrr@{}}
\toprule
\textbf{Category} & \textbf{Value 1} & \textbf{Value 2} & \textbf{Total} \\
\midrule
Item A & 10.5 & 20.3 & 30.8 \\
Item B & 15.2 & 25.1 & 40.3 \\
Item C & 22.1 & 12.7 & 34.8 \\
\bottomrule
\end{tabular}
\caption{Sample Data Table}
\label{tab:sample}
\end{table}

\section{Lists}
\subsection{Bullet Points}
\begin{itemize}
  \item First item
  \item Second item
  \item Third item with sub-items:
    \begin{itemize}
      \item Sub-item 1
      \item Sub-item 2
    \end{itemize}
\end{itemize}

\subsection{Numbered List}
\begin{enumerate}
  \item First step
  \item Second step
  \item Third step
\end{enumerate}

\section{Citations and References}
When you use the citation search feature, references will be automatically managed. For example, you might cite a paper like this: \cite{example2024}. All citations will be automatically added to your \texttt{cite.bib} file.

\section{Conclusion}
This template demonstrates basic \LaTeX{} features including mathematical formulas, tables, and lists. The bibliography system is now set up to work seamlessly with Bibby AI's citation management. Feel free to modify it for your needs!

\bibliographystyle{plain}
\bibliography{cite}

\end{document}\documentclass{article}
\usepackage[utf8]{inputenc}
\usepackage{booktabs}
\usepackage{amsmath}
\usepackage{natbib}

\title{My Document}
\author{Your Name}
\date{\today}

\begin{document}

\maketitle

\section{Introduction}
This is a sample document demonstrating various \LaTeX{} features including mathematical equations, tables, and formatting. You can cite papers using the citation search feature, and they will be automatically added to your \texttt{cite.bib} file.

\section{Mathematical Expressions}
Here are some example equations:

\subsection{Quadratic Formula}
The quadratic formula is given by:
\[x = \frac{-b \pm \sqrt{b^2 - 4ac}}{2a}\]

\subsection{Matrix Example}
A sample matrix:
\[
\begin{pmatrix}
a_{11} & a_{12} & a_{13} \\
a_{21} & a_{22} & a_{23} \\
a_{31} & a_{32} & a_{33}
\end{pmatrix}
\]

\section{Tables}
Below is a sample table showing data:

\begin{table}[h]
\centering
\begin{tabular}{@{}lrrr@{}}
\toprule
\textbf{Category} & \textbf{Value 1} & \textbf{Value 2} & \textbf{Total} \\
\midrule
Item A & 10.5 & 20.3 & 30.8 \\
Item B & 15.2 & 25.1 & 40.3 \\
Item C & 22.1 & 12.7 & 34.8 \\
\bottomrule
\end{tabular}
\caption{Sample Data Table}
\label{tab:sample}
\end{table}

\section{Lists}
\subsection{Bullet Points}
\begin{itemize}
  \item First item
  \item Second item
  \item Third item with sub-items:
    \begin{itemize}
      \item Sub-item 1
      \item Sub-item 2
    \end{itemize}
\end{itemize}

\subsection{Numbered List}
\begin{enumerate}
  \item First step
  \item Second step
  \item Third step
\end{enumerate}

\section{Citations and References}
When you use the citation search feature, references will be automatically managed. For example, you might cite a paper like this: \cite{example2024}. All citations will be automatically added to your \texttt{cite.bib} file.

\section{Conclusion}
This template demonstrates basic \LaTeX{} features including mathematical formulas, tables, and lists. The bibliography system is now set up to work seamlessly with Bibby AI's citation management. Feel free to modify it for your needs!

\bibliographystyle{plain}
\bibliography{cite}

\end{document}\documentclass{article}
\usepackage[utf8]{inputenc}
\usepackage{booktabs}
\usepackage{amsmath}
\usepackage{natbib}

\title{My Document}
\author{Your Name}
\date{\today}

\begin{document}

\maketitle

\section{Introduction}
This is a sample document demonstrating various \LaTeX{} features including mathematical equations, tables, and formatting. You can cite papers using the citation search feature, and they will be automatically added to your \texttt{cite.bib} file.

\section{Mathematical Expressions}
Here are some example equations:

\subsection{Quadratic Formula}
The quadratic formula is given by:
\[x = \frac{-b \pm \sqrt{b^2 - 4ac}}{2a}\]

\subsection{Matrix Example}
A sample matrix:
\[
\begin{pmatrix}
a_{11} & a_{12} & a_{13} \\
a_{21} & a_{22} & a_{23} \\
a_{31} & a_{32} & a_{33}
\end{pmatrix}
\]

\section{Tables}
Below is a sample table showing data:

\begin{table}[h]
\centering
\begin{tabular}{@{}lrrr@{}}
\toprule
\textbf{Category} & \textbf{Value 1} & \textbf{Value 2} & \textbf{Total} \\
\midrule
Item A & 10.5 & 20.3 & 30.8 \\
Item B & 15.2 & 25.1 & 40.3 \\
Item C & 22.1 & 12.7 & 34.8 \\
\bottomrule
\end{tabular}
\caption{Sample Data Table}
\label{tab:sample}
\end{table}

\section{Lists}
\subsection{Bullet Points}
\begin{itemize}
  \item First item
  \item Second item
  \item Third item with sub-items:
    \begin{itemize}
      \item Sub-item 1
      \item Sub-item 2
    \end{itemize}
\end{itemize}

\subsection{Numbered List}
\begin{enumerate}
  \item First step
  \item Second step
  \item Third step
\end{enumerate}

\section{Citations and References}
When you use the citation search feature, references will be automatically managed. For example, you might cite a paper like this: \cite{example2024}. All citations will be automatically added to your \texttt{cite.bib} file.

\section{Conclusion}
This template demonstrates basic \LaTeX{} features including mathematical formulas, tables, and lists. The bibliography system is now set up to work seamlessly with Bibby AI's citation management. Feel free to modify it for your needs!

\bibliographystyle{plain}
\bibliography{cite}

\end{document}\documentclass{article}
\usepackage[utf8]{inputenc}
\usepackage{booktabs}
\usepackage{amsmath}
\usepackage{natbib}

\title{My Document}
\author{Your Name}
\date{\today}

\begin{document}

\maketitle

\section{Introduction}
This is a sample document demonstrating various \LaTeX{} features including mathematical equations, tables, and formatting. You can cite papers using the citation search feature, and they will be automatically added to your \texttt{cite.bib} file.

\section{Mathematical Expressions}
Here are some example equations:

\subsection{Quadratic Formula}
The quadratic formula is given by:
\[x = \frac{-b \pm \sqrt{b^2 - 4ac}}{2a}\]

\subsection{Matrix Example}
A sample matrix:
\[
\begin{pmatrix}
a_{11} & a_{12} & a_{13} \\
a_{21} & a_{22} & a_{23} \\
a_{31} & a_{32} & a_{33}
\end{pmatrix}
\]

\section{Tables}
Below is a sample table showing data:

\begin{table}[h]
\centering
\begin{tabular}{@{}lrrr@{}}
\toprule
\textbf{Category} & \textbf{Value 1} & \textbf{Value 2} & \textbf{Total} \\
\midrule
Item A & 10.5 & 20.3 & 30.8 \\
Item B & 15.2 & 25.1 & 40.3 \\
Item C & 22.1 & 12.7 & 34.8 \\
\bottomrule
\end{tabular}
\caption{Sample Data Table}
\label{tab:sample}
\end{table}

\section{Lists}
\subsection{Bullet Points}
\begin{itemize}
  \item First item
  \item Second item
  \item Third item with sub-items:
    \begin{itemize}
      \item Sub-item 1
      \item Sub-item 2
    \end{itemize}
\end{itemize}

\subsection{Numbered List}
\begin{enumerate}
  \item First step
  \item Second step
  \item Third step
\end{enumerate}

\section{Citations and References}
When you use the citation search feature, references will be automatically managed. For example, you might cite a paper like this: \cite{example2024}. All citations will be automatically added to your \texttt{cite.bib} file.

\section{Conclusion}
This template demonstrates basic \LaTeX{} features including mathematical formulas, tables, and lists. The bibliography system is now set up to work seamlessly with Bibby AI's citation management. Feel free to modify it for your needs!

\bibliographystyle{plain}
\bibliography{cite}

\end{document}\documentclass{article}
\usepackage[utf8]{inputenc}
\usepackage{booktabs}
\usepackage{amsmath}
\usepackage{natbib}

\title{My Document}
\author{Your Name}
\date{\today}

\begin{document}

\maketitle

\section{Introduction}
This is a sample document demonstrating various \LaTeX{} features including mathematical equations, tables, and formatting. You can cite papers using the citation search feature, and they will be automatically added to your \texttt{cite.bib} file.

\section{Mathematical Expressions}
Here are some example equations:

\subsection{Quadratic Formula}
The quadratic formula is given by:
\[x = \frac{-b \pm \sqrt{b^2 - 4ac}}{2a}\]

\subsection{Matrix Example}
A sample matrix:
\[
\begin{pmatrix}
a_{11} & a_{12} & a_{13} \\
a_{21} & a_{22} & a_{23} \\
a_{31} & a_{32} & a_{33}
\end{pmatrix}
\]

\section{Tables}
Below is a sample table showing data:

\begin{table}[h]
\centering
\begin{tabular}{@{}lrrr@{}}
\toprule
\textbf{Category} & \textbf{Value 1} & \textbf{Value 2} & \textbf{Total} \\
\midrule
Item A & 10.5 & 20.3 & 30.8 \\
Item B & 15.2 & 25.1 & 40.3 \\
Item C & 22.1 & 12.7 & 34.8 \\
\bottomrule
\end{tabular}
\caption{Sample Data Table}
\label{tab:sample}
\end{table}

\section{Lists}
\subsection{Bullet Points}
\begin{itemize}
  \item First item
  \item Second item
  \item Third item with sub-items:
    \begin{itemize}
      \item Sub-item 1
      \item Sub-item 2
    \end{itemize}
\end{itemize}

\subsection{Numbered List}
\begin{enumerate}
  \item First step
  \item Second step
  \item Third step
\end{enumerate}

\section{Citations and References}
When you use the citation search feature, references will be automatically managed. For example, you might cite a paper like this: \cite{example2024}. All citations will be automatically added to your \texttt{cite.bib} file.

\section{Conclusion}
This template demonstrates basic \LaTeX{} features including mathematical formulas, tables, and lists. The bibliography system is now set up to work seamlessly with Bibby AI's citation management. Feel free to modify it for your needs!

\bibliographystyle{plain}
\bibliography{cite}

\end{document}\documentclass{article}
\usepackage[utf8]{inputenc}
\usepackage{booktabs}
\usepackage{amsmath}
\usepackage{natbib}

\title{My Document}
\author{Your Name}
\date{\today}

\begin{document}

\maketitle

\section{Introduction}
This is a sample document demonstrating various \LaTeX{} features including mathematical equations, tables, and formatting. You can cite papers using the citation search feature, and they will be automatically added to your \texttt{cite.bib} file.

\section{Mathematical Expressions}
Here are some example equations:

\subsection{Quadratic Formula}
The quadratic formula is given by:
\[x = \frac{-b \pm \sqrt{b^2 - 4ac}}{2a}\]

\subsection{Matrix Example}
A sample matrix:
\[
\begin{pmatrix}
a_{11} & a_{12} & a_{13} \\
a_{21} & a_{22} & a_{23} \\
a_{31} & a_{32} & a_{33}
\end{pmatrix}
\]

\section{Tables}
Below is a sample table showing data:

\begin{table}[h]
\centering
\begin{tabular}{@{}lrrr@{}}
\toprule
\textbf{Category} & \textbf{Value 1} & \textbf{Value 2} & \textbf{Total} \\
\midrule
Item A & 10.5 & 20.3 & 30.8 \\
Item B & 15.2 & 25.1 & 40.3 \\
Item C & 22.1 & 12.7 & 34.8 \\
\bottomrule
\end{tabular}
\caption{Sample Data Table}
\label{tab:sample}
\end{table}

\section{Lists}
\subsection{Bullet Points}
\begin{itemize}
  \item First item
  \item Second item
  \item Third item with sub-items:
    \begin{itemize}
      \item Sub-item 1
      \item Sub-item 2
    \end{itemize}
\end{itemize}

\subsection{Numbered List}
\begin{enumerate}
  \item First step
  \item Second step
  \item Third step
\end{enumerate}

\section{Citations and References}
When you use the citation search feature, references will be automatically managed. For example, you might cite a paper like this: \cite{example2024}. All citations will be automatically added to your \texttt{cite.bib} file.

\section{Conclusion}
This template demonstrates basic \LaTeX{} features including mathematical formulas, tables, and lists. The bibliography system is now set up to work seamlessly with Bibby AI's citation management. Feel free to modify it for your needs!

\bibliographystyle{plain}
\bibliography{cite}

\end{document}\documentclass{article}
\usepackage[utf8]{inputenc}
\usepackage{booktabs}
\usepackage{amsmath}
\usepackage{natbib}

\title{My Document}
\author{Your Name}
\date{\today}

\begin{document}

\maketitle

\section{Introduction}
This is a sample document demonstrating various \LaTeX{} features including mathematical equations, tables, and formatting. You can cite papers using the citation search feature, and they will be automatically added to your \texttt{cite.bib} file.

\section{Mathematical Expressions}
Here are some example equations:

\subsection{Quadratic Formula}
The quadratic formula is given by:
\[x = \frac{-b \pm \sqrt{b^2 - 4ac}}{2a}\]

\subsection{Matrix Example}
A sample matrix:
\[
\begin{pmatrix}
a_{11} & a_{12} & a_{13} \\
a_{21} & a_{22} & a_{23} \\
a_{31} & a_{32} & a_{33}
\end{pmatrix}
\]

\section{Tables}
Below is a sample table showing data:

\begin{table}[h]
\centering
\begin{tabular}{@{}lrrr@{}}
\toprule
\textbf{Category} & \textbf{Value 1} & \textbf{Value 2} & \textbf{Total} \\
\midrule
Item A & 10.5 & 20.3 & 30.8 \\
Item B & 15.2 & 25.1 & 40.3 \\
Item C & 22.1 & 12.7 & 34.8 \\
\bottomrule
\end{tabular}
\caption{Sample Data Table}
\label{tab:sample}
\end{table}

\section{Lists}
\subsection{Bullet Points}
\begin{itemize}
  \item First item
  \item Second item
  \item Third item with sub-items:
    \begin{itemize}
      \item Sub-item 1
      \item Sub-item 2
    \end{itemize}
\end{itemize}

\subsection{Numbered List}
\begin{enumerate}
  \item First step
  \item Second step
  \item Third step
\end{enumerate}

\section{Citations and References}
When you use the citation search feature, references will be automatically managed. For example, you might cite a paper like this: \cite{example2024}. All citations will be automatically added to your \texttt{cite.bib} file.

\section{Conclusion}
This template demonstrates basic \LaTeX{} features including mathematical formulas, tables, and lists. The bibliography system is now set up to work seamlessly with Bibby AI's citation management. Feel free to modify it for your needs!

\bibliographystyle{plain}
\bibliography{cite}

\end{document}\documentclass{article}
\usepackage[utf8]{inputenc}
\usepackage{booktabs}
\usepackage{amsmath}
\usepackage{natbib}

\title{My Document}
\author{Your Name}
\date{\today}

\begin{document}

\maketitle

\section{Introduction}
This is a sample document demonstrating various \LaTeX{} features including mathematical equations, tables, and formatting. You can cite papers using the citation search feature, and they will be automatically added to your \texttt{cite.bib} file.

\section{Mathematical Expressions}
Here are some example equations:

\subsection{Quadratic Formula}
The quadratic formula is given by:
\[x = \frac{-b \pm \sqrt{b^2 - 4ac}}{2a}\]

\subsection{Matrix Example}
A sample matrix:
\[
\begin{pmatrix}
a_{11} & a_{12} & a_{13} \\
a_{21} & a_{22} & a_{23} \\
a_{31} & a_{32} & a_{33}
\end{pmatrix}
\]

\section{Tables}
Below is a sample table showing data:

\begin{table}[h]
\centering
\begin{tabular}{@{}lrrr@{}}
\toprule
\textbf{Category} & \textbf{Value 1} & \textbf{Value 2} & \textbf{Total} \\
\midrule
Item A & 10.5 & 20.3 & 30.8 \\
Item B & 15.2 & 25.1 & 40.3 \\
Item C & 22.1 & 12.7 & 34.8 \\
\bottomrule
\end{tabular}
\caption{Sample Data Table}
\label{tab:sample}
\end{table}

\section{Lists}
\subsection{Bullet Points}
\begin{itemize}
  \item First item
  \item Second item
  \item Third item with sub-items:
    \begin{itemize}
      \item Sub-item 1
      \item Sub-item 2
    \end{itemize}
\end{itemize}

\subsection{Numbered List}
\begin{enumerate}
  \item First step
  \item Second step
  \item Third step
\end{enumerate}

\section{Citations and References}
When you use the citation search feature, references will be automatically managed. For example, you might cite a paper like this: \cite{example2024}. All citations will be automatically added to your \texttt{cite.bib} file.

\section{Conclusion}
This template demonstrates basic \LaTeX{} features including mathematical formulas, tables, and lists. The bibliography system is now set up to work seamlessly with Bibby AI's citation management. Feel free to modify it for your needs!

\bibliographystyle{plain}
\bibliography{cite}

\end{document}\documentclass{article}
\usepackage[utf8]{inputenc}
\usepackage{booktabs}
\usepackage{amsmath}
\usepackage{natbib}

\title{My Document}
\author{Your Name}
\date{\today}

\begin{document}

\maketitle

\section{Introduction}
This is a sample document demonstrating various \LaTeX{} features including mathematical equations, tables, and formatting. You can cite papers using the citation search feature, and they will be automatically added to your \texttt{cite.bib} file.

\section{Mathematical Expressions}
Here are some example equations:

\subsection{Quadratic Formula}
The quadratic formula is given by:
\[x = \frac{-b \pm \sqrt{b^2 - 4ac}}{2a}\]

\subsection{Matrix Example}
A sample matrix:
\[
\begin{pmatrix}
a_{11} & a_{12} & a_{13} \\
a_{21} & a_{22} & a_{23} \\
a_{31} & a_{32} & a_{33}
\end{pmatrix}
\]

\section{Tables}
Below is a sample table showing data:

\begin{table}[h]
\centering
\begin{tabular}{@{}lrrr@{}}
\toprule
\textbf{Category} & \textbf{Value 1} & \textbf{Value 2} & \textbf{Total} \\
\midrule
Item A & 10.5 & 20.3 & 30.8 \\
Item B & 15.2 & 25.1 & 40.3 \\
Item C & 22.1 & 12.7 & 34.8 \\
\bottomrule
\end{tabular}
\caption{Sample Data Table}
\label{tab:sample}
\end{table}

\section{Lists}
\subsection{Bullet Points}
\begin{itemize}
  \item First item
  \item Second item
  \item Third item with sub-items:
    \begin{itemize}
      \item Sub-item 1
      \item Sub-item 2
    \end{itemize}
\end{itemize}

\subsection{Numbered List}
\begin{enumerate}
  \item First step
  \item Second step
  \item Third step
\end{enumerate}

\section{Citations and References}
When you use the citation search feature, references will be automatically managed. For example, you might cite a paper like this: \cite{example2024}. All citations will be automatically added to your \texttt{cite.bib} file.

\section{Conclusion}
This template demonstrates basic \LaTeX{} features including mathematical formulas, tables, and lists. The bibliography system is now set up to work seamlessly with Bibby AI's citation management. Feel free to modify it for your needs!

\bibliographystyle{plain}
\bibliography{cite}

\end{document}\documentclass{article}
\usepackage[utf8]{inputenc}
\usepackage{booktabs}
\usepackage{amsmath}
\usepackage{natbib}

\title{My Document}
\author{Your Name}
\date{\today}

\begin{document}

\maketitle

\section{Introduction}
This is a sample document demonstrating various \LaTeX{} features including mathematical equations, tables, and formatting. You can cite papers using the citation search feature, and they will be automatically added to your \texttt{cite.bib} file.

\section{Mathematical Expressions}
Here are some example equations:

\subsection{Quadratic Formula}
The quadratic formula is given by:
\[x = \frac{-b \pm \sqrt{b^2 - 4ac}}{2a}\]

\subsection{Matrix Example}
A sample matrix:
\[
\begin{pmatrix}
a_{11} & a_{12} & a_{13} \\
a_{21} & a_{22} & a_{23} \\
a_{31} & a_{32} & a_{33}
\end{pmatrix}
\]

\section{Tables}
Below is a sample table showing data:

\begin{table}[h]
\centering
\begin{tabular}{@{}lrrr@{}}
\toprule
\textbf{Category} & \textbf{Value 1} & \textbf{Value 2} & \textbf{Total} \\
\midrule
Item A & 10.5 & 20.3 & 30.8 \\
Item B & 15.2 & 25.1 & 40.3 \\
Item C & 22.1 & 12.7 & 34.8 \\
\bottomrule
\end{tabular}
\caption{Sample Data Table}
\label{tab:sample}
\end{table}

\section{Lists}
\subsection{Bullet Points}
\begin{itemize}
  \item First item
  \item Second item
  \item Third item with sub-items:
    \begin{itemize}
      \item Sub-item 1
      \item Sub-item 2
    \end{itemize}
\end{itemize}

\subsection{Numbered List}
\begin{enumerate}
  \item First step
  \item Second step
  \item Third step
\end{enumerate}

\section{Citations and References}
When you use the citation search feature, references will be automatically managed. For example, you might cite a paper like this: \cite{example2024}. All citations will be automatically added to your \texttt{cite.bib} file.

\section{Conclusion}
This template demonstrates basic \LaTeX{} features including mathematical formulas, tables, and lists. The bibliography system is now set up to work seamlessly with Bibby AI's citation management. Feel free to modify it for your needs!

\bibliographystyle{plain}
\bibliography{cite}

\end{document}\documentclass{article}
\usepackage[utf8]{inputenc}
\usepackage{booktabs}
\usepackage{amsmath}
\usepackage{natbib}

\title{My Document}
\author{Your Name}
\date{\today}

\begin{document}

\maketitle

\section{Introduction}
This is a sample document demonstrating various \LaTeX{} features including mathematical equations, tables, and formatting. You can cite papers using the citation search feature, and they will be automatically added to your \texttt{cite.bib} file.

\section{Mathematical Expressions}
Here are some example equations:

\subsection{Quadratic Formula}
The quadratic formula is given by:
\[x = \frac{-b \pm \sqrt{b^2 - 4ac}}{2a}\]

\subsection{Matrix Example}
A sample matrix:
\[
\begin{pmatrix}
a_{11} & a_{12} & a_{13} \\
a_{21} & a_{22} & a_{23} \\
a_{31} & a_{32} & a_{33}
\end{pmatrix}
\]

\section{Tables}
Below is a sample table showing data:

\begin{table}[h]
\centering
\begin{tabular}{@{}lrrr@{}}
\toprule
\textbf{Category} & \textbf{Value 1} & \textbf{Value 2} & \textbf{Total} \\
\midrule
Item A & 10.5 & 20.3 & 30.8 \\
Item B & 15.2 & 25.1 & 40.3 \\
Item C & 22.1 & 12.7 & 34.8 \\
\bottomrule
\end{tabular}
\caption{Sample Data Table}
\label{tab:sample}
\end{table}

\section{Lists}
\subsection{Bullet Points}
\begin{itemize}
  \item First item
  \item Second item
  \item Third item with sub-items:
    \begin{itemize}
      \item Sub-item 1
      \item Sub-item 2
    \end{itemize}
\end{itemize}

\subsection{Numbered List}
\begin{enumerate}
  \item First step
  \item Second step
  \item Third step
\end{enumerate}

\section{Citations and References}
When you use the citation search feature, references will be automatically managed. For example, you might cite a paper like this: \cite{example2024}. All citations will be automatically added to your \texttt{cite.bib} file.

\section{Conclusion}
This template demonstrates basic \LaTeX{} features including mathematical formulas, tables, and lists. The bibliography system is now set up to work seamlessly with Bibby AI's citation management. Feel free to modify it for your needs!

\bibliographystyle{plain}
\bibliography{cite}

\end{document}\documentclass{article}
\usepackage[utf8]{inputenc}
\usepackage{booktabs}
\usepackage{amsmath}
\usepackage{natbib}

\title{My Document}
\author{Your Name}
\date{\today}

\begin{document}

\maketitle

\section{Introduction}
This is a sample document demonstrating various \LaTeX{} features including mathematical equations, tables, and formatting. You can cite papers using the citation search feature, and they will be automatically added to your \texttt{cite.bib} file.

\section{Mathematical Expressions}
Here are some example equations:

\subsection{Quadratic Formula}
The quadratic formula is given by:
\[x = \frac{-b \pm \sqrt{b^2 - 4ac}}{2a}\]

\subsection{Matrix Example}
A sample matrix:
\[
\begin{pmatrix}
a_{11} & a_{12} & a_{13} \\
a_{21} & a_{22} & a_{23} \\
a_{31} & a_{32} & a_{33}
\end{pmatrix}
\]

\section{Tables}
Below is a sample table showing data:

\begin{table}[h]
\centering
\begin{tabular}{@{}lrrr@{}}
\toprule
\textbf{Category} & \textbf{Value 1} & \textbf{Value 2} & \textbf{Total} \\
\midrule
Item A & 10.5 & 20.3 & 30.8 \\
Item B & 15.2 & 25.1 & 40.3 \\
Item C & 22.1 & 12.7 & 34.8 \\
\bottomrule
\end{tabular}
\caption{Sample Data Table}
\label{tab:sample}
\end{table}

\section{Lists}
\subsection{Bullet Points}
\begin{itemize}
  \item First item
  \item Second item
  \item Third item with sub-items:
    \begin{itemize}
      \item Sub-item 1
      \item Sub-item 2
    \end{itemize}
\end{itemize}

\subsection{Numbered List}
\begin{enumerate}
  \item First step
  \item Second step
  \item Third step
\end{enumerate}

\section{Citations and References}
When you use the citation search feature, references will be automatically managed. For example, you might cite a paper like this: \cite{example2024}. All citations will be automatically added to your \texttt{cite.bib} file.

\section{Conclusion}
This template demonstrates basic \LaTeX{} features including mathematical formulas, tables, and lists. The bibliography system is now set up to work seamlessly with Bibby AI's citation management. Feel free to modify it for your needs!

\bibliographystyle{plain}
\bibliography{cite}

\end{document}\documentclass{article}
\usepackage[utf8]{inputenc}
\usepackage{booktabs}
\usepackage{amsmath}
\usepackage{natbib}

\title{My Document}
\author{Your Name}
\date{\today}

\begin{document}

\maketitle

\section{Introduction}
This is a sample document demonstrating various \LaTeX{} features including mathematical equations, tables, and formatting. You can cite papers using the citation search feature, and they will be automatically added to your \texttt{cite.bib} file.

\section{Mathematical Expressions}
Here are some example equations:

\subsection{Quadratic Formula}
The quadratic formula is given by:
\[x = \frac{-b \pm \sqrt{b^2 - 4ac}}{2a}\]

\subsection{Matrix Example}
A sample matrix:
\[
\begin{pmatrix}
a_{11} & a_{12} & a_{13} \\
a_{21} & a_{22} & a_{23} \\
a_{31} & a_{32} & a_{33}
\end{pmatrix}
\]

\section{Tables}
Below is a sample table showing data:

\begin{table}[h]
\centering
\begin{tabular}{@{}lrrr@{}}
\toprule
\textbf{Category} & \textbf{Value 1} & \textbf{Value 2} & \textbf{Total} \\
\midrule
Item A & 10.5 & 20.3 & 30.8 \\
Item B & 15.2 & 25.1 & 40.3 \\
Item C & 22.1 & 12.7 & 34.8 \\
\bottomrule
\end{tabular}
\caption{Sample Data Table}
\label{tab:sample}
\end{table}

\section{Lists}
\subsection{Bullet Points}
\begin{itemize}
  \item First item
  \item Second item
  \item Third item with sub-items:
    \begin{itemize}
      \item Sub-item 1
      \item Sub-item 2
    \end{itemize}
\end{itemize}

\subsection{Numbered List}
\begin{enumerate}
  \item First step
  \item Second step
  \item Third step
\end{enumerate}

\section{Citations and References}
When you use the citation search feature, references will be automatically managed. For example, you might cite a paper like this: \cite{example2024}. All citations will be automatically added to your \texttt{cite.bib} file.

\section{Conclusion}
This template demonstrates basic \LaTeX{} features including mathematical formulas, tables, and lists. The bibliography system is now set up to work seamlessly with Bibby AI's citation management. Feel free to modify it for your needs!

\bibliographystyle{plain}
\bibliography{cite}

\end{document}\documentclass{article}
\usepackage[utf8]{inputenc}
\usepackage{booktabs}
\usepackage{amsmath}
\usepackage{natbib}

\title{My Document}
\author{Your Name}
\date{\today}

\begin{document}

\maketitle

\section{Introduction}
This is a sample document demonstrating various \LaTeX{} features including mathematical equations, tables, and formatting. You can cite papers using the citation search feature, and they will be automatically added to your \texttt{cite.bib} file.

\section{Mathematical Expressions}
Here are some example equations:

\subsection{Quadratic Formula}
The quadratic formula is given by:
\[x = \frac{-b \pm \sqrt{b^2 - 4ac}}{2a}\]

\subsection{Matrix Example}
A sample matrix:
\[
\begin{pmatrix}
a_{11} & a_{12} & a_{13} \\
a_{21} & a_{22} & a_{23} \\
a_{31} & a_{32} & a_{33}
\end{pmatrix}
\]

\section{Tables}
Below is a sample table showing data:

\begin{table}[h]
\centering
\begin{tabular}{@{}lrrr@{}}
\toprule
\textbf{Category} & \textbf{Value 1} & \textbf{Value 2} & \textbf{Total} \\
\midrule
Item A & 10.5 & 20.3 & 30.8 \\
Item B & 15.2 & 25.1 & 40.3 \\
Item C & 22.1 & 12.7 & 34.8 \\
\bottomrule
\end{tabular}
\caption{Sample Data Table}
\label{tab:sample}
\end{table}

\section{Lists}
\subsection{Bullet Points}
\begin{itemize}
  \item First item
  \item Second item
  \item Third item with sub-items:
    \begin{itemize}
      \item Sub-item 1
      \item Sub-item 2
    \end{itemize}
\end{itemize}

\subsection{Numbered List}
\begin{enumerate}
  \item First step
  \item Second step
  \item Third step
\end{enumerate}

\section{Citations and References}
When you use the citation search feature, references will be automatically managed. For example, you might cite a paper like this: \cite{example2024}. All citations will be automatically added to your \texttt{cite.bib} file.

\section{Conclusion}
This template demonstrates basic \LaTeX{} features including mathematical formulas, tables, and lists. The bibliography system is now set up to work seamlessly with Bibby AI's citation management. Feel free to modify it for your needs!

\bibliographystyle{plain}
\bibliography{cite}

\end{document}\documentclass{article}
\usepackage[utf8]{inputenc}
\usepackage{booktabs}
\usepackage{amsmath}
\usepackage{natbib}

\title{My Document}
\author{Your Name}
\date{\today}

\begin{document}

\maketitle

\section{Introduction}
This is a sample document demonstrating various \LaTeX{} features including mathematical equations, tables, and formatting. You can cite papers using the citation search feature, and they will be automatically added to your \texttt{cite.bib} file.

\section{Mathematical Expressions}
Here are some example equations:

\subsection{Quadratic Formula}
The quadratic formula is given by:
\[x = \frac{-b \pm \sqrt{b^2 - 4ac}}{2a}\]

\subsection{Matrix Example}
A sample matrix:
\[
\begin{pmatrix}
a_{11} & a_{12} & a_{13} \\
a_{21} & a_{22} & a_{23} \\
a_{31} & a_{32} & a_{33}
\end{pmatrix}
\]

\section{Tables}
Below is a sample table showing data:

\begin{table}[h]
\centering
\begin{tabular}{@{}lrrr@{}}
\toprule
\textbf{Category} & \textbf{Value 1} & \textbf{Value 2} & \textbf{Total} \\
\midrule
Item A & 10.5 & 20.3 & 30.8 \\
Item B & 15.2 & 25.1 & 40.3 \\
Item C & 22.1 & 12.7 & 34.8 \\
\bottomrule
\end{tabular}
\caption{Sample Data Table}
\label{tab:sample}
\end{table}

\section{Lists}
\subsection{Bullet Points}
\begin{itemize}
  \item First item
  \item Second item
  \item Third item with sub-items:
    \begin{itemize}
      \item Sub-item 1
      \item Sub-item 2
    \end{itemize}
\end{itemize}

\subsection{Numbered List}
\begin{enumerate}
  \item First step
  \item Second step
  \item Third step
\end{enumerate}

\section{Citations and References}
When you use the citation search feature, references will be automatically managed. For example, you might cite a paper like this: \cite{example2024}. All citations will be automatically added to your \texttt{cite.bib} file.

\section{Conclusion}
This template demonstrates basic \LaTeX{} features including mathematical formulas, tables, and lists. The bibliography system is now set up to work seamlessly with Bibby AI's citation management. Feel free to modify it for your needs!

\bibliographystyle{plain}
\bibliography{cite}

\end{document}\documentclass{article}
\usepackage[utf8]{inputenc}
\usepackage{booktabs}
\usepackage{amsmath}
\usepackage{natbib}

\title{My Document}
\author{Your Name}
\date{\today}

\begin{document}

\maketitle

\section{Introduction}
This is a sample document demonstrating various \LaTeX{} features including mathematical equations, tables, and formatting. You can cite papers using the citation search feature, and they will be automatically added to your \texttt{cite.bib} file.

\section{Mathematical Expressions}
Here are some example equations:

\subsection{Quadratic Formula}
The quadratic formula is given by:
\[x = \frac{-b \pm \sqrt{b^2 - 4ac}}{2a}\]

\subsection{Matrix Example}
A sample matrix:
\[
\begin{pmatrix}
a_{11} & a_{12} & a_{13} \\
a_{21} & a_{22} & a_{23} \\
a_{31} & a_{32} & a_{33}
\end{pmatrix}
\]

\section{Tables}
Below is a sample table showing data:

\begin{table}[h]
\centering
\begin{tabular}{@{}lrrr@{}}
\toprule
\textbf{Category} & \textbf{Value 1} & \textbf{Value 2} & \textbf{Total} \\
\midrule
Item A & 10.5 & 20.3 & 30.8 \\
Item B & 15.2 & 25.1 & 40.3 \\
Item C & 22.1 & 12.7 & 34.8 \\
\bottomrule
\end{tabular}
\caption{Sample Data Table}
\label{tab:sample}
\end{table}

\section{Lists}
\subsection{Bullet Points}
\begin{itemize}
  \item First item
  \item Second item
  \item Third item with sub-items:
    \begin{itemize}
      \item Sub-item 1
      \item Sub-item 2
    \end{itemize}
\end{itemize}

\subsection{Numbered List}
\begin{enumerate}
  \item First step
  \item Second step
  \item Third step
\end{enumerate}

\section{Citations and References}
When you use the citation search feature, references will be automatically managed. For example, you might cite a paper like this: \cite{example2024}. All citations will be automatically added to your \texttt{cite.bib} file.

\section{Conclusion}
This template demonstrates basic \LaTeX{} features including mathematical formulas, tables, and lists. The bibliography system is now set up to work seamlessly with Bibby AI's citation management. Feel free to modify it for your needs!

\bibliographystyle{plain}
\bibliography{cite}

\end{document}\documentclass{article}
\usepackage[utf8]{inputenc}
\usepackage{booktabs}
\usepackage{amsmath}
\usepackage{natbib}

\title{My Document}
\author{Your Name}
\date{\today}

\begin{document}

\maketitle

\section{Introduction}
This is a sample document demonstrating various \LaTeX{} features including mathematical equations, tables, and formatting. You can cite papers using the citation search feature, and they will be automatically added to your \texttt{cite.bib} file.

\section{Mathematical Expressions}
Here are some example equations:

\subsection{Quadratic Formula}
The quadratic formula is given by:
\[x = \frac{-b \pm \sqrt{b^2 - 4ac}}{2a}\]

\subsection{Matrix Example}
A sample matrix:
\[
\begin{pmatrix}
a_{11} & a_{12} & a_{13} \\
a_{21} & a_{22} & a_{23} \\
a_{31} & a_{32} & a_{33}
\end{pmatrix}
\]

\section{Tables}
Below is a sample table showing data:

\begin{table}[h]
\centering
\begin{tabular}{@{}lrrr@{}}
\toprule
\textbf{Category} & \textbf{Value 1} & \textbf{Value 2} & \textbf{Total} \\
\midrule
Item A & 10.5 & 20.3 & 30.8 \\
Item B & 15.2 & 25.1 & 40.3 \\
Item C & 22.1 & 12.7 & 34.8 \\
\bottomrule
\end{tabular}
\caption{Sample Data Table}
\label{tab:sample}
\end{table}

\section{Lists}
\subsection{Bullet Points}
\begin{itemize}
  \item First item
  \item Second item
  \item Third item with sub-items:
    \begin{itemize}
      \item Sub-item 1
      \item Sub-item 2
    \end{itemize}
\end{itemize}

\subsection{Numbered List}
\begin{enumerate}
  \item First step
  \item Second step
  \item Third step
\end{enumerate}

\section{Citations and References}
When you use the citation search feature, references will be automatically managed. For example, you might cite a paper like this: \cite{example2024}. All citations will be automatically added to your \texttt{cite.bib} file.

\section{Conclusion}
This template demonstrates basic \LaTeX{} features including mathematical formulas, tables, and lists. The bibliography system is now set up to work seamlessly with Bibby AI's citation management. Feel free to modify it for your needs!

\bibliographystyle{plain}
\bibliography{cite}

\end{document}\documentclass{article}
\usepackage[utf8]{inputenc}
\usepackage{booktabs}
\usepackage{amsmath}
\usepackage{natbib}

\title{My Document}
\author{Your Name}
\date{\today}

\begin{document}

\maketitle

\section{Introduction}
This is a sample document demonstrating various \LaTeX{} features including mathematical equations, tables, and formatting. You can cite papers using the citation search feature, and they will be automatically added to your \texttt{cite.bib} file.

\section{Mathematical Expressions}
Here are some example equations:

\subsection{Quadratic Formula}
The quadratic formula is given by:
\[x = \frac{-b \pm \sqrt{b^2 - 4ac}}{2a}\]

\subsection{Matrix Example}
A sample matrix:
\[
\begin{pmatrix}
a_{11} & a_{12} & a_{13} \\
a_{21} & a_{22} & a_{23} \\
a_{31} & a_{32} & a_{33}
\end{pmatrix}
\]

\section{Tables}
Below is a sample table showing data:

\begin{table}[h]
\centering
\begin{tabular}{@{}lrrr@{}}
\toprule
\textbf{Category} & \textbf{Value 1} & \textbf{Value 2} & \textbf{Total} \\
\midrule
Item A & 10.5 & 20.3 & 30.8 \\
Item B & 15.2 & 25.1 & 40.3 \\
Item C & 22.1 & 12.7 & 34.8 \\
\bottomrule
\end{tabular}
\caption{Sample Data Table}
\label{tab:sample}
\end{table}

\section{Lists}
\subsection{Bullet Points}
\begin{itemize}
  \item First item
  \item Second item
  \item Third item with sub-items:
    \begin{itemize}
      \item Sub-item 1
      \item Sub-item 2
    \end{itemize}
\end{itemize}

\subsection{Numbered List}
\begin{enumerate}
  \item First step
  \item Second step
  \item Third step
\end{enumerate}

\section{Citations and References}
When you use the citation search feature, references will be automatically managed. For example, you might cite a paper like this: \cite{example2024}. All citations will be automatically added to your \texttt{cite.bib} file.

\section{Conclusion}
This template demonstrates basic \LaTeX{} features including mathematical formulas, tables, and lists. The bibliography system is now set up to work seamlessly with Bibby AI's citation management. Feel free to modify it for your needs!

\bibliographystyle{plain}
\bibliography{cite}

\end{document}\documentclass{article}
\usepackage[utf8]{inputenc}
\usepackage{booktabs}
\usepackage{amsmath}
\usepackage{natbib}

\title{My Document}
\author{Your Name}
\date{\today}

\begin{document}

\maketitle

\section{Introduction}
This is a sample document demonstrating various \LaTeX{} features including mathematical equations, tables, and formatting. You can cite papers using the citation search feature, and they will be automatically added to your \texttt{cite.bib} file.

\section{Mathematical Expressions}
Here are some example equations:

\subsection{Quadratic Formula}
The quadratic formula is given by:
\[x = \frac{-b \pm \sqrt{b^2 - 4ac}}{2a}\]

\subsection{Matrix Example}
A sample matrix:
\[
\begin{pmatrix}
a_{11} & a_{12} & a_{13} \\
a_{21} & a_{22} & a_{23} \\
a_{31} & a_{32} & a_{33}
\end{pmatrix}
\]

\section{Tables}
Below is a sample table showing data:

\begin{table}[h]
\centering
\begin{tabular}{@{}lrrr@{}}
\toprule
\textbf{Category} & \textbf{Value 1} & \textbf{Value 2} & \textbf{Total} \\
\midrule
Item A & 10.5 & 20.3 & 30.8 \\
Item B & 15.2 & 25.1 & 40.3 \\
Item C & 22.1 & 12.7 & 34.8 \\
\bottomrule
\end{tabular}
\caption{Sample Data Table}
\label{tab:sample}
\end{table}

\section{Lists}
\subsection{Bullet Points}
\begin{itemize}
  \item First item
  \item Second item
  \item Third item with sub-items:
    \begin{itemize}
      \item Sub-item 1
      \item Sub-item 2
    \end{itemize}
\end{itemize}

\subsection{Numbered List}
\begin{enumerate}
  \item First step
  \item Second step
  \item Third step
\end{enumerate}

\section{Citations and References}
When you use the citation search feature, references will be automatically managed. For example, you might cite a paper like this: \cite{example2024}. All citations will be automatically added to your \texttt{cite.bib} file.

\section{Conclusion}
This template demonstrates basic \LaTeX{} features including mathematical formulas, tables, and lists. The bibliography system is now set up to work seamlessly with Bibby AI's citation management. Feel free to modify it for your needs!

\bibliographystyle{plain}
\bibliography{cite}

\end{document}\documentclass{article}
\usepackage[utf8]{inputenc}
\usepackage{booktabs}
\usepackage{amsmath}
\usepackage{natbib}

\title{My Document}
\author{Your Name}
\date{\today}

\begin{document}

\maketitle

\section{Introduction}
This is a sample document demonstrating various \LaTeX{} features including mathematical equations, tables, and formatting. You can cite papers using the citation search feature, and they will be automatically added to your \texttt{cite.bib} file.

\section{Mathematical Expressions}
Here are some example equations:

\subsection{Quadratic Formula}
The quadratic formula is given by:
\[x = \frac{-b \pm \sqrt{b^2 - 4ac}}{2a}\]

\subsection{Matrix Example}
A sample matrix:
\[
\begin{pmatrix}
a_{11} & a_{12} & a_{13} \\
a_{21} & a_{22} & a_{23} \\
a_{31} & a_{32} & a_{33}
\end{pmatrix}
\]

\section{Tables}
Below is a sample table showing data:

\begin{table}[h]
\centering
\begin{tabular}{@{}lrrr@{}}
\toprule
\textbf{Category} & \textbf{Value 1} & \textbf{Value 2} & \textbf{Total} \\
\midrule
Item A & 10.5 & 20.3 & 30.8 \\
Item B & 15.2 & 25.1 & 40.3 \\
Item C & 22.1 & 12.7 & 34.8 \\
\bottomrule
\end{tabular}
\caption{Sample Data Table}
\label{tab:sample}
\end{table}

\section{Lists}
\subsection{Bullet Points}
\begin{itemize}
  \item First item
  \item Second item
  \item Third item with sub-items:
    \begin{itemize}
      \item Sub-item 1
      \item Sub-item 2
    \end{itemize}
\end{itemize}

\subsection{Numbered List}
\begin{enumerate}
  \item First step
  \item Second step
  \item Third step
\end{enumerate}

\section{Citations and References}
When you use the citation search feature, references will be automatically managed. For example, you might cite a paper like this: \cite{example2024}. All citations will be automatically added to your \texttt{cite.bib} file.

\section{Conclusion}
This template demonstrates basic \LaTeX{} features including mathematical formulas, tables, and lists. The bibliography system is now set up to work seamlessly with Bibby AI's citation management. Feel free to modify it for your needs!

\bibliographystyle{plain}
\bibliography{cite}

\end{document}\documentclass{article}
\usepackage[utf8]{inputenc}
\usepackage{booktabs}
\usepackage{amsmath}
\usepackage{natbib}

\title{My Document}
\author{Your Name}
\date{\today}

\begin{document}

\maketitle

\section{Introduction}
This is a sample document demonstrating various \LaTeX{} features including mathematical equations, tables, and formatting. You can cite papers using the citation search feature, and they will be automatically added to your \texttt{cite.bib} file.

\section{Mathematical Expressions}
Here are some example equations:

\subsection{Quadratic Formula}
The quadratic formula is given by:
\[x = \frac{-b \pm \sqrt{b^2 - 4ac}}{2a}\]

\subsection{Matrix Example}
A sample matrix:
\[
\begin{pmatrix}
a_{11} & a_{12} & a_{13} \\
a_{21} & a_{22} & a_{23} \\
a_{31} & a_{32} & a_{33}
\end{pmatrix}
\]

\section{Tables}
Below is a sample table showing data:

\begin{table}[h]
\centering
\begin{tabular}{@{}lrrr@{}}
\toprule
\textbf{Category} & \textbf{Value 1} & \textbf{Value 2} & \textbf{Total} \\
\midrule
Item A & 10.5 & 20.3 & 30.8 \\
Item B & 15.2 & 25.1 & 40.3 \\
Item C & 22.1 & 12.7 & 34.8 \\
\bottomrule
\end{tabular}
\caption{Sample Data Table}
\label{tab:sample}
\end{table}

\section{Lists}
\subsection{Bullet Points}
\begin{itemize}
  \item First item
  \item Second item
  \item Third item with sub-items:
    \begin{itemize}
      \item Sub-item 1
      \item Sub-item 2
    \end{itemize}
\end{itemize}

\subsection{Numbered List}
\begin{enumerate}
  \item First step
  \item Second step
  \item Third step
\end{enumerate}

\section{Citations and References}
When you use the citation search feature, references will be automatically managed. For example, you might cite a paper like this: \cite{example2024}. All citations will be automatically added to your \texttt{cite.bib} file.

\section{Conclusion}
This template demonstrates basic \LaTeX{} features including mathematical formulas, tables, and lists. The bibliography system is now set up to work seamlessly with Bibby AI's citation management. Feel free to modify it for your needs!

\bibliographystyle{plain}
\bibliography{cite}

\end{document}\documentclass{article}
\usepackage[utf8]{inputenc}
\usepackage{booktabs}
\usepackage{amsmath}
\usepackage{natbib}

\title{My Document}
\author{Your Name}
\date{\today}

\begin{document}

\maketitle

\section{Introduction}
This is a sample document demonstrating various \LaTeX{} features including mathematical equations, tables, and formatting. You can cite papers using the citation search feature, and they will be automatically added to your \texttt{cite.bib} file.

\section{Mathematical Expressions}
Here are some example equations:

\subsection{Quadratic Formula}
The quadratic formula is given by:
\[x = \frac{-b \pm \sqrt{b^2 - 4ac}}{2a}\]

\subsection{Matrix Example}
A sample matrix:
\[
\begin{pmatrix}
a_{11} & a_{12} & a_{13} \\
a_{21} & a_{22} & a_{23} \\
a_{31} & a_{32} & a_{33}
\end{pmatrix}
\]

\section{Tables}
Below is a sample table showing data:

\begin{table}[h]
\centering
\begin{tabular}{@{}lrrr@{}}
\toprule
\textbf{Category} & \textbf{Value 1} & \textbf{Value 2} & \textbf{Total} \\
\midrule
Item A & 10.5 & 20.3 & 30.8 \\
Item B & 15.2 & 25.1 & 40.3 \\
Item C & 22.1 & 12.7 & 34.8 \\
\bottomrule
\end{tabular}
\caption{Sample Data Table}
\label{tab:sample}
\end{table}

\section{Lists}
\subsection{Bullet Points}
\begin{itemize}
  \item First item
  \item Second item
  \item Third item with sub-items:
    \begin{itemize}
      \item Sub-item 1
      \item Sub-item 2
    \end{itemize}
\end{itemize}

\subsection{Numbered List}
\begin{enumerate}
  \item First step
  \item Second step
  \item Third step
\end{enumerate}

\section{Citations and References}
When you use the citation search feature, references will be automatically managed. For example, you might cite a paper like this: \cite{example2024}. All citations will be automatically added to your \texttt{cite.bib} file.

\section{Conclusion}
This template demonstrates basic \LaTeX{} features including mathematical formulas, tables, and lists. The bibliography system is now set up to work seamlessly with Bibby AI's citation management. Feel free to modify it for your needs!

\bibliographystyle{plain}
\bibliography{cite}

\end{document}\documentclass{article}
\usepackage[utf8]{inputenc}
\usepackage{booktabs}
\usepackage{amsmath}
\usepackage{natbib}

\title{My Document}
\author{Your Name}
\date{\today}

\begin{document}

\maketitle

\section{Introduction}
This is a sample document demonstrating various \LaTeX{} features including mathematical equations, tables, and formatting. You can cite papers using the citation search feature, and they will be automatically added to your \texttt{cite.bib} file.

\section{Mathematical Expressions}
Here are some example equations:

\subsection{Quadratic Formula}
The quadratic formula is given by:
\[x = \frac{-b \pm \sqrt{b^2 - 4ac}}{2a}\]

\subsection{Matrix Example}
A sample matrix:
\[
\begin{pmatrix}
a_{11} & a_{12} & a_{13} \\
a_{21} & a_{22} & a_{23} \\
a_{31} & a_{32} & a_{33}
\end{pmatrix}
\]

\section{Tables}
Below is a sample table showing data:

\begin{table}[h]
\centering
\begin{tabular}{@{}lrrr@{}}
\toprule
\textbf{Category} & \textbf{Value 1} & \textbf{Value 2} & \textbf{Total} \\
\midrule
Item A & 10.5 & 20.3 & 30.8 \\
Item B & 15.2 & 25.1 & 40.3 \\
Item C & 22.1 & 12.7 & 34.8 \\
\bottomrule
\end{tabular}
\caption{Sample Data Table}
\label{tab:sample}
\end{table}

\section{Lists}
\subsection{Bullet Points}
\begin{itemize}
  \item First item
  \item Second item
  \item Third item with sub-items:
    \begin{itemize}
      \item Sub-item 1
      \item Sub-item 2
    \end{itemize}
\end{itemize}

\subsection{Numbered List}
\begin{enumerate}
  \item First step
  \item Second step
  \item Third step
\end{enumerate}

\section{Citations and References}
When you use the citation search feature, references will be automatically managed. For example, you might cite a paper like this: \cite{example2024}. All citations will be automatically added to your \texttt{cite.bib} file.

\section{Conclusion}
This template demonstrates basic \LaTeX{} features including mathematical formulas, tables, and lists. The bibliography system is now set up to work seamlessly with Bibby AI's citation management. Feel free to modify it for your needs!

\bibliographystyle{plain}
\bibliography{cite}

\end{document}\documentclass{article}
\usepackage[utf8]{inputenc}
\usepackage{booktabs}
\usepackage{amsmath}
\usepackage{natbib}

\title{My Document}
\author{Your Name}
\date{\today}

\begin{document}

\maketitle

\section{Introduction}
This is a sample document demonstrating various \LaTeX{} features including mathematical equations, tables, and formatting. You can cite papers using the citation search feature, and they will be automatically added to your \texttt{cite.bib} file.

\section{Mathematical Expressions}
Here are some example equations:

\subsection{Quadratic Formula}
The quadratic formula is given by:
\[x = \frac{-b \pm \sqrt{b^2 - 4ac}}{2a}\]

\subsection{Matrix Example}
A sample matrix:
\[
\begin{pmatrix}
a_{11} & a_{12} & a_{13} \\
a_{21} & a_{22} & a_{23} \\
a_{31} & a_{32} & a_{33}
\end{pmatrix}
\]

\section{Tables}
Below is a sample table showing data:

\begin{table}[h]
\centering
\begin{tabular}{@{}lrrr@{}}
\toprule
\textbf{Category} & \textbf{Value 1} & \textbf{Value 2} & \textbf{Total} \\
\midrule
Item A & 10.5 & 20.3 & 30.8 \\
Item B & 15.2 & 25.1 & 40.3 \\
Item C & 22.1 & 12.7 & 34.8 \\
\bottomrule
\end{tabular}
\caption{Sample Data Table}
\label{tab:sample}
\end{table}

\section{Lists}
\subsection{Bullet Points}
\begin{itemize}
  \item First item
  \item Second item
  \item Third item with sub-items:
    \begin{itemize}
      \item Sub-item 1
      \item Sub-item 2
    \end{itemize}
\end{itemize}

\subsection{Numbered List}
\begin{enumerate}
  \item First step
  \item Second step
  \item Third step
\end{enumerate}

\section{Citations and References}
When you use the citation search feature, references will be automatically managed. For example, you might cite a paper like this: \cite{example2024}. All citations will be automatically added to your \texttt{cite.bib} file.

\section{Conclusion}
This template demonstrates basic \LaTeX{} features including mathematical formulas, tables, and lists. The bibliography system is now set up to work seamlessly with Bibby AI's citation management. Feel free to modify it for your needs!

\bibliographystyle{plain}
\bibliography{cite}

\end{document}\documentclass{article}
\usepackage[utf8]{inputenc}
\usepackage{booktabs}
\usepackage{amsmath}
\usepackage{natbib}

\title{My Document}
\author{Your Name}
\date{\today}

\begin{document}

\maketitle

\section{Introduction}
This is a sample document demonstrating various \LaTeX{} features including mathematical equations, tables, and formatting. You can cite papers using the citation search feature, and they will be automatically added to your \texttt{cite.bib} file.

\section{Mathematical Expressions}
Here are some example equations:

\subsection{Quadratic Formula}
The quadratic formula is given by:
\[x = \frac{-b \pm \sqrt{b^2 - 4ac}}{2a}\]

\subsection{Matrix Example}
A sample matrix:
\[
\begin{pmatrix}
a_{11} & a_{12} & a_{13} \\
a_{21} & a_{22} & a_{23} \\
a_{31} & a_{32} & a_{33}
\end{pmatrix}
\]

\section{Tables}
Below is a sample table showing data:

\begin{table}[h]
\centering
\begin{tabular}{@{}lrrr@{}}
\toprule
\textbf{Category} & \textbf{Value 1} & \textbf{Value 2} & \textbf{Total} \\
\midrule
Item A & 10.5 & 20.3 & 30.8 \\
Item B & 15.2 & 25.1 & 40.3 \\
Item C & 22.1 & 12.7 & 34.8 \\
\bottomrule
\end{tabular}
\caption{Sample Data Table}
\label{tab:sample}
\end{table}

\section{Lists}
\subsection{Bullet Points}
\begin{itemize}
  \item First item
  \item Second item
  \item Third item with sub-items:
    \begin{itemize}
      \item Sub-item 1
      \item Sub-item 2
    \end{itemize}
\end{itemize}

\subsection{Numbered List}
\begin{enumerate}
  \item First step
  \item Second step
  \item Third step
\end{enumerate}

\section{Citations and References}
When you use the citation search feature, references will be automatically managed. For example, you might cite a paper like this: \cite{example2024}. All citations will be automatically added to your \texttt{cite.bib} file.

\section{Conclusion}
This template demonstrates basic \LaTeX{} features including mathematical formulas, tables, and lists. The bibliography system is now set up to work seamlessly with Bibby AI's citation management. Feel free to modify it for your needs!

\bibliographystyle{plain}
\bibliography{cite}

\end{document}\documentclass{article}
\usepackage[utf8]{inputenc}
\usepackage{booktabs}
\usepackage{amsmath}
\usepackage{natbib}

\title{My Document}
\author{Your Name}
\date{\today}

\begin{document}

\maketitle

\section{Introduction}
This is a sample document demonstrating various \LaTeX{} features including mathematical equations, tables, and formatting. You can cite papers using the citation search feature, and they will be automatically added to your \texttt{cite.bib} file.

\section{Mathematical Expressions}
Here are some example equations:

\subsection{Quadratic Formula}
The quadratic formula is given by:
\[x = \frac{-b \pm \sqrt{b^2 - 4ac}}{2a}\]

\subsection{Matrix Example}
A sample matrix:
\[
\begin{pmatrix}
a_{11} & a_{12} & a_{13} \\
a_{21} & a_{22} & a_{23} \\
a_{31} & a_{32} & a_{33}
\end{pmatrix}
\]

\section{Tables}
Below is a sample table showing data:

\begin{table}[h]
\centering
\begin{tabular}{@{}lrrr@{}}
\toprule
\textbf{Category} & \textbf{Value 1} & \textbf{Value 2} & \textbf{Total} \\
\midrule
Item A & 10.5 & 20.3 & 30.8 \\
Item B & 15.2 & 25.1 & 40.3 \\
Item C & 22.1 & 12.7 & 34.8 \\
\bottomrule
\end{tabular}
\caption{Sample Data Table}
\label{tab:sample}
\end{table}

\section{Lists}
\subsection{Bullet Points}
\begin{itemize}
  \item First item
  \item Second item
  \item Third item with sub-items:
    \begin{itemize}
      \item Sub-item 1
      \item Sub-item 2
    \end{itemize}
\end{itemize}

\subsection{Numbered List}
\begin{enumerate}
  \item First step
  \item Second step
  \item Third step
\end{enumerate}

\section{Citations and References}
When you use the citation search feature, references will be automatically managed. For example, you might cite a paper like this: \cite{example2024}. All citations will be automatically added to your \texttt{cite.bib} file.

\section{Conclusion}
This template demonstrates basic \LaTeX{} features including mathematical formulas, tables, and lists. The bibliography system is now set up to work seamlessly with Bibby AI's citation management. Feel free to modify it for your needs!

\bibliographystyle{plain}
\bibliography{cite}

\end{document}\documentclass{article}
\usepackage[utf8]{inputenc}
\usepackage{booktabs}
\usepackage{amsmath}
\usepackage{natbib}

\title{My Document}
\author{Your Name}
\date{\today}

\begin{document}

\maketitle

\section{Introduction}
This is a sample document demonstrating various \LaTeX{} features including mathematical equations, tables, and formatting. You can cite papers using the citation search feature, and they will be automatically added to your \texttt{cite.bib} file.

\section{Mathematical Expressions}
Here are some example equations:

\subsection{Quadratic Formula}
The quadratic formula is given by:
\[x = \frac{-b \pm \sqrt{b^2 - 4ac}}{2a}\]

\subsection{Matrix Example}
A sample matrix:
\[
\begin{pmatrix}
a_{11} & a_{12} & a_{13} \\
a_{21} & a_{22} & a_{23} \\
a_{31} & a_{32} & a_{33}
\end{pmatrix}
\]

\section{Tables}
Below is a sample table showing data:

\begin{table}[h]
\centering
\begin{tabular}{@{}lrrr@{}}
\toprule
\textbf{Category} & \textbf{Value 1} & \textbf{Value 2} & \textbf{Total} \\
\midrule
Item A & 10.5 & 20.3 & 30.8 \\
Item B & 15.2 & 25.1 & 40.3 \\
Item C & 22.1 & 12.7 & 34.8 \\
\bottomrule
\end{tabular}
\caption{Sample Data Table}
\label{tab:sample}
\end{table}

\section{Lists}
\subsection{Bullet Points}
\begin{itemize}
  \item First item
  \item Second item
  \item Third item with sub-items:
    \begin{itemize}
      \item Sub-item 1
      \item Sub-item 2
    \end{itemize}
\end{itemize}

\subsection{Numbered List}
\begin{enumerate}
  \item First step
  \item Second step
  \item Third step
\end{enumerate}

\section{Citations and References}
When you use the citation search feature, references will be automatically managed. For example, you might cite a paper like this: \cite{example2024}. All citations will be automatically added to your \texttt{cite.bib} file.

\section{Conclusion}
This template demonstrates basic \LaTeX{} features including mathematical formulas, tables, and lists. The bibliography system is now set up to work seamlessly with Bibby AI's citation management. Feel free to modify it for your needs!

\bibliographystyle{plain}
\bibliography{cite}

\end{document}\documentclass{article}
\usepackage[utf8]{inputenc}
\usepackage{booktabs}
\usepackage{amsmath}
\usepackage{natbib}

\title{My Document}
\author{Your Name}
\date{\today}

\begin{document}

\maketitle

\section{Introduction}
This is a sample document demonstrating various \LaTeX{} features including mathematical equations, tables, and formatting. You can cite papers using the citation search feature, and they will be automatically added to your \texttt{cite.bib} file.

\section{Mathematical Expressions}
Here are some example equations:

\subsection{Quadratic Formula}
The quadratic formula is given by:
\[x = \frac{-b \pm \sqrt{b^2 - 4ac}}{2a}\]

\subsection{Matrix Example}
A sample matrix:
\[
\begin{pmatrix}
a_{11} & a_{12} & a_{13} \\
a_{21} & a_{22} & a_{23} \\
a_{31} & a_{32} & a_{33}
\end{pmatrix}
\]

\section{Tables}
Below is a sample table showing data:

\begin{table}[h]
\centering
\begin{tabular}{@{}lrrr@{}}
\toprule
\textbf{Category} & \textbf{Value 1} & \textbf{Value 2} & \textbf{Total} \\
\midrule
Item A & 10.5 & 20.3 & 30.8 \\
Item B & 15.2 & 25.1 & 40.3 \\
Item C & 22.1 & 12.7 & 34.8 \\
\bottomrule
\end{tabular}
\caption{Sample Data Table}
\label{tab:sample}
\end{table}

\section{Lists}
\subsection{Bullet Points}
\begin{itemize}
  \item First item
  \item Second item
  \item Third item with sub-items:
    \begin{itemize}
      \item Sub-item 1
      \item Sub-item 2
    \end{itemize}
\end{itemize}

\subsection{Numbered List}
\begin{enumerate}
  \item First step
  \item Second step
  \item Third step
\end{enumerate}

\section{Citations and References}
When you use the citation search feature, references will be automatically managed. For example, you might cite a paper like this: \cite{example2024}. All citations will be automatically added to your \texttt{cite.bib} file.

\section{Conclusion}
This template demonstrates basic \LaTeX{} features including mathematical formulas, tables, and lists. The bibliography system is now set up to work seamlessly with Bibby AI's citation management. Feel free to modify it for your needs!

\bibliographystyle{plain}
\bibliography{cite}

\end{document}\documentclass{article}
\usepackage[utf8]{inputenc}
\usepackage{booktabs}
\usepackage{amsmath}
\usepackage{natbib}

\title{My Document}
\author{Your Name}
\date{\today}

\begin{document}

\maketitle

\section{Introduction}
This is a sample document demonstrating various \LaTeX{} features including mathematical equations, tables, and formatting. You can cite papers using the citation search feature, and they will be automatically added to your \texttt{cite.bib} file.

\section{Mathematical Expressions}
Here are some example equations:

\subsection{Quadratic Formula}
The quadratic formula is given by:
\[x = \frac{-b \pm \sqrt{b^2 - 4ac}}{2a}\]

\subsection{Matrix Example}
A sample matrix:
\[
\begin{pmatrix}
a_{11} & a_{12} & a_{13} \\
a_{21} & a_{22} & a_{23} \\
a_{31} & a_{32} & a_{33}
\end{pmatrix}
\]

\section{Tables}
Below is a sample table showing data:

\begin{table}[h]
\centering
\begin{tabular}{@{}lrrr@{}}
\toprule
\textbf{Category} & \textbf{Value 1} & \textbf{Value 2} & \textbf{Total} \\
\midrule
Item A & 10.5 & 20.3 & 30.8 \\
Item B & 15.2 & 25.1 & 40.3 \\
Item C & 22.1 & 12.7 & 34.8 \\
\bottomrule
\end{tabular}
\caption{Sample Data Table}
\label{tab:sample}
\end{table}

\section{Lists}
\subsection{Bullet Points}
\begin{itemize}
  \item First item
  \item Second item
  \item Third item with sub-items:
    \begin{itemize}
      \item Sub-item 1
      \item Sub-item 2
    \end{itemize}
\end{itemize}

\subsection{Numbered List}
\begin{enumerate}
  \item First step
  \item Second step
  \item Third step
\end{enumerate}

\section{Citations and References}
When you use the citation search feature, references will be automatically managed. For example, you might cite a paper like this: \cite{example2024}. All citations will be automatically added to your \texttt{cite.bib} file.

\section{Conclusion}
This template demonstrates basic \LaTeX{} features including mathematical formulas, tables, and lists. The bibliography system is now set up to work seamlessly with Bibby AI's citation management. Feel free to modify it for your needs!

\bibliographystyle{plain}
\bibliography{cite}

\end{document}\documentclass{article}
\usepackage[utf8]{inputenc}
\usepackage{booktabs}
\usepackage{amsmath}
\usepackage{natbib}

\title{My Document}
\author{Your Name}
\date{\today}

\begin{document}

\maketitle

\section{Introduction}
This is a sample document demonstrating various \LaTeX{} features including mathematical equations, tables, and formatting. You can cite papers using the citation search feature, and they will be automatically added to your \texttt{cite.bib} file.

\section{Mathematical Expressions}
Here are some example equations:

\subsection{Quadratic Formula}
The quadratic formula is given by:
\[x = \frac{-b \pm \sqrt{b^2 - 4ac}}{2a}\]

\subsection{Matrix Example}
A sample matrix:
\[
\begin{pmatrix}
a_{11} & a_{12} & a_{13} \\
a_{21} & a_{22} & a_{23} \\
a_{31} & a_{32} & a_{33}
\end{pmatrix}
\]

\section{Tables}
Below is a sample table showing data:

\begin{table}[h]
\centering
\begin{tabular}{@{}lrrr@{}}
\toprule
\textbf{Category} & \textbf{Value 1} & \textbf{Value 2} & \textbf{Total} \\
\midrule
Item A & 10.5 & 20.3 & 30.8 \\
Item B & 15.2 & 25.1 & 40.3 \\
Item C & 22.1 & 12.7 & 34.8 \\
\bottomrule
\end{tabular}
\caption{Sample Data Table}
\label{tab:sample}
\end{table}

\section{Lists}
\subsection{Bullet Points}
\begin{itemize}
  \item First item
  \item Second item
  \item Third item with sub-items:
    \begin{itemize}
      \item Sub-item 1
      \item Sub-item 2
    \end{itemize}
\end{itemize}

\subsection{Numbered List}
\begin{enumerate}
  \item First step
  \item Second step
  \item Third step
\end{enumerate}

\section{Citations and References}
When you use the citation search feature, references will be automatically managed. For example, you might cite a paper like this: \cite{example2024}. All citations will be automatically added to your \texttt{cite.bib} file.

\section{Conclusion}
This template demonstrates basic \LaTeX{} features including mathematical formulas, tables, and lists. The bibliography system is now set up to work seamlessly with Bibby AI's citation management. Feel free to modify it for your needs!

\bibliographystyle{plain}
\bibliography{cite}

\end{document}\documentclass{article}
\usepackage[utf8]{inputenc}
\usepackage{booktabs}
\usepackage{amsmath}
\usepackage{natbib}

\title{My Document}
\author{Your Name}
\date{\today}

\begin{document}

\maketitle

\section{Introduction}
This is a sample document demonstrating various \LaTeX{} features including mathematical equations, tables, and formatting. You can cite papers using the citation search feature, and they will be automatically added to your \texttt{cite.bib} file.

\section{Mathematical Expressions}
Here are some example equations:

\subsection{Quadratic Formula}
The quadratic formula is given by:
\[x = \frac{-b \pm \sqrt{b^2 - 4ac}}{2a}\]

\subsection{Matrix Example}
A sample matrix:
\[
\begin{pmatrix}
a_{11} & a_{12} & a_{13} \\
a_{21} & a_{22} & a_{23} \\
a_{31} & a_{32} & a_{33}
\end{pmatrix}
\]

\section{Tables}
Below is a sample table showing data:

\begin{table}[h]
\centering
\begin{tabular}{@{}lrrr@{}}
\toprule
\textbf{Category} & \textbf{Value 1} & \textbf{Value 2} & \textbf{Total} \\
\midrule
Item A & 10.5 & 20.3 & 30.8 \\
Item B & 15.2 & 25.1 & 40.3 \\
Item C & 22.1 & 12.7 & 34.8 \\
\bottomrule
\end{tabular}
\caption{Sample Data Table}
\label{tab:sample}
\end{table}

\section{Lists}
\subsection{Bullet Points}
\begin{itemize}
  \item First item
  \item Second item
  \item Third item with sub-items:
    \begin{itemize}
      \item Sub-item 1
      \item Sub-item 2
    \end{itemize}
\end{itemize}

\subsection{Numbered List}
\begin{enumerate}
  \item First step
  \item Second step
  \item Third step
\end{enumerate}

\section{Citations and References}
When you use the citation search feature, references will be automatically managed. For example, you might cite a paper like this: \cite{example2024}. All citations will be automatically added to your \texttt{cite.bib} file.

\section{Conclusion}
This template demonstrates basic \LaTeX{} features including mathematical formulas, tables, and lists. The bibliography system is now set up to work seamlessly with Bibby AI's citation management. Feel free to modify it for your needs!

\bibliographystyle{plain}
\bibliography{cite}

\end{document}\documentclass{article}
\usepackage[utf8]{inputenc}
\usepackage{booktabs}
\usepackage{amsmath}
\usepackage{natbib}

\title{My Document}
\author{Your Name}
\date{\today}

\begin{document}

\maketitle

\section{Introduction}
This is a sample document demonstrating various \LaTeX{} features including mathematical equations, tables, and formatting. You can cite papers using the citation search feature, and they will be automatically added to your \texttt{cite.bib} file.

\section{Mathematical Expressions}
Here are some example equations:

\subsection{Quadratic Formula}
The quadratic formula is given by:
\[x = \frac{-b \pm \sqrt{b^2 - 4ac}}{2a}\]

\subsection{Matrix Example}
A sample matrix:
\[
\begin{pmatrix}
a_{11} & a_{12} & a_{13} \\
a_{21} & a_{22} & a_{23} \\
a_{31} & a_{32} & a_{33}
\end{pmatrix}
\]

\section{Tables}
Below is a sample table showing data:

\begin{table}[h]
\centering
\begin{tabular}{@{}lrrr@{}}
\toprule
\textbf{Category} & \textbf{Value 1} & \textbf{Value 2} & \textbf{Total} \\
\midrule
Item A & 10.5 & 20.3 & 30.8 \\
Item B & 15.2 & 25.1 & 40.3 \\
Item C & 22.1 & 12.7 & 34.8 \\
\bottomrule
\end{tabular}
\caption{Sample Data Table}
\label{tab:sample}
\end{table}

\section{Lists}
\subsection{Bullet Points}
\begin{itemize}
  \item First item
  \item Second item
  \item Third item with sub-items:
    \begin{itemize}
      \item Sub-item 1
      \item Sub-item 2
    \end{itemize}
\end{itemize}

\subsection{Numbered List}
\begin{enumerate}
  \item First step
  \item Second step
  \item Third step
\end{enumerate}

\section{Citations and References}
When you use the citation search feature, references will be automatically managed. For example, you might cite a paper like this: \cite{example2024}. All citations will be automatically added to your \texttt{cite.bib} file.

\section{Conclusion}
This template demonstrates basic \LaTeX{} features including mathematical formulas, tables, and lists. The bibliography system is now set up to work seamlessly with Bibby AI's citation management. Feel free to modify it for your needs!

\bibliographystyle{plain}
\bibliography{cite}

\end{document}\documentclass{article}
\usepackage[utf8]{inputenc}
\usepackage{booktabs}
\usepackage{amsmath}
\usepackage{natbib}

\title{My Document}
\author{Your Name}
\date{\today}

\begin{document}

\maketitle

\section{Introduction}
This is a sample document demonstrating various \LaTeX{} features including mathematical equations, tables, and formatting. You can cite papers using the citation search feature, and they will be automatically added to your \texttt{cite.bib} file.

\section{Mathematical Expressions}
Here are some example equations:

\subsection{Quadratic Formula}
The quadratic formula is given by:
\[x = \frac{-b \pm \sqrt{b^2 - 4ac}}{2a}\]

\subsection{Matrix Example}
A sample matrix:
\[
\begin{pmatrix}
a_{11} & a_{12} & a_{13} \\
a_{21} & a_{22} & a_{23} \\
a_{31} & a_{32} & a_{33}
\end{pmatrix}
\]

\section{Tables}
Below is a sample table showing data:

\begin{table}[h]
\centering
\begin{tabular}{@{}lrrr@{}}
\toprule
\textbf{Category} & \textbf{Value 1} & \textbf{Value 2} & \textbf{Total} \\
\midrule
Item A & 10.5 & 20.3 & 30.8 \\
Item B & 15.2 & 25.1 & 40.3 \\
Item C & 22.1 & 12.7 & 34.8 \\
\bottomrule
\end{tabular}
\caption{Sample Data Table}
\label{tab:sample}
\end{table}

\section{Lists}
\subsection{Bullet Points}
\begin{itemize}
  \item First item
  \item Second item
  \item Third item with sub-items:
    \begin{itemize}
      \item Sub-item 1
      \item Sub-item 2
    \end{itemize}
\end{itemize}

\subsection{Numbered List}
\begin{enumerate}
  \item First step
  \item Second step
  \item Third step
\end{enumerate}

\section{Citations and References}
When you use the citation search feature, references will be automatically managed. For example, you might cite a paper like this: \cite{example2024}. All citations will be automatically added to your \texttt{cite.bib} file.

\section{Conclusion}
This template demonstrates basic \LaTeX{} features including mathematical formulas, tables, and lists. The bibliography system is now set up to work seamlessly with Bibby AI's citation management. Feel free to modify it for your needs!

\bibliographystyle{plain}
\bibliography{cite}

\end{document}\documentclass{article}
\usepackage[utf8]{inputenc}
\usepackage{booktabs}
\usepackage{amsmath}
\usepackage{natbib}

\title{My Document}
\author{Your Name}
\date{\today}

\begin{document}

\maketitle

\section{Introduction}
This is a sample document demonstrating various \LaTeX{} features including mathematical equations, tables, and formatting. You can cite papers using the citation search feature, and they will be automatically added to your \texttt{cite.bib} file.

\section{Mathematical Expressions}
Here are some example equations:

\subsection{Quadratic Formula}
The quadratic formula is given by:
\[x = \frac{-b \pm \sqrt{b^2 - 4ac}}{2a}\]

\subsection{Matrix Example}
A sample matrix:
\[
\begin{pmatrix}
a_{11} & a_{12} & a_{13} \\
a_{21} & a_{22} & a_{23} \\
a_{31} & a_{32} & a_{33}
\end{pmatrix}
\]

\section{Tables}
Below is a sample table showing data:

\begin{table}[h]
\centering
\begin{tabular}{@{}lrrr@{}}
\toprule
\textbf{Category} & \textbf{Value 1} & \textbf{Value 2} & \textbf{Total} \\
\midrule
Item A & 10.5 & 20.3 & 30.8 \\
Item B & 15.2 & 25.1 & 40.3 \\
Item C & 22.1 & 12.7 & 34.8 \\
\bottomrule
\end{tabular}
\caption{Sample Data Table}
\label{tab:sample}
\end{table}

\section{Lists}
\subsection{Bullet Points}
\begin{itemize}
  \item First item
  \item Second item
  \item Third item with sub-items:
    \begin{itemize}
      \item Sub-item 1
      \item Sub-item 2
    \end{itemize}
\end{itemize}

\subsection{Numbered List}
\begin{enumerate}
  \item First step
  \item Second step
  \item Third step
\end{enumerate}

\section{Citations and References}
When you use the citation search feature, references will be automatically managed. For example, you might cite a paper like this: \cite{example2024}. All citations will be automatically added to your \texttt{cite.bib} file.

\section{Conclusion}
This template demonstrates basic \LaTeX{} features including mathematical formulas, tables, and lists. The bibliography system is now set up to work seamlessly with Bibby AI's citation management. Feel free to modify it for your needs!

\bibliographystyle{plain}
\bibliography{cite}

\end{document}\documentclass{article}
\usepackage[utf8]{inputenc}
\usepackage{booktabs}
\usepackage{amsmath}
\usepackage{natbib}

\title{My Document}
\author{Your Name}
\date{\today}

\begin{document}

\maketitle

\section{Introduction}
This is a sample document demonstrating various \LaTeX{} features including mathematical equations, tables, and formatting. You can cite papers using the citation search feature, and they will be automatically added to your \texttt{cite.bib} file.

\section{Mathematical Expressions}
Here are some example equations:

\subsection{Quadratic Formula}
The quadratic formula is given by:
\[x = \frac{-b \pm \sqrt{b^2 - 4ac}}{2a}\]

\subsection{Matrix Example}
A sample matrix:
\[
\begin{pmatrix}
a_{11} & a_{12} & a_{13} \\
a_{21} & a_{22} & a_{23} \\
a_{31} & a_{32} & a_{33}
\end{pmatrix}
\]

\section{Tables}
Below is a sample table showing data:

\begin{table}[h]
\centering
\begin{tabular}{@{}lrrr@{}}
\toprule
\textbf{Category} & \textbf{Value 1} & \textbf{Value 2} & \textbf{Total} \\
\midrule
Item A & 10.5 & 20.3 & 30.8 \\
Item B & 15.2 & 25.1 & 40.3 \\
Item C & 22.1 & 12.7 & 34.8 \\
\bottomrule
\end{tabular}
\caption{Sample Data Table}
\label{tab:sample}
\end{table}

\section{Lists}
\subsection{Bullet Points}
\begin{itemize}
  \item First item
  \item Second item
  \item Third item with sub-items:
    \begin{itemize}
      \item Sub-item 1
      \item Sub-item 2
    \end{itemize}
\end{itemize}

\subsection{Numbered List}
\begin{enumerate}
  \item First step
  \item Second step
  \item Third step
\end{enumerate}

\section{Citations and References}
When you use the citation search feature, references will be automatically managed. For example, you might cite a paper like this: \cite{example2024}. All citations will be automatically added to your \texttt{cite.bib} file.

\section{Conclusion}
This template demonstrates basic \LaTeX{} features including mathematical formulas, tables, and lists. The bibliography system is now set up to work seamlessly with Bibby AI's citation management. Feel free to modify it for your needs!

\bibliographystyle{plain}
\bibliography{cite}

\end{document}\documentclass{article}
\usepackage[utf8]{inputenc}
\usepackage{booktabs}
\usepackage{amsmath}
\usepackage{natbib}

\title{My Document}
\author{Your Name}
\date{\today}

\begin{document}

\maketitle

\section{Introduction}
This is a sample document demonstrating various \LaTeX{} features including mathematical equations, tables, and formatting. You can cite papers using the citation search feature, and they will be automatically added to your \texttt{cite.bib} file.

\section{Mathematical Expressions}
Here are some example equations:

\subsection{Quadratic Formula}
The quadratic formula is given by:
\[x = \frac{-b \pm \sqrt{b^2 - 4ac}}{2a}\]

\subsection{Matrix Example}
A sample matrix:
\[
\begin{pmatrix}
a_{11} & a_{12} & a_{13} \\
a_{21} & a_{22} & a_{23} \\
a_{31} & a_{32} & a_{33}
\end{pmatrix}
\]

\section{Tables}
Below is a sample table showing data:

\begin{table}[h]
\centering
\begin{tabular}{@{}lrrr@{}}
\toprule
\textbf{Category} & \textbf{Value 1} & \textbf{Value 2} & \textbf{Total} \\
\midrule
Item A & 10.5 & 20.3 & 30.8 \\
Item B & 15.2 & 25.1 & 40.3 \\
Item C & 22.1 & 12.7 & 34.8 \\
\bottomrule
\end{tabular}
\caption{Sample Data Table}
\label{tab:sample}
\end{table}

\section{Lists}
\subsection{Bullet Points}
\begin{itemize}
  \item First item
  \item Second item
  \item Third item with sub-items:
    \begin{itemize}
      \item Sub-item 1
      \item Sub-item 2
    \end{itemize}
\end{itemize}

\subsection{Numbered List}
\begin{enumerate}
  \item First step
  \item Second step
  \item Third step
\end{enumerate}

\section{Citations and References}
When you use the citation search feature, references will be automatically managed. For example, you might cite a paper like this: \cite{example2024}. All citations will be automatically added to your \texttt{cite.bib} file.

\section{Conclusion}
This template demonstrates basic \LaTeX{} features including mathematical formulas, tables, and lists. The bibliography system is now set up to work seamlessly with Bibby AI's citation management. Feel free to modify it for your needs!

\bibliographystyle{plain}
\bibliography{cite}

\end{document}\documentclass{article}
\usepackage[utf8]{inputenc}
\usepackage{booktabs}
\usepackage{amsmath}
\usepackage{natbib}

\title{My Document}
\author{Your Name}
\date{\today}

\begin{document}

\maketitle

\section{Introduction}
This is a sample document demonstrating various \LaTeX{} features including mathematical equations, tables, and formatting. You can cite papers using the citation search feature, and they will be automatically added to your \texttt{cite.bib} file.

\section{Mathematical Expressions}
Here are some example equations:

\subsection{Quadratic Formula}
The quadratic formula is given by:
\[x = \frac{-b \pm \sqrt{b^2 - 4ac}}{2a}\]

\subsection{Matrix Example}
A sample matrix:
\[
\begin{pmatrix}
a_{11} & a_{12} & a_{13} \\
a_{21} & a_{22} & a_{23} \\
a_{31} & a_{32} & a_{33}
\end{pmatrix}
\]

\section{Tables}
Below is a sample table showing data:

\begin{table}[h]
\centering
\begin{tabular}{@{}lrrr@{}}
\toprule
\textbf{Category} & \textbf{Value 1} & \textbf{Value 2} & \textbf{Total} \\
\midrule
Item A & 10.5 & 20.3 & 30.8 \\
Item B & 15.2 & 25.1 & 40.3 \\
Item C & 22.1 & 12.7 & 34.8 \\
\bottomrule
\end{tabular}
\caption{Sample Data Table}
\label{tab:sample}
\end{table}

\section{Lists}
\subsection{Bullet Points}
\begin{itemize}
  \item First item
  \item Second item
  \item Third item with sub-items:
    \begin{itemize}
      \item Sub-item 1
      \item Sub-item 2
    \end{itemize}
\end{itemize}

\subsection{Numbered List}
\begin{enumerate}
  \item First step
  \item Second step
  \item Third step
\end{enumerate}

\section{Citations and References}
When you use the citation search feature, references will be automatically managed. For example, you might cite a paper like this: \cite{example2024}. All citations will be automatically added to your \texttt{cite.bib} file.

\section{Conclusion}
This template demonstrates basic \LaTeX{} features including mathematical formulas, tables, and lists. The bibliography system is now set up to work seamlessly with Bibby AI's citation management. Feel free to modify it for your needs!

\bibliographystyle{plain}
\bibliography{cite}

\end{document}\documentclass{article}
\usepackage[utf8]{inputenc}
\usepackage{booktabs}
\usepackage{amsmath}
\usepackage{natbib}

\title{My Document}
\author{Your Name}
\date{\today}

\begin{document}

\maketitle

\section{Introduction}
This is a sample document demonstrating various \LaTeX{} features including mathematical equations, tables, and formatting. You can cite papers using the citation search feature, and they will be automatically added to your \texttt{cite.bib} file.

\section{Mathematical Expressions}
Here are some example equations:

\subsection{Quadratic Formula}
The quadratic formula is given by:
\[x = \frac{-b \pm \sqrt{b^2 - 4ac}}{2a}\]

\subsection{Matrix Example}
A sample matrix:
\[
\begin{pmatrix}
a_{11} & a_{12} & a_{13} \\
a_{21} & a_{22} & a_{23} \\
a_{31} & a_{32} & a_{33}
\end{pmatrix}
\]

\section{Tables}
Below is a sample table showing data:

\begin{table}[h]
\centering
\begin{tabular}{@{}lrrr@{}}
\toprule
\textbf{Category} & \textbf{Value 1} & \textbf{Value 2} & \textbf{Total} \\
\midrule
Item A & 10.5 & 20.3 & 30.8 \\
Item B & 15.2 & 25.1 & 40.3 \\
Item C & 22.1 & 12.7 & 34.8 \\
\bottomrule
\end{tabular}
\caption{Sample Data Table}
\label{tab:sample}
\end{table}

\section{Lists}
\subsection{Bullet Points}
\begin{itemize}
  \item First item
  \item Second item
  \item Third item with sub-items:
    \begin{itemize}
      \item Sub-item 1
      \item Sub-item 2
    \end{itemize}
\end{itemize}

\subsection{Numbered List}
\begin{enumerate}
  \item First step
  \item Second step
  \item Third step
\end{enumerate}

\section{Citations and References}
When you use the citation search feature, references will be automatically managed. For example, you might cite a paper like this: \cite{example2024}. All citations will be automatically added to your \texttt{cite.bib} file.

\section{Conclusion}
This template demonstrates basic \LaTeX{} features including mathematical formulas, tables, and lists. The bibliography system is now set up to work seamlessly with Bibby AI's citation management. Feel free to modify it for your needs!

\bibliographystyle{plain}
\bibliography{cite}

\end{document}\documentclass{article}
\usepackage[utf8]{inputenc}
\usepackage{booktabs}
\usepackage{amsmath}
\usepackage{natbib}

\title{My Document}
\author{Your Name}
\date{\today}

\begin{document}

\maketitle

\section{Introduction}
This is a sample document demonstrating various \LaTeX{} features including mathematical equations, tables, and formatting. You can cite papers using the citation search feature, and they will be automatically added to your \texttt{cite.bib} file.

\section{Mathematical Expressions}
Here are some example equations:

\subsection{Quadratic Formula}
The quadratic formula is given by:
\[x = \frac{-b \pm \sqrt{b^2 - 4ac}}{2a}\]

\subsection{Matrix Example}
A sample matrix:
\[
\begin{pmatrix}
a_{11} & a_{12} & a_{13} \\
a_{21} & a_{22} & a_{23} \\
a_{31} & a_{32} & a_{33}
\end{pmatrix}
\]

\section{Tables}
Below is a sample table showing data:

\begin{table}[h]
\centering
\begin{tabular}{@{}lrrr@{}}
\toprule
\textbf{Category} & \textbf{Value 1} & \textbf{Value 2} & \textbf{Total} \\
\midrule
Item A & 10.5 & 20.3 & 30.8 \\
Item B & 15.2 & 25.1 & 40.3 \\
Item C & 22.1 & 12.7 & 34.8 \\
\bottomrule
\end{tabular}
\caption{Sample Data Table}
\label{tab:sample}
\end{table}

\section{Lists}
\subsection{Bullet Points}
\begin{itemize}
  \item First item
  \item Second item
  \item Third item with sub-items:
    \begin{itemize}
      \item Sub-item 1
      \item Sub-item 2
    \end{itemize}
\end{itemize}

\subsection{Numbered List}
\begin{enumerate}
  \item First step
  \item Second step
  \item Third step
\end{enumerate}

\section{Citations and References}
When you use the citation search feature, references will be automatically managed. For example, you might cite a paper like this: \cite{example2024}. All citations will be automatically added to your \texttt{cite.bib} file.

\section{Conclusion}
This template demonstrates basic \LaTeX{} features including mathematical formulas, tables, and lists. The bibliography system is now set up to work seamlessly with Bibby AI's citation management. Feel free to modify it for your needs!

\bibliographystyle{plain}
\bibliography{cite}

\end{document}\documentclass{article}
\usepackage[utf8]{inputenc}
\usepackage{booktabs}
\usepackage{amsmath}
\usepackage{natbib}

\title{My Document}
\author{Your Name}
\date{\today}

\begin{document}

\maketitle

\section{Introduction}
This is a sample document demonstrating various \LaTeX{} features including mathematical equations, tables, and formatting. You can cite papers using the citation search feature, and they will be automatically added to your \texttt{cite.bib} file.

\section{Mathematical Expressions}
Here are some example equations:

\subsection{Quadratic Formula}
The quadratic formula is given by:
\[x = \frac{-b \pm \sqrt{b^2 - 4ac}}{2a}\]

\subsection{Matrix Example}
A sample matrix:
\[
\begin{pmatrix}
a_{11} & a_{12} & a_{13} \\
a_{21} & a_{22} & a_{23} \\
a_{31} & a_{32} & a_{33}
\end{pmatrix}
\]

\section{Tables}
Below is a sample table showing data:

\begin{table}[h]
\centering
\begin{tabular}{@{}lrrr@{}}
\toprule
\textbf{Category} & \textbf{Value 1} & \textbf{Value 2} & \textbf{Total} \\
\midrule
Item A & 10.5 & 20.3 & 30.8 \\
Item B & 15.2 & 25.1 & 40.3 \\
Item C & 22.1 & 12.7 & 34.8 \\
\bottomrule
\end{tabular}
\caption{Sample Data Table}
\label{tab:sample}
\end{table}

\section{Lists}
\subsection{Bullet Points}
\begin{itemize}
  \item First item
  \item Second item
  \item Third item with sub-items:
    \begin{itemize}
      \item Sub-item 1
      \item Sub-item 2
    \end{itemize}
\end{itemize}

\subsection{Numbered List}
\begin{enumerate}
  \item First step
  \item Second step
  \item Third step
\end{enumerate}

\section{Citations and References}
When you use the citation search feature, references will be automatically managed. For example, you might cite a paper like this: \cite{example2024}. All citations will be automatically added to your \texttt{cite.bib} file.

\section{Conclusion}
This template demonstrates basic \LaTeX{} features including mathematical formulas, tables, and lists. The bibliography system is now set up to work seamlessly with Bibby AI's citation management. Feel free to modify it for your needs!

\bibliographystyle{plain}
\bibliography{cite}

\end{document}\documentclass{article}
\usepackage[utf8]{inputenc}
\usepackage{booktabs}
\usepackage{amsmath}
\usepackage{natbib}

\title{My Document}
\author{Your Name}
\date{\today}

\begin{document}

\maketitle

\section{Introduction}
This is a sample document demonstrating various \LaTeX{} features including mathematical equations, tables, and formatting. You can cite papers using the citation search feature, and they will be automatically added to your \texttt{cite.bib} file.

\section{Mathematical Expressions}
Here are some example equations:

\subsection{Quadratic Formula}
The quadratic formula is given by:
\[x = \frac{-b \pm \sqrt{b^2 - 4ac}}{2a}\]

\subsection{Matrix Example}
A sample matrix:
\[
\begin{pmatrix}
a_{11} & a_{12} & a_{13} \\
a_{21} & a_{22} & a_{23} \\
a_{31} & a_{32} & a_{33}
\end{pmatrix}
\]

\section{Tables}
Below is a sample table showing data:

\begin{table}[h]
\centering
\begin{tabular}{@{}lrrr@{}}
\toprule
\textbf{Category} & \textbf{Value 1} & \textbf{Value 2} & \textbf{Total} \\
\midrule
Item A & 10.5 & 20.3 & 30.8 \\
Item B & 15.2 & 25.1 & 40.3 \\
Item C & 22.1 & 12.7 & 34.8 \\
\bottomrule
\end{tabular}
\caption{Sample Data Table}
\label{tab:sample}
\end{table}

\section{Lists}
\subsection{Bullet Points}
\begin{itemize}
  \item First item
  \item Second item
  \item Third item with sub-items:
    \begin{itemize}
      \item Sub-item 1
      \item Sub-item 2
    \end{itemize}
\end{itemize}

\subsection{Numbered List}
\begin{enumerate}
  \item First step
  \item Second step
  \item Third step
\end{enumerate}

\section{Citations and References}
When you use the citation search feature, references will be automatically managed. For example, you might cite a paper like this: \cite{example2024}. All citations will be automatically added to your \texttt{cite.bib} file.

\section{Conclusion}
This template demonstrates basic \LaTeX{} features including mathematical formulas, tables, and lists. The bibliography system is now set up to work seamlessly with Bibby AI's citation management. Feel free to modify it for your needs!

\bibliographystyle{plain}
\bibliography{cite}

\end{document}\documentclass{article}
\usepackage[utf8]{inputenc}
\usepackage{booktabs}
\usepackage{amsmath}
\usepackage{natbib}

\title{My Document}
\author{Your Name}
\date{\today}

\begin{document}

\maketitle

\section{Introduction}
This is a sample document demonstrating various \LaTeX{} features including mathematical equations, tables, and formatting. You can cite papers using the citation search feature, and they will be automatically added to your \texttt{cite.bib} file.

\section{Mathematical Expressions}
Here are some example equations:

\subsection{Quadratic Formula}
The quadratic formula is given by:
\[x = \frac{-b \pm \sqrt{b^2 - 4ac}}{2a}\]

\subsection{Matrix Example}
A sample matrix:
\[
\begin{pmatrix}
a_{11} & a_{12} & a_{13} \\
a_{21} & a_{22} & a_{23} \\
a_{31} & a_{32} & a_{33}
\end{pmatrix}
\]

\section{Tables}
Below is a sample table showing data:

\begin{table}[h]
\centering
\begin{tabular}{@{}lrrr@{}}
\toprule
\textbf{Category} & \textbf{Value 1} & \textbf{Value 2} & \textbf{Total} \\
\midrule
Item A & 10.5 & 20.3 & 30.8 \\
Item B & 15.2 & 25.1 & 40.3 \\
Item C & 22.1 & 12.7 & 34.8 \\
\bottomrule
\end{tabular}
\caption{Sample Data Table}
\label{tab:sample}
\end{table}

\section{Lists}
\subsection{Bullet Points}
\begin{itemize}
  \item First item
  \item Second item
  \item Third item with sub-items:
    \begin{itemize}
      \item Sub-item 1
      \item Sub-item 2
    \end{itemize}
\end{itemize}

\subsection{Numbered List}
\begin{enumerate}
  \item First step
  \item Second step
  \item Third step
\end{enumerate}

\section{Citations and References}
When you use the citation search feature, references will be automatically managed. For example, you might cite a paper like this: \cite{example2024}. All citations will be automatically added to your \texttt{cite.bib} file.

\section{Conclusion}
This template demonstrates basic \LaTeX{} features including mathematical formulas, tables, and lists. The bibliography system is now set up to work seamlessly with Bibby AI's citation management. Feel free to modify it for your needs!

\bibliographystyle{plain}
\bibliography{cite}

\end{document}\documentclass{article}
\usepackage[utf8]{inputenc}
\usepackage{booktabs}
\usepackage{amsmath}
\usepackage{natbib}

\title{My Document}
\author{Your Name}
\date{\today}

\begin{document}

\maketitle

\section{Introduction}
This is a sample document demonstrating various \LaTeX{} features including mathematical equations, tables, and formatting. You can cite papers using the citation search feature, and they will be automatically added to your \texttt{cite.bib} file.

\section{Mathematical Expressions}
Here are some example equations:

\subsection{Quadratic Formula}
The quadratic formula is given by:
\[x = \frac{-b \pm \sqrt{b^2 - 4ac}}{2a}\]

\subsection{Matrix Example}
A sample matrix:
\[
\begin{pmatrix}
a_{11} & a_{12} & a_{13} \\
a_{21} & a_{22} & a_{23} \\
a_{31} & a_{32} & a_{33}
\end{pmatrix}
\]

\section{Tables}
Below is a sample table showing data:

\begin{table}[h]
\centering
\begin{tabular}{@{}lrrr@{}}
\toprule
\textbf{Category} & \textbf{Value 1} & \textbf{Value 2} & \textbf{Total} \\
\midrule
Item A & 10.5 & 20.3 & 30.8 \\
Item B & 15.2 & 25.1 & 40.3 \\
Item C & 22.1 & 12.7 & 34.8 \\
\bottomrule
\end{tabular}
\caption{Sample Data Table}
\label{tab:sample}
\end{table}

\section{Lists}
\subsection{Bullet Points}
\begin{itemize}
  \item First item
  \item Second item
  \item Third item with sub-items:
    \begin{itemize}
      \item Sub-item 1
      \item Sub-item 2
    \end{itemize}
\end{itemize}

\subsection{Numbered List}
\begin{enumerate}
  \item First step
  \item Second step
  \item Third step
\end{enumerate}

\section{Citations and References}
When you use the citation search feature, references will be automatically managed. For example, you might cite a paper like this: \cite{example2024}. All citations will be automatically added to your \texttt{cite.bib} file.

\section{Conclusion}
This template demonstrates basic \LaTeX{} features including mathematical formulas, tables, and lists. The bibliography system is now set up to work seamlessly with Bibby AI's citation management. Feel free to modify it for your needs!

\bibliographystyle{plain}
\bibliography{cite}

\end{document}\documentclass{article}
\usepackage[utf8]{inputenc}
\usepackage{booktabs}
\usepackage{amsmath}
\usepackage{natbib}

\title{My Document}
\author{Your Name}
\date{\today}

\begin{document}

\maketitle

\section{Introduction}
This is a sample document demonstrating various \LaTeX{} features including mathematical equations, tables, and formatting. You can cite papers using the citation search feature, and they will be automatically added to your \texttt{cite.bib} file.

\section{Mathematical Expressions}
Here are some example equations:

\subsection{Quadratic Formula}
The quadratic formula is given by:
\[x = \frac{-b \pm \sqrt{b^2 - 4ac}}{2a}\]

\subsection{Matrix Example}
A sample matrix:
\[
\begin{pmatrix}
a_{11} & a_{12} & a_{13} \\
a_{21} & a_{22} & a_{23} \\
a_{31} & a_{32} & a_{33}
\end{pmatrix}
\]

\section{Tables}
Below is a sample table showing data:

\begin{table}[h]
\centering
\begin{tabular}{@{}lrrr@{}}
\toprule
\textbf{Category} & \textbf{Value 1} & \textbf{Value 2} & \textbf{Total} \\
\midrule
Item A & 10.5 & 20.3 & 30.8 \\
Item B & 15.2 & 25.1 & 40.3 \\
Item C & 22.1 & 12.7 & 34.8 \\
\bottomrule
\end{tabular}
\caption{Sample Data Table}
\label{tab:sample}
\end{table}

\section{Lists}
\subsection{Bullet Points}
\begin{itemize}
  \item First item
  \item Second item
  \item Third item with sub-items:
    \begin{itemize}
      \item Sub-item 1
      \item Sub-item 2
    \end{itemize}
\end{itemize}

\subsection{Numbered List}
\begin{enumerate}
  \item First step
  \item Second step
  \item Third step
\end{enumerate}

\section{Citations and References}
When you use the citation search feature, references will be automatically managed. For example, you might cite a paper like this: \cite{example2024}. All citations will be automatically added to your \texttt{cite.bib} file.

\section{Conclusion}
This template demonstrates basic \LaTeX{} features including mathematical formulas, tables, and lists. The bibliography system is now set up to work seamlessly with Bibby AI's citation management. Feel free to modify it for your needs!

\bibliographystyle{plain}
\bibliography{cite}

\end{document}\documentclass{article}
\usepackage[utf8]{inputenc}
\usepackage{booktabs}
\usepackage{amsmath}
\usepackage{natbib}

\title{My Document}
\author{Your Name}
\date{\today}

\begin{document}

\maketitle

\section{Introduction}
This is a sample document demonstrating various \LaTeX{} features including mathematical equations, tables, and formatting. You can cite papers using the citation search feature, and they will be automatically added to your \texttt{cite.bib} file.

\section{Mathematical Expressions}
Here are some example equations:

\subsection{Quadratic Formula}
The quadratic formula is given by:
\[x = \frac{-b \pm \sqrt{b^2 - 4ac}}{2a}\]

\subsection{Matrix Example}
A sample matrix:
\[
\begin{pmatrix}
a_{11} & a_{12} & a_{13} \\
a_{21} & a_{22} & a_{23} \\
a_{31} & a_{32} & a_{33}
\end{pmatrix}
\]

\section{Tables}
Below is a sample table showing data:

\begin{table}[h]
\centering
\begin{tabular}{@{}lrrr@{}}
\toprule
\textbf{Category} & \textbf{Value 1} & \textbf{Value 2} & \textbf{Total} \\
\midrule
Item A & 10.5 & 20.3 & 30.8 \\
Item B & 15.2 & 25.1 & 40.3 \\
Item C & 22.1 & 12.7 & 34.8 \\
\bottomrule
\end{tabular}
\caption{Sample Data Table}
\label{tab:sample}
\end{table}

\section{Lists}
\subsection{Bullet Points}
\begin{itemize}
  \item First item
  \item Second item
  \item Third item with sub-items:
    \begin{itemize}
      \item Sub-item 1
      \item Sub-item 2
    \end{itemize}
\end{itemize}

\subsection{Numbered List}
\begin{enumerate}
  \item First step
  \item Second step
  \item Third step
\end{enumerate}

\section{Citations and References}
When you use the citation search feature, references will be automatically managed. For example, you might cite a paper like this: \cite{example2024}. All citations will be automatically added to your \texttt{cite.bib} file.

\section{Conclusion}
This template demonstrates basic \LaTeX{} features including mathematical formulas, tables, and lists. The bibliography system is now set up to work seamlessly with Bibby AI's citation management. Feel free to modify it for your needs!

\bibliographystyle{plain}
\bibliography{cite}

\end{document}\documentclass{article}
\usepackage[utf8]{inputenc}
\usepackage{booktabs}
\usepackage{amsmath}
\usepackage{natbib}

\title{My Document}
\author{Your Name}
\date{\today}

\begin{document}

\maketitle

\section{Introduction}
This is a sample document demonstrating various \LaTeX{} features including mathematical equations, tables, and formatting. You can cite papers using the citation search feature, and they will be automatically added to your \texttt{cite.bib} file.

\section{Mathematical Expressions}
Here are some example equations:

\subsection{Quadratic Formula}
The quadratic formula is given by:
\[x = \frac{-b \pm \sqrt{b^2 - 4ac}}{2a}\]

\subsection{Matrix Example}
A sample matrix:
\[
\begin{pmatrix}
a_{11} & a_{12} & a_{13} \\
a_{21} & a_{22} & a_{23} \\
a_{31} & a_{32} & a_{33}
\end{pmatrix}
\]

\section{Tables}
Below is a sample table showing data:

\begin{table}[h]
\centering
\begin{tabular}{@{}lrrr@{}}
\toprule
\textbf{Category} & \textbf{Value 1} & \textbf{Value 2} & \textbf{Total} \\
\midrule
Item A & 10.5 & 20.3 & 30.8 \\
Item B & 15.2 & 25.1 & 40.3 \\
Item C & 22.1 & 12.7 & 34.8 \\
\bottomrule
\end{tabular}
\caption{Sample Data Table}
\label{tab:sample}
\end{table}

\section{Lists}
\subsection{Bullet Points}
\begin{itemize}
  \item First item
  \item Second item
  \item Third item with sub-items:
    \begin{itemize}
      \item Sub-item 1
      \item Sub-item 2
    \end{itemize}
\end{itemize}

\subsection{Numbered List}
\begin{enumerate}
  \item First step
  \item Second step
  \item Third step
\end{enumerate}

\section{Citations and References}
When you use the citation search feature, references will be automatically managed. For example, you might cite a paper like this: \cite{example2024}. All citations will be automatically added to your \texttt{cite.bib} file.

\section{Conclusion}
This template demonstrates basic \LaTeX{} features including mathematical formulas, tables, and lists. The bibliography system is now set up to work seamlessly with Bibby AI's citation management. Feel free to modify it for your needs!

\bibliographystyle{plain}
\bibliography{cite}

\end{document}\documentclass{article}
\usepackage[utf8]{inputenc}
\usepackage{booktabs}
\usepackage{amsmath}
\usepackage{natbib}

\title{My Document}
\author{Your Name}
\date{\today}

\begin{document}

\maketitle

\section{Introduction}
This is a sample document demonstrating various \LaTeX{} features including mathematical equations, tables, and formatting. You can cite papers using the citation search feature, and they will be automatically added to your \texttt{cite.bib} file.

\section{Mathematical Expressions}
Here are some example equations:

\subsection{Quadratic Formula}
The quadratic formula is given by:
\[x = \frac{-b \pm \sqrt{b^2 - 4ac}}{2a}\]

\subsection{Matrix Example}
A sample matrix:
\[
\begin{pmatrix}
a_{11} & a_{12} & a_{13} \\
a_{21} & a_{22} & a_{23} \\
a_{31} & a_{32} & a_{33}
\end{pmatrix}
\]

\section{Tables}
Below is a sample table showing data:

\begin{table}[h]
\centering
\begin{tabular}{@{}lrrr@{}}
\toprule
\textbf{Category} & \textbf{Value 1} & \textbf{Value 2} & \textbf{Total} \\
\midrule
Item A & 10.5 & 20.3 & 30.8 \\
Item B & 15.2 & 25.1 & 40.3 \\
Item C & 22.1 & 12.7 & 34.8 \\
\bottomrule
\end{tabular}
\caption{Sample Data Table}
\label{tab:sample}
\end{table}

\section{Lists}
\subsection{Bullet Points}
\begin{itemize}
  \item First item
  \item Second item
  \item Third item with sub-items:
    \begin{itemize}
      \item Sub-item 1
      \item Sub-item 2
    \end{itemize}
\end{itemize}

\subsection{Numbered List}
\begin{enumerate}
  \item First step
  \item Second step
  \item Third step
\end{enumerate}

\section{Citations and References}
When you use the citation search feature, references will be automatically managed. For example, you might cite a paper like this: \cite{example2024}. All citations will be automatically added to your \texttt{cite.bib} file.

\section{Conclusion}
This template demonstrates basic \LaTeX{} features including mathematical formulas, tables, and lists. The bibliography system is now set up to work seamlessly with Bibby AI's citation management. Feel free to modify it for your needs!

\bibliographystyle{plain}
\bibliography{cite}

\end{document}\documentclass{article}
\usepackage[utf8]{inputenc}
\usepackage{booktabs}
\usepackage{amsmath}
\usepackage{natbib}

\title{My Document}
\author{Your Name}
\date{\today}

\begin{document}

\maketitle

\section{Introduction}
This is a sample document demonstrating various \LaTeX{} features including mathematical equations, tables, and formatting. You can cite papers using the citation search feature, and they will be automatically added to your \texttt{cite.bib} file.

\section{Mathematical Expressions}
Here are some example equations:

\subsection{Quadratic Formula}
The quadratic formula is given by:
\[x = \frac{-b \pm \sqrt{b^2 - 4ac}}{2a}\]

\subsection{Matrix Example}
A sample matrix:
\[
\begin{pmatrix}
a_{11} & a_{12} & a_{13} \\
a_{21} & a_{22} & a_{23} \\
a_{31} & a_{32} & a_{33}
\end{pmatrix}
\]

\section{Tables}
Below is a sample table showing data:

\begin{table}[h]
\centering
\begin{tabular}{@{}lrrr@{}}
\toprule
\textbf{Category} & \textbf{Value 1} & \textbf{Value 2} & \textbf{Total} \\
\midrule
Item A & 10.5 & 20.3 & 30.8 \\
Item B & 15.2 & 25.1 & 40.3 \\
Item C & 22.1 & 12.7 & 34.8 \\
\bottomrule
\end{tabular}
\caption{Sample Data Table}
\label{tab:sample}
\end{table}

\section{Lists}
\subsection{Bullet Points}
\begin{itemize}
  \item First item
  \item Second item
  \item Third item with sub-items:
    \begin{itemize}
      \item Sub-item 1
      \item Sub-item 2
    \end{itemize}
\end{itemize}

\subsection{Numbered List}
\begin{enumerate}
  \item First step
  \item Second step
  \item Third step
\end{enumerate}

\section{Citations and References}
When you use the citation search feature, references will be automatically managed. For example, you might cite a paper like this: \cite{example2024}. All citations will be automatically added to your \texttt{cite.bib} file.

\section{Conclusion}
This template demonstrates basic \LaTeX{} features including mathematical formulas, tables, and lists. The bibliography system is now set up to work seamlessly with Bibby AI's citation management. Feel free to modify it for your needs!

\bibliographystyle{plain}
\bibliography{cite}

\end{document}\documentclass{article}
\usepackage[utf8]{inputenc}
\usepackage{booktabs}
\usepackage{amsmath}
\usepackage{natbib}

\title{My Document}
\author{Your Name}
\date{\today}

\begin{document}

\maketitle

\section{Introduction}
This is a sample document demonstrating various \LaTeX{} features including mathematical equations, tables, and formatting. You can cite papers using the citation search feature, and they will be automatically added to your \texttt{cite.bib} file.

\section{Mathematical Expressions}
Here are some example equations:

\subsection{Quadratic Formula}
The quadratic formula is given by:
\[x = \frac{-b \pm \sqrt{b^2 - 4ac}}{2a}\]

\subsection{Matrix Example}
A sample matrix:
\[
\begin{pmatrix}
a_{11} & a_{12} & a_{13} \\
a_{21} & a_{22} & a_{23} \\
a_{31} & a_{32} & a_{33}
\end{pmatrix}
\]

\section{Tables}
Below is a sample table showing data:

\begin{table}[h]
\centering
\begin{tabular}{@{}lrrr@{}}
\toprule
\textbf{Category} & \textbf{Value 1} & \textbf{Value 2} & \textbf{Total} \\
\midrule
Item A & 10.5 & 20.3 & 30.8 \\
Item B & 15.2 & 25.1 & 40.3 \\
Item C & 22.1 & 12.7 & 34.8 \\
\bottomrule
\end{tabular}
\caption{Sample Data Table}
\label{tab:sample}
\end{table}

\section{Lists}
\subsection{Bullet Points}
\begin{itemize}
  \item First item
  \item Second item
  \item Third item with sub-items:
    \begin{itemize}
      \item Sub-item 1
      \item Sub-item 2
    \end{itemize}
\end{itemize}

\subsection{Numbered List}
\begin{enumerate}
  \item First step
  \item Second step
  \item Third step
\end{enumerate}

\section{Citations and References}
When you use the citation search feature, references will be automatically managed. For example, you might cite a paper like this: \cite{example2024}. All citations will be automatically added to your \texttt{cite.bib} file.

\section{Conclusion}
This template demonstrates basic \LaTeX{} features including mathematical formulas, tables, and lists. The bibliography system is now set up to work seamlessly with Bibby AI's citation management. Feel free to modify it for your needs!

\bibliographystyle{plain}
\bibliography{cite}

\end{document}\documentclass{article}
\usepackage[utf8]{inputenc}
\usepackage{booktabs}
\usepackage{amsmath}
\usepackage{natbib}

\title{My Document}
\author{Your Name}
\date{\today}

\begin{document}

\maketitle

\section{Introduction}
This is a sample document demonstrating various \LaTeX{} features including mathematical equations, tables, and formatting. You can cite papers using the citation search feature, and they will be automatically added to your \texttt{cite.bib} file.

\section{Mathematical Expressions}
Here are some example equations:

\subsection{Quadratic Formula}
The quadratic formula is given by:
\[x = \frac{-b \pm \sqrt{b^2 - 4ac}}{2a}\]

\subsection{Matrix Example}
A sample matrix:
\[
\begin{pmatrix}
a_{11} & a_{12} & a_{13} \\
a_{21} & a_{22} & a_{23} \\
a_{31} & a_{32} & a_{33}
\end{pmatrix}
\]

\section{Tables}
Below is a sample table showing data:

\begin{table}[h]
\centering
\begin{tabular}{@{}lrrr@{}}
\toprule
\textbf{Category} & \textbf{Value 1} & \textbf{Value 2} & \textbf{Total} \\
\midrule
Item A & 10.5 & 20.3 & 30.8 \\
Item B & 15.2 & 25.1 & 40.3 \\
Item C & 22.1 & 12.7 & 34.8 \\
\bottomrule
\end{tabular}
\caption{Sample Data Table}
\label{tab:sample}
\end{table}

\section{Lists}
\subsection{Bullet Points}
\begin{itemize}
  \item First item
  \item Second item
  \item Third item with sub-items:
    \begin{itemize}
      \item Sub-item 1
      \item Sub-item 2
    \end{itemize}
\end{itemize}

\subsection{Numbered List}
\begin{enumerate}
  \item First step
  \item Second step
  \item Third step
\end{enumerate}

\section{Citations and References}
When you use the citation search feature, references will be automatically managed. For example, you might cite a paper like this: \cite{example2024}. All citations will be automatically added to your \texttt{cite.bib} file.

\section{Conclusion}
This template demonstrates basic \LaTeX{} features including mathematical formulas, tables, and lists. The bibliography system is now set up to work seamlessly with Bibby AI's citation management. Feel free to modify it for your needs!

\bibliographystyle{plain}
\bibliography{cite}

\end{document}\documentclass{article}
\usepackage[utf8]{inputenc}
\usepackage{booktabs}
\usepackage{amsmath}
\usepackage{natbib}

\title{My Document}
\author{Your Name}
\date{\today}

\begin{document}

\maketitle

\section{Introduction}
This is a sample document demonstrating various \LaTeX{} features including mathematical equations, tables, and formatting. You can cite papers using the citation search feature, and they will be automatically added to your \texttt{cite.bib} file.

\section{Mathematical Expressions}
Here are some example equations:

\subsection{Quadratic Formula}
The quadratic formula is given by:
\[x = \frac{-b \pm \sqrt{b^2 - 4ac}}{2a}\]

\subsection{Matrix Example}
A sample matrix:
\[
\begin{pmatrix}
a_{11} & a_{12} & a_{13} \\
a_{21} & a_{22} & a_{23} \\
a_{31} & a_{32} & a_{33}
\end{pmatrix}
\]

\section{Tables}
Below is a sample table showing data:

\begin{table}[h]
\centering
\begin{tabular}{@{}lrrr@{}}
\toprule
\textbf{Category} & \textbf{Value 1} & \textbf{Value 2} & \textbf{Total} \\
\midrule
Item A & 10.5 & 20.3 & 30.8 \\
Item B & 15.2 & 25.1 & 40.3 \\
Item C & 22.1 & 12.7 & 34.8 \\
\bottomrule
\end{tabular}
\caption{Sample Data Table}
\label{tab:sample}
\end{table}

\section{Lists}
\subsection{Bullet Points}
\begin{itemize}
  \item First item
  \item Second item
  \item Third item with sub-items:
    \begin{itemize}
      \item Sub-item 1
      \item Sub-item 2
    \end{itemize}
\end{itemize}

\subsection{Numbered List}
\begin{enumerate}
  \item First step
  \item Second step
  \item Third step
\end{enumerate}

\section{Citations and References}
When you use the citation search feature, references will be automatically managed. For example, you might cite a paper like this: \cite{example2024}. All citations will be automatically added to your \texttt{cite.bib} file.

\section{Conclusion}
This template demonstrates basic \LaTeX{} features including mathematical formulas, tables, and lists. The bibliography system is now set up to work seamlessly with Bibby AI's citation management. Feel free to modify it for your needs!

\bibliographystyle{plain}
\bibliography{cite}

\end{document}\documentclass{article}
\usepackage[utf8]{inputenc}
\usepackage{booktabs}
\usepackage{amsmath}
\usepackage{natbib}

\title{My Document}
\author{Your Name}
\date{\today}

\begin{document}

\maketitle

\section{Introduction}
This is a sample document demonstrating various \LaTeX{} features including mathematical equations, tables, and formatting. You can cite papers using the citation search feature, and they will be automatically added to your \texttt{cite.bib} file.

\section{Mathematical Expressions}
Here are some example equations:

\subsection{Quadratic Formula}
The quadratic formula is given by:
\[x = \frac{-b \pm \sqrt{b^2 - 4ac}}{2a}\]

\subsection{Matrix Example}
A sample matrix:
\[
\begin{pmatrix}
a_{11} & a_{12} & a_{13} \\
a_{21} & a_{22} & a_{23} \\
a_{31} & a_{32} & a_{33}
\end{pmatrix}
\]

\section{Tables}
Below is a sample table showing data:

\begin{table}[h]
\centering
\begin{tabular}{@{}lrrr@{}}
\toprule
\textbf{Category} & \textbf{Value 1} & \textbf{Value 2} & \textbf{Total} \\
\midrule
Item A & 10.5 & 20.3 & 30.8 \\
Item B & 15.2 & 25.1 & 40.3 \\
Item C & 22.1 & 12.7 & 34.8 \\
\bottomrule
\end{tabular}
\caption{Sample Data Table}
\label{tab:sample}
\end{table}

\section{Lists}
\subsection{Bullet Points}
\begin{itemize}
  \item First item
  \item Second item
  \item Third item with sub-items:
    \begin{itemize}
      \item Sub-item 1
      \item Sub-item 2
    \end{itemize}
\end{itemize}

\subsection{Numbered List}
\begin{enumerate}
  \item First step
  \item Second step
  \item Third step
\end{enumerate}

\section{Citations and References}
When you use the citation search feature, references will be automatically managed. For example, you might cite a paper like this: \cite{example2024}. All citations will be automatically added to your \texttt{cite.bib} file.

\section{Conclusion}
This template demonstrates basic \LaTeX{} features including mathematical formulas, tables, and lists. The bibliography system is now set up to work seamlessly with Bibby AI's citation management. Feel free to modify it for your needs!

\bibliographystyle{plain}
\bibliography{cite}

\end{document}\documentclass{article}
\usepackage[utf8]{inputenc}
\usepackage{booktabs}
\usepackage{amsmath}
\usepackage{natbib}

\title{My Document}
\author{Your Name}
\date{\today}

\begin{document}

\maketitle

\section{Introduction}
This is a sample document demonstrating various \LaTeX{} features including mathematical equations, tables, and formatting. You can cite papers using the citation search feature, and they will be automatically added to your \texttt{cite.bib} file.

\section{Mathematical Expressions}
Here are some example equations:

\subsection{Quadratic Formula}
The quadratic formula is given by:
\[x = \frac{-b \pm \sqrt{b^2 - 4ac}}{2a}\]

\subsection{Matrix Example}
A sample matrix:
\[
\begin{pmatrix}
a_{11} & a_{12} & a_{13} \\
a_{21} & a_{22} & a_{23} \\
a_{31} & a_{32} & a_{33}
\end{pmatrix}
\]

\section{Tables}
Below is a sample table showing data:

\begin{table}[h]
\centering
\begin{tabular}{@{}lrrr@{}}
\toprule
\textbf{Category} & \textbf{Value 1} & \textbf{Value 2} & \textbf{Total} \\
\midrule
Item A & 10.5 & 20.3 & 30.8 \\
Item B & 15.2 & 25.1 & 40.3 \\
Item C & 22.1 & 12.7 & 34.8 \\
\bottomrule
\end{tabular}
\caption{Sample Data Table}
\label{tab:sample}
\end{table}

\section{Lists}
\subsection{Bullet Points}
\begin{itemize}
  \item First item
  \item Second item
  \item Third item with sub-items:
    \begin{itemize}
      \item Sub-item 1
      \item Sub-item 2
    \end{itemize}
\end{itemize}

\subsection{Numbered List}
\begin{enumerate}
  \item First step
  \item Second step
  \item Third step
\end{enumerate}

\section{Citations and References}
When you use the citation search feature, references will be automatically managed. For example, you might cite a paper like this: \cite{example2024}. All citations will be automatically added to your \texttt{cite.bib} file.

\section{Conclusion}
This template demonstrates basic \LaTeX{} features including mathematical formulas, tables, and lists. The bibliography system is now set up to work seamlessly with Bibby AI's citation management. Feel free to modify it for your needs!

\bibliographystyle{plain}
\bibliography{cite}

\end{document}\documentclass{article}
\usepackage[utf8]{inputenc}
\usepackage{booktabs}
\usepackage{amsmath}
\usepackage{natbib}

\title{My Document}
\author{Your Name}
\date{\today}

\begin{document}

\maketitle

\section{Introduction}
This is a sample document demonstrating various \LaTeX{} features including mathematical equations, tables, and formatting. You can cite papers using the citation search feature, and they will be automatically added to your \texttt{cite.bib} file.

\section{Mathematical Expressions}
Here are some example equations:

\subsection{Quadratic Formula}
The quadratic formula is given by:
\[x = \frac{-b \pm \sqrt{b^2 - 4ac}}{2a}\]

\subsection{Matrix Example}
A sample matrix:
\[
\begin{pmatrix}
a_{11} & a_{12} & a_{13} \\
a_{21} & a_{22} & a_{23} \\
a_{31} & a_{32} & a_{33}
\end{pmatrix}
\]

\section{Tables}
Below is a sample table showing data:

\begin{table}[h]
\centering
\begin{tabular}{@{}lrrr@{}}
\toprule
\textbf{Category} & \textbf{Value 1} & \textbf{Value 2} & \textbf{Total} \\
\midrule
Item A & 10.5 & 20.3 & 30.8 \\
Item B & 15.2 & 25.1 & 40.3 \\
Item C & 22.1 & 12.7 & 34.8 \\
\bottomrule
\end{tabular}
\caption{Sample Data Table}
\label{tab:sample}
\end{table}

\section{Lists}
\subsection{Bullet Points}
\begin{itemize}
  \item First item
  \item Second item
  \item Third item with sub-items:
    \begin{itemize}
      \item Sub-item 1
      \item Sub-item 2
    \end{itemize}
\end{itemize}

\subsection{Numbered List}
\begin{enumerate}
  \item First step
  \item Second step
  \item Third step
\end{enumerate}

\section{Citations and References}
When you use the citation search feature, references will be automatically managed. For example, you might cite a paper like this: \cite{example2024}. All citations will be automatically added to your \texttt{cite.bib} file.

\section{Conclusion}
This template demonstrates basic \LaTeX{} features including mathematical formulas, tables, and lists. The bibliography system is now set up to work seamlessly with Bibby AI's citation management. Feel free to modify it for your needs!

\bibliographystyle{plain}
\bibliography{cite}

\end{document}\documentclass{article}
\usepackage[utf8]{inputenc}
\usepackage{booktabs}
\usepackage{amsmath}
\usepackage{natbib}

\title{My Document}
\author{Your Name}
\date{\today}

\begin{document}

\maketitle

\section{Introduction}
This is a sample document demonstrating various \LaTeX{} features including mathematical equations, tables, and formatting. You can cite papers using the citation search feature, and they will be automatically added to your \texttt{cite.bib} file.

\section{Mathematical Expressions}
Here are some example equations:

\subsection{Quadratic Formula}
The quadratic formula is given by:
\[x = \frac{-b \pm \sqrt{b^2 - 4ac}}{2a}\]

\subsection{Matrix Example}
A sample matrix:
\[
\begin{pmatrix}
a_{11} & a_{12} & a_{13} \\
a_{21} & a_{22} & a_{23} \\
a_{31} & a_{32} & a_{33}
\end{pmatrix}
\]

\section{Tables}
Below is a sample table showing data:

\begin{table}[h]
\centering
\begin{tabular}{@{}lrrr@{}}
\toprule
\textbf{Category} & \textbf{Value 1} & \textbf{Value 2} & \textbf{Total} \\
\midrule
Item A & 10.5 & 20.3 & 30.8 \\
Item B & 15.2 & 25.1 & 40.3 \\
Item C & 22.1 & 12.7 & 34.8 \\
\bottomrule
\end{tabular}
\caption{Sample Data Table}
\label{tab:sample}
\end{table}

\section{Lists}
\subsection{Bullet Points}
\begin{itemize}
  \item First item
  \item Second item
  \item Third item with sub-items:
    \begin{itemize}
      \item Sub-item 1
      \item Sub-item 2
    \end{itemize}
\end{itemize}

\subsection{Numbered List}
\begin{enumerate}
  \item First step
  \item Second step
  \item Third step
\end{enumerate}

\section{Citations and References}
When you use the citation search feature, references will be automatically managed. For example, you might cite a paper like this: \cite{example2024}. All citations will be automatically added to your \texttt{cite.bib} file.

\section{Conclusion}
This template demonstrates basic \LaTeX{} features including mathematical formulas, tables, and lists. The bibliography system is now set up to work seamlessly with Bibby AI's citation management. Feel free to modify it for your needs!

\bibliographystyle{plain}
\bibliography{cite}

\end{document}\documentclass{article}
\usepackage[utf8]{inputenc}
\usepackage{booktabs}
\usepackage{amsmath}
\usepackage{natbib}

\title{My Document}
\author{Your Name}
\date{\today}

\begin{document}

\maketitle

\section{Introduction}
This is a sample document demonstrating various \LaTeX{} features including mathematical equations, tables, and formatting. You can cite papers using the citation search feature, and they will be automatically added to your \texttt{cite.bib} file.

\section{Mathematical Expressions}
Here are some example equations:

\subsection{Quadratic Formula}
The quadratic formula is given by:
\[x = \frac{-b \pm \sqrt{b^2 - 4ac}}{2a}\]

\subsection{Matrix Example}
A sample matrix:
\[
\begin{pmatrix}
a_{11} & a_{12} & a_{13} \\
a_{21} & a_{22} & a_{23} \\
a_{31} & a_{32} & a_{33}
\end{pmatrix}
\]

\section{Tables}
Below is a sample table showing data:

\begin{table}[h]
\centering
\begin{tabular}{@{}lrrr@{}}
\toprule
\textbf{Category} & \textbf{Value 1} & \textbf{Value 2} & \textbf{Total} \\
\midrule
Item A & 10.5 & 20.3 & 30.8 \\
Item B & 15.2 & 25.1 & 40.3 \\
Item C & 22.1 & 12.7 & 34.8 \\
\bottomrule
\end{tabular}
\caption{Sample Data Table}
\label{tab:sample}
\end{table}

\section{Lists}
\subsection{Bullet Points}
\begin{itemize}
  \item First item
  \item Second item
  \item Third item with sub-items:
    \begin{itemize}
      \item Sub-item 1
      \item Sub-item 2
    \end{itemize}
\end{itemize}

\subsection{Numbered List}
\begin{enumerate}
  \item First step
  \item Second step
  \item Third step
\end{enumerate}

\section{Citations and References}
When you use the citation search feature, references will be automatically managed. For example, you might cite a paper like this: \cite{example2024}. All citations will be automatically added to your \texttt{cite.bib} file.

\section{Conclusion}
This template demonstrates basic \LaTeX{} features including mathematical formulas, tables, and lists. The bibliography system is now set up to work seamlessly with Bibby AI's citation management. Feel free to modify it for your needs!

\bibliographystyle{plain}
\bibliography{cite}

\end{document}\documentclass{article}
\usepackage[utf8]{inputenc}
\usepackage{booktabs}
\usepackage{amsmath}
\usepackage{natbib}

\title{My Document}
\author{Your Name}
\date{\today}

\begin{document}

\maketitle

\section{Introduction}
This is a sample document demonstrating various \LaTeX{} features including mathematical equations, tables, and formatting. You can cite papers using the citation search feature, and they will be automatically added to your \texttt{cite.bib} file.

\section{Mathematical Expressions}
Here are some example equations:

\subsection{Quadratic Formula}
The quadratic formula is given by:
\[x = \frac{-b \pm \sqrt{b^2 - 4ac}}{2a}\]

\subsection{Matrix Example}
A sample matrix:
\[
\begin{pmatrix}
a_{11} & a_{12} & a_{13} \\
a_{21} & a_{22} & a_{23} \\
a_{31} & a_{32} & a_{33}
\end{pmatrix}
\]

\section{Tables}
Below is a sample table showing data:

\begin{table}[h]
\centering
\begin{tabular}{@{}lrrr@{}}
\toprule
\textbf{Category} & \textbf{Value 1} & \textbf{Value 2} & \textbf{Total} \\
\midrule
Item A & 10.5 & 20.3 & 30.8 \\
Item B & 15.2 & 25.1 & 40.3 \\
Item C & 22.1 & 12.7 & 34.8 \\
\bottomrule
\end{tabular}
\caption{Sample Data Table}
\label{tab:sample}
\end{table}

\section{Lists}
\subsection{Bullet Points}
\begin{itemize}
  \item First item
  \item Second item
  \item Third item with sub-items:
    \begin{itemize}
      \item Sub-item 1
      \item Sub-item 2
    \end{itemize}
\end{itemize}

\subsection{Numbered List}
\begin{enumerate}
  \item First step
  \item Second step
  \item Third step
\end{enumerate}

\section{Citations and References}
When you use the citation search feature, references will be automatically managed. For example, you might cite a paper like this: \cite{example2024}. All citations will be automatically added to your \texttt{cite.bib} file.

\section{Conclusion}
This template demonstrates basic \LaTeX{} features including mathematical formulas, tables, and lists. The bibliography system is now set up to work seamlessly with Bibby AI's citation management. Feel free to modify it for your needs!

\bibliographystyle{plain}
\bibliography{cite}

\end{document}\documentclass{article}
\usepackage[utf8]{inputenc}
\usepackage{booktabs}
\usepackage{amsmath}
\usepackage{natbib}

\title{My Document}
\author{Your Name}
\date{\today}

\begin{document}

\maketitle

\section{Introduction}
This is a sample document demonstrating various \LaTeX{} features including mathematical equations, tables, and formatting. You can cite papers using the citation search feature, and they will be automatically added to your \texttt{cite.bib} file.

\section{Mathematical Expressions}
Here are some example equations:

\subsection{Quadratic Formula}
The quadratic formula is given by:
\[x = \frac{-b \pm \sqrt{b^2 - 4ac}}{2a}\]

\subsection{Matrix Example}
A sample matrix:
\[
\begin{pmatrix}
a_{11} & a_{12} & a_{13} \\
a_{21} & a_{22} & a_{23} \\
a_{31} & a_{32} & a_{33}
\end{pmatrix}
\]

\section{Tables}
Below is a sample table showing data:

\begin{table}[h]
\centering
\begin{tabular}{@{}lrrr@{}}
\toprule
\textbf{Category} & \textbf{Value 1} & \textbf{Value 2} & \textbf{Total} \\
\midrule
Item A & 10.5 & 20.3 & 30.8 \\
Item B & 15.2 & 25.1 & 40.3 \\
Item C & 22.1 & 12.7 & 34.8 \\
\bottomrule
\end{tabular}
\caption{Sample Data Table}
\label{tab:sample}
\end{table}

\section{Lists}
\subsection{Bullet Points}
\begin{itemize}
  \item First item
  \item Second item
  \item Third item with sub-items:
    \begin{itemize}
      \item Sub-item 1
      \item Sub-item 2
    \end{itemize}
\end{itemize}

\subsection{Numbered List}
\begin{enumerate}
  \item First step
  \item Second step
  \item Third step
\end{enumerate}

\section{Citations and References}
When you use the citation search feature, references will be automatically managed. For example, you might cite a paper like this: \cite{example2024}. All citations will be automatically added to your \texttt{cite.bib} file.

\section{Conclusion}
This template demonstrates basic \LaTeX{} features including mathematical formulas, tables, and lists. The bibliography system is now set up to work seamlessly with Bibby AI's citation management. Feel free to modify it for your needs!

\bibliographystyle{plain}
\bibliography{cite}

\end{document}\documentclass{article}
\usepackage[utf8]{inputenc}
\usepackage{booktabs}
\usepackage{amsmath}
\usepackage{natbib}

\title{My Document}
\author{Your Name}
\date{\today}

\begin{document}

\maketitle

\section{Introduction}
This is a sample document demonstrating various \LaTeX{} features including mathematical equations, tables, and formatting. You can cite papers using the citation search feature, and they will be automatically added to your \texttt{cite.bib} file.

\section{Mathematical Expressions}
Here are some example equations:

\subsection{Quadratic Formula}
The quadratic formula is given by:
\[x = \frac{-b \pm \sqrt{b^2 - 4ac}}{2a}\]

\subsection{Matrix Example}
A sample matrix:
\[
\begin{pmatrix}
a_{11} & a_{12} & a_{13} \\
a_{21} & a_{22} & a_{23} \\
a_{31} & a_{32} & a_{33}
\end{pmatrix}
\]

\section{Tables}
Below is a sample table showing data:

\begin{table}[h]
\centering
\begin{tabular}{@{}lrrr@{}}
\toprule
\textbf{Category} & \textbf{Value 1} & \textbf{Value 2} & \textbf{Total} \\
\midrule
Item A & 10.5 & 20.3 & 30.8 \\
Item B & 15.2 & 25.1 & 40.3 \\
Item C & 22.1 & 12.7 & 34.8 \\
\bottomrule
\end{tabular}
\caption{Sample Data Table}
\label{tab:sample}
\end{table}

\section{Lists}
\subsection{Bullet Points}
\begin{itemize}
  \item First item
  \item Second item
  \item Third item with sub-items:
    \begin{itemize}
      \item Sub-item 1
      \item Sub-item 2
    \end{itemize}
\end{itemize}

\subsection{Numbered List}
\begin{enumerate}
  \item First step
  \item Second step
  \item Third step
\end{enumerate}

\section{Citations and References}
When you use the citation search feature, references will be automatically managed. For example, you might cite a paper like this: \cite{example2024}. All citations will be automatically added to your \texttt{cite.bib} file.

\section{Conclusion}
This template demonstrates basic \LaTeX{} features including mathematical formulas, tables, and lists. The bibliography system is now set up to work seamlessly with Bibby AI's citation management. Feel free to modify it for your needs!

\bibliographystyle{plain}
\bibliography{cite}

\end{document}\documentclass{article}
\usepackage[utf8]{inputenc}
\usepackage{booktabs}
\usepackage{amsmath}
\usepackage{natbib}

\title{My Document}
\author{Your Name}
\date{\today}

\begin{document}

\maketitle

\section{Introduction}
This is a sample document demonstrating various \LaTeX{} features including mathematical equations, tables, and formatting. You can cite papers using the citation search feature, and they will be automatically added to your \texttt{cite.bib} file.

\section{Mathematical Expressions}
Here are some example equations:

\subsection{Quadratic Formula}
The quadratic formula is given by:
\[x = \frac{-b \pm \sqrt{b^2 - 4ac}}{2a}\]

\subsection{Matrix Example}
A sample matrix:
\[
\begin{pmatrix}
a_{11} & a_{12} & a_{13} \\
a_{21} & a_{22} & a_{23} \\
a_{31} & a_{32} & a_{33}
\end{pmatrix}
\]

\section{Tables}
Below is a sample table showing data:

\begin{table}[h]
\centering
\begin{tabular}{@{}lrrr@{}}
\toprule
\textbf{Category} & \textbf{Value 1} & \textbf{Value 2} & \textbf{Total} \\
\midrule
Item A & 10.5 & 20.3 & 30.8 \\
Item B & 15.2 & 25.1 & 40.3 \\
Item C & 22.1 & 12.7 & 34.8 \\
\bottomrule
\end{tabular}
\caption{Sample Data Table}
\label{tab:sample}
\end{table}

\section{Lists}
\subsection{Bullet Points}
\begin{itemize}
  \item First item
  \item Second item
  \item Third item with sub-items:
    \begin{itemize}
      \item Sub-item 1
      \item Sub-item 2
    \end{itemize}
\end{itemize}

\subsection{Numbered List}
\begin{enumerate}
  \item First step
  \item Second step
  \item Third step
\end{enumerate}

\section{Citations and References}
When you use the citation search feature, references will be automatically managed. For example, you might cite a paper like this: \cite{example2024}. All citations will be automatically added to your \texttt{cite.bib} file.

\section{Conclusion}
This template demonstrates basic \LaTeX{} features including mathematical formulas, tables, and lists. The bibliography system is now set up to work seamlessly with Bibby AI's citation management. Feel free to modify it for your needs!

\bibliographystyle{plain}
\bibliography{cite}

\end{document}\documentclass{article}
\usepackage[utf8]{inputenc}
\usepackage{booktabs}
\usepackage{amsmath}
\usepackage{natbib}

\title{My Document}
\author{Your Name}
\date{\today}

\begin{document}

\maketitle

\section{Introduction}
This is a sample document demonstrating various \LaTeX{} features including mathematical equations, tables, and formatting. You can cite papers using the citation search feature, and they will be automatically added to your \texttt{cite.bib} file.

\section{Mathematical Expressions}
Here are some example equations:

\subsection{Quadratic Formula}
The quadratic formula is given by:
\[x = \frac{-b \pm \sqrt{b^2 - 4ac}}{2a}\]

\subsection{Matrix Example}
A sample matrix:
\[
\begin{pmatrix}
a_{11} & a_{12} & a_{13} \\
a_{21} & a_{22} & a_{23} \\
a_{31} & a_{32} & a_{33}
\end{pmatrix}
\]

\section{Tables}
Below is a sample table showing data:

\begin{table}[h]
\centering
\begin{tabular}{@{}lrrr@{}}
\toprule
\textbf{Category} & \textbf{Value 1} & \textbf{Value 2} & \textbf{Total} \\
\midrule
Item A & 10.5 & 20.3 & 30.8 \\
Item B & 15.2 & 25.1 & 40.3 \\
Item C & 22.1 & 12.7 & 34.8 \\
\bottomrule
\end{tabular}
\caption{Sample Data Table}
\label{tab:sample}
\end{table}

\section{Lists}
\subsection{Bullet Points}
\begin{itemize}
  \item First item
  \item Second item
  \item Third item with sub-items:
    \begin{itemize}
      \item Sub-item 1
      \item Sub-item 2
    \end{itemize}
\end{itemize}

\subsection{Numbered List}
\begin{enumerate}
  \item First step
  \item Second step
  \item Third step
\end{enumerate}

\section{Citations and References}
When you use the citation search feature, references will be automatically managed. For example, you might cite a paper like this: \cite{example2024}. All citations will be automatically added to your \texttt{cite.bib} file.

\section{Conclusion}
This template demonstrates basic \LaTeX{} features including mathematical formulas, tables, and lists. The bibliography system is now set up to work seamlessly with Bibby AI's citation management. Feel free to modify it for your needs!

\bibliographystyle{plain}
\bibliography{cite}

\end{document}\documentclass{article}
\usepackage[utf8]{inputenc}
\usepackage{booktabs}
\usepackage{amsmath}
\usepackage{natbib}

\title{My Document}
\author{Your Name}
\date{\today}

\begin{document}

\maketitle

\section{Introduction}
This is a sample document demonstrating various \LaTeX{} features including mathematical equations, tables, and formatting. You can cite papers using the citation search feature, and they will be automatically added to your \texttt{cite.bib} file.

\section{Mathematical Expressions}
Here are some example equations:

\subsection{Quadratic Formula}
The quadratic formula is given by:
\[x = \frac{-b \pm \sqrt{b^2 - 4ac}}{2a}\]

\subsection{Matrix Example}
A sample matrix:
\[
\begin{pmatrix}
a_{11} & a_{12} & a_{13} \\
a_{21} & a_{22} & a_{23} \\
a_{31} & a_{32} & a_{33}
\end{pmatrix}
\]

\section{Tables}
Below is a sample table showing data:

\begin{table}[h]
\centering
\begin{tabular}{@{}lrrr@{}}
\toprule
\textbf{Category} & \textbf{Value 1} & \textbf{Value 2} & \textbf{Total} \\
\midrule
Item A & 10.5 & 20.3 & 30.8 \\
Item B & 15.2 & 25.1 & 40.3 \\
Item C & 22.1 & 12.7 & 34.8 \\
\bottomrule
\end{tabular}
\caption{Sample Data Table}
\label{tab:sample}
\end{table}

\section{Lists}
\subsection{Bullet Points}
\begin{itemize}
  \item First item
  \item Second item
  \item Third item with sub-items:
    \begin{itemize}
      \item Sub-item 1
      \item Sub-item 2
    \end{itemize}
\end{itemize}

\subsection{Numbered List}
\begin{enumerate}
  \item First step
  \item Second step
  \item Third step
\end{enumerate}

\section{Citations and References}
When you use the citation search feature, references will be automatically managed. For example, you might cite a paper like this: \cite{example2024}. All citations will be automatically added to your \texttt{cite.bib} file.

\section{Conclusion}
This template demonstrates basic \LaTeX{} features including mathematical formulas, tables, and lists. The bibliography system is now set up to work seamlessly with Bibby AI's citation management. Feel free to modify it for your needs!

\bibliographystyle{plain}
\bibliography{cite}

\end{document}\documentclass{article}
\usepackage[utf8]{inputenc}
\usepackage{booktabs}
\usepackage{amsmath}
\usepackage{natbib}

\title{My Document}
\author{Your Name}
\date{\today}

\begin{document}

\maketitle

\section{Introduction}
This is a sample document demonstrating various \LaTeX{} features including mathematical equations, tables, and formatting. You can cite papers using the citation search feature, and they will be automatically added to your \texttt{cite.bib} file.

\section{Mathematical Expressions}
Here are some example equations:

\subsection{Quadratic Formula}
The quadratic formula is given by:
\[x = \frac{-b \pm \sqrt{b^2 - 4ac}}{2a}\]

\subsection{Matrix Example}
A sample matrix:
\[
\begin{pmatrix}
a_{11} & a_{12} & a_{13} \\
a_{21} & a_{22} & a_{23} \\
a_{31} & a_{32} & a_{33}
\end{pmatrix}
\]

\section{Tables}
Below is a sample table showing data:

\begin{table}[h]
\centering
\begin{tabular}{@{}lrrr@{}}
\toprule
\textbf{Category} & \textbf{Value 1} & \textbf{Value 2} & \textbf{Total} \\
\midrule
Item A & 10.5 & 20.3 & 30.8 \\
Item B & 15.2 & 25.1 & 40.3 \\
Item C & 22.1 & 12.7 & 34.8 \\
\bottomrule
\end{tabular}
\caption{Sample Data Table}
\label{tab:sample}
\end{table}

\section{Lists}
\subsection{Bullet Points}
\begin{itemize}
  \item First item
  \item Second item
  \item Third item with sub-items:
    \begin{itemize}
      \item Sub-item 1
      \item Sub-item 2
    \end{itemize}
\end{itemize}

\subsection{Numbered List}
\begin{enumerate}
  \item First step
  \item Second step
  \item Third step
\end{enumerate}

\section{Citations and References}
When you use the citation search feature, references will be automatically managed. For example, you might cite a paper like this: \cite{example2024}. All citations will be automatically added to your \texttt{cite.bib} file.

\section{Conclusion}
This template demonstrates basic \LaTeX{} features including mathematical formulas, tables, and lists. The bibliography system is now set up to work seamlessly with Bibby AI's citation management. Feel free to modify it for your needs!

\bibliographystyle{plain}
\bibliography{cite}

\end{document}\documentclass{article}
\usepackage[utf8]{inputenc}
\usepackage{booktabs}
\usepackage{amsmath}
\usepackage{natbib}

\title{My Document}
\author{Your Name}
\date{\today}

\begin{document}

\maketitle

\section{Introduction}
This is a sample document demonstrating various \LaTeX{} features including mathematical equations, tables, and formatting. You can cite papers using the citation search feature, and they will be automatically added to your \texttt{cite.bib} file.

\section{Mathematical Expressions}
Here are some example equations:

\subsection{Quadratic Formula}
The quadratic formula is given by:
\[x = \frac{-b \pm \sqrt{b^2 - 4ac}}{2a}\]

\subsection{Matrix Example}
A sample matrix:
\[
\begin{pmatrix}
a_{11} & a_{12} & a_{13} \\
a_{21} & a_{22} & a_{23} \\
a_{31} & a_{32} & a_{33}
\end{pmatrix}
\]

\section{Tables}
Below is a sample table showing data:

\begin{table}[h]
\centering
\begin{tabular}{@{}lrrr@{}}
\toprule
\textbf{Category} & \textbf{Value 1} & \textbf{Value 2} & \textbf{Total} \\
\midrule
Item A & 10.5 & 20.3 & 30.8 \\
Item B & 15.2 & 25.1 & 40.3 \\
Item C & 22.1 & 12.7 & 34.8 \\
\bottomrule
\end{tabular}
\caption{Sample Data Table}
\label{tab:sample}
\end{table}

\section{Lists}
\subsection{Bullet Points}
\begin{itemize}
  \item First item
  \item Second item
  \item Third item with sub-items:
    \begin{itemize}
      \item Sub-item 1
      \item Sub-item 2
    \end{itemize}
\end{itemize}

\subsection{Numbered List}
\begin{enumerate}
  \item First step
  \item Second step
  \item Third step
\end{enumerate}

\section{Citations and References}
When you use the citation search feature, references will be automatically managed. For example, you might cite a paper like this: \cite{example2024}. All citations will be automatically added to your \texttt{cite.bib} file.

\section{Conclusion}
This template demonstrates basic \LaTeX{} features including mathematical formulas, tables, and lists. The bibliography system is now set up to work seamlessly with Bibby AI's citation management. Feel free to modify it for your needs!

\bibliographystyle{plain}
\bibliography{cite}

\end{document}\documentclass{article}
\usepackage[utf8]{inputenc}
\usepackage{booktabs}
\usepackage{amsmath}
\usepackage{natbib}

\title{My Document}
\author{Your Name}
\date{\today}

\begin{document}

\maketitle

\section{Introduction}
This is a sample document demonstrating various \LaTeX{} features including mathematical equations, tables, and formatting. You can cite papers using the citation search feature, and they will be automatically added to your \texttt{cite.bib} file.

\section{Mathematical Expressions}
Here are some example equations:

\subsection{Quadratic Formula}
The quadratic formula is given by:
\[x = \frac{-b \pm \sqrt{b^2 - 4ac}}{2a}\]

\subsection{Matrix Example}
A sample matrix:
\[
\begin{pmatrix}
a_{11} & a_{12} & a_{13} \\
a_{21} & a_{22} & a_{23} \\
a_{31} & a_{32} & a_{33}
\end{pmatrix}
\]

\section{Tables}
Below is a sample table showing data:

\begin{table}[h]
\centering
\begin{tabular}{@{}lrrr@{}}
\toprule
\textbf{Category} & \textbf{Value 1} & \textbf{Value 2} & \textbf{Total} \\
\midrule
Item A & 10.5 & 20.3 & 30.8 \\
Item B & 15.2 & 25.1 & 40.3 \\
Item C & 22.1 & 12.7 & 34.8 \\
\bottomrule
\end{tabular}
\caption{Sample Data Table}
\label{tab:sample}
\end{table}

\section{Lists}
\subsection{Bullet Points}
\begin{itemize}
  \item First item
  \item Second item
  \item Third item with sub-items:
    \begin{itemize}
      \item Sub-item 1
      \item Sub-item 2
    \end{itemize}
\end{itemize}

\subsection{Numbered List}
\begin{enumerate}
  \item First step
  \item Second step
  \item Third step
\end{enumerate}

\section{Citations and References}
When you use the citation search feature, references will be automatically managed. For example, you might cite a paper like this: \cite{example2024}. All citations will be automatically added to your \texttt{cite.bib} file.

\section{Conclusion}
This template demonstrates basic \LaTeX{} features including mathematical formulas, tables, and lists. The bibliography system is now set up to work seamlessly with Bibby AI's citation management. Feel free to modify it for your needs!

\bibliographystyle{plain}
\bibliography{cite}

\end{document}\documentclass{article}
\usepackage[utf8]{inputenc}
\usepackage{booktabs}
\usepackage{amsmath}
\usepackage{natbib}

\title{My Document}
\author{Your Name}
\date{\today}

\begin{document}

\maketitle

\section{Introduction}
This is a sample document demonstrating various \LaTeX{} features including mathematical equations, tables, and formatting. You can cite papers using the citation search feature, and they will be automatically added to your \texttt{cite.bib} file.

\section{Mathematical Expressions}
Here are some example equations:

\subsection{Quadratic Formula}
The quadratic formula is given by:
\[x = \frac{-b \pm \sqrt{b^2 - 4ac}}{2a}\]

\subsection{Matrix Example}
A sample matrix:
\[
\begin{pmatrix}
a_{11} & a_{12} & a_{13} \\
a_{21} & a_{22} & a_{23} \\
a_{31} & a_{32} & a_{33}
\end{pmatrix}
\]

\section{Tables}
Below is a sample table showing data:

\begin{table}[h]
\centering
\begin{tabular}{@{}lrrr@{}}
\toprule
\textbf{Category} & \textbf{Value 1} & \textbf{Value 2} & \textbf{Total} \\
\midrule
Item A & 10.5 & 20.3 & 30.8 \\
Item B & 15.2 & 25.1 & 40.3 \\
Item C & 22.1 & 12.7 & 34.8 \\
\bottomrule
\end{tabular}
\caption{Sample Data Table}
\label{tab:sample}
\end{table}

\section{Lists}
\subsection{Bullet Points}
\begin{itemize}
  \item First item
  \item Second item
  \item Third item with sub-items:
    \begin{itemize}
      \item Sub-item 1
      \item Sub-item 2
    \end{itemize}
\end{itemize}

\subsection{Numbered List}
\begin{enumerate}
  \item First step
  \item Second step
  \item Third step
\end{enumerate}

\section{Citations and References}
When you use the citation search feature, references will be automatically managed. For example, you might cite a paper like this: \cite{example2024}. All citations will be automatically added to your \texttt{cite.bib} file.

\section{Conclusion}
This template demonstrates basic \LaTeX{} features including mathematical formulas, tables, and lists. The bibliography system is now set up to work seamlessly with Bibby AI's citation management. Feel free to modify it for your needs!

\bibliographystyle{plain}
\bibliography{cite}

\end{document}\documentclass{article}
\usepackage[utf8]{inputenc}
\usepackage{booktabs}
\usepackage{amsmath}
\usepackage{natbib}

\title{My Document}
\author{Your Name}
\date{\today}

\begin{document}

\maketitle

\section{Introduction}
This is a sample document demonstrating various \LaTeX{} features including mathematical equations, tables, and formatting. You can cite papers using the citation search feature, and they will be automatically added to your \texttt{cite.bib} file.

\section{Mathematical Expressions}
Here are some example equations:

\subsection{Quadratic Formula}
The quadratic formula is given by:
\[x = \frac{-b \pm \sqrt{b^2 - 4ac}}{2a}\]

\subsection{Matrix Example}
A sample matrix:
\[
\begin{pmatrix}
a_{11} & a_{12} & a_{13} \\
a_{21} & a_{22} & a_{23} \\
a_{31} & a_{32} & a_{33}
\end{pmatrix}
\]

\section{Tables}
Below is a sample table showing data:

\begin{table}[h]
\centering
\begin{tabular}{@{}lrrr@{}}
\toprule
\textbf{Category} & \textbf{Value 1} & \textbf{Value 2} & \textbf{Total} \\
\midrule
Item A & 10.5 & 20.3 & 30.8 \\
Item B & 15.2 & 25.1 & 40.3 \\
Item C & 22.1 & 12.7 & 34.8 \\
\bottomrule
\end{tabular}
\caption{Sample Data Table}
\label{tab:sample}
\end{table}

\section{Lists}
\subsection{Bullet Points}
\begin{itemize}
  \item First item
  \item Second item
  \item Third item with sub-items:
    \begin{itemize}
      \item Sub-item 1
      \item Sub-item 2
    \end{itemize}
\end{itemize}

\subsection{Numbered List}
\begin{enumerate}
  \item First step
  \item Second step
  \item Third step
\end{enumerate}

\section{Citations and References}
When you use the citation search feature, references will be automatically managed. For example, you might cite a paper like this: \cite{example2024}. All citations will be automatically added to your \texttt{cite.bib} file.

\section{Conclusion}
This template demonstrates basic \LaTeX{} features including mathematical formulas, tables, and lists. The bibliography system is now set up to work seamlessly with Bibby AI's citation management. Feel free to modify it for your needs!

\bibliographystyle{plain}
\bibliography{cite}

\end{document}\documentclass{article}
\usepackage[utf8]{inputenc}
\usepackage{booktabs}
\usepackage{amsmath}
\usepackage{natbib}

\title{My Document}
\author{Your Name}
\date{\today}

\begin{document}

\maketitle

\section{Introduction}
This is a sample document demonstrating various \LaTeX{} features including mathematical equations, tables, and formatting. You can cite papers using the citation search feature, and they will be automatically added to your \texttt{cite.bib} file.

\section{Mathematical Expressions}
Here are some example equations:

\subsection{Quadratic Formula}
The quadratic formula is given by:
\[x = \frac{-b \pm \sqrt{b^2 - 4ac}}{2a}\]

\subsection{Matrix Example}
A sample matrix:
\[
\begin{pmatrix}
a_{11} & a_{12} & a_{13} \\
a_{21} & a_{22} & a_{23} \\
a_{31} & a_{32} & a_{33}
\end{pmatrix}
\]

\section{Tables}
Below is a sample table showing data:

\begin{table}[h]
\centering
\begin{tabular}{@{}lrrr@{}}
\toprule
\textbf{Category} & \textbf{Value 1} & \textbf{Value 2} & \textbf{Total} \\
\midrule
Item A & 10.5 & 20.3 & 30.8 \\
Item B & 15.2 & 25.1 & 40.3 \\
Item C & 22.1 & 12.7 & 34.8 \\
\bottomrule
\end{tabular}
\caption{Sample Data Table}
\label{tab:sample}
\end{table}

\section{Lists}
\subsection{Bullet Points}
\begin{itemize}
  \item First item
  \item Second item
  \item Third item with sub-items:
    \begin{itemize}
      \item Sub-item 1
      \item Sub-item 2
    \end{itemize}
\end{itemize}

\subsection{Numbered List}
\begin{enumerate}
  \item First step
  \item Second step
  \item Third step
\end{enumerate}

\section{Citations and References}
When you use the citation search feature, references will be automatically managed. For example, you might cite a paper like this: \cite{example2024}. All citations will be automatically added to your \texttt{cite.bib} file.

\section{Conclusion}
This template demonstrates basic \LaTeX{} features including mathematical formulas, tables, and lists. The bibliography system is now set up to work seamlessly with Bibby AI's citation management. Feel free to modify it for your needs!

\bibliographystyle{plain}
\bibliography{cite}

\end{document}\documentclass{article}
\usepackage[utf8]{inputenc}
\usepackage{booktabs}
\usepackage{amsmath}
\usepackage{natbib}

\title{My Document}
\author{Your Name}
\date{\today}

\begin{document}

\maketitle

\section{Introduction}
This is a sample document demonstrating various \LaTeX{} features including mathematical equations, tables, and formatting. You can cite papers using the citation search feature, and they will be automatically added to your \texttt{cite.bib} file.

\section{Mathematical Expressions}
Here are some example equations:

\subsection{Quadratic Formula}
The quadratic formula is given by:
\[x = \frac{-b \pm \sqrt{b^2 - 4ac}}{2a}\]

\subsection{Matrix Example}
A sample matrix:
\[
\begin{pmatrix}
a_{11} & a_{12} & a_{13} \\
a_{21} & a_{22} & a_{23} \\
a_{31} & a_{32} & a_{33}
\end{pmatrix}
\]

\section{Tables}
Below is a sample table showing data:

\begin{table}[h]
\centering
\begin{tabular}{@{}lrrr@{}}
\toprule
\textbf{Category} & \textbf{Value 1} & \textbf{Value 2} & \textbf{Total} \\
\midrule
Item A & 10.5 & 20.3 & 30.8 \\
Item B & 15.2 & 25.1 & 40.3 \\
Item C & 22.1 & 12.7 & 34.8 \\
\bottomrule
\end{tabular}
\caption{Sample Data Table}
\label{tab:sample}
\end{table}

\section{Lists}
\subsection{Bullet Points}
\begin{itemize}
  \item First item
  \item Second item
  \item Third item with sub-items:
    \begin{itemize}
      \item Sub-item 1
      \item Sub-item 2
    \end{itemize}
\end{itemize}

\subsection{Numbered List}
\begin{enumerate}
  \item First step
  \item Second step
  \item Third step
\end{enumerate}

\section{Citations and References}
When you use the citation search feature, references will be automatically managed. For example, you might cite a paper like this: \cite{example2024}. All citations will be automatically added to your \texttt{cite.bib} file.

\section{Conclusion}
This template demonstrates basic \LaTeX{} features including mathematical formulas, tables, and lists. The bibliography system is now set up to work seamlessly with Bibby AI's citation management. Feel free to modify it for your needs!

\bibliographystyle{plain}
\bibliography{cite}

\end{document}\documentclass{article}
\usepackage[utf8]{inputenc}
\usepackage{booktabs}
\usepackage{amsmath}
\usepackage{natbib}

\title{My Document}
\author{Your Name}
\date{\today}

\begin{document}

\maketitle

\section{Introduction}
This is a sample document demonstrating various \LaTeX{} features including mathematical equations, tables, and formatting. You can cite papers using the citation search feature, and they will be automatically added to your \texttt{cite.bib} file.

\section{Mathematical Expressions}
Here are some example equations:

\subsection{Quadratic Formula}
The quadratic formula is given by:
\[x = \frac{-b \pm \sqrt{b^2 - 4ac}}{2a}\]

\subsection{Matrix Example}
A sample matrix:
\[
\begin{pmatrix}
a_{11} & a_{12} & a_{13} \\
a_{21} & a_{22} & a_{23} \\
a_{31} & a_{32} & a_{33}
\end{pmatrix}
\]

\section{Tables}
Below is a sample table showing data:

\begin{table}[h]
\centering
\begin{tabular}{@{}lrrr@{}}
\toprule
\textbf{Category} & \textbf{Value 1} & \textbf{Value 2} & \textbf{Total} \\
\midrule
Item A & 10.5 & 20.3 & 30.8 \\
Item B & 15.2 & 25.1 & 40.3 \\
Item C & 22.1 & 12.7 & 34.8 \\
\bottomrule
\end{tabular}
\caption{Sample Data Table}
\label{tab:sample}
\end{table}

\section{Lists}
\subsection{Bullet Points}
\begin{itemize}
  \item First item
  \item Second item
  \item Third item with sub-items:
    \begin{itemize}
      \item Sub-item 1
      \item Sub-item 2
    \end{itemize}
\end{itemize}

\subsection{Numbered List}
\begin{enumerate}
  \item First step
  \item Second step
  \item Third step
\end{enumerate}

\section{Citations and References}
When you use the citation search feature, references will be automatically managed. For example, you might cite a paper like this: \cite{example2024}. All citations will be automatically added to your \texttt{cite.bib} file.

\section{Conclusion}
This template demonstrates basic \LaTeX{} features including mathematical formulas, tables, and lists. The bibliography system is now set up to work seamlessly with Bibby AI's citation management. Feel free to modify it for your needs!

\bibliographystyle{plain}
\bibliography{cite}

\end{document}\documentclass{article}
\usepackage[utf8]{inputenc}
\usepackage{booktabs}
\usepackage{amsmath}
\usepackage{natbib}

\title{My Document}
\author{Your Name}
\date{\today}

\begin{document}

\maketitle

\section{Introduction}
This is a sample document demonstrating various \LaTeX{} features including mathematical equations, tables, and formatting. You can cite papers using the citation search feature, and they will be automatically added to your \texttt{cite.bib} file.

\section{Mathematical Expressions}
Here are some example equations:

\subsection{Quadratic Formula}
The quadratic formula is given by:
\[x = \frac{-b \pm \sqrt{b^2 - 4ac}}{2a}\]

\subsection{Matrix Example}
A sample matrix:
\[
\begin{pmatrix}
a_{11} & a_{12} & a_{13} \\
a_{21} & a_{22} & a_{23} \\
a_{31} & a_{32} & a_{33}
\end{pmatrix}
\]

\section{Tables}
Below is a sample table showing data:

\begin{table}[h]
\centering
\begin{tabular}{@{}lrrr@{}}
\toprule
\textbf{Category} & \textbf{Value 1} & \textbf{Value 2} & \textbf{Total} \\
\midrule
Item A & 10.5 & 20.3 & 30.8 \\
Item B & 15.2 & 25.1 & 40.3 \\
Item C & 22.1 & 12.7 & 34.8 \\
\bottomrule
\end{tabular}
\caption{Sample Data Table}
\label{tab:sample}
\end{table}

\section{Lists}
\subsection{Bullet Points}
\begin{itemize}
  \item First item
  \item Second item
  \item Third item with sub-items:
    \begin{itemize}
      \item Sub-item 1
      \item Sub-item 2
    \end{itemize}
\end{itemize}

\subsection{Numbered List}
\begin{enumerate}
  \item First step
  \item Second step
  \item Third step
\end{enumerate}

\section{Citations and References}
When you use the citation search feature, references will be automatically managed. For example, you might cite a paper like this: \cite{example2024}. All citations will be automatically added to your \texttt{cite.bib} file.

\section{Conclusion}
This template demonstrates basic \LaTeX{} features including mathematical formulas, tables, and lists. The bibliography system is now set up to work seamlessly with Bibby AI's citation management. Feel free to modify it for your needs!

\bibliographystyle{plain}
\bibliography{cite}

\end{document}\documentclass{article}
\usepackage[utf8]{inputenc}
\usepackage{booktabs}
\usepackage{amsmath}
\usepackage{natbib}

\title{My Document}
\author{Your Name}
\date{\today}

\begin{document}

\maketitle

\section{Introduction}
This is a sample document demonstrating various \LaTeX{} features including mathematical equations, tables, and formatting. You can cite papers using the citation search feature, and they will be automatically added to your \texttt{cite.bib} file.

\section{Mathematical Expressions}
Here are some example equations:

\subsection{Quadratic Formula}
The quadratic formula is given by:
\[x = \frac{-b \pm \sqrt{b^2 - 4ac}}{2a}\]

\subsection{Matrix Example}
A sample matrix:
\[
\begin{pmatrix}
a_{11} & a_{12} & a_{13} \\
a_{21} & a_{22} & a_{23} \\
a_{31} & a_{32} & a_{33}
\end{pmatrix}
\]

\section{Tables}
Below is a sample table showing data:

\begin{table}[h]
\centering
\begin{tabular}{@{}lrrr@{}}
\toprule
\textbf{Category} & \textbf{Value 1} & \textbf{Value 2} & \textbf{Total} \\
\midrule
Item A & 10.5 & 20.3 & 30.8 \\
Item B & 15.2 & 25.1 & 40.3 \\
Item C & 22.1 & 12.7 & 34.8 \\
\bottomrule
\end{tabular}
\caption{Sample Data Table}
\label{tab:sample}
\end{table}

\section{Lists}
\subsection{Bullet Points}
\begin{itemize}
  \item First item
  \item Second item
  \item Third item with sub-items:
    \begin{itemize}
      \item Sub-item 1
      \item Sub-item 2
    \end{itemize}
\end{itemize}

\subsection{Numbered List}
\begin{enumerate}
  \item First step
  \item Second step
  \item Third step
\end{enumerate}

\section{Citations and References}
When you use the citation search feature, references will be automatically managed. For example, you might cite a paper like this: \cite{example2024}. All citations will be automatically added to your \texttt{cite.bib} file.

\section{Conclusion}
This template demonstrates basic \LaTeX{} features including mathematical formulas, tables, and lists. The bibliography system is now set up to work seamlessly with Bibby AI's citation management. Feel free to modify it for your needs!

\bibliographystyle{plain}
\bibliography{cite}

\end{document}\documentclass{article}
\usepackage[utf8]{inputenc}
\usepackage{booktabs}
\usepackage{amsmath}
\usepackage{natbib}

\title{My Document}
\author{Your Name}
\date{\today}

\begin{document}

\maketitle

\section{Introduction}
This is a sample document demonstrating various \LaTeX{} features including mathematical equations, tables, and formatting. You can cite papers using the citation search feature, and they will be automatically added to your \texttt{cite.bib} file.

\section{Mathematical Expressions}
Here are some example equations:

\subsection{Quadratic Formula}
The quadratic formula is given by:
\[x = \frac{-b \pm \sqrt{b^2 - 4ac}}{2a}\]

\subsection{Matrix Example}
A sample matrix:
\[
\begin{pmatrix}
a_{11} & a_{12} & a_{13} \\
a_{21} & a_{22} & a_{23} \\
a_{31} & a_{32} & a_{33}
\end{pmatrix}
\]

\section{Tables}
Below is a sample table showing data:

\begin{table}[h]
\centering
\begin{tabular}{@{}lrrr@{}}
\toprule
\textbf{Category} & \textbf{Value 1} & \textbf{Value 2} & \textbf{Total} \\
\midrule
Item A & 10.5 & 20.3 & 30.8 \\
Item B & 15.2 & 25.1 & 40.3 \\
Item C & 22.1 & 12.7 & 34.8 \\
\bottomrule
\end{tabular}
\caption{Sample Data Table}
\label{tab:sample}
\end{table}

\section{Lists}
\subsection{Bullet Points}
\begin{itemize}
  \item First item
  \item Second item
  \item Third item with sub-items:
    \begin{itemize}
      \item Sub-item 1
      \item Sub-item 2
    \end{itemize}
\end{itemize}

\subsection{Numbered List}
\begin{enumerate}
  \item First step
  \item Second step
  \item Third step
\end{enumerate}

\section{Citations and References}
When you use the citation search feature, references will be automatically managed. For example, you might cite a paper like this: \cite{example2024}. All citations will be automatically added to your \texttt{cite.bib} file.

\section{Conclusion}
This template demonstrates basic \LaTeX{} features including mathematical formulas, tables, and lists. The bibliography system is now set up to work seamlessly with Bibby AI's citation management. Feel free to modify it for your needs!

\bibliographystyle{plain}
\bibliography{cite}

\end{document}\documentclass{article}
\usepackage[utf8]{inputenc}
\usepackage{booktabs}
\usepackage{amsmath}
\usepackage{natbib}

\title{My Document}
\author{Your Name}
\date{\today}

\begin{document}

\maketitle

\section{Introduction}
This is a sample document demonstrating various \LaTeX{} features including mathematical equations, tables, and formatting. You can cite papers using the citation search feature, and they will be automatically added to your \texttt{cite.bib} file.

\section{Mathematical Expressions}
Here are some example equations:

\subsection{Quadratic Formula}
The quadratic formula is given by:
\[x = \frac{-b \pm \sqrt{b^2 - 4ac}}{2a}\]

\subsection{Matrix Example}
A sample matrix:
\[
\begin{pmatrix}
a_{11} & a_{12} & a_{13} \\
a_{21} & a_{22} & a_{23} \\
a_{31} & a_{32} & a_{33}
\end{pmatrix}
\]

\section{Tables}
Below is a sample table showing data:

\begin{table}[h]
\centering
\begin{tabular}{@{}lrrr@{}}
\toprule
\textbf{Category} & \textbf{Value 1} & \textbf{Value 2} & \textbf{Total} \\
\midrule
Item A & 10.5 & 20.3 & 30.8 \\
Item B & 15.2 & 25.1 & 40.3 \\
Item C & 22.1 & 12.7 & 34.8 \\
\bottomrule
\end{tabular}
\caption{Sample Data Table}
\label{tab:sample}
\end{table}

\section{Lists}
\subsection{Bullet Points}
\begin{itemize}
  \item First item
  \item Second item
  \item Third item with sub-items:
    \begin{itemize}
      \item Sub-item 1
      \item Sub-item 2
    \end{itemize}
\end{itemize}

\subsection{Numbered List}
\begin{enumerate}
  \item First step
  \item Second step
  \item Third step
\end{enumerate}

\section{Citations and References}
When you use the citation search feature, references will be automatically managed. For example, you might cite a paper like this: \cite{example2024}. All citations will be automatically added to your \texttt{cite.bib} file.

\section{Conclusion}
This template demonstrates basic \LaTeX{} features including mathematical formulas, tables, and lists. The bibliography system is now set up to work seamlessly with Bibby AI's citation management. Feel free to modify it for your needs!

\bibliographystyle{plain}
\bibliography{cite}

\end{document}\documentclass{article}
\usepackage[utf8]{inputenc}
\usepackage{booktabs}
\usepackage{amsmath}
\usepackage{natbib}

\title{My Document}
\author{Your Name}
\date{\today}

\begin{document}

\maketitle

\section{Introduction}
This is a sample document demonstrating various \LaTeX{} features including mathematical equations, tables, and formatting. You can cite papers using the citation search feature, and they will be automatically added to your \texttt{cite.bib} file.

\section{Mathematical Expressions}
Here are some example equations:

\subsection{Quadratic Formula}
The quadratic formula is given by:
\[x = \frac{-b \pm \sqrt{b^2 - 4ac}}{2a}\]

\subsection{Matrix Example}
A sample matrix:
\[
\begin{pmatrix}
a_{11} & a_{12} & a_{13} \\
a_{21} & a_{22} & a_{23} \\
a_{31} & a_{32} & a_{33}
\end{pmatrix}
\]

\section{Tables}
Below is a sample table showing data:

\begin{table}[h]
\centering
\begin{tabular}{@{}lrrr@{}}
\toprule
\textbf{Category} & \textbf{Value 1} & \textbf{Value 2} & \textbf{Total} \\
\midrule
Item A & 10.5 & 20.3 & 30.8 \\
Item B & 15.2 & 25.1 & 40.3 \\
Item C & 22.1 & 12.7 & 34.8 \\
\bottomrule
\end{tabular}
\caption{Sample Data Table}
\label{tab:sample}
\end{table}

\section{Lists}
\subsection{Bullet Points}
\begin{itemize}
  \item First item
  \item Second item
  \item Third item with sub-items:
    \begin{itemize}
      \item Sub-item 1
      \item Sub-item 2
    \end{itemize}
\end{itemize}

\subsection{Numbered List}
\begin{enumerate}
  \item First step
  \item Second step
  \item Third step
\end{enumerate}

\section{Citations and References}
When you use the citation search feature, references will be automatically managed. For example, you might cite a paper like this: \cite{example2024}. All citations will be automatically added to your \texttt{cite.bib} file.

\section{Conclusion}
This template demonstrates basic \LaTeX{} features including mathematical formulas, tables, and lists. The bibliography system is now set up to work seamlessly with Bibby AI's citation management. Feel free to modify it for your needs!

\bibliographystyle{plain}
\bibliography{cite}

\end{document}\documentclass{article}
\usepackage[utf8]{inputenc}
\usepackage{booktabs}
\usepackage{amsmath}
\usepackage{natbib}

\title{My Document}
\author{Your Name}
\date{\today}

\begin{document}

\maketitle

\section{Introduction}
This is a sample document demonstrating various \LaTeX{} features including mathematical equations, tables, and formatting. You can cite papers using the citation search feature, and they will be automatically added to your \texttt{cite.bib} file.

\section{Mathematical Expressions}
Here are some example equations:

\subsection{Quadratic Formula}
The quadratic formula is given by:
\[x = \frac{-b \pm \sqrt{b^2 - 4ac}}{2a}\]

\subsection{Matrix Example}
A sample matrix:
\[
\begin{pmatrix}
a_{11} & a_{12} & a_{13} \\
a_{21} & a_{22} & a_{23} \\
a_{31} & a_{32} & a_{33}
\end{pmatrix}
\]

\section{Tables}
Below is a sample table showing data:

\begin{table}[h]
\centering
\begin{tabular}{@{}lrrr@{}}
\toprule
\textbf{Category} & \textbf{Value 1} & \textbf{Value 2} & \textbf{Total} \\
\midrule
Item A & 10.5 & 20.3 & 30.8 \\
Item B & 15.2 & 25.1 & 40.3 \\
Item C & 22.1 & 12.7 & 34.8 \\
\bottomrule
\end{tabular}
\caption{Sample Data Table}
\label{tab:sample}
\end{table}

\section{Lists}
\subsection{Bullet Points}
\begin{itemize}
  \item First item
  \item Second item
  \item Third item with sub-items:
    \begin{itemize}
      \item Sub-item 1
      \item Sub-item 2
    \end{itemize}
\end{itemize}

\subsection{Numbered List}
\begin{enumerate}
  \item First step
  \item Second step
  \item Third step
\end{enumerate}

\section{Citations and References}
When you use the citation search feature, references will be automatically managed. For example, you might cite a paper like this: \cite{example2024}. All citations will be automatically added to your \texttt{cite.bib} file.

\section{Conclusion}
This template demonstrates basic \LaTeX{} features including mathematical formulas, tables, and lists. The bibliography system is now set up to work seamlessly with Bibby AI's citation management. Feel free to modify it for your needs!

\bibliographystyle{plain}
\bibliography{cite}

\end{document}\documentclass{article}
\usepackage[utf8]{inputenc}
\usepackage{booktabs}
\usepackage{amsmath}
\usepackage{natbib}

\title{My Document}
\author{Your Name}
\date{\today}

\begin{document}

\maketitle

\section{Introduction}
This is a sample document demonstrating various \LaTeX{} features including mathematical equations, tables, and formatting. You can cite papers using the citation search feature, and they will be automatically added to your \texttt{cite.bib} file.

\section{Mathematical Expressions}
Here are some example equations:

\subsection{Quadratic Formula}
The quadratic formula is given by:
\[x = \frac{-b \pm \sqrt{b^2 - 4ac}}{2a}\]

\subsection{Matrix Example}
A sample matrix:
\[
\begin{pmatrix}
a_{11} & a_{12} & a_{13} \\
a_{21} & a_{22} & a_{23} \\
a_{31} & a_{32} & a_{33}
\end{pmatrix}
\]

\section{Tables}
Below is a sample table showing data:

\begin{table}[h]
\centering
\begin{tabular}{@{}lrrr@{}}
\toprule
\textbf{Category} & \textbf{Value 1} & \textbf{Value 2} & \textbf{Total} \\
\midrule
Item A & 10.5 & 20.3 & 30.8 \\
Item B & 15.2 & 25.1 & 40.3 \\
Item C & 22.1 & 12.7 & 34.8 \\
\bottomrule
\end{tabular}
\caption{Sample Data Table}
\label{tab:sample}
\end{table}

\section{Lists}
\subsection{Bullet Points}
\begin{itemize}
  \item First item
  \item Second item
  \item Third item with sub-items:
    \begin{itemize}
      \item Sub-item 1
      \item Sub-item 2
    \end{itemize}
\end{itemize}

\subsection{Numbered List}
\begin{enumerate}
  \item First step
  \item Second step
  \item Third step
\end{enumerate}

\section{Citations and References}
When you use the citation search feature, references will be automatically managed. For example, you might cite a paper like this: \cite{example2024}. All citations will be automatically added to your \texttt{cite.bib} file.

\section{Conclusion}
This template demonstrates basic \LaTeX{} features including mathematical formulas, tables, and lists. The bibliography system is now set up to work seamlessly with Bibby AI's citation management. Feel free to modify it for your needs!

\bibliographystyle{plain}
\bibliography{cite}

\end{document}\documentclass{article}
\usepackage[utf8]{inputenc}
\usepackage{booktabs}
\usepackage{amsmath}
\usepackage{natbib}

\title{My Document}
\author{Your Name}
\date{\today}

\begin{document}

\maketitle

\section{Introduction}
This is a sample document demonstrating various \LaTeX{} features including mathematical equations, tables, and formatting. You can cite papers using the citation search feature, and they will be automatically added to your \texttt{cite.bib} file.

\section{Mathematical Expressions}
Here are some example equations:

\subsection{Quadratic Formula}
The quadratic formula is given by:
\[x = \frac{-b \pm \sqrt{b^2 - 4ac}}{2a}\]

\subsection{Matrix Example}
A sample matrix:
\[
\begin{pmatrix}
a_{11} & a_{12} & a_{13} \\
a_{21} & a_{22} & a_{23} \\
a_{31} & a_{32} & a_{33}
\end{pmatrix}
\]

\section{Tables}
Below is a sample table showing data:

\begin{table}[h]
\centering
\begin{tabular}{@{}lrrr@{}}
\toprule
\textbf{Category} & \textbf{Value 1} & \textbf{Value 2} & \textbf{Total} \\
\midrule
Item A & 10.5 & 20.3 & 30.8 \\
Item B & 15.2 & 25.1 & 40.3 \\
Item C & 22.1 & 12.7 & 34.8 \\
\bottomrule
\end{tabular}
\caption{Sample Data Table}
\label{tab:sample}
\end{table}

\section{Lists}
\subsection{Bullet Points}
\begin{itemize}
  \item First item
  \item Second item
  \item Third item with sub-items:
    \begin{itemize}
      \item Sub-item 1
      \item Sub-item 2
    \end{itemize}
\end{itemize}

\subsection{Numbered List}
\begin{enumerate}
  \item First step
  \item Second step
  \item Third step
\end{enumerate}

\section{Citations and References}
When you use the citation search feature, references will be automatically managed. For example, you might cite a paper like this: \cite{example2024}. All citations will be automatically added to your \texttt{cite.bib} file.

\section{Conclusion}
This template demonstrates basic \LaTeX{} features including mathematical formulas, tables, and lists. The bibliography system is now set up to work seamlessly with Bibby AI's citation management. Feel free to modify it for your needs!

\bibliographystyle{plain}
\bibliography{cite}

\end{document}\documentclass{article}
\usepackage[utf8]{inputenc}
\usepackage{booktabs}
\usepackage{amsmath}
\usepackage{natbib}

\title{My Document}
\author{Your Name}
\date{\today}

\begin{document}

\maketitle

\section{Introduction}
This is a sample document demonstrating various \LaTeX{} features including mathematical equations, tables, and formatting. You can cite papers using the citation search feature, and they will be automatically added to your \texttt{cite.bib} file.

\section{Mathematical Expressions}
Here are some example equations:

\subsection{Quadratic Formula}
The quadratic formula is given by:
\[x = \frac{-b \pm \sqrt{b^2 - 4ac}}{2a}\]

\subsection{Matrix Example}
A sample matrix:
\[
\begin{pmatrix}
a_{11} & a_{12} & a_{13} \\
a_{21} & a_{22} & a_{23} \\
a_{31} & a_{32} & a_{33}
\end{pmatrix}
\]

\section{Tables}
Below is a sample table showing data:

\begin{table}[h]
\centering
\begin{tabular}{@{}lrrr@{}}
\toprule
\textbf{Category} & \textbf{Value 1} & \textbf{Value 2} & \textbf{Total} \\
\midrule
Item A & 10.5 & 20.3 & 30.8 \\
Item B & 15.2 & 25.1 & 40.3 \\
Item C & 22.1 & 12.7 & 34.8 \\
\bottomrule
\end{tabular}
\caption{Sample Data Table}
\label{tab:sample}
\end{table}

\section{Lists}
\subsection{Bullet Points}
\begin{itemize}
  \item First item
  \item Second item
  \item Third item with sub-items:
    \begin{itemize}
      \item Sub-item 1
      \item Sub-item 2
    \end{itemize}
\end{itemize}

\subsection{Numbered List}
\begin{enumerate}
  \item First step
  \item Second step
  \item Third step
\end{enumerate}

\section{Citations and References}
When you use the citation search feature, references will be automatically managed. For example, you might cite a paper like this: \cite{example2024}. All citations will be automatically added to your \texttt{cite.bib} file.

\section{Conclusion}
This template demonstrates basic \LaTeX{} features including mathematical formulas, tables, and lists. The bibliography system is now set up to work seamlessly with Bibby AI's citation management. Feel free to modify it for your needs!

\bibliographystyle{plain}
\bibliography{cite}

\end{document}\documentclass{article}
\usepackage[utf8]{inputenc}
\usepackage{booktabs}
\usepackage{amsmath}
\usepackage{natbib}

\title{My Document}
\author{Your Name}
\date{\today}

\begin{document}

\maketitle

\section{Introduction}
This is a sample document demonstrating various \LaTeX{} features including mathematical equations, tables, and formatting. You can cite papers using the citation search feature, and they will be automatically added to your \texttt{cite.bib} file.

\section{Mathematical Expressions}
Here are some example equations:

\subsection{Quadratic Formula}
The quadratic formula is given by:
\[x = \frac{-b \pm \sqrt{b^2 - 4ac}}{2a}\]

\subsection{Matrix Example}
A sample matrix:
\[
\begin{pmatrix}
a_{11} & a_{12} & a_{13} \\
a_{21} & a_{22} & a_{23} \\
a_{31} & a_{32} & a_{33}
\end{pmatrix}
\]

\section{Tables}
Below is a sample table showing data:

\begin{table}[h]
\centering
\begin{tabular}{@{}lrrr@{}}
\toprule
\textbf{Category} & \textbf{Value 1} & \textbf{Value 2} & \textbf{Total} \\
\midrule
Item A & 10.5 & 20.3 & 30.8 \\
Item B & 15.2 & 25.1 & 40.3 \\
Item C & 22.1 & 12.7 & 34.8 \\
\bottomrule
\end{tabular}
\caption{Sample Data Table}
\label{tab:sample}
\end{table}

\section{Lists}
\subsection{Bullet Points}
\begin{itemize}
  \item First item
  \item Second item
  \item Third item with sub-items:
    \begin{itemize}
      \item Sub-item 1
      \item Sub-item 2
    \end{itemize}
\end{itemize}

\subsection{Numbered List}
\begin{enumerate}
  \item First step
  \item Second step
  \item Third step
\end{enumerate}

\section{Citations and References}
When you use the citation search feature, references will be automatically managed. For example, you might cite a paper like this: \cite{example2024}. All citations will be automatically added to your \texttt{cite.bib} file.

\section{Conclusion}
This template demonstrates basic \LaTeX{} features including mathematical formulas, tables, and lists. The bibliography system is now set up to work seamlessly with Bibby AI's citation management. Feel free to modify it for your needs!

\bibliographystyle{plain}
\bibliography{cite}

\end{document}\documentclass{article}
\usepackage[utf8]{inputenc}
\usepackage{booktabs}
\usepackage{amsmath}
\usepackage{natbib}

\title{My Document}
\author{Your Name}
\date{\today}

\begin{document}

\maketitle

\section{Introduction}
This is a sample document demonstrating various \LaTeX{} features including mathematical equations, tables, and formatting. You can cite papers using the citation search feature, and they will be automatically added to your \texttt{cite.bib} file.

\section{Mathematical Expressions}
Here are some example equations:

\subsection{Quadratic Formula}
The quadratic formula is given by:
\[x = \frac{-b \pm \sqrt{b^2 - 4ac}}{2a}\]

\subsection{Matrix Example}
A sample matrix:
\[
\begin{pmatrix}
a_{11} & a_{12} & a_{13} \\
a_{21} & a_{22} & a_{23} \\
a_{31} & a_{32} & a_{33}
\end{pmatrix}
\]

\section{Tables}
Below is a sample table showing data:

\begin{table}[h]
\centering
\begin{tabular}{@{}lrrr@{}}
\toprule
\textbf{Category} & \textbf{Value 1} & \textbf{Value 2} & \textbf{Total} \\
\midrule
Item A & 10.5 & 20.3 & 30.8 \\
Item B & 15.2 & 25.1 & 40.3 \\
Item C & 22.1 & 12.7 & 34.8 \\
\bottomrule
\end{tabular}
\caption{Sample Data Table}
\label{tab:sample}
\end{table}

\section{Lists}
\subsection{Bullet Points}
\begin{itemize}
  \item First item
  \item Second item
  \item Third item with sub-items:
    \begin{itemize}
      \item Sub-item 1
      \item Sub-item 2
    \end{itemize}
\end{itemize}

\subsection{Numbered List}
\begin{enumerate}
  \item First step
  \item Second step
  \item Third step
\end{enumerate}

\section{Citations and References}
When you use the citation search feature, references will be automatically managed. For example, you might cite a paper like this: \cite{example2024}. All citations will be automatically added to your \texttt{cite.bib} file.

\section{Conclusion}
This template demonstrates basic \LaTeX{} features including mathematical formulas, tables, and lists. The bibliography system is now set up to work seamlessly with Bibby AI's citation management. Feel free to modify it for your needs!

\bibliographystyle{plain}
\bibliography{cite}

\end{document}\documentclass{article}
\usepackage[utf8]{inputenc}
\usepackage{booktabs}
\usepackage{amsmath}
\usepackage{natbib}

\title{My Document}
\author{Your Name}
\date{\today}

\begin{document}

\maketitle

\section{Introduction}
This is a sample document demonstrating various \LaTeX{} features including mathematical equations, tables, and formatting. You can cite papers using the citation search feature, and they will be automatically added to your \texttt{cite.bib} file.

\section{Mathematical Expressions}
Here are some example equations:

\subsection{Quadratic Formula}
The quadratic formula is given by:
\[x = \frac{-b \pm \sqrt{b^2 - 4ac}}{2a}\]

\subsection{Matrix Example}
A sample matrix:
\[
\begin{pmatrix}
a_{11} & a_{12} & a_{13} \\
a_{21} & a_{22} & a_{23} \\
a_{31} & a_{32} & a_{33}
\end{pmatrix}
\]

\section{Tables}
Below is a sample table showing data:

\begin{table}[h]
\centering
\begin{tabular}{@{}lrrr@{}}
\toprule
\textbf{Category} & \textbf{Value 1} & \textbf{Value 2} & \textbf{Total} \\
\midrule
Item A & 10.5 & 20.3 & 30.8 \\
Item B & 15.2 & 25.1 & 40.3 \\
Item C & 22.1 & 12.7 & 34.8 \\
\bottomrule
\end{tabular}
\caption{Sample Data Table}
\label{tab:sample}
\end{table}

\section{Lists}
\subsection{Bullet Points}
\begin{itemize}
  \item First item
  \item Second item
  \item Third item with sub-items:
    \begin{itemize}
      \item Sub-item 1
      \item Sub-item 2
    \end{itemize}
\end{itemize}

\subsection{Numbered List}
\begin{enumerate}
  \item First step
  \item Second step
  \item Third step
\end{enumerate}

\section{Citations and References}
When you use the citation search feature, references will be automatically managed. For example, you might cite a paper like this: \cite{example2024}. All citations will be automatically added to your \texttt{cite.bib} file.

\section{Conclusion}
This template demonstrates basic \LaTeX{} features including mathematical formulas, tables, and lists. The bibliography system is now set up to work seamlessly with Bibby AI's citation management. Feel free to modify it for your needs!

\bibliographystyle{plain}
\bibliography{cite}

\end{document}\documentclass{article}
\usepackage[utf8]{inputenc}
\usepackage{booktabs}
\usepackage{amsmath}
\usepackage{natbib}

\title{My Document}
\author{Your Name}
\date{\today}

\begin{document}

\maketitle

\section{Introduction}
This is a sample document demonstrating various \LaTeX{} features including mathematical equations, tables, and formatting. You can cite papers using the citation search feature, and they will be automatically added to your \texttt{cite.bib} file.

\section{Mathematical Expressions}
Here are some example equations:

\subsection{Quadratic Formula}
The quadratic formula is given by:
\[x = \frac{-b \pm \sqrt{b^2 - 4ac}}{2a}\]

\subsection{Matrix Example}
A sample matrix:
\[
\begin{pmatrix}
a_{11} & a_{12} & a_{13} \\
a_{21} & a_{22} & a_{23} \\
a_{31} & a_{32} & a_{33}
\end{pmatrix}
\]

\section{Tables}
Below is a sample table showing data:

\begin{table}[h]
\centering
\begin{tabular}{@{}lrrr@{}}
\toprule
\textbf{Category} & \textbf{Value 1} & \textbf{Value 2} & \textbf{Total} \\
\midrule
Item A & 10.5 & 20.3 & 30.8 \\
Item B & 15.2 & 25.1 & 40.3 \\
Item C & 22.1 & 12.7 & 34.8 \\
\bottomrule
\end{tabular}
\caption{Sample Data Table}
\label{tab:sample}
\end{table}

\section{Lists}
\subsection{Bullet Points}
\begin{itemize}
  \item First item
  \item Second item
  \item Third item with sub-items:
    \begin{itemize}
      \item Sub-item 1
      \item Sub-item 2
    \end{itemize}
\end{itemize}

\subsection{Numbered List}
\begin{enumerate}
  \item First step
  \item Second step
  \item Third step
\end{enumerate}

\section{Citations and References}
When you use the citation search feature, references will be automatically managed. For example, you might cite a paper like this: \cite{example2024}. All citations will be automatically added to your \texttt{cite.bib} file.

\section{Conclusion}
This template demonstrates basic \LaTeX{} features including mathematical formulas, tables, and lists. The bibliography system is now set up to work seamlessly with Bibby AI's citation management. Feel free to modify it for your needs!

\bibliographystyle{plain}
\bibliography{cite}

\end{document}\documentclass{article}
\usepackage[utf8]{inputenc}
\usepackage{booktabs}
\usepackage{amsmath}
\usepackage{natbib}

\title{My Document}
\author{Your Name}
\date{\today}

\begin{document}

\maketitle

\section{Introduction}
This is a sample document demonstrating various \LaTeX{} features including mathematical equations, tables, and formatting. You can cite papers using the citation search feature, and they will be automatically added to your \texttt{cite.bib} file.

\section{Mathematical Expressions}
Here are some example equations:

\subsection{Quadratic Formula}
The quadratic formula is given by:
\[x = \frac{-b \pm \sqrt{b^2 - 4ac}}{2a}\]

\subsection{Matrix Example}
A sample matrix:
\[
\begin{pmatrix}
a_{11} & a_{12} & a_{13} \\
a_{21} & a_{22} & a_{23} \\
a_{31} & a_{32} & a_{33}
\end{pmatrix}
\]

\section{Tables}
Below is a sample table showing data:

\begin{table}[h]
\centering
\begin{tabular}{@{}lrrr@{}}
\toprule
\textbf{Category} & \textbf{Value 1} & \textbf{Value 2} & \textbf{Total} \\
\midrule
Item A & 10.5 & 20.3 & 30.8 \\
Item B & 15.2 & 25.1 & 40.3 \\
Item C & 22.1 & 12.7 & 34.8 \\
\bottomrule
\end{tabular}
\caption{Sample Data Table}
\label{tab:sample}
\end{table}

\section{Lists}
\subsection{Bullet Points}
\begin{itemize}
  \item First item
  \item Second item
  \item Third item with sub-items:
    \begin{itemize}
      \item Sub-item 1
      \item Sub-item 2
    \end{itemize}
\end{itemize}

\subsection{Numbered List}
\begin{enumerate}
  \item First step
  \item Second step
  \item Third step
\end{enumerate}

\section{Citations and References}
When you use the citation search feature, references will be automatically managed. For example, you might cite a paper like this: \cite{example2024}. All citations will be automatically added to your \texttt{cite.bib} file.

\section{Conclusion}
This template demonstrates basic \LaTeX{} features including mathematical formulas, tables, and lists. The bibliography system is now set up to work seamlessly with Bibby AI's citation management. Feel free to modify it for your needs!

\bibliographystyle{plain}
\bibliography{cite}

\end{document}\documentclass{article}
\usepackage[utf8]{inputenc}
\usepackage{booktabs}
\usepackage{amsmath}
\usepackage{natbib}

\title{My Document}
\author{Your Name}
\date{\today}

\begin{document}

\maketitle

\section{Introduction}
This is a sample document demonstrating various \LaTeX{} features including mathematical equations, tables, and formatting. You can cite papers using the citation search feature, and they will be automatically added to your \texttt{cite.bib} file.

\section{Mathematical Expressions}
Here are some example equations:

\subsection{Quadratic Formula}
The quadratic formula is given by:
\[x = \frac{-b \pm \sqrt{b^2 - 4ac}}{2a}\]

\subsection{Matrix Example}
A sample matrix:
\[
\begin{pmatrix}
a_{11} & a_{12} & a_{13} \\
a_{21} & a_{22} & a_{23} \\
a_{31} & a_{32} & a_{33}
\end{pmatrix}
\]

\section{Tables}
Below is a sample table showing data:

\begin{table}[h]
\centering
\begin{tabular}{@{}lrrr@{}}
\toprule
\textbf{Category} & \textbf{Value 1} & \textbf{Value 2} & \textbf{Total} \\
\midrule
Item A & 10.5 & 20.3 & 30.8 \\
Item B & 15.2 & 25.1 & 40.3 \\
Item C & 22.1 & 12.7 & 34.8 \\
\bottomrule
\end{tabular}
\caption{Sample Data Table}
\label{tab:sample}
\end{table}

\section{Lists}
\subsection{Bullet Points}
\begin{itemize}
  \item First item
  \item Second item
  \item Third item with sub-items:
    \begin{itemize}
      \item Sub-item 1
      \item Sub-item 2
    \end{itemize}
\end{itemize}

\subsection{Numbered List}
\begin{enumerate}
  \item First step
  \item Second step
  \item Third step
\end{enumerate}

\section{Citations and References}
When you use the citation search feature, references will be automatically managed. For example, you might cite a paper like this: \cite{example2024}. All citations will be automatically added to your \texttt{cite.bib} file.

\section{Conclusion}
This template demonstrates basic \LaTeX{} features including mathematical formulas, tables, and lists. The bibliography system is now set up to work seamlessly with Bibby AI's citation management. Feel free to modify it for your needs!

\bibliographystyle{plain}
\bibliography{cite}

\end{document}\documentclass{article}
\usepackage[utf8]{inputenc}
\usepackage{booktabs}
\usepackage{amsmath}
\usepackage{natbib}

\title{My Document}
\author{Your Name}
\date{\today}

\begin{document}

\maketitle

\section{Introduction}
This is a sample document demonstrating various \LaTeX{} features including mathematical equations, tables, and formatting. You can cite papers using the citation search feature, and they will be automatically added to your \texttt{cite.bib} file.

\section{Mathematical Expressions}
Here are some example equations:

\subsection{Quadratic Formula}
The quadratic formula is given by:
\[x = \frac{-b \pm \sqrt{b^2 - 4ac}}{2a}\]

\subsection{Matrix Example}
A sample matrix:
\[
\begin{pmatrix}
a_{11} & a_{12} & a_{13} \\
a_{21} & a_{22} & a_{23} \\
a_{31} & a_{32} & a_{33}
\end{pmatrix}
\]

\section{Tables}
Below is a sample table showing data:

\begin{table}[h]
\centering
\begin{tabular}{@{}lrrr@{}}
\toprule
\textbf{Category} & \textbf{Value 1} & \textbf{Value 2} & \textbf{Total} \\
\midrule
Item A & 10.5 & 20.3 & 30.8 \\
Item B & 15.2 & 25.1 & 40.3 \\
Item C & 22.1 & 12.7 & 34.8 \\
\bottomrule
\end{tabular}
\caption{Sample Data Table}
\label{tab:sample}
\end{table}

\section{Lists}
\subsection{Bullet Points}
\begin{itemize}
  \item First item
  \item Second item
  \item Third item with sub-items:
    \begin{itemize}
      \item Sub-item 1
      \item Sub-item 2
    \end{itemize}
\end{itemize}

\subsection{Numbered List}
\begin{enumerate}
  \item First step
  \item Second step
  \item Third step
\end{enumerate}

\section{Citations and References}
When you use the citation search feature, references will be automatically managed. For example, you might cite a paper like this: \cite{example2024}. All citations will be automatically added to your \texttt{cite.bib} file.

\section{Conclusion}
This template demonstrates basic \LaTeX{} features including mathematical formulas, tables, and lists. The bibliography system is now set up to work seamlessly with Bibby AI's citation management. Feel free to modify it for your needs!

\bibliographystyle{plain}
\bibliography{cite}

\end{document}\documentclass{article}
\usepackage[utf8]{inputenc}
\usepackage{booktabs}
\usepackage{amsmath}
\usepackage{natbib}

\title{My Document}
\author{Your Name}
\date{\today}

\begin{document}

\maketitle

\section{Introduction}
This is a sample document demonstrating various \LaTeX{} features including mathematical equations, tables, and formatting. You can cite papers using the citation search feature, and they will be automatically added to your \texttt{cite.bib} file.

\section{Mathematical Expressions}
Here are some example equations:

\subsection{Quadratic Formula}
The quadratic formula is given by:
\[x = \frac{-b \pm \sqrt{b^2 - 4ac}}{2a}\]

\subsection{Matrix Example}
A sample matrix:
\[
\begin{pmatrix}
a_{11} & a_{12} & a_{13} \\
a_{21} & a_{22} & a_{23} \\
a_{31} & a_{32} & a_{33}
\end{pmatrix}
\]

\section{Tables}
Below is a sample table showing data:

\begin{table}[h]
\centering
\begin{tabular}{@{}lrrr@{}}
\toprule
\textbf{Category} & \textbf{Value 1} & \textbf{Value 2} & \textbf{Total} \\
\midrule
Item A & 10.5 & 20.3 & 30.8 \\
Item B & 15.2 & 25.1 & 40.3 \\
Item C & 22.1 & 12.7 & 34.8 \\
\bottomrule
\end{tabular}
\caption{Sample Data Table}
\label{tab:sample}
\end{table}

\section{Lists}
\subsection{Bullet Points}
\begin{itemize}
  \item First item
  \item Second item
  \item Third item with sub-items:
    \begin{itemize}
      \item Sub-item 1
      \item Sub-item 2
    \end{itemize}
\end{itemize}

\subsection{Numbered List}
\begin{enumerate}
  \item First step
  \item Second step
  \item Third step
\end{enumerate}

\section{Citations and References}
When you use the citation search feature, references will be automatically managed. For example, you might cite a paper like this: \cite{example2024}. All citations will be automatically added to your \texttt{cite.bib} file.

\section{Conclusion}
This template demonstrates basic \LaTeX{} features including mathematical formulas, tables, and lists. The bibliography system is now set up to work seamlessly with Bibby AI's citation management. Feel free to modify it for your needs!

\bibliographystyle{plain}
\bibliography{cite}

\end{document}\documentclass{article}
\usepackage[utf8]{inputenc}
\usepackage{booktabs}
\usepackage{amsmath}
\usepackage{natbib}

\title{My Document}
\author{Your Name}
\date{\today}

\begin{document}

\maketitle

\section{Introduction}
This is a sample document demonstrating various \LaTeX{} features including mathematical equations, tables, and formatting. You can cite papers using the citation search feature, and they will be automatically added to your \texttt{cite.bib} file.

\section{Mathematical Expressions}
Here are some example equations:

\subsection{Quadratic Formula}
The quadratic formula is given by:
\[x = \frac{-b \pm \sqrt{b^2 - 4ac}}{2a}\]

\subsection{Matrix Example}
A sample matrix:
\[
\begin{pmatrix}
a_{11} & a_{12} & a_{13} \\
a_{21} & a_{22} & a_{23} \\
a_{31} & a_{32} & a_{33}
\end{pmatrix}
\]

\section{Tables}
Below is a sample table showing data:

\begin{table}[h]
\centering
\begin{tabular}{@{}lrrr@{}}
\toprule
\textbf{Category} & \textbf{Value 1} & \textbf{Value 2} & \textbf{Total} \\
\midrule
Item A & 10.5 & 20.3 & 30.8 \\
Item B & 15.2 & 25.1 & 40.3 \\
Item C & 22.1 & 12.7 & 34.8 \\
\bottomrule
\end{tabular}
\caption{Sample Data Table}
\label{tab:sample}
\end{table}

\section{Lists}
\subsection{Bullet Points}
\begin{itemize}
  \item First item
  \item Second item
  \item Third item with sub-items:
    \begin{itemize}
      \item Sub-item 1
      \item Sub-item 2
    \end{itemize}
\end{itemize}

\subsection{Numbered List}
\begin{enumerate}
  \item First step
  \item Second step
  \item Third step
\end{enumerate}

\section{Citations and References}
When you use the citation search feature, references will be automatically managed. For example, you might cite a paper like this: \cite{example2024}. All citations will be automatically added to your \texttt{cite.bib} file.

\section{Conclusion}
This template demonstrates basic \LaTeX{} features including mathematical formulas, tables, and lists. The bibliography system is now set up to work seamlessly with Bibby AI's citation management. Feel free to modify it for your needs!

\bibliographystyle{plain}
\bibliography{cite}

\end{document}\documentclass{article}
\usepackage[utf8]{inputenc}
\usepackage{booktabs}
\usepackage{amsmath}
\usepackage{natbib}

\title{My Document}
\author{Your Name}
\date{\today}

\begin{document}

\maketitle

\section{Introduction}
This is a sample document demonstrating various \LaTeX{} features including mathematical equations, tables, and formatting. You can cite papers using the citation search feature, and they will be automatically added to your \texttt{cite.bib} file.

\section{Mathematical Expressions}
Here are some example equations:

\subsection{Quadratic Formula}
The quadratic formula is given by:
\[x = \frac{-b \pm \sqrt{b^2 - 4ac}}{2a}\]

\subsection{Matrix Example}
A sample matrix:
\[
\begin{pmatrix}
a_{11} & a_{12} & a_{13} \\
a_{21} & a_{22} & a_{23} \\
a_{31} & a_{32} & a_{33}
\end{pmatrix}
\]

\section{Tables}
Below is a sample table showing data:

\begin{table}[h]
\centering
\begin{tabular}{@{}lrrr@{}}
\toprule
\textbf{Category} & \textbf{Value 1} & \textbf{Value 2} & \textbf{Total} \\
\midrule
Item A & 10.5 & 20.3 & 30.8 \\
Item B & 15.2 & 25.1 & 40.3 \\
Item C & 22.1 & 12.7 & 34.8 \\
\bottomrule
\end{tabular}
\caption{Sample Data Table}
\label{tab:sample}
\end{table}

\section{Lists}
\subsection{Bullet Points}
\begin{itemize}
  \item First item
  \item Second item
  \item Third item with sub-items:
    \begin{itemize}
      \item Sub-item 1
      \item Sub-item 2
    \end{itemize}
\end{itemize}

\subsection{Numbered List}
\begin{enumerate}
  \item First step
  \item Second step
  \item Third step
\end{enumerate}

\section{Citations and References}
When you use the citation search feature, references will be automatically managed. For example, you might cite a paper like this: \cite{example2024}. All citations will be automatically added to your \texttt{cite.bib} file.

\section{Conclusion}
This template demonstrates basic \LaTeX{} features including mathematical formulas, tables, and lists. The bibliography system is now set up to work seamlessly with Bibby AI's citation management. Feel free to modify it for your needs!

\bibliographystyle{plain}
\bibliography{cite}

\end{document}\documentclass{article}
\usepackage[utf8]{inputenc}
\usepackage{booktabs}
\usepackage{amsmath}
\usepackage{natbib}

\title{My Document}
\author{Your Name}
\date{\today}

\begin{document}

\maketitle

\section{Introduction}
This is a sample document demonstrating various \LaTeX{} features including mathematical equations, tables, and formatting. You can cite papers using the citation search feature, and they will be automatically added to your \texttt{cite.bib} file.

\section{Mathematical Expressions}
Here are some example equations:

\subsection{Quadratic Formula}
The quadratic formula is given by:
\[x = \frac{-b \pm \sqrt{b^2 - 4ac}}{2a}\]

\subsection{Matrix Example}
A sample matrix:
\[
\begin{pmatrix}
a_{11} & a_{12} & a_{13} \\
a_{21} & a_{22} & a_{23} \\
a_{31} & a_{32} & a_{33}
\end{pmatrix}
\]

\section{Tables}
Below is a sample table showing data:

\begin{table}[h]
\centering
\begin{tabular}{@{}lrrr@{}}
\toprule
\textbf{Category} & \textbf{Value 1} & \textbf{Value 2} & \textbf{Total} \\
\midrule
Item A & 10.5 & 20.3 & 30.8 \\
Item B & 15.2 & 25.1 & 40.3 \\
Item C & 22.1 & 12.7 & 34.8 \\
\bottomrule
\end{tabular}
\caption{Sample Data Table}
\label{tab:sample}
\end{table}

\section{Lists}
\subsection{Bullet Points}
\begin{itemize}
  \item First item
  \item Second item
  \item Third item with sub-items:
    \begin{itemize}
      \item Sub-item 1
      \item Sub-item 2
    \end{itemize}
\end{itemize}

\subsection{Numbered List}
\begin{enumerate}
  \item First step
  \item Second step
  \item Third step
\end{enumerate}

\section{Citations and References}
When you use the citation search feature, references will be automatically managed. For example, you might cite a paper like this: \cite{example2024}. All citations will be automatically added to your \texttt{cite.bib} file.

\section{Conclusion}
This template demonstrates basic \LaTeX{} features including mathematical formulas, tables, and lists. The bibliography system is now set up to work seamlessly with Bibby AI's citation management. Feel free to modify it for your needs!

\bibliographystyle{plain}
\bibliography{cite}

\end{document}\documentclass{article}
\usepackage[utf8]{inputenc}
\usepackage{booktabs}
\usepackage{amsmath}
\usepackage{natbib}

\title{My Document}
\author{Your Name}
\date{\today}

\begin{document}

\maketitle

\section{Introduction}
This is a sample document demonstrating various \LaTeX{} features including mathematical equations, tables, and formatting. You can cite papers using the citation search feature, and they will be automatically added to your \texttt{cite.bib} file.

\section{Mathematical Expressions}
Here are some example equations:

\subsection{Quadratic Formula}
The quadratic formula is given by:
\[x = \frac{-b \pm \sqrt{b^2 - 4ac}}{2a}\]

\subsection{Matrix Example}
A sample matrix:
\[
\begin{pmatrix}
a_{11} & a_{12} & a_{13} \\
a_{21} & a_{22} & a_{23} \\
a_{31} & a_{32} & a_{33}
\end{pmatrix}
\]

\section{Tables}
Below is a sample table showing data:

\begin{table}[h]
\centering
\begin{tabular}{@{}lrrr@{}}
\toprule
\textbf{Category} & \textbf{Value 1} & \textbf{Value 2} & \textbf{Total} \\
\midrule
Item A & 10.5 & 20.3 & 30.8 \\
Item B & 15.2 & 25.1 & 40.3 \\
Item C & 22.1 & 12.7 & 34.8 \\
\bottomrule
\end{tabular}
\caption{Sample Data Table}
\label{tab:sample}
\end{table}

\section{Lists}
\subsection{Bullet Points}
\begin{itemize}
  \item First item
  \item Second item
  \item Third item with sub-items:
    \begin{itemize}
      \item Sub-item 1
      \item Sub-item 2
    \end{itemize}
\end{itemize}

\subsection{Numbered List}
\begin{enumerate}
  \item First step
  \item Second step
  \item Third step
\end{enumerate}

\section{Citations and References}
When you use the citation search feature, references will be automatically managed. For example, you might cite a paper like this: \cite{example2024}. All citations will be automatically added to your \texttt{cite.bib} file.

\section{Conclusion}
This template demonstrates basic \LaTeX{} features including mathematical formulas, tables, and lists. The bibliography system is now set up to work seamlessly with Bibby AI's citation management. Feel free to modify it for your needs!

\bibliographystyle{plain}
\bibliography{cite}

\end{document}\documentclass{article}
\usepackage[utf8]{inputenc}
\usepackage{booktabs}
\usepackage{amsmath}
\usepackage{natbib}

\title{My Document}
\author{Your Name}
\date{\today}

\begin{document}

\maketitle

\section{Introduction}
This is a sample document demonstrating various \LaTeX{} features including mathematical equations, tables, and formatting. You can cite papers using the citation search feature, and they will be automatically added to your \texttt{cite.bib} file.

\section{Mathematical Expressions}
Here are some example equations:

\subsection{Quadratic Formula}
The quadratic formula is given by:
\[x = \frac{-b \pm \sqrt{b^2 - 4ac}}{2a}\]

\subsection{Matrix Example}
A sample matrix:
\[
\begin{pmatrix}
a_{11} & a_{12} & a_{13} \\
a_{21} & a_{22} & a_{23} \\
a_{31} & a_{32} & a_{33}
\end{pmatrix}
\]

\section{Tables}
Below is a sample table showing data:

\begin{table}[h]
\centering
\begin{tabular}{@{}lrrr@{}}
\toprule
\textbf{Category} & \textbf{Value 1} & \textbf{Value 2} & \textbf{Total} \\
\midrule
Item A & 10.5 & 20.3 & 30.8 \\
Item B & 15.2 & 25.1 & 40.3 \\
Item C & 22.1 & 12.7 & 34.8 \\
\bottomrule
\end{tabular}
\caption{Sample Data Table}
\label{tab:sample}
\end{table}

\section{Lists}
\subsection{Bullet Points}
\begin{itemize}
  \item First item
  \item Second item
  \item Third item with sub-items:
    \begin{itemize}
      \item Sub-item 1
      \item Sub-item 2
    \end{itemize}
\end{itemize}

\subsection{Numbered List}
\begin{enumerate}
  \item First step
  \item Second step
  \item Third step
\end{enumerate}

\section{Citations and References}
When you use the citation search feature, references will be automatically managed. For example, you might cite a paper like this: \cite{example2024}. All citations will be automatically added to your \texttt{cite.bib} file.

\section{Conclusion}
This template demonstrates basic \LaTeX{} features including mathematical formulas, tables, and lists. The bibliography system is now set up to work seamlessly with Bibby AI's citation management. Feel free to modify it for your needs!

\bibliographystyle{plain}
\bibliography{cite}

\end{document}\documentclass{article}
\usepackage[utf8]{inputenc}
\usepackage{booktabs}
\usepackage{amsmath}
\usepackage{natbib}

\title{My Document}
\author{Your Name}
\date{\today}

\begin{document}

\maketitle

\section{Introduction}
This is a sample document demonstrating various \LaTeX{} features including mathematical equations, tables, and formatting. You can cite papers using the citation search feature, and they will be automatically added to your \texttt{cite.bib} file.

\section{Mathematical Expressions}
Here are some example equations:

\subsection{Quadratic Formula}
The quadratic formula is given by:
\[x = \frac{-b \pm \sqrt{b^2 - 4ac}}{2a}\]

\subsection{Matrix Example}
A sample matrix:
\[
\begin{pmatrix}
a_{11} & a_{12} & a_{13} \\
a_{21} & a_{22} & a_{23} \\
a_{31} & a_{32} & a_{33}
\end{pmatrix}
\]

\section{Tables}
Below is a sample table showing data:

\begin{table}[h]
\centering
\begin{tabular}{@{}lrrr@{}}
\toprule
\textbf{Category} & \textbf{Value 1} & \textbf{Value 2} & \textbf{Total} \\
\midrule
Item A & 10.5 & 20.3 & 30.8 \\
Item B & 15.2 & 25.1 & 40.3 \\
Item C & 22.1 & 12.7 & 34.8 \\
\bottomrule
\end{tabular}
\caption{Sample Data Table}
\label{tab:sample}
\end{table}

\section{Lists}
\subsection{Bullet Points}
\begin{itemize}
  \item First item
  \item Second item
  \item Third item with sub-items:
    \begin{itemize}
      \item Sub-item 1
      \item Sub-item 2
    \end{itemize}
\end{itemize}

\subsection{Numbered List}
\begin{enumerate}
  \item First step
  \item Second step
  \item Third step
\end{enumerate}

\section{Citations and References}
When you use the citation search feature, references will be automatically managed. For example, you might cite a paper like this: \cite{example2024}. All citations will be automatically added to your \texttt{cite.bib} file.

\section{Conclusion}
This template demonstrates basic \LaTeX{} features including mathematical formulas, tables, and lists. The bibliography system is now set up to work seamlessly with Bibby AI's citation management. Feel free to modify it for your needs!

\bibliographystyle{plain}
\bibliography{cite}

\end{document}\documentclass{article}
\usepackage[utf8]{inputenc}
\usepackage{booktabs}
\usepackage{amsmath}
\usepackage{natbib}

\title{My Document}
\author{Your Name}
\date{\today}

\begin{document}

\maketitle

\section{Introduction}
This is a sample document demonstrating various \LaTeX{} features including mathematical equations, tables, and formatting. You can cite papers using the citation search feature, and they will be automatically added to your \texttt{cite.bib} file.

\section{Mathematical Expressions}
Here are some example equations:

\subsection{Quadratic Formula}
The quadratic formula is given by:
\[x = \frac{-b \pm \sqrt{b^2 - 4ac}}{2a}\]

\subsection{Matrix Example}
A sample matrix:
\[
\begin{pmatrix}
a_{11} & a_{12} & a_{13} \\
a_{21} & a_{22} & a_{23} \\
a_{31} & a_{32} & a_{33}
\end{pmatrix}
\]

\section{Tables}
Below is a sample table showing data:

\begin{table}[h]
\centering
\begin{tabular}{@{}lrrr@{}}
\toprule
\textbf{Category} & \textbf{Value 1} & \textbf{Value 2} & \textbf{Total} \\
\midrule
Item A & 10.5 & 20.3 & 30.8 \\
Item B & 15.2 & 25.1 & 40.3 \\
Item C & 22.1 & 12.7 & 34.8 \\
\bottomrule
\end{tabular}
\caption{Sample Data Table}
\label{tab:sample}
\end{table}

\section{Lists}
\subsection{Bullet Points}
\begin{itemize}
  \item First item
  \item Second item
  \item Third item with sub-items:
    \begin{itemize}
      \item Sub-item 1
      \item Sub-item 2
    \end{itemize}
\end{itemize}

\subsection{Numbered List}
\begin{enumerate}
  \item First step
  \item Second step
  \item Third step
\end{enumerate}

\section{Citations and References}
When you use the citation search feature, references will be automatically managed. For example, you might cite a paper like this: \cite{example2024}. All citations will be automatically added to your \texttt{cite.bib} file.

\section{Conclusion}
This template demonstrates basic \LaTeX{} features including mathematical formulas, tables, and lists. The bibliography system is now set up to work seamlessly with Bibby AI's citation management. Feel free to modify it for your needs!

\bibliographystyle{plain}
\bibliography{cite}

\end{document}\documentclass{article}
\usepackage[utf8]{inputenc}
\usepackage{booktabs}
\usepackage{amsmath}
\usepackage{natbib}

\title{My Document}
\author{Your Name}
\date{\today}

\begin{document}

\maketitle

\section{Introduction}
This is a sample document demonstrating various \LaTeX{} features including mathematical equations, tables, and formatting. You can cite papers using the citation search feature, and they will be automatically added to your \texttt{cite.bib} file.

\section{Mathematical Expressions}
Here are some example equations:

\subsection{Quadratic Formula}
The quadratic formula is given by:
\[x = \frac{-b \pm \sqrt{b^2 - 4ac}}{2a}\]

\subsection{Matrix Example}
A sample matrix:
\[
\begin{pmatrix}
a_{11} & a_{12} & a_{13} \\
a_{21} & a_{22} & a_{23} \\
a_{31} & a_{32} & a_{33}
\end{pmatrix}
\]

\section{Tables}
Below is a sample table showing data:

\begin{table}[h]
\centering
\begin{tabular}{@{}lrrr@{}}
\toprule
\textbf{Category} & \textbf{Value 1} & \textbf{Value 2} & \textbf{Total} \\
\midrule
Item A & 10.5 & 20.3 & 30.8 \\
Item B & 15.2 & 25.1 & 40.3 \\
Item C & 22.1 & 12.7 & 34.8 \\
\bottomrule
\end{tabular}
\caption{Sample Data Table}
\label{tab:sample}
\end{table}

\section{Lists}
\subsection{Bullet Points}
\begin{itemize}
  \item First item
  \item Second item
  \item Third item with sub-items:
    \begin{itemize}
      \item Sub-item 1
      \item Sub-item 2
    \end{itemize}
\end{itemize}

\subsection{Numbered List}
\begin{enumerate}
  \item First step
  \item Second step
  \item Third step
\end{enumerate}

\section{Citations and References}
When you use the citation search feature, references will be automatically managed. For example, you might cite a paper like this: \cite{example2024}. All citations will be automatically added to your \texttt{cite.bib} file.

\section{Conclusion}
This template demonstrates basic \LaTeX{} features including mathematical formulas, tables, and lists. The bibliography system is now set up to work seamlessly with Bibby AI's citation management. Feel free to modify it for your needs!

\bibliographystyle{plain}
\bibliography{cite}

\end{document}\documentclass{article}
\usepackage[utf8]{inputenc}
\usepackage{booktabs}
\usepackage{amsmath}
\usepackage{natbib}

\title{My Document}
\author{Your Name}
\date{\today}

\begin{document}

\maketitle

\section{Introduction}
This is a sample document demonstrating various \LaTeX{} features including mathematical equations, tables, and formatting. You can cite papers using the citation search feature, and they will be automatically added to your \texttt{cite.bib} file.

\section{Mathematical Expressions}
Here are some example equations:

\subsection{Quadratic Formula}
The quadratic formula is given by:
\[x = \frac{-b \pm \sqrt{b^2 - 4ac}}{2a}\]

\subsection{Matrix Example}
A sample matrix:
\[
\begin{pmatrix}
a_{11} & a_{12} & a_{13} \\
a_{21} & a_{22} & a_{23} \\
a_{31} & a_{32} & a_{33}
\end{pmatrix}
\]

\section{Tables}
Below is a sample table showing data:

\begin{table}[h]
\centering
\begin{tabular}{@{}lrrr@{}}
\toprule
\textbf{Category} & \textbf{Value 1} & \textbf{Value 2} & \textbf{Total} \\
\midrule
Item A & 10.5 & 20.3 & 30.8 \\
Item B & 15.2 & 25.1 & 40.3 \\
Item C & 22.1 & 12.7 & 34.8 \\
\bottomrule
\end{tabular}
\caption{Sample Data Table}
\label{tab:sample}
\end{table}

\section{Lists}
\subsection{Bullet Points}
\begin{itemize}
  \item First item
  \item Second item
  \item Third item with sub-items:
    \begin{itemize}
      \item Sub-item 1
      \item Sub-item 2
    \end{itemize}
\end{itemize}

\subsection{Numbered List}
\begin{enumerate}
  \item First step
  \item Second step
  \item Third step
\end{enumerate}

\section{Citations and References}
When you use the citation search feature, references will be automatically managed. For example, you might cite a paper like this: \cite{example2024}. All citations will be automatically added to your \texttt{cite.bib} file.

\section{Conclusion}
This template demonstrates basic \LaTeX{} features including mathematical formulas, tables, and lists. The bibliography system is now set up to work seamlessly with Bibby AI's citation management. Feel free to modify it for your needs!

\bibliographystyle{plain}
\bibliography{cite}

\end{document}\documentclass{article}
\usepackage[utf8]{inputenc}
\usepackage{booktabs}
\usepackage{amsmath}
\usepackage{natbib}

\title{My Document}
\author{Your Name}
\date{\today}

\begin{document}

\maketitle

\section{Introduction}
This is a sample document demonstrating various \LaTeX{} features including mathematical equations, tables, and formatting. You can cite papers using the citation search feature, and they will be automatically added to your \texttt{cite.bib} file.

\section{Mathematical Expressions}
Here are some example equations:

\subsection{Quadratic Formula}
The quadratic formula is given by:
\[x = \frac{-b \pm \sqrt{b^2 - 4ac}}{2a}\]

\subsection{Matrix Example}
A sample matrix:
\[
\begin{pmatrix}
a_{11} & a_{12} & a_{13} \\
a_{21} & a_{22} & a_{23} \\
a_{31} & a_{32} & a_{33}
\end{pmatrix}
\]

\section{Tables}
Below is a sample table showing data:

\begin{table}[h]
\centering
\begin{tabular}{@{}lrrr@{}}
\toprule
\textbf{Category} & \textbf{Value 1} & \textbf{Value 2} & \textbf{Total} \\
\midrule
Item A & 10.5 & 20.3 & 30.8 \\
Item B & 15.2 & 25.1 & 40.3 \\
Item C & 22.1 & 12.7 & 34.8 \\
\bottomrule
\end{tabular}
\caption{Sample Data Table}
\label{tab:sample}
\end{table}

\section{Lists}
\subsection{Bullet Points}
\begin{itemize}
  \item First item
  \item Second item
  \item Third item with sub-items:
    \begin{itemize}
      \item Sub-item 1
      \item Sub-item 2
    \end{itemize}
\end{itemize}

\subsection{Numbered List}
\begin{enumerate}
  \item First step
  \item Second step
  \item Third step
\end{enumerate}

\section{Citations and References}
When you use the citation search feature, references will be automatically managed. For example, you might cite a paper like this: \cite{example2024}. All citations will be automatically added to your \texttt{cite.bib} file.

\section{Conclusion}
This template demonstrates basic \LaTeX{} features including mathematical formulas, tables, and lists. The bibliography system is now set up to work seamlessly with Bibby AI's citation management. Feel free to modify it for your needs!

\bibliographystyle{plain}
\bibliography{cite}

\end{document}\documentclass{article}
\usepackage[utf8]{inputenc}
\usepackage{booktabs}
\usepackage{amsmath}
\usepackage{natbib}

\title{My Document}
\author{Your Name}
\date{\today}

\begin{document}

\maketitle

\section{Introduction}
This is a sample document demonstrating various \LaTeX{} features including mathematical equations, tables, and formatting. You can cite papers using the citation search feature, and they will be automatically added to your \texttt{cite.bib} file.

\section{Mathematical Expressions}
Here are some example equations:

\subsection{Quadratic Formula}
The quadratic formula is given by:
\[x = \frac{-b \pm \sqrt{b^2 - 4ac}}{2a}\]

\subsection{Matrix Example}
A sample matrix:
\[
\begin{pmatrix}
a_{11} & a_{12} & a_{13} \\
a_{21} & a_{22} & a_{23} \\
a_{31} & a_{32} & a_{33}
\end{pmatrix}
\]

\section{Tables}
Below is a sample table showing data:

\begin{table}[h]
\centering
\begin{tabular}{@{}lrrr@{}}
\toprule
\textbf{Category} & \textbf{Value 1} & \textbf{Value 2} & \textbf{Total} \\
\midrule
Item A & 10.5 & 20.3 & 30.8 \\
Item B & 15.2 & 25.1 & 40.3 \\
Item C & 22.1 & 12.7 & 34.8 \\
\bottomrule
\end{tabular}
\caption{Sample Data Table}
\label{tab:sample}
\end{table}

\section{Lists}
\subsection{Bullet Points}
\begin{itemize}
  \item First item
  \item Second item
  \item Third item with sub-items:
    \begin{itemize}
      \item Sub-item 1
      \item Sub-item 2
    \end{itemize}
\end{itemize}

\subsection{Numbered List}
\begin{enumerate}
  \item First step
  \item Second step
  \item Third step
\end{enumerate}

\section{Citations and References}
When you use the citation search feature, references will be automatically managed. For example, you might cite a paper like this: \cite{example2024}. All citations will be automatically added to your \texttt{cite.bib} file.

\section{Conclusion}
This template demonstrates basic \LaTeX{} features including mathematical formulas, tables, and lists. The bibliography system is now set up to work seamlessly with Bibby AI's citation management. Feel free to modify it for your needs!

\bibliographystyle{plain}
\bibliography{cite}

\end{document}\documentclass{article}
\usepackage[utf8]{inputenc}
\usepackage{booktabs}
\usepackage{amsmath}
\usepackage{natbib}

\title{My Document}
\author{Your Name}
\date{\today}

\begin{document}

\maketitle

\section{Introduction}
This is a sample document demonstrating various \LaTeX{} features including mathematical equations, tables, and formatting. You can cite papers using the citation search feature, and they will be automatically added to your \texttt{cite.bib} file.

\section{Mathematical Expressions}
Here are some example equations:

\subsection{Quadratic Formula}
The quadratic formula is given by:
\[x = \frac{-b \pm \sqrt{b^2 - 4ac}}{2a}\]

\subsection{Matrix Example}
A sample matrix:
\[
\begin{pmatrix}
a_{11} & a_{12} & a_{13} \\
a_{21} & a_{22} & a_{23} \\
a_{31} & a_{32} & a_{33}
\end{pmatrix}
\]

\section{Tables}
Below is a sample table showing data:

\begin{table}[h]
\centering
\begin{tabular}{@{}lrrr@{}}
\toprule
\textbf{Category} & \textbf{Value 1} & \textbf{Value 2} & \textbf{Total} \\
\midrule
Item A & 10.5 & 20.3 & 30.8 \\
Item B & 15.2 & 25.1 & 40.3 \\
Item C & 22.1 & 12.7 & 34.8 \\
\bottomrule
\end{tabular}
\caption{Sample Data Table}
\label{tab:sample}
\end{table}

\section{Lists}
\subsection{Bullet Points}
\begin{itemize}
  \item First item
  \item Second item
  \item Third item with sub-items:
    \begin{itemize}
      \item Sub-item 1
      \item Sub-item 2
    \end{itemize}
\end{itemize}

\subsection{Numbered List}
\begin{enumerate}
  \item First step
  \item Second step
  \item Third step
\end{enumerate}

\section{Citations and References}
When you use the citation search feature, references will be automatically managed. For example, you might cite a paper like this: \cite{example2024}. All citations will be automatically added to your \texttt{cite.bib} file.

\section{Conclusion}
This template demonstrates basic \LaTeX{} features including mathematical formulas, tables, and lists. The bibliography system is now set up to work seamlessly with Bibby AI's citation management. Feel free to modify it for your needs!

\bibliographystyle{plain}
\bibliography{cite}

\end{document}\documentclass{article}
\usepackage[utf8]{inputenc}
\usepackage{booktabs}
\usepackage{amsmath}
\usepackage{natbib}

\title{My Document}
\author{Your Name}
\date{\today}

\begin{document}

\maketitle

\section{Introduction}
This is a sample document demonstrating various \LaTeX{} features including mathematical equations, tables, and formatting. You can cite papers using the citation search feature, and they will be automatically added to your \texttt{cite.bib} file.

\section{Mathematical Expressions}
Here are some example equations:

\subsection{Quadratic Formula}
The quadratic formula is given by:
\[x = \frac{-b \pm \sqrt{b^2 - 4ac}}{2a}\]

\subsection{Matrix Example}
A sample matrix:
\[
\begin{pmatrix}
a_{11} & a_{12} & a_{13} \\
a_{21} & a_{22} & a_{23} \\
a_{31} & a_{32} & a_{33}
\end{pmatrix}
\]

\section{Tables}
Below is a sample table showing data:

\begin{table}[h]
\centering
\begin{tabular}{@{}lrrr@{}}
\toprule
\textbf{Category} & \textbf{Value 1} & \textbf{Value 2} & \textbf{Total} \\
\midrule
Item A & 10.5 & 20.3 & 30.8 \\
Item B & 15.2 & 25.1 & 40.3 \\
Item C & 22.1 & 12.7 & 34.8 \\
\bottomrule
\end{tabular}
\caption{Sample Data Table}
\label{tab:sample}
\end{table}

\section{Lists}
\subsection{Bullet Points}
\begin{itemize}
  \item First item
  \item Second item
  \item Third item with sub-items:
    \begin{itemize}
      \item Sub-item 1
      \item Sub-item 2
    \end{itemize}
\end{itemize}

\subsection{Numbered List}
\begin{enumerate}
  \item First step
  \item Second step
  \item Third step
\end{enumerate}

\section{Citations and References}
When you use the citation search feature, references will be automatically managed. For example, you might cite a paper like this: \cite{example2024}. All citations will be automatically added to your \texttt{cite.bib} file.

\section{Conclusion}
This template demonstrates basic \LaTeX{} features including mathematical formulas, tables, and lists. The bibliography system is now set up to work seamlessly with Bibby AI's citation management. Feel free to modify it for your needs!

\bibliographystyle{plain}
\bibliography{cite}

\end{document}\documentclass{article}
\usepackage[utf8]{inputenc}
\usepackage{booktabs}
\usepackage{amsmath}
\usepackage{natbib}

\title{My Document}
\author{Your Name}
\date{\today}

\begin{document}

\maketitle

\section{Introduction}
This is a sample document demonstrating various \LaTeX{} features including mathematical equations, tables, and formatting. You can cite papers using the citation search feature, and they will be automatically added to your \texttt{cite.bib} file.

\section{Mathematical Expressions}
Here are some example equations:

\subsection{Quadratic Formula}
The quadratic formula is given by:
\[x = \frac{-b \pm \sqrt{b^2 - 4ac}}{2a}\]

\subsection{Matrix Example}
A sample matrix:
\[
\begin{pmatrix}
a_{11} & a_{12} & a_{13} \\
a_{21} & a_{22} & a_{23} \\
a_{31} & a_{32} & a_{33}
\end{pmatrix}
\]

\section{Tables}
Below is a sample table showing data:

\begin{table}[h]
\centering
\begin{tabular}{@{}lrrr@{}}
\toprule
\textbf{Category} & \textbf{Value 1} & \textbf{Value 2} & \textbf{Total} \\
\midrule
Item A & 10.5 & 20.3 & 30.8 \\
Item B & 15.2 & 25.1 & 40.3 \\
Item C & 22.1 & 12.7 & 34.8 \\
\bottomrule
\end{tabular}
\caption{Sample Data Table}
\label{tab:sample}
\end{table}

\section{Lists}
\subsection{Bullet Points}
\begin{itemize}
  \item First item
  \item Second item
  \item Third item with sub-items:
    \begin{itemize}
      \item Sub-item 1
      \item Sub-item 2
    \end{itemize}
\end{itemize}

\subsection{Numbered List}
\begin{enumerate}
  \item First step
  \item Second step
  \item Third step
\end{enumerate}

\section{Citations and References}
When you use the citation search feature, references will be automatically managed. For example, you might cite a paper like this: \cite{example2024}. All citations will be automatically added to your \texttt{cite.bib} file.

\section{Conclusion}
This template demonstrates basic \LaTeX{} features including mathematical formulas, tables, and lists. The bibliography system is now set up to work seamlessly with Bibby AI's citation management. Feel free to modify it for your needs!

\bibliographystyle{plain}
\bibliography{cite}

\end{document}\documentclass{article}
\usepackage[utf8]{inputenc}
\usepackage{booktabs}
\usepackage{amsmath}
\usepackage{natbib}

\title{My Document}
\author{Your Name}
\date{\today}

\begin{document}

\maketitle

\section{Introduction}
This is a sample document demonstrating various \LaTeX{} features including mathematical equations, tables, and formatting. You can cite papers using the citation search feature, and they will be automatically added to your \texttt{cite.bib} file.

\section{Mathematical Expressions}
Here are some example equations:

\subsection{Quadratic Formula}
The quadratic formula is given by:
\[x = \frac{-b \pm \sqrt{b^2 - 4ac}}{2a}\]

\subsection{Matrix Example}
A sample matrix:
\[
\begin{pmatrix}
a_{11} & a_{12} & a_{13} \\
a_{21} & a_{22} & a_{23} \\
a_{31} & a_{32} & a_{33}
\end{pmatrix}
\]

\section{Tables}
Below is a sample table showing data:

\begin{table}[h]
\centering
\begin{tabular}{@{}lrrr@{}}
\toprule
\textbf{Category} & \textbf{Value 1} & \textbf{Value 2} & \textbf{Total} \\
\midrule
Item A & 10.5 & 20.3 & 30.8 \\
Item B & 15.2 & 25.1 & 40.3 \\
Item C & 22.1 & 12.7 & 34.8 \\
\bottomrule
\end{tabular}
\caption{Sample Data Table}
\label{tab:sample}
\end{table}

\section{Lists}
\subsection{Bullet Points}
\begin{itemize}
  \item First item
  \item Second item
  \item Third item with sub-items:
    \begin{itemize}
      \item Sub-item 1
      \item Sub-item 2
    \end{itemize}
\end{itemize}

\subsection{Numbered List}
\begin{enumerate}
  \item First step
  \item Second step
  \item Third step
\end{enumerate}

\section{Citations and References}
When you use the citation search feature, references will be automatically managed. For example, you might cite a paper like this: \cite{example2024}. All citations will be automatically added to your \texttt{cite.bib} file.

\section{Conclusion}
This template demonstrates basic \LaTeX{} features including mathematical formulas, tables, and lists. The bibliography system is now set up to work seamlessly with Bibby AI's citation management. Feel free to modify it for your needs!

\bibliographystyle{plain}
\bibliography{cite}

\end{document}\documentclass{article}
\usepackage[utf8]{inputenc}
\usepackage{booktabs}
\usepackage{amsmath}
\usepackage{natbib}

\title{My Document}
\author{Your Name}
\date{\today}

\begin{document}

\maketitle

\section{Introduction}
This is a sample document demonstrating various \LaTeX{} features including mathematical equations, tables, and formatting. You can cite papers using the citation search feature, and they will be automatically added to your \texttt{cite.bib} file.

\section{Mathematical Expressions}
Here are some example equations:

\subsection{Quadratic Formula}
The quadratic formula is given by:
\[x = \frac{-b \pm \sqrt{b^2 - 4ac}}{2a}\]

\subsection{Matrix Example}
A sample matrix:
\[
\begin{pmatrix}
a_{11} & a_{12} & a_{13} \\
a_{21} & a_{22} & a_{23} \\
a_{31} & a_{32} & a_{33}
\end{pmatrix}
\]

\section{Tables}
Below is a sample table showing data:

\begin{table}[h]
\centering
\begin{tabular}{@{}lrrr@{}}
\toprule
\textbf{Category} & \textbf{Value 1} & \textbf{Value 2} & \textbf{Total} \\
\midrule
Item A & 10.5 & 20.3 & 30.8 \\
Item B & 15.2 & 25.1 & 40.3 \\
Item C & 22.1 & 12.7 & 34.8 \\
\bottomrule
\end{tabular}
\caption{Sample Data Table}
\label{tab:sample}
\end{table}

\section{Lists}
\subsection{Bullet Points}
\begin{itemize}
  \item First item
  \item Second item
  \item Third item with sub-items:
    \begin{itemize}
      \item Sub-item 1
      \item Sub-item 2
    \end{itemize}
\end{itemize}

\subsection{Numbered List}
\begin{enumerate}
  \item First step
  \item Second step
  \item Third step
\end{enumerate}

\section{Citations and References}
When you use the citation search feature, references will be automatically managed. For example, you might cite a paper like this: \cite{example2024}. All citations will be automatically added to your \texttt{cite.bib} file.

\section{Conclusion}
This template demonstrates basic \LaTeX{} features including mathematical formulas, tables, and lists. The bibliography system is now set up to work seamlessly with Bibby AI's citation management. Feel free to modify it for your needs!

\bibliographystyle{plain}
\bibliography{cite}

\end{document}\documentclass{article}
\usepackage[utf8]{inputenc}
\usepackage{booktabs}
\usepackage{amsmath}
\usepackage{natbib}

\title{My Document}
\author{Your Name}
\date{\today}

\begin{document}

\maketitle

\section{Introduction}
This is a sample document demonstrating various \LaTeX{} features including mathematical equations, tables, and formatting. You can cite papers using the citation search feature, and they will be automatically added to your \texttt{cite.bib} file.

\section{Mathematical Expressions}
Here are some example equations:

\subsection{Quadratic Formula}
The quadratic formula is given by:
\[x = \frac{-b \pm \sqrt{b^2 - 4ac}}{2a}\]

\subsection{Matrix Example}
A sample matrix:
\[
\begin{pmatrix}
a_{11} & a_{12} & a_{13} \\
a_{21} & a_{22} & a_{23} \\
a_{31} & a_{32} & a_{33}
\end{pmatrix}
\]

\section{Tables}
Below is a sample table showing data:

\begin{table}[h]
\centering
\begin{tabular}{@{}lrrr@{}}
\toprule
\textbf{Category} & \textbf{Value 1} & \textbf{Value 2} & \textbf{Total} \\
\midrule
Item A & 10.5 & 20.3 & 30.8 \\
Item B & 15.2 & 25.1 & 40.3 \\
Item C & 22.1 & 12.7 & 34.8 \\
\bottomrule
\end{tabular}
\caption{Sample Data Table}
\label{tab:sample}
\end{table}

\section{Lists}
\subsection{Bullet Points}
\begin{itemize}
  \item First item
  \item Second item
  \item Third item with sub-items:
    \begin{itemize}
      \item Sub-item 1
      \item Sub-item 2
    \end{itemize}
\end{itemize}

\subsection{Numbered List}
\begin{enumerate}
  \item First step
  \item Second step
  \item Third step
\end{enumerate}

\section{Citations and References}
When you use the citation search feature, references will be automatically managed. For example, you might cite a paper like this: \cite{example2024}. All citations will be automatically added to your \texttt{cite.bib} file.

\section{Conclusion}
This template demonstrates basic \LaTeX{} features including mathematical formulas, tables, and lists. The bibliography system is now set up to work seamlessly with Bibby AI's citation management. Feel free to modify it for your needs!

\bibliographystyle{plain}
\bibliography{cite}

\end{document}\documentclass{article}
\usepackage[utf8]{inputenc}
\usepackage{booktabs}
\usepackage{amsmath}
\usepackage{natbib}

\title{My Document}
\author{Your Name}
\date{\today}

\begin{document}

\maketitle

\section{Introduction}
This is a sample document demonstrating various \LaTeX{} features including mathematical equations, tables, and formatting. You can cite papers using the citation search feature, and they will be automatically added to your \texttt{cite.bib} file.

\section{Mathematical Expressions}
Here are some example equations:

\subsection{Quadratic Formula}
The quadratic formula is given by:
\[x = \frac{-b \pm \sqrt{b^2 - 4ac}}{2a}\]

\subsection{Matrix Example}
A sample matrix:
\[
\begin{pmatrix}
a_{11} & a_{12} & a_{13} \\
a_{21} & a_{22} & a_{23} \\
a_{31} & a_{32} & a_{33}
\end{pmatrix}
\]

\section{Tables}
Below is a sample table showing data:

\begin{table}[h]
\centering
\begin{tabular}{@{}lrrr@{}}
\toprule
\textbf{Category} & \textbf{Value 1} & \textbf{Value 2} & \textbf{Total} \\
\midrule
Item A & 10.5 & 20.3 & 30.8 \\
Item B & 15.2 & 25.1 & 40.3 \\
Item C & 22.1 & 12.7 & 34.8 \\
\bottomrule
\end{tabular}
\caption{Sample Data Table}
\label{tab:sample}
\end{table}

\section{Lists}
\subsection{Bullet Points}
\begin{itemize}
  \item First item
  \item Second item
  \item Third item with sub-items:
    \begin{itemize}
      \item Sub-item 1
      \item Sub-item 2
    \end{itemize}
\end{itemize}

\subsection{Numbered List}
\begin{enumerate}
  \item First step
  \item Second step
  \item Third step
\end{enumerate}

\section{Citations and References}
When you use the citation search feature, references will be automatically managed. For example, you might cite a paper like this: \cite{example2024}. All citations will be automatically added to your \texttt{cite.bib} file.

\section{Conclusion}
This template demonstrates basic \LaTeX{} features including mathematical formulas, tables, and lists. The bibliography system is now set up to work seamlessly with Bibby AI's citation management. Feel free to modify it for your needs!

\bibliographystyle{plain}
\bibliography{cite}

\end{document}\documentclass{article}
\usepackage[utf8]{inputenc}
\usepackage{booktabs}
\usepackage{amsmath}
\usepackage{natbib}

\title{My Document}
\author{Your Name}
\date{\today}

\begin{document}

\maketitle

\section{Introduction}
This is a sample document demonstrating various \LaTeX{} features including mathematical equations, tables, and formatting. You can cite papers using the citation search feature, and they will be automatically added to your \texttt{cite.bib} file.

\section{Mathematical Expressions}
Here are some example equations:

\subsection{Quadratic Formula}
The quadratic formula is given by:
\[x = \frac{-b \pm \sqrt{b^2 - 4ac}}{2a}\]

\subsection{Matrix Example}
A sample matrix:
\[
\begin{pmatrix}
a_{11} & a_{12} & a_{13} \\
a_{21} & a_{22} & a_{23} \\
a_{31} & a_{32} & a_{33}
\end{pmatrix}
\]

\section{Tables}
Below is a sample table showing data:

\begin{table}[h]
\centering
\begin{tabular}{@{}lrrr@{}}
\toprule
\textbf{Category} & \textbf{Value 1} & \textbf{Value 2} & \textbf{Total} \\
\midrule
Item A & 10.5 & 20.3 & 30.8 \\
Item B & 15.2 & 25.1 & 40.3 \\
Item C & 22.1 & 12.7 & 34.8 \\
\bottomrule
\end{tabular}
\caption{Sample Data Table}
\label{tab:sample}
\end{table}

\section{Lists}
\subsection{Bullet Points}
\begin{itemize}
  \item First item
  \item Second item
  \item Third item with sub-items:
    \begin{itemize}
      \item Sub-item 1
      \item Sub-item 2
    \end{itemize}
\end{itemize}

\subsection{Numbered List}
\begin{enumerate}
  \item First step
  \item Second step
  \item Third step
\end{enumerate}

\section{Citations and References}
When you use the citation search feature, references will be automatically managed. For example, you might cite a paper like this: \cite{example2024}. All citations will be automatically added to your \texttt{cite.bib} file.

\section{Conclusion}
This template demonstrates basic \LaTeX{} features including mathematical formulas, tables, and lists. The bibliography system is now set up to work seamlessly with Bibby AI's citation management. Feel free to modify it for your needs!

\bibliographystyle{plain}
\bibliography{cite}

\end{document}\documentclass{article}
\usepackage[utf8]{inputenc}
\usepackage{booktabs}
\usepackage{amsmath}
\usepackage{natbib}

\title{My Document}
\author{Your Name}
\date{\today}

\begin{document}

\maketitle

\section{Introduction}
This is a sample document demonstrating various \LaTeX{} features including mathematical equations, tables, and formatting. You can cite papers using the citation search feature, and they will be automatically added to your \texttt{cite.bib} file.

\section{Mathematical Expressions}
Here are some example equations:

\subsection{Quadratic Formula}
The quadratic formula is given by:
\[x = \frac{-b \pm \sqrt{b^2 - 4ac}}{2a}\]

\subsection{Matrix Example}
A sample matrix:
\[
\begin{pmatrix}
a_{11} & a_{12} & a_{13} \\
a_{21} & a_{22} & a_{23} \\
a_{31} & a_{32} & a_{33}
\end{pmatrix}
\]

\section{Tables}
Below is a sample table showing data:

\begin{table}[h]
\centering
\begin{tabular}{@{}lrrr@{}}
\toprule
\textbf{Category} & \textbf{Value 1} & \textbf{Value 2} & \textbf{Total} \\
\midrule
Item A & 10.5 & 20.3 & 30.8 \\
Item B & 15.2 & 25.1 & 40.3 \\
Item C & 22.1 & 12.7 & 34.8 \\
\bottomrule
\end{tabular}
\caption{Sample Data Table}
\label{tab:sample}
\end{table}

\section{Lists}
\subsection{Bullet Points}
\begin{itemize}
  \item First item
  \item Second item
  \item Third item with sub-items:
    \begin{itemize}
      \item Sub-item 1
      \item Sub-item 2
    \end{itemize}
\end{itemize}

\subsection{Numbered List}
\begin{enumerate}
  \item First step
  \item Second step
  \item Third step
\end{enumerate}

\section{Citations and References}
When you use the citation search feature, references will be automatically managed. For example, you might cite a paper like this: \cite{example2024}. All citations will be automatically added to your \texttt{cite.bib} file.

\section{Conclusion}
This template demonstrates basic \LaTeX{} features including mathematical formulas, tables, and lists. The bibliography system is now set up to work seamlessly with Bibby AI's citation management. Feel free to modify it for your needs!

\bibliographystyle{plain}
\bibliography{cite}

\end{document}\documentclass{article}
\usepackage[utf8]{inputenc}
\usepackage{booktabs}
\usepackage{amsmath}
\usepackage{natbib}

\title{My Document}
\author{Your Name}
\date{\today}

\begin{document}

\maketitle

\section{Introduction}
This is a sample document demonstrating various \LaTeX{} features including mathematical equations, tables, and formatting. You can cite papers using the citation search feature, and they will be automatically added to your \texttt{cite.bib} file.

\section{Mathematical Expressions}
Here are some example equations:

\subsection{Quadratic Formula}
The quadratic formula is given by:
\[x = \frac{-b \pm \sqrt{b^2 - 4ac}}{2a}\]

\subsection{Matrix Example}
A sample matrix:
\[
\begin{pmatrix}
a_{11} & a_{12} & a_{13} \\
a_{21} & a_{22} & a_{23} \\
a_{31} & a_{32} & a_{33}
\end{pmatrix}
\]

\section{Tables}
Below is a sample table showing data:

\begin{table}[h]
\centering
\begin{tabular}{@{}lrrr@{}}
\toprule
\textbf{Category} & \textbf{Value 1} & \textbf{Value 2} & \textbf{Total} \\
\midrule
Item A & 10.5 & 20.3 & 30.8 \\
Item B & 15.2 & 25.1 & 40.3 \\
Item C & 22.1 & 12.7 & 34.8 \\
\bottomrule
\end{tabular}
\caption{Sample Data Table}
\label{tab:sample}
\end{table}

\section{Lists}
\subsection{Bullet Points}
\begin{itemize}
  \item First item
  \item Second item
  \item Third item with sub-items:
    \begin{itemize}
      \item Sub-item 1
      \item Sub-item 2
    \end{itemize}
\end{itemize}

\subsection{Numbered List}
\begin{enumerate}
  \item First step
  \item Second step
  \item Third step
\end{enumerate}

\section{Citations and References}
When you use the citation search feature, references will be automatically managed. For example, you might cite a paper like this: \cite{example2024}. All citations will be automatically added to your \texttt{cite.bib} file.

\section{Conclusion}
This template demonstrates basic \LaTeX{} features including mathematical formulas, tables, and lists. The bibliography system is now set up to work seamlessly with Bibby AI's citation management. Feel free to modify it for your needs!

\bibliographystyle{plain}
\bibliography{cite}

\end{document}\documentclass{article}
\usepackage[utf8]{inputenc}
\usepackage{booktabs}
\usepackage{amsmath}
\usepackage{natbib}

\title{My Document}
\author{Your Name}
\date{\today}

\begin{document}

\maketitle

\section{Introduction}
This is a sample document demonstrating various \LaTeX{} features including mathematical equations, tables, and formatting. You can cite papers using the citation search feature, and they will be automatically added to your \texttt{cite.bib} file.

\section{Mathematical Expressions}
Here are some example equations:

\subsection{Quadratic Formula}
The quadratic formula is given by:
\[x = \frac{-b \pm \sqrt{b^2 - 4ac}}{2a}\]

\subsection{Matrix Example}
A sample matrix:
\[
\begin{pmatrix}
a_{11} & a_{12} & a_{13} \\
a_{21} & a_{22} & a_{23} \\
a_{31} & a_{32} & a_{33}
\end{pmatrix}
\]

\section{Tables}
Below is a sample table showing data:

\begin{table}[h]
\centering
\begin{tabular}{@{}lrrr@{}}
\toprule
\textbf{Category} & \textbf{Value 1} & \textbf{Value 2} & \textbf{Total} \\
\midrule
Item A & 10.5 & 20.3 & 30.8 \\
Item B & 15.2 & 25.1 & 40.3 \\
Item C & 22.1 & 12.7 & 34.8 \\
\bottomrule
\end{tabular}
\caption{Sample Data Table}
\label{tab:sample}
\end{table}

\section{Lists}
\subsection{Bullet Points}
\begin{itemize}
  \item First item
  \item Second item
  \item Third item with sub-items:
    \begin{itemize}
      \item Sub-item 1
      \item Sub-item 2
    \end{itemize}
\end{itemize}

\subsection{Numbered List}
\begin{enumerate}
  \item First step
  \item Second step
  \item Third step
\end{enumerate}

\section{Citations and References}
When you use the citation search feature, references will be automatically managed. For example, you might cite a paper like this: \cite{example2024}. All citations will be automatically added to your \texttt{cite.bib} file.

\section{Conclusion}
This template demonstrates basic \LaTeX{} features including mathematical formulas, tables, and lists. The bibliography system is now set up to work seamlessly with Bibby AI's citation management. Feel free to modify it for your needs!

\bibliographystyle{plain}
\bibliography{cite}

\end{document}\documentclass{article}
\usepackage[utf8]{inputenc}
\usepackage{booktabs}
\usepackage{amsmath}
\usepackage{natbib}

\title{My Document}
\author{Your Name}
\date{\today}

\begin{document}

\maketitle

\section{Introduction}
This is a sample document demonstrating various \LaTeX{} features including mathematical equations, tables, and formatting. You can cite papers using the citation search feature, and they will be automatically added to your \texttt{cite.bib} file.

\section{Mathematical Expressions}
Here are some example equations:

\subsection{Quadratic Formula}
The quadratic formula is given by:
\[x = \frac{-b \pm \sqrt{b^2 - 4ac}}{2a}\]

\subsection{Matrix Example}
A sample matrix:
\[
\begin{pmatrix}
a_{11} & a_{12} & a_{13} \\
a_{21} & a_{22} & a_{23} \\
a_{31} & a_{32} & a_{33}
\end{pmatrix}
\]

\section{Tables}
Below is a sample table showing data:

\begin{table}[h]
\centering
\begin{tabular}{@{}lrrr@{}}
\toprule
\textbf{Category} & \textbf{Value 1} & \textbf{Value 2} & \textbf{Total} \\
\midrule
Item A & 10.5 & 20.3 & 30.8 \\
Item B & 15.2 & 25.1 & 40.3 \\
Item C & 22.1 & 12.7 & 34.8 \\
\bottomrule
\end{tabular}
\caption{Sample Data Table}
\label{tab:sample}
\end{table}

\section{Lists}
\subsection{Bullet Points}
\begin{itemize}
  \item First item
  \item Second item
  \item Third item with sub-items:
    \begin{itemize}
      \item Sub-item 1
      \item Sub-item 2
    \end{itemize}
\end{itemize}

\subsection{Numbered List}
\begin{enumerate}
  \item First step
  \item Second step
  \item Third step
\end{enumerate}

\section{Citations and References}
When you use the citation search feature, references will be automatically managed. For example, you might cite a paper like this: \cite{example2024}. All citations will be automatically added to your \texttt{cite.bib} file.

\section{Conclusion}
This template demonstrates basic \LaTeX{} features including mathematical formulas, tables, and lists. The bibliography system is now set up to work seamlessly with Bibby AI's citation management. Feel free to modify it for your needs!

\bibliographystyle{plain}
\bibliography{cite}

\end{document}\documentclass{article}
\usepackage[utf8]{inputenc}
\usepackage{booktabs}
\usepackage{amsmath}
\usepackage{natbib}

\title{My Document}
\author{Your Name}
\date{\today}

\begin{document}

\maketitle

\section{Introduction}
This is a sample document demonstrating various \LaTeX{} features including mathematical equations, tables, and formatting. You can cite papers using the citation search feature, and they will be automatically added to your \texttt{cite.bib} file.

\section{Mathematical Expressions}
Here are some example equations:

\subsection{Quadratic Formula}
The quadratic formula is given by:
\[x = \frac{-b \pm \sqrt{b^2 - 4ac}}{2a}\]

\subsection{Matrix Example}
A sample matrix:
\[
\begin{pmatrix}
a_{11} & a_{12} & a_{13} \\
a_{21} & a_{22} & a_{23} \\
a_{31} & a_{32} & a_{33}
\end{pmatrix}
\]

\section{Tables}
Below is a sample table showing data:

\begin{table}[h]
\centering
\begin{tabular}{@{}lrrr@{}}
\toprule
\textbf{Category} & \textbf{Value 1} & \textbf{Value 2} & \textbf{Total} \\
\midrule
Item A & 10.5 & 20.3 & 30.8 \\
Item B & 15.2 & 25.1 & 40.3 \\
Item C & 22.1 & 12.7 & 34.8 \\
\bottomrule
\end{tabular}
\caption{Sample Data Table}
\label{tab:sample}
\end{table}

\section{Lists}
\subsection{Bullet Points}
\begin{itemize}
  \item First item
  \item Second item
  \item Third item with sub-items:
    \begin{itemize}
      \item Sub-item 1
      \item Sub-item 2
    \end{itemize}
\end{itemize}

\subsection{Numbered List}
\begin{enumerate}
  \item First step
  \item Second step
  \item Third step
\end{enumerate}

\section{Citations and References}
When you use the citation search feature, references will be automatically managed. For example, you might cite a paper like this: \cite{example2024}. All citations will be automatically added to your \texttt{cite.bib} file.

\section{Conclusion}
This template demonstrates basic \LaTeX{} features including mathematical formulas, tables, and lists. The bibliography system is now set up to work seamlessly with Bibby AI's citation management. Feel free to modify it for your needs!

\bibliographystyle{plain}
\bibliography{cite}

\end{document}\documentclass{article}
\usepackage[utf8]{inputenc}
\usepackage{booktabs}
\usepackage{amsmath}
\usepackage{natbib}

\title{My Document}
\author{Your Name}
\date{\today}

\begin{document}

\maketitle

\section{Introduction}
This is a sample document demonstrating various \LaTeX{} features including mathematical equations, tables, and formatting. You can cite papers using the citation search feature, and they will be automatically added to your \texttt{cite.bib} file.

\section{Mathematical Expressions}
Here are some example equations:

\subsection{Quadratic Formula}
The quadratic formula is given by:
\[x = \frac{-b \pm \sqrt{b^2 - 4ac}}{2a}\]

\subsection{Matrix Example}
A sample matrix:
\[
\begin{pmatrix}
a_{11} & a_{12} & a_{13} \\
a_{21} & a_{22} & a_{23} \\
a_{31} & a_{32} & a_{33}
\end{pmatrix}
\]

\section{Tables}
Below is a sample table showing data:

\begin{table}[h]
\centering
\begin{tabular}{@{}lrrr@{}}
\toprule
\textbf{Category} & \textbf{Value 1} & \textbf{Value 2} & \textbf{Total} \\
\midrule
Item A & 10.5 & 20.3 & 30.8 \\
Item B & 15.2 & 25.1 & 40.3 \\
Item C & 22.1 & 12.7 & 34.8 \\
\bottomrule
\end{tabular}
\caption{Sample Data Table}
\label{tab:sample}
\end{table}

\section{Lists}
\subsection{Bullet Points}
\begin{itemize}
  \item First item
  \item Second item
  \item Third item with sub-items:
    \begin{itemize}
      \item Sub-item 1
      \item Sub-item 2
    \end{itemize}
\end{itemize}

\subsection{Numbered List}
\begin{enumerate}
  \item First step
  \item Second step
  \item Third step
\end{enumerate}

\section{Citations and References}
When you use the citation search feature, references will be automatically managed. For example, you might cite a paper like this: \cite{example2024}. All citations will be automatically added to your \texttt{cite.bib} file.

\section{Conclusion}
This template demonstrates basic \LaTeX{} features including mathematical formulas, tables, and lists. The bibliography system is now set up to work seamlessly with Bibby AI's citation management. Feel free to modify it for your needs!

\bibliographystyle{plain}
\bibliography{cite}

\end{document}\documentclass{article}
\usepackage[utf8]{inputenc}
\usepackage{booktabs}
\usepackage{amsmath}
\usepackage{natbib}

\title{My Document}
\author{Your Name}
\date{\today}

\begin{document}

\maketitle

\section{Introduction}
This is a sample document demonstrating various \LaTeX{} features including mathematical equations, tables, and formatting. You can cite papers using the citation search feature, and they will be automatically added to your \texttt{cite.bib} file.

\section{Mathematical Expressions}
Here are some example equations:

\subsection{Quadratic Formula}
The quadratic formula is given by:
\[x = \frac{-b \pm \sqrt{b^2 - 4ac}}{2a}\]

\subsection{Matrix Example}
A sample matrix:
\[
\begin{pmatrix}
a_{11} & a_{12} & a_{13} \\
a_{21} & a_{22} & a_{23} \\
a_{31} & a_{32} & a_{33}
\end{pmatrix}
\]

\section{Tables}
Below is a sample table showing data:

\begin{table}[h]
\centering
\begin{tabular}{@{}lrrr@{}}
\toprule
\textbf{Category} & \textbf{Value 1} & \textbf{Value 2} & \textbf{Total} \\
\midrule
Item A & 10.5 & 20.3 & 30.8 \\
Item B & 15.2 & 25.1 & 40.3 \\
Item C & 22.1 & 12.7 & 34.8 \\
\bottomrule
\end{tabular}
\caption{Sample Data Table}
\label{tab:sample}
\end{table}

\section{Lists}
\subsection{Bullet Points}
\begin{itemize}
  \item First item
  \item Second item
  \item Third item with sub-items:
    \begin{itemize}
      \item Sub-item 1
      \item Sub-item 2
    \end{itemize}
\end{itemize}

\subsection{Numbered List}
\begin{enumerate}
  \item First step
  \item Second step
  \item Third step
\end{enumerate}

\section{Citations and References}
When you use the citation search feature, references will be automatically managed. For example, you might cite a paper like this: \cite{example2024}. All citations will be automatically added to your \texttt{cite.bib} file.

\section{Conclusion}
This template demonstrates basic \LaTeX{} features including mathematical formulas, tables, and lists. The bibliography system is now set up to work seamlessly with Bibby AI's citation management. Feel free to modify it for your needs!

\bibliographystyle{plain}
\bibliography{cite}

\end{document}\documentclass{article}
\usepackage[utf8]{inputenc}
\usepackage{booktabs}
\usepackage{amsmath}
\usepackage{natbib}

\title{My Document}
\author{Your Name}
\date{\today}

\begin{document}

\maketitle

\section{Introduction}
This is a sample document demonstrating various \LaTeX{} features including mathematical equations, tables, and formatting. You can cite papers using the citation search feature, and they will be automatically added to your \texttt{cite.bib} file.

\section{Mathematical Expressions}
Here are some example equations:

\subsection{Quadratic Formula}
The quadratic formula is given by:
\[x = \frac{-b \pm \sqrt{b^2 - 4ac}}{2a}\]

\subsection{Matrix Example}
A sample matrix:
\[
\begin{pmatrix}
a_{11} & a_{12} & a_{13} \\
a_{21} & a_{22} & a_{23} \\
a_{31} & a_{32} & a_{33}
\end{pmatrix}
\]

\section{Tables}
Below is a sample table showing data:

\begin{table}[h]
\centering
\begin{tabular}{@{}lrrr@{}}
\toprule
\textbf{Category} & \textbf{Value 1} & \textbf{Value 2} & \textbf{Total} \\
\midrule
Item A & 10.5 & 20.3 & 30.8 \\
Item B & 15.2 & 25.1 & 40.3 \\
Item C & 22.1 & 12.7 & 34.8 \\
\bottomrule
\end{tabular}
\caption{Sample Data Table}
\label{tab:sample}
\end{table}

\section{Lists}
\subsection{Bullet Points}
\begin{itemize}
  \item First item
  \item Second item
  \item Third item with sub-items:
    \begin{itemize}
      \item Sub-item 1
      \item Sub-item 2
    \end{itemize}
\end{itemize}

\subsection{Numbered List}
\begin{enumerate}
  \item First step
  \item Second step
  \item Third step
\end{enumerate}

\section{Citations and References}
When you use the citation search feature, references will be automatically managed. For example, you might cite a paper like this: \cite{example2024}. All citations will be automatically added to your \texttt{cite.bib} file.

\section{Conclusion}
This template demonstrates basic \LaTeX{} features including mathematical formulas, tables, and lists. The bibliography system is now set up to work seamlessly with Bibby AI's citation management. Feel free to modify it for your needs!

\bibliographystyle{plain}
\bibliography{cite}

\end{document}\documentclass{article}
\usepackage[utf8]{inputenc}
\usepackage{booktabs}
\usepackage{amsmath}
\usepackage{natbib}

\title{My Document}
\author{Your Name}
\date{\today}

\begin{document}

\maketitle

\section{Introduction}
This is a sample document demonstrating various \LaTeX{} features including mathematical equations, tables, and formatting. You can cite papers using the citation search feature, and they will be automatically added to your \texttt{cite.bib} file.

\section{Mathematical Expressions}
Here are some example equations:

\subsection{Quadratic Formula}
The quadratic formula is given by:
\[x = \frac{-b \pm \sqrt{b^2 - 4ac}}{2a}\]

\subsection{Matrix Example}
A sample matrix:
\[
\begin{pmatrix}
a_{11} & a_{12} & a_{13} \\
a_{21} & a_{22} & a_{23} \\
a_{31} & a_{32} & a_{33}
\end{pmatrix}
\]

\section{Tables}
Below is a sample table showing data:

\begin{table}[h]
\centering
\begin{tabular}{@{}lrrr@{}}
\toprule
\textbf{Category} & \textbf{Value 1} & \textbf{Value 2} & \textbf{Total} \\
\midrule
Item A & 10.5 & 20.3 & 30.8 \\
Item B & 15.2 & 25.1 & 40.3 \\
Item C & 22.1 & 12.7 & 34.8 \\
\bottomrule
\end{tabular}
\caption{Sample Data Table}
\label{tab:sample}
\end{table}

\section{Lists}
\subsection{Bullet Points}
\begin{itemize}
  \item First item
  \item Second item
  \item Third item with sub-items:
    \begin{itemize}
      \item Sub-item 1
      \item Sub-item 2
    \end{itemize}
\end{itemize}

\subsection{Numbered List}
\begin{enumerate}
  \item First step
  \item Second step
  \item Third step
\end{enumerate}

\section{Citations and References}
When you use the citation search feature, references will be automatically managed. For example, you might cite a paper like this: \cite{example2024}. All citations will be automatically added to your \texttt{cite.bib} file.

\section{Conclusion}
This template demonstrates basic \LaTeX{} features including mathematical formulas, tables, and lists. The bibliography system is now set up to work seamlessly with Bibby AI's citation management. Feel free to modify it for your needs!

\bibliographystyle{plain}
\bibliography{cite}

\end{document}\documentclass{article}
\usepackage[utf8]{inputenc}
\usepackage{booktabs}
\usepackage{amsmath}
\usepackage{natbib}

\title{My Document}
\author{Your Name}
\date{\today}

\begin{document}

\maketitle

\section{Introduction}
This is a sample document demonstrating various \LaTeX{} features including mathematical equations, tables, and formatting. You can cite papers using the citation search feature, and they will be automatically added to your \texttt{cite.bib} file.

\section{Mathematical Expressions}
Here are some example equations:

\subsection{Quadratic Formula}
The quadratic formula is given by:
\[x = \frac{-b \pm \sqrt{b^2 - 4ac}}{2a}\]

\subsection{Matrix Example}
A sample matrix:
\[
\begin{pmatrix}
a_{11} & a_{12} & a_{13} \\
a_{21} & a_{22} & a_{23} \\
a_{31} & a_{32} & a_{33}
\end{pmatrix}
\]

\section{Tables}
Below is a sample table showing data:

\begin{table}[h]
\centering
\begin{tabular}{@{}lrrr@{}}
\toprule
\textbf{Category} & \textbf{Value 1} & \textbf{Value 2} & \textbf{Total} \\
\midrule
Item A & 10.5 & 20.3 & 30.8 \\
Item B & 15.2 & 25.1 & 40.3 \\
Item C & 22.1 & 12.7 & 34.8 \\
\bottomrule
\end{tabular}
\caption{Sample Data Table}
\label{tab:sample}
\end{table}

\section{Lists}
\subsection{Bullet Points}
\begin{itemize}
  \item First item
  \item Second item
  \item Third item with sub-items:
    \begin{itemize}
      \item Sub-item 1
      \item Sub-item 2
    \end{itemize}
\end{itemize}

\subsection{Numbered List}
\begin{enumerate}
  \item First step
  \item Second step
  \item Third step
\end{enumerate}

\section{Citations and References}
When you use the citation search feature, references will be automatically managed. For example, you might cite a paper like this: \cite{example2024}. All citations will be automatically added to your \texttt{cite.bib} file.

\section{Conclusion}
This template demonstrates basic \LaTeX{} features including mathematical formulas, tables, and lists. The bibliography system is now set up to work seamlessly with Bibby AI's citation management. Feel free to modify it for your needs!

\bibliographystyle{plain}
\bibliography{cite}

\end{document}\documentclass{article}
\usepackage[utf8]{inputenc}
\usepackage{booktabs}
\usepackage{amsmath}
\usepackage{natbib}

\title{My Document}
\author{Your Name}
\date{\today}

\begin{document}

\maketitle

\section{Introduction}
This is a sample document demonstrating various \LaTeX{} features including mathematical equations, tables, and formatting. You can cite papers using the citation search feature, and they will be automatically added to your \texttt{cite.bib} file.

\section{Mathematical Expressions}
Here are some example equations:

\subsection{Quadratic Formula}
The quadratic formula is given by:
\[x = \frac{-b \pm \sqrt{b^2 - 4ac}}{2a}\]

\subsection{Matrix Example}
A sample matrix:
\[
\begin{pmatrix}
a_{11} & a_{12} & a_{13} \\
a_{21} & a_{22} & a_{23} \\
a_{31} & a_{32} & a_{33}
\end{pmatrix}
\]

\section{Tables}
Below is a sample table showing data:

\begin{table}[h]
\centering
\begin{tabular}{@{}lrrr@{}}
\toprule
\textbf{Category} & \textbf{Value 1} & \textbf{Value 2} & \textbf{Total} \\
\midrule
Item A & 10.5 & 20.3 & 30.8 \\
Item B & 15.2 & 25.1 & 40.3 \\
Item C & 22.1 & 12.7 & 34.8 \\
\bottomrule
\end{tabular}
\caption{Sample Data Table}
\label{tab:sample}
\end{table}

\section{Lists}
\subsection{Bullet Points}
\begin{itemize}
  \item First item
  \item Second item
  \item Third item with sub-items:
    \begin{itemize}
      \item Sub-item 1
      \item Sub-item 2
    \end{itemize}
\end{itemize}

\subsection{Numbered List}
\begin{enumerate}
  \item First step
  \item Second step
  \item Third step
\end{enumerate}

\section{Citations and References}
When you use the citation search feature, references will be automatically managed. For example, you might cite a paper like this: \cite{example2024}. All citations will be automatically added to your \texttt{cite.bib} file.

\section{Conclusion}
This template demonstrates basic \LaTeX{} features including mathematical formulas, tables, and lists. The bibliography system is now set up to work seamlessly with Bibby AI's citation management. Feel free to modify it for your needs!

\bibliographystyle{plain}
\bibliography{cite}

\end{document}\documentclass{article}
\usepackage[utf8]{inputenc}
\usepackage{booktabs}
\usepackage{amsmath}
\usepackage{natbib}

\title{My Document}
\author{Your Name}
\date{\today}

\begin{document}

\maketitle

\section{Introduction}
This is a sample document demonstrating various \LaTeX{} features including mathematical equations, tables, and formatting. You can cite papers using the citation search feature, and they will be automatically added to your \texttt{cite.bib} file.

\section{Mathematical Expressions}
Here are some example equations:

\subsection{Quadratic Formula}
The quadratic formula is given by:
\[x = \frac{-b \pm \sqrt{b^2 - 4ac}}{2a}\]

\subsection{Matrix Example}
A sample matrix:
\[
\begin{pmatrix}
a_{11} & a_{12} & a_{13} \\
a_{21} & a_{22} & a_{23} \\
a_{31} & a_{32} & a_{33}
\end{pmatrix}
\]

\section{Tables}
Below is a sample table showing data:

\begin{table}[h]
\centering
\begin{tabular}{@{}lrrr@{}}
\toprule
\textbf{Category} & \textbf{Value 1} & \textbf{Value 2} & \textbf{Total} \\
\midrule
Item A & 10.5 & 20.3 & 30.8 \\
Item B & 15.2 & 25.1 & 40.3 \\
Item C & 22.1 & 12.7 & 34.8 \\
\bottomrule
\end{tabular}
\caption{Sample Data Table}
\label{tab:sample}
\end{table}

\section{Lists}
\subsection{Bullet Points}
\begin{itemize}
  \item First item
  \item Second item
  \item Third item with sub-items:
    \begin{itemize}
      \item Sub-item 1
      \item Sub-item 2
    \end{itemize}
\end{itemize}

\subsection{Numbered List}
\begin{enumerate}
  \item First step
  \item Second step
  \item Third step
\end{enumerate}

\section{Citations and References}
When you use the citation search feature, references will be automatically managed. For example, you might cite a paper like this: \cite{example2024}. All citations will be automatically added to your \texttt{cite.bib} file.

\section{Conclusion}
This template demonstrates basic \LaTeX{} features including mathematical formulas, tables, and lists. The bibliography system is now set up to work seamlessly with Bibby AI's citation management. Feel free to modify it for your needs!

\bibliographystyle{plain}
\bibliography{cite}

\end{document}\documentclass{article}
\usepackage[utf8]{inputenc}
\usepackage{booktabs}
\usepackage{amsmath}
\usepackage{natbib}

\title{My Document}
\author{Your Name}
\date{\today}

\begin{document}

\maketitle

\section{Introduction}
This is a sample document demonstrating various \LaTeX{} features including mathematical equations, tables, and formatting. You can cite papers using the citation search feature, and they will be automatically added to your \texttt{cite.bib} file.

\section{Mathematical Expressions}
Here are some example equations:

\subsection{Quadratic Formula}
The quadratic formula is given by:
\[x = \frac{-b \pm \sqrt{b^2 - 4ac}}{2a}\]

\subsection{Matrix Example}
A sample matrix:
\[
\begin{pmatrix}
a_{11} & a_{12} & a_{13} \\
a_{21} & a_{22} & a_{23} \\
a_{31} & a_{32} & a_{33}
\end{pmatrix}
\]

\section{Tables}
Below is a sample table showing data:

\begin{table}[h]
\centering
\begin{tabular}{@{}lrrr@{}}
\toprule
\textbf{Category} & \textbf{Value 1} & \textbf{Value 2} & \textbf{Total} \\
\midrule
Item A & 10.5 & 20.3 & 30.8 \\
Item B & 15.2 & 25.1 & 40.3 \\
Item C & 22.1 & 12.7 & 34.8 \\
\bottomrule
\end{tabular}
\caption{Sample Data Table}
\label{tab:sample}
\end{table}

\section{Lists}
\subsection{Bullet Points}
\begin{itemize}
  \item First item
  \item Second item
  \item Third item with sub-items:
    \begin{itemize}
      \item Sub-item 1
      \item Sub-item 2
    \end{itemize}
\end{itemize}

\subsection{Numbered List}
\begin{enumerate}
  \item First step
  \item Second step
  \item Third step
\end{enumerate}

\section{Citations and References}
When you use the citation search feature, references will be automatically managed. For example, you might cite a paper like this: \cite{example2024}. All citations will be automatically added to your \texttt{cite.bib} file.

\section{Conclusion}
This template demonstrates basic \LaTeX{} features including mathematical formulas, tables, and lists. The bibliography system is now set up to work seamlessly with Bibby AI's citation management. Feel free to modify it for your needs!

\bibliographystyle{plain}
\bibliography{cite}

\end{document}\documentclass{article}
\usepackage[utf8]{inputenc}
\usepackage{booktabs}
\usepackage{amsmath}
\usepackage{natbib}

\title{My Document}
\author{Your Name}
\date{\today}

\begin{document}

\maketitle

\section{Introduction}
This is a sample document demonstrating various \LaTeX{} features including mathematical equations, tables, and formatting. You can cite papers using the citation search feature, and they will be automatically added to your \texttt{cite.bib} file.

\section{Mathematical Expressions}
Here are some example equations:

\subsection{Quadratic Formula}
The quadratic formula is given by:
\[x = \frac{-b \pm \sqrt{b^2 - 4ac}}{2a}\]

\subsection{Matrix Example}
A sample matrix:
\[
\begin{pmatrix}
a_{11} & a_{12} & a_{13} \\
a_{21} & a_{22} & a_{23} \\
a_{31} & a_{32} & a_{33}
\end{pmatrix}
\]

\section{Tables}
Below is a sample table showing data:

\begin{table}[h]
\centering
\begin{tabular}{@{}lrrr@{}}
\toprule
\textbf{Category} & \textbf{Value 1} & \textbf{Value 2} & \textbf{Total} \\
\midrule
Item A & 10.5 & 20.3 & 30.8 \\
Item B & 15.2 & 25.1 & 40.3 \\
Item C & 22.1 & 12.7 & 34.8 \\
\bottomrule
\end{tabular}
\caption{Sample Data Table}
\label{tab:sample}
\end{table}

\section{Lists}
\subsection{Bullet Points}
\begin{itemize}
  \item First item
  \item Second item
  \item Third item with sub-items:
    \begin{itemize}
      \item Sub-item 1
      \item Sub-item 2
    \end{itemize}
\end{itemize}

\subsection{Numbered List}
\begin{enumerate}
  \item First step
  \item Second step
  \item Third step
\end{enumerate}

\section{Citations and References}
When you use the citation search feature, references will be automatically managed. For example, you might cite a paper like this: \cite{example2024}. All citations will be automatically added to your \texttt{cite.bib} file.

\section{Conclusion}
This template demonstrates basic \LaTeX{} features including mathematical formulas, tables, and lists. The bibliography system is now set up to work seamlessly with Bibby AI's citation management. Feel free to modify it for your needs!

\bibliographystyle{plain}
\bibliography{cite}

\end{document}\documentclass{article}
\usepackage[utf8]{inputenc}
\usepackage{booktabs}
\usepackage{amsmath}
\usepackage{natbib}

\title{My Document}
\author{Your Name}
\date{\today}

\begin{document}

\maketitle

\section{Introduction}
This is a sample document demonstrating various \LaTeX{} features including mathematical equations, tables, and formatting. You can cite papers using the citation search feature, and they will be automatically added to your \texttt{cite.bib} file.

\section{Mathematical Expressions}
Here are some example equations:

\subsection{Quadratic Formula}
The quadratic formula is given by:
\[x = \frac{-b \pm \sqrt{b^2 - 4ac}}{2a}\]

\subsection{Matrix Example}
A sample matrix:
\[
\begin{pmatrix}
a_{11} & a_{12} & a_{13} \\
a_{21} & a_{22} & a_{23} \\
a_{31} & a_{32} & a_{33}
\end{pmatrix}
\]

\section{Tables}
Below is a sample table showing data:

\begin{table}[h]
\centering
\begin{tabular}{@{}lrrr@{}}
\toprule
\textbf{Category} & \textbf{Value 1} & \textbf{Value 2} & \textbf{Total} \\
\midrule
Item A & 10.5 & 20.3 & 30.8 \\
Item B & 15.2 & 25.1 & 40.3 \\
Item C & 22.1 & 12.7 & 34.8 \\
\bottomrule
\end{tabular}
\caption{Sample Data Table}
\label{tab:sample}
\end{table}

\section{Lists}
\subsection{Bullet Points}
\begin{itemize}
  \item First item
  \item Second item
  \item Third item with sub-items:
    \begin{itemize}
      \item Sub-item 1
      \item Sub-item 2
    \end{itemize}
\end{itemize}

\subsection{Numbered List}
\begin{enumerate}
  \item First step
  \item Second step
  \item Third step
\end{enumerate}

\section{Citations and References}
When you use the citation search feature, references will be automatically managed. For example, you might cite a paper like this: \cite{example2024}. All citations will be automatically added to your \texttt{cite.bib} file.

\section{Conclusion}
This template demonstrates basic \LaTeX{} features including mathematical formulas, tables, and lists. The bibliography system is now set up to work seamlessly with Bibby AI's citation management. Feel free to modify it for your needs!

\bibliographystyle{plain}
\bibliography{cite}

\end{document}\documentclass{article}
\usepackage[utf8]{inputenc}
\usepackage{booktabs}
\usepackage{amsmath}
\usepackage{natbib}

\title{My Document}
\author{Your Name}
\date{\today}

\begin{document}

\maketitle

\section{Introduction}
This is a sample document demonstrating various \LaTeX{} features including mathematical equations, tables, and formatting. You can cite papers using the citation search feature, and they will be automatically added to your \texttt{cite.bib} file.

\section{Mathematical Expressions}
Here are some example equations:

\subsection{Quadratic Formula}
The quadratic formula is given by:
\[x = \frac{-b \pm \sqrt{b^2 - 4ac}}{2a}\]

\subsection{Matrix Example}
A sample matrix:
\[
\begin{pmatrix}
a_{11} & a_{12} & a_{13} \\
a_{21} & a_{22} & a_{23} \\
a_{31} & a_{32} & a_{33}
\end{pmatrix}
\]

\section{Tables}
Below is a sample table showing data:

\begin{table}[h]
\centering
\begin{tabular}{@{}lrrr@{}}
\toprule
\textbf{Category} & \textbf{Value 1} & \textbf{Value 2} & \textbf{Total} \\
\midrule
Item A & 10.5 & 20.3 & 30.8 \\
Item B & 15.2 & 25.1 & 40.3 \\
Item C & 22.1 & 12.7 & 34.8 \\
\bottomrule
\end{tabular}
\caption{Sample Data Table}
\label{tab:sample}
\end{table}

\section{Lists}
\subsection{Bullet Points}
\begin{itemize}
  \item First item
  \item Second item
  \item Third item with sub-items:
    \begin{itemize}
      \item Sub-item 1
      \item Sub-item 2
    \end{itemize}
\end{itemize}

\subsection{Numbered List}
\begin{enumerate}
  \item First step
  \item Second step
  \item Third step
\end{enumerate}

\section{Citations and References}
When you use the citation search feature, references will be automatically managed. For example, you might cite a paper like this: \cite{example2024}. All citations will be automatically added to your \texttt{cite.bib} file.

\section{Conclusion}
This template demonstrates basic \LaTeX{} features including mathematical formulas, tables, and lists. The bibliography system is now set up to work seamlessly with Bibby AI's citation management. Feel free to modify it for your needs!

\bibliographystyle{plain}
\bibliography{cite}

\end{document}\documentclass{article}
\usepackage[utf8]{inputenc}
\usepackage{booktabs}
\usepackage{amsmath}
\usepackage{natbib}

\title{My Document}
\author{Your Name}
\date{\today}

\begin{document}

\maketitle

\section{Introduction}
This is a sample document demonstrating various \LaTeX{} features including mathematical equations, tables, and formatting. You can cite papers using the citation search feature, and they will be automatically added to your \texttt{cite.bib} file.

\section{Mathematical Expressions}
Here are some example equations:

\subsection{Quadratic Formula}
The quadratic formula is given by:
\[x = \frac{-b \pm \sqrt{b^2 - 4ac}}{2a}\]

\subsection{Matrix Example}
A sample matrix:
\[
\begin{pmatrix}
a_{11} & a_{12} & a_{13} \\
a_{21} & a_{22} & a_{23} \\
a_{31} & a_{32} & a_{33}
\end{pmatrix}
\]

\section{Tables}
Below is a sample table showing data:

\begin{table}[h]
\centering
\begin{tabular}{@{}lrrr@{}}
\toprule
\textbf{Category} & \textbf{Value 1} & \textbf{Value 2} & \textbf{Total} \\
\midrule
Item A & 10.5 & 20.3 & 30.8 \\
Item B & 15.2 & 25.1 & 40.3 \\
Item C & 22.1 & 12.7 & 34.8 \\
\bottomrule
\end{tabular}
\caption{Sample Data Table}
\label{tab:sample}
\end{table}

\section{Lists}
\subsection{Bullet Points}
\begin{itemize}
  \item First item
  \item Second item
  \item Third item with sub-items:
    \begin{itemize}
      \item Sub-item 1
      \item Sub-item 2
    \end{itemize}
\end{itemize}

\subsection{Numbered List}
\begin{enumerate}
  \item First step
  \item Second step
  \item Third step
\end{enumerate}

\section{Citations and References}
When you use the citation search feature, references will be automatically managed. For example, you might cite a paper like this: \cite{example2024}. All citations will be automatically added to your \texttt{cite.bib} file.

\section{Conclusion}
This template demonstrates basic \LaTeX{} features including mathematical formulas, tables, and lists. The bibliography system is now set up to work seamlessly with Bibby AI's citation management. Feel free to modify it for your needs!

\bibliographystyle{plain}
\bibliography{cite}

\end{document}\documentclass{article}
\usepackage[utf8]{inputenc}
\usepackage{booktabs}
\usepackage{amsmath}
\usepackage{natbib}

\title{My Document}
\author{Your Name}
\date{\today}

\begin{document}

\maketitle

\section{Introduction}
This is a sample document demonstrating various \LaTeX{} features including mathematical equations, tables, and formatting. You can cite papers using the citation search feature, and they will be automatically added to your \texttt{cite.bib} file.

\section{Mathematical Expressions}
Here are some example equations:

\subsection{Quadratic Formula}
The quadratic formula is given by:
\[x = \frac{-b \pm \sqrt{b^2 - 4ac}}{2a}\]

\subsection{Matrix Example}
A sample matrix:
\[
\begin{pmatrix}
a_{11} & a_{12} & a_{13} \\
a_{21} & a_{22} & a_{23} \\
a_{31} & a_{32} & a_{33}
\end{pmatrix}
\]

\section{Tables}
Below is a sample table showing data:

\begin{table}[h]
\centering
\begin{tabular}{@{}lrrr@{}}
\toprule
\textbf{Category} & \textbf{Value 1} & \textbf{Value 2} & \textbf{Total} \\
\midrule
Item A & 10.5 & 20.3 & 30.8 \\
Item B & 15.2 & 25.1 & 40.3 \\
Item C & 22.1 & 12.7 & 34.8 \\
\bottomrule
\end{tabular}
\caption{Sample Data Table}
\label{tab:sample}
\end{table}

\section{Lists}
\subsection{Bullet Points}
\begin{itemize}
  \item First item
  \item Second item
  \item Third item with sub-items:
    \begin{itemize}
      \item Sub-item 1
      \item Sub-item 2
    \end{itemize}
\end{itemize}

\subsection{Numbered List}
\begin{enumerate}
  \item First step
  \item Second step
  \item Third step
\end{enumerate}

\section{Citations and References}
When you use the citation search feature, references will be automatically managed. For example, you might cite a paper like this: \cite{example2024}. All citations will be automatically added to your \texttt{cite.bib} file.

\section{Conclusion}
This template demonstrates basic \LaTeX{} features including mathematical formulas, tables, and lists. The bibliography system is now set up to work seamlessly with Bibby AI's citation management. Feel free to modify it for your needs!

\bibliographystyle{plain}
\bibliography{cite}

\end{document}\documentclass{article}
\usepackage[utf8]{inputenc}
\usepackage{booktabs}
\usepackage{amsmath}
\usepackage{natbib}

\title{My Document}
\author{Your Name}
\date{\today}

\begin{document}

\maketitle

\section{Introduction}
This is a sample document demonstrating various \LaTeX{} features including mathematical equations, tables, and formatting. You can cite papers using the citation search feature, and they will be automatically added to your \texttt{cite.bib} file.

\section{Mathematical Expressions}
Here are some example equations:

\subsection{Quadratic Formula}
The quadratic formula is given by:
\[x = \frac{-b \pm \sqrt{b^2 - 4ac}}{2a}\]

\subsection{Matrix Example}
A sample matrix:
\[
\begin{pmatrix}
a_{11} & a_{12} & a_{13} \\
a_{21} & a_{22} & a_{23} \\
a_{31} & a_{32} & a_{33}
\end{pmatrix}
\]

\section{Tables}
Below is a sample table showing data:

\begin{table}[h]
\centering
\begin{tabular}{@{}lrrr@{}}
\toprule
\textbf{Category} & \textbf{Value 1} & \textbf{Value 2} & \textbf{Total} \\
\midrule
Item A & 10.5 & 20.3 & 30.8 \\
Item B & 15.2 & 25.1 & 40.3 \\
Item C & 22.1 & 12.7 & 34.8 \\
\bottomrule
\end{tabular}
\caption{Sample Data Table}
\label{tab:sample}
\end{table}

\section{Lists}
\subsection{Bullet Points}
\begin{itemize}
  \item First item
  \item Second item
  \item Third item with sub-items:
    \begin{itemize}
      \item Sub-item 1
      \item Sub-item 2
    \end{itemize}
\end{itemize}

\subsection{Numbered List}
\begin{enumerate}
  \item First step
  \item Second step
  \item Third step
\end{enumerate}

\section{Citations and References}
When you use the citation search feature, references will be automatically managed. For example, you might cite a paper like this: \cite{example2024}. All citations will be automatically added to your \texttt{cite.bib} file.

\section{Conclusion}
This template demonstrates basic \LaTeX{} features including mathematical formulas, tables, and lists. The bibliography system is now set up to work seamlessly with Bibby AI's citation management. Feel free to modify it for your needs!

\bibliographystyle{plain}
\bibliography{cite}

\end{document}\documentclass{article}
\usepackage[utf8]{inputenc}
\usepackage{booktabs}
\usepackage{amsmath}
\usepackage{natbib}

\title{My Document}
\author{Your Name}
\date{\today}

\begin{document}

\maketitle

\section{Introduction}
This is a sample document demonstrating various \LaTeX{} features including mathematical equations, tables, and formatting. You can cite papers using the citation search feature, and they will be automatically added to your \texttt{cite.bib} file.

\section{Mathematical Expressions}
Here are some example equations:

\subsection{Quadratic Formula}
The quadratic formula is given by:
\[x = \frac{-b \pm \sqrt{b^2 - 4ac}}{2a}\]

\subsection{Matrix Example}
A sample matrix:
\[
\begin{pmatrix}
a_{11} & a_{12} & a_{13} \\
a_{21} & a_{22} & a_{23} \\
a_{31} & a_{32} & a_{33}
\end{pmatrix}
\]

\section{Tables}
Below is a sample table showing data:

\begin{table}[h]
\centering
\begin{tabular}{@{}lrrr@{}}
\toprule
\textbf{Category} & \textbf{Value 1} & \textbf{Value 2} & \textbf{Total} \\
\midrule
Item A & 10.5 & 20.3 & 30.8 \\
Item B & 15.2 & 25.1 & 40.3 \\
Item C & 22.1 & 12.7 & 34.8 \\
\bottomrule
\end{tabular}
\caption{Sample Data Table}
\label{tab:sample}
\end{table}

\section{Lists}
\subsection{Bullet Points}
\begin{itemize}
  \item First item
  \item Second item
  \item Third item with sub-items:
    \begin{itemize}
      \item Sub-item 1
      \item Sub-item 2
    \end{itemize}
\end{itemize}

\subsection{Numbered List}
\begin{enumerate}
  \item First step
  \item Second step
  \item Third step
\end{enumerate}

\section{Citations and References}
When you use the citation search feature, references will be automatically managed. For example, you might cite a paper like this: \cite{example2024}. All citations will be automatically added to your \texttt{cite.bib} file.

\section{Conclusion}
This template demonstrates basic \LaTeX{} features including mathematical formulas, tables, and lists. The bibliography system is now set up to work seamlessly with Bibby AI's citation management. Feel free to modify it for your needs!

\bibliographystyle{plain}
\bibliography{cite}

\end{document}\documentclass{article}
\usepackage[utf8]{inputenc}
\usepackage{booktabs}
\usepackage{amsmath}
\usepackage{natbib}

\title{My Document}
\author{Your Name}
\date{\today}

\begin{document}

\maketitle

\section{Introduction}
This is a sample document demonstrating various \LaTeX{} features including mathematical equations, tables, and formatting. You can cite papers using the citation search feature, and they will be automatically added to your \texttt{cite.bib} file.

\section{Mathematical Expressions}
Here are some example equations:

\subsection{Quadratic Formula}
The quadratic formula is given by:
\[x = \frac{-b \pm \sqrt{b^2 - 4ac}}{2a}\]

\subsection{Matrix Example}
A sample matrix:
\[
\begin{pmatrix}
a_{11} & a_{12} & a_{13} \\
a_{21} & a_{22} & a_{23} \\
a_{31} & a_{32} & a_{33}
\end{pmatrix}
\]

\section{Tables}
Below is a sample table showing data:

\begin{table}[h]
\centering
\begin{tabular}{@{}lrrr@{}}
\toprule
\textbf{Category} & \textbf{Value 1} & \textbf{Value 2} & \textbf{Total} \\
\midrule
Item A & 10.5 & 20.3 & 30.8 \\
Item B & 15.2 & 25.1 & 40.3 \\
Item C & 22.1 & 12.7 & 34.8 \\
\bottomrule
\end{tabular}
\caption{Sample Data Table}
\label{tab:sample}
\end{table}

\section{Lists}
\subsection{Bullet Points}
\begin{itemize}
  \item First item
  \item Second item
  \item Third item with sub-items:
    \begin{itemize}
      \item Sub-item 1
      \item Sub-item 2
    \end{itemize}
\end{itemize}

\subsection{Numbered List}
\begin{enumerate}
  \item First step
  \item Second step
  \item Third step
\end{enumerate}

\section{Citations and References}
When you use the citation search feature, references will be automatically managed. For example, you might cite a paper like this: \cite{example2024}. All citations will be automatically added to your \texttt{cite.bib} file.

\section{Conclusion}
This template demonstrates basic \LaTeX{} features including mathematical formulas, tables, and lists. The bibliography system is now set up to work seamlessly with Bibby AI's citation management. Feel free to modify it for your needs!

\bibliographystyle{plain}
\bibliography{cite}

\end{document}\documentclass{article}
\usepackage[utf8]{inputenc}
\usepackage{booktabs}
\usepackage{amsmath}
\usepackage{natbib}

\title{My Document}
\author{Your Name}
\date{\today}

\begin{document}

\maketitle

\section{Introduction}
This is a sample document demonstrating various \LaTeX{} features including mathematical equations, tables, and formatting. You can cite papers using the citation search feature, and they will be automatically added to your \texttt{cite.bib} file.

\section{Mathematical Expressions}
Here are some example equations:

\subsection{Quadratic Formula}
The quadratic formula is given by:
\[x = \frac{-b \pm \sqrt{b^2 - 4ac}}{2a}\]

\subsection{Matrix Example}
A sample matrix:
\[
\begin{pmatrix}
a_{11} & a_{12} & a_{13} \\
a_{21} & a_{22} & a_{23} \\
a_{31} & a_{32} & a_{33}
\end{pmatrix}
\]

\section{Tables}
Below is a sample table showing data:

\begin{table}[h]
\centering
\begin{tabular}{@{}lrrr@{}}
\toprule
\textbf{Category} & \textbf{Value 1} & \textbf{Value 2} & \textbf{Total} \\
\midrule
Item A & 10.5 & 20.3 & 30.8 \\
Item B & 15.2 & 25.1 & 40.3 \\
Item C & 22.1 & 12.7 & 34.8 \\
\bottomrule
\end{tabular}
\caption{Sample Data Table}
\label{tab:sample}
\end{table}

\section{Lists}
\subsection{Bullet Points}
\begin{itemize}
  \item First item
  \item Second item
  \item Third item with sub-items:
    \begin{itemize}
      \item Sub-item 1
      \item Sub-item 2
    \end{itemize}
\end{itemize}

\subsection{Numbered List}
\begin{enumerate}
  \item First step
  \item Second step
  \item Third step
\end{enumerate}

\section{Citations and References}
When you use the citation search feature, references will be automatically managed. For example, you might cite a paper like this: \cite{example2024}. All citations will be automatically added to your \texttt{cite.bib} file.

\section{Conclusion}
This template demonstrates basic \LaTeX{} features including mathematical formulas, tables, and lists. The bibliography system is now set up to work seamlessly with Bibby AI's citation management. Feel free to modify it for your needs!

\bibliographystyle{plain}
\bibliography{cite}

\end{document}\documentclass{article}
\usepackage[utf8]{inputenc}
\usepackage{booktabs}
\usepackage{amsmath}
\usepackage{natbib}

\title{My Document}
\author{Your Name}
\date{\today}

\begin{document}

\maketitle

\section{Introduction}
This is a sample document demonstrating various \LaTeX{} features including mathematical equations, tables, and formatting. You can cite papers using the citation search feature, and they will be automatically added to your \texttt{cite.bib} file.

\section{Mathematical Expressions}
Here are some example equations:

\subsection{Quadratic Formula}
The quadratic formula is given by:
\[x = \frac{-b \pm \sqrt{b^2 - 4ac}}{2a}\]

\subsection{Matrix Example}
A sample matrix:
\[
\begin{pmatrix}
a_{11} & a_{12} & a_{13} \\
a_{21} & a_{22} & a_{23} \\
a_{31} & a_{32} & a_{33}
\end{pmatrix}
\]

\section{Tables}
Below is a sample table showing data:

\begin{table}[h]
\centering
\begin{tabular}{@{}lrrr@{}}
\toprule
\textbf{Category} & \textbf{Value 1} & \textbf{Value 2} & \textbf{Total} \\
\midrule
Item A & 10.5 & 20.3 & 30.8 \\
Item B & 15.2 & 25.1 & 40.3 \\
Item C & 22.1 & 12.7 & 34.8 \\
\bottomrule
\end{tabular}
\caption{Sample Data Table}
\label{tab:sample}
\end{table}

\section{Lists}
\subsection{Bullet Points}
\begin{itemize}
  \item First item
  \item Second item
  \item Third item with sub-items:
    \begin{itemize}
      \item Sub-item 1
      \item Sub-item 2
    \end{itemize}
\end{itemize}

\subsection{Numbered List}
\begin{enumerate}
  \item First step
  \item Second step
  \item Third step
\end{enumerate}

\section{Citations and References}
When you use the citation search feature, references will be automatically managed. For example, you might cite a paper like this: \cite{example2024}. All citations will be automatically added to your \texttt{cite.bib} file.

\section{Conclusion}
This template demonstrates basic \LaTeX{} features including mathematical formulas, tables, and lists. The bibliography system is now set up to work seamlessly with Bibby AI's citation management. Feel free to modify it for your needs!

\bibliographystyle{plain}
\bibliography{cite}

\end{document}\documentclass{article}
\usepackage[utf8]{inputenc}
\usepackage{booktabs}
\usepackage{amsmath}
\usepackage{natbib}

\title{My Document}
\author{Your Name}
\date{\today}

\begin{document}

\maketitle

\section{Introduction}
This is a sample document demonstrating various \LaTeX{} features including mathematical equations, tables, and formatting. You can cite papers using the citation search feature, and they will be automatically added to your \texttt{cite.bib} file.

\section{Mathematical Expressions}
Here are some example equations:

\subsection{Quadratic Formula}
The quadratic formula is given by:
\[x = \frac{-b \pm \sqrt{b^2 - 4ac}}{2a}\]

\subsection{Matrix Example}
A sample matrix:
\[
\begin{pmatrix}
a_{11} & a_{12} & a_{13} \\
a_{21} & a_{22} & a_{23} \\
a_{31} & a_{32} & a_{33}
\end{pmatrix}
\]

\section{Tables}
Below is a sample table showing data:

\begin{table}[h]
\centering
\begin{tabular}{@{}lrrr@{}}
\toprule
\textbf{Category} & \textbf{Value 1} & \textbf{Value 2} & \textbf{Total} \\
\midrule
Item A & 10.5 & 20.3 & 30.8 \\
Item B & 15.2 & 25.1 & 40.3 \\
Item C & 22.1 & 12.7 & 34.8 \\
\bottomrule
\end{tabular}
\caption{Sample Data Table}
\label{tab:sample}
\end{table}

\section{Lists}
\subsection{Bullet Points}
\begin{itemize}
  \item First item
  \item Second item
  \item Third item with sub-items:
    \begin{itemize}
      \item Sub-item 1
      \item Sub-item 2
    \end{itemize}
\end{itemize}

\subsection{Numbered List}
\begin{enumerate}
  \item First step
  \item Second step
  \item Third step
\end{enumerate}

\section{Citations and References}
When you use the citation search feature, references will be automatically managed. For example, you might cite a paper like this: \cite{example2024}. All citations will be automatically added to your \texttt{cite.bib} file.

\section{Conclusion}
This template demonstrates basic \LaTeX{} features including mathematical formulas, tables, and lists. The bibliography system is now set up to work seamlessly with Bibby AI's citation management. Feel free to modify it for your needs!

\bibliographystyle{plain}
\bibliography{cite}

\end{document}\documentclass{article}
\usepackage[utf8]{inputenc}
\usepackage{booktabs}
\usepackage{amsmath}
\usepackage{natbib}

\title{My Document}
\author{Your Name}
\date{\today}

\begin{document}

\maketitle

\section{Introduction}
This is a sample document demonstrating various \LaTeX{} features including mathematical equations, tables, and formatting. You can cite papers using the citation search feature, and they will be automatically added to your \texttt{cite.bib} file.

\section{Mathematical Expressions}
Here are some example equations:

\subsection{Quadratic Formula}
The quadratic formula is given by:
\[x = \frac{-b \pm \sqrt{b^2 - 4ac}}{2a}\]

\subsection{Matrix Example}
A sample matrix:
\[
\begin{pmatrix}
a_{11} & a_{12} & a_{13} \\
a_{21} & a_{22} & a_{23} \\
a_{31} & a_{32} & a_{33}
\end{pmatrix}
\]

\section{Tables}
Below is a sample table showing data:

\begin{table}[h]
\centering
\begin{tabular}{@{}lrrr@{}}
\toprule
\textbf{Category} & \textbf{Value 1} & \textbf{Value 2} & \textbf{Total} \\
\midrule
Item A & 10.5 & 20.3 & 30.8 \\
Item B & 15.2 & 25.1 & 40.3 \\
Item C & 22.1 & 12.7 & 34.8 \\
\bottomrule
\end{tabular}
\caption{Sample Data Table}
\label{tab:sample}
\end{table}

\section{Lists}
\subsection{Bullet Points}
\begin{itemize}
  \item First item
  \item Second item
  \item Third item with sub-items:
    \begin{itemize}
      \item Sub-item 1
      \item Sub-item 2
    \end{itemize}
\end{itemize}

\subsection{Numbered List}
\begin{enumerate}
  \item First step
  \item Second step
  \item Third step
\end{enumerate}

\section{Citations and References}
When you use the citation search feature, references will be automatically managed. For example, you might cite a paper like this: \cite{example2024}. All citations will be automatically added to your \texttt{cite.bib} file.

\section{Conclusion}
This template demonstrates basic \LaTeX{} features including mathematical formulas, tables, and lists. The bibliography system is now set up to work seamlessly with Bibby AI's citation management. Feel free to modify it for your needs!

\bibliographystyle{plain}
\bibliography{cite}

\end{document}\documentclass{article}
\usepackage[utf8]{inputenc}
\usepackage{booktabs}
\usepackage{amsmath}
\usepackage{natbib}

\title{My Document}
\author{Your Name}
\date{\today}

\begin{document}

\maketitle

\section{Introduction}
This is a sample document demonstrating various \LaTeX{} features including mathematical equations, tables, and formatting. You can cite papers using the citation search feature, and they will be automatically added to your \texttt{cite.bib} file.

\section{Mathematical Expressions}
Here are some example equations:

\subsection{Quadratic Formula}
The quadratic formula is given by:
\[x = \frac{-b \pm \sqrt{b^2 - 4ac}}{2a}\]

\subsection{Matrix Example}
A sample matrix:
\[
\begin{pmatrix}
a_{11} & a_{12} & a_{13} \\
a_{21} & a_{22} & a_{23} \\
a_{31} & a_{32} & a_{33}
\end{pmatrix}
\]

\section{Tables}
Below is a sample table showing data:

\begin{table}[h]
\centering
\begin{tabular}{@{}lrrr@{}}
\toprule
\textbf{Category} & \textbf{Value 1} & \textbf{Value 2} & \textbf{Total} \\
\midrule
Item A & 10.5 & 20.3 & 30.8 \\
Item B & 15.2 & 25.1 & 40.3 \\
Item C & 22.1 & 12.7 & 34.8 \\
\bottomrule
\end{tabular}
\caption{Sample Data Table}
\label{tab:sample}
\end{table}

\section{Lists}
\subsection{Bullet Points}
\begin{itemize}
  \item First item
  \item Second item
  \item Third item with sub-items:
    \begin{itemize}
      \item Sub-item 1
      \item Sub-item 2
    \end{itemize}
\end{itemize}

\subsection{Numbered List}
\begin{enumerate}
  \item First step
  \item Second step
  \item Third step
\end{enumerate}

\section{Citations and References}
When you use the citation search feature, references will be automatically managed. For example, you might cite a paper like this: \cite{example2024}. All citations will be automatically added to your \texttt{cite.bib} file.

\section{Conclusion}
This template demonstrates basic \LaTeX{} features including mathematical formulas, tables, and lists. The bibliography system is now set up to work seamlessly with Bibby AI's citation management. Feel free to modify it for your needs!

\bibliographystyle{plain}
\bibliography{cite}

\end{document}\documentclass{article}
\usepackage[utf8]{inputenc}
\usepackage{booktabs}
\usepackage{amsmath}
\usepackage{natbib}

\title{My Document}
\author{Your Name}
\date{\today}

\begin{document}

\maketitle

\section{Introduction}
This is a sample document demonstrating various \LaTeX{} features including mathematical equations, tables, and formatting. You can cite papers using the citation search feature, and they will be automatically added to your \texttt{cite.bib} file.

\section{Mathematical Expressions}
Here are some example equations:

\subsection{Quadratic Formula}
The quadratic formula is given by:
\[x = \frac{-b \pm \sqrt{b^2 - 4ac}}{2a}\]

\subsection{Matrix Example}
A sample matrix:
\[
\begin{pmatrix}
a_{11} & a_{12} & a_{13} \\
a_{21} & a_{22} & a_{23} \\
a_{31} & a_{32} & a_{33}
\end{pmatrix}
\]

\section{Tables}
Below is a sample table showing data:

\begin{table}[h]
\centering
\begin{tabular}{@{}lrrr@{}}
\toprule
\textbf{Category} & \textbf{Value 1} & \textbf{Value 2} & \textbf{Total} \\
\midrule
Item A & 10.5 & 20.3 & 30.8 \\
Item B & 15.2 & 25.1 & 40.3 \\
Item C & 22.1 & 12.7 & 34.8 \\
\bottomrule
\end{tabular}
\caption{Sample Data Table}
\label{tab:sample}
\end{table}

\section{Lists}
\subsection{Bullet Points}
\begin{itemize}
  \item First item
  \item Second item
  \item Third item with sub-items:
    \begin{itemize}
      \item Sub-item 1
      \item Sub-item 2
    \end{itemize}
\end{itemize}

\subsection{Numbered List}
\begin{enumerate}
  \item First step
  \item Second step
  \item Third step
\end{enumerate}

\section{Citations and References}
When you use the citation search feature, references will be automatically managed. For example, you might cite a paper like this: \cite{example2024}. All citations will be automatically added to your \texttt{cite.bib} file.

\section{Conclusion}
This template demonstrates basic \LaTeX{} features including mathematical formulas, tables, and lists. The bibliography system is now set up to work seamlessly with Bibby AI's citation management. Feel free to modify it for your needs!

\bibliographystyle{plain}
\bibliography{cite}

\end{document}\documentclass{article}
\usepackage[utf8]{inputenc}
\usepackage{booktabs}
\usepackage{amsmath}
\usepackage{natbib}

\title{My Document}
\author{Your Name}
\date{\today}

\begin{document}

\maketitle

\section{Introduction}
This is a sample document demonstrating various \LaTeX{} features including mathematical equations, tables, and formatting. You can cite papers using the citation search feature, and they will be automatically added to your \texttt{cite.bib} file.

\section{Mathematical Expressions}
Here are some example equations:

\subsection{Quadratic Formula}
The quadratic formula is given by:
\[x = \frac{-b \pm \sqrt{b^2 - 4ac}}{2a}\]

\subsection{Matrix Example}
A sample matrix:
\[
\begin{pmatrix}
a_{11} & a_{12} & a_{13} \\
a_{21} & a_{22} & a_{23} \\
a_{31} & a_{32} & a_{33}
\end{pmatrix}
\]

\section{Tables}
Below is a sample table showing data:

\begin{table}[h]
\centering
\begin{tabular}{@{}lrrr@{}}
\toprule
\textbf{Category} & \textbf{Value 1} & \textbf{Value 2} & \textbf{Total} \\
\midrule
Item A & 10.5 & 20.3 & 30.8 \\
Item B & 15.2 & 25.1 & 40.3 \\
Item C & 22.1 & 12.7 & 34.8 \\
\bottomrule
\end{tabular}
\caption{Sample Data Table}
\label{tab:sample}
\end{table}

\section{Lists}
\subsection{Bullet Points}
\begin{itemize}
  \item First item
  \item Second item
  \item Third item with sub-items:
    \begin{itemize}
      \item Sub-item 1
      \item Sub-item 2
    \end{itemize}
\end{itemize}

\subsection{Numbered List}
\begin{enumerate}
  \item First step
  \item Second step
  \item Third step
\end{enumerate}

\section{Citations and References}
When you use the citation search feature, references will be automatically managed. For example, you might cite a paper like this: \cite{example2024}. All citations will be automatically added to your \texttt{cite.bib} file.

\section{Conclusion}
This template demonstrates basic \LaTeX{} features including mathematical formulas, tables, and lists. The bibliography system is now set up to work seamlessly with Bibby AI's citation management. Feel free to modify it for your needs!

\bibliographystyle{plain}
\bibliography{cite}

\end{document}\documentclass{article}
\usepackage[utf8]{inputenc}
\usepackage{booktabs}
\usepackage{amsmath}
\usepackage{natbib}

\title{My Document}
\author{Your Name}
\date{\today}

\begin{document}

\maketitle

\section{Introduction}
This is a sample document demonstrating various \LaTeX{} features including mathematical equations, tables, and formatting. You can cite papers using the citation search feature, and they will be automatically added to your \texttt{cite.bib} file.

\section{Mathematical Expressions}
Here are some example equations:

\subsection{Quadratic Formula}
The quadratic formula is given by:
\[x = \frac{-b \pm \sqrt{b^2 - 4ac}}{2a}\]

\subsection{Matrix Example}
A sample matrix:
\[
\begin{pmatrix}
a_{11} & a_{12} & a_{13} \\
a_{21} & a_{22} & a_{23} \\
a_{31} & a_{32} & a_{33}
\end{pmatrix}
\]

\section{Tables}
Below is a sample table showing data:

\begin{table}[h]
\centering
\begin{tabular}{@{}lrrr@{}}
\toprule
\textbf{Category} & \textbf{Value 1} & \textbf{Value 2} & \textbf{Total} \\
\midrule
Item A & 10.5 & 20.3 & 30.8 \\
Item B & 15.2 & 25.1 & 40.3 \\
Item C & 22.1 & 12.7 & 34.8 \\
\bottomrule
\end{tabular}
\caption{Sample Data Table}
\label{tab:sample}
\end{table}

\section{Lists}
\subsection{Bullet Points}
\begin{itemize}
  \item First item
  \item Second item
  \item Third item with sub-items:
    \begin{itemize}
      \item Sub-item 1
      \item Sub-item 2
    \end{itemize}
\end{itemize}

\subsection{Numbered List}
\begin{enumerate}
  \item First step
  \item Second step
  \item Third step
\end{enumerate}

\section{Citations and References}
When you use the citation search feature, references will be automatically managed. For example, you might cite a paper like this: \cite{example2024}. All citations will be automatically added to your \texttt{cite.bib} file.

\section{Conclusion}
This template demonstrates basic \LaTeX{} features including mathematical formulas, tables, and lists. The bibliography system is now set up to work seamlessly with Bibby AI's citation management. Feel free to modify it for your needs!

\bibliographystyle{plain}
\bibliography{cite}

\end{document}\documentclass{article}
\usepackage[utf8]{inputenc}
\usepackage{booktabs}
\usepackage{amsmath}
\usepackage{natbib}

\title{My Document}
\author{Your Name}
\date{\today}

\begin{document}

\maketitle

\section{Introduction}
This is a sample document demonstrating various \LaTeX{} features including mathematical equations, tables, and formatting. You can cite papers using the citation search feature, and they will be automatically added to your \texttt{cite.bib} file.

\section{Mathematical Expressions}
Here are some example equations:

\subsection{Quadratic Formula}
The quadratic formula is given by:
\[x = \frac{-b \pm \sqrt{b^2 - 4ac}}{2a}\]

\subsection{Matrix Example}
A sample matrix:
\[
\begin{pmatrix}
a_{11} & a_{12} & a_{13} \\
a_{21} & a_{22} & a_{23} \\
a_{31} & a_{32} & a_{33}
\end{pmatrix}
\]

\section{Tables}
Below is a sample table showing data:

\begin{table}[h]
\centering
\begin{tabular}{@{}lrrr@{}}
\toprule
\textbf{Category} & \textbf{Value 1} & \textbf{Value 2} & \textbf{Total} \\
\midrule
Item A & 10.5 & 20.3 & 30.8 \\
Item B & 15.2 & 25.1 & 40.3 \\
Item C & 22.1 & 12.7 & 34.8 \\
\bottomrule
\end{tabular}
\caption{Sample Data Table}
\label{tab:sample}
\end{table}

\section{Lists}
\subsection{Bullet Points}
\begin{itemize}
  \item First item
  \item Second item
  \item Third item with sub-items:
    \begin{itemize}
      \item Sub-item 1
      \item Sub-item 2
    \end{itemize}
\end{itemize}

\subsection{Numbered List}
\begin{enumerate}
  \item First step
  \item Second step
  \item Third step
\end{enumerate}

\section{Citations and References}
When you use the citation search feature, references will be automatically managed. For example, you might cite a paper like this: \cite{example2024}. All citations will be automatically added to your \texttt{cite.bib} file.

\section{Conclusion}
This template demonstrates basic \LaTeX{} features including mathematical formulas, tables, and lists. The bibliography system is now set up to work seamlessly with Bibby AI's citation management. Feel free to modify it for your needs!

\bibliographystyle{plain}
\bibliography{cite}

\end{document}\documentclass{article}
\usepackage[utf8]{inputenc}
\usepackage{booktabs}
\usepackage{amsmath}
\usepackage{natbib}

\title{My Document}
\author{Your Name}
\date{\today}

\begin{document}

\maketitle

\section{Introduction}
This is a sample document demonstrating various \LaTeX{} features including mathematical equations, tables, and formatting. You can cite papers using the citation search feature, and they will be automatically added to your \texttt{cite.bib} file.

\section{Mathematical Expressions}
Here are some example equations:

\subsection{Quadratic Formula}
The quadratic formula is given by:
\[x = \frac{-b \pm \sqrt{b^2 - 4ac}}{2a}\]

\subsection{Matrix Example}
A sample matrix:
\[
\begin{pmatrix}
a_{11} & a_{12} & a_{13} \\
a_{21} & a_{22} & a_{23} \\
a_{31} & a_{32} & a_{33}
\end{pmatrix}
\]

\section{Tables}
Below is a sample table showing data:

\begin{table}[h]
\centering
\begin{tabular}{@{}lrrr@{}}
\toprule
\textbf{Category} & \textbf{Value 1} & \textbf{Value 2} & \textbf{Total} \\
\midrule
Item A & 10.5 & 20.3 & 30.8 \\
Item B & 15.2 & 25.1 & 40.3 \\
Item C & 22.1 & 12.7 & 34.8 \\
\bottomrule
\end{tabular}
\caption{Sample Data Table}
\label{tab:sample}
\end{table}

\section{Lists}
\subsection{Bullet Points}
\begin{itemize}
  \item First item
  \item Second item
  \item Third item with sub-items:
    \begin{itemize}
      \item Sub-item 1
      \item Sub-item 2
    \end{itemize}
\end{itemize}

\subsection{Numbered List}
\begin{enumerate}
  \item First step
  \item Second step
  \item Third step
\end{enumerate}

\section{Citations and References}
When you use the citation search feature, references will be automatically managed. For example, you might cite a paper like this: \cite{example2024}. All citations will be automatically added to your \texttt{cite.bib} file.

\section{Conclusion}
This template demonstrates basic \LaTeX{} features including mathematical formulas, tables, and lists. The bibliography system is now set up to work seamlessly with Bibby AI's citation management. Feel free to modify it for your needs!

\bibliographystyle{plain}
\bibliography{cite}

\end{document}\documentclass{article}
\usepackage[utf8]{inputenc}
\usepackage{booktabs}
\usepackage{amsmath}
\usepackage{natbib}

\title{My Document}
\author{Your Name}
\date{\today}

\begin{document}

\maketitle

\section{Introduction}
This is a sample document demonstrating various \LaTeX{} features including mathematical equations, tables, and formatting. You can cite papers using the citation search feature, and they will be automatically added to your \texttt{cite.bib} file.

\section{Mathematical Expressions}
Here are some example equations:

\subsection{Quadratic Formula}
The quadratic formula is given by:
\[x = \frac{-b \pm \sqrt{b^2 - 4ac}}{2a}\]

\subsection{Matrix Example}
A sample matrix:
\[
\begin{pmatrix}
a_{11} & a_{12} & a_{13} \\
a_{21} & a_{22} & a_{23} \\
a_{31} & a_{32} & a_{33}
\end{pmatrix}
\]

\section{Tables}
Below is a sample table showing data:

\begin{table}[h]
\centering
\begin{tabular}{@{}lrrr@{}}
\toprule
\textbf{Category} & \textbf{Value 1} & \textbf{Value 2} & \textbf{Total} \\
\midrule
Item A & 10.5 & 20.3 & 30.8 \\
Item B & 15.2 & 25.1 & 40.3 \\
Item C & 22.1 & 12.7 & 34.8 \\
\bottomrule
\end{tabular}
\caption{Sample Data Table}
\label{tab:sample}
\end{table}

\section{Lists}
\subsection{Bullet Points}
\begin{itemize}
  \item First item
  \item Second item
  \item Third item with sub-items:
    \begin{itemize}
      \item Sub-item 1
      \item Sub-item 2
    \end{itemize}
\end{itemize}

\subsection{Numbered List}
\begin{enumerate}
  \item First step
  \item Second step
  \item Third step
\end{enumerate}

\section{Citations and References}
When you use the citation search feature, references will be automatically managed. For example, you might cite a paper like this: \cite{example2024}. All citations will be automatically added to your \texttt{cite.bib} file.

\section{Conclusion}
This template demonstrates basic \LaTeX{} features including mathematical formulas, tables, and lists. The bibliography system is now set up to work seamlessly with Bibby AI's citation management. Feel free to modify it for your needs!

\bibliographystyle{plain}
\bibliography{cite}

\end{document}\documentclass{article}
\usepackage[utf8]{inputenc}
\usepackage{booktabs}
\usepackage{amsmath}
\usepackage{natbib}

\title{My Document}
\author{Your Name}
\date{\today}

\begin{document}

\maketitle

\section{Introduction}
This is a sample document demonstrating various \LaTeX{} features including mathematical equations, tables, and formatting. You can cite papers using the citation search feature, and they will be automatically added to your \texttt{cite.bib} file.

\section{Mathematical Expressions}
Here are some example equations:

\subsection{Quadratic Formula}
The quadratic formula is given by:
\[x = \frac{-b \pm \sqrt{b^2 - 4ac}}{2a}\]

\subsection{Matrix Example}
A sample matrix:
\[
\begin{pmatrix}
a_{11} & a_{12} & a_{13} \\
a_{21} & a_{22} & a_{23} \\
a_{31} & a_{32} & a_{33}
\end{pmatrix}
\]

\section{Tables}
Below is a sample table showing data:

\begin{table}[h]
\centering
\begin{tabular}{@{}lrrr@{}}
\toprule
\textbf{Category} & \textbf{Value 1} & \textbf{Value 2} & \textbf{Total} \\
\midrule
Item A & 10.5 & 20.3 & 30.8 \\
Item B & 15.2 & 25.1 & 40.3 \\
Item C & 22.1 & 12.7 & 34.8 \\
\bottomrule
\end{tabular}
\caption{Sample Data Table}
\label{tab:sample}
\end{table}

\section{Lists}
\subsection{Bullet Points}
\begin{itemize}
  \item First item
  \item Second item
  \item Third item with sub-items:
    \begin{itemize}
      \item Sub-item 1
      \item Sub-item 2
    \end{itemize}
\end{itemize}

\subsection{Numbered List}
\begin{enumerate}
  \item First step
  \item Second step
  \item Third step
\end{enumerate}

\section{Citations and References}
When you use the citation search feature, references will be automatically managed. For example, you might cite a paper like this: \cite{example2024}. All citations will be automatically added to your \texttt{cite.bib} file.

\section{Conclusion}
This template demonstrates basic \LaTeX{} features including mathematical formulas, tables, and lists. The bibliography system is now set up to work seamlessly with Bibby AI's citation management. Feel free to modify it for your needs!

\bibliographystyle{plain}
\bibliography{cite}

\end{document}\documentclass{article}
\usepackage[utf8]{inputenc}
\usepackage{booktabs}
\usepackage{amsmath}
\usepackage{natbib}

\title{My Document}
\author{Your Name}
\date{\today}

\begin{document}

\maketitle

\section{Introduction}
This is a sample document demonstrating various \LaTeX{} features including mathematical equations, tables, and formatting. You can cite papers using the citation search feature, and they will be automatically added to your \texttt{cite.bib} file.

\section{Mathematical Expressions}
Here are some example equations:

\subsection{Quadratic Formula}
The quadratic formula is given by:
\[x = \frac{-b \pm \sqrt{b^2 - 4ac}}{2a}\]

\subsection{Matrix Example}
A sample matrix:
\[
\begin{pmatrix}
a_{11} & a_{12} & a_{13} \\
a_{21} & a_{22} & a_{23} \\
a_{31} & a_{32} & a_{33}
\end{pmatrix}
\]

\section{Tables}
Below is a sample table showing data:

\begin{table}[h]
\centering
\begin{tabular}{@{}lrrr@{}}
\toprule
\textbf{Category} & \textbf{Value 1} & \textbf{Value 2} & \textbf{Total} \\
\midrule
Item A & 10.5 & 20.3 & 30.8 \\
Item B & 15.2 & 25.1 & 40.3 \\
Item C & 22.1 & 12.7 & 34.8 \\
\bottomrule
\end{tabular}
\caption{Sample Data Table}
\label{tab:sample}
\end{table}

\section{Lists}
\subsection{Bullet Points}
\begin{itemize}
  \item First item
  \item Second item
  \item Third item with sub-items:
    \begin{itemize}
      \item Sub-item 1
      \item Sub-item 2
    \end{itemize}
\end{itemize}

\subsection{Numbered List}
\begin{enumerate}
  \item First step
  \item Second step
  \item Third step
\end{enumerate}

\section{Citations and References}
When you use the citation search feature, references will be automatically managed. For example, you might cite a paper like this: \cite{example2024}. All citations will be automatically added to your \texttt{cite.bib} file.

\section{Conclusion}
This template demonstrates basic \LaTeX{} features including mathematical formulas, tables, and lists. The bibliography system is now set up to work seamlessly with Bibby AI's citation management. Feel free to modify it for your needs!

\bibliographystyle{plain}
\bibliography{cite}

\end{document}\documentclass{article}
\usepackage[utf8]{inputenc}
\usepackage{booktabs}
\usepackage{amsmath}
\usepackage{natbib}

\title{My Document}
\author{Your Name}
\date{\today}

\begin{document}

\maketitle

\section{Introduction}
This is a sample document demonstrating various \LaTeX{} features including mathematical equations, tables, and formatting. You can cite papers using the citation search feature, and they will be automatically added to your \texttt{cite.bib} file.

\section{Mathematical Expressions}
Here are some example equations:

\subsection{Quadratic Formula}
The quadratic formula is given by:
\[x = \frac{-b \pm \sqrt{b^2 - 4ac}}{2a}\]

\subsection{Matrix Example}
A sample matrix:
\[
\begin{pmatrix}
a_{11} & a_{12} & a_{13} \\
a_{21} & a_{22} & a_{23} \\
a_{31} & a_{32} & a_{33}
\end{pmatrix}
\]

\section{Tables}
Below is a sample table showing data:

\begin{table}[h]
\centering
\begin{tabular}{@{}lrrr@{}}
\toprule
\textbf{Category} & \textbf{Value 1} & \textbf{Value 2} & \textbf{Total} \\
\midrule
Item A & 10.5 & 20.3 & 30.8 \\
Item B & 15.2 & 25.1 & 40.3 \\
Item C & 22.1 & 12.7 & 34.8 \\
\bottomrule
\end{tabular}
\caption{Sample Data Table}
\label{tab:sample}
\end{table}

\section{Lists}
\subsection{Bullet Points}
\begin{itemize}
  \item First item
  \item Second item
  \item Third item with sub-items:
    \begin{itemize}
      \item Sub-item 1
      \item Sub-item 2
    \end{itemize}
\end{itemize}

\subsection{Numbered List}
\begin{enumerate}
  \item First step
  \item Second step
  \item Third step
\end{enumerate}

\section{Citations and References}
When you use the citation search feature, references will be automatically managed. For example, you might cite a paper like this: \cite{example2024}. All citations will be automatically added to your \texttt{cite.bib} file.

\section{Conclusion}
This template demonstrates basic \LaTeX{} features including mathematical formulas, tables, and lists. The bibliography system is now set up to work seamlessly with Bibby AI's citation management. Feel free to modify it for your needs!

\bibliographystyle{plain}
\bibliography{cite}

\end{document}\documentclass{article}
\usepackage[utf8]{inputenc}
\usepackage{booktabs}
\usepackage{amsmath}
\usepackage{natbib}

\title{My Document}
\author{Your Name}
\date{\today}

\begin{document}

\maketitle

\section{Introduction}
This is a sample document demonstrating various \LaTeX{} features including mathematical equations, tables, and formatting. You can cite papers using the citation search feature, and they will be automatically added to your \texttt{cite.bib} file.

\section{Mathematical Expressions}
Here are some example equations:

\subsection{Quadratic Formula}
The quadratic formula is given by:
\[x = \frac{-b \pm \sqrt{b^2 - 4ac}}{2a}\]

\subsection{Matrix Example}
A sample matrix:
\[
\begin{pmatrix}
a_{11} & a_{12} & a_{13} \\
a_{21} & a_{22} & a_{23} \\
a_{31} & a_{32} & a_{33}
\end{pmatrix}
\]

\section{Tables}
Below is a sample table showing data:

\begin{table}[h]
\centering
\begin{tabular}{@{}lrrr@{}}
\toprule
\textbf{Category} & \textbf{Value 1} & \textbf{Value 2} & \textbf{Total} \\
\midrule
Item A & 10.5 & 20.3 & 30.8 \\
Item B & 15.2 & 25.1 & 40.3 \\
Item C & 22.1 & 12.7 & 34.8 \\
\bottomrule
\end{tabular}
\caption{Sample Data Table}
\label{tab:sample}
\end{table}

\section{Lists}
\subsection{Bullet Points}
\begin{itemize}
  \item First item
  \item Second item
  \item Third item with sub-items:
    \begin{itemize}
      \item Sub-item 1
      \item Sub-item 2
    \end{itemize}
\end{itemize}

\subsection{Numbered List}
\begin{enumerate}
  \item First step
  \item Second step
  \item Third step
\end{enumerate}

\section{Citations and References}
When you use the citation search feature, references will be automatically managed. For example, you might cite a paper like this: \cite{example2024}. All citations will be automatically added to your \texttt{cite.bib} file.

\section{Conclusion}
This template demonstrates basic \LaTeX{} features including mathematical formulas, tables, and lists. The bibliography system is now set up to work seamlessly with Bibby AI's citation management. Feel free to modify it for your needs!

\bibliographystyle{plain}
\bibliography{cite}

\end{document}\documentclass{article}
\usepackage[utf8]{inputenc}
\usepackage{booktabs}
\usepackage{amsmath}
\usepackage{natbib}

\title{My Document}
\author{Your Name}
\date{\today}

\begin{document}

\maketitle

\section{Introduction}
This is a sample document demonstrating various \LaTeX{} features including mathematical equations, tables, and formatting. You can cite papers using the citation search feature, and they will be automatically added to your \texttt{cite.bib} file.

\section{Mathematical Expressions}
Here are some example equations:

\subsection{Quadratic Formula}
The quadratic formula is given by:
\[x = \frac{-b \pm \sqrt{b^2 - 4ac}}{2a}\]

\subsection{Matrix Example}
A sample matrix:
\[
\begin{pmatrix}
a_{11} & a_{12} & a_{13} \\
a_{21} & a_{22} & a_{23} \\
a_{31} & a_{32} & a_{33}
\end{pmatrix}
\]

\section{Tables}
Below is a sample table showing data:

\begin{table}[h]
\centering
\begin{tabular}{@{}lrrr@{}}
\toprule
\textbf{Category} & \textbf{Value 1} & \textbf{Value 2} & \textbf{Total} \\
\midrule
Item A & 10.5 & 20.3 & 30.8 \\
Item B & 15.2 & 25.1 & 40.3 \\
Item C & 22.1 & 12.7 & 34.8 \\
\bottomrule
\end{tabular}
\caption{Sample Data Table}
\label{tab:sample}
\end{table}

\section{Lists}
\subsection{Bullet Points}
\begin{itemize}
  \item First item
  \item Second item
  \item Third item with sub-items:
    \begin{itemize}
      \item Sub-item 1
      \item Sub-item 2
    \end{itemize}
\end{itemize}

\subsection{Numbered List}
\begin{enumerate}
  \item First step
  \item Second step
  \item Third step
\end{enumerate}

\section{Citations and References}
When you use the citation search feature, references will be automatically managed. For example, you might cite a paper like this: \cite{example2024}. All citations will be automatically added to your \texttt{cite.bib} file.

\section{Conclusion}
This template demonstrates basic \LaTeX{} features including mathematical formulas, tables, and lists. The bibliography system is now set up to work seamlessly with Bibby AI's citation management. Feel free to modify it for your needs!

\bibliographystyle{plain}
\bibliography{cite}

\end{document}\documentclass{article}
\usepackage[utf8]{inputenc}
\usepackage{booktabs}
\usepackage{amsmath}
\usepackage{natbib}

\title{My Document}
\author{Your Name}
\date{\today}

\begin{document}

\maketitle

\section{Introduction}
This is a sample document demonstrating various \LaTeX{} features including mathematical equations, tables, and formatting. You can cite papers using the citation search feature, and they will be automatically added to your \texttt{cite.bib} file.

\section{Mathematical Expressions}
Here are some example equations:

\subsection{Quadratic Formula}
The quadratic formula is given by:
\[x = \frac{-b \pm \sqrt{b^2 - 4ac}}{2a}\]

\subsection{Matrix Example}
A sample matrix:
\[
\begin{pmatrix}
a_{11} & a_{12} & a_{13} \\
a_{21} & a_{22} & a_{23} \\
a_{31} & a_{32} & a_{33}
\end{pmatrix}
\]

\section{Tables}
Below is a sample table showing data:

\begin{table}[h]
\centering
\begin{tabular}{@{}lrrr@{}}
\toprule
\textbf{Category} & \textbf{Value 1} & \textbf{Value 2} & \textbf{Total} \\
\midrule
Item A & 10.5 & 20.3 & 30.8 \\
Item B & 15.2 & 25.1 & 40.3 \\
Item C & 22.1 & 12.7 & 34.8 \\
\bottomrule
\end{tabular}
\caption{Sample Data Table}
\label{tab:sample}
\end{table}

\section{Lists}
\subsection{Bullet Points}
\begin{itemize}
  \item First item
  \item Second item
  \item Third item with sub-items:
    \begin{itemize}
      \item Sub-item 1
      \item Sub-item 2
    \end{itemize}
\end{itemize}

\subsection{Numbered List}
\begin{enumerate}
  \item First step
  \item Second step
  \item Third step
\end{enumerate}

\section{Citations and References}
When you use the citation search feature, references will be automatically managed. For example, you might cite a paper like this: \cite{example2024}. All citations will be automatically added to your \texttt{cite.bib} file.

\section{Conclusion}
This template demonstrates basic \LaTeX{} features including mathematical formulas, tables, and lists. The bibliography system is now set up to work seamlessly with Bibby AI's citation management. Feel free to modify it for your needs!

\bibliographystyle{plain}
\bibliography{cite}

\end{document}\documentclass{article}
\usepackage[utf8]{inputenc}
\usepackage{booktabs}
\usepackage{amsmath}
\usepackage{natbib}

\title{My Document}
\author{Your Name}
\date{\today}

\begin{document}

\maketitle

\section{Introduction}
This is a sample document demonstrating various \LaTeX{} features including mathematical equations, tables, and formatting. You can cite papers using the citation search feature, and they will be automatically added to your \texttt{cite.bib} file.

\section{Mathematical Expressions}
Here are some example equations:

\subsection{Quadratic Formula}
The quadratic formula is given by:
\[x = \frac{-b \pm \sqrt{b^2 - 4ac}}{2a}\]

\subsection{Matrix Example}
A sample matrix:
\[
\begin{pmatrix}
a_{11} & a_{12} & a_{13} \\
a_{21} & a_{22} & a_{23} \\
a_{31} & a_{32} & a_{33}
\end{pmatrix}
\]

\section{Tables}
Below is a sample table showing data:

\begin{table}[h]
\centering
\begin{tabular}{@{}lrrr@{}}
\toprule
\textbf{Category} & \textbf{Value 1} & \textbf{Value 2} & \textbf{Total} \\
\midrule
Item A & 10.5 & 20.3 & 30.8 \\
Item B & 15.2 & 25.1 & 40.3 \\
Item C & 22.1 & 12.7 & 34.8 \\
\bottomrule
\end{tabular}
\caption{Sample Data Table}
\label{tab:sample}
\end{table}

\section{Lists}
\subsection{Bullet Points}
\begin{itemize}
  \item First item
  \item Second item
  \item Third item with sub-items:
    \begin{itemize}
      \item Sub-item 1
      \item Sub-item 2
    \end{itemize}
\end{itemize}

\subsection{Numbered List}
\begin{enumerate}
  \item First step
  \item Second step
  \item Third step
\end{enumerate}

\section{Citations and References}
When you use the citation search feature, references will be automatically managed. For example, you might cite a paper like this: \cite{example2024}. All citations will be automatically added to your \texttt{cite.bib} file.

\section{Conclusion}
This template demonstrates basic \LaTeX{} features including mathematical formulas, tables, and lists. The bibliography system is now set up to work seamlessly with Bibby AI's citation management. Feel free to modify it for your needs!

\bibliographystyle{plain}
\bibliography{cite}

\end{document}\documentclass{article}
\usepackage[utf8]{inputenc}
\usepackage{booktabs}
\usepackage{amsmath}
\usepackage{natbib}

\title{My Document}
\author{Your Name}
\date{\today}

\begin{document}

\maketitle

\section{Introduction}
This is a sample document demonstrating various \LaTeX{} features including mathematical equations, tables, and formatting. You can cite papers using the citation search feature, and they will be automatically added to your \texttt{cite.bib} file.

\section{Mathematical Expressions}
Here are some example equations:

\subsection{Quadratic Formula}
The quadratic formula is given by:
\[x = \frac{-b \pm \sqrt{b^2 - 4ac}}{2a}\]

\subsection{Matrix Example}
A sample matrix:
\[
\begin{pmatrix}
a_{11} & a_{12} & a_{13} \\
a_{21} & a_{22} & a_{23} \\
a_{31} & a_{32} & a_{33}
\end{pmatrix}
\]

\section{Tables}
Below is a sample table showing data:

\begin{table}[h]
\centering
\begin{tabular}{@{}lrrr@{}}
\toprule
\textbf{Category} & \textbf{Value 1} & \textbf{Value 2} & \textbf{Total} \\
\midrule
Item A & 10.5 & 20.3 & 30.8 \\
Item B & 15.2 & 25.1 & 40.3 \\
Item C & 22.1 & 12.7 & 34.8 \\
\bottomrule
\end{tabular}
\caption{Sample Data Table}
\label{tab:sample}
\end{table}

\section{Lists}
\subsection{Bullet Points}
\begin{itemize}
  \item First item
  \item Second item
  \item Third item with sub-items:
    \begin{itemize}
      \item Sub-item 1
      \item Sub-item 2
    \end{itemize}
\end{itemize}

\subsection{Numbered List}
\begin{enumerate}
  \item First step
  \item Second step
  \item Third step
\end{enumerate}

\section{Citations and References}
When you use the citation search feature, references will be automatically managed. For example, you might cite a paper like this: \cite{example2024}. All citations will be automatically added to your \texttt{cite.bib} file.

\section{Conclusion}
This template demonstrates basic \LaTeX{} features including mathematical formulas, tables, and lists. The bibliography system is now set up to work seamlessly with Bibby AI's citation management. Feel free to modify it for your needs!

\bibliographystyle{plain}
\bibliography{cite}

\end{document}\documentclass{article}
\usepackage[utf8]{inputenc}
\usepackage{booktabs}
\usepackage{amsmath}
\usepackage{natbib}

\title{My Document}
\author{Your Name}
\date{\today}

\begin{document}

\maketitle

\section{Introduction}
This is a sample document demonstrating various \LaTeX{} features including mathematical equations, tables, and formatting. You can cite papers using the citation search feature, and they will be automatically added to your \texttt{cite.bib} file.

\section{Mathematical Expressions}
Here are some example equations:

\subsection{Quadratic Formula}
The quadratic formula is given by:
\[x = \frac{-b \pm \sqrt{b^2 - 4ac}}{2a}\]

\subsection{Matrix Example}
A sample matrix:
\[
\begin{pmatrix}
a_{11} & a_{12} & a_{13} \\
a_{21} & a_{22} & a_{23} \\
a_{31} & a_{32} & a_{33}
\end{pmatrix}
\]

\section{Tables}
Below is a sample table showing data:

\begin{table}[h]
\centering
\begin{tabular}{@{}lrrr@{}}
\toprule
\textbf{Category} & \textbf{Value 1} & \textbf{Value 2} & \textbf{Total} \\
\midrule
Item A & 10.5 & 20.3 & 30.8 \\
Item B & 15.2 & 25.1 & 40.3 \\
Item C & 22.1 & 12.7 & 34.8 \\
\bottomrule
\end{tabular}
\caption{Sample Data Table}
\label{tab:sample}
\end{table}

\section{Lists}
\subsection{Bullet Points}
\begin{itemize}
  \item First item
  \item Second item
  \item Third item with sub-items:
    \begin{itemize}
      \item Sub-item 1
      \item Sub-item 2
    \end{itemize}
\end{itemize}

\subsection{Numbered List}
\begin{enumerate}
  \item First step
  \item Second step
  \item Third step
\end{enumerate}

\section{Citations and References}
When you use the citation search feature, references will be automatically managed. For example, you might cite a paper like this: \cite{example2024}. All citations will be automatically added to your \texttt{cite.bib} file.

\section{Conclusion}
This template demonstrates basic \LaTeX{} features including mathematical formulas, tables, and lists. The bibliography system is now set up to work seamlessly with Bibby AI's citation management. Feel free to modify it for your needs!

\bibliographystyle{plain}
\bibliography{cite}

\end{document}\documentclass{article}
\usepackage[utf8]{inputenc}
\usepackage{booktabs}
\usepackage{amsmath}
\usepackage{natbib}

\title{My Document}
\author{Your Name}
\date{\today}

\begin{document}

\maketitle

\section{Introduction}
This is a sample document demonstrating various \LaTeX{} features including mathematical equations, tables, and formatting. You can cite papers using the citation search feature, and they will be automatically added to your \texttt{cite.bib} file.

\section{Mathematical Expressions}
Here are some example equations:

\subsection{Quadratic Formula}
The quadratic formula is given by:
\[x = \frac{-b \pm \sqrt{b^2 - 4ac}}{2a}\]

\subsection{Matrix Example}
A sample matrix:
\[
\begin{pmatrix}
a_{11} & a_{12} & a_{13} \\
a_{21} & a_{22} & a_{23} \\
a_{31} & a_{32} & a_{33}
\end{pmatrix}
\]

\section{Tables}
Below is a sample table showing data:

\begin{table}[h]
\centering
\begin{tabular}{@{}lrrr@{}}
\toprule
\textbf{Category} & \textbf{Value 1} & \textbf{Value 2} & \textbf{Total} \\
\midrule
Item A & 10.5 & 20.3 & 30.8 \\
Item B & 15.2 & 25.1 & 40.3 \\
Item C & 22.1 & 12.7 & 34.8 \\
\bottomrule
\end{tabular}
\caption{Sample Data Table}
\label{tab:sample}
\end{table}

\section{Lists}
\subsection{Bullet Points}
\begin{itemize}
  \item First item
  \item Second item
  \item Third item with sub-items:
    \begin{itemize}
      \item Sub-item 1
      \item Sub-item 2
    \end{itemize}
\end{itemize}

\subsection{Numbered List}
\begin{enumerate}
  \item First step
  \item Second step
  \item Third step
\end{enumerate}

\section{Citations and References}
When you use the citation search feature, references will be automatically managed. For example, you might cite a paper like this: \cite{example2024}. All citations will be automatically added to your \texttt{cite.bib} file.

\section{Conclusion}
This template demonstrates basic \LaTeX{} features including mathematical formulas, tables, and lists. The bibliography system is now set up to work seamlessly with Bibby AI's citation management. Feel free to modify it for your needs!

\bibliographystyle{plain}
\bibliography{cite}

\end{document}\documentclass{article}
\usepackage[utf8]{inputenc}
\usepackage{booktabs}
\usepackage{amsmath}
\usepackage{natbib}

\title{My Document}
\author{Your Name}
\date{\today}

\begin{document}

\maketitle

\section{Introduction}
This is a sample document demonstrating various \LaTeX{} features including mathematical equations, tables, and formatting. You can cite papers using the citation search feature, and they will be automatically added to your \texttt{cite.bib} file.

\section{Mathematical Expressions}
Here are some example equations:

\subsection{Quadratic Formula}
The quadratic formula is given by:
\[x = \frac{-b \pm \sqrt{b^2 - 4ac}}{2a}\]

\subsection{Matrix Example}
A sample matrix:
\[
\begin{pmatrix}
a_{11} & a_{12} & a_{13} \\
a_{21} & a_{22} & a_{23} \\
a_{31} & a_{32} & a_{33}
\end{pmatrix}
\]

\section{Tables}
Below is a sample table showing data:

\begin{table}[h]
\centering
\begin{tabular}{@{}lrrr@{}}
\toprule
\textbf{Category} & \textbf{Value 1} & \textbf{Value 2} & \textbf{Total} \\
\midrule
Item A & 10.5 & 20.3 & 30.8 \\
Item B & 15.2 & 25.1 & 40.3 \\
Item C & 22.1 & 12.7 & 34.8 \\
\bottomrule
\end{tabular}
\caption{Sample Data Table}
\label{tab:sample}
\end{table}

\section{Lists}
\subsection{Bullet Points}
\begin{itemize}
  \item First item
  \item Second item
  \item Third item with sub-items:
    \begin{itemize}
      \item Sub-item 1
      \item Sub-item 2
    \end{itemize}
\end{itemize}

\subsection{Numbered List}
\begin{enumerate}
  \item First step
  \item Second step
  \item Third step
\end{enumerate}

\section{Citations and References}
When you use the citation search feature, references will be automatically managed. For example, you might cite a paper like this: \cite{example2024}. All citations will be automatically added to your \texttt{cite.bib} file.

\section{Conclusion}
This template demonstrates basic \LaTeX{} features including mathematical formulas, tables, and lists. The bibliography system is now set up to work seamlessly with Bibby AI's citation management. Feel free to modify it for your needs!

\bibliographystyle{plain}
\bibliography{cite}

\end{document}\documentclass{article}
\usepackage[utf8]{inputenc}
\usepackage{booktabs}
\usepackage{amsmath}
\usepackage{natbib}

\title{My Document}
\author{Your Name}
\date{\today}

\begin{document}

\maketitle

\section{Introduction}
This is a sample document demonstrating various \LaTeX{} features including mathematical equations, tables, and formatting. You can cite papers using the citation search feature, and they will be automatically added to your \texttt{cite.bib} file.

\section{Mathematical Expressions}
Here are some example equations:

\subsection{Quadratic Formula}
The quadratic formula is given by:
\[x = \frac{-b \pm \sqrt{b^2 - 4ac}}{2a}\]

\subsection{Matrix Example}
A sample matrix:
\[
\begin{pmatrix}
a_{11} & a_{12} & a_{13} \\
a_{21} & a_{22} & a_{23} \\
a_{31} & a_{32} & a_{33}
\end{pmatrix}
\]

\section{Tables}
Below is a sample table showing data:

\begin{table}[h]
\centering
\begin{tabular}{@{}lrrr@{}}
\toprule
\textbf{Category} & \textbf{Value 1} & \textbf{Value 2} & \textbf{Total} \\
\midrule
Item A & 10.5 & 20.3 & 30.8 \\
Item B & 15.2 & 25.1 & 40.3 \\
Item C & 22.1 & 12.7 & 34.8 \\
\bottomrule
\end{tabular}
\caption{Sample Data Table}
\label{tab:sample}
\end{table}

\section{Lists}
\subsection{Bullet Points}
\begin{itemize}
  \item First item
  \item Second item
  \item Third item with sub-items:
    \begin{itemize}
      \item Sub-item 1
      \item Sub-item 2
    \end{itemize}
\end{itemize}

\subsection{Numbered List}
\begin{enumerate}
  \item First step
  \item Second step
  \item Third step
\end{enumerate}

\section{Citations and References}
When you use the citation search feature, references will be automatically managed. For example, you might cite a paper like this: \cite{example2024}. All citations will be automatically added to your \texttt{cite.bib} file.

\section{Conclusion}
This template demonstrates basic \LaTeX{} features including mathematical formulas, tables, and lists. The bibliography system is now set up to work seamlessly with Bibby AI's citation management. Feel free to modify it for your needs!

\bibliographystyle{plain}
\bibliography{cite}

\end{document}\documentclass{article}
\usepackage[utf8]{inputenc}
\usepackage{booktabs}
\usepackage{amsmath}
\usepackage{natbib}

\title{My Document}
\author{Your Name}
\date{\today}

\begin{document}

\maketitle

\section{Introduction}
This is a sample document demonstrating various \LaTeX{} features including mathematical equations, tables, and formatting. You can cite papers using the citation search feature, and they will be automatically added to your \texttt{cite.bib} file.

\section{Mathematical Expressions}
Here are some example equations:

\subsection{Quadratic Formula}
The quadratic formula is given by:
\[x = \frac{-b \pm \sqrt{b^2 - 4ac}}{2a}\]

\subsection{Matrix Example}
A sample matrix:
\[
\begin{pmatrix}
a_{11} & a_{12} & a_{13} \\
a_{21} & a_{22} & a_{23} \\
a_{31} & a_{32} & a_{33}
\end{pmatrix}
\]

\section{Tables}
Below is a sample table showing data:

\begin{table}[h]
\centering
\begin{tabular}{@{}lrrr@{}}
\toprule
\textbf{Category} & \textbf{Value 1} & \textbf{Value 2} & \textbf{Total} \\
\midrule
Item A & 10.5 & 20.3 & 30.8 \\
Item B & 15.2 & 25.1 & 40.3 \\
Item C & 22.1 & 12.7 & 34.8 \\
\bottomrule
\end{tabular}
\caption{Sample Data Table}
\label{tab:sample}
\end{table}

\section{Lists}
\subsection{Bullet Points}
\begin{itemize}
  \item First item
  \item Second item
  \item Third item with sub-items:
    \begin{itemize}
      \item Sub-item 1
      \item Sub-item 2
    \end{itemize}
\end{itemize}

\subsection{Numbered List}
\begin{enumerate}
  \item First step
  \item Second step
  \item Third step
\end{enumerate}

\section{Citations and References}
When you use the citation search feature, references will be automatically managed. For example, you might cite a paper like this: \cite{example2024}. All citations will be automatically added to your \texttt{cite.bib} file.

\section{Conclusion}
This template demonstrates basic \LaTeX{} features including mathematical formulas, tables, and lists. The bibliography system is now set up to work seamlessly with Bibby AI's citation management. Feel free to modify it for your needs!

\bibliographystyle{plain}
\bibliography{cite}

\end{document}\documentclass{article}
\usepackage[utf8]{inputenc}
\usepackage{booktabs}
\usepackage{amsmath}
\usepackage{natbib}

\title{My Document}
\author{Your Name}
\date{\today}

\begin{document}

\maketitle

\section{Introduction}
This is a sample document demonstrating various \LaTeX{} features including mathematical equations, tables, and formatting. You can cite papers using the citation search feature, and they will be automatically added to your \texttt{cite.bib} file.

\section{Mathematical Expressions}
Here are some example equations:

\subsection{Quadratic Formula}
The quadratic formula is given by:
\[x = \frac{-b \pm \sqrt{b^2 - 4ac}}{2a}\]

\subsection{Matrix Example}
A sample matrix:
\[
\begin{pmatrix}
a_{11} & a_{12} & a_{13} \\
a_{21} & a_{22} & a_{23} \\
a_{31} & a_{32} & a_{33}
\end{pmatrix}
\]

\section{Tables}
Below is a sample table showing data:

\begin{table}[h]
\centering
\begin{tabular}{@{}lrrr@{}}
\toprule
\textbf{Category} & \textbf{Value 1} & \textbf{Value 2} & \textbf{Total} \\
\midrule
Item A & 10.5 & 20.3 & 30.8 \\
Item B & 15.2 & 25.1 & 40.3 \\
Item C & 22.1 & 12.7 & 34.8 \\
\bottomrule
\end{tabular}
\caption{Sample Data Table}
\label{tab:sample}
\end{table}

\section{Lists}
\subsection{Bullet Points}
\begin{itemize}
  \item First item
  \item Second item
  \item Third item with sub-items:
    \begin{itemize}
      \item Sub-item 1
      \item Sub-item 2
    \end{itemize}
\end{itemize}

\subsection{Numbered List}
\begin{enumerate}
  \item First step
  \item Second step
  \item Third step
\end{enumerate}

\section{Citations and References}
When you use the citation search feature, references will be automatically managed. For example, you might cite a paper like this: \cite{example2024}. All citations will be automatically added to your \texttt{cite.bib} file.

\section{Conclusion}
This template demonstrates basic \LaTeX{} features including mathematical formulas, tables, and lists. The bibliography system is now set up to work seamlessly with Bibby AI's citation management. Feel free to modify it for your needs!

\bibliographystyle{plain}
\bibliography{cite}

\end{document}\documentclass{article}
\usepackage[utf8]{inputenc}
\usepackage{booktabs}
\usepackage{amsmath}
\usepackage{natbib}

\title{My Document}
\author{Your Name}
\date{\today}

\begin{document}

\maketitle

\section{Introduction}
This is a sample document demonstrating various \LaTeX{} features including mathematical equations, tables, and formatting. You can cite papers using the citation search feature, and they will be automatically added to your \texttt{cite.bib} file.

\section{Mathematical Expressions}
Here are some example equations:

\subsection{Quadratic Formula}
The quadratic formula is given by:
\[x = \frac{-b \pm \sqrt{b^2 - 4ac}}{2a}\]

\subsection{Matrix Example}
A sample matrix:
\[
\begin{pmatrix}
a_{11} & a_{12} & a_{13} \\
a_{21} & a_{22} & a_{23} \\
a_{31} & a_{32} & a_{33}
\end{pmatrix}
\]

\section{Tables}
Below is a sample table showing data:

\begin{table}[h]
\centering
\begin{tabular}{@{}lrrr@{}}
\toprule
\textbf{Category} & \textbf{Value 1} & \textbf{Value 2} & \textbf{Total} \\
\midrule
Item A & 10.5 & 20.3 & 30.8 \\
Item B & 15.2 & 25.1 & 40.3 \\
Item C & 22.1 & 12.7 & 34.8 \\
\bottomrule
\end{tabular}
\caption{Sample Data Table}
\label{tab:sample}
\end{table}

\section{Lists}
\subsection{Bullet Points}
\begin{itemize}
  \item First item
  \item Second item
  \item Third item with sub-items:
    \begin{itemize}
      \item Sub-item 1
      \item Sub-item 2
    \end{itemize}
\end{itemize}

\subsection{Numbered List}
\begin{enumerate}
  \item First step
  \item Second step
  \item Third step
\end{enumerate}

\section{Citations and References}
When you use the citation search feature, references will be automatically managed. For example, you might cite a paper like this: \cite{example2024}. All citations will be automatically added to your \texttt{cite.bib} file.

\section{Conclusion}
This template demonstrates basic \LaTeX{} features including mathematical formulas, tables, and lists. The bibliography system is now set up to work seamlessly with Bibby AI's citation management. Feel free to modify it for your needs!

\bibliographystyle{plain}
\bibliography{cite}

\end{document}\documentclass{article}
\usepackage[utf8]{inputenc}
\usepackage{booktabs}
\usepackage{amsmath}
\usepackage{natbib}

\title{My Document}
\author{Your Name}
\date{\today}

\begin{document}

\maketitle

\section{Introduction}
This is a sample document demonstrating various \LaTeX{} features including mathematical equations, tables, and formatting. You can cite papers using the citation search feature, and they will be automatically added to your \texttt{cite.bib} file.

\section{Mathematical Expressions}
Here are some example equations:

\subsection{Quadratic Formula}
The quadratic formula is given by:
\[x = \frac{-b \pm \sqrt{b^2 - 4ac}}{2a}\]

\subsection{Matrix Example}
A sample matrix:
\[
\begin{pmatrix}
a_{11} & a_{12} & a_{13} \\
a_{21} & a_{22} & a_{23} \\
a_{31} & a_{32} & a_{33}
\end{pmatrix}
\]

\section{Tables}
Below is a sample table showing data:

\begin{table}[h]
\centering
\begin{tabular}{@{}lrrr@{}}
\toprule
\textbf{Category} & \textbf{Value 1} & \textbf{Value 2} & \textbf{Total} \\
\midrule
Item A & 10.5 & 20.3 & 30.8 \\
Item B & 15.2 & 25.1 & 40.3 \\
Item C & 22.1 & 12.7 & 34.8 \\
\bottomrule
\end{tabular}
\caption{Sample Data Table}
\label{tab:sample}
\end{table}

\section{Lists}
\subsection{Bullet Points}
\begin{itemize}
  \item First item
  \item Second item
  \item Third item with sub-items:
    \begin{itemize}
      \item Sub-item 1
      \item Sub-item 2
    \end{itemize}
\end{itemize}

\subsection{Numbered List}
\begin{enumerate}
  \item First step
  \item Second step
  \item Third step
\end{enumerate}

\section{Citations and References}
When you use the citation search feature, references will be automatically managed. For example, you might cite a paper like this: \cite{example2024}. All citations will be automatically added to your \texttt{cite.bib} file.

\section{Conclusion}
This template demonstrates basic \LaTeX{} features including mathematical formulas, tables, and lists. The bibliography system is now set up to work seamlessly with Bibby AI's citation management. Feel free to modify it for your needs!

\bibliographystyle{plain}
\bibliography{cite}

\end{document}\documentclass{article}
\usepackage[utf8]{inputenc}
\usepackage{booktabs}
\usepackage{amsmath}
\usepackage{natbib}

\title{My Document}
\author{Your Name}
\date{\today}

\begin{document}

\maketitle

\section{Introduction}
This is a sample document demonstrating various \LaTeX{} features including mathematical equations, tables, and formatting. You can cite papers using the citation search feature, and they will be automatically added to your \texttt{cite.bib} file.

\section{Mathematical Expressions}
Here are some example equations:

\subsection{Quadratic Formula}
The quadratic formula is given by:
\[x = \frac{-b \pm \sqrt{b^2 - 4ac}}{2a}\]

\subsection{Matrix Example}
A sample matrix:
\[
\begin{pmatrix}
a_{11} & a_{12} & a_{13} \\
a_{21} & a_{22} & a_{23} \\
a_{31} & a_{32} & a_{33}
\end{pmatrix}
\]

\section{Tables}
Below is a sample table showing data:

\begin{table}[h]
\centering
\begin{tabular}{@{}lrrr@{}}
\toprule
\textbf{Category} & \textbf{Value 1} & \textbf{Value 2} & \textbf{Total} \\
\midrule
Item A & 10.5 & 20.3 & 30.8 \\
Item B & 15.2 & 25.1 & 40.3 \\
Item C & 22.1 & 12.7 & 34.8 \\
\bottomrule
\end{tabular}
\caption{Sample Data Table}
\label{tab:sample}
\end{table}

\section{Lists}
\subsection{Bullet Points}
\begin{itemize}
  \item First item
  \item Second item
  \item Third item with sub-items:
    \begin{itemize}
      \item Sub-item 1
      \item Sub-item 2
    \end{itemize}
\end{itemize}

\subsection{Numbered List}
\begin{enumerate}
  \item First step
  \item Second step
  \item Third step
\end{enumerate}

\section{Citations and References}
When you use the citation search feature, references will be automatically managed. For example, you might cite a paper like this: \cite{example2024}. All citations will be automatically added to your \texttt{cite.bib} file.

\section{Conclusion}
This template demonstrates basic \LaTeX{} features including mathematical formulas, tables, and lists. The bibliography system is now set up to work seamlessly with Bibby AI's citation management. Feel free to modify it for your needs!

\bibliographystyle{plain}
\bibliography{cite}

\end{document}\documentclass{article}
\usepackage[utf8]{inputenc}
\usepackage{booktabs}
\usepackage{amsmath}
\usepackage{natbib}

\title{My Document}
\author{Your Name}
\date{\today}

\begin{document}

\maketitle

\section{Introduction}
This is a sample document demonstrating various \LaTeX{} features including mathematical equations, tables, and formatting. You can cite papers using the citation search feature, and they will be automatically added to your \texttt{cite.bib} file.

\section{Mathematical Expressions}
Here are some example equations:

\subsection{Quadratic Formula}
The quadratic formula is given by:
\[x = \frac{-b \pm \sqrt{b^2 - 4ac}}{2a}\]

\subsection{Matrix Example}
A sample matrix:
\[
\begin{pmatrix}
a_{11} & a_{12} & a_{13} \\
a_{21} & a_{22} & a_{23} \\
a_{31} & a_{32} & a_{33}
\end{pmatrix}
\]

\section{Tables}
Below is a sample table showing data:

\begin{table}[h]
\centering
\begin{tabular}{@{}lrrr@{}}
\toprule
\textbf{Category} & \textbf{Value 1} & \textbf{Value 2} & \textbf{Total} \\
\midrule
Item A & 10.5 & 20.3 & 30.8 \\
Item B & 15.2 & 25.1 & 40.3 \\
Item C & 22.1 & 12.7 & 34.8 \\
\bottomrule
\end{tabular}
\caption{Sample Data Table}
\label{tab:sample}
\end{table}

\section{Lists}
\subsection{Bullet Points}
\begin{itemize}
  \item First item
  \item Second item
  \item Third item with sub-items:
    \begin{itemize}
      \item Sub-item 1
      \item Sub-item 2
    \end{itemize}
\end{itemize}

\subsection{Numbered List}
\begin{enumerate}
  \item First step
  \item Second step
  \item Third step
\end{enumerate}

\section{Citations and References}
When you use the citation search feature, references will be automatically managed. For example, you might cite a paper like this: \cite{example2024}. All citations will be automatically added to your \texttt{cite.bib} file.

\section{Conclusion}
This template demonstrates basic \LaTeX{} features including mathematical formulas, tables, and lists. The bibliography system is now set up to work seamlessly with Bibby AI's citation management. Feel free to modify it for your needs!

\bibliographystyle{plain}
\bibliography{cite}

\end{document}\documentclass{article}
\usepackage[utf8]{inputenc}
\usepackage{booktabs}
\usepackage{amsmath}
\usepackage{natbib}

\title{My Document}
\author{Your Name}
\date{\today}

\begin{document}

\maketitle

\section{Introduction}
This is a sample document demonstrating various \LaTeX{} features including mathematical equations, tables, and formatting. You can cite papers using the citation search feature, and they will be automatically added to your \texttt{cite.bib} file.

\section{Mathematical Expressions}
Here are some example equations:

\subsection{Quadratic Formula}
The quadratic formula is given by:
\[x = \frac{-b \pm \sqrt{b^2 - 4ac}}{2a}\]

\subsection{Matrix Example}
A sample matrix:
\[
\begin{pmatrix}
a_{11} & a_{12} & a_{13} \\
a_{21} & a_{22} & a_{23} \\
a_{31} & a_{32} & a_{33}
\end{pmatrix}
\]

\section{Tables}
Below is a sample table showing data:

\begin{table}[h]
\centering
\begin{tabular}{@{}lrrr@{}}
\toprule
\textbf{Category} & \textbf{Value 1} & \textbf{Value 2} & \textbf{Total} \\
\midrule
Item A & 10.5 & 20.3 & 30.8 \\
Item B & 15.2 & 25.1 & 40.3 \\
Item C & 22.1 & 12.7 & 34.8 \\
\bottomrule
\end{tabular}
\caption{Sample Data Table}
\label{tab:sample}
\end{table}

\section{Lists}
\subsection{Bullet Points}
\begin{itemize}
  \item First item
  \item Second item
  \item Third item with sub-items:
    \begin{itemize}
      \item Sub-item 1
      \item Sub-item 2
    \end{itemize}
\end{itemize}

\subsection{Numbered List}
\begin{enumerate}
  \item First step
  \item Second step
  \item Third step
\end{enumerate}

\section{Citations and References}
When you use the citation search feature, references will be automatically managed. For example, you might cite a paper like this: \cite{example2024}. All citations will be automatically added to your \texttt{cite.bib} file.

\section{Conclusion}
This template demonstrates basic \LaTeX{} features including mathematical formulas, tables, and lists. The bibliography system is now set up to work seamlessly with Bibby AI's citation management. Feel free to modify it for your needs!

\bibliographystyle{plain}
\bibliography{cite}

\end{document}\documentclass{article}
\usepackage[utf8]{inputenc}
\usepackage{booktabs}
\usepackage{amsmath}
\usepackage{natbib}

\title{My Document}
\author{Your Name}
\date{\today}

\begin{document}

\maketitle

\section{Introduction}
This is a sample document demonstrating various \LaTeX{} features including mathematical equations, tables, and formatting. You can cite papers using the citation search feature, and they will be automatically added to your \texttt{cite.bib} file.

\section{Mathematical Expressions}
Here are some example equations:

\subsection{Quadratic Formula}
The quadratic formula is given by:
\[x = \frac{-b \pm \sqrt{b^2 - 4ac}}{2a}\]

\subsection{Matrix Example}
A sample matrix:
\[
\begin{pmatrix}
a_{11} & a_{12} & a_{13} \\
a_{21} & a_{22} & a_{23} \\
a_{31} & a_{32} & a_{33}
\end{pmatrix}
\]

\section{Tables}
Below is a sample table showing data:

\begin{table}[h]
\centering
\begin{tabular}{@{}lrrr@{}}
\toprule
\textbf{Category} & \textbf{Value 1} & \textbf{Value 2} & \textbf{Total} \\
\midrule
Item A & 10.5 & 20.3 & 30.8 \\
Item B & 15.2 & 25.1 & 40.3 \\
Item C & 22.1 & 12.7 & 34.8 \\
\bottomrule
\end{tabular}
\caption{Sample Data Table}
\label{tab:sample}
\end{table}

\section{Lists}
\subsection{Bullet Points}
\begin{itemize}
  \item First item
  \item Second item
  \item Third item with sub-items:
    \begin{itemize}
      \item Sub-item 1
      \item Sub-item 2
    \end{itemize}
\end{itemize}

\subsection{Numbered List}
\begin{enumerate}
  \item First step
  \item Second step
  \item Third step
\end{enumerate}

\section{Citations and References}
When you use the citation search feature, references will be automatically managed. For example, you might cite a paper like this: \cite{example2024}. All citations will be automatically added to your \texttt{cite.bib} file.

\section{Conclusion}
This template demonstrates basic \LaTeX{} features including mathematical formulas, tables, and lists. The bibliography system is now set up to work seamlessly with Bibby AI's citation management. Feel free to modify it for your needs!

\bibliographystyle{plain}
\bibliography{cite}

\end{document}\documentclass{article}
\usepackage[utf8]{inputenc}
\usepackage{booktabs}
\usepackage{amsmath}
\usepackage{natbib}

\title{My Document}
\author{Your Name}
\date{\today}

\begin{document}

\maketitle

\section{Introduction}
This is a sample document demonstrating various \LaTeX{} features including mathematical equations, tables, and formatting. You can cite papers using the citation search feature, and they will be automatically added to your \texttt{cite.bib} file.

\section{Mathematical Expressions}
Here are some example equations:

\subsection{Quadratic Formula}
The quadratic formula is given by:
\[x = \frac{-b \pm \sqrt{b^2 - 4ac}}{2a}\]

\subsection{Matrix Example}
A sample matrix:
\[
\begin{pmatrix}
a_{11} & a_{12} & a_{13} \\
a_{21} & a_{22} & a_{23} \\
a_{31} & a_{32} & a_{33}
\end{pmatrix}
\]

\section{Tables}
Below is a sample table showing data:

\begin{table}[h]
\centering
\begin{tabular}{@{}lrrr@{}}
\toprule
\textbf{Category} & \textbf{Value 1} & \textbf{Value 2} & \textbf{Total} \\
\midrule
Item A & 10.5 & 20.3 & 30.8 \\
Item B & 15.2 & 25.1 & 40.3 \\
Item C & 22.1 & 12.7 & 34.8 \\
\bottomrule
\end{tabular}
\caption{Sample Data Table}
\label{tab:sample}
\end{table}

\section{Lists}
\subsection{Bullet Points}
\begin{itemize}
  \item First item
  \item Second item
  \item Third item with sub-items:
    \begin{itemize}
      \item Sub-item 1
      \item Sub-item 2
    \end{itemize}
\end{itemize}

\subsection{Numbered List}
\begin{enumerate}
  \item First step
  \item Second step
  \item Third step
\end{enumerate}

\section{Citations and References}
When you use the citation search feature, references will be automatically managed. For example, you might cite a paper like this: \cite{example2024}. All citations will be automatically added to your \texttt{cite.bib} file.

\section{Conclusion}
This template demonstrates basic \LaTeX{} features including mathematical formulas, tables, and lists. The bibliography system is now set up to work seamlessly with Bibby AI's citation management. Feel free to modify it for your needs!

\bibliographystyle{plain}
\bibliography{cite}

\end{document}\documentclass{article}
\usepackage[utf8]{inputenc}
\usepackage{booktabs}
\usepackage{amsmath}
\usepackage{natbib}

\title{My Document}
\author{Your Name}
\date{\today}

\begin{document}

\maketitle

\section{Introduction}
This is a sample document demonstrating various \LaTeX{} features including mathematical equations, tables, and formatting. You can cite papers using the citation search feature, and they will be automatically added to your \texttt{cite.bib} file.

\section{Mathematical Expressions}
Here are some example equations:

\subsection{Quadratic Formula}
The quadratic formula is given by:
\[x = \frac{-b \pm \sqrt{b^2 - 4ac}}{2a}\]

\subsection{Matrix Example}
A sample matrix:
\[
\begin{pmatrix}
a_{11} & a_{12} & a_{13} \\
a_{21} & a_{22} & a_{23} \\
a_{31} & a_{32} & a_{33}
\end{pmatrix}
\]

\section{Tables}
Below is a sample table showing data:

\begin{table}[h]
\centering
\begin{tabular}{@{}lrrr@{}}
\toprule
\textbf{Category} & \textbf{Value 1} & \textbf{Value 2} & \textbf{Total} \\
\midrule
Item A & 10.5 & 20.3 & 30.8 \\
Item B & 15.2 & 25.1 & 40.3 \\
Item C & 22.1 & 12.7 & 34.8 \\
\bottomrule
\end{tabular}
\caption{Sample Data Table}
\label{tab:sample}
\end{table}

\section{Lists}
\subsection{Bullet Points}
\begin{itemize}
  \item First item
  \item Second item
  \item Third item with sub-items:
    \begin{itemize}
      \item Sub-item 1
      \item Sub-item 2
    \end{itemize}
\end{itemize}

\subsection{Numbered List}
\begin{enumerate}
  \item First step
  \item Second step
  \item Third step
\end{enumerate}

\section{Citations and References}
When you use the citation search feature, references will be automatically managed. For example, you might cite a paper like this: \cite{example2024}. All citations will be automatically added to your \texttt{cite.bib} file.

\section{Conclusion}
This template demonstrates basic \LaTeX{} features including mathematical formulas, tables, and lists. The bibliography system is now set up to work seamlessly with Bibby AI's citation management. Feel free to modify it for your needs!

\bibliographystyle{plain}
\bibliography{cite}

\end{document}\documentclass{article}
\usepackage[utf8]{inputenc}
\usepackage{booktabs}
\usepackage{amsmath}
\usepackage{natbib}

\title{My Document}
\author{Your Name}
\date{\today}

\begin{document}

\maketitle

\section{Introduction}
This is a sample document demonstrating various \LaTeX{} features including mathematical equations, tables, and formatting. You can cite papers using the citation search feature, and they will be automatically added to your \texttt{cite.bib} file.

\section{Mathematical Expressions}
Here are some example equations:

\subsection{Quadratic Formula}
The quadratic formula is given by:
\[x = \frac{-b \pm \sqrt{b^2 - 4ac}}{2a}\]

\subsection{Matrix Example}
A sample matrix:
\[
\begin{pmatrix}
a_{11} & a_{12} & a_{13} \\
a_{21} & a_{22} & a_{23} \\
a_{31} & a_{32} & a_{33}
\end{pmatrix}
\]

\section{Tables}
Below is a sample table showing data:

\begin{table}[h]
\centering
\begin{tabular}{@{}lrrr@{}}
\toprule
\textbf{Category} & \textbf{Value 1} & \textbf{Value 2} & \textbf{Total} \\
\midrule
Item A & 10.5 & 20.3 & 30.8 \\
Item B & 15.2 & 25.1 & 40.3 \\
Item C & 22.1 & 12.7 & 34.8 \\
\bottomrule
\end{tabular}
\caption{Sample Data Table}
\label{tab:sample}
\end{table}

\section{Lists}
\subsection{Bullet Points}
\begin{itemize}
  \item First item
  \item Second item
  \item Third item with sub-items:
    \begin{itemize}
      \item Sub-item 1
      \item Sub-item 2
    \end{itemize}
\end{itemize}

\subsection{Numbered List}
\begin{enumerate}
  \item First step
  \item Second step
  \item Third step
\end{enumerate}

\section{Citations and References}
When you use the citation search feature, references will be automatically managed. For example, you might cite a paper like this: \cite{example2024}. All citations will be automatically added to your \texttt{cite.bib} file.

\section{Conclusion}
This template demonstrates basic \LaTeX{} features including mathematical formulas, tables, and lists. The bibliography system is now set up to work seamlessly with Bibby AI's citation management. Feel free to modify it for your needs!

\bibliographystyle{plain}
\bibliography{cite}

\end{document}\documentclass{article}
\usepackage[utf8]{inputenc}
\usepackage{booktabs}
\usepackage{amsmath}
\usepackage{natbib}

\title{My Document}
\author{Your Name}
\date{\today}

\begin{document}

\maketitle

\section{Introduction}
This is a sample document demonstrating various \LaTeX{} features including mathematical equations, tables, and formatting. You can cite papers using the citation search feature, and they will be automatically added to your \texttt{cite.bib} file.

\section{Mathematical Expressions}
Here are some example equations:

\subsection{Quadratic Formula}
The quadratic formula is given by:
\[x = \frac{-b \pm \sqrt{b^2 - 4ac}}{2a}\]

\subsection{Matrix Example}
A sample matrix:
\[
\begin{pmatrix}
a_{11} & a_{12} & a_{13} \\
a_{21} & a_{22} & a_{23} \\
a_{31} & a_{32} & a_{33}
\end{pmatrix}
\]

\section{Tables}
Below is a sample table showing data:

\begin{table}[h]
\centering
\begin{tabular}{@{}lrrr@{}}
\toprule
\textbf{Category} & \textbf{Value 1} & \textbf{Value 2} & \textbf{Total} \\
\midrule
Item A & 10.5 & 20.3 & 30.8 \\
Item B & 15.2 & 25.1 & 40.3 \\
Item C & 22.1 & 12.7 & 34.8 \\
\bottomrule
\end{tabular}
\caption{Sample Data Table}
\label{tab:sample}
\end{table}

\section{Lists}
\subsection{Bullet Points}
\begin{itemize}
  \item First item
  \item Second item
  \item Third item with sub-items:
    \begin{itemize}
      \item Sub-item 1
      \item Sub-item 2
    \end{itemize}
\end{itemize}

\subsection{Numbered List}
\begin{enumerate}
  \item First step
  \item Second step
  \item Third step
\end{enumerate}

\section{Citations and References}
When you use the citation search feature, references will be automatically managed. For example, you might cite a paper like this: \cite{example2024}. All citations will be automatically added to your \texttt{cite.bib} file.

\section{Conclusion}
This template demonstrates basic \LaTeX{} features including mathematical formulas, tables, and lists. The bibliography system is now set up to work seamlessly with Bibby AI's citation management. Feel free to modify it for your needs!

\bibliographystyle{plain}
\bibliography{cite}

\end{document}\documentclass{article}
\usepackage[utf8]{inputenc}
\usepackage{booktabs}
\usepackage{amsmath}
\usepackage{natbib}

\title{My Document}
\author{Your Name}
\date{\today}

\begin{document}

\maketitle

\section{Introduction}
This is a sample document demonstrating various \LaTeX{} features including mathematical equations, tables, and formatting. You can cite papers using the citation search feature, and they will be automatically added to your \texttt{cite.bib} file.

\section{Mathematical Expressions}
Here are some example equations:

\subsection{Quadratic Formula}
The quadratic formula is given by:
\[x = \frac{-b \pm \sqrt{b^2 - 4ac}}{2a}\]

\subsection{Matrix Example}
A sample matrix:
\[
\begin{pmatrix}
a_{11} & a_{12} & a_{13} \\
a_{21} & a_{22} & a_{23} \\
a_{31} & a_{32} & a_{33}
\end{pmatrix}
\]

\section{Tables}
Below is a sample table showing data:

\begin{table}[h]
\centering
\begin{tabular}{@{}lrrr@{}}
\toprule
\textbf{Category} & \textbf{Value 1} & \textbf{Value 2} & \textbf{Total} \\
\midrule
Item A & 10.5 & 20.3 & 30.8 \\
Item B & 15.2 & 25.1 & 40.3 \\
Item C & 22.1 & 12.7 & 34.8 \\
\bottomrule
\end{tabular}
\caption{Sample Data Table}
\label{tab:sample}
\end{table}

\section{Lists}
\subsection{Bullet Points}
\begin{itemize}
  \item First item
  \item Second item
  \item Third item with sub-items:
    \begin{itemize}
      \item Sub-item 1
      \item Sub-item 2
    \end{itemize}
\end{itemize}

\subsection{Numbered List}
\begin{enumerate}
  \item First step
  \item Second step
  \item Third step
\end{enumerate}

\section{Citations and References}
When you use the citation search feature, references will be automatically managed. For example, you might cite a paper like this: \cite{example2024}. All citations will be automatically added to your \texttt{cite.bib} file.

\section{Conclusion}
This template demonstrates basic \LaTeX{} features including mathematical formulas, tables, and lists. The bibliography system is now set up to work seamlessly with Bibby AI's citation management. Feel free to modify it for your needs!

\bibliographystyle{plain}
\bibliography{cite}

\end{document}\documentclass{article}
\usepackage[utf8]{inputenc}
\usepackage{booktabs}
\usepackage{amsmath}
\usepackage{natbib}

\title{My Document}
\author{Your Name}
\date{\today}

\begin{document}

\maketitle

\section{Introduction}
This is a sample document demonstrating various \LaTeX{} features including mathematical equations, tables, and formatting. You can cite papers using the citation search feature, and they will be automatically added to your \texttt{cite.bib} file.

\section{Mathematical Expressions}
Here are some example equations:

\subsection{Quadratic Formula}
The quadratic formula is given by:
\[x = \frac{-b \pm \sqrt{b^2 - 4ac}}{2a}\]

\subsection{Matrix Example}
A sample matrix:
\[
\begin{pmatrix}
a_{11} & a_{12} & a_{13} \\
a_{21} & a_{22} & a_{23} \\
a_{31} & a_{32} & a_{33}
\end{pmatrix}
\]

\section{Tables}
Below is a sample table showing data:

\begin{table}[h]
\centering
\begin{tabular}{@{}lrrr@{}}
\toprule
\textbf{Category} & \textbf{Value 1} & \textbf{Value 2} & \textbf{Total} \\
\midrule
Item A & 10.5 & 20.3 & 30.8 \\
Item B & 15.2 & 25.1 & 40.3 \\
Item C & 22.1 & 12.7 & 34.8 \\
\bottomrule
\end{tabular}
\caption{Sample Data Table}
\label{tab:sample}
\end{table}

\section{Lists}
\subsection{Bullet Points}
\begin{itemize}
  \item First item
  \item Second item
  \item Third item with sub-items:
    \begin{itemize}
      \item Sub-item 1
      \item Sub-item 2
    \end{itemize}
\end{itemize}

\subsection{Numbered List}
\begin{enumerate}
  \item First step
  \item Second step
  \item Third step
\end{enumerate}

\section{Citations and References}
When you use the citation search feature, references will be automatically managed. For example, you might cite a paper like this: \cite{example2024}. All citations will be automatically added to your \texttt{cite.bib} file.

\section{Conclusion}
This template demonstrates basic \LaTeX{} features including mathematical formulas, tables, and lists. The bibliography system is now set up to work seamlessly with Bibby AI's citation management. Feel free to modify it for your needs!

\bibliographystyle{plain}
\bibliography{cite}

\end{document}\documentclass{article}
\usepackage[utf8]{inputenc}
\usepackage{booktabs}
\usepackage{amsmath}
\usepackage{natbib}

\title{My Document}
\author{Your Name}
\date{\today}

\begin{document}

\maketitle

\section{Introduction}
This is a sample document demonstrating various \LaTeX{} features including mathematical equations, tables, and formatting. You can cite papers using the citation search feature, and they will be automatically added to your \texttt{cite.bib} file.

\section{Mathematical Expressions}
Here are some example equations:

\subsection{Quadratic Formula}
The quadratic formula is given by:
\[x = \frac{-b \pm \sqrt{b^2 - 4ac}}{2a}\]

\subsection{Matrix Example}
A sample matrix:
\[
\begin{pmatrix}
a_{11} & a_{12} & a_{13} \\
a_{21} & a_{22} & a_{23} \\
a_{31} & a_{32} & a_{33}
\end{pmatrix}
\]

\section{Tables}
Below is a sample table showing data:

\begin{table}[h]
\centering
\begin{tabular}{@{}lrrr@{}}
\toprule
\textbf{Category} & \textbf{Value 1} & \textbf{Value 2} & \textbf{Total} \\
\midrule
Item A & 10.5 & 20.3 & 30.8 \\
Item B & 15.2 & 25.1 & 40.3 \\
Item C & 22.1 & 12.7 & 34.8 \\
\bottomrule
\end{tabular}
\caption{Sample Data Table}
\label{tab:sample}
\end{table}

\section{Lists}
\subsection{Bullet Points}
\begin{itemize}
  \item First item
  \item Second item
  \item Third item with sub-items:
    \begin{itemize}
      \item Sub-item 1
      \item Sub-item 2
    \end{itemize}
\end{itemize}

\subsection{Numbered List}
\begin{enumerate}
  \item First step
  \item Second step
  \item Third step
\end{enumerate}

\section{Citations and References}
When you use the citation search feature, references will be automatically managed. For example, you might cite a paper like this: \cite{example2024}. All citations will be automatically added to your \texttt{cite.bib} file.

\section{Conclusion}
This template demonstrates basic \LaTeX{} features including mathematical formulas, tables, and lists. The bibliography system is now set up to work seamlessly with Bibby AI's citation management. Feel free to modify it for your needs!

\bibliographystyle{plain}
\bibliography{cite}

\end{document}\documentclass{article}
\usepackage[utf8]{inputenc}
\usepackage{booktabs}
\usepackage{amsmath}
\usepackage{natbib}

\title{My Document}
\author{Your Name}
\date{\today}

\begin{document}

\maketitle

\section{Introduction}
This is a sample document demonstrating various \LaTeX{} features including mathematical equations, tables, and formatting. You can cite papers using the citation search feature, and they will be automatically added to your \texttt{cite.bib} file.

\section{Mathematical Expressions}
Here are some example equations:

\subsection{Quadratic Formula}
The quadratic formula is given by:
\[x = \frac{-b \pm \sqrt{b^2 - 4ac}}{2a}\]

\subsection{Matrix Example}
A sample matrix:
\[
\begin{pmatrix}
a_{11} & a_{12} & a_{13} \\
a_{21} & a_{22} & a_{23} \\
a_{31} & a_{32} & a_{33}
\end{pmatrix}
\]

\section{Tables}
Below is a sample table showing data:

\begin{table}[h]
\centering
\begin{tabular}{@{}lrrr@{}}
\toprule
\textbf{Category} & \textbf{Value 1} & \textbf{Value 2} & \textbf{Total} \\
\midrule
Item A & 10.5 & 20.3 & 30.8 \\
Item B & 15.2 & 25.1 & 40.3 \\
Item C & 22.1 & 12.7 & 34.8 \\
\bottomrule
\end{tabular}
\caption{Sample Data Table}
\label{tab:sample}
\end{table}

\section{Lists}
\subsection{Bullet Points}
\begin{itemize}
  \item First item
  \item Second item
  \item Third item with sub-items:
    \begin{itemize}
      \item Sub-item 1
      \item Sub-item 2
    \end{itemize}
\end{itemize}

\subsection{Numbered List}
\begin{enumerate}
  \item First step
  \item Second step
  \item Third step
\end{enumerate}

\section{Citations and References}
When you use the citation search feature, references will be automatically managed. For example, you might cite a paper like this: \cite{example2024}. All citations will be automatically added to your \texttt{cite.bib} file.

\section{Conclusion}
This template demonstrates basic \LaTeX{} features including mathematical formulas, tables, and lists. The bibliography system is now set up to work seamlessly with Bibby AI's citation management. Feel free to modify it for your needs!

\bibliographystyle{plain}
\bibliography{cite}

\end{document}\documentclass{article}
\usepackage[utf8]{inputenc}
\usepackage{booktabs}
\usepackage{amsmath}
\usepackage{natbib}

\title{My Document}
\author{Your Name}
\date{\today}

\begin{document}

\maketitle

\section{Introduction}
This is a sample document demonstrating various \LaTeX{} features including mathematical equations, tables, and formatting. You can cite papers using the citation search feature, and they will be automatically added to your \texttt{cite.bib} file.

\section{Mathematical Expressions}
Here are some example equations:

\subsection{Quadratic Formula}
The quadratic formula is given by:
\[x = \frac{-b \pm \sqrt{b^2 - 4ac}}{2a}\]

\subsection{Matrix Example}
A sample matrix:
\[
\begin{pmatrix}
a_{11} & a_{12} & a_{13} \\
a_{21} & a_{22} & a_{23} \\
a_{31} & a_{32} & a_{33}
\end{pmatrix}
\]

\section{Tables}
Below is a sample table showing data:

\begin{table}[h]
\centering
\begin{tabular}{@{}lrrr@{}}
\toprule
\textbf{Category} & \textbf{Value 1} & \textbf{Value 2} & \textbf{Total} \\
\midrule
Item A & 10.5 & 20.3 & 30.8 \\
Item B & 15.2 & 25.1 & 40.3 \\
Item C & 22.1 & 12.7 & 34.8 \\
\bottomrule
\end{tabular}
\caption{Sample Data Table}
\label{tab:sample}
\end{table}

\section{Lists}
\subsection{Bullet Points}
\begin{itemize}
  \item First item
  \item Second item
  \item Third item with sub-items:
    \begin{itemize}
      \item Sub-item 1
      \item Sub-item 2
    \end{itemize}
\end{itemize}

\subsection{Numbered List}
\begin{enumerate}
  \item First step
  \item Second step
  \item Third step
\end{enumerate}

\section{Citations and References}
When you use the citation search feature, references will be automatically managed. For example, you might cite a paper like this: \cite{example2024}. All citations will be automatically added to your \texttt{cite.bib} file.

\section{Conclusion}
This template demonstrates basic \LaTeX{} features including mathematical formulas, tables, and lists. The bibliography system is now set up to work seamlessly with Bibby AI's citation management. Feel free to modify it for your needs!

\bibliographystyle{plain}
\bibliography{cite}

\end{document}\documentclass{article}
\usepackage[utf8]{inputenc}
\usepackage{booktabs}
\usepackage{amsmath}
\usepackage{natbib}

\title{My Document}
\author{Your Name}
\date{\today}

\begin{document}

\maketitle

\section{Introduction}
This is a sample document demonstrating various \LaTeX{} features including mathematical equations, tables, and formatting. You can cite papers using the citation search feature, and they will be automatically added to your \texttt{cite.bib} file.

\section{Mathematical Expressions}
Here are some example equations:

\subsection{Quadratic Formula}
The quadratic formula is given by:
\[x = \frac{-b \pm \sqrt{b^2 - 4ac}}{2a}\]

\subsection{Matrix Example}
A sample matrix:
\[
\begin{pmatrix}
a_{11} & a_{12} & a_{13} \\
a_{21} & a_{22} & a_{23} \\
a_{31} & a_{32} & a_{33}
\end{pmatrix}
\]

\section{Tables}
Below is a sample table showing data:

\begin{table}[h]
\centering
\begin{tabular}{@{}lrrr@{}}
\toprule
\textbf{Category} & \textbf{Value 1} & \textbf{Value 2} & \textbf{Total} \\
\midrule
Item A & 10.5 & 20.3 & 30.8 \\
Item B & 15.2 & 25.1 & 40.3 \\
Item C & 22.1 & 12.7 & 34.8 \\
\bottomrule
\end{tabular}
\caption{Sample Data Table}
\label{tab:sample}
\end{table}

\section{Lists}
\subsection{Bullet Points}
\begin{itemize}
  \item First item
  \item Second item
  \item Third item with sub-items:
    \begin{itemize}
      \item Sub-item 1
      \item Sub-item 2
    \end{itemize}
\end{itemize}

\subsection{Numbered List}
\begin{enumerate}
  \item First step
  \item Second step
  \item Third step
\end{enumerate}

\section{Citations and References}
When you use the citation search feature, references will be automatically managed. For example, you might cite a paper like this: \cite{example2024}. All citations will be automatically added to your \texttt{cite.bib} file.

\section{Conclusion}
This template demonstrates basic \LaTeX{} features including mathematical formulas, tables, and lists. The bibliography system is now set up to work seamlessly with Bibby AI's citation management. Feel free to modify it for your needs!

\bibliographystyle{plain}
\bibliography{cite}

\end{document}\documentclass{article}
\usepackage[utf8]{inputenc}
\usepackage{booktabs}
\usepackage{amsmath}
\usepackage{natbib}

\title{My Document}
\author{Your Name}
\date{\today}

\begin{document}

\maketitle

\section{Introduction}
This is a sample document demonstrating various \LaTeX{} features including mathematical equations, tables, and formatting. You can cite papers using the citation search feature, and they will be automatically added to your \texttt{cite.bib} file.

\section{Mathematical Expressions}
Here are some example equations:

\subsection{Quadratic Formula}
The quadratic formula is given by:
\[x = \frac{-b \pm \sqrt{b^2 - 4ac}}{2a}\]

\subsection{Matrix Example}
A sample matrix:
\[
\begin{pmatrix}
a_{11} & a_{12} & a_{13} \\
a_{21} & a_{22} & a_{23} \\
a_{31} & a_{32} & a_{33}
\end{pmatrix}
\]

\section{Tables}
Below is a sample table showing data:

\begin{table}[h]
\centering
\begin{tabular}{@{}lrrr@{}}
\toprule
\textbf{Category} & \textbf{Value 1} & \textbf{Value 2} & \textbf{Total} \\
\midrule
Item A & 10.5 & 20.3 & 30.8 \\
Item B & 15.2 & 25.1 & 40.3 \\
Item C & 22.1 & 12.7 & 34.8 \\
\bottomrule
\end{tabular}
\caption{Sample Data Table}
\label{tab:sample}
\end{table}

\section{Lists}
\subsection{Bullet Points}
\begin{itemize}
  \item First item
  \item Second item
  \item Third item with sub-items:
    \begin{itemize}
      \item Sub-item 1
      \item Sub-item 2
    \end{itemize}
\end{itemize}

\subsection{Numbered List}
\begin{enumerate}
  \item First step
  \item Second step
  \item Third step
\end{enumerate}

\section{Citations and References}
When you use the citation search feature, references will be automatically managed. For example, you might cite a paper like this: \cite{example2024}. All citations will be automatically added to your \texttt{cite.bib} file.

\section{Conclusion}
This template demonstrates basic \LaTeX{} features including mathematical formulas, tables, and lists. The bibliography system is now set up to work seamlessly with Bibby AI's citation management. Feel free to modify it for your needs!

\bibliographystyle{plain}
\bibliography{cite}

\end{document}\documentclass{article}
\usepackage[utf8]{inputenc}
\usepackage{booktabs}
\usepackage{amsmath}
\usepackage{natbib}

\title{My Document}
\author{Your Name}
\date{\today}

\begin{document}

\maketitle

\section{Introduction}
This is a sample document demonstrating various \LaTeX{} features including mathematical equations, tables, and formatting. You can cite papers using the citation search feature, and they will be automatically added to your \texttt{cite.bib} file.

\section{Mathematical Expressions}
Here are some example equations:

\subsection{Quadratic Formula}
The quadratic formula is given by:
\[x = \frac{-b \pm \sqrt{b^2 - 4ac}}{2a}\]

\subsection{Matrix Example}
A sample matrix:
\[
\begin{pmatrix}
a_{11} & a_{12} & a_{13} \\
a_{21} & a_{22} & a_{23} \\
a_{31} & a_{32} & a_{33}
\end{pmatrix}
\]

\section{Tables}
Below is a sample table showing data:

\begin{table}[h]
\centering
\begin{tabular}{@{}lrrr@{}}
\toprule
\textbf{Category} & \textbf{Value 1} & \textbf{Value 2} & \textbf{Total} \\
\midrule
Item A & 10.5 & 20.3 & 30.8 \\
Item B & 15.2 & 25.1 & 40.3 \\
Item C & 22.1 & 12.7 & 34.8 \\
\bottomrule
\end{tabular}
\caption{Sample Data Table}
\label{tab:sample}
\end{table}

\section{Lists}
\subsection{Bullet Points}
\begin{itemize}
  \item First item
  \item Second item
  \item Third item with sub-items:
    \begin{itemize}
      \item Sub-item 1
      \item Sub-item 2
    \end{itemize}
\end{itemize}

\subsection{Numbered List}
\begin{enumerate}
  \item First step
  \item Second step
  \item Third step
\end{enumerate}

\section{Citations and References}
When you use the citation search feature, references will be automatically managed. For example, you might cite a paper like this: \cite{example2024}. All citations will be automatically added to your \texttt{cite.bib} file.

\section{Conclusion}
This template demonstrates basic \LaTeX{} features including mathematical formulas, tables, and lists. The bibliography system is now set up to work seamlessly with Bibby AI's citation management. Feel free to modify it for your needs!

\bibliographystyle{plain}
\bibliography{cite}

\end{document}\documentclass{article}
\usepackage[utf8]{inputenc}
\usepackage{booktabs}
\usepackage{amsmath}
\usepackage{natbib}

\title{My Document}
\author{Your Name}
\date{\today}

\begin{document}

\maketitle

\section{Introduction}
This is a sample document demonstrating various \LaTeX{} features including mathematical equations, tables, and formatting. You can cite papers using the citation search feature, and they will be automatically added to your \texttt{cite.bib} file.

\section{Mathematical Expressions}
Here are some example equations:

\subsection{Quadratic Formula}
The quadratic formula is given by:
\[x = \frac{-b \pm \sqrt{b^2 - 4ac}}{2a}\]

\subsection{Matrix Example}
A sample matrix:
\[
\begin{pmatrix}
a_{11} & a_{12} & a_{13} \\
a_{21} & a_{22} & a_{23} \\
a_{31} & a_{32} & a_{33}
\end{pmatrix}
\]

\section{Tables}
Below is a sample table showing data:

\begin{table}[h]
\centering
\begin{tabular}{@{}lrrr@{}}
\toprule
\textbf{Category} & \textbf{Value 1} & \textbf{Value 2} & \textbf{Total} \\
\midrule
Item A & 10.5 & 20.3 & 30.8 \\
Item B & 15.2 & 25.1 & 40.3 \\
Item C & 22.1 & 12.7 & 34.8 \\
\bottomrule
\end{tabular}
\caption{Sample Data Table}
\label{tab:sample}
\end{table}

\section{Lists}
\subsection{Bullet Points}
\begin{itemize}
  \item First item
  \item Second item
  \item Third item with sub-items:
    \begin{itemize}
      \item Sub-item 1
      \item Sub-item 2
    \end{itemize}
\end{itemize}

\subsection{Numbered List}
\begin{enumerate}
  \item First step
  \item Second step
  \item Third step
\end{enumerate}

\section{Citations and References}
When you use the citation search feature, references will be automatically managed. For example, you might cite a paper like this: \cite{example2024}. All citations will be automatically added to your \texttt{cite.bib} file.

\section{Conclusion}
This template demonstrates basic \LaTeX{} features including mathematical formulas, tables, and lists. The bibliography system is now set up to work seamlessly with Bibby AI's citation management. Feel free to modify it for your needs!

\bibliographystyle{plain}
\bibliography{cite}

\end{document}\documentclass{article}
\usepackage[utf8]{inputenc}
\usepackage{booktabs}
\usepackage{amsmath}
\usepackage{natbib}

\title{My Document}
\author{Your Name}
\date{\today}

\begin{document}

\maketitle

\section{Introduction}
This is a sample document demonstrating various \LaTeX{} features including mathematical equations, tables, and formatting. You can cite papers using the citation search feature, and they will be automatically added to your \texttt{cite.bib} file.

\section{Mathematical Expressions}
Here are some example equations:

\subsection{Quadratic Formula}
The quadratic formula is given by:
\[x = \frac{-b \pm \sqrt{b^2 - 4ac}}{2a}\]

\subsection{Matrix Example}
A sample matrix:
\[
\begin{pmatrix}
a_{11} & a_{12} & a_{13} \\
a_{21} & a_{22} & a_{23} \\
a_{31} & a_{32} & a_{33}
\end{pmatrix}
\]

\section{Tables}
Below is a sample table showing data:

\begin{table}[h]
\centering
\begin{tabular}{@{}lrrr@{}}
\toprule
\textbf{Category} & \textbf{Value 1} & \textbf{Value 2} & \textbf{Total} \\
\midrule
Item A & 10.5 & 20.3 & 30.8 \\
Item B & 15.2 & 25.1 & 40.3 \\
Item C & 22.1 & 12.7 & 34.8 \\
\bottomrule
\end{tabular}
\caption{Sample Data Table}
\label{tab:sample}
\end{table}

\section{Lists}
\subsection{Bullet Points}
\begin{itemize}
  \item First item
  \item Second item
  \item Third item with sub-items:
    \begin{itemize}
      \item Sub-item 1
      \item Sub-item 2
    \end{itemize}
\end{itemize}

\subsection{Numbered List}
\begin{enumerate}
  \item First step
  \item Second step
  \item Third step
\end{enumerate}

\section{Citations and References}
When you use the citation search feature, references will be automatically managed. For example, you might cite a paper like this: \cite{example2024}. All citations will be automatically added to your \texttt{cite.bib} file.

\section{Conclusion}
This template demonstrates basic \LaTeX{} features including mathematical formulas, tables, and lists. The bibliography system is now set up to work seamlessly with Bibby AI's citation management. Feel free to modify it for your needs!

\bibliographystyle{plain}
\bibliography{cite}

\end{document}\documentclass{article}
\usepackage[utf8]{inputenc}
\usepackage{booktabs}
\usepackage{amsmath}
\usepackage{natbib}

\title{My Document}
\author{Your Name}
\date{\today}

\begin{document}

\maketitle

\section{Introduction}
This is a sample document demonstrating various \LaTeX{} features including mathematical equations, tables, and formatting. You can cite papers using the citation search feature, and they will be automatically added to your \texttt{cite.bib} file.

\section{Mathematical Expressions}
Here are some example equations:

\subsection{Quadratic Formula}
The quadratic formula is given by:
\[x = \frac{-b \pm \sqrt{b^2 - 4ac}}{2a}\]

\subsection{Matrix Example}
A sample matrix:
\[
\begin{pmatrix}
a_{11} & a_{12} & a_{13} \\
a_{21} & a_{22} & a_{23} \\
a_{31} & a_{32} & a_{33}
\end{pmatrix}
\]

\section{Tables}
Below is a sample table showing data:

\begin{table}[h]
\centering
\begin{tabular}{@{}lrrr@{}}
\toprule
\textbf{Category} & \textbf{Value 1} & \textbf{Value 2} & \textbf{Total} \\
\midrule
Item A & 10.5 & 20.3 & 30.8 \\
Item B & 15.2 & 25.1 & 40.3 \\
Item C & 22.1 & 12.7 & 34.8 \\
\bottomrule
\end{tabular}
\caption{Sample Data Table}
\label{tab:sample}
\end{table}

\section{Lists}
\subsection{Bullet Points}
\begin{itemize}
  \item First item
  \item Second item
  \item Third item with sub-items:
    \begin{itemize}
      \item Sub-item 1
      \item Sub-item 2
    \end{itemize}
\end{itemize}

\subsection{Numbered List}
\begin{enumerate}
  \item First step
  \item Second step
  \item Third step
\end{enumerate}

\section{Citations and References}
When you use the citation search feature, references will be automatically managed. For example, you might cite a paper like this: \cite{example2024}. All citations will be automatically added to your \texttt{cite.bib} file.

\section{Conclusion}
This template demonstrates basic \LaTeX{} features including mathematical formulas, tables, and lists. The bibliography system is now set up to work seamlessly with Bibby AI's citation management. Feel free to modify it for your needs!

\bibliographystyle{plain}
\bibliography{cite}

\end{document}\documentclass{article}
\usepackage[utf8]{inputenc}
\usepackage{booktabs}
\usepackage{amsmath}
\usepackage{natbib}

\title{My Document}
\author{Your Name}
\date{\today}

\begin{document}

\maketitle

\section{Introduction}
This is a sample document demonstrating various \LaTeX{} features including mathematical equations, tables, and formatting. You can cite papers using the citation search feature, and they will be automatically added to your \texttt{cite.bib} file.

\section{Mathematical Expressions}
Here are some example equations:

\subsection{Quadratic Formula}
The quadratic formula is given by:
\[x = \frac{-b \pm \sqrt{b^2 - 4ac}}{2a}\]

\subsection{Matrix Example}
A sample matrix:
\[
\begin{pmatrix}
a_{11} & a_{12} & a_{13} \\
a_{21} & a_{22} & a_{23} \\
a_{31} & a_{32} & a_{33}
\end{pmatrix}
\]

\section{Tables}
Below is a sample table showing data:

\begin{table}[h]
\centering
\begin{tabular}{@{}lrrr@{}}
\toprule
\textbf{Category} & \textbf{Value 1} & \textbf{Value 2} & \textbf{Total} \\
\midrule
Item A & 10.5 & 20.3 & 30.8 \\
Item B & 15.2 & 25.1 & 40.3 \\
Item C & 22.1 & 12.7 & 34.8 \\
\bottomrule
\end{tabular}
\caption{Sample Data Table}
\label{tab:sample}
\end{table}

\section{Lists}
\subsection{Bullet Points}
\begin{itemize}
  \item First item
  \item Second item
  \item Third item with sub-items:
    \begin{itemize}
      \item Sub-item 1
      \item Sub-item 2
    \end{itemize}
\end{itemize}

\subsection{Numbered List}
\begin{enumerate}
  \item First step
  \item Second step
  \item Third step
\end{enumerate}

\section{Citations and References}
When you use the citation search feature, references will be automatically managed. For example, you might cite a paper like this: \cite{example2024}. All citations will be automatically added to your \texttt{cite.bib} file.

\section{Conclusion}
This template demonstrates basic \LaTeX{} features including mathematical formulas, tables, and lists. The bibliography system is now set up to work seamlessly with Bibby AI's citation management. Feel free to modify it for your needs!

\bibliographystyle{plain}
\bibliography{cite}

\end{document}\documentclass{article}
\usepackage[utf8]{inputenc}
\usepackage{booktabs}
\usepackage{amsmath}
\usepackage{natbib}

\title{My Document}
\author{Your Name}
\date{\today}

\begin{document}

\maketitle

\section{Introduction}
This is a sample document demonstrating various \LaTeX{} features including mathematical equations, tables, and formatting. You can cite papers using the citation search feature, and they will be automatically added to your \texttt{cite.bib} file.

\section{Mathematical Expressions}
Here are some example equations:

\subsection{Quadratic Formula}
The quadratic formula is given by:
\[x = \frac{-b \pm \sqrt{b^2 - 4ac}}{2a}\]

\subsection{Matrix Example}
A sample matrix:
\[
\begin{pmatrix}
a_{11} & a_{12} & a_{13} \\
a_{21} & a_{22} & a_{23} \\
a_{31} & a_{32} & a_{33}
\end{pmatrix}
\]

\section{Tables}
Below is a sample table showing data:

\begin{table}[h]
\centering
\begin{tabular}{@{}lrrr@{}}
\toprule
\textbf{Category} & \textbf{Value 1} & \textbf{Value 2} & \textbf{Total} \\
\midrule
Item A & 10.5 & 20.3 & 30.8 \\
Item B & 15.2 & 25.1 & 40.3 \\
Item C & 22.1 & 12.7 & 34.8 \\
\bottomrule
\end{tabular}
\caption{Sample Data Table}
\label{tab:sample}
\end{table}

\section{Lists}
\subsection{Bullet Points}
\begin{itemize}
  \item First item
  \item Second item
  \item Third item with sub-items:
    \begin{itemize}
      \item Sub-item 1
      \item Sub-item 2
    \end{itemize}
\end{itemize}

\subsection{Numbered List}
\begin{enumerate}
  \item First step
  \item Second step
  \item Third step
\end{enumerate}

\section{Citations and References}
When you use the citation search feature, references will be automatically managed. For example, you might cite a paper like this: \cite{example2024}. All citations will be automatically added to your \texttt{cite.bib} file.

\section{Conclusion}
This template demonstrates basic \LaTeX{} features including mathematical formulas, tables, and lists. The bibliography system is now set up to work seamlessly with Bibby AI's citation management. Feel free to modify it for your needs!

\bibliographystyle{plain}
\bibliography{cite}

\end{document}\documentclass{article}
\usepackage[utf8]{inputenc}
\usepackage{booktabs}
\usepackage{amsmath}
\usepackage{natbib}

\title{My Document}
\author{Your Name}
\date{\today}

\begin{document}

\maketitle

\section{Introduction}
This is a sample document demonstrating various \LaTeX{} features including mathematical equations, tables, and formatting. You can cite papers using the citation search feature, and they will be automatically added to your \texttt{cite.bib} file.

\section{Mathematical Expressions}
Here are some example equations:

\subsection{Quadratic Formula}
The quadratic formula is given by:
\[x = \frac{-b \pm \sqrt{b^2 - 4ac}}{2a}\]

\subsection{Matrix Example}
A sample matrix:
\[
\begin{pmatrix}
a_{11} & a_{12} & a_{13} \\
a_{21} & a_{22} & a_{23} \\
a_{31} & a_{32} & a_{33}
\end{pmatrix}
\]

\section{Tables}
Below is a sample table showing data:

\begin{table}[h]
\centering
\begin{tabular}{@{}lrrr@{}}
\toprule
\textbf{Category} & \textbf{Value 1} & \textbf{Value 2} & \textbf{Total} \\
\midrule
Item A & 10.5 & 20.3 & 30.8 \\
Item B & 15.2 & 25.1 & 40.3 \\
Item C & 22.1 & 12.7 & 34.8 \\
\bottomrule
\end{tabular}
\caption{Sample Data Table}
\label{tab:sample}
\end{table}

\section{Lists}
\subsection{Bullet Points}
\begin{itemize}
  \item First item
  \item Second item
  \item Third item with sub-items:
    \begin{itemize}
      \item Sub-item 1
      \item Sub-item 2
    \end{itemize}
\end{itemize}

\subsection{Numbered List}
\begin{enumerate}
  \item First step
  \item Second step
  \item Third step
\end{enumerate}

\section{Citations and References}
When you use the citation search feature, references will be automatically managed. For example, you might cite a paper like this: \cite{example2024}. All citations will be automatically added to your \texttt{cite.bib} file.

\section{Conclusion}
This template demonstrates basic \LaTeX{} features including mathematical formulas, tables, and lists. The bibliography system is now set up to work seamlessly with Bibby AI's citation management. Feel free to modify it for your needs!

\bibliographystyle{plain}
\bibliography{cite}

\end{document}\documentclass{article}
\usepackage[utf8]{inputenc}
\usepackage{booktabs}
\usepackage{amsmath}
\usepackage{natbib}

\title{My Document}
\author{Your Name}
\date{\today}

\begin{document}

\maketitle

\section{Introduction}
This is a sample document demonstrating various \LaTeX{} features including mathematical equations, tables, and formatting. You can cite papers using the citation search feature, and they will be automatically added to your \texttt{cite.bib} file.

\section{Mathematical Expressions}
Here are some example equations:

\subsection{Quadratic Formula}
The quadratic formula is given by:
\[x = \frac{-b \pm \sqrt{b^2 - 4ac}}{2a}\]

\subsection{Matrix Example}
A sample matrix:
\[
\begin{pmatrix}
a_{11} & a_{12} & a_{13} \\
a_{21} & a_{22} & a_{23} \\
a_{31} & a_{32} & a_{33}
\end{pmatrix}
\]

\section{Tables}
Below is a sample table showing data:

\begin{table}[h]
\centering
\begin{tabular}{@{}lrrr@{}}
\toprule
\textbf{Category} & \textbf{Value 1} & \textbf{Value 2} & \textbf{Total} \\
\midrule
Item A & 10.5 & 20.3 & 30.8 \\
Item B & 15.2 & 25.1 & 40.3 \\
Item C & 22.1 & 12.7 & 34.8 \\
\bottomrule
\end{tabular}
\caption{Sample Data Table}
\label{tab:sample}
\end{table}

\section{Lists}
\subsection{Bullet Points}
\begin{itemize}
  \item First item
  \item Second item
  \item Third item with sub-items:
    \begin{itemize}
      \item Sub-item 1
      \item Sub-item 2
    \end{itemize}
\end{itemize}

\subsection{Numbered List}
\begin{enumerate}
  \item First step
  \item Second step
  \item Third step
\end{enumerate}

\section{Citations and References}
When you use the citation search feature, references will be automatically managed. For example, you might cite a paper like this: \cite{example2024}. All citations will be automatically added to your \texttt{cite.bib} file.

\section{Conclusion}
This template demonstrates basic \LaTeX{} features including mathematical formulas, tables, and lists. The bibliography system is now set up to work seamlessly with Bibby AI's citation management. Feel free to modify it for your needs!

\bibliographystyle{plain}
\bibliography{cite}

\end{document}\documentclass{article}
\usepackage[utf8]{inputenc}
\usepackage{booktabs}
\usepackage{amsmath}
\usepackage{natbib}

\title{My Document}
\author{Your Name}
\date{\today}

\begin{document}

\maketitle

\section{Introduction}
This is a sample document demonstrating various \LaTeX{} features including mathematical equations, tables, and formatting. You can cite papers using the citation search feature, and they will be automatically added to your \texttt{cite.bib} file.

\section{Mathematical Expressions}
Here are some example equations:

\subsection{Quadratic Formula}
The quadratic formula is given by:
\[x = \frac{-b \pm \sqrt{b^2 - 4ac}}{2a}\]

\subsection{Matrix Example}
A sample matrix:
\[
\begin{pmatrix}
a_{11} & a_{12} & a_{13} \\
a_{21} & a_{22} & a_{23} \\
a_{31} & a_{32} & a_{33}
\end{pmatrix}
\]

\section{Tables}
Below is a sample table showing data:

\begin{table}[h]
\centering
\begin{tabular}{@{}lrrr@{}}
\toprule
\textbf{Category} & \textbf{Value 1} & \textbf{Value 2} & \textbf{Total} \\
\midrule
Item A & 10.5 & 20.3 & 30.8 \\
Item B & 15.2 & 25.1 & 40.3 \\
Item C & 22.1 & 12.7 & 34.8 \\
\bottomrule
\end{tabular}
\caption{Sample Data Table}
\label{tab:sample}
\end{table}

\section{Lists}
\subsection{Bullet Points}
\begin{itemize}
  \item First item
  \item Second item
  \item Third item with sub-items:
    \begin{itemize}
      \item Sub-item 1
      \item Sub-item 2
    \end{itemize}
\end{itemize}

\subsection{Numbered List}
\begin{enumerate}
  \item First step
  \item Second step
  \item Third step
\end{enumerate}

\section{Citations and References}
When you use the citation search feature, references will be automatically managed. For example, you might cite a paper like this: \cite{example2024}. All citations will be automatically added to your \texttt{cite.bib} file.

\section{Conclusion}
This template demonstrates basic \LaTeX{} features including mathematical formulas, tables, and lists. The bibliography system is now set up to work seamlessly with Bibby AI's citation management. Feel free to modify it for your needs!

\bibliographystyle{plain}
\bibliography{cite}

\end{document}\documentclass{article}
\usepackage[utf8]{inputenc}
\usepackage{booktabs}
\usepackage{amsmath}
\usepackage{natbib}

\title{My Document}
\author{Your Name}
\date{\today}

\begin{document}

\maketitle

\section{Introduction}
This is a sample document demonstrating various \LaTeX{} features including mathematical equations, tables, and formatting. You can cite papers using the citation search feature, and they will be automatically added to your \texttt{cite.bib} file.

\section{Mathematical Expressions}
Here are some example equations:

\subsection{Quadratic Formula}
The quadratic formula is given by:
\[x = \frac{-b \pm \sqrt{b^2 - 4ac}}{2a}\]

\subsection{Matrix Example}
A sample matrix:
\[
\begin{pmatrix}
a_{11} & a_{12} & a_{13} \\
a_{21} & a_{22} & a_{23} \\
a_{31} & a_{32} & a_{33}
\end{pmatrix}
\]

\section{Tables}
Below is a sample table showing data:

\begin{table}[h]
\centering
\begin{tabular}{@{}lrrr@{}}
\toprule
\textbf{Category} & \textbf{Value 1} & \textbf{Value 2} & \textbf{Total} \\
\midrule
Item A & 10.5 & 20.3 & 30.8 \\
Item B & 15.2 & 25.1 & 40.3 \\
Item C & 22.1 & 12.7 & 34.8 \\
\bottomrule
\end{tabular}
\caption{Sample Data Table}
\label{tab:sample}
\end{table}

\section{Lists}
\subsection{Bullet Points}
\begin{itemize}
  \item First item
  \item Second item
  \item Third item with sub-items:
    \begin{itemize}
      \item Sub-item 1
      \item Sub-item 2
    \end{itemize}
\end{itemize}

\subsection{Numbered List}
\begin{enumerate}
  \item First step
  \item Second step
  \item Third step
\end{enumerate}

\section{Citations and References}
When you use the citation search feature, references will be automatically managed. For example, you might cite a paper like this: \cite{example2024}. All citations will be automatically added to your \texttt{cite.bib} file.

\section{Conclusion}
This template demonstrates basic \LaTeX{} features including mathematical formulas, tables, and lists. The bibliography system is now set up to work seamlessly with Bibby AI's citation management. Feel free to modify it for your needs!

\bibliographystyle{plain}
\bibliography{cite}

\end{document}\documentclass{article}
\usepackage[utf8]{inputenc}
\usepackage{booktabs}
\usepackage{amsmath}
\usepackage{natbib}

\title{My Document}
\author{Your Name}
\date{\today}

\begin{document}

\maketitle

\section{Introduction}
This is a sample document demonstrating various \LaTeX{} features including mathematical equations, tables, and formatting. You can cite papers using the citation search feature, and they will be automatically added to your \texttt{cite.bib} file.

\section{Mathematical Expressions}
Here are some example equations:

\subsection{Quadratic Formula}
The quadratic formula is given by:
\[x = \frac{-b \pm \sqrt{b^2 - 4ac}}{2a}\]

\subsection{Matrix Example}
A sample matrix:
\[
\begin{pmatrix}
a_{11} & a_{12} & a_{13} \\
a_{21} & a_{22} & a_{23} \\
a_{31} & a_{32} & a_{33}
\end{pmatrix}
\]

\section{Tables}
Below is a sample table showing data:

\begin{table}[h]
\centering
\begin{tabular}{@{}lrrr@{}}
\toprule
\textbf{Category} & \textbf{Value 1} & \textbf{Value 2} & \textbf{Total} \\
\midrule
Item A & 10.5 & 20.3 & 30.8 \\
Item B & 15.2 & 25.1 & 40.3 \\
Item C & 22.1 & 12.7 & 34.8 \\
\bottomrule
\end{tabular}
\caption{Sample Data Table}
\label{tab:sample}
\end{table}

\section{Lists}
\subsection{Bullet Points}
\begin{itemize}
  \item First item
  \item Second item
  \item Third item with sub-items:
    \begin{itemize}
      \item Sub-item 1
      \item Sub-item 2
    \end{itemize}
\end{itemize}

\subsection{Numbered List}
\begin{enumerate}
  \item First step
  \item Second step
  \item Third step
\end{enumerate}

\section{Citations and References}
When you use the citation search feature, references will be automatically managed. For example, you might cite a paper like this: \cite{example2024}. All citations will be automatically added to your \texttt{cite.bib} file.

\section{Conclusion}
This template demonstrates basic \LaTeX{} features including mathematical formulas, tables, and lists. The bibliography system is now set up to work seamlessly with Bibby AI's citation management. Feel free to modify it for your needs!

\bibliographystyle{plain}
\bibliography{cite}

\end{document}\documentclass{article}
\usepackage[utf8]{inputenc}
\usepackage{booktabs}
\usepackage{amsmath}
\usepackage{natbib}

\title{My Document}
\author{Your Name}
\date{\today}

\begin{document}

\maketitle

\section{Introduction}
This is a sample document demonstrating various \LaTeX{} features including mathematical equations, tables, and formatting. You can cite papers using the citation search feature, and they will be automatically added to your \texttt{cite.bib} file.

\section{Mathematical Expressions}
Here are some example equations:

\subsection{Quadratic Formula}
The quadratic formula is given by:
\[x = \frac{-b \pm \sqrt{b^2 - 4ac}}{2a}\]

\subsection{Matrix Example}
A sample matrix:
\[
\begin{pmatrix}
a_{11} & a_{12} & a_{13} \\
a_{21} & a_{22} & a_{23} \\
a_{31} & a_{32} & a_{33}
\end{pmatrix}
\]

\section{Tables}
Below is a sample table showing data:

\begin{table}[h]
\centering
\begin{tabular}{@{}lrrr@{}}
\toprule
\textbf{Category} & \textbf{Value 1} & \textbf{Value 2} & \textbf{Total} \\
\midrule
Item A & 10.5 & 20.3 & 30.8 \\
Item B & 15.2 & 25.1 & 40.3 \\
Item C & 22.1 & 12.7 & 34.8 \\
\bottomrule
\end{tabular}
\caption{Sample Data Table}
\label{tab:sample}
\end{table}

\section{Lists}
\subsection{Bullet Points}
\begin{itemize}
  \item First item
  \item Second item
  \item Third item with sub-items:
    \begin{itemize}
      \item Sub-item 1
      \item Sub-item 2
    \end{itemize}
\end{itemize}

\subsection{Numbered List}
\begin{enumerate}
  \item First step
  \item Second step
  \item Third step
\end{enumerate}

\section{Citations and References}
When you use the citation search feature, references will be automatically managed. For example, you might cite a paper like this: \cite{example2024}. All citations will be automatically added to your \texttt{cite.bib} file.

\section{Conclusion}
This template demonstrates basic \LaTeX{} features including mathematical formulas, tables, and lists. The bibliography system is now set up to work seamlessly with Bibby AI's citation management. Feel free to modify it for your needs!

\bibliographystyle{plain}
\bibliography{cite}

\end{document}\documentclass{article}
\usepackage[utf8]{inputenc}
\usepackage{booktabs}
\usepackage{amsmath}
\usepackage{natbib}

\title{My Document}
\author{Your Name}
\date{\today}

\begin{document}

\maketitle

\section{Introduction}
This is a sample document demonstrating various \LaTeX{} features including mathematical equations, tables, and formatting. You can cite papers using the citation search feature, and they will be automatically added to your \texttt{cite.bib} file.

\section{Mathematical Expressions}
Here are some example equations:

\subsection{Quadratic Formula}
The quadratic formula is given by:
\[x = \frac{-b \pm \sqrt{b^2 - 4ac}}{2a}\]

\subsection{Matrix Example}
A sample matrix:
\[
\begin{pmatrix}
a_{11} & a_{12} & a_{13} \\
a_{21} & a_{22} & a_{23} \\
a_{31} & a_{32} & a_{33}
\end{pmatrix}
\]

\section{Tables}
Below is a sample table showing data:

\begin{table}[h]
\centering
\begin{tabular}{@{}lrrr@{}}
\toprule
\textbf{Category} & \textbf{Value 1} & \textbf{Value 2} & \textbf{Total} \\
\midrule
Item A & 10.5 & 20.3 & 30.8 \\
Item B & 15.2 & 25.1 & 40.3 \\
Item C & 22.1 & 12.7 & 34.8 \\
\bottomrule
\end{tabular}
\caption{Sample Data Table}
\label{tab:sample}
\end{table}

\section{Lists}
\subsection{Bullet Points}
\begin{itemize}
  \item First item
  \item Second item
  \item Third item with sub-items:
    \begin{itemize}
      \item Sub-item 1
      \item Sub-item 2
    \end{itemize}
\end{itemize}

\subsection{Numbered List}
\begin{enumerate}
  \item First step
  \item Second step
  \item Third step
\end{enumerate}

\section{Citations and References}
When you use the citation search feature, references will be automatically managed. For example, you might cite a paper like this: \cite{example2024}. All citations will be automatically added to your \texttt{cite.bib} file.

\section{Conclusion}
This template demonstrates basic \LaTeX{} features including mathematical formulas, tables, and lists. The bibliography system is now set up to work seamlessly with Bibby AI's citation management. Feel free to modify it for your needs!

\bibliographystyle{plain}
\bibliography{cite}

\end{document}\documentclass{article}
\usepackage[utf8]{inputenc}
\usepackage{booktabs}
\usepackage{amsmath}
\usepackage{natbib}

\title{My Document}
\author{Your Name}
\date{\today}

\begin{document}

\maketitle

\section{Introduction}
This is a sample document demonstrating various \LaTeX{} features including mathematical equations, tables, and formatting. You can cite papers using the citation search feature, and they will be automatically added to your \texttt{cite.bib} file.

\section{Mathematical Expressions}
Here are some example equations:

\subsection{Quadratic Formula}
The quadratic formula is given by:
\[x = \frac{-b \pm \sqrt{b^2 - 4ac}}{2a}\]

\subsection{Matrix Example}
A sample matrix:
\[
\begin{pmatrix}
a_{11} & a_{12} & a_{13} \\
a_{21} & a_{22} & a_{23} \\
a_{31} & a_{32} & a_{33}
\end{pmatrix}
\]

\section{Tables}
Below is a sample table showing data:

\begin{table}[h]
\centering
\begin{tabular}{@{}lrrr@{}}
\toprule
\textbf{Category} & \textbf{Value 1} & \textbf{Value 2} & \textbf{Total} \\
\midrule
Item A & 10.5 & 20.3 & 30.8 \\
Item B & 15.2 & 25.1 & 40.3 \\
Item C & 22.1 & 12.7 & 34.8 \\
\bottomrule
\end{tabular}
\caption{Sample Data Table}
\label{tab:sample}
\end{table}

\section{Lists}
\subsection{Bullet Points}
\begin{itemize}
  \item First item
  \item Second item
  \item Third item with sub-items:
    \begin{itemize}
      \item Sub-item 1
      \item Sub-item 2
    \end{itemize}
\end{itemize}

\subsection{Numbered List}
\begin{enumerate}
  \item First step
  \item Second step
  \item Third step
\end{enumerate}

\section{Citations and References}
When you use the citation search feature, references will be automatically managed. For example, you might cite a paper like this: \cite{example2024}. All citations will be automatically added to your \texttt{cite.bib} file.

\section{Conclusion}
This template demonstrates basic \LaTeX{} features including mathematical formulas, tables, and lists. The bibliography system is now set up to work seamlessly with Bibby AI's citation management. Feel free to modify it for your needs!

\bibliographystyle{plain}
\bibliography{cite}

\end{document}\documentclass{article}
\usepackage[utf8]{inputenc}
\usepackage{booktabs}
\usepackage{amsmath}
\usepackage{natbib}

\title{My Document}
\author{Your Name}
\date{\today}

\begin{document}

\maketitle

\section{Introduction}
This is a sample document demonstrating various \LaTeX{} features including mathematical equations, tables, and formatting. You can cite papers using the citation search feature, and they will be automatically added to your \texttt{cite.bib} file.

\section{Mathematical Expressions}
Here are some example equations:

\subsection{Quadratic Formula}
The quadratic formula is given by:
\[x = \frac{-b \pm \sqrt{b^2 - 4ac}}{2a}\]

\subsection{Matrix Example}
A sample matrix:
\[
\begin{pmatrix}
a_{11} & a_{12} & a_{13} \\
a_{21} & a_{22} & a_{23} \\
a_{31} & a_{32} & a_{33}
\end{pmatrix}
\]

\section{Tables}
Below is a sample table showing data:

\begin{table}[h]
\centering
\begin{tabular}{@{}lrrr@{}}
\toprule
\textbf{Category} & \textbf{Value 1} & \textbf{Value 2} & \textbf{Total} \\
\midrule
Item A & 10.5 & 20.3 & 30.8 \\
Item B & 15.2 & 25.1 & 40.3 \\
Item C & 22.1 & 12.7 & 34.8 \\
\bottomrule
\end{tabular}
\caption{Sample Data Table}
\label{tab:sample}
\end{table}

\section{Lists}
\subsection{Bullet Points}
\begin{itemize}
  \item First item
  \item Second item
  \item Third item with sub-items:
    \begin{itemize}
      \item Sub-item 1
      \item Sub-item 2
    \end{itemize}
\end{itemize}

\subsection{Numbered List}
\begin{enumerate}
  \item First step
  \item Second step
  \item Third step
\end{enumerate}

\section{Citations and References}
When you use the citation search feature, references will be automatically managed. For example, you might cite a paper like this: \cite{example2024}. All citations will be automatically added to your \texttt{cite.bib} file.

\section{Conclusion}
This template demonstrates basic \LaTeX{} features including mathematical formulas, tables, and lists. The bibliography system is now set up to work seamlessly with Bibby AI's citation management. Feel free to modify it for your needs!

\bibliographystyle{plain}
\bibliography{cite}

\end{document}\documentclass{article}
\usepackage[utf8]{inputenc}
\usepackage{booktabs}
\usepackage{amsmath}
\usepackage{natbib}

\title{My Document}
\author{Your Name}
\date{\today}

\begin{document}

\maketitle

\section{Introduction}
This is a sample document demonstrating various \LaTeX{} features including mathematical equations, tables, and formatting. You can cite papers using the citation search feature, and they will be automatically added to your \texttt{cite.bib} file.

\section{Mathematical Expressions}
Here are some example equations:

\subsection{Quadratic Formula}
The quadratic formula is given by:
\[x = \frac{-b \pm \sqrt{b^2 - 4ac}}{2a}\]

\subsection{Matrix Example}
A sample matrix:
\[
\begin{pmatrix}
a_{11} & a_{12} & a_{13} \\
a_{21} & a_{22} & a_{23} \\
a_{31} & a_{32} & a_{33}
\end{pmatrix}
\]

\section{Tables}
Below is a sample table showing data:

\begin{table}[h]
\centering
\begin{tabular}{@{}lrrr@{}}
\toprule
\textbf{Category} & \textbf{Value 1} & \textbf{Value 2} & \textbf{Total} \\
\midrule
Item A & 10.5 & 20.3 & 30.8 \\
Item B & 15.2 & 25.1 & 40.3 \\
Item C & 22.1 & 12.7 & 34.8 \\
\bottomrule
\end{tabular}
\caption{Sample Data Table}
\label{tab:sample}
\end{table}

\section{Lists}
\subsection{Bullet Points}
\begin{itemize}
  \item First item
  \item Second item
  \item Third item with sub-items:
    \begin{itemize}
      \item Sub-item 1
      \item Sub-item 2
    \end{itemize}
\end{itemize}

\subsection{Numbered List}
\begin{enumerate}
  \item First step
  \item Second step
  \item Third step
\end{enumerate}

\section{Citations and References}
When you use the citation search feature, references will be automatically managed. For example, you might cite a paper like this: \cite{example2024}. All citations will be automatically added to your \texttt{cite.bib} file.

\section{Conclusion}
This template demonstrates basic \LaTeX{} features including mathematical formulas, tables, and lists. The bibliography system is now set up to work seamlessly with Bibby AI's citation management. Feel free to modify it for your needs!

\bibliographystyle{plain}
\bibliography{cite}

\end{document}\documentclass{article}
\usepackage[utf8]{inputenc}
\usepackage{booktabs}
\usepackage{amsmath}
\usepackage{natbib}

\title{My Document}
\author{Your Name}
\date{\today}

\begin{document}

\maketitle

\section{Introduction}
This is a sample document demonstrating various \LaTeX{} features including mathematical equations, tables, and formatting. You can cite papers using the citation search feature, and they will be automatically added to your \texttt{cite.bib} file.

\section{Mathematical Expressions}
Here are some example equations:

\subsection{Quadratic Formula}
The quadratic formula is given by:
\[x = \frac{-b \pm \sqrt{b^2 - 4ac}}{2a}\]

\subsection{Matrix Example}
A sample matrix:
\[
\begin{pmatrix}
a_{11} & a_{12} & a_{13} \\
a_{21} & a_{22} & a_{23} \\
a_{31} & a_{32} & a_{33}
\end{pmatrix}
\]

\section{Tables}
Below is a sample table showing data:

\begin{table}[h]
\centering
\begin{tabular}{@{}lrrr@{}}
\toprule
\textbf{Category} & \textbf{Value 1} & \textbf{Value 2} & \textbf{Total} \\
\midrule
Item A & 10.5 & 20.3 & 30.8 \\
Item B & 15.2 & 25.1 & 40.3 \\
Item C & 22.1 & 12.7 & 34.8 \\
\bottomrule
\end{tabular}
\caption{Sample Data Table}
\label{tab:sample}
\end{table}

\section{Lists}
\subsection{Bullet Points}
\begin{itemize}
  \item First item
  \item Second item
  \item Third item with sub-items:
    \begin{itemize}
      \item Sub-item 1
      \item Sub-item 2
    \end{itemize}
\end{itemize}

\subsection{Numbered List}
\begin{enumerate}
  \item First step
  \item Second step
  \item Third step
\end{enumerate}

\section{Citations and References}
When you use the citation search feature, references will be automatically managed. For example, you might cite a paper like this: \cite{example2024}. All citations will be automatically added to your \texttt{cite.bib} file.

\section{Conclusion}
This template demonstrates basic \LaTeX{} features including mathematical formulas, tables, and lists. The bibliography system is now set up to work seamlessly with Bibby AI's citation management. Feel free to modify it for your needs!

\bibliographystyle{plain}
\bibliography{cite}

\end{document}\documentclass{article}
\usepackage[utf8]{inputenc}
\usepackage{booktabs}
\usepackage{amsmath}
\usepackage{natbib}

\title{My Document}
\author{Your Name}
\date{\today}

\begin{document}

\maketitle

\section{Introduction}
This is a sample document demonstrating various \LaTeX{} features including mathematical equations, tables, and formatting. You can cite papers using the citation search feature, and they will be automatically added to your \texttt{cite.bib} file.

\section{Mathematical Expressions}
Here are some example equations:

\subsection{Quadratic Formula}
The quadratic formula is given by:
\[x = \frac{-b \pm \sqrt{b^2 - 4ac}}{2a}\]

\subsection{Matrix Example}
A sample matrix:
\[
\begin{pmatrix}
a_{11} & a_{12} & a_{13} \\
a_{21} & a_{22} & a_{23} \\
a_{31} & a_{32} & a_{33}
\end{pmatrix}
\]

\section{Tables}
Below is a sample table showing data:

\begin{table}[h]
\centering
\begin{tabular}{@{}lrrr@{}}
\toprule
\textbf{Category} & \textbf{Value 1} & \textbf{Value 2} & \textbf{Total} \\
\midrule
Item A & 10.5 & 20.3 & 30.8 \\
Item B & 15.2 & 25.1 & 40.3 \\
Item C & 22.1 & 12.7 & 34.8 \\
\bottomrule
\end{tabular}
\caption{Sample Data Table}
\label{tab:sample}
\end{table}

\section{Lists}
\subsection{Bullet Points}
\begin{itemize}
  \item First item
  \item Second item
  \item Third item with sub-items:
    \begin{itemize}
      \item Sub-item 1
      \item Sub-item 2
    \end{itemize}
\end{itemize}

\subsection{Numbered List}
\begin{enumerate}
  \item First step
  \item Second step
  \item Third step
\end{enumerate}

\section{Citations and References}
When you use the citation search feature, references will be automatically managed. For example, you might cite a paper like this: \cite{example2024}. All citations will be automatically added to your \texttt{cite.bib} file.

\section{Conclusion}
This template demonstrates basic \LaTeX{} features including mathematical formulas, tables, and lists. The bibliography system is now set up to work seamlessly with Bibby AI's citation management. Feel free to modify it for your needs!

\bibliographystyle{plain}
\bibliography{cite}

\end{document}\documentclass{article}
\usepackage[utf8]{inputenc}
\usepackage{booktabs}
\usepackage{amsmath}
\usepackage{natbib}

\title{My Document}
\author{Your Name}
\date{\today}

\begin{document}

\maketitle

\section{Introduction}
This is a sample document demonstrating various \LaTeX{} features including mathematical equations, tables, and formatting. You can cite papers using the citation search feature, and they will be automatically added to your \texttt{cite.bib} file.

\section{Mathematical Expressions}
Here are some example equations:

\subsection{Quadratic Formula}
The quadratic formula is given by:
\[x = \frac{-b \pm \sqrt{b^2 - 4ac}}{2a}\]

\subsection{Matrix Example}
A sample matrix:
\[
\begin{pmatrix}
a_{11} & a_{12} & a_{13} \\
a_{21} & a_{22} & a_{23} \\
a_{31} & a_{32} & a_{33}
\end{pmatrix}
\]

\section{Tables}
Below is a sample table showing data:

\begin{table}[h]
\centering
\begin{tabular}{@{}lrrr@{}}
\toprule
\textbf{Category} & \textbf{Value 1} & \textbf{Value 2} & \textbf{Total} \\
\midrule
Item A & 10.5 & 20.3 & 30.8 \\
Item B & 15.2 & 25.1 & 40.3 \\
Item C & 22.1 & 12.7 & 34.8 \\
\bottomrule
\end{tabular}
\caption{Sample Data Table}
\label{tab:sample}
\end{table}

\section{Lists}
\subsection{Bullet Points}
\begin{itemize}
  \item First item
  \item Second item
  \item Third item with sub-items:
    \begin{itemize}
      \item Sub-item 1
      \item Sub-item 2
    \end{itemize}
\end{itemize}

\subsection{Numbered List}
\begin{enumerate}
  \item First step
  \item Second step
  \item Third step
\end{enumerate}

\section{Citations and References}
When you use the citation search feature, references will be automatically managed. For example, you might cite a paper like this: \cite{example2024}. All citations will be automatically added to your \texttt{cite.bib} file.

\section{Conclusion}
This template demonstrates basic \LaTeX{} features including mathematical formulas, tables, and lists. The bibliography system is now set up to work seamlessly with Bibby AI's citation management. Feel free to modify it for your needs!

\bibliographystyle{plain}
\bibliography{cite}

\end{document}\documentclass{article}
\usepackage[utf8]{inputenc}
\usepackage{booktabs}
\usepackage{amsmath}
\usepackage{natbib}

\title{My Document}
\author{Your Name}
\date{\today}

\begin{document}

\maketitle

\section{Introduction}
This is a sample document demonstrating various \LaTeX{} features including mathematical equations, tables, and formatting. You can cite papers using the citation search feature, and they will be automatically added to your \texttt{cite.bib} file.

\section{Mathematical Expressions}
Here are some example equations:

\subsection{Quadratic Formula}
The quadratic formula is given by:
\[x = \frac{-b \pm \sqrt{b^2 - 4ac}}{2a}\]

\subsection{Matrix Example}
A sample matrix:
\[
\begin{pmatrix}
a_{11} & a_{12} & a_{13} \\
a_{21} & a_{22} & a_{23} \\
a_{31} & a_{32} & a_{33}
\end{pmatrix}
\]

\section{Tables}
Below is a sample table showing data:

\begin{table}[h]
\centering
\begin{tabular}{@{}lrrr@{}}
\toprule
\textbf{Category} & \textbf{Value 1} & \textbf{Value 2} & \textbf{Total} \\
\midrule
Item A & 10.5 & 20.3 & 30.8 \\
Item B & 15.2 & 25.1 & 40.3 \\
Item C & 22.1 & 12.7 & 34.8 \\
\bottomrule
\end{tabular}
\caption{Sample Data Table}
\label{tab:sample}
\end{table}

\section{Lists}
\subsection{Bullet Points}
\begin{itemize}
  \item First item
  \item Second item
  \item Third item with sub-items:
    \begin{itemize}
      \item Sub-item 1
      \item Sub-item 2
    \end{itemize}
\end{itemize}

\subsection{Numbered List}
\begin{enumerate}
  \item First step
  \item Second step
  \item Third step
\end{enumerate}

\section{Citations and References}
When you use the citation search feature, references will be automatically managed. For example, you might cite a paper like this: \cite{example2024}. All citations will be automatically added to your \texttt{cite.bib} file.

\section{Conclusion}
This template demonstrates basic \LaTeX{} features including mathematical formulas, tables, and lists. The bibliography system is now set up to work seamlessly with Bibby AI's citation management. Feel free to modify it for your needs!

\bibliographystyle{plain}
\bibliography{cite}

\end{document}\documentclass{article}
\usepackage[utf8]{inputenc}
\usepackage{booktabs}
\usepackage{amsmath}
\usepackage{natbib}

\title{My Document}
\author{Your Name}
\date{\today}

\begin{document}

\maketitle

\section{Introduction}
This is a sample document demonstrating various \LaTeX{} features including mathematical equations, tables, and formatting. You can cite papers using the citation search feature, and they will be automatically added to your \texttt{cite.bib} file.

\section{Mathematical Expressions}
Here are some example equations:

\subsection{Quadratic Formula}
The quadratic formula is given by:
\[x = \frac{-b \pm \sqrt{b^2 - 4ac}}{2a}\]

\subsection{Matrix Example}
A sample matrix:
\[
\begin{pmatrix}
a_{11} & a_{12} & a_{13} \\
a_{21} & a_{22} & a_{23} \\
a_{31} & a_{32} & a_{33}
\end{pmatrix}
\]

\section{Tables}
Below is a sample table showing data:

\begin{table}[h]
\centering
\begin{tabular}{@{}lrrr@{}}
\toprule
\textbf{Category} & \textbf{Value 1} & \textbf{Value 2} & \textbf{Total} \\
\midrule
Item A & 10.5 & 20.3 & 30.8 \\
Item B & 15.2 & 25.1 & 40.3 \\
Item C & 22.1 & 12.7 & 34.8 \\
\bottomrule
\end{tabular}
\caption{Sample Data Table}
\label{tab:sample}
\end{table}

\section{Lists}
\subsection{Bullet Points}
\begin{itemize}
  \item First item
  \item Second item
  \item Third item with sub-items:
    \begin{itemize}
      \item Sub-item 1
      \item Sub-item 2
    \end{itemize}
\end{itemize}

\subsection{Numbered List}
\begin{enumerate}
  \item First step
  \item Second step
  \item Third step
\end{enumerate}

\section{Citations and References}
When you use the citation search feature, references will be automatically managed. For example, you might cite a paper like this: \cite{example2024}. All citations will be automatically added to your \texttt{cite.bib} file.

\section{Conclusion}
This template demonstrates basic \LaTeX{} features including mathematical formulas, tables, and lists. The bibliography system is now set up to work seamlessly with Bibby AI's citation management. Feel free to modify it for your needs!

\bibliographystyle{plain}
\bibliography{cite}

\end{document}\documentclass{article}
\usepackage[utf8]{inputenc}
\usepackage{booktabs}
\usepackage{amsmath}
\usepackage{natbib}

\title{My Document}
\author{Your Name}
\date{\today}

\begin{document}

\maketitle

\section{Introduction}
This is a sample document demonstrating various \LaTeX{} features including mathematical equations, tables, and formatting. You can cite papers using the citation search feature, and they will be automatically added to your \texttt{cite.bib} file.

\section{Mathematical Expressions}
Here are some example equations:

\subsection{Quadratic Formula}
The quadratic formula is given by:
\[x = \frac{-b \pm \sqrt{b^2 - 4ac}}{2a}\]

\subsection{Matrix Example}
A sample matrix:
\[
\begin{pmatrix}
a_{11} & a_{12} & a_{13} \\
a_{21} & a_{22} & a_{23} \\
a_{31} & a_{32} & a_{33}
\end{pmatrix}
\]

\section{Tables}
Below is a sample table showing data:

\begin{table}[h]
\centering
\begin{tabular}{@{}lrrr@{}}
\toprule
\textbf{Category} & \textbf{Value 1} & \textbf{Value 2} & \textbf{Total} \\
\midrule
Item A & 10.5 & 20.3 & 30.8 \\
Item B & 15.2 & 25.1 & 40.3 \\
Item C & 22.1 & 12.7 & 34.8 \\
\bottomrule
\end{tabular}
\caption{Sample Data Table}
\label{tab:sample}
\end{table}

\section{Lists}
\subsection{Bullet Points}
\begin{itemize}
  \item First item
  \item Second item
  \item Third item with sub-items:
    \begin{itemize}
      \item Sub-item 1
      \item Sub-item 2
    \end{itemize}
\end{itemize}

\subsection{Numbered List}
\begin{enumerate}
  \item First step
  \item Second step
  \item Third step
\end{enumerate}

\section{Citations and References}
When you use the citation search feature, references will be automatically managed. For example, you might cite a paper like this: \cite{example2024}. All citations will be automatically added to your \texttt{cite.bib} file.

\section{Conclusion}
This template demonstrates basic \LaTeX{} features including mathematical formulas, tables, and lists. The bibliography system is now set up to work seamlessly with Bibby AI's citation management. Feel free to modify it for your needs!

\bibliographystyle{plain}
\bibliography{cite}

\end{document}\documentclass{article}
\usepackage[utf8]{inputenc}
\usepackage{booktabs}
\usepackage{amsmath}
\usepackage{natbib}

\title{My Document}
\author{Your Name}
\date{\today}

\begin{document}

\maketitle

\section{Introduction}
This is a sample document demonstrating various \LaTeX{} features including mathematical equations, tables, and formatting. You can cite papers using the citation search feature, and they will be automatically added to your \texttt{cite.bib} file.

\section{Mathematical Expressions}
Here are some example equations:

\subsection{Quadratic Formula}
The quadratic formula is given by:
\[x = \frac{-b \pm \sqrt{b^2 - 4ac}}{2a}\]

\subsection{Matrix Example}
A sample matrix:
\[
\begin{pmatrix}
a_{11} & a_{12} & a_{13} \\
a_{21} & a_{22} & a_{23} \\
a_{31} & a_{32} & a_{33}
\end{pmatrix}
\]

\section{Tables}
Below is a sample table showing data:

\begin{table}[h]
\centering
\begin{tabular}{@{}lrrr@{}}
\toprule
\textbf{Category} & \textbf{Value 1} & \textbf{Value 2} & \textbf{Total} \\
\midrule
Item A & 10.5 & 20.3 & 30.8 \\
Item B & 15.2 & 25.1 & 40.3 \\
Item C & 22.1 & 12.7 & 34.8 \\
\bottomrule
\end{tabular}
\caption{Sample Data Table}
\label{tab:sample}
\end{table}

\section{Lists}
\subsection{Bullet Points}
\begin{itemize}
  \item First item
  \item Second item
  \item Third item with sub-items:
    \begin{itemize}
      \item Sub-item 1
      \item Sub-item 2
    \end{itemize}
\end{itemize}

\subsection{Numbered List}
\begin{enumerate}
  \item First step
  \item Second step
  \item Third step
\end{enumerate}

\section{Citations and References}
When you use the citation search feature, references will be automatically managed. For example, you might cite a paper like this: \cite{example2024}. All citations will be automatically added to your \texttt{cite.bib} file.

\section{Conclusion}
This template demonstrates basic \LaTeX{} features including mathematical formulas, tables, and lists. The bibliography system is now set up to work seamlessly with Bibby AI's citation management. Feel free to modify it for your needs!

\bibliographystyle{plain}
\bibliography{cite}

\end{document}\documentclass{article}
\usepackage[utf8]{inputenc}
\usepackage{booktabs}
\usepackage{amsmath}
\usepackage{natbib}

\title{My Document}
\author{Your Name}
\date{\today}

\begin{document}

\maketitle

\section{Introduction}
This is a sample document demonstrating various \LaTeX{} features including mathematical equations, tables, and formatting. You can cite papers using the citation search feature, and they will be automatically added to your \texttt{cite.bib} file.

\section{Mathematical Expressions}
Here are some example equations:

\subsection{Quadratic Formula}
The quadratic formula is given by:
\[x = \frac{-b \pm \sqrt{b^2 - 4ac}}{2a}\]

\subsection{Matrix Example}
A sample matrix:
\[
\begin{pmatrix}
a_{11} & a_{12} & a_{13} \\
a_{21} & a_{22} & a_{23} \\
a_{31} & a_{32} & a_{33}
\end{pmatrix}
\]

\section{Tables}
Below is a sample table showing data:

\begin{table}[h]
\centering
\begin{tabular}{@{}lrrr@{}}
\toprule
\textbf{Category} & \textbf{Value 1} & \textbf{Value 2} & \textbf{Total} \\
\midrule
Item A & 10.5 & 20.3 & 30.8 \\
Item B & 15.2 & 25.1 & 40.3 \\
Item C & 22.1 & 12.7 & 34.8 \\
\bottomrule
\end{tabular}
\caption{Sample Data Table}
\label{tab:sample}
\end{table}

\section{Lists}
\subsection{Bullet Points}
\begin{itemize}
  \item First item
  \item Second item
  \item Third item with sub-items:
    \begin{itemize}
      \item Sub-item 1
      \item Sub-item 2
    \end{itemize}
\end{itemize}

\subsection{Numbered List}
\begin{enumerate}
  \item First step
  \item Second step
  \item Third step
\end{enumerate}

\section{Citations and References}
When you use the citation search feature, references will be automatically managed. For example, you might cite a paper like this: \cite{example2024}. All citations will be automatically added to your \texttt{cite.bib} file.

\section{Conclusion}
This template demonstrates basic \LaTeX{} features including mathematical formulas, tables, and lists. The bibliography system is now set up to work seamlessly with Bibby AI's citation management. Feel free to modify it for your needs!

\bibliographystyle{plain}
\bibliography{cite}

\end{document}\documentclass{article}
\usepackage[utf8]{inputenc}
\usepackage{booktabs}
\usepackage{amsmath}
\usepackage{natbib}

\title{My Document}
\author{Your Name}
\date{\today}

\begin{document}

\maketitle

\section{Introduction}
This is a sample document demonstrating various \LaTeX{} features including mathematical equations, tables, and formatting. You can cite papers using the citation search feature, and they will be automatically added to your \texttt{cite.bib} file.

\section{Mathematical Expressions}
Here are some example equations:

\subsection{Quadratic Formula}
The quadratic formula is given by:
\[x = \frac{-b \pm \sqrt{b^2 - 4ac}}{2a}\]

\subsection{Matrix Example}
A sample matrix:
\[
\begin{pmatrix}
a_{11} & a_{12} & a_{13} \\
a_{21} & a_{22} & a_{23} \\
a_{31} & a_{32} & a_{33}
\end{pmatrix}
\]

\section{Tables}
Below is a sample table showing data:

\begin{table}[h]
\centering
\begin{tabular}{@{}lrrr@{}}
\toprule
\textbf{Category} & \textbf{Value 1} & \textbf{Value 2} & \textbf{Total} \\
\midrule
Item A & 10.5 & 20.3 & 30.8 \\
Item B & 15.2 & 25.1 & 40.3 \\
Item C & 22.1 & 12.7 & 34.8 \\
\bottomrule
\end{tabular}
\caption{Sample Data Table}
\label{tab:sample}
\end{table}

\section{Lists}
\subsection{Bullet Points}
\begin{itemize}
  \item First item
  \item Second item
  \item Third item with sub-items:
    \begin{itemize}
      \item Sub-item 1
      \item Sub-item 2
    \end{itemize}
\end{itemize}

\subsection{Numbered List}
\begin{enumerate}
  \item First step
  \item Second step
  \item Third step
\end{enumerate}

\section{Citations and References}
When you use the citation search feature, references will be automatically managed. For example, you might cite a paper like this: \cite{example2024}. All citations will be automatically added to your \texttt{cite.bib} file.

\section{Conclusion}
This template demonstrates basic \LaTeX{} features including mathematical formulas, tables, and lists. The bibliography system is now set up to work seamlessly with Bibby AI's citation management. Feel free to modify it for your needs!

\bibliographystyle{plain}
\bibliography{cite}

\end{document}\documentclass{article}
\usepackage[utf8]{inputenc}
\usepackage{booktabs}
\usepackage{amsmath}
\usepackage{natbib}

\title{My Document}
\author{Your Name}
\date{\today}

\begin{document}

\maketitle

\section{Introduction}
This is a sample document demonstrating various \LaTeX{} features including mathematical equations, tables, and formatting. You can cite papers using the citation search feature, and they will be automatically added to your \texttt{cite.bib} file.

\section{Mathematical Expressions}
Here are some example equations:

\subsection{Quadratic Formula}
The quadratic formula is given by:
\[x = \frac{-b \pm \sqrt{b^2 - 4ac}}{2a}\]

\subsection{Matrix Example}
A sample matrix:
\[
\begin{pmatrix}
a_{11} & a_{12} & a_{13} \\
a_{21} & a_{22} & a_{23} \\
a_{31} & a_{32} & a_{33}
\end{pmatrix}
\]

\section{Tables}
Below is a sample table showing data:

\begin{table}[h]
\centering
\begin{tabular}{@{}lrrr@{}}
\toprule
\textbf{Category} & \textbf{Value 1} & \textbf{Value 2} & \textbf{Total} \\
\midrule
Item A & 10.5 & 20.3 & 30.8 \\
Item B & 15.2 & 25.1 & 40.3 \\
Item C & 22.1 & 12.7 & 34.8 \\
\bottomrule
\end{tabular}
\caption{Sample Data Table}
\label{tab:sample}
\end{table}

\section{Lists}
\subsection{Bullet Points}
\begin{itemize}
  \item First item
  \item Second item
  \item Third item with sub-items:
    \begin{itemize}
      \item Sub-item 1
      \item Sub-item 2
    \end{itemize}
\end{itemize}

\subsection{Numbered List}
\begin{enumerate}
  \item First step
  \item Second step
  \item Third step
\end{enumerate}

\section{Citations and References}
When you use the citation search feature, references will be automatically managed. For example, you might cite a paper like this: \cite{example2024}. All citations will be automatically added to your \texttt{cite.bib} file.

\section{Conclusion}
This template demonstrates basic \LaTeX{} features including mathematical formulas, tables, and lists. The bibliography system is now set up to work seamlessly with Bibby AI's citation management. Feel free to modify it for your needs!

\bibliographystyle{plain}
\bibliography{cite}

\end{document}\documentclass{article}
\usepackage[utf8]{inputenc}
\usepackage{booktabs}
\usepackage{amsmath}
\usepackage{natbib}

\title{My Document}
\author{Your Name}
\date{\today}

\begin{document}

\maketitle

\section{Introduction}
This is a sample document demonstrating various \LaTeX{} features including mathematical equations, tables, and formatting. You can cite papers using the citation search feature, and they will be automatically added to your \texttt{cite.bib} file.

\section{Mathematical Expressions}
Here are some example equations:

\subsection{Quadratic Formula}
The quadratic formula is given by:
\[x = \frac{-b \pm \sqrt{b^2 - 4ac}}{2a}\]

\subsection{Matrix Example}
A sample matrix:
\[
\begin{pmatrix}
a_{11} & a_{12} & a_{13} \\
a_{21} & a_{22} & a_{23} \\
a_{31} & a_{32} & a_{33}
\end{pmatrix}
\]

\section{Tables}
Below is a sample table showing data:

\begin{table}[h]
\centering
\begin{tabular}{@{}lrrr@{}}
\toprule
\textbf{Category} & \textbf{Value 1} & \textbf{Value 2} & \textbf{Total} \\
\midrule
Item A & 10.5 & 20.3 & 30.8 \\
Item B & 15.2 & 25.1 & 40.3 \\
Item C & 22.1 & 12.7 & 34.8 \\
\bottomrule
\end{tabular}
\caption{Sample Data Table}
\label{tab:sample}
\end{table}

\section{Lists}
\subsection{Bullet Points}
\begin{itemize}
  \item First item
  \item Second item
  \item Third item with sub-items:
    \begin{itemize}
      \item Sub-item 1
      \item Sub-item 2
    \end{itemize}
\end{itemize}

\subsection{Numbered List}
\begin{enumerate}
  \item First step
  \item Second step
  \item Third step
\end{enumerate}

\section{Citations and References}
When you use the citation search feature, references will be automatically managed. For example, you might cite a paper like this: \cite{example2024}. All citations will be automatically added to your \texttt{cite.bib} file.

\section{Conclusion}
This template demonstrates basic \LaTeX{} features including mathematical formulas, tables, and lists. The bibliography system is now set up to work seamlessly with Bibby AI's citation management. Feel free to modify it for your needs!

\bibliographystyle{plain}
\bibliography{cite}

\end{document}\documentclass{article}
\usepackage[utf8]{inputenc}
\usepackage{booktabs}
\usepackage{amsmath}
\usepackage{natbib}

\title{My Document}
\author{Your Name}
\date{\today}

\begin{document}

\maketitle

\section{Introduction}
This is a sample document demonstrating various \LaTeX{} features including mathematical equations, tables, and formatting. You can cite papers using the citation search feature, and they will be automatically added to your \texttt{cite.bib} file.

\section{Mathematical Expressions}
Here are some example equations:

\subsection{Quadratic Formula}
The quadratic formula is given by:
\[x = \frac{-b \pm \sqrt{b^2 - 4ac}}{2a}\]

\subsection{Matrix Example}
A sample matrix:
\[
\begin{pmatrix}
a_{11} & a_{12} & a_{13} \\
a_{21} & a_{22} & a_{23} \\
a_{31} & a_{32} & a_{33}
\end{pmatrix}
\]

\section{Tables}
Below is a sample table showing data:

\begin{table}[h]
\centering
\begin{tabular}{@{}lrrr@{}}
\toprule
\textbf{Category} & \textbf{Value 1} & \textbf{Value 2} & \textbf{Total} \\
\midrule
Item A & 10.5 & 20.3 & 30.8 \\
Item B & 15.2 & 25.1 & 40.3 \\
Item C & 22.1 & 12.7 & 34.8 \\
\bottomrule
\end{tabular}
\caption{Sample Data Table}
\label{tab:sample}
\end{table}

\section{Lists}
\subsection{Bullet Points}
\begin{itemize}
  \item First item
  \item Second item
  \item Third item with sub-items:
    \begin{itemize}
      \item Sub-item 1
      \item Sub-item 2
    \end{itemize}
\end{itemize}

\subsection{Numbered List}
\begin{enumerate}
  \item First step
  \item Second step
  \item Third step
\end{enumerate}

\section{Citations and References}
When you use the citation search feature, references will be automatically managed. For example, you might cite a paper like this: \cite{example2024}. All citations will be automatically added to your \texttt{cite.bib} file.

\section{Conclusion}
This template demonstrates basic \LaTeX{} features including mathematical formulas, tables, and lists. The bibliography system is now set up to work seamlessly with Bibby AI's citation management. Feel free to modify it for your needs!

\bibliographystyle{plain}
\bibliography{cite}

\end{document}\documentclass{article}
\usepackage[utf8]{inputenc}
\usepackage{booktabs}
\usepackage{amsmath}
\usepackage{natbib}

\title{My Document}
\author{Your Name}
\date{\today}

\begin{document}

\maketitle

\section{Introduction}
This is a sample document demonstrating various \LaTeX{} features including mathematical equations, tables, and formatting. You can cite papers using the citation search feature, and they will be automatically added to your \texttt{cite.bib} file.

\section{Mathematical Expressions}
Here are some example equations:

\subsection{Quadratic Formula}
The quadratic formula is given by:
\[x = \frac{-b \pm \sqrt{b^2 - 4ac}}{2a}\]

\subsection{Matrix Example}
A sample matrix:
\[
\begin{pmatrix}
a_{11} & a_{12} & a_{13} \\
a_{21} & a_{22} & a_{23} \\
a_{31} & a_{32} & a_{33}
\end{pmatrix}
\]

\section{Tables}
Below is a sample table showing data:

\begin{table}[h]
\centering
\begin{tabular}{@{}lrrr@{}}
\toprule
\textbf{Category} & \textbf{Value 1} & \textbf{Value 2} & \textbf{Total} \\
\midrule
Item A & 10.5 & 20.3 & 30.8 \\
Item B & 15.2 & 25.1 & 40.3 \\
Item C & 22.1 & 12.7 & 34.8 \\
\bottomrule
\end{tabular}
\caption{Sample Data Table}
\label{tab:sample}
\end{table}

\section{Lists}
\subsection{Bullet Points}
\begin{itemize}
  \item First item
  \item Second item
  \item Third item with sub-items:
    \begin{itemize}
      \item Sub-item 1
      \item Sub-item 2
    \end{itemize}
\end{itemize}

\subsection{Numbered List}
\begin{enumerate}
  \item First step
  \item Second step
  \item Third step
\end{enumerate}

\section{Citations and References}
When you use the citation search feature, references will be automatically managed. For example, you might cite a paper like this: \cite{example2024}. All citations will be automatically added to your \texttt{cite.bib} file.

\section{Conclusion}
This template demonstrates basic \LaTeX{} features including mathematical formulas, tables, and lists. The bibliography system is now set up to work seamlessly with Bibby AI's citation management. Feel free to modify it for your needs!

\bibliographystyle{plain}
\bibliography{cite}

\end{document}\documentclass{article}
\usepackage[utf8]{inputenc}
\usepackage{booktabs}
\usepackage{amsmath}
\usepackage{natbib}

\title{My Document}
\author{Your Name}
\date{\today}

\begin{document}

\maketitle

\section{Introduction}
This is a sample document demonstrating various \LaTeX{} features including mathematical equations, tables, and formatting. You can cite papers using the citation search feature, and they will be automatically added to your \texttt{cite.bib} file.

\section{Mathematical Expressions}
Here are some example equations:

\subsection{Quadratic Formula}
The quadratic formula is given by:
\[x = \frac{-b \pm \sqrt{b^2 - 4ac}}{2a}\]

\subsection{Matrix Example}
A sample matrix:
\[
\begin{pmatrix}
a_{11} & a_{12} & a_{13} \\
a_{21} & a_{22} & a_{23} \\
a_{31} & a_{32} & a_{33}
\end{pmatrix}
\]

\section{Tables}
Below is a sample table showing data:

\begin{table}[h]
\centering
\begin{tabular}{@{}lrrr@{}}
\toprule
\textbf{Category} & \textbf{Value 1} & \textbf{Value 2} & \textbf{Total} \\
\midrule
Item A & 10.5 & 20.3 & 30.8 \\
Item B & 15.2 & 25.1 & 40.3 \\
Item C & 22.1 & 12.7 & 34.8 \\
\bottomrule
\end{tabular}
\caption{Sample Data Table}
\label{tab:sample}
\end{table}

\section{Lists}
\subsection{Bullet Points}
\begin{itemize}
  \item First item
  \item Second item
  \item Third item with sub-items:
    \begin{itemize}
      \item Sub-item 1
      \item Sub-item 2
    \end{itemize}
\end{itemize}

\subsection{Numbered List}
\begin{enumerate}
  \item First step
  \item Second step
  \item Third step
\end{enumerate}

\section{Citations and References}
When you use the citation search feature, references will be automatically managed. For example, you might cite a paper like this: \cite{example2024}. All citations will be automatically added to your \texttt{cite.bib} file.

\section{Conclusion}
This template demonstrates basic \LaTeX{} features including mathematical formulas, tables, and lists. The bibliography system is now set up to work seamlessly with Bibby AI's citation management. Feel free to modify it for your needs!

\bibliographystyle{plain}
\bibliography{cite}

\end{document}\documentclass{article}
\usepackage[utf8]{inputenc}
\usepackage{booktabs}
\usepackage{amsmath}
\usepackage{natbib}

\title{My Document}
\author{Your Name}
\date{\today}

\begin{document}

\maketitle

\section{Introduction}
This is a sample document demonstrating various \LaTeX{} features including mathematical equations, tables, and formatting. You can cite papers using the citation search feature, and they will be automatically added to your \texttt{cite.bib} file.

\section{Mathematical Expressions}
Here are some example equations:

\subsection{Quadratic Formula}
The quadratic formula is given by:
\[x = \frac{-b \pm \sqrt{b^2 - 4ac}}{2a}\]

\subsection{Matrix Example}
A sample matrix:
\[
\begin{pmatrix}
a_{11} & a_{12} & a_{13} \\
a_{21} & a_{22} & a_{23} \\
a_{31} & a_{32} & a_{33}
\end{pmatrix}
\]

\section{Tables}
Below is a sample table showing data:

\begin{table}[h]
\centering
\begin{tabular}{@{}lrrr@{}}
\toprule
\textbf{Category} & \textbf{Value 1} & \textbf{Value 2} & \textbf{Total} \\
\midrule
Item A & 10.5 & 20.3 & 30.8 \\
Item B & 15.2 & 25.1 & 40.3 \\
Item C & 22.1 & 12.7 & 34.8 \\
\bottomrule
\end{tabular}
\caption{Sample Data Table}
\label{tab:sample}
\end{table}

\section{Lists}
\subsection{Bullet Points}
\begin{itemize}
  \item First item
  \item Second item
  \item Third item with sub-items:
    \begin{itemize}
      \item Sub-item 1
      \item Sub-item 2
    \end{itemize}
\end{itemize}

\subsection{Numbered List}
\begin{enumerate}
  \item First step
  \item Second step
  \item Third step
\end{enumerate}

\section{Citations and References}
When you use the citation search feature, references will be automatically managed. For example, you might cite a paper like this: \cite{example2024}. All citations will be automatically added to your \texttt{cite.bib} file.

\section{Conclusion}
This template demonstrates basic \LaTeX{} features including mathematical formulas, tables, and lists. The bibliography system is now set up to work seamlessly with Bibby AI's citation management. Feel free to modify it for your needs!

\bibliographystyle{plain}
\bibliography{cite}

\end{document}\documentclass{article}
\usepackage[utf8]{inputenc}
\usepackage{booktabs}
\usepackage{amsmath}
\usepackage{natbib}

\title{My Document}
\author{Your Name}
\date{\today}

\begin{document}

\maketitle

\section{Introduction}
This is a sample document demonstrating various \LaTeX{} features including mathematical equations, tables, and formatting. You can cite papers using the citation search feature, and they will be automatically added to your \texttt{cite.bib} file.

\section{Mathematical Expressions}
Here are some example equations:

\subsection{Quadratic Formula}
The quadratic formula is given by:
\[x = \frac{-b \pm \sqrt{b^2 - 4ac}}{2a}\]

\subsection{Matrix Example}
A sample matrix:
\[
\begin{pmatrix}
a_{11} & a_{12} & a_{13} \\
a_{21} & a_{22} & a_{23} \\
a_{31} & a_{32} & a_{33}
\end{pmatrix}
\]

\section{Tables}
Below is a sample table showing data:

\begin{table}[h]
\centering
\begin{tabular}{@{}lrrr@{}}
\toprule
\textbf{Category} & \textbf{Value 1} & \textbf{Value 2} & \textbf{Total} \\
\midrule
Item A & 10.5 & 20.3 & 30.8 \\
Item B & 15.2 & 25.1 & 40.3 \\
Item C & 22.1 & 12.7 & 34.8 \\
\bottomrule
\end{tabular}
\caption{Sample Data Table}
\label{tab:sample}
\end{table}

\section{Lists}
\subsection{Bullet Points}
\begin{itemize}
  \item First item
  \item Second item
  \item Third item with sub-items:
    \begin{itemize}
      \item Sub-item 1
      \item Sub-item 2
    \end{itemize}
\end{itemize}

\subsection{Numbered List}
\begin{enumerate}
  \item First step
  \item Second step
  \item Third step
\end{enumerate}

\section{Citations and References}
When you use the citation search feature, references will be automatically managed. For example, you might cite a paper like this: \cite{example2024}. All citations will be automatically added to your \texttt{cite.bib} file.

\section{Conclusion}
This template demonstrates basic \LaTeX{} features including mathematical formulas, tables, and lists. The bibliography system is now set up to work seamlessly with Bibby AI's citation management. Feel free to modify it for your needs!

\bibliographystyle{plain}
\bibliography{cite}

\end{document}\documentclass{article}
\usepackage[utf8]{inputenc}
\usepackage{booktabs}
\usepackage{amsmath}
\usepackage{natbib}

\title{My Document}
\author{Your Name}
\date{\today}

\begin{document}

\maketitle

\section{Introduction}
This is a sample document demonstrating various \LaTeX{} features including mathematical equations, tables, and formatting. You can cite papers using the citation search feature, and they will be automatically added to your \texttt{cite.bib} file.

\section{Mathematical Expressions}
Here are some example equations:

\subsection{Quadratic Formula}
The quadratic formula is given by:
\[x = \frac{-b \pm \sqrt{b^2 - 4ac}}{2a}\]

\subsection{Matrix Example}
A sample matrix:
\[
\begin{pmatrix}
a_{11} & a_{12} & a_{13} \\
a_{21} & a_{22} & a_{23} \\
a_{31} & a_{32} & a_{33}
\end{pmatrix}
\]

\section{Tables}
Below is a sample table showing data:

\begin{table}[h]
\centering
\begin{tabular}{@{}lrrr@{}}
\toprule
\textbf{Category} & \textbf{Value 1} & \textbf{Value 2} & \textbf{Total} \\
\midrule
Item A & 10.5 & 20.3 & 30.8 \\
Item B & 15.2 & 25.1 & 40.3 \\
Item C & 22.1 & 12.7 & 34.8 \\
\bottomrule
\end{tabular}
\caption{Sample Data Table}
\label{tab:sample}
\end{table}

\section{Lists}
\subsection{Bullet Points}
\begin{itemize}
  \item First item
  \item Second item
  \item Third item with sub-items:
    \begin{itemize}
      \item Sub-item 1
      \item Sub-item 2
    \end{itemize}
\end{itemize}

\subsection{Numbered List}
\begin{enumerate}
  \item First step
  \item Second step
  \item Third step
\end{enumerate}

\section{Citations and References}
When you use the citation search feature, references will be automatically managed. For example, you might cite a paper like this: \cite{example2024}. All citations will be automatically added to your \texttt{cite.bib} file.

\section{Conclusion}
This template demonstrates basic \LaTeX{} features including mathematical formulas, tables, and lists. The bibliography system is now set up to work seamlessly with Bibby AI's citation management. Feel free to modify it for your needs!

\bibliographystyle{plain}
\bibliography{cite}

\end{document}\documentclass{article}
\usepackage[utf8]{inputenc}
\usepackage{booktabs}
\usepackage{amsmath}
\usepackage{natbib}

\title{My Document}
\author{Your Name}
\date{\today}

\begin{document}

\maketitle

\section{Introduction}
This is a sample document demonstrating various \LaTeX{} features including mathematical equations, tables, and formatting. You can cite papers using the citation search feature, and they will be automatically added to your \texttt{cite.bib} file.

\section{Mathematical Expressions}
Here are some example equations:

\subsection{Quadratic Formula}
The quadratic formula is given by:
\[x = \frac{-b \pm \sqrt{b^2 - 4ac}}{2a}\]

\subsection{Matrix Example}
A sample matrix:
\[
\begin{pmatrix}
a_{11} & a_{12} & a_{13} \\
a_{21} & a_{22} & a_{23} \\
a_{31} & a_{32} & a_{33}
\end{pmatrix}
\]

\section{Tables}
Below is a sample table showing data:

\begin{table}[h]
\centering
\begin{tabular}{@{}lrrr@{}}
\toprule
\textbf{Category} & \textbf{Value 1} & \textbf{Value 2} & \textbf{Total} \\
\midrule
Item A & 10.5 & 20.3 & 30.8 \\
Item B & 15.2 & 25.1 & 40.3 \\
Item C & 22.1 & 12.7 & 34.8 \\
\bottomrule
\end{tabular}
\caption{Sample Data Table}
\label{tab:sample}
\end{table}

\section{Lists}
\subsection{Bullet Points}
\begin{itemize}
  \item First item
  \item Second item
  \item Third item with sub-items:
    \begin{itemize}
      \item Sub-item 1
      \item Sub-item 2
    \end{itemize}
\end{itemize}

\subsection{Numbered List}
\begin{enumerate}
  \item First step
  \item Second step
  \item Third step
\end{enumerate}

\section{Citations and References}
When you use the citation search feature, references will be automatically managed. For example, you might cite a paper like this: \cite{example2024}. All citations will be automatically added to your \texttt{cite.bib} file.

\section{Conclusion}
This template demonstrates basic \LaTeX{} features including mathematical formulas, tables, and lists. The bibliography system is now set up to work seamlessly with Bibby AI's citation management. Feel free to modify it for your needs!

\bibliographystyle{plain}
\bibliography{cite}

\end{document}\documentclass{article}
\usepackage[utf8]{inputenc}
\usepackage{booktabs}
\usepackage{amsmath}
\usepackage{natbib}

\title{My Document}
\author{Your Name}
\date{\today}

\begin{document}

\maketitle

\section{Introduction}
This is a sample document demonstrating various \LaTeX{} features including mathematical equations, tables, and formatting. You can cite papers using the citation search feature, and they will be automatically added to your \texttt{cite.bib} file.

\section{Mathematical Expressions}
Here are some example equations:

\subsection{Quadratic Formula}
The quadratic formula is given by:
\[x = \frac{-b \pm \sqrt{b^2 - 4ac}}{2a}\]

\subsection{Matrix Example}
A sample matrix:
\[
\begin{pmatrix}
a_{11} & a_{12} & a_{13} \\
a_{21} & a_{22} & a_{23} \\
a_{31} & a_{32} & a_{33}
\end{pmatrix}
\]

\section{Tables}
Below is a sample table showing data:

\begin{table}[h]
\centering
\begin{tabular}{@{}lrrr@{}}
\toprule
\textbf{Category} & \textbf{Value 1} & \textbf{Value 2} & \textbf{Total} \\
\midrule
Item A & 10.5 & 20.3 & 30.8 \\
Item B & 15.2 & 25.1 & 40.3 \\
Item C & 22.1 & 12.7 & 34.8 \\
\bottomrule
\end{tabular}
\caption{Sample Data Table}
\label{tab:sample}
\end{table}

\section{Lists}
\subsection{Bullet Points}
\begin{itemize}
  \item First item
  \item Second item
  \item Third item with sub-items:
    \begin{itemize}
      \item Sub-item 1
      \item Sub-item 2
    \end{itemize}
\end{itemize}

\subsection{Numbered List}
\begin{enumerate}
  \item First step
  \item Second step
  \item Third step
\end{enumerate}

\section{Citations and References}
When you use the citation search feature, references will be automatically managed. For example, you might cite a paper like this: \cite{example2024}. All citations will be automatically added to your \texttt{cite.bib} file.

\section{Conclusion}
This template demonstrates basic \LaTeX{} features including mathematical formulas, tables, and lists. The bibliography system is now set up to work seamlessly with Bibby AI's citation management. Feel free to modify it for your needs!

\bibliographystyle{plain}
\bibliography{cite}

\end{document}\documentclass{article}
\usepackage[utf8]{inputenc}
\usepackage{booktabs}
\usepackage{amsmath}
\usepackage{natbib}

\title{My Document}
\author{Your Name}
\date{\today}

\begin{document}

\maketitle

\section{Introduction}
This is a sample document demonstrating various \LaTeX{} features including mathematical equations, tables, and formatting. You can cite papers using the citation search feature, and they will be automatically added to your \texttt{cite.bib} file.

\section{Mathematical Expressions}
Here are some example equations:

\subsection{Quadratic Formula}
The quadratic formula is given by:
\[x = \frac{-b \pm \sqrt{b^2 - 4ac}}{2a}\]

\subsection{Matrix Example}
A sample matrix:
\[
\begin{pmatrix}
a_{11} & a_{12} & a_{13} \\
a_{21} & a_{22} & a_{23} \\
a_{31} & a_{32} & a_{33}
\end{pmatrix}
\]

\section{Tables}
Below is a sample table showing data:

\begin{table}[h]
\centering
\begin{tabular}{@{}lrrr@{}}
\toprule
\textbf{Category} & \textbf{Value 1} & \textbf{Value 2} & \textbf{Total} \\
\midrule
Item A & 10.5 & 20.3 & 30.8 \\
Item B & 15.2 & 25.1 & 40.3 \\
Item C & 22.1 & 12.7 & 34.8 \\
\bottomrule
\end{tabular}
\caption{Sample Data Table}
\label{tab:sample}
\end{table}

\section{Lists}
\subsection{Bullet Points}
\begin{itemize}
  \item First item
  \item Second item
  \item Third item with sub-items:
    \begin{itemize}
      \item Sub-item 1
      \item Sub-item 2
    \end{itemize}
\end{itemize}

\subsection{Numbered List}
\begin{enumerate}
  \item First step
  \item Second step
  \item Third step
\end{enumerate}

\section{Citations and References}
When you use the citation search feature, references will be automatically managed. For example, you might cite a paper like this: \cite{example2024}. All citations will be automatically added to your \texttt{cite.bib} file.

\section{Conclusion}
This template demonstrates basic \LaTeX{} features including mathematical formulas, tables, and lists. The bibliography system is now set up to work seamlessly with Bibby AI's citation management. Feel free to modify it for your needs!

\bibliographystyle{plain}
\bibliography{cite}

\end{document}\documentclass{article}
\usepackage[utf8]{inputenc}
\usepackage{booktabs}
\usepackage{amsmath}
\usepackage{natbib}

\title{My Document}
\author{Your Name}
\date{\today}

\begin{document}

\maketitle

\section{Introduction}
This is a sample document demonstrating various \LaTeX{} features including mathematical equations, tables, and formatting. You can cite papers using the citation search feature, and they will be automatically added to your \texttt{cite.bib} file.

\section{Mathematical Expressions}
Here are some example equations:

\subsection{Quadratic Formula}
The quadratic formula is given by:
\[x = \frac{-b \pm \sqrt{b^2 - 4ac}}{2a}\]

\subsection{Matrix Example}
A sample matrix:
\[
\begin{pmatrix}
a_{11} & a_{12} & a_{13} \\
a_{21} & a_{22} & a_{23} \\
a_{31} & a_{32} & a_{33}
\end{pmatrix}
\]

\section{Tables}
Below is a sample table showing data:

\begin{table}[h]
\centering
\begin{tabular}{@{}lrrr@{}}
\toprule
\textbf{Category} & \textbf{Value 1} & \textbf{Value 2} & \textbf{Total} \\
\midrule
Item A & 10.5 & 20.3 & 30.8 \\
Item B & 15.2 & 25.1 & 40.3 \\
Item C & 22.1 & 12.7 & 34.8 \\
\bottomrule
\end{tabular}
\caption{Sample Data Table}
\label{tab:sample}
\end{table}

\section{Lists}
\subsection{Bullet Points}
\begin{itemize}
  \item First item
  \item Second item
  \item Third item with sub-items:
    \begin{itemize}
      \item Sub-item 1
      \item Sub-item 2
    \end{itemize}
\end{itemize}

\subsection{Numbered List}
\begin{enumerate}
  \item First step
  \item Second step
  \item Third step
\end{enumerate}

\section{Citations and References}
When you use the citation search feature, references will be automatically managed. For example, you might cite a paper like this: \cite{example2024}. All citations will be automatically added to your \texttt{cite.bib} file.

\section{Conclusion}
This template demonstrates basic \LaTeX{} features including mathematical formulas, tables, and lists. The bibliography system is now set up to work seamlessly with Bibby AI's citation management. Feel free to modify it for your needs!

\bibliographystyle{plain}
\bibliography{cite}

\end{document}\documentclass{article}
\usepackage[utf8]{inputenc}
\usepackage{booktabs}
\usepackage{amsmath}
\usepackage{natbib}

\title{My Document}
\author{Your Name}
\date{\today}

\begin{document}

\maketitle

\section{Introduction}
This is a sample document demonstrating various \LaTeX{} features including mathematical equations, tables, and formatting. You can cite papers using the citation search feature, and they will be automatically added to your \texttt{cite.bib} file.

\section{Mathematical Expressions}
Here are some example equations:

\subsection{Quadratic Formula}
The quadratic formula is given by:
\[x = \frac{-b \pm \sqrt{b^2 - 4ac}}{2a}\]

\subsection{Matrix Example}
A sample matrix:
\[
\begin{pmatrix}
a_{11} & a_{12} & a_{13} \\
a_{21} & a_{22} & a_{23} \\
a_{31} & a_{32} & a_{33}
\end{pmatrix}
\]

\section{Tables}
Below is a sample table showing data:

\begin{table}[h]
\centering
\begin{tabular}{@{}lrrr@{}}
\toprule
\textbf{Category} & \textbf{Value 1} & \textbf{Value 2} & \textbf{Total} \\
\midrule
Item A & 10.5 & 20.3 & 30.8 \\
Item B & 15.2 & 25.1 & 40.3 \\
Item C & 22.1 & 12.7 & 34.8 \\
\bottomrule
\end{tabular}
\caption{Sample Data Table}
\label{tab:sample}
\end{table}

\section{Lists}
\subsection{Bullet Points}
\begin{itemize}
  \item First item
  \item Second item
  \item Third item with sub-items:
    \begin{itemize}
      \item Sub-item 1
      \item Sub-item 2
    \end{itemize}
\end{itemize}

\subsection{Numbered List}
\begin{enumerate}
  \item First step
  \item Second step
  \item Third step
\end{enumerate}

\section{Citations and References}
When you use the citation search feature, references will be automatically managed. For example, you might cite a paper like this: \cite{example2024}. All citations will be automatically added to your \texttt{cite.bib} file.

\section{Conclusion}
This template demonstrates basic \LaTeX{} features including mathematical formulas, tables, and lists. The bibliography system is now set up to work seamlessly with Bibby AI's citation management. Feel free to modify it for your needs!

\bibliographystyle{plain}
\bibliography{cite}

\end{document}\documentclass{article}
\usepackage[utf8]{inputenc}
\usepackage{booktabs}
\usepackage{amsmath}
\usepackage{natbib}

\title{My Document}
\author{Your Name}
\date{\today}

\begin{document}

\maketitle

\section{Introduction}
This is a sample document demonstrating various \LaTeX{} features including mathematical equations, tables, and formatting. You can cite papers using the citation search feature, and they will be automatically added to your \texttt{cite.bib} file.

\section{Mathematical Expressions}
Here are some example equations:

\subsection{Quadratic Formula}
The quadratic formula is given by:
\[x = \frac{-b \pm \sqrt{b^2 - 4ac}}{2a}\]

\subsection{Matrix Example}
A sample matrix:
\[
\begin{pmatrix}
a_{11} & a_{12} & a_{13} \\
a_{21} & a_{22} & a_{23} \\
a_{31} & a_{32} & a_{33}
\end{pmatrix}
\]

\section{Tables}
Below is a sample table showing data:

\begin{table}[h]
\centering
\begin{tabular}{@{}lrrr@{}}
\toprule
\textbf{Category} & \textbf{Value 1} & \textbf{Value 2} & \textbf{Total} \\
\midrule
Item A & 10.5 & 20.3 & 30.8 \\
Item B & 15.2 & 25.1 & 40.3 \\
Item C & 22.1 & 12.7 & 34.8 \\
\bottomrule
\end{tabular}
\caption{Sample Data Table}
\label{tab:sample}
\end{table}

\section{Lists}
\subsection{Bullet Points}
\begin{itemize}
  \item First item
  \item Second item
  \item Third item with sub-items:
    \begin{itemize}
      \item Sub-item 1
      \item Sub-item 2
    \end{itemize}
\end{itemize}

\subsection{Numbered List}
\begin{enumerate}
  \item First step
  \item Second step
  \item Third step
\end{enumerate}

\section{Citations and References}
When you use the citation search feature, references will be automatically managed. For example, you might cite a paper like this: \cite{example2024}. All citations will be automatically added to your \texttt{cite.bib} file.

\section{Conclusion}
This template demonstrates basic \LaTeX{} features including mathematical formulas, tables, and lists. The bibliography system is now set up to work seamlessly with Bibby AI's citation management. Feel free to modify it for your needs!

\bibliographystyle{plain}
\bibliography{cite}

\end{document}\documentclass{article}
\usepackage[utf8]{inputenc}
\usepackage{booktabs}
\usepackage{amsmath}
\usepackage{natbib}

\title{My Document}
\author{Your Name}
\date{\today}

\begin{document}

\maketitle

\section{Introduction}
This is a sample document demonstrating various \LaTeX{} features including mathematical equations, tables, and formatting. You can cite papers using the citation search feature, and they will be automatically added to your \texttt{cite.bib} file.

\section{Mathematical Expressions}
Here are some example equations:

\subsection{Quadratic Formula}
The quadratic formula is given by:
\[x = \frac{-b \pm \sqrt{b^2 - 4ac}}{2a}\]

\subsection{Matrix Example}
A sample matrix:
\[
\begin{pmatrix}
a_{11} & a_{12} & a_{13} \\
a_{21} & a_{22} & a_{23} \\
a_{31} & a_{32} & a_{33}
\end{pmatrix}
\]

\section{Tables}
Below is a sample table showing data:

\begin{table}[h]
\centering
\begin{tabular}{@{}lrrr@{}}
\toprule
\textbf{Category} & \textbf{Value 1} & \textbf{Value 2} & \textbf{Total} \\
\midrule
Item A & 10.5 & 20.3 & 30.8 \\
Item B & 15.2 & 25.1 & 40.3 \\
Item C & 22.1 & 12.7 & 34.8 \\
\bottomrule
\end{tabular}
\caption{Sample Data Table}
\label{tab:sample}
\end{table}

\section{Lists}
\subsection{Bullet Points}
\begin{itemize}
  \item First item
  \item Second item
  \item Third item with sub-items:
    \begin{itemize}
      \item Sub-item 1
      \item Sub-item 2
    \end{itemize}
\end{itemize}

\subsection{Numbered List}
\begin{enumerate}
  \item First step
  \item Second step
  \item Third step
\end{enumerate}

\section{Citations and References}
When you use the citation search feature, references will be automatically managed. For example, you might cite a paper like this: \cite{example2024}. All citations will be automatically added to your \texttt{cite.bib} file.

\section{Conclusion}
This template demonstrates basic \LaTeX{} features including mathematical formulas, tables, and lists. The bibliography system is now set up to work seamlessly with Bibby AI's citation management. Feel free to modify it for your needs!

\bibliographystyle{plain}
\bibliography{cite}

\end{document}\documentclass{article}
\usepackage[utf8]{inputenc}
\usepackage{booktabs}
\usepackage{amsmath}
\usepackage{natbib}

\title{My Document}
\author{Your Name}
\date{\today}

\begin{document}

\maketitle

\section{Introduction}
This is a sample document demonstrating various \LaTeX{} features including mathematical equations, tables, and formatting. You can cite papers using the citation search feature, and they will be automatically added to your \texttt{cite.bib} file.

\section{Mathematical Expressions}
Here are some example equations:

\subsection{Quadratic Formula}
The quadratic formula is given by:
\[x = \frac{-b \pm \sqrt{b^2 - 4ac}}{2a}\]

\subsection{Matrix Example}
A sample matrix:
\[
\begin{pmatrix}
a_{11} & a_{12} & a_{13} \\
a_{21} & a_{22} & a_{23} \\
a_{31} & a_{32} & a_{33}
\end{pmatrix}
\]

\section{Tables}
Below is a sample table showing data:

\begin{table}[h]
\centering
\begin{tabular}{@{}lrrr@{}}
\toprule
\textbf{Category} & \textbf{Value 1} & \textbf{Value 2} & \textbf{Total} \\
\midrule
Item A & 10.5 & 20.3 & 30.8 \\
Item B & 15.2 & 25.1 & 40.3 \\
Item C & 22.1 & 12.7 & 34.8 \\
\bottomrule
\end{tabular}
\caption{Sample Data Table}
\label{tab:sample}
\end{table}

\section{Lists}
\subsection{Bullet Points}
\begin{itemize}
  \item First item
  \item Second item
  \item Third item with sub-items:
    \begin{itemize}
      \item Sub-item 1
      \item Sub-item 2
    \end{itemize}
\end{itemize}

\subsection{Numbered List}
\begin{enumerate}
  \item First step
  \item Second step
  \item Third step
\end{enumerate}

\section{Citations and References}
When you use the citation search feature, references will be automatically managed. For example, you might cite a paper like this: \cite{example2024}. All citations will be automatically added to your \texttt{cite.bib} file.

\section{Conclusion}
This template demonstrates basic \LaTeX{} features including mathematical formulas, tables, and lists. The bibliography system is now set up to work seamlessly with Bibby AI's citation management. Feel free to modify it for your needs!

\bibliographystyle{plain}
\bibliography{cite}

\end{document}\documentclass{article}
\usepackage[utf8]{inputenc}
\usepackage{booktabs}
\usepackage{amsmath}
\usepackage{natbib}

\title{My Document}
\author{Your Name}
\date{\today}

\begin{document}

\maketitle

\section{Introduction}
This is a sample document demonstrating various \LaTeX{} features including mathematical equations, tables, and formatting. You can cite papers using the citation search feature, and they will be automatically added to your \texttt{cite.bib} file.

\section{Mathematical Expressions}
Here are some example equations:

\subsection{Quadratic Formula}
The quadratic formula is given by:
\[x = \frac{-b \pm \sqrt{b^2 - 4ac}}{2a}\]

\subsection{Matrix Example}
A sample matrix:
\[
\begin{pmatrix}
a_{11} & a_{12} & a_{13} \\
a_{21} & a_{22} & a_{23} \\
a_{31} & a_{32} & a_{33}
\end{pmatrix}
\]

\section{Tables}
Below is a sample table showing data:

\begin{table}[h]
\centering
\begin{tabular}{@{}lrrr@{}}
\toprule
\textbf{Category} & \textbf{Value 1} & \textbf{Value 2} & \textbf{Total} \\
\midrule
Item A & 10.5 & 20.3 & 30.8 \\
Item B & 15.2 & 25.1 & 40.3 \\
Item C & 22.1 & 12.7 & 34.8 \\
\bottomrule
\end{tabular}
\caption{Sample Data Table}
\label{tab:sample}
\end{table}

\section{Lists}
\subsection{Bullet Points}
\begin{itemize}
  \item First item
  \item Second item
  \item Third item with sub-items:
    \begin{itemize}
      \item Sub-item 1
      \item Sub-item 2
    \end{itemize}
\end{itemize}

\subsection{Numbered List}
\begin{enumerate}
  \item First step
  \item Second step
  \item Third step
\end{enumerate}

\section{Citations and References}
When you use the citation search feature, references will be automatically managed. For example, you might cite a paper like this: \cite{example2024}. All citations will be automatically added to your \texttt{cite.bib} file.

\section{Conclusion}
This template demonstrates basic \LaTeX{} features including mathematical formulas, tables, and lists. The bibliography system is now set up to work seamlessly with Bibby AI's citation management. Feel free to modify it for your needs!

\bibliographystyle{plain}
\bibliography{cite}

\end{document}\documentclass{article}
\usepackage[utf8]{inputenc}
\usepackage{booktabs}
\usepackage{amsmath}
\usepackage{natbib}

\title{My Document}
\author{Your Name}
\date{\today}

\begin{document}

\maketitle

\section{Introduction}
This is a sample document demonstrating various \LaTeX{} features including mathematical equations, tables, and formatting. You can cite papers using the citation search feature, and they will be automatically added to your \texttt{cite.bib} file.

\section{Mathematical Expressions}
Here are some example equations:

\subsection{Quadratic Formula}
The quadratic formula is given by:
\[x = \frac{-b \pm \sqrt{b^2 - 4ac}}{2a}\]

\subsection{Matrix Example}
A sample matrix:
\[
\begin{pmatrix}
a_{11} & a_{12} & a_{13} \\
a_{21} & a_{22} & a_{23} \\
a_{31} & a_{32} & a_{33}
\end{pmatrix}
\]

\section{Tables}
Below is a sample table showing data:

\begin{table}[h]
\centering
\begin{tabular}{@{}lrrr@{}}
\toprule
\textbf{Category} & \textbf{Value 1} & \textbf{Value 2} & \textbf{Total} \\
\midrule
Item A & 10.5 & 20.3 & 30.8 \\
Item B & 15.2 & 25.1 & 40.3 \\
Item C & 22.1 & 12.7 & 34.8 \\
\bottomrule
\end{tabular}
\caption{Sample Data Table}
\label{tab:sample}
\end{table}

\section{Lists}
\subsection{Bullet Points}
\begin{itemize}
  \item First item
  \item Second item
  \item Third item with sub-items:
    \begin{itemize}
      \item Sub-item 1
      \item Sub-item 2
    \end{itemize}
\end{itemize}

\subsection{Numbered List}
\begin{enumerate}
  \item First step
  \item Second step
  \item Third step
\end{enumerate}

\section{Citations and References}
When you use the citation search feature, references will be automatically managed. For example, you might cite a paper like this: \cite{example2024}. All citations will be automatically added to your \texttt{cite.bib} file.

\section{Conclusion}
This template demonstrates basic \LaTeX{} features including mathematical formulas, tables, and lists. The bibliography system is now set up to work seamlessly with Bibby AI's citation management. Feel free to modify it for your needs!

\bibliographystyle{plain}
\bibliography{cite}

\end{document}\documentclass{article}
\usepackage[utf8]{inputenc}
\usepackage{booktabs}
\usepackage{amsmath}
\usepackage{natbib}

\title{My Document}
\author{Your Name}
\date{\today}

\begin{document}

\maketitle

\section{Introduction}
This is a sample document demonstrating various \LaTeX{} features including mathematical equations, tables, and formatting. You can cite papers using the citation search feature, and they will be automatically added to your \texttt{cite.bib} file.

\section{Mathematical Expressions}
Here are some example equations:

\subsection{Quadratic Formula}
The quadratic formula is given by:
\[x = \frac{-b \pm \sqrt{b^2 - 4ac}}{2a}\]

\subsection{Matrix Example}
A sample matrix:
\[
\begin{pmatrix}
a_{11} & a_{12} & a_{13} \\
a_{21} & a_{22} & a_{23} \\
a_{31} & a_{32} & a_{33}
\end{pmatrix}
\]

\section{Tables}
Below is a sample table showing data:

\begin{table}[h]
\centering
\begin{tabular}{@{}lrrr@{}}
\toprule
\textbf{Category} & \textbf{Value 1} & \textbf{Value 2} & \textbf{Total} \\
\midrule
Item A & 10.5 & 20.3 & 30.8 \\
Item B & 15.2 & 25.1 & 40.3 \\
Item C & 22.1 & 12.7 & 34.8 \\
\bottomrule
\end{tabular}
\caption{Sample Data Table}
\label{tab:sample}
\end{table}

\section{Lists}
\subsection{Bullet Points}
\begin{itemize}
  \item First item
  \item Second item
  \item Third item with sub-items:
    \begin{itemize}
      \item Sub-item 1
      \item Sub-item 2
    \end{itemize}
\end{itemize}

\subsection{Numbered List}
\begin{enumerate}
  \item First step
  \item Second step
  \item Third step
\end{enumerate}

\section{Citations and References}
When you use the citation search feature, references will be automatically managed. For example, you might cite a paper like this: \cite{example2024}. All citations will be automatically added to your \texttt{cite.bib} file.

\section{Conclusion}
This template demonstrates basic \LaTeX{} features including mathematical formulas, tables, and lists. The bibliography system is now set up to work seamlessly with Bibby AI's citation management. Feel free to modify it for your needs!

\bibliographystyle{plain}
\bibliography{cite}

\end{document}\documentclass{article}
\usepackage[utf8]{inputenc}
\usepackage{booktabs}
\usepackage{amsmath}
\usepackage{natbib}

\title{My Document}
\author{Your Name}
\date{\today}

\begin{document}

\maketitle

\section{Introduction}
This is a sample document demonstrating various \LaTeX{} features including mathematical equations, tables, and formatting. You can cite papers using the citation search feature, and they will be automatically added to your \texttt{cite.bib} file.

\section{Mathematical Expressions}
Here are some example equations:

\subsection{Quadratic Formula}
The quadratic formula is given by:
\[x = \frac{-b \pm \sqrt{b^2 - 4ac}}{2a}\]

\subsection{Matrix Example}
A sample matrix:
\[
\begin{pmatrix}
a_{11} & a_{12} & a_{13} \\
a_{21} & a_{22} & a_{23} \\
a_{31} & a_{32} & a_{33}
\end{pmatrix}
\]

\section{Tables}
Below is a sample table showing data:

\begin{table}[h]
\centering
\begin{tabular}{@{}lrrr@{}}
\toprule
\textbf{Category} & \textbf{Value 1} & \textbf{Value 2} & \textbf{Total} \\
\midrule
Item A & 10.5 & 20.3 & 30.8 \\
Item B & 15.2 & 25.1 & 40.3 \\
Item C & 22.1 & 12.7 & 34.8 \\
\bottomrule
\end{tabular}
\caption{Sample Data Table}
\label{tab:sample}
\end{table}

\section{Lists}
\subsection{Bullet Points}
\begin{itemize}
  \item First item
  \item Second item
  \item Third item with sub-items:
    \begin{itemize}
      \item Sub-item 1
      \item Sub-item 2
    \end{itemize}
\end{itemize}

\subsection{Numbered List}
\begin{enumerate}
  \item First step
  \item Second step
  \item Third step
\end{enumerate}

\section{Citations and References}
When you use the citation search feature, references will be automatically managed. For example, you might cite a paper like this: \cite{example2024}. All citations will be automatically added to your \texttt{cite.bib} file.

\section{Conclusion}
This template demonstrates basic \LaTeX{} features including mathematical formulas, tables, and lists. The bibliography system is now set up to work seamlessly with Bibby AI's citation management. Feel free to modify it for your needs!

\bibliographystyle{plain}
\bibliography{cite}

\end{document}\documentclass{article}
\usepackage[utf8]{inputenc}
\usepackage{booktabs}
\usepackage{amsmath}
\usepackage{natbib}

\title{My Document}
\author{Your Name}
\date{\today}

\begin{document}

\maketitle

\section{Introduction}
This is a sample document demonstrating various \LaTeX{} features including mathematical equations, tables, and formatting. You can cite papers using the citation search feature, and they will be automatically added to your \texttt{cite.bib} file.

\section{Mathematical Expressions}
Here are some example equations:

\subsection{Quadratic Formula}
The quadratic formula is given by:
\[x = \frac{-b \pm \sqrt{b^2 - 4ac}}{2a}\]

\subsection{Matrix Example}
A sample matrix:
\[
\begin{pmatrix}
a_{11} & a_{12} & a_{13} \\
a_{21} & a_{22} & a_{23} \\
a_{31} & a_{32} & a_{33}
\end{pmatrix}
\]

\section{Tables}
Below is a sample table showing data:

\begin{table}[h]
\centering
\begin{tabular}{@{}lrrr@{}}
\toprule
\textbf{Category} & \textbf{Value 1} & \textbf{Value 2} & \textbf{Total} \\
\midrule
Item A & 10.5 & 20.3 & 30.8 \\
Item B & 15.2 & 25.1 & 40.3 \\
Item C & 22.1 & 12.7 & 34.8 \\
\bottomrule
\end{tabular}
\caption{Sample Data Table}
\label{tab:sample}
\end{table}

\section{Lists}
\subsection{Bullet Points}
\begin{itemize}
  \item First item
  \item Second item
  \item Third item with sub-items:
    \begin{itemize}
      \item Sub-item 1
      \item Sub-item 2
    \end{itemize}
\end{itemize}

\subsection{Numbered List}
\begin{enumerate}
  \item First step
  \item Second step
  \item Third step
\end{enumerate}

\section{Citations and References}
When you use the citation search feature, references will be automatically managed. For example, you might cite a paper like this: \cite{example2024}. All citations will be automatically added to your \texttt{cite.bib} file.

\section{Conclusion}
This template demonstrates basic \LaTeX{} features including mathematical formulas, tables, and lists. The bibliography system is now set up to work seamlessly with Bibby AI's citation management. Feel free to modify it for your needs!

\bibliographystyle{plain}
\bibliography{cite}

\end{document}\documentclass{article}
\usepackage[utf8]{inputenc}
\usepackage{booktabs}
\usepackage{amsmath}
\usepackage{natbib}

\title{My Document}
\author{Your Name}
\date{\today}

\begin{document}

\maketitle

\section{Introduction}
This is a sample document demonstrating various \LaTeX{} features including mathematical equations, tables, and formatting. You can cite papers using the citation search feature, and they will be automatically added to your \texttt{cite.bib} file.

\section{Mathematical Expressions}
Here are some example equations:

\subsection{Quadratic Formula}
The quadratic formula is given by:
\[x = \frac{-b \pm \sqrt{b^2 - 4ac}}{2a}\]

\subsection{Matrix Example}
A sample matrix:
\[
\begin{pmatrix}
a_{11} & a_{12} & a_{13} \\
a_{21} & a_{22} & a_{23} \\
a_{31} & a_{32} & a_{33}
\end{pmatrix}
\]

\section{Tables}
Below is a sample table showing data:

\begin{table}[h]
\centering
\begin{tabular}{@{}lrrr@{}}
\toprule
\textbf{Category} & \textbf{Value 1} & \textbf{Value 2} & \textbf{Total} \\
\midrule
Item A & 10.5 & 20.3 & 30.8 \\
Item B & 15.2 & 25.1 & 40.3 \\
Item C & 22.1 & 12.7 & 34.8 \\
\bottomrule
\end{tabular}
\caption{Sample Data Table}
\label{tab:sample}
\end{table}

\section{Lists}
\subsection{Bullet Points}
\begin{itemize}
  \item First item
  \item Second item
  \item Third item with sub-items:
    \begin{itemize}
      \item Sub-item 1
      \item Sub-item 2
    \end{itemize}
\end{itemize}

\subsection{Numbered List}
\begin{enumerate}
  \item First step
  \item Second step
  \item Third step
\end{enumerate}

\section{Citations and References}
When you use the citation search feature, references will be automatically managed. For example, you might cite a paper like this: \cite{example2024}. All citations will be automatically added to your \texttt{cite.bib} file.

\section{Conclusion}
This template demonstrates basic \LaTeX{} features including mathematical formulas, tables, and lists. The bibliography system is now set up to work seamlessly with Bibby AI's citation management. Feel free to modify it for your needs!

\bibliographystyle{plain}
\bibliography{cite}

\end{document}\documentclass{article}
\usepackage[utf8]{inputenc}
\usepackage{booktabs}
\usepackage{amsmath}
\usepackage{natbib}

\title{My Document}
\author{Your Name}
\date{\today}

\begin{document}

\maketitle

\section{Introduction}
This is a sample document demonstrating various \LaTeX{} features including mathematical equations, tables, and formatting. You can cite papers using the citation search feature, and they will be automatically added to your \texttt{cite.bib} file.

\section{Mathematical Expressions}
Here are some example equations:

\subsection{Quadratic Formula}
The quadratic formula is given by:
\[x = \frac{-b \pm \sqrt{b^2 - 4ac}}{2a}\]

\subsection{Matrix Example}
A sample matrix:
\[
\begin{pmatrix}
a_{11} & a_{12} & a_{13} \\
a_{21} & a_{22} & a_{23} \\
a_{31} & a_{32} & a_{33}
\end{pmatrix}
\]

\section{Tables}
Below is a sample table showing data:

\begin{table}[h]
\centering
\begin{tabular}{@{}lrrr@{}}
\toprule
\textbf{Category} & \textbf{Value 1} & \textbf{Value 2} & \textbf{Total} \\
\midrule
Item A & 10.5 & 20.3 & 30.8 \\
Item B & 15.2 & 25.1 & 40.3 \\
Item C & 22.1 & 12.7 & 34.8 \\
\bottomrule
\end{tabular}
\caption{Sample Data Table}
\label{tab:sample}
\end{table}

\section{Lists}
\subsection{Bullet Points}
\begin{itemize}
  \item First item
  \item Second item
  \item Third item with sub-items:
    \begin{itemize}
      \item Sub-item 1
      \item Sub-item 2
    \end{itemize}
\end{itemize}

\subsection{Numbered List}
\begin{enumerate}
  \item First step
  \item Second step
  \item Third step
\end{enumerate}

\section{Citations and References}
When you use the citation search feature, references will be automatically managed. For example, you might cite a paper like this: \cite{example2024}. All citations will be automatically added to your \texttt{cite.bib} file.

\section{Conclusion}
This template demonstrates basic \LaTeX{} features including mathematical formulas, tables, and lists. The bibliography system is now set up to work seamlessly with Bibby AI's citation management. Feel free to modify it for your needs!

\bibliographystyle{plain}
\bibliography{cite}

\end{document}\documentclass{article}
\usepackage[utf8]{inputenc}
\usepackage{booktabs}
\usepackage{amsmath}
\usepackage{natbib}

\title{My Document}
\author{Your Name}
\date{\today}

\begin{document}

\maketitle

\section{Introduction}
This is a sample document demonstrating various \LaTeX{} features including mathematical equations, tables, and formatting. You can cite papers using the citation search feature, and they will be automatically added to your \texttt{cite.bib} file.

\section{Mathematical Expressions}
Here are some example equations:

\subsection{Quadratic Formula}
The quadratic formula is given by:
\[x = \frac{-b \pm \sqrt{b^2 - 4ac}}{2a}\]

\subsection{Matrix Example}
A sample matrix:
\[
\begin{pmatrix}
a_{11} & a_{12} & a_{13} \\
a_{21} & a_{22} & a_{23} \\
a_{31} & a_{32} & a_{33}
\end{pmatrix}
\]

\section{Tables}
Below is a sample table showing data:

\begin{table}[h]
\centering
\begin{tabular}{@{}lrrr@{}}
\toprule
\textbf{Category} & \textbf{Value 1} & \textbf{Value 2} & \textbf{Total} \\
\midrule
Item A & 10.5 & 20.3 & 30.8 \\
Item B & 15.2 & 25.1 & 40.3 \\
Item C & 22.1 & 12.7 & 34.8 \\
\bottomrule
\end{tabular}
\caption{Sample Data Table}
\label{tab:sample}
\end{table}

\section{Lists}
\subsection{Bullet Points}
\begin{itemize}
  \item First item
  \item Second item
  \item Third item with sub-items:
    \begin{itemize}
      \item Sub-item 1
      \item Sub-item 2
    \end{itemize}
\end{itemize}

\subsection{Numbered List}
\begin{enumerate}
  \item First step
  \item Second step
  \item Third step
\end{enumerate}

\section{Citations and References}
When you use the citation search feature, references will be automatically managed. For example, you might cite a paper like this: \cite{example2024}. All citations will be automatically added to your \texttt{cite.bib} file.

\section{Conclusion}
This template demonstrates basic \LaTeX{} features including mathematical formulas, tables, and lists. The bibliography system is now set up to work seamlessly with Bibby AI's citation management. Feel free to modify it for your needs!

\bibliographystyle{plain}
\bibliography{cite}

\end{document}\documentclass{article}
\usepackage[utf8]{inputenc}
\usepackage{booktabs}
\usepackage{amsmath}
\usepackage{natbib}

\title{My Document}
\author{Your Name}
\date{\today}

\begin{document}

\maketitle

\section{Introduction}
This is a sample document demonstrating various \LaTeX{} features including mathematical equations, tables, and formatting. You can cite papers using the citation search feature, and they will be automatically added to your \texttt{cite.bib} file.

\section{Mathematical Expressions}
Here are some example equations:

\subsection{Quadratic Formula}
The quadratic formula is given by:
\[x = \frac{-b \pm \sqrt{b^2 - 4ac}}{2a}\]

\subsection{Matrix Example}
A sample matrix:
\[
\begin{pmatrix}
a_{11} & a_{12} & a_{13} \\
a_{21} & a_{22} & a_{23} \\
a_{31} & a_{32} & a_{33}
\end{pmatrix}
\]

\section{Tables}
Below is a sample table showing data:

\begin{table}[h]
\centering
\begin{tabular}{@{}lrrr@{}}
\toprule
\textbf{Category} & \textbf{Value 1} & \textbf{Value 2} & \textbf{Total} \\
\midrule
Item A & 10.5 & 20.3 & 30.8 \\
Item B & 15.2 & 25.1 & 40.3 \\
Item C & 22.1 & 12.7 & 34.8 \\
\bottomrule
\end{tabular}
\caption{Sample Data Table}
\label{tab:sample}
\end{table}

\section{Lists}
\subsection{Bullet Points}
\begin{itemize}
  \item First item
  \item Second item
  \item Third item with sub-items:
    \begin{itemize}
      \item Sub-item 1
      \item Sub-item 2
    \end{itemize}
\end{itemize}

\subsection{Numbered List}
\begin{enumerate}
  \item First step
  \item Second step
  \item Third step
\end{enumerate}

\section{Citations and References}
When you use the citation search feature, references will be automatically managed. For example, you might cite a paper like this: \cite{example2024}. All citations will be automatically added to your \texttt{cite.bib} file.

\section{Conclusion}
This template demonstrates basic \LaTeX{} features including mathematical formulas, tables, and lists. The bibliography system is now set up to work seamlessly with Bibby AI's citation management. Feel free to modify it for your needs!

\bibliographystyle{plain}
\bibliography{cite}

\end{document}\documentclass{article}
\usepackage[utf8]{inputenc}
\usepackage{booktabs}
\usepackage{amsmath}
\usepackage{natbib}

\title{My Document}
\author{Your Name}
\date{\today}

\begin{document}

\maketitle

\section{Introduction}
This is a sample document demonstrating various \LaTeX{} features including mathematical equations, tables, and formatting. You can cite papers using the citation search feature, and they will be automatically added to your \texttt{cite.bib} file.

\section{Mathematical Expressions}
Here are some example equations:

\subsection{Quadratic Formula}
The quadratic formula is given by:
\[x = \frac{-b \pm \sqrt{b^2 - 4ac}}{2a}\]

\subsection{Matrix Example}
A sample matrix:
\[
\begin{pmatrix}
a_{11} & a_{12} & a_{13} \\
a_{21} & a_{22} & a_{23} \\
a_{31} & a_{32} & a_{33}
\end{pmatrix}
\]

\section{Tables}
Below is a sample table showing data:

\begin{table}[h]
\centering
\begin{tabular}{@{}lrrr@{}}
\toprule
\textbf{Category} & \textbf{Value 1} & \textbf{Value 2} & \textbf{Total} \\
\midrule
Item A & 10.5 & 20.3 & 30.8 \\
Item B & 15.2 & 25.1 & 40.3 \\
Item C & 22.1 & 12.7 & 34.8 \\
\bottomrule
\end{tabular}
\caption{Sample Data Table}
\label{tab:sample}
\end{table}

\section{Lists}
\subsection{Bullet Points}
\begin{itemize}
  \item First item
  \item Second item
  \item Third item with sub-items:
    \begin{itemize}
      \item Sub-item 1
      \item Sub-item 2
    \end{itemize}
\end{itemize}

\subsection{Numbered List}
\begin{enumerate}
  \item First step
  \item Second step
  \item Third step
\end{enumerate}

\section{Citations and References}
When you use the citation search feature, references will be automatically managed. For example, you might cite a paper like this: \cite{example2024}. All citations will be automatically added to your \texttt{cite.bib} file.

\section{Conclusion}
This template demonstrates basic \LaTeX{} features including mathematical formulas, tables, and lists. The bibliography system is now set up to work seamlessly with Bibby AI's citation management. Feel free to modify it for your needs!

\bibliographystyle{plain}
\bibliography{cite}

\end{document}\documentclass{article}
\usepackage[utf8]{inputenc}
\usepackage{booktabs}
\usepackage{amsmath}
\usepackage{natbib}

\title{My Document}
\author{Your Name}
\date{\today}

\begin{document}

\maketitle

\section{Introduction}
This is a sample document demonstrating various \LaTeX{} features including mathematical equations, tables, and formatting. You can cite papers using the citation search feature, and they will be automatically added to your \texttt{cite.bib} file.

\section{Mathematical Expressions}
Here are some example equations:

\subsection{Quadratic Formula}
The quadratic formula is given by:
\[x = \frac{-b \pm \sqrt{b^2 - 4ac}}{2a}\]

\subsection{Matrix Example}
A sample matrix:
\[
\begin{pmatrix}
a_{11} & a_{12} & a_{13} \\
a_{21} & a_{22} & a_{23} \\
a_{31} & a_{32} & a_{33}
\end{pmatrix}
\]

\section{Tables}
Below is a sample table showing data:

\begin{table}[h]
\centering
\begin{tabular}{@{}lrrr@{}}
\toprule
\textbf{Category} & \textbf{Value 1} & \textbf{Value 2} & \textbf{Total} \\
\midrule
Item A & 10.5 & 20.3 & 30.8 \\
Item B & 15.2 & 25.1 & 40.3 \\
Item C & 22.1 & 12.7 & 34.8 \\
\bottomrule
\end{tabular}
\caption{Sample Data Table}
\label{tab:sample}
\end{table}

\section{Lists}
\subsection{Bullet Points}
\begin{itemize}
  \item First item
  \item Second item
  \item Third item with sub-items:
    \begin{itemize}
      \item Sub-item 1
      \item Sub-item 2
    \end{itemize}
\end{itemize}

\subsection{Numbered List}
\begin{enumerate}
  \item First step
  \item Second step
  \item Third step
\end{enumerate}

\section{Citations and References}
When you use the citation search feature, references will be automatically managed. For example, you might cite a paper like this: \cite{example2024}. All citations will be automatically added to your \texttt{cite.bib} file.

\section{Conclusion}
This template demonstrates basic \LaTeX{} features including mathematical formulas, tables, and lists. The bibliography system is now set up to work seamlessly with Bibby AI's citation management. Feel free to modify it for your needs!

\bibliographystyle{plain}
\bibliography{cite}

\end{document}\documentclass{article}
\usepackage[utf8]{inputenc}
\usepackage{booktabs}
\usepackage{amsmath}
\usepackage{natbib}

\title{My Document}
\author{Your Name}
\date{\today}

\begin{document}

\maketitle

\section{Introduction}
This is a sample document demonstrating various \LaTeX{} features including mathematical equations, tables, and formatting. You can cite papers using the citation search feature, and they will be automatically added to your \texttt{cite.bib} file.

\section{Mathematical Expressions}
Here are some example equations:

\subsection{Quadratic Formula}
The quadratic formula is given by:
\[x = \frac{-b \pm \sqrt{b^2 - 4ac}}{2a}\]

\subsection{Matrix Example}
A sample matrix:
\[
\begin{pmatrix}
a_{11} & a_{12} & a_{13} \\
a_{21} & a_{22} & a_{23} \\
a_{31} & a_{32} & a_{33}
\end{pmatrix}
\]

\section{Tables}
Below is a sample table showing data:

\begin{table}[h]
\centering
\begin{tabular}{@{}lrrr@{}}
\toprule
\textbf{Category} & \textbf{Value 1} & \textbf{Value 2} & \textbf{Total} \\
\midrule
Item A & 10.5 & 20.3 & 30.8 \\
Item B & 15.2 & 25.1 & 40.3 \\
Item C & 22.1 & 12.7 & 34.8 \\
\bottomrule
\end{tabular}
\caption{Sample Data Table}
\label{tab:sample}
\end{table}

\section{Lists}
\subsection{Bullet Points}
\begin{itemize}
  \item First item
  \item Second item
  \item Third item with sub-items:
    \begin{itemize}
      \item Sub-item 1
      \item Sub-item 2
    \end{itemize}
\end{itemize}

\subsection{Numbered List}
\begin{enumerate}
  \item First step
  \item Second step
  \item Third step
\end{enumerate}

\section{Citations and References}
When you use the citation search feature, references will be automatically managed. For example, you might cite a paper like this: \cite{example2024}. All citations will be automatically added to your \texttt{cite.bib} file.

\section{Conclusion}
This template demonstrates basic \LaTeX{} features including mathematical formulas, tables, and lists. The bibliography system is now set up to work seamlessly with Bibby AI's citation management. Feel free to modify it for your needs!

\bibliographystyle{plain}
\bibliography{cite}

\end{document}\documentclass{article}
\usepackage[utf8]{inputenc}
\usepackage{booktabs}
\usepackage{amsmath}
\usepackage{natbib}

\title{My Document}
\author{Your Name}
\date{\today}

\begin{document}

\maketitle

\section{Introduction}
This is a sample document demonstrating various \LaTeX{} features including mathematical equations, tables, and formatting. You can cite papers using the citation search feature, and they will be automatically added to your \texttt{cite.bib} file.

\section{Mathematical Expressions}
Here are some example equations:

\subsection{Quadratic Formula}
The quadratic formula is given by:
\[x = \frac{-b \pm \sqrt{b^2 - 4ac}}{2a}\]

\subsection{Matrix Example}
A sample matrix:
\[
\begin{pmatrix}
a_{11} & a_{12} & a_{13} \\
a_{21} & a_{22} & a_{23} \\
a_{31} & a_{32} & a_{33}
\end{pmatrix}
\]

\section{Tables}
Below is a sample table showing data:

\begin{table}[h]
\centering
\begin{tabular}{@{}lrrr@{}}
\toprule
\textbf{Category} & \textbf{Value 1} & \textbf{Value 2} & \textbf{Total} \\
\midrule
Item A & 10.5 & 20.3 & 30.8 \\
Item B & 15.2 & 25.1 & 40.3 \\
Item C & 22.1 & 12.7 & 34.8 \\
\bottomrule
\end{tabular}
\caption{Sample Data Table}
\label{tab:sample}
\end{table}

\section{Lists}
\subsection{Bullet Points}
\begin{itemize}
  \item First item
  \item Second item
  \item Third item with sub-items:
    \begin{itemize}
      \item Sub-item 1
      \item Sub-item 2
    \end{itemize}
\end{itemize}

\subsection{Numbered List}
\begin{enumerate}
  \item First step
  \item Second step
  \item Third step
\end{enumerate}

\section{Citations and References}
When you use the citation search feature, references will be automatically managed. For example, you might cite a paper like this: \cite{example2024}. All citations will be automatically added to your \texttt{cite.bib} file.

\section{Conclusion}
This template demonstrates basic \LaTeX{} features including mathematical formulas, tables, and lists. The bibliography system is now set up to work seamlessly with Bibby AI's citation management. Feel free to modify it for your needs!

\bibliographystyle{plain}
\bibliography{cite}

\end{document}\documentclass{article}
\usepackage[utf8]{inputenc}
\usepackage{booktabs}
\usepackage{amsmath}
\usepackage{natbib}

\title{My Document}
\author{Your Name}
\date{\today}

\begin{document}

\maketitle

\section{Introduction}
This is a sample document demonstrating various \LaTeX{} features including mathematical equations, tables, and formatting. You can cite papers using the citation search feature, and they will be automatically added to your \texttt{cite.bib} file.

\section{Mathematical Expressions}
Here are some example equations:

\subsection{Quadratic Formula}
The quadratic formula is given by:
\[x = \frac{-b \pm \sqrt{b^2 - 4ac}}{2a}\]

\subsection{Matrix Example}
A sample matrix:
\[
\begin{pmatrix}
a_{11} & a_{12} & a_{13} \\
a_{21} & a_{22} & a_{23} \\
a_{31} & a_{32} & a_{33}
\end{pmatrix}
\]

\section{Tables}
Below is a sample table showing data:

\begin{table}[h]
\centering
\begin{tabular}{@{}lrrr@{}}
\toprule
\textbf{Category} & \textbf{Value 1} & \textbf{Value 2} & \textbf{Total} \\
\midrule
Item A & 10.5 & 20.3 & 30.8 \\
Item B & 15.2 & 25.1 & 40.3 \\
Item C & 22.1 & 12.7 & 34.8 \\
\bottomrule
\end{tabular}
\caption{Sample Data Table}
\label{tab:sample}
\end{table}

\section{Lists}
\subsection{Bullet Points}
\begin{itemize}
  \item First item
  \item Second item
  \item Third item with sub-items:
    \begin{itemize}
      \item Sub-item 1
      \item Sub-item 2
    \end{itemize}
\end{itemize}

\subsection{Numbered List}
\begin{enumerate}
  \item First step
  \item Second step
  \item Third step
\end{enumerate}

\section{Citations and References}
When you use the citation search feature, references will be automatically managed. For example, you might cite a paper like this: \cite{example2024}. All citations will be automatically added to your \texttt{cite.bib} file.

\section{Conclusion}
This template demonstrates basic \LaTeX{} features including mathematical formulas, tables, and lists. The bibliography system is now set up to work seamlessly with Bibby AI's citation management. Feel free to modify it for your needs!

\bibliographystyle{plain}
\bibliography{cite}

\end{document}\documentclass{article}
\usepackage[utf8]{inputenc}
\usepackage{booktabs}
\usepackage{amsmath}
\usepackage{natbib}

\title{My Document}
\author{Your Name}
\date{\today}

\begin{document}

\maketitle

\section{Introduction}
This is a sample document demonstrating various \LaTeX{} features including mathematical equations, tables, and formatting. You can cite papers using the citation search feature, and they will be automatically added to your \texttt{cite.bib} file.

\section{Mathematical Expressions}
Here are some example equations:

\subsection{Quadratic Formula}
The quadratic formula is given by:
\[x = \frac{-b \pm \sqrt{b^2 - 4ac}}{2a}\]

\subsection{Matrix Example}
A sample matrix:
\[
\begin{pmatrix}
a_{11} & a_{12} & a_{13} \\
a_{21} & a_{22} & a_{23} \\
a_{31} & a_{32} & a_{33}
\end{pmatrix}
\]

\section{Tables}
Below is a sample table showing data:

\begin{table}[h]
\centering
\begin{tabular}{@{}lrrr@{}}
\toprule
\textbf{Category} & \textbf{Value 1} & \textbf{Value 2} & \textbf{Total} \\
\midrule
Item A & 10.5 & 20.3 & 30.8 \\
Item B & 15.2 & 25.1 & 40.3 \\
Item C & 22.1 & 12.7 & 34.8 \\
\bottomrule
\end{tabular}
\caption{Sample Data Table}
\label{tab:sample}
\end{table}

\section{Lists}
\subsection{Bullet Points}
\begin{itemize}
  \item First item
  \item Second item
  \item Third item with sub-items:
    \begin{itemize}
      \item Sub-item 1
      \item Sub-item 2
    \end{itemize}
\end{itemize}

\subsection{Numbered List}
\begin{enumerate}
  \item First step
  \item Second step
  \item Third step
\end{enumerate}

\section{Citations and References}
When you use the citation search feature, references will be automatically managed. For example, you might cite a paper like this: \cite{example2024}. All citations will be automatically added to your \texttt{cite.bib} file.

\section{Conclusion}
This template demonstrates basic \LaTeX{} features including mathematical formulas, tables, and lists. The bibliography system is now set up to work seamlessly with Bibby AI's citation management. Feel free to modify it for your needs!

\bibliographystyle{plain}
\bibliography{cite}

\end{document}\documentclass{article}
\usepackage[utf8]{inputenc}
\usepackage{booktabs}
\usepackage{amsmath}
\usepackage{natbib}

\title{My Document}
\author{Your Name}
\date{\today}

\begin{document}

\maketitle

\section{Introduction}
This is a sample document demonstrating various \LaTeX{} features including mathematical equations, tables, and formatting. You can cite papers using the citation search feature, and they will be automatically added to your \texttt{cite.bib} file.

\section{Mathematical Expressions}
Here are some example equations:

\subsection{Quadratic Formula}
The quadratic formula is given by:
\[x = \frac{-b \pm \sqrt{b^2 - 4ac}}{2a}\]

\subsection{Matrix Example}
A sample matrix:
\[
\begin{pmatrix}
a_{11} & a_{12} & a_{13} \\
a_{21} & a_{22} & a_{23} \\
a_{31} & a_{32} & a_{33}
\end{pmatrix}
\]

\section{Tables}
Below is a sample table showing data:

\begin{table}[h]
\centering
\begin{tabular}{@{}lrrr@{}}
\toprule
\textbf{Category} & \textbf{Value 1} & \textbf{Value 2} & \textbf{Total} \\
\midrule
Item A & 10.5 & 20.3 & 30.8 \\
Item B & 15.2 & 25.1 & 40.3 \\
Item C & 22.1 & 12.7 & 34.8 \\
\bottomrule
\end{tabular}
\caption{Sample Data Table}
\label{tab:sample}
\end{table}

\section{Lists}
\subsection{Bullet Points}
\begin{itemize}
  \item First item
  \item Second item
  \item Third item with sub-items:
    \begin{itemize}
      \item Sub-item 1
      \item Sub-item 2
    \end{itemize}
\end{itemize}

\subsection{Numbered List}
\begin{enumerate}
  \item First step
  \item Second step
  \item Third step
\end{enumerate}

\section{Citations and References}
When you use the citation search feature, references will be automatically managed. For example, you might cite a paper like this: \cite{example2024}. All citations will be automatically added to your \texttt{cite.bib} file.

\section{Conclusion}
This template demonstrates basic \LaTeX{} features including mathematical formulas, tables, and lists. The bibliography system is now set up to work seamlessly with Bibby AI's citation management. Feel free to modify it for your needs!

\bibliographystyle{plain}
\bibliography{cite}

\end{document}\documentclass{article}
\usepackage[utf8]{inputenc}
\usepackage{booktabs}
\usepackage{amsmath}
\usepackage{natbib}

\title{My Document}
\author{Your Name}
\date{\today}

\begin{document}

\maketitle

\section{Introduction}
This is a sample document demonstrating various \LaTeX{} features including mathematical equations, tables, and formatting. You can cite papers using the citation search feature, and they will be automatically added to your \texttt{cite.bib} file.

\section{Mathematical Expressions}
Here are some example equations:

\subsection{Quadratic Formula}
The quadratic formula is given by:
\[x = \frac{-b \pm \sqrt{b^2 - 4ac}}{2a}\]

\subsection{Matrix Example}
A sample matrix:
\[
\begin{pmatrix}
a_{11} & a_{12} & a_{13} \\
a_{21} & a_{22} & a_{23} \\
a_{31} & a_{32} & a_{33}
\end{pmatrix}
\]

\section{Tables}
Below is a sample table showing data:

\begin{table}[h]
\centering
\begin{tabular}{@{}lrrr@{}}
\toprule
\textbf{Category} & \textbf{Value 1} & \textbf{Value 2} & \textbf{Total} \\
\midrule
Item A & 10.5 & 20.3 & 30.8 \\
Item B & 15.2 & 25.1 & 40.3 \\
Item C & 22.1 & 12.7 & 34.8 \\
\bottomrule
\end{tabular}
\caption{Sample Data Table}
\label{tab:sample}
\end{table}

\section{Lists}
\subsection{Bullet Points}
\begin{itemize}
  \item First item
  \item Second item
  \item Third item with sub-items:
    \begin{itemize}
      \item Sub-item 1
      \item Sub-item 2
    \end{itemize}
\end{itemize}

\subsection{Numbered List}
\begin{enumerate}
  \item First step
  \item Second step
  \item Third step
\end{enumerate}

\section{Citations and References}
When you use the citation search feature, references will be automatically managed. For example, you might cite a paper like this: \cite{example2024}. All citations will be automatically added to your \texttt{cite.bib} file.

\section{Conclusion}
This template demonstrates basic \LaTeX{} features including mathematical formulas, tables, and lists. The bibliography system is now set up to work seamlessly with Bibby AI's citation management. Feel free to modify it for your needs!

\bibliographystyle{plain}
\bibliography{cite}

\end{document}\documentclass{article}
\usepackage[utf8]{inputenc}
\usepackage{booktabs}
\usepackage{amsmath}
\usepackage{natbib}

\title{My Document}
\author{Your Name}
\date{\today}

\begin{document}

\maketitle

\section{Introduction}
This is a sample document demonstrating various \LaTeX{} features including mathematical equations, tables, and formatting. You can cite papers using the citation search feature, and they will be automatically added to your \texttt{cite.bib} file.

\section{Mathematical Expressions}
Here are some example equations:

\subsection{Quadratic Formula}
The quadratic formula is given by:
\[x = \frac{-b \pm \sqrt{b^2 - 4ac}}{2a}\]

\subsection{Matrix Example}
A sample matrix:
\[
\begin{pmatrix}
a_{11} & a_{12} & a_{13} \\
a_{21} & a_{22} & a_{23} \\
a_{31} & a_{32} & a_{33}
\end{pmatrix}
\]

\section{Tables}
Below is a sample table showing data:

\begin{table}[h]
\centering
\begin{tabular}{@{}lrrr@{}}
\toprule
\textbf{Category} & \textbf{Value 1} & \textbf{Value 2} & \textbf{Total} \\
\midrule
Item A & 10.5 & 20.3 & 30.8 \\
Item B & 15.2 & 25.1 & 40.3 \\
Item C & 22.1 & 12.7 & 34.8 \\
\bottomrule
\end{tabular}
\caption{Sample Data Table}
\label{tab:sample}
\end{table}

\section{Lists}
\subsection{Bullet Points}
\begin{itemize}
  \item First item
  \item Second item
  \item Third item with sub-items:
    \begin{itemize}
      \item Sub-item 1
      \item Sub-item 2
    \end{itemize}
\end{itemize}

\subsection{Numbered List}
\begin{enumerate}
  \item First step
  \item Second step
  \item Third step
\end{enumerate}

\section{Citations and References}
When you use the citation search feature, references will be automatically managed. For example, you might cite a paper like this: \cite{example2024}. All citations will be automatically added to your \texttt{cite.bib} file.

\section{Conclusion}
This template demonstrates basic \LaTeX{} features including mathematical formulas, tables, and lists. The bibliography system is now set up to work seamlessly with Bibby AI's citation management. Feel free to modify it for your needs!

\bibliographystyle{plain}
\bibliography{cite}

\end{document}\documentclass{article}
\usepackage[utf8]{inputenc}
\usepackage{booktabs}
\usepackage{amsmath}
\usepackage{natbib}

\title{My Document}
\author{Your Name}
\date{\today}

\begin{document}

\maketitle

\section{Introduction}
This is a sample document demonstrating various \LaTeX{} features including mathematical equations, tables, and formatting. You can cite papers using the citation search feature, and they will be automatically added to your \texttt{cite.bib} file.

\section{Mathematical Expressions}
Here are some example equations:

\subsection{Quadratic Formula}
The quadratic formula is given by:
\[x = \frac{-b \pm \sqrt{b^2 - 4ac}}{2a}\]

\subsection{Matrix Example}
A sample matrix:
\[
\begin{pmatrix}
a_{11} & a_{12} & a_{13} \\
a_{21} & a_{22} & a_{23} \\
a_{31} & a_{32} & a_{33}
\end{pmatrix}
\]

\section{Tables}
Below is a sample table showing data:

\begin{table}[h]
\centering
\begin{tabular}{@{}lrrr@{}}
\toprule
\textbf{Category} & \textbf{Value 1} & \textbf{Value 2} & \textbf{Total} \\
\midrule
Item A & 10.5 & 20.3 & 30.8 \\
Item B & 15.2 & 25.1 & 40.3 \\
Item C & 22.1 & 12.7 & 34.8 \\
\bottomrule
\end{tabular}
\caption{Sample Data Table}
\label{tab:sample}
\end{table}

\section{Lists}
\subsection{Bullet Points}
\begin{itemize}
  \item First item
  \item Second item
  \item Third item with sub-items:
    \begin{itemize}
      \item Sub-item 1
      \item Sub-item 2
    \end{itemize}
\end{itemize}

\subsection{Numbered List}
\begin{enumerate}
  \item First step
  \item Second step
  \item Third step
\end{enumerate}

\section{Citations and References}
When you use the citation search feature, references will be automatically managed. For example, you might cite a paper like this: \cite{example2024}. All citations will be automatically added to your \texttt{cite.bib} file.

\section{Conclusion}
This template demonstrates basic \LaTeX{} features including mathematical formulas, tables, and lists. The bibliography system is now set up to work seamlessly with Bibby AI's citation management. Feel free to modify it for your needs!

\bibliographystyle{plain}
\bibliography{cite}

\end{document}\documentclass{article}
\usepackage[utf8]{inputenc}
\usepackage{booktabs}
\usepackage{amsmath}
\usepackage{natbib}

\title{My Document}
\author{Your Name}
\date{\today}

\begin{document}

\maketitle

\section{Introduction}
This is a sample document demonstrating various \LaTeX{} features including mathematical equations, tables, and formatting. You can cite papers using the citation search feature, and they will be automatically added to your \texttt{cite.bib} file.

\section{Mathematical Expressions}
Here are some example equations:

\subsection{Quadratic Formula}
The quadratic formula is given by:
\[x = \frac{-b \pm \sqrt{b^2 - 4ac}}{2a}\]

\subsection{Matrix Example}
A sample matrix:
\[
\begin{pmatrix}
a_{11} & a_{12} & a_{13} \\
a_{21} & a_{22} & a_{23} \\
a_{31} & a_{32} & a_{33}
\end{pmatrix}
\]

\section{Tables}
Below is a sample table showing data:

\begin{table}[h]
\centering
\begin{tabular}{@{}lrrr@{}}
\toprule
\textbf{Category} & \textbf{Value 1} & \textbf{Value 2} & \textbf{Total} \\
\midrule
Item A & 10.5 & 20.3 & 30.8 \\
Item B & 15.2 & 25.1 & 40.3 \\
Item C & 22.1 & 12.7 & 34.8 \\
\bottomrule
\end{tabular}
\caption{Sample Data Table}
\label{tab:sample}
\end{table}

\section{Lists}
\subsection{Bullet Points}
\begin{itemize}
  \item First item
  \item Second item
  \item Third item with sub-items:
    \begin{itemize}
      \item Sub-item 1
      \item Sub-item 2
    \end{itemize}
\end{itemize}

\subsection{Numbered List}
\begin{enumerate}
  \item First step
  \item Second step
  \item Third step
\end{enumerate}

\section{Citations and References}
When you use the citation search feature, references will be automatically managed. For example, you might cite a paper like this: \cite{example2024}. All citations will be automatically added to your \texttt{cite.bib} file.

\section{Conclusion}
This template demonstrates basic \LaTeX{} features including mathematical formulas, tables, and lists. The bibliography system is now set up to work seamlessly with Bibby AI's citation management. Feel free to modify it for your needs!

\bibliographystyle{plain}
\bibliography{cite}

\end{document}\documentclass{article}
\usepackage[utf8]{inputenc}
\usepackage{booktabs}
\usepackage{amsmath}
\usepackage{natbib}

\title{My Document}
\author{Your Name}
\date{\today}

\begin{document}

\maketitle

\section{Introduction}
This is a sample document demonstrating various \LaTeX{} features including mathematical equations, tables, and formatting. You can cite papers using the citation search feature, and they will be automatically added to your \texttt{cite.bib} file.

\section{Mathematical Expressions}
Here are some example equations:

\subsection{Quadratic Formula}
The quadratic formula is given by:
\[x = \frac{-b \pm \sqrt{b^2 - 4ac}}{2a}\]

\subsection{Matrix Example}
A sample matrix:
\[
\begin{pmatrix}
a_{11} & a_{12} & a_{13} \\
a_{21} & a_{22} & a_{23} \\
a_{31} & a_{32} & a_{33}
\end{pmatrix}
\]

\section{Tables}
Below is a sample table showing data:

\begin{table}[h]
\centering
\begin{tabular}{@{}lrrr@{}}
\toprule
\textbf{Category} & \textbf{Value 1} & \textbf{Value 2} & \textbf{Total} \\
\midrule
Item A & 10.5 & 20.3 & 30.8 \\
Item B & 15.2 & 25.1 & 40.3 \\
Item C & 22.1 & 12.7 & 34.8 \\
\bottomrule
\end{tabular}
\caption{Sample Data Table}
\label{tab:sample}
\end{table}

\section{Lists}
\subsection{Bullet Points}
\begin{itemize}
  \item First item
  \item Second item
  \item Third item with sub-items:
    \begin{itemize}
      \item Sub-item 1
      \item Sub-item 2
    \end{itemize}
\end{itemize}

\subsection{Numbered List}
\begin{enumerate}
  \item First step
  \item Second step
  \item Third step
\end{enumerate}

\section{Citations and References}
When you use the citation search feature, references will be automatically managed. For example, you might cite a paper like this: \cite{example2024}. All citations will be automatically added to your \texttt{cite.bib} file.

\section{Conclusion}
This template demonstrates basic \LaTeX{} features including mathematical formulas, tables, and lists. The bibliography system is now set up to work seamlessly with Bibby AI's citation management. Feel free to modify it for your needs!

\bibliographystyle{plain}
\bibliography{cite}

\end{document}\documentclass{article}
\usepackage[utf8]{inputenc}
\usepackage{booktabs}
\usepackage{amsmath}
\usepackage{natbib}

\title{My Document}
\author{Your Name}
\date{\today}

\begin{document}

\maketitle

\section{Introduction}
This is a sample document demonstrating various \LaTeX{} features including mathematical equations, tables, and formatting. You can cite papers using the citation search feature, and they will be automatically added to your \texttt{cite.bib} file.

\section{Mathematical Expressions}
Here are some example equations:

\subsection{Quadratic Formula}
The quadratic formula is given by:
\[x = \frac{-b \pm \sqrt{b^2 - 4ac}}{2a}\]

\subsection{Matrix Example}
A sample matrix:
\[
\begin{pmatrix}
a_{11} & a_{12} & a_{13} \\
a_{21} & a_{22} & a_{23} \\
a_{31} & a_{32} & a_{33}
\end{pmatrix}
\]

\section{Tables}
Below is a sample table showing data:

\begin{table}[h]
\centering
\begin{tabular}{@{}lrrr@{}}
\toprule
\textbf{Category} & \textbf{Value 1} & \textbf{Value 2} & \textbf{Total} \\
\midrule
Item A & 10.5 & 20.3 & 30.8 \\
Item B & 15.2 & 25.1 & 40.3 \\
Item C & 22.1 & 12.7 & 34.8 \\
\bottomrule
\end{tabular}
\caption{Sample Data Table}
\label{tab:sample}
\end{table}

\section{Lists}
\subsection{Bullet Points}
\begin{itemize}
  \item First item
  \item Second item
  \item Third item with sub-items:
    \begin{itemize}
      \item Sub-item 1
      \item Sub-item 2
    \end{itemize}
\end{itemize}

\subsection{Numbered List}
\begin{enumerate}
  \item First step
  \item Second step
  \item Third step
\end{enumerate}

\section{Citations and References}
When you use the citation search feature, references will be automatically managed. For example, you might cite a paper like this: \cite{example2024}. All citations will be automatically added to your \texttt{cite.bib} file.

\section{Conclusion}
This template demonstrates basic \LaTeX{} features including mathematical formulas, tables, and lists. The bibliography system is now set up to work seamlessly with Bibby AI's citation management. Feel free to modify it for your needs!

\bibliographystyle{plain}
\bibliography{cite}

\end{document}\documentclass{article}
\usepackage[utf8]{inputenc}
\usepackage{booktabs}
\usepackage{amsmath}
\usepackage{natbib}

\title{My Document}
\author{Your Name}
\date{\today}

\begin{document}

\maketitle

\section{Introduction}
This is a sample document demonstrating various \LaTeX{} features including mathematical equations, tables, and formatting. You can cite papers using the citation search feature, and they will be automatically added to your \texttt{cite.bib} file.

\section{Mathematical Expressions}
Here are some example equations:

\subsection{Quadratic Formula}
The quadratic formula is given by:
\[x = \frac{-b \pm \sqrt{b^2 - 4ac}}{2a}\]

\subsection{Matrix Example}
A sample matrix:
\[
\begin{pmatrix}
a_{11} & a_{12} & a_{13} \\
a_{21} & a_{22} & a_{23} \\
a_{31} & a_{32} & a_{33}
\end{pmatrix}
\]

\section{Tables}
Below is a sample table showing data:

\begin{table}[h]
\centering
\begin{tabular}{@{}lrrr@{}}
\toprule
\textbf{Category} & \textbf{Value 1} & \textbf{Value 2} & \textbf{Total} \\
\midrule
Item A & 10.5 & 20.3 & 30.8 \\
Item B & 15.2 & 25.1 & 40.3 \\
Item C & 22.1 & 12.7 & 34.8 \\
\bottomrule
\end{tabular}
\caption{Sample Data Table}
\label{tab:sample}
\end{table}

\section{Lists}
\subsection{Bullet Points}
\begin{itemize}
  \item First item
  \item Second item
  \item Third item with sub-items:
    \begin{itemize}
      \item Sub-item 1
      \item Sub-item 2
    \end{itemize}
\end{itemize}

\subsection{Numbered List}
\begin{enumerate}
  \item First step
  \item Second step
  \item Third step
\end{enumerate}

\section{Citations and References}
When you use the citation search feature, references will be automatically managed. For example, you might cite a paper like this: \cite{example2024}. All citations will be automatically added to your \texttt{cite.bib} file.

\section{Conclusion}
This template demonstrates basic \LaTeX{} features including mathematical formulas, tables, and lists. The bibliography system is now set up to work seamlessly with Bibby AI's citation management. Feel free to modify it for your needs!

\bibliographystyle{plain}
\bibliography{cite}

\end{document}\documentclass{article}
\usepackage[utf8]{inputenc}
\usepackage{booktabs}
\usepackage{amsmath}
\usepackage{natbib}

\title{My Document}
\author{Your Name}
\date{\today}

\begin{document}

\maketitle

\section{Introduction}
This is a sample document demonstrating various \LaTeX{} features including mathematical equations, tables, and formatting. You can cite papers using the citation search feature, and they will be automatically added to your \texttt{cite.bib} file.

\section{Mathematical Expressions}
Here are some example equations:

\subsection{Quadratic Formula}
The quadratic formula is given by:
\[x = \frac{-b \pm \sqrt{b^2 - 4ac}}{2a}\]

\subsection{Matrix Example}
A sample matrix:
\[
\begin{pmatrix}
a_{11} & a_{12} & a_{13} \\
a_{21} & a_{22} & a_{23} \\
a_{31} & a_{32} & a_{33}
\end{pmatrix}
\]

\section{Tables}
Below is a sample table showing data:

\begin{table}[h]
\centering
\begin{tabular}{@{}lrrr@{}}
\toprule
\textbf{Category} & \textbf{Value 1} & \textbf{Value 2} & \textbf{Total} \\
\midrule
Item A & 10.5 & 20.3 & 30.8 \\
Item B & 15.2 & 25.1 & 40.3 \\
Item C & 22.1 & 12.7 & 34.8 \\
\bottomrule
\end{tabular}
\caption{Sample Data Table}
\label{tab:sample}
\end{table}

\section{Lists}
\subsection{Bullet Points}
\begin{itemize}
  \item First item
  \item Second item
  \item Third item with sub-items:
    \begin{itemize}
      \item Sub-item 1
      \item Sub-item 2
    \end{itemize}
\end{itemize}

\subsection{Numbered List}
\begin{enumerate}
  \item First step
  \item Second step
  \item Third step
\end{enumerate}

\section{Citations and References}
When you use the citation search feature, references will be automatically managed. For example, you might cite a paper like this: \cite{example2024}. All citations will be automatically added to your \texttt{cite.bib} file.

\section{Conclusion}
This template demonstrates basic \LaTeX{} features including mathematical formulas, tables, and lists. The bibliography system is now set up to work seamlessly with Bibby AI's citation management. Feel free to modify it for your needs!

\bibliographystyle{plain}
\bibliography{cite}

\end{document}\documentclass{article}
\usepackage[utf8]{inputenc}
\usepackage{booktabs}
\usepackage{amsmath}
\usepackage{natbib}

\title{My Document}
\author{Your Name}
\date{\today}

\begin{document}

\maketitle

\section{Introduction}
This is a sample document demonstrating various \LaTeX{} features including mathematical equations, tables, and formatting. You can cite papers using the citation search feature, and they will be automatically added to your \texttt{cite.bib} file.

\section{Mathematical Expressions}
Here are some example equations:

\subsection{Quadratic Formula}
The quadratic formula is given by:
\[x = \frac{-b \pm \sqrt{b^2 - 4ac}}{2a}\]

\subsection{Matrix Example}
A sample matrix:
\[
\begin{pmatrix}
a_{11} & a_{12} & a_{13} \\
a_{21} & a_{22} & a_{23} \\
a_{31} & a_{32} & a_{33}
\end{pmatrix}
\]

\section{Tables}
Below is a sample table showing data:

\begin{table}[h]
\centering
\begin{tabular}{@{}lrrr@{}}
\toprule
\textbf{Category} & \textbf{Value 1} & \textbf{Value 2} & \textbf{Total} \\
\midrule
Item A & 10.5 & 20.3 & 30.8 \\
Item B & 15.2 & 25.1 & 40.3 \\
Item C & 22.1 & 12.7 & 34.8 \\
\bottomrule
\end{tabular}
\caption{Sample Data Table}
\label{tab:sample}
\end{table}

\section{Lists}
\subsection{Bullet Points}
\begin{itemize}
  \item First item
  \item Second item
  \item Third item with sub-items:
    \begin{itemize}
      \item Sub-item 1
      \item Sub-item 2
    \end{itemize}
\end{itemize}

\subsection{Numbered List}
\begin{enumerate}
  \item First step
  \item Second step
  \item Third step
\end{enumerate}

\section{Citations and References}
When you use the citation search feature, references will be automatically managed. For example, you might cite a paper like this: \cite{example2024}. All citations will be automatically added to your \texttt{cite.bib} file.

\section{Conclusion}
This template demonstrates basic \LaTeX{} features including mathematical formulas, tables, and lists. The bibliography system is now set up to work seamlessly with Bibby AI's citation management. Feel free to modify it for your needs!

\bibliographystyle{plain}
\bibliography{cite}

\end{document}\documentclass{article}
\usepackage[utf8]{inputenc}
\usepackage{booktabs}
\usepackage{amsmath}
\usepackage{natbib}

\title{My Document}
\author{Your Name}
\date{\today}

\begin{document}

\maketitle

\section{Introduction}
This is a sample document demonstrating various \LaTeX{} features including mathematical equations, tables, and formatting. You can cite papers using the citation search feature, and they will be automatically added to your \texttt{cite.bib} file.

\section{Mathematical Expressions}
Here are some example equations:

\subsection{Quadratic Formula}
The quadratic formula is given by:
\[x = \frac{-b \pm \sqrt{b^2 - 4ac}}{2a}\]

\subsection{Matrix Example}
A sample matrix:
\[
\begin{pmatrix}
a_{11} & a_{12} & a_{13} \\
a_{21} & a_{22} & a_{23} \\
a_{31} & a_{32} & a_{33}
\end{pmatrix}
\]

\section{Tables}
Below is a sample table showing data:

\begin{table}[h]
\centering
\begin{tabular}{@{}lrrr@{}}
\toprule
\textbf{Category} & \textbf{Value 1} & \textbf{Value 2} & \textbf{Total} \\
\midrule
Item A & 10.5 & 20.3 & 30.8 \\
Item B & 15.2 & 25.1 & 40.3 \\
Item C & 22.1 & 12.7 & 34.8 \\
\bottomrule
\end{tabular}
\caption{Sample Data Table}
\label{tab:sample}
\end{table}

\section{Lists}
\subsection{Bullet Points}
\begin{itemize}
  \item First item
  \item Second item
  \item Third item with sub-items:
    \begin{itemize}
      \item Sub-item 1
      \item Sub-item 2
    \end{itemize}
\end{itemize}

\subsection{Numbered List}
\begin{enumerate}
  \item First step
  \item Second step
  \item Third step
\end{enumerate}

\section{Citations and References}
When you use the citation search feature, references will be automatically managed. For example, you might cite a paper like this: \cite{example2024}. All citations will be automatically added to your \texttt{cite.bib} file.

\section{Conclusion}
This template demonstrates basic \LaTeX{} features including mathematical formulas, tables, and lists. The bibliography system is now set up to work seamlessly with Bibby AI's citation management. Feel free to modify it for your needs!

\bibliographystyle{plain}
\bibliography{cite}

\end{document}\documentclass{article}
\usepackage[utf8]{inputenc}
\usepackage{booktabs}
\usepackage{amsmath}
\usepackage{natbib}

\title{My Document}
\author{Your Name}
\date{\today}

\begin{document}

\maketitle

\section{Introduction}
This is a sample document demonstrating various \LaTeX{} features including mathematical equations, tables, and formatting. You can cite papers using the citation search feature, and they will be automatically added to your \texttt{cite.bib} file.

\section{Mathematical Expressions}
Here are some example equations:

\subsection{Quadratic Formula}
The quadratic formula is given by:
\[x = \frac{-b \pm \sqrt{b^2 - 4ac}}{2a}\]

\subsection{Matrix Example}
A sample matrix:
\[
\begin{pmatrix}
a_{11} & a_{12} & a_{13} \\
a_{21} & a_{22} & a_{23} \\
a_{31} & a_{32} & a_{33}
\end{pmatrix}
\]

\section{Tables}
Below is a sample table showing data:

\begin{table}[h]
\centering
\begin{tabular}{@{}lrrr@{}}
\toprule
\textbf{Category} & \textbf{Value 1} & \textbf{Value 2} & \textbf{Total} \\
\midrule
Item A & 10.5 & 20.3 & 30.8 \\
Item B & 15.2 & 25.1 & 40.3 \\
Item C & 22.1 & 12.7 & 34.8 \\
\bottomrule
\end{tabular}
\caption{Sample Data Table}
\label{tab:sample}
\end{table}

\section{Lists}
\subsection{Bullet Points}
\begin{itemize}
  \item First item
  \item Second item
  \item Third item with sub-items:
    \begin{itemize}
      \item Sub-item 1
      \item Sub-item 2
    \end{itemize}
\end{itemize}

\subsection{Numbered List}
\begin{enumerate}
  \item First step
  \item Second step
  \item Third step
\end{enumerate}

\section{Citations and References}
When you use the citation search feature, references will be automatically managed. For example, you might cite a paper like this: \cite{example2024}. All citations will be automatically added to your \texttt{cite.bib} file.

\section{Conclusion}
This template demonstrates basic \LaTeX{} features including mathematical formulas, tables, and lists. The bibliography system is now set up to work seamlessly with Bibby AI's citation management. Feel free to modify it for your needs!

\bibliographystyle{plain}
\bibliography{cite}

\end{document}\documentclass{article}
\usepackage[utf8]{inputenc}
\usepackage{booktabs}
\usepackage{amsmath}
\usepackage{natbib}

\title{My Document}
\author{Your Name}
\date{\today}

\begin{document}

\maketitle

\section{Introduction}
This is a sample document demonstrating various \LaTeX{} features including mathematical equations, tables, and formatting. You can cite papers using the citation search feature, and they will be automatically added to your \texttt{cite.bib} file.

\section{Mathematical Expressions}
Here are some example equations:

\subsection{Quadratic Formula}
The quadratic formula is given by:
\[x = \frac{-b \pm \sqrt{b^2 - 4ac}}{2a}\]

\subsection{Matrix Example}
A sample matrix:
\[
\begin{pmatrix}
a_{11} & a_{12} & a_{13} \\
a_{21} & a_{22} & a_{23} \\
a_{31} & a_{32} & a_{33}
\end{pmatrix}
\]

\section{Tables}
Below is a sample table showing data:

\begin{table}[h]
\centering
\begin{tabular}{@{}lrrr@{}}
\toprule
\textbf{Category} & \textbf{Value 1} & \textbf{Value 2} & \textbf{Total} \\
\midrule
Item A & 10.5 & 20.3 & 30.8 \\
Item B & 15.2 & 25.1 & 40.3 \\
Item C & 22.1 & 12.7 & 34.8 \\
\bottomrule
\end{tabular}
\caption{Sample Data Table}
\label{tab:sample}
\end{table}

\section{Lists}
\subsection{Bullet Points}
\begin{itemize}
  \item First item
  \item Second item
  \item Third item with sub-items:
    \begin{itemize}
      \item Sub-item 1
      \item Sub-item 2
    \end{itemize}
\end{itemize}

\subsection{Numbered List}
\begin{enumerate}
  \item First step
  \item Second step
  \item Third step
\end{enumerate}

\section{Citations and References}
When you use the citation search feature, references will be automatically managed. For example, you might cite a paper like this: \cite{example2024}. All citations will be automatically added to your \texttt{cite.bib} file.

\section{Conclusion}
This template demonstrates basic \LaTeX{} features including mathematical formulas, tables, and lists. The bibliography system is now set up to work seamlessly with Bibby AI's citation management. Feel free to modify it for your needs!

\bibliographystyle{plain}
\bibliography{cite}

\end{document}\documentclass{article}
\usepackage[utf8]{inputenc}
\usepackage{booktabs}
\usepackage{amsmath}
\usepackage{natbib}

\title{My Document}
\author{Your Name}
\date{\today}

\begin{document}

\maketitle

\section{Introduction}
This is a sample document demonstrating various \LaTeX{} features including mathematical equations, tables, and formatting. You can cite papers using the citation search feature, and they will be automatically added to your \texttt{cite.bib} file.

\section{Mathematical Expressions}
Here are some example equations:

\subsection{Quadratic Formula}
The quadratic formula is given by:
\[x = \frac{-b \pm \sqrt{b^2 - 4ac}}{2a}\]

\subsection{Matrix Example}
A sample matrix:
\[
\begin{pmatrix}
a_{11} & a_{12} & a_{13} \\
a_{21} & a_{22} & a_{23} \\
a_{31} & a_{32} & a_{33}
\end{pmatrix}
\]

\section{Tables}
Below is a sample table showing data:

\begin{table}[h]
\centering
\begin{tabular}{@{}lrrr@{}}
\toprule
\textbf{Category} & \textbf{Value 1} & \textbf{Value 2} & \textbf{Total} \\
\midrule
Item A & 10.5 & 20.3 & 30.8 \\
Item B & 15.2 & 25.1 & 40.3 \\
Item C & 22.1 & 12.7 & 34.8 \\
\bottomrule
\end{tabular}
\caption{Sample Data Table}
\label{tab:sample}
\end{table}

\section{Lists}
\subsection{Bullet Points}
\begin{itemize}
  \item First item
  \item Second item
  \item Third item with sub-items:
    \begin{itemize}
      \item Sub-item 1
      \item Sub-item 2
    \end{itemize}
\end{itemize}

\subsection{Numbered List}
\begin{enumerate}
  \item First step
  \item Second step
  \item Third step
\end{enumerate}

\section{Citations and References}
When you use the citation search feature, references will be automatically managed. For example, you might cite a paper like this: \cite{example2024}. All citations will be automatically added to your \texttt{cite.bib} file.

\section{Conclusion}
This template demonstrates basic \LaTeX{} features including mathematical formulas, tables, and lists. The bibliography system is now set up to work seamlessly with Bibby AI's citation management. Feel free to modify it for your needs!

\bibliographystyle{plain}
\bibliography{cite}

\end{document}\documentclass{article}
\usepackage[utf8]{inputenc}
\usepackage{booktabs}
\usepackage{amsmath}
\usepackage{natbib}

\title{My Document}
\author{Your Name}
\date{\today}

\begin{document}

\maketitle

\section{Introduction}
This is a sample document demonstrating various \LaTeX{} features including mathematical equations, tables, and formatting. You can cite papers using the citation search feature, and they will be automatically added to your \texttt{cite.bib} file.

\section{Mathematical Expressions}
Here are some example equations:

\subsection{Quadratic Formula}
The quadratic formula is given by:
\[x = \frac{-b \pm \sqrt{b^2 - 4ac}}{2a}\]

\subsection{Matrix Example}
A sample matrix:
\[
\begin{pmatrix}
a_{11} & a_{12} & a_{13} \\
a_{21} & a_{22} & a_{23} \\
a_{31} & a_{32} & a_{33}
\end{pmatrix}
\]

\section{Tables}
Below is a sample table showing data:

\begin{table}[h]
\centering
\begin{tabular}{@{}lrrr@{}}
\toprule
\textbf{Category} & \textbf{Value 1} & \textbf{Value 2} & \textbf{Total} \\
\midrule
Item A & 10.5 & 20.3 & 30.8 \\
Item B & 15.2 & 25.1 & 40.3 \\
Item C & 22.1 & 12.7 & 34.8 \\
\bottomrule
\end{tabular}
\caption{Sample Data Table}
\label{tab:sample}
\end{table}

\section{Lists}
\subsection{Bullet Points}
\begin{itemize}
  \item First item
  \item Second item
  \item Third item with sub-items:
    \begin{itemize}
      \item Sub-item 1
      \item Sub-item 2
    \end{itemize}
\end{itemize}

\subsection{Numbered List}
\begin{enumerate}
  \item First step
  \item Second step
  \item Third step
\end{enumerate}

\section{Citations and References}
When you use the citation search feature, references will be automatically managed. For example, you might cite a paper like this: \cite{example2024}. All citations will be automatically added to your \texttt{cite.bib} file.

\section{Conclusion}
This template demonstrates basic \LaTeX{} features including mathematical formulas, tables, and lists. The bibliography system is now set up to work seamlessly with Bibby AI's citation management. Feel free to modify it for your needs!

\bibliographystyle{plain}
\bibliography{cite}

\end{document}\documentclass{article}
\usepackage[utf8]{inputenc}
\usepackage{booktabs}
\usepackage{amsmath}
\usepackage{natbib}

\title{My Document}
\author{Your Name}
\date{\today}

\begin{document}

\maketitle

\section{Introduction}
This is a sample document demonstrating various \LaTeX{} features including mathematical equations, tables, and formatting. You can cite papers using the citation search feature, and they will be automatically added to your \texttt{cite.bib} file.

\section{Mathematical Expressions}
Here are some example equations:

\subsection{Quadratic Formula}
The quadratic formula is given by:
\[x = \frac{-b \pm \sqrt{b^2 - 4ac}}{2a}\]

\subsection{Matrix Example}
A sample matrix:
\[
\begin{pmatrix}
a_{11} & a_{12} & a_{13} \\
a_{21} & a_{22} & a_{23} \\
a_{31} & a_{32} & a_{33}
\end{pmatrix}
\]

\section{Tables}
Below is a sample table showing data:

\begin{table}[h]
\centering
\begin{tabular}{@{}lrrr@{}}
\toprule
\textbf{Category} & \textbf{Value 1} & \textbf{Value 2} & \textbf{Total} \\
\midrule
Item A & 10.5 & 20.3 & 30.8 \\
Item B & 15.2 & 25.1 & 40.3 \\
Item C & 22.1 & 12.7 & 34.8 \\
\bottomrule
\end{tabular}
\caption{Sample Data Table}
\label{tab:sample}
\end{table}

\section{Lists}
\subsection{Bullet Points}
\begin{itemize}
  \item First item
  \item Second item
  \item Third item with sub-items:
    \begin{itemize}
      \item Sub-item 1
      \item Sub-item 2
    \end{itemize}
\end{itemize}

\subsection{Numbered List}
\begin{enumerate}
  \item First step
  \item Second step
  \item Third step
\end{enumerate}

\section{Citations and References}
When you use the citation search feature, references will be automatically managed. For example, you might cite a paper like this: \cite{example2024}. All citations will be automatically added to your \texttt{cite.bib} file.

\section{Conclusion}
This template demonstrates basic \LaTeX{} features including mathematical formulas, tables, and lists. The bibliography system is now set up to work seamlessly with Bibby AI's citation management. Feel free to modify it for your needs!

\bibliographystyle{plain}
\bibliography{cite}

\end{document}\documentclass{article}
\usepackage[utf8]{inputenc}
\usepackage{booktabs}
\usepackage{amsmath}
\usepackage{natbib}

\title{My Document}
\author{Your Name}
\date{\today}

\begin{document}

\maketitle

\section{Introduction}
This is a sample document demonstrating various \LaTeX{} features including mathematical equations, tables, and formatting. You can cite papers using the citation search feature, and they will be automatically added to your \texttt{cite.bib} file.

\section{Mathematical Expressions}
Here are some example equations:

\subsection{Quadratic Formula}
The quadratic formula is given by:
\[x = \frac{-b \pm \sqrt{b^2 - 4ac}}{2a}\]

\subsection{Matrix Example}
A sample matrix:
\[
\begin{pmatrix}
a_{11} & a_{12} & a_{13} \\
a_{21} & a_{22} & a_{23} \\
a_{31} & a_{32} & a_{33}
\end{pmatrix}
\]

\section{Tables}
Below is a sample table showing data:

\begin{table}[h]
\centering
\begin{tabular}{@{}lrrr@{}}
\toprule
\textbf{Category} & \textbf{Value 1} & \textbf{Value 2} & \textbf{Total} \\
\midrule
Item A & 10.5 & 20.3 & 30.8 \\
Item B & 15.2 & 25.1 & 40.3 \\
Item C & 22.1 & 12.7 & 34.8 \\
\bottomrule
\end{tabular}
\caption{Sample Data Table}
\label{tab:sample}
\end{table}

\section{Lists}
\subsection{Bullet Points}
\begin{itemize}
  \item First item
  \item Second item
  \item Third item with sub-items:
    \begin{itemize}
      \item Sub-item 1
      \item Sub-item 2
    \end{itemize}
\end{itemize}

\subsection{Numbered List}
\begin{enumerate}
  \item First step
  \item Second step
  \item Third step
\end{enumerate}

\section{Citations and References}
When you use the citation search feature, references will be automatically managed. For example, you might cite a paper like this: \cite{example2024}. All citations will be automatically added to your \texttt{cite.bib} file.

\section{Conclusion}
This template demonstrates basic \LaTeX{} features including mathematical formulas, tables, and lists. The bibliography system is now set up to work seamlessly with Bibby AI's citation management. Feel free to modify it for your needs!

\bibliographystyle{plain}
\bibliography{cite}

\end{document}\documentclass{article}
\usepackage[utf8]{inputenc}
\usepackage{booktabs}
\usepackage{amsmath}
\usepackage{natbib}

\title{My Document}
\author{Your Name}
\date{\today}

\begin{document}

\maketitle

\section{Introduction}
This is a sample document demonstrating various \LaTeX{} features including mathematical equations, tables, and formatting. You can cite papers using the citation search feature, and they will be automatically added to your \texttt{cite.bib} file.

\section{Mathematical Expressions}
Here are some example equations:

\subsection{Quadratic Formula}
The quadratic formula is given by:
\[x = \frac{-b \pm \sqrt{b^2 - 4ac}}{2a}\]

\subsection{Matrix Example}
A sample matrix:
\[
\begin{pmatrix}
a_{11} & a_{12} & a_{13} \\
a_{21} & a_{22} & a_{23} \\
a_{31} & a_{32} & a_{33}
\end{pmatrix}
\]

\section{Tables}
Below is a sample table showing data:

\begin{table}[h]
\centering
\begin{tabular}{@{}lrrr@{}}
\toprule
\textbf{Category} & \textbf{Value 1} & \textbf{Value 2} & \textbf{Total} \\
\midrule
Item A & 10.5 & 20.3 & 30.8 \\
Item B & 15.2 & 25.1 & 40.3 \\
Item C & 22.1 & 12.7 & 34.8 \\
\bottomrule
\end{tabular}
\caption{Sample Data Table}
\label{tab:sample}
\end{table}

\section{Lists}
\subsection{Bullet Points}
\begin{itemize}
  \item First item
  \item Second item
  \item Third item with sub-items:
    \begin{itemize}
      \item Sub-item 1
      \item Sub-item 2
    \end{itemize}
\end{itemize}

\subsection{Numbered List}
\begin{enumerate}
  \item First step
  \item Second step
  \item Third step
\end{enumerate}

\section{Citations and References}
When you use the citation search feature, references will be automatically managed. For example, you might cite a paper like this: \cite{example2024}. All citations will be automatically added to your \texttt{cite.bib} file.

\section{Conclusion}
This template demonstrates basic \LaTeX{} features including mathematical formulas, tables, and lists. The bibliography system is now set up to work seamlessly with Bibby AI's citation management. Feel free to modify it for your needs!

\bibliographystyle{plain}
\bibliography{cite}

\end{document}\documentclass{article}
\usepackage[utf8]{inputenc}
\usepackage{booktabs}
\usepackage{amsmath}
\usepackage{natbib}

\title{My Document}
\author{Your Name}
\date{\today}

\begin{document}

\maketitle

\section{Introduction}
This is a sample document demonstrating various \LaTeX{} features including mathematical equations, tables, and formatting. You can cite papers using the citation search feature, and they will be automatically added to your \texttt{cite.bib} file.

\section{Mathematical Expressions}
Here are some example equations:

\subsection{Quadratic Formula}
The quadratic formula is given by:
\[x = \frac{-b \pm \sqrt{b^2 - 4ac}}{2a}\]

\subsection{Matrix Example}
A sample matrix:
\[
\begin{pmatrix}
a_{11} & a_{12} & a_{13} \\
a_{21} & a_{22} & a_{23} \\
a_{31} & a_{32} & a_{33}
\end{pmatrix}
\]

\section{Tables}
Below is a sample table showing data:

\begin{table}[h]
\centering
\begin{tabular}{@{}lrrr@{}}
\toprule
\textbf{Category} & \textbf{Value 1} & \textbf{Value 2} & \textbf{Total} \\
\midrule
Item A & 10.5 & 20.3 & 30.8 \\
Item B & 15.2 & 25.1 & 40.3 \\
Item C & 22.1 & 12.7 & 34.8 \\
\bottomrule
\end{tabular}
\caption{Sample Data Table}
\label{tab:sample}
\end{table}

\section{Lists}
\subsection{Bullet Points}
\begin{itemize}
  \item First item
  \item Second item
  \item Third item with sub-items:
    \begin{itemize}
      \item Sub-item 1
      \item Sub-item 2
    \end{itemize}
\end{itemize}

\subsection{Numbered List}
\begin{enumerate}
  \item First step
  \item Second step
  \item Third step
\end{enumerate}

\section{Citations and References}
When you use the citation search feature, references will be automatically managed. For example, you might cite a paper like this: \cite{example2024}. All citations will be automatically added to your \texttt{cite.bib} file.

\section{Conclusion}
This template demonstrates basic \LaTeX{} features including mathematical formulas, tables, and lists. The bibliography system is now set up to work seamlessly with Bibby AI's citation management. Feel free to modify it for your needs!

\bibliographystyle{plain}
\bibliography{cite}

\end{document}\documentclass{article}
\usepackage[utf8]{inputenc}
\usepackage{booktabs}
\usepackage{amsmath}
\usepackage{natbib}

\title{My Document}
\author{Your Name}
\date{\today}

\begin{document}

\maketitle

\section{Introduction}
This is a sample document demonstrating various \LaTeX{} features including mathematical equations, tables, and formatting. You can cite papers using the citation search feature, and they will be automatically added to your \texttt{cite.bib} file.

\section{Mathematical Expressions}
Here are some example equations:

\subsection{Quadratic Formula}
The quadratic formula is given by:
\[x = \frac{-b \pm \sqrt{b^2 - 4ac}}{2a}\]

\subsection{Matrix Example}
A sample matrix:
\[
\begin{pmatrix}
a_{11} & a_{12} & a_{13} \\
a_{21} & a_{22} & a_{23} \\
a_{31} & a_{32} & a_{33}
\end{pmatrix}
\]

\section{Tables}
Below is a sample table showing data:

\begin{table}[h]
\centering
\begin{tabular}{@{}lrrr@{}}
\toprule
\textbf{Category} & \textbf{Value 1} & \textbf{Value 2} & \textbf{Total} \\
\midrule
Item A & 10.5 & 20.3 & 30.8 \\
Item B & 15.2 & 25.1 & 40.3 \\
Item C & 22.1 & 12.7 & 34.8 \\
\bottomrule
\end{tabular}
\caption{Sample Data Table}
\label{tab:sample}
\end{table}

\section{Lists}
\subsection{Bullet Points}
\begin{itemize}
  \item First item
  \item Second item
  \item Third item with sub-items:
    \begin{itemize}
      \item Sub-item 1
      \item Sub-item 2
    \end{itemize}
\end{itemize}

\subsection{Numbered List}
\begin{enumerate}
  \item First step
  \item Second step
  \item Third step
\end{enumerate}

\section{Citations and References}
When you use the citation search feature, references will be automatically managed. For example, you might cite a paper like this: \cite{example2024}. All citations will be automatically added to your \texttt{cite.bib} file.

\section{Conclusion}
This template demonstrates basic \LaTeX{} features including mathematical formulas, tables, and lists. The bibliography system is now set up to work seamlessly with Bibby AI's citation management. Feel free to modify it for your needs!

\bibliographystyle{plain}
\bibliography{cite}

\end{document}\documentclass{article}
\usepackage[utf8]{inputenc}
\usepackage{booktabs}
\usepackage{amsmath}
\usepackage{natbib}

\title{My Document}
\author{Your Name}
\date{\today}

\begin{document}

\maketitle

\section{Introduction}
This is a sample document demonstrating various \LaTeX{} features including mathematical equations, tables, and formatting. You can cite papers using the citation search feature, and they will be automatically added to your \texttt{cite.bib} file.

\section{Mathematical Expressions}
Here are some example equations:

\subsection{Quadratic Formula}
The quadratic formula is given by:
\[x = \frac{-b \pm \sqrt{b^2 - 4ac}}{2a}\]

\subsection{Matrix Example}
A sample matrix:
\[
\begin{pmatrix}
a_{11} & a_{12} & a_{13} \\
a_{21} & a_{22} & a_{23} \\
a_{31} & a_{32} & a_{33}
\end{pmatrix}
\]

\section{Tables}
Below is a sample table showing data:

\begin{table}[h]
\centering
\begin{tabular}{@{}lrrr@{}}
\toprule
\textbf{Category} & \textbf{Value 1} & \textbf{Value 2} & \textbf{Total} \\
\midrule
Item A & 10.5 & 20.3 & 30.8 \\
Item B & 15.2 & 25.1 & 40.3 \\
Item C & 22.1 & 12.7 & 34.8 \\
\bottomrule
\end{tabular}
\caption{Sample Data Table}
\label{tab:sample}
\end{table}

\section{Lists}
\subsection{Bullet Points}
\begin{itemize}
  \item First item
  \item Second item
  \item Third item with sub-items:
    \begin{itemize}
      \item Sub-item 1
      \item Sub-item 2
    \end{itemize}
\end{itemize}

\subsection{Numbered List}
\begin{enumerate}
  \item First step
  \item Second step
  \item Third step
\end{enumerate}

\section{Citations and References}
When you use the citation search feature, references will be automatically managed. For example, you might cite a paper like this: \cite{example2024}. All citations will be automatically added to your \texttt{cite.bib} file.

\section{Conclusion}
This template demonstrates basic \LaTeX{} features including mathematical formulas, tables, and lists. The bibliography system is now set up to work seamlessly with Bibby AI's citation management. Feel free to modify it for your needs!

\bibliographystyle{plain}
\bibliography{cite}

\end{document}\documentclass{article}
\usepackage[utf8]{inputenc}
\usepackage{booktabs}
\usepackage{amsmath}
\usepackage{natbib}

\title{My Document}
\author{Your Name}
\date{\today}

\begin{document}

\maketitle

\section{Introduction}
This is a sample document demonstrating various \LaTeX{} features including mathematical equations, tables, and formatting. You can cite papers using the citation search feature, and they will be automatically added to your \texttt{cite.bib} file.

\section{Mathematical Expressions}
Here are some example equations:

\subsection{Quadratic Formula}
The quadratic formula is given by:
\[x = \frac{-b \pm \sqrt{b^2 - 4ac}}{2a}\]

\subsection{Matrix Example}
A sample matrix:
\[
\begin{pmatrix}
a_{11} & a_{12} & a_{13} \\
a_{21} & a_{22} & a_{23} \\
a_{31} & a_{32} & a_{33}
\end{pmatrix}
\]

\section{Tables}
Below is a sample table showing data:

\begin{table}[h]
\centering
\begin{tabular}{@{}lrrr@{}}
\toprule
\textbf{Category} & \textbf{Value 1} & \textbf{Value 2} & \textbf{Total} \\
\midrule
Item A & 10.5 & 20.3 & 30.8 \\
Item B & 15.2 & 25.1 & 40.3 \\
Item C & 22.1 & 12.7 & 34.8 \\
\bottomrule
\end{tabular}
\caption{Sample Data Table}
\label{tab:sample}
\end{table}

\section{Lists}
\subsection{Bullet Points}
\begin{itemize}
  \item First item
  \item Second item
  \item Third item with sub-items:
    \begin{itemize}
      \item Sub-item 1
      \item Sub-item 2
    \end{itemize}
\end{itemize}

\subsection{Numbered List}
\begin{enumerate}
  \item First step
  \item Second step
  \item Third step
\end{enumerate}

\section{Citations and References}
When you use the citation search feature, references will be automatically managed. For example, you might cite a paper like this: \cite{example2024}. All citations will be automatically added to your \texttt{cite.bib} file.

\section{Conclusion}
This template demonstrates basic \LaTeX{} features including mathematical formulas, tables, and lists. The bibliography system is now set up to work seamlessly with Bibby AI's citation management. Feel free to modify it for your needs!

\bibliographystyle{plain}
\bibliography{cite}

\end{document}\documentclass{article}
\usepackage[utf8]{inputenc}
\usepackage{booktabs}
\usepackage{amsmath}
\usepackage{natbib}

\title{My Document}
\author{Your Name}
\date{\today}

\begin{document}

\maketitle

\section{Introduction}
This is a sample document demonstrating various \LaTeX{} features including mathematical equations, tables, and formatting. You can cite papers using the citation search feature, and they will be automatically added to your \texttt{cite.bib} file.

\section{Mathematical Expressions}
Here are some example equations:

\subsection{Quadratic Formula}
The quadratic formula is given by:
\[x = \frac{-b \pm \sqrt{b^2 - 4ac}}{2a}\]

\subsection{Matrix Example}
A sample matrix:
\[
\begin{pmatrix}
a_{11} & a_{12} & a_{13} \\
a_{21} & a_{22} & a_{23} \\
a_{31} & a_{32} & a_{33}
\end{pmatrix}
\]

\section{Tables}
Below is a sample table showing data:

\begin{table}[h]
\centering
\begin{tabular}{@{}lrrr@{}}
\toprule
\textbf{Category} & \textbf{Value 1} & \textbf{Value 2} & \textbf{Total} \\
\midrule
Item A & 10.5 & 20.3 & 30.8 \\
Item B & 15.2 & 25.1 & 40.3 \\
Item C & 22.1 & 12.7 & 34.8 \\
\bottomrule
\end{tabular}
\caption{Sample Data Table}
\label{tab:sample}
\end{table}

\section{Lists}
\subsection{Bullet Points}
\begin{itemize}
  \item First item
  \item Second item
  \item Third item with sub-items:
    \begin{itemize}
      \item Sub-item 1
      \item Sub-item 2
    \end{itemize}
\end{itemize}

\subsection{Numbered List}
\begin{enumerate}
  \item First step
  \item Second step
  \item Third step
\end{enumerate}

\section{Citations and References}
When you use the citation search feature, references will be automatically managed. For example, you might cite a paper like this: \cite{example2024}. All citations will be automatically added to your \texttt{cite.bib} file.

\section{Conclusion}
This template demonstrates basic \LaTeX{} features including mathematical formulas, tables, and lists. The bibliography system is now set up to work seamlessly with Bibby AI's citation management. Feel free to modify it for your needs!

\bibliographystyle{plain}
\bibliography{cite}

\end{document}\documentclass{article}
\usepackage[utf8]{inputenc}
\usepackage{booktabs}
\usepackage{amsmath}
\usepackage{natbib}

\title{My Document}
\author{Your Name}
\date{\today}

\begin{document}

\maketitle

\section{Introduction}
This is a sample document demonstrating various \LaTeX{} features including mathematical equations, tables, and formatting. You can cite papers using the citation search feature, and they will be automatically added to your \texttt{cite.bib} file.

\section{Mathematical Expressions}
Here are some example equations:

\subsection{Quadratic Formula}
The quadratic formula is given by:
\[x = \frac{-b \pm \sqrt{b^2 - 4ac}}{2a}\]

\subsection{Matrix Example}
A sample matrix:
\[
\begin{pmatrix}
a_{11} & a_{12} & a_{13} \\
a_{21} & a_{22} & a_{23} \\
a_{31} & a_{32} & a_{33}
\end{pmatrix}
\]

\section{Tables}
Below is a sample table showing data:

\begin{table}[h]
\centering
\begin{tabular}{@{}lrrr@{}}
\toprule
\textbf{Category} & \textbf{Value 1} & \textbf{Value 2} & \textbf{Total} \\
\midrule
Item A & 10.5 & 20.3 & 30.8 \\
Item B & 15.2 & 25.1 & 40.3 \\
Item C & 22.1 & 12.7 & 34.8 \\
\bottomrule
\end{tabular}
\caption{Sample Data Table}
\label{tab:sample}
\end{table}

\section{Lists}
\subsection{Bullet Points}
\begin{itemize}
  \item First item
  \item Second item
  \item Third item with sub-items:
    \begin{itemize}
      \item Sub-item 1
      \item Sub-item 2
    \end{itemize}
\end{itemize}

\subsection{Numbered List}
\begin{enumerate}
  \item First step
  \item Second step
  \item Third step
\end{enumerate}

\section{Citations and References}
When you use the citation search feature, references will be automatically managed. For example, you might cite a paper like this: \cite{example2024}. All citations will be automatically added to your \texttt{cite.bib} file.

\section{Conclusion}
This template demonstrates basic \LaTeX{} features including mathematical formulas, tables, and lists. The bibliography system is now set up to work seamlessly with Bibby AI's citation management. Feel free to modify it for your needs!

\bibliographystyle{plain}
\bibliography{cite}

\end{document}\documentclass{article}
\usepackage[utf8]{inputenc}
\usepackage{booktabs}
\usepackage{amsmath}
\usepackage{natbib}

\title{My Document}
\author{Your Name}
\date{\today}

\begin{document}

\maketitle

\section{Introduction}
This is a sample document demonstrating various \LaTeX{} features including mathematical equations, tables, and formatting. You can cite papers using the citation search feature, and they will be automatically added to your \texttt{cite.bib} file.

\section{Mathematical Expressions}
Here are some example equations:

\subsection{Quadratic Formula}
The quadratic formula is given by:
\[x = \frac{-b \pm \sqrt{b^2 - 4ac}}{2a}\]

\subsection{Matrix Example}
A sample matrix:
\[
\begin{pmatrix}
a_{11} & a_{12} & a_{13} \\
a_{21} & a_{22} & a_{23} \\
a_{31} & a_{32} & a_{33}
\end{pmatrix}
\]

\section{Tables}
Below is a sample table showing data:

\begin{table}[h]
\centering
\begin{tabular}{@{}lrrr@{}}
\toprule
\textbf{Category} & \textbf{Value 1} & \textbf{Value 2} & \textbf{Total} \\
\midrule
Item A & 10.5 & 20.3 & 30.8 \\
Item B & 15.2 & 25.1 & 40.3 \\
Item C & 22.1 & 12.7 & 34.8 \\
\bottomrule
\end{tabular}
\caption{Sample Data Table}
\label{tab:sample}
\end{table}

\section{Lists}
\subsection{Bullet Points}
\begin{itemize}
  \item First item
  \item Second item
  \item Third item with sub-items:
    \begin{itemize}
      \item Sub-item 1
      \item Sub-item 2
    \end{itemize}
\end{itemize}

\subsection{Numbered List}
\begin{enumerate}
  \item First step
  \item Second step
  \item Third step
\end{enumerate}

\section{Citations and References}
When you use the citation search feature, references will be automatically managed. For example, you might cite a paper like this: \cite{example2024}. All citations will be automatically added to your \texttt{cite.bib} file.

\section{Conclusion}
This template demonstrates basic \LaTeX{} features including mathematical formulas, tables, and lists. The bibliography system is now set up to work seamlessly with Bibby AI's citation management. Feel free to modify it for your needs!

\bibliographystyle{plain}
\bibliography{cite}

\end{document}\documentclass{article}
\usepackage[utf8]{inputenc}
\usepackage{booktabs}
\usepackage{amsmath}
\usepackage{natbib}

\title{My Document}
\author{Your Name}
\date{\today}

\begin{document}

\maketitle

\section{Introduction}
This is a sample document demonstrating various \LaTeX{} features including mathematical equations, tables, and formatting. You can cite papers using the citation search feature, and they will be automatically added to your \texttt{cite.bib} file.

\section{Mathematical Expressions}
Here are some example equations:

\subsection{Quadratic Formula}
The quadratic formula is given by:
\[x = \frac{-b \pm \sqrt{b^2 - 4ac}}{2a}\]

\subsection{Matrix Example}
A sample matrix:
\[
\begin{pmatrix}
a_{11} & a_{12} & a_{13} \\
a_{21} & a_{22} & a_{23} \\
a_{31} & a_{32} & a_{33}
\end{pmatrix}
\]

\section{Tables}
Below is a sample table showing data:

\begin{table}[h]
\centering
\begin{tabular}{@{}lrrr@{}}
\toprule
\textbf{Category} & \textbf{Value 1} & \textbf{Value 2} & \textbf{Total} \\
\midrule
Item A & 10.5 & 20.3 & 30.8 \\
Item B & 15.2 & 25.1 & 40.3 \\
Item C & 22.1 & 12.7 & 34.8 \\
\bottomrule
\end{tabular}
\caption{Sample Data Table}
\label{tab:sample}
\end{table}

\section{Lists}
\subsection{Bullet Points}
\begin{itemize}
  \item First item
  \item Second item
  \item Third item with sub-items:
    \begin{itemize}
      \item Sub-item 1
      \item Sub-item 2
    \end{itemize}
\end{itemize}

\subsection{Numbered List}
\begin{enumerate}
  \item First step
  \item Second step
  \item Third step
\end{enumerate}

\section{Citations and References}
When you use the citation search feature, references will be automatically managed. For example, you might cite a paper like this: \cite{example2024}. All citations will be automatically added to your \texttt{cite.bib} file.

\section{Conclusion}
This template demonstrates basic \LaTeX{} features including mathematical formulas, tables, and lists. The bibliography system is now set up to work seamlessly with Bibby AI's citation management. Feel free to modify it for your needs!

\bibliographystyle{plain}
\bibliography{cite}

\end{document}\documentclass{article}
\usepackage[utf8]{inputenc}
\usepackage{booktabs}
\usepackage{amsmath}
\usepackage{natbib}

\title{My Document}
\author{Your Name}
\date{\today}

\begin{document}

\maketitle

\section{Introduction}
This is a sample document demonstrating various \LaTeX{} features including mathematical equations, tables, and formatting. You can cite papers using the citation search feature, and they will be automatically added to your \texttt{cite.bib} file.

\section{Mathematical Expressions}
Here are some example equations:

\subsection{Quadratic Formula}
The quadratic formula is given by:
\[x = \frac{-b \pm \sqrt{b^2 - 4ac}}{2a}\]

\subsection{Matrix Example}
A sample matrix:
\[
\begin{pmatrix}
a_{11} & a_{12} & a_{13} \\
a_{21} & a_{22} & a_{23} \\
a_{31} & a_{32} & a_{33}
\end{pmatrix}
\]

\section{Tables}
Below is a sample table showing data:

\begin{table}[h]
\centering
\begin{tabular}{@{}lrrr@{}}
\toprule
\textbf{Category} & \textbf{Value 1} & \textbf{Value 2} & \textbf{Total} \\
\midrule
Item A & 10.5 & 20.3 & 30.8 \\
Item B & 15.2 & 25.1 & 40.3 \\
Item C & 22.1 & 12.7 & 34.8 \\
\bottomrule
\end{tabular}
\caption{Sample Data Table}
\label{tab:sample}
\end{table}

\section{Lists}
\subsection{Bullet Points}
\begin{itemize}
  \item First item
  \item Second item
  \item Third item with sub-items:
    \begin{itemize}
      \item Sub-item 1
      \item Sub-item 2
    \end{itemize}
\end{itemize}

\subsection{Numbered List}
\begin{enumerate}
  \item First step
  \item Second step
  \item Third step
\end{enumerate}

\section{Citations and References}
When you use the citation search feature, references will be automatically managed. For example, you might cite a paper like this: \cite{example2024}. All citations will be automatically added to your \texttt{cite.bib} file.

\section{Conclusion}
This template demonstrates basic \LaTeX{} features including mathematical formulas, tables, and lists. The bibliography system is now set up to work seamlessly with Bibby AI's citation management. Feel free to modify it for your needs!

\bibliographystyle{plain}
\bibliography{cite}

\end{document}\documentclass{article}
\usepackage[utf8]{inputenc}
\usepackage{booktabs}
\usepackage{amsmath}
\usepackage{natbib}

\title{My Document}
\author{Your Name}
\date{\today}

\begin{document}

\maketitle

\section{Introduction}
This is a sample document demonstrating various \LaTeX{} features including mathematical equations, tables, and formatting. You can cite papers using the citation search feature, and they will be automatically added to your \texttt{cite.bib} file.

\section{Mathematical Expressions}
Here are some example equations:

\subsection{Quadratic Formula}
The quadratic formula is given by:
\[x = \frac{-b \pm \sqrt{b^2 - 4ac}}{2a}\]

\subsection{Matrix Example}
A sample matrix:
\[
\begin{pmatrix}
a_{11} & a_{12} & a_{13} \\
a_{21} & a_{22} & a_{23} \\
a_{31} & a_{32} & a_{33}
\end{pmatrix}
\]

\section{Tables}
Below is a sample table showing data:

\begin{table}[h]
\centering
\begin{tabular}{@{}lrrr@{}}
\toprule
\textbf{Category} & \textbf{Value 1} & \textbf{Value 2} & \textbf{Total} \\
\midrule
Item A & 10.5 & 20.3 & 30.8 \\
Item B & 15.2 & 25.1 & 40.3 \\
Item C & 22.1 & 12.7 & 34.8 \\
\bottomrule
\end{tabular}
\caption{Sample Data Table}
\label{tab:sample}
\end{table}

\section{Lists}
\subsection{Bullet Points}
\begin{itemize}
  \item First item
  \item Second item
  \item Third item with sub-items:
    \begin{itemize}
      \item Sub-item 1
      \item Sub-item 2
    \end{itemize}
\end{itemize}

\subsection{Numbered List}
\begin{enumerate}
  \item First step
  \item Second step
  \item Third step
\end{enumerate}

\section{Citations and References}
When you use the citation search feature, references will be automatically managed. For example, you might cite a paper like this: \cite{example2024}. All citations will be automatically added to your \texttt{cite.bib} file.

\section{Conclusion}
This template demonstrates basic \LaTeX{} features including mathematical formulas, tables, and lists. The bibliography system is now set up to work seamlessly with Bibby AI's citation management. Feel free to modify it for your needs!

\bibliographystyle{plain}
\bibliography{cite}

\end{document}\documentclass{article}
\usepackage[utf8]{inputenc}
\usepackage{booktabs}
\usepackage{amsmath}
\usepackage{natbib}

\title{My Document}
\author{Your Name}
\date{\today}

\begin{document}

\maketitle

\section{Introduction}
This is a sample document demonstrating various \LaTeX{} features including mathematical equations, tables, and formatting. You can cite papers using the citation search feature, and they will be automatically added to your \texttt{cite.bib} file.

\section{Mathematical Expressions}
Here are some example equations:

\subsection{Quadratic Formula}
The quadratic formula is given by:
\[x = \frac{-b \pm \sqrt{b^2 - 4ac}}{2a}\]

\subsection{Matrix Example}
A sample matrix:
\[
\begin{pmatrix}
a_{11} & a_{12} & a_{13} \\
a_{21} & a_{22} & a_{23} \\
a_{31} & a_{32} & a_{33}
\end{pmatrix}
\]

\section{Tables}
Below is a sample table showing data:

\begin{table}[h]
\centering
\begin{tabular}{@{}lrrr@{}}
\toprule
\textbf{Category} & \textbf{Value 1} & \textbf{Value 2} & \textbf{Total} \\
\midrule
Item A & 10.5 & 20.3 & 30.8 \\
Item B & 15.2 & 25.1 & 40.3 \\
Item C & 22.1 & 12.7 & 34.8 \\
\bottomrule
\end{tabular}
\caption{Sample Data Table}
\label{tab:sample}
\end{table}

\section{Lists}
\subsection{Bullet Points}
\begin{itemize}
  \item First item
  \item Second item
  \item Third item with sub-items:
    \begin{itemize}
      \item Sub-item 1
      \item Sub-item 2
    \end{itemize}
\end{itemize}

\subsection{Numbered List}
\begin{enumerate}
  \item First step
  \item Second step
  \item Third step
\end{enumerate}

\section{Citations and References}
When you use the citation search feature, references will be automatically managed. For example, you might cite a paper like this: \cite{example2024}. All citations will be automatically added to your \texttt{cite.bib} file.

\section{Conclusion}
This template demonstrates basic \LaTeX{} features including mathematical formulas, tables, and lists. The bibliography system is now set up to work seamlessly with Bibby AI's citation management. Feel free to modify it for your needs!

\bibliographystyle{plain}
\bibliography{cite}

\end{document}\documentclass{article}
\usepackage[utf8]{inputenc}
\usepackage{booktabs}
\usepackage{amsmath}
\usepackage{natbib}

\title{My Document}
\author{Your Name}
\date{\today}

\begin{document}

\maketitle

\section{Introduction}
This is a sample document demonstrating various \LaTeX{} features including mathematical equations, tables, and formatting. You can cite papers using the citation search feature, and they will be automatically added to your \texttt{cite.bib} file.

\section{Mathematical Expressions}
Here are some example equations:

\subsection{Quadratic Formula}
The quadratic formula is given by:
\[x = \frac{-b \pm \sqrt{b^2 - 4ac}}{2a}\]

\subsection{Matrix Example}
A sample matrix:
\[
\begin{pmatrix}
a_{11} & a_{12} & a_{13} \\
a_{21} & a_{22} & a_{23} \\
a_{31} & a_{32} & a_{33}
\end{pmatrix}
\]

\section{Tables}
Below is a sample table showing data:

\begin{table}[h]
\centering
\begin{tabular}{@{}lrrr@{}}
\toprule
\textbf{Category} & \textbf{Value 1} & \textbf{Value 2} & \textbf{Total} \\
\midrule
Item A & 10.5 & 20.3 & 30.8 \\
Item B & 15.2 & 25.1 & 40.3 \\
Item C & 22.1 & 12.7 & 34.8 \\
\bottomrule
\end{tabular}
\caption{Sample Data Table}
\label{tab:sample}
\end{table}

\section{Lists}
\subsection{Bullet Points}
\begin{itemize}
  \item First item
  \item Second item
  \item Third item with sub-items:
    \begin{itemize}
      \item Sub-item 1
      \item Sub-item 2
    \end{itemize}
\end{itemize}

\subsection{Numbered List}
\begin{enumerate}
  \item First step
  \item Second step
  \item Third step
\end{enumerate}

\section{Citations and References}
When you use the citation search feature, references will be automatically managed. For example, you might cite a paper like this: \cite{example2024}. All citations will be automatically added to your \texttt{cite.bib} file.

\section{Conclusion}
This template demonstrates basic \LaTeX{} features including mathematical formulas, tables, and lists. The bibliography system is now set up to work seamlessly with Bibby AI's citation management. Feel free to modify it for your needs!

\bibliographystyle{plain}
\bibliography{cite}

\end{document}\documentclass{article}
\usepackage[utf8]{inputenc}
\usepackage{booktabs}
\usepackage{amsmath}
\usepackage{natbib}

\title{My Document}
\author{Your Name}
\date{\today}

\begin{document}

\maketitle

\section{Introduction}
This is a sample document demonstrating various \LaTeX{} features including mathematical equations, tables, and formatting. You can cite papers using the citation search feature, and they will be automatically added to your \texttt{cite.bib} file.

\section{Mathematical Expressions}
Here are some example equations:

\subsection{Quadratic Formula}
The quadratic formula is given by:
\[x = \frac{-b \pm \sqrt{b^2 - 4ac}}{2a}\]

\subsection{Matrix Example}
A sample matrix:
\[
\begin{pmatrix}
a_{11} & a_{12} & a_{13} \\
a_{21} & a_{22} & a_{23} \\
a_{31} & a_{32} & a_{33}
\end{pmatrix}
\]

\section{Tables}
Below is a sample table showing data:

\begin{table}[h]
\centering
\begin{tabular}{@{}lrrr@{}}
\toprule
\textbf{Category} & \textbf{Value 1} & \textbf{Value 2} & \textbf{Total} \\
\midrule
Item A & 10.5 & 20.3 & 30.8 \\
Item B & 15.2 & 25.1 & 40.3 \\
Item C & 22.1 & 12.7 & 34.8 \\
\bottomrule
\end{tabular}
\caption{Sample Data Table}
\label{tab:sample}
\end{table}

\section{Lists}
\subsection{Bullet Points}
\begin{itemize}
  \item First item
  \item Second item
  \item Third item with sub-items:
    \begin{itemize}
      \item Sub-item 1
      \item Sub-item 2
    \end{itemize}
\end{itemize}

\subsection{Numbered List}
\begin{enumerate}
  \item First step
  \item Second step
  \item Third step
\end{enumerate}

\section{Citations and References}
When you use the citation search feature, references will be automatically managed. For example, you might cite a paper like this: \cite{example2024}. All citations will be automatically added to your \texttt{cite.bib} file.

\section{Conclusion}
This template demonstrates basic \LaTeX{} features including mathematical formulas, tables, and lists. The bibliography system is now set up to work seamlessly with Bibby AI's citation management. Feel free to modify it for your needs!

\bibliographystyle{plain}
\bibliography{cite}

\end{document}\documentclass{article}
\usepackage[utf8]{inputenc}
\usepackage{booktabs}
\usepackage{amsmath}
\usepackage{natbib}

\title{My Document}
\author{Your Name}
\date{\today}

\begin{document}

\maketitle

\section{Introduction}
This is a sample document demonstrating various \LaTeX{} features including mathematical equations, tables, and formatting. You can cite papers using the citation search feature, and they will be automatically added to your \texttt{cite.bib} file.

\section{Mathematical Expressions}
Here are some example equations:

\subsection{Quadratic Formula}
The quadratic formula is given by:
\[x = \frac{-b \pm \sqrt{b^2 - 4ac}}{2a}\]

\subsection{Matrix Example}
A sample matrix:
\[
\begin{pmatrix}
a_{11} & a_{12} & a_{13} \\
a_{21} & a_{22} & a_{23} \\
a_{31} & a_{32} & a_{33}
\end{pmatrix}
\]

\section{Tables}
Below is a sample table showing data:

\begin{table}[h]
\centering
\begin{tabular}{@{}lrrr@{}}
\toprule
\textbf{Category} & \textbf{Value 1} & \textbf{Value 2} & \textbf{Total} \\
\midrule
Item A & 10.5 & 20.3 & 30.8 \\
Item B & 15.2 & 25.1 & 40.3 \\
Item C & 22.1 & 12.7 & 34.8 \\
\bottomrule
\end{tabular}
\caption{Sample Data Table}
\label{tab:sample}
\end{table}

\section{Lists}
\subsection{Bullet Points}
\begin{itemize}
  \item First item
  \item Second item
  \item Third item with sub-items:
    \begin{itemize}
      \item Sub-item 1
      \item Sub-item 2
    \end{itemize}
\end{itemize}

\subsection{Numbered List}
\begin{enumerate}
  \item First step
  \item Second step
  \item Third step
\end{enumerate}

\section{Citations and References}
When you use the citation search feature, references will be automatically managed. For example, you might cite a paper like this: \cite{example2024}. All citations will be automatically added to your \texttt{cite.bib} file.

\section{Conclusion}
This template demonstrates basic \LaTeX{} features including mathematical formulas, tables, and lists. The bibliography system is now set up to work seamlessly with Bibby AI's citation management. Feel free to modify it for your needs!

\bibliographystyle{plain}
\bibliography{cite}

\end{document}\documentclass{article}
\usepackage[utf8]{inputenc}
\usepackage{booktabs}
\usepackage{amsmath}
\usepackage{natbib}

\title{My Document}
\author{Your Name}
\date{\today}

\begin{document}

\maketitle

\section{Introduction}
This is a sample document demonstrating various \LaTeX{} features including mathematical equations, tables, and formatting. You can cite papers using the citation search feature, and they will be automatically added to your \texttt{cite.bib} file.

\section{Mathematical Expressions}
Here are some example equations:

\subsection{Quadratic Formula}
The quadratic formula is given by:
\[x = \frac{-b \pm \sqrt{b^2 - 4ac}}{2a}\]

\subsection{Matrix Example}
A sample matrix:
\[
\begin{pmatrix}
a_{11} & a_{12} & a_{13} \\
a_{21} & a_{22} & a_{23} \\
a_{31} & a_{32} & a_{33}
\end{pmatrix}
\]

\section{Tables}
Below is a sample table showing data:

\begin{table}[h]
\centering
\begin{tabular}{@{}lrrr@{}}
\toprule
\textbf{Category} & \textbf{Value 1} & \textbf{Value 2} & \textbf{Total} \\
\midrule
Item A & 10.5 & 20.3 & 30.8 \\
Item B & 15.2 & 25.1 & 40.3 \\
Item C & 22.1 & 12.7 & 34.8 \\
\bottomrule
\end{tabular}
\caption{Sample Data Table}
\label{tab:sample}
\end{table}

\section{Lists}
\subsection{Bullet Points}
\begin{itemize}
  \item First item
  \item Second item
  \item Third item with sub-items:
    \begin{itemize}
      \item Sub-item 1
      \item Sub-item 2
    \end{itemize}
\end{itemize}

\subsection{Numbered List}
\begin{enumerate}
  \item First step
  \item Second step
  \item Third step
\end{enumerate}

\section{Citations and References}
When you use the citation search feature, references will be automatically managed. For example, you might cite a paper like this: \cite{example2024}. All citations will be automatically added to your \texttt{cite.bib} file.

\section{Conclusion}
This template demonstrates basic \LaTeX{} features including mathematical formulas, tables, and lists. The bibliography system is now set up to work seamlessly with Bibby AI's citation management. Feel free to modify it for your needs!

\bibliographystyle{plain}
\bibliography{cite}

\end{document}\documentclass{article}
\usepackage[utf8]{inputenc}
\usepackage{booktabs}
\usepackage{amsmath}
\usepackage{natbib}

\title{My Document}
\author{Your Name}
\date{\today}

\begin{document}

\maketitle

\section{Introduction}
This is a sample document demonstrating various \LaTeX{} features including mathematical equations, tables, and formatting. You can cite papers using the citation search feature, and they will be automatically added to your \texttt{cite.bib} file.

\section{Mathematical Expressions}
Here are some example equations:

\subsection{Quadratic Formula}
The quadratic formula is given by:
\[x = \frac{-b \pm \sqrt{b^2 - 4ac}}{2a}\]

\subsection{Matrix Example}
A sample matrix:
\[
\begin{pmatrix}
a_{11} & a_{12} & a_{13} \\
a_{21} & a_{22} & a_{23} \\
a_{31} & a_{32} & a_{33}
\end{pmatrix}
\]

\section{Tables}
Below is a sample table showing data:

\begin{table}[h]
\centering
\begin{tabular}{@{}lrrr@{}}
\toprule
\textbf{Category} & \textbf{Value 1} & \textbf{Value 2} & \textbf{Total} \\
\midrule
Item A & 10.5 & 20.3 & 30.8 \\
Item B & 15.2 & 25.1 & 40.3 \\
Item C & 22.1 & 12.7 & 34.8 \\
\bottomrule
\end{tabular}
\caption{Sample Data Table}
\label{tab:sample}
\end{table}

\section{Lists}
\subsection{Bullet Points}
\begin{itemize}
  \item First item
  \item Second item
  \item Third item with sub-items:
    \begin{itemize}
      \item Sub-item 1
      \item Sub-item 2
    \end{itemize}
\end{itemize}

\subsection{Numbered List}
\begin{enumerate}
  \item First step
  \item Second step
  \item Third step
\end{enumerate}

\section{Citations and References}
When you use the citation search feature, references will be automatically managed. For example, you might cite a paper like this: \cite{example2024}. All citations will be automatically added to your \texttt{cite.bib} file.

\section{Conclusion}
This template demonstrates basic \LaTeX{} features including mathematical formulas, tables, and lists. The bibliography system is now set up to work seamlessly with Bibby AI's citation management. Feel free to modify it for your needs!

\bibliographystyle{plain}
\bibliography{cite}

\end{document}\documentclass{article}
\usepackage[utf8]{inputenc}
\usepackage{booktabs}
\usepackage{amsmath}
\usepackage{natbib}

\title{My Document}
\author{Your Name}
\date{\today}

\begin{document}

\maketitle

\section{Introduction}
This is a sample document demonstrating various \LaTeX{} features including mathematical equations, tables, and formatting. You can cite papers using the citation search feature, and they will be automatically added to your \texttt{cite.bib} file.

\section{Mathematical Expressions}
Here are some example equations:

\subsection{Quadratic Formula}
The quadratic formula is given by:
\[x = \frac{-b \pm \sqrt{b^2 - 4ac}}{2a}\]

\subsection{Matrix Example}
A sample matrix:
\[
\begin{pmatrix}
a_{11} & a_{12} & a_{13} \\
a_{21} & a_{22} & a_{23} \\
a_{31} & a_{32} & a_{33}
\end{pmatrix}
\]

\section{Tables}
Below is a sample table showing data:

\begin{table}[h]
\centering
\begin{tabular}{@{}lrrr@{}}
\toprule
\textbf{Category} & \textbf{Value 1} & \textbf{Value 2} & \textbf{Total} \\
\midrule
Item A & 10.5 & 20.3 & 30.8 \\
Item B & 15.2 & 25.1 & 40.3 \\
Item C & 22.1 & 12.7 & 34.8 \\
\bottomrule
\end{tabular}
\caption{Sample Data Table}
\label{tab:sample}
\end{table}

\section{Lists}
\subsection{Bullet Points}
\begin{itemize}
  \item First item
  \item Second item
  \item Third item with sub-items:
    \begin{itemize}
      \item Sub-item 1
      \item Sub-item 2
    \end{itemize}
\end{itemize}

\subsection{Numbered List}
\begin{enumerate}
  \item First step
  \item Second step
  \item Third step
\end{enumerate}

\section{Citations and References}
When you use the citation search feature, references will be automatically managed. For example, you might cite a paper like this: \cite{example2024}. All citations will be automatically added to your \texttt{cite.bib} file.

\section{Conclusion}
This template demonstrates basic \LaTeX{} features including mathematical formulas, tables, and lists. The bibliography system is now set up to work seamlessly with Bibby AI's citation management. Feel free to modify it for your needs!

\bibliographystyle{plain}
\bibliography{cite}

\end{document}\documentclass{article}
\usepackage[utf8]{inputenc}
\usepackage{booktabs}
\usepackage{amsmath}
\usepackage{natbib}

\title{My Document}
\author{Your Name}
\date{\today}

\begin{document}

\maketitle

\section{Introduction}
This is a sample document demonstrating various \LaTeX{} features including mathematical equations, tables, and formatting. You can cite papers using the citation search feature, and they will be automatically added to your \texttt{cite.bib} file.

\section{Mathematical Expressions}
Here are some example equations:

\subsection{Quadratic Formula}
The quadratic formula is given by:
\[x = \frac{-b \pm \sqrt{b^2 - 4ac}}{2a}\]

\subsection{Matrix Example}
A sample matrix:
\[
\begin{pmatrix}
a_{11} & a_{12} & a_{13} \\
a_{21} & a_{22} & a_{23} \\
a_{31} & a_{32} & a_{33}
\end{pmatrix}
\]

\section{Tables}
Below is a sample table showing data:

\begin{table}[h]
\centering
\begin{tabular}{@{}lrrr@{}}
\toprule
\textbf{Category} & \textbf{Value 1} & \textbf{Value 2} & \textbf{Total} \\
\midrule
Item A & 10.5 & 20.3 & 30.8 \\
Item B & 15.2 & 25.1 & 40.3 \\
Item C & 22.1 & 12.7 & 34.8 \\
\bottomrule
\end{tabular}
\caption{Sample Data Table}
\label{tab:sample}
\end{table}

\section{Lists}
\subsection{Bullet Points}
\begin{itemize}
  \item First item
  \item Second item
  \item Third item with sub-items:
    \begin{itemize}
      \item Sub-item 1
      \item Sub-item 2
    \end{itemize}
\end{itemize}

\subsection{Numbered List}
\begin{enumerate}
  \item First step
  \item Second step
  \item Third step
\end{enumerate}

\section{Citations and References}
When you use the citation search feature, references will be automatically managed. For example, you might cite a paper like this: \cite{example2024}. All citations will be automatically added to your \texttt{cite.bib} file.

\section{Conclusion}
This template demonstrates basic \LaTeX{} features including mathematical formulas, tables, and lists. The bibliography system is now set up to work seamlessly with Bibby AI's citation management. Feel free to modify it for your needs!

\bibliographystyle{plain}
\bibliography{cite}

\end{document}\documentclass{article}
\usepackage[utf8]{inputenc}
\usepackage{booktabs}
\usepackage{amsmath}
\usepackage{natbib}

\title{My Document}
\author{Your Name}
\date{\today}

\begin{document}

\maketitle

\section{Introduction}
This is a sample document demonstrating various \LaTeX{} features including mathematical equations, tables, and formatting. You can cite papers using the citation search feature, and they will be automatically added to your \texttt{cite.bib} file.

\section{Mathematical Expressions}
Here are some example equations:

\subsection{Quadratic Formula}
The quadratic formula is given by:
\[x = \frac{-b \pm \sqrt{b^2 - 4ac}}{2a}\]

\subsection{Matrix Example}
A sample matrix:
\[
\begin{pmatrix}
a_{11} & a_{12} & a_{13} \\
a_{21} & a_{22} & a_{23} \\
a_{31} & a_{32} & a_{33}
\end{pmatrix}
\]

\section{Tables}
Below is a sample table showing data:

\begin{table}[h]
\centering
\begin{tabular}{@{}lrrr@{}}
\toprule
\textbf{Category} & \textbf{Value 1} & \textbf{Value 2} & \textbf{Total} \\
\midrule
Item A & 10.5 & 20.3 & 30.8 \\
Item B & 15.2 & 25.1 & 40.3 \\
Item C & 22.1 & 12.7 & 34.8 \\
\bottomrule
\end{tabular}
\caption{Sample Data Table}
\label{tab:sample}
\end{table}

\section{Lists}
\subsection{Bullet Points}
\begin{itemize}
  \item First item
  \item Second item
  \item Third item with sub-items:
    \begin{itemize}
      \item Sub-item 1
      \item Sub-item 2
    \end{itemize}
\end{itemize}

\subsection{Numbered List}
\begin{enumerate}
  \item First step
  \item Second step
  \item Third step
\end{enumerate}

\section{Citations and References}
When you use the citation search feature, references will be automatically managed. For example, you might cite a paper like this: \cite{example2024}. All citations will be automatically added to your \texttt{cite.bib} file.

\section{Conclusion}
This template demonstrates basic \LaTeX{} features including mathematical formulas, tables, and lists. The bibliography system is now set up to work seamlessly with Bibby AI's citation management. Feel free to modify it for your needs!

\bibliographystyle{plain}
\bibliography{cite}

\end{document}\documentclass{article}
\usepackage[utf8]{inputenc}
\usepackage{booktabs}
\usepackage{amsmath}
\usepackage{natbib}

\title{My Document}
\author{Your Name}
\date{\today}

\begin{document}

\maketitle

\section{Introduction}
This is a sample document demonstrating various \LaTeX{} features including mathematical equations, tables, and formatting. You can cite papers using the citation search feature, and they will be automatically added to your \texttt{cite.bib} file.

\section{Mathematical Expressions}
Here are some example equations:

\subsection{Quadratic Formula}
The quadratic formula is given by:
\[x = \frac{-b \pm \sqrt{b^2 - 4ac}}{2a}\]

\subsection{Matrix Example}
A sample matrix:
\[
\begin{pmatrix}
a_{11} & a_{12} & a_{13} \\
a_{21} & a_{22} & a_{23} \\
a_{31} & a_{32} & a_{33}
\end{pmatrix}
\]

\section{Tables}
Below is a sample table showing data:

\begin{table}[h]
\centering
\begin{tabular}{@{}lrrr@{}}
\toprule
\textbf{Category} & \textbf{Value 1} & \textbf{Value 2} & \textbf{Total} \\
\midrule
Item A & 10.5 & 20.3 & 30.8 \\
Item B & 15.2 & 25.1 & 40.3 \\
Item C & 22.1 & 12.7 & 34.8 \\
\bottomrule
\end{tabular}
\caption{Sample Data Table}
\label{tab:sample}
\end{table}

\section{Lists}
\subsection{Bullet Points}
\begin{itemize}
  \item First item
  \item Second item
  \item Third item with sub-items:
    \begin{itemize}
      \item Sub-item 1
      \item Sub-item 2
    \end{itemize}
\end{itemize}

\subsection{Numbered List}
\begin{enumerate}
  \item First step
  \item Second step
  \item Third step
\end{enumerate}

\section{Citations and References}
When you use the citation search feature, references will be automatically managed. For example, you might cite a paper like this: \cite{example2024}. All citations will be automatically added to your \texttt{cite.bib} file.

\section{Conclusion}
This template demonstrates basic \LaTeX{} features including mathematical formulas, tables, and lists. The bibliography system is now set up to work seamlessly with Bibby AI's citation management. Feel free to modify it for your needs!

\bibliographystyle{plain}
\bibliography{cite}

\end{document}\documentclass{article}
\usepackage[utf8]{inputenc}
\usepackage{booktabs}
\usepackage{amsmath}
\usepackage{natbib}

\title{My Document}
\author{Your Name}
\date{\today}

\begin{document}

\maketitle

\section{Introduction}
This is a sample document demonstrating various \LaTeX{} features including mathematical equations, tables, and formatting. You can cite papers using the citation search feature, and they will be automatically added to your \texttt{cite.bib} file.

\section{Mathematical Expressions}
Here are some example equations:

\subsection{Quadratic Formula}
The quadratic formula is given by:
\[x = \frac{-b \pm \sqrt{b^2 - 4ac}}{2a}\]

\subsection{Matrix Example}
A sample matrix:
\[
\begin{pmatrix}
a_{11} & a_{12} & a_{13} \\
a_{21} & a_{22} & a_{23} \\
a_{31} & a_{32} & a_{33}
\end{pmatrix}
\]

\section{Tables}
Below is a sample table showing data:

\begin{table}[h]
\centering
\begin{tabular}{@{}lrrr@{}}
\toprule
\textbf{Category} & \textbf{Value 1} & \textbf{Value 2} & \textbf{Total} \\
\midrule
Item A & 10.5 & 20.3 & 30.8 \\
Item B & 15.2 & 25.1 & 40.3 \\
Item C & 22.1 & 12.7 & 34.8 \\
\bottomrule
\end{tabular}
\caption{Sample Data Table}
\label{tab:sample}
\end{table}

\section{Lists}
\subsection{Bullet Points}
\begin{itemize}
  \item First item
  \item Second item
  \item Third item with sub-items:
    \begin{itemize}
      \item Sub-item 1
      \item Sub-item 2
    \end{itemize}
\end{itemize}

\subsection{Numbered List}
\begin{enumerate}
  \item First step
  \item Second step
  \item Third step
\end{enumerate}

\section{Citations and References}
When you use the citation search feature, references will be automatically managed. For example, you might cite a paper like this: \cite{example2024}. All citations will be automatically added to your \texttt{cite.bib} file.

\section{Conclusion}
This template demonstrates basic \LaTeX{} features including mathematical formulas, tables, and lists. The bibliography system is now set up to work seamlessly with Bibby AI's citation management. Feel free to modify it for your needs!

\bibliographystyle{plain}
\bibliography{cite}

\end{document}\documentclass{article}
\usepackage[utf8]{inputenc}
\usepackage{booktabs}
\usepackage{amsmath}
\usepackage{natbib}

\title{My Document}
\author{Your Name}
\date{\today}

\begin{document}

\maketitle

\section{Introduction}
This is a sample document demonstrating various \LaTeX{} features including mathematical equations, tables, and formatting. You can cite papers using the citation search feature, and they will be automatically added to your \texttt{cite.bib} file.

\section{Mathematical Expressions}
Here are some example equations:

\subsection{Quadratic Formula}
The quadratic formula is given by:
\[x = \frac{-b \pm \sqrt{b^2 - 4ac}}{2a}\]

\subsection{Matrix Example}
A sample matrix:
\[
\begin{pmatrix}
a_{11} & a_{12} & a_{13} \\
a_{21} & a_{22} & a_{23} \\
a_{31} & a_{32} & a_{33}
\end{pmatrix}
\]

\section{Tables}
Below is a sample table showing data:

\begin{table}[h]
\centering
\begin{tabular}{@{}lrrr@{}}
\toprule
\textbf{Category} & \textbf{Value 1} & \textbf{Value 2} & \textbf{Total} \\
\midrule
Item A & 10.5 & 20.3 & 30.8 \\
Item B & 15.2 & 25.1 & 40.3 \\
Item C & 22.1 & 12.7 & 34.8 \\
\bottomrule
\end{tabular}
\caption{Sample Data Table}
\label{tab:sample}
\end{table}

\section{Lists}
\subsection{Bullet Points}
\begin{itemize}
  \item First item
  \item Second item
  \item Third item with sub-items:
    \begin{itemize}
      \item Sub-item 1
      \item Sub-item 2
    \end{itemize}
\end{itemize}

\subsection{Numbered List}
\begin{enumerate}
  \item First step
  \item Second step
  \item Third step
\end{enumerate}

\section{Citations and References}
When you use the citation search feature, references will be automatically managed. For example, you might cite a paper like this: \cite{example2024}. All citations will be automatically added to your \texttt{cite.bib} file.

\section{Conclusion}
This template demonstrates basic \LaTeX{} features including mathematical formulas, tables, and lists. The bibliography system is now set up to work seamlessly with Bibby AI's citation management. Feel free to modify it for your needs!

\bibliographystyle{plain}
\bibliography{cite}

\end{document}\documentclass{article}
\usepackage[utf8]{inputenc}
\usepackage{booktabs}
\usepackage{amsmath}
\usepackage{natbib}

\title{My Document}
\author{Your Name}
\date{\today}

\begin{document}

\maketitle

\section{Introduction}
This is a sample document demonstrating various \LaTeX{} features including mathematical equations, tables, and formatting. You can cite papers using the citation search feature, and they will be automatically added to your \texttt{cite.bib} file.

\section{Mathematical Expressions}
Here are some example equations:

\subsection{Quadratic Formula}
The quadratic formula is given by:
\[x = \frac{-b \pm \sqrt{b^2 - 4ac}}{2a}\]

\subsection{Matrix Example}
A sample matrix:
\[
\begin{pmatrix}
a_{11} & a_{12} & a_{13} \\
a_{21} & a_{22} & a_{23} \\
a_{31} & a_{32} & a_{33}
\end{pmatrix}
\]

\section{Tables}
Below is a sample table showing data:

\begin{table}[h]
\centering
\begin{tabular}{@{}lrrr@{}}
\toprule
\textbf{Category} & \textbf{Value 1} & \textbf{Value 2} & \textbf{Total} \\
\midrule
Item A & 10.5 & 20.3 & 30.8 \\
Item B & 15.2 & 25.1 & 40.3 \\
Item C & 22.1 & 12.7 & 34.8 \\
\bottomrule
\end{tabular}
\caption{Sample Data Table}
\label{tab:sample}
\end{table}

\section{Lists}
\subsection{Bullet Points}
\begin{itemize}
  \item First item
  \item Second item
  \item Third item with sub-items:
    \begin{itemize}
      \item Sub-item 1
      \item Sub-item 2
    \end{itemize}
\end{itemize}

\subsection{Numbered List}
\begin{enumerate}
  \item First step
  \item Second step
  \item Third step
\end{enumerate}

\section{Citations and References}
When you use the citation search feature, references will be automatically managed. For example, you might cite a paper like this: \cite{example2024}. All citations will be automatically added to your \texttt{cite.bib} file.

\section{Conclusion}
This template demonstrates basic \LaTeX{} features including mathematical formulas, tables, and lists. The bibliography system is now set up to work seamlessly with Bibby AI's citation management. Feel free to modify it for your needs!

\bibliographystyle{plain}
\bibliography{cite}

\end{document}\documentclass{article}
\usepackage[utf8]{inputenc}
\usepackage{booktabs}
\usepackage{amsmath}
\usepackage{natbib}

\title{My Document}
\author{Your Name}
\date{\today}

\begin{document}

\maketitle

\section{Introduction}
This is a sample document demonstrating various \LaTeX{} features including mathematical equations, tables, and formatting. You can cite papers using the citation search feature, and they will be automatically added to your \texttt{cite.bib} file.

\section{Mathematical Expressions}
Here are some example equations:

\subsection{Quadratic Formula}
The quadratic formula is given by:
\[x = \frac{-b \pm \sqrt{b^2 - 4ac}}{2a}\]

\subsection{Matrix Example}
A sample matrix:
\[
\begin{pmatrix}
a_{11} & a_{12} & a_{13} \\
a_{21} & a_{22} & a_{23} \\
a_{31} & a_{32} & a_{33}
\end{pmatrix}
\]

\section{Tables}
Below is a sample table showing data:

\begin{table}[h]
\centering
\begin{tabular}{@{}lrrr@{}}
\toprule
\textbf{Category} & \textbf{Value 1} & \textbf{Value 2} & \textbf{Total} \\
\midrule
Item A & 10.5 & 20.3 & 30.8 \\
Item B & 15.2 & 25.1 & 40.3 \\
Item C & 22.1 & 12.7 & 34.8 \\
\bottomrule
\end{tabular}
\caption{Sample Data Table}
\label{tab:sample}
\end{table}

\section{Lists}
\subsection{Bullet Points}
\begin{itemize}
  \item First item
  \item Second item
  \item Third item with sub-items:
    \begin{itemize}
      \item Sub-item 1
      \item Sub-item 2
    \end{itemize}
\end{itemize}

\subsection{Numbered List}
\begin{enumerate}
  \item First step
  \item Second step
  \item Third step
\end{enumerate}

\section{Citations and References}
When you use the citation search feature, references will be automatically managed. For example, you might cite a paper like this: \cite{example2024}. All citations will be automatically added to your \texttt{cite.bib} file.

\section{Conclusion}
This template demonstrates basic \LaTeX{} features including mathematical formulas, tables, and lists. The bibliography system is now set up to work seamlessly with Bibby AI's citation management. Feel free to modify it for your needs!

\bibliographystyle{plain}
\bibliography{cite}

\end{document}\documentclass{article}
\usepackage[utf8]{inputenc}
\usepackage{booktabs}
\usepackage{amsmath}
\usepackage{natbib}

\title{My Document}
\author{Your Name}
\date{\today}

\begin{document}

\maketitle

\section{Introduction}
This is a sample document demonstrating various \LaTeX{} features including mathematical equations, tables, and formatting. You can cite papers using the citation search feature, and they will be automatically added to your \texttt{cite.bib} file.

\section{Mathematical Expressions}
Here are some example equations:

\subsection{Quadratic Formula}
The quadratic formula is given by:
\[x = \frac{-b \pm \sqrt{b^2 - 4ac}}{2a}\]

\subsection{Matrix Example}
A sample matrix:
\[
\begin{pmatrix}
a_{11} & a_{12} & a_{13} \\
a_{21} & a_{22} & a_{23} \\
a_{31} & a_{32} & a_{33}
\end{pmatrix}
\]

\section{Tables}
Below is a sample table showing data:

\begin{table}[h]
\centering
\begin{tabular}{@{}lrrr@{}}
\toprule
\textbf{Category} & \textbf{Value 1} & \textbf{Value 2} & \textbf{Total} \\
\midrule
Item A & 10.5 & 20.3 & 30.8 \\
Item B & 15.2 & 25.1 & 40.3 \\
Item C & 22.1 & 12.7 & 34.8 \\
\bottomrule
\end{tabular}
\caption{Sample Data Table}
\label{tab:sample}
\end{table}

\section{Lists}
\subsection{Bullet Points}
\begin{itemize}
  \item First item
  \item Second item
  \item Third item with sub-items:
    \begin{itemize}
      \item Sub-item 1
      \item Sub-item 2
    \end{itemize}
\end{itemize}

\subsection{Numbered List}
\begin{enumerate}
  \item First step
  \item Second step
  \item Third step
\end{enumerate}

\section{Citations and References}
When you use the citation search feature, references will be automatically managed. For example, you might cite a paper like this: \cite{example2024}. All citations will be automatically added to your \texttt{cite.bib} file.

\section{Conclusion}
This template demonstrates basic \LaTeX{} features including mathematical formulas, tables, and lists. The bibliography system is now set up to work seamlessly with Bibby AI's citation management. Feel free to modify it for your needs!

\bibliographystyle{plain}
\bibliography{cite}

\end{document}\documentclass{article}
\usepackage[utf8]{inputenc}
\usepackage{booktabs}
\usepackage{amsmath}
\usepackage{natbib}

\title{My Document}
\author{Your Name}
\date{\today}

\begin{document}

\maketitle

\section{Introduction}
This is a sample document demonstrating various \LaTeX{} features including mathematical equations, tables, and formatting. You can cite papers using the citation search feature, and they will be automatically added to your \texttt{cite.bib} file.

\section{Mathematical Expressions}
Here are some example equations:

\subsection{Quadratic Formula}
The quadratic formula is given by:
\[x = \frac{-b \pm \sqrt{b^2 - 4ac}}{2a}\]

\subsection{Matrix Example}
A sample matrix:
\[
\begin{pmatrix}
a_{11} & a_{12} & a_{13} \\
a_{21} & a_{22} & a_{23} \\
a_{31} & a_{32} & a_{33}
\end{pmatrix}
\]

\section{Tables}
Below is a sample table showing data:

\begin{table}[h]
\centering
\begin{tabular}{@{}lrrr@{}}
\toprule
\textbf{Category} & \textbf{Value 1} & \textbf{Value 2} & \textbf{Total} \\
\midrule
Item A & 10.5 & 20.3 & 30.8 \\
Item B & 15.2 & 25.1 & 40.3 \\
Item C & 22.1 & 12.7 & 34.8 \\
\bottomrule
\end{tabular}
\caption{Sample Data Table}
\label{tab:sample}
\end{table}

\section{Lists}
\subsection{Bullet Points}
\begin{itemize}
  \item First item
  \item Second item
  \item Third item with sub-items:
    \begin{itemize}
      \item Sub-item 1
      \item Sub-item 2
    \end{itemize}
\end{itemize}

\subsection{Numbered List}
\begin{enumerate}
  \item First step
  \item Second step
  \item Third step
\end{enumerate}

\section{Citations and References}
When you use the citation search feature, references will be automatically managed. For example, you might cite a paper like this: \cite{example2024}. All citations will be automatically added to your \texttt{cite.bib} file.

\section{Conclusion}
This template demonstrates basic \LaTeX{} features including mathematical formulas, tables, and lists. The bibliography system is now set up to work seamlessly with Bibby AI's citation management. Feel free to modify it for your needs!

\bibliographystyle{plain}
\bibliography{cite}

\end{document}\documentclass{article}
\usepackage[utf8]{inputenc}
\usepackage{booktabs}
\usepackage{amsmath}
\usepackage{natbib}

\title{My Document}
\author{Your Name}
\date{\today}

\begin{document}

\maketitle

\section{Introduction}
This is a sample document demonstrating various \LaTeX{} features including mathematical equations, tables, and formatting. You can cite papers using the citation search feature, and they will be automatically added to your \texttt{cite.bib} file.

\section{Mathematical Expressions}
Here are some example equations:

\subsection{Quadratic Formula}
The quadratic formula is given by:
\[x = \frac{-b \pm \sqrt{b^2 - 4ac}}{2a}\]

\subsection{Matrix Example}
A sample matrix:
\[
\begin{pmatrix}
a_{11} & a_{12} & a_{13} \\
a_{21} & a_{22} & a_{23} \\
a_{31} & a_{32} & a_{33}
\end{pmatrix}
\]

\section{Tables}
Below is a sample table showing data:

\begin{table}[h]
\centering
\begin{tabular}{@{}lrrr@{}}
\toprule
\textbf{Category} & \textbf{Value 1} & \textbf{Value 2} & \textbf{Total} \\
\midrule
Item A & 10.5 & 20.3 & 30.8 \\
Item B & 15.2 & 25.1 & 40.3 \\
Item C & 22.1 & 12.7 & 34.8 \\
\bottomrule
\end{tabular}
\caption{Sample Data Table}
\label{tab:sample}
\end{table}

\section{Lists}
\subsection{Bullet Points}
\begin{itemize}
  \item First item
  \item Second item
  \item Third item with sub-items:
    \begin{itemize}
      \item Sub-item 1
      \item Sub-item 2
    \end{itemize}
\end{itemize}

\subsection{Numbered List}
\begin{enumerate}
  \item First step
  \item Second step
  \item Third step
\end{enumerate}

\section{Citations and References}
When you use the citation search feature, references will be automatically managed. For example, you might cite a paper like this: \cite{example2024}. All citations will be automatically added to your \texttt{cite.bib} file.

\section{Conclusion}
This template demonstrates basic \LaTeX{} features including mathematical formulas, tables, and lists. The bibliography system is now set up to work seamlessly with Bibby AI's citation management. Feel free to modify it for your needs!

\bibliographystyle{plain}
\bibliography{cite}

\end{document}\documentclass{article}
\usepackage[utf8]{inputenc}
\usepackage{booktabs}
\usepackage{amsmath}
\usepackage{natbib}

\title{My Document}
\author{Your Name}
\date{\today}

\begin{document}

\maketitle

\section{Introduction}
This is a sample document demonstrating various \LaTeX{} features including mathematical equations, tables, and formatting. You can cite papers using the citation search feature, and they will be automatically added to your \texttt{cite.bib} file.

\section{Mathematical Expressions}
Here are some example equations:

\subsection{Quadratic Formula}
The quadratic formula is given by:
\[x = \frac{-b \pm \sqrt{b^2 - 4ac}}{2a}\]

\subsection{Matrix Example}
A sample matrix:
\[
\begin{pmatrix}
a_{11} & a_{12} & a_{13} \\
a_{21} & a_{22} & a_{23} \\
a_{31} & a_{32} & a_{33}
\end{pmatrix}
\]

\section{Tables}
Below is a sample table showing data:

\begin{table}[h]
\centering
\begin{tabular}{@{}lrrr@{}}
\toprule
\textbf{Category} & \textbf{Value 1} & \textbf{Value 2} & \textbf{Total} \\
\midrule
Item A & 10.5 & 20.3 & 30.8 \\
Item B & 15.2 & 25.1 & 40.3 \\
Item C & 22.1 & 12.7 & 34.8 \\
\bottomrule
\end{tabular}
\caption{Sample Data Table}
\label{tab:sample}
\end{table}

\section{Lists}
\subsection{Bullet Points}
\begin{itemize}
  \item First item
  \item Second item
  \item Third item with sub-items:
    \begin{itemize}
      \item Sub-item 1
      \item Sub-item 2
    \end{itemize}
\end{itemize}

\subsection{Numbered List}
\begin{enumerate}
  \item First step
  \item Second step
  \item Third step
\end{enumerate}

\section{Citations and References}
When you use the citation search feature, references will be automatically managed. For example, you might cite a paper like this: \cite{example2024}. All citations will be automatically added to your \texttt{cite.bib} file.

\section{Conclusion}
This template demonstrates basic \LaTeX{} features including mathematical formulas, tables, and lists. The bibliography system is now set up to work seamlessly with Bibby AI's citation management. Feel free to modify it for your needs!

\bibliographystyle{plain}
\bibliography{cite}

\end{document}\documentclass{article}
\usepackage[utf8]{inputenc}
\usepackage{booktabs}
\usepackage{amsmath}
\usepackage{natbib}

\title{My Document}
\author{Your Name}
\date{\today}

\begin{document}

\maketitle

\section{Introduction}
This is a sample document demonstrating various \LaTeX{} features including mathematical equations, tables, and formatting. You can cite papers using the citation search feature, and they will be automatically added to your \texttt{cite.bib} file.

\section{Mathematical Expressions}
Here are some example equations:

\subsection{Quadratic Formula}
The quadratic formula is given by:
\[x = \frac{-b \pm \sqrt{b^2 - 4ac}}{2a}\]

\subsection{Matrix Example}
A sample matrix:
\[
\begin{pmatrix}
a_{11} & a_{12} & a_{13} \\
a_{21} & a_{22} & a_{23} \\
a_{31} & a_{32} & a_{33}
\end{pmatrix}
\]

\section{Tables}
Below is a sample table showing data:

\begin{table}[h]
\centering
\begin{tabular}{@{}lrrr@{}}
\toprule
\textbf{Category} & \textbf{Value 1} & \textbf{Value 2} & \textbf{Total} \\
\midrule
Item A & 10.5 & 20.3 & 30.8 \\
Item B & 15.2 & 25.1 & 40.3 \\
Item C & 22.1 & 12.7 & 34.8 \\
\bottomrule
\end{tabular}
\caption{Sample Data Table}
\label{tab:sample}
\end{table}

\section{Lists}
\subsection{Bullet Points}
\begin{itemize}
  \item First item
  \item Second item
  \item Third item with sub-items:
    \begin{itemize}
      \item Sub-item 1
      \item Sub-item 2
    \end{itemize}
\end{itemize}

\subsection{Numbered List}
\begin{enumerate}
  \item First step
  \item Second step
  \item Third step
\end{enumerate}

\section{Citations and References}
When you use the citation search feature, references will be automatically managed. For example, you might cite a paper like this: \cite{example2024}. All citations will be automatically added to your \texttt{cite.bib} file.

\section{Conclusion}
This template demonstrates basic \LaTeX{} features including mathematical formulas, tables, and lists. The bibliography system is now set up to work seamlessly with Bibby AI's citation management. Feel free to modify it for your needs!

\bibliographystyle{plain}
\bibliography{cite}

\end{document}\documentclass{article}
\usepackage[utf8]{inputenc}
\usepackage{booktabs}
\usepackage{amsmath}
\usepackage{natbib}

\title{My Document}
\author{Your Name}
\date{\today}

\begin{document}

\maketitle

\section{Introduction}
This is a sample document demonstrating various \LaTeX{} features including mathematical equations, tables, and formatting. You can cite papers using the citation search feature, and they will be automatically added to your \texttt{cite.bib} file.

\section{Mathematical Expressions}
Here are some example equations:

\subsection{Quadratic Formula}
The quadratic formula is given by:
\[x = \frac{-b \pm \sqrt{b^2 - 4ac}}{2a}\]

\subsection{Matrix Example}
A sample matrix:
\[
\begin{pmatrix}
a_{11} & a_{12} & a_{13} \\
a_{21} & a_{22} & a_{23} \\
a_{31} & a_{32} & a_{33}
\end{pmatrix}
\]

\section{Tables}
Below is a sample table showing data:

\begin{table}[h]
\centering
\begin{tabular}{@{}lrrr@{}}
\toprule
\textbf{Category} & \textbf{Value 1} & \textbf{Value 2} & \textbf{Total} \\
\midrule
Item A & 10.5 & 20.3 & 30.8 \\
Item B & 15.2 & 25.1 & 40.3 \\
Item C & 22.1 & 12.7 & 34.8 \\
\bottomrule
\end{tabular}
\caption{Sample Data Table}
\label{tab:sample}
\end{table}

\section{Lists}
\subsection{Bullet Points}
\begin{itemize}
  \item First item
  \item Second item
  \item Third item with sub-items:
    \begin{itemize}
      \item Sub-item 1
      \item Sub-item 2
    \end{itemize}
\end{itemize}

\subsection{Numbered List}
\begin{enumerate}
  \item First step
  \item Second step
  \item Third step
\end{enumerate}

\section{Citations and References}
When you use the citation search feature, references will be automatically managed. For example, you might cite a paper like this: \cite{example2024}. All citations will be automatically added to your \texttt{cite.bib} file.

\section{Conclusion}
This template demonstrates basic \LaTeX{} features including mathematical formulas, tables, and lists. The bibliography system is now set up to work seamlessly with Bibby AI's citation management. Feel free to modify it for your needs!

\bibliographystyle{plain}
\bibliography{cite}

\end{document}\documentclass{article}
\usepackage[utf8]{inputenc}
\usepackage{booktabs}
\usepackage{amsmath}
\usepackage{natbib}

\title{My Document}
\author{Your Name}
\date{\today}

\begin{document}

\maketitle

\section{Introduction}
This is a sample document demonstrating various \LaTeX{} features including mathematical equations, tables, and formatting. You can cite papers using the citation search feature, and they will be automatically added to your \texttt{cite.bib} file.

\section{Mathematical Expressions}
Here are some example equations:

\subsection{Quadratic Formula}
The quadratic formula is given by:
\[x = \frac{-b \pm \sqrt{b^2 - 4ac}}{2a}\]

\subsection{Matrix Example}
A sample matrix:
\[
\begin{pmatrix}
a_{11} & a_{12} & a_{13} \\
a_{21} & a_{22} & a_{23} \\
a_{31} & a_{32} & a_{33}
\end{pmatrix}
\]

\section{Tables}
Below is a sample table showing data:

\begin{table}[h]
\centering
\begin{tabular}{@{}lrrr@{}}
\toprule
\textbf{Category} & \textbf{Value 1} & \textbf{Value 2} & \textbf{Total} \\
\midrule
Item A & 10.5 & 20.3 & 30.8 \\
Item B & 15.2 & 25.1 & 40.3 \\
Item C & 22.1 & 12.7 & 34.8 \\
\bottomrule
\end{tabular}
\caption{Sample Data Table}
\label{tab:sample}
\end{table}

\section{Lists}
\subsection{Bullet Points}
\begin{itemize}
  \item First item
  \item Second item
  \item Third item with sub-items:
    \begin{itemize}
      \item Sub-item 1
      \item Sub-item 2
    \end{itemize}
\end{itemize}

\subsection{Numbered List}
\begin{enumerate}
  \item First step
  \item Second step
  \item Third step
\end{enumerate}

\section{Citations and References}
When you use the citation search feature, references will be automatically managed. For example, you might cite a paper like this: \cite{example2024}. All citations will be automatically added to your \texttt{cite.bib} file.

\section{Conclusion}
This template demonstrates basic \LaTeX{} features including mathematical formulas, tables, and lists. The bibliography system is now set up to work seamlessly with Bibby AI's citation management. Feel free to modify it for your needs!

\bibliographystyle{plain}
\bibliography{cite}

\end{document}\documentclass{article}
\usepackage[utf8]{inputenc}
\usepackage{booktabs}
\usepackage{amsmath}
\usepackage{natbib}

\title{My Document}
\author{Your Name}
\date{\today}

\begin{document}

\maketitle

\section{Introduction}
This is a sample document demonstrating various \LaTeX{} features including mathematical equations, tables, and formatting. You can cite papers using the citation search feature, and they will be automatically added to your \texttt{cite.bib} file.

\section{Mathematical Expressions}
Here are some example equations:

\subsection{Quadratic Formula}
The quadratic formula is given by:
\[x = \frac{-b \pm \sqrt{b^2 - 4ac}}{2a}\]

\subsection{Matrix Example}
A sample matrix:
\[
\begin{pmatrix}
a_{11} & a_{12} & a_{13} \\
a_{21} & a_{22} & a_{23} \\
a_{31} & a_{32} & a_{33}
\end{pmatrix}
\]

\section{Tables}
Below is a sample table showing data:

\begin{table}[h]
\centering
\begin{tabular}{@{}lrrr@{}}
\toprule
\textbf{Category} & \textbf{Value 1} & \textbf{Value 2} & \textbf{Total} \\
\midrule
Item A & 10.5 & 20.3 & 30.8 \\
Item B & 15.2 & 25.1 & 40.3 \\
Item C & 22.1 & 12.7 & 34.8 \\
\bottomrule
\end{tabular}
\caption{Sample Data Table}
\label{tab:sample}
\end{table}

\section{Lists}
\subsection{Bullet Points}
\begin{itemize}
  \item First item
  \item Second item
  \item Third item with sub-items:
    \begin{itemize}
      \item Sub-item 1
      \item Sub-item 2
    \end{itemize}
\end{itemize}

\subsection{Numbered List}
\begin{enumerate}
  \item First step
  \item Second step
  \item Third step
\end{enumerate}

\section{Citations and References}
When you use the citation search feature, references will be automatically managed. For example, you might cite a paper like this: \cite{example2024}. All citations will be automatically added to your \texttt{cite.bib} file.

\section{Conclusion}
This template demonstrates basic \LaTeX{} features including mathematical formulas, tables, and lists. The bibliography system is now set up to work seamlessly with Bibby AI's citation management. Feel free to modify it for your needs!

\bibliographystyle{plain}
\bibliography{cite}

\end{document}\documentclass{article}
\usepackage[utf8]{inputenc}
\usepackage{booktabs}
\usepackage{amsmath}
\usepackage{natbib}

\title{My Document}
\author{Your Name}
\date{\today}

\begin{document}

\maketitle

\section{Introduction}
This is a sample document demonstrating various \LaTeX{} features including mathematical equations, tables, and formatting. You can cite papers using the citation search feature, and they will be automatically added to your \texttt{cite.bib} file.

\section{Mathematical Expressions}
Here are some example equations:

\subsection{Quadratic Formula}
The quadratic formula is given by:
\[x = \frac{-b \pm \sqrt{b^2 - 4ac}}{2a}\]

\subsection{Matrix Example}
A sample matrix:
\[
\begin{pmatrix}
a_{11} & a_{12} & a_{13} \\
a_{21} & a_{22} & a_{23} \\
a_{31} & a_{32} & a_{33}
\end{pmatrix}
\]

\section{Tables}
Below is a sample table showing data:

\begin{table}[h]
\centering
\begin{tabular}{@{}lrrr@{}}
\toprule
\textbf{Category} & \textbf{Value 1} & \textbf{Value 2} & \textbf{Total} \\
\midrule
Item A & 10.5 & 20.3 & 30.8 \\
Item B & 15.2 & 25.1 & 40.3 \\
Item C & 22.1 & 12.7 & 34.8 \\
\bottomrule
\end{tabular}
\caption{Sample Data Table}
\label{tab:sample}
\end{table}

\section{Lists}
\subsection{Bullet Points}
\begin{itemize}
  \item First item
  \item Second item
  \item Third item with sub-items:
    \begin{itemize}
      \item Sub-item 1
      \item Sub-item 2
    \end{itemize}
\end{itemize}

\subsection{Numbered List}
\begin{enumerate}
  \item First step
  \item Second step
  \item Third step
\end{enumerate}

\section{Citations and References}
When you use the citation search feature, references will be automatically managed. For example, you might cite a paper like this: \cite{example2024}. All citations will be automatically added to your \texttt{cite.bib} file.

\section{Conclusion}
This template demonstrates basic \LaTeX{} features including mathematical formulas, tables, and lists. The bibliography system is now set up to work seamlessly with Bibby AI's citation management. Feel free to modify it for your needs!

\bibliographystyle{plain}
\bibliography{cite}

\end{document}\documentclass{article}
\usepackage[utf8]{inputenc}
\usepackage{booktabs}
\usepackage{amsmath}
\usepackage{natbib}

\title{My Document}
\author{Your Name}
\date{\today}

\begin{document}

\maketitle

\section{Introduction}
This is a sample document demonstrating various \LaTeX{} features including mathematical equations, tables, and formatting. You can cite papers using the citation search feature, and they will be automatically added to your \texttt{cite.bib} file.

\section{Mathematical Expressions}
Here are some example equations:

\subsection{Quadratic Formula}
The quadratic formula is given by:
\[x = \frac{-b \pm \sqrt{b^2 - 4ac}}{2a}\]

\subsection{Matrix Example}
A sample matrix:
\[
\begin{pmatrix}
a_{11} & a_{12} & a_{13} \\
a_{21} & a_{22} & a_{23} \\
a_{31} & a_{32} & a_{33}
\end{pmatrix}
\]

\section{Tables}
Below is a sample table showing data:

\begin{table}[h]
\centering
\begin{tabular}{@{}lrrr@{}}
\toprule
\textbf{Category} & \textbf{Value 1} & \textbf{Value 2} & \textbf{Total} \\
\midrule
Item A & 10.5 & 20.3 & 30.8 \\
Item B & 15.2 & 25.1 & 40.3 \\
Item C & 22.1 & 12.7 & 34.8 \\
\bottomrule
\end{tabular}
\caption{Sample Data Table}
\label{tab:sample}
\end{table}

\section{Lists}
\subsection{Bullet Points}
\begin{itemize}
  \item First item
  \item Second item
  \item Third item with sub-items:
    \begin{itemize}
      \item Sub-item 1
      \item Sub-item 2
    \end{itemize}
\end{itemize}

\subsection{Numbered List}
\begin{enumerate}
  \item First step
  \item Second step
  \item Third step
\end{enumerate}

\section{Citations and References}
When you use the citation search feature, references will be automatically managed. For example, you might cite a paper like this: \cite{example2024}. All citations will be automatically added to your \texttt{cite.bib} file.

\section{Conclusion}
This template demonstrates basic \LaTeX{} features including mathematical formulas, tables, and lists. The bibliography system is now set up to work seamlessly with Bibby AI's citation management. Feel free to modify it for your needs!

\bibliographystyle{plain}
\bibliography{cite}

\end{document}\documentclass{article}
\usepackage[utf8]{inputenc}
\usepackage{booktabs}
\usepackage{amsmath}
\usepackage{natbib}

\title{My Document}
\author{Your Name}
\date{\today}

\begin{document}

\maketitle

\section{Introduction}
This is a sample document demonstrating various \LaTeX{} features including mathematical equations, tables, and formatting. You can cite papers using the citation search feature, and they will be automatically added to your \texttt{cite.bib} file.

\section{Mathematical Expressions}
Here are some example equations:

\subsection{Quadratic Formula}
The quadratic formula is given by:
\[x = \frac{-b \pm \sqrt{b^2 - 4ac}}{2a}\]

\subsection{Matrix Example}
A sample matrix:
\[
\begin{pmatrix}
a_{11} & a_{12} & a_{13} \\
a_{21} & a_{22} & a_{23} \\
a_{31} & a_{32} & a_{33}
\end{pmatrix}
\]

\section{Tables}
Below is a sample table showing data:

\begin{table}[h]
\centering
\begin{tabular}{@{}lrrr@{}}
\toprule
\textbf{Category} & \textbf{Value 1} & \textbf{Value 2} & \textbf{Total} \\
\midrule
Item A & 10.5 & 20.3 & 30.8 \\
Item B & 15.2 & 25.1 & 40.3 \\
Item C & 22.1 & 12.7 & 34.8 \\
\bottomrule
\end{tabular}
\caption{Sample Data Table}
\label{tab:sample}
\end{table}

\section{Lists}
\subsection{Bullet Points}
\begin{itemize}
  \item First item
  \item Second item
  \item Third item with sub-items:
    \begin{itemize}
      \item Sub-item 1
      \item Sub-item 2
    \end{itemize}
\end{itemize}

\subsection{Numbered List}
\begin{enumerate}
  \item First step
  \item Second step
  \item Third step
\end{enumerate}

\section{Citations and References}
When you use the citation search feature, references will be automatically managed. For example, you might cite a paper like this: \cite{example2024}. All citations will be automatically added to your \texttt{cite.bib} file.

\section{Conclusion}
This template demonstrates basic \LaTeX{} features including mathematical formulas, tables, and lists. The bibliography system is now set up to work seamlessly with Bibby AI's citation management. Feel free to modify it for your needs!

\bibliographystyle{plain}
\bibliography{cite}

\end{document}\documentclass{article}
\usepackage[utf8]{inputenc}
\usepackage{booktabs}
\usepackage{amsmath}
\usepackage{natbib}

\title{My Document}
\author{Your Name}
\date{\today}

\begin{document}

\maketitle

\section{Introduction}
This is a sample document demonstrating various \LaTeX{} features including mathematical equations, tables, and formatting. You can cite papers using the citation search feature, and they will be automatically added to your \texttt{cite.bib} file.

\section{Mathematical Expressions}
Here are some example equations:

\subsection{Quadratic Formula}
The quadratic formula is given by:
\[x = \frac{-b \pm \sqrt{b^2 - 4ac}}{2a}\]

\subsection{Matrix Example}
A sample matrix:
\[
\begin{pmatrix}
a_{11} & a_{12} & a_{13} \\
a_{21} & a_{22} & a_{23} \\
a_{31} & a_{32} & a_{33}
\end{pmatrix}
\]

\section{Tables}
Below is a sample table showing data:

\begin{table}[h]
\centering
\begin{tabular}{@{}lrrr@{}}
\toprule
\textbf{Category} & \textbf{Value 1} & \textbf{Value 2} & \textbf{Total} \\
\midrule
Item A & 10.5 & 20.3 & 30.8 \\
Item B & 15.2 & 25.1 & 40.3 \\
Item C & 22.1 & 12.7 & 34.8 \\
\bottomrule
\end{tabular}
\caption{Sample Data Table}
\label{tab:sample}
\end{table}

\section{Lists}
\subsection{Bullet Points}
\begin{itemize}
  \item First item
  \item Second item
  \item Third item with sub-items:
    \begin{itemize}
      \item Sub-item 1
      \item Sub-item 2
    \end{itemize}
\end{itemize}

\subsection{Numbered List}
\begin{enumerate}
  \item First step
  \item Second step
  \item Third step
\end{enumerate}

\section{Citations and References}
When you use the citation search feature, references will be automatically managed. For example, you might cite a paper like this: \cite{example2024}. All citations will be automatically added to your \texttt{cite.bib} file.

\section{Conclusion}
This template demonstrates basic \LaTeX{} features including mathematical formulas, tables, and lists. The bibliography system is now set up to work seamlessly with Bibby AI's citation management. Feel free to modify it for your needs!

\bibliographystyle{plain}
\bibliography{cite}

\end{document}\documentclass{article}
\usepackage[utf8]{inputenc}
\usepackage{booktabs}
\usepackage{amsmath}
\usepackage{natbib}

\title{My Document}
\author{Your Name}
\date{\today}

\begin{document}

\maketitle

\section{Introduction}
This is a sample document demonstrating various \LaTeX{} features including mathematical equations, tables, and formatting. You can cite papers using the citation search feature, and they will be automatically added to your \texttt{cite.bib} file.

\section{Mathematical Expressions}
Here are some example equations:

\subsection{Quadratic Formula}
The quadratic formula is given by:
\[x = \frac{-b \pm \sqrt{b^2 - 4ac}}{2a}\]

\subsection{Matrix Example}
A sample matrix:
\[
\begin{pmatrix}
a_{11} & a_{12} & a_{13} \\
a_{21} & a_{22} & a_{23} \\
a_{31} & a_{32} & a_{33}
\end{pmatrix}
\]

\section{Tables}
Below is a sample table showing data:

\begin{table}[h]
\centering
\begin{tabular}{@{}lrrr@{}}
\toprule
\textbf{Category} & \textbf{Value 1} & \textbf{Value 2} & \textbf{Total} \\
\midrule
Item A & 10.5 & 20.3 & 30.8 \\
Item B & 15.2 & 25.1 & 40.3 \\
Item C & 22.1 & 12.7 & 34.8 \\
\bottomrule
\end{tabular}
\caption{Sample Data Table}
\label{tab:sample}
\end{table}

\section{Lists}
\subsection{Bullet Points}
\begin{itemize}
  \item First item
  \item Second item
  \item Third item with sub-items:
    \begin{itemize}
      \item Sub-item 1
      \item Sub-item 2
    \end{itemize}
\end{itemize}

\subsection{Numbered List}
\begin{enumerate}
  \item First step
  \item Second step
  \item Third step
\end{enumerate}

\section{Citations and References}
When you use the citation search feature, references will be automatically managed. For example, you might cite a paper like this: \cite{example2024}. All citations will be automatically added to your \texttt{cite.bib} file.

\section{Conclusion}
This template demonstrates basic \LaTeX{} features including mathematical formulas, tables, and lists. The bibliography system is now set up to work seamlessly with Bibby AI's citation management. Feel free to modify it for your needs!

\bibliographystyle{plain}
\bibliography{cite}

\end{document}\documentclass{article}
\usepackage[utf8]{inputenc}
\usepackage{booktabs}
\usepackage{amsmath}
\usepackage{natbib}

\title{My Document}
\author{Your Name}
\date{\today}

\begin{document}

\maketitle

\section{Introduction}
This is a sample document demonstrating various \LaTeX{} features including mathematical equations, tables, and formatting. You can cite papers using the citation search feature, and they will be automatically added to your \texttt{cite.bib} file.

\section{Mathematical Expressions}
Here are some example equations:

\subsection{Quadratic Formula}
The quadratic formula is given by:
\[x = \frac{-b \pm \sqrt{b^2 - 4ac}}{2a}\]

\subsection{Matrix Example}
A sample matrix:
\[
\begin{pmatrix}
a_{11} & a_{12} & a_{13} \\
a_{21} & a_{22} & a_{23} \\
a_{31} & a_{32} & a_{33}
\end{pmatrix}
\]

\section{Tables}
Below is a sample table showing data:

\begin{table}[h]
\centering
\begin{tabular}{@{}lrrr@{}}
\toprule
\textbf{Category} & \textbf{Value 1} & \textbf{Value 2} & \textbf{Total} \\
\midrule
Item A & 10.5 & 20.3 & 30.8 \\
Item B & 15.2 & 25.1 & 40.3 \\
Item C & 22.1 & 12.7 & 34.8 \\
\bottomrule
\end{tabular}
\caption{Sample Data Table}
\label{tab:sample}
\end{table}

\section{Lists}
\subsection{Bullet Points}
\begin{itemize}
  \item First item
  \item Second item
  \item Third item with sub-items:
    \begin{itemize}
      \item Sub-item 1
      \item Sub-item 2
    \end{itemize}
\end{itemize}

\subsection{Numbered List}
\begin{enumerate}
  \item First step
  \item Second step
  \item Third step
\end{enumerate}

\section{Citations and References}
When you use the citation search feature, references will be automatically managed. For example, you might cite a paper like this: \cite{example2024}. All citations will be automatically added to your \texttt{cite.bib} file.

\section{Conclusion}
This template demonstrates basic \LaTeX{} features including mathematical formulas, tables, and lists. The bibliography system is now set up to work seamlessly with Bibby AI's citation management. Feel free to modify it for your needs!

\bibliographystyle{plain}
\bibliography{cite}

\end{document}\documentclass{article}
\usepackage[utf8]{inputenc}
\usepackage{booktabs}
\usepackage{amsmath}
\usepackage{natbib}

\title{My Document}
\author{Your Name}
\date{\today}

\begin{document}

\maketitle

\section{Introduction}
This is a sample document demonstrating various \LaTeX{} features including mathematical equations, tables, and formatting. You can cite papers using the citation search feature, and they will be automatically added to your \texttt{cite.bib} file.

\section{Mathematical Expressions}
Here are some example equations:

\subsection{Quadratic Formula}
The quadratic formula is given by:
\[x = \frac{-b \pm \sqrt{b^2 - 4ac}}{2a}\]

\subsection{Matrix Example}
A sample matrix:
\[
\begin{pmatrix}
a_{11} & a_{12} & a_{13} \\
a_{21} & a_{22} & a_{23} \\
a_{31} & a_{32} & a_{33}
\end{pmatrix}
\]

\section{Tables}
Below is a sample table showing data:

\begin{table}[h]
\centering
\begin{tabular}{@{}lrrr@{}}
\toprule
\textbf{Category} & \textbf{Value 1} & \textbf{Value 2} & \textbf{Total} \\
\midrule
Item A & 10.5 & 20.3 & 30.8 \\
Item B & 15.2 & 25.1 & 40.3 \\
Item C & 22.1 & 12.7 & 34.8 \\
\bottomrule
\end{tabular}
\caption{Sample Data Table}
\label{tab:sample}
\end{table}

\section{Lists}
\subsection{Bullet Points}
\begin{itemize}
  \item First item
  \item Second item
  \item Third item with sub-items:
    \begin{itemize}
      \item Sub-item 1
      \item Sub-item 2
    \end{itemize}
\end{itemize}

\subsection{Numbered List}
\begin{enumerate}
  \item First step
  \item Second step
  \item Third step
\end{enumerate}

\section{Citations and References}
When you use the citation search feature, references will be automatically managed. For example, you might cite a paper like this: \cite{example2024}. All citations will be automatically added to your \texttt{cite.bib} file.

\section{Conclusion}
This template demonstrates basic \LaTeX{} features including mathematical formulas, tables, and lists. The bibliography system is now set up to work seamlessly with Bibby AI's citation management. Feel free to modify it for your needs!

\bibliographystyle{plain}
\bibliography{cite}

\end{document}\documentclass{article}
\usepackage[utf8]{inputenc}
\usepackage{booktabs}
\usepackage{amsmath}
\usepackage{natbib}

\title{My Document}
\author{Your Name}
\date{\today}

\begin{document}

\maketitle

\section{Introduction}
This is a sample document demonstrating various \LaTeX{} features including mathematical equations, tables, and formatting. You can cite papers using the citation search feature, and they will be automatically added to your \texttt{cite.bib} file.

\section{Mathematical Expressions}
Here are some example equations:

\subsection{Quadratic Formula}
The quadratic formula is given by:
\[x = \frac{-b \pm \sqrt{b^2 - 4ac}}{2a}\]

\subsection{Matrix Example}
A sample matrix:
\[
\begin{pmatrix}
a_{11} & a_{12} & a_{13} \\
a_{21} & a_{22} & a_{23} \\
a_{31} & a_{32} & a_{33}
\end{pmatrix}
\]

\section{Tables}
Below is a sample table showing data:

\begin{table}[h]
\centering
\begin{tabular}{@{}lrrr@{}}
\toprule
\textbf{Category} & \textbf{Value 1} & \textbf{Value 2} & \textbf{Total} \\
\midrule
Item A & 10.5 & 20.3 & 30.8 \\
Item B & 15.2 & 25.1 & 40.3 \\
Item C & 22.1 & 12.7 & 34.8 \\
\bottomrule
\end{tabular}
\caption{Sample Data Table}
\label{tab:sample}
\end{table}

\section{Lists}
\subsection{Bullet Points}
\begin{itemize}
  \item First item
  \item Second item
  \item Third item with sub-items:
    \begin{itemize}
      \item Sub-item 1
      \item Sub-item 2
    \end{itemize}
\end{itemize}

\subsection{Numbered List}
\begin{enumerate}
  \item First step
  \item Second step
  \item Third step
\end{enumerate}

\section{Citations and References}
When you use the citation search feature, references will be automatically managed. For example, you might cite a paper like this: \cite{example2024}. All citations will be automatically added to your \texttt{cite.bib} file.

\section{Conclusion}
This template demonstrates basic \LaTeX{} features including mathematical formulas, tables, and lists. The bibliography system is now set up to work seamlessly with Bibby AI's citation management. Feel free to modify it for your needs!

\bibliographystyle{plain}
\bibliography{cite}

\end{document}\documentclass{article}
\usepackage[utf8]{inputenc}
\usepackage{booktabs}
\usepackage{amsmath}
\usepackage{natbib}

\title{My Document}
\author{Your Name}
\date{\today}

\begin{document}

\maketitle

\section{Introduction}
This is a sample document demonstrating various \LaTeX{} features including mathematical equations, tables, and formatting. You can cite papers using the citation search feature, and they will be automatically added to your \texttt{cite.bib} file.

\section{Mathematical Expressions}
Here are some example equations:

\subsection{Quadratic Formula}
The quadratic formula is given by:
\[x = \frac{-b \pm \sqrt{b^2 - 4ac}}{2a}\]

\subsection{Matrix Example}
A sample matrix:
\[
\begin{pmatrix}
a_{11} & a_{12} & a_{13} \\
a_{21} & a_{22} & a_{23} \\
a_{31} & a_{32} & a_{33}
\end{pmatrix}
\]

\section{Tables}
Below is a sample table showing data:

\begin{table}[h]
\centering
\begin{tabular}{@{}lrrr@{}}
\toprule
\textbf{Category} & \textbf{Value 1} & \textbf{Value 2} & \textbf{Total} \\
\midrule
Item A & 10.5 & 20.3 & 30.8 \\
Item B & 15.2 & 25.1 & 40.3 \\
Item C & 22.1 & 12.7 & 34.8 \\
\bottomrule
\end{tabular}
\caption{Sample Data Table}
\label{tab:sample}
\end{table}

\section{Lists}
\subsection{Bullet Points}
\begin{itemize}
  \item First item
  \item Second item
  \item Third item with sub-items:
    \begin{itemize}
      \item Sub-item 1
      \item Sub-item 2
    \end{itemize}
\end{itemize}

\subsection{Numbered List}
\begin{enumerate}
  \item First step
  \item Second step
  \item Third step
\end{enumerate}

\section{Citations and References}
When you use the citation search feature, references will be automatically managed. For example, you might cite a paper like this: \cite{example2024}. All citations will be automatically added to your \texttt{cite.bib} file.

\section{Conclusion}
This template demonstrates basic \LaTeX{} features including mathematical formulas, tables, and lists. The bibliography system is now set up to work seamlessly with Bibby AI's citation management. Feel free to modify it for your needs!

\bibliographystyle{plain}
\bibliography{cite}

\end{document}\documentclass{article}
\usepackage[utf8]{inputenc}
\usepackage{booktabs}
\usepackage{amsmath}
\usepackage{natbib}

\title{My Document}
\author{Your Name}
\date{\today}

\begin{document}

\maketitle

\section{Introduction}
This is a sample document demonstrating various \LaTeX{} features including mathematical equations, tables, and formatting. You can cite papers using the citation search feature, and they will be automatically added to your \texttt{cite.bib} file.

\section{Mathematical Expressions}
Here are some example equations:

\subsection{Quadratic Formula}
The quadratic formula is given by:
\[x = \frac{-b \pm \sqrt{b^2 - 4ac}}{2a}\]

\subsection{Matrix Example}
A sample matrix:
\[
\begin{pmatrix}
a_{11} & a_{12} & a_{13} \\
a_{21} & a_{22} & a_{23} \\
a_{31} & a_{32} & a_{33}
\end{pmatrix}
\]

\section{Tables}
Below is a sample table showing data:

\begin{table}[h]
\centering
\begin{tabular}{@{}lrrr@{}}
\toprule
\textbf{Category} & \textbf{Value 1} & \textbf{Value 2} & \textbf{Total} \\
\midrule
Item A & 10.5 & 20.3 & 30.8 \\
Item B & 15.2 & 25.1 & 40.3 \\
Item C & 22.1 & 12.7 & 34.8 \\
\bottomrule
\end{tabular}
\caption{Sample Data Table}
\label{tab:sample}
\end{table}

\section{Lists}
\subsection{Bullet Points}
\begin{itemize}
  \item First item
  \item Second item
  \item Third item with sub-items:
    \begin{itemize}
      \item Sub-item 1
      \item Sub-item 2
    \end{itemize}
\end{itemize}

\subsection{Numbered List}
\begin{enumerate}
  \item First step
  \item Second step
  \item Third step
\end{enumerate}

\section{Citations and References}
When you use the citation search feature, references will be automatically managed. For example, you might cite a paper like this: \cite{example2024}. All citations will be automatically added to your \texttt{cite.bib} file.

\section{Conclusion}
This template demonstrates basic \LaTeX{} features including mathematical formulas, tables, and lists. The bibliography system is now set up to work seamlessly with Bibby AI's citation management. Feel free to modify it for your needs!

\bibliographystyle{plain}
\bibliography{cite}

\end{document}\documentclass{article}
\usepackage[utf8]{inputenc}
\usepackage{booktabs}
\usepackage{amsmath}
\usepackage{natbib}

\title{My Document}
\author{Your Name}
\date{\today}

\begin{document}

\maketitle

\section{Introduction}
This is a sample document demonstrating various \LaTeX{} features including mathematical equations, tables, and formatting. You can cite papers using the citation search feature, and they will be automatically added to your \texttt{cite.bib} file.

\section{Mathematical Expressions}
Here are some example equations:

\subsection{Quadratic Formula}
The quadratic formula is given by:
\[x = \frac{-b \pm \sqrt{b^2 - 4ac}}{2a}\]

\subsection{Matrix Example}
A sample matrix:
\[
\begin{pmatrix}
a_{11} & a_{12} & a_{13} \\
a_{21} & a_{22} & a_{23} \\
a_{31} & a_{32} & a_{33}
\end{pmatrix}
\]

\section{Tables}
Below is a sample table showing data:

\begin{table}[h]
\centering
\begin{tabular}{@{}lrrr@{}}
\toprule
\textbf{Category} & \textbf{Value 1} & \textbf{Value 2} & \textbf{Total} \\
\midrule
Item A & 10.5 & 20.3 & 30.8 \\
Item B & 15.2 & 25.1 & 40.3 \\
Item C & 22.1 & 12.7 & 34.8 \\
\bottomrule
\end{tabular}
\caption{Sample Data Table}
\label{tab:sample}
\end{table}

\section{Lists}
\subsection{Bullet Points}
\begin{itemize}
  \item First item
  \item Second item
  \item Third item with sub-items:
    \begin{itemize}
      \item Sub-item 1
      \item Sub-item 2
    \end{itemize}
\end{itemize}

\subsection{Numbered List}
\begin{enumerate}
  \item First step
  \item Second step
  \item Third step
\end{enumerate}

\section{Citations and References}
When you use the citation search feature, references will be automatically managed. For example, you might cite a paper like this: \cite{example2024}. All citations will be automatically added to your \texttt{cite.bib} file.

\section{Conclusion}
This template demonstrates basic \LaTeX{} features including mathematical formulas, tables, and lists. The bibliography system is now set up to work seamlessly with Bibby AI's citation management. Feel free to modify it for your needs!

\bibliographystyle{plain}
\bibliography{cite}

\end{document}\documentclass{article}
\usepackage[utf8]{inputenc}
\usepackage{booktabs}
\usepackage{amsmath}
\usepackage{natbib}

\title{My Document}
\author{Your Name}
\date{\today}

\begin{document}

\maketitle

\section{Introduction}
This is a sample document demonstrating various \LaTeX{} features including mathematical equations, tables, and formatting. You can cite papers using the citation search feature, and they will be automatically added to your \texttt{cite.bib} file.

\section{Mathematical Expressions}
Here are some example equations:

\subsection{Quadratic Formula}
The quadratic formula is given by:
\[x = \frac{-b \pm \sqrt{b^2 - 4ac}}{2a}\]

\subsection{Matrix Example}
A sample matrix:
\[
\begin{pmatrix}
a_{11} & a_{12} & a_{13} \\
a_{21} & a_{22} & a_{23} \\
a_{31} & a_{32} & a_{33}
\end{pmatrix}
\]

\section{Tables}
Below is a sample table showing data:

\begin{table}[h]
\centering
\begin{tabular}{@{}lrrr@{}}
\toprule
\textbf{Category} & \textbf{Value 1} & \textbf{Value 2} & \textbf{Total} \\
\midrule
Item A & 10.5 & 20.3 & 30.8 \\
Item B & 15.2 & 25.1 & 40.3 \\
Item C & 22.1 & 12.7 & 34.8 \\
\bottomrule
\end{tabular}
\caption{Sample Data Table}
\label{tab:sample}
\end{table}

\section{Lists}
\subsection{Bullet Points}
\begin{itemize}
  \item First item
  \item Second item
  \item Third item with sub-items:
    \begin{itemize}
      \item Sub-item 1
      \item Sub-item 2
    \end{itemize}
\end{itemize}

\subsection{Numbered List}
\begin{enumerate}
  \item First step
  \item Second step
  \item Third step
\end{enumerate}

\section{Citations and References}
When you use the citation search feature, references will be automatically managed. For example, you might cite a paper like this: \cite{example2024}. All citations will be automatically added to your \texttt{cite.bib} file.

\section{Conclusion}
This template demonstrates basic \LaTeX{} features including mathematical formulas, tables, and lists. The bibliography system is now set up to work seamlessly with Bibby AI's citation management. Feel free to modify it for your needs!

\bibliographystyle{plain}
\bibliography{cite}

\end{document}\documentclass{article}
\usepackage[utf8]{inputenc}
\usepackage{booktabs}
\usepackage{amsmath}
\usepackage{natbib}

\title{My Document}
\author{Your Name}
\date{\today}

\begin{document}

\maketitle

\section{Introduction}
This is a sample document demonstrating various \LaTeX{} features including mathematical equations, tables, and formatting. You can cite papers using the citation search feature, and they will be automatically added to your \texttt{cite.bib} file.

\section{Mathematical Expressions}
Here are some example equations:

\subsection{Quadratic Formula}
The quadratic formula is given by:
\[x = \frac{-b \pm \sqrt{b^2 - 4ac}}{2a}\]

\subsection{Matrix Example}
A sample matrix:
\[
\begin{pmatrix}
a_{11} & a_{12} & a_{13} \\
a_{21} & a_{22} & a_{23} \\
a_{31} & a_{32} & a_{33}
\end{pmatrix}
\]

\section{Tables}
Below is a sample table showing data:

\begin{table}[h]
\centering
\begin{tabular}{@{}lrrr@{}}
\toprule
\textbf{Category} & \textbf{Value 1} & \textbf{Value 2} & \textbf{Total} \\
\midrule
Item A & 10.5 & 20.3 & 30.8 \\
Item B & 15.2 & 25.1 & 40.3 \\
Item C & 22.1 & 12.7 & 34.8 \\
\bottomrule
\end{tabular}
\caption{Sample Data Table}
\label{tab:sample}
\end{table}

\section{Lists}
\subsection{Bullet Points}
\begin{itemize}
  \item First item
  \item Second item
  \item Third item with sub-items:
    \begin{itemize}
      \item Sub-item 1
      \item Sub-item 2
    \end{itemize}
\end{itemize}

\subsection{Numbered List}
\begin{enumerate}
  \item First step
  \item Second step
  \item Third step
\end{enumerate}

\section{Citations and References}
When you use the citation search feature, references will be automatically managed. For example, you might cite a paper like this: \cite{example2024}. All citations will be automatically added to your \texttt{cite.bib} file.

\section{Conclusion}
This template demonstrates basic \LaTeX{} features including mathematical formulas, tables, and lists. The bibliography system is now set up to work seamlessly with Bibby AI's citation management. Feel free to modify it for your needs!

\bibliographystyle{plain}
\bibliography{cite}

\end{document}\documentclass{article}
\usepackage[utf8]{inputenc}
\usepackage{booktabs}
\usepackage{amsmath}
\usepackage{natbib}

\title{My Document}
\author{Your Name}
\date{\today}

\begin{document}

\maketitle

\section{Introduction}
This is a sample document demonstrating various \LaTeX{} features including mathematical equations, tables, and formatting. You can cite papers using the citation search feature, and they will be automatically added to your \texttt{cite.bib} file.

\section{Mathematical Expressions}
Here are some example equations:

\subsection{Quadratic Formula}
The quadratic formula is given by:
\[x = \frac{-b \pm \sqrt{b^2 - 4ac}}{2a}\]

\subsection{Matrix Example}
A sample matrix:
\[
\begin{pmatrix}
a_{11} & a_{12} & a_{13} \\
a_{21} & a_{22} & a_{23} \\
a_{31} & a_{32} & a_{33}
\end{pmatrix}
\]

\section{Tables}
Below is a sample table showing data:

\begin{table}[h]
\centering
\begin{tabular}{@{}lrrr@{}}
\toprule
\textbf{Category} & \textbf{Value 1} & \textbf{Value 2} & \textbf{Total} \\
\midrule
Item A & 10.5 & 20.3 & 30.8 \\
Item B & 15.2 & 25.1 & 40.3 \\
Item C & 22.1 & 12.7 & 34.8 \\
\bottomrule
\end{tabular}
\caption{Sample Data Table}
\label{tab:sample}
\end{table}

\section{Lists}
\subsection{Bullet Points}
\begin{itemize}
  \item First item
  \item Second item
  \item Third item with sub-items:
    \begin{itemize}
      \item Sub-item 1
      \item Sub-item 2
    \end{itemize}
\end{itemize}

\subsection{Numbered List}
\begin{enumerate}
  \item First step
  \item Second step
  \item Third step
\end{enumerate}

\section{Citations and References}
When you use the citation search feature, references will be automatically managed. For example, you might cite a paper like this: \cite{example2024}. All citations will be automatically added to your \texttt{cite.bib} file.

\section{Conclusion}
This template demonstrates basic \LaTeX{} features including mathematical formulas, tables, and lists. The bibliography system is now set up to work seamlessly with Bibby AI's citation management. Feel free to modify it for your needs!

\bibliographystyle{plain}
\bibliography{cite}

\end{document}\documentclass{article}
\usepackage[utf8]{inputenc}
\usepackage{booktabs}
\usepackage{amsmath}
\usepackage{natbib}

\title{My Document}
\author{Your Name}
\date{\today}

\begin{document}

\maketitle

\section{Introduction}
This is a sample document demonstrating various \LaTeX{} features including mathematical equations, tables, and formatting. You can cite papers using the citation search feature, and they will be automatically added to your \texttt{cite.bib} file.

\section{Mathematical Expressions}
Here are some example equations:

\subsection{Quadratic Formula}
The quadratic formula is given by:
\[x = \frac{-b \pm \sqrt{b^2 - 4ac}}{2a}\]

\subsection{Matrix Example}
A sample matrix:
\[
\begin{pmatrix}
a_{11} & a_{12} & a_{13} \\
a_{21} & a_{22} & a_{23} \\
a_{31} & a_{32} & a_{33}
\end{pmatrix}
\]

\section{Tables}
Below is a sample table showing data:

\begin{table}[h]
\centering
\begin{tabular}{@{}lrrr@{}}
\toprule
\textbf{Category} & \textbf{Value 1} & \textbf{Value 2} & \textbf{Total} \\
\midrule
Item A & 10.5 & 20.3 & 30.8 \\
Item B & 15.2 & 25.1 & 40.3 \\
Item C & 22.1 & 12.7 & 34.8 \\
\bottomrule
\end{tabular}
\caption{Sample Data Table}
\label{tab:sample}
\end{table}

\section{Lists}
\subsection{Bullet Points}
\begin{itemize}
  \item First item
  \item Second item
  \item Third item with sub-items:
    \begin{itemize}
      \item Sub-item 1
      \item Sub-item 2
    \end{itemize}
\end{itemize}

\subsection{Numbered List}
\begin{enumerate}
  \item First step
  \item Second step
  \item Third step
\end{enumerate}

\section{Citations and References}
When you use the citation search feature, references will be automatically managed. For example, you might cite a paper like this: \cite{example2024}. All citations will be automatically added to your \texttt{cite.bib} file.

\section{Conclusion}
This template demonstrates basic \LaTeX{} features including mathematical formulas, tables, and lists. The bibliography system is now set up to work seamlessly with Bibby AI's citation management. Feel free to modify it for your needs!

\bibliographystyle{plain}
\bibliography{cite}

\end{document}\documentclass{article}
\usepackage[utf8]{inputenc}
\usepackage{booktabs}
\usepackage{amsmath}
\usepackage{natbib}

\title{My Document}
\author{Your Name}
\date{\today}

\begin{document}

\maketitle

\section{Introduction}
This is a sample document demonstrating various \LaTeX{} features including mathematical equations, tables, and formatting. You can cite papers using the citation search feature, and they will be automatically added to your \texttt{cite.bib} file.

\section{Mathematical Expressions}
Here are some example equations:

\subsection{Quadratic Formula}
The quadratic formula is given by:
\[x = \frac{-b \pm \sqrt{b^2 - 4ac}}{2a}\]

\subsection{Matrix Example}
A sample matrix:
\[
\begin{pmatrix}
a_{11} & a_{12} & a_{13} \\
a_{21} & a_{22} & a_{23} \\
a_{31} & a_{32} & a_{33}
\end{pmatrix}
\]

\section{Tables}
Below is a sample table showing data:

\begin{table}[h]
\centering
\begin{tabular}{@{}lrrr@{}}
\toprule
\textbf{Category} & \textbf{Value 1} & \textbf{Value 2} & \textbf{Total} \\
\midrule
Item A & 10.5 & 20.3 & 30.8 \\
Item B & 15.2 & 25.1 & 40.3 \\
Item C & 22.1 & 12.7 & 34.8 \\
\bottomrule
\end{tabular}
\caption{Sample Data Table}
\label{tab:sample}
\end{table}

\section{Lists}
\subsection{Bullet Points}
\begin{itemize}
  \item First item
  \item Second item
  \item Third item with sub-items:
    \begin{itemize}
      \item Sub-item 1
      \item Sub-item 2
    \end{itemize}
\end{itemize}

\subsection{Numbered List}
\begin{enumerate}
  \item First step
  \item Second step
  \item Third step
\end{enumerate}

\section{Citations and References}
When you use the citation search feature, references will be automatically managed. For example, you might cite a paper like this: \cite{example2024}. All citations will be automatically added to your \texttt{cite.bib} file.

\section{Conclusion}
This template demonstrates basic \LaTeX{} features including mathematical formulas, tables, and lists. The bibliography system is now set up to work seamlessly with Bibby AI's citation management. Feel free to modify it for your needs!

\bibliographystyle{plain}
\bibliography{cite}

\end{document}\documentclass{article}
\usepackage[utf8]{inputenc}
\usepackage{booktabs}
\usepackage{amsmath}
\usepackage{natbib}

\title{My Document}
\author{Your Name}
\date{\today}

\begin{document}

\maketitle

\section{Introduction}
This is a sample document demonstrating various \LaTeX{} features including mathematical equations, tables, and formatting. You can cite papers using the citation search feature, and they will be automatically added to your \texttt{cite.bib} file.

\section{Mathematical Expressions}
Here are some example equations:

\subsection{Quadratic Formula}
The quadratic formula is given by:
\[x = \frac{-b \pm \sqrt{b^2 - 4ac}}{2a}\]

\subsection{Matrix Example}
A sample matrix:
\[
\begin{pmatrix}
a_{11} & a_{12} & a_{13} \\
a_{21} & a_{22} & a_{23} \\
a_{31} & a_{32} & a_{33}
\end{pmatrix}
\]

\section{Tables}
Below is a sample table showing data:

\begin{table}[h]
\centering
\begin{tabular}{@{}lrrr@{}}
\toprule
\textbf{Category} & \textbf{Value 1} & \textbf{Value 2} & \textbf{Total} \\
\midrule
Item A & 10.5 & 20.3 & 30.8 \\
Item B & 15.2 & 25.1 & 40.3 \\
Item C & 22.1 & 12.7 & 34.8 \\
\bottomrule
\end{tabular}
\caption{Sample Data Table}
\label{tab:sample}
\end{table}

\section{Lists}
\subsection{Bullet Points}
\begin{itemize}
  \item First item
  \item Second item
  \item Third item with sub-items:
    \begin{itemize}
      \item Sub-item 1
      \item Sub-item 2
    \end{itemize}
\end{itemize}

\subsection{Numbered List}
\begin{enumerate}
  \item First step
  \item Second step
  \item Third step
\end{enumerate}

\section{Citations and References}
When you use the citation search feature, references will be automatically managed. For example, you might cite a paper like this: \cite{example2024}. All citations will be automatically added to your \texttt{cite.bib} file.

\section{Conclusion}
This template demonstrates basic \LaTeX{} features including mathematical formulas, tables, and lists. The bibliography system is now set up to work seamlessly with Bibby AI's citation management. Feel free to modify it for your needs!

\bibliographystyle{plain}
\bibliography{cite}

\end{document}\documentclass{article}
\usepackage[utf8]{inputenc}
\usepackage{booktabs}
\usepackage{amsmath}
\usepackage{natbib}

\title{My Document}
\author{Your Name}
\date{\today}

\begin{document}

\maketitle

\section{Introduction}
This is a sample document demonstrating various \LaTeX{} features including mathematical equations, tables, and formatting. You can cite papers using the citation search feature, and they will be automatically added to your \texttt{cite.bib} file.

\section{Mathematical Expressions}
Here are some example equations:

\subsection{Quadratic Formula}
The quadratic formula is given by:
\[x = \frac{-b \pm \sqrt{b^2 - 4ac}}{2a}\]

\subsection{Matrix Example}
A sample matrix:
\[
\begin{pmatrix}
a_{11} & a_{12} & a_{13} \\
a_{21} & a_{22} & a_{23} \\
a_{31} & a_{32} & a_{33}
\end{pmatrix}
\]

\section{Tables}
Below is a sample table showing data:

\begin{table}[h]
\centering
\begin{tabular}{@{}lrrr@{}}
\toprule
\textbf{Category} & \textbf{Value 1} & \textbf{Value 2} & \textbf{Total} \\
\midrule
Item A & 10.5 & 20.3 & 30.8 \\
Item B & 15.2 & 25.1 & 40.3 \\
Item C & 22.1 & 12.7 & 34.8 \\
\bottomrule
\end{tabular}
\caption{Sample Data Table}
\label{tab:sample}
\end{table}

\section{Lists}
\subsection{Bullet Points}
\begin{itemize}
  \item First item
  \item Second item
  \item Third item with sub-items:
    \begin{itemize}
      \item Sub-item 1
      \item Sub-item 2
    \end{itemize}
\end{itemize}

\subsection{Numbered List}
\begin{enumerate}
  \item First step
  \item Second step
  \item Third step
\end{enumerate}

\section{Citations and References}
When you use the citation search feature, references will be automatically managed. For example, you might cite a paper like this: \cite{example2024}. All citations will be automatically added to your \texttt{cite.bib} file.

\section{Conclusion}
This template demonstrates basic \LaTeX{} features including mathematical formulas, tables, and lists. The bibliography system is now set up to work seamlessly with Bibby AI's citation management. Feel free to modify it for your needs!

\bibliographystyle{plain}
\bibliography{cite}

\end{document}\documentclass{article}
\usepackage[utf8]{inputenc}
\usepackage{booktabs}
\usepackage{amsmath}
\usepackage{natbib}

\title{My Document}
\author{Your Name}
\date{\today}

\begin{document}

\maketitle

\section{Introduction}
This is a sample document demonstrating various \LaTeX{} features including mathematical equations, tables, and formatting. You can cite papers using the citation search feature, and they will be automatically added to your \texttt{cite.bib} file.

\section{Mathematical Expressions}
Here are some example equations:

\subsection{Quadratic Formula}
The quadratic formula is given by:
\[x = \frac{-b \pm \sqrt{b^2 - 4ac}}{2a}\]

\subsection{Matrix Example}
A sample matrix:
\[
\begin{pmatrix}
a_{11} & a_{12} & a_{13} \\
a_{21} & a_{22} & a_{23} \\
a_{31} & a_{32} & a_{33}
\end{pmatrix}
\]

\section{Tables}
Below is a sample table showing data:

\begin{table}[h]
\centering
\begin{tabular}{@{}lrrr@{}}
\toprule
\textbf{Category} & \textbf{Value 1} & \textbf{Value 2} & \textbf{Total} \\
\midrule
Item A & 10.5 & 20.3 & 30.8 \\
Item B & 15.2 & 25.1 & 40.3 \\
Item C & 22.1 & 12.7 & 34.8 \\
\bottomrule
\end{tabular}
\caption{Sample Data Table}
\label{tab:sample}
\end{table}

\section{Lists}
\subsection{Bullet Points}
\begin{itemize}
  \item First item
  \item Second item
  \item Third item with sub-items:
    \begin{itemize}
      \item Sub-item 1
      \item Sub-item 2
    \end{itemize}
\end{itemize}

\subsection{Numbered List}
\begin{enumerate}
  \item First step
  \item Second step
  \item Third step
\end{enumerate}

\section{Citations and References}
When you use the citation search feature, references will be automatically managed. For example, you might cite a paper like this: \cite{example2024}. All citations will be automatically added to your \texttt{cite.bib} file.

\section{Conclusion}
This template demonstrates basic \LaTeX{} features including mathematical formulas, tables, and lists. The bibliography system is now set up to work seamlessly with Bibby AI's citation management. Feel free to modify it for your needs!

\bibliographystyle{plain}
\bibliography{cite}

\end{document}\documentclass{article}
\usepackage[utf8]{inputenc}
\usepackage{booktabs}
\usepackage{amsmath}
\usepackage{natbib}

\title{My Document}
\author{Your Name}
\date{\today}

\begin{document}

\maketitle

\section{Introduction}
This is a sample document demonstrating various \LaTeX{} features including mathematical equations, tables, and formatting. You can cite papers using the citation search feature, and they will be automatically added to your \texttt{cite.bib} file.

\section{Mathematical Expressions}
Here are some example equations:

\subsection{Quadratic Formula}
The quadratic formula is given by:
\[x = \frac{-b \pm \sqrt{b^2 - 4ac}}{2a}\]

\subsection{Matrix Example}
A sample matrix:
\[
\begin{pmatrix}
a_{11} & a_{12} & a_{13} \\
a_{21} & a_{22} & a_{23} \\
a_{31} & a_{32} & a_{33}
\end{pmatrix}
\]

\section{Tables}
Below is a sample table showing data:

\begin{table}[h]
\centering
\begin{tabular}{@{}lrrr@{}}
\toprule
\textbf{Category} & \textbf{Value 1} & \textbf{Value 2} & \textbf{Total} \\
\midrule
Item A & 10.5 & 20.3 & 30.8 \\
Item B & 15.2 & 25.1 & 40.3 \\
Item C & 22.1 & 12.7 & 34.8 \\
\bottomrule
\end{tabular}
\caption{Sample Data Table}
\label{tab:sample}
\end{table}

\section{Lists}
\subsection{Bullet Points}
\begin{itemize}
  \item First item
  \item Second item
  \item Third item with sub-items:
    \begin{itemize}
      \item Sub-item 1
      \item Sub-item 2
    \end{itemize}
\end{itemize}

\subsection{Numbered List}
\begin{enumerate}
  \item First step
  \item Second step
  \item Third step
\end{enumerate}

\section{Citations and References}
When you use the citation search feature, references will be automatically managed. For example, you might cite a paper like this: \cite{example2024}. All citations will be automatically added to your \texttt{cite.bib} file.

\section{Conclusion}
This template demonstrates basic \LaTeX{} features including mathematical formulas, tables, and lists. The bibliography system is now set up to work seamlessly with Bibby AI's citation management. Feel free to modify it for your needs!

\bibliographystyle{plain}
\bibliography{cite}

\end{document}\documentclass{article}
\usepackage[utf8]{inputenc}
\usepackage{booktabs}
\usepackage{amsmath}
\usepackage{natbib}

\title{My Document}
\author{Your Name}
\date{\today}

\begin{document}

\maketitle

\section{Introduction}
This is a sample document demonstrating various \LaTeX{} features including mathematical equations, tables, and formatting. You can cite papers using the citation search feature, and they will be automatically added to your \texttt{cite.bib} file.

\section{Mathematical Expressions}
Here are some example equations:

\subsection{Quadratic Formula}
The quadratic formula is given by:
\[x = \frac{-b \pm \sqrt{b^2 - 4ac}}{2a}\]

\subsection{Matrix Example}
A sample matrix:
\[
\begin{pmatrix}
a_{11} & a_{12} & a_{13} \\
a_{21} & a_{22} & a_{23} \\
a_{31} & a_{32} & a_{33}
\end{pmatrix}
\]

\section{Tables}
Below is a sample table showing data:

\begin{table}[h]
\centering
\begin{tabular}{@{}lrrr@{}}
\toprule
\textbf{Category} & \textbf{Value 1} & \textbf{Value 2} & \textbf{Total} \\
\midrule
Item A & 10.5 & 20.3 & 30.8 \\
Item B & 15.2 & 25.1 & 40.3 \\
Item C & 22.1 & 12.7 & 34.8 \\
\bottomrule
\end{tabular}
\caption{Sample Data Table}
\label{tab:sample}
\end{table}

\section{Lists}
\subsection{Bullet Points}
\begin{itemize}
  \item First item
  \item Second item
  \item Third item with sub-items:
    \begin{itemize}
      \item Sub-item 1
      \item Sub-item 2
    \end{itemize}
\end{itemize}

\subsection{Numbered List}
\begin{enumerate}
  \item First step
  \item Second step
  \item Third step
\end{enumerate}

\section{Citations and References}
When you use the citation search feature, references will be automatically managed. For example, you might cite a paper like this: \cite{example2024}. All citations will be automatically added to your \texttt{cite.bib} file.

\section{Conclusion}
This template demonstrates basic \LaTeX{} features including mathematical formulas, tables, and lists. The bibliography system is now set up to work seamlessly with Bibby AI's citation management. Feel free to modify it for your needs!

\bibliographystyle{plain}
\bibliography{cite}

\end{document}\documentclass{article}
\usepackage[utf8]{inputenc}
\usepackage{booktabs}
\usepackage{amsmath}
\usepackage{natbib}

\title{My Document}
\author{Your Name}
\date{\today}

\begin{document}

\maketitle

\section{Introduction}
This is a sample document demonstrating various \LaTeX{} features including mathematical equations, tables, and formatting. You can cite papers using the citation search feature, and they will be automatically added to your \texttt{cite.bib} file.

\section{Mathematical Expressions}
Here are some example equations:

\subsection{Quadratic Formula}
The quadratic formula is given by:
\[x = \frac{-b \pm \sqrt{b^2 - 4ac}}{2a}\]

\subsection{Matrix Example}
A sample matrix:
\[
\begin{pmatrix}
a_{11} & a_{12} & a_{13} \\
a_{21} & a_{22} & a_{23} \\
a_{31} & a_{32} & a_{33}
\end{pmatrix}
\]

\section{Tables}
Below is a sample table showing data:

\begin{table}[h]
\centering
\begin{tabular}{@{}lrrr@{}}
\toprule
\textbf{Category} & \textbf{Value 1} & \textbf{Value 2} & \textbf{Total} \\
\midrule
Item A & 10.5 & 20.3 & 30.8 \\
Item B & 15.2 & 25.1 & 40.3 \\
Item C & 22.1 & 12.7 & 34.8 \\
\bottomrule
\end{tabular}
\caption{Sample Data Table}
\label{tab:sample}
\end{table}

\section{Lists}
\subsection{Bullet Points}
\begin{itemize}
  \item First item
  \item Second item
  \item Third item with sub-items:
    \begin{itemize}
      \item Sub-item 1
      \item Sub-item 2
    \end{itemize}
\end{itemize}

\subsection{Numbered List}
\begin{enumerate}
  \item First step
  \item Second step
  \item Third step
\end{enumerate}

\section{Citations and References}
When you use the citation search feature, references will be automatically managed. For example, you might cite a paper like this: \cite{example2024}. All citations will be automatically added to your \texttt{cite.bib} file.

\section{Conclusion}
This template demonstrates basic \LaTeX{} features including mathematical formulas, tables, and lists. The bibliography system is now set up to work seamlessly with Bibby AI's citation management. Feel free to modify it for your needs!

\bibliographystyle{plain}
\bibliography{cite}

\end{document}\documentclass{article}
\usepackage[utf8]{inputenc}
\usepackage{booktabs}
\usepackage{amsmath}
\usepackage{natbib}

\title{My Document}
\author{Your Name}
\date{\today}

\begin{document}

\maketitle

\section{Introduction}
This is a sample document demonstrating various \LaTeX{} features including mathematical equations, tables, and formatting. You can cite papers using the citation search feature, and they will be automatically added to your \texttt{cite.bib} file.

\section{Mathematical Expressions}
Here are some example equations:

\subsection{Quadratic Formula}
The quadratic formula is given by:
\[x = \frac{-b \pm \sqrt{b^2 - 4ac}}{2a}\]

\subsection{Matrix Example}
A sample matrix:
\[
\begin{pmatrix}
a_{11} & a_{12} & a_{13} \\
a_{21} & a_{22} & a_{23} \\
a_{31} & a_{32} & a_{33}
\end{pmatrix}
\]

\section{Tables}
Below is a sample table showing data:

\begin{table}[h]
\centering
\begin{tabular}{@{}lrrr@{}}
\toprule
\textbf{Category} & \textbf{Value 1} & \textbf{Value 2} & \textbf{Total} \\
\midrule
Item A & 10.5 & 20.3 & 30.8 \\
Item B & 15.2 & 25.1 & 40.3 \\
Item C & 22.1 & 12.7 & 34.8 \\
\bottomrule
\end{tabular}
\caption{Sample Data Table}
\label{tab:sample}
\end{table}

\section{Lists}
\subsection{Bullet Points}
\begin{itemize}
  \item First item
  \item Second item
  \item Third item with sub-items:
    \begin{itemize}
      \item Sub-item 1
      \item Sub-item 2
    \end{itemize}
\end{itemize}

\subsection{Numbered List}
\begin{enumerate}
  \item First step
  \item Second step
  \item Third step
\end{enumerate}

\section{Citations and References}
When you use the citation search feature, references will be automatically managed. For example, you might cite a paper like this: \cite{example2024}. All citations will be automatically added to your \texttt{cite.bib} file.

\section{Conclusion}
This template demonstrates basic \LaTeX{} features including mathematical formulas, tables, and lists. The bibliography system is now set up to work seamlessly with Bibby AI's citation management. Feel free to modify it for your needs!

\bibliographystyle{plain}
\bibliography{cite}

\end{document}\documentclass{article}
\usepackage[utf8]{inputenc}
\usepackage{booktabs}
\usepackage{amsmath}
\usepackage{natbib}

\title{My Document}
\author{Your Name}
\date{\today}

\begin{document}

\maketitle

\section{Introduction}
This is a sample document demonstrating various \LaTeX{} features including mathematical equations, tables, and formatting. You can cite papers using the citation search feature, and they will be automatically added to your \texttt{cite.bib} file.

\section{Mathematical Expressions}
Here are some example equations:

\subsection{Quadratic Formula}
The quadratic formula is given by:
\[x = \frac{-b \pm \sqrt{b^2 - 4ac}}{2a}\]

\subsection{Matrix Example}
A sample matrix:
\[
\begin{pmatrix}
a_{11} & a_{12} & a_{13} \\
a_{21} & a_{22} & a_{23} \\
a_{31} & a_{32} & a_{33}
\end{pmatrix}
\]

\section{Tables}
Below is a sample table showing data:

\begin{table}[h]
\centering
\begin{tabular}{@{}lrrr@{}}
\toprule
\textbf{Category} & \textbf{Value 1} & \textbf{Value 2} & \textbf{Total} \\
\midrule
Item A & 10.5 & 20.3 & 30.8 \\
Item B & 15.2 & 25.1 & 40.3 \\
Item C & 22.1 & 12.7 & 34.8 \\
\bottomrule
\end{tabular}
\caption{Sample Data Table}
\label{tab:sample}
\end{table}

\section{Lists}
\subsection{Bullet Points}
\begin{itemize}
  \item First item
  \item Second item
  \item Third item with sub-items:
    \begin{itemize}
      \item Sub-item 1
      \item Sub-item 2
    \end{itemize}
\end{itemize}

\subsection{Numbered List}
\begin{enumerate}
  \item First step
  \item Second step
  \item Third step
\end{enumerate}

\section{Citations and References}
When you use the citation search feature, references will be automatically managed. For example, you might cite a paper like this: \cite{example2024}. All citations will be automatically added to your \texttt{cite.bib} file.

\section{Conclusion}
This template demonstrates basic \LaTeX{} features including mathematical formulas, tables, and lists. The bibliography system is now set up to work seamlessly with Bibby AI's citation management. Feel free to modify it for your needs!

\bibliographystyle{plain}
\bibliography{cite}

\end{document}\documentclass{article}
\usepackage[utf8]{inputenc}
\usepackage{booktabs}
\usepackage{amsmath}
\usepackage{natbib}

\title{My Document}
\author{Your Name}
\date{\today}

\begin{document}

\maketitle

\section{Introduction}
This is a sample document demonstrating various \LaTeX{} features including mathematical equations, tables, and formatting. You can cite papers using the citation search feature, and they will be automatically added to your \texttt{cite.bib} file.

\section{Mathematical Expressions}
Here are some example equations:

\subsection{Quadratic Formula}
The quadratic formula is given by:
\[x = \frac{-b \pm \sqrt{b^2 - 4ac}}{2a}\]

\subsection{Matrix Example}
A sample matrix:
\[
\begin{pmatrix}
a_{11} & a_{12} & a_{13} \\
a_{21} & a_{22} & a_{23} \\
a_{31} & a_{32} & a_{33}
\end{pmatrix}
\]

\section{Tables}
Below is a sample table showing data:

\begin{table}[h]
\centering
\begin{tabular}{@{}lrrr@{}}
\toprule
\textbf{Category} & \textbf{Value 1} & \textbf{Value 2} & \textbf{Total} \\
\midrule
Item A & 10.5 & 20.3 & 30.8 \\
Item B & 15.2 & 25.1 & 40.3 \\
Item C & 22.1 & 12.7 & 34.8 \\
\bottomrule
\end{tabular}
\caption{Sample Data Table}
\label{tab:sample}
\end{table}

\section{Lists}
\subsection{Bullet Points}
\begin{itemize}
  \item First item
  \item Second item
  \item Third item with sub-items:
    \begin{itemize}
      \item Sub-item 1
      \item Sub-item 2
    \end{itemize}
\end{itemize}

\subsection{Numbered List}
\begin{enumerate}
  \item First step
  \item Second step
  \item Third step
\end{enumerate}

\section{Citations and References}
When you use the citation search feature, references will be automatically managed. For example, you might cite a paper like this: \cite{example2024}. All citations will be automatically added to your \texttt{cite.bib} file.

\section{Conclusion}
This template demonstrates basic \LaTeX{} features including mathematical formulas, tables, and lists. The bibliography system is now set up to work seamlessly with Bibby AI's citation management. Feel free to modify it for your needs!

\bibliographystyle{plain}
\bibliography{cite}

\end{document}\documentclass{article}
\usepackage[utf8]{inputenc}
\usepackage{booktabs}
\usepackage{amsmath}
\usepackage{natbib}

\title{My Document}
\author{Your Name}
\date{\today}

\begin{document}

\maketitle

\section{Introduction}
This is a sample document demonstrating various \LaTeX{} features including mathematical equations, tables, and formatting. You can cite papers using the citation search feature, and they will be automatically added to your \texttt{cite.bib} file.

\section{Mathematical Expressions}
Here are some example equations:

\subsection{Quadratic Formula}
The quadratic formula is given by:
\[x = \frac{-b \pm \sqrt{b^2 - 4ac}}{2a}\]

\subsection{Matrix Example}
A sample matrix:
\[
\begin{pmatrix}
a_{11} & a_{12} & a_{13} \\
a_{21} & a_{22} & a_{23} \\
a_{31} & a_{32} & a_{33}
\end{pmatrix}
\]

\section{Tables}
Below is a sample table showing data:

\begin{table}[h]
\centering
\begin{tabular}{@{}lrrr@{}}
\toprule
\textbf{Category} & \textbf{Value 1} & \textbf{Value 2} & \textbf{Total} \\
\midrule
Item A & 10.5 & 20.3 & 30.8 \\
Item B & 15.2 & 25.1 & 40.3 \\
Item C & 22.1 & 12.7 & 34.8 \\
\bottomrule
\end{tabular}
\caption{Sample Data Table}
\label{tab:sample}
\end{table}

\section{Lists}
\subsection{Bullet Points}
\begin{itemize}
  \item First item
  \item Second item
  \item Third item with sub-items:
    \begin{itemize}
      \item Sub-item 1
      \item Sub-item 2
    \end{itemize}
\end{itemize}

\subsection{Numbered List}
\begin{enumerate}
  \item First step
  \item Second step
  \item Third step
\end{enumerate}

\section{Citations and References}
When you use the citation search feature, references will be automatically managed. For example, you might cite a paper like this: \cite{example2024}. All citations will be automatically added to your \texttt{cite.bib} file.

\section{Conclusion}
This template demonstrates basic \LaTeX{} features including mathematical formulas, tables, and lists. The bibliography system is now set up to work seamlessly with Bibby AI's citation management. Feel free to modify it for your needs!

\bibliographystyle{plain}
\bibliography{cite}

\end{document}\documentclass{article}
\usepackage[utf8]{inputenc}
\usepackage{booktabs}
\usepackage{amsmath}
\usepackage{natbib}

\title{My Document}
\author{Your Name}
\date{\today}

\begin{document}

\maketitle

\section{Introduction}
This is a sample document demonstrating various \LaTeX{} features including mathematical equations, tables, and formatting. You can cite papers using the citation search feature, and they will be automatically added to your \texttt{cite.bib} file.

\section{Mathematical Expressions}
Here are some example equations:

\subsection{Quadratic Formula}
The quadratic formula is given by:
\[x = \frac{-b \pm \sqrt{b^2 - 4ac}}{2a}\]

\subsection{Matrix Example}
A sample matrix:
\[
\begin{pmatrix}
a_{11} & a_{12} & a_{13} \\
a_{21} & a_{22} & a_{23} \\
a_{31} & a_{32} & a_{33}
\end{pmatrix}
\]

\section{Tables}
Below is a sample table showing data:

\begin{table}[h]
\centering
\begin{tabular}{@{}lrrr@{}}
\toprule
\textbf{Category} & \textbf{Value 1} & \textbf{Value 2} & \textbf{Total} \\
\midrule
Item A & 10.5 & 20.3 & 30.8 \\
Item B & 15.2 & 25.1 & 40.3 \\
Item C & 22.1 & 12.7 & 34.8 \\
\bottomrule
\end{tabular}
\caption{Sample Data Table}
\label{tab:sample}
\end{table}

\section{Lists}
\subsection{Bullet Points}
\begin{itemize}
  \item First item
  \item Second item
  \item Third item with sub-items:
    \begin{itemize}
      \item Sub-item 1
      \item Sub-item 2
    \end{itemize}
\end{itemize}

\subsection{Numbered List}
\begin{enumerate}
  \item First step
  \item Second step
  \item Third step
\end{enumerate}

\section{Citations and References}
When you use the citation search feature, references will be automatically managed. For example, you might cite a paper like this: \cite{example2024}. All citations will be automatically added to your \texttt{cite.bib} file.

\section{Conclusion}
This template demonstrates basic \LaTeX{} features including mathematical formulas, tables, and lists. The bibliography system is now set up to work seamlessly with Bibby AI's citation management. Feel free to modify it for your needs!

\bibliographystyle{plain}
\bibliography{cite}

\end{document}\documentclass{article}
\usepackage[utf8]{inputenc}
\usepackage{booktabs}
\usepackage{amsmath}
\usepackage{natbib}

\title{My Document}
\author{Your Name}
\date{\today}

\begin{document}

\maketitle

\section{Introduction}
This is a sample document demonstrating various \LaTeX{} features including mathematical equations, tables, and formatting. You can cite papers using the citation search feature, and they will be automatically added to your \texttt{cite.bib} file.

\section{Mathematical Expressions}
Here are some example equations:

\subsection{Quadratic Formula}
The quadratic formula is given by:
\[x = \frac{-b \pm \sqrt{b^2 - 4ac}}{2a}\]

\subsection{Matrix Example}
A sample matrix:
\[
\begin{pmatrix}
a_{11} & a_{12} & a_{13} \\
a_{21} & a_{22} & a_{23} \\
a_{31} & a_{32} & a_{33}
\end{pmatrix}
\]

\section{Tables}
Below is a sample table showing data:

\begin{table}[h]
\centering
\begin{tabular}{@{}lrrr@{}}
\toprule
\textbf{Category} & \textbf{Value 1} & \textbf{Value 2} & \textbf{Total} \\
\midrule
Item A & 10.5 & 20.3 & 30.8 \\
Item B & 15.2 & 25.1 & 40.3 \\
Item C & 22.1 & 12.7 & 34.8 \\
\bottomrule
\end{tabular}
\caption{Sample Data Table}
\label{tab:sample}
\end{table}

\section{Lists}
\subsection{Bullet Points}
\begin{itemize}
  \item First item
  \item Second item
  \item Third item with sub-items:
    \begin{itemize}
      \item Sub-item 1
      \item Sub-item 2
    \end{itemize}
\end{itemize}

\subsection{Numbered List}
\begin{enumerate}
  \item First step
  \item Second step
  \item Third step
\end{enumerate}

\section{Citations and References}
When you use the citation search feature, references will be automatically managed. For example, you might cite a paper like this: \cite{example2024}. All citations will be automatically added to your \texttt{cite.bib} file.

\section{Conclusion}
This template demonstrates basic \LaTeX{} features including mathematical formulas, tables, and lists. The bibliography system is now set up to work seamlessly with Bibby AI's citation management. Feel free to modify it for your needs!

\bibliographystyle{plain}
\bibliography{cite}

\end{document}\documentclass{article}
\usepackage[utf8]{inputenc}
\usepackage{booktabs}
\usepackage{amsmath}
\usepackage{natbib}

\title{My Document}
\author{Your Name}
\date{\today}

\begin{document}

\maketitle

\section{Introduction}
This is a sample document demonstrating various \LaTeX{} features including mathematical equations, tables, and formatting. You can cite papers using the citation search feature, and they will be automatically added to your \texttt{cite.bib} file.

\section{Mathematical Expressions}
Here are some example equations:

\subsection{Quadratic Formula}
The quadratic formula is given by:
\[x = \frac{-b \pm \sqrt{b^2 - 4ac}}{2a}\]

\subsection{Matrix Example}
A sample matrix:
\[
\begin{pmatrix}
a_{11} & a_{12} & a_{13} \\
a_{21} & a_{22} & a_{23} \\
a_{31} & a_{32} & a_{33}
\end{pmatrix}
\]

\section{Tables}
Below is a sample table showing data:

\begin{table}[h]
\centering
\begin{tabular}{@{}lrrr@{}}
\toprule
\textbf{Category} & \textbf{Value 1} & \textbf{Value 2} & \textbf{Total} \\
\midrule
Item A & 10.5 & 20.3 & 30.8 \\
Item B & 15.2 & 25.1 & 40.3 \\
Item C & 22.1 & 12.7 & 34.8 \\
\bottomrule
\end{tabular}
\caption{Sample Data Table}
\label{tab:sample}
\end{table}

\section{Lists}
\subsection{Bullet Points}
\begin{itemize}
  \item First item
  \item Second item
  \item Third item with sub-items:
    \begin{itemize}
      \item Sub-item 1
      \item Sub-item 2
    \end{itemize}
\end{itemize}

\subsection{Numbered List}
\begin{enumerate}
  \item First step
  \item Second step
  \item Third step
\end{enumerate}

\section{Citations and References}
When you use the citation search feature, references will be automatically managed. For example, you might cite a paper like this: \cite{example2024}. All citations will be automatically added to your \texttt{cite.bib} file.

\section{Conclusion}
This template demonstrates basic \LaTeX{} features including mathematical formulas, tables, and lists. The bibliography system is now set up to work seamlessly with Bibby AI's citation management. Feel free to modify it for your needs!

\bibliographystyle{plain}
\bibliography{cite}

\end{document}\documentclass{article}
\usepackage[utf8]{inputenc}
\usepackage{booktabs}
\usepackage{amsmath}
\usepackage{natbib}

\title{My Document}
\author{Your Name}
\date{\today}

\begin{document}

\maketitle

\section{Introduction}
This is a sample document demonstrating various \LaTeX{} features including mathematical equations, tables, and formatting. You can cite papers using the citation search feature, and they will be automatically added to your \texttt{cite.bib} file.

\section{Mathematical Expressions}
Here are some example equations:

\subsection{Quadratic Formula}
The quadratic formula is given by:
\[x = \frac{-b \pm \sqrt{b^2 - 4ac}}{2a}\]

\subsection{Matrix Example}
A sample matrix:
\[
\begin{pmatrix}
a_{11} & a_{12} & a_{13} \\
a_{21} & a_{22} & a_{23} \\
a_{31} & a_{32} & a_{33}
\end{pmatrix}
\]

\section{Tables}
Below is a sample table showing data:

\begin{table}[h]
\centering
\begin{tabular}{@{}lrrr@{}}
\toprule
\textbf{Category} & \textbf{Value 1} & \textbf{Value 2} & \textbf{Total} \\
\midrule
Item A & 10.5 & 20.3 & 30.8 \\
Item B & 15.2 & 25.1 & 40.3 \\
Item C & 22.1 & 12.7 & 34.8 \\
\bottomrule
\end{tabular}
\caption{Sample Data Table}
\label{tab:sample}
\end{table}

\section{Lists}
\subsection{Bullet Points}
\begin{itemize}
  \item First item
  \item Second item
  \item Third item with sub-items:
    \begin{itemize}
      \item Sub-item 1
      \item Sub-item 2
    \end{itemize}
\end{itemize}

\subsection{Numbered List}
\begin{enumerate}
  \item First step
  \item Second step
  \item Third step
\end{enumerate}

\section{Citations and References}
When you use the citation search feature, references will be automatically managed. For example, you might cite a paper like this: \cite{example2024}. All citations will be automatically added to your \texttt{cite.bib} file.

\section{Conclusion}
This template demonstrates basic \LaTeX{} features including mathematical formulas, tables, and lists. The bibliography system is now set up to work seamlessly with Bibby AI's citation management. Feel free to modify it for your needs!

\bibliographystyle{plain}
\bibliography{cite}

\end{document}\documentclass{article}
\usepackage[utf8]{inputenc}
\usepackage{booktabs}
\usepackage{amsmath}
\usepackage{natbib}

\title{My Document}
\author{Your Name}
\date{\today}

\begin{document}

\maketitle

\section{Introduction}
This is a sample document demonstrating various \LaTeX{} features including mathematical equations, tables, and formatting. You can cite papers using the citation search feature, and they will be automatically added to your \texttt{cite.bib} file.

\section{Mathematical Expressions}
Here are some example equations:

\subsection{Quadratic Formula}
The quadratic formula is given by:
\[x = \frac{-b \pm \sqrt{b^2 - 4ac}}{2a}\]

\subsection{Matrix Example}
A sample matrix:
\[
\begin{pmatrix}
a_{11} & a_{12} & a_{13} \\
a_{21} & a_{22} & a_{23} \\
a_{31} & a_{32} & a_{33}
\end{pmatrix}
\]

\section{Tables}
Below is a sample table showing data:

\begin{table}[h]
\centering
\begin{tabular}{@{}lrrr@{}}
\toprule
\textbf{Category} & \textbf{Value 1} & \textbf{Value 2} & \textbf{Total} \\
\midrule
Item A & 10.5 & 20.3 & 30.8 \\
Item B & 15.2 & 25.1 & 40.3 \\
Item C & 22.1 & 12.7 & 34.8 \\
\bottomrule
\end{tabular}
\caption{Sample Data Table}
\label{tab:sample}
\end{table}

\section{Lists}
\subsection{Bullet Points}
\begin{itemize}
  \item First item
  \item Second item
  \item Third item with sub-items:
    \begin{itemize}
      \item Sub-item 1
      \item Sub-item 2
    \end{itemize}
\end{itemize}

\subsection{Numbered List}
\begin{enumerate}
  \item First step
  \item Second step
  \item Third step
\end{enumerate}

\section{Citations and References}
When you use the citation search feature, references will be automatically managed. For example, you might cite a paper like this: \cite{example2024}. All citations will be automatically added to your \texttt{cite.bib} file.

\section{Conclusion}
This template demonstrates basic \LaTeX{} features including mathematical formulas, tables, and lists. The bibliography system is now set up to work seamlessly with Bibby AI's citation management. Feel free to modify it for your needs!

\bibliographystyle{plain}
\bibliography{cite}

\end{document}\documentclass{article}
\usepackage[utf8]{inputenc}
\usepackage{booktabs}
\usepackage{amsmath}
\usepackage{natbib}

\title{My Document}
\author{Your Name}
\date{\today}

\begin{document}

\maketitle

\section{Introduction}
This is a sample document demonstrating various \LaTeX{} features including mathematical equations, tables, and formatting. You can cite papers using the citation search feature, and they will be automatically added to your \texttt{cite.bib} file.

\section{Mathematical Expressions}
Here are some example equations:

\subsection{Quadratic Formula}
The quadratic formula is given by:
\[x = \frac{-b \pm \sqrt{b^2 - 4ac}}{2a}\]

\subsection{Matrix Example}
A sample matrix:
\[
\begin{pmatrix}
a_{11} & a_{12} & a_{13} \\
a_{21} & a_{22} & a_{23} \\
a_{31} & a_{32} & a_{33}
\end{pmatrix}
\]

\section{Tables}
Below is a sample table showing data:

\begin{table}[h]
\centering
\begin{tabular}{@{}lrrr@{}}
\toprule
\textbf{Category} & \textbf{Value 1} & \textbf{Value 2} & \textbf{Total} \\
\midrule
Item A & 10.5 & 20.3 & 30.8 \\
Item B & 15.2 & 25.1 & 40.3 \\
Item C & 22.1 & 12.7 & 34.8 \\
\bottomrule
\end{tabular}
\caption{Sample Data Table}
\label{tab:sample}
\end{table}

\section{Lists}
\subsection{Bullet Points}
\begin{itemize}
  \item First item
  \item Second item
  \item Third item with sub-items:
    \begin{itemize}
      \item Sub-item 1
      \item Sub-item 2
    \end{itemize}
\end{itemize}

\subsection{Numbered List}
\begin{enumerate}
  \item First step
  \item Second step
  \item Third step
\end{enumerate}

\section{Citations and References}
When you use the citation search feature, references will be automatically managed. For example, you might cite a paper like this: \cite{example2024}. All citations will be automatically added to your \texttt{cite.bib} file.

\section{Conclusion}
This template demonstrates basic \LaTeX{} features including mathematical formulas, tables, and lists. The bibliography system is now set up to work seamlessly with Bibby AI's citation management. Feel free to modify it for your needs!

\bibliographystyle{plain}
\bibliography{cite}

\end{document}\documentclass{article}
\usepackage[utf8]{inputenc}
\usepackage{booktabs}
\usepackage{amsmath}
\usepackage{natbib}

\title{My Document}
\author{Your Name}
\date{\today}

\begin{document}

\maketitle

\section{Introduction}
This is a sample document demonstrating various \LaTeX{} features including mathematical equations, tables, and formatting. You can cite papers using the citation search feature, and they will be automatically added to your \texttt{cite.bib} file.

\section{Mathematical Expressions}
Here are some example equations:

\subsection{Quadratic Formula}
The quadratic formula is given by:
\[x = \frac{-b \pm \sqrt{b^2 - 4ac}}{2a}\]

\subsection{Matrix Example}
A sample matrix:
\[
\begin{pmatrix}
a_{11} & a_{12} & a_{13} \\
a_{21} & a_{22} & a_{23} \\
a_{31} & a_{32} & a_{33}
\end{pmatrix}
\]

\section{Tables}
Below is a sample table showing data:

\begin{table}[h]
\centering
\begin{tabular}{@{}lrrr@{}}
\toprule
\textbf{Category} & \textbf{Value 1} & \textbf{Value 2} & \textbf{Total} \\
\midrule
Item A & 10.5 & 20.3 & 30.8 \\
Item B & 15.2 & 25.1 & 40.3 \\
Item C & 22.1 & 12.7 & 34.8 \\
\bottomrule
\end{tabular}
\caption{Sample Data Table}
\label{tab:sample}
\end{table}

\section{Lists}
\subsection{Bullet Points}
\begin{itemize}
  \item First item
  \item Second item
  \item Third item with sub-items:
    \begin{itemize}
      \item Sub-item 1
      \item Sub-item 2
    \end{itemize}
\end{itemize}

\subsection{Numbered List}
\begin{enumerate}
  \item First step
  \item Second step
  \item Third step
\end{enumerate}

\section{Citations and References}
When you use the citation search feature, references will be automatically managed. For example, you might cite a paper like this: \cite{example2024}. All citations will be automatically added to your \texttt{cite.bib} file.

\section{Conclusion}
This template demonstrates basic \LaTeX{} features including mathematical formulas, tables, and lists. The bibliography system is now set up to work seamlessly with Bibby AI's citation management. Feel free to modify it for your needs!

\bibliographystyle{plain}
\bibliography{cite}

\end{document}\documentclass{article}
\usepackage[utf8]{inputenc}
\usepackage{booktabs}
\usepackage{amsmath}
\usepackage{natbib}

\title{My Document}
\author{Your Name}
\date{\today}

\begin{document}

\maketitle

\section{Introduction}
This is a sample document demonstrating various \LaTeX{} features including mathematical equations, tables, and formatting. You can cite papers using the citation search feature, and they will be automatically added to your \texttt{cite.bib} file.

\section{Mathematical Expressions}
Here are some example equations:

\subsection{Quadratic Formula}
The quadratic formula is given by:
\[x = \frac{-b \pm \sqrt{b^2 - 4ac}}{2a}\]

\subsection{Matrix Example}
A sample matrix:
\[
\begin{pmatrix}
a_{11} & a_{12} & a_{13} \\
a_{21} & a_{22} & a_{23} \\
a_{31} & a_{32} & a_{33}
\end{pmatrix}
\]

\section{Tables}
Below is a sample table showing data:

\begin{table}[h]
\centering
\begin{tabular}{@{}lrrr@{}}
\toprule
\textbf{Category} & \textbf{Value 1} & \textbf{Value 2} & \textbf{Total} \\
\midrule
Item A & 10.5 & 20.3 & 30.8 \\
Item B & 15.2 & 25.1 & 40.3 \\
Item C & 22.1 & 12.7 & 34.8 \\
\bottomrule
\end{tabular}
\caption{Sample Data Table}
\label{tab:sample}
\end{table}

\section{Lists}
\subsection{Bullet Points}
\begin{itemize}
  \item First item
  \item Second item
  \item Third item with sub-items:
    \begin{itemize}
      \item Sub-item 1
      \item Sub-item 2
    \end{itemize}
\end{itemize}

\subsection{Numbered List}
\begin{enumerate}
  \item First step
  \item Second step
  \item Third step
\end{enumerate}

\section{Citations and References}
When you use the citation search feature, references will be automatically managed. For example, you might cite a paper like this: \cite{example2024}. All citations will be automatically added to your \texttt{cite.bib} file.

\section{Conclusion}
This template demonstrates basic \LaTeX{} features including mathematical formulas, tables, and lists. The bibliography system is now set up to work seamlessly with Bibby AI's citation management. Feel free to modify it for your needs!

\bibliographystyle{plain}
\bibliography{cite}

\end{document}\documentclass{article}
\usepackage[utf8]{inputenc}
\usepackage{booktabs}
\usepackage{amsmath}
\usepackage{natbib}

\title{My Document}
\author{Your Name}
\date{\today}

\begin{document}

\maketitle

\section{Introduction}
This is a sample document demonstrating various \LaTeX{} features including mathematical equations, tables, and formatting. You can cite papers using the citation search feature, and they will be automatically added to your \texttt{cite.bib} file.

\section{Mathematical Expressions}
Here are some example equations:

\subsection{Quadratic Formula}
The quadratic formula is given by:
\[x = \frac{-b \pm \sqrt{b^2 - 4ac}}{2a}\]

\subsection{Matrix Example}
A sample matrix:
\[
\begin{pmatrix}
a_{11} & a_{12} & a_{13} \\
a_{21} & a_{22} & a_{23} \\
a_{31} & a_{32} & a_{33}
\end{pmatrix}
\]

\section{Tables}
Below is a sample table showing data:

\begin{table}[h]
\centering
\begin{tabular}{@{}lrrr@{}}
\toprule
\textbf{Category} & \textbf{Value 1} & \textbf{Value 2} & \textbf{Total} \\
\midrule
Item A & 10.5 & 20.3 & 30.8 \\
Item B & 15.2 & 25.1 & 40.3 \\
Item C & 22.1 & 12.7 & 34.8 \\
\bottomrule
\end{tabular}
\caption{Sample Data Table}
\label{tab:sample}
\end{table}

\section{Lists}
\subsection{Bullet Points}
\begin{itemize}
  \item First item
  \item Second item
  \item Third item with sub-items:
    \begin{itemize}
      \item Sub-item 1
      \item Sub-item 2
    \end{itemize}
\end{itemize}

\subsection{Numbered List}
\begin{enumerate}
  \item First step
  \item Second step
  \item Third step
\end{enumerate}

\section{Citations and References}
When you use the citation search feature, references will be automatically managed. For example, you might cite a paper like this: \cite{example2024}. All citations will be automatically added to your \texttt{cite.bib} file.

\section{Conclusion}
This template demonstrates basic \LaTeX{} features including mathematical formulas, tables, and lists. The bibliography system is now set up to work seamlessly with Bibby AI's citation management. Feel free to modify it for your needs!

\bibliographystyle{plain}
\bibliography{cite}

\end{document}\documentclass{article}
\usepackage[utf8]{inputenc}
\usepackage{booktabs}
\usepackage{amsmath}
\usepackage{natbib}

\title{My Document}
\author{Your Name}
\date{\today}

\begin{document}

\maketitle

\section{Introduction}
This is a sample document demonstrating various \LaTeX{} features including mathematical equations, tables, and formatting. You can cite papers using the citation search feature, and they will be automatically added to your \texttt{cite.bib} file.

\section{Mathematical Expressions}
Here are some example equations:

\subsection{Quadratic Formula}
The quadratic formula is given by:
\[x = \frac{-b \pm \sqrt{b^2 - 4ac}}{2a}\]

\subsection{Matrix Example}
A sample matrix:
\[
\begin{pmatrix}
a_{11} & a_{12} & a_{13} \\
a_{21} & a_{22} & a_{23} \\
a_{31} & a_{32} & a_{33}
\end{pmatrix}
\]

\section{Tables}
Below is a sample table showing data:

\begin{table}[h]
\centering
\begin{tabular}{@{}lrrr@{}}
\toprule
\textbf{Category} & \textbf{Value 1} & \textbf{Value 2} & \textbf{Total} \\
\midrule
Item A & 10.5 & 20.3 & 30.8 \\
Item B & 15.2 & 25.1 & 40.3 \\
Item C & 22.1 & 12.7 & 34.8 \\
\bottomrule
\end{tabular}
\caption{Sample Data Table}
\label{tab:sample}
\end{table}

\section{Lists}
\subsection{Bullet Points}
\begin{itemize}
  \item First item
  \item Second item
  \item Third item with sub-items:
    \begin{itemize}
      \item Sub-item 1
      \item Sub-item 2
    \end{itemize}
\end{itemize}

\subsection{Numbered List}
\begin{enumerate}
  \item First step
  \item Second step
  \item Third step
\end{enumerate}

\section{Citations and References}
When you use the citation search feature, references will be automatically managed. For example, you might cite a paper like this: \cite{example2024}. All citations will be automatically added to your \texttt{cite.bib} file.

\section{Conclusion}
This template demonstrates basic \LaTeX{} features including mathematical formulas, tables, and lists. The bibliography system is now set up to work seamlessly with Bibby AI's citation management. Feel free to modify it for your needs!

\bibliographystyle{plain}
\bibliography{cite}

\end{document}\documentclass{article}
\usepackage[utf8]{inputenc}
\usepackage{booktabs}
\usepackage{amsmath}
\usepackage{natbib}

\title{My Document}
\author{Your Name}
\date{\today}

\begin{document}

\maketitle

\section{Introduction}
This is a sample document demonstrating various \LaTeX{} features including mathematical equations, tables, and formatting. You can cite papers using the citation search feature, and they will be automatically added to your \texttt{cite.bib} file.

\section{Mathematical Expressions}
Here are some example equations:

\subsection{Quadratic Formula}
The quadratic formula is given by:
\[x = \frac{-b \pm \sqrt{b^2 - 4ac}}{2a}\]

\subsection{Matrix Example}
A sample matrix:
\[
\begin{pmatrix}
a_{11} & a_{12} & a_{13} \\
a_{21} & a_{22} & a_{23} \\
a_{31} & a_{32} & a_{33}
\end{pmatrix}
\]

\section{Tables}
Below is a sample table showing data:

\begin{table}[h]
\centering
\begin{tabular}{@{}lrrr@{}}
\toprule
\textbf{Category} & \textbf{Value 1} & \textbf{Value 2} & \textbf{Total} \\
\midrule
Item A & 10.5 & 20.3 & 30.8 \\
Item B & 15.2 & 25.1 & 40.3 \\
Item C & 22.1 & 12.7 & 34.8 \\
\bottomrule
\end{tabular}
\caption{Sample Data Table}
\label{tab:sample}
\end{table}

\section{Lists}
\subsection{Bullet Points}
\begin{itemize}
  \item First item
  \item Second item
  \item Third item with sub-items:
    \begin{itemize}
      \item Sub-item 1
      \item Sub-item 2
    \end{itemize}
\end{itemize}

\subsection{Numbered List}
\begin{enumerate}
  \item First step
  \item Second step
  \item Third step
\end{enumerate}

\section{Citations and References}
When you use the citation search feature, references will be automatically managed. For example, you might cite a paper like this: \cite{example2024}. All citations will be automatically added to your \texttt{cite.bib} file.

\section{Conclusion}
This template demonstrates basic \LaTeX{} features including mathematical formulas, tables, and lists. The bibliography system is now set up to work seamlessly with Bibby AI's citation management. Feel free to modify it for your needs!

\bibliographystyle{plain}
\bibliography{cite}

\end{document}\documentclass{article}
\usepackage[utf8]{inputenc}
\usepackage{booktabs}
\usepackage{amsmath}
\usepackage{natbib}

\title{My Document}
\author{Your Name}
\date{\today}

\begin{document}

\maketitle

\section{Introduction}
This is a sample document demonstrating various \LaTeX{} features including mathematical equations, tables, and formatting. You can cite papers using the citation search feature, and they will be automatically added to your \texttt{cite.bib} file.

\section{Mathematical Expressions}
Here are some example equations:

\subsection{Quadratic Formula}
The quadratic formula is given by:
\[x = \frac{-b \pm \sqrt{b^2 - 4ac}}{2a}\]

\subsection{Matrix Example}
A sample matrix:
\[
\begin{pmatrix}
a_{11} & a_{12} & a_{13} \\
a_{21} & a_{22} & a_{23} \\
a_{31} & a_{32} & a_{33}
\end{pmatrix}
\]

\section{Tables}
Below is a sample table showing data:

\begin{table}[h]
\centering
\begin{tabular}{@{}lrrr@{}}
\toprule
\textbf{Category} & \textbf{Value 1} & \textbf{Value 2} & \textbf{Total} \\
\midrule
Item A & 10.5 & 20.3 & 30.8 \\
Item B & 15.2 & 25.1 & 40.3 \\
Item C & 22.1 & 12.7 & 34.8 \\
\bottomrule
\end{tabular}
\caption{Sample Data Table}
\label{tab:sample}
\end{table}

\section{Lists}
\subsection{Bullet Points}
\begin{itemize}
  \item First item
  \item Second item
  \item Third item with sub-items:
    \begin{itemize}
      \item Sub-item 1
      \item Sub-item 2
    \end{itemize}
\end{itemize}

\subsection{Numbered List}
\begin{enumerate}
  \item First step
  \item Second step
  \item Third step
\end{enumerate}

\section{Citations and References}
When you use the citation search feature, references will be automatically managed. For example, you might cite a paper like this: \cite{example2024}. All citations will be automatically added to your \texttt{cite.bib} file.

\section{Conclusion}
This template demonstrates basic \LaTeX{} features including mathematical formulas, tables, and lists. The bibliography system is now set up to work seamlessly with Bibby AI's citation management. Feel free to modify it for your needs!

\bibliographystyle{plain}
\bibliography{cite}

\end{document}\documentclass{article}
\usepackage[utf8]{inputenc}
\usepackage{booktabs}
\usepackage{amsmath}
\usepackage{natbib}

\title{My Document}
\author{Your Name}
\date{\today}

\begin{document}

\maketitle

\section{Introduction}
This is a sample document demonstrating various \LaTeX{} features including mathematical equations, tables, and formatting. You can cite papers using the citation search feature, and they will be automatically added to your \texttt{cite.bib} file.

\section{Mathematical Expressions}
Here are some example equations:

\subsection{Quadratic Formula}
The quadratic formula is given by:
\[x = \frac{-b \pm \sqrt{b^2 - 4ac}}{2a}\]

\subsection{Matrix Example}
A sample matrix:
\[
\begin{pmatrix}
a_{11} & a_{12} & a_{13} \\
a_{21} & a_{22} & a_{23} \\
a_{31} & a_{32} & a_{33}
\end{pmatrix}
\]

\section{Tables}
Below is a sample table showing data:

\begin{table}[h]
\centering
\begin{tabular}{@{}lrrr@{}}
\toprule
\textbf{Category} & \textbf{Value 1} & \textbf{Value 2} & \textbf{Total} \\
\midrule
Item A & 10.5 & 20.3 & 30.8 \\
Item B & 15.2 & 25.1 & 40.3 \\
Item C & 22.1 & 12.7 & 34.8 \\
\bottomrule
\end{tabular}
\caption{Sample Data Table}
\label{tab:sample}
\end{table}

\section{Lists}
\subsection{Bullet Points}
\begin{itemize}
  \item First item
  \item Second item
  \item Third item with sub-items:
    \begin{itemize}
      \item Sub-item 1
      \item Sub-item 2
    \end{itemize}
\end{itemize}

\subsection{Numbered List}
\begin{enumerate}
  \item First step
  \item Second step
  \item Third step
\end{enumerate}

\section{Citations and References}
When you use the citation search feature, references will be automatically managed. For example, you might cite a paper like this: \cite{example2024}. All citations will be automatically added to your \texttt{cite.bib} file.

\section{Conclusion}
This template demonstrates basic \LaTeX{} features including mathematical formulas, tables, and lists. The bibliography system is now set up to work seamlessly with Bibby AI's citation management. Feel free to modify it for your needs!

\bibliographystyle{plain}
\bibliography{cite}

\end{document}\documentclass{article}
\usepackage[utf8]{inputenc}
\usepackage{booktabs}
\usepackage{amsmath}
\usepackage{natbib}

\title{My Document}
\author{Your Name}
\date{\today}

\begin{document}

\maketitle

\section{Introduction}
This is a sample document demonstrating various \LaTeX{} features including mathematical equations, tables, and formatting. You can cite papers using the citation search feature, and they will be automatically added to your \texttt{cite.bib} file.

\section{Mathematical Expressions}
Here are some example equations:

\subsection{Quadratic Formula}
The quadratic formula is given by:
\[x = \frac{-b \pm \sqrt{b^2 - 4ac}}{2a}\]

\subsection{Matrix Example}
A sample matrix:
\[
\begin{pmatrix}
a_{11} & a_{12} & a_{13} \\
a_{21} & a_{22} & a_{23} \\
a_{31} & a_{32} & a_{33}
\end{pmatrix}
\]

\section{Tables}
Below is a sample table showing data:

\begin{table}[h]
\centering
\begin{tabular}{@{}lrrr@{}}
\toprule
\textbf{Category} & \textbf{Value 1} & \textbf{Value 2} & \textbf{Total} \\
\midrule
Item A & 10.5 & 20.3 & 30.8 \\
Item B & 15.2 & 25.1 & 40.3 \\
Item C & 22.1 & 12.7 & 34.8 \\
\bottomrule
\end{tabular}
\caption{Sample Data Table}
\label{tab:sample}
\end{table}

\section{Lists}
\subsection{Bullet Points}
\begin{itemize}
  \item First item
  \item Second item
  \item Third item with sub-items:
    \begin{itemize}
      \item Sub-item 1
      \item Sub-item 2
    \end{itemize}
\end{itemize}

\subsection{Numbered List}
\begin{enumerate}
  \item First step
  \item Second step
  \item Third step
\end{enumerate}

\section{Citations and References}
When you use the citation search feature, references will be automatically managed. For example, you might cite a paper like this: \cite{example2024}. All citations will be automatically added to your \texttt{cite.bib} file.

\section{Conclusion}
This template demonstrates basic \LaTeX{} features including mathematical formulas, tables, and lists. The bibliography system is now set up to work seamlessly with Bibby AI's citation management. Feel free to modify it for your needs!

\bibliographystyle{plain}
\bibliography{cite}

\end{document}\documentclass{article}
\usepackage[utf8]{inputenc}
\usepackage{booktabs}
\usepackage{amsmath}
\usepackage{natbib}

\title{My Document}
\author{Your Name}
\date{\today}

\begin{document}

\maketitle

\section{Introduction}
This is a sample document demonstrating various \LaTeX{} features including mathematical equations, tables, and formatting. You can cite papers using the citation search feature, and they will be automatically added to your \texttt{cite.bib} file.

\section{Mathematical Expressions}
Here are some example equations:

\subsection{Quadratic Formula}
The quadratic formula is given by:
\[x = \frac{-b \pm \sqrt{b^2 - 4ac}}{2a}\]

\subsection{Matrix Example}
A sample matrix:
\[
\begin{pmatrix}
a_{11} & a_{12} & a_{13} \\
a_{21} & a_{22} & a_{23} \\
a_{31} & a_{32} & a_{33}
\end{pmatrix}
\]

\section{Tables}
Below is a sample table showing data:

\begin{table}[h]
\centering
\begin{tabular}{@{}lrrr@{}}
\toprule
\textbf{Category} & \textbf{Value 1} & \textbf{Value 2} & \textbf{Total} \\
\midrule
Item A & 10.5 & 20.3 & 30.8 \\
Item B & 15.2 & 25.1 & 40.3 \\
Item C & 22.1 & 12.7 & 34.8 \\
\bottomrule
\end{tabular}
\caption{Sample Data Table}
\label{tab:sample}
\end{table}

\section{Lists}
\subsection{Bullet Points}
\begin{itemize}
  \item First item
  \item Second item
  \item Third item with sub-items:
    \begin{itemize}
      \item Sub-item 1
      \item Sub-item 2
    \end{itemize}
\end{itemize}

\subsection{Numbered List}
\begin{enumerate}
  \item First step
  \item Second step
  \item Third step
\end{enumerate}

\section{Citations and References}
When you use the citation search feature, references will be automatically managed. For example, you might cite a paper like this: \cite{example2024}. All citations will be automatically added to your \texttt{cite.bib} file.

\section{Conclusion}
This template demonstrates basic \LaTeX{} features including mathematical formulas, tables, and lists. The bibliography system is now set up to work seamlessly with Bibby AI's citation management. Feel free to modify it for your needs!

\bibliographystyle{plain}
\bibliography{cite}

\end{document}\documentclass{article}
\usepackage[utf8]{inputenc}
\usepackage{booktabs}
\usepackage{amsmath}
\usepackage{natbib}

\title{My Document}
\author{Your Name}
\date{\today}

\begin{document}

\maketitle

\section{Introduction}
This is a sample document demonstrating various \LaTeX{} features including mathematical equations, tables, and formatting. You can cite papers using the citation search feature, and they will be automatically added to your \texttt{cite.bib} file.

\section{Mathematical Expressions}
Here are some example equations:

\subsection{Quadratic Formula}
The quadratic formula is given by:
\[x = \frac{-b \pm \sqrt{b^2 - 4ac}}{2a}\]

\subsection{Matrix Example}
A sample matrix:
\[
\begin{pmatrix}
a_{11} & a_{12} & a_{13} \\
a_{21} & a_{22} & a_{23} \\
a_{31} & a_{32} & a_{33}
\end{pmatrix}
\]

\section{Tables}
Below is a sample table showing data:

\begin{table}[h]
\centering
\begin{tabular}{@{}lrrr@{}}
\toprule
\textbf{Category} & \textbf{Value 1} & \textbf{Value 2} & \textbf{Total} \\
\midrule
Item A & 10.5 & 20.3 & 30.8 \\
Item B & 15.2 & 25.1 & 40.3 \\
Item C & 22.1 & 12.7 & 34.8 \\
\bottomrule
\end{tabular}
\caption{Sample Data Table}
\label{tab:sample}
\end{table}

\section{Lists}
\subsection{Bullet Points}
\begin{itemize}
  \item First item
  \item Second item
  \item Third item with sub-items:
    \begin{itemize}
      \item Sub-item 1
      \item Sub-item 2
    \end{itemize}
\end{itemize}

\subsection{Numbered List}
\begin{enumerate}
  \item First step
  \item Second step
  \item Third step
\end{enumerate}

\section{Citations and References}
When you use the citation search feature, references will be automatically managed. For example, you might cite a paper like this: \cite{example2024}. All citations will be automatically added to your \texttt{cite.bib} file.

\section{Conclusion}
This template demonstrates basic \LaTeX{} features including mathematical formulas, tables, and lists. The bibliography system is now set up to work seamlessly with Bibby AI's citation management. Feel free to modify it for your needs!

\bibliographystyle{plain}
\bibliography{cite}

\end{document}\documentclass{article}
\usepackage[utf8]{inputenc}
\usepackage{booktabs}
\usepackage{amsmath}
\usepackage{natbib}

\title{My Document}
\author{Your Name}
\date{\today}

\begin{document}

\maketitle

\section{Introduction}
This is a sample document demonstrating various \LaTeX{} features including mathematical equations, tables, and formatting. You can cite papers using the citation search feature, and they will be automatically added to your \texttt{cite.bib} file.

\section{Mathematical Expressions}
Here are some example equations:

\subsection{Quadratic Formula}
The quadratic formula is given by:
\[x = \frac{-b \pm \sqrt{b^2 - 4ac}}{2a}\]

\subsection{Matrix Example}
A sample matrix:
\[
\begin{pmatrix}
a_{11} & a_{12} & a_{13} \\
a_{21} & a_{22} & a_{23} \\
a_{31} & a_{32} & a_{33}
\end{pmatrix}
\]

\section{Tables}
Below is a sample table showing data:

\begin{table}[h]
\centering
\begin{tabular}{@{}lrrr@{}}
\toprule
\textbf{Category} & \textbf{Value 1} & \textbf{Value 2} & \textbf{Total} \\
\midrule
Item A & 10.5 & 20.3 & 30.8 \\
Item B & 15.2 & 25.1 & 40.3 \\
Item C & 22.1 & 12.7 & 34.8 \\
\bottomrule
\end{tabular}
\caption{Sample Data Table}
\label{tab:sample}
\end{table}

\section{Lists}
\subsection{Bullet Points}
\begin{itemize}
  \item First item
  \item Second item
  \item Third item with sub-items:
    \begin{itemize}
      \item Sub-item 1
      \item Sub-item 2
    \end{itemize}
\end{itemize}

\subsection{Numbered List}
\begin{enumerate}
  \item First step
  \item Second step
  \item Third step
\end{enumerate}

\section{Citations and References}
When you use the citation search feature, references will be automatically managed. For example, you might cite a paper like this: \cite{example2024}. All citations will be automatically added to your \texttt{cite.bib} file.

\section{Conclusion}
This template demonstrates basic \LaTeX{} features including mathematical formulas, tables, and lists. The bibliography system is now set up to work seamlessly with Bibby AI's citation management. Feel free to modify it for your needs!

\bibliographystyle{plain}
\bibliography{cite}

\end{document}\documentclass{article}
\usepackage[utf8]{inputenc}
\usepackage{booktabs}
\usepackage{amsmath}
\usepackage{natbib}

\title{My Document}
\author{Your Name}
\date{\today}

\begin{document}

\maketitle

\section{Introduction}
This is a sample document demonstrating various \LaTeX{} features including mathematical equations, tables, and formatting. You can cite papers using the citation search feature, and they will be automatically added to your \texttt{cite.bib} file.

\section{Mathematical Expressions}
Here are some example equations:

\subsection{Quadratic Formula}
The quadratic formula is given by:
\[x = \frac{-b \pm \sqrt{b^2 - 4ac}}{2a}\]

\subsection{Matrix Example}
A sample matrix:
\[
\begin{pmatrix}
a_{11} & a_{12} & a_{13} \\
a_{21} & a_{22} & a_{23} \\
a_{31} & a_{32} & a_{33}
\end{pmatrix}
\]

\section{Tables}
Below is a sample table showing data:

\begin{table}[h]
\centering
\begin{tabular}{@{}lrrr@{}}
\toprule
\textbf{Category} & \textbf{Value 1} & \textbf{Value 2} & \textbf{Total} \\
\midrule
Item A & 10.5 & 20.3 & 30.8 \\
Item B & 15.2 & 25.1 & 40.3 \\
Item C & 22.1 & 12.7 & 34.8 \\
\bottomrule
\end{tabular}
\caption{Sample Data Table}
\label{tab:sample}
\end{table}

\section{Lists}
\subsection{Bullet Points}
\begin{itemize}
  \item First item
  \item Second item
  \item Third item with sub-items:
    \begin{itemize}
      \item Sub-item 1
      \item Sub-item 2
    \end{itemize}
\end{itemize}

\subsection{Numbered List}
\begin{enumerate}
  \item First step
  \item Second step
  \item Third step
\end{enumerate}

\section{Citations and References}
When you use the citation search feature, references will be automatically managed. For example, you might cite a paper like this: \cite{example2024}. All citations will be automatically added to your \texttt{cite.bib} file.

\section{Conclusion}
This template demonstrates basic \LaTeX{} features including mathematical formulas, tables, and lists. The bibliography system is now set up to work seamlessly with Bibby AI's citation management. Feel free to modify it for your needs!

\bibliographystyle{plain}
\bibliography{cite}

\end{document}\documentclass{article}
\usepackage[utf8]{inputenc}
\usepackage{booktabs}
\usepackage{amsmath}
\usepackage{natbib}

\title{My Document}
\author{Your Name}
\date{\today}

\begin{document}

\maketitle

\section{Introduction}
This is a sample document demonstrating various \LaTeX{} features including mathematical equations, tables, and formatting. You can cite papers using the citation search feature, and they will be automatically added to your \texttt{cite.bib} file.

\section{Mathematical Expressions}
Here are some example equations:

\subsection{Quadratic Formula}
The quadratic formula is given by:
\[x = \frac{-b \pm \sqrt{b^2 - 4ac}}{2a}\]

\subsection{Matrix Example}
A sample matrix:
\[
\begin{pmatrix}
a_{11} & a_{12} & a_{13} \\
a_{21} & a_{22} & a_{23} \\
a_{31} & a_{32} & a_{33}
\end{pmatrix}
\]

\section{Tables}
Below is a sample table showing data:

\begin{table}[h]
\centering
\begin{tabular}{@{}lrrr@{}}
\toprule
\textbf{Category} & \textbf{Value 1} & \textbf{Value 2} & \textbf{Total} \\
\midrule
Item A & 10.5 & 20.3 & 30.8 \\
Item B & 15.2 & 25.1 & 40.3 \\
Item C & 22.1 & 12.7 & 34.8 \\
\bottomrule
\end{tabular}
\caption{Sample Data Table}
\label{tab:sample}
\end{table}

\section{Lists}
\subsection{Bullet Points}
\begin{itemize}
  \item First item
  \item Second item
  \item Third item with sub-items:
    \begin{itemize}
      \item Sub-item 1
      \item Sub-item 2
    \end{itemize}
\end{itemize}

\subsection{Numbered List}
\begin{enumerate}
  \item First step
  \item Second step
  \item Third step
\end{enumerate}

\section{Citations and References}
When you use the citation search feature, references will be automatically managed. For example, you might cite a paper like this: \cite{example2024}. All citations will be automatically added to your \texttt{cite.bib} file.

\section{Conclusion}
This template demonstrates basic \LaTeX{} features including mathematical formulas, tables, and lists. The bibliography system is now set up to work seamlessly with Bibby AI's citation management. Feel free to modify it for your needs!

\bibliographystyle{plain}
\bibliography{cite}

\end{document}\documentclass{article}
\usepackage[utf8]{inputenc}
\usepackage{booktabs}
\usepackage{amsmath}
\usepackage{natbib}

\title{My Document}
\author{Your Name}
\date{\today}

\begin{document}

\maketitle

\section{Introduction}
This is a sample document demonstrating various \LaTeX{} features including mathematical equations, tables, and formatting. You can cite papers using the citation search feature, and they will be automatically added to your \texttt{cite.bib} file.

\section{Mathematical Expressions}
Here are some example equations:

\subsection{Quadratic Formula}
The quadratic formula is given by:
\[x = \frac{-b \pm \sqrt{b^2 - 4ac}}{2a}\]

\subsection{Matrix Example}
A sample matrix:
\[
\begin{pmatrix}
a_{11} & a_{12} & a_{13} \\
a_{21} & a_{22} & a_{23} \\
a_{31} & a_{32} & a_{33}
\end{pmatrix}
\]

\section{Tables}
Below is a sample table showing data:

\begin{table}[h]
\centering
\begin{tabular}{@{}lrrr@{}}
\toprule
\textbf{Category} & \textbf{Value 1} & \textbf{Value 2} & \textbf{Total} \\
\midrule
Item A & 10.5 & 20.3 & 30.8 \\
Item B & 15.2 & 25.1 & 40.3 \\
Item C & 22.1 & 12.7 & 34.8 \\
\bottomrule
\end{tabular}
\caption{Sample Data Table}
\label{tab:sample}
\end{table}

\section{Lists}
\subsection{Bullet Points}
\begin{itemize}
  \item First item
  \item Second item
  \item Third item with sub-items:
    \begin{itemize}
      \item Sub-item 1
      \item Sub-item 2
    \end{itemize}
\end{itemize}

\subsection{Numbered List}
\begin{enumerate}
  \item First step
  \item Second step
  \item Third step
\end{enumerate}

\section{Citations and References}
When you use the citation search feature, references will be automatically managed. For example, you might cite a paper like this: \cite{example2024}. All citations will be automatically added to your \texttt{cite.bib} file.

\section{Conclusion}
This template demonstrates basic \LaTeX{} features including mathematical formulas, tables, and lists. The bibliography system is now set up to work seamlessly with Bibby AI's citation management. Feel free to modify it for your needs!

\bibliographystyle{plain}
\bibliography{cite}

\end{document}\documentclass{article}
\usepackage[utf8]{inputenc}
\usepackage{booktabs}
\usepackage{amsmath}
\usepackage{natbib}

\title{My Document}
\author{Your Name}
\date{\today}

\begin{document}

\maketitle

\section{Introduction}
This is a sample document demonstrating various \LaTeX{} features including mathematical equations, tables, and formatting. You can cite papers using the citation search feature, and they will be automatically added to your \texttt{cite.bib} file.

\section{Mathematical Expressions}
Here are some example equations:

\subsection{Quadratic Formula}
The quadratic formula is given by:
\[x = \frac{-b \pm \sqrt{b^2 - 4ac}}{2a}\]

\subsection{Matrix Example}
A sample matrix:
\[
\begin{pmatrix}
a_{11} & a_{12} & a_{13} \\
a_{21} & a_{22} & a_{23} \\
a_{31} & a_{32} & a_{33}
\end{pmatrix}
\]

\section{Tables}
Below is a sample table showing data:

\begin{table}[h]
\centering
\begin{tabular}{@{}lrrr@{}}
\toprule
\textbf{Category} & \textbf{Value 1} & \textbf{Value 2} & \textbf{Total} \\
\midrule
Item A & 10.5 & 20.3 & 30.8 \\
Item B & 15.2 & 25.1 & 40.3 \\
Item C & 22.1 & 12.7 & 34.8 \\
\bottomrule
\end{tabular}
\caption{Sample Data Table}
\label{tab:sample}
\end{table}

\section{Lists}
\subsection{Bullet Points}
\begin{itemize}
  \item First item
  \item Second item
  \item Third item with sub-items:
    \begin{itemize}
      \item Sub-item 1
      \item Sub-item 2
    \end{itemize}
\end{itemize}

\subsection{Numbered List}
\begin{enumerate}
  \item First step
  \item Second step
  \item Third step
\end{enumerate}

\section{Citations and References}
When you use the citation search feature, references will be automatically managed. For example, you might cite a paper like this: \cite{example2024}. All citations will be automatically added to your \texttt{cite.bib} file.

\section{Conclusion}
This template demonstrates basic \LaTeX{} features including mathematical formulas, tables, and lists. The bibliography system is now set up to work seamlessly with Bibby AI's citation management. Feel free to modify it for your needs!

\bibliographystyle{plain}
\bibliography{cite}

\end{document}\documentclass{article}
\usepackage[utf8]{inputenc}
\usepackage{booktabs}
\usepackage{amsmath}
\usepackage{natbib}

\title{My Document}
\author{Your Name}
\date{\today}

\begin{document}

\maketitle

\section{Introduction}
This is a sample document demonstrating various \LaTeX{} features including mathematical equations, tables, and formatting. You can cite papers using the citation search feature, and they will be automatically added to your \texttt{cite.bib} file.

\section{Mathematical Expressions}
Here are some example equations:

\subsection{Quadratic Formula}
The quadratic formula is given by:
\[x = \frac{-b \pm \sqrt{b^2 - 4ac}}{2a}\]

\subsection{Matrix Example}
A sample matrix:
\[
\begin{pmatrix}
a_{11} & a_{12} & a_{13} \\
a_{21} & a_{22} & a_{23} \\
a_{31} & a_{32} & a_{33}
\end{pmatrix}
\]

\section{Tables}
Below is a sample table showing data:

\begin{table}[h]
\centering
\begin{tabular}{@{}lrrr@{}}
\toprule
\textbf{Category} & \textbf{Value 1} & \textbf{Value 2} & \textbf{Total} \\
\midrule
Item A & 10.5 & 20.3 & 30.8 \\
Item B & 15.2 & 25.1 & 40.3 \\
Item C & 22.1 & 12.7 & 34.8 \\
\bottomrule
\end{tabular}
\caption{Sample Data Table}
\label{tab:sample}
\end{table}

\section{Lists}
\subsection{Bullet Points}
\begin{itemize}
  \item First item
  \item Second item
  \item Third item with sub-items:
    \begin{itemize}
      \item Sub-item 1
      \item Sub-item 2
    \end{itemize}
\end{itemize}

\subsection{Numbered List}
\begin{enumerate}
  \item First step
  \item Second step
  \item Third step
\end{enumerate}

\section{Citations and References}
When you use the citation search feature, references will be automatically managed. For example, you might cite a paper like this: \cite{example2024}. All citations will be automatically added to your \texttt{cite.bib} file.

\section{Conclusion}
This template demonstrates basic \LaTeX{} features including mathematical formulas, tables, and lists. The bibliography system is now set up to work seamlessly with Bibby AI's citation management. Feel free to modify it for your needs!

\bibliographystyle{plain}
\bibliography{cite}

\end{document}\documentclass{article}
\usepackage[utf8]{inputenc}
\usepackage{booktabs}
\usepackage{amsmath}
\usepackage{natbib}

\title{My Document}
\author{Your Name}
\date{\today}

\begin{document}

\maketitle

\section{Introduction}
This is a sample document demonstrating various \LaTeX{} features including mathematical equations, tables, and formatting. You can cite papers using the citation search feature, and they will be automatically added to your \texttt{cite.bib} file.

\section{Mathematical Expressions}
Here are some example equations:

\subsection{Quadratic Formula}
The quadratic formula is given by:
\[x = \frac{-b \pm \sqrt{b^2 - 4ac}}{2a}\]

\subsection{Matrix Example}
A sample matrix:
\[
\begin{pmatrix}
a_{11} & a_{12} & a_{13} \\
a_{21} & a_{22} & a_{23} \\
a_{31} & a_{32} & a_{33}
\end{pmatrix}
\]

\section{Tables}
Below is a sample table showing data:

\begin{table}[h]
\centering
\begin{tabular}{@{}lrrr@{}}
\toprule
\textbf{Category} & \textbf{Value 1} & \textbf{Value 2} & \textbf{Total} \\
\midrule
Item A & 10.5 & 20.3 & 30.8 \\
Item B & 15.2 & 25.1 & 40.3 \\
Item C & 22.1 & 12.7 & 34.8 \\
\bottomrule
\end{tabular}
\caption{Sample Data Table}
\label{tab:sample}
\end{table}

\section{Lists}
\subsection{Bullet Points}
\begin{itemize}
  \item First item
  \item Second item
  \item Third item with sub-items:
    \begin{itemize}
      \item Sub-item 1
      \item Sub-item 2
    \end{itemize}
\end{itemize}

\subsection{Numbered List}
\begin{enumerate}
  \item First step
  \item Second step
  \item Third step
\end{enumerate}

\section{Citations and References}
When you use the citation search feature, references will be automatically managed. For example, you might cite a paper like this: \cite{example2024}. All citations will be automatically added to your \texttt{cite.bib} file.

\section{Conclusion}
This template demonstrates basic \LaTeX{} features including mathematical formulas, tables, and lists. The bibliography system is now set up to work seamlessly with Bibby AI's citation management. Feel free to modify it for your needs!

\bibliographystyle{plain}
\bibliography{cite}

\end{document}\documentclass{article}
\usepackage[utf8]{inputenc}
\usepackage{booktabs}
\usepackage{amsmath}
\usepackage{natbib}

\title{My Document}
\author{Your Name}
\date{\today}

\begin{document}

\maketitle

\section{Introduction}
This is a sample document demonstrating various \LaTeX{} features including mathematical equations, tables, and formatting. You can cite papers using the citation search feature, and they will be automatically added to your \texttt{cite.bib} file.

\section{Mathematical Expressions}
Here are some example equations:

\subsection{Quadratic Formula}
The quadratic formula is given by:
\[x = \frac{-b \pm \sqrt{b^2 - 4ac}}{2a}\]

\subsection{Matrix Example}
A sample matrix:
\[
\begin{pmatrix}
a_{11} & a_{12} & a_{13} \\
a_{21} & a_{22} & a_{23} \\
a_{31} & a_{32} & a_{33}
\end{pmatrix}
\]

\section{Tables}
Below is a sample table showing data:

\begin{table}[h]
\centering
\begin{tabular}{@{}lrrr@{}}
\toprule
\textbf{Category} & \textbf{Value 1} & \textbf{Value 2} & \textbf{Total} \\
\midrule
Item A & 10.5 & 20.3 & 30.8 \\
Item B & 15.2 & 25.1 & 40.3 \\
Item C & 22.1 & 12.7 & 34.8 \\
\bottomrule
\end{tabular}
\caption{Sample Data Table}
\label{tab:sample}
\end{table}

\section{Lists}
\subsection{Bullet Points}
\begin{itemize}
  \item First item
  \item Second item
  \item Third item with sub-items:
    \begin{itemize}
      \item Sub-item 1
      \item Sub-item 2
    \end{itemize}
\end{itemize}

\subsection{Numbered List}
\begin{enumerate}
  \item First step
  \item Second step
  \item Third step
\end{enumerate}

\section{Citations and References}
When you use the citation search feature, references will be automatically managed. For example, you might cite a paper like this: \cite{example2024}. All citations will be automatically added to your \texttt{cite.bib} file.

\section{Conclusion}
This template demonstrates basic \LaTeX{} features including mathematical formulas, tables, and lists. The bibliography system is now set up to work seamlessly with Bibby AI's citation management. Feel free to modify it for your needs!

\bibliographystyle{plain}
\bibliography{cite}

\end{document}\documentclass{article}
\usepackage[utf8]{inputenc}
\usepackage{booktabs}
\usepackage{amsmath}
\usepackage{natbib}

\title{My Document}
\author{Your Name}
\date{\today}

\begin{document}

\maketitle

\section{Introduction}
This is a sample document demonstrating various \LaTeX{} features including mathematical equations, tables, and formatting. You can cite papers using the citation search feature, and they will be automatically added to your \texttt{cite.bib} file.

\section{Mathematical Expressions}
Here are some example equations:

\subsection{Quadratic Formula}
The quadratic formula is given by:
\[x = \frac{-b \pm \sqrt{b^2 - 4ac}}{2a}\]

\subsection{Matrix Example}
A sample matrix:
\[
\begin{pmatrix}
a_{11} & a_{12} & a_{13} \\
a_{21} & a_{22} & a_{23} \\
a_{31} & a_{32} & a_{33}
\end{pmatrix}
\]

\section{Tables}
Below is a sample table showing data:

\begin{table}[h]
\centering
\begin{tabular}{@{}lrrr@{}}
\toprule
\textbf{Category} & \textbf{Value 1} & \textbf{Value 2} & \textbf{Total} \\
\midrule
Item A & 10.5 & 20.3 & 30.8 \\
Item B & 15.2 & 25.1 & 40.3 \\
Item C & 22.1 & 12.7 & 34.8 \\
\bottomrule
\end{tabular}
\caption{Sample Data Table}
\label{tab:sample}
\end{table}

\section{Lists}
\subsection{Bullet Points}
\begin{itemize}
  \item First item
  \item Second item
  \item Third item with sub-items:
    \begin{itemize}
      \item Sub-item 1
      \item Sub-item 2
    \end{itemize}
\end{itemize}

\subsection{Numbered List}
\begin{enumerate}
  \item First step
  \item Second step
  \item Third step
\end{enumerate}

\section{Citations and References}
When you use the citation search feature, references will be automatically managed. For example, you might cite a paper like this: \cite{example2024}. All citations will be automatically added to your \texttt{cite.bib} file.

\section{Conclusion}
This template demonstrates basic \LaTeX{} features including mathematical formulas, tables, and lists. The bibliography system is now set up to work seamlessly with Bibby AI's citation management. Feel free to modify it for your needs!

\bibliographystyle{plain}
\bibliography{cite}

\end{document}\documentclass{article}
\usepackage[utf8]{inputenc}
\usepackage{booktabs}
\usepackage{amsmath}
\usepackage{natbib}

\title{My Document}
\author{Your Name}
\date{\today}

\begin{document}

\maketitle

\section{Introduction}
This is a sample document demonstrating various \LaTeX{} features including mathematical equations, tables, and formatting. You can cite papers using the citation search feature, and they will be automatically added to your \texttt{cite.bib} file.

\section{Mathematical Expressions}
Here are some example equations:

\subsection{Quadratic Formula}
The quadratic formula is given by:
\[x = \frac{-b \pm \sqrt{b^2 - 4ac}}{2a}\]

\subsection{Matrix Example}
A sample matrix:
\[
\begin{pmatrix}
a_{11} & a_{12} & a_{13} \\
a_{21} & a_{22} & a_{23} \\
a_{31} & a_{32} & a_{33}
\end{pmatrix}
\]

\section{Tables}
Below is a sample table showing data:

\begin{table}[h]
\centering
\begin{tabular}{@{}lrrr@{}}
\toprule
\textbf{Category} & \textbf{Value 1} & \textbf{Value 2} & \textbf{Total} \\
\midrule
Item A & 10.5 & 20.3 & 30.8 \\
Item B & 15.2 & 25.1 & 40.3 \\
Item C & 22.1 & 12.7 & 34.8 \\
\bottomrule
\end{tabular}
\caption{Sample Data Table}
\label{tab:sample}
\end{table}

\section{Lists}
\subsection{Bullet Points}
\begin{itemize}
  \item First item
  \item Second item
  \item Third item with sub-items:
    \begin{itemize}
      \item Sub-item 1
      \item Sub-item 2
    \end{itemize}
\end{itemize}

\subsection{Numbered List}
\begin{enumerate}
  \item First step
  \item Second step
  \item Third step
\end{enumerate}

\section{Citations and References}
When you use the citation search feature, references will be automatically managed. For example, you might cite a paper like this: \cite{example2024}. All citations will be automatically added to your \texttt{cite.bib} file.

\section{Conclusion}
This template demonstrates basic \LaTeX{} features including mathematical formulas, tables, and lists. The bibliography system is now set up to work seamlessly with Bibby AI's citation management. Feel free to modify it for your needs!

\bibliographystyle{plain}
\bibliography{cite}

\end{document}\documentclass{article}
\usepackage[utf8]{inputenc}
\usepackage{booktabs}
\usepackage{amsmath}
\usepackage{natbib}

\title{My Document}
\author{Your Name}
\date{\today}

\begin{document}

\maketitle

\section{Introduction}
This is a sample document demonstrating various \LaTeX{} features including mathematical equations, tables, and formatting. You can cite papers using the citation search feature, and they will be automatically added to your \texttt{cite.bib} file.

\section{Mathematical Expressions}
Here are some example equations:

\subsection{Quadratic Formula}
The quadratic formula is given by:
\[x = \frac{-b \pm \sqrt{b^2 - 4ac}}{2a}\]

\subsection{Matrix Example}
A sample matrix:
\[
\begin{pmatrix}
a_{11} & a_{12} & a_{13} \\
a_{21} & a_{22} & a_{23} \\
a_{31} & a_{32} & a_{33}
\end{pmatrix}
\]

\section{Tables}
Below is a sample table showing data:

\begin{table}[h]
\centering
\begin{tabular}{@{}lrrr@{}}
\toprule
\textbf{Category} & \textbf{Value 1} & \textbf{Value 2} & \textbf{Total} \\
\midrule
Item A & 10.5 & 20.3 & 30.8 \\
Item B & 15.2 & 25.1 & 40.3 \\
Item C & 22.1 & 12.7 & 34.8 \\
\bottomrule
\end{tabular}
\caption{Sample Data Table}
\label{tab:sample}
\end{table}

\section{Lists}
\subsection{Bullet Points}
\begin{itemize}
  \item First item
  \item Second item
  \item Third item with sub-items:
    \begin{itemize}
      \item Sub-item 1
      \item Sub-item 2
    \end{itemize}
\end{itemize}

\subsection{Numbered List}
\begin{enumerate}
  \item First step
  \item Second step
  \item Third step
\end{enumerate}

\section{Citations and References}
When you use the citation search feature, references will be automatically managed. For example, you might cite a paper like this: \cite{example2024}. All citations will be automatically added to your \texttt{cite.bib} file.

\section{Conclusion}
This template demonstrates basic \LaTeX{} features including mathematical formulas, tables, and lists. The bibliography system is now set up to work seamlessly with Bibby AI's citation management. Feel free to modify it for your needs!

\bibliographystyle{plain}
\bibliography{cite}

\end{document}\documentclass{article}
\usepackage[utf8]{inputenc}
\usepackage{booktabs}
\usepackage{amsmath}
\usepackage{natbib}

\title{My Document}
\author{Your Name}
\date{\today}

\begin{document}

\maketitle

\section{Introduction}
This is a sample document demonstrating various \LaTeX{} features including mathematical equations, tables, and formatting. You can cite papers using the citation search feature, and they will be automatically added to your \texttt{cite.bib} file.

\section{Mathematical Expressions}
Here are some example equations:

\subsection{Quadratic Formula}
The quadratic formula is given by:
\[x = \frac{-b \pm \sqrt{b^2 - 4ac}}{2a}\]

\subsection{Matrix Example}
A sample matrix:
\[
\begin{pmatrix}
a_{11} & a_{12} & a_{13} \\
a_{21} & a_{22} & a_{23} \\
a_{31} & a_{32} & a_{33}
\end{pmatrix}
\]

\section{Tables}
Below is a sample table showing data:

\begin{table}[h]
\centering
\begin{tabular}{@{}lrrr@{}}
\toprule
\textbf{Category} & \textbf{Value 1} & \textbf{Value 2} & \textbf{Total} \\
\midrule
Item A & 10.5 & 20.3 & 30.8 \\
Item B & 15.2 & 25.1 & 40.3 \\
Item C & 22.1 & 12.7 & 34.8 \\
\bottomrule
\end{tabular}
\caption{Sample Data Table}
\label{tab:sample}
\end{table}

\section{Lists}
\subsection{Bullet Points}
\begin{itemize}
  \item First item
  \item Second item
  \item Third item with sub-items:
    \begin{itemize}
      \item Sub-item 1
      \item Sub-item 2
    \end{itemize}
\end{itemize}

\subsection{Numbered List}
\begin{enumerate}
  \item First step
  \item Second step
  \item Third step
\end{enumerate}

\section{Citations and References}
When you use the citation search feature, references will be automatically managed. For example, you might cite a paper like this: \cite{example2024}. All citations will be automatically added to your \texttt{cite.bib} file.

\section{Conclusion}
This template demonstrates basic \LaTeX{} features including mathematical formulas, tables, and lists. The bibliography system is now set up to work seamlessly with Bibby AI's citation management. Feel free to modify it for your needs!

\bibliographystyle{plain}
\bibliography{cite}

\end{document}\documentclass{article}
\usepackage[utf8]{inputenc}
\usepackage{booktabs}
\usepackage{amsmath}
\usepackage{natbib}

\title{My Document}
\author{Your Name}
\date{\today}

\begin{document}

\maketitle

\section{Introduction}
This is a sample document demonstrating various \LaTeX{} features including mathematical equations, tables, and formatting. You can cite papers using the citation search feature, and they will be automatically added to your \texttt{cite.bib} file.

\section{Mathematical Expressions}
Here are some example equations:

\subsection{Quadratic Formula}
The quadratic formula is given by:
\[x = \frac{-b \pm \sqrt{b^2 - 4ac}}{2a}\]

\subsection{Matrix Example}
A sample matrix:
\[
\begin{pmatrix}
a_{11} & a_{12} & a_{13} \\
a_{21} & a_{22} & a_{23} \\
a_{31} & a_{32} & a_{33}
\end{pmatrix}
\]

\section{Tables}
Below is a sample table showing data:

\begin{table}[h]
\centering
\begin{tabular}{@{}lrrr@{}}
\toprule
\textbf{Category} & \textbf{Value 1} & \textbf{Value 2} & \textbf{Total} \\
\midrule
Item A & 10.5 & 20.3 & 30.8 \\
Item B & 15.2 & 25.1 & 40.3 \\
Item C & 22.1 & 12.7 & 34.8 \\
\bottomrule
\end{tabular}
\caption{Sample Data Table}
\label{tab:sample}
\end{table}

\section{Lists}
\subsection{Bullet Points}
\begin{itemize}
  \item First item
  \item Second item
  \item Third item with sub-items:
    \begin{itemize}
      \item Sub-item 1
      \item Sub-item 2
    \end{itemize}
\end{itemize}

\subsection{Numbered List}
\begin{enumerate}
  \item First step
  \item Second step
  \item Third step
\end{enumerate}

\section{Citations and References}
When you use the citation search feature, references will be automatically managed. For example, you might cite a paper like this: \cite{example2024}. All citations will be automatically added to your \texttt{cite.bib} file.

\section{Conclusion}
This template demonstrates basic \LaTeX{} features including mathematical formulas, tables, and lists. The bibliography system is now set up to work seamlessly with Bibby AI's citation management. Feel free to modify it for your needs!

\bibliographystyle{plain}
\bibliography{cite}

\end{document}\documentclass{article}
\usepackage[utf8]{inputenc}
\usepackage{booktabs}
\usepackage{amsmath}
\usepackage{natbib}

\title{My Document}
\author{Your Name}
\date{\today}

\begin{document}

\maketitle

\section{Introduction}
This is a sample document demonstrating various \LaTeX{} features including mathematical equations, tables, and formatting. You can cite papers using the citation search feature, and they will be automatically added to your \texttt{cite.bib} file.

\section{Mathematical Expressions}
Here are some example equations:

\subsection{Quadratic Formula}
The quadratic formula is given by:
\[x = \frac{-b \pm \sqrt{b^2 - 4ac}}{2a}\]

\subsection{Matrix Example}
A sample matrix:
\[
\begin{pmatrix}
a_{11} & a_{12} & a_{13} \\
a_{21} & a_{22} & a_{23} \\
a_{31} & a_{32} & a_{33}
\end{pmatrix}
\]

\section{Tables}
Below is a sample table showing data:

\begin{table}[h]
\centering
\begin{tabular}{@{}lrrr@{}}
\toprule
\textbf{Category} & \textbf{Value 1} & \textbf{Value 2} & \textbf{Total} \\
\midrule
Item A & 10.5 & 20.3 & 30.8 \\
Item B & 15.2 & 25.1 & 40.3 \\
Item C & 22.1 & 12.7 & 34.8 \\
\bottomrule
\end{tabular}
\caption{Sample Data Table}
\label{tab:sample}
\end{table}

\section{Lists}
\subsection{Bullet Points}
\begin{itemize}
  \item First item
  \item Second item
  \item Third item with sub-items:
    \begin{itemize}
      \item Sub-item 1
      \item Sub-item 2
    \end{itemize}
\end{itemize}

\subsection{Numbered List}
\begin{enumerate}
  \item First step
  \item Second step
  \item Third step
\end{enumerate}

\section{Citations and References}
When you use the citation search feature, references will be automatically managed. For example, you might cite a paper like this: \cite{example2024}. All citations will be automatically added to your \texttt{cite.bib} file.

\section{Conclusion}
This template demonstrates basic \LaTeX{} features including mathematical formulas, tables, and lists. The bibliography system is now set up to work seamlessly with Bibby AI's citation management. Feel free to modify it for your needs!

\bibliographystyle{plain}
\bibliography{cite}

\end{document}\documentclass{article}
\usepackage[utf8]{inputenc}
\usepackage{booktabs}
\usepackage{amsmath}
\usepackage{natbib}

\title{My Document}
\author{Your Name}
\date{\today}

\begin{document}

\maketitle

\section{Introduction}
This is a sample document demonstrating various \LaTeX{} features including mathematical equations, tables, and formatting. You can cite papers using the citation search feature, and they will be automatically added to your \texttt{cite.bib} file.

\section{Mathematical Expressions}
Here are some example equations:

\subsection{Quadratic Formula}
The quadratic formula is given by:
\[x = \frac{-b \pm \sqrt{b^2 - 4ac}}{2a}\]

\subsection{Matrix Example}
A sample matrix:
\[
\begin{pmatrix}
a_{11} & a_{12} & a_{13} \\
a_{21} & a_{22} & a_{23} \\
a_{31} & a_{32} & a_{33}
\end{pmatrix}
\]

\section{Tables}
Below is a sample table showing data:

\begin{table}[h]
\centering
\begin{tabular}{@{}lrrr@{}}
\toprule
\textbf{Category} & \textbf{Value 1} & \textbf{Value 2} & \textbf{Total} \\
\midrule
Item A & 10.5 & 20.3 & 30.8 \\
Item B & 15.2 & 25.1 & 40.3 \\
Item C & 22.1 & 12.7 & 34.8 \\
\bottomrule
\end{tabular}
\caption{Sample Data Table}
\label{tab:sample}
\end{table}

\section{Lists}
\subsection{Bullet Points}
\begin{itemize}
  \item First item
  \item Second item
  \item Third item with sub-items:
    \begin{itemize}
      \item Sub-item 1
      \item Sub-item 2
    \end{itemize}
\end{itemize}

\subsection{Numbered List}
\begin{enumerate}
  \item First step
  \item Second step
  \item Third step
\end{enumerate}

\section{Citations and References}
When you use the citation search feature, references will be automatically managed. For example, you might cite a paper like this: \cite{example2024}. All citations will be automatically added to your \texttt{cite.bib} file.

\section{Conclusion}
This template demonstrates basic \LaTeX{} features including mathematical formulas, tables, and lists. The bibliography system is now set up to work seamlessly with Bibby AI's citation management. Feel free to modify it for your needs!

\bibliographystyle{plain}
\bibliography{cite}

\end{document}\documentclass{article}
\usepackage[utf8]{inputenc}
\usepackage{booktabs}
\usepackage{amsmath}
\usepackage{natbib}

\title{My Document}
\author{Your Name}
\date{\today}

\begin{document}

\maketitle

\section{Introduction}
This is a sample document demonstrating various \LaTeX{} features including mathematical equations, tables, and formatting. You can cite papers using the citation search feature, and they will be automatically added to your \texttt{cite.bib} file.

\section{Mathematical Expressions}
Here are some example equations:

\subsection{Quadratic Formula}
The quadratic formula is given by:
\[x = \frac{-b \pm \sqrt{b^2 - 4ac}}{2a}\]

\subsection{Matrix Example}
A sample matrix:
\[
\begin{pmatrix}
a_{11} & a_{12} & a_{13} \\
a_{21} & a_{22} & a_{23} \\
a_{31} & a_{32} & a_{33}
\end{pmatrix}
\]

\section{Tables}
Below is a sample table showing data:

\begin{table}[h]
\centering
\begin{tabular}{@{}lrrr@{}}
\toprule
\textbf{Category} & \textbf{Value 1} & \textbf{Value 2} & \textbf{Total} \\
\midrule
Item A & 10.5 & 20.3 & 30.8 \\
Item B & 15.2 & 25.1 & 40.3 \\
Item C & 22.1 & 12.7 & 34.8 \\
\bottomrule
\end{tabular}
\caption{Sample Data Table}
\label{tab:sample}
\end{table}

\section{Lists}
\subsection{Bullet Points}
\begin{itemize}
  \item First item
  \item Second item
  \item Third item with sub-items:
    \begin{itemize}
      \item Sub-item 1
      \item Sub-item 2
    \end{itemize}
\end{itemize}

\subsection{Numbered List}
\begin{enumerate}
  \item First step
  \item Second step
  \item Third step
\end{enumerate}

\section{Citations and References}
When you use the citation search feature, references will be automatically managed. For example, you might cite a paper like this: \cite{example2024}. All citations will be automatically added to your \texttt{cite.bib} file.

\section{Conclusion}
This template demonstrates basic \LaTeX{} features including mathematical formulas, tables, and lists. The bibliography system is now set up to work seamlessly with Bibby AI's citation management. Feel free to modify it for your needs!

\bibliographystyle{plain}
\bibliography{cite}

\end{document}\documentclass{article}
\usepackage[utf8]{inputenc}
\usepackage{booktabs}
\usepackage{amsmath}
\usepackage{natbib}

\title{My Document}
\author{Your Name}
\date{\today}

\begin{document}

\maketitle

\section{Introduction}
This is a sample document demonstrating various \LaTeX{} features including mathematical equations, tables, and formatting. You can cite papers using the citation search feature, and they will be automatically added to your \texttt{cite.bib} file.

\section{Mathematical Expressions}
Here are some example equations:

\subsection{Quadratic Formula}
The quadratic formula is given by:
\[x = \frac{-b \pm \sqrt{b^2 - 4ac}}{2a}\]

\subsection{Matrix Example}
A sample matrix:
\[
\begin{pmatrix}
a_{11} & a_{12} & a_{13} \\
a_{21} & a_{22} & a_{23} \\
a_{31} & a_{32} & a_{33}
\end{pmatrix}
\]

\section{Tables}
Below is a sample table showing data:

\begin{table}[h]
\centering
\begin{tabular}{@{}lrrr@{}}
\toprule
\textbf{Category} & \textbf{Value 1} & \textbf{Value 2} & \textbf{Total} \\
\midrule
Item A & 10.5 & 20.3 & 30.8 \\
Item B & 15.2 & 25.1 & 40.3 \\
Item C & 22.1 & 12.7 & 34.8 \\
\bottomrule
\end{tabular}
\caption{Sample Data Table}
\label{tab:sample}
\end{table}

\section{Lists}
\subsection{Bullet Points}
\begin{itemize}
  \item First item
  \item Second item
  \item Third item with sub-items:
    \begin{itemize}
      \item Sub-item 1
      \item Sub-item 2
    \end{itemize}
\end{itemize}

\subsection{Numbered List}
\begin{enumerate}
  \item First step
  \item Second step
  \item Third step
\end{enumerate}

\section{Citations and References}
When you use the citation search feature, references will be automatically managed. For example, you might cite a paper like this: \cite{example2024}. All citations will be automatically added to your \texttt{cite.bib} file.

\section{Conclusion}
This template demonstrates basic \LaTeX{} features including mathematical formulas, tables, and lists. The bibliography system is now set up to work seamlessly with Bibby AI's citation management. Feel free to modify it for your needs!

\bibliographystyle{plain}
\bibliography{cite}

\end{document}\documentclass{article}
\usepackage[utf8]{inputenc}
\usepackage{booktabs}
\usepackage{amsmath}
\usepackage{natbib}

\title{My Document}
\author{Your Name}
\date{\today}

\begin{document}

\maketitle

\section{Introduction}
This is a sample document demonstrating various \LaTeX{} features including mathematical equations, tables, and formatting. You can cite papers using the citation search feature, and they will be automatically added to your \texttt{cite.bib} file.

\section{Mathematical Expressions}
Here are some example equations:

\subsection{Quadratic Formula}
The quadratic formula is given by:
\[x = \frac{-b \pm \sqrt{b^2 - 4ac}}{2a}\]

\subsection{Matrix Example}
A sample matrix:
\[
\begin{pmatrix}
a_{11} & a_{12} & a_{13} \\
a_{21} & a_{22} & a_{23} \\
a_{31} & a_{32} & a_{33}
\end{pmatrix}
\]

\section{Tables}
Below is a sample table showing data:

\begin{table}[h]
\centering
\begin{tabular}{@{}lrrr@{}}
\toprule
\textbf{Category} & \textbf{Value 1} & \textbf{Value 2} & \textbf{Total} \\
\midrule
Item A & 10.5 & 20.3 & 30.8 \\
Item B & 15.2 & 25.1 & 40.3 \\
Item C & 22.1 & 12.7 & 34.8 \\
\bottomrule
\end{tabular}
\caption{Sample Data Table}
\label{tab:sample}
\end{table}

\section{Lists}
\subsection{Bullet Points}
\begin{itemize}
  \item First item
  \item Second item
  \item Third item with sub-items:
    \begin{itemize}
      \item Sub-item 1
      \item Sub-item 2
    \end{itemize}
\end{itemize}

\subsection{Numbered List}
\begin{enumerate}
  \item First step
  \item Second step
  \item Third step
\end{enumerate}

\section{Citations and References}
When you use the citation search feature, references will be automatically managed. For example, you might cite a paper like this: \cite{example2024}. All citations will be automatically added to your \texttt{cite.bib} file.

\section{Conclusion}
This template demonstrates basic \LaTeX{} features including mathematical formulas, tables, and lists. The bibliography system is now set up to work seamlessly with Bibby AI's citation management. Feel free to modify it for your needs!

\bibliographystyle{plain}
\bibliography{cite}

\end{document}\documentclass{article}
\usepackage[utf8]{inputenc}
\usepackage{booktabs}
\usepackage{amsmath}
\usepackage{natbib}

\title{My Document}
\author{Your Name}
\date{\today}

\begin{document}

\maketitle

\section{Introduction}
This is a sample document demonstrating various \LaTeX{} features including mathematical equations, tables, and formatting. You can cite papers using the citation search feature, and they will be automatically added to your \texttt{cite.bib} file.

\section{Mathematical Expressions}
Here are some example equations:

\subsection{Quadratic Formula}
The quadratic formula is given by:
\[x = \frac{-b \pm \sqrt{b^2 - 4ac}}{2a}\]

\subsection{Matrix Example}
A sample matrix:
\[
\begin{pmatrix}
a_{11} & a_{12} & a_{13} \\
a_{21} & a_{22} & a_{23} \\
a_{31} & a_{32} & a_{33}
\end{pmatrix}
\]

\section{Tables}
Below is a sample table showing data:

\begin{table}[h]
\centering
\begin{tabular}{@{}lrrr@{}}
\toprule
\textbf{Category} & \textbf{Value 1} & \textbf{Value 2} & \textbf{Total} \\
\midrule
Item A & 10.5 & 20.3 & 30.8 \\
Item B & 15.2 & 25.1 & 40.3 \\
Item C & 22.1 & 12.7 & 34.8 \\
\bottomrule
\end{tabular}
\caption{Sample Data Table}
\label{tab:sample}
\end{table}

\section{Lists}
\subsection{Bullet Points}
\begin{itemize}
  \item First item
  \item Second item
  \item Third item with sub-items:
    \begin{itemize}
      \item Sub-item 1
      \item Sub-item 2
    \end{itemize}
\end{itemize}

\subsection{Numbered List}
\begin{enumerate}
  \item First step
  \item Second step
  \item Third step
\end{enumerate}

\section{Citations and References}
When you use the citation search feature, references will be automatically managed. For example, you might cite a paper like this: \cite{example2024}. All citations will be automatically added to your \texttt{cite.bib} file.

\section{Conclusion}
This template demonstrates basic \LaTeX{} features including mathematical formulas, tables, and lists. The bibliography system is now set up to work seamlessly with Bibby AI's citation management. Feel free to modify it for your needs!

\bibliographystyle{plain}
\bibliography{cite}

\end{document}\documentclass{article}
\usepackage[utf8]{inputenc}
\usepackage{booktabs}
\usepackage{amsmath}
\usepackage{natbib}

\title{My Document}
\author{Your Name}
\date{\today}

\begin{document}

\maketitle

\section{Introduction}
This is a sample document demonstrating various \LaTeX{} features including mathematical equations, tables, and formatting. You can cite papers using the citation search feature, and they will be automatically added to your \texttt{cite.bib} file.

\section{Mathematical Expressions}
Here are some example equations:

\subsection{Quadratic Formula}
The quadratic formula is given by:
\[x = \frac{-b \pm \sqrt{b^2 - 4ac}}{2a}\]

\subsection{Matrix Example}
A sample matrix:
\[
\begin{pmatrix}
a_{11} & a_{12} & a_{13} \\
a_{21} & a_{22} & a_{23} \\
a_{31} & a_{32} & a_{33}
\end{pmatrix}
\]

\section{Tables}
Below is a sample table showing data:

\begin{table}[h]
\centering
\begin{tabular}{@{}lrrr@{}}
\toprule
\textbf{Category} & \textbf{Value 1} & \textbf{Value 2} & \textbf{Total} \\
\midrule
Item A & 10.5 & 20.3 & 30.8 \\
Item B & 15.2 & 25.1 & 40.3 \\
Item C & 22.1 & 12.7 & 34.8 \\
\bottomrule
\end{tabular}
\caption{Sample Data Table}
\label{tab:sample}
\end{table}

\section{Lists}
\subsection{Bullet Points}
\begin{itemize}
  \item First item
  \item Second item
  \item Third item with sub-items:
    \begin{itemize}
      \item Sub-item 1
      \item Sub-item 2
    \end{itemize}
\end{itemize}

\subsection{Numbered List}
\begin{enumerate}
  \item First step
  \item Second step
  \item Third step
\end{enumerate}

\section{Citations and References}
When you use the citation search feature, references will be automatically managed. For example, you might cite a paper like this: \cite{example2024}. All citations will be automatically added to your \texttt{cite.bib} file.

\section{Conclusion}
This template demonstrates basic \LaTeX{} features including mathematical formulas, tables, and lists. The bibliography system is now set up to work seamlessly with Bibby AI's citation management. Feel free to modify it for your needs!

\bibliographystyle{plain}
\bibliography{cite}

\end{document}\documentclass{article}
\usepackage[utf8]{inputenc}
\usepackage{booktabs}
\usepackage{amsmath}
\usepackage{natbib}

\title{My Document}
\author{Your Name}
\date{\today}

\begin{document}

\maketitle

\section{Introduction}
This is a sample document demonstrating various \LaTeX{} features including mathematical equations, tables, and formatting. You can cite papers using the citation search feature, and they will be automatically added to your \texttt{cite.bib} file.

\section{Mathematical Expressions}
Here are some example equations:

\subsection{Quadratic Formula}
The quadratic formula is given by:
\[x = \frac{-b \pm \sqrt{b^2 - 4ac}}{2a}\]

\subsection{Matrix Example}
A sample matrix:
\[
\begin{pmatrix}
a_{11} & a_{12} & a_{13} \\
a_{21} & a_{22} & a_{23} \\
a_{31} & a_{32} & a_{33}
\end{pmatrix}
\]

\section{Tables}
Below is a sample table showing data:

\begin{table}[h]
\centering
\begin{tabular}{@{}lrrr@{}}
\toprule
\textbf{Category} & \textbf{Value 1} & \textbf{Value 2} & \textbf{Total} \\
\midrule
Item A & 10.5 & 20.3 & 30.8 \\
Item B & 15.2 & 25.1 & 40.3 \\
Item C & 22.1 & 12.7 & 34.8 \\
\bottomrule
\end{tabular}
\caption{Sample Data Table}
\label{tab:sample}
\end{table}

\section{Lists}
\subsection{Bullet Points}
\begin{itemize}
  \item First item
  \item Second item
  \item Third item with sub-items:
    \begin{itemize}
      \item Sub-item 1
      \item Sub-item 2
    \end{itemize}
\end{itemize}

\subsection{Numbered List}
\begin{enumerate}
  \item First step
  \item Second step
  \item Third step
\end{enumerate}

\section{Citations and References}
When you use the citation search feature, references will be automatically managed. For example, you might cite a paper like this: \cite{example2024}. All citations will be automatically added to your \texttt{cite.bib} file.

\section{Conclusion}
This template demonstrates basic \LaTeX{} features including mathematical formulas, tables, and lists. The bibliography system is now set up to work seamlessly with Bibby AI's citation management. Feel free to modify it for your needs!

\bibliographystyle{plain}
\bibliography{cite}

\end{document}\documentclass{article}
\usepackage[utf8]{inputenc}
\usepackage{booktabs}
\usepackage{amsmath}
\usepackage{natbib}

\title{My Document}
\author{Your Name}
\date{\today}

\begin{document}

\maketitle

\section{Introduction}
This is a sample document demonstrating various \LaTeX{} features including mathematical equations, tables, and formatting. You can cite papers using the citation search feature, and they will be automatically added to your \texttt{cite.bib} file.

\section{Mathematical Expressions}
Here are some example equations:

\subsection{Quadratic Formula}
The quadratic formula is given by:
\[x = \frac{-b \pm \sqrt{b^2 - 4ac}}{2a}\]

\subsection{Matrix Example}
A sample matrix:
\[
\begin{pmatrix}
a_{11} & a_{12} & a_{13} \\
a_{21} & a_{22} & a_{23} \\
a_{31} & a_{32} & a_{33}
\end{pmatrix}
\]

\section{Tables}
Below is a sample table showing data:

\begin{table}[h]
\centering
\begin{tabular}{@{}lrrr@{}}
\toprule
\textbf{Category} & \textbf{Value 1} & \textbf{Value 2} & \textbf{Total} \\
\midrule
Item A & 10.5 & 20.3 & 30.8 \\
Item B & 15.2 & 25.1 & 40.3 \\
Item C & 22.1 & 12.7 & 34.8 \\
\bottomrule
\end{tabular}
\caption{Sample Data Table}
\label{tab:sample}
\end{table}

\section{Lists}
\subsection{Bullet Points}
\begin{itemize}
  \item First item
  \item Second item
  \item Third item with sub-items:
    \begin{itemize}
      \item Sub-item 1
      \item Sub-item 2
    \end{itemize}
\end{itemize}

\subsection{Numbered List}
\begin{enumerate}
  \item First step
  \item Second step
  \item Third step
\end{enumerate}

\section{Citations and References}
When you use the citation search feature, references will be automatically managed. For example, you might cite a paper like this: \cite{example2024}. All citations will be automatically added to your \texttt{cite.bib} file.

\section{Conclusion}
This template demonstrates basic \LaTeX{} features including mathematical formulas, tables, and lists. The bibliography system is now set up to work seamlessly with Bibby AI's citation management. Feel free to modify it for your needs!

\bibliographystyle{plain}
\bibliography{cite}

\end{document}\documentclass{article}
\usepackage[utf8]{inputenc}
\usepackage{booktabs}
\usepackage{amsmath}
\usepackage{natbib}

\title{My Document}
\author{Your Name}
\date{\today}

\begin{document}

\maketitle

\section{Introduction}
This is a sample document demonstrating various \LaTeX{} features including mathematical equations, tables, and formatting. You can cite papers using the citation search feature, and they will be automatically added to your \texttt{cite.bib} file.

\section{Mathematical Expressions}
Here are some example equations:

\subsection{Quadratic Formula}
The quadratic formula is given by:
\[x = \frac{-b \pm \sqrt{b^2 - 4ac}}{2a}\]

\subsection{Matrix Example}
A sample matrix:
\[
\begin{pmatrix}
a_{11} & a_{12} & a_{13} \\
a_{21} & a_{22} & a_{23} \\
a_{31} & a_{32} & a_{33}
\end{pmatrix}
\]

\section{Tables}
Below is a sample table showing data:

\begin{table}[h]
\centering
\begin{tabular}{@{}lrrr@{}}
\toprule
\textbf{Category} & \textbf{Value 1} & \textbf{Value 2} & \textbf{Total} \\
\midrule
Item A & 10.5 & 20.3 & 30.8 \\
Item B & 15.2 & 25.1 & 40.3 \\
Item C & 22.1 & 12.7 & 34.8 \\
\bottomrule
\end{tabular}
\caption{Sample Data Table}
\label{tab:sample}
\end{table}

\section{Lists}
\subsection{Bullet Points}
\begin{itemize}
  \item First item
  \item Second item
  \item Third item with sub-items:
    \begin{itemize}
      \item Sub-item 1
      \item Sub-item 2
    \end{itemize}
\end{itemize}

\subsection{Numbered List}
\begin{enumerate}
  \item First step
  \item Second step
  \item Third step
\end{enumerate}

\section{Citations and References}
When you use the citation search feature, references will be automatically managed. For example, you might cite a paper like this: \cite{example2024}. All citations will be automatically added to your \texttt{cite.bib} file.

\section{Conclusion}
This template demonstrates basic \LaTeX{} features including mathematical formulas, tables, and lists. The bibliography system is now set up to work seamlessly with Bibby AI's citation management. Feel free to modify it for your needs!

\bibliographystyle{plain}
\bibliography{cite}

\end{document}\documentclass{article}
\usepackage[utf8]{inputenc}
\usepackage{booktabs}
\usepackage{amsmath}
\usepackage{natbib}

\title{My Document}
\author{Your Name}
\date{\today}

\begin{document}

\maketitle

\section{Introduction}
This is a sample document demonstrating various \LaTeX{} features including mathematical equations, tables, and formatting. You can cite papers using the citation search feature, and they will be automatically added to your \texttt{cite.bib} file.

\section{Mathematical Expressions}
Here are some example equations:

\subsection{Quadratic Formula}
The quadratic formula is given by:
\[x = \frac{-b \pm \sqrt{b^2 - 4ac}}{2a}\]

\subsection{Matrix Example}
A sample matrix:
\[
\begin{pmatrix}
a_{11} & a_{12} & a_{13} \\
a_{21} & a_{22} & a_{23} \\
a_{31} & a_{32} & a_{33}
\end{pmatrix}
\]

\section{Tables}
Below is a sample table showing data:

\begin{table}[h]
\centering
\begin{tabular}{@{}lrrr@{}}
\toprule
\textbf{Category} & \textbf{Value 1} & \textbf{Value 2} & \textbf{Total} \\
\midrule
Item A & 10.5 & 20.3 & 30.8 \\
Item B & 15.2 & 25.1 & 40.3 \\
Item C & 22.1 & 12.7 & 34.8 \\
\bottomrule
\end{tabular}
\caption{Sample Data Table}
\label{tab:sample}
\end{table}

\section{Lists}
\subsection{Bullet Points}
\begin{itemize}
  \item First item
  \item Second item
  \item Third item with sub-items:
    \begin{itemize}
      \item Sub-item 1
      \item Sub-item 2
    \end{itemize}
\end{itemize}

\subsection{Numbered List}
\begin{enumerate}
  \item First step
  \item Second step
  \item Third step
\end{enumerate}

\section{Citations and References}
When you use the citation search feature, references will be automatically managed. For example, you might cite a paper like this: \cite{example2024}. All citations will be automatically added to your \texttt{cite.bib} file.

\section{Conclusion}
This template demonstrates basic \LaTeX{} features including mathematical formulas, tables, and lists. The bibliography system is now set up to work seamlessly with Bibby AI's citation management. Feel free to modify it for your needs!

\bibliographystyle{plain}
\bibliography{cite}

\end{document}\documentclass{article}
\usepackage[utf8]{inputenc}
\usepackage{booktabs}
\usepackage{amsmath}
\usepackage{natbib}

\title{My Document}
\author{Your Name}
\date{\today}

\begin{document}

\maketitle

\section{Introduction}
This is a sample document demonstrating various \LaTeX{} features including mathematical equations, tables, and formatting. You can cite papers using the citation search feature, and they will be automatically added to your \texttt{cite.bib} file.

\section{Mathematical Expressions}
Here are some example equations:

\subsection{Quadratic Formula}
The quadratic formula is given by:
\[x = \frac{-b \pm \sqrt{b^2 - 4ac}}{2a}\]

\subsection{Matrix Example}
A sample matrix:
\[
\begin{pmatrix}
a_{11} & a_{12} & a_{13} \\
a_{21} & a_{22} & a_{23} \\
a_{31} & a_{32} & a_{33}
\end{pmatrix}
\]

\section{Tables}
Below is a sample table showing data:

\begin{table}[h]
\centering
\begin{tabular}{@{}lrrr@{}}
\toprule
\textbf{Category} & \textbf{Value 1} & \textbf{Value 2} & \textbf{Total} \\
\midrule
Item A & 10.5 & 20.3 & 30.8 \\
Item B & 15.2 & 25.1 & 40.3 \\
Item C & 22.1 & 12.7 & 34.8 \\
\bottomrule
\end{tabular}
\caption{Sample Data Table}
\label{tab:sample}
\end{table}

\section{Lists}
\subsection{Bullet Points}
\begin{itemize}
  \item First item
  \item Second item
  \item Third item with sub-items:
    \begin{itemize}
      \item Sub-item 1
      \item Sub-item 2
    \end{itemize}
\end{itemize}

\subsection{Numbered List}
\begin{enumerate}
  \item First step
  \item Second step
  \item Third step
\end{enumerate}

\section{Citations and References}
When you use the citation search feature, references will be automatically managed. For example, you might cite a paper like this: \cite{example2024}. All citations will be automatically added to your \texttt{cite.bib} file.

\section{Conclusion}
This template demonstrates basic \LaTeX{} features including mathematical formulas, tables, and lists. The bibliography system is now set up to work seamlessly with Bibby AI's citation management. Feel free to modify it for your needs!

\bibliographystyle{plain}
\bibliography{cite}

\end{document}\documentclass{article}
\usepackage[utf8]{inputenc}
\usepackage{booktabs}
\usepackage{amsmath}
\usepackage{natbib}

\title{My Document}
\author{Your Name}
\date{\today}

\begin{document}

\maketitle

\section{Introduction}
This is a sample document demonstrating various \LaTeX{} features including mathematical equations, tables, and formatting. You can cite papers using the citation search feature, and they will be automatically added to your \texttt{cite.bib} file.

\section{Mathematical Expressions}
Here are some example equations:

\subsection{Quadratic Formula}
The quadratic formula is given by:
\[x = \frac{-b \pm \sqrt{b^2 - 4ac}}{2a}\]

\subsection{Matrix Example}
A sample matrix:
\[
\begin{pmatrix}
a_{11} & a_{12} & a_{13} \\
a_{21} & a_{22} & a_{23} \\
a_{31} & a_{32} & a_{33}
\end{pmatrix}
\]

\section{Tables}
Below is a sample table showing data:

\begin{table}[h]
\centering
\begin{tabular}{@{}lrrr@{}}
\toprule
\textbf{Category} & \textbf{Value 1} & \textbf{Value 2} & \textbf{Total} \\
\midrule
Item A & 10.5 & 20.3 & 30.8 \\
Item B & 15.2 & 25.1 & 40.3 \\
Item C & 22.1 & 12.7 & 34.8 \\
\bottomrule
\end{tabular}
\caption{Sample Data Table}
\label{tab:sample}
\end{table}

\section{Lists}
\subsection{Bullet Points}
\begin{itemize}
  \item First item
  \item Second item
  \item Third item with sub-items:
    \begin{itemize}
      \item Sub-item 1
      \item Sub-item 2
    \end{itemize}
\end{itemize}

\subsection{Numbered List}
\begin{enumerate}
  \item First step
  \item Second step
  \item Third step
\end{enumerate}

\section{Citations and References}
When you use the citation search feature, references will be automatically managed. For example, you might cite a paper like this: \cite{example2024}. All citations will be automatically added to your \texttt{cite.bib} file.

\section{Conclusion}
This template demonstrates basic \LaTeX{} features including mathematical formulas, tables, and lists. The bibliography system is now set up to work seamlessly with Bibby AI's citation management. Feel free to modify it for your needs!

\bibliographystyle{plain}
\bibliography{cite}

\end{document}\documentclass{article}
\usepackage[utf8]{inputenc}
\usepackage{booktabs}
\usepackage{amsmath}
\usepackage{natbib}

\title{My Document}
\author{Your Name}
\date{\today}

\begin{document}

\maketitle

\section{Introduction}
This is a sample document demonstrating various \LaTeX{} features including mathematical equations, tables, and formatting. You can cite papers using the citation search feature, and they will be automatically added to your \texttt{cite.bib} file.

\section{Mathematical Expressions}
Here are some example equations:

\subsection{Quadratic Formula}
The quadratic formula is given by:
\[x = \frac{-b \pm \sqrt{b^2 - 4ac}}{2a}\]

\subsection{Matrix Example}
A sample matrix:
\[
\begin{pmatrix}
a_{11} & a_{12} & a_{13} \\
a_{21} & a_{22} & a_{23} \\
a_{31} & a_{32} & a_{33}
\end{pmatrix}
\]

\section{Tables}
Below is a sample table showing data:

\begin{table}[h]
\centering
\begin{tabular}{@{}lrrr@{}}
\toprule
\textbf{Category} & \textbf{Value 1} & \textbf{Value 2} & \textbf{Total} \\
\midrule
Item A & 10.5 & 20.3 & 30.8 \\
Item B & 15.2 & 25.1 & 40.3 \\
Item C & 22.1 & 12.7 & 34.8 \\
\bottomrule
\end{tabular}
\caption{Sample Data Table}
\label{tab:sample}
\end{table}

\section{Lists}
\subsection{Bullet Points}
\begin{itemize}
  \item First item
  \item Second item
  \item Third item with sub-items:
    \begin{itemize}
      \item Sub-item 1
      \item Sub-item 2
    \end{itemize}
\end{itemize}

\subsection{Numbered List}
\begin{enumerate}
  \item First step
  \item Second step
  \item Third step
\end{enumerate}

\section{Citations and References}
When you use the citation search feature, references will be automatically managed. For example, you might cite a paper like this: \cite{example2024}. All citations will be automatically added to your \texttt{cite.bib} file.

\section{Conclusion}
This template demonstrates basic \LaTeX{} features including mathematical formulas, tables, and lists. The bibliography system is now set up to work seamlessly with Bibby AI's citation management. Feel free to modify it for your needs!

\bibliographystyle{plain}
\bibliography{cite}

\end{document}\documentclass{article}
\usepackage[utf8]{inputenc}
\usepackage{booktabs}
\usepackage{amsmath}
\usepackage{natbib}

\title{My Document}
\author{Your Name}
\date{\today}

\begin{document}

\maketitle

\section{Introduction}
This is a sample document demonstrating various \LaTeX{} features including mathematical equations, tables, and formatting. You can cite papers using the citation search feature, and they will be automatically added to your \texttt{cite.bib} file.

\section{Mathematical Expressions}
Here are some example equations:

\subsection{Quadratic Formula}
The quadratic formula is given by:
\[x = \frac{-b \pm \sqrt{b^2 - 4ac}}{2a}\]

\subsection{Matrix Example}
A sample matrix:
\[
\begin{pmatrix}
a_{11} & a_{12} & a_{13} \\
a_{21} & a_{22} & a_{23} \\
a_{31} & a_{32} & a_{33}
\end{pmatrix}
\]

\section{Tables}
Below is a sample table showing data:

\begin{table}[h]
\centering
\begin{tabular}{@{}lrrr@{}}
\toprule
\textbf{Category} & \textbf{Value 1} & \textbf{Value 2} & \textbf{Total} \\
\midrule
Item A & 10.5 & 20.3 & 30.8 \\
Item B & 15.2 & 25.1 & 40.3 \\
Item C & 22.1 & 12.7 & 34.8 \\
\bottomrule
\end{tabular}
\caption{Sample Data Table}
\label{tab:sample}
\end{table}

\section{Lists}
\subsection{Bullet Points}
\begin{itemize}
  \item First item
  \item Second item
  \item Third item with sub-items:
    \begin{itemize}
      \item Sub-item 1
      \item Sub-item 2
    \end{itemize}
\end{itemize}

\subsection{Numbered List}
\begin{enumerate}
  \item First step
  \item Second step
  \item Third step
\end{enumerate}

\section{Citations and References}
When you use the citation search feature, references will be automatically managed. For example, you might cite a paper like this: \cite{example2024}. All citations will be automatically added to your \texttt{cite.bib} file.

\section{Conclusion}
This template demonstrates basic \LaTeX{} features including mathematical formulas, tables, and lists. The bibliography system is now set up to work seamlessly with Bibby AI's citation management. Feel free to modify it for your needs!

\bibliographystyle{plain}
\bibliography{cite}

\end{document}\documentclass{article}
\usepackage[utf8]{inputenc}
\usepackage{booktabs}
\usepackage{amsmath}
\usepackage{natbib}

\title{My Document}
\author{Your Name}
\date{\today}

\begin{document}

\maketitle

\section{Introduction}
This is a sample document demonstrating various \LaTeX{} features including mathematical equations, tables, and formatting. You can cite papers using the citation search feature, and they will be automatically added to your \texttt{cite.bib} file.

\section{Mathematical Expressions}
Here are some example equations:

\subsection{Quadratic Formula}
The quadratic formula is given by:
\[x = \frac{-b \pm \sqrt{b^2 - 4ac}}{2a}\]

\subsection{Matrix Example}
A sample matrix:
\[
\begin{pmatrix}
a_{11} & a_{12} & a_{13} \\
a_{21} & a_{22} & a_{23} \\
a_{31} & a_{32} & a_{33}
\end{pmatrix}
\]

\section{Tables}
Below is a sample table showing data:

\begin{table}[h]
\centering
\begin{tabular}{@{}lrrr@{}}
\toprule
\textbf{Category} & \textbf{Value 1} & \textbf{Value 2} & \textbf{Total} \\
\midrule
Item A & 10.5 & 20.3 & 30.8 \\
Item B & 15.2 & 25.1 & 40.3 \\
Item C & 22.1 & 12.7 & 34.8 \\
\bottomrule
\end{tabular}
\caption{Sample Data Table}
\label{tab:sample}
\end{table}

\section{Lists}
\subsection{Bullet Points}
\begin{itemize}
  \item First item
  \item Second item
  \item Third item with sub-items:
    \begin{itemize}
      \item Sub-item 1
      \item Sub-item 2
    \end{itemize}
\end{itemize}

\subsection{Numbered List}
\begin{enumerate}
  \item First step
  \item Second step
  \item Third step
\end{enumerate}

\section{Citations and References}
When you use the citation search feature, references will be automatically managed. For example, you might cite a paper like this: \cite{example2024}. All citations will be automatically added to your \texttt{cite.bib} file.

\section{Conclusion}
This template demonstrates basic \LaTeX{} features including mathematical formulas, tables, and lists. The bibliography system is now set up to work seamlessly with Bibby AI's citation management. Feel free to modify it for your needs!

\bibliographystyle{plain}
\bibliography{cite}

\end{document}\documentclass{article}
\usepackage[utf8]{inputenc}
\usepackage{booktabs}
\usepackage{amsmath}
\usepackage{natbib}

\title{My Document}
\author{Your Name}
\date{\today}

\begin{document}

\maketitle

\section{Introduction}
This is a sample document demonstrating various \LaTeX{} features including mathematical equations, tables, and formatting. You can cite papers using the citation search feature, and they will be automatically added to your \texttt{cite.bib} file.

\section{Mathematical Expressions}
Here are some example equations:

\subsection{Quadratic Formula}
The quadratic formula is given by:
\[x = \frac{-b \pm \sqrt{b^2 - 4ac}}{2a}\]

\subsection{Matrix Example}
A sample matrix:
\[
\begin{pmatrix}
a_{11} & a_{12} & a_{13} \\
a_{21} & a_{22} & a_{23} \\
a_{31} & a_{32} & a_{33}
\end{pmatrix}
\]

\section{Tables}
Below is a sample table showing data:

\begin{table}[h]
\centering
\begin{tabular}{@{}lrrr@{}}
\toprule
\textbf{Category} & \textbf{Value 1} & \textbf{Value 2} & \textbf{Total} \\
\midrule
Item A & 10.5 & 20.3 & 30.8 \\
Item B & 15.2 & 25.1 & 40.3 \\
Item C & 22.1 & 12.7 & 34.8 \\
\bottomrule
\end{tabular}
\caption{Sample Data Table}
\label{tab:sample}
\end{table}

\section{Lists}
\subsection{Bullet Points}
\begin{itemize}
  \item First item
  \item Second item
  \item Third item with sub-items:
    \begin{itemize}
      \item Sub-item 1
      \item Sub-item 2
    \end{itemize}
\end{itemize}

\subsection{Numbered List}
\begin{enumerate}
  \item First step
  \item Second step
  \item Third step
\end{enumerate}

\section{Citations and References}
When you use the citation search feature, references will be automatically managed. For example, you might cite a paper like this: \cite{example2024}. All citations will be automatically added to your \texttt{cite.bib} file.

\section{Conclusion}
This template demonstrates basic \LaTeX{} features including mathematical formulas, tables, and lists. The bibliography system is now set up to work seamlessly with Bibby AI's citation management. Feel free to modify it for your needs!

\bibliographystyle{plain}
\bibliography{cite}

\end{document}\documentclass{article}
\usepackage[utf8]{inputenc}
\usepackage{booktabs}
\usepackage{amsmath}
\usepackage{natbib}

\title{My Document}
\author{Your Name}
\date{\today}

\begin{document}

\maketitle

\section{Introduction}
This is a sample document demonstrating various \LaTeX{} features including mathematical equations, tables, and formatting. You can cite papers using the citation search feature, and they will be automatically added to your \texttt{cite.bib} file.

\section{Mathematical Expressions}
Here are some example equations:

\subsection{Quadratic Formula}
The quadratic formula is given by:
\[x = \frac{-b \pm \sqrt{b^2 - 4ac}}{2a}\]

\subsection{Matrix Example}
A sample matrix:
\[
\begin{pmatrix}
a_{11} & a_{12} & a_{13} \\
a_{21} & a_{22} & a_{23} \\
a_{31} & a_{32} & a_{33}
\end{pmatrix}
\]

\section{Tables}
Below is a sample table showing data:

\begin{table}[h]
\centering
\begin{tabular}{@{}lrrr@{}}
\toprule
\textbf{Category} & \textbf{Value 1} & \textbf{Value 2} & \textbf{Total} \\
\midrule
Item A & 10.5 & 20.3 & 30.8 \\
Item B & 15.2 & 25.1 & 40.3 \\
Item C & 22.1 & 12.7 & 34.8 \\
\bottomrule
\end{tabular}
\caption{Sample Data Table}
\label{tab:sample}
\end{table}

\section{Lists}
\subsection{Bullet Points}
\begin{itemize}
  \item First item
  \item Second item
  \item Third item with sub-items:
    \begin{itemize}
      \item Sub-item 1
      \item Sub-item 2
    \end{itemize}
\end{itemize}

\subsection{Numbered List}
\begin{enumerate}
  \item First step
  \item Second step
  \item Third step
\end{enumerate}

\section{Citations and References}
When you use the citation search feature, references will be automatically managed. For example, you might cite a paper like this: \cite{example2024}. All citations will be automatically added to your \texttt{cite.bib} file.

\section{Conclusion}
This template demonstrates basic \LaTeX{} features including mathematical formulas, tables, and lists. The bibliography system is now set up to work seamlessly with Bibby AI's citation management. Feel free to modify it for your needs!

\bibliographystyle{plain}
\bibliography{cite}

\end{document}\documentclass{article}
\usepackage[utf8]{inputenc}
\usepackage{booktabs}
\usepackage{amsmath}
\usepackage{natbib}

\title{My Document}
\author{Your Name}
\date{\today}

\begin{document}

\maketitle

\section{Introduction}
This is a sample document demonstrating various \LaTeX{} features including mathematical equations, tables, and formatting. You can cite papers using the citation search feature, and they will be automatically added to your \texttt{cite.bib} file.

\section{Mathematical Expressions}
Here are some example equations:

\subsection{Quadratic Formula}
The quadratic formula is given by:
\[x = \frac{-b \pm \sqrt{b^2 - 4ac}}{2a}\]

\subsection{Matrix Example}
A sample matrix:
\[
\begin{pmatrix}
a_{11} & a_{12} & a_{13} \\
a_{21} & a_{22} & a_{23} \\
a_{31} & a_{32} & a_{33}
\end{pmatrix}
\]

\section{Tables}
Below is a sample table showing data:

\begin{table}[h]
\centering
\begin{tabular}{@{}lrrr@{}}
\toprule
\textbf{Category} & \textbf{Value 1} & \textbf{Value 2} & \textbf{Total} \\
\midrule
Item A & 10.5 & 20.3 & 30.8 \\
Item B & 15.2 & 25.1 & 40.3 \\
Item C & 22.1 & 12.7 & 34.8 \\
\bottomrule
\end{tabular}
\caption{Sample Data Table}
\label{tab:sample}
\end{table}

\section{Lists}
\subsection{Bullet Points}
\begin{itemize}
  \item First item
  \item Second item
  \item Third item with sub-items:
    \begin{itemize}
      \item Sub-item 1
      \item Sub-item 2
    \end{itemize}
\end{itemize}

\subsection{Numbered List}
\begin{enumerate}
  \item First step
  \item Second step
  \item Third step
\end{enumerate}

\section{Citations and References}
When you use the citation search feature, references will be automatically managed. For example, you might cite a paper like this: \cite{example2024}. All citations will be automatically added to your \texttt{cite.bib} file.

\section{Conclusion}
This template demonstrates basic \LaTeX{} features including mathematical formulas, tables, and lists. The bibliography system is now set up to work seamlessly with Bibby AI's citation management. Feel free to modify it for your needs!

\bibliographystyle{plain}
\bibliography{cite}

\end{document}\documentclass{article}
\usepackage[utf8]{inputenc}
\usepackage{booktabs}
\usepackage{amsmath}
\usepackage{natbib}

\title{My Document}
\author{Your Name}
\date{\today}

\begin{document}

\maketitle

\section{Introduction}
This is a sample document demonstrating various \LaTeX{} features including mathematical equations, tables, and formatting. You can cite papers using the citation search feature, and they will be automatically added to your \texttt{cite.bib} file.

\section{Mathematical Expressions}
Here are some example equations:

\subsection{Quadratic Formula}
The quadratic formula is given by:
\[x = \frac{-b \pm \sqrt{b^2 - 4ac}}{2a}\]

\subsection{Matrix Example}
A sample matrix:
\[
\begin{pmatrix}
a_{11} & a_{12} & a_{13} \\
a_{21} & a_{22} & a_{23} \\
a_{31} & a_{32} & a_{33}
\end{pmatrix}
\]

\section{Tables}
Below is a sample table showing data:

\begin{table}[h]
\centering
\begin{tabular}{@{}lrrr@{}}
\toprule
\textbf{Category} & \textbf{Value 1} & \textbf{Value 2} & \textbf{Total} \\
\midrule
Item A & 10.5 & 20.3 & 30.8 \\
Item B & 15.2 & 25.1 & 40.3 \\
Item C & 22.1 & 12.7 & 34.8 \\
\bottomrule
\end{tabular}
\caption{Sample Data Table}
\label{tab:sample}
\end{table}

\section{Lists}
\subsection{Bullet Points}
\begin{itemize}
  \item First item
  \item Second item
  \item Third item with sub-items:
    \begin{itemize}
      \item Sub-item 1
      \item Sub-item 2
    \end{itemize}
\end{itemize}

\subsection{Numbered List}
\begin{enumerate}
  \item First step
  \item Second step
  \item Third step
\end{enumerate}

\section{Citations and References}
When you use the citation search feature, references will be automatically managed. For example, you might cite a paper like this: \cite{example2024}. All citations will be automatically added to your \texttt{cite.bib} file.

\section{Conclusion}
This template demonstrates basic \LaTeX{} features including mathematical formulas, tables, and lists. The bibliography system is now set up to work seamlessly with Bibby AI's citation management. Feel free to modify it for your needs!

\bibliographystyle{plain}
\bibliography{cite}

\end{document}\documentclass{article}
\usepackage[utf8]{inputenc}
\usepackage{booktabs}
\usepackage{amsmath}
\usepackage{natbib}

\title{My Document}
\author{Your Name}
\date{\today}

\begin{document}

\maketitle

\section{Introduction}
This is a sample document demonstrating various \LaTeX{} features including mathematical equations, tables, and formatting. You can cite papers using the citation search feature, and they will be automatically added to your \texttt{cite.bib} file.

\section{Mathematical Expressions}
Here are some example equations:

\subsection{Quadratic Formula}
The quadratic formula is given by:
\[x = \frac{-b \pm \sqrt{b^2 - 4ac}}{2a}\]

\subsection{Matrix Example}
A sample matrix:
\[
\begin{pmatrix}
a_{11} & a_{12} & a_{13} \\
a_{21} & a_{22} & a_{23} \\
a_{31} & a_{32} & a_{33}
\end{pmatrix}
\]

\section{Tables}
Below is a sample table showing data:

\begin{table}[h]
\centering
\begin{tabular}{@{}lrrr@{}}
\toprule
\textbf{Category} & \textbf{Value 1} & \textbf{Value 2} & \textbf{Total} \\
\midrule
Item A & 10.5 & 20.3 & 30.8 \\
Item B & 15.2 & 25.1 & 40.3 \\
Item C & 22.1 & 12.7 & 34.8 \\
\bottomrule
\end{tabular}
\caption{Sample Data Table}
\label{tab:sample}
\end{table}

\section{Lists}
\subsection{Bullet Points}
\begin{itemize}
  \item First item
  \item Second item
  \item Third item with sub-items:
    \begin{itemize}
      \item Sub-item 1
      \item Sub-item 2
    \end{itemize}
\end{itemize}

\subsection{Numbered List}
\begin{enumerate}
  \item First step
  \item Second step
  \item Third step
\end{enumerate}

\section{Citations and References}
When you use the citation search feature, references will be automatically managed. For example, you might cite a paper like this: \cite{example2024}. All citations will be automatically added to your \texttt{cite.bib} file.

\section{Conclusion}
This template demonstrates basic \LaTeX{} features including mathematical formulas, tables, and lists. The bibliography system is now set up to work seamlessly with Bibby AI's citation management. Feel free to modify it for your needs!

\bibliographystyle{plain}
\bibliography{cite}

\end{document}\documentclass{article}
\usepackage[utf8]{inputenc}
\usepackage{booktabs}
\usepackage{amsmath}
\usepackage{natbib}

\title{My Document}
\author{Your Name}
\date{\today}

\begin{document}

\maketitle

\section{Introduction}
This is a sample document demonstrating various \LaTeX{} features including mathematical equations, tables, and formatting. You can cite papers using the citation search feature, and they will be automatically added to your \texttt{cite.bib} file.

\section{Mathematical Expressions}
Here are some example equations:

\subsection{Quadratic Formula}
The quadratic formula is given by:
\[x = \frac{-b \pm \sqrt{b^2 - 4ac}}{2a}\]

\subsection{Matrix Example}
A sample matrix:
\[
\begin{pmatrix}
a_{11} & a_{12} & a_{13} \\
a_{21} & a_{22} & a_{23} \\
a_{31} & a_{32} & a_{33}
\end{pmatrix}
\]

\section{Tables}
Below is a sample table showing data:

\begin{table}[h]
\centering
\begin{tabular}{@{}lrrr@{}}
\toprule
\textbf{Category} & \textbf{Value 1} & \textbf{Value 2} & \textbf{Total} \\
\midrule
Item A & 10.5 & 20.3 & 30.8 \\
Item B & 15.2 & 25.1 & 40.3 \\
Item C & 22.1 & 12.7 & 34.8 \\
\bottomrule
\end{tabular}
\caption{Sample Data Table}
\label{tab:sample}
\end{table}

\section{Lists}
\subsection{Bullet Points}
\begin{itemize}
  \item First item
  \item Second item
  \item Third item with sub-items:
    \begin{itemize}
      \item Sub-item 1
      \item Sub-item 2
    \end{itemize}
\end{itemize}

\subsection{Numbered List}
\begin{enumerate}
  \item First step
  \item Second step
  \item Third step
\end{enumerate}

\section{Citations and References}
When you use the citation search feature, references will be automatically managed. For example, you might cite a paper like this: \cite{example2024}. All citations will be automatically added to your \texttt{cite.bib} file.

\section{Conclusion}
This template demonstrates basic \LaTeX{} features including mathematical formulas, tables, and lists. The bibliography system is now set up to work seamlessly with Bibby AI's citation management. Feel free to modify it for your needs!

\bibliographystyle{plain}
\bibliography{cite}

\end{document}\documentclass{article}
\usepackage[utf8]{inputenc}
\usepackage{booktabs}
\usepackage{amsmath}
\usepackage{natbib}

\title{My Document}
\author{Your Name}
\date{\today}

\begin{document}

\maketitle

\section{Introduction}
This is a sample document demonstrating various \LaTeX{} features including mathematical equations, tables, and formatting. You can cite papers using the citation search feature, and they will be automatically added to your \texttt{cite.bib} file.

\section{Mathematical Expressions}
Here are some example equations:

\subsection{Quadratic Formula}
The quadratic formula is given by:
\[x = \frac{-b \pm \sqrt{b^2 - 4ac}}{2a}\]

\subsection{Matrix Example}
A sample matrix:
\[
\begin{pmatrix}
a_{11} & a_{12} & a_{13} \\
a_{21} & a_{22} & a_{23} \\
a_{31} & a_{32} & a_{33}
\end{pmatrix}
\]

\section{Tables}
Below is a sample table showing data:

\begin{table}[h]
\centering
\begin{tabular}{@{}lrrr@{}}
\toprule
\textbf{Category} & \textbf{Value 1} & \textbf{Value 2} & \textbf{Total} \\
\midrule
Item A & 10.5 & 20.3 & 30.8 \\
Item B & 15.2 & 25.1 & 40.3 \\
Item C & 22.1 & 12.7 & 34.8 \\
\bottomrule
\end{tabular}
\caption{Sample Data Table}
\label{tab:sample}
\end{table}

\section{Lists}
\subsection{Bullet Points}
\begin{itemize}
  \item First item
  \item Second item
  \item Third item with sub-items:
    \begin{itemize}
      \item Sub-item 1
      \item Sub-item 2
    \end{itemize}
\end{itemize}

\subsection{Numbered List}
\begin{enumerate}
  \item First step
  \item Second step
  \item Third step
\end{enumerate}

\section{Citations and References}
When you use the citation search feature, references will be automatically managed. For example, you might cite a paper like this: \cite{example2024}. All citations will be automatically added to your \texttt{cite.bib} file.

\section{Conclusion}
This template demonstrates basic \LaTeX{} features including mathematical formulas, tables, and lists. The bibliography system is now set up to work seamlessly with Bibby AI's citation management. Feel free to modify it for your needs!

\bibliographystyle{plain}
\bibliography{cite}

\end{document}\documentclass{article}
\usepackage[utf8]{inputenc}
\usepackage{booktabs}
\usepackage{amsmath}
\usepackage{natbib}

\title{My Document}
\author{Your Name}
\date{\today}

\begin{document}

\maketitle

\section{Introduction}
This is a sample document demonstrating various \LaTeX{} features including mathematical equations, tables, and formatting. You can cite papers using the citation search feature, and they will be automatically added to your \texttt{cite.bib} file.

\section{Mathematical Expressions}
Here are some example equations:

\subsection{Quadratic Formula}
The quadratic formula is given by:
\[x = \frac{-b \pm \sqrt{b^2 - 4ac}}{2a}\]

\subsection{Matrix Example}
A sample matrix:
\[
\begin{pmatrix}
a_{11} & a_{12} & a_{13} \\
a_{21} & a_{22} & a_{23} \\
a_{31} & a_{32} & a_{33}
\end{pmatrix}
\]

\section{Tables}
Below is a sample table showing data:

\begin{table}[h]
\centering
\begin{tabular}{@{}lrrr@{}}
\toprule
\textbf{Category} & \textbf{Value 1} & \textbf{Value 2} & \textbf{Total} \\
\midrule
Item A & 10.5 & 20.3 & 30.8 \\
Item B & 15.2 & 25.1 & 40.3 \\
Item C & 22.1 & 12.7 & 34.8 \\
\bottomrule
\end{tabular}
\caption{Sample Data Table}
\label{tab:sample}
\end{table}

\section{Lists}
\subsection{Bullet Points}
\begin{itemize}
  \item First item
  \item Second item
  \item Third item with sub-items:
    \begin{itemize}
      \item Sub-item 1
      \item Sub-item 2
    \end{itemize}
\end{itemize}

\subsection{Numbered List}
\begin{enumerate}
  \item First step
  \item Second step
  \item Third step
\end{enumerate}

\section{Citations and References}
When you use the citation search feature, references will be automatically managed. For example, you might cite a paper like this: \cite{example2024}. All citations will be automatically added to your \texttt{cite.bib} file.

\section{Conclusion}
This template demonstrates basic \LaTeX{} features including mathematical formulas, tables, and lists. The bibliography system is now set up to work seamlessly with Bibby AI's citation management. Feel free to modify it for your needs!

\bibliographystyle{plain}
\bibliography{cite}

\end{document}\documentclass{article}
\usepackage[utf8]{inputenc}
\usepackage{booktabs}
\usepackage{amsmath}
\usepackage{natbib}

\title{My Document}
\author{Your Name}
\date{\today}

\begin{document}

\maketitle

\section{Introduction}
This is a sample document demonstrating various \LaTeX{} features including mathematical equations, tables, and formatting. You can cite papers using the citation search feature, and they will be automatically added to your \texttt{cite.bib} file.

\section{Mathematical Expressions}
Here are some example equations:

\subsection{Quadratic Formula}
The quadratic formula is given by:
\[x = \frac{-b \pm \sqrt{b^2 - 4ac}}{2a}\]

\subsection{Matrix Example}
A sample matrix:
\[
\begin{pmatrix}
a_{11} & a_{12} & a_{13} \\
a_{21} & a_{22} & a_{23} \\
a_{31} & a_{32} & a_{33}
\end{pmatrix}
\]

\section{Tables}
Below is a sample table showing data:

\begin{table}[h]
\centering
\begin{tabular}{@{}lrrr@{}}
\toprule
\textbf{Category} & \textbf{Value 1} & \textbf{Value 2} & \textbf{Total} \\
\midrule
Item A & 10.5 & 20.3 & 30.8 \\
Item B & 15.2 & 25.1 & 40.3 \\
Item C & 22.1 & 12.7 & 34.8 \\
\bottomrule
\end{tabular}
\caption{Sample Data Table}
\label{tab:sample}
\end{table}

\section{Lists}
\subsection{Bullet Points}
\begin{itemize}
  \item First item
  \item Second item
  \item Third item with sub-items:
    \begin{itemize}
      \item Sub-item 1
      \item Sub-item 2
    \end{itemize}
\end{itemize}

\subsection{Numbered List}
\begin{enumerate}
  \item First step
  \item Second step
  \item Third step
\end{enumerate}

\section{Citations and References}
When you use the citation search feature, references will be automatically managed. For example, you might cite a paper like this: \cite{example2024}. All citations will be automatically added to your \texttt{cite.bib} file.

\section{Conclusion}
This template demonstrates basic \LaTeX{} features including mathematical formulas, tables, and lists. The bibliography system is now set up to work seamlessly with Bibby AI's citation management. Feel free to modify it for your needs!

\bibliographystyle{plain}
\bibliography{cite}

\end{document}\documentclass{article}
\usepackage[utf8]{inputenc}
\usepackage{booktabs}
\usepackage{amsmath}
\usepackage{natbib}

\title{My Document}
\author{Your Name}
\date{\today}

\begin{document}

\maketitle

\section{Introduction}
This is a sample document demonstrating various \LaTeX{} features including mathematical equations, tables, and formatting. You can cite papers using the citation search feature, and they will be automatically added to your \texttt{cite.bib} file.

\section{Mathematical Expressions}
Here are some example equations:

\subsection{Quadratic Formula}
The quadratic formula is given by:
\[x = \frac{-b \pm \sqrt{b^2 - 4ac}}{2a}\]

\subsection{Matrix Example}
A sample matrix:
\[
\begin{pmatrix}
a_{11} & a_{12} & a_{13} \\
a_{21} & a_{22} & a_{23} \\
a_{31} & a_{32} & a_{33}
\end{pmatrix}
\]

\section{Tables}
Below is a sample table showing data:

\begin{table}[h]
\centering
\begin{tabular}{@{}lrrr@{}}
\toprule
\textbf{Category} & \textbf{Value 1} & \textbf{Value 2} & \textbf{Total} \\
\midrule
Item A & 10.5 & 20.3 & 30.8 \\
Item B & 15.2 & 25.1 & 40.3 \\
Item C & 22.1 & 12.7 & 34.8 \\
\bottomrule
\end{tabular}
\caption{Sample Data Table}
\label{tab:sample}
\end{table}

\section{Lists}
\subsection{Bullet Points}
\begin{itemize}
  \item First item
  \item Second item
  \item Third item with sub-items:
    \begin{itemize}
      \item Sub-item 1
      \item Sub-item 2
    \end{itemize}
\end{itemize}

\subsection{Numbered List}
\begin{enumerate}
  \item First step
  \item Second step
  \item Third step
\end{enumerate}

\section{Citations and References}
When you use the citation search feature, references will be automatically managed. For example, you might cite a paper like this: \cite{example2024}. All citations will be automatically added to your \texttt{cite.bib} file.

\section{Conclusion}
This template demonstrates basic \LaTeX{} features including mathematical formulas, tables, and lists. The bibliography system is now set up to work seamlessly with Bibby AI's citation management. Feel free to modify it for your needs!

\bibliographystyle{plain}
\bibliography{cite}

\end{document}\documentclass{article}
\usepackage[utf8]{inputenc}
\usepackage{booktabs}
\usepackage{amsmath}
\usepackage{natbib}

\title{My Document}
\author{Your Name}
\date{\today}

\begin{document}

\maketitle

\section{Introduction}
This is a sample document demonstrating various \LaTeX{} features including mathematical equations, tables, and formatting. You can cite papers using the citation search feature, and they will be automatically added to your \texttt{cite.bib} file.

\section{Mathematical Expressions}
Here are some example equations:

\subsection{Quadratic Formula}
The quadratic formula is given by:
\[x = \frac{-b \pm \sqrt{b^2 - 4ac}}{2a}\]

\subsection{Matrix Example}
A sample matrix:
\[
\begin{pmatrix}
a_{11} & a_{12} & a_{13} \\
a_{21} & a_{22} & a_{23} \\
a_{31} & a_{32} & a_{33}
\end{pmatrix}
\]

\section{Tables}
Below is a sample table showing data:

\begin{table}[h]
\centering
\begin{tabular}{@{}lrrr@{}}
\toprule
\textbf{Category} & \textbf{Value 1} & \textbf{Value 2} & \textbf{Total} \\
\midrule
Item A & 10.5 & 20.3 & 30.8 \\
Item B & 15.2 & 25.1 & 40.3 \\
Item C & 22.1 & 12.7 & 34.8 \\
\bottomrule
\end{tabular}
\caption{Sample Data Table}
\label{tab:sample}
\end{table}

\section{Lists}
\subsection{Bullet Points}
\begin{itemize}
  \item First item
  \item Second item
  \item Third item with sub-items:
    \begin{itemize}
      \item Sub-item 1
      \item Sub-item 2
    \end{itemize}
\end{itemize}

\subsection{Numbered List}
\begin{enumerate}
  \item First step
  \item Second step
  \item Third step
\end{enumerate}

\section{Citations and References}
When you use the citation search feature, references will be automatically managed. For example, you might cite a paper like this: \cite{example2024}. All citations will be automatically added to your \texttt{cite.bib} file.

\section{Conclusion}
This template demonstrates basic \LaTeX{} features including mathematical formulas, tables, and lists. The bibliography system is now set up to work seamlessly with Bibby AI's citation management. Feel free to modify it for your needs!

\bibliographystyle{plain}
\bibliography{cite}

\end{document}\documentclass{article}
\usepackage[utf8]{inputenc}
\usepackage{booktabs}
\usepackage{amsmath}
\usepackage{natbib}

\title{My Document}
\author{Your Name}
\date{\today}

\begin{document}

\maketitle

\section{Introduction}
This is a sample document demonstrating various \LaTeX{} features including mathematical equations, tables, and formatting. You can cite papers using the citation search feature, and they will be automatically added to your \texttt{cite.bib} file.

\section{Mathematical Expressions}
Here are some example equations:

\subsection{Quadratic Formula}
The quadratic formula is given by:
\[x = \frac{-b \pm \sqrt{b^2 - 4ac}}{2a}\]

\subsection{Matrix Example}
A sample matrix:
\[
\begin{pmatrix}
a_{11} & a_{12} & a_{13} \\
a_{21} & a_{22} & a_{23} \\
a_{31} & a_{32} & a_{33}
\end{pmatrix}
\]

\section{Tables}
Below is a sample table showing data:

\begin{table}[h]
\centering
\begin{tabular}{@{}lrrr@{}}
\toprule
\textbf{Category} & \textbf{Value 1} & \textbf{Value 2} & \textbf{Total} \\
\midrule
Item A & 10.5 & 20.3 & 30.8 \\
Item B & 15.2 & 25.1 & 40.3 \\
Item C & 22.1 & 12.7 & 34.8 \\
\bottomrule
\end{tabular}
\caption{Sample Data Table}
\label{tab:sample}
\end{table}

\section{Lists}
\subsection{Bullet Points}
\begin{itemize}
  \item First item
  \item Second item
  \item Third item with sub-items:
    \begin{itemize}
      \item Sub-item 1
      \item Sub-item 2
    \end{itemize}
\end{itemize}

\subsection{Numbered List}
\begin{enumerate}
  \item First step
  \item Second step
  \item Third step
\end{enumerate}

\section{Citations and References}
When you use the citation search feature, references will be automatically managed. For example, you might cite a paper like this: \cite{example2024}. All citations will be automatically added to your \texttt{cite.bib} file.

\section{Conclusion}
This template demonstrates basic \LaTeX{} features including mathematical formulas, tables, and lists. The bibliography system is now set up to work seamlessly with Bibby AI's citation management. Feel free to modify it for your needs!

\bibliographystyle{plain}
\bibliography{cite}

\end{document}\documentclass{article}
\usepackage[utf8]{inputenc}
\usepackage{booktabs}
\usepackage{amsmath}
\usepackage{natbib}

\title{My Document}
\author{Your Name}
\date{\today}

\begin{document}

\maketitle

\section{Introduction}
This is a sample document demonstrating various \LaTeX{} features including mathematical equations, tables, and formatting. You can cite papers using the citation search feature, and they will be automatically added to your \texttt{cite.bib} file.

\section{Mathematical Expressions}
Here are some example equations:

\subsection{Quadratic Formula}
The quadratic formula is given by:
\[x = \frac{-b \pm \sqrt{b^2 - 4ac}}{2a}\]

\subsection{Matrix Example}
A sample matrix:
\[
\begin{pmatrix}
a_{11} & a_{12} & a_{13} \\
a_{21} & a_{22} & a_{23} \\
a_{31} & a_{32} & a_{33}
\end{pmatrix}
\]

\section{Tables}
Below is a sample table showing data:

\begin{table}[h]
\centering
\begin{tabular}{@{}lrrr@{}}
\toprule
\textbf{Category} & \textbf{Value 1} & \textbf{Value 2} & \textbf{Total} \\
\midrule
Item A & 10.5 & 20.3 & 30.8 \\
Item B & 15.2 & 25.1 & 40.3 \\
Item C & 22.1 & 12.7 & 34.8 \\
\bottomrule
\end{tabular}
\caption{Sample Data Table}
\label{tab:sample}
\end{table}

\section{Lists}
\subsection{Bullet Points}
\begin{itemize}
  \item First item
  \item Second item
  \item Third item with sub-items:
    \begin{itemize}
      \item Sub-item 1
      \item Sub-item 2
    \end{itemize}
\end{itemize}

\subsection{Numbered List}
\begin{enumerate}
  \item First step
  \item Second step
  \item Third step
\end{enumerate}

\section{Citations and References}
When you use the citation search feature, references will be automatically managed. For example, you might cite a paper like this: \cite{example2024}. All citations will be automatically added to your \texttt{cite.bib} file.

\section{Conclusion}
This template demonstrates basic \LaTeX{} features including mathematical formulas, tables, and lists. The bibliography system is now set up to work seamlessly with Bibby AI's citation management. Feel free to modify it for your needs!

\bibliographystyle{plain}
\bibliography{cite}

\end{document}\documentclass{article}
\usepackage[utf8]{inputenc}
\usepackage{booktabs}
\usepackage{amsmath}
\usepackage{natbib}

\title{My Document}
\author{Your Name}
\date{\today}

\begin{document}

\maketitle

\section{Introduction}
This is a sample document demonstrating various \LaTeX{} features including mathematical equations, tables, and formatting. You can cite papers using the citation search feature, and they will be automatically added to your \texttt{cite.bib} file.

\section{Mathematical Expressions}
Here are some example equations:

\subsection{Quadratic Formula}
The quadratic formula is given by:
\[x = \frac{-b \pm \sqrt{b^2 - 4ac}}{2a}\]

\subsection{Matrix Example}
A sample matrix:
\[
\begin{pmatrix}
a_{11} & a_{12} & a_{13} \\
a_{21} & a_{22} & a_{23} \\
a_{31} & a_{32} & a_{33}
\end{pmatrix}
\]

\section{Tables}
Below is a sample table showing data:

\begin{table}[h]
\centering
\begin{tabular}{@{}lrrr@{}}
\toprule
\textbf{Category} & \textbf{Value 1} & \textbf{Value 2} & \textbf{Total} \\
\midrule
Item A & 10.5 & 20.3 & 30.8 \\
Item B & 15.2 & 25.1 & 40.3 \\
Item C & 22.1 & 12.7 & 34.8 \\
\bottomrule
\end{tabular}
\caption{Sample Data Table}
\label{tab:sample}
\end{table}

\section{Lists}
\subsection{Bullet Points}
\begin{itemize}
  \item First item
  \item Second item
  \item Third item with sub-items:
    \begin{itemize}
      \item Sub-item 1
      \item Sub-item 2
    \end{itemize}
\end{itemize}

\subsection{Numbered List}
\begin{enumerate}
  \item First step
  \item Second step
  \item Third step
\end{enumerate}

\section{Citations and References}
When you use the citation search feature, references will be automatically managed. For example, you might cite a paper like this: \cite{example2024}. All citations will be automatically added to your \texttt{cite.bib} file.

\section{Conclusion}
This template demonstrates basic \LaTeX{} features including mathematical formulas, tables, and lists. The bibliography system is now set up to work seamlessly with Bibby AI's citation management. Feel free to modify it for your needs!

\bibliographystyle{plain}
\bibliography{cite}

\end{document}\documentclass{article}
\usepackage[utf8]{inputenc}
\usepackage{booktabs}
\usepackage{amsmath}
\usepackage{natbib}

\title{My Document}
\author{Your Name}
\date{\today}

\begin{document}

\maketitle

\section{Introduction}
This is a sample document demonstrating various \LaTeX{} features including mathematical equations, tables, and formatting. You can cite papers using the citation search feature, and they will be automatically added to your \texttt{cite.bib} file.

\section{Mathematical Expressions}
Here are some example equations:

\subsection{Quadratic Formula}
The quadratic formula is given by:
\[x = \frac{-b \pm \sqrt{b^2 - 4ac}}{2a}\]

\subsection{Matrix Example}
A sample matrix:
\[
\begin{pmatrix}
a_{11} & a_{12} & a_{13} \\
a_{21} & a_{22} & a_{23} \\
a_{31} & a_{32} & a_{33}
\end{pmatrix}
\]

\section{Tables}
Below is a sample table showing data:

\begin{table}[h]
\centering
\begin{tabular}{@{}lrrr@{}}
\toprule
\textbf{Category} & \textbf{Value 1} & \textbf{Value 2} & \textbf{Total} \\
\midrule
Item A & 10.5 & 20.3 & 30.8 \\
Item B & 15.2 & 25.1 & 40.3 \\
Item C & 22.1 & 12.7 & 34.8 \\
\bottomrule
\end{tabular}
\caption{Sample Data Table}
\label{tab:sample}
\end{table}

\section{Lists}
\subsection{Bullet Points}
\begin{itemize}
  \item First item
  \item Second item
  \item Third item with sub-items:
    \begin{itemize}
      \item Sub-item 1
      \item Sub-item 2
    \end{itemize}
\end{itemize}

\subsection{Numbered List}
\begin{enumerate}
  \item First step
  \item Second step
  \item Third step
\end{enumerate}

\section{Citations and References}
When you use the citation search feature, references will be automatically managed. For example, you might cite a paper like this: \cite{example2024}. All citations will be automatically added to your \texttt{cite.bib} file.

\section{Conclusion}
This template demonstrates basic \LaTeX{} features including mathematical formulas, tables, and lists. The bibliography system is now set up to work seamlessly with Bibby AI's citation management. Feel free to modify it for your needs!

\bibliographystyle{plain}
\bibliography{cite}

\end{document}\documentclass{article}
\usepackage[utf8]{inputenc}
\usepackage{booktabs}
\usepackage{amsmath}
\usepackage{natbib}

\title{My Document}
\author{Your Name}
\date{\today}

\begin{document}

\maketitle

\section{Introduction}
This is a sample document demonstrating various \LaTeX{} features including mathematical equations, tables, and formatting. You can cite papers using the citation search feature, and they will be automatically added to your \texttt{cite.bib} file.

\section{Mathematical Expressions}
Here are some example equations:

\subsection{Quadratic Formula}
The quadratic formula is given by:
\[x = \frac{-b \pm \sqrt{b^2 - 4ac}}{2a}\]

\subsection{Matrix Example}
A sample matrix:
\[
\begin{pmatrix}
a_{11} & a_{12} & a_{13} \\
a_{21} & a_{22} & a_{23} \\
a_{31} & a_{32} & a_{33}
\end{pmatrix}
\]

\section{Tables}
Below is a sample table showing data:

\begin{table}[h]
\centering
\begin{tabular}{@{}lrrr@{}}
\toprule
\textbf{Category} & \textbf{Value 1} & \textbf{Value 2} & \textbf{Total} \\
\midrule
Item A & 10.5 & 20.3 & 30.8 \\
Item B & 15.2 & 25.1 & 40.3 \\
Item C & 22.1 & 12.7 & 34.8 \\
\bottomrule
\end{tabular}
\caption{Sample Data Table}
\label{tab:sample}
\end{table}

\section{Lists}
\subsection{Bullet Points}
\begin{itemize}
  \item First item
  \item Second item
  \item Third item with sub-items:
    \begin{itemize}
      \item Sub-item 1
      \item Sub-item 2
    \end{itemize}
\end{itemize}

\subsection{Numbered List}
\begin{enumerate}
  \item First step
  \item Second step
  \item Third step
\end{enumerate}

\section{Citations and References}
When you use the citation search feature, references will be automatically managed. For example, you might cite a paper like this: \cite{example2024}. All citations will be automatically added to your \texttt{cite.bib} file.

\section{Conclusion}
This template demonstrates basic \LaTeX{} features including mathematical formulas, tables, and lists. The bibliography system is now set up to work seamlessly with Bibby AI's citation management. Feel free to modify it for your needs!

\bibliographystyle{plain}
\bibliography{cite}

\end{document}\documentclass{article}
\usepackage[utf8]{inputenc}
\usepackage{booktabs}
\usepackage{amsmath}
\usepackage{natbib}

\title{My Document}
\author{Your Name}
\date{\today}

\begin{document}

\maketitle

\section{Introduction}
This is a sample document demonstrating various \LaTeX{} features including mathematical equations, tables, and formatting. You can cite papers using the citation search feature, and they will be automatically added to your \texttt{cite.bib} file.

\section{Mathematical Expressions}
Here are some example equations:

\subsection{Quadratic Formula}
The quadratic formula is given by:
\[x = \frac{-b \pm \sqrt{b^2 - 4ac}}{2a}\]

\subsection{Matrix Example}
A sample matrix:
\[
\begin{pmatrix}
a_{11} & a_{12} & a_{13} \\
a_{21} & a_{22} & a_{23} \\
a_{31} & a_{32} & a_{33}
\end{pmatrix}
\]

\section{Tables}
Below is a sample table showing data:

\begin{table}[h]
\centering
\begin{tabular}{@{}lrrr@{}}
\toprule
\textbf{Category} & \textbf{Value 1} & \textbf{Value 2} & \textbf{Total} \\
\midrule
Item A & 10.5 & 20.3 & 30.8 \\
Item B & 15.2 & 25.1 & 40.3 \\
Item C & 22.1 & 12.7 & 34.8 \\
\bottomrule
\end{tabular}
\caption{Sample Data Table}
\label{tab:sample}
\end{table}

\section{Lists}
\subsection{Bullet Points}
\begin{itemize}
  \item First item
  \item Second item
  \item Third item with sub-items:
    \begin{itemize}
      \item Sub-item 1
      \item Sub-item 2
    \end{itemize}
\end{itemize}

\subsection{Numbered List}
\begin{enumerate}
  \item First step
  \item Second step
  \item Third step
\end{enumerate}

\section{Citations and References}
When you use the citation search feature, references will be automatically managed. For example, you might cite a paper like this: \cite{example2024}. All citations will be automatically added to your \texttt{cite.bib} file.

\section{Conclusion}
This template demonstrates basic \LaTeX{} features including mathematical formulas, tables, and lists. The bibliography system is now set up to work seamlessly with Bibby AI's citation management. Feel free to modify it for your needs!

\bibliographystyle{plain}
\bibliography{cite}

\end{document}\documentclass{article}
\usepackage[utf8]{inputenc}
\usepackage{booktabs}
\usepackage{amsmath}
\usepackage{natbib}

\title{My Document}
\author{Your Name}
\date{\today}

\begin{document}

\maketitle

\section{Introduction}
This is a sample document demonstrating various \LaTeX{} features including mathematical equations, tables, and formatting. You can cite papers using the citation search feature, and they will be automatically added to your \texttt{cite.bib} file.

\section{Mathematical Expressions}
Here are some example equations:

\subsection{Quadratic Formula}
The quadratic formula is given by:
\[x = \frac{-b \pm \sqrt{b^2 - 4ac}}{2a}\]

\subsection{Matrix Example}
A sample matrix:
\[
\begin{pmatrix}
a_{11} & a_{12} & a_{13} \\
a_{21} & a_{22} & a_{23} \\
a_{31} & a_{32} & a_{33}
\end{pmatrix}
\]

\section{Tables}
Below is a sample table showing data:

\begin{table}[h]
\centering
\begin{tabular}{@{}lrrr@{}}
\toprule
\textbf{Category} & \textbf{Value 1} & \textbf{Value 2} & \textbf{Total} \\
\midrule
Item A & 10.5 & 20.3 & 30.8 \\
Item B & 15.2 & 25.1 & 40.3 \\
Item C & 22.1 & 12.7 & 34.8 \\
\bottomrule
\end{tabular}
\caption{Sample Data Table}
\label{tab:sample}
\end{table}

\section{Lists}
\subsection{Bullet Points}
\begin{itemize}
  \item First item
  \item Second item
  \item Third item with sub-items:
    \begin{itemize}
      \item Sub-item 1
      \item Sub-item 2
    \end{itemize}
\end{itemize}

\subsection{Numbered List}
\begin{enumerate}
  \item First step
  \item Second step
  \item Third step
\end{enumerate}

\section{Citations and References}
When you use the citation search feature, references will be automatically managed. For example, you might cite a paper like this: \cite{example2024}. All citations will be automatically added to your \texttt{cite.bib} file.

\section{Conclusion}
This template demonstrates basic \LaTeX{} features including mathematical formulas, tables, and lists. The bibliography system is now set up to work seamlessly with Bibby AI's citation management. Feel free to modify it for your needs!

\bibliographystyle{plain}
\bibliography{cite}

\end{document}\documentclass{article}
\usepackage[utf8]{inputenc}
\usepackage{booktabs}
\usepackage{amsmath}
\usepackage{natbib}

\title{My Document}
\author{Your Name}
\date{\today}

\begin{document}

\maketitle

\section{Introduction}
This is a sample document demonstrating various \LaTeX{} features including mathematical equations, tables, and formatting. You can cite papers using the citation search feature, and they will be automatically added to your \texttt{cite.bib} file.

\section{Mathematical Expressions}
Here are some example equations:

\subsection{Quadratic Formula}
The quadratic formula is given by:
\[x = \frac{-b \pm \sqrt{b^2 - 4ac}}{2a}\]

\subsection{Matrix Example}
A sample matrix:
\[
\begin{pmatrix}
a_{11} & a_{12} & a_{13} \\
a_{21} & a_{22} & a_{23} \\
a_{31} & a_{32} & a_{33}
\end{pmatrix}
\]

\section{Tables}
Below is a sample table showing data:

\begin{table}[h]
\centering
\begin{tabular}{@{}lrrr@{}}
\toprule
\textbf{Category} & \textbf{Value 1} & \textbf{Value 2} & \textbf{Total} \\
\midrule
Item A & 10.5 & 20.3 & 30.8 \\
Item B & 15.2 & 25.1 & 40.3 \\
Item C & 22.1 & 12.7 & 34.8 \\
\bottomrule
\end{tabular}
\caption{Sample Data Table}
\label{tab:sample}
\end{table}

\section{Lists}
\subsection{Bullet Points}
\begin{itemize}
  \item First item
  \item Second item
  \item Third item with sub-items:
    \begin{itemize}
      \item Sub-item 1
      \item Sub-item 2
    \end{itemize}
\end{itemize}

\subsection{Numbered List}
\begin{enumerate}
  \item First step
  \item Second step
  \item Third step
\end{enumerate}

\section{Citations and References}
When you use the citation search feature, references will be automatically managed. For example, you might cite a paper like this: \cite{example2024}. All citations will be automatically added to your \texttt{cite.bib} file.

\section{Conclusion}
This template demonstrates basic \LaTeX{} features including mathematical formulas, tables, and lists. The bibliography system is now set up to work seamlessly with Bibby AI's citation management. Feel free to modify it for your needs!

\bibliographystyle{plain}
\bibliography{cite}

\end{document}\documentclass{article}
\usepackage[utf8]{inputenc}
\usepackage{booktabs}
\usepackage{amsmath}
\usepackage{natbib}

\title{My Document}
\author{Your Name}
\date{\today}

\begin{document}

\maketitle

\section{Introduction}
This is a sample document demonstrating various \LaTeX{} features including mathematical equations, tables, and formatting. You can cite papers using the citation search feature, and they will be automatically added to your \texttt{cite.bib} file.

\section{Mathematical Expressions}
Here are some example equations:

\subsection{Quadratic Formula}
The quadratic formula is given by:
\[x = \frac{-b \pm \sqrt{b^2 - 4ac}}{2a}\]

\subsection{Matrix Example}
A sample matrix:
\[
\begin{pmatrix}
a_{11} & a_{12} & a_{13} \\
a_{21} & a_{22} & a_{23} \\
a_{31} & a_{32} & a_{33}
\end{pmatrix}
\]

\section{Tables}
Below is a sample table showing data:

\begin{table}[h]
\centering
\begin{tabular}{@{}lrrr@{}}
\toprule
\textbf{Category} & \textbf{Value 1} & \textbf{Value 2} & \textbf{Total} \\
\midrule
Item A & 10.5 & 20.3 & 30.8 \\
Item B & 15.2 & 25.1 & 40.3 \\
Item C & 22.1 & 12.7 & 34.8 \\
\bottomrule
\end{tabular}
\caption{Sample Data Table}
\label{tab:sample}
\end{table}

\section{Lists}
\subsection{Bullet Points}
\begin{itemize}
  \item First item
  \item Second item
  \item Third item with sub-items:
    \begin{itemize}
      \item Sub-item 1
      \item Sub-item 2
    \end{itemize}
\end{itemize}

\subsection{Numbered List}
\begin{enumerate}
  \item First step
  \item Second step
  \item Third step
\end{enumerate}

\section{Citations and References}
When you use the citation search feature, references will be automatically managed. For example, you might cite a paper like this: \cite{example2024}. All citations will be automatically added to your \texttt{cite.bib} file.

\section{Conclusion}
This template demonstrates basic \LaTeX{} features including mathematical formulas, tables, and lists. The bibliography system is now set up to work seamlessly with Bibby AI's citation management. Feel free to modify it for your needs!

\bibliographystyle{plain}
\bibliography{cite}

\end{document}\documentclass{article}
\usepackage[utf8]{inputenc}
\usepackage{booktabs}
\usepackage{amsmath}
\usepackage{natbib}

\title{My Document}
\author{Your Name}
\date{\today}

\begin{document}

\maketitle

\section{Introduction}
This is a sample document demonstrating various \LaTeX{} features including mathematical equations, tables, and formatting. You can cite papers using the citation search feature, and they will be automatically added to your \texttt{cite.bib} file.

\section{Mathematical Expressions}
Here are some example equations:

\subsection{Quadratic Formula}
The quadratic formula is given by:
\[x = \frac{-b \pm \sqrt{b^2 - 4ac}}{2a}\]

\subsection{Matrix Example}
A sample matrix:
\[
\begin{pmatrix}
a_{11} & a_{12} & a_{13} \\
a_{21} & a_{22} & a_{23} \\
a_{31} & a_{32} & a_{33}
\end{pmatrix}
\]

\section{Tables}
Below is a sample table showing data:

\begin{table}[h]
\centering
\begin{tabular}{@{}lrrr@{}}
\toprule
\textbf{Category} & \textbf{Value 1} & \textbf{Value 2} & \textbf{Total} \\
\midrule
Item A & 10.5 & 20.3 & 30.8 \\
Item B & 15.2 & 25.1 & 40.3 \\
Item C & 22.1 & 12.7 & 34.8 \\
\bottomrule
\end{tabular}
\caption{Sample Data Table}
\label{tab:sample}
\end{table}

\section{Lists}
\subsection{Bullet Points}
\begin{itemize}
  \item First item
  \item Second item
  \item Third item with sub-items:
    \begin{itemize}
      \item Sub-item 1
      \item Sub-item 2
    \end{itemize}
\end{itemize}

\subsection{Numbered List}
\begin{enumerate}
  \item First step
  \item Second step
  \item Third step
\end{enumerate}

\section{Citations and References}
When you use the citation search feature, references will be automatically managed. For example, you might cite a paper like this: \cite{example2024}. All citations will be automatically added to your \texttt{cite.bib} file.

\section{Conclusion}
This template demonstrates basic \LaTeX{} features including mathematical formulas, tables, and lists. The bibliography system is now set up to work seamlessly with Bibby AI's citation management. Feel free to modify it for your needs!

\bibliographystyle{plain}
\bibliography{cite}

\end{document}\documentclass{article}
\usepackage[utf8]{inputenc}
\usepackage{booktabs}
\usepackage{amsmath}
\usepackage{natbib}

\title{My Document}
\author{Your Name}
\date{\today}

\begin{document}

\maketitle

\section{Introduction}
This is a sample document demonstrating various \LaTeX{} features including mathematical equations, tables, and formatting. You can cite papers using the citation search feature, and they will be automatically added to your \texttt{cite.bib} file.

\section{Mathematical Expressions}
Here are some example equations:

\subsection{Quadratic Formula}
The quadratic formula is given by:
\[x = \frac{-b \pm \sqrt{b^2 - 4ac}}{2a}\]

\subsection{Matrix Example}
A sample matrix:
\[
\begin{pmatrix}
a_{11} & a_{12} & a_{13} \\
a_{21} & a_{22} & a_{23} \\
a_{31} & a_{32} & a_{33}
\end{pmatrix}
\]

\section{Tables}
Below is a sample table showing data:

\begin{table}[h]
\centering
\begin{tabular}{@{}lrrr@{}}
\toprule
\textbf{Category} & \textbf{Value 1} & \textbf{Value 2} & \textbf{Total} \\
\midrule
Item A & 10.5 & 20.3 & 30.8 \\
Item B & 15.2 & 25.1 & 40.3 \\
Item C & 22.1 & 12.7 & 34.8 \\
\bottomrule
\end{tabular}
\caption{Sample Data Table}
\label{tab:sample}
\end{table}

\section{Lists}
\subsection{Bullet Points}
\begin{itemize}
  \item First item
  \item Second item
  \item Third item with sub-items:
    \begin{itemize}
      \item Sub-item 1
      \item Sub-item 2
    \end{itemize}
\end{itemize}

\subsection{Numbered List}
\begin{enumerate}
  \item First step
  \item Second step
  \item Third step
\end{enumerate}

\section{Citations and References}
When you use the citation search feature, references will be automatically managed. For example, you might cite a paper like this: \cite{example2024}. All citations will be automatically added to your \texttt{cite.bib} file.

\section{Conclusion}
This template demonstrates basic \LaTeX{} features including mathematical formulas, tables, and lists. The bibliography system is now set up to work seamlessly with Bibby AI's citation management. Feel free to modify it for your needs!

\bibliographystyle{plain}
\bibliography{cite}

\end{document}\documentclass{article}
\usepackage[utf8]{inputenc}
\usepackage{booktabs}
\usepackage{amsmath}
\usepackage{natbib}

\title{My Document}
\author{Your Name}
\date{\today}

\begin{document}

\maketitle

\section{Introduction}
This is a sample document demonstrating various \LaTeX{} features including mathematical equations, tables, and formatting. You can cite papers using the citation search feature, and they will be automatically added to your \texttt{cite.bib} file.

\section{Mathematical Expressions}
Here are some example equations:

\subsection{Quadratic Formula}
The quadratic formula is given by:
\[x = \frac{-b \pm \sqrt{b^2 - 4ac}}{2a}\]

\subsection{Matrix Example}
A sample matrix:
\[
\begin{pmatrix}
a_{11} & a_{12} & a_{13} \\
a_{21} & a_{22} & a_{23} \\
a_{31} & a_{32} & a_{33}
\end{pmatrix}
\]

\section{Tables}
Below is a sample table showing data:

\begin{table}[h]
\centering
\begin{tabular}{@{}lrrr@{}}
\toprule
\textbf{Category} & \textbf{Value 1} & \textbf{Value 2} & \textbf{Total} \\
\midrule
Item A & 10.5 & 20.3 & 30.8 \\
Item B & 15.2 & 25.1 & 40.3 \\
Item C & 22.1 & 12.7 & 34.8 \\
\bottomrule
\end{tabular}
\caption{Sample Data Table}
\label{tab:sample}
\end{table}

\section{Lists}
\subsection{Bullet Points}
\begin{itemize}
  \item First item
  \item Second item
  \item Third item with sub-items:
    \begin{itemize}
      \item Sub-item 1
      \item Sub-item 2
    \end{itemize}
\end{itemize}

\subsection{Numbered List}
\begin{enumerate}
  \item First step
  \item Second step
  \item Third step
\end{enumerate}

\section{Citations and References}
When you use the citation search feature, references will be automatically managed. For example, you might cite a paper like this: \cite{example2024}. All citations will be automatically added to your \texttt{cite.bib} file.

\section{Conclusion}
This template demonstrates basic \LaTeX{} features including mathematical formulas, tables, and lists. The bibliography system is now set up to work seamlessly with Bibby AI's citation management. Feel free to modify it for your needs!

\bibliographystyle{plain}
\bibliography{cite}

\end{document}\documentclass{article}
\usepackage[utf8]{inputenc}
\usepackage{booktabs}
\usepackage{amsmath}
\usepackage{natbib}

\title{My Document}
\author{Your Name}
\date{\today}

\begin{document}

\maketitle

\section{Introduction}
This is a sample document demonstrating various \LaTeX{} features including mathematical equations, tables, and formatting. You can cite papers using the citation search feature, and they will be automatically added to your \texttt{cite.bib} file.

\section{Mathematical Expressions}
Here are some example equations:

\subsection{Quadratic Formula}
The quadratic formula is given by:
\[x = \frac{-b \pm \sqrt{b^2 - 4ac}}{2a}\]

\subsection{Matrix Example}
A sample matrix:
\[
\begin{pmatrix}
a_{11} & a_{12} & a_{13} \\
a_{21} & a_{22} & a_{23} \\
a_{31} & a_{32} & a_{33}
\end{pmatrix}
\]

\section{Tables}
Below is a sample table showing data:

\begin{table}[h]
\centering
\begin{tabular}{@{}lrrr@{}}
\toprule
\textbf{Category} & \textbf{Value 1} & \textbf{Value 2} & \textbf{Total} \\
\midrule
Item A & 10.5 & 20.3 & 30.8 \\
Item B & 15.2 & 25.1 & 40.3 \\
Item C & 22.1 & 12.7 & 34.8 \\
\bottomrule
\end{tabular}
\caption{Sample Data Table}
\label{tab:sample}
\end{table}

\section{Lists}
\subsection{Bullet Points}
\begin{itemize}
  \item First item
  \item Second item
  \item Third item with sub-items:
    \begin{itemize}
      \item Sub-item 1
      \item Sub-item 2
    \end{itemize}
\end{itemize}

\subsection{Numbered List}
\begin{enumerate}
  \item First step
  \item Second step
  \item Third step
\end{enumerate}

\section{Citations and References}
When you use the citation search feature, references will be automatically managed. For example, you might cite a paper like this: \cite{example2024}. All citations will be automatically added to your \texttt{cite.bib} file.

\section{Conclusion}
This template demonstrates basic \LaTeX{} features including mathematical formulas, tables, and lists. The bibliography system is now set up to work seamlessly with Bibby AI's citation management. Feel free to modify it for your needs!

\bibliographystyle{plain}
\bibliography{cite}

\end{document}\documentclass{article}
\usepackage[utf8]{inputenc}
\usepackage{booktabs}
\usepackage{amsmath}
\usepackage{natbib}

\title{My Document}
\author{Your Name}
\date{\today}

\begin{document}

\maketitle

\section{Introduction}
This is a sample document demonstrating various \LaTeX{} features including mathematical equations, tables, and formatting. You can cite papers using the citation search feature, and they will be automatically added to your \texttt{cite.bib} file.

\section{Mathematical Expressions}
Here are some example equations:

\subsection{Quadratic Formula}
The quadratic formula is given by:
\[x = \frac{-b \pm \sqrt{b^2 - 4ac}}{2a}\]

\subsection{Matrix Example}
A sample matrix:
\[
\begin{pmatrix}
a_{11} & a_{12} & a_{13} \\
a_{21} & a_{22} & a_{23} \\
a_{31} & a_{32} & a_{33}
\end{pmatrix}
\]

\section{Tables}
Below is a sample table showing data:

\begin{table}[h]
\centering
\begin{tabular}{@{}lrrr@{}}
\toprule
\textbf{Category} & \textbf{Value 1} & \textbf{Value 2} & \textbf{Total} \\
\midrule
Item A & 10.5 & 20.3 & 30.8 \\
Item B & 15.2 & 25.1 & 40.3 \\
Item C & 22.1 & 12.7 & 34.8 \\
\bottomrule
\end{tabular}
\caption{Sample Data Table}
\label{tab:sample}
\end{table}

\section{Lists}
\subsection{Bullet Points}
\begin{itemize}
  \item First item
  \item Second item
  \item Third item with sub-items:
    \begin{itemize}
      \item Sub-item 1
      \item Sub-item 2
    \end{itemize}
\end{itemize}

\subsection{Numbered List}
\begin{enumerate}
  \item First step
  \item Second step
  \item Third step
\end{enumerate}

\section{Citations and References}
When you use the citation search feature, references will be automatically managed. For example, you might cite a paper like this: \cite{example2024}. All citations will be automatically added to your \texttt{cite.bib} file.

\section{Conclusion}
This template demonstrates basic \LaTeX{} features including mathematical formulas, tables, and lists. The bibliography system is now set up to work seamlessly with Bibby AI's citation management. Feel free to modify it for your needs!

\bibliographystyle{plain}
\bibliography{cite}

\end{document}\documentclass{article}
\usepackage[utf8]{inputenc}
\usepackage{booktabs}
\usepackage{amsmath}
\usepackage{natbib}

\title{My Document}
\author{Your Name}
\date{\today}

\begin{document}

\maketitle

\section{Introduction}
This is a sample document demonstrating various \LaTeX{} features including mathematical equations, tables, and formatting. You can cite papers using the citation search feature, and they will be automatically added to your \texttt{cite.bib} file.

\section{Mathematical Expressions}
Here are some example equations:

\subsection{Quadratic Formula}
The quadratic formula is given by:
\[x = \frac{-b \pm \sqrt{b^2 - 4ac}}{2a}\]

\subsection{Matrix Example}
A sample matrix:
\[
\begin{pmatrix}
a_{11} & a_{12} & a_{13} \\
a_{21} & a_{22} & a_{23} \\
a_{31} & a_{32} & a_{33}
\end{pmatrix}
\]

\section{Tables}
Below is a sample table showing data:

\begin{table}[h]
\centering
\begin{tabular}{@{}lrrr@{}}
\toprule
\textbf{Category} & \textbf{Value 1} & \textbf{Value 2} & \textbf{Total} \\
\midrule
Item A & 10.5 & 20.3 & 30.8 \\
Item B & 15.2 & 25.1 & 40.3 \\
Item C & 22.1 & 12.7 & 34.8 \\
\bottomrule
\end{tabular}
\caption{Sample Data Table}
\label{tab:sample}
\end{table}

\section{Lists}
\subsection{Bullet Points}
\begin{itemize}
  \item First item
  \item Second item
  \item Third item with sub-items:
    \begin{itemize}
      \item Sub-item 1
      \item Sub-item 2
    \end{itemize}
\end{itemize}

\subsection{Numbered List}
\begin{enumerate}
  \item First step
  \item Second step
  \item Third step
\end{enumerate}

\section{Citations and References}
When you use the citation search feature, references will be automatically managed. For example, you might cite a paper like this: \cite{example2024}. All citations will be automatically added to your \texttt{cite.bib} file.

\section{Conclusion}
This template demonstrates basic \LaTeX{} features including mathematical formulas, tables, and lists. The bibliography system is now set up to work seamlessly with Bibby AI's citation management. Feel free to modify it for your needs!

\bibliographystyle{plain}
\bibliography{cite}

\end{document}\documentclass{article}
\usepackage[utf8]{inputenc}
\usepackage{booktabs}
\usepackage{amsmath}
\usepackage{natbib}

\title{My Document}
\author{Your Name}
\date{\today}

\begin{document}

\maketitle

\section{Introduction}
This is a sample document demonstrating various \LaTeX{} features including mathematical equations, tables, and formatting. You can cite papers using the citation search feature, and they will be automatically added to your \texttt{cite.bib} file.

\section{Mathematical Expressions}
Here are some example equations:

\subsection{Quadratic Formula}
The quadratic formula is given by:
\[x = \frac{-b \pm \sqrt{b^2 - 4ac}}{2a}\]

\subsection{Matrix Example}
A sample matrix:
\[
\begin{pmatrix}
a_{11} & a_{12} & a_{13} \\
a_{21} & a_{22} & a_{23} \\
a_{31} & a_{32} & a_{33}
\end{pmatrix}
\]

\section{Tables}
Below is a sample table showing data:

\begin{table}[h]
\centering
\begin{tabular}{@{}lrrr@{}}
\toprule
\textbf{Category} & \textbf{Value 1} & \textbf{Value 2} & \textbf{Total} \\
\midrule
Item A & 10.5 & 20.3 & 30.8 \\
Item B & 15.2 & 25.1 & 40.3 \\
Item C & 22.1 & 12.7 & 34.8 \\
\bottomrule
\end{tabular}
\caption{Sample Data Table}
\label{tab:sample}
\end{table}

\section{Lists}
\subsection{Bullet Points}
\begin{itemize}
  \item First item
  \item Second item
  \item Third item with sub-items:
    \begin{itemize}
      \item Sub-item 1
      \item Sub-item 2
    \end{itemize}
\end{itemize}

\subsection{Numbered List}
\begin{enumerate}
  \item First step
  \item Second step
  \item Third step
\end{enumerate}

\section{Citations and References}
When you use the citation search feature, references will be automatically managed. For example, you might cite a paper like this: \cite{example2024}. All citations will be automatically added to your \texttt{cite.bib} file.

\section{Conclusion}
This template demonstrates basic \LaTeX{} features including mathematical formulas, tables, and lists. The bibliography system is now set up to work seamlessly with Bibby AI's citation management. Feel free to modify it for your needs!

\bibliographystyle{plain}
\bibliography{cite}

\end{document}\documentclass{article}
\usepackage[utf8]{inputenc}
\usepackage{booktabs}
\usepackage{amsmath}
\usepackage{natbib}

\title{My Document}
\author{Your Name}
\date{\today}

\begin{document}

\maketitle

\section{Introduction}
This is a sample document demonstrating various \LaTeX{} features including mathematical equations, tables, and formatting. You can cite papers using the citation search feature, and they will be automatically added to your \texttt{cite.bib} file.

\section{Mathematical Expressions}
Here are some example equations:

\subsection{Quadratic Formula}
The quadratic formula is given by:
\[x = \frac{-b \pm \sqrt{b^2 - 4ac}}{2a}\]

\subsection{Matrix Example}
A sample matrix:
\[
\begin{pmatrix}
a_{11} & a_{12} & a_{13} \\
a_{21} & a_{22} & a_{23} \\
a_{31} & a_{32} & a_{33}
\end{pmatrix}
\]

\section{Tables}
Below is a sample table showing data:

\begin{table}[h]
\centering
\begin{tabular}{@{}lrrr@{}}
\toprule
\textbf{Category} & \textbf{Value 1} & \textbf{Value 2} & \textbf{Total} \\
\midrule
Item A & 10.5 & 20.3 & 30.8 \\
Item B & 15.2 & 25.1 & 40.3 \\
Item C & 22.1 & 12.7 & 34.8 \\
\bottomrule
\end{tabular}
\caption{Sample Data Table}
\label{tab:sample}
\end{table}

\section{Lists}
\subsection{Bullet Points}
\begin{itemize}
  \item First item
  \item Second item
  \item Third item with sub-items:
    \begin{itemize}
      \item Sub-item 1
      \item Sub-item 2
    \end{itemize}
\end{itemize}

\subsection{Numbered List}
\begin{enumerate}
  \item First step
  \item Second step
  \item Third step
\end{enumerate}

\section{Citations and References}
When you use the citation search feature, references will be automatically managed. For example, you might cite a paper like this: \cite{example2024}. All citations will be automatically added to your \texttt{cite.bib} file.

\section{Conclusion}
This template demonstrates basic \LaTeX{} features including mathematical formulas, tables, and lists. The bibliography system is now set up to work seamlessly with Bibby AI's citation management. Feel free to modify it for your needs!

\bibliographystyle{plain}
\bibliography{cite}

\end{document}\documentclass{article}
\usepackage[utf8]{inputenc}
\usepackage{booktabs}
\usepackage{amsmath}
\usepackage{natbib}

\title{My Document}
\author{Your Name}
\date{\today}

\begin{document}

\maketitle

\section{Introduction}
This is a sample document demonstrating various \LaTeX{} features including mathematical equations, tables, and formatting. You can cite papers using the citation search feature, and they will be automatically added to your \texttt{cite.bib} file.

\section{Mathematical Expressions}
Here are some example equations:

\subsection{Quadratic Formula}
The quadratic formula is given by:
\[x = \frac{-b \pm \sqrt{b^2 - 4ac}}{2a}\]

\subsection{Matrix Example}
A sample matrix:
\[
\begin{pmatrix}
a_{11} & a_{12} & a_{13} \\
a_{21} & a_{22} & a_{23} \\
a_{31} & a_{32} & a_{33}
\end{pmatrix}
\]

\section{Tables}
Below is a sample table showing data:

\begin{table}[h]
\centering
\begin{tabular}{@{}lrrr@{}}
\toprule
\textbf{Category} & \textbf{Value 1} & \textbf{Value 2} & \textbf{Total} \\
\midrule
Item A & 10.5 & 20.3 & 30.8 \\
Item B & 15.2 & 25.1 & 40.3 \\
Item C & 22.1 & 12.7 & 34.8 \\
\bottomrule
\end{tabular}
\caption{Sample Data Table}
\label{tab:sample}
\end{table}

\section{Lists}
\subsection{Bullet Points}
\begin{itemize}
  \item First item
  \item Second item
  \item Third item with sub-items:
    \begin{itemize}
      \item Sub-item 1
      \item Sub-item 2
    \end{itemize}
\end{itemize}

\subsection{Numbered List}
\begin{enumerate}
  \item First step
  \item Second step
  \item Third step
\end{enumerate}

\section{Citations and References}
When you use the citation search feature, references will be automatically managed. For example, you might cite a paper like this: \cite{example2024}. All citations will be automatically added to your \texttt{cite.bib} file.

\section{Conclusion}
This template demonstrates basic \LaTeX{} features including mathematical formulas, tables, and lists. The bibliography system is now set up to work seamlessly with Bibby AI's citation management. Feel free to modify it for your needs!

\bibliographystyle{plain}
\bibliography{cite}

\end{document}\documentclass{article}
\usepackage[utf8]{inputenc}
\usepackage{booktabs}
\usepackage{amsmath}
\usepackage{natbib}

\title{My Document}
\author{Your Name}
\date{\today}

\begin{document}

\maketitle

\section{Introduction}
This is a sample document demonstrating various \LaTeX{} features including mathematical equations, tables, and formatting. You can cite papers using the citation search feature, and they will be automatically added to your \texttt{cite.bib} file.

\section{Mathematical Expressions}
Here are some example equations:

\subsection{Quadratic Formula}
The quadratic formula is given by:
\[x = \frac{-b \pm \sqrt{b^2 - 4ac}}{2a}\]

\subsection{Matrix Example}
A sample matrix:
\[
\begin{pmatrix}
a_{11} & a_{12} & a_{13} \\
a_{21} & a_{22} & a_{23} \\
a_{31} & a_{32} & a_{33}
\end{pmatrix}
\]

\section{Tables}
Below is a sample table showing data:

\begin{table}[h]
\centering
\begin{tabular}{@{}lrrr@{}}
\toprule
\textbf{Category} & \textbf{Value 1} & \textbf{Value 2} & \textbf{Total} \\
\midrule
Item A & 10.5 & 20.3 & 30.8 \\
Item B & 15.2 & 25.1 & 40.3 \\
Item C & 22.1 & 12.7 & 34.8 \\
\bottomrule
\end{tabular}
\caption{Sample Data Table}
\label{tab:sample}
\end{table}

\section{Lists}
\subsection{Bullet Points}
\begin{itemize}
  \item First item
  \item Second item
  \item Third item with sub-items:
    \begin{itemize}
      \item Sub-item 1
      \item Sub-item 2
    \end{itemize}
\end{itemize}

\subsection{Numbered List}
\begin{enumerate}
  \item First step
  \item Second step
  \item Third step
\end{enumerate}

\section{Citations and References}
When you use the citation search feature, references will be automatically managed. For example, you might cite a paper like this: \cite{example2024}. All citations will be automatically added to your \texttt{cite.bib} file.

\section{Conclusion}
This template demonstrates basic \LaTeX{} features including mathematical formulas, tables, and lists. The bibliography system is now set up to work seamlessly with Bibby AI's citation management. Feel free to modify it for your needs!

\bibliographystyle{plain}
\bibliography{cite}

\end{document}\documentclass{article}
\usepackage[utf8]{inputenc}
\usepackage{booktabs}
\usepackage{amsmath}
\usepackage{natbib}

\title{My Document}
\author{Your Name}
\date{\today}

\begin{document}

\maketitle

\section{Introduction}
This is a sample document demonstrating various \LaTeX{} features including mathematical equations, tables, and formatting. You can cite papers using the citation search feature, and they will be automatically added to your \texttt{cite.bib} file.

\section{Mathematical Expressions}
Here are some example equations:

\subsection{Quadratic Formula}
The quadratic formula is given by:
\[x = \frac{-b \pm \sqrt{b^2 - 4ac}}{2a}\]

\subsection{Matrix Example}
A sample matrix:
\[
\begin{pmatrix}
a_{11} & a_{12} & a_{13} \\
a_{21} & a_{22} & a_{23} \\
a_{31} & a_{32} & a_{33}
\end{pmatrix}
\]

\section{Tables}
Below is a sample table showing data:

\begin{table}[h]
\centering
\begin{tabular}{@{}lrrr@{}}
\toprule
\textbf{Category} & \textbf{Value 1} & \textbf{Value 2} & \textbf{Total} \\
\midrule
Item A & 10.5 & 20.3 & 30.8 \\
Item B & 15.2 & 25.1 & 40.3 \\
Item C & 22.1 & 12.7 & 34.8 \\
\bottomrule
\end{tabular}
\caption{Sample Data Table}
\label{tab:sample}
\end{table}

\section{Lists}
\subsection{Bullet Points}
\begin{itemize}
  \item First item
  \item Second item
  \item Third item with sub-items:
    \begin{itemize}
      \item Sub-item 1
      \item Sub-item 2
    \end{itemize}
\end{itemize}

\subsection{Numbered List}
\begin{enumerate}
  \item First step
  \item Second step
  \item Third step
\end{enumerate}

\section{Citations and References}
When you use the citation search feature, references will be automatically managed. For example, you might cite a paper like this: \cite{example2024}. All citations will be automatically added to your \texttt{cite.bib} file.

\section{Conclusion}
This template demonstrates basic \LaTeX{} features including mathematical formulas, tables, and lists. The bibliography system is now set up to work seamlessly with Bibby AI's citation management. Feel free to modify it for your needs!

\bibliographystyle{plain}
\bibliography{cite}

\end{document}\documentclass{article}
\usepackage[utf8]{inputenc}
\usepackage{booktabs}
\usepackage{amsmath}
\usepackage{natbib}

\title{My Document}
\author{Your Name}
\date{\today}

\begin{document}

\maketitle

\section{Introduction}
This is a sample document demonstrating various \LaTeX{} features including mathematical equations, tables, and formatting. You can cite papers using the citation search feature, and they will be automatically added to your \texttt{cite.bib} file.

\section{Mathematical Expressions}
Here are some example equations:

\subsection{Quadratic Formula}
The quadratic formula is given by:
\[x = \frac{-b \pm \sqrt{b^2 - 4ac}}{2a}\]

\subsection{Matrix Example}
A sample matrix:
\[
\begin{pmatrix}
a_{11} & a_{12} & a_{13} \\
a_{21} & a_{22} & a_{23} \\
a_{31} & a_{32} & a_{33}
\end{pmatrix}
\]

\section{Tables}
Below is a sample table showing data:

\begin{table}[h]
\centering
\begin{tabular}{@{}lrrr@{}}
\toprule
\textbf{Category} & \textbf{Value 1} & \textbf{Value 2} & \textbf{Total} \\
\midrule
Item A & 10.5 & 20.3 & 30.8 \\
Item B & 15.2 & 25.1 & 40.3 \\
Item C & 22.1 & 12.7 & 34.8 \\
\bottomrule
\end{tabular}
\caption{Sample Data Table}
\label{tab:sample}
\end{table}

\section{Lists}
\subsection{Bullet Points}
\begin{itemize}
  \item First item
  \item Second item
  \item Third item with sub-items:
    \begin{itemize}
      \item Sub-item 1
      \item Sub-item 2
    \end{itemize}
\end{itemize}

\subsection{Numbered List}
\begin{enumerate}
  \item First step
  \item Second step
  \item Third step
\end{enumerate}

\section{Citations and References}
When you use the citation search feature, references will be automatically managed. For example, you might cite a paper like this: \cite{example2024}. All citations will be automatically added to your \texttt{cite.bib} file.

\section{Conclusion}
This template demonstrates basic \LaTeX{} features including mathematical formulas, tables, and lists. The bibliography system is now set up to work seamlessly with Bibby AI's citation management. Feel free to modify it for your needs!

\bibliographystyle{plain}
\bibliography{cite}

\end{document}\documentclass{article}
\usepackage[utf8]{inputenc}
\usepackage{booktabs}
\usepackage{amsmath}
\usepackage{natbib}

\title{My Document}
\author{Your Name}
\date{\today}

\begin{document}

\maketitle

\section{Introduction}
This is a sample document demonstrating various \LaTeX{} features including mathematical equations, tables, and formatting. You can cite papers using the citation search feature, and they will be automatically added to your \texttt{cite.bib} file.

\section{Mathematical Expressions}
Here are some example equations:

\subsection{Quadratic Formula}
The quadratic formula is given by:
\[x = \frac{-b \pm \sqrt{b^2 - 4ac}}{2a}\]

\subsection{Matrix Example}
A sample matrix:
\[
\begin{pmatrix}
a_{11} & a_{12} & a_{13} \\
a_{21} & a_{22} & a_{23} \\
a_{31} & a_{32} & a_{33}
\end{pmatrix}
\]

\section{Tables}
Below is a sample table showing data:

\begin{table}[h]
\centering
\begin{tabular}{@{}lrrr@{}}
\toprule
\textbf{Category} & \textbf{Value 1} & \textbf{Value 2} & \textbf{Total} \\
\midrule
Item A & 10.5 & 20.3 & 30.8 \\
Item B & 15.2 & 25.1 & 40.3 \\
Item C & 22.1 & 12.7 & 34.8 \\
\bottomrule
\end{tabular}
\caption{Sample Data Table}
\label{tab:sample}
\end{table}

\section{Lists}
\subsection{Bullet Points}
\begin{itemize}
  \item First item
  \item Second item
  \item Third item with sub-items:
    \begin{itemize}
      \item Sub-item 1
      \item Sub-item 2
    \end{itemize}
\end{itemize}

\subsection{Numbered List}
\begin{enumerate}
  \item First step
  \item Second step
  \item Third step
\end{enumerate}

\section{Citations and References}
When you use the citation search feature, references will be automatically managed. For example, you might cite a paper like this: \cite{example2024}. All citations will be automatically added to your \texttt{cite.bib} file.

\section{Conclusion}
This template demonstrates basic \LaTeX{} features including mathematical formulas, tables, and lists. The bibliography system is now set up to work seamlessly with Bibby AI's citation management. Feel free to modify it for your needs!

\bibliographystyle{plain}
\bibliography{cite}

\end{document}\documentclass{article}
\usepackage[utf8]{inputenc}
\usepackage{booktabs}
\usepackage{amsmath}
\usepackage{natbib}

\title{My Document}
\author{Your Name}
\date{\today}

\begin{document}

\maketitle

\section{Introduction}
This is a sample document demonstrating various \LaTeX{} features including mathematical equations, tables, and formatting. You can cite papers using the citation search feature, and they will be automatically added to your \texttt{cite.bib} file.

\section{Mathematical Expressions}
Here are some example equations:

\subsection{Quadratic Formula}
The quadratic formula is given by:
\[x = \frac{-b \pm \sqrt{b^2 - 4ac}}{2a}\]

\subsection{Matrix Example}
A sample matrix:
\[
\begin{pmatrix}
a_{11} & a_{12} & a_{13} \\
a_{21} & a_{22} & a_{23} \\
a_{31} & a_{32} & a_{33}
\end{pmatrix}
\]

\section{Tables}
Below is a sample table showing data:

\begin{table}[h]
\centering
\begin{tabular}{@{}lrrr@{}}
\toprule
\textbf{Category} & \textbf{Value 1} & \textbf{Value 2} & \textbf{Total} \\
\midrule
Item A & 10.5 & 20.3 & 30.8 \\
Item B & 15.2 & 25.1 & 40.3 \\
Item C & 22.1 & 12.7 & 34.8 \\
\bottomrule
\end{tabular}
\caption{Sample Data Table}
\label{tab:sample}
\end{table}

\section{Lists}
\subsection{Bullet Points}
\begin{itemize}
  \item First item
  \item Second item
  \item Third item with sub-items:
    \begin{itemize}
      \item Sub-item 1
      \item Sub-item 2
    \end{itemize}
\end{itemize}

\subsection{Numbered List}
\begin{enumerate}
  \item First step
  \item Second step
  \item Third step
\end{enumerate}

\section{Citations and References}
When you use the citation search feature, references will be automatically managed. For example, you might cite a paper like this: \cite{example2024}. All citations will be automatically added to your \texttt{cite.bib} file.

\section{Conclusion}
This template demonstrates basic \LaTeX{} features including mathematical formulas, tables, and lists. The bibliography system is now set up to work seamlessly with Bibby AI's citation management. Feel free to modify it for your needs!

\bibliographystyle{plain}
\bibliography{cite}

\end{document}\documentclass{article}
\usepackage[utf8]{inputenc}
\usepackage{booktabs}
\usepackage{amsmath}
\usepackage{natbib}

\title{My Document}
\author{Your Name}
\date{\today}

\begin{document}

\maketitle

\section{Introduction}
This is a sample document demonstrating various \LaTeX{} features including mathematical equations, tables, and formatting. You can cite papers using the citation search feature, and they will be automatically added to your \texttt{cite.bib} file.

\section{Mathematical Expressions}
Here are some example equations:

\subsection{Quadratic Formula}
The quadratic formula is given by:
\[x = \frac{-b \pm \sqrt{b^2 - 4ac}}{2a}\]

\subsection{Matrix Example}
A sample matrix:
\[
\begin{pmatrix}
a_{11} & a_{12} & a_{13} \\
a_{21} & a_{22} & a_{23} \\
a_{31} & a_{32} & a_{33}
\end{pmatrix}
\]

\section{Tables}
Below is a sample table showing data:

\begin{table}[h]
\centering
\begin{tabular}{@{}lrrr@{}}
\toprule
\textbf{Category} & \textbf{Value 1} & \textbf{Value 2} & \textbf{Total} \\
\midrule
Item A & 10.5 & 20.3 & 30.8 \\
Item B & 15.2 & 25.1 & 40.3 \\
Item C & 22.1 & 12.7 & 34.8 \\
\bottomrule
\end{tabular}
\caption{Sample Data Table}
\label{tab:sample}
\end{table}

\section{Lists}
\subsection{Bullet Points}
\begin{itemize}
  \item First item
  \item Second item
  \item Third item with sub-items:
    \begin{itemize}
      \item Sub-item 1
      \item Sub-item 2
    \end{itemize}
\end{itemize}

\subsection{Numbered List}
\begin{enumerate}
  \item First step
  \item Second step
  \item Third step
\end{enumerate}

\section{Citations and References}
When you use the citation search feature, references will be automatically managed. For example, you might cite a paper like this: \cite{example2024}. All citations will be automatically added to your \texttt{cite.bib} file.

\section{Conclusion}
This template demonstrates basic \LaTeX{} features including mathematical formulas, tables, and lists. The bibliography system is now set up to work seamlessly with Bibby AI's citation management. Feel free to modify it for your needs!

\bibliographystyle{plain}
\bibliography{cite}

\end{document}\documentclass{article}
\usepackage[utf8]{inputenc}
\usepackage{booktabs}
\usepackage{amsmath}
\usepackage{natbib}

\title{My Document}
\author{Your Name}
\date{\today}

\begin{document}

\maketitle

\section{Introduction}
This is a sample document demonstrating various \LaTeX{} features including mathematical equations, tables, and formatting. You can cite papers using the citation search feature, and they will be automatically added to your \texttt{cite.bib} file.

\section{Mathematical Expressions}
Here are some example equations:

\subsection{Quadratic Formula}
The quadratic formula is given by:
\[x = \frac{-b \pm \sqrt{b^2 - 4ac}}{2a}\]

\subsection{Matrix Example}
A sample matrix:
\[
\begin{pmatrix}
a_{11} & a_{12} & a_{13} \\
a_{21} & a_{22} & a_{23} \\
a_{31} & a_{32} & a_{33}
\end{pmatrix}
\]

\section{Tables}
Below is a sample table showing data:

\begin{table}[h]
\centering
\begin{tabular}{@{}lrrr@{}}
\toprule
\textbf{Category} & \textbf{Value 1} & \textbf{Value 2} & \textbf{Total} \\
\midrule
Item A & 10.5 & 20.3 & 30.8 \\
Item B & 15.2 & 25.1 & 40.3 \\
Item C & 22.1 & 12.7 & 34.8 \\
\bottomrule
\end{tabular}
\caption{Sample Data Table}
\label{tab:sample}
\end{table}

\section{Lists}
\subsection{Bullet Points}
\begin{itemize}
  \item First item
  \item Second item
  \item Third item with sub-items:
    \begin{itemize}
      \item Sub-item 1
      \item Sub-item 2
    \end{itemize}
\end{itemize}

\subsection{Numbered List}
\begin{enumerate}
  \item First step
  \item Second step
  \item Third step
\end{enumerate}

\section{Citations and References}
When you use the citation search feature, references will be automatically managed. For example, you might cite a paper like this: \cite{example2024}. All citations will be automatically added to your \texttt{cite.bib} file.

\section{Conclusion}
This template demonstrates basic \LaTeX{} features including mathematical formulas, tables, and lists. The bibliography system is now set up to work seamlessly with Bibby AI's citation management. Feel free to modify it for your needs!

\bibliographystyle{plain}
\bibliography{cite}

\end{document}\documentclass{article}
\usepackage[utf8]{inputenc}
\usepackage{booktabs}
\usepackage{amsmath}
\usepackage{natbib}

\title{My Document}
\author{Your Name}
\date{\today}

\begin{document}

\maketitle

\section{Introduction}
This is a sample document demonstrating various \LaTeX{} features including mathematical equations, tables, and formatting. You can cite papers using the citation search feature, and they will be automatically added to your \texttt{cite.bib} file.

\section{Mathematical Expressions}
Here are some example equations:

\subsection{Quadratic Formula}
The quadratic formula is given by:
\[x = \frac{-b \pm \sqrt{b^2 - 4ac}}{2a}\]

\subsection{Matrix Example}
A sample matrix:
\[
\begin{pmatrix}
a_{11} & a_{12} & a_{13} \\
a_{21} & a_{22} & a_{23} \\
a_{31} & a_{32} & a_{33}
\end{pmatrix}
\]

\section{Tables}
Below is a sample table showing data:

\begin{table}[h]
\centering
\begin{tabular}{@{}lrrr@{}}
\toprule
\textbf{Category} & \textbf{Value 1} & \textbf{Value 2} & \textbf{Total} \\
\midrule
Item A & 10.5 & 20.3 & 30.8 \\
Item B & 15.2 & 25.1 & 40.3 \\
Item C & 22.1 & 12.7 & 34.8 \\
\bottomrule
\end{tabular}
\caption{Sample Data Table}
\label{tab:sample}
\end{table}

\section{Lists}
\subsection{Bullet Points}
\begin{itemize}
  \item First item
  \item Second item
  \item Third item with sub-items:
    \begin{itemize}
      \item Sub-item 1
      \item Sub-item 2
    \end{itemize}
\end{itemize}

\subsection{Numbered List}
\begin{enumerate}
  \item First step
  \item Second step
  \item Third step
\end{enumerate}

\section{Citations and References}
When you use the citation search feature, references will be automatically managed. For example, you might cite a paper like this: \cite{example2024}. All citations will be automatically added to your \texttt{cite.bib} file.

\section{Conclusion}
This template demonstrates basic \LaTeX{} features including mathematical formulas, tables, and lists. The bibliography system is now set up to work seamlessly with Bibby AI's citation management. Feel free to modify it for your needs!

\bibliographystyle{plain}
\bibliography{cite}

\end{document}\documentclass{article}
\usepackage[utf8]{inputenc}
\usepackage{booktabs}
\usepackage{amsmath}
\usepackage{natbib}

\title{My Document}
\author{Your Name}
\date{\today}

\begin{document}

\maketitle

\section{Introduction}
This is a sample document demonstrating various \LaTeX{} features including mathematical equations, tables, and formatting. You can cite papers using the citation search feature, and they will be automatically added to your \texttt{cite.bib} file.

\section{Mathematical Expressions}
Here are some example equations:

\subsection{Quadratic Formula}
The quadratic formula is given by:
\[x = \frac{-b \pm \sqrt{b^2 - 4ac}}{2a}\]

\subsection{Matrix Example}
A sample matrix:
\[
\begin{pmatrix}
a_{11} & a_{12} & a_{13} \\
a_{21} & a_{22} & a_{23} \\
a_{31} & a_{32} & a_{33}
\end{pmatrix}
\]

\section{Tables}
Below is a sample table showing data:

\begin{table}[h]
\centering
\begin{tabular}{@{}lrrr@{}}
\toprule
\textbf{Category} & \textbf{Value 1} & \textbf{Value 2} & \textbf{Total} \\
\midrule
Item A & 10.5 & 20.3 & 30.8 \\
Item B & 15.2 & 25.1 & 40.3 \\
Item C & 22.1 & 12.7 & 34.8 \\
\bottomrule
\end{tabular}
\caption{Sample Data Table}
\label{tab:sample}
\end{table}

\section{Lists}
\subsection{Bullet Points}
\begin{itemize}
  \item First item
  \item Second item
  \item Third item with sub-items:
    \begin{itemize}
      \item Sub-item 1
      \item Sub-item 2
    \end{itemize}
\end{itemize}

\subsection{Numbered List}
\begin{enumerate}
  \item First step
  \item Second step
  \item Third step
\end{enumerate}

\section{Citations and References}
When you use the citation search feature, references will be automatically managed. For example, you might cite a paper like this: \cite{example2024}. All citations will be automatically added to your \texttt{cite.bib} file.

\section{Conclusion}
This template demonstrates basic \LaTeX{} features including mathematical formulas, tables, and lists. The bibliography system is now set up to work seamlessly with Bibby AI's citation management. Feel free to modify it for your needs!

\bibliographystyle{plain}
\bibliography{cite}

\end{document}\documentclass{article}
\usepackage[utf8]{inputenc}
\usepackage{booktabs}
\usepackage{amsmath}
\usepackage{natbib}

\title{My Document}
\author{Your Name}
\date{\today}

\begin{document}

\maketitle

\section{Introduction}
This is a sample document demonstrating various \LaTeX{} features including mathematical equations, tables, and formatting. You can cite papers using the citation search feature, and they will be automatically added to your \texttt{cite.bib} file.

\section{Mathematical Expressions}
Here are some example equations:

\subsection{Quadratic Formula}
The quadratic formula is given by:
\[x = \frac{-b \pm \sqrt{b^2 - 4ac}}{2a}\]

\subsection{Matrix Example}
A sample matrix:
\[
\begin{pmatrix}
a_{11} & a_{12} & a_{13} \\
a_{21} & a_{22} & a_{23} \\
a_{31} & a_{32} & a_{33}
\end{pmatrix}
\]

\section{Tables}
Below is a sample table showing data:

\begin{table}[h]
\centering
\begin{tabular}{@{}lrrr@{}}
\toprule
\textbf{Category} & \textbf{Value 1} & \textbf{Value 2} & \textbf{Total} \\
\midrule
Item A & 10.5 & 20.3 & 30.8 \\
Item B & 15.2 & 25.1 & 40.3 \\
Item C & 22.1 & 12.7 & 34.8 \\
\bottomrule
\end{tabular}
\caption{Sample Data Table}
\label{tab:sample}
\end{table}

\section{Lists}
\subsection{Bullet Points}
\begin{itemize}
  \item First item
  \item Second item
  \item Third item with sub-items:
    \begin{itemize}
      \item Sub-item 1
      \item Sub-item 2
    \end{itemize}
\end{itemize}

\subsection{Numbered List}
\begin{enumerate}
  \item First step
  \item Second step
  \item Third step
\end{enumerate}

\section{Citations and References}
When you use the citation search feature, references will be automatically managed. For example, you might cite a paper like this: \cite{example2024}. All citations will be automatically added to your \texttt{cite.bib} file.

\section{Conclusion}
This template demonstrates basic \LaTeX{} features including mathematical formulas, tables, and lists. The bibliography system is now set up to work seamlessly with Bibby AI's citation management. Feel free to modify it for your needs!

\bibliographystyle{plain}
\bibliography{cite}

\end{document}\documentclass{article}
\usepackage[utf8]{inputenc}
\usepackage{booktabs}
\usepackage{amsmath}
\usepackage{natbib}

\title{My Document}
\author{Your Name}
\date{\today}

\begin{document}

\maketitle

\section{Introduction}
This is a sample document demonstrating various \LaTeX{} features including mathematical equations, tables, and formatting. You can cite papers using the citation search feature, and they will be automatically added to your \texttt{cite.bib} file.

\section{Mathematical Expressions}
Here are some example equations:

\subsection{Quadratic Formula}
The quadratic formula is given by:
\[x = \frac{-b \pm \sqrt{b^2 - 4ac}}{2a}\]

\subsection{Matrix Example}
A sample matrix:
\[
\begin{pmatrix}
a_{11} & a_{12} & a_{13} \\
a_{21} & a_{22} & a_{23} \\
a_{31} & a_{32} & a_{33}
\end{pmatrix}
\]

\section{Tables}
Below is a sample table showing data:

\begin{table}[h]
\centering
\begin{tabular}{@{}lrrr@{}}
\toprule
\textbf{Category} & \textbf{Value 1} & \textbf{Value 2} & \textbf{Total} \\
\midrule
Item A & 10.5 & 20.3 & 30.8 \\
Item B & 15.2 & 25.1 & 40.3 \\
Item C & 22.1 & 12.7 & 34.8 \\
\bottomrule
\end{tabular}
\caption{Sample Data Table}
\label{tab:sample}
\end{table}

\section{Lists}
\subsection{Bullet Points}
\begin{itemize}
  \item First item
  \item Second item
  \item Third item with sub-items:
    \begin{itemize}
      \item Sub-item 1
      \item Sub-item 2
    \end{itemize}
\end{itemize}

\subsection{Numbered List}
\begin{enumerate}
  \item First step
  \item Second step
  \item Third step
\end{enumerate}

\section{Citations and References}
When you use the citation search feature, references will be automatically managed. For example, you might cite a paper like this: \cite{example2024}. All citations will be automatically added to your \texttt{cite.bib} file.

\section{Conclusion}
This template demonstrates basic \LaTeX{} features including mathematical formulas, tables, and lists. The bibliography system is now set up to work seamlessly with Bibby AI's citation management. Feel free to modify it for your needs!

\bibliographystyle{plain}
\bibliography{cite}

\end{document}\documentclass{article}
\usepackage[utf8]{inputenc}
\usepackage{booktabs}
\usepackage{amsmath}
\usepackage{natbib}

\title{My Document}
\author{Your Name}
\date{\today}

\begin{document}

\maketitle

\section{Introduction}
This is a sample document demonstrating various \LaTeX{} features including mathematical equations, tables, and formatting. You can cite papers using the citation search feature, and they will be automatically added to your \texttt{cite.bib} file.

\section{Mathematical Expressions}
Here are some example equations:

\subsection{Quadratic Formula}
The quadratic formula is given by:
\[x = \frac{-b \pm \sqrt{b^2 - 4ac}}{2a}\]

\subsection{Matrix Example}
A sample matrix:
\[
\begin{pmatrix}
a_{11} & a_{12} & a_{13} \\
a_{21} & a_{22} & a_{23} \\
a_{31} & a_{32} & a_{33}
\end{pmatrix}
\]

\section{Tables}
Below is a sample table showing data:

\begin{table}[h]
\centering
\begin{tabular}{@{}lrrr@{}}
\toprule
\textbf{Category} & \textbf{Value 1} & \textbf{Value 2} & \textbf{Total} \\
\midrule
Item A & 10.5 & 20.3 & 30.8 \\
Item B & 15.2 & 25.1 & 40.3 \\
Item C & 22.1 & 12.7 & 34.8 \\
\bottomrule
\end{tabular}
\caption{Sample Data Table}
\label{tab:sample}
\end{table}

\section{Lists}
\subsection{Bullet Points}
\begin{itemize}
  \item First item
  \item Second item
  \item Third item with sub-items:
    \begin{itemize}
      \item Sub-item 1
      \item Sub-item 2
    \end{itemize}
\end{itemize}

\subsection{Numbered List}
\begin{enumerate}
  \item First step
  \item Second step
  \item Third step
\end{enumerate}

\section{Citations and References}
When you use the citation search feature, references will be automatically managed. For example, you might cite a paper like this: \cite{example2024}. All citations will be automatically added to your \texttt{cite.bib} file.

\section{Conclusion}
This template demonstrates basic \LaTeX{} features including mathematical formulas, tables, and lists. The bibliography system is now set up to work seamlessly with Bibby AI's citation management. Feel free to modify it for your needs!

\bibliographystyle{plain}
\bibliography{cite}

\end{document}\documentclass{article}
\usepackage[utf8]{inputenc}
\usepackage{booktabs}
\usepackage{amsmath}
\usepackage{natbib}

\title{My Document}
\author{Your Name}
\date{\today}

\begin{document}

\maketitle

\section{Introduction}
This is a sample document demonstrating various \LaTeX{} features including mathematical equations, tables, and formatting. You can cite papers using the citation search feature, and they will be automatically added to your \texttt{cite.bib} file.

\section{Mathematical Expressions}
Here are some example equations:

\subsection{Quadratic Formula}
The quadratic formula is given by:
\[x = \frac{-b \pm \sqrt{b^2 - 4ac}}{2a}\]

\subsection{Matrix Example}
A sample matrix:
\[
\begin{pmatrix}
a_{11} & a_{12} & a_{13} \\
a_{21} & a_{22} & a_{23} \\
a_{31} & a_{32} & a_{33}
\end{pmatrix}
\]

\section{Tables}
Below is a sample table showing data:

\begin{table}[h]
\centering
\begin{tabular}{@{}lrrr@{}}
\toprule
\textbf{Category} & \textbf{Value 1} & \textbf{Value 2} & \textbf{Total} \\
\midrule
Item A & 10.5 & 20.3 & 30.8 \\
Item B & 15.2 & 25.1 & 40.3 \\
Item C & 22.1 & 12.7 & 34.8 \\
\bottomrule
\end{tabular}
\caption{Sample Data Table}
\label{tab:sample}
\end{table}

\section{Lists}
\subsection{Bullet Points}
\begin{itemize}
  \item First item
  \item Second item
  \item Third item with sub-items:
    \begin{itemize}
      \item Sub-item 1
      \item Sub-item 2
    \end{itemize}
\end{itemize}

\subsection{Numbered List}
\begin{enumerate}
  \item First step
  \item Second step
  \item Third step
\end{enumerate}

\section{Citations and References}
When you use the citation search feature, references will be automatically managed. For example, you might cite a paper like this: \cite{example2024}. All citations will be automatically added to your \texttt{cite.bib} file.

\section{Conclusion}
This template demonstrates basic \LaTeX{} features including mathematical formulas, tables, and lists. The bibliography system is now set up to work seamlessly with Bibby AI's citation management. Feel free to modify it for your needs!

\bibliographystyle{plain}
\bibliography{cite}

\end{document}\documentclass{article}
\usepackage[utf8]{inputenc}
\usepackage{booktabs}
\usepackage{amsmath}
\usepackage{natbib}

\title{My Document}
\author{Your Name}
\date{\today}

\begin{document}

\maketitle

\section{Introduction}
This is a sample document demonstrating various \LaTeX{} features including mathematical equations, tables, and formatting. You can cite papers using the citation search feature, and they will be automatically added to your \texttt{cite.bib} file.

\section{Mathematical Expressions}
Here are some example equations:

\subsection{Quadratic Formula}
The quadratic formula is given by:
\[x = \frac{-b \pm \sqrt{b^2 - 4ac}}{2a}\]

\subsection{Matrix Example}
A sample matrix:
\[
\begin{pmatrix}
a_{11} & a_{12} & a_{13} \\
a_{21} & a_{22} & a_{23} \\
a_{31} & a_{32} & a_{33}
\end{pmatrix}
\]

\section{Tables}
Below is a sample table showing data:

\begin{table}[h]
\centering
\begin{tabular}{@{}lrrr@{}}
\toprule
\textbf{Category} & \textbf{Value 1} & \textbf{Value 2} & \textbf{Total} \\
\midrule
Item A & 10.5 & 20.3 & 30.8 \\
Item B & 15.2 & 25.1 & 40.3 \\
Item C & 22.1 & 12.7 & 34.8 \\
\bottomrule
\end{tabular}
\caption{Sample Data Table}
\label{tab:sample}
\end{table}

\section{Lists}
\subsection{Bullet Points}
\begin{itemize}
  \item First item
  \item Second item
  \item Third item with sub-items:
    \begin{itemize}
      \item Sub-item 1
      \item Sub-item 2
    \end{itemize}
\end{itemize}

\subsection{Numbered List}
\begin{enumerate}
  \item First step
  \item Second step
  \item Third step
\end{enumerate}

\section{Citations and References}
When you use the citation search feature, references will be automatically managed. For example, you might cite a paper like this: \cite{example2024}. All citations will be automatically added to your \texttt{cite.bib} file.

\section{Conclusion}
This template demonstrates basic \LaTeX{} features including mathematical formulas, tables, and lists. The bibliography system is now set up to work seamlessly with Bibby AI's citation management. Feel free to modify it for your needs!

\bibliographystyle{plain}
\bibliography{cite}

\end{document}\documentclass{article}
\usepackage[utf8]{inputenc}
\usepackage{booktabs}
\usepackage{amsmath}
\usepackage{natbib}

\title{My Document}
\author{Your Name}
\date{\today}

\begin{document}

\maketitle

\section{Introduction}
This is a sample document demonstrating various \LaTeX{} features including mathematical equations, tables, and formatting. You can cite papers using the citation search feature, and they will be automatically added to your \texttt{cite.bib} file.

\section{Mathematical Expressions}
Here are some example equations:

\subsection{Quadratic Formula}
The quadratic formula is given by:
\[x = \frac{-b \pm \sqrt{b^2 - 4ac}}{2a}\]

\subsection{Matrix Example}
A sample matrix:
\[
\begin{pmatrix}
a_{11} & a_{12} & a_{13} \\
a_{21} & a_{22} & a_{23} \\
a_{31} & a_{32} & a_{33}
\end{pmatrix}
\]

\section{Tables}
Below is a sample table showing data:

\begin{table}[h]
\centering
\begin{tabular}{@{}lrrr@{}}
\toprule
\textbf{Category} & \textbf{Value 1} & \textbf{Value 2} & \textbf{Total} \\
\midrule
Item A & 10.5 & 20.3 & 30.8 \\
Item B & 15.2 & 25.1 & 40.3 \\
Item C & 22.1 & 12.7 & 34.8 \\
\bottomrule
\end{tabular}
\caption{Sample Data Table}
\label{tab:sample}
\end{table}

\section{Lists}
\subsection{Bullet Points}
\begin{itemize}
  \item First item
  \item Second item
  \item Third item with sub-items:
    \begin{itemize}
      \item Sub-item 1
      \item Sub-item 2
    \end{itemize}
\end{itemize}

\subsection{Numbered List}
\begin{enumerate}
  \item First step
  \item Second step
  \item Third step
\end{enumerate}

\section{Citations and References}
When you use the citation search feature, references will be automatically managed. For example, you might cite a paper like this: \cite{example2024}. All citations will be automatically added to your \texttt{cite.bib} file.

\section{Conclusion}
This template demonstrates basic \LaTeX{} features including mathematical formulas, tables, and lists. The bibliography system is now set up to work seamlessly with Bibby AI's citation management. Feel free to modify it for your needs!

\bibliographystyle{plain}
\bibliography{cite}

\end{document}\documentclass{article}
\usepackage[utf8]{inputenc}
\usepackage{booktabs}
\usepackage{amsmath}
\usepackage{natbib}

\title{My Document}
\author{Your Name}
\date{\today}

\begin{document}

\maketitle

\section{Introduction}
This is a sample document demonstrating various \LaTeX{} features including mathematical equations, tables, and formatting. You can cite papers using the citation search feature, and they will be automatically added to your \texttt{cite.bib} file.

\section{Mathematical Expressions}
Here are some example equations:

\subsection{Quadratic Formula}
The quadratic formula is given by:
\[x = \frac{-b \pm \sqrt{b^2 - 4ac}}{2a}\]

\subsection{Matrix Example}
A sample matrix:
\[
\begin{pmatrix}
a_{11} & a_{12} & a_{13} \\
a_{21} & a_{22} & a_{23} \\
a_{31} & a_{32} & a_{33}
\end{pmatrix}
\]

\section{Tables}
Below is a sample table showing data:

\begin{table}[h]
\centering
\begin{tabular}{@{}lrrr@{}}
\toprule
\textbf{Category} & \textbf{Value 1} & \textbf{Value 2} & \textbf{Total} \\
\midrule
Item A & 10.5 & 20.3 & 30.8 \\
Item B & 15.2 & 25.1 & 40.3 \\
Item C & 22.1 & 12.7 & 34.8 \\
\bottomrule
\end{tabular}
\caption{Sample Data Table}
\label{tab:sample}
\end{table}

\section{Lists}
\subsection{Bullet Points}
\begin{itemize}
  \item First item
  \item Second item
  \item Third item with sub-items:
    \begin{itemize}
      \item Sub-item 1
      \item Sub-item 2
    \end{itemize}
\end{itemize}

\subsection{Numbered List}
\begin{enumerate}
  \item First step
  \item Second step
  \item Third step
\end{enumerate}

\section{Citations and References}
When you use the citation search feature, references will be automatically managed. For example, you might cite a paper like this: \cite{example2024}. All citations will be automatically added to your \texttt{cite.bib} file.

\section{Conclusion}
This template demonstrates basic \LaTeX{} features including mathematical formulas, tables, and lists. The bibliography system is now set up to work seamlessly with Bibby AI's citation management. Feel free to modify it for your needs!

\bibliographystyle{plain}
\bibliography{cite}

\end{document}\documentclass{article}
\usepackage[utf8]{inputenc}
\usepackage{booktabs}
\usepackage{amsmath}
\usepackage{natbib}

\title{My Document}
\author{Your Name}
\date{\today}

\begin{document}

\maketitle

\section{Introduction}
This is a sample document demonstrating various \LaTeX{} features including mathematical equations, tables, and formatting. You can cite papers using the citation search feature, and they will be automatically added to your \texttt{cite.bib} file.

\section{Mathematical Expressions}
Here are some example equations:

\subsection{Quadratic Formula}
The quadratic formula is given by:
\[x = \frac{-b \pm \sqrt{b^2 - 4ac}}{2a}\]

\subsection{Matrix Example}
A sample matrix:
\[
\begin{pmatrix}
a_{11} & a_{12} & a_{13} \\
a_{21} & a_{22} & a_{23} \\
a_{31} & a_{32} & a_{33}
\end{pmatrix}
\]

\section{Tables}
Below is a sample table showing data:

\begin{table}[h]
\centering
\begin{tabular}{@{}lrrr@{}}
\toprule
\textbf{Category} & \textbf{Value 1} & \textbf{Value 2} & \textbf{Total} \\
\midrule
Item A & 10.5 & 20.3 & 30.8 \\
Item B & 15.2 & 25.1 & 40.3 \\
Item C & 22.1 & 12.7 & 34.8 \\
\bottomrule
\end{tabular}
\caption{Sample Data Table}
\label{tab:sample}
\end{table}

\section{Lists}
\subsection{Bullet Points}
\begin{itemize}
  \item First item
  \item Second item
  \item Third item with sub-items:
    \begin{itemize}
      \item Sub-item 1
      \item Sub-item 2
    \end{itemize}
\end{itemize}

\subsection{Numbered List}
\begin{enumerate}
  \item First step
  \item Second step
  \item Third step
\end{enumerate}

\section{Citations and References}
When you use the citation search feature, references will be automatically managed. For example, you might cite a paper like this: \cite{example2024}. All citations will be automatically added to your \texttt{cite.bib} file.

\section{Conclusion}
This template demonstrates basic \LaTeX{} features including mathematical formulas, tables, and lists. The bibliography system is now set up to work seamlessly with Bibby AI's citation management. Feel free to modify it for your needs!

\bibliographystyle{plain}
\bibliography{cite}

\end{document}\documentclass{article}
\usepackage[utf8]{inputenc}
\usepackage{booktabs}
\usepackage{amsmath}
\usepackage{natbib}

\title{My Document}
\author{Your Name}
\date{\today}

\begin{document}

\maketitle

\section{Introduction}
This is a sample document demonstrating various \LaTeX{} features including mathematical equations, tables, and formatting. You can cite papers using the citation search feature, and they will be automatically added to your \texttt{cite.bib} file.

\section{Mathematical Expressions}
Here are some example equations:

\subsection{Quadratic Formula}
The quadratic formula is given by:
\[x = \frac{-b \pm \sqrt{b^2 - 4ac}}{2a}\]

\subsection{Matrix Example}
A sample matrix:
\[
\begin{pmatrix}
a_{11} & a_{12} & a_{13} \\
a_{21} & a_{22} & a_{23} \\
a_{31} & a_{32} & a_{33}
\end{pmatrix}
\]

\section{Tables}
Below is a sample table showing data:

\begin{table}[h]
\centering
\begin{tabular}{@{}lrrr@{}}
\toprule
\textbf{Category} & \textbf{Value 1} & \textbf{Value 2} & \textbf{Total} \\
\midrule
Item A & 10.5 & 20.3 & 30.8 \\
Item B & 15.2 & 25.1 & 40.3 \\
Item C & 22.1 & 12.7 & 34.8 \\
\bottomrule
\end{tabular}
\caption{Sample Data Table}
\label{tab:sample}
\end{table}

\section{Lists}
\subsection{Bullet Points}
\begin{itemize}
  \item First item
  \item Second item
  \item Third item with sub-items:
    \begin{itemize}
      \item Sub-item 1
      \item Sub-item 2
    \end{itemize}
\end{itemize}

\subsection{Numbered List}
\begin{enumerate}
  \item First step
  \item Second step
  \item Third step
\end{enumerate}

\section{Citations and References}
When you use the citation search feature, references will be automatically managed. For example, you might cite a paper like this: \cite{example2024}. All citations will be automatically added to your \texttt{cite.bib} file.

\section{Conclusion}
This template demonstrates basic \LaTeX{} features including mathematical formulas, tables, and lists. The bibliography system is now set up to work seamlessly with Bibby AI's citation management. Feel free to modify it for your needs!

\bibliographystyle{plain}
\bibliography{cite}

\end{document}\documentclass{article}
\usepackage[utf8]{inputenc}
\usepackage{booktabs}
\usepackage{amsmath}
\usepackage{natbib}

\title{My Document}
\author{Your Name}
\date{\today}

\begin{document}

\maketitle

\section{Introduction}
This is a sample document demonstrating various \LaTeX{} features including mathematical equations, tables, and formatting. You can cite papers using the citation search feature, and they will be automatically added to your \texttt{cite.bib} file.

\section{Mathematical Expressions}
Here are some example equations:

\subsection{Quadratic Formula}
The quadratic formula is given by:
\[x = \frac{-b \pm \sqrt{b^2 - 4ac}}{2a}\]

\subsection{Matrix Example}
A sample matrix:
\[
\begin{pmatrix}
a_{11} & a_{12} & a_{13} \\
a_{21} & a_{22} & a_{23} \\
a_{31} & a_{32} & a_{33}
\end{pmatrix}
\]

\section{Tables}
Below is a sample table showing data:

\begin{table}[h]
\centering
\begin{tabular}{@{}lrrr@{}}
\toprule
\textbf{Category} & \textbf{Value 1} & \textbf{Value 2} & \textbf{Total} \\
\midrule
Item A & 10.5 & 20.3 & 30.8 \\
Item B & 15.2 & 25.1 & 40.3 \\
Item C & 22.1 & 12.7 & 34.8 \\
\bottomrule
\end{tabular}
\caption{Sample Data Table}
\label{tab:sample}
\end{table}

\section{Lists}
\subsection{Bullet Points}
\begin{itemize}
  \item First item
  \item Second item
  \item Third item with sub-items:
    \begin{itemize}
      \item Sub-item 1
      \item Sub-item 2
    \end{itemize}
\end{itemize}

\subsection{Numbered List}
\begin{enumerate}
  \item First step
  \item Second step
  \item Third step
\end{enumerate}

\section{Citations and References}
When you use the citation search feature, references will be automatically managed. For example, you might cite a paper like this: \cite{example2024}. All citations will be automatically added to your \texttt{cite.bib} file.

\section{Conclusion}
This template demonstrates basic \LaTeX{} features including mathematical formulas, tables, and lists. The bibliography system is now set up to work seamlessly with Bibby AI's citation management. Feel free to modify it for your needs!

\bibliographystyle{plain}
\bibliography{cite}

\end{document}\documentclass{article}
\usepackage[utf8]{inputenc}
\usepackage{booktabs}
\usepackage{amsmath}
\usepackage{natbib}

\title{My Document}
\author{Your Name}
\date{\today}

\begin{document}

\maketitle

\section{Introduction}
This is a sample document demonstrating various \LaTeX{} features including mathematical equations, tables, and formatting. You can cite papers using the citation search feature, and they will be automatically added to your \texttt{cite.bib} file.

\section{Mathematical Expressions}
Here are some example equations:

\subsection{Quadratic Formula}
The quadratic formula is given by:
\[x = \frac{-b \pm \sqrt{b^2 - 4ac}}{2a}\]

\subsection{Matrix Example}
A sample matrix:
\[
\begin{pmatrix}
a_{11} & a_{12} & a_{13} \\
a_{21} & a_{22} & a_{23} \\
a_{31} & a_{32} & a_{33}
\end{pmatrix}
\]

\section{Tables}
Below is a sample table showing data:

\begin{table}[h]
\centering
\begin{tabular}{@{}lrrr@{}}
\toprule
\textbf{Category} & \textbf{Value 1} & \textbf{Value 2} & \textbf{Total} \\
\midrule
Item A & 10.5 & 20.3 & 30.8 \\
Item B & 15.2 & 25.1 & 40.3 \\
Item C & 22.1 & 12.7 & 34.8 \\
\bottomrule
\end{tabular}
\caption{Sample Data Table}
\label{tab:sample}
\end{table}

\section{Lists}
\subsection{Bullet Points}
\begin{itemize}
  \item First item
  \item Second item
  \item Third item with sub-items:
    \begin{itemize}
      \item Sub-item 1
      \item Sub-item 2
    \end{itemize}
\end{itemize}

\subsection{Numbered List}
\begin{enumerate}
  \item First step
  \item Second step
  \item Third step
\end{enumerate}

\section{Citations and References}
When you use the citation search feature, references will be automatically managed. For example, you might cite a paper like this: \cite{example2024}. All citations will be automatically added to your \texttt{cite.bib} file.

\section{Conclusion}
This template demonstrates basic \LaTeX{} features including mathematical formulas, tables, and lists. The bibliography system is now set up to work seamlessly with Bibby AI's citation management. Feel free to modify it for your needs!

\bibliographystyle{plain}
\bibliography{cite}

\end{document}\documentclass{article}
\usepackage[utf8]{inputenc}
\usepackage{booktabs}
\usepackage{amsmath}
\usepackage{natbib}

\title{My Document}
\author{Your Name}
\date{\today}

\begin{document}

\maketitle

\section{Introduction}
This is a sample document demonstrating various \LaTeX{} features including mathematical equations, tables, and formatting. You can cite papers using the citation search feature, and they will be automatically added to your \texttt{cite.bib} file.

\section{Mathematical Expressions}
Here are some example equations:

\subsection{Quadratic Formula}
The quadratic formula is given by:
\[x = \frac{-b \pm \sqrt{b^2 - 4ac}}{2a}\]

\subsection{Matrix Example}
A sample matrix:
\[
\begin{pmatrix}
a_{11} & a_{12} & a_{13} \\
a_{21} & a_{22} & a_{23} \\
a_{31} & a_{32} & a_{33}
\end{pmatrix}
\]

\section{Tables}
Below is a sample table showing data:

\begin{table}[h]
\centering
\begin{tabular}{@{}lrrr@{}}
\toprule
\textbf{Category} & \textbf{Value 1} & \textbf{Value 2} & \textbf{Total} \\
\midrule
Item A & 10.5 & 20.3 & 30.8 \\
Item B & 15.2 & 25.1 & 40.3 \\
Item C & 22.1 & 12.7 & 34.8 \\
\bottomrule
\end{tabular}
\caption{Sample Data Table}
\label{tab:sample}
\end{table}

\section{Lists}
\subsection{Bullet Points}
\begin{itemize}
  \item First item
  \item Second item
  \item Third item with sub-items:
    \begin{itemize}
      \item Sub-item 1
      \item Sub-item 2
    \end{itemize}
\end{itemize}

\subsection{Numbered List}
\begin{enumerate}
  \item First step
  \item Second step
  \item Third step
\end{enumerate}

\section{Citations and References}
When you use the citation search feature, references will be automatically managed. For example, you might cite a paper like this: \cite{example2024}. All citations will be automatically added to your \texttt{cite.bib} file.

\section{Conclusion}
This template demonstrates basic \LaTeX{} features including mathematical formulas, tables, and lists. The bibliography system is now set up to work seamlessly with Bibby AI's citation management. Feel free to modify it for your needs!

\bibliographystyle{plain}
\bibliography{cite}

\end{document}\documentclass{article}
\usepackage[utf8]{inputenc}
\usepackage{booktabs}
\usepackage{amsmath}
\usepackage{natbib}

\title{My Document}
\author{Your Name}
\date{\today}

\begin{document}

\maketitle

\section{Introduction}
This is a sample document demonstrating various \LaTeX{} features including mathematical equations, tables, and formatting. You can cite papers using the citation search feature, and they will be automatically added to your \texttt{cite.bib} file.

\section{Mathematical Expressions}
Here are some example equations:

\subsection{Quadratic Formula}
The quadratic formula is given by:
\[x = \frac{-b \pm \sqrt{b^2 - 4ac}}{2a}\]

\subsection{Matrix Example}
A sample matrix:
\[
\begin{pmatrix}
a_{11} & a_{12} & a_{13} \\
a_{21} & a_{22} & a_{23} \\
a_{31} & a_{32} & a_{33}
\end{pmatrix}
\]

\section{Tables}
Below is a sample table showing data:

\begin{table}[h]
\centering
\begin{tabular}{@{}lrrr@{}}
\toprule
\textbf{Category} & \textbf{Value 1} & \textbf{Value 2} & \textbf{Total} \\
\midrule
Item A & 10.5 & 20.3 & 30.8 \\
Item B & 15.2 & 25.1 & 40.3 \\
Item C & 22.1 & 12.7 & 34.8 \\
\bottomrule
\end{tabular}
\caption{Sample Data Table}
\label{tab:sample}
\end{table}

\section{Lists}
\subsection{Bullet Points}
\begin{itemize}
  \item First item
  \item Second item
  \item Third item with sub-items:
    \begin{itemize}
      \item Sub-item 1
      \item Sub-item 2
    \end{itemize}
\end{itemize}

\subsection{Numbered List}
\begin{enumerate}
  \item First step
  \item Second step
  \item Third step
\end{enumerate}

\section{Citations and References}
When you use the citation search feature, references will be automatically managed. For example, you might cite a paper like this: \cite{example2024}. All citations will be automatically added to your \texttt{cite.bib} file.

\section{Conclusion}
This template demonstrates basic \LaTeX{} features including mathematical formulas, tables, and lists. The bibliography system is now set up to work seamlessly with Bibby AI's citation management. Feel free to modify it for your needs!

\bibliographystyle{plain}
\bibliography{cite}

\end{document}\documentclass{article}
\usepackage[utf8]{inputenc}
\usepackage{booktabs}
\usepackage{amsmath}
\usepackage{natbib}

\title{My Document}
\author{Your Name}
\date{\today}

\begin{document}

\maketitle

\section{Introduction}
This is a sample document demonstrating various \LaTeX{} features including mathematical equations, tables, and formatting. You can cite papers using the citation search feature, and they will be automatically added to your \texttt{cite.bib} file.

\section{Mathematical Expressions}
Here are some example equations:

\subsection{Quadratic Formula}
The quadratic formula is given by:
\[x = \frac{-b \pm \sqrt{b^2 - 4ac}}{2a}\]

\subsection{Matrix Example}
A sample matrix:
\[
\begin{pmatrix}
a_{11} & a_{12} & a_{13} \\
a_{21} & a_{22} & a_{23} \\
a_{31} & a_{32} & a_{33}
\end{pmatrix}
\]

\section{Tables}
Below is a sample table showing data:

\begin{table}[h]
\centering
\begin{tabular}{@{}lrrr@{}}
\toprule
\textbf{Category} & \textbf{Value 1} & \textbf{Value 2} & \textbf{Total} \\
\midrule
Item A & 10.5 & 20.3 & 30.8 \\
Item B & 15.2 & 25.1 & 40.3 \\
Item C & 22.1 & 12.7 & 34.8 \\
\bottomrule
\end{tabular}
\caption{Sample Data Table}
\label{tab:sample}
\end{table}

\section{Lists}
\subsection{Bullet Points}
\begin{itemize}
  \item First item
  \item Second item
  \item Third item with sub-items:
    \begin{itemize}
      \item Sub-item 1
      \item Sub-item 2
    \end{itemize}
\end{itemize}

\subsection{Numbered List}
\begin{enumerate}
  \item First step
  \item Second step
  \item Third step
\end{enumerate}

\section{Citations and References}
When you use the citation search feature, references will be automatically managed. For example, you might cite a paper like this: \cite{example2024}. All citations will be automatically added to your \texttt{cite.bib} file.

\section{Conclusion}
This template demonstrates basic \LaTeX{} features including mathematical formulas, tables, and lists. The bibliography system is now set up to work seamlessly with Bibby AI's citation management. Feel free to modify it for your needs!

\bibliographystyle{plain}
\bibliography{cite}

\end{document}\documentclass{article}
\usepackage[utf8]{inputenc}
\usepackage{booktabs}
\usepackage{amsmath}
\usepackage{natbib}

\title{My Document}
\author{Your Name}
\date{\today}

\begin{document}

\maketitle

\section{Introduction}
This is a sample document demonstrating various \LaTeX{} features including mathematical equations, tables, and formatting. You can cite papers using the citation search feature, and they will be automatically added to your \texttt{cite.bib} file.

\section{Mathematical Expressions}
Here are some example equations:

\subsection{Quadratic Formula}
The quadratic formula is given by:
\[x = \frac{-b \pm \sqrt{b^2 - 4ac}}{2a}\]

\subsection{Matrix Example}
A sample matrix:
\[
\begin{pmatrix}
a_{11} & a_{12} & a_{13} \\
a_{21} & a_{22} & a_{23} \\
a_{31} & a_{32} & a_{33}
\end{pmatrix}
\]

\section{Tables}
Below is a sample table showing data:

\begin{table}[h]
\centering
\begin{tabular}{@{}lrrr@{}}
\toprule
\textbf{Category} & \textbf{Value 1} & \textbf{Value 2} & \textbf{Total} \\
\midrule
Item A & 10.5 & 20.3 & 30.8 \\
Item B & 15.2 & 25.1 & 40.3 \\
Item C & 22.1 & 12.7 & 34.8 \\
\bottomrule
\end{tabular}
\caption{Sample Data Table}
\label{tab:sample}
\end{table}

\section{Lists}
\subsection{Bullet Points}
\begin{itemize}
  \item First item
  \item Second item
  \item Third item with sub-items:
    \begin{itemize}
      \item Sub-item 1
      \item Sub-item 2
    \end{itemize}
\end{itemize}

\subsection{Numbered List}
\begin{enumerate}
  \item First step
  \item Second step
  \item Third step
\end{enumerate}

\section{Citations and References}
When you use the citation search feature, references will be automatically managed. For example, you might cite a paper like this: \cite{example2024}. All citations will be automatically added to your \texttt{cite.bib} file.

\section{Conclusion}
This template demonstrates basic \LaTeX{} features including mathematical formulas, tables, and lists. The bibliography system is now set up to work seamlessly with Bibby AI's citation management. Feel free to modify it for your needs!

\bibliographystyle{plain}
\bibliography{cite}

\end{document}\documentclass{article}
\usepackage[utf8]{inputenc}
\usepackage{booktabs}
\usepackage{amsmath}
\usepackage{natbib}

\title{My Document}
\author{Your Name}
\date{\today}

\begin{document}

\maketitle

\section{Introduction}
This is a sample document demonstrating various \LaTeX{} features including mathematical equations, tables, and formatting. You can cite papers using the citation search feature, and they will be automatically added to your \texttt{cite.bib} file.

\section{Mathematical Expressions}
Here are some example equations:

\subsection{Quadratic Formula}
The quadratic formula is given by:
\[x = \frac{-b \pm \sqrt{b^2 - 4ac}}{2a}\]

\subsection{Matrix Example}
A sample matrix:
\[
\begin{pmatrix}
a_{11} & a_{12} & a_{13} \\
a_{21} & a_{22} & a_{23} \\
a_{31} & a_{32} & a_{33}
\end{pmatrix}
\]

\section{Tables}
Below is a sample table showing data:

\begin{table}[h]
\centering
\begin{tabular}{@{}lrrr@{}}
\toprule
\textbf{Category} & \textbf{Value 1} & \textbf{Value 2} & \textbf{Total} \\
\midrule
Item A & 10.5 & 20.3 & 30.8 \\
Item B & 15.2 & 25.1 & 40.3 \\
Item C & 22.1 & 12.7 & 34.8 \\
\bottomrule
\end{tabular}
\caption{Sample Data Table}
\label{tab:sample}
\end{table}

\section{Lists}
\subsection{Bullet Points}
\begin{itemize}
  \item First item
  \item Second item
  \item Third item with sub-items:
    \begin{itemize}
      \item Sub-item 1
      \item Sub-item 2
    \end{itemize}
\end{itemize}

\subsection{Numbered List}
\begin{enumerate}
  \item First step
  \item Second step
  \item Third step
\end{enumerate}

\section{Citations and References}
When you use the citation search feature, references will be automatically managed. For example, you might cite a paper like this: \cite{example2024}. All citations will be automatically added to your \texttt{cite.bib} file.

\section{Conclusion}
This template demonstrates basic \LaTeX{} features including mathematical formulas, tables, and lists. The bibliography system is now set up to work seamlessly with Bibby AI's citation management. Feel free to modify it for your needs!

\bibliographystyle{plain}
\bibliography{cite}

\end{document}\documentclass{article}
\usepackage[utf8]{inputenc}
\usepackage{booktabs}
\usepackage{amsmath}
\usepackage{natbib}

\title{My Document}
\author{Your Name}
\date{\today}

\begin{document}

\maketitle

\section{Introduction}
This is a sample document demonstrating various \LaTeX{} features including mathematical equations, tables, and formatting. You can cite papers using the citation search feature, and they will be automatically added to your \texttt{cite.bib} file.

\section{Mathematical Expressions}
Here are some example equations:

\subsection{Quadratic Formula}
The quadratic formula is given by:
\[x = \frac{-b \pm \sqrt{b^2 - 4ac}}{2a}\]

\subsection{Matrix Example}
A sample matrix:
\[
\begin{pmatrix}
a_{11} & a_{12} & a_{13} \\
a_{21} & a_{22} & a_{23} \\
a_{31} & a_{32} & a_{33}
\end{pmatrix}
\]

\section{Tables}
Below is a sample table showing data:

\begin{table}[h]
\centering
\begin{tabular}{@{}lrrr@{}}
\toprule
\textbf{Category} & \textbf{Value 1} & \textbf{Value 2} & \textbf{Total} \\
\midrule
Item A & 10.5 & 20.3 & 30.8 \\
Item B & 15.2 & 25.1 & 40.3 \\
Item C & 22.1 & 12.7 & 34.8 \\
\bottomrule
\end{tabular}
\caption{Sample Data Table}
\label{tab:sample}
\end{table}

\section{Lists}
\subsection{Bullet Points}
\begin{itemize}
  \item First item
  \item Second item
  \item Third item with sub-items:
    \begin{itemize}
      \item Sub-item 1
      \item Sub-item 2
    \end{itemize}
\end{itemize}

\subsection{Numbered List}
\begin{enumerate}
  \item First step
  \item Second step
  \item Third step
\end{enumerate}

\section{Citations and References}
When you use the citation search feature, references will be automatically managed. For example, you might cite a paper like this: \cite{example2024}. All citations will be automatically added to your \texttt{cite.bib} file.

\section{Conclusion}
This template demonstrates basic \LaTeX{} features including mathematical formulas, tables, and lists. The bibliography system is now set up to work seamlessly with Bibby AI's citation management. Feel free to modify it for your needs!

\bibliographystyle{plain}
\bibliography{cite}

\end{document}\documentclass{article}
\usepackage[utf8]{inputenc}
\usepackage{booktabs}
\usepackage{amsmath}
\usepackage{natbib}

\title{My Document}
\author{Your Name}
\date{\today}

\begin{document}

\maketitle

\section{Introduction}
This is a sample document demonstrating various \LaTeX{} features including mathematical equations, tables, and formatting. You can cite papers using the citation search feature, and they will be automatically added to your \texttt{cite.bib} file.

\section{Mathematical Expressions}
Here are some example equations:

\subsection{Quadratic Formula}
The quadratic formula is given by:
\[x = \frac{-b \pm \sqrt{b^2 - 4ac}}{2a}\]

\subsection{Matrix Example}
A sample matrix:
\[
\begin{pmatrix}
a_{11} & a_{12} & a_{13} \\
a_{21} & a_{22} & a_{23} \\
a_{31} & a_{32} & a_{33}
\end{pmatrix}
\]

\section{Tables}
Below is a sample table showing data:

\begin{table}[h]
\centering
\begin{tabular}{@{}lrrr@{}}
\toprule
\textbf{Category} & \textbf{Value 1} & \textbf{Value 2} & \textbf{Total} \\
\midrule
Item A & 10.5 & 20.3 & 30.8 \\
Item B & 15.2 & 25.1 & 40.3 \\
Item C & 22.1 & 12.7 & 34.8 \\
\bottomrule
\end{tabular}
\caption{Sample Data Table}
\label{tab:sample}
\end{table}

\section{Lists}
\subsection{Bullet Points}
\begin{itemize}
  \item First item
  \item Second item
  \item Third item with sub-items:
    \begin{itemize}
      \item Sub-item 1
      \item Sub-item 2
    \end{itemize}
\end{itemize}

\subsection{Numbered List}
\begin{enumerate}
  \item First step
  \item Second step
  \item Third step
\end{enumerate}

\section{Citations and References}
When you use the citation search feature, references will be automatically managed. For example, you might cite a paper like this: \cite{example2024}. All citations will be automatically added to your \texttt{cite.bib} file.

\section{Conclusion}
This template demonstrates basic \LaTeX{} features including mathematical formulas, tables, and lists. The bibliography system is now set up to work seamlessly with Bibby AI's citation management. Feel free to modify it for your needs!

\bibliographystyle{plain}
\bibliography{cite}

\end{document}\documentclass{article}
\usepackage[utf8]{inputenc}
\usepackage{booktabs}
\usepackage{amsmath}
\usepackage{natbib}

\title{My Document}
\author{Your Name}
\date{\today}

\begin{document}

\maketitle

\section{Introduction}
This is a sample document demonstrating various \LaTeX{} features including mathematical equations, tables, and formatting. You can cite papers using the citation search feature, and they will be automatically added to your \texttt{cite.bib} file.

\section{Mathematical Expressions}
Here are some example equations:

\subsection{Quadratic Formula}
The quadratic formula is given by:
\[x = \frac{-b \pm \sqrt{b^2 - 4ac}}{2a}\]

\subsection{Matrix Example}
A sample matrix:
\[
\begin{pmatrix}
a_{11} & a_{12} & a_{13} \\
a_{21} & a_{22} & a_{23} \\
a_{31} & a_{32} & a_{33}
\end{pmatrix}
\]

\section{Tables}
Below is a sample table showing data:

\begin{table}[h]
\centering
\begin{tabular}{@{}lrrr@{}}
\toprule
\textbf{Category} & \textbf{Value 1} & \textbf{Value 2} & \textbf{Total} \\
\midrule
Item A & 10.5 & 20.3 & 30.8 \\
Item B & 15.2 & 25.1 & 40.3 \\
Item C & 22.1 & 12.7 & 34.8 \\
\bottomrule
\end{tabular}
\caption{Sample Data Table}
\label{tab:sample}
\end{table}

\section{Lists}
\subsection{Bullet Points}
\begin{itemize}
  \item First item
  \item Second item
  \item Third item with sub-items:
    \begin{itemize}
      \item Sub-item 1
      \item Sub-item 2
    \end{itemize}
\end{itemize}

\subsection{Numbered List}
\begin{enumerate}
  \item First step
  \item Second step
  \item Third step
\end{enumerate}

\section{Citations and References}
When you use the citation search feature, references will be automatically managed. For example, you might cite a paper like this: \cite{example2024}. All citations will be automatically added to your \texttt{cite.bib} file.

\section{Conclusion}
This template demonstrates basic \LaTeX{} features including mathematical formulas, tables, and lists. The bibliography system is now set up to work seamlessly with Bibby AI's citation management. Feel free to modify it for your needs!

\bibliographystyle{plain}
\bibliography{cite}

\end{document}\documentclass{article}
\usepackage[utf8]{inputenc}
\usepackage{booktabs}
\usepackage{amsmath}
\usepackage{natbib}

\title{My Document}
\author{Your Name}
\date{\today}

\begin{document}

\maketitle

\section{Introduction}
This is a sample document demonstrating various \LaTeX{} features including mathematical equations, tables, and formatting. You can cite papers using the citation search feature, and they will be automatically added to your \texttt{cite.bib} file.

\section{Mathematical Expressions}
Here are some example equations:

\subsection{Quadratic Formula}
The quadratic formula is given by:
\[x = \frac{-b \pm \sqrt{b^2 - 4ac}}{2a}\]

\subsection{Matrix Example}
A sample matrix:
\[
\begin{pmatrix}
a_{11} & a_{12} & a_{13} \\
a_{21} & a_{22} & a_{23} \\
a_{31} & a_{32} & a_{33}
\end{pmatrix}
\]

\section{Tables}
Below is a sample table showing data:

\begin{table}[h]
\centering
\begin{tabular}{@{}lrrr@{}}
\toprule
\textbf{Category} & \textbf{Value 1} & \textbf{Value 2} & \textbf{Total} \\
\midrule
Item A & 10.5 & 20.3 & 30.8 \\
Item B & 15.2 & 25.1 & 40.3 \\
Item C & 22.1 & 12.7 & 34.8 \\
\bottomrule
\end{tabular}
\caption{Sample Data Table}
\label{tab:sample}
\end{table}

\section{Lists}
\subsection{Bullet Points}
\begin{itemize}
  \item First item
  \item Second item
  \item Third item with sub-items:
    \begin{itemize}
      \item Sub-item 1
      \item Sub-item 2
    \end{itemize}
\end{itemize}

\subsection{Numbered List}
\begin{enumerate}
  \item First step
  \item Second step
  \item Third step
\end{enumerate}

\section{Citations and References}
When you use the citation search feature, references will be automatically managed. For example, you might cite a paper like this: \cite{example2024}. All citations will be automatically added to your \texttt{cite.bib} file.

\section{Conclusion}
This template demonstrates basic \LaTeX{} features including mathematical formulas, tables, and lists. The bibliography system is now set up to work seamlessly with Bibby AI's citation management. Feel free to modify it for your needs!

\bibliographystyle{plain}
\bibliography{cite}

\end{document}\documentclass{article}
\usepackage[utf8]{inputenc}
\usepackage{booktabs}
\usepackage{amsmath}
\usepackage{natbib}

\title{My Document}
\author{Your Name}
\date{\today}

\begin{document}

\maketitle

\section{Introduction}
This is a sample document demonstrating various \LaTeX{} features including mathematical equations, tables, and formatting. You can cite papers using the citation search feature, and they will be automatically added to your \texttt{cite.bib} file.

\section{Mathematical Expressions}
Here are some example equations:

\subsection{Quadratic Formula}
The quadratic formula is given by:
\[x = \frac{-b \pm \sqrt{b^2 - 4ac}}{2a}\]

\subsection{Matrix Example}
A sample matrix:
\[
\begin{pmatrix}
a_{11} & a_{12} & a_{13} \\
a_{21} & a_{22} & a_{23} \\
a_{31} & a_{32} & a_{33}
\end{pmatrix}
\]

\section{Tables}
Below is a sample table showing data:

\begin{table}[h]
\centering
\begin{tabular}{@{}lrrr@{}}
\toprule
\textbf{Category} & \textbf{Value 1} & \textbf{Value 2} & \textbf{Total} \\
\midrule
Item A & 10.5 & 20.3 & 30.8 \\
Item B & 15.2 & 25.1 & 40.3 \\
Item C & 22.1 & 12.7 & 34.8 \\
\bottomrule
\end{tabular}
\caption{Sample Data Table}
\label{tab:sample}
\end{table}

\section{Lists}
\subsection{Bullet Points}
\begin{itemize}
  \item First item
  \item Second item
  \item Third item with sub-items:
    \begin{itemize}
      \item Sub-item 1
      \item Sub-item 2
    \end{itemize}
\end{itemize}

\subsection{Numbered List}
\begin{enumerate}
  \item First step
  \item Second step
  \item Third step
\end{enumerate}

\section{Citations and References}
When you use the citation search feature, references will be automatically managed. For example, you might cite a paper like this: \cite{example2024}. All citations will be automatically added to your \texttt{cite.bib} file.

\section{Conclusion}
This template demonstrates basic \LaTeX{} features including mathematical formulas, tables, and lists. The bibliography system is now set up to work seamlessly with Bibby AI's citation management. Feel free to modify it for your needs!

\bibliographystyle{plain}
\bibliography{cite}

\end{document}\documentclass{article}
\usepackage[utf8]{inputenc}
\usepackage{booktabs}
\usepackage{amsmath}
\usepackage{natbib}

\title{My Document}
\author{Your Name}
\date{\today}

\begin{document}

\maketitle

\section{Introduction}
This is a sample document demonstrating various \LaTeX{} features including mathematical equations, tables, and formatting. You can cite papers using the citation search feature, and they will be automatically added to your \texttt{cite.bib} file.

\section{Mathematical Expressions}
Here are some example equations:

\subsection{Quadratic Formula}
The quadratic formula is given by:
\[x = \frac{-b \pm \sqrt{b^2 - 4ac}}{2a}\]

\subsection{Matrix Example}
A sample matrix:
\[
\begin{pmatrix}
a_{11} & a_{12} & a_{13} \\
a_{21} & a_{22} & a_{23} \\
a_{31} & a_{32} & a_{33}
\end{pmatrix}
\]

\section{Tables}
Below is a sample table showing data:

\begin{table}[h]
\centering
\begin{tabular}{@{}lrrr@{}}
\toprule
\textbf{Category} & \textbf{Value 1} & \textbf{Value 2} & \textbf{Total} \\
\midrule
Item A & 10.5 & 20.3 & 30.8 \\
Item B & 15.2 & 25.1 & 40.3 \\
Item C & 22.1 & 12.7 & 34.8 \\
\bottomrule
\end{tabular}
\caption{Sample Data Table}
\label{tab:sample}
\end{table}

\section{Lists}
\subsection{Bullet Points}
\begin{itemize}
  \item First item
  \item Second item
  \item Third item with sub-items:
    \begin{itemize}
      \item Sub-item 1
      \item Sub-item 2
    \end{itemize}
\end{itemize}

\subsection{Numbered List}
\begin{enumerate}
  \item First step
  \item Second step
  \item Third step
\end{enumerate}

\section{Citations and References}
When you use the citation search feature, references will be automatically managed. For example, you might cite a paper like this: \cite{example2024}. All citations will be automatically added to your \texttt{cite.bib} file.

\section{Conclusion}
This template demonstrates basic \LaTeX{} features including mathematical formulas, tables, and lists. The bibliography system is now set up to work seamlessly with Bibby AI's citation management. Feel free to modify it for your needs!

\bibliographystyle{plain}
\bibliography{cite}

\end{document}\documentclass{article}
\usepackage[utf8]{inputenc}
\usepackage{booktabs}
\usepackage{amsmath}
\usepackage{natbib}

\title{My Document}
\author{Your Name}
\date{\today}

\begin{document}

\maketitle

\section{Introduction}
This is a sample document demonstrating various \LaTeX{} features including mathematical equations, tables, and formatting. You can cite papers using the citation search feature, and they will be automatically added to your \texttt{cite.bib} file.

\section{Mathematical Expressions}
Here are some example equations:

\subsection{Quadratic Formula}
The quadratic formula is given by:
\[x = \frac{-b \pm \sqrt{b^2 - 4ac}}{2a}\]

\subsection{Matrix Example}
A sample matrix:
\[
\begin{pmatrix}
a_{11} & a_{12} & a_{13} \\
a_{21} & a_{22} & a_{23} \\
a_{31} & a_{32} & a_{33}
\end{pmatrix}
\]

\section{Tables}
Below is a sample table showing data:

\begin{table}[h]
\centering
\begin{tabular}{@{}lrrr@{}}
\toprule
\textbf{Category} & \textbf{Value 1} & \textbf{Value 2} & \textbf{Total} \\
\midrule
Item A & 10.5 & 20.3 & 30.8 \\
Item B & 15.2 & 25.1 & 40.3 \\
Item C & 22.1 & 12.7 & 34.8 \\
\bottomrule
\end{tabular}
\caption{Sample Data Table}
\label{tab:sample}
\end{table}

\section{Lists}
\subsection{Bullet Points}
\begin{itemize}
  \item First item
  \item Second item
  \item Third item with sub-items:
    \begin{itemize}
      \item Sub-item 1
      \item Sub-item 2
    \end{itemize}
\end{itemize}

\subsection{Numbered List}
\begin{enumerate}
  \item First step
  \item Second step
  \item Third step
\end{enumerate}

\section{Citations and References}
When you use the citation search feature, references will be automatically managed. For example, you might cite a paper like this: \cite{example2024}. All citations will be automatically added to your \texttt{cite.bib} file.

\section{Conclusion}
This template demonstrates basic \LaTeX{} features including mathematical formulas, tables, and lists. The bibliography system is now set up to work seamlessly with Bibby AI's citation management. Feel free to modify it for your needs!

\bibliographystyle{plain}
\bibliography{cite}

\end{document}\documentclass{article}
\usepackage[utf8]{inputenc}
\usepackage{booktabs}
\usepackage{amsmath}
\usepackage{natbib}

\title{My Document}
\author{Your Name}
\date{\today}

\begin{document}

\maketitle

\section{Introduction}
This is a sample document demonstrating various \LaTeX{} features including mathematical equations, tables, and formatting. You can cite papers using the citation search feature, and they will be automatically added to your \texttt{cite.bib} file.

\section{Mathematical Expressions}
Here are some example equations:

\subsection{Quadratic Formula}
The quadratic formula is given by:
\[x = \frac{-b \pm \sqrt{b^2 - 4ac}}{2a}\]

\subsection{Matrix Example}
A sample matrix:
\[
\begin{pmatrix}
a_{11} & a_{12} & a_{13} \\
a_{21} & a_{22} & a_{23} \\
a_{31} & a_{32} & a_{33}
\end{pmatrix}
\]

\section{Tables}
Below is a sample table showing data:

\begin{table}[h]
\centering
\begin{tabular}{@{}lrrr@{}}
\toprule
\textbf{Category} & \textbf{Value 1} & \textbf{Value 2} & \textbf{Total} \\
\midrule
Item A & 10.5 & 20.3 & 30.8 \\
Item B & 15.2 & 25.1 & 40.3 \\
Item C & 22.1 & 12.7 & 34.8 \\
\bottomrule
\end{tabular}
\caption{Sample Data Table}
\label{tab:sample}
\end{table}

\section{Lists}
\subsection{Bullet Points}
\begin{itemize}
  \item First item
  \item Second item
  \item Third item with sub-items:
    \begin{itemize}
      \item Sub-item 1
      \item Sub-item 2
    \end{itemize}
\end{itemize}

\subsection{Numbered List}
\begin{enumerate}
  \item First step
  \item Second step
  \item Third step
\end{enumerate}

\section{Citations and References}
When you use the citation search feature, references will be automatically managed. For example, you might cite a paper like this: \cite{example2024}. All citations will be automatically added to your \texttt{cite.bib} file.

\section{Conclusion}
This template demonstrates basic \LaTeX{} features including mathematical formulas, tables, and lists. The bibliography system is now set up to work seamlessly with Bibby AI's citation management. Feel free to modify it for your needs!

\bibliographystyle{plain}
\bibliography{cite}

\end{document}\documentclass{article}
\usepackage[utf8]{inputenc}
\usepackage{booktabs}
\usepackage{amsmath}
\usepackage{natbib}

\title{My Document}
\author{Your Name}
\date{\today}

\begin{document}

\maketitle

\section{Introduction}
This is a sample document demonstrating various \LaTeX{} features including mathematical equations, tables, and formatting. You can cite papers using the citation search feature, and they will be automatically added to your \texttt{cite.bib} file.

\section{Mathematical Expressions}
Here are some example equations:

\subsection{Quadratic Formula}
The quadratic formula is given by:
\[x = \frac{-b \pm \sqrt{b^2 - 4ac}}{2a}\]

\subsection{Matrix Example}
A sample matrix:
\[
\begin{pmatrix}
a_{11} & a_{12} & a_{13} \\
a_{21} & a_{22} & a_{23} \\
a_{31} & a_{32} & a_{33}
\end{pmatrix}
\]

\section{Tables}
Below is a sample table showing data:

\begin{table}[h]
\centering
\begin{tabular}{@{}lrrr@{}}
\toprule
\textbf{Category} & \textbf{Value 1} & \textbf{Value 2} & \textbf{Total} \\
\midrule
Item A & 10.5 & 20.3 & 30.8 \\
Item B & 15.2 & 25.1 & 40.3 \\
Item C & 22.1 & 12.7 & 34.8 \\
\bottomrule
\end{tabular}
\caption{Sample Data Table}
\label{tab:sample}
\end{table}

\section{Lists}
\subsection{Bullet Points}
\begin{itemize}
  \item First item
  \item Second item
  \item Third item with sub-items:
    \begin{itemize}
      \item Sub-item 1
      \item Sub-item 2
    \end{itemize}
\end{itemize}

\subsection{Numbered List}
\begin{enumerate}
  \item First step
  \item Second step
  \item Third step
\end{enumerate}

\section{Citations and References}
When you use the citation search feature, references will be automatically managed. For example, you might cite a paper like this: \cite{example2024}. All citations will be automatically added to your \texttt{cite.bib} file.

\section{Conclusion}
This template demonstrates basic \LaTeX{} features including mathematical formulas, tables, and lists. The bibliography system is now set up to work seamlessly with Bibby AI's citation management. Feel free to modify it for your needs!

\bibliographystyle{plain}
\bibliography{cite}

\end{document}\documentclass{article}
\usepackage[utf8]{inputenc}
\usepackage{booktabs}
\usepackage{amsmath}
\usepackage{natbib}

\title{My Document}
\author{Your Name}
\date{\today}

\begin{document}

\maketitle

\section{Introduction}
This is a sample document demonstrating various \LaTeX{} features including mathematical equations, tables, and formatting. You can cite papers using the citation search feature, and they will be automatically added to your \texttt{cite.bib} file.

\section{Mathematical Expressions}
Here are some example equations:

\subsection{Quadratic Formula}
The quadratic formula is given by:
\[x = \frac{-b \pm \sqrt{b^2 - 4ac}}{2a}\]

\subsection{Matrix Example}
A sample matrix:
\[
\begin{pmatrix}
a_{11} & a_{12} & a_{13} \\
a_{21} & a_{22} & a_{23} \\
a_{31} & a_{32} & a_{33}
\end{pmatrix}
\]

\section{Tables}
Below is a sample table showing data:

\begin{table}[h]
\centering
\begin{tabular}{@{}lrrr@{}}
\toprule
\textbf{Category} & \textbf{Value 1} & \textbf{Value 2} & \textbf{Total} \\
\midrule
Item A & 10.5 & 20.3 & 30.8 \\
Item B & 15.2 & 25.1 & 40.3 \\
Item C & 22.1 & 12.7 & 34.8 \\
\bottomrule
\end{tabular}
\caption{Sample Data Table}
\label{tab:sample}
\end{table}

\section{Lists}
\subsection{Bullet Points}
\begin{itemize}
  \item First item
  \item Second item
  \item Third item with sub-items:
    \begin{itemize}
      \item Sub-item 1
      \item Sub-item 2
    \end{itemize}
\end{itemize}

\subsection{Numbered List}
\begin{enumerate}
  \item First step
  \item Second step
  \item Third step
\end{enumerate}

\section{Citations and References}
When you use the citation search feature, references will be automatically managed. For example, you might cite a paper like this: \cite{example2024}. All citations will be automatically added to your \texttt{cite.bib} file.

\section{Conclusion}
This template demonstrates basic \LaTeX{} features including mathematical formulas, tables, and lists. The bibliography system is now set up to work seamlessly with Bibby AI's citation management. Feel free to modify it for your needs!

\bibliographystyle{plain}
\bibliography{cite}

\end{document}\documentclass{article}
\usepackage[utf8]{inputenc}
\usepackage{booktabs}
\usepackage{amsmath}
\usepackage{natbib}

\title{My Document}
\author{Your Name}
\date{\today}

\begin{document}

\maketitle

\section{Introduction}
This is a sample document demonstrating various \LaTeX{} features including mathematical equations, tables, and formatting. You can cite papers using the citation search feature, and they will be automatically added to your \texttt{cite.bib} file.

\section{Mathematical Expressions}
Here are some example equations:

\subsection{Quadratic Formula}
The quadratic formula is given by:
\[x = \frac{-b \pm \sqrt{b^2 - 4ac}}{2a}\]

\subsection{Matrix Example}
A sample matrix:
\[
\begin{pmatrix}
a_{11} & a_{12} & a_{13} \\
a_{21} & a_{22} & a_{23} \\
a_{31} & a_{32} & a_{33}
\end{pmatrix}
\]

\section{Tables}
Below is a sample table showing data:

\begin{table}[h]
\centering
\begin{tabular}{@{}lrrr@{}}
\toprule
\textbf{Category} & \textbf{Value 1} & \textbf{Value 2} & \textbf{Total} \\
\midrule
Item A & 10.5 & 20.3 & 30.8 \\
Item B & 15.2 & 25.1 & 40.3 \\
Item C & 22.1 & 12.7 & 34.8 \\
\bottomrule
\end{tabular}
\caption{Sample Data Table}
\label{tab:sample}
\end{table}

\section{Lists}
\subsection{Bullet Points}
\begin{itemize}
  \item First item
  \item Second item
  \item Third item with sub-items:
    \begin{itemize}
      \item Sub-item 1
      \item Sub-item 2
    \end{itemize}
\end{itemize}

\subsection{Numbered List}
\begin{enumerate}
  \item First step
  \item Second step
  \item Third step
\end{enumerate}

\section{Citations and References}
When you use the citation search feature, references will be automatically managed. For example, you might cite a paper like this: \cite{example2024}. All citations will be automatically added to your \texttt{cite.bib} file.

\section{Conclusion}
This template demonstrates basic \LaTeX{} features including mathematical formulas, tables, and lists. The bibliography system is now set up to work seamlessly with Bibby AI's citation management. Feel free to modify it for your needs!

\bibliographystyle{plain}
\bibliography{cite}

\end{document}\documentclass{article}
\usepackage[utf8]{inputenc}
\usepackage{booktabs}
\usepackage{amsmath}
\usepackage{natbib}

\title{My Document}
\author{Your Name}
\date{\today}

\begin{document}

\maketitle

\section{Introduction}
This is a sample document demonstrating various \LaTeX{} features including mathematical equations, tables, and formatting. You can cite papers using the citation search feature, and they will be automatically added to your \texttt{cite.bib} file.

\section{Mathematical Expressions}
Here are some example equations:

\subsection{Quadratic Formula}
The quadratic formula is given by:
\[x = \frac{-b \pm \sqrt{b^2 - 4ac}}{2a}\]

\subsection{Matrix Example}
A sample matrix:
\[
\begin{pmatrix}
a_{11} & a_{12} & a_{13} \\
a_{21} & a_{22} & a_{23} \\
a_{31} & a_{32} & a_{33}
\end{pmatrix}
\]

\section{Tables}
Below is a sample table showing data:

\begin{table}[h]
\centering
\begin{tabular}{@{}lrrr@{}}
\toprule
\textbf{Category} & \textbf{Value 1} & \textbf{Value 2} & \textbf{Total} \\
\midrule
Item A & 10.5 & 20.3 & 30.8 \\
Item B & 15.2 & 25.1 & 40.3 \\
Item C & 22.1 & 12.7 & 34.8 \\
\bottomrule
\end{tabular}
\caption{Sample Data Table}
\label{tab:sample}
\end{table}

\section{Lists}
\subsection{Bullet Points}
\begin{itemize}
  \item First item
  \item Second item
  \item Third item with sub-items:
    \begin{itemize}
      \item Sub-item 1
      \item Sub-item 2
    \end{itemize}
\end{itemize}

\subsection{Numbered List}
\begin{enumerate}
  \item First step
  \item Second step
  \item Third step
\end{enumerate}

\section{Citations and References}
When you use the citation search feature, references will be automatically managed. For example, you might cite a paper like this: \cite{example2024}. All citations will be automatically added to your \texttt{cite.bib} file.

\section{Conclusion}
This template demonstrates basic \LaTeX{} features including mathematical formulas, tables, and lists. The bibliography system is now set up to work seamlessly with Bibby AI's citation management. Feel free to modify it for your needs!

\bibliographystyle{plain}
\bibliography{cite}

\end{document}\documentclass{article}
\usepackage[utf8]{inputenc}
\usepackage{booktabs}
\usepackage{amsmath}
\usepackage{natbib}

\title{My Document}
\author{Your Name}
\date{\today}

\begin{document}

\maketitle

\section{Introduction}
This is a sample document demonstrating various \LaTeX{} features including mathematical equations, tables, and formatting. You can cite papers using the citation search feature, and they will be automatically added to your \texttt{cite.bib} file.

\section{Mathematical Expressions}
Here are some example equations:

\subsection{Quadratic Formula}
The quadratic formula is given by:
\[x = \frac{-b \pm \sqrt{b^2 - 4ac}}{2a}\]

\subsection{Matrix Example}
A sample matrix:
\[
\begin{pmatrix}
a_{11} & a_{12} & a_{13} \\
a_{21} & a_{22} & a_{23} \\
a_{31} & a_{32} & a_{33}
\end{pmatrix}
\]

\section{Tables}
Below is a sample table showing data:

\begin{table}[h]
\centering
\begin{tabular}{@{}lrrr@{}}
\toprule
\textbf{Category} & \textbf{Value 1} & \textbf{Value 2} & \textbf{Total} \\
\midrule
Item A & 10.5 & 20.3 & 30.8 \\
Item B & 15.2 & 25.1 & 40.3 \\
Item C & 22.1 & 12.7 & 34.8 \\
\bottomrule
\end{tabular}
\caption{Sample Data Table}
\label{tab:sample}
\end{table}

\section{Lists}
\subsection{Bullet Points}
\begin{itemize}
  \item First item
  \item Second item
  \item Third item with sub-items:
    \begin{itemize}
      \item Sub-item 1
      \item Sub-item 2
    \end{itemize}
\end{itemize}

\subsection{Numbered List}
\begin{enumerate}
  \item First step
  \item Second step
  \item Third step
\end{enumerate}

\section{Citations and References}
When you use the citation search feature, references will be automatically managed. For example, you might cite a paper like this: \cite{example2024}. All citations will be automatically added to your \texttt{cite.bib} file.

\section{Conclusion}
This template demonstrates basic \LaTeX{} features including mathematical formulas, tables, and lists. The bibliography system is now set up to work seamlessly with Bibby AI's citation management. Feel free to modify it for your needs!

\bibliographystyle{plain}
\bibliography{cite}

\end{document}\documentclass{article}
\usepackage[utf8]{inputenc}
\usepackage{booktabs}
\usepackage{amsmath}
\usepackage{natbib}

\title{My Document}
\author{Your Name}
\date{\today}

\begin{document}

\maketitle

\section{Introduction}
This is a sample document demonstrating various \LaTeX{} features including mathematical equations, tables, and formatting. You can cite papers using the citation search feature, and they will be automatically added to your \texttt{cite.bib} file.

\section{Mathematical Expressions}
Here are some example equations:

\subsection{Quadratic Formula}
The quadratic formula is given by:
\[x = \frac{-b \pm \sqrt{b^2 - 4ac}}{2a}\]

\subsection{Matrix Example}
A sample matrix:
\[
\begin{pmatrix}
a_{11} & a_{12} & a_{13} \\
a_{21} & a_{22} & a_{23} \\
a_{31} & a_{32} & a_{33}
\end{pmatrix}
\]

\section{Tables}
Below is a sample table showing data:

\begin{table}[h]
\centering
\begin{tabular}{@{}lrrr@{}}
\toprule
\textbf{Category} & \textbf{Value 1} & \textbf{Value 2} & \textbf{Total} \\
\midrule
Item A & 10.5 & 20.3 & 30.8 \\
Item B & 15.2 & 25.1 & 40.3 \\
Item C & 22.1 & 12.7 & 34.8 \\
\bottomrule
\end{tabular}
\caption{Sample Data Table}
\label{tab:sample}
\end{table}

\section{Lists}
\subsection{Bullet Points}
\begin{itemize}
  \item First item
  \item Second item
  \item Third item with sub-items:
    \begin{itemize}
      \item Sub-item 1
      \item Sub-item 2
    \end{itemize}
\end{itemize}

\subsection{Numbered List}
\begin{enumerate}
  \item First step
  \item Second step
  \item Third step
\end{enumerate}

\section{Citations and References}
When you use the citation search feature, references will be automatically managed. For example, you might cite a paper like this: \cite{example2024}. All citations will be automatically added to your \texttt{cite.bib} file.

\section{Conclusion}
This template demonstrates basic \LaTeX{} features including mathematical formulas, tables, and lists. The bibliography system is now set up to work seamlessly with Bibby AI's citation management. Feel free to modify it for your needs!

\bibliographystyle{plain}
\bibliography{cite}

\end{document}\documentclass{article}
\usepackage[utf8]{inputenc}
\usepackage{booktabs}
\usepackage{amsmath}
\usepackage{natbib}

\title{My Document}
\author{Your Name}
\date{\today}

\begin{document}

\maketitle

\section{Introduction}
This is a sample document demonstrating various \LaTeX{} features including mathematical equations, tables, and formatting. You can cite papers using the citation search feature, and they will be automatically added to your \texttt{cite.bib} file.

\section{Mathematical Expressions}
Here are some example equations:

\subsection{Quadratic Formula}
The quadratic formula is given by:
\[x = \frac{-b \pm \sqrt{b^2 - 4ac}}{2a}\]

\subsection{Matrix Example}
A sample matrix:
\[
\begin{pmatrix}
a_{11} & a_{12} & a_{13} \\
a_{21} & a_{22} & a_{23} \\
a_{31} & a_{32} & a_{33}
\end{pmatrix}
\]

\section{Tables}
Below is a sample table showing data:

\begin{table}[h]
\centering
\begin{tabular}{@{}lrrr@{}}
\toprule
\textbf{Category} & \textbf{Value 1} & \textbf{Value 2} & \textbf{Total} \\
\midrule
Item A & 10.5 & 20.3 & 30.8 \\
Item B & 15.2 & 25.1 & 40.3 \\
Item C & 22.1 & 12.7 & 34.8 \\
\bottomrule
\end{tabular}
\caption{Sample Data Table}
\label{tab:sample}
\end{table}

\section{Lists}
\subsection{Bullet Points}
\begin{itemize}
  \item First item
  \item Second item
  \item Third item with sub-items:
    \begin{itemize}
      \item Sub-item 1
      \item Sub-item 2
    \end{itemize}
\end{itemize}

\subsection{Numbered List}
\begin{enumerate}
  \item First step
  \item Second step
  \item Third step
\end{enumerate}

\section{Citations and References}
When you use the citation search feature, references will be automatically managed. For example, you might cite a paper like this: \cite{example2024}. All citations will be automatically added to your \texttt{cite.bib} file.

\section{Conclusion}
This template demonstrates basic \LaTeX{} features including mathematical formulas, tables, and lists. The bibliography system is now set up to work seamlessly with Bibby AI's citation management. Feel free to modify it for your needs!

\bibliographystyle{plain}
\bibliography{cite}

\end{document}\documentclass{article}
\usepackage[utf8]{inputenc}
\usepackage{booktabs}
\usepackage{amsmath}
\usepackage{natbib}

\title{My Document}
\author{Your Name}
\date{\today}

\begin{document}

\maketitle

\section{Introduction}
This is a sample document demonstrating various \LaTeX{} features including mathematical equations, tables, and formatting. You can cite papers using the citation search feature, and they will be automatically added to your \texttt{cite.bib} file.

\section{Mathematical Expressions}
Here are some example equations:

\subsection{Quadratic Formula}
The quadratic formula is given by:
\[x = \frac{-b \pm \sqrt{b^2 - 4ac}}{2a}\]

\subsection{Matrix Example}
A sample matrix:
\[
\begin{pmatrix}
a_{11} & a_{12} & a_{13} \\
a_{21} & a_{22} & a_{23} \\
a_{31} & a_{32} & a_{33}
\end{pmatrix}
\]

\section{Tables}
Below is a sample table showing data:

\begin{table}[h]
\centering
\begin{tabular}{@{}lrrr@{}}
\toprule
\textbf{Category} & \textbf{Value 1} & \textbf{Value 2} & \textbf{Total} \\
\midrule
Item A & 10.5 & 20.3 & 30.8 \\
Item B & 15.2 & 25.1 & 40.3 \\
Item C & 22.1 & 12.7 & 34.8 \\
\bottomrule
\end{tabular}
\caption{Sample Data Table}
\label{tab:sample}
\end{table}

\section{Lists}
\subsection{Bullet Points}
\begin{itemize}
  \item First item
  \item Second item
  \item Third item with sub-items:
    \begin{itemize}
      \item Sub-item 1
      \item Sub-item 2
    \end{itemize}
\end{itemize}

\subsection{Numbered List}
\begin{enumerate}
  \item First step
  \item Second step
  \item Third step
\end{enumerate}

\section{Citations and References}
When you use the citation search feature, references will be automatically managed. For example, you might cite a paper like this: \cite{example2024}. All citations will be automatically added to your \texttt{cite.bib} file.

\section{Conclusion}
This template demonstrates basic \LaTeX{} features including mathematical formulas, tables, and lists. The bibliography system is now set up to work seamlessly with Bibby AI's citation management. Feel free to modify it for your needs!

\bibliographystyle{plain}
\bibliography{cite}

\end{document}\documentclass{article}
\usepackage[utf8]{inputenc}
\usepackage{booktabs}
\usepackage{amsmath}
\usepackage{natbib}

\title{My Document}
\author{Your Name}
\date{\today}

\begin{document}

\maketitle

\section{Introduction}
This is a sample document demonstrating various \LaTeX{} features including mathematical equations, tables, and formatting. You can cite papers using the citation search feature, and they will be automatically added to your \texttt{cite.bib} file.

\section{Mathematical Expressions}
Here are some example equations:

\subsection{Quadratic Formula}
The quadratic formula is given by:
\[x = \frac{-b \pm \sqrt{b^2 - 4ac}}{2a}\]

\subsection{Matrix Example}
A sample matrix:
\[
\begin{pmatrix}
a_{11} & a_{12} & a_{13} \\
a_{21} & a_{22} & a_{23} \\
a_{31} & a_{32} & a_{33}
\end{pmatrix}
\]

\section{Tables}
Below is a sample table showing data:

\begin{table}[h]
\centering
\begin{tabular}{@{}lrrr@{}}
\toprule
\textbf{Category} & \textbf{Value 1} & \textbf{Value 2} & \textbf{Total} \\
\midrule
Item A & 10.5 & 20.3 & 30.8 \\
Item B & 15.2 & 25.1 & 40.3 \\
Item C & 22.1 & 12.7 & 34.8 \\
\bottomrule
\end{tabular}
\caption{Sample Data Table}
\label{tab:sample}
\end{table}

\section{Lists}
\subsection{Bullet Points}
\begin{itemize}
  \item First item
  \item Second item
  \item Third item with sub-items:
    \begin{itemize}
      \item Sub-item 1
      \item Sub-item 2
    \end{itemize}
\end{itemize}

\subsection{Numbered List}
\begin{enumerate}
  \item First step
  \item Second step
  \item Third step
\end{enumerate}

\section{Citations and References}
When you use the citation search feature, references will be automatically managed. For example, you might cite a paper like this: \cite{example2024}. All citations will be automatically added to your \texttt{cite.bib} file.

\section{Conclusion}
This template demonstrates basic \LaTeX{} features including mathematical formulas, tables, and lists. The bibliography system is now set up to work seamlessly with Bibby AI's citation management. Feel free to modify it for your needs!

\bibliographystyle{plain}
\bibliography{cite}

\end{document}\documentclass{article}
\usepackage[utf8]{inputenc}
\usepackage{booktabs}
\usepackage{amsmath}
\usepackage{natbib}

\title{My Document}
\author{Your Name}
\date{\today}

\begin{document}

\maketitle

\section{Introduction}
This is a sample document demonstrating various \LaTeX{} features including mathematical equations, tables, and formatting. You can cite papers using the citation search feature, and they will be automatically added to your \texttt{cite.bib} file.

\section{Mathematical Expressions}
Here are some example equations:

\subsection{Quadratic Formula}
The quadratic formula is given by:
\[x = \frac{-b \pm \sqrt{b^2 - 4ac}}{2a}\]

\subsection{Matrix Example}
A sample matrix:
\[
\begin{pmatrix}
a_{11} & a_{12} & a_{13} \\
a_{21} & a_{22} & a_{23} \\
a_{31} & a_{32} & a_{33}
\end{pmatrix}
\]

\section{Tables}
Below is a sample table showing data:

\begin{table}[h]
\centering
\begin{tabular}{@{}lrrr@{}}
\toprule
\textbf{Category} & \textbf{Value 1} & \textbf{Value 2} & \textbf{Total} \\
\midrule
Item A & 10.5 & 20.3 & 30.8 \\
Item B & 15.2 & 25.1 & 40.3 \\
Item C & 22.1 & 12.7 & 34.8 \\
\bottomrule
\end{tabular}
\caption{Sample Data Table}
\label{tab:sample}
\end{table}

\section{Lists}
\subsection{Bullet Points}
\begin{itemize}
  \item First item
  \item Second item
  \item Third item with sub-items:
    \begin{itemize}
      \item Sub-item 1
      \item Sub-item 2
    \end{itemize}
\end{itemize}

\subsection{Numbered List}
\begin{enumerate}
  \item First step
  \item Second step
  \item Third step
\end{enumerate}

\section{Citations and References}
When you use the citation search feature, references will be automatically managed. For example, you might cite a paper like this: \cite{example2024}. All citations will be automatically added to your \texttt{cite.bib} file.

\section{Conclusion}
This template demonstrates basic \LaTeX{} features including mathematical formulas, tables, and lists. The bibliography system is now set up to work seamlessly with Bibby AI's citation management. Feel free to modify it for your needs!

\bibliographystyle{plain}
\bibliography{cite}

\end{document}\documentclass{article}
\usepackage[utf8]{inputenc}
\usepackage{booktabs}
\usepackage{amsmath}
\usepackage{natbib}

\title{My Document}
\author{Your Name}
\date{\today}

\begin{document}

\maketitle

\section{Introduction}
This is a sample document demonstrating various \LaTeX{} features including mathematical equations, tables, and formatting. You can cite papers using the citation search feature, and they will be automatically added to your \texttt{cite.bib} file.

\section{Mathematical Expressions}
Here are some example equations:

\subsection{Quadratic Formula}
The quadratic formula is given by:
\[x = \frac{-b \pm \sqrt{b^2 - 4ac}}{2a}\]

\subsection{Matrix Example}
A sample matrix:
\[
\begin{pmatrix}
a_{11} & a_{12} & a_{13} \\
a_{21} & a_{22} & a_{23} \\
a_{31} & a_{32} & a_{33}
\end{pmatrix}
\]

\section{Tables}
Below is a sample table showing data:

\begin{table}[h]
\centering
\begin{tabular}{@{}lrrr@{}}
\toprule
\textbf{Category} & \textbf{Value 1} & \textbf{Value 2} & \textbf{Total} \\
\midrule
Item A & 10.5 & 20.3 & 30.8 \\
Item B & 15.2 & 25.1 & 40.3 \\
Item C & 22.1 & 12.7 & 34.8 \\
\bottomrule
\end{tabular}
\caption{Sample Data Table}
\label{tab:sample}
\end{table}

\section{Lists}
\subsection{Bullet Points}
\begin{itemize}
  \item First item
  \item Second item
  \item Third item with sub-items:
    \begin{itemize}
      \item Sub-item 1
      \item Sub-item 2
    \end{itemize}
\end{itemize}

\subsection{Numbered List}
\begin{enumerate}
  \item First step
  \item Second step
  \item Third step
\end{enumerate}

\section{Citations and References}
When you use the citation search feature, references will be automatically managed. For example, you might cite a paper like this: \cite{example2024}. All citations will be automatically added to your \texttt{cite.bib} file.

\section{Conclusion}
This template demonstrates basic \LaTeX{} features including mathematical formulas, tables, and lists. The bibliography system is now set up to work seamlessly with Bibby AI's citation management. Feel free to modify it for your needs!

\bibliographystyle{plain}
\bibliography{cite}

\end{document}\documentclass{article}
\usepackage[utf8]{inputenc}
\usepackage{booktabs}
\usepackage{amsmath}
\usepackage{natbib}

\title{My Document}
\author{Your Name}
\date{\today}

\begin{document}

\maketitle

\section{Introduction}
This is a sample document demonstrating various \LaTeX{} features including mathematical equations, tables, and formatting. You can cite papers using the citation search feature, and they will be automatically added to your \texttt{cite.bib} file.

\section{Mathematical Expressions}
Here are some example equations:

\subsection{Quadratic Formula}
The quadratic formula is given by:
\[x = \frac{-b \pm \sqrt{b^2 - 4ac}}{2a}\]

\subsection{Matrix Example}
A sample matrix:
\[
\begin{pmatrix}
a_{11} & a_{12} & a_{13} \\
a_{21} & a_{22} & a_{23} \\
a_{31} & a_{32} & a_{33}
\end{pmatrix}
\]

\section{Tables}
Below is a sample table showing data:

\begin{table}[h]
\centering
\begin{tabular}{@{}lrrr@{}}
\toprule
\textbf{Category} & \textbf{Value 1} & \textbf{Value 2} & \textbf{Total} \\
\midrule
Item A & 10.5 & 20.3 & 30.8 \\
Item B & 15.2 & 25.1 & 40.3 \\
Item C & 22.1 & 12.7 & 34.8 \\
\bottomrule
\end{tabular}
\caption{Sample Data Table}
\label{tab:sample}
\end{table}

\section{Lists}
\subsection{Bullet Points}
\begin{itemize}
  \item First item
  \item Second item
  \item Third item with sub-items:
    \begin{itemize}
      \item Sub-item 1
      \item Sub-item 2
    \end{itemize}
\end{itemize}

\subsection{Numbered List}
\begin{enumerate}
  \item First step
  \item Second step
  \item Third step
\end{enumerate}

\section{Citations and References}
When you use the citation search feature, references will be automatically managed. For example, you might cite a paper like this: \cite{example2024}. All citations will be automatically added to your \texttt{cite.bib} file.

\section{Conclusion}
This template demonstrates basic \LaTeX{} features including mathematical formulas, tables, and lists. The bibliography system is now set up to work seamlessly with Bibby AI's citation management. Feel free to modify it for your needs!

\bibliographystyle{plain}
\bibliography{cite}

\end{document}\documentclass{article}
\usepackage[utf8]{inputenc}
\usepackage{booktabs}
\usepackage{amsmath}
\usepackage{natbib}

\title{My Document}
\author{Your Name}
\date{\today}

\begin{document}

\maketitle

\section{Introduction}
This is a sample document demonstrating various \LaTeX{} features including mathematical equations, tables, and formatting. You can cite papers using the citation search feature, and they will be automatically added to your \texttt{cite.bib} file.

\section{Mathematical Expressions}
Here are some example equations:

\subsection{Quadratic Formula}
The quadratic formula is given by:
\[x = \frac{-b \pm \sqrt{b^2 - 4ac}}{2a}\]

\subsection{Matrix Example}
A sample matrix:
\[
\begin{pmatrix}
a_{11} & a_{12} & a_{13} \\
a_{21} & a_{22} & a_{23} \\
a_{31} & a_{32} & a_{33}
\end{pmatrix}
\]

\section{Tables}
Below is a sample table showing data:

\begin{table}[h]
\centering
\begin{tabular}{@{}lrrr@{}}
\toprule
\textbf{Category} & \textbf{Value 1} & \textbf{Value 2} & \textbf{Total} \\
\midrule
Item A & 10.5 & 20.3 & 30.8 \\
Item B & 15.2 & 25.1 & 40.3 \\
Item C & 22.1 & 12.7 & 34.8 \\
\bottomrule
\end{tabular}
\caption{Sample Data Table}
\label{tab:sample}
\end{table}

\section{Lists}
\subsection{Bullet Points}
\begin{itemize}
  \item First item
  \item Second item
  \item Third item with sub-items:
    \begin{itemize}
      \item Sub-item 1
      \item Sub-item 2
    \end{itemize}
\end{itemize}

\subsection{Numbered List}
\begin{enumerate}
  \item First step
  \item Second step
  \item Third step
\end{enumerate}

\section{Citations and References}
When you use the citation search feature, references will be automatically managed. For example, you might cite a paper like this: \cite{example2024}. All citations will be automatically added to your \texttt{cite.bib} file.

\section{Conclusion}
This template demonstrates basic \LaTeX{} features including mathematical formulas, tables, and lists. The bibliography system is now set up to work seamlessly with Bibby AI's citation management. Feel free to modify it for your needs!

\bibliographystyle{plain}
\bibliography{cite}

\end{document}\documentclass{article}
\usepackage[utf8]{inputenc}
\usepackage{booktabs}
\usepackage{amsmath}
\usepackage{natbib}

\title{My Document}
\author{Your Name}
\date{\today}

\begin{document}

\maketitle

\section{Introduction}
This is a sample document demonstrating various \LaTeX{} features including mathematical equations, tables, and formatting. You can cite papers using the citation search feature, and they will be automatically added to your \texttt{cite.bib} file.

\section{Mathematical Expressions}
Here are some example equations:

\subsection{Quadratic Formula}
The quadratic formula is given by:
\[x = \frac{-b \pm \sqrt{b^2 - 4ac}}{2a}\]

\subsection{Matrix Example}
A sample matrix:
\[
\begin{pmatrix}
a_{11} & a_{12} & a_{13} \\
a_{21} & a_{22} & a_{23} \\
a_{31} & a_{32} & a_{33}
\end{pmatrix}
\]

\section{Tables}
Below is a sample table showing data:

\begin{table}[h]
\centering
\begin{tabular}{@{}lrrr@{}}
\toprule
\textbf{Category} & \textbf{Value 1} & \textbf{Value 2} & \textbf{Total} \\
\midrule
Item A & 10.5 & 20.3 & 30.8 \\
Item B & 15.2 & 25.1 & 40.3 \\
Item C & 22.1 & 12.7 & 34.8 \\
\bottomrule
\end{tabular}
\caption{Sample Data Table}
\label{tab:sample}
\end{table}

\section{Lists}
\subsection{Bullet Points}
\begin{itemize}
  \item First item
  \item Second item
  \item Third item with sub-items:
    \begin{itemize}
      \item Sub-item 1
      \item Sub-item 2
    \end{itemize}
\end{itemize}

\subsection{Numbered List}
\begin{enumerate}
  \item First step
  \item Second step
  \item Third step
\end{enumerate}

\section{Citations and References}
When you use the citation search feature, references will be automatically managed. For example, you might cite a paper like this: \cite{example2024}. All citations will be automatically added to your \texttt{cite.bib} file.

\section{Conclusion}
This template demonstrates basic \LaTeX{} features including mathematical formulas, tables, and lists. The bibliography system is now set up to work seamlessly with Bibby AI's citation management. Feel free to modify it for your needs!

\bibliographystyle{plain}
\bibliography{cite}

\end{document}\documentclass{article}
\usepackage[utf8]{inputenc}
\usepackage{booktabs}
\usepackage{amsmath}
\usepackage{natbib}

\title{My Document}
\author{Your Name}
\date{\today}

\begin{document}

\maketitle

\section{Introduction}
This is a sample document demonstrating various \LaTeX{} features including mathematical equations, tables, and formatting. You can cite papers using the citation search feature, and they will be automatically added to your \texttt{cite.bib} file.

\section{Mathematical Expressions}
Here are some example equations:

\subsection{Quadratic Formula}
The quadratic formula is given by:
\[x = \frac{-b \pm \sqrt{b^2 - 4ac}}{2a}\]

\subsection{Matrix Example}
A sample matrix:
\[
\begin{pmatrix}
a_{11} & a_{12} & a_{13} \\
a_{21} & a_{22} & a_{23} \\
a_{31} & a_{32} & a_{33}
\end{pmatrix}
\]

\section{Tables}
Below is a sample table showing data:

\begin{table}[h]
\centering
\begin{tabular}{@{}lrrr@{}}
\toprule
\textbf{Category} & \textbf{Value 1} & \textbf{Value 2} & \textbf{Total} \\
\midrule
Item A & 10.5 & 20.3 & 30.8 \\
Item B & 15.2 & 25.1 & 40.3 \\
Item C & 22.1 & 12.7 & 34.8 \\
\bottomrule
\end{tabular}
\caption{Sample Data Table}
\label{tab:sample}
\end{table}

\section{Lists}
\subsection{Bullet Points}
\begin{itemize}
  \item First item
  \item Second item
  \item Third item with sub-items:
    \begin{itemize}
      \item Sub-item 1
      \item Sub-item 2
    \end{itemize}
\end{itemize}

\subsection{Numbered List}
\begin{enumerate}
  \item First step
  \item Second step
  \item Third step
\end{enumerate}

\section{Citations and References}
When you use the citation search feature, references will be automatically managed. For example, you might cite a paper like this: \cite{example2024}. All citations will be automatically added to your \texttt{cite.bib} file.

\section{Conclusion}
This template demonstrates basic \LaTeX{} features including mathematical formulas, tables, and lists. The bibliography system is now set up to work seamlessly with Bibby AI's citation management. Feel free to modify it for your needs!

\bibliographystyle{plain}
\bibliography{cite}

\end{document}\documentclass{article}
\usepackage[utf8]{inputenc}
\usepackage{booktabs}
\usepackage{amsmath}
\usepackage{natbib}

\title{My Document}
\author{Your Name}
\date{\today}

\begin{document}

\maketitle

\section{Introduction}
This is a sample document demonstrating various \LaTeX{} features including mathematical equations, tables, and formatting. You can cite papers using the citation search feature, and they will be automatically added to your \texttt{cite.bib} file.

\section{Mathematical Expressions}
Here are some example equations:

\subsection{Quadratic Formula}
The quadratic formula is given by:
\[x = \frac{-b \pm \sqrt{b^2 - 4ac}}{2a}\]

\subsection{Matrix Example}
A sample matrix:
\[
\begin{pmatrix}
a_{11} & a_{12} & a_{13} \\
a_{21} & a_{22} & a_{23} \\
a_{31} & a_{32} & a_{33}
\end{pmatrix}
\]

\section{Tables}
Below is a sample table showing data:

\begin{table}[h]
\centering
\begin{tabular}{@{}lrrr@{}}
\toprule
\textbf{Category} & \textbf{Value 1} & \textbf{Value 2} & \textbf{Total} \\
\midrule
Item A & 10.5 & 20.3 & 30.8 \\
Item B & 15.2 & 25.1 & 40.3 \\
Item C & 22.1 & 12.7 & 34.8 \\
\bottomrule
\end{tabular}
\caption{Sample Data Table}
\label{tab:sample}
\end{table}

\section{Lists}
\subsection{Bullet Points}
\begin{itemize}
  \item First item
  \item Second item
  \item Third item with sub-items:
    \begin{itemize}
      \item Sub-item 1
      \item Sub-item 2
    \end{itemize}
\end{itemize}

\subsection{Numbered List}
\begin{enumerate}
  \item First step
  \item Second step
  \item Third step
\end{enumerate}

\section{Citations and References}
When you use the citation search feature, references will be automatically managed. For example, you might cite a paper like this: \cite{example2024}. All citations will be automatically added to your \texttt{cite.bib} file.

\section{Conclusion}
This template demonstrates basic \LaTeX{} features including mathematical formulas, tables, and lists. The bibliography system is now set up to work seamlessly with Bibby AI's citation management. Feel free to modify it for your needs!

\bibliographystyle{plain}
\bibliography{cite}

\end{document}\documentclass{article}
\usepackage[utf8]{inputenc}
\usepackage{booktabs}
\usepackage{amsmath}
\usepackage{natbib}

\title{My Document}
\author{Your Name}
\date{\today}

\begin{document}

\maketitle

\section{Introduction}
This is a sample document demonstrating various \LaTeX{} features including mathematical equations, tables, and formatting. You can cite papers using the citation search feature, and they will be automatically added to your \texttt{cite.bib} file.

\section{Mathematical Expressions}
Here are some example equations:

\subsection{Quadratic Formula}
The quadratic formula is given by:
\[x = \frac{-b \pm \sqrt{b^2 - 4ac}}{2a}\]

\subsection{Matrix Example}
A sample matrix:
\[
\begin{pmatrix}
a_{11} & a_{12} & a_{13} \\
a_{21} & a_{22} & a_{23} \\
a_{31} & a_{32} & a_{33}
\end{pmatrix}
\]

\section{Tables}
Below is a sample table showing data:

\begin{table}[h]
\centering
\begin{tabular}{@{}lrrr@{}}
\toprule
\textbf{Category} & \textbf{Value 1} & \textbf{Value 2} & \textbf{Total} \\
\midrule
Item A & 10.5 & 20.3 & 30.8 \\
Item B & 15.2 & 25.1 & 40.3 \\
Item C & 22.1 & 12.7 & 34.8 \\
\bottomrule
\end{tabular}
\caption{Sample Data Table}
\label{tab:sample}
\end{table}

\section{Lists}
\subsection{Bullet Points}
\begin{itemize}
  \item First item
  \item Second item
  \item Third item with sub-items:
    \begin{itemize}
      \item Sub-item 1
      \item Sub-item 2
    \end{itemize}
\end{itemize}

\subsection{Numbered List}
\begin{enumerate}
  \item First step
  \item Second step
  \item Third step
\end{enumerate}

\section{Citations and References}
When you use the citation search feature, references will be automatically managed. For example, you might cite a paper like this: \cite{example2024}. All citations will be automatically added to your \texttt{cite.bib} file.

\section{Conclusion}
This template demonstrates basic \LaTeX{} features including mathematical formulas, tables, and lists. The bibliography system is now set up to work seamlessly with Bibby AI's citation management. Feel free to modify it for your needs!

\bibliographystyle{plain}
\bibliography{cite}

\end{document}\documentclass{article}
\usepackage[utf8]{inputenc}
\usepackage{booktabs}
\usepackage{amsmath}
\usepackage{natbib}

\title{My Document}
\author{Your Name}
\date{\today}

\begin{document}

\maketitle

\section{Introduction}
This is a sample document demonstrating various \LaTeX{} features including mathematical equations, tables, and formatting. You can cite papers using the citation search feature, and they will be automatically added to your \texttt{cite.bib} file.

\section{Mathematical Expressions}
Here are some example equations:

\subsection{Quadratic Formula}
The quadratic formula is given by:
\[x = \frac{-b \pm \sqrt{b^2 - 4ac}}{2a}\]

\subsection{Matrix Example}
A sample matrix:
\[
\begin{pmatrix}
a_{11} & a_{12} & a_{13} \\
a_{21} & a_{22} & a_{23} \\
a_{31} & a_{32} & a_{33}
\end{pmatrix}
\]

\section{Tables}
Below is a sample table showing data:

\begin{table}[h]
\centering
\begin{tabular}{@{}lrrr@{}}
\toprule
\textbf{Category} & \textbf{Value 1} & \textbf{Value 2} & \textbf{Total} \\
\midrule
Item A & 10.5 & 20.3 & 30.8 \\
Item B & 15.2 & 25.1 & 40.3 \\
Item C & 22.1 & 12.7 & 34.8 \\
\bottomrule
\end{tabular}
\caption{Sample Data Table}
\label{tab:sample}
\end{table}

\section{Lists}
\subsection{Bullet Points}
\begin{itemize}
  \item First item
  \item Second item
  \item Third item with sub-items:
    \begin{itemize}
      \item Sub-item 1
      \item Sub-item 2
    \end{itemize}
\end{itemize}

\subsection{Numbered List}
\begin{enumerate}
  \item First step
  \item Second step
  \item Third step
\end{enumerate}

\section{Citations and References}
When you use the citation search feature, references will be automatically managed. For example, you might cite a paper like this: \cite{example2024}. All citations will be automatically added to your \texttt{cite.bib} file.

\section{Conclusion}
This template demonstrates basic \LaTeX{} features including mathematical formulas, tables, and lists. The bibliography system is now set up to work seamlessly with Bibby AI's citation management. Feel free to modify it for your needs!

\bibliographystyle{plain}
\bibliography{cite}

\end{document}\documentclass{article}
\usepackage[utf8]{inputenc}
\usepackage{booktabs}
\usepackage{amsmath}
\usepackage{natbib}

\title{My Document}
\author{Your Name}
\date{\today}

\begin{document}

\maketitle

\section{Introduction}
This is a sample document demonstrating various \LaTeX{} features including mathematical equations, tables, and formatting. You can cite papers using the citation search feature, and they will be automatically added to your \texttt{cite.bib} file.

\section{Mathematical Expressions}
Here are some example equations:

\subsection{Quadratic Formula}
The quadratic formula is given by:
\[x = \frac{-b \pm \sqrt{b^2 - 4ac}}{2a}\]

\subsection{Matrix Example}
A sample matrix:
\[
\begin{pmatrix}
a_{11} & a_{12} & a_{13} \\
a_{21} & a_{22} & a_{23} \\
a_{31} & a_{32} & a_{33}
\end{pmatrix}
\]

\section{Tables}
Below is a sample table showing data:

\begin{table}[h]
\centering
\begin{tabular}{@{}lrrr@{}}
\toprule
\textbf{Category} & \textbf{Value 1} & \textbf{Value 2} & \textbf{Total} \\
\midrule
Item A & 10.5 & 20.3 & 30.8 \\
Item B & 15.2 & 25.1 & 40.3 \\
Item C & 22.1 & 12.7 & 34.8 \\
\bottomrule
\end{tabular}
\caption{Sample Data Table}
\label{tab:sample}
\end{table}

\section{Lists}
\subsection{Bullet Points}
\begin{itemize}
  \item First item
  \item Second item
  \item Third item with sub-items:
    \begin{itemize}
      \item Sub-item 1
      \item Sub-item 2
    \end{itemize}
\end{itemize}

\subsection{Numbered List}
\begin{enumerate}
  \item First step
  \item Second step
  \item Third step
\end{enumerate}

\section{Citations and References}
When you use the citation search feature, references will be automatically managed. For example, you might cite a paper like this: \cite{example2024}. All citations will be automatically added to your \texttt{cite.bib} file.

\section{Conclusion}
This template demonstrates basic \LaTeX{} features including mathematical formulas, tables, and lists. The bibliography system is now set up to work seamlessly with Bibby AI's citation management. Feel free to modify it for your needs!

\bibliographystyle{plain}
\bibliography{cite}

\end{document}\documentclass{article}
\usepackage[utf8]{inputenc}
\usepackage{booktabs}
\usepackage{amsmath}
\usepackage{natbib}

\title{My Document}
\author{Your Name}
\date{\today}

\begin{document}

\maketitle

\section{Introduction}
This is a sample document demonstrating various \LaTeX{} features including mathematical equations, tables, and formatting. You can cite papers using the citation search feature, and they will be automatically added to your \texttt{cite.bib} file.

\section{Mathematical Expressions}
Here are some example equations:

\subsection{Quadratic Formula}
The quadratic formula is given by:
\[x = \frac{-b \pm \sqrt{b^2 - 4ac}}{2a}\]

\subsection{Matrix Example}
A sample matrix:
\[
\begin{pmatrix}
a_{11} & a_{12} & a_{13} \\
a_{21} & a_{22} & a_{23} \\
a_{31} & a_{32} & a_{33}
\end{pmatrix}
\]

\section{Tables}
Below is a sample table showing data:

\begin{table}[h]
\centering
\begin{tabular}{@{}lrrr@{}}
\toprule
\textbf{Category} & \textbf{Value 1} & \textbf{Value 2} & \textbf{Total} \\
\midrule
Item A & 10.5 & 20.3 & 30.8 \\
Item B & 15.2 & 25.1 & 40.3 \\
Item C & 22.1 & 12.7 & 34.8 \\
\bottomrule
\end{tabular}
\caption{Sample Data Table}
\label{tab:sample}
\end{table}

\section{Lists}
\subsection{Bullet Points}
\begin{itemize}
  \item First item
  \item Second item
  \item Third item with sub-items:
    \begin{itemize}
      \item Sub-item 1
      \item Sub-item 2
    \end{itemize}
\end{itemize}

\subsection{Numbered List}
\begin{enumerate}
  \item First step
  \item Second step
  \item Third step
\end{enumerate}

\section{Citations and References}
When you use the citation search feature, references will be automatically managed. For example, you might cite a paper like this: \cite{example2024}. All citations will be automatically added to your \texttt{cite.bib} file.

\section{Conclusion}
This template demonstrates basic \LaTeX{} features including mathematical formulas, tables, and lists. The bibliography system is now set up to work seamlessly with Bibby AI's citation management. Feel free to modify it for your needs!

\bibliographystyle{plain}
\bibliography{cite}

\end{document}\documentclass{article}
\usepackage[utf8]{inputenc}
\usepackage{booktabs}
\usepackage{amsmath}
\usepackage{natbib}

\title{My Document}
\author{Your Name}
\date{\today}

\begin{document}

\maketitle

\section{Introduction}
This is a sample document demonstrating various \LaTeX{} features including mathematical equations, tables, and formatting. You can cite papers using the citation search feature, and they will be automatically added to your \texttt{cite.bib} file.

\section{Mathematical Expressions}
Here are some example equations:

\subsection{Quadratic Formula}
The quadratic formula is given by:
\[x = \frac{-b \pm \sqrt{b^2 - 4ac}}{2a}\]

\subsection{Matrix Example}
A sample matrix:
\[
\begin{pmatrix}
a_{11} & a_{12} & a_{13} \\
a_{21} & a_{22} & a_{23} \\
a_{31} & a_{32} & a_{33}
\end{pmatrix}
\]

\section{Tables}
Below is a sample table showing data:

\begin{table}[h]
\centering
\begin{tabular}{@{}lrrr@{}}
\toprule
\textbf{Category} & \textbf{Value 1} & \textbf{Value 2} & \textbf{Total} \\
\midrule
Item A & 10.5 & 20.3 & 30.8 \\
Item B & 15.2 & 25.1 & 40.3 \\
Item C & 22.1 & 12.7 & 34.8 \\
\bottomrule
\end{tabular}
\caption{Sample Data Table}
\label{tab:sample}
\end{table}

\section{Lists}
\subsection{Bullet Points}
\begin{itemize}
  \item First item
  \item Second item
  \item Third item with sub-items:
    \begin{itemize}
      \item Sub-item 1
      \item Sub-item 2
    \end{itemize}
\end{itemize}

\subsection{Numbered List}
\begin{enumerate}
  \item First step
  \item Second step
  \item Third step
\end{enumerate}

\section{Citations and References}
When you use the citation search feature, references will be automatically managed. For example, you might cite a paper like this: \cite{example2024}. All citations will be automatically added to your \texttt{cite.bib} file.

\section{Conclusion}
This template demonstrates basic \LaTeX{} features including mathematical formulas, tables, and lists. The bibliography system is now set up to work seamlessly with Bibby AI's citation management. Feel free to modify it for your needs!

\bibliographystyle{plain}
\bibliography{cite}

\end{document}\documentclass{article}
\usepackage[utf8]{inputenc}
\usepackage{booktabs}
\usepackage{amsmath}
\usepackage{natbib}

\title{My Document}
\author{Your Name}
\date{\today}

\begin{document}

\maketitle

\section{Introduction}
This is a sample document demonstrating various \LaTeX{} features including mathematical equations, tables, and formatting. You can cite papers using the citation search feature, and they will be automatically added to your \texttt{cite.bib} file.

\section{Mathematical Expressions}
Here are some example equations:

\subsection{Quadratic Formula}
The quadratic formula is given by:
\[x = \frac{-b \pm \sqrt{b^2 - 4ac}}{2a}\]

\subsection{Matrix Example}
A sample matrix:
\[
\begin{pmatrix}
a_{11} & a_{12} & a_{13} \\
a_{21} & a_{22} & a_{23} \\
a_{31} & a_{32} & a_{33}
\end{pmatrix}
\]

\section{Tables}
Below is a sample table showing data:

\begin{table}[h]
\centering
\begin{tabular}{@{}lrrr@{}}
\toprule
\textbf{Category} & \textbf{Value 1} & \textbf{Value 2} & \textbf{Total} \\
\midrule
Item A & 10.5 & 20.3 & 30.8 \\
Item B & 15.2 & 25.1 & 40.3 \\
Item C & 22.1 & 12.7 & 34.8 \\
\bottomrule
\end{tabular}
\caption{Sample Data Table}
\label{tab:sample}
\end{table}

\section{Lists}
\subsection{Bullet Points}
\begin{itemize}
  \item First item
  \item Second item
  \item Third item with sub-items:
    \begin{itemize}
      \item Sub-item 1
      \item Sub-item 2
    \end{itemize}
\end{itemize}

\subsection{Numbered List}
\begin{enumerate}
  \item First step
  \item Second step
  \item Third step
\end{enumerate}

\section{Citations and References}
When you use the citation search feature, references will be automatically managed. For example, you might cite a paper like this: \cite{example2024}. All citations will be automatically added to your \texttt{cite.bib} file.

\section{Conclusion}
This template demonstrates basic \LaTeX{} features including mathematical formulas, tables, and lists. The bibliography system is now set up to work seamlessly with Bibby AI's citation management. Feel free to modify it for your needs!

\bibliographystyle{plain}
\bibliography{cite}

\end{document}\documentclass{article}
\usepackage[utf8]{inputenc}
\usepackage{booktabs}
\usepackage{amsmath}
\usepackage{natbib}

\title{My Document}
\author{Your Name}
\date{\today}

\begin{document}

\maketitle

\section{Introduction}
This is a sample document demonstrating various \LaTeX{} features including mathematical equations, tables, and formatting. You can cite papers using the citation search feature, and they will be automatically added to your \texttt{cite.bib} file.

\section{Mathematical Expressions}
Here are some example equations:

\subsection{Quadratic Formula}
The quadratic formula is given by:
\[x = \frac{-b \pm \sqrt{b^2 - 4ac}}{2a}\]

\subsection{Matrix Example}
A sample matrix:
\[
\begin{pmatrix}
a_{11} & a_{12} & a_{13} \\
a_{21} & a_{22} & a_{23} \\
a_{31} & a_{32} & a_{33}
\end{pmatrix}
\]

\section{Tables}
Below is a sample table showing data:

\begin{table}[h]
\centering
\begin{tabular}{@{}lrrr@{}}
\toprule
\textbf{Category} & \textbf{Value 1} & \textbf{Value 2} & \textbf{Total} \\
\midrule
Item A & 10.5 & 20.3 & 30.8 \\
Item B & 15.2 & 25.1 & 40.3 \\
Item C & 22.1 & 12.7 & 34.8 \\
\bottomrule
\end{tabular}
\caption{Sample Data Table}
\label{tab:sample}
\end{table}

\section{Lists}
\subsection{Bullet Points}
\begin{itemize}
  \item First item
  \item Second item
  \item Third item with sub-items:
    \begin{itemize}
      \item Sub-item 1
      \item Sub-item 2
    \end{itemize}
\end{itemize}

\subsection{Numbered List}
\begin{enumerate}
  \item First step
  \item Second step
  \item Third step
\end{enumerate}

\section{Citations and References}
When you use the citation search feature, references will be automatically managed. For example, you might cite a paper like this: \cite{example2024}. All citations will be automatically added to your \texttt{cite.bib} file.

\section{Conclusion}
This template demonstrates basic \LaTeX{} features including mathematical formulas, tables, and lists. The bibliography system is now set up to work seamlessly with Bibby AI's citation management. Feel free to modify it for your needs!

\bibliographystyle{plain}
\bibliography{cite}

\end{document}\documentclass{article}
\usepackage[utf8]{inputenc}
\usepackage{booktabs}
\usepackage{amsmath}
\usepackage{natbib}

\title{My Document}
\author{Your Name}
\date{\today}

\begin{document}

\maketitle

\section{Introduction}
This is a sample document demonstrating various \LaTeX{} features including mathematical equations, tables, and formatting. You can cite papers using the citation search feature, and they will be automatically added to your \texttt{cite.bib} file.

\section{Mathematical Expressions}
Here are some example equations:

\subsection{Quadratic Formula}
The quadratic formula is given by:
\[x = \frac{-b \pm \sqrt{b^2 - 4ac}}{2a}\]

\subsection{Matrix Example}
A sample matrix:
\[
\begin{pmatrix}
a_{11} & a_{12} & a_{13} \\
a_{21} & a_{22} & a_{23} \\
a_{31} & a_{32} & a_{33}
\end{pmatrix}
\]

\section{Tables}
Below is a sample table showing data:

\begin{table}[h]
\centering
\begin{tabular}{@{}lrrr@{}}
\toprule
\textbf{Category} & \textbf{Value 1} & \textbf{Value 2} & \textbf{Total} \\
\midrule
Item A & 10.5 & 20.3 & 30.8 \\
Item B & 15.2 & 25.1 & 40.3 \\
Item C & 22.1 & 12.7 & 34.8 \\
\bottomrule
\end{tabular}
\caption{Sample Data Table}
\label{tab:sample}
\end{table}

\section{Lists}
\subsection{Bullet Points}
\begin{itemize}
  \item First item
  \item Second item
  \item Third item with sub-items:
    \begin{itemize}
      \item Sub-item 1
      \item Sub-item 2
    \end{itemize}
\end{itemize}

\subsection{Numbered List}
\begin{enumerate}
  \item First step
  \item Second step
  \item Third step
\end{enumerate}

\section{Citations and References}
When you use the citation search feature, references will be automatically managed. For example, you might cite a paper like this: \cite{example2024}. All citations will be automatically added to your \texttt{cite.bib} file.

\section{Conclusion}
This template demonstrates basic \LaTeX{} features including mathematical formulas, tables, and lists. The bibliography system is now set up to work seamlessly with Bibby AI's citation management. Feel free to modify it for your needs!

\bibliographystyle{plain}
\bibliography{cite}

\end{document}\documentclass{article}
\usepackage[utf8]{inputenc}
\usepackage{booktabs}
\usepackage{amsmath}
\usepackage{natbib}

\title{My Document}
\author{Your Name}
\date{\today}

\begin{document}

\maketitle

\section{Introduction}
This is a sample document demonstrating various \LaTeX{} features including mathematical equations, tables, and formatting. You can cite papers using the citation search feature, and they will be automatically added to your \texttt{cite.bib} file.

\section{Mathematical Expressions}
Here are some example equations:

\subsection{Quadratic Formula}
The quadratic formula is given by:
\[x = \frac{-b \pm \sqrt{b^2 - 4ac}}{2a}\]

\subsection{Matrix Example}
A sample matrix:
\[
\begin{pmatrix}
a_{11} & a_{12} & a_{13} \\
a_{21} & a_{22} & a_{23} \\
a_{31} & a_{32} & a_{33}
\end{pmatrix}
\]

\section{Tables}
Below is a sample table showing data:

\begin{table}[h]
\centering
\begin{tabular}{@{}lrrr@{}}
\toprule
\textbf{Category} & \textbf{Value 1} & \textbf{Value 2} & \textbf{Total} \\
\midrule
Item A & 10.5 & 20.3 & 30.8 \\
Item B & 15.2 & 25.1 & 40.3 \\
Item C & 22.1 & 12.7 & 34.8 \\
\bottomrule
\end{tabular}
\caption{Sample Data Table}
\label{tab:sample}
\end{table}

\section{Lists}
\subsection{Bullet Points}
\begin{itemize}
  \item First item
  \item Second item
  \item Third item with sub-items:
    \begin{itemize}
      \item Sub-item 1
      \item Sub-item 2
    \end{itemize}
\end{itemize}

\subsection{Numbered List}
\begin{enumerate}
  \item First step
  \item Second step
  \item Third step
\end{enumerate}

\section{Citations and References}
When you use the citation search feature, references will be automatically managed. For example, you might cite a paper like this: \cite{example2024}. All citations will be automatically added to your \texttt{cite.bib} file.

\section{Conclusion}
This template demonstrates basic \LaTeX{} features including mathematical formulas, tables, and lists. The bibliography system is now set up to work seamlessly with Bibby AI's citation management. Feel free to modify it for your needs!

\bibliographystyle{plain}
\bibliography{cite}

\end{document}\documentclass{article}
\usepackage[utf8]{inputenc}
\usepackage{booktabs}
\usepackage{amsmath}
\usepackage{natbib}

\title{My Document}
\author{Your Name}
\date{\today}

\begin{document}

\maketitle

\section{Introduction}
This is a sample document demonstrating various \LaTeX{} features including mathematical equations, tables, and formatting. You can cite papers using the citation search feature, and they will be automatically added to your \texttt{cite.bib} file.

\section{Mathematical Expressions}
Here are some example equations:

\subsection{Quadratic Formula}
The quadratic formula is given by:
\[x = \frac{-b \pm \sqrt{b^2 - 4ac}}{2a}\]

\subsection{Matrix Example}
A sample matrix:
\[
\begin{pmatrix}
a_{11} & a_{12} & a_{13} \\
a_{21} & a_{22} & a_{23} \\
a_{31} & a_{32} & a_{33}
\end{pmatrix}
\]

\section{Tables}
Below is a sample table showing data:

\begin{table}[h]
\centering
\begin{tabular}{@{}lrrr@{}}
\toprule
\textbf{Category} & \textbf{Value 1} & \textbf{Value 2} & \textbf{Total} \\
\midrule
Item A & 10.5 & 20.3 & 30.8 \\
Item B & 15.2 & 25.1 & 40.3 \\
Item C & 22.1 & 12.7 & 34.8 \\
\bottomrule
\end{tabular}
\caption{Sample Data Table}
\label{tab:sample}
\end{table}

\section{Lists}
\subsection{Bullet Points}
\begin{itemize}
  \item First item
  \item Second item
  \item Third item with sub-items:
    \begin{itemize}
      \item Sub-item 1
      \item Sub-item 2
    \end{itemize}
\end{itemize}

\subsection{Numbered List}
\begin{enumerate}
  \item First step
  \item Second step
  \item Third step
\end{enumerate}

\section{Citations and References}
When you use the citation search feature, references will be automatically managed. For example, you might cite a paper like this: \cite{example2024}. All citations will be automatically added to your \texttt{cite.bib} file.

\section{Conclusion}
This template demonstrates basic \LaTeX{} features including mathematical formulas, tables, and lists. The bibliography system is now set up to work seamlessly with Bibby AI's citation management. Feel free to modify it for your needs!

\bibliographystyle{plain}
\bibliography{cite}

\end{document}\documentclass{article}
\usepackage[utf8]{inputenc}
\usepackage{booktabs}
\usepackage{amsmath}
\usepackage{natbib}

\title{My Document}
\author{Your Name}
\date{\today}

\begin{document}

\maketitle

\section{Introduction}
This is a sample document demonstrating various \LaTeX{} features including mathematical equations, tables, and formatting. You can cite papers using the citation search feature, and they will be automatically added to your \texttt{cite.bib} file.

\section{Mathematical Expressions}
Here are some example equations:

\subsection{Quadratic Formula}
The quadratic formula is given by:
\[x = \frac{-b \pm \sqrt{b^2 - 4ac}}{2a}\]

\subsection{Matrix Example}
A sample matrix:
\[
\begin{pmatrix}
a_{11} & a_{12} & a_{13} \\
a_{21} & a_{22} & a_{23} \\
a_{31} & a_{32} & a_{33}
\end{pmatrix}
\]

\section{Tables}
Below is a sample table showing data:

\begin{table}[h]
\centering
\begin{tabular}{@{}lrrr@{}}
\toprule
\textbf{Category} & \textbf{Value 1} & \textbf{Value 2} & \textbf{Total} \\
\midrule
Item A & 10.5 & 20.3 & 30.8 \\
Item B & 15.2 & 25.1 & 40.3 \\
Item C & 22.1 & 12.7 & 34.8 \\
\bottomrule
\end{tabular}
\caption{Sample Data Table}
\label{tab:sample}
\end{table}

\section{Lists}
\subsection{Bullet Points}
\begin{itemize}
  \item First item
  \item Second item
  \item Third item with sub-items:
    \begin{itemize}
      \item Sub-item 1
      \item Sub-item 2
    \end{itemize}
\end{itemize}

\subsection{Numbered List}
\begin{enumerate}
  \item First step
  \item Second step
  \item Third step
\end{enumerate}

\section{Citations and References}
When you use the citation search feature, references will be automatically managed. For example, you might cite a paper like this: \cite{example2024}. All citations will be automatically added to your \texttt{cite.bib} file.

\section{Conclusion}
This template demonstrates basic \LaTeX{} features including mathematical formulas, tables, and lists. The bibliography system is now set up to work seamlessly with Bibby AI's citation management. Feel free to modify it for your needs!

\bibliographystyle{plain}
\bibliography{cite}

\end{document}\documentclass{article}
\usepackage[utf8]{inputenc}
\usepackage{booktabs}
\usepackage{amsmath}
\usepackage{natbib}

\title{My Document}
\author{Your Name}
\date{\today}

\begin{document}

\maketitle

\section{Introduction}
This is a sample document demonstrating various \LaTeX{} features including mathematical equations, tables, and formatting. You can cite papers using the citation search feature, and they will be automatically added to your \texttt{cite.bib} file.

\section{Mathematical Expressions}
Here are some example equations:

\subsection{Quadratic Formula}
The quadratic formula is given by:
\[x = \frac{-b \pm \sqrt{b^2 - 4ac}}{2a}\]

\subsection{Matrix Example}
A sample matrix:
\[
\begin{pmatrix}
a_{11} & a_{12} & a_{13} \\
a_{21} & a_{22} & a_{23} \\
a_{31} & a_{32} & a_{33}
\end{pmatrix}
\]

\section{Tables}
Below is a sample table showing data:

\begin{table}[h]
\centering
\begin{tabular}{@{}lrrr@{}}
\toprule
\textbf{Category} & \textbf{Value 1} & \textbf{Value 2} & \textbf{Total} \\
\midrule
Item A & 10.5 & 20.3 & 30.8 \\
Item B & 15.2 & 25.1 & 40.3 \\
Item C & 22.1 & 12.7 & 34.8 \\
\bottomrule
\end{tabular}
\caption{Sample Data Table}
\label{tab:sample}
\end{table}

\section{Lists}
\subsection{Bullet Points}
\begin{itemize}
  \item First item
  \item Second item
  \item Third item with sub-items:
    \begin{itemize}
      \item Sub-item 1
      \item Sub-item 2
    \end{itemize}
\end{itemize}

\subsection{Numbered List}
\begin{enumerate}
  \item First step
  \item Second step
  \item Third step
\end{enumerate}

\section{Citations and References}
When you use the citation search feature, references will be automatically managed. For example, you might cite a paper like this: \cite{example2024}. All citations will be automatically added to your \texttt{cite.bib} file.

\section{Conclusion}
This template demonstrates basic \LaTeX{} features including mathematical formulas, tables, and lists. The bibliography system is now set up to work seamlessly with Bibby AI's citation management. Feel free to modify it for your needs!

\bibliographystyle{plain}
\bibliography{cite}

\end{document}\documentclass{article}
\usepackage[utf8]{inputenc}
\usepackage{booktabs}
\usepackage{amsmath}
\usepackage{natbib}

\title{My Document}
\author{Your Name}
\date{\today}

\begin{document}

\maketitle

\section{Introduction}
This is a sample document demonstrating various \LaTeX{} features including mathematical equations, tables, and formatting. You can cite papers using the citation search feature, and they will be automatically added to your \texttt{cite.bib} file.

\section{Mathematical Expressions}
Here are some example equations:

\subsection{Quadratic Formula}
The quadratic formula is given by:
\[x = \frac{-b \pm \sqrt{b^2 - 4ac}}{2a}\]

\subsection{Matrix Example}
A sample matrix:
\[
\begin{pmatrix}
a_{11} & a_{12} & a_{13} \\
a_{21} & a_{22} & a_{23} \\
a_{31} & a_{32} & a_{33}
\end{pmatrix}
\]

\section{Tables}
Below is a sample table showing data:

\begin{table}[h]
\centering
\begin{tabular}{@{}lrrr@{}}
\toprule
\textbf{Category} & \textbf{Value 1} & \textbf{Value 2} & \textbf{Total} \\
\midrule
Item A & 10.5 & 20.3 & 30.8 \\
Item B & 15.2 & 25.1 & 40.3 \\
Item C & 22.1 & 12.7 & 34.8 \\
\bottomrule
\end{tabular}
\caption{Sample Data Table}
\label{tab:sample}
\end{table}

\section{Lists}
\subsection{Bullet Points}
\begin{itemize}
  \item First item
  \item Second item
  \item Third item with sub-items:
    \begin{itemize}
      \item Sub-item 1
      \item Sub-item 2
    \end{itemize}
\end{itemize}

\subsection{Numbered List}
\begin{enumerate}
  \item First step
  \item Second step
  \item Third step
\end{enumerate}

\section{Citations and References}
When you use the citation search feature, references will be automatically managed. For example, you might cite a paper like this: \cite{example2024}. All citations will be automatically added to your \texttt{cite.bib} file.

\section{Conclusion}
This template demonstrates basic \LaTeX{} features including mathematical formulas, tables, and lists. The bibliography system is now set up to work seamlessly with Bibby AI's citation management. Feel free to modify it for your needs!

\bibliographystyle{plain}
\bibliography{cite}

\end{document}\documentclass{article}
\usepackage[utf8]{inputenc}
\usepackage{booktabs}
\usepackage{amsmath}
\usepackage{natbib}

\title{My Document}
\author{Your Name}
\date{\today}

\begin{document}

\maketitle

\section{Introduction}
This is a sample document demonstrating various \LaTeX{} features including mathematical equations, tables, and formatting. You can cite papers using the citation search feature, and they will be automatically added to your \texttt{cite.bib} file.

\section{Mathematical Expressions}
Here are some example equations:

\subsection{Quadratic Formula}
The quadratic formula is given by:
\[x = \frac{-b \pm \sqrt{b^2 - 4ac}}{2a}\]

\subsection{Matrix Example}
A sample matrix:
\[
\begin{pmatrix}
a_{11} & a_{12} & a_{13} \\
a_{21} & a_{22} & a_{23} \\
a_{31} & a_{32} & a_{33}
\end{pmatrix}
\]

\section{Tables}
Below is a sample table showing data:

\begin{table}[h]
\centering
\begin{tabular}{@{}lrrr@{}}
\toprule
\textbf{Category} & \textbf{Value 1} & \textbf{Value 2} & \textbf{Total} \\
\midrule
Item A & 10.5 & 20.3 & 30.8 \\
Item B & 15.2 & 25.1 & 40.3 \\
Item C & 22.1 & 12.7 & 34.8 \\
\bottomrule
\end{tabular}
\caption{Sample Data Table}
\label{tab:sample}
\end{table}

\section{Lists}
\subsection{Bullet Points}
\begin{itemize}
  \item First item
  \item Second item
  \item Third item with sub-items:
    \begin{itemize}
      \item Sub-item 1
      \item Sub-item 2
    \end{itemize}
\end{itemize}

\subsection{Numbered List}
\begin{enumerate}
  \item First step
  \item Second step
  \item Third step
\end{enumerate}

\section{Citations and References}
When you use the citation search feature, references will be automatically managed. For example, you might cite a paper like this: \cite{example2024}. All citations will be automatically added to your \texttt{cite.bib} file.

\section{Conclusion}
This template demonstrates basic \LaTeX{} features including mathematical formulas, tables, and lists. The bibliography system is now set up to work seamlessly with Bibby AI's citation management. Feel free to modify it for your needs!

\bibliographystyle{plain}
\bibliography{cite}

\end{document}\documentclass{article}
\usepackage[utf8]{inputenc}
\usepackage{booktabs}
\usepackage{amsmath}
\usepackage{natbib}

\title{My Document}
\author{Your Name}
\date{\today}

\begin{document}

\maketitle

\section{Introduction}
This is a sample document demonstrating various \LaTeX{} features including mathematical equations, tables, and formatting. You can cite papers using the citation search feature, and they will be automatically added to your \texttt{cite.bib} file.

\section{Mathematical Expressions}
Here are some example equations:

\subsection{Quadratic Formula}
The quadratic formula is given by:
\[x = \frac{-b \pm \sqrt{b^2 - 4ac}}{2a}\]

\subsection{Matrix Example}
A sample matrix:
\[
\begin{pmatrix}
a_{11} & a_{12} & a_{13} \\
a_{21} & a_{22} & a_{23} \\
a_{31} & a_{32} & a_{33}
\end{pmatrix}
\]

\section{Tables}
Below is a sample table showing data:

\begin{table}[h]
\centering
\begin{tabular}{@{}lrrr@{}}
\toprule
\textbf{Category} & \textbf{Value 1} & \textbf{Value 2} & \textbf{Total} \\
\midrule
Item A & 10.5 & 20.3 & 30.8 \\
Item B & 15.2 & 25.1 & 40.3 \\
Item C & 22.1 & 12.7 & 34.8 \\
\bottomrule
\end{tabular}
\caption{Sample Data Table}
\label{tab:sample}
\end{table}

\section{Lists}
\subsection{Bullet Points}
\begin{itemize}
  \item First item
  \item Second item
  \item Third item with sub-items:
    \begin{itemize}
      \item Sub-item 1
      \item Sub-item 2
    \end{itemize}
\end{itemize}

\subsection{Numbered List}
\begin{enumerate}
  \item First step
  \item Second step
  \item Third step
\end{enumerate}

\section{Citations and References}
When you use the citation search feature, references will be automatically managed. For example, you might cite a paper like this: \cite{example2024}. All citations will be automatically added to your \texttt{cite.bib} file.

\section{Conclusion}
This template demonstrates basic \LaTeX{} features including mathematical formulas, tables, and lists. The bibliography system is now set up to work seamlessly with Bibby AI's citation management. Feel free to modify it for your needs!

\bibliographystyle{plain}
\bibliography{cite}

\end{document}\documentclass{article}
\usepackage[utf8]{inputenc}
\usepackage{booktabs}
\usepackage{amsmath}
\usepackage{natbib}

\title{My Document}
\author{Your Name}
\date{\today}

\begin{document}

\maketitle

\section{Introduction}
This is a sample document demonstrating various \LaTeX{} features including mathematical equations, tables, and formatting. You can cite papers using the citation search feature, and they will be automatically added to your \texttt{cite.bib} file.

\section{Mathematical Expressions}
Here are some example equations:

\subsection{Quadratic Formula}
The quadratic formula is given by:
\[x = \frac{-b \pm \sqrt{b^2 - 4ac}}{2a}\]

\subsection{Matrix Example}
A sample matrix:
\[
\begin{pmatrix}
a_{11} & a_{12} & a_{13} \\
a_{21} & a_{22} & a_{23} \\
a_{31} & a_{32} & a_{33}
\end{pmatrix}
\]

\section{Tables}
Below is a sample table showing data:

\begin{table}[h]
\centering
\begin{tabular}{@{}lrrr@{}}
\toprule
\textbf{Category} & \textbf{Value 1} & \textbf{Value 2} & \textbf{Total} \\
\midrule
Item A & 10.5 & 20.3 & 30.8 \\
Item B & 15.2 & 25.1 & 40.3 \\
Item C & 22.1 & 12.7 & 34.8 \\
\bottomrule
\end{tabular}
\caption{Sample Data Table}
\label{tab:sample}
\end{table}

\section{Lists}
\subsection{Bullet Points}
\begin{itemize}
  \item First item
  \item Second item
  \item Third item with sub-items:
    \begin{itemize}
      \item Sub-item 1
      \item Sub-item 2
    \end{itemize}
\end{itemize}

\subsection{Numbered List}
\begin{enumerate}
  \item First step
  \item Second step
  \item Third step
\end{enumerate}

\section{Citations and References}
When you use the citation search feature, references will be automatically managed. For example, you might cite a paper like this: \cite{example2024}. All citations will be automatically added to your \texttt{cite.bib} file.

\section{Conclusion}
This template demonstrates basic \LaTeX{} features including mathematical formulas, tables, and lists. The bibliography system is now set up to work seamlessly with Bibby AI's citation management. Feel free to modify it for your needs!

\bibliographystyle{plain}
\bibliography{cite}

\end{document}\documentclass{article}
\usepackage[utf8]{inputenc}
\usepackage{booktabs}
\usepackage{amsmath}
\usepackage{natbib}

\title{My Document}
\author{Your Name}
\date{\today}

\begin{document}

\maketitle

\section{Introduction}
This is a sample document demonstrating various \LaTeX{} features including mathematical equations, tables, and formatting. You can cite papers using the citation search feature, and they will be automatically added to your \texttt{cite.bib} file.

\section{Mathematical Expressions}
Here are some example equations:

\subsection{Quadratic Formula}
The quadratic formula is given by:
\[x = \frac{-b \pm \sqrt{b^2 - 4ac}}{2a}\]

\subsection{Matrix Example}
A sample matrix:
\[
\begin{pmatrix}
a_{11} & a_{12} & a_{13} \\
a_{21} & a_{22} & a_{23} \\
a_{31} & a_{32} & a_{33}
\end{pmatrix}
\]

\section{Tables}
Below is a sample table showing data:

\begin{table}[h]
\centering
\begin{tabular}{@{}lrrr@{}}
\toprule
\textbf{Category} & \textbf{Value 1} & \textbf{Value 2} & \textbf{Total} \\
\midrule
Item A & 10.5 & 20.3 & 30.8 \\
Item B & 15.2 & 25.1 & 40.3 \\
Item C & 22.1 & 12.7 & 34.8 \\
\bottomrule
\end{tabular}
\caption{Sample Data Table}
\label{tab:sample}
\end{table}

\section{Lists}
\subsection{Bullet Points}
\begin{itemize}
  \item First item
  \item Second item
  \item Third item with sub-items:
    \begin{itemize}
      \item Sub-item 1
      \item Sub-item 2
    \end{itemize}
\end{itemize}

\subsection{Numbered List}
\begin{enumerate}
  \item First step
  \item Second step
  \item Third step
\end{enumerate}

\section{Citations and References}
When you use the citation search feature, references will be automatically managed. For example, you might cite a paper like this: \cite{example2024}. All citations will be automatically added to your \texttt{cite.bib} file.

\section{Conclusion}
This template demonstrates basic \LaTeX{} features including mathematical formulas, tables, and lists. The bibliography system is now set up to work seamlessly with Bibby AI's citation management. Feel free to modify it for your needs!

\bibliographystyle{plain}
\bibliography{cite}

\end{document}\documentclass{article}
\usepackage[utf8]{inputenc}
\usepackage{booktabs}
\usepackage{amsmath}
\usepackage{natbib}

\title{My Document}
\author{Your Name}
\date{\today}

\begin{document}

\maketitle

\section{Introduction}
This is a sample document demonstrating various \LaTeX{} features including mathematical equations, tables, and formatting. You can cite papers using the citation search feature, and they will be automatically added to your \texttt{cite.bib} file.

\section{Mathematical Expressions}
Here are some example equations:

\subsection{Quadratic Formula}
The quadratic formula is given by:
\[x = \frac{-b \pm \sqrt{b^2 - 4ac}}{2a}\]

\subsection{Matrix Example}
A sample matrix:
\[
\begin{pmatrix}
a_{11} & a_{12} & a_{13} \\
a_{21} & a_{22} & a_{23} \\
a_{31} & a_{32} & a_{33}
\end{pmatrix}
\]

\section{Tables}
Below is a sample table showing data:

\begin{table}[h]
\centering
\begin{tabular}{@{}lrrr@{}}
\toprule
\textbf{Category} & \textbf{Value 1} & \textbf{Value 2} & \textbf{Total} \\
\midrule
Item A & 10.5 & 20.3 & 30.8 \\
Item B & 15.2 & 25.1 & 40.3 \\
Item C & 22.1 & 12.7 & 34.8 \\
\bottomrule
\end{tabular}
\caption{Sample Data Table}
\label{tab:sample}
\end{table}

\section{Lists}
\subsection{Bullet Points}
\begin{itemize}
  \item First item
  \item Second item
  \item Third item with sub-items:
    \begin{itemize}
      \item Sub-item 1
      \item Sub-item 2
    \end{itemize}
\end{itemize}

\subsection{Numbered List}
\begin{enumerate}
  \item First step
  \item Second step
  \item Third step
\end{enumerate}

\section{Citations and References}
When you use the citation search feature, references will be automatically managed. For example, you might cite a paper like this: \cite{example2024}. All citations will be automatically added to your \texttt{cite.bib} file.

\section{Conclusion}
This template demonstrates basic \LaTeX{} features including mathematical formulas, tables, and lists. The bibliography system is now set up to work seamlessly with Bibby AI's citation management. Feel free to modify it for your needs!

\bibliographystyle{plain}
\bibliography{cite}

\end{document}\documentclass{article}
\usepackage[utf8]{inputenc}
\usepackage{booktabs}
\usepackage{amsmath}
\usepackage{natbib}

\title{My Document}
\author{Your Name}
\date{\today}

\begin{document}

\maketitle

\section{Introduction}
This is a sample document demonstrating various \LaTeX{} features including mathematical equations, tables, and formatting. You can cite papers using the citation search feature, and they will be automatically added to your \texttt{cite.bib} file.

\section{Mathematical Expressions}
Here are some example equations:

\subsection{Quadratic Formula}
The quadratic formula is given by:
\[x = \frac{-b \pm \sqrt{b^2 - 4ac}}{2a}\]

\subsection{Matrix Example}
A sample matrix:
\[
\begin{pmatrix}
a_{11} & a_{12} & a_{13} \\
a_{21} & a_{22} & a_{23} \\
a_{31} & a_{32} & a_{33}
\end{pmatrix}
\]

\section{Tables}
Below is a sample table showing data:

\begin{table}[h]
\centering
\begin{tabular}{@{}lrrr@{}}
\toprule
\textbf{Category} & \textbf{Value 1} & \textbf{Value 2} & \textbf{Total} \\
\midrule
Item A & 10.5 & 20.3 & 30.8 \\
Item B & 15.2 & 25.1 & 40.3 \\
Item C & 22.1 & 12.7 & 34.8 \\
\bottomrule
\end{tabular}
\caption{Sample Data Table}
\label{tab:sample}
\end{table}

\section{Lists}
\subsection{Bullet Points}
\begin{itemize}
  \item First item
  \item Second item
  \item Third item with sub-items:
    \begin{itemize}
      \item Sub-item 1
      \item Sub-item 2
    \end{itemize}
\end{itemize}

\subsection{Numbered List}
\begin{enumerate}
  \item First step
  \item Second step
  \item Third step
\end{enumerate}

\section{Citations and References}
When you use the citation search feature, references will be automatically managed. For example, you might cite a paper like this: \cite{example2024}. All citations will be automatically added to your \texttt{cite.bib} file.

\section{Conclusion}
This template demonstrates basic \LaTeX{} features including mathematical formulas, tables, and lists. The bibliography system is now set up to work seamlessly with Bibby AI's citation management. Feel free to modify it for your needs!

\bibliographystyle{plain}
\bibliography{cite}

\end{document}\documentclass{article}
\usepackage[utf8]{inputenc}
\usepackage{booktabs}
\usepackage{amsmath}
\usepackage{natbib}

\title{My Document}
\author{Your Name}
\date{\today}

\begin{document}

\maketitle

\section{Introduction}
This is a sample document demonstrating various \LaTeX{} features including mathematical equations, tables, and formatting. You can cite papers using the citation search feature, and they will be automatically added to your \texttt{cite.bib} file.

\section{Mathematical Expressions}
Here are some example equations:

\subsection{Quadratic Formula}
The quadratic formula is given by:
\[x = \frac{-b \pm \sqrt{b^2 - 4ac}}{2a}\]

\subsection{Matrix Example}
A sample matrix:
\[
\begin{pmatrix}
a_{11} & a_{12} & a_{13} \\
a_{21} & a_{22} & a_{23} \\
a_{31} & a_{32} & a_{33}
\end{pmatrix}
\]

\section{Tables}
Below is a sample table showing data:

\begin{table}[h]
\centering
\begin{tabular}{@{}lrrr@{}}
\toprule
\textbf{Category} & \textbf{Value 1} & \textbf{Value 2} & \textbf{Total} \\
\midrule
Item A & 10.5 & 20.3 & 30.8 \\
Item B & 15.2 & 25.1 & 40.3 \\
Item C & 22.1 & 12.7 & 34.8 \\
\bottomrule
\end{tabular}
\caption{Sample Data Table}
\label{tab:sample}
\end{table}

\section{Lists}
\subsection{Bullet Points}
\begin{itemize}
  \item First item
  \item Second item
  \item Third item with sub-items:
    \begin{itemize}
      \item Sub-item 1
      \item Sub-item 2
    \end{itemize}
\end{itemize}

\subsection{Numbered List}
\begin{enumerate}
  \item First step
  \item Second step
  \item Third step
\end{enumerate}

\section{Citations and References}
When you use the citation search feature, references will be automatically managed. For example, you might cite a paper like this: \cite{example2024}. All citations will be automatically added to your \texttt{cite.bib} file.

\section{Conclusion}
This template demonstrates basic \LaTeX{} features including mathematical formulas, tables, and lists. The bibliography system is now set up to work seamlessly with Bibby AI's citation management. Feel free to modify it for your needs!

\bibliographystyle{plain}
\bibliography{cite}

\end{document}\documentclass{article}
\usepackage[utf8]{inputenc}
\usepackage{booktabs}
\usepackage{amsmath}
\usepackage{natbib}

\title{My Document}
\author{Your Name}
\date{\today}

\begin{document}

\maketitle

\section{Introduction}
This is a sample document demonstrating various \LaTeX{} features including mathematical equations, tables, and formatting. You can cite papers using the citation search feature, and they will be automatically added to your \texttt{cite.bib} file.

\section{Mathematical Expressions}
Here are some example equations:

\subsection{Quadratic Formula}
The quadratic formula is given by:
\[x = \frac{-b \pm \sqrt{b^2 - 4ac}}{2a}\]

\subsection{Matrix Example}
A sample matrix:
\[
\begin{pmatrix}
a_{11} & a_{12} & a_{13} \\
a_{21} & a_{22} & a_{23} \\
a_{31} & a_{32} & a_{33}
\end{pmatrix}
\]

\section{Tables}
Below is a sample table showing data:

\begin{table}[h]
\centering
\begin{tabular}{@{}lrrr@{}}
\toprule
\textbf{Category} & \textbf{Value 1} & \textbf{Value 2} & \textbf{Total} \\
\midrule
Item A & 10.5 & 20.3 & 30.8 \\
Item B & 15.2 & 25.1 & 40.3 \\
Item C & 22.1 & 12.7 & 34.8 \\
\bottomrule
\end{tabular}
\caption{Sample Data Table}
\label{tab:sample}
\end{table}

\section{Lists}
\subsection{Bullet Points}
\begin{itemize}
  \item First item
  \item Second item
  \item Third item with sub-items:
    \begin{itemize}
      \item Sub-item 1
      \item Sub-item 2
    \end{itemize}
\end{itemize}

\subsection{Numbered List}
\begin{enumerate}
  \item First step
  \item Second step
  \item Third step
\end{enumerate}

\section{Citations and References}
When you use the citation search feature, references will be automatically managed. For example, you might cite a paper like this: \cite{example2024}. All citations will be automatically added to your \texttt{cite.bib} file.

\section{Conclusion}
This template demonstrates basic \LaTeX{} features including mathematical formulas, tables, and lists. The bibliography system is now set up to work seamlessly with Bibby AI's citation management. Feel free to modify it for your needs!

\bibliographystyle{plain}
\bibliography{cite}

\end{document}\documentclass{article}
\usepackage[utf8]{inputenc}
\usepackage{booktabs}
\usepackage{amsmath}
\usepackage{natbib}

\title{My Document}
\author{Your Name}
\date{\today}

\begin{document}

\maketitle

\section{Introduction}
This is a sample document demonstrating various \LaTeX{} features including mathematical equations, tables, and formatting. You can cite papers using the citation search feature, and they will be automatically added to your \texttt{cite.bib} file.

\section{Mathematical Expressions}
Here are some example equations:

\subsection{Quadratic Formula}
The quadratic formula is given by:
\[x = \frac{-b \pm \sqrt{b^2 - 4ac}}{2a}\]

\subsection{Matrix Example}
A sample matrix:
\[
\begin{pmatrix}
a_{11} & a_{12} & a_{13} \\
a_{21} & a_{22} & a_{23} \\
a_{31} & a_{32} & a_{33}
\end{pmatrix}
\]

\section{Tables}
Below is a sample table showing data:

\begin{table}[h]
\centering
\begin{tabular}{@{}lrrr@{}}
\toprule
\textbf{Category} & \textbf{Value 1} & \textbf{Value 2} & \textbf{Total} \\
\midrule
Item A & 10.5 & 20.3 & 30.8 \\
Item B & 15.2 & 25.1 & 40.3 \\
Item C & 22.1 & 12.7 & 34.8 \\
\bottomrule
\end{tabular}
\caption{Sample Data Table}
\label{tab:sample}
\end{table}

\section{Lists}
\subsection{Bullet Points}
\begin{itemize}
  \item First item
  \item Second item
  \item Third item with sub-items:
    \begin{itemize}
      \item Sub-item 1
      \item Sub-item 2
    \end{itemize}
\end{itemize}

\subsection{Numbered List}
\begin{enumerate}
  \item First step
  \item Second step
  \item Third step
\end{enumerate}

\section{Citations and References}
When you use the citation search feature, references will be automatically managed. For example, you might cite a paper like this: \cite{example2024}. All citations will be automatically added to your \texttt{cite.bib} file.

\section{Conclusion}
This template demonstrates basic \LaTeX{} features including mathematical formulas, tables, and lists. The bibliography system is now set up to work seamlessly with Bibby AI's citation management. Feel free to modify it for your needs!

\bibliographystyle{plain}
\bibliography{cite}

\end{document}\documentclass{article}
\usepackage[utf8]{inputenc}
\usepackage{booktabs}
\usepackage{amsmath}
\usepackage{natbib}

\title{My Document}
\author{Your Name}
\date{\today}

\begin{document}

\maketitle

\section{Introduction}
This is a sample document demonstrating various \LaTeX{} features including mathematical equations, tables, and formatting. You can cite papers using the citation search feature, and they will be automatically added to your \texttt{cite.bib} file.

\section{Mathematical Expressions}
Here are some example equations:

\subsection{Quadratic Formula}
The quadratic formula is given by:
\[x = \frac{-b \pm \sqrt{b^2 - 4ac}}{2a}\]

\subsection{Matrix Example}
A sample matrix:
\[
\begin{pmatrix}
a_{11} & a_{12} & a_{13} \\
a_{21} & a_{22} & a_{23} \\
a_{31} & a_{32} & a_{33}
\end{pmatrix}
\]

\section{Tables}
Below is a sample table showing data:

\begin{table}[h]
\centering
\begin{tabular}{@{}lrrr@{}}
\toprule
\textbf{Category} & \textbf{Value 1} & \textbf{Value 2} & \textbf{Total} \\
\midrule
Item A & 10.5 & 20.3 & 30.8 \\
Item B & 15.2 & 25.1 & 40.3 \\
Item C & 22.1 & 12.7 & 34.8 \\
\bottomrule
\end{tabular}
\caption{Sample Data Table}
\label{tab:sample}
\end{table}

\section{Lists}
\subsection{Bullet Points}
\begin{itemize}
  \item First item
  \item Second item
  \item Third item with sub-items:
    \begin{itemize}
      \item Sub-item 1
      \item Sub-item 2
    \end{itemize}
\end{itemize}

\subsection{Numbered List}
\begin{enumerate}
  \item First step
  \item Second step
  \item Third step
\end{enumerate}

\section{Citations and References}
When you use the citation search feature, references will be automatically managed. For example, you might cite a paper like this: \cite{example2024}. All citations will be automatically added to your \texttt{cite.bib} file.

\section{Conclusion}
This template demonstrates basic \LaTeX{} features including mathematical formulas, tables, and lists. The bibliography system is now set up to work seamlessly with Bibby AI's citation management. Feel free to modify it for your needs!

\bibliographystyle{plain}
\bibliography{cite}

\end{document}\documentclass{article}
\usepackage[utf8]{inputenc}
\usepackage{booktabs}
\usepackage{amsmath}
\usepackage{natbib}

\title{My Document}
\author{Your Name}
\date{\today}

\begin{document}

\maketitle

\section{Introduction}
This is a sample document demonstrating various \LaTeX{} features including mathematical equations, tables, and formatting. You can cite papers using the citation search feature, and they will be automatically added to your \texttt{cite.bib} file.

\section{Mathematical Expressions}
Here are some example equations:

\subsection{Quadratic Formula}
The quadratic formula is given by:
\[x = \frac{-b \pm \sqrt{b^2 - 4ac}}{2a}\]

\subsection{Matrix Example}
A sample matrix:
\[
\begin{pmatrix}
a_{11} & a_{12} & a_{13} \\
a_{21} & a_{22} & a_{23} \\
a_{31} & a_{32} & a_{33}
\end{pmatrix}
\]

\section{Tables}
Below is a sample table showing data:

\begin{table}[h]
\centering
\begin{tabular}{@{}lrrr@{}}
\toprule
\textbf{Category} & \textbf{Value 1} & \textbf{Value 2} & \textbf{Total} \\
\midrule
Item A & 10.5 & 20.3 & 30.8 \\
Item B & 15.2 & 25.1 & 40.3 \\
Item C & 22.1 & 12.7 & 34.8 \\
\bottomrule
\end{tabular}
\caption{Sample Data Table}
\label{tab:sample}
\end{table}

\section{Lists}
\subsection{Bullet Points}
\begin{itemize}
  \item First item
  \item Second item
  \item Third item with sub-items:
    \begin{itemize}
      \item Sub-item 1
      \item Sub-item 2
    \end{itemize}
\end{itemize}

\subsection{Numbered List}
\begin{enumerate}
  \item First step
  \item Second step
  \item Third step
\end{enumerate}

\section{Citations and References}
When you use the citation search feature, references will be automatically managed. For example, you might cite a paper like this: \cite{example2024}. All citations will be automatically added to your \texttt{cite.bib} file.

\section{Conclusion}
This template demonstrates basic \LaTeX{} features including mathematical formulas, tables, and lists. The bibliography system is now set up to work seamlessly with Bibby AI's citation management. Feel free to modify it for your needs!

\bibliographystyle{plain}
\bibliography{cite}

\end{document}\documentclass{article}
\usepackage[utf8]{inputenc}
\usepackage{booktabs}
\usepackage{amsmath}
\usepackage{natbib}

\title{My Document}
\author{Your Name}
\date{\today}

\begin{document}

\maketitle

\section{Introduction}
This is a sample document demonstrating various \LaTeX{} features including mathematical equations, tables, and formatting. You can cite papers using the citation search feature, and they will be automatically added to your \texttt{cite.bib} file.

\section{Mathematical Expressions}
Here are some example equations:

\subsection{Quadratic Formula}
The quadratic formula is given by:
\[x = \frac{-b \pm \sqrt{b^2 - 4ac}}{2a}\]

\subsection{Matrix Example}
A sample matrix:
\[
\begin{pmatrix}
a_{11} & a_{12} & a_{13} \\
a_{21} & a_{22} & a_{23} \\
a_{31} & a_{32} & a_{33}
\end{pmatrix}
\]

\section{Tables}
Below is a sample table showing data:

\begin{table}[h]
\centering
\begin{tabular}{@{}lrrr@{}}
\toprule
\textbf{Category} & \textbf{Value 1} & \textbf{Value 2} & \textbf{Total} \\
\midrule
Item A & 10.5 & 20.3 & 30.8 \\
Item B & 15.2 & 25.1 & 40.3 \\
Item C & 22.1 & 12.7 & 34.8 \\
\bottomrule
\end{tabular}
\caption{Sample Data Table}
\label{tab:sample}
\end{table}

\section{Lists}
\subsection{Bullet Points}
\begin{itemize}
  \item First item
  \item Second item
  \item Third item with sub-items:
    \begin{itemize}
      \item Sub-item 1
      \item Sub-item 2
    \end{itemize}
\end{itemize}

\subsection{Numbered List}
\begin{enumerate}
  \item First step
  \item Second step
  \item Third step
\end{enumerate}

\section{Citations and References}
When you use the citation search feature, references will be automatically managed. For example, you might cite a paper like this: \cite{example2024}. All citations will be automatically added to your \texttt{cite.bib} file.

\section{Conclusion}
This template demonstrates basic \LaTeX{} features including mathematical formulas, tables, and lists. The bibliography system is now set up to work seamlessly with Bibby AI's citation management. Feel free to modify it for your needs!

\bibliographystyle{plain}
\bibliography{cite}

\end{document}\documentclass{article}
\usepackage[utf8]{inputenc}
\usepackage{booktabs}
\usepackage{amsmath}
\usepackage{natbib}

\title{My Document}
\author{Your Name}
\date{\today}

\begin{document}

\maketitle

\section{Introduction}
This is a sample document demonstrating various \LaTeX{} features including mathematical equations, tables, and formatting. You can cite papers using the citation search feature, and they will be automatically added to your \texttt{cite.bib} file.

\section{Mathematical Expressions}
Here are some example equations:

\subsection{Quadratic Formula}
The quadratic formula is given by:
\[x = \frac{-b \pm \sqrt{b^2 - 4ac}}{2a}\]

\subsection{Matrix Example}
A sample matrix:
\[
\begin{pmatrix}
a_{11} & a_{12} & a_{13} \\
a_{21} & a_{22} & a_{23} \\
a_{31} & a_{32} & a_{33}
\end{pmatrix}
\]

\section{Tables}
Below is a sample table showing data:

\begin{table}[h]
\centering
\begin{tabular}{@{}lrrr@{}}
\toprule
\textbf{Category} & \textbf{Value 1} & \textbf{Value 2} & \textbf{Total} \\
\midrule
Item A & 10.5 & 20.3 & 30.8 \\
Item B & 15.2 & 25.1 & 40.3 \\
Item C & 22.1 & 12.7 & 34.8 \\
\bottomrule
\end{tabular}
\caption{Sample Data Table}
\label{tab:sample}
\end{table}

\section{Lists}
\subsection{Bullet Points}
\begin{itemize}
  \item First item
  \item Second item
  \item Third item with sub-items:
    \begin{itemize}
      \item Sub-item 1
      \item Sub-item 2
    \end{itemize}
\end{itemize}

\subsection{Numbered List}
\begin{enumerate}
  \item First step
  \item Second step
  \item Third step
\end{enumerate}

\section{Citations and References}
When you use the citation search feature, references will be automatically managed. For example, you might cite a paper like this: \cite{example2024}. All citations will be automatically added to your \texttt{cite.bib} file.

\section{Conclusion}
This template demonstrates basic \LaTeX{} features including mathematical formulas, tables, and lists. The bibliography system is now set up to work seamlessly with Bibby AI's citation management. Feel free to modify it for your needs!

\bibliographystyle{plain}
\bibliography{cite}

\end{document}\documentclass{article}
\usepackage[utf8]{inputenc}
\usepackage{booktabs}
\usepackage{amsmath}
\usepackage{natbib}

\title{My Document}
\author{Your Name}
\date{\today}

\begin{document}

\maketitle

\section{Introduction}
This is a sample document demonstrating various \LaTeX{} features including mathematical equations, tables, and formatting. You can cite papers using the citation search feature, and they will be automatically added to your \texttt{cite.bib} file.

\section{Mathematical Expressions}
Here are some example equations:

\subsection{Quadratic Formula}
The quadratic formula is given by:
\[x = \frac{-b \pm \sqrt{b^2 - 4ac}}{2a}\]

\subsection{Matrix Example}
A sample matrix:
\[
\begin{pmatrix}
a_{11} & a_{12} & a_{13} \\
a_{21} & a_{22} & a_{23} \\
a_{31} & a_{32} & a_{33}
\end{pmatrix}
\]

\section{Tables}
Below is a sample table showing data:

\begin{table}[h]
\centering
\begin{tabular}{@{}lrrr@{}}
\toprule
\textbf{Category} & \textbf{Value 1} & \textbf{Value 2} & \textbf{Total} \\
\midrule
Item A & 10.5 & 20.3 & 30.8 \\
Item B & 15.2 & 25.1 & 40.3 \\
Item C & 22.1 & 12.7 & 34.8 \\
\bottomrule
\end{tabular}
\caption{Sample Data Table}
\label{tab:sample}
\end{table}

\section{Lists}
\subsection{Bullet Points}
\begin{itemize}
  \item First item
  \item Second item
  \item Third item with sub-items:
    \begin{itemize}
      \item Sub-item 1
      \item Sub-item 2
    \end{itemize}
\end{itemize}

\subsection{Numbered List}
\begin{enumerate}
  \item First step
  \item Second step
  \item Third step
\end{enumerate}

\section{Citations and References}
When you use the citation search feature, references will be automatically managed. For example, you might cite a paper like this: \cite{example2024}. All citations will be automatically added to your \texttt{cite.bib} file.

\section{Conclusion}
This template demonstrates basic \LaTeX{} features including mathematical formulas, tables, and lists. The bibliography system is now set up to work seamlessly with Bibby AI's citation management. Feel free to modify it for your needs!

\bibliographystyle{plain}
\bibliography{cite}

\end{document}\documentclass{article}
\usepackage[utf8]{inputenc}
\usepackage{booktabs}
\usepackage{amsmath}
\usepackage{natbib}

\title{My Document}
\author{Your Name}
\date{\today}

\begin{document}

\maketitle

\section{Introduction}
This is a sample document demonstrating various \LaTeX{} features including mathematical equations, tables, and formatting. You can cite papers using the citation search feature, and they will be automatically added to your \texttt{cite.bib} file.

\section{Mathematical Expressions}
Here are some example equations:

\subsection{Quadratic Formula}
The quadratic formula is given by:
\[x = \frac{-b \pm \sqrt{b^2 - 4ac}}{2a}\]

\subsection{Matrix Example}
A sample matrix:
\[
\begin{pmatrix}
a_{11} & a_{12} & a_{13} \\
a_{21} & a_{22} & a_{23} \\
a_{31} & a_{32} & a_{33}
\end{pmatrix}
\]

\section{Tables}
Below is a sample table showing data:

\begin{table}[h]
\centering
\begin{tabular}{@{}lrrr@{}}
\toprule
\textbf{Category} & \textbf{Value 1} & \textbf{Value 2} & \textbf{Total} \\
\midrule
Item A & 10.5 & 20.3 & 30.8 \\
Item B & 15.2 & 25.1 & 40.3 \\
Item C & 22.1 & 12.7 & 34.8 \\
\bottomrule
\end{tabular}
\caption{Sample Data Table}
\label{tab:sample}
\end{table}

\section{Lists}
\subsection{Bullet Points}
\begin{itemize}
  \item First item
  \item Second item
  \item Third item with sub-items:
    \begin{itemize}
      \item Sub-item 1
      \item Sub-item 2
    \end{itemize}
\end{itemize}

\subsection{Numbered List}
\begin{enumerate}
  \item First step
  \item Second step
  \item Third step
\end{enumerate}

\section{Citations and References}
When you use the citation search feature, references will be automatically managed. For example, you might cite a paper like this: \cite{example2024}. All citations will be automatically added to your \texttt{cite.bib} file.

\section{Conclusion}
This template demonstrates basic \LaTeX{} features including mathematical formulas, tables, and lists. The bibliography system is now set up to work seamlessly with Bibby AI's citation management. Feel free to modify it for your needs!

\bibliographystyle{plain}
\bibliography{cite}

\end{document}\documentclass{article}
\usepackage[utf8]{inputenc}
\usepackage{booktabs}
\usepackage{amsmath}
\usepackage{natbib}

\title{My Document}
\author{Your Name}
\date{\today}

\begin{document}

\maketitle

\section{Introduction}
This is a sample document demonstrating various \LaTeX{} features including mathematical equations, tables, and formatting. You can cite papers using the citation search feature, and they will be automatically added to your \texttt{cite.bib} file.

\section{Mathematical Expressions}
Here are some example equations:

\subsection{Quadratic Formula}
The quadratic formula is given by:
\[x = \frac{-b \pm \sqrt{b^2 - 4ac}}{2a}\]

\subsection{Matrix Example}
A sample matrix:
\[
\begin{pmatrix}
a_{11} & a_{12} & a_{13} \\
a_{21} & a_{22} & a_{23} \\
a_{31} & a_{32} & a_{33}
\end{pmatrix}
\]

\section{Tables}
Below is a sample table showing data:

\begin{table}[h]
\centering
\begin{tabular}{@{}lrrr@{}}
\toprule
\textbf{Category} & \textbf{Value 1} & \textbf{Value 2} & \textbf{Total} \\
\midrule
Item A & 10.5 & 20.3 & 30.8 \\
Item B & 15.2 & 25.1 & 40.3 \\
Item C & 22.1 & 12.7 & 34.8 \\
\bottomrule
\end{tabular}
\caption{Sample Data Table}
\label{tab:sample}
\end{table}

\section{Lists}
\subsection{Bullet Points}
\begin{itemize}
  \item First item
  \item Second item
  \item Third item with sub-items:
    \begin{itemize}
      \item Sub-item 1
      \item Sub-item 2
    \end{itemize}
\end{itemize}

\subsection{Numbered List}
\begin{enumerate}
  \item First step
  \item Second step
  \item Third step
\end{enumerate}

\section{Citations and References}
When you use the citation search feature, references will be automatically managed. For example, you might cite a paper like this: \cite{example2024}. All citations will be automatically added to your \texttt{cite.bib} file.

\section{Conclusion}
This template demonstrates basic \LaTeX{} features including mathematical formulas, tables, and lists. The bibliography system is now set up to work seamlessly with Bibby AI's citation management. Feel free to modify it for your needs!

\bibliographystyle{plain}
\bibliography{cite}

\end{document}\documentclass{article}
\usepackage[utf8]{inputenc}
\usepackage{booktabs}
\usepackage{amsmath}
\usepackage{natbib}

\title{My Document}
\author{Your Name}
\date{\today}

\begin{document}

\maketitle

\section{Introduction}
This is a sample document demonstrating various \LaTeX{} features including mathematical equations, tables, and formatting. You can cite papers using the citation search feature, and they will be automatically added to your \texttt{cite.bib} file.

\section{Mathematical Expressions}
Here are some example equations:

\subsection{Quadratic Formula}
The quadratic formula is given by:
\[x = \frac{-b \pm \sqrt{b^2 - 4ac}}{2a}\]

\subsection{Matrix Example}
A sample matrix:
\[
\begin{pmatrix}
a_{11} & a_{12} & a_{13} \\
a_{21} & a_{22} & a_{23} \\
a_{31} & a_{32} & a_{33}
\end{pmatrix}
\]

\section{Tables}
Below is a sample table showing data:

\begin{table}[h]
\centering
\begin{tabular}{@{}lrrr@{}}
\toprule
\textbf{Category} & \textbf{Value 1} & \textbf{Value 2} & \textbf{Total} \\
\midrule
Item A & 10.5 & 20.3 & 30.8 \\
Item B & 15.2 & 25.1 & 40.3 \\
Item C & 22.1 & 12.7 & 34.8 \\
\bottomrule
\end{tabular}
\caption{Sample Data Table}
\label{tab:sample}
\end{table}

\section{Lists}
\subsection{Bullet Points}
\begin{itemize}
  \item First item
  \item Second item
  \item Third item with sub-items:
    \begin{itemize}
      \item Sub-item 1
      \item Sub-item 2
    \end{itemize}
\end{itemize}

\subsection{Numbered List}
\begin{enumerate}
  \item First step
  \item Second step
  \item Third step
\end{enumerate}

\section{Citations and References}
When you use the citation search feature, references will be automatically managed. For example, you might cite a paper like this: \cite{example2024}. All citations will be automatically added to your \texttt{cite.bib} file.

\section{Conclusion}
This template demonstrates basic \LaTeX{} features including mathematical formulas, tables, and lists. The bibliography system is now set up to work seamlessly with Bibby AI's citation management. Feel free to modify it for your needs!

\bibliographystyle{plain}
\bibliography{cite}

\end{document}\documentclass{article}
\usepackage[utf8]{inputenc}
\usepackage{booktabs}
\usepackage{amsmath}
\usepackage{natbib}

\title{My Document}
\author{Your Name}
\date{\today}

\begin{document}

\maketitle

\section{Introduction}
This is a sample document demonstrating various \LaTeX{} features including mathematical equations, tables, and formatting. You can cite papers using the citation search feature, and they will be automatically added to your \texttt{cite.bib} file.

\section{Mathematical Expressions}
Here are some example equations:

\subsection{Quadratic Formula}
The quadratic formula is given by:
\[x = \frac{-b \pm \sqrt{b^2 - 4ac}}{2a}\]

\subsection{Matrix Example}
A sample matrix:
\[
\begin{pmatrix}
a_{11} & a_{12} & a_{13} \\
a_{21} & a_{22} & a_{23} \\
a_{31} & a_{32} & a_{33}
\end{pmatrix}
\]

\section{Tables}
Below is a sample table showing data:

\begin{table}[h]
\centering
\begin{tabular}{@{}lrrr@{}}
\toprule
\textbf{Category} & \textbf{Value 1} & \textbf{Value 2} & \textbf{Total} \\
\midrule
Item A & 10.5 & 20.3 & 30.8 \\
Item B & 15.2 & 25.1 & 40.3 \\
Item C & 22.1 & 12.7 & 34.8 \\
\bottomrule
\end{tabular}
\caption{Sample Data Table}
\label{tab:sample}
\end{table}

\section{Lists}
\subsection{Bullet Points}
\begin{itemize}
  \item First item
  \item Second item
  \item Third item with sub-items:
    \begin{itemize}
      \item Sub-item 1
      \item Sub-item 2
    \end{itemize}
\end{itemize}

\subsection{Numbered List}
\begin{enumerate}
  \item First step
  \item Second step
  \item Third step
\end{enumerate}

\section{Citations and References}
When you use the citation search feature, references will be automatically managed. For example, you might cite a paper like this: \cite{example2024}. All citations will be automatically added to your \texttt{cite.bib} file.

\section{Conclusion}
This template demonstrates basic \LaTeX{} features including mathematical formulas, tables, and lists. The bibliography system is now set up to work seamlessly with Bibby AI's citation management. Feel free to modify it for your needs!

\bibliographystyle{plain}
\bibliography{cite}

\end{document}\documentclass{article}
\usepackage[utf8]{inputenc}
\usepackage{booktabs}
\usepackage{amsmath}
\usepackage{natbib}

\title{My Document}
\author{Your Name}
\date{\today}

\begin{document}

\maketitle

\section{Introduction}
This is a sample document demonstrating various \LaTeX{} features including mathematical equations, tables, and formatting. You can cite papers using the citation search feature, and they will be automatically added to your \texttt{cite.bib} file.

\section{Mathematical Expressions}
Here are some example equations:

\subsection{Quadratic Formula}
The quadratic formula is given by:
\[x = \frac{-b \pm \sqrt{b^2 - 4ac}}{2a}\]

\subsection{Matrix Example}
A sample matrix:
\[
\begin{pmatrix}
a_{11} & a_{12} & a_{13} \\
a_{21} & a_{22} & a_{23} \\
a_{31} & a_{32} & a_{33}
\end{pmatrix}
\]

\section{Tables}
Below is a sample table showing data:

\begin{table}[h]
\centering
\begin{tabular}{@{}lrrr@{}}
\toprule
\textbf{Category} & \textbf{Value 1} & \textbf{Value 2} & \textbf{Total} \\
\midrule
Item A & 10.5 & 20.3 & 30.8 \\
Item B & 15.2 & 25.1 & 40.3 \\
Item C & 22.1 & 12.7 & 34.8 \\
\bottomrule
\end{tabular}
\caption{Sample Data Table}
\label{tab:sample}
\end{table}

\section{Lists}
\subsection{Bullet Points}
\begin{itemize}
  \item First item
  \item Second item
  \item Third item with sub-items:
    \begin{itemize}
      \item Sub-item 1
      \item Sub-item 2
    \end{itemize}
\end{itemize}

\subsection{Numbered List}
\begin{enumerate}
  \item First step
  \item Second step
  \item Third step
\end{enumerate}

\section{Citations and References}
When you use the citation search feature, references will be automatically managed. For example, you might cite a paper like this: \cite{example2024}. All citations will be automatically added to your \texttt{cite.bib} file.

\section{Conclusion}
This template demonstrates basic \LaTeX{} features including mathematical formulas, tables, and lists. The bibliography system is now set up to work seamlessly with Bibby AI's citation management. Feel free to modify it for your needs!

\bibliographystyle{plain}
\bibliography{cite}

\end{document}\documentclass{article}
\usepackage[utf8]{inputenc}
\usepackage{booktabs}
\usepackage{amsmath}
\usepackage{natbib}

\title{My Document}
\author{Your Name}
\date{\today}

\begin{document}

\maketitle

\section{Introduction}
This is a sample document demonstrating various \LaTeX{} features including mathematical equations, tables, and formatting. You can cite papers using the citation search feature, and they will be automatically added to your \texttt{cite.bib} file.

\section{Mathematical Expressions}
Here are some example equations:

\subsection{Quadratic Formula}
The quadratic formula is given by:
\[x = \frac{-b \pm \sqrt{b^2 - 4ac}}{2a}\]

\subsection{Matrix Example}
A sample matrix:
\[
\begin{pmatrix}
a_{11} & a_{12} & a_{13} \\
a_{21} & a_{22} & a_{23} \\
a_{31} & a_{32} & a_{33}
\end{pmatrix}
\]

\section{Tables}
Below is a sample table showing data:

\begin{table}[h]
\centering
\begin{tabular}{@{}lrrr@{}}
\toprule
\textbf{Category} & \textbf{Value 1} & \textbf{Value 2} & \textbf{Total} \\
\midrule
Item A & 10.5 & 20.3 & 30.8 \\
Item B & 15.2 & 25.1 & 40.3 \\
Item C & 22.1 & 12.7 & 34.8 \\
\bottomrule
\end{tabular}
\caption{Sample Data Table}
\label{tab:sample}
\end{table}

\section{Lists}
\subsection{Bullet Points}
\begin{itemize}
  \item First item
  \item Second item
  \item Third item with sub-items:
    \begin{itemize}
      \item Sub-item 1
      \item Sub-item 2
    \end{itemize}
\end{itemize}

\subsection{Numbered List}
\begin{enumerate}
  \item First step
  \item Second step
  \item Third step
\end{enumerate}

\section{Citations and References}
When you use the citation search feature, references will be automatically managed. For example, you might cite a paper like this: \cite{example2024}. All citations will be automatically added to your \texttt{cite.bib} file.

\section{Conclusion}
This template demonstrates basic \LaTeX{} features including mathematical formulas, tables, and lists. The bibliography system is now set up to work seamlessly with Bibby AI's citation management. Feel free to modify it for your needs!

\bibliographystyle{plain}
\bibliography{cite}

\end{document}\documentclass{article}
\usepackage[utf8]{inputenc}
\usepackage{booktabs}
\usepackage{amsmath}
\usepackage{natbib}

\title{My Document}
\author{Your Name}
\date{\today}

\begin{document}

\maketitle

\section{Introduction}
This is a sample document demonstrating various \LaTeX{} features including mathematical equations, tables, and formatting. You can cite papers using the citation search feature, and they will be automatically added to your \texttt{cite.bib} file.

\section{Mathematical Expressions}
Here are some example equations:

\subsection{Quadratic Formula}
The quadratic formula is given by:
\[x = \frac{-b \pm \sqrt{b^2 - 4ac}}{2a}\]

\subsection{Matrix Example}
A sample matrix:
\[
\begin{pmatrix}
a_{11} & a_{12} & a_{13} \\
a_{21} & a_{22} & a_{23} \\
a_{31} & a_{32} & a_{33}
\end{pmatrix}
\]

\section{Tables}
Below is a sample table showing data:

\begin{table}[h]
\centering
\begin{tabular}{@{}lrrr@{}}
\toprule
\textbf{Category} & \textbf{Value 1} & \textbf{Value 2} & \textbf{Total} \\
\midrule
Item A & 10.5 & 20.3 & 30.8 \\
Item B & 15.2 & 25.1 & 40.3 \\
Item C & 22.1 & 12.7 & 34.8 \\
\bottomrule
\end{tabular}
\caption{Sample Data Table}
\label{tab:sample}
\end{table}

\section{Lists}
\subsection{Bullet Points}
\begin{itemize}
  \item First item
  \item Second item
  \item Third item with sub-items:
    \begin{itemize}
      \item Sub-item 1
      \item Sub-item 2
    \end{itemize}
\end{itemize}

\subsection{Numbered List}
\begin{enumerate}
  \item First step
  \item Second step
  \item Third step
\end{enumerate}

\section{Citations and References}
When you use the citation search feature, references will be automatically managed. For example, you might cite a paper like this: \cite{example2024}. All citations will be automatically added to your \texttt{cite.bib} file.

\section{Conclusion}
This template demonstrates basic \LaTeX{} features including mathematical formulas, tables, and lists. The bibliography system is now set up to work seamlessly with Bibby AI's citation management. Feel free to modify it for your needs!

\bibliographystyle{plain}
\bibliography{cite}

\end{document}\documentclass{article}
\usepackage[utf8]{inputenc}
\usepackage{booktabs}
\usepackage{amsmath}
\usepackage{natbib}

\title{My Document}
\author{Your Name}
\date{\today}

\begin{document}

\maketitle

\section{Introduction}
This is a sample document demonstrating various \LaTeX{} features including mathematical equations, tables, and formatting. You can cite papers using the citation search feature, and they will be automatically added to your \texttt{cite.bib} file.

\section{Mathematical Expressions}
Here are some example equations:

\subsection{Quadratic Formula}
The quadratic formula is given by:
\[x = \frac{-b \pm \sqrt{b^2 - 4ac}}{2a}\]

\subsection{Matrix Example}
A sample matrix:
\[
\begin{pmatrix}
a_{11} & a_{12} & a_{13} \\
a_{21} & a_{22} & a_{23} \\
a_{31} & a_{32} & a_{33}
\end{pmatrix}
\]

\section{Tables}
Below is a sample table showing data:

\begin{table}[h]
\centering
\begin{tabular}{@{}lrrr@{}}
\toprule
\textbf{Category} & \textbf{Value 1} & \textbf{Value 2} & \textbf{Total} \\
\midrule
Item A & 10.5 & 20.3 & 30.8 \\
Item B & 15.2 & 25.1 & 40.3 \\
Item C & 22.1 & 12.7 & 34.8 \\
\bottomrule
\end{tabular}
\caption{Sample Data Table}
\label{tab:sample}
\end{table}

\section{Lists}
\subsection{Bullet Points}
\begin{itemize}
  \item First item
  \item Second item
  \item Third item with sub-items:
    \begin{itemize}
      \item Sub-item 1
      \item Sub-item 2
    \end{itemize}
\end{itemize}

\subsection{Numbered List}
\begin{enumerate}
  \item First step
  \item Second step
  \item Third step
\end{enumerate}

\section{Citations and References}
When you use the citation search feature, references will be automatically managed. For example, you might cite a paper like this: \cite{example2024}. All citations will be automatically added to your \texttt{cite.bib} file.

\section{Conclusion}
This template demonstrates basic \LaTeX{} features including mathematical formulas, tables, and lists. The bibliography system is now set up to work seamlessly with Bibby AI's citation management. Feel free to modify it for your needs!

\bibliographystyle{plain}
\bibliography{cite}

\end{document}\documentclass{article}
\usepackage[utf8]{inputenc}
\usepackage{booktabs}
\usepackage{amsmath}
\usepackage{natbib}

\title{My Document}
\author{Your Name}
\date{\today}

\begin{document}

\maketitle

\section{Introduction}
This is a sample document demonstrating various \LaTeX{} features including mathematical equations, tables, and formatting. You can cite papers using the citation search feature, and they will be automatically added to your \texttt{cite.bib} file.

\section{Mathematical Expressions}
Here are some example equations:

\subsection{Quadratic Formula}
The quadratic formula is given by:
\[x = \frac{-b \pm \sqrt{b^2 - 4ac}}{2a}\]

\subsection{Matrix Example}
A sample matrix:
\[
\begin{pmatrix}
a_{11} & a_{12} & a_{13} \\
a_{21} & a_{22} & a_{23} \\
a_{31} & a_{32} & a_{33}
\end{pmatrix}
\]

\section{Tables}
Below is a sample table showing data:

\begin{table}[h]
\centering
\begin{tabular}{@{}lrrr@{}}
\toprule
\textbf{Category} & \textbf{Value 1} & \textbf{Value 2} & \textbf{Total} \\
\midrule
Item A & 10.5 & 20.3 & 30.8 \\
Item B & 15.2 & 25.1 & 40.3 \\
Item C & 22.1 & 12.7 & 34.8 \\
\bottomrule
\end{tabular}
\caption{Sample Data Table}
\label{tab:sample}
\end{table}

\section{Lists}
\subsection{Bullet Points}
\begin{itemize}
  \item First item
  \item Second item
  \item Third item with sub-items:
    \begin{itemize}
      \item Sub-item 1
      \item Sub-item 2
    \end{itemize}
\end{itemize}

\subsection{Numbered List}
\begin{enumerate}
  \item First step
  \item Second step
  \item Third step
\end{enumerate}

\section{Citations and References}
When you use the citation search feature, references will be automatically managed. For example, you might cite a paper like this: \cite{example2024}. All citations will be automatically added to your \texttt{cite.bib} file.

\section{Conclusion}
This template demonstrates basic \LaTeX{} features including mathematical formulas, tables, and lists. The bibliography system is now set up to work seamlessly with Bibby AI's citation management. Feel free to modify it for your needs!

\bibliographystyle{plain}
\bibliography{cite}

\end{document}\documentclass{article}
\usepackage[utf8]{inputenc}
\usepackage{booktabs}
\usepackage{amsmath}
\usepackage{natbib}

\title{My Document}
\author{Your Name}
\date{\today}

\begin{document}

\maketitle

\section{Introduction}
This is a sample document demonstrating various \LaTeX{} features including mathematical equations, tables, and formatting. You can cite papers using the citation search feature, and they will be automatically added to your \texttt{cite.bib} file.

\section{Mathematical Expressions}
Here are some example equations:

\subsection{Quadratic Formula}
The quadratic formula is given by:
\[x = \frac{-b \pm \sqrt{b^2 - 4ac}}{2a}\]

\subsection{Matrix Example}
A sample matrix:
\[
\begin{pmatrix}
a_{11} & a_{12} & a_{13} \\
a_{21} & a_{22} & a_{23} \\
a_{31} & a_{32} & a_{33}
\end{pmatrix}
\]

\section{Tables}
Below is a sample table showing data:

\begin{table}[h]
\centering
\begin{tabular}{@{}lrrr@{}}
\toprule
\textbf{Category} & \textbf{Value 1} & \textbf{Value 2} & \textbf{Total} \\
\midrule
Item A & 10.5 & 20.3 & 30.8 \\
Item B & 15.2 & 25.1 & 40.3 \\
Item C & 22.1 & 12.7 & 34.8 \\
\bottomrule
\end{tabular}
\caption{Sample Data Table}
\label{tab:sample}
\end{table}

\section{Lists}
\subsection{Bullet Points}
\begin{itemize}
  \item First item
  \item Second item
  \item Third item with sub-items:
    \begin{itemize}
      \item Sub-item 1
      \item Sub-item 2
    \end{itemize}
\end{itemize}

\subsection{Numbered List}
\begin{enumerate}
  \item First step
  \item Second step
  \item Third step
\end{enumerate}

\section{Citations and References}
When you use the citation search feature, references will be automatically managed. For example, you might cite a paper like this: \cite{example2024}. All citations will be automatically added to your \texttt{cite.bib} file.

\section{Conclusion}
This template demonstrates basic \LaTeX{} features including mathematical formulas, tables, and lists. The bibliography system is now set up to work seamlessly with Bibby AI's citation management. Feel free to modify it for your needs!

\bibliographystyle{plain}
\bibliography{cite}

\end{document}\documentclass{article}
\usepackage[utf8]{inputenc}
\usepackage{booktabs}
\usepackage{amsmath}
\usepackage{natbib}

\title{My Document}
\author{Your Name}
\date{\today}

\begin{document}

\maketitle

\section{Introduction}
This is a sample document demonstrating various \LaTeX{} features including mathematical equations, tables, and formatting. You can cite papers using the citation search feature, and they will be automatically added to your \texttt{cite.bib} file.

\section{Mathematical Expressions}
Here are some example equations:

\subsection{Quadratic Formula}
The quadratic formula is given by:
\[x = \frac{-b \pm \sqrt{b^2 - 4ac}}{2a}\]

\subsection{Matrix Example}
A sample matrix:
\[
\begin{pmatrix}
a_{11} & a_{12} & a_{13} \\
a_{21} & a_{22} & a_{23} \\
a_{31} & a_{32} & a_{33}
\end{pmatrix}
\]

\section{Tables}
Below is a sample table showing data:

\begin{table}[h]
\centering
\begin{tabular}{@{}lrrr@{}}
\toprule
\textbf{Category} & \textbf{Value 1} & \textbf{Value 2} & \textbf{Total} \\
\midrule
Item A & 10.5 & 20.3 & 30.8 \\
Item B & 15.2 & 25.1 & 40.3 \\
Item C & 22.1 & 12.7 & 34.8 \\
\bottomrule
\end{tabular}
\caption{Sample Data Table}
\label{tab:sample}
\end{table}

\section{Lists}
\subsection{Bullet Points}
\begin{itemize}
  \item First item
  \item Second item
  \item Third item with sub-items:
    \begin{itemize}
      \item Sub-item 1
      \item Sub-item 2
    \end{itemize}
\end{itemize}

\subsection{Numbered List}
\begin{enumerate}
  \item First step
  \item Second step
  \item Third step
\end{enumerate}

\section{Citations and References}
When you use the citation search feature, references will be automatically managed. For example, you might cite a paper like this: \cite{example2024}. All citations will be automatically added to your \texttt{cite.bib} file.

\section{Conclusion}
This template demonstrates basic \LaTeX{} features including mathematical formulas, tables, and lists. The bibliography system is now set up to work seamlessly with Bibby AI's citation management. Feel free to modify it for your needs!

\bibliographystyle{plain}
\bibliography{cite}

\end{document}\documentclass{article}
\usepackage[utf8]{inputenc}
\usepackage{booktabs}
\usepackage{amsmath}
\usepackage{natbib}

\title{My Document}
\author{Your Name}
\date{\today}

\begin{document}

\maketitle

\section{Introduction}
This is a sample document demonstrating various \LaTeX{} features including mathematical equations, tables, and formatting. You can cite papers using the citation search feature, and they will be automatically added to your \texttt{cite.bib} file.

\section{Mathematical Expressions}
Here are some example equations:

\subsection{Quadratic Formula}
The quadratic formula is given by:
\[x = \frac{-b \pm \sqrt{b^2 - 4ac}}{2a}\]

\subsection{Matrix Example}
A sample matrix:
\[
\begin{pmatrix}
a_{11} & a_{12} & a_{13} \\
a_{21} & a_{22} & a_{23} \\
a_{31} & a_{32} & a_{33}
\end{pmatrix}
\]

\section{Tables}
Below is a sample table showing data:

\begin{table}[h]
\centering
\begin{tabular}{@{}lrrr@{}}
\toprule
\textbf{Category} & \textbf{Value 1} & \textbf{Value 2} & \textbf{Total} \\
\midrule
Item A & 10.5 & 20.3 & 30.8 \\
Item B & 15.2 & 25.1 & 40.3 \\
Item C & 22.1 & 12.7 & 34.8 \\
\bottomrule
\end{tabular}
\caption{Sample Data Table}
\label{tab:sample}
\end{table}

\section{Lists}
\subsection{Bullet Points}
\begin{itemize}
  \item First item
  \item Second item
  \item Third item with sub-items:
    \begin{itemize}
      \item Sub-item 1
      \item Sub-item 2
    \end{itemize}
\end{itemize}

\subsection{Numbered List}
\begin{enumerate}
  \item First step
  \item Second step
  \item Third step
\end{enumerate}

\section{Citations and References}
When you use the citation search feature, references will be automatically managed. For example, you might cite a paper like this: \cite{example2024}. All citations will be automatically added to your \texttt{cite.bib} file.

\section{Conclusion}
This template demonstrates basic \LaTeX{} features including mathematical formulas, tables, and lists. The bibliography system is now set up to work seamlessly with Bibby AI's citation management. Feel free to modify it for your needs!

\bibliographystyle{plain}
\bibliography{cite}

\end{document}\documentclass{article}
\usepackage[utf8]{inputenc}
\usepackage{booktabs}
\usepackage{amsmath}
\usepackage{natbib}

\title{My Document}
\author{Your Name}
\date{\today}

\begin{document}

\maketitle

\section{Introduction}
This is a sample document demonstrating various \LaTeX{} features including mathematical equations, tables, and formatting. You can cite papers using the citation search feature, and they will be automatically added to your \texttt{cite.bib} file.

\section{Mathematical Expressions}
Here are some example equations:

\subsection{Quadratic Formula}
The quadratic formula is given by:
\[x = \frac{-b \pm \sqrt{b^2 - 4ac}}{2a}\]

\subsection{Matrix Example}
A sample matrix:
\[
\begin{pmatrix}
a_{11} & a_{12} & a_{13} \\
a_{21} & a_{22} & a_{23} \\
a_{31} & a_{32} & a_{33}
\end{pmatrix}
\]

\section{Tables}
Below is a sample table showing data:

\begin{table}[h]
\centering
\begin{tabular}{@{}lrrr@{}}
\toprule
\textbf{Category} & \textbf{Value 1} & \textbf{Value 2} & \textbf{Total} \\
\midrule
Item A & 10.5 & 20.3 & 30.8 \\
Item B & 15.2 & 25.1 & 40.3 \\
Item C & 22.1 & 12.7 & 34.8 \\
\bottomrule
\end{tabular}
\caption{Sample Data Table}
\label{tab:sample}
\end{table}

\section{Lists}
\subsection{Bullet Points}
\begin{itemize}
  \item First item
  \item Second item
  \item Third item with sub-items:
    \begin{itemize}
      \item Sub-item 1
      \item Sub-item 2
    \end{itemize}
\end{itemize}

\subsection{Numbered List}
\begin{enumerate}
  \item First step
  \item Second step
  \item Third step
\end{enumerate}

\section{Citations and References}
When you use the citation search feature, references will be automatically managed. For example, you might cite a paper like this: \cite{example2024}. All citations will be automatically added to your \texttt{cite.bib} file.

\section{Conclusion}
This template demonstrates basic \LaTeX{} features including mathematical formulas, tables, and lists. The bibliography system is now set up to work seamlessly with Bibby AI's citation management. Feel free to modify it for your needs!

\bibliographystyle{plain}
\bibliography{cite}

\end{document}\documentclass{article}
\usepackage[utf8]{inputenc}
\usepackage{booktabs}
\usepackage{amsmath}
\usepackage{natbib}

\title{My Document}
\author{Your Name}
\date{\today}

\begin{document}

\maketitle

\section{Introduction}
This is a sample document demonstrating various \LaTeX{} features including mathematical equations, tables, and formatting. You can cite papers using the citation search feature, and they will be automatically added to your \texttt{cite.bib} file.

\section{Mathematical Expressions}
Here are some example equations:

\subsection{Quadratic Formula}
The quadratic formula is given by:
\[x = \frac{-b \pm \sqrt{b^2 - 4ac}}{2a}\]

\subsection{Matrix Example}
A sample matrix:
\[
\begin{pmatrix}
a_{11} & a_{12} & a_{13} \\
a_{21} & a_{22} & a_{23} \\
a_{31} & a_{32} & a_{33}
\end{pmatrix}
\]

\section{Tables}
Below is a sample table showing data:

\begin{table}[h]
\centering
\begin{tabular}{@{}lrrr@{}}
\toprule
\textbf{Category} & \textbf{Value 1} & \textbf{Value 2} & \textbf{Total} \\
\midrule
Item A & 10.5 & 20.3 & 30.8 \\
Item B & 15.2 & 25.1 & 40.3 \\
Item C & 22.1 & 12.7 & 34.8 \\
\bottomrule
\end{tabular}
\caption{Sample Data Table}
\label{tab:sample}
\end{table}

\section{Lists}
\subsection{Bullet Points}
\begin{itemize}
  \item First item
  \item Second item
  \item Third item with sub-items:
    \begin{itemize}
      \item Sub-item 1
      \item Sub-item 2
    \end{itemize}
\end{itemize}

\subsection{Numbered List}
\begin{enumerate}
  \item First step
  \item Second step
  \item Third step
\end{enumerate}

\section{Citations and References}
When you use the citation search feature, references will be automatically managed. For example, you might cite a paper like this: \cite{example2024}. All citations will be automatically added to your \texttt{cite.bib} file.

\section{Conclusion}
This template demonstrates basic \LaTeX{} features including mathematical formulas, tables, and lists. The bibliography system is now set up to work seamlessly with Bibby AI's citation management. Feel free to modify it for your needs!

\bibliographystyle{plain}
\bibliography{cite}

\end{document}\documentclass{article}
\usepackage[utf8]{inputenc}
\usepackage{booktabs}
\usepackage{amsmath}
\usepackage{natbib}

\title{My Document}
\author{Your Name}
\date{\today}

\begin{document}

\maketitle

\section{Introduction}
This is a sample document demonstrating various \LaTeX{} features including mathematical equations, tables, and formatting. You can cite papers using the citation search feature, and they will be automatically added to your \texttt{cite.bib} file.

\section{Mathematical Expressions}
Here are some example equations:

\subsection{Quadratic Formula}
The quadratic formula is given by:
\[x = \frac{-b \pm \sqrt{b^2 - 4ac}}{2a}\]

\subsection{Matrix Example}
A sample matrix:
\[
\begin{pmatrix}
a_{11} & a_{12} & a_{13} \\
a_{21} & a_{22} & a_{23} \\
a_{31} & a_{32} & a_{33}
\end{pmatrix}
\]

\section{Tables}
Below is a sample table showing data:

\begin{table}[h]
\centering
\begin{tabular}{@{}lrrr@{}}
\toprule
\textbf{Category} & \textbf{Value 1} & \textbf{Value 2} & \textbf{Total} \\
\midrule
Item A & 10.5 & 20.3 & 30.8 \\
Item B & 15.2 & 25.1 & 40.3 \\
Item C & 22.1 & 12.7 & 34.8 \\
\bottomrule
\end{tabular}
\caption{Sample Data Table}
\label{tab:sample}
\end{table}

\section{Lists}
\subsection{Bullet Points}
\begin{itemize}
  \item First item
  \item Second item
  \item Third item with sub-items:
    \begin{itemize}
      \item Sub-item 1
      \item Sub-item 2
    \end{itemize}
\end{itemize}

\subsection{Numbered List}
\begin{enumerate}
  \item First step
  \item Second step
  \item Third step
\end{enumerate}

\section{Citations and References}
When you use the citation search feature, references will be automatically managed. For example, you might cite a paper like this: \cite{example2024}. All citations will be automatically added to your \texttt{cite.bib} file.

\section{Conclusion}
This template demonstrates basic \LaTeX{} features including mathematical formulas, tables, and lists. The bibliography system is now set up to work seamlessly with Bibby AI's citation management. Feel free to modify it for your needs!

\bibliographystyle{plain}
\bibliography{cite}

\end{document}\documentclass{article}
\usepackage[utf8]{inputenc}
\usepackage{booktabs}
\usepackage{amsmath}
\usepackage{natbib}

\title{My Document}
\author{Your Name}
\date{\today}

\begin{document}

\maketitle

\section{Introduction}
This is a sample document demonstrating various \LaTeX{} features including mathematical equations, tables, and formatting. You can cite papers using the citation search feature, and they will be automatically added to your \texttt{cite.bib} file.

\section{Mathematical Expressions}
Here are some example equations:

\subsection{Quadratic Formula}
The quadratic formula is given by:
\[x = \frac{-b \pm \sqrt{b^2 - 4ac}}{2a}\]

\subsection{Matrix Example}
A sample matrix:
\[
\begin{pmatrix}
a_{11} & a_{12} & a_{13} \\
a_{21} & a_{22} & a_{23} \\
a_{31} & a_{32} & a_{33}
\end{pmatrix}
\]

\section{Tables}
Below is a sample table showing data:

\begin{table}[h]
\centering
\begin{tabular}{@{}lrrr@{}}
\toprule
\textbf{Category} & \textbf{Value 1} & \textbf{Value 2} & \textbf{Total} \\
\midrule
Item A & 10.5 & 20.3 & 30.8 \\
Item B & 15.2 & 25.1 & 40.3 \\
Item C & 22.1 & 12.7 & 34.8 \\
\bottomrule
\end{tabular}
\caption{Sample Data Table}
\label{tab:sample}
\end{table}

\section{Lists}
\subsection{Bullet Points}
\begin{itemize}
  \item First item
  \item Second item
  \item Third item with sub-items:
    \begin{itemize}
      \item Sub-item 1
      \item Sub-item 2
    \end{itemize}
\end{itemize}

\subsection{Numbered List}
\begin{enumerate}
  \item First step
  \item Second step
  \item Third step
\end{enumerate}

\section{Citations and References}
When you use the citation search feature, references will be automatically managed. For example, you might cite a paper like this: \cite{example2024}. All citations will be automatically added to your \texttt{cite.bib} file.

\section{Conclusion}
This template demonstrates basic \LaTeX{} features including mathematical formulas, tables, and lists. The bibliography system is now set up to work seamlessly with Bibby AI's citation management. Feel free to modify it for your needs!

\bibliographystyle{plain}
\bibliography{cite}

\end{document}\documentclass{article}
\usepackage[utf8]{inputenc}
\usepackage{booktabs}
\usepackage{amsmath}
\usepackage{natbib}

\title{My Document}
\author{Your Name}
\date{\today}

\begin{document}

\maketitle

\section{Introduction}
This is a sample document demonstrating various \LaTeX{} features including mathematical equations, tables, and formatting. You can cite papers using the citation search feature, and they will be automatically added to your \texttt{cite.bib} file.

\section{Mathematical Expressions}
Here are some example equations:

\subsection{Quadratic Formula}
The quadratic formula is given by:
\[x = \frac{-b \pm \sqrt{b^2 - 4ac}}{2a}\]

\subsection{Matrix Example}
A sample matrix:
\[
\begin{pmatrix}
a_{11} & a_{12} & a_{13} \\
a_{21} & a_{22} & a_{23} \\
a_{31} & a_{32} & a_{33}
\end{pmatrix}
\]

\section{Tables}
Below is a sample table showing data:

\begin{table}[h]
\centering
\begin{tabular}{@{}lrrr@{}}
\toprule
\textbf{Category} & \textbf{Value 1} & \textbf{Value 2} & \textbf{Total} \\
\midrule
Item A & 10.5 & 20.3 & 30.8 \\
Item B & 15.2 & 25.1 & 40.3 \\
Item C & 22.1 & 12.7 & 34.8 \\
\bottomrule
\end{tabular}
\caption{Sample Data Table}
\label{tab:sample}
\end{table}

\section{Lists}
\subsection{Bullet Points}
\begin{itemize}
  \item First item
  \item Second item
  \item Third item with sub-items:
    \begin{itemize}
      \item Sub-item 1
      \item Sub-item 2
    \end{itemize}
\end{itemize}

\subsection{Numbered List}
\begin{enumerate}
  \item First step
  \item Second step
  \item Third step
\end{enumerate}

\section{Citations and References}
When you use the citation search feature, references will be automatically managed. For example, you might cite a paper like this: \cite{example2024}. All citations will be automatically added to your \texttt{cite.bib} file.

\section{Conclusion}
This template demonstrates basic \LaTeX{} features including mathematical formulas, tables, and lists. The bibliography system is now set up to work seamlessly with Bibby AI's citation management. Feel free to modify it for your needs!

\bibliographystyle{plain}
\bibliography{cite}

\end{document}\documentclass{article}
\usepackage[utf8]{inputenc}
\usepackage{booktabs}
\usepackage{amsmath}
\usepackage{natbib}

\title{My Document}
\author{Your Name}
\date{\today}

\begin{document}

\maketitle

\section{Introduction}
This is a sample document demonstrating various \LaTeX{} features including mathematical equations, tables, and formatting. You can cite papers using the citation search feature, and they will be automatically added to your \texttt{cite.bib} file.

\section{Mathematical Expressions}
Here are some example equations:

\subsection{Quadratic Formula}
The quadratic formula is given by:
\[x = \frac{-b \pm \sqrt{b^2 - 4ac}}{2a}\]

\subsection{Matrix Example}
A sample matrix:
\[
\begin{pmatrix}
a_{11} & a_{12} & a_{13} \\
a_{21} & a_{22} & a_{23} \\
a_{31} & a_{32} & a_{33}
\end{pmatrix}
\]

\section{Tables}
Below is a sample table showing data:

\begin{table}[h]
\centering
\begin{tabular}{@{}lrrr@{}}
\toprule
\textbf{Category} & \textbf{Value 1} & \textbf{Value 2} & \textbf{Total} \\
\midrule
Item A & 10.5 & 20.3 & 30.8 \\
Item B & 15.2 & 25.1 & 40.3 \\
Item C & 22.1 & 12.7 & 34.8 \\
\bottomrule
\end{tabular}
\caption{Sample Data Table}
\label{tab:sample}
\end{table}

\section{Lists}
\subsection{Bullet Points}
\begin{itemize}
  \item First item
  \item Second item
  \item Third item with sub-items:
    \begin{itemize}
      \item Sub-item 1
      \item Sub-item 2
    \end{itemize}
\end{itemize}

\subsection{Numbered List}
\begin{enumerate}
  \item First step
  \item Second step
  \item Third step
\end{enumerate}

\section{Citations and References}
When you use the citation search feature, references will be automatically managed. For example, you might cite a paper like this: \cite{example2024}. All citations will be automatically added to your \texttt{cite.bib} file.

\section{Conclusion}
This template demonstrates basic \LaTeX{} features including mathematical formulas, tables, and lists. The bibliography system is now set up to work seamlessly with Bibby AI's citation management. Feel free to modify it for your needs!

\bibliographystyle{plain}
\bibliography{cite}

\end{document}\documentclass{article}
\usepackage[utf8]{inputenc}
\usepackage{booktabs}
\usepackage{amsmath}
\usepackage{natbib}

\title{My Document}
\author{Your Name}
\date{\today}

\begin{document}

\maketitle

\section{Introduction}
This is a sample document demonstrating various \LaTeX{} features including mathematical equations, tables, and formatting. You can cite papers using the citation search feature, and they will be automatically added to your \texttt{cite.bib} file.

\section{Mathematical Expressions}
Here are some example equations:

\subsection{Quadratic Formula}
The quadratic formula is given by:
\[x = \frac{-b \pm \sqrt{b^2 - 4ac}}{2a}\]

\subsection{Matrix Example}
A sample matrix:
\[
\begin{pmatrix}
a_{11} & a_{12} & a_{13} \\
a_{21} & a_{22} & a_{23} \\
a_{31} & a_{32} & a_{33}
\end{pmatrix}
\]

\section{Tables}
Below is a sample table showing data:

\begin{table}[h]
\centering
\begin{tabular}{@{}lrrr@{}}
\toprule
\textbf{Category} & \textbf{Value 1} & \textbf{Value 2} & \textbf{Total} \\
\midrule
Item A & 10.5 & 20.3 & 30.8 \\
Item B & 15.2 & 25.1 & 40.3 \\
Item C & 22.1 & 12.7 & 34.8 \\
\bottomrule
\end{tabular}
\caption{Sample Data Table}
\label{tab:sample}
\end{table}

\section{Lists}
\subsection{Bullet Points}
\begin{itemize}
  \item First item
  \item Second item
  \item Third item with sub-items:
    \begin{itemize}
      \item Sub-item 1
      \item Sub-item 2
    \end{itemize}
\end{itemize}

\subsection{Numbered List}
\begin{enumerate}
  \item First step
  \item Second step
  \item Third step
\end{enumerate}

\section{Citations and References}
When you use the citation search feature, references will be automatically managed. For example, you might cite a paper like this: \cite{example2024}. All citations will be automatically added to your \texttt{cite.bib} file.

\section{Conclusion}
This template demonstrates basic \LaTeX{} features including mathematical formulas, tables, and lists. The bibliography system is now set up to work seamlessly with Bibby AI's citation management. Feel free to modify it for your needs!

\bibliographystyle{plain}
\bibliography{cite}

\end{document}\documentclass{article}
\usepackage[utf8]{inputenc}
\usepackage{booktabs}
\usepackage{amsmath}
\usepackage{natbib}

\title{My Document}
\author{Your Name}
\date{\today}

\begin{document}

\maketitle

\section{Introduction}
This is a sample document demonstrating various \LaTeX{} features including mathematical equations, tables, and formatting. You can cite papers using the citation search feature, and they will be automatically added to your \texttt{cite.bib} file.

\section{Mathematical Expressions}
Here are some example equations:

\subsection{Quadratic Formula}
The quadratic formula is given by:
\[x = \frac{-b \pm \sqrt{b^2 - 4ac}}{2a}\]

\subsection{Matrix Example}
A sample matrix:
\[
\begin{pmatrix}
a_{11} & a_{12} & a_{13} \\
a_{21} & a_{22} & a_{23} \\
a_{31} & a_{32} & a_{33}
\end{pmatrix}
\]

\section{Tables}
Below is a sample table showing data:

\begin{table}[h]
\centering
\begin{tabular}{@{}lrrr@{}}
\toprule
\textbf{Category} & \textbf{Value 1} & \textbf{Value 2} & \textbf{Total} \\
\midrule
Item A & 10.5 & 20.3 & 30.8 \\
Item B & 15.2 & 25.1 & 40.3 \\
Item C & 22.1 & 12.7 & 34.8 \\
\bottomrule
\end{tabular}
\caption{Sample Data Table}
\label{tab:sample}
\end{table}

\section{Lists}
\subsection{Bullet Points}
\begin{itemize}
  \item First item
  \item Second item
  \item Third item with sub-items:
    \begin{itemize}
      \item Sub-item 1
      \item Sub-item 2
    \end{itemize}
\end{itemize}

\subsection{Numbered List}
\begin{enumerate}
  \item First step
  \item Second step
  \item Third step
\end{enumerate}

\section{Citations and References}
When you use the citation search feature, references will be automatically managed. For example, you might cite a paper like this: \cite{example2024}. All citations will be automatically added to your \texttt{cite.bib} file.

\section{Conclusion}
This template demonstrates basic \LaTeX{} features including mathematical formulas, tables, and lists. The bibliography system is now set up to work seamlessly with Bibby AI's citation management. Feel free to modify it for your needs!

\bibliographystyle{plain}
\bibliography{cite}

\end{document}\documentclass{article}
\usepackage[utf8]{inputenc}
\usepackage{booktabs}
\usepackage{amsmath}
\usepackage{natbib}

\title{My Document}
\author{Your Name}
\date{\today}

\begin{document}

\maketitle

\section{Introduction}
This is a sample document demonstrating various \LaTeX{} features including mathematical equations, tables, and formatting. You can cite papers using the citation search feature, and they will be automatically added to your \texttt{cite.bib} file.

\section{Mathematical Expressions}
Here are some example equations:

\subsection{Quadratic Formula}
The quadratic formula is given by:
\[x = \frac{-b \pm \sqrt{b^2 - 4ac}}{2a}\]

\subsection{Matrix Example}
A sample matrix:
\[
\begin{pmatrix}
a_{11} & a_{12} & a_{13} \\
a_{21} & a_{22} & a_{23} \\
a_{31} & a_{32} & a_{33}
\end{pmatrix}
\]

\section{Tables}
Below is a sample table showing data:

\begin{table}[h]
\centering
\begin{tabular}{@{}lrrr@{}}
\toprule
\textbf{Category} & \textbf{Value 1} & \textbf{Value 2} & \textbf{Total} \\
\midrule
Item A & 10.5 & 20.3 & 30.8 \\
Item B & 15.2 & 25.1 & 40.3 \\
Item C & 22.1 & 12.7 & 34.8 \\
\bottomrule
\end{tabular}
\caption{Sample Data Table}
\label{tab:sample}
\end{table}

\section{Lists}
\subsection{Bullet Points}
\begin{itemize}
  \item First item
  \item Second item
  \item Third item with sub-items:
    \begin{itemize}
      \item Sub-item 1
      \item Sub-item 2
    \end{itemize}
\end{itemize}

\subsection{Numbered List}
\begin{enumerate}
  \item First step
  \item Second step
  \item Third step
\end{enumerate}

\section{Citations and References}
When you use the citation search feature, references will be automatically managed. For example, you might cite a paper like this: \cite{example2024}. All citations will be automatically added to your \texttt{cite.bib} file.

\section{Conclusion}
This template demonstrates basic \LaTeX{} features including mathematical formulas, tables, and lists. The bibliography system is now set up to work seamlessly with Bibby AI's citation management. Feel free to modify it for your needs!

\bibliographystyle{plain}
\bibliography{cite}

\end{document}\documentclass{article}
\usepackage[utf8]{inputenc}
\usepackage{booktabs}
\usepackage{amsmath}
\usepackage{natbib}

\title{My Document}
\author{Your Name}
\date{\today}

\begin{document}

\maketitle

\section{Introduction}
This is a sample document demonstrating various \LaTeX{} features including mathematical equations, tables, and formatting. You can cite papers using the citation search feature, and they will be automatically added to your \texttt{cite.bib} file.

\section{Mathematical Expressions}
Here are some example equations:

\subsection{Quadratic Formula}
The quadratic formula is given by:
\[x = \frac{-b \pm \sqrt{b^2 - 4ac}}{2a}\]

\subsection{Matrix Example}
A sample matrix:
\[
\begin{pmatrix}
a_{11} & a_{12} & a_{13} \\
a_{21} & a_{22} & a_{23} \\
a_{31} & a_{32} & a_{33}
\end{pmatrix}
\]

\section{Tables}
Below is a sample table showing data:

\begin{table}[h]
\centering
\begin{tabular}{@{}lrrr@{}}
\toprule
\textbf{Category} & \textbf{Value 1} & \textbf{Value 2} & \textbf{Total} \\
\midrule
Item A & 10.5 & 20.3 & 30.8 \\
Item B & 15.2 & 25.1 & 40.3 \\
Item C & 22.1 & 12.7 & 34.8 \\
\bottomrule
\end{tabular}
\caption{Sample Data Table}
\label{tab:sample}
\end{table}

\section{Lists}
\subsection{Bullet Points}
\begin{itemize}
  \item First item
  \item Second item
  \item Third item with sub-items:
    \begin{itemize}
      \item Sub-item 1
      \item Sub-item 2
    \end{itemize}
\end{itemize}

\subsection{Numbered List}
\begin{enumerate}
  \item First step
  \item Second step
  \item Third step
\end{enumerate}

\section{Citations and References}
When you use the citation search feature, references will be automatically managed. For example, you might cite a paper like this: \cite{example2024}. All citations will be automatically added to your \texttt{cite.bib} file.

\section{Conclusion}
This template demonstrates basic \LaTeX{} features including mathematical formulas, tables, and lists. The bibliography system is now set up to work seamlessly with Bibby AI's citation management. Feel free to modify it for your needs!

\bibliographystyle{plain}
\bibliography{cite}

\end{document}\documentclass{article}
\usepackage[utf8]{inputenc}
\usepackage{booktabs}
\usepackage{amsmath}
\usepackage{natbib}

\title{My Document}
\author{Your Name}
\date{\today}

\begin{document}

\maketitle

\section{Introduction}
This is a sample document demonstrating various \LaTeX{} features including mathematical equations, tables, and formatting. You can cite papers using the citation search feature, and they will be automatically added to your \texttt{cite.bib} file.

\section{Mathematical Expressions}
Here are some example equations:

\subsection{Quadratic Formula}
The quadratic formula is given by:
\[x = \frac{-b \pm \sqrt{b^2 - 4ac}}{2a}\]

\subsection{Matrix Example}
A sample matrix:
\[
\begin{pmatrix}
a_{11} & a_{12} & a_{13} \\
a_{21} & a_{22} & a_{23} \\
a_{31} & a_{32} & a_{33}
\end{pmatrix}
\]

\section{Tables}
Below is a sample table showing data:

\begin{table}[h]
\centering
\begin{tabular}{@{}lrrr@{}}
\toprule
\textbf{Category} & \textbf{Value 1} & \textbf{Value 2} & \textbf{Total} \\
\midrule
Item A & 10.5 & 20.3 & 30.8 \\
Item B & 15.2 & 25.1 & 40.3 \\
Item C & 22.1 & 12.7 & 34.8 \\
\bottomrule
\end{tabular}
\caption{Sample Data Table}
\label{tab:sample}
\end{table}

\section{Lists}
\subsection{Bullet Points}
\begin{itemize}
  \item First item
  \item Second item
  \item Third item with sub-items:
    \begin{itemize}
      \item Sub-item 1
      \item Sub-item 2
    \end{itemize}
\end{itemize}

\subsection{Numbered List}
\begin{enumerate}
  \item First step
  \item Second step
  \item Third step
\end{enumerate}

\section{Citations and References}
When you use the citation search feature, references will be automatically managed. For example, you might cite a paper like this: \cite{example2024}. All citations will be automatically added to your \texttt{cite.bib} file.

\section{Conclusion}
This template demonstrates basic \LaTeX{} features including mathematical formulas, tables, and lists. The bibliography system is now set up to work seamlessly with Bibby AI's citation management. Feel free to modify it for your needs!

\bibliographystyle{plain}
\bibliography{cite}

\end{document}\documentclass{article}
\usepackage[utf8]{inputenc}
\usepackage{booktabs}
\usepackage{amsmath}
\usepackage{natbib}

\title{My Document}
\author{Your Name}
\date{\today}

\begin{document}

\maketitle

\section{Introduction}
This is a sample document demonstrating various \LaTeX{} features including mathematical equations, tables, and formatting. You can cite papers using the citation search feature, and they will be automatically added to your \texttt{cite.bib} file.

\section{Mathematical Expressions}
Here are some example equations:

\subsection{Quadratic Formula}
The quadratic formula is given by:
\[x = \frac{-b \pm \sqrt{b^2 - 4ac}}{2a}\]

\subsection{Matrix Example}
A sample matrix:
\[
\begin{pmatrix}
a_{11} & a_{12} & a_{13} \\
a_{21} & a_{22} & a_{23} \\
a_{31} & a_{32} & a_{33}
\end{pmatrix}
\]

\section{Tables}
Below is a sample table showing data:

\begin{table}[h]
\centering
\begin{tabular}{@{}lrrr@{}}
\toprule
\textbf{Category} & \textbf{Value 1} & \textbf{Value 2} & \textbf{Total} \\
\midrule
Item A & 10.5 & 20.3 & 30.8 \\
Item B & 15.2 & 25.1 & 40.3 \\
Item C & 22.1 & 12.7 & 34.8 \\
\bottomrule
\end{tabular}
\caption{Sample Data Table}
\label{tab:sample}
\end{table}

\section{Lists}
\subsection{Bullet Points}
\begin{itemize}
  \item First item
  \item Second item
  \item Third item with sub-items:
    \begin{itemize}
      \item Sub-item 1
      \item Sub-item 2
    \end{itemize}
\end{itemize}

\subsection{Numbered List}
\begin{enumerate}
  \item First step
  \item Second step
  \item Third step
\end{enumerate}

\section{Citations and References}
When you use the citation search feature, references will be automatically managed. For example, you might cite a paper like this: \cite{example2024}. All citations will be automatically added to your \texttt{cite.bib} file.

\section{Conclusion}
This template demonstrates basic \LaTeX{} features including mathematical formulas, tables, and lists. The bibliography system is now set up to work seamlessly with Bibby AI's citation management. Feel free to modify it for your needs!

\bibliographystyle{plain}
\bibliography{cite}

\end{document}\documentclass{article}
\usepackage[utf8]{inputenc}
\usepackage{booktabs}
\usepackage{amsmath}
\usepackage{natbib}

\title{My Document}
\author{Your Name}
\date{\today}

\begin{document}

\maketitle

\section{Introduction}
This is a sample document demonstrating various \LaTeX{} features including mathematical equations, tables, and formatting. You can cite papers using the citation search feature, and they will be automatically added to your \texttt{cite.bib} file.

\section{Mathematical Expressions}
Here are some example equations:

\subsection{Quadratic Formula}
The quadratic formula is given by:
\[x = \frac{-b \pm \sqrt{b^2 - 4ac}}{2a}\]

\subsection{Matrix Example}
A sample matrix:
\[
\begin{pmatrix}
a_{11} & a_{12} & a_{13} \\
a_{21} & a_{22} & a_{23} \\
a_{31} & a_{32} & a_{33}
\end{pmatrix}
\]

\section{Tables}
Below is a sample table showing data:

\begin{table}[h]
\centering
\begin{tabular}{@{}lrrr@{}}
\toprule
\textbf{Category} & \textbf{Value 1} & \textbf{Value 2} & \textbf{Total} \\
\midrule
Item A & 10.5 & 20.3 & 30.8 \\
Item B & 15.2 & 25.1 & 40.3 \\
Item C & 22.1 & 12.7 & 34.8 \\
\bottomrule
\end{tabular}
\caption{Sample Data Table}
\label{tab:sample}
\end{table}

\section{Lists}
\subsection{Bullet Points}
\begin{itemize}
  \item First item
  \item Second item
  \item Third item with sub-items:
    \begin{itemize}
      \item Sub-item 1
      \item Sub-item 2
    \end{itemize}
\end{itemize}

\subsection{Numbered List}
\begin{enumerate}
  \item First step
  \item Second step
  \item Third step
\end{enumerate}

\section{Citations and References}
When you use the citation search feature, references will be automatically managed. For example, you might cite a paper like this: \cite{example2024}. All citations will be automatically added to your \texttt{cite.bib} file.

\section{Conclusion}
This template demonstrates basic \LaTeX{} features including mathematical formulas, tables, and lists. The bibliography system is now set up to work seamlessly with Bibby AI's citation management. Feel free to modify it for your needs!

\bibliographystyle{plain}
\bibliography{cite}

\end{document}\documentclass{article}
\usepackage[utf8]{inputenc}
\usepackage{booktabs}
\usepackage{amsmath}
\usepackage{natbib}

\title{My Document}
\author{Your Name}
\date{\today}

\begin{document}

\maketitle

\section{Introduction}
This is a sample document demonstrating various \LaTeX{} features including mathematical equations, tables, and formatting. You can cite papers using the citation search feature, and they will be automatically added to your \texttt{cite.bib} file.

\section{Mathematical Expressions}
Here are some example equations:

\subsection{Quadratic Formula}
The quadratic formula is given by:
\[x = \frac{-b \pm \sqrt{b^2 - 4ac}}{2a}\]

\subsection{Matrix Example}
A sample matrix:
\[
\begin{pmatrix}
a_{11} & a_{12} & a_{13} \\
a_{21} & a_{22} & a_{23} \\
a_{31} & a_{32} & a_{33}
\end{pmatrix}
\]

\section{Tables}
Below is a sample table showing data:

\begin{table}[h]
\centering
\begin{tabular}{@{}lrrr@{}}
\toprule
\textbf{Category} & \textbf{Value 1} & \textbf{Value 2} & \textbf{Total} \\
\midrule
Item A & 10.5 & 20.3 & 30.8 \\
Item B & 15.2 & 25.1 & 40.3 \\
Item C & 22.1 & 12.7 & 34.8 \\
\bottomrule
\end{tabular}
\caption{Sample Data Table}
\label{tab:sample}
\end{table}

\section{Lists}
\subsection{Bullet Points}
\begin{itemize}
  \item First item
  \item Second item
  \item Third item with sub-items:
    \begin{itemize}
      \item Sub-item 1
      \item Sub-item 2
    \end{itemize}
\end{itemize}

\subsection{Numbered List}
\begin{enumerate}
  \item First step
  \item Second step
  \item Third step
\end{enumerate}

\section{Citations and References}
When you use the citation search feature, references will be automatically managed. For example, you might cite a paper like this: \cite{example2024}. All citations will be automatically added to your \texttt{cite.bib} file.

\section{Conclusion}
This template demonstrates basic \LaTeX{} features including mathematical formulas, tables, and lists. The bibliography system is now set up to work seamlessly with Bibby AI's citation management. Feel free to modify it for your needs!

\bibliographystyle{plain}
\bibliography{cite}

\end{document}\documentclass{article}
\usepackage[utf8]{inputenc}
\usepackage{booktabs}
\usepackage{amsmath}
\usepackage{natbib}

\title{My Document}
\author{Your Name}
\date{\today}

\begin{document}

\maketitle

\section{Introduction}
This is a sample document demonstrating various \LaTeX{} features including mathematical equations, tables, and formatting. You can cite papers using the citation search feature, and they will be automatically added to your \texttt{cite.bib} file.

\section{Mathematical Expressions}
Here are some example equations:

\subsection{Quadratic Formula}
The quadratic formula is given by:
\[x = \frac{-b \pm \sqrt{b^2 - 4ac}}{2a}\]

\subsection{Matrix Example}
A sample matrix:
\[
\begin{pmatrix}
a_{11} & a_{12} & a_{13} \\
a_{21} & a_{22} & a_{23} \\
a_{31} & a_{32} & a_{33}
\end{pmatrix}
\]

\section{Tables}
Below is a sample table showing data:

\begin{table}[h]
\centering
\begin{tabular}{@{}lrrr@{}}
\toprule
\textbf{Category} & \textbf{Value 1} & \textbf{Value 2} & \textbf{Total} \\
\midrule
Item A & 10.5 & 20.3 & 30.8 \\
Item B & 15.2 & 25.1 & 40.3 \\
Item C & 22.1 & 12.7 & 34.8 \\
\bottomrule
\end{tabular}
\caption{Sample Data Table}
\label{tab:sample}
\end{table}

\section{Lists}
\subsection{Bullet Points}
\begin{itemize}
  \item First item
  \item Second item
  \item Third item with sub-items:
    \begin{itemize}
      \item Sub-item 1
      \item Sub-item 2
    \end{itemize}
\end{itemize}

\subsection{Numbered List}
\begin{enumerate}
  \item First step
  \item Second step
  \item Third step
\end{enumerate}

\section{Citations and References}
When you use the citation search feature, references will be automatically managed. For example, you might cite a paper like this: \cite{example2024}. All citations will be automatically added to your \texttt{cite.bib} file.

\section{Conclusion}
This template demonstrates basic \LaTeX{} features including mathematical formulas, tables, and lists. The bibliography system is now set up to work seamlessly with Bibby AI's citation management. Feel free to modify it for your needs!

\bibliographystyle{plain}
\bibliography{cite}

\end{document}\documentclass{article}
\usepackage[utf8]{inputenc}
\usepackage{booktabs}
\usepackage{amsmath}
\usepackage{natbib}

\title{My Document}
\author{Your Name}
\date{\today}

\begin{document}

\maketitle

\section{Introduction}
This is a sample document demonstrating various \LaTeX{} features including mathematical equations, tables, and formatting. You can cite papers using the citation search feature, and they will be automatically added to your \texttt{cite.bib} file.

\section{Mathematical Expressions}
Here are some example equations:

\subsection{Quadratic Formula}
The quadratic formula is given by:
\[x = \frac{-b \pm \sqrt{b^2 - 4ac}}{2a}\]

\subsection{Matrix Example}
A sample matrix:
\[
\begin{pmatrix}
a_{11} & a_{12} & a_{13} \\
a_{21} & a_{22} & a_{23} \\
a_{31} & a_{32} & a_{33}
\end{pmatrix}
\]

\section{Tables}
Below is a sample table showing data:

\begin{table}[h]
\centering
\begin{tabular}{@{}lrrr@{}}
\toprule
\textbf{Category} & \textbf{Value 1} & \textbf{Value 2} & \textbf{Total} \\
\midrule
Item A & 10.5 & 20.3 & 30.8 \\
Item B & 15.2 & 25.1 & 40.3 \\
Item C & 22.1 & 12.7 & 34.8 \\
\bottomrule
\end{tabular}
\caption{Sample Data Table}
\label{tab:sample}
\end{table}

\section{Lists}
\subsection{Bullet Points}
\begin{itemize}
  \item First item
  \item Second item
  \item Third item with sub-items:
    \begin{itemize}
      \item Sub-item 1
      \item Sub-item 2
    \end{itemize}
\end{itemize}

\subsection{Numbered List}
\begin{enumerate}
  \item First step
  \item Second step
  \item Third step
\end{enumerate}

\section{Citations and References}
When you use the citation search feature, references will be automatically managed. For example, you might cite a paper like this: \cite{example2024}. All citations will be automatically added to your \texttt{cite.bib} file.

\section{Conclusion}
This template demonstrates basic \LaTeX{} features including mathematical formulas, tables, and lists. The bibliography system is now set up to work seamlessly with Bibby AI's citation management. Feel free to modify it for your needs!

\bibliographystyle{plain}
\bibliography{cite}

\end{document}\documentclass{article}
\usepackage[utf8]{inputenc}
\usepackage{booktabs}
\usepackage{amsmath}
\usepackage{natbib}

\title{My Document}
\author{Your Name}
\date{\today}

\begin{document}

\maketitle

\section{Introduction}
This is a sample document demonstrating various \LaTeX{} features including mathematical equations, tables, and formatting. You can cite papers using the citation search feature, and they will be automatically added to your \texttt{cite.bib} file.

\section{Mathematical Expressions}
Here are some example equations:

\subsection{Quadratic Formula}
The quadratic formula is given by:
\[x = \frac{-b \pm \sqrt{b^2 - 4ac}}{2a}\]

\subsection{Matrix Example}
A sample matrix:
\[
\begin{pmatrix}
a_{11} & a_{12} & a_{13} \\
a_{21} & a_{22} & a_{23} \\
a_{31} & a_{32} & a_{33}
\end{pmatrix}
\]

\section{Tables}
Below is a sample table showing data:

\begin{table}[h]
\centering
\begin{tabular}{@{}lrrr@{}}
\toprule
\textbf{Category} & \textbf{Value 1} & \textbf{Value 2} & \textbf{Total} \\
\midrule
Item A & 10.5 & 20.3 & 30.8 \\
Item B & 15.2 & 25.1 & 40.3 \\
Item C & 22.1 & 12.7 & 34.8 \\
\bottomrule
\end{tabular}
\caption{Sample Data Table}
\label{tab:sample}
\end{table}

\section{Lists}
\subsection{Bullet Points}
\begin{itemize}
  \item First item
  \item Second item
  \item Third item with sub-items:
    \begin{itemize}
      \item Sub-item 1
      \item Sub-item 2
    \end{itemize}
\end{itemize}

\subsection{Numbered List}
\begin{enumerate}
  \item First step
  \item Second step
  \item Third step
\end{enumerate}

\section{Citations and References}
When you use the citation search feature, references will be automatically managed. For example, you might cite a paper like this: \cite{example2024}. All citations will be automatically added to your \texttt{cite.bib} file.

\section{Conclusion}
This template demonstrates basic \LaTeX{} features including mathematical formulas, tables, and lists. The bibliography system is now set up to work seamlessly with Bibby AI's citation management. Feel free to modify it for your needs!

\bibliographystyle{plain}
\bibliography{cite}

\end{document}\documentclass{article}
\usepackage[utf8]{inputenc}
\usepackage{booktabs}
\usepackage{amsmath}
\usepackage{natbib}

\title{My Document}
\author{Your Name}
\date{\today}

\begin{document}

\maketitle

\section{Introduction}
This is a sample document demonstrating various \LaTeX{} features including mathematical equations, tables, and formatting. You can cite papers using the citation search feature, and they will be automatically added to your \texttt{cite.bib} file.

\section{Mathematical Expressions}
Here are some example equations:

\subsection{Quadratic Formula}
The quadratic formula is given by:
\[x = \frac{-b \pm \sqrt{b^2 - 4ac}}{2a}\]

\subsection{Matrix Example}
A sample matrix:
\[
\begin{pmatrix}
a_{11} & a_{12} & a_{13} \\
a_{21} & a_{22} & a_{23} \\
a_{31} & a_{32} & a_{33}
\end{pmatrix}
\]

\section{Tables}
Below is a sample table showing data:

\begin{table}[h]
\centering
\begin{tabular}{@{}lrrr@{}}
\toprule
\textbf{Category} & \textbf{Value 1} & \textbf{Value 2} & \textbf{Total} \\
\midrule
Item A & 10.5 & 20.3 & 30.8 \\
Item B & 15.2 & 25.1 & 40.3 \\
Item C & 22.1 & 12.7 & 34.8 \\
\bottomrule
\end{tabular}
\caption{Sample Data Table}
\label{tab:sample}
\end{table}

\section{Lists}
\subsection{Bullet Points}
\begin{itemize}
  \item First item
  \item Second item
  \item Third item with sub-items:
    \begin{itemize}
      \item Sub-item 1
      \item Sub-item 2
    \end{itemize}
\end{itemize}

\subsection{Numbered List}
\begin{enumerate}
  \item First step
  \item Second step
  \item Third step
\end{enumerate}

\section{Citations and References}
When you use the citation search feature, references will be automatically managed. For example, you might cite a paper like this: \cite{example2024}. All citations will be automatically added to your \texttt{cite.bib} file.

\section{Conclusion}
This template demonstrates basic \LaTeX{} features including mathematical formulas, tables, and lists. The bibliography system is now set up to work seamlessly with Bibby AI's citation management. Feel free to modify it for your needs!

\bibliographystyle{plain}
\bibliography{cite}

\end{document}\documentclass{article}
\usepackage[utf8]{inputenc}
\usepackage{booktabs}
\usepackage{amsmath}
\usepackage{natbib}

\title{My Document}
\author{Your Name}
\date{\today}

\begin{document}

\maketitle

\section{Introduction}
This is a sample document demonstrating various \LaTeX{} features including mathematical equations, tables, and formatting. You can cite papers using the citation search feature, and they will be automatically added to your \texttt{cite.bib} file.

\section{Mathematical Expressions}
Here are some example equations:

\subsection{Quadratic Formula}
The quadratic formula is given by:
\[x = \frac{-b \pm \sqrt{b^2 - 4ac}}{2a}\]

\subsection{Matrix Example}
A sample matrix:
\[
\begin{pmatrix}
a_{11} & a_{12} & a_{13} \\
a_{21} & a_{22} & a_{23} \\
a_{31} & a_{32} & a_{33}
\end{pmatrix}
\]

\section{Tables}
Below is a sample table showing data:

\begin{table}[h]
\centering
\begin{tabular}{@{}lrrr@{}}
\toprule
\textbf{Category} & \textbf{Value 1} & \textbf{Value 2} & \textbf{Total} \\
\midrule
Item A & 10.5 & 20.3 & 30.8 \\
Item B & 15.2 & 25.1 & 40.3 \\
Item C & 22.1 & 12.7 & 34.8 \\
\bottomrule
\end{tabular}
\caption{Sample Data Table}
\label{tab:sample}
\end{table}

\section{Lists}
\subsection{Bullet Points}
\begin{itemize}
  \item First item
  \item Second item
  \item Third item with sub-items:
    \begin{itemize}
      \item Sub-item 1
      \item Sub-item 2
    \end{itemize}
\end{itemize}

\subsection{Numbered List}
\begin{enumerate}
  \item First step
  \item Second step
  \item Third step
\end{enumerate}

\section{Citations and References}
When you use the citation search feature, references will be automatically managed. For example, you might cite a paper like this: \cite{example2024}. All citations will be automatically added to your \texttt{cite.bib} file.

\section{Conclusion}
This template demonstrates basic \LaTeX{} features including mathematical formulas, tables, and lists. The bibliography system is now set up to work seamlessly with Bibby AI's citation management. Feel free to modify it for your needs!

\bibliographystyle{plain}
\bibliography{cite}

\end{document}\documentclass{article}
\usepackage[utf8]{inputenc}
\usepackage{booktabs}
\usepackage{amsmath}
\usepackage{natbib}

\title{My Document}
\author{Your Name}
\date{\today}

\begin{document}

\maketitle

\section{Introduction}
This is a sample document demonstrating various \LaTeX{} features including mathematical equations, tables, and formatting. You can cite papers using the citation search feature, and they will be automatically added to your \texttt{cite.bib} file.

\section{Mathematical Expressions}
Here are some example equations:

\subsection{Quadratic Formula}
The quadratic formula is given by:
\[x = \frac{-b \pm \sqrt{b^2 - 4ac}}{2a}\]

\subsection{Matrix Example}
A sample matrix:
\[
\begin{pmatrix}
a_{11} & a_{12} & a_{13} \\
a_{21} & a_{22} & a_{23} \\
a_{31} & a_{32} & a_{33}
\end{pmatrix}
\]

\section{Tables}
Below is a sample table showing data:

\begin{table}[h]
\centering
\begin{tabular}{@{}lrrr@{}}
\toprule
\textbf{Category} & \textbf{Value 1} & \textbf{Value 2} & \textbf{Total} \\
\midrule
Item A & 10.5 & 20.3 & 30.8 \\
Item B & 15.2 & 25.1 & 40.3 \\
Item C & 22.1 & 12.7 & 34.8 \\
\bottomrule
\end{tabular}
\caption{Sample Data Table}
\label{tab:sample}
\end{table}

\section{Lists}
\subsection{Bullet Points}
\begin{itemize}
  \item First item
  \item Second item
  \item Third item with sub-items:
    \begin{itemize}
      \item Sub-item 1
      \item Sub-item 2
    \end{itemize}
\end{itemize}

\subsection{Numbered List}
\begin{enumerate}
  \item First step
  \item Second step
  \item Third step
\end{enumerate}

\section{Citations and References}
When you use the citation search feature, references will be automatically managed. For example, you might cite a paper like this: \cite{example2024}. All citations will be automatically added to your \texttt{cite.bib} file.

\section{Conclusion}
This template demonstrates basic \LaTeX{} features including mathematical formulas, tables, and lists. The bibliography system is now set up to work seamlessly with Bibby AI's citation management. Feel free to modify it for your needs!

\bibliographystyle{plain}
\bibliography{cite}

\end{document}\documentclass{article}
\usepackage[utf8]{inputenc}
\usepackage{booktabs}
\usepackage{amsmath}
\usepackage{natbib}

\title{My Document}
\author{Your Name}
\date{\today}

\begin{document}

\maketitle

\section{Introduction}
This is a sample document demonstrating various \LaTeX{} features including mathematical equations, tables, and formatting. You can cite papers using the citation search feature, and they will be automatically added to your \texttt{cite.bib} file.

\section{Mathematical Expressions}
Here are some example equations:

\subsection{Quadratic Formula}
The quadratic formula is given by:
\[x = \frac{-b \pm \sqrt{b^2 - 4ac}}{2a}\]

\subsection{Matrix Example}
A sample matrix:
\[
\begin{pmatrix}
a_{11} & a_{12} & a_{13} \\
a_{21} & a_{22} & a_{23} \\
a_{31} & a_{32} & a_{33}
\end{pmatrix}
\]

\section{Tables}
Below is a sample table showing data:

\begin{table}[h]
\centering
\begin{tabular}{@{}lrrr@{}}
\toprule
\textbf{Category} & \textbf{Value 1} & \textbf{Value 2} & \textbf{Total} \\
\midrule
Item A & 10.5 & 20.3 & 30.8 \\
Item B & 15.2 & 25.1 & 40.3 \\
Item C & 22.1 & 12.7 & 34.8 \\
\bottomrule
\end{tabular}
\caption{Sample Data Table}
\label{tab:sample}
\end{table}

\section{Lists}
\subsection{Bullet Points}
\begin{itemize}
  \item First item
  \item Second item
  \item Third item with sub-items:
    \begin{itemize}
      \item Sub-item 1
      \item Sub-item 2
    \end{itemize}
\end{itemize}

\subsection{Numbered List}
\begin{enumerate}
  \item First step
  \item Second step
  \item Third step
\end{enumerate}

\section{Citations and References}
When you use the citation search feature, references will be automatically managed. For example, you might cite a paper like this: \cite{example2024}. All citations will be automatically added to your \texttt{cite.bib} file.

\section{Conclusion}
This template demonstrates basic \LaTeX{} features including mathematical formulas, tables, and lists. The bibliography system is now set up to work seamlessly with Bibby AI's citation management. Feel free to modify it for your needs!

\bibliographystyle{plain}
\bibliography{cite}

\end{document}\documentclass{article}
\usepackage[utf8]{inputenc}
\usepackage{booktabs}
\usepackage{amsmath}
\usepackage{natbib}

\title{My Document}
\author{Your Name}
\date{\today}

\begin{document}

\maketitle

\section{Introduction}
This is a sample document demonstrating various \LaTeX{} features including mathematical equations, tables, and formatting. You can cite papers using the citation search feature, and they will be automatically added to your \texttt{cite.bib} file.

\section{Mathematical Expressions}
Here are some example equations:

\subsection{Quadratic Formula}
The quadratic formula is given by:
\[x = \frac{-b \pm \sqrt{b^2 - 4ac}}{2a}\]

\subsection{Matrix Example}
A sample matrix:
\[
\begin{pmatrix}
a_{11} & a_{12} & a_{13} \\
a_{21} & a_{22} & a_{23} \\
a_{31} & a_{32} & a_{33}
\end{pmatrix}
\]

\section{Tables}
Below is a sample table showing data:

\begin{table}[h]
\centering
\begin{tabular}{@{}lrrr@{}}
\toprule
\textbf{Category} & \textbf{Value 1} & \textbf{Value 2} & \textbf{Total} \\
\midrule
Item A & 10.5 & 20.3 & 30.8 \\
Item B & 15.2 & 25.1 & 40.3 \\
Item C & 22.1 & 12.7 & 34.8 \\
\bottomrule
\end{tabular}
\caption{Sample Data Table}
\label{tab:sample}
\end{table}

\section{Lists}
\subsection{Bullet Points}
\begin{itemize}
  \item First item
  \item Second item
  \item Third item with sub-items:
    \begin{itemize}
      \item Sub-item 1
      \item Sub-item 2
    \end{itemize}
\end{itemize}

\subsection{Numbered List}
\begin{enumerate}
  \item First step
  \item Second step
  \item Third step
\end{enumerate}

\section{Citations and References}
When you use the citation search feature, references will be automatically managed. For example, you might cite a paper like this: \cite{example2024}. All citations will be automatically added to your \texttt{cite.bib} file.

\section{Conclusion}
This template demonstrates basic \LaTeX{} features including mathematical formulas, tables, and lists. The bibliography system is now set up to work seamlessly with Bibby AI's citation management. Feel free to modify it for your needs!

\bibliographystyle{plain}
\bibliography{cite}

\end{document}\documentclass{article}
\usepackage[utf8]{inputenc}
\usepackage{booktabs}
\usepackage{amsmath}
\usepackage{natbib}

\title{My Document}
\author{Your Name}
\date{\today}

\begin{document}

\maketitle

\section{Introduction}
This is a sample document demonstrating various \LaTeX{} features including mathematical equations, tables, and formatting. You can cite papers using the citation search feature, and they will be automatically added to your \texttt{cite.bib} file.

\section{Mathematical Expressions}
Here are some example equations:

\subsection{Quadratic Formula}
The quadratic formula is given by:
\[x = \frac{-b \pm \sqrt{b^2 - 4ac}}{2a}\]

\subsection{Matrix Example}
A sample matrix:
\[
\begin{pmatrix}
a_{11} & a_{12} & a_{13} \\
a_{21} & a_{22} & a_{23} \\
a_{31} & a_{32} & a_{33}
\end{pmatrix}
\]

\section{Tables}
Below is a sample table showing data:

\begin{table}[h]
\centering
\begin{tabular}{@{}lrrr@{}}
\toprule
\textbf{Category} & \textbf{Value 1} & \textbf{Value 2} & \textbf{Total} \\
\midrule
Item A & 10.5 & 20.3 & 30.8 \\
Item B & 15.2 & 25.1 & 40.3 \\
Item C & 22.1 & 12.7 & 34.8 \\
\bottomrule
\end{tabular}
\caption{Sample Data Table}
\label{tab:sample}
\end{table}

\section{Lists}
\subsection{Bullet Points}
\begin{itemize}
  \item First item
  \item Second item
  \item Third item with sub-items:
    \begin{itemize}
      \item Sub-item 1
      \item Sub-item 2
    \end{itemize}
\end{itemize}

\subsection{Numbered List}
\begin{enumerate}
  \item First step
  \item Second step
  \item Third step
\end{enumerate}

\section{Citations and References}
When you use the citation search feature, references will be automatically managed. For example, you might cite a paper like this: \cite{example2024}. All citations will be automatically added to your \texttt{cite.bib} file.

\section{Conclusion}
This template demonstrates basic \LaTeX{} features including mathematical formulas, tables, and lists. The bibliography system is now set up to work seamlessly with Bibby AI's citation management. Feel free to modify it for your needs!

\bibliographystyle{plain}
\bibliography{cite}

\end{document}\documentclass{article}
\usepackage[utf8]{inputenc}
\usepackage{booktabs}
\usepackage{amsmath}
\usepackage{natbib}

\title{My Document}
\author{Your Name}
\date{\today}

\begin{document}

\maketitle

\section{Introduction}
This is a sample document demonstrating various \LaTeX{} features including mathematical equations, tables, and formatting. You can cite papers using the citation search feature, and they will be automatically added to your \texttt{cite.bib} file.

\section{Mathematical Expressions}
Here are some example equations:

\subsection{Quadratic Formula}
The quadratic formula is given by:
\[x = \frac{-b \pm \sqrt{b^2 - 4ac}}{2a}\]

\subsection{Matrix Example}
A sample matrix:
\[
\begin{pmatrix}
a_{11} & a_{12} & a_{13} \\
a_{21} & a_{22} & a_{23} \\
a_{31} & a_{32} & a_{33}
\end{pmatrix}
\]

\section{Tables}
Below is a sample table showing data:

\begin{table}[h]
\centering
\begin{tabular}{@{}lrrr@{}}
\toprule
\textbf{Category} & \textbf{Value 1} & \textbf{Value 2} & \textbf{Total} \\
\midrule
Item A & 10.5 & 20.3 & 30.8 \\
Item B & 15.2 & 25.1 & 40.3 \\
Item C & 22.1 & 12.7 & 34.8 \\
\bottomrule
\end{tabular}
\caption{Sample Data Table}
\label{tab:sample}
\end{table}

\section{Lists}
\subsection{Bullet Points}
\begin{itemize}
  \item First item
  \item Second item
  \item Third item with sub-items:
    \begin{itemize}
      \item Sub-item 1
      \item Sub-item 2
    \end{itemize}
\end{itemize}

\subsection{Numbered List}
\begin{enumerate}
  \item First step
  \item Second step
  \item Third step
\end{enumerate}

\section{Citations and References}
When you use the citation search feature, references will be automatically managed. For example, you might cite a paper like this: \cite{example2024}. All citations will be automatically added to your \texttt{cite.bib} file.

\section{Conclusion}
This template demonstrates basic \LaTeX{} features including mathematical formulas, tables, and lists. The bibliography system is now set up to work seamlessly with Bibby AI's citation management. Feel free to modify it for your needs!

\bibliographystyle{plain}
\bibliography{cite}

\end{document}\documentclass{article}
\usepackage[utf8]{inputenc}
\usepackage{booktabs}
\usepackage{amsmath}
\usepackage{natbib}

\title{My Document}
\author{Your Name}
\date{\today}

\begin{document}

\maketitle

\section{Introduction}
This is a sample document demonstrating various \LaTeX{} features including mathematical equations, tables, and formatting. You can cite papers using the citation search feature, and they will be automatically added to your \texttt{cite.bib} file.

\section{Mathematical Expressions}
Here are some example equations:

\subsection{Quadratic Formula}
The quadratic formula is given by:
\[x = \frac{-b \pm \sqrt{b^2 - 4ac}}{2a}\]

\subsection{Matrix Example}
A sample matrix:
\[
\begin{pmatrix}
a_{11} & a_{12} & a_{13} \\
a_{21} & a_{22} & a_{23} \\
a_{31} & a_{32} & a_{33}
\end{pmatrix}
\]

\section{Tables}
Below is a sample table showing data:

\begin{table}[h]
\centering
\begin{tabular}{@{}lrrr@{}}
\toprule
\textbf{Category} & \textbf{Value 1} & \textbf{Value 2} & \textbf{Total} \\
\midrule
Item A & 10.5 & 20.3 & 30.8 \\
Item B & 15.2 & 25.1 & 40.3 \\
Item C & 22.1 & 12.7 & 34.8 \\
\bottomrule
\end{tabular}
\caption{Sample Data Table}
\label{tab:sample}
\end{table}

\section{Lists}
\subsection{Bullet Points}
\begin{itemize}
  \item First item
  \item Second item
  \item Third item with sub-items:
    \begin{itemize}
      \item Sub-item 1
      \item Sub-item 2
    \end{itemize}
\end{itemize}

\subsection{Numbered List}
\begin{enumerate}
  \item First step
  \item Second step
  \item Third step
\end{enumerate}

\section{Citations and References}
When you use the citation search feature, references will be automatically managed. For example, you might cite a paper like this: \cite{example2024}. All citations will be automatically added to your \texttt{cite.bib} file.

\section{Conclusion}
This template demonstrates basic \LaTeX{} features including mathematical formulas, tables, and lists. The bibliography system is now set up to work seamlessly with Bibby AI's citation management. Feel free to modify it for your needs!

\bibliographystyle{plain}
\bibliography{cite}

\end{document}\documentclass{article}
\usepackage[utf8]{inputenc}
\usepackage{booktabs}
\usepackage{amsmath}
\usepackage{natbib}

\title{My Document}
\author{Your Name}
\date{\today}

\begin{document}

\maketitle

\section{Introduction}
This is a sample document demonstrating various \LaTeX{} features including mathematical equations, tables, and formatting. You can cite papers using the citation search feature, and they will be automatically added to your \texttt{cite.bib} file.

\section{Mathematical Expressions}
Here are some example equations:

\subsection{Quadratic Formula}
The quadratic formula is given by:
\[x = \frac{-b \pm \sqrt{b^2 - 4ac}}{2a}\]

\subsection{Matrix Example}
A sample matrix:
\[
\begin{pmatrix}
a_{11} & a_{12} & a_{13} \\
a_{21} & a_{22} & a_{23} \\
a_{31} & a_{32} & a_{33}
\end{pmatrix}
\]

\section{Tables}
Below is a sample table showing data:

\begin{table}[h]
\centering
\begin{tabular}{@{}lrrr@{}}
\toprule
\textbf{Category} & \textbf{Value 1} & \textbf{Value 2} & \textbf{Total} \\
\midrule
Item A & 10.5 & 20.3 & 30.8 \\
Item B & 15.2 & 25.1 & 40.3 \\
Item C & 22.1 & 12.7 & 34.8 \\
\bottomrule
\end{tabular}
\caption{Sample Data Table}
\label{tab:sample}
\end{table}

\section{Lists}
\subsection{Bullet Points}
\begin{itemize}
  \item First item
  \item Second item
  \item Third item with sub-items:
    \begin{itemize}
      \item Sub-item 1
      \item Sub-item 2
    \end{itemize}
\end{itemize}

\subsection{Numbered List}
\begin{enumerate}
  \item First step
  \item Second step
  \item Third step
\end{enumerate}

\section{Citations and References}
When you use the citation search feature, references will be automatically managed. For example, you might cite a paper like this: \cite{example2024}. All citations will be automatically added to your \texttt{cite.bib} file.

\section{Conclusion}
This template demonstrates basic \LaTeX{} features including mathematical formulas, tables, and lists. The bibliography system is now set up to work seamlessly with Bibby AI's citation management. Feel free to modify it for your needs!

\bibliographystyle{plain}
\bibliography{cite}

\end{document}\documentclass{article}
\usepackage[utf8]{inputenc}
\usepackage{booktabs}
\usepackage{amsmath}
\usepackage{natbib}

\title{My Document}
\author{Your Name}
\date{\today}

\begin{document}

\maketitle

\section{Introduction}
This is a sample document demonstrating various \LaTeX{} features including mathematical equations, tables, and formatting. You can cite papers using the citation search feature, and they will be automatically added to your \texttt{cite.bib} file.

\section{Mathematical Expressions}
Here are some example equations:

\subsection{Quadratic Formula}
The quadratic formula is given by:
\[x = \frac{-b \pm \sqrt{b^2 - 4ac}}{2a}\]

\subsection{Matrix Example}
A sample matrix:
\[
\begin{pmatrix}
a_{11} & a_{12} & a_{13} \\
a_{21} & a_{22} & a_{23} \\
a_{31} & a_{32} & a_{33}
\end{pmatrix}
\]

\section{Tables}
Below is a sample table showing data:

\begin{table}[h]
\centering
\begin{tabular}{@{}lrrr@{}}
\toprule
\textbf{Category} & \textbf{Value 1} & \textbf{Value 2} & \textbf{Total} \\
\midrule
Item A & 10.5 & 20.3 & 30.8 \\
Item B & 15.2 & 25.1 & 40.3 \\
Item C & 22.1 & 12.7 & 34.8 \\
\bottomrule
\end{tabular}
\caption{Sample Data Table}
\label{tab:sample}
\end{table}

\section{Lists}
\subsection{Bullet Points}
\begin{itemize}
  \item First item
  \item Second item
  \item Third item with sub-items:
    \begin{itemize}
      \item Sub-item 1
      \item Sub-item 2
    \end{itemize}
\end{itemize}

\subsection{Numbered List}
\begin{enumerate}
  \item First step
  \item Second step
  \item Third step
\end{enumerate}

\section{Citations and References}
When you use the citation search feature, references will be automatically managed. For example, you might cite a paper like this: \cite{example2024}. All citations will be automatically added to your \texttt{cite.bib} file.

\section{Conclusion}
This template demonstrates basic \LaTeX{} features including mathematical formulas, tables, and lists. The bibliography system is now set up to work seamlessly with Bibby AI's citation management. Feel free to modify it for your needs!

\bibliographystyle{plain}
\bibliography{cite}

\end{document}\documentclass{article}
\usepackage[utf8]{inputenc}
\usepackage{booktabs}
\usepackage{amsmath}
\usepackage{natbib}

\title{My Document}
\author{Your Name}
\date{\today}

\begin{document}

\maketitle

\section{Introduction}
This is a sample document demonstrating various \LaTeX{} features including mathematical equations, tables, and formatting. You can cite papers using the citation search feature, and they will be automatically added to your \texttt{cite.bib} file.

\section{Mathematical Expressions}
Here are some example equations:

\subsection{Quadratic Formula}
The quadratic formula is given by:
\[x = \frac{-b \pm \sqrt{b^2 - 4ac}}{2a}\]

\subsection{Matrix Example}
A sample matrix:
\[
\begin{pmatrix}
a_{11} & a_{12} & a_{13} \\
a_{21} & a_{22} & a_{23} \\
a_{31} & a_{32} & a_{33}
\end{pmatrix}
\]

\section{Tables}
Below is a sample table showing data:

\begin{table}[h]
\centering
\begin{tabular}{@{}lrrr@{}}
\toprule
\textbf{Category} & \textbf{Value 1} & \textbf{Value 2} & \textbf{Total} \\
\midrule
Item A & 10.5 & 20.3 & 30.8 \\
Item B & 15.2 & 25.1 & 40.3 \\
Item C & 22.1 & 12.7 & 34.8 \\
\bottomrule
\end{tabular}
\caption{Sample Data Table}
\label{tab:sample}
\end{table}

\section{Lists}
\subsection{Bullet Points}
\begin{itemize}
  \item First item
  \item Second item
  \item Third item with sub-items:
    \begin{itemize}
      \item Sub-item 1
      \item Sub-item 2
    \end{itemize}
\end{itemize}

\subsection{Numbered List}
\begin{enumerate}
  \item First step
  \item Second step
  \item Third step
\end{enumerate}

\section{Citations and References}
When you use the citation search feature, references will be automatically managed. For example, you might cite a paper like this: \cite{example2024}. All citations will be automatically added to your \texttt{cite.bib} file.

\section{Conclusion}
This template demonstrates basic \LaTeX{} features including mathematical formulas, tables, and lists. The bibliography system is now set up to work seamlessly with Bibby AI's citation management. Feel free to modify it for your needs!

\bibliographystyle{plain}
\bibliography{cite}

\end{document}\documentclass{article}
\usepackage[utf8]{inputenc}
\usepackage{booktabs}
\usepackage{amsmath}
\usepackage{natbib}

\title{My Document}
\author{Your Name}
\date{\today}

\begin{document}

\maketitle

\section{Introduction}
This is a sample document demonstrating various \LaTeX{} features including mathematical equations, tables, and formatting. You can cite papers using the citation search feature, and they will be automatically added to your \texttt{cite.bib} file.

\section{Mathematical Expressions}
Here are some example equations:

\subsection{Quadratic Formula}
The quadratic formula is given by:
\[x = \frac{-b \pm \sqrt{b^2 - 4ac}}{2a}\]

\subsection{Matrix Example}
A sample matrix:
\[
\begin{pmatrix}
a_{11} & a_{12} & a_{13} \\
a_{21} & a_{22} & a_{23} \\
a_{31} & a_{32} & a_{33}
\end{pmatrix}
\]

\section{Tables}
Below is a sample table showing data:

\begin{table}[h]
\centering
\begin{tabular}{@{}lrrr@{}}
\toprule
\textbf{Category} & \textbf{Value 1} & \textbf{Value 2} & \textbf{Total} \\
\midrule
Item A & 10.5 & 20.3 & 30.8 \\
Item B & 15.2 & 25.1 & 40.3 \\
Item C & 22.1 & 12.7 & 34.8 \\
\bottomrule
\end{tabular}
\caption{Sample Data Table}
\label{tab:sample}
\end{table}

\section{Lists}
\subsection{Bullet Points}
\begin{itemize}
  \item First item
  \item Second item
  \item Third item with sub-items:
    \begin{itemize}
      \item Sub-item 1
      \item Sub-item 2
    \end{itemize}
\end{itemize}

\subsection{Numbered List}
\begin{enumerate}
  \item First step
  \item Second step
  \item Third step
\end{enumerate}

\section{Citations and References}
When you use the citation search feature, references will be automatically managed. For example, you might cite a paper like this: \cite{example2024}. All citations will be automatically added to your \texttt{cite.bib} file.

\section{Conclusion}
This template demonstrates basic \LaTeX{} features including mathematical formulas, tables, and lists. The bibliography system is now set up to work seamlessly with Bibby AI's citation management. Feel free to modify it for your needs!

\bibliographystyle{plain}
\bibliography{cite}

\end{document}\documentclass{article}
\usepackage[utf8]{inputenc}
\usepackage{booktabs}
\usepackage{amsmath}
\usepackage{natbib}

\title{My Document}
\author{Your Name}
\date{\today}

\begin{document}

\maketitle

\section{Introduction}
This is a sample document demonstrating various \LaTeX{} features including mathematical equations, tables, and formatting. You can cite papers using the citation search feature, and they will be automatically added to your \texttt{cite.bib} file.

\section{Mathematical Expressions}
Here are some example equations:

\subsection{Quadratic Formula}
The quadratic formula is given by:
\[x = \frac{-b \pm \sqrt{b^2 - 4ac}}{2a}\]

\subsection{Matrix Example}
A sample matrix:
\[
\begin{pmatrix}
a_{11} & a_{12} & a_{13} \\
a_{21} & a_{22} & a_{23} \\
a_{31} & a_{32} & a_{33}
\end{pmatrix}
\]

\section{Tables}
Below is a sample table showing data:

\begin{table}[h]
\centering
\begin{tabular}{@{}lrrr@{}}
\toprule
\textbf{Category} & \textbf{Value 1} & \textbf{Value 2} & \textbf{Total} \\
\midrule
Item A & 10.5 & 20.3 & 30.8 \\
Item B & 15.2 & 25.1 & 40.3 \\
Item C & 22.1 & 12.7 & 34.8 \\
\bottomrule
\end{tabular}
\caption{Sample Data Table}
\label{tab:sample}
\end{table}

\section{Lists}
\subsection{Bullet Points}
\begin{itemize}
  \item First item
  \item Second item
  \item Third item with sub-items:
    \begin{itemize}
      \item Sub-item 1
      \item Sub-item 2
    \end{itemize}
\end{itemize}

\subsection{Numbered List}
\begin{enumerate}
  \item First step
  \item Second step
  \item Third step
\end{enumerate}

\section{Citations and References}
When you use the citation search feature, references will be automatically managed. For example, you might cite a paper like this: \cite{example2024}. All citations will be automatically added to your \texttt{cite.bib} file.

\section{Conclusion}
This template demonstrates basic \LaTeX{} features including mathematical formulas, tables, and lists. The bibliography system is now set up to work seamlessly with Bibby AI's citation management. Feel free to modify it for your needs!

\bibliographystyle{plain}
\bibliography{cite}

\end{document}\documentclass{article}
\usepackage[utf8]{inputenc}
\usepackage{booktabs}
\usepackage{amsmath}
\usepackage{natbib}

\title{My Document}
\author{Your Name}
\date{\today}

\begin{document}

\maketitle

\section{Introduction}
This is a sample document demonstrating various \LaTeX{} features including mathematical equations, tables, and formatting. You can cite papers using the citation search feature, and they will be automatically added to your \texttt{cite.bib} file.

\section{Mathematical Expressions}
Here are some example equations:

\subsection{Quadratic Formula}
The quadratic formula is given by:
\[x = \frac{-b \pm \sqrt{b^2 - 4ac}}{2a}\]

\subsection{Matrix Example}
A sample matrix:
\[
\begin{pmatrix}
a_{11} & a_{12} & a_{13} \\
a_{21} & a_{22} & a_{23} \\
a_{31} & a_{32} & a_{33}
\end{pmatrix}
\]

\section{Tables}
Below is a sample table showing data:

\begin{table}[h]
\centering
\begin{tabular}{@{}lrrr@{}}
\toprule
\textbf{Category} & \textbf{Value 1} & \textbf{Value 2} & \textbf{Total} \\
\midrule
Item A & 10.5 & 20.3 & 30.8 \\
Item B & 15.2 & 25.1 & 40.3 \\
Item C & 22.1 & 12.7 & 34.8 \\
\bottomrule
\end{tabular}
\caption{Sample Data Table}
\label{tab:sample}
\end{table}

\section{Lists}
\subsection{Bullet Points}
\begin{itemize}
  \item First item
  \item Second item
  \item Third item with sub-items:
    \begin{itemize}
      \item Sub-item 1
      \item Sub-item 2
    \end{itemize}
\end{itemize}

\subsection{Numbered List}
\begin{enumerate}
  \item First step
  \item Second step
  \item Third step
\end{enumerate}

\section{Citations and References}
When you use the citation search feature, references will be automatically managed. For example, you might cite a paper like this: \cite{example2024}. All citations will be automatically added to your \texttt{cite.bib} file.

\section{Conclusion}
This template demonstrates basic \LaTeX{} features including mathematical formulas, tables, and lists. The bibliography system is now set up to work seamlessly with Bibby AI's citation management. Feel free to modify it for your needs!

\bibliographystyle{plain}
\bibliography{cite}

\end{document}\documentclass{article}
\usepackage[utf8]{inputenc}
\usepackage{booktabs}
\usepackage{amsmath}
\usepackage{natbib}

\title{My Document}
\author{Your Name}
\date{\today}

\begin{document}

\maketitle

\section{Introduction}
This is a sample document demonstrating various \LaTeX{} features including mathematical equations, tables, and formatting. You can cite papers using the citation search feature, and they will be automatically added to your \texttt{cite.bib} file.

\section{Mathematical Expressions}
Here are some example equations:

\subsection{Quadratic Formula}
The quadratic formula is given by:
\[x = \frac{-b \pm \sqrt{b^2 - 4ac}}{2a}\]

\subsection{Matrix Example}
A sample matrix:
\[
\begin{pmatrix}
a_{11} & a_{12} & a_{13} \\
a_{21} & a_{22} & a_{23} \\
a_{31} & a_{32} & a_{33}
\end{pmatrix}
\]

\section{Tables}
Below is a sample table showing data:

\begin{table}[h]
\centering
\begin{tabular}{@{}lrrr@{}}
\toprule
\textbf{Category} & \textbf{Value 1} & \textbf{Value 2} & \textbf{Total} \\
\midrule
Item A & 10.5 & 20.3 & 30.8 \\
Item B & 15.2 & 25.1 & 40.3 \\
Item C & 22.1 & 12.7 & 34.8 \\
\bottomrule
\end{tabular}
\caption{Sample Data Table}
\label{tab:sample}
\end{table}

\section{Lists}
\subsection{Bullet Points}
\begin{itemize}
  \item First item
  \item Second item
  \item Third item with sub-items:
    \begin{itemize}
      \item Sub-item 1
      \item Sub-item 2
    \end{itemize}
\end{itemize}

\subsection{Numbered List}
\begin{enumerate}
  \item First step
  \item Second step
  \item Third step
\end{enumerate}

\section{Citations and References}
When you use the citation search feature, references will be automatically managed. For example, you might cite a paper like this: \cite{example2024}. All citations will be automatically added to your \texttt{cite.bib} file.

\section{Conclusion}
This template demonstrates basic \LaTeX{} features including mathematical formulas, tables, and lists. The bibliography system is now set up to work seamlessly with Bibby AI's citation management. Feel free to modify it for your needs!

\bibliographystyle{plain}
\bibliography{cite}

\end{document}\documentclass{article}
\usepackage[utf8]{inputenc}
\usepackage{booktabs}
\usepackage{amsmath}
\usepackage{natbib}

\title{My Document}
\author{Your Name}
\date{\today}

\begin{document}

\maketitle

\section{Introduction}
This is a sample document demonstrating various \LaTeX{} features including mathematical equations, tables, and formatting. You can cite papers using the citation search feature, and they will be automatically added to your \texttt{cite.bib} file.

\section{Mathematical Expressions}
Here are some example equations:

\subsection{Quadratic Formula}
The quadratic formula is given by:
\[x = \frac{-b \pm \sqrt{b^2 - 4ac}}{2a}\]

\subsection{Matrix Example}
A sample matrix:
\[
\begin{pmatrix}
a_{11} & a_{12} & a_{13} \\
a_{21} & a_{22} & a_{23} \\
a_{31} & a_{32} & a_{33}
\end{pmatrix}
\]

\section{Tables}
Below is a sample table showing data:

\begin{table}[h]
\centering
\begin{tabular}{@{}lrrr@{}}
\toprule
\textbf{Category} & \textbf{Value 1} & \textbf{Value 2} & \textbf{Total} \\
\midrule
Item A & 10.5 & 20.3 & 30.8 \\
Item B & 15.2 & 25.1 & 40.3 \\
Item C & 22.1 & 12.7 & 34.8 \\
\bottomrule
\end{tabular}
\caption{Sample Data Table}
\label{tab:sample}
\end{table}

\section{Lists}
\subsection{Bullet Points}
\begin{itemize}
  \item First item
  \item Second item
  \item Third item with sub-items:
    \begin{itemize}
      \item Sub-item 1
      \item Sub-item 2
    \end{itemize}
\end{itemize}

\subsection{Numbered List}
\begin{enumerate}
  \item First step
  \item Second step
  \item Third step
\end{enumerate}

\section{Citations and References}
When you use the citation search feature, references will be automatically managed. For example, you might cite a paper like this: \cite{example2024}. All citations will be automatically added to your \texttt{cite.bib} file.

\section{Conclusion}
This template demonstrates basic \LaTeX{} features including mathematical formulas, tables, and lists. The bibliography system is now set up to work seamlessly with Bibby AI's citation management. Feel free to modify it for your needs!

\bibliographystyle{plain}
\bibliography{cite}

\end{document}\documentclass{article}
\usepackage[utf8]{inputenc}
\usepackage{booktabs}
\usepackage{amsmath}
\usepackage{natbib}

\title{My Document}
\author{Your Name}
\date{\today}

\begin{document}

\maketitle

\section{Introduction}
This is a sample document demonstrating various \LaTeX{} features including mathematical equations, tables, and formatting. You can cite papers using the citation search feature, and they will be automatically added to your \texttt{cite.bib} file.

\section{Mathematical Expressions}
Here are some example equations:

\subsection{Quadratic Formula}
The quadratic formula is given by:
\[x = \frac{-b \pm \sqrt{b^2 - 4ac}}{2a}\]

\subsection{Matrix Example}
A sample matrix:
\[
\begin{pmatrix}
a_{11} & a_{12} & a_{13} \\
a_{21} & a_{22} & a_{23} \\
a_{31} & a_{32} & a_{33}
\end{pmatrix}
\]

\section{Tables}
Below is a sample table showing data:

\begin{table}[h]
\centering
\begin{tabular}{@{}lrrr@{}}
\toprule
\textbf{Category} & \textbf{Value 1} & \textbf{Value 2} & \textbf{Total} \\
\midrule
Item A & 10.5 & 20.3 & 30.8 \\
Item B & 15.2 & 25.1 & 40.3 \\
Item C & 22.1 & 12.7 & 34.8 \\
\bottomrule
\end{tabular}
\caption{Sample Data Table}
\label{tab:sample}
\end{table}

\section{Lists}
\subsection{Bullet Points}
\begin{itemize}
  \item First item
  \item Second item
  \item Third item with sub-items:
    \begin{itemize}
      \item Sub-item 1
      \item Sub-item 2
    \end{itemize}
\end{itemize}

\subsection{Numbered List}
\begin{enumerate}
  \item First step
  \item Second step
  \item Third step
\end{enumerate}

\section{Citations and References}
When you use the citation search feature, references will be automatically managed. For example, you might cite a paper like this: \cite{example2024}. All citations will be automatically added to your \texttt{cite.bib} file.

\section{Conclusion}
This template demonstrates basic \LaTeX{} features including mathematical formulas, tables, and lists. The bibliography system is now set up to work seamlessly with Bibby AI's citation management. Feel free to modify it for your needs!

\bibliographystyle{plain}
\bibliography{cite}

\end{document}\documentclass{article}
\usepackage[utf8]{inputenc}
\usepackage{booktabs}
\usepackage{amsmath}
\usepackage{natbib}

\title{My Document}
\author{Your Name}
\date{\today}

\begin{document}

\maketitle

\section{Introduction}
This is a sample document demonstrating various \LaTeX{} features including mathematical equations, tables, and formatting. You can cite papers using the citation search feature, and they will be automatically added to your \texttt{cite.bib} file.

\section{Mathematical Expressions}
Here are some example equations:

\subsection{Quadratic Formula}
The quadratic formula is given by:
\[x = \frac{-b \pm \sqrt{b^2 - 4ac}}{2a}\]

\subsection{Matrix Example}
A sample matrix:
\[
\begin{pmatrix}
a_{11} & a_{12} & a_{13} \\
a_{21} & a_{22} & a_{23} \\
a_{31} & a_{32} & a_{33}
\end{pmatrix}
\]

\section{Tables}
Below is a sample table showing data:

\begin{table}[h]
\centering
\begin{tabular}{@{}lrrr@{}}
\toprule
\textbf{Category} & \textbf{Value 1} & \textbf{Value 2} & \textbf{Total} \\
\midrule
Item A & 10.5 & 20.3 & 30.8 \\
Item B & 15.2 & 25.1 & 40.3 \\
Item C & 22.1 & 12.7 & 34.8 \\
\bottomrule
\end{tabular}
\caption{Sample Data Table}
\label{tab:sample}
\end{table}

\section{Lists}
\subsection{Bullet Points}
\begin{itemize}
  \item First item
  \item Second item
  \item Third item with sub-items:
    \begin{itemize}
      \item Sub-item 1
      \item Sub-item 2
    \end{itemize}
\end{itemize}

\subsection{Numbered List}
\begin{enumerate}
  \item First step
  \item Second step
  \item Third step
\end{enumerate}

\section{Citations and References}
When you use the citation search feature, references will be automatically managed. For example, you might cite a paper like this: \cite{example2024}. All citations will be automatically added to your \texttt{cite.bib} file.

\section{Conclusion}
This template demonstrates basic \LaTeX{} features including mathematical formulas, tables, and lists. The bibliography system is now set up to work seamlessly with Bibby AI's citation management. Feel free to modify it for your needs!

\bibliographystyle{plain}
\bibliography{cite}

\end{document}\documentclass{article}
\usepackage[utf8]{inputenc}
\usepackage{booktabs}
\usepackage{amsmath}
\usepackage{natbib}

\title{My Document}
\author{Your Name}
\date{\today}

\begin{document}

\maketitle

\section{Introduction}
This is a sample document demonstrating various \LaTeX{} features including mathematical equations, tables, and formatting. You can cite papers using the citation search feature, and they will be automatically added to your \texttt{cite.bib} file.

\section{Mathematical Expressions}
Here are some example equations:

\subsection{Quadratic Formula}
The quadratic formula is given by:
\[x = \frac{-b \pm \sqrt{b^2 - 4ac}}{2a}\]

\subsection{Matrix Example}
A sample matrix:
\[
\begin{pmatrix}
a_{11} & a_{12} & a_{13} \\
a_{21} & a_{22} & a_{23} \\
a_{31} & a_{32} & a_{33}
\end{pmatrix}
\]

\section{Tables}
Below is a sample table showing data:

\begin{table}[h]
\centering
\begin{tabular}{@{}lrrr@{}}
\toprule
\textbf{Category} & \textbf{Value 1} & \textbf{Value 2} & \textbf{Total} \\
\midrule
Item A & 10.5 & 20.3 & 30.8 \\
Item B & 15.2 & 25.1 & 40.3 \\
Item C & 22.1 & 12.7 & 34.8 \\
\bottomrule
\end{tabular}
\caption{Sample Data Table}
\label{tab:sample}
\end{table}

\section{Lists}
\subsection{Bullet Points}
\begin{itemize}
  \item First item
  \item Second item
  \item Third item with sub-items:
    \begin{itemize}
      \item Sub-item 1
      \item Sub-item 2
    \end{itemize}
\end{itemize}

\subsection{Numbered List}
\begin{enumerate}
  \item First step
  \item Second step
  \item Third step
\end{enumerate}

\section{Citations and References}
When you use the citation search feature, references will be automatically managed. For example, you might cite a paper like this: \cite{example2024}. All citations will be automatically added to your \texttt{cite.bib} file.

\section{Conclusion}
This template demonstrates basic \LaTeX{} features including mathematical formulas, tables, and lists. The bibliography system is now set up to work seamlessly with Bibby AI's citation management. Feel free to modify it for your needs!

\bibliographystyle{plain}
\bibliography{cite}

\end{document}\documentclass{article}
\usepackage[utf8]{inputenc}
\usepackage{booktabs}
\usepackage{amsmath}
\usepackage{natbib}

\title{My Document}
\author{Your Name}
\date{\today}

\begin{document}

\maketitle

\section{Introduction}
This is a sample document demonstrating various \LaTeX{} features including mathematical equations, tables, and formatting. You can cite papers using the citation search feature, and they will be automatically added to your \texttt{cite.bib} file.

\section{Mathematical Expressions}
Here are some example equations:

\subsection{Quadratic Formula}
The quadratic formula is given by:
\[x = \frac{-b \pm \sqrt{b^2 - 4ac}}{2a}\]

\subsection{Matrix Example}
A sample matrix:
\[
\begin{pmatrix}
a_{11} & a_{12} & a_{13} \\
a_{21} & a_{22} & a_{23} \\
a_{31} & a_{32} & a_{33}
\end{pmatrix}
\]

\section{Tables}
Below is a sample table showing data:

\begin{table}[h]
\centering
\begin{tabular}{@{}lrrr@{}}
\toprule
\textbf{Category} & \textbf{Value 1} & \textbf{Value 2} & \textbf{Total} \\
\midrule
Item A & 10.5 & 20.3 & 30.8 \\
Item B & 15.2 & 25.1 & 40.3 \\
Item C & 22.1 & 12.7 & 34.8 \\
\bottomrule
\end{tabular}
\caption{Sample Data Table}
\label{tab:sample}
\end{table}

\section{Lists}
\subsection{Bullet Points}
\begin{itemize}
  \item First item
  \item Second item
  \item Third item with sub-items:
    \begin{itemize}
      \item Sub-item 1
      \item Sub-item 2
    \end{itemize}
\end{itemize}

\subsection{Numbered List}
\begin{enumerate}
  \item First step
  \item Second step
  \item Third step
\end{enumerate}

\section{Citations and References}
When you use the citation search feature, references will be automatically managed. For example, you might cite a paper like this: \cite{example2024}. All citations will be automatically added to your \texttt{cite.bib} file.

\section{Conclusion}
This template demonstrates basic \LaTeX{} features including mathematical formulas, tables, and lists. The bibliography system is now set up to work seamlessly with Bibby AI's citation management. Feel free to modify it for your needs!

\bibliographystyle{plain}
\bibliography{cite}

\end{document}\documentclass{article}
\usepackage[utf8]{inputenc}
\usepackage{booktabs}
\usepackage{amsmath}
\usepackage{natbib}

\title{My Document}
\author{Your Name}
\date{\today}

\begin{document}

\maketitle

\section{Introduction}
This is a sample document demonstrating various \LaTeX{} features including mathematical equations, tables, and formatting. You can cite papers using the citation search feature, and they will be automatically added to your \texttt{cite.bib} file.

\section{Mathematical Expressions}
Here are some example equations:

\subsection{Quadratic Formula}
The quadratic formula is given by:
\[x = \frac{-b \pm \sqrt{b^2 - 4ac}}{2a}\]

\subsection{Matrix Example}
A sample matrix:
\[
\begin{pmatrix}
a_{11} & a_{12} & a_{13} \\
a_{21} & a_{22} & a_{23} \\
a_{31} & a_{32} & a_{33}
\end{pmatrix}
\]

\section{Tables}
Below is a sample table showing data:

\begin{table}[h]
\centering
\begin{tabular}{@{}lrrr@{}}
\toprule
\textbf{Category} & \textbf{Value 1} & \textbf{Value 2} & \textbf{Total} \\
\midrule
Item A & 10.5 & 20.3 & 30.8 \\
Item B & 15.2 & 25.1 & 40.3 \\
Item C & 22.1 & 12.7 & 34.8 \\
\bottomrule
\end{tabular}
\caption{Sample Data Table}
\label{tab:sample}
\end{table}

\section{Lists}
\subsection{Bullet Points}
\begin{itemize}
  \item First item
  \item Second item
  \item Third item with sub-items:
    \begin{itemize}
      \item Sub-item 1
      \item Sub-item 2
    \end{itemize}
\end{itemize}

\subsection{Numbered List}
\begin{enumerate}
  \item First step
  \item Second step
  \item Third step
\end{enumerate}

\section{Citations and References}
When you use the citation search feature, references will be automatically managed. For example, you might cite a paper like this: \cite{example2024}. All citations will be automatically added to your \texttt{cite.bib} file.

\section{Conclusion}
This template demonstrates basic \LaTeX{} features including mathematical formulas, tables, and lists. The bibliography system is now set up to work seamlessly with Bibby AI's citation management. Feel free to modify it for your needs!

\bibliographystyle{plain}
\bibliography{cite}

\end{document}\documentclass{article}
\usepackage[utf8]{inputenc}
\usepackage{booktabs}
\usepackage{amsmath}
\usepackage{natbib}

\title{My Document}
\author{Your Name}
\date{\today}

\begin{document}

\maketitle

\section{Introduction}
This is a sample document demonstrating various \LaTeX{} features including mathematical equations, tables, and formatting. You can cite papers using the citation search feature, and they will be automatically added to your \texttt{cite.bib} file.

\section{Mathematical Expressions}
Here are some example equations:

\subsection{Quadratic Formula}
The quadratic formula is given by:
\[x = \frac{-b \pm \sqrt{b^2 - 4ac}}{2a}\]

\subsection{Matrix Example}
A sample matrix:
\[
\begin{pmatrix}
a_{11} & a_{12} & a_{13} \\
a_{21} & a_{22} & a_{23} \\
a_{31} & a_{32} & a_{33}
\end{pmatrix}
\]

\section{Tables}
Below is a sample table showing data:

\begin{table}[h]
\centering
\begin{tabular}{@{}lrrr@{}}
\toprule
\textbf{Category} & \textbf{Value 1} & \textbf{Value 2} & \textbf{Total} \\
\midrule
Item A & 10.5 & 20.3 & 30.8 \\
Item B & 15.2 & 25.1 & 40.3 \\
Item C & 22.1 & 12.7 & 34.8 \\
\bottomrule
\end{tabular}
\caption{Sample Data Table}
\label{tab:sample}
\end{table}

\section{Lists}
\subsection{Bullet Points}
\begin{itemize}
  \item First item
  \item Second item
  \item Third item with sub-items:
    \begin{itemize}
      \item Sub-item 1
      \item Sub-item 2
    \end{itemize}
\end{itemize}

\subsection{Numbered List}
\begin{enumerate}
  \item First step
  \item Second step
  \item Third step
\end{enumerate}

\section{Citations and References}
When you use the citation search feature, references will be automatically managed. For example, you might cite a paper like this: \cite{example2024}. All citations will be automatically added to your \texttt{cite.bib} file.

\section{Conclusion}
This template demonstrates basic \LaTeX{} features including mathematical formulas, tables, and lists. The bibliography system is now set up to work seamlessly with Bibby AI's citation management. Feel free to modify it for your needs!

\bibliographystyle{plain}
\bibliography{cite}

\end{document}\documentclass{article}
\usepackage[utf8]{inputenc}
\usepackage{booktabs}
\usepackage{amsmath}
\usepackage{natbib}

\title{My Document}
\author{Your Name}
\date{\today}

\begin{document}

\maketitle

\section{Introduction}
This is a sample document demonstrating various \LaTeX{} features including mathematical equations, tables, and formatting. You can cite papers using the citation search feature, and they will be automatically added to your \texttt{cite.bib} file.

\section{Mathematical Expressions}
Here are some example equations:

\subsection{Quadratic Formula}
The quadratic formula is given by:
\[x = \frac{-b \pm \sqrt{b^2 - 4ac}}{2a}\]

\subsection{Matrix Example}
A sample matrix:
\[
\begin{pmatrix}
a_{11} & a_{12} & a_{13} \\
a_{21} & a_{22} & a_{23} \\
a_{31} & a_{32} & a_{33}
\end{pmatrix}
\]

\section{Tables}
Below is a sample table showing data:

\begin{table}[h]
\centering
\begin{tabular}{@{}lrrr@{}}
\toprule
\textbf{Category} & \textbf{Value 1} & \textbf{Value 2} & \textbf{Total} \\
\midrule
Item A & 10.5 & 20.3 & 30.8 \\
Item B & 15.2 & 25.1 & 40.3 \\
Item C & 22.1 & 12.7 & 34.8 \\
\bottomrule
\end{tabular}
\caption{Sample Data Table}
\label{tab:sample}
\end{table}

\section{Lists}
\subsection{Bullet Points}
\begin{itemize}
  \item First item
  \item Second item
  \item Third item with sub-items:
    \begin{itemize}
      \item Sub-item 1
      \item Sub-item 2
    \end{itemize}
\end{itemize}

\subsection{Numbered List}
\begin{enumerate}
  \item First step
  \item Second step
  \item Third step
\end{enumerate}

\section{Citations and References}
When you use the citation search feature, references will be automatically managed. For example, you might cite a paper like this: \cite{example2024}. All citations will be automatically added to your \texttt{cite.bib} file.

\section{Conclusion}
This template demonstrates basic \LaTeX{} features including mathematical formulas, tables, and lists. The bibliography system is now set up to work seamlessly with Bibby AI's citation management. Feel free to modify it for your needs!

\bibliographystyle{plain}
\bibliography{cite}

\end{document}\documentclass{article}
\usepackage[utf8]{inputenc}
\usepackage{booktabs}
\usepackage{amsmath}
\usepackage{natbib}

\title{My Document}
\author{Your Name}
\date{\today}

\begin{document}

\maketitle

\section{Introduction}
This is a sample document demonstrating various \LaTeX{} features including mathematical equations, tables, and formatting. You can cite papers using the citation search feature, and they will be automatically added to your \texttt{cite.bib} file.

\section{Mathematical Expressions}
Here are some example equations:

\subsection{Quadratic Formula}
The quadratic formula is given by:
\[x = \frac{-b \pm \sqrt{b^2 - 4ac}}{2a}\]

\subsection{Matrix Example}
A sample matrix:
\[
\begin{pmatrix}
a_{11} & a_{12} & a_{13} \\
a_{21} & a_{22} & a_{23} \\
a_{31} & a_{32} & a_{33}
\end{pmatrix}
\]

\section{Tables}
Below is a sample table showing data:

\begin{table}[h]
\centering
\begin{tabular}{@{}lrrr@{}}
\toprule
\textbf{Category} & \textbf{Value 1} & \textbf{Value 2} & \textbf{Total} \\
\midrule
Item A & 10.5 & 20.3 & 30.8 \\
Item B & 15.2 & 25.1 & 40.3 \\
Item C & 22.1 & 12.7 & 34.8 \\
\bottomrule
\end{tabular}
\caption{Sample Data Table}
\label{tab:sample}
\end{table}

\section{Lists}
\subsection{Bullet Points}
\begin{itemize}
  \item First item
  \item Second item
  \item Third item with sub-items:
    \begin{itemize}
      \item Sub-item 1
      \item Sub-item 2
    \end{itemize}
\end{itemize}

\subsection{Numbered List}
\begin{enumerate}
  \item First step
  \item Second step
  \item Third step
\end{enumerate}

\section{Citations and References}
When you use the citation search feature, references will be automatically managed. For example, you might cite a paper like this: \cite{example2024}. All citations will be automatically added to your \texttt{cite.bib} file.

\section{Conclusion}
This template demonstrates basic \LaTeX{} features including mathematical formulas, tables, and lists. The bibliography system is now set up to work seamlessly with Bibby AI's citation management. Feel free to modify it for your needs!

\bibliographystyle{plain}
\bibliography{cite}

\end{document}\documentclass{article}
\usepackage[utf8]{inputenc}
\usepackage{booktabs}
\usepackage{amsmath}
\usepackage{natbib}

\title{My Document}
\author{Your Name}
\date{\today}

\begin{document}

\maketitle

\section{Introduction}
This is a sample document demonstrating various \LaTeX{} features including mathematical equations, tables, and formatting. You can cite papers using the citation search feature, and they will be automatically added to your \texttt{cite.bib} file.

\section{Mathematical Expressions}
Here are some example equations:

\subsection{Quadratic Formula}
The quadratic formula is given by:
\[x = \frac{-b \pm \sqrt{b^2 - 4ac}}{2a}\]

\subsection{Matrix Example}
A sample matrix:
\[
\begin{pmatrix}
a_{11} & a_{12} & a_{13} \\
a_{21} & a_{22} & a_{23} \\
a_{31} & a_{32} & a_{33}
\end{pmatrix}
\]

\section{Tables}
Below is a sample table showing data:

\begin{table}[h]
\centering
\begin{tabular}{@{}lrrr@{}}
\toprule
\textbf{Category} & \textbf{Value 1} & \textbf{Value 2} & \textbf{Total} \\
\midrule
Item A & 10.5 & 20.3 & 30.8 \\
Item B & 15.2 & 25.1 & 40.3 \\
Item C & 22.1 & 12.7 & 34.8 \\
\bottomrule
\end{tabular}
\caption{Sample Data Table}
\label{tab:sample}
\end{table}

\section{Lists}
\subsection{Bullet Points}
\begin{itemize}
  \item First item
  \item Second item
  \item Third item with sub-items:
    \begin{itemize}
      \item Sub-item 1
      \item Sub-item 2
    \end{itemize}
\end{itemize}

\subsection{Numbered List}
\begin{enumerate}
  \item First step
  \item Second step
  \item Third step
\end{enumerate}

\section{Citations and References}
When you use the citation search feature, references will be automatically managed. For example, you might cite a paper like this: \cite{example2024}. All citations will be automatically added to your \texttt{cite.bib} file.

\section{Conclusion}
This template demonstrates basic \LaTeX{} features including mathematical formulas, tables, and lists. The bibliography system is now set up to work seamlessly with Bibby AI's citation management. Feel free to modify it for your needs!

\bibliographystyle{plain}
\bibliography{cite}

\end{document}\documentclass{article}
\usepackage[utf8]{inputenc}
\usepackage{booktabs}
\usepackage{amsmath}
\usepackage{natbib}

\title{My Document}
\author{Your Name}
\date{\today}

\begin{document}

\maketitle

\section{Introduction}
This is a sample document demonstrating various \LaTeX{} features including mathematical equations, tables, and formatting. You can cite papers using the citation search feature, and they will be automatically added to your \texttt{cite.bib} file.

\section{Mathematical Expressions}
Here are some example equations:

\subsection{Quadratic Formula}
The quadratic formula is given by:
\[x = \frac{-b \pm \sqrt{b^2 - 4ac}}{2a}\]

\subsection{Matrix Example}
A sample matrix:
\[
\begin{pmatrix}
a_{11} & a_{12} & a_{13} \\
a_{21} & a_{22} & a_{23} \\
a_{31} & a_{32} & a_{33}
\end{pmatrix}
\]

\section{Tables}
Below is a sample table showing data:

\begin{table}[h]
\centering
\begin{tabular}{@{}lrrr@{}}
\toprule
\textbf{Category} & \textbf{Value 1} & \textbf{Value 2} & \textbf{Total} \\
\midrule
Item A & 10.5 & 20.3 & 30.8 \\
Item B & 15.2 & 25.1 & 40.3 \\
Item C & 22.1 & 12.7 & 34.8 \\
\bottomrule
\end{tabular}
\caption{Sample Data Table}
\label{tab:sample}
\end{table}

\section{Lists}
\subsection{Bullet Points}
\begin{itemize}
  \item First item
  \item Second item
  \item Third item with sub-items:
    \begin{itemize}
      \item Sub-item 1
      \item Sub-item 2
    \end{itemize}
\end{itemize}

\subsection{Numbered List}
\begin{enumerate}
  \item First step
  \item Second step
  \item Third step
\end{enumerate}

\section{Citations and References}
When you use the citation search feature, references will be automatically managed. For example, you might cite a paper like this: \cite{example2024}. All citations will be automatically added to your \texttt{cite.bib} file.

\section{Conclusion}
This template demonstrates basic \LaTeX{} features including mathematical formulas, tables, and lists. The bibliography system is now set up to work seamlessly with Bibby AI's citation management. Feel free to modify it for your needs!

\bibliographystyle{plain}
\bibliography{cite}

\end{document}\documentclass{article}
\usepackage[utf8]{inputenc}
\usepackage{booktabs}
\usepackage{amsmath}
\usepackage{natbib}

\title{My Document}
\author{Your Name}
\date{\today}

\begin{document}

\maketitle

\section{Introduction}
This is a sample document demonstrating various \LaTeX{} features including mathematical equations, tables, and formatting. You can cite papers using the citation search feature, and they will be automatically added to your \texttt{cite.bib} file.

\section{Mathematical Expressions}
Here are some example equations:

\subsection{Quadratic Formula}
The quadratic formula is given by:
\[x = \frac{-b \pm \sqrt{b^2 - 4ac}}{2a}\]

\subsection{Matrix Example}
A sample matrix:
\[
\begin{pmatrix}
a_{11} & a_{12} & a_{13} \\
a_{21} & a_{22} & a_{23} \\
a_{31} & a_{32} & a_{33}
\end{pmatrix}
\]

\section{Tables}
Below is a sample table showing data:

\begin{table}[h]
\centering
\begin{tabular}{@{}lrrr@{}}
\toprule
\textbf{Category} & \textbf{Value 1} & \textbf{Value 2} & \textbf{Total} \\
\midrule
Item A & 10.5 & 20.3 & 30.8 \\
Item B & 15.2 & 25.1 & 40.3 \\
Item C & 22.1 & 12.7 & 34.8 \\
\bottomrule
\end{tabular}
\caption{Sample Data Table}
\label{tab:sample}
\end{table}

\section{Lists}
\subsection{Bullet Points}
\begin{itemize}
  \item First item
  \item Second item
  \item Third item with sub-items:
    \begin{itemize}
      \item Sub-item 1
      \item Sub-item 2
    \end{itemize}
\end{itemize}

\subsection{Numbered List}
\begin{enumerate}
  \item First step
  \item Second step
  \item Third step
\end{enumerate}

\section{Citations and References}
When you use the citation search feature, references will be automatically managed. For example, you might cite a paper like this: \cite{example2024}. All citations will be automatically added to your \texttt{cite.bib} file.

\section{Conclusion}
This template demonstrates basic \LaTeX{} features including mathematical formulas, tables, and lists. The bibliography system is now set up to work seamlessly with Bibby AI's citation management. Feel free to modify it for your needs!

\bibliographystyle{plain}
\bibliography{cite}

\end{document}\documentclass{article}
\usepackage[utf8]{inputenc}
\usepackage{booktabs}
\usepackage{amsmath}
\usepackage{natbib}

\title{My Document}
\author{Your Name}
\date{\today}

\begin{document}

\maketitle

\section{Introduction}
This is a sample document demonstrating various \LaTeX{} features including mathematical equations, tables, and formatting. You can cite papers using the citation search feature, and they will be automatically added to your \texttt{cite.bib} file.

\section{Mathematical Expressions}
Here are some example equations:

\subsection{Quadratic Formula}
The quadratic formula is given by:
\[x = \frac{-b \pm \sqrt{b^2 - 4ac}}{2a}\]

\subsection{Matrix Example}
A sample matrix:
\[
\begin{pmatrix}
a_{11} & a_{12} & a_{13} \\
a_{21} & a_{22} & a_{23} \\
a_{31} & a_{32} & a_{33}
\end{pmatrix}
\]

\section{Tables}
Below is a sample table showing data:

\begin{table}[h]
\centering
\begin{tabular}{@{}lrrr@{}}
\toprule
\textbf{Category} & \textbf{Value 1} & \textbf{Value 2} & \textbf{Total} \\
\midrule
Item A & 10.5 & 20.3 & 30.8 \\
Item B & 15.2 & 25.1 & 40.3 \\
Item C & 22.1 & 12.7 & 34.8 \\
\bottomrule
\end{tabular}
\caption{Sample Data Table}
\label{tab:sample}
\end{table}

\section{Lists}
\subsection{Bullet Points}
\begin{itemize}
  \item First item
  \item Second item
  \item Third item with sub-items:
    \begin{itemize}
      \item Sub-item 1
      \item Sub-item 2
    \end{itemize}
\end{itemize}

\subsection{Numbered List}
\begin{enumerate}
  \item First step
  \item Second step
  \item Third step
\end{enumerate}

\section{Citations and References}
When you use the citation search feature, references will be automatically managed. For example, you might cite a paper like this: \cite{example2024}. All citations will be automatically added to your \texttt{cite.bib} file.

\section{Conclusion}
This template demonstrates basic \LaTeX{} features including mathematical formulas, tables, and lists. The bibliography system is now set up to work seamlessly with Bibby AI's citation management. Feel free to modify it for your needs!

\bibliographystyle{plain}
\bibliography{cite}

\end{document}\documentclass{article}
\usepackage[utf8]{inputenc}
\usepackage{booktabs}
\usepackage{amsmath}
\usepackage{natbib}

\title{My Document}
\author{Your Name}
\date{\today}

\begin{document}

\maketitle

\section{Introduction}
This is a sample document demonstrating various \LaTeX{} features including mathematical equations, tables, and formatting. You can cite papers using the citation search feature, and they will be automatically added to your \texttt{cite.bib} file.

\section{Mathematical Expressions}
Here are some example equations:

\subsection{Quadratic Formula}
The quadratic formula is given by:
\[x = \frac{-b \pm \sqrt{b^2 - 4ac}}{2a}\]

\subsection{Matrix Example}
A sample matrix:
\[
\begin{pmatrix}
a_{11} & a_{12} & a_{13} \\
a_{21} & a_{22} & a_{23} \\
a_{31} & a_{32} & a_{33}
\end{pmatrix}
\]

\section{Tables}
Below is a sample table showing data:

\begin{table}[h]
\centering
\begin{tabular}{@{}lrrr@{}}
\toprule
\textbf{Category} & \textbf{Value 1} & \textbf{Value 2} & \textbf{Total} \\
\midrule
Item A & 10.5 & 20.3 & 30.8 \\
Item B & 15.2 & 25.1 & 40.3 \\
Item C & 22.1 & 12.7 & 34.8 \\
\bottomrule
\end{tabular}
\caption{Sample Data Table}
\label{tab:sample}
\end{table}

\section{Lists}
\subsection{Bullet Points}
\begin{itemize}
  \item First item
  \item Second item
  \item Third item with sub-items:
    \begin{itemize}
      \item Sub-item 1
      \item Sub-item 2
    \end{itemize}
\end{itemize}

\subsection{Numbered List}
\begin{enumerate}
  \item First step
  \item Second step
  \item Third step
\end{enumerate}

\section{Citations and References}
When you use the citation search feature, references will be automatically managed. For example, you might cite a paper like this: \cite{example2024}. All citations will be automatically added to your \texttt{cite.bib} file.

\section{Conclusion}
This template demonstrates basic \LaTeX{} features including mathematical formulas, tables, and lists. The bibliography system is now set up to work seamlessly with Bibby AI's citation management. Feel free to modify it for your needs!

\bibliographystyle{plain}
\bibliography{cite}

\end{document}\documentclass{article}
\usepackage[utf8]{inputenc}
\usepackage{booktabs}
\usepackage{amsmath}
\usepackage{natbib}

\title{My Document}
\author{Your Name}
\date{\today}

\begin{document}

\maketitle

\section{Introduction}
This is a sample document demonstrating various \LaTeX{} features including mathematical equations, tables, and formatting. You can cite papers using the citation search feature, and they will be automatically added to your \texttt{cite.bib} file.

\section{Mathematical Expressions}
Here are some example equations:

\subsection{Quadratic Formula}
The quadratic formula is given by:
\[x = \frac{-b \pm \sqrt{b^2 - 4ac}}{2a}\]

\subsection{Matrix Example}
A sample matrix:
\[
\begin{pmatrix}
a_{11} & a_{12} & a_{13} \\
a_{21} & a_{22} & a_{23} \\
a_{31} & a_{32} & a_{33}
\end{pmatrix}
\]

\section{Tables}
Below is a sample table showing data:

\begin{table}[h]
\centering
\begin{tabular}{@{}lrrr@{}}
\toprule
\textbf{Category} & \textbf{Value 1} & \textbf{Value 2} & \textbf{Total} \\
\midrule
Item A & 10.5 & 20.3 & 30.8 \\
Item B & 15.2 & 25.1 & 40.3 \\
Item C & 22.1 & 12.7 & 34.8 \\
\bottomrule
\end{tabular}
\caption{Sample Data Table}
\label{tab:sample}
\end{table}

\section{Lists}
\subsection{Bullet Points}
\begin{itemize}
  \item First item
  \item Second item
  \item Third item with sub-items:
    \begin{itemize}
      \item Sub-item 1
      \item Sub-item 2
    \end{itemize}
\end{itemize}

\subsection{Numbered List}
\begin{enumerate}
  \item First step
  \item Second step
  \item Third step
\end{enumerate}

\section{Citations and References}
When you use the citation search feature, references will be automatically managed. For example, you might cite a paper like this: \cite{example2024}. All citations will be automatically added to your \texttt{cite.bib} file.

\section{Conclusion}
This template demonstrates basic \LaTeX{} features including mathematical formulas, tables, and lists. The bibliography system is now set up to work seamlessly with Bibby AI's citation management. Feel free to modify it for your needs!

\bibliographystyle{plain}
\bibliography{cite}

\end{document}\documentclass{article}
\usepackage[utf8]{inputenc}
\usepackage{booktabs}
\usepackage{amsmath}
\usepackage{natbib}

\title{My Document}
\author{Your Name}
\date{\today}

\begin{document}

\maketitle

\section{Introduction}
This is a sample document demonstrating various \LaTeX{} features including mathematical equations, tables, and formatting. You can cite papers using the citation search feature, and they will be automatically added to your \texttt{cite.bib} file.

\section{Mathematical Expressions}
Here are some example equations:

\subsection{Quadratic Formula}
The quadratic formula is given by:
\[x = \frac{-b \pm \sqrt{b^2 - 4ac}}{2a}\]

\subsection{Matrix Example}
A sample matrix:
\[
\begin{pmatrix}
a_{11} & a_{12} & a_{13} \\
a_{21} & a_{22} & a_{23} \\
a_{31} & a_{32} & a_{33}
\end{pmatrix}
\]

\section{Tables}
Below is a sample table showing data:

\begin{table}[h]
\centering
\begin{tabular}{@{}lrrr@{}}
\toprule
\textbf{Category} & \textbf{Value 1} & \textbf{Value 2} & \textbf{Total} \\
\midrule
Item A & 10.5 & 20.3 & 30.8 \\
Item B & 15.2 & 25.1 & 40.3 \\
Item C & 22.1 & 12.7 & 34.8 \\
\bottomrule
\end{tabular}
\caption{Sample Data Table}
\label{tab:sample}
\end{table}

\section{Lists}
\subsection{Bullet Points}
\begin{itemize}
  \item First item
  \item Second item
  \item Third item with sub-items:
    \begin{itemize}
      \item Sub-item 1
      \item Sub-item 2
    \end{itemize}
\end{itemize}

\subsection{Numbered List}
\begin{enumerate}
  \item First step
  \item Second step
  \item Third step
\end{enumerate}

\section{Citations and References}
When you use the citation search feature, references will be automatically managed. For example, you might cite a paper like this: \cite{example2024}. All citations will be automatically added to your \texttt{cite.bib} file.

\section{Conclusion}
This template demonstrates basic \LaTeX{} features including mathematical formulas, tables, and lists. The bibliography system is now set up to work seamlessly with Bibby AI's citation management. Feel free to modify it for your needs!

\bibliographystyle{plain}
\bibliography{cite}

\end{document}\documentclass{article}
\usepackage[utf8]{inputenc}
\usepackage{booktabs}
\usepackage{amsmath}
\usepackage{natbib}

\title{My Document}
\author{Your Name}
\date{\today}

\begin{document}

\maketitle

\section{Introduction}
This is a sample document demonstrating various \LaTeX{} features including mathematical equations, tables, and formatting. You can cite papers using the citation search feature, and they will be automatically added to your \texttt{cite.bib} file.

\section{Mathematical Expressions}
Here are some example equations:

\subsection{Quadratic Formula}
The quadratic formula is given by:
\[x = \frac{-b \pm \sqrt{b^2 - 4ac}}{2a}\]

\subsection{Matrix Example}
A sample matrix:
\[
\begin{pmatrix}
a_{11} & a_{12} & a_{13} \\
a_{21} & a_{22} & a_{23} \\
a_{31} & a_{32} & a_{33}
\end{pmatrix}
\]

\section{Tables}
Below is a sample table showing data:

\begin{table}[h]
\centering
\begin{tabular}{@{}lrrr@{}}
\toprule
\textbf{Category} & \textbf{Value 1} & \textbf{Value 2} & \textbf{Total} \\
\midrule
Item A & 10.5 & 20.3 & 30.8 \\
Item B & 15.2 & 25.1 & 40.3 \\
Item C & 22.1 & 12.7 & 34.8 \\
\bottomrule
\end{tabular}
\caption{Sample Data Table}
\label{tab:sample}
\end{table}

\section{Lists}
\subsection{Bullet Points}
\begin{itemize}
  \item First item
  \item Second item
  \item Third item with sub-items:
    \begin{itemize}
      \item Sub-item 1
      \item Sub-item 2
    \end{itemize}
\end{itemize}

\subsection{Numbered List}
\begin{enumerate}
  \item First step
  \item Second step
  \item Third step
\end{enumerate}

\section{Citations and References}
When you use the citation search feature, references will be automatically managed. For example, you might cite a paper like this: \cite{example2024}. All citations will be automatically added to your \texttt{cite.bib} file.

\section{Conclusion}
This template demonstrates basic \LaTeX{} features including mathematical formulas, tables, and lists. The bibliography system is now set up to work seamlessly with Bibby AI's citation management. Feel free to modify it for your needs!

\bibliographystyle{plain}
\bibliography{cite}

\end{document}\documentclass{article}
\usepackage[utf8]{inputenc}
\usepackage{booktabs}
\usepackage{amsmath}
\usepackage{natbib}

\title{My Document}
\author{Your Name}
\date{\today}

\begin{document}

\maketitle

\section{Introduction}
This is a sample document demonstrating various \LaTeX{} features including mathematical equations, tables, and formatting. You can cite papers using the citation search feature, and they will be automatically added to your \texttt{cite.bib} file.

\section{Mathematical Expressions}
Here are some example equations:

\subsection{Quadratic Formula}
The quadratic formula is given by:
\[x = \frac{-b \pm \sqrt{b^2 - 4ac}}{2a}\]

\subsection{Matrix Example}
A sample matrix:
\[
\begin{pmatrix}
a_{11} & a_{12} & a_{13} \\
a_{21} & a_{22} & a_{23} \\
a_{31} & a_{32} & a_{33}
\end{pmatrix}
\]

\section{Tables}
Below is a sample table showing data:

\begin{table}[h]
\centering
\begin{tabular}{@{}lrrr@{}}
\toprule
\textbf{Category} & \textbf{Value 1} & \textbf{Value 2} & \textbf{Total} \\
\midrule
Item A & 10.5 & 20.3 & 30.8 \\
Item B & 15.2 & 25.1 & 40.3 \\
Item C & 22.1 & 12.7 & 34.8 \\
\bottomrule
\end{tabular}
\caption{Sample Data Table}
\label{tab:sample}
\end{table}

\section{Lists}
\subsection{Bullet Points}
\begin{itemize}
  \item First item
  \item Second item
  \item Third item with sub-items:
    \begin{itemize}
      \item Sub-item 1
      \item Sub-item 2
    \end{itemize}
\end{itemize}

\subsection{Numbered List}
\begin{enumerate}
  \item First step
  \item Second step
  \item Third step
\end{enumerate}

\section{Citations and References}
When you use the citation search feature, references will be automatically managed. For example, you might cite a paper like this: \cite{example2024}. All citations will be automatically added to your \texttt{cite.bib} file.

\section{Conclusion}
This template demonstrates basic \LaTeX{} features including mathematical formulas, tables, and lists. The bibliography system is now set up to work seamlessly with Bibby AI's citation management. Feel free to modify it for your needs!

\bibliographystyle{plain}
\bibliography{cite}

\end{document}\documentclass{article}
\usepackage[utf8]{inputenc}
\usepackage{booktabs}
\usepackage{amsmath}
\usepackage{natbib}

\title{My Document}
\author{Your Name}
\date{\today}

\begin{document}

\maketitle

\section{Introduction}
This is a sample document demonstrating various \LaTeX{} features including mathematical equations, tables, and formatting. You can cite papers using the citation search feature, and they will be automatically added to your \texttt{cite.bib} file.

\section{Mathematical Expressions}
Here are some example equations:

\subsection{Quadratic Formula}
The quadratic formula is given by:
\[x = \frac{-b \pm \sqrt{b^2 - 4ac}}{2a}\]

\subsection{Matrix Example}
A sample matrix:
\[
\begin{pmatrix}
a_{11} & a_{12} & a_{13} \\
a_{21} & a_{22} & a_{23} \\
a_{31} & a_{32} & a_{33}
\end{pmatrix}
\]

\section{Tables}
Below is a sample table showing data:

\begin{table}[h]
\centering
\begin{tabular}{@{}lrrr@{}}
\toprule
\textbf{Category} & \textbf{Value 1} & \textbf{Value 2} & \textbf{Total} \\
\midrule
Item A & 10.5 & 20.3 & 30.8 \\
Item B & 15.2 & 25.1 & 40.3 \\
Item C & 22.1 & 12.7 & 34.8 \\
\bottomrule
\end{tabular}
\caption{Sample Data Table}
\label{tab:sample}
\end{table}

\section{Lists}
\subsection{Bullet Points}
\begin{itemize}
  \item First item
  \item Second item
  \item Third item with sub-items:
    \begin{itemize}
      \item Sub-item 1
      \item Sub-item 2
    \end{itemize}
\end{itemize}

\subsection{Numbered List}
\begin{enumerate}
  \item First step
  \item Second step
  \item Third step
\end{enumerate}

\section{Citations and References}
When you use the citation search feature, references will be automatically managed. For example, you might cite a paper like this: \cite{example2024}. All citations will be automatically added to your \texttt{cite.bib} file.

\section{Conclusion}
This template demonstrates basic \LaTeX{} features including mathematical formulas, tables, and lists. The bibliography system is now set up to work seamlessly with Bibby AI's citation management. Feel free to modify it for your needs!

\bibliographystyle{plain}
\bibliography{cite}

\end{document}\documentclass{article}
\usepackage[utf8]{inputenc}
\usepackage{booktabs}
\usepackage{amsmath}
\usepackage{natbib}

\title{My Document}
\author{Your Name}
\date{\today}

\begin{document}

\maketitle

\section{Introduction}
This is a sample document demonstrating various \LaTeX{} features including mathematical equations, tables, and formatting. You can cite papers using the citation search feature, and they will be automatically added to your \texttt{cite.bib} file.

\section{Mathematical Expressions}
Here are some example equations:

\subsection{Quadratic Formula}
The quadratic formula is given by:
\[x = \frac{-b \pm \sqrt{b^2 - 4ac}}{2a}\]

\subsection{Matrix Example}
A sample matrix:
\[
\begin{pmatrix}
a_{11} & a_{12} & a_{13} \\
a_{21} & a_{22} & a_{23} \\
a_{31} & a_{32} & a_{33}
\end{pmatrix}
\]

\section{Tables}
Below is a sample table showing data:

\begin{table}[h]
\centering
\begin{tabular}{@{}lrrr@{}}
\toprule
\textbf{Category} & \textbf{Value 1} & \textbf{Value 2} & \textbf{Total} \\
\midrule
Item A & 10.5 & 20.3 & 30.8 \\
Item B & 15.2 & 25.1 & 40.3 \\
Item C & 22.1 & 12.7 & 34.8 \\
\bottomrule
\end{tabular}
\caption{Sample Data Table}
\label{tab:sample}
\end{table}

\section{Lists}
\subsection{Bullet Points}
\begin{itemize}
  \item First item
  \item Second item
  \item Third item with sub-items:
    \begin{itemize}
      \item Sub-item 1
      \item Sub-item 2
    \end{itemize}
\end{itemize}

\subsection{Numbered List}
\begin{enumerate}
  \item First step
  \item Second step
  \item Third step
\end{enumerate}

\section{Citations and References}
When you use the citation search feature, references will be automatically managed. For example, you might cite a paper like this: \cite{example2024}. All citations will be automatically added to your \texttt{cite.bib} file.

\section{Conclusion}
This template demonstrates basic \LaTeX{} features including mathematical formulas, tables, and lists. The bibliography system is now set up to work seamlessly with Bibby AI's citation management. Feel free to modify it for your needs!

\bibliographystyle{plain}
\bibliography{cite}

\end{document}\documentclass{article}
\usepackage[utf8]{inputenc}
\usepackage{booktabs}
\usepackage{amsmath}
\usepackage{natbib}

\title{My Document}
\author{Your Name}
\date{\today}

\begin{document}

\maketitle

\section{Introduction}
This is a sample document demonstrating various \LaTeX{} features including mathematical equations, tables, and formatting. You can cite papers using the citation search feature, and they will be automatically added to your \texttt{cite.bib} file.

\section{Mathematical Expressions}
Here are some example equations:

\subsection{Quadratic Formula}
The quadratic formula is given by:
\[x = \frac{-b \pm \sqrt{b^2 - 4ac}}{2a}\]

\subsection{Matrix Example}
A sample matrix:
\[
\begin{pmatrix}
a_{11} & a_{12} & a_{13} \\
a_{21} & a_{22} & a_{23} \\
a_{31} & a_{32} & a_{33}
\end{pmatrix}
\]

\section{Tables}
Below is a sample table showing data:

\begin{table}[h]
\centering
\begin{tabular}{@{}lrrr@{}}
\toprule
\textbf{Category} & \textbf{Value 1} & \textbf{Value 2} & \textbf{Total} \\
\midrule
Item A & 10.5 & 20.3 & 30.8 \\
Item B & 15.2 & 25.1 & 40.3 \\
Item C & 22.1 & 12.7 & 34.8 \\
\bottomrule
\end{tabular}
\caption{Sample Data Table}
\label{tab:sample}
\end{table}

\section{Lists}
\subsection{Bullet Points}
\begin{itemize}
  \item First item
  \item Second item
  \item Third item with sub-items:
    \begin{itemize}
      \item Sub-item 1
      \item Sub-item 2
    \end{itemize}
\end{itemize}

\subsection{Numbered List}
\begin{enumerate}
  \item First step
  \item Second step
  \item Third step
\end{enumerate}

\section{Citations and References}
When you use the citation search feature, references will be automatically managed. For example, you might cite a paper like this: \cite{example2024}. All citations will be automatically added to your \texttt{cite.bib} file.

\section{Conclusion}
This template demonstrates basic \LaTeX{} features including mathematical formulas, tables, and lists. The bibliography system is now set up to work seamlessly with Bibby AI's citation management. Feel free to modify it for your needs!

\bibliographystyle{plain}
\bibliography{cite}

\end{document}\documentclass{article}
\usepackage[utf8]{inputenc}
\usepackage{booktabs}
\usepackage{amsmath}
\usepackage{natbib}

\title{My Document}
\author{Your Name}
\date{\today}

\begin{document}

\maketitle

\section{Introduction}
This is a sample document demonstrating various \LaTeX{} features including mathematical equations, tables, and formatting. You can cite papers using the citation search feature, and they will be automatically added to your \texttt{cite.bib} file.

\section{Mathematical Expressions}
Here are some example equations:

\subsection{Quadratic Formula}
The quadratic formula is given by:
\[x = \frac{-b \pm \sqrt{b^2 - 4ac}}{2a}\]

\subsection{Matrix Example}
A sample matrix:
\[
\begin{pmatrix}
a_{11} & a_{12} & a_{13} \\
a_{21} & a_{22} & a_{23} \\
a_{31} & a_{32} & a_{33}
\end{pmatrix}
\]

\section{Tables}
Below is a sample table showing data:

\begin{table}[h]
\centering
\begin{tabular}{@{}lrrr@{}}
\toprule
\textbf{Category} & \textbf{Value 1} & \textbf{Value 2} & \textbf{Total} \\
\midrule
Item A & 10.5 & 20.3 & 30.8 \\
Item B & 15.2 & 25.1 & 40.3 \\
Item C & 22.1 & 12.7 & 34.8 \\
\bottomrule
\end{tabular}
\caption{Sample Data Table}
\label{tab:sample}
\end{table}

\section{Lists}
\subsection{Bullet Points}
\begin{itemize}
  \item First item
  \item Second item
  \item Third item with sub-items:
    \begin{itemize}
      \item Sub-item 1
      \item Sub-item 2
    \end{itemize}
\end{itemize}

\subsection{Numbered List}
\begin{enumerate}
  \item First step
  \item Second step
  \item Third step
\end{enumerate}

\section{Citations and References}
When you use the citation search feature, references will be automatically managed. For example, you might cite a paper like this: \cite{example2024}. All citations will be automatically added to your \texttt{cite.bib} file.

\section{Conclusion}
This template demonstrates basic \LaTeX{} features including mathematical formulas, tables, and lists. The bibliography system is now set up to work seamlessly with Bibby AI's citation management. Feel free to modify it for your needs!

\bibliographystyle{plain}
\bibliography{cite}

\end{document}\documentclass{article}
\usepackage[utf8]{inputenc}
\usepackage{booktabs}
\usepackage{amsmath}
\usepackage{natbib}

\title{My Document}
\author{Your Name}
\date{\today}

\begin{document}

\maketitle

\section{Introduction}
This is a sample document demonstrating various \LaTeX{} features including mathematical equations, tables, and formatting. You can cite papers using the citation search feature, and they will be automatically added to your \texttt{cite.bib} file.

\section{Mathematical Expressions}
Here are some example equations:

\subsection{Quadratic Formula}
The quadratic formula is given by:
\[x = \frac{-b \pm \sqrt{b^2 - 4ac}}{2a}\]

\subsection{Matrix Example}
A sample matrix:
\[
\begin{pmatrix}
a_{11} & a_{12} & a_{13} \\
a_{21} & a_{22} & a_{23} \\
a_{31} & a_{32} & a_{33}
\end{pmatrix}
\]

\section{Tables}
Below is a sample table showing data:

\begin{table}[h]
\centering
\begin{tabular}{@{}lrrr@{}}
\toprule
\textbf{Category} & \textbf{Value 1} & \textbf{Value 2} & \textbf{Total} \\
\midrule
Item A & 10.5 & 20.3 & 30.8 \\
Item B & 15.2 & 25.1 & 40.3 \\
Item C & 22.1 & 12.7 & 34.8 \\
\bottomrule
\end{tabular}
\caption{Sample Data Table}
\label{tab:sample}
\end{table}

\section{Lists}
\subsection{Bullet Points}
\begin{itemize}
  \item First item
  \item Second item
  \item Third item with sub-items:
    \begin{itemize}
      \item Sub-item 1
      \item Sub-item 2
    \end{itemize}
\end{itemize}

\subsection{Numbered List}
\begin{enumerate}
  \item First step
  \item Second step
  \item Third step
\end{enumerate}

\section{Citations and References}
When you use the citation search feature, references will be automatically managed. For example, you might cite a paper like this: \cite{example2024}. All citations will be automatically added to your \texttt{cite.bib} file.

\section{Conclusion}
This template demonstrates basic \LaTeX{} features including mathematical formulas, tables, and lists. The bibliography system is now set up to work seamlessly with Bibby AI's citation management. Feel free to modify it for your needs!

\bibliographystyle{plain}
\bibliography{cite}

\end{document}\documentclass{article}
\usepackage[utf8]{inputenc}
\usepackage{booktabs}
\usepackage{amsmath}
\usepackage{natbib}

\title{My Document}
\author{Your Name}
\date{\today}

\begin{document}

\maketitle

\section{Introduction}
This is a sample document demonstrating various \LaTeX{} features including mathematical equations, tables, and formatting. You can cite papers using the citation search feature, and they will be automatically added to your \texttt{cite.bib} file.

\section{Mathematical Expressions}
Here are some example equations:

\subsection{Quadratic Formula}
The quadratic formula is given by:
\[x = \frac{-b \pm \sqrt{b^2 - 4ac}}{2a}\]

\subsection{Matrix Example}
A sample matrix:
\[
\begin{pmatrix}
a_{11} & a_{12} & a_{13} \\
a_{21} & a_{22} & a_{23} \\
a_{31} & a_{32} & a_{33}
\end{pmatrix}
\]

\section{Tables}
Below is a sample table showing data:

\begin{table}[h]
\centering
\begin{tabular}{@{}lrrr@{}}
\toprule
\textbf{Category} & \textbf{Value 1} & \textbf{Value 2} & \textbf{Total} \\
\midrule
Item A & 10.5 & 20.3 & 30.8 \\
Item B & 15.2 & 25.1 & 40.3 \\
Item C & 22.1 & 12.7 & 34.8 \\
\bottomrule
\end{tabular}
\caption{Sample Data Table}
\label{tab:sample}
\end{table}

\section{Lists}
\subsection{Bullet Points}
\begin{itemize}
  \item First item
  \item Second item
  \item Third item with sub-items:
    \begin{itemize}
      \item Sub-item 1
      \item Sub-item 2
    \end{itemize}
\end{itemize}

\subsection{Numbered List}
\begin{enumerate}
  \item First step
  \item Second step
  \item Third step
\end{enumerate}

\section{Citations and References}
When you use the citation search feature, references will be automatically managed. For example, you might cite a paper like this: \cite{example2024}. All citations will be automatically added to your \texttt{cite.bib} file.

\section{Conclusion}
This template demonstrates basic \LaTeX{} features including mathematical formulas, tables, and lists. The bibliography system is now set up to work seamlessly with Bibby AI's citation management. Feel free to modify it for your needs!

\bibliographystyle{plain}
\bibliography{cite}

\end{document}\documentclass{article}
\usepackage[utf8]{inputenc}
\usepackage{booktabs}
\usepackage{amsmath}
\usepackage{natbib}

\title{My Document}
\author{Your Name}
\date{\today}

\begin{document}

\maketitle

\section{Introduction}
This is a sample document demonstrating various \LaTeX{} features including mathematical equations, tables, and formatting. You can cite papers using the citation search feature, and they will be automatically added to your \texttt{cite.bib} file.

\section{Mathematical Expressions}
Here are some example equations:

\subsection{Quadratic Formula}
The quadratic formula is given by:
\[x = \frac{-b \pm \sqrt{b^2 - 4ac}}{2a}\]

\subsection{Matrix Example}
A sample matrix:
\[
\begin{pmatrix}
a_{11} & a_{12} & a_{13} \\
a_{21} & a_{22} & a_{23} \\
a_{31} & a_{32} & a_{33}
\end{pmatrix}
\]

\section{Tables}
Below is a sample table showing data:

\begin{table}[h]
\centering
\begin{tabular}{@{}lrrr@{}}
\toprule
\textbf{Category} & \textbf{Value 1} & \textbf{Value 2} & \textbf{Total} \\
\midrule
Item A & 10.5 & 20.3 & 30.8 \\
Item B & 15.2 & 25.1 & 40.3 \\
Item C & 22.1 & 12.7 & 34.8 \\
\bottomrule
\end{tabular}
\caption{Sample Data Table}
\label{tab:sample}
\end{table}

\section{Lists}
\subsection{Bullet Points}
\begin{itemize}
  \item First item
  \item Second item
  \item Third item with sub-items:
    \begin{itemize}
      \item Sub-item 1
      \item Sub-item 2
    \end{itemize}
\end{itemize}

\subsection{Numbered List}
\begin{enumerate}
  \item First step
  \item Second step
  \item Third step
\end{enumerate}

\section{Citations and References}
When you use the citation search feature, references will be automatically managed. For example, you might cite a paper like this: \cite{example2024}. All citations will be automatically added to your \texttt{cite.bib} file.

\section{Conclusion}
This template demonstrates basic \LaTeX{} features including mathematical formulas, tables, and lists. The bibliography system is now set up to work seamlessly with Bibby AI's citation management. Feel free to modify it for your needs!

\bibliographystyle{plain}
\bibliography{cite}

\end{document}\documentclass{article}
\usepackage[utf8]{inputenc}
\usepackage{booktabs}
\usepackage{amsmath}
\usepackage{natbib}

\title{My Document}
\author{Your Name}
\date{\today}

\begin{document}

\maketitle

\section{Introduction}
This is a sample document demonstrating various \LaTeX{} features including mathematical equations, tables, and formatting. You can cite papers using the citation search feature, and they will be automatically added to your \texttt{cite.bib} file.

\section{Mathematical Expressions}
Here are some example equations:

\subsection{Quadratic Formula}
The quadratic formula is given by:
\[x = \frac{-b \pm \sqrt{b^2 - 4ac}}{2a}\]

\subsection{Matrix Example}
A sample matrix:
\[
\begin{pmatrix}
a_{11} & a_{12} & a_{13} \\
a_{21} & a_{22} & a_{23} \\
a_{31} & a_{32} & a_{33}
\end{pmatrix}
\]

\section{Tables}
Below is a sample table showing data:

\begin{table}[h]
\centering
\begin{tabular}{@{}lrrr@{}}
\toprule
\textbf{Category} & \textbf{Value 1} & \textbf{Value 2} & \textbf{Total} \\
\midrule
Item A & 10.5 & 20.3 & 30.8 \\
Item B & 15.2 & 25.1 & 40.3 \\
Item C & 22.1 & 12.7 & 34.8 \\
\bottomrule
\end{tabular}
\caption{Sample Data Table}
\label{tab:sample}
\end{table}

\section{Lists}
\subsection{Bullet Points}
\begin{itemize}
  \item First item
  \item Second item
  \item Third item with sub-items:
    \begin{itemize}
      \item Sub-item 1
      \item Sub-item 2
    \end{itemize}
\end{itemize}

\subsection{Numbered List}
\begin{enumerate}
  \item First step
  \item Second step
  \item Third step
\end{enumerate}

\section{Citations and References}
When you use the citation search feature, references will be automatically managed. For example, you might cite a paper like this: \cite{example2024}. All citations will be automatically added to your \texttt{cite.bib} file.

\section{Conclusion}
This template demonstrates basic \LaTeX{} features including mathematical formulas, tables, and lists. The bibliography system is now set up to work seamlessly with Bibby AI's citation management. Feel free to modify it for your needs!

\bibliographystyle{plain}
\bibliography{cite}

\end{document}\documentclass{article}
\usepackage[utf8]{inputenc}
\usepackage{booktabs}
\usepackage{amsmath}
\usepackage{natbib}

\title{My Document}
\author{Your Name}
\date{\today}

\begin{document}

\maketitle

\section{Introduction}
This is a sample document demonstrating various \LaTeX{} features including mathematical equations, tables, and formatting. You can cite papers using the citation search feature, and they will be automatically added to your \texttt{cite.bib} file.

\section{Mathematical Expressions}
Here are some example equations:

\subsection{Quadratic Formula}
The quadratic formula is given by:
\[x = \frac{-b \pm \sqrt{b^2 - 4ac}}{2a}\]

\subsection{Matrix Example}
A sample matrix:
\[
\begin{pmatrix}
a_{11} & a_{12} & a_{13} \\
a_{21} & a_{22} & a_{23} \\
a_{31} & a_{32} & a_{33}
\end{pmatrix}
\]

\section{Tables}
Below is a sample table showing data:

\begin{table}[h]
\centering
\begin{tabular}{@{}lrrr@{}}
\toprule
\textbf{Category} & \textbf{Value 1} & \textbf{Value 2} & \textbf{Total} \\
\midrule
Item A & 10.5 & 20.3 & 30.8 \\
Item B & 15.2 & 25.1 & 40.3 \\
Item C & 22.1 & 12.7 & 34.8 \\
\bottomrule
\end{tabular}
\caption{Sample Data Table}
\label{tab:sample}
\end{table}

\section{Lists}
\subsection{Bullet Points}
\begin{itemize}
  \item First item
  \item Second item
  \item Third item with sub-items:
    \begin{itemize}
      \item Sub-item 1
      \item Sub-item 2
    \end{itemize}
\end{itemize}

\subsection{Numbered List}
\begin{enumerate}
  \item First step
  \item Second step
  \item Third step
\end{enumerate}

\section{Citations and References}
When you use the citation search feature, references will be automatically managed. For example, you might cite a paper like this: \cite{example2024}. All citations will be automatically added to your \texttt{cite.bib} file.

\section{Conclusion}
This template demonstrates basic \LaTeX{} features including mathematical formulas, tables, and lists. The bibliography system is now set up to work seamlessly with Bibby AI's citation management. Feel free to modify it for your needs!

\bibliographystyle{plain}
\bibliography{cite}

\end{document}\documentclass{article}
\usepackage[utf8]{inputenc}
\usepackage{booktabs}
\usepackage{amsmath}
\usepackage{natbib}

\title{My Document}
\author{Your Name}
\date{\today}

\begin{document}

\maketitle

\section{Introduction}
This is a sample document demonstrating various \LaTeX{} features including mathematical equations, tables, and formatting. You can cite papers using the citation search feature, and they will be automatically added to your \texttt{cite.bib} file.

\section{Mathematical Expressions}
Here are some example equations:

\subsection{Quadratic Formula}
The quadratic formula is given by:
\[x = \frac{-b \pm \sqrt{b^2 - 4ac}}{2a}\]

\subsection{Matrix Example}
A sample matrix:
\[
\begin{pmatrix}
a_{11} & a_{12} & a_{13} \\
a_{21} & a_{22} & a_{23} \\
a_{31} & a_{32} & a_{33}
\end{pmatrix}
\]

\section{Tables}
Below is a sample table showing data:

\begin{table}[h]
\centering
\begin{tabular}{@{}lrrr@{}}
\toprule
\textbf{Category} & \textbf{Value 1} & \textbf{Value 2} & \textbf{Total} \\
\midrule
Item A & 10.5 & 20.3 & 30.8 \\
Item B & 15.2 & 25.1 & 40.3 \\
Item C & 22.1 & 12.7 & 34.8 \\
\bottomrule
\end{tabular}
\caption{Sample Data Table}
\label{tab:sample}
\end{table}

\section{Lists}
\subsection{Bullet Points}
\begin{itemize}
  \item First item
  \item Second item
  \item Third item with sub-items:
    \begin{itemize}
      \item Sub-item 1
      \item Sub-item 2
    \end{itemize}
\end{itemize}

\subsection{Numbered List}
\begin{enumerate}
  \item First step
  \item Second step
  \item Third step
\end{enumerate}

\section{Citations and References}
When you use the citation search feature, references will be automatically managed. For example, you might cite a paper like this: \cite{example2024}. All citations will be automatically added to your \texttt{cite.bib} file.

\section{Conclusion}
This template demonstrates basic \LaTeX{} features including mathematical formulas, tables, and lists. The bibliography system is now set up to work seamlessly with Bibby AI's citation management. Feel free to modify it for your needs!

\bibliographystyle{plain}
\bibliography{cite}

\end{document}\documentclass{article}
\usepackage[utf8]{inputenc}
\usepackage{booktabs}
\usepackage{amsmath}
\usepackage{natbib}

\title{My Document}
\author{Your Name}
\date{\today}

\begin{document}

\maketitle

\section{Introduction}
This is a sample document demonstrating various \LaTeX{} features including mathematical equations, tables, and formatting. You can cite papers using the citation search feature, and they will be automatically added to your \texttt{cite.bib} file.

\section{Mathematical Expressions}
Here are some example equations:

\subsection{Quadratic Formula}
The quadratic formula is given by:
\[x = \frac{-b \pm \sqrt{b^2 - 4ac}}{2a}\]

\subsection{Matrix Example}
A sample matrix:
\[
\begin{pmatrix}
a_{11} & a_{12} & a_{13} \\
a_{21} & a_{22} & a_{23} \\
a_{31} & a_{32} & a_{33}
\end{pmatrix}
\]

\section{Tables}
Below is a sample table showing data:

\begin{table}[h]
\centering
\begin{tabular}{@{}lrrr@{}}
\toprule
\textbf{Category} & \textbf{Value 1} & \textbf{Value 2} & \textbf{Total} \\
\midrule
Item A & 10.5 & 20.3 & 30.8 \\
Item B & 15.2 & 25.1 & 40.3 \\
Item C & 22.1 & 12.7 & 34.8 \\
\bottomrule
\end{tabular}
\caption{Sample Data Table}
\label{tab:sample}
\end{table}

\section{Lists}
\subsection{Bullet Points}
\begin{itemize}
  \item First item
  \item Second item
  \item Third item with sub-items:
    \begin{itemize}
      \item Sub-item 1
      \item Sub-item 2
    \end{itemize}
\end{itemize}

\subsection{Numbered List}
\begin{enumerate}
  \item First step
  \item Second step
  \item Third step
\end{enumerate}

\section{Citations and References}
When you use the citation search feature, references will be automatically managed. For example, you might cite a paper like this: \cite{example2024}. All citations will be automatically added to your \texttt{cite.bib} file.

\section{Conclusion}
This template demonstrates basic \LaTeX{} features including mathematical formulas, tables, and lists. The bibliography system is now set up to work seamlessly with Bibby AI's citation management. Feel free to modify it for your needs!

\bibliographystyle{plain}
\bibliography{cite}

\end{document}\documentclass{article}
\usepackage[utf8]{inputenc}
\usepackage{booktabs}
\usepackage{amsmath}
\usepackage{natbib}

\title{My Document}
\author{Your Name}
\date{\today}

\begin{document}

\maketitle

\section{Introduction}
This is a sample document demonstrating various \LaTeX{} features including mathematical equations, tables, and formatting. You can cite papers using the citation search feature, and they will be automatically added to your \texttt{cite.bib} file.

\section{Mathematical Expressions}
Here are some example equations:

\subsection{Quadratic Formula}
The quadratic formula is given by:
\[x = \frac{-b \pm \sqrt{b^2 - 4ac}}{2a}\]

\subsection{Matrix Example}
A sample matrix:
\[
\begin{pmatrix}
a_{11} & a_{12} & a_{13} \\
a_{21} & a_{22} & a_{23} \\
a_{31} & a_{32} & a_{33}
\end{pmatrix}
\]

\section{Tables}
Below is a sample table showing data:

\begin{table}[h]
\centering
\begin{tabular}{@{}lrrr@{}}
\toprule
\textbf{Category} & \textbf{Value 1} & \textbf{Value 2} & \textbf{Total} \\
\midrule
Item A & 10.5 & 20.3 & 30.8 \\
Item B & 15.2 & 25.1 & 40.3 \\
Item C & 22.1 & 12.7 & 34.8 \\
\bottomrule
\end{tabular}
\caption{Sample Data Table}
\label{tab:sample}
\end{table}

\section{Lists}
\subsection{Bullet Points}
\begin{itemize}
  \item First item
  \item Second item
  \item Third item with sub-items:
    \begin{itemize}
      \item Sub-item 1
      \item Sub-item 2
    \end{itemize}
\end{itemize}

\subsection{Numbered List}
\begin{enumerate}
  \item First step
  \item Second step
  \item Third step
\end{enumerate}

\section{Citations and References}
When you use the citation search feature, references will be automatically managed. For example, you might cite a paper like this: \cite{example2024}. All citations will be automatically added to your \texttt{cite.bib} file.

\section{Conclusion}
This template demonstrates basic \LaTeX{} features including mathematical formulas, tables, and lists. The bibliography system is now set up to work seamlessly with Bibby AI's citation management. Feel free to modify it for your needs!

\bibliographystyle{plain}
\bibliography{cite}

\end{document}\documentclass{article}
\usepackage[utf8]{inputenc}
\usepackage{booktabs}
\usepackage{amsmath}
\usepackage{natbib}

\title{My Document}
\author{Your Name}
\date{\today}

\begin{document}

\maketitle

\section{Introduction}
This is a sample document demonstrating various \LaTeX{} features including mathematical equations, tables, and formatting. You can cite papers using the citation search feature, and they will be automatically added to your \texttt{cite.bib} file.

\section{Mathematical Expressions}
Here are some example equations:

\subsection{Quadratic Formula}
The quadratic formula is given by:
\[x = \frac{-b \pm \sqrt{b^2 - 4ac}}{2a}\]

\subsection{Matrix Example}
A sample matrix:
\[
\begin{pmatrix}
a_{11} & a_{12} & a_{13} \\
a_{21} & a_{22} & a_{23} \\
a_{31} & a_{32} & a_{33}
\end{pmatrix}
\]

\section{Tables}
Below is a sample table showing data:

\begin{table}[h]
\centering
\begin{tabular}{@{}lrrr@{}}
\toprule
\textbf{Category} & \textbf{Value 1} & \textbf{Value 2} & \textbf{Total} \\
\midrule
Item A & 10.5 & 20.3 & 30.8 \\
Item B & 15.2 & 25.1 & 40.3 \\
Item C & 22.1 & 12.7 & 34.8 \\
\bottomrule
\end{tabular}
\caption{Sample Data Table}
\label{tab:sample}
\end{table}

\section{Lists}
\subsection{Bullet Points}
\begin{itemize}
  \item First item
  \item Second item
  \item Third item with sub-items:
    \begin{itemize}
      \item Sub-item 1
      \item Sub-item 2
    \end{itemize}
\end{itemize}

\subsection{Numbered List}
\begin{enumerate}
  \item First step
  \item Second step
  \item Third step
\end{enumerate}

\section{Citations and References}
When you use the citation search feature, references will be automatically managed. For example, you might cite a paper like this: \cite{example2024}. All citations will be automatically added to your \texttt{cite.bib} file.

\section{Conclusion}
This template demonstrates basic \LaTeX{} features including mathematical formulas, tables, and lists. The bibliography system is now set up to work seamlessly with Bibby AI's citation management. Feel free to modify it for your needs!

\bibliographystyle{plain}
\bibliography{cite}

\end{document}\documentclass{article}
\usepackage[utf8]{inputenc}
\usepackage{booktabs}
\usepackage{amsmath}
\usepackage{natbib}

\title{My Document}
\author{Your Name}
\date{\today}

\begin{document}

\maketitle

\section{Introduction}
This is a sample document demonstrating various \LaTeX{} features including mathematical equations, tables, and formatting. You can cite papers using the citation search feature, and they will be automatically added to your \texttt{cite.bib} file.

\section{Mathematical Expressions}
Here are some example equations:

\subsection{Quadratic Formula}
The quadratic formula is given by:
\[x = \frac{-b \pm \sqrt{b^2 - 4ac}}{2a}\]

\subsection{Matrix Example}
A sample matrix:
\[
\begin{pmatrix}
a_{11} & a_{12} & a_{13} \\
a_{21} & a_{22} & a_{23} \\
a_{31} & a_{32} & a_{33}
\end{pmatrix}
\]

\section{Tables}
Below is a sample table showing data:

\begin{table}[h]
\centering
\begin{tabular}{@{}lrrr@{}}
\toprule
\textbf{Category} & \textbf{Value 1} & \textbf{Value 2} & \textbf{Total} \\
\midrule
Item A & 10.5 & 20.3 & 30.8 \\
Item B & 15.2 & 25.1 & 40.3 \\
Item C & 22.1 & 12.7 & 34.8 \\
\bottomrule
\end{tabular}
\caption{Sample Data Table}
\label{tab:sample}
\end{table}

\section{Lists}
\subsection{Bullet Points}
\begin{itemize}
  \item First item
  \item Second item
  \item Third item with sub-items:
    \begin{itemize}
      \item Sub-item 1
      \item Sub-item 2
    \end{itemize}
\end{itemize}

\subsection{Numbered List}
\begin{enumerate}
  \item First step
  \item Second step
  \item Third step
\end{enumerate}

\section{Citations and References}
When you use the citation search feature, references will be automatically managed. For example, you might cite a paper like this: \cite{example2024}. All citations will be automatically added to your \texttt{cite.bib} file.

\section{Conclusion}
This template demonstrates basic \LaTeX{} features including mathematical formulas, tables, and lists. The bibliography system is now set up to work seamlessly with Bibby AI's citation management. Feel free to modify it for your needs!

\bibliographystyle{plain}
\bibliography{cite}

\end{document}\documentclass{article}
\usepackage[utf8]{inputenc}
\usepackage{booktabs}
\usepackage{amsmath}
\usepackage{natbib}

\title{My Document}
\author{Your Name}
\date{\today}

\begin{document}

\maketitle

\section{Introduction}
This is a sample document demonstrating various \LaTeX{} features including mathematical equations, tables, and formatting. You can cite papers using the citation search feature, and they will be automatically added to your \texttt{cite.bib} file.

\section{Mathematical Expressions}
Here are some example equations:

\subsection{Quadratic Formula}
The quadratic formula is given by:
\[x = \frac{-b \pm \sqrt{b^2 - 4ac}}{2a}\]

\subsection{Matrix Example}
A sample matrix:
\[
\begin{pmatrix}
a_{11} & a_{12} & a_{13} \\
a_{21} & a_{22} & a_{23} \\
a_{31} & a_{32} & a_{33}
\end{pmatrix}
\]

\section{Tables}
Below is a sample table showing data:

\begin{table}[h]
\centering
\begin{tabular}{@{}lrrr@{}}
\toprule
\textbf{Category} & \textbf{Value 1} & \textbf{Value 2} & \textbf{Total} \\
\midrule
Item A & 10.5 & 20.3 & 30.8 \\
Item B & 15.2 & 25.1 & 40.3 \\
Item C & 22.1 & 12.7 & 34.8 \\
\bottomrule
\end{tabular}
\caption{Sample Data Table}
\label{tab:sample}
\end{table}

\section{Lists}
\subsection{Bullet Points}
\begin{itemize}
  \item First item
  \item Second item
  \item Third item with sub-items:
    \begin{itemize}
      \item Sub-item 1
      \item Sub-item 2
    \end{itemize}
\end{itemize}

\subsection{Numbered List}
\begin{enumerate}
  \item First step
  \item Second step
  \item Third step
\end{enumerate}

\section{Citations and References}
When you use the citation search feature, references will be automatically managed. For example, you might cite a paper like this: \cite{example2024}. All citations will be automatically added to your \texttt{cite.bib} file.

\section{Conclusion}
This template demonstrates basic \LaTeX{} features including mathematical formulas, tables, and lists. The bibliography system is now set up to work seamlessly with Bibby AI's citation management. Feel free to modify it for your needs!

\bibliographystyle{plain}
\bibliography{cite}

\end{document}\documentclass{article}
\usepackage[utf8]{inputenc}
\usepackage{booktabs}
\usepackage{amsmath}
\usepackage{natbib}

\title{My Document}
\author{Your Name}
\date{\today}

\begin{document}

\maketitle

\section{Introduction}
This is a sample document demonstrating various \LaTeX{} features including mathematical equations, tables, and formatting. You can cite papers using the citation search feature, and they will be automatically added to your \texttt{cite.bib} file.

\section{Mathematical Expressions}
Here are some example equations:

\subsection{Quadratic Formula}
The quadratic formula is given by:
\[x = \frac{-b \pm \sqrt{b^2 - 4ac}}{2a}\]

\subsection{Matrix Example}
A sample matrix:
\[
\begin{pmatrix}
a_{11} & a_{12} & a_{13} \\
a_{21} & a_{22} & a_{23} \\
a_{31} & a_{32} & a_{33}
\end{pmatrix}
\]

\section{Tables}
Below is a sample table showing data:

\begin{table}[h]
\centering
\begin{tabular}{@{}lrrr@{}}
\toprule
\textbf{Category} & \textbf{Value 1} & \textbf{Value 2} & \textbf{Total} \\
\midrule
Item A & 10.5 & 20.3 & 30.8 \\
Item B & 15.2 & 25.1 & 40.3 \\
Item C & 22.1 & 12.7 & 34.8 \\
\bottomrule
\end{tabular}
\caption{Sample Data Table}
\label{tab:sample}
\end{table}

\section{Lists}
\subsection{Bullet Points}
\begin{itemize}
  \item First item
  \item Second item
  \item Third item with sub-items:
    \begin{itemize}
      \item Sub-item 1
      \item Sub-item 2
    \end{itemize}
\end{itemize}

\subsection{Numbered List}
\begin{enumerate}
  \item First step
  \item Second step
  \item Third step
\end{enumerate}

\section{Citations and References}
When you use the citation search feature, references will be automatically managed. For example, you might cite a paper like this: \cite{example2024}. All citations will be automatically added to your \texttt{cite.bib} file.

\section{Conclusion}
This template demonstrates basic \LaTeX{} features including mathematical formulas, tables, and lists. The bibliography system is now set up to work seamlessly with Bibby AI's citation management. Feel free to modify it for your needs!

\bibliographystyle{plain}
\bibliography{cite}

\end{document}\documentclass{article}
\usepackage[utf8]{inputenc}
\usepackage{booktabs}
\usepackage{amsmath}
\usepackage{natbib}

\title{My Document}
\author{Your Name}
\date{\today}

\begin{document}

\maketitle

\section{Introduction}
This is a sample document demonstrating various \LaTeX{} features including mathematical equations, tables, and formatting. You can cite papers using the citation search feature, and they will be automatically added to your \texttt{cite.bib} file.

\section{Mathematical Expressions}
Here are some example equations:

\subsection{Quadratic Formula}
The quadratic formula is given by:
\[x = \frac{-b \pm \sqrt{b^2 - 4ac}}{2a}\]

\subsection{Matrix Example}
A sample matrix:
\[
\begin{pmatrix}
a_{11} & a_{12} & a_{13} \\
a_{21} & a_{22} & a_{23} \\
a_{31} & a_{32} & a_{33}
\end{pmatrix}
\]

\section{Tables}
Below is a sample table showing data:

\begin{table}[h]
\centering
\begin{tabular}{@{}lrrr@{}}
\toprule
\textbf{Category} & \textbf{Value 1} & \textbf{Value 2} & \textbf{Total} \\
\midrule
Item A & 10.5 & 20.3 & 30.8 \\
Item B & 15.2 & 25.1 & 40.3 \\
Item C & 22.1 & 12.7 & 34.8 \\
\bottomrule
\end{tabular}
\caption{Sample Data Table}
\label{tab:sample}
\end{table}

\section{Lists}
\subsection{Bullet Points}
\begin{itemize}
  \item First item
  \item Second item
  \item Third item with sub-items:
    \begin{itemize}
      \item Sub-item 1
      \item Sub-item 2
    \end{itemize}
\end{itemize}

\subsection{Numbered List}
\begin{enumerate}
  \item First step
  \item Second step
  \item Third step
\end{enumerate}

\section{Citations and References}
When you use the citation search feature, references will be automatically managed. For example, you might cite a paper like this: \cite{example2024}. All citations will be automatically added to your \texttt{cite.bib} file.

\section{Conclusion}
This template demonstrates basic \LaTeX{} features including mathematical formulas, tables, and lists. The bibliography system is now set up to work seamlessly with Bibby AI's citation management. Feel free to modify it for your needs!

\bibliographystyle{plain}
\bibliography{cite}

\end{document}\documentclass{article}
\usepackage[utf8]{inputenc}
\usepackage{booktabs}
\usepackage{amsmath}
\usepackage{natbib}

\title{My Document}
\author{Your Name}
\date{\today}

\begin{document}

\maketitle

\section{Introduction}
This is a sample document demonstrating various \LaTeX{} features including mathematical equations, tables, and formatting. You can cite papers using the citation search feature, and they will be automatically added to your \texttt{cite.bib} file.

\section{Mathematical Expressions}
Here are some example equations:

\subsection{Quadratic Formula}
The quadratic formula is given by:
\[x = \frac{-b \pm \sqrt{b^2 - 4ac}}{2a}\]

\subsection{Matrix Example}
A sample matrix:
\[
\begin{pmatrix}
a_{11} & a_{12} & a_{13} \\
a_{21} & a_{22} & a_{23} \\
a_{31} & a_{32} & a_{33}
\end{pmatrix}
\]

\section{Tables}
Below is a sample table showing data:

\begin{table}[h]
\centering
\begin{tabular}{@{}lrrr@{}}
\toprule
\textbf{Category} & \textbf{Value 1} & \textbf{Value 2} & \textbf{Total} \\
\midrule
Item A & 10.5 & 20.3 & 30.8 \\
Item B & 15.2 & 25.1 & 40.3 \\
Item C & 22.1 & 12.7 & 34.8 \\
\bottomrule
\end{tabular}
\caption{Sample Data Table}
\label{tab:sample}
\end{table}

\section{Lists}
\subsection{Bullet Points}
\begin{itemize}
  \item First item
  \item Second item
  \item Third item with sub-items:
    \begin{itemize}
      \item Sub-item 1
      \item Sub-item 2
    \end{itemize}
\end{itemize}

\subsection{Numbered List}
\begin{enumerate}
  \item First step
  \item Second step
  \item Third step
\end{enumerate}

\section{Citations and References}
When you use the citation search feature, references will be automatically managed. For example, you might cite a paper like this: \cite{example2024}. All citations will be automatically added to your \texttt{cite.bib} file.

\section{Conclusion}
This template demonstrates basic \LaTeX{} features including mathematical formulas, tables, and lists. The bibliography system is now set up to work seamlessly with Bibby AI's citation management. Feel free to modify it for your needs!

\bibliographystyle{plain}
\bibliography{cite}

\end{document}\documentclass{article}
\usepackage[utf8]{inputenc}
\usepackage{booktabs}
\usepackage{amsmath}
\usepackage{natbib}

\title{My Document}
\author{Your Name}
\date{\today}

\begin{document}

\maketitle

\section{Introduction}
This is a sample document demonstrating various \LaTeX{} features including mathematical equations, tables, and formatting. You can cite papers using the citation search feature, and they will be automatically added to your \texttt{cite.bib} file.

\section{Mathematical Expressions}
Here are some example equations:

\subsection{Quadratic Formula}
The quadratic formula is given by:
\[x = \frac{-b \pm \sqrt{b^2 - 4ac}}{2a}\]

\subsection{Matrix Example}
A sample matrix:
\[
\begin{pmatrix}
a_{11} & a_{12} & a_{13} \\
a_{21} & a_{22} & a_{23} \\
a_{31} & a_{32} & a_{33}
\end{pmatrix}
\]

\section{Tables}
Below is a sample table showing data:

\begin{table}[h]
\centering
\begin{tabular}{@{}lrrr@{}}
\toprule
\textbf{Category} & \textbf{Value 1} & \textbf{Value 2} & \textbf{Total} \\
\midrule
Item A & 10.5 & 20.3 & 30.8 \\
Item B & 15.2 & 25.1 & 40.3 \\
Item C & 22.1 & 12.7 & 34.8 \\
\bottomrule
\end{tabular}
\caption{Sample Data Table}
\label{tab:sample}
\end{table}

\section{Lists}
\subsection{Bullet Points}
\begin{itemize}
  \item First item
  \item Second item
  \item Third item with sub-items:
    \begin{itemize}
      \item Sub-item 1
      \item Sub-item 2
    \end{itemize}
\end{itemize}

\subsection{Numbered List}
\begin{enumerate}
  \item First step
  \item Second step
  \item Third step
\end{enumerate}

\section{Citations and References}
When you use the citation search feature, references will be automatically managed. For example, you might cite a paper like this: \cite{example2024}. All citations will be automatically added to your \texttt{cite.bib} file.

\section{Conclusion}
This template demonstrates basic \LaTeX{} features including mathematical formulas, tables, and lists. The bibliography system is now set up to work seamlessly with Bibby AI's citation management. Feel free to modify it for your needs!

\bibliographystyle{plain}
\bibliography{cite}

\end{document}\documentclass{article}
\usepackage[utf8]{inputenc}
\usepackage{booktabs}
\usepackage{amsmath}
\usepackage{natbib}

\title{My Document}
\author{Your Name}
\date{\today}

\begin{document}

\maketitle

\section{Introduction}
This is a sample document demonstrating various \LaTeX{} features including mathematical equations, tables, and formatting. You can cite papers using the citation search feature, and they will be automatically added to your \texttt{cite.bib} file.

\section{Mathematical Expressions}
Here are some example equations:

\subsection{Quadratic Formula}
The quadratic formula is given by:
\[x = \frac{-b \pm \sqrt{b^2 - 4ac}}{2a}\]

\subsection{Matrix Example}
A sample matrix:
\[
\begin{pmatrix}
a_{11} & a_{12} & a_{13} \\
a_{21} & a_{22} & a_{23} \\
a_{31} & a_{32} & a_{33}
\end{pmatrix}
\]

\section{Tables}
Below is a sample table showing data:

\begin{table}[h]
\centering
\begin{tabular}{@{}lrrr@{}}
\toprule
\textbf{Category} & \textbf{Value 1} & \textbf{Value 2} & \textbf{Total} \\
\midrule
Item A & 10.5 & 20.3 & 30.8 \\
Item B & 15.2 & 25.1 & 40.3 \\
Item C & 22.1 & 12.7 & 34.8 \\
\bottomrule
\end{tabular}
\caption{Sample Data Table}
\label{tab:sample}
\end{table}

\section{Lists}
\subsection{Bullet Points}
\begin{itemize}
  \item First item
  \item Second item
  \item Third item with sub-items:
    \begin{itemize}
      \item Sub-item 1
      \item Sub-item 2
    \end{itemize}
\end{itemize}

\subsection{Numbered List}
\begin{enumerate}
  \item First step
  \item Second step
  \item Third step
\end{enumerate}

\section{Citations and References}
When you use the citation search feature, references will be automatically managed. For example, you might cite a paper like this: \cite{example2024}. All citations will be automatically added to your \texttt{cite.bib} file.

\section{Conclusion}
This template demonstrates basic \LaTeX{} features including mathematical formulas, tables, and lists. The bibliography system is now set up to work seamlessly with Bibby AI's citation management. Feel free to modify it for your needs!

\bibliographystyle{plain}
\bibliography{cite}

\end{document}\documentclass{article}
\usepackage[utf8]{inputenc}
\usepackage{booktabs}
\usepackage{amsmath}
\usepackage{natbib}

\title{My Document}
\author{Your Name}
\date{\today}

\begin{document}

\maketitle

\section{Introduction}
This is a sample document demonstrating various \LaTeX{} features including mathematical equations, tables, and formatting. You can cite papers using the citation search feature, and they will be automatically added to your \texttt{cite.bib} file.

\section{Mathematical Expressions}
Here are some example equations:

\subsection{Quadratic Formula}
The quadratic formula is given by:
\[x = \frac{-b \pm \sqrt{b^2 - 4ac}}{2a}\]

\subsection{Matrix Example}
A sample matrix:
\[
\begin{pmatrix}
a_{11} & a_{12} & a_{13} \\
a_{21} & a_{22} & a_{23} \\
a_{31} & a_{32} & a_{33}
\end{pmatrix}
\]

\section{Tables}
Below is a sample table showing data:

\begin{table}[h]
\centering
\begin{tabular}{@{}lrrr@{}}
\toprule
\textbf{Category} & \textbf{Value 1} & \textbf{Value 2} & \textbf{Total} \\
\midrule
Item A & 10.5 & 20.3 & 30.8 \\
Item B & 15.2 & 25.1 & 40.3 \\
Item C & 22.1 & 12.7 & 34.8 \\
\bottomrule
\end{tabular}
\caption{Sample Data Table}
\label{tab:sample}
\end{table}

\section{Lists}
\subsection{Bullet Points}
\begin{itemize}
  \item First item
  \item Second item
  \item Third item with sub-items:
    \begin{itemize}
      \item Sub-item 1
      \item Sub-item 2
    \end{itemize}
\end{itemize}

\subsection{Numbered List}
\begin{enumerate}
  \item First step
  \item Second step
  \item Third step
\end{enumerate}

\section{Citations and References}
When you use the citation search feature, references will be automatically managed. For example, you might cite a paper like this: \cite{example2024}. All citations will be automatically added to your \texttt{cite.bib} file.

\section{Conclusion}
This template demonstrates basic \LaTeX{} features including mathematical formulas, tables, and lists. The bibliography system is now set up to work seamlessly with Bibby AI's citation management. Feel free to modify it for your needs!

\bibliographystyle{plain}
\bibliography{cite}

\end{document}\documentclass{article}
\usepackage[utf8]{inputenc}
\usepackage{booktabs}
\usepackage{amsmath}
\usepackage{natbib}

\title{My Document}
\author{Your Name}
\date{\today}

\begin{document}

\maketitle

\section{Introduction}
This is a sample document demonstrating various \LaTeX{} features including mathematical equations, tables, and formatting. You can cite papers using the citation search feature, and they will be automatically added to your \texttt{cite.bib} file.

\section{Mathematical Expressions}
Here are some example equations:

\subsection{Quadratic Formula}
The quadratic formula is given by:
\[x = \frac{-b \pm \sqrt{b^2 - 4ac}}{2a}\]

\subsection{Matrix Example}
A sample matrix:
\[
\begin{pmatrix}
a_{11} & a_{12} & a_{13} \\
a_{21} & a_{22} & a_{23} \\
a_{31} & a_{32} & a_{33}
\end{pmatrix}
\]

\section{Tables}
Below is a sample table showing data:

\begin{table}[h]
\centering
\begin{tabular}{@{}lrrr@{}}
\toprule
\textbf{Category} & \textbf{Value 1} & \textbf{Value 2} & \textbf{Total} \\
\midrule
Item A & 10.5 & 20.3 & 30.8 \\
Item B & 15.2 & 25.1 & 40.3 \\
Item C & 22.1 & 12.7 & 34.8 \\
\bottomrule
\end{tabular}
\caption{Sample Data Table}
\label{tab:sample}
\end{table}

\section{Lists}
\subsection{Bullet Points}
\begin{itemize}
  \item First item
  \item Second item
  \item Third item with sub-items:
    \begin{itemize}
      \item Sub-item 1
      \item Sub-item 2
    \end{itemize}
\end{itemize}

\subsection{Numbered List}
\begin{enumerate}
  \item First step
  \item Second step
  \item Third step
\end{enumerate}

\section{Citations and References}
When you use the citation search feature, references will be automatically managed. For example, you might cite a paper like this: \cite{example2024}. All citations will be automatically added to your \texttt{cite.bib} file.

\section{Conclusion}
This template demonstrates basic \LaTeX{} features including mathematical formulas, tables, and lists. The bibliography system is now set up to work seamlessly with Bibby AI's citation management. Feel free to modify it for your needs!

\bibliographystyle{plain}
\bibliography{cite}

\end{document}\documentclass{article}
\usepackage[utf8]{inputenc}
\usepackage{booktabs}
\usepackage{amsmath}
\usepackage{natbib}

\title{My Document}
\author{Your Name}
\date{\today}

\begin{document}

\maketitle

\section{Introduction}
This is a sample document demonstrating various \LaTeX{} features including mathematical equations, tables, and formatting. You can cite papers using the citation search feature, and they will be automatically added to your \texttt{cite.bib} file.

\section{Mathematical Expressions}
Here are some example equations:

\subsection{Quadratic Formula}
The quadratic formula is given by:
\[x = \frac{-b \pm \sqrt{b^2 - 4ac}}{2a}\]

\subsection{Matrix Example}
A sample matrix:
\[
\begin{pmatrix}
a_{11} & a_{12} & a_{13} \\
a_{21} & a_{22} & a_{23} \\
a_{31} & a_{32} & a_{33}
\end{pmatrix}
\]

\section{Tables}
Below is a sample table showing data:

\begin{table}[h]
\centering
\begin{tabular}{@{}lrrr@{}}
\toprule
\textbf{Category} & \textbf{Value 1} & \textbf{Value 2} & \textbf{Total} \\
\midrule
Item A & 10.5 & 20.3 & 30.8 \\
Item B & 15.2 & 25.1 & 40.3 \\
Item C & 22.1 & 12.7 & 34.8 \\
\bottomrule
\end{tabular}
\caption{Sample Data Table}
\label{tab:sample}
\end{table}

\section{Lists}
\subsection{Bullet Points}
\begin{itemize}
  \item First item
  \item Second item
  \item Third item with sub-items:
    \begin{itemize}
      \item Sub-item 1
      \item Sub-item 2
    \end{itemize}
\end{itemize}

\subsection{Numbered List}
\begin{enumerate}
  \item First step
  \item Second step
  \item Third step
\end{enumerate}

\section{Citations and References}
When you use the citation search feature, references will be automatically managed. For example, you might cite a paper like this: \cite{example2024}. All citations will be automatically added to your \texttt{cite.bib} file.

\section{Conclusion}
This template demonstrates basic \LaTeX{} features including mathematical formulas, tables, and lists. The bibliography system is now set up to work seamlessly with Bibby AI's citation management. Feel free to modify it for your needs!

\bibliographystyle{plain}
\bibliography{cite}

\end{document}\documentclass{article}
\usepackage[utf8]{inputenc}
\usepackage{booktabs}
\usepackage{amsmath}
\usepackage{natbib}

\title{My Document}
\author{Your Name}
\date{\today}

\begin{document}

\maketitle

\section{Introduction}
This is a sample document demonstrating various \LaTeX{} features including mathematical equations, tables, and formatting. You can cite papers using the citation search feature, and they will be automatically added to your \texttt{cite.bib} file.

\section{Mathematical Expressions}
Here are some example equations:

\subsection{Quadratic Formula}
The quadratic formula is given by:
\[x = \frac{-b \pm \sqrt{b^2 - 4ac}}{2a}\]

\subsection{Matrix Example}
A sample matrix:
\[
\begin{pmatrix}
a_{11} & a_{12} & a_{13} \\
a_{21} & a_{22} & a_{23} \\
a_{31} & a_{32} & a_{33}
\end{pmatrix}
\]

\section{Tables}
Below is a sample table showing data:

\begin{table}[h]
\centering
\begin{tabular}{@{}lrrr@{}}
\toprule
\textbf{Category} & \textbf{Value 1} & \textbf{Value 2} & \textbf{Total} \\
\midrule
Item A & 10.5 & 20.3 & 30.8 \\
Item B & 15.2 & 25.1 & 40.3 \\
Item C & 22.1 & 12.7 & 34.8 \\
\bottomrule
\end{tabular}
\caption{Sample Data Table}
\label{tab:sample}
\end{table}

\section{Lists}
\subsection{Bullet Points}
\begin{itemize}
  \item First item
  \item Second item
  \item Third item with sub-items:
    \begin{itemize}
      \item Sub-item 1
      \item Sub-item 2
    \end{itemize}
\end{itemize}

\subsection{Numbered List}
\begin{enumerate}
  \item First step
  \item Second step
  \item Third step
\end{enumerate}

\section{Citations and References}
When you use the citation search feature, references will be automatically managed. For example, you might cite a paper like this: \cite{example2024}. All citations will be automatically added to your \texttt{cite.bib} file.

\section{Conclusion}
This template demonstrates basic \LaTeX{} features including mathematical formulas, tables, and lists. The bibliography system is now set up to work seamlessly with Bibby AI's citation management. Feel free to modify it for your needs!

\bibliographystyle{plain}
\bibliography{cite}

\end{document}\documentclass{article}
\usepackage[utf8]{inputenc}
\usepackage{booktabs}
\usepackage{amsmath}
\usepackage{natbib}

\title{My Document}
\author{Your Name}
\date{\today}

\begin{document}

\maketitle

\section{Introduction}
This is a sample document demonstrating various \LaTeX{} features including mathematical equations, tables, and formatting. You can cite papers using the citation search feature, and they will be automatically added to your \texttt{cite.bib} file.

\section{Mathematical Expressions}
Here are some example equations:

\subsection{Quadratic Formula}
The quadratic formula is given by:
\[x = \frac{-b \pm \sqrt{b^2 - 4ac}}{2a}\]

\subsection{Matrix Example}
A sample matrix:
\[
\begin{pmatrix}
a_{11} & a_{12} & a_{13} \\
a_{21} & a_{22} & a_{23} \\
a_{31} & a_{32} & a_{33}
\end{pmatrix}
\]

\section{Tables}
Below is a sample table showing data:

\begin{table}[h]
\centering
\begin{tabular}{@{}lrrr@{}}
\toprule
\textbf{Category} & \textbf{Value 1} & \textbf{Value 2} & \textbf{Total} \\
\midrule
Item A & 10.5 & 20.3 & 30.8 \\
Item B & 15.2 & 25.1 & 40.3 \\
Item C & 22.1 & 12.7 & 34.8 \\
\bottomrule
\end{tabular}
\caption{Sample Data Table}
\label{tab:sample}
\end{table}

\section{Lists}
\subsection{Bullet Points}
\begin{itemize}
  \item First item
  \item Second item
  \item Third item with sub-items:
    \begin{itemize}
      \item Sub-item 1
      \item Sub-item 2
    \end{itemize}
\end{itemize}

\subsection{Numbered List}
\begin{enumerate}
  \item First step
  \item Second step
  \item Third step
\end{enumerate}

\section{Citations and References}
When you use the citation search feature, references will be automatically managed. For example, you might cite a paper like this: \cite{example2024}. All citations will be automatically added to your \texttt{cite.bib} file.

\section{Conclusion}
This template demonstrates basic \LaTeX{} features including mathematical formulas, tables, and lists. The bibliography system is now set up to work seamlessly with Bibby AI's citation management. Feel free to modify it for your needs!

\bibliographystyle{plain}
\bibliography{cite}

\end{document}\documentclass{article}
\usepackage[utf8]{inputenc}
\usepackage{booktabs}
\usepackage{amsmath}
\usepackage{natbib}

\title{My Document}
\author{Your Name}
\date{\today}

\begin{document}

\maketitle

\section{Introduction}
This is a sample document demonstrating various \LaTeX{} features including mathematical equations, tables, and formatting. You can cite papers using the citation search feature, and they will be automatically added to your \texttt{cite.bib} file.

\section{Mathematical Expressions}
Here are some example equations:

\subsection{Quadratic Formula}
The quadratic formula is given by:
\[x = \frac{-b \pm \sqrt{b^2 - 4ac}}{2a}\]

\subsection{Matrix Example}
A sample matrix:
\[
\begin{pmatrix}
a_{11} & a_{12} & a_{13} \\
a_{21} & a_{22} & a_{23} \\
a_{31} & a_{32} & a_{33}
\end{pmatrix}
\]

\section{Tables}
Below is a sample table showing data:

\begin{table}[h]
\centering
\begin{tabular}{@{}lrrr@{}}
\toprule
\textbf{Category} & \textbf{Value 1} & \textbf{Value 2} & \textbf{Total} \\
\midrule
Item A & 10.5 & 20.3 & 30.8 \\
Item B & 15.2 & 25.1 & 40.3 \\
Item C & 22.1 & 12.7 & 34.8 \\
\bottomrule
\end{tabular}
\caption{Sample Data Table}
\label{tab:sample}
\end{table}

\section{Lists}
\subsection{Bullet Points}
\begin{itemize}
  \item First item
  \item Second item
  \item Third item with sub-items:
    \begin{itemize}
      \item Sub-item 1
      \item Sub-item 2
    \end{itemize}
\end{itemize}

\subsection{Numbered List}
\begin{enumerate}
  \item First step
  \item Second step
  \item Third step
\end{enumerate}

\section{Citations and References}
When you use the citation search feature, references will be automatically managed. For example, you might cite a paper like this: \cite{example2024}. All citations will be automatically added to your \texttt{cite.bib} file.

\section{Conclusion}
This template demonstrates basic \LaTeX{} features including mathematical formulas, tables, and lists. The bibliography system is now set up to work seamlessly with Bibby AI's citation management. Feel free to modify it for your needs!

\bibliographystyle{plain}
\bibliography{cite}

\end{document}\documentclass{article}
\usepackage[utf8]{inputenc}
\usepackage{booktabs}
\usepackage{amsmath}
\usepackage{natbib}

\title{My Document}
\author{Your Name}
\date{\today}

\begin{document}

\maketitle

\section{Introduction}
This is a sample document demonstrating various \LaTeX{} features including mathematical equations, tables, and formatting. You can cite papers using the citation search feature, and they will be automatically added to your \texttt{cite.bib} file.

\section{Mathematical Expressions}
Here are some example equations:

\subsection{Quadratic Formula}
The quadratic formula is given by:
\[x = \frac{-b \pm \sqrt{b^2 - 4ac}}{2a}\]

\subsection{Matrix Example}
A sample matrix:
\[
\begin{pmatrix}
a_{11} & a_{12} & a_{13} \\
a_{21} & a_{22} & a_{23} \\
a_{31} & a_{32} & a_{33}
\end{pmatrix}
\]

\section{Tables}
Below is a sample table showing data:

\begin{table}[h]
\centering
\begin{tabular}{@{}lrrr@{}}
\toprule
\textbf{Category} & \textbf{Value 1} & \textbf{Value 2} & \textbf{Total} \\
\midrule
Item A & 10.5 & 20.3 & 30.8 \\
Item B & 15.2 & 25.1 & 40.3 \\
Item C & 22.1 & 12.7 & 34.8 \\
\bottomrule
\end{tabular}
\caption{Sample Data Table}
\label{tab:sample}
\end{table}

\section{Lists}
\subsection{Bullet Points}
\begin{itemize}
  \item First item
  \item Second item
  \item Third item with sub-items:
    \begin{itemize}
      \item Sub-item 1
      \item Sub-item 2
    \end{itemize}
\end{itemize}

\subsection{Numbered List}
\begin{enumerate}
  \item First step
  \item Second step
  \item Third step
\end{enumerate}

\section{Citations and References}
When you use the citation search feature, references will be automatically managed. For example, you might cite a paper like this: \cite{example2024}. All citations will be automatically added to your \texttt{cite.bib} file.

\section{Conclusion}
This template demonstrates basic \LaTeX{} features including mathematical formulas, tables, and lists. The bibliography system is now set up to work seamlessly with Bibby AI's citation management. Feel free to modify it for your needs!

\bibliographystyle{plain}
\bibliography{cite}

\end{document}\documentclass{article}
\usepackage[utf8]{inputenc}
\usepackage{booktabs}
\usepackage{amsmath}
\usepackage{natbib}

\title{My Document}
\author{Your Name}
\date{\today}

\begin{document}

\maketitle

\section{Introduction}
This is a sample document demonstrating various \LaTeX{} features including mathematical equations, tables, and formatting. You can cite papers using the citation search feature, and they will be automatically added to your \texttt{cite.bib} file.

\section{Mathematical Expressions}
Here are some example equations:

\subsection{Quadratic Formula}
The quadratic formula is given by:
\[x = \frac{-b \pm \sqrt{b^2 - 4ac}}{2a}\]

\subsection{Matrix Example}
A sample matrix:
\[
\begin{pmatrix}
a_{11} & a_{12} & a_{13} \\
a_{21} & a_{22} & a_{23} \\
a_{31} & a_{32} & a_{33}
\end{pmatrix}
\]

\section{Tables}
Below is a sample table showing data:

\begin{table}[h]
\centering
\begin{tabular}{@{}lrrr@{}}
\toprule
\textbf{Category} & \textbf{Value 1} & \textbf{Value 2} & \textbf{Total} \\
\midrule
Item A & 10.5 & 20.3 & 30.8 \\
Item B & 15.2 & 25.1 & 40.3 \\
Item C & 22.1 & 12.7 & 34.8 \\
\bottomrule
\end{tabular}
\caption{Sample Data Table}
\label{tab:sample}
\end{table}

\section{Lists}
\subsection{Bullet Points}
\begin{itemize}
  \item First item
  \item Second item
  \item Third item with sub-items:
    \begin{itemize}
      \item Sub-item 1
      \item Sub-item 2
    \end{itemize}
\end{itemize}

\subsection{Numbered List}
\begin{enumerate}
  \item First step
  \item Second step
  \item Third step
\end{enumerate}

\section{Citations and References}
When you use the citation search feature, references will be automatically managed. For example, you might cite a paper like this: \cite{example2024}. All citations will be automatically added to your \texttt{cite.bib} file.

\section{Conclusion}
This template demonstrates basic \LaTeX{} features including mathematical formulas, tables, and lists. The bibliography system is now set up to work seamlessly with Bibby AI's citation management. Feel free to modify it for your needs!

\bibliographystyle{plain}
\bibliography{cite}

\end{document}\documentclass{article}
\usepackage[utf8]{inputenc}
\usepackage{booktabs}
\usepackage{amsmath}
\usepackage{natbib}

\title{My Document}
\author{Your Name}
\date{\today}

\begin{document}

\maketitle

\section{Introduction}
This is a sample document demonstrating various \LaTeX{} features including mathematical equations, tables, and formatting. You can cite papers using the citation search feature, and they will be automatically added to your \texttt{cite.bib} file.

\section{Mathematical Expressions}
Here are some example equations:

\subsection{Quadratic Formula}
The quadratic formula is given by:
\[x = \frac{-b \pm \sqrt{b^2 - 4ac}}{2a}\]

\subsection{Matrix Example}
A sample matrix:
\[
\begin{pmatrix}
a_{11} & a_{12} & a_{13} \\
a_{21} & a_{22} & a_{23} \\
a_{31} & a_{32} & a_{33}
\end{pmatrix}
\]

\section{Tables}
Below is a sample table showing data:

\begin{table}[h]
\centering
\begin{tabular}{@{}lrrr@{}}
\toprule
\textbf{Category} & \textbf{Value 1} & \textbf{Value 2} & \textbf{Total} \\
\midrule
Item A & 10.5 & 20.3 & 30.8 \\
Item B & 15.2 & 25.1 & 40.3 \\
Item C & 22.1 & 12.7 & 34.8 \\
\bottomrule
\end{tabular}
\caption{Sample Data Table}
\label{tab:sample}
\end{table}

\section{Lists}
\subsection{Bullet Points}
\begin{itemize}
  \item First item
  \item Second item
  \item Third item with sub-items:
    \begin{itemize}
      \item Sub-item 1
      \item Sub-item 2
    \end{itemize}
\end{itemize}

\subsection{Numbered List}
\begin{enumerate}
  \item First step
  \item Second step
  \item Third step
\end{enumerate}

\section{Citations and References}
When you use the citation search feature, references will be automatically managed. For example, you might cite a paper like this: \cite{example2024}. All citations will be automatically added to your \texttt{cite.bib} file.

\section{Conclusion}
This template demonstrates basic \LaTeX{} features including mathematical formulas, tables, and lists. The bibliography system is now set up to work seamlessly with Bibby AI's citation management. Feel free to modify it for your needs!

\bibliographystyle{plain}
\bibliography{cite}

\end{document}\documentclass{article}
\usepackage[utf8]{inputenc}
\usepackage{booktabs}
\usepackage{amsmath}
\usepackage{natbib}

\title{My Document}
\author{Your Name}
\date{\today}

\begin{document}

\maketitle

\section{Introduction}
This is a sample document demonstrating various \LaTeX{} features including mathematical equations, tables, and formatting. You can cite papers using the citation search feature, and they will be automatically added to your \texttt{cite.bib} file.

\section{Mathematical Expressions}
Here are some example equations:

\subsection{Quadratic Formula}
The quadratic formula is given by:
\[x = \frac{-b \pm \sqrt{b^2 - 4ac}}{2a}\]

\subsection{Matrix Example}
A sample matrix:
\[
\begin{pmatrix}
a_{11} & a_{12} & a_{13} \\
a_{21} & a_{22} & a_{23} \\
a_{31} & a_{32} & a_{33}
\end{pmatrix}
\]

\section{Tables}
Below is a sample table showing data:

\begin{table}[h]
\centering
\begin{tabular}{@{}lrrr@{}}
\toprule
\textbf{Category} & \textbf{Value 1} & \textbf{Value 2} & \textbf{Total} \\
\midrule
Item A & 10.5 & 20.3 & 30.8 \\
Item B & 15.2 & 25.1 & 40.3 \\
Item C & 22.1 & 12.7 & 34.8 \\
\bottomrule
\end{tabular}
\caption{Sample Data Table}
\label{tab:sample}
\end{table}

\section{Lists}
\subsection{Bullet Points}
\begin{itemize}
  \item First item
  \item Second item
  \item Third item with sub-items:
    \begin{itemize}
      \item Sub-item 1
      \item Sub-item 2
    \end{itemize}
\end{itemize}

\subsection{Numbered List}
\begin{enumerate}
  \item First step
  \item Second step
  \item Third step
\end{enumerate}

\section{Citations and References}
When you use the citation search feature, references will be automatically managed. For example, you might cite a paper like this: \cite{example2024}. All citations will be automatically added to your \texttt{cite.bib} file.

\section{Conclusion}
This template demonstrates basic \LaTeX{} features including mathematical formulas, tables, and lists. The bibliography system is now set up to work seamlessly with Bibby AI's citation management. Feel free to modify it for your needs!

\bibliographystyle{plain}
\bibliography{cite}

\end{document}\documentclass{article}
\usepackage[utf8]{inputenc}
\usepackage{booktabs}
\usepackage{amsmath}
\usepackage{natbib}

\title{My Document}
\author{Your Name}
\date{\today}

\begin{document}

\maketitle

\section{Introduction}
This is a sample document demonstrating various \LaTeX{} features including mathematical equations, tables, and formatting. You can cite papers using the citation search feature, and they will be automatically added to your \texttt{cite.bib} file.

\section{Mathematical Expressions}
Here are some example equations:

\subsection{Quadratic Formula}
The quadratic formula is given by:
\[x = \frac{-b \pm \sqrt{b^2 - 4ac}}{2a}\]

\subsection{Matrix Example}
A sample matrix:
\[
\begin{pmatrix}
a_{11} & a_{12} & a_{13} \\
a_{21} & a_{22} & a_{23} \\
a_{31} & a_{32} & a_{33}
\end{pmatrix}
\]

\section{Tables}
Below is a sample table showing data:

\begin{table}[h]
\centering
\begin{tabular}{@{}lrrr@{}}
\toprule
\textbf{Category} & \textbf{Value 1} & \textbf{Value 2} & \textbf{Total} \\
\midrule
Item A & 10.5 & 20.3 & 30.8 \\
Item B & 15.2 & 25.1 & 40.3 \\
Item C & 22.1 & 12.7 & 34.8 \\
\bottomrule
\end{tabular}
\caption{Sample Data Table}
\label{tab:sample}
\end{table}

\section{Lists}
\subsection{Bullet Points}
\begin{itemize}
  \item First item
  \item Second item
  \item Third item with sub-items:
    \begin{itemize}
      \item Sub-item 1
      \item Sub-item 2
    \end{itemize}
\end{itemize}

\subsection{Numbered List}
\begin{enumerate}
  \item First step
  \item Second step
  \item Third step
\end{enumerate}

\section{Citations and References}
When you use the citation search feature, references will be automatically managed. For example, you might cite a paper like this: \cite{example2024}. All citations will be automatically added to your \texttt{cite.bib} file.

\section{Conclusion}
This template demonstrates basic \LaTeX{} features including mathematical formulas, tables, and lists. The bibliography system is now set up to work seamlessly with Bibby AI's citation management. Feel free to modify it for your needs!

\bibliographystyle{plain}
\bibliography{cite}

\end{document}\documentclass{article}
\usepackage[utf8]{inputenc}
\usepackage{booktabs}
\usepackage{amsmath}
\usepackage{natbib}

\title{My Document}
\author{Your Name}
\date{\today}

\begin{document}

\maketitle

\section{Introduction}
This is a sample document demonstrating various \LaTeX{} features including mathematical equations, tables, and formatting. You can cite papers using the citation search feature, and they will be automatically added to your \texttt{cite.bib} file.

\section{Mathematical Expressions}
Here are some example equations:

\subsection{Quadratic Formula}
The quadratic formula is given by:
\[x = \frac{-b \pm \sqrt{b^2 - 4ac}}{2a}\]

\subsection{Matrix Example}
A sample matrix:
\[
\begin{pmatrix}
a_{11} & a_{12} & a_{13} \\
a_{21} & a_{22} & a_{23} \\
a_{31} & a_{32} & a_{33}
\end{pmatrix}
\]

\section{Tables}
Below is a sample table showing data:

\begin{table}[h]
\centering
\begin{tabular}{@{}lrrr@{}}
\toprule
\textbf{Category} & \textbf{Value 1} & \textbf{Value 2} & \textbf{Total} \\
\midrule
Item A & 10.5 & 20.3 & 30.8 \\
Item B & 15.2 & 25.1 & 40.3 \\
Item C & 22.1 & 12.7 & 34.8 \\
\bottomrule
\end{tabular}
\caption{Sample Data Table}
\label{tab:sample}
\end{table}

\section{Lists}
\subsection{Bullet Points}
\begin{itemize}
  \item First item
  \item Second item
  \item Third item with sub-items:
    \begin{itemize}
      \item Sub-item 1
      \item Sub-item 2
    \end{itemize}
\end{itemize}

\subsection{Numbered List}
\begin{enumerate}
  \item First step
  \item Second step
  \item Third step
\end{enumerate}

\section{Citations and References}
When you use the citation search feature, references will be automatically managed. For example, you might cite a paper like this: \cite{example2024}. All citations will be automatically added to your \texttt{cite.bib} file.

\section{Conclusion}
This template demonstrates basic \LaTeX{} features including mathematical formulas, tables, and lists. The bibliography system is now set up to work seamlessly with Bibby AI's citation management. Feel free to modify it for your needs!

\bibliographystyle{plain}
\bibliography{cite}

\end{document}\documentclass{article}
\usepackage[utf8]{inputenc}
\usepackage{booktabs}
\usepackage{amsmath}
\usepackage{natbib}

\title{My Document}
\author{Your Name}
\date{\today}

\begin{document}

\maketitle

\section{Introduction}
This is a sample document demonstrating various \LaTeX{} features including mathematical equations, tables, and formatting. You can cite papers using the citation search feature, and they will be automatically added to your \texttt{cite.bib} file.

\section{Mathematical Expressions}
Here are some example equations:

\subsection{Quadratic Formula}
The quadratic formula is given by:
\[x = \frac{-b \pm \sqrt{b^2 - 4ac}}{2a}\]

\subsection{Matrix Example}
A sample matrix:
\[
\begin{pmatrix}
a_{11} & a_{12} & a_{13} \\
a_{21} & a_{22} & a_{23} \\
a_{31} & a_{32} & a_{33}
\end{pmatrix}
\]

\section{Tables}
Below is a sample table showing data:

\begin{table}[h]
\centering
\begin{tabular}{@{}lrrr@{}}
\toprule
\textbf{Category} & \textbf{Value 1} & \textbf{Value 2} & \textbf{Total} \\
\midrule
Item A & 10.5 & 20.3 & 30.8 \\
Item B & 15.2 & 25.1 & 40.3 \\
Item C & 22.1 & 12.7 & 34.8 \\
\bottomrule
\end{tabular}
\caption{Sample Data Table}
\label{tab:sample}
\end{table}

\section{Lists}
\subsection{Bullet Points}
\begin{itemize}
  \item First item
  \item Second item
  \item Third item with sub-items:
    \begin{itemize}
      \item Sub-item 1
      \item Sub-item 2
    \end{itemize}
\end{itemize}

\subsection{Numbered List}
\begin{enumerate}
  \item First step
  \item Second step
  \item Third step
\end{enumerate}

\section{Citations and References}
When you use the citation search feature, references will be automatically managed. For example, you might cite a paper like this: \cite{example2024}. All citations will be automatically added to your \texttt{cite.bib} file.

\section{Conclusion}
This template demonstrates basic \LaTeX{} features including mathematical formulas, tables, and lists. The bibliography system is now set up to work seamlessly with Bibby AI's citation management. Feel free to modify it for your needs!

\bibliographystyle{plain}
\bibliography{cite}

\end{document}\documentclass{article}
\usepackage[utf8]{inputenc}
\usepackage{booktabs}
\usepackage{amsmath}
\usepackage{natbib}

\title{My Document}
\author{Your Name}
\date{\today}

\begin{document}

\maketitle

\section{Introduction}
This is a sample document demonstrating various \LaTeX{} features including mathematical equations, tables, and formatting. You can cite papers using the citation search feature, and they will be automatically added to your \texttt{cite.bib} file.

\section{Mathematical Expressions}
Here are some example equations:

\subsection{Quadratic Formula}
The quadratic formula is given by:
\[x = \frac{-b \pm \sqrt{b^2 - 4ac}}{2a}\]

\subsection{Matrix Example}
A sample matrix:
\[
\begin{pmatrix}
a_{11} & a_{12} & a_{13} \\
a_{21} & a_{22} & a_{23} \\
a_{31} & a_{32} & a_{33}
\end{pmatrix}
\]

\section{Tables}
Below is a sample table showing data:

\begin{table}[h]
\centering
\begin{tabular}{@{}lrrr@{}}
\toprule
\textbf{Category} & \textbf{Value 1} & \textbf{Value 2} & \textbf{Total} \\
\midrule
Item A & 10.5 & 20.3 & 30.8 \\
Item B & 15.2 & 25.1 & 40.3 \\
Item C & 22.1 & 12.7 & 34.8 \\
\bottomrule
\end{tabular}
\caption{Sample Data Table}
\label{tab:sample}
\end{table}

\section{Lists}
\subsection{Bullet Points}
\begin{itemize}
  \item First item
  \item Second item
  \item Third item with sub-items:
    \begin{itemize}
      \item Sub-item 1
      \item Sub-item 2
    \end{itemize}
\end{itemize}

\subsection{Numbered List}
\begin{enumerate}
  \item First step
  \item Second step
  \item Third step
\end{enumerate}

\section{Citations and References}
When you use the citation search feature, references will be automatically managed. For example, you might cite a paper like this: \cite{example2024}. All citations will be automatically added to your \texttt{cite.bib} file.

\section{Conclusion}
This template demonstrates basic \LaTeX{} features including mathematical formulas, tables, and lists. The bibliography system is now set up to work seamlessly with Bibby AI's citation management. Feel free to modify it for your needs!

\bibliographystyle{plain}
\bibliography{cite}

\end{document}\documentclass{article}
\usepackage[utf8]{inputenc}
\usepackage{booktabs}
\usepackage{amsmath}
\usepackage{natbib}

\title{My Document}
\author{Your Name}
\date{\today}

\begin{document}

\maketitle

\section{Introduction}
This is a sample document demonstrating various \LaTeX{} features including mathematical equations, tables, and formatting. You can cite papers using the citation search feature, and they will be automatically added to your \texttt{cite.bib} file.

\section{Mathematical Expressions}
Here are some example equations:

\subsection{Quadratic Formula}
The quadratic formula is given by:
\[x = \frac{-b \pm \sqrt{b^2 - 4ac}}{2a}\]

\subsection{Matrix Example}
A sample matrix:
\[
\begin{pmatrix}
a_{11} & a_{12} & a_{13} \\
a_{21} & a_{22} & a_{23} \\
a_{31} & a_{32} & a_{33}
\end{pmatrix}
\]

\section{Tables}
Below is a sample table showing data:

\begin{table}[h]
\centering
\begin{tabular}{@{}lrrr@{}}
\toprule
\textbf{Category} & \textbf{Value 1} & \textbf{Value 2} & \textbf{Total} \\
\midrule
Item A & 10.5 & 20.3 & 30.8 \\
Item B & 15.2 & 25.1 & 40.3 \\
Item C & 22.1 & 12.7 & 34.8 \\
\bottomrule
\end{tabular}
\caption{Sample Data Table}
\label{tab:sample}
\end{table}

\section{Lists}
\subsection{Bullet Points}
\begin{itemize}
  \item First item
  \item Second item
  \item Third item with sub-items:
    \begin{itemize}
      \item Sub-item 1
      \item Sub-item 2
    \end{itemize}
\end{itemize}

\subsection{Numbered List}
\begin{enumerate}
  \item First step
  \item Second step
  \item Third step
\end{enumerate}

\section{Citations and References}
When you use the citation search feature, references will be automatically managed. For example, you might cite a paper like this: \cite{example2024}. All citations will be automatically added to your \texttt{cite.bib} file.

\section{Conclusion}
This template demonstrates basic \LaTeX{} features including mathematical formulas, tables, and lists. The bibliography system is now set up to work seamlessly with Bibby AI's citation management. Feel free to modify it for your needs!

\bibliographystyle{plain}
\bibliography{cite}

\end{document}\documentclass{article}
\usepackage[utf8]{inputenc}
\usepackage{booktabs}
\usepackage{amsmath}
\usepackage{natbib}

\title{My Document}
\author{Your Name}
\date{\today}

\begin{document}

\maketitle

\section{Introduction}
This is a sample document demonstrating various \LaTeX{} features including mathematical equations, tables, and formatting. You can cite papers using the citation search feature, and they will be automatically added to your \texttt{cite.bib} file.

\section{Mathematical Expressions}
Here are some example equations:

\subsection{Quadratic Formula}
The quadratic formula is given by:
\[x = \frac{-b \pm \sqrt{b^2 - 4ac}}{2a}\]

\subsection{Matrix Example}
A sample matrix:
\[
\begin{pmatrix}
a_{11} & a_{12} & a_{13} \\
a_{21} & a_{22} & a_{23} \\
a_{31} & a_{32} & a_{33}
\end{pmatrix}
\]

\section{Tables}
Below is a sample table showing data:

\begin{table}[h]
\centering
\begin{tabular}{@{}lrrr@{}}
\toprule
\textbf{Category} & \textbf{Value 1} & \textbf{Value 2} & \textbf{Total} \\
\midrule
Item A & 10.5 & 20.3 & 30.8 \\
Item B & 15.2 & 25.1 & 40.3 \\
Item C & 22.1 & 12.7 & 34.8 \\
\bottomrule
\end{tabular}
\caption{Sample Data Table}
\label{tab:sample}
\end{table}

\section{Lists}
\subsection{Bullet Points}
\begin{itemize}
  \item First item
  \item Second item
  \item Third item with sub-items:
    \begin{itemize}
      \item Sub-item 1
      \item Sub-item 2
    \end{itemize}
\end{itemize}

\subsection{Numbered List}
\begin{enumerate}
  \item First step
  \item Second step
  \item Third step
\end{enumerate}

\section{Citations and References}
When you use the citation search feature, references will be automatically managed. For example, you might cite a paper like this: \cite{example2024}. All citations will be automatically added to your \texttt{cite.bib} file.

\section{Conclusion}
This template demonstrates basic \LaTeX{} features including mathematical formulas, tables, and lists. The bibliography system is now set up to work seamlessly with Bibby AI's citation management. Feel free to modify it for your needs!

\bibliographystyle{plain}
\bibliography{cite}

\end{document}\documentclass{article}
\usepackage[utf8]{inputenc}
\usepackage{booktabs}
\usepackage{amsmath}
\usepackage{natbib}

\title{My Document}
\author{Your Name}
\date{\today}

\begin{document}

\maketitle

\section{Introduction}
This is a sample document demonstrating various \LaTeX{} features including mathematical equations, tables, and formatting. You can cite papers using the citation search feature, and they will be automatically added to your \texttt{cite.bib} file.

\section{Mathematical Expressions}
Here are some example equations:

\subsection{Quadratic Formula}
The quadratic formula is given by:
\[x = \frac{-b \pm \sqrt{b^2 - 4ac}}{2a}\]

\subsection{Matrix Example}
A sample matrix:
\[
\begin{pmatrix}
a_{11} & a_{12} & a_{13} \\
a_{21} & a_{22} & a_{23} \\
a_{31} & a_{32} & a_{33}
\end{pmatrix}
\]

\section{Tables}
Below is a sample table showing data:

\begin{table}[h]
\centering
\begin{tabular}{@{}lrrr@{}}
\toprule
\textbf{Category} & \textbf{Value 1} & \textbf{Value 2} & \textbf{Total} \\
\midrule
Item A & 10.5 & 20.3 & 30.8 \\
Item B & 15.2 & 25.1 & 40.3 \\
Item C & 22.1 & 12.7 & 34.8 \\
\bottomrule
\end{tabular}
\caption{Sample Data Table}
\label{tab:sample}
\end{table}

\section{Lists}
\subsection{Bullet Points}
\begin{itemize}
  \item First item
  \item Second item
  \item Third item with sub-items:
    \begin{itemize}
      \item Sub-item 1
      \item Sub-item 2
    \end{itemize}
\end{itemize}

\subsection{Numbered List}
\begin{enumerate}
  \item First step
  \item Second step
  \item Third step
\end{enumerate}

\section{Citations and References}
When you use the citation search feature, references will be automatically managed. For example, you might cite a paper like this: \cite{example2024}. All citations will be automatically added to your \texttt{cite.bib} file.

\section{Conclusion}
This template demonstrates basic \LaTeX{} features including mathematical formulas, tables, and lists. The bibliography system is now set up to work seamlessly with Bibby AI's citation management. Feel free to modify it for your needs!

\bibliographystyle{plain}
\bibliography{cite}

\end{document}\documentclass{article}
\usepackage[utf8]{inputenc}
\usepackage{booktabs}
\usepackage{amsmath}
\usepackage{natbib}

\title{My Document}
\author{Your Name}
\date{\today}

\begin{document}

\maketitle

\section{Introduction}
This is a sample document demonstrating various \LaTeX{} features including mathematical equations, tables, and formatting. You can cite papers using the citation search feature, and they will be automatically added to your \texttt{cite.bib} file.

\section{Mathematical Expressions}
Here are some example equations:

\subsection{Quadratic Formula}
The quadratic formula is given by:
\[x = \frac{-b \pm \sqrt{b^2 - 4ac}}{2a}\]

\subsection{Matrix Example}
A sample matrix:
\[
\begin{pmatrix}
a_{11} & a_{12} & a_{13} \\
a_{21} & a_{22} & a_{23} \\
a_{31} & a_{32} & a_{33}
\end{pmatrix}
\]

\section{Tables}
Below is a sample table showing data:

\begin{table}[h]
\centering
\begin{tabular}{@{}lrrr@{}}
\toprule
\textbf{Category} & \textbf{Value 1} & \textbf{Value 2} & \textbf{Total} \\
\midrule
Item A & 10.5 & 20.3 & 30.8 \\
Item B & 15.2 & 25.1 & 40.3 \\
Item C & 22.1 & 12.7 & 34.8 \\
\bottomrule
\end{tabular}
\caption{Sample Data Table}
\label{tab:sample}
\end{table}

\section{Lists}
\subsection{Bullet Points}
\begin{itemize}
  \item First item
  \item Second item
  \item Third item with sub-items:
    \begin{itemize}
      \item Sub-item 1
      \item Sub-item 2
    \end{itemize}
\end{itemize}

\subsection{Numbered List}
\begin{enumerate}
  \item First step
  \item Second step
  \item Third step
\end{enumerate}

\section{Citations and References}
When you use the citation search feature, references will be automatically managed. For example, you might cite a paper like this: \cite{example2024}. All citations will be automatically added to your \texttt{cite.bib} file.

\section{Conclusion}
This template demonstrates basic \LaTeX{} features including mathematical formulas, tables, and lists. The bibliography system is now set up to work seamlessly with Bibby AI's citation management. Feel free to modify it for your needs!

\bibliographystyle{plain}
\bibliography{cite}

\end{document}\documentclass{article}
\usepackage[utf8]{inputenc}
\usepackage{booktabs}
\usepackage{amsmath}
\usepackage{natbib}

\title{My Document}
\author{Your Name}
\date{\today}

\begin{document}

\maketitle

\section{Introduction}
This is a sample document demonstrating various \LaTeX{} features including mathematical equations, tables, and formatting. You can cite papers using the citation search feature, and they will be automatically added to your \texttt{cite.bib} file.

\section{Mathematical Expressions}
Here are some example equations:

\subsection{Quadratic Formula}
The quadratic formula is given by:
\[x = \frac{-b \pm \sqrt{b^2 - 4ac}}{2a}\]

\subsection{Matrix Example}
A sample matrix:
\[
\begin{pmatrix}
a_{11} & a_{12} & a_{13} \\
a_{21} & a_{22} & a_{23} \\
a_{31} & a_{32} & a_{33}
\end{pmatrix}
\]

\section{Tables}
Below is a sample table showing data:

\begin{table}[h]
\centering
\begin{tabular}{@{}lrrr@{}}
\toprule
\textbf{Category} & \textbf{Value 1} & \textbf{Value 2} & \textbf{Total} \\
\midrule
Item A & 10.5 & 20.3 & 30.8 \\
Item B & 15.2 & 25.1 & 40.3 \\
Item C & 22.1 & 12.7 & 34.8 \\
\bottomrule
\end{tabular}
\caption{Sample Data Table}
\label{tab:sample}
\end{table}

\section{Lists}
\subsection{Bullet Points}
\begin{itemize}
  \item First item
  \item Second item
  \item Third item with sub-items:
    \begin{itemize}
      \item Sub-item 1
      \item Sub-item 2
    \end{itemize}
\end{itemize}

\subsection{Numbered List}
\begin{enumerate}
  \item First step
  \item Second step
  \item Third step
\end{enumerate}

\section{Citations and References}
When you use the citation search feature, references will be automatically managed. For example, you might cite a paper like this: \cite{example2024}. All citations will be automatically added to your \texttt{cite.bib} file.

\section{Conclusion}
This template demonstrates basic \LaTeX{} features including mathematical formulas, tables, and lists. The bibliography system is now set up to work seamlessly with Bibby AI's citation management. Feel free to modify it for your needs!

\bibliographystyle{plain}
\bibliography{cite}

\end{document}\documentclass{article}
\usepackage[utf8]{inputenc}
\usepackage{booktabs}
\usepackage{amsmath}
\usepackage{natbib}

\title{My Document}
\author{Your Name}
\date{\today}

\begin{document}

\maketitle

\section{Introduction}
This is a sample document demonstrating various \LaTeX{} features including mathematical equations, tables, and formatting. You can cite papers using the citation search feature, and they will be automatically added to your \texttt{cite.bib} file.

\section{Mathematical Expressions}
Here are some example equations:

\subsection{Quadratic Formula}
The quadratic formula is given by:
\[x = \frac{-b \pm \sqrt{b^2 - 4ac}}{2a}\]

\subsection{Matrix Example}
A sample matrix:
\[
\begin{pmatrix}
a_{11} & a_{12} & a_{13} \\
a_{21} & a_{22} & a_{23} \\
a_{31} & a_{32} & a_{33}
\end{pmatrix}
\]

\section{Tables}
Below is a sample table showing data:

\begin{table}[h]
\centering
\begin{tabular}{@{}lrrr@{}}
\toprule
\textbf{Category} & \textbf{Value 1} & \textbf{Value 2} & \textbf{Total} \\
\midrule
Item A & 10.5 & 20.3 & 30.8 \\
Item B & 15.2 & 25.1 & 40.3 \\
Item C & 22.1 & 12.7 & 34.8 \\
\bottomrule
\end{tabular}
\caption{Sample Data Table}
\label{tab:sample}
\end{table}

\section{Lists}
\subsection{Bullet Points}
\begin{itemize}
  \item First item
  \item Second item
  \item Third item with sub-items:
    \begin{itemize}
      \item Sub-item 1
      \item Sub-item 2
    \end{itemize}
\end{itemize}

\subsection{Numbered List}
\begin{enumerate}
  \item First step
  \item Second step
  \item Third step
\end{enumerate}

\section{Citations and References}
When you use the citation search feature, references will be automatically managed. For example, you might cite a paper like this: \cite{example2024}. All citations will be automatically added to your \texttt{cite.bib} file.

\section{Conclusion}
This template demonstrates basic \LaTeX{} features including mathematical formulas, tables, and lists. The bibliography system is now set up to work seamlessly with Bibby AI's citation management. Feel free to modify it for your needs!

\bibliographystyle{plain}
\bibliography{cite}

\end{document}\documentclass{article}
\usepackage[utf8]{inputenc}
\usepackage{booktabs}
\usepackage{amsmath}
\usepackage{natbib}

\title{My Document}
\author{Your Name}
\date{\today}

\begin{document}

\maketitle

\section{Introduction}
This is a sample document demonstrating various \LaTeX{} features including mathematical equations, tables, and formatting. You can cite papers using the citation search feature, and they will be automatically added to your \texttt{cite.bib} file.

\section{Mathematical Expressions}
Here are some example equations:

\subsection{Quadratic Formula}
The quadratic formula is given by:
\[x = \frac{-b \pm \sqrt{b^2 - 4ac}}{2a}\]

\subsection{Matrix Example}
A sample matrix:
\[
\begin{pmatrix}
a_{11} & a_{12} & a_{13} \\
a_{21} & a_{22} & a_{23} \\
a_{31} & a_{32} & a_{33}
\end{pmatrix}
\]

\section{Tables}
Below is a sample table showing data:

\begin{table}[h]
\centering
\begin{tabular}{@{}lrrr@{}}
\toprule
\textbf{Category} & \textbf{Value 1} & \textbf{Value 2} & \textbf{Total} \\
\midrule
Item A & 10.5 & 20.3 & 30.8 \\
Item B & 15.2 & 25.1 & 40.3 \\
Item C & 22.1 & 12.7 & 34.8 \\
\bottomrule
\end{tabular}
\caption{Sample Data Table}
\label{tab:sample}
\end{table}

\section{Lists}
\subsection{Bullet Points}
\begin{itemize}
  \item First item
  \item Second item
  \item Third item with sub-items:
    \begin{itemize}
      \item Sub-item 1
      \item Sub-item 2
    \end{itemize}
\end{itemize}

\subsection{Numbered List}
\begin{enumerate}
  \item First step
  \item Second step
  \item Third step
\end{enumerate}

\section{Citations and References}
When you use the citation search feature, references will be automatically managed. For example, you might cite a paper like this: \cite{example2024}. All citations will be automatically added to your \texttt{cite.bib} file.

\section{Conclusion}
This template demonstrates basic \LaTeX{} features including mathematical formulas, tables, and lists. The bibliography system is now set up to work seamlessly with Bibby AI's citation management. Feel free to modify it for your needs!

\bibliographystyle{plain}
\bibliography{cite}

\end{document}\documentclass{article}
\usepackage[utf8]{inputenc}
\usepackage{booktabs}
\usepackage{amsmath}
\usepackage{natbib}

\title{My Document}
\author{Your Name}
\date{\today}

\begin{document}

\maketitle

\section{Introduction}
This is a sample document demonstrating various \LaTeX{} features including mathematical equations, tables, and formatting. You can cite papers using the citation search feature, and they will be automatically added to your \texttt{cite.bib} file.

\section{Mathematical Expressions}
Here are some example equations:

\subsection{Quadratic Formula}
The quadratic formula is given by:
\[x = \frac{-b \pm \sqrt{b^2 - 4ac}}{2a}\]

\subsection{Matrix Example}
A sample matrix:
\[
\begin{pmatrix}
a_{11} & a_{12} & a_{13} \\
a_{21} & a_{22} & a_{23} \\
a_{31} & a_{32} & a_{33}
\end{pmatrix}
\]

\section{Tables}
Below is a sample table showing data:

\begin{table}[h]
\centering
\begin{tabular}{@{}lrrr@{}}
\toprule
\textbf{Category} & \textbf{Value 1} & \textbf{Value 2} & \textbf{Total} \\
\midrule
Item A & 10.5 & 20.3 & 30.8 \\
Item B & 15.2 & 25.1 & 40.3 \\
Item C & 22.1 & 12.7 & 34.8 \\
\bottomrule
\end{tabular}
\caption{Sample Data Table}
\label{tab:sample}
\end{table}

\section{Lists}
\subsection{Bullet Points}
\begin{itemize}
  \item First item
  \item Second item
  \item Third item with sub-items:
    \begin{itemize}
      \item Sub-item 1
      \item Sub-item 2
    \end{itemize}
\end{itemize}

\subsection{Numbered List}
\begin{enumerate}
  \item First step
  \item Second step
  \item Third step
\end{enumerate}

\section{Citations and References}
When you use the citation search feature, references will be automatically managed. For example, you might cite a paper like this: \cite{example2024}. All citations will be automatically added to your \texttt{cite.bib} file.

\section{Conclusion}
This template demonstrates basic \LaTeX{} features including mathematical formulas, tables, and lists. The bibliography system is now set up to work seamlessly with Bibby AI's citation management. Feel free to modify it for your needs!

\bibliographystyle{plain}
\bibliography{cite}

\end{document}\documentclass{article}
\usepackage[utf8]{inputenc}
\usepackage{booktabs}
\usepackage{amsmath}
\usepackage{natbib}

\title{My Document}
\author{Your Name}
\date{\today}

\begin{document}

\maketitle

\section{Introduction}
This is a sample document demonstrating various \LaTeX{} features including mathematical equations, tables, and formatting. You can cite papers using the citation search feature, and they will be automatically added to your \texttt{cite.bib} file.

\section{Mathematical Expressions}
Here are some example equations:

\subsection{Quadratic Formula}
The quadratic formula is given by:
\[x = \frac{-b \pm \sqrt{b^2 - 4ac}}{2a}\]

\subsection{Matrix Example}
A sample matrix:
\[
\begin{pmatrix}
a_{11} & a_{12} & a_{13} \\
a_{21} & a_{22} & a_{23} \\
a_{31} & a_{32} & a_{33}
\end{pmatrix}
\]

\section{Tables}
Below is a sample table showing data:

\begin{table}[h]
\centering
\begin{tabular}{@{}lrrr@{}}
\toprule
\textbf{Category} & \textbf{Value 1} & \textbf{Value 2} & \textbf{Total} \\
\midrule
Item A & 10.5 & 20.3 & 30.8 \\
Item B & 15.2 & 25.1 & 40.3 \\
Item C & 22.1 & 12.7 & 34.8 \\
\bottomrule
\end{tabular}
\caption{Sample Data Table}
\label{tab:sample}
\end{table}

\section{Lists}
\subsection{Bullet Points}
\begin{itemize}
  \item First item
  \item Second item
  \item Third item with sub-items:
    \begin{itemize}
      \item Sub-item 1
      \item Sub-item 2
    \end{itemize}
\end{itemize}

\subsection{Numbered List}
\begin{enumerate}
  \item First step
  \item Second step
  \item Third step
\end{enumerate}

\section{Citations and References}
When you use the citation search feature, references will be automatically managed. For example, you might cite a paper like this: \cite{example2024}. All citations will be automatically added to your \texttt{cite.bib} file.

\section{Conclusion}
This template demonstrates basic \LaTeX{} features including mathematical formulas, tables, and lists. The bibliography system is now set up to work seamlessly with Bibby AI's citation management. Feel free to modify it for your needs!

\bibliographystyle{plain}
\bibliography{cite}

\end{document}\documentclass{article}
\usepackage[utf8]{inputenc}
\usepackage{booktabs}
\usepackage{amsmath}
\usepackage{natbib}

\title{My Document}
\author{Your Name}
\date{\today}

\begin{document}

\maketitle

\section{Introduction}
This is a sample document demonstrating various \LaTeX{} features including mathematical equations, tables, and formatting. You can cite papers using the citation search feature, and they will be automatically added to your \texttt{cite.bib} file.

\section{Mathematical Expressions}
Here are some example equations:

\subsection{Quadratic Formula}
The quadratic formula is given by:
\[x = \frac{-b \pm \sqrt{b^2 - 4ac}}{2a}\]

\subsection{Matrix Example}
A sample matrix:
\[
\begin{pmatrix}
a_{11} & a_{12} & a_{13} \\
a_{21} & a_{22} & a_{23} \\
a_{31} & a_{32} & a_{33}
\end{pmatrix}
\]

\section{Tables}
Below is a sample table showing data:

\begin{table}[h]
\centering
\begin{tabular}{@{}lrrr@{}}
\toprule
\textbf{Category} & \textbf{Value 1} & \textbf{Value 2} & \textbf{Total} \\
\midrule
Item A & 10.5 & 20.3 & 30.8 \\
Item B & 15.2 & 25.1 & 40.3 \\
Item C & 22.1 & 12.7 & 34.8 \\
\bottomrule
\end{tabular}
\caption{Sample Data Table}
\label{tab:sample}
\end{table}

\section{Lists}
\subsection{Bullet Points}
\begin{itemize}
  \item First item
  \item Second item
  \item Third item with sub-items:
    \begin{itemize}
      \item Sub-item 1
      \item Sub-item 2
    \end{itemize}
\end{itemize}

\subsection{Numbered List}
\begin{enumerate}
  \item First step
  \item Second step
  \item Third step
\end{enumerate}

\section{Citations and References}
When you use the citation search feature, references will be automatically managed. For example, you might cite a paper like this: \cite{example2024}. All citations will be automatically added to your \texttt{cite.bib} file.

\section{Conclusion}
This template demonstrates basic \LaTeX{} features including mathematical formulas, tables, and lists. The bibliography system is now set up to work seamlessly with Bibby AI's citation management. Feel free to modify it for your needs!

\bibliographystyle{plain}
\bibliography{cite}

\end{document}\documentclass{article}
\usepackage[utf8]{inputenc}
\usepackage{booktabs}
\usepackage{amsmath}
\usepackage{natbib}

\title{My Document}
\author{Your Name}
\date{\today}

\begin{document}

\maketitle

\section{Introduction}
This is a sample document demonstrating various \LaTeX{} features including mathematical equations, tables, and formatting. You can cite papers using the citation search feature, and they will be automatically added to your \texttt{cite.bib} file.

\section{Mathematical Expressions}
Here are some example equations:

\subsection{Quadratic Formula}
The quadratic formula is given by:
\[x = \frac{-b \pm \sqrt{b^2 - 4ac}}{2a}\]

\subsection{Matrix Example}
A sample matrix:
\[
\begin{pmatrix}
a_{11} & a_{12} & a_{13} \\
a_{21} & a_{22} & a_{23} \\
a_{31} & a_{32} & a_{33}
\end{pmatrix}
\]

\section{Tables}
Below is a sample table showing data:

\begin{table}[h]
\centering
\begin{tabular}{@{}lrrr@{}}
\toprule
\textbf{Category} & \textbf{Value 1} & \textbf{Value 2} & \textbf{Total} \\
\midrule
Item A & 10.5 & 20.3 & 30.8 \\
Item B & 15.2 & 25.1 & 40.3 \\
Item C & 22.1 & 12.7 & 34.8 \\
\bottomrule
\end{tabular}
\caption{Sample Data Table}
\label{tab:sample}
\end{table}

\section{Lists}
\subsection{Bullet Points}
\begin{itemize}
  \item First item
  \item Second item
  \item Third item with sub-items:
    \begin{itemize}
      \item Sub-item 1
      \item Sub-item 2
    \end{itemize}
\end{itemize}

\subsection{Numbered List}
\begin{enumerate}
  \item First step
  \item Second step
  \item Third step
\end{enumerate}

\section{Citations and References}
When you use the citation search feature, references will be automatically managed. For example, you might cite a paper like this: \cite{example2024}. All citations will be automatically added to your \texttt{cite.bib} file.

\section{Conclusion}
This template demonstrates basic \LaTeX{} features including mathematical formulas, tables, and lists. The bibliography system is now set up to work seamlessly with Bibby AI's citation management. Feel free to modify it for your needs!

\bibliographystyle{plain}
\bibliography{cite}

\end{document}\documentclass{article}
\usepackage[utf8]{inputenc}
\usepackage{booktabs}
\usepackage{amsmath}
\usepackage{natbib}

\title{My Document}
\author{Your Name}
\date{\today}

\begin{document}

\maketitle

\section{Introduction}
This is a sample document demonstrating various \LaTeX{} features including mathematical equations, tables, and formatting. You can cite papers using the citation search feature, and they will be automatically added to your \texttt{cite.bib} file.

\section{Mathematical Expressions}
Here are some example equations:

\subsection{Quadratic Formula}
The quadratic formula is given by:
\[x = \frac{-b \pm \sqrt{b^2 - 4ac}}{2a}\]

\subsection{Matrix Example}
A sample matrix:
\[
\begin{pmatrix}
a_{11} & a_{12} & a_{13} \\
a_{21} & a_{22} & a_{23} \\
a_{31} & a_{32} & a_{33}
\end{pmatrix}
\]

\section{Tables}
Below is a sample table showing data:

\begin{table}[h]
\centering
\begin{tabular}{@{}lrrr@{}}
\toprule
\textbf{Category} & \textbf{Value 1} & \textbf{Value 2} & \textbf{Total} \\
\midrule
Item A & 10.5 & 20.3 & 30.8 \\
Item B & 15.2 & 25.1 & 40.3 \\
Item C & 22.1 & 12.7 & 34.8 \\
\bottomrule
\end{tabular}
\caption{Sample Data Table}
\label{tab:sample}
\end{table}

\section{Lists}
\subsection{Bullet Points}
\begin{itemize}
  \item First item
  \item Second item
  \item Third item with sub-items:
    \begin{itemize}
      \item Sub-item 1
      \item Sub-item 2
    \end{itemize}
\end{itemize}

\subsection{Numbered List}
\begin{enumerate}
  \item First step
  \item Second step
  \item Third step
\end{enumerate}

\section{Citations and References}
When you use the citation search feature, references will be automatically managed. For example, you might cite a paper like this: \cite{example2024}. All citations will be automatically added to your \texttt{cite.bib} file.

\section{Conclusion}
This template demonstrates basic \LaTeX{} features including mathematical formulas, tables, and lists. The bibliography system is now set up to work seamlessly with Bibby AI's citation management. Feel free to modify it for your needs!

\bibliographystyle{plain}
\bibliography{cite}

\end{document}\documentclass{article}
\usepackage[utf8]{inputenc}
\usepackage{booktabs}
\usepackage{amsmath}
\usepackage{natbib}

\title{My Document}
\author{Your Name}
\date{\today}

\begin{document}

\maketitle

\section{Introduction}
This is a sample document demonstrating various \LaTeX{} features including mathematical equations, tables, and formatting. You can cite papers using the citation search feature, and they will be automatically added to your \texttt{cite.bib} file.

\section{Mathematical Expressions}
Here are some example equations:

\subsection{Quadratic Formula}
The quadratic formula is given by:
\[x = \frac{-b \pm \sqrt{b^2 - 4ac}}{2a}\]

\subsection{Matrix Example}
A sample matrix:
\[
\begin{pmatrix}
a_{11} & a_{12} & a_{13} \\
a_{21} & a_{22} & a_{23} \\
a_{31} & a_{32} & a_{33}
\end{pmatrix}
\]

\section{Tables}
Below is a sample table showing data:

\begin{table}[h]
\centering
\begin{tabular}{@{}lrrr@{}}
\toprule
\textbf{Category} & \textbf{Value 1} & \textbf{Value 2} & \textbf{Total} \\
\midrule
Item A & 10.5 & 20.3 & 30.8 \\
Item B & 15.2 & 25.1 & 40.3 \\
Item C & 22.1 & 12.7 & 34.8 \\
\bottomrule
\end{tabular}
\caption{Sample Data Table}
\label{tab:sample}
\end{table}

\section{Lists}
\subsection{Bullet Points}
\begin{itemize}
  \item First item
  \item Second item
  \item Third item with sub-items:
    \begin{itemize}
      \item Sub-item 1
      \item Sub-item 2
    \end{itemize}
\end{itemize}

\subsection{Numbered List}
\begin{enumerate}
  \item First step
  \item Second step
  \item Third step
\end{enumerate}

\section{Citations and References}
When you use the citation search feature, references will be automatically managed. For example, you might cite a paper like this: \cite{example2024}. All citations will be automatically added to your \texttt{cite.bib} file.

\section{Conclusion}
This template demonstrates basic \LaTeX{} features including mathematical formulas, tables, and lists. The bibliography system is now set up to work seamlessly with Bibby AI's citation management. Feel free to modify it for your needs!

\bibliographystyle{plain}
\bibliography{cite}

\end{document}\documentclass{article}
\usepackage[utf8]{inputenc}
\usepackage{booktabs}
\usepackage{amsmath}
\usepackage{natbib}

\title{My Document}
\author{Your Name}
\date{\today}

\begin{document}

\maketitle

\section{Introduction}
This is a sample document demonstrating various \LaTeX{} features including mathematical equations, tables, and formatting. You can cite papers using the citation search feature, and they will be automatically added to your \texttt{cite.bib} file.

\section{Mathematical Expressions}
Here are some example equations:

\subsection{Quadratic Formula}
The quadratic formula is given by:
\[x = \frac{-b \pm \sqrt{b^2 - 4ac}}{2a}\]

\subsection{Matrix Example}
A sample matrix:
\[
\begin{pmatrix}
a_{11} & a_{12} & a_{13} \\
a_{21} & a_{22} & a_{23} \\
a_{31} & a_{32} & a_{33}
\end{pmatrix}
\]

\section{Tables}
Below is a sample table showing data:

\begin{table}[h]
\centering
\begin{tabular}{@{}lrrr@{}}
\toprule
\textbf{Category} & \textbf{Value 1} & \textbf{Value 2} & \textbf{Total} \\
\midrule
Item A & 10.5 & 20.3 & 30.8 \\
Item B & 15.2 & 25.1 & 40.3 \\
Item C & 22.1 & 12.7 & 34.8 \\
\bottomrule
\end{tabular}
\caption{Sample Data Table}
\label{tab:sample}
\end{table}

\section{Lists}
\subsection{Bullet Points}
\begin{itemize}
  \item First item
  \item Second item
  \item Third item with sub-items:
    \begin{itemize}
      \item Sub-item 1
      \item Sub-item 2
    \end{itemize}
\end{itemize}

\subsection{Numbered List}
\begin{enumerate}
  \item First step
  \item Second step
  \item Third step
\end{enumerate}

\section{Citations and References}
When you use the citation search feature, references will be automatically managed. For example, you might cite a paper like this: \cite{example2024}. All citations will be automatically added to your \texttt{cite.bib} file.

\section{Conclusion}
This template demonstrates basic \LaTeX{} features including mathematical formulas, tables, and lists. The bibliography system is now set up to work seamlessly with Bibby AI's citation management. Feel free to modify it for your needs!

\bibliographystyle{plain}
\bibliography{cite}

\end{document}\documentclass{article}
\usepackage[utf8]{inputenc}
\usepackage{booktabs}
\usepackage{amsmath}
\usepackage{natbib}

\title{My Document}
\author{Your Name}
\date{\today}

\begin{document}

\maketitle

\section{Introduction}
This is a sample document demonstrating various \LaTeX{} features including mathematical equations, tables, and formatting. You can cite papers using the citation search feature, and they will be automatically added to your \texttt{cite.bib} file.

\section{Mathematical Expressions}
Here are some example equations:

\subsection{Quadratic Formula}
The quadratic formula is given by:
\[x = \frac{-b \pm \sqrt{b^2 - 4ac}}{2a}\]

\subsection{Matrix Example}
A sample matrix:
\[
\begin{pmatrix}
a_{11} & a_{12} & a_{13} \\
a_{21} & a_{22} & a_{23} \\
a_{31} & a_{32} & a_{33}
\end{pmatrix}
\]

\section{Tables}
Below is a sample table showing data:

\begin{table}[h]
\centering
\begin{tabular}{@{}lrrr@{}}
\toprule
\textbf{Category} & \textbf{Value 1} & \textbf{Value 2} & \textbf{Total} \\
\midrule
Item A & 10.5 & 20.3 & 30.8 \\
Item B & 15.2 & 25.1 & 40.3 \\
Item C & 22.1 & 12.7 & 34.8 \\
\bottomrule
\end{tabular}
\caption{Sample Data Table}
\label{tab:sample}
\end{table}

\section{Lists}
\subsection{Bullet Points}
\begin{itemize}
  \item First item
  \item Second item
  \item Third item with sub-items:
    \begin{itemize}
      \item Sub-item 1
      \item Sub-item 2
    \end{itemize}
\end{itemize}

\subsection{Numbered List}
\begin{enumerate}
  \item First step
  \item Second step
  \item Third step
\end{enumerate}

\section{Citations and References}
When you use the citation search feature, references will be automatically managed. For example, you might cite a paper like this: \cite{example2024}. All citations will be automatically added to your \texttt{cite.bib} file.

\section{Conclusion}
This template demonstrates basic \LaTeX{} features including mathematical formulas, tables, and lists. The bibliography system is now set up to work seamlessly with Bibby AI's citation management. Feel free to modify it for your needs!

\bibliographystyle{plain}
\bibliography{cite}

\end{document}\documentclass{article}
\usepackage[utf8]{inputenc}
\usepackage{booktabs}
\usepackage{amsmath}
\usepackage{natbib}

\title{My Document}
\author{Your Name}
\date{\today}

\begin{document}

\maketitle

\section{Introduction}
This is a sample document demonstrating various \LaTeX{} features including mathematical equations, tables, and formatting. You can cite papers using the citation search feature, and they will be automatically added to your \texttt{cite.bib} file.

\section{Mathematical Expressions}
Here are some example equations:

\subsection{Quadratic Formula}
The quadratic formula is given by:
\[x = \frac{-b \pm \sqrt{b^2 - 4ac}}{2a}\]

\subsection{Matrix Example}
A sample matrix:
\[
\begin{pmatrix}
a_{11} & a_{12} & a_{13} \\
a_{21} & a_{22} & a_{23} \\
a_{31} & a_{32} & a_{33}
\end{pmatrix}
\]

\section{Tables}
Below is a sample table showing data:

\begin{table}[h]
\centering
\begin{tabular}{@{}lrrr@{}}
\toprule
\textbf{Category} & \textbf{Value 1} & \textbf{Value 2} & \textbf{Total} \\
\midrule
Item A & 10.5 & 20.3 & 30.8 \\
Item B & 15.2 & 25.1 & 40.3 \\
Item C & 22.1 & 12.7 & 34.8 \\
\bottomrule
\end{tabular}
\caption{Sample Data Table}
\label{tab:sample}
\end{table}

\section{Lists}
\subsection{Bullet Points}
\begin{itemize}
  \item First item
  \item Second item
  \item Third item with sub-items:
    \begin{itemize}
      \item Sub-item 1
      \item Sub-item 2
    \end{itemize}
\end{itemize}

\subsection{Numbered List}
\begin{enumerate}
  \item First step
  \item Second step
  \item Third step
\end{enumerate}

\section{Citations and References}
When you use the citation search feature, references will be automatically managed. For example, you might cite a paper like this: \cite{example2024}. All citations will be automatically added to your \texttt{cite.bib} file.

\section{Conclusion}
This template demonstrates basic \LaTeX{} features including mathematical formulas, tables, and lists. The bibliography system is now set up to work seamlessly with Bibby AI's citation management. Feel free to modify it for your needs!

\bibliographystyle{plain}
\bibliography{cite}

\end{document}\documentclass{article}
\usepackage[utf8]{inputenc}
\usepackage{booktabs}
\usepackage{amsmath}
\usepackage{natbib}

\title{My Document}
\author{Your Name}
\date{\today}

\begin{document}

\maketitle

\section{Introduction}
This is a sample document demonstrating various \LaTeX{} features including mathematical equations, tables, and formatting. You can cite papers using the citation search feature, and they will be automatically added to your \texttt{cite.bib} file.

\section{Mathematical Expressions}
Here are some example equations:

\subsection{Quadratic Formula}
The quadratic formula is given by:
\[x = \frac{-b \pm \sqrt{b^2 - 4ac}}{2a}\]

\subsection{Matrix Example}
A sample matrix:
\[
\begin{pmatrix}
a_{11} & a_{12} & a_{13} \\
a_{21} & a_{22} & a_{23} \\
a_{31} & a_{32} & a_{33}
\end{pmatrix}
\]

\section{Tables}
Below is a sample table showing data:

\begin{table}[h]
\centering
\begin{tabular}{@{}lrrr@{}}
\toprule
\textbf{Category} & \textbf{Value 1} & \textbf{Value 2} & \textbf{Total} \\
\midrule
Item A & 10.5 & 20.3 & 30.8 \\
Item B & 15.2 & 25.1 & 40.3 \\
Item C & 22.1 & 12.7 & 34.8 \\
\bottomrule
\end{tabular}
\caption{Sample Data Table}
\label{tab:sample}
\end{table}

\section{Lists}
\subsection{Bullet Points}
\begin{itemize}
  \item First item
  \item Second item
  \item Third item with sub-items:
    \begin{itemize}
      \item Sub-item 1
      \item Sub-item 2
    \end{itemize}
\end{itemize}

\subsection{Numbered List}
\begin{enumerate}
  \item First step
  \item Second step
  \item Third step
\end{enumerate}

\section{Citations and References}
When you use the citation search feature, references will be automatically managed. For example, you might cite a paper like this: \cite{example2024}. All citations will be automatically added to your \texttt{cite.bib} file.

\section{Conclusion}
This template demonstrates basic \LaTeX{} features including mathematical formulas, tables, and lists. The bibliography system is now set up to work seamlessly with Bibby AI's citation management. Feel free to modify it for your needs!

\bibliographystyle{plain}
\bibliography{cite}

\end{document}\documentclass{article}
\usepackage[utf8]{inputenc}
\usepackage{booktabs}
\usepackage{amsmath}
\usepackage{natbib}

\title{My Document}
\author{Your Name}
\date{\today}

\begin{document}

\maketitle

\section{Introduction}
This is a sample document demonstrating various \LaTeX{} features including mathematical equations, tables, and formatting. You can cite papers using the citation search feature, and they will be automatically added to your \texttt{cite.bib} file.

\section{Mathematical Expressions}
Here are some example equations:

\subsection{Quadratic Formula}
The quadratic formula is given by:
\[x = \frac{-b \pm \sqrt{b^2 - 4ac}}{2a}\]

\subsection{Matrix Example}
A sample matrix:
\[
\begin{pmatrix}
a_{11} & a_{12} & a_{13} \\
a_{21} & a_{22} & a_{23} \\
a_{31} & a_{32} & a_{33}
\end{pmatrix}
\]

\section{Tables}
Below is a sample table showing data:

\begin{table}[h]
\centering
\begin{tabular}{@{}lrrr@{}}
\toprule
\textbf{Category} & \textbf{Value 1} & \textbf{Value 2} & \textbf{Total} \\
\midrule
Item A & 10.5 & 20.3 & 30.8 \\
Item B & 15.2 & 25.1 & 40.3 \\
Item C & 22.1 & 12.7 & 34.8 \\
\bottomrule
\end{tabular}
\caption{Sample Data Table}
\label{tab:sample}
\end{table}

\section{Lists}
\subsection{Bullet Points}
\begin{itemize}
  \item First item
  \item Second item
  \item Third item with sub-items:
    \begin{itemize}
      \item Sub-item 1
      \item Sub-item 2
    \end{itemize}
\end{itemize}

\subsection{Numbered List}
\begin{enumerate}
  \item First step
  \item Second step
  \item Third step
\end{enumerate}

\section{Citations and References}
When you use the citation search feature, references will be automatically managed. For example, you might cite a paper like this: \cite{example2024}. All citations will be automatically added to your \texttt{cite.bib} file.

\section{Conclusion}
This template demonstrates basic \LaTeX{} features including mathematical formulas, tables, and lists. The bibliography system is now set up to work seamlessly with Bibby AI's citation management. Feel free to modify it for your needs!

\bibliographystyle{plain}
\bibliography{cite}

\end{document}\documentclass{article}
\usepackage[utf8]{inputenc}
\usepackage{booktabs}
\usepackage{amsmath}
\usepackage{natbib}

\title{My Document}
\author{Your Name}
\date{\today}

\begin{document}

\maketitle

\section{Introduction}
This is a sample document demonstrating various \LaTeX{} features including mathematical equations, tables, and formatting. You can cite papers using the citation search feature, and they will be automatically added to your \texttt{cite.bib} file.

\section{Mathematical Expressions}
Here are some example equations:

\subsection{Quadratic Formula}
The quadratic formula is given by:
\[x = \frac{-b \pm \sqrt{b^2 - 4ac}}{2a}\]

\subsection{Matrix Example}
A sample matrix:
\[
\begin{pmatrix}
a_{11} & a_{12} & a_{13} \\
a_{21} & a_{22} & a_{23} \\
a_{31} & a_{32} & a_{33}
\end{pmatrix}
\]

\section{Tables}
Below is a sample table showing data:

\begin{table}[h]
\centering
\begin{tabular}{@{}lrrr@{}}
\toprule
\textbf{Category} & \textbf{Value 1} & \textbf{Value 2} & \textbf{Total} \\
\midrule
Item A & 10.5 & 20.3 & 30.8 \\
Item B & 15.2 & 25.1 & 40.3 \\
Item C & 22.1 & 12.7 & 34.8 \\
\bottomrule
\end{tabular}
\caption{Sample Data Table}
\label{tab:sample}
\end{table}

\section{Lists}
\subsection{Bullet Points}
\begin{itemize}
  \item First item
  \item Second item
  \item Third item with sub-items:
    \begin{itemize}
      \item Sub-item 1
      \item Sub-item 2
    \end{itemize}
\end{itemize}

\subsection{Numbered List}
\begin{enumerate}
  \item First step
  \item Second step
  \item Third step
\end{enumerate}

\section{Citations and References}
When you use the citation search feature, references will be automatically managed. For example, you might cite a paper like this: \cite{example2024}. All citations will be automatically added to your \texttt{cite.bib} file.

\section{Conclusion}
This template demonstrates basic \LaTeX{} features including mathematical formulas, tables, and lists. The bibliography system is now set up to work seamlessly with Bibby AI's citation management. Feel free to modify it for your needs!

\bibliographystyle{plain}
\bibliography{cite}

\end{document}\documentclass{article}
\usepackage[utf8]{inputenc}
\usepackage{booktabs}
\usepackage{amsmath}
\usepackage{natbib}

\title{My Document}
\author{Your Name}
\date{\today}

\begin{document}

\maketitle

\section{Introduction}
This is a sample document demonstrating various \LaTeX{} features including mathematical equations, tables, and formatting. You can cite papers using the citation search feature, and they will be automatically added to your \texttt{cite.bib} file.

\section{Mathematical Expressions}
Here are some example equations:

\subsection{Quadratic Formula}
The quadratic formula is given by:
\[x = \frac{-b \pm \sqrt{b^2 - 4ac}}{2a}\]

\subsection{Matrix Example}
A sample matrix:
\[
\begin{pmatrix}
a_{11} & a_{12} & a_{13} \\
a_{21} & a_{22} & a_{23} \\
a_{31} & a_{32} & a_{33}
\end{pmatrix}
\]

\section{Tables}
Below is a sample table showing data:

\begin{table}[h]
\centering
\begin{tabular}{@{}lrrr@{}}
\toprule
\textbf{Category} & \textbf{Value 1} & \textbf{Value 2} & \textbf{Total} \\
\midrule
Item A & 10.5 & 20.3 & 30.8 \\
Item B & 15.2 & 25.1 & 40.3 \\
Item C & 22.1 & 12.7 & 34.8 \\
\bottomrule
\end{tabular}
\caption{Sample Data Table}
\label{tab:sample}
\end{table}

\section{Lists}
\subsection{Bullet Points}
\begin{itemize}
  \item First item
  \item Second item
  \item Third item with sub-items:
    \begin{itemize}
      \item Sub-item 1
      \item Sub-item 2
    \end{itemize}
\end{itemize}

\subsection{Numbered List}
\begin{enumerate}
  \item First step
  \item Second step
  \item Third step
\end{enumerate}

\section{Citations and References}
When you use the citation search feature, references will be automatically managed. For example, you might cite a paper like this: \cite{example2024}. All citations will be automatically added to your \texttt{cite.bib} file.

\section{Conclusion}
This template demonstrates basic \LaTeX{} features including mathematical formulas, tables, and lists. The bibliography system is now set up to work seamlessly with Bibby AI's citation management. Feel free to modify it for your needs!

\bibliographystyle{plain}
\bibliography{cite}

\end{document}\documentclass{article}
\usepackage[utf8]{inputenc}
\usepackage{booktabs}
\usepackage{amsmath}
\usepackage{natbib}

\title{My Document}
\author{Your Name}
\date{\today}

\begin{document}

\maketitle

\section{Introduction}
This is a sample document demonstrating various \LaTeX{} features including mathematical equations, tables, and formatting. You can cite papers using the citation search feature, and they will be automatically added to your \texttt{cite.bib} file.

\section{Mathematical Expressions}
Here are some example equations:

\subsection{Quadratic Formula}
The quadratic formula is given by:
\[x = \frac{-b \pm \sqrt{b^2 - 4ac}}{2a}\]

\subsection{Matrix Example}
A sample matrix:
\[
\begin{pmatrix}
a_{11} & a_{12} & a_{13} \\
a_{21} & a_{22} & a_{23} \\
a_{31} & a_{32} & a_{33}
\end{pmatrix}
\]

\section{Tables}
Below is a sample table showing data:

\begin{table}[h]
\centering
\begin{tabular}{@{}lrrr@{}}
\toprule
\textbf{Category} & \textbf{Value 1} & \textbf{Value 2} & \textbf{Total} \\
\midrule
Item A & 10.5 & 20.3 & 30.8 \\
Item B & 15.2 & 25.1 & 40.3 \\
Item C & 22.1 & 12.7 & 34.8 \\
\bottomrule
\end{tabular}
\caption{Sample Data Table}
\label{tab:sample}
\end{table}

\section{Lists}
\subsection{Bullet Points}
\begin{itemize}
  \item First item
  \item Second item
  \item Third item with sub-items:
    \begin{itemize}
      \item Sub-item 1
      \item Sub-item 2
    \end{itemize}
\end{itemize}

\subsection{Numbered List}
\begin{enumerate}
  \item First step
  \item Second step
  \item Third step
\end{enumerate}

\section{Citations and References}
When you use the citation search feature, references will be automatically managed. For example, you might cite a paper like this: \cite{example2024}. All citations will be automatically added to your \texttt{cite.bib} file.

\section{Conclusion}
This template demonstrates basic \LaTeX{} features including mathematical formulas, tables, and lists. The bibliography system is now set up to work seamlessly with Bibby AI's citation management. Feel free to modify it for your needs!

\bibliographystyle{plain}
\bibliography{cite}

\end{document}\documentclass{article}
\usepackage[utf8]{inputenc}
\usepackage{booktabs}
\usepackage{amsmath}
\usepackage{natbib}

\title{My Document}
\author{Your Name}
\date{\today}

\begin{document}

\maketitle

\section{Introduction}
This is a sample document demonstrating various \LaTeX{} features including mathematical equations, tables, and formatting. You can cite papers using the citation search feature, and they will be automatically added to your \texttt{cite.bib} file.

\section{Mathematical Expressions}
Here are some example equations:

\subsection{Quadratic Formula}
The quadratic formula is given by:
\[x = \frac{-b \pm \sqrt{b^2 - 4ac}}{2a}\]

\subsection{Matrix Example}
A sample matrix:
\[
\begin{pmatrix}
a_{11} & a_{12} & a_{13} \\
a_{21} & a_{22} & a_{23} \\
a_{31} & a_{32} & a_{33}
\end{pmatrix}
\]

\section{Tables}
Below is a sample table showing data:

\begin{table}[h]
\centering
\begin{tabular}{@{}lrrr@{}}
\toprule
\textbf{Category} & \textbf{Value 1} & \textbf{Value 2} & \textbf{Total} \\
\midrule
Item A & 10.5 & 20.3 & 30.8 \\
Item B & 15.2 & 25.1 & 40.3 \\
Item C & 22.1 & 12.7 & 34.8 \\
\bottomrule
\end{tabular}
\caption{Sample Data Table}
\label{tab:sample}
\end{table}

\section{Lists}
\subsection{Bullet Points}
\begin{itemize}
  \item First item
  \item Second item
  \item Third item with sub-items:
    \begin{itemize}
      \item Sub-item 1
      \item Sub-item 2
    \end{itemize}
\end{itemize}

\subsection{Numbered List}
\begin{enumerate}
  \item First step
  \item Second step
  \item Third step
\end{enumerate}

\section{Citations and References}
When you use the citation search feature, references will be automatically managed. For example, you might cite a paper like this: \cite{example2024}. All citations will be automatically added to your \texttt{cite.bib} file.

\section{Conclusion}
This template demonstrates basic \LaTeX{} features including mathematical formulas, tables, and lists. The bibliography system is now set up to work seamlessly with Bibby AI's citation management. Feel free to modify it for your needs!

\bibliographystyle{plain}
\bibliography{cite}

\end{document}\documentclass{article}
\usepackage[utf8]{inputenc}
\usepackage{booktabs}
\usepackage{amsmath}
\usepackage{natbib}

\title{My Document}
\author{Your Name}
\date{\today}

\begin{document}

\maketitle

\section{Introduction}
This is a sample document demonstrating various \LaTeX{} features including mathematical equations, tables, and formatting. You can cite papers using the citation search feature, and they will be automatically added to your \texttt{cite.bib} file.

\section{Mathematical Expressions}
Here are some example equations:

\subsection{Quadratic Formula}
The quadratic formula is given by:
\[x = \frac{-b \pm \sqrt{b^2 - 4ac}}{2a}\]

\subsection{Matrix Example}
A sample matrix:
\[
\begin{pmatrix}
a_{11} & a_{12} & a_{13} \\
a_{21} & a_{22} & a_{23} \\
a_{31} & a_{32} & a_{33}
\end{pmatrix}
\]

\section{Tables}
Below is a sample table showing data:

\begin{table}[h]
\centering
\begin{tabular}{@{}lrrr@{}}
\toprule
\textbf{Category} & \textbf{Value 1} & \textbf{Value 2} & \textbf{Total} \\
\midrule
Item A & 10.5 & 20.3 & 30.8 \\
Item B & 15.2 & 25.1 & 40.3 \\
Item C & 22.1 & 12.7 & 34.8 \\
\bottomrule
\end{tabular}
\caption{Sample Data Table}
\label{tab:sample}
\end{table}

\section{Lists}
\subsection{Bullet Points}
\begin{itemize}
  \item First item
  \item Second item
  \item Third item with sub-items:
    \begin{itemize}
      \item Sub-item 1
      \item Sub-item 2
    \end{itemize}
\end{itemize}

\subsection{Numbered List}
\begin{enumerate}
  \item First step
  \item Second step
  \item Third step
\end{enumerate}

\section{Citations and References}
When you use the citation search feature, references will be automatically managed. For example, you might cite a paper like this: \cite{example2024}. All citations will be automatically added to your \texttt{cite.bib} file.

\section{Conclusion}
This template demonstrates basic \LaTeX{} features including mathematical formulas, tables, and lists. The bibliography system is now set up to work seamlessly with Bibby AI's citation management. Feel free to modify it for your needs!

\bibliographystyle{plain}
\bibliography{cite}

\end{document}\documentclass{article}
\usepackage[utf8]{inputenc}
\usepackage{booktabs}
\usepackage{amsmath}
\usepackage{natbib}

\title{My Document}
\author{Your Name}
\date{\today}

\begin{document}

\maketitle

\section{Introduction}
This is a sample document demonstrating various \LaTeX{} features including mathematical equations, tables, and formatting. You can cite papers using the citation search feature, and they will be automatically added to your \texttt{cite.bib} file.

\section{Mathematical Expressions}
Here are some example equations:

\subsection{Quadratic Formula}
The quadratic formula is given by:
\[x = \frac{-b \pm \sqrt{b^2 - 4ac}}{2a}\]

\subsection{Matrix Example}
A sample matrix:
\[
\begin{pmatrix}
a_{11} & a_{12} & a_{13} \\
a_{21} & a_{22} & a_{23} \\
a_{31} & a_{32} & a_{33}
\end{pmatrix}
\]

\section{Tables}
Below is a sample table showing data:

\begin{table}[h]
\centering
\begin{tabular}{@{}lrrr@{}}
\toprule
\textbf{Category} & \textbf{Value 1} & \textbf{Value 2} & \textbf{Total} \\
\midrule
Item A & 10.5 & 20.3 & 30.8 \\
Item B & 15.2 & 25.1 & 40.3 \\
Item C & 22.1 & 12.7 & 34.8 \\
\bottomrule
\end{tabular}
\caption{Sample Data Table}
\label{tab:sample}
\end{table}

\section{Lists}
\subsection{Bullet Points}
\begin{itemize}
  \item First item
  \item Second item
  \item Third item with sub-items:
    \begin{itemize}
      \item Sub-item 1
      \item Sub-item 2
    \end{itemize}
\end{itemize}

\subsection{Numbered List}
\begin{enumerate}
  \item First step
  \item Second step
  \item Third step
\end{enumerate}

\section{Citations and References}
When you use the citation search feature, references will be automatically managed. For example, you might cite a paper like this: \cite{example2024}. All citations will be automatically added to your \texttt{cite.bib} file.

\section{Conclusion}
This template demonstrates basic \LaTeX{} features including mathematical formulas, tables, and lists. The bibliography system is now set up to work seamlessly with Bibby AI's citation management. Feel free to modify it for your needs!

\bibliographystyle{plain}
\bibliography{cite}

\end{document}\documentclass{article}
\usepackage[utf8]{inputenc}
\usepackage{booktabs}
\usepackage{amsmath}
\usepackage{natbib}

\title{My Document}
\author{Your Name}
\date{\today}

\begin{document}

\maketitle

\section{Introduction}
This is a sample document demonstrating various \LaTeX{} features including mathematical equations, tables, and formatting. You can cite papers using the citation search feature, and they will be automatically added to your \texttt{cite.bib} file.

\section{Mathematical Expressions}
Here are some example equations:

\subsection{Quadratic Formula}
The quadratic formula is given by:
\[x = \frac{-b \pm \sqrt{b^2 - 4ac}}{2a}\]

\subsection{Matrix Example}
A sample matrix:
\[
\begin{pmatrix}
a_{11} & a_{12} & a_{13} \\
a_{21} & a_{22} & a_{23} \\
a_{31} & a_{32} & a_{33}
\end{pmatrix}
\]

\section{Tables}
Below is a sample table showing data:

\begin{table}[h]
\centering
\begin{tabular}{@{}lrrr@{}}
\toprule
\textbf{Category} & \textbf{Value 1} & \textbf{Value 2} & \textbf{Total} \\
\midrule
Item A & 10.5 & 20.3 & 30.8 \\
Item B & 15.2 & 25.1 & 40.3 \\
Item C & 22.1 & 12.7 & 34.8 \\
\bottomrule
\end{tabular}
\caption{Sample Data Table}
\label{tab:sample}
\end{table}

\section{Lists}
\subsection{Bullet Points}
\begin{itemize}
  \item First item
  \item Second item
  \item Third item with sub-items:
    \begin{itemize}
      \item Sub-item 1
      \item Sub-item 2
    \end{itemize}
\end{itemize}

\subsection{Numbered List}
\begin{enumerate}
  \item First step
  \item Second step
  \item Third step
\end{enumerate}

\section{Citations and References}
When you use the citation search feature, references will be automatically managed. For example, you might cite a paper like this: \cite{example2024}. All citations will be automatically added to your \texttt{cite.bib} file.

\section{Conclusion}
This template demonstrates basic \LaTeX{} features including mathematical formulas, tables, and lists. The bibliography system is now set up to work seamlessly with Bibby AI's citation management. Feel free to modify it for your needs!

\bibliographystyle{plain}
\bibliography{cite}

\end{document}\documentclass{article}
\usepackage[utf8]{inputenc}
\usepackage{booktabs}
\usepackage{amsmath}
\usepackage{natbib}

\title{My Document}
\author{Your Name}
\date{\today}

\begin{document}

\maketitle

\section{Introduction}
This is a sample document demonstrating various \LaTeX{} features including mathematical equations, tables, and formatting. You can cite papers using the citation search feature, and they will be automatically added to your \texttt{cite.bib} file.

\section{Mathematical Expressions}
Here are some example equations:

\subsection{Quadratic Formula}
The quadratic formula is given by:
\[x = \frac{-b \pm \sqrt{b^2 - 4ac}}{2a}\]

\subsection{Matrix Example}
A sample matrix:
\[
\begin{pmatrix}
a_{11} & a_{12} & a_{13} \\
a_{21} & a_{22} & a_{23} \\
a_{31} & a_{32} & a_{33}
\end{pmatrix}
\]

\section{Tables}
Below is a sample table showing data:

\begin{table}[h]
\centering
\begin{tabular}{@{}lrrr@{}}
\toprule
\textbf{Category} & \textbf{Value 1} & \textbf{Value 2} & \textbf{Total} \\
\midrule
Item A & 10.5 & 20.3 & 30.8 \\
Item B & 15.2 & 25.1 & 40.3 \\
Item C & 22.1 & 12.7 & 34.8 \\
\bottomrule
\end{tabular}
\caption{Sample Data Table}
\label{tab:sample}
\end{table}

\section{Lists}
\subsection{Bullet Points}
\begin{itemize}
  \item First item
  \item Second item
  \item Third item with sub-items:
    \begin{itemize}
      \item Sub-item 1
      \item Sub-item 2
    \end{itemize}
\end{itemize}

\subsection{Numbered List}
\begin{enumerate}
  \item First step
  \item Second step
  \item Third step
\end{enumerate}

\section{Citations and References}
When you use the citation search feature, references will be automatically managed. For example, you might cite a paper like this: \cite{example2024}. All citations will be automatically added to your \texttt{cite.bib} file.

\section{Conclusion}
This template demonstrates basic \LaTeX{} features including mathematical formulas, tables, and lists. The bibliography system is now set up to work seamlessly with Bibby AI's citation management. Feel free to modify it for your needs!

\bibliographystyle{plain}
\bibliography{cite}

\end{document}\documentclass{article}
\usepackage[utf8]{inputenc}
\usepackage{booktabs}
\usepackage{amsmath}
\usepackage{natbib}

\title{My Document}
\author{Your Name}
\date{\today}

\begin{document}

\maketitle

\section{Introduction}
This is a sample document demonstrating various \LaTeX{} features including mathematical equations, tables, and formatting. You can cite papers using the citation search feature, and they will be automatically added to your \texttt{cite.bib} file.

\section{Mathematical Expressions}
Here are some example equations:

\subsection{Quadratic Formula}
The quadratic formula is given by:
\[x = \frac{-b \pm \sqrt{b^2 - 4ac}}{2a}\]

\subsection{Matrix Example}
A sample matrix:
\[
\begin{pmatrix}
a_{11} & a_{12} & a_{13} \\
a_{21} & a_{22} & a_{23} \\
a_{31} & a_{32} & a_{33}
\end{pmatrix}
\]

\section{Tables}
Below is a sample table showing data:

\begin{table}[h]
\centering
\begin{tabular}{@{}lrrr@{}}
\toprule
\textbf{Category} & \textbf{Value 1} & \textbf{Value 2} & \textbf{Total} \\
\midrule
Item A & 10.5 & 20.3 & 30.8 \\
Item B & 15.2 & 25.1 & 40.3 \\
Item C & 22.1 & 12.7 & 34.8 \\
\bottomrule
\end{tabular}
\caption{Sample Data Table}
\label{tab:sample}
\end{table}

\section{Lists}
\subsection{Bullet Points}
\begin{itemize}
  \item First item
  \item Second item
  \item Third item with sub-items:
    \begin{itemize}
      \item Sub-item 1
      \item Sub-item 2
    \end{itemize}
\end{itemize}

\subsection{Numbered List}
\begin{enumerate}
  \item First step
  \item Second step
  \item Third step
\end{enumerate}

\section{Citations and References}
When you use the citation search feature, references will be automatically managed. For example, you might cite a paper like this: \cite{example2024}. All citations will be automatically added to your \texttt{cite.bib} file.

\section{Conclusion}
This template demonstrates basic \LaTeX{} features including mathematical formulas, tables, and lists. The bibliography system is now set up to work seamlessly with Bibby AI's citation management. Feel free to modify it for your needs!

\bibliographystyle{plain}
\bibliography{cite}

\end{document}\documentclass{article}
\usepackage[utf8]{inputenc}
\usepackage{booktabs}
\usepackage{amsmath}
\usepackage{natbib}

\title{My Document}
\author{Your Name}
\date{\today}

\begin{document}

\maketitle

\section{Introduction}
This is a sample document demonstrating various \LaTeX{} features including mathematical equations, tables, and formatting. You can cite papers using the citation search feature, and they will be automatically added to your \texttt{cite.bib} file.

\section{Mathematical Expressions}
Here are some example equations:

\subsection{Quadratic Formula}
The quadratic formula is given by:
\[x = \frac{-b \pm \sqrt{b^2 - 4ac}}{2a}\]

\subsection{Matrix Example}
A sample matrix:
\[
\begin{pmatrix}
a_{11} & a_{12} & a_{13} \\
a_{21} & a_{22} & a_{23} \\
a_{31} & a_{32} & a_{33}
\end{pmatrix}
\]

\section{Tables}
Below is a sample table showing data:

\begin{table}[h]
\centering
\begin{tabular}{@{}lrrr@{}}
\toprule
\textbf{Category} & \textbf{Value 1} & \textbf{Value 2} & \textbf{Total} \\
\midrule
Item A & 10.5 & 20.3 & 30.8 \\
Item B & 15.2 & 25.1 & 40.3 \\
Item C & 22.1 & 12.7 & 34.8 \\
\bottomrule
\end{tabular}
\caption{Sample Data Table}
\label{tab:sample}
\end{table}

\section{Lists}
\subsection{Bullet Points}
\begin{itemize}
  \item First item
  \item Second item
  \item Third item with sub-items:
    \begin{itemize}
      \item Sub-item 1
      \item Sub-item 2
    \end{itemize}
\end{itemize}

\subsection{Numbered List}
\begin{enumerate}
  \item First step
  \item Second step
  \item Third step
\end{enumerate}

\section{Citations and References}
When you use the citation search feature, references will be automatically managed. For example, you might cite a paper like this: \cite{example2024}. All citations will be automatically added to your \texttt{cite.bib} file.

\section{Conclusion}
This template demonstrates basic \LaTeX{} features including mathematical formulas, tables, and lists. The bibliography system is now set up to work seamlessly with Bibby AI's citation management. Feel free to modify it for your needs!

\bibliographystyle{plain}
\bibliography{cite}

\end{document}\documentclass{article}
\usepackage[utf8]{inputenc}
\usepackage{booktabs}
\usepackage{amsmath}
\usepackage{natbib}

\title{My Document}
\author{Your Name}
\date{\today}

\begin{document}

\maketitle

\section{Introduction}
This is a sample document demonstrating various \LaTeX{} features including mathematical equations, tables, and formatting. You can cite papers using the citation search feature, and they will be automatically added to your \texttt{cite.bib} file.

\section{Mathematical Expressions}
Here are some example equations:

\subsection{Quadratic Formula}
The quadratic formula is given by:
\[x = \frac{-b \pm \sqrt{b^2 - 4ac}}{2a}\]

\subsection{Matrix Example}
A sample matrix:
\[
\begin{pmatrix}
a_{11} & a_{12} & a_{13} \\
a_{21} & a_{22} & a_{23} \\
a_{31} & a_{32} & a_{33}
\end{pmatrix}
\]

\section{Tables}
Below is a sample table showing data:

\begin{table}[h]
\centering
\begin{tabular}{@{}lrrr@{}}
\toprule
\textbf{Category} & \textbf{Value 1} & \textbf{Value 2} & \textbf{Total} \\
\midrule
Item A & 10.5 & 20.3 & 30.8 \\
Item B & 15.2 & 25.1 & 40.3 \\
Item C & 22.1 & 12.7 & 34.8 \\
\bottomrule
\end{tabular}
\caption{Sample Data Table}
\label{tab:sample}
\end{table}

\section{Lists}
\subsection{Bullet Points}
\begin{itemize}
  \item First item
  \item Second item
  \item Third item with sub-items:
    \begin{itemize}
      \item Sub-item 1
      \item Sub-item 2
    \end{itemize}
\end{itemize}

\subsection{Numbered List}
\begin{enumerate}
  \item First step
  \item Second step
  \item Third step
\end{enumerate}

\section{Citations and References}
When you use the citation search feature, references will be automatically managed. For example, you might cite a paper like this: \cite{example2024}. All citations will be automatically added to your \texttt{cite.bib} file.

\section{Conclusion}
This template demonstrates basic \LaTeX{} features including mathematical formulas, tables, and lists. The bibliography system is now set up to work seamlessly with Bibby AI's citation management. Feel free to modify it for your needs!

\bibliographystyle{plain}
\bibliography{cite}

\end{document}\documentclass{article}
\usepackage[utf8]{inputenc}
\usepackage{booktabs}
\usepackage{amsmath}
\usepackage{natbib}

\title{My Document}
\author{Your Name}
\date{\today}

\begin{document}

\maketitle

\section{Introduction}
This is a sample document demonstrating various \LaTeX{} features including mathematical equations, tables, and formatting. You can cite papers using the citation search feature, and they will be automatically added to your \texttt{cite.bib} file.

\section{Mathematical Expressions}
Here are some example equations:

\subsection{Quadratic Formula}
The quadratic formula is given by:
\[x = \frac{-b \pm \sqrt{b^2 - 4ac}}{2a}\]

\subsection{Matrix Example}
A sample matrix:
\[
\begin{pmatrix}
a_{11} & a_{12} & a_{13} \\
a_{21} & a_{22} & a_{23} \\
a_{31} & a_{32} & a_{33}
\end{pmatrix}
\]

\section{Tables}
Below is a sample table showing data:

\begin{table}[h]
\centering
\begin{tabular}{@{}lrrr@{}}
\toprule
\textbf{Category} & \textbf{Value 1} & \textbf{Value 2} & \textbf{Total} \\
\midrule
Item A & 10.5 & 20.3 & 30.8 \\
Item B & 15.2 & 25.1 & 40.3 \\
Item C & 22.1 & 12.7 & 34.8 \\
\bottomrule
\end{tabular}
\caption{Sample Data Table}
\label{tab:sample}
\end{table}

\section{Lists}
\subsection{Bullet Points}
\begin{itemize}
  \item First item
  \item Second item
  \item Third item with sub-items:
    \begin{itemize}
      \item Sub-item 1
      \item Sub-item 2
    \end{itemize}
\end{itemize}

\subsection{Numbered List}
\begin{enumerate}
  \item First step
  \item Second step
  \item Third step
\end{enumerate}

\section{Citations and References}
When you use the citation search feature, references will be automatically managed. For example, you might cite a paper like this: \cite{example2024}. All citations will be automatically added to your \texttt{cite.bib} file.

\section{Conclusion}
This template demonstrates basic \LaTeX{} features including mathematical formulas, tables, and lists. The bibliography system is now set up to work seamlessly with Bibby AI's citation management. Feel free to modify it for your needs!

\bibliographystyle{plain}
\bibliography{cite}

\end{document}\documentclass{article}
\usepackage[utf8]{inputenc}
\usepackage{booktabs}
\usepackage{amsmath}
\usepackage{natbib}

\title{My Document}
\author{Your Name}
\date{\today}

\begin{document}

\maketitle

\section{Introduction}
This is a sample document demonstrating various \LaTeX{} features including mathematical equations, tables, and formatting. You can cite papers using the citation search feature, and they will be automatically added to your \texttt{cite.bib} file.

\section{Mathematical Expressions}
Here are some example equations:

\subsection{Quadratic Formula}
The quadratic formula is given by:
\[x = \frac{-b \pm \sqrt{b^2 - 4ac}}{2a}\]

\subsection{Matrix Example}
A sample matrix:
\[
\begin{pmatrix}
a_{11} & a_{12} & a_{13} \\
a_{21} & a_{22} & a_{23} \\
a_{31} & a_{32} & a_{33}
\end{pmatrix}
\]

\section{Tables}
Below is a sample table showing data:

\begin{table}[h]
\centering
\begin{tabular}{@{}lrrr@{}}
\toprule
\textbf{Category} & \textbf{Value 1} & \textbf{Value 2} & \textbf{Total} \\
\midrule
Item A & 10.5 & 20.3 & 30.8 \\
Item B & 15.2 & 25.1 & 40.3 \\
Item C & 22.1 & 12.7 & 34.8 \\
\bottomrule
\end{tabular}
\caption{Sample Data Table}
\label{tab:sample}
\end{table}

\section{Lists}
\subsection{Bullet Points}
\begin{itemize}
  \item First item
  \item Second item
  \item Third item with sub-items:
    \begin{itemize}
      \item Sub-item 1
      \item Sub-item 2
    \end{itemize}
\end{itemize}

\subsection{Numbered List}
\begin{enumerate}
  \item First step
  \item Second step
  \item Third step
\end{enumerate}

\section{Citations and References}
When you use the citation search feature, references will be automatically managed. For example, you might cite a paper like this: \cite{example2024}. All citations will be automatically added to your \texttt{cite.bib} file.

\section{Conclusion}
This template demonstrates basic \LaTeX{} features including mathematical formulas, tables, and lists. The bibliography system is now set up to work seamlessly with Bibby AI's citation management. Feel free to modify it for your needs!

\bibliographystyle{plain}
\bibliography{cite}

\end{document}\documentclass{article}
\usepackage[utf8]{inputenc}
\usepackage{booktabs}
\usepackage{amsmath}
\usepackage{natbib}

\title{My Document}
\author{Your Name}
\date{\today}

\begin{document}

\maketitle

\section{Introduction}
This is a sample document demonstrating various \LaTeX{} features including mathematical equations, tables, and formatting. You can cite papers using the citation search feature, and they will be automatically added to your \texttt{cite.bib} file.

\section{Mathematical Expressions}
Here are some example equations:

\subsection{Quadratic Formula}
The quadratic formula is given by:
\[x = \frac{-b \pm \sqrt{b^2 - 4ac}}{2a}\]

\subsection{Matrix Example}
A sample matrix:
\[
\begin{pmatrix}
a_{11} & a_{12} & a_{13} \\
a_{21} & a_{22} & a_{23} \\
a_{31} & a_{32} & a_{33}
\end{pmatrix}
\]

\section{Tables}
Below is a sample table showing data:

\begin{table}[h]
\centering
\begin{tabular}{@{}lrrr@{}}
\toprule
\textbf{Category} & \textbf{Value 1} & \textbf{Value 2} & \textbf{Total} \\
\midrule
Item A & 10.5 & 20.3 & 30.8 \\
Item B & 15.2 & 25.1 & 40.3 \\
Item C & 22.1 & 12.7 & 34.8 \\
\bottomrule
\end{tabular}
\caption{Sample Data Table}
\label{tab:sample}
\end{table}

\section{Lists}
\subsection{Bullet Points}
\begin{itemize}
  \item First item
  \item Second item
  \item Third item with sub-items:
    \begin{itemize}
      \item Sub-item 1
      \item Sub-item 2
    \end{itemize}
\end{itemize}

\subsection{Numbered List}
\begin{enumerate}
  \item First step
  \item Second step
  \item Third step
\end{enumerate}

\section{Citations and References}
When you use the citation search feature, references will be automatically managed. For example, you might cite a paper like this: \cite{example2024}. All citations will be automatically added to your \texttt{cite.bib} file.

\section{Conclusion}
This template demonstrates basic \LaTeX{} features including mathematical formulas, tables, and lists. The bibliography system is now set up to work seamlessly with Bibby AI's citation management. Feel free to modify it for your needs!

\bibliographystyle{plain}
\bibliography{cite}

\end{document}\documentclass{article}
\usepackage[utf8]{inputenc}
\usepackage{booktabs}
\usepackage{amsmath}
\usepackage{natbib}

\title{My Document}
\author{Your Name}
\date{\today}

\begin{document}

\maketitle

\section{Introduction}
This is a sample document demonstrating various \LaTeX{} features including mathematical equations, tables, and formatting. You can cite papers using the citation search feature, and they will be automatically added to your \texttt{cite.bib} file.

\section{Mathematical Expressions}
Here are some example equations:

\subsection{Quadratic Formula}
The quadratic formula is given by:
\[x = \frac{-b \pm \sqrt{b^2 - 4ac}}{2a}\]

\subsection{Matrix Example}
A sample matrix:
\[
\begin{pmatrix}
a_{11} & a_{12} & a_{13} \\
a_{21} & a_{22} & a_{23} \\
a_{31} & a_{32} & a_{33}
\end{pmatrix}
\]

\section{Tables}
Below is a sample table showing data:

\begin{table}[h]
\centering
\begin{tabular}{@{}lrrr@{}}
\toprule
\textbf{Category} & \textbf{Value 1} & \textbf{Value 2} & \textbf{Total} \\
\midrule
Item A & 10.5 & 20.3 & 30.8 \\
Item B & 15.2 & 25.1 & 40.3 \\
Item C & 22.1 & 12.7 & 34.8 \\
\bottomrule
\end{tabular}
\caption{Sample Data Table}
\label{tab:sample}
\end{table}

\section{Lists}
\subsection{Bullet Points}
\begin{itemize}
  \item First item
  \item Second item
  \item Third item with sub-items:
    \begin{itemize}
      \item Sub-item 1
      \item Sub-item 2
    \end{itemize}
\end{itemize}

\subsection{Numbered List}
\begin{enumerate}
  \item First step
  \item Second step
  \item Third step
\end{enumerate}

\section{Citations and References}
When you use the citation search feature, references will be automatically managed. For example, you might cite a paper like this: \cite{example2024}. All citations will be automatically added to your \texttt{cite.bib} file.

\section{Conclusion}
This template demonstrates basic \LaTeX{} features including mathematical formulas, tables, and lists. The bibliography system is now set up to work seamlessly with Bibby AI's citation management. Feel free to modify it for your needs!

\bibliographystyle{plain}
\bibliography{cite}

\end{document}\documentclass{article}
\usepackage[utf8]{inputenc}
\usepackage{booktabs}
\usepackage{amsmath}
\usepackage{natbib}

\title{My Document}
\author{Your Name}
\date{\today}

\begin{document}

\maketitle

\section{Introduction}
This is a sample document demonstrating various \LaTeX{} features including mathematical equations, tables, and formatting. You can cite papers using the citation search feature, and they will be automatically added to your \texttt{cite.bib} file.

\section{Mathematical Expressions}
Here are some example equations:

\subsection{Quadratic Formula}
The quadratic formula is given by:
\[x = \frac{-b \pm \sqrt{b^2 - 4ac}}{2a}\]

\subsection{Matrix Example}
A sample matrix:
\[
\begin{pmatrix}
a_{11} & a_{12} & a_{13} \\
a_{21} & a_{22} & a_{23} \\
a_{31} & a_{32} & a_{33}
\end{pmatrix}
\]

\section{Tables}
Below is a sample table showing data:

\begin{table}[h]
\centering
\begin{tabular}{@{}lrrr@{}}
\toprule
\textbf{Category} & \textbf{Value 1} & \textbf{Value 2} & \textbf{Total} \\
\midrule
Item A & 10.5 & 20.3 & 30.8 \\
Item B & 15.2 & 25.1 & 40.3 \\
Item C & 22.1 & 12.7 & 34.8 \\
\bottomrule
\end{tabular}
\caption{Sample Data Table}
\label{tab:sample}
\end{table}

\section{Lists}
\subsection{Bullet Points}
\begin{itemize}
  \item First item
  \item Second item
  \item Third item with sub-items:
    \begin{itemize}
      \item Sub-item 1
      \item Sub-item 2
    \end{itemize}
\end{itemize}

\subsection{Numbered List}
\begin{enumerate}
  \item First step
  \item Second step
  \item Third step
\end{enumerate}

\section{Citations and References}
When you use the citation search feature, references will be automatically managed. For example, you might cite a paper like this: \cite{example2024}. All citations will be automatically added to your \texttt{cite.bib} file.

\section{Conclusion}
This template demonstrates basic \LaTeX{} features including mathematical formulas, tables, and lists. The bibliography system is now set up to work seamlessly with Bibby AI's citation management. Feel free to modify it for your needs!

\bibliographystyle{plain}
\bibliography{cite}

\end{document}\documentclass{article}
\usepackage[utf8]{inputenc}
\usepackage{booktabs}
\usepackage{amsmath}
\usepackage{natbib}

\title{My Document}
\author{Your Name}
\date{\today}

\begin{document}

\maketitle

\section{Introduction}
This is a sample document demonstrating various \LaTeX{} features including mathematical equations, tables, and formatting. You can cite papers using the citation search feature, and they will be automatically added to your \texttt{cite.bib} file.

\section{Mathematical Expressions}
Here are some example equations:

\subsection{Quadratic Formula}
The quadratic formula is given by:
\[x = \frac{-b \pm \sqrt{b^2 - 4ac}}{2a}\]

\subsection{Matrix Example}
A sample matrix:
\[
\begin{pmatrix}
a_{11} & a_{12} & a_{13} \\
a_{21} & a_{22} & a_{23} \\
a_{31} & a_{32} & a_{33}
\end{pmatrix}
\]

\section{Tables}
Below is a sample table showing data:

\begin{table}[h]
\centering
\begin{tabular}{@{}lrrr@{}}
\toprule
\textbf{Category} & \textbf{Value 1} & \textbf{Value 2} & \textbf{Total} \\
\midrule
Item A & 10.5 & 20.3 & 30.8 \\
Item B & 15.2 & 25.1 & 40.3 \\
Item C & 22.1 & 12.7 & 34.8 \\
\bottomrule
\end{tabular}
\caption{Sample Data Table}
\label{tab:sample}
\end{table}

\section{Lists}
\subsection{Bullet Points}
\begin{itemize}
  \item First item
  \item Second item
  \item Third item with sub-items:
    \begin{itemize}
      \item Sub-item 1
      \item Sub-item 2
    \end{itemize}
\end{itemize}

\subsection{Numbered List}
\begin{enumerate}
  \item First step
  \item Second step
  \item Third step
\end{enumerate}

\section{Citations and References}
When you use the citation search feature, references will be automatically managed. For example, you might cite a paper like this: \cite{example2024}. All citations will be automatically added to your \texttt{cite.bib} file.

\section{Conclusion}
This template demonstrates basic \LaTeX{} features including mathematical formulas, tables, and lists. The bibliography system is now set up to work seamlessly with Bibby AI's citation management. Feel free to modify it for your needs!

\bibliographystyle{plain}
\bibliography{cite}

\end{document}\documentclass{article}
\usepackage[utf8]{inputenc}
\usepackage{booktabs}
\usepackage{amsmath}
\usepackage{natbib}

\title{My Document}
\author{Your Name}
\date{\today}

\begin{document}

\maketitle

\section{Introduction}
This is a sample document demonstrating various \LaTeX{} features including mathematical equations, tables, and formatting. You can cite papers using the citation search feature, and they will be automatically added to your \texttt{cite.bib} file.

\section{Mathematical Expressions}
Here are some example equations:

\subsection{Quadratic Formula}
The quadratic formula is given by:
\[x = \frac{-b \pm \sqrt{b^2 - 4ac}}{2a}\]

\subsection{Matrix Example}
A sample matrix:
\[
\begin{pmatrix}
a_{11} & a_{12} & a_{13} \\
a_{21} & a_{22} & a_{23} \\
a_{31} & a_{32} & a_{33}
\end{pmatrix}
\]

\section{Tables}
Below is a sample table showing data:

\begin{table}[h]
\centering
\begin{tabular}{@{}lrrr@{}}
\toprule
\textbf{Category} & \textbf{Value 1} & \textbf{Value 2} & \textbf{Total} \\
\midrule
Item A & 10.5 & 20.3 & 30.8 \\
Item B & 15.2 & 25.1 & 40.3 \\
Item C & 22.1 & 12.7 & 34.8 \\
\bottomrule
\end{tabular}
\caption{Sample Data Table}
\label{tab:sample}
\end{table}

\section{Lists}
\subsection{Bullet Points}
\begin{itemize}
  \item First item
  \item Second item
  \item Third item with sub-items:
    \begin{itemize}
      \item Sub-item 1
      \item Sub-item 2
    \end{itemize}
\end{itemize}

\subsection{Numbered List}
\begin{enumerate}
  \item First step
  \item Second step
  \item Third step
\end{enumerate}

\section{Citations and References}
When you use the citation search feature, references will be automatically managed. For example, you might cite a paper like this: \cite{example2024}. All citations will be automatically added to your \texttt{cite.bib} file.

\section{Conclusion}
This template demonstrates basic \LaTeX{} features including mathematical formulas, tables, and lists. The bibliography system is now set up to work seamlessly with Bibby AI's citation management. Feel free to modify it for your needs!

\bibliographystyle{plain}
\bibliography{cite}

\end{document}\documentclass{article}
\usepackage[utf8]{inputenc}
\usepackage{booktabs}
\usepackage{amsmath}
\usepackage{natbib}

\title{My Document}
\author{Your Name}
\date{\today}

\begin{document}

\maketitle

\section{Introduction}
This is a sample document demonstrating various \LaTeX{} features including mathematical equations, tables, and formatting. You can cite papers using the citation search feature, and they will be automatically added to your \texttt{cite.bib} file.

\section{Mathematical Expressions}
Here are some example equations:

\subsection{Quadratic Formula}
The quadratic formula is given by:
\[x = \frac{-b \pm \sqrt{b^2 - 4ac}}{2a}\]

\subsection{Matrix Example}
A sample matrix:
\[
\begin{pmatrix}
a_{11} & a_{12} & a_{13} \\
a_{21} & a_{22} & a_{23} \\
a_{31} & a_{32} & a_{33}
\end{pmatrix}
\]

\section{Tables}
Below is a sample table showing data:

\begin{table}[h]
\centering
\begin{tabular}{@{}lrrr@{}}
\toprule
\textbf{Category} & \textbf{Value 1} & \textbf{Value 2} & \textbf{Total} \\
\midrule
Item A & 10.5 & 20.3 & 30.8 \\
Item B & 15.2 & 25.1 & 40.3 \\
Item C & 22.1 & 12.7 & 34.8 \\
\bottomrule
\end{tabular}
\caption{Sample Data Table}
\label{tab:sample}
\end{table}

\section{Lists}
\subsection{Bullet Points}
\begin{itemize}
  \item First item
  \item Second item
  \item Third item with sub-items:
    \begin{itemize}
      \item Sub-item 1
      \item Sub-item 2
    \end{itemize}
\end{itemize}

\subsection{Numbered List}
\begin{enumerate}
  \item First step
  \item Second step
  \item Third step
\end{enumerate}

\section{Citations and References}
When you use the citation search feature, references will be automatically managed. For example, you might cite a paper like this: \cite{example2024}. All citations will be automatically added to your \texttt{cite.bib} file.

\section{Conclusion}
This template demonstrates basic \LaTeX{} features including mathematical formulas, tables, and lists. The bibliography system is now set up to work seamlessly with Bibby AI's citation management. Feel free to modify it for your needs!

\bibliographystyle{plain}
\bibliography{cite}

\end{document}\documentclass{article}
\usepackage[utf8]{inputenc}
\usepackage{booktabs}
\usepackage{amsmath}
\usepackage{natbib}

\title{My Document}
\author{Your Name}
\date{\today}

\begin{document}

\maketitle

\section{Introduction}
This is a sample document demonstrating various \LaTeX{} features including mathematical equations, tables, and formatting. You can cite papers using the citation search feature, and they will be automatically added to your \texttt{cite.bib} file.

\section{Mathematical Expressions}
Here are some example equations:

\subsection{Quadratic Formula}
The quadratic formula is given by:
\[x = \frac{-b \pm \sqrt{b^2 - 4ac}}{2a}\]

\subsection{Matrix Example}
A sample matrix:
\[
\begin{pmatrix}
a_{11} & a_{12} & a_{13} \\
a_{21} & a_{22} & a_{23} \\
a_{31} & a_{32} & a_{33}
\end{pmatrix}
\]

\section{Tables}
Below is a sample table showing data:

\begin{table}[h]
\centering
\begin{tabular}{@{}lrrr@{}}
\toprule
\textbf{Category} & \textbf{Value 1} & \textbf{Value 2} & \textbf{Total} \\
\midrule
Item A & 10.5 & 20.3 & 30.8 \\
Item B & 15.2 & 25.1 & 40.3 \\
Item C & 22.1 & 12.7 & 34.8 \\
\bottomrule
\end{tabular}
\caption{Sample Data Table}
\label{tab:sample}
\end{table}

\section{Lists}
\subsection{Bullet Points}
\begin{itemize}
  \item First item
  \item Second item
  \item Third item with sub-items:
    \begin{itemize}
      \item Sub-item 1
      \item Sub-item 2
    \end{itemize}
\end{itemize}

\subsection{Numbered List}
\begin{enumerate}
  \item First step
  \item Second step
  \item Third step
\end{enumerate}

\section{Citations and References}
When you use the citation search feature, references will be automatically managed. For example, you might cite a paper like this: \cite{example2024}. All citations will be automatically added to your \texttt{cite.bib} file.

\section{Conclusion}
This template demonstrates basic \LaTeX{} features including mathematical formulas, tables, and lists. The bibliography system is now set up to work seamlessly with Bibby AI's citation management. Feel free to modify it for your needs!

\bibliographystyle{plain}
\bibliography{cite}

\end{document}\documentclass{article}
\usepackage[utf8]{inputenc}
\usepackage{booktabs}
\usepackage{amsmath}
\usepackage{natbib}

\title{My Document}
\author{Your Name}
\date{\today}

\begin{document}

\maketitle

\section{Introduction}
This is a sample document demonstrating various \LaTeX{} features including mathematical equations, tables, and formatting. You can cite papers using the citation search feature, and they will be automatically added to your \texttt{cite.bib} file.

\section{Mathematical Expressions}
Here are some example equations:

\subsection{Quadratic Formula}
The quadratic formula is given by:
\[x = \frac{-b \pm \sqrt{b^2 - 4ac}}{2a}\]

\subsection{Matrix Example}
A sample matrix:
\[
\begin{pmatrix}
a_{11} & a_{12} & a_{13} \\
a_{21} & a_{22} & a_{23} \\
a_{31} & a_{32} & a_{33}
\end{pmatrix}
\]

\section{Tables}
Below is a sample table showing data:

\begin{table}[h]
\centering
\begin{tabular}{@{}lrrr@{}}
\toprule
\textbf{Category} & \textbf{Value 1} & \textbf{Value 2} & \textbf{Total} \\
\midrule
Item A & 10.5 & 20.3 & 30.8 \\
Item B & 15.2 & 25.1 & 40.3 \\
Item C & 22.1 & 12.7 & 34.8 \\
\bottomrule
\end{tabular}
\caption{Sample Data Table}
\label{tab:sample}
\end{table}

\section{Lists}
\subsection{Bullet Points}
\begin{itemize}
  \item First item
  \item Second item
  \item Third item with sub-items:
    \begin{itemize}
      \item Sub-item 1
      \item Sub-item 2
    \end{itemize}
\end{itemize}

\subsection{Numbered List}
\begin{enumerate}
  \item First step
  \item Second step
  \item Third step
\end{enumerate}

\section{Citations and References}
When you use the citation search feature, references will be automatically managed. For example, you might cite a paper like this: \cite{example2024}. All citations will be automatically added to your \texttt{cite.bib} file.

\section{Conclusion}
This template demonstrates basic \LaTeX{} features including mathematical formulas, tables, and lists. The bibliography system is now set up to work seamlessly with Bibby AI's citation management. Feel free to modify it for your needs!

\bibliographystyle{plain}
\bibliography{cite}

\end{document}\documentclass{article}
\usepackage[utf8]{inputenc}
\usepackage{booktabs}
\usepackage{amsmath}
\usepackage{natbib}

\title{My Document}
\author{Your Name}
\date{\today}

\begin{document}

\maketitle

\section{Introduction}
This is a sample document demonstrating various \LaTeX{} features including mathematical equations, tables, and formatting. You can cite papers using the citation search feature, and they will be automatically added to your \texttt{cite.bib} file.

\section{Mathematical Expressions}
Here are some example equations:

\subsection{Quadratic Formula}
The quadratic formula is given by:
\[x = \frac{-b \pm \sqrt{b^2 - 4ac}}{2a}\]

\subsection{Matrix Example}
A sample matrix:
\[
\begin{pmatrix}
a_{11} & a_{12} & a_{13} \\
a_{21} & a_{22} & a_{23} \\
a_{31} & a_{32} & a_{33}
\end{pmatrix}
\]

\section{Tables}
Below is a sample table showing data:

\begin{table}[h]
\centering
\begin{tabular}{@{}lrrr@{}}
\toprule
\textbf{Category} & \textbf{Value 1} & \textbf{Value 2} & \textbf{Total} \\
\midrule
Item A & 10.5 & 20.3 & 30.8 \\
Item B & 15.2 & 25.1 & 40.3 \\
Item C & 22.1 & 12.7 & 34.8 \\
\bottomrule
\end{tabular}
\caption{Sample Data Table}
\label{tab:sample}
\end{table}

\section{Lists}
\subsection{Bullet Points}
\begin{itemize}
  \item First item
  \item Second item
  \item Third item with sub-items:
    \begin{itemize}
      \item Sub-item 1
      \item Sub-item 2
    \end{itemize}
\end{itemize}

\subsection{Numbered List}
\begin{enumerate}
  \item First step
  \item Second step
  \item Third step
\end{enumerate}

\section{Citations and References}
When you use the citation search feature, references will be automatically managed. For example, you might cite a paper like this: \cite{example2024}. All citations will be automatically added to your \texttt{cite.bib} file.

\section{Conclusion}
This template demonstrates basic \LaTeX{} features including mathematical formulas, tables, and lists. The bibliography system is now set up to work seamlessly with Bibby AI's citation management. Feel free to modify it for your needs!

\bibliographystyle{plain}
\bibliography{cite}

\end{document}\documentclass{article}
\usepackage[utf8]{inputenc}
\usepackage{booktabs}
\usepackage{amsmath}
\usepackage{natbib}

\title{My Document}
\author{Your Name}
\date{\today}

\begin{document}

\maketitle

\section{Introduction}
This is a sample document demonstrating various \LaTeX{} features including mathematical equations, tables, and formatting. You can cite papers using the citation search feature, and they will be automatically added to your \texttt{cite.bib} file.

\section{Mathematical Expressions}
Here are some example equations:

\subsection{Quadratic Formula}
The quadratic formula is given by:
\[x = \frac{-b \pm \sqrt{b^2 - 4ac}}{2a}\]

\subsection{Matrix Example}
A sample matrix:
\[
\begin{pmatrix}
a_{11} & a_{12} & a_{13} \\
a_{21} & a_{22} & a_{23} \\
a_{31} & a_{32} & a_{33}
\end{pmatrix}
\]

\section{Tables}
Below is a sample table showing data:

\begin{table}[h]
\centering
\begin{tabular}{@{}lrrr@{}}
\toprule
\textbf{Category} & \textbf{Value 1} & \textbf{Value 2} & \textbf{Total} \\
\midrule
Item A & 10.5 & 20.3 & 30.8 \\
Item B & 15.2 & 25.1 & 40.3 \\
Item C & 22.1 & 12.7 & 34.8 \\
\bottomrule
\end{tabular}
\caption{Sample Data Table}
\label{tab:sample}
\end{table}

\section{Lists}
\subsection{Bullet Points}
\begin{itemize}
  \item First item
  \item Second item
  \item Third item with sub-items:
    \begin{itemize}
      \item Sub-item 1
      \item Sub-item 2
    \end{itemize}
\end{itemize}

\subsection{Numbered List}
\begin{enumerate}
  \item First step
  \item Second step
  \item Third step
\end{enumerate}

\section{Citations and References}
When you use the citation search feature, references will be automatically managed. For example, you might cite a paper like this: \cite{example2024}. All citations will be automatically added to your \texttt{cite.bib} file.

\section{Conclusion}
This template demonstrates basic \LaTeX{} features including mathematical formulas, tables, and lists. The bibliography system is now set up to work seamlessly with Bibby AI's citation management. Feel free to modify it for your needs!

\bibliographystyle{plain}
\bibliography{cite}

\end{document}\documentclass{article}
\usepackage[utf8]{inputenc}
\usepackage{booktabs}
\usepackage{amsmath}
\usepackage{natbib}

\title{My Document}
\author{Your Name}
\date{\today}

\begin{document}

\maketitle

\section{Introduction}
This is a sample document demonstrating various \LaTeX{} features including mathematical equations, tables, and formatting. You can cite papers using the citation search feature, and they will be automatically added to your \texttt{cite.bib} file.

\section{Mathematical Expressions}
Here are some example equations:

\subsection{Quadratic Formula}
The quadratic formula is given by:
\[x = \frac{-b \pm \sqrt{b^2 - 4ac}}{2a}\]

\subsection{Matrix Example}
A sample matrix:
\[
\begin{pmatrix}
a_{11} & a_{12} & a_{13} \\
a_{21} & a_{22} & a_{23} \\
a_{31} & a_{32} & a_{33}
\end{pmatrix}
\]

\section{Tables}
Below is a sample table showing data:

\begin{table}[h]
\centering
\begin{tabular}{@{}lrrr@{}}
\toprule
\textbf{Category} & \textbf{Value 1} & \textbf{Value 2} & \textbf{Total} \\
\midrule
Item A & 10.5 & 20.3 & 30.8 \\
Item B & 15.2 & 25.1 & 40.3 \\
Item C & 22.1 & 12.7 & 34.8 \\
\bottomrule
\end{tabular}
\caption{Sample Data Table}
\label{tab:sample}
\end{table}

\section{Lists}
\subsection{Bullet Points}
\begin{itemize}
  \item First item
  \item Second item
  \item Third item with sub-items:
    \begin{itemize}
      \item Sub-item 1
      \item Sub-item 2
    \end{itemize}
\end{itemize}

\subsection{Numbered List}
\begin{enumerate}
  \item First step
  \item Second step
  \item Third step
\end{enumerate}

\section{Citations and References}
When you use the citation search feature, references will be automatically managed. For example, you might cite a paper like this: \cite{example2024}. All citations will be automatically added to your \texttt{cite.bib} file.

\section{Conclusion}
This template demonstrates basic \LaTeX{} features including mathematical formulas, tables, and lists. The bibliography system is now set up to work seamlessly with Bibby AI's citation management. Feel free to modify it for your needs!

\bibliographystyle{plain}
\bibliography{cite}

\end{document}\documentclass{article}
\usepackage[utf8]{inputenc}
\usepackage{booktabs}
\usepackage{amsmath}
\usepackage{natbib}

\title{My Document}
\author{Your Name}
\date{\today}

\begin{document}

\maketitle

\section{Introduction}
This is a sample document demonstrating various \LaTeX{} features including mathematical equations, tables, and formatting. You can cite papers using the citation search feature, and they will be automatically added to your \texttt{cite.bib} file.

\section{Mathematical Expressions}
Here are some example equations:

\subsection{Quadratic Formula}
The quadratic formula is given by:
\[x = \frac{-b \pm \sqrt{b^2 - 4ac}}{2a}\]

\subsection{Matrix Example}
A sample matrix:
\[
\begin{pmatrix}
a_{11} & a_{12} & a_{13} \\
a_{21} & a_{22} & a_{23} \\
a_{31} & a_{32} & a_{33}
\end{pmatrix}
\]

\section{Tables}
Below is a sample table showing data:

\begin{table}[h]
\centering
\begin{tabular}{@{}lrrr@{}}
\toprule
\textbf{Category} & \textbf{Value 1} & \textbf{Value 2} & \textbf{Total} \\
\midrule
Item A & 10.5 & 20.3 & 30.8 \\
Item B & 15.2 & 25.1 & 40.3 \\
Item C & 22.1 & 12.7 & 34.8 \\
\bottomrule
\end{tabular}
\caption{Sample Data Table}
\label{tab:sample}
\end{table}

\section{Lists}
\subsection{Bullet Points}
\begin{itemize}
  \item First item
  \item Second item
  \item Third item with sub-items:
    \begin{itemize}
      \item Sub-item 1
      \item Sub-item 2
    \end{itemize}
\end{itemize}

\subsection{Numbered List}
\begin{enumerate}
  \item First step
  \item Second step
  \item Third step
\end{enumerate}

\section{Citations and References}
When you use the citation search feature, references will be automatically managed. For example, you might cite a paper like this: \cite{example2024}. All citations will be automatically added to your \texttt{cite.bib} file.

\section{Conclusion}
This template demonstrates basic \LaTeX{} features including mathematical formulas, tables, and lists. The bibliography system is now set up to work seamlessly with Bibby AI's citation management. Feel free to modify it for your needs!

\bibliographystyle{plain}
\bibliography{cite}

\end{document}\documentclass{article}
\usepackage[utf8]{inputenc}
\usepackage{booktabs}
\usepackage{amsmath}
\usepackage{natbib}

\title{My Document}
\author{Your Name}
\date{\today}

\begin{document}

\maketitle

\section{Introduction}
This is a sample document demonstrating various \LaTeX{} features including mathematical equations, tables, and formatting. You can cite papers using the citation search feature, and they will be automatically added to your \texttt{cite.bib} file.

\section{Mathematical Expressions}
Here are some example equations:

\subsection{Quadratic Formula}
The quadratic formula is given by:
\[x = \frac{-b \pm \sqrt{b^2 - 4ac}}{2a}\]

\subsection{Matrix Example}
A sample matrix:
\[
\begin{pmatrix}
a_{11} & a_{12} & a_{13} \\
a_{21} & a_{22} & a_{23} \\
a_{31} & a_{32} & a_{33}
\end{pmatrix}
\]

\section{Tables}
Below is a sample table showing data:

\begin{table}[h]
\centering
\begin{tabular}{@{}lrrr@{}}
\toprule
\textbf{Category} & \textbf{Value 1} & \textbf{Value 2} & \textbf{Total} \\
\midrule
Item A & 10.5 & 20.3 & 30.8 \\
Item B & 15.2 & 25.1 & 40.3 \\
Item C & 22.1 & 12.7 & 34.8 \\
\bottomrule
\end{tabular}
\caption{Sample Data Table}
\label{tab:sample}
\end{table}

\section{Lists}
\subsection{Bullet Points}
\begin{itemize}
  \item First item
  \item Second item
  \item Third item with sub-items:
    \begin{itemize}
      \item Sub-item 1
      \item Sub-item 2
    \end{itemize}
\end{itemize}

\subsection{Numbered List}
\begin{enumerate}
  \item First step
  \item Second step
  \item Third step
\end{enumerate}

\section{Citations and References}
When you use the citation search feature, references will be automatically managed. For example, you might cite a paper like this: \cite{example2024}. All citations will be automatically added to your \texttt{cite.bib} file.

\section{Conclusion}
This template demonstrates basic \LaTeX{} features including mathematical formulas, tables, and lists. The bibliography system is now set up to work seamlessly with Bibby AI's citation management. Feel free to modify it for your needs!

\bibliographystyle{plain}
\bibliography{cite}

\end{document}\documentclass{article}
\usepackage[utf8]{inputenc}
\usepackage{booktabs}
\usepackage{amsmath}
\usepackage{natbib}

\title{My Document}
\author{Your Name}
\date{\today}

\begin{document}

\maketitle

\section{Introduction}
This is a sample document demonstrating various \LaTeX{} features including mathematical equations, tables, and formatting. You can cite papers using the citation search feature, and they will be automatically added to your \texttt{cite.bib} file.

\section{Mathematical Expressions}
Here are some example equations:

\subsection{Quadratic Formula}
The quadratic formula is given by:
\[x = \frac{-b \pm \sqrt{b^2 - 4ac}}{2a}\]

\subsection{Matrix Example}
A sample matrix:
\[
\begin{pmatrix}
a_{11} & a_{12} & a_{13} \\
a_{21} & a_{22} & a_{23} \\
a_{31} & a_{32} & a_{33}
\end{pmatrix}
\]

\section{Tables}
Below is a sample table showing data:

\begin{table}[h]
\centering
\begin{tabular}{@{}lrrr@{}}
\toprule
\textbf{Category} & \textbf{Value 1} & \textbf{Value 2} & \textbf{Total} \\
\midrule
Item A & 10.5 & 20.3 & 30.8 \\
Item B & 15.2 & 25.1 & 40.3 \\
Item C & 22.1 & 12.7 & 34.8 \\
\bottomrule
\end{tabular}
\caption{Sample Data Table}
\label{tab:sample}
\end{table}

\section{Lists}
\subsection{Bullet Points}
\begin{itemize}
  \item First item
  \item Second item
  \item Third item with sub-items:
    \begin{itemize}
      \item Sub-item 1
      \item Sub-item 2
    \end{itemize}
\end{itemize}

\subsection{Numbered List}
\begin{enumerate}
  \item First step
  \item Second step
  \item Third step
\end{enumerate}

\section{Citations and References}
When you use the citation search feature, references will be automatically managed. For example, you might cite a paper like this: \cite{example2024}. All citations will be automatically added to your \texttt{cite.bib} file.

\section{Conclusion}
This template demonstrates basic \LaTeX{} features including mathematical formulas, tables, and lists. The bibliography system is now set up to work seamlessly with Bibby AI's citation management. Feel free to modify it for your needs!

\bibliographystyle{plain}
\bibliography{cite}

\end{document}\documentclass{article}
\usepackage[utf8]{inputenc}
\usepackage{booktabs}
\usepackage{amsmath}
\usepackage{natbib}

\title{My Document}
\author{Your Name}
\date{\today}

\begin{document}

\maketitle

\section{Introduction}
This is a sample document demonstrating various \LaTeX{} features including mathematical equations, tables, and formatting. You can cite papers using the citation search feature, and they will be automatically added to your \texttt{cite.bib} file.

\section{Mathematical Expressions}
Here are some example equations:

\subsection{Quadratic Formula}
The quadratic formula is given by:
\[x = \frac{-b \pm \sqrt{b^2 - 4ac}}{2a}\]

\subsection{Matrix Example}
A sample matrix:
\[
\begin{pmatrix}
a_{11} & a_{12} & a_{13} \\
a_{21} & a_{22} & a_{23} \\
a_{31} & a_{32} & a_{33}
\end{pmatrix}
\]

\section{Tables}
Below is a sample table showing data:

\begin{table}[h]
\centering
\begin{tabular}{@{}lrrr@{}}
\toprule
\textbf{Category} & \textbf{Value 1} & \textbf{Value 2} & \textbf{Total} \\
\midrule
Item A & 10.5 & 20.3 & 30.8 \\
Item B & 15.2 & 25.1 & 40.3 \\
Item C & 22.1 & 12.7 & 34.8 \\
\bottomrule
\end{tabular}
\caption{Sample Data Table}
\label{tab:sample}
\end{table}

\section{Lists}
\subsection{Bullet Points}
\begin{itemize}
  \item First item
  \item Second item
  \item Third item with sub-items:
    \begin{itemize}
      \item Sub-item 1
      \item Sub-item 2
    \end{itemize}
\end{itemize}

\subsection{Numbered List}
\begin{enumerate}
  \item First step
  \item Second step
  \item Third step
\end{enumerate}

\section{Citations and References}
When you use the citation search feature, references will be automatically managed. For example, you might cite a paper like this: \cite{example2024}. All citations will be automatically added to your \texttt{cite.bib} file.

\section{Conclusion}
This template demonstrates basic \LaTeX{} features including mathematical formulas, tables, and lists. The bibliography system is now set up to work seamlessly with Bibby AI's citation management. Feel free to modify it for your needs!

\bibliographystyle{plain}
\bibliography{cite}

\end{document}\documentclass{article}
\usepackage[utf8]{inputenc}
\usepackage{booktabs}
\usepackage{amsmath}
\usepackage{natbib}

\title{My Document}
\author{Your Name}
\date{\today}

\begin{document}

\maketitle

\section{Introduction}
This is a sample document demonstrating various \LaTeX{} features including mathematical equations, tables, and formatting. You can cite papers using the citation search feature, and they will be automatically added to your \texttt{cite.bib} file.

\section{Mathematical Expressions}
Here are some example equations:

\subsection{Quadratic Formula}
The quadratic formula is given by:
\[x = \frac{-b \pm \sqrt{b^2 - 4ac}}{2a}\]

\subsection{Matrix Example}
A sample matrix:
\[
\begin{pmatrix}
a_{11} & a_{12} & a_{13} \\
a_{21} & a_{22} & a_{23} \\
a_{31} & a_{32} & a_{33}
\end{pmatrix}
\]

\section{Tables}
Below is a sample table showing data:

\begin{table}[h]
\centering
\begin{tabular}{@{}lrrr@{}}
\toprule
\textbf{Category} & \textbf{Value 1} & \textbf{Value 2} & \textbf{Total} \\
\midrule
Item A & 10.5 & 20.3 & 30.8 \\
Item B & 15.2 & 25.1 & 40.3 \\
Item C & 22.1 & 12.7 & 34.8 \\
\bottomrule
\end{tabular}
\caption{Sample Data Table}
\label{tab:sample}
\end{table}

\section{Lists}
\subsection{Bullet Points}
\begin{itemize}
  \item First item
  \item Second item
  \item Third item with sub-items:
    \begin{itemize}
      \item Sub-item 1
      \item Sub-item 2
    \end{itemize}
\end{itemize}

\subsection{Numbered List}
\begin{enumerate}
  \item First step
  \item Second step
  \item Third step
\end{enumerate}

\section{Citations and References}
When you use the citation search feature, references will be automatically managed. For example, you might cite a paper like this: \cite{example2024}. All citations will be automatically added to your \texttt{cite.bib} file.

\section{Conclusion}
This template demonstrates basic \LaTeX{} features including mathematical formulas, tables, and lists. The bibliography system is now set up to work seamlessly with Bibby AI's citation management. Feel free to modify it for your needs!

\bibliographystyle{plain}
\bibliography{cite}

\end{document}\documentclass{article}
\usepackage[utf8]{inputenc}
\usepackage{booktabs}
\usepackage{amsmath}
\usepackage{natbib}

\title{My Document}
\author{Your Name}
\date{\today}

\begin{document}

\maketitle

\section{Introduction}
This is a sample document demonstrating various \LaTeX{} features including mathematical equations, tables, and formatting. You can cite papers using the citation search feature, and they will be automatically added to your \texttt{cite.bib} file.

\section{Mathematical Expressions}
Here are some example equations:

\subsection{Quadratic Formula}
The quadratic formula is given by:
\[x = \frac{-b \pm \sqrt{b^2 - 4ac}}{2a}\]

\subsection{Matrix Example}
A sample matrix:
\[
\begin{pmatrix}
a_{11} & a_{12} & a_{13} \\
a_{21} & a_{22} & a_{23} \\
a_{31} & a_{32} & a_{33}
\end{pmatrix}
\]

\section{Tables}
Below is a sample table showing data:

\begin{table}[h]
\centering
\begin{tabular}{@{}lrrr@{}}
\toprule
\textbf{Category} & \textbf{Value 1} & \textbf{Value 2} & \textbf{Total} \\
\midrule
Item A & 10.5 & 20.3 & 30.8 \\
Item B & 15.2 & 25.1 & 40.3 \\
Item C & 22.1 & 12.7 & 34.8 \\
\bottomrule
\end{tabular}
\caption{Sample Data Table}
\label{tab:sample}
\end{table}

\section{Lists}
\subsection{Bullet Points}
\begin{itemize}
  \item First item
  \item Second item
  \item Third item with sub-items:
    \begin{itemize}
      \item Sub-item 1
      \item Sub-item 2
    \end{itemize}
\end{itemize}

\subsection{Numbered List}
\begin{enumerate}
  \item First step
  \item Second step
  \item Third step
\end{enumerate}

\section{Citations and References}
When you use the citation search feature, references will be automatically managed. For example, you might cite a paper like this: \cite{example2024}. All citations will be automatically added to your \texttt{cite.bib} file.

\section{Conclusion}
This template demonstrates basic \LaTeX{} features including mathematical formulas, tables, and lists. The bibliography system is now set up to work seamlessly with Bibby AI's citation management. Feel free to modify it for your needs!

\bibliographystyle{plain}
\bibliography{cite}

\end{document}\documentclass{article}
\usepackage[utf8]{inputenc}
\usepackage{booktabs}
\usepackage{amsmath}
\usepackage{natbib}

\title{My Document}
\author{Your Name}
\date{\today}

\begin{document}

\maketitle

\section{Introduction}
This is a sample document demonstrating various \LaTeX{} features including mathematical equations, tables, and formatting. You can cite papers using the citation search feature, and they will be automatically added to your \texttt{cite.bib} file.

\section{Mathematical Expressions}
Here are some example equations:

\subsection{Quadratic Formula}
The quadratic formula is given by:
\[x = \frac{-b \pm \sqrt{b^2 - 4ac}}{2a}\]

\subsection{Matrix Example}
A sample matrix:
\[
\begin{pmatrix}
a_{11} & a_{12} & a_{13} \\
a_{21} & a_{22} & a_{23} \\
a_{31} & a_{32} & a_{33}
\end{pmatrix}
\]

\section{Tables}
Below is a sample table showing data:

\begin{table}[h]
\centering
\begin{tabular}{@{}lrrr@{}}
\toprule
\textbf{Category} & \textbf{Value 1} & \textbf{Value 2} & \textbf{Total} \\
\midrule
Item A & 10.5 & 20.3 & 30.8 \\
Item B & 15.2 & 25.1 & 40.3 \\
Item C & 22.1 & 12.7 & 34.8 \\
\bottomrule
\end{tabular}
\caption{Sample Data Table}
\label{tab:sample}
\end{table}

\section{Lists}
\subsection{Bullet Points}
\begin{itemize}
  \item First item
  \item Second item
  \item Third item with sub-items:
    \begin{itemize}
      \item Sub-item 1
      \item Sub-item 2
    \end{itemize}
\end{itemize}

\subsection{Numbered List}
\begin{enumerate}
  \item First step
  \item Second step
  \item Third step
\end{enumerate}

\section{Citations and References}
When you use the citation search feature, references will be automatically managed. For example, you might cite a paper like this: \cite{example2024}. All citations will be automatically added to your \texttt{cite.bib} file.

\section{Conclusion}
This template demonstrates basic \LaTeX{} features including mathematical formulas, tables, and lists. The bibliography system is now set up to work seamlessly with Bibby AI's citation management. Feel free to modify it for your needs!

\bibliographystyle{plain}
\bibliography{cite}

\end{document}\documentclass{article}
\usepackage[utf8]{inputenc}
\usepackage{booktabs}
\usepackage{amsmath}
\usepackage{natbib}

\title{My Document}
\author{Your Name}
\date{\today}

\begin{document}

\maketitle

\section{Introduction}
This is a sample document demonstrating various \LaTeX{} features including mathematical equations, tables, and formatting. You can cite papers using the citation search feature, and they will be automatically added to your \texttt{cite.bib} file.

\section{Mathematical Expressions}
Here are some example equations:

\subsection{Quadratic Formula}
The quadratic formula is given by:
\[x = \frac{-b \pm \sqrt{b^2 - 4ac}}{2a}\]

\subsection{Matrix Example}
A sample matrix:
\[
\begin{pmatrix}
a_{11} & a_{12} & a_{13} \\
a_{21} & a_{22} & a_{23} \\
a_{31} & a_{32} & a_{33}
\end{pmatrix}
\]

\section{Tables}
Below is a sample table showing data:

\begin{table}[h]
\centering
\begin{tabular}{@{}lrrr@{}}
\toprule
\textbf{Category} & \textbf{Value 1} & \textbf{Value 2} & \textbf{Total} \\
\midrule
Item A & 10.5 & 20.3 & 30.8 \\
Item B & 15.2 & 25.1 & 40.3 \\
Item C & 22.1 & 12.7 & 34.8 \\
\bottomrule
\end{tabular}
\caption{Sample Data Table}
\label{tab:sample}
\end{table}

\section{Lists}
\subsection{Bullet Points}
\begin{itemize}
  \item First item
  \item Second item
  \item Third item with sub-items:
    \begin{itemize}
      \item Sub-item 1
      \item Sub-item 2
    \end{itemize}
\end{itemize}

\subsection{Numbered List}
\begin{enumerate}
  \item First step
  \item Second step
  \item Third step
\end{enumerate}

\section{Citations and References}
When you use the citation search feature, references will be automatically managed. For example, you might cite a paper like this: \cite{example2024}. All citations will be automatically added to your \texttt{cite.bib} file.

\section{Conclusion}
This template demonstrates basic \LaTeX{} features including mathematical formulas, tables, and lists. The bibliography system is now set up to work seamlessly with Bibby AI's citation management. Feel free to modify it for your needs!

\bibliographystyle{plain}
\bibliography{cite}

\end{document}\documentclass{article}
\usepackage[utf8]{inputenc}
\usepackage{booktabs}
\usepackage{amsmath}
\usepackage{natbib}

\title{My Document}
\author{Your Name}
\date{\today}

\begin{document}

\maketitle

\section{Introduction}
This is a sample document demonstrating various \LaTeX{} features including mathematical equations, tables, and formatting. You can cite papers using the citation search feature, and they will be automatically added to your \texttt{cite.bib} file.

\section{Mathematical Expressions}
Here are some example equations:

\subsection{Quadratic Formula}
The quadratic formula is given by:
\[x = \frac{-b \pm \sqrt{b^2 - 4ac}}{2a}\]

\subsection{Matrix Example}
A sample matrix:
\[
\begin{pmatrix}
a_{11} & a_{12} & a_{13} \\
a_{21} & a_{22} & a_{23} \\
a_{31} & a_{32} & a_{33}
\end{pmatrix}
\]

\section{Tables}
Below is a sample table showing data:

\begin{table}[h]
\centering
\begin{tabular}{@{}lrrr@{}}
\toprule
\textbf{Category} & \textbf{Value 1} & \textbf{Value 2} & \textbf{Total} \\
\midrule
Item A & 10.5 & 20.3 & 30.8 \\
Item B & 15.2 & 25.1 & 40.3 \\
Item C & 22.1 & 12.7 & 34.8 \\
\bottomrule
\end{tabular}
\caption{Sample Data Table}
\label{tab:sample}
\end{table}

\section{Lists}
\subsection{Bullet Points}
\begin{itemize}
  \item First item
  \item Second item
  \item Third item with sub-items:
    \begin{itemize}
      \item Sub-item 1
      \item Sub-item 2
    \end{itemize}
\end{itemize}

\subsection{Numbered List}
\begin{enumerate}
  \item First step
  \item Second step
  \item Third step
\end{enumerate}

\section{Citations and References}
When you use the citation search feature, references will be automatically managed. For example, you might cite a paper like this: \cite{example2024}. All citations will be automatically added to your \texttt{cite.bib} file.

\section{Conclusion}
This template demonstrates basic \LaTeX{} features including mathematical formulas, tables, and lists. The bibliography system is now set up to work seamlessly with Bibby AI's citation management. Feel free to modify it for your needs!

\bibliographystyle{plain}
\bibliography{cite}

\end{document}\documentclass{article}
\usepackage[utf8]{inputenc}
\usepackage{booktabs}
\usepackage{amsmath}
\usepackage{natbib}

\title{My Document}
\author{Your Name}
\date{\today}

\begin{document}

\maketitle

\section{Introduction}
This is a sample document demonstrating various \LaTeX{} features including mathematical equations, tables, and formatting. You can cite papers using the citation search feature, and they will be automatically added to your \texttt{cite.bib} file.

\section{Mathematical Expressions}
Here are some example equations:

\subsection{Quadratic Formula}
The quadratic formula is given by:
\[x = \frac{-b \pm \sqrt{b^2 - 4ac}}{2a}\]

\subsection{Matrix Example}
A sample matrix:
\[
\begin{pmatrix}
a_{11} & a_{12} & a_{13} \\
a_{21} & a_{22} & a_{23} \\
a_{31} & a_{32} & a_{33}
\end{pmatrix}
\]

\section{Tables}
Below is a sample table showing data:

\begin{table}[h]
\centering
\begin{tabular}{@{}lrrr@{}}
\toprule
\textbf{Category} & \textbf{Value 1} & \textbf{Value 2} & \textbf{Total} \\
\midrule
Item A & 10.5 & 20.3 & 30.8 \\
Item B & 15.2 & 25.1 & 40.3 \\
Item C & 22.1 & 12.7 & 34.8 \\
\bottomrule
\end{tabular}
\caption{Sample Data Table}
\label{tab:sample}
\end{table}

\section{Lists}
\subsection{Bullet Points}
\begin{itemize}
  \item First item
  \item Second item
  \item Third item with sub-items:
    \begin{itemize}
      \item Sub-item 1
      \item Sub-item 2
    \end{itemize}
\end{itemize}

\subsection{Numbered List}
\begin{enumerate}
  \item First step
  \item Second step
  \item Third step
\end{enumerate}

\section{Citations and References}
When you use the citation search feature, references will be automatically managed. For example, you might cite a paper like this: \cite{example2024}. All citations will be automatically added to your \texttt{cite.bib} file.

\section{Conclusion}
This template demonstrates basic \LaTeX{} features including mathematical formulas, tables, and lists. The bibliography system is now set up to work seamlessly with Bibby AI's citation management. Feel free to modify it for your needs!

\bibliographystyle{plain}
\bibliography{cite}

\end{document}\documentclass{article}
\usepackage[utf8]{inputenc}
\usepackage{booktabs}
\usepackage{amsmath}
\usepackage{natbib}

\title{My Document}
\author{Your Name}
\date{\today}

\begin{document}

\maketitle

\section{Introduction}
This is a sample document demonstrating various \LaTeX{} features including mathematical equations, tables, and formatting. You can cite papers using the citation search feature, and they will be automatically added to your \texttt{cite.bib} file.

\section{Mathematical Expressions}
Here are some example equations:

\subsection{Quadratic Formula}
The quadratic formula is given by:
\[x = \frac{-b \pm \sqrt{b^2 - 4ac}}{2a}\]

\subsection{Matrix Example}
A sample matrix:
\[
\begin{pmatrix}
a_{11} & a_{12} & a_{13} \\
a_{21} & a_{22} & a_{23} \\
a_{31} & a_{32} & a_{33}
\end{pmatrix}
\]

\section{Tables}
Below is a sample table showing data:

\begin{table}[h]
\centering
\begin{tabular}{@{}lrrr@{}}
\toprule
\textbf{Category} & \textbf{Value 1} & \textbf{Value 2} & \textbf{Total} \\
\midrule
Item A & 10.5 & 20.3 & 30.8 \\
Item B & 15.2 & 25.1 & 40.3 \\
Item C & 22.1 & 12.7 & 34.8 \\
\bottomrule
\end{tabular}
\caption{Sample Data Table}
\label{tab:sample}
\end{table}

\section{Lists}
\subsection{Bullet Points}
\begin{itemize}
  \item First item
  \item Second item
  \item Third item with sub-items:
    \begin{itemize}
      \item Sub-item 1
      \item Sub-item 2
    \end{itemize}
\end{itemize}

\subsection{Numbered List}
\begin{enumerate}
  \item First step
  \item Second step
  \item Third step
\end{enumerate}

\section{Citations and References}
When you use the citation search feature, references will be automatically managed. For example, you might cite a paper like this: \cite{example2024}. All citations will be automatically added to your \texttt{cite.bib} file.

\section{Conclusion}
This template demonstrates basic \LaTeX{} features including mathematical formulas, tables, and lists. The bibliography system is now set up to work seamlessly with Bibby AI's citation management. Feel free to modify it for your needs!

\bibliographystyle{plain}
\bibliography{cite}

\end{document}\documentclass{article}
\usepackage[utf8]{inputenc}
\usepackage{booktabs}
\usepackage{amsmath}
\usepackage{natbib}

\title{My Document}
\author{Your Name}
\date{\today}

\begin{document}

\maketitle

\section{Introduction}
This is a sample document demonstrating various \LaTeX{} features including mathematical equations, tables, and formatting. You can cite papers using the citation search feature, and they will be automatically added to your \texttt{cite.bib} file.

\section{Mathematical Expressions}
Here are some example equations:

\subsection{Quadratic Formula}
The quadratic formula is given by:
\[x = \frac{-b \pm \sqrt{b^2 - 4ac}}{2a}\]

\subsection{Matrix Example}
A sample matrix:
\[
\begin{pmatrix}
a_{11} & a_{12} & a_{13} \\
a_{21} & a_{22} & a_{23} \\
a_{31} & a_{32} & a_{33}
\end{pmatrix}
\]

\section{Tables}
Below is a sample table showing data:

\begin{table}[h]
\centering
\begin{tabular}{@{}lrrr@{}}
\toprule
\textbf{Category} & \textbf{Value 1} & \textbf{Value 2} & \textbf{Total} \\
\midrule
Item A & 10.5 & 20.3 & 30.8 \\
Item B & 15.2 & 25.1 & 40.3 \\
Item C & 22.1 & 12.7 & 34.8 \\
\bottomrule
\end{tabular}
\caption{Sample Data Table}
\label{tab:sample}
\end{table}

\section{Lists}
\subsection{Bullet Points}
\begin{itemize}
  \item First item
  \item Second item
  \item Third item with sub-items:
    \begin{itemize}
      \item Sub-item 1
      \item Sub-item 2
    \end{itemize}
\end{itemize}

\subsection{Numbered List}
\begin{enumerate}
  \item First step
  \item Second step
  \item Third step
\end{enumerate}

\section{Citations and References}
When you use the citation search feature, references will be automatically managed. For example, you might cite a paper like this: \cite{example2024}. All citations will be automatically added to your \texttt{cite.bib} file.

\section{Conclusion}
This template demonstrates basic \LaTeX{} features including mathematical formulas, tables, and lists. The bibliography system is now set up to work seamlessly with Bibby AI's citation management. Feel free to modify it for your needs!

\bibliographystyle{plain}
\bibliography{cite}

\end{document}\documentclass{article}
\usepackage[utf8]{inputenc}
\usepackage{booktabs}
\usepackage{amsmath}
\usepackage{natbib}

\title{My Document}
\author{Your Name}
\date{\today}

\begin{document}

\maketitle

\section{Introduction}
This is a sample document demonstrating various \LaTeX{} features including mathematical equations, tables, and formatting. You can cite papers using the citation search feature, and they will be automatically added to your \texttt{cite.bib} file.

\section{Mathematical Expressions}
Here are some example equations:

\subsection{Quadratic Formula}
The quadratic formula is given by:
\[x = \frac{-b \pm \sqrt{b^2 - 4ac}}{2a}\]

\subsection{Matrix Example}
A sample matrix:
\[
\begin{pmatrix}
a_{11} & a_{12} & a_{13} \\
a_{21} & a_{22} & a_{23} \\
a_{31} & a_{32} & a_{33}
\end{pmatrix}
\]

\section{Tables}
Below is a sample table showing data:

\begin{table}[h]
\centering
\begin{tabular}{@{}lrrr@{}}
\toprule
\textbf{Category} & \textbf{Value 1} & \textbf{Value 2} & \textbf{Total} \\
\midrule
Item A & 10.5 & 20.3 & 30.8 \\
Item B & 15.2 & 25.1 & 40.3 \\
Item C & 22.1 & 12.7 & 34.8 \\
\bottomrule
\end{tabular}
\caption{Sample Data Table}
\label{tab:sample}
\end{table}

\section{Lists}
\subsection{Bullet Points}
\begin{itemize}
  \item First item
  \item Second item
  \item Third item with sub-items:
    \begin{itemize}
      \item Sub-item 1
      \item Sub-item 2
    \end{itemize}
\end{itemize}

\subsection{Numbered List}
\begin{enumerate}
  \item First step
  \item Second step
  \item Third step
\end{enumerate}

\section{Citations and References}
When you use the citation search feature, references will be automatically managed. For example, you might cite a paper like this: \cite{example2024}. All citations will be automatically added to your \texttt{cite.bib} file.

\section{Conclusion}
This template demonstrates basic \LaTeX{} features including mathematical formulas, tables, and lists. The bibliography system is now set up to work seamlessly with Bibby AI's citation management. Feel free to modify it for your needs!

\bibliographystyle{plain}
\bibliography{cite}

\end{document}\documentclass{article}
\usepackage[utf8]{inputenc}
\usepackage{booktabs}
\usepackage{amsmath}
\usepackage{natbib}

\title{My Document}
\author{Your Name}
\date{\today}

\begin{document}

\maketitle

\section{Introduction}
This is a sample document demonstrating various \LaTeX{} features including mathematical equations, tables, and formatting. You can cite papers using the citation search feature, and they will be automatically added to your \texttt{cite.bib} file.

\section{Mathematical Expressions}
Here are some example equations:

\subsection{Quadratic Formula}
The quadratic formula is given by:
\[x = \frac{-b \pm \sqrt{b^2 - 4ac}}{2a}\]

\subsection{Matrix Example}
A sample matrix:
\[
\begin{pmatrix}
a_{11} & a_{12} & a_{13} \\
a_{21} & a_{22} & a_{23} \\
a_{31} & a_{32} & a_{33}
\end{pmatrix}
\]

\section{Tables}
Below is a sample table showing data:

\begin{table}[h]
\centering
\begin{tabular}{@{}lrrr@{}}
\toprule
\textbf{Category} & \textbf{Value 1} & \textbf{Value 2} & \textbf{Total} \\
\midrule
Item A & 10.5 & 20.3 & 30.8 \\
Item B & 15.2 & 25.1 & 40.3 \\
Item C & 22.1 & 12.7 & 34.8 \\
\bottomrule
\end{tabular}
\caption{Sample Data Table}
\label{tab:sample}
\end{table}

\section{Lists}
\subsection{Bullet Points}
\begin{itemize}
  \item First item
  \item Second item
  \item Third item with sub-items:
    \begin{itemize}
      \item Sub-item 1
      \item Sub-item 2
    \end{itemize}
\end{itemize}

\subsection{Numbered List}
\begin{enumerate}
  \item First step
  \item Second step
  \item Third step
\end{enumerate}

\section{Citations and References}
When you use the citation search feature, references will be automatically managed. For example, you might cite a paper like this: \cite{example2024}. All citations will be automatically added to your \texttt{cite.bib} file.

\section{Conclusion}
This template demonstrates basic \LaTeX{} features including mathematical formulas, tables, and lists. The bibliography system is now set up to work seamlessly with Bibby AI's citation management. Feel free to modify it for your needs!

\bibliographystyle{plain}
\bibliography{cite}

\end{document}\documentclass{article}
\usepackage[utf8]{inputenc}
\usepackage{booktabs}
\usepackage{amsmath}
\usepackage{natbib}

\title{My Document}
\author{Your Name}
\date{\today}

\begin{document}

\maketitle

\section{Introduction}
This is a sample document demonstrating various \LaTeX{} features including mathematical equations, tables, and formatting. You can cite papers using the citation search feature, and they will be automatically added to your \texttt{cite.bib} file.

\section{Mathematical Expressions}
Here are some example equations:

\subsection{Quadratic Formula}
The quadratic formula is given by:
\[x = \frac{-b \pm \sqrt{b^2 - 4ac}}{2a}\]

\subsection{Matrix Example}
A sample matrix:
\[
\begin{pmatrix}
a_{11} & a_{12} & a_{13} \\
a_{21} & a_{22} & a_{23} \\
a_{31} & a_{32} & a_{33}
\end{pmatrix}
\]

\section{Tables}
Below is a sample table showing data:

\begin{table}[h]
\centering
\begin{tabular}{@{}lrrr@{}}
\toprule
\textbf{Category} & \textbf{Value 1} & \textbf{Value 2} & \textbf{Total} \\
\midrule
Item A & 10.5 & 20.3 & 30.8 \\
Item B & 15.2 & 25.1 & 40.3 \\
Item C & 22.1 & 12.7 & 34.8 \\
\bottomrule
\end{tabular}
\caption{Sample Data Table}
\label{tab:sample}
\end{table}

\section{Lists}
\subsection{Bullet Points}
\begin{itemize}
  \item First item
  \item Second item
  \item Third item with sub-items:
    \begin{itemize}
      \item Sub-item 1
      \item Sub-item 2
    \end{itemize}
\end{itemize}

\subsection{Numbered List}
\begin{enumerate}
  \item First step
  \item Second step
  \item Third step
\end{enumerate}

\section{Citations and References}
When you use the citation search feature, references will be automatically managed. For example, you might cite a paper like this: \cite{example2024}. All citations will be automatically added to your \texttt{cite.bib} file.

\section{Conclusion}
This template demonstrates basic \LaTeX{} features including mathematical formulas, tables, and lists. The bibliography system is now set up to work seamlessly with Bibby AI's citation management. Feel free to modify it for your needs!

\bibliographystyle{plain}
\bibliography{cite}

\end{document}\documentclass{article}
\usepackage[utf8]{inputenc}
\usepackage{booktabs}
\usepackage{amsmath}
\usepackage{natbib}

\title{My Document}
\author{Your Name}
\date{\today}

\begin{document}

\maketitle

\section{Introduction}
This is a sample document demonstrating various \LaTeX{} features including mathematical equations, tables, and formatting. You can cite papers using the citation search feature, and they will be automatically added to your \texttt{cite.bib} file.

\section{Mathematical Expressions}
Here are some example equations:

\subsection{Quadratic Formula}
The quadratic formula is given by:
\[x = \frac{-b \pm \sqrt{b^2 - 4ac}}{2a}\]

\subsection{Matrix Example}
A sample matrix:
\[
\begin{pmatrix}
a_{11} & a_{12} & a_{13} \\
a_{21} & a_{22} & a_{23} \\
a_{31} & a_{32} & a_{33}
\end{pmatrix}
\]

\section{Tables}
Below is a sample table showing data:

\begin{table}[h]
\centering
\begin{tabular}{@{}lrrr@{}}
\toprule
\textbf{Category} & \textbf{Value 1} & \textbf{Value 2} & \textbf{Total} \\
\midrule
Item A & 10.5 & 20.3 & 30.8 \\
Item B & 15.2 & 25.1 & 40.3 \\
Item C & 22.1 & 12.7 & 34.8 \\
\bottomrule
\end{tabular}
\caption{Sample Data Table}
\label{tab:sample}
\end{table}

\section{Lists}
\subsection{Bullet Points}
\begin{itemize}
  \item First item
  \item Second item
  \item Third item with sub-items:
    \begin{itemize}
      \item Sub-item 1
      \item Sub-item 2
    \end{itemize}
\end{itemize}

\subsection{Numbered List}
\begin{enumerate}
  \item First step
  \item Second step
  \item Third step
\end{enumerate}

\section{Citations and References}
When you use the citation search feature, references will be automatically managed. For example, you might cite a paper like this: \cite{example2024}. All citations will be automatically added to your \texttt{cite.bib} file.

\section{Conclusion}
This template demonstrates basic \LaTeX{} features including mathematical formulas, tables, and lists. The bibliography system is now set up to work seamlessly with Bibby AI's citation management. Feel free to modify it for your needs!

\bibliographystyle{plain}
\bibliography{cite}

\end{document}\documentclass{article}
\usepackage[utf8]{inputenc}
\usepackage{booktabs}
\usepackage{amsmath}
\usepackage{natbib}

\title{My Document}
\author{Your Name}
\date{\today}

\begin{document}

\maketitle

\section{Introduction}
This is a sample document demonstrating various \LaTeX{} features including mathematical equations, tables, and formatting. You can cite papers using the citation search feature, and they will be automatically added to your \texttt{cite.bib} file.

\section{Mathematical Expressions}
Here are some example equations:

\subsection{Quadratic Formula}
The quadratic formula is given by:
\[x = \frac{-b \pm \sqrt{b^2 - 4ac}}{2a}\]

\subsection{Matrix Example}
A sample matrix:
\[
\begin{pmatrix}
a_{11} & a_{12} & a_{13} \\
a_{21} & a_{22} & a_{23} \\
a_{31} & a_{32} & a_{33}
\end{pmatrix}
\]

\section{Tables}
Below is a sample table showing data:

\begin{table}[h]
\centering
\begin{tabular}{@{}lrrr@{}}
\toprule
\textbf{Category} & \textbf{Value 1} & \textbf{Value 2} & \textbf{Total} \\
\midrule
Item A & 10.5 & 20.3 & 30.8 \\
Item B & 15.2 & 25.1 & 40.3 \\
Item C & 22.1 & 12.7 & 34.8 \\
\bottomrule
\end{tabular}
\caption{Sample Data Table}
\label{tab:sample}
\end{table}

\section{Lists}
\subsection{Bullet Points}
\begin{itemize}
  \item First item
  \item Second item
  \item Third item with sub-items:
    \begin{itemize}
      \item Sub-item 1
      \item Sub-item 2
    \end{itemize}
\end{itemize}

\subsection{Numbered List}
\begin{enumerate}
  \item First step
  \item Second step
  \item Third step
\end{enumerate}

\section{Citations and References}
When you use the citation search feature, references will be automatically managed. For example, you might cite a paper like this: \cite{example2024}. All citations will be automatically added to your \texttt{cite.bib} file.

\section{Conclusion}
This template demonstrates basic \LaTeX{} features including mathematical formulas, tables, and lists. The bibliography system is now set up to work seamlessly with Bibby AI's citation management. Feel free to modify it for your needs!

\bibliographystyle{plain}
\bibliography{cite}

\end{document}\documentclass{article}
\usepackage[utf8]{inputenc}
\usepackage{booktabs}
\usepackage{amsmath}
\usepackage{natbib}

\title{My Document}
\author{Your Name}
\date{\today}

\begin{document}

\maketitle

\section{Introduction}
This is a sample document demonstrating various \LaTeX{} features including mathematical equations, tables, and formatting. You can cite papers using the citation search feature, and they will be automatically added to your \texttt{cite.bib} file.

\section{Mathematical Expressions}
Here are some example equations:

\subsection{Quadratic Formula}
The quadratic formula is given by:
\[x = \frac{-b \pm \sqrt{b^2 - 4ac}}{2a}\]

\subsection{Matrix Example}
A sample matrix:
\[
\begin{pmatrix}
a_{11} & a_{12} & a_{13} \\
a_{21} & a_{22} & a_{23} \\
a_{31} & a_{32} & a_{33}
\end{pmatrix}
\]

\section{Tables}
Below is a sample table showing data:

\begin{table}[h]
\centering
\begin{tabular}{@{}lrrr@{}}
\toprule
\textbf{Category} & \textbf{Value 1} & \textbf{Value 2} & \textbf{Total} \\
\midrule
Item A & 10.5 & 20.3 & 30.8 \\
Item B & 15.2 & 25.1 & 40.3 \\
Item C & 22.1 & 12.7 & 34.8 \\
\bottomrule
\end{tabular}
\caption{Sample Data Table}
\label{tab:sample}
\end{table}

\section{Lists}
\subsection{Bullet Points}
\begin{itemize}
  \item First item
  \item Second item
  \item Third item with sub-items:
    \begin{itemize}
      \item Sub-item 1
      \item Sub-item 2
    \end{itemize}
\end{itemize}

\subsection{Numbered List}
\begin{enumerate}
  \item First step
  \item Second step
  \item Third step
\end{enumerate}

\section{Citations and References}
When you use the citation search feature, references will be automatically managed. For example, you might cite a paper like this: \cite{example2024}. All citations will be automatically added to your \texttt{cite.bib} file.

\section{Conclusion}
This template demonstrates basic \LaTeX{} features including mathematical formulas, tables, and lists. The bibliography system is now set up to work seamlessly with Bibby AI's citation management. Feel free to modify it for your needs!

\bibliographystyle{plain}
\bibliography{cite}

\end{document}\documentclass{article}
\usepackage[utf8]{inputenc}
\usepackage{booktabs}
\usepackage{amsmath}
\usepackage{natbib}

\title{My Document}
\author{Your Name}
\date{\today}

\begin{document}

\maketitle

\section{Introduction}
This is a sample document demonstrating various \LaTeX{} features including mathematical equations, tables, and formatting. You can cite papers using the citation search feature, and they will be automatically added to your \texttt{cite.bib} file.

\section{Mathematical Expressions}
Here are some example equations:

\subsection{Quadratic Formula}
The quadratic formula is given by:
\[x = \frac{-b \pm \sqrt{b^2 - 4ac}}{2a}\]

\subsection{Matrix Example}
A sample matrix:
\[
\begin{pmatrix}
a_{11} & a_{12} & a_{13} \\
a_{21} & a_{22} & a_{23} \\
a_{31} & a_{32} & a_{33}
\end{pmatrix}
\]

\section{Tables}
Below is a sample table showing data:

\begin{table}[h]
\centering
\begin{tabular}{@{}lrrr@{}}
\toprule
\textbf{Category} & \textbf{Value 1} & \textbf{Value 2} & \textbf{Total} \\
\midrule
Item A & 10.5 & 20.3 & 30.8 \\
Item B & 15.2 & 25.1 & 40.3 \\
Item C & 22.1 & 12.7 & 34.8 \\
\bottomrule
\end{tabular}
\caption{Sample Data Table}
\label{tab:sample}
\end{table}

\section{Lists}
\subsection{Bullet Points}
\begin{itemize}
  \item First item
  \item Second item
  \item Third item with sub-items:
    \begin{itemize}
      \item Sub-item 1
      \item Sub-item 2
    \end{itemize}
\end{itemize}

\subsection{Numbered List}
\begin{enumerate}
  \item First step
  \item Second step
  \item Third step
\end{enumerate}

\section{Citations and References}
When you use the citation search feature, references will be automatically managed. For example, you might cite a paper like this: \cite{example2024}. All citations will be automatically added to your \texttt{cite.bib} file.

\section{Conclusion}
This template demonstrates basic \LaTeX{} features including mathematical formulas, tables, and lists. The bibliography system is now set up to work seamlessly with Bibby AI's citation management. Feel free to modify it for your needs!

\bibliographystyle{plain}
\bibliography{cite}

\end{document}\documentclass{article}
\usepackage[utf8]{inputenc}
\usepackage{booktabs}
\usepackage{amsmath}
\usepackage{natbib}

\title{My Document}
\author{Your Name}
\date{\today}

\begin{document}

\maketitle

\section{Introduction}
This is a sample document demonstrating various \LaTeX{} features including mathematical equations, tables, and formatting. You can cite papers using the citation search feature, and they will be automatically added to your \texttt{cite.bib} file.

\section{Mathematical Expressions}
Here are some example equations:

\subsection{Quadratic Formula}
The quadratic formula is given by:
\[x = \frac{-b \pm \sqrt{b^2 - 4ac}}{2a}\]

\subsection{Matrix Example}
A sample matrix:
\[
\begin{pmatrix}
a_{11} & a_{12} & a_{13} \\
a_{21} & a_{22} & a_{23} \\
a_{31} & a_{32} & a_{33}
\end{pmatrix}
\]

\section{Tables}
Below is a sample table showing data:

\begin{table}[h]
\centering
\begin{tabular}{@{}lrrr@{}}
\toprule
\textbf{Category} & \textbf{Value 1} & \textbf{Value 2} & \textbf{Total} \\
\midrule
Item A & 10.5 & 20.3 & 30.8 \\
Item B & 15.2 & 25.1 & 40.3 \\
Item C & 22.1 & 12.7 & 34.8 \\
\bottomrule
\end{tabular}
\caption{Sample Data Table}
\label{tab:sample}
\end{table}

\section{Lists}
\subsection{Bullet Points}
\begin{itemize}
  \item First item
  \item Second item
  \item Third item with sub-items:
    \begin{itemize}
      \item Sub-item 1
      \item Sub-item 2
    \end{itemize}
\end{itemize}

\subsection{Numbered List}
\begin{enumerate}
  \item First step
  \item Second step
  \item Third step
\end{enumerate}

\section{Citations and References}
When you use the citation search feature, references will be automatically managed. For example, you might cite a paper like this: \cite{example2024}. All citations will be automatically added to your \texttt{cite.bib} file.

\section{Conclusion}
This template demonstrates basic \LaTeX{} features including mathematical formulas, tables, and lists. The bibliography system is now set up to work seamlessly with Bibby AI's citation management. Feel free to modify it for your needs!

\bibliographystyle{plain}
\bibliography{cite}

\end{document}\documentclass{article}
\usepackage[utf8]{inputenc}
\usepackage{booktabs}
\usepackage{amsmath}
\usepackage{natbib}

\title{My Document}
\author{Your Name}
\date{\today}

\begin{document}

\maketitle

\section{Introduction}
This is a sample document demonstrating various \LaTeX{} features including mathematical equations, tables, and formatting. You can cite papers using the citation search feature, and they will be automatically added to your \texttt{cite.bib} file.

\section{Mathematical Expressions}
Here are some example equations:

\subsection{Quadratic Formula}
The quadratic formula is given by:
\[x = \frac{-b \pm \sqrt{b^2 - 4ac}}{2a}\]

\subsection{Matrix Example}
A sample matrix:
\[
\begin{pmatrix}
a_{11} & a_{12} & a_{13} \\
a_{21} & a_{22} & a_{23} \\
a_{31} & a_{32} & a_{33}
\end{pmatrix}
\]

\section{Tables}
Below is a sample table showing data:

\begin{table}[h]
\centering
\begin{tabular}{@{}lrrr@{}}
\toprule
\textbf{Category} & \textbf{Value 1} & \textbf{Value 2} & \textbf{Total} \\
\midrule
Item A & 10.5 & 20.3 & 30.8 \\
Item B & 15.2 & 25.1 & 40.3 \\
Item C & 22.1 & 12.7 & 34.8 \\
\bottomrule
\end{tabular}
\caption{Sample Data Table}
\label{tab:sample}
\end{table}

\section{Lists}
\subsection{Bullet Points}
\begin{itemize}
  \item First item
  \item Second item
  \item Third item with sub-items:
    \begin{itemize}
      \item Sub-item 1
      \item Sub-item 2
    \end{itemize}
\end{itemize}

\subsection{Numbered List}
\begin{enumerate}
  \item First step
  \item Second step
  \item Third step
\end{enumerate}

\section{Citations and References}
When you use the citation search feature, references will be automatically managed. For example, you might cite a paper like this: \cite{example2024}. All citations will be automatically added to your \texttt{cite.bib} file.

\section{Conclusion}
This template demonstrates basic \LaTeX{} features including mathematical formulas, tables, and lists. The bibliography system is now set up to work seamlessly with Bibby AI's citation management. Feel free to modify it for your needs!

\bibliographystyle{plain}
\bibliography{cite}

\end{document}\documentclass{article}
\usepackage[utf8]{inputenc}
\usepackage{booktabs}
\usepackage{amsmath}
\usepackage{natbib}

\title{My Document}
\author{Your Name}
\date{\today}

\begin{document}

\maketitle

\section{Introduction}
This is a sample document demonstrating various \LaTeX{} features including mathematical equations, tables, and formatting. You can cite papers using the citation search feature, and they will be automatically added to your \texttt{cite.bib} file.

\section{Mathematical Expressions}
Here are some example equations:

\subsection{Quadratic Formula}
The quadratic formula is given by:
\[x = \frac{-b \pm \sqrt{b^2 - 4ac}}{2a}\]

\subsection{Matrix Example}
A sample matrix:
\[
\begin{pmatrix}
a_{11} & a_{12} & a_{13} \\
a_{21} & a_{22} & a_{23} \\
a_{31} & a_{32} & a_{33}
\end{pmatrix}
\]

\section{Tables}
Below is a sample table showing data:

\begin{table}[h]
\centering
\begin{tabular}{@{}lrrr@{}}
\toprule
\textbf{Category} & \textbf{Value 1} & \textbf{Value 2} & \textbf{Total} \\
\midrule
Item A & 10.5 & 20.3 & 30.8 \\
Item B & 15.2 & 25.1 & 40.3 \\
Item C & 22.1 & 12.7 & 34.8 \\
\bottomrule
\end{tabular}
\caption{Sample Data Table}
\label{tab:sample}
\end{table}

\section{Lists}
\subsection{Bullet Points}
\begin{itemize}
  \item First item
  \item Second item
  \item Third item with sub-items:
    \begin{itemize}
      \item Sub-item 1
      \item Sub-item 2
    \end{itemize}
\end{itemize}

\subsection{Numbered List}
\begin{enumerate}
  \item First step
  \item Second step
  \item Third step
\end{enumerate}

\section{Citations and References}
When you use the citation search feature, references will be automatically managed. For example, you might cite a paper like this: \cite{example2024}. All citations will be automatically added to your \texttt{cite.bib} file.

\section{Conclusion}
This template demonstrates basic \LaTeX{} features including mathematical formulas, tables, and lists. The bibliography system is now set up to work seamlessly with Bibby AI's citation management. Feel free to modify it for your needs!

\bibliographystyle{plain}
\bibliography{cite}

\end{document}\documentclass{article}
\usepackage[utf8]{inputenc}
\usepackage{booktabs}
\usepackage{amsmath}
\usepackage{natbib}

\title{My Document}
\author{Your Name}
\date{\today}

\begin{document}

\maketitle

\section{Introduction}
This is a sample document demonstrating various \LaTeX{} features including mathematical equations, tables, and formatting. You can cite papers using the citation search feature, and they will be automatically added to your \texttt{cite.bib} file.

\section{Mathematical Expressions}
Here are some example equations:

\subsection{Quadratic Formula}
The quadratic formula is given by:
\[x = \frac{-b \pm \sqrt{b^2 - 4ac}}{2a}\]

\subsection{Matrix Example}
A sample matrix:
\[
\begin{pmatrix}
a_{11} & a_{12} & a_{13} \\
a_{21} & a_{22} & a_{23} \\
a_{31} & a_{32} & a_{33}
\end{pmatrix}
\]

\section{Tables}
Below is a sample table showing data:

\begin{table}[h]
\centering
\begin{tabular}{@{}lrrr@{}}
\toprule
\textbf{Category} & \textbf{Value 1} & \textbf{Value 2} & \textbf{Total} \\
\midrule
Item A & 10.5 & 20.3 & 30.8 \\
Item B & 15.2 & 25.1 & 40.3 \\
Item C & 22.1 & 12.7 & 34.8 \\
\bottomrule
\end{tabular}
\caption{Sample Data Table}
\label{tab:sample}
\end{table}

\section{Lists}
\subsection{Bullet Points}
\begin{itemize}
  \item First item
  \item Second item
  \item Third item with sub-items:
    \begin{itemize}
      \item Sub-item 1
      \item Sub-item 2
    \end{itemize}
\end{itemize}

\subsection{Numbered List}
\begin{enumerate}
  \item First step
  \item Second step
  \item Third step
\end{enumerate}

\section{Citations and References}
When you use the citation search feature, references will be automatically managed. For example, you might cite a paper like this: \cite{example2024}. All citations will be automatically added to your \texttt{cite.bib} file.

\section{Conclusion}
This template demonstrates basic \LaTeX{} features including mathematical formulas, tables, and lists. The bibliography system is now set up to work seamlessly with Bibby AI's citation management. Feel free to modify it for your needs!

\bibliographystyle{plain}
\bibliography{cite}

\end{document}\documentclass{article}
\usepackage[utf8]{inputenc}
\usepackage{booktabs}
\usepackage{amsmath}
\usepackage{natbib}

\title{My Document}
\author{Your Name}
\date{\today}

\begin{document}

\maketitle

\section{Introduction}
This is a sample document demonstrating various \LaTeX{} features including mathematical equations, tables, and formatting. You can cite papers using the citation search feature, and they will be automatically added to your \texttt{cite.bib} file.

\section{Mathematical Expressions}
Here are some example equations:

\subsection{Quadratic Formula}
The quadratic formula is given by:
\[x = \frac{-b \pm \sqrt{b^2 - 4ac}}{2a}\]

\subsection{Matrix Example}
A sample matrix:
\[
\begin{pmatrix}
a_{11} & a_{12} & a_{13} \\
a_{21} & a_{22} & a_{23} \\
a_{31} & a_{32} & a_{33}
\end{pmatrix}
\]

\section{Tables}
Below is a sample table showing data:

\begin{table}[h]
\centering
\begin{tabular}{@{}lrrr@{}}
\toprule
\textbf{Category} & \textbf{Value 1} & \textbf{Value 2} & \textbf{Total} \\
\midrule
Item A & 10.5 & 20.3 & 30.8 \\
Item B & 15.2 & 25.1 & 40.3 \\
Item C & 22.1 & 12.7 & 34.8 \\
\bottomrule
\end{tabular}
\caption{Sample Data Table}
\label{tab:sample}
\end{table}

\section{Lists}
\subsection{Bullet Points}
\begin{itemize}
  \item First item
  \item Second item
  \item Third item with sub-items:
    \begin{itemize}
      \item Sub-item 1
      \item Sub-item 2
    \end{itemize}
\end{itemize}

\subsection{Numbered List}
\begin{enumerate}
  \item First step
  \item Second step
  \item Third step
\end{enumerate}

\section{Citations and References}
When you use the citation search feature, references will be automatically managed. For example, you might cite a paper like this: \cite{example2024}. All citations will be automatically added to your \texttt{cite.bib} file.

\section{Conclusion}
This template demonstrates basic \LaTeX{} features including mathematical formulas, tables, and lists. The bibliography system is now set up to work seamlessly with Bibby AI's citation management. Feel free to modify it for your needs!

\bibliographystyle{plain}
\bibliography{cite}

\end{document}\documentclass{article}
\usepackage[utf8]{inputenc}
\usepackage{booktabs}
\usepackage{amsmath}
\usepackage{natbib}

\title{My Document}
\author{Your Name}
\date{\today}

\begin{document}

\maketitle

\section{Introduction}
This is a sample document demonstrating various \LaTeX{} features including mathematical equations, tables, and formatting. You can cite papers using the citation search feature, and they will be automatically added to your \texttt{cite.bib} file.

\section{Mathematical Expressions}
Here are some example equations:

\subsection{Quadratic Formula}
The quadratic formula is given by:
\[x = \frac{-b \pm \sqrt{b^2 - 4ac}}{2a}\]

\subsection{Matrix Example}
A sample matrix:
\[
\begin{pmatrix}
a_{11} & a_{12} & a_{13} \\
a_{21} & a_{22} & a_{23} \\
a_{31} & a_{32} & a_{33}
\end{pmatrix}
\]

\section{Tables}
Below is a sample table showing data:

\begin{table}[h]
\centering
\begin{tabular}{@{}lrrr@{}}
\toprule
\textbf{Category} & \textbf{Value 1} & \textbf{Value 2} & \textbf{Total} \\
\midrule
Item A & 10.5 & 20.3 & 30.8 \\
Item B & 15.2 & 25.1 & 40.3 \\
Item C & 22.1 & 12.7 & 34.8 \\
\bottomrule
\end{tabular}
\caption{Sample Data Table}
\label{tab:sample}
\end{table}

\section{Lists}
\subsection{Bullet Points}
\begin{itemize}
  \item First item
  \item Second item
  \item Third item with sub-items:
    \begin{itemize}
      \item Sub-item 1
      \item Sub-item 2
    \end{itemize}
\end{itemize}

\subsection{Numbered List}
\begin{enumerate}
  \item First step
  \item Second step
  \item Third step
\end{enumerate}

\section{Citations and References}
When you use the citation search feature, references will be automatically managed. For example, you might cite a paper like this: \cite{example2024}. All citations will be automatically added to your \texttt{cite.bib} file.

\section{Conclusion}
This template demonstrates basic \LaTeX{} features including mathematical formulas, tables, and lists. The bibliography system is now set up to work seamlessly with Bibby AI's citation management. Feel free to modify it for your needs!

\bibliographystyle{plain}
\bibliography{cite}

\end{document}\documentclass{article}
\usepackage[utf8]{inputenc}
\usepackage{booktabs}
\usepackage{amsmath}
\usepackage{natbib}

\title{My Document}
\author{Your Name}
\date{\today}

\begin{document}

\maketitle

\section{Introduction}
This is a sample document demonstrating various \LaTeX{} features including mathematical equations, tables, and formatting. You can cite papers using the citation search feature, and they will be automatically added to your \texttt{cite.bib} file.

\section{Mathematical Expressions}
Here are some example equations:

\subsection{Quadratic Formula}
The quadratic formula is given by:
\[x = \frac{-b \pm \sqrt{b^2 - 4ac}}{2a}\]

\subsection{Matrix Example}
A sample matrix:
\[
\begin{pmatrix}
a_{11} & a_{12} & a_{13} \\
a_{21} & a_{22} & a_{23} \\
a_{31} & a_{32} & a_{33}
\end{pmatrix}
\]

\section{Tables}
Below is a sample table showing data:

\begin{table}[h]
\centering
\begin{tabular}{@{}lrrr@{}}
\toprule
\textbf{Category} & \textbf{Value 1} & \textbf{Value 2} & \textbf{Total} \\
\midrule
Item A & 10.5 & 20.3 & 30.8 \\
Item B & 15.2 & 25.1 & 40.3 \\
Item C & 22.1 & 12.7 & 34.8 \\
\bottomrule
\end{tabular}
\caption{Sample Data Table}
\label{tab:sample}
\end{table}

\section{Lists}
\subsection{Bullet Points}
\begin{itemize}
  \item First item
  \item Second item
  \item Third item with sub-items:
    \begin{itemize}
      \item Sub-item 1
      \item Sub-item 2
    \end{itemize}
\end{itemize}

\subsection{Numbered List}
\begin{enumerate}
  \item First step
  \item Second step
  \item Third step
\end{enumerate}

\section{Citations and References}
When you use the citation search feature, references will be automatically managed. For example, you might cite a paper like this: \cite{example2024}. All citations will be automatically added to your \texttt{cite.bib} file.

\section{Conclusion}
This template demonstrates basic \LaTeX{} features including mathematical formulas, tables, and lists. The bibliography system is now set up to work seamlessly with Bibby AI's citation management. Feel free to modify it for your needs!

\bibliographystyle{plain}
\bibliography{cite}

\end{document}\documentclass{article}
\usepackage[utf8]{inputenc}
\usepackage{booktabs}
\usepackage{amsmath}
\usepackage{natbib}

\title{My Document}
\author{Your Name}
\date{\today}

\begin{document}

\maketitle

\section{Introduction}
This is a sample document demonstrating various \LaTeX{} features including mathematical equations, tables, and formatting. You can cite papers using the citation search feature, and they will be automatically added to your \texttt{cite.bib} file.

\section{Mathematical Expressions}
Here are some example equations:

\subsection{Quadratic Formula}
The quadratic formula is given by:
\[x = \frac{-b \pm \sqrt{b^2 - 4ac}}{2a}\]

\subsection{Matrix Example}
A sample matrix:
\[
\begin{pmatrix}
a_{11} & a_{12} & a_{13} \\
a_{21} & a_{22} & a_{23} \\
a_{31} & a_{32} & a_{33}
\end{pmatrix}
\]

\section{Tables}
Below is a sample table showing data:

\begin{table}[h]
\centering
\begin{tabular}{@{}lrrr@{}}
\toprule
\textbf{Category} & \textbf{Value 1} & \textbf{Value 2} & \textbf{Total} \\
\midrule
Item A & 10.5 & 20.3 & 30.8 \\
Item B & 15.2 & 25.1 & 40.3 \\
Item C & 22.1 & 12.7 & 34.8 \\
\bottomrule
\end{tabular}
\caption{Sample Data Table}
\label{tab:sample}
\end{table}

\section{Lists}
\subsection{Bullet Points}
\begin{itemize}
  \item First item
  \item Second item
  \item Third item with sub-items:
    \begin{itemize}
      \item Sub-item 1
      \item Sub-item 2
    \end{itemize}
\end{itemize}

\subsection{Numbered List}
\begin{enumerate}
  \item First step
  \item Second step
  \item Third step
\end{enumerate}

\section{Citations and References}
When you use the citation search feature, references will be automatically managed. For example, you might cite a paper like this: \cite{example2024}. All citations will be automatically added to your \texttt{cite.bib} file.

\section{Conclusion}
This template demonstrates basic \LaTeX{} features including mathematical formulas, tables, and lists. The bibliography system is now set up to work seamlessly with Bibby AI's citation management. Feel free to modify it for your needs!

\bibliographystyle{plain}
\bibliography{cite}

\end{document}\documentclass{article}
\usepackage[utf8]{inputenc}
\usepackage{booktabs}
\usepackage{amsmath}
\usepackage{natbib}

\title{My Document}
\author{Your Name}
\date{\today}

\begin{document}

\maketitle

\section{Introduction}
This is a sample document demonstrating various \LaTeX{} features including mathematical equations, tables, and formatting. You can cite papers using the citation search feature, and they will be automatically added to your \texttt{cite.bib} file.

\section{Mathematical Expressions}
Here are some example equations:

\subsection{Quadratic Formula}
The quadratic formula is given by:
\[x = \frac{-b \pm \sqrt{b^2 - 4ac}}{2a}\]

\subsection{Matrix Example}
A sample matrix:
\[
\begin{pmatrix}
a_{11} & a_{12} & a_{13} \\
a_{21} & a_{22} & a_{23} \\
a_{31} & a_{32} & a_{33}
\end{pmatrix}
\]

\section{Tables}
Below is a sample table showing data:

\begin{table}[h]
\centering
\begin{tabular}{@{}lrrr@{}}
\toprule
\textbf{Category} & \textbf{Value 1} & \textbf{Value 2} & \textbf{Total} \\
\midrule
Item A & 10.5 & 20.3 & 30.8 \\
Item B & 15.2 & 25.1 & 40.3 \\
Item C & 22.1 & 12.7 & 34.8 \\
\bottomrule
\end{tabular}
\caption{Sample Data Table}
\label{tab:sample}
\end{table}

\section{Lists}
\subsection{Bullet Points}
\begin{itemize}
  \item First item
  \item Second item
  \item Third item with sub-items:
    \begin{itemize}
      \item Sub-item 1
      \item Sub-item 2
    \end{itemize}
\end{itemize}

\subsection{Numbered List}
\begin{enumerate}
  \item First step
  \item Second step
  \item Third step
\end{enumerate}

\section{Citations and References}
When you use the citation search feature, references will be automatically managed. For example, you might cite a paper like this: \cite{example2024}. All citations will be automatically added to your \texttt{cite.bib} file.

\section{Conclusion}
This template demonstrates basic \LaTeX{} features including mathematical formulas, tables, and lists. The bibliography system is now set up to work seamlessly with Bibby AI's citation management. Feel free to modify it for your needs!

\bibliographystyle{plain}
\bibliography{cite}

\end{document}\documentclass{article}
\usepackage[utf8]{inputenc}
\usepackage{booktabs}
\usepackage{amsmath}
\usepackage{natbib}

\title{My Document}
\author{Your Name}
\date{\today}

\begin{document}

\maketitle

\section{Introduction}
This is a sample document demonstrating various \LaTeX{} features including mathematical equations, tables, and formatting. You can cite papers using the citation search feature, and they will be automatically added to your \texttt{cite.bib} file.

\section{Mathematical Expressions}
Here are some example equations:

\subsection{Quadratic Formula}
The quadratic formula is given by:
\[x = \frac{-b \pm \sqrt{b^2 - 4ac}}{2a}\]

\subsection{Matrix Example}
A sample matrix:
\[
\begin{pmatrix}
a_{11} & a_{12} & a_{13} \\
a_{21} & a_{22} & a_{23} \\
a_{31} & a_{32} & a_{33}
\end{pmatrix}
\]

\section{Tables}
Below is a sample table showing data:

\begin{table}[h]
\centering
\begin{tabular}{@{}lrrr@{}}
\toprule
\textbf{Category} & \textbf{Value 1} & \textbf{Value 2} & \textbf{Total} \\
\midrule
Item A & 10.5 & 20.3 & 30.8 \\
Item B & 15.2 & 25.1 & 40.3 \\
Item C & 22.1 & 12.7 & 34.8 \\
\bottomrule
\end{tabular}
\caption{Sample Data Table}
\label{tab:sample}
\end{table}

\section{Lists}
\subsection{Bullet Points}
\begin{itemize}
  \item First item
  \item Second item
  \item Third item with sub-items:
    \begin{itemize}
      \item Sub-item 1
      \item Sub-item 2
    \end{itemize}
\end{itemize}

\subsection{Numbered List}
\begin{enumerate}
  \item First step
  \item Second step
  \item Third step
\end{enumerate}

\section{Citations and References}
When you use the citation search feature, references will be automatically managed. For example, you might cite a paper like this: \cite{example2024}. All citations will be automatically added to your \texttt{cite.bib} file.

\section{Conclusion}
This template demonstrates basic \LaTeX{} features including mathematical formulas, tables, and lists. The bibliography system is now set up to work seamlessly with Bibby AI's citation management. Feel free to modify it for your needs!

\bibliographystyle{plain}
\bibliography{cite}

\end{document}\documentclass{article}
\usepackage[utf8]{inputenc}
\usepackage{booktabs}
\usepackage{amsmath}
\usepackage{natbib}

\title{My Document}
\author{Your Name}
\date{\today}

\begin{document}

\maketitle

\section{Introduction}
This is a sample document demonstrating various \LaTeX{} features including mathematical equations, tables, and formatting. You can cite papers using the citation search feature, and they will be automatically added to your \texttt{cite.bib} file.

\section{Mathematical Expressions}
Here are some example equations:

\subsection{Quadratic Formula}
The quadratic formula is given by:
\[x = \frac{-b \pm \sqrt{b^2 - 4ac}}{2a}\]

\subsection{Matrix Example}
A sample matrix:
\[
\begin{pmatrix}
a_{11} & a_{12} & a_{13} \\
a_{21} & a_{22} & a_{23} \\
a_{31} & a_{32} & a_{33}
\end{pmatrix}
\]

\section{Tables}
Below is a sample table showing data:

\begin{table}[h]
\centering
\begin{tabular}{@{}lrrr@{}}
\toprule
\textbf{Category} & \textbf{Value 1} & \textbf{Value 2} & \textbf{Total} \\
\midrule
Item A & 10.5 & 20.3 & 30.8 \\
Item B & 15.2 & 25.1 & 40.3 \\
Item C & 22.1 & 12.7 & 34.8 \\
\bottomrule
\end{tabular}
\caption{Sample Data Table}
\label{tab:sample}
\end{table}

\section{Lists}
\subsection{Bullet Points}
\begin{itemize}
  \item First item
  \item Second item
  \item Third item with sub-items:
    \begin{itemize}
      \item Sub-item 1
      \item Sub-item 2
    \end{itemize}
\end{itemize}

\subsection{Numbered List}
\begin{enumerate}
  \item First step
  \item Second step
  \item Third step
\end{enumerate}

\section{Citations and References}
When you use the citation search feature, references will be automatically managed. For example, you might cite a paper like this: \cite{example2024}. All citations will be automatically added to your \texttt{cite.bib} file.

\section{Conclusion}
This template demonstrates basic \LaTeX{} features including mathematical formulas, tables, and lists. The bibliography system is now set up to work seamlessly with Bibby AI's citation management. Feel free to modify it for your needs!

\bibliographystyle{plain}
\bibliography{cite}

\end{document}\documentclass{article}
\usepackage[utf8]{inputenc}
\usepackage{booktabs}
\usepackage{amsmath}
\usepackage{natbib}

\title{My Document}
\author{Your Name}
\date{\today}

\begin{document}

\maketitle

\section{Introduction}
This is a sample document demonstrating various \LaTeX{} features including mathematical equations, tables, and formatting. You can cite papers using the citation search feature, and they will be automatically added to your \texttt{cite.bib} file.

\section{Mathematical Expressions}
Here are some example equations:

\subsection{Quadratic Formula}
The quadratic formula is given by:
\[x = \frac{-b \pm \sqrt{b^2 - 4ac}}{2a}\]

\subsection{Matrix Example}
A sample matrix:
\[
\begin{pmatrix}
a_{11} & a_{12} & a_{13} \\
a_{21} & a_{22} & a_{23} \\
a_{31} & a_{32} & a_{33}
\end{pmatrix}
\]

\section{Tables}
Below is a sample table showing data:

\begin{table}[h]
\centering
\begin{tabular}{@{}lrrr@{}}
\toprule
\textbf{Category} & \textbf{Value 1} & \textbf{Value 2} & \textbf{Total} \\
\midrule
Item A & 10.5 & 20.3 & 30.8 \\
Item B & 15.2 & 25.1 & 40.3 \\
Item C & 22.1 & 12.7 & 34.8 \\
\bottomrule
\end{tabular}
\caption{Sample Data Table}
\label{tab:sample}
\end{table}

\section{Lists}
\subsection{Bullet Points}
\begin{itemize}
  \item First item
  \item Second item
  \item Third item with sub-items:
    \begin{itemize}
      \item Sub-item 1
      \item Sub-item 2
    \end{itemize}
\end{itemize}

\subsection{Numbered List}
\begin{enumerate}
  \item First step
  \item Second step
  \item Third step
\end{enumerate}

\section{Citations and References}
When you use the citation search feature, references will be automatically managed. For example, you might cite a paper like this: \cite{example2024}. All citations will be automatically added to your \texttt{cite.bib} file.

\section{Conclusion}
This template demonstrates basic \LaTeX{} features including mathematical formulas, tables, and lists. The bibliography system is now set up to work seamlessly with Bibby AI's citation management. Feel free to modify it for your needs!

\bibliographystyle{plain}
\bibliography{cite}

\end{document}\documentclass{article}
\usepackage[utf8]{inputenc}
\usepackage{booktabs}
\usepackage{amsmath}
\usepackage{natbib}

\title{My Document}
\author{Your Name}
\date{\today}

\begin{document}

\maketitle

\section{Introduction}
This is a sample document demonstrating various \LaTeX{} features including mathematical equations, tables, and formatting. You can cite papers using the citation search feature, and they will be automatically added to your \texttt{cite.bib} file.

\section{Mathematical Expressions}
Here are some example equations:

\subsection{Quadratic Formula}
The quadratic formula is given by:
\[x = \frac{-b \pm \sqrt{b^2 - 4ac}}{2a}\]

\subsection{Matrix Example}
A sample matrix:
\[
\begin{pmatrix}
a_{11} & a_{12} & a_{13} \\
a_{21} & a_{22} & a_{23} \\
a_{31} & a_{32} & a_{33}
\end{pmatrix}
\]

\section{Tables}
Below is a sample table showing data:

\begin{table}[h]
\centering
\begin{tabular}{@{}lrrr@{}}
\toprule
\textbf{Category} & \textbf{Value 1} & \textbf{Value 2} & \textbf{Total} \\
\midrule
Item A & 10.5 & 20.3 & 30.8 \\
Item B & 15.2 & 25.1 & 40.3 \\
Item C & 22.1 & 12.7 & 34.8 \\
\bottomrule
\end{tabular}
\caption{Sample Data Table}
\label{tab:sample}
\end{table}

\section{Lists}
\subsection{Bullet Points}
\begin{itemize}
  \item First item
  \item Second item
  \item Third item with sub-items:
    \begin{itemize}
      \item Sub-item 1
      \item Sub-item 2
    \end{itemize}
\end{itemize}

\subsection{Numbered List}
\begin{enumerate}
  \item First step
  \item Second step
  \item Third step
\end{enumerate}

\section{Citations and References}
When you use the citation search feature, references will be automatically managed. For example, you might cite a paper like this: \cite{example2024}. All citations will be automatically added to your \texttt{cite.bib} file.

\section{Conclusion}
This template demonstrates basic \LaTeX{} features including mathematical formulas, tables, and lists. The bibliography system is now set up to work seamlessly with Bibby AI's citation management. Feel free to modify it for your needs!

\bibliographystyle{plain}
\bibliography{cite}

\end{document}\documentclass{article}
\usepackage[utf8]{inputenc}
\usepackage{booktabs}
\usepackage{amsmath}
\usepackage{natbib}

\title{My Document}
\author{Your Name}
\date{\today}

\begin{document}

\maketitle

\section{Introduction}
This is a sample document demonstrating various \LaTeX{} features including mathematical equations, tables, and formatting. You can cite papers using the citation search feature, and they will be automatically added to your \texttt{cite.bib} file.

\section{Mathematical Expressions}
Here are some example equations:

\subsection{Quadratic Formula}
The quadratic formula is given by:
\[x = \frac{-b \pm \sqrt{b^2 - 4ac}}{2a}\]

\subsection{Matrix Example}
A sample matrix:
\[
\begin{pmatrix}
a_{11} & a_{12} & a_{13} \\
a_{21} & a_{22} & a_{23} \\
a_{31} & a_{32} & a_{33}
\end{pmatrix}
\]

\section{Tables}
Below is a sample table showing data:

\begin{table}[h]
\centering
\begin{tabular}{@{}lrrr@{}}
\toprule
\textbf{Category} & \textbf{Value 1} & \textbf{Value 2} & \textbf{Total} \\
\midrule
Item A & 10.5 & 20.3 & 30.8 \\
Item B & 15.2 & 25.1 & 40.3 \\
Item C & 22.1 & 12.7 & 34.8 \\
\bottomrule
\end{tabular}
\caption{Sample Data Table}
\label{tab:sample}
\end{table}

\section{Lists}
\subsection{Bullet Points}
\begin{itemize}
  \item First item
  \item Second item
  \item Third item with sub-items:
    \begin{itemize}
      \item Sub-item 1
      \item Sub-item 2
    \end{itemize}
\end{itemize}

\subsection{Numbered List}
\begin{enumerate}
  \item First step
  \item Second step
  \item Third step
\end{enumerate}

\section{Citations and References}
When you use the citation search feature, references will be automatically managed. For example, you might cite a paper like this: \cite{example2024}. All citations will be automatically added to your \texttt{cite.bib} file.

\section{Conclusion}
This template demonstrates basic \LaTeX{} features including mathematical formulas, tables, and lists. The bibliography system is now set up to work seamlessly with Bibby AI's citation management. Feel free to modify it for your needs!

\bibliographystyle{plain}
\bibliography{cite}

\end{document}\documentclass{article}
\usepackage[utf8]{inputenc}
\usepackage{booktabs}
\usepackage{amsmath}
\usepackage{natbib}

\title{My Document}
\author{Your Name}
\date{\today}

\begin{document}

\maketitle

\section{Introduction}
This is a sample document demonstrating various \LaTeX{} features including mathematical equations, tables, and formatting. You can cite papers using the citation search feature, and they will be automatically added to your \texttt{cite.bib} file.

\section{Mathematical Expressions}
Here are some example equations:

\subsection{Quadratic Formula}
The quadratic formula is given by:
\[x = \frac{-b \pm \sqrt{b^2 - 4ac}}{2a}\]

\subsection{Matrix Example}
A sample matrix:
\[
\begin{pmatrix}
a_{11} & a_{12} & a_{13} \\
a_{21} & a_{22} & a_{23} \\
a_{31} & a_{32} & a_{33}
\end{pmatrix}
\]

\section{Tables}
Below is a sample table showing data:

\begin{table}[h]
\centering
\begin{tabular}{@{}lrrr@{}}
\toprule
\textbf{Category} & \textbf{Value 1} & \textbf{Value 2} & \textbf{Total} \\
\midrule
Item A & 10.5 & 20.3 & 30.8 \\
Item B & 15.2 & 25.1 & 40.3 \\
Item C & 22.1 & 12.7 & 34.8 \\
\bottomrule
\end{tabular}
\caption{Sample Data Table}
\label{tab:sample}
\end{table}

\section{Lists}
\subsection{Bullet Points}
\begin{itemize}
  \item First item
  \item Second item
  \item Third item with sub-items:
    \begin{itemize}
      \item Sub-item 1
      \item Sub-item 2
    \end{itemize}
\end{itemize}

\subsection{Numbered List}
\begin{enumerate}
  \item First step
  \item Second step
  \item Third step
\end{enumerate}

\section{Citations and References}
When you use the citation search feature, references will be automatically managed. For example, you might cite a paper like this: \cite{example2024}. All citations will be automatically added to your \texttt{cite.bib} file.

\section{Conclusion}
This template demonstrates basic \LaTeX{} features including mathematical formulas, tables, and lists. The bibliography system is now set up to work seamlessly with Bibby AI's citation management. Feel free to modify it for your needs!

\bibliographystyle{plain}
\bibliography{cite}

\end{document}\documentclass{article}
\usepackage[utf8]{inputenc}
\usepackage{booktabs}
\usepackage{amsmath}
\usepackage{natbib}

\title{My Document}
\author{Your Name}
\date{\today}

\begin{document}

\maketitle

\section{Introduction}
This is a sample document demonstrating various \LaTeX{} features including mathematical equations, tables, and formatting. You can cite papers using the citation search feature, and they will be automatically added to your \texttt{cite.bib} file.

\section{Mathematical Expressions}
Here are some example equations:

\subsection{Quadratic Formula}
The quadratic formula is given by:
\[x = \frac{-b \pm \sqrt{b^2 - 4ac}}{2a}\]

\subsection{Matrix Example}
A sample matrix:
\[
\begin{pmatrix}
a_{11} & a_{12} & a_{13} \\
a_{21} & a_{22} & a_{23} \\
a_{31} & a_{32} & a_{33}
\end{pmatrix}
\]

\section{Tables}
Below is a sample table showing data:

\begin{table}[h]
\centering
\begin{tabular}{@{}lrrr@{}}
\toprule
\textbf{Category} & \textbf{Value 1} & \textbf{Value 2} & \textbf{Total} \\
\midrule
Item A & 10.5 & 20.3 & 30.8 \\
Item B & 15.2 & 25.1 & 40.3 \\
Item C & 22.1 & 12.7 & 34.8 \\
\bottomrule
\end{tabular}
\caption{Sample Data Table}
\label{tab:sample}
\end{table}

\section{Lists}
\subsection{Bullet Points}
\begin{itemize}
  \item First item
  \item Second item
  \item Third item with sub-items:
    \begin{itemize}
      \item Sub-item 1
      \item Sub-item 2
    \end{itemize}
\end{itemize}

\subsection{Numbered List}
\begin{enumerate}
  \item First step
  \item Second step
  \item Third step
\end{enumerate}

\section{Citations and References}
When you use the citation search feature, references will be automatically managed. For example, you might cite a paper like this: \cite{example2024}. All citations will be automatically added to your \texttt{cite.bib} file.

\section{Conclusion}
This template demonstrates basic \LaTeX{} features including mathematical formulas, tables, and lists. The bibliography system is now set up to work seamlessly with Bibby AI's citation management. Feel free to modify it for your needs!

\bibliographystyle{plain}
\bibliography{cite}

\end{document}\documentclass{article}
\usepackage[utf8]{inputenc}
\usepackage{booktabs}
\usepackage{amsmath}
\usepackage{natbib}

\title{My Document}
\author{Your Name}
\date{\today}

\begin{document}

\maketitle

\section{Introduction}
This is a sample document demonstrating various \LaTeX{} features including mathematical equations, tables, and formatting. You can cite papers using the citation search feature, and they will be automatically added to your \texttt{cite.bib} file.

\section{Mathematical Expressions}
Here are some example equations:

\subsection{Quadratic Formula}
The quadratic formula is given by:
\[x = \frac{-b \pm \sqrt{b^2 - 4ac}}{2a}\]

\subsection{Matrix Example}
A sample matrix:
\[
\begin{pmatrix}
a_{11} & a_{12} & a_{13} \\
a_{21} & a_{22} & a_{23} \\
a_{31} & a_{32} & a_{33}
\end{pmatrix}
\]

\section{Tables}
Below is a sample table showing data:

\begin{table}[h]
\centering
\begin{tabular}{@{}lrrr@{}}
\toprule
\textbf{Category} & \textbf{Value 1} & \textbf{Value 2} & \textbf{Total} \\
\midrule
Item A & 10.5 & 20.3 & 30.8 \\
Item B & 15.2 & 25.1 & 40.3 \\
Item C & 22.1 & 12.7 & 34.8 \\
\bottomrule
\end{tabular}
\caption{Sample Data Table}
\label{tab:sample}
\end{table}

\section{Lists}
\subsection{Bullet Points}
\begin{itemize}
  \item First item
  \item Second item
  \item Third item with sub-items:
    \begin{itemize}
      \item Sub-item 1
      \item Sub-item 2
    \end{itemize}
\end{itemize}

\subsection{Numbered List}
\begin{enumerate}
  \item First step
  \item Second step
  \item Third step
\end{enumerate}

\section{Citations and References}
When you use the citation search feature, references will be automatically managed. For example, you might cite a paper like this: \cite{example2024}. All citations will be automatically added to your \texttt{cite.bib} file.

\section{Conclusion}
This template demonstrates basic \LaTeX{} features including mathematical formulas, tables, and lists. The bibliography system is now set up to work seamlessly with Bibby AI's citation management. Feel free to modify it for your needs!

\bibliographystyle{plain}
\bibliography{cite}

\end{document}\documentclass{article}
\usepackage[utf8]{inputenc}
\usepackage{booktabs}
\usepackage{amsmath}
\usepackage{natbib}

\title{My Document}
\author{Your Name}
\date{\today}

\begin{document}

\maketitle

\section{Introduction}
This is a sample document demonstrating various \LaTeX{} features including mathematical equations, tables, and formatting. You can cite papers using the citation search feature, and they will be automatically added to your \texttt{cite.bib} file.

\section{Mathematical Expressions}
Here are some example equations:

\subsection{Quadratic Formula}
The quadratic formula is given by:
\[x = \frac{-b \pm \sqrt{b^2 - 4ac}}{2a}\]

\subsection{Matrix Example}
A sample matrix:
\[
\begin{pmatrix}
a_{11} & a_{12} & a_{13} \\
a_{21} & a_{22} & a_{23} \\
a_{31} & a_{32} & a_{33}
\end{pmatrix}
\]

\section{Tables}
Below is a sample table showing data:

\begin{table}[h]
\centering
\begin{tabular}{@{}lrrr@{}}
\toprule
\textbf{Category} & \textbf{Value 1} & \textbf{Value 2} & \textbf{Total} \\
\midrule
Item A & 10.5 & 20.3 & 30.8 \\
Item B & 15.2 & 25.1 & 40.3 \\
Item C & 22.1 & 12.7 & 34.8 \\
\bottomrule
\end{tabular}
\caption{Sample Data Table}
\label{tab:sample}
\end{table}

\section{Lists}
\subsection{Bullet Points}
\begin{itemize}
  \item First item
  \item Second item
  \item Third item with sub-items:
    \begin{itemize}
      \item Sub-item 1
      \item Sub-item 2
    \end{itemize}
\end{itemize}

\subsection{Numbered List}
\begin{enumerate}
  \item First step
  \item Second step
  \item Third step
\end{enumerate}

\section{Citations and References}
When you use the citation search feature, references will be automatically managed. For example, you might cite a paper like this: \cite{example2024}. All citations will be automatically added to your \texttt{cite.bib} file.

\section{Conclusion}
This template demonstrates basic \LaTeX{} features including mathematical formulas, tables, and lists. The bibliography system is now set up to work seamlessly with Bibby AI's citation management. Feel free to modify it for your needs!

\bibliographystyle{plain}
\bibliography{cite}

\end{document}\documentclass{article}
\usepackage[utf8]{inputenc}
\usepackage{booktabs}
\usepackage{amsmath}
\usepackage{natbib}

\title{My Document}
\author{Your Name}
\date{\today}

\begin{document}

\maketitle

\section{Introduction}
This is a sample document demonstrating various \LaTeX{} features including mathematical equations, tables, and formatting. You can cite papers using the citation search feature, and they will be automatically added to your \texttt{cite.bib} file.

\section{Mathematical Expressions}
Here are some example equations:

\subsection{Quadratic Formula}
The quadratic formula is given by:
\[x = \frac{-b \pm \sqrt{b^2 - 4ac}}{2a}\]

\subsection{Matrix Example}
A sample matrix:
\[
\begin{pmatrix}
a_{11} & a_{12} & a_{13} \\
a_{21} & a_{22} & a_{23} \\
a_{31} & a_{32} & a_{33}
\end{pmatrix}
\]

\section{Tables}
Below is a sample table showing data:

\begin{table}[h]
\centering
\begin{tabular}{@{}lrrr@{}}
\toprule
\textbf{Category} & \textbf{Value 1} & \textbf{Value 2} & \textbf{Total} \\
\midrule
Item A & 10.5 & 20.3 & 30.8 \\
Item B & 15.2 & 25.1 & 40.3 \\
Item C & 22.1 & 12.7 & 34.8 \\
\bottomrule
\end{tabular}
\caption{Sample Data Table}
\label{tab:sample}
\end{table}

\section{Lists}
\subsection{Bullet Points}
\begin{itemize}
  \item First item
  \item Second item
  \item Third item with sub-items:
    \begin{itemize}
      \item Sub-item 1
      \item Sub-item 2
    \end{itemize}
\end{itemize}

\subsection{Numbered List}
\begin{enumerate}
  \item First step
  \item Second step
  \item Third step
\end{enumerate}

\section{Citations and References}
When you use the citation search feature, references will be automatically managed. For example, you might cite a paper like this: \cite{example2024}. All citations will be automatically added to your \texttt{cite.bib} file.

\section{Conclusion}
This template demonstrates basic \LaTeX{} features including mathematical formulas, tables, and lists. The bibliography system is now set up to work seamlessly with Bibby AI's citation management. Feel free to modify it for your needs!

\bibliographystyle{plain}
\bibliography{cite}

\end{document}\documentclass{article}
\usepackage[utf8]{inputenc}
\usepackage{booktabs}
\usepackage{amsmath}
\usepackage{natbib}

\title{My Document}
\author{Your Name}
\date{\today}

\begin{document}

\maketitle

\section{Introduction}
This is a sample document demonstrating various \LaTeX{} features including mathematical equations, tables, and formatting. You can cite papers using the citation search feature, and they will be automatically added to your \texttt{cite.bib} file.

\section{Mathematical Expressions}
Here are some example equations:

\subsection{Quadratic Formula}
The quadratic formula is given by:
\[x = \frac{-b \pm \sqrt{b^2 - 4ac}}{2a}\]

\subsection{Matrix Example}
A sample matrix:
\[
\begin{pmatrix}
a_{11} & a_{12} & a_{13} \\
a_{21} & a_{22} & a_{23} \\
a_{31} & a_{32} & a_{33}
\end{pmatrix}
\]

\section{Tables}
Below is a sample table showing data:

\begin{table}[h]
\centering
\begin{tabular}{@{}lrrr@{}}
\toprule
\textbf{Category} & \textbf{Value 1} & \textbf{Value 2} & \textbf{Total} \\
\midrule
Item A & 10.5 & 20.3 & 30.8 \\
Item B & 15.2 & 25.1 & 40.3 \\
Item C & 22.1 & 12.7 & 34.8 \\
\bottomrule
\end{tabular}
\caption{Sample Data Table}
\label{tab:sample}
\end{table}

\section{Lists}
\subsection{Bullet Points}
\begin{itemize}
  \item First item
  \item Second item
  \item Third item with sub-items:
    \begin{itemize}
      \item Sub-item 1
      \item Sub-item 2
    \end{itemize}
\end{itemize}

\subsection{Numbered List}
\begin{enumerate}
  \item First step
  \item Second step
  \item Third step
\end{enumerate}

\section{Citations and References}
When you use the citation search feature, references will be automatically managed. For example, you might cite a paper like this: \cite{example2024}. All citations will be automatically added to your \texttt{cite.bib} file.

\section{Conclusion}
This template demonstrates basic \LaTeX{} features including mathematical formulas, tables, and lists. The bibliography system is now set up to work seamlessly with Bibby AI's citation management. Feel free to modify it for your needs!

\bibliographystyle{plain}
\bibliography{cite}

\end{document}\documentclass{article}
\usepackage[utf8]{inputenc}
\usepackage{booktabs}
\usepackage{amsmath}
\usepackage{natbib}

\title{My Document}
\author{Your Name}
\date{\today}

\begin{document}

\maketitle

\section{Introduction}
This is a sample document demonstrating various \LaTeX{} features including mathematical equations, tables, and formatting. You can cite papers using the citation search feature, and they will be automatically added to your \texttt{cite.bib} file.

\section{Mathematical Expressions}
Here are some example equations:

\subsection{Quadratic Formula}
The quadratic formula is given by:
\[x = \frac{-b \pm \sqrt{b^2 - 4ac}}{2a}\]

\subsection{Matrix Example}
A sample matrix:
\[
\begin{pmatrix}
a_{11} & a_{12} & a_{13} \\
a_{21} & a_{22} & a_{23} \\
a_{31} & a_{32} & a_{33}
\end{pmatrix}
\]

\section{Tables}
Below is a sample table showing data:

\begin{table}[h]
\centering
\begin{tabular}{@{}lrrr@{}}
\toprule
\textbf{Category} & \textbf{Value 1} & \textbf{Value 2} & \textbf{Total} \\
\midrule
Item A & 10.5 & 20.3 & 30.8 \\
Item B & 15.2 & 25.1 & 40.3 \\
Item C & 22.1 & 12.7 & 34.8 \\
\bottomrule
\end{tabular}
\caption{Sample Data Table}
\label{tab:sample}
\end{table}

\section{Lists}
\subsection{Bullet Points}
\begin{itemize}
  \item First item
  \item Second item
  \item Third item with sub-items:
    \begin{itemize}
      \item Sub-item 1
      \item Sub-item 2
    \end{itemize}
\end{itemize}

\subsection{Numbered List}
\begin{enumerate}
  \item First step
  \item Second step
  \item Third step
\end{enumerate}

\section{Citations and References}
When you use the citation search feature, references will be automatically managed. For example, you might cite a paper like this: \cite{example2024}. All citations will be automatically added to your \texttt{cite.bib} file.

\section{Conclusion}
This template demonstrates basic \LaTeX{} features including mathematical formulas, tables, and lists. The bibliography system is now set up to work seamlessly with Bibby AI's citation management. Feel free to modify it for your needs!

\bibliographystyle{plain}
\bibliography{cite}

\end{document}\documentclass{article}
\usepackage[utf8]{inputenc}
\usepackage{booktabs}
\usepackage{amsmath}
\usepackage{natbib}

\title{My Document}
\author{Your Name}
\date{\today}

\begin{document}

\maketitle

\section{Introduction}
This is a sample document demonstrating various \LaTeX{} features including mathematical equations, tables, and formatting. You can cite papers using the citation search feature, and they will be automatically added to your \texttt{cite.bib} file.

\section{Mathematical Expressions}
Here are some example equations:

\subsection{Quadratic Formula}
The quadratic formula is given by:
\[x = \frac{-b \pm \sqrt{b^2 - 4ac}}{2a}\]

\subsection{Matrix Example}
A sample matrix:
\[
\begin{pmatrix}
a_{11} & a_{12} & a_{13} \\
a_{21} & a_{22} & a_{23} \\
a_{31} & a_{32} & a_{33}
\end{pmatrix}
\]

\section{Tables}
Below is a sample table showing data:

\begin{table}[h]
\centering
\begin{tabular}{@{}lrrr@{}}
\toprule
\textbf{Category} & \textbf{Value 1} & \textbf{Value 2} & \textbf{Total} \\
\midrule
Item A & 10.5 & 20.3 & 30.8 \\
Item B & 15.2 & 25.1 & 40.3 \\
Item C & 22.1 & 12.7 & 34.8 \\
\bottomrule
\end{tabular}
\caption{Sample Data Table}
\label{tab:sample}
\end{table}

\section{Lists}
\subsection{Bullet Points}
\begin{itemize}
  \item First item
  \item Second item
  \item Third item with sub-items:
    \begin{itemize}
      \item Sub-item 1
      \item Sub-item 2
    \end{itemize}
\end{itemize}

\subsection{Numbered List}
\begin{enumerate}
  \item First step
  \item Second step
  \item Third step
\end{enumerate}

\section{Citations and References}
When you use the citation search feature, references will be automatically managed. For example, you might cite a paper like this: \cite{example2024}. All citations will be automatically added to your \texttt{cite.bib} file.

\section{Conclusion}
This template demonstrates basic \LaTeX{} features including mathematical formulas, tables, and lists. The bibliography system is now set up to work seamlessly with Bibby AI's citation management. Feel free to modify it for your needs!

\bibliographystyle{plain}
\bibliography{cite}

\end{document}\documentclass{article}
\usepackage[utf8]{inputenc}
\usepackage{booktabs}
\usepackage{amsmath}
\usepackage{natbib}

\title{My Document}
\author{Your Name}
\date{\today}

\begin{document}

\maketitle

\section{Introduction}
This is a sample document demonstrating various \LaTeX{} features including mathematical equations, tables, and formatting. You can cite papers using the citation search feature, and they will be automatically added to your \texttt{cite.bib} file.

\section{Mathematical Expressions}
Here are some example equations:

\subsection{Quadratic Formula}
The quadratic formula is given by:
\[x = \frac{-b \pm \sqrt{b^2 - 4ac}}{2a}\]

\subsection{Matrix Example}
A sample matrix:
\[
\begin{pmatrix}
a_{11} & a_{12} & a_{13} \\
a_{21} & a_{22} & a_{23} \\
a_{31} & a_{32} & a_{33}
\end{pmatrix}
\]

\section{Tables}
Below is a sample table showing data:

\begin{table}[h]
\centering
\begin{tabular}{@{}lrrr@{}}
\toprule
\textbf{Category} & \textbf{Value 1} & \textbf{Value 2} & \textbf{Total} \\
\midrule
Item A & 10.5 & 20.3 & 30.8 \\
Item B & 15.2 & 25.1 & 40.3 \\
Item C & 22.1 & 12.7 & 34.8 \\
\bottomrule
\end{tabular}
\caption{Sample Data Table}
\label{tab:sample}
\end{table}

\section{Lists}
\subsection{Bullet Points}
\begin{itemize}
  \item First item
  \item Second item
  \item Third item with sub-items:
    \begin{itemize}
      \item Sub-item 1
      \item Sub-item 2
    \end{itemize}
\end{itemize}

\subsection{Numbered List}
\begin{enumerate}
  \item First step
  \item Second step
  \item Third step
\end{enumerate}

\section{Citations and References}
When you use the citation search feature, references will be automatically managed. For example, you might cite a paper like this: \cite{example2024}. All citations will be automatically added to your \texttt{cite.bib} file.

\section{Conclusion}
This template demonstrates basic \LaTeX{} features including mathematical formulas, tables, and lists. The bibliography system is now set up to work seamlessly with Bibby AI's citation management. Feel free to modify it for your needs!

\bibliographystyle{plain}
\bibliography{cite}

\end{document}\documentclass{article}
\usepackage[utf8]{inputenc}
\usepackage{booktabs}
\usepackage{amsmath}
\usepackage{natbib}

\title{My Document}
\author{Your Name}
\date{\today}

\begin{document}

\maketitle

\section{Introduction}
This is a sample document demonstrating various \LaTeX{} features including mathematical equations, tables, and formatting. You can cite papers using the citation search feature, and they will be automatically added to your \texttt{cite.bib} file.

\section{Mathematical Expressions}
Here are some example equations:

\subsection{Quadratic Formula}
The quadratic formula is given by:
\[x = \frac{-b \pm \sqrt{b^2 - 4ac}}{2a}\]

\subsection{Matrix Example}
A sample matrix:
\[
\begin{pmatrix}
a_{11} & a_{12} & a_{13} \\
a_{21} & a_{22} & a_{23} \\
a_{31} & a_{32} & a_{33}
\end{pmatrix}
\]

\section{Tables}
Below is a sample table showing data:

\begin{table}[h]
\centering
\begin{tabular}{@{}lrrr@{}}
\toprule
\textbf{Category} & \textbf{Value 1} & \textbf{Value 2} & \textbf{Total} \\
\midrule
Item A & 10.5 & 20.3 & 30.8 \\
Item B & 15.2 & 25.1 & 40.3 \\
Item C & 22.1 & 12.7 & 34.8 \\
\bottomrule
\end{tabular}
\caption{Sample Data Table}
\label{tab:sample}
\end{table}

\section{Lists}
\subsection{Bullet Points}
\begin{itemize}
  \item First item
  \item Second item
  \item Third item with sub-items:
    \begin{itemize}
      \item Sub-item 1
      \item Sub-item 2
    \end{itemize}
\end{itemize}

\subsection{Numbered List}
\begin{enumerate}
  \item First step
  \item Second step
  \item Third step
\end{enumerate}

\section{Citations and References}
When you use the citation search feature, references will be automatically managed. For example, you might cite a paper like this: \cite{example2024}. All citations will be automatically added to your \texttt{cite.bib} file.

\section{Conclusion}
This template demonstrates basic \LaTeX{} features including mathematical formulas, tables, and lists. The bibliography system is now set up to work seamlessly with Bibby AI's citation management. Feel free to modify it for your needs!

\bibliographystyle{plain}
\bibliography{cite}

\end{document}\documentclass{article}
\usepackage[utf8]{inputenc}
\usepackage{booktabs}
\usepackage{amsmath}
\usepackage{natbib}

\title{My Document}
\author{Your Name}
\date{\today}

\begin{document}

\maketitle

\section{Introduction}
This is a sample document demonstrating various \LaTeX{} features including mathematical equations, tables, and formatting. You can cite papers using the citation search feature, and they will be automatically added to your \texttt{cite.bib} file.

\section{Mathematical Expressions}
Here are some example equations:

\subsection{Quadratic Formula}
The quadratic formula is given by:
\[x = \frac{-b \pm \sqrt{b^2 - 4ac}}{2a}\]

\subsection{Matrix Example}
A sample matrix:
\[
\begin{pmatrix}
a_{11} & a_{12} & a_{13} \\
a_{21} & a_{22} & a_{23} \\
a_{31} & a_{32} & a_{33}
\end{pmatrix}
\]

\section{Tables}
Below is a sample table showing data:

\begin{table}[h]
\centering
\begin{tabular}{@{}lrrr@{}}
\toprule
\textbf{Category} & \textbf{Value 1} & \textbf{Value 2} & \textbf{Total} \\
\midrule
Item A & 10.5 & 20.3 & 30.8 \\
Item B & 15.2 & 25.1 & 40.3 \\
Item C & 22.1 & 12.7 & 34.8 \\
\bottomrule
\end{tabular}
\caption{Sample Data Table}
\label{tab:sample}
\end{table}

\section{Lists}
\subsection{Bullet Points}
\begin{itemize}
  \item First item
  \item Second item
  \item Third item with sub-items:
    \begin{itemize}
      \item Sub-item 1
      \item Sub-item 2
    \end{itemize}
\end{itemize}

\subsection{Numbered List}
\begin{enumerate}
  \item First step
  \item Second step
  \item Third step
\end{enumerate}

\section{Citations and References}
When you use the citation search feature, references will be automatically managed. For example, you might cite a paper like this: \cite{example2024}. All citations will be automatically added to your \texttt{cite.bib} file.

\section{Conclusion}
This template demonstrates basic \LaTeX{} features including mathematical formulas, tables, and lists. The bibliography system is now set up to work seamlessly with Bibby AI's citation management. Feel free to modify it for your needs!

\bibliographystyle{plain}
\bibliography{cite}

\end{document}\documentclass{article}
\usepackage[utf8]{inputenc}
\usepackage{booktabs}
\usepackage{amsmath}
\usepackage{natbib}

\title{My Document}
\author{Your Name}
\date{\today}

\begin{document}

\maketitle

\section{Introduction}
This is a sample document demonstrating various \LaTeX{} features including mathematical equations, tables, and formatting. You can cite papers using the citation search feature, and they will be automatically added to your \texttt{cite.bib} file.

\section{Mathematical Expressions}
Here are some example equations:

\subsection{Quadratic Formula}
The quadratic formula is given by:
\[x = \frac{-b \pm \sqrt{b^2 - 4ac}}{2a}\]

\subsection{Matrix Example}
A sample matrix:
\[
\begin{pmatrix}
a_{11} & a_{12} & a_{13} \\
a_{21} & a_{22} & a_{23} \\
a_{31} & a_{32} & a_{33}
\end{pmatrix}
\]

\section{Tables}
Below is a sample table showing data:

\begin{table}[h]
\centering
\begin{tabular}{@{}lrrr@{}}
\toprule
\textbf{Category} & \textbf{Value 1} & \textbf{Value 2} & \textbf{Total} \\
\midrule
Item A & 10.5 & 20.3 & 30.8 \\
Item B & 15.2 & 25.1 & 40.3 \\
Item C & 22.1 & 12.7 & 34.8 \\
\bottomrule
\end{tabular}
\caption{Sample Data Table}
\label{tab:sample}
\end{table}

\section{Lists}
\subsection{Bullet Points}
\begin{itemize}
  \item First item
  \item Second item
  \item Third item with sub-items:
    \begin{itemize}
      \item Sub-item 1
      \item Sub-item 2
    \end{itemize}
\end{itemize}

\subsection{Numbered List}
\begin{enumerate}
  \item First step
  \item Second step
  \item Third step
\end{enumerate}

\section{Citations and References}
When you use the citation search feature, references will be automatically managed. For example, you might cite a paper like this: \cite{example2024}. All citations will be automatically added to your \texttt{cite.bib} file.

\section{Conclusion}
This template demonstrates basic \LaTeX{} features including mathematical formulas, tables, and lists. The bibliography system is now set up to work seamlessly with Bibby AI's citation management. Feel free to modify it for your needs!

\bibliographystyle{plain}
\bibliography{cite}

\end{document}\documentclass{article}
\usepackage[utf8]{inputenc}
\usepackage{booktabs}
\usepackage{amsmath}
\usepackage{natbib}

\title{My Document}
\author{Your Name}
\date{\today}

\begin{document}

\maketitle

\section{Introduction}
This is a sample document demonstrating various \LaTeX{} features including mathematical equations, tables, and formatting. You can cite papers using the citation search feature, and they will be automatically added to your \texttt{cite.bib} file.

\section{Mathematical Expressions}
Here are some example equations:

\subsection{Quadratic Formula}
The quadratic formula is given by:
\[x = \frac{-b \pm \sqrt{b^2 - 4ac}}{2a}\]

\subsection{Matrix Example}
A sample matrix:
\[
\begin{pmatrix}
a_{11} & a_{12} & a_{13} \\
a_{21} & a_{22} & a_{23} \\
a_{31} & a_{32} & a_{33}
\end{pmatrix}
\]

\section{Tables}
Below is a sample table showing data:

\begin{table}[h]
\centering
\begin{tabular}{@{}lrrr@{}}
\toprule
\textbf{Category} & \textbf{Value 1} & \textbf{Value 2} & \textbf{Total} \\
\midrule
Item A & 10.5 & 20.3 & 30.8 \\
Item B & 15.2 & 25.1 & 40.3 \\
Item C & 22.1 & 12.7 & 34.8 \\
\bottomrule
\end{tabular}
\caption{Sample Data Table}
\label{tab:sample}
\end{table}

\section{Lists}
\subsection{Bullet Points}
\begin{itemize}
  \item First item
  \item Second item
  \item Third item with sub-items:
    \begin{itemize}
      \item Sub-item 1
      \item Sub-item 2
    \end{itemize}
\end{itemize}

\subsection{Numbered List}
\begin{enumerate}
  \item First step
  \item Second step
  \item Third step
\end{enumerate}

\section{Citations and References}
When you use the citation search feature, references will be automatically managed. For example, you might cite a paper like this: \cite{example2024}. All citations will be automatically added to your \texttt{cite.bib} file.

\section{Conclusion}
This template demonstrates basic \LaTeX{} features including mathematical formulas, tables, and lists. The bibliography system is now set up to work seamlessly with Bibby AI's citation management. Feel free to modify it for your needs!

\bibliographystyle{plain}
\bibliography{cite}

\end{document}\documentclass{article}
\usepackage[utf8]{inputenc}
\usepackage{booktabs}
\usepackage{amsmath}
\usepackage{natbib}

\title{My Document}
\author{Your Name}
\date{\today}

\begin{document}

\maketitle

\section{Introduction}
This is a sample document demonstrating various \LaTeX{} features including mathematical equations, tables, and formatting. You can cite papers using the citation search feature, and they will be automatically added to your \texttt{cite.bib} file.

\section{Mathematical Expressions}
Here are some example equations:

\subsection{Quadratic Formula}
The quadratic formula is given by:
\[x = \frac{-b \pm \sqrt{b^2 - 4ac}}{2a}\]

\subsection{Matrix Example}
A sample matrix:
\[
\begin{pmatrix}
a_{11} & a_{12} & a_{13} \\
a_{21} & a_{22} & a_{23} \\
a_{31} & a_{32} & a_{33}
\end{pmatrix}
\]

\section{Tables}
Below is a sample table showing data:

\begin{table}[h]
\centering
\begin{tabular}{@{}lrrr@{}}
\toprule
\textbf{Category} & \textbf{Value 1} & \textbf{Value 2} & \textbf{Total} \\
\midrule
Item A & 10.5 & 20.3 & 30.8 \\
Item B & 15.2 & 25.1 & 40.3 \\
Item C & 22.1 & 12.7 & 34.8 \\
\bottomrule
\end{tabular}
\caption{Sample Data Table}
\label{tab:sample}
\end{table}

\section{Lists}
\subsection{Bullet Points}
\begin{itemize}
  \item First item
  \item Second item
  \item Third item with sub-items:
    \begin{itemize}
      \item Sub-item 1
      \item Sub-item 2
    \end{itemize}
\end{itemize}

\subsection{Numbered List}
\begin{enumerate}
  \item First step
  \item Second step
  \item Third step
\end{enumerate}

\section{Citations and References}
When you use the citation search feature, references will be automatically managed. For example, you might cite a paper like this: \cite{example2024}. All citations will be automatically added to your \texttt{cite.bib} file.

\section{Conclusion}
This template demonstrates basic \LaTeX{} features including mathematical formulas, tables, and lists. The bibliography system is now set up to work seamlessly with Bibby AI's citation management. Feel free to modify it for your needs!

\bibliographystyle{plain}
\bibliography{cite}

\end{document}\documentclass{article}
\usepackage[utf8]{inputenc}
\usepackage{booktabs}
\usepackage{amsmath}
\usepackage{natbib}

\title{My Document}
\author{Your Name}
\date{\today}

\begin{document}

\maketitle

\section{Introduction}
This is a sample document demonstrating various \LaTeX{} features including mathematical equations, tables, and formatting. You can cite papers using the citation search feature, and they will be automatically added to your \texttt{cite.bib} file.

\section{Mathematical Expressions}
Here are some example equations:

\subsection{Quadratic Formula}
The quadratic formula is given by:
\[x = \frac{-b \pm \sqrt{b^2 - 4ac}}{2a}\]

\subsection{Matrix Example}
A sample matrix:
\[
\begin{pmatrix}
a_{11} & a_{12} & a_{13} \\
a_{21} & a_{22} & a_{23} \\
a_{31} & a_{32} & a_{33}
\end{pmatrix}
\]

\section{Tables}
Below is a sample table showing data:

\begin{table}[h]
\centering
\begin{tabular}{@{}lrrr@{}}
\toprule
\textbf{Category} & \textbf{Value 1} & \textbf{Value 2} & \textbf{Total} \\
\midrule
Item A & 10.5 & 20.3 & 30.8 \\
Item B & 15.2 & 25.1 & 40.3 \\
Item C & 22.1 & 12.7 & 34.8 \\
\bottomrule
\end{tabular}
\caption{Sample Data Table}
\label{tab:sample}
\end{table}

\section{Lists}
\subsection{Bullet Points}
\begin{itemize}
  \item First item
  \item Second item
  \item Third item with sub-items:
    \begin{itemize}
      \item Sub-item 1
      \item Sub-item 2
    \end{itemize}
\end{itemize}

\subsection{Numbered List}
\begin{enumerate}
  \item First step
  \item Second step
  \item Third step
\end{enumerate}

\section{Citations and References}
When you use the citation search feature, references will be automatically managed. For example, you might cite a paper like this: \cite{example2024}. All citations will be automatically added to your \texttt{cite.bib} file.

\section{Conclusion}
This template demonstrates basic \LaTeX{} features including mathematical formulas, tables, and lists. The bibliography system is now set up to work seamlessly with Bibby AI's citation management. Feel free to modify it for your needs!

\bibliographystyle{plain}
\bibliography{cite}

\end{document}\documentclass{article}
\usepackage[utf8]{inputenc}
\usepackage{booktabs}
\usepackage{amsmath}
\usepackage{natbib}

\title{My Document}
\author{Your Name}
\date{\today}

\begin{document}

\maketitle

\section{Introduction}
This is a sample document demonstrating various \LaTeX{} features including mathematical equations, tables, and formatting. You can cite papers using the citation search feature, and they will be automatically added to your \texttt{cite.bib} file.

\section{Mathematical Expressions}
Here are some example equations:

\subsection{Quadratic Formula}
The quadratic formula is given by:
\[x = \frac{-b \pm \sqrt{b^2 - 4ac}}{2a}\]

\subsection{Matrix Example}
A sample matrix:
\[
\begin{pmatrix}
a_{11} & a_{12} & a_{13} \\
a_{21} & a_{22} & a_{23} \\
a_{31} & a_{32} & a_{33}
\end{pmatrix}
\]

\section{Tables}
Below is a sample table showing data:

\begin{table}[h]
\centering
\begin{tabular}{@{}lrrr@{}}
\toprule
\textbf{Category} & \textbf{Value 1} & \textbf{Value 2} & \textbf{Total} \\
\midrule
Item A & 10.5 & 20.3 & 30.8 \\
Item B & 15.2 & 25.1 & 40.3 \\
Item C & 22.1 & 12.7 & 34.8 \\
\bottomrule
\end{tabular}
\caption{Sample Data Table}
\label{tab:sample}
\end{table}

\section{Lists}
\subsection{Bullet Points}
\begin{itemize}
  \item First item
  \item Second item
  \item Third item with sub-items:
    \begin{itemize}
      \item Sub-item 1
      \item Sub-item 2
    \end{itemize}
\end{itemize}

\subsection{Numbered List}
\begin{enumerate}
  \item First step
  \item Second step
  \item Third step
\end{enumerate}

\section{Citations and References}
When you use the citation search feature, references will be automatically managed. For example, you might cite a paper like this: \cite{example2024}. All citations will be automatically added to your \texttt{cite.bib} file.

\section{Conclusion}
This template demonstrates basic \LaTeX{} features including mathematical formulas, tables, and lists. The bibliography system is now set up to work seamlessly with Bibby AI's citation management. Feel free to modify it for your needs!

\bibliographystyle{plain}
\bibliography{cite}

\end{document}\documentclass{article}
\usepackage[utf8]{inputenc}
\usepackage{booktabs}
\usepackage{amsmath}
\usepackage{natbib}

\title{My Document}
\author{Your Name}
\date{\today}

\begin{document}

\maketitle

\section{Introduction}
This is a sample document demonstrating various \LaTeX{} features including mathematical equations, tables, and formatting. You can cite papers using the citation search feature, and they will be automatically added to your \texttt{cite.bib} file.

\section{Mathematical Expressions}
Here are some example equations:

\subsection{Quadratic Formula}
The quadratic formula is given by:
\[x = \frac{-b \pm \sqrt{b^2 - 4ac}}{2a}\]

\subsection{Matrix Example}
A sample matrix:
\[
\begin{pmatrix}
a_{11} & a_{12} & a_{13} \\
a_{21} & a_{22} & a_{23} \\
a_{31} & a_{32} & a_{33}
\end{pmatrix}
\]

\section{Tables}
Below is a sample table showing data:

\begin{table}[h]
\centering
\begin{tabular}{@{}lrrr@{}}
\toprule
\textbf{Category} & \textbf{Value 1} & \textbf{Value 2} & \textbf{Total} \\
\midrule
Item A & 10.5 & 20.3 & 30.8 \\
Item B & 15.2 & 25.1 & 40.3 \\
Item C & 22.1 & 12.7 & 34.8 \\
\bottomrule
\end{tabular}
\caption{Sample Data Table}
\label{tab:sample}
\end{table}

\section{Lists}
\subsection{Bullet Points}
\begin{itemize}
  \item First item
  \item Second item
  \item Third item with sub-items:
    \begin{itemize}
      \item Sub-item 1
      \item Sub-item 2
    \end{itemize}
\end{itemize}

\subsection{Numbered List}
\begin{enumerate}
  \item First step
  \item Second step
  \item Third step
\end{enumerate}

\section{Citations and References}
When you use the citation search feature, references will be automatically managed. For example, you might cite a paper like this: \cite{example2024}. All citations will be automatically added to your \texttt{cite.bib} file.

\section{Conclusion}
This template demonstrates basic \LaTeX{} features including mathematical formulas, tables, and lists. The bibliography system is now set up to work seamlessly with Bibby AI's citation management. Feel free to modify it for your needs!

\bibliographystyle{plain}
\bibliography{cite}

\end{document}\documentclass{article}
\usepackage[utf8]{inputenc}
\usepackage{booktabs}
\usepackage{amsmath}
\usepackage{natbib}

\title{My Document}
\author{Your Name}
\date{\today}

\begin{document}

\maketitle

\section{Introduction}
This is a sample document demonstrating various \LaTeX{} features including mathematical equations, tables, and formatting. You can cite papers using the citation search feature, and they will be automatically added to your \texttt{cite.bib} file.

\section{Mathematical Expressions}
Here are some example equations:

\subsection{Quadratic Formula}
The quadratic formula is given by:
\[x = \frac{-b \pm \sqrt{b^2 - 4ac}}{2a}\]

\subsection{Matrix Example}
A sample matrix:
\[
\begin{pmatrix}
a_{11} & a_{12} & a_{13} \\
a_{21} & a_{22} & a_{23} \\
a_{31} & a_{32} & a_{33}
\end{pmatrix}
\]

\section{Tables}
Below is a sample table showing data:

\begin{table}[h]
\centering
\begin{tabular}{@{}lrrr@{}}
\toprule
\textbf{Category} & \textbf{Value 1} & \textbf{Value 2} & \textbf{Total} \\
\midrule
Item A & 10.5 & 20.3 & 30.8 \\
Item B & 15.2 & 25.1 & 40.3 \\
Item C & 22.1 & 12.7 & 34.8 \\
\bottomrule
\end{tabular}
\caption{Sample Data Table}
\label{tab:sample}
\end{table}

\section{Lists}
\subsection{Bullet Points}
\begin{itemize}
  \item First item
  \item Second item
  \item Third item with sub-items:
    \begin{itemize}
      \item Sub-item 1
      \item Sub-item 2
    \end{itemize}
\end{itemize}

\subsection{Numbered List}
\begin{enumerate}
  \item First step
  \item Second step
  \item Third step
\end{enumerate}

\section{Citations and References}
When you use the citation search feature, references will be automatically managed. For example, you might cite a paper like this: \cite{example2024}. All citations will be automatically added to your \texttt{cite.bib} file.

\section{Conclusion}
This template demonstrates basic \LaTeX{} features including mathematical formulas, tables, and lists. The bibliography system is now set up to work seamlessly with Bibby AI's citation management. Feel free to modify it for your needs!

\bibliographystyle{plain}
\bibliography{cite}

\end{document}\documentclass{article}
\usepackage[utf8]{inputenc}
\usepackage{booktabs}
\usepackage{amsmath}
\usepackage{natbib}

\title{My Document}
\author{Your Name}
\date{\today}

\begin{document}

\maketitle

\section{Introduction}
This is a sample document demonstrating various \LaTeX{} features including mathematical equations, tables, and formatting. You can cite papers using the citation search feature, and they will be automatically added to your \texttt{cite.bib} file.

\section{Mathematical Expressions}
Here are some example equations:

\subsection{Quadratic Formula}
The quadratic formula is given by:
\[x = \frac{-b \pm \sqrt{b^2 - 4ac}}{2a}\]

\subsection{Matrix Example}
A sample matrix:
\[
\begin{pmatrix}
a_{11} & a_{12} & a_{13} \\
a_{21} & a_{22} & a_{23} \\
a_{31} & a_{32} & a_{33}
\end{pmatrix}
\]

\section{Tables}
Below is a sample table showing data:

\begin{table}[h]
\centering
\begin{tabular}{@{}lrrr@{}}
\toprule
\textbf{Category} & \textbf{Value 1} & \textbf{Value 2} & \textbf{Total} \\
\midrule
Item A & 10.5 & 20.3 & 30.8 \\
Item B & 15.2 & 25.1 & 40.3 \\
Item C & 22.1 & 12.7 & 34.8 \\
\bottomrule
\end{tabular}
\caption{Sample Data Table}
\label{tab:sample}
\end{table}

\section{Lists}
\subsection{Bullet Points}
\begin{itemize}
  \item First item
  \item Second item
  \item Third item with sub-items:
    \begin{itemize}
      \item Sub-item 1
      \item Sub-item 2
    \end{itemize}
\end{itemize}

\subsection{Numbered List}
\begin{enumerate}
  \item First step
  \item Second step
  \item Third step
\end{enumerate}

\section{Citations and References}
When you use the citation search feature, references will be automatically managed. For example, you might cite a paper like this: \cite{example2024}. All citations will be automatically added to your \texttt{cite.bib} file.

\section{Conclusion}
This template demonstrates basic \LaTeX{} features including mathematical formulas, tables, and lists. The bibliography system is now set up to work seamlessly with Bibby AI's citation management. Feel free to modify it for your needs!

\bibliographystyle{plain}
\bibliography{cite}

\end{document}\documentclass{article}
\usepackage[utf8]{inputenc}
\usepackage{booktabs}
\usepackage{amsmath}
\usepackage{natbib}

\title{My Document}
\author{Your Name}
\date{\today}

\begin{document}

\maketitle

\section{Introduction}
This is a sample document demonstrating various \LaTeX{} features including mathematical equations, tables, and formatting. You can cite papers using the citation search feature, and they will be automatically added to your \texttt{cite.bib} file.

\section{Mathematical Expressions}
Here are some example equations:

\subsection{Quadratic Formula}
The quadratic formula is given by:
\[x = \frac{-b \pm \sqrt{b^2 - 4ac}}{2a}\]

\subsection{Matrix Example}
A sample matrix:
\[
\begin{pmatrix}
a_{11} & a_{12} & a_{13} \\
a_{21} & a_{22} & a_{23} \\
a_{31} & a_{32} & a_{33}
\end{pmatrix}
\]

\section{Tables}
Below is a sample table showing data:

\begin{table}[h]
\centering
\begin{tabular}{@{}lrrr@{}}
\toprule
\textbf{Category} & \textbf{Value 1} & \textbf{Value 2} & \textbf{Total} \\
\midrule
Item A & 10.5 & 20.3 & 30.8 \\
Item B & 15.2 & 25.1 & 40.3 \\
Item C & 22.1 & 12.7 & 34.8 \\
\bottomrule
\end{tabular}
\caption{Sample Data Table}
\label{tab:sample}
\end{table}

\section{Lists}
\subsection{Bullet Points}
\begin{itemize}
  \item First item
  \item Second item
  \item Third item with sub-items:
    \begin{itemize}
      \item Sub-item 1
      \item Sub-item 2
    \end{itemize}
\end{itemize}

\subsection{Numbered List}
\begin{enumerate}
  \item First step
  \item Second step
  \item Third step
\end{enumerate}

\section{Citations and References}
When you use the citation search feature, references will be automatically managed. For example, you might cite a paper like this: \cite{example2024}. All citations will be automatically added to your \texttt{cite.bib} file.

\section{Conclusion}
This template demonstrates basic \LaTeX{} features including mathematical formulas, tables, and lists. The bibliography system is now set up to work seamlessly with Bibby AI's citation management. Feel free to modify it for your needs!

\bibliographystyle{plain}
\bibliography{cite}

\end{document}\documentclass{article}
\usepackage[utf8]{inputenc}
\usepackage{booktabs}
\usepackage{amsmath}
\usepackage{natbib}

\title{My Document}
\author{Your Name}
\date{\today}

\begin{document}

\maketitle

\section{Introduction}
This is a sample document demonstrating various \LaTeX{} features including mathematical equations, tables, and formatting. You can cite papers using the citation search feature, and they will be automatically added to your \texttt{cite.bib} file.

\section{Mathematical Expressions}
Here are some example equations:

\subsection{Quadratic Formula}
The quadratic formula is given by:
\[x = \frac{-b \pm \sqrt{b^2 - 4ac}}{2a}\]

\subsection{Matrix Example}
A sample matrix:
\[
\begin{pmatrix}
a_{11} & a_{12} & a_{13} \\
a_{21} & a_{22} & a_{23} \\
a_{31} & a_{32} & a_{33}
\end{pmatrix}
\]

\section{Tables}
Below is a sample table showing data:

\begin{table}[h]
\centering
\begin{tabular}{@{}lrrr@{}}
\toprule
\textbf{Category} & \textbf{Value 1} & \textbf{Value 2} & \textbf{Total} \\
\midrule
Item A & 10.5 & 20.3 & 30.8 \\
Item B & 15.2 & 25.1 & 40.3 \\
Item C & 22.1 & 12.7 & 34.8 \\
\bottomrule
\end{tabular}
\caption{Sample Data Table}
\label{tab:sample}
\end{table}

\section{Lists}
\subsection{Bullet Points}
\begin{itemize}
  \item First item
  \item Second item
  \item Third item with sub-items:
    \begin{itemize}
      \item Sub-item 1
      \item Sub-item 2
    \end{itemize}
\end{itemize}

\subsection{Numbered List}
\begin{enumerate}
  \item First step
  \item Second step
  \item Third step
\end{enumerate}

\section{Citations and References}
When you use the citation search feature, references will be automatically managed. For example, you might cite a paper like this: \cite{example2024}. All citations will be automatically added to your \texttt{cite.bib} file.

\section{Conclusion}
This template demonstrates basic \LaTeX{} features including mathematical formulas, tables, and lists. The bibliography system is now set up to work seamlessly with Bibby AI's citation management. Feel free to modify it for your needs!

\bibliographystyle{plain}
\bibliography{cite}

\end{document}\documentclass{article}
\usepackage[utf8]{inputenc}
\usepackage{booktabs}
\usepackage{amsmath}
\usepackage{natbib}

\title{My Document}
\author{Your Name}
\date{\today}

\begin{document}

\maketitle

\section{Introduction}
This is a sample document demonstrating various \LaTeX{} features including mathematical equations, tables, and formatting. You can cite papers using the citation search feature, and they will be automatically added to your \texttt{cite.bib} file.

\section{Mathematical Expressions}
Here are some example equations:

\subsection{Quadratic Formula}
The quadratic formula is given by:
\[x = \frac{-b \pm \sqrt{b^2 - 4ac}}{2a}\]

\subsection{Matrix Example}
A sample matrix:
\[
\begin{pmatrix}
a_{11} & a_{12} & a_{13} \\
a_{21} & a_{22} & a_{23} \\
a_{31} & a_{32} & a_{33}
\end{pmatrix}
\]

\section{Tables}
Below is a sample table showing data:

\begin{table}[h]
\centering
\begin{tabular}{@{}lrrr@{}}
\toprule
\textbf{Category} & \textbf{Value 1} & \textbf{Value 2} & \textbf{Total} \\
\midrule
Item A & 10.5 & 20.3 & 30.8 \\
Item B & 15.2 & 25.1 & 40.3 \\
Item C & 22.1 & 12.7 & 34.8 \\
\bottomrule
\end{tabular}
\caption{Sample Data Table}
\label{tab:sample}
\end{table}

\section{Lists}
\subsection{Bullet Points}
\begin{itemize}
  \item First item
  \item Second item
  \item Third item with sub-items:
    \begin{itemize}
      \item Sub-item 1
      \item Sub-item 2
    \end{itemize}
\end{itemize}

\subsection{Numbered List}
\begin{enumerate}
  \item First step
  \item Second step
  \item Third step
\end{enumerate}

\section{Citations and References}
When you use the citation search feature, references will be automatically managed. For example, you might cite a paper like this: \cite{example2024}. All citations will be automatically added to your \texttt{cite.bib} file.

\section{Conclusion}
This template demonstrates basic \LaTeX{} features including mathematical formulas, tables, and lists. The bibliography system is now set up to work seamlessly with Bibby AI's citation management. Feel free to modify it for your needs!

\bibliographystyle{plain}
\bibliography{cite}

\end{document}\documentclass{article}
\usepackage[utf8]{inputenc}
\usepackage{booktabs}
\usepackage{amsmath}
\usepackage{natbib}

\title{My Document}
\author{Your Name}
\date{\today}

\begin{document}

\maketitle

\section{Introduction}
This is a sample document demonstrating various \LaTeX{} features including mathematical equations, tables, and formatting. You can cite papers using the citation search feature, and they will be automatically added to your \texttt{cite.bib} file.

\section{Mathematical Expressions}
Here are some example equations:

\subsection{Quadratic Formula}
The quadratic formula is given by:
\[x = \frac{-b \pm \sqrt{b^2 - 4ac}}{2a}\]

\subsection{Matrix Example}
A sample matrix:
\[
\begin{pmatrix}
a_{11} & a_{12} & a_{13} \\
a_{21} & a_{22} & a_{23} \\
a_{31} & a_{32} & a_{33}
\end{pmatrix}
\]

\section{Tables}
Below is a sample table showing data:

\begin{table}[h]
\centering
\begin{tabular}{@{}lrrr@{}}
\toprule
\textbf{Category} & \textbf{Value 1} & \textbf{Value 2} & \textbf{Total} \\
\midrule
Item A & 10.5 & 20.3 & 30.8 \\
Item B & 15.2 & 25.1 & 40.3 \\
Item C & 22.1 & 12.7 & 34.8 \\
\bottomrule
\end{tabular}
\caption{Sample Data Table}
\label{tab:sample}
\end{table}

\section{Lists}
\subsection{Bullet Points}
\begin{itemize}
  \item First item
  \item Second item
  \item Third item with sub-items:
    \begin{itemize}
      \item Sub-item 1
      \item Sub-item 2
    \end{itemize}
\end{itemize}

\subsection{Numbered List}
\begin{enumerate}
  \item First step
  \item Second step
  \item Third step
\end{enumerate}

\section{Citations and References}
When you use the citation search feature, references will be automatically managed. For example, you might cite a paper like this: \cite{example2024}. All citations will be automatically added to your \texttt{cite.bib} file.

\section{Conclusion}
This template demonstrates basic \LaTeX{} features including mathematical formulas, tables, and lists. The bibliography system is now set up to work seamlessly with Bibby AI's citation management. Feel free to modify it for your needs!

\bibliographystyle{plain}
\bibliography{cite}

\end{document}\documentclass{article}
\usepackage[utf8]{inputenc}
\usepackage{booktabs}
\usepackage{amsmath}
\usepackage{natbib}

\title{My Document}
\author{Your Name}
\date{\today}

\begin{document}

\maketitle

\section{Introduction}
This is a sample document demonstrating various \LaTeX{} features including mathematical equations, tables, and formatting. You can cite papers using the citation search feature, and they will be automatically added to your \texttt{cite.bib} file.

\section{Mathematical Expressions}
Here are some example equations:

\subsection{Quadratic Formula}
The quadratic formula is given by:
\[x = \frac{-b \pm \sqrt{b^2 - 4ac}}{2a}\]

\subsection{Matrix Example}
A sample matrix:
\[
\begin{pmatrix}
a_{11} & a_{12} & a_{13} \\
a_{21} & a_{22} & a_{23} \\
a_{31} & a_{32} & a_{33}
\end{pmatrix}
\]

\section{Tables}
Below is a sample table showing data:

\begin{table}[h]
\centering
\begin{tabular}{@{}lrrr@{}}
\toprule
\textbf{Category} & \textbf{Value 1} & \textbf{Value 2} & \textbf{Total} \\
\midrule
Item A & 10.5 & 20.3 & 30.8 \\
Item B & 15.2 & 25.1 & 40.3 \\
Item C & 22.1 & 12.7 & 34.8 \\
\bottomrule
\end{tabular}
\caption{Sample Data Table}
\label{tab:sample}
\end{table}

\section{Lists}
\subsection{Bullet Points}
\begin{itemize}
  \item First item
  \item Second item
  \item Third item with sub-items:
    \begin{itemize}
      \item Sub-item 1
      \item Sub-item 2
    \end{itemize}
\end{itemize}

\subsection{Numbered List}
\begin{enumerate}
  \item First step
  \item Second step
  \item Third step
\end{enumerate}

\section{Citations and References}
When you use the citation search feature, references will be automatically managed. For example, you might cite a paper like this: \cite{example2024}. All citations will be automatically added to your \texttt{cite.bib} file.

\section{Conclusion}
This template demonstrates basic \LaTeX{} features including mathematical formulas, tables, and lists. The bibliography system is now set up to work seamlessly with Bibby AI's citation management. Feel free to modify it for your needs!

\bibliographystyle{plain}
\bibliography{cite}

\end{document}\documentclass{article}
\usepackage[utf8]{inputenc}
\usepackage{booktabs}
\usepackage{amsmath}
\usepackage{natbib}

\title{My Document}
\author{Your Name}
\date{\today}

\begin{document}

\maketitle

\section{Introduction}
This is a sample document demonstrating various \LaTeX{} features including mathematical equations, tables, and formatting. You can cite papers using the citation search feature, and they will be automatically added to your \texttt{cite.bib} file.

\section{Mathematical Expressions}
Here are some example equations:

\subsection{Quadratic Formula}
The quadratic formula is given by:
\[x = \frac{-b \pm \sqrt{b^2 - 4ac}}{2a}\]

\subsection{Matrix Example}
A sample matrix:
\[
\begin{pmatrix}
a_{11} & a_{12} & a_{13} \\
a_{21} & a_{22} & a_{23} \\
a_{31} & a_{32} & a_{33}
\end{pmatrix}
\]

\section{Tables}
Below is a sample table showing data:

\begin{table}[h]
\centering
\begin{tabular}{@{}lrrr@{}}
\toprule
\textbf{Category} & \textbf{Value 1} & \textbf{Value 2} & \textbf{Total} \\
\midrule
Item A & 10.5 & 20.3 & 30.8 \\
Item B & 15.2 & 25.1 & 40.3 \\
Item C & 22.1 & 12.7 & 34.8 \\
\bottomrule
\end{tabular}
\caption{Sample Data Table}
\label{tab:sample}
\end{table}

\section{Lists}
\subsection{Bullet Points}
\begin{itemize}
  \item First item
  \item Second item
  \item Third item with sub-items:
    \begin{itemize}
      \item Sub-item 1
      \item Sub-item 2
    \end{itemize}
\end{itemize}

\subsection{Numbered List}
\begin{enumerate}
  \item First step
  \item Second step
  \item Third step
\end{enumerate}

\section{Citations and References}
When you use the citation search feature, references will be automatically managed. For example, you might cite a paper like this: \cite{example2024}. All citations will be automatically added to your \texttt{cite.bib} file.

\section{Conclusion}
This template demonstrates basic \LaTeX{} features including mathematical formulas, tables, and lists. The bibliography system is now set up to work seamlessly with Bibby AI's citation management. Feel free to modify it for your needs!

\bibliographystyle{plain}
\bibliography{cite}

\end{document}\documentclass{article}
\usepackage[utf8]{inputenc}
\usepackage{booktabs}
\usepackage{amsmath}
\usepackage{natbib}

\title{My Document}
\author{Your Name}
\date{\today}

\begin{document}

\maketitle

\section{Introduction}
This is a sample document demonstrating various \LaTeX{} features including mathematical equations, tables, and formatting. You can cite papers using the citation search feature, and they will be automatically added to your \texttt{cite.bib} file.

\section{Mathematical Expressions}
Here are some example equations:

\subsection{Quadratic Formula}
The quadratic formula is given by:
\[x = \frac{-b \pm \sqrt{b^2 - 4ac}}{2a}\]

\subsection{Matrix Example}
A sample matrix:
\[
\begin{pmatrix}
a_{11} & a_{12} & a_{13} \\
a_{21} & a_{22} & a_{23} \\
a_{31} & a_{32} & a_{33}
\end{pmatrix}
\]

\section{Tables}
Below is a sample table showing data:

\begin{table}[h]
\centering
\begin{tabular}{@{}lrrr@{}}
\toprule
\textbf{Category} & \textbf{Value 1} & \textbf{Value 2} & \textbf{Total} \\
\midrule
Item A & 10.5 & 20.3 & 30.8 \\
Item B & 15.2 & 25.1 & 40.3 \\
Item C & 22.1 & 12.7 & 34.8 \\
\bottomrule
\end{tabular}
\caption{Sample Data Table}
\label{tab:sample}
\end{table}

\section{Lists}
\subsection{Bullet Points}
\begin{itemize}
  \item First item
  \item Second item
  \item Third item with sub-items:
    \begin{itemize}
      \item Sub-item 1
      \item Sub-item 2
    \end{itemize}
\end{itemize}

\subsection{Numbered List}
\begin{enumerate}
  \item First step
  \item Second step
  \item Third step
\end{enumerate}

\section{Citations and References}
When you use the citation search feature, references will be automatically managed. For example, you might cite a paper like this: \cite{example2024}. All citations will be automatically added to your \texttt{cite.bib} file.

\section{Conclusion}
This template demonstrates basic \LaTeX{} features including mathematical formulas, tables, and lists. The bibliography system is now set up to work seamlessly with Bibby AI's citation management. Feel free to modify it for your needs!

\bibliographystyle{plain}
\bibliography{cite}

\end{document}\documentclass{article}
\usepackage[utf8]{inputenc}
\usepackage{booktabs}
\usepackage{amsmath}
\usepackage{natbib}

\title{My Document}
\author{Your Name}
\date{\today}

\begin{document}

\maketitle

\section{Introduction}
This is a sample document demonstrating various \LaTeX{} features including mathematical equations, tables, and formatting. You can cite papers using the citation search feature, and they will be automatically added to your \texttt{cite.bib} file.

\section{Mathematical Expressions}
Here are some example equations:

\subsection{Quadratic Formula}
The quadratic formula is given by:
\[x = \frac{-b \pm \sqrt{b^2 - 4ac}}{2a}\]

\subsection{Matrix Example}
A sample matrix:
\[
\begin{pmatrix}
a_{11} & a_{12} & a_{13} \\
a_{21} & a_{22} & a_{23} \\
a_{31} & a_{32} & a_{33}
\end{pmatrix}
\]

\section{Tables}
Below is a sample table showing data:

\begin{table}[h]
\centering
\begin{tabular}{@{}lrrr@{}}
\toprule
\textbf{Category} & \textbf{Value 1} & \textbf{Value 2} & \textbf{Total} \\
\midrule
Item A & 10.5 & 20.3 & 30.8 \\
Item B & 15.2 & 25.1 & 40.3 \\
Item C & 22.1 & 12.7 & 34.8 \\
\bottomrule
\end{tabular}
\caption{Sample Data Table}
\label{tab:sample}
\end{table}

\section{Lists}
\subsection{Bullet Points}
\begin{itemize}
  \item First item
  \item Second item
  \item Third item with sub-items:
    \begin{itemize}
      \item Sub-item 1
      \item Sub-item 2
    \end{itemize}
\end{itemize}

\subsection{Numbered List}
\begin{enumerate}
  \item First step
  \item Second step
  \item Third step
\end{enumerate}

\section{Citations and References}
When you use the citation search feature, references will be automatically managed. For example, you might cite a paper like this: \cite{example2024}. All citations will be automatically added to your \texttt{cite.bib} file.

\section{Conclusion}
This template demonstrates basic \LaTeX{} features including mathematical formulas, tables, and lists. The bibliography system is now set up to work seamlessly with Bibby AI's citation management. Feel free to modify it for your needs!

\bibliographystyle{plain}
\bibliography{cite}

\end{document}\documentclass{article}
\usepackage[utf8]{inputenc}
\usepackage{booktabs}
\usepackage{amsmath}
\usepackage{natbib}

\title{My Document}
\author{Your Name}
\date{\today}

\begin{document}

\maketitle

\section{Introduction}
This is a sample document demonstrating various \LaTeX{} features including mathematical equations, tables, and formatting. You can cite papers using the citation search feature, and they will be automatically added to your \texttt{cite.bib} file.

\section{Mathematical Expressions}
Here are some example equations:

\subsection{Quadratic Formula}
The quadratic formula is given by:
\[x = \frac{-b \pm \sqrt{b^2 - 4ac}}{2a}\]

\subsection{Matrix Example}
A sample matrix:
\[
\begin{pmatrix}
a_{11} & a_{12} & a_{13} \\
a_{21} & a_{22} & a_{23} \\
a_{31} & a_{32} & a_{33}
\end{pmatrix}
\]

\section{Tables}
Below is a sample table showing data:

\begin{table}[h]
\centering
\begin{tabular}{@{}lrrr@{}}
\toprule
\textbf{Category} & \textbf{Value 1} & \textbf{Value 2} & \textbf{Total} \\
\midrule
Item A & 10.5 & 20.3 & 30.8 \\
Item B & 15.2 & 25.1 & 40.3 \\
Item C & 22.1 & 12.7 & 34.8 \\
\bottomrule
\end{tabular}
\caption{Sample Data Table}
\label{tab:sample}
\end{table}

\section{Lists}
\subsection{Bullet Points}
\begin{itemize}
  \item First item
  \item Second item
  \item Third item with sub-items:
    \begin{itemize}
      \item Sub-item 1
      \item Sub-item 2
    \end{itemize}
\end{itemize}

\subsection{Numbered List}
\begin{enumerate}
  \item First step
  \item Second step
  \item Third step
\end{enumerate}

\section{Citations and References}
When you use the citation search feature, references will be automatically managed. For example, you might cite a paper like this: \cite{example2024}. All citations will be automatically added to your \texttt{cite.bib} file.

\section{Conclusion}
This template demonstrates basic \LaTeX{} features including mathematical formulas, tables, and lists. The bibliography system is now set up to work seamlessly with Bibby AI's citation management. Feel free to modify it for your needs!

\bibliographystyle{plain}
\bibliography{cite}

\end{document}\documentclass{article}
\usepackage[utf8]{inputenc}
\usepackage{booktabs}
\usepackage{amsmath}
\usepackage{natbib}

\title{My Document}
\author{Your Name}
\date{\today}

\begin{document}

\maketitle

\section{Introduction}
This is a sample document demonstrating various \LaTeX{} features including mathematical equations, tables, and formatting. You can cite papers using the citation search feature, and they will be automatically added to your \texttt{cite.bib} file.

\section{Mathematical Expressions}
Here are some example equations:

\subsection{Quadratic Formula}
The quadratic formula is given by:
\[x = \frac{-b \pm \sqrt{b^2 - 4ac}}{2a}\]

\subsection{Matrix Example}
A sample matrix:
\[
\begin{pmatrix}
a_{11} & a_{12} & a_{13} \\
a_{21} & a_{22} & a_{23} \\
a_{31} & a_{32} & a_{33}
\end{pmatrix}
\]

\section{Tables}
Below is a sample table showing data:

\begin{table}[h]
\centering
\begin{tabular}{@{}lrrr@{}}
\toprule
\textbf{Category} & \textbf{Value 1} & \textbf{Value 2} & \textbf{Total} \\
\midrule
Item A & 10.5 & 20.3 & 30.8 \\
Item B & 15.2 & 25.1 & 40.3 \\
Item C & 22.1 & 12.7 & 34.8 \\
\bottomrule
\end{tabular}
\caption{Sample Data Table}
\label{tab:sample}
\end{table}

\section{Lists}
\subsection{Bullet Points}
\begin{itemize}
  \item First item
  \item Second item
  \item Third item with sub-items:
    \begin{itemize}
      \item Sub-item 1
      \item Sub-item 2
    \end{itemize}
\end{itemize}

\subsection{Numbered List}
\begin{enumerate}
  \item First step
  \item Second step
  \item Third step
\end{enumerate}

\section{Citations and References}
When you use the citation search feature, references will be automatically managed. For example, you might cite a paper like this: \cite{example2024}. All citations will be automatically added to your \texttt{cite.bib} file.

\section{Conclusion}
This template demonstrates basic \LaTeX{} features including mathematical formulas, tables, and lists. The bibliography system is now set up to work seamlessly with Bibby AI's citation management. Feel free to modify it for your needs!

\bibliographystyle{plain}
\bibliography{cite}

\end{document}\documentclass{article}
\usepackage[utf8]{inputenc}
\usepackage{booktabs}
\usepackage{amsmath}
\usepackage{natbib}

\title{My Document}
\author{Your Name}
\date{\today}

\begin{document}

\maketitle

\section{Introduction}
This is a sample document demonstrating various \LaTeX{} features including mathematical equations, tables, and formatting. You can cite papers using the citation search feature, and they will be automatically added to your \texttt{cite.bib} file.

\section{Mathematical Expressions}
Here are some example equations:

\subsection{Quadratic Formula}
The quadratic formula is given by:
\[x = \frac{-b \pm \sqrt{b^2 - 4ac}}{2a}\]

\subsection{Matrix Example}
A sample matrix:
\[
\begin{pmatrix}
a_{11} & a_{12} & a_{13} \\
a_{21} & a_{22} & a_{23} \\
a_{31} & a_{32} & a_{33}
\end{pmatrix}
\]

\section{Tables}
Below is a sample table showing data:

\begin{table}[h]
\centering
\begin{tabular}{@{}lrrr@{}}
\toprule
\textbf{Category} & \textbf{Value 1} & \textbf{Value 2} & \textbf{Total} \\
\midrule
Item A & 10.5 & 20.3 & 30.8 \\
Item B & 15.2 & 25.1 & 40.3 \\
Item C & 22.1 & 12.7 & 34.8 \\
\bottomrule
\end{tabular}
\caption{Sample Data Table}
\label{tab:sample}
\end{table}

\section{Lists}
\subsection{Bullet Points}
\begin{itemize}
  \item First item
  \item Second item
  \item Third item with sub-items:
    \begin{itemize}
      \item Sub-item 1
      \item Sub-item 2
    \end{itemize}
\end{itemize}

\subsection{Numbered List}
\begin{enumerate}
  \item First step
  \item Second step
  \item Third step
\end{enumerate}

\section{Citations and References}
When you use the citation search feature, references will be automatically managed. For example, you might cite a paper like this: \cite{example2024}. All citations will be automatically added to your \texttt{cite.bib} file.

\section{Conclusion}
This template demonstrates basic \LaTeX{} features including mathematical formulas, tables, and lists. The bibliography system is now set up to work seamlessly with Bibby AI's citation management. Feel free to modify it for your needs!

\bibliographystyle{plain}
\bibliography{cite}

\end{document}\documentclass{article}
\usepackage[utf8]{inputenc}
\usepackage{booktabs}
\usepackage{amsmath}
\usepackage{natbib}

\title{My Document}
\author{Your Name}
\date{\today}

\begin{document}

\maketitle

\section{Introduction}
This is a sample document demonstrating various \LaTeX{} features including mathematical equations, tables, and formatting. You can cite papers using the citation search feature, and they will be automatically added to your \texttt{cite.bib} file.

\section{Mathematical Expressions}
Here are some example equations:

\subsection{Quadratic Formula}
The quadratic formula is given by:
\[x = \frac{-b \pm \sqrt{b^2 - 4ac}}{2a}\]

\subsection{Matrix Example}
A sample matrix:
\[
\begin{pmatrix}
a_{11} & a_{12} & a_{13} \\
a_{21} & a_{22} & a_{23} \\
a_{31} & a_{32} & a_{33}
\end{pmatrix}
\]

\section{Tables}
Below is a sample table showing data:

\begin{table}[h]
\centering
\begin{tabular}{@{}lrrr@{}}
\toprule
\textbf{Category} & \textbf{Value 1} & \textbf{Value 2} & \textbf{Total} \\
\midrule
Item A & 10.5 & 20.3 & 30.8 \\
Item B & 15.2 & 25.1 & 40.3 \\
Item C & 22.1 & 12.7 & 34.8 \\
\bottomrule
\end{tabular}
\caption{Sample Data Table}
\label{tab:sample}
\end{table}

\section{Lists}
\subsection{Bullet Points}
\begin{itemize}
  \item First item
  \item Second item
  \item Third item with sub-items:
    \begin{itemize}
      \item Sub-item 1
      \item Sub-item 2
    \end{itemize}
\end{itemize}

\subsection{Numbered List}
\begin{enumerate}
  \item First step
  \item Second step
  \item Third step
\end{enumerate}

\section{Citations and References}
When you use the citation search feature, references will be automatically managed. For example, you might cite a paper like this: \cite{example2024}. All citations will be automatically added to your \texttt{cite.bib} file.

\section{Conclusion}
This template demonstrates basic \LaTeX{} features including mathematical formulas, tables, and lists. The bibliography system is now set up to work seamlessly with Bibby AI's citation management. Feel free to modify it for your needs!

\bibliographystyle{plain}
\bibliography{cite}

\end{document}\documentclass{article}
\usepackage[utf8]{inputenc}
\usepackage{booktabs}
\usepackage{amsmath}
\usepackage{natbib}

\title{My Document}
\author{Your Name}
\date{\today}

\begin{document}

\maketitle

\section{Introduction}
This is a sample document demonstrating various \LaTeX{} features including mathematical equations, tables, and formatting. You can cite papers using the citation search feature, and they will be automatically added to your \texttt{cite.bib} file.

\section{Mathematical Expressions}
Here are some example equations:

\subsection{Quadratic Formula}
The quadratic formula is given by:
\[x = \frac{-b \pm \sqrt{b^2 - 4ac}}{2a}\]

\subsection{Matrix Example}
A sample matrix:
\[
\begin{pmatrix}
a_{11} & a_{12} & a_{13} \\
a_{21} & a_{22} & a_{23} \\
a_{31} & a_{32} & a_{33}
\end{pmatrix}
\]

\section{Tables}
Below is a sample table showing data:

\begin{table}[h]
\centering
\begin{tabular}{@{}lrrr@{}}
\toprule
\textbf{Category} & \textbf{Value 1} & \textbf{Value 2} & \textbf{Total} \\
\midrule
Item A & 10.5 & 20.3 & 30.8 \\
Item B & 15.2 & 25.1 & 40.3 \\
Item C & 22.1 & 12.7 & 34.8 \\
\bottomrule
\end{tabular}
\caption{Sample Data Table}
\label{tab:sample}
\end{table}

\section{Lists}
\subsection{Bullet Points}
\begin{itemize}
  \item First item
  \item Second item
  \item Third item with sub-items:
    \begin{itemize}
      \item Sub-item 1
      \item Sub-item 2
    \end{itemize}
\end{itemize}

\subsection{Numbered List}
\begin{enumerate}
  \item First step
  \item Second step
  \item Third step
\end{enumerate}

\section{Citations and References}
When you use the citation search feature, references will be automatically managed. For example, you might cite a paper like this: \cite{example2024}. All citations will be automatically added to your \texttt{cite.bib} file.

\section{Conclusion}
This template demonstrates basic \LaTeX{} features including mathematical formulas, tables, and lists. The bibliography system is now set up to work seamlessly with Bibby AI's citation management. Feel free to modify it for your needs!

\bibliographystyle{plain}
\bibliography{cite}

\end{document}\documentclass{article}
\usepackage[utf8]{inputenc}
\usepackage{booktabs}
\usepackage{amsmath}
\usepackage{natbib}

\title{My Document}
\author{Your Name}
\date{\today}

\begin{document}

\maketitle

\section{Introduction}
This is a sample document demonstrating various \LaTeX{} features including mathematical equations, tables, and formatting. You can cite papers using the citation search feature, and they will be automatically added to your \texttt{cite.bib} file.

\section{Mathematical Expressions}
Here are some example equations:

\subsection{Quadratic Formula}
The quadratic formula is given by:
\[x = \frac{-b \pm \sqrt{b^2 - 4ac}}{2a}\]

\subsection{Matrix Example}
A sample matrix:
\[
\begin{pmatrix}
a_{11} & a_{12} & a_{13} \\
a_{21} & a_{22} & a_{23} \\
a_{31} & a_{32} & a_{33}
\end{pmatrix}
\]

\section{Tables}
Below is a sample table showing data:

\begin{table}[h]
\centering
\begin{tabular}{@{}lrrr@{}}
\toprule
\textbf{Category} & \textbf{Value 1} & \textbf{Value 2} & \textbf{Total} \\
\midrule
Item A & 10.5 & 20.3 & 30.8 \\
Item B & 15.2 & 25.1 & 40.3 \\
Item C & 22.1 & 12.7 & 34.8 \\
\bottomrule
\end{tabular}
\caption{Sample Data Table}
\label{tab:sample}
\end{table}

\section{Lists}
\subsection{Bullet Points}
\begin{itemize}
  \item First item
  \item Second item
  \item Third item with sub-items:
    \begin{itemize}
      \item Sub-item 1
      \item Sub-item 2
    \end{itemize}
\end{itemize}

\subsection{Numbered List}
\begin{enumerate}
  \item First step
  \item Second step
  \item Third step
\end{enumerate}

\section{Citations and References}
When you use the citation search feature, references will be automatically managed. For example, you might cite a paper like this: \cite{example2024}. All citations will be automatically added to your \texttt{cite.bib} file.

\section{Conclusion}
This template demonstrates basic \LaTeX{} features including mathematical formulas, tables, and lists. The bibliography system is now set up to work seamlessly with Bibby AI's citation management. Feel free to modify it for your needs!

\bibliographystyle{plain}
\bibliography{cite}

\end{document}\documentclass{article}
\usepackage[utf8]{inputenc}
\usepackage{booktabs}
\usepackage{amsmath}
\usepackage{natbib}

\title{My Document}
\author{Your Name}
\date{\today}

\begin{document}

\maketitle

\section{Introduction}
This is a sample document demonstrating various \LaTeX{} features including mathematical equations, tables, and formatting. You can cite papers using the citation search feature, and they will be automatically added to your \texttt{cite.bib} file.

\section{Mathematical Expressions}
Here are some example equations:

\subsection{Quadratic Formula}
The quadratic formula is given by:
\[x = \frac{-b \pm \sqrt{b^2 - 4ac}}{2a}\]

\subsection{Matrix Example}
A sample matrix:
\[
\begin{pmatrix}
a_{11} & a_{12} & a_{13} \\
a_{21} & a_{22} & a_{23} \\
a_{31} & a_{32} & a_{33}
\end{pmatrix}
\]

\section{Tables}
Below is a sample table showing data:

\begin{table}[h]
\centering
\begin{tabular}{@{}lrrr@{}}
\toprule
\textbf{Category} & \textbf{Value 1} & \textbf{Value 2} & \textbf{Total} \\
\midrule
Item A & 10.5 & 20.3 & 30.8 \\
Item B & 15.2 & 25.1 & 40.3 \\
Item C & 22.1 & 12.7 & 34.8 \\
\bottomrule
\end{tabular}
\caption{Sample Data Table}
\label{tab:sample}
\end{table}

\section{Lists}
\subsection{Bullet Points}
\begin{itemize}
  \item First item
  \item Second item
  \item Third item with sub-items:
    \begin{itemize}
      \item Sub-item 1
      \item Sub-item 2
    \end{itemize}
\end{itemize}

\subsection{Numbered List}
\begin{enumerate}
  \item First step
  \item Second step
  \item Third step
\end{enumerate}

\section{Citations and References}
When you use the citation search feature, references will be automatically managed. For example, you might cite a paper like this: \cite{example2024}. All citations will be automatically added to your \texttt{cite.bib} file.

\section{Conclusion}
This template demonstrates basic \LaTeX{} features including mathematical formulas, tables, and lists. The bibliography system is now set up to work seamlessly with Bibby AI's citation management. Feel free to modify it for your needs!

\bibliographystyle{plain}
\bibliography{cite}

\end{document}\documentclass{article}
\usepackage[utf8]{inputenc}
\usepackage{booktabs}
\usepackage{amsmath}
\usepackage{natbib}

\title{My Document}
\author{Your Name}
\date{\today}

\begin{document}

\maketitle

\section{Introduction}
This is a sample document demonstrating various \LaTeX{} features including mathematical equations, tables, and formatting. You can cite papers using the citation search feature, and they will be automatically added to your \texttt{cite.bib} file.

\section{Mathematical Expressions}
Here are some example equations:

\subsection{Quadratic Formula}
The quadratic formula is given by:
\[x = \frac{-b \pm \sqrt{b^2 - 4ac}}{2a}\]

\subsection{Matrix Example}
A sample matrix:
\[
\begin{pmatrix}
a_{11} & a_{12} & a_{13} \\
a_{21} & a_{22} & a_{23} \\
a_{31} & a_{32} & a_{33}
\end{pmatrix}
\]

\section{Tables}
Below is a sample table showing data:

\begin{table}[h]
\centering
\begin{tabular}{@{}lrrr@{}}
\toprule
\textbf{Category} & \textbf{Value 1} & \textbf{Value 2} & \textbf{Total} \\
\midrule
Item A & 10.5 & 20.3 & 30.8 \\
Item B & 15.2 & 25.1 & 40.3 \\
Item C & 22.1 & 12.7 & 34.8 \\
\bottomrule
\end{tabular}
\caption{Sample Data Table}
\label{tab:sample}
\end{table}

\section{Lists}
\subsection{Bullet Points}
\begin{itemize}
  \item First item
  \item Second item
  \item Third item with sub-items:
    \begin{itemize}
      \item Sub-item 1
      \item Sub-item 2
    \end{itemize}
\end{itemize}

\subsection{Numbered List}
\begin{enumerate}
  \item First step
  \item Second step
  \item Third step
\end{enumerate}

\section{Citations and References}
When you use the citation search feature, references will be automatically managed. For example, you might cite a paper like this: \cite{example2024}. All citations will be automatically added to your \texttt{cite.bib} file.

\section{Conclusion}
This template demonstrates basic \LaTeX{} features including mathematical formulas, tables, and lists. The bibliography system is now set up to work seamlessly with Bibby AI's citation management. Feel free to modify it for your needs!

\bibliographystyle{plain}
\bibliography{cite}

\end{document}\documentclass{article}
\usepackage[utf8]{inputenc}
\usepackage{booktabs}
\usepackage{amsmath}
\usepackage{natbib}

\title{My Document}
\author{Your Name}
\date{\today}

\begin{document}

\maketitle

\section{Introduction}
This is a sample document demonstrating various \LaTeX{} features including mathematical equations, tables, and formatting. You can cite papers using the citation search feature, and they will be automatically added to your \texttt{cite.bib} file.

\section{Mathematical Expressions}
Here are some example equations:

\subsection{Quadratic Formula}
The quadratic formula is given by:
\[x = \frac{-b \pm \sqrt{b^2 - 4ac}}{2a}\]

\subsection{Matrix Example}
A sample matrix:
\[
\begin{pmatrix}
a_{11} & a_{12} & a_{13} \\
a_{21} & a_{22} & a_{23} \\
a_{31} & a_{32} & a_{33}
\end{pmatrix}
\]

\section{Tables}
Below is a sample table showing data:

\begin{table}[h]
\centering
\begin{tabular}{@{}lrrr@{}}
\toprule
\textbf{Category} & \textbf{Value 1} & \textbf{Value 2} & \textbf{Total} \\
\midrule
Item A & 10.5 & 20.3 & 30.8 \\
Item B & 15.2 & 25.1 & 40.3 \\
Item C & 22.1 & 12.7 & 34.8 \\
\bottomrule
\end{tabular}
\caption{Sample Data Table}
\label{tab:sample}
\end{table}

\section{Lists}
\subsection{Bullet Points}
\begin{itemize}
  \item First item
  \item Second item
  \item Third item with sub-items:
    \begin{itemize}
      \item Sub-item 1
      \item Sub-item 2
    \end{itemize}
\end{itemize}

\subsection{Numbered List}
\begin{enumerate}
  \item First step
  \item Second step
  \item Third step
\end{enumerate}

\section{Citations and References}
When you use the citation search feature, references will be automatically managed. For example, you might cite a paper like this: \cite{example2024}. All citations will be automatically added to your \texttt{cite.bib} file.

\section{Conclusion}
This template demonstrates basic \LaTeX{} features including mathematical formulas, tables, and lists. The bibliography system is now set up to work seamlessly with Bibby AI's citation management. Feel free to modify it for your needs!

\bibliographystyle{plain}
\bibliography{cite}

\end{document}\documentclass{article}
\usepackage[utf8]{inputenc}
\usepackage{booktabs}
\usepackage{amsmath}
\usepackage{natbib}

\title{My Document}
\author{Your Name}
\date{\today}

\begin{document}

\maketitle

\section{Introduction}
This is a sample document demonstrating various \LaTeX{} features including mathematical equations, tables, and formatting. You can cite papers using the citation search feature, and they will be automatically added to your \texttt{cite.bib} file.

\section{Mathematical Expressions}
Here are some example equations:

\subsection{Quadratic Formula}
The quadratic formula is given by:
\[x = \frac{-b \pm \sqrt{b^2 - 4ac}}{2a}\]

\subsection{Matrix Example}
A sample matrix:
\[
\begin{pmatrix}
a_{11} & a_{12} & a_{13} \\
a_{21} & a_{22} & a_{23} \\
a_{31} & a_{32} & a_{33}
\end{pmatrix}
\]

\section{Tables}
Below is a sample table showing data:

\begin{table}[h]
\centering
\begin{tabular}{@{}lrrr@{}}
\toprule
\textbf{Category} & \textbf{Value 1} & \textbf{Value 2} & \textbf{Total} \\
\midrule
Item A & 10.5 & 20.3 & 30.8 \\
Item B & 15.2 & 25.1 & 40.3 \\
Item C & 22.1 & 12.7 & 34.8 \\
\bottomrule
\end{tabular}
\caption{Sample Data Table}
\label{tab:sample}
\end{table}

\section{Lists}
\subsection{Bullet Points}
\begin{itemize}
  \item First item
  \item Second item
  \item Third item with sub-items:
    \begin{itemize}
      \item Sub-item 1
      \item Sub-item 2
    \end{itemize}
\end{itemize}

\subsection{Numbered List}
\begin{enumerate}
  \item First step
  \item Second step
  \item Third step
\end{enumerate}

\section{Citations and References}
When you use the citation search feature, references will be automatically managed. For example, you might cite a paper like this: \cite{example2024}. All citations will be automatically added to your \texttt{cite.bib} file.

\section{Conclusion}
This template demonstrates basic \LaTeX{} features including mathematical formulas, tables, and lists. The bibliography system is now set up to work seamlessly with Bibby AI's citation management. Feel free to modify it for your needs!

\bibliographystyle{plain}
\bibliography{cite}

\end{document}\documentclass{article}
\usepackage[utf8]{inputenc}
\usepackage{booktabs}
\usepackage{amsmath}
\usepackage{natbib}

\title{My Document}
\author{Your Name}
\date{\today}

\begin{document}

\maketitle

\section{Introduction}
This is a sample document demonstrating various \LaTeX{} features including mathematical equations, tables, and formatting. You can cite papers using the citation search feature, and they will be automatically added to your \texttt{cite.bib} file.

\section{Mathematical Expressions}
Here are some example equations:

\subsection{Quadratic Formula}
The quadratic formula is given by:
\[x = \frac{-b \pm \sqrt{b^2 - 4ac}}{2a}\]

\subsection{Matrix Example}
A sample matrix:
\[
\begin{pmatrix}
a_{11} & a_{12} & a_{13} \\
a_{21} & a_{22} & a_{23} \\
a_{31} & a_{32} & a_{33}
\end{pmatrix}
\]

\section{Tables}
Below is a sample table showing data:

\begin{table}[h]
\centering
\begin{tabular}{@{}lrrr@{}}
\toprule
\textbf{Category} & \textbf{Value 1} & \textbf{Value 2} & \textbf{Total} \\
\midrule
Item A & 10.5 & 20.3 & 30.8 \\
Item B & 15.2 & 25.1 & 40.3 \\
Item C & 22.1 & 12.7 & 34.8 \\
\bottomrule
\end{tabular}
\caption{Sample Data Table}
\label{tab:sample}
\end{table}

\section{Lists}
\subsection{Bullet Points}
\begin{itemize}
  \item First item
  \item Second item
  \item Third item with sub-items:
    \begin{itemize}
      \item Sub-item 1
      \item Sub-item 2
    \end{itemize}
\end{itemize}

\subsection{Numbered List}
\begin{enumerate}
  \item First step
  \item Second step
  \item Third step
\end{enumerate}

\section{Citations and References}
When you use the citation search feature, references will be automatically managed. For example, you might cite a paper like this: \cite{example2024}. All citations will be automatically added to your \texttt{cite.bib} file.

\section{Conclusion}
This template demonstrates basic \LaTeX{} features including mathematical formulas, tables, and lists. The bibliography system is now set up to work seamlessly with Bibby AI's citation management. Feel free to modify it for your needs!

\bibliographystyle{plain}
\bibliography{cite}

\end{document}\documentclass{article}
\usepackage[utf8]{inputenc}
\usepackage{booktabs}
\usepackage{amsmath}
\usepackage{natbib}

\title{My Document}
\author{Your Name}
\date{\today}

\begin{document}

\maketitle

\section{Introduction}
This is a sample document demonstrating various \LaTeX{} features including mathematical equations, tables, and formatting. You can cite papers using the citation search feature, and they will be automatically added to your \texttt{cite.bib} file.

\section{Mathematical Expressions}
Here are some example equations:

\subsection{Quadratic Formula}
The quadratic formula is given by:
\[x = \frac{-b \pm \sqrt{b^2 - 4ac}}{2a}\]

\subsection{Matrix Example}
A sample matrix:
\[
\begin{pmatrix}
a_{11} & a_{12} & a_{13} \\
a_{21} & a_{22} & a_{23} \\
a_{31} & a_{32} & a_{33}
\end{pmatrix}
\]

\section{Tables}
Below is a sample table showing data:

\begin{table}[h]
\centering
\begin{tabular}{@{}lrrr@{}}
\toprule
\textbf{Category} & \textbf{Value 1} & \textbf{Value 2} & \textbf{Total} \\
\midrule
Item A & 10.5 & 20.3 & 30.8 \\
Item B & 15.2 & 25.1 & 40.3 \\
Item C & 22.1 & 12.7 & 34.8 \\
\bottomrule
\end{tabular}
\caption{Sample Data Table}
\label{tab:sample}
\end{table}

\section{Lists}
\subsection{Bullet Points}
\begin{itemize}
  \item First item
  \item Second item
  \item Third item with sub-items:
    \begin{itemize}
      \item Sub-item 1
      \item Sub-item 2
    \end{itemize}
\end{itemize}

\subsection{Numbered List}
\begin{enumerate}
  \item First step
  \item Second step
  \item Third step
\end{enumerate}

\section{Citations and References}
When you use the citation search feature, references will be automatically managed. For example, you might cite a paper like this: \cite{example2024}. All citations will be automatically added to your \texttt{cite.bib} file.

\section{Conclusion}
This template demonstrates basic \LaTeX{} features including mathematical formulas, tables, and lists. The bibliography system is now set up to work seamlessly with Bibby AI's citation management. Feel free to modify it for your needs!

\bibliographystyle{plain}
\bibliography{cite}

\end{document}\documentclass{article}
\usepackage[utf8]{inputenc}
\usepackage{booktabs}
\usepackage{amsmath}
\usepackage{natbib}

\title{My Document}
\author{Your Name}
\date{\today}

\begin{document}

\maketitle

\section{Introduction}
This is a sample document demonstrating various \LaTeX{} features including mathematical equations, tables, and formatting. You can cite papers using the citation search feature, and they will be automatically added to your \texttt{cite.bib} file.

\section{Mathematical Expressions}
Here are some example equations:

\subsection{Quadratic Formula}
The quadratic formula is given by:
\[x = \frac{-b \pm \sqrt{b^2 - 4ac}}{2a}\]

\subsection{Matrix Example}
A sample matrix:
\[
\begin{pmatrix}
a_{11} & a_{12} & a_{13} \\
a_{21} & a_{22} & a_{23} \\
a_{31} & a_{32} & a_{33}
\end{pmatrix}
\]

\section{Tables}
Below is a sample table showing data:

\begin{table}[h]
\centering
\begin{tabular}{@{}lrrr@{}}
\toprule
\textbf{Category} & \textbf{Value 1} & \textbf{Value 2} & \textbf{Total} \\
\midrule
Item A & 10.5 & 20.3 & 30.8 \\
Item B & 15.2 & 25.1 & 40.3 \\
Item C & 22.1 & 12.7 & 34.8 \\
\bottomrule
\end{tabular}
\caption{Sample Data Table}
\label{tab:sample}
\end{table}

\section{Lists}
\subsection{Bullet Points}
\begin{itemize}
  \item First item
  \item Second item
  \item Third item with sub-items:
    \begin{itemize}
      \item Sub-item 1
      \item Sub-item 2
    \end{itemize}
\end{itemize}

\subsection{Numbered List}
\begin{enumerate}
  \item First step
  \item Second step
  \item Third step
\end{enumerate}

\section{Citations and References}
When you use the citation search feature, references will be automatically managed. For example, you might cite a paper like this: \cite{example2024}. All citations will be automatically added to your \texttt{cite.bib} file.

\section{Conclusion}
This template demonstrates basic \LaTeX{} features including mathematical formulas, tables, and lists. The bibliography system is now set up to work seamlessly with Bibby AI's citation management. Feel free to modify it for your needs!

\bibliographystyle{plain}
\bibliography{cite}

\end{document}\documentclass{article}
\usepackage[utf8]{inputenc}
\usepackage{booktabs}
\usepackage{amsmath}
\usepackage{natbib}

\title{My Document}
\author{Your Name}
\date{\today}

\begin{document}

\maketitle

\section{Introduction}
This is a sample document demonstrating various \LaTeX{} features including mathematical equations, tables, and formatting. You can cite papers using the citation search feature, and they will be automatically added to your \texttt{cite.bib} file.

\section{Mathematical Expressions}
Here are some example equations:

\subsection{Quadratic Formula}
The quadratic formula is given by:
\[x = \frac{-b \pm \sqrt{b^2 - 4ac}}{2a}\]

\subsection{Matrix Example}
A sample matrix:
\[
\begin{pmatrix}
a_{11} & a_{12} & a_{13} \\
a_{21} & a_{22} & a_{23} \\
a_{31} & a_{32} & a_{33}
\end{pmatrix}
\]

\section{Tables}
Below is a sample table showing data:

\begin{table}[h]
\centering
\begin{tabular}{@{}lrrr@{}}
\toprule
\textbf{Category} & \textbf{Value 1} & \textbf{Value 2} & \textbf{Total} \\
\midrule
Item A & 10.5 & 20.3 & 30.8 \\
Item B & 15.2 & 25.1 & 40.3 \\
Item C & 22.1 & 12.7 & 34.8 \\
\bottomrule
\end{tabular}
\caption{Sample Data Table}
\label{tab:sample}
\end{table}

\section{Lists}
\subsection{Bullet Points}
\begin{itemize}
  \item First item
  \item Second item
  \item Third item with sub-items:
    \begin{itemize}
      \item Sub-item 1
      \item Sub-item 2
    \end{itemize}
\end{itemize}

\subsection{Numbered List}
\begin{enumerate}
  \item First step
  \item Second step
  \item Third step
\end{enumerate}

\section{Citations and References}
When you use the citation search feature, references will be automatically managed. For example, you might cite a paper like this: \cite{example2024}. All citations will be automatically added to your \texttt{cite.bib} file.

\section{Conclusion}
This template demonstrates basic \LaTeX{} features including mathematical formulas, tables, and lists. The bibliography system is now set up to work seamlessly with Bibby AI's citation management. Feel free to modify it for your needs!

\bibliographystyle{plain}
\bibliography{cite}

\end{document}\documentclass{article}
\usepackage[utf8]{inputenc}
\usepackage{booktabs}
\usepackage{amsmath}
\usepackage{natbib}

\title{My Document}
\author{Your Name}
\date{\today}

\begin{document}

\maketitle

\section{Introduction}
This is a sample document demonstrating various \LaTeX{} features including mathematical equations, tables, and formatting. You can cite papers using the citation search feature, and they will be automatically added to your \texttt{cite.bib} file.

\section{Mathematical Expressions}
Here are some example equations:

\subsection{Quadratic Formula}
The quadratic formula is given by:
\[x = \frac{-b \pm \sqrt{b^2 - 4ac}}{2a}\]

\subsection{Matrix Example}
A sample matrix:
\[
\begin{pmatrix}
a_{11} & a_{12} & a_{13} \\
a_{21} & a_{22} & a_{23} \\
a_{31} & a_{32} & a_{33}
\end{pmatrix}
\]

\section{Tables}
Below is a sample table showing data:

\begin{table}[h]
\centering
\begin{tabular}{@{}lrrr@{}}
\toprule
\textbf{Category} & \textbf{Value 1} & \textbf{Value 2} & \textbf{Total} \\
\midrule
Item A & 10.5 & 20.3 & 30.8 \\
Item B & 15.2 & 25.1 & 40.3 \\
Item C & 22.1 & 12.7 & 34.8 \\
\bottomrule
\end{tabular}
\caption{Sample Data Table}
\label{tab:sample}
\end{table}

\section{Lists}
\subsection{Bullet Points}
\begin{itemize}
  \item First item
  \item Second item
  \item Third item with sub-items:
    \begin{itemize}
      \item Sub-item 1
      \item Sub-item 2
    \end{itemize}
\end{itemize}

\subsection{Numbered List}
\begin{enumerate}
  \item First step
  \item Second step
  \item Third step
\end{enumerate}

\section{Citations and References}
When you use the citation search feature, references will be automatically managed. For example, you might cite a paper like this: \cite{example2024}. All citations will be automatically added to your \texttt{cite.bib} file.

\section{Conclusion}
This template demonstrates basic \LaTeX{} features including mathematical formulas, tables, and lists. The bibliography system is now set up to work seamlessly with Bibby AI's citation management. Feel free to modify it for your needs!

\bibliographystyle{plain}
\bibliography{cite}

\end{document}\documentclass{article}
\usepackage[utf8]{inputenc}
\usepackage{booktabs}
\usepackage{amsmath}
\usepackage{natbib}

\title{My Document}
\author{Your Name}
\date{\today}

\begin{document}

\maketitle

\section{Introduction}
This is a sample document demonstrating various \LaTeX{} features including mathematical equations, tables, and formatting. You can cite papers using the citation search feature, and they will be automatically added to your \texttt{cite.bib} file.

\section{Mathematical Expressions}
Here are some example equations:

\subsection{Quadratic Formula}
The quadratic formula is given by:
\[x = \frac{-b \pm \sqrt{b^2 - 4ac}}{2a}\]

\subsection{Matrix Example}
A sample matrix:
\[
\begin{pmatrix}
a_{11} & a_{12} & a_{13} \\
a_{21} & a_{22} & a_{23} \\
a_{31} & a_{32} & a_{33}
\end{pmatrix}
\]

\section{Tables}
Below is a sample table showing data:

\begin{table}[h]
\centering
\begin{tabular}{@{}lrrr@{}}
\toprule
\textbf{Category} & \textbf{Value 1} & \textbf{Value 2} & \textbf{Total} \\
\midrule
Item A & 10.5 & 20.3 & 30.8 \\
Item B & 15.2 & 25.1 & 40.3 \\
Item C & 22.1 & 12.7 & 34.8 \\
\bottomrule
\end{tabular}
\caption{Sample Data Table}
\label{tab:sample}
\end{table}

\section{Lists}
\subsection{Bullet Points}
\begin{itemize}
  \item First item
  \item Second item
  \item Third item with sub-items:
    \begin{itemize}
      \item Sub-item 1
      \item Sub-item 2
    \end{itemize}
\end{itemize}

\subsection{Numbered List}
\begin{enumerate}
  \item First step
  \item Second step
  \item Third step
\end{enumerate}

\section{Citations and References}
When you use the citation search feature, references will be automatically managed. For example, you might cite a paper like this: \cite{example2024}. All citations will be automatically added to your \texttt{cite.bib} file.

\section{Conclusion}
This template demonstrates basic \LaTeX{} features including mathematical formulas, tables, and lists. The bibliography system is now set up to work seamlessly with Bibby AI's citation management. Feel free to modify it for your needs!

\bibliographystyle{plain}
\bibliography{cite}

\end{document}\documentclass{article}
\usepackage[utf8]{inputenc}
\usepackage{booktabs}
\usepackage{amsmath}
\usepackage{natbib}

\title{My Document}
\author{Your Name}
\date{\today}

\begin{document}

\maketitle

\section{Introduction}
This is a sample document demonstrating various \LaTeX{} features including mathematical equations, tables, and formatting. You can cite papers using the citation search feature, and they will be automatically added to your \texttt{cite.bib} file.

\section{Mathematical Expressions}
Here are some example equations:

\subsection{Quadratic Formula}
The quadratic formula is given by:
\[x = \frac{-b \pm \sqrt{b^2 - 4ac}}{2a}\]

\subsection{Matrix Example}
A sample matrix:
\[
\begin{pmatrix}
a_{11} & a_{12} & a_{13} \\
a_{21} & a_{22} & a_{23} \\
a_{31} & a_{32} & a_{33}
\end{pmatrix}
\]

\section{Tables}
Below is a sample table showing data:

\begin{table}[h]
\centering
\begin{tabular}{@{}lrrr@{}}
\toprule
\textbf{Category} & \textbf{Value 1} & \textbf{Value 2} & \textbf{Total} \\
\midrule
Item A & 10.5 & 20.3 & 30.8 \\
Item B & 15.2 & 25.1 & 40.3 \\
Item C & 22.1 & 12.7 & 34.8 \\
\bottomrule
\end{tabular}
\caption{Sample Data Table}
\label{tab:sample}
\end{table}

\section{Lists}
\subsection{Bullet Points}
\begin{itemize}
  \item First item
  \item Second item
  \item Third item with sub-items:
    \begin{itemize}
      \item Sub-item 1
      \item Sub-item 2
    \end{itemize}
\end{itemize}

\subsection{Numbered List}
\begin{enumerate}
  \item First step
  \item Second step
  \item Third step
\end{enumerate}

\section{Citations and References}
When you use the citation search feature, references will be automatically managed. For example, you might cite a paper like this: \cite{example2024}. All citations will be automatically added to your \texttt{cite.bib} file.

\section{Conclusion}
This template demonstrates basic \LaTeX{} features including mathematical formulas, tables, and lists. The bibliography system is now set up to work seamlessly with Bibby AI's citation management. Feel free to modify it for your needs!

\bibliographystyle{plain}
\bibliography{cite}

\end{document}\documentclass{article}
\usepackage[utf8]{inputenc}
\usepackage{booktabs}
\usepackage{amsmath}
\usepackage{natbib}

\title{My Document}
\author{Your Name}
\date{\today}

\begin{document}

\maketitle

\section{Introduction}
This is a sample document demonstrating various \LaTeX{} features including mathematical equations, tables, and formatting. You can cite papers using the citation search feature, and they will be automatically added to your \texttt{cite.bib} file.

\section{Mathematical Expressions}
Here are some example equations:

\subsection{Quadratic Formula}
The quadratic formula is given by:
\[x = \frac{-b \pm \sqrt{b^2 - 4ac}}{2a}\]

\subsection{Matrix Example}
A sample matrix:
\[
\begin{pmatrix}
a_{11} & a_{12} & a_{13} \\
a_{21} & a_{22} & a_{23} \\
a_{31} & a_{32} & a_{33}
\end{pmatrix}
\]

\section{Tables}
Below is a sample table showing data:

\begin{table}[h]
\centering
\begin{tabular}{@{}lrrr@{}}
\toprule
\textbf{Category} & \textbf{Value 1} & \textbf{Value 2} & \textbf{Total} \\
\midrule
Item A & 10.5 & 20.3 & 30.8 \\
Item B & 15.2 & 25.1 & 40.3 \\
Item C & 22.1 & 12.7 & 34.8 \\
\bottomrule
\end{tabular}
\caption{Sample Data Table}
\label{tab:sample}
\end{table}

\section{Lists}
\subsection{Bullet Points}
\begin{itemize}
  \item First item
  \item Second item
  \item Third item with sub-items:
    \begin{itemize}
      \item Sub-item 1
      \item Sub-item 2
    \end{itemize}
\end{itemize}

\subsection{Numbered List}
\begin{enumerate}
  \item First step
  \item Second step
  \item Third step
\end{enumerate}

\section{Citations and References}
When you use the citation search feature, references will be automatically managed. For example, you might cite a paper like this: \cite{example2024}. All citations will be automatically added to your \texttt{cite.bib} file.

\section{Conclusion}
This template demonstrates basic \LaTeX{} features including mathematical formulas, tables, and lists. The bibliography system is now set up to work seamlessly with Bibby AI's citation management. Feel free to modify it for your needs!

\bibliographystyle{plain}
\bibliography{cite}

\end{document}\documentclass{article}
\usepackage[utf8]{inputenc}
\usepackage{booktabs}
\usepackage{amsmath}
\usepackage{natbib}

\title{My Document}
\author{Your Name}
\date{\today}

\begin{document}

\maketitle

\section{Introduction}
This is a sample document demonstrating various \LaTeX{} features including mathematical equations, tables, and formatting. You can cite papers using the citation search feature, and they will be automatically added to your \texttt{cite.bib} file.

\section{Mathematical Expressions}
Here are some example equations:

\subsection{Quadratic Formula}
The quadratic formula is given by:
\[x = \frac{-b \pm \sqrt{b^2 - 4ac}}{2a}\]

\subsection{Matrix Example}
A sample matrix:
\[
\begin{pmatrix}
a_{11} & a_{12} & a_{13} \\
a_{21} & a_{22} & a_{23} \\
a_{31} & a_{32} & a_{33}
\end{pmatrix}
\]

\section{Tables}
Below is a sample table showing data:

\begin{table}[h]
\centering
\begin{tabular}{@{}lrrr@{}}
\toprule
\textbf{Category} & \textbf{Value 1} & \textbf{Value 2} & \textbf{Total} \\
\midrule
Item A & 10.5 & 20.3 & 30.8 \\
Item B & 15.2 & 25.1 & 40.3 \\
Item C & 22.1 & 12.7 & 34.8 \\
\bottomrule
\end{tabular}
\caption{Sample Data Table}
\label{tab:sample}
\end{table}

\section{Lists}
\subsection{Bullet Points}
\begin{itemize}
  \item First item
  \item Second item
  \item Third item with sub-items:
    \begin{itemize}
      \item Sub-item 1
      \item Sub-item 2
    \end{itemize}
\end{itemize}

\subsection{Numbered List}
\begin{enumerate}
  \item First step
  \item Second step
  \item Third step
\end{enumerate}

\section{Citations and References}
When you use the citation search feature, references will be automatically managed. For example, you might cite a paper like this: \cite{example2024}. All citations will be automatically added to your \texttt{cite.bib} file.

\section{Conclusion}
This template demonstrates basic \LaTeX{} features including mathematical formulas, tables, and lists. The bibliography system is now set up to work seamlessly with Bibby AI's citation management. Feel free to modify it for your needs!

\bibliographystyle{plain}
\bibliography{cite}

\end{document}\documentclass{article}
\usepackage[utf8]{inputenc}
\usepackage{booktabs}
\usepackage{amsmath}
\usepackage{natbib}

\title{My Document}
\author{Your Name}
\date{\today}

\begin{document}

\maketitle

\section{Introduction}
This is a sample document demonstrating various \LaTeX{} features including mathematical equations, tables, and formatting. You can cite papers using the citation search feature, and they will be automatically added to your \texttt{cite.bib} file.

\section{Mathematical Expressions}
Here are some example equations:

\subsection{Quadratic Formula}
The quadratic formula is given by:
\[x = \frac{-b \pm \sqrt{b^2 - 4ac}}{2a}\]

\subsection{Matrix Example}
A sample matrix:
\[
\begin{pmatrix}
a_{11} & a_{12} & a_{13} \\
a_{21} & a_{22} & a_{23} \\
a_{31} & a_{32} & a_{33}
\end{pmatrix}
\]

\section{Tables}
Below is a sample table showing data:

\begin{table}[h]
\centering
\begin{tabular}{@{}lrrr@{}}
\toprule
\textbf{Category} & \textbf{Value 1} & \textbf{Value 2} & \textbf{Total} \\
\midrule
Item A & 10.5 & 20.3 & 30.8 \\
Item B & 15.2 & 25.1 & 40.3 \\
Item C & 22.1 & 12.7 & 34.8 \\
\bottomrule
\end{tabular}
\caption{Sample Data Table}
\label{tab:sample}
\end{table}

\section{Lists}
\subsection{Bullet Points}
\begin{itemize}
  \item First item
  \item Second item
  \item Third item with sub-items:
    \begin{itemize}
      \item Sub-item 1
      \item Sub-item 2
    \end{itemize}
\end{itemize}

\subsection{Numbered List}
\begin{enumerate}
  \item First step
  \item Second step
  \item Third step
\end{enumerate}

\section{Citations and References}
When you use the citation search feature, references will be automatically managed. For example, you might cite a paper like this: \cite{example2024}. All citations will be automatically added to your \texttt{cite.bib} file.

\section{Conclusion}
This template demonstrates basic \LaTeX{} features including mathematical formulas, tables, and lists. The bibliography system is now set up to work seamlessly with Bibby AI's citation management. Feel free to modify it for your needs!

\bibliographystyle{plain}
\bibliography{cite}

\end{document}\documentclass{article}
\usepackage[utf8]{inputenc}
\usepackage{booktabs}
\usepackage{amsmath}
\usepackage{natbib}

\title{My Document}
\author{Your Name}
\date{\today}

\begin{document}

\maketitle

\section{Introduction}
This is a sample document demonstrating various \LaTeX{} features including mathematical equations, tables, and formatting. You can cite papers using the citation search feature, and they will be automatically added to your \texttt{cite.bib} file.

\section{Mathematical Expressions}
Here are some example equations:

\subsection{Quadratic Formula}
The quadratic formula is given by:
\[x = \frac{-b \pm \sqrt{b^2 - 4ac}}{2a}\]

\subsection{Matrix Example}
A sample matrix:
\[
\begin{pmatrix}
a_{11} & a_{12} & a_{13} \\
a_{21} & a_{22} & a_{23} \\
a_{31} & a_{32} & a_{33}
\end{pmatrix}
\]

\section{Tables}
Below is a sample table showing data:

\begin{table}[h]
\centering
\begin{tabular}{@{}lrrr@{}}
\toprule
\textbf{Category} & \textbf{Value 1} & \textbf{Value 2} & \textbf{Total} \\
\midrule
Item A & 10.5 & 20.3 & 30.8 \\
Item B & 15.2 & 25.1 & 40.3 \\
Item C & 22.1 & 12.7 & 34.8 \\
\bottomrule
\end{tabular}
\caption{Sample Data Table}
\label{tab:sample}
\end{table}

\section{Lists}
\subsection{Bullet Points}
\begin{itemize}
  \item First item
  \item Second item
  \item Third item with sub-items:
    \begin{itemize}
      \item Sub-item 1
      \item Sub-item 2
    \end{itemize}
\end{itemize}

\subsection{Numbered List}
\begin{enumerate}
  \item First step
  \item Second step
  \item Third step
\end{enumerate}

\section{Citations and References}
When you use the citation search feature, references will be automatically managed. For example, you might cite a paper like this: \cite{example2024}. All citations will be automatically added to your \texttt{cite.bib} file.

\section{Conclusion}
This template demonstrates basic \LaTeX{} features including mathematical formulas, tables, and lists. The bibliography system is now set up to work seamlessly with Bibby AI's citation management. Feel free to modify it for your needs!

\bibliographystyle{plain}
\bibliography{cite}

\end{document}\documentclass{article}
\usepackage[utf8]{inputenc}
\usepackage{booktabs}
\usepackage{amsmath}
\usepackage{natbib}

\title{My Document}
\author{Your Name}
\date{\today}

\begin{document}

\maketitle

\section{Introduction}
This is a sample document demonstrating various \LaTeX{} features including mathematical equations, tables, and formatting. You can cite papers using the citation search feature, and they will be automatically added to your \texttt{cite.bib} file.

\section{Mathematical Expressions}
Here are some example equations:

\subsection{Quadratic Formula}
The quadratic formula is given by:
\[x = \frac{-b \pm \sqrt{b^2 - 4ac}}{2a}\]

\subsection{Matrix Example}
A sample matrix:
\[
\begin{pmatrix}
a_{11} & a_{12} & a_{13} \\
a_{21} & a_{22} & a_{23} \\
a_{31} & a_{32} & a_{33}
\end{pmatrix}
\]

\section{Tables}
Below is a sample table showing data:

\begin{table}[h]
\centering
\begin{tabular}{@{}lrrr@{}}
\toprule
\textbf{Category} & \textbf{Value 1} & \textbf{Value 2} & \textbf{Total} \\
\midrule
Item A & 10.5 & 20.3 & 30.8 \\
Item B & 15.2 & 25.1 & 40.3 \\
Item C & 22.1 & 12.7 & 34.8 \\
\bottomrule
\end{tabular}
\caption{Sample Data Table}
\label{tab:sample}
\end{table}

\section{Lists}
\subsection{Bullet Points}
\begin{itemize}
  \item First item
  \item Second item
  \item Third item with sub-items:
    \begin{itemize}
      \item Sub-item 1
      \item Sub-item 2
    \end{itemize}
\end{itemize}

\subsection{Numbered List}
\begin{enumerate}
  \item First step
  \item Second step
  \item Third step
\end{enumerate}

\section{Citations and References}
When you use the citation search feature, references will be automatically managed. For example, you might cite a paper like this: \cite{example2024}. All citations will be automatically added to your \texttt{cite.bib} file.

\section{Conclusion}
This template demonstrates basic \LaTeX{} features including mathematical formulas, tables, and lists. The bibliography system is now set up to work seamlessly with Bibby AI's citation management. Feel free to modify it for your needs!

\bibliographystyle{plain}
\bibliography{cite}

\end{document}\documentclass{article}
\usepackage[utf8]{inputenc}
\usepackage{booktabs}
\usepackage{amsmath}
\usepackage{natbib}

\title{My Document}
\author{Your Name}
\date{\today}

\begin{document}

\maketitle

\section{Introduction}
This is a sample document demonstrating various \LaTeX{} features including mathematical equations, tables, and formatting. You can cite papers using the citation search feature, and they will be automatically added to your \texttt{cite.bib} file.

\section{Mathematical Expressions}
Here are some example equations:

\subsection{Quadratic Formula}
The quadratic formula is given by:
\[x = \frac{-b \pm \sqrt{b^2 - 4ac}}{2a}\]

\subsection{Matrix Example}
A sample matrix:
\[
\begin{pmatrix}
a_{11} & a_{12} & a_{13} \\
a_{21} & a_{22} & a_{23} \\
a_{31} & a_{32} & a_{33}
\end{pmatrix}
\]

\section{Tables}
Below is a sample table showing data:

\begin{table}[h]
\centering
\begin{tabular}{@{}lrrr@{}}
\toprule
\textbf{Category} & \textbf{Value 1} & \textbf{Value 2} & \textbf{Total} \\
\midrule
Item A & 10.5 & 20.3 & 30.8 \\
Item B & 15.2 & 25.1 & 40.3 \\
Item C & 22.1 & 12.7 & 34.8 \\
\bottomrule
\end{tabular}
\caption{Sample Data Table}
\label{tab:sample}
\end{table}

\section{Lists}
\subsection{Bullet Points}
\begin{itemize}
  \item First item
  \item Second item
  \item Third item with sub-items:
    \begin{itemize}
      \item Sub-item 1
      \item Sub-item 2
    \end{itemize}
\end{itemize}

\subsection{Numbered List}
\begin{enumerate}
  \item First step
  \item Second step
  \item Third step
\end{enumerate}

\section{Citations and References}
When you use the citation search feature, references will be automatically managed. For example, you might cite a paper like this: \cite{example2024}. All citations will be automatically added to your \texttt{cite.bib} file.

\section{Conclusion}
This template demonstrates basic \LaTeX{} features including mathematical formulas, tables, and lists. The bibliography system is now set up to work seamlessly with Bibby AI's citation management. Feel free to modify it for your needs!

\bibliographystyle{plain}
\bibliography{cite}

\end{document}\documentclass{article}
\usepackage[utf8]{inputenc}
\usepackage{booktabs}
\usepackage{amsmath}
\usepackage{natbib}

\title{My Document}
\author{Your Name}
\date{\today}

\begin{document}

\maketitle

\section{Introduction}
This is a sample document demonstrating various \LaTeX{} features including mathematical equations, tables, and formatting. You can cite papers using the citation search feature, and they will be automatically added to your \texttt{cite.bib} file.

\section{Mathematical Expressions}
Here are some example equations:

\subsection{Quadratic Formula}
The quadratic formula is given by:
\[x = \frac{-b \pm \sqrt{b^2 - 4ac}}{2a}\]

\subsection{Matrix Example}
A sample matrix:
\[
\begin{pmatrix}
a_{11} & a_{12} & a_{13} \\
a_{21} & a_{22} & a_{23} \\
a_{31} & a_{32} & a_{33}
\end{pmatrix}
\]

\section{Tables}
Below is a sample table showing data:

\begin{table}[h]
\centering
\begin{tabular}{@{}lrrr@{}}
\toprule
\textbf{Category} & \textbf{Value 1} & \textbf{Value 2} & \textbf{Total} \\
\midrule
Item A & 10.5 & 20.3 & 30.8 \\
Item B & 15.2 & 25.1 & 40.3 \\
Item C & 22.1 & 12.7 & 34.8 \\
\bottomrule
\end{tabular}
\caption{Sample Data Table}
\label{tab:sample}
\end{table}

\section{Lists}
\subsection{Bullet Points}
\begin{itemize}
  \item First item
  \item Second item
  \item Third item with sub-items:
    \begin{itemize}
      \item Sub-item 1
      \item Sub-item 2
    \end{itemize}
\end{itemize}

\subsection{Numbered List}
\begin{enumerate}
  \item First step
  \item Second step
  \item Third step
\end{enumerate}

\section{Citations and References}
When you use the citation search feature, references will be automatically managed. For example, you might cite a paper like this: \cite{example2024}. All citations will be automatically added to your \texttt{cite.bib} file.

\section{Conclusion}
This template demonstrates basic \LaTeX{} features including mathematical formulas, tables, and lists. The bibliography system is now set up to work seamlessly with Bibby AI's citation management. Feel free to modify it for your needs!

\bibliographystyle{plain}
\bibliography{cite}

\end{document}\documentclass{article}
\usepackage[utf8]{inputenc}
\usepackage{booktabs}
\usepackage{amsmath}
\usepackage{natbib}

\title{My Document}
\author{Your Name}
\date{\today}

\begin{document}

\maketitle

\section{Introduction}
This is a sample document demonstrating various \LaTeX{} features including mathematical equations, tables, and formatting. You can cite papers using the citation search feature, and they will be automatically added to your \texttt{cite.bib} file.

\section{Mathematical Expressions}
Here are some example equations:

\subsection{Quadratic Formula}
The quadratic formula is given by:
\[x = \frac{-b \pm \sqrt{b^2 - 4ac}}{2a}\]

\subsection{Matrix Example}
A sample matrix:
\[
\begin{pmatrix}
a_{11} & a_{12} & a_{13} \\
a_{21} & a_{22} & a_{23} \\
a_{31} & a_{32} & a_{33}
\end{pmatrix}
\]

\section{Tables}
Below is a sample table showing data:

\begin{table}[h]
\centering
\begin{tabular}{@{}lrrr@{}}
\toprule
\textbf{Category} & \textbf{Value 1} & \textbf{Value 2} & \textbf{Total} \\
\midrule
Item A & 10.5 & 20.3 & 30.8 \\
Item B & 15.2 & 25.1 & 40.3 \\
Item C & 22.1 & 12.7 & 34.8 \\
\bottomrule
\end{tabular}
\caption{Sample Data Table}
\label{tab:sample}
\end{table}

\section{Lists}
\subsection{Bullet Points}
\begin{itemize}
  \item First item
  \item Second item
  \item Third item with sub-items:
    \begin{itemize}
      \item Sub-item 1
      \item Sub-item 2
    \end{itemize}
\end{itemize}

\subsection{Numbered List}
\begin{enumerate}
  \item First step
  \item Second step
  \item Third step
\end{enumerate}

\section{Citations and References}
When you use the citation search feature, references will be automatically managed. For example, you might cite a paper like this: \cite{example2024}. All citations will be automatically added to your \texttt{cite.bib} file.

\section{Conclusion}
This template demonstrates basic \LaTeX{} features including mathematical formulas, tables, and lists. The bibliography system is now set up to work seamlessly with Bibby AI's citation management. Feel free to modify it for your needs!

\bibliographystyle{plain}
\bibliography{cite}

\end{document}\documentclass{article}
\usepackage[utf8]{inputenc}
\usepackage{booktabs}
\usepackage{amsmath}
\usepackage{natbib}

\title{My Document}
\author{Your Name}
\date{\today}

\begin{document}

\maketitle

\section{Introduction}
This is a sample document demonstrating various \LaTeX{} features including mathematical equations, tables, and formatting. You can cite papers using the citation search feature, and they will be automatically added to your \texttt{cite.bib} file.

\section{Mathematical Expressions}
Here are some example equations:

\subsection{Quadratic Formula}
The quadratic formula is given by:
\[x = \frac{-b \pm \sqrt{b^2 - 4ac}}{2a}\]

\subsection{Matrix Example}
A sample matrix:
\[
\begin{pmatrix}
a_{11} & a_{12} & a_{13} \\
a_{21} & a_{22} & a_{23} \\
a_{31} & a_{32} & a_{33}
\end{pmatrix}
\]

\section{Tables}
Below is a sample table showing data:

\begin{table}[h]
\centering
\begin{tabular}{@{}lrrr@{}}
\toprule
\textbf{Category} & \textbf{Value 1} & \textbf{Value 2} & \textbf{Total} \\
\midrule
Item A & 10.5 & 20.3 & 30.8 \\
Item B & 15.2 & 25.1 & 40.3 \\
Item C & 22.1 & 12.7 & 34.8 \\
\bottomrule
\end{tabular}
\caption{Sample Data Table}
\label{tab:sample}
\end{table}

\section{Lists}
\subsection{Bullet Points}
\begin{itemize}
  \item First item
  \item Second item
  \item Third item with sub-items:
    \begin{itemize}
      \item Sub-item 1
      \item Sub-item 2
    \end{itemize}
\end{itemize}

\subsection{Numbered List}
\begin{enumerate}
  \item First step
  \item Second step
  \item Third step
\end{enumerate}

\section{Citations and References}
When you use the citation search feature, references will be automatically managed. For example, you might cite a paper like this: \cite{example2024}. All citations will be automatically added to your \texttt{cite.bib} file.

\section{Conclusion}
This template demonstrates basic \LaTeX{} features including mathematical formulas, tables, and lists. The bibliography system is now set up to work seamlessly with Bibby AI's citation management. Feel free to modify it for your needs!

\bibliographystyle{plain}
\bibliography{cite}

\end{document}\documentclass{article}
\usepackage[utf8]{inputenc}
\usepackage{booktabs}
\usepackage{amsmath}
\usepackage{natbib}

\title{My Document}
\author{Your Name}
\date{\today}

\begin{document}

\maketitle

\section{Introduction}
This is a sample document demonstrating various \LaTeX{} features including mathematical equations, tables, and formatting. You can cite papers using the citation search feature, and they will be automatically added to your \texttt{cite.bib} file.

\section{Mathematical Expressions}
Here are some example equations:

\subsection{Quadratic Formula}
The quadratic formula is given by:
\[x = \frac{-b \pm \sqrt{b^2 - 4ac}}{2a}\]

\subsection{Matrix Example}
A sample matrix:
\[
\begin{pmatrix}
a_{11} & a_{12} & a_{13} \\
a_{21} & a_{22} & a_{23} \\
a_{31} & a_{32} & a_{33}
\end{pmatrix}
\]

\section{Tables}
Below is a sample table showing data:

\begin{table}[h]
\centering
\begin{tabular}{@{}lrrr@{}}
\toprule
\textbf{Category} & \textbf{Value 1} & \textbf{Value 2} & \textbf{Total} \\
\midrule
Item A & 10.5 & 20.3 & 30.8 \\
Item B & 15.2 & 25.1 & 40.3 \\
Item C & 22.1 & 12.7 & 34.8 \\
\bottomrule
\end{tabular}
\caption{Sample Data Table}
\label{tab:sample}
\end{table}

\section{Lists}
\subsection{Bullet Points}
\begin{itemize}
  \item First item
  \item Second item
  \item Third item with sub-items:
    \begin{itemize}
      \item Sub-item 1
      \item Sub-item 2
    \end{itemize}
\end{itemize}

\subsection{Numbered List}
\begin{enumerate}
  \item First step
  \item Second step
  \item Third step
\end{enumerate}

\section{Citations and References}
When you use the citation search feature, references will be automatically managed. For example, you might cite a paper like this: \cite{example2024}. All citations will be automatically added to your \texttt{cite.bib} file.

\section{Conclusion}
This template demonstrates basic \LaTeX{} features including mathematical formulas, tables, and lists. The bibliography system is now set up to work seamlessly with Bibby AI's citation management. Feel free to modify it for your needs!

\bibliographystyle{plain}
\bibliography{cite}

\end{document}\documentclass{article}
\usepackage[utf8]{inputenc}
\usepackage{booktabs}
\usepackage{amsmath}
\usepackage{natbib}

\title{My Document}
\author{Your Name}
\date{\today}

\begin{document}

\maketitle

\section{Introduction}
This is a sample document demonstrating various \LaTeX{} features including mathematical equations, tables, and formatting. You can cite papers using the citation search feature, and they will be automatically added to your \texttt{cite.bib} file.

\section{Mathematical Expressions}
Here are some example equations:

\subsection{Quadratic Formula}
The quadratic formula is given by:
\[x = \frac{-b \pm \sqrt{b^2 - 4ac}}{2a}\]

\subsection{Matrix Example}
A sample matrix:
\[
\begin{pmatrix}
a_{11} & a_{12} & a_{13} \\
a_{21} & a_{22} & a_{23} \\
a_{31} & a_{32} & a_{33}
\end{pmatrix}
\]

\section{Tables}
Below is a sample table showing data:

\begin{table}[h]
\centering
\begin{tabular}{@{}lrrr@{}}
\toprule
\textbf{Category} & \textbf{Value 1} & \textbf{Value 2} & \textbf{Total} \\
\midrule
Item A & 10.5 & 20.3 & 30.8 \\
Item B & 15.2 & 25.1 & 40.3 \\
Item C & 22.1 & 12.7 & 34.8 \\
\bottomrule
\end{tabular}
\caption{Sample Data Table}
\label{tab:sample}
\end{table}

\section{Lists}
\subsection{Bullet Points}
\begin{itemize}
  \item First item
  \item Second item
  \item Third item with sub-items:
    \begin{itemize}
      \item Sub-item 1
      \item Sub-item 2
    \end{itemize}
\end{itemize}

\subsection{Numbered List}
\begin{enumerate}
  \item First step
  \item Second step
  \item Third step
\end{enumerate}

\section{Citations and References}
When you use the citation search feature, references will be automatically managed. For example, you might cite a paper like this: \cite{example2024}. All citations will be automatically added to your \texttt{cite.bib} file.

\section{Conclusion}
This template demonstrates basic \LaTeX{} features including mathematical formulas, tables, and lists. The bibliography system is now set up to work seamlessly with Bibby AI's citation management. Feel free to modify it for your needs!

\bibliographystyle{plain}
\bibliography{cite}

\end{document}\documentclass{article}
\usepackage[utf8]{inputenc}
\usepackage{booktabs}
\usepackage{amsmath}
\usepackage{natbib}

\title{My Document}
\author{Your Name}
\date{\today}

\begin{document}

\maketitle

\section{Introduction}
This is a sample document demonstrating various \LaTeX{} features including mathematical equations, tables, and formatting. You can cite papers using the citation search feature, and they will be automatically added to your \texttt{cite.bib} file.

\section{Mathematical Expressions}
Here are some example equations:

\subsection{Quadratic Formula}
The quadratic formula is given by:
\[x = \frac{-b \pm \sqrt{b^2 - 4ac}}{2a}\]

\subsection{Matrix Example}
A sample matrix:
\[
\begin{pmatrix}
a_{11} & a_{12} & a_{13} \\
a_{21} & a_{22} & a_{23} \\
a_{31} & a_{32} & a_{33}
\end{pmatrix}
\]

\section{Tables}
Below is a sample table showing data:

\begin{table}[h]
\centering
\begin{tabular}{@{}lrrr@{}}
\toprule
\textbf{Category} & \textbf{Value 1} & \textbf{Value 2} & \textbf{Total} \\
\midrule
Item A & 10.5 & 20.3 & 30.8 \\
Item B & 15.2 & 25.1 & 40.3 \\
Item C & 22.1 & 12.7 & 34.8 \\
\bottomrule
\end{tabular}
\caption{Sample Data Table}
\label{tab:sample}
\end{table}

\section{Lists}
\subsection{Bullet Points}
\begin{itemize}
  \item First item
  \item Second item
  \item Third item with sub-items:
    \begin{itemize}
      \item Sub-item 1
      \item Sub-item 2
    \end{itemize}
\end{itemize}

\subsection{Numbered List}
\begin{enumerate}
  \item First step
  \item Second step
  \item Third step
\end{enumerate}

\section{Citations and References}
When you use the citation search feature, references will be automatically managed. For example, you might cite a paper like this: \cite{example2024}. All citations will be automatically added to your \texttt{cite.bib} file.

\section{Conclusion}
This template demonstrates basic \LaTeX{} features including mathematical formulas, tables, and lists. The bibliography system is now set up to work seamlessly with Bibby AI's citation management. Feel free to modify it for your needs!

\bibliographystyle{plain}
\bibliography{cite}

\end{document}\documentclass{article}
\usepackage[utf8]{inputenc}
\usepackage{booktabs}
\usepackage{amsmath}
\usepackage{natbib}

\title{My Document}
\author{Your Name}
\date{\today}

\begin{document}

\maketitle

\section{Introduction}
This is a sample document demonstrating various \LaTeX{} features including mathematical equations, tables, and formatting. You can cite papers using the citation search feature, and they will be automatically added to your \texttt{cite.bib} file.

\section{Mathematical Expressions}
Here are some example equations:

\subsection{Quadratic Formula}
The quadratic formula is given by:
\[x = \frac{-b \pm \sqrt{b^2 - 4ac}}{2a}\]

\subsection{Matrix Example}
A sample matrix:
\[
\begin{pmatrix}
a_{11} & a_{12} & a_{13} \\
a_{21} & a_{22} & a_{23} \\
a_{31} & a_{32} & a_{33}
\end{pmatrix}
\]

\section{Tables}
Below is a sample table showing data:

\begin{table}[h]
\centering
\begin{tabular}{@{}lrrr@{}}
\toprule
\textbf{Category} & \textbf{Value 1} & \textbf{Value 2} & \textbf{Total} \\
\midrule
Item A & 10.5 & 20.3 & 30.8 \\
Item B & 15.2 & 25.1 & 40.3 \\
Item C & 22.1 & 12.7 & 34.8 \\
\bottomrule
\end{tabular}
\caption{Sample Data Table}
\label{tab:sample}
\end{table}

\section{Lists}
\subsection{Bullet Points}
\begin{itemize}
  \item First item
  \item Second item
  \item Third item with sub-items:
    \begin{itemize}
      \item Sub-item 1
      \item Sub-item 2
    \end{itemize}
\end{itemize}

\subsection{Numbered List}
\begin{enumerate}
  \item First step
  \item Second step
  \item Third step
\end{enumerate}

\section{Citations and References}
When you use the citation search feature, references will be automatically managed. For example, you might cite a paper like this: \cite{example2024}. All citations will be automatically added to your \texttt{cite.bib} file.

\section{Conclusion}
This template demonstrates basic \LaTeX{} features including mathematical formulas, tables, and lists. The bibliography system is now set up to work seamlessly with Bibby AI's citation management. Feel free to modify it for your needs!

\bibliographystyle{plain}
\bibliography{cite}

\end{document}\documentclass{article}
\usepackage[utf8]{inputenc}
\usepackage{booktabs}
\usepackage{amsmath}
\usepackage{natbib}

\title{My Document}
\author{Your Name}
\date{\today}

\begin{document}

\maketitle

\section{Introduction}
This is a sample document demonstrating various \LaTeX{} features including mathematical equations, tables, and formatting. You can cite papers using the citation search feature, and they will be automatically added to your \texttt{cite.bib} file.

\section{Mathematical Expressions}
Here are some example equations:

\subsection{Quadratic Formula}
The quadratic formula is given by:
\[x = \frac{-b \pm \sqrt{b^2 - 4ac}}{2a}\]

\subsection{Matrix Example}
A sample matrix:
\[
\begin{pmatrix}
a_{11} & a_{12} & a_{13} \\
a_{21} & a_{22} & a_{23} \\
a_{31} & a_{32} & a_{33}
\end{pmatrix}
\]

\section{Tables}
Below is a sample table showing data:

\begin{table}[h]
\centering
\begin{tabular}{@{}lrrr@{}}
\toprule
\textbf{Category} & \textbf{Value 1} & \textbf{Value 2} & \textbf{Total} \\
\midrule
Item A & 10.5 & 20.3 & 30.8 \\
Item B & 15.2 & 25.1 & 40.3 \\
Item C & 22.1 & 12.7 & 34.8 \\
\bottomrule
\end{tabular}
\caption{Sample Data Table}
\label{tab:sample}
\end{table}

\section{Lists}
\subsection{Bullet Points}
\begin{itemize}
  \item First item
  \item Second item
  \item Third item with sub-items:
    \begin{itemize}
      \item Sub-item 1
      \item Sub-item 2
    \end{itemize}
\end{itemize}

\subsection{Numbered List}
\begin{enumerate}
  \item First step
  \item Second step
  \item Third step
\end{enumerate}

\section{Citations and References}
When you use the citation search feature, references will be automatically managed. For example, you might cite a paper like this: \cite{example2024}. All citations will be automatically added to your \texttt{cite.bib} file.

\section{Conclusion}
This template demonstrates basic \LaTeX{} features including mathematical formulas, tables, and lists. The bibliography system is now set up to work seamlessly with Bibby AI's citation management. Feel free to modify it for your needs!

\bibliographystyle{plain}
\bibliography{cite}

\end{document}\documentclass{article}
\usepackage[utf8]{inputenc}
\usepackage{booktabs}
\usepackage{amsmath}
\usepackage{natbib}

\title{My Document}
\author{Your Name}
\date{\today}

\begin{document}

\maketitle

\section{Introduction}
This is a sample document demonstrating various \LaTeX{} features including mathematical equations, tables, and formatting. You can cite papers using the citation search feature, and they will be automatically added to your \texttt{cite.bib} file.

\section{Mathematical Expressions}
Here are some example equations:

\subsection{Quadratic Formula}
The quadratic formula is given by:
\[x = \frac{-b \pm \sqrt{b^2 - 4ac}}{2a}\]

\subsection{Matrix Example}
A sample matrix:
\[
\begin{pmatrix}
a_{11} & a_{12} & a_{13} \\
a_{21} & a_{22} & a_{23} \\
a_{31} & a_{32} & a_{33}
\end{pmatrix}
\]

\section{Tables}
Below is a sample table showing data:

\begin{table}[h]
\centering
\begin{tabular}{@{}lrrr@{}}
\toprule
\textbf{Category} & \textbf{Value 1} & \textbf{Value 2} & \textbf{Total} \\
\midrule
Item A & 10.5 & 20.3 & 30.8 \\
Item B & 15.2 & 25.1 & 40.3 \\
Item C & 22.1 & 12.7 & 34.8 \\
\bottomrule
\end{tabular}
\caption{Sample Data Table}
\label{tab:sample}
\end{table}

\section{Lists}
\subsection{Bullet Points}
\begin{itemize}
  \item First item
  \item Second item
  \item Third item with sub-items:
    \begin{itemize}
      \item Sub-item 1
      \item Sub-item 2
    \end{itemize}
\end{itemize}

\subsection{Numbered List}
\begin{enumerate}
  \item First step
  \item Second step
  \item Third step
\end{enumerate}

\section{Citations and References}
When you use the citation search feature, references will be automatically managed. For example, you might cite a paper like this: \cite{example2024}. All citations will be automatically added to your \texttt{cite.bib} file.

\section{Conclusion}
This template demonstrates basic \LaTeX{} features including mathematical formulas, tables, and lists. The bibliography system is now set up to work seamlessly with Bibby AI's citation management. Feel free to modify it for your needs!

\bibliographystyle{plain}
\bibliography{cite}

\end{document}\documentclass{article}
\usepackage[utf8]{inputenc}
\usepackage{booktabs}
\usepackage{amsmath}
\usepackage{natbib}

\title{My Document}
\author{Your Name}
\date{\today}

\begin{document}

\maketitle

\section{Introduction}
This is a sample document demonstrating various \LaTeX{} features including mathematical equations, tables, and formatting. You can cite papers using the citation search feature, and they will be automatically added to your \texttt{cite.bib} file.

\section{Mathematical Expressions}
Here are some example equations:

\subsection{Quadratic Formula}
The quadratic formula is given by:
\[x = \frac{-b \pm \sqrt{b^2 - 4ac}}{2a}\]

\subsection{Matrix Example}
A sample matrix:
\[
\begin{pmatrix}
a_{11} & a_{12} & a_{13} \\
a_{21} & a_{22} & a_{23} \\
a_{31} & a_{32} & a_{33}
\end{pmatrix}
\]

\section{Tables}
Below is a sample table showing data:

\begin{table}[h]
\centering
\begin{tabular}{@{}lrrr@{}}
\toprule
\textbf{Category} & \textbf{Value 1} & \textbf{Value 2} & \textbf{Total} \\
\midrule
Item A & 10.5 & 20.3 & 30.8 \\
Item B & 15.2 & 25.1 & 40.3 \\
Item C & 22.1 & 12.7 & 34.8 \\
\bottomrule
\end{tabular}
\caption{Sample Data Table}
\label{tab:sample}
\end{table}

\section{Lists}
\subsection{Bullet Points}
\begin{itemize}
  \item First item
  \item Second item
  \item Third item with sub-items:
    \begin{itemize}
      \item Sub-item 1
      \item Sub-item 2
    \end{itemize}
\end{itemize}

\subsection{Numbered List}
\begin{enumerate}
  \item First step
  \item Second step
  \item Third step
\end{enumerate}

\section{Citations and References}
When you use the citation search feature, references will be automatically managed. For example, you might cite a paper like this: \cite{example2024}. All citations will be automatically added to your \texttt{cite.bib} file.

\section{Conclusion}
This template demonstrates basic \LaTeX{} features including mathematical formulas, tables, and lists. The bibliography system is now set up to work seamlessly with Bibby AI's citation management. Feel free to modify it for your needs!

\bibliographystyle{plain}
\bibliography{cite}

\end{document}\documentclass{article}
\usepackage[utf8]{inputenc}
\usepackage{booktabs}
\usepackage{amsmath}
\usepackage{natbib}

\title{My Document}
\author{Your Name}
\date{\today}

\begin{document}

\maketitle

\section{Introduction}
This is a sample document demonstrating various \LaTeX{} features including mathematical equations, tables, and formatting. You can cite papers using the citation search feature, and they will be automatically added to your \texttt{cite.bib} file.

\section{Mathematical Expressions}
Here are some example equations:

\subsection{Quadratic Formula}
The quadratic formula is given by:
\[x = \frac{-b \pm \sqrt{b^2 - 4ac}}{2a}\]

\subsection{Matrix Example}
A sample matrix:
\[
\begin{pmatrix}
a_{11} & a_{12} & a_{13} \\
a_{21} & a_{22} & a_{23} \\
a_{31} & a_{32} & a_{33}
\end{pmatrix}
\]

\section{Tables}
Below is a sample table showing data:

\begin{table}[h]
\centering
\begin{tabular}{@{}lrrr@{}}
\toprule
\textbf{Category} & \textbf{Value 1} & \textbf{Value 2} & \textbf{Total} \\
\midrule
Item A & 10.5 & 20.3 & 30.8 \\
Item B & 15.2 & 25.1 & 40.3 \\
Item C & 22.1 & 12.7 & 34.8 \\
\bottomrule
\end{tabular}
\caption{Sample Data Table}
\label{tab:sample}
\end{table}

\section{Lists}
\subsection{Bullet Points}
\begin{itemize}
  \item First item
  \item Second item
  \item Third item with sub-items:
    \begin{itemize}
      \item Sub-item 1
      \item Sub-item 2
    \end{itemize}
\end{itemize}

\subsection{Numbered List}
\begin{enumerate}
  \item First step
  \item Second step
  \item Third step
\end{enumerate}

\section{Citations and References}
When you use the citation search feature, references will be automatically managed. For example, you might cite a paper like this: \cite{example2024}. All citations will be automatically added to your \texttt{cite.bib} file.

\section{Conclusion}
This template demonstrates basic \LaTeX{} features including mathematical formulas, tables, and lists. The bibliography system is now set up to work seamlessly with Bibby AI's citation management. Feel free to modify it for your needs!

\bibliographystyle{plain}
\bibliography{cite}

\end{document}\documentclass{article}
\usepackage[utf8]{inputenc}
\usepackage{booktabs}
\usepackage{amsmath}
\usepackage{natbib}

\title{My Document}
\author{Your Name}
\date{\today}

\begin{document}

\maketitle

\section{Introduction}
This is a sample document demonstrating various \LaTeX{} features including mathematical equations, tables, and formatting. You can cite papers using the citation search feature, and they will be automatically added to your \texttt{cite.bib} file.

\section{Mathematical Expressions}
Here are some example equations:

\subsection{Quadratic Formula}
The quadratic formula is given by:
\[x = \frac{-b \pm \sqrt{b^2 - 4ac}}{2a}\]

\subsection{Matrix Example}
A sample matrix:
\[
\begin{pmatrix}
a_{11} & a_{12} & a_{13} \\
a_{21} & a_{22} & a_{23} \\
a_{31} & a_{32} & a_{33}
\end{pmatrix}
\]

\section{Tables}
Below is a sample table showing data:

\begin{table}[h]
\centering
\begin{tabular}{@{}lrrr@{}}
\toprule
\textbf{Category} & \textbf{Value 1} & \textbf{Value 2} & \textbf{Total} \\
\midrule
Item A & 10.5 & 20.3 & 30.8 \\
Item B & 15.2 & 25.1 & 40.3 \\
Item C & 22.1 & 12.7 & 34.8 \\
\bottomrule
\end{tabular}
\caption{Sample Data Table}
\label{tab:sample}
\end{table}

\section{Lists}
\subsection{Bullet Points}
\begin{itemize}
  \item First item
  \item Second item
  \item Third item with sub-items:
    \begin{itemize}
      \item Sub-item 1
      \item Sub-item 2
    \end{itemize}
\end{itemize}

\subsection{Numbered List}
\begin{enumerate}
  \item First step
  \item Second step
  \item Third step
\end{enumerate}

\section{Citations and References}
When you use the citation search feature, references will be automatically managed. For example, you might cite a paper like this: \cite{example2024}. All citations will be automatically added to your \texttt{cite.bib} file.

\section{Conclusion}
This template demonstrates basic \LaTeX{} features including mathematical formulas, tables, and lists. The bibliography system is now set up to work seamlessly with Bibby AI's citation management. Feel free to modify it for your needs!

\bibliographystyle{plain}
\bibliography{cite}

\end{document}\documentclass{article}
\usepackage[utf8]{inputenc}
\usepackage{booktabs}
\usepackage{amsmath}
\usepackage{natbib}

\title{My Document}
\author{Your Name}
\date{\today}

\begin{document}

\maketitle

\section{Introduction}
This is a sample document demonstrating various \LaTeX{} features including mathematical equations, tables, and formatting. You can cite papers using the citation search feature, and they will be automatically added to your \texttt{cite.bib} file.

\section{Mathematical Expressions}
Here are some example equations:

\subsection{Quadratic Formula}
The quadratic formula is given by:
\[x = \frac{-b \pm \sqrt{b^2 - 4ac}}{2a}\]

\subsection{Matrix Example}
A sample matrix:
\[
\begin{pmatrix}
a_{11} & a_{12} & a_{13} \\
a_{21} & a_{22} & a_{23} \\
a_{31} & a_{32} & a_{33}
\end{pmatrix}
\]

\section{Tables}
Below is a sample table showing data:

\begin{table}[h]
\centering
\begin{tabular}{@{}lrrr@{}}
\toprule
\textbf{Category} & \textbf{Value 1} & \textbf{Value 2} & \textbf{Total} \\
\midrule
Item A & 10.5 & 20.3 & 30.8 \\
Item B & 15.2 & 25.1 & 40.3 \\
Item C & 22.1 & 12.7 & 34.8 \\
\bottomrule
\end{tabular}
\caption{Sample Data Table}
\label{tab:sample}
\end{table}

\section{Lists}
\subsection{Bullet Points}
\begin{itemize}
  \item First item
  \item Second item
  \item Third item with sub-items:
    \begin{itemize}
      \item Sub-item 1
      \item Sub-item 2
    \end{itemize}
\end{itemize}

\subsection{Numbered List}
\begin{enumerate}
  \item First step
  \item Second step
  \item Third step
\end{enumerate}

\section{Citations and References}
When you use the citation search feature, references will be automatically managed. For example, you might cite a paper like this: \cite{example2024}. All citations will be automatically added to your \texttt{cite.bib} file.

\section{Conclusion}
This template demonstrates basic \LaTeX{} features including mathematical formulas, tables, and lists. The bibliography system is now set up to work seamlessly with Bibby AI's citation management. Feel free to modify it for your needs!

\bibliographystyle{plain}
\bibliography{cite}

\end{document}\documentclass{article}
\usepackage[utf8]{inputenc}
\usepackage{booktabs}
\usepackage{amsmath}
\usepackage{natbib}

\title{My Document}
\author{Your Name}
\date{\today}

\begin{document}

\maketitle

\section{Introduction}
This is a sample document demonstrating various \LaTeX{} features including mathematical equations, tables, and formatting. You can cite papers using the citation search feature, and they will be automatically added to your \texttt{cite.bib} file.

\section{Mathematical Expressions}
Here are some example equations:

\subsection{Quadratic Formula}
The quadratic formula is given by:
\[x = \frac{-b \pm \sqrt{b^2 - 4ac}}{2a}\]

\subsection{Matrix Example}
A sample matrix:
\[
\begin{pmatrix}
a_{11} & a_{12} & a_{13} \\
a_{21} & a_{22} & a_{23} \\
a_{31} & a_{32} & a_{33}
\end{pmatrix}
\]

\section{Tables}
Below is a sample table showing data:

\begin{table}[h]
\centering
\begin{tabular}{@{}lrrr@{}}
\toprule
\textbf{Category} & \textbf{Value 1} & \textbf{Value 2} & \textbf{Total} \\
\midrule
Item A & 10.5 & 20.3 & 30.8 \\
Item B & 15.2 & 25.1 & 40.3 \\
Item C & 22.1 & 12.7 & 34.8 \\
\bottomrule
\end{tabular}
\caption{Sample Data Table}
\label{tab:sample}
\end{table}

\section{Lists}
\subsection{Bullet Points}
\begin{itemize}
  \item First item
  \item Second item
  \item Third item with sub-items:
    \begin{itemize}
      \item Sub-item 1
      \item Sub-item 2
    \end{itemize}
\end{itemize}

\subsection{Numbered List}
\begin{enumerate}
  \item First step
  \item Second step
  \item Third step
\end{enumerate}

\section{Citations and References}
When you use the citation search feature, references will be automatically managed. For example, you might cite a paper like this: \cite{example2024}. All citations will be automatically added to your \texttt{cite.bib} file.

\section{Conclusion}
This template demonstrates basic \LaTeX{} features including mathematical formulas, tables, and lists. The bibliography system is now set up to work seamlessly with Bibby AI's citation management. Feel free to modify it for your needs!

\bibliographystyle{plain}
\bibliography{cite}

\end{document}\documentclass{article}
\usepackage[utf8]{inputenc}
\usepackage{booktabs}
\usepackage{amsmath}
\usepackage{natbib}

\title{My Document}
\author{Your Name}
\date{\today}

\begin{document}

\maketitle

\section{Introduction}
This is a sample document demonstrating various \LaTeX{} features including mathematical equations, tables, and formatting. You can cite papers using the citation search feature, and they will be automatically added to your \texttt{cite.bib} file.

\section{Mathematical Expressions}
Here are some example equations:

\subsection{Quadratic Formula}
The quadratic formula is given by:
\[x = \frac{-b \pm \sqrt{b^2 - 4ac}}{2a}\]

\subsection{Matrix Example}
A sample matrix:
\[
\begin{pmatrix}
a_{11} & a_{12} & a_{13} \\
a_{21} & a_{22} & a_{23} \\
a_{31} & a_{32} & a_{33}
\end{pmatrix}
\]

\section{Tables}
Below is a sample table showing data:

\begin{table}[h]
\centering
\begin{tabular}{@{}lrrr@{}}
\toprule
\textbf{Category} & \textbf{Value 1} & \textbf{Value 2} & \textbf{Total} \\
\midrule
Item A & 10.5 & 20.3 & 30.8 \\
Item B & 15.2 & 25.1 & 40.3 \\
Item C & 22.1 & 12.7 & 34.8 \\
\bottomrule
\end{tabular}
\caption{Sample Data Table}
\label{tab:sample}
\end{table}

\section{Lists}
\subsection{Bullet Points}
\begin{itemize}
  \item First item
  \item Second item
  \item Third item with sub-items:
    \begin{itemize}
      \item Sub-item 1
      \item Sub-item 2
    \end{itemize}
\end{itemize}

\subsection{Numbered List}
\begin{enumerate}
  \item First step
  \item Second step
  \item Third step
\end{enumerate}

\section{Citations and References}
When you use the citation search feature, references will be automatically managed. For example, you might cite a paper like this: \cite{example2024}. All citations will be automatically added to your \texttt{cite.bib} file.

\section{Conclusion}
This template demonstrates basic \LaTeX{} features including mathematical formulas, tables, and lists. The bibliography system is now set up to work seamlessly with Bibby AI's citation management. Feel free to modify it for your needs!

\bibliographystyle{plain}
\bibliography{cite}

\end{document}\documentclass{article}
\usepackage[utf8]{inputenc}
\usepackage{booktabs}
\usepackage{amsmath}
\usepackage{natbib}

\title{My Document}
\author{Your Name}
\date{\today}

\begin{document}

\maketitle

\section{Introduction}
This is a sample document demonstrating various \LaTeX{} features including mathematical equations, tables, and formatting. You can cite papers using the citation search feature, and they will be automatically added to your \texttt{cite.bib} file.

\section{Mathematical Expressions}
Here are some example equations:

\subsection{Quadratic Formula}
The quadratic formula is given by:
\[x = \frac{-b \pm \sqrt{b^2 - 4ac}}{2a}\]

\subsection{Matrix Example}
A sample matrix:
\[
\begin{pmatrix}
a_{11} & a_{12} & a_{13} \\
a_{21} & a_{22} & a_{23} \\
a_{31} & a_{32} & a_{33}
\end{pmatrix}
\]

\section{Tables}
Below is a sample table showing data:

\begin{table}[h]
\centering
\begin{tabular}{@{}lrrr@{}}
\toprule
\textbf{Category} & \textbf{Value 1} & \textbf{Value 2} & \textbf{Total} \\
\midrule
Item A & 10.5 & 20.3 & 30.8 \\
Item B & 15.2 & 25.1 & 40.3 \\
Item C & 22.1 & 12.7 & 34.8 \\
\bottomrule
\end{tabular}
\caption{Sample Data Table}
\label{tab:sample}
\end{table}

\section{Lists}
\subsection{Bullet Points}
\begin{itemize}
  \item First item
  \item Second item
  \item Third item with sub-items:
    \begin{itemize}
      \item Sub-item 1
      \item Sub-item 2
    \end{itemize}
\end{itemize}

\subsection{Numbered List}
\begin{enumerate}
  \item First step
  \item Second step
  \item Third step
\end{enumerate}

\section{Citations and References}
When you use the citation search feature, references will be automatically managed. For example, you might cite a paper like this: \cite{example2024}. All citations will be automatically added to your \texttt{cite.bib} file.

\section{Conclusion}
This template demonstrates basic \LaTeX{} features including mathematical formulas, tables, and lists. The bibliography system is now set up to work seamlessly with Bibby AI's citation management. Feel free to modify it for your needs!

\bibliographystyle{plain}
\bibliography{cite}

\end{document}\documentclass{article}
\usepackage[utf8]{inputenc}
\usepackage{booktabs}
\usepackage{amsmath}
\usepackage{natbib}

\title{My Document}
\author{Your Name}
\date{\today}

\begin{document}

\maketitle

\section{Introduction}
This is a sample document demonstrating various \LaTeX{} features including mathematical equations, tables, and formatting. You can cite papers using the citation search feature, and they will be automatically added to your \texttt{cite.bib} file.

\section{Mathematical Expressions}
Here are some example equations:

\subsection{Quadratic Formula}
The quadratic formula is given by:
\[x = \frac{-b \pm \sqrt{b^2 - 4ac}}{2a}\]

\subsection{Matrix Example}
A sample matrix:
\[
\begin{pmatrix}
a_{11} & a_{12} & a_{13} \\
a_{21} & a_{22} & a_{23} \\
a_{31} & a_{32} & a_{33}
\end{pmatrix}
\]

\section{Tables}
Below is a sample table showing data:

\begin{table}[h]
\centering
\begin{tabular}{@{}lrrr@{}}
\toprule
\textbf{Category} & \textbf{Value 1} & \textbf{Value 2} & \textbf{Total} \\
\midrule
Item A & 10.5 & 20.3 & 30.8 \\
Item B & 15.2 & 25.1 & 40.3 \\
Item C & 22.1 & 12.7 & 34.8 \\
\bottomrule
\end{tabular}
\caption{Sample Data Table}
\label{tab:sample}
\end{table}

\section{Lists}
\subsection{Bullet Points}
\begin{itemize}
  \item First item
  \item Second item
  \item Third item with sub-items:
    \begin{itemize}
      \item Sub-item 1
      \item Sub-item 2
    \end{itemize}
\end{itemize}

\subsection{Numbered List}
\begin{enumerate}
  \item First step
  \item Second step
  \item Third step
\end{enumerate}

\section{Citations and References}
When you use the citation search feature, references will be automatically managed. For example, you might cite a paper like this: \cite{example2024}. All citations will be automatically added to your \texttt{cite.bib} file.

\section{Conclusion}
This template demonstrates basic \LaTeX{} features including mathematical formulas, tables, and lists. The bibliography system is now set up to work seamlessly with Bibby AI's citation management. Feel free to modify it for your needs!

\bibliographystyle{plain}
\bibliography{cite}

\end{document}\documentclass{article}
\usepackage[utf8]{inputenc}
\usepackage{booktabs}
\usepackage{amsmath}
\usepackage{natbib}

\title{My Document}
\author{Your Name}
\date{\today}

\begin{document}

\maketitle

\section{Introduction}
This is a sample document demonstrating various \LaTeX{} features including mathematical equations, tables, and formatting. You can cite papers using the citation search feature, and they will be automatically added to your \texttt{cite.bib} file.

\section{Mathematical Expressions}
Here are some example equations:

\subsection{Quadratic Formula}
The quadratic formula is given by:
\[x = \frac{-b \pm \sqrt{b^2 - 4ac}}{2a}\]

\subsection{Matrix Example}
A sample matrix:
\[
\begin{pmatrix}
a_{11} & a_{12} & a_{13} \\
a_{21} & a_{22} & a_{23} \\
a_{31} & a_{32} & a_{33}
\end{pmatrix}
\]

\section{Tables}
Below is a sample table showing data:

\begin{table}[h]
\centering
\begin{tabular}{@{}lrrr@{}}
\toprule
\textbf{Category} & \textbf{Value 1} & \textbf{Value 2} & \textbf{Total} \\
\midrule
Item A & 10.5 & 20.3 & 30.8 \\
Item B & 15.2 & 25.1 & 40.3 \\
Item C & 22.1 & 12.7 & 34.8 \\
\bottomrule
\end{tabular}
\caption{Sample Data Table}
\label{tab:sample}
\end{table}

\section{Lists}
\subsection{Bullet Points}
\begin{itemize}
  \item First item
  \item Second item
  \item Third item with sub-items:
    \begin{itemize}
      \item Sub-item 1
      \item Sub-item 2
    \end{itemize}
\end{itemize}

\subsection{Numbered List}
\begin{enumerate}
  \item First step
  \item Second step
  \item Third step
\end{enumerate}

\section{Citations and References}
When you use the citation search feature, references will be automatically managed. For example, you might cite a paper like this: \cite{example2024}. All citations will be automatically added to your \texttt{cite.bib} file.

\section{Conclusion}
This template demonstrates basic \LaTeX{} features including mathematical formulas, tables, and lists. The bibliography system is now set up to work seamlessly with Bibby AI's citation management. Feel free to modify it for your needs!

\bibliographystyle{plain}
\bibliography{cite}

\end{document}\documentclass{article}
\usepackage[utf8]{inputenc}
\usepackage{booktabs}
\usepackage{amsmath}
\usepackage{natbib}

\title{My Document}
\author{Your Name}
\date{\today}

\begin{document}

\maketitle

\section{Introduction}
This is a sample document demonstrating various \LaTeX{} features including mathematical equations, tables, and formatting. You can cite papers using the citation search feature, and they will be automatically added to your \texttt{cite.bib} file.

\section{Mathematical Expressions}
Here are some example equations:

\subsection{Quadratic Formula}
The quadratic formula is given by:
\[x = \frac{-b \pm \sqrt{b^2 - 4ac}}{2a}\]

\subsection{Matrix Example}
A sample matrix:
\[
\begin{pmatrix}
a_{11} & a_{12} & a_{13} \\
a_{21} & a_{22} & a_{23} \\
a_{31} & a_{32} & a_{33}
\end{pmatrix}
\]

\section{Tables}
Below is a sample table showing data:

\begin{table}[h]
\centering
\begin{tabular}{@{}lrrr@{}}
\toprule
\textbf{Category} & \textbf{Value 1} & \textbf{Value 2} & \textbf{Total} \\
\midrule
Item A & 10.5 & 20.3 & 30.8 \\
Item B & 15.2 & 25.1 & 40.3 \\
Item C & 22.1 & 12.7 & 34.8 \\
\bottomrule
\end{tabular}
\caption{Sample Data Table}
\label{tab:sample}
\end{table}

\section{Lists}
\subsection{Bullet Points}
\begin{itemize}
  \item First item
  \item Second item
  \item Third item with sub-items:
    \begin{itemize}
      \item Sub-item 1
      \item Sub-item 2
    \end{itemize}
\end{itemize}

\subsection{Numbered List}
\begin{enumerate}
  \item First step
  \item Second step
  \item Third step
\end{enumerate}

\section{Citations and References}
When you use the citation search feature, references will be automatically managed. For example, you might cite a paper like this: \cite{example2024}. All citations will be automatically added to your \texttt{cite.bib} file.

\section{Conclusion}
This template demonstrates basic \LaTeX{} features including mathematical formulas, tables, and lists. The bibliography system is now set up to work seamlessly with Bibby AI's citation management. Feel free to modify it for your needs!

\bibliographystyle{plain}
\bibliography{cite}

\end{document}\documentclass{article}
\usepackage[utf8]{inputenc}
\usepackage{booktabs}
\usepackage{amsmath}
\usepackage{natbib}

\title{My Document}
\author{Your Name}
\date{\today}

\begin{document}

\maketitle

\section{Introduction}
This is a sample document demonstrating various \LaTeX{} features including mathematical equations, tables, and formatting. You can cite papers using the citation search feature, and they will be automatically added to your \texttt{cite.bib} file.

\section{Mathematical Expressions}
Here are some example equations:

\subsection{Quadratic Formula}
The quadratic formula is given by:
\[x = \frac{-b \pm \sqrt{b^2 - 4ac}}{2a}\]

\subsection{Matrix Example}
A sample matrix:
\[
\begin{pmatrix}
a_{11} & a_{12} & a_{13} \\
a_{21} & a_{22} & a_{23} \\
a_{31} & a_{32} & a_{33}
\end{pmatrix}
\]

\section{Tables}
Below is a sample table showing data:

\begin{table}[h]
\centering
\begin{tabular}{@{}lrrr@{}}
\toprule
\textbf{Category} & \textbf{Value 1} & \textbf{Value 2} & \textbf{Total} \\
\midrule
Item A & 10.5 & 20.3 & 30.8 \\
Item B & 15.2 & 25.1 & 40.3 \\
Item C & 22.1 & 12.7 & 34.8 \\
\bottomrule
\end{tabular}
\caption{Sample Data Table}
\label{tab:sample}
\end{table}

\section{Lists}
\subsection{Bullet Points}
\begin{itemize}
  \item First item
  \item Second item
  \item Third item with sub-items:
    \begin{itemize}
      \item Sub-item 1
      \item Sub-item 2
    \end{itemize}
\end{itemize}

\subsection{Numbered List}
\begin{enumerate}
  \item First step
  \item Second step
  \item Third step
\end{enumerate}

\section{Citations and References}
When you use the citation search feature, references will be automatically managed. For example, you might cite a paper like this: \cite{example2024}. All citations will be automatically added to your \texttt{cite.bib} file.

\section{Conclusion}
This template demonstrates basic \LaTeX{} features including mathematical formulas, tables, and lists. The bibliography system is now set up to work seamlessly with Bibby AI's citation management. Feel free to modify it for your needs!

\bibliographystyle{plain}
\bibliography{cite}

\end{document}\documentclass{article}
\usepackage[utf8]{inputenc}
\usepackage{booktabs}
\usepackage{amsmath}
\usepackage{natbib}

\title{My Document}
\author{Your Name}
\date{\today}

\begin{document}

\maketitle

\section{Introduction}
This is a sample document demonstrating various \LaTeX{} features including mathematical equations, tables, and formatting. You can cite papers using the citation search feature, and they will be automatically added to your \texttt{cite.bib} file.

\section{Mathematical Expressions}
Here are some example equations:

\subsection{Quadratic Formula}
The quadratic formula is given by:
\[x = \frac{-b \pm \sqrt{b^2 - 4ac}}{2a}\]

\subsection{Matrix Example}
A sample matrix:
\[
\begin{pmatrix}
a_{11} & a_{12} & a_{13} \\
a_{21} & a_{22} & a_{23} \\
a_{31} & a_{32} & a_{33}
\end{pmatrix}
\]

\section{Tables}
Below is a sample table showing data:

\begin{table}[h]
\centering
\begin{tabular}{@{}lrrr@{}}
\toprule
\textbf{Category} & \textbf{Value 1} & \textbf{Value 2} & \textbf{Total} \\
\midrule
Item A & 10.5 & 20.3 & 30.8 \\
Item B & 15.2 & 25.1 & 40.3 \\
Item C & 22.1 & 12.7 & 34.8 \\
\bottomrule
\end{tabular}
\caption{Sample Data Table}
\label{tab:sample}
\end{table}

\section{Lists}
\subsection{Bullet Points}
\begin{itemize}
  \item First item
  \item Second item
  \item Third item with sub-items:
    \begin{itemize}
      \item Sub-item 1
      \item Sub-item 2
    \end{itemize}
\end{itemize}

\subsection{Numbered List}
\begin{enumerate}
  \item First step
  \item Second step
  \item Third step
\end{enumerate}

\section{Citations and References}
When you use the citation search feature, references will be automatically managed. For example, you might cite a paper like this: \cite{example2024}. All citations will be automatically added to your \texttt{cite.bib} file.

\section{Conclusion}
This template demonstrates basic \LaTeX{} features including mathematical formulas, tables, and lists. The bibliography system is now set up to work seamlessly with Bibby AI's citation management. Feel free to modify it for your needs!

\bibliographystyle{plain}
\bibliography{cite}

\end{document}\documentclass{article}
\usepackage[utf8]{inputenc}
\usepackage{booktabs}
\usepackage{amsmath}
\usepackage{natbib}

\title{My Document}
\author{Your Name}
\date{\today}

\begin{document}

\maketitle

\section{Introduction}
This is a sample document demonstrating various \LaTeX{} features including mathematical equations, tables, and formatting. You can cite papers using the citation search feature, and they will be automatically added to your \texttt{cite.bib} file.

\section{Mathematical Expressions}
Here are some example equations:

\subsection{Quadratic Formula}
The quadratic formula is given by:
\[x = \frac{-b \pm \sqrt{b^2 - 4ac}}{2a}\]

\subsection{Matrix Example}
A sample matrix:
\[
\begin{pmatrix}
a_{11} & a_{12} & a_{13} \\
a_{21} & a_{22} & a_{23} \\
a_{31} & a_{32} & a_{33}
\end{pmatrix}
\]

\section{Tables}
Below is a sample table showing data:

\begin{table}[h]
\centering
\begin{tabular}{@{}lrrr@{}}
\toprule
\textbf{Category} & \textbf{Value 1} & \textbf{Value 2} & \textbf{Total} \\
\midrule
Item A & 10.5 & 20.3 & 30.8 \\
Item B & 15.2 & 25.1 & 40.3 \\
Item C & 22.1 & 12.7 & 34.8 \\
\bottomrule
\end{tabular}
\caption{Sample Data Table}
\label{tab:sample}
\end{table}

\section{Lists}
\subsection{Bullet Points}
\begin{itemize}
  \item First item
  \item Second item
  \item Third item with sub-items:
    \begin{itemize}
      \item Sub-item 1
      \item Sub-item 2
    \end{itemize}
\end{itemize}

\subsection{Numbered List}
\begin{enumerate}
  \item First step
  \item Second step
  \item Third step
\end{enumerate}

\section{Citations and References}
When you use the citation search feature, references will be automatically managed. For example, you might cite a paper like this: \cite{example2024}. All citations will be automatically added to your \texttt{cite.bib} file.

\section{Conclusion}
This template demonstrates basic \LaTeX{} features including mathematical formulas, tables, and lists. The bibliography system is now set up to work seamlessly with Bibby AI's citation management. Feel free to modify it for your needs!

\bibliographystyle{plain}
\bibliography{cite}

\end{document}\documentclass{article}
\usepackage[utf8]{inputenc}
\usepackage{booktabs}
\usepackage{amsmath}
\usepackage{natbib}

\title{My Document}
\author{Your Name}
\date{\today}

\begin{document}

\maketitle

\section{Introduction}
This is a sample document demonstrating various \LaTeX{} features including mathematical equations, tables, and formatting. You can cite papers using the citation search feature, and they will be automatically added to your \texttt{cite.bib} file.

\section{Mathematical Expressions}
Here are some example equations:

\subsection{Quadratic Formula}
The quadratic formula is given by:
\[x = \frac{-b \pm \sqrt{b^2 - 4ac}}{2a}\]

\subsection{Matrix Example}
A sample matrix:
\[
\begin{pmatrix}
a_{11} & a_{12} & a_{13} \\
a_{21} & a_{22} & a_{23} \\
a_{31} & a_{32} & a_{33}
\end{pmatrix}
\]

\section{Tables}
Below is a sample table showing data:

\begin{table}[h]
\centering
\begin{tabular}{@{}lrrr@{}}
\toprule
\textbf{Category} & \textbf{Value 1} & \textbf{Value 2} & \textbf{Total} \\
\midrule
Item A & 10.5 & 20.3 & 30.8 \\
Item B & 15.2 & 25.1 & 40.3 \\
Item C & 22.1 & 12.7 & 34.8 \\
\bottomrule
\end{tabular}
\caption{Sample Data Table}
\label{tab:sample}
\end{table}

\section{Lists}
\subsection{Bullet Points}
\begin{itemize}
  \item First item
  \item Second item
  \item Third item with sub-items:
    \begin{itemize}
      \item Sub-item 1
      \item Sub-item 2
    \end{itemize}
\end{itemize}

\subsection{Numbered List}
\begin{enumerate}
  \item First step
  \item Second step
  \item Third step
\end{enumerate}

\section{Citations and References}
When you use the citation search feature, references will be automatically managed. For example, you might cite a paper like this: \cite{example2024}. All citations will be automatically added to your \texttt{cite.bib} file.

\section{Conclusion}
This template demonstrates basic \LaTeX{} features including mathematical formulas, tables, and lists. The bibliography system is now set up to work seamlessly with Bibby AI's citation management. Feel free to modify it for your needs!

\bibliographystyle{plain}
\bibliography{cite}

\end{document}\documentclass{article}
\usepackage[utf8]{inputenc}
\usepackage{booktabs}
\usepackage{amsmath}
\usepackage{natbib}

\title{My Document}
\author{Your Name}
\date{\today}

\begin{document}

\maketitle

\section{Introduction}
This is a sample document demonstrating various \LaTeX{} features including mathematical equations, tables, and formatting. You can cite papers using the citation search feature, and they will be automatically added to your \texttt{cite.bib} file.

\section{Mathematical Expressions}
Here are some example equations:

\subsection{Quadratic Formula}
The quadratic formula is given by:
\[x = \frac{-b \pm \sqrt{b^2 - 4ac}}{2a}\]

\subsection{Matrix Example}
A sample matrix:
\[
\begin{pmatrix}
a_{11} & a_{12} & a_{13} \\
a_{21} & a_{22} & a_{23} \\
a_{31} & a_{32} & a_{33}
\end{pmatrix}
\]

\section{Tables}
Below is a sample table showing data:

\begin{table}[h]
\centering
\begin{tabular}{@{}lrrr@{}}
\toprule
\textbf{Category} & \textbf{Value 1} & \textbf{Value 2} & \textbf{Total} \\
\midrule
Item A & 10.5 & 20.3 & 30.8 \\
Item B & 15.2 & 25.1 & 40.3 \\
Item C & 22.1 & 12.7 & 34.8 \\
\bottomrule
\end{tabular}
\caption{Sample Data Table}
\label{tab:sample}
\end{table}

\section{Lists}
\subsection{Bullet Points}
\begin{itemize}
  \item First item
  \item Second item
  \item Third item with sub-items:
    \begin{itemize}
      \item Sub-item 1
      \item Sub-item 2
    \end{itemize}
\end{itemize}

\subsection{Numbered List}
\begin{enumerate}
  \item First step
  \item Second step
  \item Third step
\end{enumerate}

\section{Citations and References}
When you use the citation search feature, references will be automatically managed. For example, you might cite a paper like this: \cite{example2024}. All citations will be automatically added to your \texttt{cite.bib} file.

\section{Conclusion}
This template demonstrates basic \LaTeX{} features including mathematical formulas, tables, and lists. The bibliography system is now set up to work seamlessly with Bibby AI's citation management. Feel free to modify it for your needs!

\bibliographystyle{plain}
\bibliography{cite}

\end{document}\documentclass{article}
\usepackage[utf8]{inputenc}
\usepackage{booktabs}
\usepackage{amsmath}
\usepackage{natbib}

\title{My Document}
\author{Your Name}
\date{\today}

\begin{document}

\maketitle

\section{Introduction}
This is a sample document demonstrating various \LaTeX{} features including mathematical equations, tables, and formatting. You can cite papers using the citation search feature, and they will be automatically added to your \texttt{cite.bib} file.

\section{Mathematical Expressions}
Here are some example equations:

\subsection{Quadratic Formula}
The quadratic formula is given by:
\[x = \frac{-b \pm \sqrt{b^2 - 4ac}}{2a}\]

\subsection{Matrix Example}
A sample matrix:
\[
\begin{pmatrix}
a_{11} & a_{12} & a_{13} \\
a_{21} & a_{22} & a_{23} \\
a_{31} & a_{32} & a_{33}
\end{pmatrix}
\]

\section{Tables}
Below is a sample table showing data:

\begin{table}[h]
\centering
\begin{tabular}{@{}lrrr@{}}
\toprule
\textbf{Category} & \textbf{Value 1} & \textbf{Value 2} & \textbf{Total} \\
\midrule
Item A & 10.5 & 20.3 & 30.8 \\
Item B & 15.2 & 25.1 & 40.3 \\
Item C & 22.1 & 12.7 & 34.8 \\
\bottomrule
\end{tabular}
\caption{Sample Data Table}
\label{tab:sample}
\end{table}

\section{Lists}
\subsection{Bullet Points}
\begin{itemize}
  \item First item
  \item Second item
  \item Third item with sub-items:
    \begin{itemize}
      \item Sub-item 1
      \item Sub-item 2
    \end{itemize}
\end{itemize}

\subsection{Numbered List}
\begin{enumerate}
  \item First step
  \item Second step
  \item Third step
\end{enumerate}

\section{Citations and References}
When you use the citation search feature, references will be automatically managed. For example, you might cite a paper like this: \cite{example2024}. All citations will be automatically added to your \texttt{cite.bib} file.

\section{Conclusion}
This template demonstrates basic \LaTeX{} features including mathematical formulas, tables, and lists. The bibliography system is now set up to work seamlessly with Bibby AI's citation management. Feel free to modify it for your needs!

\bibliographystyle{plain}
\bibliography{cite}

\end{document}\documentclass{article}
\usepackage[utf8]{inputenc}
\usepackage{booktabs}
\usepackage{amsmath}
\usepackage{natbib}

\title{My Document}
\author{Your Name}
\date{\today}

\begin{document}

\maketitle

\section{Introduction}
This is a sample document demonstrating various \LaTeX{} features including mathematical equations, tables, and formatting. You can cite papers using the citation search feature, and they will be automatically added to your \texttt{cite.bib} file.

\section{Mathematical Expressions}
Here are some example equations:

\subsection{Quadratic Formula}
The quadratic formula is given by:
\[x = \frac{-b \pm \sqrt{b^2 - 4ac}}{2a}\]

\subsection{Matrix Example}
A sample matrix:
\[
\begin{pmatrix}
a_{11} & a_{12} & a_{13} \\
a_{21} & a_{22} & a_{23} \\
a_{31} & a_{32} & a_{33}
\end{pmatrix}
\]

\section{Tables}
Below is a sample table showing data:

\begin{table}[h]
\centering
\begin{tabular}{@{}lrrr@{}}
\toprule
\textbf{Category} & \textbf{Value 1} & \textbf{Value 2} & \textbf{Total} \\
\midrule
Item A & 10.5 & 20.3 & 30.8 \\
Item B & 15.2 & 25.1 & 40.3 \\
Item C & 22.1 & 12.7 & 34.8 \\
\bottomrule
\end{tabular}
\caption{Sample Data Table}
\label{tab:sample}
\end{table}

\section{Lists}
\subsection{Bullet Points}
\begin{itemize}
  \item First item
  \item Second item
  \item Third item with sub-items:
    \begin{itemize}
      \item Sub-item 1
      \item Sub-item 2
    \end{itemize}
\end{itemize}

\subsection{Numbered List}
\begin{enumerate}
  \item First step
  \item Second step
  \item Third step
\end{enumerate}

\section{Citations and References}
When you use the citation search feature, references will be automatically managed. For example, you might cite a paper like this: \cite{example2024}. All citations will be automatically added to your \texttt{cite.bib} file.

\section{Conclusion}
This template demonstrates basic \LaTeX{} features including mathematical formulas, tables, and lists. The bibliography system is now set up to work seamlessly with Bibby AI's citation management. Feel free to modify it for your needs!

\bibliographystyle{plain}
\bibliography{cite}

\end{document}\documentclass{article}
\usepackage[utf8]{inputenc}
\usepackage{booktabs}
\usepackage{amsmath}
\usepackage{natbib}

\title{My Document}
\author{Your Name}
\date{\today}

\begin{document}

\maketitle

\section{Introduction}
This is a sample document demonstrating various \LaTeX{} features including mathematical equations, tables, and formatting. You can cite papers using the citation search feature, and they will be automatically added to your \texttt{cite.bib} file.

\section{Mathematical Expressions}
Here are some example equations:

\subsection{Quadratic Formula}
The quadratic formula is given by:
\[x = \frac{-b \pm \sqrt{b^2 - 4ac}}{2a}\]

\subsection{Matrix Example}
A sample matrix:
\[
\begin{pmatrix}
a_{11} & a_{12} & a_{13} \\
a_{21} & a_{22} & a_{23} \\
a_{31} & a_{32} & a_{33}
\end{pmatrix}
\]

\section{Tables}
Below is a sample table showing data:

\begin{table}[h]
\centering
\begin{tabular}{@{}lrrr@{}}
\toprule
\textbf{Category} & \textbf{Value 1} & \textbf{Value 2} & \textbf{Total} \\
\midrule
Item A & 10.5 & 20.3 & 30.8 \\
Item B & 15.2 & 25.1 & 40.3 \\
Item C & 22.1 & 12.7 & 34.8 \\
\bottomrule
\end{tabular}
\caption{Sample Data Table}
\label{tab:sample}
\end{table}

\section{Lists}
\subsection{Bullet Points}
\begin{itemize}
  \item First item
  \item Second item
  \item Third item with sub-items:
    \begin{itemize}
      \item Sub-item 1
      \item Sub-item 2
    \end{itemize}
\end{itemize}

\subsection{Numbered List}
\begin{enumerate}
  \item First step
  \item Second step
  \item Third step
\end{enumerate}

\section{Citations and References}
When you use the citation search feature, references will be automatically managed. For example, you might cite a paper like this: \cite{example2024}. All citations will be automatically added to your \texttt{cite.bib} file.

\section{Conclusion}
This template demonstrates basic \LaTeX{} features including mathematical formulas, tables, and lists. The bibliography system is now set up to work seamlessly with Bibby AI's citation management. Feel free to modify it for your needs!

\bibliographystyle{plain}
\bibliography{cite}

\end{document}\documentclass{article}
\usepackage[utf8]{inputenc}
\usepackage{booktabs}
\usepackage{amsmath}
\usepackage{natbib}

\title{My Document}
\author{Your Name}
\date{\today}

\begin{document}

\maketitle

\section{Introduction}
This is a sample document demonstrating various \LaTeX{} features including mathematical equations, tables, and formatting. You can cite papers using the citation search feature, and they will be automatically added to your \texttt{cite.bib} file.

\section{Mathematical Expressions}
Here are some example equations:

\subsection{Quadratic Formula}
The quadratic formula is given by:
\[x = \frac{-b \pm \sqrt{b^2 - 4ac}}{2a}\]

\subsection{Matrix Example}
A sample matrix:
\[
\begin{pmatrix}
a_{11} & a_{12} & a_{13} \\
a_{21} & a_{22} & a_{23} \\
a_{31} & a_{32} & a_{33}
\end{pmatrix}
\]

\section{Tables}
Below is a sample table showing data:

\begin{table}[h]
\centering
\begin{tabular}{@{}lrrr@{}}
\toprule
\textbf{Category} & \textbf{Value 1} & \textbf{Value 2} & \textbf{Total} \\
\midrule
Item A & 10.5 & 20.3 & 30.8 \\
Item B & 15.2 & 25.1 & 40.3 \\
Item C & 22.1 & 12.7 & 34.8 \\
\bottomrule
\end{tabular}
\caption{Sample Data Table}
\label{tab:sample}
\end{table}

\section{Lists}
\subsection{Bullet Points}
\begin{itemize}
  \item First item
  \item Second item
  \item Third item with sub-items:
    \begin{itemize}
      \item Sub-item 1
      \item Sub-item 2
    \end{itemize}
\end{itemize}

\subsection{Numbered List}
\begin{enumerate}
  \item First step
  \item Second step
  \item Third step
\end{enumerate}

\section{Citations and References}
When you use the citation search feature, references will be automatically managed. For example, you might cite a paper like this: \cite{example2024}. All citations will be automatically added to your \texttt{cite.bib} file.

\section{Conclusion}
This template demonstrates basic \LaTeX{} features including mathematical formulas, tables, and lists. The bibliography system is now set up to work seamlessly with Bibby AI's citation management. Feel free to modify it for your needs!

\bibliographystyle{plain}
\bibliography{cite}

\end{document}\documentclass{article}
\usepackage[utf8]{inputenc}
\usepackage{booktabs}
\usepackage{amsmath}
\usepackage{natbib}

\title{My Document}
\author{Your Name}
\date{\today}

\begin{document}

\maketitle

\section{Introduction}
This is a sample document demonstrating various \LaTeX{} features including mathematical equations, tables, and formatting. You can cite papers using the citation search feature, and they will be automatically added to your \texttt{cite.bib} file.

\section{Mathematical Expressions}
Here are some example equations:

\subsection{Quadratic Formula}
The quadratic formula is given by:
\[x = \frac{-b \pm \sqrt{b^2 - 4ac}}{2a}\]

\subsection{Matrix Example}
A sample matrix:
\[
\begin{pmatrix}
a_{11} & a_{12} & a_{13} \\
a_{21} & a_{22} & a_{23} \\
a_{31} & a_{32} & a_{33}
\end{pmatrix}
\]

\section{Tables}
Below is a sample table showing data:

\begin{table}[h]
\centering
\begin{tabular}{@{}lrrr@{}}
\toprule
\textbf{Category} & \textbf{Value 1} & \textbf{Value 2} & \textbf{Total} \\
\midrule
Item A & 10.5 & 20.3 & 30.8 \\
Item B & 15.2 & 25.1 & 40.3 \\
Item C & 22.1 & 12.7 & 34.8 \\
\bottomrule
\end{tabular}
\caption{Sample Data Table}
\label{tab:sample}
\end{table}

\section{Lists}
\subsection{Bullet Points}
\begin{itemize}
  \item First item
  \item Second item
  \item Third item with sub-items:
    \begin{itemize}
      \item Sub-item 1
      \item Sub-item 2
    \end{itemize}
\end{itemize}

\subsection{Numbered List}
\begin{enumerate}
  \item First step
  \item Second step
  \item Third step
\end{enumerate}

\section{Citations and References}
When you use the citation search feature, references will be automatically managed. For example, you might cite a paper like this: \cite{example2024}. All citations will be automatically added to your \texttt{cite.bib} file.

\section{Conclusion}
This template demonstrates basic \LaTeX{} features including mathematical formulas, tables, and lists. The bibliography system is now set up to work seamlessly with Bibby AI's citation management. Feel free to modify it for your needs!

\bibliographystyle{plain}
\bibliography{cite}

\end{document}\documentclass{article}
\usepackage[utf8]{inputenc}
\usepackage{booktabs}
\usepackage{amsmath}
\usepackage{natbib}

\title{My Document}
\author{Your Name}
\date{\today}

\begin{document}

\maketitle

\section{Introduction}
This is a sample document demonstrating various \LaTeX{} features including mathematical equations, tables, and formatting. You can cite papers using the citation search feature, and they will be automatically added to your \texttt{cite.bib} file.

\section{Mathematical Expressions}
Here are some example equations:

\subsection{Quadratic Formula}
The quadratic formula is given by:
\[x = \frac{-b \pm \sqrt{b^2 - 4ac}}{2a}\]

\subsection{Matrix Example}
A sample matrix:
\[
\begin{pmatrix}
a_{11} & a_{12} & a_{13} \\
a_{21} & a_{22} & a_{23} \\
a_{31} & a_{32} & a_{33}
\end{pmatrix}
\]

\section{Tables}
Below is a sample table showing data:

\begin{table}[h]
\centering
\begin{tabular}{@{}lrrr@{}}
\toprule
\textbf{Category} & \textbf{Value 1} & \textbf{Value 2} & \textbf{Total} \\
\midrule
Item A & 10.5 & 20.3 & 30.8 \\
Item B & 15.2 & 25.1 & 40.3 \\
Item C & 22.1 & 12.7 & 34.8 \\
\bottomrule
\end{tabular}
\caption{Sample Data Table}
\label{tab:sample}
\end{table}

\section{Lists}
\subsection{Bullet Points}
\begin{itemize}
  \item First item
  \item Second item
  \item Third item with sub-items:
    \begin{itemize}
      \item Sub-item 1
      \item Sub-item 2
    \end{itemize}
\end{itemize}

\subsection{Numbered List}
\begin{enumerate}
  \item First step
  \item Second step
  \item Third step
\end{enumerate}

\section{Citations and References}
When you use the citation search feature, references will be automatically managed. For example, you might cite a paper like this: \cite{example2024}. All citations will be automatically added to your \texttt{cite.bib} file.

\section{Conclusion}
This template demonstrates basic \LaTeX{} features including mathematical formulas, tables, and lists. The bibliography system is now set up to work seamlessly with Bibby AI's citation management. Feel free to modify it for your needs!

\bibliographystyle{plain}
\bibliography{cite}

\end{document}\documentclass{article}
\usepackage[utf8]{inputenc}
\usepackage{booktabs}
\usepackage{amsmath}
\usepackage{natbib}

\title{My Document}
\author{Your Name}
\date{\today}

\begin{document}

\maketitle

\section{Introduction}
This is a sample document demonstrating various \LaTeX{} features including mathematical equations, tables, and formatting. You can cite papers using the citation search feature, and they will be automatically added to your \texttt{cite.bib} file.

\section{Mathematical Expressions}
Here are some example equations:

\subsection{Quadratic Formula}
The quadratic formula is given by:
\[x = \frac{-b \pm \sqrt{b^2 - 4ac}}{2a}\]

\subsection{Matrix Example}
A sample matrix:
\[
\begin{pmatrix}
a_{11} & a_{12} & a_{13} \\
a_{21} & a_{22} & a_{23} \\
a_{31} & a_{32} & a_{33}
\end{pmatrix}
\]

\section{Tables}
Below is a sample table showing data:

\begin{table}[h]
\centering
\begin{tabular}{@{}lrrr@{}}
\toprule
\textbf{Category} & \textbf{Value 1} & \textbf{Value 2} & \textbf{Total} \\
\midrule
Item A & 10.5 & 20.3 & 30.8 \\
Item B & 15.2 & 25.1 & 40.3 \\
Item C & 22.1 & 12.7 & 34.8 \\
\bottomrule
\end{tabular}
\caption{Sample Data Table}
\label{tab:sample}
\end{table}

\section{Lists}
\subsection{Bullet Points}
\begin{itemize}
  \item First item
  \item Second item
  \item Third item with sub-items:
    \begin{itemize}
      \item Sub-item 1
      \item Sub-item 2
    \end{itemize}
\end{itemize}

\subsection{Numbered List}
\begin{enumerate}
  \item First step
  \item Second step
  \item Third step
\end{enumerate}

\section{Citations and References}
When you use the citation search feature, references will be automatically managed. For example, you might cite a paper like this: \cite{example2024}. All citations will be automatically added to your \texttt{cite.bib} file.

\section{Conclusion}
This template demonstrates basic \LaTeX{} features including mathematical formulas, tables, and lists. The bibliography system is now set up to work seamlessly with Bibby AI's citation management. Feel free to modify it for your needs!

\bibliographystyle{plain}
\bibliography{cite}

\end{document}\documentclass{article}
\usepackage[utf8]{inputenc}
\usepackage{booktabs}
\usepackage{amsmath}
\usepackage{natbib}

\title{My Document}
\author{Your Name}
\date{\today}

\begin{document}

\maketitle

\section{Introduction}
This is a sample document demonstrating various \LaTeX{} features including mathematical equations, tables, and formatting. You can cite papers using the citation search feature, and they will be automatically added to your \texttt{cite.bib} file.

\section{Mathematical Expressions}
Here are some example equations:

\subsection{Quadratic Formula}
The quadratic formula is given by:
\[x = \frac{-b \pm \sqrt{b^2 - 4ac}}{2a}\]

\subsection{Matrix Example}
A sample matrix:
\[
\begin{pmatrix}
a_{11} & a_{12} & a_{13} \\
a_{21} & a_{22} & a_{23} \\
a_{31} & a_{32} & a_{33}
\end{pmatrix}
\]

\section{Tables}
Below is a sample table showing data:

\begin{table}[h]
\centering
\begin{tabular}{@{}lrrr@{}}
\toprule
\textbf{Category} & \textbf{Value 1} & \textbf{Value 2} & \textbf{Total} \\
\midrule
Item A & 10.5 & 20.3 & 30.8 \\
Item B & 15.2 & 25.1 & 40.3 \\
Item C & 22.1 & 12.7 & 34.8 \\
\bottomrule
\end{tabular}
\caption{Sample Data Table}
\label{tab:sample}
\end{table}

\section{Lists}
\subsection{Bullet Points}
\begin{itemize}
  \item First item
  \item Second item
  \item Third item with sub-items:
    \begin{itemize}
      \item Sub-item 1
      \item Sub-item 2
    \end{itemize}
\end{itemize}

\subsection{Numbered List}
\begin{enumerate}
  \item First step
  \item Second step
  \item Third step
\end{enumerate}

\section{Citations and References}
When you use the citation search feature, references will be automatically managed. For example, you might cite a paper like this: \cite{example2024}. All citations will be automatically added to your \texttt{cite.bib} file.

\section{Conclusion}
This template demonstrates basic \LaTeX{} features including mathematical formulas, tables, and lists. The bibliography system is now set up to work seamlessly with Bibby AI's citation management. Feel free to modify it for your needs!

\bibliographystyle{plain}
\bibliography{cite}

\end{document}\documentclass{article}
\usepackage[utf8]{inputenc}
\usepackage{booktabs}
\usepackage{amsmath}
\usepackage{natbib}

\title{My Document}
\author{Your Name}
\date{\today}

\begin{document}

\maketitle

\section{Introduction}
This is a sample document demonstrating various \LaTeX{} features including mathematical equations, tables, and formatting. You can cite papers using the citation search feature, and they will be automatically added to your \texttt{cite.bib} file.

\section{Mathematical Expressions}
Here are some example equations:

\subsection{Quadratic Formula}
The quadratic formula is given by:
\[x = \frac{-b \pm \sqrt{b^2 - 4ac}}{2a}\]

\subsection{Matrix Example}
A sample matrix:
\[
\begin{pmatrix}
a_{11} & a_{12} & a_{13} \\
a_{21} & a_{22} & a_{23} \\
a_{31} & a_{32} & a_{33}
\end{pmatrix}
\]

\section{Tables}
Below is a sample table showing data:

\begin{table}[h]
\centering
\begin{tabular}{@{}lrrr@{}}
\toprule
\textbf{Category} & \textbf{Value 1} & \textbf{Value 2} & \textbf{Total} \\
\midrule
Item A & 10.5 & 20.3 & 30.8 \\
Item B & 15.2 & 25.1 & 40.3 \\
Item C & 22.1 & 12.7 & 34.8 \\
\bottomrule
\end{tabular}
\caption{Sample Data Table}
\label{tab:sample}
\end{table}

\section{Lists}
\subsection{Bullet Points}
\begin{itemize}
  \item First item
  \item Second item
  \item Third item with sub-items:
    \begin{itemize}
      \item Sub-item 1
      \item Sub-item 2
    \end{itemize}
\end{itemize}

\subsection{Numbered List}
\begin{enumerate}
  \item First step
  \item Second step
  \item Third step
\end{enumerate}

\section{Citations and References}
When you use the citation search feature, references will be automatically managed. For example, you might cite a paper like this: \cite{example2024}. All citations will be automatically added to your \texttt{cite.bib} file.

\section{Conclusion}
This template demonstrates basic \LaTeX{} features including mathematical formulas, tables, and lists. The bibliography system is now set up to work seamlessly with Bibby AI's citation management. Feel free to modify it for your needs!

\bibliographystyle{plain}
\bibliography{cite}

\end{document}\documentclass{article}
\usepackage[utf8]{inputenc}
\usepackage{booktabs}
\usepackage{amsmath}
\usepackage{natbib}

\title{My Document}
\author{Your Name}
\date{\today}

\begin{document}

\maketitle

\section{Introduction}
This is a sample document demonstrating various \LaTeX{} features including mathematical equations, tables, and formatting. You can cite papers using the citation search feature, and they will be automatically added to your \texttt{cite.bib} file.

\section{Mathematical Expressions}
Here are some example equations:

\subsection{Quadratic Formula}
The quadratic formula is given by:
\[x = \frac{-b \pm \sqrt{b^2 - 4ac}}{2a}\]

\subsection{Matrix Example}
A sample matrix:
\[
\begin{pmatrix}
a_{11} & a_{12} & a_{13} \\
a_{21} & a_{22} & a_{23} \\
a_{31} & a_{32} & a_{33}
\end{pmatrix}
\]

\section{Tables}
Below is a sample table showing data:

\begin{table}[h]
\centering
\begin{tabular}{@{}lrrr@{}}
\toprule
\textbf{Category} & \textbf{Value 1} & \textbf{Value 2} & \textbf{Total} \\
\midrule
Item A & 10.5 & 20.3 & 30.8 \\
Item B & 15.2 & 25.1 & 40.3 \\
Item C & 22.1 & 12.7 & 34.8 \\
\bottomrule
\end{tabular}
\caption{Sample Data Table}
\label{tab:sample}
\end{table}

\section{Lists}
\subsection{Bullet Points}
\begin{itemize}
  \item First item
  \item Second item
  \item Third item with sub-items:
    \begin{itemize}
      \item Sub-item 1
      \item Sub-item 2
    \end{itemize}
\end{itemize}

\subsection{Numbered List}
\begin{enumerate}
  \item First step
  \item Second step
  \item Third step
\end{enumerate}

\section{Citations and References}
When you use the citation search feature, references will be automatically managed. For example, you might cite a paper like this: \cite{example2024}. All citations will be automatically added to your \texttt{cite.bib} file.

\section{Conclusion}
This template demonstrates basic \LaTeX{} features including mathematical formulas, tables, and lists. The bibliography system is now set up to work seamlessly with Bibby AI's citation management. Feel free to modify it for your needs!

\bibliographystyle{plain}
\bibliography{cite}

\end{document}\documentclass{article}
\usepackage[utf8]{inputenc}
\usepackage{booktabs}
\usepackage{amsmath}
\usepackage{natbib}

\title{My Document}
\author{Your Name}
\date{\today}

\begin{document}

\maketitle

\section{Introduction}
This is a sample document demonstrating various \LaTeX{} features including mathematical equations, tables, and formatting. You can cite papers using the citation search feature, and they will be automatically added to your \texttt{cite.bib} file.

\section{Mathematical Expressions}
Here are some example equations:

\subsection{Quadratic Formula}
The quadratic formula is given by:
\[x = \frac{-b \pm \sqrt{b^2 - 4ac}}{2a}\]

\subsection{Matrix Example}
A sample matrix:
\[
\begin{pmatrix}
a_{11} & a_{12} & a_{13} \\
a_{21} & a_{22} & a_{23} \\
a_{31} & a_{32} & a_{33}
\end{pmatrix}
\]

\section{Tables}
Below is a sample table showing data:

\begin{table}[h]
\centering
\begin{tabular}{@{}lrrr@{}}
\toprule
\textbf{Category} & \textbf{Value 1} & \textbf{Value 2} & \textbf{Total} \\
\midrule
Item A & 10.5 & 20.3 & 30.8 \\
Item B & 15.2 & 25.1 & 40.3 \\
Item C & 22.1 & 12.7 & 34.8 \\
\bottomrule
\end{tabular}
\caption{Sample Data Table}
\label{tab:sample}
\end{table}

\section{Lists}
\subsection{Bullet Points}
\begin{itemize}
  \item First item
  \item Second item
  \item Third item with sub-items:
    \begin{itemize}
      \item Sub-item 1
      \item Sub-item 2
    \end{itemize}
\end{itemize}

\subsection{Numbered List}
\begin{enumerate}
  \item First step
  \item Second step
  \item Third step
\end{enumerate}

\section{Citations and References}
When you use the citation search feature, references will be automatically managed. For example, you might cite a paper like this: \cite{example2024}. All citations will be automatically added to your \texttt{cite.bib} file.

\section{Conclusion}
This template demonstrates basic \LaTeX{} features including mathematical formulas, tables, and lists. The bibliography system is now set up to work seamlessly with Bibby AI's citation management. Feel free to modify it for your needs!

\bibliographystyle{plain}
\bibliography{cite}

\end{document}\documentclass{article}
\usepackage[utf8]{inputenc}
\usepackage{booktabs}
\usepackage{amsmath}
\usepackage{natbib}

\title{My Document}
\author{Your Name}
\date{\today}

\begin{document}

\maketitle

\section{Introduction}
This is a sample document demonstrating various \LaTeX{} features including mathematical equations, tables, and formatting. You can cite papers using the citation search feature, and they will be automatically added to your \texttt{cite.bib} file.

\section{Mathematical Expressions}
Here are some example equations:

\subsection{Quadratic Formula}
The quadratic formula is given by:
\[x = \frac{-b \pm \sqrt{b^2 - 4ac}}{2a}\]

\subsection{Matrix Example}
A sample matrix:
\[
\begin{pmatrix}
a_{11} & a_{12} & a_{13} \\
a_{21} & a_{22} & a_{23} \\
a_{31} & a_{32} & a_{33}
\end{pmatrix}
\]

\section{Tables}
Below is a sample table showing data:

\begin{table}[h]
\centering
\begin{tabular}{@{}lrrr@{}}
\toprule
\textbf{Category} & \textbf{Value 1} & \textbf{Value 2} & \textbf{Total} \\
\midrule
Item A & 10.5 & 20.3 & 30.8 \\
Item B & 15.2 & 25.1 & 40.3 \\
Item C & 22.1 & 12.7 & 34.8 \\
\bottomrule
\end{tabular}
\caption{Sample Data Table}
\label{tab:sample}
\end{table}

\section{Lists}
\subsection{Bullet Points}
\begin{itemize}
  \item First item
  \item Second item
  \item Third item with sub-items:
    \begin{itemize}
      \item Sub-item 1
      \item Sub-item 2
    \end{itemize}
\end{itemize}

\subsection{Numbered List}
\begin{enumerate}
  \item First step
  \item Second step
  \item Third step
\end{enumerate}

\section{Citations and References}
When you use the citation search feature, references will be automatically managed. For example, you might cite a paper like this: \cite{example2024}. All citations will be automatically added to your \texttt{cite.bib} file.

\section{Conclusion}
This template demonstrates basic \LaTeX{} features including mathematical formulas, tables, and lists. The bibliography system is now set up to work seamlessly with Bibby AI's citation management. Feel free to modify it for your needs!

\bibliographystyle{plain}
\bibliography{cite}

\end{document}\documentclass{article}
\usepackage[utf8]{inputenc}
\usepackage{booktabs}
\usepackage{amsmath}
\usepackage{natbib}

\title{My Document}
\author{Your Name}
\date{\today}

\begin{document}

\maketitle

\section{Introduction}
This is a sample document demonstrating various \LaTeX{} features including mathematical equations, tables, and formatting. You can cite papers using the citation search feature, and they will be automatically added to your \texttt{cite.bib} file.

\section{Mathematical Expressions}
Here are some example equations:

\subsection{Quadratic Formula}
The quadratic formula is given by:
\[x = \frac{-b \pm \sqrt{b^2 - 4ac}}{2a}\]

\subsection{Matrix Example}
A sample matrix:
\[
\begin{pmatrix}
a_{11} & a_{12} & a_{13} \\
a_{21} & a_{22} & a_{23} \\
a_{31} & a_{32} & a_{33}
\end{pmatrix}
\]

\section{Tables}
Below is a sample table showing data:

\begin{table}[h]
\centering
\begin{tabular}{@{}lrrr@{}}
\toprule
\textbf{Category} & \textbf{Value 1} & \textbf{Value 2} & \textbf{Total} \\
\midrule
Item A & 10.5 & 20.3 & 30.8 \\
Item B & 15.2 & 25.1 & 40.3 \\
Item C & 22.1 & 12.7 & 34.8 \\
\bottomrule
\end{tabular}
\caption{Sample Data Table}
\label{tab:sample}
\end{table}

\section{Lists}
\subsection{Bullet Points}
\begin{itemize}
  \item First item
  \item Second item
  \item Third item with sub-items:
    \begin{itemize}
      \item Sub-item 1
      \item Sub-item 2
    \end{itemize}
\end{itemize}

\subsection{Numbered List}
\begin{enumerate}
  \item First step
  \item Second step
  \item Third step
\end{enumerate}

\section{Citations and References}
When you use the citation search feature, references will be automatically managed. For example, you might cite a paper like this: \cite{example2024}. All citations will be automatically added to your \texttt{cite.bib} file.

\section{Conclusion}
This template demonstrates basic \LaTeX{} features including mathematical formulas, tables, and lists. The bibliography system is now set up to work seamlessly with Bibby AI's citation management. Feel free to modify it for your needs!

\bibliographystyle{plain}
\bibliography{cite}

\end{document}\documentclass{article}
\usepackage[utf8]{inputenc}
\usepackage{booktabs}
\usepackage{amsmath}
\usepackage{natbib}

\title{My Document}
\author{Your Name}
\date{\today}

\begin{document}

\maketitle

\section{Introduction}
This is a sample document demonstrating various \LaTeX{} features including mathematical equations, tables, and formatting. You can cite papers using the citation search feature, and they will be automatically added to your \texttt{cite.bib} file.

\section{Mathematical Expressions}
Here are some example equations:

\subsection{Quadratic Formula}
The quadratic formula is given by:
\[x = \frac{-b \pm \sqrt{b^2 - 4ac}}{2a}\]

\subsection{Matrix Example}
A sample matrix:
\[
\begin{pmatrix}
a_{11} & a_{12} & a_{13} \\
a_{21} & a_{22} & a_{23} \\
a_{31} & a_{32} & a_{33}
\end{pmatrix}
\]

\section{Tables}
Below is a sample table showing data:

\begin{table}[h]
\centering
\begin{tabular}{@{}lrrr@{}}
\toprule
\textbf{Category} & \textbf{Value 1} & \textbf{Value 2} & \textbf{Total} \\
\midrule
Item A & 10.5 & 20.3 & 30.8 \\
Item B & 15.2 & 25.1 & 40.3 \\
Item C & 22.1 & 12.7 & 34.8 \\
\bottomrule
\end{tabular}
\caption{Sample Data Table}
\label{tab:sample}
\end{table}

\section{Lists}
\subsection{Bullet Points}
\begin{itemize}
  \item First item
  \item Second item
  \item Third item with sub-items:
    \begin{itemize}
      \item Sub-item 1
      \item Sub-item 2
    \end{itemize}
\end{itemize}

\subsection{Numbered List}
\begin{enumerate}
  \item First step
  \item Second step
  \item Third step
\end{enumerate}

\section{Citations and References}
When you use the citation search feature, references will be automatically managed. For example, you might cite a paper like this: \cite{example2024}. All citations will be automatically added to your \texttt{cite.bib} file.

\section{Conclusion}
This template demonstrates basic \LaTeX{} features including mathematical formulas, tables, and lists. The bibliography system is now set up to work seamlessly with Bibby AI's citation management. Feel free to modify it for your needs!

\bibliographystyle{plain}
\bibliography{cite}

\end{document}\documentclass{article}
\usepackage[utf8]{inputenc}
\usepackage{booktabs}
\usepackage{amsmath}
\usepackage{natbib}

\title{My Document}
\author{Your Name}
\date{\today}

\begin{document}

\maketitle

\section{Introduction}
This is a sample document demonstrating various \LaTeX{} features including mathematical equations, tables, and formatting. You can cite papers using the citation search feature, and they will be automatically added to your \texttt{cite.bib} file.

\section{Mathematical Expressions}
Here are some example equations:

\subsection{Quadratic Formula}
The quadratic formula is given by:
\[x = \frac{-b \pm \sqrt{b^2 - 4ac}}{2a}\]

\subsection{Matrix Example}
A sample matrix:
\[
\begin{pmatrix}
a_{11} & a_{12} & a_{13} \\
a_{21} & a_{22} & a_{23} \\
a_{31} & a_{32} & a_{33}
\end{pmatrix}
\]

\section{Tables}
Below is a sample table showing data:

\begin{table}[h]
\centering
\begin{tabular}{@{}lrrr@{}}
\toprule
\textbf{Category} & \textbf{Value 1} & \textbf{Value 2} & \textbf{Total} \\
\midrule
Item A & 10.5 & 20.3 & 30.8 \\
Item B & 15.2 & 25.1 & 40.3 \\
Item C & 22.1 & 12.7 & 34.8 \\
\bottomrule
\end{tabular}
\caption{Sample Data Table}
\label{tab:sample}
\end{table}

\section{Lists}
\subsection{Bullet Points}
\begin{itemize}
  \item First item
  \item Second item
  \item Third item with sub-items:
    \begin{itemize}
      \item Sub-item 1
      \item Sub-item 2
    \end{itemize}
\end{itemize}

\subsection{Numbered List}
\begin{enumerate}
  \item First step
  \item Second step
  \item Third step
\end{enumerate}

\section{Citations and References}
When you use the citation search feature, references will be automatically managed. For example, you might cite a paper like this: \cite{example2024}. All citations will be automatically added to your \texttt{cite.bib} file.

\section{Conclusion}
This template demonstrates basic \LaTeX{} features including mathematical formulas, tables, and lists. The bibliography system is now set up to work seamlessly with Bibby AI's citation management. Feel free to modify it for your needs!

\bibliographystyle{plain}
\bibliography{cite}

\end{document}\documentclass{article}
\usepackage[utf8]{inputenc}
\usepackage{booktabs}
\usepackage{amsmath}
\usepackage{natbib}

\title{My Document}
\author{Your Name}
\date{\today}

\begin{document}

\maketitle

\section{Introduction}
This is a sample document demonstrating various \LaTeX{} features including mathematical equations, tables, and formatting. You can cite papers using the citation search feature, and they will be automatically added to your \texttt{cite.bib} file.

\section{Mathematical Expressions}
Here are some example equations:

\subsection{Quadratic Formula}
The quadratic formula is given by:
\[x = \frac{-b \pm \sqrt{b^2 - 4ac}}{2a}\]

\subsection{Matrix Example}
A sample matrix:
\[
\begin{pmatrix}
a_{11} & a_{12} & a_{13} \\
a_{21} & a_{22} & a_{23} \\
a_{31} & a_{32} & a_{33}
\end{pmatrix}
\]

\section{Tables}
Below is a sample table showing data:

\begin{table}[h]
\centering
\begin{tabular}{@{}lrrr@{}}
\toprule
\textbf{Category} & \textbf{Value 1} & \textbf{Value 2} & \textbf{Total} \\
\midrule
Item A & 10.5 & 20.3 & 30.8 \\
Item B & 15.2 & 25.1 & 40.3 \\
Item C & 22.1 & 12.7 & 34.8 \\
\bottomrule
\end{tabular}
\caption{Sample Data Table}
\label{tab:sample}
\end{table}

\section{Lists}
\subsection{Bullet Points}
\begin{itemize}
  \item First item
  \item Second item
  \item Third item with sub-items:
    \begin{itemize}
      \item Sub-item 1
      \item Sub-item 2
    \end{itemize}
\end{itemize}

\subsection{Numbered List}
\begin{enumerate}
  \item First step
  \item Second step
  \item Third step
\end{enumerate}

\section{Citations and References}
When you use the citation search feature, references will be automatically managed. For example, you might cite a paper like this: \cite{example2024}. All citations will be automatically added to your \texttt{cite.bib} file.

\section{Conclusion}
This template demonstrates basic \LaTeX{} features including mathematical formulas, tables, and lists. The bibliography system is now set up to work seamlessly with Bibby AI's citation management. Feel free to modify it for your needs!

\bibliographystyle{plain}
\bibliography{cite}

\end{document}\documentclass{article}
\usepackage[utf8]{inputenc}
\usepackage{booktabs}
\usepackage{amsmath}
\usepackage{natbib}

\title{My Document}
\author{Your Name}
\date{\today}

\begin{document}

\maketitle

\section{Introduction}
This is a sample document demonstrating various \LaTeX{} features including mathematical equations, tables, and formatting. You can cite papers using the citation search feature, and they will be automatically added to your \texttt{cite.bib} file.

\section{Mathematical Expressions}
Here are some example equations:

\subsection{Quadratic Formula}
The quadratic formula is given by:
\[x = \frac{-b \pm \sqrt{b^2 - 4ac}}{2a}\]

\subsection{Matrix Example}
A sample matrix:
\[
\begin{pmatrix}
a_{11} & a_{12} & a_{13} \\
a_{21} & a_{22} & a_{23} \\
a_{31} & a_{32} & a_{33}
\end{pmatrix}
\]

\section{Tables}
Below is a sample table showing data:

\begin{table}[h]
\centering
\begin{tabular}{@{}lrrr@{}}
\toprule
\textbf{Category} & \textbf{Value 1} & \textbf{Value 2} & \textbf{Total} \\
\midrule
Item A & 10.5 & 20.3 & 30.8 \\
Item B & 15.2 & 25.1 & 40.3 \\
Item C & 22.1 & 12.7 & 34.8 \\
\bottomrule
\end{tabular}
\caption{Sample Data Table}
\label{tab:sample}
\end{table}

\section{Lists}
\subsection{Bullet Points}
\begin{itemize}
  \item First item
  \item Second item
  \item Third item with sub-items:
    \begin{itemize}
      \item Sub-item 1
      \item Sub-item 2
    \end{itemize}
\end{itemize}

\subsection{Numbered List}
\begin{enumerate}
  \item First step
  \item Second step
  \item Third step
\end{enumerate}

\section{Citations and References}
When you use the citation search feature, references will be automatically managed. For example, you might cite a paper like this: \cite{example2024}. All citations will be automatically added to your \texttt{cite.bib} file.

\section{Conclusion}
This template demonstrates basic \LaTeX{} features including mathematical formulas, tables, and lists. The bibliography system is now set up to work seamlessly with Bibby AI's citation management. Feel free to modify it for your needs!

\bibliographystyle{plain}
\bibliography{cite}

\end{document}\documentclass{article}
\usepackage[utf8]{inputenc}
\usepackage{booktabs}
\usepackage{amsmath}
\usepackage{natbib}

\title{My Document}
\author{Your Name}
\date{\today}

\begin{document}

\maketitle

\section{Introduction}
This is a sample document demonstrating various \LaTeX{} features including mathematical equations, tables, and formatting. You can cite papers using the citation search feature, and they will be automatically added to your \texttt{cite.bib} file.

\section{Mathematical Expressions}
Here are some example equations:

\subsection{Quadratic Formula}
The quadratic formula is given by:
\[x = \frac{-b \pm \sqrt{b^2 - 4ac}}{2a}\]

\subsection{Matrix Example}
A sample matrix:
\[
\begin{pmatrix}
a_{11} & a_{12} & a_{13} \\
a_{21} & a_{22} & a_{23} \\
a_{31} & a_{32} & a_{33}
\end{pmatrix}
\]

\section{Tables}
Below is a sample table showing data:

\begin{table}[h]
\centering
\begin{tabular}{@{}lrrr@{}}
\toprule
\textbf{Category} & \textbf{Value 1} & \textbf{Value 2} & \textbf{Total} \\
\midrule
Item A & 10.5 & 20.3 & 30.8 \\
Item B & 15.2 & 25.1 & 40.3 \\
Item C & 22.1 & 12.7 & 34.8 \\
\bottomrule
\end{tabular}
\caption{Sample Data Table}
\label{tab:sample}
\end{table}

\section{Lists}
\subsection{Bullet Points}
\begin{itemize}
  \item First item
  \item Second item
  \item Third item with sub-items:
    \begin{itemize}
      \item Sub-item 1
      \item Sub-item 2
    \end{itemize}
\end{itemize}

\subsection{Numbered List}
\begin{enumerate}
  \item First step
  \item Second step
  \item Third step
\end{enumerate}

\section{Citations and References}
When you use the citation search feature, references will be automatically managed. For example, you might cite a paper like this: \cite{example2024}. All citations will be automatically added to your \texttt{cite.bib} file.

\section{Conclusion}
This template demonstrates basic \LaTeX{} features including mathematical formulas, tables, and lists. The bibliography system is now set up to work seamlessly with Bibby AI's citation management. Feel free to modify it for your needs!

\bibliographystyle{plain}
\bibliography{cite}

\end{document}\documentclass{article}
\usepackage[utf8]{inputenc}
\usepackage{booktabs}
\usepackage{amsmath}
\usepackage{natbib}

\title{My Document}
\author{Your Name}
\date{\today}

\begin{document}

\maketitle

\section{Introduction}
This is a sample document demonstrating various \LaTeX{} features including mathematical equations, tables, and formatting. You can cite papers using the citation search feature, and they will be automatically added to your \texttt{cite.bib} file.

\section{Mathematical Expressions}
Here are some example equations:

\subsection{Quadratic Formula}
The quadratic formula is given by:
\[x = \frac{-b \pm \sqrt{b^2 - 4ac}}{2a}\]

\subsection{Matrix Example}
A sample matrix:
\[
\begin{pmatrix}
a_{11} & a_{12} & a_{13} \\
a_{21} & a_{22} & a_{23} \\
a_{31} & a_{32} & a_{33}
\end{pmatrix}
\]

\section{Tables}
Below is a sample table showing data:

\begin{table}[h]
\centering
\begin{tabular}{@{}lrrr@{}}
\toprule
\textbf{Category} & \textbf{Value 1} & \textbf{Value 2} & \textbf{Total} \\
\midrule
Item A & 10.5 & 20.3 & 30.8 \\
Item B & 15.2 & 25.1 & 40.3 \\
Item C & 22.1 & 12.7 & 34.8 \\
\bottomrule
\end{tabular}
\caption{Sample Data Table}
\label{tab:sample}
\end{table}

\section{Lists}
\subsection{Bullet Points}
\begin{itemize}
  \item First item
  \item Second item
  \item Third item with sub-items:
    \begin{itemize}
      \item Sub-item 1
      \item Sub-item 2
    \end{itemize}
\end{itemize}

\subsection{Numbered List}
\begin{enumerate}
  \item First step
  \item Second step
  \item Third step
\end{enumerate}

\section{Citations and References}
When you use the citation search feature, references will be automatically managed. For example, you might cite a paper like this: \cite{example2024}. All citations will be automatically added to your \texttt{cite.bib} file.

\section{Conclusion}
This template demonstrates basic \LaTeX{} features including mathematical formulas, tables, and lists. The bibliography system is now set up to work seamlessly with Bibby AI's citation management. Feel free to modify it for your needs!

\bibliographystyle{plain}
\bibliography{cite}

\end{document}\documentclass{article}
\usepackage[utf8]{inputenc}
\usepackage{booktabs}
\usepackage{amsmath}
\usepackage{natbib}

\title{My Document}
\author{Your Name}
\date{\today}

\begin{document}

\maketitle

\section{Introduction}
This is a sample document demonstrating various \LaTeX{} features including mathematical equations, tables, and formatting. You can cite papers using the citation search feature, and they will be automatically added to your \texttt{cite.bib} file.

\section{Mathematical Expressions}
Here are some example equations:

\subsection{Quadratic Formula}
The quadratic formula is given by:
\[x = \frac{-b \pm \sqrt{b^2 - 4ac}}{2a}\]

\subsection{Matrix Example}
A sample matrix:
\[
\begin{pmatrix}
a_{11} & a_{12} & a_{13} \\
a_{21} & a_{22} & a_{23} \\
a_{31} & a_{32} & a_{33}
\end{pmatrix}
\]

\section{Tables}
Below is a sample table showing data:

\begin{table}[h]
\centering
\begin{tabular}{@{}lrrr@{}}
\toprule
\textbf{Category} & \textbf{Value 1} & \textbf{Value 2} & \textbf{Total} \\
\midrule
Item A & 10.5 & 20.3 & 30.8 \\
Item B & 15.2 & 25.1 & 40.3 \\
Item C & 22.1 & 12.7 & 34.8 \\
\bottomrule
\end{tabular}
\caption{Sample Data Table}
\label{tab:sample}
\end{table}

\section{Lists}
\subsection{Bullet Points}
\begin{itemize}
  \item First item
  \item Second item
  \item Third item with sub-items:
    \begin{itemize}
      \item Sub-item 1
      \item Sub-item 2
    \end{itemize}
\end{itemize}

\subsection{Numbered List}
\begin{enumerate}
  \item First step
  \item Second step
  \item Third step
\end{enumerate}

\section{Citations and References}
When you use the citation search feature, references will be automatically managed. For example, you might cite a paper like this: \cite{example2024}. All citations will be automatically added to your \texttt{cite.bib} file.

\section{Conclusion}
This template demonstrates basic \LaTeX{} features including mathematical formulas, tables, and lists. The bibliography system is now set up to work seamlessly with Bibby AI's citation management. Feel free to modify it for your needs!

\bibliographystyle{plain}
\bibliography{cite}

\end{document}\documentclass{article}
\usepackage[utf8]{inputenc}
\usepackage{booktabs}
\usepackage{amsmath}
\usepackage{natbib}

\title{My Document}
\author{Your Name}
\date{\today}

\begin{document}

\maketitle

\section{Introduction}
This is a sample document demonstrating various \LaTeX{} features including mathematical equations, tables, and formatting. You can cite papers using the citation search feature, and they will be automatically added to your \texttt{cite.bib} file.

\section{Mathematical Expressions}
Here are some example equations:

\subsection{Quadratic Formula}
The quadratic formula is given by:
\[x = \frac{-b \pm \sqrt{b^2 - 4ac}}{2a}\]

\subsection{Matrix Example}
A sample matrix:
\[
\begin{pmatrix}
a_{11} & a_{12} & a_{13} \\
a_{21} & a_{22} & a_{23} \\
a_{31} & a_{32} & a_{33}
\end{pmatrix}
\]

\section{Tables}
Below is a sample table showing data:

\begin{table}[h]
\centering
\begin{tabular}{@{}lrrr@{}}
\toprule
\textbf{Category} & \textbf{Value 1} & \textbf{Value 2} & \textbf{Total} \\
\midrule
Item A & 10.5 & 20.3 & 30.8 \\
Item B & 15.2 & 25.1 & 40.3 \\
Item C & 22.1 & 12.7 & 34.8 \\
\bottomrule
\end{tabular}
\caption{Sample Data Table}
\label{tab:sample}
\end{table}

\section{Lists}
\subsection{Bullet Points}
\begin{itemize}
  \item First item
  \item Second item
  \item Third item with sub-items:
    \begin{itemize}
      \item Sub-item 1
      \item Sub-item 2
    \end{itemize}
\end{itemize}

\subsection{Numbered List}
\begin{enumerate}
  \item First step
  \item Second step
  \item Third step
\end{enumerate}

\section{Citations and References}
When you use the citation search feature, references will be automatically managed. For example, you might cite a paper like this: \cite{example2024}. All citations will be automatically added to your \texttt{cite.bib} file.

\section{Conclusion}
This template demonstrates basic \LaTeX{} features including mathematical formulas, tables, and lists. The bibliography system is now set up to work seamlessly with Bibby AI's citation management. Feel free to modify it for your needs!

\bibliographystyle{plain}
\bibliography{cite}

\end{document}\documentclass{article}
\usepackage[utf8]{inputenc}
\usepackage{booktabs}
\usepackage{amsmath}
\usepackage{natbib}

\title{My Document}
\author{Your Name}
\date{\today}

\begin{document}

\maketitle

\section{Introduction}
This is a sample document demonstrating various \LaTeX{} features including mathematical equations, tables, and formatting. You can cite papers using the citation search feature, and they will be automatically added to your \texttt{cite.bib} file.

\section{Mathematical Expressions}
Here are some example equations:

\subsection{Quadratic Formula}
The quadratic formula is given by:
\[x = \frac{-b \pm \sqrt{b^2 - 4ac}}{2a}\]

\subsection{Matrix Example}
A sample matrix:
\[
\begin{pmatrix}
a_{11} & a_{12} & a_{13} \\
a_{21} & a_{22} & a_{23} \\
a_{31} & a_{32} & a_{33}
\end{pmatrix}
\]

\section{Tables}
Below is a sample table showing data:

\begin{table}[h]
\centering
\begin{tabular}{@{}lrrr@{}}
\toprule
\textbf{Category} & \textbf{Value 1} & \textbf{Value 2} & \textbf{Total} \\
\midrule
Item A & 10.5 & 20.3 & 30.8 \\
Item B & 15.2 & 25.1 & 40.3 \\
Item C & 22.1 & 12.7 & 34.8 \\
\bottomrule
\end{tabular}
\caption{Sample Data Table}
\label{tab:sample}
\end{table}

\section{Lists}
\subsection{Bullet Points}
\begin{itemize}
  \item First item
  \item Second item
  \item Third item with sub-items:
    \begin{itemize}
      \item Sub-item 1
      \item Sub-item 2
    \end{itemize}
\end{itemize}

\subsection{Numbered List}
\begin{enumerate}
  \item First step
  \item Second step
  \item Third step
\end{enumerate}

\section{Citations and References}
When you use the citation search feature, references will be automatically managed. For example, you might cite a paper like this: \cite{example2024}. All citations will be automatically added to your \texttt{cite.bib} file.

\section{Conclusion}
This template demonstrates basic \LaTeX{} features including mathematical formulas, tables, and lists. The bibliography system is now set up to work seamlessly with Bibby AI's citation management. Feel free to modify it for your needs!

\bibliographystyle{plain}
\bibliography{cite}

\end{document}\documentclass{article}
\usepackage[utf8]{inputenc}
\usepackage{booktabs}
\usepackage{amsmath}
\usepackage{natbib}

\title{My Document}
\author{Your Name}
\date{\today}

\begin{document}

\maketitle

\section{Introduction}
This is a sample document demonstrating various \LaTeX{} features including mathematical equations, tables, and formatting. You can cite papers using the citation search feature, and they will be automatically added to your \texttt{cite.bib} file.

\section{Mathematical Expressions}
Here are some example equations:

\subsection{Quadratic Formula}
The quadratic formula is given by:
\[x = \frac{-b \pm \sqrt{b^2 - 4ac}}{2a}\]

\subsection{Matrix Example}
A sample matrix:
\[
\begin{pmatrix}
a_{11} & a_{12} & a_{13} \\
a_{21} & a_{22} & a_{23} \\
a_{31} & a_{32} & a_{33}
\end{pmatrix}
\]

\section{Tables}
Below is a sample table showing data:

\begin{table}[h]
\centering
\begin{tabular}{@{}lrrr@{}}
\toprule
\textbf{Category} & \textbf{Value 1} & \textbf{Value 2} & \textbf{Total} \\
\midrule
Item A & 10.5 & 20.3 & 30.8 \\
Item B & 15.2 & 25.1 & 40.3 \\
Item C & 22.1 & 12.7 & 34.8 \\
\bottomrule
\end{tabular}
\caption{Sample Data Table}
\label{tab:sample}
\end{table}

\section{Lists}
\subsection{Bullet Points}
\begin{itemize}
  \item First item
  \item Second item
  \item Third item with sub-items:
    \begin{itemize}
      \item Sub-item 1
      \item Sub-item 2
    \end{itemize}
\end{itemize}

\subsection{Numbered List}
\begin{enumerate}
  \item First step
  \item Second step
  \item Third step
\end{enumerate}

\section{Citations and References}
When you use the citation search feature, references will be automatically managed. For example, you might cite a paper like this: \cite{example2024}. All citations will be automatically added to your \texttt{cite.bib} file.

\section{Conclusion}
This template demonstrates basic \LaTeX{} features including mathematical formulas, tables, and lists. The bibliography system is now set up to work seamlessly with Bibby AI's citation management. Feel free to modify it for your needs!

\bibliographystyle{plain}
\bibliography{cite}

\end{document}\documentclass{article}
\usepackage[utf8]{inputenc}
\usepackage{booktabs}
\usepackage{amsmath}
\usepackage{natbib}

\title{My Document}
\author{Your Name}
\date{\today}

\begin{document}

\maketitle

\section{Introduction}
This is a sample document demonstrating various \LaTeX{} features including mathematical equations, tables, and formatting. You can cite papers using the citation search feature, and they will be automatically added to your \texttt{cite.bib} file.

\section{Mathematical Expressions}
Here are some example equations:

\subsection{Quadratic Formula}
The quadratic formula is given by:
\[x = \frac{-b \pm \sqrt{b^2 - 4ac}}{2a}\]

\subsection{Matrix Example}
A sample matrix:
\[
\begin{pmatrix}
a_{11} & a_{12} & a_{13} \\
a_{21} & a_{22} & a_{23} \\
a_{31} & a_{32} & a_{33}
\end{pmatrix}
\]

\section{Tables}
Below is a sample table showing data:

\begin{table}[h]
\centering
\begin{tabular}{@{}lrrr@{}}
\toprule
\textbf{Category} & \textbf{Value 1} & \textbf{Value 2} & \textbf{Total} \\
\midrule
Item A & 10.5 & 20.3 & 30.8 \\
Item B & 15.2 & 25.1 & 40.3 \\
Item C & 22.1 & 12.7 & 34.8 \\
\bottomrule
\end{tabular}
\caption{Sample Data Table}
\label{tab:sample}
\end{table}

\section{Lists}
\subsection{Bullet Points}
\begin{itemize}
  \item First item
  \item Second item
  \item Third item with sub-items:
    \begin{itemize}
      \item Sub-item 1
      \item Sub-item 2
    \end{itemize}
\end{itemize}

\subsection{Numbered List}
\begin{enumerate}
  \item First step
  \item Second step
  \item Third step
\end{enumerate}

\section{Citations and References}
When you use the citation search feature, references will be automatically managed. For example, you might cite a paper like this: \cite{example2024}. All citations will be automatically added to your \texttt{cite.bib} file.

\section{Conclusion}
This template demonstrates basic \LaTeX{} features including mathematical formulas, tables, and lists. The bibliography system is now set up to work seamlessly with Bibby AI's citation management. Feel free to modify it for your needs!

\bibliographystyle{plain}
\bibliography{cite}

\end{document}\documentclass{article}
\usepackage[utf8]{inputenc}
\usepackage{booktabs}
\usepackage{amsmath}
\usepackage{natbib}

\title{My Document}
\author{Your Name}
\date{\today}

\begin{document}

\maketitle

\section{Introduction}
This is a sample document demonstrating various \LaTeX{} features including mathematical equations, tables, and formatting. You can cite papers using the citation search feature, and they will be automatically added to your \texttt{cite.bib} file.

\section{Mathematical Expressions}
Here are some example equations:

\subsection{Quadratic Formula}
The quadratic formula is given by:
\[x = \frac{-b \pm \sqrt{b^2 - 4ac}}{2a}\]

\subsection{Matrix Example}
A sample matrix:
\[
\begin{pmatrix}
a_{11} & a_{12} & a_{13} \\
a_{21} & a_{22} & a_{23} \\
a_{31} & a_{32} & a_{33}
\end{pmatrix}
\]

\section{Tables}
Below is a sample table showing data:

\begin{table}[h]
\centering
\begin{tabular}{@{}lrrr@{}}
\toprule
\textbf{Category} & \textbf{Value 1} & \textbf{Value 2} & \textbf{Total} \\
\midrule
Item A & 10.5 & 20.3 & 30.8 \\
Item B & 15.2 & 25.1 & 40.3 \\
Item C & 22.1 & 12.7 & 34.8 \\
\bottomrule
\end{tabular}
\caption{Sample Data Table}
\label{tab:sample}
\end{table}

\section{Lists}
\subsection{Bullet Points}
\begin{itemize}
  \item First item
  \item Second item
  \item Third item with sub-items:
    \begin{itemize}
      \item Sub-item 1
      \item Sub-item 2
    \end{itemize}
\end{itemize}

\subsection{Numbered List}
\begin{enumerate}
  \item First step
  \item Second step
  \item Third step
\end{enumerate}

\section{Citations and References}
When you use the citation search feature, references will be automatically managed. For example, you might cite a paper like this: \cite{example2024}. All citations will be automatically added to your \texttt{cite.bib} file.

\section{Conclusion}
This template demonstrates basic \LaTeX{} features including mathematical formulas, tables, and lists. The bibliography system is now set up to work seamlessly with Bibby AI's citation management. Feel free to modify it for your needs!

\bibliographystyle{plain}
\bibliography{cite}

\end{document}\documentclass{article}
\usepackage[utf8]{inputenc}
\usepackage{booktabs}
\usepackage{amsmath}
\usepackage{natbib}

\title{My Document}
\author{Your Name}
\date{\today}

\begin{document}

\maketitle

\section{Introduction}
This is a sample document demonstrating various \LaTeX{} features including mathematical equations, tables, and formatting. You can cite papers using the citation search feature, and they will be automatically added to your \texttt{cite.bib} file.

\section{Mathematical Expressions}
Here are some example equations:

\subsection{Quadratic Formula}
The quadratic formula is given by:
\[x = \frac{-b \pm \sqrt{b^2 - 4ac}}{2a}\]

\subsection{Matrix Example}
A sample matrix:
\[
\begin{pmatrix}
a_{11} & a_{12} & a_{13} \\
a_{21} & a_{22} & a_{23} \\
a_{31} & a_{32} & a_{33}
\end{pmatrix}
\]

\section{Tables}
Below is a sample table showing data:

\begin{table}[h]
\centering
\begin{tabular}{@{}lrrr@{}}
\toprule
\textbf{Category} & \textbf{Value 1} & \textbf{Value 2} & \textbf{Total} \\
\midrule
Item A & 10.5 & 20.3 & 30.8 \\
Item B & 15.2 & 25.1 & 40.3 \\
Item C & 22.1 & 12.7 & 34.8 \\
\bottomrule
\end{tabular}
\caption{Sample Data Table}
\label{tab:sample}
\end{table}

\section{Lists}
\subsection{Bullet Points}
\begin{itemize}
  \item First item
  \item Second item
  \item Third item with sub-items:
    \begin{itemize}
      \item Sub-item 1
      \item Sub-item 2
    \end{itemize}
\end{itemize}

\subsection{Numbered List}
\begin{enumerate}
  \item First step
  \item Second step
  \item Third step
\end{enumerate}

\section{Citations and References}
When you use the citation search feature, references will be automatically managed. For example, you might cite a paper like this: \cite{example2024}. All citations will be automatically added to your \texttt{cite.bib} file.

\section{Conclusion}
This template demonstrates basic \LaTeX{} features including mathematical formulas, tables, and lists. The bibliography system is now set up to work seamlessly with Bibby AI's citation management. Feel free to modify it for your needs!

\bibliographystyle{plain}
\bibliography{cite}

\end{document}\documentclass{article}
\usepackage[utf8]{inputenc}
\usepackage{booktabs}
\usepackage{amsmath}
\usepackage{natbib}

\title{My Document}
\author{Your Name}
\date{\today}

\begin{document}

\maketitle

\section{Introduction}
This is a sample document demonstrating various \LaTeX{} features including mathematical equations, tables, and formatting. You can cite papers using the citation search feature, and they will be automatically added to your \texttt{cite.bib} file.

\section{Mathematical Expressions}
Here are some example equations:

\subsection{Quadratic Formula}
The quadratic formula is given by:
\[x = \frac{-b \pm \sqrt{b^2 - 4ac}}{2a}\]

\subsection{Matrix Example}
A sample matrix:
\[
\begin{pmatrix}
a_{11} & a_{12} & a_{13} \\
a_{21} & a_{22} & a_{23} \\
a_{31} & a_{32} & a_{33}
\end{pmatrix}
\]

\section{Tables}
Below is a sample table showing data:

\begin{table}[h]
\centering
\begin{tabular}{@{}lrrr@{}}
\toprule
\textbf{Category} & \textbf{Value 1} & \textbf{Value 2} & \textbf{Total} \\
\midrule
Item A & 10.5 & 20.3 & 30.8 \\
Item B & 15.2 & 25.1 & 40.3 \\
Item C & 22.1 & 12.7 & 34.8 \\
\bottomrule
\end{tabular}
\caption{Sample Data Table}
\label{tab:sample}
\end{table}

\section{Lists}
\subsection{Bullet Points}
\begin{itemize}
  \item First item
  \item Second item
  \item Third item with sub-items:
    \begin{itemize}
      \item Sub-item 1
      \item Sub-item 2
    \end{itemize}
\end{itemize}

\subsection{Numbered List}
\begin{enumerate}
  \item First step
  \item Second step
  \item Third step
\end{enumerate}

\section{Citations and References}
When you use the citation search feature, references will be automatically managed. For example, you might cite a paper like this: \cite{example2024}. All citations will be automatically added to your \texttt{cite.bib} file.

\section{Conclusion}
This template demonstrates basic \LaTeX{} features including mathematical formulas, tables, and lists. The bibliography system is now set up to work seamlessly with Bibby AI's citation management. Feel free to modify it for your needs!

\bibliographystyle{plain}
\bibliography{cite}

\end{document}\documentclass{article}
\usepackage[utf8]{inputenc}
\usepackage{booktabs}
\usepackage{amsmath}
\usepackage{natbib}

\title{My Document}
\author{Your Name}
\date{\today}

\begin{document}

\maketitle

\section{Introduction}
This is a sample document demonstrating various \LaTeX{} features including mathematical equations, tables, and formatting. You can cite papers using the citation search feature, and they will be automatically added to your \texttt{cite.bib} file.

\section{Mathematical Expressions}
Here are some example equations:

\subsection{Quadratic Formula}
The quadratic formula is given by:
\[x = \frac{-b \pm \sqrt{b^2 - 4ac}}{2a}\]

\subsection{Matrix Example}
A sample matrix:
\[
\begin{pmatrix}
a_{11} & a_{12} & a_{13} \\
a_{21} & a_{22} & a_{23} \\
a_{31} & a_{32} & a_{33}
\end{pmatrix}
\]

\section{Tables}
Below is a sample table showing data:

\begin{table}[h]
\centering
\begin{tabular}{@{}lrrr@{}}
\toprule
\textbf{Category} & \textbf{Value 1} & \textbf{Value 2} & \textbf{Total} \\
\midrule
Item A & 10.5 & 20.3 & 30.8 \\
Item B & 15.2 & 25.1 & 40.3 \\
Item C & 22.1 & 12.7 & 34.8 \\
\bottomrule
\end{tabular}
\caption{Sample Data Table}
\label{tab:sample}
\end{table}

\section{Lists}
\subsection{Bullet Points}
\begin{itemize}
  \item First item
  \item Second item
  \item Third item with sub-items:
    \begin{itemize}
      \item Sub-item 1
      \item Sub-item 2
    \end{itemize}
\end{itemize}

\subsection{Numbered List}
\begin{enumerate}
  \item First step
  \item Second step
  \item Third step
\end{enumerate}

\section{Citations and References}
When you use the citation search feature, references will be automatically managed. For example, you might cite a paper like this: \cite{example2024}. All citations will be automatically added to your \texttt{cite.bib} file.

\section{Conclusion}
This template demonstrates basic \LaTeX{} features including mathematical formulas, tables, and lists. The bibliography system is now set up to work seamlessly with Bibby AI's citation management. Feel free to modify it for your needs!

\bibliographystyle{plain}
\bibliography{cite}

\end{document}\documentclass{article}
\usepackage[utf8]{inputenc}
\usepackage{booktabs}
\usepackage{amsmath}
\usepackage{natbib}

\title{My Document}
\author{Your Name}
\date{\today}

\begin{document}

\maketitle

\section{Introduction}
This is a sample document demonstrating various \LaTeX{} features including mathematical equations, tables, and formatting. You can cite papers using the citation search feature, and they will be automatically added to your \texttt{cite.bib} file.

\section{Mathematical Expressions}
Here are some example equations:

\subsection{Quadratic Formula}
The quadratic formula is given by:
\[x = \frac{-b \pm \sqrt{b^2 - 4ac}}{2a}\]

\subsection{Matrix Example}
A sample matrix:
\[
\begin{pmatrix}
a_{11} & a_{12} & a_{13} \\
a_{21} & a_{22} & a_{23} \\
a_{31} & a_{32} & a_{33}
\end{pmatrix}
\]

\section{Tables}
Below is a sample table showing data:

\begin{table}[h]
\centering
\begin{tabular}{@{}lrrr@{}}
\toprule
\textbf{Category} & \textbf{Value 1} & \textbf{Value 2} & \textbf{Total} \\
\midrule
Item A & 10.5 & 20.3 & 30.8 \\
Item B & 15.2 & 25.1 & 40.3 \\
Item C & 22.1 & 12.7 & 34.8 \\
\bottomrule
\end{tabular}
\caption{Sample Data Table}
\label{tab:sample}
\end{table}

\section{Lists}
\subsection{Bullet Points}
\begin{itemize}
  \item First item
  \item Second item
  \item Third item with sub-items:
    \begin{itemize}
      \item Sub-item 1
      \item Sub-item 2
    \end{itemize}
\end{itemize}

\subsection{Numbered List}
\begin{enumerate}
  \item First step
  \item Second step
  \item Third step
\end{enumerate}

\section{Citations and References}
When you use the citation search feature, references will be automatically managed. For example, you might cite a paper like this: \cite{example2024}. All citations will be automatically added to your \texttt{cite.bib} file.

\section{Conclusion}
This template demonstrates basic \LaTeX{} features including mathematical formulas, tables, and lists. The bibliography system is now set up to work seamlessly with Bibby AI's citation management. Feel free to modify it for your needs!

\bibliographystyle{plain}
\bibliography{cite}

\end{document}\documentclass{article}
\usepackage[utf8]{inputenc}
\usepackage{booktabs}
\usepackage{amsmath}
\usepackage{natbib}

\title{My Document}
\author{Your Name}
\date{\today}

\begin{document}

\maketitle

\section{Introduction}
This is a sample document demonstrating various \LaTeX{} features including mathematical equations, tables, and formatting. You can cite papers using the citation search feature, and they will be automatically added to your \texttt{cite.bib} file.

\section{Mathematical Expressions}
Here are some example equations:

\subsection{Quadratic Formula}
The quadratic formula is given by:
\[x = \frac{-b \pm \sqrt{b^2 - 4ac}}{2a}\]

\subsection{Matrix Example}
A sample matrix:
\[
\begin{pmatrix}
a_{11} & a_{12} & a_{13} \\
a_{21} & a_{22} & a_{23} \\
a_{31} & a_{32} & a_{33}
\end{pmatrix}
\]

\section{Tables}
Below is a sample table showing data:

\begin{table}[h]
\centering
\begin{tabular}{@{}lrrr@{}}
\toprule
\textbf{Category} & \textbf{Value 1} & \textbf{Value 2} & \textbf{Total} \\
\midrule
Item A & 10.5 & 20.3 & 30.8 \\
Item B & 15.2 & 25.1 & 40.3 \\
Item C & 22.1 & 12.7 & 34.8 \\
\bottomrule
\end{tabular}
\caption{Sample Data Table}
\label{tab:sample}
\end{table}

\section{Lists}
\subsection{Bullet Points}
\begin{itemize}
  \item First item
  \item Second item
  \item Third item with sub-items:
    \begin{itemize}
      \item Sub-item 1
      \item Sub-item 2
    \end{itemize}
\end{itemize}

\subsection{Numbered List}
\begin{enumerate}
  \item First step
  \item Second step
  \item Third step
\end{enumerate}

\section{Citations and References}
When you use the citation search feature, references will be automatically managed. For example, you might cite a paper like this: \cite{example2024}. All citations will be automatically added to your \texttt{cite.bib} file.

\section{Conclusion}
This template demonstrates basic \LaTeX{} features including mathematical formulas, tables, and lists. The bibliography system is now set up to work seamlessly with Bibby AI's citation management. Feel free to modify it for your needs!

\bibliographystyle{plain}
\bibliography{cite}

\end{document}\documentclass{article}
\usepackage[utf8]{inputenc}
\usepackage{booktabs}
\usepackage{amsmath}
\usepackage{natbib}

\title{My Document}
\author{Your Name}
\date{\today}

\begin{document}

\maketitle

\section{Introduction}
This is a sample document demonstrating various \LaTeX{} features including mathematical equations, tables, and formatting. You can cite papers using the citation search feature, and they will be automatically added to your \texttt{cite.bib} file.

\section{Mathematical Expressions}
Here are some example equations:

\subsection{Quadratic Formula}
The quadratic formula is given by:
\[x = \frac{-b \pm \sqrt{b^2 - 4ac}}{2a}\]

\subsection{Matrix Example}
A sample matrix:
\[
\begin{pmatrix}
a_{11} & a_{12} & a_{13} \\
a_{21} & a_{22} & a_{23} \\
a_{31} & a_{32} & a_{33}
\end{pmatrix}
\]

\section{Tables}
Below is a sample table showing data:

\begin{table}[h]
\centering
\begin{tabular}{@{}lrrr@{}}
\toprule
\textbf{Category} & \textbf{Value 1} & \textbf{Value 2} & \textbf{Total} \\
\midrule
Item A & 10.5 & 20.3 & 30.8 \\
Item B & 15.2 & 25.1 & 40.3 \\
Item C & 22.1 & 12.7 & 34.8 \\
\bottomrule
\end{tabular}
\caption{Sample Data Table}
\label{tab:sample}
\end{table}

\section{Lists}
\subsection{Bullet Points}
\begin{itemize}
  \item First item
  \item Second item
  \item Third item with sub-items:
    \begin{itemize}
      \item Sub-item 1
      \item Sub-item 2
    \end{itemize}
\end{itemize}

\subsection{Numbered List}
\begin{enumerate}
  \item First step
  \item Second step
  \item Third step
\end{enumerate}

\section{Citations and References}
When you use the citation search feature, references will be automatically managed. For example, you might cite a paper like this: \cite{example2024}. All citations will be automatically added to your \texttt{cite.bib} file.

\section{Conclusion}
This template demonstrates basic \LaTeX{} features including mathematical formulas, tables, and lists. The bibliography system is now set up to work seamlessly with Bibby AI's citation management. Feel free to modify it for your needs!

\bibliographystyle{plain}
\bibliography{cite}

\end{document}\documentclass{article}
\usepackage[utf8]{inputenc}
\usepackage{booktabs}
\usepackage{amsmath}
\usepackage{natbib}

\title{My Document}
\author{Your Name}
\date{\today}

\begin{document}

\maketitle

\section{Introduction}
This is a sample document demonstrating various \LaTeX{} features including mathematical equations, tables, and formatting. You can cite papers using the citation search feature, and they will be automatically added to your \texttt{cite.bib} file.

\section{Mathematical Expressions}
Here are some example equations:

\subsection{Quadratic Formula}
The quadratic formula is given by:
\[x = \frac{-b \pm \sqrt{b^2 - 4ac}}{2a}\]

\subsection{Matrix Example}
A sample matrix:
\[
\begin{pmatrix}
a_{11} & a_{12} & a_{13} \\
a_{21} & a_{22} & a_{23} \\
a_{31} & a_{32} & a_{33}
\end{pmatrix}
\]

\section{Tables}
Below is a sample table showing data:

\begin{table}[h]
\centering
\begin{tabular}{@{}lrrr@{}}
\toprule
\textbf{Category} & \textbf{Value 1} & \textbf{Value 2} & \textbf{Total} \\
\midrule
Item A & 10.5 & 20.3 & 30.8 \\
Item B & 15.2 & 25.1 & 40.3 \\
Item C & 22.1 & 12.7 & 34.8 \\
\bottomrule
\end{tabular}
\caption{Sample Data Table}
\label{tab:sample}
\end{table}

\section{Lists}
\subsection{Bullet Points}
\begin{itemize}
  \item First item
  \item Second item
  \item Third item with sub-items:
    \begin{itemize}
      \item Sub-item 1
      \item Sub-item 2
    \end{itemize}
\end{itemize}

\subsection{Numbered List}
\begin{enumerate}
  \item First step
  \item Second step
  \item Third step
\end{enumerate}

\section{Citations and References}
When you use the citation search feature, references will be automatically managed. For example, you might cite a paper like this: \cite{example2024}. All citations will be automatically added to your \texttt{cite.bib} file.

\section{Conclusion}
This template demonstrates basic \LaTeX{} features including mathematical formulas, tables, and lists. The bibliography system is now set up to work seamlessly with Bibby AI's citation management. Feel free to modify it for your needs!

\bibliographystyle{plain}
\bibliography{cite}

\end{document}\documentclass{article}
\usepackage[utf8]{inputenc}
\usepackage{booktabs}
\usepackage{amsmath}
\usepackage{natbib}

\title{My Document}
\author{Your Name}
\date{\today}

\begin{document}

\maketitle

\section{Introduction}
This is a sample document demonstrating various \LaTeX{} features including mathematical equations, tables, and formatting. You can cite papers using the citation search feature, and they will be automatically added to your \texttt{cite.bib} file.

\section{Mathematical Expressions}
Here are some example equations:

\subsection{Quadratic Formula}
The quadratic formula is given by:
\[x = \frac{-b \pm \sqrt{b^2 - 4ac}}{2a}\]

\subsection{Matrix Example}
A sample matrix:
\[
\begin{pmatrix}
a_{11} & a_{12} & a_{13} \\
a_{21} & a_{22} & a_{23} \\
a_{31} & a_{32} & a_{33}
\end{pmatrix}
\]

\section{Tables}
Below is a sample table showing data:

\begin{table}[h]
\centering
\begin{tabular}{@{}lrrr@{}}
\toprule
\textbf{Category} & \textbf{Value 1} & \textbf{Value 2} & \textbf{Total} \\
\midrule
Item A & 10.5 & 20.3 & 30.8 \\
Item B & 15.2 & 25.1 & 40.3 \\
Item C & 22.1 & 12.7 & 34.8 \\
\bottomrule
\end{tabular}
\caption{Sample Data Table}
\label{tab:sample}
\end{table}

\section{Lists}
\subsection{Bullet Points}
\begin{itemize}
  \item First item
  \item Second item
  \item Third item with sub-items:
    \begin{itemize}
      \item Sub-item 1
      \item Sub-item 2
    \end{itemize}
\end{itemize}

\subsection{Numbered List}
\begin{enumerate}
  \item First step
  \item Second step
  \item Third step
\end{enumerate}

\section{Citations and References}
When you use the citation search feature, references will be automatically managed. For example, you might cite a paper like this: \cite{example2024}. All citations will be automatically added to your \texttt{cite.bib} file.

\section{Conclusion}
This template demonstrates basic \LaTeX{} features including mathematical formulas, tables, and lists. The bibliography system is now set up to work seamlessly with Bibby AI's citation management. Feel free to modify it for your needs!

\bibliographystyle{plain}
\bibliography{cite}

\end{document}\documentclass{article}
\usepackage[utf8]{inputenc}
\usepackage{booktabs}
\usepackage{amsmath}
\usepackage{natbib}

\title{My Document}
\author{Your Name}
\date{\today}

\begin{document}

\maketitle

\section{Introduction}
This is a sample document demonstrating various \LaTeX{} features including mathematical equations, tables, and formatting. You can cite papers using the citation search feature, and they will be automatically added to your \texttt{cite.bib} file.

\section{Mathematical Expressions}
Here are some example equations:

\subsection{Quadratic Formula}
The quadratic formula is given by:
\[x = \frac{-b \pm \sqrt{b^2 - 4ac}}{2a}\]

\subsection{Matrix Example}
A sample matrix:
\[
\begin{pmatrix}
a_{11} & a_{12} & a_{13} \\
a_{21} & a_{22} & a_{23} \\
a_{31} & a_{32} & a_{33}
\end{pmatrix}
\]

\section{Tables}
Below is a sample table showing data:

\begin{table}[h]
\centering
\begin{tabular}{@{}lrrr@{}}
\toprule
\textbf{Category} & \textbf{Value 1} & \textbf{Value 2} & \textbf{Total} \\
\midrule
Item A & 10.5 & 20.3 & 30.8 \\
Item B & 15.2 & 25.1 & 40.3 \\
Item C & 22.1 & 12.7 & 34.8 \\
\bottomrule
\end{tabular}
\caption{Sample Data Table}
\label{tab:sample}
\end{table}

\section{Lists}
\subsection{Bullet Points}
\begin{itemize}
  \item First item
  \item Second item
  \item Third item with sub-items:
    \begin{itemize}
      \item Sub-item 1
      \item Sub-item 2
    \end{itemize}
\end{itemize}

\subsection{Numbered List}
\begin{enumerate}
  \item First step
  \item Second step
  \item Third step
\end{enumerate}

\section{Citations and References}
When you use the citation search feature, references will be automatically managed. For example, you might cite a paper like this: \cite{example2024}. All citations will be automatically added to your \texttt{cite.bib} file.

\section{Conclusion}
This template demonstrates basic \LaTeX{} features including mathematical formulas, tables, and lists. The bibliography system is now set up to work seamlessly with Bibby AI's citation management. Feel free to modify it for your needs!

\bibliographystyle{plain}
\bibliography{cite}

\end{document}\documentclass{article}
\usepackage[utf8]{inputenc}
\usepackage{booktabs}
\usepackage{amsmath}
\usepackage{natbib}

\title{My Document}
\author{Your Name}
\date{\today}

\begin{document}

\maketitle

\section{Introduction}
This is a sample document demonstrating various \LaTeX{} features including mathematical equations, tables, and formatting. You can cite papers using the citation search feature, and they will be automatically added to your \texttt{cite.bib} file.

\section{Mathematical Expressions}
Here are some example equations:

\subsection{Quadratic Formula}
The quadratic formula is given by:
\[x = \frac{-b \pm \sqrt{b^2 - 4ac}}{2a}\]

\subsection{Matrix Example}
A sample matrix:
\[
\begin{pmatrix}
a_{11} & a_{12} & a_{13} \\
a_{21} & a_{22} & a_{23} \\
a_{31} & a_{32} & a_{33}
\end{pmatrix}
\]

\section{Tables}
Below is a sample table showing data:

\begin{table}[h]
\centering
\begin{tabular}{@{}lrrr@{}}
\toprule
\textbf{Category} & \textbf{Value 1} & \textbf{Value 2} & \textbf{Total} \\
\midrule
Item A & 10.5 & 20.3 & 30.8 \\
Item B & 15.2 & 25.1 & 40.3 \\
Item C & 22.1 & 12.7 & 34.8 \\
\bottomrule
\end{tabular}
\caption{Sample Data Table}
\label{tab:sample}
\end{table}

\section{Lists}
\subsection{Bullet Points}
\begin{itemize}
  \item First item
  \item Second item
  \item Third item with sub-items:
    \begin{itemize}
      \item Sub-item 1
      \item Sub-item 2
    \end{itemize}
\end{itemize}

\subsection{Numbered List}
\begin{enumerate}
  \item First step
  \item Second step
  \item Third step
\end{enumerate}

\section{Citations and References}
When you use the citation search feature, references will be automatically managed. For example, you might cite a paper like this: \cite{example2024}. All citations will be automatically added to your \texttt{cite.bib} file.

\section{Conclusion}
This template demonstrates basic \LaTeX{} features including mathematical formulas, tables, and lists. The bibliography system is now set up to work seamlessly with Bibby AI's citation management. Feel free to modify it for your needs!

\bibliographystyle{plain}
\bibliography{cite}

\end{document}\documentclass{article}
\usepackage[utf8]{inputenc}
\usepackage{booktabs}
\usepackage{amsmath}
\usepackage{natbib}

\title{My Document}
\author{Your Name}
\date{\today}

\begin{document}

\maketitle

\section{Introduction}
This is a sample document demonstrating various \LaTeX{} features including mathematical equations, tables, and formatting. You can cite papers using the citation search feature, and they will be automatically added to your \texttt{cite.bib} file.

\section{Mathematical Expressions}
Here are some example equations:

\subsection{Quadratic Formula}
The quadratic formula is given by:
\[x = \frac{-b \pm \sqrt{b^2 - 4ac}}{2a}\]

\subsection{Matrix Example}
A sample matrix:
\[
\begin{pmatrix}
a_{11} & a_{12} & a_{13} \\
a_{21} & a_{22} & a_{23} \\
a_{31} & a_{32} & a_{33}
\end{pmatrix}
\]

\section{Tables}
Below is a sample table showing data:

\begin{table}[h]
\centering
\begin{tabular}{@{}lrrr@{}}
\toprule
\textbf{Category} & \textbf{Value 1} & \textbf{Value 2} & \textbf{Total} \\
\midrule
Item A & 10.5 & 20.3 & 30.8 \\
Item B & 15.2 & 25.1 & 40.3 \\
Item C & 22.1 & 12.7 & 34.8 \\
\bottomrule
\end{tabular}
\caption{Sample Data Table}
\label{tab:sample}
\end{table}

\section{Lists}
\subsection{Bullet Points}
\begin{itemize}
  \item First item
  \item Second item
  \item Third item with sub-items:
    \begin{itemize}
      \item Sub-item 1
      \item Sub-item 2
    \end{itemize}
\end{itemize}

\subsection{Numbered List}
\begin{enumerate}
  \item First step
  \item Second step
  \item Third step
\end{enumerate}

\section{Citations and References}
When you use the citation search feature, references will be automatically managed. For example, you might cite a paper like this: \cite{example2024}. All citations will be automatically added to your \texttt{cite.bib} file.

\section{Conclusion}
This template demonstrates basic \LaTeX{} features including mathematical formulas, tables, and lists. The bibliography system is now set up to work seamlessly with Bibby AI's citation management. Feel free to modify it for your needs!

\bibliographystyle{plain}
\bibliography{cite}

\end{document}\documentclass{article}
\usepackage[utf8]{inputenc}
\usepackage{booktabs}
\usepackage{amsmath}
\usepackage{natbib}

\title{My Document}
\author{Your Name}
\date{\today}

\begin{document}

\maketitle

\section{Introduction}
This is a sample document demonstrating various \LaTeX{} features including mathematical equations, tables, and formatting. You can cite papers using the citation search feature, and they will be automatically added to your \texttt{cite.bib} file.

\section{Mathematical Expressions}
Here are some example equations:

\subsection{Quadratic Formula}
The quadratic formula is given by:
\[x = \frac{-b \pm \sqrt{b^2 - 4ac}}{2a}\]

\subsection{Matrix Example}
A sample matrix:
\[
\begin{pmatrix}
a_{11} & a_{12} & a_{13} \\
a_{21} & a_{22} & a_{23} \\
a_{31} & a_{32} & a_{33}
\end{pmatrix}
\]

\section{Tables}
Below is a sample table showing data:

\begin{table}[h]
\centering
\begin{tabular}{@{}lrrr@{}}
\toprule
\textbf{Category} & \textbf{Value 1} & \textbf{Value 2} & \textbf{Total} \\
\midrule
Item A & 10.5 & 20.3 & 30.8 \\
Item B & 15.2 & 25.1 & 40.3 \\
Item C & 22.1 & 12.7 & 34.8 \\
\bottomrule
\end{tabular}
\caption{Sample Data Table}
\label{tab:sample}
\end{table}

\section{Lists}
\subsection{Bullet Points}
\begin{itemize}
  \item First item
  \item Second item
  \item Third item with sub-items:
    \begin{itemize}
      \item Sub-item 1
      \item Sub-item 2
    \end{itemize}
\end{itemize}

\subsection{Numbered List}
\begin{enumerate}
  \item First step
  \item Second step
  \item Third step
\end{enumerate}

\section{Citations and References}
When you use the citation search feature, references will be automatically managed. For example, you might cite a paper like this: \cite{example2024}. All citations will be automatically added to your \texttt{cite.bib} file.

\section{Conclusion}
This template demonstrates basic \LaTeX{} features including mathematical formulas, tables, and lists. The bibliography system is now set up to work seamlessly with Bibby AI's citation management. Feel free to modify it for your needs!

\bibliographystyle{plain}
\bibliography{cite}

\end{document}\documentclass{article}
\usepackage[utf8]{inputenc}
\usepackage{booktabs}
\usepackage{amsmath}
\usepackage{natbib}

\title{My Document}
\author{Your Name}
\date{\today}

\begin{document}

\maketitle

\section{Introduction}
This is a sample document demonstrating various \LaTeX{} features including mathematical equations, tables, and formatting. You can cite papers using the citation search feature, and they will be automatically added to your \texttt{cite.bib} file.

\section{Mathematical Expressions}
Here are some example equations:

\subsection{Quadratic Formula}
The quadratic formula is given by:
\[x = \frac{-b \pm \sqrt{b^2 - 4ac}}{2a}\]

\subsection{Matrix Example}
A sample matrix:
\[
\begin{pmatrix}
a_{11} & a_{12} & a_{13} \\
a_{21} & a_{22} & a_{23} \\
a_{31} & a_{32} & a_{33}
\end{pmatrix}
\]

\section{Tables}
Below is a sample table showing data:

\begin{table}[h]
\centering
\begin{tabular}{@{}lrrr@{}}
\toprule
\textbf{Category} & \textbf{Value 1} & \textbf{Value 2} & \textbf{Total} \\
\midrule
Item A & 10.5 & 20.3 & 30.8 \\
Item B & 15.2 & 25.1 & 40.3 \\
Item C & 22.1 & 12.7 & 34.8 \\
\bottomrule
\end{tabular}
\caption{Sample Data Table}
\label{tab:sample}
\end{table}

\section{Lists}
\subsection{Bullet Points}
\begin{itemize}
  \item First item
  \item Second item
  \item Third item with sub-items:
    \begin{itemize}
      \item Sub-item 1
      \item Sub-item 2
    \end{itemize}
\end{itemize}

\subsection{Numbered List}
\begin{enumerate}
  \item First step
  \item Second step
  \item Third step
\end{enumerate}

\section{Citations and References}
When you use the citation search feature, references will be automatically managed. For example, you might cite a paper like this: \cite{example2024}. All citations will be automatically added to your \texttt{cite.bib} file.

\section{Conclusion}
This template demonstrates basic \LaTeX{} features including mathematical formulas, tables, and lists. The bibliography system is now set up to work seamlessly with Bibby AI's citation management. Feel free to modify it for your needs!

\bibliographystyle{plain}
\bibliography{cite}

\end{document}\documentclass{article}
\usepackage[utf8]{inputenc}
\usepackage{booktabs}
\usepackage{amsmath}
\usepackage{natbib}

\title{My Document}
\author{Your Name}
\date{\today}

\begin{document}

\maketitle

\section{Introduction}
This is a sample document demonstrating various \LaTeX{} features including mathematical equations, tables, and formatting. You can cite papers using the citation search feature, and they will be automatically added to your \texttt{cite.bib} file.

\section{Mathematical Expressions}
Here are some example equations:

\subsection{Quadratic Formula}
The quadratic formula is given by:
\[x = \frac{-b \pm \sqrt{b^2 - 4ac}}{2a}\]

\subsection{Matrix Example}
A sample matrix:
\[
\begin{pmatrix}
a_{11} & a_{12} & a_{13} \\
a_{21} & a_{22} & a_{23} \\
a_{31} & a_{32} & a_{33}
\end{pmatrix}
\]

\section{Tables}
Below is a sample table showing data:

\begin{table}[h]
\centering
\begin{tabular}{@{}lrrr@{}}
\toprule
\textbf{Category} & \textbf{Value 1} & \textbf{Value 2} & \textbf{Total} \\
\midrule
Item A & 10.5 & 20.3 & 30.8 \\
Item B & 15.2 & 25.1 & 40.3 \\
Item C & 22.1 & 12.7 & 34.8 \\
\bottomrule
\end{tabular}
\caption{Sample Data Table}
\label{tab:sample}
\end{table}

\section{Lists}
\subsection{Bullet Points}
\begin{itemize}
  \item First item
  \item Second item
  \item Third item with sub-items:
    \begin{itemize}
      \item Sub-item 1
      \item Sub-item 2
    \end{itemize}
\end{itemize}

\subsection{Numbered List}
\begin{enumerate}
  \item First step
  \item Second step
  \item Third step
\end{enumerate}

\section{Citations and References}
When you use the citation search feature, references will be automatically managed. For example, you might cite a paper like this: \cite{example2024}. All citations will be automatically added to your \texttt{cite.bib} file.

\section{Conclusion}
This template demonstrates basic \LaTeX{} features including mathematical formulas, tables, and lists. The bibliography system is now set up to work seamlessly with Bibby AI's citation management. Feel free to modify it for your needs!

\bibliographystyle{plain}
\bibliography{cite}

\end{document}\documentclass{article}
\usepackage[utf8]{inputenc}
\usepackage{booktabs}
\usepackage{amsmath}
\usepackage{natbib}

\title{My Document}
\author{Your Name}
\date{\today}

\begin{document}

\maketitle

\section{Introduction}
This is a sample document demonstrating various \LaTeX{} features including mathematical equations, tables, and formatting. You can cite papers using the citation search feature, and they will be automatically added to your \texttt{cite.bib} file.

\section{Mathematical Expressions}
Here are some example equations:

\subsection{Quadratic Formula}
The quadratic formula is given by:
\[x = \frac{-b \pm \sqrt{b^2 - 4ac}}{2a}\]

\subsection{Matrix Example}
A sample matrix:
\[
\begin{pmatrix}
a_{11} & a_{12} & a_{13} \\
a_{21} & a_{22} & a_{23} \\
a_{31} & a_{32} & a_{33}
\end{pmatrix}
\]

\section{Tables}
Below is a sample table showing data:

\begin{table}[h]
\centering
\begin{tabular}{@{}lrrr@{}}
\toprule
\textbf{Category} & \textbf{Value 1} & \textbf{Value 2} & \textbf{Total} \\
\midrule
Item A & 10.5 & 20.3 & 30.8 \\
Item B & 15.2 & 25.1 & 40.3 \\
Item C & 22.1 & 12.7 & 34.8 \\
\bottomrule
\end{tabular}
\caption{Sample Data Table}
\label{tab:sample}
\end{table}

\section{Lists}
\subsection{Bullet Points}
\begin{itemize}
  \item First item
  \item Second item
  \item Third item with sub-items:
    \begin{itemize}
      \item Sub-item 1
      \item Sub-item 2
    \end{itemize}
\end{itemize}

\subsection{Numbered List}
\begin{enumerate}
  \item First step
  \item Second step
  \item Third step
\end{enumerate}

\section{Citations and References}
When you use the citation search feature, references will be automatically managed. For example, you might cite a paper like this: \cite{example2024}. All citations will be automatically added to your \texttt{cite.bib} file.

\section{Conclusion}
This template demonstrates basic \LaTeX{} features including mathematical formulas, tables, and lists. The bibliography system is now set up to work seamlessly with Bibby AI's citation management. Feel free to modify it for your needs!

\bibliographystyle{plain}
\bibliography{cite}

\end{document}\documentclass{article}
\usepackage[utf8]{inputenc}
\usepackage{booktabs}
\usepackage{amsmath}
\usepackage{natbib}

\title{My Document}
\author{Your Name}
\date{\today}

\begin{document}

\maketitle

\section{Introduction}
This is a sample document demonstrating various \LaTeX{} features including mathematical equations, tables, and formatting. You can cite papers using the citation search feature, and they will be automatically added to your \texttt{cite.bib} file.

\section{Mathematical Expressions}
Here are some example equations:

\subsection{Quadratic Formula}
The quadratic formula is given by:
\[x = \frac{-b \pm \sqrt{b^2 - 4ac}}{2a}\]

\subsection{Matrix Example}
A sample matrix:
\[
\begin{pmatrix}
a_{11} & a_{12} & a_{13} \\
a_{21} & a_{22} & a_{23} \\
a_{31} & a_{32} & a_{33}
\end{pmatrix}
\]

\section{Tables}
Below is a sample table showing data:

\begin{table}[h]
\centering
\begin{tabular}{@{}lrrr@{}}
\toprule
\textbf{Category} & \textbf{Value 1} & \textbf{Value 2} & \textbf{Total} \\
\midrule
Item A & 10.5 & 20.3 & 30.8 \\
Item B & 15.2 & 25.1 & 40.3 \\
Item C & 22.1 & 12.7 & 34.8 \\
\bottomrule
\end{tabular}
\caption{Sample Data Table}
\label{tab:sample}
\end{table}

\section{Lists}
\subsection{Bullet Points}
\begin{itemize}
  \item First item
  \item Second item
  \item Third item with sub-items:
    \begin{itemize}
      \item Sub-item 1
      \item Sub-item 2
    \end{itemize}
\end{itemize}

\subsection{Numbered List}
\begin{enumerate}
  \item First step
  \item Second step
  \item Third step
\end{enumerate}

\section{Citations and References}
When you use the citation search feature, references will be automatically managed. For example, you might cite a paper like this: \cite{example2024}. All citations will be automatically added to your \texttt{cite.bib} file.

\section{Conclusion}
This template demonstrates basic \LaTeX{} features including mathematical formulas, tables, and lists. The bibliography system is now set up to work seamlessly with Bibby AI's citation management. Feel free to modify it for your needs!

\bibliographystyle{plain}
\bibliography{cite}

\end{document}\documentclass{article}
\usepackage[utf8]{inputenc}
\usepackage{booktabs}
\usepackage{amsmath}
\usepackage{natbib}

\title{My Document}
\author{Your Name}
\date{\today}

\begin{document}

\maketitle

\section{Introduction}
This is a sample document demonstrating various \LaTeX{} features including mathematical equations, tables, and formatting. You can cite papers using the citation search feature, and they will be automatically added to your \texttt{cite.bib} file.

\section{Mathematical Expressions}
Here are some example equations:

\subsection{Quadratic Formula}
The quadratic formula is given by:
\[x = \frac{-b \pm \sqrt{b^2 - 4ac}}{2a}\]

\subsection{Matrix Example}
A sample matrix:
\[
\begin{pmatrix}
a_{11} & a_{12} & a_{13} \\
a_{21} & a_{22} & a_{23} \\
a_{31} & a_{32} & a_{33}
\end{pmatrix}
\]

\section{Tables}
Below is a sample table showing data:

\begin{table}[h]
\centering
\begin{tabular}{@{}lrrr@{}}
\toprule
\textbf{Category} & \textbf{Value 1} & \textbf{Value 2} & \textbf{Total} \\
\midrule
Item A & 10.5 & 20.3 & 30.8 \\
Item B & 15.2 & 25.1 & 40.3 \\
Item C & 22.1 & 12.7 & 34.8 \\
\bottomrule
\end{tabular}
\caption{Sample Data Table}
\label{tab:sample}
\end{table}

\section{Lists}
\subsection{Bullet Points}
\begin{itemize}
  \item First item
  \item Second item
  \item Third item with sub-items:
    \begin{itemize}
      \item Sub-item 1
      \item Sub-item 2
    \end{itemize}
\end{itemize}

\subsection{Numbered List}
\begin{enumerate}
  \item First step
  \item Second step
  \item Third step
\end{enumerate}

\section{Citations and References}
When you use the citation search feature, references will be automatically managed. For example, you might cite a paper like this: \cite{example2024}. All citations will be automatically added to your \texttt{cite.bib} file.

\section{Conclusion}
This template demonstrates basic \LaTeX{} features including mathematical formulas, tables, and lists. The bibliography system is now set up to work seamlessly with Bibby AI's citation management. Feel free to modify it for your needs!

\bibliographystyle{plain}
\bibliography{cite}

\end{document}\documentclass{article}
\usepackage[utf8]{inputenc}
\usepackage{booktabs}
\usepackage{amsmath}
\usepackage{natbib}

\title{My Document}
\author{Your Name}
\date{\today}

\begin{document}

\maketitle

\section{Introduction}
This is a sample document demonstrating various \LaTeX{} features including mathematical equations, tables, and formatting. You can cite papers using the citation search feature, and they will be automatically added to your \texttt{cite.bib} file.

\section{Mathematical Expressions}
Here are some example equations:

\subsection{Quadratic Formula}
The quadratic formula is given by:
\[x = \frac{-b \pm \sqrt{b^2 - 4ac}}{2a}\]

\subsection{Matrix Example}
A sample matrix:
\[
\begin{pmatrix}
a_{11} & a_{12} & a_{13} \\
a_{21} & a_{22} & a_{23} \\
a_{31} & a_{32} & a_{33}
\end{pmatrix}
\]

\section{Tables}
Below is a sample table showing data:

\begin{table}[h]
\centering
\begin{tabular}{@{}lrrr@{}}
\toprule
\textbf{Category} & \textbf{Value 1} & \textbf{Value 2} & \textbf{Total} \\
\midrule
Item A & 10.5 & 20.3 & 30.8 \\
Item B & 15.2 & 25.1 & 40.3 \\
Item C & 22.1 & 12.7 & 34.8 \\
\bottomrule
\end{tabular}
\caption{Sample Data Table}
\label{tab:sample}
\end{table}

\section{Lists}
\subsection{Bullet Points}
\begin{itemize}
  \item First item
  \item Second item
  \item Third item with sub-items:
    \begin{itemize}
      \item Sub-item 1
      \item Sub-item 2
    \end{itemize}
\end{itemize}

\subsection{Numbered List}
\begin{enumerate}
  \item First step
  \item Second step
  \item Third step
\end{enumerate}

\section{Citations and References}
When you use the citation search feature, references will be automatically managed. For example, you might cite a paper like this: \cite{example2024}. All citations will be automatically added to your \texttt{cite.bib} file.

\section{Conclusion}
This template demonstrates basic \LaTeX{} features including mathematical formulas, tables, and lists. The bibliography system is now set up to work seamlessly with Bibby AI's citation management. Feel free to modify it for your needs!

\bibliographystyle{plain}
\bibliography{cite}

\end{document}\documentclass{article}
\usepackage[utf8]{inputenc}
\usepackage{booktabs}
\usepackage{amsmath}
\usepackage{natbib}

\title{My Document}
\author{Your Name}
\date{\today}

\begin{document}

\maketitle

\section{Introduction}
This is a sample document demonstrating various \LaTeX{} features including mathematical equations, tables, and formatting. You can cite papers using the citation search feature, and they will be automatically added to your \texttt{cite.bib} file.

\section{Mathematical Expressions}
Here are some example equations:

\subsection{Quadratic Formula}
The quadratic formula is given by:
\[x = \frac{-b \pm \sqrt{b^2 - 4ac}}{2a}\]

\subsection{Matrix Example}
A sample matrix:
\[
\begin{pmatrix}
a_{11} & a_{12} & a_{13} \\
a_{21} & a_{22} & a_{23} \\
a_{31} & a_{32} & a_{33}
\end{pmatrix}
\]

\section{Tables}
Below is a sample table showing data:

\begin{table}[h]
\centering
\begin{tabular}{@{}lrrr@{}}
\toprule
\textbf{Category} & \textbf{Value 1} & \textbf{Value 2} & \textbf{Total} \\
\midrule
Item A & 10.5 & 20.3 & 30.8 \\
Item B & 15.2 & 25.1 & 40.3 \\
Item C & 22.1 & 12.7 & 34.8 \\
\bottomrule
\end{tabular}
\caption{Sample Data Table}
\label{tab:sample}
\end{table}

\section{Lists}
\subsection{Bullet Points}
\begin{itemize}
  \item First item
  \item Second item
  \item Third item with sub-items:
    \begin{itemize}
      \item Sub-item 1
      \item Sub-item 2
    \end{itemize}
\end{itemize}

\subsection{Numbered List}
\begin{enumerate}
  \item First step
  \item Second step
  \item Third step
\end{enumerate}

\section{Citations and References}
When you use the citation search feature, references will be automatically managed. For example, you might cite a paper like this: \cite{example2024}. All citations will be automatically added to your \texttt{cite.bib} file.

\section{Conclusion}
This template demonstrates basic \LaTeX{} features including mathematical formulas, tables, and lists. The bibliography system is now set up to work seamlessly with Bibby AI's citation management. Feel free to modify it for your needs!

\bibliographystyle{plain}
\bibliography{cite}

\end{document}\documentclass{article}
\usepackage[utf8]{inputenc}
\usepackage{booktabs}
\usepackage{amsmath}
\usepackage{natbib}

\title{My Document}
\author{Your Name}
\date{\today}

\begin{document}

\maketitle

\section{Introduction}
This is a sample document demonstrating various \LaTeX{} features including mathematical equations, tables, and formatting. You can cite papers using the citation search feature, and they will be automatically added to your \texttt{cite.bib} file.

\section{Mathematical Expressions}
Here are some example equations:

\subsection{Quadratic Formula}
The quadratic formula is given by:
\[x = \frac{-b \pm \sqrt{b^2 - 4ac}}{2a}\]

\subsection{Matrix Example}
A sample matrix:
\[
\begin{pmatrix}
a_{11} & a_{12} & a_{13} \\
a_{21} & a_{22} & a_{23} \\
a_{31} & a_{32} & a_{33}
\end{pmatrix}
\]

\section{Tables}
Below is a sample table showing data:

\begin{table}[h]
\centering
\begin{tabular}{@{}lrrr@{}}
\toprule
\textbf{Category} & \textbf{Value 1} & \textbf{Value 2} & \textbf{Total} \\
\midrule
Item A & 10.5 & 20.3 & 30.8 \\
Item B & 15.2 & 25.1 & 40.3 \\
Item C & 22.1 & 12.7 & 34.8 \\
\bottomrule
\end{tabular}
\caption{Sample Data Table}
\label{tab:sample}
\end{table}

\section{Lists}
\subsection{Bullet Points}
\begin{itemize}
  \item First item
  \item Second item
  \item Third item with sub-items:
    \begin{itemize}
      \item Sub-item 1
      \item Sub-item 2
    \end{itemize}
\end{itemize}

\subsection{Numbered List}
\begin{enumerate}
  \item First step
  \item Second step
  \item Third step
\end{enumerate}

\section{Citations and References}
When you use the citation search feature, references will be automatically managed. For example, you might cite a paper like this: \cite{example2024}. All citations will be automatically added to your \texttt{cite.bib} file.

\section{Conclusion}
This template demonstrates basic \LaTeX{} features including mathematical formulas, tables, and lists. The bibliography system is now set up to work seamlessly with Bibby AI's citation management. Feel free to modify it for your needs!

\bibliographystyle{plain}
\bibliography{cite}

\end{document}\documentclass{article}
\usepackage[utf8]{inputenc}
\usepackage{booktabs}
\usepackage{amsmath}
\usepackage{natbib}

\title{My Document}
\author{Your Name}
\date{\today}

\begin{document}

\maketitle

\section{Introduction}
This is a sample document demonstrating various \LaTeX{} features including mathematical equations, tables, and formatting. You can cite papers using the citation search feature, and they will be automatically added to your \texttt{cite.bib} file.

\section{Mathematical Expressions}
Here are some example equations:

\subsection{Quadratic Formula}
The quadratic formula is given by:
\[x = \frac{-b \pm \sqrt{b^2 - 4ac}}{2a}\]

\subsection{Matrix Example}
A sample matrix:
\[
\begin{pmatrix}
a_{11} & a_{12} & a_{13} \\
a_{21} & a_{22} & a_{23} \\
a_{31} & a_{32} & a_{33}
\end{pmatrix}
\]

\section{Tables}
Below is a sample table showing data:

\begin{table}[h]
\centering
\begin{tabular}{@{}lrrr@{}}
\toprule
\textbf{Category} & \textbf{Value 1} & \textbf{Value 2} & \textbf{Total} \\
\midrule
Item A & 10.5 & 20.3 & 30.8 \\
Item B & 15.2 & 25.1 & 40.3 \\
Item C & 22.1 & 12.7 & 34.8 \\
\bottomrule
\end{tabular}
\caption{Sample Data Table}
\label{tab:sample}
\end{table}

\section{Lists}
\subsection{Bullet Points}
\begin{itemize}
  \item First item
  \item Second item
  \item Third item with sub-items:
    \begin{itemize}
      \item Sub-item 1
      \item Sub-item 2
    \end{itemize}
\end{itemize}

\subsection{Numbered List}
\begin{enumerate}
  \item First step
  \item Second step
  \item Third step
\end{enumerate}

\section{Citations and References}
When you use the citation search feature, references will be automatically managed. For example, you might cite a paper like this: \cite{example2024}. All citations will be automatically added to your \texttt{cite.bib} file.

\section{Conclusion}
This template demonstrates basic \LaTeX{} features including mathematical formulas, tables, and lists. The bibliography system is now set up to work seamlessly with Bibby AI's citation management. Feel free to modify it for your needs!

\bibliographystyle{plain}
\bibliography{cite}

\end{document}\documentclass{article}
\usepackage[utf8]{inputenc}
\usepackage{booktabs}
\usepackage{amsmath}
\usepackage{natbib}

\title{My Document}
\author{Your Name}
\date{\today}

\begin{document}

\maketitle

\section{Introduction}
This is a sample document demonstrating various \LaTeX{} features including mathematical equations, tables, and formatting. You can cite papers using the citation search feature, and they will be automatically added to your \texttt{cite.bib} file.

\section{Mathematical Expressions}
Here are some example equations:

\subsection{Quadratic Formula}
The quadratic formula is given by:
\[x = \frac{-b \pm \sqrt{b^2 - 4ac}}{2a}\]

\subsection{Matrix Example}
A sample matrix:
\[
\begin{pmatrix}
a_{11} & a_{12} & a_{13} \\
a_{21} & a_{22} & a_{23} \\
a_{31} & a_{32} & a_{33}
\end{pmatrix}
\]

\section{Tables}
Below is a sample table showing data:

\begin{table}[h]
\centering
\begin{tabular}{@{}lrrr@{}}
\toprule
\textbf{Category} & \textbf{Value 1} & \textbf{Value 2} & \textbf{Total} \\
\midrule
Item A & 10.5 & 20.3 & 30.8 \\
Item B & 15.2 & 25.1 & 40.3 \\
Item C & 22.1 & 12.7 & 34.8 \\
\bottomrule
\end{tabular}
\caption{Sample Data Table}
\label{tab:sample}
\end{table}

\section{Lists}
\subsection{Bullet Points}
\begin{itemize}
  \item First item
  \item Second item
  \item Third item with sub-items:
    \begin{itemize}
      \item Sub-item 1
      \item Sub-item 2
    \end{itemize}
\end{itemize}

\subsection{Numbered List}
\begin{enumerate}
  \item First step
  \item Second step
  \item Third step
\end{enumerate}

\section{Citations and References}
When you use the citation search feature, references will be automatically managed. For example, you might cite a paper like this: \cite{example2024}. All citations will be automatically added to your \texttt{cite.bib} file.

\section{Conclusion}
This template demonstrates basic \LaTeX{} features including mathematical formulas, tables, and lists. The bibliography system is now set up to work seamlessly with Bibby AI's citation management. Feel free to modify it for your needs!

\bibliographystyle{plain}
\bibliography{cite}

\end{document}\documentclass{article}
\usepackage[utf8]{inputenc}
\usepackage{booktabs}
\usepackage{amsmath}
\usepackage{natbib}

\title{My Document}
\author{Your Name}
\date{\today}

\begin{document}

\maketitle

\section{Introduction}
This is a sample document demonstrating various \LaTeX{} features including mathematical equations, tables, and formatting. You can cite papers using the citation search feature, and they will be automatically added to your \texttt{cite.bib} file.

\section{Mathematical Expressions}
Here are some example equations:

\subsection{Quadratic Formula}
The quadratic formula is given by:
\[x = \frac{-b \pm \sqrt{b^2 - 4ac}}{2a}\]

\subsection{Matrix Example}
A sample matrix:
\[
\begin{pmatrix}
a_{11} & a_{12} & a_{13} \\
a_{21} & a_{22} & a_{23} \\
a_{31} & a_{32} & a_{33}
\end{pmatrix}
\]

\section{Tables}
Below is a sample table showing data:

\begin{table}[h]
\centering
\begin{tabular}{@{}lrrr@{}}
\toprule
\textbf{Category} & \textbf{Value 1} & \textbf{Value 2} & \textbf{Total} \\
\midrule
Item A & 10.5 & 20.3 & 30.8 \\
Item B & 15.2 & 25.1 & 40.3 \\
Item C & 22.1 & 12.7 & 34.8 \\
\bottomrule
\end{tabular}
\caption{Sample Data Table}
\label{tab:sample}
\end{table}

\section{Lists}
\subsection{Bullet Points}
\begin{itemize}
  \item First item
  \item Second item
  \item Third item with sub-items:
    \begin{itemize}
      \item Sub-item 1
      \item Sub-item 2
    \end{itemize}
\end{itemize}

\subsection{Numbered List}
\begin{enumerate}
  \item First step
  \item Second step
  \item Third step
\end{enumerate}

\section{Citations and References}
When you use the citation search feature, references will be automatically managed. For example, you might cite a paper like this: \cite{example2024}. All citations will be automatically added to your \texttt{cite.bib} file.

\section{Conclusion}
This template demonstrates basic \LaTeX{} features including mathematical formulas, tables, and lists. The bibliography system is now set up to work seamlessly with Bibby AI's citation management. Feel free to modify it for your needs!

\bibliographystyle{plain}
\bibliography{cite}

\end{document}\documentclass{article}
\usepackage[utf8]{inputenc}
\usepackage{booktabs}
\usepackage{amsmath}
\usepackage{natbib}

\title{My Document}
\author{Your Name}
\date{\today}

\begin{document}

\maketitle

\section{Introduction}
This is a sample document demonstrating various \LaTeX{} features including mathematical equations, tables, and formatting. You can cite papers using the citation search feature, and they will be automatically added to your \texttt{cite.bib} file.

\section{Mathematical Expressions}
Here are some example equations:

\subsection{Quadratic Formula}
The quadratic formula is given by:
\[x = \frac{-b \pm \sqrt{b^2 - 4ac}}{2a}\]

\subsection{Matrix Example}
A sample matrix:
\[
\begin{pmatrix}
a_{11} & a_{12} & a_{13} \\
a_{21} & a_{22} & a_{23} \\
a_{31} & a_{32} & a_{33}
\end{pmatrix}
\]

\section{Tables}
Below is a sample table showing data:

\begin{table}[h]
\centering
\begin{tabular}{@{}lrrr@{}}
\toprule
\textbf{Category} & \textbf{Value 1} & \textbf{Value 2} & \textbf{Total} \\
\midrule
Item A & 10.5 & 20.3 & 30.8 \\
Item B & 15.2 & 25.1 & 40.3 \\
Item C & 22.1 & 12.7 & 34.8 \\
\bottomrule
\end{tabular}
\caption{Sample Data Table}
\label{tab:sample}
\end{table}

\section{Lists}
\subsection{Bullet Points}
\begin{itemize}
  \item First item
  \item Second item
  \item Third item with sub-items:
    \begin{itemize}
      \item Sub-item 1
      \item Sub-item 2
    \end{itemize}
\end{itemize}

\subsection{Numbered List}
\begin{enumerate}
  \item First step
  \item Second step
  \item Third step
\end{enumerate}

\section{Citations and References}
When you use the citation search feature, references will be automatically managed. For example, you might cite a paper like this: \cite{example2024}. All citations will be automatically added to your \texttt{cite.bib} file.

\section{Conclusion}
This template demonstrates basic \LaTeX{} features including mathematical formulas, tables, and lists. The bibliography system is now set up to work seamlessly with Bibby AI's citation management. Feel free to modify it for your needs!

\bibliographystyle{plain}
\bibliography{cite}

\end{document}\documentclass{article}
\usepackage[utf8]{inputenc}
\usepackage{booktabs}
\usepackage{amsmath}
\usepackage{natbib}

\title{My Document}
\author{Your Name}
\date{\today}

\begin{document}

\maketitle

\section{Introduction}
This is a sample document demonstrating various \LaTeX{} features including mathematical equations, tables, and formatting. You can cite papers using the citation search feature, and they will be automatically added to your \texttt{cite.bib} file.

\section{Mathematical Expressions}
Here are some example equations:

\subsection{Quadratic Formula}
The quadratic formula is given by:
\[x = \frac{-b \pm \sqrt{b^2 - 4ac}}{2a}\]

\subsection{Matrix Example}
A sample matrix:
\[
\begin{pmatrix}
a_{11} & a_{12} & a_{13} \\
a_{21} & a_{22} & a_{23} \\
a_{31} & a_{32} & a_{33}
\end{pmatrix}
\]

\section{Tables}
Below is a sample table showing data:

\begin{table}[h]
\centering
\begin{tabular}{@{}lrrr@{}}
\toprule
\textbf{Category} & \textbf{Value 1} & \textbf{Value 2} & \textbf{Total} \\
\midrule
Item A & 10.5 & 20.3 & 30.8 \\
Item B & 15.2 & 25.1 & 40.3 \\
Item C & 22.1 & 12.7 & 34.8 \\
\bottomrule
\end{tabular}
\caption{Sample Data Table}
\label{tab:sample}
\end{table}

\section{Lists}
\subsection{Bullet Points}
\begin{itemize}
  \item First item
  \item Second item
  \item Third item with sub-items:
    \begin{itemize}
      \item Sub-item 1
      \item Sub-item 2
    \end{itemize}
\end{itemize}

\subsection{Numbered List}
\begin{enumerate}
  \item First step
  \item Second step
  \item Third step
\end{enumerate}

\section{Citations and References}
When you use the citation search feature, references will be automatically managed. For example, you might cite a paper like this: \cite{example2024}. All citations will be automatically added to your \texttt{cite.bib} file.

\section{Conclusion}
This template demonstrates basic \LaTeX{} features including mathematical formulas, tables, and lists. The bibliography system is now set up to work seamlessly with Bibby AI's citation management. Feel free to modify it for your needs!

\bibliographystyle{plain}
\bibliography{cite}

\end{document}\documentclass{article}
\usepackage[utf8]{inputenc}
\usepackage{booktabs}
\usepackage{amsmath}
\usepackage{natbib}

\title{My Document}
\author{Your Name}
\date{\today}

\begin{document}

\maketitle

\section{Introduction}
This is a sample document demonstrating various \LaTeX{} features including mathematical equations, tables, and formatting. You can cite papers using the citation search feature, and they will be automatically added to your \texttt{cite.bib} file.

\section{Mathematical Expressions}
Here are some example equations:

\subsection{Quadratic Formula}
The quadratic formula is given by:
\[x = \frac{-b \pm \sqrt{b^2 - 4ac}}{2a}\]

\subsection{Matrix Example}
A sample matrix:
\[
\begin{pmatrix}
a_{11} & a_{12} & a_{13} \\
a_{21} & a_{22} & a_{23} \\
a_{31} & a_{32} & a_{33}
\end{pmatrix}
\]

\section{Tables}
Below is a sample table showing data:

\begin{table}[h]
\centering
\begin{tabular}{@{}lrrr@{}}
\toprule
\textbf{Category} & \textbf{Value 1} & \textbf{Value 2} & \textbf{Total} \\
\midrule
Item A & 10.5 & 20.3 & 30.8 \\
Item B & 15.2 & 25.1 & 40.3 \\
Item C & 22.1 & 12.7 & 34.8 \\
\bottomrule
\end{tabular}
\caption{Sample Data Table}
\label{tab:sample}
\end{table}

\section{Lists}
\subsection{Bullet Points}
\begin{itemize}
  \item First item
  \item Second item
  \item Third item with sub-items:
    \begin{itemize}
      \item Sub-item 1
      \item Sub-item 2
    \end{itemize}
\end{itemize}

\subsection{Numbered List}
\begin{enumerate}
  \item First step
  \item Second step
  \item Third step
\end{enumerate}

\section{Citations and References}
When you use the citation search feature, references will be automatically managed. For example, you might cite a paper like this: \cite{example2024}. All citations will be automatically added to your \texttt{cite.bib} file.

\section{Conclusion}
This template demonstrates basic \LaTeX{} features including mathematical formulas, tables, and lists. The bibliography system is now set up to work seamlessly with Bibby AI's citation management. Feel free to modify it for your needs!

\bibliographystyle{plain}
\bibliography{cite}

\end{document}\documentclass{article}
\usepackage[utf8]{inputenc}
\usepackage{booktabs}
\usepackage{amsmath}
\usepackage{natbib}

\title{My Document}
\author{Your Name}
\date{\today}

\begin{document}

\maketitle

\section{Introduction}
This is a sample document demonstrating various \LaTeX{} features including mathematical equations, tables, and formatting. You can cite papers using the citation search feature, and they will be automatically added to your \texttt{cite.bib} file.

\section{Mathematical Expressions}
Here are some example equations:

\subsection{Quadratic Formula}
The quadratic formula is given by:
\[x = \frac{-b \pm \sqrt{b^2 - 4ac}}{2a}\]

\subsection{Matrix Example}
A sample matrix:
\[
\begin{pmatrix}
a_{11} & a_{12} & a_{13} \\
a_{21} & a_{22} & a_{23} \\
a_{31} & a_{32} & a_{33}
\end{pmatrix}
\]

\section{Tables}
Below is a sample table showing data:

\begin{table}[h]
\centering
\begin{tabular}{@{}lrrr@{}}
\toprule
\textbf{Category} & \textbf{Value 1} & \textbf{Value 2} & \textbf{Total} \\
\midrule
Item A & 10.5 & 20.3 & 30.8 \\
Item B & 15.2 & 25.1 & 40.3 \\
Item C & 22.1 & 12.7 & 34.8 \\
\bottomrule
\end{tabular}
\caption{Sample Data Table}
\label{tab:sample}
\end{table}

\section{Lists}
\subsection{Bullet Points}
\begin{itemize}
  \item First item
  \item Second item
  \item Third item with sub-items:
    \begin{itemize}
      \item Sub-item 1
      \item Sub-item 2
    \end{itemize}
\end{itemize}

\subsection{Numbered List}
\begin{enumerate}
  \item First step
  \item Second step
  \item Third step
\end{enumerate}

\section{Citations and References}
When you use the citation search feature, references will be automatically managed. For example, you might cite a paper like this: \cite{example2024}. All citations will be automatically added to your \texttt{cite.bib} file.

\section{Conclusion}
This template demonstrates basic \LaTeX{} features including mathematical formulas, tables, and lists. The bibliography system is now set up to work seamlessly with Bibby AI's citation management. Feel free to modify it for your needs!

\bibliographystyle{plain}
\bibliography{cite}

\end{document}\documentclass{article}
\usepackage[utf8]{inputenc}
\usepackage{booktabs}
\usepackage{amsmath}
\usepackage{natbib}

\title{My Document}
\author{Your Name}
\date{\today}

\begin{document}

\maketitle

\section{Introduction}
This is a sample document demonstrating various \LaTeX{} features including mathematical equations, tables, and formatting. You can cite papers using the citation search feature, and they will be automatically added to your \texttt{cite.bib} file.

\section{Mathematical Expressions}
Here are some example equations:

\subsection{Quadratic Formula}
The quadratic formula is given by:
\[x = \frac{-b \pm \sqrt{b^2 - 4ac}}{2a}\]

\subsection{Matrix Example}
A sample matrix:
\[
\begin{pmatrix}
a_{11} & a_{12} & a_{13} \\
a_{21} & a_{22} & a_{23} \\
a_{31} & a_{32} & a_{33}
\end{pmatrix}
\]

\section{Tables}
Below is a sample table showing data:

\begin{table}[h]
\centering
\begin{tabular}{@{}lrrr@{}}
\toprule
\textbf{Category} & \textbf{Value 1} & \textbf{Value 2} & \textbf{Total} \\
\midrule
Item A & 10.5 & 20.3 & 30.8 \\
Item B & 15.2 & 25.1 & 40.3 \\
Item C & 22.1 & 12.7 & 34.8 \\
\bottomrule
\end{tabular}
\caption{Sample Data Table}
\label{tab:sample}
\end{table}

\section{Lists}
\subsection{Bullet Points}
\begin{itemize}
  \item First item
  \item Second item
  \item Third item with sub-items:
    \begin{itemize}
      \item Sub-item 1
      \item Sub-item 2
    \end{itemize}
\end{itemize}

\subsection{Numbered List}
\begin{enumerate}
  \item First step
  \item Second step
  \item Third step
\end{enumerate}

\section{Citations and References}
When you use the citation search feature, references will be automatically managed. For example, you might cite a paper like this: \cite{example2024}. All citations will be automatically added to your \texttt{cite.bib} file.

\section{Conclusion}
This template demonstrates basic \LaTeX{} features including mathematical formulas, tables, and lists. The bibliography system is now set up to work seamlessly with Bibby AI's citation management. Feel free to modify it for your needs!

\bibliographystyle{plain}
\bibliography{cite}

\end{document}\documentclass{article}
\usepackage[utf8]{inputenc}
\usepackage{booktabs}
\usepackage{amsmath}
\usepackage{natbib}

\title{My Document}
\author{Your Name}
\date{\today}

\begin{document}

\maketitle

\section{Introduction}
This is a sample document demonstrating various \LaTeX{} features including mathematical equations, tables, and formatting. You can cite papers using the citation search feature, and they will be automatically added to your \texttt{cite.bib} file.

\section{Mathematical Expressions}
Here are some example equations:

\subsection{Quadratic Formula}
The quadratic formula is given by:
\[x = \frac{-b \pm \sqrt{b^2 - 4ac}}{2a}\]

\subsection{Matrix Example}
A sample matrix:
\[
\begin{pmatrix}
a_{11} & a_{12} & a_{13} \\
a_{21} & a_{22} & a_{23} \\
a_{31} & a_{32} & a_{33}
\end{pmatrix}
\]

\section{Tables}
Below is a sample table showing data:

\begin{table}[h]
\centering
\begin{tabular}{@{}lrrr@{}}
\toprule
\textbf{Category} & \textbf{Value 1} & \textbf{Value 2} & \textbf{Total} \\
\midrule
Item A & 10.5 & 20.3 & 30.8 \\
Item B & 15.2 & 25.1 & 40.3 \\
Item C & 22.1 & 12.7 & 34.8 \\
\bottomrule
\end{tabular}
\caption{Sample Data Table}
\label{tab:sample}
\end{table}

\section{Lists}
\subsection{Bullet Points}
\begin{itemize}
  \item First item
  \item Second item
  \item Third item with sub-items:
    \begin{itemize}
      \item Sub-item 1
      \item Sub-item 2
    \end{itemize}
\end{itemize}

\subsection{Numbered List}
\begin{enumerate}
  \item First step
  \item Second step
  \item Third step
\end{enumerate}

\section{Citations and References}
When you use the citation search feature, references will be automatically managed. For example, you might cite a paper like this: \cite{example2024}. All citations will be automatically added to your \texttt{cite.bib} file.

\section{Conclusion}
This template demonstrates basic \LaTeX{} features including mathematical formulas, tables, and lists. The bibliography system is now set up to work seamlessly with Bibby AI's citation management. Feel free to modify it for your needs!

\bibliographystyle{plain}
\bibliography{cite}

\end{document}\documentclass{article}
\usepackage[utf8]{inputenc}
\usepackage{booktabs}
\usepackage{amsmath}
\usepackage{natbib}

\title{My Document}
\author{Your Name}
\date{\today}

\begin{document}

\maketitle

\section{Introduction}
This is a sample document demonstrating various \LaTeX{} features including mathematical equations, tables, and formatting. You can cite papers using the citation search feature, and they will be automatically added to your \texttt{cite.bib} file.

\section{Mathematical Expressions}
Here are some example equations:

\subsection{Quadratic Formula}
The quadratic formula is given by:
\[x = \frac{-b \pm \sqrt{b^2 - 4ac}}{2a}\]

\subsection{Matrix Example}
A sample matrix:
\[
\begin{pmatrix}
a_{11} & a_{12} & a_{13} \\
a_{21} & a_{22} & a_{23} \\
a_{31} & a_{32} & a_{33}
\end{pmatrix}
\]

\section{Tables}
Below is a sample table showing data:

\begin{table}[h]
\centering
\begin{tabular}{@{}lrrr@{}}
\toprule
\textbf{Category} & \textbf{Value 1} & \textbf{Value 2} & \textbf{Total} \\
\midrule
Item A & 10.5 & 20.3 & 30.8 \\
Item B & 15.2 & 25.1 & 40.3 \\
Item C & 22.1 & 12.7 & 34.8 \\
\bottomrule
\end{tabular}
\caption{Sample Data Table}
\label{tab:sample}
\end{table}

\section{Lists}
\subsection{Bullet Points}
\begin{itemize}
  \item First item
  \item Second item
  \item Third item with sub-items:
    \begin{itemize}
      \item Sub-item 1
      \item Sub-item 2
    \end{itemize}
\end{itemize}

\subsection{Numbered List}
\begin{enumerate}
  \item First step
  \item Second step
  \item Third step
\end{enumerate}

\section{Citations and References}
When you use the citation search feature, references will be automatically managed. For example, you might cite a paper like this: \cite{example2024}. All citations will be automatically added to your \texttt{cite.bib} file.

\section{Conclusion}
This template demonstrates basic \LaTeX{} features including mathematical formulas, tables, and lists. The bibliography system is now set up to work seamlessly with Bibby AI's citation management. Feel free to modify it for your needs!

\bibliographystyle{plain}
\bibliography{cite}

\end{document}\documentclass{article}
\usepackage[utf8]{inputenc}
\usepackage{booktabs}
\usepackage{amsmath}
\usepackage{natbib}

\title{My Document}
\author{Your Name}
\date{\today}

\begin{document}

\maketitle

\section{Introduction}
This is a sample document demonstrating various \LaTeX{} features including mathematical equations, tables, and formatting. You can cite papers using the citation search feature, and they will be automatically added to your \texttt{cite.bib} file.

\section{Mathematical Expressions}
Here are some example equations:

\subsection{Quadratic Formula}
The quadratic formula is given by:
\[x = \frac{-b \pm \sqrt{b^2 - 4ac}}{2a}\]

\subsection{Matrix Example}
A sample matrix:
\[
\begin{pmatrix}
a_{11} & a_{12} & a_{13} \\
a_{21} & a_{22} & a_{23} \\
a_{31} & a_{32} & a_{33}
\end{pmatrix}
\]

\section{Tables}
Below is a sample table showing data:

\begin{table}[h]
\centering
\begin{tabular}{@{}lrrr@{}}
\toprule
\textbf{Category} & \textbf{Value 1} & \textbf{Value 2} & \textbf{Total} \\
\midrule
Item A & 10.5 & 20.3 & 30.8 \\
Item B & 15.2 & 25.1 & 40.3 \\
Item C & 22.1 & 12.7 & 34.8 \\
\bottomrule
\end{tabular}
\caption{Sample Data Table}
\label{tab:sample}
\end{table}

\section{Lists}
\subsection{Bullet Points}
\begin{itemize}
  \item First item
  \item Second item
  \item Third item with sub-items:
    \begin{itemize}
      \item Sub-item 1
      \item Sub-item 2
    \end{itemize}
\end{itemize}

\subsection{Numbered List}
\begin{enumerate}
  \item First step
  \item Second step
  \item Third step
\end{enumerate}

\section{Citations and References}
When you use the citation search feature, references will be automatically managed. For example, you might cite a paper like this: \cite{example2024}. All citations will be automatically added to your \texttt{cite.bib} file.

\section{Conclusion}
This template demonstrates basic \LaTeX{} features including mathematical formulas, tables, and lists. The bibliography system is now set up to work seamlessly with Bibby AI's citation management. Feel free to modify it for your needs!

\bibliographystyle{plain}
\bibliography{cite}

\end{document}\documentclass{article}
\usepackage[utf8]{inputenc}
\usepackage{booktabs}
\usepackage{amsmath}
\usepackage{natbib}

\title{My Document}
\author{Your Name}
\date{\today}

\begin{document}

\maketitle

\section{Introduction}
This is a sample document demonstrating various \LaTeX{} features including mathematical equations, tables, and formatting. You can cite papers using the citation search feature, and they will be automatically added to your \texttt{cite.bib} file.

\section{Mathematical Expressions}
Here are some example equations:

\subsection{Quadratic Formula}
The quadratic formula is given by:
\[x = \frac{-b \pm \sqrt{b^2 - 4ac}}{2a}\]

\subsection{Matrix Example}
A sample matrix:
\[
\begin{pmatrix}
a_{11} & a_{12} & a_{13} \\
a_{21} & a_{22} & a_{23} \\
a_{31} & a_{32} & a_{33}
\end{pmatrix}
\]

\section{Tables}
Below is a sample table showing data:

\begin{table}[h]
\centering
\begin{tabular}{@{}lrrr@{}}
\toprule
\textbf{Category} & \textbf{Value 1} & \textbf{Value 2} & \textbf{Total} \\
\midrule
Item A & 10.5 & 20.3 & 30.8 \\
Item B & 15.2 & 25.1 & 40.3 \\
Item C & 22.1 & 12.7 & 34.8 \\
\bottomrule
\end{tabular}
\caption{Sample Data Table}
\label{tab:sample}
\end{table}

\section{Lists}
\subsection{Bullet Points}
\begin{itemize}
  \item First item
  \item Second item
  \item Third item with sub-items:
    \begin{itemize}
      \item Sub-item 1
      \item Sub-item 2
    \end{itemize}
\end{itemize}

\subsection{Numbered List}
\begin{enumerate}
  \item First step
  \item Second step
  \item Third step
\end{enumerate}

\section{Citations and References}
When you use the citation search feature, references will be automatically managed. For example, you might cite a paper like this: \cite{example2024}. All citations will be automatically added to your \texttt{cite.bib} file.

\section{Conclusion}
This template demonstrates basic \LaTeX{} features including mathematical formulas, tables, and lists. The bibliography system is now set up to work seamlessly with Bibby AI's citation management. Feel free to modify it for your needs!

\bibliographystyle{plain}
\bibliography{cite}

\end{document}\documentclass{article}
\usepackage[utf8]{inputenc}
\usepackage{booktabs}
\usepackage{amsmath}
\usepackage{natbib}

\title{My Document}
\author{Your Name}
\date{\today}

\begin{document}

\maketitle

\section{Introduction}
This is a sample document demonstrating various \LaTeX{} features including mathematical equations, tables, and formatting. You can cite papers using the citation search feature, and they will be automatically added to your \texttt{cite.bib} file.

\section{Mathematical Expressions}
Here are some example equations:

\subsection{Quadratic Formula}
The quadratic formula is given by:
\[x = \frac{-b \pm \sqrt{b^2 - 4ac}}{2a}\]

\subsection{Matrix Example}
A sample matrix:
\[
\begin{pmatrix}
a_{11} & a_{12} & a_{13} \\
a_{21} & a_{22} & a_{23} \\
a_{31} & a_{32} & a_{33}
\end{pmatrix}
\]

\section{Tables}
Below is a sample table showing data:

\begin{table}[h]
\centering
\begin{tabular}{@{}lrrr@{}}
\toprule
\textbf{Category} & \textbf{Value 1} & \textbf{Value 2} & \textbf{Total} \\
\midrule
Item A & 10.5 & 20.3 & 30.8 \\
Item B & 15.2 & 25.1 & 40.3 \\
Item C & 22.1 & 12.7 & 34.8 \\
\bottomrule
\end{tabular}
\caption{Sample Data Table}
\label{tab:sample}
\end{table}

\section{Lists}
\subsection{Bullet Points}
\begin{itemize}
  \item First item
  \item Second item
  \item Third item with sub-items:
    \begin{itemize}
      \item Sub-item 1
      \item Sub-item 2
    \end{itemize}
\end{itemize}

\subsection{Numbered List}
\begin{enumerate}
  \item First step
  \item Second step
  \item Third step
\end{enumerate}

\section{Citations and References}
When you use the citation search feature, references will be automatically managed. For example, you might cite a paper like this: \cite{example2024}. All citations will be automatically added to your \texttt{cite.bib} file.

\section{Conclusion}
This template demonstrates basic \LaTeX{} features including mathematical formulas, tables, and lists. The bibliography system is now set up to work seamlessly with Bibby AI's citation management. Feel free to modify it for your needs!

\bibliographystyle{plain}
\bibliography{cite}

\end{document}\documentclass{article}
\usepackage[utf8]{inputenc}
\usepackage{booktabs}
\usepackage{amsmath}
\usepackage{natbib}

\title{My Document}
\author{Your Name}
\date{\today}

\begin{document}

\maketitle

\section{Introduction}
This is a sample document demonstrating various \LaTeX{} features including mathematical equations, tables, and formatting. You can cite papers using the citation search feature, and they will be automatically added to your \texttt{cite.bib} file.

\section{Mathematical Expressions}
Here are some example equations:

\subsection{Quadratic Formula}
The quadratic formula is given by:
\[x = \frac{-b \pm \sqrt{b^2 - 4ac}}{2a}\]

\subsection{Matrix Example}
A sample matrix:
\[
\begin{pmatrix}
a_{11} & a_{12} & a_{13} \\
a_{21} & a_{22} & a_{23} \\
a_{31} & a_{32} & a_{33}
\end{pmatrix}
\]

\section{Tables}
Below is a sample table showing data:

\begin{table}[h]
\centering
\begin{tabular}{@{}lrrr@{}}
\toprule
\textbf{Category} & \textbf{Value 1} & \textbf{Value 2} & \textbf{Total} \\
\midrule
Item A & 10.5 & 20.3 & 30.8 \\
Item B & 15.2 & 25.1 & 40.3 \\
Item C & 22.1 & 12.7 & 34.8 \\
\bottomrule
\end{tabular}
\caption{Sample Data Table}
\label{tab:sample}
\end{table}

\section{Lists}
\subsection{Bullet Points}
\begin{itemize}
  \item First item
  \item Second item
  \item Third item with sub-items:
    \begin{itemize}
      \item Sub-item 1
      \item Sub-item 2
    \end{itemize}
\end{itemize}

\subsection{Numbered List}
\begin{enumerate}
  \item First step
  \item Second step
  \item Third step
\end{enumerate}

\section{Citations and References}
When you use the citation search feature, references will be automatically managed. For example, you might cite a paper like this: \cite{example2024}. All citations will be automatically added to your \texttt{cite.bib} file.

\section{Conclusion}
This template demonstrates basic \LaTeX{} features including mathematical formulas, tables, and lists. The bibliography system is now set up to work seamlessly with Bibby AI's citation management. Feel free to modify it for your needs!

\bibliographystyle{plain}
\bibliography{cite}

\end{document}\documentclass{article}
\usepackage[utf8]{inputenc}
\usepackage{booktabs}
\usepackage{amsmath}
\usepackage{natbib}

\title{My Document}
\author{Your Name}
\date{\today}

\begin{document}

\maketitle

\section{Introduction}
This is a sample document demonstrating various \LaTeX{} features including mathematical equations, tables, and formatting. You can cite papers using the citation search feature, and they will be automatically added to your \texttt{cite.bib} file.

\section{Mathematical Expressions}
Here are some example equations:

\subsection{Quadratic Formula}
The quadratic formula is given by:
\[x = \frac{-b \pm \sqrt{b^2 - 4ac}}{2a}\]

\subsection{Matrix Example}
A sample matrix:
\[
\begin{pmatrix}
a_{11} & a_{12} & a_{13} \\
a_{21} & a_{22} & a_{23} \\
a_{31} & a_{32} & a_{33}
\end{pmatrix}
\]

\section{Tables}
Below is a sample table showing data:

\begin{table}[h]
\centering
\begin{tabular}{@{}lrrr@{}}
\toprule
\textbf{Category} & \textbf{Value 1} & \textbf{Value 2} & \textbf{Total} \\
\midrule
Item A & 10.5 & 20.3 & 30.8 \\
Item B & 15.2 & 25.1 & 40.3 \\
Item C & 22.1 & 12.7 & 34.8 \\
\bottomrule
\end{tabular}
\caption{Sample Data Table}
\label{tab:sample}
\end{table}

\section{Lists}
\subsection{Bullet Points}
\begin{itemize}
  \item First item
  \item Second item
  \item Third item with sub-items:
    \begin{itemize}
      \item Sub-item 1
      \item Sub-item 2
    \end{itemize}
\end{itemize}

\subsection{Numbered List}
\begin{enumerate}
  \item First step
  \item Second step
  \item Third step
\end{enumerate}

\section{Citations and References}
When you use the citation search feature, references will be automatically managed. For example, you might cite a paper like this: \cite{example2024}. All citations will be automatically added to your \texttt{cite.bib} file.

\section{Conclusion}
This template demonstrates basic \LaTeX{} features including mathematical formulas, tables, and lists. The bibliography system is now set up to work seamlessly with Bibby AI's citation management. Feel free to modify it for your needs!

\bibliographystyle{plain}
\bibliography{cite}

\end{document}\documentclass{article}
\usepackage[utf8]{inputenc}
\usepackage{booktabs}
\usepackage{amsmath}
\usepackage{natbib}

\title{My Document}
\author{Your Name}
\date{\today}

\begin{document}

\maketitle

\section{Introduction}
This is a sample document demonstrating various \LaTeX{} features including mathematical equations, tables, and formatting. You can cite papers using the citation search feature, and they will be automatically added to your \texttt{cite.bib} file.

\section{Mathematical Expressions}
Here are some example equations:

\subsection{Quadratic Formula}
The quadratic formula is given by:
\[x = \frac{-b \pm \sqrt{b^2 - 4ac}}{2a}\]

\subsection{Matrix Example}
A sample matrix:
\[
\begin{pmatrix}
a_{11} & a_{12} & a_{13} \\
a_{21} & a_{22} & a_{23} \\
a_{31} & a_{32} & a_{33}
\end{pmatrix}
\]

\section{Tables}
Below is a sample table showing data:

\begin{table}[h]
\centering
\begin{tabular}{@{}lrrr@{}}
\toprule
\textbf{Category} & \textbf{Value 1} & \textbf{Value 2} & \textbf{Total} \\
\midrule
Item A & 10.5 & 20.3 & 30.8 \\
Item B & 15.2 & 25.1 & 40.3 \\
Item C & 22.1 & 12.7 & 34.8 \\
\bottomrule
\end{tabular}
\caption{Sample Data Table}
\label{tab:sample}
\end{table}

\section{Lists}
\subsection{Bullet Points}
\begin{itemize}
  \item First item
  \item Second item
  \item Third item with sub-items:
    \begin{itemize}
      \item Sub-item 1
      \item Sub-item 2
    \end{itemize}
\end{itemize}

\subsection{Numbered List}
\begin{enumerate}
  \item First step
  \item Second step
  \item Third step
\end{enumerate}

\section{Citations and References}
When you use the citation search feature, references will be automatically managed. For example, you might cite a paper like this: \cite{example2024}. All citations will be automatically added to your \texttt{cite.bib} file.

\section{Conclusion}
This template demonstrates basic \LaTeX{} features including mathematical formulas, tables, and lists. The bibliography system is now set up to work seamlessly with Bibby AI's citation management. Feel free to modify it for your needs!

\bibliographystyle{plain}
\bibliography{cite}

\end{document}\documentclass{article}
\usepackage[utf8]{inputenc}
\usepackage{booktabs}
\usepackage{amsmath}
\usepackage{natbib}

\title{My Document}
\author{Your Name}
\date{\today}

\begin{document}

\maketitle

\section{Introduction}
This is a sample document demonstrating various \LaTeX{} features including mathematical equations, tables, and formatting. You can cite papers using the citation search feature, and they will be automatically added to your \texttt{cite.bib} file.

\section{Mathematical Expressions}
Here are some example equations:

\subsection{Quadratic Formula}
The quadratic formula is given by:
\[x = \frac{-b \pm \sqrt{b^2 - 4ac}}{2a}\]

\subsection{Matrix Example}
A sample matrix:
\[
\begin{pmatrix}
a_{11} & a_{12} & a_{13} \\
a_{21} & a_{22} & a_{23} \\
a_{31} & a_{32} & a_{33}
\end{pmatrix}
\]

\section{Tables}
Below is a sample table showing data:

\begin{table}[h]
\centering
\begin{tabular}{@{}lrrr@{}}
\toprule
\textbf{Category} & \textbf{Value 1} & \textbf{Value 2} & \textbf{Total} \\
\midrule
Item A & 10.5 & 20.3 & 30.8 \\
Item B & 15.2 & 25.1 & 40.3 \\
Item C & 22.1 & 12.7 & 34.8 \\
\bottomrule
\end{tabular}
\caption{Sample Data Table}
\label{tab:sample}
\end{table}

\section{Lists}
\subsection{Bullet Points}
\begin{itemize}
  \item First item
  \item Second item
  \item Third item with sub-items:
    \begin{itemize}
      \item Sub-item 1
      \item Sub-item 2
    \end{itemize}
\end{itemize}

\subsection{Numbered List}
\begin{enumerate}
  \item First step
  \item Second step
  \item Third step
\end{enumerate}

\section{Citations and References}
When you use the citation search feature, references will be automatically managed. For example, you might cite a paper like this: \cite{example2024}. All citations will be automatically added to your \texttt{cite.bib} file.

\section{Conclusion}
This template demonstrates basic \LaTeX{} features including mathematical formulas, tables, and lists. The bibliography system is now set up to work seamlessly with Bibby AI's citation management. Feel free to modify it for your needs!

\bibliographystyle{plain}
\bibliography{cite}

\end{document}\documentclass{article}
\usepackage[utf8]{inputenc}
\usepackage{booktabs}
\usepackage{amsmath}
\usepackage{natbib}

\title{My Document}
\author{Your Name}
\date{\today}

\begin{document}

\maketitle

\section{Introduction}
This is a sample document demonstrating various \LaTeX{} features including mathematical equations, tables, and formatting. You can cite papers using the citation search feature, and they will be automatically added to your \texttt{cite.bib} file.

\section{Mathematical Expressions}
Here are some example equations:

\subsection{Quadratic Formula}
The quadratic formula is given by:
\[x = \frac{-b \pm \sqrt{b^2 - 4ac}}{2a}\]

\subsection{Matrix Example}
A sample matrix:
\[
\begin{pmatrix}
a_{11} & a_{12} & a_{13} \\
a_{21} & a_{22} & a_{23} \\
a_{31} & a_{32} & a_{33}
\end{pmatrix}
\]

\section{Tables}
Below is a sample table showing data:

\begin{table}[h]
\centering
\begin{tabular}{@{}lrrr@{}}
\toprule
\textbf{Category} & \textbf{Value 1} & \textbf{Value 2} & \textbf{Total} \\
\midrule
Item A & 10.5 & 20.3 & 30.8 \\
Item B & 15.2 & 25.1 & 40.3 \\
Item C & 22.1 & 12.7 & 34.8 \\
\bottomrule
\end{tabular}
\caption{Sample Data Table}
\label{tab:sample}
\end{table}

\section{Lists}
\subsection{Bullet Points}
\begin{itemize}
  \item First item
  \item Second item
  \item Third item with sub-items:
    \begin{itemize}
      \item Sub-item 1
      \item Sub-item 2
    \end{itemize}
\end{itemize}

\subsection{Numbered List}
\begin{enumerate}
  \item First step
  \item Second step
  \item Third step
\end{enumerate}

\section{Citations and References}
When you use the citation search feature, references will be automatically managed. For example, you might cite a paper like this: \cite{example2024}. All citations will be automatically added to your \texttt{cite.bib} file.

\section{Conclusion}
This template demonstrates basic \LaTeX{} features including mathematical formulas, tables, and lists. The bibliography system is now set up to work seamlessly with Bibby AI's citation management. Feel free to modify it for your needs!

\bibliographystyle{plain}
\bibliography{cite}

\end{document}\documentclass{article}
\usepackage[utf8]{inputenc}
\usepackage{booktabs}
\usepackage{amsmath}
\usepackage{natbib}

\title{My Document}
\author{Your Name}
\date{\today}

\begin{document}

\maketitle

\section{Introduction}
This is a sample document demonstrating various \LaTeX{} features including mathematical equations, tables, and formatting. You can cite papers using the citation search feature, and they will be automatically added to your \texttt{cite.bib} file.

\section{Mathematical Expressions}
Here are some example equations:

\subsection{Quadratic Formula}
The quadratic formula is given by:
\[x = \frac{-b \pm \sqrt{b^2 - 4ac}}{2a}\]

\subsection{Matrix Example}
A sample matrix:
\[
\begin{pmatrix}
a_{11} & a_{12} & a_{13} \\
a_{21} & a_{22} & a_{23} \\
a_{31} & a_{32} & a_{33}
\end{pmatrix}
\]

\section{Tables}
Below is a sample table showing data:

\begin{table}[h]
\centering
\begin{tabular}{@{}lrrr@{}}
\toprule
\textbf{Category} & \textbf{Value 1} & \textbf{Value 2} & \textbf{Total} \\
\midrule
Item A & 10.5 & 20.3 & 30.8 \\
Item B & 15.2 & 25.1 & 40.3 \\
Item C & 22.1 & 12.7 & 34.8 \\
\bottomrule
\end{tabular}
\caption{Sample Data Table}
\label{tab:sample}
\end{table}

\section{Lists}
\subsection{Bullet Points}
\begin{itemize}
  \item First item
  \item Second item
  \item Third item with sub-items:
    \begin{itemize}
      \item Sub-item 1
      \item Sub-item 2
    \end{itemize}
\end{itemize}

\subsection{Numbered List}
\begin{enumerate}
  \item First step
  \item Second step
  \item Third step
\end{enumerate}

\section{Citations and References}
When you use the citation search feature, references will be automatically managed. For example, you might cite a paper like this: \cite{example2024}. All citations will be automatically added to your \texttt{cite.bib} file.

\section{Conclusion}
This template demonstrates basic \LaTeX{} features including mathematical formulas, tables, and lists. The bibliography system is now set up to work seamlessly with Bibby AI's citation management. Feel free to modify it for your needs!

\bibliographystyle{plain}
\bibliography{cite}

\end{document}\documentclass{article}
\usepackage[utf8]{inputenc}
\usepackage{booktabs}
\usepackage{amsmath}
\usepackage{natbib}

\title{My Document}
\author{Your Name}
\date{\today}

\begin{document}

\maketitle

\section{Introduction}
This is a sample document demonstrating various \LaTeX{} features including mathematical equations, tables, and formatting. You can cite papers using the citation search feature, and they will be automatically added to your \texttt{cite.bib} file.

\section{Mathematical Expressions}
Here are some example equations:

\subsection{Quadratic Formula}
The quadratic formula is given by:
\[x = \frac{-b \pm \sqrt{b^2 - 4ac}}{2a}\]

\subsection{Matrix Example}
A sample matrix:
\[
\begin{pmatrix}
a_{11} & a_{12} & a_{13} \\
a_{21} & a_{22} & a_{23} \\
a_{31} & a_{32} & a_{33}
\end{pmatrix}
\]

\section{Tables}
Below is a sample table showing data:

\begin{table}[h]
\centering
\begin{tabular}{@{}lrrr@{}}
\toprule
\textbf{Category} & \textbf{Value 1} & \textbf{Value 2} & \textbf{Total} \\
\midrule
Item A & 10.5 & 20.3 & 30.8 \\
Item B & 15.2 & 25.1 & 40.3 \\
Item C & 22.1 & 12.7 & 34.8 \\
\bottomrule
\end{tabular}
\caption{Sample Data Table}
\label{tab:sample}
\end{table}

\section{Lists}
\subsection{Bullet Points}
\begin{itemize}
  \item First item
  \item Second item
  \item Third item with sub-items:
    \begin{itemize}
      \item Sub-item 1
      \item Sub-item 2
    \end{itemize}
\end{itemize}

\subsection{Numbered List}
\begin{enumerate}
  \item First step
  \item Second step
  \item Third step
\end{enumerate}

\section{Citations and References}
When you use the citation search feature, references will be automatically managed. For example, you might cite a paper like this: \cite{example2024}. All citations will be automatically added to your \texttt{cite.bib} file.

\section{Conclusion}
This template demonstrates basic \LaTeX{} features including mathematical formulas, tables, and lists. The bibliography system is now set up to work seamlessly with Bibby AI's citation management. Feel free to modify it for your needs!

\bibliographystyle{plain}
\bibliography{cite}

\end{document}\documentclass{article}
\usepackage[utf8]{inputenc}
\usepackage{booktabs}
\usepackage{amsmath}
\usepackage{natbib}

\title{My Document}
\author{Your Name}
\date{\today}

\begin{document}

\maketitle

\section{Introduction}
This is a sample document demonstrating various \LaTeX{} features including mathematical equations, tables, and formatting. You can cite papers using the citation search feature, and they will be automatically added to your \texttt{cite.bib} file.

\section{Mathematical Expressions}
Here are some example equations:

\subsection{Quadratic Formula}
The quadratic formula is given by:
\[x = \frac{-b \pm \sqrt{b^2 - 4ac}}{2a}\]

\subsection{Matrix Example}
A sample matrix:
\[
\begin{pmatrix}
a_{11} & a_{12} & a_{13} \\
a_{21} & a_{22} & a_{23} \\
a_{31} & a_{32} & a_{33}
\end{pmatrix}
\]

\section{Tables}
Below is a sample table showing data:

\begin{table}[h]
\centering
\begin{tabular}{@{}lrrr@{}}
\toprule
\textbf{Category} & \textbf{Value 1} & \textbf{Value 2} & \textbf{Total} \\
\midrule
Item A & 10.5 & 20.3 & 30.8 \\
Item B & 15.2 & 25.1 & 40.3 \\
Item C & 22.1 & 12.7 & 34.8 \\
\bottomrule
\end{tabular}
\caption{Sample Data Table}
\label{tab:sample}
\end{table}

\section{Lists}
\subsection{Bullet Points}
\begin{itemize}
  \item First item
  \item Second item
  \item Third item with sub-items:
    \begin{itemize}
      \item Sub-item 1
      \item Sub-item 2
    \end{itemize}
\end{itemize}

\subsection{Numbered List}
\begin{enumerate}
  \item First step
  \item Second step
  \item Third step
\end{enumerate}

\section{Citations and References}
When you use the citation search feature, references will be automatically managed. For example, you might cite a paper like this: \cite{example2024}. All citations will be automatically added to your \texttt{cite.bib} file.

\section{Conclusion}
This template demonstrates basic \LaTeX{} features including mathematical formulas, tables, and lists. The bibliography system is now set up to work seamlessly with Bibby AI's citation management. Feel free to modify it for your needs!

\bibliographystyle{plain}
\bibliography{cite}

\end{document}\documentclass{article}
\usepackage[utf8]{inputenc}
\usepackage{booktabs}
\usepackage{amsmath}
\usepackage{natbib}

\title{My Document}
\author{Your Name}
\date{\today}

\begin{document}

\maketitle

\section{Introduction}
This is a sample document demonstrating various \LaTeX{} features including mathematical equations, tables, and formatting. You can cite papers using the citation search feature, and they will be automatically added to your \texttt{cite.bib} file.

\section{Mathematical Expressions}
Here are some example equations:

\subsection{Quadratic Formula}
The quadratic formula is given by:
\[x = \frac{-b \pm \sqrt{b^2 - 4ac}}{2a}\]

\subsection{Matrix Example}
A sample matrix:
\[
\begin{pmatrix}
a_{11} & a_{12} & a_{13} \\
a_{21} & a_{22} & a_{23} \\
a_{31} & a_{32} & a_{33}
\end{pmatrix}
\]

\section{Tables}
Below is a sample table showing data:

\begin{table}[h]
\centering
\begin{tabular}{@{}lrrr@{}}
\toprule
\textbf{Category} & \textbf{Value 1} & \textbf{Value 2} & \textbf{Total} \\
\midrule
Item A & 10.5 & 20.3 & 30.8 \\
Item B & 15.2 & 25.1 & 40.3 \\
Item C & 22.1 & 12.7 & 34.8 \\
\bottomrule
\end{tabular}
\caption{Sample Data Table}
\label{tab:sample}
\end{table}

\section{Lists}
\subsection{Bullet Points}
\begin{itemize}
  \item First item
  \item Second item
  \item Third item with sub-items:
    \begin{itemize}
      \item Sub-item 1
      \item Sub-item 2
    \end{itemize}
\end{itemize}

\subsection{Numbered List}
\begin{enumerate}
  \item First step
  \item Second step
  \item Third step
\end{enumerate}

\section{Citations and References}
When you use the citation search feature, references will be automatically managed. For example, you might cite a paper like this: \cite{example2024}. All citations will be automatically added to your \texttt{cite.bib} file.

\section{Conclusion}
This template demonstrates basic \LaTeX{} features including mathematical formulas, tables, and lists. The bibliography system is now set up to work seamlessly with Bibby AI's citation management. Feel free to modify it for your needs!

\bibliographystyle{plain}
\bibliography{cite}

\end{document}\documentclass{article}
\usepackage[utf8]{inputenc}
\usepackage{booktabs}
\usepackage{amsmath}
\usepackage{natbib}

\title{My Document}
\author{Your Name}
\date{\today}

\begin{document}

\maketitle

\section{Introduction}
This is a sample document demonstrating various \LaTeX{} features including mathematical equations, tables, and formatting. You can cite papers using the citation search feature, and they will be automatically added to your \texttt{cite.bib} file.

\section{Mathematical Expressions}
Here are some example equations:

\subsection{Quadratic Formula}
The quadratic formula is given by:
\[x = \frac{-b \pm \sqrt{b^2 - 4ac}}{2a}\]

\subsection{Matrix Example}
A sample matrix:
\[
\begin{pmatrix}
a_{11} & a_{12} & a_{13} \\
a_{21} & a_{22} & a_{23} \\
a_{31} & a_{32} & a_{33}
\end{pmatrix}
\]

\section{Tables}
Below is a sample table showing data:

\begin{table}[h]
\centering
\begin{tabular}{@{}lrrr@{}}
\toprule
\textbf{Category} & \textbf{Value 1} & \textbf{Value 2} & \textbf{Total} \\
\midrule
Item A & 10.5 & 20.3 & 30.8 \\
Item B & 15.2 & 25.1 & 40.3 \\
Item C & 22.1 & 12.7 & 34.8 \\
\bottomrule
\end{tabular}
\caption{Sample Data Table}
\label{tab:sample}
\end{table}

\section{Lists}
\subsection{Bullet Points}
\begin{itemize}
  \item First item
  \item Second item
  \item Third item with sub-items:
    \begin{itemize}
      \item Sub-item 1
      \item Sub-item 2
    \end{itemize}
\end{itemize}

\subsection{Numbered List}
\begin{enumerate}
  \item First step
  \item Second step
  \item Third step
\end{enumerate}

\section{Citations and References}
When you use the citation search feature, references will be automatically managed. For example, you might cite a paper like this: \cite{example2024}. All citations will be automatically added to your \texttt{cite.bib} file.

\section{Conclusion}
This template demonstrates basic \LaTeX{} features including mathematical formulas, tables, and lists. The bibliography system is now set up to work seamlessly with Bibby AI's citation management. Feel free to modify it for your needs!

\bibliographystyle{plain}
\bibliography{cite}

\end{document}\documentclass{article}
\usepackage[utf8]{inputenc}
\usepackage{booktabs}
\usepackage{amsmath}
\usepackage{natbib}

\title{My Document}
\author{Your Name}
\date{\today}

\begin{document}

\maketitle

\section{Introduction}
This is a sample document demonstrating various \LaTeX{} features including mathematical equations, tables, and formatting. You can cite papers using the citation search feature, and they will be automatically added to your \texttt{cite.bib} file.

\section{Mathematical Expressions}
Here are some example equations:

\subsection{Quadratic Formula}
The quadratic formula is given by:
\[x = \frac{-b \pm \sqrt{b^2 - 4ac}}{2a}\]

\subsection{Matrix Example}
A sample matrix:
\[
\begin{pmatrix}
a_{11} & a_{12} & a_{13} \\
a_{21} & a_{22} & a_{23} \\
a_{31} & a_{32} & a_{33}
\end{pmatrix}
\]

\section{Tables}
Below is a sample table showing data:

\begin{table}[h]
\centering
\begin{tabular}{@{}lrrr@{}}
\toprule
\textbf{Category} & \textbf{Value 1} & \textbf{Value 2} & \textbf{Total} \\
\midrule
Item A & 10.5 & 20.3 & 30.8 \\
Item B & 15.2 & 25.1 & 40.3 \\
Item C & 22.1 & 12.7 & 34.8 \\
\bottomrule
\end{tabular}
\caption{Sample Data Table}
\label{tab:sample}
\end{table}

\section{Lists}
\subsection{Bullet Points}
\begin{itemize}
  \item First item
  \item Second item
  \item Third item with sub-items:
    \begin{itemize}
      \item Sub-item 1
      \item Sub-item 2
    \end{itemize}
\end{itemize}

\subsection{Numbered List}
\begin{enumerate}
  \item First step
  \item Second step
  \item Third step
\end{enumerate}

\section{Citations and References}
When you use the citation search feature, references will be automatically managed. For example, you might cite a paper like this: \cite{example2024}. All citations will be automatically added to your \texttt{cite.bib} file.

\section{Conclusion}
This template demonstrates basic \LaTeX{} features including mathematical formulas, tables, and lists. The bibliography system is now set up to work seamlessly with Bibby AI's citation management. Feel free to modify it for your needs!

\bibliographystyle{plain}
\bibliography{cite}

\end{document}\documentclass{article}
\usepackage[utf8]{inputenc}
\usepackage{booktabs}
\usepackage{amsmath}
\usepackage{natbib}

\title{My Document}
\author{Your Name}
\date{\today}

\begin{document}

\maketitle

\section{Introduction}
This is a sample document demonstrating various \LaTeX{} features including mathematical equations, tables, and formatting. You can cite papers using the citation search feature, and they will be automatically added to your \texttt{cite.bib} file.

\section{Mathematical Expressions}
Here are some example equations:

\subsection{Quadratic Formula}
The quadratic formula is given by:
\[x = \frac{-b \pm \sqrt{b^2 - 4ac}}{2a}\]

\subsection{Matrix Example}
A sample matrix:
\[
\begin{pmatrix}
a_{11} & a_{12} & a_{13} \\
a_{21} & a_{22} & a_{23} \\
a_{31} & a_{32} & a_{33}
\end{pmatrix}
\]

\section{Tables}
Below is a sample table showing data:

\begin{table}[h]
\centering
\begin{tabular}{@{}lrrr@{}}
\toprule
\textbf{Category} & \textbf{Value 1} & \textbf{Value 2} & \textbf{Total} \\
\midrule
Item A & 10.5 & 20.3 & 30.8 \\
Item B & 15.2 & 25.1 & 40.3 \\
Item C & 22.1 & 12.7 & 34.8 \\
\bottomrule
\end{tabular}
\caption{Sample Data Table}
\label{tab:sample}
\end{table}

\section{Lists}
\subsection{Bullet Points}
\begin{itemize}
  \item First item
  \item Second item
  \item Third item with sub-items:
    \begin{itemize}
      \item Sub-item 1
      \item Sub-item 2
    \end{itemize}
\end{itemize}

\subsection{Numbered List}
\begin{enumerate}
  \item First step
  \item Second step
  \item Third step
\end{enumerate}

\section{Citations and References}
When you use the citation search feature, references will be automatically managed. For example, you might cite a paper like this: \cite{example2024}. All citations will be automatically added to your \texttt{cite.bib} file.

\section{Conclusion}
This template demonstrates basic \LaTeX{} features including mathematical formulas, tables, and lists. The bibliography system is now set up to work seamlessly with Bibby AI's citation management. Feel free to modify it for your needs!

\bibliographystyle{plain}
\bibliography{cite}

\end{document}\documentclass{article}
\usepackage[utf8]{inputenc}
\usepackage{booktabs}
\usepackage{amsmath}
\usepackage{natbib}

\title{My Document}
\author{Your Name}
\date{\today}

\begin{document}

\maketitle

\section{Introduction}
This is a sample document demonstrating various \LaTeX{} features including mathematical equations, tables, and formatting. You can cite papers using the citation search feature, and they will be automatically added to your \texttt{cite.bib} file.

\section{Mathematical Expressions}
Here are some example equations:

\subsection{Quadratic Formula}
The quadratic formula is given by:
\[x = \frac{-b \pm \sqrt{b^2 - 4ac}}{2a}\]

\subsection{Matrix Example}
A sample matrix:
\[
\begin{pmatrix}
a_{11} & a_{12} & a_{13} \\
a_{21} & a_{22} & a_{23} \\
a_{31} & a_{32} & a_{33}
\end{pmatrix}
\]

\section{Tables}
Below is a sample table showing data:

\begin{table}[h]
\centering
\begin{tabular}{@{}lrrr@{}}
\toprule
\textbf{Category} & \textbf{Value 1} & \textbf{Value 2} & \textbf{Total} \\
\midrule
Item A & 10.5 & 20.3 & 30.8 \\
Item B & 15.2 & 25.1 & 40.3 \\
Item C & 22.1 & 12.7 & 34.8 \\
\bottomrule
\end{tabular}
\caption{Sample Data Table}
\label{tab:sample}
\end{table}

\section{Lists}
\subsection{Bullet Points}
\begin{itemize}
  \item First item
  \item Second item
  \item Third item with sub-items:
    \begin{itemize}
      \item Sub-item 1
      \item Sub-item 2
    \end{itemize}
\end{itemize}

\subsection{Numbered List}
\begin{enumerate}
  \item First step
  \item Second step
  \item Third step
\end{enumerate}

\section{Citations and References}
When you use the citation search feature, references will be automatically managed. For example, you might cite a paper like this: \cite{example2024}. All citations will be automatically added to your \texttt{cite.bib} file.

\section{Conclusion}
This template demonstrates basic \LaTeX{} features including mathematical formulas, tables, and lists. The bibliography system is now set up to work seamlessly with Bibby AI's citation management. Feel free to modify it for your needs!

\bibliographystyle{plain}
\bibliography{cite}

\end{document}\documentclass{article}
\usepackage[utf8]{inputenc}
\usepackage{booktabs}
\usepackage{amsmath}
\usepackage{natbib}

\title{My Document}
\author{Your Name}
\date{\today}

\begin{document}

\maketitle

\section{Introduction}
This is a sample document demonstrating various \LaTeX{} features including mathematical equations, tables, and formatting. You can cite papers using the citation search feature, and they will be automatically added to your \texttt{cite.bib} file.

\section{Mathematical Expressions}
Here are some example equations:

\subsection{Quadratic Formula}
The quadratic formula is given by:
\[x = \frac{-b \pm \sqrt{b^2 - 4ac}}{2a}\]

\subsection{Matrix Example}
A sample matrix:
\[
\begin{pmatrix}
a_{11} & a_{12} & a_{13} \\
a_{21} & a_{22} & a_{23} \\
a_{31} & a_{32} & a_{33}
\end{pmatrix}
\]

\section{Tables}
Below is a sample table showing data:

\begin{table}[h]
\centering
\begin{tabular}{@{}lrrr@{}}
\toprule
\textbf{Category} & \textbf{Value 1} & \textbf{Value 2} & \textbf{Total} \\
\midrule
Item A & 10.5 & 20.3 & 30.8 \\
Item B & 15.2 & 25.1 & 40.3 \\
Item C & 22.1 & 12.7 & 34.8 \\
\bottomrule
\end{tabular}
\caption{Sample Data Table}
\label{tab:sample}
\end{table}

\section{Lists}
\subsection{Bullet Points}
\begin{itemize}
  \item First item
  \item Second item
  \item Third item with sub-items:
    \begin{itemize}
      \item Sub-item 1
      \item Sub-item 2
    \end{itemize}
\end{itemize}

\subsection{Numbered List}
\begin{enumerate}
  \item First step
  \item Second step
  \item Third step
\end{enumerate}

\section{Citations and References}
When you use the citation search feature, references will be automatically managed. For example, you might cite a paper like this: \cite{example2024}. All citations will be automatically added to your \texttt{cite.bib} file.

\section{Conclusion}
This template demonstrates basic \LaTeX{} features including mathematical formulas, tables, and lists. The bibliography system is now set up to work seamlessly with Bibby AI's citation management. Feel free to modify it for your needs!

\bibliographystyle{plain}
\bibliography{cite}

\end{document}\documentclass{article}
\usepackage[utf8]{inputenc}
\usepackage{booktabs}
\usepackage{amsmath}
\usepackage{natbib}

\title{My Document}
\author{Your Name}
\date{\today}

\begin{document}

\maketitle

\section{Introduction}
This is a sample document demonstrating various \LaTeX{} features including mathematical equations, tables, and formatting. You can cite papers using the citation search feature, and they will be automatically added to your \texttt{cite.bib} file.

\section{Mathematical Expressions}
Here are some example equations:

\subsection{Quadratic Formula}
The quadratic formula is given by:
\[x = \frac{-b \pm \sqrt{b^2 - 4ac}}{2a}\]

\subsection{Matrix Example}
A sample matrix:
\[
\begin{pmatrix}
a_{11} & a_{12} & a_{13} \\
a_{21} & a_{22} & a_{23} \\
a_{31} & a_{32} & a_{33}
\end{pmatrix}
\]

\section{Tables}
Below is a sample table showing data:

\begin{table}[h]
\centering
\begin{tabular}{@{}lrrr@{}}
\toprule
\textbf{Category} & \textbf{Value 1} & \textbf{Value 2} & \textbf{Total} \\
\midrule
Item A & 10.5 & 20.3 & 30.8 \\
Item B & 15.2 & 25.1 & 40.3 \\
Item C & 22.1 & 12.7 & 34.8 \\
\bottomrule
\end{tabular}
\caption{Sample Data Table}
\label{tab:sample}
\end{table}

\section{Lists}
\subsection{Bullet Points}
\begin{itemize}
  \item First item
  \item Second item
  \item Third item with sub-items:
    \begin{itemize}
      \item Sub-item 1
      \item Sub-item 2
    \end{itemize}
\end{itemize}

\subsection{Numbered List}
\begin{enumerate}
  \item First step
  \item Second step
  \item Third step
\end{enumerate}

\section{Citations and References}
When you use the citation search feature, references will be automatically managed. For example, you might cite a paper like this: \cite{example2024}. All citations will be automatically added to your \texttt{cite.bib} file.

\section{Conclusion}
This template demonstrates basic \LaTeX{} features including mathematical formulas, tables, and lists. The bibliography system is now set up to work seamlessly with Bibby AI's citation management. Feel free to modify it for your needs!

\bibliographystyle{plain}
\bibliography{cite}

\end{document}\documentclass{article}
\usepackage[utf8]{inputenc}
\usepackage{booktabs}
\usepackage{amsmath}
\usepackage{natbib}

\title{My Document}
\author{Your Name}
\date{\today}

\begin{document}

\maketitle

\section{Introduction}
This is a sample document demonstrating various \LaTeX{} features including mathematical equations, tables, and formatting. You can cite papers using the citation search feature, and they will be automatically added to your \texttt{cite.bib} file.

\section{Mathematical Expressions}
Here are some example equations:

\subsection{Quadratic Formula}
The quadratic formula is given by:
\[x = \frac{-b \pm \sqrt{b^2 - 4ac}}{2a}\]

\subsection{Matrix Example}
A sample matrix:
\[
\begin{pmatrix}
a_{11} & a_{12} & a_{13} \\
a_{21} & a_{22} & a_{23} \\
a_{31} & a_{32} & a_{33}
\end{pmatrix}
\]

\section{Tables}
Below is a sample table showing data:

\begin{table}[h]
\centering
\begin{tabular}{@{}lrrr@{}}
\toprule
\textbf{Category} & \textbf{Value 1} & \textbf{Value 2} & \textbf{Total} \\
\midrule
Item A & 10.5 & 20.3 & 30.8 \\
Item B & 15.2 & 25.1 & 40.3 \\
Item C & 22.1 & 12.7 & 34.8 \\
\bottomrule
\end{tabular}
\caption{Sample Data Table}
\label{tab:sample}
\end{table}

\section{Lists}
\subsection{Bullet Points}
\begin{itemize}
  \item First item
  \item Second item
  \item Third item with sub-items:
    \begin{itemize}
      \item Sub-item 1
      \item Sub-item 2
    \end{itemize}
\end{itemize}

\subsection{Numbered List}
\begin{enumerate}
  \item First step
  \item Second step
  \item Third step
\end{enumerate}

\section{Citations and References}
When you use the citation search feature, references will be automatically managed. For example, you might cite a paper like this: \cite{example2024}. All citations will be automatically added to your \texttt{cite.bib} file.

\section{Conclusion}
This template demonstrates basic \LaTeX{} features including mathematical formulas, tables, and lists. The bibliography system is now set up to work seamlessly with Bibby AI's citation management. Feel free to modify it for your needs!

\bibliographystyle{plain}
\bibliography{cite}

\end{document}\documentclass{article}
\usepackage[utf8]{inputenc}
\usepackage{booktabs}
\usepackage{amsmath}
\usepackage{natbib}

\title{My Document}
\author{Your Name}
\date{\today}

\begin{document}

\maketitle

\section{Introduction}
This is a sample document demonstrating various \LaTeX{} features including mathematical equations, tables, and formatting. You can cite papers using the citation search feature, and they will be automatically added to your \texttt{cite.bib} file.

\section{Mathematical Expressions}
Here are some example equations:

\subsection{Quadratic Formula}
The quadratic formula is given by:
\[x = \frac{-b \pm \sqrt{b^2 - 4ac}}{2a}\]

\subsection{Matrix Example}
A sample matrix:
\[
\begin{pmatrix}
a_{11} & a_{12} & a_{13} \\
a_{21} & a_{22} & a_{23} \\
a_{31} & a_{32} & a_{33}
\end{pmatrix}
\]

\section{Tables}
Below is a sample table showing data:

\begin{table}[h]
\centering
\begin{tabular}{@{}lrrr@{}}
\toprule
\textbf{Category} & \textbf{Value 1} & \textbf{Value 2} & \textbf{Total} \\
\midrule
Item A & 10.5 & 20.3 & 30.8 \\
Item B & 15.2 & 25.1 & 40.3 \\
Item C & 22.1 & 12.7 & 34.8 \\
\bottomrule
\end{tabular}
\caption{Sample Data Table}
\label{tab:sample}
\end{table}

\section{Lists}
\subsection{Bullet Points}
\begin{itemize}
  \item First item
  \item Second item
  \item Third item with sub-items:
    \begin{itemize}
      \item Sub-item 1
      \item Sub-item 2
    \end{itemize}
\end{itemize}

\subsection{Numbered List}
\begin{enumerate}
  \item First step
  \item Second step
  \item Third step
\end{enumerate}

\section{Citations and References}
When you use the citation search feature, references will be automatically managed. For example, you might cite a paper like this: \cite{example2024}. All citations will be automatically added to your \texttt{cite.bib} file.

\section{Conclusion}
This template demonstrates basic \LaTeX{} features including mathematical formulas, tables, and lists. The bibliography system is now set up to work seamlessly with Bibby AI's citation management. Feel free to modify it for your needs!

\bibliographystyle{plain}
\bibliography{cite}

\end{document}\documentclass{article}
\usepackage[utf8]{inputenc}
\usepackage{booktabs}
\usepackage{amsmath}
\usepackage{natbib}

\title{My Document}
\author{Your Name}
\date{\today}

\begin{document}

\maketitle

\section{Introduction}
This is a sample document demonstrating various \LaTeX{} features including mathematical equations, tables, and formatting. You can cite papers using the citation search feature, and they will be automatically added to your \texttt{cite.bib} file.

\section{Mathematical Expressions}
Here are some example equations:

\subsection{Quadratic Formula}
The quadratic formula is given by:
\[x = \frac{-b \pm \sqrt{b^2 - 4ac}}{2a}\]

\subsection{Matrix Example}
A sample matrix:
\[
\begin{pmatrix}
a_{11} & a_{12} & a_{13} \\
a_{21} & a_{22} & a_{23} \\
a_{31} & a_{32} & a_{33}
\end{pmatrix}
\]

\section{Tables}
Below is a sample table showing data:

\begin{table}[h]
\centering
\begin{tabular}{@{}lrrr@{}}
\toprule
\textbf{Category} & \textbf{Value 1} & \textbf{Value 2} & \textbf{Total} \\
\midrule
Item A & 10.5 & 20.3 & 30.8 \\
Item B & 15.2 & 25.1 & 40.3 \\
Item C & 22.1 & 12.7 & 34.8 \\
\bottomrule
\end{tabular}
\caption{Sample Data Table}
\label{tab:sample}
\end{table}

\section{Lists}
\subsection{Bullet Points}
\begin{itemize}
  \item First item
  \item Second item
  \item Third item with sub-items:
    \begin{itemize}
      \item Sub-item 1
      \item Sub-item 2
    \end{itemize}
\end{itemize}

\subsection{Numbered List}
\begin{enumerate}
  \item First step
  \item Second step
  \item Third step
\end{enumerate}

\section{Citations and References}
When you use the citation search feature, references will be automatically managed. For example, you might cite a paper like this: \cite{example2024}. All citations will be automatically added to your \texttt{cite.bib} file.

\section{Conclusion}
This template demonstrates basic \LaTeX{} features including mathematical formulas, tables, and lists. The bibliography system is now set up to work seamlessly with Bibby AI's citation management. Feel free to modify it for your needs!

\bibliographystyle{plain}
\bibliography{cite}

\end{document}\documentclass{article}
\usepackage[utf8]{inputenc}
\usepackage{booktabs}
\usepackage{amsmath}
\usepackage{natbib}

\title{My Document}
\author{Your Name}
\date{\today}

\begin{document}

\maketitle

\section{Introduction}
This is a sample document demonstrating various \LaTeX{} features including mathematical equations, tables, and formatting. You can cite papers using the citation search feature, and they will be automatically added to your \texttt{cite.bib} file.

\section{Mathematical Expressions}
Here are some example equations:

\subsection{Quadratic Formula}
The quadratic formula is given by:
\[x = \frac{-b \pm \sqrt{b^2 - 4ac}}{2a}\]

\subsection{Matrix Example}
A sample matrix:
\[
\begin{pmatrix}
a_{11} & a_{12} & a_{13} \\
a_{21} & a_{22} & a_{23} \\
a_{31} & a_{32} & a_{33}
\end{pmatrix}
\]

\section{Tables}
Below is a sample table showing data:

\begin{table}[h]
\centering
\begin{tabular}{@{}lrrr@{}}
\toprule
\textbf{Category} & \textbf{Value 1} & \textbf{Value 2} & \textbf{Total} \\
\midrule
Item A & 10.5 & 20.3 & 30.8 \\
Item B & 15.2 & 25.1 & 40.3 \\
Item C & 22.1 & 12.7 & 34.8 \\
\bottomrule
\end{tabular}
\caption{Sample Data Table}
\label{tab:sample}
\end{table}

\section{Lists}
\subsection{Bullet Points}
\begin{itemize}
  \item First item
  \item Second item
  \item Third item with sub-items:
    \begin{itemize}
      \item Sub-item 1
      \item Sub-item 2
    \end{itemize}
\end{itemize}

\subsection{Numbered List}
\begin{enumerate}
  \item First step
  \item Second step
  \item Third step
\end{enumerate}

\section{Citations and References}
When you use the citation search feature, references will be automatically managed. For example, you might cite a paper like this: \cite{example2024}. All citations will be automatically added to your \texttt{cite.bib} file.

\section{Conclusion}
This template demonstrates basic \LaTeX{} features including mathematical formulas, tables, and lists. The bibliography system is now set up to work seamlessly with Bibby AI's citation management. Feel free to modify it for your needs!

\bibliographystyle{plain}
\bibliography{cite}

\end{document}\documentclass{article}
\usepackage[utf8]{inputenc}
\usepackage{booktabs}
\usepackage{amsmath}
\usepackage{natbib}

\title{My Document}
\author{Your Name}
\date{\today}

\begin{document}

\maketitle

\section{Introduction}
This is a sample document demonstrating various \LaTeX{} features including mathematical equations, tables, and formatting. You can cite papers using the citation search feature, and they will be automatically added to your \texttt{cite.bib} file.

\section{Mathematical Expressions}
Here are some example equations:

\subsection{Quadratic Formula}
The quadratic formula is given by:
\[x = \frac{-b \pm \sqrt{b^2 - 4ac}}{2a}\]

\subsection{Matrix Example}
A sample matrix:
\[
\begin{pmatrix}
a_{11} & a_{12} & a_{13} \\
a_{21} & a_{22} & a_{23} \\
a_{31} & a_{32} & a_{33}
\end{pmatrix}
\]

\section{Tables}
Below is a sample table showing data:

\begin{table}[h]
\centering
\begin{tabular}{@{}lrrr@{}}
\toprule
\textbf{Category} & \textbf{Value 1} & \textbf{Value 2} & \textbf{Total} \\
\midrule
Item A & 10.5 & 20.3 & 30.8 \\
Item B & 15.2 & 25.1 & 40.3 \\
Item C & 22.1 & 12.7 & 34.8 \\
\bottomrule
\end{tabular}
\caption{Sample Data Table}
\label{tab:sample}
\end{table}

\section{Lists}
\subsection{Bullet Points}
\begin{itemize}
  \item First item
  \item Second item
  \item Third item with sub-items:
    \begin{itemize}
      \item Sub-item 1
      \item Sub-item 2
    \end{itemize}
\end{itemize}

\subsection{Numbered List}
\begin{enumerate}
  \item First step
  \item Second step
  \item Third step
\end{enumerate}

\section{Citations and References}
When you use the citation search feature, references will be automatically managed. For example, you might cite a paper like this: \cite{example2024}. All citations will be automatically added to your \texttt{cite.bib} file.

\section{Conclusion}
This template demonstrates basic \LaTeX{} features including mathematical formulas, tables, and lists. The bibliography system is now set up to work seamlessly with Bibby AI's citation management. Feel free to modify it for your needs!

\bibliographystyle{plain}
\bibliography{cite}

\end{document}\documentclass{article}
\usepackage[utf8]{inputenc}
\usepackage{booktabs}
\usepackage{amsmath}
\usepackage{natbib}

\title{My Document}
\author{Your Name}
\date{\today}

\begin{document}

\maketitle

\section{Introduction}
This is a sample document demonstrating various \LaTeX{} features including mathematical equations, tables, and formatting. You can cite papers using the citation search feature, and they will be automatically added to your \texttt{cite.bib} file.

\section{Mathematical Expressions}
Here are some example equations:

\subsection{Quadratic Formula}
The quadratic formula is given by:
\[x = \frac{-b \pm \sqrt{b^2 - 4ac}}{2a}\]

\subsection{Matrix Example}
A sample matrix:
\[
\begin{pmatrix}
a_{11} & a_{12} & a_{13} \\
a_{21} & a_{22} & a_{23} \\
a_{31} & a_{32} & a_{33}
\end{pmatrix}
\]

\section{Tables}
Below is a sample table showing data:

\begin{table}[h]
\centering
\begin{tabular}{@{}lrrr@{}}
\toprule
\textbf{Category} & \textbf{Value 1} & \textbf{Value 2} & \textbf{Total} \\
\midrule
Item A & 10.5 & 20.3 & 30.8 \\
Item B & 15.2 & 25.1 & 40.3 \\
Item C & 22.1 & 12.7 & 34.8 \\
\bottomrule
\end{tabular}
\caption{Sample Data Table}
\label{tab:sample}
\end{table}

\section{Lists}
\subsection{Bullet Points}
\begin{itemize}
  \item First item
  \item Second item
  \item Third item with sub-items:
    \begin{itemize}
      \item Sub-item 1
      \item Sub-item 2
    \end{itemize}
\end{itemize}

\subsection{Numbered List}
\begin{enumerate}
  \item First step
  \item Second step
  \item Third step
\end{enumerate}

\section{Citations and References}
When you use the citation search feature, references will be automatically managed. For example, you might cite a paper like this: \cite{example2024}. All citations will be automatically added to your \texttt{cite.bib} file.

\section{Conclusion}
This template demonstrates basic \LaTeX{} features including mathematical formulas, tables, and lists. The bibliography system is now set up to work seamlessly with Bibby AI's citation management. Feel free to modify it for your needs!

\bibliographystyle{plain}
\bibliography{cite}

\end{document}\documentclass{article}
\usepackage[utf8]{inputenc}
\usepackage{booktabs}
\usepackage{amsmath}
\usepackage{natbib}

\title{My Document}
\author{Your Name}
\date{\today}

\begin{document}

\maketitle

\section{Introduction}
This is a sample document demonstrating various \LaTeX{} features including mathematical equations, tables, and formatting. You can cite papers using the citation search feature, and they will be automatically added to your \texttt{cite.bib} file.

\section{Mathematical Expressions}
Here are some example equations:

\subsection{Quadratic Formula}
The quadratic formula is given by:
\[x = \frac{-b \pm \sqrt{b^2 - 4ac}}{2a}\]

\subsection{Matrix Example}
A sample matrix:
\[
\begin{pmatrix}
a_{11} & a_{12} & a_{13} \\
a_{21} & a_{22} & a_{23} \\
a_{31} & a_{32} & a_{33}
\end{pmatrix}
\]

\section{Tables}
Below is a sample table showing data:

\begin{table}[h]
\centering
\begin{tabular}{@{}lrrr@{}}
\toprule
\textbf{Category} & \textbf{Value 1} & \textbf{Value 2} & \textbf{Total} \\
\midrule
Item A & 10.5 & 20.3 & 30.8 \\
Item B & 15.2 & 25.1 & 40.3 \\
Item C & 22.1 & 12.7 & 34.8 \\
\bottomrule
\end{tabular}
\caption{Sample Data Table}
\label{tab:sample}
\end{table}

\section{Lists}
\subsection{Bullet Points}
\begin{itemize}
  \item First item
  \item Second item
  \item Third item with sub-items:
    \begin{itemize}
      \item Sub-item 1
      \item Sub-item 2
    \end{itemize}
\end{itemize}

\subsection{Numbered List}
\begin{enumerate}
  \item First step
  \item Second step
  \item Third step
\end{enumerate}

\section{Citations and References}
When you use the citation search feature, references will be automatically managed. For example, you might cite a paper like this: \cite{example2024}. All citations will be automatically added to your \texttt{cite.bib} file.

\section{Conclusion}
This template demonstrates basic \LaTeX{} features including mathematical formulas, tables, and lists. The bibliography system is now set up to work seamlessly with Bibby AI's citation management. Feel free to modify it for your needs!

\bibliographystyle{plain}
\bibliography{cite}

\end{document}\documentclass{article}
\usepackage[utf8]{inputenc}
\usepackage{booktabs}
\usepackage{amsmath}
\usepackage{natbib}

\title{My Document}
\author{Your Name}
\date{\today}

\begin{document}

\maketitle

\section{Introduction}
This is a sample document demonstrating various \LaTeX{} features including mathematical equations, tables, and formatting. You can cite papers using the citation search feature, and they will be automatically added to your \texttt{cite.bib} file.

\section{Mathematical Expressions}
Here are some example equations:

\subsection{Quadratic Formula}
The quadratic formula is given by:
\[x = \frac{-b \pm \sqrt{b^2 - 4ac}}{2a}\]

\subsection{Matrix Example}
A sample matrix:
\[
\begin{pmatrix}
a_{11} & a_{12} & a_{13} \\
a_{21} & a_{22} & a_{23} \\
a_{31} & a_{32} & a_{33}
\end{pmatrix}
\]

\section{Tables}
Below is a sample table showing data:

\begin{table}[h]
\centering
\begin{tabular}{@{}lrrr@{}}
\toprule
\textbf{Category} & \textbf{Value 1} & \textbf{Value 2} & \textbf{Total} \\
\midrule
Item A & 10.5 & 20.3 & 30.8 \\
Item B & 15.2 & 25.1 & 40.3 \\
Item C & 22.1 & 12.7 & 34.8 \\
\bottomrule
\end{tabular}
\caption{Sample Data Table}
\label{tab:sample}
\end{table}

\section{Lists}
\subsection{Bullet Points}
\begin{itemize}
  \item First item
  \item Second item
  \item Third item with sub-items:
    \begin{itemize}
      \item Sub-item 1
      \item Sub-item 2
    \end{itemize}
\end{itemize}

\subsection{Numbered List}
\begin{enumerate}
  \item First step
  \item Second step
  \item Third step
\end{enumerate}

\section{Citations and References}
When you use the citation search feature, references will be automatically managed. For example, you might cite a paper like this: \cite{example2024}. All citations will be automatically added to your \texttt{cite.bib} file.

\section{Conclusion}
This template demonstrates basic \LaTeX{} features including mathematical formulas, tables, and lists. The bibliography system is now set up to work seamlessly with Bibby AI's citation management. Feel free to modify it for your needs!

\bibliographystyle{plain}
\bibliography{cite}

\end{document}\documentclass{article}
\usepackage[utf8]{inputenc}
\usepackage{booktabs}
\usepackage{amsmath}
\usepackage{natbib}

\title{My Document}
\author{Your Name}
\date{\today}

\begin{document}

\maketitle

\section{Introduction}
This is a sample document demonstrating various \LaTeX{} features including mathematical equations, tables, and formatting. You can cite papers using the citation search feature, and they will be automatically added to your \texttt{cite.bib} file.

\section{Mathematical Expressions}
Here are some example equations:

\subsection{Quadratic Formula}
The quadratic formula is given by:
\[x = \frac{-b \pm \sqrt{b^2 - 4ac}}{2a}\]

\subsection{Matrix Example}
A sample matrix:
\[
\begin{pmatrix}
a_{11} & a_{12} & a_{13} \\
a_{21} & a_{22} & a_{23} \\
a_{31} & a_{32} & a_{33}
\end{pmatrix}
\]

\section{Tables}
Below is a sample table showing data:

\begin{table}[h]
\centering
\begin{tabular}{@{}lrrr@{}}
\toprule
\textbf{Category} & \textbf{Value 1} & \textbf{Value 2} & \textbf{Total} \\
\midrule
Item A & 10.5 & 20.3 & 30.8 \\
Item B & 15.2 & 25.1 & 40.3 \\
Item C & 22.1 & 12.7 & 34.8 \\
\bottomrule
\end{tabular}
\caption{Sample Data Table}
\label{tab:sample}
\end{table}

\section{Lists}
\subsection{Bullet Points}
\begin{itemize}
  \item First item
  \item Second item
  \item Third item with sub-items:
    \begin{itemize}
      \item Sub-item 1
      \item Sub-item 2
    \end{itemize}
\end{itemize}

\subsection{Numbered List}
\begin{enumerate}
  \item First step
  \item Second step
  \item Third step
\end{enumerate}

\section{Citations and References}
When you use the citation search feature, references will be automatically managed. For example, you might cite a paper like this: \cite{example2024}. All citations will be automatically added to your \texttt{cite.bib} file.

\section{Conclusion}
This template demonstrates basic \LaTeX{} features including mathematical formulas, tables, and lists. The bibliography system is now set up to work seamlessly with Bibby AI's citation management. Feel free to modify it for your needs!

\bibliographystyle{plain}
\bibliography{cite}

\end{document}\documentclass{article}
\usepackage[utf8]{inputenc}
\usepackage{booktabs}
\usepackage{amsmath}
\usepackage{natbib}

\title{My Document}
\author{Your Name}
\date{\today}

\begin{document}

\maketitle

\section{Introduction}
This is a sample document demonstrating various \LaTeX{} features including mathematical equations, tables, and formatting. You can cite papers using the citation search feature, and they will be automatically added to your \texttt{cite.bib} file.

\section{Mathematical Expressions}
Here are some example equations:

\subsection{Quadratic Formula}
The quadratic formula is given by:
\[x = \frac{-b \pm \sqrt{b^2 - 4ac}}{2a}\]

\subsection{Matrix Example}
A sample matrix:
\[
\begin{pmatrix}
a_{11} & a_{12} & a_{13} \\
a_{21} & a_{22} & a_{23} \\
a_{31} & a_{32} & a_{33}
\end{pmatrix}
\]

\section{Tables}
Below is a sample table showing data:

\begin{table}[h]
\centering
\begin{tabular}{@{}lrrr@{}}
\toprule
\textbf{Category} & \textbf{Value 1} & \textbf{Value 2} & \textbf{Total} \\
\midrule
Item A & 10.5 & 20.3 & 30.8 \\
Item B & 15.2 & 25.1 & 40.3 \\
Item C & 22.1 & 12.7 & 34.8 \\
\bottomrule
\end{tabular}
\caption{Sample Data Table}
\label{tab:sample}
\end{table}

\section{Lists}
\subsection{Bullet Points}
\begin{itemize}
  \item First item
  \item Second item
  \item Third item with sub-items:
    \begin{itemize}
      \item Sub-item 1
      \item Sub-item 2
    \end{itemize}
\end{itemize}

\subsection{Numbered List}
\begin{enumerate}
  \item First step
  \item Second step
  \item Third step
\end{enumerate}

\section{Citations and References}
When you use the citation search feature, references will be automatically managed. For example, you might cite a paper like this: \cite{example2024}. All citations will be automatically added to your \texttt{cite.bib} file.

\section{Conclusion}
This template demonstrates basic \LaTeX{} features including mathematical formulas, tables, and lists. The bibliography system is now set up to work seamlessly with Bibby AI's citation management. Feel free to modify it for your needs!

\bibliographystyle{plain}
\bibliography{cite}

\end{document}\documentclass{article}
\usepackage[utf8]{inputenc}
\usepackage{booktabs}
\usepackage{amsmath}
\usepackage{natbib}

\title{My Document}
\author{Your Name}
\date{\today}

\begin{document}

\maketitle

\section{Introduction}
This is a sample document demonstrating various \LaTeX{} features including mathematical equations, tables, and formatting. You can cite papers using the citation search feature, and they will be automatically added to your \texttt{cite.bib} file.

\section{Mathematical Expressions}
Here are some example equations:

\subsection{Quadratic Formula}
The quadratic formula is given by:
\[x = \frac{-b \pm \sqrt{b^2 - 4ac}}{2a}\]

\subsection{Matrix Example}
A sample matrix:
\[
\begin{pmatrix}
a_{11} & a_{12} & a_{13} \\
a_{21} & a_{22} & a_{23} \\
a_{31} & a_{32} & a_{33}
\end{pmatrix}
\]

\section{Tables}
Below is a sample table showing data:

\begin{table}[h]
\centering
\begin{tabular}{@{}lrrr@{}}
\toprule
\textbf{Category} & \textbf{Value 1} & \textbf{Value 2} & \textbf{Total} \\
\midrule
Item A & 10.5 & 20.3 & 30.8 \\
Item B & 15.2 & 25.1 & 40.3 \\
Item C & 22.1 & 12.7 & 34.8 \\
\bottomrule
\end{tabular}
\caption{Sample Data Table}
\label{tab:sample}
\end{table}

\section{Lists}
\subsection{Bullet Points}
\begin{itemize}
  \item First item
  \item Second item
  \item Third item with sub-items:
    \begin{itemize}
      \item Sub-item 1
      \item Sub-item 2
    \end{itemize}
\end{itemize}

\subsection{Numbered List}
\begin{enumerate}
  \item First step
  \item Second step
  \item Third step
\end{enumerate}

\section{Citations and References}
When you use the citation search feature, references will be automatically managed. For example, you might cite a paper like this: \cite{example2024}. All citations will be automatically added to your \texttt{cite.bib} file.

\section{Conclusion}
This template demonstrates basic \LaTeX{} features including mathematical formulas, tables, and lists. The bibliography system is now set up to work seamlessly with Bibby AI's citation management. Feel free to modify it for your needs!

\bibliographystyle{plain}
\bibliography{cite}

\end{document}\documentclass{article}
\usepackage[utf8]{inputenc}
\usepackage{booktabs}
\usepackage{amsmath}
\usepackage{natbib}

\title{My Document}
\author{Your Name}
\date{\today}

\begin{document}

\maketitle

\section{Introduction}
This is a sample document demonstrating various \LaTeX{} features including mathematical equations, tables, and formatting. You can cite papers using the citation search feature, and they will be automatically added to your \texttt{cite.bib} file.

\section{Mathematical Expressions}
Here are some example equations:

\subsection{Quadratic Formula}
The quadratic formula is given by:
\[x = \frac{-b \pm \sqrt{b^2 - 4ac}}{2a}\]

\subsection{Matrix Example}
A sample matrix:
\[
\begin{pmatrix}
a_{11} & a_{12} & a_{13} \\
a_{21} & a_{22} & a_{23} \\
a_{31} & a_{32} & a_{33}
\end{pmatrix}
\]

\section{Tables}
Below is a sample table showing data:

\begin{table}[h]
\centering
\begin{tabular}{@{}lrrr@{}}
\toprule
\textbf{Category} & \textbf{Value 1} & \textbf{Value 2} & \textbf{Total} \\
\midrule
Item A & 10.5 & 20.3 & 30.8 \\
Item B & 15.2 & 25.1 & 40.3 \\
Item C & 22.1 & 12.7 & 34.8 \\
\bottomrule
\end{tabular}
\caption{Sample Data Table}
\label{tab:sample}
\end{table}

\section{Lists}
\subsection{Bullet Points}
\begin{itemize}
  \item First item
  \item Second item
  \item Third item with sub-items:
    \begin{itemize}
      \item Sub-item 1
      \item Sub-item 2
    \end{itemize}
\end{itemize}

\subsection{Numbered List}
\begin{enumerate}
  \item First step
  \item Second step
  \item Third step
\end{enumerate}

\section{Citations and References}
When you use the citation search feature, references will be automatically managed. For example, you might cite a paper like this: \cite{example2024}. All citations will be automatically added to your \texttt{cite.bib} file.

\section{Conclusion}
This template demonstrates basic \LaTeX{} features including mathematical formulas, tables, and lists. The bibliography system is now set up to work seamlessly with Bibby AI's citation management. Feel free to modify it for your needs!

\bibliographystyle{plain}
\bibliography{cite}

\end{document}\documentclass{article}
\usepackage[utf8]{inputenc}
\usepackage{booktabs}
\usepackage{amsmath}
\usepackage{natbib}

\title{My Document}
\author{Your Name}
\date{\today}

\begin{document}

\maketitle

\section{Introduction}
This is a sample document demonstrating various \LaTeX{} features including mathematical equations, tables, and formatting. You can cite papers using the citation search feature, and they will be automatically added to your \texttt{cite.bib} file.

\section{Mathematical Expressions}
Here are some example equations:

\subsection{Quadratic Formula}
The quadratic formula is given by:
\[x = \frac{-b \pm \sqrt{b^2 - 4ac}}{2a}\]

\subsection{Matrix Example}
A sample matrix:
\[
\begin{pmatrix}
a_{11} & a_{12} & a_{13} \\
a_{21} & a_{22} & a_{23} \\
a_{31} & a_{32} & a_{33}
\end{pmatrix}
\]

\section{Tables}
Below is a sample table showing data:

\begin{table}[h]
\centering
\begin{tabular}{@{}lrrr@{}}
\toprule
\textbf{Category} & \textbf{Value 1} & \textbf{Value 2} & \textbf{Total} \\
\midrule
Item A & 10.5 & 20.3 & 30.8 \\
Item B & 15.2 & 25.1 & 40.3 \\
Item C & 22.1 & 12.7 & 34.8 \\
\bottomrule
\end{tabular}
\caption{Sample Data Table}
\label{tab:sample}
\end{table}

\section{Lists}
\subsection{Bullet Points}
\begin{itemize}
  \item First item
  \item Second item
  \item Third item with sub-items:
    \begin{itemize}
      \item Sub-item 1
      \item Sub-item 2
    \end{itemize}
\end{itemize}

\subsection{Numbered List}
\begin{enumerate}
  \item First step
  \item Second step
  \item Third step
\end{enumerate}

\section{Citations and References}
When you use the citation search feature, references will be automatically managed. For example, you might cite a paper like this: \cite{example2024}. All citations will be automatically added to your \texttt{cite.bib} file.

\section{Conclusion}
This template demonstrates basic \LaTeX{} features including mathematical formulas, tables, and lists. The bibliography system is now set up to work seamlessly with Bibby AI's citation management. Feel free to modify it for your needs!

\bibliographystyle{plain}
\bibliography{cite}

\end{document}\documentclass{article}
\usepackage[utf8]{inputenc}
\usepackage{booktabs}
\usepackage{amsmath}
\usepackage{natbib}

\title{My Document}
\author{Your Name}
\date{\today}

\begin{document}

\maketitle

\section{Introduction}
This is a sample document demonstrating various \LaTeX{} features including mathematical equations, tables, and formatting. You can cite papers using the citation search feature, and they will be automatically added to your \texttt{cite.bib} file.

\section{Mathematical Expressions}
Here are some example equations:

\subsection{Quadratic Formula}
The quadratic formula is given by:
\[x = \frac{-b \pm \sqrt{b^2 - 4ac}}{2a}\]

\subsection{Matrix Example}
A sample matrix:
\[
\begin{pmatrix}
a_{11} & a_{12} & a_{13} \\
a_{21} & a_{22} & a_{23} \\
a_{31} & a_{32} & a_{33}
\end{pmatrix}
\]

\section{Tables}
Below is a sample table showing data:

\begin{table}[h]
\centering
\begin{tabular}{@{}lrrr@{}}
\toprule
\textbf{Category} & \textbf{Value 1} & \textbf{Value 2} & \textbf{Total} \\
\midrule
Item A & 10.5 & 20.3 & 30.8 \\
Item B & 15.2 & 25.1 & 40.3 \\
Item C & 22.1 & 12.7 & 34.8 \\
\bottomrule
\end{tabular}
\caption{Sample Data Table}
\label{tab:sample}
\end{table}

\section{Lists}
\subsection{Bullet Points}
\begin{itemize}
  \item First item
  \item Second item
  \item Third item with sub-items:
    \begin{itemize}
      \item Sub-item 1
      \item Sub-item 2
    \end{itemize}
\end{itemize}

\subsection{Numbered List}
\begin{enumerate}
  \item First step
  \item Second step
  \item Third step
\end{enumerate}

\section{Citations and References}
When you use the citation search feature, references will be automatically managed. For example, you might cite a paper like this: \cite{example2024}. All citations will be automatically added to your \texttt{cite.bib} file.

\section{Conclusion}
This template demonstrates basic \LaTeX{} features including mathematical formulas, tables, and lists. The bibliography system is now set up to work seamlessly with Bibby AI's citation management. Feel free to modify it for your needs!

\bibliographystyle{plain}
\bibliography{cite}

\end{document}\documentclass{article}
\usepackage[utf8]{inputenc}
\usepackage{booktabs}
\usepackage{amsmath}
\usepackage{natbib}

\title{My Document}
\author{Your Name}
\date{\today}

\begin{document}

\maketitle

\section{Introduction}
This is a sample document demonstrating various \LaTeX{} features including mathematical equations, tables, and formatting. You can cite papers using the citation search feature, and they will be automatically added to your \texttt{cite.bib} file.

\section{Mathematical Expressions}
Here are some example equations:

\subsection{Quadratic Formula}
The quadratic formula is given by:
\[x = \frac{-b \pm \sqrt{b^2 - 4ac}}{2a}\]

\subsection{Matrix Example}
A sample matrix:
\[
\begin{pmatrix}
a_{11} & a_{12} & a_{13} \\
a_{21} & a_{22} & a_{23} \\
a_{31} & a_{32} & a_{33}
\end{pmatrix}
\]

\section{Tables}
Below is a sample table showing data:

\begin{table}[h]
\centering
\begin{tabular}{@{}lrrr@{}}
\toprule
\textbf{Category} & \textbf{Value 1} & \textbf{Value 2} & \textbf{Total} \\
\midrule
Item A & 10.5 & 20.3 & 30.8 \\
Item B & 15.2 & 25.1 & 40.3 \\
Item C & 22.1 & 12.7 & 34.8 \\
\bottomrule
\end{tabular}
\caption{Sample Data Table}
\label{tab:sample}
\end{table}

\section{Lists}
\subsection{Bullet Points}
\begin{itemize}
  \item First item
  \item Second item
  \item Third item with sub-items:
    \begin{itemize}
      \item Sub-item 1
      \item Sub-item 2
    \end{itemize}
\end{itemize}

\subsection{Numbered List}
\begin{enumerate}
  \item First step
  \item Second step
  \item Third step
\end{enumerate}

\section{Citations and References}
When you use the citation search feature, references will be automatically managed. For example, you might cite a paper like this: \cite{example2024}. All citations will be automatically added to your \texttt{cite.bib} file.

\section{Conclusion}
This template demonstrates basic \LaTeX{} features including mathematical formulas, tables, and lists. The bibliography system is now set up to work seamlessly with Bibby AI's citation management. Feel free to modify it for your needs!

\bibliographystyle{plain}
\bibliography{cite}

\end{document}\documentclass{article}
\usepackage[utf8]{inputenc}
\usepackage{booktabs}
\usepackage{amsmath}
\usepackage{natbib}

\title{My Document}
\author{Your Name}
\date{\today}

\begin{document}

\maketitle

\section{Introduction}
This is a sample document demonstrating various \LaTeX{} features including mathematical equations, tables, and formatting. You can cite papers using the citation search feature, and they will be automatically added to your \texttt{cite.bib} file.

\section{Mathematical Expressions}
Here are some example equations:

\subsection{Quadratic Formula}
The quadratic formula is given by:
\[x = \frac{-b \pm \sqrt{b^2 - 4ac}}{2a}\]

\subsection{Matrix Example}
A sample matrix:
\[
\begin{pmatrix}
a_{11} & a_{12} & a_{13} \\
a_{21} & a_{22} & a_{23} \\
a_{31} & a_{32} & a_{33}
\end{pmatrix}
\]

\section{Tables}
Below is a sample table showing data:

\begin{table}[h]
\centering
\begin{tabular}{@{}lrrr@{}}
\toprule
\textbf{Category} & \textbf{Value 1} & \textbf{Value 2} & \textbf{Total} \\
\midrule
Item A & 10.5 & 20.3 & 30.8 \\
Item B & 15.2 & 25.1 & 40.3 \\
Item C & 22.1 & 12.7 & 34.8 \\
\bottomrule
\end{tabular}
\caption{Sample Data Table}
\label{tab:sample}
\end{table}

\section{Lists}
\subsection{Bullet Points}
\begin{itemize}
  \item First item
  \item Second item
  \item Third item with sub-items:
    \begin{itemize}
      \item Sub-item 1
      \item Sub-item 2
    \end{itemize}
\end{itemize}

\subsection{Numbered List}
\begin{enumerate}
  \item First step
  \item Second step
  \item Third step
\end{enumerate}

\section{Citations and References}
When you use the citation search feature, references will be automatically managed. For example, you might cite a paper like this: \cite{example2024}. All citations will be automatically added to your \texttt{cite.bib} file.

\section{Conclusion}
This template demonstrates basic \LaTeX{} features including mathematical formulas, tables, and lists. The bibliography system is now set up to work seamlessly with Bibby AI's citation management. Feel free to modify it for your needs!

\bibliographystyle{plain}
\bibliography{cite}

\end{document}\documentclass{article}
\usepackage[utf8]{inputenc}
\usepackage{booktabs}
\usepackage{amsmath}
\usepackage{natbib}

\title{My Document}
\author{Your Name}
\date{\today}

\begin{document}

\maketitle

\section{Introduction}
This is a sample document demonstrating various \LaTeX{} features including mathematical equations, tables, and formatting. You can cite papers using the citation search feature, and they will be automatically added to your \texttt{cite.bib} file.

\section{Mathematical Expressions}
Here are some example equations:

\subsection{Quadratic Formula}
The quadratic formula is given by:
\[x = \frac{-b \pm \sqrt{b^2 - 4ac}}{2a}\]

\subsection{Matrix Example}
A sample matrix:
\[
\begin{pmatrix}
a_{11} & a_{12} & a_{13} \\
a_{21} & a_{22} & a_{23} \\
a_{31} & a_{32} & a_{33}
\end{pmatrix}
\]

\section{Tables}
Below is a sample table showing data:

\begin{table}[h]
\centering
\begin{tabular}{@{}lrrr@{}}
\toprule
\textbf{Category} & \textbf{Value 1} & \textbf{Value 2} & \textbf{Total} \\
\midrule
Item A & 10.5 & 20.3 & 30.8 \\
Item B & 15.2 & 25.1 & 40.3 \\
Item C & 22.1 & 12.7 & 34.8 \\
\bottomrule
\end{tabular}
\caption{Sample Data Table}
\label{tab:sample}
\end{table}

\section{Lists}
\subsection{Bullet Points}
\begin{itemize}
  \item First item
  \item Second item
  \item Third item with sub-items:
    \begin{itemize}
      \item Sub-item 1
      \item Sub-item 2
    \end{itemize}
\end{itemize}

\subsection{Numbered List}
\begin{enumerate}
  \item First step
  \item Second step
  \item Third step
\end{enumerate}

\section{Citations and References}
When you use the citation search feature, references will be automatically managed. For example, you might cite a paper like this: \cite{example2024}. All citations will be automatically added to your \texttt{cite.bib} file.

\section{Conclusion}
This template demonstrates basic \LaTeX{} features including mathematical formulas, tables, and lists. The bibliography system is now set up to work seamlessly with Bibby AI's citation management. Feel free to modify it for your needs!

\bibliographystyle{plain}
\bibliography{cite}

\end{document}\documentclass{article}
\usepackage[utf8]{inputenc}
\usepackage{booktabs}
\usepackage{amsmath}
\usepackage{natbib}

\title{My Document}
\author{Your Name}
\date{\today}

\begin{document}

\maketitle

\section{Introduction}
This is a sample document demonstrating various \LaTeX{} features including mathematical equations, tables, and formatting. You can cite papers using the citation search feature, and they will be automatically added to your \texttt{cite.bib} file.

\section{Mathematical Expressions}
Here are some example equations:

\subsection{Quadratic Formula}
The quadratic formula is given by:
\[x = \frac{-b \pm \sqrt{b^2 - 4ac}}{2a}\]

\subsection{Matrix Example}
A sample matrix:
\[
\begin{pmatrix}
a_{11} & a_{12} & a_{13} \\
a_{21} & a_{22} & a_{23} \\
a_{31} & a_{32} & a_{33}
\end{pmatrix}
\]

\section{Tables}
Below is a sample table showing data:

\begin{table}[h]
\centering
\begin{tabular}{@{}lrrr@{}}
\toprule
\textbf{Category} & \textbf{Value 1} & \textbf{Value 2} & \textbf{Total} \\
\midrule
Item A & 10.5 & 20.3 & 30.8 \\
Item B & 15.2 & 25.1 & 40.3 \\
Item C & 22.1 & 12.7 & 34.8 \\
\bottomrule
\end{tabular}
\caption{Sample Data Table}
\label{tab:sample}
\end{table}

\section{Lists}
\subsection{Bullet Points}
\begin{itemize}
  \item First item
  \item Second item
  \item Third item with sub-items:
    \begin{itemize}
      \item Sub-item 1
      \item Sub-item 2
    \end{itemize}
\end{itemize}

\subsection{Numbered List}
\begin{enumerate}
  \item First step
  \item Second step
  \item Third step
\end{enumerate}

\section{Citations and References}
When you use the citation search feature, references will be automatically managed. For example, you might cite a paper like this: \cite{example2024}. All citations will be automatically added to your \texttt{cite.bib} file.

\section{Conclusion}
This template demonstrates basic \LaTeX{} features including mathematical formulas, tables, and lists. The bibliography system is now set up to work seamlessly with Bibby AI's citation management. Feel free to modify it for your needs!

\bibliographystyle{plain}
\bibliography{cite}

\end{document}\documentclass{article}
\usepackage[utf8]{inputenc}
\usepackage{booktabs}
\usepackage{amsmath}
\usepackage{natbib}

\title{My Document}
\author{Your Name}
\date{\today}

\begin{document}

\maketitle

\section{Introduction}
This is a sample document demonstrating various \LaTeX{} features including mathematical equations, tables, and formatting. You can cite papers using the citation search feature, and they will be automatically added to your \texttt{cite.bib} file.

\section{Mathematical Expressions}
Here are some example equations:

\subsection{Quadratic Formula}
The quadratic formula is given by:
\[x = \frac{-b \pm \sqrt{b^2 - 4ac}}{2a}\]

\subsection{Matrix Example}
A sample matrix:
\[
\begin{pmatrix}
a_{11} & a_{12} & a_{13} \\
a_{21} & a_{22} & a_{23} \\
a_{31} & a_{32} & a_{33}
\end{pmatrix}
\]

\section{Tables}
Below is a sample table showing data:

\begin{table}[h]
\centering
\begin{tabular}{@{}lrrr@{}}
\toprule
\textbf{Category} & \textbf{Value 1} & \textbf{Value 2} & \textbf{Total} \\
\midrule
Item A & 10.5 & 20.3 & 30.8 \\
Item B & 15.2 & 25.1 & 40.3 \\
Item C & 22.1 & 12.7 & 34.8 \\
\bottomrule
\end{tabular}
\caption{Sample Data Table}
\label{tab:sample}
\end{table}

\section{Lists}
\subsection{Bullet Points}
\begin{itemize}
  \item First item
  \item Second item
  \item Third item with sub-items:
    \begin{itemize}
      \item Sub-item 1
      \item Sub-item 2
    \end{itemize}
\end{itemize}

\subsection{Numbered List}
\begin{enumerate}
  \item First step
  \item Second step
  \item Third step
\end{enumerate}

\section{Citations and References}
When you use the citation search feature, references will be automatically managed. For example, you might cite a paper like this: \cite{example2024}. All citations will be automatically added to your \texttt{cite.bib} file.

\section{Conclusion}
This template demonstrates basic \LaTeX{} features including mathematical formulas, tables, and lists. The bibliography system is now set up to work seamlessly with Bibby AI's citation management. Feel free to modify it for your needs!

\bibliographystyle{plain}
\bibliography{cite}

\end{document}\documentclass{article}
\usepackage[utf8]{inputenc}
\usepackage{booktabs}
\usepackage{amsmath}
\usepackage{natbib}

\title{My Document}
\author{Your Name}
\date{\today}

\begin{document}

\maketitle

\section{Introduction}
This is a sample document demonstrating various \LaTeX{} features including mathematical equations, tables, and formatting. You can cite papers using the citation search feature, and they will be automatically added to your \texttt{cite.bib} file.

\section{Mathematical Expressions}
Here are some example equations:

\subsection{Quadratic Formula}
The quadratic formula is given by:
\[x = \frac{-b \pm \sqrt{b^2 - 4ac}}{2a}\]

\subsection{Matrix Example}
A sample matrix:
\[
\begin{pmatrix}
a_{11} & a_{12} & a_{13} \\
a_{21} & a_{22} & a_{23} \\
a_{31} & a_{32} & a_{33}
\end{pmatrix}
\]

\section{Tables}
Below is a sample table showing data:

\begin{table}[h]
\centering
\begin{tabular}{@{}lrrr@{}}
\toprule
\textbf{Category} & \textbf{Value 1} & \textbf{Value 2} & \textbf{Total} \\
\midrule
Item A & 10.5 & 20.3 & 30.8 \\
Item B & 15.2 & 25.1 & 40.3 \\
Item C & 22.1 & 12.7 & 34.8 \\
\bottomrule
\end{tabular}
\caption{Sample Data Table}
\label{tab:sample}
\end{table}

\section{Lists}
\subsection{Bullet Points}
\begin{itemize}
  \item First item
  \item Second item
  \item Third item with sub-items:
    \begin{itemize}
      \item Sub-item 1
      \item Sub-item 2
    \end{itemize}
\end{itemize}

\subsection{Numbered List}
\begin{enumerate}
  \item First step
  \item Second step
  \item Third step
\end{enumerate}

\section{Citations and References}
When you use the citation search feature, references will be automatically managed. For example, you might cite a paper like this: \cite{example2024}. All citations will be automatically added to your \texttt{cite.bib} file.

\section{Conclusion}
This template demonstrates basic \LaTeX{} features including mathematical formulas, tables, and lists. The bibliography system is now set up to work seamlessly with Bibby AI's citation management. Feel free to modify it for your needs!

\bibliographystyle{plain}
\bibliography{cite}

\end{document}\documentclass{article}
\usepackage[utf8]{inputenc}
\usepackage{booktabs}
\usepackage{amsmath}
\usepackage{natbib}

\title{My Document}
\author{Your Name}
\date{\today}

\begin{document}

\maketitle

\section{Introduction}
This is a sample document demonstrating various \LaTeX{} features including mathematical equations, tables, and formatting. You can cite papers using the citation search feature, and they will be automatically added to your \texttt{cite.bib} file.

\section{Mathematical Expressions}
Here are some example equations:

\subsection{Quadratic Formula}
The quadratic formula is given by:
\[x = \frac{-b \pm \sqrt{b^2 - 4ac}}{2a}\]

\subsection{Matrix Example}
A sample matrix:
\[
\begin{pmatrix}
a_{11} & a_{12} & a_{13} \\
a_{21} & a_{22} & a_{23} \\
a_{31} & a_{32} & a_{33}
\end{pmatrix}
\]

\section{Tables}
Below is a sample table showing data:

\begin{table}[h]
\centering
\begin{tabular}{@{}lrrr@{}}
\toprule
\textbf{Category} & \textbf{Value 1} & \textbf{Value 2} & \textbf{Total} \\
\midrule
Item A & 10.5 & 20.3 & 30.8 \\
Item B & 15.2 & 25.1 & 40.3 \\
Item C & 22.1 & 12.7 & 34.8 \\
\bottomrule
\end{tabular}
\caption{Sample Data Table}
\label{tab:sample}
\end{table}

\section{Lists}
\subsection{Bullet Points}
\begin{itemize}
  \item First item
  \item Second item
  \item Third item with sub-items:
    \begin{itemize}
      \item Sub-item 1
      \item Sub-item 2
    \end{itemize}
\end{itemize}

\subsection{Numbered List}
\begin{enumerate}
  \item First step
  \item Second step
  \item Third step
\end{enumerate}

\section{Citations and References}
When you use the citation search feature, references will be automatically managed. For example, you might cite a paper like this: \cite{example2024}. All citations will be automatically added to your \texttt{cite.bib} file.

\section{Conclusion}
This template demonstrates basic \LaTeX{} features including mathematical formulas, tables, and lists. The bibliography system is now set up to work seamlessly with Bibby AI's citation management. Feel free to modify it for your needs!

\bibliographystyle{plain}
\bibliography{cite}

\end{document}\documentclass{article}
\usepackage[utf8]{inputenc}
\usepackage{booktabs}
\usepackage{amsmath}
\usepackage{natbib}

\title{My Document}
\author{Your Name}
\date{\today}

\begin{document}

\maketitle

\section{Introduction}
This is a sample document demonstrating various \LaTeX{} features including mathematical equations, tables, and formatting. You can cite papers using the citation search feature, and they will be automatically added to your \texttt{cite.bib} file.

\section{Mathematical Expressions}
Here are some example equations:

\subsection{Quadratic Formula}
The quadratic formula is given by:
\[x = \frac{-b \pm \sqrt{b^2 - 4ac}}{2a}\]

\subsection{Matrix Example}
A sample matrix:
\[
\begin{pmatrix}
a_{11} & a_{12} & a_{13} \\
a_{21} & a_{22} & a_{23} \\
a_{31} & a_{32} & a_{33}
\end{pmatrix}
\]

\section{Tables}
Below is a sample table showing data:

\begin{table}[h]
\centering
\begin{tabular}{@{}lrrr@{}}
\toprule
\textbf{Category} & \textbf{Value 1} & \textbf{Value 2} & \textbf{Total} \\
\midrule
Item A & 10.5 & 20.3 & 30.8 \\
Item B & 15.2 & 25.1 & 40.3 \\
Item C & 22.1 & 12.7 & 34.8 \\
\bottomrule
\end{tabular}
\caption{Sample Data Table}
\label{tab:sample}
\end{table}

\section{Lists}
\subsection{Bullet Points}
\begin{itemize}
  \item First item
  \item Second item
  \item Third item with sub-items:
    \begin{itemize}
      \item Sub-item 1
      \item Sub-item 2
    \end{itemize}
\end{itemize}

\subsection{Numbered List}
\begin{enumerate}
  \item First step
  \item Second step
  \item Third step
\end{enumerate}

\section{Citations and References}
When you use the citation search feature, references will be automatically managed. For example, you might cite a paper like this: \cite{example2024}. All citations will be automatically added to your \texttt{cite.bib} file.

\section{Conclusion}
This template demonstrates basic \LaTeX{} features including mathematical formulas, tables, and lists. The bibliography system is now set up to work seamlessly with Bibby AI's citation management. Feel free to modify it for your needs!

\bibliographystyle{plain}
\bibliography{cite}

\end{document}\documentclass{article}
\usepackage[utf8]{inputenc}
\usepackage{booktabs}
\usepackage{amsmath}
\usepackage{natbib}

\title{My Document}
\author{Your Name}
\date{\today}

\begin{document}

\maketitle

\section{Introduction}
This is a sample document demonstrating various \LaTeX{} features including mathematical equations, tables, and formatting. You can cite papers using the citation search feature, and they will be automatically added to your \texttt{cite.bib} file.

\section{Mathematical Expressions}
Here are some example equations:

\subsection{Quadratic Formula}
The quadratic formula is given by:
\[x = \frac{-b \pm \sqrt{b^2 - 4ac}}{2a}\]

\subsection{Matrix Example}
A sample matrix:
\[
\begin{pmatrix}
a_{11} & a_{12} & a_{13} \\
a_{21} & a_{22} & a_{23} \\
a_{31} & a_{32} & a_{33}
\end{pmatrix}
\]

\section{Tables}
Below is a sample table showing data:

\begin{table}[h]
\centering
\begin{tabular}{@{}lrrr@{}}
\toprule
\textbf{Category} & \textbf{Value 1} & \textbf{Value 2} & \textbf{Total} \\
\midrule
Item A & 10.5 & 20.3 & 30.8 \\
Item B & 15.2 & 25.1 & 40.3 \\
Item C & 22.1 & 12.7 & 34.8 \\
\bottomrule
\end{tabular}
\caption{Sample Data Table}
\label{tab:sample}
\end{table}

\section{Lists}
\subsection{Bullet Points}
\begin{itemize}
  \item First item
  \item Second item
  \item Third item with sub-items:
    \begin{itemize}
      \item Sub-item 1
      \item Sub-item 2
    \end{itemize}
\end{itemize}

\subsection{Numbered List}
\begin{enumerate}
  \item First step
  \item Second step
  \item Third step
\end{enumerate}

\section{Citations and References}
When you use the citation search feature, references will be automatically managed. For example, you might cite a paper like this: \cite{example2024}. All citations will be automatically added to your \texttt{cite.bib} file.

\section{Conclusion}
This template demonstrates basic \LaTeX{} features including mathematical formulas, tables, and lists. The bibliography system is now set up to work seamlessly with Bibby AI's citation management. Feel free to modify it for your needs!

\bibliographystyle{plain}
\bibliography{cite}

\end{document}\documentclass{article}
\usepackage[utf8]{inputenc}
\usepackage{booktabs}
\usepackage{amsmath}
\usepackage{natbib}

\title{My Document}
\author{Your Name}
\date{\today}

\begin{document}

\maketitle

\section{Introduction}
This is a sample document demonstrating various \LaTeX{} features including mathematical equations, tables, and formatting. You can cite papers using the citation search feature, and they will be automatically added to your \texttt{cite.bib} file.

\section{Mathematical Expressions}
Here are some example equations:

\subsection{Quadratic Formula}
The quadratic formula is given by:
\[x = \frac{-b \pm \sqrt{b^2 - 4ac}}{2a}\]

\subsection{Matrix Example}
A sample matrix:
\[
\begin{pmatrix}
a_{11} & a_{12} & a_{13} \\
a_{21} & a_{22} & a_{23} \\
a_{31} & a_{32} & a_{33}
\end{pmatrix}
\]

\section{Tables}
Below is a sample table showing data:

\begin{table}[h]
\centering
\begin{tabular}{@{}lrrr@{}}
\toprule
\textbf{Category} & \textbf{Value 1} & \textbf{Value 2} & \textbf{Total} \\
\midrule
Item A & 10.5 & 20.3 & 30.8 \\
Item B & 15.2 & 25.1 & 40.3 \\
Item C & 22.1 & 12.7 & 34.8 \\
\bottomrule
\end{tabular}
\caption{Sample Data Table}
\label{tab:sample}
\end{table}

\section{Lists}
\subsection{Bullet Points}
\begin{itemize}
  \item First item
  \item Second item
  \item Third item with sub-items:
    \begin{itemize}
      \item Sub-item 1
      \item Sub-item 2
    \end{itemize}
\end{itemize}

\subsection{Numbered List}
\begin{enumerate}
  \item First step
  \item Second step
  \item Third step
\end{enumerate}

\section{Citations and References}
When you use the citation search feature, references will be automatically managed. For example, you might cite a paper like this: \cite{example2024}. All citations will be automatically added to your \texttt{cite.bib} file.

\section{Conclusion}
This template demonstrates basic \LaTeX{} features including mathematical formulas, tables, and lists. The bibliography system is now set up to work seamlessly with Bibby AI's citation management. Feel free to modify it for your needs!

\bibliographystyle{plain}
\bibliography{cite}

\end{document}\documentclass{article}
\usepackage[utf8]{inputenc}
\usepackage{booktabs}
\usepackage{amsmath}
\usepackage{natbib}

\title{My Document}
\author{Your Name}
\date{\today}

\begin{document}

\maketitle

\section{Introduction}
This is a sample document demonstrating various \LaTeX{} features including mathematical equations, tables, and formatting. You can cite papers using the citation search feature, and they will be automatically added to your \texttt{cite.bib} file.

\section{Mathematical Expressions}
Here are some example equations:

\subsection{Quadratic Formula}
The quadratic formula is given by:
\[x = \frac{-b \pm \sqrt{b^2 - 4ac}}{2a}\]

\subsection{Matrix Example}
A sample matrix:
\[
\begin{pmatrix}
a_{11} & a_{12} & a_{13} \\
a_{21} & a_{22} & a_{23} \\
a_{31} & a_{32} & a_{33}
\end{pmatrix}
\]

\section{Tables}
Below is a sample table showing data:

\begin{table}[h]
\centering
\begin{tabular}{@{}lrrr@{}}
\toprule
\textbf{Category} & \textbf{Value 1} & \textbf{Value 2} & \textbf{Total} \\
\midrule
Item A & 10.5 & 20.3 & 30.8 \\
Item B & 15.2 & 25.1 & 40.3 \\
Item C & 22.1 & 12.7 & 34.8 \\
\bottomrule
\end{tabular}
\caption{Sample Data Table}
\label{tab:sample}
\end{table}

\section{Lists}
\subsection{Bullet Points}
\begin{itemize}
  \item First item
  \item Second item
  \item Third item with sub-items:
    \begin{itemize}
      \item Sub-item 1
      \item Sub-item 2
    \end{itemize}
\end{itemize}

\subsection{Numbered List}
\begin{enumerate}
  \item First step
  \item Second step
  \item Third step
\end{enumerate}

\section{Citations and References}
When you use the citation search feature, references will be automatically managed. For example, you might cite a paper like this: \cite{example2024}. All citations will be automatically added to your \texttt{cite.bib} file.

\section{Conclusion}
This template demonstrates basic \LaTeX{} features including mathematical formulas, tables, and lists. The bibliography system is now set up to work seamlessly with Bibby AI's citation management. Feel free to modify it for your needs!

\bibliographystyle{plain}
\bibliography{cite}

\end{document}\documentclass{article}
\usepackage[utf8]{inputenc}
\usepackage{booktabs}
\usepackage{amsmath}
\usepackage{natbib}

\title{My Document}
\author{Your Name}
\date{\today}

\begin{document}

\maketitle

\section{Introduction}
This is a sample document demonstrating various \LaTeX{} features including mathematical equations, tables, and formatting. You can cite papers using the citation search feature, and they will be automatically added to your \texttt{cite.bib} file.

\section{Mathematical Expressions}
Here are some example equations:

\subsection{Quadratic Formula}
The quadratic formula is given by:
\[x = \frac{-b \pm \sqrt{b^2 - 4ac}}{2a}\]

\subsection{Matrix Example}
A sample matrix:
\[
\begin{pmatrix}
a_{11} & a_{12} & a_{13} \\
a_{21} & a_{22} & a_{23} \\
a_{31} & a_{32} & a_{33}
\end{pmatrix}
\]

\section{Tables}
Below is a sample table showing data:

\begin{table}[h]
\centering
\begin{tabular}{@{}lrrr@{}}
\toprule
\textbf{Category} & \textbf{Value 1} & \textbf{Value 2} & \textbf{Total} \\
\midrule
Item A & 10.5 & 20.3 & 30.8 \\
Item B & 15.2 & 25.1 & 40.3 \\
Item C & 22.1 & 12.7 & 34.8 \\
\bottomrule
\end{tabular}
\caption{Sample Data Table}
\label{tab:sample}
\end{table}

\section{Lists}
\subsection{Bullet Points}
\begin{itemize}
  \item First item
  \item Second item
  \item Third item with sub-items:
    \begin{itemize}
      \item Sub-item 1
      \item Sub-item 2
    \end{itemize}
\end{itemize}

\subsection{Numbered List}
\begin{enumerate}
  \item First step
  \item Second step
  \item Third step
\end{enumerate}

\section{Citations and References}
When you use the citation search feature, references will be automatically managed. For example, you might cite a paper like this: \cite{example2024}. All citations will be automatically added to your \texttt{cite.bib} file.

\section{Conclusion}
This template demonstrates basic \LaTeX{} features including mathematical formulas, tables, and lists. The bibliography system is now set up to work seamlessly with Bibby AI's citation management. Feel free to modify it for your needs!

\bibliographystyle{plain}
\bibliography{cite}

\end{document}\documentclass{article}
\usepackage[utf8]{inputenc}
\usepackage{booktabs}
\usepackage{amsmath}
\usepackage{natbib}

\title{My Document}
\author{Your Name}
\date{\today}

\begin{document}

\maketitle

\section{Introduction}
This is a sample document demonstrating various \LaTeX{} features including mathematical equations, tables, and formatting. You can cite papers using the citation search feature, and they will be automatically added to your \texttt{cite.bib} file.

\section{Mathematical Expressions}
Here are some example equations:

\subsection{Quadratic Formula}
The quadratic formula is given by:
\[x = \frac{-b \pm \sqrt{b^2 - 4ac}}{2a}\]

\subsection{Matrix Example}
A sample matrix:
\[
\begin{pmatrix}
a_{11} & a_{12} & a_{13} \\
a_{21} & a_{22} & a_{23} \\
a_{31} & a_{32} & a_{33}
\end{pmatrix}
\]

\section{Tables}
Below is a sample table showing data:

\begin{table}[h]
\centering
\begin{tabular}{@{}lrrr@{}}
\toprule
\textbf{Category} & \textbf{Value 1} & \textbf{Value 2} & \textbf{Total} \\
\midrule
Item A & 10.5 & 20.3 & 30.8 \\
Item B & 15.2 & 25.1 & 40.3 \\
Item C & 22.1 & 12.7 & 34.8 \\
\bottomrule
\end{tabular}
\caption{Sample Data Table}
\label{tab:sample}
\end{table}

\section{Lists}
\subsection{Bullet Points}
\begin{itemize}
  \item First item
  \item Second item
  \item Third item with sub-items:
    \begin{itemize}
      \item Sub-item 1
      \item Sub-item 2
    \end{itemize}
\end{itemize}

\subsection{Numbered List}
\begin{enumerate}
  \item First step
  \item Second step
  \item Third step
\end{enumerate}

\section{Citations and References}
When you use the citation search feature, references will be automatically managed. For example, you might cite a paper like this: \cite{example2024}. All citations will be automatically added to your \texttt{cite.bib} file.

\section{Conclusion}
This template demonstrates basic \LaTeX{} features including mathematical formulas, tables, and lists. The bibliography system is now set up to work seamlessly with Bibby AI's citation management. Feel free to modify it for your needs!

\bibliographystyle{plain}
\bibliography{cite}

\end{document}\documentclass{article}
\usepackage[utf8]{inputenc}
\usepackage{booktabs}
\usepackage{amsmath}
\usepackage{natbib}

\title{My Document}
\author{Your Name}
\date{\today}

\begin{document}

\maketitle

\section{Introduction}
This is a sample document demonstrating various \LaTeX{} features including mathematical equations, tables, and formatting. You can cite papers using the citation search feature, and they will be automatically added to your \texttt{cite.bib} file.

\section{Mathematical Expressions}
Here are some example equations:

\subsection{Quadratic Formula}
The quadratic formula is given by:
\[x = \frac{-b \pm \sqrt{b^2 - 4ac}}{2a}\]

\subsection{Matrix Example}
A sample matrix:
\[
\begin{pmatrix}
a_{11} & a_{12} & a_{13} \\
a_{21} & a_{22} & a_{23} \\
a_{31} & a_{32} & a_{33}
\end{pmatrix}
\]

\section{Tables}
Below is a sample table showing data:

\begin{table}[h]
\centering
\begin{tabular}{@{}lrrr@{}}
\toprule
\textbf{Category} & \textbf{Value 1} & \textbf{Value 2} & \textbf{Total} \\
\midrule
Item A & 10.5 & 20.3 & 30.8 \\
Item B & 15.2 & 25.1 & 40.3 \\
Item C & 22.1 & 12.7 & 34.8 \\
\bottomrule
\end{tabular}
\caption{Sample Data Table}
\label{tab:sample}
\end{table}

\section{Lists}
\subsection{Bullet Points}
\begin{itemize}
  \item First item
  \item Second item
  \item Third item with sub-items:
    \begin{itemize}
      \item Sub-item 1
      \item Sub-item 2
    \end{itemize}
\end{itemize}

\subsection{Numbered List}
\begin{enumerate}
  \item First step
  \item Second step
  \item Third step
\end{enumerate}

\section{Citations and References}
When you use the citation search feature, references will be automatically managed. For example, you might cite a paper like this: \cite{example2024}. All citations will be automatically added to your \texttt{cite.bib} file.

\section{Conclusion}
This template demonstrates basic \LaTeX{} features including mathematical formulas, tables, and lists. The bibliography system is now set up to work seamlessly with Bibby AI's citation management. Feel free to modify it for your needs!

\bibliographystyle{plain}
\bibliography{cite}

\end{document}\documentclass{article}
\usepackage[utf8]{inputenc}
\usepackage{booktabs}
\usepackage{amsmath}
\usepackage{natbib}

\title{My Document}
\author{Your Name}
\date{\today}

\begin{document}

\maketitle

\section{Introduction}
This is a sample document demonstrating various \LaTeX{} features including mathematical equations, tables, and formatting. You can cite papers using the citation search feature, and they will be automatically added to your \texttt{cite.bib} file.

\section{Mathematical Expressions}
Here are some example equations:

\subsection{Quadratic Formula}
The quadratic formula is given by:
\[x = \frac{-b \pm \sqrt{b^2 - 4ac}}{2a}\]

\subsection{Matrix Example}
A sample matrix:
\[
\begin{pmatrix}
a_{11} & a_{12} & a_{13} \\
a_{21} & a_{22} & a_{23} \\
a_{31} & a_{32} & a_{33}
\end{pmatrix}
\]

\section{Tables}
Below is a sample table showing data:

\begin{table}[h]
\centering
\begin{tabular}{@{}lrrr@{}}
\toprule
\textbf{Category} & \textbf{Value 1} & \textbf{Value 2} & \textbf{Total} \\
\midrule
Item A & 10.5 & 20.3 & 30.8 \\
Item B & 15.2 & 25.1 & 40.3 \\
Item C & 22.1 & 12.7 & 34.8 \\
\bottomrule
\end{tabular}
\caption{Sample Data Table}
\label{tab:sample}
\end{table}

\section{Lists}
\subsection{Bullet Points}
\begin{itemize}
  \item First item
  \item Second item
  \item Third item with sub-items:
    \begin{itemize}
      \item Sub-item 1
      \item Sub-item 2
    \end{itemize}
\end{itemize}

\subsection{Numbered List}
\begin{enumerate}
  \item First step
  \item Second step
  \item Third step
\end{enumerate}

\section{Citations and References}
When you use the citation search feature, references will be automatically managed. For example, you might cite a paper like this: \cite{example2024}. All citations will be automatically added to your \texttt{cite.bib} file.

\section{Conclusion}
This template demonstrates basic \LaTeX{} features including mathematical formulas, tables, and lists. The bibliography system is now set up to work seamlessly with Bibby AI's citation management. Feel free to modify it for your needs!

\bibliographystyle{plain}
\bibliography{cite}

\end{document}\documentclass{article}
\usepackage[utf8]{inputenc}
\usepackage{booktabs}
\usepackage{amsmath}
\usepackage{natbib}

\title{My Document}
\author{Your Name}
\date{\today}

\begin{document}

\maketitle

\section{Introduction}
This is a sample document demonstrating various \LaTeX{} features including mathematical equations, tables, and formatting. You can cite papers using the citation search feature, and they will be automatically added to your \texttt{cite.bib} file.

\section{Mathematical Expressions}
Here are some example equations:

\subsection{Quadratic Formula}
The quadratic formula is given by:
\[x = \frac{-b \pm \sqrt{b^2 - 4ac}}{2a}\]

\subsection{Matrix Example}
A sample matrix:
\[
\begin{pmatrix}
a_{11} & a_{12} & a_{13} \\
a_{21} & a_{22} & a_{23} \\
a_{31} & a_{32} & a_{33}
\end{pmatrix}
\]

\section{Tables}
Below is a sample table showing data:

\begin{table}[h]
\centering
\begin{tabular}{@{}lrrr@{}}
\toprule
\textbf{Category} & \textbf{Value 1} & \textbf{Value 2} & \textbf{Total} \\
\midrule
Item A & 10.5 & 20.3 & 30.8 \\
Item B & 15.2 & 25.1 & 40.3 \\
Item C & 22.1 & 12.7 & 34.8 \\
\bottomrule
\end{tabular}
\caption{Sample Data Table}
\label{tab:sample}
\end{table}

\section{Lists}
\subsection{Bullet Points}
\begin{itemize}
  \item First item
  \item Second item
  \item Third item with sub-items:
    \begin{itemize}
      \item Sub-item 1
      \item Sub-item 2
    \end{itemize}
\end{itemize}

\subsection{Numbered List}
\begin{enumerate}
  \item First step
  \item Second step
  \item Third step
\end{enumerate}

\section{Citations and References}
When you use the citation search feature, references will be automatically managed. For example, you might cite a paper like this: \cite{example2024}. All citations will be automatically added to your \texttt{cite.bib} file.

\section{Conclusion}
This template demonstrates basic \LaTeX{} features including mathematical formulas, tables, and lists. The bibliography system is now set up to work seamlessly with Bibby AI's citation management. Feel free to modify it for your needs!

\bibliographystyle{plain}
\bibliography{cite}

\end{document}\documentclass{article}
\usepackage[utf8]{inputenc}
\usepackage{booktabs}
\usepackage{amsmath}
\usepackage{natbib}

\title{My Document}
\author{Your Name}
\date{\today}

\begin{document}

\maketitle

\section{Introduction}
This is a sample document demonstrating various \LaTeX{} features including mathematical equations, tables, and formatting. You can cite papers using the citation search feature, and they will be automatically added to your \texttt{cite.bib} file.

\section{Mathematical Expressions}
Here are some example equations:

\subsection{Quadratic Formula}
The quadratic formula is given by:
\[x = \frac{-b \pm \sqrt{b^2 - 4ac}}{2a}\]

\subsection{Matrix Example}
A sample matrix:
\[
\begin{pmatrix}
a_{11} & a_{12} & a_{13} \\
a_{21} & a_{22} & a_{23} \\
a_{31} & a_{32} & a_{33}
\end{pmatrix}
\]

\section{Tables}
Below is a sample table showing data:

\begin{table}[h]
\centering
\begin{tabular}{@{}lrrr@{}}
\toprule
\textbf{Category} & \textbf{Value 1} & \textbf{Value 2} & \textbf{Total} \\
\midrule
Item A & 10.5 & 20.3 & 30.8 \\
Item B & 15.2 & 25.1 & 40.3 \\
Item C & 22.1 & 12.7 & 34.8 \\
\bottomrule
\end{tabular}
\caption{Sample Data Table}
\label{tab:sample}
\end{table}

\section{Lists}
\subsection{Bullet Points}
\begin{itemize}
  \item First item
  \item Second item
  \item Third item with sub-items:
    \begin{itemize}
      \item Sub-item 1
      \item Sub-item 2
    \end{itemize}
\end{itemize}

\subsection{Numbered List}
\begin{enumerate}
  \item First step
  \item Second step
  \item Third step
\end{enumerate}

\section{Citations and References}
When you use the citation search feature, references will be automatically managed. For example, you might cite a paper like this: \cite{example2024}. All citations will be automatically added to your \texttt{cite.bib} file.

\section{Conclusion}
This template demonstrates basic \LaTeX{} features including mathematical formulas, tables, and lists. The bibliography system is now set up to work seamlessly with Bibby AI's citation management. Feel free to modify it for your needs!

\bibliographystyle{plain}
\bibliography{cite}

\end{document}\documentclass{article}
\usepackage[utf8]{inputenc}
\usepackage{booktabs}
\usepackage{amsmath}
\usepackage{natbib}

\title{My Document}
\author{Your Name}
\date{\today}

\begin{document}

\maketitle

\section{Introduction}
This is a sample document demonstrating various \LaTeX{} features including mathematical equations, tables, and formatting. You can cite papers using the citation search feature, and they will be automatically added to your \texttt{cite.bib} file.

\section{Mathematical Expressions}
Here are some example equations:

\subsection{Quadratic Formula}
The quadratic formula is given by:
\[x = \frac{-b \pm \sqrt{b^2 - 4ac}}{2a}\]

\subsection{Matrix Example}
A sample matrix:
\[
\begin{pmatrix}
a_{11} & a_{12} & a_{13} \\
a_{21} & a_{22} & a_{23} \\
a_{31} & a_{32} & a_{33}
\end{pmatrix}
\]

\section{Tables}
Below is a sample table showing data:

\begin{table}[h]
\centering
\begin{tabular}{@{}lrrr@{}}
\toprule
\textbf{Category} & \textbf{Value 1} & \textbf{Value 2} & \textbf{Total} \\
\midrule
Item A & 10.5 & 20.3 & 30.8 \\
Item B & 15.2 & 25.1 & 40.3 \\
Item C & 22.1 & 12.7 & 34.8 \\
\bottomrule
\end{tabular}
\caption{Sample Data Table}
\label{tab:sample}
\end{table}

\section{Lists}
\subsection{Bullet Points}
\begin{itemize}
  \item First item
  \item Second item
  \item Third item with sub-items:
    \begin{itemize}
      \item Sub-item 1
      \item Sub-item 2
    \end{itemize}
\end{itemize}

\subsection{Numbered List}
\begin{enumerate}
  \item First step
  \item Second step
  \item Third step
\end{enumerate}

\section{Citations and References}
When you use the citation search feature, references will be automatically managed. For example, you might cite a paper like this: \cite{example2024}. All citations will be automatically added to your \texttt{cite.bib} file.

\section{Conclusion}
This template demonstrates basic \LaTeX{} features including mathematical formulas, tables, and lists. The bibliography system is now set up to work seamlessly with Bibby AI's citation management. Feel free to modify it for your needs!

\bibliographystyle{plain}
\bibliography{cite}

\end{document}\documentclass{article}
\usepackage[utf8]{inputenc}
\usepackage{booktabs}
\usepackage{amsmath}
\usepackage{natbib}

\title{My Document}
\author{Your Name}
\date{\today}

\begin{document}

\maketitle

\section{Introduction}
This is a sample document demonstrating various \LaTeX{} features including mathematical equations, tables, and formatting. You can cite papers using the citation search feature, and they will be automatically added to your \texttt{cite.bib} file.

\section{Mathematical Expressions}
Here are some example equations:

\subsection{Quadratic Formula}
The quadratic formula is given by:
\[x = \frac{-b \pm \sqrt{b^2 - 4ac}}{2a}\]

\subsection{Matrix Example}
A sample matrix:
\[
\begin{pmatrix}
a_{11} & a_{12} & a_{13} \\
a_{21} & a_{22} & a_{23} \\
a_{31} & a_{32} & a_{33}
\end{pmatrix}
\]

\section{Tables}
Below is a sample table showing data:

\begin{table}[h]
\centering
\begin{tabular}{@{}lrrr@{}}
\toprule
\textbf{Category} & \textbf{Value 1} & \textbf{Value 2} & \textbf{Total} \\
\midrule
Item A & 10.5 & 20.3 & 30.8 \\
Item B & 15.2 & 25.1 & 40.3 \\
Item C & 22.1 & 12.7 & 34.8 \\
\bottomrule
\end{tabular}
\caption{Sample Data Table}
\label{tab:sample}
\end{table}

\section{Lists}
\subsection{Bullet Points}
\begin{itemize}
  \item First item
  \item Second item
  \item Third item with sub-items:
    \begin{itemize}
      \item Sub-item 1
      \item Sub-item 2
    \end{itemize}
\end{itemize}

\subsection{Numbered List}
\begin{enumerate}
  \item First step
  \item Second step
  \item Third step
\end{enumerate}

\section{Citations and References}
When you use the citation search feature, references will be automatically managed. For example, you might cite a paper like this: \cite{example2024}. All citations will be automatically added to your \texttt{cite.bib} file.

\section{Conclusion}
This template demonstrates basic \LaTeX{} features including mathematical formulas, tables, and lists. The bibliography system is now set up to work seamlessly with Bibby AI's citation management. Feel free to modify it for your needs!

\bibliographystyle{plain}
\bibliography{cite}

\end{document}\documentclass{article}
\usepackage[utf8]{inputenc}
\usepackage{booktabs}
\usepackage{amsmath}
\usepackage{natbib}

\title{My Document}
\author{Your Name}
\date{\today}

\begin{document}

\maketitle

\section{Introduction}
This is a sample document demonstrating various \LaTeX{} features including mathematical equations, tables, and formatting. You can cite papers using the citation search feature, and they will be automatically added to your \texttt{cite.bib} file.

\section{Mathematical Expressions}
Here are some example equations:

\subsection{Quadratic Formula}
The quadratic formula is given by:
\[x = \frac{-b \pm \sqrt{b^2 - 4ac}}{2a}\]

\subsection{Matrix Example}
A sample matrix:
\[
\begin{pmatrix}
a_{11} & a_{12} & a_{13} \\
a_{21} & a_{22} & a_{23} \\
a_{31} & a_{32} & a_{33}
\end{pmatrix}
\]

\section{Tables}
Below is a sample table showing data:

\begin{table}[h]
\centering
\begin{tabular}{@{}lrrr@{}}
\toprule
\textbf{Category} & \textbf{Value 1} & \textbf{Value 2} & \textbf{Total} \\
\midrule
Item A & 10.5 & 20.3 & 30.8 \\
Item B & 15.2 & 25.1 & 40.3 \\
Item C & 22.1 & 12.7 & 34.8 \\
\bottomrule
\end{tabular}
\caption{Sample Data Table}
\label{tab:sample}
\end{table}

\section{Lists}
\subsection{Bullet Points}
\begin{itemize}
  \item First item
  \item Second item
  \item Third item with sub-items:
    \begin{itemize}
      \item Sub-item 1
      \item Sub-item 2
    \end{itemize}
\end{itemize}

\subsection{Numbered List}
\begin{enumerate}
  \item First step
  \item Second step
  \item Third step
\end{enumerate}

\section{Citations and References}
When you use the citation search feature, references will be automatically managed. For example, you might cite a paper like this: \cite{example2024}. All citations will be automatically added to your \texttt{cite.bib} file.

\section{Conclusion}
This template demonstrates basic \LaTeX{} features including mathematical formulas, tables, and lists. The bibliography system is now set up to work seamlessly with Bibby AI's citation management. Feel free to modify it for your needs!

\bibliographystyle{plain}
\bibliography{cite}

\end{document}\documentclass{article}
\usepackage[utf8]{inputenc}
\usepackage{booktabs}
\usepackage{amsmath}
\usepackage{natbib}

\title{My Document}
\author{Your Name}
\date{\today}

\begin{document}

\maketitle

\section{Introduction}
This is a sample document demonstrating various \LaTeX{} features including mathematical equations, tables, and formatting. You can cite papers using the citation search feature, and they will be automatically added to your \texttt{cite.bib} file.

\section{Mathematical Expressions}
Here are some example equations:

\subsection{Quadratic Formula}
The quadratic formula is given by:
\[x = \frac{-b \pm \sqrt{b^2 - 4ac}}{2a}\]

\subsection{Matrix Example}
A sample matrix:
\[
\begin{pmatrix}
a_{11} & a_{12} & a_{13} \\
a_{21} & a_{22} & a_{23} \\
a_{31} & a_{32} & a_{33}
\end{pmatrix}
\]

\section{Tables}
Below is a sample table showing data:

\begin{table}[h]
\centering
\begin{tabular}{@{}lrrr@{}}
\toprule
\textbf{Category} & \textbf{Value 1} & \textbf{Value 2} & \textbf{Total} \\
\midrule
Item A & 10.5 & 20.3 & 30.8 \\
Item B & 15.2 & 25.1 & 40.3 \\
Item C & 22.1 & 12.7 & 34.8 \\
\bottomrule
\end{tabular}
\caption{Sample Data Table}
\label{tab:sample}
\end{table}

\section{Lists}
\subsection{Bullet Points}
\begin{itemize}
  \item First item
  \item Second item
  \item Third item with sub-items:
    \begin{itemize}
      \item Sub-item 1
      \item Sub-item 2
    \end{itemize}
\end{itemize}

\subsection{Numbered List}
\begin{enumerate}
  \item First step
  \item Second step
  \item Third step
\end{enumerate}

\section{Citations and References}
When you use the citation search feature, references will be automatically managed. For example, you might cite a paper like this: \cite{example2024}. All citations will be automatically added to your \texttt{cite.bib} file.

\section{Conclusion}
This template demonstrates basic \LaTeX{} features including mathematical formulas, tables, and lists. The bibliography system is now set up to work seamlessly with Bibby AI's citation management. Feel free to modify it for your needs!

\bibliographystyle{plain}
\bibliography{cite}

\end{document}\documentclass{article}
\usepackage[utf8]{inputenc}
\usepackage{booktabs}
\usepackage{amsmath}
\usepackage{natbib}

\title{My Document}
\author{Your Name}
\date{\today}

\begin{document}

\maketitle

\section{Introduction}
This is a sample document demonstrating various \LaTeX{} features including mathematical equations, tables, and formatting. You can cite papers using the citation search feature, and they will be automatically added to your \texttt{cite.bib} file.

\section{Mathematical Expressions}
Here are some example equations:

\subsection{Quadratic Formula}
The quadratic formula is given by:
\[x = \frac{-b \pm \sqrt{b^2 - 4ac}}{2a}\]

\subsection{Matrix Example}
A sample matrix:
\[
\begin{pmatrix}
a_{11} & a_{12} & a_{13} \\
a_{21} & a_{22} & a_{23} \\
a_{31} & a_{32} & a_{33}
\end{pmatrix}
\]

\section{Tables}
Below is a sample table showing data:

\begin{table}[h]
\centering
\begin{tabular}{@{}lrrr@{}}
\toprule
\textbf{Category} & \textbf{Value 1} & \textbf{Value 2} & \textbf{Total} \\
\midrule
Item A & 10.5 & 20.3 & 30.8 \\
Item B & 15.2 & 25.1 & 40.3 \\
Item C & 22.1 & 12.7 & 34.8 \\
\bottomrule
\end{tabular}
\caption{Sample Data Table}
\label{tab:sample}
\end{table}

\section{Lists}
\subsection{Bullet Points}
\begin{itemize}
  \item First item
  \item Second item
  \item Third item with sub-items:
    \begin{itemize}
      \item Sub-item 1
      \item Sub-item 2
    \end{itemize}
\end{itemize}

\subsection{Numbered List}
\begin{enumerate}
  \item First step
  \item Second step
  \item Third step
\end{enumerate}

\section{Citations and References}
When you use the citation search feature, references will be automatically managed. For example, you might cite a paper like this: \cite{example2024}. All citations will be automatically added to your \texttt{cite.bib} file.

\section{Conclusion}
This template demonstrates basic \LaTeX{} features including mathematical formulas, tables, and lists. The bibliography system is now set up to work seamlessly with Bibby AI's citation management. Feel free to modify it for your needs!

\bibliographystyle{plain}
\bibliography{cite}

\end{document}\documentclass{article}
\usepackage[utf8]{inputenc}
\usepackage{booktabs}
\usepackage{amsmath}
\usepackage{natbib}

\title{My Document}
\author{Your Name}
\date{\today}

\begin{document}

\maketitle

\section{Introduction}
This is a sample document demonstrating various \LaTeX{} features including mathematical equations, tables, and formatting. You can cite papers using the citation search feature, and they will be automatically added to your \texttt{cite.bib} file.

\section{Mathematical Expressions}
Here are some example equations:

\subsection{Quadratic Formula}
The quadratic formula is given by:
\[x = \frac{-b \pm \sqrt{b^2 - 4ac}}{2a}\]

\subsection{Matrix Example}
A sample matrix:
\[
\begin{pmatrix}
a_{11} & a_{12} & a_{13} \\
a_{21} & a_{22} & a_{23} \\
a_{31} & a_{32} & a_{33}
\end{pmatrix}
\]

\section{Tables}
Below is a sample table showing data:

\begin{table}[h]
\centering
\begin{tabular}{@{}lrrr@{}}
\toprule
\textbf{Category} & \textbf{Value 1} & \textbf{Value 2} & \textbf{Total} \\
\midrule
Item A & 10.5 & 20.3 & 30.8 \\
Item B & 15.2 & 25.1 & 40.3 \\
Item C & 22.1 & 12.7 & 34.8 \\
\bottomrule
\end{tabular}
\caption{Sample Data Table}
\label{tab:sample}
\end{table}

\section{Lists}
\subsection{Bullet Points}
\begin{itemize}
  \item First item
  \item Second item
  \item Third item with sub-items:
    \begin{itemize}
      \item Sub-item 1
      \item Sub-item 2
    \end{itemize}
\end{itemize}

\subsection{Numbered List}
\begin{enumerate}
  \item First step
  \item Second step
  \item Third step
\end{enumerate}

\section{Citations and References}
When you use the citation search feature, references will be automatically managed. For example, you might cite a paper like this: \cite{example2024}. All citations will be automatically added to your \texttt{cite.bib} file.

\section{Conclusion}
This template demonstrates basic \LaTeX{} features including mathematical formulas, tables, and lists. The bibliography system is now set up to work seamlessly with Bibby AI's citation management. Feel free to modify it for your needs!

\bibliographystyle{plain}
\bibliography{cite}

\end{document}\documentclass{article}
\usepackage[utf8]{inputenc}
\usepackage{booktabs}
\usepackage{amsmath}
\usepackage{natbib}

\title{My Document}
\author{Your Name}
\date{\today}

\begin{document}

\maketitle

\section{Introduction}
This is a sample document demonstrating various \LaTeX{} features including mathematical equations, tables, and formatting. You can cite papers using the citation search feature, and they will be automatically added to your \texttt{cite.bib} file.

\section{Mathematical Expressions}
Here are some example equations:

\subsection{Quadratic Formula}
The quadratic formula is given by:
\[x = \frac{-b \pm \sqrt{b^2 - 4ac}}{2a}\]

\subsection{Matrix Example}
A sample matrix:
\[
\begin{pmatrix}
a_{11} & a_{12} & a_{13} \\
a_{21} & a_{22} & a_{23} \\
a_{31} & a_{32} & a_{33}
\end{pmatrix}
\]

\section{Tables}
Below is a sample table showing data:

\begin{table}[h]
\centering
\begin{tabular}{@{}lrrr@{}}
\toprule
\textbf{Category} & \textbf{Value 1} & \textbf{Value 2} & \textbf{Total} \\
\midrule
Item A & 10.5 & 20.3 & 30.8 \\
Item B & 15.2 & 25.1 & 40.3 \\
Item C & 22.1 & 12.7 & 34.8 \\
\bottomrule
\end{tabular}
\caption{Sample Data Table}
\label{tab:sample}
\end{table}

\section{Lists}
\subsection{Bullet Points}
\begin{itemize}
  \item First item
  \item Second item
  \item Third item with sub-items:
    \begin{itemize}
      \item Sub-item 1
      \item Sub-item 2
    \end{itemize}
\end{itemize}

\subsection{Numbered List}
\begin{enumerate}
  \item First step
  \item Second step
  \item Third step
\end{enumerate}

\section{Citations and References}
When you use the citation search feature, references will be automatically managed. For example, you might cite a paper like this: \cite{example2024}. All citations will be automatically added to your \texttt{cite.bib} file.

\section{Conclusion}
This template demonstrates basic \LaTeX{} features including mathematical formulas, tables, and lists. The bibliography system is now set up to work seamlessly with Bibby AI's citation management. Feel free to modify it for your needs!

\bibliographystyle{plain}
\bibliography{cite}

\end{document}\documentclass{article}
\usepackage[utf8]{inputenc}
\usepackage{booktabs}
\usepackage{amsmath}
\usepackage{natbib}

\title{My Document}
\author{Your Name}
\date{\today}

\begin{document}

\maketitle

\section{Introduction}
This is a sample document demonstrating various \LaTeX{} features including mathematical equations, tables, and formatting. You can cite papers using the citation search feature, and they will be automatically added to your \texttt{cite.bib} file.

\section{Mathematical Expressions}
Here are some example equations:

\subsection{Quadratic Formula}
The quadratic formula is given by:
\[x = \frac{-b \pm \sqrt{b^2 - 4ac}}{2a}\]

\subsection{Matrix Example}
A sample matrix:
\[
\begin{pmatrix}
a_{11} & a_{12} & a_{13} \\
a_{21} & a_{22} & a_{23} \\
a_{31} & a_{32} & a_{33}
\end{pmatrix}
\]

\section{Tables}
Below is a sample table showing data:

\begin{table}[h]
\centering
\begin{tabular}{@{}lrrr@{}}
\toprule
\textbf{Category} & \textbf{Value 1} & \textbf{Value 2} & \textbf{Total} \\
\midrule
Item A & 10.5 & 20.3 & 30.8 \\
Item B & 15.2 & 25.1 & 40.3 \\
Item C & 22.1 & 12.7 & 34.8 \\
\bottomrule
\end{tabular}
\caption{Sample Data Table}
\label{tab:sample}
\end{table}

\section{Lists}
\subsection{Bullet Points}
\begin{itemize}
  \item First item
  \item Second item
  \item Third item with sub-items:
    \begin{itemize}
      \item Sub-item 1
      \item Sub-item 2
    \end{itemize}
\end{itemize}

\subsection{Numbered List}
\begin{enumerate}
  \item First step
  \item Second step
  \item Third step
\end{enumerate}

\section{Citations and References}
When you use the citation search feature, references will be automatically managed. For example, you might cite a paper like this: \cite{example2024}. All citations will be automatically added to your \texttt{cite.bib} file.

\section{Conclusion}
This template demonstrates basic \LaTeX{} features including mathematical formulas, tables, and lists. The bibliography system is now set up to work seamlessly with Bibby AI's citation management. Feel free to modify it for your needs!

\bibliographystyle{plain}
\bibliography{cite}

\end{document}\documentclass{article}
\usepackage[utf8]{inputenc}
\usepackage{booktabs}
\usepackage{amsmath}
\usepackage{natbib}

\title{My Document}
\author{Your Name}
\date{\today}

\begin{document}

\maketitle

\section{Introduction}
This is a sample document demonstrating various \LaTeX{} features including mathematical equations, tables, and formatting. You can cite papers using the citation search feature, and they will be automatically added to your \texttt{cite.bib} file.

\section{Mathematical Expressions}
Here are some example equations:

\subsection{Quadratic Formula}
The quadratic formula is given by:
\[x = \frac{-b \pm \sqrt{b^2 - 4ac}}{2a}\]

\subsection{Matrix Example}
A sample matrix:
\[
\begin{pmatrix}
a_{11} & a_{12} & a_{13} \\
a_{21} & a_{22} & a_{23} \\
a_{31} & a_{32} & a_{33}
\end{pmatrix}
\]

\section{Tables}
Below is a sample table showing data:

\begin{table}[h]
\centering
\begin{tabular}{@{}lrrr@{}}
\toprule
\textbf{Category} & \textbf{Value 1} & \textbf{Value 2} & \textbf{Total} \\
\midrule
Item A & 10.5 & 20.3 & 30.8 \\
Item B & 15.2 & 25.1 & 40.3 \\
Item C & 22.1 & 12.7 & 34.8 \\
\bottomrule
\end{tabular}
\caption{Sample Data Table}
\label{tab:sample}
\end{table}

\section{Lists}
\subsection{Bullet Points}
\begin{itemize}
  \item First item
  \item Second item
  \item Third item with sub-items:
    \begin{itemize}
      \item Sub-item 1
      \item Sub-item 2
    \end{itemize}
\end{itemize}

\subsection{Numbered List}
\begin{enumerate}
  \item First step
  \item Second step
  \item Third step
\end{enumerate}

\section{Citations and References}
When you use the citation search feature, references will be automatically managed. For example, you might cite a paper like this: \cite{example2024}. All citations will be automatically added to your \texttt{cite.bib} file.

\section{Conclusion}
This template demonstrates basic \LaTeX{} features including mathematical formulas, tables, and lists. The bibliography system is now set up to work seamlessly with Bibby AI's citation management. Feel free to modify it for your needs!

\bibliographystyle{plain}
\bibliography{cite}

\end{document}\documentclass{article}
\usepackage[utf8]{inputenc}
\usepackage{booktabs}
\usepackage{amsmath}
\usepackage{natbib}

\title{My Document}
\author{Your Name}
\date{\today}

\begin{document}

\maketitle

\section{Introduction}
This is a sample document demonstrating various \LaTeX{} features including mathematical equations, tables, and formatting. You can cite papers using the citation search feature, and they will be automatically added to your \texttt{cite.bib} file.

\section{Mathematical Expressions}
Here are some example equations:

\subsection{Quadratic Formula}
The quadratic formula is given by:
\[x = \frac{-b \pm \sqrt{b^2 - 4ac}}{2a}\]

\subsection{Matrix Example}
A sample matrix:
\[
\begin{pmatrix}
a_{11} & a_{12} & a_{13} \\
a_{21} & a_{22} & a_{23} \\
a_{31} & a_{32} & a_{33}
\end{pmatrix}
\]

\section{Tables}
Below is a sample table showing data:

\begin{table}[h]
\centering
\begin{tabular}{@{}lrrr@{}}
\toprule
\textbf{Category} & \textbf{Value 1} & \textbf{Value 2} & \textbf{Total} \\
\midrule
Item A & 10.5 & 20.3 & 30.8 \\
Item B & 15.2 & 25.1 & 40.3 \\
Item C & 22.1 & 12.7 & 34.8 \\
\bottomrule
\end{tabular}
\caption{Sample Data Table}
\label{tab:sample}
\end{table}

\section{Lists}
\subsection{Bullet Points}
\begin{itemize}
  \item First item
  \item Second item
  \item Third item with sub-items:
    \begin{itemize}
      \item Sub-item 1
      \item Sub-item 2
    \end{itemize}
\end{itemize}

\subsection{Numbered List}
\begin{enumerate}
  \item First step
  \item Second step
  \item Third step
\end{enumerate}

\section{Citations and References}
When you use the citation search feature, references will be automatically managed. For example, you might cite a paper like this: \cite{example2024}. All citations will be automatically added to your \texttt{cite.bib} file.

\section{Conclusion}
This template demonstrates basic \LaTeX{} features including mathematical formulas, tables, and lists. The bibliography system is now set up to work seamlessly with Bibby AI's citation management. Feel free to modify it for your needs!

\bibliographystyle{plain}
\bibliography{cite}

\end{document}\documentclass{article}
\usepackage[utf8]{inputenc}
\usepackage{booktabs}
\usepackage{amsmath}
\usepackage{natbib}

\title{My Document}
\author{Your Name}
\date{\today}

\begin{document}

\maketitle

\section{Introduction}
This is a sample document demonstrating various \LaTeX{} features including mathematical equations, tables, and formatting. You can cite papers using the citation search feature, and they will be automatically added to your \texttt{cite.bib} file.

\section{Mathematical Expressions}
Here are some example equations:

\subsection{Quadratic Formula}
The quadratic formula is given by:
\[x = \frac{-b \pm \sqrt{b^2 - 4ac}}{2a}\]

\subsection{Matrix Example}
A sample matrix:
\[
\begin{pmatrix}
a_{11} & a_{12} & a_{13} \\
a_{21} & a_{22} & a_{23} \\
a_{31} & a_{32} & a_{33}
\end{pmatrix}
\]

\section{Tables}
Below is a sample table showing data:

\begin{table}[h]
\centering
\begin{tabular}{@{}lrrr@{}}
\toprule
\textbf{Category} & \textbf{Value 1} & \textbf{Value 2} & \textbf{Total} \\
\midrule
Item A & 10.5 & 20.3 & 30.8 \\
Item B & 15.2 & 25.1 & 40.3 \\
Item C & 22.1 & 12.7 & 34.8 \\
\bottomrule
\end{tabular}
\caption{Sample Data Table}
\label{tab:sample}
\end{table}

\section{Lists}
\subsection{Bullet Points}
\begin{itemize}
  \item First item
  \item Second item
  \item Third item with sub-items:
    \begin{itemize}
      \item Sub-item 1
      \item Sub-item 2
    \end{itemize}
\end{itemize}

\subsection{Numbered List}
\begin{enumerate}
  \item First step
  \item Second step
  \item Third step
\end{enumerate}

\section{Citations and References}
When you use the citation search feature, references will be automatically managed. For example, you might cite a paper like this: \cite{example2024}. All citations will be automatically added to your \texttt{cite.bib} file.

\section{Conclusion}
This template demonstrates basic \LaTeX{} features including mathematical formulas, tables, and lists. The bibliography system is now set up to work seamlessly with Bibby AI's citation management. Feel free to modify it for your needs!

\bibliographystyle{plain}
\bibliography{cite}

\end{document}\documentclass{article}
\usepackage[utf8]{inputenc}
\usepackage{booktabs}
\usepackage{amsmath}
\usepackage{natbib}

\title{My Document}
\author{Your Name}
\date{\today}

\begin{document}

\maketitle

\section{Introduction}
This is a sample document demonstrating various \LaTeX{} features including mathematical equations, tables, and formatting. You can cite papers using the citation search feature, and they will be automatically added to your \texttt{cite.bib} file.

\section{Mathematical Expressions}
Here are some example equations:

\subsection{Quadratic Formula}
The quadratic formula is given by:
\[x = \frac{-b \pm \sqrt{b^2 - 4ac}}{2a}\]

\subsection{Matrix Example}
A sample matrix:
\[
\begin{pmatrix}
a_{11} & a_{12} & a_{13} \\
a_{21} & a_{22} & a_{23} \\
a_{31} & a_{32} & a_{33}
\end{pmatrix}
\]

\section{Tables}
Below is a sample table showing data:

\begin{table}[h]
\centering
\begin{tabular}{@{}lrrr@{}}
\toprule
\textbf{Category} & \textbf{Value 1} & \textbf{Value 2} & \textbf{Total} \\
\midrule
Item A & 10.5 & 20.3 & 30.8 \\
Item B & 15.2 & 25.1 & 40.3 \\
Item C & 22.1 & 12.7 & 34.8 \\
\bottomrule
\end{tabular}
\caption{Sample Data Table}
\label{tab:sample}
\end{table}

\section{Lists}
\subsection{Bullet Points}
\begin{itemize}
  \item First item
  \item Second item
  \item Third item with sub-items:
    \begin{itemize}
      \item Sub-item 1
      \item Sub-item 2
    \end{itemize}
\end{itemize}

\subsection{Numbered List}
\begin{enumerate}
  \item First step
  \item Second step
  \item Third step
\end{enumerate}

\section{Citations and References}
When you use the citation search feature, references will be automatically managed. For example, you might cite a paper like this: \cite{example2024}. All citations will be automatically added to your \texttt{cite.bib} file.

\section{Conclusion}
This template demonstrates basic \LaTeX{} features including mathematical formulas, tables, and lists. The bibliography system is now set up to work seamlessly with Bibby AI's citation management. Feel free to modify it for your needs!

\bibliographystyle{plain}
\bibliography{cite}

\end{document}\documentclass{article}
\usepackage[utf8]{inputenc}
\usepackage{booktabs}
\usepackage{amsmath}
\usepackage{natbib}

\title{My Document}
\author{Your Name}
\date{\today}

\begin{document}

\maketitle

\section{Introduction}
This is a sample document demonstrating various \LaTeX{} features including mathematical equations, tables, and formatting. You can cite papers using the citation search feature, and they will be automatically added to your \texttt{cite.bib} file.

\section{Mathematical Expressions}
Here are some example equations:

\subsection{Quadratic Formula}
The quadratic formula is given by:
\[x = \frac{-b \pm \sqrt{b^2 - 4ac}}{2a}\]

\subsection{Matrix Example}
A sample matrix:
\[
\begin{pmatrix}
a_{11} & a_{12} & a_{13} \\
a_{21} & a_{22} & a_{23} \\
a_{31} & a_{32} & a_{33}
\end{pmatrix}
\]

\section{Tables}
Below is a sample table showing data:

\begin{table}[h]
\centering
\begin{tabular}{@{}lrrr@{}}
\toprule
\textbf{Category} & \textbf{Value 1} & \textbf{Value 2} & \textbf{Total} \\
\midrule
Item A & 10.5 & 20.3 & 30.8 \\
Item B & 15.2 & 25.1 & 40.3 \\
Item C & 22.1 & 12.7 & 34.8 \\
\bottomrule
\end{tabular}
\caption{Sample Data Table}
\label{tab:sample}
\end{table}

\section{Lists}
\subsection{Bullet Points}
\begin{itemize}
  \item First item
  \item Second item
  \item Third item with sub-items:
    \begin{itemize}
      \item Sub-item 1
      \item Sub-item 2
    \end{itemize}
\end{itemize}

\subsection{Numbered List}
\begin{enumerate}
  \item First step
  \item Second step
  \item Third step
\end{enumerate}

\section{Citations and References}
When you use the citation search feature, references will be automatically managed. For example, you might cite a paper like this: \cite{example2024}. All citations will be automatically added to your \texttt{cite.bib} file.

\section{Conclusion}
This template demonstrates basic \LaTeX{} features including mathematical formulas, tables, and lists. The bibliography system is now set up to work seamlessly with Bibby AI's citation management. Feel free to modify it for your needs!

\bibliographystyle{plain}
\bibliography{cite}

\end{document}\documentclass{article}
\usepackage[utf8]{inputenc}
\usepackage{booktabs}
\usepackage{amsmath}
\usepackage{natbib}

\title{My Document}
\author{Your Name}
\date{\today}

\begin{document}

\maketitle

\section{Introduction}
This is a sample document demonstrating various \LaTeX{} features including mathematical equations, tables, and formatting. You can cite papers using the citation search feature, and they will be automatically added to your \texttt{cite.bib} file.

\section{Mathematical Expressions}
Here are some example equations:

\subsection{Quadratic Formula}
The quadratic formula is given by:
\[x = \frac{-b \pm \sqrt{b^2 - 4ac}}{2a}\]

\subsection{Matrix Example}
A sample matrix:
\[
\begin{pmatrix}
a_{11} & a_{12} & a_{13} \\
a_{21} & a_{22} & a_{23} \\
a_{31} & a_{32} & a_{33}
\end{pmatrix}
\]

\section{Tables}
Below is a sample table showing data:

\begin{table}[h]
\centering
\begin{tabular}{@{}lrrr@{}}
\toprule
\textbf{Category} & \textbf{Value 1} & \textbf{Value 2} & \textbf{Total} \\
\midrule
Item A & 10.5 & 20.3 & 30.8 \\
Item B & 15.2 & 25.1 & 40.3 \\
Item C & 22.1 & 12.7 & 34.8 \\
\bottomrule
\end{tabular}
\caption{Sample Data Table}
\label{tab:sample}
\end{table}

\section{Lists}
\subsection{Bullet Points}
\begin{itemize}
  \item First item
  \item Second item
  \item Third item with sub-items:
    \begin{itemize}
      \item Sub-item 1
      \item Sub-item 2
    \end{itemize}
\end{itemize}

\subsection{Numbered List}
\begin{enumerate}
  \item First step
  \item Second step
  \item Third step
\end{enumerate}

\section{Citations and References}
When you use the citation search feature, references will be automatically managed. For example, you might cite a paper like this: \cite{example2024}. All citations will be automatically added to your \texttt{cite.bib} file.

\section{Conclusion}
This template demonstrates basic \LaTeX{} features including mathematical formulas, tables, and lists. The bibliography system is now set up to work seamlessly with Bibby AI's citation management. Feel free to modify it for your needs!

\bibliographystyle{plain}
\bibliography{cite}

\end{document}\documentclass{article}
\usepackage[utf8]{inputenc}
\usepackage{booktabs}
\usepackage{amsmath}
\usepackage{natbib}

\title{My Document}
\author{Your Name}
\date{\today}

\begin{document}

\maketitle

\section{Introduction}
This is a sample document demonstrating various \LaTeX{} features including mathematical equations, tables, and formatting. You can cite papers using the citation search feature, and they will be automatically added to your \texttt{cite.bib} file.

\section{Mathematical Expressions}
Here are some example equations:

\subsection{Quadratic Formula}
The quadratic formula is given by:
\[x = \frac{-b \pm \sqrt{b^2 - 4ac}}{2a}\]

\subsection{Matrix Example}
A sample matrix:
\[
\begin{pmatrix}
a_{11} & a_{12} & a_{13} \\
a_{21} & a_{22} & a_{23} \\
a_{31} & a_{32} & a_{33}
\end{pmatrix}
\]

\section{Tables}
Below is a sample table showing data:

\begin{table}[h]
\centering
\begin{tabular}{@{}lrrr@{}}
\toprule
\textbf{Category} & \textbf{Value 1} & \textbf{Value 2} & \textbf{Total} \\
\midrule
Item A & 10.5 & 20.3 & 30.8 \\
Item B & 15.2 & 25.1 & 40.3 \\
Item C & 22.1 & 12.7 & 34.8 \\
\bottomrule
\end{tabular}
\caption{Sample Data Table}
\label{tab:sample}
\end{table}

\section{Lists}
\subsection{Bullet Points}
\begin{itemize}
  \item First item
  \item Second item
  \item Third item with sub-items:
    \begin{itemize}
      \item Sub-item 1
      \item Sub-item 2
    \end{itemize}
\end{itemize}

\subsection{Numbered List}
\begin{enumerate}
  \item First step
  \item Second step
  \item Third step
\end{enumerate}

\section{Citations and References}
When you use the citation search feature, references will be automatically managed. For example, you might cite a paper like this: \cite{example2024}. All citations will be automatically added to your \texttt{cite.bib} file.

\section{Conclusion}
This template demonstrates basic \LaTeX{} features including mathematical formulas, tables, and lists. The bibliography system is now set up to work seamlessly with Bibby AI's citation management. Feel free to modify it for your needs!

\bibliographystyle{plain}
\bibliography{cite}

\end{document}\documentclass{article}
\usepackage[utf8]{inputenc}
\usepackage{booktabs}
\usepackage{amsmath}
\usepackage{natbib}

\title{My Document}
\author{Your Name}
\date{\today}

\begin{document}

\maketitle

\section{Introduction}
This is a sample document demonstrating various \LaTeX{} features including mathematical equations, tables, and formatting. You can cite papers using the citation search feature, and they will be automatically added to your \texttt{cite.bib} file.

\section{Mathematical Expressions}
Here are some example equations:

\subsection{Quadratic Formula}
The quadratic formula is given by:
\[x = \frac{-b \pm \sqrt{b^2 - 4ac}}{2a}\]

\subsection{Matrix Example}
A sample matrix:
\[
\begin{pmatrix}
a_{11} & a_{12} & a_{13} \\
a_{21} & a_{22} & a_{23} \\
a_{31} & a_{32} & a_{33}
\end{pmatrix}
\]

\section{Tables}
Below is a sample table showing data:

\begin{table}[h]
\centering
\begin{tabular}{@{}lrrr@{}}
\toprule
\textbf{Category} & \textbf{Value 1} & \textbf{Value 2} & \textbf{Total} \\
\midrule
Item A & 10.5 & 20.3 & 30.8 \\
Item B & 15.2 & 25.1 & 40.3 \\
Item C & 22.1 & 12.7 & 34.8 \\
\bottomrule
\end{tabular}
\caption{Sample Data Table}
\label{tab:sample}
\end{table}

\section{Lists}
\subsection{Bullet Points}
\begin{itemize}
  \item First item
  \item Second item
  \item Third item with sub-items:
    \begin{itemize}
      \item Sub-item 1
      \item Sub-item 2
    \end{itemize}
\end{itemize}

\subsection{Numbered List}
\begin{enumerate}
  \item First step
  \item Second step
  \item Third step
\end{enumerate}

\section{Citations and References}
When you use the citation search feature, references will be automatically managed. For example, you might cite a paper like this: \cite{example2024}. All citations will be automatically added to your \texttt{cite.bib} file.

\section{Conclusion}
This template demonstrates basic \LaTeX{} features including mathematical formulas, tables, and lists. The bibliography system is now set up to work seamlessly with Bibby AI's citation management. Feel free to modify it for your needs!

\bibliographystyle{plain}
\bibliography{cite}

\end{document}\documentclass{article}
\usepackage[utf8]{inputenc}
\usepackage{booktabs}
\usepackage{amsmath}
\usepackage{natbib}

\title{My Document}
\author{Your Name}
\date{\today}

\begin{document}

\maketitle

\section{Introduction}
This is a sample document demonstrating various \LaTeX{} features including mathematical equations, tables, and formatting. You can cite papers using the citation search feature, and they will be automatically added to your \texttt{cite.bib} file.

\section{Mathematical Expressions}
Here are some example equations:

\subsection{Quadratic Formula}
The quadratic formula is given by:
\[x = \frac{-b \pm \sqrt{b^2 - 4ac}}{2a}\]

\subsection{Matrix Example}
A sample matrix:
\[
\begin{pmatrix}
a_{11} & a_{12} & a_{13} \\
a_{21} & a_{22} & a_{23} \\
a_{31} & a_{32} & a_{33}
\end{pmatrix}
\]

\section{Tables}
Below is a sample table showing data:

\begin{table}[h]
\centering
\begin{tabular}{@{}lrrr@{}}
\toprule
\textbf{Category} & \textbf{Value 1} & \textbf{Value 2} & \textbf{Total} \\
\midrule
Item A & 10.5 & 20.3 & 30.8 \\
Item B & 15.2 & 25.1 & 40.3 \\
Item C & 22.1 & 12.7 & 34.8 \\
\bottomrule
\end{tabular}
\caption{Sample Data Table}
\label{tab:sample}
\end{table}

\section{Lists}
\subsection{Bullet Points}
\begin{itemize}
  \item First item
  \item Second item
  \item Third item with sub-items:
    \begin{itemize}
      \item Sub-item 1
      \item Sub-item 2
    \end{itemize}
\end{itemize}

\subsection{Numbered List}
\begin{enumerate}
  \item First step
  \item Second step
  \item Third step
\end{enumerate}

\section{Citations and References}
When you use the citation search feature, references will be automatically managed. For example, you might cite a paper like this: \cite{example2024}. All citations will be automatically added to your \texttt{cite.bib} file.

\section{Conclusion}
This template demonstrates basic \LaTeX{} features including mathematical formulas, tables, and lists. The bibliography system is now set up to work seamlessly with Bibby AI's citation management. Feel free to modify it for your needs!

\bibliographystyle{plain}
\bibliography{cite}

\end{document}\documentclass{article}
\usepackage[utf8]{inputenc}
\usepackage{booktabs}
\usepackage{amsmath}
\usepackage{natbib}

\title{My Document}
\author{Your Name}
\date{\today}

\begin{document}

\maketitle

\section{Introduction}
This is a sample document demonstrating various \LaTeX{} features including mathematical equations, tables, and formatting. You can cite papers using the citation search feature, and they will be automatically added to your \texttt{cite.bib} file.

\section{Mathematical Expressions}
Here are some example equations:

\subsection{Quadratic Formula}
The quadratic formula is given by:
\[x = \frac{-b \pm \sqrt{b^2 - 4ac}}{2a}\]

\subsection{Matrix Example}
A sample matrix:
\[
\begin{pmatrix}
a_{11} & a_{12} & a_{13} \\
a_{21} & a_{22} & a_{23} \\
a_{31} & a_{32} & a_{33}
\end{pmatrix}
\]

\section{Tables}
Below is a sample table showing data:

\begin{table}[h]
\centering
\begin{tabular}{@{}lrrr@{}}
\toprule
\textbf{Category} & \textbf{Value 1} & \textbf{Value 2} & \textbf{Total} \\
\midrule
Item A & 10.5 & 20.3 & 30.8 \\
Item B & 15.2 & 25.1 & 40.3 \\
Item C & 22.1 & 12.7 & 34.8 \\
\bottomrule
\end{tabular}
\caption{Sample Data Table}
\label{tab:sample}
\end{table}

\section{Lists}
\subsection{Bullet Points}
\begin{itemize}
  \item First item
  \item Second item
  \item Third item with sub-items:
    \begin{itemize}
      \item Sub-item 1
      \item Sub-item 2
    \end{itemize}
\end{itemize}

\subsection{Numbered List}
\begin{enumerate}
  \item First step
  \item Second step
  \item Third step
\end{enumerate}

\section{Citations and References}
When you use the citation search feature, references will be automatically managed. For example, you might cite a paper like this: \cite{example2024}. All citations will be automatically added to your \texttt{cite.bib} file.

\section{Conclusion}
This template demonstrates basic \LaTeX{} features including mathematical formulas, tables, and lists. The bibliography system is now set up to work seamlessly with Bibby AI's citation management. Feel free to modify it for your needs!

\bibliographystyle{plain}
\bibliography{cite}

\end{document}\documentclass{article}
\usepackage[utf8]{inputenc}
\usepackage{booktabs}
\usepackage{amsmath}
\usepackage{natbib}

\title{My Document}
\author{Your Name}
\date{\today}

\begin{document}

\maketitle

\section{Introduction}
This is a sample document demonstrating various \LaTeX{} features including mathematical equations, tables, and formatting. You can cite papers using the citation search feature, and they will be automatically added to your \texttt{cite.bib} file.

\section{Mathematical Expressions}
Here are some example equations:

\subsection{Quadratic Formula}
The quadratic formula is given by:
\[x = \frac{-b \pm \sqrt{b^2 - 4ac}}{2a}\]

\subsection{Matrix Example}
A sample matrix:
\[
\begin{pmatrix}
a_{11} & a_{12} & a_{13} \\
a_{21} & a_{22} & a_{23} \\
a_{31} & a_{32} & a_{33}
\end{pmatrix}
\]

\section{Tables}
Below is a sample table showing data:

\begin{table}[h]
\centering
\begin{tabular}{@{}lrrr@{}}
\toprule
\textbf{Category} & \textbf{Value 1} & \textbf{Value 2} & \textbf{Total} \\
\midrule
Item A & 10.5 & 20.3 & 30.8 \\
Item B & 15.2 & 25.1 & 40.3 \\
Item C & 22.1 & 12.7 & 34.8 \\
\bottomrule
\end{tabular}
\caption{Sample Data Table}
\label{tab:sample}
\end{table}

\section{Lists}
\subsection{Bullet Points}
\begin{itemize}
  \item First item
  \item Second item
  \item Third item with sub-items:
    \begin{itemize}
      \item Sub-item 1
      \item Sub-item 2
    \end{itemize}
\end{itemize}

\subsection{Numbered List}
\begin{enumerate}
  \item First step
  \item Second step
  \item Third step
\end{enumerate}

\section{Citations and References}
When you use the citation search feature, references will be automatically managed. For example, you might cite a paper like this: \cite{example2024}. All citations will be automatically added to your \texttt{cite.bib} file.

\section{Conclusion}
This template demonstrates basic \LaTeX{} features including mathematical formulas, tables, and lists. The bibliography system is now set up to work seamlessly with Bibby AI's citation management. Feel free to modify it for your needs!

\bibliographystyle{plain}
\bibliography{cite}

\end{document}\documentclass{article}
\usepackage[utf8]{inputenc}
\usepackage{booktabs}
\usepackage{amsmath}
\usepackage{natbib}

\title{My Document}
\author{Your Name}
\date{\today}

\begin{document}

\maketitle

\section{Introduction}
This is a sample document demonstrating various \LaTeX{} features including mathematical equations, tables, and formatting. You can cite papers using the citation search feature, and they will be automatically added to your \texttt{cite.bib} file.

\section{Mathematical Expressions}
Here are some example equations:

\subsection{Quadratic Formula}
The quadratic formula is given by:
\[x = \frac{-b \pm \sqrt{b^2 - 4ac}}{2a}\]

\subsection{Matrix Example}
A sample matrix:
\[
\begin{pmatrix}
a_{11} & a_{12} & a_{13} \\
a_{21} & a_{22} & a_{23} \\
a_{31} & a_{32} & a_{33}
\end{pmatrix}
\]

\section{Tables}
Below is a sample table showing data:

\begin{table}[h]
\centering
\begin{tabular}{@{}lrrr@{}}
\toprule
\textbf{Category} & \textbf{Value 1} & \textbf{Value 2} & \textbf{Total} \\
\midrule
Item A & 10.5 & 20.3 & 30.8 \\
Item B & 15.2 & 25.1 & 40.3 \\
Item C & 22.1 & 12.7 & 34.8 \\
\bottomrule
\end{tabular}
\caption{Sample Data Table}
\label{tab:sample}
\end{table}

\section{Lists}
\subsection{Bullet Points}
\begin{itemize}
  \item First item
  \item Second item
  \item Third item with sub-items:
    \begin{itemize}
      \item Sub-item 1
      \item Sub-item 2
    \end{itemize}
\end{itemize}

\subsection{Numbered List}
\begin{enumerate}
  \item First step
  \item Second step
  \item Third step
\end{enumerate}

\section{Citations and References}
When you use the citation search feature, references will be automatically managed. For example, you might cite a paper like this: \cite{example2024}. All citations will be automatically added to your \texttt{cite.bib} file.

\section{Conclusion}
This template demonstrates basic \LaTeX{} features including mathematical formulas, tables, and lists. The bibliography system is now set up to work seamlessly with Bibby AI's citation management. Feel free to modify it for your needs!

\bibliographystyle{plain}
\bibliography{cite}

\end{document}\documentclass{article}
\usepackage[utf8]{inputenc}
\usepackage{booktabs}
\usepackage{amsmath}
\usepackage{natbib}

\title{My Document}
\author{Your Name}
\date{\today}

\begin{document}

\maketitle

\section{Introduction}
This is a sample document demonstrating various \LaTeX{} features including mathematical equations, tables, and formatting. You can cite papers using the citation search feature, and they will be automatically added to your \texttt{cite.bib} file.

\section{Mathematical Expressions}
Here are some example equations:

\subsection{Quadratic Formula}
The quadratic formula is given by:
\[x = \frac{-b \pm \sqrt{b^2 - 4ac}}{2a}\]

\subsection{Matrix Example}
A sample matrix:
\[
\begin{pmatrix}
a_{11} & a_{12} & a_{13} \\
a_{21} & a_{22} & a_{23} \\
a_{31} & a_{32} & a_{33}
\end{pmatrix}
\]

\section{Tables}
Below is a sample table showing data:

\begin{table}[h]
\centering
\begin{tabular}{@{}lrrr@{}}
\toprule
\textbf{Category} & \textbf{Value 1} & \textbf{Value 2} & \textbf{Total} \\
\midrule
Item A & 10.5 & 20.3 & 30.8 \\
Item B & 15.2 & 25.1 & 40.3 \\
Item C & 22.1 & 12.7 & 34.8 \\
\bottomrule
\end{tabular}
\caption{Sample Data Table}
\label{tab:sample}
\end{table}

\section{Lists}
\subsection{Bullet Points}
\begin{itemize}
  \item First item
  \item Second item
  \item Third item with sub-items:
    \begin{itemize}
      \item Sub-item 1
      \item Sub-item 2
    \end{itemize}
\end{itemize}

\subsection{Numbered List}
\begin{enumerate}
  \item First step
  \item Second step
  \item Third step
\end{enumerate}

\section{Citations and References}
When you use the citation search feature, references will be automatically managed. For example, you might cite a paper like this: \cite{example2024}. All citations will be automatically added to your \texttt{cite.bib} file.

\section{Conclusion}
This template demonstrates basic \LaTeX{} features including mathematical formulas, tables, and lists. The bibliography system is now set up to work seamlessly with Bibby AI's citation management. Feel free to modify it for your needs!

\bibliographystyle{plain}
\bibliography{cite}

\end{document}\documentclass{article}
\usepackage[utf8]{inputenc}
\usepackage{booktabs}
\usepackage{amsmath}
\usepackage{natbib}

\title{My Document}
\author{Your Name}
\date{\today}

\begin{document}

\maketitle

\section{Introduction}
This is a sample document demonstrating various \LaTeX{} features including mathematical equations, tables, and formatting. You can cite papers using the citation search feature, and they will be automatically added to your \texttt{cite.bib} file.

\section{Mathematical Expressions}
Here are some example equations:

\subsection{Quadratic Formula}
The quadratic formula is given by:
\[x = \frac{-b \pm \sqrt{b^2 - 4ac}}{2a}\]

\subsection{Matrix Example}
A sample matrix:
\[
\begin{pmatrix}
a_{11} & a_{12} & a_{13} \\
a_{21} & a_{22} & a_{23} \\
a_{31} & a_{32} & a_{33}
\end{pmatrix}
\]

\section{Tables}
Below is a sample table showing data:

\begin{table}[h]
\centering
\begin{tabular}{@{}lrrr@{}}
\toprule
\textbf{Category} & \textbf{Value 1} & \textbf{Value 2} & \textbf{Total} \\
\midrule
Item A & 10.5 & 20.3 & 30.8 \\
Item B & 15.2 & 25.1 & 40.3 \\
Item C & 22.1 & 12.7 & 34.8 \\
\bottomrule
\end{tabular}
\caption{Sample Data Table}
\label{tab:sample}
\end{table}

\section{Lists}
\subsection{Bullet Points}
\begin{itemize}
  \item First item
  \item Second item
  \item Third item with sub-items:
    \begin{itemize}
      \item Sub-item 1
      \item Sub-item 2
    \end{itemize}
\end{itemize}

\subsection{Numbered List}
\begin{enumerate}
  \item First step
  \item Second step
  \item Third step
\end{enumerate}

\section{Citations and References}
When you use the citation search feature, references will be automatically managed. For example, you might cite a paper like this: \cite{example2024}. All citations will be automatically added to your \texttt{cite.bib} file.

\section{Conclusion}
This template demonstrates basic \LaTeX{} features including mathematical formulas, tables, and lists. The bibliography system is now set up to work seamlessly with Bibby AI's citation management. Feel free to modify it for your needs!

\bibliographystyle{plain}
\bibliography{cite}

\end{document}\documentclass{article}
\usepackage[utf8]{inputenc}
\usepackage{booktabs}
\usepackage{amsmath}
\usepackage{natbib}

\title{My Document}
\author{Your Name}
\date{\today}

\begin{document}

\maketitle

\section{Introduction}
This is a sample document demonstrating various \LaTeX{} features including mathematical equations, tables, and formatting. You can cite papers using the citation search feature, and they will be automatically added to your \texttt{cite.bib} file.

\section{Mathematical Expressions}
Here are some example equations:

\subsection{Quadratic Formula}
The quadratic formula is given by:
\[x = \frac{-b \pm \sqrt{b^2 - 4ac}}{2a}\]

\subsection{Matrix Example}
A sample matrix:
\[
\begin{pmatrix}
a_{11} & a_{12} & a_{13} \\
a_{21} & a_{22} & a_{23} \\
a_{31} & a_{32} & a_{33}
\end{pmatrix}
\]

\section{Tables}
Below is a sample table showing data:

\begin{table}[h]
\centering
\begin{tabular}{@{}lrrr@{}}
\toprule
\textbf{Category} & \textbf{Value 1} & \textbf{Value 2} & \textbf{Total} \\
\midrule
Item A & 10.5 & 20.3 & 30.8 \\
Item B & 15.2 & 25.1 & 40.3 \\
Item C & 22.1 & 12.7 & 34.8 \\
\bottomrule
\end{tabular}
\caption{Sample Data Table}
\label{tab:sample}
\end{table}

\section{Lists}
\subsection{Bullet Points}
\begin{itemize}
  \item First item
  \item Second item
  \item Third item with sub-items:
    \begin{itemize}
      \item Sub-item 1
      \item Sub-item 2
    \end{itemize}
\end{itemize}

\subsection{Numbered List}
\begin{enumerate}
  \item First step
  \item Second step
  \item Third step
\end{enumerate}

\section{Citations and References}
When you use the citation search feature, references will be automatically managed. For example, you might cite a paper like this: \cite{example2024}. All citations will be automatically added to your \texttt{cite.bib} file.

\section{Conclusion}
This template demonstrates basic \LaTeX{} features including mathematical formulas, tables, and lists. The bibliography system is now set up to work seamlessly with Bibby AI's citation management. Feel free to modify it for your needs!

\bibliographystyle{plain}
\bibliography{cite}

\end{document}\documentclass{article}
\usepackage[utf8]{inputenc}
\usepackage{booktabs}
\usepackage{amsmath}
\usepackage{natbib}

\title{My Document}
\author{Your Name}
\date{\today}

\begin{document}

\maketitle

\section{Introduction}
This is a sample document demonstrating various \LaTeX{} features including mathematical equations, tables, and formatting. You can cite papers using the citation search feature, and they will be automatically added to your \texttt{cite.bib} file.

\section{Mathematical Expressions}
Here are some example equations:

\subsection{Quadratic Formula}
The quadratic formula is given by:
\[x = \frac{-b \pm \sqrt{b^2 - 4ac}}{2a}\]

\subsection{Matrix Example}
A sample matrix:
\[
\begin{pmatrix}
a_{11} & a_{12} & a_{13} \\
a_{21} & a_{22} & a_{23} \\
a_{31} & a_{32} & a_{33}
\end{pmatrix}
\]

\section{Tables}
Below is a sample table showing data:

\begin{table}[h]
\centering
\begin{tabular}{@{}lrrr@{}}
\toprule
\textbf{Category} & \textbf{Value 1} & \textbf{Value 2} & \textbf{Total} \\
\midrule
Item A & 10.5 & 20.3 & 30.8 \\
Item B & 15.2 & 25.1 & 40.3 \\
Item C & 22.1 & 12.7 & 34.8 \\
\bottomrule
\end{tabular}
\caption{Sample Data Table}
\label{tab:sample}
\end{table}

\section{Lists}
\subsection{Bullet Points}
\begin{itemize}
  \item First item
  \item Second item
  \item Third item with sub-items:
    \begin{itemize}
      \item Sub-item 1
      \item Sub-item 2
    \end{itemize}
\end{itemize}

\subsection{Numbered List}
\begin{enumerate}
  \item First step
  \item Second step
  \item Third step
\end{enumerate}

\section{Citations and References}
When you use the citation search feature, references will be automatically managed. For example, you might cite a paper like this: \cite{example2024}. All citations will be automatically added to your \texttt{cite.bib} file.

\section{Conclusion}
This template demonstrates basic \LaTeX{} features including mathematical formulas, tables, and lists. The bibliography system is now set up to work seamlessly with Bibby AI's citation management. Feel free to modify it for your needs!

\bibliographystyle{plain}
\bibliography{cite}

\end{document}\documentclass{article}
\usepackage[utf8]{inputenc}
\usepackage{booktabs}
\usepackage{amsmath}
\usepackage{natbib}

\title{My Document}
\author{Your Name}
\date{\today}

\begin{document}

\maketitle

\section{Introduction}
This is a sample document demonstrating various \LaTeX{} features including mathematical equations, tables, and formatting. You can cite papers using the citation search feature, and they will be automatically added to your \texttt{cite.bib} file.

\section{Mathematical Expressions}
Here are some example equations:

\subsection{Quadratic Formula}
The quadratic formula is given by:
\[x = \frac{-b \pm \sqrt{b^2 - 4ac}}{2a}\]

\subsection{Matrix Example}
A sample matrix:
\[
\begin{pmatrix}
a_{11} & a_{12} & a_{13} \\
a_{21} & a_{22} & a_{23} \\
a_{31} & a_{32} & a_{33}
\end{pmatrix}
\]

\section{Tables}
Below is a sample table showing data:

\begin{table}[h]
\centering
\begin{tabular}{@{}lrrr@{}}
\toprule
\textbf{Category} & \textbf{Value 1} & \textbf{Value 2} & \textbf{Total} \\
\midrule
Item A & 10.5 & 20.3 & 30.8 \\
Item B & 15.2 & 25.1 & 40.3 \\
Item C & 22.1 & 12.7 & 34.8 \\
\bottomrule
\end{tabular}
\caption{Sample Data Table}
\label{tab:sample}
\end{table}

\section{Lists}
\subsection{Bullet Points}
\begin{itemize}
  \item First item
  \item Second item
  \item Third item with sub-items:
    \begin{itemize}
      \item Sub-item 1
      \item Sub-item 2
    \end{itemize}
\end{itemize}

\subsection{Numbered List}
\begin{enumerate}
  \item First step
  \item Second step
  \item Third step
\end{enumerate}

\section{Citations and References}
When you use the citation search feature, references will be automatically managed. For example, you might cite a paper like this: \cite{example2024}. All citations will be automatically added to your \texttt{cite.bib} file.

\section{Conclusion}
This template demonstrates basic \LaTeX{} features including mathematical formulas, tables, and lists. The bibliography system is now set up to work seamlessly with Bibby AI's citation management. Feel free to modify it for your needs!

\bibliographystyle{plain}
\bibliography{cite}

\end{document}\documentclass{article}
\usepackage[utf8]{inputenc}
\usepackage{booktabs}
\usepackage{amsmath}
\usepackage{natbib}

\title{My Document}
\author{Your Name}
\date{\today}

\begin{document}

\maketitle

\section{Introduction}
This is a sample document demonstrating various \LaTeX{} features including mathematical equations, tables, and formatting. You can cite papers using the citation search feature, and they will be automatically added to your \texttt{cite.bib} file.

\section{Mathematical Expressions}
Here are some example equations:

\subsection{Quadratic Formula}
The quadratic formula is given by:
\[x = \frac{-b \pm \sqrt{b^2 - 4ac}}{2a}\]

\subsection{Matrix Example}
A sample matrix:
\[
\begin{pmatrix}
a_{11} & a_{12} & a_{13} \\
a_{21} & a_{22} & a_{23} \\
a_{31} & a_{32} & a_{33}
\end{pmatrix}
\]

\section{Tables}
Below is a sample table showing data:

\begin{table}[h]
\centering
\begin{tabular}{@{}lrrr@{}}
\toprule
\textbf{Category} & \textbf{Value 1} & \textbf{Value 2} & \textbf{Total} \\
\midrule
Item A & 10.5 & 20.3 & 30.8 \\
Item B & 15.2 & 25.1 & 40.3 \\
Item C & 22.1 & 12.7 & 34.8 \\
\bottomrule
\end{tabular}
\caption{Sample Data Table}
\label{tab:sample}
\end{table}

\section{Lists}
\subsection{Bullet Points}
\begin{itemize}
  \item First item
  \item Second item
  \item Third item with sub-items:
    \begin{itemize}
      \item Sub-item 1
      \item Sub-item 2
    \end{itemize}
\end{itemize}

\subsection{Numbered List}
\begin{enumerate}
  \item First step
  \item Second step
  \item Third step
\end{enumerate}

\section{Citations and References}
When you use the citation search feature, references will be automatically managed. For example, you might cite a paper like this: \cite{example2024}. All citations will be automatically added to your \texttt{cite.bib} file.

\section{Conclusion}
This template demonstrates basic \LaTeX{} features including mathematical formulas, tables, and lists. The bibliography system is now set up to work seamlessly with Bibby AI's citation management. Feel free to modify it for your needs!

\bibliographystyle{plain}
\bibliography{cite}

\end{document}\documentclass{article}
\usepackage[utf8]{inputenc}
\usepackage{booktabs}
\usepackage{amsmath}
\usepackage{natbib}

\title{My Document}
\author{Your Name}
\date{\today}

\begin{document}

\maketitle

\section{Introduction}
This is a sample document demonstrating various \LaTeX{} features including mathematical equations, tables, and formatting. You can cite papers using the citation search feature, and they will be automatically added to your \texttt{cite.bib} file.

\section{Mathematical Expressions}
Here are some example equations:

\subsection{Quadratic Formula}
The quadratic formula is given by:
\[x = \frac{-b \pm \sqrt{b^2 - 4ac}}{2a}\]

\subsection{Matrix Example}
A sample matrix:
\[
\begin{pmatrix}
a_{11} & a_{12} & a_{13} \\
a_{21} & a_{22} & a_{23} \\
a_{31} & a_{32} & a_{33}
\end{pmatrix}
\]

\section{Tables}
Below is a sample table showing data:

\begin{table}[h]
\centering
\begin{tabular}{@{}lrrr@{}}
\toprule
\textbf{Category} & \textbf{Value 1} & \textbf{Value 2} & \textbf{Total} \\
\midrule
Item A & 10.5 & 20.3 & 30.8 \\
Item B & 15.2 & 25.1 & 40.3 \\
Item C & 22.1 & 12.7 & 34.8 \\
\bottomrule
\end{tabular}
\caption{Sample Data Table}
\label{tab:sample}
\end{table}

\section{Lists}
\subsection{Bullet Points}
\begin{itemize}
  \item First item
  \item Second item
  \item Third item with sub-items:
    \begin{itemize}
      \item Sub-item 1
      \item Sub-item 2
    \end{itemize}
\end{itemize}

\subsection{Numbered List}
\begin{enumerate}
  \item First step
  \item Second step
  \item Third step
\end{enumerate}

\section{Citations and References}
When you use the citation search feature, references will be automatically managed. For example, you might cite a paper like this: \cite{example2024}. All citations will be automatically added to your \texttt{cite.bib} file.

\section{Conclusion}
This template demonstrates basic \LaTeX{} features including mathematical formulas, tables, and lists. The bibliography system is now set up to work seamlessly with Bibby AI's citation management. Feel free to modify it for your needs!

\bibliographystyle{plain}
\bibliography{cite}

\end{document}\documentclass{article}
\usepackage[utf8]{inputenc}
\usepackage{booktabs}
\usepackage{amsmath}
\usepackage{natbib}

\title{My Document}
\author{Your Name}
\date{\today}

\begin{document}

\maketitle

\section{Introduction}
This is a sample document demonstrating various \LaTeX{} features including mathematical equations, tables, and formatting. You can cite papers using the citation search feature, and they will be automatically added to your \texttt{cite.bib} file.

\section{Mathematical Expressions}
Here are some example equations:

\subsection{Quadratic Formula}
The quadratic formula is given by:
\[x = \frac{-b \pm \sqrt{b^2 - 4ac}}{2a}\]

\subsection{Matrix Example}
A sample matrix:
\[
\begin{pmatrix}
a_{11} & a_{12} & a_{13} \\
a_{21} & a_{22} & a_{23} \\
a_{31} & a_{32} & a_{33}
\end{pmatrix}
\]

\section{Tables}
Below is a sample table showing data:

\begin{table}[h]
\centering
\begin{tabular}{@{}lrrr@{}}
\toprule
\textbf{Category} & \textbf{Value 1} & \textbf{Value 2} & \textbf{Total} \\
\midrule
Item A & 10.5 & 20.3 & 30.8 \\
Item B & 15.2 & 25.1 & 40.3 \\
Item C & 22.1 & 12.7 & 34.8 \\
\bottomrule
\end{tabular}
\caption{Sample Data Table}
\label{tab:sample}
\end{table}

\section{Lists}
\subsection{Bullet Points}
\begin{itemize}
  \item First item
  \item Second item
  \item Third item with sub-items:
    \begin{itemize}
      \item Sub-item 1
      \item Sub-item 2
    \end{itemize}
\end{itemize}

\subsection{Numbered List}
\begin{enumerate}
  \item First step
  \item Second step
  \item Third step
\end{enumerate}

\section{Citations and References}
When you use the citation search feature, references will be automatically managed. For example, you might cite a paper like this: \cite{example2024}. All citations will be automatically added to your \texttt{cite.bib} file.

\section{Conclusion}
This template demonstrates basic \LaTeX{} features including mathematical formulas, tables, and lists. The bibliography system is now set up to work seamlessly with Bibby AI's citation management. Feel free to modify it for your needs!

\bibliographystyle{plain}
\bibliography{cite}

\end{document}\documentclass{article}
\usepackage[utf8]{inputenc}
\usepackage{booktabs}
\usepackage{amsmath}
\usepackage{natbib}

\title{My Document}
\author{Your Name}
\date{\today}

\begin{document}

\maketitle

\section{Introduction}
This is a sample document demonstrating various \LaTeX{} features including mathematical equations, tables, and formatting. You can cite papers using the citation search feature, and they will be automatically added to your \texttt{cite.bib} file.

\section{Mathematical Expressions}
Here are some example equations:

\subsection{Quadratic Formula}
The quadratic formula is given by:
\[x = \frac{-b \pm \sqrt{b^2 - 4ac}}{2a}\]

\subsection{Matrix Example}
A sample matrix:
\[
\begin{pmatrix}
a_{11} & a_{12} & a_{13} \\
a_{21} & a_{22} & a_{23} \\
a_{31} & a_{32} & a_{33}
\end{pmatrix}
\]

\section{Tables}
Below is a sample table showing data:

\begin{table}[h]
\centering
\begin{tabular}{@{}lrrr@{}}
\toprule
\textbf{Category} & \textbf{Value 1} & \textbf{Value 2} & \textbf{Total} \\
\midrule
Item A & 10.5 & 20.3 & 30.8 \\
Item B & 15.2 & 25.1 & 40.3 \\
Item C & 22.1 & 12.7 & 34.8 \\
\bottomrule
\end{tabular}
\caption{Sample Data Table}
\label{tab:sample}
\end{table}

\section{Lists}
\subsection{Bullet Points}
\begin{itemize}
  \item First item
  \item Second item
  \item Third item with sub-items:
    \begin{itemize}
      \item Sub-item 1
      \item Sub-item 2
    \end{itemize}
\end{itemize}

\subsection{Numbered List}
\begin{enumerate}
  \item First step
  \item Second step
  \item Third step
\end{enumerate}

\section{Citations and References}
When you use the citation search feature, references will be automatically managed. For example, you might cite a paper like this: \cite{example2024}. All citations will be automatically added to your \texttt{cite.bib} file.

\section{Conclusion}
This template demonstrates basic \LaTeX{} features including mathematical formulas, tables, and lists. The bibliography system is now set up to work seamlessly with Bibby AI's citation management. Feel free to modify it for your needs!

\bibliographystyle{plain}
\bibliography{cite}

\end{document}\documentclass{article}
\usepackage[utf8]{inputenc}
\usepackage{booktabs}
\usepackage{amsmath}
\usepackage{natbib}

\title{My Document}
\author{Your Name}
\date{\today}

\begin{document}

\maketitle

\section{Introduction}
This is a sample document demonstrating various \LaTeX{} features including mathematical equations, tables, and formatting. You can cite papers using the citation search feature, and they will be automatically added to your \texttt{cite.bib} file.

\section{Mathematical Expressions}
Here are some example equations:

\subsection{Quadratic Formula}
The quadratic formula is given by:
\[x = \frac{-b \pm \sqrt{b^2 - 4ac}}{2a}\]

\subsection{Matrix Example}
A sample matrix:
\[
\begin{pmatrix}
a_{11} & a_{12} & a_{13} \\
a_{21} & a_{22} & a_{23} \\
a_{31} & a_{32} & a_{33}
\end{pmatrix}
\]

\section{Tables}
Below is a sample table showing data:

\begin{table}[h]
\centering
\begin{tabular}{@{}lrrr@{}}
\toprule
\textbf{Category} & \textbf{Value 1} & \textbf{Value 2} & \textbf{Total} \\
\midrule
Item A & 10.5 & 20.3 & 30.8 \\
Item B & 15.2 & 25.1 & 40.3 \\
Item C & 22.1 & 12.7 & 34.8 \\
\bottomrule
\end{tabular}
\caption{Sample Data Table}
\label{tab:sample}
\end{table}

\section{Lists}
\subsection{Bullet Points}
\begin{itemize}
  \item First item
  \item Second item
  \item Third item with sub-items:
    \begin{itemize}
      \item Sub-item 1
      \item Sub-item 2
    \end{itemize}
\end{itemize}

\subsection{Numbered List}
\begin{enumerate}
  \item First step
  \item Second step
  \item Third step
\end{enumerate}

\section{Citations and References}
When you use the citation search feature, references will be automatically managed. For example, you might cite a paper like this: \cite{example2024}. All citations will be automatically added to your \texttt{cite.bib} file.

\section{Conclusion}
This template demonstrates basic \LaTeX{} features including mathematical formulas, tables, and lists. The bibliography system is now set up to work seamlessly with Bibby AI's citation management. Feel free to modify it for your needs!

\bibliographystyle{plain}
\bibliography{cite}

\end{document}\documentclass{article}
\usepackage[utf8]{inputenc}
\usepackage{booktabs}
\usepackage{amsmath}
\usepackage{natbib}

\title{My Document}
\author{Your Name}
\date{\today}

\begin{document}

\maketitle

\section{Introduction}
This is a sample document demonstrating various \LaTeX{} features including mathematical equations, tables, and formatting. You can cite papers using the citation search feature, and they will be automatically added to your \texttt{cite.bib} file.

\section{Mathematical Expressions}
Here are some example equations:

\subsection{Quadratic Formula}
The quadratic formula is given by:
\[x = \frac{-b \pm \sqrt{b^2 - 4ac}}{2a}\]

\subsection{Matrix Example}
A sample matrix:
\[
\begin{pmatrix}
a_{11} & a_{12} & a_{13} \\
a_{21} & a_{22} & a_{23} \\
a_{31} & a_{32} & a_{33}
\end{pmatrix}
\]

\section{Tables}
Below is a sample table showing data:

\begin{table}[h]
\centering
\begin{tabular}{@{}lrrr@{}}
\toprule
\textbf{Category} & \textbf{Value 1} & \textbf{Value 2} & \textbf{Total} \\
\midrule
Item A & 10.5 & 20.3 & 30.8 \\
Item B & 15.2 & 25.1 & 40.3 \\
Item C & 22.1 & 12.7 & 34.8 \\
\bottomrule
\end{tabular}
\caption{Sample Data Table}
\label{tab:sample}
\end{table}

\section{Lists}
\subsection{Bullet Points}
\begin{itemize}
  \item First item
  \item Second item
  \item Third item with sub-items:
    \begin{itemize}
      \item Sub-item 1
      \item Sub-item 2
    \end{itemize}
\end{itemize}

\subsection{Numbered List}
\begin{enumerate}
  \item First step
  \item Second step
  \item Third step
\end{enumerate}

\section{Citations and References}
When you use the citation search feature, references will be automatically managed. For example, you might cite a paper like this: \cite{example2024}. All citations will be automatically added to your \texttt{cite.bib} file.

\section{Conclusion}
This template demonstrates basic \LaTeX{} features including mathematical formulas, tables, and lists. The bibliography system is now set up to work seamlessly with Bibby AI's citation management. Feel free to modify it for your needs!

\bibliographystyle{plain}
\bibliography{cite}

\end{document}\documentclass{article}
\usepackage[utf8]{inputenc}
\usepackage{booktabs}
\usepackage{amsmath}
\usepackage{natbib}

\title{My Document}
\author{Your Name}
\date{\today}

\begin{document}

\maketitle

\section{Introduction}
This is a sample document demonstrating various \LaTeX{} features including mathematical equations, tables, and formatting. You can cite papers using the citation search feature, and they will be automatically added to your \texttt{cite.bib} file.

\section{Mathematical Expressions}
Here are some example equations:

\subsection{Quadratic Formula}
The quadratic formula is given by:
\[x = \frac{-b \pm \sqrt{b^2 - 4ac}}{2a}\]

\subsection{Matrix Example}
A sample matrix:
\[
\begin{pmatrix}
a_{11} & a_{12} & a_{13} \\
a_{21} & a_{22} & a_{23} \\
a_{31} & a_{32} & a_{33}
\end{pmatrix}
\]

\section{Tables}
Below is a sample table showing data:

\begin{table}[h]
\centering
\begin{tabular}{@{}lrrr@{}}
\toprule
\textbf{Category} & \textbf{Value 1} & \textbf{Value 2} & \textbf{Total} \\
\midrule
Item A & 10.5 & 20.3 & 30.8 \\
Item B & 15.2 & 25.1 & 40.3 \\
Item C & 22.1 & 12.7 & 34.8 \\
\bottomrule
\end{tabular}
\caption{Sample Data Table}
\label{tab:sample}
\end{table}

\section{Lists}
\subsection{Bullet Points}
\begin{itemize}
  \item First item
  \item Second item
  \item Third item with sub-items:
    \begin{itemize}
      \item Sub-item 1
      \item Sub-item 2
    \end{itemize}
\end{itemize}

\subsection{Numbered List}
\begin{enumerate}
  \item First step
  \item Second step
  \item Third step
\end{enumerate}

\section{Citations and References}
When you use the citation search feature, references will be automatically managed. For example, you might cite a paper like this: \cite{example2024}. All citations will be automatically added to your \texttt{cite.bib} file.

\section{Conclusion}
This template demonstrates basic \LaTeX{} features including mathematical formulas, tables, and lists. The bibliography system is now set up to work seamlessly with Bibby AI's citation management. Feel free to modify it for your needs!

\bibliographystyle{plain}
\bibliography{cite}

\end{document}\documentclass{article}
\usepackage[utf8]{inputenc}
\usepackage{booktabs}
\usepackage{amsmath}
\usepackage{natbib}

\title{My Document}
\author{Your Name}
\date{\today}

\begin{document}

\maketitle

\section{Introduction}
This is a sample document demonstrating various \LaTeX{} features including mathematical equations, tables, and formatting. You can cite papers using the citation search feature, and they will be automatically added to your \texttt{cite.bib} file.

\section{Mathematical Expressions}
Here are some example equations:

\subsection{Quadratic Formula}
The quadratic formula is given by:
\[x = \frac{-b \pm \sqrt{b^2 - 4ac}}{2a}\]

\subsection{Matrix Example}
A sample matrix:
\[
\begin{pmatrix}
a_{11} & a_{12} & a_{13} \\
a_{21} & a_{22} & a_{23} \\
a_{31} & a_{32} & a_{33}
\end{pmatrix}
\]

\section{Tables}
Below is a sample table showing data:

\begin{table}[h]
\centering
\begin{tabular}{@{}lrrr@{}}
\toprule
\textbf{Category} & \textbf{Value 1} & \textbf{Value 2} & \textbf{Total} \\
\midrule
Item A & 10.5 & 20.3 & 30.8 \\
Item B & 15.2 & 25.1 & 40.3 \\
Item C & 22.1 & 12.7 & 34.8 \\
\bottomrule
\end{tabular}
\caption{Sample Data Table}
\label{tab:sample}
\end{table}

\section{Lists}
\subsection{Bullet Points}
\begin{itemize}
  \item First item
  \item Second item
  \item Third item with sub-items:
    \begin{itemize}
      \item Sub-item 1
      \item Sub-item 2
    \end{itemize}
\end{itemize}

\subsection{Numbered List}
\begin{enumerate}
  \item First step
  \item Second step
  \item Third step
\end{enumerate}

\section{Citations and References}
When you use the citation search feature, references will be automatically managed. For example, you might cite a paper like this: \cite{example2024}. All citations will be automatically added to your \texttt{cite.bib} file.

\section{Conclusion}
This template demonstrates basic \LaTeX{} features including mathematical formulas, tables, and lists. The bibliography system is now set up to work seamlessly with Bibby AI's citation management. Feel free to modify it for your needs!

\bibliographystyle{plain}
\bibliography{cite}

\end{document}\documentclass{article}
\usepackage[utf8]{inputenc}
\usepackage{booktabs}
\usepackage{amsmath}
\usepackage{natbib}

\title{My Document}
\author{Your Name}
\date{\today}

\begin{document}

\maketitle

\section{Introduction}
This is a sample document demonstrating various \LaTeX{} features including mathematical equations, tables, and formatting. You can cite papers using the citation search feature, and they will be automatically added to your \texttt{cite.bib} file.

\section{Mathematical Expressions}
Here are some example equations:

\subsection{Quadratic Formula}
The quadratic formula is given by:
\[x = \frac{-b \pm \sqrt{b^2 - 4ac}}{2a}\]

\subsection{Matrix Example}
A sample matrix:
\[
\begin{pmatrix}
a_{11} & a_{12} & a_{13} \\
a_{21} & a_{22} & a_{23} \\
a_{31} & a_{32} & a_{33}
\end{pmatrix}
\]

\section{Tables}
Below is a sample table showing data:

\begin{table}[h]
\centering
\begin{tabular}{@{}lrrr@{}}
\toprule
\textbf{Category} & \textbf{Value 1} & \textbf{Value 2} & \textbf{Total} \\
\midrule
Item A & 10.5 & 20.3 & 30.8 \\
Item B & 15.2 & 25.1 & 40.3 \\
Item C & 22.1 & 12.7 & 34.8 \\
\bottomrule
\end{tabular}
\caption{Sample Data Table}
\label{tab:sample}
\end{table}

\section{Lists}
\subsection{Bullet Points}
\begin{itemize}
  \item First item
  \item Second item
  \item Third item with sub-items:
    \begin{itemize}
      \item Sub-item 1
      \item Sub-item 2
    \end{itemize}
\end{itemize}

\subsection{Numbered List}
\begin{enumerate}
  \item First step
  \item Second step
  \item Third step
\end{enumerate}

\section{Citations and References}
When you use the citation search feature, references will be automatically managed. For example, you might cite a paper like this: \cite{example2024}. All citations will be automatically added to your \texttt{cite.bib} file.

\section{Conclusion}
This template demonstrates basic \LaTeX{} features including mathematical formulas, tables, and lists. The bibliography system is now set up to work seamlessly with Bibby AI's citation management. Feel free to modify it for your needs!

\bibliographystyle{plain}
\bibliography{cite}

\end{document}\documentclass{article}
\usepackage[utf8]{inputenc}
\usepackage{booktabs}
\usepackage{amsmath}
\usepackage{natbib}

\title{My Document}
\author{Your Name}
\date{\today}

\begin{document}

\maketitle

\section{Introduction}
This is a sample document demonstrating various \LaTeX{} features including mathematical equations, tables, and formatting. You can cite papers using the citation search feature, and they will be automatically added to your \texttt{cite.bib} file.

\section{Mathematical Expressions}
Here are some example equations:

\subsection{Quadratic Formula}
The quadratic formula is given by:
\[x = \frac{-b \pm \sqrt{b^2 - 4ac}}{2a}\]

\subsection{Matrix Example}
A sample matrix:
\[
\begin{pmatrix}
a_{11} & a_{12} & a_{13} \\
a_{21} & a_{22} & a_{23} \\
a_{31} & a_{32} & a_{33}
\end{pmatrix}
\]

\section{Tables}
Below is a sample table showing data:

\begin{table}[h]
\centering
\begin{tabular}{@{}lrrr@{}}
\toprule
\textbf{Category} & \textbf{Value 1} & \textbf{Value 2} & \textbf{Total} \\
\midrule
Item A & 10.5 & 20.3 & 30.8 \\
Item B & 15.2 & 25.1 & 40.3 \\
Item C & 22.1 & 12.7 & 34.8 \\
\bottomrule
\end{tabular}
\caption{Sample Data Table}
\label{tab:sample}
\end{table}

\section{Lists}
\subsection{Bullet Points}
\begin{itemize}
  \item First item
  \item Second item
  \item Third item with sub-items:
    \begin{itemize}
      \item Sub-item 1
      \item Sub-item 2
    \end{itemize}
\end{itemize}

\subsection{Numbered List}
\begin{enumerate}
  \item First step
  \item Second step
  \item Third step
\end{enumerate}

\section{Citations and References}
When you use the citation search feature, references will be automatically managed. For example, you might cite a paper like this: \cite{example2024}. All citations will be automatically added to your \texttt{cite.bib} file.

\section{Conclusion}
This template demonstrates basic \LaTeX{} features including mathematical formulas, tables, and lists. The bibliography system is now set up to work seamlessly with Bibby AI's citation management. Feel free to modify it for your needs!

\bibliographystyle{plain}
\bibliography{cite}

\end{document}\documentclass{article}
\usepackage[utf8]{inputenc}
\usepackage{booktabs}
\usepackage{amsmath}
\usepackage{natbib}

\title{My Document}
\author{Your Name}
\date{\today}

\begin{document}

\maketitle

\section{Introduction}
This is a sample document demonstrating various \LaTeX{} features including mathematical equations, tables, and formatting. You can cite papers using the citation search feature, and they will be automatically added to your \texttt{cite.bib} file.

\section{Mathematical Expressions}
Here are some example equations:

\subsection{Quadratic Formula}
The quadratic formula is given by:
\[x = \frac{-b \pm \sqrt{b^2 - 4ac}}{2a}\]

\subsection{Matrix Example}
A sample matrix:
\[
\begin{pmatrix}
a_{11} & a_{12} & a_{13} \\
a_{21} & a_{22} & a_{23} \\
a_{31} & a_{32} & a_{33}
\end{pmatrix}
\]

\section{Tables}
Below is a sample table showing data:

\begin{table}[h]
\centering
\begin{tabular}{@{}lrrr@{}}
\toprule
\textbf{Category} & \textbf{Value 1} & \textbf{Value 2} & \textbf{Total} \\
\midrule
Item A & 10.5 & 20.3 & 30.8 \\
Item B & 15.2 & 25.1 & 40.3 \\
Item C & 22.1 & 12.7 & 34.8 \\
\bottomrule
\end{tabular}
\caption{Sample Data Table}
\label{tab:sample}
\end{table}

\section{Lists}
\subsection{Bullet Points}
\begin{itemize}
  \item First item
  \item Second item
  \item Third item with sub-items:
    \begin{itemize}
      \item Sub-item 1
      \item Sub-item 2
    \end{itemize}
\end{itemize}

\subsection{Numbered List}
\begin{enumerate}
  \item First step
  \item Second step
  \item Third step
\end{enumerate}

\section{Citations and References}
When you use the citation search feature, references will be automatically managed. For example, you might cite a paper like this: \cite{example2024}. All citations will be automatically added to your \texttt{cite.bib} file.

\section{Conclusion}
This template demonstrates basic \LaTeX{} features including mathematical formulas, tables, and lists. The bibliography system is now set up to work seamlessly with Bibby AI's citation management. Feel free to modify it for your needs!

\bibliographystyle{plain}
\bibliography{cite}

\end{document}\documentclass{article}
\usepackage[utf8]{inputenc}
\usepackage{booktabs}
\usepackage{amsmath}
\usepackage{natbib}

\title{My Document}
\author{Your Name}
\date{\today}

\begin{document}

\maketitle

\section{Introduction}
This is a sample document demonstrating various \LaTeX{} features including mathematical equations, tables, and formatting. You can cite papers using the citation search feature, and they will be automatically added to your \texttt{cite.bib} file.

\section{Mathematical Expressions}
Here are some example equations:

\subsection{Quadratic Formula}
The quadratic formula is given by:
\[x = \frac{-b \pm \sqrt{b^2 - 4ac}}{2a}\]

\subsection{Matrix Example}
A sample matrix:
\[
\begin{pmatrix}
a_{11} & a_{12} & a_{13} \\
a_{21} & a_{22} & a_{23} \\
a_{31} & a_{32} & a_{33}
\end{pmatrix}
\]

\section{Tables}
Below is a sample table showing data:

\begin{table}[h]
\centering
\begin{tabular}{@{}lrrr@{}}
\toprule
\textbf{Category} & \textbf{Value 1} & \textbf{Value 2} & \textbf{Total} \\
\midrule
Item A & 10.5 & 20.3 & 30.8 \\
Item B & 15.2 & 25.1 & 40.3 \\
Item C & 22.1 & 12.7 & 34.8 \\
\bottomrule
\end{tabular}
\caption{Sample Data Table}
\label{tab:sample}
\end{table}

\section{Lists}
\subsection{Bullet Points}
\begin{itemize}
  \item First item
  \item Second item
  \item Third item with sub-items:
    \begin{itemize}
      \item Sub-item 1
      \item Sub-item 2
    \end{itemize}
\end{itemize}

\subsection{Numbered List}
\begin{enumerate}
  \item First step
  \item Second step
  \item Third step
\end{enumerate}

\section{Citations and References}
When you use the citation search feature, references will be automatically managed. For example, you might cite a paper like this: \cite{example2024}. All citations will be automatically added to your \texttt{cite.bib} file.

\section{Conclusion}
This template demonstrates basic \LaTeX{} features including mathematical formulas, tables, and lists. The bibliography system is now set up to work seamlessly with Bibby AI's citation management. Feel free to modify it for your needs!

\bibliographystyle{plain}
\bibliography{cite}

\end{document}\documentclass{article}
\usepackage[utf8]{inputenc}
\usepackage{booktabs}
\usepackage{amsmath}
\usepackage{natbib}

\title{My Document}
\author{Your Name}
\date{\today}

\begin{document}

\maketitle

\section{Introduction}
This is a sample document demonstrating various \LaTeX{} features including mathematical equations, tables, and formatting. You can cite papers using the citation search feature, and they will be automatically added to your \texttt{cite.bib} file.

\section{Mathematical Expressions}
Here are some example equations:

\subsection{Quadratic Formula}
The quadratic formula is given by:
\[x = \frac{-b \pm \sqrt{b^2 - 4ac}}{2a}\]

\subsection{Matrix Example}
A sample matrix:
\[
\begin{pmatrix}
a_{11} & a_{12} & a_{13} \\
a_{21} & a_{22} & a_{23} \\
a_{31} & a_{32} & a_{33}
\end{pmatrix}
\]

\section{Tables}
Below is a sample table showing data:

\begin{table}[h]
\centering
\begin{tabular}{@{}lrrr@{}}
\toprule
\textbf{Category} & \textbf{Value 1} & \textbf{Value 2} & \textbf{Total} \\
\midrule
Item A & 10.5 & 20.3 & 30.8 \\
Item B & 15.2 & 25.1 & 40.3 \\
Item C & 22.1 & 12.7 & 34.8 \\
\bottomrule
\end{tabular}
\caption{Sample Data Table}
\label{tab:sample}
\end{table}

\section{Lists}
\subsection{Bullet Points}
\begin{itemize}
  \item First item
  \item Second item
  \item Third item with sub-items:
    \begin{itemize}
      \item Sub-item 1
      \item Sub-item 2
    \end{itemize}
\end{itemize}

\subsection{Numbered List}
\begin{enumerate}
  \item First step
  \item Second step
  \item Third step
\end{enumerate}

\section{Citations and References}
When you use the citation search feature, references will be automatically managed. For example, you might cite a paper like this: \cite{example2024}. All citations will be automatically added to your \texttt{cite.bib} file.

\section{Conclusion}
This template demonstrates basic \LaTeX{} features including mathematical formulas, tables, and lists. The bibliography system is now set up to work seamlessly with Bibby AI's citation management. Feel free to modify it for your needs!

\bibliographystyle{plain}
\bibliography{cite}

\end{document}\documentclass{article}
\usepackage[utf8]{inputenc}
\usepackage{booktabs}
\usepackage{amsmath}
\usepackage{natbib}

\title{My Document}
\author{Your Name}
\date{\today}

\begin{document}

\maketitle

\section{Introduction}
This is a sample document demonstrating various \LaTeX{} features including mathematical equations, tables, and formatting. You can cite papers using the citation search feature, and they will be automatically added to your \texttt{cite.bib} file.

\section{Mathematical Expressions}
Here are some example equations:

\subsection{Quadratic Formula}
The quadratic formula is given by:
\[x = \frac{-b \pm \sqrt{b^2 - 4ac}}{2a}\]

\subsection{Matrix Example}
A sample matrix:
\[
\begin{pmatrix}
a_{11} & a_{12} & a_{13} \\
a_{21} & a_{22} & a_{23} \\
a_{31} & a_{32} & a_{33}
\end{pmatrix}
\]

\section{Tables}
Below is a sample table showing data:

\begin{table}[h]
\centering
\begin{tabular}{@{}lrrr@{}}
\toprule
\textbf{Category} & \textbf{Value 1} & \textbf{Value 2} & \textbf{Total} \\
\midrule
Item A & 10.5 & 20.3 & 30.8 \\
Item B & 15.2 & 25.1 & 40.3 \\
Item C & 22.1 & 12.7 & 34.8 \\
\bottomrule
\end{tabular}
\caption{Sample Data Table}
\label{tab:sample}
\end{table}

\section{Lists}
\subsection{Bullet Points}
\begin{itemize}
  \item First item
  \item Second item
  \item Third item with sub-items:
    \begin{itemize}
      \item Sub-item 1
      \item Sub-item 2
    \end{itemize}
\end{itemize}

\subsection{Numbered List}
\begin{enumerate}
  \item First step
  \item Second step
  \item Third step
\end{enumerate}

\section{Citations and References}
When you use the citation search feature, references will be automatically managed. For example, you might cite a paper like this: \cite{example2024}. All citations will be automatically added to your \texttt{cite.bib} file.

\section{Conclusion}
This template demonstrates basic \LaTeX{} features including mathematical formulas, tables, and lists. The bibliography system is now set up to work seamlessly with Bibby AI's citation management. Feel free to modify it for your needs!

\bibliographystyle{plain}
\bibliography{cite}

\end{document}\documentclass{article}
\usepackage[utf8]{inputenc}
\usepackage{booktabs}
\usepackage{amsmath}
\usepackage{natbib}

\title{My Document}
\author{Your Name}
\date{\today}

\begin{document}

\maketitle

\section{Introduction}
This is a sample document demonstrating various \LaTeX{} features including mathematical equations, tables, and formatting. You can cite papers using the citation search feature, and they will be automatically added to your \texttt{cite.bib} file.

\section{Mathematical Expressions}
Here are some example equations:

\subsection{Quadratic Formula}
The quadratic formula is given by:
\[x = \frac{-b \pm \sqrt{b^2 - 4ac}}{2a}\]

\subsection{Matrix Example}
A sample matrix:
\[
\begin{pmatrix}
a_{11} & a_{12} & a_{13} \\
a_{21} & a_{22} & a_{23} \\
a_{31} & a_{32} & a_{33}
\end{pmatrix}
\]

\section{Tables}
Below is a sample table showing data:

\begin{table}[h]
\centering
\begin{tabular}{@{}lrrr@{}}
\toprule
\textbf{Category} & \textbf{Value 1} & \textbf{Value 2} & \textbf{Total} \\
\midrule
Item A & 10.5 & 20.3 & 30.8 \\
Item B & 15.2 & 25.1 & 40.3 \\
Item C & 22.1 & 12.7 & 34.8 \\
\bottomrule
\end{tabular}
\caption{Sample Data Table}
\label{tab:sample}
\end{table}

\section{Lists}
\subsection{Bullet Points}
\begin{itemize}
  \item First item
  \item Second item
  \item Third item with sub-items:
    \begin{itemize}
      \item Sub-item 1
      \item Sub-item 2
    \end{itemize}
\end{itemize}

\subsection{Numbered List}
\begin{enumerate}
  \item First step
  \item Second step
  \item Third step
\end{enumerate}

\section{Citations and References}
When you use the citation search feature, references will be automatically managed. For example, you might cite a paper like this: \cite{example2024}. All citations will be automatically added to your \texttt{cite.bib} file.

\section{Conclusion}
This template demonstrates basic \LaTeX{} features including mathematical formulas, tables, and lists. The bibliography system is now set up to work seamlessly with Bibby AI's citation management. Feel free to modify it for your needs!

\bibliographystyle{plain}
\bibliography{cite}

\end{document}\documentclass{article}
\usepackage[utf8]{inputenc}
\usepackage{booktabs}
\usepackage{amsmath}
\usepackage{natbib}

\title{My Document}
\author{Your Name}
\date{\today}

\begin{document}

\maketitle

\section{Introduction}
This is a sample document demonstrating various \LaTeX{} features including mathematical equations, tables, and formatting. You can cite papers using the citation search feature, and they will be automatically added to your \texttt{cite.bib} file.

\section{Mathematical Expressions}
Here are some example equations:

\subsection{Quadratic Formula}
The quadratic formula is given by:
\[x = \frac{-b \pm \sqrt{b^2 - 4ac}}{2a}\]

\subsection{Matrix Example}
A sample matrix:
\[
\begin{pmatrix}
a_{11} & a_{12} & a_{13} \\
a_{21} & a_{22} & a_{23} \\
a_{31} & a_{32} & a_{33}
\end{pmatrix}
\]

\section{Tables}
Below is a sample table showing data:

\begin{table}[h]
\centering
\begin{tabular}{@{}lrrr@{}}
\toprule
\textbf{Category} & \textbf{Value 1} & \textbf{Value 2} & \textbf{Total} \\
\midrule
Item A & 10.5 & 20.3 & 30.8 \\
Item B & 15.2 & 25.1 & 40.3 \\
Item C & 22.1 & 12.7 & 34.8 \\
\bottomrule
\end{tabular}
\caption{Sample Data Table}
\label{tab:sample}
\end{table}

\section{Lists}
\subsection{Bullet Points}
\begin{itemize}
  \item First item
  \item Second item
  \item Third item with sub-items:
    \begin{itemize}
      \item Sub-item 1
      \item Sub-item 2
    \end{itemize}
\end{itemize}

\subsection{Numbered List}
\begin{enumerate}
  \item First step
  \item Second step
  \item Third step
\end{enumerate}

\section{Citations and References}
When you use the citation search feature, references will be automatically managed. For example, you might cite a paper like this: \cite{example2024}. All citations will be automatically added to your \texttt{cite.bib} file.

\section{Conclusion}
This template demonstrates basic \LaTeX{} features including mathematical formulas, tables, and lists. The bibliography system is now set up to work seamlessly with Bibby AI's citation management. Feel free to modify it for your needs!

\bibliographystyle{plain}
\bibliography{cite}

\end{document}\documentclass{article}
\usepackage[utf8]{inputenc}
\usepackage{booktabs}
\usepackage{amsmath}
\usepackage{natbib}

\title{My Document}
\author{Your Name}
\date{\today}

\begin{document}

\maketitle

\section{Introduction}
This is a sample document demonstrating various \LaTeX{} features including mathematical equations, tables, and formatting. You can cite papers using the citation search feature, and they will be automatically added to your \texttt{cite.bib} file.

\section{Mathematical Expressions}
Here are some example equations:

\subsection{Quadratic Formula}
The quadratic formula is given by:
\[x = \frac{-b \pm \sqrt{b^2 - 4ac}}{2a}\]

\subsection{Matrix Example}
A sample matrix:
\[
\begin{pmatrix}
a_{11} & a_{12} & a_{13} \\
a_{21} & a_{22} & a_{23} \\
a_{31} & a_{32} & a_{33}
\end{pmatrix}
\]

\section{Tables}
Below is a sample table showing data:

\begin{table}[h]
\centering
\begin{tabular}{@{}lrrr@{}}
\toprule
\textbf{Category} & \textbf{Value 1} & \textbf{Value 2} & \textbf{Total} \\
\midrule
Item A & 10.5 & 20.3 & 30.8 \\
Item B & 15.2 & 25.1 & 40.3 \\
Item C & 22.1 & 12.7 & 34.8 \\
\bottomrule
\end{tabular}
\caption{Sample Data Table}
\label{tab:sample}
\end{table}

\section{Lists}
\subsection{Bullet Points}
\begin{itemize}
  \item First item
  \item Second item
  \item Third item with sub-items:
    \begin{itemize}
      \item Sub-item 1
      \item Sub-item 2
    \end{itemize}
\end{itemize}

\subsection{Numbered List}
\begin{enumerate}
  \item First step
  \item Second step
  \item Third step
\end{enumerate}

\section{Citations and References}
When you use the citation search feature, references will be automatically managed. For example, you might cite a paper like this: \cite{example2024}. All citations will be automatically added to your \texttt{cite.bib} file.

\section{Conclusion}
This template demonstrates basic \LaTeX{} features including mathematical formulas, tables, and lists. The bibliography system is now set up to work seamlessly with Bibby AI's citation management. Feel free to modify it for your needs!

\bibliographystyle{plain}
\bibliography{cite}

\end{document}\documentclass{article}
\usepackage[utf8]{inputenc}
\usepackage{booktabs}
\usepackage{amsmath}
\usepackage{natbib}

\title{My Document}
\author{Your Name}
\date{\today}

\begin{document}

\maketitle

\section{Introduction}
This is a sample document demonstrating various \LaTeX{} features including mathematical equations, tables, and formatting. You can cite papers using the citation search feature, and they will be automatically added to your \texttt{cite.bib} file.

\section{Mathematical Expressions}
Here are some example equations:

\subsection{Quadratic Formula}
The quadratic formula is given by:
\[x = \frac{-b \pm \sqrt{b^2 - 4ac}}{2a}\]

\subsection{Matrix Example}
A sample matrix:
\[
\begin{pmatrix}
a_{11} & a_{12} & a_{13} \\
a_{21} & a_{22} & a_{23} \\
a_{31} & a_{32} & a_{33}
\end{pmatrix}
\]

\section{Tables}
Below is a sample table showing data:

\begin{table}[h]
\centering
\begin{tabular}{@{}lrrr@{}}
\toprule
\textbf{Category} & \textbf{Value 1} & \textbf{Value 2} & \textbf{Total} \\
\midrule
Item A & 10.5 & 20.3 & 30.8 \\
Item B & 15.2 & 25.1 & 40.3 \\
Item C & 22.1 & 12.7 & 34.8 \\
\bottomrule
\end{tabular}
\caption{Sample Data Table}
\label{tab:sample}
\end{table}

\section{Lists}
\subsection{Bullet Points}
\begin{itemize}
  \item First item
  \item Second item
  \item Third item with sub-items:
    \begin{itemize}
      \item Sub-item 1
      \item Sub-item 2
    \end{itemize}
\end{itemize}

\subsection{Numbered List}
\begin{enumerate}
  \item First step
  \item Second step
  \item Third step
\end{enumerate}

\section{Citations and References}
When you use the citation search feature, references will be automatically managed. For example, you might cite a paper like this: \cite{example2024}. All citations will be automatically added to your \texttt{cite.bib} file.

\section{Conclusion}
This template demonstrates basic \LaTeX{} features including mathematical formulas, tables, and lists. The bibliography system is now set up to work seamlessly with Bibby AI's citation management. Feel free to modify it for your needs!

\bibliographystyle{plain}
\bibliography{cite}

\end{document}\documentclass{article}
\usepackage[utf8]{inputenc}
\usepackage{booktabs}
\usepackage{amsmath}
\usepackage{natbib}

\title{My Document}
\author{Your Name}
\date{\today}

\begin{document}

\maketitle

\section{Introduction}
This is a sample document demonstrating various \LaTeX{} features including mathematical equations, tables, and formatting. You can cite papers using the citation search feature, and they will be automatically added to your \texttt{cite.bib} file.

\section{Mathematical Expressions}
Here are some example equations:

\subsection{Quadratic Formula}
The quadratic formula is given by:
\[x = \frac{-b \pm \sqrt{b^2 - 4ac}}{2a}\]

\subsection{Matrix Example}
A sample matrix:
\[
\begin{pmatrix}
a_{11} & a_{12} & a_{13} \\
a_{21} & a_{22} & a_{23} \\
a_{31} & a_{32} & a_{33}
\end{pmatrix}
\]

\section{Tables}
Below is a sample table showing data:

\begin{table}[h]
\centering
\begin{tabular}{@{}lrrr@{}}
\toprule
\textbf{Category} & \textbf{Value 1} & \textbf{Value 2} & \textbf{Total} \\
\midrule
Item A & 10.5 & 20.3 & 30.8 \\
Item B & 15.2 & 25.1 & 40.3 \\
Item C & 22.1 & 12.7 & 34.8 \\
\bottomrule
\end{tabular}
\caption{Sample Data Table}
\label{tab:sample}
\end{table}

\section{Lists}
\subsection{Bullet Points}
\begin{itemize}
  \item First item
  \item Second item
  \item Third item with sub-items:
    \begin{itemize}
      \item Sub-item 1
      \item Sub-item 2
    \end{itemize}
\end{itemize}

\subsection{Numbered List}
\begin{enumerate}
  \item First step
  \item Second step
  \item Third step
\end{enumerate}

\section{Citations and References}
When you use the citation search feature, references will be automatically managed. For example, you might cite a paper like this: \cite{example2024}. All citations will be automatically added to your \texttt{cite.bib} file.

\section{Conclusion}
This template demonstrates basic \LaTeX{} features including mathematical formulas, tables, and lists. The bibliography system is now set up to work seamlessly with Bibby AI's citation management. Feel free to modify it for your needs!

\bibliographystyle{plain}
\bibliography{cite}

\end{document}\documentclass{article}
\usepackage[utf8]{inputenc}
\usepackage{booktabs}
\usepackage{amsmath}
\usepackage{natbib}

\title{My Document}
\author{Your Name}
\date{\today}

\begin{document}

\maketitle

\section{Introduction}
This is a sample document demonstrating various \LaTeX{} features including mathematical equations, tables, and formatting. You can cite papers using the citation search feature, and they will be automatically added to your \texttt{cite.bib} file.

\section{Mathematical Expressions}
Here are some example equations:

\subsection{Quadratic Formula}
The quadratic formula is given by:
\[x = \frac{-b \pm \sqrt{b^2 - 4ac}}{2a}\]

\subsection{Matrix Example}
A sample matrix:
\[
\begin{pmatrix}
a_{11} & a_{12} & a_{13} \\
a_{21} & a_{22} & a_{23} \\
a_{31} & a_{32} & a_{33}
\end{pmatrix}
\]

\section{Tables}
Below is a sample table showing data:

\begin{table}[h]
\centering
\begin{tabular}{@{}lrrr@{}}
\toprule
\textbf{Category} & \textbf{Value 1} & \textbf{Value 2} & \textbf{Total} \\
\midrule
Item A & 10.5 & 20.3 & 30.8 \\
Item B & 15.2 & 25.1 & 40.3 \\
Item C & 22.1 & 12.7 & 34.8 \\
\bottomrule
\end{tabular}
\caption{Sample Data Table}
\label{tab:sample}
\end{table}

\section{Lists}
\subsection{Bullet Points}
\begin{itemize}
  \item First item
  \item Second item
  \item Third item with sub-items:
    \begin{itemize}
      \item Sub-item 1
      \item Sub-item 2
    \end{itemize}
\end{itemize}

\subsection{Numbered List}
\begin{enumerate}
  \item First step
  \item Second step
  \item Third step
\end{enumerate}

\section{Citations and References}
When you use the citation search feature, references will be automatically managed. For example, you might cite a paper like this: \cite{example2024}. All citations will be automatically added to your \texttt{cite.bib} file.

\section{Conclusion}
This template demonstrates basic \LaTeX{} features including mathematical formulas, tables, and lists. The bibliography system is now set up to work seamlessly with Bibby AI's citation management. Feel free to modify it for your needs!

\bibliographystyle{plain}
\bibliography{cite}

\end{document}\documentclass{article}
\usepackage[utf8]{inputenc}
\usepackage{booktabs}
\usepackage{amsmath}
\usepackage{natbib}

\title{My Document}
\author{Your Name}
\date{\today}

\begin{document}

\maketitle

\section{Introduction}
This is a sample document demonstrating various \LaTeX{} features including mathematical equations, tables, and formatting. You can cite papers using the citation search feature, and they will be automatically added to your \texttt{cite.bib} file.

\section{Mathematical Expressions}
Here are some example equations:

\subsection{Quadratic Formula}
The quadratic formula is given by:
\[x = \frac{-b \pm \sqrt{b^2 - 4ac}}{2a}\]

\subsection{Matrix Example}
A sample matrix:
\[
\begin{pmatrix}
a_{11} & a_{12} & a_{13} \\
a_{21} & a_{22} & a_{23} \\
a_{31} & a_{32} & a_{33}
\end{pmatrix}
\]

\section{Tables}
Below is a sample table showing data:

\begin{table}[h]
\centering
\begin{tabular}{@{}lrrr@{}}
\toprule
\textbf{Category} & \textbf{Value 1} & \textbf{Value 2} & \textbf{Total} \\
\midrule
Item A & 10.5 & 20.3 & 30.8 \\
Item B & 15.2 & 25.1 & 40.3 \\
Item C & 22.1 & 12.7 & 34.8 \\
\bottomrule
\end{tabular}
\caption{Sample Data Table}
\label{tab:sample}
\end{table}

\section{Lists}
\subsection{Bullet Points}
\begin{itemize}
  \item First item
  \item Second item
  \item Third item with sub-items:
    \begin{itemize}
      \item Sub-item 1
      \item Sub-item 2
    \end{itemize}
\end{itemize}

\subsection{Numbered List}
\begin{enumerate}
  \item First step
  \item Second step
  \item Third step
\end{enumerate}

\section{Citations and References}
When you use the citation search feature, references will be automatically managed. For example, you might cite a paper like this: \cite{example2024}. All citations will be automatically added to your \texttt{cite.bib} file.

\section{Conclusion}
This template demonstrates basic \LaTeX{} features including mathematical formulas, tables, and lists. The bibliography system is now set up to work seamlessly with Bibby AI's citation management. Feel free to modify it for your needs!

\bibliographystyle{plain}
\bibliography{cite}

\end{document}\documentclass{article}
\usepackage[utf8]{inputenc}
\usepackage{booktabs}
\usepackage{amsmath}
\usepackage{natbib}

\title{My Document}
\author{Your Name}
\date{\today}

\begin{document}

\maketitle

\section{Introduction}
This is a sample document demonstrating various \LaTeX{} features including mathematical equations, tables, and formatting. You can cite papers using the citation search feature, and they will be automatically added to your \texttt{cite.bib} file.

\section{Mathematical Expressions}
Here are some example equations:

\subsection{Quadratic Formula}
The quadratic formula is given by:
\[x = \frac{-b \pm \sqrt{b^2 - 4ac}}{2a}\]

\subsection{Matrix Example}
A sample matrix:
\[
\begin{pmatrix}
a_{11} & a_{12} & a_{13} \\
a_{21} & a_{22} & a_{23} \\
a_{31} & a_{32} & a_{33}
\end{pmatrix}
\]

\section{Tables}
Below is a sample table showing data:

\begin{table}[h]
\centering
\begin{tabular}{@{}lrrr@{}}
\toprule
\textbf{Category} & \textbf{Value 1} & \textbf{Value 2} & \textbf{Total} \\
\midrule
Item A & 10.5 & 20.3 & 30.8 \\
Item B & 15.2 & 25.1 & 40.3 \\
Item C & 22.1 & 12.7 & 34.8 \\
\bottomrule
\end{tabular}
\caption{Sample Data Table}
\label{tab:sample}
\end{table}

\section{Lists}
\subsection{Bullet Points}
\begin{itemize}
  \item First item
  \item Second item
  \item Third item with sub-items:
    \begin{itemize}
      \item Sub-item 1
      \item Sub-item 2
    \end{itemize}
\end{itemize}

\subsection{Numbered List}
\begin{enumerate}
  \item First step
  \item Second step
  \item Third step
\end{enumerate}

\section{Citations and References}
When you use the citation search feature, references will be automatically managed. For example, you might cite a paper like this: \cite{example2024}. All citations will be automatically added to your \texttt{cite.bib} file.

\section{Conclusion}
This template demonstrates basic \LaTeX{} features including mathematical formulas, tables, and lists. The bibliography system is now set up to work seamlessly with Bibby AI's citation management. Feel free to modify it for your needs!

\bibliographystyle{plain}
\bibliography{cite}

\end{document}\documentclass{article}
\usepackage[utf8]{inputenc}
\usepackage{booktabs}
\usepackage{amsmath}
\usepackage{natbib}

\title{My Document}
\author{Your Name}
\date{\today}

\begin{document}

\maketitle

\section{Introduction}
This is a sample document demonstrating various \LaTeX{} features including mathematical equations, tables, and formatting. You can cite papers using the citation search feature, and they will be automatically added to your \texttt{cite.bib} file.

\section{Mathematical Expressions}
Here are some example equations:

\subsection{Quadratic Formula}
The quadratic formula is given by:
\[x = \frac{-b \pm \sqrt{b^2 - 4ac}}{2a}\]

\subsection{Matrix Example}
A sample matrix:
\[
\begin{pmatrix}
a_{11} & a_{12} & a_{13} \\
a_{21} & a_{22} & a_{23} \\
a_{31} & a_{32} & a_{33}
\end{pmatrix}
\]

\section{Tables}
Below is a sample table showing data:

\begin{table}[h]
\centering
\begin{tabular}{@{}lrrr@{}}
\toprule
\textbf{Category} & \textbf{Value 1} & \textbf{Value 2} & \textbf{Total} \\
\midrule
Item A & 10.5 & 20.3 & 30.8 \\
Item B & 15.2 & 25.1 & 40.3 \\
Item C & 22.1 & 12.7 & 34.8 \\
\bottomrule
\end{tabular}
\caption{Sample Data Table}
\label{tab:sample}
\end{table}

\section{Lists}
\subsection{Bullet Points}
\begin{itemize}
  \item First item
  \item Second item
  \item Third item with sub-items:
    \begin{itemize}
      \item Sub-item 1
      \item Sub-item 2
    \end{itemize}
\end{itemize}

\subsection{Numbered List}
\begin{enumerate}
  \item First step
  \item Second step
  \item Third step
\end{enumerate}

\section{Citations and References}
When you use the citation search feature, references will be automatically managed. For example, you might cite a paper like this: \cite{example2024}. All citations will be automatically added to your \texttt{cite.bib} file.

\section{Conclusion}
This template demonstrates basic \LaTeX{} features including mathematical formulas, tables, and lists. The bibliography system is now set up to work seamlessly with Bibby AI's citation management. Feel free to modify it for your needs!

\bibliographystyle{plain}
\bibliography{cite}

\end{document}\documentclass{article}
\usepackage[utf8]{inputenc}
\usepackage{booktabs}
\usepackage{amsmath}
\usepackage{natbib}

\title{My Document}
\author{Your Name}
\date{\today}

\begin{document}

\maketitle

\section{Introduction}
This is a sample document demonstrating various \LaTeX{} features including mathematical equations, tables, and formatting. You can cite papers using the citation search feature, and they will be automatically added to your \texttt{cite.bib} file.

\section{Mathematical Expressions}
Here are some example equations:

\subsection{Quadratic Formula}
The quadratic formula is given by:
\[x = \frac{-b \pm \sqrt{b^2 - 4ac}}{2a}\]

\subsection{Matrix Example}
A sample matrix:
\[
\begin{pmatrix}
a_{11} & a_{12} & a_{13} \\
a_{21} & a_{22} & a_{23} \\
a_{31} & a_{32} & a_{33}
\end{pmatrix}
\]

\section{Tables}
Below is a sample table showing data:

\begin{table}[h]
\centering
\begin{tabular}{@{}lrrr@{}}
\toprule
\textbf{Category} & \textbf{Value 1} & \textbf{Value 2} & \textbf{Total} \\
\midrule
Item A & 10.5 & 20.3 & 30.8 \\
Item B & 15.2 & 25.1 & 40.3 \\
Item C & 22.1 & 12.7 & 34.8 \\
\bottomrule
\end{tabular}
\caption{Sample Data Table}
\label{tab:sample}
\end{table}

\section{Lists}
\subsection{Bullet Points}
\begin{itemize}
  \item First item
  \item Second item
  \item Third item with sub-items:
    \begin{itemize}
      \item Sub-item 1
      \item Sub-item 2
    \end{itemize}
\end{itemize}

\subsection{Numbered List}
\begin{enumerate}
  \item First step
  \item Second step
  \item Third step
\end{enumerate}

\section{Citations and References}
When you use the citation search feature, references will be automatically managed. For example, you might cite a paper like this: \cite{example2024}. All citations will be automatically added to your \texttt{cite.bib} file.

\section{Conclusion}
This template demonstrates basic \LaTeX{} features including mathematical formulas, tables, and lists. The bibliography system is now set up to work seamlessly with Bibby AI's citation management. Feel free to modify it for your needs!

\bibliographystyle{plain}
\bibliography{cite}

\end{document}\documentclass{article}
\usepackage[utf8]{inputenc}
\usepackage{booktabs}
\usepackage{amsmath}
\usepackage{natbib}

\title{My Document}
\author{Your Name}
\date{\today}

\begin{document}

\maketitle

\section{Introduction}
This is a sample document demonstrating various \LaTeX{} features including mathematical equations, tables, and formatting. You can cite papers using the citation search feature, and they will be automatically added to your \texttt{cite.bib} file.

\section{Mathematical Expressions}
Here are some example equations:

\subsection{Quadratic Formula}
The quadratic formula is given by:
\[x = \frac{-b \pm \sqrt{b^2 - 4ac}}{2a}\]

\subsection{Matrix Example}
A sample matrix:
\[
\begin{pmatrix}
a_{11} & a_{12} & a_{13} \\
a_{21} & a_{22} & a_{23} \\
a_{31} & a_{32} & a_{33}
\end{pmatrix}
\]

\section{Tables}
Below is a sample table showing data:

\begin{table}[h]
\centering
\begin{tabular}{@{}lrrr@{}}
\toprule
\textbf{Category} & \textbf{Value 1} & \textbf{Value 2} & \textbf{Total} \\
\midrule
Item A & 10.5 & 20.3 & 30.8 \\
Item B & 15.2 & 25.1 & 40.3 \\
Item C & 22.1 & 12.7 & 34.8 \\
\bottomrule
\end{tabular}
\caption{Sample Data Table}
\label{tab:sample}
\end{table}

\section{Lists}
\subsection{Bullet Points}
\begin{itemize}
  \item First item
  \item Second item
  \item Third item with sub-items:
    \begin{itemize}
      \item Sub-item 1
      \item Sub-item 2
    \end{itemize}
\end{itemize}

\subsection{Numbered List}
\begin{enumerate}
  \item First step
  \item Second step
  \item Third step
\end{enumerate}

\section{Citations and References}
When you use the citation search feature, references will be automatically managed. For example, you might cite a paper like this: \cite{example2024}. All citations will be automatically added to your \texttt{cite.bib} file.

\section{Conclusion}
This template demonstrates basic \LaTeX{} features including mathematical formulas, tables, and lists. The bibliography system is now set up to work seamlessly with Bibby AI's citation management. Feel free to modify it for your needs!

\bibliographystyle{plain}
\bibliography{cite}

\end{document}\documentclass{article}
\usepackage[utf8]{inputenc}
\usepackage{booktabs}
\usepackage{amsmath}
\usepackage{natbib}

\title{My Document}
\author{Your Name}
\date{\today}

\begin{document}

\maketitle

\section{Introduction}
This is a sample document demonstrating various \LaTeX{} features including mathematical equations, tables, and formatting. You can cite papers using the citation search feature, and they will be automatically added to your \texttt{cite.bib} file.

\section{Mathematical Expressions}
Here are some example equations:

\subsection{Quadratic Formula}
The quadratic formula is given by:
\[x = \frac{-b \pm \sqrt{b^2 - 4ac}}{2a}\]

\subsection{Matrix Example}
A sample matrix:
\[
\begin{pmatrix}
a_{11} & a_{12} & a_{13} \\
a_{21} & a_{22} & a_{23} \\
a_{31} & a_{32} & a_{33}
\end{pmatrix}
\]

\section{Tables}
Below is a sample table showing data:

\begin{table}[h]
\centering
\begin{tabular}{@{}lrrr@{}}
\toprule
\textbf{Category} & \textbf{Value 1} & \textbf{Value 2} & \textbf{Total} \\
\midrule
Item A & 10.5 & 20.3 & 30.8 \\
Item B & 15.2 & 25.1 & 40.3 \\
Item C & 22.1 & 12.7 & 34.8 \\
\bottomrule
\end{tabular}
\caption{Sample Data Table}
\label{tab:sample}
\end{table}

\section{Lists}
\subsection{Bullet Points}
\begin{itemize}
  \item First item
  \item Second item
  \item Third item with sub-items:
    \begin{itemize}
      \item Sub-item 1
      \item Sub-item 2
    \end{itemize}
\end{itemize}

\subsection{Numbered List}
\begin{enumerate}
  \item First step
  \item Second step
  \item Third step
\end{enumerate}

\section{Citations and References}
When you use the citation search feature, references will be automatically managed. For example, you might cite a paper like this: \cite{example2024}. All citations will be automatically added to your \texttt{cite.bib} file.

\section{Conclusion}
This template demonstrates basic \LaTeX{} features including mathematical formulas, tables, and lists. The bibliography system is now set up to work seamlessly with Bibby AI's citation management. Feel free to modify it for your needs!

\bibliographystyle{plain}
\bibliography{cite}

\end{document}\documentclass{article}
\usepackage[utf8]{inputenc}
\usepackage{booktabs}
\usepackage{amsmath}
\usepackage{natbib}

\title{My Document}
\author{Your Name}
\date{\today}

\begin{document}

\maketitle

\section{Introduction}
This is a sample document demonstrating various \LaTeX{} features including mathematical equations, tables, and formatting. You can cite papers using the citation search feature, and they will be automatically added to your \texttt{cite.bib} file.

\section{Mathematical Expressions}
Here are some example equations:

\subsection{Quadratic Formula}
The quadratic formula is given by:
\[x = \frac{-b \pm \sqrt{b^2 - 4ac}}{2a}\]

\subsection{Matrix Example}
A sample matrix:
\[
\begin{pmatrix}
a_{11} & a_{12} & a_{13} \\
a_{21} & a_{22} & a_{23} \\
a_{31} & a_{32} & a_{33}
\end{pmatrix}
\]

\section{Tables}
Below is a sample table showing data:

\begin{table}[h]
\centering
\begin{tabular}{@{}lrrr@{}}
\toprule
\textbf{Category} & \textbf{Value 1} & \textbf{Value 2} & \textbf{Total} \\
\midrule
Item A & 10.5 & 20.3 & 30.8 \\
Item B & 15.2 & 25.1 & 40.3 \\
Item C & 22.1 & 12.7 & 34.8 \\
\bottomrule
\end{tabular}
\caption{Sample Data Table}
\label{tab:sample}
\end{table}

\section{Lists}
\subsection{Bullet Points}
\begin{itemize}
  \item First item
  \item Second item
  \item Third item with sub-items:
    \begin{itemize}
      \item Sub-item 1
      \item Sub-item 2
    \end{itemize}
\end{itemize}

\subsection{Numbered List}
\begin{enumerate}
  \item First step
  \item Second step
  \item Third step
\end{enumerate}

\section{Citations and References}
When you use the citation search feature, references will be automatically managed. For example, you might cite a paper like this: \cite{example2024}. All citations will be automatically added to your \texttt{cite.bib} file.

\section{Conclusion}
This template demonstrates basic \LaTeX{} features including mathematical formulas, tables, and lists. The bibliography system is now set up to work seamlessly with Bibby AI's citation management. Feel free to modify it for your needs!

\bibliographystyle{plain}
\bibliography{cite}

\end{document}\documentclass{article}
\usepackage[utf8]{inputenc}
\usepackage{booktabs}
\usepackage{amsmath}
\usepackage{natbib}

\title{My Document}
\author{Your Name}
\date{\today}

\begin{document}

\maketitle

\section{Introduction}
This is a sample document demonstrating various \LaTeX{} features including mathematical equations, tables, and formatting. You can cite papers using the citation search feature, and they will be automatically added to your \texttt{cite.bib} file.

\section{Mathematical Expressions}
Here are some example equations:

\subsection{Quadratic Formula}
The quadratic formula is given by:
\[x = \frac{-b \pm \sqrt{b^2 - 4ac}}{2a}\]

\subsection{Matrix Example}
A sample matrix:
\[
\begin{pmatrix}
a_{11} & a_{12} & a_{13} \\
a_{21} & a_{22} & a_{23} \\
a_{31} & a_{32} & a_{33}
\end{pmatrix}
\]

\section{Tables}
Below is a sample table showing data:

\begin{table}[h]
\centering
\begin{tabular}{@{}lrrr@{}}
\toprule
\textbf{Category} & \textbf{Value 1} & \textbf{Value 2} & \textbf{Total} \\
\midrule
Item A & 10.5 & 20.3 & 30.8 \\
Item B & 15.2 & 25.1 & 40.3 \\
Item C & 22.1 & 12.7 & 34.8 \\
\bottomrule
\end{tabular}
\caption{Sample Data Table}
\label{tab:sample}
\end{table}

\section{Lists}
\subsection{Bullet Points}
\begin{itemize}
  \item First item
  \item Second item
  \item Third item with sub-items:
    \begin{itemize}
      \item Sub-item 1
      \item Sub-item 2
    \end{itemize}
\end{itemize}

\subsection{Numbered List}
\begin{enumerate}
  \item First step
  \item Second step
  \item Third step
\end{enumerate}

\section{Citations and References}
When you use the citation search feature, references will be automatically managed. For example, you might cite a paper like this: \cite{example2024}. All citations will be automatically added to your \texttt{cite.bib} file.

\section{Conclusion}
This template demonstrates basic \LaTeX{} features including mathematical formulas, tables, and lists. The bibliography system is now set up to work seamlessly with Bibby AI's citation management. Feel free to modify it for your needs!

\bibliographystyle{plain}
\bibliography{cite}

\end{document}\documentclass{article}
\usepackage[utf8]{inputenc}
\usepackage{booktabs}
\usepackage{amsmath}
\usepackage{natbib}

\title{My Document}
\author{Your Name}
\date{\today}

\begin{document}

\maketitle

\section{Introduction}
This is a sample document demonstrating various \LaTeX{} features including mathematical equations, tables, and formatting. You can cite papers using the citation search feature, and they will be automatically added to your \texttt{cite.bib} file.

\section{Mathematical Expressions}
Here are some example equations:

\subsection{Quadratic Formula}
The quadratic formula is given by:
\[x = \frac{-b \pm \sqrt{b^2 - 4ac}}{2a}\]

\subsection{Matrix Example}
A sample matrix:
\[
\begin{pmatrix}
a_{11} & a_{12} & a_{13} \\
a_{21} & a_{22} & a_{23} \\
a_{31} & a_{32} & a_{33}
\end{pmatrix}
\]

\section{Tables}
Below is a sample table showing data:

\begin{table}[h]
\centering
\begin{tabular}{@{}lrrr@{}}
\toprule
\textbf{Category} & \textbf{Value 1} & \textbf{Value 2} & \textbf{Total} \\
\midrule
Item A & 10.5 & 20.3 & 30.8 \\
Item B & 15.2 & 25.1 & 40.3 \\
Item C & 22.1 & 12.7 & 34.8 \\
\bottomrule
\end{tabular}
\caption{Sample Data Table}
\label{tab:sample}
\end{table}

\section{Lists}
\subsection{Bullet Points}
\begin{itemize}
  \item First item
  \item Second item
  \item Third item with sub-items:
    \begin{itemize}
      \item Sub-item 1
      \item Sub-item 2
    \end{itemize}
\end{itemize}

\subsection{Numbered List}
\begin{enumerate}
  \item First step
  \item Second step
  \item Third step
\end{enumerate}

\section{Citations and References}
When you use the citation search feature, references will be automatically managed. For example, you might cite a paper like this: \cite{example2024}. All citations will be automatically added to your \texttt{cite.bib} file.

\section{Conclusion}
This template demonstrates basic \LaTeX{} features including mathematical formulas, tables, and lists. The bibliography system is now set up to work seamlessly with Bibby AI's citation management. Feel free to modify it for your needs!

\bibliographystyle{plain}
\bibliography{cite}

\end{document}\documentclass{article}
\usepackage[utf8]{inputenc}
\usepackage{booktabs}
\usepackage{amsmath}
\usepackage{natbib}

\title{My Document}
\author{Your Name}
\date{\today}

\begin{document}

\maketitle

\section{Introduction}
This is a sample document demonstrating various \LaTeX{} features including mathematical equations, tables, and formatting. You can cite papers using the citation search feature, and they will be automatically added to your \texttt{cite.bib} file.

\section{Mathematical Expressions}
Here are some example equations:

\subsection{Quadratic Formula}
The quadratic formula is given by:
\[x = \frac{-b \pm \sqrt{b^2 - 4ac}}{2a}\]

\subsection{Matrix Example}
A sample matrix:
\[
\begin{pmatrix}
a_{11} & a_{12} & a_{13} \\
a_{21} & a_{22} & a_{23} \\
a_{31} & a_{32} & a_{33}
\end{pmatrix}
\]

\section{Tables}
Below is a sample table showing data:

\begin{table}[h]
\centering
\begin{tabular}{@{}lrrr@{}}
\toprule
\textbf{Category} & \textbf{Value 1} & \textbf{Value 2} & \textbf{Total} \\
\midrule
Item A & 10.5 & 20.3 & 30.8 \\
Item B & 15.2 & 25.1 & 40.3 \\
Item C & 22.1 & 12.7 & 34.8 \\
\bottomrule
\end{tabular}
\caption{Sample Data Table}
\label{tab:sample}
\end{table}

\section{Lists}
\subsection{Bullet Points}
\begin{itemize}
  \item First item
  \item Second item
  \item Third item with sub-items:
    \begin{itemize}
      \item Sub-item 1
      \item Sub-item 2
    \end{itemize}
\end{itemize}

\subsection{Numbered List}
\begin{enumerate}
  \item First step
  \item Second step
  \item Third step
\end{enumerate}

\section{Citations and References}
When you use the citation search feature, references will be automatically managed. For example, you might cite a paper like this: \cite{example2024}. All citations will be automatically added to your \texttt{cite.bib} file.

\section{Conclusion}
This template demonstrates basic \LaTeX{} features including mathematical formulas, tables, and lists. The bibliography system is now set up to work seamlessly with Bibby AI's citation management. Feel free to modify it for your needs!

\bibliographystyle{plain}
\bibliography{cite}

\end{document}\documentclass{article}
\usepackage[utf8]{inputenc}
\usepackage{booktabs}
\usepackage{amsmath}
\usepackage{natbib}

\title{My Document}
\author{Your Name}
\date{\today}

\begin{document}

\maketitle

\section{Introduction}
This is a sample document demonstrating various \LaTeX{} features including mathematical equations, tables, and formatting. You can cite papers using the citation search feature, and they will be automatically added to your \texttt{cite.bib} file.

\section{Mathematical Expressions}
Here are some example equations:

\subsection{Quadratic Formula}
The quadratic formula is given by:
\[x = \frac{-b \pm \sqrt{b^2 - 4ac}}{2a}\]

\subsection{Matrix Example}
A sample matrix:
\[
\begin{pmatrix}
a_{11} & a_{12} & a_{13} \\
a_{21} & a_{22} & a_{23} \\
a_{31} & a_{32} & a_{33}
\end{pmatrix}
\]

\section{Tables}
Below is a sample table showing data:

\begin{table}[h]
\centering
\begin{tabular}{@{}lrrr@{}}
\toprule
\textbf{Category} & \textbf{Value 1} & \textbf{Value 2} & \textbf{Total} \\
\midrule
Item A & 10.5 & 20.3 & 30.8 \\
Item B & 15.2 & 25.1 & 40.3 \\
Item C & 22.1 & 12.7 & 34.8 \\
\bottomrule
\end{tabular}
\caption{Sample Data Table}
\label{tab:sample}
\end{table}

\section{Lists}
\subsection{Bullet Points}
\begin{itemize}
  \item First item
  \item Second item
  \item Third item with sub-items:
    \begin{itemize}
      \item Sub-item 1
      \item Sub-item 2
    \end{itemize}
\end{itemize}

\subsection{Numbered List}
\begin{enumerate}
  \item First step
  \item Second step
  \item Third step
\end{enumerate}

\section{Citations and References}
When you use the citation search feature, references will be automatically managed. For example, you might cite a paper like this: \cite{example2024}. All citations will be automatically added to your \texttt{cite.bib} file.

\section{Conclusion}
This template demonstrates basic \LaTeX{} features including mathematical formulas, tables, and lists. The bibliography system is now set up to work seamlessly with Bibby AI's citation management. Feel free to modify it for your needs!

\bibliographystyle{plain}
\bibliography{cite}

\end{document}\documentclass{article}
\usepackage[utf8]{inputenc}
\usepackage{booktabs}
\usepackage{amsmath}
\usepackage{natbib}

\title{My Document}
\author{Your Name}
\date{\today}

\begin{document}

\maketitle

\section{Introduction}
This is a sample document demonstrating various \LaTeX{} features including mathematical equations, tables, and formatting. You can cite papers using the citation search feature, and they will be automatically added to your \texttt{cite.bib} file.

\section{Mathematical Expressions}
Here are some example equations:

\subsection{Quadratic Formula}
The quadratic formula is given by:
\[x = \frac{-b \pm \sqrt{b^2 - 4ac}}{2a}\]

\subsection{Matrix Example}
A sample matrix:
\[
\begin{pmatrix}
a_{11} & a_{12} & a_{13} \\
a_{21} & a_{22} & a_{23} \\
a_{31} & a_{32} & a_{33}
\end{pmatrix}
\]

\section{Tables}
Below is a sample table showing data:

\begin{table}[h]
\centering
\begin{tabular}{@{}lrrr@{}}
\toprule
\textbf{Category} & \textbf{Value 1} & \textbf{Value 2} & \textbf{Total} \\
\midrule
Item A & 10.5 & 20.3 & 30.8 \\
Item B & 15.2 & 25.1 & 40.3 \\
Item C & 22.1 & 12.7 & 34.8 \\
\bottomrule
\end{tabular}
\caption{Sample Data Table}
\label{tab:sample}
\end{table}

\section{Lists}
\subsection{Bullet Points}
\begin{itemize}
  \item First item
  \item Second item
  \item Third item with sub-items:
    \begin{itemize}
      \item Sub-item 1
      \item Sub-item 2
    \end{itemize}
\end{itemize}

\subsection{Numbered List}
\begin{enumerate}
  \item First step
  \item Second step
  \item Third step
\end{enumerate}

\section{Citations and References}
When you use the citation search feature, references will be automatically managed. For example, you might cite a paper like this: \cite{example2024}. All citations will be automatically added to your \texttt{cite.bib} file.

\section{Conclusion}
This template demonstrates basic \LaTeX{} features including mathematical formulas, tables, and lists. The bibliography system is now set up to work seamlessly with Bibby AI's citation management. Feel free to modify it for your needs!

\bibliographystyle{plain}
\bibliography{cite}

\end{document}\documentclass{article}
\usepackage[utf8]{inputenc}
\usepackage{booktabs}
\usepackage{amsmath}
\usepackage{natbib}

\title{My Document}
\author{Your Name}
\date{\today}

\begin{document}

\maketitle

\section{Introduction}
This is a sample document demonstrating various \LaTeX{} features including mathematical equations, tables, and formatting. You can cite papers using the citation search feature, and they will be automatically added to your \texttt{cite.bib} file.

\section{Mathematical Expressions}
Here are some example equations:

\subsection{Quadratic Formula}
The quadratic formula is given by:
\[x = \frac{-b \pm \sqrt{b^2 - 4ac}}{2a}\]

\subsection{Matrix Example}
A sample matrix:
\[
\begin{pmatrix}
a_{11} & a_{12} & a_{13} \\
a_{21} & a_{22} & a_{23} \\
a_{31} & a_{32} & a_{33}
\end{pmatrix}
\]

\section{Tables}
Below is a sample table showing data:

\begin{table}[h]
\centering
\begin{tabular}{@{}lrrr@{}}
\toprule
\textbf{Category} & \textbf{Value 1} & \textbf{Value 2} & \textbf{Total} \\
\midrule
Item A & 10.5 & 20.3 & 30.8 \\
Item B & 15.2 & 25.1 & 40.3 \\
Item C & 22.1 & 12.7 & 34.8 \\
\bottomrule
\end{tabular}
\caption{Sample Data Table}
\label{tab:sample}
\end{table}

\section{Lists}
\subsection{Bullet Points}
\begin{itemize}
  \item First item
  \item Second item
  \item Third item with sub-items:
    \begin{itemize}
      \item Sub-item 1
      \item Sub-item 2
    \end{itemize}
\end{itemize}

\subsection{Numbered List}
\begin{enumerate}
  \item First step
  \item Second step
  \item Third step
\end{enumerate}

\section{Citations and References}
When you use the citation search feature, references will be automatically managed. For example, you might cite a paper like this: \cite{example2024}. All citations will be automatically added to your \texttt{cite.bib} file.

\section{Conclusion}
This template demonstrates basic \LaTeX{} features including mathematical formulas, tables, and lists. The bibliography system is now set up to work seamlessly with Bibby AI's citation management. Feel free to modify it for your needs!

\bibliographystyle{plain}
\bibliography{cite}

\end{document}\documentclass{article}
\usepackage[utf8]{inputenc}
\usepackage{booktabs}
\usepackage{amsmath}
\usepackage{natbib}

\title{My Document}
\author{Your Name}
\date{\today}

\begin{document}

\maketitle

\section{Introduction}
This is a sample document demonstrating various \LaTeX{} features including mathematical equations, tables, and formatting. You can cite papers using the citation search feature, and they will be automatically added to your \texttt{cite.bib} file.

\section{Mathematical Expressions}
Here are some example equations:

\subsection{Quadratic Formula}
The quadratic formula is given by:
\[x = \frac{-b \pm \sqrt{b^2 - 4ac}}{2a}\]

\subsection{Matrix Example}
A sample matrix:
\[
\begin{pmatrix}
a_{11} & a_{12} & a_{13} \\
a_{21} & a_{22} & a_{23} \\
a_{31} & a_{32} & a_{33}
\end{pmatrix}
\]

\section{Tables}
Below is a sample table showing data:

\begin{table}[h]
\centering
\begin{tabular}{@{}lrrr@{}}
\toprule
\textbf{Category} & \textbf{Value 1} & \textbf{Value 2} & \textbf{Total} \\
\midrule
Item A & 10.5 & 20.3 & 30.8 \\
Item B & 15.2 & 25.1 & 40.3 \\
Item C & 22.1 & 12.7 & 34.8 \\
\bottomrule
\end{tabular}
\caption{Sample Data Table}
\label{tab:sample}
\end{table}

\section{Lists}
\subsection{Bullet Points}
\begin{itemize}
  \item First item
  \item Second item
  \item Third item with sub-items:
    \begin{itemize}
      \item Sub-item 1
      \item Sub-item 2
    \end{itemize}
\end{itemize}

\subsection{Numbered List}
\begin{enumerate}
  \item First step
  \item Second step
  \item Third step
\end{enumerate}

\section{Citations and References}
When you use the citation search feature, references will be automatically managed. For example, you might cite a paper like this: \cite{example2024}. All citations will be automatically added to your \texttt{cite.bib} file.

\section{Conclusion}
This template demonstrates basic \LaTeX{} features including mathematical formulas, tables, and lists. The bibliography system is now set up to work seamlessly with Bibby AI's citation management. Feel free to modify it for your needs!

\bibliographystyle{plain}
\bibliography{cite}

\end{document}\documentclass{article}
\usepackage[utf8]{inputenc}
\usepackage{booktabs}
\usepackage{amsmath}
\usepackage{natbib}

\title{My Document}
\author{Your Name}
\date{\today}

\begin{document}

\maketitle

\section{Introduction}
This is a sample document demonstrating various \LaTeX{} features including mathematical equations, tables, and formatting. You can cite papers using the citation search feature, and they will be automatically added to your \texttt{cite.bib} file.

\section{Mathematical Expressions}
Here are some example equations:

\subsection{Quadratic Formula}
The quadratic formula is given by:
\[x = \frac{-b \pm \sqrt{b^2 - 4ac}}{2a}\]

\subsection{Matrix Example}
A sample matrix:
\[
\begin{pmatrix}
a_{11} & a_{12} & a_{13} \\
a_{21} & a_{22} & a_{23} \\
a_{31} & a_{32} & a_{33}
\end{pmatrix}
\]

\section{Tables}
Below is a sample table showing data:

\begin{table}[h]
\centering
\begin{tabular}{@{}lrrr@{}}
\toprule
\textbf{Category} & \textbf{Value 1} & \textbf{Value 2} & \textbf{Total} \\
\midrule
Item A & 10.5 & 20.3 & 30.8 \\
Item B & 15.2 & 25.1 & 40.3 \\
Item C & 22.1 & 12.7 & 34.8 \\
\bottomrule
\end{tabular}
\caption{Sample Data Table}
\label{tab:sample}
\end{table}

\section{Lists}
\subsection{Bullet Points}
\begin{itemize}
  \item First item
  \item Second item
  \item Third item with sub-items:
    \begin{itemize}
      \item Sub-item 1
      \item Sub-item 2
    \end{itemize}
\end{itemize}

\subsection{Numbered List}
\begin{enumerate}
  \item First step
  \item Second step
  \item Third step
\end{enumerate}

\section{Citations and References}
When you use the citation search feature, references will be automatically managed. For example, you might cite a paper like this: \cite{example2024}. All citations will be automatically added to your \texttt{cite.bib} file.

\section{Conclusion}
This template demonstrates basic \LaTeX{} features including mathematical formulas, tables, and lists. The bibliography system is now set up to work seamlessly with Bibby AI's citation management. Feel free to modify it for your needs!

\bibliographystyle{plain}
\bibliography{cite}

\end{document}\documentclass{article}
\usepackage[utf8]{inputenc}
\usepackage{booktabs}
\usepackage{amsmath}
\usepackage{natbib}

\title{My Document}
\author{Your Name}
\date{\today}

\begin{document}

\maketitle

\section{Introduction}
This is a sample document demonstrating various \LaTeX{} features including mathematical equations, tables, and formatting. You can cite papers using the citation search feature, and they will be automatically added to your \texttt{cite.bib} file.

\section{Mathematical Expressions}
Here are some example equations:

\subsection{Quadratic Formula}
The quadratic formula is given by:
\[x = \frac{-b \pm \sqrt{b^2 - 4ac}}{2a}\]

\subsection{Matrix Example}
A sample matrix:
\[
\begin{pmatrix}
a_{11} & a_{12} & a_{13} \\
a_{21} & a_{22} & a_{23} \\
a_{31} & a_{32} & a_{33}
\end{pmatrix}
\]

\section{Tables}
Below is a sample table showing data:

\begin{table}[h]
\centering
\begin{tabular}{@{}lrrr@{}}
\toprule
\textbf{Category} & \textbf{Value 1} & \textbf{Value 2} & \textbf{Total} \\
\midrule
Item A & 10.5 & 20.3 & 30.8 \\
Item B & 15.2 & 25.1 & 40.3 \\
Item C & 22.1 & 12.7 & 34.8 \\
\bottomrule
\end{tabular}
\caption{Sample Data Table}
\label{tab:sample}
\end{table}

\section{Lists}
\subsection{Bullet Points}
\begin{itemize}
  \item First item
  \item Second item
  \item Third item with sub-items:
    \begin{itemize}
      \item Sub-item 1
      \item Sub-item 2
    \end{itemize}
\end{itemize}

\subsection{Numbered List}
\begin{enumerate}
  \item First step
  \item Second step
  \item Third step
\end{enumerate}

\section{Citations and References}
When you use the citation search feature, references will be automatically managed. For example, you might cite a paper like this: \cite{example2024}. All citations will be automatically added to your \texttt{cite.bib} file.

\section{Conclusion}
This template demonstrates basic \LaTeX{} features including mathematical formulas, tables, and lists. The bibliography system is now set up to work seamlessly with Bibby AI's citation management. Feel free to modify it for your needs!

\bibliographystyle{plain}
\bibliography{cite}

\end{document}\documentclass{article}
\usepackage[utf8]{inputenc}
\usepackage{booktabs}
\usepackage{amsmath}
\usepackage{natbib}

\title{My Document}
\author{Your Name}
\date{\today}

\begin{document}

\maketitle

\section{Introduction}
This is a sample document demonstrating various \LaTeX{} features including mathematical equations, tables, and formatting. You can cite papers using the citation search feature, and they will be automatically added to your \texttt{cite.bib} file.

\section{Mathematical Expressions}
Here are some example equations:

\subsection{Quadratic Formula}
The quadratic formula is given by:
\[x = \frac{-b \pm \sqrt{b^2 - 4ac}}{2a}\]

\subsection{Matrix Example}
A sample matrix:
\[
\begin{pmatrix}
a_{11} & a_{12} & a_{13} \\
a_{21} & a_{22} & a_{23} \\
a_{31} & a_{32} & a_{33}
\end{pmatrix}
\]

\section{Tables}
Below is a sample table showing data:

\begin{table}[h]
\centering
\begin{tabular}{@{}lrrr@{}}
\toprule
\textbf{Category} & \textbf{Value 1} & \textbf{Value 2} & \textbf{Total} \\
\midrule
Item A & 10.5 & 20.3 & 30.8 \\
Item B & 15.2 & 25.1 & 40.3 \\
Item C & 22.1 & 12.7 & 34.8 \\
\bottomrule
\end{tabular}
\caption{Sample Data Table}
\label{tab:sample}
\end{table}

\section{Lists}
\subsection{Bullet Points}
\begin{itemize}
  \item First item
  \item Second item
  \item Third item with sub-items:
    \begin{itemize}
      \item Sub-item 1
      \item Sub-item 2
    \end{itemize}
\end{itemize}

\subsection{Numbered List}
\begin{enumerate}
  \item First step
  \item Second step
  \item Third step
\end{enumerate}

\section{Citations and References}
When you use the citation search feature, references will be automatically managed. For example, you might cite a paper like this: \cite{example2024}. All citations will be automatically added to your \texttt{cite.bib} file.

\section{Conclusion}
This template demonstrates basic \LaTeX{} features including mathematical formulas, tables, and lists. The bibliography system is now set up to work seamlessly with Bibby AI's citation management. Feel free to modify it for your needs!

\bibliographystyle{plain}
\bibliography{cite}

\end{document}\documentclass{article}
\usepackage[utf8]{inputenc}
\usepackage{booktabs}
\usepackage{amsmath}
\usepackage{natbib}

\title{My Document}
\author{Your Name}
\date{\today}

\begin{document}

\maketitle

\section{Introduction}
This is a sample document demonstrating various \LaTeX{} features including mathematical equations, tables, and formatting. You can cite papers using the citation search feature, and they will be automatically added to your \texttt{cite.bib} file.

\section{Mathematical Expressions}
Here are some example equations:

\subsection{Quadratic Formula}
The quadratic formula is given by:
\[x = \frac{-b \pm \sqrt{b^2 - 4ac}}{2a}\]

\subsection{Matrix Example}
A sample matrix:
\[
\begin{pmatrix}
a_{11} & a_{12} & a_{13} \\
a_{21} & a_{22} & a_{23} \\
a_{31} & a_{32} & a_{33}
\end{pmatrix}
\]

\section{Tables}
Below is a sample table showing data:

\begin{table}[h]
\centering
\begin{tabular}{@{}lrrr@{}}
\toprule
\textbf{Category} & \textbf{Value 1} & \textbf{Value 2} & \textbf{Total} \\
\midrule
Item A & 10.5 & 20.3 & 30.8 \\
Item B & 15.2 & 25.1 & 40.3 \\
Item C & 22.1 & 12.7 & 34.8 \\
\bottomrule
\end{tabular}
\caption{Sample Data Table}
\label{tab:sample}
\end{table}

\section{Lists}
\subsection{Bullet Points}
\begin{itemize}
  \item First item
  \item Second item
  \item Third item with sub-items:
    \begin{itemize}
      \item Sub-item 1
      \item Sub-item 2
    \end{itemize}
\end{itemize}

\subsection{Numbered List}
\begin{enumerate}
  \item First step
  \item Second step
  \item Third step
\end{enumerate}

\section{Citations and References}
When you use the citation search feature, references will be automatically managed. For example, you might cite a paper like this: \cite{example2024}. All citations will be automatically added to your \texttt{cite.bib} file.

\section{Conclusion}
This template demonstrates basic \LaTeX{} features including mathematical formulas, tables, and lists. The bibliography system is now set up to work seamlessly with Bibby AI's citation management. Feel free to modify it for your needs!

\bibliographystyle{plain}
\bibliography{cite}

\end{document}\documentclass{article}
\usepackage[utf8]{inputenc}
\usepackage{booktabs}
\usepackage{amsmath}
\usepackage{natbib}

\title{My Document}
\author{Your Name}
\date{\today}

\begin{document}

\maketitle

\section{Introduction}
This is a sample document demonstrating various \LaTeX{} features including mathematical equations, tables, and formatting. You can cite papers using the citation search feature, and they will be automatically added to your \texttt{cite.bib} file.

\section{Mathematical Expressions}
Here are some example equations:

\subsection{Quadratic Formula}
The quadratic formula is given by:
\[x = \frac{-b \pm \sqrt{b^2 - 4ac}}{2a}\]

\subsection{Matrix Example}
A sample matrix:
\[
\begin{pmatrix}
a_{11} & a_{12} & a_{13} \\
a_{21} & a_{22} & a_{23} \\
a_{31} & a_{32} & a_{33}
\end{pmatrix}
\]

\section{Tables}
Below is a sample table showing data:

\begin{table}[h]
\centering
\begin{tabular}{@{}lrrr@{}}
\toprule
\textbf{Category} & \textbf{Value 1} & \textbf{Value 2} & \textbf{Total} \\
\midrule
Item A & 10.5 & 20.3 & 30.8 \\
Item B & 15.2 & 25.1 & 40.3 \\
Item C & 22.1 & 12.7 & 34.8 \\
\bottomrule
\end{tabular}
\caption{Sample Data Table}
\label{tab:sample}
\end{table}

\section{Lists}
\subsection{Bullet Points}
\begin{itemize}
  \item First item
  \item Second item
  \item Third item with sub-items:
    \begin{itemize}
      \item Sub-item 1
      \item Sub-item 2
    \end{itemize}
\end{itemize}

\subsection{Numbered List}
\begin{enumerate}
  \item First step
  \item Second step
  \item Third step
\end{enumerate}

\section{Citations and References}
When you use the citation search feature, references will be automatically managed. For example, you might cite a paper like this: \cite{example2024}. All citations will be automatically added to your \texttt{cite.bib} file.

\section{Conclusion}
This template demonstrates basic \LaTeX{} features including mathematical formulas, tables, and lists. The bibliography system is now set up to work seamlessly with Bibby AI's citation management. Feel free to modify it for your needs!

\bibliographystyle{plain}
\bibliography{cite}

\end{document}\documentclass{article}
\usepackage[utf8]{inputenc}
\usepackage{booktabs}
\usepackage{amsmath}
\usepackage{natbib}

\title{My Document}
\author{Your Name}
\date{\today}

\begin{document}

\maketitle

\section{Introduction}
This is a sample document demonstrating various \LaTeX{} features including mathematical equations, tables, and formatting. You can cite papers using the citation search feature, and they will be automatically added to your \texttt{cite.bib} file.

\section{Mathematical Expressions}
Here are some example equations:

\subsection{Quadratic Formula}
The quadratic formula is given by:
\[x = \frac{-b \pm \sqrt{b^2 - 4ac}}{2a}\]

\subsection{Matrix Example}
A sample matrix:
\[
\begin{pmatrix}
a_{11} & a_{12} & a_{13} \\
a_{21} & a_{22} & a_{23} \\
a_{31} & a_{32} & a_{33}
\end{pmatrix}
\]

\section{Tables}
Below is a sample table showing data:

\begin{table}[h]
\centering
\begin{tabular}{@{}lrrr@{}}
\toprule
\textbf{Category} & \textbf{Value 1} & \textbf{Value 2} & \textbf{Total} \\
\midrule
Item A & 10.5 & 20.3 & 30.8 \\
Item B & 15.2 & 25.1 & 40.3 \\
Item C & 22.1 & 12.7 & 34.8 \\
\bottomrule
\end{tabular}
\caption{Sample Data Table}
\label{tab:sample}
\end{table}

\section{Lists}
\subsection{Bullet Points}
\begin{itemize}
  \item First item
  \item Second item
  \item Third item with sub-items:
    \begin{itemize}
      \item Sub-item 1
      \item Sub-item 2
    \end{itemize}
\end{itemize}

\subsection{Numbered List}
\begin{enumerate}
  \item First step
  \item Second step
  \item Third step
\end{enumerate}

\section{Citations and References}
When you use the citation search feature, references will be automatically managed. For example, you might cite a paper like this: \cite{example2024}. All citations will be automatically added to your \texttt{cite.bib} file.

\section{Conclusion}
This template demonstrates basic \LaTeX{} features including mathematical formulas, tables, and lists. The bibliography system is now set up to work seamlessly with Bibby AI's citation management. Feel free to modify it for your needs!

\bibliographystyle{plain}
\bibliography{cite}

\end{document}\documentclass{article}
\usepackage[utf8]{inputenc}
\usepackage{booktabs}
\usepackage{amsmath}
\usepackage{natbib}

\title{My Document}
\author{Your Name}
\date{\today}

\begin{document}

\maketitle

\section{Introduction}
This is a sample document demonstrating various \LaTeX{} features including mathematical equations, tables, and formatting. You can cite papers using the citation search feature, and they will be automatically added to your \texttt{cite.bib} file.

\section{Mathematical Expressions}
Here are some example equations:

\subsection{Quadratic Formula}
The quadratic formula is given by:
\[x = \frac{-b \pm \sqrt{b^2 - 4ac}}{2a}\]

\subsection{Matrix Example}
A sample matrix:
\[
\begin{pmatrix}
a_{11} & a_{12} & a_{13} \\
a_{21} & a_{22} & a_{23} \\
a_{31} & a_{32} & a_{33}
\end{pmatrix}
\]

\section{Tables}
Below is a sample table showing data:

\begin{table}[h]
\centering
\begin{tabular}{@{}lrrr@{}}
\toprule
\textbf{Category} & \textbf{Value 1} & \textbf{Value 2} & \textbf{Total} \\
\midrule
Item A & 10.5 & 20.3 & 30.8 \\
Item B & 15.2 & 25.1 & 40.3 \\
Item C & 22.1 & 12.7 & 34.8 \\
\bottomrule
\end{tabular}
\caption{Sample Data Table}
\label{tab:sample}
\end{table}

\section{Lists}
\subsection{Bullet Points}
\begin{itemize}
  \item First item
  \item Second item
  \item Third item with sub-items:
    \begin{itemize}
      \item Sub-item 1
      \item Sub-item 2
    \end{itemize}
\end{itemize}

\subsection{Numbered List}
\begin{enumerate}
  \item First step
  \item Second step
  \item Third step
\end{enumerate}

\section{Citations and References}
When you use the citation search feature, references will be automatically managed. For example, you might cite a paper like this: \cite{example2024}. All citations will be automatically added to your \texttt{cite.bib} file.

\section{Conclusion}
This template demonstrates basic \LaTeX{} features including mathematical formulas, tables, and lists. The bibliography system is now set up to work seamlessly with Bibby AI's citation management. Feel free to modify it for your needs!

\bibliographystyle{plain}
\bibliography{cite}

\end{document}\documentclass{article}
\usepackage[utf8]{inputenc}
\usepackage{booktabs}
\usepackage{amsmath}
\usepackage{natbib}

\title{My Document}
\author{Your Name}
\date{\today}

\begin{document}

\maketitle

\section{Introduction}
This is a sample document demonstrating various \LaTeX{} features including mathematical equations, tables, and formatting. You can cite papers using the citation search feature, and they will be automatically added to your \texttt{cite.bib} file.

\section{Mathematical Expressions}
Here are some example equations:

\subsection{Quadratic Formula}
The quadratic formula is given by:
\[x = \frac{-b \pm \sqrt{b^2 - 4ac}}{2a}\]

\subsection{Matrix Example}
A sample matrix:
\[
\begin{pmatrix}
a_{11} & a_{12} & a_{13} \\
a_{21} & a_{22} & a_{23} \\
a_{31} & a_{32} & a_{33}
\end{pmatrix}
\]

\section{Tables}
Below is a sample table showing data:

\begin{table}[h]
\centering
\begin{tabular}{@{}lrrr@{}}
\toprule
\textbf{Category} & \textbf{Value 1} & \textbf{Value 2} & \textbf{Total} \\
\midrule
Item A & 10.5 & 20.3 & 30.8 \\
Item B & 15.2 & 25.1 & 40.3 \\
Item C & 22.1 & 12.7 & 34.8 \\
\bottomrule
\end{tabular}
\caption{Sample Data Table}
\label{tab:sample}
\end{table}

\section{Lists}
\subsection{Bullet Points}
\begin{itemize}
  \item First item
  \item Second item
  \item Third item with sub-items:
    \begin{itemize}
      \item Sub-item 1
      \item Sub-item 2
    \end{itemize}
\end{itemize}

\subsection{Numbered List}
\begin{enumerate}
  \item First step
  \item Second step
  \item Third step
\end{enumerate}

\section{Citations and References}
When you use the citation search feature, references will be automatically managed. For example, you might cite a paper like this: \cite{example2024}. All citations will be automatically added to your \texttt{cite.bib} file.

\section{Conclusion}
This template demonstrates basic \LaTeX{} features including mathematical formulas, tables, and lists. The bibliography system is now set up to work seamlessly with Bibby AI's citation management. Feel free to modify it for your needs!

\bibliographystyle{plain}
\bibliography{cite}

\end{document}\documentclass{article}
\usepackage[utf8]{inputenc}
\usepackage{booktabs}
\usepackage{amsmath}
\usepackage{natbib}

\title{My Document}
\author{Your Name}
\date{\today}

\begin{document}

\maketitle

\section{Introduction}
This is a sample document demonstrating various \LaTeX{} features including mathematical equations, tables, and formatting. You can cite papers using the citation search feature, and they will be automatically added to your \texttt{cite.bib} file.

\section{Mathematical Expressions}
Here are some example equations:

\subsection{Quadratic Formula}
The quadratic formula is given by:
\[x = \frac{-b \pm \sqrt{b^2 - 4ac}}{2a}\]

\subsection{Matrix Example}
A sample matrix:
\[
\begin{pmatrix}
a_{11} & a_{12} & a_{13} \\
a_{21} & a_{22} & a_{23} \\
a_{31} & a_{32} & a_{33}
\end{pmatrix}
\]

\section{Tables}
Below is a sample table showing data:

\begin{table}[h]
\centering
\begin{tabular}{@{}lrrr@{}}
\toprule
\textbf{Category} & \textbf{Value 1} & \textbf{Value 2} & \textbf{Total} \\
\midrule
Item A & 10.5 & 20.3 & 30.8 \\
Item B & 15.2 & 25.1 & 40.3 \\
Item C & 22.1 & 12.7 & 34.8 \\
\bottomrule
\end{tabular}
\caption{Sample Data Table}
\label{tab:sample}
\end{table}

\section{Lists}
\subsection{Bullet Points}
\begin{itemize}
  \item First item
  \item Second item
  \item Third item with sub-items:
    \begin{itemize}
      \item Sub-item 1
      \item Sub-item 2
    \end{itemize}
\end{itemize}

\subsection{Numbered List}
\begin{enumerate}
  \item First step
  \item Second step
  \item Third step
\end{enumerate}

\section{Citations and References}
When you use the citation search feature, references will be automatically managed. For example, you might cite a paper like this: \cite{example2024}. All citations will be automatically added to your \texttt{cite.bib} file.

\section{Conclusion}
This template demonstrates basic \LaTeX{} features including mathematical formulas, tables, and lists. The bibliography system is now set up to work seamlessly with Bibby AI's citation management. Feel free to modify it for your needs!

\bibliographystyle{plain}
\bibliography{cite}

\end{document}\documentclass{article}
\usepackage[utf8]{inputenc}
\usepackage{booktabs}
\usepackage{amsmath}
\usepackage{natbib}

\title{My Document}
\author{Your Name}
\date{\today}

\begin{document}

\maketitle

\section{Introduction}
This is a sample document demonstrating various \LaTeX{} features including mathematical equations, tables, and formatting. You can cite papers using the citation search feature, and they will be automatically added to your \texttt{cite.bib} file.

\section{Mathematical Expressions}
Here are some example equations:

\subsection{Quadratic Formula}
The quadratic formula is given by:
\[x = \frac{-b \pm \sqrt{b^2 - 4ac}}{2a}\]

\subsection{Matrix Example}
A sample matrix:
\[
\begin{pmatrix}
a_{11} & a_{12} & a_{13} \\
a_{21} & a_{22} & a_{23} \\
a_{31} & a_{32} & a_{33}
\end{pmatrix}
\]

\section{Tables}
Below is a sample table showing data:

\begin{table}[h]
\centering
\begin{tabular}{@{}lrrr@{}}
\toprule
\textbf{Category} & \textbf{Value 1} & \textbf{Value 2} & \textbf{Total} \\
\midrule
Item A & 10.5 & 20.3 & 30.8 \\
Item B & 15.2 & 25.1 & 40.3 \\
Item C & 22.1 & 12.7 & 34.8 \\
\bottomrule
\end{tabular}
\caption{Sample Data Table}
\label{tab:sample}
\end{table}

\section{Lists}
\subsection{Bullet Points}
\begin{itemize}
  \item First item
  \item Second item
  \item Third item with sub-items:
    \begin{itemize}
      \item Sub-item 1
      \item Sub-item 2
    \end{itemize}
\end{itemize}

\subsection{Numbered List}
\begin{enumerate}
  \item First step
  \item Second step
  \item Third step
\end{enumerate}

\section{Citations and References}
When you use the citation search feature, references will be automatically managed. For example, you might cite a paper like this: \cite{example2024}. All citations will be automatically added to your \texttt{cite.bib} file.

\section{Conclusion}
This template demonstrates basic \LaTeX{} features including mathematical formulas, tables, and lists. The bibliography system is now set up to work seamlessly with Bibby AI's citation management. Feel free to modify it for your needs!

\bibliographystyle{plain}
\bibliography{cite}

\end{document}\documentclass{article}
\usepackage[utf8]{inputenc}
\usepackage{booktabs}
\usepackage{amsmath}
\usepackage{natbib}

\title{My Document}
\author{Your Name}
\date{\today}

\begin{document}

\maketitle

\section{Introduction}
This is a sample document demonstrating various \LaTeX{} features including mathematical equations, tables, and formatting. You can cite papers using the citation search feature, and they will be automatically added to your \texttt{cite.bib} file.

\section{Mathematical Expressions}
Here are some example equations:

\subsection{Quadratic Formula}
The quadratic formula is given by:
\[x = \frac{-b \pm \sqrt{b^2 - 4ac}}{2a}\]

\subsection{Matrix Example}
A sample matrix:
\[
\begin{pmatrix}
a_{11} & a_{12} & a_{13} \\
a_{21} & a_{22} & a_{23} \\
a_{31} & a_{32} & a_{33}
\end{pmatrix}
\]

\section{Tables}
Below is a sample table showing data:

\begin{table}[h]
\centering
\begin{tabular}{@{}lrrr@{}}
\toprule
\textbf{Category} & \textbf{Value 1} & \textbf{Value 2} & \textbf{Total} \\
\midrule
Item A & 10.5 & 20.3 & 30.8 \\
Item B & 15.2 & 25.1 & 40.3 \\
Item C & 22.1 & 12.7 & 34.8 \\
\bottomrule
\end{tabular}
\caption{Sample Data Table}
\label{tab:sample}
\end{table}

\section{Lists}
\subsection{Bullet Points}
\begin{itemize}
  \item First item
  \item Second item
  \item Third item with sub-items:
    \begin{itemize}
      \item Sub-item 1
      \item Sub-item 2
    \end{itemize}
\end{itemize}

\subsection{Numbered List}
\begin{enumerate}
  \item First step
  \item Second step
  \item Third step
\end{enumerate}

\section{Citations and References}
When you use the citation search feature, references will be automatically managed. For example, you might cite a paper like this: \cite{example2024}. All citations will be automatically added to your \texttt{cite.bib} file.

\section{Conclusion}
This template demonstrates basic \LaTeX{} features including mathematical formulas, tables, and lists. The bibliography system is now set up to work seamlessly with Bibby AI's citation management. Feel free to modify it for your needs!

\bibliographystyle{plain}
\bibliography{cite}

\end{document}\documentclass{article}
\usepackage[utf8]{inputenc}
\usepackage{booktabs}
\usepackage{amsmath}
\usepackage{natbib}

\title{My Document}
\author{Your Name}
\date{\today}

\begin{document}

\maketitle

\section{Introduction}
This is a sample document demonstrating various \LaTeX{} features including mathematical equations, tables, and formatting. You can cite papers using the citation search feature, and they will be automatically added to your \texttt{cite.bib} file.

\section{Mathematical Expressions}
Here are some example equations:

\subsection{Quadratic Formula}
The quadratic formula is given by:
\[x = \frac{-b \pm \sqrt{b^2 - 4ac}}{2a}\]

\subsection{Matrix Example}
A sample matrix:
\[
\begin{pmatrix}
a_{11} & a_{12} & a_{13} \\
a_{21} & a_{22} & a_{23} \\
a_{31} & a_{32} & a_{33}
\end{pmatrix}
\]

\section{Tables}
Below is a sample table showing data:

\begin{table}[h]
\centering
\begin{tabular}{@{}lrrr@{}}
\toprule
\textbf{Category} & \textbf{Value 1} & \textbf{Value 2} & \textbf{Total} \\
\midrule
Item A & 10.5 & 20.3 & 30.8 \\
Item B & 15.2 & 25.1 & 40.3 \\
Item C & 22.1 & 12.7 & 34.8 \\
\bottomrule
\end{tabular}
\caption{Sample Data Table}
\label{tab:sample}
\end{table}

\section{Lists}
\subsection{Bullet Points}
\begin{itemize}
  \item First item
  \item Second item
  \item Third item with sub-items:
    \begin{itemize}
      \item Sub-item 1
      \item Sub-item 2
    \end{itemize}
\end{itemize}

\subsection{Numbered List}
\begin{enumerate}
  \item First step
  \item Second step
  \item Third step
\end{enumerate}

\section{Citations and References}
When you use the citation search feature, references will be automatically managed. For example, you might cite a paper like this: \cite{example2024}. All citations will be automatically added to your \texttt{cite.bib} file.

\section{Conclusion}
This template demonstrates basic \LaTeX{} features including mathematical formulas, tables, and lists. The bibliography system is now set up to work seamlessly with Bibby AI's citation management. Feel free to modify it for your needs!

\bibliographystyle{plain}
\bibliography{cite}

\end{document}\documentclass{article}
\usepackage[utf8]{inputenc}
\usepackage{booktabs}
\usepackage{amsmath}
\usepackage{natbib}

\title{My Document}
\author{Your Name}
\date{\today}

\begin{document}

\maketitle

\section{Introduction}
This is a sample document demonstrating various \LaTeX{} features including mathematical equations, tables, and formatting. You can cite papers using the citation search feature, and they will be automatically added to your \texttt{cite.bib} file.

\section{Mathematical Expressions}
Here are some example equations:

\subsection{Quadratic Formula}
The quadratic formula is given by:
\[x = \frac{-b \pm \sqrt{b^2 - 4ac}}{2a}\]

\subsection{Matrix Example}
A sample matrix:
\[
\begin{pmatrix}
a_{11} & a_{12} & a_{13} \\
a_{21} & a_{22} & a_{23} \\
a_{31} & a_{32} & a_{33}
\end{pmatrix}
\]

\section{Tables}
Below is a sample table showing data:

\begin{table}[h]
\centering
\begin{tabular}{@{}lrrr@{}}
\toprule
\textbf{Category} & \textbf{Value 1} & \textbf{Value 2} & \textbf{Total} \\
\midrule
Item A & 10.5 & 20.3 & 30.8 \\
Item B & 15.2 & 25.1 & 40.3 \\
Item C & 22.1 & 12.7 & 34.8 \\
\bottomrule
\end{tabular}
\caption{Sample Data Table}
\label{tab:sample}
\end{table}

\section{Lists}
\subsection{Bullet Points}
\begin{itemize}
  \item First item
  \item Second item
  \item Third item with sub-items:
    \begin{itemize}
      \item Sub-item 1
      \item Sub-item 2
    \end{itemize}
\end{itemize}

\subsection{Numbered List}
\begin{enumerate}
  \item First step
  \item Second step
  \item Third step
\end{enumerate}

\section{Citations and References}
When you use the citation search feature, references will be automatically managed. For example, you might cite a paper like this: \cite{example2024}. All citations will be automatically added to your \texttt{cite.bib} file.

\section{Conclusion}
This template demonstrates basic \LaTeX{} features including mathematical formulas, tables, and lists. The bibliography system is now set up to work seamlessly with Bibby AI's citation management. Feel free to modify it for your needs!

\bibliographystyle{plain}
\bibliography{cite}

\end{document}\documentclass{article}
\usepackage[utf8]{inputenc}
\usepackage{booktabs}
\usepackage{amsmath}
\usepackage{natbib}

\title{My Document}
\author{Your Name}
\date{\today}

\begin{document}

\maketitle

\section{Introduction}
This is a sample document demonstrating various \LaTeX{} features including mathematical equations, tables, and formatting. You can cite papers using the citation search feature, and they will be automatically added to your \texttt{cite.bib} file.

\section{Mathematical Expressions}
Here are some example equations:

\subsection{Quadratic Formula}
The quadratic formula is given by:
\[x = \frac{-b \pm \sqrt{b^2 - 4ac}}{2a}\]

\subsection{Matrix Example}
A sample matrix:
\[
\begin{pmatrix}
a_{11} & a_{12} & a_{13} \\
a_{21} & a_{22} & a_{23} \\
a_{31} & a_{32} & a_{33}
\end{pmatrix}
\]

\section{Tables}
Below is a sample table showing data:

\begin{table}[h]
\centering
\begin{tabular}{@{}lrrr@{}}
\toprule
\textbf{Category} & \textbf{Value 1} & \textbf{Value 2} & \textbf{Total} \\
\midrule
Item A & 10.5 & 20.3 & 30.8 \\
Item B & 15.2 & 25.1 & 40.3 \\
Item C & 22.1 & 12.7 & 34.8 \\
\bottomrule
\end{tabular}
\caption{Sample Data Table}
\label{tab:sample}
\end{table}

\section{Lists}
\subsection{Bullet Points}
\begin{itemize}
  \item First item
  \item Second item
  \item Third item with sub-items:
    \begin{itemize}
      \item Sub-item 1
      \item Sub-item 2
    \end{itemize}
\end{itemize}

\subsection{Numbered List}
\begin{enumerate}
  \item First step
  \item Second step
  \item Third step
\end{enumerate}

\section{Citations and References}
When you use the citation search feature, references will be automatically managed. For example, you might cite a paper like this: \cite{example2024}. All citations will be automatically added to your \texttt{cite.bib} file.

\section{Conclusion}
This template demonstrates basic \LaTeX{} features including mathematical formulas, tables, and lists. The bibliography system is now set up to work seamlessly with Bibby AI's citation management. Feel free to modify it for your needs!

\bibliographystyle{plain}
\bibliography{cite}

\end{document}\documentclass{article}
\usepackage[utf8]{inputenc}
\usepackage{booktabs}
\usepackage{amsmath}
\usepackage{natbib}

\title{My Document}
\author{Your Name}
\date{\today}

\begin{document}

\maketitle

\section{Introduction}
This is a sample document demonstrating various \LaTeX{} features including mathematical equations, tables, and formatting. You can cite papers using the citation search feature, and they will be automatically added to your \texttt{cite.bib} file.

\section{Mathematical Expressions}
Here are some example equations:

\subsection{Quadratic Formula}
The quadratic formula is given by:
\[x = \frac{-b \pm \sqrt{b^2 - 4ac}}{2a}\]

\subsection{Matrix Example}
A sample matrix:
\[
\begin{pmatrix}
a_{11} & a_{12} & a_{13} \\
a_{21} & a_{22} & a_{23} \\
a_{31} & a_{32} & a_{33}
\end{pmatrix}
\]

\section{Tables}
Below is a sample table showing data:

\begin{table}[h]
\centering
\begin{tabular}{@{}lrrr@{}}
\toprule
\textbf{Category} & \textbf{Value 1} & \textbf{Value 2} & \textbf{Total} \\
\midrule
Item A & 10.5 & 20.3 & 30.8 \\
Item B & 15.2 & 25.1 & 40.3 \\
Item C & 22.1 & 12.7 & 34.8 \\
\bottomrule
\end{tabular}
\caption{Sample Data Table}
\label{tab:sample}
\end{table}

\section{Lists}
\subsection{Bullet Points}
\begin{itemize}
  \item First item
  \item Second item
  \item Third item with sub-items:
    \begin{itemize}
      \item Sub-item 1
      \item Sub-item 2
    \end{itemize}
\end{itemize}

\subsection{Numbered List}
\begin{enumerate}
  \item First step
  \item Second step
  \item Third step
\end{enumerate}

\section{Citations and References}
When you use the citation search feature, references will be automatically managed. For example, you might cite a paper like this: \cite{example2024}. All citations will be automatically added to your \texttt{cite.bib} file.

\section{Conclusion}
This template demonstrates basic \LaTeX{} features including mathematical formulas, tables, and lists. The bibliography system is now set up to work seamlessly with Bibby AI's citation management. Feel free to modify it for your needs!

\bibliographystyle{plain}
\bibliography{cite}

\end{document}\documentclass{article}
\usepackage[utf8]{inputenc}
\usepackage{booktabs}
\usepackage{amsmath}
\usepackage{natbib}

\title{My Document}
\author{Your Name}
\date{\today}

\begin{document}

\maketitle

\section{Introduction}
This is a sample document demonstrating various \LaTeX{} features including mathematical equations, tables, and formatting. You can cite papers using the citation search feature, and they will be automatically added to your \texttt{cite.bib} file.

\section{Mathematical Expressions}
Here are some example equations:

\subsection{Quadratic Formula}
The quadratic formula is given by:
\[x = \frac{-b \pm \sqrt{b^2 - 4ac}}{2a}\]

\subsection{Matrix Example}
A sample matrix:
\[
\begin{pmatrix}
a_{11} & a_{12} & a_{13} \\
a_{21} & a_{22} & a_{23} \\
a_{31} & a_{32} & a_{33}
\end{pmatrix}
\]

\section{Tables}
Below is a sample table showing data:

\begin{table}[h]
\centering
\begin{tabular}{@{}lrrr@{}}
\toprule
\textbf{Category} & \textbf{Value 1} & \textbf{Value 2} & \textbf{Total} \\
\midrule
Item A & 10.5 & 20.3 & 30.8 \\
Item B & 15.2 & 25.1 & 40.3 \\
Item C & 22.1 & 12.7 & 34.8 \\
\bottomrule
\end{tabular}
\caption{Sample Data Table}
\label{tab:sample}
\end{table}

\section{Lists}
\subsection{Bullet Points}
\begin{itemize}
  \item First item
  \item Second item
  \item Third item with sub-items:
    \begin{itemize}
      \item Sub-item 1
      \item Sub-item 2
    \end{itemize}
\end{itemize}

\subsection{Numbered List}
\begin{enumerate}
  \item First step
  \item Second step
  \item Third step
\end{enumerate}

\section{Citations and References}
When you use the citation search feature, references will be automatically managed. For example, you might cite a paper like this: \cite{example2024}. All citations will be automatically added to your \texttt{cite.bib} file.

\section{Conclusion}
This template demonstrates basic \LaTeX{} features including mathematical formulas, tables, and lists. The bibliography system is now set up to work seamlessly with Bibby AI's citation management. Feel free to modify it for your needs!

\bibliographystyle{plain}
\bibliography{cite}

\end{document}\documentclass{article}
\usepackage[utf8]{inputenc}
\usepackage{booktabs}
\usepackage{amsmath}
\usepackage{natbib}

\title{My Document}
\author{Your Name}
\date{\today}

\begin{document}

\maketitle

\section{Introduction}
This is a sample document demonstrating various \LaTeX{} features including mathematical equations, tables, and formatting. You can cite papers using the citation search feature, and they will be automatically added to your \texttt{cite.bib} file.

\section{Mathematical Expressions}
Here are some example equations:

\subsection{Quadratic Formula}
The quadratic formula is given by:
\[x = \frac{-b \pm \sqrt{b^2 - 4ac}}{2a}\]

\subsection{Matrix Example}
A sample matrix:
\[
\begin{pmatrix}
a_{11} & a_{12} & a_{13} \\
a_{21} & a_{22} & a_{23} \\
a_{31} & a_{32} & a_{33}
\end{pmatrix}
\]

\section{Tables}
Below is a sample table showing data:

\begin{table}[h]
\centering
\begin{tabular}{@{}lrrr@{}}
\toprule
\textbf{Category} & \textbf{Value 1} & \textbf{Value 2} & \textbf{Total} \\
\midrule
Item A & 10.5 & 20.3 & 30.8 \\
Item B & 15.2 & 25.1 & 40.3 \\
Item C & 22.1 & 12.7 & 34.8 \\
\bottomrule
\end{tabular}
\caption{Sample Data Table}
\label{tab:sample}
\end{table}

\section{Lists}
\subsection{Bullet Points}
\begin{itemize}
  \item First item
  \item Second item
  \item Third item with sub-items:
    \begin{itemize}
      \item Sub-item 1
      \item Sub-item 2
    \end{itemize}
\end{itemize}

\subsection{Numbered List}
\begin{enumerate}
  \item First step
  \item Second step
  \item Third step
\end{enumerate}

\section{Citations and References}
When you use the citation search feature, references will be automatically managed. For example, you might cite a paper like this: \cite{example2024}. All citations will be automatically added to your \texttt{cite.bib} file.

\section{Conclusion}
This template demonstrates basic \LaTeX{} features including mathematical formulas, tables, and lists. The bibliography system is now set up to work seamlessly with Bibby AI's citation management. Feel free to modify it for your needs!

\bibliographystyle{plain}
\bibliography{cite}

\end{document}\documentclass{article}
\usepackage[utf8]{inputenc}
\usepackage{booktabs}
\usepackage{amsmath}
\usepackage{natbib}

\title{My Document}
\author{Your Name}
\date{\today}

\begin{document}

\maketitle

\section{Introduction}
This is a sample document demonstrating various \LaTeX{} features including mathematical equations, tables, and formatting. You can cite papers using the citation search feature, and they will be automatically added to your \texttt{cite.bib} file.

\section{Mathematical Expressions}
Here are some example equations:

\subsection{Quadratic Formula}
The quadratic formula is given by:
\[x = \frac{-b \pm \sqrt{b^2 - 4ac}}{2a}\]

\subsection{Matrix Example}
A sample matrix:
\[
\begin{pmatrix}
a_{11} & a_{12} & a_{13} \\
a_{21} & a_{22} & a_{23} \\
a_{31} & a_{32} & a_{33}
\end{pmatrix}
\]

\section{Tables}
Below is a sample table showing data:

\begin{table}[h]
\centering
\begin{tabular}{@{}lrrr@{}}
\toprule
\textbf{Category} & \textbf{Value 1} & \textbf{Value 2} & \textbf{Total} \\
\midrule
Item A & 10.5 & 20.3 & 30.8 \\
Item B & 15.2 & 25.1 & 40.3 \\
Item C & 22.1 & 12.7 & 34.8 \\
\bottomrule
\end{tabular}
\caption{Sample Data Table}
\label{tab:sample}
\end{table}

\section{Lists}
\subsection{Bullet Points}
\begin{itemize}
  \item First item
  \item Second item
  \item Third item with sub-items:
    \begin{itemize}
      \item Sub-item 1
      \item Sub-item 2
    \end{itemize}
\end{itemize}

\subsection{Numbered List}
\begin{enumerate}
  \item First step
  \item Second step
  \item Third step
\end{enumerate}

\section{Citations and References}
When you use the citation search feature, references will be automatically managed. For example, you might cite a paper like this: \cite{example2024}. All citations will be automatically added to your \texttt{cite.bib} file.

\section{Conclusion}
This template demonstrates basic \LaTeX{} features including mathematical formulas, tables, and lists. The bibliography system is now set up to work seamlessly with Bibby AI's citation management. Feel free to modify it for your needs!

\bibliographystyle{plain}
\bibliography{cite}

\end{document}\documentclass{article}
\usepackage[utf8]{inputenc}
\usepackage{booktabs}
\usepackage{amsmath}
\usepackage{natbib}

\title{My Document}
\author{Your Name}
\date{\today}

\begin{document}

\maketitle

\section{Introduction}
This is a sample document demonstrating various \LaTeX{} features including mathematical equations, tables, and formatting. You can cite papers using the citation search feature, and they will be automatically added to your \texttt{cite.bib} file.

\section{Mathematical Expressions}
Here are some example equations:

\subsection{Quadratic Formula}
The quadratic formula is given by:
\[x = \frac{-b \pm \sqrt{b^2 - 4ac}}{2a}\]

\subsection{Matrix Example}
A sample matrix:
\[
\begin{pmatrix}
a_{11} & a_{12} & a_{13} \\
a_{21} & a_{22} & a_{23} \\
a_{31} & a_{32} & a_{33}
\end{pmatrix}
\]

\section{Tables}
Below is a sample table showing data:

\begin{table}[h]
\centering
\begin{tabular}{@{}lrrr@{}}
\toprule
\textbf{Category} & \textbf{Value 1} & \textbf{Value 2} & \textbf{Total} \\
\midrule
Item A & 10.5 & 20.3 & 30.8 \\
Item B & 15.2 & 25.1 & 40.3 \\
Item C & 22.1 & 12.7 & 34.8 \\
\bottomrule
\end{tabular}
\caption{Sample Data Table}
\label{tab:sample}
\end{table}

\section{Lists}
\subsection{Bullet Points}
\begin{itemize}
  \item First item
  \item Second item
  \item Third item with sub-items:
    \begin{itemize}
      \item Sub-item 1
      \item Sub-item 2
    \end{itemize}
\end{itemize}

\subsection{Numbered List}
\begin{enumerate}
  \item First step
  \item Second step
  \item Third step
\end{enumerate}

\section{Citations and References}
When you use the citation search feature, references will be automatically managed. For example, you might cite a paper like this: \cite{example2024}. All citations will be automatically added to your \texttt{cite.bib} file.

\section{Conclusion}
This template demonstrates basic \LaTeX{} features including mathematical formulas, tables, and lists. The bibliography system is now set up to work seamlessly with Bibby AI's citation management. Feel free to modify it for your needs!

\bibliographystyle{plain}
\bibliography{cite}

\end{document}\documentclass{article}
\usepackage[utf8]{inputenc}
\usepackage{booktabs}
\usepackage{amsmath}
\usepackage{natbib}

\title{My Document}
\author{Your Name}
\date{\today}

\begin{document}

\maketitle

\section{Introduction}
This is a sample document demonstrating various \LaTeX{} features including mathematical equations, tables, and formatting. You can cite papers using the citation search feature, and they will be automatically added to your \texttt{cite.bib} file.

\section{Mathematical Expressions}
Here are some example equations:

\subsection{Quadratic Formula}
The quadratic formula is given by:
\[x = \frac{-b \pm \sqrt{b^2 - 4ac}}{2a}\]

\subsection{Matrix Example}
A sample matrix:
\[
\begin{pmatrix}
a_{11} & a_{12} & a_{13} \\
a_{21} & a_{22} & a_{23} \\
a_{31} & a_{32} & a_{33}
\end{pmatrix}
\]

\section{Tables}
Below is a sample table showing data:

\begin{table}[h]
\centering
\begin{tabular}{@{}lrrr@{}}
\toprule
\textbf{Category} & \textbf{Value 1} & \textbf{Value 2} & \textbf{Total} \\
\midrule
Item A & 10.5 & 20.3 & 30.8 \\
Item B & 15.2 & 25.1 & 40.3 \\
Item C & 22.1 & 12.7 & 34.8 \\
\bottomrule
\end{tabular}
\caption{Sample Data Table}
\label{tab:sample}
\end{table}

\section{Lists}
\subsection{Bullet Points}
\begin{itemize}
  \item First item
  \item Second item
  \item Third item with sub-items:
    \begin{itemize}
      \item Sub-item 1
      \item Sub-item 2
    \end{itemize}
\end{itemize}

\subsection{Numbered List}
\begin{enumerate}
  \item First step
  \item Second step
  \item Third step
\end{enumerate}

\section{Citations and References}
When you use the citation search feature, references will be automatically managed. For example, you might cite a paper like this: \cite{example2024}. All citations will be automatically added to your \texttt{cite.bib} file.

\section{Conclusion}
This template demonstrates basic \LaTeX{} features including mathematical formulas, tables, and lists. The bibliography system is now set up to work seamlessly with Bibby AI's citation management. Feel free to modify it for your needs!

\bibliographystyle{plain}
\bibliography{cite}

\end{document}\documentclass{article}
\usepackage[utf8]{inputenc}
\usepackage{booktabs}
\usepackage{amsmath}
\usepackage{natbib}

\title{My Document}
\author{Your Name}
\date{\today}

\begin{document}

\maketitle

\section{Introduction}
This is a sample document demonstrating various \LaTeX{} features including mathematical equations, tables, and formatting. You can cite papers using the citation search feature, and they will be automatically added to your \texttt{cite.bib} file.

\section{Mathematical Expressions}
Here are some example equations:

\subsection{Quadratic Formula}
The quadratic formula is given by:
\[x = \frac{-b \pm \sqrt{b^2 - 4ac}}{2a}\]

\subsection{Matrix Example}
A sample matrix:
\[
\begin{pmatrix}
a_{11} & a_{12} & a_{13} \\
a_{21} & a_{22} & a_{23} \\
a_{31} & a_{32} & a_{33}
\end{pmatrix}
\]

\section{Tables}
Below is a sample table showing data:

\begin{table}[h]
\centering
\begin{tabular}{@{}lrrr@{}}
\toprule
\textbf{Category} & \textbf{Value 1} & \textbf{Value 2} & \textbf{Total} \\
\midrule
Item A & 10.5 & 20.3 & 30.8 \\
Item B & 15.2 & 25.1 & 40.3 \\
Item C & 22.1 & 12.7 & 34.8 \\
\bottomrule
\end{tabular}
\caption{Sample Data Table}
\label{tab:sample}
\end{table}

\section{Lists}
\subsection{Bullet Points}
\begin{itemize}
  \item First item
  \item Second item
  \item Third item with sub-items:
    \begin{itemize}
      \item Sub-item 1
      \item Sub-item 2
    \end{itemize}
\end{itemize}

\subsection{Numbered List}
\begin{enumerate}
  \item First step
  \item Second step
  \item Third step
\end{enumerate}

\section{Citations and References}
When you use the citation search feature, references will be automatically managed. For example, you might cite a paper like this: \cite{example2024}. All citations will be automatically added to your \texttt{cite.bib} file.

\section{Conclusion}
This template demonstrates basic \LaTeX{} features including mathematical formulas, tables, and lists. The bibliography system is now set up to work seamlessly with Bibby AI's citation management. Feel free to modify it for your needs!

\bibliographystyle{plain}
\bibliography{cite}

\end{document}\documentclass{article}
\usepackage[utf8]{inputenc}
\usepackage{booktabs}
\usepackage{amsmath}
\usepackage{natbib}

\title{My Document}
\author{Your Name}
\date{\today}

\begin{document}

\maketitle

\section{Introduction}
This is a sample document demonstrating various \LaTeX{} features including mathematical equations, tables, and formatting. You can cite papers using the citation search feature, and they will be automatically added to your \texttt{cite.bib} file.

\section{Mathematical Expressions}
Here are some example equations:

\subsection{Quadratic Formula}
The quadratic formula is given by:
\[x = \frac{-b \pm \sqrt{b^2 - 4ac}}{2a}\]

\subsection{Matrix Example}
A sample matrix:
\[
\begin{pmatrix}
a_{11} & a_{12} & a_{13} \\
a_{21} & a_{22} & a_{23} \\
a_{31} & a_{32} & a_{33}
\end{pmatrix}
\]

\section{Tables}
Below is a sample table showing data:

\begin{table}[h]
\centering
\begin{tabular}{@{}lrrr@{}}
\toprule
\textbf{Category} & \textbf{Value 1} & \textbf{Value 2} & \textbf{Total} \\
\midrule
Item A & 10.5 & 20.3 & 30.8 \\
Item B & 15.2 & 25.1 & 40.3 \\
Item C & 22.1 & 12.7 & 34.8 \\
\bottomrule
\end{tabular}
\caption{Sample Data Table}
\label{tab:sample}
\end{table}

\section{Lists}
\subsection{Bullet Points}
\begin{itemize}
  \item First item
  \item Second item
  \item Third item with sub-items:
    \begin{itemize}
      \item Sub-item 1
      \item Sub-item 2
    \end{itemize}
\end{itemize}

\subsection{Numbered List}
\begin{enumerate}
  \item First step
  \item Second step
  \item Third step
\end{enumerate}

\section{Citations and References}
When you use the citation search feature, references will be automatically managed. For example, you might cite a paper like this: \cite{example2024}. All citations will be automatically added to your \texttt{cite.bib} file.

\section{Conclusion}
This template demonstrates basic \LaTeX{} features including mathematical formulas, tables, and lists. The bibliography system is now set up to work seamlessly with Bibby AI's citation management. Feel free to modify it for your needs!

\bibliographystyle{plain}
\bibliography{cite}

\end{document}\documentclass{article}
\usepackage[utf8]{inputenc}
\usepackage{booktabs}
\usepackage{amsmath}
\usepackage{natbib}

\title{My Document}
\author{Your Name}
\date{\today}

\begin{document}

\maketitle

\section{Introduction}
This is a sample document demonstrating various \LaTeX{} features including mathematical equations, tables, and formatting. You can cite papers using the citation search feature, and they will be automatically added to your \texttt{cite.bib} file.

\section{Mathematical Expressions}
Here are some example equations:

\subsection{Quadratic Formula}
The quadratic formula is given by:
\[x = \frac{-b \pm \sqrt{b^2 - 4ac}}{2a}\]

\subsection{Matrix Example}
A sample matrix:
\[
\begin{pmatrix}
a_{11} & a_{12} & a_{13} \\
a_{21} & a_{22} & a_{23} \\
a_{31} & a_{32} & a_{33}
\end{pmatrix}
\]

\section{Tables}
Below is a sample table showing data:

\begin{table}[h]
\centering
\begin{tabular}{@{}lrrr@{}}
\toprule
\textbf{Category} & \textbf{Value 1} & \textbf{Value 2} & \textbf{Total} \\
\midrule
Item A & 10.5 & 20.3 & 30.8 \\
Item B & 15.2 & 25.1 & 40.3 \\
Item C & 22.1 & 12.7 & 34.8 \\
\bottomrule
\end{tabular}
\caption{Sample Data Table}
\label{tab:sample}
\end{table}

\section{Lists}
\subsection{Bullet Points}
\begin{itemize}
  \item First item
  \item Second item
  \item Third item with sub-items:
    \begin{itemize}
      \item Sub-item 1
      \item Sub-item 2
    \end{itemize}
\end{itemize}

\subsection{Numbered List}
\begin{enumerate}
  \item First step
  \item Second step
  \item Third step
\end{enumerate}

\section{Citations and References}
When you use the citation search feature, references will be automatically managed. For example, you might cite a paper like this: \cite{example2024}. All citations will be automatically added to your \texttt{cite.bib} file.

\section{Conclusion}
This template demonstrates basic \LaTeX{} features including mathematical formulas, tables, and lists. The bibliography system is now set up to work seamlessly with Bibby AI's citation management. Feel free to modify it for your needs!

\bibliographystyle{plain}
\bibliography{cite}

\end{document}\documentclass{article}
\usepackage[utf8]{inputenc}
\usepackage{booktabs}
\usepackage{amsmath}
\usepackage{natbib}

\title{My Document}
\author{Your Name}
\date{\today}

\begin{document}

\maketitle

\section{Introduction}
This is a sample document demonstrating various \LaTeX{} features including mathematical equations, tables, and formatting. You can cite papers using the citation search feature, and they will be automatically added to your \texttt{cite.bib} file.

\section{Mathematical Expressions}
Here are some example equations:

\subsection{Quadratic Formula}
The quadratic formula is given by:
\[x = \frac{-b \pm \sqrt{b^2 - 4ac}}{2a}\]

\subsection{Matrix Example}
A sample matrix:
\[
\begin{pmatrix}
a_{11} & a_{12} & a_{13} \\
a_{21} & a_{22} & a_{23} \\
a_{31} & a_{32} & a_{33}
\end{pmatrix}
\]

\section{Tables}
Below is a sample table showing data:

\begin{table}[h]
\centering
\begin{tabular}{@{}lrrr@{}}
\toprule
\textbf{Category} & \textbf{Value 1} & \textbf{Value 2} & \textbf{Total} \\
\midrule
Item A & 10.5 & 20.3 & 30.8 \\
Item B & 15.2 & 25.1 & 40.3 \\
Item C & 22.1 & 12.7 & 34.8 \\
\bottomrule
\end{tabular}
\caption{Sample Data Table}
\label{tab:sample}
\end{table}

\section{Lists}
\subsection{Bullet Points}
\begin{itemize}
  \item First item
  \item Second item
  \item Third item with sub-items:
    \begin{itemize}
      \item Sub-item 1
      \item Sub-item 2
    \end{itemize}
\end{itemize}

\subsection{Numbered List}
\begin{enumerate}
  \item First step
  \item Second step
  \item Third step
\end{enumerate}

\section{Citations and References}
When you use the citation search feature, references will be automatically managed. For example, you might cite a paper like this: \cite{example2024}. All citations will be automatically added to your \texttt{cite.bib} file.

\section{Conclusion}
This template demonstrates basic \LaTeX{} features including mathematical formulas, tables, and lists. The bibliography system is now set up to work seamlessly with Bibby AI's citation management. Feel free to modify it for your needs!

\bibliographystyle{plain}
\bibliography{cite}

\end{document}\documentclass{article}
\usepackage[utf8]{inputenc}
\usepackage{booktabs}
\usepackage{amsmath}
\usepackage{natbib}

\title{My Document}
\author{Your Name}
\date{\today}

\begin{document}

\maketitle

\section{Introduction}
This is a sample document demonstrating various \LaTeX{} features including mathematical equations, tables, and formatting. You can cite papers using the citation search feature, and they will be automatically added to your \texttt{cite.bib} file.

\section{Mathematical Expressions}
Here are some example equations:

\subsection{Quadratic Formula}
The quadratic formula is given by:
\[x = \frac{-b \pm \sqrt{b^2 - 4ac}}{2a}\]

\subsection{Matrix Example}
A sample matrix:
\[
\begin{pmatrix}
a_{11} & a_{12} & a_{13} \\
a_{21} & a_{22} & a_{23} \\
a_{31} & a_{32} & a_{33}
\end{pmatrix}
\]

\section{Tables}
Below is a sample table showing data:

\begin{table}[h]
\centering
\begin{tabular}{@{}lrrr@{}}
\toprule
\textbf{Category} & \textbf{Value 1} & \textbf{Value 2} & \textbf{Total} \\
\midrule
Item A & 10.5 & 20.3 & 30.8 \\
Item B & 15.2 & 25.1 & 40.3 \\
Item C & 22.1 & 12.7 & 34.8 \\
\bottomrule
\end{tabular}
\caption{Sample Data Table}
\label{tab:sample}
\end{table}

\section{Lists}
\subsection{Bullet Points}
\begin{itemize}
  \item First item
  \item Second item
  \item Third item with sub-items:
    \begin{itemize}
      \item Sub-item 1
      \item Sub-item 2
    \end{itemize}
\end{itemize}

\subsection{Numbered List}
\begin{enumerate}
  \item First step
  \item Second step
  \item Third step
\end{enumerate}

\section{Citations and References}
When you use the citation search feature, references will be automatically managed. For example, you might cite a paper like this: \cite{example2024}. All citations will be automatically added to your \texttt{cite.bib} file.

\section{Conclusion}
This template demonstrates basic \LaTeX{} features including mathematical formulas, tables, and lists. The bibliography system is now set up to work seamlessly with Bibby AI's citation management. Feel free to modify it for your needs!

\bibliographystyle{plain}
\bibliography{cite}

\end{document}\documentclass{article}
\usepackage[utf8]{inputenc}
\usepackage{booktabs}
\usepackage{amsmath}
\usepackage{natbib}

\title{My Document}
\author{Your Name}
\date{\today}

\begin{document}

\maketitle

\section{Introduction}
This is a sample document demonstrating various \LaTeX{} features including mathematical equations, tables, and formatting. You can cite papers using the citation search feature, and they will be automatically added to your \texttt{cite.bib} file.

\section{Mathematical Expressions}
Here are some example equations:

\subsection{Quadratic Formula}
The quadratic formula is given by:
\[x = \frac{-b \pm \sqrt{b^2 - 4ac}}{2a}\]

\subsection{Matrix Example}
A sample matrix:
\[
\begin{pmatrix}
a_{11} & a_{12} & a_{13} \\
a_{21} & a_{22} & a_{23} \\
a_{31} & a_{32} & a_{33}
\end{pmatrix}
\]

\section{Tables}
Below is a sample table showing data:

\begin{table}[h]
\centering
\begin{tabular}{@{}lrrr@{}}
\toprule
\textbf{Category} & \textbf{Value 1} & \textbf{Value 2} & \textbf{Total} \\
\midrule
Item A & 10.5 & 20.3 & 30.8 \\
Item B & 15.2 & 25.1 & 40.3 \\
Item C & 22.1 & 12.7 & 34.8 \\
\bottomrule
\end{tabular}
\caption{Sample Data Table}
\label{tab:sample}
\end{table}

\section{Lists}
\subsection{Bullet Points}
\begin{itemize}
  \item First item
  \item Second item
  \item Third item with sub-items:
    \begin{itemize}
      \item Sub-item 1
      \item Sub-item 2
    \end{itemize}
\end{itemize}

\subsection{Numbered List}
\begin{enumerate}
  \item First step
  \item Second step
  \item Third step
\end{enumerate}

\section{Citations and References}
When you use the citation search feature, references will be automatically managed. For example, you might cite a paper like this: \cite{example2024}. All citations will be automatically added to your \texttt{cite.bib} file.

\section{Conclusion}
This template demonstrates basic \LaTeX{} features including mathematical formulas, tables, and lists. The bibliography system is now set up to work seamlessly with Bibby AI's citation management. Feel free to modify it for your needs!

\bibliographystyle{plain}
\bibliography{cite}

\end{document}\documentclass{article}
\usepackage[utf8]{inputenc}
\usepackage{booktabs}
\usepackage{amsmath}
\usepackage{natbib}

\title{My Document}
\author{Your Name}
\date{\today}

\begin{document}

\maketitle

\section{Introduction}
This is a sample document demonstrating various \LaTeX{} features including mathematical equations, tables, and formatting. You can cite papers using the citation search feature, and they will be automatically added to your \texttt{cite.bib} file.

\section{Mathematical Expressions}
Here are some example equations:

\subsection{Quadratic Formula}
The quadratic formula is given by:
\[x = \frac{-b \pm \sqrt{b^2 - 4ac}}{2a}\]

\subsection{Matrix Example}
A sample matrix:
\[
\begin{pmatrix}
a_{11} & a_{12} & a_{13} \\
a_{21} & a_{22} & a_{23} \\
a_{31} & a_{32} & a_{33}
\end{pmatrix}
\]

\section{Tables}
Below is a sample table showing data:

\begin{table}[h]
\centering
\begin{tabular}{@{}lrrr@{}}
\toprule
\textbf{Category} & \textbf{Value 1} & \textbf{Value 2} & \textbf{Total} \\
\midrule
Item A & 10.5 & 20.3 & 30.8 \\
Item B & 15.2 & 25.1 & 40.3 \\
Item C & 22.1 & 12.7 & 34.8 \\
\bottomrule
\end{tabular}
\caption{Sample Data Table}
\label{tab:sample}
\end{table}

\section{Lists}
\subsection{Bullet Points}
\begin{itemize}
  \item First item
  \item Second item
  \item Third item with sub-items:
    \begin{itemize}
      \item Sub-item 1
      \item Sub-item 2
    \end{itemize}
\end{itemize}

\subsection{Numbered List}
\begin{enumerate}
  \item First step
  \item Second step
  \item Third step
\end{enumerate}

\section{Citations and References}
When you use the citation search feature, references will be automatically managed. For example, you might cite a paper like this: \cite{example2024}. All citations will be automatically added to your \texttt{cite.bib} file.

\section{Conclusion}
This template demonstrates basic \LaTeX{} features including mathematical formulas, tables, and lists. The bibliography system is now set up to work seamlessly with Bibby AI's citation management. Feel free to modify it for your needs!

\bibliographystyle{plain}
\bibliography{cite}

\end{document}\documentclass{article}
\usepackage[utf8]{inputenc}
\usepackage{booktabs}
\usepackage{amsmath}
\usepackage{natbib}

\title{My Document}
\author{Your Name}
\date{\today}

\begin{document}

\maketitle

\section{Introduction}
This is a sample document demonstrating various \LaTeX{} features including mathematical equations, tables, and formatting. You can cite papers using the citation search feature, and they will be automatically added to your \texttt{cite.bib} file.

\section{Mathematical Expressions}
Here are some example equations:

\subsection{Quadratic Formula}
The quadratic formula is given by:
\[x = \frac{-b \pm \sqrt{b^2 - 4ac}}{2a}\]

\subsection{Matrix Example}
A sample matrix:
\[
\begin{pmatrix}
a_{11} & a_{12} & a_{13} \\
a_{21} & a_{22} & a_{23} \\
a_{31} & a_{32} & a_{33}
\end{pmatrix}
\]

\section{Tables}
Below is a sample table showing data:

\begin{table}[h]
\centering
\begin{tabular}{@{}lrrr@{}}
\toprule
\textbf{Category} & \textbf{Value 1} & \textbf{Value 2} & \textbf{Total} \\
\midrule
Item A & 10.5 & 20.3 & 30.8 \\
Item B & 15.2 & 25.1 & 40.3 \\
Item C & 22.1 & 12.7 & 34.8 \\
\bottomrule
\end{tabular}
\caption{Sample Data Table}
\label{tab:sample}
\end{table}

\section{Lists}
\subsection{Bullet Points}
\begin{itemize}
  \item First item
  \item Second item
  \item Third item with sub-items:
    \begin{itemize}
      \item Sub-item 1
      \item Sub-item 2
    \end{itemize}
\end{itemize}

\subsection{Numbered List}
\begin{enumerate}
  \item First step
  \item Second step
  \item Third step
\end{enumerate}

\section{Citations and References}
When you use the citation search feature, references will be automatically managed. For example, you might cite a paper like this: \cite{example2024}. All citations will be automatically added to your \texttt{cite.bib} file.

\section{Conclusion}
This template demonstrates basic \LaTeX{} features including mathematical formulas, tables, and lists. The bibliography system is now set up to work seamlessly with Bibby AI's citation management. Feel free to modify it for your needs!

\bibliographystyle{plain}
\bibliography{cite}

\end{document}\documentclass{article}
\usepackage[utf8]{inputenc}
\usepackage{booktabs}
\usepackage{amsmath}
\usepackage{natbib}

\title{My Document}
\author{Your Name}
\date{\today}

\begin{document}

\maketitle

\section{Introduction}
This is a sample document demonstrating various \LaTeX{} features including mathematical equations, tables, and formatting. You can cite papers using the citation search feature, and they will be automatically added to your \texttt{cite.bib} file.

\section{Mathematical Expressions}
Here are some example equations:

\subsection{Quadratic Formula}
The quadratic formula is given by:
\[x = \frac{-b \pm \sqrt{b^2 - 4ac}}{2a}\]

\subsection{Matrix Example}
A sample matrix:
\[
\begin{pmatrix}
a_{11} & a_{12} & a_{13} \\
a_{21} & a_{22} & a_{23} \\
a_{31} & a_{32} & a_{33}
\end{pmatrix}
\]

\section{Tables}
Below is a sample table showing data:

\begin{table}[h]
\centering
\begin{tabular}{@{}lrrr@{}}
\toprule
\textbf{Category} & \textbf{Value 1} & \textbf{Value 2} & \textbf{Total} \\
\midrule
Item A & 10.5 & 20.3 & 30.8 \\
Item B & 15.2 & 25.1 & 40.3 \\
Item C & 22.1 & 12.7 & 34.8 \\
\bottomrule
\end{tabular}
\caption{Sample Data Table}
\label{tab:sample}
\end{table}

\section{Lists}
\subsection{Bullet Points}
\begin{itemize}
  \item First item
  \item Second item
  \item Third item with sub-items:
    \begin{itemize}
      \item Sub-item 1
      \item Sub-item 2
    \end{itemize}
\end{itemize}

\subsection{Numbered List}
\begin{enumerate}
  \item First step
  \item Second step
  \item Third step
\end{enumerate}

\section{Citations and References}
When you use the citation search feature, references will be automatically managed. For example, you might cite a paper like this: \cite{example2024}. All citations will be automatically added to your \texttt{cite.bib} file.

\section{Conclusion}
This template demonstrates basic \LaTeX{} features including mathematical formulas, tables, and lists. The bibliography system is now set up to work seamlessly with Bibby AI's citation management. Feel free to modify it for your needs!

\bibliographystyle{plain}
\bibliography{cite}

\end{document}\documentclass{article}
\usepackage[utf8]{inputenc}
\usepackage{booktabs}
\usepackage{amsmath}
\usepackage{natbib}

\title{My Document}
\author{Your Name}
\date{\today}

\begin{document}

\maketitle

\section{Introduction}
This is a sample document demonstrating various \LaTeX{} features including mathematical equations, tables, and formatting. You can cite papers using the citation search feature, and they will be automatically added to your \texttt{cite.bib} file.

\section{Mathematical Expressions}
Here are some example equations:

\subsection{Quadratic Formula}
The quadratic formula is given by:
\[x = \frac{-b \pm \sqrt{b^2 - 4ac}}{2a}\]

\subsection{Matrix Example}
A sample matrix:
\[
\begin{pmatrix}
a_{11} & a_{12} & a_{13} \\
a_{21} & a_{22} & a_{23} \\
a_{31} & a_{32} & a_{33}
\end{pmatrix}
\]

\section{Tables}
Below is a sample table showing data:

\begin{table}[h]
\centering
\begin{tabular}{@{}lrrr@{}}
\toprule
\textbf{Category} & \textbf{Value 1} & \textbf{Value 2} & \textbf{Total} \\
\midrule
Item A & 10.5 & 20.3 & 30.8 \\
Item B & 15.2 & 25.1 & 40.3 \\
Item C & 22.1 & 12.7 & 34.8 \\
\bottomrule
\end{tabular}
\caption{Sample Data Table}
\label{tab:sample}
\end{table}

\section{Lists}
\subsection{Bullet Points}
\begin{itemize}
  \item First item
  \item Second item
  \item Third item with sub-items:
    \begin{itemize}
      \item Sub-item 1
      \item Sub-item 2
    \end{itemize}
\end{itemize}

\subsection{Numbered List}
\begin{enumerate}
  \item First step
  \item Second step
  \item Third step
\end{enumerate}

\section{Citations and References}
When you use the citation search feature, references will be automatically managed. For example, you might cite a paper like this: \cite{example2024}. All citations will be automatically added to your \texttt{cite.bib} file.

\section{Conclusion}
This template demonstrates basic \LaTeX{} features including mathematical formulas, tables, and lists. The bibliography system is now set up to work seamlessly with Bibby AI's citation management. Feel free to modify it for your needs!

\bibliographystyle{plain}
\bibliography{cite}

\end{document}\documentclass{article}
\usepackage[utf8]{inputenc}
\usepackage{booktabs}
\usepackage{amsmath}
\usepackage{natbib}

\title{My Document}
\author{Your Name}
\date{\today}

\begin{document}

\maketitle

\section{Introduction}
This is a sample document demonstrating various \LaTeX{} features including mathematical equations, tables, and formatting. You can cite papers using the citation search feature, and they will be automatically added to your \texttt{cite.bib} file.

\section{Mathematical Expressions}
Here are some example equations:

\subsection{Quadratic Formula}
The quadratic formula is given by:
\[x = \frac{-b \pm \sqrt{b^2 - 4ac}}{2a}\]

\subsection{Matrix Example}
A sample matrix:
\[
\begin{pmatrix}
a_{11} & a_{12} & a_{13} \\
a_{21} & a_{22} & a_{23} \\
a_{31} & a_{32} & a_{33}
\end{pmatrix}
\]

\section{Tables}
Below is a sample table showing data:

\begin{table}[h]
\centering
\begin{tabular}{@{}lrrr@{}}
\toprule
\textbf{Category} & \textbf{Value 1} & \textbf{Value 2} & \textbf{Total} \\
\midrule
Item A & 10.5 & 20.3 & 30.8 \\
Item B & 15.2 & 25.1 & 40.3 \\
Item C & 22.1 & 12.7 & 34.8 \\
\bottomrule
\end{tabular}
\caption{Sample Data Table}
\label{tab:sample}
\end{table}

\section{Lists}
\subsection{Bullet Points}
\begin{itemize}
  \item First item
  \item Second item
  \item Third item with sub-items:
    \begin{itemize}
      \item Sub-item 1
      \item Sub-item 2
    \end{itemize}
\end{itemize}

\subsection{Numbered List}
\begin{enumerate}
  \item First step
  \item Second step
  \item Third step
\end{enumerate}

\section{Citations and References}
When you use the citation search feature, references will be automatically managed. For example, you might cite a paper like this: \cite{example2024}. All citations will be automatically added to your \texttt{cite.bib} file.

\section{Conclusion}
This template demonstrates basic \LaTeX{} features including mathematical formulas, tables, and lists. The bibliography system is now set up to work seamlessly with Bibby AI's citation management. Feel free to modify it for your needs!

\bibliographystyle{plain}
\bibliography{cite}

\end{document}\documentclass{article}
\usepackage[utf8]{inputenc}
\usepackage{booktabs}
\usepackage{amsmath}
\usepackage{natbib}

\title{My Document}
\author{Your Name}
\date{\today}

\begin{document}

\maketitle

\section{Introduction}
This is a sample document demonstrating various \LaTeX{} features including mathematical equations, tables, and formatting. You can cite papers using the citation search feature, and they will be automatically added to your \texttt{cite.bib} file.

\section{Mathematical Expressions}
Here are some example equations:

\subsection{Quadratic Formula}
The quadratic formula is given by:
\[x = \frac{-b \pm \sqrt{b^2 - 4ac}}{2a}\]

\subsection{Matrix Example}
A sample matrix:
\[
\begin{pmatrix}
a_{11} & a_{12} & a_{13} \\
a_{21} & a_{22} & a_{23} \\
a_{31} & a_{32} & a_{33}
\end{pmatrix}
\]

\section{Tables}
Below is a sample table showing data:

\begin{table}[h]
\centering
\begin{tabular}{@{}lrrr@{}}
\toprule
\textbf{Category} & \textbf{Value 1} & \textbf{Value 2} & \textbf{Total} \\
\midrule
Item A & 10.5 & 20.3 & 30.8 \\
Item B & 15.2 & 25.1 & 40.3 \\
Item C & 22.1 & 12.7 & 34.8 \\
\bottomrule
\end{tabular}
\caption{Sample Data Table}
\label{tab:sample}
\end{table}

\section{Lists}
\subsection{Bullet Points}
\begin{itemize}
  \item First item
  \item Second item
  \item Third item with sub-items:
    \begin{itemize}
      \item Sub-item 1
      \item Sub-item 2
    \end{itemize}
\end{itemize}

\subsection{Numbered List}
\begin{enumerate}
  \item First step
  \item Second step
  \item Third step
\end{enumerate}

\section{Citations and References}
When you use the citation search feature, references will be automatically managed. For example, you might cite a paper like this: \cite{example2024}. All citations will be automatically added to your \texttt{cite.bib} file.

\section{Conclusion}
This template demonstrates basic \LaTeX{} features including mathematical formulas, tables, and lists. The bibliography system is now set up to work seamlessly with Bibby AI's citation management. Feel free to modify it for your needs!

\bibliographystyle{plain}
\bibliography{cite}

\end{document}\documentclass{article}
\usepackage[utf8]{inputenc}
\usepackage{booktabs}
\usepackage{amsmath}
\usepackage{natbib}

\title{My Document}
\author{Your Name}
\date{\today}

\begin{document}

\maketitle

\section{Introduction}
This is a sample document demonstrating various \LaTeX{} features including mathematical equations, tables, and formatting. You can cite papers using the citation search feature, and they will be automatically added to your \texttt{cite.bib} file.

\section{Mathematical Expressions}
Here are some example equations:

\subsection{Quadratic Formula}
The quadratic formula is given by:
\[x = \frac{-b \pm \sqrt{b^2 - 4ac}}{2a}\]

\subsection{Matrix Example}
A sample matrix:
\[
\begin{pmatrix}
a_{11} & a_{12} & a_{13} \\
a_{21} & a_{22} & a_{23} \\
a_{31} & a_{32} & a_{33}
\end{pmatrix}
\]

\section{Tables}
Below is a sample table showing data:

\begin{table}[h]
\centering
\begin{tabular}{@{}lrrr@{}}
\toprule
\textbf{Category} & \textbf{Value 1} & \textbf{Value 2} & \textbf{Total} \\
\midrule
Item A & 10.5 & 20.3 & 30.8 \\
Item B & 15.2 & 25.1 & 40.3 \\
Item C & 22.1 & 12.7 & 34.8 \\
\bottomrule
\end{tabular}
\caption{Sample Data Table}
\label{tab:sample}
\end{table}

\section{Lists}
\subsection{Bullet Points}
\begin{itemize}
  \item First item
  \item Second item
  \item Third item with sub-items:
    \begin{itemize}
      \item Sub-item 1
      \item Sub-item 2
    \end{itemize}
\end{itemize}

\subsection{Numbered List}
\begin{enumerate}
  \item First step
  \item Second step
  \item Third step
\end{enumerate}

\section{Citations and References}
When you use the citation search feature, references will be automatically managed. For example, you might cite a paper like this: \cite{example2024}. All citations will be automatically added to your \texttt{cite.bib} file.

\section{Conclusion}
This template demonstrates basic \LaTeX{} features including mathematical formulas, tables, and lists. The bibliography system is now set up to work seamlessly with Bibby AI's citation management. Feel free to modify it for your needs!

\bibliographystyle{plain}
\bibliography{cite}

\end{document}\documentclass{article}
\usepackage[utf8]{inputenc}
\usepackage{booktabs}
\usepackage{amsmath}
\usepackage{natbib}

\title{My Document}
\author{Your Name}
\date{\today}

\begin{document}

\maketitle

\section{Introduction}
This is a sample document demonstrating various \LaTeX{} features including mathematical equations, tables, and formatting. You can cite papers using the citation search feature, and they will be automatically added to your \texttt{cite.bib} file.

\section{Mathematical Expressions}
Here are some example equations:

\subsection{Quadratic Formula}
The quadratic formula is given by:
\[x = \frac{-b \pm \sqrt{b^2 - 4ac}}{2a}\]

\subsection{Matrix Example}
A sample matrix:
\[
\begin{pmatrix}
a_{11} & a_{12} & a_{13} \\
a_{21} & a_{22} & a_{23} \\
a_{31} & a_{32} & a_{33}
\end{pmatrix}
\]

\section{Tables}
Below is a sample table showing data:

\begin{table}[h]
\centering
\begin{tabular}{@{}lrrr@{}}
\toprule
\textbf{Category} & \textbf{Value 1} & \textbf{Value 2} & \textbf{Total} \\
\midrule
Item A & 10.5 & 20.3 & 30.8 \\
Item B & 15.2 & 25.1 & 40.3 \\
Item C & 22.1 & 12.7 & 34.8 \\
\bottomrule
\end{tabular}
\caption{Sample Data Table}
\label{tab:sample}
\end{table}

\section{Lists}
\subsection{Bullet Points}
\begin{itemize}
  \item First item
  \item Second item
  \item Third item with sub-items:
    \begin{itemize}
      \item Sub-item 1
      \item Sub-item 2
    \end{itemize}
\end{itemize}

\subsection{Numbered List}
\begin{enumerate}
  \item First step
  \item Second step
  \item Third step
\end{enumerate}

\section{Citations and References}
When you use the citation search feature, references will be automatically managed. For example, you might cite a paper like this: \cite{example2024}. All citations will be automatically added to your \texttt{cite.bib} file.

\section{Conclusion}
This template demonstrates basic \LaTeX{} features including mathematical formulas, tables, and lists. The bibliography system is now set up to work seamlessly with Bibby AI's citation management. Feel free to modify it for your needs!

\bibliographystyle{plain}
\bibliography{cite}

\end{document}\documentclass{article}
\usepackage[utf8]{inputenc}
\usepackage{booktabs}
\usepackage{amsmath}
\usepackage{natbib}

\title{My Document}
\author{Your Name}
\date{\today}

\begin{document}

\maketitle

\section{Introduction}
This is a sample document demonstrating various \LaTeX{} features including mathematical equations, tables, and formatting. You can cite papers using the citation search feature, and they will be automatically added to your \texttt{cite.bib} file.

\section{Mathematical Expressions}
Here are some example equations:

\subsection{Quadratic Formula}
The quadratic formula is given by:
\[x = \frac{-b \pm \sqrt{b^2 - 4ac}}{2a}\]

\subsection{Matrix Example}
A sample matrix:
\[
\begin{pmatrix}
a_{11} & a_{12} & a_{13} \\
a_{21} & a_{22} & a_{23} \\
a_{31} & a_{32} & a_{33}
\end{pmatrix}
\]

\section{Tables}
Below is a sample table showing data:

\begin{table}[h]
\centering
\begin{tabular}{@{}lrrr@{}}
\toprule
\textbf{Category} & \textbf{Value 1} & \textbf{Value 2} & \textbf{Total} \\
\midrule
Item A & 10.5 & 20.3 & 30.8 \\
Item B & 15.2 & 25.1 & 40.3 \\
Item C & 22.1 & 12.7 & 34.8 \\
\bottomrule
\end{tabular}
\caption{Sample Data Table}
\label{tab:sample}
\end{table}

\section{Lists}
\subsection{Bullet Points}
\begin{itemize}
  \item First item
  \item Second item
  \item Third item with sub-items:
    \begin{itemize}
      \item Sub-item 1
      \item Sub-item 2
    \end{itemize}
\end{itemize}

\subsection{Numbered List}
\begin{enumerate}
  \item First step
  \item Second step
  \item Third step
\end{enumerate}

\section{Citations and References}
When you use the citation search feature, references will be automatically managed. For example, you might cite a paper like this: \cite{example2024}. All citations will be automatically added to your \texttt{cite.bib} file.

\section{Conclusion}
This template demonstrates basic \LaTeX{} features including mathematical formulas, tables, and lists. The bibliography system is now set up to work seamlessly with Bibby AI's citation management. Feel free to modify it for your needs!

\bibliographystyle{plain}
\bibliography{cite}

\end{document}\documentclass{article}
\usepackage[utf8]{inputenc}
\usepackage{booktabs}
\usepackage{amsmath}
\usepackage{natbib}

\title{My Document}
\author{Your Name}
\date{\today}

\begin{document}

\maketitle

\section{Introduction}
This is a sample document demonstrating various \LaTeX{} features including mathematical equations, tables, and formatting. You can cite papers using the citation search feature, and they will be automatically added to your \texttt{cite.bib} file.

\section{Mathematical Expressions}
Here are some example equations:

\subsection{Quadratic Formula}
The quadratic formula is given by:
\[x = \frac{-b \pm \sqrt{b^2 - 4ac}}{2a}\]

\subsection{Matrix Example}
A sample matrix:
\[
\begin{pmatrix}
a_{11} & a_{12} & a_{13} \\
a_{21} & a_{22} & a_{23} \\
a_{31} & a_{32} & a_{33}
\end{pmatrix}
\]

\section{Tables}
Below is a sample table showing data:

\begin{table}[h]
\centering
\begin{tabular}{@{}lrrr@{}}
\toprule
\textbf{Category} & \textbf{Value 1} & \textbf{Value 2} & \textbf{Total} \\
\midrule
Item A & 10.5 & 20.3 & 30.8 \\
Item B & 15.2 & 25.1 & 40.3 \\
Item C & 22.1 & 12.7 & 34.8 \\
\bottomrule
\end{tabular}
\caption{Sample Data Table}
\label{tab:sample}
\end{table}

\section{Lists}
\subsection{Bullet Points}
\begin{itemize}
  \item First item
  \item Second item
  \item Third item with sub-items:
    \begin{itemize}
      \item Sub-item 1
      \item Sub-item 2
    \end{itemize}
\end{itemize}

\subsection{Numbered List}
\begin{enumerate}
  \item First step
  \item Second step
  \item Third step
\end{enumerate}

\section{Citations and References}
When you use the citation search feature, references will be automatically managed. For example, you might cite a paper like this: \cite{example2024}. All citations will be automatically added to your \texttt{cite.bib} file.

\section{Conclusion}
This template demonstrates basic \LaTeX{} features including mathematical formulas, tables, and lists. The bibliography system is now set up to work seamlessly with Bibby AI's citation management. Feel free to modify it for your needs!

\bibliographystyle{plain}
\bibliography{cite}

\end{document}\documentclass{article}
\usepackage[utf8]{inputenc}
\usepackage{booktabs}
\usepackage{amsmath}
\usepackage{natbib}

\title{My Document}
\author{Your Name}
\date{\today}

\begin{document}

\maketitle

\section{Introduction}
This is a sample document demonstrating various \LaTeX{} features including mathematical equations, tables, and formatting. You can cite papers using the citation search feature, and they will be automatically added to your \texttt{cite.bib} file.

\section{Mathematical Expressions}
Here are some example equations:

\subsection{Quadratic Formula}
The quadratic formula is given by:
\[x = \frac{-b \pm \sqrt{b^2 - 4ac}}{2a}\]

\subsection{Matrix Example}
A sample matrix:
\[
\begin{pmatrix}
a_{11} & a_{12} & a_{13} \\
a_{21} & a_{22} & a_{23} \\
a_{31} & a_{32} & a_{33}
\end{pmatrix}
\]

\section{Tables}
Below is a sample table showing data:

\begin{table}[h]
\centering
\begin{tabular}{@{}lrrr@{}}
\toprule
\textbf{Category} & \textbf{Value 1} & \textbf{Value 2} & \textbf{Total} \\
\midrule
Item A & 10.5 & 20.3 & 30.8 \\
Item B & 15.2 & 25.1 & 40.3 \\
Item C & 22.1 & 12.7 & 34.8 \\
\bottomrule
\end{tabular}
\caption{Sample Data Table}
\label{tab:sample}
\end{table}

\section{Lists}
\subsection{Bullet Points}
\begin{itemize}
  \item First item
  \item Second item
  \item Third item with sub-items:
    \begin{itemize}
      \item Sub-item 1
      \item Sub-item 2
    \end{itemize}
\end{itemize}

\subsection{Numbered List}
\begin{enumerate}
  \item First step
  \item Second step
  \item Third step
\end{enumerate}

\section{Citations and References}
When you use the citation search feature, references will be automatically managed. For example, you might cite a paper like this: \cite{example2024}. All citations will be automatically added to your \texttt{cite.bib} file.

\section{Conclusion}
This template demonstrates basic \LaTeX{} features including mathematical formulas, tables, and lists. The bibliography system is now set up to work seamlessly with Bibby AI's citation management. Feel free to modify it for your needs!

\bibliographystyle{plain}
\bibliography{cite}

\end{document}\documentclass{article}
\usepackage[utf8]{inputenc}
\usepackage{booktabs}
\usepackage{amsmath}
\usepackage{natbib}

\title{My Document}
\author{Your Name}
\date{\today}

\begin{document}

\maketitle

\section{Introduction}
This is a sample document demonstrating various \LaTeX{} features including mathematical equations, tables, and formatting. You can cite papers using the citation search feature, and they will be automatically added to your \texttt{cite.bib} file.

\section{Mathematical Expressions}
Here are some example equations:

\subsection{Quadratic Formula}
The quadratic formula is given by:
\[x = \frac{-b \pm \sqrt{b^2 - 4ac}}{2a}\]

\subsection{Matrix Example}
A sample matrix:
\[
\begin{pmatrix}
a_{11} & a_{12} & a_{13} \\
a_{21} & a_{22} & a_{23} \\
a_{31} & a_{32} & a_{33}
\end{pmatrix}
\]

\section{Tables}
Below is a sample table showing data:

\begin{table}[h]
\centering
\begin{tabular}{@{}lrrr@{}}
\toprule
\textbf{Category} & \textbf{Value 1} & \textbf{Value 2} & \textbf{Total} \\
\midrule
Item A & 10.5 & 20.3 & 30.8 \\
Item B & 15.2 & 25.1 & 40.3 \\
Item C & 22.1 & 12.7 & 34.8 \\
\bottomrule
\end{tabular}
\caption{Sample Data Table}
\label{tab:sample}
\end{table}

\section{Lists}
\subsection{Bullet Points}
\begin{itemize}
  \item First item
  \item Second item
  \item Third item with sub-items:
    \begin{itemize}
      \item Sub-item 1
      \item Sub-item 2
    \end{itemize}
\end{itemize}

\subsection{Numbered List}
\begin{enumerate}
  \item First step
  \item Second step
  \item Third step
\end{enumerate}

\section{Citations and References}
When you use the citation search feature, references will be automatically managed. For example, you might cite a paper like this: \cite{example2024}. All citations will be automatically added to your \texttt{cite.bib} file.

\section{Conclusion}
This template demonstrates basic \LaTeX{} features including mathematical formulas, tables, and lists. The bibliography system is now set up to work seamlessly with Bibby AI's citation management. Feel free to modify it for your needs!

\bibliographystyle{plain}
\bibliography{cite}

\end{document}\documentclass{article}
\usepackage[utf8]{inputenc}
\usepackage{booktabs}
\usepackage{amsmath}
\usepackage{natbib}

\title{My Document}
\author{Your Name}
\date{\today}

\begin{document}

\maketitle

\section{Introduction}
This is a sample document demonstrating various \LaTeX{} features including mathematical equations, tables, and formatting. You can cite papers using the citation search feature, and they will be automatically added to your \texttt{cite.bib} file.

\section{Mathematical Expressions}
Here are some example equations:

\subsection{Quadratic Formula}
The quadratic formula is given by:
\[x = \frac{-b \pm \sqrt{b^2 - 4ac}}{2a}\]

\subsection{Matrix Example}
A sample matrix:
\[
\begin{pmatrix}
a_{11} & a_{12} & a_{13} \\
a_{21} & a_{22} & a_{23} \\
a_{31} & a_{32} & a_{33}
\end{pmatrix}
\]

\section{Tables}
Below is a sample table showing data:

\begin{table}[h]
\centering
\begin{tabular}{@{}lrrr@{}}
\toprule
\textbf{Category} & \textbf{Value 1} & \textbf{Value 2} & \textbf{Total} \\
\midrule
Item A & 10.5 & 20.3 & 30.8 \\
Item B & 15.2 & 25.1 & 40.3 \\
Item C & 22.1 & 12.7 & 34.8 \\
\bottomrule
\end{tabular}
\caption{Sample Data Table}
\label{tab:sample}
\end{table}

\section{Lists}
\subsection{Bullet Points}
\begin{itemize}
  \item First item
  \item Second item
  \item Third item with sub-items:
    \begin{itemize}
      \item Sub-item 1
      \item Sub-item 2
    \end{itemize}
\end{itemize}

\subsection{Numbered List}
\begin{enumerate}
  \item First step
  \item Second step
  \item Third step
\end{enumerate}

\section{Citations and References}
When you use the citation search feature, references will be automatically managed. For example, you might cite a paper like this: \cite{example2024}. All citations will be automatically added to your \texttt{cite.bib} file.

\section{Conclusion}
This template demonstrates basic \LaTeX{} features including mathematical formulas, tables, and lists. The bibliography system is now set up to work seamlessly with Bibby AI's citation management. Feel free to modify it for your needs!

\bibliographystyle{plain}
\bibliography{cite}

\end{document}\documentclass{article}
\usepackage[utf8]{inputenc}
\usepackage{booktabs}
\usepackage{amsmath}
\usepackage{natbib}

\title{My Document}
\author{Your Name}
\date{\today}

\begin{document}

\maketitle

\section{Introduction}
This is a sample document demonstrating various \LaTeX{} features including mathematical equations, tables, and formatting. You can cite papers using the citation search feature, and they will be automatically added to your \texttt{cite.bib} file.

\section{Mathematical Expressions}
Here are some example equations:

\subsection{Quadratic Formula}
The quadratic formula is given by:
\[x = \frac{-b \pm \sqrt{b^2 - 4ac}}{2a}\]

\subsection{Matrix Example}
A sample matrix:
\[
\begin{pmatrix}
a_{11} & a_{12} & a_{13} \\
a_{21} & a_{22} & a_{23} \\
a_{31} & a_{32} & a_{33}
\end{pmatrix}
\]

\section{Tables}
Below is a sample table showing data:

\begin{table}[h]
\centering
\begin{tabular}{@{}lrrr@{}}
\toprule
\textbf{Category} & \textbf{Value 1} & \textbf{Value 2} & \textbf{Total} \\
\midrule
Item A & 10.5 & 20.3 & 30.8 \\
Item B & 15.2 & 25.1 & 40.3 \\
Item C & 22.1 & 12.7 & 34.8 \\
\bottomrule
\end{tabular}
\caption{Sample Data Table}
\label{tab:sample}
\end{table}

\section{Lists}
\subsection{Bullet Points}
\begin{itemize}
  \item First item
  \item Second item
  \item Third item with sub-items:
    \begin{itemize}
      \item Sub-item 1
      \item Sub-item 2
    \end{itemize}
\end{itemize}

\subsection{Numbered List}
\begin{enumerate}
  \item First step
  \item Second step
  \item Third step
\end{enumerate}

\section{Citations and References}
When you use the citation search feature, references will be automatically managed. For example, you might cite a paper like this: \cite{example2024}. All citations will be automatically added to your \texttt{cite.bib} file.

\section{Conclusion}
This template demonstrates basic \LaTeX{} features including mathematical formulas, tables, and lists. The bibliography system is now set up to work seamlessly with Bibby AI's citation management. Feel free to modify it for your needs!

\bibliographystyle{plain}
\bibliography{cite}

\end{document}\documentclass{article}
\usepackage[utf8]{inputenc}
\usepackage{booktabs}
\usepackage{amsmath}
\usepackage{natbib}

\title{My Document}
\author{Your Name}
\date{\today}

\begin{document}

\maketitle

\section{Introduction}
This is a sample document demonstrating various \LaTeX{} features including mathematical equations, tables, and formatting. You can cite papers using the citation search feature, and they will be automatically added to your \texttt{cite.bib} file.

\section{Mathematical Expressions}
Here are some example equations:

\subsection{Quadratic Formula}
The quadratic formula is given by:
\[x = \frac{-b \pm \sqrt{b^2 - 4ac}}{2a}\]

\subsection{Matrix Example}
A sample matrix:
\[
\begin{pmatrix}
a_{11} & a_{12} & a_{13} \\
a_{21} & a_{22} & a_{23} \\
a_{31} & a_{32} & a_{33}
\end{pmatrix}
\]

\section{Tables}
Below is a sample table showing data:

\begin{table}[h]
\centering
\begin{tabular}{@{}lrrr@{}}
\toprule
\textbf{Category} & \textbf{Value 1} & \textbf{Value 2} & \textbf{Total} \\
\midrule
Item A & 10.5 & 20.3 & 30.8 \\
Item B & 15.2 & 25.1 & 40.3 \\
Item C & 22.1 & 12.7 & 34.8 \\
\bottomrule
\end{tabular}
\caption{Sample Data Table}
\label{tab:sample}
\end{table}

\section{Lists}
\subsection{Bullet Points}
\begin{itemize}
  \item First item
  \item Second item
  \item Third item with sub-items:
    \begin{itemize}
      \item Sub-item 1
      \item Sub-item 2
    \end{itemize}
\end{itemize}

\subsection{Numbered List}
\begin{enumerate}
  \item First step
  \item Second step
  \item Third step
\end{enumerate}

\section{Citations and References}
When you use the citation search feature, references will be automatically managed. For example, you might cite a paper like this: \cite{example2024}. All citations will be automatically added to your \texttt{cite.bib} file.

\section{Conclusion}
This template demonstrates basic \LaTeX{} features including mathematical formulas, tables, and lists. The bibliography system is now set up to work seamlessly with Bibby AI's citation management. Feel free to modify it for your needs!

\bibliographystyle{plain}
\bibliography{cite}

\end{document}\documentclass{article}
\usepackage[utf8]{inputenc}
\usepackage{booktabs}
\usepackage{amsmath}
\usepackage{natbib}

\title{My Document}
\author{Your Name}
\date{\today}

\begin{document}

\maketitle

\section{Introduction}
This is a sample document demonstrating various \LaTeX{} features including mathematical equations, tables, and formatting. You can cite papers using the citation search feature, and they will be automatically added to your \texttt{cite.bib} file.

\section{Mathematical Expressions}
Here are some example equations:

\subsection{Quadratic Formula}
The quadratic formula is given by:
\[x = \frac{-b \pm \sqrt{b^2 - 4ac}}{2a}\]

\subsection{Matrix Example}
A sample matrix:
\[
\begin{pmatrix}
a_{11} & a_{12} & a_{13} \\
a_{21} & a_{22} & a_{23} \\
a_{31} & a_{32} & a_{33}
\end{pmatrix}
\]

\section{Tables}
Below is a sample table showing data:

\begin{table}[h]
\centering
\begin{tabular}{@{}lrrr@{}}
\toprule
\textbf{Category} & \textbf{Value 1} & \textbf{Value 2} & \textbf{Total} \\
\midrule
Item A & 10.5 & 20.3 & 30.8 \\
Item B & 15.2 & 25.1 & 40.3 \\
Item C & 22.1 & 12.7 & 34.8 \\
\bottomrule
\end{tabular}
\caption{Sample Data Table}
\label{tab:sample}
\end{table}

\section{Lists}
\subsection{Bullet Points}
\begin{itemize}
  \item First item
  \item Second item
  \item Third item with sub-items:
    \begin{itemize}
      \item Sub-item 1
      \item Sub-item 2
    \end{itemize}
\end{itemize}

\subsection{Numbered List}
\begin{enumerate}
  \item First step
  \item Second step
  \item Third step
\end{enumerate}

\section{Citations and References}
When you use the citation search feature, references will be automatically managed. For example, you might cite a paper like this: \cite{example2024}. All citations will be automatically added to your \texttt{cite.bib} file.

\section{Conclusion}
This template demonstrates basic \LaTeX{} features including mathematical formulas, tables, and lists. The bibliography system is now set up to work seamlessly with Bibby AI's citation management. Feel free to modify it for your needs!

\bibliographystyle{plain}
\bibliography{cite}

\end{document}\documentclass{article}
\usepackage[utf8]{inputenc}
\usepackage{booktabs}
\usepackage{amsmath}
\usepackage{natbib}

\title{My Document}
\author{Your Name}
\date{\today}

\begin{document}

\maketitle

\section{Introduction}
This is a sample document demonstrating various \LaTeX{} features including mathematical equations, tables, and formatting. You can cite papers using the citation search feature, and they will be automatically added to your \texttt{cite.bib} file.

\section{Mathematical Expressions}
Here are some example equations:

\subsection{Quadratic Formula}
The quadratic formula is given by:
\[x = \frac{-b \pm \sqrt{b^2 - 4ac}}{2a}\]

\subsection{Matrix Example}
A sample matrix:
\[
\begin{pmatrix}
a_{11} & a_{12} & a_{13} \\
a_{21} & a_{22} & a_{23} \\
a_{31} & a_{32} & a_{33}
\end{pmatrix}
\]

\section{Tables}
Below is a sample table showing data:

\begin{table}[h]
\centering
\begin{tabular}{@{}lrrr@{}}
\toprule
\textbf{Category} & \textbf{Value 1} & \textbf{Value 2} & \textbf{Total} \\
\midrule
Item A & 10.5 & 20.3 & 30.8 \\
Item B & 15.2 & 25.1 & 40.3 \\
Item C & 22.1 & 12.7 & 34.8 \\
\bottomrule
\end{tabular}
\caption{Sample Data Table}
\label{tab:sample}
\end{table}

\section{Lists}
\subsection{Bullet Points}
\begin{itemize}
  \item First item
  \item Second item
  \item Third item with sub-items:
    \begin{itemize}
      \item Sub-item 1
      \item Sub-item 2
    \end{itemize}
\end{itemize}

\subsection{Numbered List}
\begin{enumerate}
  \item First step
  \item Second step
  \item Third step
\end{enumerate}

\section{Citations and References}
When you use the citation search feature, references will be automatically managed. For example, you might cite a paper like this: \cite{example2024}. All citations will be automatically added to your \texttt{cite.bib} file.

\section{Conclusion}
This template demonstrates basic \LaTeX{} features including mathematical formulas, tables, and lists. The bibliography system is now set up to work seamlessly with Bibby AI's citation management. Feel free to modify it for your needs!

\bibliographystyle{plain}
\bibliography{cite}

\end{document}\documentclass{article}
\usepackage[utf8]{inputenc}
\usepackage{booktabs}
\usepackage{amsmath}
\usepackage{natbib}

\title{My Document}
\author{Your Name}
\date{\today}

\begin{document}

\maketitle

\section{Introduction}
This is a sample document demonstrating various \LaTeX{} features including mathematical equations, tables, and formatting. You can cite papers using the citation search feature, and they will be automatically added to your \texttt{cite.bib} file.

\section{Mathematical Expressions}
Here are some example equations:

\subsection{Quadratic Formula}
The quadratic formula is given by:
\[x = \frac{-b \pm \sqrt{b^2 - 4ac}}{2a}\]

\subsection{Matrix Example}
A sample matrix:
\[
\begin{pmatrix}
a_{11} & a_{12} & a_{13} \\
a_{21} & a_{22} & a_{23} \\
a_{31} & a_{32} & a_{33}
\end{pmatrix}
\]

\section{Tables}
Below is a sample table showing data:

\begin{table}[h]
\centering
\begin{tabular}{@{}lrrr@{}}
\toprule
\textbf{Category} & \textbf{Value 1} & \textbf{Value 2} & \textbf{Total} \\
\midrule
Item A & 10.5 & 20.3 & 30.8 \\
Item B & 15.2 & 25.1 & 40.3 \\
Item C & 22.1 & 12.7 & 34.8 \\
\bottomrule
\end{tabular}
\caption{Sample Data Table}
\label{tab:sample}
\end{table}

\section{Lists}
\subsection{Bullet Points}
\begin{itemize}
  \item First item
  \item Second item
  \item Third item with sub-items:
    \begin{itemize}
      \item Sub-item 1
      \item Sub-item 2
    \end{itemize}
\end{itemize}

\subsection{Numbered List}
\begin{enumerate}
  \item First step
  \item Second step
  \item Third step
\end{enumerate}

\section{Citations and References}
When you use the citation search feature, references will be automatically managed. For example, you might cite a paper like this: \cite{example2024}. All citations will be automatically added to your \texttt{cite.bib} file.

\section{Conclusion}
This template demonstrates basic \LaTeX{} features including mathematical formulas, tables, and lists. The bibliography system is now set up to work seamlessly with Bibby AI's citation management. Feel free to modify it for your needs!

\bibliographystyle{plain}
\bibliography{cite}

\end{document}\documentclass{article}
\usepackage[utf8]{inputenc}
\usepackage{booktabs}
\usepackage{amsmath}
\usepackage{natbib}

\title{My Document}
\author{Your Name}
\date{\today}

\begin{document}

\maketitle

\section{Introduction}
This is a sample document demonstrating various \LaTeX{} features including mathematical equations, tables, and formatting. You can cite papers using the citation search feature, and they will be automatically added to your \texttt{cite.bib} file.

\section{Mathematical Expressions}
Here are some example equations:

\subsection{Quadratic Formula}
The quadratic formula is given by:
\[x = \frac{-b \pm \sqrt{b^2 - 4ac}}{2a}\]

\subsection{Matrix Example}
A sample matrix:
\[
\begin{pmatrix}
a_{11} & a_{12} & a_{13} \\
a_{21} & a_{22} & a_{23} \\
a_{31} & a_{32} & a_{33}
\end{pmatrix}
\]

\section{Tables}
Below is a sample table showing data:

\begin{table}[h]
\centering
\begin{tabular}{@{}lrrr@{}}
\toprule
\textbf{Category} & \textbf{Value 1} & \textbf{Value 2} & \textbf{Total} \\
\midrule
Item A & 10.5 & 20.3 & 30.8 \\
Item B & 15.2 & 25.1 & 40.3 \\
Item C & 22.1 & 12.7 & 34.8 \\
\bottomrule
\end{tabular}
\caption{Sample Data Table}
\label{tab:sample}
\end{table}

\section{Lists}
\subsection{Bullet Points}
\begin{itemize}
  \item First item
  \item Second item
  \item Third item with sub-items:
    \begin{itemize}
      \item Sub-item 1
      \item Sub-item 2
    \end{itemize}
\end{itemize}

\subsection{Numbered List}
\begin{enumerate}
  \item First step
  \item Second step
  \item Third step
\end{enumerate}

\section{Citations and References}
When you use the citation search feature, references will be automatically managed. For example, you might cite a paper like this: \cite{example2024}. All citations will be automatically added to your \texttt{cite.bib} file.

\section{Conclusion}
This template demonstrates basic \LaTeX{} features including mathematical formulas, tables, and lists. The bibliography system is now set up to work seamlessly with Bibby AI's citation management. Feel free to modify it for your needs!

\bibliographystyle{plain}
\bibliography{cite}

\end{document}\documentclass{article}
\usepackage[utf8]{inputenc}
\usepackage{booktabs}
\usepackage{amsmath}
\usepackage{natbib}

\title{My Document}
\author{Your Name}
\date{\today}

\begin{document}

\maketitle

\section{Introduction}
This is a sample document demonstrating various \LaTeX{} features including mathematical equations, tables, and formatting. You can cite papers using the citation search feature, and they will be automatically added to your \texttt{cite.bib} file.

\section{Mathematical Expressions}
Here are some example equations:

\subsection{Quadratic Formula}
The quadratic formula is given by:
\[x = \frac{-b \pm \sqrt{b^2 - 4ac}}{2a}\]

\subsection{Matrix Example}
A sample matrix:
\[
\begin{pmatrix}
a_{11} & a_{12} & a_{13} \\
a_{21} & a_{22} & a_{23} \\
a_{31} & a_{32} & a_{33}
\end{pmatrix}
\]

\section{Tables}
Below is a sample table showing data:

\begin{table}[h]
\centering
\begin{tabular}{@{}lrrr@{}}
\toprule
\textbf{Category} & \textbf{Value 1} & \textbf{Value 2} & \textbf{Total} \\
\midrule
Item A & 10.5 & 20.3 & 30.8 \\
Item B & 15.2 & 25.1 & 40.3 \\
Item C & 22.1 & 12.7 & 34.8 \\
\bottomrule
\end{tabular}
\caption{Sample Data Table}
\label{tab:sample}
\end{table}

\section{Lists}
\subsection{Bullet Points}
\begin{itemize}
  \item First item
  \item Second item
  \item Third item with sub-items:
    \begin{itemize}
      \item Sub-item 1
      \item Sub-item 2
    \end{itemize}
\end{itemize}

\subsection{Numbered List}
\begin{enumerate}
  \item First step
  \item Second step
  \item Third step
\end{enumerate}

\section{Citations and References}
When you use the citation search feature, references will be automatically managed. For example, you might cite a paper like this: \cite{example2024}. All citations will be automatically added to your \texttt{cite.bib} file.

\section{Conclusion}
This template demonstrates basic \LaTeX{} features including mathematical formulas, tables, and lists. The bibliography system is now set up to work seamlessly with Bibby AI's citation management. Feel free to modify it for your needs!

\bibliographystyle{plain}
\bibliography{cite}

\end{document}\documentclass{article}
\usepackage[utf8]{inputenc}
\usepackage{booktabs}
\usepackage{amsmath}
\usepackage{natbib}

\title{My Document}
\author{Your Name}
\date{\today}

\begin{document}

\maketitle

\section{Introduction}
This is a sample document demonstrating various \LaTeX{} features including mathematical equations, tables, and formatting. You can cite papers using the citation search feature, and they will be automatically added to your \texttt{cite.bib} file.

\section{Mathematical Expressions}
Here are some example equations:

\subsection{Quadratic Formula}
The quadratic formula is given by:
\[x = \frac{-b \pm \sqrt{b^2 - 4ac}}{2a}\]

\subsection{Matrix Example}
A sample matrix:
\[
\begin{pmatrix}
a_{11} & a_{12} & a_{13} \\
a_{21} & a_{22} & a_{23} \\
a_{31} & a_{32} & a_{33}
\end{pmatrix}
\]

\section{Tables}
Below is a sample table showing data:

\begin{table}[h]
\centering
\begin{tabular}{@{}lrrr@{}}
\toprule
\textbf{Category} & \textbf{Value 1} & \textbf{Value 2} & \textbf{Total} \\
\midrule
Item A & 10.5 & 20.3 & 30.8 \\
Item B & 15.2 & 25.1 & 40.3 \\
Item C & 22.1 & 12.7 & 34.8 \\
\bottomrule
\end{tabular}
\caption{Sample Data Table}
\label{tab:sample}
\end{table}

\section{Lists}
\subsection{Bullet Points}
\begin{itemize}
  \item First item
  \item Second item
  \item Third item with sub-items:
    \begin{itemize}
      \item Sub-item 1
      \item Sub-item 2
    \end{itemize}
\end{itemize}

\subsection{Numbered List}
\begin{enumerate}
  \item First step
  \item Second step
  \item Third step
\end{enumerate}

\section{Citations and References}
When you use the citation search feature, references will be automatically managed. For example, you might cite a paper like this: \cite{example2024}. All citations will be automatically added to your \texttt{cite.bib} file.

\section{Conclusion}
This template demonstrates basic \LaTeX{} features including mathematical formulas, tables, and lists. The bibliography system is now set up to work seamlessly with Bibby AI's citation management. Feel free to modify it for your needs!

\bibliographystyle{plain}
\bibliography{cite}

\end{document}\documentclass{article}
\usepackage[utf8]{inputenc}
\usepackage{booktabs}
\usepackage{amsmath}
\usepackage{natbib}

\title{My Document}
\author{Your Name}
\date{\today}

\begin{document}

\maketitle

\section{Introduction}
This is a sample document demonstrating various \LaTeX{} features including mathematical equations, tables, and formatting. You can cite papers using the citation search feature, and they will be automatically added to your \texttt{cite.bib} file.

\section{Mathematical Expressions}
Here are some example equations:

\subsection{Quadratic Formula}
The quadratic formula is given by:
\[x = \frac{-b \pm \sqrt{b^2 - 4ac}}{2a}\]

\subsection{Matrix Example}
A sample matrix:
\[
\begin{pmatrix}
a_{11} & a_{12} & a_{13} \\
a_{21} & a_{22} & a_{23} \\
a_{31} & a_{32} & a_{33}
\end{pmatrix}
\]

\section{Tables}
Below is a sample table showing data:

\begin{table}[h]
\centering
\begin{tabular}{@{}lrrr@{}}
\toprule
\textbf{Category} & \textbf{Value 1} & \textbf{Value 2} & \textbf{Total} \\
\midrule
Item A & 10.5 & 20.3 & 30.8 \\
Item B & 15.2 & 25.1 & 40.3 \\
Item C & 22.1 & 12.7 & 34.8 \\
\bottomrule
\end{tabular}
\caption{Sample Data Table}
\label{tab:sample}
\end{table}

\section{Lists}
\subsection{Bullet Points}
\begin{itemize}
  \item First item
  \item Second item
  \item Third item with sub-items:
    \begin{itemize}
      \item Sub-item 1
      \item Sub-item 2
    \end{itemize}
\end{itemize}

\subsection{Numbered List}
\begin{enumerate}
  \item First step
  \item Second step
  \item Third step
\end{enumerate}

\section{Citations and References}
When you use the citation search feature, references will be automatically managed. For example, you might cite a paper like this: \cite{example2024}. All citations will be automatically added to your \texttt{cite.bib} file.

\section{Conclusion}
This template demonstrates basic \LaTeX{} features including mathematical formulas, tables, and lists. The bibliography system is now set up to work seamlessly with Bibby AI's citation management. Feel free to modify it for your needs!

\bibliographystyle{plain}
\bibliography{cite}

\end{document}\documentclass{article}
\usepackage[utf8]{inputenc}
\usepackage{booktabs}
\usepackage{amsmath}
\usepackage{natbib}

\title{My Document}
\author{Your Name}
\date{\today}

\begin{document}

\maketitle

\section{Introduction}
This is a sample document demonstrating various \LaTeX{} features including mathematical equations, tables, and formatting. You can cite papers using the citation search feature, and they will be automatically added to your \texttt{cite.bib} file.

\section{Mathematical Expressions}
Here are some example equations:

\subsection{Quadratic Formula}
The quadratic formula is given by:
\[x = \frac{-b \pm \sqrt{b^2 - 4ac}}{2a}\]

\subsection{Matrix Example}
A sample matrix:
\[
\begin{pmatrix}
a_{11} & a_{12} & a_{13} \\
a_{21} & a_{22} & a_{23} \\
a_{31} & a_{32} & a_{33}
\end{pmatrix}
\]

\section{Tables}
Below is a sample table showing data:

\begin{table}[h]
\centering
\begin{tabular}{@{}lrrr@{}}
\toprule
\textbf{Category} & \textbf{Value 1} & \textbf{Value 2} & \textbf{Total} \\
\midrule
Item A & 10.5 & 20.3 & 30.8 \\
Item B & 15.2 & 25.1 & 40.3 \\
Item C & 22.1 & 12.7 & 34.8 \\
\bottomrule
\end{tabular}
\caption{Sample Data Table}
\label{tab:sample}
\end{table}

\section{Lists}
\subsection{Bullet Points}
\begin{itemize}
  \item First item
  \item Second item
  \item Third item with sub-items:
    \begin{itemize}
      \item Sub-item 1
      \item Sub-item 2
    \end{itemize}
\end{itemize}

\subsection{Numbered List}
\begin{enumerate}
  \item First step
  \item Second step
  \item Third step
\end{enumerate}

\section{Citations and References}
When you use the citation search feature, references will be automatically managed. For example, you might cite a paper like this: \cite{example2024}. All citations will be automatically added to your \texttt{cite.bib} file.

\section{Conclusion}
This template demonstrates basic \LaTeX{} features including mathematical formulas, tables, and lists. The bibliography system is now set up to work seamlessly with Bibby AI's citation management. Feel free to modify it for your needs!

\bibliographystyle{plain}
\bibliography{cite}

\end{document}\documentclass{article}
\usepackage[utf8]{inputenc}
\usepackage{booktabs}
\usepackage{amsmath}
\usepackage{natbib}

\title{My Document}
\author{Your Name}
\date{\today}

\begin{document}

\maketitle

\section{Introduction}
This is a sample document demonstrating various \LaTeX{} features including mathematical equations, tables, and formatting. You can cite papers using the citation search feature, and they will be automatically added to your \texttt{cite.bib} file.

\section{Mathematical Expressions}
Here are some example equations:

\subsection{Quadratic Formula}
The quadratic formula is given by:
\[x = \frac{-b \pm \sqrt{b^2 - 4ac}}{2a}\]

\subsection{Matrix Example}
A sample matrix:
\[
\begin{pmatrix}
a_{11} & a_{12} & a_{13} \\
a_{21} & a_{22} & a_{23} \\
a_{31} & a_{32} & a_{33}
\end{pmatrix}
\]

\section{Tables}
Below is a sample table showing data:

\begin{table}[h]
\centering
\begin{tabular}{@{}lrrr@{}}
\toprule
\textbf{Category} & \textbf{Value 1} & \textbf{Value 2} & \textbf{Total} \\
\midrule
Item A & 10.5 & 20.3 & 30.8 \\
Item B & 15.2 & 25.1 & 40.3 \\
Item C & 22.1 & 12.7 & 34.8 \\
\bottomrule
\end{tabular}
\caption{Sample Data Table}
\label{tab:sample}
\end{table}

\section{Lists}
\subsection{Bullet Points}
\begin{itemize}
  \item First item
  \item Second item
  \item Third item with sub-items:
    \begin{itemize}
      \item Sub-item 1
      \item Sub-item 2
    \end{itemize}
\end{itemize}

\subsection{Numbered List}
\begin{enumerate}
  \item First step
  \item Second step
  \item Third step
\end{enumerate}

\section{Citations and References}
When you use the citation search feature, references will be automatically managed. For example, you might cite a paper like this: \cite{example2024}. All citations will be automatically added to your \texttt{cite.bib} file.

\section{Conclusion}
This template demonstrates basic \LaTeX{} features including mathematical formulas, tables, and lists. The bibliography system is now set up to work seamlessly with Bibby AI's citation management. Feel free to modify it for your needs!

\bibliographystyle{plain}
\bibliography{cite}

\end{document}\documentclass{article}
\usepackage[utf8]{inputenc}
\usepackage{booktabs}
\usepackage{amsmath}
\usepackage{natbib}

\title{My Document}
\author{Your Name}
\date{\today}

\begin{document}

\maketitle

\section{Introduction}
This is a sample document demonstrating various \LaTeX{} features including mathematical equations, tables, and formatting. You can cite papers using the citation search feature, and they will be automatically added to your \texttt{cite.bib} file.

\section{Mathematical Expressions}
Here are some example equations:

\subsection{Quadratic Formula}
The quadratic formula is given by:
\[x = \frac{-b \pm \sqrt{b^2 - 4ac}}{2a}\]

\subsection{Matrix Example}
A sample matrix:
\[
\begin{pmatrix}
a_{11} & a_{12} & a_{13} \\
a_{21} & a_{22} & a_{23} \\
a_{31} & a_{32} & a_{33}
\end{pmatrix}
\]

\section{Tables}
Below is a sample table showing data:

\begin{table}[h]
\centering
\begin{tabular}{@{}lrrr@{}}
\toprule
\textbf{Category} & \textbf{Value 1} & \textbf{Value 2} & \textbf{Total} \\
\midrule
Item A & 10.5 & 20.3 & 30.8 \\
Item B & 15.2 & 25.1 & 40.3 \\
Item C & 22.1 & 12.7 & 34.8 \\
\bottomrule
\end{tabular}
\caption{Sample Data Table}
\label{tab:sample}
\end{table}

\section{Lists}
\subsection{Bullet Points}
\begin{itemize}
  \item First item
  \item Second item
  \item Third item with sub-items:
    \begin{itemize}
      \item Sub-item 1
      \item Sub-item 2
    \end{itemize}
\end{itemize}

\subsection{Numbered List}
\begin{enumerate}
  \item First step
  \item Second step
  \item Third step
\end{enumerate}

\section{Citations and References}
When you use the citation search feature, references will be automatically managed. For example, you might cite a paper like this: \cite{example2024}. All citations will be automatically added to your \texttt{cite.bib} file.

\section{Conclusion}
This template demonstrates basic \LaTeX{} features including mathematical formulas, tables, and lists. The bibliography system is now set up to work seamlessly with Bibby AI's citation management. Feel free to modify it for your needs!

\bibliographystyle{plain}
\bibliography{cite}

\end{document}\documentclass{article}
\usepackage[utf8]{inputenc}
\usepackage{booktabs}
\usepackage{amsmath}
\usepackage{natbib}

\title{My Document}
\author{Your Name}
\date{\today}

\begin{document}

\maketitle

\section{Introduction}
This is a sample document demonstrating various \LaTeX{} features including mathematical equations, tables, and formatting. You can cite papers using the citation search feature, and they will be automatically added to your \texttt{cite.bib} file.

\section{Mathematical Expressions}
Here are some example equations:

\subsection{Quadratic Formula}
The quadratic formula is given by:
\[x = \frac{-b \pm \sqrt{b^2 - 4ac}}{2a}\]

\subsection{Matrix Example}
A sample matrix:
\[
\begin{pmatrix}
a_{11} & a_{12} & a_{13} \\
a_{21} & a_{22} & a_{23} \\
a_{31} & a_{32} & a_{33}
\end{pmatrix}
\]

\section{Tables}
Below is a sample table showing data:

\begin{table}[h]
\centering
\begin{tabular}{@{}lrrr@{}}
\toprule
\textbf{Category} & \textbf{Value 1} & \textbf{Value 2} & \textbf{Total} \\
\midrule
Item A & 10.5 & 20.3 & 30.8 \\
Item B & 15.2 & 25.1 & 40.3 \\
Item C & 22.1 & 12.7 & 34.8 \\
\bottomrule
\end{tabular}
\caption{Sample Data Table}
\label{tab:sample}
\end{table}

\section{Lists}
\subsection{Bullet Points}
\begin{itemize}
  \item First item
  \item Second item
  \item Third item with sub-items:
    \begin{itemize}
      \item Sub-item 1
      \item Sub-item 2
    \end{itemize}
\end{itemize}

\subsection{Numbered List}
\begin{enumerate}
  \item First step
  \item Second step
  \item Third step
\end{enumerate}

\section{Citations and References}
When you use the citation search feature, references will be automatically managed. For example, you might cite a paper like this: \cite{example2024}. All citations will be automatically added to your \texttt{cite.bib} file.

\section{Conclusion}
This template demonstrates basic \LaTeX{} features including mathematical formulas, tables, and lists. The bibliography system is now set up to work seamlessly with Bibby AI's citation management. Feel free to modify it for your needs!

\bibliographystyle{plain}
\bibliography{cite}

\end{document}\documentclass{article}
\usepackage[utf8]{inputenc}
\usepackage{booktabs}
\usepackage{amsmath}
\usepackage{natbib}

\title{My Document}
\author{Your Name}
\date{\today}

\begin{document}

\maketitle

\section{Introduction}
This is a sample document demonstrating various \LaTeX{} features including mathematical equations, tables, and formatting. You can cite papers using the citation search feature, and they will be automatically added to your \texttt{cite.bib} file.

\section{Mathematical Expressions}
Here are some example equations:

\subsection{Quadratic Formula}
The quadratic formula is given by:
\[x = \frac{-b \pm \sqrt{b^2 - 4ac}}{2a}\]

\subsection{Matrix Example}
A sample matrix:
\[
\begin{pmatrix}
a_{11} & a_{12} & a_{13} \\
a_{21} & a_{22} & a_{23} \\
a_{31} & a_{32} & a_{33}
\end{pmatrix}
\]

\section{Tables}
Below is a sample table showing data:

\begin{table}[h]
\centering
\begin{tabular}{@{}lrrr@{}}
\toprule
\textbf{Category} & \textbf{Value 1} & \textbf{Value 2} & \textbf{Total} \\
\midrule
Item A & 10.5 & 20.3 & 30.8 \\
Item B & 15.2 & 25.1 & 40.3 \\
Item C & 22.1 & 12.7 & 34.8 \\
\bottomrule
\end{tabular}
\caption{Sample Data Table}
\label{tab:sample}
\end{table}

\section{Lists}
\subsection{Bullet Points}
\begin{itemize}
  \item First item
  \item Second item
  \item Third item with sub-items:
    \begin{itemize}
      \item Sub-item 1
      \item Sub-item 2
    \end{itemize}
\end{itemize}

\subsection{Numbered List}
\begin{enumerate}
  \item First step
  \item Second step
  \item Third step
\end{enumerate}

\section{Citations and References}
When you use the citation search feature, references will be automatically managed. For example, you might cite a paper like this: \cite{example2024}. All citations will be automatically added to your \texttt{cite.bib} file.

\section{Conclusion}
This template demonstrates basic \LaTeX{} features including mathematical formulas, tables, and lists. The bibliography system is now set up to work seamlessly with Bibby AI's citation management. Feel free to modify it for your needs!

\bibliographystyle{plain}
\bibliography{cite}

\end{document}\documentclass{article}
\usepackage[utf8]{inputenc}
\usepackage{booktabs}
\usepackage{amsmath}
\usepackage{natbib}

\title{My Document}
\author{Your Name}
\date{\today}

\begin{document}

\maketitle

\section{Introduction}
This is a sample document demonstrating various \LaTeX{} features including mathematical equations, tables, and formatting. You can cite papers using the citation search feature, and they will be automatically added to your \texttt{cite.bib} file.

\section{Mathematical Expressions}
Here are some example equations:

\subsection{Quadratic Formula}
The quadratic formula is given by:
\[x = \frac{-b \pm \sqrt{b^2 - 4ac}}{2a}\]

\subsection{Matrix Example}
A sample matrix:
\[
\begin{pmatrix}
a_{11} & a_{12} & a_{13} \\
a_{21} & a_{22} & a_{23} \\
a_{31} & a_{32} & a_{33}
\end{pmatrix}
\]

\section{Tables}
Below is a sample table showing data:

\begin{table}[h]
\centering
\begin{tabular}{@{}lrrr@{}}
\toprule
\textbf{Category} & \textbf{Value 1} & \textbf{Value 2} & \textbf{Total} \\
\midrule
Item A & 10.5 & 20.3 & 30.8 \\
Item B & 15.2 & 25.1 & 40.3 \\
Item C & 22.1 & 12.7 & 34.8 \\
\bottomrule
\end{tabular}
\caption{Sample Data Table}
\label{tab:sample}
\end{table}

\section{Lists}
\subsection{Bullet Points}
\begin{itemize}
  \item First item
  \item Second item
  \item Third item with sub-items:
    \begin{itemize}
      \item Sub-item 1
      \item Sub-item 2
    \end{itemize}
\end{itemize}

\subsection{Numbered List}
\begin{enumerate}
  \item First step
  \item Second step
  \item Third step
\end{enumerate}

\section{Citations and References}
When you use the citation search feature, references will be automatically managed. For example, you might cite a paper like this: \cite{example2024}. All citations will be automatically added to your \texttt{cite.bib} file.

\section{Conclusion}
This template demonstrates basic \LaTeX{} features including mathematical formulas, tables, and lists. The bibliography system is now set up to work seamlessly with Bibby AI's citation management. Feel free to modify it for your needs!

\bibliographystyle{plain}
\bibliography{cite}

\end{document}\documentclass{article}
\usepackage[utf8]{inputenc}
\usepackage{booktabs}
\usepackage{amsmath}
\usepackage{natbib}

\title{My Document}
\author{Your Name}
\date{\today}

\begin{document}

\maketitle

\section{Introduction}
This is a sample document demonstrating various \LaTeX{} features including mathematical equations, tables, and formatting. You can cite papers using the citation search feature, and they will be automatically added to your \texttt{cite.bib} file.

\section{Mathematical Expressions}
Here are some example equations:

\subsection{Quadratic Formula}
The quadratic formula is given by:
\[x = \frac{-b \pm \sqrt{b^2 - 4ac}}{2a}\]

\subsection{Matrix Example}
A sample matrix:
\[
\begin{pmatrix}
a_{11} & a_{12} & a_{13} \\
a_{21} & a_{22} & a_{23} \\
a_{31} & a_{32} & a_{33}
\end{pmatrix}
\]

\section{Tables}
Below is a sample table showing data:

\begin{table}[h]
\centering
\begin{tabular}{@{}lrrr@{}}
\toprule
\textbf{Category} & \textbf{Value 1} & \textbf{Value 2} & \textbf{Total} \\
\midrule
Item A & 10.5 & 20.3 & 30.8 \\
Item B & 15.2 & 25.1 & 40.3 \\
Item C & 22.1 & 12.7 & 34.8 \\
\bottomrule
\end{tabular}
\caption{Sample Data Table}
\label{tab:sample}
\end{table}

\section{Lists}
\subsection{Bullet Points}
\begin{itemize}
  \item First item
  \item Second item
  \item Third item with sub-items:
    \begin{itemize}
      \item Sub-item 1
      \item Sub-item 2
    \end{itemize}
\end{itemize}

\subsection{Numbered List}
\begin{enumerate}
  \item First step
  \item Second step
  \item Third step
\end{enumerate}

\section{Citations and References}
When you use the citation search feature, references will be automatically managed. For example, you might cite a paper like this: \cite{example2024}. All citations will be automatically added to your \texttt{cite.bib} file.

\section{Conclusion}
This template demonstrates basic \LaTeX{} features including mathematical formulas, tables, and lists. The bibliography system is now set up to work seamlessly with Bibby AI's citation management. Feel free to modify it for your needs!

\bibliographystyle{plain}
\bibliography{cite}

\end{document}\documentclass{article}
\usepackage[utf8]{inputenc}
\usepackage{booktabs}
\usepackage{amsmath}
\usepackage{natbib}

\title{My Document}
\author{Your Name}
\date{\today}

\begin{document}

\maketitle

\section{Introduction}
This is a sample document demonstrating various \LaTeX{} features including mathematical equations, tables, and formatting. You can cite papers using the citation search feature, and they will be automatically added to your \texttt{cite.bib} file.

\section{Mathematical Expressions}
Here are some example equations:

\subsection{Quadratic Formula}
The quadratic formula is given by:
\[x = \frac{-b \pm \sqrt{b^2 - 4ac}}{2a}\]

\subsection{Matrix Example}
A sample matrix:
\[
\begin{pmatrix}
a_{11} & a_{12} & a_{13} \\
a_{21} & a_{22} & a_{23} \\
a_{31} & a_{32} & a_{33}
\end{pmatrix}
\]

\section{Tables}
Below is a sample table showing data:

\begin{table}[h]
\centering
\begin{tabular}{@{}lrrr@{}}
\toprule
\textbf{Category} & \textbf{Value 1} & \textbf{Value 2} & \textbf{Total} \\
\midrule
Item A & 10.5 & 20.3 & 30.8 \\
Item B & 15.2 & 25.1 & 40.3 \\
Item C & 22.1 & 12.7 & 34.8 \\
\bottomrule
\end{tabular}
\caption{Sample Data Table}
\label{tab:sample}
\end{table}

\section{Lists}
\subsection{Bullet Points}
\begin{itemize}
  \item First item
  \item Second item
  \item Third item with sub-items:
    \begin{itemize}
      \item Sub-item 1
      \item Sub-item 2
    \end{itemize}
\end{itemize}

\subsection{Numbered List}
\begin{enumerate}
  \item First step
  \item Second step
  \item Third step
\end{enumerate}

\section{Citations and References}
When you use the citation search feature, references will be automatically managed. For example, you might cite a paper like this: \cite{example2024}. All citations will be automatically added to your \texttt{cite.bib} file.

\section{Conclusion}
This template demonstrates basic \LaTeX{} features including mathematical formulas, tables, and lists. The bibliography system is now set up to work seamlessly with Bibby AI's citation management. Feel free to modify it for your needs!

\bibliographystyle{plain}
\bibliography{cite}

\end{document}\documentclass{article}
\usepackage[utf8]{inputenc}
\usepackage{booktabs}
\usepackage{amsmath}
\usepackage{natbib}

\title{My Document}
\author{Your Name}
\date{\today}

\begin{document}

\maketitle

\section{Introduction}
This is a sample document demonstrating various \LaTeX{} features including mathematical equations, tables, and formatting. You can cite papers using the citation search feature, and they will be automatically added to your \texttt{cite.bib} file.

\section{Mathematical Expressions}
Here are some example equations:

\subsection{Quadratic Formula}
The quadratic formula is given by:
\[x = \frac{-b \pm \sqrt{b^2 - 4ac}}{2a}\]

\subsection{Matrix Example}
A sample matrix:
\[
\begin{pmatrix}
a_{11} & a_{12} & a_{13} \\
a_{21} & a_{22} & a_{23} \\
a_{31} & a_{32} & a_{33}
\end{pmatrix}
\]

\section{Tables}
Below is a sample table showing data:

\begin{table}[h]
\centering
\begin{tabular}{@{}lrrr@{}}
\toprule
\textbf{Category} & \textbf{Value 1} & \textbf{Value 2} & \textbf{Total} \\
\midrule
Item A & 10.5 & 20.3 & 30.8 \\
Item B & 15.2 & 25.1 & 40.3 \\
Item C & 22.1 & 12.7 & 34.8 \\
\bottomrule
\end{tabular}
\caption{Sample Data Table}
\label{tab:sample}
\end{table}

\section{Lists}
\subsection{Bullet Points}
\begin{itemize}
  \item First item
  \item Second item
  \item Third item with sub-items:
    \begin{itemize}
      \item Sub-item 1
      \item Sub-item 2
    \end{itemize}
\end{itemize}

\subsection{Numbered List}
\begin{enumerate}
  \item First step
  \item Second step
  \item Third step
\end{enumerate}

\section{Citations and References}
When you use the citation search feature, references will be automatically managed. For example, you might cite a paper like this: \cite{example2024}. All citations will be automatically added to your \texttt{cite.bib} file.

\section{Conclusion}
This template demonstrates basic \LaTeX{} features including mathematical formulas, tables, and lists. The bibliography system is now set up to work seamlessly with Bibby AI's citation management. Feel free to modify it for your needs!

\bibliographystyle{plain}
\bibliography{cite}

\end{document}\documentclass{article}
\usepackage[utf8]{inputenc}
\usepackage{booktabs}
\usepackage{amsmath}
\usepackage{natbib}

\title{My Document}
\author{Your Name}
\date{\today}

\begin{document}

\maketitle

\section{Introduction}
This is a sample document demonstrating various \LaTeX{} features including mathematical equations, tables, and formatting. You can cite papers using the citation search feature, and they will be automatically added to your \texttt{cite.bib} file.

\section{Mathematical Expressions}
Here are some example equations:

\subsection{Quadratic Formula}
The quadratic formula is given by:
\[x = \frac{-b \pm \sqrt{b^2 - 4ac}}{2a}\]

\subsection{Matrix Example}
A sample matrix:
\[
\begin{pmatrix}
a_{11} & a_{12} & a_{13} \\
a_{21} & a_{22} & a_{23} \\
a_{31} & a_{32} & a_{33}
\end{pmatrix}
\]

\section{Tables}
Below is a sample table showing data:

\begin{table}[h]
\centering
\begin{tabular}{@{}lrrr@{}}
\toprule
\textbf{Category} & \textbf{Value 1} & \textbf{Value 2} & \textbf{Total} \\
\midrule
Item A & 10.5 & 20.3 & 30.8 \\
Item B & 15.2 & 25.1 & 40.3 \\
Item C & 22.1 & 12.7 & 34.8 \\
\bottomrule
\end{tabular}
\caption{Sample Data Table}
\label{tab:sample}
\end{table}

\section{Lists}
\subsection{Bullet Points}
\begin{itemize}
  \item First item
  \item Second item
  \item Third item with sub-items:
    \begin{itemize}
      \item Sub-item 1
      \item Sub-item 2
    \end{itemize}
\end{itemize}

\subsection{Numbered List}
\begin{enumerate}
  \item First step
  \item Second step
  \item Third step
\end{enumerate}

\section{Citations and References}
When you use the citation search feature, references will be automatically managed. For example, you might cite a paper like this: \cite{example2024}. All citations will be automatically added to your \texttt{cite.bib} file.

\section{Conclusion}
This template demonstrates basic \LaTeX{} features including mathematical formulas, tables, and lists. The bibliography system is now set up to work seamlessly with Bibby AI's citation management. Feel free to modify it for your needs!

\bibliographystyle{plain}
\bibliography{cite}

\end{document}\documentclass{article}
\usepackage[utf8]{inputenc}
\usepackage{booktabs}
\usepackage{amsmath}
\usepackage{natbib}

\title{My Document}
\author{Your Name}
\date{\today}

\begin{document}

\maketitle

\section{Introduction}
This is a sample document demonstrating various \LaTeX{} features including mathematical equations, tables, and formatting. You can cite papers using the citation search feature, and they will be automatically added to your \texttt{cite.bib} file.

\section{Mathematical Expressions}
Here are some example equations:

\subsection{Quadratic Formula}
The quadratic formula is given by:
\[x = \frac{-b \pm \sqrt{b^2 - 4ac}}{2a}\]

\subsection{Matrix Example}
A sample matrix:
\[
\begin{pmatrix}
a_{11} & a_{12} & a_{13} \\
a_{21} & a_{22} & a_{23} \\
a_{31} & a_{32} & a_{33}
\end{pmatrix}
\]

\section{Tables}
Below is a sample table showing data:

\begin{table}[h]
\centering
\begin{tabular}{@{}lrrr@{}}
\toprule
\textbf{Category} & \textbf{Value 1} & \textbf{Value 2} & \textbf{Total} \\
\midrule
Item A & 10.5 & 20.3 & 30.8 \\
Item B & 15.2 & 25.1 & 40.3 \\
Item C & 22.1 & 12.7 & 34.8 \\
\bottomrule
\end{tabular}
\caption{Sample Data Table}
\label{tab:sample}
\end{table}

\section{Lists}
\subsection{Bullet Points}
\begin{itemize}
  \item First item
  \item Second item
  \item Third item with sub-items:
    \begin{itemize}
      \item Sub-item 1
      \item Sub-item 2
    \end{itemize}
\end{itemize}

\subsection{Numbered List}
\begin{enumerate}
  \item First step
  \item Second step
  \item Third step
\end{enumerate}

\section{Citations and References}
When you use the citation search feature, references will be automatically managed. For example, you might cite a paper like this: \cite{example2024}. All citations will be automatically added to your \texttt{cite.bib} file.

\section{Conclusion}
This template demonstrates basic \LaTeX{} features including mathematical formulas, tables, and lists. The bibliography system is now set up to work seamlessly with Bibby AI's citation management. Feel free to modify it for your needs!

\bibliographystyle{plain}
\bibliography{cite}

\end{document}\documentclass{article}
\usepackage[utf8]{inputenc}
\usepackage{booktabs}
\usepackage{amsmath}
\usepackage{natbib}

\title{My Document}
\author{Your Name}
\date{\today}

\begin{document}

\maketitle

\section{Introduction}
This is a sample document demonstrating various \LaTeX{} features including mathematical equations, tables, and formatting. You can cite papers using the citation search feature, and they will be automatically added to your \texttt{cite.bib} file.

\section{Mathematical Expressions}
Here are some example equations:

\subsection{Quadratic Formula}
The quadratic formula is given by:
\[x = \frac{-b \pm \sqrt{b^2 - 4ac}}{2a}\]

\subsection{Matrix Example}
A sample matrix:
\[
\begin{pmatrix}
a_{11} & a_{12} & a_{13} \\
a_{21} & a_{22} & a_{23} \\
a_{31} & a_{32} & a_{33}
\end{pmatrix}
\]

\section{Tables}
Below is a sample table showing data:

\begin{table}[h]
\centering
\begin{tabular}{@{}lrrr@{}}
\toprule
\textbf{Category} & \textbf{Value 1} & \textbf{Value 2} & \textbf{Total} \\
\midrule
Item A & 10.5 & 20.3 & 30.8 \\
Item B & 15.2 & 25.1 & 40.3 \\
Item C & 22.1 & 12.7 & 34.8 \\
\bottomrule
\end{tabular}
\caption{Sample Data Table}
\label{tab:sample}
\end{table}

\section{Lists}
\subsection{Bullet Points}
\begin{itemize}
  \item First item
  \item Second item
  \item Third item with sub-items:
    \begin{itemize}
      \item Sub-item 1
      \item Sub-item 2
    \end{itemize}
\end{itemize}

\subsection{Numbered List}
\begin{enumerate}
  \item First step
  \item Second step
  \item Third step
\end{enumerate}

\section{Citations and References}
When you use the citation search feature, references will be automatically managed. For example, you might cite a paper like this: \cite{example2024}. All citations will be automatically added to your \texttt{cite.bib} file.

\section{Conclusion}
This template demonstrates basic \LaTeX{} features including mathematical formulas, tables, and lists. The bibliography system is now set up to work seamlessly with Bibby AI's citation management. Feel free to modify it for your needs!

\bibliographystyle{plain}
\bibliography{cite}

\end{document}\documentclass{article}
\usepackage[utf8]{inputenc}
\usepackage{booktabs}
\usepackage{amsmath}
\usepackage{natbib}

\title{My Document}
\author{Your Name}
\date{\today}

\begin{document}

\maketitle

\section{Introduction}
This is a sample document demonstrating various \LaTeX{} features including mathematical equations, tables, and formatting. You can cite papers using the citation search feature, and they will be automatically added to your \texttt{cite.bib} file.

\section{Mathematical Expressions}
Here are some example equations:

\subsection{Quadratic Formula}
The quadratic formula is given by:
\[x = \frac{-b \pm \sqrt{b^2 - 4ac}}{2a}\]

\subsection{Matrix Example}
A sample matrix:
\[
\begin{pmatrix}
a_{11} & a_{12} & a_{13} \\
a_{21} & a_{22} & a_{23} \\
a_{31} & a_{32} & a_{33}
\end{pmatrix}
\]

\section{Tables}
Below is a sample table showing data:

\begin{table}[h]
\centering
\begin{tabular}{@{}lrrr@{}}
\toprule
\textbf{Category} & \textbf{Value 1} & \textbf{Value 2} & \textbf{Total} \\
\midrule
Item A & 10.5 & 20.3 & 30.8 \\
Item B & 15.2 & 25.1 & 40.3 \\
Item C & 22.1 & 12.7 & 34.8 \\
\bottomrule
\end{tabular}
\caption{Sample Data Table}
\label{tab:sample}
\end{table}

\section{Lists}
\subsection{Bullet Points}
\begin{itemize}
  \item First item
  \item Second item
  \item Third item with sub-items:
    \begin{itemize}
      \item Sub-item 1
      \item Sub-item 2
    \end{itemize}
\end{itemize}

\subsection{Numbered List}
\begin{enumerate}
  \item First step
  \item Second step
  \item Third step
\end{enumerate}

\section{Citations and References}
When you use the citation search feature, references will be automatically managed. For example, you might cite a paper like this: \cite{example2024}. All citations will be automatically added to your \texttt{cite.bib} file.

\section{Conclusion}
This template demonstrates basic \LaTeX{} features including mathematical formulas, tables, and lists. The bibliography system is now set up to work seamlessly with Bibby AI's citation management. Feel free to modify it for your needs!

\bibliographystyle{plain}
\bibliography{cite}

\end{document}\documentclass{article}
\usepackage[utf8]{inputenc}
\usepackage{booktabs}
\usepackage{amsmath}
\usepackage{natbib}

\title{My Document}
\author{Your Name}
\date{\today}

\begin{document}

\maketitle

\section{Introduction}
This is a sample document demonstrating various \LaTeX{} features including mathematical equations, tables, and formatting. You can cite papers using the citation search feature, and they will be automatically added to your \texttt{cite.bib} file.

\section{Mathematical Expressions}
Here are some example equations:

\subsection{Quadratic Formula}
The quadratic formula is given by:
\[x = \frac{-b \pm \sqrt{b^2 - 4ac}}{2a}\]

\subsection{Matrix Example}
A sample matrix:
\[
\begin{pmatrix}
a_{11} & a_{12} & a_{13} \\
a_{21} & a_{22} & a_{23} \\
a_{31} & a_{32} & a_{33}
\end{pmatrix}
\]

\section{Tables}
Below is a sample table showing data:

\begin{table}[h]
\centering
\begin{tabular}{@{}lrrr@{}}
\toprule
\textbf{Category} & \textbf{Value 1} & \textbf{Value 2} & \textbf{Total} \\
\midrule
Item A & 10.5 & 20.3 & 30.8 \\
Item B & 15.2 & 25.1 & 40.3 \\
Item C & 22.1 & 12.7 & 34.8 \\
\bottomrule
\end{tabular}
\caption{Sample Data Table}
\label{tab:sample}
\end{table}

\section{Lists}
\subsection{Bullet Points}
\begin{itemize}
  \item First item
  \item Second item
  \item Third item with sub-items:
    \begin{itemize}
      \item Sub-item 1
      \item Sub-item 2
    \end{itemize}
\end{itemize}

\subsection{Numbered List}
\begin{enumerate}
  \item First step
  \item Second step
  \item Third step
\end{enumerate}

\section{Citations and References}
When you use the citation search feature, references will be automatically managed. For example, you might cite a paper like this: \cite{example2024}. All citations will be automatically added to your \texttt{cite.bib} file.

\section{Conclusion}
This template demonstrates basic \LaTeX{} features including mathematical formulas, tables, and lists. The bibliography system is now set up to work seamlessly with Bibby AI's citation management. Feel free to modify it for your needs!

\bibliographystyle{plain}
\bibliography{cite}

\end{document}\documentclass{article}
\usepackage[utf8]{inputenc}
\usepackage{booktabs}
\usepackage{amsmath}
\usepackage{natbib}

\title{My Document}
\author{Your Name}
\date{\today}

\begin{document}

\maketitle

\section{Introduction}
This is a sample document demonstrating various \LaTeX{} features including mathematical equations, tables, and formatting. You can cite papers using the citation search feature, and they will be automatically added to your \texttt{cite.bib} file.

\section{Mathematical Expressions}
Here are some example equations:

\subsection{Quadratic Formula}
The quadratic formula is given by:
\[x = \frac{-b \pm \sqrt{b^2 - 4ac}}{2a}\]

\subsection{Matrix Example}
A sample matrix:
\[
\begin{pmatrix}
a_{11} & a_{12} & a_{13} \\
a_{21} & a_{22} & a_{23} \\
a_{31} & a_{32} & a_{33}
\end{pmatrix}
\]

\section{Tables}
Below is a sample table showing data:

\begin{table}[h]
\centering
\begin{tabular}{@{}lrrr@{}}
\toprule
\textbf{Category} & \textbf{Value 1} & \textbf{Value 2} & \textbf{Total} \\
\midrule
Item A & 10.5 & 20.3 & 30.8 \\
Item B & 15.2 & 25.1 & 40.3 \\
Item C & 22.1 & 12.7 & 34.8 \\
\bottomrule
\end{tabular}
\caption{Sample Data Table}
\label{tab:sample}
\end{table}

\section{Lists}
\subsection{Bullet Points}
\begin{itemize}
  \item First item
  \item Second item
  \item Third item with sub-items:
    \begin{itemize}
      \item Sub-item 1
      \item Sub-item 2
    \end{itemize}
\end{itemize}

\subsection{Numbered List}
\begin{enumerate}
  \item First step
  \item Second step
  \item Third step
\end{enumerate}

\section{Citations and References}
When you use the citation search feature, references will be automatically managed. For example, you might cite a paper like this: \cite{example2024}. All citations will be automatically added to your \texttt{cite.bib} file.

\section{Conclusion}
This template demonstrates basic \LaTeX{} features including mathematical formulas, tables, and lists. The bibliography system is now set up to work seamlessly with Bibby AI's citation management. Feel free to modify it for your needs!

\bibliographystyle{plain}
\bibliography{cite}

\end{document}\documentclass{article}
\usepackage[utf8]{inputenc}
\usepackage{booktabs}
\usepackage{amsmath}
\usepackage{natbib}

\title{My Document}
\author{Your Name}
\date{\today}

\begin{document}

\maketitle

\section{Introduction}
This is a sample document demonstrating various \LaTeX{} features including mathematical equations, tables, and formatting. You can cite papers using the citation search feature, and they will be automatically added to your \texttt{cite.bib} file.

\section{Mathematical Expressions}
Here are some example equations:

\subsection{Quadratic Formula}
The quadratic formula is given by:
\[x = \frac{-b \pm \sqrt{b^2 - 4ac}}{2a}\]

\subsection{Matrix Example}
A sample matrix:
\[
\begin{pmatrix}
a_{11} & a_{12} & a_{13} \\
a_{21} & a_{22} & a_{23} \\
a_{31} & a_{32} & a_{33}
\end{pmatrix}
\]

\section{Tables}
Below is a sample table showing data:

\begin{table}[h]
\centering
\begin{tabular}{@{}lrrr@{}}
\toprule
\textbf{Category} & \textbf{Value 1} & \textbf{Value 2} & \textbf{Total} \\
\midrule
Item A & 10.5 & 20.3 & 30.8 \\
Item B & 15.2 & 25.1 & 40.3 \\
Item C & 22.1 & 12.7 & 34.8 \\
\bottomrule
\end{tabular}
\caption{Sample Data Table}
\label{tab:sample}
\end{table}

\section{Lists}
\subsection{Bullet Points}
\begin{itemize}
  \item First item
  \item Second item
  \item Third item with sub-items:
    \begin{itemize}
      \item Sub-item 1
      \item Sub-item 2
    \end{itemize}
\end{itemize}

\subsection{Numbered List}
\begin{enumerate}
  \item First step
  \item Second step
  \item Third step
\end{enumerate}

\section{Citations and References}
When you use the citation search feature, references will be automatically managed. For example, you might cite a paper like this: \cite{example2024}. All citations will be automatically added to your \texttt{cite.bib} file.

\section{Conclusion}
This template demonstrates basic \LaTeX{} features including mathematical formulas, tables, and lists. The bibliography system is now set up to work seamlessly with Bibby AI's citation management. Feel free to modify it for your needs!

\bibliographystyle{plain}
\bibliography{cite}

\end{document}\documentclass{article}
\usepackage[utf8]{inputenc}
\usepackage{booktabs}
\usepackage{amsmath}
\usepackage{natbib}

\title{My Document}
\author{Your Name}
\date{\today}

\begin{document}

\maketitle

\section{Introduction}
This is a sample document demonstrating various \LaTeX{} features including mathematical equations, tables, and formatting. You can cite papers using the citation search feature, and they will be automatically added to your \texttt{cite.bib} file.

\section{Mathematical Expressions}
Here are some example equations:

\subsection{Quadratic Formula}
The quadratic formula is given by:
\[x = \frac{-b \pm \sqrt{b^2 - 4ac}}{2a}\]

\subsection{Matrix Example}
A sample matrix:
\[
\begin{pmatrix}
a_{11} & a_{12} & a_{13} \\
a_{21} & a_{22} & a_{23} \\
a_{31} & a_{32} & a_{33}
\end{pmatrix}
\]

\section{Tables}
Below is a sample table showing data:

\begin{table}[h]
\centering
\begin{tabular}{@{}lrrr@{}}
\toprule
\textbf{Category} & \textbf{Value 1} & \textbf{Value 2} & \textbf{Total} \\
\midrule
Item A & 10.5 & 20.3 & 30.8 \\
Item B & 15.2 & 25.1 & 40.3 \\
Item C & 22.1 & 12.7 & 34.8 \\
\bottomrule
\end{tabular}
\caption{Sample Data Table}
\label{tab:sample}
\end{table}

\section{Lists}
\subsection{Bullet Points}
\begin{itemize}
  \item First item
  \item Second item
  \item Third item with sub-items:
    \begin{itemize}
      \item Sub-item 1
      \item Sub-item 2
    \end{itemize}
\end{itemize}

\subsection{Numbered List}
\begin{enumerate}
  \item First step
  \item Second step
  \item Third step
\end{enumerate}

\section{Citations and References}
When you use the citation search feature, references will be automatically managed. For example, you might cite a paper like this: \cite{example2024}. All citations will be automatically added to your \texttt{cite.bib} file.

\section{Conclusion}
This template demonstrates basic \LaTeX{} features including mathematical formulas, tables, and lists. The bibliography system is now set up to work seamlessly with Bibby AI's citation management. Feel free to modify it for your needs!

\bibliographystyle{plain}
\bibliography{cite}

\end{document}\documentclass{article}
\usepackage[utf8]{inputenc}
\usepackage{booktabs}
\usepackage{amsmath}
\usepackage{natbib}

\title{My Document}
\author{Your Name}
\date{\today}

\begin{document}

\maketitle

\section{Introduction}
This is a sample document demonstrating various \LaTeX{} features including mathematical equations, tables, and formatting. You can cite papers using the citation search feature, and they will be automatically added to your \texttt{cite.bib} file.

\section{Mathematical Expressions}
Here are some example equations:

\subsection{Quadratic Formula}
The quadratic formula is given by:
\[x = \frac{-b \pm \sqrt{b^2 - 4ac}}{2a}\]

\subsection{Matrix Example}
A sample matrix:
\[
\begin{pmatrix}
a_{11} & a_{12} & a_{13} \\
a_{21} & a_{22} & a_{23} \\
a_{31} & a_{32} & a_{33}
\end{pmatrix}
\]

\section{Tables}
Below is a sample table showing data:

\begin{table}[h]
\centering
\begin{tabular}{@{}lrrr@{}}
\toprule
\textbf{Category} & \textbf{Value 1} & \textbf{Value 2} & \textbf{Total} \\
\midrule
Item A & 10.5 & 20.3 & 30.8 \\
Item B & 15.2 & 25.1 & 40.3 \\
Item C & 22.1 & 12.7 & 34.8 \\
\bottomrule
\end{tabular}
\caption{Sample Data Table}
\label{tab:sample}
\end{table}

\section{Lists}
\subsection{Bullet Points}
\begin{itemize}
  \item First item
  \item Second item
  \item Third item with sub-items:
    \begin{itemize}
      \item Sub-item 1
      \item Sub-item 2
    \end{itemize}
\end{itemize}

\subsection{Numbered List}
\begin{enumerate}
  \item First step
  \item Second step
  \item Third step
\end{enumerate}

\section{Citations and References}
When you use the citation search feature, references will be automatically managed. For example, you might cite a paper like this: \cite{example2024}. All citations will be automatically added to your \texttt{cite.bib} file.

\section{Conclusion}
This template demonstrates basic \LaTeX{} features including mathematical formulas, tables, and lists. The bibliography system is now set up to work seamlessly with Bibby AI's citation management. Feel free to modify it for your needs!

\bibliographystyle{plain}
\bibliography{cite}

\end{document}\documentclass{article}
\usepackage[utf8]{inputenc}
\usepackage{booktabs}
\usepackage{amsmath}
\usepackage{natbib}

\title{My Document}
\author{Your Name}
\date{\today}

\begin{document}

\maketitle

\section{Introduction}
This is a sample document demonstrating various \LaTeX{} features including mathematical equations, tables, and formatting. You can cite papers using the citation search feature, and they will be automatically added to your \texttt{cite.bib} file.

\section{Mathematical Expressions}
Here are some example equations:

\subsection{Quadratic Formula}
The quadratic formula is given by:
\[x = \frac{-b \pm \sqrt{b^2 - 4ac}}{2a}\]

\subsection{Matrix Example}
A sample matrix:
\[
\begin{pmatrix}
a_{11} & a_{12} & a_{13} \\
a_{21} & a_{22} & a_{23} \\
a_{31} & a_{32} & a_{33}
\end{pmatrix}
\]

\section{Tables}
Below is a sample table showing data:

\begin{table}[h]
\centering
\begin{tabular}{@{}lrrr@{}}
\toprule
\textbf{Category} & \textbf{Value 1} & \textbf{Value 2} & \textbf{Total} \\
\midrule
Item A & 10.5 & 20.3 & 30.8 \\
Item B & 15.2 & 25.1 & 40.3 \\
Item C & 22.1 & 12.7 & 34.8 \\
\bottomrule
\end{tabular}
\caption{Sample Data Table}
\label{tab:sample}
\end{table}

\section{Lists}
\subsection{Bullet Points}
\begin{itemize}
  \item First item
  \item Second item
  \item Third item with sub-items:
    \begin{itemize}
      \item Sub-item 1
      \item Sub-item 2
    \end{itemize}
\end{itemize}

\subsection{Numbered List}
\begin{enumerate}
  \item First step
  \item Second step
  \item Third step
\end{enumerate}

\section{Citations and References}
When you use the citation search feature, references will be automatically managed. For example, you might cite a paper like this: \cite{example2024}. All citations will be automatically added to your \texttt{cite.bib} file.

\section{Conclusion}
This template demonstrates basic \LaTeX{} features including mathematical formulas, tables, and lists. The bibliography system is now set up to work seamlessly with Bibby AI's citation management. Feel free to modify it for your needs!

\bibliographystyle{plain}
\bibliography{cite}

\end{document}\documentclass{article}
\usepackage[utf8]{inputenc}
\usepackage{booktabs}
\usepackage{amsmath}
\usepackage{natbib}

\title{My Document}
\author{Your Name}
\date{\today}

\begin{document}

\maketitle

\section{Introduction}
This is a sample document demonstrating various \LaTeX{} features including mathematical equations, tables, and formatting. You can cite papers using the citation search feature, and they will be automatically added to your \texttt{cite.bib} file.

\section{Mathematical Expressions}
Here are some example equations:

\subsection{Quadratic Formula}
The quadratic formula is given by:
\[x = \frac{-b \pm \sqrt{b^2 - 4ac}}{2a}\]

\subsection{Matrix Example}
A sample matrix:
\[
\begin{pmatrix}
a_{11} & a_{12} & a_{13} \\
a_{21} & a_{22} & a_{23} \\
a_{31} & a_{32} & a_{33}
\end{pmatrix}
\]

\section{Tables}
Below is a sample table showing data:

\begin{table}[h]
\centering
\begin{tabular}{@{}lrrr@{}}
\toprule
\textbf{Category} & \textbf{Value 1} & \textbf{Value 2} & \textbf{Total} \\
\midrule
Item A & 10.5 & 20.3 & 30.8 \\
Item B & 15.2 & 25.1 & 40.3 \\
Item C & 22.1 & 12.7 & 34.8 \\
\bottomrule
\end{tabular}
\caption{Sample Data Table}
\label{tab:sample}
\end{table}

\section{Lists}
\subsection{Bullet Points}
\begin{itemize}
  \item First item
  \item Second item
  \item Third item with sub-items:
    \begin{itemize}
      \item Sub-item 1
      \item Sub-item 2
    \end{itemize}
\end{itemize}

\subsection{Numbered List}
\begin{enumerate}
  \item First step
  \item Second step
  \item Third step
\end{enumerate}

\section{Citations and References}
When you use the citation search feature, references will be automatically managed. For example, you might cite a paper like this: \cite{example2024}. All citations will be automatically added to your \texttt{cite.bib} file.

\section{Conclusion}
This template demonstrates basic \LaTeX{} features including mathematical formulas, tables, and lists. The bibliography system is now set up to work seamlessly with Bibby AI's citation management. Feel free to modify it for your needs!

\bibliographystyle{plain}
\bibliography{cite}

\end{document}\documentclass{article}
\usepackage[utf8]{inputenc}
\usepackage{booktabs}
\usepackage{amsmath}
\usepackage{natbib}

\title{My Document}
\author{Your Name}
\date{\today}

\begin{document}

\maketitle

\section{Introduction}
This is a sample document demonstrating various \LaTeX{} features including mathematical equations, tables, and formatting. You can cite papers using the citation search feature, and they will be automatically added to your \texttt{cite.bib} file.

\section{Mathematical Expressions}
Here are some example equations:

\subsection{Quadratic Formula}
The quadratic formula is given by:
\[x = \frac{-b \pm \sqrt{b^2 - 4ac}}{2a}\]

\subsection{Matrix Example}
A sample matrix:
\[
\begin{pmatrix}
a_{11} & a_{12} & a_{13} \\
a_{21} & a_{22} & a_{23} \\
a_{31} & a_{32} & a_{33}
\end{pmatrix}
\]

\section{Tables}
Below is a sample table showing data:

\begin{table}[h]
\centering
\begin{tabular}{@{}lrrr@{}}
\toprule
\textbf{Category} & \textbf{Value 1} & \textbf{Value 2} & \textbf{Total} \\
\midrule
Item A & 10.5 & 20.3 & 30.8 \\
Item B & 15.2 & 25.1 & 40.3 \\
Item C & 22.1 & 12.7 & 34.8 \\
\bottomrule
\end{tabular}
\caption{Sample Data Table}
\label{tab:sample}
\end{table}

\section{Lists}
\subsection{Bullet Points}
\begin{itemize}
  \item First item
  \item Second item
  \item Third item with sub-items:
    \begin{itemize}
      \item Sub-item 1
      \item Sub-item 2
    \end{itemize}
\end{itemize}

\subsection{Numbered List}
\begin{enumerate}
  \item First step
  \item Second step
  \item Third step
\end{enumerate}

\section{Citations and References}
When you use the citation search feature, references will be automatically managed. For example, you might cite a paper like this: \cite{example2024}. All citations will be automatically added to your \texttt{cite.bib} file.

\section{Conclusion}
This template demonstrates basic \LaTeX{} features including mathematical formulas, tables, and lists. The bibliography system is now set up to work seamlessly with Bibby AI's citation management. Feel free to modify it for your needs!

\bibliographystyle{plain}
\bibliography{cite}

\end{document}\documentclass{article}
\usepackage[utf8]{inputenc}
\usepackage{booktabs}
\usepackage{amsmath}
\usepackage{natbib}

\title{My Document}
\author{Your Name}
\date{\today}

\begin{document}

\maketitle

\section{Introduction}
This is a sample document demonstrating various \LaTeX{} features including mathematical equations, tables, and formatting. You can cite papers using the citation search feature, and they will be automatically added to your \texttt{cite.bib} file.

\section{Mathematical Expressions}
Here are some example equations:

\subsection{Quadratic Formula}
The quadratic formula is given by:
\[x = \frac{-b \pm \sqrt{b^2 - 4ac}}{2a}\]

\subsection{Matrix Example}
A sample matrix:
\[
\begin{pmatrix}
a_{11} & a_{12} & a_{13} \\
a_{21} & a_{22} & a_{23} \\
a_{31} & a_{32} & a_{33}
\end{pmatrix}
\]

\section{Tables}
Below is a sample table showing data:

\begin{table}[h]
\centering
\begin{tabular}{@{}lrrr@{}}
\toprule
\textbf{Category} & \textbf{Value 1} & \textbf{Value 2} & \textbf{Total} \\
\midrule
Item A & 10.5 & 20.3 & 30.8 \\
Item B & 15.2 & 25.1 & 40.3 \\
Item C & 22.1 & 12.7 & 34.8 \\
\bottomrule
\end{tabular}
\caption{Sample Data Table}
\label{tab:sample}
\end{table}

\section{Lists}
\subsection{Bullet Points}
\begin{itemize}
  \item First item
  \item Second item
  \item Third item with sub-items:
    \begin{itemize}
      \item Sub-item 1
      \item Sub-item 2
    \end{itemize}
\end{itemize}

\subsection{Numbered List}
\begin{enumerate}
  \item First step
  \item Second step
  \item Third step
\end{enumerate}

\section{Citations and References}
When you use the citation search feature, references will be automatically managed. For example, you might cite a paper like this: \cite{example2024}. All citations will be automatically added to your \texttt{cite.bib} file.

\section{Conclusion}
This template demonstrates basic \LaTeX{} features including mathematical formulas, tables, and lists. The bibliography system is now set up to work seamlessly with Bibby AI's citation management. Feel free to modify it for your needs!

\bibliographystyle{plain}
\bibliography{cite}

\end{document}\documentclass{article}
\usepackage[utf8]{inputenc}
\usepackage{booktabs}
\usepackage{amsmath}
\usepackage{natbib}

\title{My Document}
\author{Your Name}
\date{\today}

\begin{document}

\maketitle

\section{Introduction}
This is a sample document demonstrating various \LaTeX{} features including mathematical equations, tables, and formatting. You can cite papers using the citation search feature, and they will be automatically added to your \texttt{cite.bib} file.

\section{Mathematical Expressions}
Here are some example equations:

\subsection{Quadratic Formula}
The quadratic formula is given by:
\[x = \frac{-b \pm \sqrt{b^2 - 4ac}}{2a}\]

\subsection{Matrix Example}
A sample matrix:
\[
\begin{pmatrix}
a_{11} & a_{12} & a_{13} \\
a_{21} & a_{22} & a_{23} \\
a_{31} & a_{32} & a_{33}
\end{pmatrix}
\]

\section{Tables}
Below is a sample table showing data:

\begin{table}[h]
\centering
\begin{tabular}{@{}lrrr@{}}
\toprule
\textbf{Category} & \textbf{Value 1} & \textbf{Value 2} & \textbf{Total} \\
\midrule
Item A & 10.5 & 20.3 & 30.8 \\
Item B & 15.2 & 25.1 & 40.3 \\
Item C & 22.1 & 12.7 & 34.8 \\
\bottomrule
\end{tabular}
\caption{Sample Data Table}
\label{tab:sample}
\end{table}

\section{Lists}
\subsection{Bullet Points}
\begin{itemize}
  \item First item
  \item Second item
  \item Third item with sub-items:
    \begin{itemize}
      \item Sub-item 1
      \item Sub-item 2
    \end{itemize}
\end{itemize}

\subsection{Numbered List}
\begin{enumerate}
  \item First step
  \item Second step
  \item Third step
\end{enumerate}

\section{Citations and References}
When you use the citation search feature, references will be automatically managed. For example, you might cite a paper like this: \cite{example2024}. All citations will be automatically added to your \texttt{cite.bib} file.

\section{Conclusion}
This template demonstrates basic \LaTeX{} features including mathematical formulas, tables, and lists. The bibliography system is now set up to work seamlessly with Bibby AI's citation management. Feel free to modify it for your needs!

\bibliographystyle{plain}
\bibliography{cite}

\end{document}\documentclass{article}
\usepackage[utf8]{inputenc}
\usepackage{booktabs}
\usepackage{amsmath}
\usepackage{natbib}

\title{My Document}
\author{Your Name}
\date{\today}

\begin{document}

\maketitle

\section{Introduction}
This is a sample document demonstrating various \LaTeX{} features including mathematical equations, tables, and formatting. You can cite papers using the citation search feature, and they will be automatically added to your \texttt{cite.bib} file.

\section{Mathematical Expressions}
Here are some example equations:

\subsection{Quadratic Formula}
The quadratic formula is given by:
\[x = \frac{-b \pm \sqrt{b^2 - 4ac}}{2a}\]

\subsection{Matrix Example}
A sample matrix:
\[
\begin{pmatrix}
a_{11} & a_{12} & a_{13} \\
a_{21} & a_{22} & a_{23} \\
a_{31} & a_{32} & a_{33}
\end{pmatrix}
\]

\section{Tables}
Below is a sample table showing data:

\begin{table}[h]
\centering
\begin{tabular}{@{}lrrr@{}}
\toprule
\textbf{Category} & \textbf{Value 1} & \textbf{Value 2} & \textbf{Total} \\
\midrule
Item A & 10.5 & 20.3 & 30.8 \\
Item B & 15.2 & 25.1 & 40.3 \\
Item C & 22.1 & 12.7 & 34.8 \\
\bottomrule
\end{tabular}
\caption{Sample Data Table}
\label{tab:sample}
\end{table}

\section{Lists}
\subsection{Bullet Points}
\begin{itemize}
  \item First item
  \item Second item
  \item Third item with sub-items:
    \begin{itemize}
      \item Sub-item 1
      \item Sub-item 2
    \end{itemize}
\end{itemize}

\subsection{Numbered List}
\begin{enumerate}
  \item First step
  \item Second step
  \item Third step
\end{enumerate}

\section{Citations and References}
When you use the citation search feature, references will be automatically managed. For example, you might cite a paper like this: \cite{example2024}. All citations will be automatically added to your \texttt{cite.bib} file.

\section{Conclusion}
This template demonstrates basic \LaTeX{} features including mathematical formulas, tables, and lists. The bibliography system is now set up to work seamlessly with Bibby AI's citation management. Feel free to modify it for your needs!

\bibliographystyle{plain}
\bibliography{cite}

\end{document}\documentclass{article}
\usepackage[utf8]{inputenc}
\usepackage{booktabs}
\usepackage{amsmath}
\usepackage{natbib}

\title{My Document}
\author{Your Name}
\date{\today}

\begin{document}

\maketitle

\section{Introduction}
This is a sample document demonstrating various \LaTeX{} features including mathematical equations, tables, and formatting. You can cite papers using the citation search feature, and they will be automatically added to your \texttt{cite.bib} file.

\section{Mathematical Expressions}
Here are some example equations:

\subsection{Quadratic Formula}
The quadratic formula is given by:
\[x = \frac{-b \pm \sqrt{b^2 - 4ac}}{2a}\]

\subsection{Matrix Example}
A sample matrix:
\[
\begin{pmatrix}
a_{11} & a_{12} & a_{13} \\
a_{21} & a_{22} & a_{23} \\
a_{31} & a_{32} & a_{33}
\end{pmatrix}
\]

\section{Tables}
Below is a sample table showing data:

\begin{table}[h]
\centering
\begin{tabular}{@{}lrrr@{}}
\toprule
\textbf{Category} & \textbf{Value 1} & \textbf{Value 2} & \textbf{Total} \\
\midrule
Item A & 10.5 & 20.3 & 30.8 \\
Item B & 15.2 & 25.1 & 40.3 \\
Item C & 22.1 & 12.7 & 34.8 \\
\bottomrule
\end{tabular}
\caption{Sample Data Table}
\label{tab:sample}
\end{table}

\section{Lists}
\subsection{Bullet Points}
\begin{itemize}
  \item First item
  \item Second item
  \item Third item with sub-items:
    \begin{itemize}
      \item Sub-item 1
      \item Sub-item 2
    \end{itemize}
\end{itemize}

\subsection{Numbered List}
\begin{enumerate}
  \item First step
  \item Second step
  \item Third step
\end{enumerate}

\section{Citations and References}
When you use the citation search feature, references will be automatically managed. For example, you might cite a paper like this: \cite{example2024}. All citations will be automatically added to your \texttt{cite.bib} file.

\section{Conclusion}
This template demonstrates basic \LaTeX{} features including mathematical formulas, tables, and lists. The bibliography system is now set up to work seamlessly with Bibby AI's citation management. Feel free to modify it for your needs!

\bibliographystyle{plain}
\bibliography{cite}

\end{document}\documentclass{article}
\usepackage[utf8]{inputenc}
\usepackage{booktabs}
\usepackage{amsmath}
\usepackage{natbib}

\title{My Document}
\author{Your Name}
\date{\today}

\begin{document}

\maketitle

\section{Introduction}
This is a sample document demonstrating various \LaTeX{} features including mathematical equations, tables, and formatting. You can cite papers using the citation search feature, and they will be automatically added to your \texttt{cite.bib} file.

\section{Mathematical Expressions}
Here are some example equations:

\subsection{Quadratic Formula}
The quadratic formula is given by:
\[x = \frac{-b \pm \sqrt{b^2 - 4ac}}{2a}\]

\subsection{Matrix Example}
A sample matrix:
\[
\begin{pmatrix}
a_{11} & a_{12} & a_{13} \\
a_{21} & a_{22} & a_{23} \\
a_{31} & a_{32} & a_{33}
\end{pmatrix}
\]

\section{Tables}
Below is a sample table showing data:

\begin{table}[h]
\centering
\begin{tabular}{@{}lrrr@{}}
\toprule
\textbf{Category} & \textbf{Value 1} & \textbf{Value 2} & \textbf{Total} \\
\midrule
Item A & 10.5 & 20.3 & 30.8 \\
Item B & 15.2 & 25.1 & 40.3 \\
Item C & 22.1 & 12.7 & 34.8 \\
\bottomrule
\end{tabular}
\caption{Sample Data Table}
\label{tab:sample}
\end{table}

\section{Lists}
\subsection{Bullet Points}
\begin{itemize}
  \item First item
  \item Second item
  \item Third item with sub-items:
    \begin{itemize}
      \item Sub-item 1
      \item Sub-item 2
    \end{itemize}
\end{itemize}

\subsection{Numbered List}
\begin{enumerate}
  \item First step
  \item Second step
  \item Third step
\end{enumerate}

\section{Citations and References}
When you use the citation search feature, references will be automatically managed. For example, you might cite a paper like this: \cite{example2024}. All citations will be automatically added to your \texttt{cite.bib} file.

\section{Conclusion}
This template demonstrates basic \LaTeX{} features including mathematical formulas, tables, and lists. The bibliography system is now set up to work seamlessly with Bibby AI's citation management. Feel free to modify it for your needs!

\bibliographystyle{plain}
\bibliography{cite}

\end{document}\documentclass{article}
\usepackage[utf8]{inputenc}
\usepackage{booktabs}
\usepackage{amsmath}
\usepackage{natbib}

\title{My Document}
\author{Your Name}
\date{\today}

\begin{document}

\maketitle

\section{Introduction}
This is a sample document demonstrating various \LaTeX{} features including mathematical equations, tables, and formatting. You can cite papers using the citation search feature, and they will be automatically added to your \texttt{cite.bib} file.

\section{Mathematical Expressions}
Here are some example equations:

\subsection{Quadratic Formula}
The quadratic formula is given by:
\[x = \frac{-b \pm \sqrt{b^2 - 4ac}}{2a}\]

\subsection{Matrix Example}
A sample matrix:
\[
\begin{pmatrix}
a_{11} & a_{12} & a_{13} \\
a_{21} & a_{22} & a_{23} \\
a_{31} & a_{32} & a_{33}
\end{pmatrix}
\]

\section{Tables}
Below is a sample table showing data:

\begin{table}[h]
\centering
\begin{tabular}{@{}lrrr@{}}
\toprule
\textbf{Category} & \textbf{Value 1} & \textbf{Value 2} & \textbf{Total} \\
\midrule
Item A & 10.5 & 20.3 & 30.8 \\
Item B & 15.2 & 25.1 & 40.3 \\
Item C & 22.1 & 12.7 & 34.8 \\
\bottomrule
\end{tabular}
\caption{Sample Data Table}
\label{tab:sample}
\end{table}

\section{Lists}
\subsection{Bullet Points}
\begin{itemize}
  \item First item
  \item Second item
  \item Third item with sub-items:
    \begin{itemize}
      \item Sub-item 1
      \item Sub-item 2
    \end{itemize}
\end{itemize}

\subsection{Numbered List}
\begin{enumerate}
  \item First step
  \item Second step
  \item Third step
\end{enumerate}

\section{Citations and References}
When you use the citation search feature, references will be automatically managed. For example, you might cite a paper like this: \cite{example2024}. All citations will be automatically added to your \texttt{cite.bib} file.

\section{Conclusion}
This template demonstrates basic \LaTeX{} features including mathematical formulas, tables, and lists. The bibliography system is now set up to work seamlessly with Bibby AI's citation management. Feel free to modify it for your needs!

\bibliographystyle{plain}
\bibliography{cite}

\end{document}\documentclass{article}
\usepackage[utf8]{inputenc}
\usepackage{booktabs}
\usepackage{amsmath}
\usepackage{natbib}

\title{My Document}
\author{Your Name}
\date{\today}

\begin{document}

\maketitle

\section{Introduction}
This is a sample document demonstrating various \LaTeX{} features including mathematical equations, tables, and formatting. You can cite papers using the citation search feature, and they will be automatically added to your \texttt{cite.bib} file.

\section{Mathematical Expressions}
Here are some example equations:

\subsection{Quadratic Formula}
The quadratic formula is given by:
\[x = \frac{-b \pm \sqrt{b^2 - 4ac}}{2a}\]

\subsection{Matrix Example}
A sample matrix:
\[
\begin{pmatrix}
a_{11} & a_{12} & a_{13} \\
a_{21} & a_{22} & a_{23} \\
a_{31} & a_{32} & a_{33}
\end{pmatrix}
\]

\section{Tables}
Below is a sample table showing data:

\begin{table}[h]
\centering
\begin{tabular}{@{}lrrr@{}}
\toprule
\textbf{Category} & \textbf{Value 1} & \textbf{Value 2} & \textbf{Total} \\
\midrule
Item A & 10.5 & 20.3 & 30.8 \\
Item B & 15.2 & 25.1 & 40.3 \\
Item C & 22.1 & 12.7 & 34.8 \\
\bottomrule
\end{tabular}
\caption{Sample Data Table}
\label{tab:sample}
\end{table}

\section{Lists}
\subsection{Bullet Points}
\begin{itemize}
  \item First item
  \item Second item
  \item Third item with sub-items:
    \begin{itemize}
      \item Sub-item 1
      \item Sub-item 2
    \end{itemize}
\end{itemize}

\subsection{Numbered List}
\begin{enumerate}
  \item First step
  \item Second step
  \item Third step
\end{enumerate}

\section{Citations and References}
When you use the citation search feature, references will be automatically managed. For example, you might cite a paper like this: \cite{example2024}. All citations will be automatically added to your \texttt{cite.bib} file.

\section{Conclusion}
This template demonstrates basic \LaTeX{} features including mathematical formulas, tables, and lists. The bibliography system is now set up to work seamlessly with Bibby AI's citation management. Feel free to modify it for your needs!

\bibliographystyle{plain}
\bibliography{cite}

\end{document}\documentclass{article}
\usepackage[utf8]{inputenc}
\usepackage{booktabs}
\usepackage{amsmath}
\usepackage{natbib}

\title{My Document}
\author{Your Name}
\date{\today}

\begin{document}

\maketitle

\section{Introduction}
This is a sample document demonstrating various \LaTeX{} features including mathematical equations, tables, and formatting. You can cite papers using the citation search feature, and they will be automatically added to your \texttt{cite.bib} file.

\section{Mathematical Expressions}
Here are some example equations:

\subsection{Quadratic Formula}
The quadratic formula is given by:
\[x = \frac{-b \pm \sqrt{b^2 - 4ac}}{2a}\]

\subsection{Matrix Example}
A sample matrix:
\[
\begin{pmatrix}
a_{11} & a_{12} & a_{13} \\
a_{21} & a_{22} & a_{23} \\
a_{31} & a_{32} & a_{33}
\end{pmatrix}
\]

\section{Tables}
Below is a sample table showing data:

\begin{table}[h]
\centering
\begin{tabular}{@{}lrrr@{}}
\toprule
\textbf{Category} & \textbf{Value 1} & \textbf{Value 2} & \textbf{Total} \\
\midrule
Item A & 10.5 & 20.3 & 30.8 \\
Item B & 15.2 & 25.1 & 40.3 \\
Item C & 22.1 & 12.7 & 34.8 \\
\bottomrule
\end{tabular}
\caption{Sample Data Table}
\label{tab:sample}
\end{table}

\section{Lists}
\subsection{Bullet Points}
\begin{itemize}
  \item First item
  \item Second item
  \item Third item with sub-items:
    \begin{itemize}
      \item Sub-item 1
      \item Sub-item 2
    \end{itemize}
\end{itemize}

\subsection{Numbered List}
\begin{enumerate}
  \item First step
  \item Second step
  \item Third step
\end{enumerate}

\section{Citations and References}
When you use the citation search feature, references will be automatically managed. For example, you might cite a paper like this: \cite{example2024}. All citations will be automatically added to your \texttt{cite.bib} file.

\section{Conclusion}
This template demonstrates basic \LaTeX{} features including mathematical formulas, tables, and lists. The bibliography system is now set up to work seamlessly with Bibby AI's citation management. Feel free to modify it for your needs!

\bibliographystyle{plain}
\bibliography{cite}

\end{document}\documentclass{article}
\usepackage[utf8]{inputenc}
\usepackage{booktabs}
\usepackage{amsmath}
\usepackage{natbib}

\title{My Document}
\author{Your Name}
\date{\today}

\begin{document}

\maketitle

\section{Introduction}
This is a sample document demonstrating various \LaTeX{} features including mathematical equations, tables, and formatting. You can cite papers using the citation search feature, and they will be automatically added to your \texttt{cite.bib} file.

\section{Mathematical Expressions}
Here are some example equations:

\subsection{Quadratic Formula}
The quadratic formula is given by:
\[x = \frac{-b \pm \sqrt{b^2 - 4ac}}{2a}\]

\subsection{Matrix Example}
A sample matrix:
\[
\begin{pmatrix}
a_{11} & a_{12} & a_{13} \\
a_{21} & a_{22} & a_{23} \\
a_{31} & a_{32} & a_{33}
\end{pmatrix}
\]

\section{Tables}
Below is a sample table showing data:

\begin{table}[h]
\centering
\begin{tabular}{@{}lrrr@{}}
\toprule
\textbf{Category} & \textbf{Value 1} & \textbf{Value 2} & \textbf{Total} \\
\midrule
Item A & 10.5 & 20.3 & 30.8 \\
Item B & 15.2 & 25.1 & 40.3 \\
Item C & 22.1 & 12.7 & 34.8 \\
\bottomrule
\end{tabular}
\caption{Sample Data Table}
\label{tab:sample}
\end{table}

\section{Lists}
\subsection{Bullet Points}
\begin{itemize}
  \item First item
  \item Second item
  \item Third item with sub-items:
    \begin{itemize}
      \item Sub-item 1
      \item Sub-item 2
    \end{itemize}
\end{itemize}

\subsection{Numbered List}
\begin{enumerate}
  \item First step
  \item Second step
  \item Third step
\end{enumerate}

\section{Citations and References}
When you use the citation search feature, references will be automatically managed. For example, you might cite a paper like this: \cite{example2024}. All citations will be automatically added to your \texttt{cite.bib} file.

\section{Conclusion}
This template demonstrates basic \LaTeX{} features including mathematical formulas, tables, and lists. The bibliography system is now set up to work seamlessly with Bibby AI's citation management. Feel free to modify it for your needs!

\bibliographystyle{plain}
\bibliography{cite}

\end{document}\documentclass{article}
\usepackage[utf8]{inputenc}
\usepackage{booktabs}
\usepackage{amsmath}
\usepackage{natbib}

\title{My Document}
\author{Your Name}
\date{\today}

\begin{document}

\maketitle

\section{Introduction}
This is a sample document demonstrating various \LaTeX{} features including mathematical equations, tables, and formatting. You can cite papers using the citation search feature, and they will be automatically added to your \texttt{cite.bib} file.

\section{Mathematical Expressions}
Here are some example equations:

\subsection{Quadratic Formula}
The quadratic formula is given by:
\[x = \frac{-b \pm \sqrt{b^2 - 4ac}}{2a}\]

\subsection{Matrix Example}
A sample matrix:
\[
\begin{pmatrix}
a_{11} & a_{12} & a_{13} \\
a_{21} & a_{22} & a_{23} \\
a_{31} & a_{32} & a_{33}
\end{pmatrix}
\]

\section{Tables}
Below is a sample table showing data:

\begin{table}[h]
\centering
\begin{tabular}{@{}lrrr@{}}
\toprule
\textbf{Category} & \textbf{Value 1} & \textbf{Value 2} & \textbf{Total} \\
\midrule
Item A & 10.5 & 20.3 & 30.8 \\
Item B & 15.2 & 25.1 & 40.3 \\
Item C & 22.1 & 12.7 & 34.8 \\
\bottomrule
\end{tabular}
\caption{Sample Data Table}
\label{tab:sample}
\end{table}

\section{Lists}
\subsection{Bullet Points}
\begin{itemize}
  \item First item
  \item Second item
  \item Third item with sub-items:
    \begin{itemize}
      \item Sub-item 1
      \item Sub-item 2
    \end{itemize}
\end{itemize}

\subsection{Numbered List}
\begin{enumerate}
  \item First step
  \item Second step
  \item Third step
\end{enumerate}

\section{Citations and References}
When you use the citation search feature, references will be automatically managed. For example, you might cite a paper like this: \cite{example2024}. All citations will be automatically added to your \texttt{cite.bib} file.

\section{Conclusion}
This template demonstrates basic \LaTeX{} features including mathematical formulas, tables, and lists. The bibliography system is now set up to work seamlessly with Bibby AI's citation management. Feel free to modify it for your needs!

\bibliographystyle{plain}
\bibliography{cite}

\end{document}\documentclass{article}
\usepackage[utf8]{inputenc}
\usepackage{booktabs}
\usepackage{amsmath}
\usepackage{natbib}

\title{My Document}
\author{Your Name}
\date{\today}

\begin{document}

\maketitle

\section{Introduction}
This is a sample document demonstrating various \LaTeX{} features including mathematical equations, tables, and formatting. You can cite papers using the citation search feature, and they will be automatically added to your \texttt{cite.bib} file.

\section{Mathematical Expressions}
Here are some example equations:

\subsection{Quadratic Formula}
The quadratic formula is given by:
\[x = \frac{-b \pm \sqrt{b^2 - 4ac}}{2a}\]

\subsection{Matrix Example}
A sample matrix:
\[
\begin{pmatrix}
a_{11} & a_{12} & a_{13} \\
a_{21} & a_{22} & a_{23} \\
a_{31} & a_{32} & a_{33}
\end{pmatrix}
\]

\section{Tables}
Below is a sample table showing data:

\begin{table}[h]
\centering
\begin{tabular}{@{}lrrr@{}}
\toprule
\textbf{Category} & \textbf{Value 1} & \textbf{Value 2} & \textbf{Total} \\
\midrule
Item A & 10.5 & 20.3 & 30.8 \\
Item B & 15.2 & 25.1 & 40.3 \\
Item C & 22.1 & 12.7 & 34.8 \\
\bottomrule
\end{tabular}
\caption{Sample Data Table}
\label{tab:sample}
\end{table}

\section{Lists}
\subsection{Bullet Points}
\begin{itemize}
  \item First item
  \item Second item
  \item Third item with sub-items:
    \begin{itemize}
      \item Sub-item 1
      \item Sub-item 2
    \end{itemize}
\end{itemize}

\subsection{Numbered List}
\begin{enumerate}
  \item First step
  \item Second step
  \item Third step
\end{enumerate}

\section{Citations and References}
When you use the citation search feature, references will be automatically managed. For example, you might cite a paper like this: \cite{example2024}. All citations will be automatically added to your \texttt{cite.bib} file.

\section{Conclusion}
This template demonstrates basic \LaTeX{} features including mathematical formulas, tables, and lists. The bibliography system is now set up to work seamlessly with Bibby AI's citation management. Feel free to modify it for your needs!

\bibliographystyle{plain}
\bibliography{cite}

\end{document}\documentclass{article}
\usepackage[utf8]{inputenc}
\usepackage{booktabs}
\usepackage{amsmath}
\usepackage{natbib}

\title{My Document}
\author{Your Name}
\date{\today}

\begin{document}

\maketitle

\section{Introduction}
This is a sample document demonstrating various \LaTeX{} features including mathematical equations, tables, and formatting. You can cite papers using the citation search feature, and they will be automatically added to your \texttt{cite.bib} file.

\section{Mathematical Expressions}
Here are some example equations:

\subsection{Quadratic Formula}
The quadratic formula is given by:
\[x = \frac{-b \pm \sqrt{b^2 - 4ac}}{2a}\]

\subsection{Matrix Example}
A sample matrix:
\[
\begin{pmatrix}
a_{11} & a_{12} & a_{13} \\
a_{21} & a_{22} & a_{23} \\
a_{31} & a_{32} & a_{33}
\end{pmatrix}
\]

\section{Tables}
Below is a sample table showing data:

\begin{table}[h]
\centering
\begin{tabular}{@{}lrrr@{}}
\toprule
\textbf{Category} & \textbf{Value 1} & \textbf{Value 2} & \textbf{Total} \\
\midrule
Item A & 10.5 & 20.3 & 30.8 \\
Item B & 15.2 & 25.1 & 40.3 \\
Item C & 22.1 & 12.7 & 34.8 \\
\bottomrule
\end{tabular}
\caption{Sample Data Table}
\label{tab:sample}
\end{table}

\section{Lists}
\subsection{Bullet Points}
\begin{itemize}
  \item First item
  \item Second item
  \item Third item with sub-items:
    \begin{itemize}
      \item Sub-item 1
      \item Sub-item 2
    \end{itemize}
\end{itemize}

\subsection{Numbered List}
\begin{enumerate}
  \item First step
  \item Second step
  \item Third step
\end{enumerate}

\section{Citations and References}
When you use the citation search feature, references will be automatically managed. For example, you might cite a paper like this: \cite{example2024}. All citations will be automatically added to your \texttt{cite.bib} file.

\section{Conclusion}
This template demonstrates basic \LaTeX{} features including mathematical formulas, tables, and lists. The bibliography system is now set up to work seamlessly with Bibby AI's citation management. Feel free to modify it for your needs!

\bibliographystyle{plain}
\bibliography{cite}

\end{document}\documentclass{article}
\usepackage[utf8]{inputenc}
\usepackage{booktabs}
\usepackage{amsmath}
\usepackage{natbib}

\title{My Document}
\author{Your Name}
\date{\today}

\begin{document}

\maketitle

\section{Introduction}
This is a sample document demonstrating various \LaTeX{} features including mathematical equations, tables, and formatting. You can cite papers using the citation search feature, and they will be automatically added to your \texttt{cite.bib} file.

\section{Mathematical Expressions}
Here are some example equations:

\subsection{Quadratic Formula}
The quadratic formula is given by:
\[x = \frac{-b \pm \sqrt{b^2 - 4ac}}{2a}\]

\subsection{Matrix Example}
A sample matrix:
\[
\begin{pmatrix}
a_{11} & a_{12} & a_{13} \\
a_{21} & a_{22} & a_{23} \\
a_{31} & a_{32} & a_{33}
\end{pmatrix}
\]

\section{Tables}
Below is a sample table showing data:

\begin{table}[h]
\centering
\begin{tabular}{@{}lrrr@{}}
\toprule
\textbf{Category} & \textbf{Value 1} & \textbf{Value 2} & \textbf{Total} \\
\midrule
Item A & 10.5 & 20.3 & 30.8 \\
Item B & 15.2 & 25.1 & 40.3 \\
Item C & 22.1 & 12.7 & 34.8 \\
\bottomrule
\end{tabular}
\caption{Sample Data Table}
\label{tab:sample}
\end{table}

\section{Lists}
\subsection{Bullet Points}
\begin{itemize}
  \item First item
  \item Second item
  \item Third item with sub-items:
    \begin{itemize}
      \item Sub-item 1
      \item Sub-item 2
    \end{itemize}
\end{itemize}

\subsection{Numbered List}
\begin{enumerate}
  \item First step
  \item Second step
  \item Third step
\end{enumerate}

\section{Citations and References}
When you use the citation search feature, references will be automatically managed. For example, you might cite a paper like this: \cite{example2024}. All citations will be automatically added to your \texttt{cite.bib} file.

\section{Conclusion}
This template demonstrates basic \LaTeX{} features including mathematical formulas, tables, and lists. The bibliography system is now set up to work seamlessly with Bibby AI's citation management. Feel free to modify it for your needs!

\bibliographystyle{plain}
\bibliography{cite}

\end{document}\documentclass{article}
\usepackage[utf8]{inputenc}
\usepackage{booktabs}
\usepackage{amsmath}
\usepackage{natbib}

\title{My Document}
\author{Your Name}
\date{\today}

\begin{document}

\maketitle

\section{Introduction}
This is a sample document demonstrating various \LaTeX{} features including mathematical equations, tables, and formatting. You can cite papers using the citation search feature, and they will be automatically added to your \texttt{cite.bib} file.

\section{Mathematical Expressions}
Here are some example equations:

\subsection{Quadratic Formula}
The quadratic formula is given by:
\[x = \frac{-b \pm \sqrt{b^2 - 4ac}}{2a}\]

\subsection{Matrix Example}
A sample matrix:
\[
\begin{pmatrix}
a_{11} & a_{12} & a_{13} \\
a_{21} & a_{22} & a_{23} \\
a_{31} & a_{32} & a_{33}
\end{pmatrix}
\]

\section{Tables}
Below is a sample table showing data:

\begin{table}[h]
\centering
\begin{tabular}{@{}lrrr@{}}
\toprule
\textbf{Category} & \textbf{Value 1} & \textbf{Value 2} & \textbf{Total} \\
\midrule
Item A & 10.5 & 20.3 & 30.8 \\
Item B & 15.2 & 25.1 & 40.3 \\
Item C & 22.1 & 12.7 & 34.8 \\
\bottomrule
\end{tabular}
\caption{Sample Data Table}
\label{tab:sample}
\end{table}

\section{Lists}
\subsection{Bullet Points}
\begin{itemize}
  \item First item
  \item Second item
  \item Third item with sub-items:
    \begin{itemize}
      \item Sub-item 1
      \item Sub-item 2
    \end{itemize}
\end{itemize}

\subsection{Numbered List}
\begin{enumerate}
  \item First step
  \item Second step
  \item Third step
\end{enumerate}

\section{Citations and References}
When you use the citation search feature, references will be automatically managed. For example, you might cite a paper like this: \cite{example2024}. All citations will be automatically added to your \texttt{cite.bib} file.

\section{Conclusion}
This template demonstrates basic \LaTeX{} features including mathematical formulas, tables, and lists. The bibliography system is now set up to work seamlessly with Bibby AI's citation management. Feel free to modify it for your needs!

\bibliographystyle{plain}
\bibliography{cite}

\end{document}\documentclass{article}
\usepackage[utf8]{inputenc}
\usepackage{booktabs}
\usepackage{amsmath}
\usepackage{natbib}

\title{My Document}
\author{Your Name}
\date{\today}

\begin{document}

\maketitle

\section{Introduction}
This is a sample document demonstrating various \LaTeX{} features including mathematical equations, tables, and formatting. You can cite papers using the citation search feature, and they will be automatically added to your \texttt{cite.bib} file.

\section{Mathematical Expressions}
Here are some example equations:

\subsection{Quadratic Formula}
The quadratic formula is given by:
\[x = \frac{-b \pm \sqrt{b^2 - 4ac}}{2a}\]

\subsection{Matrix Example}
A sample matrix:
\[
\begin{pmatrix}
a_{11} & a_{12} & a_{13} \\
a_{21} & a_{22} & a_{23} \\
a_{31} & a_{32} & a_{33}
\end{pmatrix}
\]

\section{Tables}
Below is a sample table showing data:

\begin{table}[h]
\centering
\begin{tabular}{@{}lrrr@{}}
\toprule
\textbf{Category} & \textbf{Value 1} & \textbf{Value 2} & \textbf{Total} \\
\midrule
Item A & 10.5 & 20.3 & 30.8 \\
Item B & 15.2 & 25.1 & 40.3 \\
Item C & 22.1 & 12.7 & 34.8 \\
\bottomrule
\end{tabular}
\caption{Sample Data Table}
\label{tab:sample}
\end{table}

\section{Lists}
\subsection{Bullet Points}
\begin{itemize}
  \item First item
  \item Second item
  \item Third item with sub-items:
    \begin{itemize}
      \item Sub-item 1
      \item Sub-item 2
    \end{itemize}
\end{itemize}

\subsection{Numbered List}
\begin{enumerate}
  \item First step
  \item Second step
  \item Third step
\end{enumerate}

\section{Citations and References}
When you use the citation search feature, references will be automatically managed. For example, you might cite a paper like this: \cite{example2024}. All citations will be automatically added to your \texttt{cite.bib} file.

\section{Conclusion}
This template demonstrates basic \LaTeX{} features including mathematical formulas, tables, and lists. The bibliography system is now set up to work seamlessly with Bibby AI's citation management. Feel free to modify it for your needs!

\bibliographystyle{plain}
\bibliography{cite}

\end{document}\documentclass{article}
\usepackage[utf8]{inputenc}
\usepackage{booktabs}
\usepackage{amsmath}
\usepackage{natbib}

\title{My Document}
\author{Your Name}
\date{\today}

\begin{document}

\maketitle

\section{Introduction}
This is a sample document demonstrating various \LaTeX{} features including mathematical equations, tables, and formatting. You can cite papers using the citation search feature, and they will be automatically added to your \texttt{cite.bib} file.

\section{Mathematical Expressions}
Here are some example equations:

\subsection{Quadratic Formula}
The quadratic formula is given by:
\[x = \frac{-b \pm \sqrt{b^2 - 4ac}}{2a}\]

\subsection{Matrix Example}
A sample matrix:
\[
\begin{pmatrix}
a_{11} & a_{12} & a_{13} \\
a_{21} & a_{22} & a_{23} \\
a_{31} & a_{32} & a_{33}
\end{pmatrix}
\]

\section{Tables}
Below is a sample table showing data:

\begin{table}[h]
\centering
\begin{tabular}{@{}lrrr@{}}
\toprule
\textbf{Category} & \textbf{Value 1} & \textbf{Value 2} & \textbf{Total} \\
\midrule
Item A & 10.5 & 20.3 & 30.8 \\
Item B & 15.2 & 25.1 & 40.3 \\
Item C & 22.1 & 12.7 & 34.8 \\
\bottomrule
\end{tabular}
\caption{Sample Data Table}
\label{tab:sample}
\end{table}

\section{Lists}
\subsection{Bullet Points}
\begin{itemize}
  \item First item
  \item Second item
  \item Third item with sub-items:
    \begin{itemize}
      \item Sub-item 1
      \item Sub-item 2
    \end{itemize}
\end{itemize}

\subsection{Numbered List}
\begin{enumerate}
  \item First step
  \item Second step
  \item Third step
\end{enumerate}

\section{Citations and References}
When you use the citation search feature, references will be automatically managed. For example, you might cite a paper like this: \cite{example2024}. All citations will be automatically added to your \texttt{cite.bib} file.

\section{Conclusion}
This template demonstrates basic \LaTeX{} features including mathematical formulas, tables, and lists. The bibliography system is now set up to work seamlessly with Bibby AI's citation management. Feel free to modify it for your needs!

\bibliographystyle{plain}
\bibliography{cite}

\end{document}\documentclass{article}
\usepackage[utf8]{inputenc}
\usepackage{booktabs}
\usepackage{amsmath}
\usepackage{natbib}

\title{My Document}
\author{Your Name}
\date{\today}

\begin{document}

\maketitle

\section{Introduction}
This is a sample document demonstrating various \LaTeX{} features including mathematical equations, tables, and formatting. You can cite papers using the citation search feature, and they will be automatically added to your \texttt{cite.bib} file.

\section{Mathematical Expressions}
Here are some example equations:

\subsection{Quadratic Formula}
The quadratic formula is given by:
\[x = \frac{-b \pm \sqrt{b^2 - 4ac}}{2a}\]

\subsection{Matrix Example}
A sample matrix:
\[
\begin{pmatrix}
a_{11} & a_{12} & a_{13} \\
a_{21} & a_{22} & a_{23} \\
a_{31} & a_{32} & a_{33}
\end{pmatrix}
\]

\section{Tables}
Below is a sample table showing data:

\begin{table}[h]
\centering
\begin{tabular}{@{}lrrr@{}}
\toprule
\textbf{Category} & \textbf{Value 1} & \textbf{Value 2} & \textbf{Total} \\
\midrule
Item A & 10.5 & 20.3 & 30.8 \\
Item B & 15.2 & 25.1 & 40.3 \\
Item C & 22.1 & 12.7 & 34.8 \\
\bottomrule
\end{tabular}
\caption{Sample Data Table}
\label{tab:sample}
\end{table}

\section{Lists}
\subsection{Bullet Points}
\begin{itemize}
  \item First item
  \item Second item
  \item Third item with sub-items:
    \begin{itemize}
      \item Sub-item 1
      \item Sub-item 2
    \end{itemize}
\end{itemize}

\subsection{Numbered List}
\begin{enumerate}
  \item First step
  \item Second step
  \item Third step
\end{enumerate}

\section{Citations and References}
When you use the citation search feature, references will be automatically managed. For example, you might cite a paper like this: \cite{example2024}. All citations will be automatically added to your \texttt{cite.bib} file.

\section{Conclusion}
This template demonstrates basic \LaTeX{} features including mathematical formulas, tables, and lists. The bibliography system is now set up to work seamlessly with Bibby AI's citation management. Feel free to modify it for your needs!

\bibliographystyle{plain}
\bibliography{cite}

\end{document}\documentclass{article}
\usepackage[utf8]{inputenc}
\usepackage{booktabs}
\usepackage{amsmath}
\usepackage{natbib}

\title{My Document}
\author{Your Name}
\date{\today}

\begin{document}

\maketitle

\section{Introduction}
This is a sample document demonstrating various \LaTeX{} features including mathematical equations, tables, and formatting. You can cite papers using the citation search feature, and they will be automatically added to your \texttt{cite.bib} file.

\section{Mathematical Expressions}
Here are some example equations:

\subsection{Quadratic Formula}
The quadratic formula is given by:
\[x = \frac{-b \pm \sqrt{b^2 - 4ac}}{2a}\]

\subsection{Matrix Example}
A sample matrix:
\[
\begin{pmatrix}
a_{11} & a_{12} & a_{13} \\
a_{21} & a_{22} & a_{23} \\
a_{31} & a_{32} & a_{33}
\end{pmatrix}
\]

\section{Tables}
Below is a sample table showing data:

\begin{table}[h]
\centering
\begin{tabular}{@{}lrrr@{}}
\toprule
\textbf{Category} & \textbf{Value 1} & \textbf{Value 2} & \textbf{Total} \\
\midrule
Item A & 10.5 & 20.3 & 30.8 \\
Item B & 15.2 & 25.1 & 40.3 \\
Item C & 22.1 & 12.7 & 34.8 \\
\bottomrule
\end{tabular}
\caption{Sample Data Table}
\label{tab:sample}
\end{table}

\section{Lists}
\subsection{Bullet Points}
\begin{itemize}
  \item First item
  \item Second item
  \item Third item with sub-items:
    \begin{itemize}
      \item Sub-item 1
      \item Sub-item 2
    \end{itemize}
\end{itemize}

\subsection{Numbered List}
\begin{enumerate}
  \item First step
  \item Second step
  \item Third step
\end{enumerate}

\section{Citations and References}
When you use the citation search feature, references will be automatically managed. For example, you might cite a paper like this: \cite{example2024}. All citations will be automatically added to your \texttt{cite.bib} file.

\section{Conclusion}
This template demonstrates basic \LaTeX{} features including mathematical formulas, tables, and lists. The bibliography system is now set up to work seamlessly with Bibby AI's citation management. Feel free to modify it for your needs!

\bibliographystyle{plain}
\bibliography{cite}

\end{document}\documentclass{article}
\usepackage[utf8]{inputenc}
\usepackage{booktabs}
\usepackage{amsmath}
\usepackage{natbib}

\title{My Document}
\author{Your Name}
\date{\today}

\begin{document}

\maketitle

\section{Introduction}
This is a sample document demonstrating various \LaTeX{} features including mathematical equations, tables, and formatting. You can cite papers using the citation search feature, and they will be automatically added to your \texttt{cite.bib} file.

\section{Mathematical Expressions}
Here are some example equations:

\subsection{Quadratic Formula}
The quadratic formula is given by:
\[x = \frac{-b \pm \sqrt{b^2 - 4ac}}{2a}\]

\subsection{Matrix Example}
A sample matrix:
\[
\begin{pmatrix}
a_{11} & a_{12} & a_{13} \\
a_{21} & a_{22} & a_{23} \\
a_{31} & a_{32} & a_{33}
\end{pmatrix}
\]

\section{Tables}
Below is a sample table showing data:

\begin{table}[h]
\centering
\begin{tabular}{@{}lrrr@{}}
\toprule
\textbf{Category} & \textbf{Value 1} & \textbf{Value 2} & \textbf{Total} \\
\midrule
Item A & 10.5 & 20.3 & 30.8 \\
Item B & 15.2 & 25.1 & 40.3 \\
Item C & 22.1 & 12.7 & 34.8 \\
\bottomrule
\end{tabular}
\caption{Sample Data Table}
\label{tab:sample}
\end{table}

\section{Lists}
\subsection{Bullet Points}
\begin{itemize}
  \item First item
  \item Second item
  \item Third item with sub-items:
    \begin{itemize}
      \item Sub-item 1
      \item Sub-item 2
    \end{itemize}
\end{itemize}

\subsection{Numbered List}
\begin{enumerate}
  \item First step
  \item Second step
  \item Third step
\end{enumerate}

\section{Citations and References}
When you use the citation search feature, references will be automatically managed. For example, you might cite a paper like this: \cite{example2024}. All citations will be automatically added to your \texttt{cite.bib} file.

\section{Conclusion}
This template demonstrates basic \LaTeX{} features including mathematical formulas, tables, and lists. The bibliography system is now set up to work seamlessly with Bibby AI's citation management. Feel free to modify it for your needs!

\bibliographystyle{plain}
\bibliography{cite}

\end{document}\documentclass{article}
\usepackage[utf8]{inputenc}
\usepackage{booktabs}
\usepackage{amsmath}
\usepackage{natbib}

\title{My Document}
\author{Your Name}
\date{\today}

\begin{document}

\maketitle

\section{Introduction}
This is a sample document demonstrating various \LaTeX{} features including mathematical equations, tables, and formatting. You can cite papers using the citation search feature, and they will be automatically added to your \texttt{cite.bib} file.

\section{Mathematical Expressions}
Here are some example equations:

\subsection{Quadratic Formula}
The quadratic formula is given by:
\[x = \frac{-b \pm \sqrt{b^2 - 4ac}}{2a}\]

\subsection{Matrix Example}
A sample matrix:
\[
\begin{pmatrix}
a_{11} & a_{12} & a_{13} \\
a_{21} & a_{22} & a_{23} \\
a_{31} & a_{32} & a_{33}
\end{pmatrix}
\]

\section{Tables}
Below is a sample table showing data:

\begin{table}[h]
\centering
\begin{tabular}{@{}lrrr@{}}
\toprule
\textbf{Category} & \textbf{Value 1} & \textbf{Value 2} & \textbf{Total} \\
\midrule
Item A & 10.5 & 20.3 & 30.8 \\
Item B & 15.2 & 25.1 & 40.3 \\
Item C & 22.1 & 12.7 & 34.8 \\
\bottomrule
\end{tabular}
\caption{Sample Data Table}
\label{tab:sample}
\end{table}

\section{Lists}
\subsection{Bullet Points}
\begin{itemize}
  \item First item
  \item Second item
  \item Third item with sub-items:
    \begin{itemize}
      \item Sub-item 1
      \item Sub-item 2
    \end{itemize}
\end{itemize}

\subsection{Numbered List}
\begin{enumerate}
  \item First step
  \item Second step
  \item Third step
\end{enumerate}

\section{Citations and References}
When you use the citation search feature, references will be automatically managed. For example, you might cite a paper like this: \cite{example2024}. All citations will be automatically added to your \texttt{cite.bib} file.

\section{Conclusion}
This template demonstrates basic \LaTeX{} features including mathematical formulas, tables, and lists. The bibliography system is now set up to work seamlessly with Bibby AI's citation management. Feel free to modify it for your needs!

\bibliographystyle{plain}
\bibliography{cite}

\end{document}\documentclass{article}
\usepackage[utf8]{inputenc}
\usepackage{booktabs}
\usepackage{amsmath}
\usepackage{natbib}

\title{My Document}
\author{Your Name}
\date{\today}

\begin{document}

\maketitle

\section{Introduction}
This is a sample document demonstrating various \LaTeX{} features including mathematical equations, tables, and formatting. You can cite papers using the citation search feature, and they will be automatically added to your \texttt{cite.bib} file.

\section{Mathematical Expressions}
Here are some example equations:

\subsection{Quadratic Formula}
The quadratic formula is given by:
\[x = \frac{-b \pm \sqrt{b^2 - 4ac}}{2a}\]

\subsection{Matrix Example}
A sample matrix:
\[
\begin{pmatrix}
a_{11} & a_{12} & a_{13} \\
a_{21} & a_{22} & a_{23} \\
a_{31} & a_{32} & a_{33}
\end{pmatrix}
\]

\section{Tables}
Below is a sample table showing data:

\begin{table}[h]
\centering
\begin{tabular}{@{}lrrr@{}}
\toprule
\textbf{Category} & \textbf{Value 1} & \textbf{Value 2} & \textbf{Total} \\
\midrule
Item A & 10.5 & 20.3 & 30.8 \\
Item B & 15.2 & 25.1 & 40.3 \\
Item C & 22.1 & 12.7 & 34.8 \\
\bottomrule
\end{tabular}
\caption{Sample Data Table}
\label{tab:sample}
\end{table}

\section{Lists}
\subsection{Bullet Points}
\begin{itemize}
  \item First item
  \item Second item
  \item Third item with sub-items:
    \begin{itemize}
      \item Sub-item 1
      \item Sub-item 2
    \end{itemize}
\end{itemize}

\subsection{Numbered List}
\begin{enumerate}
  \item First step
  \item Second step
  \item Third step
\end{enumerate}

\section{Citations and References}
When you use the citation search feature, references will be automatically managed. For example, you might cite a paper like this: \cite{example2024}. All citations will be automatically added to your \texttt{cite.bib} file.

\section{Conclusion}
This template demonstrates basic \LaTeX{} features including mathematical formulas, tables, and lists. The bibliography system is now set up to work seamlessly with Bibby AI's citation management. Feel free to modify it for your needs!

\bibliographystyle{plain}
\bibliography{cite}

\end{document}\documentclass{article}
\usepackage[utf8]{inputenc}
\usepackage{booktabs}
\usepackage{amsmath}
\usepackage{natbib}

\title{My Document}
\author{Your Name}
\date{\today}

\begin{document}

\maketitle

\section{Introduction}
This is a sample document demonstrating various \LaTeX{} features including mathematical equations, tables, and formatting. You can cite papers using the citation search feature, and they will be automatically added to your \texttt{cite.bib} file.

\section{Mathematical Expressions}
Here are some example equations:

\subsection{Quadratic Formula}
The quadratic formula is given by:
\[x = \frac{-b \pm \sqrt{b^2 - 4ac}}{2a}\]

\subsection{Matrix Example}
A sample matrix:
\[
\begin{pmatrix}
a_{11} & a_{12} & a_{13} \\
a_{21} & a_{22} & a_{23} \\
a_{31} & a_{32} & a_{33}
\end{pmatrix}
\]

\section{Tables}
Below is a sample table showing data:

\begin{table}[h]
\centering
\begin{tabular}{@{}lrrr@{}}
\toprule
\textbf{Category} & \textbf{Value 1} & \textbf{Value 2} & \textbf{Total} \\
\midrule
Item A & 10.5 & 20.3 & 30.8 \\
Item B & 15.2 & 25.1 & 40.3 \\
Item C & 22.1 & 12.7 & 34.8 \\
\bottomrule
\end{tabular}
\caption{Sample Data Table}
\label{tab:sample}
\end{table}

\section{Lists}
\subsection{Bullet Points}
\begin{itemize}
  \item First item
  \item Second item
  \item Third item with sub-items:
    \begin{itemize}
      \item Sub-item 1
      \item Sub-item 2
    \end{itemize}
\end{itemize}

\subsection{Numbered List}
\begin{enumerate}
  \item First step
  \item Second step
  \item Third step
\end{enumerate}

\section{Citations and References}
When you use the citation search feature, references will be automatically managed. For example, you might cite a paper like this: \cite{example2024}. All citations will be automatically added to your \texttt{cite.bib} file.

\section{Conclusion}
This template demonstrates basic \LaTeX{} features including mathematical formulas, tables, and lists. The bibliography system is now set up to work seamlessly with Bibby AI's citation management. Feel free to modify it for your needs!

\bibliographystyle{plain}
\bibliography{cite}

\end{document}\documentclass{article}
\usepackage[utf8]{inputenc}
\usepackage{booktabs}
\usepackage{amsmath}
\usepackage{natbib}

\title{My Document}
\author{Your Name}
\date{\today}

\begin{document}

\maketitle

\section{Introduction}
This is a sample document demonstrating various \LaTeX{} features including mathematical equations, tables, and formatting. You can cite papers using the citation search feature, and they will be automatically added to your \texttt{cite.bib} file.

\section{Mathematical Expressions}
Here are some example equations:

\subsection{Quadratic Formula}
The quadratic formula is given by:
\[x = \frac{-b \pm \sqrt{b^2 - 4ac}}{2a}\]

\subsection{Matrix Example}
A sample matrix:
\[
\begin{pmatrix}
a_{11} & a_{12} & a_{13} \\
a_{21} & a_{22} & a_{23} \\
a_{31} & a_{32} & a_{33}
\end{pmatrix}
\]

\section{Tables}
Below is a sample table showing data:

\begin{table}[h]
\centering
\begin{tabular}{@{}lrrr@{}}
\toprule
\textbf{Category} & \textbf{Value 1} & \textbf{Value 2} & \textbf{Total} \\
\midrule
Item A & 10.5 & 20.3 & 30.8 \\
Item B & 15.2 & 25.1 & 40.3 \\
Item C & 22.1 & 12.7 & 34.8 \\
\bottomrule
\end{tabular}
\caption{Sample Data Table}
\label{tab:sample}
\end{table}

\section{Lists}
\subsection{Bullet Points}
\begin{itemize}
  \item First item
  \item Second item
  \item Third item with sub-items:
    \begin{itemize}
      \item Sub-item 1
      \item Sub-item 2
    \end{itemize}
\end{itemize}

\subsection{Numbered List}
\begin{enumerate}
  \item First step
  \item Second step
  \item Third step
\end{enumerate}

\section{Citations and References}
When you use the citation search feature, references will be automatically managed. For example, you might cite a paper like this: \cite{example2024}. All citations will be automatically added to your \texttt{cite.bib} file.

\section{Conclusion}
This template demonstrates basic \LaTeX{} features including mathematical formulas, tables, and lists. The bibliography system is now set up to work seamlessly with Bibby AI's citation management. Feel free to modify it for your needs!

\bibliographystyle{plain}
\bibliography{cite}

\end{document}\documentclass{article}
\usepackage[utf8]{inputenc}
\usepackage{booktabs}
\usepackage{amsmath}
\usepackage{natbib}

\title{My Document}
\author{Your Name}
\date{\today}

\begin{document}

\maketitle

\section{Introduction}
This is a sample document demonstrating various \LaTeX{} features including mathematical equations, tables, and formatting. You can cite papers using the citation search feature, and they will be automatically added to your \texttt{cite.bib} file.

\section{Mathematical Expressions}
Here are some example equations:

\subsection{Quadratic Formula}
The quadratic formula is given by:
\[x = \frac{-b \pm \sqrt{b^2 - 4ac}}{2a}\]

\subsection{Matrix Example}
A sample matrix:
\[
\begin{pmatrix}
a_{11} & a_{12} & a_{13} \\
a_{21} & a_{22} & a_{23} \\
a_{31} & a_{32} & a_{33}
\end{pmatrix}
\]

\section{Tables}
Below is a sample table showing data:

\begin{table}[h]
\centering
\begin{tabular}{@{}lrrr@{}}
\toprule
\textbf{Category} & \textbf{Value 1} & \textbf{Value 2} & \textbf{Total} \\
\midrule
Item A & 10.5 & 20.3 & 30.8 \\
Item B & 15.2 & 25.1 & 40.3 \\
Item C & 22.1 & 12.7 & 34.8 \\
\bottomrule
\end{tabular}
\caption{Sample Data Table}
\label{tab:sample}
\end{table}

\section{Lists}
\subsection{Bullet Points}
\begin{itemize}
  \item First item
  \item Second item
  \item Third item with sub-items:
    \begin{itemize}
      \item Sub-item 1
      \item Sub-item 2
    \end{itemize}
\end{itemize}

\subsection{Numbered List}
\begin{enumerate}
  \item First step
  \item Second step
  \item Third step
\end{enumerate}

\section{Citations and References}
When you use the citation search feature, references will be automatically managed. For example, you might cite a paper like this: \cite{example2024}. All citations will be automatically added to your \texttt{cite.bib} file.

\section{Conclusion}
This template demonstrates basic \LaTeX{} features including mathematical formulas, tables, and lists. The bibliography system is now set up to work seamlessly with Bibby AI's citation management. Feel free to modify it for your needs!

\bibliographystyle{plain}
\bibliography{cite}

\end{document}\documentclass{article}
\usepackage[utf8]{inputenc}
\usepackage{booktabs}
\usepackage{amsmath}
\usepackage{natbib}

\title{My Document}
\author{Your Name}
\date{\today}

\begin{document}

\maketitle

\section{Introduction}
This is a sample document demonstrating various \LaTeX{} features including mathematical equations, tables, and formatting. You can cite papers using the citation search feature, and they will be automatically added to your \texttt{cite.bib} file.

\section{Mathematical Expressions}
Here are some example equations:

\subsection{Quadratic Formula}
The quadratic formula is given by:
\[x = \frac{-b \pm \sqrt{b^2 - 4ac}}{2a}\]

\subsection{Matrix Example}
A sample matrix:
\[
\begin{pmatrix}
a_{11} & a_{12} & a_{13} \\
a_{21} & a_{22} & a_{23} \\
a_{31} & a_{32} & a_{33}
\end{pmatrix}
\]

\section{Tables}
Below is a sample table showing data:

\begin{table}[h]
\centering
\begin{tabular}{@{}lrrr@{}}
\toprule
\textbf{Category} & \textbf{Value 1} & \textbf{Value 2} & \textbf{Total} \\
\midrule
Item A & 10.5 & 20.3 & 30.8 \\
Item B & 15.2 & 25.1 & 40.3 \\
Item C & 22.1 & 12.7 & 34.8 \\
\bottomrule
\end{tabular}
\caption{Sample Data Table}
\label{tab:sample}
\end{table}

\section{Lists}
\subsection{Bullet Points}
\begin{itemize}
  \item First item
  \item Second item
  \item Third item with sub-items:
    \begin{itemize}
      \item Sub-item 1
      \item Sub-item 2
    \end{itemize}
\end{itemize}

\subsection{Numbered List}
\begin{enumerate}
  \item First step
  \item Second step
  \item Third step
\end{enumerate}

\section{Citations and References}
When you use the citation search feature, references will be automatically managed. For example, you might cite a paper like this: \cite{example2024}. All citations will be automatically added to your \texttt{cite.bib} file.

\section{Conclusion}
This template demonstrates basic \LaTeX{} features including mathematical formulas, tables, and lists. The bibliography system is now set up to work seamlessly with Bibby AI's citation management. Feel free to modify it for your needs!

\bibliographystyle{plain}
\bibliography{cite}

\end{document}\documentclass{article}
\usepackage[utf8]{inputenc}
\usepackage{booktabs}
\usepackage{amsmath}
\usepackage{natbib}

\title{My Document}
\author{Your Name}
\date{\today}

\begin{document}

\maketitle

\section{Introduction}
This is a sample document demonstrating various \LaTeX{} features including mathematical equations, tables, and formatting. You can cite papers using the citation search feature, and they will be automatically added to your \texttt{cite.bib} file.

\section{Mathematical Expressions}
Here are some example equations:

\subsection{Quadratic Formula}
The quadratic formula is given by:
\[x = \frac{-b \pm \sqrt{b^2 - 4ac}}{2a}\]

\subsection{Matrix Example}
A sample matrix:
\[
\begin{pmatrix}
a_{11} & a_{12} & a_{13} \\
a_{21} & a_{22} & a_{23} \\
a_{31} & a_{32} & a_{33}
\end{pmatrix}
\]

\section{Tables}
Below is a sample table showing data:

\begin{table}[h]
\centering
\begin{tabular}{@{}lrrr@{}}
\toprule
\textbf{Category} & \textbf{Value 1} & \textbf{Value 2} & \textbf{Total} \\
\midrule
Item A & 10.5 & 20.3 & 30.8 \\
Item B & 15.2 & 25.1 & 40.3 \\
Item C & 22.1 & 12.7 & 34.8 \\
\bottomrule
\end{tabular}
\caption{Sample Data Table}
\label{tab:sample}
\end{table}

\section{Lists}
\subsection{Bullet Points}
\begin{itemize}
  \item First item
  \item Second item
  \item Third item with sub-items:
    \begin{itemize}
      \item Sub-item 1
      \item Sub-item 2
    \end{itemize}
\end{itemize}

\subsection{Numbered List}
\begin{enumerate}
  \item First step
  \item Second step
  \item Third step
\end{enumerate}

\section{Citations and References}
When you use the citation search feature, references will be automatically managed. For example, you might cite a paper like this: \cite{example2024}. All citations will be automatically added to your \texttt{cite.bib} file.

\section{Conclusion}
This template demonstrates basic \LaTeX{} features including mathematical formulas, tables, and lists. The bibliography system is now set up to work seamlessly with Bibby AI's citation management. Feel free to modify it for your needs!

\bibliographystyle{plain}
\bibliography{cite}

\end{document}\documentclass{article}
\usepackage[utf8]{inputenc}
\usepackage{booktabs}
\usepackage{amsmath}
\usepackage{natbib}

\title{My Document}
\author{Your Name}
\date{\today}

\begin{document}

\maketitle

\section{Introduction}
This is a sample document demonstrating various \LaTeX{} features including mathematical equations, tables, and formatting. You can cite papers using the citation search feature, and they will be automatically added to your \texttt{cite.bib} file.

\section{Mathematical Expressions}
Here are some example equations:

\subsection{Quadratic Formula}
The quadratic formula is given by:
\[x = \frac{-b \pm \sqrt{b^2 - 4ac}}{2a}\]

\subsection{Matrix Example}
A sample matrix:
\[
\begin{pmatrix}
a_{11} & a_{12} & a_{13} \\
a_{21} & a_{22} & a_{23} \\
a_{31} & a_{32} & a_{33}
\end{pmatrix}
\]

\section{Tables}
Below is a sample table showing data:

\begin{table}[h]
\centering
\begin{tabular}{@{}lrrr@{}}
\toprule
\textbf{Category} & \textbf{Value 1} & \textbf{Value 2} & \textbf{Total} \\
\midrule
Item A & 10.5 & 20.3 & 30.8 \\
Item B & 15.2 & 25.1 & 40.3 \\
Item C & 22.1 & 12.7 & 34.8 \\
\bottomrule
\end{tabular}
\caption{Sample Data Table}
\label{tab:sample}
\end{table}

\section{Lists}
\subsection{Bullet Points}
\begin{itemize}
  \item First item
  \item Second item
  \item Third item with sub-items:
    \begin{itemize}
      \item Sub-item 1
      \item Sub-item 2
    \end{itemize}
\end{itemize}

\subsection{Numbered List}
\begin{enumerate}
  \item First step
  \item Second step
  \item Third step
\end{enumerate}

\section{Citations and References}
When you use the citation search feature, references will be automatically managed. For example, you might cite a paper like this: \cite{example2024}. All citations will be automatically added to your \texttt{cite.bib} file.

\section{Conclusion}
This template demonstrates basic \LaTeX{} features including mathematical formulas, tables, and lists. The bibliography system is now set up to work seamlessly with Bibby AI's citation management. Feel free to modify it for your needs!

\bibliographystyle{plain}
\bibliography{cite}

\end{document}\documentclass{article}
\usepackage[utf8]{inputenc}
\usepackage{booktabs}
\usepackage{amsmath}
\usepackage{natbib}

\title{My Document}
\author{Your Name}
\date{\today}

\begin{document}

\maketitle

\section{Introduction}
This is a sample document demonstrating various \LaTeX{} features including mathematical equations, tables, and formatting. You can cite papers using the citation search feature, and they will be automatically added to your \texttt{cite.bib} file.

\section{Mathematical Expressions}
Here are some example equations:

\subsection{Quadratic Formula}
The quadratic formula is given by:
\[x = \frac{-b \pm \sqrt{b^2 - 4ac}}{2a}\]

\subsection{Matrix Example}
A sample matrix:
\[
\begin{pmatrix}
a_{11} & a_{12} & a_{13} \\
a_{21} & a_{22} & a_{23} \\
a_{31} & a_{32} & a_{33}
\end{pmatrix}
\]

\section{Tables}
Below is a sample table showing data:

\begin{table}[h]
\centering
\begin{tabular}{@{}lrrr@{}}
\toprule
\textbf{Category} & \textbf{Value 1} & \textbf{Value 2} & \textbf{Total} \\
\midrule
Item A & 10.5 & 20.3 & 30.8 \\
Item B & 15.2 & 25.1 & 40.3 \\
Item C & 22.1 & 12.7 & 34.8 \\
\bottomrule
\end{tabular}
\caption{Sample Data Table}
\label{tab:sample}
\end{table}

\section{Lists}
\subsection{Bullet Points}
\begin{itemize}
  \item First item
  \item Second item
  \item Third item with sub-items:
    \begin{itemize}
      \item Sub-item 1
      \item Sub-item 2
    \end{itemize}
\end{itemize}

\subsection{Numbered List}
\begin{enumerate}
  \item First step
  \item Second step
  \item Third step
\end{enumerate}

\section{Citations and References}
When you use the citation search feature, references will be automatically managed. For example, you might cite a paper like this: \cite{example2024}. All citations will be automatically added to your \texttt{cite.bib} file.

\section{Conclusion}
This template demonstrates basic \LaTeX{} features including mathematical formulas, tables, and lists. The bibliography system is now set up to work seamlessly with Bibby AI's citation management. Feel free to modify it for your needs!

\bibliographystyle{plain}
\bibliography{cite}

\end{document}\documentclass{article}
\usepackage[utf8]{inputenc}
\usepackage{booktabs}
\usepackage{amsmath}
\usepackage{natbib}

\title{My Document}
\author{Your Name}
\date{\today}

\begin{document}

\maketitle

\section{Introduction}
This is a sample document demonstrating various \LaTeX{} features including mathematical equations, tables, and formatting. You can cite papers using the citation search feature, and they will be automatically added to your \texttt{cite.bib} file.

\section{Mathematical Expressions}
Here are some example equations:

\subsection{Quadratic Formula}
The quadratic formula is given by:
\[x = \frac{-b \pm \sqrt{b^2 - 4ac}}{2a}\]

\subsection{Matrix Example}
A sample matrix:
\[
\begin{pmatrix}
a_{11} & a_{12} & a_{13} \\
a_{21} & a_{22} & a_{23} \\
a_{31} & a_{32} & a_{33}
\end{pmatrix}
\]

\section{Tables}
Below is a sample table showing data:

\begin{table}[h]
\centering
\begin{tabular}{@{}lrrr@{}}
\toprule
\textbf{Category} & \textbf{Value 1} & \textbf{Value 2} & \textbf{Total} \\
\midrule
Item A & 10.5 & 20.3 & 30.8 \\
Item B & 15.2 & 25.1 & 40.3 \\
Item C & 22.1 & 12.7 & 34.8 \\
\bottomrule
\end{tabular}
\caption{Sample Data Table}
\label{tab:sample}
\end{table}

\section{Lists}
\subsection{Bullet Points}
\begin{itemize}
  \item First item
  \item Second item
  \item Third item with sub-items:
    \begin{itemize}
      \item Sub-item 1
      \item Sub-item 2
    \end{itemize}
\end{itemize}

\subsection{Numbered List}
\begin{enumerate}
  \item First step
  \item Second step
  \item Third step
\end{enumerate}

\section{Citations and References}
When you use the citation search feature, references will be automatically managed. For example, you might cite a paper like this: \cite{example2024}. All citations will be automatically added to your \texttt{cite.bib} file.

\section{Conclusion}
This template demonstrates basic \LaTeX{} features including mathematical formulas, tables, and lists. The bibliography system is now set up to work seamlessly with Bibby AI's citation management. Feel free to modify it for your needs!

\bibliographystyle{plain}
\bibliography{cite}

\end{document}\documentclass{article}
\usepackage[utf8]{inputenc}
\usepackage{booktabs}
\usepackage{amsmath}
\usepackage{natbib}

\title{My Document}
\author{Your Name}
\date{\today}

\begin{document}

\maketitle

\section{Introduction}
This is a sample document demonstrating various \LaTeX{} features including mathematical equations, tables, and formatting. You can cite papers using the citation search feature, and they will be automatically added to your \texttt{cite.bib} file.

\section{Mathematical Expressions}
Here are some example equations:

\subsection{Quadratic Formula}
The quadratic formula is given by:
\[x = \frac{-b \pm \sqrt{b^2 - 4ac}}{2a}\]

\subsection{Matrix Example}
A sample matrix:
\[
\begin{pmatrix}
a_{11} & a_{12} & a_{13} \\
a_{21} & a_{22} & a_{23} \\
a_{31} & a_{32} & a_{33}
\end{pmatrix}
\]

\section{Tables}
Below is a sample table showing data:

\begin{table}[h]
\centering
\begin{tabular}{@{}lrrr@{}}
\toprule
\textbf{Category} & \textbf{Value 1} & \textbf{Value 2} & \textbf{Total} \\
\midrule
Item A & 10.5 & 20.3 & 30.8 \\
Item B & 15.2 & 25.1 & 40.3 \\
Item C & 22.1 & 12.7 & 34.8 \\
\bottomrule
\end{tabular}
\caption{Sample Data Table}
\label{tab:sample}
\end{table}

\section{Lists}
\subsection{Bullet Points}
\begin{itemize}
  \item First item
  \item Second item
  \item Third item with sub-items:
    \begin{itemize}
      \item Sub-item 1
      \item Sub-item 2
    \end{itemize}
\end{itemize}

\subsection{Numbered List}
\begin{enumerate}
  \item First step
  \item Second step
  \item Third step
\end{enumerate}

\section{Citations and References}
When you use the citation search feature, references will be automatically managed. For example, you might cite a paper like this: \cite{example2024}. All citations will be automatically added to your \texttt{cite.bib} file.

\section{Conclusion}
This template demonstrates basic \LaTeX{} features including mathematical formulas, tables, and lists. The bibliography system is now set up to work seamlessly with Bibby AI's citation management. Feel free to modify it for your needs!

\bibliographystyle{plain}
\bibliography{cite}

\end{document}\documentclass{article}
\usepackage[utf8]{inputenc}
\usepackage{booktabs}
\usepackage{amsmath}
\usepackage{natbib}

\title{My Document}
\author{Your Name}
\date{\today}

\begin{document}

\maketitle

\section{Introduction}
This is a sample document demonstrating various \LaTeX{} features including mathematical equations, tables, and formatting. You can cite papers using the citation search feature, and they will be automatically added to your \texttt{cite.bib} file.

\section{Mathematical Expressions}
Here are some example equations:

\subsection{Quadratic Formula}
The quadratic formula is given by:
\[x = \frac{-b \pm \sqrt{b^2 - 4ac}}{2a}\]

\subsection{Matrix Example}
A sample matrix:
\[
\begin{pmatrix}
a_{11} & a_{12} & a_{13} \\
a_{21} & a_{22} & a_{23} \\
a_{31} & a_{32} & a_{33}
\end{pmatrix}
\]

\section{Tables}
Below is a sample table showing data:

\begin{table}[h]
\centering
\begin{tabular}{@{}lrrr@{}}
\toprule
\textbf{Category} & \textbf{Value 1} & \textbf{Value 2} & \textbf{Total} \\
\midrule
Item A & 10.5 & 20.3 & 30.8 \\
Item B & 15.2 & 25.1 & 40.3 \\
Item C & 22.1 & 12.7 & 34.8 \\
\bottomrule
\end{tabular}
\caption{Sample Data Table}
\label{tab:sample}
\end{table}

\section{Lists}
\subsection{Bullet Points}
\begin{itemize}
  \item First item
  \item Second item
  \item Third item with sub-items:
    \begin{itemize}
      \item Sub-item 1
      \item Sub-item 2
    \end{itemize}
\end{itemize}

\subsection{Numbered List}
\begin{enumerate}
  \item First step
  \item Second step
  \item Third step
\end{enumerate}

\section{Citations and References}
When you use the citation search feature, references will be automatically managed. For example, you might cite a paper like this: \cite{example2024}. All citations will be automatically added to your \texttt{cite.bib} file.

\section{Conclusion}
This template demonstrates basic \LaTeX{} features including mathematical formulas, tables, and lists. The bibliography system is now set up to work seamlessly with Bibby AI's citation management. Feel free to modify it for your needs!

\bibliographystyle{plain}
\bibliography{cite}

\end{document}\documentclass{article}
\usepackage[utf8]{inputenc}
\usepackage{booktabs}
\usepackage{amsmath}
\usepackage{natbib}

\title{My Document}
\author{Your Name}
\date{\today}

\begin{document}

\maketitle

\section{Introduction}
This is a sample document demonstrating various \LaTeX{} features including mathematical equations, tables, and formatting. You can cite papers using the citation search feature, and they will be automatically added to your \texttt{cite.bib} file.

\section{Mathematical Expressions}
Here are some example equations:

\subsection{Quadratic Formula}
The quadratic formula is given by:
\[x = \frac{-b \pm \sqrt{b^2 - 4ac}}{2a}\]

\subsection{Matrix Example}
A sample matrix:
\[
\begin{pmatrix}
a_{11} & a_{12} & a_{13} \\
a_{21} & a_{22} & a_{23} \\
a_{31} & a_{32} & a_{33}
\end{pmatrix}
\]

\section{Tables}
Below is a sample table showing data:

\begin{table}[h]
\centering
\begin{tabular}{@{}lrrr@{}}
\toprule
\textbf{Category} & \textbf{Value 1} & \textbf{Value 2} & \textbf{Total} \\
\midrule
Item A & 10.5 & 20.3 & 30.8 \\
Item B & 15.2 & 25.1 & 40.3 \\
Item C & 22.1 & 12.7 & 34.8 \\
\bottomrule
\end{tabular}
\caption{Sample Data Table}
\label{tab:sample}
\end{table}

\section{Lists}
\subsection{Bullet Points}
\begin{itemize}
  \item First item
  \item Second item
  \item Third item with sub-items:
    \begin{itemize}
      \item Sub-item 1
      \item Sub-item 2
    \end{itemize}
\end{itemize}

\subsection{Numbered List}
\begin{enumerate}
  \item First step
  \item Second step
  \item Third step
\end{enumerate}

\section{Citations and References}
When you use the citation search feature, references will be automatically managed. For example, you might cite a paper like this: \cite{example2024}. All citations will be automatically added to your \texttt{cite.bib} file.

\section{Conclusion}
This template demonstrates basic \LaTeX{} features including mathematical formulas, tables, and lists. The bibliography system is now set up to work seamlessly with Bibby AI's citation management. Feel free to modify it for your needs!

\bibliographystyle{plain}
\bibliography{cite}

\end{document}\documentclass{article}
\usepackage[utf8]{inputenc}
\usepackage{booktabs}
\usepackage{amsmath}
\usepackage{natbib}

\title{My Document}
\author{Your Name}
\date{\today}

\begin{document}

\maketitle

\section{Introduction}
This is a sample document demonstrating various \LaTeX{} features including mathematical equations, tables, and formatting. You can cite papers using the citation search feature, and they will be automatically added to your \texttt{cite.bib} file.

\section{Mathematical Expressions}
Here are some example equations:

\subsection{Quadratic Formula}
The quadratic formula is given by:
\[x = \frac{-b \pm \sqrt{b^2 - 4ac}}{2a}\]

\subsection{Matrix Example}
A sample matrix:
\[
\begin{pmatrix}
a_{11} & a_{12} & a_{13} \\
a_{21} & a_{22} & a_{23} \\
a_{31} & a_{32} & a_{33}
\end{pmatrix}
\]

\section{Tables}
Below is a sample table showing data:

\begin{table}[h]
\centering
\begin{tabular}{@{}lrrr@{}}
\toprule
\textbf{Category} & \textbf{Value 1} & \textbf{Value 2} & \textbf{Total} \\
\midrule
Item A & 10.5 & 20.3 & 30.8 \\
Item B & 15.2 & 25.1 & 40.3 \\
Item C & 22.1 & 12.7 & 34.8 \\
\bottomrule
\end{tabular}
\caption{Sample Data Table}
\label{tab:sample}
\end{table}

\section{Lists}
\subsection{Bullet Points}
\begin{itemize}
  \item First item
  \item Second item
  \item Third item with sub-items:
    \begin{itemize}
      \item Sub-item 1
      \item Sub-item 2
    \end{itemize}
\end{itemize}

\subsection{Numbered List}
\begin{enumerate}
  \item First step
  \item Second step
  \item Third step
\end{enumerate}

\section{Citations and References}
When you use the citation search feature, references will be automatically managed. For example, you might cite a paper like this: \cite{example2024}. All citations will be automatically added to your \texttt{cite.bib} file.

\section{Conclusion}
This template demonstrates basic \LaTeX{} features including mathematical formulas, tables, and lists. The bibliography system is now set up to work seamlessly with Bibby AI's citation management. Feel free to modify it for your needs!

\bibliographystyle{plain}
\bibliography{cite}

\end{document}\documentclass{article}
\usepackage[utf8]{inputenc}
\usepackage{booktabs}
\usepackage{amsmath}
\usepackage{natbib}

\title{My Document}
\author{Your Name}
\date{\today}

\begin{document}

\maketitle

\section{Introduction}
This is a sample document demonstrating various \LaTeX{} features including mathematical equations, tables, and formatting. You can cite papers using the citation search feature, and they will be automatically added to your \texttt{cite.bib} file.

\section{Mathematical Expressions}
Here are some example equations:

\subsection{Quadratic Formula}
The quadratic formula is given by:
\[x = \frac{-b \pm \sqrt{b^2 - 4ac}}{2a}\]

\subsection{Matrix Example}
A sample matrix:
\[
\begin{pmatrix}
a_{11} & a_{12} & a_{13} \\
a_{21} & a_{22} & a_{23} \\
a_{31} & a_{32} & a_{33}
\end{pmatrix}
\]

\section{Tables}
Below is a sample table showing data:

\begin{table}[h]
\centering
\begin{tabular}{@{}lrrr@{}}
\toprule
\textbf{Category} & \textbf{Value 1} & \textbf{Value 2} & \textbf{Total} \\
\midrule
Item A & 10.5 & 20.3 & 30.8 \\
Item B & 15.2 & 25.1 & 40.3 \\
Item C & 22.1 & 12.7 & 34.8 \\
\bottomrule
\end{tabular}
\caption{Sample Data Table}
\label{tab:sample}
\end{table}

\section{Lists}
\subsection{Bullet Points}
\begin{itemize}
  \item First item
  \item Second item
  \item Third item with sub-items:
    \begin{itemize}
      \item Sub-item 1
      \item Sub-item 2
    \end{itemize}
\end{itemize}

\subsection{Numbered List}
\begin{enumerate}
  \item First step
  \item Second step
  \item Third step
\end{enumerate}

\section{Citations and References}
When you use the citation search feature, references will be automatically managed. For example, you might cite a paper like this: \cite{example2024}. All citations will be automatically added to your \texttt{cite.bib} file.

\section{Conclusion}
This template demonstrates basic \LaTeX{} features including mathematical formulas, tables, and lists. The bibliography system is now set up to work seamlessly with Bibby AI's citation management. Feel free to modify it for your needs!

\bibliographystyle{plain}
\bibliography{cite}

\end{document}\documentclass{article}
\usepackage[utf8]{inputenc}
\usepackage{booktabs}
\usepackage{amsmath}
\usepackage{natbib}

\title{My Document}
\author{Your Name}
\date{\today}

\begin{document}

\maketitle

\section{Introduction}
This is a sample document demonstrating various \LaTeX{} features including mathematical equations, tables, and formatting. You can cite papers using the citation search feature, and they will be automatically added to your \texttt{cite.bib} file.

\section{Mathematical Expressions}
Here are some example equations:

\subsection{Quadratic Formula}
The quadratic formula is given by:
\[x = \frac{-b \pm \sqrt{b^2 - 4ac}}{2a}\]

\subsection{Matrix Example}
A sample matrix:
\[
\begin{pmatrix}
a_{11} & a_{12} & a_{13} \\
a_{21} & a_{22} & a_{23} \\
a_{31} & a_{32} & a_{33}
\end{pmatrix}
\]

\section{Tables}
Below is a sample table showing data:

\begin{table}[h]
\centering
\begin{tabular}{@{}lrrr@{}}
\toprule
\textbf{Category} & \textbf{Value 1} & \textbf{Value 2} & \textbf{Total} \\
\midrule
Item A & 10.5 & 20.3 & 30.8 \\
Item B & 15.2 & 25.1 & 40.3 \\
Item C & 22.1 & 12.7 & 34.8 \\
\bottomrule
\end{tabular}
\caption{Sample Data Table}
\label{tab:sample}
\end{table}

\section{Lists}
\subsection{Bullet Points}
\begin{itemize}
  \item First item
  \item Second item
  \item Third item with sub-items:
    \begin{itemize}
      \item Sub-item 1
      \item Sub-item 2
    \end{itemize}
\end{itemize}

\subsection{Numbered List}
\begin{enumerate}
  \item First step
  \item Second step
  \item Third step
\end{enumerate}

\section{Citations and References}
When you use the citation search feature, references will be automatically managed. For example, you might cite a paper like this: \cite{example2024}. All citations will be automatically added to your \texttt{cite.bib} file.

\section{Conclusion}
This template demonstrates basic \LaTeX{} features including mathematical formulas, tables, and lists. The bibliography system is now set up to work seamlessly with Bibby AI's citation management. Feel free to modify it for your needs!

\bibliographystyle{plain}
\bibliography{cite}

\end{document}\documentclass{article}
\usepackage[utf8]{inputenc}
\usepackage{booktabs}
\usepackage{amsmath}
\usepackage{natbib}

\title{My Document}
\author{Your Name}
\date{\today}

\begin{document}

\maketitle

\section{Introduction}
This is a sample document demonstrating various \LaTeX{} features including mathematical equations, tables, and formatting. You can cite papers using the citation search feature, and they will be automatically added to your \texttt{cite.bib} file.

\section{Mathematical Expressions}
Here are some example equations:

\subsection{Quadratic Formula}
The quadratic formula is given by:
\[x = \frac{-b \pm \sqrt{b^2 - 4ac}}{2a}\]

\subsection{Matrix Example}
A sample matrix:
\[
\begin{pmatrix}
a_{11} & a_{12} & a_{13} \\
a_{21} & a_{22} & a_{23} \\
a_{31} & a_{32} & a_{33}
\end{pmatrix}
\]

\section{Tables}
Below is a sample table showing data:

\begin{table}[h]
\centering
\begin{tabular}{@{}lrrr@{}}
\toprule
\textbf{Category} & \textbf{Value 1} & \textbf{Value 2} & \textbf{Total} \\
\midrule
Item A & 10.5 & 20.3 & 30.8 \\
Item B & 15.2 & 25.1 & 40.3 \\
Item C & 22.1 & 12.7 & 34.8 \\
\bottomrule
\end{tabular}
\caption{Sample Data Table}
\label{tab:sample}
\end{table}

\section{Lists}
\subsection{Bullet Points}
\begin{itemize}
  \item First item
  \item Second item
  \item Third item with sub-items:
    \begin{itemize}
      \item Sub-item 1
      \item Sub-item 2
    \end{itemize}
\end{itemize}

\subsection{Numbered List}
\begin{enumerate}
  \item First step
  \item Second step
  \item Third step
\end{enumerate}

\section{Citations and References}
When you use the citation search feature, references will be automatically managed. For example, you might cite a paper like this: \cite{example2024}. All citations will be automatically added to your \texttt{cite.bib} file.

\section{Conclusion}
This template demonstrates basic \LaTeX{} features including mathematical formulas, tables, and lists. The bibliography system is now set up to work seamlessly with Bibby AI's citation management. Feel free to modify it for your needs!

\bibliographystyle{plain}
\bibliography{cite}

\end{document}\documentclass{article}
\usepackage[utf8]{inputenc}
\usepackage{booktabs}
\usepackage{amsmath}
\usepackage{natbib}

\title{My Document}
\author{Your Name}
\date{\today}

\begin{document}

\maketitle

\section{Introduction}
This is a sample document demonstrating various \LaTeX{} features including mathematical equations, tables, and formatting. You can cite papers using the citation search feature, and they will be automatically added to your \texttt{cite.bib} file.

\section{Mathematical Expressions}
Here are some example equations:

\subsection{Quadratic Formula}
The quadratic formula is given by:
\[x = \frac{-b \pm \sqrt{b^2 - 4ac}}{2a}\]

\subsection{Matrix Example}
A sample matrix:
\[
\begin{pmatrix}
a_{11} & a_{12} & a_{13} \\
a_{21} & a_{22} & a_{23} \\
a_{31} & a_{32} & a_{33}
\end{pmatrix}
\]

\section{Tables}
Below is a sample table showing data:

\begin{table}[h]
\centering
\begin{tabular}{@{}lrrr@{}}
\toprule
\textbf{Category} & \textbf{Value 1} & \textbf{Value 2} & \textbf{Total} \\
\midrule
Item A & 10.5 & 20.3 & 30.8 \\
Item B & 15.2 & 25.1 & 40.3 \\
Item C & 22.1 & 12.7 & 34.8 \\
\bottomrule
\end{tabular}
\caption{Sample Data Table}
\label{tab:sample}
\end{table}

\section{Lists}
\subsection{Bullet Points}
\begin{itemize}
  \item First item
  \item Second item
  \item Third item with sub-items:
    \begin{itemize}
      \item Sub-item 1
      \item Sub-item 2
    \end{itemize}
\end{itemize}

\subsection{Numbered List}
\begin{enumerate}
  \item First step
  \item Second step
  \item Third step
\end{enumerate}

\section{Citations and References}
When you use the citation search feature, references will be automatically managed. For example, you might cite a paper like this: \cite{example2024}. All citations will be automatically added to your \texttt{cite.bib} file.

\section{Conclusion}
This template demonstrates basic \LaTeX{} features including mathematical formulas, tables, and lists. The bibliography system is now set up to work seamlessly with Bibby AI's citation management. Feel free to modify it for your needs!

\bibliographystyle{plain}
\bibliography{cite}

\end{document}\documentclass{article}
\usepackage[utf8]{inputenc}
\usepackage{booktabs}
\usepackage{amsmath}
\usepackage{natbib}

\title{My Document}
\author{Your Name}
\date{\today}

\begin{document}

\maketitle

\section{Introduction}
This is a sample document demonstrating various \LaTeX{} features including mathematical equations, tables, and formatting. You can cite papers using the citation search feature, and they will be automatically added to your \texttt{cite.bib} file.

\section{Mathematical Expressions}
Here are some example equations:

\subsection{Quadratic Formula}
The quadratic formula is given by:
\[x = \frac{-b \pm \sqrt{b^2 - 4ac}}{2a}\]

\subsection{Matrix Example}
A sample matrix:
\[
\begin{pmatrix}
a_{11} & a_{12} & a_{13} \\
a_{21} & a_{22} & a_{23} \\
a_{31} & a_{32} & a_{33}
\end{pmatrix}
\]

\section{Tables}
Below is a sample table showing data:

\begin{table}[h]
\centering
\begin{tabular}{@{}lrrr@{}}
\toprule
\textbf{Category} & \textbf{Value 1} & \textbf{Value 2} & \textbf{Total} \\
\midrule
Item A & 10.5 & 20.3 & 30.8 \\
Item B & 15.2 & 25.1 & 40.3 \\
Item C & 22.1 & 12.7 & 34.8 \\
\bottomrule
\end{tabular}
\caption{Sample Data Table}
\label{tab:sample}
\end{table}

\section{Lists}
\subsection{Bullet Points}
\begin{itemize}
  \item First item
  \item Second item
  \item Third item with sub-items:
    \begin{itemize}
      \item Sub-item 1
      \item Sub-item 2
    \end{itemize}
\end{itemize}

\subsection{Numbered List}
\begin{enumerate}
  \item First step
  \item Second step
  \item Third step
\end{enumerate}

\section{Citations and References}
When you use the citation search feature, references will be automatically managed. For example, you might cite a paper like this: \cite{example2024}. All citations will be automatically added to your \texttt{cite.bib} file.

\section{Conclusion}
This template demonstrates basic \LaTeX{} features including mathematical formulas, tables, and lists. The bibliography system is now set up to work seamlessly with Bibby AI's citation management. Feel free to modify it for your needs!

\bibliographystyle{plain}
\bibliography{cite}

\end{document}\documentclass{article}
\usepackage[utf8]{inputenc}
\usepackage{booktabs}
\usepackage{amsmath}
\usepackage{natbib}

\title{My Document}
\author{Your Name}
\date{\today}

\begin{document}

\maketitle

\section{Introduction}
This is a sample document demonstrating various \LaTeX{} features including mathematical equations, tables, and formatting. You can cite papers using the citation search feature, and they will be automatically added to your \texttt{cite.bib} file.

\section{Mathematical Expressions}
Here are some example equations:

\subsection{Quadratic Formula}
The quadratic formula is given by:
\[x = \frac{-b \pm \sqrt{b^2 - 4ac}}{2a}\]

\subsection{Matrix Example}
A sample matrix:
\[
\begin{pmatrix}
a_{11} & a_{12} & a_{13} \\
a_{21} & a_{22} & a_{23} \\
a_{31} & a_{32} & a_{33}
\end{pmatrix}
\]

\section{Tables}
Below is a sample table showing data:

\begin{table}[h]
\centering
\begin{tabular}{@{}lrrr@{}}
\toprule
\textbf{Category} & \textbf{Value 1} & \textbf{Value 2} & \textbf{Total} \\
\midrule
Item A & 10.5 & 20.3 & 30.8 \\
Item B & 15.2 & 25.1 & 40.3 \\
Item C & 22.1 & 12.7 & 34.8 \\
\bottomrule
\end{tabular}
\caption{Sample Data Table}
\label{tab:sample}
\end{table}

\section{Lists}
\subsection{Bullet Points}
\begin{itemize}
  \item First item
  \item Second item
  \item Third item with sub-items:
    \begin{itemize}
      \item Sub-item 1
      \item Sub-item 2
    \end{itemize}
\end{itemize}

\subsection{Numbered List}
\begin{enumerate}
  \item First step
  \item Second step
  \item Third step
\end{enumerate}

\section{Citations and References}
When you use the citation search feature, references will be automatically managed. For example, you might cite a paper like this: \cite{example2024}. All citations will be automatically added to your \texttt{cite.bib} file.

\section{Conclusion}
This template demonstrates basic \LaTeX{} features including mathematical formulas, tables, and lists. The bibliography system is now set up to work seamlessly with Bibby AI's citation management. Feel free to modify it for your needs!

\bibliographystyle{plain}
\bibliography{cite}

\end{document}\documentclass{article}
\usepackage[utf8]{inputenc}
\usepackage{booktabs}
\usepackage{amsmath}
\usepackage{natbib}

\title{My Document}
\author{Your Name}
\date{\today}

\begin{document}

\maketitle

\section{Introduction}
This is a sample document demonstrating various \LaTeX{} features including mathematical equations, tables, and formatting. You can cite papers using the citation search feature, and they will be automatically added to your \texttt{cite.bib} file.

\section{Mathematical Expressions}
Here are some example equations:

\subsection{Quadratic Formula}
The quadratic formula is given by:
\[x = \frac{-b \pm \sqrt{b^2 - 4ac}}{2a}\]

\subsection{Matrix Example}
A sample matrix:
\[
\begin{pmatrix}
a_{11} & a_{12} & a_{13} \\
a_{21} & a_{22} & a_{23} \\
a_{31} & a_{32} & a_{33}
\end{pmatrix}
\]

\section{Tables}
Below is a sample table showing data:

\begin{table}[h]
\centering
\begin{tabular}{@{}lrrr@{}}
\toprule
\textbf{Category} & \textbf{Value 1} & \textbf{Value 2} & \textbf{Total} \\
\midrule
Item A & 10.5 & 20.3 & 30.8 \\
Item B & 15.2 & 25.1 & 40.3 \\
Item C & 22.1 & 12.7 & 34.8 \\
\bottomrule
\end{tabular}
\caption{Sample Data Table}
\label{tab:sample}
\end{table}

\section{Lists}
\subsection{Bullet Points}
\begin{itemize}
  \item First item
  \item Second item
  \item Third item with sub-items:
    \begin{itemize}
      \item Sub-item 1
      \item Sub-item 2
    \end{itemize}
\end{itemize}

\subsection{Numbered List}
\begin{enumerate}
  \item First step
  \item Second step
  \item Third step
\end{enumerate}

\section{Citations and References}
When you use the citation search feature, references will be automatically managed. For example, you might cite a paper like this: \cite{example2024}. All citations will be automatically added to your \texttt{cite.bib} file.

\section{Conclusion}
This template demonstrates basic \LaTeX{} features including mathematical formulas, tables, and lists. The bibliography system is now set up to work seamlessly with Bibby AI's citation management. Feel free to modify it for your needs!

\bibliographystyle{plain}
\bibliography{cite}

\end{document}\documentclass{article}
\usepackage[utf8]{inputenc}
\usepackage{booktabs}
\usepackage{amsmath}
\usepackage{natbib}

\title{My Document}
\author{Your Name}
\date{\today}

\begin{document}

\maketitle

\section{Introduction}
This is a sample document demonstrating various \LaTeX{} features including mathematical equations, tables, and formatting. You can cite papers using the citation search feature, and they will be automatically added to your \texttt{cite.bib} file.

\section{Mathematical Expressions}
Here are some example equations:

\subsection{Quadratic Formula}
The quadratic formula is given by:
\[x = \frac{-b \pm \sqrt{b^2 - 4ac}}{2a}\]

\subsection{Matrix Example}
A sample matrix:
\[
\begin{pmatrix}
a_{11} & a_{12} & a_{13} \\
a_{21} & a_{22} & a_{23} \\
a_{31} & a_{32} & a_{33}
\end{pmatrix}
\]

\section{Tables}
Below is a sample table showing data:

\begin{table}[h]
\centering
\begin{tabular}{@{}lrrr@{}}
\toprule
\textbf{Category} & \textbf{Value 1} & \textbf{Value 2} & \textbf{Total} \\
\midrule
Item A & 10.5 & 20.3 & 30.8 \\
Item B & 15.2 & 25.1 & 40.3 \\
Item C & 22.1 & 12.7 & 34.8 \\
\bottomrule
\end{tabular}
\caption{Sample Data Table}
\label{tab:sample}
\end{table}

\section{Lists}
\subsection{Bullet Points}
\begin{itemize}
  \item First item
  \item Second item
  \item Third item with sub-items:
    \begin{itemize}
      \item Sub-item 1
      \item Sub-item 2
    \end{itemize}
\end{itemize}

\subsection{Numbered List}
\begin{enumerate}
  \item First step
  \item Second step
  \item Third step
\end{enumerate}

\section{Citations and References}
When you use the citation search feature, references will be automatically managed. For example, you might cite a paper like this: \cite{example2024}. All citations will be automatically added to your \texttt{cite.bib} file.

\section{Conclusion}
This template demonstrates basic \LaTeX{} features including mathematical formulas, tables, and lists. The bibliography system is now set up to work seamlessly with Bibby AI's citation management. Feel free to modify it for your needs!

\bibliographystyle{plain}
\bibliography{cite}

\end{document}\documentclass{article}
\usepackage[utf8]{inputenc}
\usepackage{booktabs}
\usepackage{amsmath}
\usepackage{natbib}

\title{My Document}
\author{Your Name}
\date{\today}

\begin{document}

\maketitle

\section{Introduction}
This is a sample document demonstrating various \LaTeX{} features including mathematical equations, tables, and formatting. You can cite papers using the citation search feature, and they will be automatically added to your \texttt{cite.bib} file.

\section{Mathematical Expressions}
Here are some example equations:

\subsection{Quadratic Formula}
The quadratic formula is given by:
\[x = \frac{-b \pm \sqrt{b^2 - 4ac}}{2a}\]

\subsection{Matrix Example}
A sample matrix:
\[
\begin{pmatrix}
a_{11} & a_{12} & a_{13} \\
a_{21} & a_{22} & a_{23} \\
a_{31} & a_{32} & a_{33}
\end{pmatrix}
\]

\section{Tables}
Below is a sample table showing data:

\begin{table}[h]
\centering
\begin{tabular}{@{}lrrr@{}}
\toprule
\textbf{Category} & \textbf{Value 1} & \textbf{Value 2} & \textbf{Total} \\
\midrule
Item A & 10.5 & 20.3 & 30.8 \\
Item B & 15.2 & 25.1 & 40.3 \\
Item C & 22.1 & 12.7 & 34.8 \\
\bottomrule
\end{tabular}
\caption{Sample Data Table}
\label{tab:sample}
\end{table}

\section{Lists}
\subsection{Bullet Points}
\begin{itemize}
  \item First item
  \item Second item
  \item Third item with sub-items:
    \begin{itemize}
      \item Sub-item 1
      \item Sub-item 2
    \end{itemize}
\end{itemize}

\subsection{Numbered List}
\begin{enumerate}
  \item First step
  \item Second step
  \item Third step
\end{enumerate}

\section{Citations and References}
When you use the citation search feature, references will be automatically managed. For example, you might cite a paper like this: \cite{example2024}. All citations will be automatically added to your \texttt{cite.bib} file.

\section{Conclusion}
This template demonstrates basic \LaTeX{} features including mathematical formulas, tables, and lists. The bibliography system is now set up to work seamlessly with Bibby AI's citation management. Feel free to modify it for your needs!

\bibliographystyle{plain}
\bibliography{cite}

\end{document}\documentclass{article}
\usepackage[utf8]{inputenc}
\usepackage{booktabs}
\usepackage{amsmath}
\usepackage{natbib}

\title{My Document}
\author{Your Name}
\date{\today}

\begin{document}

\maketitle

\section{Introduction}
This is a sample document demonstrating various \LaTeX{} features including mathematical equations, tables, and formatting. You can cite papers using the citation search feature, and they will be automatically added to your \texttt{cite.bib} file.

\section{Mathematical Expressions}
Here are some example equations:

\subsection{Quadratic Formula}
The quadratic formula is given by:
\[x = \frac{-b \pm \sqrt{b^2 - 4ac}}{2a}\]

\subsection{Matrix Example}
A sample matrix:
\[
\begin{pmatrix}
a_{11} & a_{12} & a_{13} \\
a_{21} & a_{22} & a_{23} \\
a_{31} & a_{32} & a_{33}
\end{pmatrix}
\]

\section{Tables}
Below is a sample table showing data:

\begin{table}[h]
\centering
\begin{tabular}{@{}lrrr@{}}
\toprule
\textbf{Category} & \textbf{Value 1} & \textbf{Value 2} & \textbf{Total} \\
\midrule
Item A & 10.5 & 20.3 & 30.8 \\
Item B & 15.2 & 25.1 & 40.3 \\
Item C & 22.1 & 12.7 & 34.8 \\
\bottomrule
\end{tabular}
\caption{Sample Data Table}
\label{tab:sample}
\end{table}

\section{Lists}
\subsection{Bullet Points}
\begin{itemize}
  \item First item
  \item Second item
  \item Third item with sub-items:
    \begin{itemize}
      \item Sub-item 1
      \item Sub-item 2
    \end{itemize}
\end{itemize}

\subsection{Numbered List}
\begin{enumerate}
  \item First step
  \item Second step
  \item Third step
\end{enumerate}

\section{Citations and References}
When you use the citation search feature, references will be automatically managed. For example, you might cite a paper like this: \cite{example2024}. All citations will be automatically added to your \texttt{cite.bib} file.

\section{Conclusion}
This template demonstrates basic \LaTeX{} features including mathematical formulas, tables, and lists. The bibliography system is now set up to work seamlessly with Bibby AI's citation management. Feel free to modify it for your needs!

\bibliographystyle{plain}
\bibliography{cite}

\end{document}\documentclass{article}
\usepackage[utf8]{inputenc}
\usepackage{booktabs}
\usepackage{amsmath}
\usepackage{natbib}

\title{My Document}
\author{Your Name}
\date{\today}

\begin{document}

\maketitle

\section{Introduction}
This is a sample document demonstrating various \LaTeX{} features including mathematical equations, tables, and formatting. You can cite papers using the citation search feature, and they will be automatically added to your \texttt{cite.bib} file.

\section{Mathematical Expressions}
Here are some example equations:

\subsection{Quadratic Formula}
The quadratic formula is given by:
\[x = \frac{-b \pm \sqrt{b^2 - 4ac}}{2a}\]

\subsection{Matrix Example}
A sample matrix:
\[
\begin{pmatrix}
a_{11} & a_{12} & a_{13} \\
a_{21} & a_{22} & a_{23} \\
a_{31} & a_{32} & a_{33}
\end{pmatrix}
\]

\section{Tables}
Below is a sample table showing data:

\begin{table}[h]
\centering
\begin{tabular}{@{}lrrr@{}}
\toprule
\textbf{Category} & \textbf{Value 1} & \textbf{Value 2} & \textbf{Total} \\
\midrule
Item A & 10.5 & 20.3 & 30.8 \\
Item B & 15.2 & 25.1 & 40.3 \\
Item C & 22.1 & 12.7 & 34.8 \\
\bottomrule
\end{tabular}
\caption{Sample Data Table}
\label{tab:sample}
\end{table}

\section{Lists}
\subsection{Bullet Points}
\begin{itemize}
  \item First item
  \item Second item
  \item Third item with sub-items:
    \begin{itemize}
      \item Sub-item 1
      \item Sub-item 2
    \end{itemize}
\end{itemize}

\subsection{Numbered List}
\begin{enumerate}
  \item First step
  \item Second step
  \item Third step
\end{enumerate}

\section{Citations and References}
When you use the citation search feature, references will be automatically managed. For example, you might cite a paper like this: \cite{example2024}. All citations will be automatically added to your \texttt{cite.bib} file.

\section{Conclusion}
This template demonstrates basic \LaTeX{} features including mathematical formulas, tables, and lists. The bibliography system is now set up to work seamlessly with Bibby AI's citation management. Feel free to modify it for your needs!

\bibliographystyle{plain}
\bibliography{cite}

\end{document}\documentclass{article}
\usepackage[utf8]{inputenc}
\usepackage{booktabs}
\usepackage{amsmath}
\usepackage{natbib}

\title{My Document}
\author{Your Name}
\date{\today}

\begin{document}

\maketitle

\section{Introduction}
This is a sample document demonstrating various \LaTeX{} features including mathematical equations, tables, and formatting. You can cite papers using the citation search feature, and they will be automatically added to your \texttt{cite.bib} file.

\section{Mathematical Expressions}
Here are some example equations:

\subsection{Quadratic Formula}
The quadratic formula is given by:
\[x = \frac{-b \pm \sqrt{b^2 - 4ac}}{2a}\]

\subsection{Matrix Example}
A sample matrix:
\[
\begin{pmatrix}
a_{11} & a_{12} & a_{13} \\
a_{21} & a_{22} & a_{23} \\
a_{31} & a_{32} & a_{33}
\end{pmatrix}
\]

\section{Tables}
Below is a sample table showing data:

\begin{table}[h]
\centering
\begin{tabular}{@{}lrrr@{}}
\toprule
\textbf{Category} & \textbf{Value 1} & \textbf{Value 2} & \textbf{Total} \\
\midrule
Item A & 10.5 & 20.3 & 30.8 \\
Item B & 15.2 & 25.1 & 40.3 \\
Item C & 22.1 & 12.7 & 34.8 \\
\bottomrule
\end{tabular}
\caption{Sample Data Table}
\label{tab:sample}
\end{table}

\section{Lists}
\subsection{Bullet Points}
\begin{itemize}
  \item First item
  \item Second item
  \item Third item with sub-items:
    \begin{itemize}
      \item Sub-item 1
      \item Sub-item 2
    \end{itemize}
\end{itemize}

\subsection{Numbered List}
\begin{enumerate}
  \item First step
  \item Second step
  \item Third step
\end{enumerate}

\section{Citations and References}
When you use the citation search feature, references will be automatically managed. For example, you might cite a paper like this: \cite{example2024}. All citations will be automatically added to your \texttt{cite.bib} file.

\section{Conclusion}
This template demonstrates basic \LaTeX{} features including mathematical formulas, tables, and lists. The bibliography system is now set up to work seamlessly with Bibby AI's citation management. Feel free to modify it for your needs!

\bibliographystyle{plain}
\bibliography{cite}

\end{document}\documentclass{article}
\usepackage[utf8]{inputenc}
\usepackage{booktabs}
\usepackage{amsmath}
\usepackage{natbib}

\title{My Document}
\author{Your Name}
\date{\today}

\begin{document}

\maketitle

\section{Introduction}
This is a sample document demonstrating various \LaTeX{} features including mathematical equations, tables, and formatting. You can cite papers using the citation search feature, and they will be automatically added to your \texttt{cite.bib} file.

\section{Mathematical Expressions}
Here are some example equations:

\subsection{Quadratic Formula}
The quadratic formula is given by:
\[x = \frac{-b \pm \sqrt{b^2 - 4ac}}{2a}\]

\subsection{Matrix Example}
A sample matrix:
\[
\begin{pmatrix}
a_{11} & a_{12} & a_{13} \\
a_{21} & a_{22} & a_{23} \\
a_{31} & a_{32} & a_{33}
\end{pmatrix}
\]

\section{Tables}
Below is a sample table showing data:

\begin{table}[h]
\centering
\begin{tabular}{@{}lrrr@{}}
\toprule
\textbf{Category} & \textbf{Value 1} & \textbf{Value 2} & \textbf{Total} \\
\midrule
Item A & 10.5 & 20.3 & 30.8 \\
Item B & 15.2 & 25.1 & 40.3 \\
Item C & 22.1 & 12.7 & 34.8 \\
\bottomrule
\end{tabular}
\caption{Sample Data Table}
\label{tab:sample}
\end{table}

\section{Lists}
\subsection{Bullet Points}
\begin{itemize}
  \item First item
  \item Second item
  \item Third item with sub-items:
    \begin{itemize}
      \item Sub-item 1
      \item Sub-item 2
    \end{itemize}
\end{itemize}

\subsection{Numbered List}
\begin{enumerate}
  \item First step
  \item Second step
  \item Third step
\end{enumerate}

\section{Citations and References}
When you use the citation search feature, references will be automatically managed. For example, you might cite a paper like this: \cite{example2024}. All citations will be automatically added to your \texttt{cite.bib} file.

\section{Conclusion}
This template demonstrates basic \LaTeX{} features including mathematical formulas, tables, and lists. The bibliography system is now set up to work seamlessly with Bibby AI's citation management. Feel free to modify it for your needs!

\bibliographystyle{plain}
\bibliography{cite}

\end{document}\documentclass{article}
\usepackage[utf8]{inputenc}
\usepackage{booktabs}
\usepackage{amsmath}
\usepackage{natbib}

\title{My Document}
\author{Your Name}
\date{\today}

\begin{document}

\maketitle

\section{Introduction}
This is a sample document demonstrating various \LaTeX{} features including mathematical equations, tables, and formatting. You can cite papers using the citation search feature, and they will be automatically added to your \texttt{cite.bib} file.

\section{Mathematical Expressions}
Here are some example equations:

\subsection{Quadratic Formula}
The quadratic formula is given by:
\[x = \frac{-b \pm \sqrt{b^2 - 4ac}}{2a}\]

\subsection{Matrix Example}
A sample matrix:
\[
\begin{pmatrix}
a_{11} & a_{12} & a_{13} \\
a_{21} & a_{22} & a_{23} \\
a_{31} & a_{32} & a_{33}
\end{pmatrix}
\]

\section{Tables}
Below is a sample table showing data:

\begin{table}[h]
\centering
\begin{tabular}{@{}lrrr@{}}
\toprule
\textbf{Category} & \textbf{Value 1} & \textbf{Value 2} & \textbf{Total} \\
\midrule
Item A & 10.5 & 20.3 & 30.8 \\
Item B & 15.2 & 25.1 & 40.3 \\
Item C & 22.1 & 12.7 & 34.8 \\
\bottomrule
\end{tabular}
\caption{Sample Data Table}
\label{tab:sample}
\end{table}

\section{Lists}
\subsection{Bullet Points}
\begin{itemize}
  \item First item
  \item Second item
  \item Third item with sub-items:
    \begin{itemize}
      \item Sub-item 1
      \item Sub-item 2
    \end{itemize}
\end{itemize}

\subsection{Numbered List}
\begin{enumerate}
  \item First step
  \item Second step
  \item Third step
\end{enumerate}

\section{Citations and References}
When you use the citation search feature, references will be automatically managed. For example, you might cite a paper like this: \cite{example2024}. All citations will be automatically added to your \texttt{cite.bib} file.

\section{Conclusion}
This template demonstrates basic \LaTeX{} features including mathematical formulas, tables, and lists. The bibliography system is now set up to work seamlessly with Bibby AI's citation management. Feel free to modify it for your needs!

\bibliographystyle{plain}
\bibliography{cite}

\end{document}\documentclass{article}
\usepackage[utf8]{inputenc}
\usepackage{booktabs}
\usepackage{amsmath}
\usepackage{natbib}

\title{My Document}
\author{Your Name}
\date{\today}

\begin{document}

\maketitle

\section{Introduction}
This is a sample document demonstrating various \LaTeX{} features including mathematical equations, tables, and formatting. You can cite papers using the citation search feature, and they will be automatically added to your \texttt{cite.bib} file.

\section{Mathematical Expressions}
Here are some example equations:

\subsection{Quadratic Formula}
The quadratic formula is given by:
\[x = \frac{-b \pm \sqrt{b^2 - 4ac}}{2a}\]

\subsection{Matrix Example}
A sample matrix:
\[
\begin{pmatrix}
a_{11} & a_{12} & a_{13} \\
a_{21} & a_{22} & a_{23} \\
a_{31} & a_{32} & a_{33}
\end{pmatrix}
\]

\section{Tables}
Below is a sample table showing data:

\begin{table}[h]
\centering
\begin{tabular}{@{}lrrr@{}}
\toprule
\textbf{Category} & \textbf{Value 1} & \textbf{Value 2} & \textbf{Total} \\
\midrule
Item A & 10.5 & 20.3 & 30.8 \\
Item B & 15.2 & 25.1 & 40.3 \\
Item C & 22.1 & 12.7 & 34.8 \\
\bottomrule
\end{tabular}
\caption{Sample Data Table}
\label{tab:sample}
\end{table}

\section{Lists}
\subsection{Bullet Points}
\begin{itemize}
  \item First item
  \item Second item
  \item Third item with sub-items:
    \begin{itemize}
      \item Sub-item 1
      \item Sub-item 2
    \end{itemize}
\end{itemize}

\subsection{Numbered List}
\begin{enumerate}
  \item First step
  \item Second step
  \item Third step
\end{enumerate}

\section{Citations and References}
When you use the citation search feature, references will be automatically managed. For example, you might cite a paper like this: \cite{example2024}. All citations will be automatically added to your \texttt{cite.bib} file.

\section{Conclusion}
This template demonstrates basic \LaTeX{} features including mathematical formulas, tables, and lists. The bibliography system is now set up to work seamlessly with Bibby AI's citation management. Feel free to modify it for your needs!

\bibliographystyle{plain}
\bibliography{cite}

\end{document}\documentclass{article}
\usepackage[utf8]{inputenc}
\usepackage{booktabs}
\usepackage{amsmath}
\usepackage{natbib}

\title{My Document}
\author{Your Name}
\date{\today}

\begin{document}

\maketitle

\section{Introduction}
This is a sample document demonstrating various \LaTeX{} features including mathematical equations, tables, and formatting. You can cite papers using the citation search feature, and they will be automatically added to your \texttt{cite.bib} file.

\section{Mathematical Expressions}
Here are some example equations:

\subsection{Quadratic Formula}
The quadratic formula is given by:
\[x = \frac{-b \pm \sqrt{b^2 - 4ac}}{2a}\]

\subsection{Matrix Example}
A sample matrix:
\[
\begin{pmatrix}
a_{11} & a_{12} & a_{13} \\
a_{21} & a_{22} & a_{23} \\
a_{31} & a_{32} & a_{33}
\end{pmatrix}
\]

\section{Tables}
Below is a sample table showing data:

\begin{table}[h]
\centering
\begin{tabular}{@{}lrrr@{}}
\toprule
\textbf{Category} & \textbf{Value 1} & \textbf{Value 2} & \textbf{Total} \\
\midrule
Item A & 10.5 & 20.3 & 30.8 \\
Item B & 15.2 & 25.1 & 40.3 \\
Item C & 22.1 & 12.7 & 34.8 \\
\bottomrule
\end{tabular}
\caption{Sample Data Table}
\label{tab:sample}
\end{table}

\section{Lists}
\subsection{Bullet Points}
\begin{itemize}
  \item First item
  \item Second item
  \item Third item with sub-items:
    \begin{itemize}
      \item Sub-item 1
      \item Sub-item 2
    \end{itemize}
\end{itemize}

\subsection{Numbered List}
\begin{enumerate}
  \item First step
  \item Second step
  \item Third step
\end{enumerate}

\section{Citations and References}
When you use the citation search feature, references will be automatically managed. For example, you might cite a paper like this: \cite{example2024}. All citations will be automatically added to your \texttt{cite.bib} file.

\section{Conclusion}
This template demonstrates basic \LaTeX{} features including mathematical formulas, tables, and lists. The bibliography system is now set up to work seamlessly with Bibby AI's citation management. Feel free to modify it for your needs!

\bibliographystyle{plain}
\bibliography{cite}

\end{document}\documentclass{article}
\usepackage[utf8]{inputenc}
\usepackage{booktabs}
\usepackage{amsmath}
\usepackage{natbib}

\title{My Document}
\author{Your Name}
\date{\today}

\begin{document}

\maketitle

\section{Introduction}
This is a sample document demonstrating various \LaTeX{} features including mathematical equations, tables, and formatting. You can cite papers using the citation search feature, and they will be automatically added to your \texttt{cite.bib} file.

\section{Mathematical Expressions}
Here are some example equations:

\subsection{Quadratic Formula}
The quadratic formula is given by:
\[x = \frac{-b \pm \sqrt{b^2 - 4ac}}{2a}\]

\subsection{Matrix Example}
A sample matrix:
\[
\begin{pmatrix}
a_{11} & a_{12} & a_{13} \\
a_{21} & a_{22} & a_{23} \\
a_{31} & a_{32} & a_{33}
\end{pmatrix}
\]

\section{Tables}
Below is a sample table showing data:

\begin{table}[h]
\centering
\begin{tabular}{@{}lrrr@{}}
\toprule
\textbf{Category} & \textbf{Value 1} & \textbf{Value 2} & \textbf{Total} \\
\midrule
Item A & 10.5 & 20.3 & 30.8 \\
Item B & 15.2 & 25.1 & 40.3 \\
Item C & 22.1 & 12.7 & 34.8 \\
\bottomrule
\end{tabular}
\caption{Sample Data Table}
\label{tab:sample}
\end{table}

\section{Lists}
\subsection{Bullet Points}
\begin{itemize}
  \item First item
  \item Second item
  \item Third item with sub-items:
    \begin{itemize}
      \item Sub-item 1
      \item Sub-item 2
    \end{itemize}
\end{itemize}

\subsection{Numbered List}
\begin{enumerate}
  \item First step
  \item Second step
  \item Third step
\end{enumerate}

\section{Citations and References}
When you use the citation search feature, references will be automatically managed. For example, you might cite a paper like this: \cite{example2024}. All citations will be automatically added to your \texttt{cite.bib} file.

\section{Conclusion}
This template demonstrates basic \LaTeX{} features including mathematical formulas, tables, and lists. The bibliography system is now set up to work seamlessly with Bibby AI's citation management. Feel free to modify it for your needs!

\bibliographystyle{plain}
\bibliography{cite}

\end{document}\documentclass{article}
\usepackage[utf8]{inputenc}
\usepackage{booktabs}
\usepackage{amsmath}
\usepackage{natbib}

\title{My Document}
\author{Your Name}
\date{\today}

\begin{document}

\maketitle

\section{Introduction}
This is a sample document demonstrating various \LaTeX{} features including mathematical equations, tables, and formatting. You can cite papers using the citation search feature, and they will be automatically added to your \texttt{cite.bib} file.

\section{Mathematical Expressions}
Here are some example equations:

\subsection{Quadratic Formula}
The quadratic formula is given by:
\[x = \frac{-b \pm \sqrt{b^2 - 4ac}}{2a}\]

\subsection{Matrix Example}
A sample matrix:
\[
\begin{pmatrix}
a_{11} & a_{12} & a_{13} \\
a_{21} & a_{22} & a_{23} \\
a_{31} & a_{32} & a_{33}
\end{pmatrix}
\]

\section{Tables}
Below is a sample table showing data:

\begin{table}[h]
\centering
\begin{tabular}{@{}lrrr@{}}
\toprule
\textbf{Category} & \textbf{Value 1} & \textbf{Value 2} & \textbf{Total} \\
\midrule
Item A & 10.5 & 20.3 & 30.8 \\
Item B & 15.2 & 25.1 & 40.3 \\
Item C & 22.1 & 12.7 & 34.8 \\
\bottomrule
\end{tabular}
\caption{Sample Data Table}
\label{tab:sample}
\end{table}

\section{Lists}
\subsection{Bullet Points}
\begin{itemize}
  \item First item
  \item Second item
  \item Third item with sub-items:
    \begin{itemize}
      \item Sub-item 1
      \item Sub-item 2
    \end{itemize}
\end{itemize}

\subsection{Numbered List}
\begin{enumerate}
  \item First step
  \item Second step
  \item Third step
\end{enumerate}

\section{Citations and References}
When you use the citation search feature, references will be automatically managed. For example, you might cite a paper like this: \cite{example2024}. All citations will be automatically added to your \texttt{cite.bib} file.

\section{Conclusion}
This template demonstrates basic \LaTeX{} features including mathematical formulas, tables, and lists. The bibliography system is now set up to work seamlessly with Bibby AI's citation management. Feel free to modify it for your needs!

\bibliographystyle{plain}
\bibliography{cite}

\end{document}\documentclass{article}
\usepackage[utf8]{inputenc}
\usepackage{booktabs}
\usepackage{amsmath}
\usepackage{natbib}

\title{My Document}
\author{Your Name}
\date{\today}

\begin{document}

\maketitle

\section{Introduction}
This is a sample document demonstrating various \LaTeX{} features including mathematical equations, tables, and formatting. You can cite papers using the citation search feature, and they will be automatically added to your \texttt{cite.bib} file.

\section{Mathematical Expressions}
Here are some example equations:

\subsection{Quadratic Formula}
The quadratic formula is given by:
\[x = \frac{-b \pm \sqrt{b^2 - 4ac}}{2a}\]

\subsection{Matrix Example}
A sample matrix:
\[
\begin{pmatrix}
a_{11} & a_{12} & a_{13} \\
a_{21} & a_{22} & a_{23} \\
a_{31} & a_{32} & a_{33}
\end{pmatrix}
\]

\section{Tables}
Below is a sample table showing data:

\begin{table}[h]
\centering
\begin{tabular}{@{}lrrr@{}}
\toprule
\textbf{Category} & \textbf{Value 1} & \textbf{Value 2} & \textbf{Total} \\
\midrule
Item A & 10.5 & 20.3 & 30.8 \\
Item B & 15.2 & 25.1 & 40.3 \\
Item C & 22.1 & 12.7 & 34.8 \\
\bottomrule
\end{tabular}
\caption{Sample Data Table}
\label{tab:sample}
\end{table}

\section{Lists}
\subsection{Bullet Points}
\begin{itemize}
  \item First item
  \item Second item
  \item Third item with sub-items:
    \begin{itemize}
      \item Sub-item 1
      \item Sub-item 2
    \end{itemize}
\end{itemize}

\subsection{Numbered List}
\begin{enumerate}
  \item First step
  \item Second step
  \item Third step
\end{enumerate}

\section{Citations and References}
When you use the citation search feature, references will be automatically managed. For example, you might cite a paper like this: \cite{example2024}. All citations will be automatically added to your \texttt{cite.bib} file.

\section{Conclusion}
This template demonstrates basic \LaTeX{} features including mathematical formulas, tables, and lists. The bibliography system is now set up to work seamlessly with Bibby AI's citation management. Feel free to modify it for your needs!

\bibliographystyle{plain}
\bibliography{cite}

\end{document}\documentclass{article}
\usepackage[utf8]{inputenc}
\usepackage{booktabs}
\usepackage{amsmath}
\usepackage{natbib}

\title{My Document}
\author{Your Name}
\date{\today}

\begin{document}

\maketitle

\section{Introduction}
This is a sample document demonstrating various \LaTeX{} features including mathematical equations, tables, and formatting. You can cite papers using the citation search feature, and they will be automatically added to your \texttt{cite.bib} file.

\section{Mathematical Expressions}
Here are some example equations:

\subsection{Quadratic Formula}
The quadratic formula is given by:
\[x = \frac{-b \pm \sqrt{b^2 - 4ac}}{2a}\]

\subsection{Matrix Example}
A sample matrix:
\[
\begin{pmatrix}
a_{11} & a_{12} & a_{13} \\
a_{21} & a_{22} & a_{23} \\
a_{31} & a_{32} & a_{33}
\end{pmatrix}
\]

\section{Tables}
Below is a sample table showing data:

\begin{table}[h]
\centering
\begin{tabular}{@{}lrrr@{}}
\toprule
\textbf{Category} & \textbf{Value 1} & \textbf{Value 2} & \textbf{Total} \\
\midrule
Item A & 10.5 & 20.3 & 30.8 \\
Item B & 15.2 & 25.1 & 40.3 \\
Item C & 22.1 & 12.7 & 34.8 \\
\bottomrule
\end{tabular}
\caption{Sample Data Table}
\label{tab:sample}
\end{table}

\section{Lists}
\subsection{Bullet Points}
\begin{itemize}
  \item First item
  \item Second item
  \item Third item with sub-items:
    \begin{itemize}
      \item Sub-item 1
      \item Sub-item 2
    \end{itemize}
\end{itemize}

\subsection{Numbered List}
\begin{enumerate}
  \item First step
  \item Second step
  \item Third step
\end{enumerate}

\section{Citations and References}
When you use the citation search feature, references will be automatically managed. For example, you might cite a paper like this: \cite{example2024}. All citations will be automatically added to your \texttt{cite.bib} file.

\section{Conclusion}
This template demonstrates basic \LaTeX{} features including mathematical formulas, tables, and lists. The bibliography system is now set up to work seamlessly with Bibby AI's citation management. Feel free to modify it for your needs!

\bibliographystyle{plain}
\bibliography{cite}

\end{document}\documentclass{article}
\usepackage[utf8]{inputenc}
\usepackage{booktabs}
\usepackage{amsmath}
\usepackage{natbib}

\title{My Document}
\author{Your Name}
\date{\today}

\begin{document}

\maketitle

\section{Introduction}
This is a sample document demonstrating various \LaTeX{} features including mathematical equations, tables, and formatting. You can cite papers using the citation search feature, and they will be automatically added to your \texttt{cite.bib} file.

\section{Mathematical Expressions}
Here are some example equations:

\subsection{Quadratic Formula}
The quadratic formula is given by:
\[x = \frac{-b \pm \sqrt{b^2 - 4ac}}{2a}\]

\subsection{Matrix Example}
A sample matrix:
\[
\begin{pmatrix}
a_{11} & a_{12} & a_{13} \\
a_{21} & a_{22} & a_{23} \\
a_{31} & a_{32} & a_{33}
\end{pmatrix}
\]

\section{Tables}
Below is a sample table showing data:

\begin{table}[h]
\centering
\begin{tabular}{@{}lrrr@{}}
\toprule
\textbf{Category} & \textbf{Value 1} & \textbf{Value 2} & \textbf{Total} \\
\midrule
Item A & 10.5 & 20.3 & 30.8 \\
Item B & 15.2 & 25.1 & 40.3 \\
Item C & 22.1 & 12.7 & 34.8 \\
\bottomrule
\end{tabular}
\caption{Sample Data Table}
\label{tab:sample}
\end{table}

\section{Lists}
\subsection{Bullet Points}
\begin{itemize}
  \item First item
  \item Second item
  \item Third item with sub-items:
    \begin{itemize}
      \item Sub-item 1
      \item Sub-item 2
    \end{itemize}
\end{itemize}

\subsection{Numbered List}
\begin{enumerate}
  \item First step
  \item Second step
  \item Third step
\end{enumerate}

\section{Citations and References}
When you use the citation search feature, references will be automatically managed. For example, you might cite a paper like this: \cite{example2024}. All citations will be automatically added to your \texttt{cite.bib} file.

\section{Conclusion}
This template demonstrates basic \LaTeX{} features including mathematical formulas, tables, and lists. The bibliography system is now set up to work seamlessly with Bibby AI's citation management. Feel free to modify it for your needs!

\bibliographystyle{plain}
\bibliography{cite}

\end{document}\documentclass{article}
\usepackage[utf8]{inputenc}
\usepackage{booktabs}
\usepackage{amsmath}
\usepackage{natbib}

\title{My Document}
\author{Your Name}
\date{\today}

\begin{document}

\maketitle

\section{Introduction}
This is a sample document demonstrating various \LaTeX{} features including mathematical equations, tables, and formatting. You can cite papers using the citation search feature, and they will be automatically added to your \texttt{cite.bib} file.

\section{Mathematical Expressions}
Here are some example equations:

\subsection{Quadratic Formula}
The quadratic formula is given by:
\[x = \frac{-b \pm \sqrt{b^2 - 4ac}}{2a}\]

\subsection{Matrix Example}
A sample matrix:
\[
\begin{pmatrix}
a_{11} & a_{12} & a_{13} \\
a_{21} & a_{22} & a_{23} \\
a_{31} & a_{32} & a_{33}
\end{pmatrix}
\]

\section{Tables}
Below is a sample table showing data:

\begin{table}[h]
\centering
\begin{tabular}{@{}lrrr@{}}
\toprule
\textbf{Category} & \textbf{Value 1} & \textbf{Value 2} & \textbf{Total} \\
\midrule
Item A & 10.5 & 20.3 & 30.8 \\
Item B & 15.2 & 25.1 & 40.3 \\
Item C & 22.1 & 12.7 & 34.8 \\
\bottomrule
\end{tabular}
\caption{Sample Data Table}
\label{tab:sample}
\end{table}

\section{Lists}
\subsection{Bullet Points}
\begin{itemize}
  \item First item
  \item Second item
  \item Third item with sub-items:
    \begin{itemize}
      \item Sub-item 1
      \item Sub-item 2
    \end{itemize}
\end{itemize}

\subsection{Numbered List}
\begin{enumerate}
  \item First step
  \item Second step
  \item Third step
\end{enumerate}

\section{Citations and References}
When you use the citation search feature, references will be automatically managed. For example, you might cite a paper like this: \cite{example2024}. All citations will be automatically added to your \texttt{cite.bib} file.

\section{Conclusion}
This template demonstrates basic \LaTeX{} features including mathematical formulas, tables, and lists. The bibliography system is now set up to work seamlessly with Bibby AI's citation management. Feel free to modify it for your needs!

\bibliographystyle{plain}
\bibliography{cite}

\end{document}\documentclass{article}
\usepackage[utf8]{inputenc}
\usepackage{booktabs}
\usepackage{amsmath}
\usepackage{natbib}

\title{My Document}
\author{Your Name}
\date{\today}

\begin{document}

\maketitle

\section{Introduction}
This is a sample document demonstrating various \LaTeX{} features including mathematical equations, tables, and formatting. You can cite papers using the citation search feature, and they will be automatically added to your \texttt{cite.bib} file.

\section{Mathematical Expressions}
Here are some example equations:

\subsection{Quadratic Formula}
The quadratic formula is given by:
\[x = \frac{-b \pm \sqrt{b^2 - 4ac}}{2a}\]

\subsection{Matrix Example}
A sample matrix:
\[
\begin{pmatrix}
a_{11} & a_{12} & a_{13} \\
a_{21} & a_{22} & a_{23} \\
a_{31} & a_{32} & a_{33}
\end{pmatrix}
\]

\section{Tables}
Below is a sample table showing data:

\begin{table}[h]
\centering
\begin{tabular}{@{}lrrr@{}}
\toprule
\textbf{Category} & \textbf{Value 1} & \textbf{Value 2} & \textbf{Total} \\
\midrule
Item A & 10.5 & 20.3 & 30.8 \\
Item B & 15.2 & 25.1 & 40.3 \\
Item C & 22.1 & 12.7 & 34.8 \\
\bottomrule
\end{tabular}
\caption{Sample Data Table}
\label{tab:sample}
\end{table}

\section{Lists}
\subsection{Bullet Points}
\begin{itemize}
  \item First item
  \item Second item
  \item Third item with sub-items:
    \begin{itemize}
      \item Sub-item 1
      \item Sub-item 2
    \end{itemize}
\end{itemize}

\subsection{Numbered List}
\begin{enumerate}
  \item First step
  \item Second step
  \item Third step
\end{enumerate}

\section{Citations and References}
When you use the citation search feature, references will be automatically managed. For example, you might cite a paper like this: \cite{example2024}. All citations will be automatically added to your \texttt{cite.bib} file.

\section{Conclusion}
This template demonstrates basic \LaTeX{} features including mathematical formulas, tables, and lists. The bibliography system is now set up to work seamlessly with Bibby AI's citation management. Feel free to modify it for your needs!

\bibliographystyle{plain}
\bibliography{cite}

\end{document}\documentclass{article}
\usepackage[utf8]{inputenc}
\usepackage{booktabs}
\usepackage{amsmath}
\usepackage{natbib}

\title{My Document}
\author{Your Name}
\date{\today}

\begin{document}

\maketitle

\section{Introduction}
This is a sample document demonstrating various \LaTeX{} features including mathematical equations, tables, and formatting. You can cite papers using the citation search feature, and they will be automatically added to your \texttt{cite.bib} file.

\section{Mathematical Expressions}
Here are some example equations:

\subsection{Quadratic Formula}
The quadratic formula is given by:
\[x = \frac{-b \pm \sqrt{b^2 - 4ac}}{2a}\]

\subsection{Matrix Example}
A sample matrix:
\[
\begin{pmatrix}
a_{11} & a_{12} & a_{13} \\
a_{21} & a_{22} & a_{23} \\
a_{31} & a_{32} & a_{33}
\end{pmatrix}
\]

\section{Tables}
Below is a sample table showing data:

\begin{table}[h]
\centering
\begin{tabular}{@{}lrrr@{}}
\toprule
\textbf{Category} & \textbf{Value 1} & \textbf{Value 2} & \textbf{Total} \\
\midrule
Item A & 10.5 & 20.3 & 30.8 \\
Item B & 15.2 & 25.1 & 40.3 \\
Item C & 22.1 & 12.7 & 34.8 \\
\bottomrule
\end{tabular}
\caption{Sample Data Table}
\label{tab:sample}
\end{table}

\section{Lists}
\subsection{Bullet Points}
\begin{itemize}
  \item First item
  \item Second item
  \item Third item with sub-items:
    \begin{itemize}
      \item Sub-item 1
      \item Sub-item 2
    \end{itemize}
\end{itemize}

\subsection{Numbered List}
\begin{enumerate}
  \item First step
  \item Second step
  \item Third step
\end{enumerate}

\section{Citations and References}
When you use the citation search feature, references will be automatically managed. For example, you might cite a paper like this: \cite{example2024}. All citations will be automatically added to your \texttt{cite.bib} file.

\section{Conclusion}
This template demonstrates basic \LaTeX{} features including mathematical formulas, tables, and lists. The bibliography system is now set up to work seamlessly with Bibby AI's citation management. Feel free to modify it for your needs!

\bibliographystyle{plain}
\bibliography{cite}

\end{document}\documentclass{article}
\usepackage[utf8]{inputenc}
\usepackage{booktabs}
\usepackage{amsmath}
\usepackage{natbib}

\title{My Document}
\author{Your Name}
\date{\today}

\begin{document}

\maketitle

\section{Introduction}
This is a sample document demonstrating various \LaTeX{} features including mathematical equations, tables, and formatting. You can cite papers using the citation search feature, and they will be automatically added to your \texttt{cite.bib} file.

\section{Mathematical Expressions}
Here are some example equations:

\subsection{Quadratic Formula}
The quadratic formula is given by:
\[x = \frac{-b \pm \sqrt{b^2 - 4ac}}{2a}\]

\subsection{Matrix Example}
A sample matrix:
\[
\begin{pmatrix}
a_{11} & a_{12} & a_{13} \\
a_{21} & a_{22} & a_{23} \\
a_{31} & a_{32} & a_{33}
\end{pmatrix}
\]

\section{Tables}
Below is a sample table showing data:

\begin{table}[h]
\centering
\begin{tabular}{@{}lrrr@{}}
\toprule
\textbf{Category} & \textbf{Value 1} & \textbf{Value 2} & \textbf{Total} \\
\midrule
Item A & 10.5 & 20.3 & 30.8 \\
Item B & 15.2 & 25.1 & 40.3 \\
Item C & 22.1 & 12.7 & 34.8 \\
\bottomrule
\end{tabular}
\caption{Sample Data Table}
\label{tab:sample}
\end{table}

\section{Lists}
\subsection{Bullet Points}
\begin{itemize}
  \item First item
  \item Second item
  \item Third item with sub-items:
    \begin{itemize}
      \item Sub-item 1
      \item Sub-item 2
    \end{itemize}
\end{itemize}

\subsection{Numbered List}
\begin{enumerate}
  \item First step
  \item Second step
  \item Third step
\end{enumerate}

\section{Citations and References}
When you use the citation search feature, references will be automatically managed. For example, you might cite a paper like this: \cite{example2024}. All citations will be automatically added to your \texttt{cite.bib} file.

\section{Conclusion}
This template demonstrates basic \LaTeX{} features including mathematical formulas, tables, and lists. The bibliography system is now set up to work seamlessly with Bibby AI's citation management. Feel free to modify it for your needs!

\bibliographystyle{plain}
\bibliography{cite}

\end{document}\documentclass{article}
\usepackage[utf8]{inputenc}
\usepackage{booktabs}
\usepackage{amsmath}
\usepackage{natbib}

\title{My Document}
\author{Your Name}
\date{\today}

\begin{document}

\maketitle

\section{Introduction}
This is a sample document demonstrating various \LaTeX{} features including mathematical equations, tables, and formatting. You can cite papers using the citation search feature, and they will be automatically added to your \texttt{cite.bib} file.

\section{Mathematical Expressions}
Here are some example equations:

\subsection{Quadratic Formula}
The quadratic formula is given by:
\[x = \frac{-b \pm \sqrt{b^2 - 4ac}}{2a}\]

\subsection{Matrix Example}
A sample matrix:
\[
\begin{pmatrix}
a_{11} & a_{12} & a_{13} \\
a_{21} & a_{22} & a_{23} \\
a_{31} & a_{32} & a_{33}
\end{pmatrix}
\]

\section{Tables}
Below is a sample table showing data:

\begin{table}[h]
\centering
\begin{tabular}{@{}lrrr@{}}
\toprule
\textbf{Category} & \textbf{Value 1} & \textbf{Value 2} & \textbf{Total} \\
\midrule
Item A & 10.5 & 20.3 & 30.8 \\
Item B & 15.2 & 25.1 & 40.3 \\
Item C & 22.1 & 12.7 & 34.8 \\
\bottomrule
\end{tabular}
\caption{Sample Data Table}
\label{tab:sample}
\end{table}

\section{Lists}
\subsection{Bullet Points}
\begin{itemize}
  \item First item
  \item Second item
  \item Third item with sub-items:
    \begin{itemize}
      \item Sub-item 1
      \item Sub-item 2
    \end{itemize}
\end{itemize}

\subsection{Numbered List}
\begin{enumerate}
  \item First step
  \item Second step
  \item Third step
\end{enumerate}

\section{Citations and References}
When you use the citation search feature, references will be automatically managed. For example, you might cite a paper like this: \cite{example2024}. All citations will be automatically added to your \texttt{cite.bib} file.

\section{Conclusion}
This template demonstrates basic \LaTeX{} features including mathematical formulas, tables, and lists. The bibliography system is now set up to work seamlessly with Bibby AI's citation management. Feel free to modify it for your needs!

\bibliographystyle{plain}
\bibliography{cite}

\end{document}\documentclass{article}
\usepackage[utf8]{inputenc}
\usepackage{booktabs}
\usepackage{amsmath}
\usepackage{natbib}

\title{My Document}
\author{Your Name}
\date{\today}

\begin{document}

\maketitle

\section{Introduction}
This is a sample document demonstrating various \LaTeX{} features including mathematical equations, tables, and formatting. You can cite papers using the citation search feature, and they will be automatically added to your \texttt{cite.bib} file.

\section{Mathematical Expressions}
Here are some example equations:

\subsection{Quadratic Formula}
The quadratic formula is given by:
\[x = \frac{-b \pm \sqrt{b^2 - 4ac}}{2a}\]

\subsection{Matrix Example}
A sample matrix:
\[
\begin{pmatrix}
a_{11} & a_{12} & a_{13} \\
a_{21} & a_{22} & a_{23} \\
a_{31} & a_{32} & a_{33}
\end{pmatrix}
\]

\section{Tables}
Below is a sample table showing data:

\begin{table}[h]
\centering
\begin{tabular}{@{}lrrr@{}}
\toprule
\textbf{Category} & \textbf{Value 1} & \textbf{Value 2} & \textbf{Total} \\
\midrule
Item A & 10.5 & 20.3 & 30.8 \\
Item B & 15.2 & 25.1 & 40.3 \\
Item C & 22.1 & 12.7 & 34.8 \\
\bottomrule
\end{tabular}
\caption{Sample Data Table}
\label{tab:sample}
\end{table}

\section{Lists}
\subsection{Bullet Points}
\begin{itemize}
  \item First item
  \item Second item
  \item Third item with sub-items:
    \begin{itemize}
      \item Sub-item 1
      \item Sub-item 2
    \end{itemize}
\end{itemize}

\subsection{Numbered List}
\begin{enumerate}
  \item First step
  \item Second step
  \item Third step
\end{enumerate}

\section{Citations and References}
When you use the citation search feature, references will be automatically managed. For example, you might cite a paper like this: \cite{example2024}. All citations will be automatically added to your \texttt{cite.bib} file.

\section{Conclusion}
This template demonstrates basic \LaTeX{} features including mathematical formulas, tables, and lists. The bibliography system is now set up to work seamlessly with Bibby AI's citation management. Feel free to modify it for your needs!

\bibliographystyle{plain}
\bibliography{cite}

\end{document}\documentclass{article}
\usepackage[utf8]{inputenc}
\usepackage{booktabs}
\usepackage{amsmath}
\usepackage{natbib}

\title{My Document}
\author{Your Name}
\date{\today}

\begin{document}

\maketitle

\section{Introduction}
This is a sample document demonstrating various \LaTeX{} features including mathematical equations, tables, and formatting. You can cite papers using the citation search feature, and they will be automatically added to your \texttt{cite.bib} file.

\section{Mathematical Expressions}
Here are some example equations:

\subsection{Quadratic Formula}
The quadratic formula is given by:
\[x = \frac{-b \pm \sqrt{b^2 - 4ac}}{2a}\]

\subsection{Matrix Example}
A sample matrix:
\[
\begin{pmatrix}
a_{11} & a_{12} & a_{13} \\
a_{21} & a_{22} & a_{23} \\
a_{31} & a_{32} & a_{33}
\end{pmatrix}
\]

\section{Tables}
Below is a sample table showing data:

\begin{table}[h]
\centering
\begin{tabular}{@{}lrrr@{}}
\toprule
\textbf{Category} & \textbf{Value 1} & \textbf{Value 2} & \textbf{Total} \\
\midrule
Item A & 10.5 & 20.3 & 30.8 \\
Item B & 15.2 & 25.1 & 40.3 \\
Item C & 22.1 & 12.7 & 34.8 \\
\bottomrule
\end{tabular}
\caption{Sample Data Table}
\label{tab:sample}
\end{table}

\section{Lists}
\subsection{Bullet Points}
\begin{itemize}
  \item First item
  \item Second item
  \item Third item with sub-items:
    \begin{itemize}
      \item Sub-item 1
      \item Sub-item 2
    \end{itemize}
\end{itemize}

\subsection{Numbered List}
\begin{enumerate}
  \item First step
  \item Second step
  \item Third step
\end{enumerate}

\section{Citations and References}
When you use the citation search feature, references will be automatically managed. For example, you might cite a paper like this: \cite{example2024}. All citations will be automatically added to your \texttt{cite.bib} file.

\section{Conclusion}
This template demonstrates basic \LaTeX{} features including mathematical formulas, tables, and lists. The bibliography system is now set up to work seamlessly with Bibby AI's citation management. Feel free to modify it for your needs!

\bibliographystyle{plain}
\bibliography{cite}

\end{document}\documentclass{article}
\usepackage[utf8]{inputenc}
\usepackage{booktabs}
\usepackage{amsmath}
\usepackage{natbib}

\title{My Document}
\author{Your Name}
\date{\today}

\begin{document}

\maketitle

\section{Introduction}
This is a sample document demonstrating various \LaTeX{} features including mathematical equations, tables, and formatting. You can cite papers using the citation search feature, and they will be automatically added to your \texttt{cite.bib} file.

\section{Mathematical Expressions}
Here are some example equations:

\subsection{Quadratic Formula}
The quadratic formula is given by:
\[x = \frac{-b \pm \sqrt{b^2 - 4ac}}{2a}\]

\subsection{Matrix Example}
A sample matrix:
\[
\begin{pmatrix}
a_{11} & a_{12} & a_{13} \\
a_{21} & a_{22} & a_{23} \\
a_{31} & a_{32} & a_{33}
\end{pmatrix}
\]

\section{Tables}
Below is a sample table showing data:

\begin{table}[h]
\centering
\begin{tabular}{@{}lrrr@{}}
\toprule
\textbf{Category} & \textbf{Value 1} & \textbf{Value 2} & \textbf{Total} \\
\midrule
Item A & 10.5 & 20.3 & 30.8 \\
Item B & 15.2 & 25.1 & 40.3 \\
Item C & 22.1 & 12.7 & 34.8 \\
\bottomrule
\end{tabular}
\caption{Sample Data Table}
\label{tab:sample}
\end{table}

\section{Lists}
\subsection{Bullet Points}
\begin{itemize}
  \item First item
  \item Second item
  \item Third item with sub-items:
    \begin{itemize}
      \item Sub-item 1
      \item Sub-item 2
    \end{itemize}
\end{itemize}

\subsection{Numbered List}
\begin{enumerate}
  \item First step
  \item Second step
  \item Third step
\end{enumerate}

\section{Citations and References}
When you use the citation search feature, references will be automatically managed. For example, you might cite a paper like this: \cite{example2024}. All citations will be automatically added to your \texttt{cite.bib} file.

\section{Conclusion}
This template demonstrates basic \LaTeX{} features including mathematical formulas, tables, and lists. The bibliography system is now set up to work seamlessly with Bibby AI's citation management. Feel free to modify it for your needs!

\bibliographystyle{plain}
\bibliography{cite}

\end{document}\documentclass{article}
\usepackage[utf8]{inputenc}
\usepackage{booktabs}
\usepackage{amsmath}
\usepackage{natbib}

\title{My Document}
\author{Your Name}
\date{\today}

\begin{document}

\maketitle

\section{Introduction}
This is a sample document demonstrating various \LaTeX{} features including mathematical equations, tables, and formatting. You can cite papers using the citation search feature, and they will be automatically added to your \texttt{cite.bib} file.

\section{Mathematical Expressions}
Here are some example equations:

\subsection{Quadratic Formula}
The quadratic formula is given by:
\[x = \frac{-b \pm \sqrt{b^2 - 4ac}}{2a}\]

\subsection{Matrix Example}
A sample matrix:
\[
\begin{pmatrix}
a_{11} & a_{12} & a_{13} \\
a_{21} & a_{22} & a_{23} \\
a_{31} & a_{32} & a_{33}
\end{pmatrix}
\]

\section{Tables}
Below is a sample table showing data:

\begin{table}[h]
\centering
\begin{tabular}{@{}lrrr@{}}
\toprule
\textbf{Category} & \textbf{Value 1} & \textbf{Value 2} & \textbf{Total} \\
\midrule
Item A & 10.5 & 20.3 & 30.8 \\
Item B & 15.2 & 25.1 & 40.3 \\
Item C & 22.1 & 12.7 & 34.8 \\
\bottomrule
\end{tabular}
\caption{Sample Data Table}
\label{tab:sample}
\end{table}

\section{Lists}
\subsection{Bullet Points}
\begin{itemize}
  \item First item
  \item Second item
  \item Third item with sub-items:
    \begin{itemize}
      \item Sub-item 1
      \item Sub-item 2
    \end{itemize}
\end{itemize}

\subsection{Numbered List}
\begin{enumerate}
  \item First step
  \item Second step
  \item Third step
\end{enumerate}

\section{Citations and References}
When you use the citation search feature, references will be automatically managed. For example, you might cite a paper like this: \cite{example2024}. All citations will be automatically added to your \texttt{cite.bib} file.

\section{Conclusion}
This template demonstrates basic \LaTeX{} features including mathematical formulas, tables, and lists. The bibliography system is now set up to work seamlessly with Bibby AI's citation management. Feel free to modify it for your needs!

\bibliographystyle{plain}
\bibliography{cite}

\end{document}\documentclass{article}
\usepackage[utf8]{inputenc}
\usepackage{booktabs}
\usepackage{amsmath}
\usepackage{natbib}

\title{My Document}
\author{Your Name}
\date{\today}

\begin{document}

\maketitle

\section{Introduction}
This is a sample document demonstrating various \LaTeX{} features including mathematical equations, tables, and formatting. You can cite papers using the citation search feature, and they will be automatically added to your \texttt{cite.bib} file.

\section{Mathematical Expressions}
Here are some example equations:

\subsection{Quadratic Formula}
The quadratic formula is given by:
\[x = \frac{-b \pm \sqrt{b^2 - 4ac}}{2a}\]

\subsection{Matrix Example}
A sample matrix:
\[
\begin{pmatrix}
a_{11} & a_{12} & a_{13} \\
a_{21} & a_{22} & a_{23} \\
a_{31} & a_{32} & a_{33}
\end{pmatrix}
\]

\section{Tables}
Below is a sample table showing data:

\begin{table}[h]
\centering
\begin{tabular}{@{}lrrr@{}}
\toprule
\textbf{Category} & \textbf{Value 1} & \textbf{Value 2} & \textbf{Total} \\
\midrule
Item A & 10.5 & 20.3 & 30.8 \\
Item B & 15.2 & 25.1 & 40.3 \\
Item C & 22.1 & 12.7 & 34.8 \\
\bottomrule
\end{tabular}
\caption{Sample Data Table}
\label{tab:sample}
\end{table}

\section{Lists}
\subsection{Bullet Points}
\begin{itemize}
  \item First item
  \item Second item
  \item Third item with sub-items:
    \begin{itemize}
      \item Sub-item 1
      \item Sub-item 2
    \end{itemize}
\end{itemize}

\subsection{Numbered List}
\begin{enumerate}
  \item First step
  \item Second step
  \item Third step
\end{enumerate}

\section{Citations and References}
When you use the citation search feature, references will be automatically managed. For example, you might cite a paper like this: \cite{example2024}. All citations will be automatically added to your \texttt{cite.bib} file.

\section{Conclusion}
This template demonstrates basic \LaTeX{} features including mathematical formulas, tables, and lists. The bibliography system is now set up to work seamlessly with Bibby AI's citation management. Feel free to modify it for your needs!

\bibliographystyle{plain}
\bibliography{cite}

\end{document}\documentclass{article}
\usepackage[utf8]{inputenc}
\usepackage{booktabs}
\usepackage{amsmath}
\usepackage{natbib}

\title{My Document}
\author{Your Name}
\date{\today}

\begin{document}

\maketitle

\section{Introduction}
This is a sample document demonstrating various \LaTeX{} features including mathematical equations, tables, and formatting. You can cite papers using the citation search feature, and they will be automatically added to your \texttt{cite.bib} file.

\section{Mathematical Expressions}
Here are some example equations:

\subsection{Quadratic Formula}
The quadratic formula is given by:
\[x = \frac{-b \pm \sqrt{b^2 - 4ac}}{2a}\]

\subsection{Matrix Example}
A sample matrix:
\[
\begin{pmatrix}
a_{11} & a_{12} & a_{13} \\
a_{21} & a_{22} & a_{23} \\
a_{31} & a_{32} & a_{33}
\end{pmatrix}
\]

\section{Tables}
Below is a sample table showing data:

\begin{table}[h]
\centering
\begin{tabular}{@{}lrrr@{}}
\toprule
\textbf{Category} & \textbf{Value 1} & \textbf{Value 2} & \textbf{Total} \\
\midrule
Item A & 10.5 & 20.3 & 30.8 \\
Item B & 15.2 & 25.1 & 40.3 \\
Item C & 22.1 & 12.7 & 34.8 \\
\bottomrule
\end{tabular}
\caption{Sample Data Table}
\label{tab:sample}
\end{table}

\section{Lists}
\subsection{Bullet Points}
\begin{itemize}
  \item First item
  \item Second item
  \item Third item with sub-items:
    \begin{itemize}
      \item Sub-item 1
      \item Sub-item 2
    \end{itemize}
\end{itemize}

\subsection{Numbered List}
\begin{enumerate}
  \item First step
  \item Second step
  \item Third step
\end{enumerate}

\section{Citations and References}
When you use the citation search feature, references will be automatically managed. For example, you might cite a paper like this: \cite{example2024}. All citations will be automatically added to your \texttt{cite.bib} file.

\section{Conclusion}
This template demonstrates basic \LaTeX{} features including mathematical formulas, tables, and lists. The bibliography system is now set up to work seamlessly with Bibby AI's citation management. Feel free to modify it for your needs!

\bibliographystyle{plain}
\bibliography{cite}

\end{document}\documentclass{article}
\usepackage[utf8]{inputenc}
\usepackage{booktabs}
\usepackage{amsmath}
\usepackage{natbib}

\title{My Document}
\author{Your Name}
\date{\today}

\begin{document}

\maketitle

\section{Introduction}
This is a sample document demonstrating various \LaTeX{} features including mathematical equations, tables, and formatting. You can cite papers using the citation search feature, and they will be automatically added to your \texttt{cite.bib} file.

\section{Mathematical Expressions}
Here are some example equations:

\subsection{Quadratic Formula}
The quadratic formula is given by:
\[x = \frac{-b \pm \sqrt{b^2 - 4ac}}{2a}\]

\subsection{Matrix Example}
A sample matrix:
\[
\begin{pmatrix}
a_{11} & a_{12} & a_{13} \\
a_{21} & a_{22} & a_{23} \\
a_{31} & a_{32} & a_{33}
\end{pmatrix}
\]

\section{Tables}
Below is a sample table showing data:

\begin{table}[h]
\centering
\begin{tabular}{@{}lrrr@{}}
\toprule
\textbf{Category} & \textbf{Value 1} & \textbf{Value 2} & \textbf{Total} \\
\midrule
Item A & 10.5 & 20.3 & 30.8 \\
Item B & 15.2 & 25.1 & 40.3 \\
Item C & 22.1 & 12.7 & 34.8 \\
\bottomrule
\end{tabular}
\caption{Sample Data Table}
\label{tab:sample}
\end{table}

\section{Lists}
\subsection{Bullet Points}
\begin{itemize}
  \item First item
  \item Second item
  \item Third item with sub-items:
    \begin{itemize}
      \item Sub-item 1
      \item Sub-item 2
    \end{itemize}
\end{itemize}

\subsection{Numbered List}
\begin{enumerate}
  \item First step
  \item Second step
  \item Third step
\end{enumerate}

\section{Citations and References}
When you use the citation search feature, references will be automatically managed. For example, you might cite a paper like this: \cite{example2024}. All citations will be automatically added to your \texttt{cite.bib} file.

\section{Conclusion}
This template demonstrates basic \LaTeX{} features including mathematical formulas, tables, and lists. The bibliography system is now set up to work seamlessly with Bibby AI's citation management. Feel free to modify it for your needs!

\bibliographystyle{plain}
\bibliography{cite}

\end{document}\documentclass{article}
\usepackage[utf8]{inputenc}
\usepackage{booktabs}
\usepackage{amsmath}
\usepackage{natbib}

\title{My Document}
\author{Your Name}
\date{\today}

\begin{document}

\maketitle

\section{Introduction}
This is a sample document demonstrating various \LaTeX{} features including mathematical equations, tables, and formatting. You can cite papers using the citation search feature, and they will be automatically added to your \texttt{cite.bib} file.

\section{Mathematical Expressions}
Here are some example equations:

\subsection{Quadratic Formula}
The quadratic formula is given by:
\[x = \frac{-b \pm \sqrt{b^2 - 4ac}}{2a}\]

\subsection{Matrix Example}
A sample matrix:
\[
\begin{pmatrix}
a_{11} & a_{12} & a_{13} \\
a_{21} & a_{22} & a_{23} \\
a_{31} & a_{32} & a_{33}
\end{pmatrix}
\]

\section{Tables}
Below is a sample table showing data:

\begin{table}[h]
\centering
\begin{tabular}{@{}lrrr@{}}
\toprule
\textbf{Category} & \textbf{Value 1} & \textbf{Value 2} & \textbf{Total} \\
\midrule
Item A & 10.5 & 20.3 & 30.8 \\
Item B & 15.2 & 25.1 & 40.3 \\
Item C & 22.1 & 12.7 & 34.8 \\
\bottomrule
\end{tabular}
\caption{Sample Data Table}
\label{tab:sample}
\end{table}

\section{Lists}
\subsection{Bullet Points}
\begin{itemize}
  \item First item
  \item Second item
  \item Third item with sub-items:
    \begin{itemize}
      \item Sub-item 1
      \item Sub-item 2
    \end{itemize}
\end{itemize}

\subsection{Numbered List}
\begin{enumerate}
  \item First step
  \item Second step
  \item Third step
\end{enumerate}

\section{Citations and References}
When you use the citation search feature, references will be automatically managed. For example, you might cite a paper like this: \cite{example2024}. All citations will be automatically added to your \texttt{cite.bib} file.

\section{Conclusion}
This template demonstrates basic \LaTeX{} features including mathematical formulas, tables, and lists. The bibliography system is now set up to work seamlessly with Bibby AI's citation management. Feel free to modify it for your needs!

\bibliographystyle{plain}
\bibliography{cite}

\end{document}\documentclass{article}
\usepackage[utf8]{inputenc}
\usepackage{booktabs}
\usepackage{amsmath}
\usepackage{natbib}

\title{My Document}
\author{Your Name}
\date{\today}

\begin{document}

\maketitle

\section{Introduction}
This is a sample document demonstrating various \LaTeX{} features including mathematical equations, tables, and formatting. You can cite papers using the citation search feature, and they will be automatically added to your \texttt{cite.bib} file.

\section{Mathematical Expressions}
Here are some example equations:

\subsection{Quadratic Formula}
The quadratic formula is given by:
\[x = \frac{-b \pm \sqrt{b^2 - 4ac}}{2a}\]

\subsection{Matrix Example}
A sample matrix:
\[
\begin{pmatrix}
a_{11} & a_{12} & a_{13} \\
a_{21} & a_{22} & a_{23} \\
a_{31} & a_{32} & a_{33}
\end{pmatrix}
\]

\section{Tables}
Below is a sample table showing data:

\begin{table}[h]
\centering
\begin{tabular}{@{}lrrr@{}}
\toprule
\textbf{Category} & \textbf{Value 1} & \textbf{Value 2} & \textbf{Total} \\
\midrule
Item A & 10.5 & 20.3 & 30.8 \\
Item B & 15.2 & 25.1 & 40.3 \\
Item C & 22.1 & 12.7 & 34.8 \\
\bottomrule
\end{tabular}
\caption{Sample Data Table}
\label{tab:sample}
\end{table}

\section{Lists}
\subsection{Bullet Points}
\begin{itemize}
  \item First item
  \item Second item
  \item Third item with sub-items:
    \begin{itemize}
      \item Sub-item 1
      \item Sub-item 2
    \end{itemize}
\end{itemize}

\subsection{Numbered List}
\begin{enumerate}
  \item First step
  \item Second step
  \item Third step
\end{enumerate}

\section{Citations and References}
When you use the citation search feature, references will be automatically managed. For example, you might cite a paper like this: \cite{example2024}. All citations will be automatically added to your \texttt{cite.bib} file.

\section{Conclusion}
This template demonstrates basic \LaTeX{} features including mathematical formulas, tables, and lists. The bibliography system is now set up to work seamlessly with Bibby AI's citation management. Feel free to modify it for your needs!

\bibliographystyle{plain}
\bibliography{cite}

\end{document}\documentclass{article}
\usepackage[utf8]{inputenc}
\usepackage{booktabs}
\usepackage{amsmath}
\usepackage{natbib}

\title{My Document}
\author{Your Name}
\date{\today}

\begin{document}

\maketitle

\section{Introduction}
This is a sample document demonstrating various \LaTeX{} features including mathematical equations, tables, and formatting. You can cite papers using the citation search feature, and they will be automatically added to your \texttt{cite.bib} file.

\section{Mathematical Expressions}
Here are some example equations:

\subsection{Quadratic Formula}
The quadratic formula is given by:
\[x = \frac{-b \pm \sqrt{b^2 - 4ac}}{2a}\]

\subsection{Matrix Example}
A sample matrix:
\[
\begin{pmatrix}
a_{11} & a_{12} & a_{13} \\
a_{21} & a_{22} & a_{23} \\
a_{31} & a_{32} & a_{33}
\end{pmatrix}
\]

\section{Tables}
Below is a sample table showing data:

\begin{table}[h]
\centering
\begin{tabular}{@{}lrrr@{}}
\toprule
\textbf{Category} & \textbf{Value 1} & \textbf{Value 2} & \textbf{Total} \\
\midrule
Item A & 10.5 & 20.3 & 30.8 \\
Item B & 15.2 & 25.1 & 40.3 \\
Item C & 22.1 & 12.7 & 34.8 \\
\bottomrule
\end{tabular}
\caption{Sample Data Table}
\label{tab:sample}
\end{table}

\section{Lists}
\subsection{Bullet Points}
\begin{itemize}
  \item First item
  \item Second item
  \item Third item with sub-items:
    \begin{itemize}
      \item Sub-item 1
      \item Sub-item 2
    \end{itemize}
\end{itemize}

\subsection{Numbered List}
\begin{enumerate}
  \item First step
  \item Second step
  \item Third step
\end{enumerate}

\section{Citations and References}
When you use the citation search feature, references will be automatically managed. For example, you might cite a paper like this: \cite{example2024}. All citations will be automatically added to your \texttt{cite.bib} file.

\section{Conclusion}
This template demonstrates basic \LaTeX{} features including mathematical formulas, tables, and lists. The bibliography system is now set up to work seamlessly with Bibby AI's citation management. Feel free to modify it for your needs!

\bibliographystyle{plain}
\bibliography{cite}

\end{document}\documentclass{article}
\usepackage[utf8]{inputenc}
\usepackage{booktabs}
\usepackage{amsmath}
\usepackage{natbib}

\title{My Document}
\author{Your Name}
\date{\today}

\begin{document}

\maketitle

\section{Introduction}
This is a sample document demonstrating various \LaTeX{} features including mathematical equations, tables, and formatting. You can cite papers using the citation search feature, and they will be automatically added to your \texttt{cite.bib} file.

\section{Mathematical Expressions}
Here are some example equations:

\subsection{Quadratic Formula}
The quadratic formula is given by:
\[x = \frac{-b \pm \sqrt{b^2 - 4ac}}{2a}\]

\subsection{Matrix Example}
A sample matrix:
\[
\begin{pmatrix}
a_{11} & a_{12} & a_{13} \\
a_{21} & a_{22} & a_{23} \\
a_{31} & a_{32} & a_{33}
\end{pmatrix}
\]

\section{Tables}
Below is a sample table showing data:

\begin{table}[h]
\centering
\begin{tabular}{@{}lrrr@{}}
\toprule
\textbf{Category} & \textbf{Value 1} & \textbf{Value 2} & \textbf{Total} \\
\midrule
Item A & 10.5 & 20.3 & 30.8 \\
Item B & 15.2 & 25.1 & 40.3 \\
Item C & 22.1 & 12.7 & 34.8 \\
\bottomrule
\end{tabular}
\caption{Sample Data Table}
\label{tab:sample}
\end{table}

\section{Lists}
\subsection{Bullet Points}
\begin{itemize}
  \item First item
  \item Second item
  \item Third item with sub-items:
    \begin{itemize}
      \item Sub-item 1
      \item Sub-item 2
    \end{itemize}
\end{itemize}

\subsection{Numbered List}
\begin{enumerate}
  \item First step
  \item Second step
  \item Third step
\end{enumerate}

\section{Citations and References}
When you use the citation search feature, references will be automatically managed. For example, you might cite a paper like this: \cite{example2024}. All citations will be automatically added to your \texttt{cite.bib} file.

\section{Conclusion}
This template demonstrates basic \LaTeX{} features including mathematical formulas, tables, and lists. The bibliography system is now set up to work seamlessly with Bibby AI's citation management. Feel free to modify it for your needs!

\bibliographystyle{plain}
\bibliography{cite}

\end{document}\documentclass{article}
\usepackage[utf8]{inputenc}
\usepackage{booktabs}
\usepackage{amsmath}
\usepackage{natbib}

\title{My Document}
\author{Your Name}
\date{\today}

\begin{document}

\maketitle

\section{Introduction}
This is a sample document demonstrating various \LaTeX{} features including mathematical equations, tables, and formatting. You can cite papers using the citation search feature, and they will be automatically added to your \texttt{cite.bib} file.

\section{Mathematical Expressions}
Here are some example equations:

\subsection{Quadratic Formula}
The quadratic formula is given by:
\[x = \frac{-b \pm \sqrt{b^2 - 4ac}}{2a}\]

\subsection{Matrix Example}
A sample matrix:
\[
\begin{pmatrix}
a_{11} & a_{12} & a_{13} \\
a_{21} & a_{22} & a_{23} \\
a_{31} & a_{32} & a_{33}
\end{pmatrix}
\]

\section{Tables}
Below is a sample table showing data:

\begin{table}[h]
\centering
\begin{tabular}{@{}lrrr@{}}
\toprule
\textbf{Category} & \textbf{Value 1} & \textbf{Value 2} & \textbf{Total} \\
\midrule
Item A & 10.5 & 20.3 & 30.8 \\
Item B & 15.2 & 25.1 & 40.3 \\
Item C & 22.1 & 12.7 & 34.8 \\
\bottomrule
\end{tabular}
\caption{Sample Data Table}
\label{tab:sample}
\end{table}

\section{Lists}
\subsection{Bullet Points}
\begin{itemize}
  \item First item
  \item Second item
  \item Third item with sub-items:
    \begin{itemize}
      \item Sub-item 1
      \item Sub-item 2
    \end{itemize}
\end{itemize}

\subsection{Numbered List}
\begin{enumerate}
  \item First step
  \item Second step
  \item Third step
\end{enumerate}

\section{Citations and References}
When you use the citation search feature, references will be automatically managed. For example, you might cite a paper like this: \cite{example2024}. All citations will be automatically added to your \texttt{cite.bib} file.

\section{Conclusion}
This template demonstrates basic \LaTeX{} features including mathematical formulas, tables, and lists. The bibliography system is now set up to work seamlessly with Bibby AI's citation management. Feel free to modify it for your needs!

\bibliographystyle{plain}
\bibliography{cite}

\end{document}\documentclass{article}
\usepackage[utf8]{inputenc}
\usepackage{booktabs}
\usepackage{amsmath}
\usepackage{natbib}

\title{My Document}
\author{Your Name}
\date{\today}

\begin{document}

\maketitle

\section{Introduction}
This is a sample document demonstrating various \LaTeX{} features including mathematical equations, tables, and formatting. You can cite papers using the citation search feature, and they will be automatically added to your \texttt{cite.bib} file.

\section{Mathematical Expressions}
Here are some example equations:

\subsection{Quadratic Formula}
The quadratic formula is given by:
\[x = \frac{-b \pm \sqrt{b^2 - 4ac}}{2a}\]

\subsection{Matrix Example}
A sample matrix:
\[
\begin{pmatrix}
a_{11} & a_{12} & a_{13} \\
a_{21} & a_{22} & a_{23} \\
a_{31} & a_{32} & a_{33}
\end{pmatrix}
\]

\section{Tables}
Below is a sample table showing data:

\begin{table}[h]
\centering
\begin{tabular}{@{}lrrr@{}}
\toprule
\textbf{Category} & \textbf{Value 1} & \textbf{Value 2} & \textbf{Total} \\
\midrule
Item A & 10.5 & 20.3 & 30.8 \\
Item B & 15.2 & 25.1 & 40.3 \\
Item C & 22.1 & 12.7 & 34.8 \\
\bottomrule
\end{tabular}
\caption{Sample Data Table}
\label{tab:sample}
\end{table}

\section{Lists}
\subsection{Bullet Points}
\begin{itemize}
  \item First item
  \item Second item
  \item Third item with sub-items:
    \begin{itemize}
      \item Sub-item 1
      \item Sub-item 2
    \end{itemize}
\end{itemize}

\subsection{Numbered List}
\begin{enumerate}
  \item First step
  \item Second step
  \item Third step
\end{enumerate}

\section{Citations and References}
When you use the citation search feature, references will be automatically managed. For example, you might cite a paper like this: \cite{example2024}. All citations will be automatically added to your \texttt{cite.bib} file.

\section{Conclusion}
This template demonstrates basic \LaTeX{} features including mathematical formulas, tables, and lists. The bibliography system is now set up to work seamlessly with Bibby AI's citation management. Feel free to modify it for your needs!

\bibliographystyle{plain}
\bibliography{cite}

\end{document}\documentclass{article}
\usepackage[utf8]{inputenc}
\usepackage{booktabs}
\usepackage{amsmath}
\usepackage{natbib}

\title{My Document}
\author{Your Name}
\date{\today}

\begin{document}

\maketitle

\section{Introduction}
This is a sample document demonstrating various \LaTeX{} features including mathematical equations, tables, and formatting. You can cite papers using the citation search feature, and they will be automatically added to your \texttt{cite.bib} file.

\section{Mathematical Expressions}
Here are some example equations:

\subsection{Quadratic Formula}
The quadratic formula is given by:
\[x = \frac{-b \pm \sqrt{b^2 - 4ac}}{2a}\]

\subsection{Matrix Example}
A sample matrix:
\[
\begin{pmatrix}
a_{11} & a_{12} & a_{13} \\
a_{21} & a_{22} & a_{23} \\
a_{31} & a_{32} & a_{33}
\end{pmatrix}
\]

\section{Tables}
Below is a sample table showing data:

\begin{table}[h]
\centering
\begin{tabular}{@{}lrrr@{}}
\toprule
\textbf{Category} & \textbf{Value 1} & \textbf{Value 2} & \textbf{Total} \\
\midrule
Item A & 10.5 & 20.3 & 30.8 \\
Item B & 15.2 & 25.1 & 40.3 \\
Item C & 22.1 & 12.7 & 34.8 \\
\bottomrule
\end{tabular}
\caption{Sample Data Table}
\label{tab:sample}
\end{table}

\section{Lists}
\subsection{Bullet Points}
\begin{itemize}
  \item First item
  \item Second item
  \item Third item with sub-items:
    \begin{itemize}
      \item Sub-item 1
      \item Sub-item 2
    \end{itemize}
\end{itemize}

\subsection{Numbered List}
\begin{enumerate}
  \item First step
  \item Second step
  \item Third step
\end{enumerate}

\section{Citations and References}
When you use the citation search feature, references will be automatically managed. For example, you might cite a paper like this: \cite{example2024}. All citations will be automatically added to your \texttt{cite.bib} file.

\section{Conclusion}
This template demonstrates basic \LaTeX{} features including mathematical formulas, tables, and lists. The bibliography system is now set up to work seamlessly with Bibby AI's citation management. Feel free to modify it for your needs!

\bibliographystyle{plain}
\bibliography{cite}

\end{document}\documentclass{article}
\usepackage[utf8]{inputenc}
\usepackage{booktabs}
\usepackage{amsmath}
\usepackage{natbib}

\title{My Document}
\author{Your Name}
\date{\today}

\begin{document}

\maketitle

\section{Introduction}
This is a sample document demonstrating various \LaTeX{} features including mathematical equations, tables, and formatting. You can cite papers using the citation search feature, and they will be automatically added to your \texttt{cite.bib} file.

\section{Mathematical Expressions}
Here are some example equations:

\subsection{Quadratic Formula}
The quadratic formula is given by:
\[x = \frac{-b \pm \sqrt{b^2 - 4ac}}{2a}\]

\subsection{Matrix Example}
A sample matrix:
\[
\begin{pmatrix}
a_{11} & a_{12} & a_{13} \\
a_{21} & a_{22} & a_{23} \\
a_{31} & a_{32} & a_{33}
\end{pmatrix}
\]

\section{Tables}
Below is a sample table showing data:

\begin{table}[h]
\centering
\begin{tabular}{@{}lrrr@{}}
\toprule
\textbf{Category} & \textbf{Value 1} & \textbf{Value 2} & \textbf{Total} \\
\midrule
Item A & 10.5 & 20.3 & 30.8 \\
Item B & 15.2 & 25.1 & 40.3 \\
Item C & 22.1 & 12.7 & 34.8 \\
\bottomrule
\end{tabular}
\caption{Sample Data Table}
\label{tab:sample}
\end{table}

\section{Lists}
\subsection{Bullet Points}
\begin{itemize}
  \item First item
  \item Second item
  \item Third item with sub-items:
    \begin{itemize}
      \item Sub-item 1
      \item Sub-item 2
    \end{itemize}
\end{itemize}

\subsection{Numbered List}
\begin{enumerate}
  \item First step
  \item Second step
  \item Third step
\end{enumerate}

\section{Citations and References}
When you use the citation search feature, references will be automatically managed. For example, you might cite a paper like this: \cite{example2024}. All citations will be automatically added to your \texttt{cite.bib} file.

\section{Conclusion}
This template demonstrates basic \LaTeX{} features including mathematical formulas, tables, and lists. The bibliography system is now set up to work seamlessly with Bibby AI's citation management. Feel free to modify it for your needs!

\bibliographystyle{plain}
\bibliography{cite}

\end{document}\documentclass{article}
\usepackage[utf8]{inputenc}
\usepackage{booktabs}
\usepackage{amsmath}
\usepackage{natbib}

\title{My Document}
\author{Your Name}
\date{\today}

\begin{document}

\maketitle

\section{Introduction}
This is a sample document demonstrating various \LaTeX{} features including mathematical equations, tables, and formatting. You can cite papers using the citation search feature, and they will be automatically added to your \texttt{cite.bib} file.

\section{Mathematical Expressions}
Here are some example equations:

\subsection{Quadratic Formula}
The quadratic formula is given by:
\[x = \frac{-b \pm \sqrt{b^2 - 4ac}}{2a}\]

\subsection{Matrix Example}
A sample matrix:
\[
\begin{pmatrix}
a_{11} & a_{12} & a_{13} \\
a_{21} & a_{22} & a_{23} \\
a_{31} & a_{32} & a_{33}
\end{pmatrix}
\]

\section{Tables}
Below is a sample table showing data:

\begin{table}[h]
\centering
\begin{tabular}{@{}lrrr@{}}
\toprule
\textbf{Category} & \textbf{Value 1} & \textbf{Value 2} & \textbf{Total} \\
\midrule
Item A & 10.5 & 20.3 & 30.8 \\
Item B & 15.2 & 25.1 & 40.3 \\
Item C & 22.1 & 12.7 & 34.8 \\
\bottomrule
\end{tabular}
\caption{Sample Data Table}
\label{tab:sample}
\end{table}

\section{Lists}
\subsection{Bullet Points}
\begin{itemize}
  \item First item
  \item Second item
  \item Third item with sub-items:
    \begin{itemize}
      \item Sub-item 1
      \item Sub-item 2
    \end{itemize}
\end{itemize}

\subsection{Numbered List}
\begin{enumerate}
  \item First step
  \item Second step
  \item Third step
\end{enumerate}

\section{Citations and References}
When you use the citation search feature, references will be automatically managed. For example, you might cite a paper like this: \cite{example2024}. All citations will be automatically added to your \texttt{cite.bib} file.

\section{Conclusion}
This template demonstrates basic \LaTeX{} features including mathematical formulas, tables, and lists. The bibliography system is now set up to work seamlessly with Bibby AI's citation management. Feel free to modify it for your needs!

\bibliographystyle{plain}
\bibliography{cite}

\end{document}\documentclass{article}
\usepackage[utf8]{inputenc}
\usepackage{booktabs}
\usepackage{amsmath}
\usepackage{natbib}

\title{My Document}
\author{Your Name}
\date{\today}

\begin{document}

\maketitle

\section{Introduction}
This is a sample document demonstrating various \LaTeX{} features including mathematical equations, tables, and formatting. You can cite papers using the citation search feature, and they will be automatically added to your \texttt{cite.bib} file.

\section{Mathematical Expressions}
Here are some example equations:

\subsection{Quadratic Formula}
The quadratic formula is given by:
\[x = \frac{-b \pm \sqrt{b^2 - 4ac}}{2a}\]

\subsection{Matrix Example}
A sample matrix:
\[
\begin{pmatrix}
a_{11} & a_{12} & a_{13} \\
a_{21} & a_{22} & a_{23} \\
a_{31} & a_{32} & a_{33}
\end{pmatrix}
\]

\section{Tables}
Below is a sample table showing data:

\begin{table}[h]
\centering
\begin{tabular}{@{}lrrr@{}}
\toprule
\textbf{Category} & \textbf{Value 1} & \textbf{Value 2} & \textbf{Total} \\
\midrule
Item A & 10.5 & 20.3 & 30.8 \\
Item B & 15.2 & 25.1 & 40.3 \\
Item C & 22.1 & 12.7 & 34.8 \\
\bottomrule
\end{tabular}
\caption{Sample Data Table}
\label{tab:sample}
\end{table}

\section{Lists}
\subsection{Bullet Points}
\begin{itemize}
  \item First item
  \item Second item
  \item Third item with sub-items:
    \begin{itemize}
      \item Sub-item 1
      \item Sub-item 2
    \end{itemize}
\end{itemize}

\subsection{Numbered List}
\begin{enumerate}
  \item First step
  \item Second step
  \item Third step
\end{enumerate}

\section{Citations and References}
When you use the citation search feature, references will be automatically managed. For example, you might cite a paper like this: \cite{example2024}. All citations will be automatically added to your \texttt{cite.bib} file.

\section{Conclusion}
This template demonstrates basic \LaTeX{} features including mathematical formulas, tables, and lists. The bibliography system is now set up to work seamlessly with Bibby AI's citation management. Feel free to modify it for your needs!

\bibliographystyle{plain}
\bibliography{cite}

\end{document}\documentclass{article}
\usepackage[utf8]{inputenc}
\usepackage{booktabs}
\usepackage{amsmath}
\usepackage{natbib}

\title{My Document}
\author{Your Name}
\date{\today}

\begin{document}

\maketitle

\section{Introduction}
This is a sample document demonstrating various \LaTeX{} features including mathematical equations, tables, and formatting. You can cite papers using the citation search feature, and they will be automatically added to your \texttt{cite.bib} file.

\section{Mathematical Expressions}
Here are some example equations:

\subsection{Quadratic Formula}
The quadratic formula is given by:
\[x = \frac{-b \pm \sqrt{b^2 - 4ac}}{2a}\]

\subsection{Matrix Example}
A sample matrix:
\[
\begin{pmatrix}
a_{11} & a_{12} & a_{13} \\
a_{21} & a_{22} & a_{23} \\
a_{31} & a_{32} & a_{33}
\end{pmatrix}
\]

\section{Tables}
Below is a sample table showing data:

\begin{table}[h]
\centering
\begin{tabular}{@{}lrrr@{}}
\toprule
\textbf{Category} & \textbf{Value 1} & \textbf{Value 2} & \textbf{Total} \\
\midrule
Item A & 10.5 & 20.3 & 30.8 \\
Item B & 15.2 & 25.1 & 40.3 \\
Item C & 22.1 & 12.7 & 34.8 \\
\bottomrule
\end{tabular}
\caption{Sample Data Table}
\label{tab:sample}
\end{table}

\section{Lists}
\subsection{Bullet Points}
\begin{itemize}
  \item First item
  \item Second item
  \item Third item with sub-items:
    \begin{itemize}
      \item Sub-item 1
      \item Sub-item 2
    \end{itemize}
\end{itemize}

\subsection{Numbered List}
\begin{enumerate}
  \item First step
  \item Second step
  \item Third step
\end{enumerate}

\section{Citations and References}
When you use the citation search feature, references will be automatically managed. For example, you might cite a paper like this: \cite{example2024}. All citations will be automatically added to your \texttt{cite.bib} file.

\section{Conclusion}
This template demonstrates basic \LaTeX{} features including mathematical formulas, tables, and lists. The bibliography system is now set up to work seamlessly with Bibby AI's citation management. Feel free to modify it for your needs!

\bibliographystyle{plain}
\bibliography{cite}

\end{document}\documentclass{article}
\usepackage[utf8]{inputenc}
\usepackage{booktabs}
\usepackage{amsmath}
\usepackage{natbib}

\title{My Document}
\author{Your Name}
\date{\today}

\begin{document}

\maketitle

\section{Introduction}
This is a sample document demonstrating various \LaTeX{} features including mathematical equations, tables, and formatting. You can cite papers using the citation search feature, and they will be automatically added to your \texttt{cite.bib} file.

\section{Mathematical Expressions}
Here are some example equations:

\subsection{Quadratic Formula}
The quadratic formula is given by:
\[x = \frac{-b \pm \sqrt{b^2 - 4ac}}{2a}\]

\subsection{Matrix Example}
A sample matrix:
\[
\begin{pmatrix}
a_{11} & a_{12} & a_{13} \\
a_{21} & a_{22} & a_{23} \\
a_{31} & a_{32} & a_{33}
\end{pmatrix}
\]

\section{Tables}
Below is a sample table showing data:

\begin{table}[h]
\centering
\begin{tabular}{@{}lrrr@{}}
\toprule
\textbf{Category} & \textbf{Value 1} & \textbf{Value 2} & \textbf{Total} \\
\midrule
Item A & 10.5 & 20.3 & 30.8 \\
Item B & 15.2 & 25.1 & 40.3 \\
Item C & 22.1 & 12.7 & 34.8 \\
\bottomrule
\end{tabular}
\caption{Sample Data Table}
\label{tab:sample}
\end{table}

\section{Lists}
\subsection{Bullet Points}
\begin{itemize}
  \item First item
  \item Second item
  \item Third item with sub-items:
    \begin{itemize}
      \item Sub-item 1
      \item Sub-item 2
    \end{itemize}
\end{itemize}

\subsection{Numbered List}
\begin{enumerate}
  \item First step
  \item Second step
  \item Third step
\end{enumerate}

\section{Citations and References}
When you use the citation search feature, references will be automatically managed. For example, you might cite a paper like this: \cite{example2024}. All citations will be automatically added to your \texttt{cite.bib} file.

\section{Conclusion}
This template demonstrates basic \LaTeX{} features including mathematical formulas, tables, and lists. The bibliography system is now set up to work seamlessly with Bibby AI's citation management. Feel free to modify it for your needs!

\bibliographystyle{plain}
\bibliography{cite}

\end{document}\documentclass{article}
\usepackage[utf8]{inputenc}
\usepackage{booktabs}
\usepackage{amsmath}
\usepackage{natbib}

\title{My Document}
\author{Your Name}
\date{\today}

\begin{document}

\maketitle

\section{Introduction}
This is a sample document demonstrating various \LaTeX{} features including mathematical equations, tables, and formatting. You can cite papers using the citation search feature, and they will be automatically added to your \texttt{cite.bib} file.

\section{Mathematical Expressions}
Here are some example equations:

\subsection{Quadratic Formula}
The quadratic formula is given by:
\[x = \frac{-b \pm \sqrt{b^2 - 4ac}}{2a}\]

\subsection{Matrix Example}
A sample matrix:
\[
\begin{pmatrix}
a_{11} & a_{12} & a_{13} \\
a_{21} & a_{22} & a_{23} \\
a_{31} & a_{32} & a_{33}
\end{pmatrix}
\]

\section{Tables}
Below is a sample table showing data:

\begin{table}[h]
\centering
\begin{tabular}{@{}lrrr@{}}
\toprule
\textbf{Category} & \textbf{Value 1} & \textbf{Value 2} & \textbf{Total} \\
\midrule
Item A & 10.5 & 20.3 & 30.8 \\
Item B & 15.2 & 25.1 & 40.3 \\
Item C & 22.1 & 12.7 & 34.8 \\
\bottomrule
\end{tabular}
\caption{Sample Data Table}
\label{tab:sample}
\end{table}

\section{Lists}
\subsection{Bullet Points}
\begin{itemize}
  \item First item
  \item Second item
  \item Third item with sub-items:
    \begin{itemize}
      \item Sub-item 1
      \item Sub-item 2
    \end{itemize}
\end{itemize}

\subsection{Numbered List}
\begin{enumerate}
  \item First step
  \item Second step
  \item Third step
\end{enumerate}

\section{Citations and References}
When you use the citation search feature, references will be automatically managed. For example, you might cite a paper like this: \cite{example2024}. All citations will be automatically added to your \texttt{cite.bib} file.

\section{Conclusion}
This template demonstrates basic \LaTeX{} features including mathematical formulas, tables, and lists. The bibliography system is now set up to work seamlessly with Bibby AI's citation management. Feel free to modify it for your needs!

\bibliographystyle{plain}
\bibliography{cite}

\end{document}\documentclass{article}
\usepackage[utf8]{inputenc}
\usepackage{booktabs}
\usepackage{amsmath}
\usepackage{natbib}

\title{My Document}
\author{Your Name}
\date{\today}

\begin{document}

\maketitle

\section{Introduction}
This is a sample document demonstrating various \LaTeX{} features including mathematical equations, tables, and formatting. You can cite papers using the citation search feature, and they will be automatically added to your \texttt{cite.bib} file.

\section{Mathematical Expressions}
Here are some example equations:

\subsection{Quadratic Formula}
The quadratic formula is given by:
\[x = \frac{-b \pm \sqrt{b^2 - 4ac}}{2a}\]

\subsection{Matrix Example}
A sample matrix:
\[
\begin{pmatrix}
a_{11} & a_{12} & a_{13} \\
a_{21} & a_{22} & a_{23} \\
a_{31} & a_{32} & a_{33}
\end{pmatrix}
\]

\section{Tables}
Below is a sample table showing data:

\begin{table}[h]
\centering
\begin{tabular}{@{}lrrr@{}}
\toprule
\textbf{Category} & \textbf{Value 1} & \textbf{Value 2} & \textbf{Total} \\
\midrule
Item A & 10.5 & 20.3 & 30.8 \\
Item B & 15.2 & 25.1 & 40.3 \\
Item C & 22.1 & 12.7 & 34.8 \\
\bottomrule
\end{tabular}
\caption{Sample Data Table}
\label{tab:sample}
\end{table}

\section{Lists}
\subsection{Bullet Points}
\begin{itemize}
  \item First item
  \item Second item
  \item Third item with sub-items:
    \begin{itemize}
      \item Sub-item 1
      \item Sub-item 2
    \end{itemize}
\end{itemize}

\subsection{Numbered List}
\begin{enumerate}
  \item First step
  \item Second step
  \item Third step
\end{enumerate}

\section{Citations and References}
When you use the citation search feature, references will be automatically managed. For example, you might cite a paper like this: \cite{example2024}. All citations will be automatically added to your \texttt{cite.bib} file.

\section{Conclusion}
This template demonstrates basic \LaTeX{} features including mathematical formulas, tables, and lists. The bibliography system is now set up to work seamlessly with Bibby AI's citation management. Feel free to modify it for your needs!

\bibliographystyle{plain}
\bibliography{cite}

\end{document}\documentclass{article}
\usepackage[utf8]{inputenc}
\usepackage{booktabs}
\usepackage{amsmath}
\usepackage{natbib}

\title{My Document}
\author{Your Name}
\date{\today}

\begin{document}

\maketitle

\section{Introduction}
This is a sample document demonstrating various \LaTeX{} features including mathematical equations, tables, and formatting. You can cite papers using the citation search feature, and they will be automatically added to your \texttt{cite.bib} file.

\section{Mathematical Expressions}
Here are some example equations:

\subsection{Quadratic Formula}
The quadratic formula is given by:
\[x = \frac{-b \pm \sqrt{b^2 - 4ac}}{2a}\]

\subsection{Matrix Example}
A sample matrix:
\[
\begin{pmatrix}
a_{11} & a_{12} & a_{13} \\
a_{21} & a_{22} & a_{23} \\
a_{31} & a_{32} & a_{33}
\end{pmatrix}
\]

\section{Tables}
Below is a sample table showing data:

\begin{table}[h]
\centering
\begin{tabular}{@{}lrrr@{}}
\toprule
\textbf{Category} & \textbf{Value 1} & \textbf{Value 2} & \textbf{Total} \\
\midrule
Item A & 10.5 & 20.3 & 30.8 \\
Item B & 15.2 & 25.1 & 40.3 \\
Item C & 22.1 & 12.7 & 34.8 \\
\bottomrule
\end{tabular}
\caption{Sample Data Table}
\label{tab:sample}
\end{table}

\section{Lists}
\subsection{Bullet Points}
\begin{itemize}
  \item First item
  \item Second item
  \item Third item with sub-items:
    \begin{itemize}
      \item Sub-item 1
      \item Sub-item 2
    \end{itemize}
\end{itemize}

\subsection{Numbered List}
\begin{enumerate}
  \item First step
  \item Second step
  \item Third step
\end{enumerate}

\section{Citations and References}
When you use the citation search feature, references will be automatically managed. For example, you might cite a paper like this: \cite{example2024}. All citations will be automatically added to your \texttt{cite.bib} file.

\section{Conclusion}
This template demonstrates basic \LaTeX{} features including mathematical formulas, tables, and lists. The bibliography system is now set up to work seamlessly with Bibby AI's citation management. Feel free to modify it for your needs!

\bibliographystyle{plain}
\bibliography{cite}

\end{document}\documentclass{article}
\usepackage[utf8]{inputenc}
\usepackage{booktabs}
\usepackage{amsmath}
\usepackage{natbib}

\title{My Document}
\author{Your Name}
\date{\today}

\begin{document}

\maketitle

\section{Introduction}
This is a sample document demonstrating various \LaTeX{} features including mathematical equations, tables, and formatting. You can cite papers using the citation search feature, and they will be automatically added to your \texttt{cite.bib} file.

\section{Mathematical Expressions}
Here are some example equations:

\subsection{Quadratic Formula}
The quadratic formula is given by:
\[x = \frac{-b \pm \sqrt{b^2 - 4ac}}{2a}\]

\subsection{Matrix Example}
A sample matrix:
\[
\begin{pmatrix}
a_{11} & a_{12} & a_{13} \\
a_{21} & a_{22} & a_{23} \\
a_{31} & a_{32} & a_{33}
\end{pmatrix}
\]

\section{Tables}
Below is a sample table showing data:

\begin{table}[h]
\centering
\begin{tabular}{@{}lrrr@{}}
\toprule
\textbf{Category} & \textbf{Value 1} & \textbf{Value 2} & \textbf{Total} \\
\midrule
Item A & 10.5 & 20.3 & 30.8 \\
Item B & 15.2 & 25.1 & 40.3 \\
Item C & 22.1 & 12.7 & 34.8 \\
\bottomrule
\end{tabular}
\caption{Sample Data Table}
\label{tab:sample}
\end{table}

\section{Lists}
\subsection{Bullet Points}
\begin{itemize}
  \item First item
  \item Second item
  \item Third item with sub-items:
    \begin{itemize}
      \item Sub-item 1
      \item Sub-item 2
    \end{itemize}
\end{itemize}

\subsection{Numbered List}
\begin{enumerate}
  \item First step
  \item Second step
  \item Third step
\end{enumerate}

\section{Citations and References}
When you use the citation search feature, references will be automatically managed. For example, you might cite a paper like this: \cite{example2024}. All citations will be automatically added to your \texttt{cite.bib} file.

\section{Conclusion}
This template demonstrates basic \LaTeX{} features including mathematical formulas, tables, and lists. The bibliography system is now set up to work seamlessly with Bibby AI's citation management. Feel free to modify it for your needs!

\bibliographystyle{plain}
\bibliography{cite}

\end{document}\documentclass{article}
\usepackage[utf8]{inputenc}
\usepackage{booktabs}
\usepackage{amsmath}
\usepackage{natbib}

\title{My Document}
\author{Your Name}
\date{\today}

\begin{document}

\maketitle

\section{Introduction}
This is a sample document demonstrating various \LaTeX{} features including mathematical equations, tables, and formatting. You can cite papers using the citation search feature, and they will be automatically added to your \texttt{cite.bib} file.

\section{Mathematical Expressions}
Here are some example equations:

\subsection{Quadratic Formula}
The quadratic formula is given by:
\[x = \frac{-b \pm \sqrt{b^2 - 4ac}}{2a}\]

\subsection{Matrix Example}
A sample matrix:
\[
\begin{pmatrix}
a_{11} & a_{12} & a_{13} \\
a_{21} & a_{22} & a_{23} \\
a_{31} & a_{32} & a_{33}
\end{pmatrix}
\]

\section{Tables}
Below is a sample table showing data:

\begin{table}[h]
\centering
\begin{tabular}{@{}lrrr@{}}
\toprule
\textbf{Category} & \textbf{Value 1} & \textbf{Value 2} & \textbf{Total} \\
\midrule
Item A & 10.5 & 20.3 & 30.8 \\
Item B & 15.2 & 25.1 & 40.3 \\
Item C & 22.1 & 12.7 & 34.8 \\
\bottomrule
\end{tabular}
\caption{Sample Data Table}
\label{tab:sample}
\end{table}

\section{Lists}
\subsection{Bullet Points}
\begin{itemize}
  \item First item
  \item Second item
  \item Third item with sub-items:
    \begin{itemize}
      \item Sub-item 1
      \item Sub-item 2
    \end{itemize}
\end{itemize}

\subsection{Numbered List}
\begin{enumerate}
  \item First step
  \item Second step
  \item Third step
\end{enumerate}

\section{Citations and References}
When you use the citation search feature, references will be automatically managed. For example, you might cite a paper like this: \cite{example2024}. All citations will be automatically added to your \texttt{cite.bib} file.

\section{Conclusion}
This template demonstrates basic \LaTeX{} features including mathematical formulas, tables, and lists. The bibliography system is now set up to work seamlessly with Bibby AI's citation management. Feel free to modify it for your needs!

\bibliographystyle{plain}
\bibliography{cite}

\end{document}\documentclass{article}
\usepackage[utf8]{inputenc}
\usepackage{booktabs}
\usepackage{amsmath}
\usepackage{natbib}

\title{My Document}
\author{Your Name}
\date{\today}

\begin{document}

\maketitle

\section{Introduction}
This is a sample document demonstrating various \LaTeX{} features including mathematical equations, tables, and formatting. You can cite papers using the citation search feature, and they will be automatically added to your \texttt{cite.bib} file.

\section{Mathematical Expressions}
Here are some example equations:

\subsection{Quadratic Formula}
The quadratic formula is given by:
\[x = \frac{-b \pm \sqrt{b^2 - 4ac}}{2a}\]

\subsection{Matrix Example}
A sample matrix:
\[
\begin{pmatrix}
a_{11} & a_{12} & a_{13} \\
a_{21} & a_{22} & a_{23} \\
a_{31} & a_{32} & a_{33}
\end{pmatrix}
\]

\section{Tables}
Below is a sample table showing data:

\begin{table}[h]
\centering
\begin{tabular}{@{}lrrr@{}}
\toprule
\textbf{Category} & \textbf{Value 1} & \textbf{Value 2} & \textbf{Total} \\
\midrule
Item A & 10.5 & 20.3 & 30.8 \\
Item B & 15.2 & 25.1 & 40.3 \\
Item C & 22.1 & 12.7 & 34.8 \\
\bottomrule
\end{tabular}
\caption{Sample Data Table}
\label{tab:sample}
\end{table}

\section{Lists}
\subsection{Bullet Points}
\begin{itemize}
  \item First item
  \item Second item
  \item Third item with sub-items:
    \begin{itemize}
      \item Sub-item 1
      \item Sub-item 2
    \end{itemize}
\end{itemize}

\subsection{Numbered List}
\begin{enumerate}
  \item First step
  \item Second step
  \item Third step
\end{enumerate}

\section{Citations and References}
When you use the citation search feature, references will be automatically managed. For example, you might cite a paper like this: \cite{example2024}. All citations will be automatically added to your \texttt{cite.bib} file.

\section{Conclusion}
This template demonstrates basic \LaTeX{} features including mathematical formulas, tables, and lists. The bibliography system is now set up to work seamlessly with Bibby AI's citation management. Feel free to modify it for your needs!

\bibliographystyle{plain}
\bibliography{cite}

\end{document}\documentclass{article}
\usepackage[utf8]{inputenc}
\usepackage{booktabs}
\usepackage{amsmath}
\usepackage{natbib}

\title{My Document}
\author{Your Name}
\date{\today}

\begin{document}

\maketitle

\section{Introduction}
This is a sample document demonstrating various \LaTeX{} features including mathematical equations, tables, and formatting. You can cite papers using the citation search feature, and they will be automatically added to your \texttt{cite.bib} file.

\section{Mathematical Expressions}
Here are some example equations:

\subsection{Quadratic Formula}
The quadratic formula is given by:
\[x = \frac{-b \pm \sqrt{b^2 - 4ac}}{2a}\]

\subsection{Matrix Example}
A sample matrix:
\[
\begin{pmatrix}
a_{11} & a_{12} & a_{13} \\
a_{21} & a_{22} & a_{23} \\
a_{31} & a_{32} & a_{33}
\end{pmatrix}
\]

\section{Tables}
Below is a sample table showing data:

\begin{table}[h]
\centering
\begin{tabular}{@{}lrrr@{}}
\toprule
\textbf{Category} & \textbf{Value 1} & \textbf{Value 2} & \textbf{Total} \\
\midrule
Item A & 10.5 & 20.3 & 30.8 \\
Item B & 15.2 & 25.1 & 40.3 \\
Item C & 22.1 & 12.7 & 34.8 \\
\bottomrule
\end{tabular}
\caption{Sample Data Table}
\label{tab:sample}
\end{table}

\section{Lists}
\subsection{Bullet Points}
\begin{itemize}
  \item First item
  \item Second item
  \item Third item with sub-items:
    \begin{itemize}
      \item Sub-item 1
      \item Sub-item 2
    \end{itemize}
\end{itemize}

\subsection{Numbered List}
\begin{enumerate}
  \item First step
  \item Second step
  \item Third step
\end{enumerate}

\section{Citations and References}
When you use the citation search feature, references will be automatically managed. For example, you might cite a paper like this: \cite{example2024}. All citations will be automatically added to your \texttt{cite.bib} file.

\section{Conclusion}
This template demonstrates basic \LaTeX{} features including mathematical formulas, tables, and lists. The bibliography system is now set up to work seamlessly with Bibby AI's citation management. Feel free to modify it for your needs!

\bibliographystyle{plain}
\bibliography{cite}

\end{document}\documentclass{article}
\usepackage[utf8]{inputenc}
\usepackage{booktabs}
\usepackage{amsmath}
\usepackage{natbib}

\title{My Document}
\author{Your Name}
\date{\today}

\begin{document}

\maketitle

\section{Introduction}
This is a sample document demonstrating various \LaTeX{} features including mathematical equations, tables, and formatting. You can cite papers using the citation search feature, and they will be automatically added to your \texttt{cite.bib} file.

\section{Mathematical Expressions}
Here are some example equations:

\subsection{Quadratic Formula}
The quadratic formula is given by:
\[x = \frac{-b \pm \sqrt{b^2 - 4ac}}{2a}\]

\subsection{Matrix Example}
A sample matrix:
\[
\begin{pmatrix}
a_{11} & a_{12} & a_{13} \\
a_{21} & a_{22} & a_{23} \\
a_{31} & a_{32} & a_{33}
\end{pmatrix}
\]

\section{Tables}
Below is a sample table showing data:

\begin{table}[h]
\centering
\begin{tabular}{@{}lrrr@{}}
\toprule
\textbf{Category} & \textbf{Value 1} & \textbf{Value 2} & \textbf{Total} \\
\midrule
Item A & 10.5 & 20.3 & 30.8 \\
Item B & 15.2 & 25.1 & 40.3 \\
Item C & 22.1 & 12.7 & 34.8 \\
\bottomrule
\end{tabular}
\caption{Sample Data Table}
\label{tab:sample}
\end{table}

\section{Lists}
\subsection{Bullet Points}
\begin{itemize}
  \item First item
  \item Second item
  \item Third item with sub-items:
    \begin{itemize}
      \item Sub-item 1
      \item Sub-item 2
    \end{itemize}
\end{itemize}

\subsection{Numbered List}
\begin{enumerate}
  \item First step
  \item Second step
  \item Third step
\end{enumerate}

\section{Citations and References}
When you use the citation search feature, references will be automatically managed. For example, you might cite a paper like this: \cite{example2024}. All citations will be automatically added to your \texttt{cite.bib} file.

\section{Conclusion}
This template demonstrates basic \LaTeX{} features including mathematical formulas, tables, and lists. The bibliography system is now set up to work seamlessly with Bibby AI's citation management. Feel free to modify it for your needs!

\bibliographystyle{plain}
\bibliography{cite}

\end{document}\documentclass{article}
\usepackage[utf8]{inputenc}
\usepackage{booktabs}
\usepackage{amsmath}
\usepackage{natbib}

\title{My Document}
\author{Your Name}
\date{\today}

\begin{document}

\maketitle

\section{Introduction}
This is a sample document demonstrating various \LaTeX{} features including mathematical equations, tables, and formatting. You can cite papers using the citation search feature, and they will be automatically added to your \texttt{cite.bib} file.

\section{Mathematical Expressions}
Here are some example equations:

\subsection{Quadratic Formula}
The quadratic formula is given by:
\[x = \frac{-b \pm \sqrt{b^2 - 4ac}}{2a}\]

\subsection{Matrix Example}
A sample matrix:
\[
\begin{pmatrix}
a_{11} & a_{12} & a_{13} \\
a_{21} & a_{22} & a_{23} \\
a_{31} & a_{32} & a_{33}
\end{pmatrix}
\]

\section{Tables}
Below is a sample table showing data:

\begin{table}[h]
\centering
\begin{tabular}{@{}lrrr@{}}
\toprule
\textbf{Category} & \textbf{Value 1} & \textbf{Value 2} & \textbf{Total} \\
\midrule
Item A & 10.5 & 20.3 & 30.8 \\
Item B & 15.2 & 25.1 & 40.3 \\
Item C & 22.1 & 12.7 & 34.8 \\
\bottomrule
\end{tabular}
\caption{Sample Data Table}
\label{tab:sample}
\end{table}

\section{Lists}
\subsection{Bullet Points}
\begin{itemize}
  \item First item
  \item Second item
  \item Third item with sub-items:
    \begin{itemize}
      \item Sub-item 1
      \item Sub-item 2
    \end{itemize}
\end{itemize}

\subsection{Numbered List}
\begin{enumerate}
  \item First step
  \item Second step
  \item Third step
\end{enumerate}

\section{Citations and References}
When you use the citation search feature, references will be automatically managed. For example, you might cite a paper like this: \cite{example2024}. All citations will be automatically added to your \texttt{cite.bib} file.

\section{Conclusion}
This template demonstrates basic \LaTeX{} features including mathematical formulas, tables, and lists. The bibliography system is now set up to work seamlessly with Bibby AI's citation management. Feel free to modify it for your needs!

\bibliographystyle{plain}
\bibliography{cite}

\end{document}\documentclass{article}
\usepackage[utf8]{inputenc}
\usepackage{booktabs}
\usepackage{amsmath}
\usepackage{natbib}

\title{My Document}
\author{Your Name}
\date{\today}

\begin{document}

\maketitle

\section{Introduction}
This is a sample document demonstrating various \LaTeX{} features including mathematical equations, tables, and formatting. You can cite papers using the citation search feature, and they will be automatically added to your \texttt{cite.bib} file.

\section{Mathematical Expressions}
Here are some example equations:

\subsection{Quadratic Formula}
The quadratic formula is given by:
\[x = \frac{-b \pm \sqrt{b^2 - 4ac}}{2a}\]

\subsection{Matrix Example}
A sample matrix:
\[
\begin{pmatrix}
a_{11} & a_{12} & a_{13} \\
a_{21} & a_{22} & a_{23} \\
a_{31} & a_{32} & a_{33}
\end{pmatrix}
\]

\section{Tables}
Below is a sample table showing data:

\begin{table}[h]
\centering
\begin{tabular}{@{}lrrr@{}}
\toprule
\textbf{Category} & \textbf{Value 1} & \textbf{Value 2} & \textbf{Total} \\
\midrule
Item A & 10.5 & 20.3 & 30.8 \\
Item B & 15.2 & 25.1 & 40.3 \\
Item C & 22.1 & 12.7 & 34.8 \\
\bottomrule
\end{tabular}
\caption{Sample Data Table}
\label{tab:sample}
\end{table}

\section{Lists}
\subsection{Bullet Points}
\begin{itemize}
  \item First item
  \item Second item
  \item Third item with sub-items:
    \begin{itemize}
      \item Sub-item 1
      \item Sub-item 2
    \end{itemize}
\end{itemize}

\subsection{Numbered List}
\begin{enumerate}
  \item First step
  \item Second step
  \item Third step
\end{enumerate}

\section{Citations and References}
When you use the citation search feature, references will be automatically managed. For example, you might cite a paper like this: \cite{example2024}. All citations will be automatically added to your \texttt{cite.bib} file.

\section{Conclusion}
This template demonstrates basic \LaTeX{} features including mathematical formulas, tables, and lists. The bibliography system is now set up to work seamlessly with Bibby AI's citation management. Feel free to modify it for your needs!

\bibliographystyle{plain}
\bibliography{cite}

\end{document}\documentclass{article}
\usepackage[utf8]{inputenc}
\usepackage{booktabs}
\usepackage{amsmath}
\usepackage{natbib}

\title{My Document}
\author{Your Name}
\date{\today}

\begin{document}

\maketitle

\section{Introduction}
This is a sample document demonstrating various \LaTeX{} features including mathematical equations, tables, and formatting. You can cite papers using the citation search feature, and they will be automatically added to your \texttt{cite.bib} file.

\section{Mathematical Expressions}
Here are some example equations:

\subsection{Quadratic Formula}
The quadratic formula is given by:
\[x = \frac{-b \pm \sqrt{b^2 - 4ac}}{2a}\]

\subsection{Matrix Example}
A sample matrix:
\[
\begin{pmatrix}
a_{11} & a_{12} & a_{13} \\
a_{21} & a_{22} & a_{23} \\
a_{31} & a_{32} & a_{33}
\end{pmatrix}
\]

\section{Tables}
Below is a sample table showing data:

\begin{table}[h]
\centering
\begin{tabular}{@{}lrrr@{}}
\toprule
\textbf{Category} & \textbf{Value 1} & \textbf{Value 2} & \textbf{Total} \\
\midrule
Item A & 10.5 & 20.3 & 30.8 \\
Item B & 15.2 & 25.1 & 40.3 \\
Item C & 22.1 & 12.7 & 34.8 \\
\bottomrule
\end{tabular}
\caption{Sample Data Table}
\label{tab:sample}
\end{table}

\section{Lists}
\subsection{Bullet Points}
\begin{itemize}
  \item First item
  \item Second item
  \item Third item with sub-items:
    \begin{itemize}
      \item Sub-item 1
      \item Sub-item 2
    \end{itemize}
\end{itemize}

\subsection{Numbered List}
\begin{enumerate}
  \item First step
  \item Second step
  \item Third step
\end{enumerate}

\section{Citations and References}
When you use the citation search feature, references will be automatically managed. For example, you might cite a paper like this: \cite{example2024}. All citations will be automatically added to your \texttt{cite.bib} file.

\section{Conclusion}
This template demonstrates basic \LaTeX{} features including mathematical formulas, tables, and lists. The bibliography system is now set up to work seamlessly with Bibby AI's citation management. Feel free to modify it for your needs!

\bibliographystyle{plain}
\bibliography{cite}

\end{document}\documentclass{article}
\usepackage[utf8]{inputenc}
\usepackage{booktabs}
\usepackage{amsmath}
\usepackage{natbib}

\title{My Document}
\author{Your Name}
\date{\today}

\begin{document}

\maketitle

\section{Introduction}
This is a sample document demonstrating various \LaTeX{} features including mathematical equations, tables, and formatting. You can cite papers using the citation search feature, and they will be automatically added to your \texttt{cite.bib} file.

\section{Mathematical Expressions}
Here are some example equations:

\subsection{Quadratic Formula}
The quadratic formula is given by:
\[x = \frac{-b \pm \sqrt{b^2 - 4ac}}{2a}\]

\subsection{Matrix Example}
A sample matrix:
\[
\begin{pmatrix}
a_{11} & a_{12} & a_{13} \\
a_{21} & a_{22} & a_{23} \\
a_{31} & a_{32} & a_{33}
\end{pmatrix}
\]

\section{Tables}
Below is a sample table showing data:

\begin{table}[h]
\centering
\begin{tabular}{@{}lrrr@{}}
\toprule
\textbf{Category} & \textbf{Value 1} & \textbf{Value 2} & \textbf{Total} \\
\midrule
Item A & 10.5 & 20.3 & 30.8 \\
Item B & 15.2 & 25.1 & 40.3 \\
Item C & 22.1 & 12.7 & 34.8 \\
\bottomrule
\end{tabular}
\caption{Sample Data Table}
\label{tab:sample}
\end{table}

\section{Lists}
\subsection{Bullet Points}
\begin{itemize}
  \item First item
  \item Second item
  \item Third item with sub-items:
    \begin{itemize}
      \item Sub-item 1
      \item Sub-item 2
    \end{itemize}
\end{itemize}

\subsection{Numbered List}
\begin{enumerate}
  \item First step
  \item Second step
  \item Third step
\end{enumerate}

\section{Citations and References}
When you use the citation search feature, references will be automatically managed. For example, you might cite a paper like this: \cite{example2024}. All citations will be automatically added to your \texttt{cite.bib} file.

\section{Conclusion}
This template demonstrates basic \LaTeX{} features including mathematical formulas, tables, and lists. The bibliography system is now set up to work seamlessly with Bibby AI's citation management. Feel free to modify it for your needs!

\bibliographystyle{plain}
\bibliography{cite}

\end{document}\documentclass{article}
\usepackage[utf8]{inputenc}
\usepackage{booktabs}
\usepackage{amsmath}
\usepackage{natbib}

\title{My Document}
\author{Your Name}
\date{\today}

\begin{document}

\maketitle

\section{Introduction}
This is a sample document demonstrating various \LaTeX{} features including mathematical equations, tables, and formatting. You can cite papers using the citation search feature, and they will be automatically added to your \texttt{cite.bib} file.

\section{Mathematical Expressions}
Here are some example equations:

\subsection{Quadratic Formula}
The quadratic formula is given by:
\[x = \frac{-b \pm \sqrt{b^2 - 4ac}}{2a}\]

\subsection{Matrix Example}
A sample matrix:
\[
\begin{pmatrix}
a_{11} & a_{12} & a_{13} \\
a_{21} & a_{22} & a_{23} \\
a_{31} & a_{32} & a_{33}
\end{pmatrix}
\]

\section{Tables}
Below is a sample table showing data:

\begin{table}[h]
\centering
\begin{tabular}{@{}lrrr@{}}
\toprule
\textbf{Category} & \textbf{Value 1} & \textbf{Value 2} & \textbf{Total} \\
\midrule
Item A & 10.5 & 20.3 & 30.8 \\
Item B & 15.2 & 25.1 & 40.3 \\
Item C & 22.1 & 12.7 & 34.8 \\
\bottomrule
\end{tabular}
\caption{Sample Data Table}
\label{tab:sample}
\end{table}

\section{Lists}
\subsection{Bullet Points}
\begin{itemize}
  \item First item
  \item Second item
  \item Third item with sub-items:
    \begin{itemize}
      \item Sub-item 1
      \item Sub-item 2
    \end{itemize}
\end{itemize}

\subsection{Numbered List}
\begin{enumerate}
  \item First step
  \item Second step
  \item Third step
\end{enumerate}

\section{Citations and References}
When you use the citation search feature, references will be automatically managed. For example, you might cite a paper like this: \cite{example2024}. All citations will be automatically added to your \texttt{cite.bib} file.

\section{Conclusion}
This template demonstrates basic \LaTeX{} features including mathematical formulas, tables, and lists. The bibliography system is now set up to work seamlessly with Bibby AI's citation management. Feel free to modify it for your needs!

\bibliographystyle{plain}
\bibliography{cite}

\end{document}\documentclass{article}
\usepackage[utf8]{inputenc}
\usepackage{booktabs}
\usepackage{amsmath}
\usepackage{natbib}

\title{My Document}
\author{Your Name}
\date{\today}

\begin{document}

\maketitle

\section{Introduction}
This is a sample document demonstrating various \LaTeX{} features including mathematical equations, tables, and formatting. You can cite papers using the citation search feature, and they will be automatically added to your \texttt{cite.bib} file.

\section{Mathematical Expressions}
Here are some example equations:

\subsection{Quadratic Formula}
The quadratic formula is given by:
\[x = \frac{-b \pm \sqrt{b^2 - 4ac}}{2a}\]

\subsection{Matrix Example}
A sample matrix:
\[
\begin{pmatrix}
a_{11} & a_{12} & a_{13} \\
a_{21} & a_{22} & a_{23} \\
a_{31} & a_{32} & a_{33}
\end{pmatrix}
\]

\section{Tables}
Below is a sample table showing data:

\begin{table}[h]
\centering
\begin{tabular}{@{}lrrr@{}}
\toprule
\textbf{Category} & \textbf{Value 1} & \textbf{Value 2} & \textbf{Total} \\
\midrule
Item A & 10.5 & 20.3 & 30.8 \\
Item B & 15.2 & 25.1 & 40.3 \\
Item C & 22.1 & 12.7 & 34.8 \\
\bottomrule
\end{tabular}
\caption{Sample Data Table}
\label{tab:sample}
\end{table}

\section{Lists}
\subsection{Bullet Points}
\begin{itemize}
  \item First item
  \item Second item
  \item Third item with sub-items:
    \begin{itemize}
      \item Sub-item 1
      \item Sub-item 2
    \end{itemize}
\end{itemize}

\subsection{Numbered List}
\begin{enumerate}
  \item First step
  \item Second step
  \item Third step
\end{enumerate}

\section{Citations and References}
When you use the citation search feature, references will be automatically managed. For example, you might cite a paper like this: \cite{example2024}. All citations will be automatically added to your \texttt{cite.bib} file.

\section{Conclusion}
This template demonstrates basic \LaTeX{} features including mathematical formulas, tables, and lists. The bibliography system is now set up to work seamlessly with Bibby AI's citation management. Feel free to modify it for your needs!

\bibliographystyle{plain}
\bibliography{cite}

\end{document}\documentclass{article}
\usepackage[utf8]{inputenc}
\usepackage{booktabs}
\usepackage{amsmath}
\usepackage{natbib}

\title{My Document}
\author{Your Name}
\date{\today}

\begin{document}

\maketitle

\section{Introduction}
This is a sample document demonstrating various \LaTeX{} features including mathematical equations, tables, and formatting. You can cite papers using the citation search feature, and they will be automatically added to your \texttt{cite.bib} file.

\section{Mathematical Expressions}
Here are some example equations:

\subsection{Quadratic Formula}
The quadratic formula is given by:
\[x = \frac{-b \pm \sqrt{b^2 - 4ac}}{2a}\]

\subsection{Matrix Example}
A sample matrix:
\[
\begin{pmatrix}
a_{11} & a_{12} & a_{13} \\
a_{21} & a_{22} & a_{23} \\
a_{31} & a_{32} & a_{33}
\end{pmatrix}
\]

\section{Tables}
Below is a sample table showing data:

\begin{table}[h]
\centering
\begin{tabular}{@{}lrrr@{}}
\toprule
\textbf{Category} & \textbf{Value 1} & \textbf{Value 2} & \textbf{Total} \\
\midrule
Item A & 10.5 & 20.3 & 30.8 \\
Item B & 15.2 & 25.1 & 40.3 \\
Item C & 22.1 & 12.7 & 34.8 \\
\bottomrule
\end{tabular}
\caption{Sample Data Table}
\label{tab:sample}
\end{table}

\section{Lists}
\subsection{Bullet Points}
\begin{itemize}
  \item First item
  \item Second item
  \item Third item with sub-items:
    \begin{itemize}
      \item Sub-item 1
      \item Sub-item 2
    \end{itemize}
\end{itemize}

\subsection{Numbered List}
\begin{enumerate}
  \item First step
  \item Second step
  \item Third step
\end{enumerate}

\section{Citations and References}
When you use the citation search feature, references will be automatically managed. For example, you might cite a paper like this: \cite{example2024}. All citations will be automatically added to your \texttt{cite.bib} file.

\section{Conclusion}
This template demonstrates basic \LaTeX{} features including mathematical formulas, tables, and lists. The bibliography system is now set up to work seamlessly with Bibby AI's citation management. Feel free to modify it for your needs!

\bibliographystyle{plain}
\bibliography{cite}

\end{document}\documentclass{article}
\usepackage[utf8]{inputenc}
\usepackage{booktabs}
\usepackage{amsmath}
\usepackage{natbib}

\title{My Document}
\author{Your Name}
\date{\today}

\begin{document}

\maketitle

\section{Introduction}
This is a sample document demonstrating various \LaTeX{} features including mathematical equations, tables, and formatting. You can cite papers using the citation search feature, and they will be automatically added to your \texttt{cite.bib} file.

\section{Mathematical Expressions}
Here are some example equations:

\subsection{Quadratic Formula}
The quadratic formula is given by:
\[x = \frac{-b \pm \sqrt{b^2 - 4ac}}{2a}\]

\subsection{Matrix Example}
A sample matrix:
\[
\begin{pmatrix}
a_{11} & a_{12} & a_{13} \\
a_{21} & a_{22} & a_{23} \\
a_{31} & a_{32} & a_{33}
\end{pmatrix}
\]

\section{Tables}
Below is a sample table showing data:

\begin{table}[h]
\centering
\begin{tabular}{@{}lrrr@{}}
\toprule
\textbf{Category} & \textbf{Value 1} & \textbf{Value 2} & \textbf{Total} \\
\midrule
Item A & 10.5 & 20.3 & 30.8 \\
Item B & 15.2 & 25.1 & 40.3 \\
Item C & 22.1 & 12.7 & 34.8 \\
\bottomrule
\end{tabular}
\caption{Sample Data Table}
\label{tab:sample}
\end{table}

\section{Lists}
\subsection{Bullet Points}
\begin{itemize}
  \item First item
  \item Second item
  \item Third item with sub-items:
    \begin{itemize}
      \item Sub-item 1
      \item Sub-item 2
    \end{itemize}
\end{itemize}

\subsection{Numbered List}
\begin{enumerate}
  \item First step
  \item Second step
  \item Third step
\end{enumerate}

\section{Citations and References}
When you use the citation search feature, references will be automatically managed. For example, you might cite a paper like this: \cite{example2024}. All citations will be automatically added to your \texttt{cite.bib} file.

\section{Conclusion}
This template demonstrates basic \LaTeX{} features including mathematical formulas, tables, and lists. The bibliography system is now set up to work seamlessly with Bibby AI's citation management. Feel free to modify it for your needs!

\bibliographystyle{plain}
\bibliography{cite}

\end{document}\documentclass{article}
\usepackage[utf8]{inputenc}
\usepackage{booktabs}
\usepackage{amsmath}
\usepackage{natbib}

\title{My Document}
\author{Your Name}
\date{\today}

\begin{document}

\maketitle

\section{Introduction}
This is a sample document demonstrating various \LaTeX{} features including mathematical equations, tables, and formatting. You can cite papers using the citation search feature, and they will be automatically added to your \texttt{cite.bib} file.

\section{Mathematical Expressions}
Here are some example equations:

\subsection{Quadratic Formula}
The quadratic formula is given by:
\[x = \frac{-b \pm \sqrt{b^2 - 4ac}}{2a}\]

\subsection{Matrix Example}
A sample matrix:
\[
\begin{pmatrix}
a_{11} & a_{12} & a_{13} \\
a_{21} & a_{22} & a_{23} \\
a_{31} & a_{32} & a_{33}
\end{pmatrix}
\]

\section{Tables}
Below is a sample table showing data:

\begin{table}[h]
\centering
\begin{tabular}{@{}lrrr@{}}
\toprule
\textbf{Category} & \textbf{Value 1} & \textbf{Value 2} & \textbf{Total} \\
\midrule
Item A & 10.5 & 20.3 & 30.8 \\
Item B & 15.2 & 25.1 & 40.3 \\
Item C & 22.1 & 12.7 & 34.8 \\
\bottomrule
\end{tabular}
\caption{Sample Data Table}
\label{tab:sample}
\end{table}

\section{Lists}
\subsection{Bullet Points}
\begin{itemize}
  \item First item
  \item Second item
  \item Third item with sub-items:
    \begin{itemize}
      \item Sub-item 1
      \item Sub-item 2
    \end{itemize}
\end{itemize}

\subsection{Numbered List}
\begin{enumerate}
  \item First step
  \item Second step
  \item Third step
\end{enumerate}

\section{Citations and References}
When you use the citation search feature, references will be automatically managed. For example, you might cite a paper like this: \cite{example2024}. All citations will be automatically added to your \texttt{cite.bib} file.

\section{Conclusion}
This template demonstrates basic \LaTeX{} features including mathematical formulas, tables, and lists. The bibliography system is now set up to work seamlessly with Bibby AI's citation management. Feel free to modify it for your needs!

\bibliographystyle{plain}
\bibliography{cite}

\end{document}\documentclass{article}
\usepackage[utf8]{inputenc}
\usepackage{booktabs}
\usepackage{amsmath}
\usepackage{natbib}

\title{My Document}
\author{Your Name}
\date{\today}

\begin{document}

\maketitle

\section{Introduction}
This is a sample document demonstrating various \LaTeX{} features including mathematical equations, tables, and formatting. You can cite papers using the citation search feature, and they will be automatically added to your \texttt{cite.bib} file.

\section{Mathematical Expressions}
Here are some example equations:

\subsection{Quadratic Formula}
The quadratic formula is given by:
\[x = \frac{-b \pm \sqrt{b^2 - 4ac}}{2a}\]

\subsection{Matrix Example}
A sample matrix:
\[
\begin{pmatrix}
a_{11} & a_{12} & a_{13} \\
a_{21} & a_{22} & a_{23} \\
a_{31} & a_{32} & a_{33}
\end{pmatrix}
\]

\section{Tables}
Below is a sample table showing data:

\begin{table}[h]
\centering
\begin{tabular}{@{}lrrr@{}}
\toprule
\textbf{Category} & \textbf{Value 1} & \textbf{Value 2} & \textbf{Total} \\
\midrule
Item A & 10.5 & 20.3 & 30.8 \\
Item B & 15.2 & 25.1 & 40.3 \\
Item C & 22.1 & 12.7 & 34.8 \\
\bottomrule
\end{tabular}
\caption{Sample Data Table}
\label{tab:sample}
\end{table}

\section{Lists}
\subsection{Bullet Points}
\begin{itemize}
  \item First item
  \item Second item
  \item Third item with sub-items:
    \begin{itemize}
      \item Sub-item 1
      \item Sub-item 2
    \end{itemize}
\end{itemize}

\subsection{Numbered List}
\begin{enumerate}
  \item First step
  \item Second step
  \item Third step
\end{enumerate}

\section{Citations and References}
When you use the citation search feature, references will be automatically managed. For example, you might cite a paper like this: \cite{example2024}. All citations will be automatically added to your \texttt{cite.bib} file.

\section{Conclusion}
This template demonstrates basic \LaTeX{} features including mathematical formulas, tables, and lists. The bibliography system is now set up to work seamlessly with Bibby AI's citation management. Feel free to modify it for your needs!

\bibliographystyle{plain}
\bibliography{cite}

\end{document}\documentclass{article}
\usepackage[utf8]{inputenc}
\usepackage{booktabs}
\usepackage{amsmath}
\usepackage{natbib}

\title{My Document}
\author{Your Name}
\date{\today}

\begin{document}

\maketitle

\section{Introduction}
This is a sample document demonstrating various \LaTeX{} features including mathematical equations, tables, and formatting. You can cite papers using the citation search feature, and they will be automatically added to your \texttt{cite.bib} file.

\section{Mathematical Expressions}
Here are some example equations:

\subsection{Quadratic Formula}
The quadratic formula is given by:
\[x = \frac{-b \pm \sqrt{b^2 - 4ac}}{2a}\]

\subsection{Matrix Example}
A sample matrix:
\[
\begin{pmatrix}
a_{11} & a_{12} & a_{13} \\
a_{21} & a_{22} & a_{23} \\
a_{31} & a_{32} & a_{33}
\end{pmatrix}
\]

\section{Tables}
Below is a sample table showing data:

\begin{table}[h]
\centering
\begin{tabular}{@{}lrrr@{}}
\toprule
\textbf{Category} & \textbf{Value 1} & \textbf{Value 2} & \textbf{Total} \\
\midrule
Item A & 10.5 & 20.3 & 30.8 \\
Item B & 15.2 & 25.1 & 40.3 \\
Item C & 22.1 & 12.7 & 34.8 \\
\bottomrule
\end{tabular}
\caption{Sample Data Table}
\label{tab:sample}
\end{table}

\section{Lists}
\subsection{Bullet Points}
\begin{itemize}
  \item First item
  \item Second item
  \item Third item with sub-items:
    \begin{itemize}
      \item Sub-item 1
      \item Sub-item 2
    \end{itemize}
\end{itemize}

\subsection{Numbered List}
\begin{enumerate}
  \item First step
  \item Second step
  \item Third step
\end{enumerate}

\section{Citations and References}
When you use the citation search feature, references will be automatically managed. For example, you might cite a paper like this: \cite{example2024}. All citations will be automatically added to your \texttt{cite.bib} file.

\section{Conclusion}
This template demonstrates basic \LaTeX{} features including mathematical formulas, tables, and lists. The bibliography system is now set up to work seamlessly with Bibby AI's citation management. Feel free to modify it for your needs!

\bibliographystyle{plain}
\bibliography{cite}

\end{document}\documentclass{article}
\usepackage[utf8]{inputenc}
\usepackage{booktabs}
\usepackage{amsmath}
\usepackage{natbib}

\title{My Document}
\author{Your Name}
\date{\today}

\begin{document}

\maketitle

\section{Introduction}
This is a sample document demonstrating various \LaTeX{} features including mathematical equations, tables, and formatting. You can cite papers using the citation search feature, and they will be automatically added to your \texttt{cite.bib} file.

\section{Mathematical Expressions}
Here are some example equations:

\subsection{Quadratic Formula}
The quadratic formula is given by:
\[x = \frac{-b \pm \sqrt{b^2 - 4ac}}{2a}\]

\subsection{Matrix Example}
A sample matrix:
\[
\begin{pmatrix}
a_{11} & a_{12} & a_{13} \\
a_{21} & a_{22} & a_{23} \\
a_{31} & a_{32} & a_{33}
\end{pmatrix}
\]

\section{Tables}
Below is a sample table showing data:

\begin{table}[h]
\centering
\begin{tabular}{@{}lrrr@{}}
\toprule
\textbf{Category} & \textbf{Value 1} & \textbf{Value 2} & \textbf{Total} \\
\midrule
Item A & 10.5 & 20.3 & 30.8 \\
Item B & 15.2 & 25.1 & 40.3 \\
Item C & 22.1 & 12.7 & 34.8 \\
\bottomrule
\end{tabular}
\caption{Sample Data Table}
\label{tab:sample}
\end{table}

\section{Lists}
\subsection{Bullet Points}
\begin{itemize}
  \item First item
  \item Second item
  \item Third item with sub-items:
    \begin{itemize}
      \item Sub-item 1
      \item Sub-item 2
    \end{itemize}
\end{itemize}

\subsection{Numbered List}
\begin{enumerate}
  \item First step
  \item Second step
  \item Third step
\end{enumerate}

\section{Citations and References}
When you use the citation search feature, references will be automatically managed. For example, you might cite a paper like this: \cite{example2024}. All citations will be automatically added to your \texttt{cite.bib} file.

\section{Conclusion}
This template demonstrates basic \LaTeX{} features including mathematical formulas, tables, and lists. The bibliography system is now set up to work seamlessly with Bibby AI's citation management. Feel free to modify it for your needs!

\bibliographystyle{plain}
\bibliography{cite}

\end{document}\documentclass{article}
\usepackage[utf8]{inputenc}
\usepackage{booktabs}
\usepackage{amsmath}
\usepackage{natbib}

\title{My Document}
\author{Your Name}
\date{\today}

\begin{document}

\maketitle

\section{Introduction}
This is a sample document demonstrating various \LaTeX{} features including mathematical equations, tables, and formatting. You can cite papers using the citation search feature, and they will be automatically added to your \texttt{cite.bib} file.

\section{Mathematical Expressions}
Here are some example equations:

\subsection{Quadratic Formula}
The quadratic formula is given by:
\[x = \frac{-b \pm \sqrt{b^2 - 4ac}}{2a}\]

\subsection{Matrix Example}
A sample matrix:
\[
\begin{pmatrix}
a_{11} & a_{12} & a_{13} \\
a_{21} & a_{22} & a_{23} \\
a_{31} & a_{32} & a_{33}
\end{pmatrix}
\]

\section{Tables}
Below is a sample table showing data:

\begin{table}[h]
\centering
\begin{tabular}{@{}lrrr@{}}
\toprule
\textbf{Category} & \textbf{Value 1} & \textbf{Value 2} & \textbf{Total} \\
\midrule
Item A & 10.5 & 20.3 & 30.8 \\
Item B & 15.2 & 25.1 & 40.3 \\
Item C & 22.1 & 12.7 & 34.8 \\
\bottomrule
\end{tabular}
\caption{Sample Data Table}
\label{tab:sample}
\end{table}

\section{Lists}
\subsection{Bullet Points}
\begin{itemize}
  \item First item
  \item Second item
  \item Third item with sub-items:
    \begin{itemize}
      \item Sub-item 1
      \item Sub-item 2
    \end{itemize}
\end{itemize}

\subsection{Numbered List}
\begin{enumerate}
  \item First step
  \item Second step
  \item Third step
\end{enumerate}

\section{Citations and References}
When you use the citation search feature, references will be automatically managed. For example, you might cite a paper like this: \cite{example2024}. All citations will be automatically added to your \texttt{cite.bib} file.

\section{Conclusion}
This template demonstrates basic \LaTeX{} features including mathematical formulas, tables, and lists. The bibliography system is now set up to work seamlessly with Bibby AI's citation management. Feel free to modify it for your needs!

\bibliographystyle{plain}
\bibliography{cite}

\end{document}\documentclass{article}
\usepackage[utf8]{inputenc}
\usepackage{booktabs}
\usepackage{amsmath}
\usepackage{natbib}

\title{My Document}
\author{Your Name}
\date{\today}

\begin{document}

\maketitle

\section{Introduction}
This is a sample document demonstrating various \LaTeX{} features including mathematical equations, tables, and formatting. You can cite papers using the citation search feature, and they will be automatically added to your \texttt{cite.bib} file.

\section{Mathematical Expressions}
Here are some example equations:

\subsection{Quadratic Formula}
The quadratic formula is given by:
\[x = \frac{-b \pm \sqrt{b^2 - 4ac}}{2a}\]

\subsection{Matrix Example}
A sample matrix:
\[
\begin{pmatrix}
a_{11} & a_{12} & a_{13} \\
a_{21} & a_{22} & a_{23} \\
a_{31} & a_{32} & a_{33}
\end{pmatrix}
\]

\section{Tables}
Below is a sample table showing data:

\begin{table}[h]
\centering
\begin{tabular}{@{}lrrr@{}}
\toprule
\textbf{Category} & \textbf{Value 1} & \textbf{Value 2} & \textbf{Total} \\
\midrule
Item A & 10.5 & 20.3 & 30.8 \\
Item B & 15.2 & 25.1 & 40.3 \\
Item C & 22.1 & 12.7 & 34.8 \\
\bottomrule
\end{tabular}
\caption{Sample Data Table}
\label{tab:sample}
\end{table}

\section{Lists}
\subsection{Bullet Points}
\begin{itemize}
  \item First item
  \item Second item
  \item Third item with sub-items:
    \begin{itemize}
      \item Sub-item 1
      \item Sub-item 2
    \end{itemize}
\end{itemize}

\subsection{Numbered List}
\begin{enumerate}
  \item First step
  \item Second step
  \item Third step
\end{enumerate}

\section{Citations and References}
When you use the citation search feature, references will be automatically managed. For example, you might cite a paper like this: \cite{example2024}. All citations will be automatically added to your \texttt{cite.bib} file.

\section{Conclusion}
This template demonstrates basic \LaTeX{} features including mathematical formulas, tables, and lists. The bibliography system is now set up to work seamlessly with Bibby AI's citation management. Feel free to modify it for your needs!

\bibliographystyle{plain}
\bibliography{cite}

\end{document}\documentclass{article}
\usepackage[utf8]{inputenc}
\usepackage{booktabs}
\usepackage{amsmath}
\usepackage{natbib}

\title{My Document}
\author{Your Name}
\date{\today}

\begin{document}

\maketitle

\section{Introduction}
This is a sample document demonstrating various \LaTeX{} features including mathematical equations, tables, and formatting. You can cite papers using the citation search feature, and they will be automatically added to your \texttt{cite.bib} file.

\section{Mathematical Expressions}
Here are some example equations:

\subsection{Quadratic Formula}
The quadratic formula is given by:
\[x = \frac{-b \pm \sqrt{b^2 - 4ac}}{2a}\]

\subsection{Matrix Example}
A sample matrix:
\[
\begin{pmatrix}
a_{11} & a_{12} & a_{13} \\
a_{21} & a_{22} & a_{23} \\
a_{31} & a_{32} & a_{33}
\end{pmatrix}
\]

\section{Tables}
Below is a sample table showing data:

\begin{table}[h]
\centering
\begin{tabular}{@{}lrrr@{}}
\toprule
\textbf{Category} & \textbf{Value 1} & \textbf{Value 2} & \textbf{Total} \\
\midrule
Item A & 10.5 & 20.3 & 30.8 \\
Item B & 15.2 & 25.1 & 40.3 \\
Item C & 22.1 & 12.7 & 34.8 \\
\bottomrule
\end{tabular}
\caption{Sample Data Table}
\label{tab:sample}
\end{table}

\section{Lists}
\subsection{Bullet Points}
\begin{itemize}
  \item First item
  \item Second item
  \item Third item with sub-items:
    \begin{itemize}
      \item Sub-item 1
      \item Sub-item 2
    \end{itemize}
\end{itemize}

\subsection{Numbered List}
\begin{enumerate}
  \item First step
  \item Second step
  \item Third step
\end{enumerate}

\section{Citations and References}
When you use the citation search feature, references will be automatically managed. For example, you might cite a paper like this: \cite{example2024}. All citations will be automatically added to your \texttt{cite.bib} file.

\section{Conclusion}
This template demonstrates basic \LaTeX{} features including mathematical formulas, tables, and lists. The bibliography system is now set up to work seamlessly with Bibby AI's citation management. Feel free to modify it for your needs!

\bibliographystyle{plain}
\bibliography{cite}

\end{document}\documentclass{article}
\usepackage[utf8]{inputenc}
\usepackage{booktabs}
\usepackage{amsmath}
\usepackage{natbib}

\title{My Document}
\author{Your Name}
\date{\today}

\begin{document}

\maketitle

\section{Introduction}
This is a sample document demonstrating various \LaTeX{} features including mathematical equations, tables, and formatting. You can cite papers using the citation search feature, and they will be automatically added to your \texttt{cite.bib} file.

\section{Mathematical Expressions}
Here are some example equations:

\subsection{Quadratic Formula}
The quadratic formula is given by:
\[x = \frac{-b \pm \sqrt{b^2 - 4ac}}{2a}\]

\subsection{Matrix Example}
A sample matrix:
\[
\begin{pmatrix}
a_{11} & a_{12} & a_{13} \\
a_{21} & a_{22} & a_{23} \\
a_{31} & a_{32} & a_{33}
\end{pmatrix}
\]

\section{Tables}
Below is a sample table showing data:

\begin{table}[h]
\centering
\begin{tabular}{@{}lrrr@{}}
\toprule
\textbf{Category} & \textbf{Value 1} & \textbf{Value 2} & \textbf{Total} \\
\midrule
Item A & 10.5 & 20.3 & 30.8 \\
Item B & 15.2 & 25.1 & 40.3 \\
Item C & 22.1 & 12.7 & 34.8 \\
\bottomrule
\end{tabular}
\caption{Sample Data Table}
\label{tab:sample}
\end{table}

\section{Lists}
\subsection{Bullet Points}
\begin{itemize}
  \item First item
  \item Second item
  \item Third item with sub-items:
    \begin{itemize}
      \item Sub-item 1
      \item Sub-item 2
    \end{itemize}
\end{itemize}

\subsection{Numbered List}
\begin{enumerate}
  \item First step
  \item Second step
  \item Third step
\end{enumerate}

\section{Citations and References}
When you use the citation search feature, references will be automatically managed. For example, you might cite a paper like this: \cite{example2024}. All citations will be automatically added to your \texttt{cite.bib} file.

\section{Conclusion}
This template demonstrates basic \LaTeX{} features including mathematical formulas, tables, and lists. The bibliography system is now set up to work seamlessly with Bibby AI's citation management. Feel free to modify it for your needs!

\bibliographystyle{plain}
\bibliography{cite}

\end{document}\documentclass{article}
\usepackage[utf8]{inputenc}
\usepackage{booktabs}
\usepackage{amsmath}
\usepackage{natbib}

\title{My Document}
\author{Your Name}
\date{\today}

\begin{document}

\maketitle

\section{Introduction}
This is a sample document demonstrating various \LaTeX{} features including mathematical equations, tables, and formatting. You can cite papers using the citation search feature, and they will be automatically added to your \texttt{cite.bib} file.

\section{Mathematical Expressions}
Here are some example equations:

\subsection{Quadratic Formula}
The quadratic formula is given by:
\[x = \frac{-b \pm \sqrt{b^2 - 4ac}}{2a}\]

\subsection{Matrix Example}
A sample matrix:
\[
\begin{pmatrix}
a_{11} & a_{12} & a_{13} \\
a_{21} & a_{22} & a_{23} \\
a_{31} & a_{32} & a_{33}
\end{pmatrix}
\]

\section{Tables}
Below is a sample table showing data:

\begin{table}[h]
\centering
\begin{tabular}{@{}lrrr@{}}
\toprule
\textbf{Category} & \textbf{Value 1} & \textbf{Value 2} & \textbf{Total} \\
\midrule
Item A & 10.5 & 20.3 & 30.8 \\
Item B & 15.2 & 25.1 & 40.3 \\
Item C & 22.1 & 12.7 & 34.8 \\
\bottomrule
\end{tabular}
\caption{Sample Data Table}
\label{tab:sample}
\end{table}

\section{Lists}
\subsection{Bullet Points}
\begin{itemize}
  \item First item
  \item Second item
  \item Third item with sub-items:
    \begin{itemize}
      \item Sub-item 1
      \item Sub-item 2
    \end{itemize}
\end{itemize}

\subsection{Numbered List}
\begin{enumerate}
  \item First step
  \item Second step
  \item Third step
\end{enumerate}

\section{Citations and References}
When you use the citation search feature, references will be automatically managed. For example, you might cite a paper like this: \cite{example2024}. All citations will be automatically added to your \texttt{cite.bib} file.

\section{Conclusion}
This template demonstrates basic \LaTeX{} features including mathematical formulas, tables, and lists. The bibliography system is now set up to work seamlessly with Bibby AI's citation management. Feel free to modify it for your needs!

\bibliographystyle{plain}
\bibliography{cite}

\end{document}\documentclass{article}
\usepackage[utf8]{inputenc}
\usepackage{booktabs}
\usepackage{amsmath}
\usepackage{natbib}

\title{My Document}
\author{Your Name}
\date{\today}

\begin{document}

\maketitle

\section{Introduction}
This is a sample document demonstrating various \LaTeX{} features including mathematical equations, tables, and formatting. You can cite papers using the citation search feature, and they will be automatically added to your \texttt{cite.bib} file.

\section{Mathematical Expressions}
Here are some example equations:

\subsection{Quadratic Formula}
The quadratic formula is given by:
\[x = \frac{-b \pm \sqrt{b^2 - 4ac}}{2a}\]

\subsection{Matrix Example}
A sample matrix:
\[
\begin{pmatrix}
a_{11} & a_{12} & a_{13} \\
a_{21} & a_{22} & a_{23} \\
a_{31} & a_{32} & a_{33}
\end{pmatrix}
\]

\section{Tables}
Below is a sample table showing data:

\begin{table}[h]
\centering
\begin{tabular}{@{}lrrr@{}}
\toprule
\textbf{Category} & \textbf{Value 1} & \textbf{Value 2} & \textbf{Total} \\
\midrule
Item A & 10.5 & 20.3 & 30.8 \\
Item B & 15.2 & 25.1 & 40.3 \\
Item C & 22.1 & 12.7 & 34.8 \\
\bottomrule
\end{tabular}
\caption{Sample Data Table}
\label{tab:sample}
\end{table}

\section{Lists}
\subsection{Bullet Points}
\begin{itemize}
  \item First item
  \item Second item
  \item Third item with sub-items:
    \begin{itemize}
      \item Sub-item 1
      \item Sub-item 2
    \end{itemize}
\end{itemize}

\subsection{Numbered List}
\begin{enumerate}
  \item First step
  \item Second step
  \item Third step
\end{enumerate}

\section{Citations and References}
When you use the citation search feature, references will be automatically managed. For example, you might cite a paper like this: \cite{example2024}. All citations will be automatically added to your \texttt{cite.bib} file.

\section{Conclusion}
This template demonstrates basic \LaTeX{} features including mathematical formulas, tables, and lists. The bibliography system is now set up to work seamlessly with Bibby AI's citation management. Feel free to modify it for your needs!

\bibliographystyle{plain}
\bibliography{cite}

\end{document}\documentclass{article}
\usepackage[utf8]{inputenc}
\usepackage{booktabs}
\usepackage{amsmath}
\usepackage{natbib}

\title{My Document}
\author{Your Name}
\date{\today}

\begin{document}

\maketitle

\section{Introduction}
This is a sample document demonstrating various \LaTeX{} features including mathematical equations, tables, and formatting. You can cite papers using the citation search feature, and they will be automatically added to your \texttt{cite.bib} file.

\section{Mathematical Expressions}
Here are some example equations:

\subsection{Quadratic Formula}
The quadratic formula is given by:
\[x = \frac{-b \pm \sqrt{b^2 - 4ac}}{2a}\]

\subsection{Matrix Example}
A sample matrix:
\[
\begin{pmatrix}
a_{11} & a_{12} & a_{13} \\
a_{21} & a_{22} & a_{23} \\
a_{31} & a_{32} & a_{33}
\end{pmatrix}
\]

\section{Tables}
Below is a sample table showing data:

\begin{table}[h]
\centering
\begin{tabular}{@{}lrrr@{}}
\toprule
\textbf{Category} & \textbf{Value 1} & \textbf{Value 2} & \textbf{Total} \\
\midrule
Item A & 10.5 & 20.3 & 30.8 \\
Item B & 15.2 & 25.1 & 40.3 \\
Item C & 22.1 & 12.7 & 34.8 \\
\bottomrule
\end{tabular}
\caption{Sample Data Table}
\label{tab:sample}
\end{table}

\section{Lists}
\subsection{Bullet Points}
\begin{itemize}
  \item First item
  \item Second item
  \item Third item with sub-items:
    \begin{itemize}
      \item Sub-item 1
      \item Sub-item 2
    \end{itemize}
\end{itemize}

\subsection{Numbered List}
\begin{enumerate}
  \item First step
  \item Second step
  \item Third step
\end{enumerate}

\section{Citations and References}
When you use the citation search feature, references will be automatically managed. For example, you might cite a paper like this: \cite{example2024}. All citations will be automatically added to your \texttt{cite.bib} file.

\section{Conclusion}
This template demonstrates basic \LaTeX{} features including mathematical formulas, tables, and lists. The bibliography system is now set up to work seamlessly with Bibby AI's citation management. Feel free to modify it for your needs!

\bibliographystyle{plain}
\bibliography{cite}

\end{document}\documentclass{article}
\usepackage[utf8]{inputenc}
\usepackage{booktabs}
\usepackage{amsmath}
\usepackage{natbib}

\title{My Document}
\author{Your Name}
\date{\today}

\begin{document}

\maketitle

\section{Introduction}
This is a sample document demonstrating various \LaTeX{} features including mathematical equations, tables, and formatting. You can cite papers using the citation search feature, and they will be automatically added to your \texttt{cite.bib} file.

\section{Mathematical Expressions}
Here are some example equations:

\subsection{Quadratic Formula}
The quadratic formula is given by:
\[x = \frac{-b \pm \sqrt{b^2 - 4ac}}{2a}\]

\subsection{Matrix Example}
A sample matrix:
\[
\begin{pmatrix}
a_{11} & a_{12} & a_{13} \\
a_{21} & a_{22} & a_{23} \\
a_{31} & a_{32} & a_{33}
\end{pmatrix}
\]

\section{Tables}
Below is a sample table showing data:

\begin{table}[h]
\centering
\begin{tabular}{@{}lrrr@{}}
\toprule
\textbf{Category} & \textbf{Value 1} & \textbf{Value 2} & \textbf{Total} \\
\midrule
Item A & 10.5 & 20.3 & 30.8 \\
Item B & 15.2 & 25.1 & 40.3 \\
Item C & 22.1 & 12.7 & 34.8 \\
\bottomrule
\end{tabular}
\caption{Sample Data Table}
\label{tab:sample}
\end{table}

\section{Lists}
\subsection{Bullet Points}
\begin{itemize}
  \item First item
  \item Second item
  \item Third item with sub-items:
    \begin{itemize}
      \item Sub-item 1
      \item Sub-item 2
    \end{itemize}
\end{itemize}

\subsection{Numbered List}
\begin{enumerate}
  \item First step
  \item Second step
  \item Third step
\end{enumerate}

\section{Citations and References}
When you use the citation search feature, references will be automatically managed. For example, you might cite a paper like this: \cite{example2024}. All citations will be automatically added to your \texttt{cite.bib} file.

\section{Conclusion}
This template demonstrates basic \LaTeX{} features including mathematical formulas, tables, and lists. The bibliography system is now set up to work seamlessly with Bibby AI's citation management. Feel free to modify it for your needs!

\bibliographystyle{plain}
\bibliography{cite}

\end{document}\documentclass{article}
\usepackage[utf8]{inputenc}
\usepackage{booktabs}
\usepackage{amsmath}
\usepackage{natbib}

\title{My Document}
\author{Your Name}
\date{\today}

\begin{document}

\maketitle

\section{Introduction}
This is a sample document demonstrating various \LaTeX{} features including mathematical equations, tables, and formatting. You can cite papers using the citation search feature, and they will be automatically added to your \texttt{cite.bib} file.

\section{Mathematical Expressions}
Here are some example equations:

\subsection{Quadratic Formula}
The quadratic formula is given by:
\[x = \frac{-b \pm \sqrt{b^2 - 4ac}}{2a}\]

\subsection{Matrix Example}
A sample matrix:
\[
\begin{pmatrix}
a_{11} & a_{12} & a_{13} \\
a_{21} & a_{22} & a_{23} \\
a_{31} & a_{32} & a_{33}
\end{pmatrix}
\]

\section{Tables}
Below is a sample table showing data:

\begin{table}[h]
\centering
\begin{tabular}{@{}lrrr@{}}
\toprule
\textbf{Category} & \textbf{Value 1} & \textbf{Value 2} & \textbf{Total} \\
\midrule
Item A & 10.5 & 20.3 & 30.8 \\
Item B & 15.2 & 25.1 & 40.3 \\
Item C & 22.1 & 12.7 & 34.8 \\
\bottomrule
\end{tabular}
\caption{Sample Data Table}
\label{tab:sample}
\end{table}

\section{Lists}
\subsection{Bullet Points}
\begin{itemize}
  \item First item
  \item Second item
  \item Third item with sub-items:
    \begin{itemize}
      \item Sub-item 1
      \item Sub-item 2
    \end{itemize}
\end{itemize}

\subsection{Numbered List}
\begin{enumerate}
  \item First step
  \item Second step
  \item Third step
\end{enumerate}

\section{Citations and References}
When you use the citation search feature, references will be automatically managed. For example, you might cite a paper like this: \cite{example2024}. All citations will be automatically added to your \texttt{cite.bib} file.

\section{Conclusion}
This template demonstrates basic \LaTeX{} features including mathematical formulas, tables, and lists. The bibliography system is now set up to work seamlessly with Bibby AI's citation management. Feel free to modify it for your needs!

\bibliographystyle{plain}
\bibliography{cite}

\end{document}\documentclass{article}
\usepackage[utf8]{inputenc}
\usepackage{booktabs}
\usepackage{amsmath}
\usepackage{natbib}

\title{My Document}
\author{Your Name}
\date{\today}

\begin{document}

\maketitle

\section{Introduction}
This is a sample document demonstrating various \LaTeX{} features including mathematical equations, tables, and formatting. You can cite papers using the citation search feature, and they will be automatically added to your \texttt{cite.bib} file.

\section{Mathematical Expressions}
Here are some example equations:

\subsection{Quadratic Formula}
The quadratic formula is given by:
\[x = \frac{-b \pm \sqrt{b^2 - 4ac}}{2a}\]

\subsection{Matrix Example}
A sample matrix:
\[
\begin{pmatrix}
a_{11} & a_{12} & a_{13} \\
a_{21} & a_{22} & a_{23} \\
a_{31} & a_{32} & a_{33}
\end{pmatrix}
\]

\section{Tables}
Below is a sample table showing data:

\begin{table}[h]
\centering
\begin{tabular}{@{}lrrr@{}}
\toprule
\textbf{Category} & \textbf{Value 1} & \textbf{Value 2} & \textbf{Total} \\
\midrule
Item A & 10.5 & 20.3 & 30.8 \\
Item B & 15.2 & 25.1 & 40.3 \\
Item C & 22.1 & 12.7 & 34.8 \\
\bottomrule
\end{tabular}
\caption{Sample Data Table}
\label{tab:sample}
\end{table}

\section{Lists}
\subsection{Bullet Points}
\begin{itemize}
  \item First item
  \item Second item
  \item Third item with sub-items:
    \begin{itemize}
      \item Sub-item 1
      \item Sub-item 2
    \end{itemize}
\end{itemize}

\subsection{Numbered List}
\begin{enumerate}
  \item First step
  \item Second step
  \item Third step
\end{enumerate}

\section{Citations and References}
When you use the citation search feature, references will be automatically managed. For example, you might cite a paper like this: \cite{example2024}. All citations will be automatically added to your \texttt{cite.bib} file.

\section{Conclusion}
This template demonstrates basic \LaTeX{} features including mathematical formulas, tables, and lists. The bibliography system is now set up to work seamlessly with Bibby AI's citation management. Feel free to modify it for your needs!

\bibliographystyle{plain}
\bibliography{cite}

\end{document}\documentclass{article}
\usepackage[utf8]{inputenc}
\usepackage{booktabs}
\usepackage{amsmath}
\usepackage{natbib}

\title{My Document}
\author{Your Name}
\date{\today}

\begin{document}

\maketitle

\section{Introduction}
This is a sample document demonstrating various \LaTeX{} features including mathematical equations, tables, and formatting. You can cite papers using the citation search feature, and they will be automatically added to your \texttt{cite.bib} file.

\section{Mathematical Expressions}
Here are some example equations:

\subsection{Quadratic Formula}
The quadratic formula is given by:
\[x = \frac{-b \pm \sqrt{b^2 - 4ac}}{2a}\]

\subsection{Matrix Example}
A sample matrix:
\[
\begin{pmatrix}
a_{11} & a_{12} & a_{13} \\
a_{21} & a_{22} & a_{23} \\
a_{31} & a_{32} & a_{33}
\end{pmatrix}
\]

\section{Tables}
Below is a sample table showing data:

\begin{table}[h]
\centering
\begin{tabular}{@{}lrrr@{}}
\toprule
\textbf{Category} & \textbf{Value 1} & \textbf{Value 2} & \textbf{Total} \\
\midrule
Item A & 10.5 & 20.3 & 30.8 \\
Item B & 15.2 & 25.1 & 40.3 \\
Item C & 22.1 & 12.7 & 34.8 \\
\bottomrule
\end{tabular}
\caption{Sample Data Table}
\label{tab:sample}
\end{table}

\section{Lists}
\subsection{Bullet Points}
\begin{itemize}
  \item First item
  \item Second item
  \item Third item with sub-items:
    \begin{itemize}
      \item Sub-item 1
      \item Sub-item 2
    \end{itemize}
\end{itemize}

\subsection{Numbered List}
\begin{enumerate}
  \item First step
  \item Second step
  \item Third step
\end{enumerate}

\section{Citations and References}
When you use the citation search feature, references will be automatically managed. For example, you might cite a paper like this: \cite{example2024}. All citations will be automatically added to your \texttt{cite.bib} file.

\section{Conclusion}
This template demonstrates basic \LaTeX{} features including mathematical formulas, tables, and lists. The bibliography system is now set up to work seamlessly with Bibby AI's citation management. Feel free to modify it for your needs!

\bibliographystyle{plain}
\bibliography{cite}

\end{document}\documentclass{article}
\usepackage[utf8]{inputenc}
\usepackage{booktabs}
\usepackage{amsmath}
\usepackage{natbib}

\title{My Document}
\author{Your Name}
\date{\today}

\begin{document}

\maketitle

\section{Introduction}
This is a sample document demonstrating various \LaTeX{} features including mathematical equations, tables, and formatting. You can cite papers using the citation search feature, and they will be automatically added to your \texttt{cite.bib} file.

\section{Mathematical Expressions}
Here are some example equations:

\subsection{Quadratic Formula}
The quadratic formula is given by:
\[x = \frac{-b \pm \sqrt{b^2 - 4ac}}{2a}\]

\subsection{Matrix Example}
A sample matrix:
\[
\begin{pmatrix}
a_{11} & a_{12} & a_{13} \\
a_{21} & a_{22} & a_{23} \\
a_{31} & a_{32} & a_{33}
\end{pmatrix}
\]

\section{Tables}
Below is a sample table showing data:

\begin{table}[h]
\centering
\begin{tabular}{@{}lrrr@{}}
\toprule
\textbf{Category} & \textbf{Value 1} & \textbf{Value 2} & \textbf{Total} \\
\midrule
Item A & 10.5 & 20.3 & 30.8 \\
Item B & 15.2 & 25.1 & 40.3 \\
Item C & 22.1 & 12.7 & 34.8 \\
\bottomrule
\end{tabular}
\caption{Sample Data Table}
\label{tab:sample}
\end{table}

\section{Lists}
\subsection{Bullet Points}
\begin{itemize}
  \item First item
  \item Second item
  \item Third item with sub-items:
    \begin{itemize}
      \item Sub-item 1
      \item Sub-item 2
    \end{itemize}
\end{itemize}

\subsection{Numbered List}
\begin{enumerate}
  \item First step
  \item Second step
  \item Third step
\end{enumerate}

\section{Citations and References}
When you use the citation search feature, references will be automatically managed. For example, you might cite a paper like this: \cite{example2024}. All citations will be automatically added to your \texttt{cite.bib} file.

\section{Conclusion}
This template demonstrates basic \LaTeX{} features including mathematical formulas, tables, and lists. The bibliography system is now set up to work seamlessly with Bibby AI's citation management. Feel free to modify it for your needs!

\bibliographystyle{plain}
\bibliography{cite}

\end{document}\documentclass{article}
\usepackage[utf8]{inputenc}
\usepackage{booktabs}
\usepackage{amsmath}
\usepackage{natbib}

\title{My Document}
\author{Your Name}
\date{\today}

\begin{document}

\maketitle

\section{Introduction}
This is a sample document demonstrating various \LaTeX{} features including mathematical equations, tables, and formatting. You can cite papers using the citation search feature, and they will be automatically added to your \texttt{cite.bib} file.

\section{Mathematical Expressions}
Here are some example equations:

\subsection{Quadratic Formula}
The quadratic formula is given by:
\[x = \frac{-b \pm \sqrt{b^2 - 4ac}}{2a}\]

\subsection{Matrix Example}
A sample matrix:
\[
\begin{pmatrix}
a_{11} & a_{12} & a_{13} \\
a_{21} & a_{22} & a_{23} \\
a_{31} & a_{32} & a_{33}
\end{pmatrix}
\]

\section{Tables}
Below is a sample table showing data:

\begin{table}[h]
\centering
\begin{tabular}{@{}lrrr@{}}
\toprule
\textbf{Category} & \textbf{Value 1} & \textbf{Value 2} & \textbf{Total} \\
\midrule
Item A & 10.5 & 20.3 & 30.8 \\
Item B & 15.2 & 25.1 & 40.3 \\
Item C & 22.1 & 12.7 & 34.8 \\
\bottomrule
\end{tabular}
\caption{Sample Data Table}
\label{tab:sample}
\end{table}

\section{Lists}
\subsection{Bullet Points}
\begin{itemize}
  \item First item
  \item Second item
  \item Third item with sub-items:
    \begin{itemize}
      \item Sub-item 1
      \item Sub-item 2
    \end{itemize}
\end{itemize}

\subsection{Numbered List}
\begin{enumerate}
  \item First step
  \item Second step
  \item Third step
\end{enumerate}

\section{Citations and References}
When you use the citation search feature, references will be automatically managed. For example, you might cite a paper like this: \cite{example2024}. All citations will be automatically added to your \texttt{cite.bib} file.

\section{Conclusion}
This template demonstrates basic \LaTeX{} features including mathematical formulas, tables, and lists. The bibliography system is now set up to work seamlessly with Bibby AI's citation management. Feel free to modify it for your needs!

\bibliographystyle{plain}
\bibliography{cite}

\end{document}\documentclass{article}
\usepackage[utf8]{inputenc}
\usepackage{booktabs}
\usepackage{amsmath}
\usepackage{natbib}

\title{My Document}
\author{Your Name}
\date{\today}

\begin{document}

\maketitle

\section{Introduction}
This is a sample document demonstrating various \LaTeX{} features including mathematical equations, tables, and formatting. You can cite papers using the citation search feature, and they will be automatically added to your \texttt{cite.bib} file.

\section{Mathematical Expressions}
Here are some example equations:

\subsection{Quadratic Formula}
The quadratic formula is given by:
\[x = \frac{-b \pm \sqrt{b^2 - 4ac}}{2a}\]

\subsection{Matrix Example}
A sample matrix:
\[
\begin{pmatrix}
a_{11} & a_{12} & a_{13} \\
a_{21} & a_{22} & a_{23} \\
a_{31} & a_{32} & a_{33}
\end{pmatrix}
\]

\section{Tables}
Below is a sample table showing data:

\begin{table}[h]
\centering
\begin{tabular}{@{}lrrr@{}}
\toprule
\textbf{Category} & \textbf{Value 1} & \textbf{Value 2} & \textbf{Total} \\
\midrule
Item A & 10.5 & 20.3 & 30.8 \\
Item B & 15.2 & 25.1 & 40.3 \\
Item C & 22.1 & 12.7 & 34.8 \\
\bottomrule
\end{tabular}
\caption{Sample Data Table}
\label{tab:sample}
\end{table}

\section{Lists}
\subsection{Bullet Points}
\begin{itemize}
  \item First item
  \item Second item
  \item Third item with sub-items:
    \begin{itemize}
      \item Sub-item 1
      \item Sub-item 2
    \end{itemize}
\end{itemize}

\subsection{Numbered List}
\begin{enumerate}
  \item First step
  \item Second step
  \item Third step
\end{enumerate}

\section{Citations and References}
When you use the citation search feature, references will be automatically managed. For example, you might cite a paper like this: \cite{example2024}. All citations will be automatically added to your \texttt{cite.bib} file.

\section{Conclusion}
This template demonstrates basic \LaTeX{} features including mathematical formulas, tables, and lists. The bibliography system is now set up to work seamlessly with Bibby AI's citation management. Feel free to modify it for your needs!

\bibliographystyle{plain}
\bibliography{cite}

\end{document}\documentclass{article}
\usepackage[utf8]{inputenc}
\usepackage{booktabs}
\usepackage{amsmath}
\usepackage{natbib}

\title{My Document}
\author{Your Name}
\date{\today}

\begin{document}

\maketitle

\section{Introduction}
This is a sample document demonstrating various \LaTeX{} features including mathematical equations, tables, and formatting. You can cite papers using the citation search feature, and they will be automatically added to your \texttt{cite.bib} file.

\section{Mathematical Expressions}
Here are some example equations:

\subsection{Quadratic Formula}
The quadratic formula is given by:
\[x = \frac{-b \pm \sqrt{b^2 - 4ac}}{2a}\]

\subsection{Matrix Example}
A sample matrix:
\[
\begin{pmatrix}
a_{11} & a_{12} & a_{13} \\
a_{21} & a_{22} & a_{23} \\
a_{31} & a_{32} & a_{33}
\end{pmatrix}
\]

\section{Tables}
Below is a sample table showing data:

\begin{table}[h]
\centering
\begin{tabular}{@{}lrrr@{}}
\toprule
\textbf{Category} & \textbf{Value 1} & \textbf{Value 2} & \textbf{Total} \\
\midrule
Item A & 10.5 & 20.3 & 30.8 \\
Item B & 15.2 & 25.1 & 40.3 \\
Item C & 22.1 & 12.7 & 34.8 \\
\bottomrule
\end{tabular}
\caption{Sample Data Table}
\label{tab:sample}
\end{table}

\section{Lists}
\subsection{Bullet Points}
\begin{itemize}
  \item First item
  \item Second item
  \item Third item with sub-items:
    \begin{itemize}
      \item Sub-item 1
      \item Sub-item 2
    \end{itemize}
\end{itemize}

\subsection{Numbered List}
\begin{enumerate}
  \item First step
  \item Second step
  \item Third step
\end{enumerate}

\section{Citations and References}
When you use the citation search feature, references will be automatically managed. For example, you might cite a paper like this: \cite{example2024}. All citations will be automatically added to your \texttt{cite.bib} file.

\section{Conclusion}
This template demonstrates basic \LaTeX{} features including mathematical formulas, tables, and lists. The bibliography system is now set up to work seamlessly with Bibby AI's citation management. Feel free to modify it for your needs!

\bibliographystyle{plain}
\bibliography{cite}

\end{document}\documentclass{article}
\usepackage[utf8]{inputenc}
\usepackage{booktabs}
\usepackage{amsmath}
\usepackage{natbib}

\title{My Document}
\author{Your Name}
\date{\today}

\begin{document}

\maketitle

\section{Introduction}
This is a sample document demonstrating various \LaTeX{} features including mathematical equations, tables, and formatting. You can cite papers using the citation search feature, and they will be automatically added to your \texttt{cite.bib} file.

\section{Mathematical Expressions}
Here are some example equations:

\subsection{Quadratic Formula}
The quadratic formula is given by:
\[x = \frac{-b \pm \sqrt{b^2 - 4ac}}{2a}\]

\subsection{Matrix Example}
A sample matrix:
\[
\begin{pmatrix}
a_{11} & a_{12} & a_{13} \\
a_{21} & a_{22} & a_{23} \\
a_{31} & a_{32} & a_{33}
\end{pmatrix}
\]

\section{Tables}
Below is a sample table showing data:

\begin{table}[h]
\centering
\begin{tabular}{@{}lrrr@{}}
\toprule
\textbf{Category} & \textbf{Value 1} & \textbf{Value 2} & \textbf{Total} \\
\midrule
Item A & 10.5 & 20.3 & 30.8 \\
Item B & 15.2 & 25.1 & 40.3 \\
Item C & 22.1 & 12.7 & 34.8 \\
\bottomrule
\end{tabular}
\caption{Sample Data Table}
\label{tab:sample}
\end{table}

\section{Lists}
\subsection{Bullet Points}
\begin{itemize}
  \item First item
  \item Second item
  \item Third item with sub-items:
    \begin{itemize}
      \item Sub-item 1
      \item Sub-item 2
    \end{itemize}
\end{itemize}

\subsection{Numbered List}
\begin{enumerate}
  \item First step
  \item Second step
  \item Third step
\end{enumerate}

\section{Citations and References}
When you use the citation search feature, references will be automatically managed. For example, you might cite a paper like this: \cite{example2024}. All citations will be automatically added to your \texttt{cite.bib} file.

\section{Conclusion}
This template demonstrates basic \LaTeX{} features including mathematical formulas, tables, and lists. The bibliography system is now set up to work seamlessly with Bibby AI's citation management. Feel free to modify it for your needs!

\bibliographystyle{plain}
\bibliography{cite}

\end{document}\documentclass{article}
\usepackage[utf8]{inputenc}
\usepackage{booktabs}
\usepackage{amsmath}
\usepackage{natbib}

\title{My Document}
\author{Your Name}
\date{\today}

\begin{document}

\maketitle

\section{Introduction}
This is a sample document demonstrating various \LaTeX{} features including mathematical equations, tables, and formatting. You can cite papers using the citation search feature, and they will be automatically added to your \texttt{cite.bib} file.

\section{Mathematical Expressions}
Here are some example equations:

\subsection{Quadratic Formula}
The quadratic formula is given by:
\[x = \frac{-b \pm \sqrt{b^2 - 4ac}}{2a}\]

\subsection{Matrix Example}
A sample matrix:
\[
\begin{pmatrix}
a_{11} & a_{12} & a_{13} \\
a_{21} & a_{22} & a_{23} \\
a_{31} & a_{32} & a_{33}
\end{pmatrix}
\]

\section{Tables}
Below is a sample table showing data:

\begin{table}[h]
\centering
\begin{tabular}{@{}lrrr@{}}
\toprule
\textbf{Category} & \textbf{Value 1} & \textbf{Value 2} & \textbf{Total} \\
\midrule
Item A & 10.5 & 20.3 & 30.8 \\
Item B & 15.2 & 25.1 & 40.3 \\
Item C & 22.1 & 12.7 & 34.8 \\
\bottomrule
\end{tabular}
\caption{Sample Data Table}
\label{tab:sample}
\end{table}

\section{Lists}
\subsection{Bullet Points}
\begin{itemize}
  \item First item
  \item Second item
  \item Third item with sub-items:
    \begin{itemize}
      \item Sub-item 1
      \item Sub-item 2
    \end{itemize}
\end{itemize}

\subsection{Numbered List}
\begin{enumerate}
  \item First step
  \item Second step
  \item Third step
\end{enumerate}

\section{Citations and References}
When you use the citation search feature, references will be automatically managed. For example, you might cite a paper like this: \cite{example2024}. All citations will be automatically added to your \texttt{cite.bib} file.

\section{Conclusion}
This template demonstrates basic \LaTeX{} features including mathematical formulas, tables, and lists. The bibliography system is now set up to work seamlessly with Bibby AI's citation management. Feel free to modify it for your needs!

\bibliographystyle{plain}
\bibliography{cite}

\end{document}\documentclass{article}
\usepackage[utf8]{inputenc}
\usepackage{booktabs}
\usepackage{amsmath}
\usepackage{natbib}

\title{My Document}
\author{Your Name}
\date{\today}

\begin{document}

\maketitle

\section{Introduction}
This is a sample document demonstrating various \LaTeX{} features including mathematical equations, tables, and formatting. You can cite papers using the citation search feature, and they will be automatically added to your \texttt{cite.bib} file.

\section{Mathematical Expressions}
Here are some example equations:

\subsection{Quadratic Formula}
The quadratic formula is given by:
\[x = \frac{-b \pm \sqrt{b^2 - 4ac}}{2a}\]

\subsection{Matrix Example}
A sample matrix:
\[
\begin{pmatrix}
a_{11} & a_{12} & a_{13} \\
a_{21} & a_{22} & a_{23} \\
a_{31} & a_{32} & a_{33}
\end{pmatrix}
\]

\section{Tables}
Below is a sample table showing data:

\begin{table}[h]
\centering
\begin{tabular}{@{}lrrr@{}}
\toprule
\textbf{Category} & \textbf{Value 1} & \textbf{Value 2} & \textbf{Total} \\
\midrule
Item A & 10.5 & 20.3 & 30.8 \\
Item B & 15.2 & 25.1 & 40.3 \\
Item C & 22.1 & 12.7 & 34.8 \\
\bottomrule
\end{tabular}
\caption{Sample Data Table}
\label{tab:sample}
\end{table}

\section{Lists}
\subsection{Bullet Points}
\begin{itemize}
  \item First item
  \item Second item
  \item Third item with sub-items:
    \begin{itemize}
      \item Sub-item 1
      \item Sub-item 2
    \end{itemize}
\end{itemize}

\subsection{Numbered List}
\begin{enumerate}
  \item First step
  \item Second step
  \item Third step
\end{enumerate}

\section{Citations and References}
When you use the citation search feature, references will be automatically managed. For example, you might cite a paper like this: \cite{example2024}. All citations will be automatically added to your \texttt{cite.bib} file.

\section{Conclusion}
This template demonstrates basic \LaTeX{} features including mathematical formulas, tables, and lists. The bibliography system is now set up to work seamlessly with Bibby AI's citation management. Feel free to modify it for your needs!

\bibliographystyle{plain}
\bibliography{cite}

\end{document}\documentclass{article}
\usepackage[utf8]{inputenc}
\usepackage{booktabs}
\usepackage{amsmath}
\usepackage{natbib}

\title{My Document}
\author{Your Name}
\date{\today}

\begin{document}

\maketitle

\section{Introduction}
This is a sample document demonstrating various \LaTeX{} features including mathematical equations, tables, and formatting. You can cite papers using the citation search feature, and they will be automatically added to your \texttt{cite.bib} file.

\section{Mathematical Expressions}
Here are some example equations:

\subsection{Quadratic Formula}
The quadratic formula is given by:
\[x = \frac{-b \pm \sqrt{b^2 - 4ac}}{2a}\]

\subsection{Matrix Example}
A sample matrix:
\[
\begin{pmatrix}
a_{11} & a_{12} & a_{13} \\
a_{21} & a_{22} & a_{23} \\
a_{31} & a_{32} & a_{33}
\end{pmatrix}
\]

\section{Tables}
Below is a sample table showing data:

\begin{table}[h]
\centering
\begin{tabular}{@{}lrrr@{}}
\toprule
\textbf{Category} & \textbf{Value 1} & \textbf{Value 2} & \textbf{Total} \\
\midrule
Item A & 10.5 & 20.3 & 30.8 \\
Item B & 15.2 & 25.1 & 40.3 \\
Item C & 22.1 & 12.7 & 34.8 \\
\bottomrule
\end{tabular}
\caption{Sample Data Table}
\label{tab:sample}
\end{table}

\section{Lists}
\subsection{Bullet Points}
\begin{itemize}
  \item First item
  \item Second item
  \item Third item with sub-items:
    \begin{itemize}
      \item Sub-item 1
      \item Sub-item 2
    \end{itemize}
\end{itemize}

\subsection{Numbered List}
\begin{enumerate}
  \item First step
  \item Second step
  \item Third step
\end{enumerate}

\section{Citations and References}
When you use the citation search feature, references will be automatically managed. For example, you might cite a paper like this: \cite{example2024}. All citations will be automatically added to your \texttt{cite.bib} file.

\section{Conclusion}
This template demonstrates basic \LaTeX{} features including mathematical formulas, tables, and lists. The bibliography system is now set up to work seamlessly with Bibby AI's citation management. Feel free to modify it for your needs!

\bibliographystyle{plain}
\bibliography{cite}

\end{document}\documentclass{article}
\usepackage[utf8]{inputenc}
\usepackage{booktabs}
\usepackage{amsmath}
\usepackage{natbib}

\title{My Document}
\author{Your Name}
\date{\today}

\begin{document}

\maketitle

\section{Introduction}
This is a sample document demonstrating various \LaTeX{} features including mathematical equations, tables, and formatting. You can cite papers using the citation search feature, and they will be automatically added to your \texttt{cite.bib} file.

\section{Mathematical Expressions}
Here are some example equations:

\subsection{Quadratic Formula}
The quadratic formula is given by:
\[x = \frac{-b \pm \sqrt{b^2 - 4ac}}{2a}\]

\subsection{Matrix Example}
A sample matrix:
\[
\begin{pmatrix}
a_{11} & a_{12} & a_{13} \\
a_{21} & a_{22} & a_{23} \\
a_{31} & a_{32} & a_{33}
\end{pmatrix}
\]

\section{Tables}
Below is a sample table showing data:

\begin{table}[h]
\centering
\begin{tabular}{@{}lrrr@{}}
\toprule
\textbf{Category} & \textbf{Value 1} & \textbf{Value 2} & \textbf{Total} \\
\midrule
Item A & 10.5 & 20.3 & 30.8 \\
Item B & 15.2 & 25.1 & 40.3 \\
Item C & 22.1 & 12.7 & 34.8 \\
\bottomrule
\end{tabular}
\caption{Sample Data Table}
\label{tab:sample}
\end{table}

\section{Lists}
\subsection{Bullet Points}
\begin{itemize}
  \item First item
  \item Second item
  \item Third item with sub-items:
    \begin{itemize}
      \item Sub-item 1
      \item Sub-item 2
    \end{itemize}
\end{itemize}

\subsection{Numbered List}
\begin{enumerate}
  \item First step
  \item Second step
  \item Third step
\end{enumerate}

\section{Citations and References}
When you use the citation search feature, references will be automatically managed. For example, you might cite a paper like this: \cite{example2024}. All citations will be automatically added to your \texttt{cite.bib} file.

\section{Conclusion}
This template demonstrates basic \LaTeX{} features including mathematical formulas, tables, and lists. The bibliography system is now set up to work seamlessly with Bibby AI's citation management. Feel free to modify it for your needs!

\bibliographystyle{plain}
\bibliography{cite}

\end{document}\documentclass{article}
\usepackage[utf8]{inputenc}
\usepackage{booktabs}
\usepackage{amsmath}
\usepackage{natbib}

\title{My Document}
\author{Your Name}
\date{\today}

\begin{document}

\maketitle

\section{Introduction}
This is a sample document demonstrating various \LaTeX{} features including mathematical equations, tables, and formatting. You can cite papers using the citation search feature, and they will be automatically added to your \texttt{cite.bib} file.

\section{Mathematical Expressions}
Here are some example equations:

\subsection{Quadratic Formula}
The quadratic formula is given by:
\[x = \frac{-b \pm \sqrt{b^2 - 4ac}}{2a}\]

\subsection{Matrix Example}
A sample matrix:
\[
\begin{pmatrix}
a_{11} & a_{12} & a_{13} \\
a_{21} & a_{22} & a_{23} \\
a_{31} & a_{32} & a_{33}
\end{pmatrix}
\]

\section{Tables}
Below is a sample table showing data:

\begin{table}[h]
\centering
\begin{tabular}{@{}lrrr@{}}
\toprule
\textbf{Category} & \textbf{Value 1} & \textbf{Value 2} & \textbf{Total} \\
\midrule
Item A & 10.5 & 20.3 & 30.8 \\
Item B & 15.2 & 25.1 & 40.3 \\
Item C & 22.1 & 12.7 & 34.8 \\
\bottomrule
\end{tabular}
\caption{Sample Data Table}
\label{tab:sample}
\end{table}

\section{Lists}
\subsection{Bullet Points}
\begin{itemize}
  \item First item
  \item Second item
  \item Third item with sub-items:
    \begin{itemize}
      \item Sub-item 1
      \item Sub-item 2
    \end{itemize}
\end{itemize}

\subsection{Numbered List}
\begin{enumerate}
  \item First step
  \item Second step
  \item Third step
\end{enumerate}

\section{Citations and References}
When you use the citation search feature, references will be automatically managed. For example, you might cite a paper like this: \cite{example2024}. All citations will be automatically added to your \texttt{cite.bib} file.

\section{Conclusion}
This template demonstrates basic \LaTeX{} features including mathematical formulas, tables, and lists. The bibliography system is now set up to work seamlessly with Bibby AI's citation management. Feel free to modify it for your needs!

\bibliographystyle{plain}
\bibliography{cite}

\end{document}\documentclass{article}
\usepackage[utf8]{inputenc}
\usepackage{booktabs}
\usepackage{amsmath}
\usepackage{natbib}

\title{My Document}
\author{Your Name}
\date{\today}

\begin{document}

\maketitle

\section{Introduction}
This is a sample document demonstrating various \LaTeX{} features including mathematical equations, tables, and formatting. You can cite papers using the citation search feature, and they will be automatically added to your \texttt{cite.bib} file.

\section{Mathematical Expressions}
Here are some example equations:

\subsection{Quadratic Formula}
The quadratic formula is given by:
\[x = \frac{-b \pm \sqrt{b^2 - 4ac}}{2a}\]

\subsection{Matrix Example}
A sample matrix:
\[
\begin{pmatrix}
a_{11} & a_{12} & a_{13} \\
a_{21} & a_{22} & a_{23} \\
a_{31} & a_{32} & a_{33}
\end{pmatrix}
\]

\section{Tables}
Below is a sample table showing data:

\begin{table}[h]
\centering
\begin{tabular}{@{}lrrr@{}}
\toprule
\textbf{Category} & \textbf{Value 1} & \textbf{Value 2} & \textbf{Total} \\
\midrule
Item A & 10.5 & 20.3 & 30.8 \\
Item B & 15.2 & 25.1 & 40.3 \\
Item C & 22.1 & 12.7 & 34.8 \\
\bottomrule
\end{tabular}
\caption{Sample Data Table}
\label{tab:sample}
\end{table}

\section{Lists}
\subsection{Bullet Points}
\begin{itemize}
  \item First item
  \item Second item
  \item Third item with sub-items:
    \begin{itemize}
      \item Sub-item 1
      \item Sub-item 2
    \end{itemize}
\end{itemize}

\subsection{Numbered List}
\begin{enumerate}
  \item First step
  \item Second step
  \item Third step
\end{enumerate}

\section{Citations and References}
When you use the citation search feature, references will be automatically managed. For example, you might cite a paper like this: \cite{example2024}. All citations will be automatically added to your \texttt{cite.bib} file.

\section{Conclusion}
This template demonstrates basic \LaTeX{} features including mathematical formulas, tables, and lists. The bibliography system is now set up to work seamlessly with Bibby AI's citation management. Feel free to modify it for your needs!

\bibliographystyle{plain}
\bibliography{cite}

\end{document}\documentclass{article}
\usepackage[utf8]{inputenc}
\usepackage{booktabs}
\usepackage{amsmath}
\usepackage{natbib}

\title{My Document}
\author{Your Name}
\date{\today}

\begin{document}

\maketitle

\section{Introduction}
This is a sample document demonstrating various \LaTeX{} features including mathematical equations, tables, and formatting. You can cite papers using the citation search feature, and they will be automatically added to your \texttt{cite.bib} file.

\section{Mathematical Expressions}
Here are some example equations:

\subsection{Quadratic Formula}
The quadratic formula is given by:
\[x = \frac{-b \pm \sqrt{b^2 - 4ac}}{2a}\]

\subsection{Matrix Example}
A sample matrix:
\[
\begin{pmatrix}
a_{11} & a_{12} & a_{13} \\
a_{21} & a_{22} & a_{23} \\
a_{31} & a_{32} & a_{33}
\end{pmatrix}
\]

\section{Tables}
Below is a sample table showing data:

\begin{table}[h]
\centering
\begin{tabular}{@{}lrrr@{}}
\toprule
\textbf{Category} & \textbf{Value 1} & \textbf{Value 2} & \textbf{Total} \\
\midrule
Item A & 10.5 & 20.3 & 30.8 \\
Item B & 15.2 & 25.1 & 40.3 \\
Item C & 22.1 & 12.7 & 34.8 \\
\bottomrule
\end{tabular}
\caption{Sample Data Table}
\label{tab:sample}
\end{table}

\section{Lists}
\subsection{Bullet Points}
\begin{itemize}
  \item First item
  \item Second item
  \item Third item with sub-items:
    \begin{itemize}
      \item Sub-item 1
      \item Sub-item 2
    \end{itemize}
\end{itemize}

\subsection{Numbered List}
\begin{enumerate}
  \item First step
  \item Second step
  \item Third step
\end{enumerate}

\section{Citations and References}
When you use the citation search feature, references will be automatically managed. For example, you might cite a paper like this: \cite{example2024}. All citations will be automatically added to your \texttt{cite.bib} file.

\section{Conclusion}
This template demonstrates basic \LaTeX{} features including mathematical formulas, tables, and lists. The bibliography system is now set up to work seamlessly with Bibby AI's citation management. Feel free to modify it for your needs!

\bibliographystyle{plain}
\bibliography{cite}

\end{document}\documentclass{article}
\usepackage[utf8]{inputenc}
\usepackage{booktabs}
\usepackage{amsmath}
\usepackage{natbib}

\title{My Document}
\author{Your Name}
\date{\today}

\begin{document}

\maketitle

\section{Introduction}
This is a sample document demonstrating various \LaTeX{} features including mathematical equations, tables, and formatting. You can cite papers using the citation search feature, and they will be automatically added to your \texttt{cite.bib} file.

\section{Mathematical Expressions}
Here are some example equations:

\subsection{Quadratic Formula}
The quadratic formula is given by:
\[x = \frac{-b \pm \sqrt{b^2 - 4ac}}{2a}\]

\subsection{Matrix Example}
A sample matrix:
\[
\begin{pmatrix}
a_{11} & a_{12} & a_{13} \\
a_{21} & a_{22} & a_{23} \\
a_{31} & a_{32} & a_{33}
\end{pmatrix}
\]

\section{Tables}
Below is a sample table showing data:

\begin{table}[h]
\centering
\begin{tabular}{@{}lrrr@{}}
\toprule
\textbf{Category} & \textbf{Value 1} & \textbf{Value 2} & \textbf{Total} \\
\midrule
Item A & 10.5 & 20.3 & 30.8 \\
Item B & 15.2 & 25.1 & 40.3 \\
Item C & 22.1 & 12.7 & 34.8 \\
\bottomrule
\end{tabular}
\caption{Sample Data Table}
\label{tab:sample}
\end{table}

\section{Lists}
\subsection{Bullet Points}
\begin{itemize}
  \item First item
  \item Second item
  \item Third item with sub-items:
    \begin{itemize}
      \item Sub-item 1
      \item Sub-item 2
    \end{itemize}
\end{itemize}

\subsection{Numbered List}
\begin{enumerate}
  \item First step
  \item Second step
  \item Third step
\end{enumerate}

\section{Citations and References}
When you use the citation search feature, references will be automatically managed. For example, you might cite a paper like this: \cite{example2024}. All citations will be automatically added to your \texttt{cite.bib} file.

\section{Conclusion}
This template demonstrates basic \LaTeX{} features including mathematical formulas, tables, and lists. The bibliography system is now set up to work seamlessly with Bibby AI's citation management. Feel free to modify it for your needs!

\bibliographystyle{plain}
\bibliography{cite}

\end{document}\documentclass{article}
\usepackage[utf8]{inputenc}
\usepackage{booktabs}
\usepackage{amsmath}
\usepackage{natbib}

\title{My Document}
\author{Your Name}
\date{\today}

\begin{document}

\maketitle

\section{Introduction}
This is a sample document demonstrating various \LaTeX{} features including mathematical equations, tables, and formatting. You can cite papers using the citation search feature, and they will be automatically added to your \texttt{cite.bib} file.

\section{Mathematical Expressions}
Here are some example equations:

\subsection{Quadratic Formula}
The quadratic formula is given by:
\[x = \frac{-b \pm \sqrt{b^2 - 4ac}}{2a}\]

\subsection{Matrix Example}
A sample matrix:
\[
\begin{pmatrix}
a_{11} & a_{12} & a_{13} \\
a_{21} & a_{22} & a_{23} \\
a_{31} & a_{32} & a_{33}
\end{pmatrix}
\]

\section{Tables}
Below is a sample table showing data:

\begin{table}[h]
\centering
\begin{tabular}{@{}lrrr@{}}
\toprule
\textbf{Category} & \textbf{Value 1} & \textbf{Value 2} & \textbf{Total} \\
\midrule
Item A & 10.5 & 20.3 & 30.8 \\
Item B & 15.2 & 25.1 & 40.3 \\
Item C & 22.1 & 12.7 & 34.8 \\
\bottomrule
\end{tabular}
\caption{Sample Data Table}
\label{tab:sample}
\end{table}

\section{Lists}
\subsection{Bullet Points}
\begin{itemize}
  \item First item
  \item Second item
  \item Third item with sub-items:
    \begin{itemize}
      \item Sub-item 1
      \item Sub-item 2
    \end{itemize}
\end{itemize}

\subsection{Numbered List}
\begin{enumerate}
  \item First step
  \item Second step
  \item Third step
\end{enumerate}

\section{Citations and References}
When you use the citation search feature, references will be automatically managed. For example, you might cite a paper like this: \cite{example2024}. All citations will be automatically added to your \texttt{cite.bib} file.

\section{Conclusion}
This template demonstrates basic \LaTeX{} features including mathematical formulas, tables, and lists. The bibliography system is now set up to work seamlessly with Bibby AI's citation management. Feel free to modify it for your needs!

\bibliographystyle{plain}
\bibliography{cite}

\end{document}\documentclass{article}
\usepackage[utf8]{inputenc}
\usepackage{booktabs}
\usepackage{amsmath}
\usepackage{natbib}

\title{My Document}
\author{Your Name}
\date{\today}

\begin{document}

\maketitle

\section{Introduction}
This is a sample document demonstrating various \LaTeX{} features including mathematical equations, tables, and formatting. You can cite papers using the citation search feature, and they will be automatically added to your \texttt{cite.bib} file.

\section{Mathematical Expressions}
Here are some example equations:

\subsection{Quadratic Formula}
The quadratic formula is given by:
\[x = \frac{-b \pm \sqrt{b^2 - 4ac}}{2a}\]

\subsection{Matrix Example}
A sample matrix:
\[
\begin{pmatrix}
a_{11} & a_{12} & a_{13} \\
a_{21} & a_{22} & a_{23} \\
a_{31} & a_{32} & a_{33}
\end{pmatrix}
\]

\section{Tables}
Below is a sample table showing data:

\begin{table}[h]
\centering
\begin{tabular}{@{}lrrr@{}}
\toprule
\textbf{Category} & \textbf{Value 1} & \textbf{Value 2} & \textbf{Total} \\
\midrule
Item A & 10.5 & 20.3 & 30.8 \\
Item B & 15.2 & 25.1 & 40.3 \\
Item C & 22.1 & 12.7 & 34.8 \\
\bottomrule
\end{tabular}
\caption{Sample Data Table}
\label{tab:sample}
\end{table}

\section{Lists}
\subsection{Bullet Points}
\begin{itemize}
  \item First item
  \item Second item
  \item Third item with sub-items:
    \begin{itemize}
      \item Sub-item 1
      \item Sub-item 2
    \end{itemize}
\end{itemize}

\subsection{Numbered List}
\begin{enumerate}
  \item First step
  \item Second step
  \item Third step
\end{enumerate}

\section{Citations and References}
When you use the citation search feature, references will be automatically managed. For example, you might cite a paper like this: \cite{example2024}. All citations will be automatically added to your \texttt{cite.bib} file.

\section{Conclusion}
This template demonstrates basic \LaTeX{} features including mathematical formulas, tables, and lists. The bibliography system is now set up to work seamlessly with Bibby AI's citation management. Feel free to modify it for your needs!

\bibliographystyle{plain}
\bibliography{cite}

\end{document}\documentclass{article}
\usepackage[utf8]{inputenc}
\usepackage{booktabs}
\usepackage{amsmath}
\usepackage{natbib}

\title{My Document}
\author{Your Name}
\date{\today}

\begin{document}

\maketitle

\section{Introduction}
This is a sample document demonstrating various \LaTeX{} features including mathematical equations, tables, and formatting. You can cite papers using the citation search feature, and they will be automatically added to your \texttt{cite.bib} file.

\section{Mathematical Expressions}
Here are some example equations:

\subsection{Quadratic Formula}
The quadratic formula is given by:
\[x = \frac{-b \pm \sqrt{b^2 - 4ac}}{2a}\]

\subsection{Matrix Example}
A sample matrix:
\[
\begin{pmatrix}
a_{11} & a_{12} & a_{13} \\
a_{21} & a_{22} & a_{23} \\
a_{31} & a_{32} & a_{33}
\end{pmatrix}
\]

\section{Tables}
Below is a sample table showing data:

\begin{table}[h]
\centering
\begin{tabular}{@{}lrrr@{}}
\toprule
\textbf{Category} & \textbf{Value 1} & \textbf{Value 2} & \textbf{Total} \\
\midrule
Item A & 10.5 & 20.3 & 30.8 \\
Item B & 15.2 & 25.1 & 40.3 \\
Item C & 22.1 & 12.7 & 34.8 \\
\bottomrule
\end{tabular}
\caption{Sample Data Table}
\label{tab:sample}
\end{table}

\section{Lists}
\subsection{Bullet Points}
\begin{itemize}
  \item First item
  \item Second item
  \item Third item with sub-items:
    \begin{itemize}
      \item Sub-item 1
      \item Sub-item 2
    \end{itemize}
\end{itemize}

\subsection{Numbered List}
\begin{enumerate}
  \item First step
  \item Second step
  \item Third step
\end{enumerate}

\section{Citations and References}
When you use the citation search feature, references will be automatically managed. For example, you might cite a paper like this: \cite{example2024}. All citations will be automatically added to your \texttt{cite.bib} file.

\section{Conclusion}
This template demonstrates basic \LaTeX{} features including mathematical formulas, tables, and lists. The bibliography system is now set up to work seamlessly with Bibby AI's citation management. Feel free to modify it for your needs!

\bibliographystyle{plain}
\bibliography{cite}

\end{document}\documentclass{article}
\usepackage[utf8]{inputenc}
\usepackage{booktabs}
\usepackage{amsmath}
\usepackage{natbib}

\title{My Document}
\author{Your Name}
\date{\today}

\begin{document}

\maketitle

\section{Introduction}
This is a sample document demonstrating various \LaTeX{} features including mathematical equations, tables, and formatting. You can cite papers using the citation search feature, and they will be automatically added to your \texttt{cite.bib} file.

\section{Mathematical Expressions}
Here are some example equations:

\subsection{Quadratic Formula}
The quadratic formula is given by:
\[x = \frac{-b \pm \sqrt{b^2 - 4ac}}{2a}\]

\subsection{Matrix Example}
A sample matrix:
\[
\begin{pmatrix}
a_{11} & a_{12} & a_{13} \\
a_{21} & a_{22} & a_{23} \\
a_{31} & a_{32} & a_{33}
\end{pmatrix}
\]

\section{Tables}
Below is a sample table showing data:

\begin{table}[h]
\centering
\begin{tabular}{@{}lrrr@{}}
\toprule
\textbf{Category} & \textbf{Value 1} & \textbf{Value 2} & \textbf{Total} \\
\midrule
Item A & 10.5 & 20.3 & 30.8 \\
Item B & 15.2 & 25.1 & 40.3 \\
Item C & 22.1 & 12.7 & 34.8 \\
\bottomrule
\end{tabular}
\caption{Sample Data Table}
\label{tab:sample}
\end{table}

\section{Lists}
\subsection{Bullet Points}
\begin{itemize}
  \item First item
  \item Second item
  \item Third item with sub-items:
    \begin{itemize}
      \item Sub-item 1
      \item Sub-item 2
    \end{itemize}
\end{itemize}

\subsection{Numbered List}
\begin{enumerate}
  \item First step
  \item Second step
  \item Third step
\end{enumerate}

\section{Citations and References}
When you use the citation search feature, references will be automatically managed. For example, you might cite a paper like this: \cite{example2024}. All citations will be automatically added to your \texttt{cite.bib} file.

\section{Conclusion}
This template demonstrates basic \LaTeX{} features including mathematical formulas, tables, and lists. The bibliography system is now set up to work seamlessly with Bibby AI's citation management. Feel free to modify it for your needs!

\bibliographystyle{plain}
\bibliography{cite}

\end{document}\documentclass{article}
\usepackage[utf8]{inputenc}
\usepackage{booktabs}
\usepackage{amsmath}
\usepackage{natbib}

\title{My Document}
\author{Your Name}
\date{\today}

\begin{document}

\maketitle

\section{Introduction}
This is a sample document demonstrating various \LaTeX{} features including mathematical equations, tables, and formatting. You can cite papers using the citation search feature, and they will be automatically added to your \texttt{cite.bib} file.

\section{Mathematical Expressions}
Here are some example equations:

\subsection{Quadratic Formula}
The quadratic formula is given by:
\[x = \frac{-b \pm \sqrt{b^2 - 4ac}}{2a}\]

\subsection{Matrix Example}
A sample matrix:
\[
\begin{pmatrix}
a_{11} & a_{12} & a_{13} \\
a_{21} & a_{22} & a_{23} \\
a_{31} & a_{32} & a_{33}
\end{pmatrix}
\]

\section{Tables}
Below is a sample table showing data:

\begin{table}[h]
\centering
\begin{tabular}{@{}lrrr@{}}
\toprule
\textbf{Category} & \textbf{Value 1} & \textbf{Value 2} & \textbf{Total} \\
\midrule
Item A & 10.5 & 20.3 & 30.8 \\
Item B & 15.2 & 25.1 & 40.3 \\
Item C & 22.1 & 12.7 & 34.8 \\
\bottomrule
\end{tabular}
\caption{Sample Data Table}
\label{tab:sample}
\end{table}

\section{Lists}
\subsection{Bullet Points}
\begin{itemize}
  \item First item
  \item Second item
  \item Third item with sub-items:
    \begin{itemize}
      \item Sub-item 1
      \item Sub-item 2
    \end{itemize}
\end{itemize}

\subsection{Numbered List}
\begin{enumerate}
  \item First step
  \item Second step
  \item Third step
\end{enumerate}

\section{Citations and References}
When you use the citation search feature, references will be automatically managed. For example, you might cite a paper like this: \cite{example2024}. All citations will be automatically added to your \texttt{cite.bib} file.

\section{Conclusion}
This template demonstrates basic \LaTeX{} features including mathematical formulas, tables, and lists. The bibliography system is now set up to work seamlessly with Bibby AI's citation management. Feel free to modify it for your needs!

\bibliographystyle{plain}
\bibliography{cite}

\end{document}\documentclass{article}
\usepackage[utf8]{inputenc}
\usepackage{booktabs}
\usepackage{amsmath}
\usepackage{natbib}

\title{My Document}
\author{Your Name}
\date{\today}

\begin{document}

\maketitle

\section{Introduction}
This is a sample document demonstrating various \LaTeX{} features including mathematical equations, tables, and formatting. You can cite papers using the citation search feature, and they will be automatically added to your \texttt{cite.bib} file.

\section{Mathematical Expressions}
Here are some example equations:

\subsection{Quadratic Formula}
The quadratic formula is given by:
\[x = \frac{-b \pm \sqrt{b^2 - 4ac}}{2a}\]

\subsection{Matrix Example}
A sample matrix:
\[
\begin{pmatrix}
a_{11} & a_{12} & a_{13} \\
a_{21} & a_{22} & a_{23} \\
a_{31} & a_{32} & a_{33}
\end{pmatrix}
\]

\section{Tables}
Below is a sample table showing data:

\begin{table}[h]
\centering
\begin{tabular}{@{}lrrr@{}}
\toprule
\textbf{Category} & \textbf{Value 1} & \textbf{Value 2} & \textbf{Total} \\
\midrule
Item A & 10.5 & 20.3 & 30.8 \\
Item B & 15.2 & 25.1 & 40.3 \\
Item C & 22.1 & 12.7 & 34.8 \\
\bottomrule
\end{tabular}
\caption{Sample Data Table}
\label{tab:sample}
\end{table}

\section{Lists}
\subsection{Bullet Points}
\begin{itemize}
  \item First item
  \item Second item
  \item Third item with sub-items:
    \begin{itemize}
      \item Sub-item 1
      \item Sub-item 2
    \end{itemize}
\end{itemize}

\subsection{Numbered List}
\begin{enumerate}
  \item First step
  \item Second step
  \item Third step
\end{enumerate}

\section{Citations and References}
When you use the citation search feature, references will be automatically managed. For example, you might cite a paper like this: \cite{example2024}. All citations will be automatically added to your \texttt{cite.bib} file.

\section{Conclusion}
This template demonstrates basic \LaTeX{} features including mathematical formulas, tables, and lists. The bibliography system is now set up to work seamlessly with Bibby AI's citation management. Feel free to modify it for your needs!

\bibliographystyle{plain}
\bibliography{cite}

\end{document}\documentclass{article}
\usepackage[utf8]{inputenc}
\usepackage{booktabs}
\usepackage{amsmath}
\usepackage{natbib}

\title{My Document}
\author{Your Name}
\date{\today}

\begin{document}

\maketitle

\section{Introduction}
This is a sample document demonstrating various \LaTeX{} features including mathematical equations, tables, and formatting. You can cite papers using the citation search feature, and they will be automatically added to your \texttt{cite.bib} file.

\section{Mathematical Expressions}
Here are some example equations:

\subsection{Quadratic Formula}
The quadratic formula is given by:
\[x = \frac{-b \pm \sqrt{b^2 - 4ac}}{2a}\]

\subsection{Matrix Example}
A sample matrix:
\[
\begin{pmatrix}
a_{11} & a_{12} & a_{13} \\
a_{21} & a_{22} & a_{23} \\
a_{31} & a_{32} & a_{33}
\end{pmatrix}
\]

\section{Tables}
Below is a sample table showing data:

\begin{table}[h]
\centering
\begin{tabular}{@{}lrrr@{}}
\toprule
\textbf{Category} & \textbf{Value 1} & \textbf{Value 2} & \textbf{Total} \\
\midrule
Item A & 10.5 & 20.3 & 30.8 \\
Item B & 15.2 & 25.1 & 40.3 \\
Item C & 22.1 & 12.7 & 34.8 \\
\bottomrule
\end{tabular}
\caption{Sample Data Table}
\label{tab:sample}
\end{table}

\section{Lists}
\subsection{Bullet Points}
\begin{itemize}
  \item First item
  \item Second item
  \item Third item with sub-items:
    \begin{itemize}
      \item Sub-item 1
      \item Sub-item 2
    \end{itemize}
\end{itemize}

\subsection{Numbered List}
\begin{enumerate}
  \item First step
  \item Second step
  \item Third step
\end{enumerate}

\section{Citations and References}
When you use the citation search feature, references will be automatically managed. For example, you might cite a paper like this: \cite{example2024}. All citations will be automatically added to your \texttt{cite.bib} file.

\section{Conclusion}
This template demonstrates basic \LaTeX{} features including mathematical formulas, tables, and lists. The bibliography system is now set up to work seamlessly with Bibby AI's citation management. Feel free to modify it for your needs!

\bibliographystyle{plain}
\bibliography{cite}

\end{document}\documentclass{article}
\usepackage[utf8]{inputenc}
\usepackage{booktabs}
\usepackage{amsmath}
\usepackage{natbib}

\title{My Document}
\author{Your Name}
\date{\today}

\begin{document}

\maketitle

\section{Introduction}
This is a sample document demonstrating various \LaTeX{} features including mathematical equations, tables, and formatting. You can cite papers using the citation search feature, and they will be automatically added to your \texttt{cite.bib} file.

\section{Mathematical Expressions}
Here are some example equations:

\subsection{Quadratic Formula}
The quadratic formula is given by:
\[x = \frac{-b \pm \sqrt{b^2 - 4ac}}{2a}\]

\subsection{Matrix Example}
A sample matrix:
\[
\begin{pmatrix}
a_{11} & a_{12} & a_{13} \\
a_{21} & a_{22} & a_{23} \\
a_{31} & a_{32} & a_{33}
\end{pmatrix}
\]

\section{Tables}
Below is a sample table showing data:

\begin{table}[h]
\centering
\begin{tabular}{@{}lrrr@{}}
\toprule
\textbf{Category} & \textbf{Value 1} & \textbf{Value 2} & \textbf{Total} \\
\midrule
Item A & 10.5 & 20.3 & 30.8 \\
Item B & 15.2 & 25.1 & 40.3 \\
Item C & 22.1 & 12.7 & 34.8 \\
\bottomrule
\end{tabular}
\caption{Sample Data Table}
\label{tab:sample}
\end{table}

\section{Lists}
\subsection{Bullet Points}
\begin{itemize}
  \item First item
  \item Second item
  \item Third item with sub-items:
    \begin{itemize}
      \item Sub-item 1
      \item Sub-item 2
    \end{itemize}
\end{itemize}

\subsection{Numbered List}
\begin{enumerate}
  \item First step
  \item Second step
  \item Third step
\end{enumerate}

\section{Citations and References}
When you use the citation search feature, references will be automatically managed. For example, you might cite a paper like this: \cite{example2024}. All citations will be automatically added to your \texttt{cite.bib} file.

\section{Conclusion}
This template demonstrates basic \LaTeX{} features including mathematical formulas, tables, and lists. The bibliography system is now set up to work seamlessly with Bibby AI's citation management. Feel free to modify it for your needs!

\bibliographystyle{plain}
\bibliography{cite}

\end{document}\documentclass{article}
\usepackage[utf8]{inputenc}
\usepackage{booktabs}
\usepackage{amsmath}
\usepackage{natbib}

\title{My Document}
\author{Your Name}
\date{\today}

\begin{document}

\maketitle

\section{Introduction}
This is a sample document demonstrating various \LaTeX{} features including mathematical equations, tables, and formatting. You can cite papers using the citation search feature, and they will be automatically added to your \texttt{cite.bib} file.

\section{Mathematical Expressions}
Here are some example equations:

\subsection{Quadratic Formula}
The quadratic formula is given by:
\[x = \frac{-b \pm \sqrt{b^2 - 4ac}}{2a}\]

\subsection{Matrix Example}
A sample matrix:
\[
\begin{pmatrix}
a_{11} & a_{12} & a_{13} \\
a_{21} & a_{22} & a_{23} \\
a_{31} & a_{32} & a_{33}
\end{pmatrix}
\]

\section{Tables}
Below is a sample table showing data:

\begin{table}[h]
\centering
\begin{tabular}{@{}lrrr@{}}
\toprule
\textbf{Category} & \textbf{Value 1} & \textbf{Value 2} & \textbf{Total} \\
\midrule
Item A & 10.5 & 20.3 & 30.8 \\
Item B & 15.2 & 25.1 & 40.3 \\
Item C & 22.1 & 12.7 & 34.8 \\
\bottomrule
\end{tabular}
\caption{Sample Data Table}
\label{tab:sample}
\end{table}

\section{Lists}
\subsection{Bullet Points}
\begin{itemize}
  \item First item
  \item Second item
  \item Third item with sub-items:
    \begin{itemize}
      \item Sub-item 1
      \item Sub-item 2
    \end{itemize}
\end{itemize}

\subsection{Numbered List}
\begin{enumerate}
  \item First step
  \item Second step
  \item Third step
\end{enumerate}

\section{Citations and References}
When you use the citation search feature, references will be automatically managed. For example, you might cite a paper like this: \cite{example2024}. All citations will be automatically added to your \texttt{cite.bib} file.

\section{Conclusion}
This template demonstrates basic \LaTeX{} features including mathematical formulas, tables, and lists. The bibliography system is now set up to work seamlessly with Bibby AI's citation management. Feel free to modify it for your needs!

\bibliographystyle{plain}
\bibliography{cite}

\end{document}\documentclass{article}
\usepackage[utf8]{inputenc}
\usepackage{booktabs}
\usepackage{amsmath}
\usepackage{natbib}

\title{My Document}
\author{Your Name}
\date{\today}

\begin{document}

\maketitle

\section{Introduction}
This is a sample document demonstrating various \LaTeX{} features including mathematical equations, tables, and formatting. You can cite papers using the citation search feature, and they will be automatically added to your \texttt{cite.bib} file.

\section{Mathematical Expressions}
Here are some example equations:

\subsection{Quadratic Formula}
The quadratic formula is given by:
\[x = \frac{-b \pm \sqrt{b^2 - 4ac}}{2a}\]

\subsection{Matrix Example}
A sample matrix:
\[
\begin{pmatrix}
a_{11} & a_{12} & a_{13} \\
a_{21} & a_{22} & a_{23} \\
a_{31} & a_{32} & a_{33}
\end{pmatrix}
\]

\section{Tables}
Below is a sample table showing data:

\begin{table}[h]
\centering
\begin{tabular}{@{}lrrr@{}}
\toprule
\textbf{Category} & \textbf{Value 1} & \textbf{Value 2} & \textbf{Total} \\
\midrule
Item A & 10.5 & 20.3 & 30.8 \\
Item B & 15.2 & 25.1 & 40.3 \\
Item C & 22.1 & 12.7 & 34.8 \\
\bottomrule
\end{tabular}
\caption{Sample Data Table}
\label{tab:sample}
\end{table}

\section{Lists}
\subsection{Bullet Points}
\begin{itemize}
  \item First item
  \item Second item
  \item Third item with sub-items:
    \begin{itemize}
      \item Sub-item 1
      \item Sub-item 2
    \end{itemize}
\end{itemize}

\subsection{Numbered List}
\begin{enumerate}
  \item First step
  \item Second step
  \item Third step
\end{enumerate}

\section{Citations and References}
When you use the citation search feature, references will be automatically managed. For example, you might cite a paper like this: \cite{example2024}. All citations will be automatically added to your \texttt{cite.bib} file.

\section{Conclusion}
This template demonstrates basic \LaTeX{} features including mathematical formulas, tables, and lists. The bibliography system is now set up to work seamlessly with Bibby AI's citation management. Feel free to modify it for your needs!

\bibliographystyle{plain}
\bibliography{cite}

\end{document}\documentclass{article}
\usepackage[utf8]{inputenc}
\usepackage{booktabs}
\usepackage{amsmath}
\usepackage{natbib}

\title{My Document}
\author{Your Name}
\date{\today}

\begin{document}

\maketitle

\section{Introduction}
This is a sample document demonstrating various \LaTeX{} features including mathematical equations, tables, and formatting. You can cite papers using the citation search feature, and they will be automatically added to your \texttt{cite.bib} file.

\section{Mathematical Expressions}
Here are some example equations:

\subsection{Quadratic Formula}
The quadratic formula is given by:
\[x = \frac{-b \pm \sqrt{b^2 - 4ac}}{2a}\]

\subsection{Matrix Example}
A sample matrix:
\[
\begin{pmatrix}
a_{11} & a_{12} & a_{13} \\
a_{21} & a_{22} & a_{23} \\
a_{31} & a_{32} & a_{33}
\end{pmatrix}
\]

\section{Tables}
Below is a sample table showing data:

\begin{table}[h]
\centering
\begin{tabular}{@{}lrrr@{}}
\toprule
\textbf{Category} & \textbf{Value 1} & \textbf{Value 2} & \textbf{Total} \\
\midrule
Item A & 10.5 & 20.3 & 30.8 \\
Item B & 15.2 & 25.1 & 40.3 \\
Item C & 22.1 & 12.7 & 34.8 \\
\bottomrule
\end{tabular}
\caption{Sample Data Table}
\label{tab:sample}
\end{table}

\section{Lists}
\subsection{Bullet Points}
\begin{itemize}
  \item First item
  \item Second item
  \item Third item with sub-items:
    \begin{itemize}
      \item Sub-item 1
      \item Sub-item 2
    \end{itemize}
\end{itemize}

\subsection{Numbered List}
\begin{enumerate}
  \item First step
  \item Second step
  \item Third step
\end{enumerate}

\section{Citations and References}
When you use the citation search feature, references will be automatically managed. For example, you might cite a paper like this: \cite{example2024}. All citations will be automatically added to your \texttt{cite.bib} file.

\section{Conclusion}
This template demonstrates basic \LaTeX{} features including mathematical formulas, tables, and lists. The bibliography system is now set up to work seamlessly with Bibby AI's citation management. Feel free to modify it for your needs!

\bibliographystyle{plain}
\bibliography{cite}

\end{document}\documentclass{article}
\usepackage[utf8]{inputenc}
\usepackage{booktabs}
\usepackage{amsmath}
\usepackage{natbib}

\title{My Document}
\author{Your Name}
\date{\today}

\begin{document}

\maketitle

\section{Introduction}
This is a sample document demonstrating various \LaTeX{} features including mathematical equations, tables, and formatting. You can cite papers using the citation search feature, and they will be automatically added to your \texttt{cite.bib} file.

\section{Mathematical Expressions}
Here are some example equations:

\subsection{Quadratic Formula}
The quadratic formula is given by:
\[x = \frac{-b \pm \sqrt{b^2 - 4ac}}{2a}\]

\subsection{Matrix Example}
A sample matrix:
\[
\begin{pmatrix}
a_{11} & a_{12} & a_{13} \\
a_{21} & a_{22} & a_{23} \\
a_{31} & a_{32} & a_{33}
\end{pmatrix}
\]

\section{Tables}
Below is a sample table showing data:

\begin{table}[h]
\centering
\begin{tabular}{@{}lrrr@{}}
\toprule
\textbf{Category} & \textbf{Value 1} & \textbf{Value 2} & \textbf{Total} \\
\midrule
Item A & 10.5 & 20.3 & 30.8 \\
Item B & 15.2 & 25.1 & 40.3 \\
Item C & 22.1 & 12.7 & 34.8 \\
\bottomrule
\end{tabular}
\caption{Sample Data Table}
\label{tab:sample}
\end{table}

\section{Lists}
\subsection{Bullet Points}
\begin{itemize}
  \item First item
  \item Second item
  \item Third item with sub-items:
    \begin{itemize}
      \item Sub-item 1
      \item Sub-item 2
    \end{itemize}
\end{itemize}

\subsection{Numbered List}
\begin{enumerate}
  \item First step
  \item Second step
  \item Third step
\end{enumerate}

\section{Citations and References}
When you use the citation search feature, references will be automatically managed. For example, you might cite a paper like this: \cite{example2024}. All citations will be automatically added to your \texttt{cite.bib} file.

\section{Conclusion}
This template demonstrates basic \LaTeX{} features including mathematical formulas, tables, and lists. The bibliography system is now set up to work seamlessly with Bibby AI's citation management. Feel free to modify it for your needs!

\bibliographystyle{plain}
\bibliography{cite}

\end{document}\documentclass{article}
\usepackage[utf8]{inputenc}
\usepackage{booktabs}
\usepackage{amsmath}
\usepackage{natbib}

\title{My Document}
\author{Your Name}
\date{\today}

\begin{document}

\maketitle

\section{Introduction}
This is a sample document demonstrating various \LaTeX{} features including mathematical equations, tables, and formatting. You can cite papers using the citation search feature, and they will be automatically added to your \texttt{cite.bib} file.

\section{Mathematical Expressions}
Here are some example equations:

\subsection{Quadratic Formula}
The quadratic formula is given by:
\[x = \frac{-b \pm \sqrt{b^2 - 4ac}}{2a}\]

\subsection{Matrix Example}
A sample matrix:
\[
\begin{pmatrix}
a_{11} & a_{12} & a_{13} \\
a_{21} & a_{22} & a_{23} \\
a_{31} & a_{32} & a_{33}
\end{pmatrix}
\]

\section{Tables}
Below is a sample table showing data:

\begin{table}[h]
\centering
\begin{tabular}{@{}lrrr@{}}
\toprule
\textbf{Category} & \textbf{Value 1} & \textbf{Value 2} & \textbf{Total} \\
\midrule
Item A & 10.5 & 20.3 & 30.8 \\
Item B & 15.2 & 25.1 & 40.3 \\
Item C & 22.1 & 12.7 & 34.8 \\
\bottomrule
\end{tabular}
\caption{Sample Data Table}
\label{tab:sample}
\end{table}

\section{Lists}
\subsection{Bullet Points}
\begin{itemize}
  \item First item
  \item Second item
  \item Third item with sub-items:
    \begin{itemize}
      \item Sub-item 1
      \item Sub-item 2
    \end{itemize}
\end{itemize}

\subsection{Numbered List}
\begin{enumerate}
  \item First step
  \item Second step
  \item Third step
\end{enumerate}

\section{Citations and References}
When you use the citation search feature, references will be automatically managed. For example, you might cite a paper like this: \cite{example2024}. All citations will be automatically added to your \texttt{cite.bib} file.

\section{Conclusion}
This template demonstrates basic \LaTeX{} features including mathematical formulas, tables, and lists. The bibliography system is now set up to work seamlessly with Bibby AI's citation management. Feel free to modify it for your needs!

\bibliographystyle{plain}
\bibliography{cite}

\end{document}\documentclass{article}
\usepackage[utf8]{inputenc}
\usepackage{booktabs}
\usepackage{amsmath}
\usepackage{natbib}

\title{My Document}
\author{Your Name}
\date{\today}

\begin{document}

\maketitle

\section{Introduction}
This is a sample document demonstrating various \LaTeX{} features including mathematical equations, tables, and formatting. You can cite papers using the citation search feature, and they will be automatically added to your \texttt{cite.bib} file.

\section{Mathematical Expressions}
Here are some example equations:

\subsection{Quadratic Formula}
The quadratic formula is given by:
\[x = \frac{-b \pm \sqrt{b^2 - 4ac}}{2a}\]

\subsection{Matrix Example}
A sample matrix:
\[
\begin{pmatrix}
a_{11} & a_{12} & a_{13} \\
a_{21} & a_{22} & a_{23} \\
a_{31} & a_{32} & a_{33}
\end{pmatrix}
\]

\section{Tables}
Below is a sample table showing data:

\begin{table}[h]
\centering
\begin{tabular}{@{}lrrr@{}}
\toprule
\textbf{Category} & \textbf{Value 1} & \textbf{Value 2} & \textbf{Total} \\
\midrule
Item A & 10.5 & 20.3 & 30.8 \\
Item B & 15.2 & 25.1 & 40.3 \\
Item C & 22.1 & 12.7 & 34.8 \\
\bottomrule
\end{tabular}
\caption{Sample Data Table}
\label{tab:sample}
\end{table}

\section{Lists}
\subsection{Bullet Points}
\begin{itemize}
  \item First item
  \item Second item
  \item Third item with sub-items:
    \begin{itemize}
      \item Sub-item 1
      \item Sub-item 2
    \end{itemize}
\end{itemize}

\subsection{Numbered List}
\begin{enumerate}
  \item First step
  \item Second step
  \item Third step
\end{enumerate}

\section{Citations and References}
When you use the citation search feature, references will be automatically managed. For example, you might cite a paper like this: \cite{example2024}. All citations will be automatically added to your \texttt{cite.bib} file.

\section{Conclusion}
This template demonstrates basic \LaTeX{} features including mathematical formulas, tables, and lists. The bibliography system is now set up to work seamlessly with Bibby AI's citation management. Feel free to modify it for your needs!

\bibliographystyle{plain}
\bibliography{cite}

\end{document}\documentclass{article}
\usepackage[utf8]{inputenc}
\usepackage{booktabs}
\usepackage{amsmath}
\usepackage{natbib}

\title{My Document}
\author{Your Name}
\date{\today}

\begin{document}

\maketitle

\section{Introduction}
This is a sample document demonstrating various \LaTeX{} features including mathematical equations, tables, and formatting. You can cite papers using the citation search feature, and they will be automatically added to your \texttt{cite.bib} file.

\section{Mathematical Expressions}
Here are some example equations:

\subsection{Quadratic Formula}
The quadratic formula is given by:
\[x = \frac{-b \pm \sqrt{b^2 - 4ac}}{2a}\]

\subsection{Matrix Example}
A sample matrix:
\[
\begin{pmatrix}
a_{11} & a_{12} & a_{13} \\
a_{21} & a_{22} & a_{23} \\
a_{31} & a_{32} & a_{33}
\end{pmatrix}
\]

\section{Tables}
Below is a sample table showing data:

\begin{table}[h]
\centering
\begin{tabular}{@{}lrrr@{}}
\toprule
\textbf{Category} & \textbf{Value 1} & \textbf{Value 2} & \textbf{Total} \\
\midrule
Item A & 10.5 & 20.3 & 30.8 \\
Item B & 15.2 & 25.1 & 40.3 \\
Item C & 22.1 & 12.7 & 34.8 \\
\bottomrule
\end{tabular}
\caption{Sample Data Table}
\label{tab:sample}
\end{table}

\section{Lists}
\subsection{Bullet Points}
\begin{itemize}
  \item First item
  \item Second item
  \item Third item with sub-items:
    \begin{itemize}
      \item Sub-item 1
      \item Sub-item 2
    \end{itemize}
\end{itemize}

\subsection{Numbered List}
\begin{enumerate}
  \item First step
  \item Second step
  \item Third step
\end{enumerate}

\section{Citations and References}
When you use the citation search feature, references will be automatically managed. For example, you might cite a paper like this: \cite{example2024}. All citations will be automatically added to your \texttt{cite.bib} file.

\section{Conclusion}
This template demonstrates basic \LaTeX{} features including mathematical formulas, tables, and lists. The bibliography system is now set up to work seamlessly with Bibby AI's citation management. Feel free to modify it for your needs!

\bibliographystyle{plain}
\bibliography{cite}

\end{document}\documentclass{article}
\usepackage[utf8]{inputenc}
\usepackage{booktabs}
\usepackage{amsmath}
\usepackage{natbib}

\title{My Document}
\author{Your Name}
\date{\today}

\begin{document}

\maketitle

\section{Introduction}
This is a sample document demonstrating various \LaTeX{} features including mathematical equations, tables, and formatting. You can cite papers using the citation search feature, and they will be automatically added to your \texttt{cite.bib} file.

\section{Mathematical Expressions}
Here are some example equations:

\subsection{Quadratic Formula}
The quadratic formula is given by:
\[x = \frac{-b \pm \sqrt{b^2 - 4ac}}{2a}\]

\subsection{Matrix Example}
A sample matrix:
\[
\begin{pmatrix}
a_{11} & a_{12} & a_{13} \\
a_{21} & a_{22} & a_{23} \\
a_{31} & a_{32} & a_{33}
\end{pmatrix}
\]

\section{Tables}
Below is a sample table showing data:

\begin{table}[h]
\centering
\begin{tabular}{@{}lrrr@{}}
\toprule
\textbf{Category} & \textbf{Value 1} & \textbf{Value 2} & \textbf{Total} \\
\midrule
Item A & 10.5 & 20.3 & 30.8 \\
Item B & 15.2 & 25.1 & 40.3 \\
Item C & 22.1 & 12.7 & 34.8 \\
\bottomrule
\end{tabular}
\caption{Sample Data Table}
\label{tab:sample}
\end{table}

\section{Lists}
\subsection{Bullet Points}
\begin{itemize}
  \item First item
  \item Second item
  \item Third item with sub-items:
    \begin{itemize}
      \item Sub-item 1
      \item Sub-item 2
    \end{itemize}
\end{itemize}

\subsection{Numbered List}
\begin{enumerate}
  \item First step
  \item Second step
  \item Third step
\end{enumerate}

\section{Citations and References}
When you use the citation search feature, references will be automatically managed. For example, you might cite a paper like this: \cite{example2024}. All citations will be automatically added to your \texttt{cite.bib} file.

\section{Conclusion}
This template demonstrates basic \LaTeX{} features including mathematical formulas, tables, and lists. The bibliography system is now set up to work seamlessly with Bibby AI's citation management. Feel free to modify it for your needs!

\bibliographystyle{plain}
\bibliography{cite}

\end{document}\documentclass{article}
\usepackage[utf8]{inputenc}
\usepackage{booktabs}
\usepackage{amsmath}
\usepackage{natbib}

\title{My Document}
\author{Your Name}
\date{\today}

\begin{document}

\maketitle

\section{Introduction}
This is a sample document demonstrating various \LaTeX{} features including mathematical equations, tables, and formatting. You can cite papers using the citation search feature, and they will be automatically added to your \texttt{cite.bib} file.

\section{Mathematical Expressions}
Here are some example equations:

\subsection{Quadratic Formula}
The quadratic formula is given by:
\[x = \frac{-b \pm \sqrt{b^2 - 4ac}}{2a}\]

\subsection{Matrix Example}
A sample matrix:
\[
\begin{pmatrix}
a_{11} & a_{12} & a_{13} \\
a_{21} & a_{22} & a_{23} \\
a_{31} & a_{32} & a_{33}
\end{pmatrix}
\]

\section{Tables}
Below is a sample table showing data:

\begin{table}[h]
\centering
\begin{tabular}{@{}lrrr@{}}
\toprule
\textbf{Category} & \textbf{Value 1} & \textbf{Value 2} & \textbf{Total} \\
\midrule
Item A & 10.5 & 20.3 & 30.8 \\
Item B & 15.2 & 25.1 & 40.3 \\
Item C & 22.1 & 12.7 & 34.8 \\
\bottomrule
\end{tabular}
\caption{Sample Data Table}
\label{tab:sample}
\end{table}

\section{Lists}
\subsection{Bullet Points}
\begin{itemize}
  \item First item
  \item Second item
  \item Third item with sub-items:
    \begin{itemize}
      \item Sub-item 1
      \item Sub-item 2
    \end{itemize}
\end{itemize}

\subsection{Numbered List}
\begin{enumerate}
  \item First step
  \item Second step
  \item Third step
\end{enumerate}

\section{Citations and References}
When you use the citation search feature, references will be automatically managed. For example, you might cite a paper like this: \cite{example2024}. All citations will be automatically added to your \texttt{cite.bib} file.

\section{Conclusion}
This template demonstrates basic \LaTeX{} features including mathematical formulas, tables, and lists. The bibliography system is now set up to work seamlessly with Bibby AI's citation management. Feel free to modify it for your needs!

\bibliographystyle{plain}
\bibliography{cite}

\end{document}\documentclass{article}
\usepackage[utf8]{inputenc}
\usepackage{booktabs}
\usepackage{amsmath}
\usepackage{natbib}

\title{My Document}
\author{Your Name}
\date{\today}

\begin{document}

\maketitle

\section{Introduction}
This is a sample document demonstrating various \LaTeX{} features including mathematical equations, tables, and formatting. You can cite papers using the citation search feature, and they will be automatically added to your \texttt{cite.bib} file.

\section{Mathematical Expressions}
Here are some example equations:

\subsection{Quadratic Formula}
The quadratic formula is given by:
\[x = \frac{-b \pm \sqrt{b^2 - 4ac}}{2a}\]

\subsection{Matrix Example}
A sample matrix:
\[
\begin{pmatrix}
a_{11} & a_{12} & a_{13} \\
a_{21} & a_{22} & a_{23} \\
a_{31} & a_{32} & a_{33}
\end{pmatrix}
\]

\section{Tables}
Below is a sample table showing data:

\begin{table}[h]
\centering
\begin{tabular}{@{}lrrr@{}}
\toprule
\textbf{Category} & \textbf{Value 1} & \textbf{Value 2} & \textbf{Total} \\
\midrule
Item A & 10.5 & 20.3 & 30.8 \\
Item B & 15.2 & 25.1 & 40.3 \\
Item C & 22.1 & 12.7 & 34.8 \\
\bottomrule
\end{tabular}
\caption{Sample Data Table}
\label{tab:sample}
\end{table}

\section{Lists}
\subsection{Bullet Points}
\begin{itemize}
  \item First item
  \item Second item
  \item Third item with sub-items:
    \begin{itemize}
      \item Sub-item 1
      \item Sub-item 2
    \end{itemize}
\end{itemize}

\subsection{Numbered List}
\begin{enumerate}
  \item First step
  \item Second step
  \item Third step
\end{enumerate}

\section{Citations and References}
When you use the citation search feature, references will be automatically managed. For example, you might cite a paper like this: \cite{example2024}. All citations will be automatically added to your \texttt{cite.bib} file.

\section{Conclusion}
This template demonstrates basic \LaTeX{} features including mathematical formulas, tables, and lists. The bibliography system is now set up to work seamlessly with Bibby AI's citation management. Feel free to modify it for your needs!

\bibliographystyle{plain}
\bibliography{cite}

\end{document}\documentclass{article}
\usepackage[utf8]{inputenc}
\usepackage{booktabs}
\usepackage{amsmath}
\usepackage{natbib}

\title{My Document}
\author{Your Name}
\date{\today}

\begin{document}

\maketitle

\section{Introduction}
This is a sample document demonstrating various \LaTeX{} features including mathematical equations, tables, and formatting. You can cite papers using the citation search feature, and they will be automatically added to your \texttt{cite.bib} file.

\section{Mathematical Expressions}
Here are some example equations:

\subsection{Quadratic Formula}
The quadratic formula is given by:
\[x = \frac{-b \pm \sqrt{b^2 - 4ac}}{2a}\]

\subsection{Matrix Example}
A sample matrix:
\[
\begin{pmatrix}
a_{11} & a_{12} & a_{13} \\
a_{21} & a_{22} & a_{23} \\
a_{31} & a_{32} & a_{33}
\end{pmatrix}
\]

\section{Tables}
Below is a sample table showing data:

\begin{table}[h]
\centering
\begin{tabular}{@{}lrrr@{}}
\toprule
\textbf{Category} & \textbf{Value 1} & \textbf{Value 2} & \textbf{Total} \\
\midrule
Item A & 10.5 & 20.3 & 30.8 \\
Item B & 15.2 & 25.1 & 40.3 \\
Item C & 22.1 & 12.7 & 34.8 \\
\bottomrule
\end{tabular}
\caption{Sample Data Table}
\label{tab:sample}
\end{table}

\section{Lists}
\subsection{Bullet Points}
\begin{itemize}
  \item First item
  \item Second item
  \item Third item with sub-items:
    \begin{itemize}
      \item Sub-item 1
      \item Sub-item 2
    \end{itemize}
\end{itemize}

\subsection{Numbered List}
\begin{enumerate}
  \item First step
  \item Second step
  \item Third step
\end{enumerate}

\section{Citations and References}
When you use the citation search feature, references will be automatically managed. For example, you might cite a paper like this: \cite{example2024}. All citations will be automatically added to your \texttt{cite.bib} file.

\section{Conclusion}
This template demonstrates basic \LaTeX{} features including mathematical formulas, tables, and lists. The bibliography system is now set up to work seamlessly with Bibby AI's citation management. Feel free to modify it for your needs!

\bibliographystyle{plain}
\bibliography{cite}

\end{document}\documentclass{article}
\usepackage[utf8]{inputenc}
\usepackage{booktabs}
\usepackage{amsmath}
\usepackage{natbib}

\title{My Document}
\author{Your Name}
\date{\today}

\begin{document}

\maketitle

\section{Introduction}
This is a sample document demonstrating various \LaTeX{} features including mathematical equations, tables, and formatting. You can cite papers using the citation search feature, and they will be automatically added to your \texttt{cite.bib} file.

\section{Mathematical Expressions}
Here are some example equations:

\subsection{Quadratic Formula}
The quadratic formula is given by:
\[x = \frac{-b \pm \sqrt{b^2 - 4ac}}{2a}\]

\subsection{Matrix Example}
A sample matrix:
\[
\begin{pmatrix}
a_{11} & a_{12} & a_{13} \\
a_{21} & a_{22} & a_{23} \\
a_{31} & a_{32} & a_{33}
\end{pmatrix}
\]

\section{Tables}
Below is a sample table showing data:

\begin{table}[h]
\centering
\begin{tabular}{@{}lrrr@{}}
\toprule
\textbf{Category} & \textbf{Value 1} & \textbf{Value 2} & \textbf{Total} \\
\midrule
Item A & 10.5 & 20.3 & 30.8 \\
Item B & 15.2 & 25.1 & 40.3 \\
Item C & 22.1 & 12.7 & 34.8 \\
\bottomrule
\end{tabular}
\caption{Sample Data Table}
\label{tab:sample}
\end{table}

\section{Lists}
\subsection{Bullet Points}
\begin{itemize}
  \item First item
  \item Second item
  \item Third item with sub-items:
    \begin{itemize}
      \item Sub-item 1
      \item Sub-item 2
    \end{itemize}
\end{itemize}

\subsection{Numbered List}
\begin{enumerate}
  \item First step
  \item Second step
  \item Third step
\end{enumerate}

\section{Citations and References}
When you use the citation search feature, references will be automatically managed. For example, you might cite a paper like this: \cite{example2024}. All citations will be automatically added to your \texttt{cite.bib} file.

\section{Conclusion}
This template demonstrates basic \LaTeX{} features including mathematical formulas, tables, and lists. The bibliography system is now set up to work seamlessly with Bibby AI's citation management. Feel free to modify it for your needs!

\bibliographystyle{plain}
\bibliography{cite}

\end{document}\documentclass{article}
\usepackage[utf8]{inputenc}
\usepackage{booktabs}
\usepackage{amsmath}
\usepackage{natbib}

\title{My Document}
\author{Your Name}
\date{\today}

\begin{document}

\maketitle

\section{Introduction}
This is a sample document demonstrating various \LaTeX{} features including mathematical equations, tables, and formatting. You can cite papers using the citation search feature, and they will be automatically added to your \texttt{cite.bib} file.

\section{Mathematical Expressions}
Here are some example equations:

\subsection{Quadratic Formula}
The quadratic formula is given by:
\[x = \frac{-b \pm \sqrt{b^2 - 4ac}}{2a}\]

\subsection{Matrix Example}
A sample matrix:
\[
\begin{pmatrix}
a_{11} & a_{12} & a_{13} \\
a_{21} & a_{22} & a_{23} \\
a_{31} & a_{32} & a_{33}
\end{pmatrix}
\]

\section{Tables}
Below is a sample table showing data:

\begin{table}[h]
\centering
\begin{tabular}{@{}lrrr@{}}
\toprule
\textbf{Category} & \textbf{Value 1} & \textbf{Value 2} & \textbf{Total} \\
\midrule
Item A & 10.5 & 20.3 & 30.8 \\
Item B & 15.2 & 25.1 & 40.3 \\
Item C & 22.1 & 12.7 & 34.8 \\
\bottomrule
\end{tabular}
\caption{Sample Data Table}
\label{tab:sample}
\end{table}

\section{Lists}
\subsection{Bullet Points}
\begin{itemize}
  \item First item
  \item Second item
  \item Third item with sub-items:
    \begin{itemize}
      \item Sub-item 1
      \item Sub-item 2
    \end{itemize}
\end{itemize}

\subsection{Numbered List}
\begin{enumerate}
  \item First step
  \item Second step
  \item Third step
\end{enumerate}

\section{Citations and References}
When you use the citation search feature, references will be automatically managed. For example, you might cite a paper like this: \cite{example2024}. All citations will be automatically added to your \texttt{cite.bib} file.

\section{Conclusion}
This template demonstrates basic \LaTeX{} features including mathematical formulas, tables, and lists. The bibliography system is now set up to work seamlessly with Bibby AI's citation management. Feel free to modify it for your needs!

\bibliographystyle{plain}
\bibliography{cite}

\end{document}\documentclass{article}
\usepackage[utf8]{inputenc}
\usepackage{booktabs}
\usepackage{amsmath}
\usepackage{natbib}

\title{My Document}
\author{Your Name}
\date{\today}

\begin{document}

\maketitle

\section{Introduction}
This is a sample document demonstrating various \LaTeX{} features including mathematical equations, tables, and formatting. You can cite papers using the citation search feature, and they will be automatically added to your \texttt{cite.bib} file.

\section{Mathematical Expressions}
Here are some example equations:

\subsection{Quadratic Formula}
The quadratic formula is given by:
\[x = \frac{-b \pm \sqrt{b^2 - 4ac}}{2a}\]

\subsection{Matrix Example}
A sample matrix:
\[
\begin{pmatrix}
a_{11} & a_{12} & a_{13} \\
a_{21} & a_{22} & a_{23} \\
a_{31} & a_{32} & a_{33}
\end{pmatrix}
\]

\section{Tables}
Below is a sample table showing data:

\begin{table}[h]
\centering
\begin{tabular}{@{}lrrr@{}}
\toprule
\textbf{Category} & \textbf{Value 1} & \textbf{Value 2} & \textbf{Total} \\
\midrule
Item A & 10.5 & 20.3 & 30.8 \\
Item B & 15.2 & 25.1 & 40.3 \\
Item C & 22.1 & 12.7 & 34.8 \\
\bottomrule
\end{tabular}
\caption{Sample Data Table}
\label{tab:sample}
\end{table}

\section{Lists}
\subsection{Bullet Points}
\begin{itemize}
  \item First item
  \item Second item
  \item Third item with sub-items:
    \begin{itemize}
      \item Sub-item 1
      \item Sub-item 2
    \end{itemize}
\end{itemize}

\subsection{Numbered List}
\begin{enumerate}
  \item First step
  \item Second step
  \item Third step
\end{enumerate}

\section{Citations and References}
When you use the citation search feature, references will be automatically managed. For example, you might cite a paper like this: \cite{example2024}. All citations will be automatically added to your \texttt{cite.bib} file.

\section{Conclusion}
This template demonstrates basic \LaTeX{} features including mathematical formulas, tables, and lists. The bibliography system is now set up to work seamlessly with Bibby AI's citation management. Feel free to modify it for your needs!

\bibliographystyle{plain}
\bibliography{cite}

\end{document}\documentclass{article}
\usepackage[utf8]{inputenc}
\usepackage{booktabs}
\usepackage{amsmath}
\usepackage{natbib}

\title{My Document}
\author{Your Name}
\date{\today}

\begin{document}

\maketitle

\section{Introduction}
This is a sample document demonstrating various \LaTeX{} features including mathematical equations, tables, and formatting. You can cite papers using the citation search feature, and they will be automatically added to your \texttt{cite.bib} file.

\section{Mathematical Expressions}
Here are some example equations:

\subsection{Quadratic Formula}
The quadratic formula is given by:
\[x = \frac{-b \pm \sqrt{b^2 - 4ac}}{2a}\]

\subsection{Matrix Example}
A sample matrix:
\[
\begin{pmatrix}
a_{11} & a_{12} & a_{13} \\
a_{21} & a_{22} & a_{23} \\
a_{31} & a_{32} & a_{33}
\end{pmatrix}
\]

\section{Tables}
Below is a sample table showing data:

\begin{table}[h]
\centering
\begin{tabular}{@{}lrrr@{}}
\toprule
\textbf{Category} & \textbf{Value 1} & \textbf{Value 2} & \textbf{Total} \\
\midrule
Item A & 10.5 & 20.3 & 30.8 \\
Item B & 15.2 & 25.1 & 40.3 \\
Item C & 22.1 & 12.7 & 34.8 \\
\bottomrule
\end{tabular}
\caption{Sample Data Table}
\label{tab:sample}
\end{table}

\section{Lists}
\subsection{Bullet Points}
\begin{itemize}
  \item First item
  \item Second item
  \item Third item with sub-items:
    \begin{itemize}
      \item Sub-item 1
      \item Sub-item 2
    \end{itemize}
\end{itemize}

\subsection{Numbered List}
\begin{enumerate}
  \item First step
  \item Second step
  \item Third step
\end{enumerate}

\section{Citations and References}
When you use the citation search feature, references will be automatically managed. For example, you might cite a paper like this: \cite{example2024}. All citations will be automatically added to your \texttt{cite.bib} file.

\section{Conclusion}
This template demonstrates basic \LaTeX{} features including mathematical formulas, tables, and lists. The bibliography system is now set up to work seamlessly with Bibby AI's citation management. Feel free to modify it for your needs!

\bibliographystyle{plain}
\bibliography{cite}

\end{document}\documentclass{article}
\usepackage[utf8]{inputenc}
\usepackage{booktabs}
\usepackage{amsmath}
\usepackage{natbib}

\title{My Document}
\author{Your Name}
\date{\today}

\begin{document}

\maketitle

\section{Introduction}
This is a sample document demonstrating various \LaTeX{} features including mathematical equations, tables, and formatting. You can cite papers using the citation search feature, and they will be automatically added to your \texttt{cite.bib} file.

\section{Mathematical Expressions}
Here are some example equations:

\subsection{Quadratic Formula}
The quadratic formula is given by:
\[x = \frac{-b \pm \sqrt{b^2 - 4ac}}{2a}\]

\subsection{Matrix Example}
A sample matrix:
\[
\begin{pmatrix}
a_{11} & a_{12} & a_{13} \\
a_{21} & a_{22} & a_{23} \\
a_{31} & a_{32} & a_{33}
\end{pmatrix}
\]

\section{Tables}
Below is a sample table showing data:

\begin{table}[h]
\centering
\begin{tabular}{@{}lrrr@{}}
\toprule
\textbf{Category} & \textbf{Value 1} & \textbf{Value 2} & \textbf{Total} \\
\midrule
Item A & 10.5 & 20.3 & 30.8 \\
Item B & 15.2 & 25.1 & 40.3 \\
Item C & 22.1 & 12.7 & 34.8 \\
\bottomrule
\end{tabular}
\caption{Sample Data Table}
\label{tab:sample}
\end{table}

\section{Lists}
\subsection{Bullet Points}
\begin{itemize}
  \item First item
  \item Second item
  \item Third item with sub-items:
    \begin{itemize}
      \item Sub-item 1
      \item Sub-item 2
    \end{itemize}
\end{itemize}

\subsection{Numbered List}
\begin{enumerate}
  \item First step
  \item Second step
  \item Third step
\end{enumerate}

\section{Citations and References}
When you use the citation search feature, references will be automatically managed. For example, you might cite a paper like this: \cite{example2024}. All citations will be automatically added to your \texttt{cite.bib} file.

\section{Conclusion}
This template demonstrates basic \LaTeX{} features including mathematical formulas, tables, and lists. The bibliography system is now set up to work seamlessly with Bibby AI's citation management. Feel free to modify it for your needs!

\bibliographystyle{plain}
\bibliography{cite}

\end{document}\documentclass{article}
\usepackage[utf8]{inputenc}
\usepackage{booktabs}
\usepackage{amsmath}
\usepackage{natbib}

\title{My Document}
\author{Your Name}
\date{\today}

\begin{document}

\maketitle

\section{Introduction}
This is a sample document demonstrating various \LaTeX{} features including mathematical equations, tables, and formatting. You can cite papers using the citation search feature, and they will be automatically added to your \texttt{cite.bib} file.

\section{Mathematical Expressions}
Here are some example equations:

\subsection{Quadratic Formula}
The quadratic formula is given by:
\[x = \frac{-b \pm \sqrt{b^2 - 4ac}}{2a}\]

\subsection{Matrix Example}
A sample matrix:
\[
\begin{pmatrix}
a_{11} & a_{12} & a_{13} \\
a_{21} & a_{22} & a_{23} \\
a_{31} & a_{32} & a_{33}
\end{pmatrix}
\]

\section{Tables}
Below is a sample table showing data:

\begin{table}[h]
\centering
\begin{tabular}{@{}lrrr@{}}
\toprule
\textbf{Category} & \textbf{Value 1} & \textbf{Value 2} & \textbf{Total} \\
\midrule
Item A & 10.5 & 20.3 & 30.8 \\
Item B & 15.2 & 25.1 & 40.3 \\
Item C & 22.1 & 12.7 & 34.8 \\
\bottomrule
\end{tabular}
\caption{Sample Data Table}
\label{tab:sample}
\end{table}

\section{Lists}
\subsection{Bullet Points}
\begin{itemize}
  \item First item
  \item Second item
  \item Third item with sub-items:
    \begin{itemize}
      \item Sub-item 1
      \item Sub-item 2
    \end{itemize}
\end{itemize}

\subsection{Numbered List}
\begin{enumerate}
  \item First step
  \item Second step
  \item Third step
\end{enumerate}

\section{Citations and References}
When you use the citation search feature, references will be automatically managed. For example, you might cite a paper like this: \cite{example2024}. All citations will be automatically added to your \texttt{cite.bib} file.

\section{Conclusion}
This template demonstrates basic \LaTeX{} features including mathematical formulas, tables, and lists. The bibliography system is now set up to work seamlessly with Bibby AI's citation management. Feel free to modify it for your needs!

\bibliographystyle{plain}
\bibliography{cite}

\end{document}\documentclass{article}
\usepackage[utf8]{inputenc}
\usepackage{booktabs}
\usepackage{amsmath}
\usepackage{natbib}

\title{My Document}
\author{Your Name}
\date{\today}

\begin{document}

\maketitle

\section{Introduction}
This is a sample document demonstrating various \LaTeX{} features including mathematical equations, tables, and formatting. You can cite papers using the citation search feature, and they will be automatically added to your \texttt{cite.bib} file.

\section{Mathematical Expressions}
Here are some example equations:

\subsection{Quadratic Formula}
The quadratic formula is given by:
\[x = \frac{-b \pm \sqrt{b^2 - 4ac}}{2a}\]

\subsection{Matrix Example}
A sample matrix:
\[
\begin{pmatrix}
a_{11} & a_{12} & a_{13} \\
a_{21} & a_{22} & a_{23} \\
a_{31} & a_{32} & a_{33}
\end{pmatrix}
\]

\section{Tables}
Below is a sample table showing data:

\begin{table}[h]
\centering
\begin{tabular}{@{}lrrr@{}}
\toprule
\textbf{Category} & \textbf{Value 1} & \textbf{Value 2} & \textbf{Total} \\
\midrule
Item A & 10.5 & 20.3 & 30.8 \\
Item B & 15.2 & 25.1 & 40.3 \\
Item C & 22.1 & 12.7 & 34.8 \\
\bottomrule
\end{tabular}
\caption{Sample Data Table}
\label{tab:sample}
\end{table}

\section{Lists}
\subsection{Bullet Points}
\begin{itemize}
  \item First item
  \item Second item
  \item Third item with sub-items:
    \begin{itemize}
      \item Sub-item 1
      \item Sub-item 2
    \end{itemize}
\end{itemize}

\subsection{Numbered List}
\begin{enumerate}
  \item First step
  \item Second step
  \item Third step
\end{enumerate}

\section{Citations and References}
When you use the citation search feature, references will be automatically managed. For example, you might cite a paper like this: \cite{example2024}. All citations will be automatically added to your \texttt{cite.bib} file.

\section{Conclusion}
This template demonstrates basic \LaTeX{} features including mathematical formulas, tables, and lists. The bibliography system is now set up to work seamlessly with Bibby AI's citation management. Feel free to modify it for your needs!

\bibliographystyle{plain}
\bibliography{cite}

\end{document}\documentclass{article}
\usepackage[utf8]{inputenc}
\usepackage{booktabs}
\usepackage{amsmath}
\usepackage{natbib}

\title{My Document}
\author{Your Name}
\date{\today}

\begin{document}

\maketitle

\section{Introduction}
This is a sample document demonstrating various \LaTeX{} features including mathematical equations, tables, and formatting. You can cite papers using the citation search feature, and they will be automatically added to your \texttt{cite.bib} file.

\section{Mathematical Expressions}
Here are some example equations:

\subsection{Quadratic Formula}
The quadratic formula is given by:
\[x = \frac{-b \pm \sqrt{b^2 - 4ac}}{2a}\]

\subsection{Matrix Example}
A sample matrix:
\[
\begin{pmatrix}
a_{11} & a_{12} & a_{13} \\
a_{21} & a_{22} & a_{23} \\
a_{31} & a_{32} & a_{33}
\end{pmatrix}
\]

\section{Tables}
Below is a sample table showing data:

\begin{table}[h]
\centering
\begin{tabular}{@{}lrrr@{}}
\toprule
\textbf{Category} & \textbf{Value 1} & \textbf{Value 2} & \textbf{Total} \\
\midrule
Item A & 10.5 & 20.3 & 30.8 \\
Item B & 15.2 & 25.1 & 40.3 \\
Item C & 22.1 & 12.7 & 34.8 \\
\bottomrule
\end{tabular}
\caption{Sample Data Table}
\label{tab:sample}
\end{table}

\section{Lists}
\subsection{Bullet Points}
\begin{itemize}
  \item First item
  \item Second item
  \item Third item with sub-items:
    \begin{itemize}
      \item Sub-item 1
      \item Sub-item 2
    \end{itemize}
\end{itemize}

\subsection{Numbered List}
\begin{enumerate}
  \item First step
  \item Second step
  \item Third step
\end{enumerate}

\section{Citations and References}
When you use the citation search feature, references will be automatically managed. For example, you might cite a paper like this: \cite{example2024}. All citations will be automatically added to your \texttt{cite.bib} file.

\section{Conclusion}
This template demonstrates basic \LaTeX{} features including mathematical formulas, tables, and lists. The bibliography system is now set up to work seamlessly with Bibby AI's citation management. Feel free to modify it for your needs!

\bibliographystyle{plain}
\bibliography{cite}

\end{document}\documentclass{article}
\usepackage[utf8]{inputenc}
\usepackage{booktabs}
\usepackage{amsmath}
\usepackage{natbib}

\title{My Document}
\author{Your Name}
\date{\today}

\begin{document}

\maketitle

\section{Introduction}
This is a sample document demonstrating various \LaTeX{} features including mathematical equations, tables, and formatting. You can cite papers using the citation search feature, and they will be automatically added to your \texttt{cite.bib} file.

\section{Mathematical Expressions}
Here are some example equations:

\subsection{Quadratic Formula}
The quadratic formula is given by:
\[x = \frac{-b \pm \sqrt{b^2 - 4ac}}{2a}\]

\subsection{Matrix Example}
A sample matrix:
\[
\begin{pmatrix}
a_{11} & a_{12} & a_{13} \\
a_{21} & a_{22} & a_{23} \\
a_{31} & a_{32} & a_{33}
\end{pmatrix}
\]

\section{Tables}
Below is a sample table showing data:

\begin{table}[h]
\centering
\begin{tabular}{@{}lrrr@{}}
\toprule
\textbf{Category} & \textbf{Value 1} & \textbf{Value 2} & \textbf{Total} \\
\midrule
Item A & 10.5 & 20.3 & 30.8 \\
Item B & 15.2 & 25.1 & 40.3 \\
Item C & 22.1 & 12.7 & 34.8 \\
\bottomrule
\end{tabular}
\caption{Sample Data Table}
\label{tab:sample}
\end{table}

\section{Lists}
\subsection{Bullet Points}
\begin{itemize}
  \item First item
  \item Second item
  \item Third item with sub-items:
    \begin{itemize}
      \item Sub-item 1
      \item Sub-item 2
    \end{itemize}
\end{itemize}

\subsection{Numbered List}
\begin{enumerate}
  \item First step
  \item Second step
  \item Third step
\end{enumerate}

\section{Citations and References}
When you use the citation search feature, references will be automatically managed. For example, you might cite a paper like this: \cite{example2024}. All citations will be automatically added to your \texttt{cite.bib} file.

\section{Conclusion}
This template demonstrates basic \LaTeX{} features including mathematical formulas, tables, and lists. The bibliography system is now set up to work seamlessly with Bibby AI's citation management. Feel free to modify it for your needs!

\bibliographystyle{plain}
\bibliography{cite}

\end{document}\documentclass{article}
\usepackage[utf8]{inputenc}
\usepackage{booktabs}
\usepackage{amsmath}
\usepackage{natbib}

\title{My Document}
\author{Your Name}
\date{\today}

\begin{document}

\maketitle

\section{Introduction}
This is a sample document demonstrating various \LaTeX{} features including mathematical equations, tables, and formatting. You can cite papers using the citation search feature, and they will be automatically added to your \texttt{cite.bib} file.

\section{Mathematical Expressions}
Here are some example equations:

\subsection{Quadratic Formula}
The quadratic formula is given by:
\[x = \frac{-b \pm \sqrt{b^2 - 4ac}}{2a}\]

\subsection{Matrix Example}
A sample matrix:
\[
\begin{pmatrix}
a_{11} & a_{12} & a_{13} \\
a_{21} & a_{22} & a_{23} \\
a_{31} & a_{32} & a_{33}
\end{pmatrix}
\]

\section{Tables}
Below is a sample table showing data:

\begin{table}[h]
\centering
\begin{tabular}{@{}lrrr@{}}
\toprule
\textbf{Category} & \textbf{Value 1} & \textbf{Value 2} & \textbf{Total} \\
\midrule
Item A & 10.5 & 20.3 & 30.8 \\
Item B & 15.2 & 25.1 & 40.3 \\
Item C & 22.1 & 12.7 & 34.8 \\
\bottomrule
\end{tabular}
\caption{Sample Data Table}
\label{tab:sample}
\end{table}

\section{Lists}
\subsection{Bullet Points}
\begin{itemize}
  \item First item
  \item Second item
  \item Third item with sub-items:
    \begin{itemize}
      \item Sub-item 1
      \item Sub-item 2
    \end{itemize}
\end{itemize}

\subsection{Numbered List}
\begin{enumerate}
  \item First step
  \item Second step
  \item Third step
\end{enumerate}

\section{Citations and References}
When you use the citation search feature, references will be automatically managed. For example, you might cite a paper like this: \cite{example2024}. All citations will be automatically added to your \texttt{cite.bib} file.

\section{Conclusion}
This template demonstrates basic \LaTeX{} features including mathematical formulas, tables, and lists. The bibliography system is now set up to work seamlessly with Bibby AI's citation management. Feel free to modify it for your needs!

\bibliographystyle{plain}
\bibliography{cite}

\end{document}\documentclass{article}
\usepackage[utf8]{inputenc}
\usepackage{booktabs}
\usepackage{amsmath}
\usepackage{natbib}

\title{My Document}
\author{Your Name}
\date{\today}

\begin{document}

\maketitle

\section{Introduction}
This is a sample document demonstrating various \LaTeX{} features including mathematical equations, tables, and formatting. You can cite papers using the citation search feature, and they will be automatically added to your \texttt{cite.bib} file.

\section{Mathematical Expressions}
Here are some example equations:

\subsection{Quadratic Formula}
The quadratic formula is given by:
\[x = \frac{-b \pm \sqrt{b^2 - 4ac}}{2a}\]

\subsection{Matrix Example}
A sample matrix:
\[
\begin{pmatrix}
a_{11} & a_{12} & a_{13} \\
a_{21} & a_{22} & a_{23} \\
a_{31} & a_{32} & a_{33}
\end{pmatrix}
\]

\section{Tables}
Below is a sample table showing data:

\begin{table}[h]
\centering
\begin{tabular}{@{}lrrr@{}}
\toprule
\textbf{Category} & \textbf{Value 1} & \textbf{Value 2} & \textbf{Total} \\
\midrule
Item A & 10.5 & 20.3 & 30.8 \\
Item B & 15.2 & 25.1 & 40.3 \\
Item C & 22.1 & 12.7 & 34.8 \\
\bottomrule
\end{tabular}
\caption{Sample Data Table}
\label{tab:sample}
\end{table}

\section{Lists}
\subsection{Bullet Points}
\begin{itemize}
  \item First item
  \item Second item
  \item Third item with sub-items:
    \begin{itemize}
      \item Sub-item 1
      \item Sub-item 2
    \end{itemize}
\end{itemize}

\subsection{Numbered List}
\begin{enumerate}
  \item First step
  \item Second step
  \item Third step
\end{enumerate}

\section{Citations and References}
When you use the citation search feature, references will be automatically managed. For example, you might cite a paper like this: \cite{example2024}. All citations will be automatically added to your \texttt{cite.bib} file.

\section{Conclusion}
This template demonstrates basic \LaTeX{} features including mathematical formulas, tables, and lists. The bibliography system is now set up to work seamlessly with Bibby AI's citation management. Feel free to modify it for your needs!

\bibliographystyle{plain}
\bibliography{cite}

\end{document}\documentclass{article}
\usepackage[utf8]{inputenc}
\usepackage{booktabs}
\usepackage{amsmath}
\usepackage{natbib}

\title{My Document}
\author{Your Name}
\date{\today}

\begin{document}

\maketitle

\section{Introduction}
This is a sample document demonstrating various \LaTeX{} features including mathematical equations, tables, and formatting. You can cite papers using the citation search feature, and they will be automatically added to your \texttt{cite.bib} file.

\section{Mathematical Expressions}
Here are some example equations:

\subsection{Quadratic Formula}
The quadratic formula is given by:
\[x = \frac{-b \pm \sqrt{b^2 - 4ac}}{2a}\]

\subsection{Matrix Example}
A sample matrix:
\[
\begin{pmatrix}
a_{11} & a_{12} & a_{13} \\
a_{21} & a_{22} & a_{23} \\
a_{31} & a_{32} & a_{33}
\end{pmatrix}
\]

\section{Tables}
Below is a sample table showing data:

\begin{table}[h]
\centering
\begin{tabular}{@{}lrrr@{}}
\toprule
\textbf{Category} & \textbf{Value 1} & \textbf{Value 2} & \textbf{Total} \\
\midrule
Item A & 10.5 & 20.3 & 30.8 \\
Item B & 15.2 & 25.1 & 40.3 \\
Item C & 22.1 & 12.7 & 34.8 \\
\bottomrule
\end{tabular}
\caption{Sample Data Table}
\label{tab:sample}
\end{table}

\section{Lists}
\subsection{Bullet Points}
\begin{itemize}
  \item First item
  \item Second item
  \item Third item with sub-items:
    \begin{itemize}
      \item Sub-item 1
      \item Sub-item 2
    \end{itemize}
\end{itemize}

\subsection{Numbered List}
\begin{enumerate}
  \item First step
  \item Second step
  \item Third step
\end{enumerate}

\section{Citations and References}
When you use the citation search feature, references will be automatically managed. For example, you might cite a paper like this: \cite{example2024}. All citations will be automatically added to your \texttt{cite.bib} file.

\section{Conclusion}
This template demonstrates basic \LaTeX{} features including mathematical formulas, tables, and lists. The bibliography system is now set up to work seamlessly with Bibby AI's citation management. Feel free to modify it for your needs!

\bibliographystyle{plain}
\bibliography{cite}

\end{document}\documentclass{article}
\usepackage[utf8]{inputenc}
\usepackage{booktabs}
\usepackage{amsmath}
\usepackage{natbib}

\title{My Document}
\author{Your Name}
\date{\today}

\begin{document}

\maketitle

\section{Introduction}
This is a sample document demonstrating various \LaTeX{} features including mathematical equations, tables, and formatting. You can cite papers using the citation search feature, and they will be automatically added to your \texttt{cite.bib} file.

\section{Mathematical Expressions}
Here are some example equations:

\subsection{Quadratic Formula}
The quadratic formula is given by:
\[x = \frac{-b \pm \sqrt{b^2 - 4ac}}{2a}\]

\subsection{Matrix Example}
A sample matrix:
\[
\begin{pmatrix}
a_{11} & a_{12} & a_{13} \\
a_{21} & a_{22} & a_{23} \\
a_{31} & a_{32} & a_{33}
\end{pmatrix}
\]

\section{Tables}
Below is a sample table showing data:

\begin{table}[h]
\centering
\begin{tabular}{@{}lrrr@{}}
\toprule
\textbf{Category} & \textbf{Value 1} & \textbf{Value 2} & \textbf{Total} \\
\midrule
Item A & 10.5 & 20.3 & 30.8 \\
Item B & 15.2 & 25.1 & 40.3 \\
Item C & 22.1 & 12.7 & 34.8 \\
\bottomrule
\end{tabular}
\caption{Sample Data Table}
\label{tab:sample}
\end{table}

\section{Lists}
\subsection{Bullet Points}
\begin{itemize}
  \item First item
  \item Second item
  \item Third item with sub-items:
    \begin{itemize}
      \item Sub-item 1
      \item Sub-item 2
    \end{itemize}
\end{itemize}

\subsection{Numbered List}
\begin{enumerate}
  \item First step
  \item Second step
  \item Third step
\end{enumerate}

\section{Citations and References}
When you use the citation search feature, references will be automatically managed. For example, you might cite a paper like this: \cite{example2024}. All citations will be automatically added to your \texttt{cite.bib} file.

\section{Conclusion}
This template demonstrates basic \LaTeX{} features including mathematical formulas, tables, and lists. The bibliography system is now set up to work seamlessly with Bibby AI's citation management. Feel free to modify it for your needs!

\bibliographystyle{plain}
\bibliography{cite}

\end{document}\documentclass{article}
\usepackage[utf8]{inputenc}
\usepackage{booktabs}
\usepackage{amsmath}
\usepackage{natbib}

\title{My Document}
\author{Your Name}
\date{\today}

\begin{document}

\maketitle

\section{Introduction}
This is a sample document demonstrating various \LaTeX{} features including mathematical equations, tables, and formatting. You can cite papers using the citation search feature, and they will be automatically added to your \texttt{cite.bib} file.

\section{Mathematical Expressions}
Here are some example equations:

\subsection{Quadratic Formula}
The quadratic formula is given by:
\[x = \frac{-b \pm \sqrt{b^2 - 4ac}}{2a}\]

\subsection{Matrix Example}
A sample matrix:
\[
\begin{pmatrix}
a_{11} & a_{12} & a_{13} \\
a_{21} & a_{22} & a_{23} \\
a_{31} & a_{32} & a_{33}
\end{pmatrix}
\]

\section{Tables}
Below is a sample table showing data:

\begin{table}[h]
\centering
\begin{tabular}{@{}lrrr@{}}
\toprule
\textbf{Category} & \textbf{Value 1} & \textbf{Value 2} & \textbf{Total} \\
\midrule
Item A & 10.5 & 20.3 & 30.8 \\
Item B & 15.2 & 25.1 & 40.3 \\
Item C & 22.1 & 12.7 & 34.8 \\
\bottomrule
\end{tabular}
\caption{Sample Data Table}
\label{tab:sample}
\end{table}

\section{Lists}
\subsection{Bullet Points}
\begin{itemize}
  \item First item
  \item Second item
  \item Third item with sub-items:
    \begin{itemize}
      \item Sub-item 1
      \item Sub-item 2
    \end{itemize}
\end{itemize}

\subsection{Numbered List}
\begin{enumerate}
  \item First step
  \item Second step
  \item Third step
\end{enumerate}

\section{Citations and References}
When you use the citation search feature, references will be automatically managed. For example, you might cite a paper like this: \cite{example2024}. All citations will be automatically added to your \texttt{cite.bib} file.

\section{Conclusion}
This template demonstrates basic \LaTeX{} features including mathematical formulas, tables, and lists. The bibliography system is now set up to work seamlessly with Bibby AI's citation management. Feel free to modify it for your needs!

\bibliographystyle{plain}
\bibliography{cite}

\end{document}\documentclass{article}
\usepackage[utf8]{inputenc}
\usepackage{booktabs}
\usepackage{amsmath}
\usepackage{natbib}

\title{My Document}
\author{Your Name}
\date{\today}

\begin{document}

\maketitle

\section{Introduction}
This is a sample document demonstrating various \LaTeX{} features including mathematical equations, tables, and formatting. You can cite papers using the citation search feature, and they will be automatically added to your \texttt{cite.bib} file.

\section{Mathematical Expressions}
Here are some example equations:

\subsection{Quadratic Formula}
The quadratic formula is given by:
\[x = \frac{-b \pm \sqrt{b^2 - 4ac}}{2a}\]

\subsection{Matrix Example}
A sample matrix:
\[
\begin{pmatrix}
a_{11} & a_{12} & a_{13} \\
a_{21} & a_{22} & a_{23} \\
a_{31} & a_{32} & a_{33}
\end{pmatrix}
\]

\section{Tables}
Below is a sample table showing data:

\begin{table}[h]
\centering
\begin{tabular}{@{}lrrr@{}}
\toprule
\textbf{Category} & \textbf{Value 1} & \textbf{Value 2} & \textbf{Total} \\
\midrule
Item A & 10.5 & 20.3 & 30.8 \\
Item B & 15.2 & 25.1 & 40.3 \\
Item C & 22.1 & 12.7 & 34.8 \\
\bottomrule
\end{tabular}
\caption{Sample Data Table}
\label{tab:sample}
\end{table}

\section{Lists}
\subsection{Bullet Points}
\begin{itemize}
  \item First item
  \item Second item
  \item Third item with sub-items:
    \begin{itemize}
      \item Sub-item 1
      \item Sub-item 2
    \end{itemize}
\end{itemize}

\subsection{Numbered List}
\begin{enumerate}
  \item First step
  \item Second step
  \item Third step
\end{enumerate}

\section{Citations and References}
When you use the citation search feature, references will be automatically managed. For example, you might cite a paper like this: \cite{example2024}. All citations will be automatically added to your \texttt{cite.bib} file.

\section{Conclusion}
This template demonstrates basic \LaTeX{} features including mathematical formulas, tables, and lists. The bibliography system is now set up to work seamlessly with Bibby AI's citation management. Feel free to modify it for your needs!

\bibliographystyle{plain}
\bibliography{cite}

\end{document}\documentclass{article}
\usepackage[utf8]{inputenc}
\usepackage{booktabs}
\usepackage{amsmath}
\usepackage{natbib}

\title{My Document}
\author{Your Name}
\date{\today}

\begin{document}

\maketitle

\section{Introduction}
This is a sample document demonstrating various \LaTeX{} features including mathematical equations, tables, and formatting. You can cite papers using the citation search feature, and they will be automatically added to your \texttt{cite.bib} file.

\section{Mathematical Expressions}
Here are some example equations:

\subsection{Quadratic Formula}
The quadratic formula is given by:
\[x = \frac{-b \pm \sqrt{b^2 - 4ac}}{2a}\]

\subsection{Matrix Example}
A sample matrix:
\[
\begin{pmatrix}
a_{11} & a_{12} & a_{13} \\
a_{21} & a_{22} & a_{23} \\
a_{31} & a_{32} & a_{33}
\end{pmatrix}
\]

\section{Tables}
Below is a sample table showing data:

\begin{table}[h]
\centering
\begin{tabular}{@{}lrrr@{}}
\toprule
\textbf{Category} & \textbf{Value 1} & \textbf{Value 2} & \textbf{Total} \\
\midrule
Item A & 10.5 & 20.3 & 30.8 \\
Item B & 15.2 & 25.1 & 40.3 \\
Item C & 22.1 & 12.7 & 34.8 \\
\bottomrule
\end{tabular}
\caption{Sample Data Table}
\label{tab:sample}
\end{table}

\section{Lists}
\subsection{Bullet Points}
\begin{itemize}
  \item First item
  \item Second item
  \item Third item with sub-items:
    \begin{itemize}
      \item Sub-item 1
      \item Sub-item 2
    \end{itemize}
\end{itemize}

\subsection{Numbered List}
\begin{enumerate}
  \item First step
  \item Second step
  \item Third step
\end{enumerate}

\section{Citations and References}
When you use the citation search feature, references will be automatically managed. For example, you might cite a paper like this: \cite{example2024}. All citations will be automatically added to your \texttt{cite.bib} file.

\section{Conclusion}
This template demonstrates basic \LaTeX{} features including mathematical formulas, tables, and lists. The bibliography system is now set up to work seamlessly with Bibby AI's citation management. Feel free to modify it for your needs!

\bibliographystyle{plain}
\bibliography{cite}

\end{document}\documentclass{article}
\usepackage[utf8]{inputenc}
\usepackage{booktabs}
\usepackage{amsmath}
\usepackage{natbib}

\title{My Document}
\author{Your Name}
\date{\today}

\begin{document}

\maketitle

\section{Introduction}
This is a sample document demonstrating various \LaTeX{} features including mathematical equations, tables, and formatting. You can cite papers using the citation search feature, and they will be automatically added to your \texttt{cite.bib} file.

\section{Mathematical Expressions}
Here are some example equations:

\subsection{Quadratic Formula}
The quadratic formula is given by:
\[x = \frac{-b \pm \sqrt{b^2 - 4ac}}{2a}\]

\subsection{Matrix Example}
A sample matrix:
\[
\begin{pmatrix}
a_{11} & a_{12} & a_{13} \\
a_{21} & a_{22} & a_{23} \\
a_{31} & a_{32} & a_{33}
\end{pmatrix}
\]

\section{Tables}
Below is a sample table showing data:

\begin{table}[h]
\centering
\begin{tabular}{@{}lrrr@{}}
\toprule
\textbf{Category} & \textbf{Value 1} & \textbf{Value 2} & \textbf{Total} \\
\midrule
Item A & 10.5 & 20.3 & 30.8 \\
Item B & 15.2 & 25.1 & 40.3 \\
Item C & 22.1 & 12.7 & 34.8 \\
\bottomrule
\end{tabular}
\caption{Sample Data Table}
\label{tab:sample}
\end{table}

\section{Lists}
\subsection{Bullet Points}
\begin{itemize}
  \item First item
  \item Second item
  \item Third item with sub-items:
    \begin{itemize}
      \item Sub-item 1
      \item Sub-item 2
    \end{itemize}
\end{itemize}

\subsection{Numbered List}
\begin{enumerate}
  \item First step
  \item Second step
  \item Third step
\end{enumerate}

\section{Citations and References}
When you use the citation search feature, references will be automatically managed. For example, you might cite a paper like this: \cite{example2024}. All citations will be automatically added to your \texttt{cite.bib} file.

\section{Conclusion}
This template demonstrates basic \LaTeX{} features including mathematical formulas, tables, and lists. The bibliography system is now set up to work seamlessly with Bibby AI's citation management. Feel free to modify it for your needs!

\bibliographystyle{plain}
\bibliography{cite}

\end{document}\documentclass{article}
\usepackage[utf8]{inputenc}
\usepackage{booktabs}
\usepackage{amsmath}
\usepackage{natbib}

\title{My Document}
\author{Your Name}
\date{\today}

\begin{document}

\maketitle

\section{Introduction}
This is a sample document demonstrating various \LaTeX{} features including mathematical equations, tables, and formatting. You can cite papers using the citation search feature, and they will be automatically added to your \texttt{cite.bib} file.

\section{Mathematical Expressions}
Here are some example equations:

\subsection{Quadratic Formula}
The quadratic formula is given by:
\[x = \frac{-b \pm \sqrt{b^2 - 4ac}}{2a}\]

\subsection{Matrix Example}
A sample matrix:
\[
\begin{pmatrix}
a_{11} & a_{12} & a_{13} \\
a_{21} & a_{22} & a_{23} \\
a_{31} & a_{32} & a_{33}
\end{pmatrix}
\]

\section{Tables}
Below is a sample table showing data:

\begin{table}[h]
\centering
\begin{tabular}{@{}lrrr@{}}
\toprule
\textbf{Category} & \textbf{Value 1} & \textbf{Value 2} & \textbf{Total} \\
\midrule
Item A & 10.5 & 20.3 & 30.8 \\
Item B & 15.2 & 25.1 & 40.3 \\
Item C & 22.1 & 12.7 & 34.8 \\
\bottomrule
\end{tabular}
\caption{Sample Data Table}
\label{tab:sample}
\end{table}

\section{Lists}
\subsection{Bullet Points}
\begin{itemize}
  \item First item
  \item Second item
  \item Third item with sub-items:
    \begin{itemize}
      \item Sub-item 1
      \item Sub-item 2
    \end{itemize}
\end{itemize}

\subsection{Numbered List}
\begin{enumerate}
  \item First step
  \item Second step
  \item Third step
\end{enumerate}

\section{Citations and References}
When you use the citation search feature, references will be automatically managed. For example, you might cite a paper like this: \cite{example2024}. All citations will be automatically added to your \texttt{cite.bib} file.

\section{Conclusion}
This template demonstrates basic \LaTeX{} features including mathematical formulas, tables, and lists. The bibliography system is now set up to work seamlessly with Bibby AI's citation management. Feel free to modify it for your needs!

\bibliographystyle{plain}
\bibliography{cite}

\end{document}\documentclass{article}
\usepackage[utf8]{inputenc}
\usepackage{booktabs}
\usepackage{amsmath}
\usepackage{natbib}

\title{My Document}
\author{Your Name}
\date{\today}

\begin{document}

\maketitle

\section{Introduction}
This is a sample document demonstrating various \LaTeX{} features including mathematical equations, tables, and formatting. You can cite papers using the citation search feature, and they will be automatically added to your \texttt{cite.bib} file.

\section{Mathematical Expressions}
Here are some example equations:

\subsection{Quadratic Formula}
The quadratic formula is given by:
\[x = \frac{-b \pm \sqrt{b^2 - 4ac}}{2a}\]

\subsection{Matrix Example}
A sample matrix:
\[
\begin{pmatrix}
a_{11} & a_{12} & a_{13} \\
a_{21} & a_{22} & a_{23} \\
a_{31} & a_{32} & a_{33}
\end{pmatrix}
\]

\section{Tables}
Below is a sample table showing data:

\begin{table}[h]
\centering
\begin{tabular}{@{}lrrr@{}}
\toprule
\textbf{Category} & \textbf{Value 1} & \textbf{Value 2} & \textbf{Total} \\
\midrule
Item A & 10.5 & 20.3 & 30.8 \\
Item B & 15.2 & 25.1 & 40.3 \\
Item C & 22.1 & 12.7 & 34.8 \\
\bottomrule
\end{tabular}
\caption{Sample Data Table}
\label{tab:sample}
\end{table}

\section{Lists}
\subsection{Bullet Points}
\begin{itemize}
  \item First item
  \item Second item
  \item Third item with sub-items:
    \begin{itemize}
      \item Sub-item 1
      \item Sub-item 2
    \end{itemize}
\end{itemize}

\subsection{Numbered List}
\begin{enumerate}
  \item First step
  \item Second step
  \item Third step
\end{enumerate}

\section{Citations and References}
When you use the citation search feature, references will be automatically managed. For example, you might cite a paper like this: \cite{example2024}. All citations will be automatically added to your \texttt{cite.bib} file.

\section{Conclusion}
This template demonstrates basic \LaTeX{} features including mathematical formulas, tables, and lists. The bibliography system is now set up to work seamlessly with Bibby AI's citation management. Feel free to modify it for your needs!

\bibliographystyle{plain}
\bibliography{cite}

\end{document}\documentclass{article}
\usepackage[utf8]{inputenc}
\usepackage{booktabs}
\usepackage{amsmath}
\usepackage{natbib}

\title{My Document}
\author{Your Name}
\date{\today}

\begin{document}

\maketitle

\section{Introduction}
This is a sample document demonstrating various \LaTeX{} features including mathematical equations, tables, and formatting. You can cite papers using the citation search feature, and they will be automatically added to your \texttt{cite.bib} file.

\section{Mathematical Expressions}
Here are some example equations:

\subsection{Quadratic Formula}
The quadratic formula is given by:
\[x = \frac{-b \pm \sqrt{b^2 - 4ac}}{2a}\]

\subsection{Matrix Example}
A sample matrix:
\[
\begin{pmatrix}
a_{11} & a_{12} & a_{13} \\
a_{21} & a_{22} & a_{23} \\
a_{31} & a_{32} & a_{33}
\end{pmatrix}
\]

\section{Tables}
Below is a sample table showing data:

\begin{table}[h]
\centering
\begin{tabular}{@{}lrrr@{}}
\toprule
\textbf{Category} & \textbf{Value 1} & \textbf{Value 2} & \textbf{Total} \\
\midrule
Item A & 10.5 & 20.3 & 30.8 \\
Item B & 15.2 & 25.1 & 40.3 \\
Item C & 22.1 & 12.7 & 34.8 \\
\bottomrule
\end{tabular}
\caption{Sample Data Table}
\label{tab:sample}
\end{table}

\section{Lists}
\subsection{Bullet Points}
\begin{itemize}
  \item First item
  \item Second item
  \item Third item with sub-items:
    \begin{itemize}
      \item Sub-item 1
      \item Sub-item 2
    \end{itemize}
\end{itemize}

\subsection{Numbered List}
\begin{enumerate}
  \item First step
  \item Second step
  \item Third step
\end{enumerate}

\section{Citations and References}
When you use the citation search feature, references will be automatically managed. For example, you might cite a paper like this: \cite{example2024}. All citations will be automatically added to your \texttt{cite.bib} file.

\section{Conclusion}
This template demonstrates basic \LaTeX{} features including mathematical formulas, tables, and lists. The bibliography system is now set up to work seamlessly with Bibby AI's citation management. Feel free to modify it for your needs!

\bibliographystyle{plain}
\bibliography{cite}

\end{document}\documentclass{article}
\usepackage[utf8]{inputenc}
\usepackage{booktabs}
\usepackage{amsmath}
\usepackage{natbib}

\title{My Document}
\author{Your Name}
\date{\today}

\begin{document}

\maketitle

\section{Introduction}
This is a sample document demonstrating various \LaTeX{} features including mathematical equations, tables, and formatting. You can cite papers using the citation search feature, and they will be automatically added to your \texttt{cite.bib} file.

\section{Mathematical Expressions}
Here are some example equations:

\subsection{Quadratic Formula}
The quadratic formula is given by:
\[x = \frac{-b \pm \sqrt{b^2 - 4ac}}{2a}\]

\subsection{Matrix Example}
A sample matrix:
\[
\begin{pmatrix}
a_{11} & a_{12} & a_{13} \\
a_{21} & a_{22} & a_{23} \\
a_{31} & a_{32} & a_{33}
\end{pmatrix}
\]

\section{Tables}
Below is a sample table showing data:

\begin{table}[h]
\centering
\begin{tabular}{@{}lrrr@{}}
\toprule
\textbf{Category} & \textbf{Value 1} & \textbf{Value 2} & \textbf{Total} \\
\midrule
Item A & 10.5 & 20.3 & 30.8 \\
Item B & 15.2 & 25.1 & 40.3 \\
Item C & 22.1 & 12.7 & 34.8 \\
\bottomrule
\end{tabular}
\caption{Sample Data Table}
\label{tab:sample}
\end{table}

\section{Lists}
\subsection{Bullet Points}
\begin{itemize}
  \item First item
  \item Second item
  \item Third item with sub-items:
    \begin{itemize}
      \item Sub-item 1
      \item Sub-item 2
    \end{itemize}
\end{itemize}

\subsection{Numbered List}
\begin{enumerate}
  \item First step
  \item Second step
  \item Third step
\end{enumerate}

\section{Citations and References}
When you use the citation search feature, references will be automatically managed. For example, you might cite a paper like this: \cite{example2024}. All citations will be automatically added to your \texttt{cite.bib} file.

\section{Conclusion}
This template demonstrates basic \LaTeX{} features including mathematical formulas, tables, and lists. The bibliography system is now set up to work seamlessly with Bibby AI's citation management. Feel free to modify it for your needs!

\bibliographystyle{plain}
\bibliography{cite}

\end{document}\documentclass{article}
\usepackage[utf8]{inputenc}
\usepackage{booktabs}
\usepackage{amsmath}
\usepackage{natbib}

\title{My Document}
\author{Your Name}
\date{\today}

\begin{document}

\maketitle

\section{Introduction}
This is a sample document demonstrating various \LaTeX{} features including mathematical equations, tables, and formatting. You can cite papers using the citation search feature, and they will be automatically added to your \texttt{cite.bib} file.

\section{Mathematical Expressions}
Here are some example equations:

\subsection{Quadratic Formula}
The quadratic formula is given by:
\[x = \frac{-b \pm \sqrt{b^2 - 4ac}}{2a}\]

\subsection{Matrix Example}
A sample matrix:
\[
\begin{pmatrix}
a_{11} & a_{12} & a_{13} \\
a_{21} & a_{22} & a_{23} \\
a_{31} & a_{32} & a_{33}
\end{pmatrix}
\]

\section{Tables}
Below is a sample table showing data:

\begin{table}[h]
\centering
\begin{tabular}{@{}lrrr@{}}
\toprule
\textbf{Category} & \textbf{Value 1} & \textbf{Value 2} & \textbf{Total} \\
\midrule
Item A & 10.5 & 20.3 & 30.8 \\
Item B & 15.2 & 25.1 & 40.3 \\
Item C & 22.1 & 12.7 & 34.8 \\
\bottomrule
\end{tabular}
\caption{Sample Data Table}
\label{tab:sample}
\end{table}

\section{Lists}
\subsection{Bullet Points}
\begin{itemize}
  \item First item
  \item Second item
  \item Third item with sub-items:
    \begin{itemize}
      \item Sub-item 1
      \item Sub-item 2
    \end{itemize}
\end{itemize}

\subsection{Numbered List}
\begin{enumerate}
  \item First step
  \item Second step
  \item Third step
\end{enumerate}

\section{Citations and References}
When you use the citation search feature, references will be automatically managed. For example, you might cite a paper like this: \cite{example2024}. All citations will be automatically added to your \texttt{cite.bib} file.

\section{Conclusion}
This template demonstrates basic \LaTeX{} features including mathematical formulas, tables, and lists. The bibliography system is now set up to work seamlessly with Bibby AI's citation management. Feel free to modify it for your needs!

\bibliographystyle{plain}
\bibliography{cite}

\end{document}\documentclass{article}
\usepackage[utf8]{inputenc}
\usepackage{booktabs}
\usepackage{amsmath}
\usepackage{natbib}

\title{My Document}
\author{Your Name}
\date{\today}

\begin{document}

\maketitle

\section{Introduction}
This is a sample document demonstrating various \LaTeX{} features including mathematical equations, tables, and formatting. You can cite papers using the citation search feature, and they will be automatically added to your \texttt{cite.bib} file.

\section{Mathematical Expressions}
Here are some example equations:

\subsection{Quadratic Formula}
The quadratic formula is given by:
\[x = \frac{-b \pm \sqrt{b^2 - 4ac}}{2a}\]

\subsection{Matrix Example}
A sample matrix:
\[
\begin{pmatrix}
a_{11} & a_{12} & a_{13} \\
a_{21} & a_{22} & a_{23} \\
a_{31} & a_{32} & a_{33}
\end{pmatrix}
\]

\section{Tables}
Below is a sample table showing data:

\begin{table}[h]
\centering
\begin{tabular}{@{}lrrr@{}}
\toprule
\textbf{Category} & \textbf{Value 1} & \textbf{Value 2} & \textbf{Total} \\
\midrule
Item A & 10.5 & 20.3 & 30.8 \\
Item B & 15.2 & 25.1 & 40.3 \\
Item C & 22.1 & 12.7 & 34.8 \\
\bottomrule
\end{tabular}
\caption{Sample Data Table}
\label{tab:sample}
\end{table}

\section{Lists}
\subsection{Bullet Points}
\begin{itemize}
  \item First item
  \item Second item
  \item Third item with sub-items:
    \begin{itemize}
      \item Sub-item 1
      \item Sub-item 2
    \end{itemize}
\end{itemize}

\subsection{Numbered List}
\begin{enumerate}
  \item First step
  \item Second step
  \item Third step
\end{enumerate}

\section{Citations and References}
When you use the citation search feature, references will be automatically managed. For example, you might cite a paper like this: \cite{example2024}. All citations will be automatically added to your \texttt{cite.bib} file.

\section{Conclusion}
This template demonstrates basic \LaTeX{} features including mathematical formulas, tables, and lists. The bibliography system is now set up to work seamlessly with Bibby AI's citation management. Feel free to modify it for your needs!

\bibliographystyle{plain}
\bibliography{cite}

\end{document}\documentclass{article}
\usepackage[utf8]{inputenc}
\usepackage{booktabs}
\usepackage{amsmath}
\usepackage{natbib}

\title{My Document}
\author{Your Name}
\date{\today}

\begin{document}

\maketitle

\section{Introduction}
This is a sample document demonstrating various \LaTeX{} features including mathematical equations, tables, and formatting. You can cite papers using the citation search feature, and they will be automatically added to your \texttt{cite.bib} file.

\section{Mathematical Expressions}
Here are some example equations:

\subsection{Quadratic Formula}
The quadratic formula is given by:
\[x = \frac{-b \pm \sqrt{b^2 - 4ac}}{2a}\]

\subsection{Matrix Example}
A sample matrix:
\[
\begin{pmatrix}
a_{11} & a_{12} & a_{13} \\
a_{21} & a_{22} & a_{23} \\
a_{31} & a_{32} & a_{33}
\end{pmatrix}
\]

\section{Tables}
Below is a sample table showing data:

\begin{table}[h]
\centering
\begin{tabular}{@{}lrrr@{}}
\toprule
\textbf{Category} & \textbf{Value 1} & \textbf{Value 2} & \textbf{Total} \\
\midrule
Item A & 10.5 & 20.3 & 30.8 \\
Item B & 15.2 & 25.1 & 40.3 \\
Item C & 22.1 & 12.7 & 34.8 \\
\bottomrule
\end{tabular}
\caption{Sample Data Table}
\label{tab:sample}
\end{table}

\section{Lists}
\subsection{Bullet Points}
\begin{itemize}
  \item First item
  \item Second item
  \item Third item with sub-items:
    \begin{itemize}
      \item Sub-item 1
      \item Sub-item 2
    \end{itemize}
\end{itemize}

\subsection{Numbered List}
\begin{enumerate}
  \item First step
  \item Second step
  \item Third step
\end{enumerate}

\section{Citations and References}
When you use the citation search feature, references will be automatically managed. For example, you might cite a paper like this: \cite{example2024}. All citations will be automatically added to your \texttt{cite.bib} file.

\section{Conclusion}
This template demonstrates basic \LaTeX{} features including mathematical formulas, tables, and lists. The bibliography system is now set up to work seamlessly with Bibby AI's citation management. Feel free to modify it for your needs!

\bibliographystyle{plain}
\bibliography{cite}

\end{document}\documentclass{article}
\usepackage[utf8]{inputenc}
\usepackage{booktabs}
\usepackage{amsmath}
\usepackage{natbib}

\title{My Document}
\author{Your Name}
\date{\today}

\begin{document}

\maketitle

\section{Introduction}
This is a sample document demonstrating various \LaTeX{} features including mathematical equations, tables, and formatting. You can cite papers using the citation search feature, and they will be automatically added to your \texttt{cite.bib} file.

\section{Mathematical Expressions}
Here are some example equations:

\subsection{Quadratic Formula}
The quadratic formula is given by:
\[x = \frac{-b \pm \sqrt{b^2 - 4ac}}{2a}\]

\subsection{Matrix Example}
A sample matrix:
\[
\begin{pmatrix}
a_{11} & a_{12} & a_{13} \\
a_{21} & a_{22} & a_{23} \\
a_{31} & a_{32} & a_{33}
\end{pmatrix}
\]

\section{Tables}
Below is a sample table showing data:

\begin{table}[h]
\centering
\begin{tabular}{@{}lrrr@{}}
\toprule
\textbf{Category} & \textbf{Value 1} & \textbf{Value 2} & \textbf{Total} \\
\midrule
Item A & 10.5 & 20.3 & 30.8 \\
Item B & 15.2 & 25.1 & 40.3 \\
Item C & 22.1 & 12.7 & 34.8 \\
\bottomrule
\end{tabular}
\caption{Sample Data Table}
\label{tab:sample}
\end{table}

\section{Lists}
\subsection{Bullet Points}
\begin{itemize}
  \item First item
  \item Second item
  \item Third item with sub-items:
    \begin{itemize}
      \item Sub-item 1
      \item Sub-item 2
    \end{itemize}
\end{itemize}

\subsection{Numbered List}
\begin{enumerate}
  \item First step
  \item Second step
  \item Third step
\end{enumerate}

\section{Citations and References}
When you use the citation search feature, references will be automatically managed. For example, you might cite a paper like this: \cite{example2024}. All citations will be automatically added to your \texttt{cite.bib} file.

\section{Conclusion}
This template demonstrates basic \LaTeX{} features including mathematical formulas, tables, and lists. The bibliography system is now set up to work seamlessly with Bibby AI's citation management. Feel free to modify it for your needs!

\bibliographystyle{plain}
\bibliography{cite}

\end{document}\documentclass{article}
\usepackage[utf8]{inputenc}
\usepackage{booktabs}
\usepackage{amsmath}
\usepackage{natbib}

\title{My Document}
\author{Your Name}
\date{\today}

\begin{document}

\maketitle

\section{Introduction}
This is a sample document demonstrating various \LaTeX{} features including mathematical equations, tables, and formatting. You can cite papers using the citation search feature, and they will be automatically added to your \texttt{cite.bib} file.

\section{Mathematical Expressions}
Here are some example equations:

\subsection{Quadratic Formula}
The quadratic formula is given by:
\[x = \frac{-b \pm \sqrt{b^2 - 4ac}}{2a}\]

\subsection{Matrix Example}
A sample matrix:
\[
\begin{pmatrix}
a_{11} & a_{12} & a_{13} \\
a_{21} & a_{22} & a_{23} \\
a_{31} & a_{32} & a_{33}
\end{pmatrix}
\]

\section{Tables}
Below is a sample table showing data:

\begin{table}[h]
\centering
\begin{tabular}{@{}lrrr@{}}
\toprule
\textbf{Category} & \textbf{Value 1} & \textbf{Value 2} & \textbf{Total} \\
\midrule
Item A & 10.5 & 20.3 & 30.8 \\
Item B & 15.2 & 25.1 & 40.3 \\
Item C & 22.1 & 12.7 & 34.8 \\
\bottomrule
\end{tabular}
\caption{Sample Data Table}
\label{tab:sample}
\end{table}

\section{Lists}
\subsection{Bullet Points}
\begin{itemize}
  \item First item
  \item Second item
  \item Third item with sub-items:
    \begin{itemize}
      \item Sub-item 1
      \item Sub-item 2
    \end{itemize}
\end{itemize}

\subsection{Numbered List}
\begin{enumerate}
  \item First step
  \item Second step
  \item Third step
\end{enumerate}

\section{Citations and References}
When you use the citation search feature, references will be automatically managed. For example, you might cite a paper like this: \cite{example2024}. All citations will be automatically added to your \texttt{cite.bib} file.

\section{Conclusion}
This template demonstrates basic \LaTeX{} features including mathematical formulas, tables, and lists. The bibliography system is now set up to work seamlessly with Bibby AI's citation management. Feel free to modify it for your needs!

\bibliographystyle{plain}
\bibliography{cite}

\end{document}\documentclass{article}
\usepackage[utf8]{inputenc}
\usepackage{booktabs}
\usepackage{amsmath}
\usepackage{natbib}

\title{My Document}
\author{Your Name}
\date{\today}

\begin{document}

\maketitle

\section{Introduction}
This is a sample document demonstrating various \LaTeX{} features including mathematical equations, tables, and formatting. You can cite papers using the citation search feature, and they will be automatically added to your \texttt{cite.bib} file.

\section{Mathematical Expressions}
Here are some example equations:

\subsection{Quadratic Formula}
The quadratic formula is given by:
\[x = \frac{-b \pm \sqrt{b^2 - 4ac}}{2a}\]

\subsection{Matrix Example}
A sample matrix:
\[
\begin{pmatrix}
a_{11} & a_{12} & a_{13} \\
a_{21} & a_{22} & a_{23} \\
a_{31} & a_{32} & a_{33}
\end{pmatrix}
\]

\section{Tables}
Below is a sample table showing data:

\begin{table}[h]
\centering
\begin{tabular}{@{}lrrr@{}}
\toprule
\textbf{Category} & \textbf{Value 1} & \textbf{Value 2} & \textbf{Total} \\
\midrule
Item A & 10.5 & 20.3 & 30.8 \\
Item B & 15.2 & 25.1 & 40.3 \\
Item C & 22.1 & 12.7 & 34.8 \\
\bottomrule
\end{tabular}
\caption{Sample Data Table}
\label{tab:sample}
\end{table}

\section{Lists}
\subsection{Bullet Points}
\begin{itemize}
  \item First item
  \item Second item
  \item Third item with sub-items:
    \begin{itemize}
      \item Sub-item 1
      \item Sub-item 2
    \end{itemize}
\end{itemize}

\subsection{Numbered List}
\begin{enumerate}
  \item First step
  \item Second step
  \item Third step
\end{enumerate}

\section{Citations and References}
When you use the citation search feature, references will be automatically managed. For example, you might cite a paper like this: \cite{example2024}. All citations will be automatically added to your \texttt{cite.bib} file.

\section{Conclusion}
This template demonstrates basic \LaTeX{} features including mathematical formulas, tables, and lists. The bibliography system is now set up to work seamlessly with Bibby AI's citation management. Feel free to modify it for your needs!

\bibliographystyle{plain}
\bibliography{cite}

\end{document}\documentclass{article}
\usepackage[utf8]{inputenc}
\usepackage{booktabs}
\usepackage{amsmath}
\usepackage{natbib}

\title{My Document}
\author{Your Name}
\date{\today}

\begin{document}

\maketitle

\section{Introduction}
This is a sample document demonstrating various \LaTeX{} features including mathematical equations, tables, and formatting. You can cite papers using the citation search feature, and they will be automatically added to your \texttt{cite.bib} file.

\section{Mathematical Expressions}
Here are some example equations:

\subsection{Quadratic Formula}
The quadratic formula is given by:
\[x = \frac{-b \pm \sqrt{b^2 - 4ac}}{2a}\]

\subsection{Matrix Example}
A sample matrix:
\[
\begin{pmatrix}
a_{11} & a_{12} & a_{13} \\
a_{21} & a_{22} & a_{23} \\
a_{31} & a_{32} & a_{33}
\end{pmatrix}
\]

\section{Tables}
Below is a sample table showing data:

\begin{table}[h]
\centering
\begin{tabular}{@{}lrrr@{}}
\toprule
\textbf{Category} & \textbf{Value 1} & \textbf{Value 2} & \textbf{Total} \\
\midrule
Item A & 10.5 & 20.3 & 30.8 \\
Item B & 15.2 & 25.1 & 40.3 \\
Item C & 22.1 & 12.7 & 34.8 \\
\bottomrule
\end{tabular}
\caption{Sample Data Table}
\label{tab:sample}
\end{table}

\section{Lists}
\subsection{Bullet Points}
\begin{itemize}
  \item First item
  \item Second item
  \item Third item with sub-items:
    \begin{itemize}
      \item Sub-item 1
      \item Sub-item 2
    \end{itemize}
\end{itemize}

\subsection{Numbered List}
\begin{enumerate}
  \item First step
  \item Second step
  \item Third step
\end{enumerate}

\section{Citations and References}
When you use the citation search feature, references will be automatically managed. For example, you might cite a paper like this: \cite{example2024}. All citations will be automatically added to your \texttt{cite.bib} file.

\section{Conclusion}
This template demonstrates basic \LaTeX{} features including mathematical formulas, tables, and lists. The bibliography system is now set up to work seamlessly with Bibby AI's citation management. Feel free to modify it for your needs!

\bibliographystyle{plain}
\bibliography{cite}

\end{document}\documentclass{article}
\usepackage[utf8]{inputenc}
\usepackage{booktabs}
\usepackage{amsmath}
\usepackage{natbib}

\title{My Document}
\author{Your Name}
\date{\today}

\begin{document}

\maketitle

\section{Introduction}
This is a sample document demonstrating various \LaTeX{} features including mathematical equations, tables, and formatting. You can cite papers using the citation search feature, and they will be automatically added to your \texttt{cite.bib} file.

\section{Mathematical Expressions}
Here are some example equations:

\subsection{Quadratic Formula}
The quadratic formula is given by:
\[x = \frac{-b \pm \sqrt{b^2 - 4ac}}{2a}\]

\subsection{Matrix Example}
A sample matrix:
\[
\begin{pmatrix}
a_{11} & a_{12} & a_{13} \\
a_{21} & a_{22} & a_{23} \\
a_{31} & a_{32} & a_{33}
\end{pmatrix}
\]

\section{Tables}
Below is a sample table showing data:

\begin{table}[h]
\centering
\begin{tabular}{@{}lrrr@{}}
\toprule
\textbf{Category} & \textbf{Value 1} & \textbf{Value 2} & \textbf{Total} \\
\midrule
Item A & 10.5 & 20.3 & 30.8 \\
Item B & 15.2 & 25.1 & 40.3 \\
Item C & 22.1 & 12.7 & 34.8 \\
\bottomrule
\end{tabular}
\caption{Sample Data Table}
\label{tab:sample}
\end{table}

\section{Lists}
\subsection{Bullet Points}
\begin{itemize}
  \item First item
  \item Second item
  \item Third item with sub-items:
    \begin{itemize}
      \item Sub-item 1
      \item Sub-item 2
    \end{itemize}
\end{itemize}

\subsection{Numbered List}
\begin{enumerate}
  \item First step
  \item Second step
  \item Third step
\end{enumerate}

\section{Citations and References}
When you use the citation search feature, references will be automatically managed. For example, you might cite a paper like this: \cite{example2024}. All citations will be automatically added to your \texttt{cite.bib} file.

\section{Conclusion}
This template demonstrates basic \LaTeX{} features including mathematical formulas, tables, and lists. The bibliography system is now set up to work seamlessly with Bibby AI's citation management. Feel free to modify it for your needs!

\bibliographystyle{plain}
\bibliography{cite}

\end{document}\documentclass{article}
\usepackage[utf8]{inputenc}
\usepackage{booktabs}
\usepackage{amsmath}
\usepackage{natbib}

\title{My Document}
\author{Your Name}
\date{\today}

\begin{document}

\maketitle

\section{Introduction}
This is a sample document demonstrating various \LaTeX{} features including mathematical equations, tables, and formatting. You can cite papers using the citation search feature, and they will be automatically added to your \texttt{cite.bib} file.

\section{Mathematical Expressions}
Here are some example equations:

\subsection{Quadratic Formula}
The quadratic formula is given by:
\[x = \frac{-b \pm \sqrt{b^2 - 4ac}}{2a}\]

\subsection{Matrix Example}
A sample matrix:
\[
\begin{pmatrix}
a_{11} & a_{12} & a_{13} \\
a_{21} & a_{22} & a_{23} \\
a_{31} & a_{32} & a_{33}
\end{pmatrix}
\]

\section{Tables}
Below is a sample table showing data:

\begin{table}[h]
\centering
\begin{tabular}{@{}lrrr@{}}
\toprule
\textbf{Category} & \textbf{Value 1} & \textbf{Value 2} & \textbf{Total} \\
\midrule
Item A & 10.5 & 20.3 & 30.8 \\
Item B & 15.2 & 25.1 & 40.3 \\
Item C & 22.1 & 12.7 & 34.8 \\
\bottomrule
\end{tabular}
\caption{Sample Data Table}
\label{tab:sample}
\end{table}

\section{Lists}
\subsection{Bullet Points}
\begin{itemize}
  \item First item
  \item Second item
  \item Third item with sub-items:
    \begin{itemize}
      \item Sub-item 1
      \item Sub-item 2
    \end{itemize}
\end{itemize}

\subsection{Numbered List}
\begin{enumerate}
  \item First step
  \item Second step
  \item Third step
\end{enumerate}

\section{Citations and References}
When you use the citation search feature, references will be automatically managed. For example, you might cite a paper like this: \cite{example2024}. All citations will be automatically added to your \texttt{cite.bib} file.

\section{Conclusion}
This template demonstrates basic \LaTeX{} features including mathematical formulas, tables, and lists. The bibliography system is now set up to work seamlessly with Bibby AI's citation management. Feel free to modify it for your needs!

\bibliographystyle{plain}
\bibliography{cite}

\end{document}\documentclass{article}
\usepackage[utf8]{inputenc}
\usepackage{booktabs}
\usepackage{amsmath}
\usepackage{natbib}

\title{My Document}
\author{Your Name}
\date{\today}

\begin{document}

\maketitle

\section{Introduction}
This is a sample document demonstrating various \LaTeX{} features including mathematical equations, tables, and formatting. You can cite papers using the citation search feature, and they will be automatically added to your \texttt{cite.bib} file.

\section{Mathematical Expressions}
Here are some example equations:

\subsection{Quadratic Formula}
The quadratic formula is given by:
\[x = \frac{-b \pm \sqrt{b^2 - 4ac}}{2a}\]

\subsection{Matrix Example}
A sample matrix:
\[
\begin{pmatrix}
a_{11} & a_{12} & a_{13} \\
a_{21} & a_{22} & a_{23} \\
a_{31} & a_{32} & a_{33}
\end{pmatrix}
\]

\section{Tables}
Below is a sample table showing data:

\begin{table}[h]
\centering
\begin{tabular}{@{}lrrr@{}}
\toprule
\textbf{Category} & \textbf{Value 1} & \textbf{Value 2} & \textbf{Total} \\
\midrule
Item A & 10.5 & 20.3 & 30.8 \\
Item B & 15.2 & 25.1 & 40.3 \\
Item C & 22.1 & 12.7 & 34.8 \\
\bottomrule
\end{tabular}
\caption{Sample Data Table}
\label{tab:sample}
\end{table}

\section{Lists}
\subsection{Bullet Points}
\begin{itemize}
  \item First item
  \item Second item
  \item Third item with sub-items:
    \begin{itemize}
      \item Sub-item 1
      \item Sub-item 2
    \end{itemize}
\end{itemize}

\subsection{Numbered List}
\begin{enumerate}
  \item First step
  \item Second step
  \item Third step
\end{enumerate}

\section{Citations and References}
When you use the citation search feature, references will be automatically managed. For example, you might cite a paper like this: \cite{example2024}. All citations will be automatically added to your \texttt{cite.bib} file.

\section{Conclusion}
This template demonstrates basic \LaTeX{} features including mathematical formulas, tables, and lists. The bibliography system is now set up to work seamlessly with Bibby AI's citation management. Feel free to modify it for your needs!

\bibliographystyle{plain}
\bibliography{cite}

\end{document}\documentclass{article}
\usepackage[utf8]{inputenc}
\usepackage{booktabs}
\usepackage{amsmath}
\usepackage{natbib}

\title{My Document}
\author{Your Name}
\date{\today}

\begin{document}

\maketitle

\section{Introduction}
This is a sample document demonstrating various \LaTeX{} features including mathematical equations, tables, and formatting. You can cite papers using the citation search feature, and they will be automatically added to your \texttt{cite.bib} file.

\section{Mathematical Expressions}
Here are some example equations:

\subsection{Quadratic Formula}
The quadratic formula is given by:
\[x = \frac{-b \pm \sqrt{b^2 - 4ac}}{2a}\]

\subsection{Matrix Example}
A sample matrix:
\[
\begin{pmatrix}
a_{11} & a_{12} & a_{13} \\
a_{21} & a_{22} & a_{23} \\
a_{31} & a_{32} & a_{33}
\end{pmatrix}
\]

\section{Tables}
Below is a sample table showing data:

\begin{table}[h]
\centering
\begin{tabular}{@{}lrrr@{}}
\toprule
\textbf{Category} & \textbf{Value 1} & \textbf{Value 2} & \textbf{Total} \\
\midrule
Item A & 10.5 & 20.3 & 30.8 \\
Item B & 15.2 & 25.1 & 40.3 \\
Item C & 22.1 & 12.7 & 34.8 \\
\bottomrule
\end{tabular}
\caption{Sample Data Table}
\label{tab:sample}
\end{table}

\section{Lists}
\subsection{Bullet Points}
\begin{itemize}
  \item First item
  \item Second item
  \item Third item with sub-items:
    \begin{itemize}
      \item Sub-item 1
      \item Sub-item 2
    \end{itemize}
\end{itemize}

\subsection{Numbered List}
\begin{enumerate}
  \item First step
  \item Second step
  \item Third step
\end{enumerate}

\section{Citations and References}
When you use the citation search feature, references will be automatically managed. For example, you might cite a paper like this: \cite{example2024}. All citations will be automatically added to your \texttt{cite.bib} file.

\section{Conclusion}
This template demonstrates basic \LaTeX{} features including mathematical formulas, tables, and lists. The bibliography system is now set up to work seamlessly with Bibby AI's citation management. Feel free to modify it for your needs!

\bibliographystyle{plain}
\bibliography{cite}

\end{document}\documentclass{article}
\usepackage[utf8]{inputenc}
\usepackage{booktabs}
\usepackage{amsmath}
\usepackage{natbib}

\title{My Document}
\author{Your Name}
\date{\today}

\begin{document}

\maketitle

\section{Introduction}
This is a sample document demonstrating various \LaTeX{} features including mathematical equations, tables, and formatting. You can cite papers using the citation search feature, and they will be automatically added to your \texttt{cite.bib} file.

\section{Mathematical Expressions}
Here are some example equations:

\subsection{Quadratic Formula}
The quadratic formula is given by:
\[x = \frac{-b \pm \sqrt{b^2 - 4ac}}{2a}\]

\subsection{Matrix Example}
A sample matrix:
\[
\begin{pmatrix}
a_{11} & a_{12} & a_{13} \\
a_{21} & a_{22} & a_{23} \\
a_{31} & a_{32} & a_{33}
\end{pmatrix}
\]

\section{Tables}
Below is a sample table showing data:

\begin{table}[h]
\centering
\begin{tabular}{@{}lrrr@{}}
\toprule
\textbf{Category} & \textbf{Value 1} & \textbf{Value 2} & \textbf{Total} \\
\midrule
Item A & 10.5 & 20.3 & 30.8 \\
Item B & 15.2 & 25.1 & 40.3 \\
Item C & 22.1 & 12.7 & 34.8 \\
\bottomrule
\end{tabular}
\caption{Sample Data Table}
\label{tab:sample}
\end{table}

\section{Lists}
\subsection{Bullet Points}
\begin{itemize}
  \item First item
  \item Second item
  \item Third item with sub-items:
    \begin{itemize}
      \item Sub-item 1
      \item Sub-item 2
    \end{itemize}
\end{itemize}

\subsection{Numbered List}
\begin{enumerate}
  \item First step
  \item Second step
  \item Third step
\end{enumerate}

\section{Citations and References}
When you use the citation search feature, references will be automatically managed. For example, you might cite a paper like this: \cite{example2024}. All citations will be automatically added to your \texttt{cite.bib} file.

\section{Conclusion}
This template demonstrates basic \LaTeX{} features including mathematical formulas, tables, and lists. The bibliography system is now set up to work seamlessly with Bibby AI's citation management. Feel free to modify it for your needs!

\bibliographystyle{plain}
\bibliography{cite}

\end{document}\documentclass{article}
\usepackage[utf8]{inputenc}
\usepackage{booktabs}
\usepackage{amsmath}
\usepackage{natbib}

\title{My Document}
\author{Your Name}
\date{\today}

\begin{document}

\maketitle

\section{Introduction}
This is a sample document demonstrating various \LaTeX{} features including mathematical equations, tables, and formatting. You can cite papers using the citation search feature, and they will be automatically added to your \texttt{cite.bib} file.

\section{Mathematical Expressions}
Here are some example equations:

\subsection{Quadratic Formula}
The quadratic formula is given by:
\[x = \frac{-b \pm \sqrt{b^2 - 4ac}}{2a}\]

\subsection{Matrix Example}
A sample matrix:
\[
\begin{pmatrix}
a_{11} & a_{12} & a_{13} \\
a_{21} & a_{22} & a_{23} \\
a_{31} & a_{32} & a_{33}
\end{pmatrix}
\]

\section{Tables}
Below is a sample table showing data:

\begin{table}[h]
\centering
\begin{tabular}{@{}lrrr@{}}
\toprule
\textbf{Category} & \textbf{Value 1} & \textbf{Value 2} & \textbf{Total} \\
\midrule
Item A & 10.5 & 20.3 & 30.8 \\
Item B & 15.2 & 25.1 & 40.3 \\
Item C & 22.1 & 12.7 & 34.8 \\
\bottomrule
\end{tabular}
\caption{Sample Data Table}
\label{tab:sample}
\end{table}

\section{Lists}
\subsection{Bullet Points}
\begin{itemize}
  \item First item
  \item Second item
  \item Third item with sub-items:
    \begin{itemize}
      \item Sub-item 1
      \item Sub-item 2
    \end{itemize}
\end{itemize}

\subsection{Numbered List}
\begin{enumerate}
  \item First step
  \item Second step
  \item Third step
\end{enumerate}

\section{Citations and References}
When you use the citation search feature, references will be automatically managed. For example, you might cite a paper like this: \cite{example2024}. All citations will be automatically added to your \texttt{cite.bib} file.

\section{Conclusion}
This template demonstrates basic \LaTeX{} features including mathematical formulas, tables, and lists. The bibliography system is now set up to work seamlessly with Bibby AI's citation management. Feel free to modify it for your needs!

\bibliographystyle{plain}
\bibliography{cite}

\end{document}\documentclass{article}
\usepackage[utf8]{inputenc}
\usepackage{booktabs}
\usepackage{amsmath}
\usepackage{natbib}

\title{My Document}
\author{Your Name}
\date{\today}

\begin{document}

\maketitle

\section{Introduction}
This is a sample document demonstrating various \LaTeX{} features including mathematical equations, tables, and formatting. You can cite papers using the citation search feature, and they will be automatically added to your \texttt{cite.bib} file.

\section{Mathematical Expressions}
Here are some example equations:

\subsection{Quadratic Formula}
The quadratic formula is given by:
\[x = \frac{-b \pm \sqrt{b^2 - 4ac}}{2a}\]

\subsection{Matrix Example}
A sample matrix:
\[
\begin{pmatrix}
a_{11} & a_{12} & a_{13} \\
a_{21} & a_{22} & a_{23} \\
a_{31} & a_{32} & a_{33}
\end{pmatrix}
\]

\section{Tables}
Below is a sample table showing data:

\begin{table}[h]
\centering
\begin{tabular}{@{}lrrr@{}}
\toprule
\textbf{Category} & \textbf{Value 1} & \textbf{Value 2} & \textbf{Total} \\
\midrule
Item A & 10.5 & 20.3 & 30.8 \\
Item B & 15.2 & 25.1 & 40.3 \\
Item C & 22.1 & 12.7 & 34.8 \\
\bottomrule
\end{tabular}
\caption{Sample Data Table}
\label{tab:sample}
\end{table}

\section{Lists}
\subsection{Bullet Points}
\begin{itemize}
  \item First item
  \item Second item
  \item Third item with sub-items:
    \begin{itemize}
      \item Sub-item 1
      \item Sub-item 2
    \end{itemize}
\end{itemize}

\subsection{Numbered List}
\begin{enumerate}
  \item First step
  \item Second step
  \item Third step
\end{enumerate}

\section{Citations and References}
When you use the citation search feature, references will be automatically managed. For example, you might cite a paper like this: \cite{example2024}. All citations will be automatically added to your \texttt{cite.bib} file.

\section{Conclusion}
This template demonstrates basic \LaTeX{} features including mathematical formulas, tables, and lists. The bibliography system is now set up to work seamlessly with Bibby AI's citation management. Feel free to modify it for your needs!

\bibliographystyle{plain}
\bibliography{cite}

\end{document}\documentclass{article}
\usepackage[utf8]{inputenc}
\usepackage{booktabs}
\usepackage{amsmath}
\usepackage{natbib}

\title{My Document}
\author{Your Name}
\date{\today}

\begin{document}

\maketitle

\section{Introduction}
This is a sample document demonstrating various \LaTeX{} features including mathematical equations, tables, and formatting. You can cite papers using the citation search feature, and they will be automatically added to your \texttt{cite.bib} file.

\section{Mathematical Expressions}
Here are some example equations:

\subsection{Quadratic Formula}
The quadratic formula is given by:
\[x = \frac{-b \pm \sqrt{b^2 - 4ac}}{2a}\]

\subsection{Matrix Example}
A sample matrix:
\[
\begin{pmatrix}
a_{11} & a_{12} & a_{13} \\
a_{21} & a_{22} & a_{23} \\
a_{31} & a_{32} & a_{33}
\end{pmatrix}
\]

\section{Tables}
Below is a sample table showing data:

\begin{table}[h]
\centering
\begin{tabular}{@{}lrrr@{}}
\toprule
\textbf{Category} & \textbf{Value 1} & \textbf{Value 2} & \textbf{Total} \\
\midrule
Item A & 10.5 & 20.3 & 30.8 \\
Item B & 15.2 & 25.1 & 40.3 \\
Item C & 22.1 & 12.7 & 34.8 \\
\bottomrule
\end{tabular}
\caption{Sample Data Table}
\label{tab:sample}
\end{table}

\section{Lists}
\subsection{Bullet Points}
\begin{itemize}
  \item First item
  \item Second item
  \item Third item with sub-items:
    \begin{itemize}
      \item Sub-item 1
      \item Sub-item 2
    \end{itemize}
\end{itemize}

\subsection{Numbered List}
\begin{enumerate}
  \item First step
  \item Second step
  \item Third step
\end{enumerate}

\section{Citations and References}
When you use the citation search feature, references will be automatically managed. For example, you might cite a paper like this: \cite{example2024}. All citations will be automatically added to your \texttt{cite.bib} file.

\section{Conclusion}
This template demonstrates basic \LaTeX{} features including mathematical formulas, tables, and lists. The bibliography system is now set up to work seamlessly with Bibby AI's citation management. Feel free to modify it for your needs!

\bibliographystyle{plain}
\bibliography{cite}

\end{document}\documentclass{article}
\usepackage[utf8]{inputenc}
\usepackage{booktabs}
\usepackage{amsmath}
\usepackage{natbib}

\title{My Document}
\author{Your Name}
\date{\today}

\begin{document}

\maketitle

\section{Introduction}
This is a sample document demonstrating various \LaTeX{} features including mathematical equations, tables, and formatting. You can cite papers using the citation search feature, and they will be automatically added to your \texttt{cite.bib} file.

\section{Mathematical Expressions}
Here are some example equations:

\subsection{Quadratic Formula}
The quadratic formula is given by:
\[x = \frac{-b \pm \sqrt{b^2 - 4ac}}{2a}\]

\subsection{Matrix Example}
A sample matrix:
\[
\begin{pmatrix}
a_{11} & a_{12} & a_{13} \\
a_{21} & a_{22} & a_{23} \\
a_{31} & a_{32} & a_{33}
\end{pmatrix}
\]

\section{Tables}
Below is a sample table showing data:

\begin{table}[h]
\centering
\begin{tabular}{@{}lrrr@{}}
\toprule
\textbf{Category} & \textbf{Value 1} & \textbf{Value 2} & \textbf{Total} \\
\midrule
Item A & 10.5 & 20.3 & 30.8 \\
Item B & 15.2 & 25.1 & 40.3 \\
Item C & 22.1 & 12.7 & 34.8 \\
\bottomrule
\end{tabular}
\caption{Sample Data Table}
\label{tab:sample}
\end{table}

\section{Lists}
\subsection{Bullet Points}
\begin{itemize}
  \item First item
  \item Second item
  \item Third item with sub-items:
    \begin{itemize}
      \item Sub-item 1
      \item Sub-item 2
    \end{itemize}
\end{itemize}

\subsection{Numbered List}
\begin{enumerate}
  \item First step
  \item Second step
  \item Third step
\end{enumerate}

\section{Citations and References}
When you use the citation search feature, references will be automatically managed. For example, you might cite a paper like this: \cite{example2024}. All citations will be automatically added to your \texttt{cite.bib} file.

\section{Conclusion}
This template demonstrates basic \LaTeX{} features including mathematical formulas, tables, and lists. The bibliography system is now set up to work seamlessly with Bibby AI's citation management. Feel free to modify it for your needs!

\bibliographystyle{plain}
\bibliography{cite}

\end{document}\documentclass{article}
\usepackage[utf8]{inputenc}
\usepackage{booktabs}
\usepackage{amsmath}
\usepackage{natbib}

\title{My Document}
\author{Your Name}
\date{\today}

\begin{document}

\maketitle

\section{Introduction}
This is a sample document demonstrating various \LaTeX{} features including mathematical equations, tables, and formatting. You can cite papers using the citation search feature, and they will be automatically added to your \texttt{cite.bib} file.

\section{Mathematical Expressions}
Here are some example equations:

\subsection{Quadratic Formula}
The quadratic formula is given by:
\[x = \frac{-b \pm \sqrt{b^2 - 4ac}}{2a}\]

\subsection{Matrix Example}
A sample matrix:
\[
\begin{pmatrix}
a_{11} & a_{12} & a_{13} \\
a_{21} & a_{22} & a_{23} \\
a_{31} & a_{32} & a_{33}
\end{pmatrix}
\]

\section{Tables}
Below is a sample table showing data:

\begin{table}[h]
\centering
\begin{tabular}{@{}lrrr@{}}
\toprule
\textbf{Category} & \textbf{Value 1} & \textbf{Value 2} & \textbf{Total} \\
\midrule
Item A & 10.5 & 20.3 & 30.8 \\
Item B & 15.2 & 25.1 & 40.3 \\
Item C & 22.1 & 12.7 & 34.8 \\
\bottomrule
\end{tabular}
\caption{Sample Data Table}
\label{tab:sample}
\end{table}

\section{Lists}
\subsection{Bullet Points}
\begin{itemize}
  \item First item
  \item Second item
  \item Third item with sub-items:
    \begin{itemize}
      \item Sub-item 1
      \item Sub-item 2
    \end{itemize}
\end{itemize}

\subsection{Numbered List}
\begin{enumerate}
  \item First step
  \item Second step
  \item Third step
\end{enumerate}

\section{Citations and References}
When you use the citation search feature, references will be automatically managed. For example, you might cite a paper like this: \cite{example2024}. All citations will be automatically added to your \texttt{cite.bib} file.

\section{Conclusion}
This template demonstrates basic \LaTeX{} features including mathematical formulas, tables, and lists. The bibliography system is now set up to work seamlessly with Bibby AI's citation management. Feel free to modify it for your needs!

\bibliographystyle{plain}
\bibliography{cite}

\end{document}\documentclass{article}
\usepackage[utf8]{inputenc}
\usepackage{booktabs}
\usepackage{amsmath}
\usepackage{natbib}

\title{My Document}
\author{Your Name}
\date{\today}

\begin{document}

\maketitle

\section{Introduction}
This is a sample document demonstrating various \LaTeX{} features including mathematical equations, tables, and formatting. You can cite papers using the citation search feature, and they will be automatically added to your \texttt{cite.bib} file.

\section{Mathematical Expressions}
Here are some example equations:

\subsection{Quadratic Formula}
The quadratic formula is given by:
\[x = \frac{-b \pm \sqrt{b^2 - 4ac}}{2a}\]

\subsection{Matrix Example}
A sample matrix:
\[
\begin{pmatrix}
a_{11} & a_{12} & a_{13} \\
a_{21} & a_{22} & a_{23} \\
a_{31} & a_{32} & a_{33}
\end{pmatrix}
\]

\section{Tables}
Below is a sample table showing data:

\begin{table}[h]
\centering
\begin{tabular}{@{}lrrr@{}}
\toprule
\textbf{Category} & \textbf{Value 1} & \textbf{Value 2} & \textbf{Total} \\
\midrule
Item A & 10.5 & 20.3 & 30.8 \\
Item B & 15.2 & 25.1 & 40.3 \\
Item C & 22.1 & 12.7 & 34.8 \\
\bottomrule
\end{tabular}
\caption{Sample Data Table}
\label{tab:sample}
\end{table}

\section{Lists}
\subsection{Bullet Points}
\begin{itemize}
  \item First item
  \item Second item
  \item Third item with sub-items:
    \begin{itemize}
      \item Sub-item 1
      \item Sub-item 2
    \end{itemize}
\end{itemize}

\subsection{Numbered List}
\begin{enumerate}
  \item First step
  \item Second step
  \item Third step
\end{enumerate}

\section{Citations and References}
When you use the citation search feature, references will be automatically managed. For example, you might cite a paper like this: \cite{example2024}. All citations will be automatically added to your \texttt{cite.bib} file.

\section{Conclusion}
This template demonstrates basic \LaTeX{} features including mathematical formulas, tables, and lists. The bibliography system is now set up to work seamlessly with Bibby AI's citation management. Feel free to modify it for your needs!

\bibliographystyle{plain}
\bibliography{cite}

\end{document}\documentclass{article}
\usepackage[utf8]{inputenc}
\usepackage{booktabs}
\usepackage{amsmath}
\usepackage{natbib}

\title{My Document}
\author{Your Name}
\date{\today}

\begin{document}

\maketitle

\section{Introduction}
This is a sample document demonstrating various \LaTeX{} features including mathematical equations, tables, and formatting. You can cite papers using the citation search feature, and they will be automatically added to your \texttt{cite.bib} file.

\section{Mathematical Expressions}
Here are some example equations:

\subsection{Quadratic Formula}
The quadratic formula is given by:
\[x = \frac{-b \pm \sqrt{b^2 - 4ac}}{2a}\]

\subsection{Matrix Example}
A sample matrix:
\[
\begin{pmatrix}
a_{11} & a_{12} & a_{13} \\
a_{21} & a_{22} & a_{23} \\
a_{31} & a_{32} & a_{33}
\end{pmatrix}
\]

\section{Tables}
Below is a sample table showing data:

\begin{table}[h]
\centering
\begin{tabular}{@{}lrrr@{}}
\toprule
\textbf{Category} & \textbf{Value 1} & \textbf{Value 2} & \textbf{Total} \\
\midrule
Item A & 10.5 & 20.3 & 30.8 \\
Item B & 15.2 & 25.1 & 40.3 \\
Item C & 22.1 & 12.7 & 34.8 \\
\bottomrule
\end{tabular}
\caption{Sample Data Table}
\label{tab:sample}
\end{table}

\section{Lists}
\subsection{Bullet Points}
\begin{itemize}
  \item First item
  \item Second item
  \item Third item with sub-items:
    \begin{itemize}
      \item Sub-item 1
      \item Sub-item 2
    \end{itemize}
\end{itemize}

\subsection{Numbered List}
\begin{enumerate}
  \item First step
  \item Second step
  \item Third step
\end{enumerate}

\section{Citations and References}
When you use the citation search feature, references will be automatically managed. For example, you might cite a paper like this: \cite{example2024}. All citations will be automatically added to your \texttt{cite.bib} file.

\section{Conclusion}
This template demonstrates basic \LaTeX{} features including mathematical formulas, tables, and lists. The bibliography system is now set up to work seamlessly with Bibby AI's citation management. Feel free to modify it for your needs!

\bibliographystyle{plain}
\bibliography{cite}

\end{document}\documentclass{article}
\usepackage[utf8]{inputenc}
\usepackage{booktabs}
\usepackage{amsmath}
\usepackage{natbib}

\title{My Document}
\author{Your Name}
\date{\today}

\begin{document}

\maketitle

\section{Introduction}
This is a sample document demonstrating various \LaTeX{} features including mathematical equations, tables, and formatting. You can cite papers using the citation search feature, and they will be automatically added to your \texttt{cite.bib} file.

\section{Mathematical Expressions}
Here are some example equations:

\subsection{Quadratic Formula}
The quadratic formula is given by:
\[x = \frac{-b \pm \sqrt{b^2 - 4ac}}{2a}\]

\subsection{Matrix Example}
A sample matrix:
\[
\begin{pmatrix}
a_{11} & a_{12} & a_{13} \\
a_{21} & a_{22} & a_{23} \\
a_{31} & a_{32} & a_{33}
\end{pmatrix}
\]

\section{Tables}
Below is a sample table showing data:

\begin{table}[h]
\centering
\begin{tabular}{@{}lrrr@{}}
\toprule
\textbf{Category} & \textbf{Value 1} & \textbf{Value 2} & \textbf{Total} \\
\midrule
Item A & 10.5 & 20.3 & 30.8 \\
Item B & 15.2 & 25.1 & 40.3 \\
Item C & 22.1 & 12.7 & 34.8 \\
\bottomrule
\end{tabular}
\caption{Sample Data Table}
\label{tab:sample}
\end{table}

\section{Lists}
\subsection{Bullet Points}
\begin{itemize}
  \item First item
  \item Second item
  \item Third item with sub-items:
    \begin{itemize}
      \item Sub-item 1
      \item Sub-item 2
    \end{itemize}
\end{itemize}

\subsection{Numbered List}
\begin{enumerate}
  \item First step
  \item Second step
  \item Third step
\end{enumerate}

\section{Citations and References}
When you use the citation search feature, references will be automatically managed. For example, you might cite a paper like this: \cite{example2024}. All citations will be automatically added to your \texttt{cite.bib} file.

\section{Conclusion}
This template demonstrates basic \LaTeX{} features including mathematical formulas, tables, and lists. The bibliography system is now set up to work seamlessly with Bibby AI's citation management. Feel free to modify it for your needs!

\bibliographystyle{plain}
\bibliography{cite}

\end{document}\documentclass{article}
\usepackage[utf8]{inputenc}
\usepackage{booktabs}
\usepackage{amsmath}
\usepackage{natbib}

\title{My Document}
\author{Your Name}
\date{\today}

\begin{document}

\maketitle

\section{Introduction}
This is a sample document demonstrating various \LaTeX{} features including mathematical equations, tables, and formatting. You can cite papers using the citation search feature, and they will be automatically added to your \texttt{cite.bib} file.

\section{Mathematical Expressions}
Here are some example equations:

\subsection{Quadratic Formula}
The quadratic formula is given by:
\[x = \frac{-b \pm \sqrt{b^2 - 4ac}}{2a}\]

\subsection{Matrix Example}
A sample matrix:
\[
\begin{pmatrix}
a_{11} & a_{12} & a_{13} \\
a_{21} & a_{22} & a_{23} \\
a_{31} & a_{32} & a_{33}
\end{pmatrix}
\]

\section{Tables}
Below is a sample table showing data:

\begin{table}[h]
\centering
\begin{tabular}{@{}lrrr@{}}
\toprule
\textbf{Category} & \textbf{Value 1} & \textbf{Value 2} & \textbf{Total} \\
\midrule
Item A & 10.5 & 20.3 & 30.8 \\
Item B & 15.2 & 25.1 & 40.3 \\
Item C & 22.1 & 12.7 & 34.8 \\
\bottomrule
\end{tabular}
\caption{Sample Data Table}
\label{tab:sample}
\end{table}

\section{Lists}
\subsection{Bullet Points}
\begin{itemize}
  \item First item
  \item Second item
  \item Third item with sub-items:
    \begin{itemize}
      \item Sub-item 1
      \item Sub-item 2
    \end{itemize}
\end{itemize}

\subsection{Numbered List}
\begin{enumerate}
  \item First step
  \item Second step
  \item Third step
\end{enumerate}

\section{Citations and References}
When you use the citation search feature, references will be automatically managed. For example, you might cite a paper like this: \cite{example2024}. All citations will be automatically added to your \texttt{cite.bib} file.

\section{Conclusion}
This template demonstrates basic \LaTeX{} features including mathematical formulas, tables, and lists. The bibliography system is now set up to work seamlessly with Bibby AI's citation management. Feel free to modify it for your needs!

\bibliographystyle{plain}
\bibliography{cite}

\end{document}\documentclass{article}
\usepackage[utf8]{inputenc}
\usepackage{booktabs}
\usepackage{amsmath}
\usepackage{natbib}

\title{My Document}
\author{Your Name}
\date{\today}

\begin{document}

\maketitle

\section{Introduction}
This is a sample document demonstrating various \LaTeX{} features including mathematical equations, tables, and formatting. You can cite papers using the citation search feature, and they will be automatically added to your \texttt{cite.bib} file.

\section{Mathematical Expressions}
Here are some example equations:

\subsection{Quadratic Formula}
The quadratic formula is given by:
\[x = \frac{-b \pm \sqrt{b^2 - 4ac}}{2a}\]

\subsection{Matrix Example}
A sample matrix:
\[
\begin{pmatrix}
a_{11} & a_{12} & a_{13} \\
a_{21} & a_{22} & a_{23} \\
a_{31} & a_{32} & a_{33}
\end{pmatrix}
\]

\section{Tables}
Below is a sample table showing data:

\begin{table}[h]
\centering
\begin{tabular}{@{}lrrr@{}}
\toprule
\textbf{Category} & \textbf{Value 1} & \textbf{Value 2} & \textbf{Total} \\
\midrule
Item A & 10.5 & 20.3 & 30.8 \\
Item B & 15.2 & 25.1 & 40.3 \\
Item C & 22.1 & 12.7 & 34.8 \\
\bottomrule
\end{tabular}
\caption{Sample Data Table}
\label{tab:sample}
\end{table}

\section{Lists}
\subsection{Bullet Points}
\begin{itemize}
  \item First item
  \item Second item
  \item Third item with sub-items:
    \begin{itemize}
      \item Sub-item 1
      \item Sub-item 2
    \end{itemize}
\end{itemize}

\subsection{Numbered List}
\begin{enumerate}
  \item First step
  \item Second step
  \item Third step
\end{enumerate}

\section{Citations and References}
When you use the citation search feature, references will be automatically managed. For example, you might cite a paper like this: \cite{example2024}. All citations will be automatically added to your \texttt{cite.bib} file.

\section{Conclusion}
This template demonstrates basic \LaTeX{} features including mathematical formulas, tables, and lists. The bibliography system is now set up to work seamlessly with Bibby AI's citation management. Feel free to modify it for your needs!

\bibliographystyle{plain}
\bibliography{cite}

\end{document}\documentclass{article}
\usepackage[utf8]{inputenc}
\usepackage{booktabs}
\usepackage{amsmath}
\usepackage{natbib}

\title{My Document}
\author{Your Name}
\date{\today}

\begin{document}

\maketitle

\section{Introduction}
This is a sample document demonstrating various \LaTeX{} features including mathematical equations, tables, and formatting. You can cite papers using the citation search feature, and they will be automatically added to your \texttt{cite.bib} file.

\section{Mathematical Expressions}
Here are some example equations:

\subsection{Quadratic Formula}
The quadratic formula is given by:
\[x = \frac{-b \pm \sqrt{b^2 - 4ac}}{2a}\]

\subsection{Matrix Example}
A sample matrix:
\[
\begin{pmatrix}
a_{11} & a_{12} & a_{13} \\
a_{21} & a_{22} & a_{23} \\
a_{31} & a_{32} & a_{33}
\end{pmatrix}
\]

\section{Tables}
Below is a sample table showing data:

\begin{table}[h]
\centering
\begin{tabular}{@{}lrrr@{}}
\toprule
\textbf{Category} & \textbf{Value 1} & \textbf{Value 2} & \textbf{Total} \\
\midrule
Item A & 10.5 & 20.3 & 30.8 \\
Item B & 15.2 & 25.1 & 40.3 \\
Item C & 22.1 & 12.7 & 34.8 \\
\bottomrule
\end{tabular}
\caption{Sample Data Table}
\label{tab:sample}
\end{table}

\section{Lists}
\subsection{Bullet Points}
\begin{itemize}
  \item First item
  \item Second item
  \item Third item with sub-items:
    \begin{itemize}
      \item Sub-item 1
      \item Sub-item 2
    \end{itemize}
\end{itemize}

\subsection{Numbered List}
\begin{enumerate}
  \item First step
  \item Second step
  \item Third step
\end{enumerate}

\section{Citations and References}
When you use the citation search feature, references will be automatically managed. For example, you might cite a paper like this: \cite{example2024}. All citations will be automatically added to your \texttt{cite.bib} file.

\section{Conclusion}
This template demonstrates basic \LaTeX{} features including mathematical formulas, tables, and lists. The bibliography system is now set up to work seamlessly with Bibby AI's citation management. Feel free to modify it for your needs!

\bibliographystyle{plain}
\bibliography{cite}

\end{document}\documentclass{article}
\usepackage[utf8]{inputenc}
\usepackage{booktabs}
\usepackage{amsmath}
\usepackage{natbib}

\title{My Document}
\author{Your Name}
\date{\today}

\begin{document}

\maketitle

\section{Introduction}
This is a sample document demonstrating various \LaTeX{} features including mathematical equations, tables, and formatting. You can cite papers using the citation search feature, and they will be automatically added to your \texttt{cite.bib} file.

\section{Mathematical Expressions}
Here are some example equations:

\subsection{Quadratic Formula}
The quadratic formula is given by:
\[x = \frac{-b \pm \sqrt{b^2 - 4ac}}{2a}\]

\subsection{Matrix Example}
A sample matrix:
\[
\begin{pmatrix}
a_{11} & a_{12} & a_{13} \\
a_{21} & a_{22} & a_{23} \\
a_{31} & a_{32} & a_{33}
\end{pmatrix}
\]

\section{Tables}
Below is a sample table showing data:

\begin{table}[h]
\centering
\begin{tabular}{@{}lrrr@{}}
\toprule
\textbf{Category} & \textbf{Value 1} & \textbf{Value 2} & \textbf{Total} \\
\midrule
Item A & 10.5 & 20.3 & 30.8 \\
Item B & 15.2 & 25.1 & 40.3 \\
Item C & 22.1 & 12.7 & 34.8 \\
\bottomrule
\end{tabular}
\caption{Sample Data Table}
\label{tab:sample}
\end{table}

\section{Lists}
\subsection{Bullet Points}
\begin{itemize}
  \item First item
  \item Second item
  \item Third item with sub-items:
    \begin{itemize}
      \item Sub-item 1
      \item Sub-item 2
    \end{itemize}
\end{itemize}

\subsection{Numbered List}
\begin{enumerate}
  \item First step
  \item Second step
  \item Third step
\end{enumerate}

\section{Citations and References}
When you use the citation search feature, references will be automatically managed. For example, you might cite a paper like this: \cite{example2024}. All citations will be automatically added to your \texttt{cite.bib} file.

\section{Conclusion}
This template demonstrates basic \LaTeX{} features including mathematical formulas, tables, and lists. The bibliography system is now set up to work seamlessly with Bibby AI's citation management. Feel free to modify it for your needs!

\bibliographystyle{plain}
\bibliography{cite}

\end{document}\documentclass{article}
\usepackage[utf8]{inputenc}
\usepackage{booktabs}
\usepackage{amsmath}
\usepackage{natbib}

\title{My Document}
\author{Your Name}
\date{\today}

\begin{document}

\maketitle

\section{Introduction}
This is a sample document demonstrating various \LaTeX{} features including mathematical equations, tables, and formatting. You can cite papers using the citation search feature, and they will be automatically added to your \texttt{cite.bib} file.

\section{Mathematical Expressions}
Here are some example equations:

\subsection{Quadratic Formula}
The quadratic formula is given by:
\[x = \frac{-b \pm \sqrt{b^2 - 4ac}}{2a}\]

\subsection{Matrix Example}
A sample matrix:
\[
\begin{pmatrix}
a_{11} & a_{12} & a_{13} \\
a_{21} & a_{22} & a_{23} \\
a_{31} & a_{32} & a_{33}
\end{pmatrix}
\]

\section{Tables}
Below is a sample table showing data:

\begin{table}[h]
\centering
\begin{tabular}{@{}lrrr@{}}
\toprule
\textbf{Category} & \textbf{Value 1} & \textbf{Value 2} & \textbf{Total} \\
\midrule
Item A & 10.5 & 20.3 & 30.8 \\
Item B & 15.2 & 25.1 & 40.3 \\
Item C & 22.1 & 12.7 & 34.8 \\
\bottomrule
\end{tabular}
\caption{Sample Data Table}
\label{tab:sample}
\end{table}

\section{Lists}
\subsection{Bullet Points}
\begin{itemize}
  \item First item
  \item Second item
  \item Third item with sub-items:
    \begin{itemize}
      \item Sub-item 1
      \item Sub-item 2
    \end{itemize}
\end{itemize}

\subsection{Numbered List}
\begin{enumerate}
  \item First step
  \item Second step
  \item Third step
\end{enumerate}

\section{Citations and References}
When you use the citation search feature, references will be automatically managed. For example, you might cite a paper like this: \cite{example2024}. All citations will be automatically added to your \texttt{cite.bib} file.

\section{Conclusion}
This template demonstrates basic \LaTeX{} features including mathematical formulas, tables, and lists. The bibliography system is now set up to work seamlessly with Bibby AI's citation management. Feel free to modify it for your needs!

\bibliographystyle{plain}
\bibliography{cite}

\end{document}\documentclass{article}
\usepackage[utf8]{inputenc}
\usepackage{booktabs}
\usepackage{amsmath}
\usepackage{natbib}

\title{My Document}
\author{Your Name}
\date{\today}

\begin{document}

\maketitle

\section{Introduction}
This is a sample document demonstrating various \LaTeX{} features including mathematical equations, tables, and formatting. You can cite papers using the citation search feature, and they will be automatically added to your \texttt{cite.bib} file.

\section{Mathematical Expressions}
Here are some example equations:

\subsection{Quadratic Formula}
The quadratic formula is given by:
\[x = \frac{-b \pm \sqrt{b^2 - 4ac}}{2a}\]

\subsection{Matrix Example}
A sample matrix:
\[
\begin{pmatrix}
a_{11} & a_{12} & a_{13} \\
a_{21} & a_{22} & a_{23} \\
a_{31} & a_{32} & a_{33}
\end{pmatrix}
\]

\section{Tables}
Below is a sample table showing data:

\begin{table}[h]
\centering
\begin{tabular}{@{}lrrr@{}}
\toprule
\textbf{Category} & \textbf{Value 1} & \textbf{Value 2} & \textbf{Total} \\
\midrule
Item A & 10.5 & 20.3 & 30.8 \\
Item B & 15.2 & 25.1 & 40.3 \\
Item C & 22.1 & 12.7 & 34.8 \\
\bottomrule
\end{tabular}
\caption{Sample Data Table}
\label{tab:sample}
\end{table}

\section{Lists}
\subsection{Bullet Points}
\begin{itemize}
  \item First item
  \item Second item
  \item Third item with sub-items:
    \begin{itemize}
      \item Sub-item 1
      \item Sub-item 2
    \end{itemize}
\end{itemize}

\subsection{Numbered List}
\begin{enumerate}
  \item First step
  \item Second step
  \item Third step
\end{enumerate}

\section{Citations and References}
When you use the citation search feature, references will be automatically managed. For example, you might cite a paper like this: \cite{example2024}. All citations will be automatically added to your \texttt{cite.bib} file.

\section{Conclusion}
This template demonstrates basic \LaTeX{} features including mathematical formulas, tables, and lists. The bibliography system is now set up to work seamlessly with Bibby AI's citation management. Feel free to modify it for your needs!

\bibliographystyle{plain}
\bibliography{cite}

\end{document}\documentclass{article}
\usepackage[utf8]{inputenc}
\usepackage{booktabs}
\usepackage{amsmath}
\usepackage{natbib}

\title{My Document}
\author{Your Name}
\date{\today}

\begin{document}

\maketitle

\section{Introduction}
This is a sample document demonstrating various \LaTeX{} features including mathematical equations, tables, and formatting. You can cite papers using the citation search feature, and they will be automatically added to your \texttt{cite.bib} file.

\section{Mathematical Expressions}
Here are some example equations:

\subsection{Quadratic Formula}
The quadratic formula is given by:
\[x = \frac{-b \pm \sqrt{b^2 - 4ac}}{2a}\]

\subsection{Matrix Example}
A sample matrix:
\[
\begin{pmatrix}
a_{11} & a_{12} & a_{13} \\
a_{21} & a_{22} & a_{23} \\
a_{31} & a_{32} & a_{33}
\end{pmatrix}
\]

\section{Tables}
Below is a sample table showing data:

\begin{table}[h]
\centering
\begin{tabular}{@{}lrrr@{}}
\toprule
\textbf{Category} & \textbf{Value 1} & \textbf{Value 2} & \textbf{Total} \\
\midrule
Item A & 10.5 & 20.3 & 30.8 \\
Item B & 15.2 & 25.1 & 40.3 \\
Item C & 22.1 & 12.7 & 34.8 \\
\bottomrule
\end{tabular}
\caption{Sample Data Table}
\label{tab:sample}
\end{table}

\section{Lists}
\subsection{Bullet Points}
\begin{itemize}
  \item First item
  \item Second item
  \item Third item with sub-items:
    \begin{itemize}
      \item Sub-item 1
      \item Sub-item 2
    \end{itemize}
\end{itemize}

\subsection{Numbered List}
\begin{enumerate}
  \item First step
  \item Second step
  \item Third step
\end{enumerate}

\section{Citations and References}
When you use the citation search feature, references will be automatically managed. For example, you might cite a paper like this: \cite{example2024}. All citations will be automatically added to your \texttt{cite.bib} file.

\section{Conclusion}
This template demonstrates basic \LaTeX{} features including mathematical formulas, tables, and lists. The bibliography system is now set up to work seamlessly with Bibby AI's citation management. Feel free to modify it for your needs!

\bibliographystyle{plain}
\bibliography{cite}

\end{document}\documentclass{article}
\usepackage[utf8]{inputenc}
\usepackage{booktabs}
\usepackage{amsmath}
\usepackage{natbib}

\title{My Document}
\author{Your Name}
\date{\today}

\begin{document}

\maketitle

\section{Introduction}
This is a sample document demonstrating various \LaTeX{} features including mathematical equations, tables, and formatting. You can cite papers using the citation search feature, and they will be automatically added to your \texttt{cite.bib} file.

\section{Mathematical Expressions}
Here are some example equations:

\subsection{Quadratic Formula}
The quadratic formula is given by:
\[x = \frac{-b \pm \sqrt{b^2 - 4ac}}{2a}\]

\subsection{Matrix Example}
A sample matrix:
\[
\begin{pmatrix}
a_{11} & a_{12} & a_{13} \\
a_{21} & a_{22} & a_{23} \\
a_{31} & a_{32} & a_{33}
\end{pmatrix}
\]

\section{Tables}
Below is a sample table showing data:

\begin{table}[h]
\centering
\begin{tabular}{@{}lrrr@{}}
\toprule
\textbf{Category} & \textbf{Value 1} & \textbf{Value 2} & \textbf{Total} \\
\midrule
Item A & 10.5 & 20.3 & 30.8 \\
Item B & 15.2 & 25.1 & 40.3 \\
Item C & 22.1 & 12.7 & 34.8 \\
\bottomrule
\end{tabular}
\caption{Sample Data Table}
\label{tab:sample}
\end{table}

\section{Lists}
\subsection{Bullet Points}
\begin{itemize}
  \item First item
  \item Second item
  \item Third item with sub-items:
    \begin{itemize}
      \item Sub-item 1
      \item Sub-item 2
    \end{itemize}
\end{itemize}

\subsection{Numbered List}
\begin{enumerate}
  \item First step
  \item Second step
  \item Third step
\end{enumerate}

\section{Citations and References}
When you use the citation search feature, references will be automatically managed. For example, you might cite a paper like this: \cite{example2024}. All citations will be automatically added to your \texttt{cite.bib} file.

\section{Conclusion}
This template demonstrates basic \LaTeX{} features including mathematical formulas, tables, and lists. The bibliography system is now set up to work seamlessly with Bibby AI's citation management. Feel free to modify it for your needs!

\bibliographystyle{plain}
\bibliography{cite}

\end{document}\documentclass{article}
\usepackage[utf8]{inputenc}
\usepackage{booktabs}
\usepackage{amsmath}
\usepackage{natbib}

\title{My Document}
\author{Your Name}
\date{\today}

\begin{document}

\maketitle

\section{Introduction}
This is a sample document demonstrating various \LaTeX{} features including mathematical equations, tables, and formatting. You can cite papers using the citation search feature, and they will be automatically added to your \texttt{cite.bib} file.

\section{Mathematical Expressions}
Here are some example equations:

\subsection{Quadratic Formula}
The quadratic formula is given by:
\[x = \frac{-b \pm \sqrt{b^2 - 4ac}}{2a}\]

\subsection{Matrix Example}
A sample matrix:
\[
\begin{pmatrix}
a_{11} & a_{12} & a_{13} \\
a_{21} & a_{22} & a_{23} \\
a_{31} & a_{32} & a_{33}
\end{pmatrix}
\]

\section{Tables}
Below is a sample table showing data:

\begin{table}[h]
\centering
\begin{tabular}{@{}lrrr@{}}
\toprule
\textbf{Category} & \textbf{Value 1} & \textbf{Value 2} & \textbf{Total} \\
\midrule
Item A & 10.5 & 20.3 & 30.8 \\
Item B & 15.2 & 25.1 & 40.3 \\
Item C & 22.1 & 12.7 & 34.8 \\
\bottomrule
\end{tabular}
\caption{Sample Data Table}
\label{tab:sample}
\end{table}

\section{Lists}
\subsection{Bullet Points}
\begin{itemize}
  \item First item
  \item Second item
  \item Third item with sub-items:
    \begin{itemize}
      \item Sub-item 1
      \item Sub-item 2
    \end{itemize}
\end{itemize}

\subsection{Numbered List}
\begin{enumerate}
  \item First step
  \item Second step
  \item Third step
\end{enumerate}

\section{Citations and References}
When you use the citation search feature, references will be automatically managed. For example, you might cite a paper like this: \cite{example2024}. All citations will be automatically added to your \texttt{cite.bib} file.

\section{Conclusion}
This template demonstrates basic \LaTeX{} features including mathematical formulas, tables, and lists. The bibliography system is now set up to work seamlessly with Bibby AI's citation management. Feel free to modify it for your needs!

\bibliographystyle{plain}
\bibliography{cite}

\end{document}\documentclass{article}
\usepackage[utf8]{inputenc}
\usepackage{booktabs}
\usepackage{amsmath}
\usepackage{natbib}

\title{My Document}
\author{Your Name}
\date{\today}

\begin{document}

\maketitle

\section{Introduction}
This is a sample document demonstrating various \LaTeX{} features including mathematical equations, tables, and formatting. You can cite papers using the citation search feature, and they will be automatically added to your \texttt{cite.bib} file.

\section{Mathematical Expressions}
Here are some example equations:

\subsection{Quadratic Formula}
The quadratic formula is given by:
\[x = \frac{-b \pm \sqrt{b^2 - 4ac}}{2a}\]

\subsection{Matrix Example}
A sample matrix:
\[
\begin{pmatrix}
a_{11} & a_{12} & a_{13} \\
a_{21} & a_{22} & a_{23} \\
a_{31} & a_{32} & a_{33}
\end{pmatrix}
\]

\section{Tables}
Below is a sample table showing data:

\begin{table}[h]
\centering
\begin{tabular}{@{}lrrr@{}}
\toprule
\textbf{Category} & \textbf{Value 1} & \textbf{Value 2} & \textbf{Total} \\
\midrule
Item A & 10.5 & 20.3 & 30.8 \\
Item B & 15.2 & 25.1 & 40.3 \\
Item C & 22.1 & 12.7 & 34.8 \\
\bottomrule
\end{tabular}
\caption{Sample Data Table}
\label{tab:sample}
\end{table}

\section{Lists}
\subsection{Bullet Points}
\begin{itemize}
  \item First item
  \item Second item
  \item Third item with sub-items:
    \begin{itemize}
      \item Sub-item 1
      \item Sub-item 2
    \end{itemize}
\end{itemize}

\subsection{Numbered List}
\begin{enumerate}
  \item First step
  \item Second step
  \item Third step
\end{enumerate}

\section{Citations and References}
When you use the citation search feature, references will be automatically managed. For example, you might cite a paper like this: \cite{example2024}. All citations will be automatically added to your \texttt{cite.bib} file.

\section{Conclusion}
This template demonstrates basic \LaTeX{} features including mathematical formulas, tables, and lists. The bibliography system is now set up to work seamlessly with Bibby AI's citation management. Feel free to modify it for your needs!

\bibliographystyle{plain}
\bibliography{cite}

\end{document}\documentclass{article}
\usepackage[utf8]{inputenc}
\usepackage{booktabs}
\usepackage{amsmath}
\usepackage{natbib}

\title{My Document}
\author{Your Name}
\date{\today}

\begin{document}

\maketitle

\section{Introduction}
This is a sample document demonstrating various \LaTeX{} features including mathematical equations, tables, and formatting. You can cite papers using the citation search feature, and they will be automatically added to your \texttt{cite.bib} file.

\section{Mathematical Expressions}
Here are some example equations:

\subsection{Quadratic Formula}
The quadratic formula is given by:
\[x = \frac{-b \pm \sqrt{b^2 - 4ac}}{2a}\]

\subsection{Matrix Example}
A sample matrix:
\[
\begin{pmatrix}
a_{11} & a_{12} & a_{13} \\
a_{21} & a_{22} & a_{23} \\
a_{31} & a_{32} & a_{33}
\end{pmatrix}
\]

\section{Tables}
Below is a sample table showing data:

\begin{table}[h]
\centering
\begin{tabular}{@{}lrrr@{}}
\toprule
\textbf{Category} & \textbf{Value 1} & \textbf{Value 2} & \textbf{Total} \\
\midrule
Item A & 10.5 & 20.3 & 30.8 \\
Item B & 15.2 & 25.1 & 40.3 \\
Item C & 22.1 & 12.7 & 34.8 \\
\bottomrule
\end{tabular}
\caption{Sample Data Table}
\label{tab:sample}
\end{table}

\section{Lists}
\subsection{Bullet Points}
\begin{itemize}
  \item First item
  \item Second item
  \item Third item with sub-items:
    \begin{itemize}
      \item Sub-item 1
      \item Sub-item 2
    \end{itemize}
\end{itemize}

\subsection{Numbered List}
\begin{enumerate}
  \item First step
  \item Second step
  \item Third step
\end{enumerate}

\section{Citations and References}
When you use the citation search feature, references will be automatically managed. For example, you might cite a paper like this: \cite{example2024}. All citations will be automatically added to your \texttt{cite.bib} file.

\section{Conclusion}
This template demonstrates basic \LaTeX{} features including mathematical formulas, tables, and lists. The bibliography system is now set up to work seamlessly with Bibby AI's citation management. Feel free to modify it for your needs!

\bibliographystyle{plain}
\bibliography{cite}

\end{document}\documentclass{article}
\usepackage[utf8]{inputenc}
\usepackage{booktabs}
\usepackage{amsmath}
\usepackage{natbib}

\title{My Document}
\author{Your Name}
\date{\today}

\begin{document}

\maketitle

\section{Introduction}
This is a sample document demonstrating various \LaTeX{} features including mathematical equations, tables, and formatting. You can cite papers using the citation search feature, and they will be automatically added to your \texttt{cite.bib} file.

\section{Mathematical Expressions}
Here are some example equations:

\subsection{Quadratic Formula}
The quadratic formula is given by:
\[x = \frac{-b \pm \sqrt{b^2 - 4ac}}{2a}\]

\subsection{Matrix Example}
A sample matrix:
\[
\begin{pmatrix}
a_{11} & a_{12} & a_{13} \\
a_{21} & a_{22} & a_{23} \\
a_{31} & a_{32} & a_{33}
\end{pmatrix}
\]

\section{Tables}
Below is a sample table showing data:

\begin{table}[h]
\centering
\begin{tabular}{@{}lrrr@{}}
\toprule
\textbf{Category} & \textbf{Value 1} & \textbf{Value 2} & \textbf{Total} \\
\midrule
Item A & 10.5 & 20.3 & 30.8 \\
Item B & 15.2 & 25.1 & 40.3 \\
Item C & 22.1 & 12.7 & 34.8 \\
\bottomrule
\end{tabular}
\caption{Sample Data Table}
\label{tab:sample}
\end{table}

\section{Lists}
\subsection{Bullet Points}
\begin{itemize}
  \item First item
  \item Second item
  \item Third item with sub-items:
    \begin{itemize}
      \item Sub-item 1
      \item Sub-item 2
    \end{itemize}
\end{itemize}

\subsection{Numbered List}
\begin{enumerate}
  \item First step
  \item Second step
  \item Third step
\end{enumerate}

\section{Citations and References}
When you use the citation search feature, references will be automatically managed. For example, you might cite a paper like this: \cite{example2024}. All citations will be automatically added to your \texttt{cite.bib} file.

\section{Conclusion}
This template demonstrates basic \LaTeX{} features including mathematical formulas, tables, and lists. The bibliography system is now set up to work seamlessly with Bibby AI's citation management. Feel free to modify it for your needs!

\bibliographystyle{plain}
\bibliography{cite}

\end{document}\documentclass{article}
\usepackage[utf8]{inputenc}
\usepackage{booktabs}
\usepackage{amsmath}
\usepackage{natbib}

\title{My Document}
\author{Your Name}
\date{\today}

\begin{document}

\maketitle

\section{Introduction}
This is a sample document demonstrating various \LaTeX{} features including mathematical equations, tables, and formatting. You can cite papers using the citation search feature, and they will be automatically added to your \texttt{cite.bib} file.

\section{Mathematical Expressions}
Here are some example equations:

\subsection{Quadratic Formula}
The quadratic formula is given by:
\[x = \frac{-b \pm \sqrt{b^2 - 4ac}}{2a}\]

\subsection{Matrix Example}
A sample matrix:
\[
\begin{pmatrix}
a_{11} & a_{12} & a_{13} \\
a_{21} & a_{22} & a_{23} \\
a_{31} & a_{32} & a_{33}
\end{pmatrix}
\]

\section{Tables}
Below is a sample table showing data:

\begin{table}[h]
\centering
\begin{tabular}{@{}lrrr@{}}
\toprule
\textbf{Category} & \textbf{Value 1} & \textbf{Value 2} & \textbf{Total} \\
\midrule
Item A & 10.5 & 20.3 & 30.8 \\
Item B & 15.2 & 25.1 & 40.3 \\
Item C & 22.1 & 12.7 & 34.8 \\
\bottomrule
\end{tabular}
\caption{Sample Data Table}
\label{tab:sample}
\end{table}

\section{Lists}
\subsection{Bullet Points}
\begin{itemize}
  \item First item
  \item Second item
  \item Third item with sub-items:
    \begin{itemize}
      \item Sub-item 1
      \item Sub-item 2
    \end{itemize}
\end{itemize}

\subsection{Numbered List}
\begin{enumerate}
  \item First step
  \item Second step
  \item Third step
\end{enumerate}

\section{Citations and References}
When you use the citation search feature, references will be automatically managed. For example, you might cite a paper like this: \cite{example2024}. All citations will be automatically added to your \texttt{cite.bib} file.

\section{Conclusion}
This template demonstrates basic \LaTeX{} features including mathematical formulas, tables, and lists. The bibliography system is now set up to work seamlessly with Bibby AI's citation management. Feel free to modify it for your needs!

\bibliographystyle{plain}
\bibliography{cite}

\end{document}\documentclass{article}
\usepackage[utf8]{inputenc}
\usepackage{booktabs}
\usepackage{amsmath}
\usepackage{natbib}

\title{My Document}
\author{Your Name}
\date{\today}

\begin{document}

\maketitle

\section{Introduction}
This is a sample document demonstrating various \LaTeX{} features including mathematical equations, tables, and formatting. You can cite papers using the citation search feature, and they will be automatically added to your \texttt{cite.bib} file.

\section{Mathematical Expressions}
Here are some example equations:

\subsection{Quadratic Formula}
The quadratic formula is given by:
\[x = \frac{-b \pm \sqrt{b^2 - 4ac}}{2a}\]

\subsection{Matrix Example}
A sample matrix:
\[
\begin{pmatrix}
a_{11} & a_{12} & a_{13} \\
a_{21} & a_{22} & a_{23} \\
a_{31} & a_{32} & a_{33}
\end{pmatrix}
\]

\section{Tables}
Below is a sample table showing data:

\begin{table}[h]
\centering
\begin{tabular}{@{}lrrr@{}}
\toprule
\textbf{Category} & \textbf{Value 1} & \textbf{Value 2} & \textbf{Total} \\
\midrule
Item A & 10.5 & 20.3 & 30.8 \\
Item B & 15.2 & 25.1 & 40.3 \\
Item C & 22.1 & 12.7 & 34.8 \\
\bottomrule
\end{tabular}
\caption{Sample Data Table}
\label{tab:sample}
\end{table}

\section{Lists}
\subsection{Bullet Points}
\begin{itemize}
  \item First item
  \item Second item
  \item Third item with sub-items:
    \begin{itemize}
      \item Sub-item 1
      \item Sub-item 2
    \end{itemize}
\end{itemize}

\subsection{Numbered List}
\begin{enumerate}
  \item First step
  \item Second step
  \item Third step
\end{enumerate}

\section{Citations and References}
When you use the citation search feature, references will be automatically managed. For example, you might cite a paper like this: \cite{example2024}. All citations will be automatically added to your \texttt{cite.bib} file.

\section{Conclusion}
This template demonstrates basic \LaTeX{} features including mathematical formulas, tables, and lists. The bibliography system is now set up to work seamlessly with Bibby AI's citation management. Feel free to modify it for your needs!

\bibliographystyle{plain}
\bibliography{cite}

\end{document}\documentclass{article}
\usepackage[utf8]{inputenc}
\usepackage{booktabs}
\usepackage{amsmath}
\usepackage{natbib}

\title{My Document}
\author{Your Name}
\date{\today}

\begin{document}

\maketitle

\section{Introduction}
This is a sample document demonstrating various \LaTeX{} features including mathematical equations, tables, and formatting. You can cite papers using the citation search feature, and they will be automatically added to your \texttt{cite.bib} file.

\section{Mathematical Expressions}
Here are some example equations:

\subsection{Quadratic Formula}
The quadratic formula is given by:
\[x = \frac{-b \pm \sqrt{b^2 - 4ac}}{2a}\]

\subsection{Matrix Example}
A sample matrix:
\[
\begin{pmatrix}
a_{11} & a_{12} & a_{13} \\
a_{21} & a_{22} & a_{23} \\
a_{31} & a_{32} & a_{33}
\end{pmatrix}
\]

\section{Tables}
Below is a sample table showing data:

\begin{table}[h]
\centering
\begin{tabular}{@{}lrrr@{}}
\toprule
\textbf{Category} & \textbf{Value 1} & \textbf{Value 2} & \textbf{Total} \\
\midrule
Item A & 10.5 & 20.3 & 30.8 \\
Item B & 15.2 & 25.1 & 40.3 \\
Item C & 22.1 & 12.7 & 34.8 \\
\bottomrule
\end{tabular}
\caption{Sample Data Table}
\label{tab:sample}
\end{table}

\section{Lists}
\subsection{Bullet Points}
\begin{itemize}
  \item First item
  \item Second item
  \item Third item with sub-items:
    \begin{itemize}
      \item Sub-item 1
      \item Sub-item 2
    \end{itemize}
\end{itemize}

\subsection{Numbered List}
\begin{enumerate}
  \item First step
  \item Second step
  \item Third step
\end{enumerate}

\section{Citations and References}
When you use the citation search feature, references will be automatically managed. For example, you might cite a paper like this: \cite{example2024}. All citations will be automatically added to your \texttt{cite.bib} file.

\section{Conclusion}
This template demonstrates basic \LaTeX{} features including mathematical formulas, tables, and lists. The bibliography system is now set up to work seamlessly with Bibby AI's citation management. Feel free to modify it for your needs!

\bibliographystyle{plain}
\bibliography{cite}

\end{document}\documentclass{article}
\usepackage[utf8]{inputenc}
\usepackage{booktabs}
\usepackage{amsmath}
\usepackage{natbib}

\title{My Document}
\author{Your Name}
\date{\today}

\begin{document}

\maketitle

\section{Introduction}
This is a sample document demonstrating various \LaTeX{} features including mathematical equations, tables, and formatting. You can cite papers using the citation search feature, and they will be automatically added to your \texttt{cite.bib} file.

\section{Mathematical Expressions}
Here are some example equations:

\subsection{Quadratic Formula}
The quadratic formula is given by:
\[x = \frac{-b \pm \sqrt{b^2 - 4ac}}{2a}\]

\subsection{Matrix Example}
A sample matrix:
\[
\begin{pmatrix}
a_{11} & a_{12} & a_{13} \\
a_{21} & a_{22} & a_{23} \\
a_{31} & a_{32} & a_{33}
\end{pmatrix}
\]

\section{Tables}
Below is a sample table showing data:

\begin{table}[h]
\centering
\begin{tabular}{@{}lrrr@{}}
\toprule
\textbf{Category} & \textbf{Value 1} & \textbf{Value 2} & \textbf{Total} \\
\midrule
Item A & 10.5 & 20.3 & 30.8 \\
Item B & 15.2 & 25.1 & 40.3 \\
Item C & 22.1 & 12.7 & 34.8 \\
\bottomrule
\end{tabular}
\caption{Sample Data Table}
\label{tab:sample}
\end{table}

\section{Lists}
\subsection{Bullet Points}
\begin{itemize}
  \item First item
  \item Second item
  \item Third item with sub-items:
    \begin{itemize}
      \item Sub-item 1
      \item Sub-item 2
    \end{itemize}
\end{itemize}

\subsection{Numbered List}
\begin{enumerate}
  \item First step
  \item Second step
  \item Third step
\end{enumerate}

\section{Citations and References}
When you use the citation search feature, references will be automatically managed. For example, you might cite a paper like this: \cite{example2024}. All citations will be automatically added to your \texttt{cite.bib} file.

\section{Conclusion}
This template demonstrates basic \LaTeX{} features including mathematical formulas, tables, and lists. The bibliography system is now set up to work seamlessly with Bibby AI's citation management. Feel free to modify it for your needs!

\bibliographystyle{plain}
\bibliography{cite}

\end{document}\documentclass{article}
\usepackage[utf8]{inputenc}
\usepackage{booktabs}
\usepackage{amsmath}
\usepackage{natbib}

\title{My Document}
\author{Your Name}
\date{\today}

\begin{document}

\maketitle

\section{Introduction}
This is a sample document demonstrating various \LaTeX{} features including mathematical equations, tables, and formatting. You can cite papers using the citation search feature, and they will be automatically added to your \texttt{cite.bib} file.

\section{Mathematical Expressions}
Here are some example equations:

\subsection{Quadratic Formula}
The quadratic formula is given by:
\[x = \frac{-b \pm \sqrt{b^2 - 4ac}}{2a}\]

\subsection{Matrix Example}
A sample matrix:
\[
\begin{pmatrix}
a_{11} & a_{12} & a_{13} \\
a_{21} & a_{22} & a_{23} \\
a_{31} & a_{32} & a_{33}
\end{pmatrix}
\]

\section{Tables}
Below is a sample table showing data:

\begin{table}[h]
\centering
\begin{tabular}{@{}lrrr@{}}
\toprule
\textbf{Category} & \textbf{Value 1} & \textbf{Value 2} & \textbf{Total} \\
\midrule
Item A & 10.5 & 20.3 & 30.8 \\
Item B & 15.2 & 25.1 & 40.3 \\
Item C & 22.1 & 12.7 & 34.8 \\
\bottomrule
\end{tabular}
\caption{Sample Data Table}
\label{tab:sample}
\end{table}

\section{Lists}
\subsection{Bullet Points}
\begin{itemize}
  \item First item
  \item Second item
  \item Third item with sub-items:
    \begin{itemize}
      \item Sub-item 1
      \item Sub-item 2
    \end{itemize}
\end{itemize}

\subsection{Numbered List}
\begin{enumerate}
  \item First step
  \item Second step
  \item Third step
\end{enumerate}

\section{Citations and References}
When you use the citation search feature, references will be automatically managed. For example, you might cite a paper like this: \cite{example2024}. All citations will be automatically added to your \texttt{cite.bib} file.

\section{Conclusion}
This template demonstrates basic \LaTeX{} features including mathematical formulas, tables, and lists. The bibliography system is now set up to work seamlessly with Bibby AI's citation management. Feel free to modify it for your needs!

\bibliographystyle{plain}
\bibliography{cite}

\end{document}\documentclass{article}
\usepackage[utf8]{inputenc}
\usepackage{booktabs}
\usepackage{amsmath}
\usepackage{natbib}

\title{My Document}
\author{Your Name}
\date{\today}

\begin{document}

\maketitle

\section{Introduction}
This is a sample document demonstrating various \LaTeX{} features including mathematical equations, tables, and formatting. You can cite papers using the citation search feature, and they will be automatically added to your \texttt{cite.bib} file.

\section{Mathematical Expressions}
Here are some example equations:

\subsection{Quadratic Formula}
The quadratic formula is given by:
\[x = \frac{-b \pm \sqrt{b^2 - 4ac}}{2a}\]

\subsection{Matrix Example}
A sample matrix:
\[
\begin{pmatrix}
a_{11} & a_{12} & a_{13} \\
a_{21} & a_{22} & a_{23} \\
a_{31} & a_{32} & a_{33}
\end{pmatrix}
\]

\section{Tables}
Below is a sample table showing data:

\begin{table}[h]
\centering
\begin{tabular}{@{}lrrr@{}}
\toprule
\textbf{Category} & \textbf{Value 1} & \textbf{Value 2} & \textbf{Total} \\
\midrule
Item A & 10.5 & 20.3 & 30.8 \\
Item B & 15.2 & 25.1 & 40.3 \\
Item C & 22.1 & 12.7 & 34.8 \\
\bottomrule
\end{tabular}
\caption{Sample Data Table}
\label{tab:sample}
\end{table}

\section{Lists}
\subsection{Bullet Points}
\begin{itemize}
  \item First item
  \item Second item
  \item Third item with sub-items:
    \begin{itemize}
      \item Sub-item 1
      \item Sub-item 2
    \end{itemize}
\end{itemize}

\subsection{Numbered List}
\begin{enumerate}
  \item First step
  \item Second step
  \item Third step
\end{enumerate}

\section{Citations and References}
When you use the citation search feature, references will be automatically managed. For example, you might cite a paper like this: \cite{example2024}. All citations will be automatically added to your \texttt{cite.bib} file.

\section{Conclusion}
This template demonstrates basic \LaTeX{} features including mathematical formulas, tables, and lists. The bibliography system is now set up to work seamlessly with Bibby AI's citation management. Feel free to modify it for your needs!

\bibliographystyle{plain}
\bibliography{cite}

\end{document}\documentclass{article}
\usepackage[utf8]{inputenc}
\usepackage{booktabs}
\usepackage{amsmath}
\usepackage{natbib}

\title{My Document}
\author{Your Name}
\date{\today}

\begin{document}

\maketitle

\section{Introduction}
This is a sample document demonstrating various \LaTeX{} features including mathematical equations, tables, and formatting. You can cite papers using the citation search feature, and they will be automatically added to your \texttt{cite.bib} file.

\section{Mathematical Expressions}
Here are some example equations:

\subsection{Quadratic Formula}
The quadratic formula is given by:
\[x = \frac{-b \pm \sqrt{b^2 - 4ac}}{2a}\]

\subsection{Matrix Example}
A sample matrix:
\[
\begin{pmatrix}
a_{11} & a_{12} & a_{13} \\
a_{21} & a_{22} & a_{23} \\
a_{31} & a_{32} & a_{33}
\end{pmatrix}
\]

\section{Tables}
Below is a sample table showing data:

\begin{table}[h]
\centering
\begin{tabular}{@{}lrrr@{}}
\toprule
\textbf{Category} & \textbf{Value 1} & \textbf{Value 2} & \textbf{Total} \\
\midrule
Item A & 10.5 & 20.3 & 30.8 \\
Item B & 15.2 & 25.1 & 40.3 \\
Item C & 22.1 & 12.7 & 34.8 \\
\bottomrule
\end{tabular}
\caption{Sample Data Table}
\label{tab:sample}
\end{table}

\section{Lists}
\subsection{Bullet Points}
\begin{itemize}
  \item First item
  \item Second item
  \item Third item with sub-items:
    \begin{itemize}
      \item Sub-item 1
      \item Sub-item 2
    \end{itemize}
\end{itemize}

\subsection{Numbered List}
\begin{enumerate}
  \item First step
  \item Second step
  \item Third step
\end{enumerate}

\section{Citations and References}
When you use the citation search feature, references will be automatically managed. For example, you might cite a paper like this: \cite{example2024}. All citations will be automatically added to your \texttt{cite.bib} file.

\section{Conclusion}
This template demonstrates basic \LaTeX{} features including mathematical formulas, tables, and lists. The bibliography system is now set up to work seamlessly with Bibby AI's citation management. Feel free to modify it for your needs!

\bibliographystyle{plain}
\bibliography{cite}

\end{document}\documentclass{article}
\usepackage[utf8]{inputenc}
\usepackage{booktabs}
\usepackage{amsmath}
\usepackage{natbib}

\title{My Document}
\author{Your Name}
\date{\today}

\begin{document}

\maketitle

\section{Introduction}
This is a sample document demonstrating various \LaTeX{} features including mathematical equations, tables, and formatting. You can cite papers using the citation search feature, and they will be automatically added to your \texttt{cite.bib} file.

\section{Mathematical Expressions}
Here are some example equations:

\subsection{Quadratic Formula}
The quadratic formula is given by:
\[x = \frac{-b \pm \sqrt{b^2 - 4ac}}{2a}\]

\subsection{Matrix Example}
A sample matrix:
\[
\begin{pmatrix}
a_{11} & a_{12} & a_{13} \\
a_{21} & a_{22} & a_{23} \\
a_{31} & a_{32} & a_{33}
\end{pmatrix}
\]

\section{Tables}
Below is a sample table showing data:

\begin{table}[h]
\centering
\begin{tabular}{@{}lrrr@{}}
\toprule
\textbf{Category} & \textbf{Value 1} & \textbf{Value 2} & \textbf{Total} \\
\midrule
Item A & 10.5 & 20.3 & 30.8 \\
Item B & 15.2 & 25.1 & 40.3 \\
Item C & 22.1 & 12.7 & 34.8 \\
\bottomrule
\end{tabular}
\caption{Sample Data Table}
\label{tab:sample}
\end{table}

\section{Lists}
\subsection{Bullet Points}
\begin{itemize}
  \item First item
  \item Second item
  \item Third item with sub-items:
    \begin{itemize}
      \item Sub-item 1
      \item Sub-item 2
    \end{itemize}
\end{itemize}

\subsection{Numbered List}
\begin{enumerate}
  \item First step
  \item Second step
  \item Third step
\end{enumerate}

\section{Citations and References}
When you use the citation search feature, references will be automatically managed. For example, you might cite a paper like this: \cite{example2024}. All citations will be automatically added to your \texttt{cite.bib} file.

\section{Conclusion}
This template demonstrates basic \LaTeX{} features including mathematical formulas, tables, and lists. The bibliography system is now set up to work seamlessly with Bibby AI's citation management. Feel free to modify it for your needs!

\bibliographystyle{plain}
\bibliography{cite}

\end{document}\documentclass{article}
\usepackage[utf8]{inputenc}
\usepackage{booktabs}
\usepackage{amsmath}
\usepackage{natbib}

\title{My Document}
\author{Your Name}
\date{\today}

\begin{document}

\maketitle

\section{Introduction}
This is a sample document demonstrating various \LaTeX{} features including mathematical equations, tables, and formatting. You can cite papers using the citation search feature, and they will be automatically added to your \texttt{cite.bib} file.

\section{Mathematical Expressions}
Here are some example equations:

\subsection{Quadratic Formula}
The quadratic formula is given by:
\[x = \frac{-b \pm \sqrt{b^2 - 4ac}}{2a}\]

\subsection{Matrix Example}
A sample matrix:
\[
\begin{pmatrix}
a_{11} & a_{12} & a_{13} \\
a_{21} & a_{22} & a_{23} \\
a_{31} & a_{32} & a_{33}
\end{pmatrix}
\]

\section{Tables}
Below is a sample table showing data:

\begin{table}[h]
\centering
\begin{tabular}{@{}lrrr@{}}
\toprule
\textbf{Category} & \textbf{Value 1} & \textbf{Value 2} & \textbf{Total} \\
\midrule
Item A & 10.5 & 20.3 & 30.8 \\
Item B & 15.2 & 25.1 & 40.3 \\
Item C & 22.1 & 12.7 & 34.8 \\
\bottomrule
\end{tabular}
\caption{Sample Data Table}
\label{tab:sample}
\end{table}

\section{Lists}
\subsection{Bullet Points}
\begin{itemize}
  \item First item
  \item Second item
  \item Third item with sub-items:
    \begin{itemize}
      \item Sub-item 1
      \item Sub-item 2
    \end{itemize}
\end{itemize}

\subsection{Numbered List}
\begin{enumerate}
  \item First step
  \item Second step
  \item Third step
\end{enumerate}

\section{Citations and References}
When you use the citation search feature, references will be automatically managed. For example, you might cite a paper like this: \cite{example2024}. All citations will be automatically added to your \texttt{cite.bib} file.

\section{Conclusion}
This template demonstrates basic \LaTeX{} features including mathematical formulas, tables, and lists. The bibliography system is now set up to work seamlessly with Bibby AI's citation management. Feel free to modify it for your needs!

\bibliographystyle{plain}
\bibliography{cite}

\end{document}\documentclass{article}
\usepackage[utf8]{inputenc}
\usepackage{booktabs}
\usepackage{amsmath}
\usepackage{natbib}

\title{My Document}
\author{Your Name}
\date{\today}

\begin{document}

\maketitle

\section{Introduction}
This is a sample document demonstrating various \LaTeX{} features including mathematical equations, tables, and formatting. You can cite papers using the citation search feature, and they will be automatically added to your \texttt{cite.bib} file.

\section{Mathematical Expressions}
Here are some example equations:

\subsection{Quadratic Formula}
The quadratic formula is given by:
\[x = \frac{-b \pm \sqrt{b^2 - 4ac}}{2a}\]

\subsection{Matrix Example}
A sample matrix:
\[
\begin{pmatrix}
a_{11} & a_{12} & a_{13} \\
a_{21} & a_{22} & a_{23} \\
a_{31} & a_{32} & a_{33}
\end{pmatrix}
\]

\section{Tables}
Below is a sample table showing data:

\begin{table}[h]
\centering
\begin{tabular}{@{}lrrr@{}}
\toprule
\textbf{Category} & \textbf{Value 1} & \textbf{Value 2} & \textbf{Total} \\
\midrule
Item A & 10.5 & 20.3 & 30.8 \\
Item B & 15.2 & 25.1 & 40.3 \\
Item C & 22.1 & 12.7 & 34.8 \\
\bottomrule
\end{tabular}
\caption{Sample Data Table}
\label{tab:sample}
\end{table}

\section{Lists}
\subsection{Bullet Points}
\begin{itemize}
  \item First item
  \item Second item
  \item Third item with sub-items:
    \begin{itemize}
      \item Sub-item 1
      \item Sub-item 2
    \end{itemize}
\end{itemize}

\subsection{Numbered List}
\begin{enumerate}
  \item First step
  \item Second step
  \item Third step
\end{enumerate}

\section{Citations and References}
When you use the citation search feature, references will be automatically managed. For example, you might cite a paper like this: \cite{example2024}. All citations will be automatically added to your \texttt{cite.bib} file.

\section{Conclusion}
This template demonstrates basic \LaTeX{} features including mathematical formulas, tables, and lists. The bibliography system is now set up to work seamlessly with Bibby AI's citation management. Feel free to modify it for your needs!

\bibliographystyle{plain}
\bibliography{cite}

\end{document}\documentclass{article}
\usepackage[utf8]{inputenc}
\usepackage{booktabs}
\usepackage{amsmath}
\usepackage{natbib}

\title{My Document}
\author{Your Name}
\date{\today}

\begin{document}

\maketitle

\section{Introduction}
This is a sample document demonstrating various \LaTeX{} features including mathematical equations, tables, and formatting. You can cite papers using the citation search feature, and they will be automatically added to your \texttt{cite.bib} file.

\section{Mathematical Expressions}
Here are some example equations:

\subsection{Quadratic Formula}
The quadratic formula is given by:
\[x = \frac{-b \pm \sqrt{b^2 - 4ac}}{2a}\]

\subsection{Matrix Example}
A sample matrix:
\[
\begin{pmatrix}
a_{11} & a_{12} & a_{13} \\
a_{21} & a_{22} & a_{23} \\
a_{31} & a_{32} & a_{33}
\end{pmatrix}
\]

\section{Tables}
Below is a sample table showing data:

\begin{table}[h]
\centering
\begin{tabular}{@{}lrrr@{}}
\toprule
\textbf{Category} & \textbf{Value 1} & \textbf{Value 2} & \textbf{Total} \\
\midrule
Item A & 10.5 & 20.3 & 30.8 \\
Item B & 15.2 & 25.1 & 40.3 \\
Item C & 22.1 & 12.7 & 34.8 \\
\bottomrule
\end{tabular}
\caption{Sample Data Table}
\label{tab:sample}
\end{table}

\section{Lists}
\subsection{Bullet Points}
\begin{itemize}
  \item First item
  \item Second item
  \item Third item with sub-items:
    \begin{itemize}
      \item Sub-item 1
      \item Sub-item 2
    \end{itemize}
\end{itemize}

\subsection{Numbered List}
\begin{enumerate}
  \item First step
  \item Second step
  \item Third step
\end{enumerate}

\section{Citations and References}
When you use the citation search feature, references will be automatically managed. For example, you might cite a paper like this: \cite{example2024}. All citations will be automatically added to your \texttt{cite.bib} file.

\section{Conclusion}
This template demonstrates basic \LaTeX{} features including mathematical formulas, tables, and lists. The bibliography system is now set up to work seamlessly with Bibby AI's citation management. Feel free to modify it for your needs!

\bibliographystyle{plain}
\bibliography{cite}

\end{document}\documentclass{article}
\usepackage[utf8]{inputenc}
\usepackage{booktabs}
\usepackage{amsmath}
\usepackage{natbib}

\title{My Document}
\author{Your Name}
\date{\today}

\begin{document}

\maketitle

\section{Introduction}
This is a sample document demonstrating various \LaTeX{} features including mathematical equations, tables, and formatting. You can cite papers using the citation search feature, and they will be automatically added to your \texttt{cite.bib} file.

\section{Mathematical Expressions}
Here are some example equations:

\subsection{Quadratic Formula}
The quadratic formula is given by:
\[x = \frac{-b \pm \sqrt{b^2 - 4ac}}{2a}\]

\subsection{Matrix Example}
A sample matrix:
\[
\begin{pmatrix}
a_{11} & a_{12} & a_{13} \\
a_{21} & a_{22} & a_{23} \\
a_{31} & a_{32} & a_{33}
\end{pmatrix}
\]

\section{Tables}
Below is a sample table showing data:

\begin{table}[h]
\centering
\begin{tabular}{@{}lrrr@{}}
\toprule
\textbf{Category} & \textbf{Value 1} & \textbf{Value 2} & \textbf{Total} \\
\midrule
Item A & 10.5 & 20.3 & 30.8 \\
Item B & 15.2 & 25.1 & 40.3 \\
Item C & 22.1 & 12.7 & 34.8 \\
\bottomrule
\end{tabular}
\caption{Sample Data Table}
\label{tab:sample}
\end{table}

\section{Lists}
\subsection{Bullet Points}
\begin{itemize}
  \item First item
  \item Second item
  \item Third item with sub-items:
    \begin{itemize}
      \item Sub-item 1
      \item Sub-item 2
    \end{itemize}
\end{itemize}

\subsection{Numbered List}
\begin{enumerate}
  \item First step
  \item Second step
  \item Third step
\end{enumerate}

\section{Citations and References}
When you use the citation search feature, references will be automatically managed. For example, you might cite a paper like this: \cite{example2024}. All citations will be automatically added to your \texttt{cite.bib} file.

\section{Conclusion}
This template demonstrates basic \LaTeX{} features including mathematical formulas, tables, and lists. The bibliography system is now set up to work seamlessly with Bibby AI's citation management. Feel free to modify it for your needs!

\bibliographystyle{plain}
\bibliography{cite}

\end{document}\documentclass{article}
\usepackage[utf8]{inputenc}
\usepackage{booktabs}
\usepackage{amsmath}
\usepackage{natbib}

\title{My Document}
\author{Your Name}
\date{\today}

\begin{document}

\maketitle

\section{Introduction}
This is a sample document demonstrating various \LaTeX{} features including mathematical equations, tables, and formatting. You can cite papers using the citation search feature, and they will be automatically added to your \texttt{cite.bib} file.

\section{Mathematical Expressions}
Here are some example equations:

\subsection{Quadratic Formula}
The quadratic formula is given by:
\[x = \frac{-b \pm \sqrt{b^2 - 4ac}}{2a}\]

\subsection{Matrix Example}
A sample matrix:
\[
\begin{pmatrix}
a_{11} & a_{12} & a_{13} \\
a_{21} & a_{22} & a_{23} \\
a_{31} & a_{32} & a_{33}
\end{pmatrix}
\]

\section{Tables}
Below is a sample table showing data:

\begin{table}[h]
\centering
\begin{tabular}{@{}lrrr@{}}
\toprule
\textbf{Category} & \textbf{Value 1} & \textbf{Value 2} & \textbf{Total} \\
\midrule
Item A & 10.5 & 20.3 & 30.8 \\
Item B & 15.2 & 25.1 & 40.3 \\
Item C & 22.1 & 12.7 & 34.8 \\
\bottomrule
\end{tabular}
\caption{Sample Data Table}
\label{tab:sample}
\end{table}

\section{Lists}
\subsection{Bullet Points}
\begin{itemize}
  \item First item
  \item Second item
  \item Third item with sub-items:
    \begin{itemize}
      \item Sub-item 1
      \item Sub-item 2
    \end{itemize}
\end{itemize}

\subsection{Numbered List}
\begin{enumerate}
  \item First step
  \item Second step
  \item Third step
\end{enumerate}

\section{Citations and References}
When you use the citation search feature, references will be automatically managed. For example, you might cite a paper like this: \cite{example2024}. All citations will be automatically added to your \texttt{cite.bib} file.

\section{Conclusion}
This template demonstrates basic \LaTeX{} features including mathematical formulas, tables, and lists. The bibliography system is now set up to work seamlessly with Bibby AI's citation management. Feel free to modify it for your needs!

\bibliographystyle{plain}
\bibliography{cite}

\end{document}\documentclass{article}
\usepackage[utf8]{inputenc}
\usepackage{booktabs}
\usepackage{amsmath}
\usepackage{natbib}

\title{My Document}
\author{Your Name}
\date{\today}

\begin{document}

\maketitle

\section{Introduction}
This is a sample document demonstrating various \LaTeX{} features including mathematical equations, tables, and formatting. You can cite papers using the citation search feature, and they will be automatically added to your \texttt{cite.bib} file.

\section{Mathematical Expressions}
Here are some example equations:

\subsection{Quadratic Formula}
The quadratic formula is given by:
\[x = \frac{-b \pm \sqrt{b^2 - 4ac}}{2a}\]

\subsection{Matrix Example}
A sample matrix:
\[
\begin{pmatrix}
a_{11} & a_{12} & a_{13} \\
a_{21} & a_{22} & a_{23} \\
a_{31} & a_{32} & a_{33}
\end{pmatrix}
\]

\section{Tables}
Below is a sample table showing data:

\begin{table}[h]
\centering
\begin{tabular}{@{}lrrr@{}}
\toprule
\textbf{Category} & \textbf{Value 1} & \textbf{Value 2} & \textbf{Total} \\
\midrule
Item A & 10.5 & 20.3 & 30.8 \\
Item B & 15.2 & 25.1 & 40.3 \\
Item C & 22.1 & 12.7 & 34.8 \\
\bottomrule
\end{tabular}
\caption{Sample Data Table}
\label{tab:sample}
\end{table}

\section{Lists}
\subsection{Bullet Points}
\begin{itemize}
  \item First item
  \item Second item
  \item Third item with sub-items:
    \begin{itemize}
      \item Sub-item 1
      \item Sub-item 2
    \end{itemize}
\end{itemize}

\subsection{Numbered List}
\begin{enumerate}
  \item First step
  \item Second step
  \item Third step
\end{enumerate}

\section{Citations and References}
When you use the citation search feature, references will be automatically managed. For example, you might cite a paper like this: \cite{example2024}. All citations will be automatically added to your \texttt{cite.bib} file.

\section{Conclusion}
This template demonstrates basic \LaTeX{} features including mathematical formulas, tables, and lists. The bibliography system is now set up to work seamlessly with Bibby AI's citation management. Feel free to modify it for your needs!

\bibliographystyle{plain}
\bibliography{cite}

\end{document}\documentclass{article}
\usepackage[utf8]{inputenc}
\usepackage{booktabs}
\usepackage{amsmath}
\usepackage{natbib}

\title{My Document}
\author{Your Name}
\date{\today}

\begin{document}

\maketitle

\section{Introduction}
This is a sample document demonstrating various \LaTeX{} features including mathematical equations, tables, and formatting. You can cite papers using the citation search feature, and they will be automatically added to your \texttt{cite.bib} file.

\section{Mathematical Expressions}
Here are some example equations:

\subsection{Quadratic Formula}
The quadratic formula is given by:
\[x = \frac{-b \pm \sqrt{b^2 - 4ac}}{2a}\]

\subsection{Matrix Example}
A sample matrix:
\[
\begin{pmatrix}
a_{11} & a_{12} & a_{13} \\
a_{21} & a_{22} & a_{23} \\
a_{31} & a_{32} & a_{33}
\end{pmatrix}
\]

\section{Tables}
Below is a sample table showing data:

\begin{table}[h]
\centering
\begin{tabular}{@{}lrrr@{}}
\toprule
\textbf{Category} & \textbf{Value 1} & \textbf{Value 2} & \textbf{Total} \\
\midrule
Item A & 10.5 & 20.3 & 30.8 \\
Item B & 15.2 & 25.1 & 40.3 \\
Item C & 22.1 & 12.7 & 34.8 \\
\bottomrule
\end{tabular}
\caption{Sample Data Table}
\label{tab:sample}
\end{table}

\section{Lists}
\subsection{Bullet Points}
\begin{itemize}
  \item First item
  \item Second item
  \item Third item with sub-items:
    \begin{itemize}
      \item Sub-item 1
      \item Sub-item 2
    \end{itemize}
\end{itemize}

\subsection{Numbered List}
\begin{enumerate}
  \item First step
  \item Second step
  \item Third step
\end{enumerate}

\section{Citations and References}
When you use the citation search feature, references will be automatically managed. For example, you might cite a paper like this: \cite{example2024}. All citations will be automatically added to your \texttt{cite.bib} file.

\section{Conclusion}
This template demonstrates basic \LaTeX{} features including mathematical formulas, tables, and lists. The bibliography system is now set up to work seamlessly with Bibby AI's citation management. Feel free to modify it for your needs!

\bibliographystyle{plain}
\bibliography{cite}

\end{document}\documentclass{article}
\usepackage[utf8]{inputenc}
\usepackage{booktabs}
\usepackage{amsmath}
\usepackage{natbib}

\title{My Document}
\author{Your Name}
\date{\today}

\begin{document}

\maketitle

\section{Introduction}
This is a sample document demonstrating various \LaTeX{} features including mathematical equations, tables, and formatting. You can cite papers using the citation search feature, and they will be automatically added to your \texttt{cite.bib} file.

\section{Mathematical Expressions}
Here are some example equations:

\subsection{Quadratic Formula}
The quadratic formula is given by:
\[x = \frac{-b \pm \sqrt{b^2 - 4ac}}{2a}\]

\subsection{Matrix Example}
A sample matrix:
\[
\begin{pmatrix}
a_{11} & a_{12} & a_{13} \\
a_{21} & a_{22} & a_{23} \\
a_{31} & a_{32} & a_{33}
\end{pmatrix}
\]

\section{Tables}
Below is a sample table showing data:

\begin{table}[h]
\centering
\begin{tabular}{@{}lrrr@{}}
\toprule
\textbf{Category} & \textbf{Value 1} & \textbf{Value 2} & \textbf{Total} \\
\midrule
Item A & 10.5 & 20.3 & 30.8 \\
Item B & 15.2 & 25.1 & 40.3 \\
Item C & 22.1 & 12.7 & 34.8 \\
\bottomrule
\end{tabular}
\caption{Sample Data Table}
\label{tab:sample}
\end{table}

\section{Lists}
\subsection{Bullet Points}
\begin{itemize}
  \item First item
  \item Second item
  \item Third item with sub-items:
    \begin{itemize}
      \item Sub-item 1
      \item Sub-item 2
    \end{itemize}
\end{itemize}

\subsection{Numbered List}
\begin{enumerate}
  \item First step
  \item Second step
  \item Third step
\end{enumerate}

\section{Citations and References}
When you use the citation search feature, references will be automatically managed. For example, you might cite a paper like this: \cite{example2024}. All citations will be automatically added to your \texttt{cite.bib} file.

\section{Conclusion}
This template demonstrates basic \LaTeX{} features including mathematical formulas, tables, and lists. The bibliography system is now set up to work seamlessly with Bibby AI's citation management. Feel free to modify it for your needs!

\bibliographystyle{plain}
\bibliography{cite}

\end{document}\documentclass{article}
\usepackage[utf8]{inputenc}
\usepackage{booktabs}
\usepackage{amsmath}
\usepackage{natbib}

\title{My Document}
\author{Your Name}
\date{\today}

\begin{document}

\maketitle

\section{Introduction}
This is a sample document demonstrating various \LaTeX{} features including mathematical equations, tables, and formatting. You can cite papers using the citation search feature, and they will be automatically added to your \texttt{cite.bib} file.

\section{Mathematical Expressions}
Here are some example equations:

\subsection{Quadratic Formula}
The quadratic formula is given by:
\[x = \frac{-b \pm \sqrt{b^2 - 4ac}}{2a}\]

\subsection{Matrix Example}
A sample matrix:
\[
\begin{pmatrix}
a_{11} & a_{12} & a_{13} \\
a_{21} & a_{22} & a_{23} \\
a_{31} & a_{32} & a_{33}
\end{pmatrix}
\]

\section{Tables}
Below is a sample table showing data:

\begin{table}[h]
\centering
\begin{tabular}{@{}lrrr@{}}
\toprule
\textbf{Category} & \textbf{Value 1} & \textbf{Value 2} & \textbf{Total} \\
\midrule
Item A & 10.5 & 20.3 & 30.8 \\
Item B & 15.2 & 25.1 & 40.3 \\
Item C & 22.1 & 12.7 & 34.8 \\
\bottomrule
\end{tabular}
\caption{Sample Data Table}
\label{tab:sample}
\end{table}

\section{Lists}
\subsection{Bullet Points}
\begin{itemize}
  \item First item
  \item Second item
  \item Third item with sub-items:
    \begin{itemize}
      \item Sub-item 1
      \item Sub-item 2
    \end{itemize}
\end{itemize}

\subsection{Numbered List}
\begin{enumerate}
  \item First step
  \item Second step
  \item Third step
\end{enumerate}

\section{Citations and References}
When you use the citation search feature, references will be automatically managed. For example, you might cite a paper like this: \cite{example2024}. All citations will be automatically added to your \texttt{cite.bib} file.

\section{Conclusion}
This template demonstrates basic \LaTeX{} features including mathematical formulas, tables, and lists. The bibliography system is now set up to work seamlessly with Bibby AI's citation management. Feel free to modify it for your needs!

\bibliographystyle{plain}
\bibliography{cite}

\end{document}\documentclass{article}
\usepackage[utf8]{inputenc}
\usepackage{booktabs}
\usepackage{amsmath}
\usepackage{natbib}

\title{My Document}
\author{Your Name}
\date{\today}

\begin{document}

\maketitle

\section{Introduction}
This is a sample document demonstrating various \LaTeX{} features including mathematical equations, tables, and formatting. You can cite papers using the citation search feature, and they will be automatically added to your \texttt{cite.bib} file.

\section{Mathematical Expressions}
Here are some example equations:

\subsection{Quadratic Formula}
The quadratic formula is given by:
\[x = \frac{-b \pm \sqrt{b^2 - 4ac}}{2a}\]

\subsection{Matrix Example}
A sample matrix:
\[
\begin{pmatrix}
a_{11} & a_{12} & a_{13} \\
a_{21} & a_{22} & a_{23} \\
a_{31} & a_{32} & a_{33}
\end{pmatrix}
\]

\section{Tables}
Below is a sample table showing data:

\begin{table}[h]
\centering
\begin{tabular}{@{}lrrr@{}}
\toprule
\textbf{Category} & \textbf{Value 1} & \textbf{Value 2} & \textbf{Total} \\
\midrule
Item A & 10.5 & 20.3 & 30.8 \\
Item B & 15.2 & 25.1 & 40.3 \\
Item C & 22.1 & 12.7 & 34.8 \\
\bottomrule
\end{tabular}
\caption{Sample Data Table}
\label{tab:sample}
\end{table}

\section{Lists}
\subsection{Bullet Points}
\begin{itemize}
  \item First item
  \item Second item
  \item Third item with sub-items:
    \begin{itemize}
      \item Sub-item 1
      \item Sub-item 2
    \end{itemize}
\end{itemize}

\subsection{Numbered List}
\begin{enumerate}
  \item First step
  \item Second step
  \item Third step
\end{enumerate}

\section{Citations and References}
When you use the citation search feature, references will be automatically managed. For example, you might cite a paper like this: \cite{example2024}. All citations will be automatically added to your \texttt{cite.bib} file.

\section{Conclusion}
This template demonstrates basic \LaTeX{} features including mathematical formulas, tables, and lists. The bibliography system is now set up to work seamlessly with Bibby AI's citation management. Feel free to modify it for your needs!

\bibliographystyle{plain}
\bibliography{cite}

\end{document}\documentclass{article}
\usepackage[utf8]{inputenc}
\usepackage{booktabs}
\usepackage{amsmath}
\usepackage{natbib}

\title{My Document}
\author{Your Name}
\date{\today}

\begin{document}

\maketitle

\section{Introduction}
This is a sample document demonstrating various \LaTeX{} features including mathematical equations, tables, and formatting. You can cite papers using the citation search feature, and they will be automatically added to your \texttt{cite.bib} file.

\section{Mathematical Expressions}
Here are some example equations:

\subsection{Quadratic Formula}
The quadratic formula is given by:
\[x = \frac{-b \pm \sqrt{b^2 - 4ac}}{2a}\]

\subsection{Matrix Example}
A sample matrix:
\[
\begin{pmatrix}
a_{11} & a_{12} & a_{13} \\
a_{21} & a_{22} & a_{23} \\
a_{31} & a_{32} & a_{33}
\end{pmatrix}
\]

\section{Tables}
Below is a sample table showing data:

\begin{table}[h]
\centering
\begin{tabular}{@{}lrrr@{}}
\toprule
\textbf{Category} & \textbf{Value 1} & \textbf{Value 2} & \textbf{Total} \\
\midrule
Item A & 10.5 & 20.3 & 30.8 \\
Item B & 15.2 & 25.1 & 40.3 \\
Item C & 22.1 & 12.7 & 34.8 \\
\bottomrule
\end{tabular}
\caption{Sample Data Table}
\label{tab:sample}
\end{table}

\section{Lists}
\subsection{Bullet Points}
\begin{itemize}
  \item First item
  \item Second item
  \item Third item with sub-items:
    \begin{itemize}
      \item Sub-item 1
      \item Sub-item 2
    \end{itemize}
\end{itemize}

\subsection{Numbered List}
\begin{enumerate}
  \item First step
  \item Second step
  \item Third step
\end{enumerate}

\section{Citations and References}
When you use the citation search feature, references will be automatically managed. For example, you might cite a paper like this: \cite{example2024}. All citations will be automatically added to your \texttt{cite.bib} file.

\section{Conclusion}
This template demonstrates basic \LaTeX{} features including mathematical formulas, tables, and lists. The bibliography system is now set up to work seamlessly with Bibby AI's citation management. Feel free to modify it for your needs!

\bibliographystyle{plain}
\bibliography{cite}

\end{document}\documentclass{article}
\usepackage[utf8]{inputenc}
\usepackage{booktabs}
\usepackage{amsmath}
\usepackage{natbib}

\title{My Document}
\author{Your Name}
\date{\today}

\begin{document}

\maketitle

\section{Introduction}
This is a sample document demonstrating various \LaTeX{} features including mathematical equations, tables, and formatting. You can cite papers using the citation search feature, and they will be automatically added to your \texttt{cite.bib} file.

\section{Mathematical Expressions}
Here are some example equations:

\subsection{Quadratic Formula}
The quadratic formula is given by:
\[x = \frac{-b \pm \sqrt{b^2 - 4ac}}{2a}\]

\subsection{Matrix Example}
A sample matrix:
\[
\begin{pmatrix}
a_{11} & a_{12} & a_{13} \\
a_{21} & a_{22} & a_{23} \\
a_{31} & a_{32} & a_{33}
\end{pmatrix}
\]

\section{Tables}
Below is a sample table showing data:

\begin{table}[h]
\centering
\begin{tabular}{@{}lrrr@{}}
\toprule
\textbf{Category} & \textbf{Value 1} & \textbf{Value 2} & \textbf{Total} \\
\midrule
Item A & 10.5 & 20.3 & 30.8 \\
Item B & 15.2 & 25.1 & 40.3 \\
Item C & 22.1 & 12.7 & 34.8 \\
\bottomrule
\end{tabular}
\caption{Sample Data Table}
\label{tab:sample}
\end{table}

\section{Lists}
\subsection{Bullet Points}
\begin{itemize}
  \item First item
  \item Second item
  \item Third item with sub-items:
    \begin{itemize}
      \item Sub-item 1
      \item Sub-item 2
    \end{itemize}
\end{itemize}

\subsection{Numbered List}
\begin{enumerate}
  \item First step
  \item Second step
  \item Third step
\end{enumerate}

\section{Citations and References}
When you use the citation search feature, references will be automatically managed. For example, you might cite a paper like this: \cite{example2024}. All citations will be automatically added to your \texttt{cite.bib} file.

\section{Conclusion}
This template demonstrates basic \LaTeX{} features including mathematical formulas, tables, and lists. The bibliography system is now set up to work seamlessly with Bibby AI's citation management. Feel free to modify it for your needs!

\bibliographystyle{plain}
\bibliography{cite}

\end{document}\documentclass{article}
\usepackage[utf8]{inputenc}
\usepackage{booktabs}
\usepackage{amsmath}
\usepackage{natbib}

\title{My Document}
\author{Your Name}
\date{\today}

\begin{document}

\maketitle

\section{Introduction}
This is a sample document demonstrating various \LaTeX{} features including mathematical equations, tables, and formatting. You can cite papers using the citation search feature, and they will be automatically added to your \texttt{cite.bib} file.

\section{Mathematical Expressions}
Here are some example equations:

\subsection{Quadratic Formula}
The quadratic formula is given by:
\[x = \frac{-b \pm \sqrt{b^2 - 4ac}}{2a}\]

\subsection{Matrix Example}
A sample matrix:
\[
\begin{pmatrix}
a_{11} & a_{12} & a_{13} \\
a_{21} & a_{22} & a_{23} \\
a_{31} & a_{32} & a_{33}
\end{pmatrix}
\]

\section{Tables}
Below is a sample table showing data:

\begin{table}[h]
\centering
\begin{tabular}{@{}lrrr@{}}
\toprule
\textbf{Category} & \textbf{Value 1} & \textbf{Value 2} & \textbf{Total} \\
\midrule
Item A & 10.5 & 20.3 & 30.8 \\
Item B & 15.2 & 25.1 & 40.3 \\
Item C & 22.1 & 12.7 & 34.8 \\
\bottomrule
\end{tabular}
\caption{Sample Data Table}
\label{tab:sample}
\end{table}

\section{Lists}
\subsection{Bullet Points}
\begin{itemize}
  \item First item
  \item Second item
  \item Third item with sub-items:
    \begin{itemize}
      \item Sub-item 1
      \item Sub-item 2
    \end{itemize}
\end{itemize}

\subsection{Numbered List}
\begin{enumerate}
  \item First step
  \item Second step
  \item Third step
\end{enumerate}

\section{Citations and References}
When you use the citation search feature, references will be automatically managed. For example, you might cite a paper like this: \cite{example2024}. All citations will be automatically added to your \texttt{cite.bib} file.

\section{Conclusion}
This template demonstrates basic \LaTeX{} features including mathematical formulas, tables, and lists. The bibliography system is now set up to work seamlessly with Bibby AI's citation management. Feel free to modify it for your needs!

\bibliographystyle{plain}
\bibliography{cite}

\end{document}\documentclass{article}
\usepackage[utf8]{inputenc}
\usepackage{booktabs}
\usepackage{amsmath}
\usepackage{natbib}

\title{My Document}
\author{Your Name}
\date{\today}

\begin{document}

\maketitle

\section{Introduction}
This is a sample document demonstrating various \LaTeX{} features including mathematical equations, tables, and formatting. You can cite papers using the citation search feature, and they will be automatically added to your \texttt{cite.bib} file.

\section{Mathematical Expressions}
Here are some example equations:

\subsection{Quadratic Formula}
The quadratic formula is given by:
\[x = \frac{-b \pm \sqrt{b^2 - 4ac}}{2a}\]

\subsection{Matrix Example}
A sample matrix:
\[
\begin{pmatrix}
a_{11} & a_{12} & a_{13} \\
a_{21} & a_{22} & a_{23} \\
a_{31} & a_{32} & a_{33}
\end{pmatrix}
\]

\section{Tables}
Below is a sample table showing data:

\begin{table}[h]
\centering
\begin{tabular}{@{}lrrr@{}}
\toprule
\textbf{Category} & \textbf{Value 1} & \textbf{Value 2} & \textbf{Total} \\
\midrule
Item A & 10.5 & 20.3 & 30.8 \\
Item B & 15.2 & 25.1 & 40.3 \\
Item C & 22.1 & 12.7 & 34.8 \\
\bottomrule
\end{tabular}
\caption{Sample Data Table}
\label{tab:sample}
\end{table}

\section{Lists}
\subsection{Bullet Points}
\begin{itemize}
  \item First item
  \item Second item
  \item Third item with sub-items:
    \begin{itemize}
      \item Sub-item 1
      \item Sub-item 2
    \end{itemize}
\end{itemize}

\subsection{Numbered List}
\begin{enumerate}
  \item First step
  \item Second step
  \item Third step
\end{enumerate}

\section{Citations and References}
When you use the citation search feature, references will be automatically managed. For example, you might cite a paper like this: \cite{example2024}. All citations will be automatically added to your \texttt{cite.bib} file.

\section{Conclusion}
This template demonstrates basic \LaTeX{} features including mathematical formulas, tables, and lists. The bibliography system is now set up to work seamlessly with Bibby AI's citation management. Feel free to modify it for your needs!

\bibliographystyle{plain}
\bibliography{cite}

\end{document}\documentclass{article}
\usepackage[utf8]{inputenc}
\usepackage{booktabs}
\usepackage{amsmath}
\usepackage{natbib}

\title{My Document}
\author{Your Name}
\date{\today}

\begin{document}

\maketitle

\section{Introduction}
This is a sample document demonstrating various \LaTeX{} features including mathematical equations, tables, and formatting. You can cite papers using the citation search feature, and they will be automatically added to your \texttt{cite.bib} file.

\section{Mathematical Expressions}
Here are some example equations:

\subsection{Quadratic Formula}
The quadratic formula is given by:
\[x = \frac{-b \pm \sqrt{b^2 - 4ac}}{2a}\]

\subsection{Matrix Example}
A sample matrix:
\[
\begin{pmatrix}
a_{11} & a_{12} & a_{13} \\
a_{21} & a_{22} & a_{23} \\
a_{31} & a_{32} & a_{33}
\end{pmatrix}
\]

\section{Tables}
Below is a sample table showing data:

\begin{table}[h]
\centering
\begin{tabular}{@{}lrrr@{}}
\toprule
\textbf{Category} & \textbf{Value 1} & \textbf{Value 2} & \textbf{Total} \\
\midrule
Item A & 10.5 & 20.3 & 30.8 \\
Item B & 15.2 & 25.1 & 40.3 \\
Item C & 22.1 & 12.7 & 34.8 \\
\bottomrule
\end{tabular}
\caption{Sample Data Table}
\label{tab:sample}
\end{table}

\section{Lists}
\subsection{Bullet Points}
\begin{itemize}
  \item First item
  \item Second item
  \item Third item with sub-items:
    \begin{itemize}
      \item Sub-item 1
      \item Sub-item 2
    \end{itemize}
\end{itemize}

\subsection{Numbered List}
\begin{enumerate}
  \item First step
  \item Second step
  \item Third step
\end{enumerate}

\section{Citations and References}
When you use the citation search feature, references will be automatically managed. For example, you might cite a paper like this: \cite{example2024}. All citations will be automatically added to your \texttt{cite.bib} file.

\section{Conclusion}
This template demonstrates basic \LaTeX{} features including mathematical formulas, tables, and lists. The bibliography system is now set up to work seamlessly with Bibby AI's citation management. Feel free to modify it for your needs!

\bibliographystyle{plain}
\bibliography{cite}

\end{document}\documentclass{article}
\usepackage[utf8]{inputenc}
\usepackage{booktabs}
\usepackage{amsmath}
\usepackage{natbib}

\title{My Document}
\author{Your Name}
\date{\today}

\begin{document}

\maketitle

\section{Introduction}
This is a sample document demonstrating various \LaTeX{} features including mathematical equations, tables, and formatting. You can cite papers using the citation search feature, and they will be automatically added to your \texttt{cite.bib} file.

\section{Mathematical Expressions}
Here are some example equations:

\subsection{Quadratic Formula}
The quadratic formula is given by:
\[x = \frac{-b \pm \sqrt{b^2 - 4ac}}{2a}\]

\subsection{Matrix Example}
A sample matrix:
\[
\begin{pmatrix}
a_{11} & a_{12} & a_{13} \\
a_{21} & a_{22} & a_{23} \\
a_{31} & a_{32} & a_{33}
\end{pmatrix}
\]

\section{Tables}
Below is a sample table showing data:

\begin{table}[h]
\centering
\begin{tabular}{@{}lrrr@{}}
\toprule
\textbf{Category} & \textbf{Value 1} & \textbf{Value 2} & \textbf{Total} \\
\midrule
Item A & 10.5 & 20.3 & 30.8 \\
Item B & 15.2 & 25.1 & 40.3 \\
Item C & 22.1 & 12.7 & 34.8 \\
\bottomrule
\end{tabular}
\caption{Sample Data Table}
\label{tab:sample}
\end{table}

\section{Lists}
\subsection{Bullet Points}
\begin{itemize}
  \item First item
  \item Second item
  \item Third item with sub-items:
    \begin{itemize}
      \item Sub-item 1
      \item Sub-item 2
    \end{itemize}
\end{itemize}

\subsection{Numbered List}
\begin{enumerate}
  \item First step
  \item Second step
  \item Third step
\end{enumerate}

\section{Citations and References}
When you use the citation search feature, references will be automatically managed. For example, you might cite a paper like this: \cite{example2024}. All citations will be automatically added to your \texttt{cite.bib} file.

\section{Conclusion}
This template demonstrates basic \LaTeX{} features including mathematical formulas, tables, and lists. The bibliography system is now set up to work seamlessly with Bibby AI's citation management. Feel free to modify it for your needs!

\bibliographystyle{plain}
\bibliography{cite}

\end{document}\documentclass{article}
\usepackage[utf8]{inputenc}
\usepackage{booktabs}
\usepackage{amsmath}
\usepackage{natbib}

\title{My Document}
\author{Your Name}
\date{\today}

\begin{document}

\maketitle

\section{Introduction}
This is a sample document demonstrating various \LaTeX{} features including mathematical equations, tables, and formatting. You can cite papers using the citation search feature, and they will be automatically added to your \texttt{cite.bib} file.

\section{Mathematical Expressions}
Here are some example equations:

\subsection{Quadratic Formula}
The quadratic formula is given by:
\[x = \frac{-b \pm \sqrt{b^2 - 4ac}}{2a}\]

\subsection{Matrix Example}
A sample matrix:
\[
\begin{pmatrix}
a_{11} & a_{12} & a_{13} \\
a_{21} & a_{22} & a_{23} \\
a_{31} & a_{32} & a_{33}
\end{pmatrix}
\]

\section{Tables}
Below is a sample table showing data:

\begin{table}[h]
\centering
\begin{tabular}{@{}lrrr@{}}
\toprule
\textbf{Category} & \textbf{Value 1} & \textbf{Value 2} & \textbf{Total} \\
\midrule
Item A & 10.5 & 20.3 & 30.8 \\
Item B & 15.2 & 25.1 & 40.3 \\
Item C & 22.1 & 12.7 & 34.8 \\
\bottomrule
\end{tabular}
\caption{Sample Data Table}
\label{tab:sample}
\end{table}

\section{Lists}
\subsection{Bullet Points}
\begin{itemize}
  \item First item
  \item Second item
  \item Third item with sub-items:
    \begin{itemize}
      \item Sub-item 1
      \item Sub-item 2
    \end{itemize}
\end{itemize}

\subsection{Numbered List}
\begin{enumerate}
  \item First step
  \item Second step
  \item Third step
\end{enumerate}

\section{Citations and References}
When you use the citation search feature, references will be automatically managed. For example, you might cite a paper like this: \cite{example2024}. All citations will be automatically added to your \texttt{cite.bib} file.

\section{Conclusion}
This template demonstrates basic \LaTeX{} features including mathematical formulas, tables, and lists. The bibliography system is now set up to work seamlessly with Bibby AI's citation management. Feel free to modify it for your needs!

\bibliographystyle{plain}
\bibliography{cite}

\end{document}\documentclass{article}
\usepackage[utf8]{inputenc}
\usepackage{booktabs}
\usepackage{amsmath}
\usepackage{natbib}

\title{My Document}
\author{Your Name}
\date{\today}

\begin{document}

\maketitle

\section{Introduction}
This is a sample document demonstrating various \LaTeX{} features including mathematical equations, tables, and formatting. You can cite papers using the citation search feature, and they will be automatically added to your \texttt{cite.bib} file.

\section{Mathematical Expressions}
Here are some example equations:

\subsection{Quadratic Formula}
The quadratic formula is given by:
\[x = \frac{-b \pm \sqrt{b^2 - 4ac}}{2a}\]

\subsection{Matrix Example}
A sample matrix:
\[
\begin{pmatrix}
a_{11} & a_{12} & a_{13} \\
a_{21} & a_{22} & a_{23} \\
a_{31} & a_{32} & a_{33}
\end{pmatrix}
\]

\section{Tables}
Below is a sample table showing data:

\begin{table}[h]
\centering
\begin{tabular}{@{}lrrr@{}}
\toprule
\textbf{Category} & \textbf{Value 1} & \textbf{Value 2} & \textbf{Total} \\
\midrule
Item A & 10.5 & 20.3 & 30.8 \\
Item B & 15.2 & 25.1 & 40.3 \\
Item C & 22.1 & 12.7 & 34.8 \\
\bottomrule
\end{tabular}
\caption{Sample Data Table}
\label{tab:sample}
\end{table}

\section{Lists}
\subsection{Bullet Points}
\begin{itemize}
  \item First item
  \item Second item
  \item Third item with sub-items:
    \begin{itemize}
      \item Sub-item 1
      \item Sub-item 2
    \end{itemize}
\end{itemize}

\subsection{Numbered List}
\begin{enumerate}
  \item First step
  \item Second step
  \item Third step
\end{enumerate}

\section{Citations and References}
When you use the citation search feature, references will be automatically managed. For example, you might cite a paper like this: \cite{example2024}. All citations will be automatically added to your \texttt{cite.bib} file.

\section{Conclusion}
This template demonstrates basic \LaTeX{} features including mathematical formulas, tables, and lists. The bibliography system is now set up to work seamlessly with Bibby AI's citation management. Feel free to modify it for your needs!

\bibliographystyle{plain}
\bibliography{cite}

\end{document}\documentclass{article}
\usepackage[utf8]{inputenc}
\usepackage{booktabs}
\usepackage{amsmath}
\usepackage{natbib}

\title{My Document}
\author{Your Name}
\date{\today}

\begin{document}

\maketitle

\section{Introduction}
This is a sample document demonstrating various \LaTeX{} features including mathematical equations, tables, and formatting. You can cite papers using the citation search feature, and they will be automatically added to your \texttt{cite.bib} file.

\section{Mathematical Expressions}
Here are some example equations:

\subsection{Quadratic Formula}
The quadratic formula is given by:
\[x = \frac{-b \pm \sqrt{b^2 - 4ac}}{2a}\]

\subsection{Matrix Example}
A sample matrix:
\[
\begin{pmatrix}
a_{11} & a_{12} & a_{13} \\
a_{21} & a_{22} & a_{23} \\
a_{31} & a_{32} & a_{33}
\end{pmatrix}
\]

\section{Tables}
Below is a sample table showing data:

\begin{table}[h]
\centering
\begin{tabular}{@{}lrrr@{}}
\toprule
\textbf{Category} & \textbf{Value 1} & \textbf{Value 2} & \textbf{Total} \\
\midrule
Item A & 10.5 & 20.3 & 30.8 \\
Item B & 15.2 & 25.1 & 40.3 \\
Item C & 22.1 & 12.7 & 34.8 \\
\bottomrule
\end{tabular}
\caption{Sample Data Table}
\label{tab:sample}
\end{table}

\section{Lists}
\subsection{Bullet Points}
\begin{itemize}
  \item First item
  \item Second item
  \item Third item with sub-items:
    \begin{itemize}
      \item Sub-item 1
      \item Sub-item 2
    \end{itemize}
\end{itemize}

\subsection{Numbered List}
\begin{enumerate}
  \item First step
  \item Second step
  \item Third step
\end{enumerate}

\section{Citations and References}
When you use the citation search feature, references will be automatically managed. For example, you might cite a paper like this: \cite{example2024}. All citations will be automatically added to your \texttt{cite.bib} file.

\section{Conclusion}
This template demonstrates basic \LaTeX{} features including mathematical formulas, tables, and lists. The bibliography system is now set up to work seamlessly with Bibby AI's citation management. Feel free to modify it for your needs!

\bibliographystyle{plain}
\bibliography{cite}

\end{document}\documentclass{article}
\usepackage[utf8]{inputenc}
\usepackage{booktabs}
\usepackage{amsmath}
\usepackage{natbib}

\title{My Document}
\author{Your Name}
\date{\today}

\begin{document}

\maketitle

\section{Introduction}
This is a sample document demonstrating various \LaTeX{} features including mathematical equations, tables, and formatting. You can cite papers using the citation search feature, and they will be automatically added to your \texttt{cite.bib} file.

\section{Mathematical Expressions}
Here are some example equations:

\subsection{Quadratic Formula}
The quadratic formula is given by:
\[x = \frac{-b \pm \sqrt{b^2 - 4ac}}{2a}\]

\subsection{Matrix Example}
A sample matrix:
\[
\begin{pmatrix}
a_{11} & a_{12} & a_{13} \\
a_{21} & a_{22} & a_{23} \\
a_{31} & a_{32} & a_{33}
\end{pmatrix}
\]

\section{Tables}
Below is a sample table showing data:

\begin{table}[h]
\centering
\begin{tabular}{@{}lrrr@{}}
\toprule
\textbf{Category} & \textbf{Value 1} & \textbf{Value 2} & \textbf{Total} \\
\midrule
Item A & 10.5 & 20.3 & 30.8 \\
Item B & 15.2 & 25.1 & 40.3 \\
Item C & 22.1 & 12.7 & 34.8 \\
\bottomrule
\end{tabular}
\caption{Sample Data Table}
\label{tab:sample}
\end{table}

\section{Lists}
\subsection{Bullet Points}
\begin{itemize}
  \item First item
  \item Second item
  \item Third item with sub-items:
    \begin{itemize}
      \item Sub-item 1
      \item Sub-item 2
    \end{itemize}
\end{itemize}

\subsection{Numbered List}
\begin{enumerate}
  \item First step
  \item Second step
  \item Third step
\end{enumerate}

\section{Citations and References}
When you use the citation search feature, references will be automatically managed. For example, you might cite a paper like this: \cite{example2024}. All citations will be automatically added to your \texttt{cite.bib} file.

\section{Conclusion}
This template demonstrates basic \LaTeX{} features including mathematical formulas, tables, and lists. The bibliography system is now set up to work seamlessly with Bibby AI's citation management. Feel free to modify it for your needs!

\bibliographystyle{plain}
\bibliography{cite}

\end{document}\documentclass{article}
\usepackage[utf8]{inputenc}
\usepackage{booktabs}
\usepackage{amsmath}
\usepackage{natbib}

\title{My Document}
\author{Your Name}
\date{\today}

\begin{document}

\maketitle

\section{Introduction}
This is a sample document demonstrating various \LaTeX{} features including mathematical equations, tables, and formatting. You can cite papers using the citation search feature, and they will be automatically added to your \texttt{cite.bib} file.

\section{Mathematical Expressions}
Here are some example equations:

\subsection{Quadratic Formula}
The quadratic formula is given by:
\[x = \frac{-b \pm \sqrt{b^2 - 4ac}}{2a}\]

\subsection{Matrix Example}
A sample matrix:
\[
\begin{pmatrix}
a_{11} & a_{12} & a_{13} \\
a_{21} & a_{22} & a_{23} \\
a_{31} & a_{32} & a_{33}
\end{pmatrix}
\]

\section{Tables}
Below is a sample table showing data:

\begin{table}[h]
\centering
\begin{tabular}{@{}lrrr@{}}
\toprule
\textbf{Category} & \textbf{Value 1} & \textbf{Value 2} & \textbf{Total} \\
\midrule
Item A & 10.5 & 20.3 & 30.8 \\
Item B & 15.2 & 25.1 & 40.3 \\
Item C & 22.1 & 12.7 & 34.8 \\
\bottomrule
\end{tabular}
\caption{Sample Data Table}
\label{tab:sample}
\end{table}

\section{Lists}
\subsection{Bullet Points}
\begin{itemize}
  \item First item
  \item Second item
  \item Third item with sub-items:
    \begin{itemize}
      \item Sub-item 1
      \item Sub-item 2
    \end{itemize}
\end{itemize}

\subsection{Numbered List}
\begin{enumerate}
  \item First step
  \item Second step
  \item Third step
\end{enumerate}

\section{Citations and References}
When you use the citation search feature, references will be automatically managed. For example, you might cite a paper like this: \cite{example2024}. All citations will be automatically added to your \texttt{cite.bib} file.

\section{Conclusion}
This template demonstrates basic \LaTeX{} features including mathematical formulas, tables, and lists. The bibliography system is now set up to work seamlessly with Bibby AI's citation management. Feel free to modify it for your needs!

\bibliographystyle{plain}
\bibliography{cite}

\end{document}\documentclass{article}
\usepackage[utf8]{inputenc}
\usepackage{booktabs}
\usepackage{amsmath}
\usepackage{natbib}

\title{My Document}
\author{Your Name}
\date{\today}

\begin{document}

\maketitle

\section{Introduction}
This is a sample document demonstrating various \LaTeX{} features including mathematical equations, tables, and formatting. You can cite papers using the citation search feature, and they will be automatically added to your \texttt{cite.bib} file.

\section{Mathematical Expressions}
Here are some example equations:

\subsection{Quadratic Formula}
The quadratic formula is given by:
\[x = \frac{-b \pm \sqrt{b^2 - 4ac}}{2a}\]

\subsection{Matrix Example}
A sample matrix:
\[
\begin{pmatrix}
a_{11} & a_{12} & a_{13} \\
a_{21} & a_{22} & a_{23} \\
a_{31} & a_{32} & a_{33}
\end{pmatrix}
\]

\section{Tables}
Below is a sample table showing data:

\begin{table}[h]
\centering
\begin{tabular}{@{}lrrr@{}}
\toprule
\textbf{Category} & \textbf{Value 1} & \textbf{Value 2} & \textbf{Total} \\
\midrule
Item A & 10.5 & 20.3 & 30.8 \\
Item B & 15.2 & 25.1 & 40.3 \\
Item C & 22.1 & 12.7 & 34.8 \\
\bottomrule
\end{tabular}
\caption{Sample Data Table}
\label{tab:sample}
\end{table}

\section{Lists}
\subsection{Bullet Points}
\begin{itemize}
  \item First item
  \item Second item
  \item Third item with sub-items:
    \begin{itemize}
      \item Sub-item 1
      \item Sub-item 2
    \end{itemize}
\end{itemize}

\subsection{Numbered List}
\begin{enumerate}
  \item First step
  \item Second step
  \item Third step
\end{enumerate}

\section{Citations and References}
When you use the citation search feature, references will be automatically managed. For example, you might cite a paper like this: \cite{example2024}. All citations will be automatically added to your \texttt{cite.bib} file.

\section{Conclusion}
This template demonstrates basic \LaTeX{} features including mathematical formulas, tables, and lists. The bibliography system is now set up to work seamlessly with Bibby AI's citation management. Feel free to modify it for your needs!

\bibliographystyle{plain}
\bibliography{cite}

\end{document}\documentclass{article}
\usepackage[utf8]{inputenc}
\usepackage{booktabs}
\usepackage{amsmath}
\usepackage{natbib}

\title{My Document}
\author{Your Name}
\date{\today}

\begin{document}

\maketitle

\section{Introduction}
This is a sample document demonstrating various \LaTeX{} features including mathematical equations, tables, and formatting. You can cite papers using the citation search feature, and they will be automatically added to your \texttt{cite.bib} file.

\section{Mathematical Expressions}
Here are some example equations:

\subsection{Quadratic Formula}
The quadratic formula is given by:
\[x = \frac{-b \pm \sqrt{b^2 - 4ac}}{2a}\]

\subsection{Matrix Example}
A sample matrix:
\[
\begin{pmatrix}
a_{11} & a_{12} & a_{13} \\
a_{21} & a_{22} & a_{23} \\
a_{31} & a_{32} & a_{33}
\end{pmatrix}
\]

\section{Tables}
Below is a sample table showing data:

\begin{table}[h]
\centering
\begin{tabular}{@{}lrrr@{}}
\toprule
\textbf{Category} & \textbf{Value 1} & \textbf{Value 2} & \textbf{Total} \\
\midrule
Item A & 10.5 & 20.3 & 30.8 \\
Item B & 15.2 & 25.1 & 40.3 \\
Item C & 22.1 & 12.7 & 34.8 \\
\bottomrule
\end{tabular}
\caption{Sample Data Table}
\label{tab:sample}
\end{table}

\section{Lists}
\subsection{Bullet Points}
\begin{itemize}
  \item First item
  \item Second item
  \item Third item with sub-items:
    \begin{itemize}
      \item Sub-item 1
      \item Sub-item 2
    \end{itemize}
\end{itemize}

\subsection{Numbered List}
\begin{enumerate}
  \item First step
  \item Second step
  \item Third step
\end{enumerate}

\section{Citations and References}
When you use the citation search feature, references will be automatically managed. For example, you might cite a paper like this: \cite{example2024}. All citations will be automatically added to your \texttt{cite.bib} file.

\section{Conclusion}
This template demonstrates basic \LaTeX{} features including mathematical formulas, tables, and lists. The bibliography system is now set up to work seamlessly with Bibby AI's citation management. Feel free to modify it for your needs!

\bibliographystyle{plain}
\bibliography{cite}

\end{document}\documentclass{article}
\usepackage[utf8]{inputenc}
\usepackage{booktabs}
\usepackage{amsmath}
\usepackage{natbib}

\title{My Document}
\author{Your Name}
\date{\today}

\begin{document}

\maketitle

\section{Introduction}
This is a sample document demonstrating various \LaTeX{} features including mathematical equations, tables, and formatting. You can cite papers using the citation search feature, and they will be automatically added to your \texttt{cite.bib} file.

\section{Mathematical Expressions}
Here are some example equations:

\subsection{Quadratic Formula}
The quadratic formula is given by:
\[x = \frac{-b \pm \sqrt{b^2 - 4ac}}{2a}\]

\subsection{Matrix Example}
A sample matrix:
\[
\begin{pmatrix}
a_{11} & a_{12} & a_{13} \\
a_{21} & a_{22} & a_{23} \\
a_{31} & a_{32} & a_{33}
\end{pmatrix}
\]

\section{Tables}
Below is a sample table showing data:

\begin{table}[h]
\centering
\begin{tabular}{@{}lrrr@{}}
\toprule
\textbf{Category} & \textbf{Value 1} & \textbf{Value 2} & \textbf{Total} \\
\midrule
Item A & 10.5 & 20.3 & 30.8 \\
Item B & 15.2 & 25.1 & 40.3 \\
Item C & 22.1 & 12.7 & 34.8 \\
\bottomrule
\end{tabular}
\caption{Sample Data Table}
\label{tab:sample}
\end{table}

\section{Lists}
\subsection{Bullet Points}
\begin{itemize}
  \item First item
  \item Second item
  \item Third item with sub-items:
    \begin{itemize}
      \item Sub-item 1
      \item Sub-item 2
    \end{itemize}
\end{itemize}

\subsection{Numbered List}
\begin{enumerate}
  \item First step
  \item Second step
  \item Third step
\end{enumerate}

\section{Citations and References}
When you use the citation search feature, references will be automatically managed. For example, you might cite a paper like this: \cite{example2024}. All citations will be automatically added to your \texttt{cite.bib} file.

\section{Conclusion}
This template demonstrates basic \LaTeX{} features including mathematical formulas, tables, and lists. The bibliography system is now set up to work seamlessly with Bibby AI's citation management. Feel free to modify it for your needs!

\bibliographystyle{plain}
\bibliography{cite}

\end{document}\documentclass{article}
\usepackage[utf8]{inputenc}
\usepackage{booktabs}
\usepackage{amsmath}
\usepackage{natbib}

\title{My Document}
\author{Your Name}
\date{\today}

\begin{document}

\maketitle

\section{Introduction}
This is a sample document demonstrating various \LaTeX{} features including mathematical equations, tables, and formatting. You can cite papers using the citation search feature, and they will be automatically added to your \texttt{cite.bib} file.

\section{Mathematical Expressions}
Here are some example equations:

\subsection{Quadratic Formula}
The quadratic formula is given by:
\[x = \frac{-b \pm \sqrt{b^2 - 4ac}}{2a}\]

\subsection{Matrix Example}
A sample matrix:
\[
\begin{pmatrix}
a_{11} & a_{12} & a_{13} \\
a_{21} & a_{22} & a_{23} \\
a_{31} & a_{32} & a_{33}
\end{pmatrix}
\]

\section{Tables}
Below is a sample table showing data:

\begin{table}[h]
\centering
\begin{tabular}{@{}lrrr@{}}
\toprule
\textbf{Category} & \textbf{Value 1} & \textbf{Value 2} & \textbf{Total} \\
\midrule
Item A & 10.5 & 20.3 & 30.8 \\
Item B & 15.2 & 25.1 & 40.3 \\
Item C & 22.1 & 12.7 & 34.8 \\
\bottomrule
\end{tabular}
\caption{Sample Data Table}
\label{tab:sample}
\end{table}

\section{Lists}
\subsection{Bullet Points}
\begin{itemize}
  \item First item
  \item Second item
  \item Third item with sub-items:
    \begin{itemize}
      \item Sub-item 1
      \item Sub-item 2
    \end{itemize}
\end{itemize}

\subsection{Numbered List}
\begin{enumerate}
  \item First step
  \item Second step
  \item Third step
\end{enumerate}

\section{Citations and References}
When you use the citation search feature, references will be automatically managed. For example, you might cite a paper like this: \cite{example2024}. All citations will be automatically added to your \texttt{cite.bib} file.

\section{Conclusion}
This template demonstrates basic \LaTeX{} features including mathematical formulas, tables, and lists. The bibliography system is now set up to work seamlessly with Bibby AI's citation management. Feel free to modify it for your needs!

\bibliographystyle{plain}
\bibliography{cite}

\end{document}\documentclass{article}
\usepackage[utf8]{inputenc}
\usepackage{booktabs}
\usepackage{amsmath}
\usepackage{natbib}

\title{My Document}
\author{Your Name}
\date{\today}

\begin{document}

\maketitle

\section{Introduction}
This is a sample document demonstrating various \LaTeX{} features including mathematical equations, tables, and formatting. You can cite papers using the citation search feature, and they will be automatically added to your \texttt{cite.bib} file.

\section{Mathematical Expressions}
Here are some example equations:

\subsection{Quadratic Formula}
The quadratic formula is given by:
\[x = \frac{-b \pm \sqrt{b^2 - 4ac}}{2a}\]

\subsection{Matrix Example}
A sample matrix:
\[
\begin{pmatrix}
a_{11} & a_{12} & a_{13} \\
a_{21} & a_{22} & a_{23} \\
a_{31} & a_{32} & a_{33}
\end{pmatrix}
\]

\section{Tables}
Below is a sample table showing data:

\begin{table}[h]
\centering
\begin{tabular}{@{}lrrr@{}}
\toprule
\textbf{Category} & \textbf{Value 1} & \textbf{Value 2} & \textbf{Total} \\
\midrule
Item A & 10.5 & 20.3 & 30.8 \\
Item B & 15.2 & 25.1 & 40.3 \\
Item C & 22.1 & 12.7 & 34.8 \\
\bottomrule
\end{tabular}
\caption{Sample Data Table}
\label{tab:sample}
\end{table}

\section{Lists}
\subsection{Bullet Points}
\begin{itemize}
  \item First item
  \item Second item
  \item Third item with sub-items:
    \begin{itemize}
      \item Sub-item 1
      \item Sub-item 2
    \end{itemize}
\end{itemize}

\subsection{Numbered List}
\begin{enumerate}
  \item First step
  \item Second step
  \item Third step
\end{enumerate}

\section{Citations and References}
When you use the citation search feature, references will be automatically managed. For example, you might cite a paper like this: \cite{example2024}. All citations will be automatically added to your \texttt{cite.bib} file.

\section{Conclusion}
This template demonstrates basic \LaTeX{} features including mathematical formulas, tables, and lists. The bibliography system is now set up to work seamlessly with Bibby AI's citation management. Feel free to modify it for your needs!

\bibliographystyle{plain}
\bibliography{cite}

\end{document}\documentclass{article}
\usepackage[utf8]{inputenc}
\usepackage{booktabs}
\usepackage{amsmath}
\usepackage{natbib}

\title{My Document}
\author{Your Name}
\date{\today}

\begin{document}

\maketitle

\section{Introduction}
This is a sample document demonstrating various \LaTeX{} features including mathematical equations, tables, and formatting. You can cite papers using the citation search feature, and they will be automatically added to your \texttt{cite.bib} file.

\section{Mathematical Expressions}
Here are some example equations:

\subsection{Quadratic Formula}
The quadratic formula is given by:
\[x = \frac{-b \pm \sqrt{b^2 - 4ac}}{2a}\]

\subsection{Matrix Example}
A sample matrix:
\[
\begin{pmatrix}
a_{11} & a_{12} & a_{13} \\
a_{21} & a_{22} & a_{23} \\
a_{31} & a_{32} & a_{33}
\end{pmatrix}
\]

\section{Tables}
Below is a sample table showing data:

\begin{table}[h]
\centering
\begin{tabular}{@{}lrrr@{}}
\toprule
\textbf{Category} & \textbf{Value 1} & \textbf{Value 2} & \textbf{Total} \\
\midrule
Item A & 10.5 & 20.3 & 30.8 \\
Item B & 15.2 & 25.1 & 40.3 \\
Item C & 22.1 & 12.7 & 34.8 \\
\bottomrule
\end{tabular}
\caption{Sample Data Table}
\label{tab:sample}
\end{table}

\section{Lists}
\subsection{Bullet Points}
\begin{itemize}
  \item First item
  \item Second item
  \item Third item with sub-items:
    \begin{itemize}
      \item Sub-item 1
      \item Sub-item 2
    \end{itemize}
\end{itemize}

\subsection{Numbered List}
\begin{enumerate}
  \item First step
  \item Second step
  \item Third step
\end{enumerate}

\section{Citations and References}
When you use the citation search feature, references will be automatically managed. For example, you might cite a paper like this: \cite{example2024}. All citations will be automatically added to your \texttt{cite.bib} file.

\section{Conclusion}
This template demonstrates basic \LaTeX{} features including mathematical formulas, tables, and lists. The bibliography system is now set up to work seamlessly with Bibby AI's citation management. Feel free to modify it for your needs!

\bibliographystyle{plain}
\bibliography{cite}

\end{document}\documentclass{article}
\usepackage[utf8]{inputenc}
\usepackage{booktabs}
\usepackage{amsmath}
\usepackage{natbib}

\title{My Document}
\author{Your Name}
\date{\today}

\begin{document}

\maketitle

\section{Introduction}
This is a sample document demonstrating various \LaTeX{} features including mathematical equations, tables, and formatting. You can cite papers using the citation search feature, and they will be automatically added to your \texttt{cite.bib} file.

\section{Mathematical Expressions}
Here are some example equations:

\subsection{Quadratic Formula}
The quadratic formula is given by:
\[x = \frac{-b \pm \sqrt{b^2 - 4ac}}{2a}\]

\subsection{Matrix Example}
A sample matrix:
\[
\begin{pmatrix}
a_{11} & a_{12} & a_{13} \\
a_{21} & a_{22} & a_{23} \\
a_{31} & a_{32} & a_{33}
\end{pmatrix}
\]

\section{Tables}
Below is a sample table showing data:

\begin{table}[h]
\centering
\begin{tabular}{@{}lrrr@{}}
\toprule
\textbf{Category} & \textbf{Value 1} & \textbf{Value 2} & \textbf{Total} \\
\midrule
Item A & 10.5 & 20.3 & 30.8 \\
Item B & 15.2 & 25.1 & 40.3 \\
Item C & 22.1 & 12.7 & 34.8 \\
\bottomrule
\end{tabular}
\caption{Sample Data Table}
\label{tab:sample}
\end{table}

\section{Lists}
\subsection{Bullet Points}
\begin{itemize}
  \item First item
  \item Second item
  \item Third item with sub-items:
    \begin{itemize}
      \item Sub-item 1
      \item Sub-item 2
    \end{itemize}
\end{itemize}

\subsection{Numbered List}
\begin{enumerate}
  \item First step
  \item Second step
  \item Third step
\end{enumerate}

\section{Citations and References}
When you use the citation search feature, references will be automatically managed. For example, you might cite a paper like this: \cite{example2024}. All citations will be automatically added to your \texttt{cite.bib} file.

\section{Conclusion}
This template demonstrates basic \LaTeX{} features including mathematical formulas, tables, and lists. The bibliography system is now set up to work seamlessly with Bibby AI's citation management. Feel free to modify it for your needs!

\bibliographystyle{plain}
\bibliography{cite}

\end{document}\documentclass{article}
\usepackage[utf8]{inputenc}
\usepackage{booktabs}
\usepackage{amsmath}
\usepackage{natbib}

\title{My Document}
\author{Your Name}
\date{\today}

\begin{document}

\maketitle

\section{Introduction}
This is a sample document demonstrating various \LaTeX{} features including mathematical equations, tables, and formatting. You can cite papers using the citation search feature, and they will be automatically added to your \texttt{cite.bib} file.

\section{Mathematical Expressions}
Here are some example equations:

\subsection{Quadratic Formula}
The quadratic formula is given by:
\[x = \frac{-b \pm \sqrt{b^2 - 4ac}}{2a}\]

\subsection{Matrix Example}
A sample matrix:
\[
\begin{pmatrix}
a_{11} & a_{12} & a_{13} \\
a_{21} & a_{22} & a_{23} \\
a_{31} & a_{32} & a_{33}
\end{pmatrix}
\]

\section{Tables}
Below is a sample table showing data:

\begin{table}[h]
\centering
\begin{tabular}{@{}lrrr@{}}
\toprule
\textbf{Category} & \textbf{Value 1} & \textbf{Value 2} & \textbf{Total} \\
\midrule
Item A & 10.5 & 20.3 & 30.8 \\
Item B & 15.2 & 25.1 & 40.3 \\
Item C & 22.1 & 12.7 & 34.8 \\
\bottomrule
\end{tabular}
\caption{Sample Data Table}
\label{tab:sample}
\end{table}

\section{Lists}
\subsection{Bullet Points}
\begin{itemize}
  \item First item
  \item Second item
  \item Third item with sub-items:
    \begin{itemize}
      \item Sub-item 1
      \item Sub-item 2
    \end{itemize}
\end{itemize}

\subsection{Numbered List}
\begin{enumerate}
  \item First step
  \item Second step
  \item Third step
\end{enumerate}

\section{Citations and References}
When you use the citation search feature, references will be automatically managed. For example, you might cite a paper like this: \cite{example2024}. All citations will be automatically added to your \texttt{cite.bib} file.

\section{Conclusion}
This template demonstrates basic \LaTeX{} features including mathematical formulas, tables, and lists. The bibliography system is now set up to work seamlessly with Bibby AI's citation management. Feel free to modify it for your needs!

\bibliographystyle{plain}
\bibliography{cite}

\end{document}\documentclass{article}
\usepackage[utf8]{inputenc}
\usepackage{booktabs}
\usepackage{amsmath}
\usepackage{natbib}

\title{My Document}
\author{Your Name}
\date{\today}

\begin{document}

\maketitle

\section{Introduction}
This is a sample document demonstrating various \LaTeX{} features including mathematical equations, tables, and formatting. You can cite papers using the citation search feature, and they will be automatically added to your \texttt{cite.bib} file.

\section{Mathematical Expressions}
Here are some example equations:

\subsection{Quadratic Formula}
The quadratic formula is given by:
\[x = \frac{-b \pm \sqrt{b^2 - 4ac}}{2a}\]

\subsection{Matrix Example}
A sample matrix:
\[
\begin{pmatrix}
a_{11} & a_{12} & a_{13} \\
a_{21} & a_{22} & a_{23} \\
a_{31} & a_{32} & a_{33}
\end{pmatrix}
\]

\section{Tables}
Below is a sample table showing data:

\begin{table}[h]
\centering
\begin{tabular}{@{}lrrr@{}}
\toprule
\textbf{Category} & \textbf{Value 1} & \textbf{Value 2} & \textbf{Total} \\
\midrule
Item A & 10.5 & 20.3 & 30.8 \\
Item B & 15.2 & 25.1 & 40.3 \\
Item C & 22.1 & 12.7 & 34.8 \\
\bottomrule
\end{tabular}
\caption{Sample Data Table}
\label{tab:sample}
\end{table}

\section{Lists}
\subsection{Bullet Points}
\begin{itemize}
  \item First item
  \item Second item
  \item Third item with sub-items:
    \begin{itemize}
      \item Sub-item 1
      \item Sub-item 2
    \end{itemize}
\end{itemize}

\subsection{Numbered List}
\begin{enumerate}
  \item First step
  \item Second step
  \item Third step
\end{enumerate}

\section{Citations and References}
When you use the citation search feature, references will be automatically managed. For example, you might cite a paper like this: \cite{example2024}. All citations will be automatically added to your \texttt{cite.bib} file.

\section{Conclusion}
This template demonstrates basic \LaTeX{} features including mathematical formulas, tables, and lists. The bibliography system is now set up to work seamlessly with Bibby AI's citation management. Feel free to modify it for your needs!

\bibliographystyle{plain}
\bibliography{cite}

\end{document}\documentclass{article}
\usepackage[utf8]{inputenc}
\usepackage{booktabs}
\usepackage{amsmath}
\usepackage{natbib}

\title{My Document}
\author{Your Name}
\date{\today}

\begin{document}

\maketitle

\section{Introduction}
This is a sample document demonstrating various \LaTeX{} features including mathematical equations, tables, and formatting. You can cite papers using the citation search feature, and they will be automatically added to your \texttt{cite.bib} file.

\section{Mathematical Expressions}
Here are some example equations:

\subsection{Quadratic Formula}
The quadratic formula is given by:
\[x = \frac{-b \pm \sqrt{b^2 - 4ac}}{2a}\]

\subsection{Matrix Example}
A sample matrix:
\[
\begin{pmatrix}
a_{11} & a_{12} & a_{13} \\
a_{21} & a_{22} & a_{23} \\
a_{31} & a_{32} & a_{33}
\end{pmatrix}
\]

\section{Tables}
Below is a sample table showing data:

\begin{table}[h]
\centering
\begin{tabular}{@{}lrrr@{}}
\toprule
\textbf{Category} & \textbf{Value 1} & \textbf{Value 2} & \textbf{Total} \\
\midrule
Item A & 10.5 & 20.3 & 30.8 \\
Item B & 15.2 & 25.1 & 40.3 \\
Item C & 22.1 & 12.7 & 34.8 \\
\bottomrule
\end{tabular}
\caption{Sample Data Table}
\label{tab:sample}
\end{table}

\section{Lists}
\subsection{Bullet Points}
\begin{itemize}
  \item First item
  \item Second item
  \item Third item with sub-items:
    \begin{itemize}
      \item Sub-item 1
      \item Sub-item 2
    \end{itemize}
\end{itemize}

\subsection{Numbered List}
\begin{enumerate}
  \item First step
  \item Second step
  \item Third step
\end{enumerate}

\section{Citations and References}
When you use the citation search feature, references will be automatically managed. For example, you might cite a paper like this: \cite{example2024}. All citations will be automatically added to your \texttt{cite.bib} file.

\section{Conclusion}
This template demonstrates basic \LaTeX{} features including mathematical formulas, tables, and lists. The bibliography system is now set up to work seamlessly with Bibby AI's citation management. Feel free to modify it for your needs!

\bibliographystyle{plain}
\bibliography{cite}

\end{document}\documentclass{article}
\usepackage[utf8]{inputenc}
\usepackage{booktabs}
\usepackage{amsmath}
\usepackage{natbib}

\title{My Document}
\author{Your Name}
\date{\today}

\begin{document}

\maketitle

\section{Introduction}
This is a sample document demonstrating various \LaTeX{} features including mathematical equations, tables, and formatting. You can cite papers using the citation search feature, and they will be automatically added to your \texttt{cite.bib} file.

\section{Mathematical Expressions}
Here are some example equations:

\subsection{Quadratic Formula}
The quadratic formula is given by:
\[x = \frac{-b \pm \sqrt{b^2 - 4ac}}{2a}\]

\subsection{Matrix Example}
A sample matrix:
\[
\begin{pmatrix}
a_{11} & a_{12} & a_{13} \\
a_{21} & a_{22} & a_{23} \\
a_{31} & a_{32} & a_{33}
\end{pmatrix}
\]

\section{Tables}
Below is a sample table showing data:

\begin{table}[h]
\centering
\begin{tabular}{@{}lrrr@{}}
\toprule
\textbf{Category} & \textbf{Value 1} & \textbf{Value 2} & \textbf{Total} \\
\midrule
Item A & 10.5 & 20.3 & 30.8 \\
Item B & 15.2 & 25.1 & 40.3 \\
Item C & 22.1 & 12.7 & 34.8 \\
\bottomrule
\end{tabular}
\caption{Sample Data Table}
\label{tab:sample}
\end{table}

\section{Lists}
\subsection{Bullet Points}
\begin{itemize}
  \item First item
  \item Second item
  \item Third item with sub-items:
    \begin{itemize}
      \item Sub-item 1
      \item Sub-item 2
    \end{itemize}
\end{itemize}

\subsection{Numbered List}
\begin{enumerate}
  \item First step
  \item Second step
  \item Third step
\end{enumerate}

\section{Citations and References}
When you use the citation search feature, references will be automatically managed. For example, you might cite a paper like this: \cite{example2024}. All citations will be automatically added to your \texttt{cite.bib} file.

\section{Conclusion}
This template demonstrates basic \LaTeX{} features including mathematical formulas, tables, and lists. The bibliography system is now set up to work seamlessly with Bibby AI's citation management. Feel free to modify it for your needs!

\bibliographystyle{plain}
\bibliography{cite}

\end{document}\documentclass{article}
\usepackage[utf8]{inputenc}
\usepackage{booktabs}
\usepackage{amsmath}
\usepackage{natbib}

\title{My Document}
\author{Your Name}
\date{\today}

\begin{document}

\maketitle

\section{Introduction}
This is a sample document demonstrating various \LaTeX{} features including mathematical equations, tables, and formatting. You can cite papers using the citation search feature, and they will be automatically added to your \texttt{cite.bib} file.

\section{Mathematical Expressions}
Here are some example equations:

\subsection{Quadratic Formula}
The quadratic formula is given by:
\[x = \frac{-b \pm \sqrt{b^2 - 4ac}}{2a}\]

\subsection{Matrix Example}
A sample matrix:
\[
\begin{pmatrix}
a_{11} & a_{12} & a_{13} \\
a_{21} & a_{22} & a_{23} \\
a_{31} & a_{32} & a_{33}
\end{pmatrix}
\]

\section{Tables}
Below is a sample table showing data:

\begin{table}[h]
\centering
\begin{tabular}{@{}lrrr@{}}
\toprule
\textbf{Category} & \textbf{Value 1} & \textbf{Value 2} & \textbf{Total} \\
\midrule
Item A & 10.5 & 20.3 & 30.8 \\
Item B & 15.2 & 25.1 & 40.3 \\
Item C & 22.1 & 12.7 & 34.8 \\
\bottomrule
\end{tabular}
\caption{Sample Data Table}
\label{tab:sample}
\end{table}

\section{Lists}
\subsection{Bullet Points}
\begin{itemize}
  \item First item
  \item Second item
  \item Third item with sub-items:
    \begin{itemize}
      \item Sub-item 1
      \item Sub-item 2
    \end{itemize}
\end{itemize}

\subsection{Numbered List}
\begin{enumerate}
  \item First step
  \item Second step
  \item Third step
\end{enumerate}

\section{Citations and References}
When you use the citation search feature, references will be automatically managed. For example, you might cite a paper like this: \cite{example2024}. All citations will be automatically added to your \texttt{cite.bib} file.

\section{Conclusion}
This template demonstrates basic \LaTeX{} features including mathematical formulas, tables, and lists. The bibliography system is now set up to work seamlessly with Bibby AI's citation management. Feel free to modify it for your needs!

\bibliographystyle{plain}
\bibliography{cite}

\end{document}\documentclass{article}
\usepackage[utf8]{inputenc}
\usepackage{booktabs}
\usepackage{amsmath}
\usepackage{natbib}

\title{My Document}
\author{Your Name}
\date{\today}

\begin{document}

\maketitle

\section{Introduction}
This is a sample document demonstrating various \LaTeX{} features including mathematical equations, tables, and formatting. You can cite papers using the citation search feature, and they will be automatically added to your \texttt{cite.bib} file.

\section{Mathematical Expressions}
Here are some example equations:

\subsection{Quadratic Formula}
The quadratic formula is given by:
\[x = \frac{-b \pm \sqrt{b^2 - 4ac}}{2a}\]

\subsection{Matrix Example}
A sample matrix:
\[
\begin{pmatrix}
a_{11} & a_{12} & a_{13} \\
a_{21} & a_{22} & a_{23} \\
a_{31} & a_{32} & a_{33}
\end{pmatrix}
\]

\section{Tables}
Below is a sample table showing data:

\begin{table}[h]
\centering
\begin{tabular}{@{}lrrr@{}}
\toprule
\textbf{Category} & \textbf{Value 1} & \textbf{Value 2} & \textbf{Total} \\
\midrule
Item A & 10.5 & 20.3 & 30.8 \\
Item B & 15.2 & 25.1 & 40.3 \\
Item C & 22.1 & 12.7 & 34.8 \\
\bottomrule
\end{tabular}
\caption{Sample Data Table}
\label{tab:sample}
\end{table}

\section{Lists}
\subsection{Bullet Points}
\begin{itemize}
  \item First item
  \item Second item
  \item Third item with sub-items:
    \begin{itemize}
      \item Sub-item 1
      \item Sub-item 2
    \end{itemize}
\end{itemize}

\subsection{Numbered List}
\begin{enumerate}
  \item First step
  \item Second step
  \item Third step
\end{enumerate}

\section{Citations and References}
When you use the citation search feature, references will be automatically managed. For example, you might cite a paper like this: \cite{example2024}. All citations will be automatically added to your \texttt{cite.bib} file.

\section{Conclusion}
This template demonstrates basic \LaTeX{} features including mathematical formulas, tables, and lists. The bibliography system is now set up to work seamlessly with Bibby AI's citation management. Feel free to modify it for your needs!

\bibliographystyle{plain}
\bibliography{cite}

\end{document}\documentclass{article}
\usepackage[utf8]{inputenc}
\usepackage{booktabs}
\usepackage{amsmath}
\usepackage{natbib}

\title{My Document}
\author{Your Name}
\date{\today}

\begin{document}

\maketitle

\section{Introduction}
This is a sample document demonstrating various \LaTeX{} features including mathematical equations, tables, and formatting. You can cite papers using the citation search feature, and they will be automatically added to your \texttt{cite.bib} file.

\section{Mathematical Expressions}
Here are some example equations:

\subsection{Quadratic Formula}
The quadratic formula is given by:
\[x = \frac{-b \pm \sqrt{b^2 - 4ac}}{2a}\]

\subsection{Matrix Example}
A sample matrix:
\[
\begin{pmatrix}
a_{11} & a_{12} & a_{13} \\
a_{21} & a_{22} & a_{23} \\
a_{31} & a_{32} & a_{33}
\end{pmatrix}
\]

\section{Tables}
Below is a sample table showing data:

\begin{table}[h]
\centering
\begin{tabular}{@{}lrrr@{}}
\toprule
\textbf{Category} & \textbf{Value 1} & \textbf{Value 2} & \textbf{Total} \\
\midrule
Item A & 10.5 & 20.3 & 30.8 \\
Item B & 15.2 & 25.1 & 40.3 \\
Item C & 22.1 & 12.7 & 34.8 \\
\bottomrule
\end{tabular}
\caption{Sample Data Table}
\label{tab:sample}
\end{table}

\section{Lists}
\subsection{Bullet Points}
\begin{itemize}
  \item First item
  \item Second item
  \item Third item with sub-items:
    \begin{itemize}
      \item Sub-item 1
      \item Sub-item 2
    \end{itemize}
\end{itemize}

\subsection{Numbered List}
\begin{enumerate}
  \item First step
  \item Second step
  \item Third step
\end{enumerate}

\section{Citations and References}
When you use the citation search feature, references will be automatically managed. For example, you might cite a paper like this: \cite{example2024}. All citations will be automatically added to your \texttt{cite.bib} file.

\section{Conclusion}
This template demonstrates basic \LaTeX{} features including mathematical formulas, tables, and lists. The bibliography system is now set up to work seamlessly with Bibby AI's citation management. Feel free to modify it for your needs!

\bibliographystyle{plain}
\bibliography{cite}

\end{document}\documentclass{article}
\usepackage[utf8]{inputenc}
\usepackage{booktabs}
\usepackage{amsmath}
\usepackage{natbib}

\title{My Document}
\author{Your Name}
\date{\today}

\begin{document}

\maketitle

\section{Introduction}
This is a sample document demonstrating various \LaTeX{} features including mathematical equations, tables, and formatting. You can cite papers using the citation search feature, and they will be automatically added to your \texttt{cite.bib} file.

\section{Mathematical Expressions}
Here are some example equations:

\subsection{Quadratic Formula}
The quadratic formula is given by:
\[x = \frac{-b \pm \sqrt{b^2 - 4ac}}{2a}\]

\subsection{Matrix Example}
A sample matrix:
\[
\begin{pmatrix}
a_{11} & a_{12} & a_{13} \\
a_{21} & a_{22} & a_{23} \\
a_{31} & a_{32} & a_{33}
\end{pmatrix}
\]

\section{Tables}
Below is a sample table showing data:

\begin{table}[h]
\centering
\begin{tabular}{@{}lrrr@{}}
\toprule
\textbf{Category} & \textbf{Value 1} & \textbf{Value 2} & \textbf{Total} \\
\midrule
Item A & 10.5 & 20.3 & 30.8 \\
Item B & 15.2 & 25.1 & 40.3 \\
Item C & 22.1 & 12.7 & 34.8 \\
\bottomrule
\end{tabular}
\caption{Sample Data Table}
\label{tab:sample}
\end{table}

\section{Lists}
\subsection{Bullet Points}
\begin{itemize}
  \item First item
  \item Second item
  \item Third item with sub-items:
    \begin{itemize}
      \item Sub-item 1
      \item Sub-item 2
    \end{itemize}
\end{itemize}

\subsection{Numbered List}
\begin{enumerate}
  \item First step
  \item Second step
  \item Third step
\end{enumerate}

\section{Citations and References}
When you use the citation search feature, references will be automatically managed. For example, you might cite a paper like this: \cite{example2024}. All citations will be automatically added to your \texttt{cite.bib} file.

\section{Conclusion}
This template demonstrates basic \LaTeX{} features including mathematical formulas, tables, and lists. The bibliography system is now set up to work seamlessly with Bibby AI's citation management. Feel free to modify it for your needs!

\bibliographystyle{plain}
\bibliography{cite}

\end{document}\documentclass{article}
\usepackage[utf8]{inputenc}
\usepackage{booktabs}
\usepackage{amsmath}
\usepackage{natbib}

\title{My Document}
\author{Your Name}
\date{\today}

\begin{document}

\maketitle

\section{Introduction}
This is a sample document demonstrating various \LaTeX{} features including mathematical equations, tables, and formatting. You can cite papers using the citation search feature, and they will be automatically added to your \texttt{cite.bib} file.

\section{Mathematical Expressions}
Here are some example equations:

\subsection{Quadratic Formula}
The quadratic formula is given by:
\[x = \frac{-b \pm \sqrt{b^2 - 4ac}}{2a}\]

\subsection{Matrix Example}
A sample matrix:
\[
\begin{pmatrix}
a_{11} & a_{12} & a_{13} \\
a_{21} & a_{22} & a_{23} \\
a_{31} & a_{32} & a_{33}
\end{pmatrix}
\]

\section{Tables}
Below is a sample table showing data:

\begin{table}[h]
\centering
\begin{tabular}{@{}lrrr@{}}
\toprule
\textbf{Category} & \textbf{Value 1} & \textbf{Value 2} & \textbf{Total} \\
\midrule
Item A & 10.5 & 20.3 & 30.8 \\
Item B & 15.2 & 25.1 & 40.3 \\
Item C & 22.1 & 12.7 & 34.8 \\
\bottomrule
\end{tabular}
\caption{Sample Data Table}
\label{tab:sample}
\end{table}

\section{Lists}
\subsection{Bullet Points}
\begin{itemize}
  \item First item
  \item Second item
  \item Third item with sub-items:
    \begin{itemize}
      \item Sub-item 1
      \item Sub-item 2
    \end{itemize}
\end{itemize}

\subsection{Numbered List}
\begin{enumerate}
  \item First step
  \item Second step
  \item Third step
\end{enumerate}

\section{Citations and References}
When you use the citation search feature, references will be automatically managed. For example, you might cite a paper like this: \cite{example2024}. All citations will be automatically added to your \texttt{cite.bib} file.

\section{Conclusion}
This template demonstrates basic \LaTeX{} features including mathematical formulas, tables, and lists. The bibliography system is now set up to work seamlessly with Bibby AI's citation management. Feel free to modify it for your needs!

\bibliographystyle{plain}
\bibliography{cite}

\end{document}\documentclass{article}
\usepackage[utf8]{inputenc}
\usepackage{booktabs}
\usepackage{amsmath}
\usepackage{natbib}

\title{My Document}
\author{Your Name}
\date{\today}

\begin{document}

\maketitle

\section{Introduction}
This is a sample document demonstrating various \LaTeX{} features including mathematical equations, tables, and formatting. You can cite papers using the citation search feature, and they will be automatically added to your \texttt{cite.bib} file.

\section{Mathematical Expressions}
Here are some example equations:

\subsection{Quadratic Formula}
The quadratic formula is given by:
\[x = \frac{-b \pm \sqrt{b^2 - 4ac}}{2a}\]

\subsection{Matrix Example}
A sample matrix:
\[
\begin{pmatrix}
a_{11} & a_{12} & a_{13} \\
a_{21} & a_{22} & a_{23} \\
a_{31} & a_{32} & a_{33}
\end{pmatrix}
\]

\section{Tables}
Below is a sample table showing data:

\begin{table}[h]
\centering
\begin{tabular}{@{}lrrr@{}}
\toprule
\textbf{Category} & \textbf{Value 1} & \textbf{Value 2} & \textbf{Total} \\
\midrule
Item A & 10.5 & 20.3 & 30.8 \\
Item B & 15.2 & 25.1 & 40.3 \\
Item C & 22.1 & 12.7 & 34.8 \\
\bottomrule
\end{tabular}
\caption{Sample Data Table}
\label{tab:sample}
\end{table}

\section{Lists}
\subsection{Bullet Points}
\begin{itemize}
  \item First item
  \item Second item
  \item Third item with sub-items:
    \begin{itemize}
      \item Sub-item 1
      \item Sub-item 2
    \end{itemize}
\end{itemize}

\subsection{Numbered List}
\begin{enumerate}
  \item First step
  \item Second step
  \item Third step
\end{enumerate}

\section{Citations and References}
When you use the citation search feature, references will be automatically managed. For example, you might cite a paper like this: \cite{example2024}. All citations will be automatically added to your \texttt{cite.bib} file.

\section{Conclusion}
This template demonstrates basic \LaTeX{} features including mathematical formulas, tables, and lists. The bibliography system is now set up to work seamlessly with Bibby AI's citation management. Feel free to modify it for your needs!

\bibliographystyle{plain}
\bibliography{cite}

\end{document}\documentclass{article}
\usepackage[utf8]{inputenc}
\usepackage{booktabs}
\usepackage{amsmath}
\usepackage{natbib}

\title{My Document}
\author{Your Name}
\date{\today}

\begin{document}

\maketitle

\section{Introduction}
This is a sample document demonstrating various \LaTeX{} features including mathematical equations, tables, and formatting. You can cite papers using the citation search feature, and they will be automatically added to your \texttt{cite.bib} file.

\section{Mathematical Expressions}
Here are some example equations:

\subsection{Quadratic Formula}
The quadratic formula is given by:
\[x = \frac{-b \pm \sqrt{b^2 - 4ac}}{2a}\]

\subsection{Matrix Example}
A sample matrix:
\[
\begin{pmatrix}
a_{11} & a_{12} & a_{13} \\
a_{21} & a_{22} & a_{23} \\
a_{31} & a_{32} & a_{33}
\end{pmatrix}
\]

\section{Tables}
Below is a sample table showing data:

\begin{table}[h]
\centering
\begin{tabular}{@{}lrrr@{}}
\toprule
\textbf{Category} & \textbf{Value 1} & \textbf{Value 2} & \textbf{Total} \\
\midrule
Item A & 10.5 & 20.3 & 30.8 \\
Item B & 15.2 & 25.1 & 40.3 \\
Item C & 22.1 & 12.7 & 34.8 \\
\bottomrule
\end{tabular}
\caption{Sample Data Table}
\label{tab:sample}
\end{table}

\section{Lists}
\subsection{Bullet Points}
\begin{itemize}
  \item First item
  \item Second item
  \item Third item with sub-items:
    \begin{itemize}
      \item Sub-item 1
      \item Sub-item 2
    \end{itemize}
\end{itemize}

\subsection{Numbered List}
\begin{enumerate}
  \item First step
  \item Second step
  \item Third step
\end{enumerate}

\section{Citations and References}
When you use the citation search feature, references will be automatically managed. For example, you might cite a paper like this: \cite{example2024}. All citations will be automatically added to your \texttt{cite.bib} file.

\section{Conclusion}
This template demonstrates basic \LaTeX{} features including mathematical formulas, tables, and lists. The bibliography system is now set up to work seamlessly with Bibby AI's citation management. Feel free to modify it for your needs!

\bibliographystyle{plain}
\bibliography{cite}

\end{document}\documentclass{article}
\usepackage[utf8]{inputenc}
\usepackage{booktabs}
\usepackage{amsmath}
\usepackage{natbib}

\title{My Document}
\author{Your Name}
\date{\today}

\begin{document}

\maketitle

\section{Introduction}
This is a sample document demonstrating various \LaTeX{} features including mathematical equations, tables, and formatting. You can cite papers using the citation search feature, and they will be automatically added to your \texttt{cite.bib} file.

\section{Mathematical Expressions}
Here are some example equations:

\subsection{Quadratic Formula}
The quadratic formula is given by:
\[x = \frac{-b \pm \sqrt{b^2 - 4ac}}{2a}\]

\subsection{Matrix Example}
A sample matrix:
\[
\begin{pmatrix}
a_{11} & a_{12} & a_{13} \\
a_{21} & a_{22} & a_{23} \\
a_{31} & a_{32} & a_{33}
\end{pmatrix}
\]

\section{Tables}
Below is a sample table showing data:

\begin{table}[h]
\centering
\begin{tabular}{@{}lrrr@{}}
\toprule
\textbf{Category} & \textbf{Value 1} & \textbf{Value 2} & \textbf{Total} \\
\midrule
Item A & 10.5 & 20.3 & 30.8 \\
Item B & 15.2 & 25.1 & 40.3 \\
Item C & 22.1 & 12.7 & 34.8 \\
\bottomrule
\end{tabular}
\caption{Sample Data Table}
\label{tab:sample}
\end{table}

\section{Lists}
\subsection{Bullet Points}
\begin{itemize}
  \item First item
  \item Second item
  \item Third item with sub-items:
    \begin{itemize}
      \item Sub-item 1
      \item Sub-item 2
    \end{itemize}
\end{itemize}

\subsection{Numbered List}
\begin{enumerate}
  \item First step
  \item Second step
  \item Third step
\end{enumerate}

\section{Citations and References}
When you use the citation search feature, references will be automatically managed. For example, you might cite a paper like this: \cite{example2024}. All citations will be automatically added to your \texttt{cite.bib} file.

\section{Conclusion}
This template demonstrates basic \LaTeX{} features including mathematical formulas, tables, and lists. The bibliography system is now set up to work seamlessly with Bibby AI's citation management. Feel free to modify it for your needs!

\bibliographystyle{plain}
\bibliography{cite}

\end{document}\documentclass{article}
\usepackage[utf8]{inputenc}
\usepackage{booktabs}
\usepackage{amsmath}
\usepackage{natbib}

\title{My Document}
\author{Your Name}
\date{\today}

\begin{document}

\maketitle

\section{Introduction}
This is a sample document demonstrating various \LaTeX{} features including mathematical equations, tables, and formatting. You can cite papers using the citation search feature, and they will be automatically added to your \texttt{cite.bib} file.

\section{Mathematical Expressions}
Here are some example equations:

\subsection{Quadratic Formula}
The quadratic formula is given by:
\[x = \frac{-b \pm \sqrt{b^2 - 4ac}}{2a}\]

\subsection{Matrix Example}
A sample matrix:
\[
\begin{pmatrix}
a_{11} & a_{12} & a_{13} \\
a_{21} & a_{22} & a_{23} \\
a_{31} & a_{32} & a_{33}
\end{pmatrix}
\]

\section{Tables}
Below is a sample table showing data:

\begin{table}[h]
\centering
\begin{tabular}{@{}lrrr@{}}
\toprule
\textbf{Category} & \textbf{Value 1} & \textbf{Value 2} & \textbf{Total} \\
\midrule
Item A & 10.5 & 20.3 & 30.8 \\
Item B & 15.2 & 25.1 & 40.3 \\
Item C & 22.1 & 12.7 & 34.8 \\
\bottomrule
\end{tabular}
\caption{Sample Data Table}
\label{tab:sample}
\end{table}

\section{Lists}
\subsection{Bullet Points}
\begin{itemize}
  \item First item
  \item Second item
  \item Third item with sub-items:
    \begin{itemize}
      \item Sub-item 1
      \item Sub-item 2
    \end{itemize}
\end{itemize}

\subsection{Numbered List}
\begin{enumerate}
  \item First step
  \item Second step
  \item Third step
\end{enumerate}

\section{Citations and References}
When you use the citation search feature, references will be automatically managed. For example, you might cite a paper like this: \cite{example2024}. All citations will be automatically added to your \texttt{cite.bib} file.

\section{Conclusion}
This template demonstrates basic \LaTeX{} features including mathematical formulas, tables, and lists. The bibliography system is now set up to work seamlessly with Bibby AI's citation management. Feel free to modify it for your needs!

\bibliographystyle{plain}
\bibliography{cite}

\end{document}\documentclass{article}
\usepackage[utf8]{inputenc}
\usepackage{booktabs}
\usepackage{amsmath}
\usepackage{natbib}

\title{My Document}
\author{Your Name}
\date{\today}

\begin{document}

\maketitle

\section{Introduction}
This is a sample document demonstrating various \LaTeX{} features including mathematical equations, tables, and formatting. You can cite papers using the citation search feature, and they will be automatically added to your \texttt{cite.bib} file.

\section{Mathematical Expressions}
Here are some example equations:

\subsection{Quadratic Formula}
The quadratic formula is given by:
\[x = \frac{-b \pm \sqrt{b^2 - 4ac}}{2a}\]

\subsection{Matrix Example}
A sample matrix:
\[
\begin{pmatrix}
a_{11} & a_{12} & a_{13} \\
a_{21} & a_{22} & a_{23} \\
a_{31} & a_{32} & a_{33}
\end{pmatrix}
\]

\section{Tables}
Below is a sample table showing data:

\begin{table}[h]
\centering
\begin{tabular}{@{}lrrr@{}}
\toprule
\textbf{Category} & \textbf{Value 1} & \textbf{Value 2} & \textbf{Total} \\
\midrule
Item A & 10.5 & 20.3 & 30.8 \\
Item B & 15.2 & 25.1 & 40.3 \\
Item C & 22.1 & 12.7 & 34.8 \\
\bottomrule
\end{tabular}
\caption{Sample Data Table}
\label{tab:sample}
\end{table}

\section{Lists}
\subsection{Bullet Points}
\begin{itemize}
  \item First item
  \item Second item
  \item Third item with sub-items:
    \begin{itemize}
      \item Sub-item 1
      \item Sub-item 2
    \end{itemize}
\end{itemize}

\subsection{Numbered List}
\begin{enumerate}
  \item First step
  \item Second step
  \item Third step
\end{enumerate}

\section{Citations and References}
When you use the citation search feature, references will be automatically managed. For example, you might cite a paper like this: \cite{example2024}. All citations will be automatically added to your \texttt{cite.bib} file.

\section{Conclusion}
This template demonstrates basic \LaTeX{} features including mathematical formulas, tables, and lists. The bibliography system is now set up to work seamlessly with Bibby AI's citation management. Feel free to modify it for your needs!

\bibliographystyle{plain}
\bibliography{cite}

\end{document}\documentclass{article}
\usepackage[utf8]{inputenc}
\usepackage{booktabs}
\usepackage{amsmath}
\usepackage{natbib}

\title{My Document}
\author{Your Name}
\date{\today}

\begin{document}

\maketitle

\section{Introduction}
This is a sample document demonstrating various \LaTeX{} features including mathematical equations, tables, and formatting. You can cite papers using the citation search feature, and they will be automatically added to your \texttt{cite.bib} file.

\section{Mathematical Expressions}
Here are some example equations:

\subsection{Quadratic Formula}
The quadratic formula is given by:
\[x = \frac{-b \pm \sqrt{b^2 - 4ac}}{2a}\]

\subsection{Matrix Example}
A sample matrix:
\[
\begin{pmatrix}
a_{11} & a_{12} & a_{13} \\
a_{21} & a_{22} & a_{23} \\
a_{31} & a_{32} & a_{33}
\end{pmatrix}
\]

\section{Tables}
Below is a sample table showing data:

\begin{table}[h]
\centering
\begin{tabular}{@{}lrrr@{}}
\toprule
\textbf{Category} & \textbf{Value 1} & \textbf{Value 2} & \textbf{Total} \\
\midrule
Item A & 10.5 & 20.3 & 30.8 \\
Item B & 15.2 & 25.1 & 40.3 \\
Item C & 22.1 & 12.7 & 34.8 \\
\bottomrule
\end{tabular}
\caption{Sample Data Table}
\label{tab:sample}
\end{table}

\section{Lists}
\subsection{Bullet Points}
\begin{itemize}
  \item First item
  \item Second item
  \item Third item with sub-items:
    \begin{itemize}
      \item Sub-item 1
      \item Sub-item 2
    \end{itemize}
\end{itemize}

\subsection{Numbered List}
\begin{enumerate}
  \item First step
  \item Second step
  \item Third step
\end{enumerate}

\section{Citations and References}
When you use the citation search feature, references will be automatically managed. For example, you might cite a paper like this: \cite{example2024}. All citations will be automatically added to your \texttt{cite.bib} file.

\section{Conclusion}
This template demonstrates basic \LaTeX{} features including mathematical formulas, tables, and lists. The bibliography system is now set up to work seamlessly with Bibby AI's citation management. Feel free to modify it for your needs!

\bibliographystyle{plain}
\bibliography{cite}

\end{document}\documentclass{article}
\usepackage[utf8]{inputenc}
\usepackage{booktabs}
\usepackage{amsmath}
\usepackage{natbib}

\title{My Document}
\author{Your Name}
\date{\today}

\begin{document}

\maketitle

\section{Introduction}
This is a sample document demonstrating various \LaTeX{} features including mathematical equations, tables, and formatting. You can cite papers using the citation search feature, and they will be automatically added to your \texttt{cite.bib} file.

\section{Mathematical Expressions}
Here are some example equations:

\subsection{Quadratic Formula}
The quadratic formula is given by:
\[x = \frac{-b \pm \sqrt{b^2 - 4ac}}{2a}\]

\subsection{Matrix Example}
A sample matrix:
\[
\begin{pmatrix}
a_{11} & a_{12} & a_{13} \\
a_{21} & a_{22} & a_{23} \\
a_{31} & a_{32} & a_{33}
\end{pmatrix}
\]

\section{Tables}
Below is a sample table showing data:

\begin{table}[h]
\centering
\begin{tabular}{@{}lrrr@{}}
\toprule
\textbf{Category} & \textbf{Value 1} & \textbf{Value 2} & \textbf{Total} \\
\midrule
Item A & 10.5 & 20.3 & 30.8 \\
Item B & 15.2 & 25.1 & 40.3 \\
Item C & 22.1 & 12.7 & 34.8 \\
\bottomrule
\end{tabular}
\caption{Sample Data Table}
\label{tab:sample}
\end{table}

\section{Lists}
\subsection{Bullet Points}
\begin{itemize}
  \item First item
  \item Second item
  \item Third item with sub-items:
    \begin{itemize}
      \item Sub-item 1
      \item Sub-item 2
    \end{itemize}
\end{itemize}

\subsection{Numbered List}
\begin{enumerate}
  \item First step
  \item Second step
  \item Third step
\end{enumerate}

\section{Citations and References}
When you use the citation search feature, references will be automatically managed. For example, you might cite a paper like this: \cite{example2024}. All citations will be automatically added to your \texttt{cite.bib} file.

\section{Conclusion}
This template demonstrates basic \LaTeX{} features including mathematical formulas, tables, and lists. The bibliography system is now set up to work seamlessly with Bibby AI's citation management. Feel free to modify it for your needs!

\bibliographystyle{plain}
\bibliography{cite}

\end{document}\documentclass{article}
\usepackage[utf8]{inputenc}
\usepackage{booktabs}
\usepackage{amsmath}
\usepackage{natbib}

\title{My Document}
\author{Your Name}
\date{\today}

\begin{document}

\maketitle

\section{Introduction}
This is a sample document demonstrating various \LaTeX{} features including mathematical equations, tables, and formatting. You can cite papers using the citation search feature, and they will be automatically added to your \texttt{cite.bib} file.

\section{Mathematical Expressions}
Here are some example equations:

\subsection{Quadratic Formula}
The quadratic formula is given by:
\[x = \frac{-b \pm \sqrt{b^2 - 4ac}}{2a}\]

\subsection{Matrix Example}
A sample matrix:
\[
\begin{pmatrix}
a_{11} & a_{12} & a_{13} \\
a_{21} & a_{22} & a_{23} \\
a_{31} & a_{32} & a_{33}
\end{pmatrix}
\]

\section{Tables}
Below is a sample table showing data:

\begin{table}[h]
\centering
\begin{tabular}{@{}lrrr@{}}
\toprule
\textbf{Category} & \textbf{Value 1} & \textbf{Value 2} & \textbf{Total} \\
\midrule
Item A & 10.5 & 20.3 & 30.8 \\
Item B & 15.2 & 25.1 & 40.3 \\
Item C & 22.1 & 12.7 & 34.8 \\
\bottomrule
\end{tabular}
\caption{Sample Data Table}
\label{tab:sample}
\end{table}

\section{Lists}
\subsection{Bullet Points}
\begin{itemize}
  \item First item
  \item Second item
  \item Third item with sub-items:
    \begin{itemize}
      \item Sub-item 1
      \item Sub-item 2
    \end{itemize}
\end{itemize}

\subsection{Numbered List}
\begin{enumerate}
  \item First step
  \item Second step
  \item Third step
\end{enumerate}

\section{Citations and References}
When you use the citation search feature, references will be automatically managed. For example, you might cite a paper like this: \cite{example2024}. All citations will be automatically added to your \texttt{cite.bib} file.

\section{Conclusion}
This template demonstrates basic \LaTeX{} features including mathematical formulas, tables, and lists. The bibliography system is now set up to work seamlessly with Bibby AI's citation management. Feel free to modify it for your needs!

\bibliographystyle{plain}
\bibliography{cite}

\end{document}\documentclass{article}
\usepackage[utf8]{inputenc}
\usepackage{booktabs}
\usepackage{amsmath}
\usepackage{natbib}

\title{My Document}
\author{Your Name}
\date{\today}

\begin{document}

\maketitle

\section{Introduction}
This is a sample document demonstrating various \LaTeX{} features including mathematical equations, tables, and formatting. You can cite papers using the citation search feature, and they will be automatically added to your \texttt{cite.bib} file.

\section{Mathematical Expressions}
Here are some example equations:

\subsection{Quadratic Formula}
The quadratic formula is given by:
\[x = \frac{-b \pm \sqrt{b^2 - 4ac}}{2a}\]

\subsection{Matrix Example}
A sample matrix:
\[
\begin{pmatrix}
a_{11} & a_{12} & a_{13} \\
a_{21} & a_{22} & a_{23} \\
a_{31} & a_{32} & a_{33}
\end{pmatrix}
\]

\section{Tables}
Below is a sample table showing data:

\begin{table}[h]
\centering
\begin{tabular}{@{}lrrr@{}}
\toprule
\textbf{Category} & \textbf{Value 1} & \textbf{Value 2} & \textbf{Total} \\
\midrule
Item A & 10.5 & 20.3 & 30.8 \\
Item B & 15.2 & 25.1 & 40.3 \\
Item C & 22.1 & 12.7 & 34.8 \\
\bottomrule
\end{tabular}
\caption{Sample Data Table}
\label{tab:sample}
\end{table}

\section{Lists}
\subsection{Bullet Points}
\begin{itemize}
  \item First item
  \item Second item
  \item Third item with sub-items:
    \begin{itemize}
      \item Sub-item 1
      \item Sub-item 2
    \end{itemize}
\end{itemize}

\subsection{Numbered List}
\begin{enumerate}
  \item First step
  \item Second step
  \item Third step
\end{enumerate}

\section{Citations and References}
When you use the citation search feature, references will be automatically managed. For example, you might cite a paper like this: \cite{example2024}. All citations will be automatically added to your \texttt{cite.bib} file.

\section{Conclusion}
This template demonstrates basic \LaTeX{} features including mathematical formulas, tables, and lists. The bibliography system is now set up to work seamlessly with Bibby AI's citation management. Feel free to modify it for your needs!

\bibliographystyle{plain}
\bibliography{cite}

\end{document}\documentclass{article}
\usepackage[utf8]{inputenc}
\usepackage{booktabs}
\usepackage{amsmath}
\usepackage{natbib}

\title{My Document}
\author{Your Name}
\date{\today}

\begin{document}

\maketitle

\section{Introduction}
This is a sample document demonstrating various \LaTeX{} features including mathematical equations, tables, and formatting. You can cite papers using the citation search feature, and they will be automatically added to your \texttt{cite.bib} file.

\section{Mathematical Expressions}
Here are some example equations:

\subsection{Quadratic Formula}
The quadratic formula is given by:
\[x = \frac{-b \pm \sqrt{b^2 - 4ac}}{2a}\]

\subsection{Matrix Example}
A sample matrix:
\[
\begin{pmatrix}
a_{11} & a_{12} & a_{13} \\
a_{21} & a_{22} & a_{23} \\
a_{31} & a_{32} & a_{33}
\end{pmatrix}
\]

\section{Tables}
Below is a sample table showing data:

\begin{table}[h]
\centering
\begin{tabular}{@{}lrrr@{}}
\toprule
\textbf{Category} & \textbf{Value 1} & \textbf{Value 2} & \textbf{Total} \\
\midrule
Item A & 10.5 & 20.3 & 30.8 \\
Item B & 15.2 & 25.1 & 40.3 \\
Item C & 22.1 & 12.7 & 34.8 \\
\bottomrule
\end{tabular}
\caption{Sample Data Table}
\label{tab:sample}
\end{table}

\section{Lists}
\subsection{Bullet Points}
\begin{itemize}
  \item First item
  \item Second item
  \item Third item with sub-items:
    \begin{itemize}
      \item Sub-item 1
      \item Sub-item 2
    \end{itemize}
\end{itemize}

\subsection{Numbered List}
\begin{enumerate}
  \item First step
  \item Second step
  \item Third step
\end{enumerate}

\section{Citations and References}
When you use the citation search feature, references will be automatically managed. For example, you might cite a paper like this: \cite{example2024}. All citations will be automatically added to your \texttt{cite.bib} file.

\section{Conclusion}
This template demonstrates basic \LaTeX{} features including mathematical formulas, tables, and lists. The bibliography system is now set up to work seamlessly with Bibby AI's citation management. Feel free to modify it for your needs!

\bibliographystyle{plain}
\bibliography{cite}

\end{document}\documentclass{article}
\usepackage[utf8]{inputenc}
\usepackage{booktabs}
\usepackage{amsmath}
\usepackage{natbib}

\title{My Document}
\author{Your Name}
\date{\today}

\begin{document}

\maketitle

\section{Introduction}
This is a sample document demonstrating various \LaTeX{} features including mathematical equations, tables, and formatting. You can cite papers using the citation search feature, and they will be automatically added to your \texttt{cite.bib} file.

\section{Mathematical Expressions}
Here are some example equations:

\subsection{Quadratic Formula}
The quadratic formula is given by:
\[x = \frac{-b \pm \sqrt{b^2 - 4ac}}{2a}\]

\subsection{Matrix Example}
A sample matrix:
\[
\begin{pmatrix}
a_{11} & a_{12} & a_{13} \\
a_{21} & a_{22} & a_{23} \\
a_{31} & a_{32} & a_{33}
\end{pmatrix}
\]

\section{Tables}
Below is a sample table showing data:

\begin{table}[h]
\centering
\begin{tabular}{@{}lrrr@{}}
\toprule
\textbf{Category} & \textbf{Value 1} & \textbf{Value 2} & \textbf{Total} \\
\midrule
Item A & 10.5 & 20.3 & 30.8 \\
Item B & 15.2 & 25.1 & 40.3 \\
Item C & 22.1 & 12.7 & 34.8 \\
\bottomrule
\end{tabular}
\caption{Sample Data Table}
\label{tab:sample}
\end{table}

\section{Lists}
\subsection{Bullet Points}
\begin{itemize}
  \item First item
  \item Second item
  \item Third item with sub-items:
    \begin{itemize}
      \item Sub-item 1
      \item Sub-item 2
    \end{itemize}
\end{itemize}

\subsection{Numbered List}
\begin{enumerate}
  \item First step
  \item Second step
  \item Third step
\end{enumerate}

\section{Citations and References}
When you use the citation search feature, references will be automatically managed. For example, you might cite a paper like this: \cite{example2024}. All citations will be automatically added to your \texttt{cite.bib} file.

\section{Conclusion}
This template demonstrates basic \LaTeX{} features including mathematical formulas, tables, and lists. The bibliography system is now set up to work seamlessly with Bibby AI's citation management. Feel free to modify it for your needs!

\bibliographystyle{plain}
\bibliography{cite}

\end{document}\documentclass{article}
\usepackage[utf8]{inputenc}
\usepackage{booktabs}
\usepackage{amsmath}
\usepackage{natbib}

\title{My Document}
\author{Your Name}
\date{\today}

\begin{document}

\maketitle

\section{Introduction}
This is a sample document demonstrating various \LaTeX{} features including mathematical equations, tables, and formatting. You can cite papers using the citation search feature, and they will be automatically added to your \texttt{cite.bib} file.

\section{Mathematical Expressions}
Here are some example equations:

\subsection{Quadratic Formula}
The quadratic formula is given by:
\[x = \frac{-b \pm \sqrt{b^2 - 4ac}}{2a}\]

\subsection{Matrix Example}
A sample matrix:
\[
\begin{pmatrix}
a_{11} & a_{12} & a_{13} \\
a_{21} & a_{22} & a_{23} \\
a_{31} & a_{32} & a_{33}
\end{pmatrix}
\]

\section{Tables}
Below is a sample table showing data:

\begin{table}[h]
\centering
\begin{tabular}{@{}lrrr@{}}
\toprule
\textbf{Category} & \textbf{Value 1} & \textbf{Value 2} & \textbf{Total} \\
\midrule
Item A & 10.5 & 20.3 & 30.8 \\
Item B & 15.2 & 25.1 & 40.3 \\
Item C & 22.1 & 12.7 & 34.8 \\
\bottomrule
\end{tabular}
\caption{Sample Data Table}
\label{tab:sample}
\end{table}

\section{Lists}
\subsection{Bullet Points}
\begin{itemize}
  \item First item
  \item Second item
  \item Third item with sub-items:
    \begin{itemize}
      \item Sub-item 1
      \item Sub-item 2
    \end{itemize}
\end{itemize}

\subsection{Numbered List}
\begin{enumerate}
  \item First step
  \item Second step
  \item Third step
\end{enumerate}

\section{Citations and References}
When you use the citation search feature, references will be automatically managed. For example, you might cite a paper like this: \cite{example2024}. All citations will be automatically added to your \texttt{cite.bib} file.

\section{Conclusion}
This template demonstrates basic \LaTeX{} features including mathematical formulas, tables, and lists. The bibliography system is now set up to work seamlessly with Bibby AI's citation management. Feel free to modify it for your needs!

\bibliographystyle{plain}
\bibliography{cite}

\end{document}\documentclass{article}
\usepackage[utf8]{inputenc}
\usepackage{booktabs}
\usepackage{amsmath}
\usepackage{natbib}

\title{My Document}
\author{Your Name}
\date{\today}

\begin{document}

\maketitle

\section{Introduction}
This is a sample document demonstrating various \LaTeX{} features including mathematical equations, tables, and formatting. You can cite papers using the citation search feature, and they will be automatically added to your \texttt{cite.bib} file.

\section{Mathematical Expressions}
Here are some example equations:

\subsection{Quadratic Formula}
The quadratic formula is given by:
\[x = \frac{-b \pm \sqrt{b^2 - 4ac}}{2a}\]

\subsection{Matrix Example}
A sample matrix:
\[
\begin{pmatrix}
a_{11} & a_{12} & a_{13} \\
a_{21} & a_{22} & a_{23} \\
a_{31} & a_{32} & a_{33}
\end{pmatrix}
\]

\section{Tables}
Below is a sample table showing data:

\begin{table}[h]
\centering
\begin{tabular}{@{}lrrr@{}}
\toprule
\textbf{Category} & \textbf{Value 1} & \textbf{Value 2} & \textbf{Total} \\
\midrule
Item A & 10.5 & 20.3 & 30.8 \\
Item B & 15.2 & 25.1 & 40.3 \\
Item C & 22.1 & 12.7 & 34.8 \\
\bottomrule
\end{tabular}
\caption{Sample Data Table}
\label{tab:sample}
\end{table}

\section{Lists}
\subsection{Bullet Points}
\begin{itemize}
  \item First item
  \item Second item
  \item Third item with sub-items:
    \begin{itemize}
      \item Sub-item 1
      \item Sub-item 2
    \end{itemize}
\end{itemize}

\subsection{Numbered List}
\begin{enumerate}
  \item First step
  \item Second step
  \item Third step
\end{enumerate}

\section{Citations and References}
When you use the citation search feature, references will be automatically managed. For example, you might cite a paper like this: \cite{example2024}. All citations will be automatically added to your \texttt{cite.bib} file.

\section{Conclusion}
This template demonstrates basic \LaTeX{} features including mathematical formulas, tables, and lists. The bibliography system is now set up to work seamlessly with Bibby AI's citation management. Feel free to modify it for your needs!

\bibliographystyle{plain}
\bibliography{cite}

\end{document}\documentclass{article}
\usepackage[utf8]{inputenc}
\usepackage{booktabs}
\usepackage{amsmath}
\usepackage{natbib}

\title{My Document}
\author{Your Name}
\date{\today}

\begin{document}

\maketitle

\section{Introduction}
This is a sample document demonstrating various \LaTeX{} features including mathematical equations, tables, and formatting. You can cite papers using the citation search feature, and they will be automatically added to your \texttt{cite.bib} file.

\section{Mathematical Expressions}
Here are some example equations:

\subsection{Quadratic Formula}
The quadratic formula is given by:
\[x = \frac{-b \pm \sqrt{b^2 - 4ac}}{2a}\]

\subsection{Matrix Example}
A sample matrix:
\[
\begin{pmatrix}
a_{11} & a_{12} & a_{13} \\
a_{21} & a_{22} & a_{23} \\
a_{31} & a_{32} & a_{33}
\end{pmatrix}
\]

\section{Tables}
Below is a sample table showing data:

\begin{table}[h]
\centering
\begin{tabular}{@{}lrrr@{}}
\toprule
\textbf{Category} & \textbf{Value 1} & \textbf{Value 2} & \textbf{Total} \\
\midrule
Item A & 10.5 & 20.3 & 30.8 \\
Item B & 15.2 & 25.1 & 40.3 \\
Item C & 22.1 & 12.7 & 34.8 \\
\bottomrule
\end{tabular}
\caption{Sample Data Table}
\label{tab:sample}
\end{table}

\section{Lists}
\subsection{Bullet Points}
\begin{itemize}
  \item First item
  \item Second item
  \item Third item with sub-items:
    \begin{itemize}
      \item Sub-item 1
      \item Sub-item 2
    \end{itemize}
\end{itemize}

\subsection{Numbered List}
\begin{enumerate}
  \item First step
  \item Second step
  \item Third step
\end{enumerate}

\section{Citations and References}
When you use the citation search feature, references will be automatically managed. For example, you might cite a paper like this: \cite{example2024}. All citations will be automatically added to your \texttt{cite.bib} file.

\section{Conclusion}
This template demonstrates basic \LaTeX{} features including mathematical formulas, tables, and lists. The bibliography system is now set up to work seamlessly with Bibby AI's citation management. Feel free to modify it for your needs!

\bibliographystyle{plain}
\bibliography{cite}

\end{document}\documentclass{article}
\usepackage[utf8]{inputenc}
\usepackage{booktabs}
\usepackage{amsmath}
\usepackage{natbib}

\title{My Document}
\author{Your Name}
\date{\today}

\begin{document}

\maketitle

\section{Introduction}
This is a sample document demonstrating various \LaTeX{} features including mathematical equations, tables, and formatting. You can cite papers using the citation search feature, and they will be automatically added to your \texttt{cite.bib} file.

\section{Mathematical Expressions}
Here are some example equations:

\subsection{Quadratic Formula}
The quadratic formula is given by:
\[x = \frac{-b \pm \sqrt{b^2 - 4ac}}{2a}\]

\subsection{Matrix Example}
A sample matrix:
\[
\begin{pmatrix}
a_{11} & a_{12} & a_{13} \\
a_{21} & a_{22} & a_{23} \\
a_{31} & a_{32} & a_{33}
\end{pmatrix}
\]

\section{Tables}
Below is a sample table showing data:

\begin{table}[h]
\centering
\begin{tabular}{@{}lrrr@{}}
\toprule
\textbf{Category} & \textbf{Value 1} & \textbf{Value 2} & \textbf{Total} \\
\midrule
Item A & 10.5 & 20.3 & 30.8 \\
Item B & 15.2 & 25.1 & 40.3 \\
Item C & 22.1 & 12.7 & 34.8 \\
\bottomrule
\end{tabular}
\caption{Sample Data Table}
\label{tab:sample}
\end{table}

\section{Lists}
\subsection{Bullet Points}
\begin{itemize}
  \item First item
  \item Second item
  \item Third item with sub-items:
    \begin{itemize}
      \item Sub-item 1
      \item Sub-item 2
    \end{itemize}
\end{itemize}

\subsection{Numbered List}
\begin{enumerate}
  \item First step
  \item Second step
  \item Third step
\end{enumerate}

\section{Citations and References}
When you use the citation search feature, references will be automatically managed. For example, you might cite a paper like this: \cite{example2024}. All citations will be automatically added to your \texttt{cite.bib} file.

\section{Conclusion}
This template demonstrates basic \LaTeX{} features including mathematical formulas, tables, and lists. The bibliography system is now set up to work seamlessly with Bibby AI's citation management. Feel free to modify it for your needs!

\bibliographystyle{plain}
\bibliography{cite}

\end{document}\documentclass{article}
\usepackage[utf8]{inputenc}
\usepackage{booktabs}
\usepackage{amsmath}
\usepackage{natbib}

\title{My Document}
\author{Your Name}
\date{\today}

\begin{document}

\maketitle

\section{Introduction}
This is a sample document demonstrating various \LaTeX{} features including mathematical equations, tables, and formatting. You can cite papers using the citation search feature, and they will be automatically added to your \texttt{cite.bib} file.

\section{Mathematical Expressions}
Here are some example equations:

\subsection{Quadratic Formula}
The quadratic formula is given by:
\[x = \frac{-b \pm \sqrt{b^2 - 4ac}}{2a}\]

\subsection{Matrix Example}
A sample matrix:
\[
\begin{pmatrix}
a_{11} & a_{12} & a_{13} \\
a_{21} & a_{22} & a_{23} \\
a_{31} & a_{32} & a_{33}
\end{pmatrix}
\]

\section{Tables}
Below is a sample table showing data:

\begin{table}[h]
\centering
\begin{tabular}{@{}lrrr@{}}
\toprule
\textbf{Category} & \textbf{Value 1} & \textbf{Value 2} & \textbf{Total} \\
\midrule
Item A & 10.5 & 20.3 & 30.8 \\
Item B & 15.2 & 25.1 & 40.3 \\
Item C & 22.1 & 12.7 & 34.8 \\
\bottomrule
\end{tabular}
\caption{Sample Data Table}
\label{tab:sample}
\end{table}

\section{Lists}
\subsection{Bullet Points}
\begin{itemize}
  \item First item
  \item Second item
  \item Third item with sub-items:
    \begin{itemize}
      \item Sub-item 1
      \item Sub-item 2
    \end{itemize}
\end{itemize}

\subsection{Numbered List}
\begin{enumerate}
  \item First step
  \item Second step
  \item Third step
\end{enumerate}

\section{Citations and References}
When you use the citation search feature, references will be automatically managed. For example, you might cite a paper like this: \cite{example2024}. All citations will be automatically added to your \texttt{cite.bib} file.

\section{Conclusion}
This template demonstrates basic \LaTeX{} features including mathematical formulas, tables, and lists. The bibliography system is now set up to work seamlessly with Bibby AI's citation management. Feel free to modify it for your needs!

\bibliographystyle{plain}
\bibliography{cite}

\end{document}\documentclass{article}
\usepackage[utf8]{inputenc}
\usepackage{booktabs}
\usepackage{amsmath}
\usepackage{natbib}

\title{My Document}
\author{Your Name}
\date{\today}

\begin{document}

\maketitle

\section{Introduction}
This is a sample document demonstrating various \LaTeX{} features including mathematical equations, tables, and formatting. You can cite papers using the citation search feature, and they will be automatically added to your \texttt{cite.bib} file.

\section{Mathematical Expressions}
Here are some example equations:

\subsection{Quadratic Formula}
The quadratic formula is given by:
\[x = \frac{-b \pm \sqrt{b^2 - 4ac}}{2a}\]

\subsection{Matrix Example}
A sample matrix:
\[
\begin{pmatrix}
a_{11} & a_{12} & a_{13} \\
a_{21} & a_{22} & a_{23} \\
a_{31} & a_{32} & a_{33}
\end{pmatrix}
\]

\section{Tables}
Below is a sample table showing data:

\begin{table}[h]
\centering
\begin{tabular}{@{}lrrr@{}}
\toprule
\textbf{Category} & \textbf{Value 1} & \textbf{Value 2} & \textbf{Total} \\
\midrule
Item A & 10.5 & 20.3 & 30.8 \\
Item B & 15.2 & 25.1 & 40.3 \\
Item C & 22.1 & 12.7 & 34.8 \\
\bottomrule
\end{tabular}
\caption{Sample Data Table}
\label{tab:sample}
\end{table}

\section{Lists}
\subsection{Bullet Points}
\begin{itemize}
  \item First item
  \item Second item
  \item Third item with sub-items:
    \begin{itemize}
      \item Sub-item 1
      \item Sub-item 2
    \end{itemize}
\end{itemize}

\subsection{Numbered List}
\begin{enumerate}
  \item First step
  \item Second step
  \item Third step
\end{enumerate}

\section{Citations and References}
When you use the citation search feature, references will be automatically managed. For example, you might cite a paper like this: \cite{example2024}. All citations will be automatically added to your \texttt{cite.bib} file.

\section{Conclusion}
This template demonstrates basic \LaTeX{} features including mathematical formulas, tables, and lists. The bibliography system is now set up to work seamlessly with Bibby AI's citation management. Feel free to modify it for your needs!

\bibliographystyle{plain}
\bibliography{cite}

\end{document}\documentclass{article}
\usepackage[utf8]{inputenc}
\usepackage{booktabs}
\usepackage{amsmath}
\usepackage{natbib}

\title{My Document}
\author{Your Name}
\date{\today}

\begin{document}

\maketitle

\section{Introduction}
This is a sample document demonstrating various \LaTeX{} features including mathematical equations, tables, and formatting. You can cite papers using the citation search feature, and they will be automatically added to your \texttt{cite.bib} file.

\section{Mathematical Expressions}
Here are some example equations:

\subsection{Quadratic Formula}
The quadratic formula is given by:
\[x = \frac{-b \pm \sqrt{b^2 - 4ac}}{2a}\]

\subsection{Matrix Example}
A sample matrix:
\[
\begin{pmatrix}
a_{11} & a_{12} & a_{13} \\
a_{21} & a_{22} & a_{23} \\
a_{31} & a_{32} & a_{33}
\end{pmatrix}
\]

\section{Tables}
Below is a sample table showing data:

\begin{table}[h]
\centering
\begin{tabular}{@{}lrrr@{}}
\toprule
\textbf{Category} & \textbf{Value 1} & \textbf{Value 2} & \textbf{Total} \\
\midrule
Item A & 10.5 & 20.3 & 30.8 \\
Item B & 15.2 & 25.1 & 40.3 \\
Item C & 22.1 & 12.7 & 34.8 \\
\bottomrule
\end{tabular}
\caption{Sample Data Table}
\label{tab:sample}
\end{table}

\section{Lists}
\subsection{Bullet Points}
\begin{itemize}
  \item First item
  \item Second item
  \item Third item with sub-items:
    \begin{itemize}
      \item Sub-item 1
      \item Sub-item 2
    \end{itemize}
\end{itemize}

\subsection{Numbered List}
\begin{enumerate}
  \item First step
  \item Second step
  \item Third step
\end{enumerate}

\section{Citations and References}
When you use the citation search feature, references will be automatically managed. For example, you might cite a paper like this: \cite{example2024}. All citations will be automatically added to your \texttt{cite.bib} file.

\section{Conclusion}
This template demonstrates basic \LaTeX{} features including mathematical formulas, tables, and lists. The bibliography system is now set up to work seamlessly with Bibby AI's citation management. Feel free to modify it for your needs!

\bibliographystyle{plain}
\bibliography{cite}

\end{document}\documentclass{article}
\usepackage[utf8]{inputenc}
\usepackage{booktabs}
\usepackage{amsmath}
\usepackage{natbib}

\title{My Document}
\author{Your Name}
\date{\today}

\begin{document}

\maketitle

\section{Introduction}
This is a sample document demonstrating various \LaTeX{} features including mathematical equations, tables, and formatting. You can cite papers using the citation search feature, and they will be automatically added to your \texttt{cite.bib} file.

\section{Mathematical Expressions}
Here are some example equations:

\subsection{Quadratic Formula}
The quadratic formula is given by:
\[x = \frac{-b \pm \sqrt{b^2 - 4ac}}{2a}\]

\subsection{Matrix Example}
A sample matrix:
\[
\begin{pmatrix}
a_{11} & a_{12} & a_{13} \\
a_{21} & a_{22} & a_{23} \\
a_{31} & a_{32} & a_{33}
\end{pmatrix}
\]

\section{Tables}
Below is a sample table showing data:

\begin{table}[h]
\centering
\begin{tabular}{@{}lrrr@{}}
\toprule
\textbf{Category} & \textbf{Value 1} & \textbf{Value 2} & \textbf{Total} \\
\midrule
Item A & 10.5 & 20.3 & 30.8 \\
Item B & 15.2 & 25.1 & 40.3 \\
Item C & 22.1 & 12.7 & 34.8 \\
\bottomrule
\end{tabular}
\caption{Sample Data Table}
\label{tab:sample}
\end{table}

\section{Lists}
\subsection{Bullet Points}
\begin{itemize}
  \item First item
  \item Second item
  \item Third item with sub-items:
    \begin{itemize}
      \item Sub-item 1
      \item Sub-item 2
    \end{itemize}
\end{itemize}

\subsection{Numbered List}
\begin{enumerate}
  \item First step
  \item Second step
  \item Third step
\end{enumerate}

\section{Citations and References}
When you use the citation search feature, references will be automatically managed. For example, you might cite a paper like this: \cite{example2024}. All citations will be automatically added to your \texttt{cite.bib} file.

\section{Conclusion}
This template demonstrates basic \LaTeX{} features including mathematical formulas, tables, and lists. The bibliography system is now set up to work seamlessly with Bibby AI's citation management. Feel free to modify it for your needs!

\bibliographystyle{plain}
\bibliography{cite}

\end{document}\documentclass{article}
\usepackage[utf8]{inputenc}
\usepackage{booktabs}
\usepackage{amsmath}
\usepackage{natbib}

\title{My Document}
\author{Your Name}
\date{\today}

\begin{document}

\maketitle

\section{Introduction}
This is a sample document demonstrating various \LaTeX{} features including mathematical equations, tables, and formatting. You can cite papers using the citation search feature, and they will be automatically added to your \texttt{cite.bib} file.

\section{Mathematical Expressions}
Here are some example equations:

\subsection{Quadratic Formula}
The quadratic formula is given by:
\[x = \frac{-b \pm \sqrt{b^2 - 4ac}}{2a}\]

\subsection{Matrix Example}
A sample matrix:
\[
\begin{pmatrix}
a_{11} & a_{12} & a_{13} \\
a_{21} & a_{22} & a_{23} \\
a_{31} & a_{32} & a_{33}
\end{pmatrix}
\]

\section{Tables}
Below is a sample table showing data:

\begin{table}[h]
\centering
\begin{tabular}{@{}lrrr@{}}
\toprule
\textbf{Category} & \textbf{Value 1} & \textbf{Value 2} & \textbf{Total} \\
\midrule
Item A & 10.5 & 20.3 & 30.8 \\
Item B & 15.2 & 25.1 & 40.3 \\
Item C & 22.1 & 12.7 & 34.8 \\
\bottomrule
\end{tabular}
\caption{Sample Data Table}
\label{tab:sample}
\end{table}

\section{Lists}
\subsection{Bullet Points}
\begin{itemize}
  \item First item
  \item Second item
  \item Third item with sub-items:
    \begin{itemize}
      \item Sub-item 1
      \item Sub-item 2
    \end{itemize}
\end{itemize}

\subsection{Numbered List}
\begin{enumerate}
  \item First step
  \item Second step
  \item Third step
\end{enumerate}

\section{Citations and References}
When you use the citation search feature, references will be automatically managed. For example, you might cite a paper like this: \cite{example2024}. All citations will be automatically added to your \texttt{cite.bib} file.

\section{Conclusion}
This template demonstrates basic \LaTeX{} features including mathematical formulas, tables, and lists. The bibliography system is now set up to work seamlessly with Bibby AI's citation management. Feel free to modify it for your needs!

\bibliographystyle{plain}
\bibliography{cite}

\end{document}\documentclass{article}
\usepackage[utf8]{inputenc}
\usepackage{booktabs}
\usepackage{amsmath}
\usepackage{natbib}

\title{My Document}
\author{Your Name}
\date{\today}

\begin{document}

\maketitle

\section{Introduction}
This is a sample document demonstrating various \LaTeX{} features including mathematical equations, tables, and formatting. You can cite papers using the citation search feature, and they will be automatically added to your \texttt{cite.bib} file.

\section{Mathematical Expressions}
Here are some example equations:

\subsection{Quadratic Formula}
The quadratic formula is given by:
\[x = \frac{-b \pm \sqrt{b^2 - 4ac}}{2a}\]

\subsection{Matrix Example}
A sample matrix:
\[
\begin{pmatrix}
a_{11} & a_{12} & a_{13} \\
a_{21} & a_{22} & a_{23} \\
a_{31} & a_{32} & a_{33}
\end{pmatrix}
\]

\section{Tables}
Below is a sample table showing data:

\begin{table}[h]
\centering
\begin{tabular}{@{}lrrr@{}}
\toprule
\textbf{Category} & \textbf{Value 1} & \textbf{Value 2} & \textbf{Total} \\
\midrule
Item A & 10.5 & 20.3 & 30.8 \\
Item B & 15.2 & 25.1 & 40.3 \\
Item C & 22.1 & 12.7 & 34.8 \\
\bottomrule
\end{tabular}
\caption{Sample Data Table}
\label{tab:sample}
\end{table}

\section{Lists}
\subsection{Bullet Points}
\begin{itemize}
  \item First item
  \item Second item
  \item Third item with sub-items:
    \begin{itemize}
      \item Sub-item 1
      \item Sub-item 2
    \end{itemize}
\end{itemize}

\subsection{Numbered List}
\begin{enumerate}
  \item First step
  \item Second step
  \item Third step
\end{enumerate}

\section{Citations and References}
When you use the citation search feature, references will be automatically managed. For example, you might cite a paper like this: \cite{example2024}. All citations will be automatically added to your \texttt{cite.bib} file.

\section{Conclusion}
This template demonstrates basic \LaTeX{} features including mathematical formulas, tables, and lists. The bibliography system is now set up to work seamlessly with Bibby AI's citation management. Feel free to modify it for your needs!

\bibliographystyle{plain}
\bibliography{cite}

\end{document}\documentclass{article}
\usepackage[utf8]{inputenc}
\usepackage{booktabs}
\usepackage{amsmath}
\usepackage{natbib}

\title{My Document}
\author{Your Name}
\date{\today}

\begin{document}

\maketitle

\section{Introduction}
This is a sample document demonstrating various \LaTeX{} features including mathematical equations, tables, and formatting. You can cite papers using the citation search feature, and they will be automatically added to your \texttt{cite.bib} file.

\section{Mathematical Expressions}
Here are some example equations:

\subsection{Quadratic Formula}
The quadratic formula is given by:
\[x = \frac{-b \pm \sqrt{b^2 - 4ac}}{2a}\]

\subsection{Matrix Example}
A sample matrix:
\[
\begin{pmatrix}
a_{11} & a_{12} & a_{13} \\
a_{21} & a_{22} & a_{23} \\
a_{31} & a_{32} & a_{33}
\end{pmatrix}
\]

\section{Tables}
Below is a sample table showing data:

\begin{table}[h]
\centering
\begin{tabular}{@{}lrrr@{}}
\toprule
\textbf{Category} & \textbf{Value 1} & \textbf{Value 2} & \textbf{Total} \\
\midrule
Item A & 10.5 & 20.3 & 30.8 \\
Item B & 15.2 & 25.1 & 40.3 \\
Item C & 22.1 & 12.7 & 34.8 \\
\bottomrule
\end{tabular}
\caption{Sample Data Table}
\label{tab:sample}
\end{table}

\section{Lists}
\subsection{Bullet Points}
\begin{itemize}
  \item First item
  \item Second item
  \item Third item with sub-items:
    \begin{itemize}
      \item Sub-item 1
      \item Sub-item 2
    \end{itemize}
\end{itemize}

\subsection{Numbered List}
\begin{enumerate}
  \item First step
  \item Second step
  \item Third step
\end{enumerate}

\section{Citations and References}
When you use the citation search feature, references will be automatically managed. For example, you might cite a paper like this: \cite{example2024}. All citations will be automatically added to your \texttt{cite.bib} file.

\section{Conclusion}
This template demonstrates basic \LaTeX{} features including mathematical formulas, tables, and lists. The bibliography system is now set up to work seamlessly with Bibby AI's citation management. Feel free to modify it for your needs!

\bibliographystyle{plain}
\bibliography{cite}

\end{document}\documentclass{article}
\usepackage[utf8]{inputenc}
\usepackage{booktabs}
\usepackage{amsmath}
\usepackage{natbib}

\title{My Document}
\author{Your Name}
\date{\today}

\begin{document}

\maketitle

\section{Introduction}
This is a sample document demonstrating various \LaTeX{} features including mathematical equations, tables, and formatting. You can cite papers using the citation search feature, and they will be automatically added to your \texttt{cite.bib} file.

\section{Mathematical Expressions}
Here are some example equations:

\subsection{Quadratic Formula}
The quadratic formula is given by:
\[x = \frac{-b \pm \sqrt{b^2 - 4ac}}{2a}\]

\subsection{Matrix Example}
A sample matrix:
\[
\begin{pmatrix}
a_{11} & a_{12} & a_{13} \\
a_{21} & a_{22} & a_{23} \\
a_{31} & a_{32} & a_{33}
\end{pmatrix}
\]

\section{Tables}
Below is a sample table showing data:

\begin{table}[h]
\centering
\begin{tabular}{@{}lrrr@{}}
\toprule
\textbf{Category} & \textbf{Value 1} & \textbf{Value 2} & \textbf{Total} \\
\midrule
Item A & 10.5 & 20.3 & 30.8 \\
Item B & 15.2 & 25.1 & 40.3 \\
Item C & 22.1 & 12.7 & 34.8 \\
\bottomrule
\end{tabular}
\caption{Sample Data Table}
\label{tab:sample}
\end{table}

\section{Lists}
\subsection{Bullet Points}
\begin{itemize}
  \item First item
  \item Second item
  \item Third item with sub-items:
    \begin{itemize}
      \item Sub-item 1
      \item Sub-item 2
    \end{itemize}
\end{itemize}

\subsection{Numbered List}
\begin{enumerate}
  \item First step
  \item Second step
  \item Third step
\end{enumerate}

\section{Citations and References}
When you use the citation search feature, references will be automatically managed. For example, you might cite a paper like this: \cite{example2024}. All citations will be automatically added to your \texttt{cite.bib} file.

\section{Conclusion}
This template demonstrates basic \LaTeX{} features including mathematical formulas, tables, and lists. The bibliography system is now set up to work seamlessly with Bibby AI's citation management. Feel free to modify it for your needs!

\bibliographystyle{plain}
\bibliography{cite}

\end{document}\documentclass{article}
\usepackage[utf8]{inputenc}
\usepackage{booktabs}
\usepackage{amsmath}
\usepackage{natbib}

\title{My Document}
\author{Your Name}
\date{\today}

\begin{document}

\maketitle

\section{Introduction}
This is a sample document demonstrating various \LaTeX{} features including mathematical equations, tables, and formatting. You can cite papers using the citation search feature, and they will be automatically added to your \texttt{cite.bib} file.

\section{Mathematical Expressions}
Here are some example equations:

\subsection{Quadratic Formula}
The quadratic formula is given by:
\[x = \frac{-b \pm \sqrt{b^2 - 4ac}}{2a}\]

\subsection{Matrix Example}
A sample matrix:
\[
\begin{pmatrix}
a_{11} & a_{12} & a_{13} \\
a_{21} & a_{22} & a_{23} \\
a_{31} & a_{32} & a_{33}
\end{pmatrix}
\]

\section{Tables}
Below is a sample table showing data:

\begin{table}[h]
\centering
\begin{tabular}{@{}lrrr@{}}
\toprule
\textbf{Category} & \textbf{Value 1} & \textbf{Value 2} & \textbf{Total} \\
\midrule
Item A & 10.5 & 20.3 & 30.8 \\
Item B & 15.2 & 25.1 & 40.3 \\
Item C & 22.1 & 12.7 & 34.8 \\
\bottomrule
\end{tabular}
\caption{Sample Data Table}
\label{tab:sample}
\end{table}

\section{Lists}
\subsection{Bullet Points}
\begin{itemize}
  \item First item
  \item Second item
  \item Third item with sub-items:
    \begin{itemize}
      \item Sub-item 1
      \item Sub-item 2
    \end{itemize}
\end{itemize}

\subsection{Numbered List}
\begin{enumerate}
  \item First step
  \item Second step
  \item Third step
\end{enumerate}

\section{Citations and References}
When you use the citation search feature, references will be automatically managed. For example, you might cite a paper like this: \cite{example2024}. All citations will be automatically added to your \texttt{cite.bib} file.

\section{Conclusion}
This template demonstrates basic \LaTeX{} features including mathematical formulas, tables, and lists. The bibliography system is now set up to work seamlessly with Bibby AI's citation management. Feel free to modify it for your needs!

\bibliographystyle{plain}
\bibliography{cite}

\end{document}\documentclass{article}
\usepackage[utf8]{inputenc}
\usepackage{booktabs}
\usepackage{amsmath}
\usepackage{natbib}

\title{My Document}
\author{Your Name}
\date{\today}

\begin{document}

\maketitle

\section{Introduction}
This is a sample document demonstrating various \LaTeX{} features including mathematical equations, tables, and formatting. You can cite papers using the citation search feature, and they will be automatically added to your \texttt{cite.bib} file.

\section{Mathematical Expressions}
Here are some example equations:

\subsection{Quadratic Formula}
The quadratic formula is given by:
\[x = \frac{-b \pm \sqrt{b^2 - 4ac}}{2a}\]

\subsection{Matrix Example}
A sample matrix:
\[
\begin{pmatrix}
a_{11} & a_{12} & a_{13} \\
a_{21} & a_{22} & a_{23} \\
a_{31} & a_{32} & a_{33}
\end{pmatrix}
\]

\section{Tables}
Below is a sample table showing data:

\begin{table}[h]
\centering
\begin{tabular}{@{}lrrr@{}}
\toprule
\textbf{Category} & \textbf{Value 1} & \textbf{Value 2} & \textbf{Total} \\
\midrule
Item A & 10.5 & 20.3 & 30.8 \\
Item B & 15.2 & 25.1 & 40.3 \\
Item C & 22.1 & 12.7 & 34.8 \\
\bottomrule
\end{tabular}
\caption{Sample Data Table}
\label{tab:sample}
\end{table}

\section{Lists}
\subsection{Bullet Points}
\begin{itemize}
  \item First item
  \item Second item
  \item Third item with sub-items:
    \begin{itemize}
      \item Sub-item 1
      \item Sub-item 2
    \end{itemize}
\end{itemize}

\subsection{Numbered List}
\begin{enumerate}
  \item First step
  \item Second step
  \item Third step
\end{enumerate}

\section{Citations and References}
When you use the citation search feature, references will be automatically managed. For example, you might cite a paper like this: \cite{example2024}. All citations will be automatically added to your \texttt{cite.bib} file.

\section{Conclusion}
This template demonstrates basic \LaTeX{} features including mathematical formulas, tables, and lists. The bibliography system is now set up to work seamlessly with Bibby AI's citation management. Feel free to modify it for your needs!

\bibliographystyle{plain}
\bibliography{cite}

\end{document}\documentclass{article}
\usepackage[utf8]{inputenc}
\usepackage{booktabs}
\usepackage{amsmath}
\usepackage{natbib}

\title{My Document}
\author{Your Name}
\date{\today}

\begin{document}

\maketitle

\section{Introduction}
This is a sample document demonstrating various \LaTeX{} features including mathematical equations, tables, and formatting. You can cite papers using the citation search feature, and they will be automatically added to your \texttt{cite.bib} file.

\section{Mathematical Expressions}
Here are some example equations:

\subsection{Quadratic Formula}
The quadratic formula is given by:
\[x = \frac{-b \pm \sqrt{b^2 - 4ac}}{2a}\]

\subsection{Matrix Example}
A sample matrix:
\[
\begin{pmatrix}
a_{11} & a_{12} & a_{13} \\
a_{21} & a_{22} & a_{23} \\
a_{31} & a_{32} & a_{33}
\end{pmatrix}
\]

\section{Tables}
Below is a sample table showing data:

\begin{table}[h]
\centering
\begin{tabular}{@{}lrrr@{}}
\toprule
\textbf{Category} & \textbf{Value 1} & \textbf{Value 2} & \textbf{Total} \\
\midrule
Item A & 10.5 & 20.3 & 30.8 \\
Item B & 15.2 & 25.1 & 40.3 \\
Item C & 22.1 & 12.7 & 34.8 \\
\bottomrule
\end{tabular}
\caption{Sample Data Table}
\label{tab:sample}
\end{table}

\section{Lists}
\subsection{Bullet Points}
\begin{itemize}
  \item First item
  \item Second item
  \item Third item with sub-items:
    \begin{itemize}
      \item Sub-item 1
      \item Sub-item 2
    \end{itemize}
\end{itemize}

\subsection{Numbered List}
\begin{enumerate}
  \item First step
  \item Second step
  \item Third step
\end{enumerate}

\section{Citations and References}
When you use the citation search feature, references will be automatically managed. For example, you might cite a paper like this: \cite{example2024}. All citations will be automatically added to your \texttt{cite.bib} file.

\section{Conclusion}
This template demonstrates basic \LaTeX{} features including mathematical formulas, tables, and lists. The bibliography system is now set up to work seamlessly with Bibby AI's citation management. Feel free to modify it for your needs!

\bibliographystyle{plain}
\bibliography{cite}

\end{document}\documentclass{article}
\usepackage[utf8]{inputenc}
\usepackage{booktabs}
\usepackage{amsmath}
\usepackage{natbib}

\title{My Document}
\author{Your Name}
\date{\today}

\begin{document}

\maketitle

\section{Introduction}
This is a sample document demonstrating various \LaTeX{} features including mathematical equations, tables, and formatting. You can cite papers using the citation search feature, and they will be automatically added to your \texttt{cite.bib} file.

\section{Mathematical Expressions}
Here are some example equations:

\subsection{Quadratic Formula}
The quadratic formula is given by:
\[x = \frac{-b \pm \sqrt{b^2 - 4ac}}{2a}\]

\subsection{Matrix Example}
A sample matrix:
\[
\begin{pmatrix}
a_{11} & a_{12} & a_{13} \\
a_{21} & a_{22} & a_{23} \\
a_{31} & a_{32} & a_{33}
\end{pmatrix}
\]

\section{Tables}
Below is a sample table showing data:

\begin{table}[h]
\centering
\begin{tabular}{@{}lrrr@{}}
\toprule
\textbf{Category} & \textbf{Value 1} & \textbf{Value 2} & \textbf{Total} \\
\midrule
Item A & 10.5 & 20.3 & 30.8 \\
Item B & 15.2 & 25.1 & 40.3 \\
Item C & 22.1 & 12.7 & 34.8 \\
\bottomrule
\end{tabular}
\caption{Sample Data Table}
\label{tab:sample}
\end{table}

\section{Lists}
\subsection{Bullet Points}
\begin{itemize}
  \item First item
  \item Second item
  \item Third item with sub-items:
    \begin{itemize}
      \item Sub-item 1
      \item Sub-item 2
    \end{itemize}
\end{itemize}

\subsection{Numbered List}
\begin{enumerate}
  \item First step
  \item Second step
  \item Third step
\end{enumerate}

\section{Citations and References}
When you use the citation search feature, references will be automatically managed. For example, you might cite a paper like this: \cite{example2024}. All citations will be automatically added to your \texttt{cite.bib} file.

\section{Conclusion}
This template demonstrates basic \LaTeX{} features including mathematical formulas, tables, and lists. The bibliography system is now set up to work seamlessly with Bibby AI's citation management. Feel free to modify it for your needs!

\bibliographystyle{plain}
\bibliography{cite}

\end{document}\documentclass{article}
\usepackage[utf8]{inputenc}
\usepackage{booktabs}
\usepackage{amsmath}
\usepackage{natbib}

\title{My Document}
\author{Your Name}
\date{\today}

\begin{document}

\maketitle

\section{Introduction}
This is a sample document demonstrating various \LaTeX{} features including mathematical equations, tables, and formatting. You can cite papers using the citation search feature, and they will be automatically added to your \texttt{cite.bib} file.

\section{Mathematical Expressions}
Here are some example equations:

\subsection{Quadratic Formula}
The quadratic formula is given by:
\[x = \frac{-b \pm \sqrt{b^2 - 4ac}}{2a}\]

\subsection{Matrix Example}
A sample matrix:
\[
\begin{pmatrix}
a_{11} & a_{12} & a_{13} \\
a_{21} & a_{22} & a_{23} \\
a_{31} & a_{32} & a_{33}
\end{pmatrix}
\]

\section{Tables}
Below is a sample table showing data:

\begin{table}[h]
\centering
\begin{tabular}{@{}lrrr@{}}
\toprule
\textbf{Category} & \textbf{Value 1} & \textbf{Value 2} & \textbf{Total} \\
\midrule
Item A & 10.5 & 20.3 & 30.8 \\
Item B & 15.2 & 25.1 & 40.3 \\
Item C & 22.1 & 12.7 & 34.8 \\
\bottomrule
\end{tabular}
\caption{Sample Data Table}
\label{tab:sample}
\end{table}

\section{Lists}
\subsection{Bullet Points}
\begin{itemize}
  \item First item
  \item Second item
  \item Third item with sub-items:
    \begin{itemize}
      \item Sub-item 1
      \item Sub-item 2
    \end{itemize}
\end{itemize}

\subsection{Numbered List}
\begin{enumerate}
  \item First step
  \item Second step
  \item Third step
\end{enumerate}

\section{Citations and References}
When you use the citation search feature, references will be automatically managed. For example, you might cite a paper like this: \cite{example2024}. All citations will be automatically added to your \texttt{cite.bib} file.

\section{Conclusion}
This template demonstrates basic \LaTeX{} features including mathematical formulas, tables, and lists. The bibliography system is now set up to work seamlessly with Bibby AI's citation management. Feel free to modify it for your needs!

\bibliographystyle{plain}
\bibliography{cite}

\end{document}\documentclass{article}
\usepackage[utf8]{inputenc}
\usepackage{booktabs}
\usepackage{amsmath}
\usepackage{natbib}

\title{My Document}
\author{Your Name}
\date{\today}

\begin{document}

\maketitle

\section{Introduction}
This is a sample document demonstrating various \LaTeX{} features including mathematical equations, tables, and formatting. You can cite papers using the citation search feature, and they will be automatically added to your \texttt{cite.bib} file.

\section{Mathematical Expressions}
Here are some example equations:

\subsection{Quadratic Formula}
The quadratic formula is given by:
\[x = \frac{-b \pm \sqrt{b^2 - 4ac}}{2a}\]

\subsection{Matrix Example}
A sample matrix:
\[
\begin{pmatrix}
a_{11} & a_{12} & a_{13} \\
a_{21} & a_{22} & a_{23} \\
a_{31} & a_{32} & a_{33}
\end{pmatrix}
\]

\section{Tables}
Below is a sample table showing data:

\begin{table}[h]
\centering
\begin{tabular}{@{}lrrr@{}}
\toprule
\textbf{Category} & \textbf{Value 1} & \textbf{Value 2} & \textbf{Total} \\
\midrule
Item A & 10.5 & 20.3 & 30.8 \\
Item B & 15.2 & 25.1 & 40.3 \\
Item C & 22.1 & 12.7 & 34.8 \\
\bottomrule
\end{tabular}
\caption{Sample Data Table}
\label{tab:sample}
\end{table}

\section{Lists}
\subsection{Bullet Points}
\begin{itemize}
  \item First item
  \item Second item
  \item Third item with sub-items:
    \begin{itemize}
      \item Sub-item 1
      \item Sub-item 2
    \end{itemize}
\end{itemize}

\subsection{Numbered List}
\begin{enumerate}
  \item First step
  \item Second step
  \item Third step
\end{enumerate}

\section{Citations and References}
When you use the citation search feature, references will be automatically managed. For example, you might cite a paper like this: \cite{example2024}. All citations will be automatically added to your \texttt{cite.bib} file.

\section{Conclusion}
This template demonstrates basic \LaTeX{} features including mathematical formulas, tables, and lists. The bibliography system is now set up to work seamlessly with Bibby AI's citation management. Feel free to modify it for your needs!

\bibliographystyle{plain}
\bibliography{cite}

\end{document}\documentclass{article}
\usepackage[utf8]{inputenc}
\usepackage{booktabs}
\usepackage{amsmath}
\usepackage{natbib}

\title{My Document}
\author{Your Name}
\date{\today}

\begin{document}

\maketitle

\section{Introduction}
This is a sample document demonstrating various \LaTeX{} features including mathematical equations, tables, and formatting. You can cite papers using the citation search feature, and they will be automatically added to your \texttt{cite.bib} file.

\section{Mathematical Expressions}
Here are some example equations:

\subsection{Quadratic Formula}
The quadratic formula is given by:
\[x = \frac{-b \pm \sqrt{b^2 - 4ac}}{2a}\]

\subsection{Matrix Example}
A sample matrix:
\[
\begin{pmatrix}
a_{11} & a_{12} & a_{13} \\
a_{21} & a_{22} & a_{23} \\
a_{31} & a_{32} & a_{33}
\end{pmatrix}
\]

\section{Tables}
Below is a sample table showing data:

\begin{table}[h]
\centering
\begin{tabular}{@{}lrrr@{}}
\toprule
\textbf{Category} & \textbf{Value 1} & \textbf{Value 2} & \textbf{Total} \\
\midrule
Item A & 10.5 & 20.3 & 30.8 \\
Item B & 15.2 & 25.1 & 40.3 \\
Item C & 22.1 & 12.7 & 34.8 \\
\bottomrule
\end{tabular}
\caption{Sample Data Table}
\label{tab:sample}
\end{table}

\section{Lists}
\subsection{Bullet Points}
\begin{itemize}
  \item First item
  \item Second item
  \item Third item with sub-items:
    \begin{itemize}
      \item Sub-item 1
      \item Sub-item 2
    \end{itemize}
\end{itemize}

\subsection{Numbered List}
\begin{enumerate}
  \item First step
  \item Second step
  \item Third step
\end{enumerate}

\section{Citations and References}
When you use the citation search feature, references will be automatically managed. For example, you might cite a paper like this: \cite{example2024}. All citations will be automatically added to your \texttt{cite.bib} file.

\section{Conclusion}
This template demonstrates basic \LaTeX{} features including mathematical formulas, tables, and lists. The bibliography system is now set up to work seamlessly with Bibby AI's citation management. Feel free to modify it for your needs!

\bibliographystyle{plain}
\bibliography{cite}

\end{document}\documentclass{article}
\usepackage[utf8]{inputenc}
\usepackage{booktabs}
\usepackage{amsmath}
\usepackage{natbib}

\title{My Document}
\author{Your Name}
\date{\today}

\begin{document}

\maketitle

\section{Introduction}
This is a sample document demonstrating various \LaTeX{} features including mathematical equations, tables, and formatting. You can cite papers using the citation search feature, and they will be automatically added to your \texttt{cite.bib} file.

\section{Mathematical Expressions}
Here are some example equations:

\subsection{Quadratic Formula}
The quadratic formula is given by:
\[x = \frac{-b \pm \sqrt{b^2 - 4ac}}{2a}\]

\subsection{Matrix Example}
A sample matrix:
\[
\begin{pmatrix}
a_{11} & a_{12} & a_{13} \\
a_{21} & a_{22} & a_{23} \\
a_{31} & a_{32} & a_{33}
\end{pmatrix}
\]

\section{Tables}
Below is a sample table showing data:

\begin{table}[h]
\centering
\begin{tabular}{@{}lrrr@{}}
\toprule
\textbf{Category} & \textbf{Value 1} & \textbf{Value 2} & \textbf{Total} \\
\midrule
Item A & 10.5 & 20.3 & 30.8 \\
Item B & 15.2 & 25.1 & 40.3 \\
Item C & 22.1 & 12.7 & 34.8 \\
\bottomrule
\end{tabular}
\caption{Sample Data Table}
\label{tab:sample}
\end{table}

\section{Lists}
\subsection{Bullet Points}
\begin{itemize}
  \item First item
  \item Second item
  \item Third item with sub-items:
    \begin{itemize}
      \item Sub-item 1
      \item Sub-item 2
    \end{itemize}
\end{itemize}

\subsection{Numbered List}
\begin{enumerate}
  \item First step
  \item Second step
  \item Third step
\end{enumerate}

\section{Citations and References}
When you use the citation search feature, references will be automatically managed. For example, you might cite a paper like this: \cite{example2024}. All citations will be automatically added to your \texttt{cite.bib} file.

\section{Conclusion}
This template demonstrates basic \LaTeX{} features including mathematical formulas, tables, and lists. The bibliography system is now set up to work seamlessly with Bibby AI's citation management. Feel free to modify it for your needs!

\bibliographystyle{plain}
\bibliography{cite}

\end{document}\documentclass{article}
\usepackage[utf8]{inputenc}
\usepackage{booktabs}
\usepackage{amsmath}
\usepackage{natbib}

\title{My Document}
\author{Your Name}
\date{\today}

\begin{document}

\maketitle

\section{Introduction}
This is a sample document demonstrating various \LaTeX{} features including mathematical equations, tables, and formatting. You can cite papers using the citation search feature, and they will be automatically added to your \texttt{cite.bib} file.

\section{Mathematical Expressions}
Here are some example equations:

\subsection{Quadratic Formula}
The quadratic formula is given by:
\[x = \frac{-b \pm \sqrt{b^2 - 4ac}}{2a}\]

\subsection{Matrix Example}
A sample matrix:
\[
\begin{pmatrix}
a_{11} & a_{12} & a_{13} \\
a_{21} & a_{22} & a_{23} \\
a_{31} & a_{32} & a_{33}
\end{pmatrix}
\]

\section{Tables}
Below is a sample table showing data:

\begin{table}[h]
\centering
\begin{tabular}{@{}lrrr@{}}
\toprule
\textbf{Category} & \textbf{Value 1} & \textbf{Value 2} & \textbf{Total} \\
\midrule
Item A & 10.5 & 20.3 & 30.8 \\
Item B & 15.2 & 25.1 & 40.3 \\
Item C & 22.1 & 12.7 & 34.8 \\
\bottomrule
\end{tabular}
\caption{Sample Data Table}
\label{tab:sample}
\end{table}

\section{Lists}
\subsection{Bullet Points}
\begin{itemize}
  \item First item
  \item Second item
  \item Third item with sub-items:
    \begin{itemize}
      \item Sub-item 1
      \item Sub-item 2
    \end{itemize}
\end{itemize}

\subsection{Numbered List}
\begin{enumerate}
  \item First step
  \item Second step
  \item Third step
\end{enumerate}

\section{Citations and References}
When you use the citation search feature, references will be automatically managed. For example, you might cite a paper like this: \cite{example2024}. All citations will be automatically added to your \texttt{cite.bib} file.

\section{Conclusion}
This template demonstrates basic \LaTeX{} features including mathematical formulas, tables, and lists. The bibliography system is now set up to work seamlessly with Bibby AI's citation management. Feel free to modify it for your needs!

\bibliographystyle{plain}
\bibliography{cite}

\end{document}\documentclass{article}
\usepackage[utf8]{inputenc}
\usepackage{booktabs}
\usepackage{amsmath}
\usepackage{natbib}

\title{My Document}
\author{Your Name}
\date{\today}

\begin{document}

\maketitle

\section{Introduction}
This is a sample document demonstrating various \LaTeX{} features including mathematical equations, tables, and formatting. You can cite papers using the citation search feature, and they will be automatically added to your \texttt{cite.bib} file.

\section{Mathematical Expressions}
Here are some example equations:

\subsection{Quadratic Formula}
The quadratic formula is given by:
\[x = \frac{-b \pm \sqrt{b^2 - 4ac}}{2a}\]

\subsection{Matrix Example}
A sample matrix:
\[
\begin{pmatrix}
a_{11} & a_{12} & a_{13} \\
a_{21} & a_{22} & a_{23} \\
a_{31} & a_{32} & a_{33}
\end{pmatrix}
\]

\section{Tables}
Below is a sample table showing data:

\begin{table}[h]
\centering
\begin{tabular}{@{}lrrr@{}}
\toprule
\textbf{Category} & \textbf{Value 1} & \textbf{Value 2} & \textbf{Total} \\
\midrule
Item A & 10.5 & 20.3 & 30.8 \\
Item B & 15.2 & 25.1 & 40.3 \\
Item C & 22.1 & 12.7 & 34.8 \\
\bottomrule
\end{tabular}
\caption{Sample Data Table}
\label{tab:sample}
\end{table}

\section{Lists}
\subsection{Bullet Points}
\begin{itemize}
  \item First item
  \item Second item
  \item Third item with sub-items:
    \begin{itemize}
      \item Sub-item 1
      \item Sub-item 2
    \end{itemize}
\end{itemize}

\subsection{Numbered List}
\begin{enumerate}
  \item First step
  \item Second step
  \item Third step
\end{enumerate}

\section{Citations and References}
When you use the citation search feature, references will be automatically managed. For example, you might cite a paper like this: \cite{example2024}. All citations will be automatically added to your \texttt{cite.bib} file.

\section{Conclusion}
This template demonstrates basic \LaTeX{} features including mathematical formulas, tables, and lists. The bibliography system is now set up to work seamlessly with Bibby AI's citation management. Feel free to modify it for your needs!

\bibliographystyle{plain}
\bibliography{cite}

\end{document}\documentclass{article}
\usepackage[utf8]{inputenc}
\usepackage{booktabs}
\usepackage{amsmath}
\usepackage{natbib}

\title{My Document}
\author{Your Name}
\date{\today}

\begin{document}

\maketitle

\section{Introduction}
This is a sample document demonstrating various \LaTeX{} features including mathematical equations, tables, and formatting. You can cite papers using the citation search feature, and they will be automatically added to your \texttt{cite.bib} file.

\section{Mathematical Expressions}
Here are some example equations:

\subsection{Quadratic Formula}
The quadratic formula is given by:
\[x = \frac{-b \pm \sqrt{b^2 - 4ac}}{2a}\]

\subsection{Matrix Example}
A sample matrix:
\[
\begin{pmatrix}
a_{11} & a_{12} & a_{13} \\
a_{21} & a_{22} & a_{23} \\
a_{31} & a_{32} & a_{33}
\end{pmatrix}
\]

\section{Tables}
Below is a sample table showing data:

\begin{table}[h]
\centering
\begin{tabular}{@{}lrrr@{}}
\toprule
\textbf{Category} & \textbf{Value 1} & \textbf{Value 2} & \textbf{Total} \\
\midrule
Item A & 10.5 & 20.3 & 30.8 \\
Item B & 15.2 & 25.1 & 40.3 \\
Item C & 22.1 & 12.7 & 34.8 \\
\bottomrule
\end{tabular}
\caption{Sample Data Table}
\label{tab:sample}
\end{table}

\section{Lists}
\subsection{Bullet Points}
\begin{itemize}
  \item First item
  \item Second item
  \item Third item with sub-items:
    \begin{itemize}
      \item Sub-item 1
      \item Sub-item 2
    \end{itemize}
\end{itemize}

\subsection{Numbered List}
\begin{enumerate}
  \item First step
  \item Second step
  \item Third step
\end{enumerate}

\section{Citations and References}
When you use the citation search feature, references will be automatically managed. For example, you might cite a paper like this: \cite{example2024}. All citations will be automatically added to your \texttt{cite.bib} file.

\section{Conclusion}
This template demonstrates basic \LaTeX{} features including mathematical formulas, tables, and lists. The bibliography system is now set up to work seamlessly with Bibby AI's citation management. Feel free to modify it for your needs!

\bibliographystyle{plain}
\bibliography{cite}

\end{document}\documentclass{article}
\usepackage[utf8]{inputenc}
\usepackage{booktabs}
\usepackage{amsmath}
\usepackage{natbib}

\title{My Document}
\author{Your Name}
\date{\today}

\begin{document}

\maketitle

\section{Introduction}
This is a sample document demonstrating various \LaTeX{} features including mathematical equations, tables, and formatting. You can cite papers using the citation search feature, and they will be automatically added to your \texttt{cite.bib} file.

\section{Mathematical Expressions}
Here are some example equations:

\subsection{Quadratic Formula}
The quadratic formula is given by:
\[x = \frac{-b \pm \sqrt{b^2 - 4ac}}{2a}\]

\subsection{Matrix Example}
A sample matrix:
\[
\begin{pmatrix}
a_{11} & a_{12} & a_{13} \\
a_{21} & a_{22} & a_{23} \\
a_{31} & a_{32} & a_{33}
\end{pmatrix}
\]

\section{Tables}
Below is a sample table showing data:

\begin{table}[h]
\centering
\begin{tabular}{@{}lrrr@{}}
\toprule
\textbf{Category} & \textbf{Value 1} & \textbf{Value 2} & \textbf{Total} \\
\midrule
Item A & 10.5 & 20.3 & 30.8 \\
Item B & 15.2 & 25.1 & 40.3 \\
Item C & 22.1 & 12.7 & 34.8 \\
\bottomrule
\end{tabular}
\caption{Sample Data Table}
\label{tab:sample}
\end{table}

\section{Lists}
\subsection{Bullet Points}
\begin{itemize}
  \item First item
  \item Second item
  \item Third item with sub-items:
    \begin{itemize}
      \item Sub-item 1
      \item Sub-item 2
    \end{itemize}
\end{itemize}

\subsection{Numbered List}
\begin{enumerate}
  \item First step
  \item Second step
  \item Third step
\end{enumerate}

\section{Citations and References}
When you use the citation search feature, references will be automatically managed. For example, you might cite a paper like this: \cite{example2024}. All citations will be automatically added to your \texttt{cite.bib} file.

\section{Conclusion}
This template demonstrates basic \LaTeX{} features including mathematical formulas, tables, and lists. The bibliography system is now set up to work seamlessly with Bibby AI's citation management. Feel free to modify it for your needs!

\bibliographystyle{plain}
\bibliography{cite}

\end{document}\documentclass{article}
\usepackage[utf8]{inputenc}
\usepackage{booktabs}
\usepackage{amsmath}
\usepackage{natbib}

\title{My Document}
\author{Your Name}
\date{\today}

\begin{document}

\maketitle

\section{Introduction}
This is a sample document demonstrating various \LaTeX{} features including mathematical equations, tables, and formatting. You can cite papers using the citation search feature, and they will be automatically added to your \texttt{cite.bib} file.

\section{Mathematical Expressions}
Here are some example equations:

\subsection{Quadratic Formula}
The quadratic formula is given by:
\[x = \frac{-b \pm \sqrt{b^2 - 4ac}}{2a}\]

\subsection{Matrix Example}
A sample matrix:
\[
\begin{pmatrix}
a_{11} & a_{12} & a_{13} \\
a_{21} & a_{22} & a_{23} \\
a_{31} & a_{32} & a_{33}
\end{pmatrix}
\]

\section{Tables}
Below is a sample table showing data:

\begin{table}[h]
\centering
\begin{tabular}{@{}lrrr@{}}
\toprule
\textbf{Category} & \textbf{Value 1} & \textbf{Value 2} & \textbf{Total} \\
\midrule
Item A & 10.5 & 20.3 & 30.8 \\
Item B & 15.2 & 25.1 & 40.3 \\
Item C & 22.1 & 12.7 & 34.8 \\
\bottomrule
\end{tabular}
\caption{Sample Data Table}
\label{tab:sample}
\end{table}

\section{Lists}
\subsection{Bullet Points}
\begin{itemize}
  \item First item
  \item Second item
  \item Third item with sub-items:
    \begin{itemize}
      \item Sub-item 1
      \item Sub-item 2
    \end{itemize}
\end{itemize}

\subsection{Numbered List}
\begin{enumerate}
  \item First step
  \item Second step
  \item Third step
\end{enumerate}

\section{Citations and References}
When you use the citation search feature, references will be automatically managed. For example, you might cite a paper like this: \cite{example2024}. All citations will be automatically added to your \texttt{cite.bib} file.

\section{Conclusion}
This template demonstrates basic \LaTeX{} features including mathematical formulas, tables, and lists. The bibliography system is now set up to work seamlessly with Bibby AI's citation management. Feel free to modify it for your needs!

\bibliographystyle{plain}
\bibliography{cite}

\end{document}\documentclass{article}
\usepackage[utf8]{inputenc}
\usepackage{booktabs}
\usepackage{amsmath}
\usepackage{natbib}

\title{My Document}
\author{Your Name}
\date{\today}

\begin{document}

\maketitle

\section{Introduction}
This is a sample document demonstrating various \LaTeX{} features including mathematical equations, tables, and formatting. You can cite papers using the citation search feature, and they will be automatically added to your \texttt{cite.bib} file.

\section{Mathematical Expressions}
Here are some example equations:

\subsection{Quadratic Formula}
The quadratic formula is given by:
\[x = \frac{-b \pm \sqrt{b^2 - 4ac}}{2a}\]

\subsection{Matrix Example}
A sample matrix:
\[
\begin{pmatrix}
a_{11} & a_{12} & a_{13} \\
a_{21} & a_{22} & a_{23} \\
a_{31} & a_{32} & a_{33}
\end{pmatrix}
\]

\section{Tables}
Below is a sample table showing data:

\begin{table}[h]
\centering
\begin{tabular}{@{}lrrr@{}}
\toprule
\textbf{Category} & \textbf{Value 1} & \textbf{Value 2} & \textbf{Total} \\
\midrule
Item A & 10.5 & 20.3 & 30.8 \\
Item B & 15.2 & 25.1 & 40.3 \\
Item C & 22.1 & 12.7 & 34.8 \\
\bottomrule
\end{tabular}
\caption{Sample Data Table}
\label{tab:sample}
\end{table}

\section{Lists}
\subsection{Bullet Points}
\begin{itemize}
  \item First item
  \item Second item
  \item Third item with sub-items:
    \begin{itemize}
      \item Sub-item 1
      \item Sub-item 2
    \end{itemize}
\end{itemize}

\subsection{Numbered List}
\begin{enumerate}
  \item First step
  \item Second step
  \item Third step
\end{enumerate}

\section{Citations and References}
When you use the citation search feature, references will be automatically managed. For example, you might cite a paper like this: \cite{example2024}. All citations will be automatically added to your \texttt{cite.bib} file.

\section{Conclusion}
This template demonstrates basic \LaTeX{} features including mathematical formulas, tables, and lists. The bibliography system is now set up to work seamlessly with Bibby AI's citation management. Feel free to modify it for your needs!

\bibliographystyle{plain}
\bibliography{cite}

\end{document}\documentclass{article}
\usepackage[utf8]{inputenc}
\usepackage{booktabs}
\usepackage{amsmath}
\usepackage{natbib}

\title{My Document}
\author{Your Name}
\date{\today}

\begin{document}

\maketitle

\section{Introduction}
This is a sample document demonstrating various \LaTeX{} features including mathematical equations, tables, and formatting. You can cite papers using the citation search feature, and they will be automatically added to your \texttt{cite.bib} file.

\section{Mathematical Expressions}
Here are some example equations:

\subsection{Quadratic Formula}
The quadratic formula is given by:
\[x = \frac{-b \pm \sqrt{b^2 - 4ac}}{2a}\]

\subsection{Matrix Example}
A sample matrix:
\[
\begin{pmatrix}
a_{11} & a_{12} & a_{13} \\
a_{21} & a_{22} & a_{23} \\
a_{31} & a_{32} & a_{33}
\end{pmatrix}
\]

\section{Tables}
Below is a sample table showing data:

\begin{table}[h]
\centering
\begin{tabular}{@{}lrrr@{}}
\toprule
\textbf{Category} & \textbf{Value 1} & \textbf{Value 2} & \textbf{Total} \\
\midrule
Item A & 10.5 & 20.3 & 30.8 \\
Item B & 15.2 & 25.1 & 40.3 \\
Item C & 22.1 & 12.7 & 34.8 \\
\bottomrule
\end{tabular}
\caption{Sample Data Table}
\label{tab:sample}
\end{table}

\section{Lists}
\subsection{Bullet Points}
\begin{itemize}
  \item First item
  \item Second item
  \item Third item with sub-items:
    \begin{itemize}
      \item Sub-item 1
      \item Sub-item 2
    \end{itemize}
\end{itemize}

\subsection{Numbered List}
\begin{enumerate}
  \item First step
  \item Second step
  \item Third step
\end{enumerate}

\section{Citations and References}
When you use the citation search feature, references will be automatically managed. For example, you might cite a paper like this: \cite{example2024}. All citations will be automatically added to your \texttt{cite.bib} file.

\section{Conclusion}
This template demonstrates basic \LaTeX{} features including mathematical formulas, tables, and lists. The bibliography system is now set up to work seamlessly with Bibby AI's citation management. Feel free to modify it for your needs!

\bibliographystyle{plain}
\bibliography{cite}

\end{document}\documentclass{article}
\usepackage[utf8]{inputenc}
\usepackage{booktabs}
\usepackage{amsmath}
\usepackage{natbib}

\title{My Document}
\author{Your Name}
\date{\today}

\begin{document}

\maketitle

\section{Introduction}
This is a sample document demonstrating various \LaTeX{} features including mathematical equations, tables, and formatting. You can cite papers using the citation search feature, and they will be automatically added to your \texttt{cite.bib} file.

\section{Mathematical Expressions}
Here are some example equations:

\subsection{Quadratic Formula}
The quadratic formula is given by:
\[x = \frac{-b \pm \sqrt{b^2 - 4ac}}{2a}\]

\subsection{Matrix Example}
A sample matrix:
\[
\begin{pmatrix}
a_{11} & a_{12} & a_{13} \\
a_{21} & a_{22} & a_{23} \\
a_{31} & a_{32} & a_{33}
\end{pmatrix}
\]

\section{Tables}
Below is a sample table showing data:

\begin{table}[h]
\centering
\begin{tabular}{@{}lrrr@{}}
\toprule
\textbf{Category} & \textbf{Value 1} & \textbf{Value 2} & \textbf{Total} \\
\midrule
Item A & 10.5 & 20.3 & 30.8 \\
Item B & 15.2 & 25.1 & 40.3 \\
Item C & 22.1 & 12.7 & 34.8 \\
\bottomrule
\end{tabular}
\caption{Sample Data Table}
\label{tab:sample}
\end{table}

\section{Lists}
\subsection{Bullet Points}
\begin{itemize}
  \item First item
  \item Second item
  \item Third item with sub-items:
    \begin{itemize}
      \item Sub-item 1
      \item Sub-item 2
    \end{itemize}
\end{itemize}

\subsection{Numbered List}
\begin{enumerate}
  \item First step
  \item Second step
  \item Third step
\end{enumerate}

\section{Citations and References}
When you use the citation search feature, references will be automatically managed. For example, you might cite a paper like this: \cite{example2024}. All citations will be automatically added to your \texttt{cite.bib} file.

\section{Conclusion}
This template demonstrates basic \LaTeX{} features including mathematical formulas, tables, and lists. The bibliography system is now set up to work seamlessly with Bibby AI's citation management. Feel free to modify it for your needs!

\bibliographystyle{plain}
\bibliography{cite}

\end{document}\documentclass{article}
\usepackage[utf8]{inputenc}
\usepackage{booktabs}
\usepackage{amsmath}
\usepackage{natbib}

\title{My Document}
\author{Your Name}
\date{\today}

\begin{document}

\maketitle

\section{Introduction}
This is a sample document demonstrating various \LaTeX{} features including mathematical equations, tables, and formatting. You can cite papers using the citation search feature, and they will be automatically added to your \texttt{cite.bib} file.

\section{Mathematical Expressions}
Here are some example equations:

\subsection{Quadratic Formula}
The quadratic formula is given by:
\[x = \frac{-b \pm \sqrt{b^2 - 4ac}}{2a}\]

\subsection{Matrix Example}
A sample matrix:
\[
\begin{pmatrix}
a_{11} & a_{12} & a_{13} \\
a_{21} & a_{22} & a_{23} \\
a_{31} & a_{32} & a_{33}
\end{pmatrix}
\]

\section{Tables}
Below is a sample table showing data:

\begin{table}[h]
\centering
\begin{tabular}{@{}lrrr@{}}
\toprule
\textbf{Category} & \textbf{Value 1} & \textbf{Value 2} & \textbf{Total} \\
\midrule
Item A & 10.5 & 20.3 & 30.8 \\
Item B & 15.2 & 25.1 & 40.3 \\
Item C & 22.1 & 12.7 & 34.8 \\
\bottomrule
\end{tabular}
\caption{Sample Data Table}
\label{tab:sample}
\end{table}

\section{Lists}
\subsection{Bullet Points}
\begin{itemize}
  \item First item
  \item Second item
  \item Third item with sub-items:
    \begin{itemize}
      \item Sub-item 1
      \item Sub-item 2
    \end{itemize}
\end{itemize}

\subsection{Numbered List}
\begin{enumerate}
  \item First step
  \item Second step
  \item Third step
\end{enumerate}

\section{Citations and References}
When you use the citation search feature, references will be automatically managed. For example, you might cite a paper like this: \cite{example2024}. All citations will be automatically added to your \texttt{cite.bib} file.

\section{Conclusion}
This template demonstrates basic \LaTeX{} features including mathematical formulas, tables, and lists. The bibliography system is now set up to work seamlessly with Bibby AI's citation management. Feel free to modify it for your needs!

\bibliographystyle{plain}
\bibliography{cite}

\end{document}\documentclass{article}
\usepackage[utf8]{inputenc}
\usepackage{booktabs}
\usepackage{amsmath}
\usepackage{natbib}

\title{My Document}
\author{Your Name}
\date{\today}

\begin{document}

\maketitle

\section{Introduction}
This is a sample document demonstrating various \LaTeX{} features including mathematical equations, tables, and formatting. You can cite papers using the citation search feature, and they will be automatically added to your \texttt{cite.bib} file.

\section{Mathematical Expressions}
Here are some example equations:

\subsection{Quadratic Formula}
The quadratic formula is given by:
\[x = \frac{-b \pm \sqrt{b^2 - 4ac}}{2a}\]

\subsection{Matrix Example}
A sample matrix:
\[
\begin{pmatrix}
a_{11} & a_{12} & a_{13} \\
a_{21} & a_{22} & a_{23} \\
a_{31} & a_{32} & a_{33}
\end{pmatrix}
\]

\section{Tables}
Below is a sample table showing data:

\begin{table}[h]
\centering
\begin{tabular}{@{}lrrr@{}}
\toprule
\textbf{Category} & \textbf{Value 1} & \textbf{Value 2} & \textbf{Total} \\
\midrule
Item A & 10.5 & 20.3 & 30.8 \\
Item B & 15.2 & 25.1 & 40.3 \\
Item C & 22.1 & 12.7 & 34.8 \\
\bottomrule
\end{tabular}
\caption{Sample Data Table}
\label{tab:sample}
\end{table}

\section{Lists}
\subsection{Bullet Points}
\begin{itemize}
  \item First item
  \item Second item
  \item Third item with sub-items:
    \begin{itemize}
      \item Sub-item 1
      \item Sub-item 2
    \end{itemize}
\end{itemize}

\subsection{Numbered List}
\begin{enumerate}
  \item First step
  \item Second step
  \item Third step
\end{enumerate}

\section{Citations and References}
When you use the citation search feature, references will be automatically managed. For example, you might cite a paper like this: \cite{example2024}. All citations will be automatically added to your \texttt{cite.bib} file.

\section{Conclusion}
This template demonstrates basic \LaTeX{} features including mathematical formulas, tables, and lists. The bibliography system is now set up to work seamlessly with Bibby AI's citation management. Feel free to modify it for your needs!

\bibliographystyle{plain}
\bibliography{cite}

\end{document}\documentclass{article}
\usepackage[utf8]{inputenc}
\usepackage{booktabs}
\usepackage{amsmath}
\usepackage{natbib}

\title{My Document}
\author{Your Name}
\date{\today}

\begin{document}

\maketitle

\section{Introduction}
This is a sample document demonstrating various \LaTeX{} features including mathematical equations, tables, and formatting. You can cite papers using the citation search feature, and they will be automatically added to your \texttt{cite.bib} file.

\section{Mathematical Expressions}
Here are some example equations:

\subsection{Quadratic Formula}
The quadratic formula is given by:
\[x = \frac{-b \pm \sqrt{b^2 - 4ac}}{2a}\]

\subsection{Matrix Example}
A sample matrix:
\[
\begin{pmatrix}
a_{11} & a_{12} & a_{13} \\
a_{21} & a_{22} & a_{23} \\
a_{31} & a_{32} & a_{33}
\end{pmatrix}
\]

\section{Tables}
Below is a sample table showing data:

\begin{table}[h]
\centering
\begin{tabular}{@{}lrrr@{}}
\toprule
\textbf{Category} & \textbf{Value 1} & \textbf{Value 2} & \textbf{Total} \\
\midrule
Item A & 10.5 & 20.3 & 30.8 \\
Item B & 15.2 & 25.1 & 40.3 \\
Item C & 22.1 & 12.7 & 34.8 \\
\bottomrule
\end{tabular}
\caption{Sample Data Table}
\label{tab:sample}
\end{table}

\section{Lists}
\subsection{Bullet Points}
\begin{itemize}
  \item First item
  \item Second item
  \item Third item with sub-items:
    \begin{itemize}
      \item Sub-item 1
      \item Sub-item 2
    \end{itemize}
\end{itemize}

\subsection{Numbered List}
\begin{enumerate}
  \item First step
  \item Second step
  \item Third step
\end{enumerate}

\section{Citations and References}
When you use the citation search feature, references will be automatically managed. For example, you might cite a paper like this: \cite{example2024}. All citations will be automatically added to your \texttt{cite.bib} file.

\section{Conclusion}
This template demonstrates basic \LaTeX{} features including mathematical formulas, tables, and lists. The bibliography system is now set up to work seamlessly with Bibby AI's citation management. Feel free to modify it for your needs!

\bibliographystyle{plain}
\bibliography{cite}

\end{document}\documentclass{article}
\usepackage[utf8]{inputenc}
\usepackage{booktabs}
\usepackage{amsmath}
\usepackage{natbib}

\title{My Document}
\author{Your Name}
\date{\today}

\begin{document}

\maketitle

\section{Introduction}
This is a sample document demonstrating various \LaTeX{} features including mathematical equations, tables, and formatting. You can cite papers using the citation search feature, and they will be automatically added to your \texttt{cite.bib} file.

\section{Mathematical Expressions}
Here are some example equations:

\subsection{Quadratic Formula}
The quadratic formula is given by:
\[x = \frac{-b \pm \sqrt{b^2 - 4ac}}{2a}\]

\subsection{Matrix Example}
A sample matrix:
\[
\begin{pmatrix}
a_{11} & a_{12} & a_{13} \\
a_{21} & a_{22} & a_{23} \\
a_{31} & a_{32} & a_{33}
\end{pmatrix}
\]

\section{Tables}
Below is a sample table showing data:

\begin{table}[h]
\centering
\begin{tabular}{@{}lrrr@{}}
\toprule
\textbf{Category} & \textbf{Value 1} & \textbf{Value 2} & \textbf{Total} \\
\midrule
Item A & 10.5 & 20.3 & 30.8 \\
Item B & 15.2 & 25.1 & 40.3 \\
Item C & 22.1 & 12.7 & 34.8 \\
\bottomrule
\end{tabular}
\caption{Sample Data Table}
\label{tab:sample}
\end{table}

\section{Lists}
\subsection{Bullet Points}
\begin{itemize}
  \item First item
  \item Second item
  \item Third item with sub-items:
    \begin{itemize}
      \item Sub-item 1
      \item Sub-item 2
    \end{itemize}
\end{itemize}

\subsection{Numbered List}
\begin{enumerate}
  \item First step
  \item Second step
  \item Third step
\end{enumerate}

\section{Citations and References}
When you use the citation search feature, references will be automatically managed. For example, you might cite a paper like this: \cite{example2024}. All citations will be automatically added to your \texttt{cite.bib} file.

\section{Conclusion}
This template demonstrates basic \LaTeX{} features including mathematical formulas, tables, and lists. The bibliography system is now set up to work seamlessly with Bibby AI's citation management. Feel free to modify it for your needs!

\bibliographystyle{plain}
\bibliography{cite}

\end{document}\documentclass{article}
\usepackage[utf8]{inputenc}
\usepackage{booktabs}
\usepackage{amsmath}
\usepackage{natbib}

\title{My Document}
\author{Your Name}
\date{\today}

\begin{document}

\maketitle

\section{Introduction}
This is a sample document demonstrating various \LaTeX{} features including mathematical equations, tables, and formatting. You can cite papers using the citation search feature, and they will be automatically added to your \texttt{cite.bib} file.

\section{Mathematical Expressions}
Here are some example equations:

\subsection{Quadratic Formula}
The quadratic formula is given by:
\[x = \frac{-b \pm \sqrt{b^2 - 4ac}}{2a}\]

\subsection{Matrix Example}
A sample matrix:
\[
\begin{pmatrix}
a_{11} & a_{12} & a_{13} \\
a_{21} & a_{22} & a_{23} \\
a_{31} & a_{32} & a_{33}
\end{pmatrix}
\]

\section{Tables}
Below is a sample table showing data:

\begin{table}[h]
\centering
\begin{tabular}{@{}lrrr@{}}
\toprule
\textbf{Category} & \textbf{Value 1} & \textbf{Value 2} & \textbf{Total} \\
\midrule
Item A & 10.5 & 20.3 & 30.8 \\
Item B & 15.2 & 25.1 & 40.3 \\
Item C & 22.1 & 12.7 & 34.8 \\
\bottomrule
\end{tabular}
\caption{Sample Data Table}
\label{tab:sample}
\end{table}

\section{Lists}
\subsection{Bullet Points}
\begin{itemize}
  \item First item
  \item Second item
  \item Third item with sub-items:
    \begin{itemize}
      \item Sub-item 1
      \item Sub-item 2
    \end{itemize}
\end{itemize}

\subsection{Numbered List}
\begin{enumerate}
  \item First step
  \item Second step
  \item Third step
\end{enumerate}

\section{Citations and References}
When you use the citation search feature, references will be automatically managed. For example, you might cite a paper like this: \cite{example2024}. All citations will be automatically added to your \texttt{cite.bib} file.

\section{Conclusion}
This template demonstrates basic \LaTeX{} features including mathematical formulas, tables, and lists. The bibliography system is now set up to work seamlessly with Bibby AI's citation management. Feel free to modify it for your needs!

\bibliographystyle{plain}
\bibliography{cite}

\end{document}\documentclass{article}
\usepackage[utf8]{inputenc}
\usepackage{booktabs}
\usepackage{amsmath}
\usepackage{natbib}

\title{My Document}
\author{Your Name}
\date{\today}

\begin{document}

\maketitle

\section{Introduction}
This is a sample document demonstrating various \LaTeX{} features including mathematical equations, tables, and formatting. You can cite papers using the citation search feature, and they will be automatically added to your \texttt{cite.bib} file.

\section{Mathematical Expressions}
Here are some example equations:

\subsection{Quadratic Formula}
The quadratic formula is given by:
\[x = \frac{-b \pm \sqrt{b^2 - 4ac}}{2a}\]

\subsection{Matrix Example}
A sample matrix:
\[
\begin{pmatrix}
a_{11} & a_{12} & a_{13} \\
a_{21} & a_{22} & a_{23} \\
a_{31} & a_{32} & a_{33}
\end{pmatrix}
\]

\section{Tables}
Below is a sample table showing data:

\begin{table}[h]
\centering
\begin{tabular}{@{}lrrr@{}}
\toprule
\textbf{Category} & \textbf{Value 1} & \textbf{Value 2} & \textbf{Total} \\
\midrule
Item A & 10.5 & 20.3 & 30.8 \\
Item B & 15.2 & 25.1 & 40.3 \\
Item C & 22.1 & 12.7 & 34.8 \\
\bottomrule
\end{tabular}
\caption{Sample Data Table}
\label{tab:sample}
\end{table}

\section{Lists}
\subsection{Bullet Points}
\begin{itemize}
  \item First item
  \item Second item
  \item Third item with sub-items:
    \begin{itemize}
      \item Sub-item 1
      \item Sub-item 2
    \end{itemize}
\end{itemize}

\subsection{Numbered List}
\begin{enumerate}
  \item First step
  \item Second step
  \item Third step
\end{enumerate}

\section{Citations and References}
When you use the citation search feature, references will be automatically managed. For example, you might cite a paper like this: \cite{example2024}. All citations will be automatically added to your \texttt{cite.bib} file.

\section{Conclusion}
This template demonstrates basic \LaTeX{} features including mathematical formulas, tables, and lists. The bibliography system is now set up to work seamlessly with Bibby AI's citation management. Feel free to modify it for your needs!

\bibliographystyle{plain}
\bibliography{cite}

\end{document}\documentclass{article}
\usepackage[utf8]{inputenc}
\usepackage{booktabs}
\usepackage{amsmath}
\usepackage{natbib}

\title{My Document}
\author{Your Name}
\date{\today}

\begin{document}

\maketitle

\section{Introduction}
This is a sample document demonstrating various \LaTeX{} features including mathematical equations, tables, and formatting. You can cite papers using the citation search feature, and they will be automatically added to your \texttt{cite.bib} file.

\section{Mathematical Expressions}
Here are some example equations:

\subsection{Quadratic Formula}
The quadratic formula is given by:
\[x = \frac{-b \pm \sqrt{b^2 - 4ac}}{2a}\]

\subsection{Matrix Example}
A sample matrix:
\[
\begin{pmatrix}
a_{11} & a_{12} & a_{13} \\
a_{21} & a_{22} & a_{23} \\
a_{31} & a_{32} & a_{33}
\end{pmatrix}
\]

\section{Tables}
Below is a sample table showing data:

\begin{table}[h]
\centering
\begin{tabular}{@{}lrrr@{}}
\toprule
\textbf{Category} & \textbf{Value 1} & \textbf{Value 2} & \textbf{Total} \\
\midrule
Item A & 10.5 & 20.3 & 30.8 \\
Item B & 15.2 & 25.1 & 40.3 \\
Item C & 22.1 & 12.7 & 34.8 \\
\bottomrule
\end{tabular}
\caption{Sample Data Table}
\label{tab:sample}
\end{table}

\section{Lists}
\subsection{Bullet Points}
\begin{itemize}
  \item First item
  \item Second item
  \item Third item with sub-items:
    \begin{itemize}
      \item Sub-item 1
      \item Sub-item 2
    \end{itemize}
\end{itemize}

\subsection{Numbered List}
\begin{enumerate}
  \item First step
  \item Second step
  \item Third step
\end{enumerate}

\section{Citations and References}
When you use the citation search feature, references will be automatically managed. For example, you might cite a paper like this: \cite{example2024}. All citations will be automatically added to your \texttt{cite.bib} file.

\section{Conclusion}
This template demonstrates basic \LaTeX{} features including mathematical formulas, tables, and lists. The bibliography system is now set up to work seamlessly with Bibby AI's citation management. Feel free to modify it for your needs!

\bibliographystyle{plain}
\bibliography{cite}

\end{document}\documentclass{article}
\usepackage[utf8]{inputenc}
\usepackage{booktabs}
\usepackage{amsmath}
\usepackage{natbib}

\title{My Document}
\author{Your Name}
\date{\today}

\begin{document}

\maketitle

\section{Introduction}
This is a sample document demonstrating various \LaTeX{} features including mathematical equations, tables, and formatting. You can cite papers using the citation search feature, and they will be automatically added to your \texttt{cite.bib} file.

\section{Mathematical Expressions}
Here are some example equations:

\subsection{Quadratic Formula}
The quadratic formula is given by:
\[x = \frac{-b \pm \sqrt{b^2 - 4ac}}{2a}\]

\subsection{Matrix Example}
A sample matrix:
\[
\begin{pmatrix}
a_{11} & a_{12} & a_{13} \\
a_{21} & a_{22} & a_{23} \\
a_{31} & a_{32} & a_{33}
\end{pmatrix}
\]

\section{Tables}
Below is a sample table showing data:

\begin{table}[h]
\centering
\begin{tabular}{@{}lrrr@{}}
\toprule
\textbf{Category} & \textbf{Value 1} & \textbf{Value 2} & \textbf{Total} \\
\midrule
Item A & 10.5 & 20.3 & 30.8 \\
Item B & 15.2 & 25.1 & 40.3 \\
Item C & 22.1 & 12.7 & 34.8 \\
\bottomrule
\end{tabular}
\caption{Sample Data Table}
\label{tab:sample}
\end{table}

\section{Lists}
\subsection{Bullet Points}
\begin{itemize}
  \item First item
  \item Second item
  \item Third item with sub-items:
    \begin{itemize}
      \item Sub-item 1
      \item Sub-item 2
    \end{itemize}
\end{itemize}

\subsection{Numbered List}
\begin{enumerate}
  \item First step
  \item Second step
  \item Third step
\end{enumerate}

\section{Citations and References}
When you use the citation search feature, references will be automatically managed. For example, you might cite a paper like this: \cite{example2024}. All citations will be automatically added to your \texttt{cite.bib} file.

\section{Conclusion}
This template demonstrates basic \LaTeX{} features including mathematical formulas, tables, and lists. The bibliography system is now set up to work seamlessly with Bibby AI's citation management. Feel free to modify it for your needs!

\bibliographystyle{plain}
\bibliography{cite}

\end{document}\documentclass{article}
\usepackage[utf8]{inputenc}
\usepackage{booktabs}
\usepackage{amsmath}
\usepackage{natbib}

\title{My Document}
\author{Your Name}
\date{\today}

\begin{document}

\maketitle

\section{Introduction}
This is a sample document demonstrating various \LaTeX{} features including mathematical equations, tables, and formatting. You can cite papers using the citation search feature, and they will be automatically added to your \texttt{cite.bib} file.

\section{Mathematical Expressions}
Here are some example equations:

\subsection{Quadratic Formula}
The quadratic formula is given by:
\[x = \frac{-b \pm \sqrt{b^2 - 4ac}}{2a}\]

\subsection{Matrix Example}
A sample matrix:
\[
\begin{pmatrix}
a_{11} & a_{12} & a_{13} \\
a_{21} & a_{22} & a_{23} \\
a_{31} & a_{32} & a_{33}
\end{pmatrix}
\]

\section{Tables}
Below is a sample table showing data:

\begin{table}[h]
\centering
\begin{tabular}{@{}lrrr@{}}
\toprule
\textbf{Category} & \textbf{Value 1} & \textbf{Value 2} & \textbf{Total} \\
\midrule
Item A & 10.5 & 20.3 & 30.8 \\
Item B & 15.2 & 25.1 & 40.3 \\
Item C & 22.1 & 12.7 & 34.8 \\
\bottomrule
\end{tabular}
\caption{Sample Data Table}
\label{tab:sample}
\end{table}

\section{Lists}
\subsection{Bullet Points}
\begin{itemize}
  \item First item
  \item Second item
  \item Third item with sub-items:
    \begin{itemize}
      \item Sub-item 1
      \item Sub-item 2
    \end{itemize}
\end{itemize}

\subsection{Numbered List}
\begin{enumerate}
  \item First step
  \item Second step
  \item Third step
\end{enumerate}

\section{Citations and References}
When you use the citation search feature, references will be automatically managed. For example, you might cite a paper like this: \cite{example2024}. All citations will be automatically added to your \texttt{cite.bib} file.

\section{Conclusion}
This template demonstrates basic \LaTeX{} features including mathematical formulas, tables, and lists. The bibliography system is now set up to work seamlessly with Bibby AI's citation management. Feel free to modify it for your needs!

\bibliographystyle{plain}
\bibliography{cite}

\end{document}\documentclass{article}
\usepackage[utf8]{inputenc}
\usepackage{booktabs}
\usepackage{amsmath}
\usepackage{natbib}

\title{My Document}
\author{Your Name}
\date{\today}

\begin{document}

\maketitle

\section{Introduction}
This is a sample document demonstrating various \LaTeX{} features including mathematical equations, tables, and formatting. You can cite papers using the citation search feature, and they will be automatically added to your \texttt{cite.bib} file.

\section{Mathematical Expressions}
Here are some example equations:

\subsection{Quadratic Formula}
The quadratic formula is given by:
\[x = \frac{-b \pm \sqrt{b^2 - 4ac}}{2a}\]

\subsection{Matrix Example}
A sample matrix:
\[
\begin{pmatrix}
a_{11} & a_{12} & a_{13} \\
a_{21} & a_{22} & a_{23} \\
a_{31} & a_{32} & a_{33}
\end{pmatrix}
\]

\section{Tables}
Below is a sample table showing data:

\begin{table}[h]
\centering
\begin{tabular}{@{}lrrr@{}}
\toprule
\textbf{Category} & \textbf{Value 1} & \textbf{Value 2} & \textbf{Total} \\
\midrule
Item A & 10.5 & 20.3 & 30.8 \\
Item B & 15.2 & 25.1 & 40.3 \\
Item C & 22.1 & 12.7 & 34.8 \\
\bottomrule
\end{tabular}
\caption{Sample Data Table}
\label{tab:sample}
\end{table}

\section{Lists}
\subsection{Bullet Points}
\begin{itemize}
  \item First item
  \item Second item
  \item Third item with sub-items:
    \begin{itemize}
      \item Sub-item 1
      \item Sub-item 2
    \end{itemize}
\end{itemize}

\subsection{Numbered List}
\begin{enumerate}
  \item First step
  \item Second step
  \item Third step
\end{enumerate}

\section{Citations and References}
When you use the citation search feature, references will be automatically managed. For example, you might cite a paper like this: \cite{example2024}. All citations will be automatically added to your \texttt{cite.bib} file.

\section{Conclusion}
This template demonstrates basic \LaTeX{} features including mathematical formulas, tables, and lists. The bibliography system is now set up to work seamlessly with Bibby AI's citation management. Feel free to modify it for your needs!

\bibliographystyle{plain}
\bibliography{cite}

\end{document}\documentclass{article}
\usepackage[utf8]{inputenc}
\usepackage{booktabs}
\usepackage{amsmath}
\usepackage{natbib}

\title{My Document}
\author{Your Name}
\date{\today}

\begin{document}

\maketitle

\section{Introduction}
This is a sample document demonstrating various \LaTeX{} features including mathematical equations, tables, and formatting. You can cite papers using the citation search feature, and they will be automatically added to your \texttt{cite.bib} file.

\section{Mathematical Expressions}
Here are some example equations:

\subsection{Quadratic Formula}
The quadratic formula is given by:
\[x = \frac{-b \pm \sqrt{b^2 - 4ac}}{2a}\]

\subsection{Matrix Example}
A sample matrix:
\[
\begin{pmatrix}
a_{11} & a_{12} & a_{13} \\
a_{21} & a_{22} & a_{23} \\
a_{31} & a_{32} & a_{33}
\end{pmatrix}
\]

\section{Tables}
Below is a sample table showing data:

\begin{table}[h]
\centering
\begin{tabular}{@{}lrrr@{}}
\toprule
\textbf{Category} & \textbf{Value 1} & \textbf{Value 2} & \textbf{Total} \\
\midrule
Item A & 10.5 & 20.3 & 30.8 \\
Item B & 15.2 & 25.1 & 40.3 \\
Item C & 22.1 & 12.7 & 34.8 \\
\bottomrule
\end{tabular}
\caption{Sample Data Table}
\label{tab:sample}
\end{table}

\section{Lists}
\subsection{Bullet Points}
\begin{itemize}
  \item First item
  \item Second item
  \item Third item with sub-items:
    \begin{itemize}
      \item Sub-item 1
      \item Sub-item 2
    \end{itemize}
\end{itemize}

\subsection{Numbered List}
\begin{enumerate}
  \item First step
  \item Second step
  \item Third step
\end{enumerate}

\section{Citations and References}
When you use the citation search feature, references will be automatically managed. For example, you might cite a paper like this: \cite{example2024}. All citations will be automatically added to your \texttt{cite.bib} file.

\section{Conclusion}
This template demonstrates basic \LaTeX{} features including mathematical formulas, tables, and lists. The bibliography system is now set up to work seamlessly with Bibby AI's citation management. Feel free to modify it for your needs!

\bibliographystyle{plain}
\bibliography{cite}

\end{document}\documentclass{article}
\usepackage[utf8]{inputenc}
\usepackage{booktabs}
\usepackage{amsmath}
\usepackage{natbib}

\title{My Document}
\author{Your Name}
\date{\today}

\begin{document}

\maketitle

\section{Introduction}
This is a sample document demonstrating various \LaTeX{} features including mathematical equations, tables, and formatting. You can cite papers using the citation search feature, and they will be automatically added to your \texttt{cite.bib} file.

\section{Mathematical Expressions}
Here are some example equations:

\subsection{Quadratic Formula}
The quadratic formula is given by:
\[x = \frac{-b \pm \sqrt{b^2 - 4ac}}{2a}\]

\subsection{Matrix Example}
A sample matrix:
\[
\begin{pmatrix}
a_{11} & a_{12} & a_{13} \\
a_{21} & a_{22} & a_{23} \\
a_{31} & a_{32} & a_{33}
\end{pmatrix}
\]

\section{Tables}
Below is a sample table showing data:

\begin{table}[h]
\centering
\begin{tabular}{@{}lrrr@{}}
\toprule
\textbf{Category} & \textbf{Value 1} & \textbf{Value 2} & \textbf{Total} \\
\midrule
Item A & 10.5 & 20.3 & 30.8 \\
Item B & 15.2 & 25.1 & 40.3 \\
Item C & 22.1 & 12.7 & 34.8 \\
\bottomrule
\end{tabular}
\caption{Sample Data Table}
\label{tab:sample}
\end{table}

\section{Lists}
\subsection{Bullet Points}
\begin{itemize}
  \item First item
  \item Second item
  \item Third item with sub-items:
    \begin{itemize}
      \item Sub-item 1
      \item Sub-item 2
    \end{itemize}
\end{itemize}

\subsection{Numbered List}
\begin{enumerate}
  \item First step
  \item Second step
  \item Third step
\end{enumerate}

\section{Citations and References}
When you use the citation search feature, references will be automatically managed. For example, you might cite a paper like this: \cite{example2024}. All citations will be automatically added to your \texttt{cite.bib} file.

\section{Conclusion}
This template demonstrates basic \LaTeX{} features including mathematical formulas, tables, and lists. The bibliography system is now set up to work seamlessly with Bibby AI's citation management. Feel free to modify it for your needs!

\bibliographystyle{plain}
\bibliography{cite}

\end{document}\documentclass{article}
\usepackage[utf8]{inputenc}
\usepackage{booktabs}
\usepackage{amsmath}
\usepackage{natbib}

\title{My Document}
\author{Your Name}
\date{\today}

\begin{document}

\maketitle

\section{Introduction}
This is a sample document demonstrating various \LaTeX{} features including mathematical equations, tables, and formatting. You can cite papers using the citation search feature, and they will be automatically added to your \texttt{cite.bib} file.

\section{Mathematical Expressions}
Here are some example equations:

\subsection{Quadratic Formula}
The quadratic formula is given by:
\[x = \frac{-b \pm \sqrt{b^2 - 4ac}}{2a}\]

\subsection{Matrix Example}
A sample matrix:
\[
\begin{pmatrix}
a_{11} & a_{12} & a_{13} \\
a_{21} & a_{22} & a_{23} \\
a_{31} & a_{32} & a_{33}
\end{pmatrix}
\]

\section{Tables}
Below is a sample table showing data:

\begin{table}[h]
\centering
\begin{tabular}{@{}lrrr@{}}
\toprule
\textbf{Category} & \textbf{Value 1} & \textbf{Value 2} & \textbf{Total} \\
\midrule
Item A & 10.5 & 20.3 & 30.8 \\
Item B & 15.2 & 25.1 & 40.3 \\
Item C & 22.1 & 12.7 & 34.8 \\
\bottomrule
\end{tabular}
\caption{Sample Data Table}
\label{tab:sample}
\end{table}

\section{Lists}
\subsection{Bullet Points}
\begin{itemize}
  \item First item
  \item Second item
  \item Third item with sub-items:
    \begin{itemize}
      \item Sub-item 1
      \item Sub-item 2
    \end{itemize}
\end{itemize}

\subsection{Numbered List}
\begin{enumerate}
  \item First step
  \item Second step
  \item Third step
\end{enumerate}

\section{Citations and References}
When you use the citation search feature, references will be automatically managed. For example, you might cite a paper like this: \cite{example2024}. All citations will be automatically added to your \texttt{cite.bib} file.

\section{Conclusion}
This template demonstrates basic \LaTeX{} features including mathematical formulas, tables, and lists. The bibliography system is now set up to work seamlessly with Bibby AI's citation management. Feel free to modify it for your needs!

\bibliographystyle{plain}
\bibliography{cite}

\end{document}\documentclass{article}
\usepackage[utf8]{inputenc}
\usepackage{booktabs}
\usepackage{amsmath}
\usepackage{natbib}

\title{My Document}
\author{Your Name}
\date{\today}

\begin{document}

\maketitle

\section{Introduction}
This is a sample document demonstrating various \LaTeX{} features including mathematical equations, tables, and formatting. You can cite papers using the citation search feature, and they will be automatically added to your \texttt{cite.bib} file.

\section{Mathematical Expressions}
Here are some example equations:

\subsection{Quadratic Formula}
The quadratic formula is given by:
\[x = \frac{-b \pm \sqrt{b^2 - 4ac}}{2a}\]

\subsection{Matrix Example}
A sample matrix:
\[
\begin{pmatrix}
a_{11} & a_{12} & a_{13} \\
a_{21} & a_{22} & a_{23} \\
a_{31} & a_{32} & a_{33}
\end{pmatrix}
\]

\section{Tables}
Below is a sample table showing data:

\begin{table}[h]
\centering
\begin{tabular}{@{}lrrr@{}}
\toprule
\textbf{Category} & \textbf{Value 1} & \textbf{Value 2} & \textbf{Total} \\
\midrule
Item A & 10.5 & 20.3 & 30.8 \\
Item B & 15.2 & 25.1 & 40.3 \\
Item C & 22.1 & 12.7 & 34.8 \\
\bottomrule
\end{tabular}
\caption{Sample Data Table}
\label{tab:sample}
\end{table}

\section{Lists}
\subsection{Bullet Points}
\begin{itemize}
  \item First item
  \item Second item
  \item Third item with sub-items:
    \begin{itemize}
      \item Sub-item 1
      \item Sub-item 2
    \end{itemize}
\end{itemize}

\subsection{Numbered List}
\begin{enumerate}
  \item First step
  \item Second step
  \item Third step
\end{enumerate}

\section{Citations and References}
When you use the citation search feature, references will be automatically managed. For example, you might cite a paper like this: \cite{example2024}. All citations will be automatically added to your \texttt{cite.bib} file.

\section{Conclusion}
This template demonstrates basic \LaTeX{} features including mathematical formulas, tables, and lists. The bibliography system is now set up to work seamlessly with Bibby AI's citation management. Feel free to modify it for your needs!

\bibliographystyle{plain}
\bibliography{cite}

\end{document}\documentclass{article}
\usepackage[utf8]{inputenc}
\usepackage{booktabs}
\usepackage{amsmath}
\usepackage{natbib}

\title{My Document}
\author{Your Name}
\date{\today}

\begin{document}

\maketitle

\section{Introduction}
This is a sample document demonstrating various \LaTeX{} features including mathematical equations, tables, and formatting. You can cite papers using the citation search feature, and they will be automatically added to your \texttt{cite.bib} file.

\section{Mathematical Expressions}
Here are some example equations:

\subsection{Quadratic Formula}
The quadratic formula is given by:
\[x = \frac{-b \pm \sqrt{b^2 - 4ac}}{2a}\]

\subsection{Matrix Example}
A sample matrix:
\[
\begin{pmatrix}
a_{11} & a_{12} & a_{13} \\
a_{21} & a_{22} & a_{23} \\
a_{31} & a_{32} & a_{33}
\end{pmatrix}
\]

\section{Tables}
Below is a sample table showing data:

\begin{table}[h]
\centering
\begin{tabular}{@{}lrrr@{}}
\toprule
\textbf{Category} & \textbf{Value 1} & \textbf{Value 2} & \textbf{Total} \\
\midrule
Item A & 10.5 & 20.3 & 30.8 \\
Item B & 15.2 & 25.1 & 40.3 \\
Item C & 22.1 & 12.7 & 34.8 \\
\bottomrule
\end{tabular}
\caption{Sample Data Table}
\label{tab:sample}
\end{table}

\section{Lists}
\subsection{Bullet Points}
\begin{itemize}
  \item First item
  \item Second item
  \item Third item with sub-items:
    \begin{itemize}
      \item Sub-item 1
      \item Sub-item 2
    \end{itemize}
\end{itemize}

\subsection{Numbered List}
\begin{enumerate}
  \item First step
  \item Second step
  \item Third step
\end{enumerate}

\section{Citations and References}
When you use the citation search feature, references will be automatically managed. For example, you might cite a paper like this: \cite{example2024}. All citations will be automatically added to your \texttt{cite.bib} file.

\section{Conclusion}
This template demonstrates basic \LaTeX{} features including mathematical formulas, tables, and lists. The bibliography system is now set up to work seamlessly with Bibby AI's citation management. Feel free to modify it for your needs!

\bibliographystyle{plain}
\bibliography{cite}

\end{document}\documentclass{article}
\usepackage[utf8]{inputenc}
\usepackage{booktabs}
\usepackage{amsmath}
\usepackage{natbib}

\title{My Document}
\author{Your Name}
\date{\today}

\begin{document}

\maketitle

\section{Introduction}
This is a sample document demonstrating various \LaTeX{} features including mathematical equations, tables, and formatting. You can cite papers using the citation search feature, and they will be automatically added to your \texttt{cite.bib} file.

\section{Mathematical Expressions}
Here are some example equations:

\subsection{Quadratic Formula}
The quadratic formula is given by:
\[x = \frac{-b \pm \sqrt{b^2 - 4ac}}{2a}\]

\subsection{Matrix Example}
A sample matrix:
\[
\begin{pmatrix}
a_{11} & a_{12} & a_{13} \\
a_{21} & a_{22} & a_{23} \\
a_{31} & a_{32} & a_{33}
\end{pmatrix}
\]

\section{Tables}
Below is a sample table showing data:

\begin{table}[h]
\centering
\begin{tabular}{@{}lrrr@{}}
\toprule
\textbf{Category} & \textbf{Value 1} & \textbf{Value 2} & \textbf{Total} \\
\midrule
Item A & 10.5 & 20.3 & 30.8 \\
Item B & 15.2 & 25.1 & 40.3 \\
Item C & 22.1 & 12.7 & 34.8 \\
\bottomrule
\end{tabular}
\caption{Sample Data Table}
\label{tab:sample}
\end{table}

\section{Lists}
\subsection{Bullet Points}
\begin{itemize}
  \item First item
  \item Second item
  \item Third item with sub-items:
    \begin{itemize}
      \item Sub-item 1
      \item Sub-item 2
    \end{itemize}
\end{itemize}

\subsection{Numbered List}
\begin{enumerate}
  \item First step
  \item Second step
  \item Third step
\end{enumerate}

\section{Citations and References}
When you use the citation search feature, references will be automatically managed. For example, you might cite a paper like this: \cite{example2024}. All citations will be automatically added to your \texttt{cite.bib} file.

\section{Conclusion}
This template demonstrates basic \LaTeX{} features including mathematical formulas, tables, and lists. The bibliography system is now set up to work seamlessly with Bibby AI's citation management. Feel free to modify it for your needs!

\bibliographystyle{plain}
\bibliography{cite}

\end{document}\documentclass{article}
\usepackage[utf8]{inputenc}
\usepackage{booktabs}
\usepackage{amsmath}
\usepackage{natbib}

\title{My Document}
\author{Your Name}
\date{\today}

\begin{document}

\maketitle

\section{Introduction}
This is a sample document demonstrating various \LaTeX{} features including mathematical equations, tables, and formatting. You can cite papers using the citation search feature, and they will be automatically added to your \texttt{cite.bib} file.

\section{Mathematical Expressions}
Here are some example equations:

\subsection{Quadratic Formula}
The quadratic formula is given by:
\[x = \frac{-b \pm \sqrt{b^2 - 4ac}}{2a}\]

\subsection{Matrix Example}
A sample matrix:
\[
\begin{pmatrix}
a_{11} & a_{12} & a_{13} \\
a_{21} & a_{22} & a_{23} \\
a_{31} & a_{32} & a_{33}
\end{pmatrix}
\]

\section{Tables}
Below is a sample table showing data:

\begin{table}[h]
\centering
\begin{tabular}{@{}lrrr@{}}
\toprule
\textbf{Category} & \textbf{Value 1} & \textbf{Value 2} & \textbf{Total} \\
\midrule
Item A & 10.5 & 20.3 & 30.8 \\
Item B & 15.2 & 25.1 & 40.3 \\
Item C & 22.1 & 12.7 & 34.8 \\
\bottomrule
\end{tabular}
\caption{Sample Data Table}
\label{tab:sample}
\end{table}

\section{Lists}
\subsection{Bullet Points}
\begin{itemize}
  \item First item
  \item Second item
  \item Third item with sub-items:
    \begin{itemize}
      \item Sub-item 1
      \item Sub-item 2
    \end{itemize}
\end{itemize}

\subsection{Numbered List}
\begin{enumerate}
  \item First step
  \item Second step
  \item Third step
\end{enumerate}

\section{Citations and References}
When you use the citation search feature, references will be automatically managed. For example, you might cite a paper like this: \cite{example2024}. All citations will be automatically added to your \texttt{cite.bib} file.

\section{Conclusion}
This template demonstrates basic \LaTeX{} features including mathematical formulas, tables, and lists. The bibliography system is now set up to work seamlessly with Bibby AI's citation management. Feel free to modify it for your needs!

\bibliographystyle{plain}
\bibliography{cite}

\end{document}\documentclass{article}
\usepackage[utf8]{inputenc}
\usepackage{booktabs}
\usepackage{amsmath}
\usepackage{natbib}

\title{My Document}
\author{Your Name}
\date{\today}

\begin{document}

\maketitle

\section{Introduction}
This is a sample document demonstrating various \LaTeX{} features including mathematical equations, tables, and formatting. You can cite papers using the citation search feature, and they will be automatically added to your \texttt{cite.bib} file.

\section{Mathematical Expressions}
Here are some example equations:

\subsection{Quadratic Formula}
The quadratic formula is given by:
\[x = \frac{-b \pm \sqrt{b^2 - 4ac}}{2a}\]

\subsection{Matrix Example}
A sample matrix:
\[
\begin{pmatrix}
a_{11} & a_{12} & a_{13} \\
a_{21} & a_{22} & a_{23} \\
a_{31} & a_{32} & a_{33}
\end{pmatrix}
\]

\section{Tables}
Below is a sample table showing data:

\begin{table}[h]
\centering
\begin{tabular}{@{}lrrr@{}}
\toprule
\textbf{Category} & \textbf{Value 1} & \textbf{Value 2} & \textbf{Total} \\
\midrule
Item A & 10.5 & 20.3 & 30.8 \\
Item B & 15.2 & 25.1 & 40.3 \\
Item C & 22.1 & 12.7 & 34.8 \\
\bottomrule
\end{tabular}
\caption{Sample Data Table}
\label{tab:sample}
\end{table}

\section{Lists}
\subsection{Bullet Points}
\begin{itemize}
  \item First item
  \item Second item
  \item Third item with sub-items:
    \begin{itemize}
      \item Sub-item 1
      \item Sub-item 2
    \end{itemize}
\end{itemize}

\subsection{Numbered List}
\begin{enumerate}
  \item First step
  \item Second step
  \item Third step
\end{enumerate}

\section{Citations and References}
When you use the citation search feature, references will be automatically managed. For example, you might cite a paper like this: \cite{example2024}. All citations will be automatically added to your \texttt{cite.bib} file.

\section{Conclusion}
This template demonstrates basic \LaTeX{} features including mathematical formulas, tables, and lists. The bibliography system is now set up to work seamlessly with Bibby AI's citation management. Feel free to modify it for your needs!

\bibliographystyle{plain}
\bibliography{cite}

\end{document}\documentclass{article}
\usepackage[utf8]{inputenc}
\usepackage{booktabs}
\usepackage{amsmath}
\usepackage{natbib}

\title{My Document}
\author{Your Name}
\date{\today}

\begin{document}

\maketitle

\section{Introduction}
This is a sample document demonstrating various \LaTeX{} features including mathematical equations, tables, and formatting. You can cite papers using the citation search feature, and they will be automatically added to your \texttt{cite.bib} file.

\section{Mathematical Expressions}
Here are some example equations:

\subsection{Quadratic Formula}
The quadratic formula is given by:
\[x = \frac{-b \pm \sqrt{b^2 - 4ac}}{2a}\]

\subsection{Matrix Example}
A sample matrix:
\[
\begin{pmatrix}
a_{11} & a_{12} & a_{13} \\
a_{21} & a_{22} & a_{23} \\
a_{31} & a_{32} & a_{33}
\end{pmatrix}
\]

\section{Tables}
Below is a sample table showing data:

\begin{table}[h]
\centering
\begin{tabular}{@{}lrrr@{}}
\toprule
\textbf{Category} & \textbf{Value 1} & \textbf{Value 2} & \textbf{Total} \\
\midrule
Item A & 10.5 & 20.3 & 30.8 \\
Item B & 15.2 & 25.1 & 40.3 \\
Item C & 22.1 & 12.7 & 34.8 \\
\bottomrule
\end{tabular}
\caption{Sample Data Table}
\label{tab:sample}
\end{table}

\section{Lists}
\subsection{Bullet Points}
\begin{itemize}
  \item First item
  \item Second item
  \item Third item with sub-items:
    \begin{itemize}
      \item Sub-item 1
      \item Sub-item 2
    \end{itemize}
\end{itemize}

\subsection{Numbered List}
\begin{enumerate}
  \item First step
  \item Second step
  \item Third step
\end{enumerate}

\section{Citations and References}
When you use the citation search feature, references will be automatically managed. For example, you might cite a paper like this: \cite{example2024}. All citations will be automatically added to your \texttt{cite.bib} file.

\section{Conclusion}
This template demonstrates basic \LaTeX{} features including mathematical formulas, tables, and lists. The bibliography system is now set up to work seamlessly with Bibby AI's citation management. Feel free to modify it for your needs!

\bibliographystyle{plain}
\bibliography{cite}

\end{document}\documentclass{article}
\usepackage[utf8]{inputenc}
\usepackage{booktabs}
\usepackage{amsmath}
\usepackage{natbib}

\title{My Document}
\author{Your Name}
\date{\today}

\begin{document}

\maketitle

\section{Introduction}
This is a sample document demonstrating various \LaTeX{} features including mathematical equations, tables, and formatting. You can cite papers using the citation search feature, and they will be automatically added to your \texttt{cite.bib} file.

\section{Mathematical Expressions}
Here are some example equations:

\subsection{Quadratic Formula}
The quadratic formula is given by:
\[x = \frac{-b \pm \sqrt{b^2 - 4ac}}{2a}\]

\subsection{Matrix Example}
A sample matrix:
\[
\begin{pmatrix}
a_{11} & a_{12} & a_{13} \\
a_{21} & a_{22} & a_{23} \\
a_{31} & a_{32} & a_{33}
\end{pmatrix}
\]

\section{Tables}
Below is a sample table showing data:

\begin{table}[h]
\centering
\begin{tabular}{@{}lrrr@{}}
\toprule
\textbf{Category} & \textbf{Value 1} & \textbf{Value 2} & \textbf{Total} \\
\midrule
Item A & 10.5 & 20.3 & 30.8 \\
Item B & 15.2 & 25.1 & 40.3 \\
Item C & 22.1 & 12.7 & 34.8 \\
\bottomrule
\end{tabular}
\caption{Sample Data Table}
\label{tab:sample}
\end{table}

\section{Lists}
\subsection{Bullet Points}
\begin{itemize}
  \item First item
  \item Second item
  \item Third item with sub-items:
    \begin{itemize}
      \item Sub-item 1
      \item Sub-item 2
    \end{itemize}
\end{itemize}

\subsection{Numbered List}
\begin{enumerate}
  \item First step
  \item Second step
  \item Third step
\end{enumerate}

\section{Citations and References}
When you use the citation search feature, references will be automatically managed. For example, you might cite a paper like this: \cite{example2024}. All citations will be automatically added to your \texttt{cite.bib} file.

\section{Conclusion}
This template demonstrates basic \LaTeX{} features including mathematical formulas, tables, and lists. The bibliography system is now set up to work seamlessly with Bibby AI's citation management. Feel free to modify it for your needs!

\bibliographystyle{plain}
\bibliography{cite}

\end{document}\documentclass{article}
\usepackage[utf8]{inputenc}
\usepackage{booktabs}
\usepackage{amsmath}
\usepackage{natbib}

\title{My Document}
\author{Your Name}
\date{\today}

\begin{document}

\maketitle

\section{Introduction}
This is a sample document demonstrating various \LaTeX{} features including mathematical equations, tables, and formatting. You can cite papers using the citation search feature, and they will be automatically added to your \texttt{cite.bib} file.

\section{Mathematical Expressions}
Here are some example equations:

\subsection{Quadratic Formula}
The quadratic formula is given by:
\[x = \frac{-b \pm \sqrt{b^2 - 4ac}}{2a}\]

\subsection{Matrix Example}
A sample matrix:
\[
\begin{pmatrix}
a_{11} & a_{12} & a_{13} \\
a_{21} & a_{22} & a_{23} \\
a_{31} & a_{32} & a_{33}
\end{pmatrix}
\]

\section{Tables}
Below is a sample table showing data:

\begin{table}[h]
\centering
\begin{tabular}{@{}lrrr@{}}
\toprule
\textbf{Category} & \textbf{Value 1} & \textbf{Value 2} & \textbf{Total} \\
\midrule
Item A & 10.5 & 20.3 & 30.8 \\
Item B & 15.2 & 25.1 & 40.3 \\
Item C & 22.1 & 12.7 & 34.8 \\
\bottomrule
\end{tabular}
\caption{Sample Data Table}
\label{tab:sample}
\end{table}

\section{Lists}
\subsection{Bullet Points}
\begin{itemize}
  \item First item
  \item Second item
  \item Third item with sub-items:
    \begin{itemize}
      \item Sub-item 1
      \item Sub-item 2
    \end{itemize}
\end{itemize}

\subsection{Numbered List}
\begin{enumerate}
  \item First step
  \item Second step
  \item Third step
\end{enumerate}

\section{Citations and References}
When you use the citation search feature, references will be automatically managed. For example, you might cite a paper like this: \cite{example2024}. All citations will be automatically added to your \texttt{cite.bib} file.

\section{Conclusion}
This template demonstrates basic \LaTeX{} features including mathematical formulas, tables, and lists. The bibliography system is now set up to work seamlessly with Bibby AI's citation management. Feel free to modify it for your needs!

\bibliographystyle{plain}
\bibliography{cite}

\end{document}\documentclass{article}
\usepackage[utf8]{inputenc}
\usepackage{booktabs}
\usepackage{amsmath}
\usepackage{natbib}

\title{My Document}
\author{Your Name}
\date{\today}

\begin{document}

\maketitle

\section{Introduction}
This is a sample document demonstrating various \LaTeX{} features including mathematical equations, tables, and formatting. You can cite papers using the citation search feature, and they will be automatically added to your \texttt{cite.bib} file.

\section{Mathematical Expressions}
Here are some example equations:

\subsection{Quadratic Formula}
The quadratic formula is given by:
\[x = \frac{-b \pm \sqrt{b^2 - 4ac}}{2a}\]

\subsection{Matrix Example}
A sample matrix:
\[
\begin{pmatrix}
a_{11} & a_{12} & a_{13} \\
a_{21} & a_{22} & a_{23} \\
a_{31} & a_{32} & a_{33}
\end{pmatrix}
\]

\section{Tables}
Below is a sample table showing data:

\begin{table}[h]
\centering
\begin{tabular}{@{}lrrr@{}}
\toprule
\textbf{Category} & \textbf{Value 1} & \textbf{Value 2} & \textbf{Total} \\
\midrule
Item A & 10.5 & 20.3 & 30.8 \\
Item B & 15.2 & 25.1 & 40.3 \\
Item C & 22.1 & 12.7 & 34.8 \\
\bottomrule
\end{tabular}
\caption{Sample Data Table}
\label{tab:sample}
\end{table}

\section{Lists}
\subsection{Bullet Points}
\begin{itemize}
  \item First item
  \item Second item
  \item Third item with sub-items:
    \begin{itemize}
      \item Sub-item 1
      \item Sub-item 2
    \end{itemize}
\end{itemize}

\subsection{Numbered List}
\begin{enumerate}
  \item First step
  \item Second step
  \item Third step
\end{enumerate}

\section{Citations and References}
When you use the citation search feature, references will be automatically managed. For example, you might cite a paper like this: \cite{example2024}. All citations will be automatically added to your \texttt{cite.bib} file.

\section{Conclusion}
This template demonstrates basic \LaTeX{} features including mathematical formulas, tables, and lists. The bibliography system is now set up to work seamlessly with Bibby AI's citation management. Feel free to modify it for your needs!

\bibliographystyle{plain}
\bibliography{cite}

\end{document}\documentclass{article}
\usepackage[utf8]{inputenc}
\usepackage{booktabs}
\usepackage{amsmath}
\usepackage{natbib}

\title{My Document}
\author{Your Name}
\date{\today}

\begin{document}

\maketitle

\section{Introduction}
This is a sample document demonstrating various \LaTeX{} features including mathematical equations, tables, and formatting. You can cite papers using the citation search feature, and they will be automatically added to your \texttt{cite.bib} file.

\section{Mathematical Expressions}
Here are some example equations:

\subsection{Quadratic Formula}
The quadratic formula is given by:
\[x = \frac{-b \pm \sqrt{b^2 - 4ac}}{2a}\]

\subsection{Matrix Example}
A sample matrix:
\[
\begin{pmatrix}
a_{11} & a_{12} & a_{13} \\
a_{21} & a_{22} & a_{23} \\
a_{31} & a_{32} & a_{33}
\end{pmatrix}
\]

\section{Tables}
Below is a sample table showing data:

\begin{table}[h]
\centering
\begin{tabular}{@{}lrrr@{}}
\toprule
\textbf{Category} & \textbf{Value 1} & \textbf{Value 2} & \textbf{Total} \\
\midrule
Item A & 10.5 & 20.3 & 30.8 \\
Item B & 15.2 & 25.1 & 40.3 \\
Item C & 22.1 & 12.7 & 34.8 \\
\bottomrule
\end{tabular}
\caption{Sample Data Table}
\label{tab:sample}
\end{table}

\section{Lists}
\subsection{Bullet Points}
\begin{itemize}
  \item First item
  \item Second item
  \item Third item with sub-items:
    \begin{itemize}
      \item Sub-item 1
      \item Sub-item 2
    \end{itemize}
\end{itemize}

\subsection{Numbered List}
\begin{enumerate}
  \item First step
  \item Second step
  \item Third step
\end{enumerate}

\section{Citations and References}
When you use the citation search feature, references will be automatically managed. For example, you might cite a paper like this: \cite{example2024}. All citations will be automatically added to your \texttt{cite.bib} file.

\section{Conclusion}
This template demonstrates basic \LaTeX{} features including mathematical formulas, tables, and lists. The bibliography system is now set up to work seamlessly with Bibby AI's citation management. Feel free to modify it for your needs!

\bibliographystyle{plain}
\bibliography{cite}

\end{document}\documentclass{article}
\usepackage[utf8]{inputenc}
\usepackage{booktabs}
\usepackage{amsmath}
\usepackage{natbib}

\title{My Document}
\author{Your Name}
\date{\today}

\begin{document}

\maketitle

\section{Introduction}
This is a sample document demonstrating various \LaTeX{} features including mathematical equations, tables, and formatting. You can cite papers using the citation search feature, and they will be automatically added to your \texttt{cite.bib} file.

\section{Mathematical Expressions}
Here are some example equations:

\subsection{Quadratic Formula}
The quadratic formula is given by:
\[x = \frac{-b \pm \sqrt{b^2 - 4ac}}{2a}\]

\subsection{Matrix Example}
A sample matrix:
\[
\begin{pmatrix}
a_{11} & a_{12} & a_{13} \\
a_{21} & a_{22} & a_{23} \\
a_{31} & a_{32} & a_{33}
\end{pmatrix}
\]

\section{Tables}
Below is a sample table showing data:

\begin{table}[h]
\centering
\begin{tabular}{@{}lrrr@{}}
\toprule
\textbf{Category} & \textbf{Value 1} & \textbf{Value 2} & \textbf{Total} \\
\midrule
Item A & 10.5 & 20.3 & 30.8 \\
Item B & 15.2 & 25.1 & 40.3 \\
Item C & 22.1 & 12.7 & 34.8 \\
\bottomrule
\end{tabular}
\caption{Sample Data Table}
\label{tab:sample}
\end{table}

\section{Lists}
\subsection{Bullet Points}
\begin{itemize}
  \item First item
  \item Second item
  \item Third item with sub-items:
    \begin{itemize}
      \item Sub-item 1
      \item Sub-item 2
    \end{itemize}
\end{itemize}

\subsection{Numbered List}
\begin{enumerate}
  \item First step
  \item Second step
  \item Third step
\end{enumerate}

\section{Citations and References}
When you use the citation search feature, references will be automatically managed. For example, you might cite a paper like this: \cite{example2024}. All citations will be automatically added to your \texttt{cite.bib} file.

\section{Conclusion}
This template demonstrates basic \LaTeX{} features including mathematical formulas, tables, and lists. The bibliography system is now set up to work seamlessly with Bibby AI's citation management. Feel free to modify it for your needs!

\bibliographystyle{plain}
\bibliography{cite}

\end{document}\documentclass{article}
\usepackage[utf8]{inputenc}
\usepackage{booktabs}
\usepackage{amsmath}
\usepackage{natbib}

\title{My Document}
\author{Your Name}
\date{\today}

\begin{document}

\maketitle

\section{Introduction}
This is a sample document demonstrating various \LaTeX{} features including mathematical equations, tables, and formatting. You can cite papers using the citation search feature, and they will be automatically added to your \texttt{cite.bib} file.

\section{Mathematical Expressions}
Here are some example equations:

\subsection{Quadratic Formula}
The quadratic formula is given by:
\[x = \frac{-b \pm \sqrt{b^2 - 4ac}}{2a}\]

\subsection{Matrix Example}
A sample matrix:
\[
\begin{pmatrix}
a_{11} & a_{12} & a_{13} \\
a_{21} & a_{22} & a_{23} \\
a_{31} & a_{32} & a_{33}
\end{pmatrix}
\]

\section{Tables}
Below is a sample table showing data:

\begin{table}[h]
\centering
\begin{tabular}{@{}lrrr@{}}
\toprule
\textbf{Category} & \textbf{Value 1} & \textbf{Value 2} & \textbf{Total} \\
\midrule
Item A & 10.5 & 20.3 & 30.8 \\
Item B & 15.2 & 25.1 & 40.3 \\
Item C & 22.1 & 12.7 & 34.8 \\
\bottomrule
\end{tabular}
\caption{Sample Data Table}
\label{tab:sample}
\end{table}

\section{Lists}
\subsection{Bullet Points}
\begin{itemize}
  \item First item
  \item Second item
  \item Third item with sub-items:
    \begin{itemize}
      \item Sub-item 1
      \item Sub-item 2
    \end{itemize}
\end{itemize}

\subsection{Numbered List}
\begin{enumerate}
  \item First step
  \item Second step
  \item Third step
\end{enumerate}

\section{Citations and References}
When you use the citation search feature, references will be automatically managed. For example, you might cite a paper like this: \cite{example2024}. All citations will be automatically added to your \texttt{cite.bib} file.

\section{Conclusion}
This template demonstrates basic \LaTeX{} features including mathematical formulas, tables, and lists. The bibliography system is now set up to work seamlessly with Bibby AI's citation management. Feel free to modify it for your needs!

\bibliographystyle{plain}
\bibliography{cite}

\end{document}\documentclass{article}
\usepackage[utf8]{inputenc}
\usepackage{booktabs}
\usepackage{amsmath}
\usepackage{natbib}

\title{My Document}
\author{Your Name}
\date{\today}

\begin{document}

\maketitle

\section{Introduction}
This is a sample document demonstrating various \LaTeX{} features including mathematical equations, tables, and formatting. You can cite papers using the citation search feature, and they will be automatically added to your \texttt{cite.bib} file.

\section{Mathematical Expressions}
Here are some example equations:

\subsection{Quadratic Formula}
The quadratic formula is given by:
\[x = \frac{-b \pm \sqrt{b^2 - 4ac}}{2a}\]

\subsection{Matrix Example}
A sample matrix:
\[
\begin{pmatrix}
a_{11} & a_{12} & a_{13} \\
a_{21} & a_{22} & a_{23} \\
a_{31} & a_{32} & a_{33}
\end{pmatrix}
\]

\section{Tables}
Below is a sample table showing data:

\begin{table}[h]
\centering
\begin{tabular}{@{}lrrr@{}}
\toprule
\textbf{Category} & \textbf{Value 1} & \textbf{Value 2} & \textbf{Total} \\
\midrule
Item A & 10.5 & 20.3 & 30.8 \\
Item B & 15.2 & 25.1 & 40.3 \\
Item C & 22.1 & 12.7 & 34.8 \\
\bottomrule
\end{tabular}
\caption{Sample Data Table}
\label{tab:sample}
\end{table}

\section{Lists}
\subsection{Bullet Points}
\begin{itemize}
  \item First item
  \item Second item
  \item Third item with sub-items:
    \begin{itemize}
      \item Sub-item 1
      \item Sub-item 2
    \end{itemize}
\end{itemize}

\subsection{Numbered List}
\begin{enumerate}
  \item First step
  \item Second step
  \item Third step
\end{enumerate}

\section{Citations and References}
When you use the citation search feature, references will be automatically managed. For example, you might cite a paper like this: \cite{example2024}. All citations will be automatically added to your \texttt{cite.bib} file.

\section{Conclusion}
This template demonstrates basic \LaTeX{} features including mathematical formulas, tables, and lists. The bibliography system is now set up to work seamlessly with Bibby AI's citation management. Feel free to modify it for your needs!

\bibliographystyle{plain}
\bibliography{cite}

\end{document}\documentclass{article}
\usepackage[utf8]{inputenc}
\usepackage{booktabs}
\usepackage{amsmath}
\usepackage{natbib}

\title{My Document}
\author{Your Name}
\date{\today}

\begin{document}

\maketitle

\section{Introduction}
This is a sample document demonstrating various \LaTeX{} features including mathematical equations, tables, and formatting. You can cite papers using the citation search feature, and they will be automatically added to your \texttt{cite.bib} file.

\section{Mathematical Expressions}
Here are some example equations:

\subsection{Quadratic Formula}
The quadratic formula is given by:
\[x = \frac{-b \pm \sqrt{b^2 - 4ac}}{2a}\]

\subsection{Matrix Example}
A sample matrix:
\[
\begin{pmatrix}
a_{11} & a_{12} & a_{13} \\
a_{21} & a_{22} & a_{23} \\
a_{31} & a_{32} & a_{33}
\end{pmatrix}
\]

\section{Tables}
Below is a sample table showing data:

\begin{table}[h]
\centering
\begin{tabular}{@{}lrrr@{}}
\toprule
\textbf{Category} & \textbf{Value 1} & \textbf{Value 2} & \textbf{Total} \\
\midrule
Item A & 10.5 & 20.3 & 30.8 \\
Item B & 15.2 & 25.1 & 40.3 \\
Item C & 22.1 & 12.7 & 34.8 \\
\bottomrule
\end{tabular}
\caption{Sample Data Table}
\label{tab:sample}
\end{table}

\section{Lists}
\subsection{Bullet Points}
\begin{itemize}
  \item First item
  \item Second item
  \item Third item with sub-items:
    \begin{itemize}
      \item Sub-item 1
      \item Sub-item 2
    \end{itemize}
\end{itemize}

\subsection{Numbered List}
\begin{enumerate}
  \item First step
  \item Second step
  \item Third step
\end{enumerate}

\section{Citations and References}
When you use the citation search feature, references will be automatically managed. For example, you might cite a paper like this: \cite{example2024}. All citations will be automatically added to your \texttt{cite.bib} file.

\section{Conclusion}
This template demonstrates basic \LaTeX{} features including mathematical formulas, tables, and lists. The bibliography system is now set up to work seamlessly with Bibby AI's citation management. Feel free to modify it for your needs!

\bibliographystyle{plain}
\bibliography{cite}

\end{document}\documentclass{article}
\usepackage[utf8]{inputenc}
\usepackage{booktabs}
\usepackage{amsmath}
\usepackage{natbib}

\title{My Document}
\author{Your Name}
\date{\today}

\begin{document}

\maketitle

\section{Introduction}
This is a sample document demonstrating various \LaTeX{} features including mathematical equations, tables, and formatting. You can cite papers using the citation search feature, and they will be automatically added to your \texttt{cite.bib} file.

\section{Mathematical Expressions}
Here are some example equations:

\subsection{Quadratic Formula}
The quadratic formula is given by:
\[x = \frac{-b \pm \sqrt{b^2 - 4ac}}{2a}\]

\subsection{Matrix Example}
A sample matrix:
\[
\begin{pmatrix}
a_{11} & a_{12} & a_{13} \\
a_{21} & a_{22} & a_{23} \\
a_{31} & a_{32} & a_{33}
\end{pmatrix}
\]

\section{Tables}
Below is a sample table showing data:

\begin{table}[h]
\centering
\begin{tabular}{@{}lrrr@{}}
\toprule
\textbf{Category} & \textbf{Value 1} & \textbf{Value 2} & \textbf{Total} \\
\midrule
Item A & 10.5 & 20.3 & 30.8 \\
Item B & 15.2 & 25.1 & 40.3 \\
Item C & 22.1 & 12.7 & 34.8 \\
\bottomrule
\end{tabular}
\caption{Sample Data Table}
\label{tab:sample}
\end{table}

\section{Lists}
\subsection{Bullet Points}
\begin{itemize}
  \item First item
  \item Second item
  \item Third item with sub-items:
    \begin{itemize}
      \item Sub-item 1
      \item Sub-item 2
    \end{itemize}
\end{itemize}

\subsection{Numbered List}
\begin{enumerate}
  \item First step
  \item Second step
  \item Third step
\end{enumerate}

\section{Citations and References}
When you use the citation search feature, references will be automatically managed. For example, you might cite a paper like this: \cite{example2024}. All citations will be automatically added to your \texttt{cite.bib} file.

\section{Conclusion}
This template demonstrates basic \LaTeX{} features including mathematical formulas, tables, and lists. The bibliography system is now set up to work seamlessly with Bibby AI's citation management. Feel free to modify it for your needs!

\bibliographystyle{plain}
\bibliography{cite}

\end{document}\documentclass{article}
\usepackage[utf8]{inputenc}
\usepackage{booktabs}
\usepackage{amsmath}
\usepackage{natbib}

\title{My Document}
\author{Your Name}
\date{\today}

\begin{document}

\maketitle

\section{Introduction}
This is a sample document demonstrating various \LaTeX{} features including mathematical equations, tables, and formatting. You can cite papers using the citation search feature, and they will be automatically added to your \texttt{cite.bib} file.

\section{Mathematical Expressions}
Here are some example equations:

\subsection{Quadratic Formula}
The quadratic formula is given by:
\[x = \frac{-b \pm \sqrt{b^2 - 4ac}}{2a}\]

\subsection{Matrix Example}
A sample matrix:
\[
\begin{pmatrix}
a_{11} & a_{12} & a_{13} \\
a_{21} & a_{22} & a_{23} \\
a_{31} & a_{32} & a_{33}
\end{pmatrix}
\]

\section{Tables}
Below is a sample table showing data:

\begin{table}[h]
\centering
\begin{tabular}{@{}lrrr@{}}
\toprule
\textbf{Category} & \textbf{Value 1} & \textbf{Value 2} & \textbf{Total} \\
\midrule
Item A & 10.5 & 20.3 & 30.8 \\
Item B & 15.2 & 25.1 & 40.3 \\
Item C & 22.1 & 12.7 & 34.8 \\
\bottomrule
\end{tabular}
\caption{Sample Data Table}
\label{tab:sample}
\end{table}

\section{Lists}
\subsection{Bullet Points}
\begin{itemize}
  \item First item
  \item Second item
  \item Third item with sub-items:
    \begin{itemize}
      \item Sub-item 1
      \item Sub-item 2
    \end{itemize}
\end{itemize}

\subsection{Numbered List}
\begin{enumerate}
  \item First step
  \item Second step
  \item Third step
\end{enumerate}

\section{Citations and References}
When you use the citation search feature, references will be automatically managed. For example, you might cite a paper like this: \cite{example2024}. All citations will be automatically added to your \texttt{cite.bib} file.

\section{Conclusion}
This template demonstrates basic \LaTeX{} features including mathematical formulas, tables, and lists. The bibliography system is now set up to work seamlessly with Bibby AI's citation management. Feel free to modify it for your needs!

\bibliographystyle{plain}
\bibliography{cite}

\end{document}\documentclass{article}
\usepackage[utf8]{inputenc}
\usepackage{booktabs}
\usepackage{amsmath}
\usepackage{natbib}

\title{My Document}
\author{Your Name}
\date{\today}

\begin{document}

\maketitle

\section{Introduction}
This is a sample document demonstrating various \LaTeX{} features including mathematical equations, tables, and formatting. You can cite papers using the citation search feature, and they will be automatically added to your \texttt{cite.bib} file.

\section{Mathematical Expressions}
Here are some example equations:

\subsection{Quadratic Formula}
The quadratic formula is given by:
\[x = \frac{-b \pm \sqrt{b^2 - 4ac}}{2a}\]

\subsection{Matrix Example}
A sample matrix:
\[
\begin{pmatrix}
a_{11} & a_{12} & a_{13} \\
a_{21} & a_{22} & a_{23} \\
a_{31} & a_{32} & a_{33}
\end{pmatrix}
\]

\section{Tables}
Below is a sample table showing data:

\begin{table}[h]
\centering
\begin{tabular}{@{}lrrr@{}}
\toprule
\textbf{Category} & \textbf{Value 1} & \textbf{Value 2} & \textbf{Total} \\
\midrule
Item A & 10.5 & 20.3 & 30.8 \\
Item B & 15.2 & 25.1 & 40.3 \\
Item C & 22.1 & 12.7 & 34.8 \\
\bottomrule
\end{tabular}
\caption{Sample Data Table}
\label{tab:sample}
\end{table}

\section{Lists}
\subsection{Bullet Points}
\begin{itemize}
  \item First item
  \item Second item
  \item Third item with sub-items:
    \begin{itemize}
      \item Sub-item 1
      \item Sub-item 2
    \end{itemize}
\end{itemize}

\subsection{Numbered List}
\begin{enumerate}
  \item First step
  \item Second step
  \item Third step
\end{enumerate}

\section{Citations and References}
When you use the citation search feature, references will be automatically managed. For example, you might cite a paper like this: \cite{example2024}. All citations will be automatically added to your \texttt{cite.bib} file.

\section{Conclusion}
This template demonstrates basic \LaTeX{} features including mathematical formulas, tables, and lists. The bibliography system is now set up to work seamlessly with Bibby AI's citation management. Feel free to modify it for your needs!

\bibliographystyle{plain}
\bibliography{cite}

\end{document}\documentclass{article}
\usepackage[utf8]{inputenc}
\usepackage{booktabs}
\usepackage{amsmath}
\usepackage{natbib}

\title{My Document}
\author{Your Name}
\date{\today}

\begin{document}

\maketitle

\section{Introduction}
This is a sample document demonstrating various \LaTeX{} features including mathematical equations, tables, and formatting. You can cite papers using the citation search feature, and they will be automatically added to your \texttt{cite.bib} file.

\section{Mathematical Expressions}
Here are some example equations:

\subsection{Quadratic Formula}
The quadratic formula is given by:
\[x = \frac{-b \pm \sqrt{b^2 - 4ac}}{2a}\]

\subsection{Matrix Example}
A sample matrix:
\[
\begin{pmatrix}
a_{11} & a_{12} & a_{13} \\
a_{21} & a_{22} & a_{23} \\
a_{31} & a_{32} & a_{33}
\end{pmatrix}
\]

\section{Tables}
Below is a sample table showing data:

\begin{table}[h]
\centering
\begin{tabular}{@{}lrrr@{}}
\toprule
\textbf{Category} & \textbf{Value 1} & \textbf{Value 2} & \textbf{Total} \\
\midrule
Item A & 10.5 & 20.3 & 30.8 \\
Item B & 15.2 & 25.1 & 40.3 \\
Item C & 22.1 & 12.7 & 34.8 \\
\bottomrule
\end{tabular}
\caption{Sample Data Table}
\label{tab:sample}
\end{table}

\section{Lists}
\subsection{Bullet Points}
\begin{itemize}
  \item First item
  \item Second item
  \item Third item with sub-items:
    \begin{itemize}
      \item Sub-item 1
      \item Sub-item 2
    \end{itemize}
\end{itemize}

\subsection{Numbered List}
\begin{enumerate}
  \item First step
  \item Second step
  \item Third step
\end{enumerate}

\section{Citations and References}
When you use the citation search feature, references will be automatically managed. For example, you might cite a paper like this: \cite{example2024}. All citations will be automatically added to your \texttt{cite.bib} file.

\section{Conclusion}
This template demonstrates basic \LaTeX{} features including mathematical formulas, tables, and lists. The bibliography system is now set up to work seamlessly with Bibby AI's citation management. Feel free to modify it for your needs!

\bibliographystyle{plain}
\bibliography{cite}

\end{document}\documentclass{article}
\usepackage[utf8]{inputenc}
\usepackage{booktabs}
\usepackage{amsmath}
\usepackage{natbib}

\title{My Document}
\author{Your Name}
\date{\today}

\begin{document}

\maketitle

\section{Introduction}
This is a sample document demonstrating various \LaTeX{} features including mathematical equations, tables, and formatting. You can cite papers using the citation search feature, and they will be automatically added to your \texttt{cite.bib} file.

\section{Mathematical Expressions}
Here are some example equations:

\subsection{Quadratic Formula}
The quadratic formula is given by:
\[x = \frac{-b \pm \sqrt{b^2 - 4ac}}{2a}\]

\subsection{Matrix Example}
A sample matrix:
\[
\begin{pmatrix}
a_{11} & a_{12} & a_{13} \\
a_{21} & a_{22} & a_{23} \\
a_{31} & a_{32} & a_{33}
\end{pmatrix}
\]

\section{Tables}
Below is a sample table showing data:

\begin{table}[h]
\centering
\begin{tabular}{@{}lrrr@{}}
\toprule
\textbf{Category} & \textbf{Value 1} & \textbf{Value 2} & \textbf{Total} \\
\midrule
Item A & 10.5 & 20.3 & 30.8 \\
Item B & 15.2 & 25.1 & 40.3 \\
Item C & 22.1 & 12.7 & 34.8 \\
\bottomrule
\end{tabular}
\caption{Sample Data Table}
\label{tab:sample}
\end{table}

\section{Lists}
\subsection{Bullet Points}
\begin{itemize}
  \item First item
  \item Second item
  \item Third item with sub-items:
    \begin{itemize}
      \item Sub-item 1
      \item Sub-item 2
    \end{itemize}
\end{itemize}

\subsection{Numbered List}
\begin{enumerate}
  \item First step
  \item Second step
  \item Third step
\end{enumerate}

\section{Citations and References}
When you use the citation search feature, references will be automatically managed. For example, you might cite a paper like this: \cite{example2024}. All citations will be automatically added to your \texttt{cite.bib} file.

\section{Conclusion}
This template demonstrates basic \LaTeX{} features including mathematical formulas, tables, and lists. The bibliography system is now set up to work seamlessly with Bibby AI's citation management. Feel free to modify it for your needs!

\bibliographystyle{plain}
\bibliography{cite}

\end{document}\documentclass{article}
\usepackage[utf8]{inputenc}
\usepackage{booktabs}
\usepackage{amsmath}
\usepackage{natbib}

\title{My Document}
\author{Your Name}
\date{\today}

\begin{document}

\maketitle

\section{Introduction}
This is a sample document demonstrating various \LaTeX{} features including mathematical equations, tables, and formatting. You can cite papers using the citation search feature, and they will be automatically added to your \texttt{cite.bib} file.

\section{Mathematical Expressions}
Here are some example equations:

\subsection{Quadratic Formula}
The quadratic formula is given by:
\[x = \frac{-b \pm \sqrt{b^2 - 4ac}}{2a}\]

\subsection{Matrix Example}
A sample matrix:
\[
\begin{pmatrix}
a_{11} & a_{12} & a_{13} \\
a_{21} & a_{22} & a_{23} \\
a_{31} & a_{32} & a_{33}
\end{pmatrix}
\]

\section{Tables}
Below is a sample table showing data:

\begin{table}[h]
\centering
\begin{tabular}{@{}lrrr@{}}
\toprule
\textbf{Category} & \textbf{Value 1} & \textbf{Value 2} & \textbf{Total} \\
\midrule
Item A & 10.5 & 20.3 & 30.8 \\
Item B & 15.2 & 25.1 & 40.3 \\
Item C & 22.1 & 12.7 & 34.8 \\
\bottomrule
\end{tabular}
\caption{Sample Data Table}
\label{tab:sample}
\end{table}

\section{Lists}
\subsection{Bullet Points}
\begin{itemize}
  \item First item
  \item Second item
  \item Third item with sub-items:
    \begin{itemize}
      \item Sub-item 1
      \item Sub-item 2
    \end{itemize}
\end{itemize}

\subsection{Numbered List}
\begin{enumerate}
  \item First step
  \item Second step
  \item Third step
\end{enumerate}

\section{Citations and References}
When you use the citation search feature, references will be automatically managed. For example, you might cite a paper like this: \cite{example2024}. All citations will be automatically added to your \texttt{cite.bib} file.

\section{Conclusion}
This template demonstrates basic \LaTeX{} features including mathematical formulas, tables, and lists. The bibliography system is now set up to work seamlessly with Bibby AI's citation management. Feel free to modify it for your needs!

\bibliographystyle{plain}
\bibliography{cite}

\end{document}\documentclass{article}
\usepackage[utf8]{inputenc}
\usepackage{booktabs}
\usepackage{amsmath}
\usepackage{natbib}

\title{My Document}
\author{Your Name}
\date{\today}

\begin{document}

\maketitle

\section{Introduction}
This is a sample document demonstrating various \LaTeX{} features including mathematical equations, tables, and formatting. You can cite papers using the citation search feature, and they will be automatically added to your \texttt{cite.bib} file.

\section{Mathematical Expressions}
Here are some example equations:

\subsection{Quadratic Formula}
The quadratic formula is given by:
\[x = \frac{-b \pm \sqrt{b^2 - 4ac}}{2a}\]

\subsection{Matrix Example}
A sample matrix:
\[
\begin{pmatrix}
a_{11} & a_{12} & a_{13} \\
a_{21} & a_{22} & a_{23} \\
a_{31} & a_{32} & a_{33}
\end{pmatrix}
\]

\section{Tables}
Below is a sample table showing data:

\begin{table}[h]
\centering
\begin{tabular}{@{}lrrr@{}}
\toprule
\textbf{Category} & \textbf{Value 1} & \textbf{Value 2} & \textbf{Total} \\
\midrule
Item A & 10.5 & 20.3 & 30.8 \\
Item B & 15.2 & 25.1 & 40.3 \\
Item C & 22.1 & 12.7 & 34.8 \\
\bottomrule
\end{tabular}
\caption{Sample Data Table}
\label{tab:sample}
\end{table}

\section{Lists}
\subsection{Bullet Points}
\begin{itemize}
  \item First item
  \item Second item
  \item Third item with sub-items:
    \begin{itemize}
      \item Sub-item 1
      \item Sub-item 2
    \end{itemize}
\end{itemize}

\subsection{Numbered List}
\begin{enumerate}
  \item First step
  \item Second step
  \item Third step
\end{enumerate}

\section{Citations and References}
When you use the citation search feature, references will be automatically managed. For example, you might cite a paper like this: \cite{example2024}. All citations will be automatically added to your \texttt{cite.bib} file.

\section{Conclusion}
This template demonstrates basic \LaTeX{} features including mathematical formulas, tables, and lists. The bibliography system is now set up to work seamlessly with Bibby AI's citation management. Feel free to modify it for your needs!

\bibliographystyle{plain}
\bibliography{cite}

\end{document}\documentclass{article}
\usepackage[utf8]{inputenc}
\usepackage{booktabs}
\usepackage{amsmath}
\usepackage{natbib}

\title{My Document}
\author{Your Name}
\date{\today}

\begin{document}

\maketitle

\section{Introduction}
This is a sample document demonstrating various \LaTeX{} features including mathematical equations, tables, and formatting. You can cite papers using the citation search feature, and they will be automatically added to your \texttt{cite.bib} file.

\section{Mathematical Expressions}
Here are some example equations:

\subsection{Quadratic Formula}
The quadratic formula is given by:
\[x = \frac{-b \pm \sqrt{b^2 - 4ac}}{2a}\]

\subsection{Matrix Example}
A sample matrix:
\[
\begin{pmatrix}
a_{11} & a_{12} & a_{13} \\
a_{21} & a_{22} & a_{23} \\
a_{31} & a_{32} & a_{33}
\end{pmatrix}
\]

\section{Tables}
Below is a sample table showing data:

\begin{table}[h]
\centering
\begin{tabular}{@{}lrrr@{}}
\toprule
\textbf{Category} & \textbf{Value 1} & \textbf{Value 2} & \textbf{Total} \\
\midrule
Item A & 10.5 & 20.3 & 30.8 \\
Item B & 15.2 & 25.1 & 40.3 \\
Item C & 22.1 & 12.7 & 34.8 \\
\bottomrule
\end{tabular}
\caption{Sample Data Table}
\label{tab:sample}
\end{table}

\section{Lists}
\subsection{Bullet Points}
\begin{itemize}
  \item First item
  \item Second item
  \item Third item with sub-items:
    \begin{itemize}
      \item Sub-item 1
      \item Sub-item 2
    \end{itemize}
\end{itemize}

\subsection{Numbered List}
\begin{enumerate}
  \item First step
  \item Second step
  \item Third step
\end{enumerate}

\section{Citations and References}
When you use the citation search feature, references will be automatically managed. For example, you might cite a paper like this: \cite{example2024}. All citations will be automatically added to your \texttt{cite.bib} file.

\section{Conclusion}
This template demonstrates basic \LaTeX{} features including mathematical formulas, tables, and lists. The bibliography system is now set up to work seamlessly with Bibby AI's citation management. Feel free to modify it for your needs!

\bibliographystyle{plain}
\bibliography{cite}

\end{document}\documentclass{article}
\usepackage[utf8]{inputenc}
\usepackage{booktabs}
\usepackage{amsmath}
\usepackage{natbib}

\title{My Document}
\author{Your Name}
\date{\today}

\begin{document}

\maketitle

\section{Introduction}
This is a sample document demonstrating various \LaTeX{} features including mathematical equations, tables, and formatting. You can cite papers using the citation search feature, and they will be automatically added to your \texttt{cite.bib} file.

\section{Mathematical Expressions}
Here are some example equations:

\subsection{Quadratic Formula}
The quadratic formula is given by:
\[x = \frac{-b \pm \sqrt{b^2 - 4ac}}{2a}\]

\subsection{Matrix Example}
A sample matrix:
\[
\begin{pmatrix}
a_{11} & a_{12} & a_{13} \\
a_{21} & a_{22} & a_{23} \\
a_{31} & a_{32} & a_{33}
\end{pmatrix}
\]

\section{Tables}
Below is a sample table showing data:

\begin{table}[h]
\centering
\begin{tabular}{@{}lrrr@{}}
\toprule
\textbf{Category} & \textbf{Value 1} & \textbf{Value 2} & \textbf{Total} \\
\midrule
Item A & 10.5 & 20.3 & 30.8 \\
Item B & 15.2 & 25.1 & 40.3 \\
Item C & 22.1 & 12.7 & 34.8 \\
\bottomrule
\end{tabular}
\caption{Sample Data Table}
\label{tab:sample}
\end{table}

\section{Lists}
\subsection{Bullet Points}
\begin{itemize}
  \item First item
  \item Second item
  \item Third item with sub-items:
    \begin{itemize}
      \item Sub-item 1
      \item Sub-item 2
    \end{itemize}
\end{itemize}

\subsection{Numbered List}
\begin{enumerate}
  \item First step
  \item Second step
  \item Third step
\end{enumerate}

\section{Citations and References}
When you use the citation search feature, references will be automatically managed. For example, you might cite a paper like this: \cite{example2024}. All citations will be automatically added to your \texttt{cite.bib} file.

\section{Conclusion}
This template demonstrates basic \LaTeX{} features including mathematical formulas, tables, and lists. The bibliography system is now set up to work seamlessly with Bibby AI's citation management. Feel free to modify it for your needs!

\bibliographystyle{plain}
\bibliography{cite}

\end{document}\documentclass{article}
\usepackage[utf8]{inputenc}
\usepackage{booktabs}
\usepackage{amsmath}
\usepackage{natbib}

\title{My Document}
\author{Your Name}
\date{\today}

\begin{document}

\maketitle

\section{Introduction}
This is a sample document demonstrating various \LaTeX{} features including mathematical equations, tables, and formatting. You can cite papers using the citation search feature, and they will be automatically added to your \texttt{cite.bib} file.

\section{Mathematical Expressions}
Here are some example equations:

\subsection{Quadratic Formula}
The quadratic formula is given by:
\[x = \frac{-b \pm \sqrt{b^2 - 4ac}}{2a}\]

\subsection{Matrix Example}
A sample matrix:
\[
\begin{pmatrix}
a_{11} & a_{12} & a_{13} \\
a_{21} & a_{22} & a_{23} \\
a_{31} & a_{32} & a_{33}
\end{pmatrix}
\]

\section{Tables}
Below is a sample table showing data:

\begin{table}[h]
\centering
\begin{tabular}{@{}lrrr@{}}
\toprule
\textbf{Category} & \textbf{Value 1} & \textbf{Value 2} & \textbf{Total} \\
\midrule
Item A & 10.5 & 20.3 & 30.8 \\
Item B & 15.2 & 25.1 & 40.3 \\
Item C & 22.1 & 12.7 & 34.8 \\
\bottomrule
\end{tabular}
\caption{Sample Data Table}
\label{tab:sample}
\end{table}

\section{Lists}
\subsection{Bullet Points}
\begin{itemize}
  \item First item
  \item Second item
  \item Third item with sub-items:
    \begin{itemize}
      \item Sub-item 1
      \item Sub-item 2
    \end{itemize}
\end{itemize}

\subsection{Numbered List}
\begin{enumerate}
  \item First step
  \item Second step
  \item Third step
\end{enumerate}

\section{Citations and References}
When you use the citation search feature, references will be automatically managed. For example, you might cite a paper like this: \cite{example2024}. All citations will be automatically added to your \texttt{cite.bib} file.

\section{Conclusion}
This template demonstrates basic \LaTeX{} features including mathematical formulas, tables, and lists. The bibliography system is now set up to work seamlessly with Bibby AI's citation management. Feel free to modify it for your needs!

\bibliographystyle{plain}
\bibliography{cite}

\end{document}\documentclass{article}
\usepackage[utf8]{inputenc}
\usepackage{booktabs}
\usepackage{amsmath}
\usepackage{natbib}

\title{My Document}
\author{Your Name}
\date{\today}

\begin{document}

\maketitle

\section{Introduction}
This is a sample document demonstrating various \LaTeX{} features including mathematical equations, tables, and formatting. You can cite papers using the citation search feature, and they will be automatically added to your \texttt{cite.bib} file.

\section{Mathematical Expressions}
Here are some example equations:

\subsection{Quadratic Formula}
The quadratic formula is given by:
\[x = \frac{-b \pm \sqrt{b^2 - 4ac}}{2a}\]

\subsection{Matrix Example}
A sample matrix:
\[
\begin{pmatrix}
a_{11} & a_{12} & a_{13} \\
a_{21} & a_{22} & a_{23} \\
a_{31} & a_{32} & a_{33}
\end{pmatrix}
\]

\section{Tables}
Below is a sample table showing data:

\begin{table}[h]
\centering
\begin{tabular}{@{}lrrr@{}}
\toprule
\textbf{Category} & \textbf{Value 1} & \textbf{Value 2} & \textbf{Total} \\
\midrule
Item A & 10.5 & 20.3 & 30.8 \\
Item B & 15.2 & 25.1 & 40.3 \\
Item C & 22.1 & 12.7 & 34.8 \\
\bottomrule
\end{tabular}
\caption{Sample Data Table}
\label{tab:sample}
\end{table}

\section{Lists}
\subsection{Bullet Points}
\begin{itemize}
  \item First item
  \item Second item
  \item Third item with sub-items:
    \begin{itemize}
      \item Sub-item 1
      \item Sub-item 2
    \end{itemize}
\end{itemize}

\subsection{Numbered List}
\begin{enumerate}
  \item First step
  \item Second step
  \item Third step
\end{enumerate}

\section{Citations and References}
When you use the citation search feature, references will be automatically managed. For example, you might cite a paper like this: \cite{example2024}. All citations will be automatically added to your \texttt{cite.bib} file.

\section{Conclusion}
This template demonstrates basic \LaTeX{} features including mathematical formulas, tables, and lists. The bibliography system is now set up to work seamlessly with Bibby AI's citation management. Feel free to modify it for your needs!

\bibliographystyle{plain}
\bibliography{cite}

\end{document}\documentclass{article}
\usepackage[utf8]{inputenc}
\usepackage{booktabs}
\usepackage{amsmath}
\usepackage{natbib}

\title{My Document}
\author{Your Name}
\date{\today}

\begin{document}

\maketitle

\section{Introduction}
This is a sample document demonstrating various \LaTeX{} features including mathematical equations, tables, and formatting. You can cite papers using the citation search feature, and they will be automatically added to your \texttt{cite.bib} file.

\section{Mathematical Expressions}
Here are some example equations:

\subsection{Quadratic Formula}
The quadratic formula is given by:
\[x = \frac{-b \pm \sqrt{b^2 - 4ac}}{2a}\]

\subsection{Matrix Example}
A sample matrix:
\[
\begin{pmatrix}
a_{11} & a_{12} & a_{13} \\
a_{21} & a_{22} & a_{23} \\
a_{31} & a_{32} & a_{33}
\end{pmatrix}
\]

\section{Tables}
Below is a sample table showing data:

\begin{table}[h]
\centering
\begin{tabular}{@{}lrrr@{}}
\toprule
\textbf{Category} & \textbf{Value 1} & \textbf{Value 2} & \textbf{Total} \\
\midrule
Item A & 10.5 & 20.3 & 30.8 \\
Item B & 15.2 & 25.1 & 40.3 \\
Item C & 22.1 & 12.7 & 34.8 \\
\bottomrule
\end{tabular}
\caption{Sample Data Table}
\label{tab:sample}
\end{table}

\section{Lists}
\subsection{Bullet Points}
\begin{itemize}
  \item First item
  \item Second item
  \item Third item with sub-items:
    \begin{itemize}
      \item Sub-item 1
      \item Sub-item 2
    \end{itemize}
\end{itemize}

\subsection{Numbered List}
\begin{enumerate}
  \item First step
  \item Second step
  \item Third step
\end{enumerate}

\section{Citations and References}
When you use the citation search feature, references will be automatically managed. For example, you might cite a paper like this: \cite{example2024}. All citations will be automatically added to your \texttt{cite.bib} file.

\section{Conclusion}
This template demonstrates basic \LaTeX{} features including mathematical formulas, tables, and lists. The bibliography system is now set up to work seamlessly with Bibby AI's citation management. Feel free to modify it for your needs!

\bibliographystyle{plain}
\bibliography{cite}

\end{document}\documentclass{article}
\usepackage[utf8]{inputenc}
\usepackage{booktabs}
\usepackage{amsmath}
\usepackage{natbib}

\title{My Document}
\author{Your Name}
\date{\today}

\begin{document}

\maketitle

\section{Introduction}
This is a sample document demonstrating various \LaTeX{} features including mathematical equations, tables, and formatting. You can cite papers using the citation search feature, and they will be automatically added to your \texttt{cite.bib} file.

\section{Mathematical Expressions}
Here are some example equations:

\subsection{Quadratic Formula}
The quadratic formula is given by:
\[x = \frac{-b \pm \sqrt{b^2 - 4ac}}{2a}\]

\subsection{Matrix Example}
A sample matrix:
\[
\begin{pmatrix}
a_{11} & a_{12} & a_{13} \\
a_{21} & a_{22} & a_{23} \\
a_{31} & a_{32} & a_{33}
\end{pmatrix}
\]

\section{Tables}
Below is a sample table showing data:

\begin{table}[h]
\centering
\begin{tabular}{@{}lrrr@{}}
\toprule
\textbf{Category} & \textbf{Value 1} & \textbf{Value 2} & \textbf{Total} \\
\midrule
Item A & 10.5 & 20.3 & 30.8 \\
Item B & 15.2 & 25.1 & 40.3 \\
Item C & 22.1 & 12.7 & 34.8 \\
\bottomrule
\end{tabular}
\caption{Sample Data Table}
\label{tab:sample}
\end{table}

\section{Lists}
\subsection{Bullet Points}
\begin{itemize}
  \item First item
  \item Second item
  \item Third item with sub-items:
    \begin{itemize}
      \item Sub-item 1
      \item Sub-item 2
    \end{itemize}
\end{itemize}

\subsection{Numbered List}
\begin{enumerate}
  \item First step
  \item Second step
  \item Third step
\end{enumerate}

\section{Citations and References}
When you use the citation search feature, references will be automatically managed. For example, you might cite a paper like this: \cite{example2024}. All citations will be automatically added to your \texttt{cite.bib} file.

\section{Conclusion}
This template demonstrates basic \LaTeX{} features including mathematical formulas, tables, and lists. The bibliography system is now set up to work seamlessly with Bibby AI's citation management. Feel free to modify it for your needs!

\bibliographystyle{plain}
\bibliography{cite}

\end{document}\documentclass{article}
\usepackage[utf8]{inputenc}
\usepackage{booktabs}
\usepackage{amsmath}
\usepackage{natbib}

\title{My Document}
\author{Your Name}
\date{\today}

\begin{document}

\maketitle

\section{Introduction}
This is a sample document demonstrating various \LaTeX{} features including mathematical equations, tables, and formatting. You can cite papers using the citation search feature, and they will be automatically added to your \texttt{cite.bib} file.

\section{Mathematical Expressions}
Here are some example equations:

\subsection{Quadratic Formula}
The quadratic formula is given by:
\[x = \frac{-b \pm \sqrt{b^2 - 4ac}}{2a}\]

\subsection{Matrix Example}
A sample matrix:
\[
\begin{pmatrix}
a_{11} & a_{12} & a_{13} \\
a_{21} & a_{22} & a_{23} \\
a_{31} & a_{32} & a_{33}
\end{pmatrix}
\]

\section{Tables}
Below is a sample table showing data:

\begin{table}[h]
\centering
\begin{tabular}{@{}lrrr@{}}
\toprule
\textbf{Category} & \textbf{Value 1} & \textbf{Value 2} & \textbf{Total} \\
\midrule
Item A & 10.5 & 20.3 & 30.8 \\
Item B & 15.2 & 25.1 & 40.3 \\
Item C & 22.1 & 12.7 & 34.8 \\
\bottomrule
\end{tabular}
\caption{Sample Data Table}
\label{tab:sample}
\end{table}

\section{Lists}
\subsection{Bullet Points}
\begin{itemize}
  \item First item
  \item Second item
  \item Third item with sub-items:
    \begin{itemize}
      \item Sub-item 1
      \item Sub-item 2
    \end{itemize}
\end{itemize}

\subsection{Numbered List}
\begin{enumerate}
  \item First step
  \item Second step
  \item Third step
\end{enumerate}

\section{Citations and References}
When you use the citation search feature, references will be automatically managed. For example, you might cite a paper like this: \cite{example2024}. All citations will be automatically added to your \texttt{cite.bib} file.

\section{Conclusion}
This template demonstrates basic \LaTeX{} features including mathematical formulas, tables, and lists. The bibliography system is now set up to work seamlessly with Bibby AI's citation management. Feel free to modify it for your needs!

\bibliographystyle{plain}
\bibliography{cite}

\end{document}\documentclass{article}
\usepackage[utf8]{inputenc}
\usepackage{booktabs}
\usepackage{amsmath}
\usepackage{natbib}

\title{My Document}
\author{Your Name}
\date{\today}

\begin{document}

\maketitle

\section{Introduction}
This is a sample document demonstrating various \LaTeX{} features including mathematical equations, tables, and formatting. You can cite papers using the citation search feature, and they will be automatically added to your \texttt{cite.bib} file.

\section{Mathematical Expressions}
Here are some example equations:

\subsection{Quadratic Formula}
The quadratic formula is given by:
\[x = \frac{-b \pm \sqrt{b^2 - 4ac}}{2a}\]

\subsection{Matrix Example}
A sample matrix:
\[
\begin{pmatrix}
a_{11} & a_{12} & a_{13} \\
a_{21} & a_{22} & a_{23} \\
a_{31} & a_{32} & a_{33}
\end{pmatrix}
\]

\section{Tables}
Below is a sample table showing data:

\begin{table}[h]
\centering
\begin{tabular}{@{}lrrr@{}}
\toprule
\textbf{Category} & \textbf{Value 1} & \textbf{Value 2} & \textbf{Total} \\
\midrule
Item A & 10.5 & 20.3 & 30.8 \\
Item B & 15.2 & 25.1 & 40.3 \\
Item C & 22.1 & 12.7 & 34.8 \\
\bottomrule
\end{tabular}
\caption{Sample Data Table}
\label{tab:sample}
\end{table}

\section{Lists}
\subsection{Bullet Points}
\begin{itemize}
  \item First item
  \item Second item
  \item Third item with sub-items:
    \begin{itemize}
      \item Sub-item 1
      \item Sub-item 2
    \end{itemize}
\end{itemize}

\subsection{Numbered List}
\begin{enumerate}
  \item First step
  \item Second step
  \item Third step
\end{enumerate}

\section{Citations and References}
When you use the citation search feature, references will be automatically managed. For example, you might cite a paper like this: \cite{example2024}. All citations will be automatically added to your \texttt{cite.bib} file.

\section{Conclusion}
This template demonstrates basic \LaTeX{} features including mathematical formulas, tables, and lists. The bibliography system is now set up to work seamlessly with Bibby AI's citation management. Feel free to modify it for your needs!

\bibliographystyle{plain}
\bibliography{cite}

\end{document}\documentclass{article}
\usepackage[utf8]{inputenc}
\usepackage{booktabs}
\usepackage{amsmath}
\usepackage{natbib}

\title{My Document}
\author{Your Name}
\date{\today}

\begin{document}

\maketitle

\section{Introduction}
This is a sample document demonstrating various \LaTeX{} features including mathematical equations, tables, and formatting. You can cite papers using the citation search feature, and they will be automatically added to your \texttt{cite.bib} file.

\section{Mathematical Expressions}
Here are some example equations:

\subsection{Quadratic Formula}
The quadratic formula is given by:
\[x = \frac{-b \pm \sqrt{b^2 - 4ac}}{2a}\]

\subsection{Matrix Example}
A sample matrix:
\[
\begin{pmatrix}
a_{11} & a_{12} & a_{13} \\
a_{21} & a_{22} & a_{23} \\
a_{31} & a_{32} & a_{33}
\end{pmatrix}
\]

\section{Tables}
Below is a sample table showing data:

\begin{table}[h]
\centering
\begin{tabular}{@{}lrrr@{}}
\toprule
\textbf{Category} & \textbf{Value 1} & \textbf{Value 2} & \textbf{Total} \\
\midrule
Item A & 10.5 & 20.3 & 30.8 \\
Item B & 15.2 & 25.1 & 40.3 \\
Item C & 22.1 & 12.7 & 34.8 \\
\bottomrule
\end{tabular}
\caption{Sample Data Table}
\label{tab:sample}
\end{table}

\section{Lists}
\subsection{Bullet Points}
\begin{itemize}
  \item First item
  \item Second item
  \item Third item with sub-items:
    \begin{itemize}
      \item Sub-item 1
      \item Sub-item 2
    \end{itemize}
\end{itemize}

\subsection{Numbered List}
\begin{enumerate}
  \item First step
  \item Second step
  \item Third step
\end{enumerate}

\section{Citations and References}
When you use the citation search feature, references will be automatically managed. For example, you might cite a paper like this: \cite{example2024}. All citations will be automatically added to your \texttt{cite.bib} file.

\section{Conclusion}
This template demonstrates basic \LaTeX{} features including mathematical formulas, tables, and lists. The bibliography system is now set up to work seamlessly with Bibby AI's citation management. Feel free to modify it for your needs!

\bibliographystyle{plain}
\bibliography{cite}

\end{document}\documentclass{article}
\usepackage[utf8]{inputenc}
\usepackage{booktabs}
\usepackage{amsmath}
\usepackage{natbib}

\title{My Document}
\author{Your Name}
\date{\today}

\begin{document}

\maketitle

\section{Introduction}
This is a sample document demonstrating various \LaTeX{} features including mathematical equations, tables, and formatting. You can cite papers using the citation search feature, and they will be automatically added to your \texttt{cite.bib} file.

\section{Mathematical Expressions}
Here are some example equations:

\subsection{Quadratic Formula}
The quadratic formula is given by:
\[x = \frac{-b \pm \sqrt{b^2 - 4ac}}{2a}\]

\subsection{Matrix Example}
A sample matrix:
\[
\begin{pmatrix}
a_{11} & a_{12} & a_{13} \\
a_{21} & a_{22} & a_{23} \\
a_{31} & a_{32} & a_{33}
\end{pmatrix}
\]

\section{Tables}
Below is a sample table showing data:

\begin{table}[h]
\centering
\begin{tabular}{@{}lrrr@{}}
\toprule
\textbf{Category} & \textbf{Value 1} & \textbf{Value 2} & \textbf{Total} \\
\midrule
Item A & 10.5 & 20.3 & 30.8 \\
Item B & 15.2 & 25.1 & 40.3 \\
Item C & 22.1 & 12.7 & 34.8 \\
\bottomrule
\end{tabular}
\caption{Sample Data Table}
\label{tab:sample}
\end{table}

\section{Lists}
\subsection{Bullet Points}
\begin{itemize}
  \item First item
  \item Second item
  \item Third item with sub-items:
    \begin{itemize}
      \item Sub-item 1
      \item Sub-item 2
    \end{itemize}
\end{itemize}

\subsection{Numbered List}
\begin{enumerate}
  \item First step
  \item Second step
  \item Third step
\end{enumerate}

\section{Citations and References}
When you use the citation search feature, references will be automatically managed. For example, you might cite a paper like this: \cite{example2024}. All citations will be automatically added to your \texttt{cite.bib} file.

\section{Conclusion}
This template demonstrates basic \LaTeX{} features including mathematical formulas, tables, and lists. The bibliography system is now set up to work seamlessly with Bibby AI's citation management. Feel free to modify it for your needs!

\bibliographystyle{plain}
\bibliography{cite}

\end{document}\documentclass{article}
\usepackage[utf8]{inputenc}
\usepackage{booktabs}
\usepackage{amsmath}
\usepackage{natbib}

\title{My Document}
\author{Your Name}
\date{\today}

\begin{document}

\maketitle

\section{Introduction}
This is a sample document demonstrating various \LaTeX{} features including mathematical equations, tables, and formatting. You can cite papers using the citation search feature, and they will be automatically added to your \texttt{cite.bib} file.

\section{Mathematical Expressions}
Here are some example equations:

\subsection{Quadratic Formula}
The quadratic formula is given by:
\[x = \frac{-b \pm \sqrt{b^2 - 4ac}}{2a}\]

\subsection{Matrix Example}
A sample matrix:
\[
\begin{pmatrix}
a_{11} & a_{12} & a_{13} \\
a_{21} & a_{22} & a_{23} \\
a_{31} & a_{32} & a_{33}
\end{pmatrix}
\]

\section{Tables}
Below is a sample table showing data:

\begin{table}[h]
\centering
\begin{tabular}{@{}lrrr@{}}
\toprule
\textbf{Category} & \textbf{Value 1} & \textbf{Value 2} & \textbf{Total} \\
\midrule
Item A & 10.5 & 20.3 & 30.8 \\
Item B & 15.2 & 25.1 & 40.3 \\
Item C & 22.1 & 12.7 & 34.8 \\
\bottomrule
\end{tabular}
\caption{Sample Data Table}
\label{tab:sample}
\end{table}

\section{Lists}
\subsection{Bullet Points}
\begin{itemize}
  \item First item
  \item Second item
  \item Third item with sub-items:
    \begin{itemize}
      \item Sub-item 1
      \item Sub-item 2
    \end{itemize}
\end{itemize}

\subsection{Numbered List}
\begin{enumerate}
  \item First step
  \item Second step
  \item Third step
\end{enumerate}

\section{Citations and References}
When you use the citation search feature, references will be automatically managed. For example, you might cite a paper like this: \cite{example2024}. All citations will be automatically added to your \texttt{cite.bib} file.

\section{Conclusion}
This template demonstrates basic \LaTeX{} features including mathematical formulas, tables, and lists. The bibliography system is now set up to work seamlessly with Bibby AI's citation management. Feel free to modify it for your needs!

\bibliographystyle{plain}
\bibliography{cite}

\end{document}\documentclass{article}
\usepackage[utf8]{inputenc}
\usepackage{booktabs}
\usepackage{amsmath}
\usepackage{natbib}

\title{My Document}
\author{Your Name}
\date{\today}

\begin{document}

\maketitle

\section{Introduction}
This is a sample document demonstrating various \LaTeX{} features including mathematical equations, tables, and formatting. You can cite papers using the citation search feature, and they will be automatically added to your \texttt{cite.bib} file.

\section{Mathematical Expressions}
Here are some example equations:

\subsection{Quadratic Formula}
The quadratic formula is given by:
\[x = \frac{-b \pm \sqrt{b^2 - 4ac}}{2a}\]

\subsection{Matrix Example}
A sample matrix:
\[
\begin{pmatrix}
a_{11} & a_{12} & a_{13} \\
a_{21} & a_{22} & a_{23} \\
a_{31} & a_{32} & a_{33}
\end{pmatrix}
\]

\section{Tables}
Below is a sample table showing data:

\begin{table}[h]
\centering
\begin{tabular}{@{}lrrr@{}}
\toprule
\textbf{Category} & \textbf{Value 1} & \textbf{Value 2} & \textbf{Total} \\
\midrule
Item A & 10.5 & 20.3 & 30.8 \\
Item B & 15.2 & 25.1 & 40.3 \\
Item C & 22.1 & 12.7 & 34.8 \\
\bottomrule
\end{tabular}
\caption{Sample Data Table}
\label{tab:sample}
\end{table}

\section{Lists}
\subsection{Bullet Points}
\begin{itemize}
  \item First item
  \item Second item
  \item Third item with sub-items:
    \begin{itemize}
      \item Sub-item 1
      \item Sub-item 2
    \end{itemize}
\end{itemize}

\subsection{Numbered List}
\begin{enumerate}
  \item First step
  \item Second step
  \item Third step
\end{enumerate}

\section{Citations and References}
When you use the citation search feature, references will be automatically managed. For example, you might cite a paper like this: \cite{example2024}. All citations will be automatically added to your \texttt{cite.bib} file.

\section{Conclusion}
This template demonstrates basic \LaTeX{} features including mathematical formulas, tables, and lists. The bibliography system is now set up to work seamlessly with Bibby AI's citation management. Feel free to modify it for your needs!

\bibliographystyle{plain}
\bibliography{cite}

\end{document}\documentclass{article}
\usepackage[utf8]{inputenc}
\usepackage{booktabs}
\usepackage{amsmath}
\usepackage{natbib}

\title{My Document}
\author{Your Name}
\date{\today}

\begin{document}

\maketitle

\section{Introduction}
This is a sample document demonstrating various \LaTeX{} features including mathematical equations, tables, and formatting. You can cite papers using the citation search feature, and they will be automatically added to your \texttt{cite.bib} file.

\section{Mathematical Expressions}
Here are some example equations:

\subsection{Quadratic Formula}
The quadratic formula is given by:
\[x = \frac{-b \pm \sqrt{b^2 - 4ac}}{2a}\]

\subsection{Matrix Example}
A sample matrix:
\[
\begin{pmatrix}
a_{11} & a_{12} & a_{13} \\
a_{21} & a_{22} & a_{23} \\
a_{31} & a_{32} & a_{33}
\end{pmatrix}
\]

\section{Tables}
Below is a sample table showing data:

\begin{table}[h]
\centering
\begin{tabular}{@{}lrrr@{}}
\toprule
\textbf{Category} & \textbf{Value 1} & \textbf{Value 2} & \textbf{Total} \\
\midrule
Item A & 10.5 & 20.3 & 30.8 \\
Item B & 15.2 & 25.1 & 40.3 \\
Item C & 22.1 & 12.7 & 34.8 \\
\bottomrule
\end{tabular}
\caption{Sample Data Table}
\label{tab:sample}
\end{table}

\section{Lists}
\subsection{Bullet Points}
\begin{itemize}
  \item First item
  \item Second item
  \item Third item with sub-items:
    \begin{itemize}
      \item Sub-item 1
      \item Sub-item 2
    \end{itemize}
\end{itemize}

\subsection{Numbered List}
\begin{enumerate}
  \item First step
  \item Second step
  \item Third step
\end{enumerate}

\section{Citations and References}
When you use the citation search feature, references will be automatically managed. For example, you might cite a paper like this: \cite{example2024}. All citations will be automatically added to your \texttt{cite.bib} file.

\section{Conclusion}
This template demonstrates basic \LaTeX{} features including mathematical formulas, tables, and lists. The bibliography system is now set up to work seamlessly with Bibby AI's citation management. Feel free to modify it for your needs!

\bibliographystyle{plain}
\bibliography{cite}

\end{document}\documentclass{article}
\usepackage[utf8]{inputenc}
\usepackage{booktabs}
\usepackage{amsmath}
\usepackage{natbib}

\title{My Document}
\author{Your Name}
\date{\today}

\begin{document}

\maketitle

\section{Introduction}
This is a sample document demonstrating various \LaTeX{} features including mathematical equations, tables, and formatting. You can cite papers using the citation search feature, and they will be automatically added to your \texttt{cite.bib} file.

\section{Mathematical Expressions}
Here are some example equations:

\subsection{Quadratic Formula}
The quadratic formula is given by:
\[x = \frac{-b \pm \sqrt{b^2 - 4ac}}{2a}\]

\subsection{Matrix Example}
A sample matrix:
\[
\begin{pmatrix}
a_{11} & a_{12} & a_{13} \\
a_{21} & a_{22} & a_{23} \\
a_{31} & a_{32} & a_{33}
\end{pmatrix}
\]

\section{Tables}
Below is a sample table showing data:

\begin{table}[h]
\centering
\begin{tabular}{@{}lrrr@{}}
\toprule
\textbf{Category} & \textbf{Value 1} & \textbf{Value 2} & \textbf{Total} \\
\midrule
Item A & 10.5 & 20.3 & 30.8 \\
Item B & 15.2 & 25.1 & 40.3 \\
Item C & 22.1 & 12.7 & 34.8 \\
\bottomrule
\end{tabular}
\caption{Sample Data Table}
\label{tab:sample}
\end{table}

\section{Lists}
\subsection{Bullet Points}
\begin{itemize}
  \item First item
  \item Second item
  \item Third item with sub-items:
    \begin{itemize}
      \item Sub-item 1
      \item Sub-item 2
    \end{itemize}
\end{itemize}

\subsection{Numbered List}
\begin{enumerate}
  \item First step
  \item Second step
  \item Third step
\end{enumerate}

\section{Citations and References}
When you use the citation search feature, references will be automatically managed. For example, you might cite a paper like this: \cite{example2024}. All citations will be automatically added to your \texttt{cite.bib} file.

\section{Conclusion}
This template demonstrates basic \LaTeX{} features including mathematical formulas, tables, and lists. The bibliography system is now set up to work seamlessly with Bibby AI's citation management. Feel free to modify it for your needs!

\bibliographystyle{plain}
\bibliography{cite}

\end{document}\documentclass{article}
\usepackage[utf8]{inputenc}
\usepackage{booktabs}
\usepackage{amsmath}
\usepackage{natbib}

\title{My Document}
\author{Your Name}
\date{\today}

\begin{document}

\maketitle

\section{Introduction}
This is a sample document demonstrating various \LaTeX{} features including mathematical equations, tables, and formatting. You can cite papers using the citation search feature, and they will be automatically added to your \texttt{cite.bib} file.

\section{Mathematical Expressions}
Here are some example equations:

\subsection{Quadratic Formula}
The quadratic formula is given by:
\[x = \frac{-b \pm \sqrt{b^2 - 4ac}}{2a}\]

\subsection{Matrix Example}
A sample matrix:
\[
\begin{pmatrix}
a_{11} & a_{12} & a_{13} \\
a_{21} & a_{22} & a_{23} \\
a_{31} & a_{32} & a_{33}
\end{pmatrix}
\]

\section{Tables}
Below is a sample table showing data:

\begin{table}[h]
\centering
\begin{tabular}{@{}lrrr@{}}
\toprule
\textbf{Category} & \textbf{Value 1} & \textbf{Value 2} & \textbf{Total} \\
\midrule
Item A & 10.5 & 20.3 & 30.8 \\
Item B & 15.2 & 25.1 & 40.3 \\
Item C & 22.1 & 12.7 & 34.8 \\
\bottomrule
\end{tabular}
\caption{Sample Data Table}
\label{tab:sample}
\end{table}

\section{Lists}
\subsection{Bullet Points}
\begin{itemize}
  \item First item
  \item Second item
  \item Third item with sub-items:
    \begin{itemize}
      \item Sub-item 1
      \item Sub-item 2
    \end{itemize}
\end{itemize}

\subsection{Numbered List}
\begin{enumerate}
  \item First step
  \item Second step
  \item Third step
\end{enumerate}

\section{Citations and References}
When you use the citation search feature, references will be automatically managed. For example, you might cite a paper like this: \cite{example2024}. All citations will be automatically added to your \texttt{cite.bib} file.

\section{Conclusion}
This template demonstrates basic \LaTeX{} features including mathematical formulas, tables, and lists. The bibliography system is now set up to work seamlessly with Bibby AI's citation management. Feel free to modify it for your needs!

\bibliographystyle{plain}
\bibliography{cite}

\end{document}\documentclass{article}
\usepackage[utf8]{inputenc}
\usepackage{booktabs}
\usepackage{amsmath}
\usepackage{natbib}

\title{My Document}
\author{Your Name}
\date{\today}

\begin{document}

\maketitle

\section{Introduction}
This is a sample document demonstrating various \LaTeX{} features including mathematical equations, tables, and formatting. You can cite papers using the citation search feature, and they will be automatically added to your \texttt{cite.bib} file.

\section{Mathematical Expressions}
Here are some example equations:

\subsection{Quadratic Formula}
The quadratic formula is given by:
\[x = \frac{-b \pm \sqrt{b^2 - 4ac}}{2a}\]

\subsection{Matrix Example}
A sample matrix:
\[
\begin{pmatrix}
a_{11} & a_{12} & a_{13} \\
a_{21} & a_{22} & a_{23} \\
a_{31} & a_{32} & a_{33}
\end{pmatrix}
\]

\section{Tables}
Below is a sample table showing data:

\begin{table}[h]
\centering
\begin{tabular}{@{}lrrr@{}}
\toprule
\textbf{Category} & \textbf{Value 1} & \textbf{Value 2} & \textbf{Total} \\
\midrule
Item A & 10.5 & 20.3 & 30.8 \\
Item B & 15.2 & 25.1 & 40.3 \\
Item C & 22.1 & 12.7 & 34.8 \\
\bottomrule
\end{tabular}
\caption{Sample Data Table}
\label{tab:sample}
\end{table}

\section{Lists}
\subsection{Bullet Points}
\begin{itemize}
  \item First item
  \item Second item
  \item Third item with sub-items:
    \begin{itemize}
      \item Sub-item 1
      \item Sub-item 2
    \end{itemize}
\end{itemize}

\subsection{Numbered List}
\begin{enumerate}
  \item First step
  \item Second step
  \item Third step
\end{enumerate}

\section{Citations and References}
When you use the citation search feature, references will be automatically managed. For example, you might cite a paper like this: \cite{example2024}. All citations will be automatically added to your \texttt{cite.bib} file.

\section{Conclusion}
This template demonstrates basic \LaTeX{} features including mathematical formulas, tables, and lists. The bibliography system is now set up to work seamlessly with Bibby AI's citation management. Feel free to modify it for your needs!

\bibliographystyle{plain}
\bibliography{cite}

\end{document}\documentclass{article}
\usepackage[utf8]{inputenc}
\usepackage{booktabs}
\usepackage{amsmath}
\usepackage{natbib}

\title{My Document}
\author{Your Name}
\date{\today}

\begin{document}

\maketitle

\section{Introduction}
This is a sample document demonstrating various \LaTeX{} features including mathematical equations, tables, and formatting. You can cite papers using the citation search feature, and they will be automatically added to your \texttt{cite.bib} file.

\section{Mathematical Expressions}
Here are some example equations:

\subsection{Quadratic Formula}
The quadratic formula is given by:
\[x = \frac{-b \pm \sqrt{b^2 - 4ac}}{2a}\]

\subsection{Matrix Example}
A sample matrix:
\[
\begin{pmatrix}
a_{11} & a_{12} & a_{13} \\
a_{21} & a_{22} & a_{23} \\
a_{31} & a_{32} & a_{33}
\end{pmatrix}
\]

\section{Tables}
Below is a sample table showing data:

\begin{table}[h]
\centering
\begin{tabular}{@{}lrrr@{}}
\toprule
\textbf{Category} & \textbf{Value 1} & \textbf{Value 2} & \textbf{Total} \\
\midrule
Item A & 10.5 & 20.3 & 30.8 \\
Item B & 15.2 & 25.1 & 40.3 \\
Item C & 22.1 & 12.7 & 34.8 \\
\bottomrule
\end{tabular}
\caption{Sample Data Table}
\label{tab:sample}
\end{table}

\section{Lists}
\subsection{Bullet Points}
\begin{itemize}
  \item First item
  \item Second item
  \item Third item with sub-items:
    \begin{itemize}
      \item Sub-item 1
      \item Sub-item 2
    \end{itemize}
\end{itemize}

\subsection{Numbered List}
\begin{enumerate}
  \item First step
  \item Second step
  \item Third step
\end{enumerate}

\section{Citations and References}
When you use the citation search feature, references will be automatically managed. For example, you might cite a paper like this: \cite{example2024}. All citations will be automatically added to your \texttt{cite.bib} file.

\section{Conclusion}
This template demonstrates basic \LaTeX{} features including mathematical formulas, tables, and lists. The bibliography system is now set up to work seamlessly with Bibby AI's citation management. Feel free to modify it for your needs!

\bibliographystyle{plain}
\bibliography{cite}

\end{document}\documentclass{article}
\usepackage[utf8]{inputenc}
\usepackage{booktabs}
\usepackage{amsmath}
\usepackage{natbib}

\title{My Document}
\author{Your Name}
\date{\today}

\begin{document}

\maketitle

\section{Introduction}
This is a sample document demonstrating various \LaTeX{} features including mathematical equations, tables, and formatting. You can cite papers using the citation search feature, and they will be automatically added to your \texttt{cite.bib} file.

\section{Mathematical Expressions}
Here are some example equations:

\subsection{Quadratic Formula}
The quadratic formula is given by:
\[x = \frac{-b \pm \sqrt{b^2 - 4ac}}{2a}\]

\subsection{Matrix Example}
A sample matrix:
\[
\begin{pmatrix}
a_{11} & a_{12} & a_{13} \\
a_{21} & a_{22} & a_{23} \\
a_{31} & a_{32} & a_{33}
\end{pmatrix}
\]

\section{Tables}
Below is a sample table showing data:

\begin{table}[h]
\centering
\begin{tabular}{@{}lrrr@{}}
\toprule
\textbf{Category} & \textbf{Value 1} & \textbf{Value 2} & \textbf{Total} \\
\midrule
Item A & 10.5 & 20.3 & 30.8 \\
Item B & 15.2 & 25.1 & 40.3 \\
Item C & 22.1 & 12.7 & 34.8 \\
\bottomrule
\end{tabular}
\caption{Sample Data Table}
\label{tab:sample}
\end{table}

\section{Lists}
\subsection{Bullet Points}
\begin{itemize}
  \item First item
  \item Second item
  \item Third item with sub-items:
    \begin{itemize}
      \item Sub-item 1
      \item Sub-item 2
    \end{itemize}
\end{itemize}

\subsection{Numbered List}
\begin{enumerate}
  \item First step
  \item Second step
  \item Third step
\end{enumerate}

\section{Citations and References}
When you use the citation search feature, references will be automatically managed. For example, you might cite a paper like this: \cite{example2024}. All citations will be automatically added to your \texttt{cite.bib} file.

\section{Conclusion}
This template demonstrates basic \LaTeX{} features including mathematical formulas, tables, and lists. The bibliography system is now set up to work seamlessly with Bibby AI's citation management. Feel free to modify it for your needs!

\bibliographystyle{plain}
\bibliography{cite}

\end{document}\documentclass{article}
\usepackage[utf8]{inputenc}
\usepackage{booktabs}
\usepackage{amsmath}
\usepackage{natbib}

\title{My Document}
\author{Your Name}
\date{\today}

\begin{document}

\maketitle

\section{Introduction}
This is a sample document demonstrating various \LaTeX{} features including mathematical equations, tables, and formatting. You can cite papers using the citation search feature, and they will be automatically added to your \texttt{cite.bib} file.

\section{Mathematical Expressions}
Here are some example equations:

\subsection{Quadratic Formula}
The quadratic formula is given by:
\[x = \frac{-b \pm \sqrt{b^2 - 4ac}}{2a}\]

\subsection{Matrix Example}
A sample matrix:
\[
\begin{pmatrix}
a_{11} & a_{12} & a_{13} \\
a_{21} & a_{22} & a_{23} \\
a_{31} & a_{32} & a_{33}
\end{pmatrix}
\]

\section{Tables}
Below is a sample table showing data:

\begin{table}[h]
\centering
\begin{tabular}{@{}lrrr@{}}
\toprule
\textbf{Category} & \textbf{Value 1} & \textbf{Value 2} & \textbf{Total} \\
\midrule
Item A & 10.5 & 20.3 & 30.8 \\
Item B & 15.2 & 25.1 & 40.3 \\
Item C & 22.1 & 12.7 & 34.8 \\
\bottomrule
\end{tabular}
\caption{Sample Data Table}
\label{tab:sample}
\end{table}

\section{Lists}
\subsection{Bullet Points}
\begin{itemize}
  \item First item
  \item Second item
  \item Third item with sub-items:
    \begin{itemize}
      \item Sub-item 1
      \item Sub-item 2
    \end{itemize}
\end{itemize}

\subsection{Numbered List}
\begin{enumerate}
  \item First step
  \item Second step
  \item Third step
\end{enumerate}

\section{Citations and References}
When you use the citation search feature, references will be automatically managed. For example, you might cite a paper like this: \cite{example2024}. All citations will be automatically added to your \texttt{cite.bib} file.

\section{Conclusion}
This template demonstrates basic \LaTeX{} features including mathematical formulas, tables, and lists. The bibliography system is now set up to work seamlessly with Bibby AI's citation management. Feel free to modify it for your needs!

\bibliographystyle{plain}
\bibliography{cite}

\end{document}\documentclass{article}
\usepackage[utf8]{inputenc}
\usepackage{booktabs}
\usepackage{amsmath}
\usepackage{natbib}

\title{My Document}
\author{Your Name}
\date{\today}

\begin{document}

\maketitle

\section{Introduction}
This is a sample document demonstrating various \LaTeX{} features including mathematical equations, tables, and formatting. You can cite papers using the citation search feature, and they will be automatically added to your \texttt{cite.bib} file.

\section{Mathematical Expressions}
Here are some example equations:

\subsection{Quadratic Formula}
The quadratic formula is given by:
\[x = \frac{-b \pm \sqrt{b^2 - 4ac}}{2a}\]

\subsection{Matrix Example}
A sample matrix:
\[
\begin{pmatrix}
a_{11} & a_{12} & a_{13} \\
a_{21} & a_{22} & a_{23} \\
a_{31} & a_{32} & a_{33}
\end{pmatrix}
\]

\section{Tables}
Below is a sample table showing data:

\begin{table}[h]
\centering
\begin{tabular}{@{}lrrr@{}}
\toprule
\textbf{Category} & \textbf{Value 1} & \textbf{Value 2} & \textbf{Total} \\
\midrule
Item A & 10.5 & 20.3 & 30.8 \\
Item B & 15.2 & 25.1 & 40.3 \\
Item C & 22.1 & 12.7 & 34.8 \\
\bottomrule
\end{tabular}
\caption{Sample Data Table}
\label{tab:sample}
\end{table}

\section{Lists}
\subsection{Bullet Points}
\begin{itemize}
  \item First item
  \item Second item
  \item Third item with sub-items:
    \begin{itemize}
      \item Sub-item 1
      \item Sub-item 2
    \end{itemize}
\end{itemize}

\subsection{Numbered List}
\begin{enumerate}
  \item First step
  \item Second step
  \item Third step
\end{enumerate}

\section{Citations and References}
When you use the citation search feature, references will be automatically managed. For example, you might cite a paper like this: \cite{example2024}. All citations will be automatically added to your \texttt{cite.bib} file.

\section{Conclusion}
This template demonstrates basic \LaTeX{} features including mathematical formulas, tables, and lists. The bibliography system is now set up to work seamlessly with Bibby AI's citation management. Feel free to modify it for your needs!

\bibliographystyle{plain}
\bibliography{cite}

\end{document}\documentclass{article}
\usepackage[utf8]{inputenc}
\usepackage{booktabs}
\usepackage{amsmath}
\usepackage{natbib}

\title{My Document}
\author{Your Name}
\date{\today}

\begin{document}

\maketitle

\section{Introduction}
This is a sample document demonstrating various \LaTeX{} features including mathematical equations, tables, and formatting. You can cite papers using the citation search feature, and they will be automatically added to your \texttt{cite.bib} file.

\section{Mathematical Expressions}
Here are some example equations:

\subsection{Quadratic Formula}
The quadratic formula is given by:
\[x = \frac{-b \pm \sqrt{b^2 - 4ac}}{2a}\]

\subsection{Matrix Example}
A sample matrix:
\[
\begin{pmatrix}
a_{11} & a_{12} & a_{13} \\
a_{21} & a_{22} & a_{23} \\
a_{31} & a_{32} & a_{33}
\end{pmatrix}
\]

\section{Tables}
Below is a sample table showing data:

\begin{table}[h]
\centering
\begin{tabular}{@{}lrrr@{}}
\toprule
\textbf{Category} & \textbf{Value 1} & \textbf{Value 2} & \textbf{Total} \\
\midrule
Item A & 10.5 & 20.3 & 30.8 \\
Item B & 15.2 & 25.1 & 40.3 \\
Item C & 22.1 & 12.7 & 34.8 \\
\bottomrule
\end{tabular}
\caption{Sample Data Table}
\label{tab:sample}
\end{table}

\section{Lists}
\subsection{Bullet Points}
\begin{itemize}
  \item First item
  \item Second item
  \item Third item with sub-items:
    \begin{itemize}
      \item Sub-item 1
      \item Sub-item 2
    \end{itemize}
\end{itemize}

\subsection{Numbered List}
\begin{enumerate}
  \item First step
  \item Second step
  \item Third step
\end{enumerate}

\section{Citations and References}
When you use the citation search feature, references will be automatically managed. For example, you might cite a paper like this: \cite{example2024}. All citations will be automatically added to your \texttt{cite.bib} file.

\section{Conclusion}
This template demonstrates basic \LaTeX{} features including mathematical formulas, tables, and lists. The bibliography system is now set up to work seamlessly with Bibby AI's citation management. Feel free to modify it for your needs!

\bibliographystyle{plain}
\bibliography{cite}

\end{document}\documentclass{article}
\usepackage[utf8]{inputenc}
\usepackage{booktabs}
\usepackage{amsmath}
\usepackage{natbib}

\title{My Document}
\author{Your Name}
\date{\today}

\begin{document}

\maketitle

\section{Introduction}
This is a sample document demonstrating various \LaTeX{} features including mathematical equations, tables, and formatting. You can cite papers using the citation search feature, and they will be automatically added to your \texttt{cite.bib} file.

\section{Mathematical Expressions}
Here are some example equations:

\subsection{Quadratic Formula}
The quadratic formula is given by:
\[x = \frac{-b \pm \sqrt{b^2 - 4ac}}{2a}\]

\subsection{Matrix Example}
A sample matrix:
\[
\begin{pmatrix}
a_{11} & a_{12} & a_{13} \\
a_{21} & a_{22} & a_{23} \\
a_{31} & a_{32} & a_{33}
\end{pmatrix}
\]

\section{Tables}
Below is a sample table showing data:

\begin{table}[h]
\centering
\begin{tabular}{@{}lrrr@{}}
\toprule
\textbf{Category} & \textbf{Value 1} & \textbf{Value 2} & \textbf{Total} \\
\midrule
Item A & 10.5 & 20.3 & 30.8 \\
Item B & 15.2 & 25.1 & 40.3 \\
Item C & 22.1 & 12.7 & 34.8 \\
\bottomrule
\end{tabular}
\caption{Sample Data Table}
\label{tab:sample}
\end{table}

\section{Lists}
\subsection{Bullet Points}
\begin{itemize}
  \item First item
  \item Second item
  \item Third item with sub-items:
    \begin{itemize}
      \item Sub-item 1
      \item Sub-item 2
    \end{itemize}
\end{itemize}

\subsection{Numbered List}
\begin{enumerate}
  \item First step
  \item Second step
  \item Third step
\end{enumerate}

\section{Citations and References}
When you use the citation search feature, references will be automatically managed. For example, you might cite a paper like this: \cite{example2024}. All citations will be automatically added to your \texttt{cite.bib} file.

\section{Conclusion}
This template demonstrates basic \LaTeX{} features including mathematical formulas, tables, and lists. The bibliography system is now set up to work seamlessly with Bibby AI's citation management. Feel free to modify it for your needs!

\bibliographystyle{plain}
\bibliography{cite}

\end{document}\documentclass{article}
\usepackage[utf8]{inputenc}
\usepackage{booktabs}
\usepackage{amsmath}
\usepackage{natbib}

\title{My Document}
\author{Your Name}
\date{\today}

\begin{document}

\maketitle

\section{Introduction}
This is a sample document demonstrating various \LaTeX{} features including mathematical equations, tables, and formatting. You can cite papers using the citation search feature, and they will be automatically added to your \texttt{cite.bib} file.

\section{Mathematical Expressions}
Here are some example equations:

\subsection{Quadratic Formula}
The quadratic formula is given by:
\[x = \frac{-b \pm \sqrt{b^2 - 4ac}}{2a}\]

\subsection{Matrix Example}
A sample matrix:
\[
\begin{pmatrix}
a_{11} & a_{12} & a_{13} \\
a_{21} & a_{22} & a_{23} \\
a_{31} & a_{32} & a_{33}
\end{pmatrix}
\]

\section{Tables}
Below is a sample table showing data:

\begin{table}[h]
\centering
\begin{tabular}{@{}lrrr@{}}
\toprule
\textbf{Category} & \textbf{Value 1} & \textbf{Value 2} & \textbf{Total} \\
\midrule
Item A & 10.5 & 20.3 & 30.8 \\
Item B & 15.2 & 25.1 & 40.3 \\
Item C & 22.1 & 12.7 & 34.8 \\
\bottomrule
\end{tabular}
\caption{Sample Data Table}
\label{tab:sample}
\end{table}

\section{Lists}
\subsection{Bullet Points}
\begin{itemize}
  \item First item
  \item Second item
  \item Third item with sub-items:
    \begin{itemize}
      \item Sub-item 1
      \item Sub-item 2
    \end{itemize}
\end{itemize}

\subsection{Numbered List}
\begin{enumerate}
  \item First step
  \item Second step
  \item Third step
\end{enumerate}

\section{Citations and References}
When you use the citation search feature, references will be automatically managed. For example, you might cite a paper like this: \cite{example2024}. All citations will be automatically added to your \texttt{cite.bib} file.

\section{Conclusion}
This template demonstrates basic \LaTeX{} features including mathematical formulas, tables, and lists. The bibliography system is now set up to work seamlessly with Bibby AI's citation management. Feel free to modify it for your needs!

\bibliographystyle{plain}
\bibliography{cite}

\end{document}\documentclass{article}
\usepackage[utf8]{inputenc}
\usepackage{booktabs}
\usepackage{amsmath}
\usepackage{natbib}

\title{My Document}
\author{Your Name}
\date{\today}

\begin{document}

\maketitle

\section{Introduction}
This is a sample document demonstrating various \LaTeX{} features including mathematical equations, tables, and formatting. You can cite papers using the citation search feature, and they will be automatically added to your \texttt{cite.bib} file.

\section{Mathematical Expressions}
Here are some example equations:

\subsection{Quadratic Formula}
The quadratic formula is given by:
\[x = \frac{-b \pm \sqrt{b^2 - 4ac}}{2a}\]

\subsection{Matrix Example}
A sample matrix:
\[
\begin{pmatrix}
a_{11} & a_{12} & a_{13} \\
a_{21} & a_{22} & a_{23} \\
a_{31} & a_{32} & a_{33}
\end{pmatrix}
\]

\section{Tables}
Below is a sample table showing data:

\begin{table}[h]
\centering
\begin{tabular}{@{}lrrr@{}}
\toprule
\textbf{Category} & \textbf{Value 1} & \textbf{Value 2} & \textbf{Total} \\
\midrule
Item A & 10.5 & 20.3 & 30.8 \\
Item B & 15.2 & 25.1 & 40.3 \\
Item C & 22.1 & 12.7 & 34.8 \\
\bottomrule
\end{tabular}
\caption{Sample Data Table}
\label{tab:sample}
\end{table}

\section{Lists}
\subsection{Bullet Points}
\begin{itemize}
  \item First item
  \item Second item
  \item Third item with sub-items:
    \begin{itemize}
      \item Sub-item 1
      \item Sub-item 2
    \end{itemize}
\end{itemize}

\subsection{Numbered List}
\begin{enumerate}
  \item First step
  \item Second step
  \item Third step
\end{enumerate}

\section{Citations and References}
When you use the citation search feature, references will be automatically managed. For example, you might cite a paper like this: \cite{example2024}. All citations will be automatically added to your \texttt{cite.bib} file.

\section{Conclusion}
This template demonstrates basic \LaTeX{} features including mathematical formulas, tables, and lists. The bibliography system is now set up to work seamlessly with Bibby AI's citation management. Feel free to modify it for your needs!

\bibliographystyle{plain}
\bibliography{cite}

\end{document}\documentclass{article}
\usepackage[utf8]{inputenc}
\usepackage{booktabs}
\usepackage{amsmath}
\usepackage{natbib}

\title{My Document}
\author{Your Name}
\date{\today}

\begin{document}

\maketitle

\section{Introduction}
This is a sample document demonstrating various \LaTeX{} features including mathematical equations, tables, and formatting. You can cite papers using the citation search feature, and they will be automatically added to your \texttt{cite.bib} file.

\section{Mathematical Expressions}
Here are some example equations:

\subsection{Quadratic Formula}
The quadratic formula is given by:
\[x = \frac{-b \pm \sqrt{b^2 - 4ac}}{2a}\]

\subsection{Matrix Example}
A sample matrix:
\[
\begin{pmatrix}
a_{11} & a_{12} & a_{13} \\
a_{21} & a_{22} & a_{23} \\
a_{31} & a_{32} & a_{33}
\end{pmatrix}
\]

\section{Tables}
Below is a sample table showing data:

\begin{table}[h]
\centering
\begin{tabular}{@{}lrrr@{}}
\toprule
\textbf{Category} & \textbf{Value 1} & \textbf{Value 2} & \textbf{Total} \\
\midrule
Item A & 10.5 & 20.3 & 30.8 \\
Item B & 15.2 & 25.1 & 40.3 \\
Item C & 22.1 & 12.7 & 34.8 \\
\bottomrule
\end{tabular}
\caption{Sample Data Table}
\label{tab:sample}
\end{table}

\section{Lists}
\subsection{Bullet Points}
\begin{itemize}
  \item First item
  \item Second item
  \item Third item with sub-items:
    \begin{itemize}
      \item Sub-item 1
      \item Sub-item 2
    \end{itemize}
\end{itemize}

\subsection{Numbered List}
\begin{enumerate}
  \item First step
  \item Second step
  \item Third step
\end{enumerate}

\section{Citations and References}
When you use the citation search feature, references will be automatically managed. For example, you might cite a paper like this: \cite{example2024}. All citations will be automatically added to your \texttt{cite.bib} file.

\section{Conclusion}
This template demonstrates basic \LaTeX{} features including mathematical formulas, tables, and lists. The bibliography system is now set up to work seamlessly with Bibby AI's citation management. Feel free to modify it for your needs!

\bibliographystyle{plain}
\bibliography{cite}

\end{document}\documentclass{article}
\usepackage[utf8]{inputenc}
\usepackage{booktabs}
\usepackage{amsmath}
\usepackage{natbib}

\title{My Document}
\author{Your Name}
\date{\today}

\begin{document}

\maketitle

\section{Introduction}
This is a sample document demonstrating various \LaTeX{} features including mathematical equations, tables, and formatting. You can cite papers using the citation search feature, and they will be automatically added to your \texttt{cite.bib} file.

\section{Mathematical Expressions}
Here are some example equations:

\subsection{Quadratic Formula}
The quadratic formula is given by:
\[x = \frac{-b \pm \sqrt{b^2 - 4ac}}{2a}\]

\subsection{Matrix Example}
A sample matrix:
\[
\begin{pmatrix}
a_{11} & a_{12} & a_{13} \\
a_{21} & a_{22} & a_{23} \\
a_{31} & a_{32} & a_{33}
\end{pmatrix}
\]

\section{Tables}
Below is a sample table showing data:

\begin{table}[h]
\centering
\begin{tabular}{@{}lrrr@{}}
\toprule
\textbf{Category} & \textbf{Value 1} & \textbf{Value 2} & \textbf{Total} \\
\midrule
Item A & 10.5 & 20.3 & 30.8 \\
Item B & 15.2 & 25.1 & 40.3 \\
Item C & 22.1 & 12.7 & 34.8 \\
\bottomrule
\end{tabular}
\caption{Sample Data Table}
\label{tab:sample}
\end{table}

\section{Lists}
\subsection{Bullet Points}
\begin{itemize}
  \item First item
  \item Second item
  \item Third item with sub-items:
    \begin{itemize}
      \item Sub-item 1
      \item Sub-item 2
    \end{itemize}
\end{itemize}

\subsection{Numbered List}
\begin{enumerate}
  \item First step
  \item Second step
  \item Third step
\end{enumerate}

\section{Citations and References}
When you use the citation search feature, references will be automatically managed. For example, you might cite a paper like this: \cite{example2024}. All citations will be automatically added to your \texttt{cite.bib} file.

\section{Conclusion}
This template demonstrates basic \LaTeX{} features including mathematical formulas, tables, and lists. The bibliography system is now set up to work seamlessly with Bibby AI's citation management. Feel free to modify it for your needs!

\bibliographystyle{plain}
\bibliography{cite}

\end{document}\documentclass{article}
\usepackage[utf8]{inputenc}
\usepackage{booktabs}
\usepackage{amsmath}
\usepackage{natbib}

\title{My Document}
\author{Your Name}
\date{\today}

\begin{document}

\maketitle

\section{Introduction}
This is a sample document demonstrating various \LaTeX{} features including mathematical equations, tables, and formatting. You can cite papers using the citation search feature, and they will be automatically added to your \texttt{cite.bib} file.

\section{Mathematical Expressions}
Here are some example equations:

\subsection{Quadratic Formula}
The quadratic formula is given by:
\[x = \frac{-b \pm \sqrt{b^2 - 4ac}}{2a}\]

\subsection{Matrix Example}
A sample matrix:
\[
\begin{pmatrix}
a_{11} & a_{12} & a_{13} \\
a_{21} & a_{22} & a_{23} \\
a_{31} & a_{32} & a_{33}
\end{pmatrix}
\]

\section{Tables}
Below is a sample table showing data:

\begin{table}[h]
\centering
\begin{tabular}{@{}lrrr@{}}
\toprule
\textbf{Category} & \textbf{Value 1} & \textbf{Value 2} & \textbf{Total} \\
\midrule
Item A & 10.5 & 20.3 & 30.8 \\
Item B & 15.2 & 25.1 & 40.3 \\
Item C & 22.1 & 12.7 & 34.8 \\
\bottomrule
\end{tabular}
\caption{Sample Data Table}
\label{tab:sample}
\end{table}

\section{Lists}
\subsection{Bullet Points}
\begin{itemize}
  \item First item
  \item Second item
  \item Third item with sub-items:
    \begin{itemize}
      \item Sub-item 1
      \item Sub-item 2
    \end{itemize}
\end{itemize}

\subsection{Numbered List}
\begin{enumerate}
  \item First step
  \item Second step
  \item Third step
\end{enumerate}

\section{Citations and References}
When you use the citation search feature, references will be automatically managed. For example, you might cite a paper like this: \cite{example2024}. All citations will be automatically added to your \texttt{cite.bib} file.

\section{Conclusion}
This template demonstrates basic \LaTeX{} features including mathematical formulas, tables, and lists. The bibliography system is now set up to work seamlessly with Bibby AI's citation management. Feel free to modify it for your needs!

\bibliographystyle{plain}
\bibliography{cite}

\end{document}\documentclass{article}
\usepackage[utf8]{inputenc}
\usepackage{booktabs}
\usepackage{amsmath}
\usepackage{natbib}

\title{My Document}
\author{Your Name}
\date{\today}

\begin{document}

\maketitle

\section{Introduction}
This is a sample document demonstrating various \LaTeX{} features including mathematical equations, tables, and formatting. You can cite papers using the citation search feature, and they will be automatically added to your \texttt{cite.bib} file.

\section{Mathematical Expressions}
Here are some example equations:

\subsection{Quadratic Formula}
The quadratic formula is given by:
\[x = \frac{-b \pm \sqrt{b^2 - 4ac}}{2a}\]

\subsection{Matrix Example}
A sample matrix:
\[
\begin{pmatrix}
a_{11} & a_{12} & a_{13} \\
a_{21} & a_{22} & a_{23} \\
a_{31} & a_{32} & a_{33}
\end{pmatrix}
\]

\section{Tables}
Below is a sample table showing data:

\begin{table}[h]
\centering
\begin{tabular}{@{}lrrr@{}}
\toprule
\textbf{Category} & \textbf{Value 1} & \textbf{Value 2} & \textbf{Total} \\
\midrule
Item A & 10.5 & 20.3 & 30.8 \\
Item B & 15.2 & 25.1 & 40.3 \\
Item C & 22.1 & 12.7 & 34.8 \\
\bottomrule
\end{tabular}
\caption{Sample Data Table}
\label{tab:sample}
\end{table}

\section{Lists}
\subsection{Bullet Points}
\begin{itemize}
  \item First item
  \item Second item
  \item Third item with sub-items:
    \begin{itemize}
      \item Sub-item 1
      \item Sub-item 2
    \end{itemize}
\end{itemize}

\subsection{Numbered List}
\begin{enumerate}
  \item First step
  \item Second step
  \item Third step
\end{enumerate}

\section{Citations and References}
When you use the citation search feature, references will be automatically managed. For example, you might cite a paper like this: \cite{example2024}. All citations will be automatically added to your \texttt{cite.bib} file.

\section{Conclusion}
This template demonstrates basic \LaTeX{} features including mathematical formulas, tables, and lists. The bibliography system is now set up to work seamlessly with Bibby AI's citation management. Feel free to modify it for your needs!

\bibliographystyle{plain}
\bibliography{cite}

\end{document}\documentclass{article}
\usepackage[utf8]{inputenc}
\usepackage{booktabs}
\usepackage{amsmath}
\usepackage{natbib}

\title{My Document}
\author{Your Name}
\date{\today}

\begin{document}

\maketitle

\section{Introduction}
This is a sample document demonstrating various \LaTeX{} features including mathematical equations, tables, and formatting. You can cite papers using the citation search feature, and they will be automatically added to your \texttt{cite.bib} file.

\section{Mathematical Expressions}
Here are some example equations:

\subsection{Quadratic Formula}
The quadratic formula is given by:
\[x = \frac{-b \pm \sqrt{b^2 - 4ac}}{2a}\]

\subsection{Matrix Example}
A sample matrix:
\[
\begin{pmatrix}
a_{11} & a_{12} & a_{13} \\
a_{21} & a_{22} & a_{23} \\
a_{31} & a_{32} & a_{33}
\end{pmatrix}
\]

\section{Tables}
Below is a sample table showing data:

\begin{table}[h]
\centering
\begin{tabular}{@{}lrrr@{}}
\toprule
\textbf{Category} & \textbf{Value 1} & \textbf{Value 2} & \textbf{Total} \\
\midrule
Item A & 10.5 & 20.3 & 30.8 \\
Item B & 15.2 & 25.1 & 40.3 \\
Item C & 22.1 & 12.7 & 34.8 \\
\bottomrule
\end{tabular}
\caption{Sample Data Table}
\label{tab:sample}
\end{table}

\section{Lists}
\subsection{Bullet Points}
\begin{itemize}
  \item First item
  \item Second item
  \item Third item with sub-items:
    \begin{itemize}
      \item Sub-item 1
      \item Sub-item 2
    \end{itemize}
\end{itemize}

\subsection{Numbered List}
\begin{enumerate}
  \item First step
  \item Second step
  \item Third step
\end{enumerate}

\section{Citations and References}
When you use the citation search feature, references will be automatically managed. For example, you might cite a paper like this: \cite{example2024}. All citations will be automatically added to your \texttt{cite.bib} file.

\section{Conclusion}
This template demonstrates basic \LaTeX{} features including mathematical formulas, tables, and lists. The bibliography system is now set up to work seamlessly with Bibby AI's citation management. Feel free to modify it for your needs!

\bibliographystyle{plain}
\bibliography{cite}

\end{document}\documentclass{article}
\usepackage[utf8]{inputenc}
\usepackage{booktabs}
\usepackage{amsmath}
\usepackage{natbib}

\title{My Document}
\author{Your Name}
\date{\today}

\begin{document}

\maketitle

\section{Introduction}
This is a sample document demonstrating various \LaTeX{} features including mathematical equations, tables, and formatting. You can cite papers using the citation search feature, and they will be automatically added to your \texttt{cite.bib} file.

\section{Mathematical Expressions}
Here are some example equations:

\subsection{Quadratic Formula}
The quadratic formula is given by:
\[x = \frac{-b \pm \sqrt{b^2 - 4ac}}{2a}\]

\subsection{Matrix Example}
A sample matrix:
\[
\begin{pmatrix}
a_{11} & a_{12} & a_{13} \\
a_{21} & a_{22} & a_{23} \\
a_{31} & a_{32} & a_{33}
\end{pmatrix}
\]

\section{Tables}
Below is a sample table showing data:

\begin{table}[h]
\centering
\begin{tabular}{@{}lrrr@{}}
\toprule
\textbf{Category} & \textbf{Value 1} & \textbf{Value 2} & \textbf{Total} \\
\midrule
Item A & 10.5 & 20.3 & 30.8 \\
Item B & 15.2 & 25.1 & 40.3 \\
Item C & 22.1 & 12.7 & 34.8 \\
\bottomrule
\end{tabular}
\caption{Sample Data Table}
\label{tab:sample}
\end{table}

\section{Lists}
\subsection{Bullet Points}
\begin{itemize}
  \item First item
  \item Second item
  \item Third item with sub-items:
    \begin{itemize}
      \item Sub-item 1
      \item Sub-item 2
    \end{itemize}
\end{itemize}

\subsection{Numbered List}
\begin{enumerate}
  \item First step
  \item Second step
  \item Third step
\end{enumerate}

\section{Citations and References}
When you use the citation search feature, references will be automatically managed. For example, you might cite a paper like this: \cite{example2024}. All citations will be automatically added to your \texttt{cite.bib} file.

\section{Conclusion}
This template demonstrates basic \LaTeX{} features including mathematical formulas, tables, and lists. The bibliography system is now set up to work seamlessly with Bibby AI's citation management. Feel free to modify it for your needs!

\bibliographystyle{plain}
\bibliography{cite}

\end{document}\documentclass{article}
\usepackage[utf8]{inputenc}
\usepackage{booktabs}
\usepackage{amsmath}
\usepackage{natbib}

\title{My Document}
\author{Your Name}
\date{\today}

\begin{document}

\maketitle

\section{Introduction}
This is a sample document demonstrating various \LaTeX{} features including mathematical equations, tables, and formatting. You can cite papers using the citation search feature, and they will be automatically added to your \texttt{cite.bib} file.

\section{Mathematical Expressions}
Here are some example equations:

\subsection{Quadratic Formula}
The quadratic formula is given by:
\[x = \frac{-b \pm \sqrt{b^2 - 4ac}}{2a}\]

\subsection{Matrix Example}
A sample matrix:
\[
\begin{pmatrix}
a_{11} & a_{12} & a_{13} \\
a_{21} & a_{22} & a_{23} \\
a_{31} & a_{32} & a_{33}
\end{pmatrix}
\]

\section{Tables}
Below is a sample table showing data:

\begin{table}[h]
\centering
\begin{tabular}{@{}lrrr@{}}
\toprule
\textbf{Category} & \textbf{Value 1} & \textbf{Value 2} & \textbf{Total} \\
\midrule
Item A & 10.5 & 20.3 & 30.8 \\
Item B & 15.2 & 25.1 & 40.3 \\
Item C & 22.1 & 12.7 & 34.8 \\
\bottomrule
\end{tabular}
\caption{Sample Data Table}
\label{tab:sample}
\end{table}

\section{Lists}
\subsection{Bullet Points}
\begin{itemize}
  \item First item
  \item Second item
  \item Third item with sub-items:
    \begin{itemize}
      \item Sub-item 1
      \item Sub-item 2
    \end{itemize}
\end{itemize}

\subsection{Numbered List}
\begin{enumerate}
  \item First step
  \item Second step
  \item Third step
\end{enumerate}

\section{Citations and References}
When you use the citation search feature, references will be automatically managed. For example, you might cite a paper like this: \cite{example2024}. All citations will be automatically added to your \texttt{cite.bib} file.

\section{Conclusion}
This template demonstrates basic \LaTeX{} features including mathematical formulas, tables, and lists. The bibliography system is now set up to work seamlessly with Bibby AI's citation management. Feel free to modify it for your needs!

\bibliographystyle{plain}
\bibliography{cite}

\end{document}\documentclass{article}
\usepackage[utf8]{inputenc}
\usepackage{booktabs}
\usepackage{amsmath}
\usepackage{natbib}

\title{My Document}
\author{Your Name}
\date{\today}

\begin{document}

\maketitle

\section{Introduction}
This is a sample document demonstrating various \LaTeX{} features including mathematical equations, tables, and formatting. You can cite papers using the citation search feature, and they will be automatically added to your \texttt{cite.bib} file.

\section{Mathematical Expressions}
Here are some example equations:

\subsection{Quadratic Formula}
The quadratic formula is given by:
\[x = \frac{-b \pm \sqrt{b^2 - 4ac}}{2a}\]

\subsection{Matrix Example}
A sample matrix:
\[
\begin{pmatrix}
a_{11} & a_{12} & a_{13} \\
a_{21} & a_{22} & a_{23} \\
a_{31} & a_{32} & a_{33}
\end{pmatrix}
\]

\section{Tables}
Below is a sample table showing data:

\begin{table}[h]
\centering
\begin{tabular}{@{}lrrr@{}}
\toprule
\textbf{Category} & \textbf{Value 1} & \textbf{Value 2} & \textbf{Total} \\
\midrule
Item A & 10.5 & 20.3 & 30.8 \\
Item B & 15.2 & 25.1 & 40.3 \\
Item C & 22.1 & 12.7 & 34.8 \\
\bottomrule
\end{tabular}
\caption{Sample Data Table}
\label{tab:sample}
\end{table}

\section{Lists}
\subsection{Bullet Points}
\begin{itemize}
  \item First item
  \item Second item
  \item Third item with sub-items:
    \begin{itemize}
      \item Sub-item 1
      \item Sub-item 2
    \end{itemize}
\end{itemize}

\subsection{Numbered List}
\begin{enumerate}
  \item First step
  \item Second step
  \item Third step
\end{enumerate}

\section{Citations and References}
When you use the citation search feature, references will be automatically managed. For example, you might cite a paper like this: \cite{example2024}. All citations will be automatically added to your \texttt{cite.bib} file.

\section{Conclusion}
This template demonstrates basic \LaTeX{} features including mathematical formulas, tables, and lists. The bibliography system is now set up to work seamlessly with Bibby AI's citation management. Feel free to modify it for your needs!

\bibliographystyle{plain}
\bibliography{cite}

\end{document}\documentclass{article}
\usepackage[utf8]{inputenc}
\usepackage{booktabs}
\usepackage{amsmath}
\usepackage{natbib}

\title{My Document}
\author{Your Name}
\date{\today}

\begin{document}

\maketitle

\section{Introduction}
This is a sample document demonstrating various \LaTeX{} features including mathematical equations, tables, and formatting. You can cite papers using the citation search feature, and they will be automatically added to your \texttt{cite.bib} file.

\section{Mathematical Expressions}
Here are some example equations:

\subsection{Quadratic Formula}
The quadratic formula is given by:
\[x = \frac{-b \pm \sqrt{b^2 - 4ac}}{2a}\]

\subsection{Matrix Example}
A sample matrix:
\[
\begin{pmatrix}
a_{11} & a_{12} & a_{13} \\
a_{21} & a_{22} & a_{23} \\
a_{31} & a_{32} & a_{33}
\end{pmatrix}
\]

\section{Tables}
Below is a sample table showing data:

\begin{table}[h]
\centering
\begin{tabular}{@{}lrrr@{}}
\toprule
\textbf{Category} & \textbf{Value 1} & \textbf{Value 2} & \textbf{Total} \\
\midrule
Item A & 10.5 & 20.3 & 30.8 \\
Item B & 15.2 & 25.1 & 40.3 \\
Item C & 22.1 & 12.7 & 34.8 \\
\bottomrule
\end{tabular}
\caption{Sample Data Table}
\label{tab:sample}
\end{table}

\section{Lists}
\subsection{Bullet Points}
\begin{itemize}
  \item First item
  \item Second item
  \item Third item with sub-items:
    \begin{itemize}
      \item Sub-item 1
      \item Sub-item 2
    \end{itemize}
\end{itemize}

\subsection{Numbered List}
\begin{enumerate}
  \item First step
  \item Second step
  \item Third step
\end{enumerate}

\section{Citations and References}
When you use the citation search feature, references will be automatically managed. For example, you might cite a paper like this: \cite{example2024}. All citations will be automatically added to your \texttt{cite.bib} file.

\section{Conclusion}
This template demonstrates basic \LaTeX{} features including mathematical formulas, tables, and lists. The bibliography system is now set up to work seamlessly with Bibby AI's citation management. Feel free to modify it for your needs!

\bibliographystyle{plain}
\bibliography{cite}

\end{document}\documentclass{article}
\usepackage[utf8]{inputenc}
\usepackage{booktabs}
\usepackage{amsmath}
\usepackage{natbib}

\title{My Document}
\author{Your Name}
\date{\today}

\begin{document}

\maketitle

\section{Introduction}
This is a sample document demonstrating various \LaTeX{} features including mathematical equations, tables, and formatting. You can cite papers using the citation search feature, and they will be automatically added to your \texttt{cite.bib} file.

\section{Mathematical Expressions}
Here are some example equations:

\subsection{Quadratic Formula}
The quadratic formula is given by:
\[x = \frac{-b \pm \sqrt{b^2 - 4ac}}{2a}\]

\subsection{Matrix Example}
A sample matrix:
\[
\begin{pmatrix}
a_{11} & a_{12} & a_{13} \\
a_{21} & a_{22} & a_{23} \\
a_{31} & a_{32} & a_{33}
\end{pmatrix}
\]

\section{Tables}
Below is a sample table showing data:

\begin{table}[h]
\centering
\begin{tabular}{@{}lrrr@{}}
\toprule
\textbf{Category} & \textbf{Value 1} & \textbf{Value 2} & \textbf{Total} \\
\midrule
Item A & 10.5 & 20.3 & 30.8 \\
Item B & 15.2 & 25.1 & 40.3 \\
Item C & 22.1 & 12.7 & 34.8 \\
\bottomrule
\end{tabular}
\caption{Sample Data Table}
\label{tab:sample}
\end{table}

\section{Lists}
\subsection{Bullet Points}
\begin{itemize}
  \item First item
  \item Second item
  \item Third item with sub-items:
    \begin{itemize}
      \item Sub-item 1
      \item Sub-item 2
    \end{itemize}
\end{itemize}

\subsection{Numbered List}
\begin{enumerate}
  \item First step
  \item Second step
  \item Third step
\end{enumerate}

\section{Citations and References}
When you use the citation search feature, references will be automatically managed. For example, you might cite a paper like this: \cite{example2024}. All citations will be automatically added to your \texttt{cite.bib} file.

\section{Conclusion}
This template demonstrates basic \LaTeX{} features including mathematical formulas, tables, and lists. The bibliography system is now set up to work seamlessly with Bibby AI's citation management. Feel free to modify it for your needs!

\bibliographystyle{plain}
\bibliography{cite}

\end{document}\documentclass{article}
\usepackage[utf8]{inputenc}
\usepackage{booktabs}
\usepackage{amsmath}
\usepackage{natbib}

\title{My Document}
\author{Your Name}
\date{\today}

\begin{document}

\maketitle

\section{Introduction}
This is a sample document demonstrating various \LaTeX{} features including mathematical equations, tables, and formatting. You can cite papers using the citation search feature, and they will be automatically added to your \texttt{cite.bib} file.

\section{Mathematical Expressions}
Here are some example equations:

\subsection{Quadratic Formula}
The quadratic formula is given by:
\[x = \frac{-b \pm \sqrt{b^2 - 4ac}}{2a}\]

\subsection{Matrix Example}
A sample matrix:
\[
\begin{pmatrix}
a_{11} & a_{12} & a_{13} \\
a_{21} & a_{22} & a_{23} \\
a_{31} & a_{32} & a_{33}
\end{pmatrix}
\]

\section{Tables}
Below is a sample table showing data:

\begin{table}[h]
\centering
\begin{tabular}{@{}lrrr@{}}
\toprule
\textbf{Category} & \textbf{Value 1} & \textbf{Value 2} & \textbf{Total} \\
\midrule
Item A & 10.5 & 20.3 & 30.8 \\
Item B & 15.2 & 25.1 & 40.3 \\
Item C & 22.1 & 12.7 & 34.8 \\
\bottomrule
\end{tabular}
\caption{Sample Data Table}
\label{tab:sample}
\end{table}

\section{Lists}
\subsection{Bullet Points}
\begin{itemize}
  \item First item
  \item Second item
  \item Third item with sub-items:
    \begin{itemize}
      \item Sub-item 1
      \item Sub-item 2
    \end{itemize}
\end{itemize}

\subsection{Numbered List}
\begin{enumerate}
  \item First step
  \item Second step
  \item Third step
\end{enumerate}

\section{Citations and References}
When you use the citation search feature, references will be automatically managed. For example, you might cite a paper like this: \cite{example2024}. All citations will be automatically added to your \texttt{cite.bib} file.

\section{Conclusion}
This template demonstrates basic \LaTeX{} features including mathematical formulas, tables, and lists. The bibliography system is now set up to work seamlessly with Bibby AI's citation management. Feel free to modify it for your needs!

\bibliographystyle{plain}
\bibliography{cite}

\end{document}\documentclass{article}
\usepackage[utf8]{inputenc}
\usepackage{booktabs}
\usepackage{amsmath}
\usepackage{natbib}

\title{My Document}
\author{Your Name}
\date{\today}

\begin{document}

\maketitle

\section{Introduction}
This is a sample document demonstrating various \LaTeX{} features including mathematical equations, tables, and formatting. You can cite papers using the citation search feature, and they will be automatically added to your \texttt{cite.bib} file.

\section{Mathematical Expressions}
Here are some example equations:

\subsection{Quadratic Formula}
The quadratic formula is given by:
\[x = \frac{-b \pm \sqrt{b^2 - 4ac}}{2a}\]

\subsection{Matrix Example}
A sample matrix:
\[
\begin{pmatrix}
a_{11} & a_{12} & a_{13} \\
a_{21} & a_{22} & a_{23} \\
a_{31} & a_{32} & a_{33}
\end{pmatrix}
\]

\section{Tables}
Below is a sample table showing data:

\begin{table}[h]
\centering
\begin{tabular}{@{}lrrr@{}}
\toprule
\textbf{Category} & \textbf{Value 1} & \textbf{Value 2} & \textbf{Total} \\
\midrule
Item A & 10.5 & 20.3 & 30.8 \\
Item B & 15.2 & 25.1 & 40.3 \\
Item C & 22.1 & 12.7 & 34.8 \\
\bottomrule
\end{tabular}
\caption{Sample Data Table}
\label{tab:sample}
\end{table}

\section{Lists}
\subsection{Bullet Points}
\begin{itemize}
  \item First item
  \item Second item
  \item Third item with sub-items:
    \begin{itemize}
      \item Sub-item 1
      \item Sub-item 2
    \end{itemize}
\end{itemize}

\subsection{Numbered List}
\begin{enumerate}
  \item First step
  \item Second step
  \item Third step
\end{enumerate}

\section{Citations and References}
When you use the citation search feature, references will be automatically managed. For example, you might cite a paper like this: \cite{example2024}. All citations will be automatically added to your \texttt{cite.bib} file.

\section{Conclusion}
This template demonstrates basic \LaTeX{} features including mathematical formulas, tables, and lists. The bibliography system is now set up to work seamlessly with Bibby AI's citation management. Feel free to modify it for your needs!

\bibliographystyle{plain}
\bibliography{cite}

\end{document}\documentclass{article}
\usepackage[utf8]{inputenc}
\usepackage{booktabs}
\usepackage{amsmath}
\usepackage{natbib}

\title{My Document}
\author{Your Name}
\date{\today}

\begin{document}

\maketitle

\section{Introduction}
This is a sample document demonstrating various \LaTeX{} features including mathematical equations, tables, and formatting. You can cite papers using the citation search feature, and they will be automatically added to your \texttt{cite.bib} file.

\section{Mathematical Expressions}
Here are some example equations:

\subsection{Quadratic Formula}
The quadratic formula is given by:
\[x = \frac{-b \pm \sqrt{b^2 - 4ac}}{2a}\]

\subsection{Matrix Example}
A sample matrix:
\[
\begin{pmatrix}
a_{11} & a_{12} & a_{13} \\
a_{21} & a_{22} & a_{23} \\
a_{31} & a_{32} & a_{33}
\end{pmatrix}
\]

\section{Tables}
Below is a sample table showing data:

\begin{table}[h]
\centering
\begin{tabular}{@{}lrrr@{}}
\toprule
\textbf{Category} & \textbf{Value 1} & \textbf{Value 2} & \textbf{Total} \\
\midrule
Item A & 10.5 & 20.3 & 30.8 \\
Item B & 15.2 & 25.1 & 40.3 \\
Item C & 22.1 & 12.7 & 34.8 \\
\bottomrule
\end{tabular}
\caption{Sample Data Table}
\label{tab:sample}
\end{table}

\section{Lists}
\subsection{Bullet Points}
\begin{itemize}
  \item First item
  \item Second item
  \item Third item with sub-items:
    \begin{itemize}
      \item Sub-item 1
      \item Sub-item 2
    \end{itemize}
\end{itemize}

\subsection{Numbered List}
\begin{enumerate}
  \item First step
  \item Second step
  \item Third step
\end{enumerate}

\section{Citations and References}
When you use the citation search feature, references will be automatically managed. For example, you might cite a paper like this: \cite{example2024}. All citations will be automatically added to your \texttt{cite.bib} file.

\section{Conclusion}
This template demonstrates basic \LaTeX{} features including mathematical formulas, tables, and lists. The bibliography system is now set up to work seamlessly with Bibby AI's citation management. Feel free to modify it for your needs!

\bibliographystyle{plain}
\bibliography{cite}

\end{document}\documentclass{article}
\usepackage[utf8]{inputenc}
\usepackage{booktabs}
\usepackage{amsmath}
\usepackage{natbib}

\title{My Document}
\author{Your Name}
\date{\today}

\begin{document}

\maketitle

\section{Introduction}
This is a sample document demonstrating various \LaTeX{} features including mathematical equations, tables, and formatting. You can cite papers using the citation search feature, and they will be automatically added to your \texttt{cite.bib} file.

\section{Mathematical Expressions}
Here are some example equations:

\subsection{Quadratic Formula}
The quadratic formula is given by:
\[x = \frac{-b \pm \sqrt{b^2 - 4ac}}{2a}\]

\subsection{Matrix Example}
A sample matrix:
\[
\begin{pmatrix}
a_{11} & a_{12} & a_{13} \\
a_{21} & a_{22} & a_{23} \\
a_{31} & a_{32} & a_{33}
\end{pmatrix}
\]

\section{Tables}
Below is a sample table showing data:

\begin{table}[h]
\centering
\begin{tabular}{@{}lrrr@{}}
\toprule
\textbf{Category} & \textbf{Value 1} & \textbf{Value 2} & \textbf{Total} \\
\midrule
Item A & 10.5 & 20.3 & 30.8 \\
Item B & 15.2 & 25.1 & 40.3 \\
Item C & 22.1 & 12.7 & 34.8 \\
\bottomrule
\end{tabular}
\caption{Sample Data Table}
\label{tab:sample}
\end{table}

\section{Lists}
\subsection{Bullet Points}
\begin{itemize}
  \item First item
  \item Second item
  \item Third item with sub-items:
    \begin{itemize}
      \item Sub-item 1
      \item Sub-item 2
    \end{itemize}
\end{itemize}

\subsection{Numbered List}
\begin{enumerate}
  \item First step
  \item Second step
  \item Third step
\end{enumerate}

\section{Citations and References}
When you use the citation search feature, references will be automatically managed. For example, you might cite a paper like this: \cite{example2024}. All citations will be automatically added to your \texttt{cite.bib} file.

\section{Conclusion}
This template demonstrates basic \LaTeX{} features including mathematical formulas, tables, and lists. The bibliography system is now set up to work seamlessly with Bibby AI's citation management. Feel free to modify it for your needs!

\bibliographystyle{plain}
\bibliography{cite}

\end{document}\documentclass{article}
\usepackage[utf8]{inputenc}
\usepackage{booktabs}
\usepackage{amsmath}
\usepackage{natbib}

\title{My Document}
\author{Your Name}
\date{\today}

\begin{document}

\maketitle

\section{Introduction}
This is a sample document demonstrating various \LaTeX{} features including mathematical equations, tables, and formatting. You can cite papers using the citation search feature, and they will be automatically added to your \texttt{cite.bib} file.

\section{Mathematical Expressions}
Here are some example equations:

\subsection{Quadratic Formula}
The quadratic formula is given by:
\[x = \frac{-b \pm \sqrt{b^2 - 4ac}}{2a}\]

\subsection{Matrix Example}
A sample matrix:
\[
\begin{pmatrix}
a_{11} & a_{12} & a_{13} \\
a_{21} & a_{22} & a_{23} \\
a_{31} & a_{32} & a_{33}
\end{pmatrix}
\]

\section{Tables}
Below is a sample table showing data:

\begin{table}[h]
\centering
\begin{tabular}{@{}lrrr@{}}
\toprule
\textbf{Category} & \textbf{Value 1} & \textbf{Value 2} & \textbf{Total} \\
\midrule
Item A & 10.5 & 20.3 & 30.8 \\
Item B & 15.2 & 25.1 & 40.3 \\
Item C & 22.1 & 12.7 & 34.8 \\
\bottomrule
\end{tabular}
\caption{Sample Data Table}
\label{tab:sample}
\end{table}

\section{Lists}
\subsection{Bullet Points}
\begin{itemize}
  \item First item
  \item Second item
  \item Third item with sub-items:
    \begin{itemize}
      \item Sub-item 1
      \item Sub-item 2
    \end{itemize}
\end{itemize}

\subsection{Numbered List}
\begin{enumerate}
  \item First step
  \item Second step
  \item Third step
\end{enumerate}

\section{Citations and References}
When you use the citation search feature, references will be automatically managed. For example, you might cite a paper like this: \cite{example2024}. All citations will be automatically added to your \texttt{cite.bib} file.

\section{Conclusion}
This template demonstrates basic \LaTeX{} features including mathematical formulas, tables, and lists. The bibliography system is now set up to work seamlessly with Bibby AI's citation management. Feel free to modify it for your needs!

\bibliographystyle{plain}
\bibliography{cite}

\end{document}\documentclass{article}
\usepackage[utf8]{inputenc}
\usepackage{booktabs}
\usepackage{amsmath}
\usepackage{natbib}

\title{My Document}
\author{Your Name}
\date{\today}

\begin{document}

\maketitle

\section{Introduction}
This is a sample document demonstrating various \LaTeX{} features including mathematical equations, tables, and formatting. You can cite papers using the citation search feature, and they will be automatically added to your \texttt{cite.bib} file.

\section{Mathematical Expressions}
Here are some example equations:

\subsection{Quadratic Formula}
The quadratic formula is given by:
\[x = \frac{-b \pm \sqrt{b^2 - 4ac}}{2a}\]

\subsection{Matrix Example}
A sample matrix:
\[
\begin{pmatrix}
a_{11} & a_{12} & a_{13} \\
a_{21} & a_{22} & a_{23} \\
a_{31} & a_{32} & a_{33}
\end{pmatrix}
\]

\section{Tables}
Below is a sample table showing data:

\begin{table}[h]
\centering
\begin{tabular}{@{}lrrr@{}}
\toprule
\textbf{Category} & \textbf{Value 1} & \textbf{Value 2} & \textbf{Total} \\
\midrule
Item A & 10.5 & 20.3 & 30.8 \\
Item B & 15.2 & 25.1 & 40.3 \\
Item C & 22.1 & 12.7 & 34.8 \\
\bottomrule
\end{tabular}
\caption{Sample Data Table}
\label{tab:sample}
\end{table}

\section{Lists}
\subsection{Bullet Points}
\begin{itemize}
  \item First item
  \item Second item
  \item Third item with sub-items:
    \begin{itemize}
      \item Sub-item 1
      \item Sub-item 2
    \end{itemize}
\end{itemize}

\subsection{Numbered List}
\begin{enumerate}
  \item First step
  \item Second step
  \item Third step
\end{enumerate}

\section{Citations and References}
When you use the citation search feature, references will be automatically managed. For example, you might cite a paper like this: \cite{example2024}. All citations will be automatically added to your \texttt{cite.bib} file.

\section{Conclusion}
This template demonstrates basic \LaTeX{} features including mathematical formulas, tables, and lists. The bibliography system is now set up to work seamlessly with Bibby AI's citation management. Feel free to modify it for your needs!

\bibliographystyle{plain}
\bibliography{cite}

\end{document}\documentclass{article}
\usepackage[utf8]{inputenc}
\usepackage{booktabs}
\usepackage{amsmath}
\usepackage{natbib}

\title{My Document}
\author{Your Name}
\date{\today}

\begin{document}

\maketitle

\section{Introduction}
This is a sample document demonstrating various \LaTeX{} features including mathematical equations, tables, and formatting. You can cite papers using the citation search feature, and they will be automatically added to your \texttt{cite.bib} file.

\section{Mathematical Expressions}
Here are some example equations:

\subsection{Quadratic Formula}
The quadratic formula is given by:
\[x = \frac{-b \pm \sqrt{b^2 - 4ac}}{2a}\]

\subsection{Matrix Example}
A sample matrix:
\[
\begin{pmatrix}
a_{11} & a_{12} & a_{13} \\
a_{21} & a_{22} & a_{23} \\
a_{31} & a_{32} & a_{33}
\end{pmatrix}
\]

\section{Tables}
Below is a sample table showing data:

\begin{table}[h]
\centering
\begin{tabular}{@{}lrrr@{}}
\toprule
\textbf{Category} & \textbf{Value 1} & \textbf{Value 2} & \textbf{Total} \\
\midrule
Item A & 10.5 & 20.3 & 30.8 \\
Item B & 15.2 & 25.1 & 40.3 \\
Item C & 22.1 & 12.7 & 34.8 \\
\bottomrule
\end{tabular}
\caption{Sample Data Table}
\label{tab:sample}
\end{table}

\section{Lists}
\subsection{Bullet Points}
\begin{itemize}
  \item First item
  \item Second item
  \item Third item with sub-items:
    \begin{itemize}
      \item Sub-item 1
      \item Sub-item 2
    \end{itemize}
\end{itemize}

\subsection{Numbered List}
\begin{enumerate}
  \item First step
  \item Second step
  \item Third step
\end{enumerate}

\section{Citations and References}
When you use the citation search feature, references will be automatically managed. For example, you might cite a paper like this: \cite{example2024}. All citations will be automatically added to your \texttt{cite.bib} file.

\section{Conclusion}
This template demonstrates basic \LaTeX{} features including mathematical formulas, tables, and lists. The bibliography system is now set up to work seamlessly with Bibby AI's citation management. Feel free to modify it for your needs!

\bibliographystyle{plain}
\bibliography{cite}

\end{document}\documentclass{article}
\usepackage[utf8]{inputenc}
\usepackage{booktabs}
\usepackage{amsmath}
\usepackage{natbib}

\title{My Document}
\author{Your Name}
\date{\today}

\begin{document}

\maketitle

\section{Introduction}
This is a sample document demonstrating various \LaTeX{} features including mathematical equations, tables, and formatting. You can cite papers using the citation search feature, and they will be automatically added to your \texttt{cite.bib} file.

\section{Mathematical Expressions}
Here are some example equations:

\subsection{Quadratic Formula}
The quadratic formula is given by:
\[x = \frac{-b \pm \sqrt{b^2 - 4ac}}{2a}\]

\subsection{Matrix Example}
A sample matrix:
\[
\begin{pmatrix}
a_{11} & a_{12} & a_{13} \\
a_{21} & a_{22} & a_{23} \\
a_{31} & a_{32} & a_{33}
\end{pmatrix}
\]

\section{Tables}
Below is a sample table showing data:

\begin{table}[h]
\centering
\begin{tabular}{@{}lrrr@{}}
\toprule
\textbf{Category} & \textbf{Value 1} & \textbf{Value 2} & \textbf{Total} \\
\midrule
Item A & 10.5 & 20.3 & 30.8 \\
Item B & 15.2 & 25.1 & 40.3 \\
Item C & 22.1 & 12.7 & 34.8 \\
\bottomrule
\end{tabular}
\caption{Sample Data Table}
\label{tab:sample}
\end{table}

\section{Lists}
\subsection{Bullet Points}
\begin{itemize}
  \item First item
  \item Second item
  \item Third item with sub-items:
    \begin{itemize}
      \item Sub-item 1
      \item Sub-item 2
    \end{itemize}
\end{itemize}

\subsection{Numbered List}
\begin{enumerate}
  \item First step
  \item Second step
  \item Third step
\end{enumerate}

\section{Citations and References}
When you use the citation search feature, references will be automatically managed. For example, you might cite a paper like this: \cite{example2024}. All citations will be automatically added to your \texttt{cite.bib} file.

\section{Conclusion}
This template demonstrates basic \LaTeX{} features including mathematical formulas, tables, and lists. The bibliography system is now set up to work seamlessly with Bibby AI's citation management. Feel free to modify it for your needs!

\bibliographystyle{plain}
\bibliography{cite}

\end{document}\documentclass{article}
\usepackage[utf8]{inputenc}
\usepackage{booktabs}
\usepackage{amsmath}
\usepackage{natbib}

\title{My Document}
\author{Your Name}
\date{\today}

\begin{document}

\maketitle

\section{Introduction}
This is a sample document demonstrating various \LaTeX{} features including mathematical equations, tables, and formatting. You can cite papers using the citation search feature, and they will be automatically added to your \texttt{cite.bib} file.

\section{Mathematical Expressions}
Here are some example equations:

\subsection{Quadratic Formula}
The quadratic formula is given by:
\[x = \frac{-b \pm \sqrt{b^2 - 4ac}}{2a}\]

\subsection{Matrix Example}
A sample matrix:
\[
\begin{pmatrix}
a_{11} & a_{12} & a_{13} \\
a_{21} & a_{22} & a_{23} \\
a_{31} & a_{32} & a_{33}
\end{pmatrix}
\]

\section{Tables}
Below is a sample table showing data:

\begin{table}[h]
\centering
\begin{tabular}{@{}lrrr@{}}
\toprule
\textbf{Category} & \textbf{Value 1} & \textbf{Value 2} & \textbf{Total} \\
\midrule
Item A & 10.5 & 20.3 & 30.8 \\
Item B & 15.2 & 25.1 & 40.3 \\
Item C & 22.1 & 12.7 & 34.8 \\
\bottomrule
\end{tabular}
\caption{Sample Data Table}
\label{tab:sample}
\end{table}

\section{Lists}
\subsection{Bullet Points}
\begin{itemize}
  \item First item
  \item Second item
  \item Third item with sub-items:
    \begin{itemize}
      \item Sub-item 1
      \item Sub-item 2
    \end{itemize}
\end{itemize}

\subsection{Numbered List}
\begin{enumerate}
  \item First step
  \item Second step
  \item Third step
\end{enumerate}

\section{Citations and References}
When you use the citation search feature, references will be automatically managed. For example, you might cite a paper like this: \cite{example2024}. All citations will be automatically added to your \texttt{cite.bib} file.

\section{Conclusion}
This template demonstrates basic \LaTeX{} features including mathematical formulas, tables, and lists. The bibliography system is now set up to work seamlessly with Bibby AI's citation management. Feel free to modify it for your needs!

\bibliographystyle{plain}
\bibliography{cite}

\end{document}\documentclass{article}
\usepackage[utf8]{inputenc}
\usepackage{booktabs}
\usepackage{amsmath}
\usepackage{natbib}

\title{My Document}
\author{Your Name}
\date{\today}

\begin{document}

\maketitle

\section{Introduction}
This is a sample document demonstrating various \LaTeX{} features including mathematical equations, tables, and formatting. You can cite papers using the citation search feature, and they will be automatically added to your \texttt{cite.bib} file.

\section{Mathematical Expressions}
Here are some example equations:

\subsection{Quadratic Formula}
The quadratic formula is given by:
\[x = \frac{-b \pm \sqrt{b^2 - 4ac}}{2a}\]

\subsection{Matrix Example}
A sample matrix:
\[
\begin{pmatrix}
a_{11} & a_{12} & a_{13} \\
a_{21} & a_{22} & a_{23} \\
a_{31} & a_{32} & a_{33}
\end{pmatrix}
\]

\section{Tables}
Below is a sample table showing data:

\begin{table}[h]
\centering
\begin{tabular}{@{}lrrr@{}}
\toprule
\textbf{Category} & \textbf{Value 1} & \textbf{Value 2} & \textbf{Total} \\
\midrule
Item A & 10.5 & 20.3 & 30.8 \\
Item B & 15.2 & 25.1 & 40.3 \\
Item C & 22.1 & 12.7 & 34.8 \\
\bottomrule
\end{tabular}
\caption{Sample Data Table}
\label{tab:sample}
\end{table}

\section{Lists}
\subsection{Bullet Points}
\begin{itemize}
  \item First item
  \item Second item
  \item Third item with sub-items:
    \begin{itemize}
      \item Sub-item 1
      \item Sub-item 2
    \end{itemize}
\end{itemize}

\subsection{Numbered List}
\begin{enumerate}
  \item First step
  \item Second step
  \item Third step
\end{enumerate}

\section{Citations and References}
When you use the citation search feature, references will be automatically managed. For example, you might cite a paper like this: \cite{example2024}. All citations will be automatically added to your \texttt{cite.bib} file.

\section{Conclusion}
This template demonstrates basic \LaTeX{} features including mathematical formulas, tables, and lists. The bibliography system is now set up to work seamlessly with Bibby AI's citation management. Feel free to modify it for your needs!

\bibliographystyle{plain}
\bibliography{cite}

\end{document}\documentclass{article}
\usepackage[utf8]{inputenc}
\usepackage{booktabs}
\usepackage{amsmath}
\usepackage{natbib}

\title{My Document}
\author{Your Name}
\date{\today}

\begin{document}

\maketitle

\section{Introduction}
This is a sample document demonstrating various \LaTeX{} features including mathematical equations, tables, and formatting. You can cite papers using the citation search feature, and they will be automatically added to your \texttt{cite.bib} file.

\section{Mathematical Expressions}
Here are some example equations:

\subsection{Quadratic Formula}
The quadratic formula is given by:
\[x = \frac{-b \pm \sqrt{b^2 - 4ac}}{2a}\]

\subsection{Matrix Example}
A sample matrix:
\[
\begin{pmatrix}
a_{11} & a_{12} & a_{13} \\
a_{21} & a_{22} & a_{23} \\
a_{31} & a_{32} & a_{33}
\end{pmatrix}
\]

\section{Tables}
Below is a sample table showing data:

\begin{table}[h]
\centering
\begin{tabular}{@{}lrrr@{}}
\toprule
\textbf{Category} & \textbf{Value 1} & \textbf{Value 2} & \textbf{Total} \\
\midrule
Item A & 10.5 & 20.3 & 30.8 \\
Item B & 15.2 & 25.1 & 40.3 \\
Item C & 22.1 & 12.7 & 34.8 \\
\bottomrule
\end{tabular}
\caption{Sample Data Table}
\label{tab:sample}
\end{table}

\section{Lists}
\subsection{Bullet Points}
\begin{itemize}
  \item First item
  \item Second item
  \item Third item with sub-items:
    \begin{itemize}
      \item Sub-item 1
      \item Sub-item 2
    \end{itemize}
\end{itemize}

\subsection{Numbered List}
\begin{enumerate}
  \item First step
  \item Second step
  \item Third step
\end{enumerate}

\section{Citations and References}
When you use the citation search feature, references will be automatically managed. For example, you might cite a paper like this: \cite{example2024}. All citations will be automatically added to your \texttt{cite.bib} file.

\section{Conclusion}
This template demonstrates basic \LaTeX{} features including mathematical formulas, tables, and lists. The bibliography system is now set up to work seamlessly with Bibby AI's citation management. Feel free to modify it for your needs!

\bibliographystyle{plain}
\bibliography{cite}

\end{document}\documentclass{article}
\usepackage[utf8]{inputenc}
\usepackage{booktabs}
\usepackage{amsmath}
\usepackage{natbib}

\title{My Document}
\author{Your Name}
\date{\today}

\begin{document}

\maketitle

\section{Introduction}
This is a sample document demonstrating various \LaTeX{} features including mathematical equations, tables, and formatting. You can cite papers using the citation search feature, and they will be automatically added to your \texttt{cite.bib} file.

\section{Mathematical Expressions}
Here are some example equations:

\subsection{Quadratic Formula}
The quadratic formula is given by:
\[x = \frac{-b \pm \sqrt{b^2 - 4ac}}{2a}\]

\subsection{Matrix Example}
A sample matrix:
\[
\begin{pmatrix}
a_{11} & a_{12} & a_{13} \\
a_{21} & a_{22} & a_{23} \\
a_{31} & a_{32} & a_{33}
\end{pmatrix}
\]

\section{Tables}
Below is a sample table showing data:

\begin{table}[h]
\centering
\begin{tabular}{@{}lrrr@{}}
\toprule
\textbf{Category} & \textbf{Value 1} & \textbf{Value 2} & \textbf{Total} \\
\midrule
Item A & 10.5 & 20.3 & 30.8 \\
Item B & 15.2 & 25.1 & 40.3 \\
Item C & 22.1 & 12.7 & 34.8 \\
\bottomrule
\end{tabular}
\caption{Sample Data Table}
\label{tab:sample}
\end{table}

\section{Lists}
\subsection{Bullet Points}
\begin{itemize}
  \item First item
  \item Second item
  \item Third item with sub-items:
    \begin{itemize}
      \item Sub-item 1
      \item Sub-item 2
    \end{itemize}
\end{itemize}

\subsection{Numbered List}
\begin{enumerate}
  \item First step
  \item Second step
  \item Third step
\end{enumerate}

\section{Citations and References}
When you use the citation search feature, references will be automatically managed. For example, you might cite a paper like this: \cite{example2024}. All citations will be automatically added to your \texttt{cite.bib} file.

\section{Conclusion}
This template demonstrates basic \LaTeX{} features including mathematical formulas, tables, and lists. The bibliography system is now set up to work seamlessly with Bibby AI's citation management. Feel free to modify it for your needs!

\bibliographystyle{plain}
\bibliography{cite}

\end{document}\documentclass{article}
\usepackage[utf8]{inputenc}
\usepackage{booktabs}
\usepackage{amsmath}
\usepackage{natbib}

\title{My Document}
\author{Your Name}
\date{\today}

\begin{document}

\maketitle

\section{Introduction}
This is a sample document demonstrating various \LaTeX{} features including mathematical equations, tables, and formatting. You can cite papers using the citation search feature, and they will be automatically added to your \texttt{cite.bib} file.

\section{Mathematical Expressions}
Here are some example equations:

\subsection{Quadratic Formula}
The quadratic formula is given by:
\[x = \frac{-b \pm \sqrt{b^2 - 4ac}}{2a}\]

\subsection{Matrix Example}
A sample matrix:
\[
\begin{pmatrix}
a_{11} & a_{12} & a_{13} \\
a_{21} & a_{22} & a_{23} \\
a_{31} & a_{32} & a_{33}
\end{pmatrix}
\]

\section{Tables}
Below is a sample table showing data:

\begin{table}[h]
\centering
\begin{tabular}{@{}lrrr@{}}
\toprule
\textbf{Category} & \textbf{Value 1} & \textbf{Value 2} & \textbf{Total} \\
\midrule
Item A & 10.5 & 20.3 & 30.8 \\
Item B & 15.2 & 25.1 & 40.3 \\
Item C & 22.1 & 12.7 & 34.8 \\
\bottomrule
\end{tabular}
\caption{Sample Data Table}
\label{tab:sample}
\end{table}

\section{Lists}
\subsection{Bullet Points}
\begin{itemize}
  \item First item
  \item Second item
  \item Third item with sub-items:
    \begin{itemize}
      \item Sub-item 1
      \item Sub-item 2
    \end{itemize}
\end{itemize}

\subsection{Numbered List}
\begin{enumerate}
  \item First step
  \item Second step
  \item Third step
\end{enumerate}

\section{Citations and References}
When you use the citation search feature, references will be automatically managed. For example, you might cite a paper like this: \cite{example2024}. All citations will be automatically added to your \texttt{cite.bib} file.

\section{Conclusion}
This template demonstrates basic \LaTeX{} features including mathematical formulas, tables, and lists. The bibliography system is now set up to work seamlessly with Bibby AI's citation management. Feel free to modify it for your needs!

\bibliographystyle{plain}
\bibliography{cite}

\end{document}\documentclass{article}
\usepackage[utf8]{inputenc}
\usepackage{booktabs}
\usepackage{amsmath}
\usepackage{natbib}

\title{My Document}
\author{Your Name}
\date{\today}

\begin{document}

\maketitle

\section{Introduction}
This is a sample document demonstrating various \LaTeX{} features including mathematical equations, tables, and formatting. You can cite papers using the citation search feature, and they will be automatically added to your \texttt{cite.bib} file.

\section{Mathematical Expressions}
Here are some example equations:

\subsection{Quadratic Formula}
The quadratic formula is given by:
\[x = \frac{-b \pm \sqrt{b^2 - 4ac}}{2a}\]

\subsection{Matrix Example}
A sample matrix:
\[
\begin{pmatrix}
a_{11} & a_{12} & a_{13} \\
a_{21} & a_{22} & a_{23} \\
a_{31} & a_{32} & a_{33}
\end{pmatrix}
\]

\section{Tables}
Below is a sample table showing data:

\begin{table}[h]
\centering
\begin{tabular}{@{}lrrr@{}}
\toprule
\textbf{Category} & \textbf{Value 1} & \textbf{Value 2} & \textbf{Total} \\
\midrule
Item A & 10.5 & 20.3 & 30.8 \\
Item B & 15.2 & 25.1 & 40.3 \\
Item C & 22.1 & 12.7 & 34.8 \\
\bottomrule
\end{tabular}
\caption{Sample Data Table}
\label{tab:sample}
\end{table}

\section{Lists}
\subsection{Bullet Points}
\begin{itemize}
  \item First item
  \item Second item
  \item Third item with sub-items:
    \begin{itemize}
      \item Sub-item 1
      \item Sub-item 2
    \end{itemize}
\end{itemize}

\subsection{Numbered List}
\begin{enumerate}
  \item First step
  \item Second step
  \item Third step
\end{enumerate}

\section{Citations and References}
When you use the citation search feature, references will be automatically managed. For example, you might cite a paper like this: \cite{example2024}. All citations will be automatically added to your \texttt{cite.bib} file.

\section{Conclusion}
This template demonstrates basic \LaTeX{} features including mathematical formulas, tables, and lists. The bibliography system is now set up to work seamlessly with Bibby AI's citation management. Feel free to modify it for your needs!

\bibliographystyle{plain}
\bibliography{cite}

\end{document}\documentclass{article}
\usepackage[utf8]{inputenc}
\usepackage{booktabs}
\usepackage{amsmath}
\usepackage{natbib}

\title{My Document}
\author{Your Name}
\date{\today}

\begin{document}

\maketitle

\section{Introduction}
This is a sample document demonstrating various \LaTeX{} features including mathematical equations, tables, and formatting. You can cite papers using the citation search feature, and they will be automatically added to your \texttt{cite.bib} file.

\section{Mathematical Expressions}
Here are some example equations:

\subsection{Quadratic Formula}
The quadratic formula is given by:
\[x = \frac{-b \pm \sqrt{b^2 - 4ac}}{2a}\]

\subsection{Matrix Example}
A sample matrix:
\[
\begin{pmatrix}
a_{11} & a_{12} & a_{13} \\
a_{21} & a_{22} & a_{23} \\
a_{31} & a_{32} & a_{33}
\end{pmatrix}
\]

\section{Tables}
Below is a sample table showing data:

\begin{table}[h]
\centering
\begin{tabular}{@{}lrrr@{}}
\toprule
\textbf{Category} & \textbf{Value 1} & \textbf{Value 2} & \textbf{Total} \\
\midrule
Item A & 10.5 & 20.3 & 30.8 \\
Item B & 15.2 & 25.1 & 40.3 \\
Item C & 22.1 & 12.7 & 34.8 \\
\bottomrule
\end{tabular}
\caption{Sample Data Table}
\label{tab:sample}
\end{table}

\section{Lists}
\subsection{Bullet Points}
\begin{itemize}
  \item First item
  \item Second item
  \item Third item with sub-items:
    \begin{itemize}
      \item Sub-item 1
      \item Sub-item 2
    \end{itemize}
\end{itemize}

\subsection{Numbered List}
\begin{enumerate}
  \item First step
  \item Second step
  \item Third step
\end{enumerate}

\section{Citations and References}
When you use the citation search feature, references will be automatically managed. For example, you might cite a paper like this: \cite{example2024}. All citations will be automatically added to your \texttt{cite.bib} file.

\section{Conclusion}
This template demonstrates basic \LaTeX{} features including mathematical formulas, tables, and lists. The bibliography system is now set up to work seamlessly with Bibby AI's citation management. Feel free to modify it for your needs!

\bibliographystyle{plain}
\bibliography{cite}

\end{document}\documentclass{article}
\usepackage[utf8]{inputenc}
\usepackage{booktabs}
\usepackage{amsmath}
\usepackage{natbib}

\title{My Document}
\author{Your Name}
\date{\today}

\begin{document}

\maketitle

\section{Introduction}
This is a sample document demonstrating various \LaTeX{} features including mathematical equations, tables, and formatting. You can cite papers using the citation search feature, and they will be automatically added to your \texttt{cite.bib} file.

\section{Mathematical Expressions}
Here are some example equations:

\subsection{Quadratic Formula}
The quadratic formula is given by:
\[x = \frac{-b \pm \sqrt{b^2 - 4ac}}{2a}\]

\subsection{Matrix Example}
A sample matrix:
\[
\begin{pmatrix}
a_{11} & a_{12} & a_{13} \\
a_{21} & a_{22} & a_{23} \\
a_{31} & a_{32} & a_{33}
\end{pmatrix}
\]

\section{Tables}
Below is a sample table showing data:

\begin{table}[h]
\centering
\begin{tabular}{@{}lrrr@{}}
\toprule
\textbf{Category} & \textbf{Value 1} & \textbf{Value 2} & \textbf{Total} \\
\midrule
Item A & 10.5 & 20.3 & 30.8 \\
Item B & 15.2 & 25.1 & 40.3 \\
Item C & 22.1 & 12.7 & 34.8 \\
\bottomrule
\end{tabular}
\caption{Sample Data Table}
\label{tab:sample}
\end{table}

\section{Lists}
\subsection{Bullet Points}
\begin{itemize}
  \item First item
  \item Second item
  \item Third item with sub-items:
    \begin{itemize}
      \item Sub-item 1
      \item Sub-item 2
    \end{itemize}
\end{itemize}

\subsection{Numbered List}
\begin{enumerate}
  \item First step
  \item Second step
  \item Third step
\end{enumerate}

\section{Citations and References}
When you use the citation search feature, references will be automatically managed. For example, you might cite a paper like this: \cite{example2024}. All citations will be automatically added to your \texttt{cite.bib} file.

\section{Conclusion}
This template demonstrates basic \LaTeX{} features including mathematical formulas, tables, and lists. The bibliography system is now set up to work seamlessly with Bibby AI's citation management. Feel free to modify it for your needs!

\bibliographystyle{plain}
\bibliography{cite}

\end{document}\documentclass{article}
\usepackage[utf8]{inputenc}
\usepackage{booktabs}
\usepackage{amsmath}
\usepackage{natbib}

\title{My Document}
\author{Your Name}
\date{\today}

\begin{document}

\maketitle

\section{Introduction}
This is a sample document demonstrating various \LaTeX{} features including mathematical equations, tables, and formatting. You can cite papers using the citation search feature, and they will be automatically added to your \texttt{cite.bib} file.

\section{Mathematical Expressions}
Here are some example equations:

\subsection{Quadratic Formula}
The quadratic formula is given by:
\[x = \frac{-b \pm \sqrt{b^2 - 4ac}}{2a}\]

\subsection{Matrix Example}
A sample matrix:
\[
\begin{pmatrix}
a_{11} & a_{12} & a_{13} \\
a_{21} & a_{22} & a_{23} \\
a_{31} & a_{32} & a_{33}
\end{pmatrix}
\]

\section{Tables}
Below is a sample table showing data:

\begin{table}[h]
\centering
\begin{tabular}{@{}lrrr@{}}
\toprule
\textbf{Category} & \textbf{Value 1} & \textbf{Value 2} & \textbf{Total} \\
\midrule
Item A & 10.5 & 20.3 & 30.8 \\
Item B & 15.2 & 25.1 & 40.3 \\
Item C & 22.1 & 12.7 & 34.8 \\
\bottomrule
\end{tabular}
\caption{Sample Data Table}
\label{tab:sample}
\end{table}

\section{Lists}
\subsection{Bullet Points}
\begin{itemize}
  \item First item
  \item Second item
  \item Third item with sub-items:
    \begin{itemize}
      \item Sub-item 1
      \item Sub-item 2
    \end{itemize}
\end{itemize}

\subsection{Numbered List}
\begin{enumerate}
  \item First step
  \item Second step
  \item Third step
\end{enumerate}

\section{Citations and References}
When you use the citation search feature, references will be automatically managed. For example, you might cite a paper like this: \cite{example2024}. All citations will be automatically added to your \texttt{cite.bib} file.

\section{Conclusion}
This template demonstrates basic \LaTeX{} features including mathematical formulas, tables, and lists. The bibliography system is now set up to work seamlessly with Bibby AI's citation management. Feel free to modify it for your needs!

\bibliographystyle{plain}
\bibliography{cite}

\end{document}\documentclass{article}
\usepackage[utf8]{inputenc}
\usepackage{booktabs}
\usepackage{amsmath}
\usepackage{natbib}

\title{My Document}
\author{Your Name}
\date{\today}

\begin{document}

\maketitle

\section{Introduction}
This is a sample document demonstrating various \LaTeX{} features including mathematical equations, tables, and formatting. You can cite papers using the citation search feature, and they will be automatically added to your \texttt{cite.bib} file.

\section{Mathematical Expressions}
Here are some example equations:

\subsection{Quadratic Formula}
The quadratic formula is given by:
\[x = \frac{-b \pm \sqrt{b^2 - 4ac}}{2a}\]

\subsection{Matrix Example}
A sample matrix:
\[
\begin{pmatrix}
a_{11} & a_{12} & a_{13} \\
a_{21} & a_{22} & a_{23} \\
a_{31} & a_{32} & a_{33}
\end{pmatrix}
\]

\section{Tables}
Below is a sample table showing data:

\begin{table}[h]
\centering
\begin{tabular}{@{}lrrr@{}}
\toprule
\textbf{Category} & \textbf{Value 1} & \textbf{Value 2} & \textbf{Total} \\
\midrule
Item A & 10.5 & 20.3 & 30.8 \\
Item B & 15.2 & 25.1 & 40.3 \\
Item C & 22.1 & 12.7 & 34.8 \\
\bottomrule
\end{tabular}
\caption{Sample Data Table}
\label{tab:sample}
\end{table}

\section{Lists}
\subsection{Bullet Points}
\begin{itemize}
  \item First item
  \item Second item
  \item Third item with sub-items:
    \begin{itemize}
      \item Sub-item 1
      \item Sub-item 2
    \end{itemize}
\end{itemize}

\subsection{Numbered List}
\begin{enumerate}
  \item First step
  \item Second step
  \item Third step
\end{enumerate}

\section{Citations and References}
When you use the citation search feature, references will be automatically managed. For example, you might cite a paper like this: \cite{example2024}. All citations will be automatically added to your \texttt{cite.bib} file.

\section{Conclusion}
This template demonstrates basic \LaTeX{} features including mathematical formulas, tables, and lists. The bibliography system is now set up to work seamlessly with Bibby AI's citation management. Feel free to modify it for your needs!

\bibliographystyle{plain}
\bibliography{cite}

\end{document}\documentclass{article}
\usepackage[utf8]{inputenc}
\usepackage{booktabs}
\usepackage{amsmath}
\usepackage{natbib}

\title{My Document}
\author{Your Name}
\date{\today}

\begin{document}

\maketitle

\section{Introduction}
This is a sample document demonstrating various \LaTeX{} features including mathematical equations, tables, and formatting. You can cite papers using the citation search feature, and they will be automatically added to your \texttt{cite.bib} file.

\section{Mathematical Expressions}
Here are some example equations:

\subsection{Quadratic Formula}
The quadratic formula is given by:
\[x = \frac{-b \pm \sqrt{b^2 - 4ac}}{2a}\]

\subsection{Matrix Example}
A sample matrix:
\[
\begin{pmatrix}
a_{11} & a_{12} & a_{13} \\
a_{21} & a_{22} & a_{23} \\
a_{31} & a_{32} & a_{33}
\end{pmatrix}
\]

\section{Tables}
Below is a sample table showing data:

\begin{table}[h]
\centering
\begin{tabular}{@{}lrrr@{}}
\toprule
\textbf{Category} & \textbf{Value 1} & \textbf{Value 2} & \textbf{Total} \\
\midrule
Item A & 10.5 & 20.3 & 30.8 \\
Item B & 15.2 & 25.1 & 40.3 \\
Item C & 22.1 & 12.7 & 34.8 \\
\bottomrule
\end{tabular}
\caption{Sample Data Table}
\label{tab:sample}
\end{table}

\section{Lists}
\subsection{Bullet Points}
\begin{itemize}
  \item First item
  \item Second item
  \item Third item with sub-items:
    \begin{itemize}
      \item Sub-item 1
      \item Sub-item 2
    \end{itemize}
\end{itemize}

\subsection{Numbered List}
\begin{enumerate}
  \item First step
  \item Second step
  \item Third step
\end{enumerate}

\section{Citations and References}
When you use the citation search feature, references will be automatically managed. For example, you might cite a paper like this: \cite{example2024}. All citations will be automatically added to your \texttt{cite.bib} file.

\section{Conclusion}
This template demonstrates basic \LaTeX{} features including mathematical formulas, tables, and lists. The bibliography system is now set up to work seamlessly with Bibby AI's citation management. Feel free to modify it for your needs!

\bibliographystyle{plain}
\bibliography{cite}

\end{document}\documentclass{article}
\usepackage[utf8]{inputenc}
\usepackage{booktabs}
\usepackage{amsmath}
\usepackage{natbib}

\title{My Document}
\author{Your Name}
\date{\today}

\begin{document}

\maketitle

\section{Introduction}
This is a sample document demonstrating various \LaTeX{} features including mathematical equations, tables, and formatting. You can cite papers using the citation search feature, and they will be automatically added to your \texttt{cite.bib} file.

\section{Mathematical Expressions}
Here are some example equations:

\subsection{Quadratic Formula}
The quadratic formula is given by:
\[x = \frac{-b \pm \sqrt{b^2 - 4ac}}{2a}\]

\subsection{Matrix Example}
A sample matrix:
\[
\begin{pmatrix}
a_{11} & a_{12} & a_{13} \\
a_{21} & a_{22} & a_{23} \\
a_{31} & a_{32} & a_{33}
\end{pmatrix}
\]

\section{Tables}
Below is a sample table showing data:

\begin{table}[h]
\centering
\begin{tabular}{@{}lrrr@{}}
\toprule
\textbf{Category} & \textbf{Value 1} & \textbf{Value 2} & \textbf{Total} \\
\midrule
Item A & 10.5 & 20.3 & 30.8 \\
Item B & 15.2 & 25.1 & 40.3 \\
Item C & 22.1 & 12.7 & 34.8 \\
\bottomrule
\end{tabular}
\caption{Sample Data Table}
\label{tab:sample}
\end{table}

\section{Lists}
\subsection{Bullet Points}
\begin{itemize}
  \item First item
  \item Second item
  \item Third item with sub-items:
    \begin{itemize}
      \item Sub-item 1
      \item Sub-item 2
    \end{itemize}
\end{itemize}

\subsection{Numbered List}
\begin{enumerate}
  \item First step
  \item Second step
  \item Third step
\end{enumerate}

\section{Citations and References}
When you use the citation search feature, references will be automatically managed. For example, you might cite a paper like this: \cite{example2024}. All citations will be automatically added to your \texttt{cite.bib} file.

\section{Conclusion}
This template demonstrates basic \LaTeX{} features including mathematical formulas, tables, and lists. The bibliography system is now set up to work seamlessly with Bibby AI's citation management. Feel free to modify it for your needs!

\bibliographystyle{plain}
\bibliography{cite}

\end{document}\documentclass{article}
\usepackage[utf8]{inputenc}
\usepackage{booktabs}
\usepackage{amsmath}
\usepackage{natbib}

\title{My Document}
\author{Your Name}
\date{\today}

\begin{document}

\maketitle

\section{Introduction}
This is a sample document demonstrating various \LaTeX{} features including mathematical equations, tables, and formatting. You can cite papers using the citation search feature, and they will be automatically added to your \texttt{cite.bib} file.

\section{Mathematical Expressions}
Here are some example equations:

\subsection{Quadratic Formula}
The quadratic formula is given by:
\[x = \frac{-b \pm \sqrt{b^2 - 4ac}}{2a}\]

\subsection{Matrix Example}
A sample matrix:
\[
\begin{pmatrix}
a_{11} & a_{12} & a_{13} \\
a_{21} & a_{22} & a_{23} \\
a_{31} & a_{32} & a_{33}
\end{pmatrix}
\]

\section{Tables}
Below is a sample table showing data:

\begin{table}[h]
\centering
\begin{tabular}{@{}lrrr@{}}
\toprule
\textbf{Category} & \textbf{Value 1} & \textbf{Value 2} & \textbf{Total} \\
\midrule
Item A & 10.5 & 20.3 & 30.8 \\
Item B & 15.2 & 25.1 & 40.3 \\
Item C & 22.1 & 12.7 & 34.8 \\
\bottomrule
\end{tabular}
\caption{Sample Data Table}
\label{tab:sample}
\end{table}

\section{Lists}
\subsection{Bullet Points}
\begin{itemize}
  \item First item
  \item Second item
  \item Third item with sub-items:
    \begin{itemize}
      \item Sub-item 1
      \item Sub-item 2
    \end{itemize}
\end{itemize}

\subsection{Numbered List}
\begin{enumerate}
  \item First step
  \item Second step
  \item Third step
\end{enumerate}

\section{Citations and References}
When you use the citation search feature, references will be automatically managed. For example, you might cite a paper like this: \cite{example2024}. All citations will be automatically added to your \texttt{cite.bib} file.

\section{Conclusion}
This template demonstrates basic \LaTeX{} features including mathematical formulas, tables, and lists. The bibliography system is now set up to work seamlessly with Bibby AI's citation management. Feel free to modify it for your needs!

\bibliographystyle{plain}
\bibliography{cite}

\end{document}\documentclass{article}
\usepackage[utf8]{inputenc}
\usepackage{booktabs}
\usepackage{amsmath}
\usepackage{natbib}

\title{My Document}
\author{Your Name}
\date{\today}

\begin{document}

\maketitle

\section{Introduction}
This is a sample document demonstrating various \LaTeX{} features including mathematical equations, tables, and formatting. You can cite papers using the citation search feature, and they will be automatically added to your \texttt{cite.bib} file.

\section{Mathematical Expressions}
Here are some example equations:

\subsection{Quadratic Formula}
The quadratic formula is given by:
\[x = \frac{-b \pm \sqrt{b^2 - 4ac}}{2a}\]

\subsection{Matrix Example}
A sample matrix:
\[
\begin{pmatrix}
a_{11} & a_{12} & a_{13} \\
a_{21} & a_{22} & a_{23} \\
a_{31} & a_{32} & a_{33}
\end{pmatrix}
\]

\section{Tables}
Below is a sample table showing data:

\begin{table}[h]
\centering
\begin{tabular}{@{}lrrr@{}}
\toprule
\textbf{Category} & \textbf{Value 1} & \textbf{Value 2} & \textbf{Total} \\
\midrule
Item A & 10.5 & 20.3 & 30.8 \\
Item B & 15.2 & 25.1 & 40.3 \\
Item C & 22.1 & 12.7 & 34.8 \\
\bottomrule
\end{tabular}
\caption{Sample Data Table}
\label{tab:sample}
\end{table}

\section{Lists}
\subsection{Bullet Points}
\begin{itemize}
  \item First item
  \item Second item
  \item Third item with sub-items:
    \begin{itemize}
      \item Sub-item 1
      \item Sub-item 2
    \end{itemize}
\end{itemize}

\subsection{Numbered List}
\begin{enumerate}
  \item First step
  \item Second step
  \item Third step
\end{enumerate}

\section{Citations and References}
When you use the citation search feature, references will be automatically managed. For example, you might cite a paper like this: \cite{example2024}. All citations will be automatically added to your \texttt{cite.bib} file.

\section{Conclusion}
This template demonstrates basic \LaTeX{} features including mathematical formulas, tables, and lists. The bibliography system is now set up to work seamlessly with Bibby AI's citation management. Feel free to modify it for your needs!

\bibliographystyle{plain}
\bibliography{cite}

\end{document}\documentclass{article}
\usepackage[utf8]{inputenc}
\usepackage{booktabs}
\usepackage{amsmath}
\usepackage{natbib}

\title{My Document}
\author{Your Name}
\date{\today}

\begin{document}

\maketitle

\section{Introduction}
This is a sample document demonstrating various \LaTeX{} features including mathematical equations, tables, and formatting. You can cite papers using the citation search feature, and they will be automatically added to your \texttt{cite.bib} file.

\section{Mathematical Expressions}
Here are some example equations:

\subsection{Quadratic Formula}
The quadratic formula is given by:
\[x = \frac{-b \pm \sqrt{b^2 - 4ac}}{2a}\]

\subsection{Matrix Example}
A sample matrix:
\[
\begin{pmatrix}
a_{11} & a_{12} & a_{13} \\
a_{21} & a_{22} & a_{23} \\
a_{31} & a_{32} & a_{33}
\end{pmatrix}
\]

\section{Tables}
Below is a sample table showing data:

\begin{table}[h]
\centering
\begin{tabular}{@{}lrrr@{}}
\toprule
\textbf{Category} & \textbf{Value 1} & \textbf{Value 2} & \textbf{Total} \\
\midrule
Item A & 10.5 & 20.3 & 30.8 \\
Item B & 15.2 & 25.1 & 40.3 \\
Item C & 22.1 & 12.7 & 34.8 \\
\bottomrule
\end{tabular}
\caption{Sample Data Table}
\label{tab:sample}
\end{table}

\section{Lists}
\subsection{Bullet Points}
\begin{itemize}
  \item First item
  \item Second item
  \item Third item with sub-items:
    \begin{itemize}
      \item Sub-item 1
      \item Sub-item 2
    \end{itemize}
\end{itemize}

\subsection{Numbered List}
\begin{enumerate}
  \item First step
  \item Second step
  \item Third step
\end{enumerate}

\section{Citations and References}
When you use the citation search feature, references will be automatically managed. For example, you might cite a paper like this: \cite{example2024}. All citations will be automatically added to your \texttt{cite.bib} file.

\section{Conclusion}
This template demonstrates basic \LaTeX{} features including mathematical formulas, tables, and lists. The bibliography system is now set up to work seamlessly with Bibby AI's citation management. Feel free to modify it for your needs!

\bibliographystyle{plain}
\bibliography{cite}

\end{document}\documentclass{article}
\usepackage[utf8]{inputenc}
\usepackage{booktabs}
\usepackage{amsmath}
\usepackage{natbib}

\title{My Document}
\author{Your Name}
\date{\today}

\begin{document}

\maketitle

\section{Introduction}
This is a sample document demonstrating various \LaTeX{} features including mathematical equations, tables, and formatting. You can cite papers using the citation search feature, and they will be automatically added to your \texttt{cite.bib} file.

\section{Mathematical Expressions}
Here are some example equations:

\subsection{Quadratic Formula}
The quadratic formula is given by:
\[x = \frac{-b \pm \sqrt{b^2 - 4ac}}{2a}\]

\subsection{Matrix Example}
A sample matrix:
\[
\begin{pmatrix}
a_{11} & a_{12} & a_{13} \\
a_{21} & a_{22} & a_{23} \\
a_{31} & a_{32} & a_{33}
\end{pmatrix}
\]

\section{Tables}
Below is a sample table showing data:

\begin{table}[h]
\centering
\begin{tabular}{@{}lrrr@{}}
\toprule
\textbf{Category} & \textbf{Value 1} & \textbf{Value 2} & \textbf{Total} \\
\midrule
Item A & 10.5 & 20.3 & 30.8 \\
Item B & 15.2 & 25.1 & 40.3 \\
Item C & 22.1 & 12.7 & 34.8 \\
\bottomrule
\end{tabular}
\caption{Sample Data Table}
\label{tab:sample}
\end{table}

\section{Lists}
\subsection{Bullet Points}
\begin{itemize}
  \item First item
  \item Second item
  \item Third item with sub-items:
    \begin{itemize}
      \item Sub-item 1
      \item Sub-item 2
    \end{itemize}
\end{itemize}

\subsection{Numbered List}
\begin{enumerate}
  \item First step
  \item Second step
  \item Third step
\end{enumerate}

\section{Citations and References}
When you use the citation search feature, references will be automatically managed. For example, you might cite a paper like this: \cite{example2024}. All citations will be automatically added to your \texttt{cite.bib} file.

\section{Conclusion}
This template demonstrates basic \LaTeX{} features including mathematical formulas, tables, and lists. The bibliography system is now set up to work seamlessly with Bibby AI's citation management. Feel free to modify it for your needs!

\bibliographystyle{plain}
\bibliography{cite}

\end{document}\documentclass{article}
\usepackage[utf8]{inputenc}
\usepackage{booktabs}
\usepackage{amsmath}
\usepackage{natbib}

\title{My Document}
\author{Your Name}
\date{\today}

\begin{document}

\maketitle

\section{Introduction}
This is a sample document demonstrating various \LaTeX{} features including mathematical equations, tables, and formatting. You can cite papers using the citation search feature, and they will be automatically added to your \texttt{cite.bib} file.

\section{Mathematical Expressions}
Here are some example equations:

\subsection{Quadratic Formula}
The quadratic formula is given by:
\[x = \frac{-b \pm \sqrt{b^2 - 4ac}}{2a}\]

\subsection{Matrix Example}
A sample matrix:
\[
\begin{pmatrix}
a_{11} & a_{12} & a_{13} \\
a_{21} & a_{22} & a_{23} \\
a_{31} & a_{32} & a_{33}
\end{pmatrix}
\]

\section{Tables}
Below is a sample table showing data:

\begin{table}[h]
\centering
\begin{tabular}{@{}lrrr@{}}
\toprule
\textbf{Category} & \textbf{Value 1} & \textbf{Value 2} & \textbf{Total} \\
\midrule
Item A & 10.5 & 20.3 & 30.8 \\
Item B & 15.2 & 25.1 & 40.3 \\
Item C & 22.1 & 12.7 & 34.8 \\
\bottomrule
\end{tabular}
\caption{Sample Data Table}
\label{tab:sample}
\end{table}

\section{Lists}
\subsection{Bullet Points}
\begin{itemize}
  \item First item
  \item Second item
  \item Third item with sub-items:
    \begin{itemize}
      \item Sub-item 1
      \item Sub-item 2
    \end{itemize}
\end{itemize}

\subsection{Numbered List}
\begin{enumerate}
  \item First step
  \item Second step
  \item Third step
\end{enumerate}

\section{Citations and References}
When you use the citation search feature, references will be automatically managed. For example, you might cite a paper like this: \cite{example2024}. All citations will be automatically added to your \texttt{cite.bib} file.

\section{Conclusion}
This template demonstrates basic \LaTeX{} features including mathematical formulas, tables, and lists. The bibliography system is now set up to work seamlessly with Bibby AI's citation management. Feel free to modify it for your needs!

\bibliographystyle{plain}
\bibliography{cite}

\end{document}\documentclass{article}
\usepackage[utf8]{inputenc}
\usepackage{booktabs}
\usepackage{amsmath}
\usepackage{natbib}

\title{My Document}
\author{Your Name}
\date{\today}

\begin{document}

\maketitle

\section{Introduction}
This is a sample document demonstrating various \LaTeX{} features including mathematical equations, tables, and formatting. You can cite papers using the citation search feature, and they will be automatically added to your \texttt{cite.bib} file.

\section{Mathematical Expressions}
Here are some example equations:

\subsection{Quadratic Formula}
The quadratic formula is given by:
\[x = \frac{-b \pm \sqrt{b^2 - 4ac}}{2a}\]

\subsection{Matrix Example}
A sample matrix:
\[
\begin{pmatrix}
a_{11} & a_{12} & a_{13} \\
a_{21} & a_{22} & a_{23} \\
a_{31} & a_{32} & a_{33}
\end{pmatrix}
\]

\section{Tables}
Below is a sample table showing data:

\begin{table}[h]
\centering
\begin{tabular}{@{}lrrr@{}}
\toprule
\textbf{Category} & \textbf{Value 1} & \textbf{Value 2} & \textbf{Total} \\
\midrule
Item A & 10.5 & 20.3 & 30.8 \\
Item B & 15.2 & 25.1 & 40.3 \\
Item C & 22.1 & 12.7 & 34.8 \\
\bottomrule
\end{tabular}
\caption{Sample Data Table}
\label{tab:sample}
\end{table}

\section{Lists}
\subsection{Bullet Points}
\begin{itemize}
  \item First item
  \item Second item
  \item Third item with sub-items:
    \begin{itemize}
      \item Sub-item 1
      \item Sub-item 2
    \end{itemize}
\end{itemize}

\subsection{Numbered List}
\begin{enumerate}
  \item First step
  \item Second step
  \item Third step
\end{enumerate}

\section{Citations and References}
When you use the citation search feature, references will be automatically managed. For example, you might cite a paper like this: \cite{example2024}. All citations will be automatically added to your \texttt{cite.bib} file.

\section{Conclusion}
This template demonstrates basic \LaTeX{} features including mathematical formulas, tables, and lists. The bibliography system is now set up to work seamlessly with Bibby AI's citation management. Feel free to modify it for your needs!

\bibliographystyle{plain}
\bibliography{cite}

\end{document}\documentclass{article}
\usepackage[utf8]{inputenc}
\usepackage{booktabs}
\usepackage{amsmath}
\usepackage{natbib}

\title{My Document}
\author{Your Name}
\date{\today}

\begin{document}

\maketitle

\section{Introduction}
This is a sample document demonstrating various \LaTeX{} features including mathematical equations, tables, and formatting. You can cite papers using the citation search feature, and they will be automatically added to your \texttt{cite.bib} file.

\section{Mathematical Expressions}
Here are some example equations:

\subsection{Quadratic Formula}
The quadratic formula is given by:
\[x = \frac{-b \pm \sqrt{b^2 - 4ac}}{2a}\]

\subsection{Matrix Example}
A sample matrix:
\[
\begin{pmatrix}
a_{11} & a_{12} & a_{13} \\
a_{21} & a_{22} & a_{23} \\
a_{31} & a_{32} & a_{33}
\end{pmatrix}
\]

\section{Tables}
Below is a sample table showing data:

\begin{table}[h]
\centering
\begin{tabular}{@{}lrrr@{}}
\toprule
\textbf{Category} & \textbf{Value 1} & \textbf{Value 2} & \textbf{Total} \\
\midrule
Item A & 10.5 & 20.3 & 30.8 \\
Item B & 15.2 & 25.1 & 40.3 \\
Item C & 22.1 & 12.7 & 34.8 \\
\bottomrule
\end{tabular}
\caption{Sample Data Table}
\label{tab:sample}
\end{table}

\section{Lists}
\subsection{Bullet Points}
\begin{itemize}
  \item First item
  \item Second item
  \item Third item with sub-items:
    \begin{itemize}
      \item Sub-item 1
      \item Sub-item 2
    \end{itemize}
\end{itemize}

\subsection{Numbered List}
\begin{enumerate}
  \item First step
  \item Second step
  \item Third step
\end{enumerate}

\section{Citations and References}
When you use the citation search feature, references will be automatically managed. For example, you might cite a paper like this: \cite{example2024}. All citations will be automatically added to your \texttt{cite.bib} file.

\section{Conclusion}
This template demonstrates basic \LaTeX{} features including mathematical formulas, tables, and lists. The bibliography system is now set up to work seamlessly with Bibby AI's citation management. Feel free to modify it for your needs!

\bibliographystyle{plain}
\bibliography{cite}

\end{document}\documentclass{article}
\usepackage[utf8]{inputenc}
\usepackage{booktabs}
\usepackage{amsmath}
\usepackage{natbib}

\title{My Document}
\author{Your Name}
\date{\today}

\begin{document}

\maketitle

\section{Introduction}
This is a sample document demonstrating various \LaTeX{} features including mathematical equations, tables, and formatting. You can cite papers using the citation search feature, and they will be automatically added to your \texttt{cite.bib} file.

\section{Mathematical Expressions}
Here are some example equations:

\subsection{Quadratic Formula}
The quadratic formula is given by:
\[x = \frac{-b \pm \sqrt{b^2 - 4ac}}{2a}\]

\subsection{Matrix Example}
A sample matrix:
\[
\begin{pmatrix}
a_{11} & a_{12} & a_{13} \\
a_{21} & a_{22} & a_{23} \\
a_{31} & a_{32} & a_{33}
\end{pmatrix}
\]

\section{Tables}
Below is a sample table showing data:

\begin{table}[h]
\centering
\begin{tabular}{@{}lrrr@{}}
\toprule
\textbf{Category} & \textbf{Value 1} & \textbf{Value 2} & \textbf{Total} \\
\midrule
Item A & 10.5 & 20.3 & 30.8 \\
Item B & 15.2 & 25.1 & 40.3 \\
Item C & 22.1 & 12.7 & 34.8 \\
\bottomrule
\end{tabular}
\caption{Sample Data Table}
\label{tab:sample}
\end{table}

\section{Lists}
\subsection{Bullet Points}
\begin{itemize}
  \item First item
  \item Second item
  \item Third item with sub-items:
    \begin{itemize}
      \item Sub-item 1
      \item Sub-item 2
    \end{itemize}
\end{itemize}

\subsection{Numbered List}
\begin{enumerate}
  \item First step
  \item Second step
  \item Third step
\end{enumerate}

\section{Citations and References}
When you use the citation search feature, references will be automatically managed. For example, you might cite a paper like this: \cite{example2024}. All citations will be automatically added to your \texttt{cite.bib} file.

\section{Conclusion}
This template demonstrates basic \LaTeX{} features including mathematical formulas, tables, and lists. The bibliography system is now set up to work seamlessly with Bibby AI's citation management. Feel free to modify it for your needs!

\bibliographystyle{plain}
\bibliography{cite}

\end{document}\documentclass{article}
\usepackage[utf8]{inputenc}
\usepackage{booktabs}
\usepackage{amsmath}
\usepackage{natbib}

\title{My Document}
\author{Your Name}
\date{\today}

\begin{document}

\maketitle

\section{Introduction}
This is a sample document demonstrating various \LaTeX{} features including mathematical equations, tables, and formatting. You can cite papers using the citation search feature, and they will be automatically added to your \texttt{cite.bib} file.

\section{Mathematical Expressions}
Here are some example equations:

\subsection{Quadratic Formula}
The quadratic formula is given by:
\[x = \frac{-b \pm \sqrt{b^2 - 4ac}}{2a}\]

\subsection{Matrix Example}
A sample matrix:
\[
\begin{pmatrix}
a_{11} & a_{12} & a_{13} \\
a_{21} & a_{22} & a_{23} \\
a_{31} & a_{32} & a_{33}
\end{pmatrix}
\]

\section{Tables}
Below is a sample table showing data:

\begin{table}[h]
\centering
\begin{tabular}{@{}lrrr@{}}
\toprule
\textbf{Category} & \textbf{Value 1} & \textbf{Value 2} & \textbf{Total} \\
\midrule
Item A & 10.5 & 20.3 & 30.8 \\
Item B & 15.2 & 25.1 & 40.3 \\
Item C & 22.1 & 12.7 & 34.8 \\
\bottomrule
\end{tabular}
\caption{Sample Data Table}
\label{tab:sample}
\end{table}

\section{Lists}
\subsection{Bullet Points}
\begin{itemize}
  \item First item
  \item Second item
  \item Third item with sub-items:
    \begin{itemize}
      \item Sub-item 1
      \item Sub-item 2
    \end{itemize}
\end{itemize}

\subsection{Numbered List}
\begin{enumerate}
  \item First step
  \item Second step
  \item Third step
\end{enumerate}

\section{Citations and References}
When you use the citation search feature, references will be automatically managed. For example, you might cite a paper like this: \cite{example2024}. All citations will be automatically added to your \texttt{cite.bib} file.

\section{Conclusion}
This template demonstrates basic \LaTeX{} features including mathematical formulas, tables, and lists. The bibliography system is now set up to work seamlessly with Bibby AI's citation management. Feel free to modify it for your needs!

\bibliographystyle{plain}
\bibliography{cite}

\end{document}\documentclass{article}
\usepackage[utf8]{inputenc}
\usepackage{booktabs}
\usepackage{amsmath}
\usepackage{natbib}

\title{My Document}
\author{Your Name}
\date{\today}

\begin{document}

\maketitle

\section{Introduction}
This is a sample document demonstrating various \LaTeX{} features including mathematical equations, tables, and formatting. You can cite papers using the citation search feature, and they will be automatically added to your \texttt{cite.bib} file.

\section{Mathematical Expressions}
Here are some example equations:

\subsection{Quadratic Formula}
The quadratic formula is given by:
\[x = \frac{-b \pm \sqrt{b^2 - 4ac}}{2a}\]

\subsection{Matrix Example}
A sample matrix:
\[
\begin{pmatrix}
a_{11} & a_{12} & a_{13} \\
a_{21} & a_{22} & a_{23} \\
a_{31} & a_{32} & a_{33}
\end{pmatrix}
\]

\section{Tables}
Below is a sample table showing data:

\begin{table}[h]
\centering
\begin{tabular}{@{}lrrr@{}}
\toprule
\textbf{Category} & \textbf{Value 1} & \textbf{Value 2} & \textbf{Total} \\
\midrule
Item A & 10.5 & 20.3 & 30.8 \\
Item B & 15.2 & 25.1 & 40.3 \\
Item C & 22.1 & 12.7 & 34.8 \\
\bottomrule
\end{tabular}
\caption{Sample Data Table}
\label{tab:sample}
\end{table}

\section{Lists}
\subsection{Bullet Points}
\begin{itemize}
  \item First item
  \item Second item
  \item Third item with sub-items:
    \begin{itemize}
      \item Sub-item 1
      \item Sub-item 2
    \end{itemize}
\end{itemize}

\subsection{Numbered List}
\begin{enumerate}
  \item First step
  \item Second step
  \item Third step
\end{enumerate}

\section{Citations and References}
When you use the citation search feature, references will be automatically managed. For example, you might cite a paper like this: \cite{example2024}. All citations will be automatically added to your \texttt{cite.bib} file.

\section{Conclusion}
This template demonstrates basic \LaTeX{} features including mathematical formulas, tables, and lists. The bibliography system is now set up to work seamlessly with Bibby AI's citation management. Feel free to modify it for your needs!

\bibliographystyle{plain}
\bibliography{cite}

\end{document}\documentclass{article}
\usepackage[utf8]{inputenc}
\usepackage{booktabs}
\usepackage{amsmath}
\usepackage{natbib}

\title{My Document}
\author{Your Name}
\date{\today}

\begin{document}

\maketitle

\section{Introduction}
This is a sample document demonstrating various \LaTeX{} features including mathematical equations, tables, and formatting. You can cite papers using the citation search feature, and they will be automatically added to your \texttt{cite.bib} file.

\section{Mathematical Expressions}
Here are some example equations:

\subsection{Quadratic Formula}
The quadratic formula is given by:
\[x = \frac{-b \pm \sqrt{b^2 - 4ac}}{2a}\]

\subsection{Matrix Example}
A sample matrix:
\[
\begin{pmatrix}
a_{11} & a_{12} & a_{13} \\
a_{21} & a_{22} & a_{23} \\
a_{31} & a_{32} & a_{33}
\end{pmatrix}
\]

\section{Tables}
Below is a sample table showing data:

\begin{table}[h]
\centering
\begin{tabular}{@{}lrrr@{}}
\toprule
\textbf{Category} & \textbf{Value 1} & \textbf{Value 2} & \textbf{Total} \\
\midrule
Item A & 10.5 & 20.3 & 30.8 \\
Item B & 15.2 & 25.1 & 40.3 \\
Item C & 22.1 & 12.7 & 34.8 \\
\bottomrule
\end{tabular}
\caption{Sample Data Table}
\label{tab:sample}
\end{table}

\section{Lists}
\subsection{Bullet Points}
\begin{itemize}
  \item First item
  \item Second item
  \item Third item with sub-items:
    \begin{itemize}
      \item Sub-item 1
      \item Sub-item 2
    \end{itemize}
\end{itemize}

\subsection{Numbered List}
\begin{enumerate}
  \item First step
  \item Second step
  \item Third step
\end{enumerate}

\section{Citations and References}
When you use the citation search feature, references will be automatically managed. For example, you might cite a paper like this: \cite{example2024}. All citations will be automatically added to your \texttt{cite.bib} file.

\section{Conclusion}
This template demonstrates basic \LaTeX{} features including mathematical formulas, tables, and lists. The bibliography system is now set up to work seamlessly with Bibby AI's citation management. Feel free to modify it for your needs!

\bibliographystyle{plain}
\bibliography{cite}

\end{document}\documentclass{article}
\usepackage[utf8]{inputenc}
\usepackage{booktabs}
\usepackage{amsmath}
\usepackage{natbib}

\title{My Document}
\author{Your Name}
\date{\today}

\begin{document}

\maketitle

\section{Introduction}
This is a sample document demonstrating various \LaTeX{} features including mathematical equations, tables, and formatting. You can cite papers using the citation search feature, and they will be automatically added to your \texttt{cite.bib} file.

\section{Mathematical Expressions}
Here are some example equations:

\subsection{Quadratic Formula}
The quadratic formula is given by:
\[x = \frac{-b \pm \sqrt{b^2 - 4ac}}{2a}\]

\subsection{Matrix Example}
A sample matrix:
\[
\begin{pmatrix}
a_{11} & a_{12} & a_{13} \\
a_{21} & a_{22} & a_{23} \\
a_{31} & a_{32} & a_{33}
\end{pmatrix}
\]

\section{Tables}
Below is a sample table showing data:

\begin{table}[h]
\centering
\begin{tabular}{@{}lrrr@{}}
\toprule
\textbf{Category} & \textbf{Value 1} & \textbf{Value 2} & \textbf{Total} \\
\midrule
Item A & 10.5 & 20.3 & 30.8 \\
Item B & 15.2 & 25.1 & 40.3 \\
Item C & 22.1 & 12.7 & 34.8 \\
\bottomrule
\end{tabular}
\caption{Sample Data Table}
\label{tab:sample}
\end{table}

\section{Lists}
\subsection{Bullet Points}
\begin{itemize}
  \item First item
  \item Second item
  \item Third item with sub-items:
    \begin{itemize}
      \item Sub-item 1
      \item Sub-item 2
    \end{itemize}
\end{itemize}

\subsection{Numbered List}
\begin{enumerate}
  \item First step
  \item Second step
  \item Third step
\end{enumerate}

\section{Citations and References}
When you use the citation search feature, references will be automatically managed. For example, you might cite a paper like this: \cite{example2024}. All citations will be automatically added to your \texttt{cite.bib} file.

\section{Conclusion}
This template demonstrates basic \LaTeX{} features including mathematical formulas, tables, and lists. The bibliography system is now set up to work seamlessly with Bibby AI's citation management. Feel free to modify it for your needs!

\bibliographystyle{plain}
\bibliography{cite}

\end{document}\documentclass{article}
\usepackage[utf8]{inputenc}
\usepackage{booktabs}
\usepackage{amsmath}
\usepackage{natbib}

\title{My Document}
\author{Your Name}
\date{\today}

\begin{document}

\maketitle

\section{Introduction}
This is a sample document demonstrating various \LaTeX{} features including mathematical equations, tables, and formatting. You can cite papers using the citation search feature, and they will be automatically added to your \texttt{cite.bib} file.

\section{Mathematical Expressions}
Here are some example equations:

\subsection{Quadratic Formula}
The quadratic formula is given by:
\[x = \frac{-b \pm \sqrt{b^2 - 4ac}}{2a}\]

\subsection{Matrix Example}
A sample matrix:
\[
\begin{pmatrix}
a_{11} & a_{12} & a_{13} \\
a_{21} & a_{22} & a_{23} \\
a_{31} & a_{32} & a_{33}
\end{pmatrix}
\]

\section{Tables}
Below is a sample table showing data:

\begin{table}[h]
\centering
\begin{tabular}{@{}lrrr@{}}
\toprule
\textbf{Category} & \textbf{Value 1} & \textbf{Value 2} & \textbf{Total} \\
\midrule
Item A & 10.5 & 20.3 & 30.8 \\
Item B & 15.2 & 25.1 & 40.3 \\
Item C & 22.1 & 12.7 & 34.8 \\
\bottomrule
\end{tabular}
\caption{Sample Data Table}
\label{tab:sample}
\end{table}

\section{Lists}
\subsection{Bullet Points}
\begin{itemize}
  \item First item
  \item Second item
  \item Third item with sub-items:
    \begin{itemize}
      \item Sub-item 1
      \item Sub-item 2
    \end{itemize}
\end{itemize}

\subsection{Numbered List}
\begin{enumerate}
  \item First step
  \item Second step
  \item Third step
\end{enumerate}

\section{Citations and References}
When you use the citation search feature, references will be automatically managed. For example, you might cite a paper like this: \cite{example2024}. All citations will be automatically added to your \texttt{cite.bib} file.

\section{Conclusion}
This template demonstrates basic \LaTeX{} features including mathematical formulas, tables, and lists. The bibliography system is now set up to work seamlessly with Bibby AI's citation management. Feel free to modify it for your needs!

\bibliographystyle{plain}
\bibliography{cite}

\end{document}\documentclass{article}
\usepackage[utf8]{inputenc}
\usepackage{booktabs}
\usepackage{amsmath}
\usepackage{natbib}

\title{My Document}
\author{Your Name}
\date{\today}

\begin{document}

\maketitle

\section{Introduction}
This is a sample document demonstrating various \LaTeX{} features including mathematical equations, tables, and formatting. You can cite papers using the citation search feature, and they will be automatically added to your \texttt{cite.bib} file.

\section{Mathematical Expressions}
Here are some example equations:

\subsection{Quadratic Formula}
The quadratic formula is given by:
\[x = \frac{-b \pm \sqrt{b^2 - 4ac}}{2a}\]

\subsection{Matrix Example}
A sample matrix:
\[
\begin{pmatrix}
a_{11} & a_{12} & a_{13} \\
a_{21} & a_{22} & a_{23} \\
a_{31} & a_{32} & a_{33}
\end{pmatrix}
\]

\section{Tables}
Below is a sample table showing data:

\begin{table}[h]
\centering
\begin{tabular}{@{}lrrr@{}}
\toprule
\textbf{Category} & \textbf{Value 1} & \textbf{Value 2} & \textbf{Total} \\
\midrule
Item A & 10.5 & 20.3 & 30.8 \\
Item B & 15.2 & 25.1 & 40.3 \\
Item C & 22.1 & 12.7 & 34.8 \\
\bottomrule
\end{tabular}
\caption{Sample Data Table}
\label{tab:sample}
\end{table}

\section{Lists}
\subsection{Bullet Points}
\begin{itemize}
  \item First item
  \item Second item
  \item Third item with sub-items:
    \begin{itemize}
      \item Sub-item 1
      \item Sub-item 2
    \end{itemize}
\end{itemize}

\subsection{Numbered List}
\begin{enumerate}
  \item First step
  \item Second step
  \item Third step
\end{enumerate}

\section{Citations and References}
When you use the citation search feature, references will be automatically managed. For example, you might cite a paper like this: \cite{example2024}. All citations will be automatically added to your \texttt{cite.bib} file.

\section{Conclusion}
This template demonstrates basic \LaTeX{} features including mathematical formulas, tables, and lists. The bibliography system is now set up to work seamlessly with Bibby AI's citation management. Feel free to modify it for your needs!

\bibliographystyle{plain}
\bibliography{cite}

\end{document}\documentclass{article}
\usepackage[utf8]{inputenc}
\usepackage{booktabs}
\usepackage{amsmath}
\usepackage{natbib}

\title{My Document}
\author{Your Name}
\date{\today}

\begin{document}

\maketitle

\section{Introduction}
This is a sample document demonstrating various \LaTeX{} features including mathematical equations, tables, and formatting. You can cite papers using the citation search feature, and they will be automatically added to your \texttt{cite.bib} file.

\section{Mathematical Expressions}
Here are some example equations:

\subsection{Quadratic Formula}
The quadratic formula is given by:
\[x = \frac{-b \pm \sqrt{b^2 - 4ac}}{2a}\]

\subsection{Matrix Example}
A sample matrix:
\[
\begin{pmatrix}
a_{11} & a_{12} & a_{13} \\
a_{21} & a_{22} & a_{23} \\
a_{31} & a_{32} & a_{33}
\end{pmatrix}
\]

\section{Tables}
Below is a sample table showing data:

\begin{table}[h]
\centering
\begin{tabular}{@{}lrrr@{}}
\toprule
\textbf{Category} & \textbf{Value 1} & \textbf{Value 2} & \textbf{Total} \\
\midrule
Item A & 10.5 & 20.3 & 30.8 \\
Item B & 15.2 & 25.1 & 40.3 \\
Item C & 22.1 & 12.7 & 34.8 \\
\bottomrule
\end{tabular}
\caption{Sample Data Table}
\label{tab:sample}
\end{table}

\section{Lists}
\subsection{Bullet Points}
\begin{itemize}
  \item First item
  \item Second item
  \item Third item with sub-items:
    \begin{itemize}
      \item Sub-item 1
      \item Sub-item 2
    \end{itemize}
\end{itemize}

\subsection{Numbered List}
\begin{enumerate}
  \item First step
  \item Second step
  \item Third step
\end{enumerate}

\section{Citations and References}
When you use the citation search feature, references will be automatically managed. For example, you might cite a paper like this: \cite{example2024}. All citations will be automatically added to your \texttt{cite.bib} file.

\section{Conclusion}
This template demonstrates basic \LaTeX{} features including mathematical formulas, tables, and lists. The bibliography system is now set up to work seamlessly with Bibby AI's citation management. Feel free to modify it for your needs!

\bibliographystyle{plain}
\bibliography{cite}

\end{document}\documentclass{article}
\usepackage[utf8]{inputenc}
\usepackage{booktabs}
\usepackage{amsmath}
\usepackage{natbib}

\title{My Document}
\author{Your Name}
\date{\today}

\begin{document}

\maketitle

\section{Introduction}
This is a sample document demonstrating various \LaTeX{} features including mathematical equations, tables, and formatting. You can cite papers using the citation search feature, and they will be automatically added to your \texttt{cite.bib} file.

\section{Mathematical Expressions}
Here are some example equations:

\subsection{Quadratic Formula}
The quadratic formula is given by:
\[x = \frac{-b \pm \sqrt{b^2 - 4ac}}{2a}\]

\subsection{Matrix Example}
A sample matrix:
\[
\begin{pmatrix}
a_{11} & a_{12} & a_{13} \\
a_{21} & a_{22} & a_{23} \\
a_{31} & a_{32} & a_{33}
\end{pmatrix}
\]

\section{Tables}
Below is a sample table showing data:

\begin{table}[h]
\centering
\begin{tabular}{@{}lrrr@{}}
\toprule
\textbf{Category} & \textbf{Value 1} & \textbf{Value 2} & \textbf{Total} \\
\midrule
Item A & 10.5 & 20.3 & 30.8 \\
Item B & 15.2 & 25.1 & 40.3 \\
Item C & 22.1 & 12.7 & 34.8 \\
\bottomrule
\end{tabular}
\caption{Sample Data Table}
\label{tab:sample}
\end{table}

\section{Lists}
\subsection{Bullet Points}
\begin{itemize}
  \item First item
  \item Second item
  \item Third item with sub-items:
    \begin{itemize}
      \item Sub-item 1
      \item Sub-item 2
    \end{itemize}
\end{itemize}

\subsection{Numbered List}
\begin{enumerate}
  \item First step
  \item Second step
  \item Third step
\end{enumerate}

\section{Citations and References}
When you use the citation search feature, references will be automatically managed. For example, you might cite a paper like this: \cite{example2024}. All citations will be automatically added to your \texttt{cite.bib} file.

\section{Conclusion}
This template demonstrates basic \LaTeX{} features including mathematical formulas, tables, and lists. The bibliography system is now set up to work seamlessly with Bibby AI's citation management. Feel free to modify it for your needs!

\bibliographystyle{plain}
\bibliography{cite}

\end{document}\documentclass{article}
\usepackage[utf8]{inputenc}
\usepackage{booktabs}
\usepackage{amsmath}
\usepackage{natbib}

\title{My Document}
\author{Your Name}
\date{\today}

\begin{document}

\maketitle

\section{Introduction}
This is a sample document demonstrating various \LaTeX{} features including mathematical equations, tables, and formatting. You can cite papers using the citation search feature, and they will be automatically added to your \texttt{cite.bib} file.

\section{Mathematical Expressions}
Here are some example equations:

\subsection{Quadratic Formula}
The quadratic formula is given by:
\[x = \frac{-b \pm \sqrt{b^2 - 4ac}}{2a}\]

\subsection{Matrix Example}
A sample matrix:
\[
\begin{pmatrix}
a_{11} & a_{12} & a_{13} \\
a_{21} & a_{22} & a_{23} \\
a_{31} & a_{32} & a_{33}
\end{pmatrix}
\]

\section{Tables}
Below is a sample table showing data:

\begin{table}[h]
\centering
\begin{tabular}{@{}lrrr@{}}
\toprule
\textbf{Category} & \textbf{Value 1} & \textbf{Value 2} & \textbf{Total} \\
\midrule
Item A & 10.5 & 20.3 & 30.8 \\
Item B & 15.2 & 25.1 & 40.3 \\
Item C & 22.1 & 12.7 & 34.8 \\
\bottomrule
\end{tabular}
\caption{Sample Data Table}
\label{tab:sample}
\end{table}

\section{Lists}
\subsection{Bullet Points}
\begin{itemize}
  \item First item
  \item Second item
  \item Third item with sub-items:
    \begin{itemize}
      \item Sub-item 1
      \item Sub-item 2
    \end{itemize}
\end{itemize}

\subsection{Numbered List}
\begin{enumerate}
  \item First step
  \item Second step
  \item Third step
\end{enumerate}

\section{Citations and References}
When you use the citation search feature, references will be automatically managed. For example, you might cite a paper like this: \cite{example2024}. All citations will be automatically added to your \texttt{cite.bib} file.

\section{Conclusion}
This template demonstrates basic \LaTeX{} features including mathematical formulas, tables, and lists. The bibliography system is now set up to work seamlessly with Bibby AI's citation management. Feel free to modify it for your needs!

\bibliographystyle{plain}
\bibliography{cite}

\end{document}\documentclass{article}
\usepackage[utf8]{inputenc}
\usepackage{booktabs}
\usepackage{amsmath}
\usepackage{natbib}

\title{My Document}
\author{Your Name}
\date{\today}

\begin{document}

\maketitle

\section{Introduction}
This is a sample document demonstrating various \LaTeX{} features including mathematical equations, tables, and formatting. You can cite papers using the citation search feature, and they will be automatically added to your \texttt{cite.bib} file.

\section{Mathematical Expressions}
Here are some example equations:

\subsection{Quadratic Formula}
The quadratic formula is given by:
\[x = \frac{-b \pm \sqrt{b^2 - 4ac}}{2a}\]

\subsection{Matrix Example}
A sample matrix:
\[
\begin{pmatrix}
a_{11} & a_{12} & a_{13} \\
a_{21} & a_{22} & a_{23} \\
a_{31} & a_{32} & a_{33}
\end{pmatrix}
\]

\section{Tables}
Below is a sample table showing data:

\begin{table}[h]
\centering
\begin{tabular}{@{}lrrr@{}}
\toprule
\textbf{Category} & \textbf{Value 1} & \textbf{Value 2} & \textbf{Total} \\
\midrule
Item A & 10.5 & 20.3 & 30.8 \\
Item B & 15.2 & 25.1 & 40.3 \\
Item C & 22.1 & 12.7 & 34.8 \\
\bottomrule
\end{tabular}
\caption{Sample Data Table}
\label{tab:sample}
\end{table}

\section{Lists}
\subsection{Bullet Points}
\begin{itemize}
  \item First item
  \item Second item
  \item Third item with sub-items:
    \begin{itemize}
      \item Sub-item 1
      \item Sub-item 2
    \end{itemize}
\end{itemize}

\subsection{Numbered List}
\begin{enumerate}
  \item First step
  \item Second step
  \item Third step
\end{enumerate}

\section{Citations and References}
When you use the citation search feature, references will be automatically managed. For example, you might cite a paper like this: \cite{example2024}. All citations will be automatically added to your \texttt{cite.bib} file.

\section{Conclusion}
This template demonstrates basic \LaTeX{} features including mathematical formulas, tables, and lists. The bibliography system is now set up to work seamlessly with Bibby AI's citation management. Feel free to modify it for your needs!

\bibliographystyle{plain}
\bibliography{cite}

\end{document}\documentclass{article}
\usepackage[utf8]{inputenc}
\usepackage{booktabs}
\usepackage{amsmath}
\usepackage{natbib}

\title{My Document}
\author{Your Name}
\date{\today}

\begin{document}

\maketitle

\section{Introduction}
This is a sample document demonstrating various \LaTeX{} features including mathematical equations, tables, and formatting. You can cite papers using the citation search feature, and they will be automatically added to your \texttt{cite.bib} file.

\section{Mathematical Expressions}
Here are some example equations:

\subsection{Quadratic Formula}
The quadratic formula is given by:
\[x = \frac{-b \pm \sqrt{b^2 - 4ac}}{2a}\]

\subsection{Matrix Example}
A sample matrix:
\[
\begin{pmatrix}
a_{11} & a_{12} & a_{13} \\
a_{21} & a_{22} & a_{23} \\
a_{31} & a_{32} & a_{33}
\end{pmatrix}
\]

\section{Tables}
Below is a sample table showing data:

\begin{table}[h]
\centering
\begin{tabular}{@{}lrrr@{}}
\toprule
\textbf{Category} & \textbf{Value 1} & \textbf{Value 2} & \textbf{Total} \\
\midrule
Item A & 10.5 & 20.3 & 30.8 \\
Item B & 15.2 & 25.1 & 40.3 \\
Item C & 22.1 & 12.7 & 34.8 \\
\bottomrule
\end{tabular}
\caption{Sample Data Table}
\label{tab:sample}
\end{table}

\section{Lists}
\subsection{Bullet Points}
\begin{itemize}
  \item First item
  \item Second item
  \item Third item with sub-items:
    \begin{itemize}
      \item Sub-item 1
      \item Sub-item 2
    \end{itemize}
\end{itemize}

\subsection{Numbered List}
\begin{enumerate}
  \item First step
  \item Second step
  \item Third step
\end{enumerate}

\section{Citations and References}
When you use the citation search feature, references will be automatically managed. For example, you might cite a paper like this: \cite{example2024}. All citations will be automatically added to your \texttt{cite.bib} file.

\section{Conclusion}
This template demonstrates basic \LaTeX{} features including mathematical formulas, tables, and lists. The bibliography system is now set up to work seamlessly with Bibby AI's citation management. Feel free to modify it for your needs!

\bibliographystyle{plain}
\bibliography{cite}

\end{document}\documentclass{article}
\usepackage[utf8]{inputenc}
\usepackage{booktabs}
\usepackage{amsmath}
\usepackage{natbib}

\title{My Document}
\author{Your Name}
\date{\today}

\begin{document}

\maketitle

\section{Introduction}
This is a sample document demonstrating various \LaTeX{} features including mathematical equations, tables, and formatting. You can cite papers using the citation search feature, and they will be automatically added to your \texttt{cite.bib} file.

\section{Mathematical Expressions}
Here are some example equations:

\subsection{Quadratic Formula}
The quadratic formula is given by:
\[x = \frac{-b \pm \sqrt{b^2 - 4ac}}{2a}\]

\subsection{Matrix Example}
A sample matrix:
\[
\begin{pmatrix}
a_{11} & a_{12} & a_{13} \\
a_{21} & a_{22} & a_{23} \\
a_{31} & a_{32} & a_{33}
\end{pmatrix}
\]

\section{Tables}
Below is a sample table showing data:

\begin{table}[h]
\centering
\begin{tabular}{@{}lrrr@{}}
\toprule
\textbf{Category} & \textbf{Value 1} & \textbf{Value 2} & \textbf{Total} \\
\midrule
Item A & 10.5 & 20.3 & 30.8 \\
Item B & 15.2 & 25.1 & 40.3 \\
Item C & 22.1 & 12.7 & 34.8 \\
\bottomrule
\end{tabular}
\caption{Sample Data Table}
\label{tab:sample}
\end{table}

\section{Lists}
\subsection{Bullet Points}
\begin{itemize}
  \item First item
  \item Second item
  \item Third item with sub-items:
    \begin{itemize}
      \item Sub-item 1
      \item Sub-item 2
    \end{itemize}
\end{itemize}

\subsection{Numbered List}
\begin{enumerate}
  \item First step
  \item Second step
  \item Third step
\end{enumerate}

\section{Citations and References}
When you use the citation search feature, references will be automatically managed. For example, you might cite a paper like this: \cite{example2024}. All citations will be automatically added to your \texttt{cite.bib} file.

\section{Conclusion}
This template demonstrates basic \LaTeX{} features including mathematical formulas, tables, and lists. The bibliography system is now set up to work seamlessly with Bibby AI's citation management. Feel free to modify it for your needs!

\bibliographystyle{plain}
\bibliography{cite}

\end{document}\documentclass{article}
\usepackage[utf8]{inputenc}
\usepackage{booktabs}
\usepackage{amsmath}
\usepackage{natbib}

\title{My Document}
\author{Your Name}
\date{\today}

\begin{document}

\maketitle

\section{Introduction}
This is a sample document demonstrating various \LaTeX{} features including mathematical equations, tables, and formatting. You can cite papers using the citation search feature, and they will be automatically added to your \texttt{cite.bib} file.

\section{Mathematical Expressions}
Here are some example equations:

\subsection{Quadratic Formula}
The quadratic formula is given by:
\[x = \frac{-b \pm \sqrt{b^2 - 4ac}}{2a}\]

\subsection{Matrix Example}
A sample matrix:
\[
\begin{pmatrix}
a_{11} & a_{12} & a_{13} \\
a_{21} & a_{22} & a_{23} \\
a_{31} & a_{32} & a_{33}
\end{pmatrix}
\]

\section{Tables}
Below is a sample table showing data:

\begin{table}[h]
\centering
\begin{tabular}{@{}lrrr@{}}
\toprule
\textbf{Category} & \textbf{Value 1} & \textbf{Value 2} & \textbf{Total} \\
\midrule
Item A & 10.5 & 20.3 & 30.8 \\
Item B & 15.2 & 25.1 & 40.3 \\
Item C & 22.1 & 12.7 & 34.8 \\
\bottomrule
\end{tabular}
\caption{Sample Data Table}
\label{tab:sample}
\end{table}

\section{Lists}
\subsection{Bullet Points}
\begin{itemize}
  \item First item
  \item Second item
  \item Third item with sub-items:
    \begin{itemize}
      \item Sub-item 1
      \item Sub-item 2
    \end{itemize}
\end{itemize}

\subsection{Numbered List}
\begin{enumerate}
  \item First step
  \item Second step
  \item Third step
\end{enumerate}

\section{Citations and References}
When you use the citation search feature, references will be automatically managed. For example, you might cite a paper like this: \cite{example2024}. All citations will be automatically added to your \texttt{cite.bib} file.

\section{Conclusion}
This template demonstrates basic \LaTeX{} features including mathematical formulas, tables, and lists. The bibliography system is now set up to work seamlessly with Bibby AI's citation management. Feel free to modify it for your needs!

\bibliographystyle{plain}
\bibliography{cite}

\end{document}\documentclass{article}
\usepackage[utf8]{inputenc}
\usepackage{booktabs}
\usepackage{amsmath}
\usepackage{natbib}

\title{My Document}
\author{Your Name}
\date{\today}

\begin{document}

\maketitle

\section{Introduction}
This is a sample document demonstrating various \LaTeX{} features including mathematical equations, tables, and formatting. You can cite papers using the citation search feature, and they will be automatically added to your \texttt{cite.bib} file.

\section{Mathematical Expressions}
Here are some example equations:

\subsection{Quadratic Formula}
The quadratic formula is given by:
\[x = \frac{-b \pm \sqrt{b^2 - 4ac}}{2a}\]

\subsection{Matrix Example}
A sample matrix:
\[
\begin{pmatrix}
a_{11} & a_{12} & a_{13} \\
a_{21} & a_{22} & a_{23} \\
a_{31} & a_{32} & a_{33}
\end{pmatrix}
\]

\section{Tables}
Below is a sample table showing data:

\begin{table}[h]
\centering
\begin{tabular}{@{}lrrr@{}}
\toprule
\textbf{Category} & \textbf{Value 1} & \textbf{Value 2} & \textbf{Total} \\
\midrule
Item A & 10.5 & 20.3 & 30.8 \\
Item B & 15.2 & 25.1 & 40.3 \\
Item C & 22.1 & 12.7 & 34.8 \\
\bottomrule
\end{tabular}
\caption{Sample Data Table}
\label{tab:sample}
\end{table}

\section{Lists}
\subsection{Bullet Points}
\begin{itemize}
  \item First item
  \item Second item
  \item Third item with sub-items:
    \begin{itemize}
      \item Sub-item 1
      \item Sub-item 2
    \end{itemize}
\end{itemize}

\subsection{Numbered List}
\begin{enumerate}
  \item First step
  \item Second step
  \item Third step
\end{enumerate}

\section{Citations and References}
When you use the citation search feature, references will be automatically managed. For example, you might cite a paper like this: \cite{example2024}. All citations will be automatically added to your \texttt{cite.bib} file.

\section{Conclusion}
This template demonstrates basic \LaTeX{} features including mathematical formulas, tables, and lists. The bibliography system is now set up to work seamlessly with Bibby AI's citation management. Feel free to modify it for your needs!

\bibliographystyle{plain}
\bibliography{cite}

\end{document}\documentclass{article}
\usepackage[utf8]{inputenc}
\usepackage{booktabs}
\usepackage{amsmath}
\usepackage{natbib}

\title{My Document}
\author{Your Name}
\date{\today}

\begin{document}

\maketitle

\section{Introduction}
This is a sample document demonstrating various \LaTeX{} features including mathematical equations, tables, and formatting. You can cite papers using the citation search feature, and they will be automatically added to your \texttt{cite.bib} file.

\section{Mathematical Expressions}
Here are some example equations:

\subsection{Quadratic Formula}
The quadratic formula is given by:
\[x = \frac{-b \pm \sqrt{b^2 - 4ac}}{2a}\]

\subsection{Matrix Example}
A sample matrix:
\[
\begin{pmatrix}
a_{11} & a_{12} & a_{13} \\
a_{21} & a_{22} & a_{23} \\
a_{31} & a_{32} & a_{33}
\end{pmatrix}
\]

\section{Tables}
Below is a sample table showing data:

\begin{table}[h]
\centering
\begin{tabular}{@{}lrrr@{}}
\toprule
\textbf{Category} & \textbf{Value 1} & \textbf{Value 2} & \textbf{Total} \\
\midrule
Item A & 10.5 & 20.3 & 30.8 \\
Item B & 15.2 & 25.1 & 40.3 \\
Item C & 22.1 & 12.7 & 34.8 \\
\bottomrule
\end{tabular}
\caption{Sample Data Table}
\label{tab:sample}
\end{table}

\section{Lists}
\subsection{Bullet Points}
\begin{itemize}
  \item First item
  \item Second item
  \item Third item with sub-items:
    \begin{itemize}
      \item Sub-item 1
      \item Sub-item 2
    \end{itemize}
\end{itemize}

\subsection{Numbered List}
\begin{enumerate}
  \item First step
  \item Second step
  \item Third step
\end{enumerate}

\section{Citations and References}
When you use the citation search feature, references will be automatically managed. For example, you might cite a paper like this: \cite{example2024}. All citations will be automatically added to your \texttt{cite.bib} file.

\section{Conclusion}
This template demonstrates basic \LaTeX{} features including mathematical formulas, tables, and lists. The bibliography system is now set up to work seamlessly with Bibby AI's citation management. Feel free to modify it for your needs!

\bibliographystyle{plain}
\bibliography{cite}

\end{document}\documentclass{article}
\usepackage[utf8]{inputenc}
\usepackage{booktabs}
\usepackage{amsmath}
\usepackage{natbib}

\title{My Document}
\author{Your Name}
\date{\today}

\begin{document}

\maketitle

\section{Introduction}
This is a sample document demonstrating various \LaTeX{} features including mathematical equations, tables, and formatting. You can cite papers using the citation search feature, and they will be automatically added to your \texttt{cite.bib} file.

\section{Mathematical Expressions}
Here are some example equations:

\subsection{Quadratic Formula}
The quadratic formula is given by:
\[x = \frac{-b \pm \sqrt{b^2 - 4ac}}{2a}\]

\subsection{Matrix Example}
A sample matrix:
\[
\begin{pmatrix}
a_{11} & a_{12} & a_{13} \\
a_{21} & a_{22} & a_{23} \\
a_{31} & a_{32} & a_{33}
\end{pmatrix}
\]

\section{Tables}
Below is a sample table showing data:

\begin{table}[h]
\centering
\begin{tabular}{@{}lrrr@{}}
\toprule
\textbf{Category} & \textbf{Value 1} & \textbf{Value 2} & \textbf{Total} \\
\midrule
Item A & 10.5 & 20.3 & 30.8 \\
Item B & 15.2 & 25.1 & 40.3 \\
Item C & 22.1 & 12.7 & 34.8 \\
\bottomrule
\end{tabular}
\caption{Sample Data Table}
\label{tab:sample}
\end{table}

\section{Lists}
\subsection{Bullet Points}
\begin{itemize}
  \item First item
  \item Second item
  \item Third item with sub-items:
    \begin{itemize}
      \item Sub-item 1
      \item Sub-item 2
    \end{itemize}
\end{itemize}

\subsection{Numbered List}
\begin{enumerate}
  \item First step
  \item Second step
  \item Third step
\end{enumerate}

\section{Citations and References}
When you use the citation search feature, references will be automatically managed. For example, you might cite a paper like this: \cite{example2024}. All citations will be automatically added to your \texttt{cite.bib} file.

\section{Conclusion}
This template demonstrates basic \LaTeX{} features including mathematical formulas, tables, and lists. The bibliography system is now set up to work seamlessly with Bibby AI's citation management. Feel free to modify it for your needs!

\bibliographystyle{plain}
\bibliography{cite}

\end{document}\documentclass{article}
\usepackage[utf8]{inputenc}
\usepackage{booktabs}
\usepackage{amsmath}
\usepackage{natbib}

\title{My Document}
\author{Your Name}
\date{\today}

\begin{document}

\maketitle

\section{Introduction}
This is a sample document demonstrating various \LaTeX{} features including mathematical equations, tables, and formatting. You can cite papers using the citation search feature, and they will be automatically added to your \texttt{cite.bib} file.

\section{Mathematical Expressions}
Here are some example equations:

\subsection{Quadratic Formula}
The quadratic formula is given by:
\[x = \frac{-b \pm \sqrt{b^2 - 4ac}}{2a}\]

\subsection{Matrix Example}
A sample matrix:
\[
\begin{pmatrix}
a_{11} & a_{12} & a_{13} \\
a_{21} & a_{22} & a_{23} \\
a_{31} & a_{32} & a_{33}
\end{pmatrix}
\]

\section{Tables}
Below is a sample table showing data:

\begin{table}[h]
\centering
\begin{tabular}{@{}lrrr@{}}
\toprule
\textbf{Category} & \textbf{Value 1} & \textbf{Value 2} & \textbf{Total} \\
\midrule
Item A & 10.5 & 20.3 & 30.8 \\
Item B & 15.2 & 25.1 & 40.3 \\
Item C & 22.1 & 12.7 & 34.8 \\
\bottomrule
\end{tabular}
\caption{Sample Data Table}
\label{tab:sample}
\end{table}

\section{Lists}
\subsection{Bullet Points}
\begin{itemize}
  \item First item
  \item Second item
  \item Third item with sub-items:
    \begin{itemize}
      \item Sub-item 1
      \item Sub-item 2
    \end{itemize}
\end{itemize}

\subsection{Numbered List}
\begin{enumerate}
  \item First step
  \item Second step
  \item Third step
\end{enumerate}

\section{Citations and References}
When you use the citation search feature, references will be automatically managed. For example, you might cite a paper like this: \cite{example2024}. All citations will be automatically added to your \texttt{cite.bib} file.

\section{Conclusion}
This template demonstrates basic \LaTeX{} features including mathematical formulas, tables, and lists. The bibliography system is now set up to work seamlessly with Bibby AI's citation management. Feel free to modify it for your needs!

\bibliographystyle{plain}
\bibliography{cite}

\end{document}\documentclass{article}
\usepackage[utf8]{inputenc}
\usepackage{booktabs}
\usepackage{amsmath}
\usepackage{natbib}

\title{My Document}
\author{Your Name}
\date{\today}

\begin{document}

\maketitle

\section{Introduction}
This is a sample document demonstrating various \LaTeX{} features including mathematical equations, tables, and formatting. You can cite papers using the citation search feature, and they will be automatically added to your \texttt{cite.bib} file.

\section{Mathematical Expressions}
Here are some example equations:

\subsection{Quadratic Formula}
The quadratic formula is given by:
\[x = \frac{-b \pm \sqrt{b^2 - 4ac}}{2a}\]

\subsection{Matrix Example}
A sample matrix:
\[
\begin{pmatrix}
a_{11} & a_{12} & a_{13} \\
a_{21} & a_{22} & a_{23} \\
a_{31} & a_{32} & a_{33}
\end{pmatrix}
\]

\section{Tables}
Below is a sample table showing data:

\begin{table}[h]
\centering
\begin{tabular}{@{}lrrr@{}}
\toprule
\textbf{Category} & \textbf{Value 1} & \textbf{Value 2} & \textbf{Total} \\
\midrule
Item A & 10.5 & 20.3 & 30.8 \\
Item B & 15.2 & 25.1 & 40.3 \\
Item C & 22.1 & 12.7 & 34.8 \\
\bottomrule
\end{tabular}
\caption{Sample Data Table}
\label{tab:sample}
\end{table}

\section{Lists}
\subsection{Bullet Points}
\begin{itemize}
  \item First item
  \item Second item
  \item Third item with sub-items:
    \begin{itemize}
      \item Sub-item 1
      \item Sub-item 2
    \end{itemize}
\end{itemize}

\subsection{Numbered List}
\begin{enumerate}
  \item First step
  \item Second step
  \item Third step
\end{enumerate}

\section{Citations and References}
When you use the citation search feature, references will be automatically managed. For example, you might cite a paper like this: \cite{example2024}. All citations will be automatically added to your \texttt{cite.bib} file.

\section{Conclusion}
This template demonstrates basic \LaTeX{} features including mathematical formulas, tables, and lists. The bibliography system is now set up to work seamlessly with Bibby AI's citation management. Feel free to modify it for your needs!

\bibliographystyle{plain}
\bibliography{cite}

\end{document}\documentclass{article}
\usepackage[utf8]{inputenc}
\usepackage{booktabs}
\usepackage{amsmath}
\usepackage{natbib}

\title{My Document}
\author{Your Name}
\date{\today}

\begin{document}

\maketitle

\section{Introduction}
This is a sample document demonstrating various \LaTeX{} features including mathematical equations, tables, and formatting. You can cite papers using the citation search feature, and they will be automatically added to your \texttt{cite.bib} file.

\section{Mathematical Expressions}
Here are some example equations:

\subsection{Quadratic Formula}
The quadratic formula is given by:
\[x = \frac{-b \pm \sqrt{b^2 - 4ac}}{2a}\]

\subsection{Matrix Example}
A sample matrix:
\[
\begin{pmatrix}
a_{11} & a_{12} & a_{13} \\
a_{21} & a_{22} & a_{23} \\
a_{31} & a_{32} & a_{33}
\end{pmatrix}
\]

\section{Tables}
Below is a sample table showing data:

\begin{table}[h]
\centering
\begin{tabular}{@{}lrrr@{}}
\toprule
\textbf{Category} & \textbf{Value 1} & \textbf{Value 2} & \textbf{Total} \\
\midrule
Item A & 10.5 & 20.3 & 30.8 \\
Item B & 15.2 & 25.1 & 40.3 \\
Item C & 22.1 & 12.7 & 34.8 \\
\bottomrule
\end{tabular}
\caption{Sample Data Table}
\label{tab:sample}
\end{table}

\section{Lists}
\subsection{Bullet Points}
\begin{itemize}
  \item First item
  \item Second item
  \item Third item with sub-items:
    \begin{itemize}
      \item Sub-item 1
      \item Sub-item 2
    \end{itemize}
\end{itemize}

\subsection{Numbered List}
\begin{enumerate}
  \item First step
  \item Second step
  \item Third step
\end{enumerate}

\section{Citations and References}
When you use the citation search feature, references will be automatically managed. For example, you might cite a paper like this: \cite{example2024}. All citations will be automatically added to your \texttt{cite.bib} file.

\section{Conclusion}
This template demonstrates basic \LaTeX{} features including mathematical formulas, tables, and lists. The bibliography system is now set up to work seamlessly with Bibby AI's citation management. Feel free to modify it for your needs!

\bibliographystyle{plain}
\bibliography{cite}

\end{document}\documentclass{article}
\usepackage[utf8]{inputenc}
\usepackage{booktabs}
\usepackage{amsmath}
\usepackage{natbib}

\title{My Document}
\author{Your Name}
\date{\today}

\begin{document}

\maketitle

\section{Introduction}
This is a sample document demonstrating various \LaTeX{} features including mathematical equations, tables, and formatting. You can cite papers using the citation search feature, and they will be automatically added to your \texttt{cite.bib} file.

\section{Mathematical Expressions}
Here are some example equations:

\subsection{Quadratic Formula}
The quadratic formula is given by:
\[x = \frac{-b \pm \sqrt{b^2 - 4ac}}{2a}\]

\subsection{Matrix Example}
A sample matrix:
\[
\begin{pmatrix}
a_{11} & a_{12} & a_{13} \\
a_{21} & a_{22} & a_{23} \\
a_{31} & a_{32} & a_{33}
\end{pmatrix}
\]

\section{Tables}
Below is a sample table showing data:

\begin{table}[h]
\centering
\begin{tabular}{@{}lrrr@{}}
\toprule
\textbf{Category} & \textbf{Value 1} & \textbf{Value 2} & \textbf{Total} \\
\midrule
Item A & 10.5 & 20.3 & 30.8 \\
Item B & 15.2 & 25.1 & 40.3 \\
Item C & 22.1 & 12.7 & 34.8 \\
\bottomrule
\end{tabular}
\caption{Sample Data Table}
\label{tab:sample}
\end{table}

\section{Lists}
\subsection{Bullet Points}
\begin{itemize}
  \item First item
  \item Second item
  \item Third item with sub-items:
    \begin{itemize}
      \item Sub-item 1
      \item Sub-item 2
    \end{itemize}
\end{itemize}

\subsection{Numbered List}
\begin{enumerate}
  \item First step
  \item Second step
  \item Third step
\end{enumerate}

\section{Citations and References}
When you use the citation search feature, references will be automatically managed. For example, you might cite a paper like this: \cite{example2024}. All citations will be automatically added to your \texttt{cite.bib} file.

\section{Conclusion}
This template demonstrates basic \LaTeX{} features including mathematical formulas, tables, and lists. The bibliography system is now set up to work seamlessly with Bibby AI's citation management. Feel free to modify it for your needs!

\bibliographystyle{plain}
\bibliography{cite}

\end{document}\documentclass{article}
\usepackage[utf8]{inputenc}
\usepackage{booktabs}
\usepackage{amsmath}
\usepackage{natbib}

\title{My Document}
\author{Your Name}
\date{\today}

\begin{document}

\maketitle

\section{Introduction}
This is a sample document demonstrating various \LaTeX{} features including mathematical equations, tables, and formatting. You can cite papers using the citation search feature, and they will be automatically added to your \texttt{cite.bib} file.

\section{Mathematical Expressions}
Here are some example equations:

\subsection{Quadratic Formula}
The quadratic formula is given by:
\[x = \frac{-b \pm \sqrt{b^2 - 4ac}}{2a}\]

\subsection{Matrix Example}
A sample matrix:
\[
\begin{pmatrix}
a_{11} & a_{12} & a_{13} \\
a_{21} & a_{22} & a_{23} \\
a_{31} & a_{32} & a_{33}
\end{pmatrix}
\]

\section{Tables}
Below is a sample table showing data:

\begin{table}[h]
\centering
\begin{tabular}{@{}lrrr@{}}
\toprule
\textbf{Category} & \textbf{Value 1} & \textbf{Value 2} & \textbf{Total} \\
\midrule
Item A & 10.5 & 20.3 & 30.8 \\
Item B & 15.2 & 25.1 & 40.3 \\
Item C & 22.1 & 12.7 & 34.8 \\
\bottomrule
\end{tabular}
\caption{Sample Data Table}
\label{tab:sample}
\end{table}

\section{Lists}
\subsection{Bullet Points}
\begin{itemize}
  \item First item
  \item Second item
  \item Third item with sub-items:
    \begin{itemize}
      \item Sub-item 1
      \item Sub-item 2
    \end{itemize}
\end{itemize}

\subsection{Numbered List}
\begin{enumerate}
  \item First step
  \item Second step
  \item Third step
\end{enumerate}

\section{Citations and References}
When you use the citation search feature, references will be automatically managed. For example, you might cite a paper like this: \cite{example2024}. All citations will be automatically added to your \texttt{cite.bib} file.

\section{Conclusion}
This template demonstrates basic \LaTeX{} features including mathematical formulas, tables, and lists. The bibliography system is now set up to work seamlessly with Bibby AI's citation management. Feel free to modify it for your needs!

\bibliographystyle{plain}
\bibliography{cite}

\end{document}\documentclass{article}
\usepackage[utf8]{inputenc}
\usepackage{booktabs}
\usepackage{amsmath}
\usepackage{natbib}

\title{My Document}
\author{Your Name}
\date{\today}

\begin{document}

\maketitle

\section{Introduction}
This is a sample document demonstrating various \LaTeX{} features including mathematical equations, tables, and formatting. You can cite papers using the citation search feature, and they will be automatically added to your \texttt{cite.bib} file.

\section{Mathematical Expressions}
Here are some example equations:

\subsection{Quadratic Formula}
The quadratic formula is given by:
\[x = \frac{-b \pm \sqrt{b^2 - 4ac}}{2a}\]

\subsection{Matrix Example}
A sample matrix:
\[
\begin{pmatrix}
a_{11} & a_{12} & a_{13} \\
a_{21} & a_{22} & a_{23} \\
a_{31} & a_{32} & a_{33}
\end{pmatrix}
\]

\section{Tables}
Below is a sample table showing data:

\begin{table}[h]
\centering
\begin{tabular}{@{}lrrr@{}}
\toprule
\textbf{Category} & \textbf{Value 1} & \textbf{Value 2} & \textbf{Total} \\
\midrule
Item A & 10.5 & 20.3 & 30.8 \\
Item B & 15.2 & 25.1 & 40.3 \\
Item C & 22.1 & 12.7 & 34.8 \\
\bottomrule
\end{tabular}
\caption{Sample Data Table}
\label{tab:sample}
\end{table}

\section{Lists}
\subsection{Bullet Points}
\begin{itemize}
  \item First item
  \item Second item
  \item Third item with sub-items:
    \begin{itemize}
      \item Sub-item 1
      \item Sub-item 2
    \end{itemize}
\end{itemize}

\subsection{Numbered List}
\begin{enumerate}
  \item First step
  \item Second step
  \item Third step
\end{enumerate}

\section{Citations and References}
When you use the citation search feature, references will be automatically managed. For example, you might cite a paper like this: \cite{example2024}. All citations will be automatically added to your \texttt{cite.bib} file.

\section{Conclusion}
This template demonstrates basic \LaTeX{} features including mathematical formulas, tables, and lists. The bibliography system is now set up to work seamlessly with Bibby AI's citation management. Feel free to modify it for your needs!

\bibliographystyle{plain}
\bibliography{cite}

\end{document}\documentclass{article}
\usepackage[utf8]{inputenc}
\usepackage{booktabs}
\usepackage{amsmath}
\usepackage{natbib}

\title{My Document}
\author{Your Name}
\date{\today}

\begin{document}

\maketitle

\section{Introduction}
This is a sample document demonstrating various \LaTeX{} features including mathematical equations, tables, and formatting. You can cite papers using the citation search feature, and they will be automatically added to your \texttt{cite.bib} file.

\section{Mathematical Expressions}
Here are some example equations:

\subsection{Quadratic Formula}
The quadratic formula is given by:
\[x = \frac{-b \pm \sqrt{b^2 - 4ac}}{2a}\]

\subsection{Matrix Example}
A sample matrix:
\[
\begin{pmatrix}
a_{11} & a_{12} & a_{13} \\
a_{21} & a_{22} & a_{23} \\
a_{31} & a_{32} & a_{33}
\end{pmatrix}
\]

\section{Tables}
Below is a sample table showing data:

\begin{table}[h]
\centering
\begin{tabular}{@{}lrrr@{}}
\toprule
\textbf{Category} & \textbf{Value 1} & \textbf{Value 2} & \textbf{Total} \\
\midrule
Item A & 10.5 & 20.3 & 30.8 \\
Item B & 15.2 & 25.1 & 40.3 \\
Item C & 22.1 & 12.7 & 34.8 \\
\bottomrule
\end{tabular}
\caption{Sample Data Table}
\label{tab:sample}
\end{table}

\section{Lists}
\subsection{Bullet Points}
\begin{itemize}
  \item First item
  \item Second item
  \item Third item with sub-items:
    \begin{itemize}
      \item Sub-item 1
      \item Sub-item 2
    \end{itemize}
\end{itemize}

\subsection{Numbered List}
\begin{enumerate}
  \item First step
  \item Second step
  \item Third step
\end{enumerate}

\section{Citations and References}
When you use the citation search feature, references will be automatically managed. For example, you might cite a paper like this: \cite{example2024}. All citations will be automatically added to your \texttt{cite.bib} file.

\section{Conclusion}
This template demonstrates basic \LaTeX{} features including mathematical formulas, tables, and lists. The bibliography system is now set up to work seamlessly with Bibby AI's citation management. Feel free to modify it for your needs!

\bibliographystyle{plain}
\bibliography{cite}

\end{document}\documentclass{article}
\usepackage[utf8]{inputenc}
\usepackage{booktabs}
\usepackage{amsmath}
\usepackage{natbib}

\title{My Document}
\author{Your Name}
\date{\today}

\begin{document}

\maketitle

\section{Introduction}
This is a sample document demonstrating various \LaTeX{} features including mathematical equations, tables, and formatting. You can cite papers using the citation search feature, and they will be automatically added to your \texttt{cite.bib} file.

\section{Mathematical Expressions}
Here are some example equations:

\subsection{Quadratic Formula}
The quadratic formula is given by:
\[x = \frac{-b \pm \sqrt{b^2 - 4ac}}{2a}\]

\subsection{Matrix Example}
A sample matrix:
\[
\begin{pmatrix}
a_{11} & a_{12} & a_{13} \\
a_{21} & a_{22} & a_{23} \\
a_{31} & a_{32} & a_{33}
\end{pmatrix}
\]

\section{Tables}
Below is a sample table showing data:

\begin{table}[h]
\centering
\begin{tabular}{@{}lrrr@{}}
\toprule
\textbf{Category} & \textbf{Value 1} & \textbf{Value 2} & \textbf{Total} \\
\midrule
Item A & 10.5 & 20.3 & 30.8 \\
Item B & 15.2 & 25.1 & 40.3 \\
Item C & 22.1 & 12.7 & 34.8 \\
\bottomrule
\end{tabular}
\caption{Sample Data Table}
\label{tab:sample}
\end{table}

\section{Lists}
\subsection{Bullet Points}
\begin{itemize}
  \item First item
  \item Second item
  \item Third item with sub-items:
    \begin{itemize}
      \item Sub-item 1
      \item Sub-item 2
    \end{itemize}
\end{itemize}

\subsection{Numbered List}
\begin{enumerate}
  \item First step
  \item Second step
  \item Third step
\end{enumerate}

\section{Citations and References}
When you use the citation search feature, references will be automatically managed. For example, you might cite a paper like this: \cite{example2024}. All citations will be automatically added to your \texttt{cite.bib} file.

\section{Conclusion}
This template demonstrates basic \LaTeX{} features including mathematical formulas, tables, and lists. The bibliography system is now set up to work seamlessly with Bibby AI's citation management. Feel free to modify it for your needs!

\bibliographystyle{plain}
\bibliography{cite}

\end{document}\documentclass{article}
\usepackage[utf8]{inputenc}
\usepackage{booktabs}
\usepackage{amsmath}
\usepackage{natbib}

\title{My Document}
\author{Your Name}
\date{\today}

\begin{document}

\maketitle

\section{Introduction}
This is a sample document demonstrating various \LaTeX{} features including mathematical equations, tables, and formatting. You can cite papers using the citation search feature, and they will be automatically added to your \texttt{cite.bib} file.

\section{Mathematical Expressions}
Here are some example equations:

\subsection{Quadratic Formula}
The quadratic formula is given by:
\[x = \frac{-b \pm \sqrt{b^2 - 4ac}}{2a}\]

\subsection{Matrix Example}
A sample matrix:
\[
\begin{pmatrix}
a_{11} & a_{12} & a_{13} \\
a_{21} & a_{22} & a_{23} \\
a_{31} & a_{32} & a_{33}
\end{pmatrix}
\]

\section{Tables}
Below is a sample table showing data:

\begin{table}[h]
\centering
\begin{tabular}{@{}lrrr@{}}
\toprule
\textbf{Category} & \textbf{Value 1} & \textbf{Value 2} & \textbf{Total} \\
\midrule
Item A & 10.5 & 20.3 & 30.8 \\
Item B & 15.2 & 25.1 & 40.3 \\
Item C & 22.1 & 12.7 & 34.8 \\
\bottomrule
\end{tabular}
\caption{Sample Data Table}
\label{tab:sample}
\end{table}

\section{Lists}
\subsection{Bullet Points}
\begin{itemize}
  \item First item
  \item Second item
  \item Third item with sub-items:
    \begin{itemize}
      \item Sub-item 1
      \item Sub-item 2
    \end{itemize}
\end{itemize}

\subsection{Numbered List}
\begin{enumerate}
  \item First step
  \item Second step
  \item Third step
\end{enumerate}

\section{Citations and References}
When you use the citation search feature, references will be automatically managed. For example, you might cite a paper like this: \cite{example2024}. All citations will be automatically added to your \texttt{cite.bib} file.

\section{Conclusion}
This template demonstrates basic \LaTeX{} features including mathematical formulas, tables, and lists. The bibliography system is now set up to work seamlessly with Bibby AI's citation management. Feel free to modify it for your needs!

\bibliographystyle{plain}
\bibliography{cite}

\end{document}\documentclass{article}
\usepackage[utf8]{inputenc}
\usepackage{booktabs}
\usepackage{amsmath}
\usepackage{natbib}

\title{My Document}
\author{Your Name}
\date{\today}

\begin{document}

\maketitle

\section{Introduction}
This is a sample document demonstrating various \LaTeX{} features including mathematical equations, tables, and formatting. You can cite papers using the citation search feature, and they will be automatically added to your \texttt{cite.bib} file.

\section{Mathematical Expressions}
Here are some example equations:

\subsection{Quadratic Formula}
The quadratic formula is given by:
\[x = \frac{-b \pm \sqrt{b^2 - 4ac}}{2a}\]

\subsection{Matrix Example}
A sample matrix:
\[
\begin{pmatrix}
a_{11} & a_{12} & a_{13} \\
a_{21} & a_{22} & a_{23} \\
a_{31} & a_{32} & a_{33}
\end{pmatrix}
\]

\section{Tables}
Below is a sample table showing data:

\begin{table}[h]
\centering
\begin{tabular}{@{}lrrr@{}}
\toprule
\textbf{Category} & \textbf{Value 1} & \textbf{Value 2} & \textbf{Total} \\
\midrule
Item A & 10.5 & 20.3 & 30.8 \\
Item B & 15.2 & 25.1 & 40.3 \\
Item C & 22.1 & 12.7 & 34.8 \\
\bottomrule
\end{tabular}
\caption{Sample Data Table}
\label{tab:sample}
\end{table}

\section{Lists}
\subsection{Bullet Points}
\begin{itemize}
  \item First item
  \item Second item
  \item Third item with sub-items:
    \begin{itemize}
      \item Sub-item 1
      \item Sub-item 2
    \end{itemize}
\end{itemize}

\subsection{Numbered List}
\begin{enumerate}
  \item First step
  \item Second step
  \item Third step
\end{enumerate}

\section{Citations and References}
When you use the citation search feature, references will be automatically managed. For example, you might cite a paper like this: \cite{example2024}. All citations will be automatically added to your \texttt{cite.bib} file.

\section{Conclusion}
This template demonstrates basic \LaTeX{} features including mathematical formulas, tables, and lists. The bibliography system is now set up to work seamlessly with Bibby AI's citation management. Feel free to modify it for your needs!

\bibliographystyle{plain}
\bibliography{cite}

\end{document}\documentclass{article}
\usepackage[utf8]{inputenc}
\usepackage{booktabs}
\usepackage{amsmath}
\usepackage{natbib}

\title{My Document}
\author{Your Name}
\date{\today}

\begin{document}

\maketitle

\section{Introduction}
This is a sample document demonstrating various \LaTeX{} features including mathematical equations, tables, and formatting. You can cite papers using the citation search feature, and they will be automatically added to your \texttt{cite.bib} file.

\section{Mathematical Expressions}
Here are some example equations:

\subsection{Quadratic Formula}
The quadratic formula is given by:
\[x = \frac{-b \pm \sqrt{b^2 - 4ac}}{2a}\]

\subsection{Matrix Example}
A sample matrix:
\[
\begin{pmatrix}
a_{11} & a_{12} & a_{13} \\
a_{21} & a_{22} & a_{23} \\
a_{31} & a_{32} & a_{33}
\end{pmatrix}
\]

\section{Tables}
Below is a sample table showing data:

\begin{table}[h]
\centering
\begin{tabular}{@{}lrrr@{}}
\toprule
\textbf{Category} & \textbf{Value 1} & \textbf{Value 2} & \textbf{Total} \\
\midrule
Item A & 10.5 & 20.3 & 30.8 \\
Item B & 15.2 & 25.1 & 40.3 \\
Item C & 22.1 & 12.7 & 34.8 \\
\bottomrule
\end{tabular}
\caption{Sample Data Table}
\label{tab:sample}
\end{table}

\section{Lists}
\subsection{Bullet Points}
\begin{itemize}
  \item First item
  \item Second item
  \item Third item with sub-items:
    \begin{itemize}
      \item Sub-item 1
      \item Sub-item 2
    \end{itemize}
\end{itemize}

\subsection{Numbered List}
\begin{enumerate}
  \item First step
  \item Second step
  \item Third step
\end{enumerate}

\section{Citations and References}
When you use the citation search feature, references will be automatically managed. For example, you might cite a paper like this: \cite{example2024}. All citations will be automatically added to your \texttt{cite.bib} file.

\section{Conclusion}
This template demonstrates basic \LaTeX{} features including mathematical formulas, tables, and lists. The bibliography system is now set up to work seamlessly with Bibby AI's citation management. Feel free to modify it for your needs!

\bibliographystyle{plain}
\bibliography{cite}

\end{document}\documentclass{article}
\usepackage[utf8]{inputenc}
\usepackage{booktabs}
\usepackage{amsmath}
\usepackage{natbib}

\title{My Document}
\author{Your Name}
\date{\today}

\begin{document}

\maketitle

\section{Introduction}
This is a sample document demonstrating various \LaTeX{} features including mathematical equations, tables, and formatting. You can cite papers using the citation search feature, and they will be automatically added to your \texttt{cite.bib} file.

\section{Mathematical Expressions}
Here are some example equations:

\subsection{Quadratic Formula}
The quadratic formula is given by:
\[x = \frac{-b \pm \sqrt{b^2 - 4ac}}{2a}\]

\subsection{Matrix Example}
A sample matrix:
\[
\begin{pmatrix}
a_{11} & a_{12} & a_{13} \\
a_{21} & a_{22} & a_{23} \\
a_{31} & a_{32} & a_{33}
\end{pmatrix}
\]

\section{Tables}
Below is a sample table showing data:

\begin{table}[h]
\centering
\begin{tabular}{@{}lrrr@{}}
\toprule
\textbf{Category} & \textbf{Value 1} & \textbf{Value 2} & \textbf{Total} \\
\midrule
Item A & 10.5 & 20.3 & 30.8 \\
Item B & 15.2 & 25.1 & 40.3 \\
Item C & 22.1 & 12.7 & 34.8 \\
\bottomrule
\end{tabular}
\caption{Sample Data Table}
\label{tab:sample}
\end{table}

\section{Lists}
\subsection{Bullet Points}
\begin{itemize}
  \item First item
  \item Second item
  \item Third item with sub-items:
    \begin{itemize}
      \item Sub-item 1
      \item Sub-item 2
    \end{itemize}
\end{itemize}

\subsection{Numbered List}
\begin{enumerate}
  \item First step
  \item Second step
  \item Third step
\end{enumerate}

\section{Citations and References}
When you use the citation search feature, references will be automatically managed. For example, you might cite a paper like this: \cite{example2024}. All citations will be automatically added to your \texttt{cite.bib} file.

\section{Conclusion}
This template demonstrates basic \LaTeX{} features including mathematical formulas, tables, and lists. The bibliography system is now set up to work seamlessly with Bibby AI's citation management. Feel free to modify it for your needs!

\bibliographystyle{plain}
\bibliography{cite}

\end{document}\documentclass{article}
\usepackage[utf8]{inputenc}
\usepackage{booktabs}
\usepackage{amsmath}
\usepackage{natbib}

\title{My Document}
\author{Your Name}
\date{\today}

\begin{document}

\maketitle

\section{Introduction}
This is a sample document demonstrating various \LaTeX{} features including mathematical equations, tables, and formatting. You can cite papers using the citation search feature, and they will be automatically added to your \texttt{cite.bib} file.

\section{Mathematical Expressions}
Here are some example equations:

\subsection{Quadratic Formula}
The quadratic formula is given by:
\[x = \frac{-b \pm \sqrt{b^2 - 4ac}}{2a}\]

\subsection{Matrix Example}
A sample matrix:
\[
\begin{pmatrix}
a_{11} & a_{12} & a_{13} \\
a_{21} & a_{22} & a_{23} \\
a_{31} & a_{32} & a_{33}
\end{pmatrix}
\]

\section{Tables}
Below is a sample table showing data:

\begin{table}[h]
\centering
\begin{tabular}{@{}lrrr@{}}
\toprule
\textbf{Category} & \textbf{Value 1} & \textbf{Value 2} & \textbf{Total} \\
\midrule
Item A & 10.5 & 20.3 & 30.8 \\
Item B & 15.2 & 25.1 & 40.3 \\
Item C & 22.1 & 12.7 & 34.8 \\
\bottomrule
\end{tabular}
\caption{Sample Data Table}
\label{tab:sample}
\end{table}

\section{Lists}
\subsection{Bullet Points}
\begin{itemize}
  \item First item
  \item Second item
  \item Third item with sub-items:
    \begin{itemize}
      \item Sub-item 1
      \item Sub-item 2
    \end{itemize}
\end{itemize}

\subsection{Numbered List}
\begin{enumerate}
  \item First step
  \item Second step
  \item Third step
\end{enumerate}

\section{Citations and References}
When you use the citation search feature, references will be automatically managed. For example, you might cite a paper like this: \cite{example2024}. All citations will be automatically added to your \texttt{cite.bib} file.

\section{Conclusion}
This template demonstrates basic \LaTeX{} features including mathematical formulas, tables, and lists. The bibliography system is now set up to work seamlessly with Bibby AI's citation management. Feel free to modify it for your needs!

\bibliographystyle{plain}
\bibliography{cite}

\end{document}\documentclass{article}
\usepackage[utf8]{inputenc}
\usepackage{booktabs}
\usepackage{amsmath}
\usepackage{natbib}

\title{My Document}
\author{Your Name}
\date{\today}

\begin{document}

\maketitle

\section{Introduction}
This is a sample document demonstrating various \LaTeX{} features including mathematical equations, tables, and formatting. You can cite papers using the citation search feature, and they will be automatically added to your \texttt{cite.bib} file.

\section{Mathematical Expressions}
Here are some example equations:

\subsection{Quadratic Formula}
The quadratic formula is given by:
\[x = \frac{-b \pm \sqrt{b^2 - 4ac}}{2a}\]

\subsection{Matrix Example}
A sample matrix:
\[
\begin{pmatrix}
a_{11} & a_{12} & a_{13} \\
a_{21} & a_{22} & a_{23} \\
a_{31} & a_{32} & a_{33}
\end{pmatrix}
\]

\section{Tables}
Below is a sample table showing data:

\begin{table}[h]
\centering
\begin{tabular}{@{}lrrr@{}}
\toprule
\textbf{Category} & \textbf{Value 1} & \textbf{Value 2} & \textbf{Total} \\
\midrule
Item A & 10.5 & 20.3 & 30.8 \\
Item B & 15.2 & 25.1 & 40.3 \\
Item C & 22.1 & 12.7 & 34.8 \\
\bottomrule
\end{tabular}
\caption{Sample Data Table}
\label{tab:sample}
\end{table}

\section{Lists}
\subsection{Bullet Points}
\begin{itemize}
  \item First item
  \item Second item
  \item Third item with sub-items:
    \begin{itemize}
      \item Sub-item 1
      \item Sub-item 2
    \end{itemize}
\end{itemize}

\subsection{Numbered List}
\begin{enumerate}
  \item First step
  \item Second step
  \item Third step
\end{enumerate}

\section{Citations and References}
When you use the citation search feature, references will be automatically managed. For example, you might cite a paper like this: \cite{example2024}. All citations will be automatically added to your \texttt{cite.bib} file.

\section{Conclusion}
This template demonstrates basic \LaTeX{} features including mathematical formulas, tables, and lists. The bibliography system is now set up to work seamlessly with Bibby AI's citation management. Feel free to modify it for your needs!

\bibliographystyle{plain}
\bibliography{cite}

\end{document}\documentclass{article}
\usepackage[utf8]{inputenc}
\usepackage{booktabs}
\usepackage{amsmath}
\usepackage{natbib}

\title{My Document}
\author{Your Name}
\date{\today}

\begin{document}

\maketitle

\section{Introduction}
This is a sample document demonstrating various \LaTeX{} features including mathematical equations, tables, and formatting. You can cite papers using the citation search feature, and they will be automatically added to your \texttt{cite.bib} file.

\section{Mathematical Expressions}
Here are some example equations:

\subsection{Quadratic Formula}
The quadratic formula is given by:
\[x = \frac{-b \pm \sqrt{b^2 - 4ac}}{2a}\]

\subsection{Matrix Example}
A sample matrix:
\[
\begin{pmatrix}
a_{11} & a_{12} & a_{13} \\
a_{21} & a_{22} & a_{23} \\
a_{31} & a_{32} & a_{33}
\end{pmatrix}
\]

\section{Tables}
Below is a sample table showing data:

\begin{table}[h]
\centering
\begin{tabular}{@{}lrrr@{}}
\toprule
\textbf{Category} & \textbf{Value 1} & \textbf{Value 2} & \textbf{Total} \\
\midrule
Item A & 10.5 & 20.3 & 30.8 \\
Item B & 15.2 & 25.1 & 40.3 \\
Item C & 22.1 & 12.7 & 34.8 \\
\bottomrule
\end{tabular}
\caption{Sample Data Table}
\label{tab:sample}
\end{table}

\section{Lists}
\subsection{Bullet Points}
\begin{itemize}
  \item First item
  \item Second item
  \item Third item with sub-items:
    \begin{itemize}
      \item Sub-item 1
      \item Sub-item 2
    \end{itemize}
\end{itemize}

\subsection{Numbered List}
\begin{enumerate}
  \item First step
  \item Second step
  \item Third step
\end{enumerate}

\section{Citations and References}
When you use the citation search feature, references will be automatically managed. For example, you might cite a paper like this: \cite{example2024}. All citations will be automatically added to your \texttt{cite.bib} file.

\section{Conclusion}
This template demonstrates basic \LaTeX{} features including mathematical formulas, tables, and lists. The bibliography system is now set up to work seamlessly with Bibby AI's citation management. Feel free to modify it for your needs!

\bibliographystyle{plain}
\bibliography{cite}

\end{document}\documentclass{article}
\usepackage[utf8]{inputenc}
\usepackage{booktabs}
\usepackage{amsmath}
\usepackage{natbib}

\title{My Document}
\author{Your Name}
\date{\today}

\begin{document}

\maketitle

\section{Introduction}
This is a sample document demonstrating various \LaTeX{} features including mathematical equations, tables, and formatting. You can cite papers using the citation search feature, and they will be automatically added to your \texttt{cite.bib} file.

\section{Mathematical Expressions}
Here are some example equations:

\subsection{Quadratic Formula}
The quadratic formula is given by:
\[x = \frac{-b \pm \sqrt{b^2 - 4ac}}{2a}\]

\subsection{Matrix Example}
A sample matrix:
\[
\begin{pmatrix}
a_{11} & a_{12} & a_{13} \\
a_{21} & a_{22} & a_{23} \\
a_{31} & a_{32} & a_{33}
\end{pmatrix}
\]

\section{Tables}
Below is a sample table showing data:

\begin{table}[h]
\centering
\begin{tabular}{@{}lrrr@{}}
\toprule
\textbf{Category} & \textbf{Value 1} & \textbf{Value 2} & \textbf{Total} \\
\midrule
Item A & 10.5 & 20.3 & 30.8 \\
Item B & 15.2 & 25.1 & 40.3 \\
Item C & 22.1 & 12.7 & 34.8 \\
\bottomrule
\end{tabular}
\caption{Sample Data Table}
\label{tab:sample}
\end{table}

\section{Lists}
\subsection{Bullet Points}
\begin{itemize}
  \item First item
  \item Second item
  \item Third item with sub-items:
    \begin{itemize}
      \item Sub-item 1
      \item Sub-item 2
    \end{itemize}
\end{itemize}

\subsection{Numbered List}
\begin{enumerate}
  \item First step
  \item Second step
  \item Third step
\end{enumerate}

\section{Citations and References}
When you use the citation search feature, references will be automatically managed. For example, you might cite a paper like this: \cite{example2024}. All citations will be automatically added to your \texttt{cite.bib} file.

\section{Conclusion}
This template demonstrates basic \LaTeX{} features including mathematical formulas, tables, and lists. The bibliography system is now set up to work seamlessly with Bibby AI's citation management. Feel free to modify it for your needs!

\bibliographystyle{plain}
\bibliography{cite}

\end{document}\documentclass{article}
\usepackage[utf8]{inputenc}
\usepackage{booktabs}
\usepackage{amsmath}
\usepackage{natbib}

\title{My Document}
\author{Your Name}
\date{\today}

\begin{document}

\maketitle

\section{Introduction}
This is a sample document demonstrating various \LaTeX{} features including mathematical equations, tables, and formatting. You can cite papers using the citation search feature, and they will be automatically added to your \texttt{cite.bib} file.

\section{Mathematical Expressions}
Here are some example equations:

\subsection{Quadratic Formula}
The quadratic formula is given by:
\[x = \frac{-b \pm \sqrt{b^2 - 4ac}}{2a}\]

\subsection{Matrix Example}
A sample matrix:
\[
\begin{pmatrix}
a_{11} & a_{12} & a_{13} \\
a_{21} & a_{22} & a_{23} \\
a_{31} & a_{32} & a_{33}
\end{pmatrix}
\]

\section{Tables}
Below is a sample table showing data:

\begin{table}[h]
\centering
\begin{tabular}{@{}lrrr@{}}
\toprule
\textbf{Category} & \textbf{Value 1} & \textbf{Value 2} & \textbf{Total} \\
\midrule
Item A & 10.5 & 20.3 & 30.8 \\
Item B & 15.2 & 25.1 & 40.3 \\
Item C & 22.1 & 12.7 & 34.8 \\
\bottomrule
\end{tabular}
\caption{Sample Data Table}
\label{tab:sample}
\end{table}

\section{Lists}
\subsection{Bullet Points}
\begin{itemize}
  \item First item
  \item Second item
  \item Third item with sub-items:
    \begin{itemize}
      \item Sub-item 1
      \item Sub-item 2
    \end{itemize}
\end{itemize}

\subsection{Numbered List}
\begin{enumerate}
  \item First step
  \item Second step
  \item Third step
\end{enumerate}

\section{Citations and References}
When you use the citation search feature, references will be automatically managed. For example, you might cite a paper like this: \cite{example2024}. All citations will be automatically added to your \texttt{cite.bib} file.

\section{Conclusion}
This template demonstrates basic \LaTeX{} features including mathematical formulas, tables, and lists. The bibliography system is now set up to work seamlessly with Bibby AI's citation management. Feel free to modify it for your needs!

\bibliographystyle{plain}
\bibliography{cite}

\end{document}\documentclass{article}
\usepackage[utf8]{inputenc}
\usepackage{booktabs}
\usepackage{amsmath}
\usepackage{natbib}

\title{My Document}
\author{Your Name}
\date{\today}

\begin{document}

\maketitle

\section{Introduction}
This is a sample document demonstrating various \LaTeX{} features including mathematical equations, tables, and formatting. You can cite papers using the citation search feature, and they will be automatically added to your \texttt{cite.bib} file.

\section{Mathematical Expressions}
Here are some example equations:

\subsection{Quadratic Formula}
The quadratic formula is given by:
\[x = \frac{-b \pm \sqrt{b^2 - 4ac}}{2a}\]

\subsection{Matrix Example}
A sample matrix:
\[
\begin{pmatrix}
a_{11} & a_{12} & a_{13} \\
a_{21} & a_{22} & a_{23} \\
a_{31} & a_{32} & a_{33}
\end{pmatrix}
\]

\section{Tables}
Below is a sample table showing data:

\begin{table}[h]
\centering
\begin{tabular}{@{}lrrr@{}}
\toprule
\textbf{Category} & \textbf{Value 1} & \textbf{Value 2} & \textbf{Total} \\
\midrule
Item A & 10.5 & 20.3 & 30.8 \\
Item B & 15.2 & 25.1 & 40.3 \\
Item C & 22.1 & 12.7 & 34.8 \\
\bottomrule
\end{tabular}
\caption{Sample Data Table}
\label{tab:sample}
\end{table}

\section{Lists}
\subsection{Bullet Points}
\begin{itemize}
  \item First item
  \item Second item
  \item Third item with sub-items:
    \begin{itemize}
      \item Sub-item 1
      \item Sub-item 2
    \end{itemize}
\end{itemize}

\subsection{Numbered List}
\begin{enumerate}
  \item First step
  \item Second step
  \item Third step
\end{enumerate}

\section{Citations and References}
When you use the citation search feature, references will be automatically managed. For example, you might cite a paper like this: \cite{example2024}. All citations will be automatically added to your \texttt{cite.bib} file.

\section{Conclusion}
This template demonstrates basic \LaTeX{} features including mathematical formulas, tables, and lists. The bibliography system is now set up to work seamlessly with Bibby AI's citation management. Feel free to modify it for your needs!

\bibliographystyle{plain}
\bibliography{cite}

\end{document}\documentclass{article}
\usepackage[utf8]{inputenc}
\usepackage{booktabs}
\usepackage{amsmath}
\usepackage{natbib}

\title{My Document}
\author{Your Name}
\date{\today}

\begin{document}

\maketitle

\section{Introduction}
This is a sample document demonstrating various \LaTeX{} features including mathematical equations, tables, and formatting. You can cite papers using the citation search feature, and they will be automatically added to your \texttt{cite.bib} file.

\section{Mathematical Expressions}
Here are some example equations:

\subsection{Quadratic Formula}
The quadratic formula is given by:
\[x = \frac{-b \pm \sqrt{b^2 - 4ac}}{2a}\]

\subsection{Matrix Example}
A sample matrix:
\[
\begin{pmatrix}
a_{11} & a_{12} & a_{13} \\
a_{21} & a_{22} & a_{23} \\
a_{31} & a_{32} & a_{33}
\end{pmatrix}
\]

\section{Tables}
Below is a sample table showing data:

\begin{table}[h]
\centering
\begin{tabular}{@{}lrrr@{}}
\toprule
\textbf{Category} & \textbf{Value 1} & \textbf{Value 2} & \textbf{Total} \\
\midrule
Item A & 10.5 & 20.3 & 30.8 \\
Item B & 15.2 & 25.1 & 40.3 \\
Item C & 22.1 & 12.7 & 34.8 \\
\bottomrule
\end{tabular}
\caption{Sample Data Table}
\label{tab:sample}
\end{table}

\section{Lists}
\subsection{Bullet Points}
\begin{itemize}
  \item First item
  \item Second item
  \item Third item with sub-items:
    \begin{itemize}
      \item Sub-item 1
      \item Sub-item 2
    \end{itemize}
\end{itemize}

\subsection{Numbered List}
\begin{enumerate}
  \item First step
  \item Second step
  \item Third step
\end{enumerate}

\section{Citations and References}
When you use the citation search feature, references will be automatically managed. For example, you might cite a paper like this: \cite{example2024}. All citations will be automatically added to your \texttt{cite.bib} file.

\section{Conclusion}
This template demonstrates basic \LaTeX{} features including mathematical formulas, tables, and lists. The bibliography system is now set up to work seamlessly with Bibby AI's citation management. Feel free to modify it for your needs!

\bibliographystyle{plain}
\bibliography{cite}

\end{document}\documentclass{article}
\usepackage[utf8]{inputenc}
\usepackage{booktabs}
\usepackage{amsmath}
\usepackage{natbib}

\title{My Document}
\author{Your Name}
\date{\today}

\begin{document}

\maketitle

\section{Introduction}
This is a sample document demonstrating various \LaTeX{} features including mathematical equations, tables, and formatting. You can cite papers using the citation search feature, and they will be automatically added to your \texttt{cite.bib} file.

\section{Mathematical Expressions}
Here are some example equations:

\subsection{Quadratic Formula}
The quadratic formula is given by:
\[x = \frac{-b \pm \sqrt{b^2 - 4ac}}{2a}\]

\subsection{Matrix Example}
A sample matrix:
\[
\begin{pmatrix}
a_{11} & a_{12} & a_{13} \\
a_{21} & a_{22} & a_{23} \\
a_{31} & a_{32} & a_{33}
\end{pmatrix}
\]

\section{Tables}
Below is a sample table showing data:

\begin{table}[h]
\centering
\begin{tabular}{@{}lrrr@{}}
\toprule
\textbf{Category} & \textbf{Value 1} & \textbf{Value 2} & \textbf{Total} \\
\midrule
Item A & 10.5 & 20.3 & 30.8 \\
Item B & 15.2 & 25.1 & 40.3 \\
Item C & 22.1 & 12.7 & 34.8 \\
\bottomrule
\end{tabular}
\caption{Sample Data Table}
\label{tab:sample}
\end{table}

\section{Lists}
\subsection{Bullet Points}
\begin{itemize}
  \item First item
  \item Second item
  \item Third item with sub-items:
    \begin{itemize}
      \item Sub-item 1
      \item Sub-item 2
    \end{itemize}
\end{itemize}

\subsection{Numbered List}
\begin{enumerate}
  \item First step
  \item Second step
  \item Third step
\end{enumerate}

\section{Citations and References}
When you use the citation search feature, references will be automatically managed. For example, you might cite a paper like this: \cite{example2024}. All citations will be automatically added to your \texttt{cite.bib} file.

\section{Conclusion}
This template demonstrates basic \LaTeX{} features including mathematical formulas, tables, and lists. The bibliography system is now set up to work seamlessly with Bibby AI's citation management. Feel free to modify it for your needs!

\bibliographystyle{plain}
\bibliography{cite}

\end{document}\documentclass{article}
\usepackage[utf8]{inputenc}
\usepackage{booktabs}
\usepackage{amsmath}
\usepackage{natbib}

\title{My Document}
\author{Your Name}
\date{\today}

\begin{document}

\maketitle

\section{Introduction}
This is a sample document demonstrating various \LaTeX{} features including mathematical equations, tables, and formatting. You can cite papers using the citation search feature, and they will be automatically added to your \texttt{cite.bib} file.

\section{Mathematical Expressions}
Here are some example equations:

\subsection{Quadratic Formula}
The quadratic formula is given by:
\[x = \frac{-b \pm \sqrt{b^2 - 4ac}}{2a}\]

\subsection{Matrix Example}
A sample matrix:
\[
\begin{pmatrix}
a_{11} & a_{12} & a_{13} \\
a_{21} & a_{22} & a_{23} \\
a_{31} & a_{32} & a_{33}
\end{pmatrix}
\]

\section{Tables}
Below is a sample table showing data:

\begin{table}[h]
\centering
\begin{tabular}{@{}lrrr@{}}
\toprule
\textbf{Category} & \textbf{Value 1} & \textbf{Value 2} & \textbf{Total} \\
\midrule
Item A & 10.5 & 20.3 & 30.8 \\
Item B & 15.2 & 25.1 & 40.3 \\
Item C & 22.1 & 12.7 & 34.8 \\
\bottomrule
\end{tabular}
\caption{Sample Data Table}
\label{tab:sample}
\end{table}

\section{Lists}
\subsection{Bullet Points}
\begin{itemize}
  \item First item
  \item Second item
  \item Third item with sub-items:
    \begin{itemize}
      \item Sub-item 1
      \item Sub-item 2
    \end{itemize}
\end{itemize}

\subsection{Numbered List}
\begin{enumerate}
  \item First step
  \item Second step
  \item Third step
\end{enumerate}

\section{Citations and References}
When you use the citation search feature, references will be automatically managed. For example, you might cite a paper like this: \cite{example2024}. All citations will be automatically added to your \texttt{cite.bib} file.

\section{Conclusion}
This template demonstrates basic \LaTeX{} features including mathematical formulas, tables, and lists. The bibliography system is now set up to work seamlessly with Bibby AI's citation management. Feel free to modify it for your needs!

\bibliographystyle{plain}
\bibliography{cite}

\end{document}\documentclass{article}
\usepackage[utf8]{inputenc}
\usepackage{booktabs}
\usepackage{amsmath}
\usepackage{natbib}

\title{My Document}
\author{Your Name}
\date{\today}

\begin{document}

\maketitle

\section{Introduction}
This is a sample document demonstrating various \LaTeX{} features including mathematical equations, tables, and formatting. You can cite papers using the citation search feature, and they will be automatically added to your \texttt{cite.bib} file.

\section{Mathematical Expressions}
Here are some example equations:

\subsection{Quadratic Formula}
The quadratic formula is given by:
\[x = \frac{-b \pm \sqrt{b^2 - 4ac}}{2a}\]

\subsection{Matrix Example}
A sample matrix:
\[
\begin{pmatrix}
a_{11} & a_{12} & a_{13} \\
a_{21} & a_{22} & a_{23} \\
a_{31} & a_{32} & a_{33}
\end{pmatrix}
\]

\section{Tables}
Below is a sample table showing data:

\begin{table}[h]
\centering
\begin{tabular}{@{}lrrr@{}}
\toprule
\textbf{Category} & \textbf{Value 1} & \textbf{Value 2} & \textbf{Total} \\
\midrule
Item A & 10.5 & 20.3 & 30.8 \\
Item B & 15.2 & 25.1 & 40.3 \\
Item C & 22.1 & 12.7 & 34.8 \\
\bottomrule
\end{tabular}
\caption{Sample Data Table}
\label{tab:sample}
\end{table}

\section{Lists}
\subsection{Bullet Points}
\begin{itemize}
  \item First item
  \item Second item
  \item Third item with sub-items:
    \begin{itemize}
      \item Sub-item 1
      \item Sub-item 2
    \end{itemize}
\end{itemize}

\subsection{Numbered List}
\begin{enumerate}
  \item First step
  \item Second step
  \item Third step
\end{enumerate}

\section{Citations and References}
When you use the citation search feature, references will be automatically managed. For example, you might cite a paper like this: \cite{example2024}. All citations will be automatically added to your \texttt{cite.bib} file.

\section{Conclusion}
This template demonstrates basic \LaTeX{} features including mathematical formulas, tables, and lists. The bibliography system is now set up to work seamlessly with Bibby AI's citation management. Feel free to modify it for your needs!

\bibliographystyle{plain}
\bibliography{cite}

\end{document}\documentclass{article}
\usepackage[utf8]{inputenc}
\usepackage{booktabs}
\usepackage{amsmath}
\usepackage{natbib}

\title{My Document}
\author{Your Name}
\date{\today}

\begin{document}

\maketitle

\section{Introduction}
This is a sample document demonstrating various \LaTeX{} features including mathematical equations, tables, and formatting. You can cite papers using the citation search feature, and they will be automatically added to your \texttt{cite.bib} file.

\section{Mathematical Expressions}
Here are some example equations:

\subsection{Quadratic Formula}
The quadratic formula is given by:
\[x = \frac{-b \pm \sqrt{b^2 - 4ac}}{2a}\]

\subsection{Matrix Example}
A sample matrix:
\[
\begin{pmatrix}
a_{11} & a_{12} & a_{13} \\
a_{21} & a_{22} & a_{23} \\
a_{31} & a_{32} & a_{33}
\end{pmatrix}
\]

\section{Tables}
Below is a sample table showing data:

\begin{table}[h]
\centering
\begin{tabular}{@{}lrrr@{}}
\toprule
\textbf{Category} & \textbf{Value 1} & \textbf{Value 2} & \textbf{Total} \\
\midrule
Item A & 10.5 & 20.3 & 30.8 \\
Item B & 15.2 & 25.1 & 40.3 \\
Item C & 22.1 & 12.7 & 34.8 \\
\bottomrule
\end{tabular}
\caption{Sample Data Table}
\label{tab:sample}
\end{table}

\section{Lists}
\subsection{Bullet Points}
\begin{itemize}
  \item First item
  \item Second item
  \item Third item with sub-items:
    \begin{itemize}
      \item Sub-item 1
      \item Sub-item 2
    \end{itemize}
\end{itemize}

\subsection{Numbered List}
\begin{enumerate}
  \item First step
  \item Second step
  \item Third step
\end{enumerate}

\section{Citations and References}
When you use the citation search feature, references will be automatically managed. For example, you might cite a paper like this: \cite{example2024}. All citations will be automatically added to your \texttt{cite.bib} file.

\section{Conclusion}
This template demonstrates basic \LaTeX{} features including mathematical formulas, tables, and lists. The bibliography system is now set up to work seamlessly with Bibby AI's citation management. Feel free to modify it for your needs!

\bibliographystyle{plain}
\bibliography{cite}

\end{document}\documentclass{article}
\usepackage[utf8]{inputenc}
\usepackage{booktabs}
\usepackage{amsmath}
\usepackage{natbib}

\title{My Document}
\author{Your Name}
\date{\today}

\begin{document}

\maketitle

\section{Introduction}
This is a sample document demonstrating various \LaTeX{} features including mathematical equations, tables, and formatting. You can cite papers using the citation search feature, and they will be automatically added to your \texttt{cite.bib} file.

\section{Mathematical Expressions}
Here are some example equations:

\subsection{Quadratic Formula}
The quadratic formula is given by:
\[x = \frac{-b \pm \sqrt{b^2 - 4ac}}{2a}\]

\subsection{Matrix Example}
A sample matrix:
\[
\begin{pmatrix}
a_{11} & a_{12} & a_{13} \\
a_{21} & a_{22} & a_{23} \\
a_{31} & a_{32} & a_{33}
\end{pmatrix}
\]

\section{Tables}
Below is a sample table showing data:

\begin{table}[h]
\centering
\begin{tabular}{@{}lrrr@{}}
\toprule
\textbf{Category} & \textbf{Value 1} & \textbf{Value 2} & \textbf{Total} \\
\midrule
Item A & 10.5 & 20.3 & 30.8 \\
Item B & 15.2 & 25.1 & 40.3 \\
Item C & 22.1 & 12.7 & 34.8 \\
\bottomrule
\end{tabular}
\caption{Sample Data Table}
\label{tab:sample}
\end{table}

\section{Lists}
\subsection{Bullet Points}
\begin{itemize}
  \item First item
  \item Second item
  \item Third item with sub-items:
    \begin{itemize}
      \item Sub-item 1
      \item Sub-item 2
    \end{itemize}
\end{itemize}

\subsection{Numbered List}
\begin{enumerate}
  \item First step
  \item Second step
  \item Third step
\end{enumerate}

\section{Citations and References}
When you use the citation search feature, references will be automatically managed. For example, you might cite a paper like this: \cite{example2024}. All citations will be automatically added to your \texttt{cite.bib} file.

\section{Conclusion}
This template demonstrates basic \LaTeX{} features including mathematical formulas, tables, and lists. The bibliography system is now set up to work seamlessly with Bibby AI's citation management. Feel free to modify it for your needs!

\bibliographystyle{plain}
\bibliography{cite}

\end{document}\documentclass{article}
\usepackage[utf8]{inputenc}
\usepackage{booktabs}
\usepackage{amsmath}
\usepackage{natbib}

\title{My Document}
\author{Your Name}
\date{\today}

\begin{document}

\maketitle

\section{Introduction}
This is a sample document demonstrating various \LaTeX{} features including mathematical equations, tables, and formatting. You can cite papers using the citation search feature, and they will be automatically added to your \texttt{cite.bib} file.

\section{Mathematical Expressions}
Here are some example equations:

\subsection{Quadratic Formula}
The quadratic formula is given by:
\[x = \frac{-b \pm \sqrt{b^2 - 4ac}}{2a}\]

\subsection{Matrix Example}
A sample matrix:
\[
\begin{pmatrix}
a_{11} & a_{12} & a_{13} \\
a_{21} & a_{22} & a_{23} \\
a_{31} & a_{32} & a_{33}
\end{pmatrix}
\]

\section{Tables}
Below is a sample table showing data:

\begin{table}[h]
\centering
\begin{tabular}{@{}lrrr@{}}
\toprule
\textbf{Category} & \textbf{Value 1} & \textbf{Value 2} & \textbf{Total} \\
\midrule
Item A & 10.5 & 20.3 & 30.8 \\
Item B & 15.2 & 25.1 & 40.3 \\
Item C & 22.1 & 12.7 & 34.8 \\
\bottomrule
\end{tabular}
\caption{Sample Data Table}
\label{tab:sample}
\end{table}

\section{Lists}
\subsection{Bullet Points}
\begin{itemize}
  \item First item
  \item Second item
  \item Third item with sub-items:
    \begin{itemize}
      \item Sub-item 1
      \item Sub-item 2
    \end{itemize}
\end{itemize}

\subsection{Numbered List}
\begin{enumerate}
  \item First step
  \item Second step
  \item Third step
\end{enumerate}

\section{Citations and References}
When you use the citation search feature, references will be automatically managed. For example, you might cite a paper like this: \cite{example2024}. All citations will be automatically added to your \texttt{cite.bib} file.

\section{Conclusion}
This template demonstrates basic \LaTeX{} features including mathematical formulas, tables, and lists. The bibliography system is now set up to work seamlessly with Bibby AI's citation management. Feel free to modify it for your needs!

\bibliographystyle{plain}
\bibliography{cite}

\end{document}\documentclass{article}
\usepackage[utf8]{inputenc}
\usepackage{booktabs}
\usepackage{amsmath}
\usepackage{natbib}

\title{My Document}
\author{Your Name}
\date{\today}

\begin{document}

\maketitle

\section{Introduction}
This is a sample document demonstrating various \LaTeX{} features including mathematical equations, tables, and formatting. You can cite papers using the citation search feature, and they will be automatically added to your \texttt{cite.bib} file.

\section{Mathematical Expressions}
Here are some example equations:

\subsection{Quadratic Formula}
The quadratic formula is given by:
\[x = \frac{-b \pm \sqrt{b^2 - 4ac}}{2a}\]

\subsection{Matrix Example}
A sample matrix:
\[
\begin{pmatrix}
a_{11} & a_{12} & a_{13} \\
a_{21} & a_{22} & a_{23} \\
a_{31} & a_{32} & a_{33}
\end{pmatrix}
\]

\section{Tables}
Below is a sample table showing data:

\begin{table}[h]
\centering
\begin{tabular}{@{}lrrr@{}}
\toprule
\textbf{Category} & \textbf{Value 1} & \textbf{Value 2} & \textbf{Total} \\
\midrule
Item A & 10.5 & 20.3 & 30.8 \\
Item B & 15.2 & 25.1 & 40.3 \\
Item C & 22.1 & 12.7 & 34.8 \\
\bottomrule
\end{tabular}
\caption{Sample Data Table}
\label{tab:sample}
\end{table}

\section{Lists}
\subsection{Bullet Points}
\begin{itemize}
  \item First item
  \item Second item
  \item Third item with sub-items:
    \begin{itemize}
      \item Sub-item 1
      \item Sub-item 2
    \end{itemize}
\end{itemize}

\subsection{Numbered List}
\begin{enumerate}
  \item First step
  \item Second step
  \item Third step
\end{enumerate}

\section{Citations and References}
When you use the citation search feature, references will be automatically managed. For example, you might cite a paper like this: \cite{example2024}. All citations will be automatically added to your \texttt{cite.bib} file.

\section{Conclusion}
This template demonstrates basic \LaTeX{} features including mathematical formulas, tables, and lists. The bibliography system is now set up to work seamlessly with Bibby AI's citation management. Feel free to modify it for your needs!

\bibliographystyle{plain}
\bibliography{cite}

\end{document}\documentclass{article}
\usepackage[utf8]{inputenc}
\usepackage{booktabs}
\usepackage{amsmath}
\usepackage{natbib}

\title{My Document}
\author{Your Name}
\date{\today}

\begin{document}

\maketitle

\section{Introduction}
This is a sample document demonstrating various \LaTeX{} features including mathematical equations, tables, and formatting. You can cite papers using the citation search feature, and they will be automatically added to your \texttt{cite.bib} file.

\section{Mathematical Expressions}
Here are some example equations:

\subsection{Quadratic Formula}
The quadratic formula is given by:
\[x = \frac{-b \pm \sqrt{b^2 - 4ac}}{2a}\]

\subsection{Matrix Example}
A sample matrix:
\[
\begin{pmatrix}
a_{11} & a_{12} & a_{13} \\
a_{21} & a_{22} & a_{23} \\
a_{31} & a_{32} & a_{33}
\end{pmatrix}
\]

\section{Tables}
Below is a sample table showing data:

\begin{table}[h]
\centering
\begin{tabular}{@{}lrrr@{}}
\toprule
\textbf{Category} & \textbf{Value 1} & \textbf{Value 2} & \textbf{Total} \\
\midrule
Item A & 10.5 & 20.3 & 30.8 \\
Item B & 15.2 & 25.1 & 40.3 \\
Item C & 22.1 & 12.7 & 34.8 \\
\bottomrule
\end{tabular}
\caption{Sample Data Table}
\label{tab:sample}
\end{table}

\section{Lists}
\subsection{Bullet Points}
\begin{itemize}
  \item First item
  \item Second item
  \item Third item with sub-items:
    \begin{itemize}
      \item Sub-item 1
      \item Sub-item 2
    \end{itemize}
\end{itemize}

\subsection{Numbered List}
\begin{enumerate}
  \item First step
  \item Second step
  \item Third step
\end{enumerate}

\section{Citations and References}
When you use the citation search feature, references will be automatically managed. For example, you might cite a paper like this: \cite{example2024}. All citations will be automatically added to your \texttt{cite.bib} file.

\section{Conclusion}
This template demonstrates basic \LaTeX{} features including mathematical formulas, tables, and lists. The bibliography system is now set up to work seamlessly with Bibby AI's citation management. Feel free to modify it for your needs!

\bibliographystyle{plain}
\bibliography{cite}

\end{document}\documentclass{article}
\usepackage[utf8]{inputenc}
\usepackage{booktabs}
\usepackage{amsmath}
\usepackage{natbib}

\title{My Document}
\author{Your Name}
\date{\today}

\begin{document}

\maketitle

\section{Introduction}
This is a sample document demonstrating various \LaTeX{} features including mathematical equations, tables, and formatting. You can cite papers using the citation search feature, and they will be automatically added to your \texttt{cite.bib} file.

\section{Mathematical Expressions}
Here are some example equations:

\subsection{Quadratic Formula}
The quadratic formula is given by:
\[x = \frac{-b \pm \sqrt{b^2 - 4ac}}{2a}\]

\subsection{Matrix Example}
A sample matrix:
\[
\begin{pmatrix}
a_{11} & a_{12} & a_{13} \\
a_{21} & a_{22} & a_{23} \\
a_{31} & a_{32} & a_{33}
\end{pmatrix}
\]

\section{Tables}
Below is a sample table showing data:

\begin{table}[h]
\centering
\begin{tabular}{@{}lrrr@{}}
\toprule
\textbf{Category} & \textbf{Value 1} & \textbf{Value 2} & \textbf{Total} \\
\midrule
Item A & 10.5 & 20.3 & 30.8 \\
Item B & 15.2 & 25.1 & 40.3 \\
Item C & 22.1 & 12.7 & 34.8 \\
\bottomrule
\end{tabular}
\caption{Sample Data Table}
\label{tab:sample}
\end{table}

\section{Lists}
\subsection{Bullet Points}
\begin{itemize}
  \item First item
  \item Second item
  \item Third item with sub-items:
    \begin{itemize}
      \item Sub-item 1
      \item Sub-item 2
    \end{itemize}
\end{itemize}

\subsection{Numbered List}
\begin{enumerate}
  \item First step
  \item Second step
  \item Third step
\end{enumerate}

\section{Citations and References}
When you use the citation search feature, references will be automatically managed. For example, you might cite a paper like this: \cite{example2024}. All citations will be automatically added to your \texttt{cite.bib} file.

\section{Conclusion}
This template demonstrates basic \LaTeX{} features including mathematical formulas, tables, and lists. The bibliography system is now set up to work seamlessly with Bibby AI's citation management. Feel free to modify it for your needs!

\bibliographystyle{plain}
\bibliography{cite}

\end{document}\documentclass{article}
\usepackage[utf8]{inputenc}
\usepackage{booktabs}
\usepackage{amsmath}
\usepackage{natbib}

\title{My Document}
\author{Your Name}
\date{\today}

\begin{document}

\maketitle

\section{Introduction}
This is a sample document demonstrating various \LaTeX{} features including mathematical equations, tables, and formatting. You can cite papers using the citation search feature, and they will be automatically added to your \texttt{cite.bib} file.

\section{Mathematical Expressions}
Here are some example equations:

\subsection{Quadratic Formula}
The quadratic formula is given by:
\[x = \frac{-b \pm \sqrt{b^2 - 4ac}}{2a}\]

\subsection{Matrix Example}
A sample matrix:
\[
\begin{pmatrix}
a_{11} & a_{12} & a_{13} \\
a_{21} & a_{22} & a_{23} \\
a_{31} & a_{32} & a_{33}
\end{pmatrix}
\]

\section{Tables}
Below is a sample table showing data:

\begin{table}[h]
\centering
\begin{tabular}{@{}lrrr@{}}
\toprule
\textbf{Category} & \textbf{Value 1} & \textbf{Value 2} & \textbf{Total} \\
\midrule
Item A & 10.5 & 20.3 & 30.8 \\
Item B & 15.2 & 25.1 & 40.3 \\
Item C & 22.1 & 12.7 & 34.8 \\
\bottomrule
\end{tabular}
\caption{Sample Data Table}
\label{tab:sample}
\end{table}

\section{Lists}
\subsection{Bullet Points}
\begin{itemize}
  \item First item
  \item Second item
  \item Third item with sub-items:
    \begin{itemize}
      \item Sub-item 1
      \item Sub-item 2
    \end{itemize}
\end{itemize}

\subsection{Numbered List}
\begin{enumerate}
  \item First step
  \item Second step
  \item Third step
\end{enumerate}

\section{Citations and References}
When you use the citation search feature, references will be automatically managed. For example, you might cite a paper like this: \cite{example2024}. All citations will be automatically added to your \texttt{cite.bib} file.

\section{Conclusion}
This template demonstrates basic \LaTeX{} features including mathematical formulas, tables, and lists. The bibliography system is now set up to work seamlessly with Bibby AI's citation management. Feel free to modify it for your needs!

\bibliographystyle{plain}
\bibliography{cite}

\end{document}\documentclass{article}
\usepackage[utf8]{inputenc}
\usepackage{booktabs}
\usepackage{amsmath}
\usepackage{natbib}

\title{My Document}
\author{Your Name}
\date{\today}

\begin{document}

\maketitle

\section{Introduction}
This is a sample document demonstrating various \LaTeX{} features including mathematical equations, tables, and formatting. You can cite papers using the citation search feature, and they will be automatically added to your \texttt{cite.bib} file.

\section{Mathematical Expressions}
Here are some example equations:

\subsection{Quadratic Formula}
The quadratic formula is given by:
\[x = \frac{-b \pm \sqrt{b^2 - 4ac}}{2a}\]

\subsection{Matrix Example}
A sample matrix:
\[
\begin{pmatrix}
a_{11} & a_{12} & a_{13} \\
a_{21} & a_{22} & a_{23} \\
a_{31} & a_{32} & a_{33}
\end{pmatrix}
\]

\section{Tables}
Below is a sample table showing data:

\begin{table}[h]
\centering
\begin{tabular}{@{}lrrr@{}}
\toprule
\textbf{Category} & \textbf{Value 1} & \textbf{Value 2} & \textbf{Total} \\
\midrule
Item A & 10.5 & 20.3 & 30.8 \\
Item B & 15.2 & 25.1 & 40.3 \\
Item C & 22.1 & 12.7 & 34.8 \\
\bottomrule
\end{tabular}
\caption{Sample Data Table}
\label{tab:sample}
\end{table}

\section{Lists}
\subsection{Bullet Points}
\begin{itemize}
  \item First item
  \item Second item
  \item Third item with sub-items:
    \begin{itemize}
      \item Sub-item 1
      \item Sub-item 2
    \end{itemize}
\end{itemize}

\subsection{Numbered List}
\begin{enumerate}
  \item First step
  \item Second step
  \item Third step
\end{enumerate}

\section{Citations and References}
When you use the citation search feature, references will be automatically managed. For example, you might cite a paper like this: \cite{example2024}. All citations will be automatically added to your \texttt{cite.bib} file.

\section{Conclusion}
This template demonstrates basic \LaTeX{} features including mathematical formulas, tables, and lists. The bibliography system is now set up to work seamlessly with Bibby AI's citation management. Feel free to modify it for your needs!

\bibliographystyle{plain}
\bibliography{cite}

\end{document}\documentclass{article}
\usepackage[utf8]{inputenc}
\usepackage{booktabs}
\usepackage{amsmath}
\usepackage{natbib}

\title{My Document}
\author{Your Name}
\date{\today}

\begin{document}

\maketitle

\section{Introduction}
This is a sample document demonstrating various \LaTeX{} features including mathematical equations, tables, and formatting. You can cite papers using the citation search feature, and they will be automatically added to your \texttt{cite.bib} file.

\section{Mathematical Expressions}
Here are some example equations:

\subsection{Quadratic Formula}
The quadratic formula is given by:
\[x = \frac{-b \pm \sqrt{b^2 - 4ac}}{2a}\]

\subsection{Matrix Example}
A sample matrix:
\[
\begin{pmatrix}
a_{11} & a_{12} & a_{13} \\
a_{21} & a_{22} & a_{23} \\
a_{31} & a_{32} & a_{33}
\end{pmatrix}
\]

\section{Tables}
Below is a sample table showing data:

\begin{table}[h]
\centering
\begin{tabular}{@{}lrrr@{}}
\toprule
\textbf{Category} & \textbf{Value 1} & \textbf{Value 2} & \textbf{Total} \\
\midrule
Item A & 10.5 & 20.3 & 30.8 \\
Item B & 15.2 & 25.1 & 40.3 \\
Item C & 22.1 & 12.7 & 34.8 \\
\bottomrule
\end{tabular}
\caption{Sample Data Table}
\label{tab:sample}
\end{table}

\section{Lists}
\subsection{Bullet Points}
\begin{itemize}
  \item First item
  \item Second item
  \item Third item with sub-items:
    \begin{itemize}
      \item Sub-item 1
      \item Sub-item 2
    \end{itemize}
\end{itemize}

\subsection{Numbered List}
\begin{enumerate}
  \item First step
  \item Second step
  \item Third step
\end{enumerate}

\section{Citations and References}
When you use the citation search feature, references will be automatically managed. For example, you might cite a paper like this: \cite{example2024}. All citations will be automatically added to your \texttt{cite.bib} file.

\section{Conclusion}
This template demonstrates basic \LaTeX{} features including mathematical formulas, tables, and lists. The bibliography system is now set up to work seamlessly with Bibby AI's citation management. Feel free to modify it for your needs!

\bibliographystyle{plain}
\bibliography{cite}

\end{document}\documentclass{article}
\usepackage[utf8]{inputenc}
\usepackage{booktabs}
\usepackage{amsmath}
\usepackage{natbib}

\title{My Document}
\author{Your Name}
\date{\today}

\begin{document}

\maketitle

\section{Introduction}
This is a sample document demonstrating various \LaTeX{} features including mathematical equations, tables, and formatting. You can cite papers using the citation search feature, and they will be automatically added to your \texttt{cite.bib} file.

\section{Mathematical Expressions}
Here are some example equations:

\subsection{Quadratic Formula}
The quadratic formula is given by:
\[x = \frac{-b \pm \sqrt{b^2 - 4ac}}{2a}\]

\subsection{Matrix Example}
A sample matrix:
\[
\begin{pmatrix}
a_{11} & a_{12} & a_{13} \\
a_{21} & a_{22} & a_{23} \\
a_{31} & a_{32} & a_{33}
\end{pmatrix}
\]

\section{Tables}
Below is a sample table showing data:

\begin{table}[h]
\centering
\begin{tabular}{@{}lrrr@{}}
\toprule
\textbf{Category} & \textbf{Value 1} & \textbf{Value 2} & \textbf{Total} \\
\midrule
Item A & 10.5 & 20.3 & 30.8 \\
Item B & 15.2 & 25.1 & 40.3 \\
Item C & 22.1 & 12.7 & 34.8 \\
\bottomrule
\end{tabular}
\caption{Sample Data Table}
\label{tab:sample}
\end{table}

\section{Lists}
\subsection{Bullet Points}
\begin{itemize}
  \item First item
  \item Second item
  \item Third item with sub-items:
    \begin{itemize}
      \item Sub-item 1
      \item Sub-item 2
    \end{itemize}
\end{itemize}

\subsection{Numbered List}
\begin{enumerate}
  \item First step
  \item Second step
  \item Third step
\end{enumerate}

\section{Citations and References}
When you use the citation search feature, references will be automatically managed. For example, you might cite a paper like this: \cite{example2024}. All citations will be automatically added to your \texttt{cite.bib} file.

\section{Conclusion}
This template demonstrates basic \LaTeX{} features including mathematical formulas, tables, and lists. The bibliography system is now set up to work seamlessly with Bibby AI's citation management. Feel free to modify it for your needs!

\bibliographystyle{plain}
\bibliography{cite}

\end{document}\documentclass{article}
\usepackage[utf8]{inputenc}
\usepackage{booktabs}
\usepackage{amsmath}
\usepackage{natbib}

\title{My Document}
\author{Your Name}
\date{\today}

\begin{document}

\maketitle

\section{Introduction}
This is a sample document demonstrating various \LaTeX{} features including mathematical equations, tables, and formatting. You can cite papers using the citation search feature, and they will be automatically added to your \texttt{cite.bib} file.

\section{Mathematical Expressions}
Here are some example equations:

\subsection{Quadratic Formula}
The quadratic formula is given by:
\[x = \frac{-b \pm \sqrt{b^2 - 4ac}}{2a}\]

\subsection{Matrix Example}
A sample matrix:
\[
\begin{pmatrix}
a_{11} & a_{12} & a_{13} \\
a_{21} & a_{22} & a_{23} \\
a_{31} & a_{32} & a_{33}
\end{pmatrix}
\]

\section{Tables}
Below is a sample table showing data:

\begin{table}[h]
\centering
\begin{tabular}{@{}lrrr@{}}
\toprule
\textbf{Category} & \textbf{Value 1} & \textbf{Value 2} & \textbf{Total} \\
\midrule
Item A & 10.5 & 20.3 & 30.8 \\
Item B & 15.2 & 25.1 & 40.3 \\
Item C & 22.1 & 12.7 & 34.8 \\
\bottomrule
\end{tabular}
\caption{Sample Data Table}
\label{tab:sample}
\end{table}

\section{Lists}
\subsection{Bullet Points}
\begin{itemize}
  \item First item
  \item Second item
  \item Third item with sub-items:
    \begin{itemize}
      \item Sub-item 1
      \item Sub-item 2
    \end{itemize}
\end{itemize}

\subsection{Numbered List}
\begin{enumerate}
  \item First step
  \item Second step
  \item Third step
\end{enumerate}

\section{Citations and References}
When you use the citation search feature, references will be automatically managed. For example, you might cite a paper like this: \cite{example2024}. All citations will be automatically added to your \texttt{cite.bib} file.

\section{Conclusion}
This template demonstrates basic \LaTeX{} features including mathematical formulas, tables, and lists. The bibliography system is now set up to work seamlessly with Bibby AI's citation management. Feel free to modify it for your needs!

\bibliographystyle{plain}
\bibliography{cite}

\end{document}\documentclass{article}
\usepackage[utf8]{inputenc}
\usepackage{booktabs}
\usepackage{amsmath}
\usepackage{natbib}

\title{My Document}
\author{Your Name}
\date{\today}

\begin{document}

\maketitle

\section{Introduction}
This is a sample document demonstrating various \LaTeX{} features including mathematical equations, tables, and formatting. You can cite papers using the citation search feature, and they will be automatically added to your \texttt{cite.bib} file.

\section{Mathematical Expressions}
Here are some example equations:

\subsection{Quadratic Formula}
The quadratic formula is given by:
\[x = \frac{-b \pm \sqrt{b^2 - 4ac}}{2a}\]

\subsection{Matrix Example}
A sample matrix:
\[
\begin{pmatrix}
a_{11} & a_{12} & a_{13} \\
a_{21} & a_{22} & a_{23} \\
a_{31} & a_{32} & a_{33}
\end{pmatrix}
\]

\section{Tables}
Below is a sample table showing data:

\begin{table}[h]
\centering
\begin{tabular}{@{}lrrr@{}}
\toprule
\textbf{Category} & \textbf{Value 1} & \textbf{Value 2} & \textbf{Total} \\
\midrule
Item A & 10.5 & 20.3 & 30.8 \\
Item B & 15.2 & 25.1 & 40.3 \\
Item C & 22.1 & 12.7 & 34.8 \\
\bottomrule
\end{tabular}
\caption{Sample Data Table}
\label{tab:sample}
\end{table}

\section{Lists}
\subsection{Bullet Points}
\begin{itemize}
  \item First item
  \item Second item
  \item Third item with sub-items:
    \begin{itemize}
      \item Sub-item 1
      \item Sub-item 2
    \end{itemize}
\end{itemize}

\subsection{Numbered List}
\begin{enumerate}
  \item First step
  \item Second step
  \item Third step
\end{enumerate}

\section{Citations and References}
When you use the citation search feature, references will be automatically managed. For example, you might cite a paper like this: \cite{example2024}. All citations will be automatically added to your \texttt{cite.bib} file.

\section{Conclusion}
This template demonstrates basic \LaTeX{} features including mathematical formulas, tables, and lists. The bibliography system is now set up to work seamlessly with Bibby AI's citation management. Feel free to modify it for your needs!

\bibliographystyle{plain}
\bibliography{cite}

\end{document}\documentclass{article}
\usepackage[utf8]{inputenc}
\usepackage{booktabs}
\usepackage{amsmath}
\usepackage{natbib}

\title{My Document}
\author{Your Name}
\date{\today}

\begin{document}

\maketitle

\section{Introduction}
This is a sample document demonstrating various \LaTeX{} features including mathematical equations, tables, and formatting. You can cite papers using the citation search feature, and they will be automatically added to your \texttt{cite.bib} file.

\section{Mathematical Expressions}
Here are some example equations:

\subsection{Quadratic Formula}
The quadratic formula is given by:
\[x = \frac{-b \pm \sqrt{b^2 - 4ac}}{2a}\]

\subsection{Matrix Example}
A sample matrix:
\[
\begin{pmatrix}
a_{11} & a_{12} & a_{13} \\
a_{21} & a_{22} & a_{23} \\
a_{31} & a_{32} & a_{33}
\end{pmatrix}
\]

\section{Tables}
Below is a sample table showing data:

\begin{table}[h]
\centering
\begin{tabular}{@{}lrrr@{}}
\toprule
\textbf{Category} & \textbf{Value 1} & \textbf{Value 2} & \textbf{Total} \\
\midrule
Item A & 10.5 & 20.3 & 30.8 \\
Item B & 15.2 & 25.1 & 40.3 \\
Item C & 22.1 & 12.7 & 34.8 \\
\bottomrule
\end{tabular}
\caption{Sample Data Table}
\label{tab:sample}
\end{table}

\section{Lists}
\subsection{Bullet Points}
\begin{itemize}
  \item First item
  \item Second item
  \item Third item with sub-items:
    \begin{itemize}
      \item Sub-item 1
      \item Sub-item 2
    \end{itemize}
\end{itemize}

\subsection{Numbered List}
\begin{enumerate}
  \item First step
  \item Second step
  \item Third step
\end{enumerate}

\section{Citations and References}
When you use the citation search feature, references will be automatically managed. For example, you might cite a paper like this: \cite{example2024}. All citations will be automatically added to your \texttt{cite.bib} file.

\section{Conclusion}
This template demonstrates basic \LaTeX{} features including mathematical formulas, tables, and lists. The bibliography system is now set up to work seamlessly with Bibby AI's citation management. Feel free to modify it for your needs!

\bibliographystyle{plain}
\bibliography{cite}

\end{document}\documentclass{article}
\usepackage[utf8]{inputenc}
\usepackage{booktabs}
\usepackage{amsmath}
\usepackage{natbib}

\title{My Document}
\author{Your Name}
\date{\today}

\begin{document}

\maketitle

\section{Introduction}
This is a sample document demonstrating various \LaTeX{} features including mathematical equations, tables, and formatting. You can cite papers using the citation search feature, and they will be automatically added to your \texttt{cite.bib} file.

\section{Mathematical Expressions}
Here are some example equations:

\subsection{Quadratic Formula}
The quadratic formula is given by:
\[x = \frac{-b \pm \sqrt{b^2 - 4ac}}{2a}\]

\subsection{Matrix Example}
A sample matrix:
\[
\begin{pmatrix}
a_{11} & a_{12} & a_{13} \\
a_{21} & a_{22} & a_{23} \\
a_{31} & a_{32} & a_{33}
\end{pmatrix}
\]

\section{Tables}
Below is a sample table showing data:

\begin{table}[h]
\centering
\begin{tabular}{@{}lrrr@{}}
\toprule
\textbf{Category} & \textbf{Value 1} & \textbf{Value 2} & \textbf{Total} \\
\midrule
Item A & 10.5 & 20.3 & 30.8 \\
Item B & 15.2 & 25.1 & 40.3 \\
Item C & 22.1 & 12.7 & 34.8 \\
\bottomrule
\end{tabular}
\caption{Sample Data Table}
\label{tab:sample}
\end{table}

\section{Lists}
\subsection{Bullet Points}
\begin{itemize}
  \item First item
  \item Second item
  \item Third item with sub-items:
    \begin{itemize}
      \item Sub-item 1
      \item Sub-item 2
    \end{itemize}
\end{itemize}

\subsection{Numbered List}
\begin{enumerate}
  \item First step
  \item Second step
  \item Third step
\end{enumerate}

\section{Citations and References}
When you use the citation search feature, references will be automatically managed. For example, you might cite a paper like this: \cite{example2024}. All citations will be automatically added to your \texttt{cite.bib} file.

\section{Conclusion}
This template demonstrates basic \LaTeX{} features including mathematical formulas, tables, and lists. The bibliography system is now set up to work seamlessly with Bibby AI's citation management. Feel free to modify it for your needs!

\bibliographystyle{plain}
\bibliography{cite}

\end{document}\documentclass{article}
\usepackage[utf8]{inputenc}
\usepackage{booktabs}
\usepackage{amsmath}
\usepackage{natbib}

\title{My Document}
\author{Your Name}
\date{\today}

\begin{document}

\maketitle

\section{Introduction}
This is a sample document demonstrating various \LaTeX{} features including mathematical equations, tables, and formatting. You can cite papers using the citation search feature, and they will be automatically added to your \texttt{cite.bib} file.

\section{Mathematical Expressions}
Here are some example equations:

\subsection{Quadratic Formula}
The quadratic formula is given by:
\[x = \frac{-b \pm \sqrt{b^2 - 4ac}}{2a}\]

\subsection{Matrix Example}
A sample matrix:
\[
\begin{pmatrix}
a_{11} & a_{12} & a_{13} \\
a_{21} & a_{22} & a_{23} \\
a_{31} & a_{32} & a_{33}
\end{pmatrix}
\]

\section{Tables}
Below is a sample table showing data:

\begin{table}[h]
\centering
\begin{tabular}{@{}lrrr@{}}
\toprule
\textbf{Category} & \textbf{Value 1} & \textbf{Value 2} & \textbf{Total} \\
\midrule
Item A & 10.5 & 20.3 & 30.8 \\
Item B & 15.2 & 25.1 & 40.3 \\
Item C & 22.1 & 12.7 & 34.8 \\
\bottomrule
\end{tabular}
\caption{Sample Data Table}
\label{tab:sample}
\end{table}

\section{Lists}
\subsection{Bullet Points}
\begin{itemize}
  \item First item
  \item Second item
  \item Third item with sub-items:
    \begin{itemize}
      \item Sub-item 1
      \item Sub-item 2
    \end{itemize}
\end{itemize}

\subsection{Numbered List}
\begin{enumerate}
  \item First step
  \item Second step
  \item Third step
\end{enumerate}

\section{Citations and References}
When you use the citation search feature, references will be automatically managed. For example, you might cite a paper like this: \cite{example2024}. All citations will be automatically added to your \texttt{cite.bib} file.

\section{Conclusion}
This template demonstrates basic \LaTeX{} features including mathematical formulas, tables, and lists. The bibliography system is now set up to work seamlessly with Bibby AI's citation management. Feel free to modify it for your needs!

\bibliographystyle{plain}
\bibliography{cite}

\end{document}\documentclass{article}
\usepackage[utf8]{inputenc}
\usepackage{booktabs}
\usepackage{amsmath}
\usepackage{natbib}

\title{My Document}
\author{Your Name}
\date{\today}

\begin{document}

\maketitle

\section{Introduction}
This is a sample document demonstrating various \LaTeX{} features including mathematical equations, tables, and formatting. You can cite papers using the citation search feature, and they will be automatically added to your \texttt{cite.bib} file.

\section{Mathematical Expressions}
Here are some example equations:

\subsection{Quadratic Formula}
The quadratic formula is given by:
\[x = \frac{-b \pm \sqrt{b^2 - 4ac}}{2a}\]

\subsection{Matrix Example}
A sample matrix:
\[
\begin{pmatrix}
a_{11} & a_{12} & a_{13} \\
a_{21} & a_{22} & a_{23} \\
a_{31} & a_{32} & a_{33}
\end{pmatrix}
\]

\section{Tables}
Below is a sample table showing data:

\begin{table}[h]
\centering
\begin{tabular}{@{}lrrr@{}}
\toprule
\textbf{Category} & \textbf{Value 1} & \textbf{Value 2} & \textbf{Total} \\
\midrule
Item A & 10.5 & 20.3 & 30.8 \\
Item B & 15.2 & 25.1 & 40.3 \\
Item C & 22.1 & 12.7 & 34.8 \\
\bottomrule
\end{tabular}
\caption{Sample Data Table}
\label{tab:sample}
\end{table}

\section{Lists}
\subsection{Bullet Points}
\begin{itemize}
  \item First item
  \item Second item
  \item Third item with sub-items:
    \begin{itemize}
      \item Sub-item 1
      \item Sub-item 2
    \end{itemize}
\end{itemize}

\subsection{Numbered List}
\begin{enumerate}
  \item First step
  \item Second step
  \item Third step
\end{enumerate}

\section{Citations and References}
When you use the citation search feature, references will be automatically managed. For example, you might cite a paper like this: \cite{example2024}. All citations will be automatically added to your \texttt{cite.bib} file.

\section{Conclusion}
This template demonstrates basic \LaTeX{} features including mathematical formulas, tables, and lists. The bibliography system is now set up to work seamlessly with Bibby AI's citation management. Feel free to modify it for your needs!

\bibliographystyle{plain}
\bibliography{cite}

\end{document}\documentclass{article}
\usepackage[utf8]{inputenc}
\usepackage{booktabs}
\usepackage{amsmath}
\usepackage{natbib}

\title{My Document}
\author{Your Name}
\date{\today}

\begin{document}

\maketitle

\section{Introduction}
This is a sample document demonstrating various \LaTeX{} features including mathematical equations, tables, and formatting. You can cite papers using the citation search feature, and they will be automatically added to your \texttt{cite.bib} file.

\section{Mathematical Expressions}
Here are some example equations:

\subsection{Quadratic Formula}
The quadratic formula is given by:
\[x = \frac{-b \pm \sqrt{b^2 - 4ac}}{2a}\]

\subsection{Matrix Example}
A sample matrix:
\[
\begin{pmatrix}
a_{11} & a_{12} & a_{13} \\
a_{21} & a_{22} & a_{23} \\
a_{31} & a_{32} & a_{33}
\end{pmatrix}
\]

\section{Tables}
Below is a sample table showing data:

\begin{table}[h]
\centering
\begin{tabular}{@{}lrrr@{}}
\toprule
\textbf{Category} & \textbf{Value 1} & \textbf{Value 2} & \textbf{Total} \\
\midrule
Item A & 10.5 & 20.3 & 30.8 \\
Item B & 15.2 & 25.1 & 40.3 \\
Item C & 22.1 & 12.7 & 34.8 \\
\bottomrule
\end{tabular}
\caption{Sample Data Table}
\label{tab:sample}
\end{table}

\section{Lists}
\subsection{Bullet Points}
\begin{itemize}
  \item First item
  \item Second item
  \item Third item with sub-items:
    \begin{itemize}
      \item Sub-item 1
      \item Sub-item 2
    \end{itemize}
\end{itemize}

\subsection{Numbered List}
\begin{enumerate}
  \item First step
  \item Second step
  \item Third step
\end{enumerate}

\section{Citations and References}
When you use the citation search feature, references will be automatically managed. For example, you might cite a paper like this: \cite{example2024}. All citations will be automatically added to your \texttt{cite.bib} file.

\section{Conclusion}
This template demonstrates basic \LaTeX{} features including mathematical formulas, tables, and lists. The bibliography system is now set up to work seamlessly with Bibby AI's citation management. Feel free to modify it for your needs!

\bibliographystyle{plain}
\bibliography{cite}

\end{document}\documentclass{article}
\usepackage[utf8]{inputenc}
\usepackage{booktabs}
\usepackage{amsmath}
\usepackage{natbib}

\title{My Document}
\author{Your Name}
\date{\today}

\begin{document}

\maketitle

\section{Introduction}
This is a sample document demonstrating various \LaTeX{} features including mathematical equations, tables, and formatting. You can cite papers using the citation search feature, and they will be automatically added to your \texttt{cite.bib} file.

\section{Mathematical Expressions}
Here are some example equations:

\subsection{Quadratic Formula}
The quadratic formula is given by:
\[x = \frac{-b \pm \sqrt{b^2 - 4ac}}{2a}\]

\subsection{Matrix Example}
A sample matrix:
\[
\begin{pmatrix}
a_{11} & a_{12} & a_{13} \\
a_{21} & a_{22} & a_{23} \\
a_{31} & a_{32} & a_{33}
\end{pmatrix}
\]

\section{Tables}
Below is a sample table showing data:

\begin{table}[h]
\centering
\begin{tabular}{@{}lrrr@{}}
\toprule
\textbf{Category} & \textbf{Value 1} & \textbf{Value 2} & \textbf{Total} \\
\midrule
Item A & 10.5 & 20.3 & 30.8 \\
Item B & 15.2 & 25.1 & 40.3 \\
Item C & 22.1 & 12.7 & 34.8 \\
\bottomrule
\end{tabular}
\caption{Sample Data Table}
\label{tab:sample}
\end{table}

\section{Lists}
\subsection{Bullet Points}
\begin{itemize}
  \item First item
  \item Second item
  \item Third item with sub-items:
    \begin{itemize}
      \item Sub-item 1
      \item Sub-item 2
    \end{itemize}
\end{itemize}

\subsection{Numbered List}
\begin{enumerate}
  \item First step
  \item Second step
  \item Third step
\end{enumerate}

\section{Citations and References}
When you use the citation search feature, references will be automatically managed. For example, you might cite a paper like this: \cite{example2024}. All citations will be automatically added to your \texttt{cite.bib} file.

\section{Conclusion}
This template demonstrates basic \LaTeX{} features including mathematical formulas, tables, and lists. The bibliography system is now set up to work seamlessly with Bibby AI's citation management. Feel free to modify it for your needs!

\bibliographystyle{plain}
\bibliography{cite}

\end{document}\documentclass{article}
\usepackage[utf8]{inputenc}
\usepackage{booktabs}
\usepackage{amsmath}
\usepackage{natbib}

\title{My Document}
\author{Your Name}
\date{\today}

\begin{document}

\maketitle

\section{Introduction}
This is a sample document demonstrating various \LaTeX{} features including mathematical equations, tables, and formatting. You can cite papers using the citation search feature, and they will be automatically added to your \texttt{cite.bib} file.

\section{Mathematical Expressions}
Here are some example equations:

\subsection{Quadratic Formula}
The quadratic formula is given by:
\[x = \frac{-b \pm \sqrt{b^2 - 4ac}}{2a}\]

\subsection{Matrix Example}
A sample matrix:
\[
\begin{pmatrix}
a_{11} & a_{12} & a_{13} \\
a_{21} & a_{22} & a_{23} \\
a_{31} & a_{32} & a_{33}
\end{pmatrix}
\]

\section{Tables}
Below is a sample table showing data:

\begin{table}[h]
\centering
\begin{tabular}{@{}lrrr@{}}
\toprule
\textbf{Category} & \textbf{Value 1} & \textbf{Value 2} & \textbf{Total} \\
\midrule
Item A & 10.5 & 20.3 & 30.8 \\
Item B & 15.2 & 25.1 & 40.3 \\
Item C & 22.1 & 12.7 & 34.8 \\
\bottomrule
\end{tabular}
\caption{Sample Data Table}
\label{tab:sample}
\end{table}

\section{Lists}
\subsection{Bullet Points}
\begin{itemize}
  \item First item
  \item Second item
  \item Third item with sub-items:
    \begin{itemize}
      \item Sub-item 1
      \item Sub-item 2
    \end{itemize}
\end{itemize}

\subsection{Numbered List}
\begin{enumerate}
  \item First step
  \item Second step
  \item Third step
\end{enumerate}

\section{Citations and References}
When you use the citation search feature, references will be automatically managed. For example, you might cite a paper like this: \cite{example2024}. All citations will be automatically added to your \texttt{cite.bib} file.

\section{Conclusion}
This template demonstrates basic \LaTeX{} features including mathematical formulas, tables, and lists. The bibliography system is now set up to work seamlessly with Bibby AI's citation management. Feel free to modify it for your needs!

\bibliographystyle{plain}
\bibliography{cite}

\end{document}\documentclass{article}
\usepackage[utf8]{inputenc}
\usepackage{booktabs}
\usepackage{amsmath}
\usepackage{natbib}

\title{My Document}
\author{Your Name}
\date{\today}

\begin{document}

\maketitle

\section{Introduction}
This is a sample document demonstrating various \LaTeX{} features including mathematical equations, tables, and formatting. You can cite papers using the citation search feature, and they will be automatically added to your \texttt{cite.bib} file.

\section{Mathematical Expressions}
Here are some example equations:

\subsection{Quadratic Formula}
The quadratic formula is given by:
\[x = \frac{-b \pm \sqrt{b^2 - 4ac}}{2a}\]

\subsection{Matrix Example}
A sample matrix:
\[
\begin{pmatrix}
a_{11} & a_{12} & a_{13} \\
a_{21} & a_{22} & a_{23} \\
a_{31} & a_{32} & a_{33}
\end{pmatrix}
\]

\section{Tables}
Below is a sample table showing data:

\begin{table}[h]
\centering
\begin{tabular}{@{}lrrr@{}}
\toprule
\textbf{Category} & \textbf{Value 1} & \textbf{Value 2} & \textbf{Total} \\
\midrule
Item A & 10.5 & 20.3 & 30.8 \\
Item B & 15.2 & 25.1 & 40.3 \\
Item C & 22.1 & 12.7 & 34.8 \\
\bottomrule
\end{tabular}
\caption{Sample Data Table}
\label{tab:sample}
\end{table}

\section{Lists}
\subsection{Bullet Points}
\begin{itemize}
  \item First item
  \item Second item
  \item Third item with sub-items:
    \begin{itemize}
      \item Sub-item 1
      \item Sub-item 2
    \end{itemize}
\end{itemize}

\subsection{Numbered List}
\begin{enumerate}
  \item First step
  \item Second step
  \item Third step
\end{enumerate}

\section{Citations and References}
When you use the citation search feature, references will be automatically managed. For example, you might cite a paper like this: \cite{example2024}. All citations will be automatically added to your \texttt{cite.bib} file.

\section{Conclusion}
This template demonstrates basic \LaTeX{} features including mathematical formulas, tables, and lists. The bibliography system is now set up to work seamlessly with Bibby AI's citation management. Feel free to modify it for your needs!

\bibliographystyle{plain}
\bibliography{cite}

\end{document}\documentclass{article}
\usepackage[utf8]{inputenc}
\usepackage{booktabs}
\usepackage{amsmath}
\usepackage{natbib}

\title{My Document}
\author{Your Name}
\date{\today}

\begin{document}

\maketitle

\section{Introduction}
This is a sample document demonstrating various \LaTeX{} features including mathematical equations, tables, and formatting. You can cite papers using the citation search feature, and they will be automatically added to your \texttt{cite.bib} file.

\section{Mathematical Expressions}
Here are some example equations:

\subsection{Quadratic Formula}
The quadratic formula is given by:
\[x = \frac{-b \pm \sqrt{b^2 - 4ac}}{2a}\]

\subsection{Matrix Example}
A sample matrix:
\[
\begin{pmatrix}
a_{11} & a_{12} & a_{13} \\
a_{21} & a_{22} & a_{23} \\
a_{31} & a_{32} & a_{33}
\end{pmatrix}
\]

\section{Tables}
Below is a sample table showing data:

\begin{table}[h]
\centering
\begin{tabular}{@{}lrrr@{}}
\toprule
\textbf{Category} & \textbf{Value 1} & \textbf{Value 2} & \textbf{Total} \\
\midrule
Item A & 10.5 & 20.3 & 30.8 \\
Item B & 15.2 & 25.1 & 40.3 \\
Item C & 22.1 & 12.7 & 34.8 \\
\bottomrule
\end{tabular}
\caption{Sample Data Table}
\label{tab:sample}
\end{table}

\section{Lists}
\subsection{Bullet Points}
\begin{itemize}
  \item First item
  \item Second item
  \item Third item with sub-items:
    \begin{itemize}
      \item Sub-item 1
      \item Sub-item 2
    \end{itemize}
\end{itemize}

\subsection{Numbered List}
\begin{enumerate}
  \item First step
  \item Second step
  \item Third step
\end{enumerate}

\section{Citations and References}
When you use the citation search feature, references will be automatically managed. For example, you might cite a paper like this: \cite{example2024}. All citations will be automatically added to your \texttt{cite.bib} file.

\section{Conclusion}
This template demonstrates basic \LaTeX{} features including mathematical formulas, tables, and lists. The bibliography system is now set up to work seamlessly with Bibby AI's citation management. Feel free to modify it for your needs!

\bibliographystyle{plain}
\bibliography{cite}

\end{document}\documentclass{article}
\usepackage[utf8]{inputenc}
\usepackage{booktabs}
\usepackage{amsmath}
\usepackage{natbib}

\title{My Document}
\author{Your Name}
\date{\today}

\begin{document}

\maketitle

\section{Introduction}
This is a sample document demonstrating various \LaTeX{} features including mathematical equations, tables, and formatting. You can cite papers using the citation search feature, and they will be automatically added to your \texttt{cite.bib} file.

\section{Mathematical Expressions}
Here are some example equations:

\subsection{Quadratic Formula}
The quadratic formula is given by:
\[x = \frac{-b \pm \sqrt{b^2 - 4ac}}{2a}\]

\subsection{Matrix Example}
A sample matrix:
\[
\begin{pmatrix}
a_{11} & a_{12} & a_{13} \\
a_{21} & a_{22} & a_{23} \\
a_{31} & a_{32} & a_{33}
\end{pmatrix}
\]

\section{Tables}
Below is a sample table showing data:

\begin{table}[h]
\centering
\begin{tabular}{@{}lrrr@{}}
\toprule
\textbf{Category} & \textbf{Value 1} & \textbf{Value 2} & \textbf{Total} \\
\midrule
Item A & 10.5 & 20.3 & 30.8 \\
Item B & 15.2 & 25.1 & 40.3 \\
Item C & 22.1 & 12.7 & 34.8 \\
\bottomrule
\end{tabular}
\caption{Sample Data Table}
\label{tab:sample}
\end{table}

\section{Lists}
\subsection{Bullet Points}
\begin{itemize}
  \item First item
  \item Second item
  \item Third item with sub-items:
    \begin{itemize}
      \item Sub-item 1
      \item Sub-item 2
    \end{itemize}
\end{itemize}

\subsection{Numbered List}
\begin{enumerate}
  \item First step
  \item Second step
  \item Third step
\end{enumerate}

\section{Citations and References}
When you use the citation search feature, references will be automatically managed. For example, you might cite a paper like this: \cite{example2024}. All citations will be automatically added to your \texttt{cite.bib} file.

\section{Conclusion}
This template demonstrates basic \LaTeX{} features including mathematical formulas, tables, and lists. The bibliography system is now set up to work seamlessly with Bibby AI's citation management. Feel free to modify it for your needs!

\bibliographystyle{plain}
\bibliography{cite}

\end{document}\documentclass{article}
\usepackage[utf8]{inputenc}
\usepackage{booktabs}
\usepackage{amsmath}
\usepackage{natbib}

\title{My Document}
\author{Your Name}
\date{\today}

\begin{document}

\maketitle

\section{Introduction}
This is a sample document demonstrating various \LaTeX{} features including mathematical equations, tables, and formatting. You can cite papers using the citation search feature, and they will be automatically added to your \texttt{cite.bib} file.

\section{Mathematical Expressions}
Here are some example equations:

\subsection{Quadratic Formula}
The quadratic formula is given by:
\[x = \frac{-b \pm \sqrt{b^2 - 4ac}}{2a}\]

\subsection{Matrix Example}
A sample matrix:
\[
\begin{pmatrix}
a_{11} & a_{12} & a_{13} \\
a_{21} & a_{22} & a_{23} \\
a_{31} & a_{32} & a_{33}
\end{pmatrix}
\]

\section{Tables}
Below is a sample table showing data:

\begin{table}[h]
\centering
\begin{tabular}{@{}lrrr@{}}
\toprule
\textbf{Category} & \textbf{Value 1} & \textbf{Value 2} & \textbf{Total} \\
\midrule
Item A & 10.5 & 20.3 & 30.8 \\
Item B & 15.2 & 25.1 & 40.3 \\
Item C & 22.1 & 12.7 & 34.8 \\
\bottomrule
\end{tabular}
\caption{Sample Data Table}
\label{tab:sample}
\end{table}

\section{Lists}
\subsection{Bullet Points}
\begin{itemize}
  \item First item
  \item Second item
  \item Third item with sub-items:
    \begin{itemize}
      \item Sub-item 1
      \item Sub-item 2
    \end{itemize}
\end{itemize}

\subsection{Numbered List}
\begin{enumerate}
  \item First step
  \item Second step
  \item Third step
\end{enumerate}

\section{Citations and References}
When you use the citation search feature, references will be automatically managed. For example, you might cite a paper like this: \cite{example2024}. All citations will be automatically added to your \texttt{cite.bib} file.

\section{Conclusion}
This template demonstrates basic \LaTeX{} features including mathematical formulas, tables, and lists. The bibliography system is now set up to work seamlessly with Bibby AI's citation management. Feel free to modify it for your needs!

\bibliographystyle{plain}
\bibliography{cite}

\end{document}\documentclass{article}
\usepackage[utf8]{inputenc}
\usepackage{booktabs}
\usepackage{amsmath}
\usepackage{natbib}

\title{My Document}
\author{Your Name}
\date{\today}

\begin{document}

\maketitle

\section{Introduction}
This is a sample document demonstrating various \LaTeX{} features including mathematical equations, tables, and formatting. You can cite papers using the citation search feature, and they will be automatically added to your \texttt{cite.bib} file.

\section{Mathematical Expressions}
Here are some example equations:

\subsection{Quadratic Formula}
The quadratic formula is given by:
\[x = \frac{-b \pm \sqrt{b^2 - 4ac}}{2a}\]

\subsection{Matrix Example}
A sample matrix:
\[
\begin{pmatrix}
a_{11} & a_{12} & a_{13} \\
a_{21} & a_{22} & a_{23} \\
a_{31} & a_{32} & a_{33}
\end{pmatrix}
\]

\section{Tables}
Below is a sample table showing data:

\begin{table}[h]
\centering
\begin{tabular}{@{}lrrr@{}}
\toprule
\textbf{Category} & \textbf{Value 1} & \textbf{Value 2} & \textbf{Total} \\
\midrule
Item A & 10.5 & 20.3 & 30.8 \\
Item B & 15.2 & 25.1 & 40.3 \\
Item C & 22.1 & 12.7 & 34.8 \\
\bottomrule
\end{tabular}
\caption{Sample Data Table}
\label{tab:sample}
\end{table}

\section{Lists}
\subsection{Bullet Points}
\begin{itemize}
  \item First item
  \item Second item
  \item Third item with sub-items:
    \begin{itemize}
      \item Sub-item 1
      \item Sub-item 2
    \end{itemize}
\end{itemize}

\subsection{Numbered List}
\begin{enumerate}
  \item First step
  \item Second step
  \item Third step
\end{enumerate}

\section{Citations and References}
When you use the citation search feature, references will be automatically managed. For example, you might cite a paper like this: \cite{example2024}. All citations will be automatically added to your \texttt{cite.bib} file.

\section{Conclusion}
This template demonstrates basic \LaTeX{} features including mathematical formulas, tables, and lists. The bibliography system is now set up to work seamlessly with Bibby AI's citation management. Feel free to modify it for your needs!

\bibliographystyle{plain}
\bibliography{cite}

\end{document}\documentclass{article}
\usepackage[utf8]{inputenc}
\usepackage{booktabs}
\usepackage{amsmath}
\usepackage{natbib}

\title{My Document}
\author{Your Name}
\date{\today}

\begin{document}

\maketitle

\section{Introduction}
This is a sample document demonstrating various \LaTeX{} features including mathematical equations, tables, and formatting. You can cite papers using the citation search feature, and they will be automatically added to your \texttt{cite.bib} file.

\section{Mathematical Expressions}
Here are some example equations:

\subsection{Quadratic Formula}
The quadratic formula is given by:
\[x = \frac{-b \pm \sqrt{b^2 - 4ac}}{2a}\]

\subsection{Matrix Example}
A sample matrix:
\[
\begin{pmatrix}
a_{11} & a_{12} & a_{13} \\
a_{21} & a_{22} & a_{23} \\
a_{31} & a_{32} & a_{33}
\end{pmatrix}
\]

\section{Tables}
Below is a sample table showing data:

\begin{table}[h]
\centering
\begin{tabular}{@{}lrrr@{}}
\toprule
\textbf{Category} & \textbf{Value 1} & \textbf{Value 2} & \textbf{Total} \\
\midrule
Item A & 10.5 & 20.3 & 30.8 \\
Item B & 15.2 & 25.1 & 40.3 \\
Item C & 22.1 & 12.7 & 34.8 \\
\bottomrule
\end{tabular}
\caption{Sample Data Table}
\label{tab:sample}
\end{table}

\section{Lists}
\subsection{Bullet Points}
\begin{itemize}
  \item First item
  \item Second item
  \item Third item with sub-items:
    \begin{itemize}
      \item Sub-item 1
      \item Sub-item 2
    \end{itemize}
\end{itemize}

\subsection{Numbered List}
\begin{enumerate}
  \item First step
  \item Second step
  \item Third step
\end{enumerate}

\section{Citations and References}
When you use the citation search feature, references will be automatically managed. For example, you might cite a paper like this: \cite{example2024}. All citations will be automatically added to your \texttt{cite.bib} file.

\section{Conclusion}
This template demonstrates basic \LaTeX{} features including mathematical formulas, tables, and lists. The bibliography system is now set up to work seamlessly with Bibby AI's citation management. Feel free to modify it for your needs!

\bibliographystyle{plain}
\bibliography{cite}

\end{document}\documentclass{article}
\usepackage[utf8]{inputenc}
\usepackage{booktabs}
\usepackage{amsmath}
\usepackage{natbib}

\title{My Document}
\author{Your Name}
\date{\today}

\begin{document}

\maketitle

\section{Introduction}
This is a sample document demonstrating various \LaTeX{} features including mathematical equations, tables, and formatting. You can cite papers using the citation search feature, and they will be automatically added to your \texttt{cite.bib} file.

\section{Mathematical Expressions}
Here are some example equations:

\subsection{Quadratic Formula}
The quadratic formula is given by:
\[x = \frac{-b \pm \sqrt{b^2 - 4ac}}{2a}\]

\subsection{Matrix Example}
A sample matrix:
\[
\begin{pmatrix}
a_{11} & a_{12} & a_{13} \\
a_{21} & a_{22} & a_{23} \\
a_{31} & a_{32} & a_{33}
\end{pmatrix}
\]

\section{Tables}
Below is a sample table showing data:

\begin{table}[h]
\centering
\begin{tabular}{@{}lrrr@{}}
\toprule
\textbf{Category} & \textbf{Value 1} & \textbf{Value 2} & \textbf{Total} \\
\midrule
Item A & 10.5 & 20.3 & 30.8 \\
Item B & 15.2 & 25.1 & 40.3 \\
Item C & 22.1 & 12.7 & 34.8 \\
\bottomrule
\end{tabular}
\caption{Sample Data Table}
\label{tab:sample}
\end{table}

\section{Lists}
\subsection{Bullet Points}
\begin{itemize}
  \item First item
  \item Second item
  \item Third item with sub-items:
    \begin{itemize}
      \item Sub-item 1
      \item Sub-item 2
    \end{itemize}
\end{itemize}

\subsection{Numbered List}
\begin{enumerate}
  \item First step
  \item Second step
  \item Third step
\end{enumerate}

\section{Citations and References}
When you use the citation search feature, references will be automatically managed. For example, you might cite a paper like this: \cite{example2024}. All citations will be automatically added to your \texttt{cite.bib} file.

\section{Conclusion}
This template demonstrates basic \LaTeX{} features including mathematical formulas, tables, and lists. The bibliography system is now set up to work seamlessly with Bibby AI's citation management. Feel free to modify it for your needs!

\bibliographystyle{plain}
\bibliography{cite}

\end{document}\documentclass{article}
\usepackage[utf8]{inputenc}
\usepackage{booktabs}
\usepackage{amsmath}
\usepackage{natbib}

\title{My Document}
\author{Your Name}
\date{\today}

\begin{document}

\maketitle

\section{Introduction}
This is a sample document demonstrating various \LaTeX{} features including mathematical equations, tables, and formatting. You can cite papers using the citation search feature, and they will be automatically added to your \texttt{cite.bib} file.

\section{Mathematical Expressions}
Here are some example equations:

\subsection{Quadratic Formula}
The quadratic formula is given by:
\[x = \frac{-b \pm \sqrt{b^2 - 4ac}}{2a}\]

\subsection{Matrix Example}
A sample matrix:
\[
\begin{pmatrix}
a_{11} & a_{12} & a_{13} \\
a_{21} & a_{22} & a_{23} \\
a_{31} & a_{32} & a_{33}
\end{pmatrix}
\]

\section{Tables}
Below is a sample table showing data:

\begin{table}[h]
\centering
\begin{tabular}{@{}lrrr@{}}
\toprule
\textbf{Category} & \textbf{Value 1} & \textbf{Value 2} & \textbf{Total} \\
\midrule
Item A & 10.5 & 20.3 & 30.8 \\
Item B & 15.2 & 25.1 & 40.3 \\
Item C & 22.1 & 12.7 & 34.8 \\
\bottomrule
\end{tabular}
\caption{Sample Data Table}
\label{tab:sample}
\end{table}

\section{Lists}
\subsection{Bullet Points}
\begin{itemize}
  \item First item
  \item Second item
  \item Third item with sub-items:
    \begin{itemize}
      \item Sub-item 1
      \item Sub-item 2
    \end{itemize}
\end{itemize}

\subsection{Numbered List}
\begin{enumerate}
  \item First step
  \item Second step
  \item Third step
\end{enumerate}

\section{Citations and References}
When you use the citation search feature, references will be automatically managed. For example, you might cite a paper like this: \cite{example2024}. All citations will be automatically added to your \texttt{cite.bib} file.

\section{Conclusion}
This template demonstrates basic \LaTeX{} features including mathematical formulas, tables, and lists. The bibliography system is now set up to work seamlessly with Bibby AI's citation management. Feel free to modify it for your needs!

\bibliographystyle{plain}
\bibliography{cite}

\end{document}\documentclass{article}
\usepackage[utf8]{inputenc}
\usepackage{booktabs}
\usepackage{amsmath}
\usepackage{natbib}

\title{My Document}
\author{Your Name}
\date{\today}

\begin{document}

\maketitle

\section{Introduction}
This is a sample document demonstrating various \LaTeX{} features including mathematical equations, tables, and formatting. You can cite papers using the citation search feature, and they will be automatically added to your \texttt{cite.bib} file.

\section{Mathematical Expressions}
Here are some example equations:

\subsection{Quadratic Formula}
The quadratic formula is given by:
\[x = \frac{-b \pm \sqrt{b^2 - 4ac}}{2a}\]

\subsection{Matrix Example}
A sample matrix:
\[
\begin{pmatrix}
a_{11} & a_{12} & a_{13} \\
a_{21} & a_{22} & a_{23} \\
a_{31} & a_{32} & a_{33}
\end{pmatrix}
\]

\section{Tables}
Below is a sample table showing data:

\begin{table}[h]
\centering
\begin{tabular}{@{}lrrr@{}}
\toprule
\textbf{Category} & \textbf{Value 1} & \textbf{Value 2} & \textbf{Total} \\
\midrule
Item A & 10.5 & 20.3 & 30.8 \\
Item B & 15.2 & 25.1 & 40.3 \\
Item C & 22.1 & 12.7 & 34.8 \\
\bottomrule
\end{tabular}
\caption{Sample Data Table}
\label{tab:sample}
\end{table}

\section{Lists}
\subsection{Bullet Points}
\begin{itemize}
  \item First item
  \item Second item
  \item Third item with sub-items:
    \begin{itemize}
      \item Sub-item 1
      \item Sub-item 2
    \end{itemize}
\end{itemize}

\subsection{Numbered List}
\begin{enumerate}
  \item First step
  \item Second step
  \item Third step
\end{enumerate}

\section{Citations and References}
When you use the citation search feature, references will be automatically managed. For example, you might cite a paper like this: \cite{example2024}. All citations will be automatically added to your \texttt{cite.bib} file.

\section{Conclusion}
This template demonstrates basic \LaTeX{} features including mathematical formulas, tables, and lists. The bibliography system is now set up to work seamlessly with Bibby AI's citation management. Feel free to modify it for your needs!

\bibliographystyle{plain}
\bibliography{cite}

\end{document}\documentclass{article}
\usepackage[utf8]{inputenc}
\usepackage{booktabs}
\usepackage{amsmath}
\usepackage{natbib}

\title{My Document}
\author{Your Name}
\date{\today}

\begin{document}

\maketitle

\section{Introduction}
This is a sample document demonstrating various \LaTeX{} features including mathematical equations, tables, and formatting. You can cite papers using the citation search feature, and they will be automatically added to your \texttt{cite.bib} file.

\section{Mathematical Expressions}
Here are some example equations:

\subsection{Quadratic Formula}
The quadratic formula is given by:
\[x = \frac{-b \pm \sqrt{b^2 - 4ac}}{2a}\]

\subsection{Matrix Example}
A sample matrix:
\[
\begin{pmatrix}
a_{11} & a_{12} & a_{13} \\
a_{21} & a_{22} & a_{23} \\
a_{31} & a_{32} & a_{33}
\end{pmatrix}
\]

\section{Tables}
Below is a sample table showing data:

\begin{table}[h]
\centering
\begin{tabular}{@{}lrrr@{}}
\toprule
\textbf{Category} & \textbf{Value 1} & \textbf{Value 2} & \textbf{Total} \\
\midrule
Item A & 10.5 & 20.3 & 30.8 \\
Item B & 15.2 & 25.1 & 40.3 \\
Item C & 22.1 & 12.7 & 34.8 \\
\bottomrule
\end{tabular}
\caption{Sample Data Table}
\label{tab:sample}
\end{table}

\section{Lists}
\subsection{Bullet Points}
\begin{itemize}
  \item First item
  \item Second item
  \item Third item with sub-items:
    \begin{itemize}
      \item Sub-item 1
      \item Sub-item 2
    \end{itemize}
\end{itemize}

\subsection{Numbered List}
\begin{enumerate}
  \item First step
  \item Second step
  \item Third step
\end{enumerate}

\section{Citations and References}
When you use the citation search feature, references will be automatically managed. For example, you might cite a paper like this: \cite{example2024}. All citations will be automatically added to your \texttt{cite.bib} file.

\section{Conclusion}
This template demonstrates basic \LaTeX{} features including mathematical formulas, tables, and lists. The bibliography system is now set up to work seamlessly with Bibby AI's citation management. Feel free to modify it for your needs!

\bibliographystyle{plain}
\bibliography{cite}

\end{document}\documentclass{article}
\usepackage[utf8]{inputenc}
\usepackage{booktabs}
\usepackage{amsmath}
\usepackage{natbib}

\title{My Document}
\author{Your Name}
\date{\today}

\begin{document}

\maketitle

\section{Introduction}
This is a sample document demonstrating various \LaTeX{} features including mathematical equations, tables, and formatting. You can cite papers using the citation search feature, and they will be automatically added to your \texttt{cite.bib} file.

\section{Mathematical Expressions}
Here are some example equations:

\subsection{Quadratic Formula}
The quadratic formula is given by:
\[x = \frac{-b \pm \sqrt{b^2 - 4ac}}{2a}\]

\subsection{Matrix Example}
A sample matrix:
\[
\begin{pmatrix}
a_{11} & a_{12} & a_{13} \\
a_{21} & a_{22} & a_{23} \\
a_{31} & a_{32} & a_{33}
\end{pmatrix}
\]

\section{Tables}
Below is a sample table showing data:

\begin{table}[h]
\centering
\begin{tabular}{@{}lrrr@{}}
\toprule
\textbf{Category} & \textbf{Value 1} & \textbf{Value 2} & \textbf{Total} \\
\midrule
Item A & 10.5 & 20.3 & 30.8 \\
Item B & 15.2 & 25.1 & 40.3 \\
Item C & 22.1 & 12.7 & 34.8 \\
\bottomrule
\end{tabular}
\caption{Sample Data Table}
\label{tab:sample}
\end{table}

\section{Lists}
\subsection{Bullet Points}
\begin{itemize}
  \item First item
  \item Second item
  \item Third item with sub-items:
    \begin{itemize}
      \item Sub-item 1
      \item Sub-item 2
    \end{itemize}
\end{itemize}

\subsection{Numbered List}
\begin{enumerate}
  \item First step
  \item Second step
  \item Third step
\end{enumerate}

\section{Citations and References}
When you use the citation search feature, references will be automatically managed. For example, you might cite a paper like this: \cite{example2024}. All citations will be automatically added to your \texttt{cite.bib} file.

\section{Conclusion}
This template demonstrates basic \LaTeX{} features including mathematical formulas, tables, and lists. The bibliography system is now set up to work seamlessly with Bibby AI's citation management. Feel free to modify it for your needs!

\bibliographystyle{plain}
\bibliography{cite}

\end{document}\documentclass{article}
\usepackage[utf8]{inputenc}
\usepackage{booktabs}
\usepackage{amsmath}
\usepackage{natbib}

\title{My Document}
\author{Your Name}
\date{\today}

\begin{document}

\maketitle

\section{Introduction}
This is a sample document demonstrating various \LaTeX{} features including mathematical equations, tables, and formatting. You can cite papers using the citation search feature, and they will be automatically added to your \texttt{cite.bib} file.

\section{Mathematical Expressions}
Here are some example equations:

\subsection{Quadratic Formula}
The quadratic formula is given by:
\[x = \frac{-b \pm \sqrt{b^2 - 4ac}}{2a}\]

\subsection{Matrix Example}
A sample matrix:
\[
\begin{pmatrix}
a_{11} & a_{12} & a_{13} \\
a_{21} & a_{22} & a_{23} \\
a_{31} & a_{32} & a_{33}
\end{pmatrix}
\]

\section{Tables}
Below is a sample table showing data:

\begin{table}[h]
\centering
\begin{tabular}{@{}lrrr@{}}
\toprule
\textbf{Category} & \textbf{Value 1} & \textbf{Value 2} & \textbf{Total} \\
\midrule
Item A & 10.5 & 20.3 & 30.8 \\
Item B & 15.2 & 25.1 & 40.3 \\
Item C & 22.1 & 12.7 & 34.8 \\
\bottomrule
\end{tabular}
\caption{Sample Data Table}
\label{tab:sample}
\end{table}

\section{Lists}
\subsection{Bullet Points}
\begin{itemize}
  \item First item
  \item Second item
  \item Third item with sub-items:
    \begin{itemize}
      \item Sub-item 1
      \item Sub-item 2
    \end{itemize}
\end{itemize}

\subsection{Numbered List}
\begin{enumerate}
  \item First step
  \item Second step
  \item Third step
\end{enumerate}

\section{Citations and References}
When you use the citation search feature, references will be automatically managed. For example, you might cite a paper like this: \cite{example2024}. All citations will be automatically added to your \texttt{cite.bib} file.

\section{Conclusion}
This template demonstrates basic \LaTeX{} features including mathematical formulas, tables, and lists. The bibliography system is now set up to work seamlessly with Bibby AI's citation management. Feel free to modify it for your needs!

\bibliographystyle{plain}
\bibliography{cite}

\end{document}\documentclass{article}
\usepackage[utf8]{inputenc}
\usepackage{booktabs}
\usepackage{amsmath}
\usepackage{natbib}

\title{My Document}
\author{Your Name}
\date{\today}

\begin{document}

\maketitle

\section{Introduction}
This is a sample document demonstrating various \LaTeX{} features including mathematical equations, tables, and formatting. You can cite papers using the citation search feature, and they will be automatically added to your \texttt{cite.bib} file.

\section{Mathematical Expressions}
Here are some example equations:

\subsection{Quadratic Formula}
The quadratic formula is given by:
\[x = \frac{-b \pm \sqrt{b^2 - 4ac}}{2a}\]

\subsection{Matrix Example}
A sample matrix:
\[
\begin{pmatrix}
a_{11} & a_{12} & a_{13} \\
a_{21} & a_{22} & a_{23} \\
a_{31} & a_{32} & a_{33}
\end{pmatrix}
\]

\section{Tables}
Below is a sample table showing data:

\begin{table}[h]
\centering
\begin{tabular}{@{}lrrr@{}}
\toprule
\textbf{Category} & \textbf{Value 1} & \textbf{Value 2} & \textbf{Total} \\
\midrule
Item A & 10.5 & 20.3 & 30.8 \\
Item B & 15.2 & 25.1 & 40.3 \\
Item C & 22.1 & 12.7 & 34.8 \\
\bottomrule
\end{tabular}
\caption{Sample Data Table}
\label{tab:sample}
\end{table}

\section{Lists}
\subsection{Bullet Points}
\begin{itemize}
  \item First item
  \item Second item
  \item Third item with sub-items:
    \begin{itemize}
      \item Sub-item 1
      \item Sub-item 2
    \end{itemize}
\end{itemize}

\subsection{Numbered List}
\begin{enumerate}
  \item First step
  \item Second step
  \item Third step
\end{enumerate}

\section{Citations and References}
When you use the citation search feature, references will be automatically managed. For example, you might cite a paper like this: \cite{example2024}. All citations will be automatically added to your \texttt{cite.bib} file.

\section{Conclusion}
This template demonstrates basic \LaTeX{} features including mathematical formulas, tables, and lists. The bibliography system is now set up to work seamlessly with Bibby AI's citation management. Feel free to modify it for your needs!

\bibliographystyle{plain}
\bibliography{cite}

\end{document}\documentclass{article}
\usepackage[utf8]{inputenc}
\usepackage{booktabs}
\usepackage{amsmath}
\usepackage{natbib}

\title{My Document}
\author{Your Name}
\date{\today}

\begin{document}

\maketitle

\section{Introduction}
This is a sample document demonstrating various \LaTeX{} features including mathematical equations, tables, and formatting. You can cite papers using the citation search feature, and they will be automatically added to your \texttt{cite.bib} file.

\section{Mathematical Expressions}
Here are some example equations:

\subsection{Quadratic Formula}
The quadratic formula is given by:
\[x = \frac{-b \pm \sqrt{b^2 - 4ac}}{2a}\]

\subsection{Matrix Example}
A sample matrix:
\[
\begin{pmatrix}
a_{11} & a_{12} & a_{13} \\
a_{21} & a_{22} & a_{23} \\
a_{31} & a_{32} & a_{33}
\end{pmatrix}
\]

\section{Tables}
Below is a sample table showing data:

\begin{table}[h]
\centering
\begin{tabular}{@{}lrrr@{}}
\toprule
\textbf{Category} & \textbf{Value 1} & \textbf{Value 2} & \textbf{Total} \\
\midrule
Item A & 10.5 & 20.3 & 30.8 \\
Item B & 15.2 & 25.1 & 40.3 \\
Item C & 22.1 & 12.7 & 34.8 \\
\bottomrule
\end{tabular}
\caption{Sample Data Table}
\label{tab:sample}
\end{table}

\section{Lists}
\subsection{Bullet Points}
\begin{itemize}
  \item First item
  \item Second item
  \item Third item with sub-items:
    \begin{itemize}
      \item Sub-item 1
      \item Sub-item 2
    \end{itemize}
\end{itemize}

\subsection{Numbered List}
\begin{enumerate}
  \item First step
  \item Second step
  \item Third step
\end{enumerate}

\section{Citations and References}
When you use the citation search feature, references will be automatically managed. For example, you might cite a paper like this: \cite{example2024}. All citations will be automatically added to your \texttt{cite.bib} file.

\section{Conclusion}
This template demonstrates basic \LaTeX{} features including mathematical formulas, tables, and lists. The bibliography system is now set up to work seamlessly with Bibby AI's citation management. Feel free to modify it for your needs!

\bibliographystyle{plain}
\bibliography{cite}

\end{document}\documentclass{article}
\usepackage[utf8]{inputenc}
\usepackage{booktabs}
\usepackage{amsmath}
\usepackage{natbib}

\title{My Document}
\author{Your Name}
\date{\today}

\begin{document}

\maketitle

\section{Introduction}
This is a sample document demonstrating various \LaTeX{} features including mathematical equations, tables, and formatting. You can cite papers using the citation search feature, and they will be automatically added to your \texttt{cite.bib} file.

\section{Mathematical Expressions}
Here are some example equations:

\subsection{Quadratic Formula}
The quadratic formula is given by:
\[x = \frac{-b \pm \sqrt{b^2 - 4ac}}{2a}\]

\subsection{Matrix Example}
A sample matrix:
\[
\begin{pmatrix}
a_{11} & a_{12} & a_{13} \\
a_{21} & a_{22} & a_{23} \\
a_{31} & a_{32} & a_{33}
\end{pmatrix}
\]

\section{Tables}
Below is a sample table showing data:

\begin{table}[h]
\centering
\begin{tabular}{@{}lrrr@{}}
\toprule
\textbf{Category} & \textbf{Value 1} & \textbf{Value 2} & \textbf{Total} \\
\midrule
Item A & 10.5 & 20.3 & 30.8 \\
Item B & 15.2 & 25.1 & 40.3 \\
Item C & 22.1 & 12.7 & 34.8 \\
\bottomrule
\end{tabular}
\caption{Sample Data Table}
\label{tab:sample}
\end{table}

\section{Lists}
\subsection{Bullet Points}
\begin{itemize}
  \item First item
  \item Second item
  \item Third item with sub-items:
    \begin{itemize}
      \item Sub-item 1
      \item Sub-item 2
    \end{itemize}
\end{itemize}

\subsection{Numbered List}
\begin{enumerate}
  \item First step
  \item Second step
  \item Third step
\end{enumerate}

\section{Citations and References}
When you use the citation search feature, references will be automatically managed. For example, you might cite a paper like this: \cite{example2024}. All citations will be automatically added to your \texttt{cite.bib} file.

\section{Conclusion}
This template demonstrates basic \LaTeX{} features including mathematical formulas, tables, and lists. The bibliography system is now set up to work seamlessly with Bibby AI's citation management. Feel free to modify it for your needs!

\bibliographystyle{plain}
\bibliography{cite}

\end{document}\documentclass{article}
\usepackage[utf8]{inputenc}
\usepackage{booktabs}
\usepackage{amsmath}
\usepackage{natbib}

\title{My Document}
\author{Your Name}
\date{\today}

\begin{document}

\maketitle

\section{Introduction}
This is a sample document demonstrating various \LaTeX{} features including mathematical equations, tables, and formatting. You can cite papers using the citation search feature, and they will be automatically added to your \texttt{cite.bib} file.

\section{Mathematical Expressions}
Here are some example equations:

\subsection{Quadratic Formula}
The quadratic formula is given by:
\[x = \frac{-b \pm \sqrt{b^2 - 4ac}}{2a}\]

\subsection{Matrix Example}
A sample matrix:
\[
\begin{pmatrix}
a_{11} & a_{12} & a_{13} \\
a_{21} & a_{22} & a_{23} \\
a_{31} & a_{32} & a_{33}
\end{pmatrix}
\]

\section{Tables}
Below is a sample table showing data:

\begin{table}[h]
\centering
\begin{tabular}{@{}lrrr@{}}
\toprule
\textbf{Category} & \textbf{Value 1} & \textbf{Value 2} & \textbf{Total} \\
\midrule
Item A & 10.5 & 20.3 & 30.8 \\
Item B & 15.2 & 25.1 & 40.3 \\
Item C & 22.1 & 12.7 & 34.8 \\
\bottomrule
\end{tabular}
\caption{Sample Data Table}
\label{tab:sample}
\end{table}

\section{Lists}
\subsection{Bullet Points}
\begin{itemize}
  \item First item
  \item Second item
  \item Third item with sub-items:
    \begin{itemize}
      \item Sub-item 1
      \item Sub-item 2
    \end{itemize}
\end{itemize}

\subsection{Numbered List}
\begin{enumerate}
  \item First step
  \item Second step
  \item Third step
\end{enumerate}

\section{Citations and References}
When you use the citation search feature, references will be automatically managed. For example, you might cite a paper like this: \cite{example2024}. All citations will be automatically added to your \texttt{cite.bib} file.

\section{Conclusion}
This template demonstrates basic \LaTeX{} features including mathematical formulas, tables, and lists. The bibliography system is now set up to work seamlessly with Bibby AI's citation management. Feel free to modify it for your needs!

\bibliographystyle{plain}
\bibliography{cite}

\end{document}\documentclass{article}
\usepackage[utf8]{inputenc}
\usepackage{booktabs}
\usepackage{amsmath}
\usepackage{natbib}

\title{My Document}
\author{Your Name}
\date{\today}

\begin{document}

\maketitle

\section{Introduction}
This is a sample document demonstrating various \LaTeX{} features including mathematical equations, tables, and formatting. You can cite papers using the citation search feature, and they will be automatically added to your \texttt{cite.bib} file.

\section{Mathematical Expressions}
Here are some example equations:

\subsection{Quadratic Formula}
The quadratic formula is given by:
\[x = \frac{-b \pm \sqrt{b^2 - 4ac}}{2a}\]

\subsection{Matrix Example}
A sample matrix:
\[
\begin{pmatrix}
a_{11} & a_{12} & a_{13} \\
a_{21} & a_{22} & a_{23} \\
a_{31} & a_{32} & a_{33}
\end{pmatrix}
\]

\section{Tables}
Below is a sample table showing data:

\begin{table}[h]
\centering
\begin{tabular}{@{}lrrr@{}}
\toprule
\textbf{Category} & \textbf{Value 1} & \textbf{Value 2} & \textbf{Total} \\
\midrule
Item A & 10.5 & 20.3 & 30.8 \\
Item B & 15.2 & 25.1 & 40.3 \\
Item C & 22.1 & 12.7 & 34.8 \\
\bottomrule
\end{tabular}
\caption{Sample Data Table}
\label{tab:sample}
\end{table}

\section{Lists}
\subsection{Bullet Points}
\begin{itemize}
  \item First item
  \item Second item
  \item Third item with sub-items:
    \begin{itemize}
      \item Sub-item 1
      \item Sub-item 2
    \end{itemize}
\end{itemize}

\subsection{Numbered List}
\begin{enumerate}
  \item First step
  \item Second step
  \item Third step
\end{enumerate}

\section{Citations and References}
When you use the citation search feature, references will be automatically managed. For example, you might cite a paper like this: \cite{example2024}. All citations will be automatically added to your \texttt{cite.bib} file.

\section{Conclusion}
This template demonstrates basic \LaTeX{} features including mathematical formulas, tables, and lists. The bibliography system is now set up to work seamlessly with Bibby AI's citation management. Feel free to modify it for your needs!

\bibliographystyle{plain}
\bibliography{cite}

\end{document}\documentclass{article}
\usepackage[utf8]{inputenc}
\usepackage{booktabs}
\usepackage{amsmath}
\usepackage{natbib}

\title{My Document}
\author{Your Name}
\date{\today}

\begin{document}

\maketitle

\section{Introduction}
This is a sample document demonstrating various \LaTeX{} features including mathematical equations, tables, and formatting. You can cite papers using the citation search feature, and they will be automatically added to your \texttt{cite.bib} file.

\section{Mathematical Expressions}
Here are some example equations:

\subsection{Quadratic Formula}
The quadratic formula is given by:
\[x = \frac{-b \pm \sqrt{b^2 - 4ac}}{2a}\]

\subsection{Matrix Example}
A sample matrix:
\[
\begin{pmatrix}
a_{11} & a_{12} & a_{13} \\
a_{21} & a_{22} & a_{23} \\
a_{31} & a_{32} & a_{33}
\end{pmatrix}
\]

\section{Tables}
Below is a sample table showing data:

\begin{table}[h]
\centering
\begin{tabular}{@{}lrrr@{}}
\toprule
\textbf{Category} & \textbf{Value 1} & \textbf{Value 2} & \textbf{Total} \\
\midrule
Item A & 10.5 & 20.3 & 30.8 \\
Item B & 15.2 & 25.1 & 40.3 \\
Item C & 22.1 & 12.7 & 34.8 \\
\bottomrule
\end{tabular}
\caption{Sample Data Table}
\label{tab:sample}
\end{table}

\section{Lists}
\subsection{Bullet Points}
\begin{itemize}
  \item First item
  \item Second item
  \item Third item with sub-items:
    \begin{itemize}
      \item Sub-item 1
      \item Sub-item 2
    \end{itemize}
\end{itemize}

\subsection{Numbered List}
\begin{enumerate}
  \item First step
  \item Second step
  \item Third step
\end{enumerate}

\section{Citations and References}
When you use the citation search feature, references will be automatically managed. For example, you might cite a paper like this: \cite{example2024}. All citations will be automatically added to your \texttt{cite.bib} file.

\section{Conclusion}
This template demonstrates basic \LaTeX{} features including mathematical formulas, tables, and lists. The bibliography system is now set up to work seamlessly with Bibby AI's citation management. Feel free to modify it for your needs!

\bibliographystyle{plain}
\bibliography{cite}

\end{document}\documentclass{article}
\usepackage[utf8]{inputenc}
\usepackage{booktabs}
\usepackage{amsmath}
\usepackage{natbib}

\title{My Document}
\author{Your Name}
\date{\today}

\begin{document}

\maketitle

\section{Introduction}
This is a sample document demonstrating various \LaTeX{} features including mathematical equations, tables, and formatting. You can cite papers using the citation search feature, and they will be automatically added to your \texttt{cite.bib} file.

\section{Mathematical Expressions}
Here are some example equations:

\subsection{Quadratic Formula}
The quadratic formula is given by:
\[x = \frac{-b \pm \sqrt{b^2 - 4ac}}{2a}\]

\subsection{Matrix Example}
A sample matrix:
\[
\begin{pmatrix}
a_{11} & a_{12} & a_{13} \\
a_{21} & a_{22} & a_{23} \\
a_{31} & a_{32} & a_{33}
\end{pmatrix}
\]

\section{Tables}
Below is a sample table showing data:

\begin{table}[h]
\centering
\begin{tabular}{@{}lrrr@{}}
\toprule
\textbf{Category} & \textbf{Value 1} & \textbf{Value 2} & \textbf{Total} \\
\midrule
Item A & 10.5 & 20.3 & 30.8 \\
Item B & 15.2 & 25.1 & 40.3 \\
Item C & 22.1 & 12.7 & 34.8 \\
\bottomrule
\end{tabular}
\caption{Sample Data Table}
\label{tab:sample}
\end{table}

\section{Lists}
\subsection{Bullet Points}
\begin{itemize}
  \item First item
  \item Second item
  \item Third item with sub-items:
    \begin{itemize}
      \item Sub-item 1
      \item Sub-item 2
    \end{itemize}
\end{itemize}

\subsection{Numbered List}
\begin{enumerate}
  \item First step
  \item Second step
  \item Third step
\end{enumerate}

\section{Citations and References}
When you use the citation search feature, references will be automatically managed. For example, you might cite a paper like this: \cite{example2024}. All citations will be automatically added to your \texttt{cite.bib} file.

\section{Conclusion}
This template demonstrates basic \LaTeX{} features including mathematical formulas, tables, and lists. The bibliography system is now set up to work seamlessly with Bibby AI's citation management. Feel free to modify it for your needs!

\bibliographystyle{plain}
\bibliography{cite}

\end{document}\documentclass{article}
\usepackage[utf8]{inputenc}
\usepackage{booktabs}
\usepackage{amsmath}
\usepackage{natbib}

\title{My Document}
\author{Your Name}
\date{\today}

\begin{document}

\maketitle

\section{Introduction}
This is a sample document demonstrating various \LaTeX{} features including mathematical equations, tables, and formatting. You can cite papers using the citation search feature, and they will be automatically added to your \texttt{cite.bib} file.

\section{Mathematical Expressions}
Here are some example equations:

\subsection{Quadratic Formula}
The quadratic formula is given by:
\[x = \frac{-b \pm \sqrt{b^2 - 4ac}}{2a}\]

\subsection{Matrix Example}
A sample matrix:
\[
\begin{pmatrix}
a_{11} & a_{12} & a_{13} \\
a_{21} & a_{22} & a_{23} \\
a_{31} & a_{32} & a_{33}
\end{pmatrix}
\]

\section{Tables}
Below is a sample table showing data:

\begin{table}[h]
\centering
\begin{tabular}{@{}lrrr@{}}
\toprule
\textbf{Category} & \textbf{Value 1} & \textbf{Value 2} & \textbf{Total} \\
\midrule
Item A & 10.5 & 20.3 & 30.8 \\
Item B & 15.2 & 25.1 & 40.3 \\
Item C & 22.1 & 12.7 & 34.8 \\
\bottomrule
\end{tabular}
\caption{Sample Data Table}
\label{tab:sample}
\end{table}

\section{Lists}
\subsection{Bullet Points}
\begin{itemize}
  \item First item
  \item Second item
  \item Third item with sub-items:
    \begin{itemize}
      \item Sub-item 1
      \item Sub-item 2
    \end{itemize}
\end{itemize}

\subsection{Numbered List}
\begin{enumerate}
  \item First step
  \item Second step
  \item Third step
\end{enumerate}

\section{Citations and References}
When you use the citation search feature, references will be automatically managed. For example, you might cite a paper like this: \cite{example2024}. All citations will be automatically added to your \texttt{cite.bib} file.

\section{Conclusion}
This template demonstrates basic \LaTeX{} features including mathematical formulas, tables, and lists. The bibliography system is now set up to work seamlessly with Bibby AI's citation management. Feel free to modify it for your needs!

\bibliographystyle{plain}
\bibliography{cite}

\end{document}\documentclass{article}
\usepackage[utf8]{inputenc}
\usepackage{booktabs}
\usepackage{amsmath}
\usepackage{natbib}

\title{My Document}
\author{Your Name}
\date{\today}

\begin{document}

\maketitle

\section{Introduction}
This is a sample document demonstrating various \LaTeX{} features including mathematical equations, tables, and formatting. You can cite papers using the citation search feature, and they will be automatically added to your \texttt{cite.bib} file.

\section{Mathematical Expressions}
Here are some example equations:

\subsection{Quadratic Formula}
The quadratic formula is given by:
\[x = \frac{-b \pm \sqrt{b^2 - 4ac}}{2a}\]

\subsection{Matrix Example}
A sample matrix:
\[
\begin{pmatrix}
a_{11} & a_{12} & a_{13} \\
a_{21} & a_{22} & a_{23} \\
a_{31} & a_{32} & a_{33}
\end{pmatrix}
\]

\section{Tables}
Below is a sample table showing data:

\begin{table}[h]
\centering
\begin{tabular}{@{}lrrr@{}}
\toprule
\textbf{Category} & \textbf{Value 1} & \textbf{Value 2} & \textbf{Total} \\
\midrule
Item A & 10.5 & 20.3 & 30.8 \\
Item B & 15.2 & 25.1 & 40.3 \\
Item C & 22.1 & 12.7 & 34.8 \\
\bottomrule
\end{tabular}
\caption{Sample Data Table}
\label{tab:sample}
\end{table}

\section{Lists}
\subsection{Bullet Points}
\begin{itemize}
  \item First item
  \item Second item
  \item Third item with sub-items:
    \begin{itemize}
      \item Sub-item 1
      \item Sub-item 2
    \end{itemize}
\end{itemize}

\subsection{Numbered List}
\begin{enumerate}
  \item First step
  \item Second step
  \item Third step
\end{enumerate}

\section{Citations and References}
When you use the citation search feature, references will be automatically managed. For example, you might cite a paper like this: \cite{example2024}. All citations will be automatically added to your \texttt{cite.bib} file.

\section{Conclusion}
This template demonstrates basic \LaTeX{} features including mathematical formulas, tables, and lists. The bibliography system is now set up to work seamlessly with Bibby AI's citation management. Feel free to modify it for your needs!

\bibliographystyle{plain}
\bibliography{cite}

\end{document}\documentclass{article}
\usepackage[utf8]{inputenc}
\usepackage{booktabs}
\usepackage{amsmath}
\usepackage{natbib}

\title{My Document}
\author{Your Name}
\date{\today}

\begin{document}

\maketitle

\section{Introduction}
This is a sample document demonstrating various \LaTeX{} features including mathematical equations, tables, and formatting. You can cite papers using the citation search feature, and they will be automatically added to your \texttt{cite.bib} file.

\section{Mathematical Expressions}
Here are some example equations:

\subsection{Quadratic Formula}
The quadratic formula is given by:
\[x = \frac{-b \pm \sqrt{b^2 - 4ac}}{2a}\]

\subsection{Matrix Example}
A sample matrix:
\[
\begin{pmatrix}
a_{11} & a_{12} & a_{13} \\
a_{21} & a_{22} & a_{23} \\
a_{31} & a_{32} & a_{33}
\end{pmatrix}
\]

\section{Tables}
Below is a sample table showing data:

\begin{table}[h]
\centering
\begin{tabular}{@{}lrrr@{}}
\toprule
\textbf{Category} & \textbf{Value 1} & \textbf{Value 2} & \textbf{Total} \\
\midrule
Item A & 10.5 & 20.3 & 30.8 \\
Item B & 15.2 & 25.1 & 40.3 \\
Item C & 22.1 & 12.7 & 34.8 \\
\bottomrule
\end{tabular}
\caption{Sample Data Table}
\label{tab:sample}
\end{table}

\section{Lists}
\subsection{Bullet Points}
\begin{itemize}
  \item First item
  \item Second item
  \item Third item with sub-items:
    \begin{itemize}
      \item Sub-item 1
      \item Sub-item 2
    \end{itemize}
\end{itemize}

\subsection{Numbered List}
\begin{enumerate}
  \item First step
  \item Second step
  \item Third step
\end{enumerate}

\section{Citations and References}
When you use the citation search feature, references will be automatically managed. For example, you might cite a paper like this: \cite{example2024}. All citations will be automatically added to your \texttt{cite.bib} file.

\section{Conclusion}
This template demonstrates basic \LaTeX{} features including mathematical formulas, tables, and lists. The bibliography system is now set up to work seamlessly with Bibby AI's citation management. Feel free to modify it for your needs!

\bibliographystyle{plain}
\bibliography{cite}

\end{document}\documentclass{article}
\usepackage[utf8]{inputenc}
\usepackage{booktabs}
\usepackage{amsmath}
\usepackage{natbib}

\title{My Document}
\author{Your Name}
\date{\today}

\begin{document}

\maketitle

\section{Introduction}
This is a sample document demonstrating various \LaTeX{} features including mathematical equations, tables, and formatting. You can cite papers using the citation search feature, and they will be automatically added to your \texttt{cite.bib} file.

\section{Mathematical Expressions}
Here are some example equations:

\subsection{Quadratic Formula}
The quadratic formula is given by:
\[x = \frac{-b \pm \sqrt{b^2 - 4ac}}{2a}\]

\subsection{Matrix Example}
A sample matrix:
\[
\begin{pmatrix}
a_{11} & a_{12} & a_{13} \\
a_{21} & a_{22} & a_{23} \\
a_{31} & a_{32} & a_{33}
\end{pmatrix}
\]

\section{Tables}
Below is a sample table showing data:

\begin{table}[h]
\centering
\begin{tabular}{@{}lrrr@{}}
\toprule
\textbf{Category} & \textbf{Value 1} & \textbf{Value 2} & \textbf{Total} \\
\midrule
Item A & 10.5 & 20.3 & 30.8 \\
Item B & 15.2 & 25.1 & 40.3 \\
Item C & 22.1 & 12.7 & 34.8 \\
\bottomrule
\end{tabular}
\caption{Sample Data Table}
\label{tab:sample}
\end{table}

\section{Lists}
\subsection{Bullet Points}
\begin{itemize}
  \item First item
  \item Second item
  \item Third item with sub-items:
    \begin{itemize}
      \item Sub-item 1
      \item Sub-item 2
    \end{itemize}
\end{itemize}

\subsection{Numbered List}
\begin{enumerate}
  \item First step
  \item Second step
  \item Third step
\end{enumerate}

\section{Citations and References}
When you use the citation search feature, references will be automatically managed. For example, you might cite a paper like this: \cite{example2024}. All citations will be automatically added to your \texttt{cite.bib} file.

\section{Conclusion}
This template demonstrates basic \LaTeX{} features including mathematical formulas, tables, and lists. The bibliography system is now set up to work seamlessly with Bibby AI's citation management. Feel free to modify it for your needs!

\bibliographystyle{plain}
\bibliography{cite}

\end{document}\documentclass{article}
\usepackage[utf8]{inputenc}
\usepackage{booktabs}
\usepackage{amsmath}
\usepackage{natbib}

\title{My Document}
\author{Your Name}
\date{\today}

\begin{document}

\maketitle

\section{Introduction}
This is a sample document demonstrating various \LaTeX{} features including mathematical equations, tables, and formatting. You can cite papers using the citation search feature, and they will be automatically added to your \texttt{cite.bib} file.

\section{Mathematical Expressions}
Here are some example equations:

\subsection{Quadratic Formula}
The quadratic formula is given by:
\[x = \frac{-b \pm \sqrt{b^2 - 4ac}}{2a}\]

\subsection{Matrix Example}
A sample matrix:
\[
\begin{pmatrix}
a_{11} & a_{12} & a_{13} \\
a_{21} & a_{22} & a_{23} \\
a_{31} & a_{32} & a_{33}
\end{pmatrix}
\]

\section{Tables}
Below is a sample table showing data:

\begin{table}[h]
\centering
\begin{tabular}{@{}lrrr@{}}
\toprule
\textbf{Category} & \textbf{Value 1} & \textbf{Value 2} & \textbf{Total} \\
\midrule
Item A & 10.5 & 20.3 & 30.8 \\
Item B & 15.2 & 25.1 & 40.3 \\
Item C & 22.1 & 12.7 & 34.8 \\
\bottomrule
\end{tabular}
\caption{Sample Data Table}
\label{tab:sample}
\end{table}

\section{Lists}
\subsection{Bullet Points}
\begin{itemize}
  \item First item
  \item Second item
  \item Third item with sub-items:
    \begin{itemize}
      \item Sub-item 1
      \item Sub-item 2
    \end{itemize}
\end{itemize}

\subsection{Numbered List}
\begin{enumerate}
  \item First step
  \item Second step
  \item Third step
\end{enumerate}

\section{Citations and References}
When you use the citation search feature, references will be automatically managed. For example, you might cite a paper like this: \cite{example2024}. All citations will be automatically added to your \texttt{cite.bib} file.

\section{Conclusion}
This template demonstrates basic \LaTeX{} features including mathematical formulas, tables, and lists. The bibliography system is now set up to work seamlessly with Bibby AI's citation management. Feel free to modify it for your needs!

\bibliographystyle{plain}
\bibliography{cite}

\end{document}\documentclass{article}
\usepackage[utf8]{inputenc}
\usepackage{booktabs}
\usepackage{amsmath}
\usepackage{natbib}

\title{My Document}
\author{Your Name}
\date{\today}

\begin{document}

\maketitle

\section{Introduction}
This is a sample document demonstrating various \LaTeX{} features including mathematical equations, tables, and formatting. You can cite papers using the citation search feature, and they will be automatically added to your \texttt{cite.bib} file.

\section{Mathematical Expressions}
Here are some example equations:

\subsection{Quadratic Formula}
The quadratic formula is given by:
\[x = \frac{-b \pm \sqrt{b^2 - 4ac}}{2a}\]

\subsection{Matrix Example}
A sample matrix:
\[
\begin{pmatrix}
a_{11} & a_{12} & a_{13} \\
a_{21} & a_{22} & a_{23} \\
a_{31} & a_{32} & a_{33}
\end{pmatrix}
\]

\section{Tables}
Below is a sample table showing data:

\begin{table}[h]
\centering
\begin{tabular}{@{}lrrr@{}}
\toprule
\textbf{Category} & \textbf{Value 1} & \textbf{Value 2} & \textbf{Total} \\
\midrule
Item A & 10.5 & 20.3 & 30.8 \\
Item B & 15.2 & 25.1 & 40.3 \\
Item C & 22.1 & 12.7 & 34.8 \\
\bottomrule
\end{tabular}
\caption{Sample Data Table}
\label{tab:sample}
\end{table}

\section{Lists}
\subsection{Bullet Points}
\begin{itemize}
  \item First item
  \item Second item
  \item Third item with sub-items:
    \begin{itemize}
      \item Sub-item 1
      \item Sub-item 2
    \end{itemize}
\end{itemize}

\subsection{Numbered List}
\begin{enumerate}
  \item First step
  \item Second step
  \item Third step
\end{enumerate}

\section{Citations and References}
When you use the citation search feature, references will be automatically managed. For example, you might cite a paper like this: \cite{example2024}. All citations will be automatically added to your \texttt{cite.bib} file.

\section{Conclusion}
This template demonstrates basic \LaTeX{} features including mathematical formulas, tables, and lists. The bibliography system is now set up to work seamlessly with Bibby AI's citation management. Feel free to modify it for your needs!

\bibliographystyle{plain}
\bibliography{cite}

\end{document}\documentclass{article}
\usepackage[utf8]{inputenc}
\usepackage{booktabs}
\usepackage{amsmath}
\usepackage{natbib}

\title{My Document}
\author{Your Name}
\date{\today}

\begin{document}

\maketitle

\section{Introduction}
This is a sample document demonstrating various \LaTeX{} features including mathematical equations, tables, and formatting. You can cite papers using the citation search feature, and they will be automatically added to your \texttt{cite.bib} file.

\section{Mathematical Expressions}
Here are some example equations:

\subsection{Quadratic Formula}
The quadratic formula is given by:
\[x = \frac{-b \pm \sqrt{b^2 - 4ac}}{2a}\]

\subsection{Matrix Example}
A sample matrix:
\[
\begin{pmatrix}
a_{11} & a_{12} & a_{13} \\
a_{21} & a_{22} & a_{23} \\
a_{31} & a_{32} & a_{33}
\end{pmatrix}
\]

\section{Tables}
Below is a sample table showing data:

\begin{table}[h]
\centering
\begin{tabular}{@{}lrrr@{}}
\toprule
\textbf{Category} & \textbf{Value 1} & \textbf{Value 2} & \textbf{Total} \\
\midrule
Item A & 10.5 & 20.3 & 30.8 \\
Item B & 15.2 & 25.1 & 40.3 \\
Item C & 22.1 & 12.7 & 34.8 \\
\bottomrule
\end{tabular}
\caption{Sample Data Table}
\label{tab:sample}
\end{table}

\section{Lists}
\subsection{Bullet Points}
\begin{itemize}
  \item First item
  \item Second item
  \item Third item with sub-items:
    \begin{itemize}
      \item Sub-item 1
      \item Sub-item 2
    \end{itemize}
\end{itemize}

\subsection{Numbered List}
\begin{enumerate}
  \item First step
  \item Second step
  \item Third step
\end{enumerate}

\section{Citations and References}
When you use the citation search feature, references will be automatically managed. For example, you might cite a paper like this: \cite{example2024}. All citations will be automatically added to your \texttt{cite.bib} file.

\section{Conclusion}
This template demonstrates basic \LaTeX{} features including mathematical formulas, tables, and lists. The bibliography system is now set up to work seamlessly with Bibby AI's citation management. Feel free to modify it for your needs!

\bibliographystyle{plain}
\bibliography{cite}

\end{document}\documentclass{article}
\usepackage[utf8]{inputenc}
\usepackage{booktabs}
\usepackage{amsmath}
\usepackage{natbib}

\title{My Document}
\author{Your Name}
\date{\today}

\begin{document}

\maketitle

\section{Introduction}
This is a sample document demonstrating various \LaTeX{} features including mathematical equations, tables, and formatting. You can cite papers using the citation search feature, and they will be automatically added to your \texttt{cite.bib} file.

\section{Mathematical Expressions}
Here are some example equations:

\subsection{Quadratic Formula}
The quadratic formula is given by:
\[x = \frac{-b \pm \sqrt{b^2 - 4ac}}{2a}\]

\subsection{Matrix Example}
A sample matrix:
\[
\begin{pmatrix}
a_{11} & a_{12} & a_{13} \\
a_{21} & a_{22} & a_{23} \\
a_{31} & a_{32} & a_{33}
\end{pmatrix}
\]

\section{Tables}
Below is a sample table showing data:

\begin{table}[h]
\centering
\begin{tabular}{@{}lrrr@{}}
\toprule
\textbf{Category} & \textbf{Value 1} & \textbf{Value 2} & \textbf{Total} \\
\midrule
Item A & 10.5 & 20.3 & 30.8 \\
Item B & 15.2 & 25.1 & 40.3 \\
Item C & 22.1 & 12.7 & 34.8 \\
\bottomrule
\end{tabular}
\caption{Sample Data Table}
\label{tab:sample}
\end{table}

\section{Lists}
\subsection{Bullet Points}
\begin{itemize}
  \item First item
  \item Second item
  \item Third item with sub-items:
    \begin{itemize}
      \item Sub-item 1
      \item Sub-item 2
    \end{itemize}
\end{itemize}

\subsection{Numbered List}
\begin{enumerate}
  \item First step
  \item Second step
  \item Third step
\end{enumerate}

\section{Citations and References}
When you use the citation search feature, references will be automatically managed. For example, you might cite a paper like this: \cite{example2024}. All citations will be automatically added to your \texttt{cite.bib} file.

\section{Conclusion}
This template demonstrates basic \LaTeX{} features including mathematical formulas, tables, and lists. The bibliography system is now set up to work seamlessly with Bibby AI's citation management. Feel free to modify it for your needs!

\bibliographystyle{plain}
\bibliography{cite}

\end{document}\documentclass{article}
\usepackage[utf8]{inputenc}
\usepackage{booktabs}
\usepackage{amsmath}
\usepackage{natbib}

\title{My Document}
\author{Your Name}
\date{\today}

\begin{document}

\maketitle

\section{Introduction}
This is a sample document demonstrating various \LaTeX{} features including mathematical equations, tables, and formatting. You can cite papers using the citation search feature, and they will be automatically added to your \texttt{cite.bib} file.

\section{Mathematical Expressions}
Here are some example equations:

\subsection{Quadratic Formula}
The quadratic formula is given by:
\[x = \frac{-b \pm \sqrt{b^2 - 4ac}}{2a}\]

\subsection{Matrix Example}
A sample matrix:
\[
\begin{pmatrix}
a_{11} & a_{12} & a_{13} \\
a_{21} & a_{22} & a_{23} \\
a_{31} & a_{32} & a_{33}
\end{pmatrix}
\]

\section{Tables}
Below is a sample table showing data:

\begin{table}[h]
\centering
\begin{tabular}{@{}lrrr@{}}
\toprule
\textbf{Category} & \textbf{Value 1} & \textbf{Value 2} & \textbf{Total} \\
\midrule
Item A & 10.5 & 20.3 & 30.8 \\
Item B & 15.2 & 25.1 & 40.3 \\
Item C & 22.1 & 12.7 & 34.8 \\
\bottomrule
\end{tabular}
\caption{Sample Data Table}
\label{tab:sample}
\end{table}

\section{Lists}
\subsection{Bullet Points}
\begin{itemize}
  \item First item
  \item Second item
  \item Third item with sub-items:
    \begin{itemize}
      \item Sub-item 1
      \item Sub-item 2
    \end{itemize}
\end{itemize}

\subsection{Numbered List}
\begin{enumerate}
  \item First step
  \item Second step
  \item Third step
\end{enumerate}

\section{Citations and References}
When you use the citation search feature, references will be automatically managed. For example, you might cite a paper like this: \cite{example2024}. All citations will be automatically added to your \texttt{cite.bib} file.

\section{Conclusion}
This template demonstrates basic \LaTeX{} features including mathematical formulas, tables, and lists. The bibliography system is now set up to work seamlessly with Bibby AI's citation management. Feel free to modify it for your needs!

\bibliographystyle{plain}
\bibliography{cite}

\end{document}\documentclass{article}
\usepackage[utf8]{inputenc}
\usepackage{booktabs}
\usepackage{amsmath}
\usepackage{natbib}

\title{My Document}
\author{Your Name}
\date{\today}

\begin{document}

\maketitle

\section{Introduction}
This is a sample document demonstrating various \LaTeX{} features including mathematical equations, tables, and formatting. You can cite papers using the citation search feature, and they will be automatically added to your \texttt{cite.bib} file.

\section{Mathematical Expressions}
Here are some example equations:

\subsection{Quadratic Formula}
The quadratic formula is given by:
\[x = \frac{-b \pm \sqrt{b^2 - 4ac}}{2a}\]

\subsection{Matrix Example}
A sample matrix:
\[
\begin{pmatrix}
a_{11} & a_{12} & a_{13} \\
a_{21} & a_{22} & a_{23} \\
a_{31} & a_{32} & a_{33}
\end{pmatrix}
\]

\section{Tables}
Below is a sample table showing data:

\begin{table}[h]
\centering
\begin{tabular}{@{}lrrr@{}}
\toprule
\textbf{Category} & \textbf{Value 1} & \textbf{Value 2} & \textbf{Total} \\
\midrule
Item A & 10.5 & 20.3 & 30.8 \\
Item B & 15.2 & 25.1 & 40.3 \\
Item C & 22.1 & 12.7 & 34.8 \\
\bottomrule
\end{tabular}
\caption{Sample Data Table}
\label{tab:sample}
\end{table}

\section{Lists}
\subsection{Bullet Points}
\begin{itemize}
  \item First item
  \item Second item
  \item Third item with sub-items:
    \begin{itemize}
      \item Sub-item 1
      \item Sub-item 2
    \end{itemize}
\end{itemize}

\subsection{Numbered List}
\begin{enumerate}
  \item First step
  \item Second step
  \item Third step
\end{enumerate}

\section{Citations and References}
When you use the citation search feature, references will be automatically managed. For example, you might cite a paper like this: \cite{example2024}. All citations will be automatically added to your \texttt{cite.bib} file.

\section{Conclusion}
This template demonstrates basic \LaTeX{} features including mathematical formulas, tables, and lists. The bibliography system is now set up to work seamlessly with Bibby AI's citation management. Feel free to modify it for your needs!

\bibliographystyle{plain}
\bibliography{cite}

\end{document}\documentclass{article}
\usepackage[utf8]{inputenc}
\usepackage{booktabs}
\usepackage{amsmath}
\usepackage{natbib}

\title{My Document}
\author{Your Name}
\date{\today}

\begin{document}

\maketitle

\section{Introduction}
This is a sample document demonstrating various \LaTeX{} features including mathematical equations, tables, and formatting. You can cite papers using the citation search feature, and they will be automatically added to your \texttt{cite.bib} file.

\section{Mathematical Expressions}
Here are some example equations:

\subsection{Quadratic Formula}
The quadratic formula is given by:
\[x = \frac{-b \pm \sqrt{b^2 - 4ac}}{2a}\]

\subsection{Matrix Example}
A sample matrix:
\[
\begin{pmatrix}
a_{11} & a_{12} & a_{13} \\
a_{21} & a_{22} & a_{23} \\
a_{31} & a_{32} & a_{33}
\end{pmatrix}
\]

\section{Tables}
Below is a sample table showing data:

\begin{table}[h]
\centering
\begin{tabular}{@{}lrrr@{}}
\toprule
\textbf{Category} & \textbf{Value 1} & \textbf{Value 2} & \textbf{Total} \\
\midrule
Item A & 10.5 & 20.3 & 30.8 \\
Item B & 15.2 & 25.1 & 40.3 \\
Item C & 22.1 & 12.7 & 34.8 \\
\bottomrule
\end{tabular}
\caption{Sample Data Table}
\label{tab:sample}
\end{table}

\section{Lists}
\subsection{Bullet Points}
\begin{itemize}
  \item First item
  \item Second item
  \item Third item with sub-items:
    \begin{itemize}
      \item Sub-item 1
      \item Sub-item 2
    \end{itemize}
\end{itemize}

\subsection{Numbered List}
\begin{enumerate}
  \item First step
  \item Second step
  \item Third step
\end{enumerate}

\section{Citations and References}
When you use the citation search feature, references will be automatically managed. For example, you might cite a paper like this: \cite{example2024}. All citations will be automatically added to your \texttt{cite.bib} file.

\section{Conclusion}
This template demonstrates basic \LaTeX{} features including mathematical formulas, tables, and lists. The bibliography system is now set up to work seamlessly with Bibby AI's citation management. Feel free to modify it for your needs!

\bibliographystyle{plain}
\bibliography{cite}

\end{document}\documentclass{article}
\usepackage[utf8]{inputenc}
\usepackage{booktabs}
\usepackage{amsmath}
\usepackage{natbib}

\title{My Document}
\author{Your Name}
\date{\today}

\begin{document}

\maketitle

\section{Introduction}
This is a sample document demonstrating various \LaTeX{} features including mathematical equations, tables, and formatting. You can cite papers using the citation search feature, and they will be automatically added to your \texttt{cite.bib} file.

\section{Mathematical Expressions}
Here are some example equations:

\subsection{Quadratic Formula}
The quadratic formula is given by:
\[x = \frac{-b \pm \sqrt{b^2 - 4ac}}{2a}\]

\subsection{Matrix Example}
A sample matrix:
\[
\begin{pmatrix}
a_{11} & a_{12} & a_{13} \\
a_{21} & a_{22} & a_{23} \\
a_{31} & a_{32} & a_{33}
\end{pmatrix}
\]

\section{Tables}
Below is a sample table showing data:

\begin{table}[h]
\centering
\begin{tabular}{@{}lrrr@{}}
\toprule
\textbf{Category} & \textbf{Value 1} & \textbf{Value 2} & \textbf{Total} \\
\midrule
Item A & 10.5 & 20.3 & 30.8 \\
Item B & 15.2 & 25.1 & 40.3 \\
Item C & 22.1 & 12.7 & 34.8 \\
\bottomrule
\end{tabular}
\caption{Sample Data Table}
\label{tab:sample}
\end{table}

\section{Lists}
\subsection{Bullet Points}
\begin{itemize}
  \item First item
  \item Second item
  \item Third item with sub-items:
    \begin{itemize}
      \item Sub-item 1
      \item Sub-item 2
    \end{itemize}
\end{itemize}

\subsection{Numbered List}
\begin{enumerate}
  \item First step
  \item Second step
  \item Third step
\end{enumerate}

\section{Citations and References}
When you use the citation search feature, references will be automatically managed. For example, you might cite a paper like this: \cite{example2024}. All citations will be automatically added to your \texttt{cite.bib} file.

\section{Conclusion}
This template demonstrates basic \LaTeX{} features including mathematical formulas, tables, and lists. The bibliography system is now set up to work seamlessly with Bibby AI's citation management. Feel free to modify it for your needs!

\bibliographystyle{plain}
\bibliography{cite}

\end{document}\documentclass{article}
\usepackage[utf8]{inputenc}
\usepackage{booktabs}
\usepackage{amsmath}
\usepackage{natbib}

\title{My Document}
\author{Your Name}
\date{\today}

\begin{document}

\maketitle

\section{Introduction}
This is a sample document demonstrating various \LaTeX{} features including mathematical equations, tables, and formatting. You can cite papers using the citation search feature, and they will be automatically added to your \texttt{cite.bib} file.

\section{Mathematical Expressions}
Here are some example equations:

\subsection{Quadratic Formula}
The quadratic formula is given by:
\[x = \frac{-b \pm \sqrt{b^2 - 4ac}}{2a}\]

\subsection{Matrix Example}
A sample matrix:
\[
\begin{pmatrix}
a_{11} & a_{12} & a_{13} \\
a_{21} & a_{22} & a_{23} \\
a_{31} & a_{32} & a_{33}
\end{pmatrix}
\]

\section{Tables}
Below is a sample table showing data:

\begin{table}[h]
\centering
\begin{tabular}{@{}lrrr@{}}
\toprule
\textbf{Category} & \textbf{Value 1} & \textbf{Value 2} & \textbf{Total} \\
\midrule
Item A & 10.5 & 20.3 & 30.8 \\
Item B & 15.2 & 25.1 & 40.3 \\
Item C & 22.1 & 12.7 & 34.8 \\
\bottomrule
\end{tabular}
\caption{Sample Data Table}
\label{tab:sample}
\end{table}

\section{Lists}
\subsection{Bullet Points}
\begin{itemize}
  \item First item
  \item Second item
  \item Third item with sub-items:
    \begin{itemize}
      \item Sub-item 1
      \item Sub-item 2
    \end{itemize}
\end{itemize}

\subsection{Numbered List}
\begin{enumerate}
  \item First step
  \item Second step
  \item Third step
\end{enumerate}

\section{Citations and References}
When you use the citation search feature, references will be automatically managed. For example, you might cite a paper like this: \cite{example2024}. All citations will be automatically added to your \texttt{cite.bib} file.

\section{Conclusion}
This template demonstrates basic \LaTeX{} features including mathematical formulas, tables, and lists. The bibliography system is now set up to work seamlessly with Bibby AI's citation management. Feel free to modify it for your needs!

\bibliographystyle{plain}
\bibliography{cite}

\end{document}\documentclass{article}
\usepackage[utf8]{inputenc}
\usepackage{booktabs}
\usepackage{amsmath}
\usepackage{natbib}

\title{My Document}
\author{Your Name}
\date{\today}

\begin{document}

\maketitle

\section{Introduction}
This is a sample document demonstrating various \LaTeX{} features including mathematical equations, tables, and formatting. You can cite papers using the citation search feature, and they will be automatically added to your \texttt{cite.bib} file.

\section{Mathematical Expressions}
Here are some example equations:

\subsection{Quadratic Formula}
The quadratic formula is given by:
\[x = \frac{-b \pm \sqrt{b^2 - 4ac}}{2a}\]

\subsection{Matrix Example}
A sample matrix:
\[
\begin{pmatrix}
a_{11} & a_{12} & a_{13} \\
a_{21} & a_{22} & a_{23} \\
a_{31} & a_{32} & a_{33}
\end{pmatrix}
\]

\section{Tables}
Below is a sample table showing data:

\begin{table}[h]
\centering
\begin{tabular}{@{}lrrr@{}}
\toprule
\textbf{Category} & \textbf{Value 1} & \textbf{Value 2} & \textbf{Total} \\
\midrule
Item A & 10.5 & 20.3 & 30.8 \\
Item B & 15.2 & 25.1 & 40.3 \\
Item C & 22.1 & 12.7 & 34.8 \\
\bottomrule
\end{tabular}
\caption{Sample Data Table}
\label{tab:sample}
\end{table}

\section{Lists}
\subsection{Bullet Points}
\begin{itemize}
  \item First item
  \item Second item
  \item Third item with sub-items:
    \begin{itemize}
      \item Sub-item 1
      \item Sub-item 2
    \end{itemize}
\end{itemize}

\subsection{Numbered List}
\begin{enumerate}
  \item First step
  \item Second step
  \item Third step
\end{enumerate}

\section{Citations and References}
When you use the citation search feature, references will be automatically managed. For example, you might cite a paper like this: \cite{example2024}. All citations will be automatically added to your \texttt{cite.bib} file.

\section{Conclusion}
This template demonstrates basic \LaTeX{} features including mathematical formulas, tables, and lists. The bibliography system is now set up to work seamlessly with Bibby AI's citation management. Feel free to modify it for your needs!

\bibliographystyle{plain}
\bibliography{cite}

\end{document}\documentclass{article}
\usepackage[utf8]{inputenc}
\usepackage{booktabs}
\usepackage{amsmath}
\usepackage{natbib}

\title{My Document}
\author{Your Name}
\date{\today}

\begin{document}

\maketitle

\section{Introduction}
This is a sample document demonstrating various \LaTeX{} features including mathematical equations, tables, and formatting. You can cite papers using the citation search feature, and they will be automatically added to your \texttt{cite.bib} file.

\section{Mathematical Expressions}
Here are some example equations:

\subsection{Quadratic Formula}
The quadratic formula is given by:
\[x = \frac{-b \pm \sqrt{b^2 - 4ac}}{2a}\]

\subsection{Matrix Example}
A sample matrix:
\[
\begin{pmatrix}
a_{11} & a_{12} & a_{13} \\
a_{21} & a_{22} & a_{23} \\
a_{31} & a_{32} & a_{33}
\end{pmatrix}
\]

\section{Tables}
Below is a sample table showing data:

\begin{table}[h]
\centering
\begin{tabular}{@{}lrrr@{}}
\toprule
\textbf{Category} & \textbf{Value 1} & \textbf{Value 2} & \textbf{Total} \\
\midrule
Item A & 10.5 & 20.3 & 30.8 \\
Item B & 15.2 & 25.1 & 40.3 \\
Item C & 22.1 & 12.7 & 34.8 \\
\bottomrule
\end{tabular}
\caption{Sample Data Table}
\label{tab:sample}
\end{table}

\section{Lists}
\subsection{Bullet Points}
\begin{itemize}
  \item First item
  \item Second item
  \item Third item with sub-items:
    \begin{itemize}
      \item Sub-item 1
      \item Sub-item 2
    \end{itemize}
\end{itemize}

\subsection{Numbered List}
\begin{enumerate}
  \item First step
  \item Second step
  \item Third step
\end{enumerate}

\section{Citations and References}
When you use the citation search feature, references will be automatically managed. For example, you might cite a paper like this: \cite{example2024}. All citations will be automatically added to your \texttt{cite.bib} file.

\section{Conclusion}
This template demonstrates basic \LaTeX{} features including mathematical formulas, tables, and lists. The bibliography system is now set up to work seamlessly with Bibby AI's citation management. Feel free to modify it for your needs!

\bibliographystyle{plain}
\bibliography{cite}

\end{document}\documentclass{article}
\usepackage[utf8]{inputenc}
\usepackage{booktabs}
\usepackage{amsmath}
\usepackage{natbib}

\title{My Document}
\author{Your Name}
\date{\today}

\begin{document}

\maketitle

\section{Introduction}
This is a sample document demonstrating various \LaTeX{} features including mathematical equations, tables, and formatting. You can cite papers using the citation search feature, and they will be automatically added to your \texttt{cite.bib} file.

\section{Mathematical Expressions}
Here are some example equations:

\subsection{Quadratic Formula}
The quadratic formula is given by:
\[x = \frac{-b \pm \sqrt{b^2 - 4ac}}{2a}\]

\subsection{Matrix Example}
A sample matrix:
\[
\begin{pmatrix}
a_{11} & a_{12} & a_{13} \\
a_{21} & a_{22} & a_{23} \\
a_{31} & a_{32} & a_{33}
\end{pmatrix}
\]

\section{Tables}
Below is a sample table showing data:

\begin{table}[h]
\centering
\begin{tabular}{@{}lrrr@{}}
\toprule
\textbf{Category} & \textbf{Value 1} & \textbf{Value 2} & \textbf{Total} \\
\midrule
Item A & 10.5 & 20.3 & 30.8 \\
Item B & 15.2 & 25.1 & 40.3 \\
Item C & 22.1 & 12.7 & 34.8 \\
\bottomrule
\end{tabular}
\caption{Sample Data Table}
\label{tab:sample}
\end{table}

\section{Lists}
\subsection{Bullet Points}
\begin{itemize}
  \item First item
  \item Second item
  \item Third item with sub-items:
    \begin{itemize}
      \item Sub-item 1
      \item Sub-item 2
    \end{itemize}
\end{itemize}

\subsection{Numbered List}
\begin{enumerate}
  \item First step
  \item Second step
  \item Third step
\end{enumerate}

\section{Citations and References}
When you use the citation search feature, references will be automatically managed. For example, you might cite a paper like this: \cite{example2024}. All citations will be automatically added to your \texttt{cite.bib} file.

\section{Conclusion}
This template demonstrates basic \LaTeX{} features including mathematical formulas, tables, and lists. The bibliography system is now set up to work seamlessly with Bibby AI's citation management. Feel free to modify it for your needs!

\bibliographystyle{plain}
\bibliography{cite}

\end{document}\documentclass{article}
\usepackage[utf8]{inputenc}
\usepackage{booktabs}
\usepackage{amsmath}
\usepackage{natbib}

\title{My Document}
\author{Your Name}
\date{\today}

\begin{document}

\maketitle

\section{Introduction}
This is a sample document demonstrating various \LaTeX{} features including mathematical equations, tables, and formatting. You can cite papers using the citation search feature, and they will be automatically added to your \texttt{cite.bib} file.

\section{Mathematical Expressions}
Here are some example equations:

\subsection{Quadratic Formula}
The quadratic formula is given by:
\[x = \frac{-b \pm \sqrt{b^2 - 4ac}}{2a}\]

\subsection{Matrix Example}
A sample matrix:
\[
\begin{pmatrix}
a_{11} & a_{12} & a_{13} \\
a_{21} & a_{22} & a_{23} \\
a_{31} & a_{32} & a_{33}
\end{pmatrix}
\]

\section{Tables}
Below is a sample table showing data:

\begin{table}[h]
\centering
\begin{tabular}{@{}lrrr@{}}
\toprule
\textbf{Category} & \textbf{Value 1} & \textbf{Value 2} & \textbf{Total} \\
\midrule
Item A & 10.5 & 20.3 & 30.8 \\
Item B & 15.2 & 25.1 & 40.3 \\
Item C & 22.1 & 12.7 & 34.8 \\
\bottomrule
\end{tabular}
\caption{Sample Data Table}
\label{tab:sample}
\end{table}

\section{Lists}
\subsection{Bullet Points}
\begin{itemize}
  \item First item
  \item Second item
  \item Third item with sub-items:
    \begin{itemize}
      \item Sub-item 1
      \item Sub-item 2
    \end{itemize}
\end{itemize}

\subsection{Numbered List}
\begin{enumerate}
  \item First step
  \item Second step
  \item Third step
\end{enumerate}

\section{Citations and References}
When you use the citation search feature, references will be automatically managed. For example, you might cite a paper like this: \cite{example2024}. All citations will be automatically added to your \texttt{cite.bib} file.

\section{Conclusion}
This template demonstrates basic \LaTeX{} features including mathematical formulas, tables, and lists. The bibliography system is now set up to work seamlessly with Bibby AI's citation management. Feel free to modify it for your needs!

\bibliographystyle{plain}
\bibliography{cite}

\end{document}\documentclass{article}
\usepackage[utf8]{inputenc}
\usepackage{booktabs}
\usepackage{amsmath}
\usepackage{natbib}

\title{My Document}
\author{Your Name}
\date{\today}

\begin{document}

\maketitle

\section{Introduction}
This is a sample document demonstrating various \LaTeX{} features including mathematical equations, tables, and formatting. You can cite papers using the citation search feature, and they will be automatically added to your \texttt{cite.bib} file.

\section{Mathematical Expressions}
Here are some example equations:

\subsection{Quadratic Formula}
The quadratic formula is given by:
\[x = \frac{-b \pm \sqrt{b^2 - 4ac}}{2a}\]

\subsection{Matrix Example}
A sample matrix:
\[
\begin{pmatrix}
a_{11} & a_{12} & a_{13} \\
a_{21} & a_{22} & a_{23} \\
a_{31} & a_{32} & a_{33}
\end{pmatrix}
\]

\section{Tables}
Below is a sample table showing data:

\begin{table}[h]
\centering
\begin{tabular}{@{}lrrr@{}}
\toprule
\textbf{Category} & \textbf{Value 1} & \textbf{Value 2} & \textbf{Total} \\
\midrule
Item A & 10.5 & 20.3 & 30.8 \\
Item B & 15.2 & 25.1 & 40.3 \\
Item C & 22.1 & 12.7 & 34.8 \\
\bottomrule
\end{tabular}
\caption{Sample Data Table}
\label{tab:sample}
\end{table}

\section{Lists}
\subsection{Bullet Points}
\begin{itemize}
  \item First item
  \item Second item
  \item Third item with sub-items:
    \begin{itemize}
      \item Sub-item 1
      \item Sub-item 2
    \end{itemize}
\end{itemize}

\subsection{Numbered List}
\begin{enumerate}
  \item First step
  \item Second step
  \item Third step
\end{enumerate}

\section{Citations and References}
When you use the citation search feature, references will be automatically managed. For example, you might cite a paper like this: \cite{example2024}. All citations will be automatically added to your \texttt{cite.bib} file.

\section{Conclusion}
This template demonstrates basic \LaTeX{} features including mathematical formulas, tables, and lists. The bibliography system is now set up to work seamlessly with Bibby AI's citation management. Feel free to modify it for your needs!

\bibliographystyle{plain}
\bibliography{cite}

\end{document}\documentclass{article}
\usepackage[utf8]{inputenc}
\usepackage{booktabs}
\usepackage{amsmath}
\usepackage{natbib}

\title{My Document}
\author{Your Name}
\date{\today}

\begin{document}

\maketitle

\section{Introduction}
This is a sample document demonstrating various \LaTeX{} features including mathematical equations, tables, and formatting. You can cite papers using the citation search feature, and they will be automatically added to your \texttt{cite.bib} file.

\section{Mathematical Expressions}
Here are some example equations:

\subsection{Quadratic Formula}
The quadratic formula is given by:
\[x = \frac{-b \pm \sqrt{b^2 - 4ac}}{2a}\]

\subsection{Matrix Example}
A sample matrix:
\[
\begin{pmatrix}
a_{11} & a_{12} & a_{13} \\
a_{21} & a_{22} & a_{23} \\
a_{31} & a_{32} & a_{33}
\end{pmatrix}
\]

\section{Tables}
Below is a sample table showing data:

\begin{table}[h]
\centering
\begin{tabular}{@{}lrrr@{}}
\toprule
\textbf{Category} & \textbf{Value 1} & \textbf{Value 2} & \textbf{Total} \\
\midrule
Item A & 10.5 & 20.3 & 30.8 \\
Item B & 15.2 & 25.1 & 40.3 \\
Item C & 22.1 & 12.7 & 34.8 \\
\bottomrule
\end{tabular}
\caption{Sample Data Table}
\label{tab:sample}
\end{table}

\section{Lists}
\subsection{Bullet Points}
\begin{itemize}
  \item First item
  \item Second item
  \item Third item with sub-items:
    \begin{itemize}
      \item Sub-item 1
      \item Sub-item 2
    \end{itemize}
\end{itemize}

\subsection{Numbered List}
\begin{enumerate}
  \item First step
  \item Second step
  \item Third step
\end{enumerate}

\section{Citations and References}
When you use the citation search feature, references will be automatically managed. For example, you might cite a paper like this: \cite{example2024}. All citations will be automatically added to your \texttt{cite.bib} file.

\section{Conclusion}
This template demonstrates basic \LaTeX{} features including mathematical formulas, tables, and lists. The bibliography system is now set up to work seamlessly with Bibby AI's citation management. Feel free to modify it for your needs!

\bibliographystyle{plain}
\bibliography{cite}

\end{document}\documentclass{article}
\usepackage[utf8]{inputenc}
\usepackage{booktabs}
\usepackage{amsmath}
\usepackage{natbib}

\title{My Document}
\author{Your Name}
\date{\today}

\begin{document}

\maketitle

\section{Introduction}
This is a sample document demonstrating various \LaTeX{} features including mathematical equations, tables, and formatting. You can cite papers using the citation search feature, and they will be automatically added to your \texttt{cite.bib} file.

\section{Mathematical Expressions}
Here are some example equations:

\subsection{Quadratic Formula}
The quadratic formula is given by:
\[x = \frac{-b \pm \sqrt{b^2 - 4ac}}{2a}\]

\subsection{Matrix Example}
A sample matrix:
\[
\begin{pmatrix}
a_{11} & a_{12} & a_{13} \\
a_{21} & a_{22} & a_{23} \\
a_{31} & a_{32} & a_{33}
\end{pmatrix}
\]

\section{Tables}
Below is a sample table showing data:

\begin{table}[h]
\centering
\begin{tabular}{@{}lrrr@{}}
\toprule
\textbf{Category} & \textbf{Value 1} & \textbf{Value 2} & \textbf{Total} \\
\midrule
Item A & 10.5 & 20.3 & 30.8 \\
Item B & 15.2 & 25.1 & 40.3 \\
Item C & 22.1 & 12.7 & 34.8 \\
\bottomrule
\end{tabular}
\caption{Sample Data Table}
\label{tab:sample}
\end{table}

\section{Lists}
\subsection{Bullet Points}
\begin{itemize}
  \item First item
  \item Second item
  \item Third item with sub-items:
    \begin{itemize}
      \item Sub-item 1
      \item Sub-item 2
    \end{itemize}
\end{itemize}

\subsection{Numbered List}
\begin{enumerate}
  \item First step
  \item Second step
  \item Third step
\end{enumerate}

\section{Citations and References}
When you use the citation search feature, references will be automatically managed. For example, you might cite a paper like this: \cite{example2024}. All citations will be automatically added to your \texttt{cite.bib} file.

\section{Conclusion}
This template demonstrates basic \LaTeX{} features including mathematical formulas, tables, and lists. The bibliography system is now set up to work seamlessly with Bibby AI's citation management. Feel free to modify it for your needs!

\bibliographystyle{plain}
\bibliography{cite}

\end{document}\documentclass{article}
\usepackage[utf8]{inputenc}
\usepackage{booktabs}
\usepackage{amsmath}
\usepackage{natbib}

\title{My Document}
\author{Your Name}
\date{\today}

\begin{document}

\maketitle

\section{Introduction}
This is a sample document demonstrating various \LaTeX{} features including mathematical equations, tables, and formatting. You can cite papers using the citation search feature, and they will be automatically added to your \texttt{cite.bib} file.

\section{Mathematical Expressions}
Here are some example equations:

\subsection{Quadratic Formula}
The quadratic formula is given by:
\[x = \frac{-b \pm \sqrt{b^2 - 4ac}}{2a}\]

\subsection{Matrix Example}
A sample matrix:
\[
\begin{pmatrix}
a_{11} & a_{12} & a_{13} \\
a_{21} & a_{22} & a_{23} \\
a_{31} & a_{32} & a_{33}
\end{pmatrix}
\]

\section{Tables}
Below is a sample table showing data:

\begin{table}[h]
\centering
\begin{tabular}{@{}lrrr@{}}
\toprule
\textbf{Category} & \textbf{Value 1} & \textbf{Value 2} & \textbf{Total} \\
\midrule
Item A & 10.5 & 20.3 & 30.8 \\
Item B & 15.2 & 25.1 & 40.3 \\
Item C & 22.1 & 12.7 & 34.8 \\
\bottomrule
\end{tabular}
\caption{Sample Data Table}
\label{tab:sample}
\end{table}

\section{Lists}
\subsection{Bullet Points}
\begin{itemize}
  \item First item
  \item Second item
  \item Third item with sub-items:
    \begin{itemize}
      \item Sub-item 1
      \item Sub-item 2
    \end{itemize}
\end{itemize}

\subsection{Numbered List}
\begin{enumerate}
  \item First step
  \item Second step
  \item Third step
\end{enumerate}

\section{Citations and References}
When you use the citation search feature, references will be automatically managed. For example, you might cite a paper like this: \cite{example2024}. All citations will be automatically added to your \texttt{cite.bib} file.

\section{Conclusion}
This template demonstrates basic \LaTeX{} features including mathematical formulas, tables, and lists. The bibliography system is now set up to work seamlessly with Bibby AI's citation management. Feel free to modify it for your needs!

\bibliographystyle{plain}
\bibliography{cite}

\end{document}\documentclass{article}
\usepackage[utf8]{inputenc}
\usepackage{booktabs}
\usepackage{amsmath}
\usepackage{natbib}

\title{My Document}
\author{Your Name}
\date{\today}

\begin{document}

\maketitle

\section{Introduction}
This is a sample document demonstrating various \LaTeX{} features including mathematical equations, tables, and formatting. You can cite papers using the citation search feature, and they will be automatically added to your \texttt{cite.bib} file.

\section{Mathematical Expressions}
Here are some example equations:

\subsection{Quadratic Formula}
The quadratic formula is given by:
\[x = \frac{-b \pm \sqrt{b^2 - 4ac}}{2a}\]

\subsection{Matrix Example}
A sample matrix:
\[
\begin{pmatrix}
a_{11} & a_{12} & a_{13} \\
a_{21} & a_{22} & a_{23} \\
a_{31} & a_{32} & a_{33}
\end{pmatrix}
\]

\section{Tables}
Below is a sample table showing data:

\begin{table}[h]
\centering
\begin{tabular}{@{}lrrr@{}}
\toprule
\textbf{Category} & \textbf{Value 1} & \textbf{Value 2} & \textbf{Total} \\
\midrule
Item A & 10.5 & 20.3 & 30.8 \\
Item B & 15.2 & 25.1 & 40.3 \\
Item C & 22.1 & 12.7 & 34.8 \\
\bottomrule
\end{tabular}
\caption{Sample Data Table}
\label{tab:sample}
\end{table}

\section{Lists}
\subsection{Bullet Points}
\begin{itemize}
  \item First item
  \item Second item
  \item Third item with sub-items:
    \begin{itemize}
      \item Sub-item 1
      \item Sub-item 2
    \end{itemize}
\end{itemize}

\subsection{Numbered List}
\begin{enumerate}
  \item First step
  \item Second step
  \item Third step
\end{enumerate}

\section{Citations and References}
When you use the citation search feature, references will be automatically managed. For example, you might cite a paper like this: \cite{example2024}. All citations will be automatically added to your \texttt{cite.bib} file.

\section{Conclusion}
This template demonstrates basic \LaTeX{} features including mathematical formulas, tables, and lists. The bibliography system is now set up to work seamlessly with Bibby AI's citation management. Feel free to modify it for your needs!

\bibliographystyle{plain}
\bibliography{cite}

\end{document}\documentclass{article}
\usepackage[utf8]{inputenc}
\usepackage{booktabs}
\usepackage{amsmath}
\usepackage{natbib}

\title{My Document}
\author{Your Name}
\date{\today}

\begin{document}

\maketitle

\section{Introduction}
This is a sample document demonstrating various \LaTeX{} features including mathematical equations, tables, and formatting. You can cite papers using the citation search feature, and they will be automatically added to your \texttt{cite.bib} file.

\section{Mathematical Expressions}
Here are some example equations:

\subsection{Quadratic Formula}
The quadratic formula is given by:
\[x = \frac{-b \pm \sqrt{b^2 - 4ac}}{2a}\]

\subsection{Matrix Example}
A sample matrix:
\[
\begin{pmatrix}
a_{11} & a_{12} & a_{13} \\
a_{21} & a_{22} & a_{23} \\
a_{31} & a_{32} & a_{33}
\end{pmatrix}
\]

\section{Tables}
Below is a sample table showing data:

\begin{table}[h]
\centering
\begin{tabular}{@{}lrrr@{}}
\toprule
\textbf{Category} & \textbf{Value 1} & \textbf{Value 2} & \textbf{Total} \\
\midrule
Item A & 10.5 & 20.3 & 30.8 \\
Item B & 15.2 & 25.1 & 40.3 \\
Item C & 22.1 & 12.7 & 34.8 \\
\bottomrule
\end{tabular}
\caption{Sample Data Table}
\label{tab:sample}
\end{table}

\section{Lists}
\subsection{Bullet Points}
\begin{itemize}
  \item First item
  \item Second item
  \item Third item with sub-items:
    \begin{itemize}
      \item Sub-item 1
      \item Sub-item 2
    \end{itemize}
\end{itemize}

\subsection{Numbered List}
\begin{enumerate}
  \item First step
  \item Second step
  \item Third step
\end{enumerate}

\section{Citations and References}
When you use the citation search feature, references will be automatically managed. For example, you might cite a paper like this: \cite{example2024}. All citations will be automatically added to your \texttt{cite.bib} file.

\section{Conclusion}
This template demonstrates basic \LaTeX{} features including mathematical formulas, tables, and lists. The bibliography system is now set up to work seamlessly with Bibby AI's citation management. Feel free to modify it for your needs!

\bibliographystyle{plain}
\bibliography{cite}

\end{document}\documentclass{article}
\usepackage[utf8]{inputenc}
\usepackage{booktabs}
\usepackage{amsmath}
\usepackage{natbib}

\title{My Document}
\author{Your Name}
\date{\today}

\begin{document}

\maketitle

\section{Introduction}
This is a sample document demonstrating various \LaTeX{} features including mathematical equations, tables, and formatting. You can cite papers using the citation search feature, and they will be automatically added to your \texttt{cite.bib} file.

\section{Mathematical Expressions}
Here are some example equations:

\subsection{Quadratic Formula}
The quadratic formula is given by:
\[x = \frac{-b \pm \sqrt{b^2 - 4ac}}{2a}\]

\subsection{Matrix Example}
A sample matrix:
\[
\begin{pmatrix}
a_{11} & a_{12} & a_{13} \\
a_{21} & a_{22} & a_{23} \\
a_{31} & a_{32} & a_{33}
\end{pmatrix}
\]

\section{Tables}
Below is a sample table showing data:

\begin{table}[h]
\centering
\begin{tabular}{@{}lrrr@{}}
\toprule
\textbf{Category} & \textbf{Value 1} & \textbf{Value 2} & \textbf{Total} \\
\midrule
Item A & 10.5 & 20.3 & 30.8 \\
Item B & 15.2 & 25.1 & 40.3 \\
Item C & 22.1 & 12.7 & 34.8 \\
\bottomrule
\end{tabular}
\caption{Sample Data Table}
\label{tab:sample}
\end{table}

\section{Lists}
\subsection{Bullet Points}
\begin{itemize}
  \item First item
  \item Second item
  \item Third item with sub-items:
    \begin{itemize}
      \item Sub-item 1
      \item Sub-item 2
    \end{itemize}
\end{itemize}

\subsection{Numbered List}
\begin{enumerate}
  \item First step
  \item Second step
  \item Third step
\end{enumerate}

\section{Citations and References}
When you use the citation search feature, references will be automatically managed. For example, you might cite a paper like this: \cite{example2024}. All citations will be automatically added to your \texttt{cite.bib} file.

\section{Conclusion}
This template demonstrates basic \LaTeX{} features including mathematical formulas, tables, and lists. The bibliography system is now set up to work seamlessly with Bibby AI's citation management. Feel free to modify it for your needs!

\bibliographystyle{plain}
\bibliography{cite}

\end{document}\documentclass{article}
\usepackage[utf8]{inputenc}
\usepackage{booktabs}
\usepackage{amsmath}
\usepackage{natbib}

\title{My Document}
\author{Your Name}
\date{\today}

\begin{document}

\maketitle

\section{Introduction}
This is a sample document demonstrating various \LaTeX{} features including mathematical equations, tables, and formatting. You can cite papers using the citation search feature, and they will be automatically added to your \texttt{cite.bib} file.

\section{Mathematical Expressions}
Here are some example equations:

\subsection{Quadratic Formula}
The quadratic formula is given by:
\[x = \frac{-b \pm \sqrt{b^2 - 4ac}}{2a}\]

\subsection{Matrix Example}
A sample matrix:
\[
\begin{pmatrix}
a_{11} & a_{12} & a_{13} \\
a_{21} & a_{22} & a_{23} \\
a_{31} & a_{32} & a_{33}
\end{pmatrix}
\]

\section{Tables}
Below is a sample table showing data:

\begin{table}[h]
\centering
\begin{tabular}{@{}lrrr@{}}
\toprule
\textbf{Category} & \textbf{Value 1} & \textbf{Value 2} & \textbf{Total} \\
\midrule
Item A & 10.5 & 20.3 & 30.8 \\
Item B & 15.2 & 25.1 & 40.3 \\
Item C & 22.1 & 12.7 & 34.8 \\
\bottomrule
\end{tabular}
\caption{Sample Data Table}
\label{tab:sample}
\end{table}

\section{Lists}
\subsection{Bullet Points}
\begin{itemize}
  \item First item
  \item Second item
  \item Third item with sub-items:
    \begin{itemize}
      \item Sub-item 1
      \item Sub-item 2
    \end{itemize}
\end{itemize}

\subsection{Numbered List}
\begin{enumerate}
  \item First step
  \item Second step
  \item Third step
\end{enumerate}

\section{Citations and References}
When you use the citation search feature, references will be automatically managed. For example, you might cite a paper like this: \cite{example2024}. All citations will be automatically added to your \texttt{cite.bib} file.

\section{Conclusion}
This template demonstrates basic \LaTeX{} features including mathematical formulas, tables, and lists. The bibliography system is now set up to work seamlessly with Bibby AI's citation management. Feel free to modify it for your needs!

\bibliographystyle{plain}
\bibliography{cite}

\end{document}\documentclass{article}
\usepackage[utf8]{inputenc}
\usepackage{booktabs}
\usepackage{amsmath}
\usepackage{natbib}

\title{My Document}
\author{Your Name}
\date{\today}

\begin{document}

\maketitle

\section{Introduction}
This is a sample document demonstrating various \LaTeX{} features including mathematical equations, tables, and formatting. You can cite papers using the citation search feature, and they will be automatically added to your \texttt{cite.bib} file.

\section{Mathematical Expressions}
Here are some example equations:

\subsection{Quadratic Formula}
The quadratic formula is given by:
\[x = \frac{-b \pm \sqrt{b^2 - 4ac}}{2a}\]

\subsection{Matrix Example}
A sample matrix:
\[
\begin{pmatrix}
a_{11} & a_{12} & a_{13} \\
a_{21} & a_{22} & a_{23} \\
a_{31} & a_{32} & a_{33}
\end{pmatrix}
\]

\section{Tables}
Below is a sample table showing data:

\begin{table}[h]
\centering
\begin{tabular}{@{}lrrr@{}}
\toprule
\textbf{Category} & \textbf{Value 1} & \textbf{Value 2} & \textbf{Total} \\
\midrule
Item A & 10.5 & 20.3 & 30.8 \\
Item B & 15.2 & 25.1 & 40.3 \\
Item C & 22.1 & 12.7 & 34.8 \\
\bottomrule
\end{tabular}
\caption{Sample Data Table}
\label{tab:sample}
\end{table}

\section{Lists}
\subsection{Bullet Points}
\begin{itemize}
  \item First item
  \item Second item
  \item Third item with sub-items:
    \begin{itemize}
      \item Sub-item 1
      \item Sub-item 2
    \end{itemize}
\end{itemize}

\subsection{Numbered List}
\begin{enumerate}
  \item First step
  \item Second step
  \item Third step
\end{enumerate}

\section{Citations and References}
When you use the citation search feature, references will be automatically managed. For example, you might cite a paper like this: \cite{example2024}. All citations will be automatically added to your \texttt{cite.bib} file.

\section{Conclusion}
This template demonstrates basic \LaTeX{} features including mathematical formulas, tables, and lists. The bibliography system is now set up to work seamlessly with Bibby AI's citation management. Feel free to modify it for your needs!

\bibliographystyle{plain}
\bibliography{cite}

\end{document}\documentclass{article}
\usepackage[utf8]{inputenc}
\usepackage{booktabs}
\usepackage{amsmath}
\usepackage{natbib}

\title{My Document}
\author{Your Name}
\date{\today}

\begin{document}

\maketitle

\section{Introduction}
This is a sample document demonstrating various \LaTeX{} features including mathematical equations, tables, and formatting. You can cite papers using the citation search feature, and they will be automatically added to your \texttt{cite.bib} file.

\section{Mathematical Expressions}
Here are some example equations:

\subsection{Quadratic Formula}
The quadratic formula is given by:
\[x = \frac{-b \pm \sqrt{b^2 - 4ac}}{2a}\]

\subsection{Matrix Example}
A sample matrix:
\[
\begin{pmatrix}
a_{11} & a_{12} & a_{13} \\
a_{21} & a_{22} & a_{23} \\
a_{31} & a_{32} & a_{33}
\end{pmatrix}
\]

\section{Tables}
Below is a sample table showing data:

\begin{table}[h]
\centering
\begin{tabular}{@{}lrrr@{}}
\toprule
\textbf{Category} & \textbf{Value 1} & \textbf{Value 2} & \textbf{Total} \\
\midrule
Item A & 10.5 & 20.3 & 30.8 \\
Item B & 15.2 & 25.1 & 40.3 \\
Item C & 22.1 & 12.7 & 34.8 \\
\bottomrule
\end{tabular}
\caption{Sample Data Table}
\label{tab:sample}
\end{table}

\section{Lists}
\subsection{Bullet Points}
\begin{itemize}
  \item First item
  \item Second item
  \item Third item with sub-items:
    \begin{itemize}
      \item Sub-item 1
      \item Sub-item 2
    \end{itemize}
\end{itemize}

\subsection{Numbered List}
\begin{enumerate}
  \item First step
  \item Second step
  \item Third step
\end{enumerate}

\section{Citations and References}
When you use the citation search feature, references will be automatically managed. For example, you might cite a paper like this: \cite{example2024}. All citations will be automatically added to your \texttt{cite.bib} file.

\section{Conclusion}
This template demonstrates basic \LaTeX{} features including mathematical formulas, tables, and lists. The bibliography system is now set up to work seamlessly with Bibby AI's citation management. Feel free to modify it for your needs!

\bibliographystyle{plain}
\bibliography{cite}

\end{document}\documentclass{article}
\usepackage[utf8]{inputenc}
\usepackage{booktabs}
\usepackage{amsmath}
\usepackage{natbib}

\title{My Document}
\author{Your Name}
\date{\today}

\begin{document}

\maketitle

\section{Introduction}
This is a sample document demonstrating various \LaTeX{} features including mathematical equations, tables, and formatting. You can cite papers using the citation search feature, and they will be automatically added to your \texttt{cite.bib} file.

\section{Mathematical Expressions}
Here are some example equations:

\subsection{Quadratic Formula}
The quadratic formula is given by:
\[x = \frac{-b \pm \sqrt{b^2 - 4ac}}{2a}\]

\subsection{Matrix Example}
A sample matrix:
\[
\begin{pmatrix}
a_{11} & a_{12} & a_{13} \\
a_{21} & a_{22} & a_{23} \\
a_{31} & a_{32} & a_{33}
\end{pmatrix}
\]

\section{Tables}
Below is a sample table showing data:

\begin{table}[h]
\centering
\begin{tabular}{@{}lrrr@{}}
\toprule
\textbf{Category} & \textbf{Value 1} & \textbf{Value 2} & \textbf{Total} \\
\midrule
Item A & 10.5 & 20.3 & 30.8 \\
Item B & 15.2 & 25.1 & 40.3 \\
Item C & 22.1 & 12.7 & 34.8 \\
\bottomrule
\end{tabular}
\caption{Sample Data Table}
\label{tab:sample}
\end{table}

\section{Lists}
\subsection{Bullet Points}
\begin{itemize}
  \item First item
  \item Second item
  \item Third item with sub-items:
    \begin{itemize}
      \item Sub-item 1
      \item Sub-item 2
    \end{itemize}
\end{itemize}

\subsection{Numbered List}
\begin{enumerate}
  \item First step
  \item Second step
  \item Third step
\end{enumerate}

\section{Citations and References}
When you use the citation search feature, references will be automatically managed. For example, you might cite a paper like this: \cite{example2024}. All citations will be automatically added to your \texttt{cite.bib} file.

\section{Conclusion}
This template demonstrates basic \LaTeX{} features including mathematical formulas, tables, and lists. The bibliography system is now set up to work seamlessly with Bibby AI's citation management. Feel free to modify it for your needs!

\bibliographystyle{plain}
\bibliography{cite}

\end{document}\documentclass{article}
\usepackage[utf8]{inputenc}
\usepackage{booktabs}
\usepackage{amsmath}
\usepackage{natbib}

\title{My Document}
\author{Your Name}
\date{\today}

\begin{document}

\maketitle

\section{Introduction}
This is a sample document demonstrating various \LaTeX{} features including mathematical equations, tables, and formatting. You can cite papers using the citation search feature, and they will be automatically added to your \texttt{cite.bib} file.

\section{Mathematical Expressions}
Here are some example equations:

\subsection{Quadratic Formula}
The quadratic formula is given by:
\[x = \frac{-b \pm \sqrt{b^2 - 4ac}}{2a}\]

\subsection{Matrix Example}
A sample matrix:
\[
\begin{pmatrix}
a_{11} & a_{12} & a_{13} \\
a_{21} & a_{22} & a_{23} \\
a_{31} & a_{32} & a_{33}
\end{pmatrix}
\]

\section{Tables}
Below is a sample table showing data:

\begin{table}[h]
\centering
\begin{tabular}{@{}lrrr@{}}
\toprule
\textbf{Category} & \textbf{Value 1} & \textbf{Value 2} & \textbf{Total} \\
\midrule
Item A & 10.5 & 20.3 & 30.8 \\
Item B & 15.2 & 25.1 & 40.3 \\
Item C & 22.1 & 12.7 & 34.8 \\
\bottomrule
\end{tabular}
\caption{Sample Data Table}
\label{tab:sample}
\end{table}

\section{Lists}
\subsection{Bullet Points}
\begin{itemize}
  \item First item
  \item Second item
  \item Third item with sub-items:
    \begin{itemize}
      \item Sub-item 1
      \item Sub-item 2
    \end{itemize}
\end{itemize}

\subsection{Numbered List}
\begin{enumerate}
  \item First step
  \item Second step
  \item Third step
\end{enumerate}

\section{Citations and References}
When you use the citation search feature, references will be automatically managed. For example, you might cite a paper like this: \cite{example2024}. All citations will be automatically added to your \texttt{cite.bib} file.

\section{Conclusion}
This template demonstrates basic \LaTeX{} features including mathematical formulas, tables, and lists. The bibliography system is now set up to work seamlessly with Bibby AI's citation management. Feel free to modify it for your needs!

\bibliographystyle{plain}
\bibliography{cite}

\end{document}\documentclass{article}
\usepackage[utf8]{inputenc}
\usepackage{booktabs}
\usepackage{amsmath}
\usepackage{natbib}

\title{My Document}
\author{Your Name}
\date{\today}

\begin{document}

\maketitle

\section{Introduction}
This is a sample document demonstrating various \LaTeX{} features including mathematical equations, tables, and formatting. You can cite papers using the citation search feature, and they will be automatically added to your \texttt{cite.bib} file.

\section{Mathematical Expressions}
Here are some example equations:

\subsection{Quadratic Formula}
The quadratic formula is given by:
\[x = \frac{-b \pm \sqrt{b^2 - 4ac}}{2a}\]

\subsection{Matrix Example}
A sample matrix:
\[
\begin{pmatrix}
a_{11} & a_{12} & a_{13} \\
a_{21} & a_{22} & a_{23} \\
a_{31} & a_{32} & a_{33}
\end{pmatrix}
\]

\section{Tables}
Below is a sample table showing data:

\begin{table}[h]
\centering
\begin{tabular}{@{}lrrr@{}}
\toprule
\textbf{Category} & \textbf{Value 1} & \textbf{Value 2} & \textbf{Total} \\
\midrule
Item A & 10.5 & 20.3 & 30.8 \\
Item B & 15.2 & 25.1 & 40.3 \\
Item C & 22.1 & 12.7 & 34.8 \\
\bottomrule
\end{tabular}
\caption{Sample Data Table}
\label{tab:sample}
\end{table}

\section{Lists}
\subsection{Bullet Points}
\begin{itemize}
  \item First item
  \item Second item
  \item Third item with sub-items:
    \begin{itemize}
      \item Sub-item 1
      \item Sub-item 2
    \end{itemize}
\end{itemize}

\subsection{Numbered List}
\begin{enumerate}
  \item First step
  \item Second step
  \item Third step
\end{enumerate}

\section{Citations and References}
When you use the citation search feature, references will be automatically managed. For example, you might cite a paper like this: \cite{example2024}. All citations will be automatically added to your \texttt{cite.bib} file.

\section{Conclusion}
This template demonstrates basic \LaTeX{} features including mathematical formulas, tables, and lists. The bibliography system is now set up to work seamlessly with Bibby AI's citation management. Feel free to modify it for your needs!

\bibliographystyle{plain}
\bibliography{cite}

\end{document}\documentclass{article}
\usepackage[utf8]{inputenc}
\usepackage{booktabs}
\usepackage{amsmath}
\usepackage{natbib}

\title{My Document}
\author{Your Name}
\date{\today}

\begin{document}

\maketitle

\section{Introduction}
This is a sample document demonstrating various \LaTeX{} features including mathematical equations, tables, and formatting. You can cite papers using the citation search feature, and they will be automatically added to your \texttt{cite.bib} file.

\section{Mathematical Expressions}
Here are some example equations:

\subsection{Quadratic Formula}
The quadratic formula is given by:
\[x = \frac{-b \pm \sqrt{b^2 - 4ac}}{2a}\]

\subsection{Matrix Example}
A sample matrix:
\[
\begin{pmatrix}
a_{11} & a_{12} & a_{13} \\
a_{21} & a_{22} & a_{23} \\
a_{31} & a_{32} & a_{33}
\end{pmatrix}
\]

\section{Tables}
Below is a sample table showing data:

\begin{table}[h]
\centering
\begin{tabular}{@{}lrrr@{}}
\toprule
\textbf{Category} & \textbf{Value 1} & \textbf{Value 2} & \textbf{Total} \\
\midrule
Item A & 10.5 & 20.3 & 30.8 \\
Item B & 15.2 & 25.1 & 40.3 \\
Item C & 22.1 & 12.7 & 34.8 \\
\bottomrule
\end{tabular}
\caption{Sample Data Table}
\label{tab:sample}
\end{table}

\section{Lists}
\subsection{Bullet Points}
\begin{itemize}
  \item First item
  \item Second item
  \item Third item with sub-items:
    \begin{itemize}
      \item Sub-item 1
      \item Sub-item 2
    \end{itemize}
\end{itemize}

\subsection{Numbered List}
\begin{enumerate}
  \item First step
  \item Second step
  \item Third step
\end{enumerate}

\section{Citations and References}
When you use the citation search feature, references will be automatically managed. For example, you might cite a paper like this: \cite{example2024}. All citations will be automatically added to your \texttt{cite.bib} file.

\section{Conclusion}
This template demonstrates basic \LaTeX{} features including mathematical formulas, tables, and lists. The bibliography system is now set up to work seamlessly with Bibby AI's citation management. Feel free to modify it for your needs!

\bibliographystyle{plain}
\bibliography{cite}

\end{document}\documentclass{article}
\usepackage[utf8]{inputenc}
\usepackage{booktabs}
\usepackage{amsmath}
\usepackage{natbib}

\title{My Document}
\author{Your Name}
\date{\today}

\begin{document}

\maketitle

\section{Introduction}
This is a sample document demonstrating various \LaTeX{} features including mathematical equations, tables, and formatting. You can cite papers using the citation search feature, and they will be automatically added to your \texttt{cite.bib} file.

\section{Mathematical Expressions}
Here are some example equations:

\subsection{Quadratic Formula}
The quadratic formula is given by:
\[x = \frac{-b \pm \sqrt{b^2 - 4ac}}{2a}\]

\subsection{Matrix Example}
A sample matrix:
\[
\begin{pmatrix}
a_{11} & a_{12} & a_{13} \\
a_{21} & a_{22} & a_{23} \\
a_{31} & a_{32} & a_{33}
\end{pmatrix}
\]

\section{Tables}
Below is a sample table showing data:

\begin{table}[h]
\centering
\begin{tabular}{@{}lrrr@{}}
\toprule
\textbf{Category} & \textbf{Value 1} & \textbf{Value 2} & \textbf{Total} \\
\midrule
Item A & 10.5 & 20.3 & 30.8 \\
Item B & 15.2 & 25.1 & 40.3 \\
Item C & 22.1 & 12.7 & 34.8 \\
\bottomrule
\end{tabular}
\caption{Sample Data Table}
\label{tab:sample}
\end{table}

\section{Lists}
\subsection{Bullet Points}
\begin{itemize}
  \item First item
  \item Second item
  \item Third item with sub-items:
    \begin{itemize}
      \item Sub-item 1
      \item Sub-item 2
    \end{itemize}
\end{itemize}

\subsection{Numbered List}
\begin{enumerate}
  \item First step
  \item Second step
  \item Third step
\end{enumerate}

\section{Citations and References}
When you use the citation search feature, references will be automatically managed. For example, you might cite a paper like this: \cite{example2024}. All citations will be automatically added to your \texttt{cite.bib} file.

\section{Conclusion}
This template demonstrates basic \LaTeX{} features including mathematical formulas, tables, and lists. The bibliography system is now set up to work seamlessly with Bibby AI's citation management. Feel free to modify it for your needs!

\bibliographystyle{plain}
\bibliography{cite}

\end{document}\documentclass{article}
\usepackage[utf8]{inputenc}
\usepackage{booktabs}
\usepackage{amsmath}
\usepackage{natbib}

\title{My Document}
\author{Your Name}
\date{\today}

\begin{document}

\maketitle

\section{Introduction}
This is a sample document demonstrating various \LaTeX{} features including mathematical equations, tables, and formatting. You can cite papers using the citation search feature, and they will be automatically added to your \texttt{cite.bib} file.

\section{Mathematical Expressions}
Here are some example equations:

\subsection{Quadratic Formula}
The quadratic formula is given by:
\[x = \frac{-b \pm \sqrt{b^2 - 4ac}}{2a}\]

\subsection{Matrix Example}
A sample matrix:
\[
\begin{pmatrix}
a_{11} & a_{12} & a_{13} \\
a_{21} & a_{22} & a_{23} \\
a_{31} & a_{32} & a_{33}
\end{pmatrix}
\]

\section{Tables}
Below is a sample table showing data:

\begin{table}[h]
\centering
\begin{tabular}{@{}lrrr@{}}
\toprule
\textbf{Category} & \textbf{Value 1} & \textbf{Value 2} & \textbf{Total} \\
\midrule
Item A & 10.5 & 20.3 & 30.8 \\
Item B & 15.2 & 25.1 & 40.3 \\
Item C & 22.1 & 12.7 & 34.8 \\
\bottomrule
\end{tabular}
\caption{Sample Data Table}
\label{tab:sample}
\end{table}

\section{Lists}
\subsection{Bullet Points}
\begin{itemize}
  \item First item
  \item Second item
  \item Third item with sub-items:
    \begin{itemize}
      \item Sub-item 1
      \item Sub-item 2
    \end{itemize}
\end{itemize}

\subsection{Numbered List}
\begin{enumerate}
  \item First step
  \item Second step
  \item Third step
\end{enumerate}

\section{Citations and References}
When you use the citation search feature, references will be automatically managed. For example, you might cite a paper like this: \cite{example2024}. All citations will be automatically added to your \texttt{cite.bib} file.

\section{Conclusion}
This template demonstrates basic \LaTeX{} features including mathematical formulas, tables, and lists. The bibliography system is now set up to work seamlessly with Bibby AI's citation management. Feel free to modify it for your needs!

\bibliographystyle{plain}
\bibliography{cite}

\end{document}\documentclass{article}
\usepackage[utf8]{inputenc}
\usepackage{booktabs}
\usepackage{amsmath}
\usepackage{natbib}

\title{My Document}
\author{Your Name}
\date{\today}

\begin{document}

\maketitle

\section{Introduction}
This is a sample document demonstrating various \LaTeX{} features including mathematical equations, tables, and formatting. You can cite papers using the citation search feature, and they will be automatically added to your \texttt{cite.bib} file.

\section{Mathematical Expressions}
Here are some example equations:

\subsection{Quadratic Formula}
The quadratic formula is given by:
\[x = \frac{-b \pm \sqrt{b^2 - 4ac}}{2a}\]

\subsection{Matrix Example}
A sample matrix:
\[
\begin{pmatrix}
a_{11} & a_{12} & a_{13} \\
a_{21} & a_{22} & a_{23} \\
a_{31} & a_{32} & a_{33}
\end{pmatrix}
\]

\section{Tables}
Below is a sample table showing data:

\begin{table}[h]
\centering
\begin{tabular}{@{}lrrr@{}}
\toprule
\textbf{Category} & \textbf{Value 1} & \textbf{Value 2} & \textbf{Total} \\
\midrule
Item A & 10.5 & 20.3 & 30.8 \\
Item B & 15.2 & 25.1 & 40.3 \\
Item C & 22.1 & 12.7 & 34.8 \\
\bottomrule
\end{tabular}
\caption{Sample Data Table}
\label{tab:sample}
\end{table}

\section{Lists}
\subsection{Bullet Points}
\begin{itemize}
  \item First item
  \item Second item
  \item Third item with sub-items:
    \begin{itemize}
      \item Sub-item 1
      \item Sub-item 2
    \end{itemize}
\end{itemize}

\subsection{Numbered List}
\begin{enumerate}
  \item First step
  \item Second step
  \item Third step
\end{enumerate}

\section{Citations and References}
When you use the citation search feature, references will be automatically managed. For example, you might cite a paper like this: \cite{example2024}. All citations will be automatically added to your \texttt{cite.bib} file.

\section{Conclusion}
This template demonstrates basic \LaTeX{} features including mathematical formulas, tables, and lists. The bibliography system is now set up to work seamlessly with Bibby AI's citation management. Feel free to modify it for your needs!

\bibliographystyle{plain}
\bibliography{cite}

\end{document}\documentclass{article}
\usepackage[utf8]{inputenc}
\usepackage{booktabs}
\usepackage{amsmath}
\usepackage{natbib}

\title{My Document}
\author{Your Name}
\date{\today}

\begin{document}

\maketitle

\section{Introduction}
This is a sample document demonstrating various \LaTeX{} features including mathematical equations, tables, and formatting. You can cite papers using the citation search feature, and they will be automatically added to your \texttt{cite.bib} file.

\section{Mathematical Expressions}
Here are some example equations:

\subsection{Quadratic Formula}
The quadratic formula is given by:
\[x = \frac{-b \pm \sqrt{b^2 - 4ac}}{2a}\]

\subsection{Matrix Example}
A sample matrix:
\[
\begin{pmatrix}
a_{11} & a_{12} & a_{13} \\
a_{21} & a_{22} & a_{23} \\
a_{31} & a_{32} & a_{33}
\end{pmatrix}
\]

\section{Tables}
Below is a sample table showing data:

\begin{table}[h]
\centering
\begin{tabular}{@{}lrrr@{}}
\toprule
\textbf{Category} & \textbf{Value 1} & \textbf{Value 2} & \textbf{Total} \\
\midrule
Item A & 10.5 & 20.3 & 30.8 \\
Item B & 15.2 & 25.1 & 40.3 \\
Item C & 22.1 & 12.7 & 34.8 \\
\bottomrule
\end{tabular}
\caption{Sample Data Table}
\label{tab:sample}
\end{table}

\section{Lists}
\subsection{Bullet Points}
\begin{itemize}
  \item First item
  \item Second item
  \item Third item with sub-items:
    \begin{itemize}
      \item Sub-item 1
      \item Sub-item 2
    \end{itemize}
\end{itemize}

\subsection{Numbered List}
\begin{enumerate}
  \item First step
  \item Second step
  \item Third step
\end{enumerate}

\section{Citations and References}
When you use the citation search feature, references will be automatically managed. For example, you might cite a paper like this: \cite{example2024}. All citations will be automatically added to your \texttt{cite.bib} file.

\section{Conclusion}
This template demonstrates basic \LaTeX{} features including mathematical formulas, tables, and lists. The bibliography system is now set up to work seamlessly with Bibby AI's citation management. Feel free to modify it for your needs!

\bibliographystyle{plain}
\bibliography{cite}

\end{document}\documentclass{article}
\usepackage[utf8]{inputenc}
\usepackage{booktabs}
\usepackage{amsmath}
\usepackage{natbib}

\title{My Document}
\author{Your Name}
\date{\today}

\begin{document}

\maketitle

\section{Introduction}
This is a sample document demonstrating various \LaTeX{} features including mathematical equations, tables, and formatting. You can cite papers using the citation search feature, and they will be automatically added to your \texttt{cite.bib} file.

\section{Mathematical Expressions}
Here are some example equations:

\subsection{Quadratic Formula}
The quadratic formula is given by:
\[x = \frac{-b \pm \sqrt{b^2 - 4ac}}{2a}\]

\subsection{Matrix Example}
A sample matrix:
\[
\begin{pmatrix}
a_{11} & a_{12} & a_{13} \\
a_{21} & a_{22} & a_{23} \\
a_{31} & a_{32} & a_{33}
\end{pmatrix}
\]

\section{Tables}
Below is a sample table showing data:

\begin{table}[h]
\centering
\begin{tabular}{@{}lrrr@{}}
\toprule
\textbf{Category} & \textbf{Value 1} & \textbf{Value 2} & \textbf{Total} \\
\midrule
Item A & 10.5 & 20.3 & 30.8 \\
Item B & 15.2 & 25.1 & 40.3 \\
Item C & 22.1 & 12.7 & 34.8 \\
\bottomrule
\end{tabular}
\caption{Sample Data Table}
\label{tab:sample}
\end{table}

\section{Lists}
\subsection{Bullet Points}
\begin{itemize}
  \item First item
  \item Second item
  \item Third item with sub-items:
    \begin{itemize}
      \item Sub-item 1
      \item Sub-item 2
    \end{itemize}
\end{itemize}

\subsection{Numbered List}
\begin{enumerate}
  \item First step
  \item Second step
  \item Third step
\end{enumerate}

\section{Citations and References}
When you use the citation search feature, references will be automatically managed. For example, you might cite a paper like this: \cite{example2024}. All citations will be automatically added to your \texttt{cite.bib} file.

\section{Conclusion}
This template demonstrates basic \LaTeX{} features including mathematical formulas, tables, and lists. The bibliography system is now set up to work seamlessly with Bibby AI's citation management. Feel free to modify it for your needs!

\bibliographystyle{plain}
\bibliography{cite}

\end{document}\documentclass{article}
\usepackage[utf8]{inputenc}
\usepackage{booktabs}
\usepackage{amsmath}
\usepackage{natbib}

\title{My Document}
\author{Your Name}
\date{\today}

\begin{document}

\maketitle

\section{Introduction}
This is a sample document demonstrating various \LaTeX{} features including mathematical equations, tables, and formatting. You can cite papers using the citation search feature, and they will be automatically added to your \texttt{cite.bib} file.

\section{Mathematical Expressions}
Here are some example equations:

\subsection{Quadratic Formula}
The quadratic formula is given by:
\[x = \frac{-b \pm \sqrt{b^2 - 4ac}}{2a}\]

\subsection{Matrix Example}
A sample matrix:
\[
\begin{pmatrix}
a_{11} & a_{12} & a_{13} \\
a_{21} & a_{22} & a_{23} \\
a_{31} & a_{32} & a_{33}
\end{pmatrix}
\]

\section{Tables}
Below is a sample table showing data:

\begin{table}[h]
\centering
\begin{tabular}{@{}lrrr@{}}
\toprule
\textbf{Category} & \textbf{Value 1} & \textbf{Value 2} & \textbf{Total} \\
\midrule
Item A & 10.5 & 20.3 & 30.8 \\
Item B & 15.2 & 25.1 & 40.3 \\
Item C & 22.1 & 12.7 & 34.8 \\
\bottomrule
\end{tabular}
\caption{Sample Data Table}
\label{tab:sample}
\end{table}

\section{Lists}
\subsection{Bullet Points}
\begin{itemize}
  \item First item
  \item Second item
  \item Third item with sub-items:
    \begin{itemize}
      \item Sub-item 1
      \item Sub-item 2
    \end{itemize}
\end{itemize}

\subsection{Numbered List}
\begin{enumerate}
  \item First step
  \item Second step
  \item Third step
\end{enumerate}

\section{Citations and References}
When you use the citation search feature, references will be automatically managed. For example, you might cite a paper like this: \cite{example2024}. All citations will be automatically added to your \texttt{cite.bib} file.

\section{Conclusion}
This template demonstrates basic \LaTeX{} features including mathematical formulas, tables, and lists. The bibliography system is now set up to work seamlessly with Bibby AI's citation management. Feel free to modify it for your needs!

\bibliographystyle{plain}
\bibliography{cite}

\end{document}\documentclass{article}
\usepackage[utf8]{inputenc}
\usepackage{booktabs}
\usepackage{amsmath}
\usepackage{natbib}

\title{My Document}
\author{Your Name}
\date{\today}

\begin{document}

\maketitle

\section{Introduction}
This is a sample document demonstrating various \LaTeX{} features including mathematical equations, tables, and formatting. You can cite papers using the citation search feature, and they will be automatically added to your \texttt{cite.bib} file.

\section{Mathematical Expressions}
Here are some example equations:

\subsection{Quadratic Formula}
The quadratic formula is given by:
\[x = \frac{-b \pm \sqrt{b^2 - 4ac}}{2a}\]

\subsection{Matrix Example}
A sample matrix:
\[
\begin{pmatrix}
a_{11} & a_{12} & a_{13} \\
a_{21} & a_{22} & a_{23} \\
a_{31} & a_{32} & a_{33}
\end{pmatrix}
\]

\section{Tables}
Below is a sample table showing data:

\begin{table}[h]
\centering
\begin{tabular}{@{}lrrr@{}}
\toprule
\textbf{Category} & \textbf{Value 1} & \textbf{Value 2} & \textbf{Total} \\
\midrule
Item A & 10.5 & 20.3 & 30.8 \\
Item B & 15.2 & 25.1 & 40.3 \\
Item C & 22.1 & 12.7 & 34.8 \\
\bottomrule
\end{tabular}
\caption{Sample Data Table}
\label{tab:sample}
\end{table}

\section{Lists}
\subsection{Bullet Points}
\begin{itemize}
  \item First item
  \item Second item
  \item Third item with sub-items:
    \begin{itemize}
      \item Sub-item 1
      \item Sub-item 2
    \end{itemize}
\end{itemize}

\subsection{Numbered List}
\begin{enumerate}
  \item First step
  \item Second step
  \item Third step
\end{enumerate}

\section{Citations and References}
When you use the citation search feature, references will be automatically managed. For example, you might cite a paper like this: \cite{example2024}. All citations will be automatically added to your \texttt{cite.bib} file.

\section{Conclusion}
This template demonstrates basic \LaTeX{} features including mathematical formulas, tables, and lists. The bibliography system is now set up to work seamlessly with Bibby AI's citation management. Feel free to modify it for your needs!

\bibliographystyle{plain}
\bibliography{cite}

\end{document}\documentclass{article}
\usepackage[utf8]{inputenc}
\usepackage{booktabs}
\usepackage{amsmath}
\usepackage{natbib}

\title{My Document}
\author{Your Name}
\date{\today}

\begin{document}

\maketitle

\section{Introduction}
This is a sample document demonstrating various \LaTeX{} features including mathematical equations, tables, and formatting. You can cite papers using the citation search feature, and they will be automatically added to your \texttt{cite.bib} file.

\section{Mathematical Expressions}
Here are some example equations:

\subsection{Quadratic Formula}
The quadratic formula is given by:
\[x = \frac{-b \pm \sqrt{b^2 - 4ac}}{2a}\]

\subsection{Matrix Example}
A sample matrix:
\[
\begin{pmatrix}
a_{11} & a_{12} & a_{13} \\
a_{21} & a_{22} & a_{23} \\
a_{31} & a_{32} & a_{33}
\end{pmatrix}
\]

\section{Tables}
Below is a sample table showing data:

\begin{table}[h]
\centering
\begin{tabular}{@{}lrrr@{}}
\toprule
\textbf{Category} & \textbf{Value 1} & \textbf{Value 2} & \textbf{Total} \\
\midrule
Item A & 10.5 & 20.3 & 30.8 \\
Item B & 15.2 & 25.1 & 40.3 \\
Item C & 22.1 & 12.7 & 34.8 \\
\bottomrule
\end{tabular}
\caption{Sample Data Table}
\label{tab:sample}
\end{table}

\section{Lists}
\subsection{Bullet Points}
\begin{itemize}
  \item First item
  \item Second item
  \item Third item with sub-items:
    \begin{itemize}
      \item Sub-item 1
      \item Sub-item 2
    \end{itemize}
\end{itemize}

\subsection{Numbered List}
\begin{enumerate}
  \item First step
  \item Second step
  \item Third step
\end{enumerate}

\section{Citations and References}
When you use the citation search feature, references will be automatically managed. For example, you might cite a paper like this: \cite{example2024}. All citations will be automatically added to your \texttt{cite.bib} file.

\section{Conclusion}
This template demonstrates basic \LaTeX{} features including mathematical formulas, tables, and lists. The bibliography system is now set up to work seamlessly with Bibby AI's citation management. Feel free to modify it for your needs!

\bibliographystyle{plain}
\bibliography{cite}

\end{document}\documentclass{article}
\usepackage[utf8]{inputenc}
\usepackage{booktabs}
\usepackage{amsmath}
\usepackage{natbib}

\title{My Document}
\author{Your Name}
\date{\today}

\begin{document}

\maketitle

\section{Introduction}
This is a sample document demonstrating various \LaTeX{} features including mathematical equations, tables, and formatting. You can cite papers using the citation search feature, and they will be automatically added to your \texttt{cite.bib} file.

\section{Mathematical Expressions}
Here are some example equations:

\subsection{Quadratic Formula}
The quadratic formula is given by:
\[x = \frac{-b \pm \sqrt{b^2 - 4ac}}{2a}\]

\subsection{Matrix Example}
A sample matrix:
\[
\begin{pmatrix}
a_{11} & a_{12} & a_{13} \\
a_{21} & a_{22} & a_{23} \\
a_{31} & a_{32} & a_{33}
\end{pmatrix}
\]

\section{Tables}
Below is a sample table showing data:

\begin{table}[h]
\centering
\begin{tabular}{@{}lrrr@{}}
\toprule
\textbf{Category} & \textbf{Value 1} & \textbf{Value 2} & \textbf{Total} \\
\midrule
Item A & 10.5 & 20.3 & 30.8 \\
Item B & 15.2 & 25.1 & 40.3 \\
Item C & 22.1 & 12.7 & 34.8 \\
\bottomrule
\end{tabular}
\caption{Sample Data Table}
\label{tab:sample}
\end{table}

\section{Lists}
\subsection{Bullet Points}
\begin{itemize}
  \item First item
  \item Second item
  \item Third item with sub-items:
    \begin{itemize}
      \item Sub-item 1
      \item Sub-item 2
    \end{itemize}
\end{itemize}

\subsection{Numbered List}
\begin{enumerate}
  \item First step
  \item Second step
  \item Third step
\end{enumerate}

\section{Citations and References}
When you use the citation search feature, references will be automatically managed. For example, you might cite a paper like this: \cite{example2024}. All citations will be automatically added to your \texttt{cite.bib} file.

\section{Conclusion}
This template demonstrates basic \LaTeX{} features including mathematical formulas, tables, and lists. The bibliography system is now set up to work seamlessly with Bibby AI's citation management. Feel free to modify it for your needs!

\bibliographystyle{plain}
\bibliography{cite}

\end{document}\documentclass{article}
\usepackage[utf8]{inputenc}
\usepackage{booktabs}
\usepackage{amsmath}
\usepackage{natbib}

\title{My Document}
\author{Your Name}
\date{\today}

\begin{document}

\maketitle

\section{Introduction}
This is a sample document demonstrating various \LaTeX{} features including mathematical equations, tables, and formatting. You can cite papers using the citation search feature, and they will be automatically added to your \texttt{cite.bib} file.

\section{Mathematical Expressions}
Here are some example equations:

\subsection{Quadratic Formula}
The quadratic formula is given by:
\[x = \frac{-b \pm \sqrt{b^2 - 4ac}}{2a}\]

\subsection{Matrix Example}
A sample matrix:
\[
\begin{pmatrix}
a_{11} & a_{12} & a_{13} \\
a_{21} & a_{22} & a_{23} \\
a_{31} & a_{32} & a_{33}
\end{pmatrix}
\]

\section{Tables}
Below is a sample table showing data:

\begin{table}[h]
\centering
\begin{tabular}{@{}lrrr@{}}
\toprule
\textbf{Category} & \textbf{Value 1} & \textbf{Value 2} & \textbf{Total} \\
\midrule
Item A & 10.5 & 20.3 & 30.8 \\
Item B & 15.2 & 25.1 & 40.3 \\
Item C & 22.1 & 12.7 & 34.8 \\
\bottomrule
\end{tabular}
\caption{Sample Data Table}
\label{tab:sample}
\end{table}

\section{Lists}
\subsection{Bullet Points}
\begin{itemize}
  \item First item
  \item Second item
  \item Third item with sub-items:
    \begin{itemize}
      \item Sub-item 1
      \item Sub-item 2
    \end{itemize}
\end{itemize}

\subsection{Numbered List}
\begin{enumerate}
  \item First step
  \item Second step
  \item Third step
\end{enumerate}

\section{Citations and References}
When you use the citation search feature, references will be automatically managed. For example, you might cite a paper like this: \cite{example2024}. All citations will be automatically added to your \texttt{cite.bib} file.

\section{Conclusion}
This template demonstrates basic \LaTeX{} features including mathematical formulas, tables, and lists. The bibliography system is now set up to work seamlessly with Bibby AI's citation management. Feel free to modify it for your needs!

\bibliographystyle{plain}
\bibliography{cite}

\end{document}\documentclass{article}
\usepackage[utf8]{inputenc}
\usepackage{booktabs}
\usepackage{amsmath}
\usepackage{natbib}

\title{My Document}
\author{Your Name}
\date{\today}

\begin{document}

\maketitle

\section{Introduction}
This is a sample document demonstrating various \LaTeX{} features including mathematical equations, tables, and formatting. You can cite papers using the citation search feature, and they will be automatically added to your \texttt{cite.bib} file.

\section{Mathematical Expressions}
Here are some example equations:

\subsection{Quadratic Formula}
The quadratic formula is given by:
\[x = \frac{-b \pm \sqrt{b^2 - 4ac}}{2a}\]

\subsection{Matrix Example}
A sample matrix:
\[
\begin{pmatrix}
a_{11} & a_{12} & a_{13} \\
a_{21} & a_{22} & a_{23} \\
a_{31} & a_{32} & a_{33}
\end{pmatrix}
\]

\section{Tables}
Below is a sample table showing data:

\begin{table}[h]
\centering
\begin{tabular}{@{}lrrr@{}}
\toprule
\textbf{Category} & \textbf{Value 1} & \textbf{Value 2} & \textbf{Total} \\
\midrule
Item A & 10.5 & 20.3 & 30.8 \\
Item B & 15.2 & 25.1 & 40.3 \\
Item C & 22.1 & 12.7 & 34.8 \\
\bottomrule
\end{tabular}
\caption{Sample Data Table}
\label{tab:sample}
\end{table}

\section{Lists}
\subsection{Bullet Points}
\begin{itemize}
  \item First item
  \item Second item
  \item Third item with sub-items:
    \begin{itemize}
      \item Sub-item 1
      \item Sub-item 2
    \end{itemize}
\end{itemize}

\subsection{Numbered List}
\begin{enumerate}
  \item First step
  \item Second step
  \item Third step
\end{enumerate}

\section{Citations and References}
When you use the citation search feature, references will be automatically managed. For example, you might cite a paper like this: \cite{example2024}. All citations will be automatically added to your \texttt{cite.bib} file.

\section{Conclusion}
This template demonstrates basic \LaTeX{} features including mathematical formulas, tables, and lists. The bibliography system is now set up to work seamlessly with Bibby AI's citation management. Feel free to modify it for your needs!

\bibliographystyle{plain}
\bibliography{cite}

\end{document}\documentclass{article}
\usepackage[utf8]{inputenc}
\usepackage{booktabs}
\usepackage{amsmath}
\usepackage{natbib}

\title{My Document}
\author{Your Name}
\date{\today}

\begin{document}

\maketitle

\section{Introduction}
This is a sample document demonstrating various \LaTeX{} features including mathematical equations, tables, and formatting. You can cite papers using the citation search feature, and they will be automatically added to your \texttt{cite.bib} file.

\section{Mathematical Expressions}
Here are some example equations:

\subsection{Quadratic Formula}
The quadratic formula is given by:
\[x = \frac{-b \pm \sqrt{b^2 - 4ac}}{2a}\]

\subsection{Matrix Example}
A sample matrix:
\[
\begin{pmatrix}
a_{11} & a_{12} & a_{13} \\
a_{21} & a_{22} & a_{23} \\
a_{31} & a_{32} & a_{33}
\end{pmatrix}
\]

\section{Tables}
Below is a sample table showing data:

\begin{table}[h]
\centering
\begin{tabular}{@{}lrrr@{}}
\toprule
\textbf{Category} & \textbf{Value 1} & \textbf{Value 2} & \textbf{Total} \\
\midrule
Item A & 10.5 & 20.3 & 30.8 \\
Item B & 15.2 & 25.1 & 40.3 \\
Item C & 22.1 & 12.7 & 34.8 \\
\bottomrule
\end{tabular}
\caption{Sample Data Table}
\label{tab:sample}
\end{table}

\section{Lists}
\subsection{Bullet Points}
\begin{itemize}
  \item First item
  \item Second item
  \item Third item with sub-items:
    \begin{itemize}
      \item Sub-item 1
      \item Sub-item 2
    \end{itemize}
\end{itemize}

\subsection{Numbered List}
\begin{enumerate}
  \item First step
  \item Second step
  \item Third step
\end{enumerate}

\section{Citations and References}
When you use the citation search feature, references will be automatically managed. For example, you might cite a paper like this: \cite{example2024}. All citations will be automatically added to your \texttt{cite.bib} file.

\section{Conclusion}
This template demonstrates basic \LaTeX{} features including mathematical formulas, tables, and lists. The bibliography system is now set up to work seamlessly with Bibby AI's citation management. Feel free to modify it for your needs!

\bibliographystyle{plain}
\bibliography{cite}

\end{document}\documentclass{article}
\usepackage[utf8]{inputenc}
\usepackage{booktabs}
\usepackage{amsmath}
\usepackage{natbib}

\title{My Document}
\author{Your Name}
\date{\today}

\begin{document}

\maketitle

\section{Introduction}
This is a sample document demonstrating various \LaTeX{} features including mathematical equations, tables, and formatting. You can cite papers using the citation search feature, and they will be automatically added to your \texttt{cite.bib} file.

\section{Mathematical Expressions}
Here are some example equations:

\subsection{Quadratic Formula}
The quadratic formula is given by:
\[x = \frac{-b \pm \sqrt{b^2 - 4ac}}{2a}\]

\subsection{Matrix Example}
A sample matrix:
\[
\begin{pmatrix}
a_{11} & a_{12} & a_{13} \\
a_{21} & a_{22} & a_{23} \\
a_{31} & a_{32} & a_{33}
\end{pmatrix}
\]

\section{Tables}
Below is a sample table showing data:

\begin{table}[h]
\centering
\begin{tabular}{@{}lrrr@{}}
\toprule
\textbf{Category} & \textbf{Value 1} & \textbf{Value 2} & \textbf{Total} \\
\midrule
Item A & 10.5 & 20.3 & 30.8 \\
Item B & 15.2 & 25.1 & 40.3 \\
Item C & 22.1 & 12.7 & 34.8 \\
\bottomrule
\end{tabular}
\caption{Sample Data Table}
\label{tab:sample}
\end{table}

\section{Lists}
\subsection{Bullet Points}
\begin{itemize}
  \item First item
  \item Second item
  \item Third item with sub-items:
    \begin{itemize}
      \item Sub-item 1
      \item Sub-item 2
    \end{itemize}
\end{itemize}

\subsection{Numbered List}
\begin{enumerate}
  \item First step
  \item Second step
  \item Third step
\end{enumerate}

\section{Citations and References}
When you use the citation search feature, references will be automatically managed. For example, you might cite a paper like this: \cite{example2024}. All citations will be automatically added to your \texttt{cite.bib} file.

\section{Conclusion}
This template demonstrates basic \LaTeX{} features including mathematical formulas, tables, and lists. The bibliography system is now set up to work seamlessly with Bibby AI's citation management. Feel free to modify it for your needs!

\bibliographystyle{plain}
\bibliography{cite}

\end{document}\documentclass{article}
\usepackage[utf8]{inputenc}
\usepackage{booktabs}
\usepackage{amsmath}
\usepackage{natbib}

\title{My Document}
\author{Your Name}
\date{\today}

\begin{document}

\maketitle

\section{Introduction}
This is a sample document demonstrating various \LaTeX{} features including mathematical equations, tables, and formatting. You can cite papers using the citation search feature, and they will be automatically added to your \texttt{cite.bib} file.

\section{Mathematical Expressions}
Here are some example equations:

\subsection{Quadratic Formula}
The quadratic formula is given by:
\[x = \frac{-b \pm \sqrt{b^2 - 4ac}}{2a}\]

\subsection{Matrix Example}
A sample matrix:
\[
\begin{pmatrix}
a_{11} & a_{12} & a_{13} \\
a_{21} & a_{22} & a_{23} \\
a_{31} & a_{32} & a_{33}
\end{pmatrix}
\]

\section{Tables}
Below is a sample table showing data:

\begin{table}[h]
\centering
\begin{tabular}{@{}lrrr@{}}
\toprule
\textbf{Category} & \textbf{Value 1} & \textbf{Value 2} & \textbf{Total} \\
\midrule
Item A & 10.5 & 20.3 & 30.8 \\
Item B & 15.2 & 25.1 & 40.3 \\
Item C & 22.1 & 12.7 & 34.8 \\
\bottomrule
\end{tabular}
\caption{Sample Data Table}
\label{tab:sample}
\end{table}

\section{Lists}
\subsection{Bullet Points}
\begin{itemize}
  \item First item
  \item Second item
  \item Third item with sub-items:
    \begin{itemize}
      \item Sub-item 1
      \item Sub-item 2
    \end{itemize}
\end{itemize}

\subsection{Numbered List}
\begin{enumerate}
  \item First step
  \item Second step
  \item Third step
\end{enumerate}

\section{Citations and References}
When you use the citation search feature, references will be automatically managed. For example, you might cite a paper like this: \cite{example2024}. All citations will be automatically added to your \texttt{cite.bib} file.

\section{Conclusion}
This template demonstrates basic \LaTeX{} features including mathematical formulas, tables, and lists. The bibliography system is now set up to work seamlessly with Bibby AI's citation management. Feel free to modify it for your needs!

\bibliographystyle{plain}
\bibliography{cite}

\end{document}\documentclass{article}
\usepackage[utf8]{inputenc}
\usepackage{booktabs}
\usepackage{amsmath}
\usepackage{natbib}

\title{My Document}
\author{Your Name}
\date{\today}

\begin{document}

\maketitle

\section{Introduction}
This is a sample document demonstrating various \LaTeX{} features including mathematical equations, tables, and formatting. You can cite papers using the citation search feature, and they will be automatically added to your \texttt{cite.bib} file.

\section{Mathematical Expressions}
Here are some example equations:

\subsection{Quadratic Formula}
The quadratic formula is given by:
\[x = \frac{-b \pm \sqrt{b^2 - 4ac}}{2a}\]

\subsection{Matrix Example}
A sample matrix:
\[
\begin{pmatrix}
a_{11} & a_{12} & a_{13} \\
a_{21} & a_{22} & a_{23} \\
a_{31} & a_{32} & a_{33}
\end{pmatrix}
\]

\section{Tables}
Below is a sample table showing data:

\begin{table}[h]
\centering
\begin{tabular}{@{}lrrr@{}}
\toprule
\textbf{Category} & \textbf{Value 1} & \textbf{Value 2} & \textbf{Total} \\
\midrule
Item A & 10.5 & 20.3 & 30.8 \\
Item B & 15.2 & 25.1 & 40.3 \\
Item C & 22.1 & 12.7 & 34.8 \\
\bottomrule
\end{tabular}
\caption{Sample Data Table}
\label{tab:sample}
\end{table}

\section{Lists}
\subsection{Bullet Points}
\begin{itemize}
  \item First item
  \item Second item
  \item Third item with sub-items:
    \begin{itemize}
      \item Sub-item 1
      \item Sub-item 2
    \end{itemize}
\end{itemize}

\subsection{Numbered List}
\begin{enumerate}
  \item First step
  \item Second step
  \item Third step
\end{enumerate}

\section{Citations and References}
When you use the citation search feature, references will be automatically managed. For example, you might cite a paper like this: \cite{example2024}. All citations will be automatically added to your \texttt{cite.bib} file.

\section{Conclusion}
This template demonstrates basic \LaTeX{} features including mathematical formulas, tables, and lists. The bibliography system is now set up to work seamlessly with Bibby AI's citation management. Feel free to modify it for your needs!

\bibliographystyle{plain}
\bibliography{cite}

\end{document}\documentclass{article}
\usepackage[utf8]{inputenc}
\usepackage{booktabs}
\usepackage{amsmath}
\usepackage{natbib}

\title{My Document}
\author{Your Name}
\date{\today}

\begin{document}

\maketitle

\section{Introduction}
This is a sample document demonstrating various \LaTeX{} features including mathematical equations, tables, and formatting. You can cite papers using the citation search feature, and they will be automatically added to your \texttt{cite.bib} file.

\section{Mathematical Expressions}
Here are some example equations:

\subsection{Quadratic Formula}
The quadratic formula is given by:
\[x = \frac{-b \pm \sqrt{b^2 - 4ac}}{2a}\]

\subsection{Matrix Example}
A sample matrix:
\[
\begin{pmatrix}
a_{11} & a_{12} & a_{13} \\
a_{21} & a_{22} & a_{23} \\
a_{31} & a_{32} & a_{33}
\end{pmatrix}
\]

\section{Tables}
Below is a sample table showing data:

\begin{table}[h]
\centering
\begin{tabular}{@{}lrrr@{}}
\toprule
\textbf{Category} & \textbf{Value 1} & \textbf{Value 2} & \textbf{Total} \\
\midrule
Item A & 10.5 & 20.3 & 30.8 \\
Item B & 15.2 & 25.1 & 40.3 \\
Item C & 22.1 & 12.7 & 34.8 \\
\bottomrule
\end{tabular}
\caption{Sample Data Table}
\label{tab:sample}
\end{table}

\section{Lists}
\subsection{Bullet Points}
\begin{itemize}
  \item First item
  \item Second item
  \item Third item with sub-items:
    \begin{itemize}
      \item Sub-item 1
      \item Sub-item 2
    \end{itemize}
\end{itemize}

\subsection{Numbered List}
\begin{enumerate}
  \item First step
  \item Second step
  \item Third step
\end{enumerate}

\section{Citations and References}
When you use the citation search feature, references will be automatically managed. For example, you might cite a paper like this: \cite{example2024}. All citations will be automatically added to your \texttt{cite.bib} file.

\section{Conclusion}
This template demonstrates basic \LaTeX{} features including mathematical formulas, tables, and lists. The bibliography system is now set up to work seamlessly with Bibby AI's citation management. Feel free to modify it for your needs!

\bibliographystyle{plain}
\bibliography{cite}

\end{document}\documentclass{article}
\usepackage[utf8]{inputenc}
\usepackage{booktabs}
\usepackage{amsmath}
\usepackage{natbib}

\title{My Document}
\author{Your Name}
\date{\today}

\begin{document}

\maketitle

\section{Introduction}
This is a sample document demonstrating various \LaTeX{} features including mathematical equations, tables, and formatting. You can cite papers using the citation search feature, and they will be automatically added to your \texttt{cite.bib} file.

\section{Mathematical Expressions}
Here are some example equations:

\subsection{Quadratic Formula}
The quadratic formula is given by:
\[x = \frac{-b \pm \sqrt{b^2 - 4ac}}{2a}\]

\subsection{Matrix Example}
A sample matrix:
\[
\begin{pmatrix}
a_{11} & a_{12} & a_{13} \\
a_{21} & a_{22} & a_{23} \\
a_{31} & a_{32} & a_{33}
\end{pmatrix}
\]

\section{Tables}
Below is a sample table showing data:

\begin{table}[h]
\centering
\begin{tabular}{@{}lrrr@{}}
\toprule
\textbf{Category} & \textbf{Value 1} & \textbf{Value 2} & \textbf{Total} \\
\midrule
Item A & 10.5 & 20.3 & 30.8 \\
Item B & 15.2 & 25.1 & 40.3 \\
Item C & 22.1 & 12.7 & 34.8 \\
\bottomrule
\end{tabular}
\caption{Sample Data Table}
\label{tab:sample}
\end{table}

\section{Lists}
\subsection{Bullet Points}
\begin{itemize}
  \item First item
  \item Second item
  \item Third item with sub-items:
    \begin{itemize}
      \item Sub-item 1
      \item Sub-item 2
    \end{itemize}
\end{itemize}

\subsection{Numbered List}
\begin{enumerate}
  \item First step
  \item Second step
  \item Third step
\end{enumerate}

\section{Citations and References}
When you use the citation search feature, references will be automatically managed. For example, you might cite a paper like this: \cite{example2024}. All citations will be automatically added to your \texttt{cite.bib} file.

\section{Conclusion}
This template demonstrates basic \LaTeX{} features including mathematical formulas, tables, and lists. The bibliography system is now set up to work seamlessly with Bibby AI's citation management. Feel free to modify it for your needs!

\bibliographystyle{plain}
\bibliography{cite}

\end{document}\documentclass{article}
\usepackage[utf8]{inputenc}
\usepackage{booktabs}
\usepackage{amsmath}
\usepackage{natbib}

\title{My Document}
\author{Your Name}
\date{\today}

\begin{document}

\maketitle

\section{Introduction}
This is a sample document demonstrating various \LaTeX{} features including mathematical equations, tables, and formatting. You can cite papers using the citation search feature, and they will be automatically added to your \texttt{cite.bib} file.

\section{Mathematical Expressions}
Here are some example equations:

\subsection{Quadratic Formula}
The quadratic formula is given by:
\[x = \frac{-b \pm \sqrt{b^2 - 4ac}}{2a}\]

\subsection{Matrix Example}
A sample matrix:
\[
\begin{pmatrix}
a_{11} & a_{12} & a_{13} \\
a_{21} & a_{22} & a_{23} \\
a_{31} & a_{32} & a_{33}
\end{pmatrix}
\]

\section{Tables}
Below is a sample table showing data:

\begin{table}[h]
\centering
\begin{tabular}{@{}lrrr@{}}
\toprule
\textbf{Category} & \textbf{Value 1} & \textbf{Value 2} & \textbf{Total} \\
\midrule
Item A & 10.5 & 20.3 & 30.8 \\
Item B & 15.2 & 25.1 & 40.3 \\
Item C & 22.1 & 12.7 & 34.8 \\
\bottomrule
\end{tabular}
\caption{Sample Data Table}
\label{tab:sample}
\end{table}

\section{Lists}
\subsection{Bullet Points}
\begin{itemize}
  \item First item
  \item Second item
  \item Third item with sub-items:
    \begin{itemize}
      \item Sub-item 1
      \item Sub-item 2
    \end{itemize}
\end{itemize}

\subsection{Numbered List}
\begin{enumerate}
  \item First step
  \item Second step
  \item Third step
\end{enumerate}

\section{Citations and References}
When you use the citation search feature, references will be automatically managed. For example, you might cite a paper like this: \cite{example2024}. All citations will be automatically added to your \texttt{cite.bib} file.

\section{Conclusion}
This template demonstrates basic \LaTeX{} features including mathematical formulas, tables, and lists. The bibliography system is now set up to work seamlessly with Bibby AI's citation management. Feel free to modify it for your needs!

\bibliographystyle{plain}
\bibliography{cite}

\end{document}\documentclass{article}
\usepackage[utf8]{inputenc}
\usepackage{booktabs}
\usepackage{amsmath}
\usepackage{natbib}

\title{My Document}
\author{Your Name}
\date{\today}

\begin{document}

\maketitle

\section{Introduction}
This is a sample document demonstrating various \LaTeX{} features including mathematical equations, tables, and formatting. You can cite papers using the citation search feature, and they will be automatically added to your \texttt{cite.bib} file.

\section{Mathematical Expressions}
Here are some example equations:

\subsection{Quadratic Formula}
The quadratic formula is given by:
\[x = \frac{-b \pm \sqrt{b^2 - 4ac}}{2a}\]

\subsection{Matrix Example}
A sample matrix:
\[
\begin{pmatrix}
a_{11} & a_{12} & a_{13} \\
a_{21} & a_{22} & a_{23} \\
a_{31} & a_{32} & a_{33}
\end{pmatrix}
\]

\section{Tables}
Below is a sample table showing data:

\begin{table}[h]
\centering
\begin{tabular}{@{}lrrr@{}}
\toprule
\textbf{Category} & \textbf{Value 1} & \textbf{Value 2} & \textbf{Total} \\
\midrule
Item A & 10.5 & 20.3 & 30.8 \\
Item B & 15.2 & 25.1 & 40.3 \\
Item C & 22.1 & 12.7 & 34.8 \\
\bottomrule
\end{tabular}
\caption{Sample Data Table}
\label{tab:sample}
\end{table}

\section{Lists}
\subsection{Bullet Points}
\begin{itemize}
  \item First item
  \item Second item
  \item Third item with sub-items:
    \begin{itemize}
      \item Sub-item 1
      \item Sub-item 2
    \end{itemize}
\end{itemize}

\subsection{Numbered List}
\begin{enumerate}
  \item First step
  \item Second step
  \item Third step
\end{enumerate}

\section{Citations and References}
When you use the citation search feature, references will be automatically managed. For example, you might cite a paper like this: \cite{example2024}. All citations will be automatically added to your \texttt{cite.bib} file.

\section{Conclusion}
This template demonstrates basic \LaTeX{} features including mathematical formulas, tables, and lists. The bibliography system is now set up to work seamlessly with Bibby AI's citation management. Feel free to modify it for your needs!

\bibliographystyle{plain}
\bibliography{cite}

\end{document}\documentclass{article}
\usepackage[utf8]{inputenc}
\usepackage{booktabs}
\usepackage{amsmath}
\usepackage{natbib}

\title{My Document}
\author{Your Name}
\date{\today}

\begin{document}

\maketitle

\section{Introduction}
This is a sample document demonstrating various \LaTeX{} features including mathematical equations, tables, and formatting. You can cite papers using the citation search feature, and they will be automatically added to your \texttt{cite.bib} file.

\section{Mathematical Expressions}
Here are some example equations:

\subsection{Quadratic Formula}
The quadratic formula is given by:
\[x = \frac{-b \pm \sqrt{b^2 - 4ac}}{2a}\]

\subsection{Matrix Example}
A sample matrix:
\[
\begin{pmatrix}
a_{11} & a_{12} & a_{13} \\
a_{21} & a_{22} & a_{23} \\
a_{31} & a_{32} & a_{33}
\end{pmatrix}
\]

\section{Tables}
Below is a sample table showing data:

\begin{table}[h]
\centering
\begin{tabular}{@{}lrrr@{}}
\toprule
\textbf{Category} & \textbf{Value 1} & \textbf{Value 2} & \textbf{Total} \\
\midrule
Item A & 10.5 & 20.3 & 30.8 \\
Item B & 15.2 & 25.1 & 40.3 \\
Item C & 22.1 & 12.7 & 34.8 \\
\bottomrule
\end{tabular}
\caption{Sample Data Table}
\label{tab:sample}
\end{table}

\section{Lists}
\subsection{Bullet Points}
\begin{itemize}
  \item First item
  \item Second item
  \item Third item with sub-items:
    \begin{itemize}
      \item Sub-item 1
      \item Sub-item 2
    \end{itemize}
\end{itemize}

\subsection{Numbered List}
\begin{enumerate}
  \item First step
  \item Second step
  \item Third step
\end{enumerate}

\section{Citations and References}
When you use the citation search feature, references will be automatically managed. For example, you might cite a paper like this: \cite{example2024}. All citations will be automatically added to your \texttt{cite.bib} file.

\section{Conclusion}
This template demonstrates basic \LaTeX{} features including mathematical formulas, tables, and lists. The bibliography system is now set up to work seamlessly with Bibby AI's citation management. Feel free to modify it for your needs!

\bibliographystyle{plain}
\bibliography{cite}

\end{document}\documentclass{article}
\usepackage[utf8]{inputenc}
\usepackage{booktabs}
\usepackage{amsmath}
\usepackage{natbib}

\title{My Document}
\author{Your Name}
\date{\today}

\begin{document}

\maketitle

\section{Introduction}
This is a sample document demonstrating various \LaTeX{} features including mathematical equations, tables, and formatting. You can cite papers using the citation search feature, and they will be automatically added to your \texttt{cite.bib} file.

\section{Mathematical Expressions}
Here are some example equations:

\subsection{Quadratic Formula}
The quadratic formula is given by:
\[x = \frac{-b \pm \sqrt{b^2 - 4ac}}{2a}\]

\subsection{Matrix Example}
A sample matrix:
\[
\begin{pmatrix}
a_{11} & a_{12} & a_{13} \\
a_{21} & a_{22} & a_{23} \\
a_{31} & a_{32} & a_{33}
\end{pmatrix}
\]

\section{Tables}
Below is a sample table showing data:

\begin{table}[h]
\centering
\begin{tabular}{@{}lrrr@{}}
\toprule
\textbf{Category} & \textbf{Value 1} & \textbf{Value 2} & \textbf{Total} \\
\midrule
Item A & 10.5 & 20.3 & 30.8 \\
Item B & 15.2 & 25.1 & 40.3 \\
Item C & 22.1 & 12.7 & 34.8 \\
\bottomrule
\end{tabular}
\caption{Sample Data Table}
\label{tab:sample}
\end{table}

\section{Lists}
\subsection{Bullet Points}
\begin{itemize}
  \item First item
  \item Second item
  \item Third item with sub-items:
    \begin{itemize}
      \item Sub-item 1
      \item Sub-item 2
    \end{itemize}
\end{itemize}

\subsection{Numbered List}
\begin{enumerate}
  \item First step
  \item Second step
  \item Third step
\end{enumerate}

\section{Citations and References}
When you use the citation search feature, references will be automatically managed. For example, you might cite a paper like this: \cite{example2024}. All citations will be automatically added to your \texttt{cite.bib} file.

\section{Conclusion}
This template demonstrates basic \LaTeX{} features including mathematical formulas, tables, and lists. The bibliography system is now set up to work seamlessly with Bibby AI's citation management. Feel free to modify it for your needs!

\bibliographystyle{plain}
\bibliography{cite}

\end{document}\documentclass{article}
\usepackage[utf8]{inputenc}
\usepackage{booktabs}
\usepackage{amsmath}
\usepackage{natbib}

\title{My Document}
\author{Your Name}
\date{\today}

\begin{document}

\maketitle

\section{Introduction}
This is a sample document demonstrating various \LaTeX{} features including mathematical equations, tables, and formatting. You can cite papers using the citation search feature, and they will be automatically added to your \texttt{cite.bib} file.

\section{Mathematical Expressions}
Here are some example equations:

\subsection{Quadratic Formula}
The quadratic formula is given by:
\[x = \frac{-b \pm \sqrt{b^2 - 4ac}}{2a}\]

\subsection{Matrix Example}
A sample matrix:
\[
\begin{pmatrix}
a_{11} & a_{12} & a_{13} \\
a_{21} & a_{22} & a_{23} \\
a_{31} & a_{32} & a_{33}
\end{pmatrix}
\]

\section{Tables}
Below is a sample table showing data:

\begin{table}[h]
\centering
\begin{tabular}{@{}lrrr@{}}
\toprule
\textbf{Category} & \textbf{Value 1} & \textbf{Value 2} & \textbf{Total} \\
\midrule
Item A & 10.5 & 20.3 & 30.8 \\
Item B & 15.2 & 25.1 & 40.3 \\
Item C & 22.1 & 12.7 & 34.8 \\
\bottomrule
\end{tabular}
\caption{Sample Data Table}
\label{tab:sample}
\end{table}

\section{Lists}
\subsection{Bullet Points}
\begin{itemize}
  \item First item
  \item Second item
  \item Third item with sub-items:
    \begin{itemize}
      \item Sub-item 1
      \item Sub-item 2
    \end{itemize}
\end{itemize}

\subsection{Numbered List}
\begin{enumerate}
  \item First step
  \item Second step
  \item Third step
\end{enumerate}

\section{Citations and References}
When you use the citation search feature, references will be automatically managed. For example, you might cite a paper like this: \cite{example2024}. All citations will be automatically added to your \texttt{cite.bib} file.

\section{Conclusion}
This template demonstrates basic \LaTeX{} features including mathematical formulas, tables, and lists. The bibliography system is now set up to work seamlessly with Bibby AI's citation management. Feel free to modify it for your needs!

\bibliographystyle{plain}
\bibliography{cite}

\end{document}\documentclass{article}
\usepackage[utf8]{inputenc}
\usepackage{booktabs}
\usepackage{amsmath}
\usepackage{natbib}

\title{My Document}
\author{Your Name}
\date{\today}

\begin{document}

\maketitle

\section{Introduction}
This is a sample document demonstrating various \LaTeX{} features including mathematical equations, tables, and formatting. You can cite papers using the citation search feature, and they will be automatically added to your \texttt{cite.bib} file.

\section{Mathematical Expressions}
Here are some example equations:

\subsection{Quadratic Formula}
The quadratic formula is given by:
\[x = \frac{-b \pm \sqrt{b^2 - 4ac}}{2a}\]

\subsection{Matrix Example}
A sample matrix:
\[
\begin{pmatrix}
a_{11} & a_{12} & a_{13} \\
a_{21} & a_{22} & a_{23} \\
a_{31} & a_{32} & a_{33}
\end{pmatrix}
\]

\section{Tables}
Below is a sample table showing data:

\begin{table}[h]
\centering
\begin{tabular}{@{}lrrr@{}}
\toprule
\textbf{Category} & \textbf{Value 1} & \textbf{Value 2} & \textbf{Total} \\
\midrule
Item A & 10.5 & 20.3 & 30.8 \\
Item B & 15.2 & 25.1 & 40.3 \\
Item C & 22.1 & 12.7 & 34.8 \\
\bottomrule
\end{tabular}
\caption{Sample Data Table}
\label{tab:sample}
\end{table}

\section{Lists}
\subsection{Bullet Points}
\begin{itemize}
  \item First item
  \item Second item
  \item Third item with sub-items:
    \begin{itemize}
      \item Sub-item 1
      \item Sub-item 2
    \end{itemize}
\end{itemize}

\subsection{Numbered List}
\begin{enumerate}
  \item First step
  \item Second step
  \item Third step
\end{enumerate}

\section{Citations and References}
When you use the citation search feature, references will be automatically managed. For example, you might cite a paper like this: \cite{example2024}. All citations will be automatically added to your \texttt{cite.bib} file.

\section{Conclusion}
This template demonstrates basic \LaTeX{} features including mathematical formulas, tables, and lists. The bibliography system is now set up to work seamlessly with Bibby AI's citation management. Feel free to modify it for your needs!

\bibliographystyle{plain}
\bibliography{cite}

\end{document}\documentclass{article}
\usepackage[utf8]{inputenc}
\usepackage{booktabs}
\usepackage{amsmath}
\usepackage{natbib}

\title{My Document}
\author{Your Name}
\date{\today}

\begin{document}

\maketitle

\section{Introduction}
This is a sample document demonstrating various \LaTeX{} features including mathematical equations, tables, and formatting. You can cite papers using the citation search feature, and they will be automatically added to your \texttt{cite.bib} file.

\section{Mathematical Expressions}
Here are some example equations:

\subsection{Quadratic Formula}
The quadratic formula is given by:
\[x = \frac{-b \pm \sqrt{b^2 - 4ac}}{2a}\]

\subsection{Matrix Example}
A sample matrix:
\[
\begin{pmatrix}
a_{11} & a_{12} & a_{13} \\
a_{21} & a_{22} & a_{23} \\
a_{31} & a_{32} & a_{33}
\end{pmatrix}
\]

\section{Tables}
Below is a sample table showing data:

\begin{table}[h]
\centering
\begin{tabular}{@{}lrrr@{}}
\toprule
\textbf{Category} & \textbf{Value 1} & \textbf{Value 2} & \textbf{Total} \\
\midrule
Item A & 10.5 & 20.3 & 30.8 \\
Item B & 15.2 & 25.1 & 40.3 \\
Item C & 22.1 & 12.7 & 34.8 \\
\bottomrule
\end{tabular}
\caption{Sample Data Table}
\label{tab:sample}
\end{table}

\section{Lists}
\subsection{Bullet Points}
\begin{itemize}
  \item First item
  \item Second item
  \item Third item with sub-items:
    \begin{itemize}
      \item Sub-item 1
      \item Sub-item 2
    \end{itemize}
\end{itemize}

\subsection{Numbered List}
\begin{enumerate}
  \item First step
  \item Second step
  \item Third step
\end{enumerate}

\section{Citations and References}
When you use the citation search feature, references will be automatically managed. For example, you might cite a paper like this: \cite{example2024}. All citations will be automatically added to your \texttt{cite.bib} file.

\section{Conclusion}
This template demonstrates basic \LaTeX{} features including mathematical formulas, tables, and lists. The bibliography system is now set up to work seamlessly with Bibby AI's citation management. Feel free to modify it for your needs!

\bibliographystyle{plain}
\bibliography{cite}

\end{document}\documentclass{article}
\usepackage[utf8]{inputenc}
\usepackage{booktabs}
\usepackage{amsmath}
\usepackage{natbib}

\title{My Document}
\author{Your Name}
\date{\today}

\begin{document}

\maketitle

\section{Introduction}
This is a sample document demonstrating various \LaTeX{} features including mathematical equations, tables, and formatting. You can cite papers using the citation search feature, and they will be automatically added to your \texttt{cite.bib} file.

\section{Mathematical Expressions}
Here are some example equations:

\subsection{Quadratic Formula}
The quadratic formula is given by:
\[x = \frac{-b \pm \sqrt{b^2 - 4ac}}{2a}\]

\subsection{Matrix Example}
A sample matrix:
\[
\begin{pmatrix}
a_{11} & a_{12} & a_{13} \\
a_{21} & a_{22} & a_{23} \\
a_{31} & a_{32} & a_{33}
\end{pmatrix}
\]

\section{Tables}
Below is a sample table showing data:

\begin{table}[h]
\centering
\begin{tabular}{@{}lrrr@{}}
\toprule
\textbf{Category} & \textbf{Value 1} & \textbf{Value 2} & \textbf{Total} \\
\midrule
Item A & 10.5 & 20.3 & 30.8 \\
Item B & 15.2 & 25.1 & 40.3 \\
Item C & 22.1 & 12.7 & 34.8 \\
\bottomrule
\end{tabular}
\caption{Sample Data Table}
\label{tab:sample}
\end{table}

\section{Lists}
\subsection{Bullet Points}
\begin{itemize}
  \item First item
  \item Second item
  \item Third item with sub-items:
    \begin{itemize}
      \item Sub-item 1
      \item Sub-item 2
    \end{itemize}
\end{itemize}

\subsection{Numbered List}
\begin{enumerate}
  \item First step
  \item Second step
  \item Third step
\end{enumerate}

\section{Citations and References}
When you use the citation search feature, references will be automatically managed. For example, you might cite a paper like this: \cite{example2024}. All citations will be automatically added to your \texttt{cite.bib} file.

\section{Conclusion}
This template demonstrates basic \LaTeX{} features including mathematical formulas, tables, and lists. The bibliography system is now set up to work seamlessly with Bibby AI's citation management. Feel free to modify it for your needs!

\bibliographystyle{plain}
\bibliography{cite}

\end{document}\documentclass{article}
\usepackage[utf8]{inputenc}
\usepackage{booktabs}
\usepackage{amsmath}
\usepackage{natbib}

\title{My Document}
\author{Your Name}
\date{\today}

\begin{document}

\maketitle

\section{Introduction}
This is a sample document demonstrating various \LaTeX{} features including mathematical equations, tables, and formatting. You can cite papers using the citation search feature, and they will be automatically added to your \texttt{cite.bib} file.

\section{Mathematical Expressions}
Here are some example equations:

\subsection{Quadratic Formula}
The quadratic formula is given by:
\[x = \frac{-b \pm \sqrt{b^2 - 4ac}}{2a}\]

\subsection{Matrix Example}
A sample matrix:
\[
\begin{pmatrix}
a_{11} & a_{12} & a_{13} \\
a_{21} & a_{22} & a_{23} \\
a_{31} & a_{32} & a_{33}
\end{pmatrix}
\]

\section{Tables}
Below is a sample table showing data:

\begin{table}[h]
\centering
\begin{tabular}{@{}lrrr@{}}
\toprule
\textbf{Category} & \textbf{Value 1} & \textbf{Value 2} & \textbf{Total} \\
\midrule
Item A & 10.5 & 20.3 & 30.8 \\
Item B & 15.2 & 25.1 & 40.3 \\
Item C & 22.1 & 12.7 & 34.8 \\
\bottomrule
\end{tabular}
\caption{Sample Data Table}
\label{tab:sample}
\end{table}

\section{Lists}
\subsection{Bullet Points}
\begin{itemize}
  \item First item
  \item Second item
  \item Third item with sub-items:
    \begin{itemize}
      \item Sub-item 1
      \item Sub-item 2
    \end{itemize}
\end{itemize}

\subsection{Numbered List}
\begin{enumerate}
  \item First step
  \item Second step
  \item Third step
\end{enumerate}

\section{Citations and References}
When you use the citation search feature, references will be automatically managed. For example, you might cite a paper like this: \cite{example2024}. All citations will be automatically added to your \texttt{cite.bib} file.

\section{Conclusion}
This template demonstrates basic \LaTeX{} features including mathematical formulas, tables, and lists. The bibliography system is now set up to work seamlessly with Bibby AI's citation management. Feel free to modify it for your needs!

\bibliographystyle{plain}
\bibliography{cite}

\end{document}\documentclass{article}
\usepackage[utf8]{inputenc}
\usepackage{booktabs}
\usepackage{amsmath}
\usepackage{natbib}

\title{My Document}
\author{Your Name}
\date{\today}

\begin{document}

\maketitle

\section{Introduction}
This is a sample document demonstrating various \LaTeX{} features including mathematical equations, tables, and formatting. You can cite papers using the citation search feature, and they will be automatically added to your \texttt{cite.bib} file.

\section{Mathematical Expressions}
Here are some example equations:

\subsection{Quadratic Formula}
The quadratic formula is given by:
\[x = \frac{-b \pm \sqrt{b^2 - 4ac}}{2a}\]

\subsection{Matrix Example}
A sample matrix:
\[
\begin{pmatrix}
a_{11} & a_{12} & a_{13} \\
a_{21} & a_{22} & a_{23} \\
a_{31} & a_{32} & a_{33}
\end{pmatrix}
\]

\section{Tables}
Below is a sample table showing data:

\begin{table}[h]
\centering
\begin{tabular}{@{}lrrr@{}}
\toprule
\textbf{Category} & \textbf{Value 1} & \textbf{Value 2} & \textbf{Total} \\
\midrule
Item A & 10.5 & 20.3 & 30.8 \\
Item B & 15.2 & 25.1 & 40.3 \\
Item C & 22.1 & 12.7 & 34.8 \\
\bottomrule
\end{tabular}
\caption{Sample Data Table}
\label{tab:sample}
\end{table}

\section{Lists}
\subsection{Bullet Points}
\begin{itemize}
  \item First item
  \item Second item
  \item Third item with sub-items:
    \begin{itemize}
      \item Sub-item 1
      \item Sub-item 2
    \end{itemize}
\end{itemize}

\subsection{Numbered List}
\begin{enumerate}
  \item First step
  \item Second step
  \item Third step
\end{enumerate}

\section{Citations and References}
When you use the citation search feature, references will be automatically managed. For example, you might cite a paper like this: \cite{example2024}. All citations will be automatically added to your \texttt{cite.bib} file.

\section{Conclusion}
This template demonstrates basic \LaTeX{} features including mathematical formulas, tables, and lists. The bibliography system is now set up to work seamlessly with Bibby AI's citation management. Feel free to modify it for your needs!

\bibliographystyle{plain}
\bibliography{cite}

\end{document}\documentclass{article}
\usepackage[utf8]{inputenc}
\usepackage{booktabs}
\usepackage{amsmath}
\usepackage{natbib}

\title{My Document}
\author{Your Name}
\date{\today}

\begin{document}

\maketitle

\section{Introduction}
This is a sample document demonstrating various \LaTeX{} features including mathematical equations, tables, and formatting. You can cite papers using the citation search feature, and they will be automatically added to your \texttt{cite.bib} file.

\section{Mathematical Expressions}
Here are some example equations:

\subsection{Quadratic Formula}
The quadratic formula is given by:
\[x = \frac{-b \pm \sqrt{b^2 - 4ac}}{2a}\]

\subsection{Matrix Example}
A sample matrix:
\[
\begin{pmatrix}
a_{11} & a_{12} & a_{13} \\
a_{21} & a_{22} & a_{23} \\
a_{31} & a_{32} & a_{33}
\end{pmatrix}
\]

\section{Tables}
Below is a sample table showing data:

\begin{table}[h]
\centering
\begin{tabular}{@{}lrrr@{}}
\toprule
\textbf{Category} & \textbf{Value 1} & \textbf{Value 2} & \textbf{Total} \\
\midrule
Item A & 10.5 & 20.3 & 30.8 \\
Item B & 15.2 & 25.1 & 40.3 \\
Item C & 22.1 & 12.7 & 34.8 \\
\bottomrule
\end{tabular}
\caption{Sample Data Table}
\label{tab:sample}
\end{table}

\section{Lists}
\subsection{Bullet Points}
\begin{itemize}
  \item First item
  \item Second item
  \item Third item with sub-items:
    \begin{itemize}
      \item Sub-item 1
      \item Sub-item 2
    \end{itemize}
\end{itemize}

\subsection{Numbered List}
\begin{enumerate}
  \item First step
  \item Second step
  \item Third step
\end{enumerate}

\section{Citations and References}
When you use the citation search feature, references will be automatically managed. For example, you might cite a paper like this: \cite{example2024}. All citations will be automatically added to your \texttt{cite.bib} file.

\section{Conclusion}
This template demonstrates basic \LaTeX{} features including mathematical formulas, tables, and lists. The bibliography system is now set up to work seamlessly with Bibby AI's citation management. Feel free to modify it for your needs!

\bibliographystyle{plain}
\bibliography{cite}

\end{document}\documentclass{article}
\usepackage[utf8]{inputenc}
\usepackage{booktabs}
\usepackage{amsmath}
\usepackage{natbib}

\title{My Document}
\author{Your Name}
\date{\today}

\begin{document}

\maketitle

\section{Introduction}
This is a sample document demonstrating various \LaTeX{} features including mathematical equations, tables, and formatting. You can cite papers using the citation search feature, and they will be automatically added to your \texttt{cite.bib} file.

\section{Mathematical Expressions}
Here are some example equations:

\subsection{Quadratic Formula}
The quadratic formula is given by:
\[x = \frac{-b \pm \sqrt{b^2 - 4ac}}{2a}\]

\subsection{Matrix Example}
A sample matrix:
\[
\begin{pmatrix}
a_{11} & a_{12} & a_{13} \\
a_{21} & a_{22} & a_{23} \\
a_{31} & a_{32} & a_{33}
\end{pmatrix}
\]

\section{Tables}
Below is a sample table showing data:

\begin{table}[h]
\centering
\begin{tabular}{@{}lrrr@{}}
\toprule
\textbf{Category} & \textbf{Value 1} & \textbf{Value 2} & \textbf{Total} \\
\midrule
Item A & 10.5 & 20.3 & 30.8 \\
Item B & 15.2 & 25.1 & 40.3 \\
Item C & 22.1 & 12.7 & 34.8 \\
\bottomrule
\end{tabular}
\caption{Sample Data Table}
\label{tab:sample}
\end{table}

\section{Lists}
\subsection{Bullet Points}
\begin{itemize}
  \item First item
  \item Second item
  \item Third item with sub-items:
    \begin{itemize}
      \item Sub-item 1
      \item Sub-item 2
    \end{itemize}
\end{itemize}

\subsection{Numbered List}
\begin{enumerate}
  \item First step
  \item Second step
  \item Third step
\end{enumerate}

\section{Citations and References}
When you use the citation search feature, references will be automatically managed. For example, you might cite a paper like this: \cite{example2024}. All citations will be automatically added to your \texttt{cite.bib} file.

\section{Conclusion}
This template demonstrates basic \LaTeX{} features including mathematical formulas, tables, and lists. The bibliography system is now set up to work seamlessly with Bibby AI's citation management. Feel free to modify it for your needs!

\bibliographystyle{plain}
\bibliography{cite}

\end{document}\documentclass{article}
\usepackage[utf8]{inputenc}
\usepackage{booktabs}
\usepackage{amsmath}
\usepackage{natbib}

\title{My Document}
\author{Your Name}
\date{\today}

\begin{document}

\maketitle

\section{Introduction}
This is a sample document demonstrating various \LaTeX{} features including mathematical equations, tables, and formatting. You can cite papers using the citation search feature, and they will be automatically added to your \texttt{cite.bib} file.

\section{Mathematical Expressions}
Here are some example equations:

\subsection{Quadratic Formula}
The quadratic formula is given by:
\[x = \frac{-b \pm \sqrt{b^2 - 4ac}}{2a}\]

\subsection{Matrix Example}
A sample matrix:
\[
\begin{pmatrix}
a_{11} & a_{12} & a_{13} \\
a_{21} & a_{22} & a_{23} \\
a_{31} & a_{32} & a_{33}
\end{pmatrix}
\]

\section{Tables}
Below is a sample table showing data:

\begin{table}[h]
\centering
\begin{tabular}{@{}lrrr@{}}
\toprule
\textbf{Category} & \textbf{Value 1} & \textbf{Value 2} & \textbf{Total} \\
\midrule
Item A & 10.5 & 20.3 & 30.8 \\
Item B & 15.2 & 25.1 & 40.3 \\
Item C & 22.1 & 12.7 & 34.8 \\
\bottomrule
\end{tabular}
\caption{Sample Data Table}
\label{tab:sample}
\end{table}

\section{Lists}
\subsection{Bullet Points}
\begin{itemize}
  \item First item
  \item Second item
  \item Third item with sub-items:
    \begin{itemize}
      \item Sub-item 1
      \item Sub-item 2
    \end{itemize}
\end{itemize}

\subsection{Numbered List}
\begin{enumerate}
  \item First step
  \item Second step
  \item Third step
\end{enumerate}

\section{Citations and References}
When you use the citation search feature, references will be automatically managed. For example, you might cite a paper like this: \cite{example2024}. All citations will be automatically added to your \texttt{cite.bib} file.

\section{Conclusion}
This template demonstrates basic \LaTeX{} features including mathematical formulas, tables, and lists. The bibliography system is now set up to work seamlessly with Bibby AI's citation management. Feel free to modify it for your needs!

\bibliographystyle{plain}
\bibliography{cite}

\end{document}\documentclass{article}
\usepackage[utf8]{inputenc}
\usepackage{booktabs}
\usepackage{amsmath}
\usepackage{natbib}

\title{My Document}
\author{Your Name}
\date{\today}

\begin{document}

\maketitle

\section{Introduction}
This is a sample document demonstrating various \LaTeX{} features including mathematical equations, tables, and formatting. You can cite papers using the citation search feature, and they will be automatically added to your \texttt{cite.bib} file.

\section{Mathematical Expressions}
Here are some example equations:

\subsection{Quadratic Formula}
The quadratic formula is given by:
\[x = \frac{-b \pm \sqrt{b^2 - 4ac}}{2a}\]

\subsection{Matrix Example}
A sample matrix:
\[
\begin{pmatrix}
a_{11} & a_{12} & a_{13} \\
a_{21} & a_{22} & a_{23} \\
a_{31} & a_{32} & a_{33}
\end{pmatrix}
\]

\section{Tables}
Below is a sample table showing data:

\begin{table}[h]
\centering
\begin{tabular}{@{}lrrr@{}}
\toprule
\textbf{Category} & \textbf{Value 1} & \textbf{Value 2} & \textbf{Total} \\
\midrule
Item A & 10.5 & 20.3 & 30.8 \\
Item B & 15.2 & 25.1 & 40.3 \\
Item C & 22.1 & 12.7 & 34.8 \\
\bottomrule
\end{tabular}
\caption{Sample Data Table}
\label{tab:sample}
\end{table}

\section{Lists}
\subsection{Bullet Points}
\begin{itemize}
  \item First item
  \item Second item
  \item Third item with sub-items:
    \begin{itemize}
      \item Sub-item 1
      \item Sub-item 2
    \end{itemize}
\end{itemize}

\subsection{Numbered List}
\begin{enumerate}
  \item First step
  \item Second step
  \item Third step
\end{enumerate}

\section{Citations and References}
When you use the citation search feature, references will be automatically managed. For example, you might cite a paper like this: \cite{example2024}. All citations will be automatically added to your \texttt{cite.bib} file.

\section{Conclusion}
This template demonstrates basic \LaTeX{} features including mathematical formulas, tables, and lists. The bibliography system is now set up to work seamlessly with Bibby AI's citation management. Feel free to modify it for your needs!

\bibliographystyle{plain}
\bibliography{cite}

\end{document}\documentclass{article}
\usepackage[utf8]{inputenc}
\usepackage{booktabs}
\usepackage{amsmath}
\usepackage{natbib}

\title{My Document}
\author{Your Name}
\date{\today}

\begin{document}

\maketitle

\section{Introduction}
This is a sample document demonstrating various \LaTeX{} features including mathematical equations, tables, and formatting. You can cite papers using the citation search feature, and they will be automatically added to your \texttt{cite.bib} file.

\section{Mathematical Expressions}
Here are some example equations:

\subsection{Quadratic Formula}
The quadratic formula is given by:
\[x = \frac{-b \pm \sqrt{b^2 - 4ac}}{2a}\]

\subsection{Matrix Example}
A sample matrix:
\[
\begin{pmatrix}
a_{11} & a_{12} & a_{13} \\
a_{21} & a_{22} & a_{23} \\
a_{31} & a_{32} & a_{33}
\end{pmatrix}
\]

\section{Tables}
Below is a sample table showing data:

\begin{table}[h]
\centering
\begin{tabular}{@{}lrrr@{}}
\toprule
\textbf{Category} & \textbf{Value 1} & \textbf{Value 2} & \textbf{Total} \\
\midrule
Item A & 10.5 & 20.3 & 30.8 \\
Item B & 15.2 & 25.1 & 40.3 \\
Item C & 22.1 & 12.7 & 34.8 \\
\bottomrule
\end{tabular}
\caption{Sample Data Table}
\label{tab:sample}
\end{table}

\section{Lists}
\subsection{Bullet Points}
\begin{itemize}
  \item First item
  \item Second item
  \item Third item with sub-items:
    \begin{itemize}
      \item Sub-item 1
      \item Sub-item 2
    \end{itemize}
\end{itemize}

\subsection{Numbered List}
\begin{enumerate}
  \item First step
  \item Second step
  \item Third step
\end{enumerate}

\section{Citations and References}
When you use the citation search feature, references will be automatically managed. For example, you might cite a paper like this: \cite{example2024}. All citations will be automatically added to your \texttt{cite.bib} file.

\section{Conclusion}
This template demonstrates basic \LaTeX{} features including mathematical formulas, tables, and lists. The bibliography system is now set up to work seamlessly with Bibby AI's citation management. Feel free to modify it for your needs!

\bibliographystyle{plain}
\bibliography{cite}

\end{document}\documentclass{article}
\usepackage[utf8]{inputenc}
\usepackage{booktabs}
\usepackage{amsmath}
\usepackage{natbib}

\title{My Document}
\author{Your Name}
\date{\today}

\begin{document}

\maketitle

\section{Introduction}
This is a sample document demonstrating various \LaTeX{} features including mathematical equations, tables, and formatting. You can cite papers using the citation search feature, and they will be automatically added to your \texttt{cite.bib} file.

\section{Mathematical Expressions}
Here are some example equations:

\subsection{Quadratic Formula}
The quadratic formula is given by:
\[x = \frac{-b \pm \sqrt{b^2 - 4ac}}{2a}\]

\subsection{Matrix Example}
A sample matrix:
\[
\begin{pmatrix}
a_{11} & a_{12} & a_{13} \\
a_{21} & a_{22} & a_{23} \\
a_{31} & a_{32} & a_{33}
\end{pmatrix}
\]

\section{Tables}
Below is a sample table showing data:

\begin{table}[h]
\centering
\begin{tabular}{@{}lrrr@{}}
\toprule
\textbf{Category} & \textbf{Value 1} & \textbf{Value 2} & \textbf{Total} \\
\midrule
Item A & 10.5 & 20.3 & 30.8 \\
Item B & 15.2 & 25.1 & 40.3 \\
Item C & 22.1 & 12.7 & 34.8 \\
\bottomrule
\end{tabular}
\caption{Sample Data Table}
\label{tab:sample}
\end{table}

\section{Lists}
\subsection{Bullet Points}
\begin{itemize}
  \item First item
  \item Second item
  \item Third item with sub-items:
    \begin{itemize}
      \item Sub-item 1
      \item Sub-item 2
    \end{itemize}
\end{itemize}

\subsection{Numbered List}
\begin{enumerate}
  \item First step
  \item Second step
  \item Third step
\end{enumerate}

\section{Citations and References}
When you use the citation search feature, references will be automatically managed. For example, you might cite a paper like this: \cite{example2024}. All citations will be automatically added to your \texttt{cite.bib} file.

\section{Conclusion}
This template demonstrates basic \LaTeX{} features including mathematical formulas, tables, and lists. The bibliography system is now set up to work seamlessly with Bibby AI's citation management. Feel free to modify it for your needs!

\bibliographystyle{plain}
\bibliography{cite}

\end{document}\documentclass{article}
\usepackage[utf8]{inputenc}
\usepackage{booktabs}
\usepackage{amsmath}
\usepackage{natbib}

\title{My Document}
\author{Your Name}
\date{\today}

\begin{document}

\maketitle

\section{Introduction}
This is a sample document demonstrating various \LaTeX{} features including mathematical equations, tables, and formatting. You can cite papers using the citation search feature, and they will be automatically added to your \texttt{cite.bib} file.

\section{Mathematical Expressions}
Here are some example equations:

\subsection{Quadratic Formula}
The quadratic formula is given by:
\[x = \frac{-b \pm \sqrt{b^2 - 4ac}}{2a}\]

\subsection{Matrix Example}
A sample matrix:
\[
\begin{pmatrix}
a_{11} & a_{12} & a_{13} \\
a_{21} & a_{22} & a_{23} \\
a_{31} & a_{32} & a_{33}
\end{pmatrix}
\]

\section{Tables}
Below is a sample table showing data:

\begin{table}[h]
\centering
\begin{tabular}{@{}lrrr@{}}
\toprule
\textbf{Category} & \textbf{Value 1} & \textbf{Value 2} & \textbf{Total} \\
\midrule
Item A & 10.5 & 20.3 & 30.8 \\
Item B & 15.2 & 25.1 & 40.3 \\
Item C & 22.1 & 12.7 & 34.8 \\
\bottomrule
\end{tabular}
\caption{Sample Data Table}
\label{tab:sample}
\end{table}

\section{Lists}
\subsection{Bullet Points}
\begin{itemize}
  \item First item
  \item Second item
  \item Third item with sub-items:
    \begin{itemize}
      \item Sub-item 1
      \item Sub-item 2
    \end{itemize}
\end{itemize}

\subsection{Numbered List}
\begin{enumerate}
  \item First step
  \item Second step
  \item Third step
\end{enumerate}

\section{Citations and References}
When you use the citation search feature, references will be automatically managed. For example, you might cite a paper like this: \cite{example2024}. All citations will be automatically added to your \texttt{cite.bib} file.

\section{Conclusion}
This template demonstrates basic \LaTeX{} features including mathematical formulas, tables, and lists. The bibliography system is now set up to work seamlessly with Bibby AI's citation management. Feel free to modify it for your needs!

\bibliographystyle{plain}
\bibliography{cite}

\end{document}\documentclass{article}
\usepackage[utf8]{inputenc}
\usepackage{booktabs}
\usepackage{amsmath}
\usepackage{natbib}

\title{My Document}
\author{Your Name}
\date{\today}

\begin{document}

\maketitle

\section{Introduction}
This is a sample document demonstrating various \LaTeX{} features including mathematical equations, tables, and formatting. You can cite papers using the citation search feature, and they will be automatically added to your \texttt{cite.bib} file.

\section{Mathematical Expressions}
Here are some example equations:

\subsection{Quadratic Formula}
The quadratic formula is given by:
\[x = \frac{-b \pm \sqrt{b^2 - 4ac}}{2a}\]

\subsection{Matrix Example}
A sample matrix:
\[
\begin{pmatrix}
a_{11} & a_{12} & a_{13} \\
a_{21} & a_{22} & a_{23} \\
a_{31} & a_{32} & a_{33}
\end{pmatrix}
\]

\section{Tables}
Below is a sample table showing data:

\begin{table}[h]
\centering
\begin{tabular}{@{}lrrr@{}}
\toprule
\textbf{Category} & \textbf{Value 1} & \textbf{Value 2} & \textbf{Total} \\
\midrule
Item A & 10.5 & 20.3 & 30.8 \\
Item B & 15.2 & 25.1 & 40.3 \\
Item C & 22.1 & 12.7 & 34.8 \\
\bottomrule
\end{tabular}
\caption{Sample Data Table}
\label{tab:sample}
\end{table}

\section{Lists}
\subsection{Bullet Points}
\begin{itemize}
  \item First item
  \item Second item
  \item Third item with sub-items:
    \begin{itemize}
      \item Sub-item 1
      \item Sub-item 2
    \end{itemize}
\end{itemize}

\subsection{Numbered List}
\begin{enumerate}
  \item First step
  \item Second step
  \item Third step
\end{enumerate}

\section{Citations and References}
When you use the citation search feature, references will be automatically managed. For example, you might cite a paper like this: \cite{example2024}. All citations will be automatically added to your \texttt{cite.bib} file.

\section{Conclusion}
This template demonstrates basic \LaTeX{} features including mathematical formulas, tables, and lists. The bibliography system is now set up to work seamlessly with Bibby AI's citation management. Feel free to modify it for your needs!

\bibliographystyle{plain}
\bibliography{cite}

\end{document}\documentclass{article}
\usepackage[utf8]{inputenc}
\usepackage{booktabs}
\usepackage{amsmath}
\usepackage{natbib}

\title{My Document}
\author{Your Name}
\date{\today}

\begin{document}

\maketitle

\section{Introduction}
This is a sample document demonstrating various \LaTeX{} features including mathematical equations, tables, and formatting. You can cite papers using the citation search feature, and they will be automatically added to your \texttt{cite.bib} file.

\section{Mathematical Expressions}
Here are some example equations:

\subsection{Quadratic Formula}
The quadratic formula is given by:
\[x = \frac{-b \pm \sqrt{b^2 - 4ac}}{2a}\]

\subsection{Matrix Example}
A sample matrix:
\[
\begin{pmatrix}
a_{11} & a_{12} & a_{13} \\
a_{21} & a_{22} & a_{23} \\
a_{31} & a_{32} & a_{33}
\end{pmatrix}
\]

\section{Tables}
Below is a sample table showing data:

\begin{table}[h]
\centering
\begin{tabular}{@{}lrrr@{}}
\toprule
\textbf{Category} & \textbf{Value 1} & \textbf{Value 2} & \textbf{Total} \\
\midrule
Item A & 10.5 & 20.3 & 30.8 \\
Item B & 15.2 & 25.1 & 40.3 \\
Item C & 22.1 & 12.7 & 34.8 \\
\bottomrule
\end{tabular}
\caption{Sample Data Table}
\label{tab:sample}
\end{table}

\section{Lists}
\subsection{Bullet Points}
\begin{itemize}
  \item First item
  \item Second item
  \item Third item with sub-items:
    \begin{itemize}
      \item Sub-item 1
      \item Sub-item 2
    \end{itemize}
\end{itemize}

\subsection{Numbered List}
\begin{enumerate}
  \item First step
  \item Second step
  \item Third step
\end{enumerate}

\section{Citations and References}
When you use the citation search feature, references will be automatically managed. For example, you might cite a paper like this: \cite{example2024}. All citations will be automatically added to your \texttt{cite.bib} file.

\section{Conclusion}
This template demonstrates basic \LaTeX{} features including mathematical formulas, tables, and lists. The bibliography system is now set up to work seamlessly with Bibby AI's citation management. Feel free to modify it for your needs!

\bibliographystyle{plain}
\bibliography{cite}

\end{document}\documentclass{article}
\usepackage[utf8]{inputenc}
\usepackage{booktabs}
\usepackage{amsmath}
\usepackage{natbib}

\title{My Document}
\author{Your Name}
\date{\today}

\begin{document}

\maketitle

\section{Introduction}
This is a sample document demonstrating various \LaTeX{} features including mathematical equations, tables, and formatting. You can cite papers using the citation search feature, and they will be automatically added to your \texttt{cite.bib} file.

\section{Mathematical Expressions}
Here are some example equations:

\subsection{Quadratic Formula}
The quadratic formula is given by:
\[x = \frac{-b \pm \sqrt{b^2 - 4ac}}{2a}\]

\subsection{Matrix Example}
A sample matrix:
\[
\begin{pmatrix}
a_{11} & a_{12} & a_{13} \\
a_{21} & a_{22} & a_{23} \\
a_{31} & a_{32} & a_{33}
\end{pmatrix}
\]

\section{Tables}
Below is a sample table showing data:

\begin{table}[h]
\centering
\begin{tabular}{@{}lrrr@{}}
\toprule
\textbf{Category} & \textbf{Value 1} & \textbf{Value 2} & \textbf{Total} \\
\midrule
Item A & 10.5 & 20.3 & 30.8 \\
Item B & 15.2 & 25.1 & 40.3 \\
Item C & 22.1 & 12.7 & 34.8 \\
\bottomrule
\end{tabular}
\caption{Sample Data Table}
\label{tab:sample}
\end{table}

\section{Lists}
\subsection{Bullet Points}
\begin{itemize}
  \item First item
  \item Second item
  \item Third item with sub-items:
    \begin{itemize}
      \item Sub-item 1
      \item Sub-item 2
    \end{itemize}
\end{itemize}

\subsection{Numbered List}
\begin{enumerate}
  \item First step
  \item Second step
  \item Third step
\end{enumerate}

\section{Citations and References}
When you use the citation search feature, references will be automatically managed. For example, you might cite a paper like this: \cite{example2024}. All citations will be automatically added to your \texttt{cite.bib} file.

\section{Conclusion}
This template demonstrates basic \LaTeX{} features including mathematical formulas, tables, and lists. The bibliography system is now set up to work seamlessly with Bibby AI's citation management. Feel free to modify it for your needs!

\bibliographystyle{plain}
\bibliography{cite}

\end{document}\documentclass{article}
\usepackage[utf8]{inputenc}
\usepackage{booktabs}
\usepackage{amsmath}
\usepackage{natbib}

\title{My Document}
\author{Your Name}
\date{\today}

\begin{document}

\maketitle

\section{Introduction}
This is a sample document demonstrating various \LaTeX{} features including mathematical equations, tables, and formatting. You can cite papers using the citation search feature, and they will be automatically added to your \texttt{cite.bib} file.

\section{Mathematical Expressions}
Here are some example equations:

\subsection{Quadratic Formula}
The quadratic formula is given by:
\[x = \frac{-b \pm \sqrt{b^2 - 4ac}}{2a}\]

\subsection{Matrix Example}
A sample matrix:
\[
\begin{pmatrix}
a_{11} & a_{12} & a_{13} \\
a_{21} & a_{22} & a_{23} \\
a_{31} & a_{32} & a_{33}
\end{pmatrix}
\]

\section{Tables}
Below is a sample table showing data:

\begin{table}[h]
\centering
\begin{tabular}{@{}lrrr@{}}
\toprule
\textbf{Category} & \textbf{Value 1} & \textbf{Value 2} & \textbf{Total} \\
\midrule
Item A & 10.5 & 20.3 & 30.8 \\
Item B & 15.2 & 25.1 & 40.3 \\
Item C & 22.1 & 12.7 & 34.8 \\
\bottomrule
\end{tabular}
\caption{Sample Data Table}
\label{tab:sample}
\end{table}

\section{Lists}
\subsection{Bullet Points}
\begin{itemize}
  \item First item
  \item Second item
  \item Third item with sub-items:
    \begin{itemize}
      \item Sub-item 1
      \item Sub-item 2
    \end{itemize}
\end{itemize}

\subsection{Numbered List}
\begin{enumerate}
  \item First step
  \item Second step
  \item Third step
\end{enumerate}

\section{Citations and References}
When you use the citation search feature, references will be automatically managed. For example, you might cite a paper like this: \cite{example2024}. All citations will be automatically added to your \texttt{cite.bib} file.

\section{Conclusion}
This template demonstrates basic \LaTeX{} features including mathematical formulas, tables, and lists. The bibliography system is now set up to work seamlessly with Bibby AI's citation management. Feel free to modify it for your needs!

\bibliographystyle{plain}
\bibliography{cite}

\end{document}\documentclass{article}
\usepackage[utf8]{inputenc}
\usepackage{booktabs}
\usepackage{amsmath}
\usepackage{natbib}

\title{My Document}
\author{Your Name}
\date{\today}

\begin{document}

\maketitle

\section{Introduction}
This is a sample document demonstrating various \LaTeX{} features including mathematical equations, tables, and formatting. You can cite papers using the citation search feature, and they will be automatically added to your \texttt{cite.bib} file.

\section{Mathematical Expressions}
Here are some example equations:

\subsection{Quadratic Formula}
The quadratic formula is given by:
\[x = \frac{-b \pm \sqrt{b^2 - 4ac}}{2a}\]

\subsection{Matrix Example}
A sample matrix:
\[
\begin{pmatrix}
a_{11} & a_{12} & a_{13} \\
a_{21} & a_{22} & a_{23} \\
a_{31} & a_{32} & a_{33}
\end{pmatrix}
\]

\section{Tables}
Below is a sample table showing data:

\begin{table}[h]
\centering
\begin{tabular}{@{}lrrr@{}}
\toprule
\textbf{Category} & \textbf{Value 1} & \textbf{Value 2} & \textbf{Total} \\
\midrule
Item A & 10.5 & 20.3 & 30.8 \\
Item B & 15.2 & 25.1 & 40.3 \\
Item C & 22.1 & 12.7 & 34.8 \\
\bottomrule
\end{tabular}
\caption{Sample Data Table}
\label{tab:sample}
\end{table}

\section{Lists}
\subsection{Bullet Points}
\begin{itemize}
  \item First item
  \item Second item
  \item Third item with sub-items:
    \begin{itemize}
      \item Sub-item 1
      \item Sub-item 2
    \end{itemize}
\end{itemize}

\subsection{Numbered List}
\begin{enumerate}
  \item First step
  \item Second step
  \item Third step
\end{enumerate}

\section{Citations and References}
When you use the citation search feature, references will be automatically managed. For example, you might cite a paper like this: \cite{example2024}. All citations will be automatically added to your \texttt{cite.bib} file.

\section{Conclusion}
This template demonstrates basic \LaTeX{} features including mathematical formulas, tables, and lists. The bibliography system is now set up to work seamlessly with Bibby AI's citation management. Feel free to modify it for your needs!

\bibliographystyle{plain}
\bibliography{cite}

\end{document}\documentclass{article}
\usepackage[utf8]{inputenc}
\usepackage{booktabs}
\usepackage{amsmath}
\usepackage{natbib}

\title{My Document}
\author{Your Name}
\date{\today}

\begin{document}

\maketitle

\section{Introduction}
This is a sample document demonstrating various \LaTeX{} features including mathematical equations, tables, and formatting. You can cite papers using the citation search feature, and they will be automatically added to your \texttt{cite.bib} file.

\section{Mathematical Expressions}
Here are some example equations:

\subsection{Quadratic Formula}
The quadratic formula is given by:
\[x = \frac{-b \pm \sqrt{b^2 - 4ac}}{2a}\]

\subsection{Matrix Example}
A sample matrix:
\[
\begin{pmatrix}
a_{11} & a_{12} & a_{13} \\
a_{21} & a_{22} & a_{23} \\
a_{31} & a_{32} & a_{33}
\end{pmatrix}
\]

\section{Tables}
Below is a sample table showing data:

\begin{table}[h]
\centering
\begin{tabular}{@{}lrrr@{}}
\toprule
\textbf{Category} & \textbf{Value 1} & \textbf{Value 2} & \textbf{Total} \\
\midrule
Item A & 10.5 & 20.3 & 30.8 \\
Item B & 15.2 & 25.1 & 40.3 \\
Item C & 22.1 & 12.7 & 34.8 \\
\bottomrule
\end{tabular}
\caption{Sample Data Table}
\label{tab:sample}
\end{table}

\section{Lists}
\subsection{Bullet Points}
\begin{itemize}
  \item First item
  \item Second item
  \item Third item with sub-items:
    \begin{itemize}
      \item Sub-item 1
      \item Sub-item 2
    \end{itemize}
\end{itemize}

\subsection{Numbered List}
\begin{enumerate}
  \item First step
  \item Second step
  \item Third step
\end{enumerate}

\section{Citations and References}
When you use the citation search feature, references will be automatically managed. For example, you might cite a paper like this: \cite{example2024}. All citations will be automatically added to your \texttt{cite.bib} file.

\section{Conclusion}
This template demonstrates basic \LaTeX{} features including mathematical formulas, tables, and lists. The bibliography system is now set up to work seamlessly with Bibby AI's citation management. Feel free to modify it for your needs!

\bibliographystyle{plain}
\bibliography{cite}

\end{document}\documentclass{article}
\usepackage[utf8]{inputenc}
\usepackage{booktabs}
\usepackage{amsmath}
\usepackage{natbib}

\title{My Document}
\author{Your Name}
\date{\today}

\begin{document}

\maketitle

\section{Introduction}
This is a sample document demonstrating various \LaTeX{} features including mathematical equations, tables, and formatting. You can cite papers using the citation search feature, and they will be automatically added to your \texttt{cite.bib} file.

\section{Mathematical Expressions}
Here are some example equations:

\subsection{Quadratic Formula}
The quadratic formula is given by:
\[x = \frac{-b \pm \sqrt{b^2 - 4ac}}{2a}\]

\subsection{Matrix Example}
A sample matrix:
\[
\begin{pmatrix}
a_{11} & a_{12} & a_{13} \\
a_{21} & a_{22} & a_{23} \\
a_{31} & a_{32} & a_{33}
\end{pmatrix}
\]

\section{Tables}
Below is a sample table showing data:

\begin{table}[h]
\centering
\begin{tabular}{@{}lrrr@{}}
\toprule
\textbf{Category} & \textbf{Value 1} & \textbf{Value 2} & \textbf{Total} \\
\midrule
Item A & 10.5 & 20.3 & 30.8 \\
Item B & 15.2 & 25.1 & 40.3 \\
Item C & 22.1 & 12.7 & 34.8 \\
\bottomrule
\end{tabular}
\caption{Sample Data Table}
\label{tab:sample}
\end{table}

\section{Lists}
\subsection{Bullet Points}
\begin{itemize}
  \item First item
  \item Second item
  \item Third item with sub-items:
    \begin{itemize}
      \item Sub-item 1
      \item Sub-item 2
    \end{itemize}
\end{itemize}

\subsection{Numbered List}
\begin{enumerate}
  \item First step
  \item Second step
  \item Third step
\end{enumerate}

\section{Citations and References}
When you use the citation search feature, references will be automatically managed. For example, you might cite a paper like this: \cite{example2024}. All citations will be automatically added to your \texttt{cite.bib} file.

\section{Conclusion}
This template demonstrates basic \LaTeX{} features including mathematical formulas, tables, and lists. The bibliography system is now set up to work seamlessly with Bibby AI's citation management. Feel free to modify it for your needs!

\bibliographystyle{plain}
\bibliography{cite}

\end{document}\documentclass{article}
\usepackage[utf8]{inputenc}
\usepackage{booktabs}
\usepackage{amsmath}
\usepackage{natbib}

\title{My Document}
\author{Your Name}
\date{\today}

\begin{document}

\maketitle

\section{Introduction}
This is a sample document demonstrating various \LaTeX{} features including mathematical equations, tables, and formatting. You can cite papers using the citation search feature, and they will be automatically added to your \texttt{cite.bib} file.

\section{Mathematical Expressions}
Here are some example equations:

\subsection{Quadratic Formula}
The quadratic formula is given by:
\[x = \frac{-b \pm \sqrt{b^2 - 4ac}}{2a}\]

\subsection{Matrix Example}
A sample matrix:
\[
\begin{pmatrix}
a_{11} & a_{12} & a_{13} \\
a_{21} & a_{22} & a_{23} \\
a_{31} & a_{32} & a_{33}
\end{pmatrix}
\]

\section{Tables}
Below is a sample table showing data:

\begin{table}[h]
\centering
\begin{tabular}{@{}lrrr@{}}
\toprule
\textbf{Category} & \textbf{Value 1} & \textbf{Value 2} & \textbf{Total} \\
\midrule
Item A & 10.5 & 20.3 & 30.8 \\
Item B & 15.2 & 25.1 & 40.3 \\
Item C & 22.1 & 12.7 & 34.8 \\
\bottomrule
\end{tabular}
\caption{Sample Data Table}
\label{tab:sample}
\end{table}

\section{Lists}
\subsection{Bullet Points}
\begin{itemize}
  \item First item
  \item Second item
  \item Third item with sub-items:
    \begin{itemize}
      \item Sub-item 1
      \item Sub-item 2
    \end{itemize}
\end{itemize}

\subsection{Numbered List}
\begin{enumerate}
  \item First step
  \item Second step
  \item Third step
\end{enumerate}

\section{Citations and References}
When you use the citation search feature, references will be automatically managed. For example, you might cite a paper like this: \cite{example2024}. All citations will be automatically added to your \texttt{cite.bib} file.

\section{Conclusion}
This template demonstrates basic \LaTeX{} features including mathematical formulas, tables, and lists. The bibliography system is now set up to work seamlessly with Bibby AI's citation management. Feel free to modify it for your needs!

\bibliographystyle{plain}
\bibliography{cite}

\end{document}\documentclass{article}
\usepackage[utf8]{inputenc}
\usepackage{booktabs}
\usepackage{amsmath}
\usepackage{natbib}

\title{My Document}
\author{Your Name}
\date{\today}

\begin{document}

\maketitle

\section{Introduction}
This is a sample document demonstrating various \LaTeX{} features including mathematical equations, tables, and formatting. You can cite papers using the citation search feature, and they will be automatically added to your \texttt{cite.bib} file.

\section{Mathematical Expressions}
Here are some example equations:

\subsection{Quadratic Formula}
The quadratic formula is given by:
\[x = \frac{-b \pm \sqrt{b^2 - 4ac}}{2a}\]

\subsection{Matrix Example}
A sample matrix:
\[
\begin{pmatrix}
a_{11} & a_{12} & a_{13} \\
a_{21} & a_{22} & a_{23} \\
a_{31} & a_{32} & a_{33}
\end{pmatrix}
\]

\section{Tables}
Below is a sample table showing data:

\begin{table}[h]
\centering
\begin{tabular}{@{}lrrr@{}}
\toprule
\textbf{Category} & \textbf{Value 1} & \textbf{Value 2} & \textbf{Total} \\
\midrule
Item A & 10.5 & 20.3 & 30.8 \\
Item B & 15.2 & 25.1 & 40.3 \\
Item C & 22.1 & 12.7 & 34.8 \\
\bottomrule
\end{tabular}
\caption{Sample Data Table}
\label{tab:sample}
\end{table}

\section{Lists}
\subsection{Bullet Points}
\begin{itemize}
  \item First item
  \item Second item
  \item Third item with sub-items:
    \begin{itemize}
      \item Sub-item 1
      \item Sub-item 2
    \end{itemize}
\end{itemize}

\subsection{Numbered List}
\begin{enumerate}
  \item First step
  \item Second step
  \item Third step
\end{enumerate}

\section{Citations and References}
When you use the citation search feature, references will be automatically managed. For example, you might cite a paper like this: \cite{example2024}. All citations will be automatically added to your \texttt{cite.bib} file.

\section{Conclusion}
This template demonstrates basic \LaTeX{} features including mathematical formulas, tables, and lists. The bibliography system is now set up to work seamlessly with Bibby AI's citation management. Feel free to modify it for your needs!

\bibliographystyle{plain}
\bibliography{cite}

\end{document}\documentclass{article}
\usepackage[utf8]{inputenc}
\usepackage{booktabs}
\usepackage{amsmath}
\usepackage{natbib}

\title{My Document}
\author{Your Name}
\date{\today}

\begin{document}

\maketitle

\section{Introduction}
This is a sample document demonstrating various \LaTeX{} features including mathematical equations, tables, and formatting. You can cite papers using the citation search feature, and they will be automatically added to your \texttt{cite.bib} file.

\section{Mathematical Expressions}
Here are some example equations:

\subsection{Quadratic Formula}
The quadratic formula is given by:
\[x = \frac{-b \pm \sqrt{b^2 - 4ac}}{2a}\]

\subsection{Matrix Example}
A sample matrix:
\[
\begin{pmatrix}
a_{11} & a_{12} & a_{13} \\
a_{21} & a_{22} & a_{23} \\
a_{31} & a_{32} & a_{33}
\end{pmatrix}
\]

\section{Tables}
Below is a sample table showing data:

\begin{table}[h]
\centering
\begin{tabular}{@{}lrrr@{}}
\toprule
\textbf{Category} & \textbf{Value 1} & \textbf{Value 2} & \textbf{Total} \\
\midrule
Item A & 10.5 & 20.3 & 30.8 \\
Item B & 15.2 & 25.1 & 40.3 \\
Item C & 22.1 & 12.7 & 34.8 \\
\bottomrule
\end{tabular}
\caption{Sample Data Table}
\label{tab:sample}
\end{table}

\section{Lists}
\subsection{Bullet Points}
\begin{itemize}
  \item First item
  \item Second item
  \item Third item with sub-items:
    \begin{itemize}
      \item Sub-item 1
      \item Sub-item 2
    \end{itemize}
\end{itemize}

\subsection{Numbered List}
\begin{enumerate}
  \item First step
  \item Second step
  \item Third step
\end{enumerate}

\section{Citations and References}
When you use the citation search feature, references will be automatically managed. For example, you might cite a paper like this: \cite{example2024}. All citations will be automatically added to your \texttt{cite.bib} file.

\section{Conclusion}
This template demonstrates basic \LaTeX{} features including mathematical formulas, tables, and lists. The bibliography system is now set up to work seamlessly with Bibby AI's citation management. Feel free to modify it for your needs!

\bibliographystyle{plain}
\bibliography{cite}

\end{document}\documentclass{article}
\usepackage[utf8]{inputenc}
\usepackage{booktabs}
\usepackage{amsmath}
\usepackage{natbib}

\title{My Document}
\author{Your Name}
\date{\today}

\begin{document}

\maketitle

\section{Introduction}
This is a sample document demonstrating various \LaTeX{} features including mathematical equations, tables, and formatting. You can cite papers using the citation search feature, and they will be automatically added to your \texttt{cite.bib} file.

\section{Mathematical Expressions}
Here are some example equations:

\subsection{Quadratic Formula}
The quadratic formula is given by:
\[x = \frac{-b \pm \sqrt{b^2 - 4ac}}{2a}\]

\subsection{Matrix Example}
A sample matrix:
\[
\begin{pmatrix}
a_{11} & a_{12} & a_{13} \\
a_{21} & a_{22} & a_{23} \\
a_{31} & a_{32} & a_{33}
\end{pmatrix}
\]

\section{Tables}
Below is a sample table showing data:

\begin{table}[h]
\centering
\begin{tabular}{@{}lrrr@{}}
\toprule
\textbf{Category} & \textbf{Value 1} & \textbf{Value 2} & \textbf{Total} \\
\midrule
Item A & 10.5 & 20.3 & 30.8 \\
Item B & 15.2 & 25.1 & 40.3 \\
Item C & 22.1 & 12.7 & 34.8 \\
\bottomrule
\end{tabular}
\caption{Sample Data Table}
\label{tab:sample}
\end{table}

\section{Lists}
\subsection{Bullet Points}
\begin{itemize}
  \item First item
  \item Second item
  \item Third item with sub-items:
    \begin{itemize}
      \item Sub-item 1
      \item Sub-item 2
    \end{itemize}
\end{itemize}

\subsection{Numbered List}
\begin{enumerate}
  \item First step
  \item Second step
  \item Third step
\end{enumerate}

\section{Citations and References}
When you use the citation search feature, references will be automatically managed. For example, you might cite a paper like this: \cite{example2024}. All citations will be automatically added to your \texttt{cite.bib} file.

\section{Conclusion}
This template demonstrates basic \LaTeX{} features including mathematical formulas, tables, and lists. The bibliography system is now set up to work seamlessly with Bibby AI's citation management. Feel free to modify it for your needs!

\bibliographystyle{plain}
\bibliography{cite}

\end{document}\documentclass{article}
\usepackage[utf8]{inputenc}
\usepackage{booktabs}
\usepackage{amsmath}
\usepackage{natbib}

\title{My Document}
\author{Your Name}
\date{\today}

\begin{document}

\maketitle

\section{Introduction}
This is a sample document demonstrating various \LaTeX{} features including mathematical equations, tables, and formatting. You can cite papers using the citation search feature, and they will be automatically added to your \texttt{cite.bib} file.

\section{Mathematical Expressions}
Here are some example equations:

\subsection{Quadratic Formula}
The quadratic formula is given by:
\[x = \frac{-b \pm \sqrt{b^2 - 4ac}}{2a}\]

\subsection{Matrix Example}
A sample matrix:
\[
\begin{pmatrix}
a_{11} & a_{12} & a_{13} \\
a_{21} & a_{22} & a_{23} \\
a_{31} & a_{32} & a_{33}
\end{pmatrix}
\]

\section{Tables}
Below is a sample table showing data:

\begin{table}[h]
\centering
\begin{tabular}{@{}lrrr@{}}
\toprule
\textbf{Category} & \textbf{Value 1} & \textbf{Value 2} & \textbf{Total} \\
\midrule
Item A & 10.5 & 20.3 & 30.8 \\
Item B & 15.2 & 25.1 & 40.3 \\
Item C & 22.1 & 12.7 & 34.8 \\
\bottomrule
\end{tabular}
\caption{Sample Data Table}
\label{tab:sample}
\end{table}

\section{Lists}
\subsection{Bullet Points}
\begin{itemize}
  \item First item
  \item Second item
  \item Third item with sub-items:
    \begin{itemize}
      \item Sub-item 1
      \item Sub-item 2
    \end{itemize}
\end{itemize}

\subsection{Numbered List}
\begin{enumerate}
  \item First step
  \item Second step
  \item Third step
\end{enumerate}

\section{Citations and References}
When you use the citation search feature, references will be automatically managed. For example, you might cite a paper like this: \cite{example2024}. All citations will be automatically added to your \texttt{cite.bib} file.

\section{Conclusion}
This template demonstrates basic \LaTeX{} features including mathematical formulas, tables, and lists. The bibliography system is now set up to work seamlessly with Bibby AI's citation management. Feel free to modify it for your needs!

\bibliographystyle{plain}
\bibliography{cite}

\end{document}\documentclass{article}
\usepackage[utf8]{inputenc}
\usepackage{booktabs}
\usepackage{amsmath}
\usepackage{natbib}

\title{My Document}
\author{Your Name}
\date{\today}

\begin{document}

\maketitle

\section{Introduction}
This is a sample document demonstrating various \LaTeX{} features including mathematical equations, tables, and formatting. You can cite papers using the citation search feature, and they will be automatically added to your \texttt{cite.bib} file.

\section{Mathematical Expressions}
Here are some example equations:

\subsection{Quadratic Formula}
The quadratic formula is given by:
\[x = \frac{-b \pm \sqrt{b^2 - 4ac}}{2a}\]

\subsection{Matrix Example}
A sample matrix:
\[
\begin{pmatrix}
a_{11} & a_{12} & a_{13} \\
a_{21} & a_{22} & a_{23} \\
a_{31} & a_{32} & a_{33}
\end{pmatrix}
\]

\section{Tables}
Below is a sample table showing data:

\begin{table}[h]
\centering
\begin{tabular}{@{}lrrr@{}}
\toprule
\textbf{Category} & \textbf{Value 1} & \textbf{Value 2} & \textbf{Total} \\
\midrule
Item A & 10.5 & 20.3 & 30.8 \\
Item B & 15.2 & 25.1 & 40.3 \\
Item C & 22.1 & 12.7 & 34.8 \\
\bottomrule
\end{tabular}
\caption{Sample Data Table}
\label{tab:sample}
\end{table}

\section{Lists}
\subsection{Bullet Points}
\begin{itemize}
  \item First item
  \item Second item
  \item Third item with sub-items:
    \begin{itemize}
      \item Sub-item 1
      \item Sub-item 2
    \end{itemize}
\end{itemize}

\subsection{Numbered List}
\begin{enumerate}
  \item First step
  \item Second step
  \item Third step
\end{enumerate}

\section{Citations and References}
When you use the citation search feature, references will be automatically managed. For example, you might cite a paper like this: \cite{example2024}. All citations will be automatically added to your \texttt{cite.bib} file.

\section{Conclusion}
This template demonstrates basic \LaTeX{} features including mathematical formulas, tables, and lists. The bibliography system is now set up to work seamlessly with Bibby AI's citation management. Feel free to modify it for your needs!

\bibliographystyle{plain}
\bibliography{cite}

\end{document}\documentclass{article}
\usepackage[utf8]{inputenc}
\usepackage{booktabs}
\usepackage{amsmath}
\usepackage{natbib}

\title{My Document}
\author{Your Name}
\date{\today}

\begin{document}

\maketitle

\section{Introduction}
This is a sample document demonstrating various \LaTeX{} features including mathematical equations, tables, and formatting. You can cite papers using the citation search feature, and they will be automatically added to your \texttt{cite.bib} file.

\section{Mathematical Expressions}
Here are some example equations:

\subsection{Quadratic Formula}
The quadratic formula is given by:
\[x = \frac{-b \pm \sqrt{b^2 - 4ac}}{2a}\]

\subsection{Matrix Example}
A sample matrix:
\[
\begin{pmatrix}
a_{11} & a_{12} & a_{13} \\
a_{21} & a_{22} & a_{23} \\
a_{31} & a_{32} & a_{33}
\end{pmatrix}
\]

\section{Tables}
Below is a sample table showing data:

\begin{table}[h]
\centering
\begin{tabular}{@{}lrrr@{}}
\toprule
\textbf{Category} & \textbf{Value 1} & \textbf{Value 2} & \textbf{Total} \\
\midrule
Item A & 10.5 & 20.3 & 30.8 \\
Item B & 15.2 & 25.1 & 40.3 \\
Item C & 22.1 & 12.7 & 34.8 \\
\bottomrule
\end{tabular}
\caption{Sample Data Table}
\label{tab:sample}
\end{table}

\section{Lists}
\subsection{Bullet Points}
\begin{itemize}
  \item First item
  \item Second item
  \item Third item with sub-items:
    \begin{itemize}
      \item Sub-item 1
      \item Sub-item 2
    \end{itemize}
\end{itemize}

\subsection{Numbered List}
\begin{enumerate}
  \item First step
  \item Second step
  \item Third step
\end{enumerate}

\section{Citations and References}
When you use the citation search feature, references will be automatically managed. For example, you might cite a paper like this: \cite{example2024}. All citations will be automatically added to your \texttt{cite.bib} file.

\section{Conclusion}
This template demonstrates basic \LaTeX{} features including mathematical formulas, tables, and lists. The bibliography system is now set up to work seamlessly with Bibby AI's citation management. Feel free to modify it for your needs!

\bibliographystyle{plain}
\bibliography{cite}

\end{document}\documentclass{article}
\usepackage[utf8]{inputenc}
\usepackage{booktabs}
\usepackage{amsmath}
\usepackage{natbib}

\title{My Document}
\author{Your Name}
\date{\today}

\begin{document}

\maketitle

\section{Introduction}
This is a sample document demonstrating various \LaTeX{} features including mathematical equations, tables, and formatting. You can cite papers using the citation search feature, and they will be automatically added to your \texttt{cite.bib} file.

\section{Mathematical Expressions}
Here are some example equations:

\subsection{Quadratic Formula}
The quadratic formula is given by:
\[x = \frac{-b \pm \sqrt{b^2 - 4ac}}{2a}\]

\subsection{Matrix Example}
A sample matrix:
\[
\begin{pmatrix}
a_{11} & a_{12} & a_{13} \\
a_{21} & a_{22} & a_{23} \\
a_{31} & a_{32} & a_{33}
\end{pmatrix}
\]

\section{Tables}
Below is a sample table showing data:

\begin{table}[h]
\centering
\begin{tabular}{@{}lrrr@{}}
\toprule
\textbf{Category} & \textbf{Value 1} & \textbf{Value 2} & \textbf{Total} \\
\midrule
Item A & 10.5 & 20.3 & 30.8 \\
Item B & 15.2 & 25.1 & 40.3 \\
Item C & 22.1 & 12.7 & 34.8 \\
\bottomrule
\end{tabular}
\caption{Sample Data Table}
\label{tab:sample}
\end{table}

\section{Lists}
\subsection{Bullet Points}
\begin{itemize}
  \item First item
  \item Second item
  \item Third item with sub-items:
    \begin{itemize}
      \item Sub-item 1
      \item Sub-item 2
    \end{itemize}
\end{itemize}

\subsection{Numbered List}
\begin{enumerate}
  \item First step
  \item Second step
  \item Third step
\end{enumerate}

\section{Citations and References}
When you use the citation search feature, references will be automatically managed. For example, you might cite a paper like this: \cite{example2024}. All citations will be automatically added to your \texttt{cite.bib} file.

\section{Conclusion}
This template demonstrates basic \LaTeX{} features including mathematical formulas, tables, and lists. The bibliography system is now set up to work seamlessly with Bibby AI's citation management. Feel free to modify it for your needs!

\bibliographystyle{plain}
\bibliography{cite}

\end{document}\documentclass{article}
\usepackage[utf8]{inputenc}
\usepackage{booktabs}
\usepackage{amsmath}
\usepackage{natbib}

\title{My Document}
\author{Your Name}
\date{\today}

\begin{document}

\maketitle

\section{Introduction}
This is a sample document demonstrating various \LaTeX{} features including mathematical equations, tables, and formatting. You can cite papers using the citation search feature, and they will be automatically added to your \texttt{cite.bib} file.

\section{Mathematical Expressions}
Here are some example equations:

\subsection{Quadratic Formula}
The quadratic formula is given by:
\[x = \frac{-b \pm \sqrt{b^2 - 4ac}}{2a}\]

\subsection{Matrix Example}
A sample matrix:
\[
\begin{pmatrix}
a_{11} & a_{12} & a_{13} \\
a_{21} & a_{22} & a_{23} \\
a_{31} & a_{32} & a_{33}
\end{pmatrix}
\]

\section{Tables}
Below is a sample table showing data:

\begin{table}[h]
\centering
\begin{tabular}{@{}lrrr@{}}
\toprule
\textbf{Category} & \textbf{Value 1} & \textbf{Value 2} & \textbf{Total} \\
\midrule
Item A & 10.5 & 20.3 & 30.8 \\
Item B & 15.2 & 25.1 & 40.3 \\
Item C & 22.1 & 12.7 & 34.8 \\
\bottomrule
\end{tabular}
\caption{Sample Data Table}
\label{tab:sample}
\end{table}

\section{Lists}
\subsection{Bullet Points}
\begin{itemize}
  \item First item
  \item Second item
  \item Third item with sub-items:
    \begin{itemize}
      \item Sub-item 1
      \item Sub-item 2
    \end{itemize}
\end{itemize}

\subsection{Numbered List}
\begin{enumerate}
  \item First step
  \item Second step
  \item Third step
\end{enumerate}

\section{Citations and References}
When you use the citation search feature, references will be automatically managed. For example, you might cite a paper like this: \cite{example2024}. All citations will be automatically added to your \texttt{cite.bib} file.

\section{Conclusion}
This template demonstrates basic \LaTeX{} features including mathematical formulas, tables, and lists. The bibliography system is now set up to work seamlessly with Bibby AI's citation management. Feel free to modify it for your needs!

\bibliographystyle{plain}
\bibliography{cite}

\end{document}\documentclass{article}
\usepackage[utf8]{inputenc}
\usepackage{booktabs}
\usepackage{amsmath}
\usepackage{natbib}

\title{My Document}
\author{Your Name}
\date{\today}

\begin{document}

\maketitle

\section{Introduction}
This is a sample document demonstrating various \LaTeX{} features including mathematical equations, tables, and formatting. You can cite papers using the citation search feature, and they will be automatically added to your \texttt{cite.bib} file.

\section{Mathematical Expressions}
Here are some example equations:

\subsection{Quadratic Formula}
The quadratic formula is given by:
\[x = \frac{-b \pm \sqrt{b^2 - 4ac}}{2a}\]

\subsection{Matrix Example}
A sample matrix:
\[
\begin{pmatrix}
a_{11} & a_{12} & a_{13} \\
a_{21} & a_{22} & a_{23} \\
a_{31} & a_{32} & a_{33}
\end{pmatrix}
\]

\section{Tables}
Below is a sample table showing data:

\begin{table}[h]
\centering
\begin{tabular}{@{}lrrr@{}}
\toprule
\textbf{Category} & \textbf{Value 1} & \textbf{Value 2} & \textbf{Total} \\
\midrule
Item A & 10.5 & 20.3 & 30.8 \\
Item B & 15.2 & 25.1 & 40.3 \\
Item C & 22.1 & 12.7 & 34.8 \\
\bottomrule
\end{tabular}
\caption{Sample Data Table}
\label{tab:sample}
\end{table}

\section{Lists}
\subsection{Bullet Points}
\begin{itemize}
  \item First item
  \item Second item
  \item Third item with sub-items:
    \begin{itemize}
      \item Sub-item 1
      \item Sub-item 2
    \end{itemize}
\end{itemize}

\subsection{Numbered List}
\begin{enumerate}
  \item First step
  \item Second step
  \item Third step
\end{enumerate}

\section{Citations and References}
When you use the citation search feature, references will be automatically managed. For example, you might cite a paper like this: \cite{example2024}. All citations will be automatically added to your \texttt{cite.bib} file.

\section{Conclusion}
This template demonstrates basic \LaTeX{} features including mathematical formulas, tables, and lists. The bibliography system is now set up to work seamlessly with Bibby AI's citation management. Feel free to modify it for your needs!

\bibliographystyle{plain}
\bibliography{cite}

\end{document}\documentclass{article}
\usepackage[utf8]{inputenc}
\usepackage{booktabs}
\usepackage{amsmath}
\usepackage{natbib}

\title{My Document}
\author{Your Name}
\date{\today}

\begin{document}

\maketitle

\section{Introduction}
This is a sample document demonstrating various \LaTeX{} features including mathematical equations, tables, and formatting. You can cite papers using the citation search feature, and they will be automatically added to your \texttt{cite.bib} file.

\section{Mathematical Expressions}
Here are some example equations:

\subsection{Quadratic Formula}
The quadratic formula is given by:
\[x = \frac{-b \pm \sqrt{b^2 - 4ac}}{2a}\]

\subsection{Matrix Example}
A sample matrix:
\[
\begin{pmatrix}
a_{11} & a_{12} & a_{13} \\
a_{21} & a_{22} & a_{23} \\
a_{31} & a_{32} & a_{33}
\end{pmatrix}
\]

\section{Tables}
Below is a sample table showing data:

\begin{table}[h]
\centering
\begin{tabular}{@{}lrrr@{}}
\toprule
\textbf{Category} & \textbf{Value 1} & \textbf{Value 2} & \textbf{Total} \\
\midrule
Item A & 10.5 & 20.3 & 30.8 \\
Item B & 15.2 & 25.1 & 40.3 \\
Item C & 22.1 & 12.7 & 34.8 \\
\bottomrule
\end{tabular}
\caption{Sample Data Table}
\label{tab:sample}
\end{table}

\section{Lists}
\subsection{Bullet Points}
\begin{itemize}
  \item First item
  \item Second item
  \item Third item with sub-items:
    \begin{itemize}
      \item Sub-item 1
      \item Sub-item 2
    \end{itemize}
\end{itemize}

\subsection{Numbered List}
\begin{enumerate}
  \item First step
  \item Second step
  \item Third step
\end{enumerate}

\section{Citations and References}
When you use the citation search feature, references will be automatically managed. For example, you might cite a paper like this: \cite{example2024}. All citations will be automatically added to your \texttt{cite.bib} file.

\section{Conclusion}
This template demonstrates basic \LaTeX{} features including mathematical formulas, tables, and lists. The bibliography system is now set up to work seamlessly with Bibby AI's citation management. Feel free to modify it for your needs!

\bibliographystyle{plain}
\bibliography{cite}

\end{document}\documentclass{article}
\usepackage[utf8]{inputenc}
\usepackage{booktabs}
\usepackage{amsmath}
\usepackage{natbib}

\title{My Document}
\author{Your Name}
\date{\today}

\begin{document}

\maketitle

\section{Introduction}
This is a sample document demonstrating various \LaTeX{} features including mathematical equations, tables, and formatting. You can cite papers using the citation search feature, and they will be automatically added to your \texttt{cite.bib} file.

\section{Mathematical Expressions}
Here are some example equations:

\subsection{Quadratic Formula}
The quadratic formula is given by:
\[x = \frac{-b \pm \sqrt{b^2 - 4ac}}{2a}\]

\subsection{Matrix Example}
A sample matrix:
\[
\begin{pmatrix}
a_{11} & a_{12} & a_{13} \\
a_{21} & a_{22} & a_{23} \\
a_{31} & a_{32} & a_{33}
\end{pmatrix}
\]

\section{Tables}
Below is a sample table showing data:

\begin{table}[h]
\centering
\begin{tabular}{@{}lrrr@{}}
\toprule
\textbf{Category} & \textbf{Value 1} & \textbf{Value 2} & \textbf{Total} \\
\midrule
Item A & 10.5 & 20.3 & 30.8 \\
Item B & 15.2 & 25.1 & 40.3 \\
Item C & 22.1 & 12.7 & 34.8 \\
\bottomrule
\end{tabular}
\caption{Sample Data Table}
\label{tab:sample}
\end{table}

\section{Lists}
\subsection{Bullet Points}
\begin{itemize}
  \item First item
  \item Second item
  \item Third item with sub-items:
    \begin{itemize}
      \item Sub-item 1
      \item Sub-item 2
    \end{itemize}
\end{itemize}

\subsection{Numbered List}
\begin{enumerate}
  \item First step
  \item Second step
  \item Third step
\end{enumerate}

\section{Citations and References}
When you use the citation search feature, references will be automatically managed. For example, you might cite a paper like this: \cite{example2024}. All citations will be automatically added to your \texttt{cite.bib} file.

\section{Conclusion}
This template demonstrates basic \LaTeX{} features including mathematical formulas, tables, and lists. The bibliography system is now set up to work seamlessly with Bibby AI's citation management. Feel free to modify it for your needs!

\bibliographystyle{plain}
\bibliography{cite}

\end{document}\documentclass{article}
\usepackage[utf8]{inputenc}
\usepackage{booktabs}
\usepackage{amsmath}
\usepackage{natbib}

\title{My Document}
\author{Your Name}
\date{\today}

\begin{document}

\maketitle

\section{Introduction}
This is a sample document demonstrating various \LaTeX{} features including mathematical equations, tables, and formatting. You can cite papers using the citation search feature, and they will be automatically added to your \texttt{cite.bib} file.

\section{Mathematical Expressions}
Here are some example equations:

\subsection{Quadratic Formula}
The quadratic formula is given by:
\[x = \frac{-b \pm \sqrt{b^2 - 4ac}}{2a}\]

\subsection{Matrix Example}
A sample matrix:
\[
\begin{pmatrix}
a_{11} & a_{12} & a_{13} \\
a_{21} & a_{22} & a_{23} \\
a_{31} & a_{32} & a_{33}
\end{pmatrix}
\]

\section{Tables}
Below is a sample table showing data:

\begin{table}[h]
\centering
\begin{tabular}{@{}lrrr@{}}
\toprule
\textbf{Category} & \textbf{Value 1} & \textbf{Value 2} & \textbf{Total} \\
\midrule
Item A & 10.5 & 20.3 & 30.8 \\
Item B & 15.2 & 25.1 & 40.3 \\
Item C & 22.1 & 12.7 & 34.8 \\
\bottomrule
\end{tabular}
\caption{Sample Data Table}
\label{tab:sample}
\end{table}

\section{Lists}
\subsection{Bullet Points}
\begin{itemize}
  \item First item
  \item Second item
  \item Third item with sub-items:
    \begin{itemize}
      \item Sub-item 1
      \item Sub-item 2
    \end{itemize}
\end{itemize}

\subsection{Numbered List}
\begin{enumerate}
  \item First step
  \item Second step
  \item Third step
\end{enumerate}

\section{Citations and References}
When you use the citation search feature, references will be automatically managed. For example, you might cite a paper like this: \cite{example2024}. All citations will be automatically added to your \texttt{cite.bib} file.

\section{Conclusion}
This template demonstrates basic \LaTeX{} features including mathematical formulas, tables, and lists. The bibliography system is now set up to work seamlessly with Bibby AI's citation management. Feel free to modify it for your needs!

\bibliographystyle{plain}
\bibliography{cite}

\end{document}\documentclass{article}
\usepackage[utf8]{inputenc}
\usepackage{booktabs}
\usepackage{amsmath}
\usepackage{natbib}

\title{My Document}
\author{Your Name}
\date{\today}

\begin{document}

\maketitle

\section{Introduction}
This is a sample document demonstrating various \LaTeX{} features including mathematical equations, tables, and formatting. You can cite papers using the citation search feature, and they will be automatically added to your \texttt{cite.bib} file.

\section{Mathematical Expressions}
Here are some example equations:

\subsection{Quadratic Formula}
The quadratic formula is given by:
\[x = \frac{-b \pm \sqrt{b^2 - 4ac}}{2a}\]

\subsection{Matrix Example}
A sample matrix:
\[
\begin{pmatrix}
a_{11} & a_{12} & a_{13} \\
a_{21} & a_{22} & a_{23} \\
a_{31} & a_{32} & a_{33}
\end{pmatrix}
\]

\section{Tables}
Below is a sample table showing data:

\begin{table}[h]
\centering
\begin{tabular}{@{}lrrr@{}}
\toprule
\textbf{Category} & \textbf{Value 1} & \textbf{Value 2} & \textbf{Total} \\
\midrule
Item A & 10.5 & 20.3 & 30.8 \\
Item B & 15.2 & 25.1 & 40.3 \\
Item C & 22.1 & 12.7 & 34.8 \\
\bottomrule
\end{tabular}
\caption{Sample Data Table}
\label{tab:sample}
\end{table}

\section{Lists}
\subsection{Bullet Points}
\begin{itemize}
  \item First item
  \item Second item
  \item Third item with sub-items:
    \begin{itemize}
      \item Sub-item 1
      \item Sub-item 2
    \end{itemize}
\end{itemize}

\subsection{Numbered List}
\begin{enumerate}
  \item First step
  \item Second step
  \item Third step
\end{enumerate}

\section{Citations and References}
When you use the citation search feature, references will be automatically managed. For example, you might cite a paper like this: \cite{example2024}. All citations will be automatically added to your \texttt{cite.bib} file.

\section{Conclusion}
This template demonstrates basic \LaTeX{} features including mathematical formulas, tables, and lists. The bibliography system is now set up to work seamlessly with Bibby AI's citation management. Feel free to modify it for your needs!

\bibliographystyle{plain}
\bibliography{cite}

\end{document}\documentclass{article}
\usepackage[utf8]{inputenc}
\usepackage{booktabs}
\usepackage{amsmath}
\usepackage{natbib}

\title{My Document}
\author{Your Name}
\date{\today}

\begin{document}

\maketitle

\section{Introduction}
This is a sample document demonstrating various \LaTeX{} features including mathematical equations, tables, and formatting. You can cite papers using the citation search feature, and they will be automatically added to your \texttt{cite.bib} file.

\section{Mathematical Expressions}
Here are some example equations:

\subsection{Quadratic Formula}
The quadratic formula is given by:
\[x = \frac{-b \pm \sqrt{b^2 - 4ac}}{2a}\]

\subsection{Matrix Example}
A sample matrix:
\[
\begin{pmatrix}
a_{11} & a_{12} & a_{13} \\
a_{21} & a_{22} & a_{23} \\
a_{31} & a_{32} & a_{33}
\end{pmatrix}
\]

\section{Tables}
Below is a sample table showing data:

\begin{table}[h]
\centering
\begin{tabular}{@{}lrrr@{}}
\toprule
\textbf{Category} & \textbf{Value 1} & \textbf{Value 2} & \textbf{Total} \\
\midrule
Item A & 10.5 & 20.3 & 30.8 \\
Item B & 15.2 & 25.1 & 40.3 \\
Item C & 22.1 & 12.7 & 34.8 \\
\bottomrule
\end{tabular}
\caption{Sample Data Table}
\label{tab:sample}
\end{table}

\section{Lists}
\subsection{Bullet Points}
\begin{itemize}
  \item First item
  \item Second item
  \item Third item with sub-items:
    \begin{itemize}
      \item Sub-item 1
      \item Sub-item 2
    \end{itemize}
\end{itemize}

\subsection{Numbered List}
\begin{enumerate}
  \item First step
  \item Second step
  \item Third step
\end{enumerate}

\section{Citations and References}
When you use the citation search feature, references will be automatically managed. For example, you might cite a paper like this: \cite{example2024}. All citations will be automatically added to your \texttt{cite.bib} file.

\section{Conclusion}
This template demonstrates basic \LaTeX{} features including mathematical formulas, tables, and lists. The bibliography system is now set up to work seamlessly with Bibby AI's citation management. Feel free to modify it for your needs!

\bibliographystyle{plain}
\bibliography{cite}

\end{document}\documentclass{article}
\usepackage[utf8]{inputenc}
\usepackage{booktabs}
\usepackage{amsmath}
\usepackage{natbib}

\title{My Document}
\author{Your Name}
\date{\today}

\begin{document}

\maketitle

\section{Introduction}
This is a sample document demonstrating various \LaTeX{} features including mathematical equations, tables, and formatting. You can cite papers using the citation search feature, and they will be automatically added to your \texttt{cite.bib} file.

\section{Mathematical Expressions}
Here are some example equations:

\subsection{Quadratic Formula}
The quadratic formula is given by:
\[x = \frac{-b \pm \sqrt{b^2 - 4ac}}{2a}\]

\subsection{Matrix Example}
A sample matrix:
\[
\begin{pmatrix}
a_{11} & a_{12} & a_{13} \\
a_{21} & a_{22} & a_{23} \\
a_{31} & a_{32} & a_{33}
\end{pmatrix}
\]

\section{Tables}
Below is a sample table showing data:

\begin{table}[h]
\centering
\begin{tabular}{@{}lrrr@{}}
\toprule
\textbf{Category} & \textbf{Value 1} & \textbf{Value 2} & \textbf{Total} \\
\midrule
Item A & 10.5 & 20.3 & 30.8 \\
Item B & 15.2 & 25.1 & 40.3 \\
Item C & 22.1 & 12.7 & 34.8 \\
\bottomrule
\end{tabular}
\caption{Sample Data Table}
\label{tab:sample}
\end{table}

\section{Lists}
\subsection{Bullet Points}
\begin{itemize}
  \item First item
  \item Second item
  \item Third item with sub-items:
    \begin{itemize}
      \item Sub-item 1
      \item Sub-item 2
    \end{itemize}
\end{itemize}

\subsection{Numbered List}
\begin{enumerate}
  \item First step
  \item Second step
  \item Third step
\end{enumerate}

\section{Citations and References}
When you use the citation search feature, references will be automatically managed. For example, you might cite a paper like this: \cite{example2024}. All citations will be automatically added to your \texttt{cite.bib} file.

\section{Conclusion}
This template demonstrates basic \LaTeX{} features including mathematical formulas, tables, and lists. The bibliography system is now set up to work seamlessly with Bibby AI's citation management. Feel free to modify it for your needs!

\bibliographystyle{plain}
\bibliography{cite}

\end{document}\documentclass{article}
\usepackage[utf8]{inputenc}
\usepackage{booktabs}
\usepackage{amsmath}
\usepackage{natbib}

\title{My Document}
\author{Your Name}
\date{\today}

\begin{document}

\maketitle

\section{Introduction}
This is a sample document demonstrating various \LaTeX{} features including mathematical equations, tables, and formatting. You can cite papers using the citation search feature, and they will be automatically added to your \texttt{cite.bib} file.

\section{Mathematical Expressions}
Here are some example equations:

\subsection{Quadratic Formula}
The quadratic formula is given by:
\[x = \frac{-b \pm \sqrt{b^2 - 4ac}}{2a}\]

\subsection{Matrix Example}
A sample matrix:
\[
\begin{pmatrix}
a_{11} & a_{12} & a_{13} \\
a_{21} & a_{22} & a_{23} \\
a_{31} & a_{32} & a_{33}
\end{pmatrix}
\]

\section{Tables}
Below is a sample table showing data:

\begin{table}[h]
\centering
\begin{tabular}{@{}lrrr@{}}
\toprule
\textbf{Category} & \textbf{Value 1} & \textbf{Value 2} & \textbf{Total} \\
\midrule
Item A & 10.5 & 20.3 & 30.8 \\
Item B & 15.2 & 25.1 & 40.3 \\
Item C & 22.1 & 12.7 & 34.8 \\
\bottomrule
\end{tabular}
\caption{Sample Data Table}
\label{tab:sample}
\end{table}

\section{Lists}
\subsection{Bullet Points}
\begin{itemize}
  \item First item
  \item Second item
  \item Third item with sub-items:
    \begin{itemize}
      \item Sub-item 1
      \item Sub-item 2
    \end{itemize}
\end{itemize}

\subsection{Numbered List}
\begin{enumerate}
  \item First step
  \item Second step
  \item Third step
\end{enumerate}

\section{Citations and References}
When you use the citation search feature, references will be automatically managed. For example, you might cite a paper like this: \cite{example2024}. All citations will be automatically added to your \texttt{cite.bib} file.

\section{Conclusion}
This template demonstrates basic \LaTeX{} features including mathematical formulas, tables, and lists. The bibliography system is now set up to work seamlessly with Bibby AI's citation management. Feel free to modify it for your needs!

\bibliographystyle{plain}
\bibliography{cite}

\end{document}\documentclass{article}
\usepackage[utf8]{inputenc}
\usepackage{booktabs}
\usepackage{amsmath}
\usepackage{natbib}

\title{My Document}
\author{Your Name}
\date{\today}

\begin{document}

\maketitle

\section{Introduction}
This is a sample document demonstrating various \LaTeX{} features including mathematical equations, tables, and formatting. You can cite papers using the citation search feature, and they will be automatically added to your \texttt{cite.bib} file.

\section{Mathematical Expressions}
Here are some example equations:

\subsection{Quadratic Formula}
The quadratic formula is given by:
\[x = \frac{-b \pm \sqrt{b^2 - 4ac}}{2a}\]

\subsection{Matrix Example}
A sample matrix:
\[
\begin{pmatrix}
a_{11} & a_{12} & a_{13} \\
a_{21} & a_{22} & a_{23} \\
a_{31} & a_{32} & a_{33}
\end{pmatrix}
\]

\section{Tables}
Below is a sample table showing data:

\begin{table}[h]
\centering
\begin{tabular}{@{}lrrr@{}}
\toprule
\textbf{Category} & \textbf{Value 1} & \textbf{Value 2} & \textbf{Total} \\
\midrule
Item A & 10.5 & 20.3 & 30.8 \\
Item B & 15.2 & 25.1 & 40.3 \\
Item C & 22.1 & 12.7 & 34.8 \\
\bottomrule
\end{tabular}
\caption{Sample Data Table}
\label{tab:sample}
\end{table}

\section{Lists}
\subsection{Bullet Points}
\begin{itemize}
  \item First item
  \item Second item
  \item Third item with sub-items:
    \begin{itemize}
      \item Sub-item 1
      \item Sub-item 2
    \end{itemize}
\end{itemize}

\subsection{Numbered List}
\begin{enumerate}
  \item First step
  \item Second step
  \item Third step
\end{enumerate}

\section{Citations and References}
When you use the citation search feature, references will be automatically managed. For example, you might cite a paper like this: \cite{example2024}. All citations will be automatically added to your \texttt{cite.bib} file.

\section{Conclusion}
This template demonstrates basic \LaTeX{} features including mathematical formulas, tables, and lists. The bibliography system is now set up to work seamlessly with Bibby AI's citation management. Feel free to modify it for your needs!

\bibliographystyle{plain}
\bibliography{cite}

\end{document}\documentclass{article}
\usepackage[utf8]{inputenc}
\usepackage{booktabs}
\usepackage{amsmath}
\usepackage{natbib}

\title{My Document}
\author{Your Name}
\date{\today}

\begin{document}

\maketitle

\section{Introduction}
This is a sample document demonstrating various \LaTeX{} features including mathematical equations, tables, and formatting. You can cite papers using the citation search feature, and they will be automatically added to your \texttt{cite.bib} file.

\section{Mathematical Expressions}
Here are some example equations:

\subsection{Quadratic Formula}
The quadratic formula is given by:
\[x = \frac{-b \pm \sqrt{b^2 - 4ac}}{2a}\]

\subsection{Matrix Example}
A sample matrix:
\[
\begin{pmatrix}
a_{11} & a_{12} & a_{13} \\
a_{21} & a_{22} & a_{23} \\
a_{31} & a_{32} & a_{33}
\end{pmatrix}
\]

\section{Tables}
Below is a sample table showing data:

\begin{table}[h]
\centering
\begin{tabular}{@{}lrrr@{}}
\toprule
\textbf{Category} & \textbf{Value 1} & \textbf{Value 2} & \textbf{Total} \\
\midrule
Item A & 10.5 & 20.3 & 30.8 \\
Item B & 15.2 & 25.1 & 40.3 \\
Item C & 22.1 & 12.7 & 34.8 \\
\bottomrule
\end{tabular}
\caption{Sample Data Table}
\label{tab:sample}
\end{table}

\section{Lists}
\subsection{Bullet Points}
\begin{itemize}
  \item First item
  \item Second item
  \item Third item with sub-items:
    \begin{itemize}
      \item Sub-item 1
      \item Sub-item 2
    \end{itemize}
\end{itemize}

\subsection{Numbered List}
\begin{enumerate}
  \item First step
  \item Second step
  \item Third step
\end{enumerate}

\section{Citations and References}
When you use the citation search feature, references will be automatically managed. For example, you might cite a paper like this: \cite{example2024}. All citations will be automatically added to your \texttt{cite.bib} file.

\section{Conclusion}
This template demonstrates basic \LaTeX{} features including mathematical formulas, tables, and lists. The bibliography system is now set up to work seamlessly with Bibby AI's citation management. Feel free to modify it for your needs!

\bibliographystyle{plain}
\bibliography{cite}

\end{document}\documentclass{article}
\usepackage[utf8]{inputenc}
\usepackage{booktabs}
\usepackage{amsmath}
\usepackage{natbib}

\title{My Document}
\author{Your Name}
\date{\today}

\begin{document}

\maketitle

\section{Introduction}
This is a sample document demonstrating various \LaTeX{} features including mathematical equations, tables, and formatting. You can cite papers using the citation search feature, and they will be automatically added to your \texttt{cite.bib} file.

\section{Mathematical Expressions}
Here are some example equations:

\subsection{Quadratic Formula}
The quadratic formula is given by:
\[x = \frac{-b \pm \sqrt{b^2 - 4ac}}{2a}\]

\subsection{Matrix Example}
A sample matrix:
\[
\begin{pmatrix}
a_{11} & a_{12} & a_{13} \\
a_{21} & a_{22} & a_{23} \\
a_{31} & a_{32} & a_{33}
\end{pmatrix}
\]

\section{Tables}
Below is a sample table showing data:

\begin{table}[h]
\centering
\begin{tabular}{@{}lrrr@{}}
\toprule
\textbf{Category} & \textbf{Value 1} & \textbf{Value 2} & \textbf{Total} \\
\midrule
Item A & 10.5 & 20.3 & 30.8 \\
Item B & 15.2 & 25.1 & 40.3 \\
Item C & 22.1 & 12.7 & 34.8 \\
\bottomrule
\end{tabular}
\caption{Sample Data Table}
\label{tab:sample}
\end{table}

\section{Lists}
\subsection{Bullet Points}
\begin{itemize}
  \item First item
  \item Second item
  \item Third item with sub-items:
    \begin{itemize}
      \item Sub-item 1
      \item Sub-item 2
    \end{itemize}
\end{itemize}

\subsection{Numbered List}
\begin{enumerate}
  \item First step
  \item Second step
  \item Third step
\end{enumerate}

\section{Citations and References}
When you use the citation search feature, references will be automatically managed. For example, you might cite a paper like this: \cite{example2024}. All citations will be automatically added to your \texttt{cite.bib} file.

\section{Conclusion}
This template demonstrates basic \LaTeX{} features including mathematical formulas, tables, and lists. The bibliography system is now set up to work seamlessly with Bibby AI's citation management. Feel free to modify it for your needs!

\bibliographystyle{plain}
\bibliography{cite}

\end{document}\documentclass{article}
\usepackage[utf8]{inputenc}
\usepackage{booktabs}
\usepackage{amsmath}
\usepackage{natbib}

\title{My Document}
\author{Your Name}
\date{\today}

\begin{document}

\maketitle

\section{Introduction}
This is a sample document demonstrating various \LaTeX{} features including mathematical equations, tables, and formatting. You can cite papers using the citation search feature, and they will be automatically added to your \texttt{cite.bib} file.

\section{Mathematical Expressions}
Here are some example equations:

\subsection{Quadratic Formula}
The quadratic formula is given by:
\[x = \frac{-b \pm \sqrt{b^2 - 4ac}}{2a}\]

\subsection{Matrix Example}
A sample matrix:
\[
\begin{pmatrix}
a_{11} & a_{12} & a_{13} \\
a_{21} & a_{22} & a_{23} \\
a_{31} & a_{32} & a_{33}
\end{pmatrix}
\]

\section{Tables}
Below is a sample table showing data:

\begin{table}[h]
\centering
\begin{tabular}{@{}lrrr@{}}
\toprule
\textbf{Category} & \textbf{Value 1} & \textbf{Value 2} & \textbf{Total} \\
\midrule
Item A & 10.5 & 20.3 & 30.8 \\
Item B & 15.2 & 25.1 & 40.3 \\
Item C & 22.1 & 12.7 & 34.8 \\
\bottomrule
\end{tabular}
\caption{Sample Data Table}
\label{tab:sample}
\end{table}

\section{Lists}
\subsection{Bullet Points}
\begin{itemize}
  \item First item
  \item Second item
  \item Third item with sub-items:
    \begin{itemize}
      \item Sub-item 1
      \item Sub-item 2
    \end{itemize}
\end{itemize}

\subsection{Numbered List}
\begin{enumerate}
  \item First step
  \item Second step
  \item Third step
\end{enumerate}

\section{Citations and References}
When you use the citation search feature, references will be automatically managed. For example, you might cite a paper like this: \cite{example2024}. All citations will be automatically added to your \texttt{cite.bib} file.

\section{Conclusion}
This template demonstrates basic \LaTeX{} features including mathematical formulas, tables, and lists. The bibliography system is now set up to work seamlessly with Bibby AI's citation management. Feel free to modify it for your needs!

\bibliographystyle{plain}
\bibliography{cite}

\end{document}\documentclass{article}
\usepackage[utf8]{inputenc}
\usepackage{booktabs}
\usepackage{amsmath}
\usepackage{natbib}

\title{My Document}
\author{Your Name}
\date{\today}

\begin{document}

\maketitle

\section{Introduction}
This is a sample document demonstrating various \LaTeX{} features including mathematical equations, tables, and formatting. You can cite papers using the citation search feature, and they will be automatically added to your \texttt{cite.bib} file.

\section{Mathematical Expressions}
Here are some example equations:

\subsection{Quadratic Formula}
The quadratic formula is given by:
\[x = \frac{-b \pm \sqrt{b^2 - 4ac}}{2a}\]

\subsection{Matrix Example}
A sample matrix:
\[
\begin{pmatrix}
a_{11} & a_{12} & a_{13} \\
a_{21} & a_{22} & a_{23} \\
a_{31} & a_{32} & a_{33}
\end{pmatrix}
\]

\section{Tables}
Below is a sample table showing data:

\begin{table}[h]
\centering
\begin{tabular}{@{}lrrr@{}}
\toprule
\textbf{Category} & \textbf{Value 1} & \textbf{Value 2} & \textbf{Total} \\
\midrule
Item A & 10.5 & 20.3 & 30.8 \\
Item B & 15.2 & 25.1 & 40.3 \\
Item C & 22.1 & 12.7 & 34.8 \\
\bottomrule
\end{tabular}
\caption{Sample Data Table}
\label{tab:sample}
\end{table}

\section{Lists}
\subsection{Bullet Points}
\begin{itemize}
  \item First item
  \item Second item
  \item Third item with sub-items:
    \begin{itemize}
      \item Sub-item 1
      \item Sub-item 2
    \end{itemize}
\end{itemize}

\subsection{Numbered List}
\begin{enumerate}
  \item First step
  \item Second step
  \item Third step
\end{enumerate}

\section{Citations and References}
When you use the citation search feature, references will be automatically managed. For example, you might cite a paper like this: \cite{example2024}. All citations will be automatically added to your \texttt{cite.bib} file.

\section{Conclusion}
This template demonstrates basic \LaTeX{} features including mathematical formulas, tables, and lists. The bibliography system is now set up to work seamlessly with Bibby AI's citation management. Feel free to modify it for your needs!

\bibliographystyle{plain}
\bibliography{cite}

\end{document}\documentclass{article}
\usepackage[utf8]{inputenc}
\usepackage{booktabs}
\usepackage{amsmath}
\usepackage{natbib}

\title{My Document}
\author{Your Name}
\date{\today}

\begin{document}

\maketitle

\section{Introduction}
This is a sample document demonstrating various \LaTeX{} features including mathematical equations, tables, and formatting. You can cite papers using the citation search feature, and they will be automatically added to your \texttt{cite.bib} file.

\section{Mathematical Expressions}
Here are some example equations:

\subsection{Quadratic Formula}
The quadratic formula is given by:
\[x = \frac{-b \pm \sqrt{b^2 - 4ac}}{2a}\]

\subsection{Matrix Example}
A sample matrix:
\[
\begin{pmatrix}
a_{11} & a_{12} & a_{13} \\
a_{21} & a_{22} & a_{23} \\
a_{31} & a_{32} & a_{33}
\end{pmatrix}
\]

\section{Tables}
Below is a sample table showing data:

\begin{table}[h]
\centering
\begin{tabular}{@{}lrrr@{}}
\toprule
\textbf{Category} & \textbf{Value 1} & \textbf{Value 2} & \textbf{Total} \\
\midrule
Item A & 10.5 & 20.3 & 30.8 \\
Item B & 15.2 & 25.1 & 40.3 \\
Item C & 22.1 & 12.7 & 34.8 \\
\bottomrule
\end{tabular}
\caption{Sample Data Table}
\label{tab:sample}
\end{table}

\section{Lists}
\subsection{Bullet Points}
\begin{itemize}
  \item First item
  \item Second item
  \item Third item with sub-items:
    \begin{itemize}
      \item Sub-item 1
      \item Sub-item 2
    \end{itemize}
\end{itemize}

\subsection{Numbered List}
\begin{enumerate}
  \item First step
  \item Second step
  \item Third step
\end{enumerate}

\section{Citations and References}
When you use the citation search feature, references will be automatically managed. For example, you might cite a paper like this: \cite{example2024}. All citations will be automatically added to your \texttt{cite.bib} file.

\section{Conclusion}
This template demonstrates basic \LaTeX{} features including mathematical formulas, tables, and lists. The bibliography system is now set up to work seamlessly with Bibby AI's citation management. Feel free to modify it for your needs!

\bibliographystyle{plain}
\bibliography{cite}

\end{document}\documentclass{article}
\usepackage[utf8]{inputenc}
\usepackage{booktabs}
\usepackage{amsmath}
\usepackage{natbib}

\title{My Document}
\author{Your Name}
\date{\today}

\begin{document}

\maketitle

\section{Introduction}
This is a sample document demonstrating various \LaTeX{} features including mathematical equations, tables, and formatting. You can cite papers using the citation search feature, and they will be automatically added to your \texttt{cite.bib} file.

\section{Mathematical Expressions}
Here are some example equations:

\subsection{Quadratic Formula}
The quadratic formula is given by:
\[x = \frac{-b \pm \sqrt{b^2 - 4ac}}{2a}\]

\subsection{Matrix Example}
A sample matrix:
\[
\begin{pmatrix}
a_{11} & a_{12} & a_{13} \\
a_{21} & a_{22} & a_{23} \\
a_{31} & a_{32} & a_{33}
\end{pmatrix}
\]

\section{Tables}
Below is a sample table showing data:

\begin{table}[h]
\centering
\begin{tabular}{@{}lrrr@{}}
\toprule
\textbf{Category} & \textbf{Value 1} & \textbf{Value 2} & \textbf{Total} \\
\midrule
Item A & 10.5 & 20.3 & 30.8 \\
Item B & 15.2 & 25.1 & 40.3 \\
Item C & 22.1 & 12.7 & 34.8 \\
\bottomrule
\end{tabular}
\caption{Sample Data Table}
\label{tab:sample}
\end{table}

\section{Lists}
\subsection{Bullet Points}
\begin{itemize}
  \item First item
  \item Second item
  \item Third item with sub-items:
    \begin{itemize}
      \item Sub-item 1
      \item Sub-item 2
    \end{itemize}
\end{itemize}

\subsection{Numbered List}
\begin{enumerate}
  \item First step
  \item Second step
  \item Third step
\end{enumerate}

\section{Citations and References}
When you use the citation search feature, references will be automatically managed. For example, you might cite a paper like this: \cite{example2024}. All citations will be automatically added to your \texttt{cite.bib} file.

\section{Conclusion}
This template demonstrates basic \LaTeX{} features including mathematical formulas, tables, and lists. The bibliography system is now set up to work seamlessly with Bibby AI's citation management. Feel free to modify it for your needs!

\bibliographystyle{plain}
\bibliography{cite}

\end{document}\documentclass{article}
\usepackage[utf8]{inputenc}
\usepackage{booktabs}
\usepackage{amsmath}
\usepackage{natbib}

\title{My Document}
\author{Your Name}
\date{\today}

\begin{document}

\maketitle

\section{Introduction}
This is a sample document demonstrating various \LaTeX{} features including mathematical equations, tables, and formatting. You can cite papers using the citation search feature, and they will be automatically added to your \texttt{cite.bib} file.

\section{Mathematical Expressions}
Here are some example equations:

\subsection{Quadratic Formula}
The quadratic formula is given by:
\[x = \frac{-b \pm \sqrt{b^2 - 4ac}}{2a}\]

\subsection{Matrix Example}
A sample matrix:
\[
\begin{pmatrix}
a_{11} & a_{12} & a_{13} \\
a_{21} & a_{22} & a_{23} \\
a_{31} & a_{32} & a_{33}
\end{pmatrix}
\]

\section{Tables}
Below is a sample table showing data:

\begin{table}[h]
\centering
\begin{tabular}{@{}lrrr@{}}
\toprule
\textbf{Category} & \textbf{Value 1} & \textbf{Value 2} & \textbf{Total} \\
\midrule
Item A & 10.5 & 20.3 & 30.8 \\
Item B & 15.2 & 25.1 & 40.3 \\
Item C & 22.1 & 12.7 & 34.8 \\
\bottomrule
\end{tabular}
\caption{Sample Data Table}
\label{tab:sample}
\end{table}

\section{Lists}
\subsection{Bullet Points}
\begin{itemize}
  \item First item
  \item Second item
  \item Third item with sub-items:
    \begin{itemize}
      \item Sub-item 1
      \item Sub-item 2
    \end{itemize}
\end{itemize}

\subsection{Numbered List}
\begin{enumerate}
  \item First step
  \item Second step
  \item Third step
\end{enumerate}

\section{Citations and References}
When you use the citation search feature, references will be automatically managed. For example, you might cite a paper like this: \cite{example2024}. All citations will be automatically added to your \texttt{cite.bib} file.

\section{Conclusion}
This template demonstrates basic \LaTeX{} features including mathematical formulas, tables, and lists. The bibliography system is now set up to work seamlessly with Bibby AI's citation management. Feel free to modify it for your needs!

\bibliographystyle{plain}
\bibliography{cite}

\end{document}\documentclass{article}
\usepackage[utf8]{inputenc}
\usepackage{booktabs}
\usepackage{amsmath}
\usepackage{natbib}

\title{My Document}
\author{Your Name}
\date{\today}

\begin{document}

\maketitle

\section{Introduction}
This is a sample document demonstrating various \LaTeX{} features including mathematical equations, tables, and formatting. You can cite papers using the citation search feature, and they will be automatically added to your \texttt{cite.bib} file.

\section{Mathematical Expressions}
Here are some example equations:

\subsection{Quadratic Formula}
The quadratic formula is given by:
\[x = \frac{-b \pm \sqrt{b^2 - 4ac}}{2a}\]

\subsection{Matrix Example}
A sample matrix:
\[
\begin{pmatrix}
a_{11} & a_{12} & a_{13} \\
a_{21} & a_{22} & a_{23} \\
a_{31} & a_{32} & a_{33}
\end{pmatrix}
\]

\section{Tables}
Below is a sample table showing data:

\begin{table}[h]
\centering
\begin{tabular}{@{}lrrr@{}}
\toprule
\textbf{Category} & \textbf{Value 1} & \textbf{Value 2} & \textbf{Total} \\
\midrule
Item A & 10.5 & 20.3 & 30.8 \\
Item B & 15.2 & 25.1 & 40.3 \\
Item C & 22.1 & 12.7 & 34.8 \\
\bottomrule
\end{tabular}
\caption{Sample Data Table}
\label{tab:sample}
\end{table}

\section{Lists}
\subsection{Bullet Points}
\begin{itemize}
  \item First item
  \item Second item
  \item Third item with sub-items:
    \begin{itemize}
      \item Sub-item 1
      \item Sub-item 2
    \end{itemize}
\end{itemize}

\subsection{Numbered List}
\begin{enumerate}
  \item First step
  \item Second step
  \item Third step
\end{enumerate}

\section{Citations and References}
When you use the citation search feature, references will be automatically managed. For example, you might cite a paper like this: \cite{example2024}. All citations will be automatically added to your \texttt{cite.bib} file.

\section{Conclusion}
This template demonstrates basic \LaTeX{} features including mathematical formulas, tables, and lists. The bibliography system is now set up to work seamlessly with Bibby AI's citation management. Feel free to modify it for your needs!

\bibliographystyle{plain}
\bibliography{cite}

\end{document}\documentclass{article}
\usepackage[utf8]{inputenc}
\usepackage{booktabs}
\usepackage{amsmath}
\usepackage{natbib}

\title{My Document}
\author{Your Name}
\date{\today}

\begin{document}

\maketitle

\section{Introduction}
This is a sample document demonstrating various \LaTeX{} features including mathematical equations, tables, and formatting. You can cite papers using the citation search feature, and they will be automatically added to your \texttt{cite.bib} file.

\section{Mathematical Expressions}
Here are some example equations:

\subsection{Quadratic Formula}
The quadratic formula is given by:
\[x = \frac{-b \pm \sqrt{b^2 - 4ac}}{2a}\]

\subsection{Matrix Example}
A sample matrix:
\[
\begin{pmatrix}
a_{11} & a_{12} & a_{13} \\
a_{21} & a_{22} & a_{23} \\
a_{31} & a_{32} & a_{33}
\end{pmatrix}
\]

\section{Tables}
Below is a sample table showing data:

\begin{table}[h]
\centering
\begin{tabular}{@{}lrrr@{}}
\toprule
\textbf{Category} & \textbf{Value 1} & \textbf{Value 2} & \textbf{Total} \\
\midrule
Item A & 10.5 & 20.3 & 30.8 \\
Item B & 15.2 & 25.1 & 40.3 \\
Item C & 22.1 & 12.7 & 34.8 \\
\bottomrule
\end{tabular}
\caption{Sample Data Table}
\label{tab:sample}
\end{table}

\section{Lists}
\subsection{Bullet Points}
\begin{itemize}
  \item First item
  \item Second item
  \item Third item with sub-items:
    \begin{itemize}
      \item Sub-item 1
      \item Sub-item 2
    \end{itemize}
\end{itemize}

\subsection{Numbered List}
\begin{enumerate}
  \item First step
  \item Second step
  \item Third step
\end{enumerate}

\section{Citations and References}
When you use the citation search feature, references will be automatically managed. For example, you might cite a paper like this: \cite{example2024}. All citations will be automatically added to your \texttt{cite.bib} file.

\section{Conclusion}
This template demonstrates basic \LaTeX{} features including mathematical formulas, tables, and lists. The bibliography system is now set up to work seamlessly with Bibby AI's citation management. Feel free to modify it for your needs!

\bibliographystyle{plain}
\bibliography{cite}

\end{document}\documentclass{article}
\usepackage[utf8]{inputenc}
\usepackage{booktabs}
\usepackage{amsmath}
\usepackage{natbib}

\title{My Document}
\author{Your Name}
\date{\today}

\begin{document}

\maketitle

\section{Introduction}
This is a sample document demonstrating various \LaTeX{} features including mathematical equations, tables, and formatting. You can cite papers using the citation search feature, and they will be automatically added to your \texttt{cite.bib} file.

\section{Mathematical Expressions}
Here are some example equations:

\subsection{Quadratic Formula}
The quadratic formula is given by:
\[x = \frac{-b \pm \sqrt{b^2 - 4ac}}{2a}\]

\subsection{Matrix Example}
A sample matrix:
\[
\begin{pmatrix}
a_{11} & a_{12} & a_{13} \\
a_{21} & a_{22} & a_{23} \\
a_{31} & a_{32} & a_{33}
\end{pmatrix}
\]

\section{Tables}
Below is a sample table showing data:

\begin{table}[h]
\centering
\begin{tabular}{@{}lrrr@{}}
\toprule
\textbf{Category} & \textbf{Value 1} & \textbf{Value 2} & \textbf{Total} \\
\midrule
Item A & 10.5 & 20.3 & 30.8 \\
Item B & 15.2 & 25.1 & 40.3 \\
Item C & 22.1 & 12.7 & 34.8 \\
\bottomrule
\end{tabular}
\caption{Sample Data Table}
\label{tab:sample}
\end{table}

\section{Lists}
\subsection{Bullet Points}
\begin{itemize}
  \item First item
  \item Second item
  \item Third item with sub-items:
    \begin{itemize}
      \item Sub-item 1
      \item Sub-item 2
    \end{itemize}
\end{itemize}

\subsection{Numbered List}
\begin{enumerate}
  \item First step
  \item Second step
  \item Third step
\end{enumerate}

\section{Citations and References}
When you use the citation search feature, references will be automatically managed. For example, you might cite a paper like this: \cite{example2024}. All citations will be automatically added to your \texttt{cite.bib} file.

\section{Conclusion}
This template demonstrates basic \LaTeX{} features including mathematical formulas, tables, and lists. The bibliography system is now set up to work seamlessly with Bibby AI's citation management. Feel free to modify it for your needs!

\bibliographystyle{plain}
\bibliography{cite}

\end{document}\documentclass{article}
\usepackage[utf8]{inputenc}
\usepackage{booktabs}
\usepackage{amsmath}
\usepackage{natbib}

\title{My Document}
\author{Your Name}
\date{\today}

\begin{document}

\maketitle

\section{Introduction}
This is a sample document demonstrating various \LaTeX{} features including mathematical equations, tables, and formatting. You can cite papers using the citation search feature, and they will be automatically added to your \texttt{cite.bib} file.

\section{Mathematical Expressions}
Here are some example equations:

\subsection{Quadratic Formula}
The quadratic formula is given by:
\[x = \frac{-b \pm \sqrt{b^2 - 4ac}}{2a}\]

\subsection{Matrix Example}
A sample matrix:
\[
\begin{pmatrix}
a_{11} & a_{12} & a_{13} \\
a_{21} & a_{22} & a_{23} \\
a_{31} & a_{32} & a_{33}
\end{pmatrix}
\]

\section{Tables}
Below is a sample table showing data:

\begin{table}[h]
\centering
\begin{tabular}{@{}lrrr@{}}
\toprule
\textbf{Category} & \textbf{Value 1} & \textbf{Value 2} & \textbf{Total} \\
\midrule
Item A & 10.5 & 20.3 & 30.8 \\
Item B & 15.2 & 25.1 & 40.3 \\
Item C & 22.1 & 12.7 & 34.8 \\
\bottomrule
\end{tabular}
\caption{Sample Data Table}
\label{tab:sample}
\end{table}

\section{Lists}
\subsection{Bullet Points}
\begin{itemize}
  \item First item
  \item Second item
  \item Third item with sub-items:
    \begin{itemize}
      \item Sub-item 1
      \item Sub-item 2
    \end{itemize}
\end{itemize}

\subsection{Numbered List}
\begin{enumerate}
  \item First step
  \item Second step
  \item Third step
\end{enumerate}

\section{Citations and References}
When you use the citation search feature, references will be automatically managed. For example, you might cite a paper like this: \cite{example2024}. All citations will be automatically added to your \texttt{cite.bib} file.

\section{Conclusion}
This template demonstrates basic \LaTeX{} features including mathematical formulas, tables, and lists. The bibliography system is now set up to work seamlessly with Bibby AI's citation management. Feel free to modify it for your needs!

\bibliographystyle{plain}
\bibliography{cite}

\end{document}\documentclass{article}
\usepackage[utf8]{inputenc}
\usepackage{booktabs}
\usepackage{amsmath}
\usepackage{natbib}

\title{My Document}
\author{Your Name}
\date{\today}

\begin{document}

\maketitle

\section{Introduction}
This is a sample document demonstrating various \LaTeX{} features including mathematical equations, tables, and formatting. You can cite papers using the citation search feature, and they will be automatically added to your \texttt{cite.bib} file.

\section{Mathematical Expressions}
Here are some example equations:

\subsection{Quadratic Formula}
The quadratic formula is given by:
\[x = \frac{-b \pm \sqrt{b^2 - 4ac}}{2a}\]

\subsection{Matrix Example}
A sample matrix:
\[
\begin{pmatrix}
a_{11} & a_{12} & a_{13} \\
a_{21} & a_{22} & a_{23} \\
a_{31} & a_{32} & a_{33}
\end{pmatrix}
\]

\section{Tables}
Below is a sample table showing data:

\begin{table}[h]
\centering
\begin{tabular}{@{}lrrr@{}}
\toprule
\textbf{Category} & \textbf{Value 1} & \textbf{Value 2} & \textbf{Total} \\
\midrule
Item A & 10.5 & 20.3 & 30.8 \\
Item B & 15.2 & 25.1 & 40.3 \\
Item C & 22.1 & 12.7 & 34.8 \\
\bottomrule
\end{tabular}
\caption{Sample Data Table}
\label{tab:sample}
\end{table}

\section{Lists}
\subsection{Bullet Points}
\begin{itemize}
  \item First item
  \item Second item
  \item Third item with sub-items:
    \begin{itemize}
      \item Sub-item 1
      \item Sub-item 2
    \end{itemize}
\end{itemize}

\subsection{Numbered List}
\begin{enumerate}
  \item First step
  \item Second step
  \item Third step
\end{enumerate}

\section{Citations and References}
When you use the citation search feature, references will be automatically managed. For example, you might cite a paper like this: \cite{example2024}. All citations will be automatically added to your \texttt{cite.bib} file.

\section{Conclusion}
This template demonstrates basic \LaTeX{} features including mathematical formulas, tables, and lists. The bibliography system is now set up to work seamlessly with Bibby AI's citation management. Feel free to modify it for your needs!

\bibliographystyle{plain}
\bibliography{cite}

\end{document}\documentclass{article}
\usepackage[utf8]{inputenc}
\usepackage{booktabs}
\usepackage{amsmath}
\usepackage{natbib}

\title{My Document}
\author{Your Name}
\date{\today}

\begin{document}

\maketitle

\section{Introduction}
This is a sample document demonstrating various \LaTeX{} features including mathematical equations, tables, and formatting. You can cite papers using the citation search feature, and they will be automatically added to your \texttt{cite.bib} file.

\section{Mathematical Expressions}
Here are some example equations:

\subsection{Quadratic Formula}
The quadratic formula is given by:
\[x = \frac{-b \pm \sqrt{b^2 - 4ac}}{2a}\]

\subsection{Matrix Example}
A sample matrix:
\[
\begin{pmatrix}
a_{11} & a_{12} & a_{13} \\
a_{21} & a_{22} & a_{23} \\
a_{31} & a_{32} & a_{33}
\end{pmatrix}
\]

\section{Tables}
Below is a sample table showing data:

\begin{table}[h]
\centering
\begin{tabular}{@{}lrrr@{}}
\toprule
\textbf{Category} & \textbf{Value 1} & \textbf{Value 2} & \textbf{Total} \\
\midrule
Item A & 10.5 & 20.3 & 30.8 \\
Item B & 15.2 & 25.1 & 40.3 \\
Item C & 22.1 & 12.7 & 34.8 \\
\bottomrule
\end{tabular}
\caption{Sample Data Table}
\label{tab:sample}
\end{table}

\section{Lists}
\subsection{Bullet Points}
\begin{itemize}
  \item First item
  \item Second item
  \item Third item with sub-items:
    \begin{itemize}
      \item Sub-item 1
      \item Sub-item 2
    \end{itemize}
\end{itemize}

\subsection{Numbered List}
\begin{enumerate}
  \item First step
  \item Second step
  \item Third step
\end{enumerate}

\section{Citations and References}
When you use the citation search feature, references will be automatically managed. For example, you might cite a paper like this: \cite{example2024}. All citations will be automatically added to your \texttt{cite.bib} file.

\section{Conclusion}
This template demonstrates basic \LaTeX{} features including mathematical formulas, tables, and lists. The bibliography system is now set up to work seamlessly with Bibby AI's citation management. Feel free to modify it for your needs!

\bibliographystyle{plain}
\bibliography{cite}

\end{document}\documentclass{article}
\usepackage[utf8]{inputenc}
\usepackage{booktabs}
\usepackage{amsmath}
\usepackage{natbib}

\title{My Document}
\author{Your Name}
\date{\today}

\begin{document}

\maketitle

\section{Introduction}
This is a sample document demonstrating various \LaTeX{} features including mathematical equations, tables, and formatting. You can cite papers using the citation search feature, and they will be automatically added to your \texttt{cite.bib} file.

\section{Mathematical Expressions}
Here are some example equations:

\subsection{Quadratic Formula}
The quadratic formula is given by:
\[x = \frac{-b \pm \sqrt{b^2 - 4ac}}{2a}\]

\subsection{Matrix Example}
A sample matrix:
\[
\begin{pmatrix}
a_{11} & a_{12} & a_{13} \\
a_{21} & a_{22} & a_{23} \\
a_{31} & a_{32} & a_{33}
\end{pmatrix}
\]

\section{Tables}
Below is a sample table showing data:

\begin{table}[h]
\centering
\begin{tabular}{@{}lrrr@{}}
\toprule
\textbf{Category} & \textbf{Value 1} & \textbf{Value 2} & \textbf{Total} \\
\midrule
Item A & 10.5 & 20.3 & 30.8 \\
Item B & 15.2 & 25.1 & 40.3 \\
Item C & 22.1 & 12.7 & 34.8 \\
\bottomrule
\end{tabular}
\caption{Sample Data Table}
\label{tab:sample}
\end{table}

\section{Lists}
\subsection{Bullet Points}
\begin{itemize}
  \item First item
  \item Second item
  \item Third item with sub-items:
    \begin{itemize}
      \item Sub-item 1
      \item Sub-item 2
    \end{itemize}
\end{itemize}

\subsection{Numbered List}
\begin{enumerate}
  \item First step
  \item Second step
  \item Third step
\end{enumerate}

\section{Citations and References}
When you use the citation search feature, references will be automatically managed. For example, you might cite a paper like this: \cite{example2024}. All citations will be automatically added to your \texttt{cite.bib} file.

\section{Conclusion}
This template demonstrates basic \LaTeX{} features including mathematical formulas, tables, and lists. The bibliography system is now set up to work seamlessly with Bibby AI's citation management. Feel free to modify it for your needs!

\bibliographystyle{plain}
\bibliography{cite}

\end{document}\documentclass{article}
\usepackage[utf8]{inputenc}
\usepackage{booktabs}
\usepackage{amsmath}
\usepackage{natbib}

\title{My Document}
\author{Your Name}
\date{\today}

\begin{document}

\maketitle

\section{Introduction}
This is a sample document demonstrating various \LaTeX{} features including mathematical equations, tables, and formatting. You can cite papers using the citation search feature, and they will be automatically added to your \texttt{cite.bib} file.

\section{Mathematical Expressions}
Here are some example equations:

\subsection{Quadratic Formula}
The quadratic formula is given by:
\[x = \frac{-b \pm \sqrt{b^2 - 4ac}}{2a}\]

\subsection{Matrix Example}
A sample matrix:
\[
\begin{pmatrix}
a_{11} & a_{12} & a_{13} \\
a_{21} & a_{22} & a_{23} \\
a_{31} & a_{32} & a_{33}
\end{pmatrix}
\]

\section{Tables}
Below is a sample table showing data:

\begin{table}[h]
\centering
\begin{tabular}{@{}lrrr@{}}
\toprule
\textbf{Category} & \textbf{Value 1} & \textbf{Value 2} & \textbf{Total} \\
\midrule
Item A & 10.5 & 20.3 & 30.8 \\
Item B & 15.2 & 25.1 & 40.3 \\
Item C & 22.1 & 12.7 & 34.8 \\
\bottomrule
\end{tabular}
\caption{Sample Data Table}
\label{tab:sample}
\end{table}

\section{Lists}
\subsection{Bullet Points}
\begin{itemize}
  \item First item
  \item Second item
  \item Third item with sub-items:
    \begin{itemize}
      \item Sub-item 1
      \item Sub-item 2
    \end{itemize}
\end{itemize}

\subsection{Numbered List}
\begin{enumerate}
  \item First step
  \item Second step
  \item Third step
\end{enumerate}

\section{Citations and References}
When you use the citation search feature, references will be automatically managed. For example, you might cite a paper like this: \cite{example2024}. All citations will be automatically added to your \texttt{cite.bib} file.

\section{Conclusion}
This template demonstrates basic \LaTeX{} features including mathematical formulas, tables, and lists. The bibliography system is now set up to work seamlessly with Bibby AI's citation management. Feel free to modify it for your needs!

\bibliographystyle{plain}
\bibliography{cite}

\end{document}\documentclass{article}
\usepackage[utf8]{inputenc}
\usepackage{booktabs}
\usepackage{amsmath}
\usepackage{natbib}

\title{My Document}
\author{Your Name}
\date{\today}

\begin{document}

\maketitle

\section{Introduction}
This is a sample document demonstrating various \LaTeX{} features including mathematical equations, tables, and formatting. You can cite papers using the citation search feature, and they will be automatically added to your \texttt{cite.bib} file.

\section{Mathematical Expressions}
Here are some example equations:

\subsection{Quadratic Formula}
The quadratic formula is given by:
\[x = \frac{-b \pm \sqrt{b^2 - 4ac}}{2a}\]

\subsection{Matrix Example}
A sample matrix:
\[
\begin{pmatrix}
a_{11} & a_{12} & a_{13} \\
a_{21} & a_{22} & a_{23} \\
a_{31} & a_{32} & a_{33}
\end{pmatrix}
\]

\section{Tables}
Below is a sample table showing data:

\begin{table}[h]
\centering
\begin{tabular}{@{}lrrr@{}}
\toprule
\textbf{Category} & \textbf{Value 1} & \textbf{Value 2} & \textbf{Total} \\
\midrule
Item A & 10.5 & 20.3 & 30.8 \\
Item B & 15.2 & 25.1 & 40.3 \\
Item C & 22.1 & 12.7 & 34.8 \\
\bottomrule
\end{tabular}
\caption{Sample Data Table}
\label{tab:sample}
\end{table}

\section{Lists}
\subsection{Bullet Points}
\begin{itemize}
  \item First item
  \item Second item
  \item Third item with sub-items:
    \begin{itemize}
      \item Sub-item 1
      \item Sub-item 2
    \end{itemize}
\end{itemize}

\subsection{Numbered List}
\begin{enumerate}
  \item First step
  \item Second step
  \item Third step
\end{enumerate}

\section{Citations and References}
When you use the citation search feature, references will be automatically managed. For example, you might cite a paper like this: \cite{example2024}. All citations will be automatically added to your \texttt{cite.bib} file.

\section{Conclusion}
This template demonstrates basic \LaTeX{} features including mathematical formulas, tables, and lists. The bibliography system is now set up to work seamlessly with Bibby AI's citation management. Feel free to modify it for your needs!

\bibliographystyle{plain}
\bibliography{cite}

\end{document}\documentclass{article}
\usepackage[utf8]{inputenc}
\usepackage{booktabs}
\usepackage{amsmath}
\usepackage{natbib}

\title{My Document}
\author{Your Name}
\date{\today}

\begin{document}

\maketitle

\section{Introduction}
This is a sample document demonstrating various \LaTeX{} features including mathematical equations, tables, and formatting. You can cite papers using the citation search feature, and they will be automatically added to your \texttt{cite.bib} file.

\section{Mathematical Expressions}
Here are some example equations:

\subsection{Quadratic Formula}
The quadratic formula is given by:
\[x = \frac{-b \pm \sqrt{b^2 - 4ac}}{2a}\]

\subsection{Matrix Example}
A sample matrix:
\[
\begin{pmatrix}
a_{11} & a_{12} & a_{13} \\
a_{21} & a_{22} & a_{23} \\
a_{31} & a_{32} & a_{33}
\end{pmatrix}
\]

\section{Tables}
Below is a sample table showing data:

\begin{table}[h]
\centering
\begin{tabular}{@{}lrrr@{}}
\toprule
\textbf{Category} & \textbf{Value 1} & \textbf{Value 2} & \textbf{Total} \\
\midrule
Item A & 10.5 & 20.3 & 30.8 \\
Item B & 15.2 & 25.1 & 40.3 \\
Item C & 22.1 & 12.7 & 34.8 \\
\bottomrule
\end{tabular}
\caption{Sample Data Table}
\label{tab:sample}
\end{table}

\section{Lists}
\subsection{Bullet Points}
\begin{itemize}
  \item First item
  \item Second item
  \item Third item with sub-items:
    \begin{itemize}
      \item Sub-item 1
      \item Sub-item 2
    \end{itemize}
\end{itemize}

\subsection{Numbered List}
\begin{enumerate}
  \item First step
  \item Second step
  \item Third step
\end{enumerate}

\section{Citations and References}
When you use the citation search feature, references will be automatically managed. For example, you might cite a paper like this: \cite{example2024}. All citations will be automatically added to your \texttt{cite.bib} file.

\section{Conclusion}
This template demonstrates basic \LaTeX{} features including mathematical formulas, tables, and lists. The bibliography system is now set up to work seamlessly with Bibby AI's citation management. Feel free to modify it for your needs!

\bibliographystyle{plain}
\bibliography{cite}

\end{document}\documentclass{article}
\usepackage[utf8]{inputenc}
\usepackage{booktabs}
\usepackage{amsmath}
\usepackage{natbib}

\title{My Document}
\author{Your Name}
\date{\today}

\begin{document}

\maketitle

\section{Introduction}
This is a sample document demonstrating various \LaTeX{} features including mathematical equations, tables, and formatting. You can cite papers using the citation search feature, and they will be automatically added to your \texttt{cite.bib} file.

\section{Mathematical Expressions}
Here are some example equations:

\subsection{Quadratic Formula}
The quadratic formula is given by:
\[x = \frac{-b \pm \sqrt{b^2 - 4ac}}{2a}\]

\subsection{Matrix Example}
A sample matrix:
\[
\begin{pmatrix}
a_{11} & a_{12} & a_{13} \\
a_{21} & a_{22} & a_{23} \\
a_{31} & a_{32} & a_{33}
\end{pmatrix}
\]

\section{Tables}
Below is a sample table showing data:

\begin{table}[h]
\centering
\begin{tabular}{@{}lrrr@{}}
\toprule
\textbf{Category} & \textbf{Value 1} & \textbf{Value 2} & \textbf{Total} \\
\midrule
Item A & 10.5 & 20.3 & 30.8 \\
Item B & 15.2 & 25.1 & 40.3 \\
Item C & 22.1 & 12.7 & 34.8 \\
\bottomrule
\end{tabular}
\caption{Sample Data Table}
\label{tab:sample}
\end{table}

\section{Lists}
\subsection{Bullet Points}
\begin{itemize}
  \item First item
  \item Second item
  \item Third item with sub-items:
    \begin{itemize}
      \item Sub-item 1
      \item Sub-item 2
    \end{itemize}
\end{itemize}

\subsection{Numbered List}
\begin{enumerate}
  \item First step
  \item Second step
  \item Third step
\end{enumerate}

\section{Citations and References}
When you use the citation search feature, references will be automatically managed. For example, you might cite a paper like this: \cite{example2024}. All citations will be automatically added to your \texttt{cite.bib} file.

\section{Conclusion}
This template demonstrates basic \LaTeX{} features including mathematical formulas, tables, and lists. The bibliography system is now set up to work seamlessly with Bibby AI's citation management. Feel free to modify it for your needs!

\bibliographystyle{plain}
\bibliography{cite}

\end{document}\documentclass{article}
\usepackage[utf8]{inputenc}
\usepackage{booktabs}
\usepackage{amsmath}
\usepackage{natbib}

\title{My Document}
\author{Your Name}
\date{\today}

\begin{document}

\maketitle

\section{Introduction}
This is a sample document demonstrating various \LaTeX{} features including mathematical equations, tables, and formatting. You can cite papers using the citation search feature, and they will be automatically added to your \texttt{cite.bib} file.

\section{Mathematical Expressions}
Here are some example equations:

\subsection{Quadratic Formula}
The quadratic formula is given by:
\[x = \frac{-b \pm \sqrt{b^2 - 4ac}}{2a}\]

\subsection{Matrix Example}
A sample matrix:
\[
\begin{pmatrix}
a_{11} & a_{12} & a_{13} \\
a_{21} & a_{22} & a_{23} \\
a_{31} & a_{32} & a_{33}
\end{pmatrix}
\]

\section{Tables}
Below is a sample table showing data:

\begin{table}[h]
\centering
\begin{tabular}{@{}lrrr@{}}
\toprule
\textbf{Category} & \textbf{Value 1} & \textbf{Value 2} & \textbf{Total} \\
\midrule
Item A & 10.5 & 20.3 & 30.8 \\
Item B & 15.2 & 25.1 & 40.3 \\
Item C & 22.1 & 12.7 & 34.8 \\
\bottomrule
\end{tabular}
\caption{Sample Data Table}
\label{tab:sample}
\end{table}

\section{Lists}
\subsection{Bullet Points}
\begin{itemize}
  \item First item
  \item Second item
  \item Third item with sub-items:
    \begin{itemize}
      \item Sub-item 1
      \item Sub-item 2
    \end{itemize}
\end{itemize}

\subsection{Numbered List}
\begin{enumerate}
  \item First step
  \item Second step
  \item Third step
\end{enumerate}

\section{Citations and References}
When you use the citation search feature, references will be automatically managed. For example, you might cite a paper like this: \cite{example2024}. All citations will be automatically added to your \texttt{cite.bib} file.

\section{Conclusion}
This template demonstrates basic \LaTeX{} features including mathematical formulas, tables, and lists. The bibliography system is now set up to work seamlessly with Bibby AI's citation management. Feel free to modify it for your needs!

\bibliographystyle{plain}
\bibliography{cite}

\end{document}\documentclass{article}
\usepackage[utf8]{inputenc}
\usepackage{booktabs}
\usepackage{amsmath}
\usepackage{natbib}

\title{My Document}
\author{Your Name}
\date{\today}

\begin{document}

\maketitle

\section{Introduction}
This is a sample document demonstrating various \LaTeX{} features including mathematical equations, tables, and formatting. You can cite papers using the citation search feature, and they will be automatically added to your \texttt{cite.bib} file.

\section{Mathematical Expressions}
Here are some example equations:

\subsection{Quadratic Formula}
The quadratic formula is given by:
\[x = \frac{-b \pm \sqrt{b^2 - 4ac}}{2a}\]

\subsection{Matrix Example}
A sample matrix:
\[
\begin{pmatrix}
a_{11} & a_{12} & a_{13} \\
a_{21} & a_{22} & a_{23} \\
a_{31} & a_{32} & a_{33}
\end{pmatrix}
\]

\section{Tables}
Below is a sample table showing data:

\begin{table}[h]
\centering
\begin{tabular}{@{}lrrr@{}}
\toprule
\textbf{Category} & \textbf{Value 1} & \textbf{Value 2} & \textbf{Total} \\
\midrule
Item A & 10.5 & 20.3 & 30.8 \\
Item B & 15.2 & 25.1 & 40.3 \\
Item C & 22.1 & 12.7 & 34.8 \\
\bottomrule
\end{tabular}
\caption{Sample Data Table}
\label{tab:sample}
\end{table}

\section{Lists}
\subsection{Bullet Points}
\begin{itemize}
  \item First item
  \item Second item
  \item Third item with sub-items:
    \begin{itemize}
      \item Sub-item 1
      \item Sub-item 2
    \end{itemize}
\end{itemize}

\subsection{Numbered List}
\begin{enumerate}
  \item First step
  \item Second step
  \item Third step
\end{enumerate}

\section{Citations and References}
When you use the citation search feature, references will be automatically managed. For example, you might cite a paper like this: \cite{example2024}. All citations will be automatically added to your \texttt{cite.bib} file.

\section{Conclusion}
This template demonstrates basic \LaTeX{} features including mathematical formulas, tables, and lists. The bibliography system is now set up to work seamlessly with Bibby AI's citation management. Feel free to modify it for your needs!

\bibliographystyle{plain}
\bibliography{cite}

\end{document}\documentclass{article}
\usepackage[utf8]{inputenc}
\usepackage{booktabs}
\usepackage{amsmath}
\usepackage{natbib}

\title{My Document}
\author{Your Name}
\date{\today}

\begin{document}

\maketitle

\section{Introduction}
This is a sample document demonstrating various \LaTeX{} features including mathematical equations, tables, and formatting. You can cite papers using the citation search feature, and they will be automatically added to your \texttt{cite.bib} file.

\section{Mathematical Expressions}
Here are some example equations:

\subsection{Quadratic Formula}
The quadratic formula is given by:
\[x = \frac{-b \pm \sqrt{b^2 - 4ac}}{2a}\]

\subsection{Matrix Example}
A sample matrix:
\[
\begin{pmatrix}
a_{11} & a_{12} & a_{13} \\
a_{21} & a_{22} & a_{23} \\
a_{31} & a_{32} & a_{33}
\end{pmatrix}
\]

\section{Tables}
Below is a sample table showing data:

\begin{table}[h]
\centering
\begin{tabular}{@{}lrrr@{}}
\toprule
\textbf{Category} & \textbf{Value 1} & \textbf{Value 2} & \textbf{Total} \\
\midrule
Item A & 10.5 & 20.3 & 30.8 \\
Item B & 15.2 & 25.1 & 40.3 \\
Item C & 22.1 & 12.7 & 34.8 \\
\bottomrule
\end{tabular}
\caption{Sample Data Table}
\label{tab:sample}
\end{table}

\section{Lists}
\subsection{Bullet Points}
\begin{itemize}
  \item First item
  \item Second item
  \item Third item with sub-items:
    \begin{itemize}
      \item Sub-item 1
      \item Sub-item 2
    \end{itemize}
\end{itemize}

\subsection{Numbered List}
\begin{enumerate}
  \item First step
  \item Second step
  \item Third step
\end{enumerate}

\section{Citations and References}
When you use the citation search feature, references will be automatically managed. For example, you might cite a paper like this: \cite{example2024}. All citations will be automatically added to your \texttt{cite.bib} file.

\section{Conclusion}
This template demonstrates basic \LaTeX{} features including mathematical formulas, tables, and lists. The bibliography system is now set up to work seamlessly with Bibby AI's citation management. Feel free to modify it for your needs!

\bibliographystyle{plain}
\bibliography{cite}

\end{document}\documentclass{article}
\usepackage[utf8]{inputenc}
\usepackage{booktabs}
\usepackage{amsmath}
\usepackage{natbib}

\title{My Document}
\author{Your Name}
\date{\today}

\begin{document}

\maketitle

\section{Introduction}
This is a sample document demonstrating various \LaTeX{} features including mathematical equations, tables, and formatting. You can cite papers using the citation search feature, and they will be automatically added to your \texttt{cite.bib} file.

\section{Mathematical Expressions}
Here are some example equations:

\subsection{Quadratic Formula}
The quadratic formula is given by:
\[x = \frac{-b \pm \sqrt{b^2 - 4ac}}{2a}\]

\subsection{Matrix Example}
A sample matrix:
\[
\begin{pmatrix}
a_{11} & a_{12} & a_{13} \\
a_{21} & a_{22} & a_{23} \\
a_{31} & a_{32} & a_{33}
\end{pmatrix}
\]

\section{Tables}
Below is a sample table showing data:

\begin{table}[h]
\centering
\begin{tabular}{@{}lrrr@{}}
\toprule
\textbf{Category} & \textbf{Value 1} & \textbf{Value 2} & \textbf{Total} \\
\midrule
Item A & 10.5 & 20.3 & 30.8 \\
Item B & 15.2 & 25.1 & 40.3 \\
Item C & 22.1 & 12.7 & 34.8 \\
\bottomrule
\end{tabular}
\caption{Sample Data Table}
\label{tab:sample}
\end{table}

\section{Lists}
\subsection{Bullet Points}
\begin{itemize}
  \item First item
  \item Second item
  \item Third item with sub-items:
    \begin{itemize}
      \item Sub-item 1
      \item Sub-item 2
    \end{itemize}
\end{itemize}

\subsection{Numbered List}
\begin{enumerate}
  \item First step
  \item Second step
  \item Third step
\end{enumerate}

\section{Citations and References}
When you use the citation search feature, references will be automatically managed. For example, you might cite a paper like this: \cite{example2024}. All citations will be automatically added to your \texttt{cite.bib} file.

\section{Conclusion}
This template demonstrates basic \LaTeX{} features including mathematical formulas, tables, and lists. The bibliography system is now set up to work seamlessly with Bibby AI's citation management. Feel free to modify it for your needs!

\bibliographystyle{plain}
\bibliography{cite}

\end{document}\documentclass{article}
\usepackage[utf8]{inputenc}
\usepackage{booktabs}
\usepackage{amsmath}
\usepackage{natbib}

\title{My Document}
\author{Your Name}
\date{\today}

\begin{document}

\maketitle

\section{Introduction}
This is a sample document demonstrating various \LaTeX{} features including mathematical equations, tables, and formatting. You can cite papers using the citation search feature, and they will be automatically added to your \texttt{cite.bib} file.

\section{Mathematical Expressions}
Here are some example equations:

\subsection{Quadratic Formula}
The quadratic formula is given by:
\[x = \frac{-b \pm \sqrt{b^2 - 4ac}}{2a}\]

\subsection{Matrix Example}
A sample matrix:
\[
\begin{pmatrix}
a_{11} & a_{12} & a_{13} \\
a_{21} & a_{22} & a_{23} \\
a_{31} & a_{32} & a_{33}
\end{pmatrix}
\]

\section{Tables}
Below is a sample table showing data:

\begin{table}[h]
\centering
\begin{tabular}{@{}lrrr@{}}
\toprule
\textbf{Category} & \textbf{Value 1} & \textbf{Value 2} & \textbf{Total} \\
\midrule
Item A & 10.5 & 20.3 & 30.8 \\
Item B & 15.2 & 25.1 & 40.3 \\
Item C & 22.1 & 12.7 & 34.8 \\
\bottomrule
\end{tabular}
\caption{Sample Data Table}
\label{tab:sample}
\end{table}

\section{Lists}
\subsection{Bullet Points}
\begin{itemize}
  \item First item
  \item Second item
  \item Third item with sub-items:
    \begin{itemize}
      \item Sub-item 1
      \item Sub-item 2
    \end{itemize}
\end{itemize}

\subsection{Numbered List}
\begin{enumerate}
  \item First step
  \item Second step
  \item Third step
\end{enumerate}

\section{Citations and References}
When you use the citation search feature, references will be automatically managed. For example, you might cite a paper like this: \cite{example2024}. All citations will be automatically added to your \texttt{cite.bib} file.

\section{Conclusion}
This template demonstrates basic \LaTeX{} features including mathematical formulas, tables, and lists. The bibliography system is now set up to work seamlessly with Bibby AI's citation management. Feel free to modify it for your needs!

\bibliographystyle{plain}
\bibliography{cite}

\end{document}\documentclass{article}
\usepackage[utf8]{inputenc}
\usepackage{booktabs}
\usepackage{amsmath}
\usepackage{natbib}

\title{My Document}
\author{Your Name}
\date{\today}

\begin{document}

\maketitle

\section{Introduction}
This is a sample document demonstrating various \LaTeX{} features including mathematical equations, tables, and formatting. You can cite papers using the citation search feature, and they will be automatically added to your \texttt{cite.bib} file.

\section{Mathematical Expressions}
Here are some example equations:

\subsection{Quadratic Formula}
The quadratic formula is given by:
\[x = \frac{-b \pm \sqrt{b^2 - 4ac}}{2a}\]

\subsection{Matrix Example}
A sample matrix:
\[
\begin{pmatrix}
a_{11} & a_{12} & a_{13} \\
a_{21} & a_{22} & a_{23} \\
a_{31} & a_{32} & a_{33}
\end{pmatrix}
\]

\section{Tables}
Below is a sample table showing data:

\begin{table}[h]
\centering
\begin{tabular}{@{}lrrr@{}}
\toprule
\textbf{Category} & \textbf{Value 1} & \textbf{Value 2} & \textbf{Total} \\
\midrule
Item A & 10.5 & 20.3 & 30.8 \\
Item B & 15.2 & 25.1 & 40.3 \\
Item C & 22.1 & 12.7 & 34.8 \\
\bottomrule
\end{tabular}
\caption{Sample Data Table}
\label{tab:sample}
\end{table}

\section{Lists}
\subsection{Bullet Points}
\begin{itemize}
  \item First item
  \item Second item
  \item Third item with sub-items:
    \begin{itemize}
      \item Sub-item 1
      \item Sub-item 2
    \end{itemize}
\end{itemize}

\subsection{Numbered List}
\begin{enumerate}
  \item First step
  \item Second step
  \item Third step
\end{enumerate}

\section{Citations and References}
When you use the citation search feature, references will be automatically managed. For example, you might cite a paper like this: \cite{example2024}. All citations will be automatically added to your \texttt{cite.bib} file.

\section{Conclusion}
This template demonstrates basic \LaTeX{} features including mathematical formulas, tables, and lists. The bibliography system is now set up to work seamlessly with Bibby AI's citation management. Feel free to modify it for your needs!

\bibliographystyle{plain}
\bibliography{cite}

\end{document}\documentclass{article}
\usepackage[utf8]{inputenc}
\usepackage{booktabs}
\usepackage{amsmath}
\usepackage{natbib}

\title{My Document}
\author{Your Name}
\date{\today}

\begin{document}

\maketitle

\section{Introduction}
This is a sample document demonstrating various \LaTeX{} features including mathematical equations, tables, and formatting. You can cite papers using the citation search feature, and they will be automatically added to your \texttt{cite.bib} file.

\section{Mathematical Expressions}
Here are some example equations:

\subsection{Quadratic Formula}
The quadratic formula is given by:
\[x = \frac{-b \pm \sqrt{b^2 - 4ac}}{2a}\]

\subsection{Matrix Example}
A sample matrix:
\[
\begin{pmatrix}
a_{11} & a_{12} & a_{13} \\
a_{21} & a_{22} & a_{23} \\
a_{31} & a_{32} & a_{33}
\end{pmatrix}
\]

\section{Tables}
Below is a sample table showing data:

\begin{table}[h]
\centering
\begin{tabular}{@{}lrrr@{}}
\toprule
\textbf{Category} & \textbf{Value 1} & \textbf{Value 2} & \textbf{Total} \\
\midrule
Item A & 10.5 & 20.3 & 30.8 \\
Item B & 15.2 & 25.1 & 40.3 \\
Item C & 22.1 & 12.7 & 34.8 \\
\bottomrule
\end{tabular}
\caption{Sample Data Table}
\label{tab:sample}
\end{table}

\section{Lists}
\subsection{Bullet Points}
\begin{itemize}
  \item First item
  \item Second item
  \item Third item with sub-items:
    \begin{itemize}
      \item Sub-item 1
      \item Sub-item 2
    \end{itemize}
\end{itemize}

\subsection{Numbered List}
\begin{enumerate}
  \item First step
  \item Second step
  \item Third step
\end{enumerate}

\section{Citations and References}
When you use the citation search feature, references will be automatically managed. For example, you might cite a paper like this: \cite{example2024}. All citations will be automatically added to your \texttt{cite.bib} file.

\section{Conclusion}
This template demonstrates basic \LaTeX{} features including mathematical formulas, tables, and lists. The bibliography system is now set up to work seamlessly with Bibby AI's citation management. Feel free to modify it for your needs!

\bibliographystyle{plain}
\bibliography{cite}

\end{document}\documentclass{article}
\usepackage[utf8]{inputenc}
\usepackage{booktabs}
\usepackage{amsmath}
\usepackage{natbib}

\title{My Document}
\author{Your Name}
\date{\today}

\begin{document}

\maketitle

\section{Introduction}
This is a sample document demonstrating various \LaTeX{} features including mathematical equations, tables, and formatting. You can cite papers using the citation search feature, and they will be automatically added to your \texttt{cite.bib} file.

\section{Mathematical Expressions}
Here are some example equations:

\subsection{Quadratic Formula}
The quadratic formula is given by:
\[x = \frac{-b \pm \sqrt{b^2 - 4ac}}{2a}\]

\subsection{Matrix Example}
A sample matrix:
\[
\begin{pmatrix}
a_{11} & a_{12} & a_{13} \\
a_{21} & a_{22} & a_{23} \\
a_{31} & a_{32} & a_{33}
\end{pmatrix}
\]

\section{Tables}
Below is a sample table showing data:

\begin{table}[h]
\centering
\begin{tabular}{@{}lrrr@{}}
\toprule
\textbf{Category} & \textbf{Value 1} & \textbf{Value 2} & \textbf{Total} \\
\midrule
Item A & 10.5 & 20.3 & 30.8 \\
Item B & 15.2 & 25.1 & 40.3 \\
Item C & 22.1 & 12.7 & 34.8 \\
\bottomrule
\end{tabular}
\caption{Sample Data Table}
\label{tab:sample}
\end{table}

\section{Lists}
\subsection{Bullet Points}
\begin{itemize}
  \item First item
  \item Second item
  \item Third item with sub-items:
    \begin{itemize}
      \item Sub-item 1
      \item Sub-item 2
    \end{itemize}
\end{itemize}

\subsection{Numbered List}
\begin{enumerate}
  \item First step
  \item Second step
  \item Third step
\end{enumerate}

\section{Citations and References}
When you use the citation search feature, references will be automatically managed. For example, you might cite a paper like this: \cite{example2024}. All citations will be automatically added to your \texttt{cite.bib} file.

\section{Conclusion}
This template demonstrates basic \LaTeX{} features including mathematical formulas, tables, and lists. The bibliography system is now set up to work seamlessly with Bibby AI's citation management. Feel free to modify it for your needs!

\bibliographystyle{plain}
\bibliography{cite}

\end{document}\documentclass{article}
\usepackage[utf8]{inputenc}
\usepackage{booktabs}
\usepackage{amsmath}
\usepackage{natbib}

\title{My Document}
\author{Your Name}
\date{\today}

\begin{document}

\maketitle

\section{Introduction}
This is a sample document demonstrating various \LaTeX{} features including mathematical equations, tables, and formatting. You can cite papers using the citation search feature, and they will be automatically added to your \texttt{cite.bib} file.

\section{Mathematical Expressions}
Here are some example equations:

\subsection{Quadratic Formula}
The quadratic formula is given by:
\[x = \frac{-b \pm \sqrt{b^2 - 4ac}}{2a}\]

\subsection{Matrix Example}
A sample matrix:
\[
\begin{pmatrix}
a_{11} & a_{12} & a_{13} \\
a_{21} & a_{22} & a_{23} \\
a_{31} & a_{32} & a_{33}
\end{pmatrix}
\]

\section{Tables}
Below is a sample table showing data:

\begin{table}[h]
\centering
\begin{tabular}{@{}lrrr@{}}
\toprule
\textbf{Category} & \textbf{Value 1} & \textbf{Value 2} & \textbf{Total} \\
\midrule
Item A & 10.5 & 20.3 & 30.8 \\
Item B & 15.2 & 25.1 & 40.3 \\
Item C & 22.1 & 12.7 & 34.8 \\
\bottomrule
\end{tabular}
\caption{Sample Data Table}
\label{tab:sample}
\end{table}

\section{Lists}
\subsection{Bullet Points}
\begin{itemize}
  \item First item
  \item Second item
  \item Third item with sub-items:
    \begin{itemize}
      \item Sub-item 1
      \item Sub-item 2
    \end{itemize}
\end{itemize}

\subsection{Numbered List}
\begin{enumerate}
  \item First step
  \item Second step
  \item Third step
\end{enumerate}

\section{Citations and References}
When you use the citation search feature, references will be automatically managed. For example, you might cite a paper like this: \cite{example2024}. All citations will be automatically added to your \texttt{cite.bib} file.

\section{Conclusion}
This template demonstrates basic \LaTeX{} features including mathematical formulas, tables, and lists. The bibliography system is now set up to work seamlessly with Bibby AI's citation management. Feel free to modify it for your needs!

\bibliographystyle{plain}
\bibliography{cite}

\end{document}\documentclass{article}
\usepackage[utf8]{inputenc}
\usepackage{booktabs}
\usepackage{amsmath}
\usepackage{natbib}

\title{My Document}
\author{Your Name}
\date{\today}

\begin{document}

\maketitle

\section{Introduction}
This is a sample document demonstrating various \LaTeX{} features including mathematical equations, tables, and formatting. You can cite papers using the citation search feature, and they will be automatically added to your \texttt{cite.bib} file.

\section{Mathematical Expressions}
Here are some example equations:

\subsection{Quadratic Formula}
The quadratic formula is given by:
\[x = \frac{-b \pm \sqrt{b^2 - 4ac}}{2a}\]

\subsection{Matrix Example}
A sample matrix:
\[
\begin{pmatrix}
a_{11} & a_{12} & a_{13} \\
a_{21} & a_{22} & a_{23} \\
a_{31} & a_{32} & a_{33}
\end{pmatrix}
\]

\section{Tables}
Below is a sample table showing data:

\begin{table}[h]
\centering
\begin{tabular}{@{}lrrr@{}}
\toprule
\textbf{Category} & \textbf{Value 1} & \textbf{Value 2} & \textbf{Total} \\
\midrule
Item A & 10.5 & 20.3 & 30.8 \\
Item B & 15.2 & 25.1 & 40.3 \\
Item C & 22.1 & 12.7 & 34.8 \\
\bottomrule
\end{tabular}
\caption{Sample Data Table}
\label{tab:sample}
\end{table}

\section{Lists}
\subsection{Bullet Points}
\begin{itemize}
  \item First item
  \item Second item
  \item Third item with sub-items:
    \begin{itemize}
      \item Sub-item 1
      \item Sub-item 2
    \end{itemize}
\end{itemize}

\subsection{Numbered List}
\begin{enumerate}
  \item First step
  \item Second step
  \item Third step
\end{enumerate}

\section{Citations and References}
When you use the citation search feature, references will be automatically managed. For example, you might cite a paper like this: \cite{example2024}. All citations will be automatically added to your \texttt{cite.bib} file.

\section{Conclusion}
This template demonstrates basic \LaTeX{} features including mathematical formulas, tables, and lists. The bibliography system is now set up to work seamlessly with Bibby AI's citation management. Feel free to modify it for your needs!

\bibliographystyle{plain}
\bibliography{cite}

\end{document}\documentclass{article}
\usepackage[utf8]{inputenc}
\usepackage{booktabs}
\usepackage{amsmath}
\usepackage{natbib}

\title{My Document}
\author{Your Name}
\date{\today}

\begin{document}

\maketitle

\section{Introduction}
This is a sample document demonstrating various \LaTeX{} features including mathematical equations, tables, and formatting. You can cite papers using the citation search feature, and they will be automatically added to your \texttt{cite.bib} file.

\section{Mathematical Expressions}
Here are some example equations:

\subsection{Quadratic Formula}
The quadratic formula is given by:
\[x = \frac{-b \pm \sqrt{b^2 - 4ac}}{2a}\]

\subsection{Matrix Example}
A sample matrix:
\[
\begin{pmatrix}
a_{11} & a_{12} & a_{13} \\
a_{21} & a_{22} & a_{23} \\
a_{31} & a_{32} & a_{33}
\end{pmatrix}
\]

\section{Tables}
Below is a sample table showing data:

\begin{table}[h]
\centering
\begin{tabular}{@{}lrrr@{}}
\toprule
\textbf{Category} & \textbf{Value 1} & \textbf{Value 2} & \textbf{Total} \\
\midrule
Item A & 10.5 & 20.3 & 30.8 \\
Item B & 15.2 & 25.1 & 40.3 \\
Item C & 22.1 & 12.7 & 34.8 \\
\bottomrule
\end{tabular}
\caption{Sample Data Table}
\label{tab:sample}
\end{table}

\section{Lists}
\subsection{Bullet Points}
\begin{itemize}
  \item First item
  \item Second item
  \item Third item with sub-items:
    \begin{itemize}
      \item Sub-item 1
      \item Sub-item 2
    \end{itemize}
\end{itemize}

\subsection{Numbered List}
\begin{enumerate}
  \item First step
  \item Second step
  \item Third step
\end{enumerate}

\section{Citations and References}
When you use the citation search feature, references will be automatically managed. For example, you might cite a paper like this: \cite{example2024}. All citations will be automatically added to your \texttt{cite.bib} file.

\section{Conclusion}
This template demonstrates basic \LaTeX{} features including mathematical formulas, tables, and lists. The bibliography system is now set up to work seamlessly with Bibby AI's citation management. Feel free to modify it for your needs!

\bibliographystyle{plain}
\bibliography{cite}

\end{document}\documentclass{article}
\usepackage[utf8]{inputenc}
\usepackage{booktabs}
\usepackage{amsmath}
\usepackage{natbib}

\title{My Document}
\author{Your Name}
\date{\today}

\begin{document}

\maketitle

\section{Introduction}
This is a sample document demonstrating various \LaTeX{} features including mathematical equations, tables, and formatting. You can cite papers using the citation search feature, and they will be automatically added to your \texttt{cite.bib} file.

\section{Mathematical Expressions}
Here are some example equations:

\subsection{Quadratic Formula}
The quadratic formula is given by:
\[x = \frac{-b \pm \sqrt{b^2 - 4ac}}{2a}\]

\subsection{Matrix Example}
A sample matrix:
\[
\begin{pmatrix}
a_{11} & a_{12} & a_{13} \\
a_{21} & a_{22} & a_{23} \\
a_{31} & a_{32} & a_{33}
\end{pmatrix}
\]

\section{Tables}
Below is a sample table showing data:

\begin{table}[h]
\centering
\begin{tabular}{@{}lrrr@{}}
\toprule
\textbf{Category} & \textbf{Value 1} & \textbf{Value 2} & \textbf{Total} \\
\midrule
Item A & 10.5 & 20.3 & 30.8 \\
Item B & 15.2 & 25.1 & 40.3 \\
Item C & 22.1 & 12.7 & 34.8 \\
\bottomrule
\end{tabular}
\caption{Sample Data Table}
\label{tab:sample}
\end{table}

\section{Lists}
\subsection{Bullet Points}
\begin{itemize}
  \item First item
  \item Second item
  \item Third item with sub-items:
    \begin{itemize}
      \item Sub-item 1
      \item Sub-item 2
    \end{itemize}
\end{itemize}

\subsection{Numbered List}
\begin{enumerate}
  \item First step
  \item Second step
  \item Third step
\end{enumerate}

\section{Citations and References}
When you use the citation search feature, references will be automatically managed. For example, you might cite a paper like this: \cite{example2024}. All citations will be automatically added to your \texttt{cite.bib} file.

\section{Conclusion}
This template demonstrates basic \LaTeX{} features including mathematical formulas, tables, and lists. The bibliography system is now set up to work seamlessly with Bibby AI's citation management. Feel free to modify it for your needs!

\bibliographystyle{plain}
\bibliography{cite}

\end{document}\documentclass{article}
\usepackage[utf8]{inputenc}
\usepackage{booktabs}
\usepackage{amsmath}
\usepackage{natbib}

\title{My Document}
\author{Your Name}
\date{\today}

\begin{document}

\maketitle

\section{Introduction}
This is a sample document demonstrating various \LaTeX{} features including mathematical equations, tables, and formatting. You can cite papers using the citation search feature, and they will be automatically added to your \texttt{cite.bib} file.

\section{Mathematical Expressions}
Here are some example equations:

\subsection{Quadratic Formula}
The quadratic formula is given by:
\[x = \frac{-b \pm \sqrt{b^2 - 4ac}}{2a}\]

\subsection{Matrix Example}
A sample matrix:
\[
\begin{pmatrix}
a_{11} & a_{12} & a_{13} \\
a_{21} & a_{22} & a_{23} \\
a_{31} & a_{32} & a_{33}
\end{pmatrix}
\]

\section{Tables}
Below is a sample table showing data:

\begin{table}[h]
\centering
\begin{tabular}{@{}lrrr@{}}
\toprule
\textbf{Category} & \textbf{Value 1} & \textbf{Value 2} & \textbf{Total} \\
\midrule
Item A & 10.5 & 20.3 & 30.8 \\
Item B & 15.2 & 25.1 & 40.3 \\
Item C & 22.1 & 12.7 & 34.8 \\
\bottomrule
\end{tabular}
\caption{Sample Data Table}
\label{tab:sample}
\end{table}

\section{Lists}
\subsection{Bullet Points}
\begin{itemize}
  \item First item
  \item Second item
  \item Third item with sub-items:
    \begin{itemize}
      \item Sub-item 1
      \item Sub-item 2
    \end{itemize}
\end{itemize}

\subsection{Numbered List}
\begin{enumerate}
  \item First step
  \item Second step
  \item Third step
\end{enumerate}

\section{Citations and References}
When you use the citation search feature, references will be automatically managed. For example, you might cite a paper like this: \cite{example2024}. All citations will be automatically added to your \texttt{cite.bib} file.

\section{Conclusion}
This template demonstrates basic \LaTeX{} features including mathematical formulas, tables, and lists. The bibliography system is now set up to work seamlessly with Bibby AI's citation management. Feel free to modify it for your needs!

\bibliographystyle{plain}
\bibliography{cite}

\end{document}\documentclass{article}
\usepackage[utf8]{inputenc}
\usepackage{booktabs}
\usepackage{amsmath}
\usepackage{natbib}

\title{My Document}
\author{Your Name}
\date{\today}

\begin{document}

\maketitle

\section{Introduction}
This is a sample document demonstrating various \LaTeX{} features including mathematical equations, tables, and formatting. You can cite papers using the citation search feature, and they will be automatically added to your \texttt{cite.bib} file.

\section{Mathematical Expressions}
Here are some example equations:

\subsection{Quadratic Formula}
The quadratic formula is given by:
\[x = \frac{-b \pm \sqrt{b^2 - 4ac}}{2a}\]

\subsection{Matrix Example}
A sample matrix:
\[
\begin{pmatrix}
a_{11} & a_{12} & a_{13} \\
a_{21} & a_{22} & a_{23} \\
a_{31} & a_{32} & a_{33}
\end{pmatrix}
\]

\section{Tables}
Below is a sample table showing data:

\begin{table}[h]
\centering
\begin{tabular}{@{}lrrr@{}}
\toprule
\textbf{Category} & \textbf{Value 1} & \textbf{Value 2} & \textbf{Total} \\
\midrule
Item A & 10.5 & 20.3 & 30.8 \\
Item B & 15.2 & 25.1 & 40.3 \\
Item C & 22.1 & 12.7 & 34.8 \\
\bottomrule
\end{tabular}
\caption{Sample Data Table}
\label{tab:sample}
\end{table}

\section{Lists}
\subsection{Bullet Points}
\begin{itemize}
  \item First item
  \item Second item
  \item Third item with sub-items:
    \begin{itemize}
      \item Sub-item 1
      \item Sub-item 2
    \end{itemize}
\end{itemize}

\subsection{Numbered List}
\begin{enumerate}
  \item First step
  \item Second step
  \item Third step
\end{enumerate}

\section{Citations and References}
When you use the citation search feature, references will be automatically managed. For example, you might cite a paper like this: \cite{example2024}. All citations will be automatically added to your \texttt{cite.bib} file.

\section{Conclusion}
This template demonstrates basic \LaTeX{} features including mathematical formulas, tables, and lists. The bibliography system is now set up to work seamlessly with Bibby AI's citation management. Feel free to modify it for your needs!

\bibliographystyle{plain}
\bibliography{cite}

\end{document}\documentclass{article}
\usepackage[utf8]{inputenc}
\usepackage{booktabs}
\usepackage{amsmath}
\usepackage{natbib}

\title{My Document}
\author{Your Name}
\date{\today}

\begin{document}

\maketitle

\section{Introduction}
This is a sample document demonstrating various \LaTeX{} features including mathematical equations, tables, and formatting. You can cite papers using the citation search feature, and they will be automatically added to your \texttt{cite.bib} file.

\section{Mathematical Expressions}
Here are some example equations:

\subsection{Quadratic Formula}
The quadratic formula is given by:
\[x = \frac{-b \pm \sqrt{b^2 - 4ac}}{2a}\]

\subsection{Matrix Example}
A sample matrix:
\[
\begin{pmatrix}
a_{11} & a_{12} & a_{13} \\
a_{21} & a_{22} & a_{23} \\
a_{31} & a_{32} & a_{33}
\end{pmatrix}
\]

\section{Tables}
Below is a sample table showing data:

\begin{table}[h]
\centering
\begin{tabular}{@{}lrrr@{}}
\toprule
\textbf{Category} & \textbf{Value 1} & \textbf{Value 2} & \textbf{Total} \\
\midrule
Item A & 10.5 & 20.3 & 30.8 \\
Item B & 15.2 & 25.1 & 40.3 \\
Item C & 22.1 & 12.7 & 34.8 \\
\bottomrule
\end{tabular}
\caption{Sample Data Table}
\label{tab:sample}
\end{table}

\section{Lists}
\subsection{Bullet Points}
\begin{itemize}
  \item First item
  \item Second item
  \item Third item with sub-items:
    \begin{itemize}
      \item Sub-item 1
      \item Sub-item 2
    \end{itemize}
\end{itemize}

\subsection{Numbered List}
\begin{enumerate}
  \item First step
  \item Second step
  \item Third step
\end{enumerate}

\section{Citations and References}
When you use the citation search feature, references will be automatically managed. For example, you might cite a paper like this: \cite{example2024}. All citations will be automatically added to your \texttt{cite.bib} file.

\section{Conclusion}
This template demonstrates basic \LaTeX{} features including mathematical formulas, tables, and lists. The bibliography system is now set up to work seamlessly with Bibby AI's citation management. Feel free to modify it for your needs!

\bibliographystyle{plain}
\bibliography{cite}

\end{document}\documentclass{article}
\usepackage[utf8]{inputenc}
\usepackage{booktabs}
\usepackage{amsmath}
\usepackage{natbib}

\title{My Document}
\author{Your Name}
\date{\today}

\begin{document}

\maketitle

\section{Introduction}
This is a sample document demonstrating various \LaTeX{} features including mathematical equations, tables, and formatting. You can cite papers using the citation search feature, and they will be automatically added to your \texttt{cite.bib} file.

\section{Mathematical Expressions}
Here are some example equations:

\subsection{Quadratic Formula}
The quadratic formula is given by:
\[x = \frac{-b \pm \sqrt{b^2 - 4ac}}{2a}\]

\subsection{Matrix Example}
A sample matrix:
\[
\begin{pmatrix}
a_{11} & a_{12} & a_{13} \\
a_{21} & a_{22} & a_{23} \\
a_{31} & a_{32} & a_{33}
\end{pmatrix}
\]

\section{Tables}
Below is a sample table showing data:

\begin{table}[h]
\centering
\begin{tabular}{@{}lrrr@{}}
\toprule
\textbf{Category} & \textbf{Value 1} & \textbf{Value 2} & \textbf{Total} \\
\midrule
Item A & 10.5 & 20.3 & 30.8 \\
Item B & 15.2 & 25.1 & 40.3 \\
Item C & 22.1 & 12.7 & 34.8 \\
\bottomrule
\end{tabular}
\caption{Sample Data Table}
\label{tab:sample}
\end{table}

\section{Lists}
\subsection{Bullet Points}
\begin{itemize}
  \item First item
  \item Second item
  \item Third item with sub-items:
    \begin{itemize}
      \item Sub-item 1
      \item Sub-item 2
    \end{itemize}
\end{itemize}

\subsection{Numbered List}
\begin{enumerate}
  \item First step
  \item Second step
  \item Third step
\end{enumerate}

\section{Citations and References}
When you use the citation search feature, references will be automatically managed. For example, you might cite a paper like this: \cite{example2024}. All citations will be automatically added to your \texttt{cite.bib} file.

\section{Conclusion}
This template demonstrates basic \LaTeX{} features including mathematical formulas, tables, and lists. The bibliography system is now set up to work seamlessly with Bibby AI's citation management. Feel free to modify it for your needs!

\bibliographystyle{plain}
\bibliography{cite}

\end{document}\documentclass{article}
\usepackage[utf8]{inputenc}
\usepackage{booktabs}
\usepackage{amsmath}
\usepackage{natbib}

\title{My Document}
\author{Your Name}
\date{\today}

\begin{document}

\maketitle

\section{Introduction}
This is a sample document demonstrating various \LaTeX{} features including mathematical equations, tables, and formatting. You can cite papers using the citation search feature, and they will be automatically added to your \texttt{cite.bib} file.

\section{Mathematical Expressions}
Here are some example equations:

\subsection{Quadratic Formula}
The quadratic formula is given by:
\[x = \frac{-b \pm \sqrt{b^2 - 4ac}}{2a}\]

\subsection{Matrix Example}
A sample matrix:
\[
\begin{pmatrix}
a_{11} & a_{12} & a_{13} \\
a_{21} & a_{22} & a_{23} \\
a_{31} & a_{32} & a_{33}
\end{pmatrix}
\]

\section{Tables}
Below is a sample table showing data:

\begin{table}[h]
\centering
\begin{tabular}{@{}lrrr@{}}
\toprule
\textbf{Category} & \textbf{Value 1} & \textbf{Value 2} & \textbf{Total} \\
\midrule
Item A & 10.5 & 20.3 & 30.8 \\
Item B & 15.2 & 25.1 & 40.3 \\
Item C & 22.1 & 12.7 & 34.8 \\
\bottomrule
\end{tabular}
\caption{Sample Data Table}
\label{tab:sample}
\end{table}

\section{Lists}
\subsection{Bullet Points}
\begin{itemize}
  \item First item
  \item Second item
  \item Third item with sub-items:
    \begin{itemize}
      \item Sub-item 1
      \item Sub-item 2
    \end{itemize}
\end{itemize}

\subsection{Numbered List}
\begin{enumerate}
  \item First step
  \item Second step
  \item Third step
\end{enumerate}

\section{Citations and References}
When you use the citation search feature, references will be automatically managed. For example, you might cite a paper like this: \cite{example2024}. All citations will be automatically added to your \texttt{cite.bib} file.

\section{Conclusion}
This template demonstrates basic \LaTeX{} features including mathematical formulas, tables, and lists. The bibliography system is now set up to work seamlessly with Bibby AI's citation management. Feel free to modify it for your needs!

\bibliographystyle{plain}
\bibliography{cite}

\end{document}\documentclass{article}
\usepackage[utf8]{inputenc}
\usepackage{booktabs}
\usepackage{amsmath}
\usepackage{natbib}

\title{My Document}
\author{Your Name}
\date{\today}

\begin{document}

\maketitle

\section{Introduction}
This is a sample document demonstrating various \LaTeX{} features including mathematical equations, tables, and formatting. You can cite papers using the citation search feature, and they will be automatically added to your \texttt{cite.bib} file.

\section{Mathematical Expressions}
Here are some example equations:

\subsection{Quadratic Formula}
The quadratic formula is given by:
\[x = \frac{-b \pm \sqrt{b^2 - 4ac}}{2a}\]

\subsection{Matrix Example}
A sample matrix:
\[
\begin{pmatrix}
a_{11} & a_{12} & a_{13} \\
a_{21} & a_{22} & a_{23} \\
a_{31} & a_{32} & a_{33}
\end{pmatrix}
\]

\section{Tables}
Below is a sample table showing data:

\begin{table}[h]
\centering
\begin{tabular}{@{}lrrr@{}}
\toprule
\textbf{Category} & \textbf{Value 1} & \textbf{Value 2} & \textbf{Total} \\
\midrule
Item A & 10.5 & 20.3 & 30.8 \\
Item B & 15.2 & 25.1 & 40.3 \\
Item C & 22.1 & 12.7 & 34.8 \\
\bottomrule
\end{tabular}
\caption{Sample Data Table}
\label{tab:sample}
\end{table}

\section{Lists}
\subsection{Bullet Points}
\begin{itemize}
  \item First item
  \item Second item
  \item Third item with sub-items:
    \begin{itemize}
      \item Sub-item 1
      \item Sub-item 2
    \end{itemize}
\end{itemize}

\subsection{Numbered List}
\begin{enumerate}
  \item First step
  \item Second step
  \item Third step
\end{enumerate}

\section{Citations and References}
When you use the citation search feature, references will be automatically managed. For example, you might cite a paper like this: \cite{example2024}. All citations will be automatically added to your \texttt{cite.bib} file.

\section{Conclusion}
This template demonstrates basic \LaTeX{} features including mathematical formulas, tables, and lists. The bibliography system is now set up to work seamlessly with Bibby AI's citation management. Feel free to modify it for your needs!

\bibliographystyle{plain}
\bibliography{cite}

\end{document}\documentclass{article}
\usepackage[utf8]{inputenc}
\usepackage{booktabs}
\usepackage{amsmath}
\usepackage{natbib}

\title{My Document}
\author{Your Name}
\date{\today}

\begin{document}

\maketitle

\section{Introduction}
This is a sample document demonstrating various \LaTeX{} features including mathematical equations, tables, and formatting. You can cite papers using the citation search feature, and they will be automatically added to your \texttt{cite.bib} file.

\section{Mathematical Expressions}
Here are some example equations:

\subsection{Quadratic Formula}
The quadratic formula is given by:
\[x = \frac{-b \pm \sqrt{b^2 - 4ac}}{2a}\]

\subsection{Matrix Example}
A sample matrix:
\[
\begin{pmatrix}
a_{11} & a_{12} & a_{13} \\
a_{21} & a_{22} & a_{23} \\
a_{31} & a_{32} & a_{33}
\end{pmatrix}
\]

\section{Tables}
Below is a sample table showing data:

\begin{table}[h]
\centering
\begin{tabular}{@{}lrrr@{}}
\toprule
\textbf{Category} & \textbf{Value 1} & \textbf{Value 2} & \textbf{Total} \\
\midrule
Item A & 10.5 & 20.3 & 30.8 \\
Item B & 15.2 & 25.1 & 40.3 \\
Item C & 22.1 & 12.7 & 34.8 \\
\bottomrule
\end{tabular}
\caption{Sample Data Table}
\label{tab:sample}
\end{table}

\section{Lists}
\subsection{Bullet Points}
\begin{itemize}
  \item First item
  \item Second item
  \item Third item with sub-items:
    \begin{itemize}
      \item Sub-item 1
      \item Sub-item 2
    \end{itemize}
\end{itemize}

\subsection{Numbered List}
\begin{enumerate}
  \item First step
  \item Second step
  \item Third step
\end{enumerate}

\section{Citations and References}
When you use the citation search feature, references will be automatically managed. For example, you might cite a paper like this: \cite{example2024}. All citations will be automatically added to your \texttt{cite.bib} file.

\section{Conclusion}
This template demonstrates basic \LaTeX{} features including mathematical formulas, tables, and lists. The bibliography system is now set up to work seamlessly with Bibby AI's citation management. Feel free to modify it for your needs!

\bibliographystyle{plain}
\bibliography{cite}

\end{document}\documentclass{article}
\usepackage[utf8]{inputenc}
\usepackage{booktabs}
\usepackage{amsmath}
\usepackage{natbib}

\title{My Document}
\author{Your Name}
\date{\today}

\begin{document}

\maketitle

\section{Introduction}
This is a sample document demonstrating various \LaTeX{} features including mathematical equations, tables, and formatting. You can cite papers using the citation search feature, and they will be automatically added to your \texttt{cite.bib} file.

\section{Mathematical Expressions}
Here are some example equations:

\subsection{Quadratic Formula}
The quadratic formula is given by:
\[x = \frac{-b \pm \sqrt{b^2 - 4ac}}{2a}\]

\subsection{Matrix Example}
A sample matrix:
\[
\begin{pmatrix}
a_{11} & a_{12} & a_{13} \\
a_{21} & a_{22} & a_{23} \\
a_{31} & a_{32} & a_{33}
\end{pmatrix}
\]

\section{Tables}
Below is a sample table showing data:

\begin{table}[h]
\centering
\begin{tabular}{@{}lrrr@{}}
\toprule
\textbf{Category} & \textbf{Value 1} & \textbf{Value 2} & \textbf{Total} \\
\midrule
Item A & 10.5 & 20.3 & 30.8 \\
Item B & 15.2 & 25.1 & 40.3 \\
Item C & 22.1 & 12.7 & 34.8 \\
\bottomrule
\end{tabular}
\caption{Sample Data Table}
\label{tab:sample}
\end{table}

\section{Lists}
\subsection{Bullet Points}
\begin{itemize}
  \item First item
  \item Second item
  \item Third item with sub-items:
    \begin{itemize}
      \item Sub-item 1
      \item Sub-item 2
    \end{itemize}
\end{itemize}

\subsection{Numbered List}
\begin{enumerate}
  \item First step
  \item Second step
  \item Third step
\end{enumerate}

\section{Citations and References}
When you use the citation search feature, references will be automatically managed. For example, you might cite a paper like this: \cite{example2024}. All citations will be automatically added to your \texttt{cite.bib} file.

\section{Conclusion}
This template demonstrates basic \LaTeX{} features including mathematical formulas, tables, and lists. The bibliography system is now set up to work seamlessly with Bibby AI's citation management. Feel free to modify it for your needs!

\bibliographystyle{plain}
\bibliography{cite}

\end{document}\documentclass{article}
\usepackage[utf8]{inputenc}
\usepackage{booktabs}
\usepackage{amsmath}
\usepackage{natbib}

\title{My Document}
\author{Your Name}
\date{\today}

\begin{document}

\maketitle

\section{Introduction}
This is a sample document demonstrating various \LaTeX{} features including mathematical equations, tables, and formatting. You can cite papers using the citation search feature, and they will be automatically added to your \texttt{cite.bib} file.

\section{Mathematical Expressions}
Here are some example equations:

\subsection{Quadratic Formula}
The quadratic formula is given by:
\[x = \frac{-b \pm \sqrt{b^2 - 4ac}}{2a}\]

\subsection{Matrix Example}
A sample matrix:
\[
\begin{pmatrix}
a_{11} & a_{12} & a_{13} \\
a_{21} & a_{22} & a_{23} \\
a_{31} & a_{32} & a_{33}
\end{pmatrix}
\]

\section{Tables}
Below is a sample table showing data:

\begin{table}[h]
\centering
\begin{tabular}{@{}lrrr@{}}
\toprule
\textbf{Category} & \textbf{Value 1} & \textbf{Value 2} & \textbf{Total} \\
\midrule
Item A & 10.5 & 20.3 & 30.8 \\
Item B & 15.2 & 25.1 & 40.3 \\
Item C & 22.1 & 12.7 & 34.8 \\
\bottomrule
\end{tabular}
\caption{Sample Data Table}
\label{tab:sample}
\end{table}

\section{Lists}
\subsection{Bullet Points}
\begin{itemize}
  \item First item
  \item Second item
  \item Third item with sub-items:
    \begin{itemize}
      \item Sub-item 1
      \item Sub-item 2
    \end{itemize}
\end{itemize}

\subsection{Numbered List}
\begin{enumerate}
  \item First step
  \item Second step
  \item Third step
\end{enumerate}

\section{Citations and References}
When you use the citation search feature, references will be automatically managed. For example, you might cite a paper like this: \cite{example2024}. All citations will be automatically added to your \texttt{cite.bib} file.

\section{Conclusion}
This template demonstrates basic \LaTeX{} features including mathematical formulas, tables, and lists. The bibliography system is now set up to work seamlessly with Bibby AI's citation management. Feel free to modify it for your needs!

\bibliographystyle{plain}
\bibliography{cite}

\end{document}\documentclass{article}
\usepackage[utf8]{inputenc}
\usepackage{booktabs}
\usepackage{amsmath}
\usepackage{natbib}

\title{My Document}
\author{Your Name}
\date{\today}

\begin{document}

\maketitle

\section{Introduction}
This is a sample document demonstrating various \LaTeX{} features including mathematical equations, tables, and formatting. You can cite papers using the citation search feature, and they will be automatically added to your \texttt{cite.bib} file.

\section{Mathematical Expressions}
Here are some example equations:

\subsection{Quadratic Formula}
The quadratic formula is given by:
\[x = \frac{-b \pm \sqrt{b^2 - 4ac}}{2a}\]

\subsection{Matrix Example}
A sample matrix:
\[
\begin{pmatrix}
a_{11} & a_{12} & a_{13} \\
a_{21} & a_{22} & a_{23} \\
a_{31} & a_{32} & a_{33}
\end{pmatrix}
\]

\section{Tables}
Below is a sample table showing data:

\begin{table}[h]
\centering
\begin{tabular}{@{}lrrr@{}}
\toprule
\textbf{Category} & \textbf{Value 1} & \textbf{Value 2} & \textbf{Total} \\
\midrule
Item A & 10.5 & 20.3 & 30.8 \\
Item B & 15.2 & 25.1 & 40.3 \\
Item C & 22.1 & 12.7 & 34.8 \\
\bottomrule
\end{tabular}
\caption{Sample Data Table}
\label{tab:sample}
\end{table}

\section{Lists}
\subsection{Bullet Points}
\begin{itemize}
  \item First item
  \item Second item
  \item Third item with sub-items:
    \begin{itemize}
      \item Sub-item 1
      \item Sub-item 2
    \end{itemize}
\end{itemize}

\subsection{Numbered List}
\begin{enumerate}
  \item First step
  \item Second step
  \item Third step
\end{enumerate}

\section{Citations and References}
When you use the citation search feature, references will be automatically managed. For example, you might cite a paper like this: \cite{example2024}. All citations will be automatically added to your \texttt{cite.bib} file.

\section{Conclusion}
This template demonstrates basic \LaTeX{} features including mathematical formulas, tables, and lists. The bibliography system is now set up to work seamlessly with Bibby AI's citation management. Feel free to modify it for your needs!

\bibliographystyle{plain}
\bibliography{cite}

\end{document}\documentclass{article}
\usepackage[utf8]{inputenc}
\usepackage{booktabs}
\usepackage{amsmath}
\usepackage{natbib}

\title{My Document}
\author{Your Name}
\date{\today}

\begin{document}

\maketitle

\section{Introduction}
This is a sample document demonstrating various \LaTeX{} features including mathematical equations, tables, and formatting. You can cite papers using the citation search feature, and they will be automatically added to your \texttt{cite.bib} file.

\section{Mathematical Expressions}
Here are some example equations:

\subsection{Quadratic Formula}
The quadratic formula is given by:
\[x = \frac{-b \pm \sqrt{b^2 - 4ac}}{2a}\]

\subsection{Matrix Example}
A sample matrix:
\[
\begin{pmatrix}
a_{11} & a_{12} & a_{13} \\
a_{21} & a_{22} & a_{23} \\
a_{31} & a_{32} & a_{33}
\end{pmatrix}
\]

\section{Tables}
Below is a sample table showing data:

\begin{table}[h]
\centering
\begin{tabular}{@{}lrrr@{}}
\toprule
\textbf{Category} & \textbf{Value 1} & \textbf{Value 2} & \textbf{Total} \\
\midrule
Item A & 10.5 & 20.3 & 30.8 \\
Item B & 15.2 & 25.1 & 40.3 \\
Item C & 22.1 & 12.7 & 34.8 \\
\bottomrule
\end{tabular}
\caption{Sample Data Table}
\label{tab:sample}
\end{table}

\section{Lists}
\subsection{Bullet Points}
\begin{itemize}
  \item First item
  \item Second item
  \item Third item with sub-items:
    \begin{itemize}
      \item Sub-item 1
      \item Sub-item 2
    \end{itemize}
\end{itemize}

\subsection{Numbered List}
\begin{enumerate}
  \item First step
  \item Second step
  \item Third step
\end{enumerate}

\section{Citations and References}
When you use the citation search feature, references will be automatically managed. For example, you might cite a paper like this: \cite{example2024}. All citations will be automatically added to your \texttt{cite.bib} file.

\section{Conclusion}
This template demonstrates basic \LaTeX{} features including mathematical formulas, tables, and lists. The bibliography system is now set up to work seamlessly with Bibby AI's citation management. Feel free to modify it for your needs!

\bibliographystyle{plain}
\bibliography{cite}

\end{document}\documentclass{article}
\usepackage[utf8]{inputenc}
\usepackage{booktabs}
\usepackage{amsmath}
\usepackage{natbib}

\title{My Document}
\author{Your Name}
\date{\today}

\begin{document}

\maketitle

\section{Introduction}
This is a sample document demonstrating various \LaTeX{} features including mathematical equations, tables, and formatting. You can cite papers using the citation search feature, and they will be automatically added to your \texttt{cite.bib} file.

\section{Mathematical Expressions}
Here are some example equations:

\subsection{Quadratic Formula}
The quadratic formula is given by:
\[x = \frac{-b \pm \sqrt{b^2 - 4ac}}{2a}\]

\subsection{Matrix Example}
A sample matrix:
\[
\begin{pmatrix}
a_{11} & a_{12} & a_{13} \\
a_{21} & a_{22} & a_{23} \\
a_{31} & a_{32} & a_{33}
\end{pmatrix}
\]

\section{Tables}
Below is a sample table showing data:

\begin{table}[h]
\centering
\begin{tabular}{@{}lrrr@{}}
\toprule
\textbf{Category} & \textbf{Value 1} & \textbf{Value 2} & \textbf{Total} \\
\midrule
Item A & 10.5 & 20.3 & 30.8 \\
Item B & 15.2 & 25.1 & 40.3 \\
Item C & 22.1 & 12.7 & 34.8 \\
\bottomrule
\end{tabular}
\caption{Sample Data Table}
\label{tab:sample}
\end{table}

\section{Lists}
\subsection{Bullet Points}
\begin{itemize}
  \item First item
  \item Second item
  \item Third item with sub-items:
    \begin{itemize}
      \item Sub-item 1
      \item Sub-item 2
    \end{itemize}
\end{itemize}

\subsection{Numbered List}
\begin{enumerate}
  \item First step
  \item Second step
  \item Third step
\end{enumerate}

\section{Citations and References}
When you use the citation search feature, references will be automatically managed. For example, you might cite a paper like this: \cite{example2024}. All citations will be automatically added to your \texttt{cite.bib} file.

\section{Conclusion}
This template demonstrates basic \LaTeX{} features including mathematical formulas, tables, and lists. The bibliography system is now set up to work seamlessly with Bibby AI's citation management. Feel free to modify it for your needs!

\bibliographystyle{plain}
\bibliography{cite}

\end{document}\documentclass{article}
\usepackage[utf8]{inputenc}
\usepackage{booktabs}
\usepackage{amsmath}
\usepackage{natbib}

\title{My Document}
\author{Your Name}
\date{\today}

\begin{document}

\maketitle

\section{Introduction}
This is a sample document demonstrating various \LaTeX{} features including mathematical equations, tables, and formatting. You can cite papers using the citation search feature, and they will be automatically added to your \texttt{cite.bib} file.

\section{Mathematical Expressions}
Here are some example equations:

\subsection{Quadratic Formula}
The quadratic formula is given by:
\[x = \frac{-b \pm \sqrt{b^2 - 4ac}}{2a}\]

\subsection{Matrix Example}
A sample matrix:
\[
\begin{pmatrix}
a_{11} & a_{12} & a_{13} \\
a_{21} & a_{22} & a_{23} \\
a_{31} & a_{32} & a_{33}
\end{pmatrix}
\]

\section{Tables}
Below is a sample table showing data:

\begin{table}[h]
\centering
\begin{tabular}{@{}lrrr@{}}
\toprule
\textbf{Category} & \textbf{Value 1} & \textbf{Value 2} & \textbf{Total} \\
\midrule
Item A & 10.5 & 20.3 & 30.8 \\
Item B & 15.2 & 25.1 & 40.3 \\
Item C & 22.1 & 12.7 & 34.8 \\
\bottomrule
\end{tabular}
\caption{Sample Data Table}
\label{tab:sample}
\end{table}

\section{Lists}
\subsection{Bullet Points}
\begin{itemize}
  \item First item
  \item Second item
  \item Third item with sub-items:
    \begin{itemize}
      \item Sub-item 1
      \item Sub-item 2
    \end{itemize}
\end{itemize}

\subsection{Numbered List}
\begin{enumerate}
  \item First step
  \item Second step
  \item Third step
\end{enumerate}

\section{Citations and References}
When you use the citation search feature, references will be automatically managed. For example, you might cite a paper like this: \cite{example2024}. All citations will be automatically added to your \texttt{cite.bib} file.

\section{Conclusion}
This template demonstrates basic \LaTeX{} features including mathematical formulas, tables, and lists. The bibliography system is now set up to work seamlessly with Bibby AI's citation management. Feel free to modify it for your needs!

\bibliographystyle{plain}
\bibliography{cite}

\end{document}\documentclass{article}
\usepackage[utf8]{inputenc}
\usepackage{booktabs}
\usepackage{amsmath}
\usepackage{natbib}

\title{My Document}
\author{Your Name}
\date{\today}

\begin{document}

\maketitle

\section{Introduction}
This is a sample document demonstrating various \LaTeX{} features including mathematical equations, tables, and formatting. You can cite papers using the citation search feature, and they will be automatically added to your \texttt{cite.bib} file.

\section{Mathematical Expressions}
Here are some example equations:

\subsection{Quadratic Formula}
The quadratic formula is given by:
\[x = \frac{-b \pm \sqrt{b^2 - 4ac}}{2a}\]

\subsection{Matrix Example}
A sample matrix:
\[
\begin{pmatrix}
a_{11} & a_{12} & a_{13} \\
a_{21} & a_{22} & a_{23} \\
a_{31} & a_{32} & a_{33}
\end{pmatrix}
\]

\section{Tables}
Below is a sample table showing data:

\begin{table}[h]
\centering
\begin{tabular}{@{}lrrr@{}}
\toprule
\textbf{Category} & \textbf{Value 1} & \textbf{Value 2} & \textbf{Total} \\
\midrule
Item A & 10.5 & 20.3 & 30.8 \\
Item B & 15.2 & 25.1 & 40.3 \\
Item C & 22.1 & 12.7 & 34.8 \\
\bottomrule
\end{tabular}
\caption{Sample Data Table}
\label{tab:sample}
\end{table}

\section{Lists}
\subsection{Bullet Points}
\begin{itemize}
  \item First item
  \item Second item
  \item Third item with sub-items:
    \begin{itemize}
      \item Sub-item 1
      \item Sub-item 2
    \end{itemize}
\end{itemize}

\subsection{Numbered List}
\begin{enumerate}
  \item First step
  \item Second step
  \item Third step
\end{enumerate}

\section{Citations and References}
When you use the citation search feature, references will be automatically managed. For example, you might cite a paper like this: \cite{example2024}. All citations will be automatically added to your \texttt{cite.bib} file.

\section{Conclusion}
This template demonstrates basic \LaTeX{} features including mathematical formulas, tables, and lists. The bibliography system is now set up to work seamlessly with Bibby AI's citation management. Feel free to modify it for your needs!

\bibliographystyle{plain}
\bibliography{cite}

\end{document}\documentclass{article}
\usepackage[utf8]{inputenc}
\usepackage{booktabs}
\usepackage{amsmath}
\usepackage{natbib}

\title{My Document}
\author{Your Name}
\date{\today}

\begin{document}

\maketitle

\section{Introduction}
This is a sample document demonstrating various \LaTeX{} features including mathematical equations, tables, and formatting. You can cite papers using the citation search feature, and they will be automatically added to your \texttt{cite.bib} file.

\section{Mathematical Expressions}
Here are some example equations:

\subsection{Quadratic Formula}
The quadratic formula is given by:
\[x = \frac{-b \pm \sqrt{b^2 - 4ac}}{2a}\]

\subsection{Matrix Example}
A sample matrix:
\[
\begin{pmatrix}
a_{11} & a_{12} & a_{13} \\
a_{21} & a_{22} & a_{23} \\
a_{31} & a_{32} & a_{33}
\end{pmatrix}
\]

\section{Tables}
Below is a sample table showing data:

\begin{table}[h]
\centering
\begin{tabular}{@{}lrrr@{}}
\toprule
\textbf{Category} & \textbf{Value 1} & \textbf{Value 2} & \textbf{Total} \\
\midrule
Item A & 10.5 & 20.3 & 30.8 \\
Item B & 15.2 & 25.1 & 40.3 \\
Item C & 22.1 & 12.7 & 34.8 \\
\bottomrule
\end{tabular}
\caption{Sample Data Table}
\label{tab:sample}
\end{table}

\section{Lists}
\subsection{Bullet Points}
\begin{itemize}
  \item First item
  \item Second item
  \item Third item with sub-items:
    \begin{itemize}
      \item Sub-item 1
      \item Sub-item 2
    \end{itemize}
\end{itemize}

\subsection{Numbered List}
\begin{enumerate}
  \item First step
  \item Second step
  \item Third step
\end{enumerate}

\section{Citations and References}
When you use the citation search feature, references will be automatically managed. For example, you might cite a paper like this: \cite{example2024}. All citations will be automatically added to your \texttt{cite.bib} file.

\section{Conclusion}
This template demonstrates basic \LaTeX{} features including mathematical formulas, tables, and lists. The bibliography system is now set up to work seamlessly with Bibby AI's citation management. Feel free to modify it for your needs!

\bibliographystyle{plain}
\bibliography{cite}

\end{document}\documentclass{article}
\usepackage[utf8]{inputenc}
\usepackage{booktabs}
\usepackage{amsmath}
\usepackage{natbib}

\title{My Document}
\author{Your Name}
\date{\today}

\begin{document}

\maketitle

\section{Introduction}
This is a sample document demonstrating various \LaTeX{} features including mathematical equations, tables, and formatting. You can cite papers using the citation search feature, and they will be automatically added to your \texttt{cite.bib} file.

\section{Mathematical Expressions}
Here are some example equations:

\subsection{Quadratic Formula}
The quadratic formula is given by:
\[x = \frac{-b \pm \sqrt{b^2 - 4ac}}{2a}\]

\subsection{Matrix Example}
A sample matrix:
\[
\begin{pmatrix}
a_{11} & a_{12} & a_{13} \\
a_{21} & a_{22} & a_{23} \\
a_{31} & a_{32} & a_{33}
\end{pmatrix}
\]

\section{Tables}
Below is a sample table showing data:

\begin{table}[h]
\centering
\begin{tabular}{@{}lrrr@{}}
\toprule
\textbf{Category} & \textbf{Value 1} & \textbf{Value 2} & \textbf{Total} \\
\midrule
Item A & 10.5 & 20.3 & 30.8 \\
Item B & 15.2 & 25.1 & 40.3 \\
Item C & 22.1 & 12.7 & 34.8 \\
\bottomrule
\end{tabular}
\caption{Sample Data Table}
\label{tab:sample}
\end{table}

\section{Lists}
\subsection{Bullet Points}
\begin{itemize}
  \item First item
  \item Second item
  \item Third item with sub-items:
    \begin{itemize}
      \item Sub-item 1
      \item Sub-item 2
    \end{itemize}
\end{itemize}

\subsection{Numbered List}
\begin{enumerate}
  \item First step
  \item Second step
  \item Third step
\end{enumerate}

\section{Citations and References}
When you use the citation search feature, references will be automatically managed. For example, you might cite a paper like this: \cite{example2024}. All citations will be automatically added to your \texttt{cite.bib} file.

\section{Conclusion}
This template demonstrates basic \LaTeX{} features including mathematical formulas, tables, and lists. The bibliography system is now set up to work seamlessly with Bibby AI's citation management. Feel free to modify it for your needs!

\bibliographystyle{plain}
\bibliography{cite}

\end{document}\documentclass{article}
\usepackage[utf8]{inputenc}
\usepackage{booktabs}
\usepackage{amsmath}
\usepackage{natbib}

\title{My Document}
\author{Your Name}
\date{\today}

\begin{document}

\maketitle

\section{Introduction}
This is a sample document demonstrating various \LaTeX{} features including mathematical equations, tables, and formatting. You can cite papers using the citation search feature, and they will be automatically added to your \texttt{cite.bib} file.

\section{Mathematical Expressions}
Here are some example equations:

\subsection{Quadratic Formula}
The quadratic formula is given by:
\[x = \frac{-b \pm \sqrt{b^2 - 4ac}}{2a}\]

\subsection{Matrix Example}
A sample matrix:
\[
\begin{pmatrix}
a_{11} & a_{12} & a_{13} \\
a_{21} & a_{22} & a_{23} \\
a_{31} & a_{32} & a_{33}
\end{pmatrix}
\]

\section{Tables}
Below is a sample table showing data:

\begin{table}[h]
\centering
\begin{tabular}{@{}lrrr@{}}
\toprule
\textbf{Category} & \textbf{Value 1} & \textbf{Value 2} & \textbf{Total} \\
\midrule
Item A & 10.5 & 20.3 & 30.8 \\
Item B & 15.2 & 25.1 & 40.3 \\
Item C & 22.1 & 12.7 & 34.8 \\
\bottomrule
\end{tabular}
\caption{Sample Data Table}
\label{tab:sample}
\end{table}

\section{Lists}
\subsection{Bullet Points}
\begin{itemize}
  \item First item
  \item Second item
  \item Third item with sub-items:
    \begin{itemize}
      \item Sub-item 1
      \item Sub-item 2
    \end{itemize}
\end{itemize}

\subsection{Numbered List}
\begin{enumerate}
  \item First step
  \item Second step
  \item Third step
\end{enumerate}

\section{Citations and References}
When you use the citation search feature, references will be automatically managed. For example, you might cite a paper like this: \cite{example2024}. All citations will be automatically added to your \texttt{cite.bib} file.

\section{Conclusion}
This template demonstrates basic \LaTeX{} features including mathematical formulas, tables, and lists. The bibliography system is now set up to work seamlessly with Bibby AI's citation management. Feel free to modify it for your needs!

\bibliographystyle{plain}
\bibliography{cite}

\end{document}\documentclass{article}
\usepackage[utf8]{inputenc}
\usepackage{booktabs}
\usepackage{amsmath}
\usepackage{natbib}

\title{My Document}
\author{Your Name}
\date{\today}

\begin{document}

\maketitle

\section{Introduction}
This is a sample document demonstrating various \LaTeX{} features including mathematical equations, tables, and formatting. You can cite papers using the citation search feature, and they will be automatically added to your \texttt{cite.bib} file.

\section{Mathematical Expressions}
Here are some example equations:

\subsection{Quadratic Formula}
The quadratic formula is given by:
\[x = \frac{-b \pm \sqrt{b^2 - 4ac}}{2a}\]

\subsection{Matrix Example}
A sample matrix:
\[
\begin{pmatrix}
a_{11} & a_{12} & a_{13} \\
a_{21} & a_{22} & a_{23} \\
a_{31} & a_{32} & a_{33}
\end{pmatrix}
\]

\section{Tables}
Below is a sample table showing data:

\begin{table}[h]
\centering
\begin{tabular}{@{}lrrr@{}}
\toprule
\textbf{Category} & \textbf{Value 1} & \textbf{Value 2} & \textbf{Total} \\
\midrule
Item A & 10.5 & 20.3 & 30.8 \\
Item B & 15.2 & 25.1 & 40.3 \\
Item C & 22.1 & 12.7 & 34.8 \\
\bottomrule
\end{tabular}
\caption{Sample Data Table}
\label{tab:sample}
\end{table}

\section{Lists}
\subsection{Bullet Points}
\begin{itemize}
  \item First item
  \item Second item
  \item Third item with sub-items:
    \begin{itemize}
      \item Sub-item 1
      \item Sub-item 2
    \end{itemize}
\end{itemize}

\subsection{Numbered List}
\begin{enumerate}
  \item First step
  \item Second step
  \item Third step
\end{enumerate}

\section{Citations and References}
When you use the citation search feature, references will be automatically managed. For example, you might cite a paper like this: \cite{example2024}. All citations will be automatically added to your \texttt{cite.bib} file.

\section{Conclusion}
This template demonstrates basic \LaTeX{} features including mathematical formulas, tables, and lists. The bibliography system is now set up to work seamlessly with Bibby AI's citation management. Feel free to modify it for your needs!

\bibliographystyle{plain}
\bibliography{cite}

\end{document}\documentclass{article}
\usepackage[utf8]{inputenc}
\usepackage{booktabs}
\usepackage{amsmath}
\usepackage{natbib}

\title{My Document}
\author{Your Name}
\date{\today}

\begin{document}

\maketitle

\section{Introduction}
This is a sample document demonstrating various \LaTeX{} features including mathematical equations, tables, and formatting. You can cite papers using the citation search feature, and they will be automatically added to your \texttt{cite.bib} file.

\section{Mathematical Expressions}
Here are some example equations:

\subsection{Quadratic Formula}
The quadratic formula is given by:
\[x = \frac{-b \pm \sqrt{b^2 - 4ac}}{2a}\]

\subsection{Matrix Example}
A sample matrix:
\[
\begin{pmatrix}
a_{11} & a_{12} & a_{13} \\
a_{21} & a_{22} & a_{23} \\
a_{31} & a_{32} & a_{33}
\end{pmatrix}
\]

\section{Tables}
Below is a sample table showing data:

\begin{table}[h]
\centering
\begin{tabular}{@{}lrrr@{}}
\toprule
\textbf{Category} & \textbf{Value 1} & \textbf{Value 2} & \textbf{Total} \\
\midrule
Item A & 10.5 & 20.3 & 30.8 \\
Item B & 15.2 & 25.1 & 40.3 \\
Item C & 22.1 & 12.7 & 34.8 \\
\bottomrule
\end{tabular}
\caption{Sample Data Table}
\label{tab:sample}
\end{table}

\section{Lists}
\subsection{Bullet Points}
\begin{itemize}
  \item First item
  \item Second item
  \item Third item with sub-items:
    \begin{itemize}
      \item Sub-item 1
      \item Sub-item 2
    \end{itemize}
\end{itemize}

\subsection{Numbered List}
\begin{enumerate}
  \item First step
  \item Second step
  \item Third step
\end{enumerate}

\section{Citations and References}
When you use the citation search feature, references will be automatically managed. For example, you might cite a paper like this: \cite{example2024}. All citations will be automatically added to your \texttt{cite.bib} file.

\section{Conclusion}
This template demonstrates basic \LaTeX{} features including mathematical formulas, tables, and lists. The bibliography system is now set up to work seamlessly with Bibby AI's citation management. Feel free to modify it for your needs!

\bibliographystyle{plain}
\bibliography{cite}

\end{document}\documentclass{article}
\usepackage[utf8]{inputenc}
\usepackage{booktabs}
\usepackage{amsmath}
\usepackage{natbib}

\title{My Document}
\author{Your Name}
\date{\today}

\begin{document}

\maketitle

\section{Introduction}
This is a sample document demonstrating various \LaTeX{} features including mathematical equations, tables, and formatting. You can cite papers using the citation search feature, and they will be automatically added to your \texttt{cite.bib} file.

\section{Mathematical Expressions}
Here are some example equations:

\subsection{Quadratic Formula}
The quadratic formula is given by:
\[x = \frac{-b \pm \sqrt{b^2 - 4ac}}{2a}\]

\subsection{Matrix Example}
A sample matrix:
\[
\begin{pmatrix}
a_{11} & a_{12} & a_{13} \\
a_{21} & a_{22} & a_{23} \\
a_{31} & a_{32} & a_{33}
\end{pmatrix}
\]

\section{Tables}
Below is a sample table showing data:

\begin{table}[h]
\centering
\begin{tabular}{@{}lrrr@{}}
\toprule
\textbf{Category} & \textbf{Value 1} & \textbf{Value 2} & \textbf{Total} \\
\midrule
Item A & 10.5 & 20.3 & 30.8 \\
Item B & 15.2 & 25.1 & 40.3 \\
Item C & 22.1 & 12.7 & 34.8 \\
\bottomrule
\end{tabular}
\caption{Sample Data Table}
\label{tab:sample}
\end{table}

\section{Lists}
\subsection{Bullet Points}
\begin{itemize}
  \item First item
  \item Second item
  \item Third item with sub-items:
    \begin{itemize}
      \item Sub-item 1
      \item Sub-item 2
    \end{itemize}
\end{itemize}

\subsection{Numbered List}
\begin{enumerate}
  \item First step
  \item Second step
  \item Third step
\end{enumerate}

\section{Citations and References}
When you use the citation search feature, references will be automatically managed. For example, you might cite a paper like this: \cite{example2024}. All citations will be automatically added to your \texttt{cite.bib} file.

\section{Conclusion}
This template demonstrates basic \LaTeX{} features including mathematical formulas, tables, and lists. The bibliography system is now set up to work seamlessly with Bibby AI's citation management. Feel free to modify it for your needs!

\bibliographystyle{plain}
\bibliography{cite}

\end{document}\documentclass{article}
\usepackage[utf8]{inputenc}
\usepackage{booktabs}
\usepackage{amsmath}
\usepackage{natbib}

\title{My Document}
\author{Your Name}
\date{\today}

\begin{document}

\maketitle

\section{Introduction}
This is a sample document demonstrating various \LaTeX{} features including mathematical equations, tables, and formatting. You can cite papers using the citation search feature, and they will be automatically added to your \texttt{cite.bib} file.

\section{Mathematical Expressions}
Here are some example equations:

\subsection{Quadratic Formula}
The quadratic formula is given by:
\[x = \frac{-b \pm \sqrt{b^2 - 4ac}}{2a}\]

\subsection{Matrix Example}
A sample matrix:
\[
\begin{pmatrix}
a_{11} & a_{12} & a_{13} \\
a_{21} & a_{22} & a_{23} \\
a_{31} & a_{32} & a_{33}
\end{pmatrix}
\]

\section{Tables}
Below is a sample table showing data:

\begin{table}[h]
\centering
\begin{tabular}{@{}lrrr@{}}
\toprule
\textbf{Category} & \textbf{Value 1} & \textbf{Value 2} & \textbf{Total} \\
\midrule
Item A & 10.5 & 20.3 & 30.8 \\
Item B & 15.2 & 25.1 & 40.3 \\
Item C & 22.1 & 12.7 & 34.8 \\
\bottomrule
\end{tabular}
\caption{Sample Data Table}
\label{tab:sample}
\end{table}

\section{Lists}
\subsection{Bullet Points}
\begin{itemize}
  \item First item
  \item Second item
  \item Third item with sub-items:
    \begin{itemize}
      \item Sub-item 1
      \item Sub-item 2
    \end{itemize}
\end{itemize}

\subsection{Numbered List}
\begin{enumerate}
  \item First step
  \item Second step
  \item Third step
\end{enumerate}

\section{Citations and References}
When you use the citation search feature, references will be automatically managed. For example, you might cite a paper like this: \cite{example2024}. All citations will be automatically added to your \texttt{cite.bib} file.

\section{Conclusion}
This template demonstrates basic \LaTeX{} features including mathematical formulas, tables, and lists. The bibliography system is now set up to work seamlessly with Bibby AI's citation management. Feel free to modify it for your needs!

\bibliographystyle{plain}
\bibliography{cite}

\end{document}\documentclass{article}
\usepackage[utf8]{inputenc}
\usepackage{booktabs}
\usepackage{amsmath}
\usepackage{natbib}

\title{My Document}
\author{Your Name}
\date{\today}

\begin{document}

\maketitle

\section{Introduction}
This is a sample document demonstrating various \LaTeX{} features including mathematical equations, tables, and formatting. You can cite papers using the citation search feature, and they will be automatically added to your \texttt{cite.bib} file.

\section{Mathematical Expressions}
Here are some example equations:

\subsection{Quadratic Formula}
The quadratic formula is given by:
\[x = \frac{-b \pm \sqrt{b^2 - 4ac}}{2a}\]

\subsection{Matrix Example}
A sample matrix:
\[
\begin{pmatrix}
a_{11} & a_{12} & a_{13} \\
a_{21} & a_{22} & a_{23} \\
a_{31} & a_{32} & a_{33}
\end{pmatrix}
\]

\section{Tables}
Below is a sample table showing data:

\begin{table}[h]
\centering
\begin{tabular}{@{}lrrr@{}}
\toprule
\textbf{Category} & \textbf{Value 1} & \textbf{Value 2} & \textbf{Total} \\
\midrule
Item A & 10.5 & 20.3 & 30.8 \\
Item B & 15.2 & 25.1 & 40.3 \\
Item C & 22.1 & 12.7 & 34.8 \\
\bottomrule
\end{tabular}
\caption{Sample Data Table}
\label{tab:sample}
\end{table}

\section{Lists}
\subsection{Bullet Points}
\begin{itemize}
  \item First item
  \item Second item
  \item Third item with sub-items:
    \begin{itemize}
      \item Sub-item 1
      \item Sub-item 2
    \end{itemize}
\end{itemize}

\subsection{Numbered List}
\begin{enumerate}
  \item First step
  \item Second step
  \item Third step
\end{enumerate}

\section{Citations and References}
When you use the citation search feature, references will be automatically managed. For example, you might cite a paper like this: \cite{example2024}. All citations will be automatically added to your \texttt{cite.bib} file.

\section{Conclusion}
This template demonstrates basic \LaTeX{} features including mathematical formulas, tables, and lists. The bibliography system is now set up to work seamlessly with Bibby AI's citation management. Feel free to modify it for your needs!

\bibliographystyle{plain}
\bibliography{cite}

\end{document}\documentclass{article}
\usepackage[utf8]{inputenc}
\usepackage{booktabs}
\usepackage{amsmath}
\usepackage{natbib}

\title{My Document}
\author{Your Name}
\date{\today}

\begin{document}

\maketitle

\section{Introduction}
This is a sample document demonstrating various \LaTeX{} features including mathematical equations, tables, and formatting. You can cite papers using the citation search feature, and they will be automatically added to your \texttt{cite.bib} file.

\section{Mathematical Expressions}
Here are some example equations:

\subsection{Quadratic Formula}
The quadratic formula is given by:
\[x = \frac{-b \pm \sqrt{b^2 - 4ac}}{2a}\]

\subsection{Matrix Example}
A sample matrix:
\[
\begin{pmatrix}
a_{11} & a_{12} & a_{13} \\
a_{21} & a_{22} & a_{23} \\
a_{31} & a_{32} & a_{33}
\end{pmatrix}
\]

\section{Tables}
Below is a sample table showing data:

\begin{table}[h]
\centering
\begin{tabular}{@{}lrrr@{}}
\toprule
\textbf{Category} & \textbf{Value 1} & \textbf{Value 2} & \textbf{Total} \\
\midrule
Item A & 10.5 & 20.3 & 30.8 \\
Item B & 15.2 & 25.1 & 40.3 \\
Item C & 22.1 & 12.7 & 34.8 \\
\bottomrule
\end{tabular}
\caption{Sample Data Table}
\label{tab:sample}
\end{table}

\section{Lists}
\subsection{Bullet Points}
\begin{itemize}
  \item First item
  \item Second item
  \item Third item with sub-items:
    \begin{itemize}
      \item Sub-item 1
      \item Sub-item 2
    \end{itemize}
\end{itemize}

\subsection{Numbered List}
\begin{enumerate}
  \item First step
  \item Second step
  \item Third step
\end{enumerate}

\section{Citations and References}
When you use the citation search feature, references will be automatically managed. For example, you might cite a paper like this: \cite{example2024}. All citations will be automatically added to your \texttt{cite.bib} file.

\section{Conclusion}
This template demonstrates basic \LaTeX{} features including mathematical formulas, tables, and lists. The bibliography system is now set up to work seamlessly with Bibby AI's citation management. Feel free to modify it for your needs!

\bibliographystyle{plain}
\bibliography{cite}

\end{document}\documentclass{article}
\usepackage[utf8]{inputenc}
\usepackage{booktabs}
\usepackage{amsmath}
\usepackage{natbib}

\title{My Document}
\author{Your Name}
\date{\today}

\begin{document}

\maketitle

\section{Introduction}
This is a sample document demonstrating various \LaTeX{} features including mathematical equations, tables, and formatting. You can cite papers using the citation search feature, and they will be automatically added to your \texttt{cite.bib} file.

\section{Mathematical Expressions}
Here are some example equations:

\subsection{Quadratic Formula}
The quadratic formula is given by:
\[x = \frac{-b \pm \sqrt{b^2 - 4ac}}{2a}\]

\subsection{Matrix Example}
A sample matrix:
\[
\begin{pmatrix}
a_{11} & a_{12} & a_{13} \\
a_{21} & a_{22} & a_{23} \\
a_{31} & a_{32} & a_{33}
\end{pmatrix}
\]

\section{Tables}
Below is a sample table showing data:

\begin{table}[h]
\centering
\begin{tabular}{@{}lrrr@{}}
\toprule
\textbf{Category} & \textbf{Value 1} & \textbf{Value 2} & \textbf{Total} \\
\midrule
Item A & 10.5 & 20.3 & 30.8 \\
Item B & 15.2 & 25.1 & 40.3 \\
Item C & 22.1 & 12.7 & 34.8 \\
\bottomrule
\end{tabular}
\caption{Sample Data Table}
\label{tab:sample}
\end{table}

\section{Lists}
\subsection{Bullet Points}
\begin{itemize}
  \item First item
  \item Second item
  \item Third item with sub-items:
    \begin{itemize}
      \item Sub-item 1
      \item Sub-item 2
    \end{itemize}
\end{itemize}

\subsection{Numbered List}
\begin{enumerate}
  \item First step
  \item Second step
  \item Third step
\end{enumerate}

\section{Citations and References}
When you use the citation search feature, references will be automatically managed. For example, you might cite a paper like this: \cite{example2024}. All citations will be automatically added to your \texttt{cite.bib} file.

\section{Conclusion}
This template demonstrates basic \LaTeX{} features including mathematical formulas, tables, and lists. The bibliography system is now set up to work seamlessly with Bibby AI's citation management. Feel free to modify it for your needs!

\bibliographystyle{plain}
\bibliography{cite}

\end{document}\documentclass{article}
\usepackage[utf8]{inputenc}
\usepackage{booktabs}
\usepackage{amsmath}
\usepackage{natbib}

\title{My Document}
\author{Your Name}
\date{\today}

\begin{document}

\maketitle

\section{Introduction}
This is a sample document demonstrating various \LaTeX{} features including mathematical equations, tables, and formatting. You can cite papers using the citation search feature, and they will be automatically added to your \texttt{cite.bib} file.

\section{Mathematical Expressions}
Here are some example equations:

\subsection{Quadratic Formula}
The quadratic formula is given by:
\[x = \frac{-b \pm \sqrt{b^2 - 4ac}}{2a}\]

\subsection{Matrix Example}
A sample matrix:
\[
\begin{pmatrix}
a_{11} & a_{12} & a_{13} \\
a_{21} & a_{22} & a_{23} \\
a_{31} & a_{32} & a_{33}
\end{pmatrix}
\]

\section{Tables}
Below is a sample table showing data:

\begin{table}[h]
\centering
\begin{tabular}{@{}lrrr@{}}
\toprule
\textbf{Category} & \textbf{Value 1} & \textbf{Value 2} & \textbf{Total} \\
\midrule
Item A & 10.5 & 20.3 & 30.8 \\
Item B & 15.2 & 25.1 & 40.3 \\
Item C & 22.1 & 12.7 & 34.8 \\
\bottomrule
\end{tabular}
\caption{Sample Data Table}
\label{tab:sample}
\end{table}

\section{Lists}
\subsection{Bullet Points}
\begin{itemize}
  \item First item
  \item Second item
  \item Third item with sub-items:
    \begin{itemize}
      \item Sub-item 1
      \item Sub-item 2
    \end{itemize}
\end{itemize}

\subsection{Numbered List}
\begin{enumerate}
  \item First step
  \item Second step
  \item Third step
\end{enumerate}

\section{Citations and References}
When you use the citation search feature, references will be automatically managed. For example, you might cite a paper like this: \cite{example2024}. All citations will be automatically added to your \texttt{cite.bib} file.

\section{Conclusion}
This template demonstrates basic \LaTeX{} features including mathematical formulas, tables, and lists. The bibliography system is now set up to work seamlessly with Bibby AI's citation management. Feel free to modify it for your needs!

\bibliographystyle{plain}
\bibliography{cite}

\end{document}\documentclass{article}
\usepackage[utf8]{inputenc}
\usepackage{booktabs}
\usepackage{amsmath}
\usepackage{natbib}

\title{My Document}
\author{Your Name}
\date{\today}

\begin{document}

\maketitle

\section{Introduction}
This is a sample document demonstrating various \LaTeX{} features including mathematical equations, tables, and formatting. You can cite papers using the citation search feature, and they will be automatically added to your \texttt{cite.bib} file.

\section{Mathematical Expressions}
Here are some example equations:

\subsection{Quadratic Formula}
The quadratic formula is given by:
\[x = \frac{-b \pm \sqrt{b^2 - 4ac}}{2a}\]

\subsection{Matrix Example}
A sample matrix:
\[
\begin{pmatrix}
a_{11} & a_{12} & a_{13} \\
a_{21} & a_{22} & a_{23} \\
a_{31} & a_{32} & a_{33}
\end{pmatrix}
\]

\section{Tables}
Below is a sample table showing data:

\begin{table}[h]
\centering
\begin{tabular}{@{}lrrr@{}}
\toprule
\textbf{Category} & \textbf{Value 1} & \textbf{Value 2} & \textbf{Total} \\
\midrule
Item A & 10.5 & 20.3 & 30.8 \\
Item B & 15.2 & 25.1 & 40.3 \\
Item C & 22.1 & 12.7 & 34.8 \\
\bottomrule
\end{tabular}
\caption{Sample Data Table}
\label{tab:sample}
\end{table}

\section{Lists}
\subsection{Bullet Points}
\begin{itemize}
  \item First item
  \item Second item
  \item Third item with sub-items:
    \begin{itemize}
      \item Sub-item 1
      \item Sub-item 2
    \end{itemize}
\end{itemize}

\subsection{Numbered List}
\begin{enumerate}
  \item First step
  \item Second step
  \item Third step
\end{enumerate}

\section{Citations and References}
When you use the citation search feature, references will be automatically managed. For example, you might cite a paper like this: \cite{example2024}. All citations will be automatically added to your \texttt{cite.bib} file.

\section{Conclusion}
This template demonstrates basic \LaTeX{} features including mathematical formulas, tables, and lists. The bibliography system is now set up to work seamlessly with Bibby AI's citation management. Feel free to modify it for your needs!

\bibliographystyle{plain}
\bibliography{cite}

\end{document}\documentclass{article}
\usepackage[utf8]{inputenc}
\usepackage{booktabs}
\usepackage{amsmath}
\usepackage{natbib}

\title{My Document}
\author{Your Name}
\date{\today}

\begin{document}

\maketitle

\section{Introduction}
This is a sample document demonstrating various \LaTeX{} features including mathematical equations, tables, and formatting. You can cite papers using the citation search feature, and they will be automatically added to your \texttt{cite.bib} file.

\section{Mathematical Expressions}
Here are some example equations:

\subsection{Quadratic Formula}
The quadratic formula is given by:
\[x = \frac{-b \pm \sqrt{b^2 - 4ac}}{2a}\]

\subsection{Matrix Example}
A sample matrix:
\[
\begin{pmatrix}
a_{11} & a_{12} & a_{13} \\
a_{21} & a_{22} & a_{23} \\
a_{31} & a_{32} & a_{33}
\end{pmatrix}
\]

\section{Tables}
Below is a sample table showing data:

\begin{table}[h]
\centering
\begin{tabular}{@{}lrrr@{}}
\toprule
\textbf{Category} & \textbf{Value 1} & \textbf{Value 2} & \textbf{Total} \\
\midrule
Item A & 10.5 & 20.3 & 30.8 \\
Item B & 15.2 & 25.1 & 40.3 \\
Item C & 22.1 & 12.7 & 34.8 \\
\bottomrule
\end{tabular}
\caption{Sample Data Table}
\label{tab:sample}
\end{table}

\section{Lists}
\subsection{Bullet Points}
\begin{itemize}
  \item First item
  \item Second item
  \item Third item with sub-items:
    \begin{itemize}
      \item Sub-item 1
      \item Sub-item 2
    \end{itemize}
\end{itemize}

\subsection{Numbered List}
\begin{enumerate}
  \item First step
  \item Second step
  \item Third step
\end{enumerate}

\section{Citations and References}
When you use the citation search feature, references will be automatically managed. For example, you might cite a paper like this: \cite{example2024}. All citations will be automatically added to your \texttt{cite.bib} file.

\section{Conclusion}
This template demonstrates basic \LaTeX{} features including mathematical formulas, tables, and lists. The bibliography system is now set up to work seamlessly with Bibby AI's citation management. Feel free to modify it for your needs!

\bibliographystyle{plain}
\bibliography{cite}

\end{document}\documentclass{article}
\usepackage[utf8]{inputenc}
\usepackage{booktabs}
\usepackage{amsmath}
\usepackage{natbib}

\title{My Document}
\author{Your Name}
\date{\today}

\begin{document}

\maketitle

\section{Introduction}
This is a sample document demonstrating various \LaTeX{} features including mathematical equations, tables, and formatting. You can cite papers using the citation search feature, and they will be automatically added to your \texttt{cite.bib} file.

\section{Mathematical Expressions}
Here are some example equations:

\subsection{Quadratic Formula}
The quadratic formula is given by:
\[x = \frac{-b \pm \sqrt{b^2 - 4ac}}{2a}\]

\subsection{Matrix Example}
A sample matrix:
\[
\begin{pmatrix}
a_{11} & a_{12} & a_{13} \\
a_{21} & a_{22} & a_{23} \\
a_{31} & a_{32} & a_{33}
\end{pmatrix}
\]

\section{Tables}
Below is a sample table showing data:

\begin{table}[h]
\centering
\begin{tabular}{@{}lrrr@{}}
\toprule
\textbf{Category} & \textbf{Value 1} & \textbf{Value 2} & \textbf{Total} \\
\midrule
Item A & 10.5 & 20.3 & 30.8 \\
Item B & 15.2 & 25.1 & 40.3 \\
Item C & 22.1 & 12.7 & 34.8 \\
\bottomrule
\end{tabular}
\caption{Sample Data Table}
\label{tab:sample}
\end{table}

\section{Lists}
\subsection{Bullet Points}
\begin{itemize}
  \item First item
  \item Second item
  \item Third item with sub-items:
    \begin{itemize}
      \item Sub-item 1
      \item Sub-item 2
    \end{itemize}
\end{itemize}

\subsection{Numbered List}
\begin{enumerate}
  \item First step
  \item Second step
  \item Third step
\end{enumerate}

\section{Citations and References}
When you use the citation search feature, references will be automatically managed. For example, you might cite a paper like this: \cite{example2024}. All citations will be automatically added to your \texttt{cite.bib} file.

\section{Conclusion}
This template demonstrates basic \LaTeX{} features including mathematical formulas, tables, and lists. The bibliography system is now set up to work seamlessly with Bibby AI's citation management. Feel free to modify it for your needs!

\bibliographystyle{plain}
\bibliography{cite}

\end{document}\documentclass{article}
\usepackage[utf8]{inputenc}
\usepackage{booktabs}
\usepackage{amsmath}
\usepackage{natbib}

\title{My Document}
\author{Your Name}
\date{\today}

\begin{document}

\maketitle

\section{Introduction}
This is a sample document demonstrating various \LaTeX{} features including mathematical equations, tables, and formatting. You can cite papers using the citation search feature, and they will be automatically added to your \texttt{cite.bib} file.

\section{Mathematical Expressions}
Here are some example equations:

\subsection{Quadratic Formula}
The quadratic formula is given by:
\[x = \frac{-b \pm \sqrt{b^2 - 4ac}}{2a}\]

\subsection{Matrix Example}
A sample matrix:
\[
\begin{pmatrix}
a_{11} & a_{12} & a_{13} \\
a_{21} & a_{22} & a_{23} \\
a_{31} & a_{32} & a_{33}
\end{pmatrix}
\]

\section{Tables}
Below is a sample table showing data:

\begin{table}[h]
\centering
\begin{tabular}{@{}lrrr@{}}
\toprule
\textbf{Category} & \textbf{Value 1} & \textbf{Value 2} & \textbf{Total} \\
\midrule
Item A & 10.5 & 20.3 & 30.8 \\
Item B & 15.2 & 25.1 & 40.3 \\
Item C & 22.1 & 12.7 & 34.8 \\
\bottomrule
\end{tabular}
\caption{Sample Data Table}
\label{tab:sample}
\end{table}

\section{Lists}
\subsection{Bullet Points}
\begin{itemize}
  \item First item
  \item Second item
  \item Third item with sub-items:
    \begin{itemize}
      \item Sub-item 1
      \item Sub-item 2
    \end{itemize}
\end{itemize}

\subsection{Numbered List}
\begin{enumerate}
  \item First step
  \item Second step
  \item Third step
\end{enumerate}

\section{Citations and References}
When you use the citation search feature, references will be automatically managed. For example, you might cite a paper like this: \cite{example2024}. All citations will be automatically added to your \texttt{cite.bib} file.

\section{Conclusion}
This template demonstrates basic \LaTeX{} features including mathematical formulas, tables, and lists. The bibliography system is now set up to work seamlessly with Bibby AI's citation management. Feel free to modify it for your needs!

\bibliographystyle{plain}
\bibliography{cite}

\end{document}\documentclass{article}
\usepackage[utf8]{inputenc}
\usepackage{booktabs}
\usepackage{amsmath}
\usepackage{natbib}

\title{My Document}
\author{Your Name}
\date{\today}

\begin{document}

\maketitle

\section{Introduction}
This is a sample document demonstrating various \LaTeX{} features including mathematical equations, tables, and formatting. You can cite papers using the citation search feature, and they will be automatically added to your \texttt{cite.bib} file.

\section{Mathematical Expressions}
Here are some example equations:

\subsection{Quadratic Formula}
The quadratic formula is given by:
\[x = \frac{-b \pm \sqrt{b^2 - 4ac}}{2a}\]

\subsection{Matrix Example}
A sample matrix:
\[
\begin{pmatrix}
a_{11} & a_{12} & a_{13} \\
a_{21} & a_{22} & a_{23} \\
a_{31} & a_{32} & a_{33}
\end{pmatrix}
\]

\section{Tables}
Below is a sample table showing data:

\begin{table}[h]
\centering
\begin{tabular}{@{}lrrr@{}}
\toprule
\textbf{Category} & \textbf{Value 1} & \textbf{Value 2} & \textbf{Total} \\
\midrule
Item A & 10.5 & 20.3 & 30.8 \\
Item B & 15.2 & 25.1 & 40.3 \\
Item C & 22.1 & 12.7 & 34.8 \\
\bottomrule
\end{tabular}
\caption{Sample Data Table}
\label{tab:sample}
\end{table}

\section{Lists}
\subsection{Bullet Points}
\begin{itemize}
  \item First item
  \item Second item
  \item Third item with sub-items:
    \begin{itemize}
      \item Sub-item 1
      \item Sub-item 2
    \end{itemize}
\end{itemize}

\subsection{Numbered List}
\begin{enumerate}
  \item First step
  \item Second step
  \item Third step
\end{enumerate}

\section{Citations and References}
When you use the citation search feature, references will be automatically managed. For example, you might cite a paper like this: \cite{example2024}. All citations will be automatically added to your \texttt{cite.bib} file.

\section{Conclusion}
This template demonstrates basic \LaTeX{} features including mathematical formulas, tables, and lists. The bibliography system is now set up to work seamlessly with Bibby AI's citation management. Feel free to modify it for your needs!

\bibliographystyle{plain}
\bibliography{cite}

\end{document}\documentclass{article}
\usepackage[utf8]{inputenc}
\usepackage{booktabs}
\usepackage{amsmath}
\usepackage{natbib}

\title{My Document}
\author{Your Name}
\date{\today}

\begin{document}

\maketitle

\section{Introduction}
This is a sample document demonstrating various \LaTeX{} features including mathematical equations, tables, and formatting. You can cite papers using the citation search feature, and they will be automatically added to your \texttt{cite.bib} file.

\section{Mathematical Expressions}
Here are some example equations:

\subsection{Quadratic Formula}
The quadratic formula is given by:
\[x = \frac{-b \pm \sqrt{b^2 - 4ac}}{2a}\]

\subsection{Matrix Example}
A sample matrix:
\[
\begin{pmatrix}
a_{11} & a_{12} & a_{13} \\
a_{21} & a_{22} & a_{23} \\
a_{31} & a_{32} & a_{33}
\end{pmatrix}
\]

\section{Tables}
Below is a sample table showing data:

\begin{table}[h]
\centering
\begin{tabular}{@{}lrrr@{}}
\toprule
\textbf{Category} & \textbf{Value 1} & \textbf{Value 2} & \textbf{Total} \\
\midrule
Item A & 10.5 & 20.3 & 30.8 \\
Item B & 15.2 & 25.1 & 40.3 \\
Item C & 22.1 & 12.7 & 34.8 \\
\bottomrule
\end{tabular}
\caption{Sample Data Table}
\label{tab:sample}
\end{table}

\section{Lists}
\subsection{Bullet Points}
\begin{itemize}
  \item First item
  \item Second item
  \item Third item with sub-items:
    \begin{itemize}
      \item Sub-item 1
      \item Sub-item 2
    \end{itemize}
\end{itemize}

\subsection{Numbered List}
\begin{enumerate}
  \item First step
  \item Second step
  \item Third step
\end{enumerate}

\section{Citations and References}
When you use the citation search feature, references will be automatically managed. For example, you might cite a paper like this: \cite{example2024}. All citations will be automatically added to your \texttt{cite.bib} file.

\section{Conclusion}
This template demonstrates basic \LaTeX{} features including mathematical formulas, tables, and lists. The bibliography system is now set up to work seamlessly with Bibby AI's citation management. Feel free to modify it for your needs!

\bibliographystyle{plain}
\bibliography{cite}

\end{document}\documentclass{article}
\usepackage[utf8]{inputenc}
\usepackage{booktabs}
\usepackage{amsmath}
\usepackage{natbib}

\title{My Document}
\author{Your Name}
\date{\today}

\begin{document}

\maketitle

\section{Introduction}
This is a sample document demonstrating various \LaTeX{} features including mathematical equations, tables, and formatting. You can cite papers using the citation search feature, and they will be automatically added to your \texttt{cite.bib} file.

\section{Mathematical Expressions}
Here are some example equations:

\subsection{Quadratic Formula}
The quadratic formula is given by:
\[x = \frac{-b \pm \sqrt{b^2 - 4ac}}{2a}\]

\subsection{Matrix Example}
A sample matrix:
\[
\begin{pmatrix}
a_{11} & a_{12} & a_{13} \\
a_{21} & a_{22} & a_{23} \\
a_{31} & a_{32} & a_{33}
\end{pmatrix}
\]

\section{Tables}
Below is a sample table showing data:

\begin{table}[h]
\centering
\begin{tabular}{@{}lrrr@{}}
\toprule
\textbf{Category} & \textbf{Value 1} & \textbf{Value 2} & \textbf{Total} \\
\midrule
Item A & 10.5 & 20.3 & 30.8 \\
Item B & 15.2 & 25.1 & 40.3 \\
Item C & 22.1 & 12.7 & 34.8 \\
\bottomrule
\end{tabular}
\caption{Sample Data Table}
\label{tab:sample}
\end{table}

\section{Lists}
\subsection{Bullet Points}
\begin{itemize}
  \item First item
  \item Second item
  \item Third item with sub-items:
    \begin{itemize}
      \item Sub-item 1
      \item Sub-item 2
    \end{itemize}
\end{itemize}

\subsection{Numbered List}
\begin{enumerate}
  \item First step
  \item Second step
  \item Third step
\end{enumerate}

\section{Citations and References}
When you use the citation search feature, references will be automatically managed. For example, you might cite a paper like this: \cite{example2024}. All citations will be automatically added to your \texttt{cite.bib} file.

\section{Conclusion}
This template demonstrates basic \LaTeX{} features including mathematical formulas, tables, and lists. The bibliography system is now set up to work seamlessly with Bibby AI's citation management. Feel free to modify it for your needs!

\bibliographystyle{plain}
\bibliography{cite}

\end{document}\documentclass{article}
\usepackage[utf8]{inputenc}
\usepackage{booktabs}
\usepackage{amsmath}
\usepackage{natbib}

\title{My Document}
\author{Your Name}
\date{\today}

\begin{document}

\maketitle

\section{Introduction}
This is a sample document demonstrating various \LaTeX{} features including mathematical equations, tables, and formatting. You can cite papers using the citation search feature, and they will be automatically added to your \texttt{cite.bib} file.

\section{Mathematical Expressions}
Here are some example equations:

\subsection{Quadratic Formula}
The quadratic formula is given by:
\[x = \frac{-b \pm \sqrt{b^2 - 4ac}}{2a}\]

\subsection{Matrix Example}
A sample matrix:
\[
\begin{pmatrix}
a_{11} & a_{12} & a_{13} \\
a_{21} & a_{22} & a_{23} \\
a_{31} & a_{32} & a_{33}
\end{pmatrix}
\]

\section{Tables}
Below is a sample table showing data:

\begin{table}[h]
\centering
\begin{tabular}{@{}lrrr@{}}
\toprule
\textbf{Category} & \textbf{Value 1} & \textbf{Value 2} & \textbf{Total} \\
\midrule
Item A & 10.5 & 20.3 & 30.8 \\
Item B & 15.2 & 25.1 & 40.3 \\
Item C & 22.1 & 12.7 & 34.8 \\
\bottomrule
\end{tabular}
\caption{Sample Data Table}
\label{tab:sample}
\end{table}

\section{Lists}
\subsection{Bullet Points}
\begin{itemize}
  \item First item
  \item Second item
  \item Third item with sub-items:
    \begin{itemize}
      \item Sub-item 1
      \item Sub-item 2
    \end{itemize}
\end{itemize}

\subsection{Numbered List}
\begin{enumerate}
  \item First step
  \item Second step
  \item Third step
\end{enumerate}

\section{Citations and References}
When you use the citation search feature, references will be automatically managed. For example, you might cite a paper like this: \cite{example2024}. All citations will be automatically added to your \texttt{cite.bib} file.

\section{Conclusion}
This template demonstrates basic \LaTeX{} features including mathematical formulas, tables, and lists. The bibliography system is now set up to work seamlessly with Bibby AI's citation management. Feel free to modify it for your needs!

\bibliographystyle{plain}
\bibliography{cite}

\end{document}\documentclass{article}
\usepackage[utf8]{inputenc}
\usepackage{booktabs}
\usepackage{amsmath}
\usepackage{natbib}

\title{My Document}
\author{Your Name}
\date{\today}

\begin{document}

\maketitle

\section{Introduction}
This is a sample document demonstrating various \LaTeX{} features including mathematical equations, tables, and formatting. You can cite papers using the citation search feature, and they will be automatically added to your \texttt{cite.bib} file.

\section{Mathematical Expressions}
Here are some example equations:

\subsection{Quadratic Formula}
The quadratic formula is given by:
\[x = \frac{-b \pm \sqrt{b^2 - 4ac}}{2a}\]

\subsection{Matrix Example}
A sample matrix:
\[
\begin{pmatrix}
a_{11} & a_{12} & a_{13} \\
a_{21} & a_{22} & a_{23} \\
a_{31} & a_{32} & a_{33}
\end{pmatrix}
\]

\section{Tables}
Below is a sample table showing data:

\begin{table}[h]
\centering
\begin{tabular}{@{}lrrr@{}}
\toprule
\textbf{Category} & \textbf{Value 1} & \textbf{Value 2} & \textbf{Total} \\
\midrule
Item A & 10.5 & 20.3 & 30.8 \\
Item B & 15.2 & 25.1 & 40.3 \\
Item C & 22.1 & 12.7 & 34.8 \\
\bottomrule
\end{tabular}
\caption{Sample Data Table}
\label{tab:sample}
\end{table}

\section{Lists}
\subsection{Bullet Points}
\begin{itemize}
  \item First item
  \item Second item
  \item Third item with sub-items:
    \begin{itemize}
      \item Sub-item 1
      \item Sub-item 2
    \end{itemize}
\end{itemize}

\subsection{Numbered List}
\begin{enumerate}
  \item First step
  \item Second step
  \item Third step
\end{enumerate}

\section{Citations and References}
When you use the citation search feature, references will be automatically managed. For example, you might cite a paper like this: \cite{example2024}. All citations will be automatically added to your \texttt{cite.bib} file.

\section{Conclusion}
This template demonstrates basic \LaTeX{} features including mathematical formulas, tables, and lists. The bibliography system is now set up to work seamlessly with Bibby AI's citation management. Feel free to modify it for your needs!

\bibliographystyle{plain}
\bibliography{cite}

\end{document}\documentclass{article}
\usepackage[utf8]{inputenc}
\usepackage{booktabs}
\usepackage{amsmath}
\usepackage{natbib}

\title{My Document}
\author{Your Name}
\date{\today}

\begin{document}

\maketitle

\section{Introduction}
This is a sample document demonstrating various \LaTeX{} features including mathematical equations, tables, and formatting. You can cite papers using the citation search feature, and they will be automatically added to your \texttt{cite.bib} file.

\section{Mathematical Expressions}
Here are some example equations:

\subsection{Quadratic Formula}
The quadratic formula is given by:
\[x = \frac{-b \pm \sqrt{b^2 - 4ac}}{2a}\]

\subsection{Matrix Example}
A sample matrix:
\[
\begin{pmatrix}
a_{11} & a_{12} & a_{13} \\
a_{21} & a_{22} & a_{23} \\
a_{31} & a_{32} & a_{33}
\end{pmatrix}
\]

\section{Tables}
Below is a sample table showing data:

\begin{table}[h]
\centering
\begin{tabular}{@{}lrrr@{}}
\toprule
\textbf{Category} & \textbf{Value 1} & \textbf{Value 2} & \textbf{Total} \\
\midrule
Item A & 10.5 & 20.3 & 30.8 \\
Item B & 15.2 & 25.1 & 40.3 \\
Item C & 22.1 & 12.7 & 34.8 \\
\bottomrule
\end{tabular}
\caption{Sample Data Table}
\label{tab:sample}
\end{table}

\section{Lists}
\subsection{Bullet Points}
\begin{itemize}
  \item First item
  \item Second item
  \item Third item with sub-items:
    \begin{itemize}
      \item Sub-item 1
      \item Sub-item 2
    \end{itemize}
\end{itemize}

\subsection{Numbered List}
\begin{enumerate}
  \item First step
  \item Second step
  \item Third step
\end{enumerate}

\section{Citations and References}
When you use the citation search feature, references will be automatically managed. For example, you might cite a paper like this: \cite{example2024}. All citations will be automatically added to your \texttt{cite.bib} file.

\section{Conclusion}
This template demonstrates basic \LaTeX{} features including mathematical formulas, tables, and lists. The bibliography system is now set up to work seamlessly with Bibby AI's citation management. Feel free to modify it for your needs!

\bibliographystyle{plain}
\bibliography{cite}

\end{document}\documentclass{article}
\usepackage[utf8]{inputenc}
\usepackage{booktabs}
\usepackage{amsmath}
\usepackage{natbib}

\title{My Document}
\author{Your Name}
\date{\today}

\begin{document}

\maketitle

\section{Introduction}
This is a sample document demonstrating various \LaTeX{} features including mathematical equations, tables, and formatting. You can cite papers using the citation search feature, and they will be automatically added to your \texttt{cite.bib} file.

\section{Mathematical Expressions}
Here are some example equations:

\subsection{Quadratic Formula}
The quadratic formula is given by:
\[x = \frac{-b \pm \sqrt{b^2 - 4ac}}{2a}\]

\subsection{Matrix Example}
A sample matrix:
\[
\begin{pmatrix}
a_{11} & a_{12} & a_{13} \\
a_{21} & a_{22} & a_{23} \\
a_{31} & a_{32} & a_{33}
\end{pmatrix}
\]

\section{Tables}
Below is a sample table showing data:

\begin{table}[h]
\centering
\begin{tabular}{@{}lrrr@{}}
\toprule
\textbf{Category} & \textbf{Value 1} & \textbf{Value 2} & \textbf{Total} \\
\midrule
Item A & 10.5 & 20.3 & 30.8 \\
Item B & 15.2 & 25.1 & 40.3 \\
Item C & 22.1 & 12.7 & 34.8 \\
\bottomrule
\end{tabular}
\caption{Sample Data Table}
\label{tab:sample}
\end{table}

\section{Lists}
\subsection{Bullet Points}
\begin{itemize}
  \item First item
  \item Second item
  \item Third item with sub-items:
    \begin{itemize}
      \item Sub-item 1
      \item Sub-item 2
    \end{itemize}
\end{itemize}

\subsection{Numbered List}
\begin{enumerate}
  \item First step
  \item Second step
  \item Third step
\end{enumerate}

\section{Citations and References}
When you use the citation search feature, references will be automatically managed. For example, you might cite a paper like this: \cite{example2024}. All citations will be automatically added to your \texttt{cite.bib} file.

\section{Conclusion}
This template demonstrates basic \LaTeX{} features including mathematical formulas, tables, and lists. The bibliography system is now set up to work seamlessly with Bibby AI's citation management. Feel free to modify it for your needs!

\bibliographystyle{plain}
\bibliography{cite}

\end{document}\documentclass{article}
\usepackage[utf8]{inputenc}
\usepackage{booktabs}
\usepackage{amsmath}
\usepackage{natbib}

\title{My Document}
\author{Your Name}
\date{\today}

\begin{document}

\maketitle

\section{Introduction}
This is a sample document demonstrating various \LaTeX{} features including mathematical equations, tables, and formatting. You can cite papers using the citation search feature, and they will be automatically added to your \texttt{cite.bib} file.

\section{Mathematical Expressions}
Here are some example equations:

\subsection{Quadratic Formula}
The quadratic formula is given by:
\[x = \frac{-b \pm \sqrt{b^2 - 4ac}}{2a}\]

\subsection{Matrix Example}
A sample matrix:
\[
\begin{pmatrix}
a_{11} & a_{12} & a_{13} \\
a_{21} & a_{22} & a_{23} \\
a_{31} & a_{32} & a_{33}
\end{pmatrix}
\]

\section{Tables}
Below is a sample table showing data:

\begin{table}[h]
\centering
\begin{tabular}{@{}lrrr@{}}
\toprule
\textbf{Category} & \textbf{Value 1} & \textbf{Value 2} & \textbf{Total} \\
\midrule
Item A & 10.5 & 20.3 & 30.8 \\
Item B & 15.2 & 25.1 & 40.3 \\
Item C & 22.1 & 12.7 & 34.8 \\
\bottomrule
\end{tabular}
\caption{Sample Data Table}
\label{tab:sample}
\end{table}

\section{Lists}
\subsection{Bullet Points}
\begin{itemize}
  \item First item
  \item Second item
  \item Third item with sub-items:
    \begin{itemize}
      \item Sub-item 1
      \item Sub-item 2
    \end{itemize}
\end{itemize}

\subsection{Numbered List}
\begin{enumerate}
  \item First step
  \item Second step
  \item Third step
\end{enumerate}

\section{Citations and References}
When you use the citation search feature, references will be automatically managed. For example, you might cite a paper like this: \cite{example2024}. All citations will be automatically added to your \texttt{cite.bib} file.

\section{Conclusion}
This template demonstrates basic \LaTeX{} features including mathematical formulas, tables, and lists. The bibliography system is now set up to work seamlessly with Bibby AI's citation management. Feel free to modify it for your needs!

\bibliographystyle{plain}
\bibliography{cite}

\end{document}\documentclass{article}
\usepackage[utf8]{inputenc}
\usepackage{booktabs}
\usepackage{amsmath}
\usepackage{natbib}

\title{My Document}
\author{Your Name}
\date{\today}

\begin{document}

\maketitle

\section{Introduction}
This is a sample document demonstrating various \LaTeX{} features including mathematical equations, tables, and formatting. You can cite papers using the citation search feature, and they will be automatically added to your \texttt{cite.bib} file.

\section{Mathematical Expressions}
Here are some example equations:

\subsection{Quadratic Formula}
The quadratic formula is given by:
\[x = \frac{-b \pm \sqrt{b^2 - 4ac}}{2a}\]

\subsection{Matrix Example}
A sample matrix:
\[
\begin{pmatrix}
a_{11} & a_{12} & a_{13} \\
a_{21} & a_{22} & a_{23} \\
a_{31} & a_{32} & a_{33}
\end{pmatrix}
\]

\section{Tables}
Below is a sample table showing data:

\begin{table}[h]
\centering
\begin{tabular}{@{}lrrr@{}}
\toprule
\textbf{Category} & \textbf{Value 1} & \textbf{Value 2} & \textbf{Total} \\
\midrule
Item A & 10.5 & 20.3 & 30.8 \\
Item B & 15.2 & 25.1 & 40.3 \\
Item C & 22.1 & 12.7 & 34.8 \\
\bottomrule
\end{tabular}
\caption{Sample Data Table}
\label{tab:sample}
\end{table}

\section{Lists}
\subsection{Bullet Points}
\begin{itemize}
  \item First item
  \item Second item
  \item Third item with sub-items:
    \begin{itemize}
      \item Sub-item 1
      \item Sub-item 2
    \end{itemize}
\end{itemize}

\subsection{Numbered List}
\begin{enumerate}
  \item First step
  \item Second step
  \item Third step
\end{enumerate}

\section{Citations and References}
When you use the citation search feature, references will be automatically managed. For example, you might cite a paper like this: \cite{example2024}. All citations will be automatically added to your \texttt{cite.bib} file.

\section{Conclusion}
This template demonstrates basic \LaTeX{} features including mathematical formulas, tables, and lists. The bibliography system is now set up to work seamlessly with Bibby AI's citation management. Feel free to modify it for your needs!

\bibliographystyle{plain}
\bibliography{cite}

\end{document}\documentclass{article}
\usepackage[utf8]{inputenc}
\usepackage{booktabs}
\usepackage{amsmath}
\usepackage{natbib}

\title{My Document}
\author{Your Name}
\date{\today}

\begin{document}

\maketitle

\section{Introduction}
This is a sample document demonstrating various \LaTeX{} features including mathematical equations, tables, and formatting. You can cite papers using the citation search feature, and they will be automatically added to your \texttt{cite.bib} file.

\section{Mathematical Expressions}
Here are some example equations:

\subsection{Quadratic Formula}
The quadratic formula is given by:
\[x = \frac{-b \pm \sqrt{b^2 - 4ac}}{2a}\]

\subsection{Matrix Example}
A sample matrix:
\[
\begin{pmatrix}
a_{11} & a_{12} & a_{13} \\
a_{21} & a_{22} & a_{23} \\
a_{31} & a_{32} & a_{33}
\end{pmatrix}
\]

\section{Tables}
Below is a sample table showing data:

\begin{table}[h]
\centering
\begin{tabular}{@{}lrrr@{}}
\toprule
\textbf{Category} & \textbf{Value 1} & \textbf{Value 2} & \textbf{Total} \\
\midrule
Item A & 10.5 & 20.3 & 30.8 \\
Item B & 15.2 & 25.1 & 40.3 \\
Item C & 22.1 & 12.7 & 34.8 \\
\bottomrule
\end{tabular}
\caption{Sample Data Table}
\label{tab:sample}
\end{table}

\section{Lists}
\subsection{Bullet Points}
\begin{itemize}
  \item First item
  \item Second item
  \item Third item with sub-items:
    \begin{itemize}
      \item Sub-item 1
      \item Sub-item 2
    \end{itemize}
\end{itemize}

\subsection{Numbered List}
\begin{enumerate}
  \item First step
  \item Second step
  \item Third step
\end{enumerate}

\section{Citations and References}
When you use the citation search feature, references will be automatically managed. For example, you might cite a paper like this: \cite{example2024}. All citations will be automatically added to your \texttt{cite.bib} file.

\section{Conclusion}
This template demonstrates basic \LaTeX{} features including mathematical formulas, tables, and lists. The bibliography system is now set up to work seamlessly with Bibby AI's citation management. Feel free to modify it for your needs!

\bibliographystyle{plain}
\bibliography{cite}

\end{document}\documentclass{article}
\usepackage[utf8]{inputenc}
\usepackage{booktabs}
\usepackage{amsmath}
\usepackage{natbib}

\title{My Document}
\author{Your Name}
\date{\today}

\begin{document}

\maketitle

\section{Introduction}
This is a sample document demonstrating various \LaTeX{} features including mathematical equations, tables, and formatting. You can cite papers using the citation search feature, and they will be automatically added to your \texttt{cite.bib} file.

\section{Mathematical Expressions}
Here are some example equations:

\subsection{Quadratic Formula}
The quadratic formula is given by:
\[x = \frac{-b \pm \sqrt{b^2 - 4ac}}{2a}\]

\subsection{Matrix Example}
A sample matrix:
\[
\begin{pmatrix}
a_{11} & a_{12} & a_{13} \\
a_{21} & a_{22} & a_{23} \\
a_{31} & a_{32} & a_{33}
\end{pmatrix}
\]

\section{Tables}
Below is a sample table showing data:

\begin{table}[h]
\centering
\begin{tabular}{@{}lrrr@{}}
\toprule
\textbf{Category} & \textbf{Value 1} & \textbf{Value 2} & \textbf{Total} \\
\midrule
Item A & 10.5 & 20.3 & 30.8 \\
Item B & 15.2 & 25.1 & 40.3 \\
Item C & 22.1 & 12.7 & 34.8 \\
\bottomrule
\end{tabular}
\caption{Sample Data Table}
\label{tab:sample}
\end{table}

\section{Lists}
\subsection{Bullet Points}
\begin{itemize}
  \item First item
  \item Second item
  \item Third item with sub-items:
    \begin{itemize}
      \item Sub-item 1
      \item Sub-item 2
    \end{itemize}
\end{itemize}

\subsection{Numbered List}
\begin{enumerate}
  \item First step
  \item Second step
  \item Third step
\end{enumerate}

\section{Citations and References}
When you use the citation search feature, references will be automatically managed. For example, you might cite a paper like this: \cite{example2024}. All citations will be automatically added to your \texttt{cite.bib} file.

\section{Conclusion}
This template demonstrates basic \LaTeX{} features including mathematical formulas, tables, and lists. The bibliography system is now set up to work seamlessly with Bibby AI's citation management. Feel free to modify it for your needs!

\bibliographystyle{plain}
\bibliography{cite}

\end{document}\documentclass{article}
\usepackage[utf8]{inputenc}
\usepackage{booktabs}
\usepackage{amsmath}
\usepackage{natbib}

\title{My Document}
\author{Your Name}
\date{\today}

\begin{document}

\maketitle

\section{Introduction}
This is a sample document demonstrating various \LaTeX{} features including mathematical equations, tables, and formatting. You can cite papers using the citation search feature, and they will be automatically added to your \texttt{cite.bib} file.

\section{Mathematical Expressions}
Here are some example equations:

\subsection{Quadratic Formula}
The quadratic formula is given by:
\[x = \frac{-b \pm \sqrt{b^2 - 4ac}}{2a}\]

\subsection{Matrix Example}
A sample matrix:
\[
\begin{pmatrix}
a_{11} & a_{12} & a_{13} \\
a_{21} & a_{22} & a_{23} \\
a_{31} & a_{32} & a_{33}
\end{pmatrix}
\]

\section{Tables}
Below is a sample table showing data:

\begin{table}[h]
\centering
\begin{tabular}{@{}lrrr@{}}
\toprule
\textbf{Category} & \textbf{Value 1} & \textbf{Value 2} & \textbf{Total} \\
\midrule
Item A & 10.5 & 20.3 & 30.8 \\
Item B & 15.2 & 25.1 & 40.3 \\
Item C & 22.1 & 12.7 & 34.8 \\
\bottomrule
\end{tabular}
\caption{Sample Data Table}
\label{tab:sample}
\end{table}

\section{Lists}
\subsection{Bullet Points}
\begin{itemize}
  \item First item
  \item Second item
  \item Third item with sub-items:
    \begin{itemize}
      \item Sub-item 1
      \item Sub-item 2
    \end{itemize}
\end{itemize}

\subsection{Numbered List}
\begin{enumerate}
  \item First step
  \item Second step
  \item Third step
\end{enumerate}

\section{Citations and References}
When you use the citation search feature, references will be automatically managed. For example, you might cite a paper like this: \cite{example2024}. All citations will be automatically added to your \texttt{cite.bib} file.

\section{Conclusion}
This template demonstrates basic \LaTeX{} features including mathematical formulas, tables, and lists. The bibliography system is now set up to work seamlessly with Bibby AI's citation management. Feel free to modify it for your needs!

\bibliographystyle{plain}
\bibliography{cite}

\end{document}\documentclass{article}
\usepackage[utf8]{inputenc}
\usepackage{booktabs}
\usepackage{amsmath}
\usepackage{natbib}

\title{My Document}
\author{Your Name}
\date{\today}

\begin{document}

\maketitle

\section{Introduction}
This is a sample document demonstrating various \LaTeX{} features including mathematical equations, tables, and formatting. You can cite papers using the citation search feature, and they will be automatically added to your \texttt{cite.bib} file.

\section{Mathematical Expressions}
Here are some example equations:

\subsection{Quadratic Formula}
The quadratic formula is given by:
\[x = \frac{-b \pm \sqrt{b^2 - 4ac}}{2a}\]

\subsection{Matrix Example}
A sample matrix:
\[
\begin{pmatrix}
a_{11} & a_{12} & a_{13} \\
a_{21} & a_{22} & a_{23} \\
a_{31} & a_{32} & a_{33}
\end{pmatrix}
\]

\section{Tables}
Below is a sample table showing data:

\begin{table}[h]
\centering
\begin{tabular}{@{}lrrr@{}}
\toprule
\textbf{Category} & \textbf{Value 1} & \textbf{Value 2} & \textbf{Total} \\
\midrule
Item A & 10.5 & 20.3 & 30.8 \\
Item B & 15.2 & 25.1 & 40.3 \\
Item C & 22.1 & 12.7 & 34.8 \\
\bottomrule
\end{tabular}
\caption{Sample Data Table}
\label{tab:sample}
\end{table}

\section{Lists}
\subsection{Bullet Points}
\begin{itemize}
  \item First item
  \item Second item
  \item Third item with sub-items:
    \begin{itemize}
      \item Sub-item 1
      \item Sub-item 2
    \end{itemize}
\end{itemize}

\subsection{Numbered List}
\begin{enumerate}
  \item First step
  \item Second step
  \item Third step
\end{enumerate}

\section{Citations and References}
When you use the citation search feature, references will be automatically managed. For example, you might cite a paper like this: \cite{example2024}. All citations will be automatically added to your \texttt{cite.bib} file.

\section{Conclusion}
This template demonstrates basic \LaTeX{} features including mathematical formulas, tables, and lists. The bibliography system is now set up to work seamlessly with Bibby AI's citation management. Feel free to modify it for your needs!

\bibliographystyle{plain}
\bibliography{cite}

\end{document}\documentclass{article}
\usepackage[utf8]{inputenc}
\usepackage{booktabs}
\usepackage{amsmath}
\usepackage{natbib}

\title{My Document}
\author{Your Name}
\date{\today}

\begin{document}

\maketitle

\section{Introduction}
This is a sample document demonstrating various \LaTeX{} features including mathematical equations, tables, and formatting. You can cite papers using the citation search feature, and they will be automatically added to your \texttt{cite.bib} file.

\section{Mathematical Expressions}
Here are some example equations:

\subsection{Quadratic Formula}
The quadratic formula is given by:
\[x = \frac{-b \pm \sqrt{b^2 - 4ac}}{2a}\]

\subsection{Matrix Example}
A sample matrix:
\[
\begin{pmatrix}
a_{11} & a_{12} & a_{13} \\
a_{21} & a_{22} & a_{23} \\
a_{31} & a_{32} & a_{33}
\end{pmatrix}
\]

\section{Tables}
Below is a sample table showing data:

\begin{table}[h]
\centering
\begin{tabular}{@{}lrrr@{}}
\toprule
\textbf{Category} & \textbf{Value 1} & \textbf{Value 2} & \textbf{Total} \\
\midrule
Item A & 10.5 & 20.3 & 30.8 \\
Item B & 15.2 & 25.1 & 40.3 \\
Item C & 22.1 & 12.7 & 34.8 \\
\bottomrule
\end{tabular}
\caption{Sample Data Table}
\label{tab:sample}
\end{table}

\section{Lists}
\subsection{Bullet Points}
\begin{itemize}
  \item First item
  \item Second item
  \item Third item with sub-items:
    \begin{itemize}
      \item Sub-item 1
      \item Sub-item 2
    \end{itemize}
\end{itemize}

\subsection{Numbered List}
\begin{enumerate}
  \item First step
  \item Second step
  \item Third step
\end{enumerate}

\section{Citations and References}
When you use the citation search feature, references will be automatically managed. For example, you might cite a paper like this: \cite{example2024}. All citations will be automatically added to your \texttt{cite.bib} file.

\section{Conclusion}
This template demonstrates basic \LaTeX{} features including mathematical formulas, tables, and lists. The bibliography system is now set up to work seamlessly with Bibby AI's citation management. Feel free to modify it for your needs!

\bibliographystyle{plain}
\bibliography{cite}

\end{document}\documentclass{article}
\usepackage[utf8]{inputenc}
\usepackage{booktabs}
\usepackage{amsmath}
\usepackage{natbib}

\title{My Document}
\author{Your Name}
\date{\today}

\begin{document}

\maketitle

\section{Introduction}
This is a sample document demonstrating various \LaTeX{} features including mathematical equations, tables, and formatting. You can cite papers using the citation search feature, and they will be automatically added to your \texttt{cite.bib} file.

\section{Mathematical Expressions}
Here are some example equations:

\subsection{Quadratic Formula}
The quadratic formula is given by:
\[x = \frac{-b \pm \sqrt{b^2 - 4ac}}{2a}\]

\subsection{Matrix Example}
A sample matrix:
\[
\begin{pmatrix}
a_{11} & a_{12} & a_{13} \\
a_{21} & a_{22} & a_{23} \\
a_{31} & a_{32} & a_{33}
\end{pmatrix}
\]

\section{Tables}
Below is a sample table showing data:

\begin{table}[h]
\centering
\begin{tabular}{@{}lrrr@{}}
\toprule
\textbf{Category} & \textbf{Value 1} & \textbf{Value 2} & \textbf{Total} \\
\midrule
Item A & 10.5 & 20.3 & 30.8 \\
Item B & 15.2 & 25.1 & 40.3 \\
Item C & 22.1 & 12.7 & 34.8 \\
\bottomrule
\end{tabular}
\caption{Sample Data Table}
\label{tab:sample}
\end{table}

\section{Lists}
\subsection{Bullet Points}
\begin{itemize}
  \item First item
  \item Second item
  \item Third item with sub-items:
    \begin{itemize}
      \item Sub-item 1
      \item Sub-item 2
    \end{itemize}
\end{itemize}

\subsection{Numbered List}
\begin{enumerate}
  \item First step
  \item Second step
  \item Third step
\end{enumerate}

\section{Citations and References}
When you use the citation search feature, references will be automatically managed. For example, you might cite a paper like this: \cite{example2024}. All citations will be automatically added to your \texttt{cite.bib} file.

\section{Conclusion}
This template demonstrates basic \LaTeX{} features including mathematical formulas, tables, and lists. The bibliography system is now set up to work seamlessly with Bibby AI's citation management. Feel free to modify it for your needs!

\bibliographystyle{plain}
\bibliography{cite}

\end{document}\documentclass{article}
\usepackage[utf8]{inputenc}
\usepackage{booktabs}
\usepackage{amsmath}
\usepackage{natbib}

\title{My Document}
\author{Your Name}
\date{\today}

\begin{document}

\maketitle

\section{Introduction}
This is a sample document demonstrating various \LaTeX{} features including mathematical equations, tables, and formatting. You can cite papers using the citation search feature, and they will be automatically added to your \texttt{cite.bib} file.

\section{Mathematical Expressions}
Here are some example equations:

\subsection{Quadratic Formula}
The quadratic formula is given by:
\[x = \frac{-b \pm \sqrt{b^2 - 4ac}}{2a}\]

\subsection{Matrix Example}
A sample matrix:
\[
\begin{pmatrix}
a_{11} & a_{12} & a_{13} \\
a_{21} & a_{22} & a_{23} \\
a_{31} & a_{32} & a_{33}
\end{pmatrix}
\]

\section{Tables}
Below is a sample table showing data:

\begin{table}[h]
\centering
\begin{tabular}{@{}lrrr@{}}
\toprule
\textbf{Category} & \textbf{Value 1} & \textbf{Value 2} & \textbf{Total} \\
\midrule
Item A & 10.5 & 20.3 & 30.8 \\
Item B & 15.2 & 25.1 & 40.3 \\
Item C & 22.1 & 12.7 & 34.8 \\
\bottomrule
\end{tabular}
\caption{Sample Data Table}
\label{tab:sample}
\end{table}

\section{Lists}
\subsection{Bullet Points}
\begin{itemize}
  \item First item
  \item Second item
  \item Third item with sub-items:
    \begin{itemize}
      \item Sub-item 1
      \item Sub-item 2
    \end{itemize}
\end{itemize}

\subsection{Numbered List}
\begin{enumerate}
  \item First step
  \item Second step
  \item Third step
\end{enumerate}

\section{Citations and References}
When you use the citation search feature, references will be automatically managed. For example, you might cite a paper like this: \cite{example2024}. All citations will be automatically added to your \texttt{cite.bib} file.

\section{Conclusion}
This template demonstrates basic \LaTeX{} features including mathematical formulas, tables, and lists. The bibliography system is now set up to work seamlessly with Bibby AI's citation management. Feel free to modify it for your needs!

\bibliographystyle{plain}
\bibliography{cite}

\end{document}\documentclass{article}
\usepackage[utf8]{inputenc}
\usepackage{booktabs}
\usepackage{amsmath}
\usepackage{natbib}

\title{My Document}
\author{Your Name}
\date{\today}

\begin{document}

\maketitle

\section{Introduction}
This is a sample document demonstrating various \LaTeX{} features including mathematical equations, tables, and formatting. You can cite papers using the citation search feature, and they will be automatically added to your \texttt{cite.bib} file.

\section{Mathematical Expressions}
Here are some example equations:

\subsection{Quadratic Formula}
The quadratic formula is given by:
\[x = \frac{-b \pm \sqrt{b^2 - 4ac}}{2a}\]

\subsection{Matrix Example}
A sample matrix:
\[
\begin{pmatrix}
a_{11} & a_{12} & a_{13} \\
a_{21} & a_{22} & a_{23} \\
a_{31} & a_{32} & a_{33}
\end{pmatrix}
\]

\section{Tables}
Below is a sample table showing data:

\begin{table}[h]
\centering
\begin{tabular}{@{}lrrr@{}}
\toprule
\textbf{Category} & \textbf{Value 1} & \textbf{Value 2} & \textbf{Total} \\
\midrule
Item A & 10.5 & 20.3 & 30.8 \\
Item B & 15.2 & 25.1 & 40.3 \\
Item C & 22.1 & 12.7 & 34.8 \\
\bottomrule
\end{tabular}
\caption{Sample Data Table}
\label{tab:sample}
\end{table}

\section{Lists}
\subsection{Bullet Points}
\begin{itemize}
  \item First item
  \item Second item
  \item Third item with sub-items:
    \begin{itemize}
      \item Sub-item 1
      \item Sub-item 2
    \end{itemize}
\end{itemize}

\subsection{Numbered List}
\begin{enumerate}
  \item First step
  \item Second step
  \item Third step
\end{enumerate}

\section{Citations and References}
When you use the citation search feature, references will be automatically managed. For example, you might cite a paper like this: \cite{example2024}. All citations will be automatically added to your \texttt{cite.bib} file.

\section{Conclusion}
This template demonstrates basic \LaTeX{} features including mathematical formulas, tables, and lists. The bibliography system is now set up to work seamlessly with Bibby AI's citation management. Feel free to modify it for your needs!

\bibliographystyle{plain}
\bibliography{cite}

\end{document}\documentclass{article}
\usepackage[utf8]{inputenc}
\usepackage{booktabs}
\usepackage{amsmath}
\usepackage{natbib}

\title{My Document}
\author{Your Name}
\date{\today}

\begin{document}

\maketitle

\section{Introduction}
This is a sample document demonstrating various \LaTeX{} features including mathematical equations, tables, and formatting. You can cite papers using the citation search feature, and they will be automatically added to your \texttt{cite.bib} file.

\section{Mathematical Expressions}
Here are some example equations:

\subsection{Quadratic Formula}
The quadratic formula is given by:
\[x = \frac{-b \pm \sqrt{b^2 - 4ac}}{2a}\]

\subsection{Matrix Example}
A sample matrix:
\[
\begin{pmatrix}
a_{11} & a_{12} & a_{13} \\
a_{21} & a_{22} & a_{23} \\
a_{31} & a_{32} & a_{33}
\end{pmatrix}
\]

\section{Tables}
Below is a sample table showing data:

\begin{table}[h]
\centering
\begin{tabular}{@{}lrrr@{}}
\toprule
\textbf{Category} & \textbf{Value 1} & \textbf{Value 2} & \textbf{Total} \\
\midrule
Item A & 10.5 & 20.3 & 30.8 \\
Item B & 15.2 & 25.1 & 40.3 \\
Item C & 22.1 & 12.7 & 34.8 \\
\bottomrule
\end{tabular}
\caption{Sample Data Table}
\label{tab:sample}
\end{table}

\section{Lists}
\subsection{Bullet Points}
\begin{itemize}
  \item First item
  \item Second item
  \item Third item with sub-items:
    \begin{itemize}
      \item Sub-item 1
      \item Sub-item 2
    \end{itemize}
\end{itemize}

\subsection{Numbered List}
\begin{enumerate}
  \item First step
  \item Second step
  \item Third step
\end{enumerate}

\section{Citations and References}
When you use the citation search feature, references will be automatically managed. For example, you might cite a paper like this: \cite{example2024}. All citations will be automatically added to your \texttt{cite.bib} file.

\section{Conclusion}
This template demonstrates basic \LaTeX{} features including mathematical formulas, tables, and lists. The bibliography system is now set up to work seamlessly with Bibby AI's citation management. Feel free to modify it for your needs!

\bibliographystyle{plain}
\bibliography{cite}

\end{document}\documentclass{article}
\usepackage[utf8]{inputenc}
\usepackage{booktabs}
\usepackage{amsmath}
\usepackage{natbib}

\title{My Document}
\author{Your Name}
\date{\today}

\begin{document}

\maketitle

\section{Introduction}
This is a sample document demonstrating various \LaTeX{} features including mathematical equations, tables, and formatting. You can cite papers using the citation search feature, and they will be automatically added to your \texttt{cite.bib} file.

\section{Mathematical Expressions}
Here are some example equations:

\subsection{Quadratic Formula}
The quadratic formula is given by:
\[x = \frac{-b \pm \sqrt{b^2 - 4ac}}{2a}\]

\subsection{Matrix Example}
A sample matrix:
\[
\begin{pmatrix}
a_{11} & a_{12} & a_{13} \\
a_{21} & a_{22} & a_{23} \\
a_{31} & a_{32} & a_{33}
\end{pmatrix}
\]

\section{Tables}
Below is a sample table showing data:

\begin{table}[h]
\centering
\begin{tabular}{@{}lrrr@{}}
\toprule
\textbf{Category} & \textbf{Value 1} & \textbf{Value 2} & \textbf{Total} \\
\midrule
Item A & 10.5 & 20.3 & 30.8 \\
Item B & 15.2 & 25.1 & 40.3 \\
Item C & 22.1 & 12.7 & 34.8 \\
\bottomrule
\end{tabular}
\caption{Sample Data Table}
\label{tab:sample}
\end{table}

\section{Lists}
\subsection{Bullet Points}
\begin{itemize}
  \item First item
  \item Second item
  \item Third item with sub-items:
    \begin{itemize}
      \item Sub-item 1
      \item Sub-item 2
    \end{itemize}
\end{itemize}

\subsection{Numbered List}
\begin{enumerate}
  \item First step
  \item Second step
  \item Third step
\end{enumerate}

\section{Citations and References}
When you use the citation search feature, references will be automatically managed. For example, you might cite a paper like this: \cite{example2024}. All citations will be automatically added to your \texttt{cite.bib} file.

\section{Conclusion}
This template demonstrates basic \LaTeX{} features including mathematical formulas, tables, and lists. The bibliography system is now set up to work seamlessly with Bibby AI's citation management. Feel free to modify it for your needs!

\bibliographystyle{plain}
\bibliography{cite}

\end{document}\documentclass{article}
\usepackage[utf8]{inputenc}
\usepackage{booktabs}
\usepackage{amsmath}
\usepackage{natbib}

\title{My Document}
\author{Your Name}
\date{\today}

\begin{document}

\maketitle

\section{Introduction}
This is a sample document demonstrating various \LaTeX{} features including mathematical equations, tables, and formatting. You can cite papers using the citation search feature, and they will be automatically added to your \texttt{cite.bib} file.

\section{Mathematical Expressions}
Here are some example equations:

\subsection{Quadratic Formula}
The quadratic formula is given by:
\[x = \frac{-b \pm \sqrt{b^2 - 4ac}}{2a}\]

\subsection{Matrix Example}
A sample matrix:
\[
\begin{pmatrix}
a_{11} & a_{12} & a_{13} \\
a_{21} & a_{22} & a_{23} \\
a_{31} & a_{32} & a_{33}
\end{pmatrix}
\]

\section{Tables}
Below is a sample table showing data:

\begin{table}[h]
\centering
\begin{tabular}{@{}lrrr@{}}
\toprule
\textbf{Category} & \textbf{Value 1} & \textbf{Value 2} & \textbf{Total} \\
\midrule
Item A & 10.5 & 20.3 & 30.8 \\
Item B & 15.2 & 25.1 & 40.3 \\
Item C & 22.1 & 12.7 & 34.8 \\
\bottomrule
\end{tabular}
\caption{Sample Data Table}
\label{tab:sample}
\end{table}

\section{Lists}
\subsection{Bullet Points}
\begin{itemize}
  \item First item
  \item Second item
  \item Third item with sub-items:
    \begin{itemize}
      \item Sub-item 1
      \item Sub-item 2
    \end{itemize}
\end{itemize}

\subsection{Numbered List}
\begin{enumerate}
  \item First step
  \item Second step
  \item Third step
\end{enumerate}

\section{Citations and References}
When you use the citation search feature, references will be automatically managed. For example, you might cite a paper like this: \cite{example2024}. All citations will be automatically added to your \texttt{cite.bib} file.

\section{Conclusion}
This template demonstrates basic \LaTeX{} features including mathematical formulas, tables, and lists. The bibliography system is now set up to work seamlessly with Bibby AI's citation management. Feel free to modify it for your needs!

\bibliographystyle{plain}
\bibliography{cite}

\end{document}\documentclass{article}
\usepackage[utf8]{inputenc}
\usepackage{booktabs}
\usepackage{amsmath}
\usepackage{natbib}

\title{My Document}
\author{Your Name}
\date{\today}

\begin{document}

\maketitle

\section{Introduction}
This is a sample document demonstrating various \LaTeX{} features including mathematical equations, tables, and formatting. You can cite papers using the citation search feature, and they will be automatically added to your \texttt{cite.bib} file.

\section{Mathematical Expressions}
Here are some example equations:

\subsection{Quadratic Formula}
The quadratic formula is given by:
\[x = \frac{-b \pm \sqrt{b^2 - 4ac}}{2a}\]

\subsection{Matrix Example}
A sample matrix:
\[
\begin{pmatrix}
a_{11} & a_{12} & a_{13} \\
a_{21} & a_{22} & a_{23} \\
a_{31} & a_{32} & a_{33}
\end{pmatrix}
\]

\section{Tables}
Below is a sample table showing data:

\begin{table}[h]
\centering
\begin{tabular}{@{}lrrr@{}}
\toprule
\textbf{Category} & \textbf{Value 1} & \textbf{Value 2} & \textbf{Total} \\
\midrule
Item A & 10.5 & 20.3 & 30.8 \\
Item B & 15.2 & 25.1 & 40.3 \\
Item C & 22.1 & 12.7 & 34.8 \\
\bottomrule
\end{tabular}
\caption{Sample Data Table}
\label{tab:sample}
\end{table}

\section{Lists}
\subsection{Bullet Points}
\begin{itemize}
  \item First item
  \item Second item
  \item Third item with sub-items:
    \begin{itemize}
      \item Sub-item 1
      \item Sub-item 2
    \end{itemize}
\end{itemize}

\subsection{Numbered List}
\begin{enumerate}
  \item First step
  \item Second step
  \item Third step
\end{enumerate}

\section{Citations and References}
When you use the citation search feature, references will be automatically managed. For example, you might cite a paper like this: \cite{example2024}. All citations will be automatically added to your \texttt{cite.bib} file.

\section{Conclusion}
This template demonstrates basic \LaTeX{} features including mathematical formulas, tables, and lists. The bibliography system is now set up to work seamlessly with Bibby AI's citation management. Feel free to modify it for your needs!

\bibliographystyle{plain}
\bibliography{cite}

\end{document}\documentclass{article}
\usepackage[utf8]{inputenc}
\usepackage{booktabs}
\usepackage{amsmath}
\usepackage{natbib}

\title{My Document}
\author{Your Name}
\date{\today}

\begin{document}

\maketitle

\section{Introduction}
This is a sample document demonstrating various \LaTeX{} features including mathematical equations, tables, and formatting. You can cite papers using the citation search feature, and they will be automatically added to your \texttt{cite.bib} file.

\section{Mathematical Expressions}
Here are some example equations:

\subsection{Quadratic Formula}
The quadratic formula is given by:
\[x = \frac{-b \pm \sqrt{b^2 - 4ac}}{2a}\]

\subsection{Matrix Example}
A sample matrix:
\[
\begin{pmatrix}
a_{11} & a_{12} & a_{13} \\
a_{21} & a_{22} & a_{23} \\
a_{31} & a_{32} & a_{33}
\end{pmatrix}
\]

\section{Tables}
Below is a sample table showing data:

\begin{table}[h]
\centering
\begin{tabular}{@{}lrrr@{}}
\toprule
\textbf{Category} & \textbf{Value 1} & \textbf{Value 2} & \textbf{Total} \\
\midrule
Item A & 10.5 & 20.3 & 30.8 \\
Item B & 15.2 & 25.1 & 40.3 \\
Item C & 22.1 & 12.7 & 34.8 \\
\bottomrule
\end{tabular}
\caption{Sample Data Table}
\label{tab:sample}
\end{table}

\section{Lists}
\subsection{Bullet Points}
\begin{itemize}
  \item First item
  \item Second item
  \item Third item with sub-items:
    \begin{itemize}
      \item Sub-item 1
      \item Sub-item 2
    \end{itemize}
\end{itemize}

\subsection{Numbered List}
\begin{enumerate}
  \item First step
  \item Second step
  \item Third step
\end{enumerate}

\section{Citations and References}
When you use the citation search feature, references will be automatically managed. For example, you might cite a paper like this: \cite{example2024}. All citations will be automatically added to your \texttt{cite.bib} file.

\section{Conclusion}
This template demonstrates basic \LaTeX{} features including mathematical formulas, tables, and lists. The bibliography system is now set up to work seamlessly with Bibby AI's citation management. Feel free to modify it for your needs!

\bibliographystyle{plain}
\bibliography{cite}

\end{document}\documentclass{article}
\usepackage[utf8]{inputenc}
\usepackage{booktabs}
\usepackage{amsmath}
\usepackage{natbib}

\title{My Document}
\author{Your Name}
\date{\today}

\begin{document}

\maketitle

\section{Introduction}
This is a sample document demonstrating various \LaTeX{} features including mathematical equations, tables, and formatting. You can cite papers using the citation search feature, and they will be automatically added to your \texttt{cite.bib} file.

\section{Mathematical Expressions}
Here are some example equations:

\subsection{Quadratic Formula}
The quadratic formula is given by:
\[x = \frac{-b \pm \sqrt{b^2 - 4ac}}{2a}\]

\subsection{Matrix Example}
A sample matrix:
\[
\begin{pmatrix}
a_{11} & a_{12} & a_{13} \\
a_{21} & a_{22} & a_{23} \\
a_{31} & a_{32} & a_{33}
\end{pmatrix}
\]

\section{Tables}
Below is a sample table showing data:

\begin{table}[h]
\centering
\begin{tabular}{@{}lrrr@{}}
\toprule
\textbf{Category} & \textbf{Value 1} & \textbf{Value 2} & \textbf{Total} \\
\midrule
Item A & 10.5 & 20.3 & 30.8 \\
Item B & 15.2 & 25.1 & 40.3 \\
Item C & 22.1 & 12.7 & 34.8 \\
\bottomrule
\end{tabular}
\caption{Sample Data Table}
\label{tab:sample}
\end{table}

\section{Lists}
\subsection{Bullet Points}
\begin{itemize}
  \item First item
  \item Second item
  \item Third item with sub-items:
    \begin{itemize}
      \item Sub-item 1
      \item Sub-item 2
    \end{itemize}
\end{itemize}

\subsection{Numbered List}
\begin{enumerate}
  \item First step
  \item Second step
  \item Third step
\end{enumerate}

\section{Citations and References}
When you use the citation search feature, references will be automatically managed. For example, you might cite a paper like this: \cite{example2024}. All citations will be automatically added to your \texttt{cite.bib} file.

\section{Conclusion}
This template demonstrates basic \LaTeX{} features including mathematical formulas, tables, and lists. The bibliography system is now set up to work seamlessly with Bibby AI's citation management. Feel free to modify it for your needs!

\bibliographystyle{plain}
\bibliography{cite}

\end{document}\documentclass{article}
\usepackage[utf8]{inputenc}
\usepackage{booktabs}
\usepackage{amsmath}
\usepackage{natbib}

\title{My Document}
\author{Your Name}
\date{\today}

\begin{document}

\maketitle

\section{Introduction}
This is a sample document demonstrating various \LaTeX{} features including mathematical equations, tables, and formatting. You can cite papers using the citation search feature, and they will be automatically added to your \texttt{cite.bib} file.

\section{Mathematical Expressions}
Here are some example equations:

\subsection{Quadratic Formula}
The quadratic formula is given by:
\[x = \frac{-b \pm \sqrt{b^2 - 4ac}}{2a}\]

\subsection{Matrix Example}
A sample matrix:
\[
\begin{pmatrix}
a_{11} & a_{12} & a_{13} \\
a_{21} & a_{22} & a_{23} \\
a_{31} & a_{32} & a_{33}
\end{pmatrix}
\]

\section{Tables}
Below is a sample table showing data:

\begin{table}[h]
\centering
\begin{tabular}{@{}lrrr@{}}
\toprule
\textbf{Category} & \textbf{Value 1} & \textbf{Value 2} & \textbf{Total} \\
\midrule
Item A & 10.5 & 20.3 & 30.8 \\
Item B & 15.2 & 25.1 & 40.3 \\
Item C & 22.1 & 12.7 & 34.8 \\
\bottomrule
\end{tabular}
\caption{Sample Data Table}
\label{tab:sample}
\end{table}

\section{Lists}
\subsection{Bullet Points}
\begin{itemize}
  \item First item
  \item Second item
  \item Third item with sub-items:
    \begin{itemize}
      \item Sub-item 1
      \item Sub-item 2
    \end{itemize}
\end{itemize}

\subsection{Numbered List}
\begin{enumerate}
  \item First step
  \item Second step
  \item Third step
\end{enumerate}

\section{Citations and References}
When you use the citation search feature, references will be automatically managed. For example, you might cite a paper like this: \cite{example2024}. All citations will be automatically added to your \texttt{cite.bib} file.

\section{Conclusion}
This template demonstrates basic \LaTeX{} features including mathematical formulas, tables, and lists. The bibliography system is now set up to work seamlessly with Bibby AI's citation management. Feel free to modify it for your needs!

\bibliographystyle{plain}
\bibliography{cite}

\end{document}\documentclass{article}
\usepackage[utf8]{inputenc}
\usepackage{booktabs}
\usepackage{amsmath}
\usepackage{natbib}

\title{My Document}
\author{Your Name}
\date{\today}

\begin{document}

\maketitle

\section{Introduction}
This is a sample document demonstrating various \LaTeX{} features including mathematical equations, tables, and formatting. You can cite papers using the citation search feature, and they will be automatically added to your \texttt{cite.bib} file.

\section{Mathematical Expressions}
Here are some example equations:

\subsection{Quadratic Formula}
The quadratic formula is given by:
\[x = \frac{-b \pm \sqrt{b^2 - 4ac}}{2a}\]

\subsection{Matrix Example}
A sample matrix:
\[
\begin{pmatrix}
a_{11} & a_{12} & a_{13} \\
a_{21} & a_{22} & a_{23} \\
a_{31} & a_{32} & a_{33}
\end{pmatrix}
\]

\section{Tables}
Below is a sample table showing data:

\begin{table}[h]
\centering
\begin{tabular}{@{}lrrr@{}}
\toprule
\textbf{Category} & \textbf{Value 1} & \textbf{Value 2} & \textbf{Total} \\
\midrule
Item A & 10.5 & 20.3 & 30.8 \\
Item B & 15.2 & 25.1 & 40.3 \\
Item C & 22.1 & 12.7 & 34.8 \\
\bottomrule
\end{tabular}
\caption{Sample Data Table}
\label{tab:sample}
\end{table}

\section{Lists}
\subsection{Bullet Points}
\begin{itemize}
  \item First item
  \item Second item
  \item Third item with sub-items:
    \begin{itemize}
      \item Sub-item 1
      \item Sub-item 2
    \end{itemize}
\end{itemize}

\subsection{Numbered List}
\begin{enumerate}
  \item First step
  \item Second step
  \item Third step
\end{enumerate}

\section{Citations and References}
When you use the citation search feature, references will be automatically managed. For example, you might cite a paper like this: \cite{example2024}. All citations will be automatically added to your \texttt{cite.bib} file.

\section{Conclusion}
This template demonstrates basic \LaTeX{} features including mathematical formulas, tables, and lists. The bibliography system is now set up to work seamlessly with Bibby AI's citation management. Feel free to modify it for your needs!

\bibliographystyle{plain}
\bibliography{cite}

\end{document}\documentclass{article}
\usepackage[utf8]{inputenc}
\usepackage{booktabs}
\usepackage{amsmath}
\usepackage{natbib}

\title{My Document}
\author{Your Name}
\date{\today}

\begin{document}

\maketitle

\section{Introduction}
This is a sample document demonstrating various \LaTeX{} features including mathematical equations, tables, and formatting. You can cite papers using the citation search feature, and they will be automatically added to your \texttt{cite.bib} file.

\section{Mathematical Expressions}
Here are some example equations:

\subsection{Quadratic Formula}
The quadratic formula is given by:
\[x = \frac{-b \pm \sqrt{b^2 - 4ac}}{2a}\]

\subsection{Matrix Example}
A sample matrix:
\[
\begin{pmatrix}
a_{11} & a_{12} & a_{13} \\
a_{21} & a_{22} & a_{23} \\
a_{31} & a_{32} & a_{33}
\end{pmatrix}
\]

\section{Tables}
Below is a sample table showing data:

\begin{table}[h]
\centering
\begin{tabular}{@{}lrrr@{}}
\toprule
\textbf{Category} & \textbf{Value 1} & \textbf{Value 2} & \textbf{Total} \\
\midrule
Item A & 10.5 & 20.3 & 30.8 \\
Item B & 15.2 & 25.1 & 40.3 \\
Item C & 22.1 & 12.7 & 34.8 \\
\bottomrule
\end{tabular}
\caption{Sample Data Table}
\label{tab:sample}
\end{table}

\section{Lists}
\subsection{Bullet Points}
\begin{itemize}
  \item First item
  \item Second item
  \item Third item with sub-items:
    \begin{itemize}
      \item Sub-item 1
      \item Sub-item 2
    \end{itemize}
\end{itemize}

\subsection{Numbered List}
\begin{enumerate}
  \item First step
  \item Second step
  \item Third step
\end{enumerate}

\section{Citations and References}
When you use the citation search feature, references will be automatically managed. For example, you might cite a paper like this: \cite{example2024}. All citations will be automatically added to your \texttt{cite.bib} file.

\section{Conclusion}
This template demonstrates basic \LaTeX{} features including mathematical formulas, tables, and lists. The bibliography system is now set up to work seamlessly with Bibby AI's citation management. Feel free to modify it for your needs!

\bibliographystyle{plain}
\bibliography{cite}

\end{document}\documentclass{article}
\usepackage[utf8]{inputenc}
\usepackage{booktabs}
\usepackage{amsmath}
\usepackage{natbib}

\title{My Document}
\author{Your Name}
\date{\today}

\begin{document}

\maketitle

\section{Introduction}
This is a sample document demonstrating various \LaTeX{} features including mathematical equations, tables, and formatting. You can cite papers using the citation search feature, and they will be automatically added to your \texttt{cite.bib} file.

\section{Mathematical Expressions}
Here are some example equations:

\subsection{Quadratic Formula}
The quadratic formula is given by:
\[x = \frac{-b \pm \sqrt{b^2 - 4ac}}{2a}\]

\subsection{Matrix Example}
A sample matrix:
\[
\begin{pmatrix}
a_{11} & a_{12} & a_{13} \\
a_{21} & a_{22} & a_{23} \\
a_{31} & a_{32} & a_{33}
\end{pmatrix}
\]

\section{Tables}
Below is a sample table showing data:

\begin{table}[h]
\centering
\begin{tabular}{@{}lrrr@{}}
\toprule
\textbf{Category} & \textbf{Value 1} & \textbf{Value 2} & \textbf{Total} \\
\midrule
Item A & 10.5 & 20.3 & 30.8 \\
Item B & 15.2 & 25.1 & 40.3 \\
Item C & 22.1 & 12.7 & 34.8 \\
\bottomrule
\end{tabular}
\caption{Sample Data Table}
\label{tab:sample}
\end{table}

\section{Lists}
\subsection{Bullet Points}
\begin{itemize}
  \item First item
  \item Second item
  \item Third item with sub-items:
    \begin{itemize}
      \item Sub-item 1
      \item Sub-item 2
    \end{itemize}
\end{itemize}

\subsection{Numbered List}
\begin{enumerate}
  \item First step
  \item Second step
  \item Third step
\end{enumerate}

\section{Citations and References}
When you use the citation search feature, references will be automatically managed. For example, you might cite a paper like this: \cite{example2024}. All citations will be automatically added to your \texttt{cite.bib} file.

\section{Conclusion}
This template demonstrates basic \LaTeX{} features including mathematical formulas, tables, and lists. The bibliography system is now set up to work seamlessly with Bibby AI's citation management. Feel free to modify it for your needs!

\bibliographystyle{plain}
\bibliography{cite}

\end{document}\documentclass{article}
\usepackage[utf8]{inputenc}
\usepackage{booktabs}
\usepackage{amsmath}
\usepackage{natbib}

\title{My Document}
\author{Your Name}
\date{\today}

\begin{document}

\maketitle

\section{Introduction}
This is a sample document demonstrating various \LaTeX{} features including mathematical equations, tables, and formatting. You can cite papers using the citation search feature, and they will be automatically added to your \texttt{cite.bib} file.

\section{Mathematical Expressions}
Here are some example equations:

\subsection{Quadratic Formula}
The quadratic formula is given by:
\[x = \frac{-b \pm \sqrt{b^2 - 4ac}}{2a}\]

\subsection{Matrix Example}
A sample matrix:
\[
\begin{pmatrix}
a_{11} & a_{12} & a_{13} \\
a_{21} & a_{22} & a_{23} \\
a_{31} & a_{32} & a_{33}
\end{pmatrix}
\]

\section{Tables}
Below is a sample table showing data:

\begin{table}[h]
\centering
\begin{tabular}{@{}lrrr@{}}
\toprule
\textbf{Category} & \textbf{Value 1} & \textbf{Value 2} & \textbf{Total} \\
\midrule
Item A & 10.5 & 20.3 & 30.8 \\
Item B & 15.2 & 25.1 & 40.3 \\
Item C & 22.1 & 12.7 & 34.8 \\
\bottomrule
\end{tabular}
\caption{Sample Data Table}
\label{tab:sample}
\end{table}

\section{Lists}
\subsection{Bullet Points}
\begin{itemize}
  \item First item
  \item Second item
  \item Third item with sub-items:
    \begin{itemize}
      \item Sub-item 1
      \item Sub-item 2
    \end{itemize}
\end{itemize}

\subsection{Numbered List}
\begin{enumerate}
  \item First step
  \item Second step
  \item Third step
\end{enumerate}

\section{Citations and References}
When you use the citation search feature, references will be automatically managed. For example, you might cite a paper like this: \cite{example2024}. All citations will be automatically added to your \texttt{cite.bib} file.

\section{Conclusion}
This template demonstrates basic \LaTeX{} features including mathematical formulas, tables, and lists. The bibliography system is now set up to work seamlessly with Bibby AI's citation management. Feel free to modify it for your needs!

\bibliographystyle{plain}
\bibliography{cite}

\end{document}\documentclass{article}
\usepackage[utf8]{inputenc}
\usepackage{booktabs}
\usepackage{amsmath}
\usepackage{natbib}

\title{My Document}
\author{Your Name}
\date{\today}

\begin{document}

\maketitle

\section{Introduction}
This is a sample document demonstrating various \LaTeX{} features including mathematical equations, tables, and formatting. You can cite papers using the citation search feature, and they will be automatically added to your \texttt{cite.bib} file.

\section{Mathematical Expressions}
Here are some example equations:

\subsection{Quadratic Formula}
The quadratic formula is given by:
\[x = \frac{-b \pm \sqrt{b^2 - 4ac}}{2a}\]

\subsection{Matrix Example}
A sample matrix:
\[
\begin{pmatrix}
a_{11} & a_{12} & a_{13} \\
a_{21} & a_{22} & a_{23} \\
a_{31} & a_{32} & a_{33}
\end{pmatrix}
\]

\section{Tables}
Below is a sample table showing data:

\begin{table}[h]
\centering
\begin{tabular}{@{}lrrr@{}}
\toprule
\textbf{Category} & \textbf{Value 1} & \textbf{Value 2} & \textbf{Total} \\
\midrule
Item A & 10.5 & 20.3 & 30.8 \\
Item B & 15.2 & 25.1 & 40.3 \\
Item C & 22.1 & 12.7 & 34.8 \\
\bottomrule
\end{tabular}
\caption{Sample Data Table}
\label{tab:sample}
\end{table}

\section{Lists}
\subsection{Bullet Points}
\begin{itemize}
  \item First item
  \item Second item
  \item Third item with sub-items:
    \begin{itemize}
      \item Sub-item 1
      \item Sub-item 2
    \end{itemize}
\end{itemize}

\subsection{Numbered List}
\begin{enumerate}
  \item First step
  \item Second step
  \item Third step
\end{enumerate}

\section{Citations and References}
When you use the citation search feature, references will be automatically managed. For example, you might cite a paper like this: \cite{example2024}. All citations will be automatically added to your \texttt{cite.bib} file.

\section{Conclusion}
This template demonstrates basic \LaTeX{} features including mathematical formulas, tables, and lists. The bibliography system is now set up to work seamlessly with Bibby AI's citation management. Feel free to modify it for your needs!

\bibliographystyle{plain}
\bibliography{cite}

\end{document}\documentclass{article}
\usepackage[utf8]{inputenc}
\usepackage{booktabs}
\usepackage{amsmath}
\usepackage{natbib}

\title{My Document}
\author{Your Name}
\date{\today}

\begin{document}

\maketitle

\section{Introduction}
This is a sample document demonstrating various \LaTeX{} features including mathematical equations, tables, and formatting. You can cite papers using the citation search feature, and they will be automatically added to your \texttt{cite.bib} file.

\section{Mathematical Expressions}
Here are some example equations:

\subsection{Quadratic Formula}
The quadratic formula is given by:
\[x = \frac{-b \pm \sqrt{b^2 - 4ac}}{2a}\]

\subsection{Matrix Example}
A sample matrix:
\[
\begin{pmatrix}
a_{11} & a_{12} & a_{13} \\
a_{21} & a_{22} & a_{23} \\
a_{31} & a_{32} & a_{33}
\end{pmatrix}
\]

\section{Tables}
Below is a sample table showing data:

\begin{table}[h]
\centering
\begin{tabular}{@{}lrrr@{}}
\toprule
\textbf{Category} & \textbf{Value 1} & \textbf{Value 2} & \textbf{Total} \\
\midrule
Item A & 10.5 & 20.3 & 30.8 \\
Item B & 15.2 & 25.1 & 40.3 \\
Item C & 22.1 & 12.7 & 34.8 \\
\bottomrule
\end{tabular}
\caption{Sample Data Table}
\label{tab:sample}
\end{table}

\section{Lists}
\subsection{Bullet Points}
\begin{itemize}
  \item First item
  \item Second item
  \item Third item with sub-items:
    \begin{itemize}
      \item Sub-item 1
      \item Sub-item 2
    \end{itemize}
\end{itemize}

\subsection{Numbered List}
\begin{enumerate}
  \item First step
  \item Second step
  \item Third step
\end{enumerate}

\section{Citations and References}
When you use the citation search feature, references will be automatically managed. For example, you might cite a paper like this: \cite{example2024}. All citations will be automatically added to your \texttt{cite.bib} file.

\section{Conclusion}
This template demonstrates basic \LaTeX{} features including mathematical formulas, tables, and lists. The bibliography system is now set up to work seamlessly with Bibby AI's citation management. Feel free to modify it for your needs!

\bibliographystyle{plain}
\bibliography{cite}

\end{document}\documentclass{article}
\usepackage[utf8]{inputenc}
\usepackage{booktabs}
\usepackage{amsmath}
\usepackage{natbib}

\title{My Document}
\author{Your Name}
\date{\today}

\begin{document}

\maketitle

\section{Introduction}
This is a sample document demonstrating various \LaTeX{} features including mathematical equations, tables, and formatting. You can cite papers using the citation search feature, and they will be automatically added to your \texttt{cite.bib} file.

\section{Mathematical Expressions}
Here are some example equations:

\subsection{Quadratic Formula}
The quadratic formula is given by:
\[x = \frac{-b \pm \sqrt{b^2 - 4ac}}{2a}\]

\subsection{Matrix Example}
A sample matrix:
\[
\begin{pmatrix}
a_{11} & a_{12} & a_{13} \\
a_{21} & a_{22} & a_{23} \\
a_{31} & a_{32} & a_{33}
\end{pmatrix}
\]

\section{Tables}
Below is a sample table showing data:

\begin{table}[h]
\centering
\begin{tabular}{@{}lrrr@{}}
\toprule
\textbf{Category} & \textbf{Value 1} & \textbf{Value 2} & \textbf{Total} \\
\midrule
Item A & 10.5 & 20.3 & 30.8 \\
Item B & 15.2 & 25.1 & 40.3 \\
Item C & 22.1 & 12.7 & 34.8 \\
\bottomrule
\end{tabular}
\caption{Sample Data Table}
\label{tab:sample}
\end{table}

\section{Lists}
\subsection{Bullet Points}
\begin{itemize}
  \item First item
  \item Second item
  \item Third item with sub-items:
    \begin{itemize}
      \item Sub-item 1
      \item Sub-item 2
    \end{itemize}
\end{itemize}

\subsection{Numbered List}
\begin{enumerate}
  \item First step
  \item Second step
  \item Third step
\end{enumerate}

\section{Citations and References}
When you use the citation search feature, references will be automatically managed. For example, you might cite a paper like this: \cite{example2024}. All citations will be automatically added to your \texttt{cite.bib} file.

\section{Conclusion}
This template demonstrates basic \LaTeX{} features including mathematical formulas, tables, and lists. The bibliography system is now set up to work seamlessly with Bibby AI's citation management. Feel free to modify it for your needs!

\bibliographystyle{plain}
\bibliography{cite}

\end{document}\documentclass{article}
\usepackage[utf8]{inputenc}
\usepackage{booktabs}
\usepackage{amsmath}
\usepackage{natbib}

\title{My Document}
\author{Your Name}
\date{\today}

\begin{document}

\maketitle

\section{Introduction}
This is a sample document demonstrating various \LaTeX{} features including mathematical equations, tables, and formatting. You can cite papers using the citation search feature, and they will be automatically added to your \texttt{cite.bib} file.

\section{Mathematical Expressions}
Here are some example equations:

\subsection{Quadratic Formula}
The quadratic formula is given by:
\[x = \frac{-b \pm \sqrt{b^2 - 4ac}}{2a}\]

\subsection{Matrix Example}
A sample matrix:
\[
\begin{pmatrix}
a_{11} & a_{12} & a_{13} \\
a_{21} & a_{22} & a_{23} \\
a_{31} & a_{32} & a_{33}
\end{pmatrix}
\]

\section{Tables}
Below is a sample table showing data:

\begin{table}[h]
\centering
\begin{tabular}{@{}lrrr@{}}
\toprule
\textbf{Category} & \textbf{Value 1} & \textbf{Value 2} & \textbf{Total} \\
\midrule
Item A & 10.5 & 20.3 & 30.8 \\
Item B & 15.2 & 25.1 & 40.3 \\
Item C & 22.1 & 12.7 & 34.8 \\
\bottomrule
\end{tabular}
\caption{Sample Data Table}
\label{tab:sample}
\end{table}

\section{Lists}
\subsection{Bullet Points}
\begin{itemize}
  \item First item
  \item Second item
  \item Third item with sub-items:
    \begin{itemize}
      \item Sub-item 1
      \item Sub-item 2
    \end{itemize}
\end{itemize}

\subsection{Numbered List}
\begin{enumerate}
  \item First step
  \item Second step
  \item Third step
\end{enumerate}

\section{Citations and References}
When you use the citation search feature, references will be automatically managed. For example, you might cite a paper like this: \cite{example2024}. All citations will be automatically added to your \texttt{cite.bib} file.

\section{Conclusion}
This template demonstrates basic \LaTeX{} features including mathematical formulas, tables, and lists. The bibliography system is now set up to work seamlessly with Bibby AI's citation management. Feel free to modify it for your needs!

\bibliographystyle{plain}
\bibliography{cite}

\end{document}\documentclass{article}
\usepackage[utf8]{inputenc}
\usepackage{booktabs}
\usepackage{amsmath}
\usepackage{natbib}

\title{My Document}
\author{Your Name}
\date{\today}

\begin{document}

\maketitle

\section{Introduction}
This is a sample document demonstrating various \LaTeX{} features including mathematical equations, tables, and formatting. You can cite papers using the citation search feature, and they will be automatically added to your \texttt{cite.bib} file.

\section{Mathematical Expressions}
Here are some example equations:

\subsection{Quadratic Formula}
The quadratic formula is given by:
\[x = \frac{-b \pm \sqrt{b^2 - 4ac}}{2a}\]

\subsection{Matrix Example}
A sample matrix:
\[
\begin{pmatrix}
a_{11} & a_{12} & a_{13} \\
a_{21} & a_{22} & a_{23} \\
a_{31} & a_{32} & a_{33}
\end{pmatrix}
\]

\section{Tables}
Below is a sample table showing data:

\begin{table}[h]
\centering
\begin{tabular}{@{}lrrr@{}}
\toprule
\textbf{Category} & \textbf{Value 1} & \textbf{Value 2} & \textbf{Total} \\
\midrule
Item A & 10.5 & 20.3 & 30.8 \\
Item B & 15.2 & 25.1 & 40.3 \\
Item C & 22.1 & 12.7 & 34.8 \\
\bottomrule
\end{tabular}
\caption{Sample Data Table}
\label{tab:sample}
\end{table}

\section{Lists}
\subsection{Bullet Points}
\begin{itemize}
  \item First item
  \item Second item
  \item Third item with sub-items:
    \begin{itemize}
      \item Sub-item 1
      \item Sub-item 2
    \end{itemize}
\end{itemize}

\subsection{Numbered List}
\begin{enumerate}
  \item First step
  \item Second step
  \item Third step
\end{enumerate}

\section{Citations and References}
When you use the citation search feature, references will be automatically managed. For example, you might cite a paper like this: \cite{example2024}. All citations will be automatically added to your \texttt{cite.bib} file.

\section{Conclusion}
This template demonstrates basic \LaTeX{} features including mathematical formulas, tables, and lists. The bibliography system is now set up to work seamlessly with Bibby AI's citation management. Feel free to modify it for your needs!

\bibliographystyle{plain}
\bibliography{cite}

\end{document}\documentclass{article}
\usepackage[utf8]{inputenc}
\usepackage{booktabs}
\usepackage{amsmath}
\usepackage{natbib}

\title{My Document}
\author{Your Name}
\date{\today}

\begin{document}

\maketitle

\section{Introduction}
This is a sample document demonstrating various \LaTeX{} features including mathematical equations, tables, and formatting. You can cite papers using the citation search feature, and they will be automatically added to your \texttt{cite.bib} file.

\section{Mathematical Expressions}
Here are some example equations:

\subsection{Quadratic Formula}
The quadratic formula is given by:
\[x = \frac{-b \pm \sqrt{b^2 - 4ac}}{2a}\]

\subsection{Matrix Example}
A sample matrix:
\[
\begin{pmatrix}
a_{11} & a_{12} & a_{13} \\
a_{21} & a_{22} & a_{23} \\
a_{31} & a_{32} & a_{33}
\end{pmatrix}
\]

\section{Tables}
Below is a sample table showing data:

\begin{table}[h]
\centering
\begin{tabular}{@{}lrrr@{}}
\toprule
\textbf{Category} & \textbf{Value 1} & \textbf{Value 2} & \textbf{Total} \\
\midrule
Item A & 10.5 & 20.3 & 30.8 \\
Item B & 15.2 & 25.1 & 40.3 \\
Item C & 22.1 & 12.7 & 34.8 \\
\bottomrule
\end{tabular}
\caption{Sample Data Table}
\label{tab:sample}
\end{table}

\section{Lists}
\subsection{Bullet Points}
\begin{itemize}
  \item First item
  \item Second item
  \item Third item with sub-items:
    \begin{itemize}
      \item Sub-item 1
      \item Sub-item 2
    \end{itemize}
\end{itemize}

\subsection{Numbered List}
\begin{enumerate}
  \item First step
  \item Second step
  \item Third step
\end{enumerate}

\section{Citations and References}
When you use the citation search feature, references will be automatically managed. For example, you might cite a paper like this: \cite{example2024}. All citations will be automatically added to your \texttt{cite.bib} file.

\section{Conclusion}
This template demonstrates basic \LaTeX{} features including mathematical formulas, tables, and lists. The bibliography system is now set up to work seamlessly with Bibby AI's citation management. Feel free to modify it for your needs!

\bibliographystyle{plain}
\bibliography{cite}

\end{document}\documentclass{article}
\usepackage[utf8]{inputenc}
\usepackage{booktabs}
\usepackage{amsmath}
\usepackage{natbib}

\title{My Document}
\author{Your Name}
\date{\today}

\begin{document}

\maketitle

\section{Introduction}
This is a sample document demonstrating various \LaTeX{} features including mathematical equations, tables, and formatting. You can cite papers using the citation search feature, and they will be automatically added to your \texttt{cite.bib} file.

\section{Mathematical Expressions}
Here are some example equations:

\subsection{Quadratic Formula}
The quadratic formula is given by:
\[x = \frac{-b \pm \sqrt{b^2 - 4ac}}{2a}\]

\subsection{Matrix Example}
A sample matrix:
\[
\begin{pmatrix}
a_{11} & a_{12} & a_{13} \\
a_{21} & a_{22} & a_{23} \\
a_{31} & a_{32} & a_{33}
\end{pmatrix}
\]

\section{Tables}
Below is a sample table showing data:

\begin{table}[h]
\centering
\begin{tabular}{@{}lrrr@{}}
\toprule
\textbf{Category} & \textbf{Value 1} & \textbf{Value 2} & \textbf{Total} \\
\midrule
Item A & 10.5 & 20.3 & 30.8 \\
Item B & 15.2 & 25.1 & 40.3 \\
Item C & 22.1 & 12.7 & 34.8 \\
\bottomrule
\end{tabular}
\caption{Sample Data Table}
\label{tab:sample}
\end{table}

\section{Lists}
\subsection{Bullet Points}
\begin{itemize}
  \item First item
  \item Second item
  \item Third item with sub-items:
    \begin{itemize}
      \item Sub-item 1
      \item Sub-item 2
    \end{itemize}
\end{itemize}

\subsection{Numbered List}
\begin{enumerate}
  \item First step
  \item Second step
  \item Third step
\end{enumerate}

\section{Citations and References}
When you use the citation search feature, references will be automatically managed. For example, you might cite a paper like this: \cite{example2024}. All citations will be automatically added to your \texttt{cite.bib} file.

\section{Conclusion}
This template demonstrates basic \LaTeX{} features including mathematical formulas, tables, and lists. The bibliography system is now set up to work seamlessly with Bibby AI's citation management. Feel free to modify it for your needs!

\bibliographystyle{plain}
\bibliography{cite}

\end{document}\documentclass{article}
\usepackage[utf8]{inputenc}
\usepackage{booktabs}
\usepackage{amsmath}
\usepackage{natbib}

\title{My Document}
\author{Your Name}
\date{\today}

\begin{document}

\maketitle

\section{Introduction}
This is a sample document demonstrating various \LaTeX{} features including mathematical equations, tables, and formatting. You can cite papers using the citation search feature, and they will be automatically added to your \texttt{cite.bib} file.

\section{Mathematical Expressions}
Here are some example equations:

\subsection{Quadratic Formula}
The quadratic formula is given by:
\[x = \frac{-b \pm \sqrt{b^2 - 4ac}}{2a}\]

\subsection{Matrix Example}
A sample matrix:
\[
\begin{pmatrix}
a_{11} & a_{12} & a_{13} \\
a_{21} & a_{22} & a_{23} \\
a_{31} & a_{32} & a_{33}
\end{pmatrix}
\]

\section{Tables}
Below is a sample table showing data:

\begin{table}[h]
\centering
\begin{tabular}{@{}lrrr@{}}
\toprule
\textbf{Category} & \textbf{Value 1} & \textbf{Value 2} & \textbf{Total} \\
\midrule
Item A & 10.5 & 20.3 & 30.8 \\
Item B & 15.2 & 25.1 & 40.3 \\
Item C & 22.1 & 12.7 & 34.8 \\
\bottomrule
\end{tabular}
\caption{Sample Data Table}
\label{tab:sample}
\end{table}

\section{Lists}
\subsection{Bullet Points}
\begin{itemize}
  \item First item
  \item Second item
  \item Third item with sub-items:
    \begin{itemize}
      \item Sub-item 1
      \item Sub-item 2
    \end{itemize}
\end{itemize}

\subsection{Numbered List}
\begin{enumerate}
  \item First step
  \item Second step
  \item Third step
\end{enumerate}

\section{Citations and References}
When you use the citation search feature, references will be automatically managed. For example, you might cite a paper like this: \cite{example2024}. All citations will be automatically added to your \texttt{cite.bib} file.

\section{Conclusion}
This template demonstrates basic \LaTeX{} features including mathematical formulas, tables, and lists. The bibliography system is now set up to work seamlessly with Bibby AI's citation management. Feel free to modify it for your needs!

\bibliographystyle{plain}
\bibliography{cite}

\end{document}\documentclass{article}
\usepackage[utf8]{inputenc}
\usepackage{booktabs}
\usepackage{amsmath}
\usepackage{natbib}

\title{My Document}
\author{Your Name}
\date{\today}

\begin{document}

\maketitle

\section{Introduction}
This is a sample document demonstrating various \LaTeX{} features including mathematical equations, tables, and formatting. You can cite papers using the citation search feature, and they will be automatically added to your \texttt{cite.bib} file.

\section{Mathematical Expressions}
Here are some example equations:

\subsection{Quadratic Formula}
The quadratic formula is given by:
\[x = \frac{-b \pm \sqrt{b^2 - 4ac}}{2a}\]

\subsection{Matrix Example}
A sample matrix:
\[
\begin{pmatrix}
a_{11} & a_{12} & a_{13} \\
a_{21} & a_{22} & a_{23} \\
a_{31} & a_{32} & a_{33}
\end{pmatrix}
\]

\section{Tables}
Below is a sample table showing data:

\begin{table}[h]
\centering
\begin{tabular}{@{}lrrr@{}}
\toprule
\textbf{Category} & \textbf{Value 1} & \textbf{Value 2} & \textbf{Total} \\
\midrule
Item A & 10.5 & 20.3 & 30.8 \\
Item B & 15.2 & 25.1 & 40.3 \\
Item C & 22.1 & 12.7 & 34.8 \\
\bottomrule
\end{tabular}
\caption{Sample Data Table}
\label{tab:sample}
\end{table}

\section{Lists}
\subsection{Bullet Points}
\begin{itemize}
  \item First item
  \item Second item
  \item Third item with sub-items:
    \begin{itemize}
      \item Sub-item 1
      \item Sub-item 2
    \end{itemize}
\end{itemize}

\subsection{Numbered List}
\begin{enumerate}
  \item First step
  \item Second step
  \item Third step
\end{enumerate}

\section{Citations and References}
When you use the citation search feature, references will be automatically managed. For example, you might cite a paper like this: \cite{example2024}. All citations will be automatically added to your \texttt{cite.bib} file.

\section{Conclusion}
This template demonstrates basic \LaTeX{} features including mathematical formulas, tables, and lists. The bibliography system is now set up to work seamlessly with Bibby AI's citation management. Feel free to modify it for your needs!

\bibliographystyle{plain}
\bibliography{cite}

\end{document}\documentclass{article}
\usepackage[utf8]{inputenc}
\usepackage{booktabs}
\usepackage{amsmath}
\usepackage{natbib}

\title{My Document}
\author{Your Name}
\date{\today}

\begin{document}

\maketitle

\section{Introduction}
This is a sample document demonstrating various \LaTeX{} features including mathematical equations, tables, and formatting. You can cite papers using the citation search feature, and they will be automatically added to your \texttt{cite.bib} file.

\section{Mathematical Expressions}
Here are some example equations:

\subsection{Quadratic Formula}
The quadratic formula is given by:
\[x = \frac{-b \pm \sqrt{b^2 - 4ac}}{2a}\]

\subsection{Matrix Example}
A sample matrix:
\[
\begin{pmatrix}
a_{11} & a_{12} & a_{13} \\
a_{21} & a_{22} & a_{23} \\
a_{31} & a_{32} & a_{33}
\end{pmatrix}
\]

\section{Tables}
Below is a sample table showing data:

\begin{table}[h]
\centering
\begin{tabular}{@{}lrrr@{}}
\toprule
\textbf{Category} & \textbf{Value 1} & \textbf{Value 2} & \textbf{Total} \\
\midrule
Item A & 10.5 & 20.3 & 30.8 \\
Item B & 15.2 & 25.1 & 40.3 \\
Item C & 22.1 & 12.7 & 34.8 \\
\bottomrule
\end{tabular}
\caption{Sample Data Table}
\label{tab:sample}
\end{table}

\section{Lists}
\subsection{Bullet Points}
\begin{itemize}
  \item First item
  \item Second item
  \item Third item with sub-items:
    \begin{itemize}
      \item Sub-item 1
      \item Sub-item 2
    \end{itemize}
\end{itemize}

\subsection{Numbered List}
\begin{enumerate}
  \item First step
  \item Second step
  \item Third step
\end{enumerate}

\section{Citations and References}
When you use the citation search feature, references will be automatically managed. For example, you might cite a paper like this: \cite{example2024}. All citations will be automatically added to your \texttt{cite.bib} file.

\section{Conclusion}
This template demonstrates basic \LaTeX{} features including mathematical formulas, tables, and lists. The bibliography system is now set up to work seamlessly with Bibby AI's citation management. Feel free to modify it for your needs!

\bibliographystyle{plain}
\bibliography{cite}

\end{document}\documentclass{article}
\usepackage[utf8]{inputenc}
\usepackage{booktabs}
\usepackage{amsmath}
\usepackage{natbib}

\title{My Document}
\author{Your Name}
\date{\today}

\begin{document}

\maketitle

\section{Introduction}
This is a sample document demonstrating various \LaTeX{} features including mathematical equations, tables, and formatting. You can cite papers using the citation search feature, and they will be automatically added to your \texttt{cite.bib} file.

\section{Mathematical Expressions}
Here are some example equations:

\subsection{Quadratic Formula}
The quadratic formula is given by:
\[x = \frac{-b \pm \sqrt{b^2 - 4ac}}{2a}\]

\subsection{Matrix Example}
A sample matrix:
\[
\begin{pmatrix}
a_{11} & a_{12} & a_{13} \\
a_{21} & a_{22} & a_{23} \\
a_{31} & a_{32} & a_{33}
\end{pmatrix}
\]

\section{Tables}
Below is a sample table showing data:

\begin{table}[h]
\centering
\begin{tabular}{@{}lrrr@{}}
\toprule
\textbf{Category} & \textbf{Value 1} & \textbf{Value 2} & \textbf{Total} \\
\midrule
Item A & 10.5 & 20.3 & 30.8 \\
Item B & 15.2 & 25.1 & 40.3 \\
Item C & 22.1 & 12.7 & 34.8 \\
\bottomrule
\end{tabular}
\caption{Sample Data Table}
\label{tab:sample}
\end{table}

\section{Lists}
\subsection{Bullet Points}
\begin{itemize}
  \item First item
  \item Second item
  \item Third item with sub-items:
    \begin{itemize}
      \item Sub-item 1
      \item Sub-item 2
    \end{itemize}
\end{itemize}

\subsection{Numbered List}
\begin{enumerate}
  \item First step
  \item Second step
  \item Third step
\end{enumerate}

\section{Citations and References}
When you use the citation search feature, references will be automatically managed. For example, you might cite a paper like this: \cite{example2024}. All citations will be automatically added to your \texttt{cite.bib} file.

\section{Conclusion}
This template demonstrates basic \LaTeX{} features including mathematical formulas, tables, and lists. The bibliography system is now set up to work seamlessly with Bibby AI's citation management. Feel free to modify it for your needs!

\bibliographystyle{plain}
\bibliography{cite}

\end{document}\documentclass{article}
\usepackage[utf8]{inputenc}
\usepackage{booktabs}
\usepackage{amsmath}
\usepackage{natbib}

\title{My Document}
\author{Your Name}
\date{\today}

\begin{document}

\maketitle

\section{Introduction}
This is a sample document demonstrating various \LaTeX{} features including mathematical equations, tables, and formatting. You can cite papers using the citation search feature, and they will be automatically added to your \texttt{cite.bib} file.

\section{Mathematical Expressions}
Here are some example equations:

\subsection{Quadratic Formula}
The quadratic formula is given by:
\[x = \frac{-b \pm \sqrt{b^2 - 4ac}}{2a}\]

\subsection{Matrix Example}
A sample matrix:
\[
\begin{pmatrix}
a_{11} & a_{12} & a_{13} \\
a_{21} & a_{22} & a_{23} \\
a_{31} & a_{32} & a_{33}
\end{pmatrix}
\]

\section{Tables}
Below is a sample table showing data:

\begin{table}[h]
\centering
\begin{tabular}{@{}lrrr@{}}
\toprule
\textbf{Category} & \textbf{Value 1} & \textbf{Value 2} & \textbf{Total} \\
\midrule
Item A & 10.5 & 20.3 & 30.8 \\
Item B & 15.2 & 25.1 & 40.3 \\
Item C & 22.1 & 12.7 & 34.8 \\
\bottomrule
\end{tabular}
\caption{Sample Data Table}
\label{tab:sample}
\end{table}

\section{Lists}
\subsection{Bullet Points}
\begin{itemize}
  \item First item
  \item Second item
  \item Third item with sub-items:
    \begin{itemize}
      \item Sub-item 1
      \item Sub-item 2
    \end{itemize}
\end{itemize}

\subsection{Numbered List}
\begin{enumerate}
  \item First step
  \item Second step
  \item Third step
\end{enumerate}

\section{Citations and References}
When you use the citation search feature, references will be automatically managed. For example, you might cite a paper like this: \cite{example2024}. All citations will be automatically added to your \texttt{cite.bib} file.

\section{Conclusion}
This template demonstrates basic \LaTeX{} features including mathematical formulas, tables, and lists. The bibliography system is now set up to work seamlessly with Bibby AI's citation management. Feel free to modify it for your needs!

\bibliographystyle{plain}
\bibliography{cite}

\end{document}\documentclass{article}
\usepackage[utf8]{inputenc}
\usepackage{booktabs}
\usepackage{amsmath}
\usepackage{natbib}

\title{My Document}
\author{Your Name}
\date{\today}

\begin{document}

\maketitle

\section{Introduction}
This is a sample document demonstrating various \LaTeX{} features including mathematical equations, tables, and formatting. You can cite papers using the citation search feature, and they will be automatically added to your \texttt{cite.bib} file.

\section{Mathematical Expressions}
Here are some example equations:

\subsection{Quadratic Formula}
The quadratic formula is given by:
\[x = \frac{-b \pm \sqrt{b^2 - 4ac}}{2a}\]

\subsection{Matrix Example}
A sample matrix:
\[
\begin{pmatrix}
a_{11} & a_{12} & a_{13} \\
a_{21} & a_{22} & a_{23} \\
a_{31} & a_{32} & a_{33}
\end{pmatrix}
\]

\section{Tables}
Below is a sample table showing data:

\begin{table}[h]
\centering
\begin{tabular}{@{}lrrr@{}}
\toprule
\textbf{Category} & \textbf{Value 1} & \textbf{Value 2} & \textbf{Total} \\
\midrule
Item A & 10.5 & 20.3 & 30.8 \\
Item B & 15.2 & 25.1 & 40.3 \\
Item C & 22.1 & 12.7 & 34.8 \\
\bottomrule
\end{tabular}
\caption{Sample Data Table}
\label{tab:sample}
\end{table}

\section{Lists}
\subsection{Bullet Points}
\begin{itemize}
  \item First item
  \item Second item
  \item Third item with sub-items:
    \begin{itemize}
      \item Sub-item 1
      \item Sub-item 2
    \end{itemize}
\end{itemize}

\subsection{Numbered List}
\begin{enumerate}
  \item First step
  \item Second step
  \item Third step
\end{enumerate}

\section{Citations and References}
When you use the citation search feature, references will be automatically managed. For example, you might cite a paper like this: \cite{example2024}. All citations will be automatically added to your \texttt{cite.bib} file.

\section{Conclusion}
This template demonstrates basic \LaTeX{} features including mathematical formulas, tables, and lists. The bibliography system is now set up to work seamlessly with Bibby AI's citation management. Feel free to modify it for your needs!

\bibliographystyle{plain}
\bibliography{cite}

\end{document}\documentclass{article}
\usepackage[utf8]{inputenc}
\usepackage{booktabs}
\usepackage{amsmath}
\usepackage{natbib}

\title{My Document}
\author{Your Name}
\date{\today}

\begin{document}

\maketitle

\section{Introduction}
This is a sample document demonstrating various \LaTeX{} features including mathematical equations, tables, and formatting. You can cite papers using the citation search feature, and they will be automatically added to your \texttt{cite.bib} file.

\section{Mathematical Expressions}
Here are some example equations:

\subsection{Quadratic Formula}
The quadratic formula is given by:
\[x = \frac{-b \pm \sqrt{b^2 - 4ac}}{2a}\]

\subsection{Matrix Example}
A sample matrix:
\[
\begin{pmatrix}
a_{11} & a_{12} & a_{13} \\
a_{21} & a_{22} & a_{23} \\
a_{31} & a_{32} & a_{33}
\end{pmatrix}
\]

\section{Tables}
Below is a sample table showing data:

\begin{table}[h]
\centering
\begin{tabular}{@{}lrrr@{}}
\toprule
\textbf{Category} & \textbf{Value 1} & \textbf{Value 2} & \textbf{Total} \\
\midrule
Item A & 10.5 & 20.3 & 30.8 \\
Item B & 15.2 & 25.1 & 40.3 \\
Item C & 22.1 & 12.7 & 34.8 \\
\bottomrule
\end{tabular}
\caption{Sample Data Table}
\label{tab:sample}
\end{table}

\section{Lists}
\subsection{Bullet Points}
\begin{itemize}
  \item First item
  \item Second item
  \item Third item with sub-items:
    \begin{itemize}
      \item Sub-item 1
      \item Sub-item 2
    \end{itemize}
\end{itemize}

\subsection{Numbered List}
\begin{enumerate}
  \item First step
  \item Second step
  \item Third step
\end{enumerate}

\section{Citations and References}
When you use the citation search feature, references will be automatically managed. For example, you might cite a paper like this: \cite{example2024}. All citations will be automatically added to your \texttt{cite.bib} file.

\section{Conclusion}
This template demonstrates basic \LaTeX{} features including mathematical formulas, tables, and lists. The bibliography system is now set up to work seamlessly with Bibby AI's citation management. Feel free to modify it for your needs!

\bibliographystyle{plain}
\bibliography{cite}

\end{document}\documentclass{article}
\usepackage[utf8]{inputenc}
\usepackage{booktabs}
\usepackage{amsmath}
\usepackage{natbib}

\title{My Document}
\author{Your Name}
\date{\today}

\begin{document}

\maketitle

\section{Introduction}
This is a sample document demonstrating various \LaTeX{} features including mathematical equations, tables, and formatting. You can cite papers using the citation search feature, and they will be automatically added to your \texttt{cite.bib} file.

\section{Mathematical Expressions}
Here are some example equations:

\subsection{Quadratic Formula}
The quadratic formula is given by:
\[x = \frac{-b \pm \sqrt{b^2 - 4ac}}{2a}\]

\subsection{Matrix Example}
A sample matrix:
\[
\begin{pmatrix}
a_{11} & a_{12} & a_{13} \\
a_{21} & a_{22} & a_{23} \\
a_{31} & a_{32} & a_{33}
\end{pmatrix}
\]

\section{Tables}
Below is a sample table showing data:

\begin{table}[h]
\centering
\begin{tabular}{@{}lrrr@{}}
\toprule
\textbf{Category} & \textbf{Value 1} & \textbf{Value 2} & \textbf{Total} \\
\midrule
Item A & 10.5 & 20.3 & 30.8 \\
Item B & 15.2 & 25.1 & 40.3 \\
Item C & 22.1 & 12.7 & 34.8 \\
\bottomrule
\end{tabular}
\caption{Sample Data Table}
\label{tab:sample}
\end{table}

\section{Lists}
\subsection{Bullet Points}
\begin{itemize}
  \item First item
  \item Second item
  \item Third item with sub-items:
    \begin{itemize}
      \item Sub-item 1
      \item Sub-item 2
    \end{itemize}
\end{itemize}

\subsection{Numbered List}
\begin{enumerate}
  \item First step
  \item Second step
  \item Third step
\end{enumerate}

\section{Citations and References}
When you use the citation search feature, references will be automatically managed. For example, you might cite a paper like this: \cite{example2024}. All citations will be automatically added to your \texttt{cite.bib} file.

\section{Conclusion}
This template demonstrates basic \LaTeX{} features including mathematical formulas, tables, and lists. The bibliography system is now set up to work seamlessly with Bibby AI's citation management. Feel free to modify it for your needs!

\bibliographystyle{plain}
\bibliography{cite}

\end{document}\documentclass{article}
\usepackage[utf8]{inputenc}
\usepackage{booktabs}
\usepackage{amsmath}
\usepackage{natbib}

\title{My Document}
\author{Your Name}
\date{\today}

\begin{document}

\maketitle

\section{Introduction}
This is a sample document demonstrating various \LaTeX{} features including mathematical equations, tables, and formatting. You can cite papers using the citation search feature, and they will be automatically added to your \texttt{cite.bib} file.

\section{Mathematical Expressions}
Here are some example equations:

\subsection{Quadratic Formula}
The quadratic formula is given by:
\[x = \frac{-b \pm \sqrt{b^2 - 4ac}}{2a}\]

\subsection{Matrix Example}
A sample matrix:
\[
\begin{pmatrix}
a_{11} & a_{12} & a_{13} \\
a_{21} & a_{22} & a_{23} \\
a_{31} & a_{32} & a_{33}
\end{pmatrix}
\]

\section{Tables}
Below is a sample table showing data:

\begin{table}[h]
\centering
\begin{tabular}{@{}lrrr@{}}
\toprule
\textbf{Category} & \textbf{Value 1} & \textbf{Value 2} & \textbf{Total} \\
\midrule
Item A & 10.5 & 20.3 & 30.8 \\
Item B & 15.2 & 25.1 & 40.3 \\
Item C & 22.1 & 12.7 & 34.8 \\
\bottomrule
\end{tabular}
\caption{Sample Data Table}
\label{tab:sample}
\end{table}

\section{Lists}
\subsection{Bullet Points}
\begin{itemize}
  \item First item
  \item Second item
  \item Third item with sub-items:
    \begin{itemize}
      \item Sub-item 1
      \item Sub-item 2
    \end{itemize}
\end{itemize}

\subsection{Numbered List}
\begin{enumerate}
  \item First step
  \item Second step
  \item Third step
\end{enumerate}

\section{Citations and References}
When you use the citation search feature, references will be automatically managed. For example, you might cite a paper like this: \cite{example2024}. All citations will be automatically added to your \texttt{cite.bib} file.

\section{Conclusion}
This template demonstrates basic \LaTeX{} features including mathematical formulas, tables, and lists. The bibliography system is now set up to work seamlessly with Bibby AI's citation management. Feel free to modify it for your needs!

\bibliographystyle{plain}
\bibliography{cite}

\end{document}\documentclass{article}
\usepackage[utf8]{inputenc}
\usepackage{booktabs}
\usepackage{amsmath}
\usepackage{natbib}

\title{My Document}
\author{Your Name}
\date{\today}

\begin{document}

\maketitle

\section{Introduction}
This is a sample document demonstrating various \LaTeX{} features including mathematical equations, tables, and formatting. You can cite papers using the citation search feature, and they will be automatically added to your \texttt{cite.bib} file.

\section{Mathematical Expressions}
Here are some example equations:

\subsection{Quadratic Formula}
The quadratic formula is given by:
\[x = \frac{-b \pm \sqrt{b^2 - 4ac}}{2a}\]

\subsection{Matrix Example}
A sample matrix:
\[
\begin{pmatrix}
a_{11} & a_{12} & a_{13} \\
a_{21} & a_{22} & a_{23} \\
a_{31} & a_{32} & a_{33}
\end{pmatrix}
\]

\section{Tables}
Below is a sample table showing data:

\begin{table}[h]
\centering
\begin{tabular}{@{}lrrr@{}}
\toprule
\textbf{Category} & \textbf{Value 1} & \textbf{Value 2} & \textbf{Total} \\
\midrule
Item A & 10.5 & 20.3 & 30.8 \\
Item B & 15.2 & 25.1 & 40.3 \\
Item C & 22.1 & 12.7 & 34.8 \\
\bottomrule
\end{tabular}
\caption{Sample Data Table}
\label{tab:sample}
\end{table}

\section{Lists}
\subsection{Bullet Points}
\begin{itemize}
  \item First item
  \item Second item
  \item Third item with sub-items:
    \begin{itemize}
      \item Sub-item 1
      \item Sub-item 2
    \end{itemize}
\end{itemize}

\subsection{Numbered List}
\begin{enumerate}
  \item First step
  \item Second step
  \item Third step
\end{enumerate}

\section{Citations and References}
When you use the citation search feature, references will be automatically managed. For example, you might cite a paper like this: \cite{example2024}. All citations will be automatically added to your \texttt{cite.bib} file.

\section{Conclusion}
This template demonstrates basic \LaTeX{} features including mathematical formulas, tables, and lists. The bibliography system is now set up to work seamlessly with Bibby AI's citation management. Feel free to modify it for your needs!

\bibliographystyle{plain}
\bibliography{cite}

\end{document}\documentclass{article}
\usepackage[utf8]{inputenc}
\usepackage{booktabs}
\usepackage{amsmath}
\usepackage{natbib}

\title{My Document}
\author{Your Name}
\date{\today}

\begin{document}

\maketitle

\section{Introduction}
This is a sample document demonstrating various \LaTeX{} features including mathematical equations, tables, and formatting. You can cite papers using the citation search feature, and they will be automatically added to your \texttt{cite.bib} file.

\section{Mathematical Expressions}
Here are some example equations:

\subsection{Quadratic Formula}
The quadratic formula is given by:
\[x = \frac{-b \pm \sqrt{b^2 - 4ac}}{2a}\]

\subsection{Matrix Example}
A sample matrix:
\[
\begin{pmatrix}
a_{11} & a_{12} & a_{13} \\
a_{21} & a_{22} & a_{23} \\
a_{31} & a_{32} & a_{33}
\end{pmatrix}
\]

\section{Tables}
Below is a sample table showing data:

\begin{table}[h]
\centering
\begin{tabular}{@{}lrrr@{}}
\toprule
\textbf{Category} & \textbf{Value 1} & \textbf{Value 2} & \textbf{Total} \\
\midrule
Item A & 10.5 & 20.3 & 30.8 \\
Item B & 15.2 & 25.1 & 40.3 \\
Item C & 22.1 & 12.7 & 34.8 \\
\bottomrule
\end{tabular}
\caption{Sample Data Table}
\label{tab:sample}
\end{table}

\section{Lists}
\subsection{Bullet Points}
\begin{itemize}
  \item First item
  \item Second item
  \item Third item with sub-items:
    \begin{itemize}
      \item Sub-item 1
      \item Sub-item 2
    \end{itemize}
\end{itemize}

\subsection{Numbered List}
\begin{enumerate}
  \item First step
  \item Second step
  \item Third step
\end{enumerate}

\section{Citations and References}
When you use the citation search feature, references will be automatically managed. For example, you might cite a paper like this: \cite{example2024}. All citations will be automatically added to your \texttt{cite.bib} file.

\section{Conclusion}
This template demonstrates basic \LaTeX{} features including mathematical formulas, tables, and lists. The bibliography system is now set up to work seamlessly with Bibby AI's citation management. Feel free to modify it for your needs!

\bibliographystyle{plain}
\bibliography{cite}

\end{document}\documentclass{article}
\usepackage[utf8]{inputenc}
\usepackage{booktabs}
\usepackage{amsmath}
\usepackage{natbib}

\title{My Document}
\author{Your Name}
\date{\today}

\begin{document}

\maketitle

\section{Introduction}
This is a sample document demonstrating various \LaTeX{} features including mathematical equations, tables, and formatting. You can cite papers using the citation search feature, and they will be automatically added to your \texttt{cite.bib} file.

\section{Mathematical Expressions}
Here are some example equations:

\subsection{Quadratic Formula}
The quadratic formula is given by:
\[x = \frac{-b \pm \sqrt{b^2 - 4ac}}{2a}\]

\subsection{Matrix Example}
A sample matrix:
\[
\begin{pmatrix}
a_{11} & a_{12} & a_{13} \\
a_{21} & a_{22} & a_{23} \\
a_{31} & a_{32} & a_{33}
\end{pmatrix}
\]

\section{Tables}
Below is a sample table showing data:

\begin{table}[h]
\centering
\begin{tabular}{@{}lrrr@{}}
\toprule
\textbf{Category} & \textbf{Value 1} & \textbf{Value 2} & \textbf{Total} \\
\midrule
Item A & 10.5 & 20.3 & 30.8 \\
Item B & 15.2 & 25.1 & 40.3 \\
Item C & 22.1 & 12.7 & 34.8 \\
\bottomrule
\end{tabular}
\caption{Sample Data Table}
\label{tab:sample}
\end{table}

\section{Lists}
\subsection{Bullet Points}
\begin{itemize}
  \item First item
  \item Second item
  \item Third item with sub-items:
    \begin{itemize}
      \item Sub-item 1
      \item Sub-item 2
    \end{itemize}
\end{itemize}

\subsection{Numbered List}
\begin{enumerate}
  \item First step
  \item Second step
  \item Third step
\end{enumerate}

\section{Citations and References}
When you use the citation search feature, references will be automatically managed. For example, you might cite a paper like this: \cite{example2024}. All citations will be automatically added to your \texttt{cite.bib} file.

\section{Conclusion}
This template demonstrates basic \LaTeX{} features including mathematical formulas, tables, and lists. The bibliography system is now set up to work seamlessly with Bibby AI's citation management. Feel free to modify it for your needs!

\bibliographystyle{plain}
\bibliography{cite}

\end{document}\documentclass{article}
\usepackage[utf8]{inputenc}
\usepackage{booktabs}
\usepackage{amsmath}
\usepackage{natbib}

\title{My Document}
\author{Your Name}
\date{\today}

\begin{document}

\maketitle

\section{Introduction}
This is a sample document demonstrating various \LaTeX{} features including mathematical equations, tables, and formatting. You can cite papers using the citation search feature, and they will be automatically added to your \texttt{cite.bib} file.

\section{Mathematical Expressions}
Here are some example equations:

\subsection{Quadratic Formula}
The quadratic formula is given by:
\[x = \frac{-b \pm \sqrt{b^2 - 4ac}}{2a}\]

\subsection{Matrix Example}
A sample matrix:
\[
\begin{pmatrix}
a_{11} & a_{12} & a_{13} \\
a_{21} & a_{22} & a_{23} \\
a_{31} & a_{32} & a_{33}
\end{pmatrix}
\]

\section{Tables}
Below is a sample table showing data:

\begin{table}[h]
\centering
\begin{tabular}{@{}lrrr@{}}
\toprule
\textbf{Category} & \textbf{Value 1} & \textbf{Value 2} & \textbf{Total} \\
\midrule
Item A & 10.5 & 20.3 & 30.8 \\
Item B & 15.2 & 25.1 & 40.3 \\
Item C & 22.1 & 12.7 & 34.8 \\
\bottomrule
\end{tabular}
\caption{Sample Data Table}
\label{tab:sample}
\end{table}

\section{Lists}
\subsection{Bullet Points}
\begin{itemize}
  \item First item
  \item Second item
  \item Third item with sub-items:
    \begin{itemize}
      \item Sub-item 1
      \item Sub-item 2
    \end{itemize}
\end{itemize}

\subsection{Numbered List}
\begin{enumerate}
  \item First step
  \item Second step
  \item Third step
\end{enumerate}

\section{Citations and References}
When you use the citation search feature, references will be automatically managed. For example, you might cite a paper like this: \cite{example2024}. All citations will be automatically added to your \texttt{cite.bib} file.

\section{Conclusion}
This template demonstrates basic \LaTeX{} features including mathematical formulas, tables, and lists. The bibliography system is now set up to work seamlessly with Bibby AI's citation management. Feel free to modify it for your needs!

\bibliographystyle{plain}
\bibliography{cite}

\end{document}\documentclass{article}
\usepackage[utf8]{inputenc}
\usepackage{booktabs}
\usepackage{amsmath}
\usepackage{natbib}

\title{My Document}
\author{Your Name}
\date{\today}

\begin{document}

\maketitle

\section{Introduction}
This is a sample document demonstrating various \LaTeX{} features including mathematical equations, tables, and formatting. You can cite papers using the citation search feature, and they will be automatically added to your \texttt{cite.bib} file.

\section{Mathematical Expressions}
Here are some example equations:

\subsection{Quadratic Formula}
The quadratic formula is given by:
\[x = \frac{-b \pm \sqrt{b^2 - 4ac}}{2a}\]

\subsection{Matrix Example}
A sample matrix:
\[
\begin{pmatrix}
a_{11} & a_{12} & a_{13} \\
a_{21} & a_{22} & a_{23} \\
a_{31} & a_{32} & a_{33}
\end{pmatrix}
\]

\section{Tables}
Below is a sample table showing data:

\begin{table}[h]
\centering
\begin{tabular}{@{}lrrr@{}}
\toprule
\textbf{Category} & \textbf{Value 1} & \textbf{Value 2} & \textbf{Total} \\
\midrule
Item A & 10.5 & 20.3 & 30.8 \\
Item B & 15.2 & 25.1 & 40.3 \\
Item C & 22.1 & 12.7 & 34.8 \\
\bottomrule
\end{tabular}
\caption{Sample Data Table}
\label{tab:sample}
\end{table}

\section{Lists}
\subsection{Bullet Points}
\begin{itemize}
  \item First item
  \item Second item
  \item Third item with sub-items:
    \begin{itemize}
      \item Sub-item 1
      \item Sub-item 2
    \end{itemize}
\end{itemize}

\subsection{Numbered List}
\begin{enumerate}
  \item First step
  \item Second step
  \item Third step
\end{enumerate}

\section{Citations and References}
When you use the citation search feature, references will be automatically managed. For example, you might cite a paper like this: \cite{example2024}. All citations will be automatically added to your \texttt{cite.bib} file.

\section{Conclusion}
This template demonstrates basic \LaTeX{} features including mathematical formulas, tables, and lists. The bibliography system is now set up to work seamlessly with Bibby AI's citation management. Feel free to modify it for your needs!

\bibliographystyle{plain}
\bibliography{cite}

\end{document}\documentclass{article}
\usepackage[utf8]{inputenc}
\usepackage{booktabs}
\usepackage{amsmath}
\usepackage{natbib}

\title{My Document}
\author{Your Name}
\date{\today}

\begin{document}

\maketitle

\section{Introduction}
This is a sample document demonstrating various \LaTeX{} features including mathematical equations, tables, and formatting. You can cite papers using the citation search feature, and they will be automatically added to your \texttt{cite.bib} file.

\section{Mathematical Expressions}
Here are some example equations:

\subsection{Quadratic Formula}
The quadratic formula is given by:
\[x = \frac{-b \pm \sqrt{b^2 - 4ac}}{2a}\]

\subsection{Matrix Example}
A sample matrix:
\[
\begin{pmatrix}
a_{11} & a_{12} & a_{13} \\
a_{21} & a_{22} & a_{23} \\
a_{31} & a_{32} & a_{33}
\end{pmatrix}
\]

\section{Tables}
Below is a sample table showing data:

\begin{table}[h]
\centering
\begin{tabular}{@{}lrrr@{}}
\toprule
\textbf{Category} & \textbf{Value 1} & \textbf{Value 2} & \textbf{Total} \\
\midrule
Item A & 10.5 & 20.3 & 30.8 \\
Item B & 15.2 & 25.1 & 40.3 \\
Item C & 22.1 & 12.7 & 34.8 \\
\bottomrule
\end{tabular}
\caption{Sample Data Table}
\label{tab:sample}
\end{table}

\section{Lists}
\subsection{Bullet Points}
\begin{itemize}
  \item First item
  \item Second item
  \item Third item with sub-items:
    \begin{itemize}
      \item Sub-item 1
      \item Sub-item 2
    \end{itemize}
\end{itemize}

\subsection{Numbered List}
\begin{enumerate}
  \item First step
  \item Second step
  \item Third step
\end{enumerate}

\section{Citations and References}
When you use the citation search feature, references will be automatically managed. For example, you might cite a paper like this: \cite{example2024}. All citations will be automatically added to your \texttt{cite.bib} file.

\section{Conclusion}
This template demonstrates basic \LaTeX{} features including mathematical formulas, tables, and lists. The bibliography system is now set up to work seamlessly with Bibby AI's citation management. Feel free to modify it for your needs!

\bibliographystyle{plain}
\bibliography{cite}

\end{document}\documentclass{article}
\usepackage[utf8]{inputenc}
\usepackage{booktabs}
\usepackage{amsmath}
\usepackage{natbib}

\title{My Document}
\author{Your Name}
\date{\today}

\begin{document}

\maketitle

\section{Introduction}
This is a sample document demonstrating various \LaTeX{} features including mathematical equations, tables, and formatting. You can cite papers using the citation search feature, and they will be automatically added to your \texttt{cite.bib} file.

\section{Mathematical Expressions}
Here are some example equations:

\subsection{Quadratic Formula}
The quadratic formula is given by:
\[x = \frac{-b \pm \sqrt{b^2 - 4ac}}{2a}\]

\subsection{Matrix Example}
A sample matrix:
\[
\begin{pmatrix}
a_{11} & a_{12} & a_{13} \\
a_{21} & a_{22} & a_{23} \\
a_{31} & a_{32} & a_{33}
\end{pmatrix}
\]

\section{Tables}
Below is a sample table showing data:

\begin{table}[h]
\centering
\begin{tabular}{@{}lrrr@{}}
\toprule
\textbf{Category} & \textbf{Value 1} & \textbf{Value 2} & \textbf{Total} \\
\midrule
Item A & 10.5 & 20.3 & 30.8 \\
Item B & 15.2 & 25.1 & 40.3 \\
Item C & 22.1 & 12.7 & 34.8 \\
\bottomrule
\end{tabular}
\caption{Sample Data Table}
\label{tab:sample}
\end{table}

\section{Lists}
\subsection{Bullet Points}
\begin{itemize}
  \item First item
  \item Second item
  \item Third item with sub-items:
    \begin{itemize}
      \item Sub-item 1
      \item Sub-item 2
    \end{itemize}
\end{itemize}

\subsection{Numbered List}
\begin{enumerate}
  \item First step
  \item Second step
  \item Third step
\end{enumerate}

\section{Citations and References}
When you use the citation search feature, references will be automatically managed. For example, you might cite a paper like this: \cite{example2024}. All citations will be automatically added to your \texttt{cite.bib} file.

\section{Conclusion}
This template demonstrates basic \LaTeX{} features including mathematical formulas, tables, and lists. The bibliography system is now set up to work seamlessly with Bibby AI's citation management. Feel free to modify it for your needs!

\bibliographystyle{plain}
\bibliography{cite}

\end{document}\documentclass{article}
\usepackage[utf8]{inputenc}
\usepackage{booktabs}
\usepackage{amsmath}
\usepackage{natbib}

\title{My Document}
\author{Your Name}
\date{\today}

\begin{document}

\maketitle

\section{Introduction}
This is a sample document demonstrating various \LaTeX{} features including mathematical equations, tables, and formatting. You can cite papers using the citation search feature, and they will be automatically added to your \texttt{cite.bib} file.

\section{Mathematical Expressions}
Here are some example equations:

\subsection{Quadratic Formula}
The quadratic formula is given by:
\[x = \frac{-b \pm \sqrt{b^2 - 4ac}}{2a}\]

\subsection{Matrix Example}
A sample matrix:
\[
\begin{pmatrix}
a_{11} & a_{12} & a_{13} \\
a_{21} & a_{22} & a_{23} \\
a_{31} & a_{32} & a_{33}
\end{pmatrix}
\]

\section{Tables}
Below is a sample table showing data:

\begin{table}[h]
\centering
\begin{tabular}{@{}lrrr@{}}
\toprule
\textbf{Category} & \textbf{Value 1} & \textbf{Value 2} & \textbf{Total} \\
\midrule
Item A & 10.5 & 20.3 & 30.8 \\
Item B & 15.2 & 25.1 & 40.3 \\
Item C & 22.1 & 12.7 & 34.8 \\
\bottomrule
\end{tabular}
\caption{Sample Data Table}
\label{tab:sample}
\end{table}

\section{Lists}
\subsection{Bullet Points}
\begin{itemize}
  \item First item
  \item Second item
  \item Third item with sub-items:
    \begin{itemize}
      \item Sub-item 1
      \item Sub-item 2
    \end{itemize}
\end{itemize}

\subsection{Numbered List}
\begin{enumerate}
  \item First step
  \item Second step
  \item Third step
\end{enumerate}

\section{Citations and References}
When you use the citation search feature, references will be automatically managed. For example, you might cite a paper like this: \cite{example2024}. All citations will be automatically added to your \texttt{cite.bib} file.

\section{Conclusion}
This template demonstrates basic \LaTeX{} features including mathematical formulas, tables, and lists. The bibliography system is now set up to work seamlessly with Bibby AI's citation management. Feel free to modify it for your needs!

\bibliographystyle{plain}
\bibliography{cite}

\end{document}\documentclass{article}
\usepackage[utf8]{inputenc}
\usepackage{booktabs}
\usepackage{amsmath}
\usepackage{natbib}

\title{My Document}
\author{Your Name}
\date{\today}

\begin{document}

\maketitle

\section{Introduction}
This is a sample document demonstrating various \LaTeX{} features including mathematical equations, tables, and formatting. You can cite papers using the citation search feature, and they will be automatically added to your \texttt{cite.bib} file.

\section{Mathematical Expressions}
Here are some example equations:

\subsection{Quadratic Formula}
The quadratic formula is given by:
\[x = \frac{-b \pm \sqrt{b^2 - 4ac}}{2a}\]

\subsection{Matrix Example}
A sample matrix:
\[
\begin{pmatrix}
a_{11} & a_{12} & a_{13} \\
a_{21} & a_{22} & a_{23} \\
a_{31} & a_{32} & a_{33}
\end{pmatrix}
\]

\section{Tables}
Below is a sample table showing data:

\begin{table}[h]
\centering
\begin{tabular}{@{}lrrr@{}}
\toprule
\textbf{Category} & \textbf{Value 1} & \textbf{Value 2} & \textbf{Total} \\
\midrule
Item A & 10.5 & 20.3 & 30.8 \\
Item B & 15.2 & 25.1 & 40.3 \\
Item C & 22.1 & 12.7 & 34.8 \\
\bottomrule
\end{tabular}
\caption{Sample Data Table}
\label{tab:sample}
\end{table}

\section{Lists}
\subsection{Bullet Points}
\begin{itemize}
  \item First item
  \item Second item
  \item Third item with sub-items:
    \begin{itemize}
      \item Sub-item 1
      \item Sub-item 2
    \end{itemize}
\end{itemize}

\subsection{Numbered List}
\begin{enumerate}
  \item First step
  \item Second step
  \item Third step
\end{enumerate}

\section{Citations and References}
When you use the citation search feature, references will be automatically managed. For example, you might cite a paper like this: \cite{example2024}. All citations will be automatically added to your \texttt{cite.bib} file.

\section{Conclusion}
This template demonstrates basic \LaTeX{} features including mathematical formulas, tables, and lists. The bibliography system is now set up to work seamlessly with Bibby AI's citation management. Feel free to modify it for your needs!

\bibliographystyle{plain}
\bibliography{cite}

\end{document}\documentclass{article}
\usepackage[utf8]{inputenc}
\usepackage{booktabs}
\usepackage{amsmath}
\usepackage{natbib}

\title{My Document}
\author{Your Name}
\date{\today}

\begin{document}

\maketitle

\section{Introduction}
This is a sample document demonstrating various \LaTeX{} features including mathematical equations, tables, and formatting. You can cite papers using the citation search feature, and they will be automatically added to your \texttt{cite.bib} file.

\section{Mathematical Expressions}
Here are some example equations:

\subsection{Quadratic Formula}
The quadratic formula is given by:
\[x = \frac{-b \pm \sqrt{b^2 - 4ac}}{2a}\]

\subsection{Matrix Example}
A sample matrix:
\[
\begin{pmatrix}
a_{11} & a_{12} & a_{13} \\
a_{21} & a_{22} & a_{23} \\
a_{31} & a_{32} & a_{33}
\end{pmatrix}
\]

\section{Tables}
Below is a sample table showing data:

\begin{table}[h]
\centering
\begin{tabular}{@{}lrrr@{}}
\toprule
\textbf{Category} & \textbf{Value 1} & \textbf{Value 2} & \textbf{Total} \\
\midrule
Item A & 10.5 & 20.3 & 30.8 \\
Item B & 15.2 & 25.1 & 40.3 \\
Item C & 22.1 & 12.7 & 34.8 \\
\bottomrule
\end{tabular}
\caption{Sample Data Table}
\label{tab:sample}
\end{table}

\section{Lists}
\subsection{Bullet Points}
\begin{itemize}
  \item First item
  \item Second item
  \item Third item with sub-items:
    \begin{itemize}
      \item Sub-item 1
      \item Sub-item 2
    \end{itemize}
\end{itemize}

\subsection{Numbered List}
\begin{enumerate}
  \item First step
  \item Second step
  \item Third step
\end{enumerate}

\section{Citations and References}
When you use the citation search feature, references will be automatically managed. For example, you might cite a paper like this: \cite{example2024}. All citations will be automatically added to your \texttt{cite.bib} file.

\section{Conclusion}
This template demonstrates basic \LaTeX{} features including mathematical formulas, tables, and lists. The bibliography system is now set up to work seamlessly with Bibby AI's citation management. Feel free to modify it for your needs!

\bibliographystyle{plain}
\bibliography{cite}

\end{document}\documentclass{article}
\usepackage[utf8]{inputenc}
\usepackage{booktabs}
\usepackage{amsmath}
\usepackage{natbib}

\title{My Document}
\author{Your Name}
\date{\today}

\begin{document}

\maketitle

\section{Introduction}
This is a sample document demonstrating various \LaTeX{} features including mathematical equations, tables, and formatting. You can cite papers using the citation search feature, and they will be automatically added to your \texttt{cite.bib} file.

\section{Mathematical Expressions}
Here are some example equations:

\subsection{Quadratic Formula}
The quadratic formula is given by:
\[x = \frac{-b \pm \sqrt{b^2 - 4ac}}{2a}\]

\subsection{Matrix Example}
A sample matrix:
\[
\begin{pmatrix}
a_{11} & a_{12} & a_{13} \\
a_{21} & a_{22} & a_{23} \\
a_{31} & a_{32} & a_{33}
\end{pmatrix}
\]

\section{Tables}
Below is a sample table showing data:

\begin{table}[h]
\centering
\begin{tabular}{@{}lrrr@{}}
\toprule
\textbf{Category} & \textbf{Value 1} & \textbf{Value 2} & \textbf{Total} \\
\midrule
Item A & 10.5 & 20.3 & 30.8 \\
Item B & 15.2 & 25.1 & 40.3 \\
Item C & 22.1 & 12.7 & 34.8 \\
\bottomrule
\end{tabular}
\caption{Sample Data Table}
\label{tab:sample}
\end{table}

\section{Lists}
\subsection{Bullet Points}
\begin{itemize}
  \item First item
  \item Second item
  \item Third item with sub-items:
    \begin{itemize}
      \item Sub-item 1
      \item Sub-item 2
    \end{itemize}
\end{itemize}

\subsection{Numbered List}
\begin{enumerate}
  \item First step
  \item Second step
  \item Third step
\end{enumerate}

\section{Citations and References}
When you use the citation search feature, references will be automatically managed. For example, you might cite a paper like this: \cite{example2024}. All citations will be automatically added to your \texttt{cite.bib} file.

\section{Conclusion}
This template demonstrates basic \LaTeX{} features including mathematical formulas, tables, and lists. The bibliography system is now set up to work seamlessly with Bibby AI's citation management. Feel free to modify it for your needs!

\bibliographystyle{plain}
\bibliography{cite}

\end{document}\documentclass{article}
\usepackage[utf8]{inputenc}
\usepackage{booktabs}
\usepackage{amsmath}
\usepackage{natbib}

\title{My Document}
\author{Your Name}
\date{\today}

\begin{document}

\maketitle

\section{Introduction}
This is a sample document demonstrating various \LaTeX{} features including mathematical equations, tables, and formatting. You can cite papers using the citation search feature, and they will be automatically added to your \texttt{cite.bib} file.

\section{Mathematical Expressions}
Here are some example equations:

\subsection{Quadratic Formula}
The quadratic formula is given by:
\[x = \frac{-b \pm \sqrt{b^2 - 4ac}}{2a}\]

\subsection{Matrix Example}
A sample matrix:
\[
\begin{pmatrix}
a_{11} & a_{12} & a_{13} \\
a_{21} & a_{22} & a_{23} \\
a_{31} & a_{32} & a_{33}
\end{pmatrix}
\]

\section{Tables}
Below is a sample table showing data:

\begin{table}[h]
\centering
\begin{tabular}{@{}lrrr@{}}
\toprule
\textbf{Category} & \textbf{Value 1} & \textbf{Value 2} & \textbf{Total} \\
\midrule
Item A & 10.5 & 20.3 & 30.8 \\
Item B & 15.2 & 25.1 & 40.3 \\
Item C & 22.1 & 12.7 & 34.8 \\
\bottomrule
\end{tabular}
\caption{Sample Data Table}
\label{tab:sample}
\end{table}

\section{Lists}
\subsection{Bullet Points}
\begin{itemize}
  \item First item
  \item Second item
  \item Third item with sub-items:
    \begin{itemize}
      \item Sub-item 1
      \item Sub-item 2
    \end{itemize}
\end{itemize}

\subsection{Numbered List}
\begin{enumerate}
  \item First step
  \item Second step
  \item Third step
\end{enumerate}

\section{Citations and References}
When you use the citation search feature, references will be automatically managed. For example, you might cite a paper like this: \cite{example2024}. All citations will be automatically added to your \texttt{cite.bib} file.

\section{Conclusion}
This template demonstrates basic \LaTeX{} features including mathematical formulas, tables, and lists. The bibliography system is now set up to work seamlessly with Bibby AI's citation management. Feel free to modify it for your needs!

\bibliographystyle{plain}
\bibliography{cite}

\end{document}\documentclass{article}
\usepackage[utf8]{inputenc}
\usepackage{booktabs}
\usepackage{amsmath}
\usepackage{natbib}

\title{My Document}
\author{Your Name}
\date{\today}

\begin{document}

\maketitle

\section{Introduction}
This is a sample document demonstrating various \LaTeX{} features including mathematical equations, tables, and formatting. You can cite papers using the citation search feature, and they will be automatically added to your \texttt{cite.bib} file.

\section{Mathematical Expressions}
Here are some example equations:

\subsection{Quadratic Formula}
The quadratic formula is given by:
\[x = \frac{-b \pm \sqrt{b^2 - 4ac}}{2a}\]

\subsection{Matrix Example}
A sample matrix:
\[
\begin{pmatrix}
a_{11} & a_{12} & a_{13} \\
a_{21} & a_{22} & a_{23} \\
a_{31} & a_{32} & a_{33}
\end{pmatrix}
\]

\section{Tables}
Below is a sample table showing data:

\begin{table}[h]
\centering
\begin{tabular}{@{}lrrr@{}}
\toprule
\textbf{Category} & \textbf{Value 1} & \textbf{Value 2} & \textbf{Total} \\
\midrule
Item A & 10.5 & 20.3 & 30.8 \\
Item B & 15.2 & 25.1 & 40.3 \\
Item C & 22.1 & 12.7 & 34.8 \\
\bottomrule
\end{tabular}
\caption{Sample Data Table}
\label{tab:sample}
\end{table}

\section{Lists}
\subsection{Bullet Points}
\begin{itemize}
  \item First item
  \item Second item
  \item Third item with sub-items:
    \begin{itemize}
      \item Sub-item 1
      \item Sub-item 2
    \end{itemize}
\end{itemize}

\subsection{Numbered List}
\begin{enumerate}
  \item First step
  \item Second step
  \item Third step
\end{enumerate}

\section{Citations and References}
When you use the citation search feature, references will be automatically managed. For example, you might cite a paper like this: \cite{example2024}. All citations will be automatically added to your \texttt{cite.bib} file.

\section{Conclusion}
This template demonstrates basic \LaTeX{} features including mathematical formulas, tables, and lists. The bibliography system is now set up to work seamlessly with Bibby AI's citation management. Feel free to modify it for your needs!

\bibliographystyle{plain}
\bibliography{cite}

\end{document}\documentclass{article}
\usepackage[utf8]{inputenc}
\usepackage{booktabs}
\usepackage{amsmath}
\usepackage{natbib}

\title{My Document}
\author{Your Name}
\date{\today}

\begin{document}

\maketitle

\section{Introduction}
This is a sample document demonstrating various \LaTeX{} features including mathematical equations, tables, and formatting. You can cite papers using the citation search feature, and they will be automatically added to your \texttt{cite.bib} file.

\section{Mathematical Expressions}
Here are some example equations:

\subsection{Quadratic Formula}
The quadratic formula is given by:
\[x = \frac{-b \pm \sqrt{b^2 - 4ac}}{2a}\]

\subsection{Matrix Example}
A sample matrix:
\[
\begin{pmatrix}
a_{11} & a_{12} & a_{13} \\
a_{21} & a_{22} & a_{23} \\
a_{31} & a_{32} & a_{33}
\end{pmatrix}
\]

\section{Tables}
Below is a sample table showing data:

\begin{table}[h]
\centering
\begin{tabular}{@{}lrrr@{}}
\toprule
\textbf{Category} & \textbf{Value 1} & \textbf{Value 2} & \textbf{Total} \\
\midrule
Item A & 10.5 & 20.3 & 30.8 \\
Item B & 15.2 & 25.1 & 40.3 \\
Item C & 22.1 & 12.7 & 34.8 \\
\bottomrule
\end{tabular}
\caption{Sample Data Table}
\label{tab:sample}
\end{table}

\section{Lists}
\subsection{Bullet Points}
\begin{itemize}
  \item First item
  \item Second item
  \item Third item with sub-items:
    \begin{itemize}
      \item Sub-item 1
      \item Sub-item 2
    \end{itemize}
\end{itemize}

\subsection{Numbered List}
\begin{enumerate}
  \item First step
  \item Second step
  \item Third step
\end{enumerate}

\section{Citations and References}
When you use the citation search feature, references will be automatically managed. For example, you might cite a paper like this: \cite{example2024}. All citations will be automatically added to your \texttt{cite.bib} file.

\section{Conclusion}
This template demonstrates basic \LaTeX{} features including mathematical formulas, tables, and lists. The bibliography system is now set up to work seamlessly with Bibby AI's citation management. Feel free to modify it for your needs!

\bibliographystyle{plain}
\bibliography{cite}

\end{document}\documentclass{article}
\usepackage[utf8]{inputenc}
\usepackage{booktabs}
\usepackage{amsmath}
\usepackage{natbib}

\title{My Document}
\author{Your Name}
\date{\today}

\begin{document}

\maketitle

\section{Introduction}
This is a sample document demonstrating various \LaTeX{} features including mathematical equations, tables, and formatting. You can cite papers using the citation search feature, and they will be automatically added to your \texttt{cite.bib} file.

\section{Mathematical Expressions}
Here are some example equations:

\subsection{Quadratic Formula}
The quadratic formula is given by:
\[x = \frac{-b \pm \sqrt{b^2 - 4ac}}{2a}\]

\subsection{Matrix Example}
A sample matrix:
\[
\begin{pmatrix}
a_{11} & a_{12} & a_{13} \\
a_{21} & a_{22} & a_{23} \\
a_{31} & a_{32} & a_{33}
\end{pmatrix}
\]

\section{Tables}
Below is a sample table showing data:

\begin{table}[h]
\centering
\begin{tabular}{@{}lrrr@{}}
\toprule
\textbf{Category} & \textbf{Value 1} & \textbf{Value 2} & \textbf{Total} \\
\midrule
Item A & 10.5 & 20.3 & 30.8 \\
Item B & 15.2 & 25.1 & 40.3 \\
Item C & 22.1 & 12.7 & 34.8 \\
\bottomrule
\end{tabular}
\caption{Sample Data Table}
\label{tab:sample}
\end{table}

\section{Lists}
\subsection{Bullet Points}
\begin{itemize}
  \item First item
  \item Second item
  \item Third item with sub-items:
    \begin{itemize}
      \item Sub-item 1
      \item Sub-item 2
    \end{itemize}
\end{itemize}

\subsection{Numbered List}
\begin{enumerate}
  \item First step
  \item Second step
  \item Third step
\end{enumerate}

\section{Citations and References}
When you use the citation search feature, references will be automatically managed. For example, you might cite a paper like this: \cite{example2024}. All citations will be automatically added to your \texttt{cite.bib} file.

\section{Conclusion}
This template demonstrates basic \LaTeX{} features including mathematical formulas, tables, and lists. The bibliography system is now set up to work seamlessly with Bibby AI's citation management. Feel free to modify it for your needs!

\bibliographystyle{plain}
\bibliography{cite}

\end{document}\documentclass{article}
\usepackage[utf8]{inputenc}
\usepackage{booktabs}
\usepackage{amsmath}
\usepackage{natbib}

\title{My Document}
\author{Your Name}
\date{\today}

\begin{document}

\maketitle

\section{Introduction}
This is a sample document demonstrating various \LaTeX{} features including mathematical equations, tables, and formatting. You can cite papers using the citation search feature, and they will be automatically added to your \texttt{cite.bib} file.

\section{Mathematical Expressions}
Here are some example equations:

\subsection{Quadratic Formula}
The quadratic formula is given by:
\[x = \frac{-b \pm \sqrt{b^2 - 4ac}}{2a}\]

\subsection{Matrix Example}
A sample matrix:
\[
\begin{pmatrix}
a_{11} & a_{12} & a_{13} \\
a_{21} & a_{22} & a_{23} \\
a_{31} & a_{32} & a_{33}
\end{pmatrix}
\]

\section{Tables}
Below is a sample table showing data:

\begin{table}[h]
\centering
\begin{tabular}{@{}lrrr@{}}
\toprule
\textbf{Category} & \textbf{Value 1} & \textbf{Value 2} & \textbf{Total} \\
\midrule
Item A & 10.5 & 20.3 & 30.8 \\
Item B & 15.2 & 25.1 & 40.3 \\
Item C & 22.1 & 12.7 & 34.8 \\
\bottomrule
\end{tabular}
\caption{Sample Data Table}
\label{tab:sample}
\end{table}

\section{Lists}
\subsection{Bullet Points}
\begin{itemize}
  \item First item
  \item Second item
  \item Third item with sub-items:
    \begin{itemize}
      \item Sub-item 1
      \item Sub-item 2
    \end{itemize}
\end{itemize}

\subsection{Numbered List}
\begin{enumerate}
  \item First step
  \item Second step
  \item Third step
\end{enumerate}

\section{Citations and References}
When you use the citation search feature, references will be automatically managed. For example, you might cite a paper like this: \cite{example2024}. All citations will be automatically added to your \texttt{cite.bib} file.

\section{Conclusion}
This template demonstrates basic \LaTeX{} features including mathematical formulas, tables, and lists. The bibliography system is now set up to work seamlessly with Bibby AI's citation management. Feel free to modify it for your needs!

\bibliographystyle{plain}
\bibliography{cite}

\end{document}\documentclass{article}
\usepackage[utf8]{inputenc}
\usepackage{booktabs}
\usepackage{amsmath}
\usepackage{natbib}

\title{My Document}
\author{Your Name}
\date{\today}

\begin{document}

\maketitle

\section{Introduction}
This is a sample document demonstrating various \LaTeX{} features including mathematical equations, tables, and formatting. You can cite papers using the citation search feature, and they will be automatically added to your \texttt{cite.bib} file.

\section{Mathematical Expressions}
Here are some example equations:

\subsection{Quadratic Formula}
The quadratic formula is given by:
\[x = \frac{-b \pm \sqrt{b^2 - 4ac}}{2a}\]

\subsection{Matrix Example}
A sample matrix:
\[
\begin{pmatrix}
a_{11} & a_{12} & a_{13} \\
a_{21} & a_{22} & a_{23} \\
a_{31} & a_{32} & a_{33}
\end{pmatrix}
\]

\section{Tables}
Below is a sample table showing data:

\begin{table}[h]
\centering
\begin{tabular}{@{}lrrr@{}}
\toprule
\textbf{Category} & \textbf{Value 1} & \textbf{Value 2} & \textbf{Total} \\
\midrule
Item A & 10.5 & 20.3 & 30.8 \\
Item B & 15.2 & 25.1 & 40.3 \\
Item C & 22.1 & 12.7 & 34.8 \\
\bottomrule
\end{tabular}
\caption{Sample Data Table}
\label{tab:sample}
\end{table}

\section{Lists}
\subsection{Bullet Points}
\begin{itemize}
  \item First item
  \item Second item
  \item Third item with sub-items:
    \begin{itemize}
      \item Sub-item 1
      \item Sub-item 2
    \end{itemize}
\end{itemize}

\subsection{Numbered List}
\begin{enumerate}
  \item First step
  \item Second step
  \item Third step
\end{enumerate}

\section{Citations and References}
When you use the citation search feature, references will be automatically managed. For example, you might cite a paper like this: \cite{example2024}. All citations will be automatically added to your \texttt{cite.bib} file.

\section{Conclusion}
This template demonstrates basic \LaTeX{} features including mathematical formulas, tables, and lists. The bibliography system is now set up to work seamlessly with Bibby AI's citation management. Feel free to modify it for your needs!

\bibliographystyle{plain}
\bibliography{cite}

\end{document}\documentclass{article}
\usepackage[utf8]{inputenc}
\usepackage{booktabs}
\usepackage{amsmath}
\usepackage{natbib}

\title{My Document}
\author{Your Name}
\date{\today}

\begin{document}

\maketitle

\section{Introduction}
This is a sample document demonstrating various \LaTeX{} features including mathematical equations, tables, and formatting. You can cite papers using the citation search feature, and they will be automatically added to your \texttt{cite.bib} file.

\section{Mathematical Expressions}
Here are some example equations:

\subsection{Quadratic Formula}
The quadratic formula is given by:
\[x = \frac{-b \pm \sqrt{b^2 - 4ac}}{2a}\]

\subsection{Matrix Example}
A sample matrix:
\[
\begin{pmatrix}
a_{11} & a_{12} & a_{13} \\
a_{21} & a_{22} & a_{23} \\
a_{31} & a_{32} & a_{33}
\end{pmatrix}
\]

\section{Tables}
Below is a sample table showing data:

\begin{table}[h]
\centering
\begin{tabular}{@{}lrrr@{}}
\toprule
\textbf{Category} & \textbf{Value 1} & \textbf{Value 2} & \textbf{Total} \\
\midrule
Item A & 10.5 & 20.3 & 30.8 \\
Item B & 15.2 & 25.1 & 40.3 \\
Item C & 22.1 & 12.7 & 34.8 \\
\bottomrule
\end{tabular}
\caption{Sample Data Table}
\label{tab:sample}
\end{table}

\section{Lists}
\subsection{Bullet Points}
\begin{itemize}
  \item First item
  \item Second item
  \item Third item with sub-items:
    \begin{itemize}
      \item Sub-item 1
      \item Sub-item 2
    \end{itemize}
\end{itemize}

\subsection{Numbered List}
\begin{enumerate}
  \item First step
  \item Second step
  \item Third step
\end{enumerate}

\section{Citations and References}
When you use the citation search feature, references will be automatically managed. For example, you might cite a paper like this: \cite{example2024}. All citations will be automatically added to your \texttt{cite.bib} file.

\section{Conclusion}
This template demonstrates basic \LaTeX{} features including mathematical formulas, tables, and lists. The bibliography system is now set up to work seamlessly with Bibby AI's citation management. Feel free to modify it for your needs!

\bibliographystyle{plain}
\bibliography{cite}

\end{document}\documentclass{article}
\usepackage[utf8]{inputenc}
\usepackage{booktabs}
\usepackage{amsmath}
\usepackage{natbib}

\title{My Document}
\author{Your Name}
\date{\today}

\begin{document}

\maketitle

\section{Introduction}
This is a sample document demonstrating various \LaTeX{} features including mathematical equations, tables, and formatting. You can cite papers using the citation search feature, and they will be automatically added to your \texttt{cite.bib} file.

\section{Mathematical Expressions}
Here are some example equations:

\subsection{Quadratic Formula}
The quadratic formula is given by:
\[x = \frac{-b \pm \sqrt{b^2 - 4ac}}{2a}\]

\subsection{Matrix Example}
A sample matrix:
\[
\begin{pmatrix}
a_{11} & a_{12} & a_{13} \\
a_{21} & a_{22} & a_{23} \\
a_{31} & a_{32} & a_{33}
\end{pmatrix}
\]

\section{Tables}
Below is a sample table showing data:

\begin{table}[h]
\centering
\begin{tabular}{@{}lrrr@{}}
\toprule
\textbf{Category} & \textbf{Value 1} & \textbf{Value 2} & \textbf{Total} \\
\midrule
Item A & 10.5 & 20.3 & 30.8 \\
Item B & 15.2 & 25.1 & 40.3 \\
Item C & 22.1 & 12.7 & 34.8 \\
\bottomrule
\end{tabular}
\caption{Sample Data Table}
\label{tab:sample}
\end{table}

\section{Lists}
\subsection{Bullet Points}
\begin{itemize}
  \item First item
  \item Second item
  \item Third item with sub-items:
    \begin{itemize}
      \item Sub-item 1
      \item Sub-item 2
    \end{itemize}
\end{itemize}

\subsection{Numbered List}
\begin{enumerate}
  \item First step
  \item Second step
  \item Third step
\end{enumerate}

\section{Citations and References}
When you use the citation search feature, references will be automatically managed. For example, you might cite a paper like this: \cite{example2024}. All citations will be automatically added to your \texttt{cite.bib} file.

\section{Conclusion}
This template demonstrates basic \LaTeX{} features including mathematical formulas, tables, and lists. The bibliography system is now set up to work seamlessly with Bibby AI's citation management. Feel free to modify it for your needs!

\bibliographystyle{plain}
\bibliography{cite}

\end{document}\documentclass{article}
\usepackage[utf8]{inputenc}
\usepackage{booktabs}
\usepackage{amsmath}
\usepackage{natbib}

\title{My Document}
\author{Your Name}
\date{\today}

\begin{document}

\maketitle

\section{Introduction}
This is a sample document demonstrating various \LaTeX{} features including mathematical equations, tables, and formatting. You can cite papers using the citation search feature, and they will be automatically added to your \texttt{cite.bib} file.

\section{Mathematical Expressions}
Here are some example equations:

\subsection{Quadratic Formula}
The quadratic formula is given by:
\[x = \frac{-b \pm \sqrt{b^2 - 4ac}}{2a}\]

\subsection{Matrix Example}
A sample matrix:
\[
\begin{pmatrix}
a_{11} & a_{12} & a_{13} \\
a_{21} & a_{22} & a_{23} \\
a_{31} & a_{32} & a_{33}
\end{pmatrix}
\]

\section{Tables}
Below is a sample table showing data:

\begin{table}[h]
\centering
\begin{tabular}{@{}lrrr@{}}
\toprule
\textbf{Category} & \textbf{Value 1} & \textbf{Value 2} & \textbf{Total} \\
\midrule
Item A & 10.5 & 20.3 & 30.8 \\
Item B & 15.2 & 25.1 & 40.3 \\
Item C & 22.1 & 12.7 & 34.8 \\
\bottomrule
\end{tabular}
\caption{Sample Data Table}
\label{tab:sample}
\end{table}

\section{Lists}
\subsection{Bullet Points}
\begin{itemize}
  \item First item
  \item Second item
  \item Third item with sub-items:
    \begin{itemize}
      \item Sub-item 1
      \item Sub-item 2
    \end{itemize}
\end{itemize}

\subsection{Numbered List}
\begin{enumerate}
  \item First step
  \item Second step
  \item Third step
\end{enumerate}

\section{Citations and References}
When you use the citation search feature, references will be automatically managed. For example, you might cite a paper like this: \cite{example2024}. All citations will be automatically added to your \texttt{cite.bib} file.

\section{Conclusion}
This template demonstrates basic \LaTeX{} features including mathematical formulas, tables, and lists. The bibliography system is now set up to work seamlessly with Bibby AI's citation management. Feel free to modify it for your needs!

\bibliographystyle{plain}
\bibliography{cite}

\end{document}\documentclass{article}
\usepackage[utf8]{inputenc}
\usepackage{booktabs}
\usepackage{amsmath}
\usepackage{natbib}

\title{My Document}
\author{Your Name}
\date{\today}

\begin{document}

\maketitle

\section{Introduction}
This is a sample document demonstrating various \LaTeX{} features including mathematical equations, tables, and formatting. You can cite papers using the citation search feature, and they will be automatically added to your \texttt{cite.bib} file.

\section{Mathematical Expressions}
Here are some example equations:

\subsection{Quadratic Formula}
The quadratic formula is given by:
\[x = \frac{-b \pm \sqrt{b^2 - 4ac}}{2a}\]

\subsection{Matrix Example}
A sample matrix:
\[
\begin{pmatrix}
a_{11} & a_{12} & a_{13} \\
a_{21} & a_{22} & a_{23} \\
a_{31} & a_{32} & a_{33}
\end{pmatrix}
\]

\section{Tables}
Below is a sample table showing data:

\begin{table}[h]
\centering
\begin{tabular}{@{}lrrr@{}}
\toprule
\textbf{Category} & \textbf{Value 1} & \textbf{Value 2} & \textbf{Total} \\
\midrule
Item A & 10.5 & 20.3 & 30.8 \\
Item B & 15.2 & 25.1 & 40.3 \\
Item C & 22.1 & 12.7 & 34.8 \\
\bottomrule
\end{tabular}
\caption{Sample Data Table}
\label{tab:sample}
\end{table}

\section{Lists}
\subsection{Bullet Points}
\begin{itemize}
  \item First item
  \item Second item
  \item Third item with sub-items:
    \begin{itemize}
      \item Sub-item 1
      \item Sub-item 2
    \end{itemize}
\end{itemize}

\subsection{Numbered List}
\begin{enumerate}
  \item First step
  \item Second step
  \item Third step
\end{enumerate}

\section{Citations and References}
When you use the citation search feature, references will be automatically managed. For example, you might cite a paper like this: \cite{example2024}. All citations will be automatically added to your \texttt{cite.bib} file.

\section{Conclusion}
This template demonstrates basic \LaTeX{} features including mathematical formulas, tables, and lists. The bibliography system is now set up to work seamlessly with Bibby AI's citation management. Feel free to modify it for your needs!

\bibliographystyle{plain}
\bibliography{cite}

\end{document}\documentclass{article}
\usepackage[utf8]{inputenc}
\usepackage{booktabs}
\usepackage{amsmath}
\usepackage{natbib}

\title{My Document}
\author{Your Name}
\date{\today}

\begin{document}

\maketitle

\section{Introduction}
This is a sample document demonstrating various \LaTeX{} features including mathematical equations, tables, and formatting. You can cite papers using the citation search feature, and they will be automatically added to your \texttt{cite.bib} file.

\section{Mathematical Expressions}
Here are some example equations:

\subsection{Quadratic Formula}
The quadratic formula is given by:
\[x = \frac{-b \pm \sqrt{b^2 - 4ac}}{2a}\]

\subsection{Matrix Example}
A sample matrix:
\[
\begin{pmatrix}
a_{11} & a_{12} & a_{13} \\
a_{21} & a_{22} & a_{23} \\
a_{31} & a_{32} & a_{33}
\end{pmatrix}
\]

\section{Tables}
Below is a sample table showing data:

\begin{table}[h]
\centering
\begin{tabular}{@{}lrrr@{}}
\toprule
\textbf{Category} & \textbf{Value 1} & \textbf{Value 2} & \textbf{Total} \\
\midrule
Item A & 10.5 & 20.3 & 30.8 \\
Item B & 15.2 & 25.1 & 40.3 \\
Item C & 22.1 & 12.7 & 34.8 \\
\bottomrule
\end{tabular}
\caption{Sample Data Table}
\label{tab:sample}
\end{table}

\section{Lists}
\subsection{Bullet Points}
\begin{itemize}
  \item First item
  \item Second item
  \item Third item with sub-items:
    \begin{itemize}
      \item Sub-item 1
      \item Sub-item 2
    \end{itemize}
\end{itemize}

\subsection{Numbered List}
\begin{enumerate}
  \item First step
  \item Second step
  \item Third step
\end{enumerate}

\section{Citations and References}
When you use the citation search feature, references will be automatically managed. For example, you might cite a paper like this: \cite{example2024}. All citations will be automatically added to your \texttt{cite.bib} file.

\section{Conclusion}
This template demonstrates basic \LaTeX{} features including mathematical formulas, tables, and lists. The bibliography system is now set up to work seamlessly with Bibby AI's citation management. Feel free to modify it for your needs!

\bibliographystyle{plain}
\bibliography{cite}

\end{document}\documentclass{article}
\usepackage[utf8]{inputenc}
\usepackage{booktabs}
\usepackage{amsmath}
\usepackage{natbib}

\title{My Document}
\author{Your Name}
\date{\today}

\begin{document}

\maketitle

\section{Introduction}
This is a sample document demonstrating various \LaTeX{} features including mathematical equations, tables, and formatting. You can cite papers using the citation search feature, and they will be automatically added to your \texttt{cite.bib} file.

\section{Mathematical Expressions}
Here are some example equations:

\subsection{Quadratic Formula}
The quadratic formula is given by:
\[x = \frac{-b \pm \sqrt{b^2 - 4ac}}{2a}\]

\subsection{Matrix Example}
A sample matrix:
\[
\begin{pmatrix}
a_{11} & a_{12} & a_{13} \\
a_{21} & a_{22} & a_{23} \\
a_{31} & a_{32} & a_{33}
\end{pmatrix}
\]

\section{Tables}
Below is a sample table showing data:

\begin{table}[h]
\centering
\begin{tabular}{@{}lrrr@{}}
\toprule
\textbf{Category} & \textbf{Value 1} & \textbf{Value 2} & \textbf{Total} \\
\midrule
Item A & 10.5 & 20.3 & 30.8 \\
Item B & 15.2 & 25.1 & 40.3 \\
Item C & 22.1 & 12.7 & 34.8 \\
\bottomrule
\end{tabular}
\caption{Sample Data Table}
\label{tab:sample}
\end{table}

\section{Lists}
\subsection{Bullet Points}
\begin{itemize}
  \item First item
  \item Second item
  \item Third item with sub-items:
    \begin{itemize}
      \item Sub-item 1
      \item Sub-item 2
    \end{itemize}
\end{itemize}

\subsection{Numbered List}
\begin{enumerate}
  \item First step
  \item Second step
  \item Third step
\end{enumerate}

\section{Citations and References}
When you use the citation search feature, references will be automatically managed. For example, you might cite a paper like this: \cite{example2024}. All citations will be automatically added to your \texttt{cite.bib} file.

\section{Conclusion}
This template demonstrates basic \LaTeX{} features including mathematical formulas, tables, and lists. The bibliography system is now set up to work seamlessly with Bibby AI's citation management. Feel free to modify it for your needs!

\bibliographystyle{plain}
\bibliography{cite}

\end{document}\documentclass{article}
\usepackage[utf8]{inputenc}
\usepackage{booktabs}
\usepackage{amsmath}
\usepackage{natbib}

\title{My Document}
\author{Your Name}
\date{\today}

\begin{document}

\maketitle

\section{Introduction}
This is a sample document demonstrating various \LaTeX{} features including mathematical equations, tables, and formatting. You can cite papers using the citation search feature, and they will be automatically added to your \texttt{cite.bib} file.

\section{Mathematical Expressions}
Here are some example equations:

\subsection{Quadratic Formula}
The quadratic formula is given by:
\[x = \frac{-b \pm \sqrt{b^2 - 4ac}}{2a}\]

\subsection{Matrix Example}
A sample matrix:
\[
\begin{pmatrix}
a_{11} & a_{12} & a_{13} \\
a_{21} & a_{22} & a_{23} \\
a_{31} & a_{32} & a_{33}
\end{pmatrix}
\]

\section{Tables}
Below is a sample table showing data:

\begin{table}[h]
\centering
\begin{tabular}{@{}lrrr@{}}
\toprule
\textbf{Category} & \textbf{Value 1} & \textbf{Value 2} & \textbf{Total} \\
\midrule
Item A & 10.5 & 20.3 & 30.8 \\
Item B & 15.2 & 25.1 & 40.3 \\
Item C & 22.1 & 12.7 & 34.8 \\
\bottomrule
\end{tabular}
\caption{Sample Data Table}
\label{tab:sample}
\end{table}

\section{Lists}
\subsection{Bullet Points}
\begin{itemize}
  \item First item
  \item Second item
  \item Third item with sub-items:
    \begin{itemize}
      \item Sub-item 1
      \item Sub-item 2
    \end{itemize}
\end{itemize}

\subsection{Numbered List}
\begin{enumerate}
  \item First step
  \item Second step
  \item Third step
\end{enumerate}

\section{Citations and References}
When you use the citation search feature, references will be automatically managed. For example, you might cite a paper like this: \cite{example2024}. All citations will be automatically added to your \texttt{cite.bib} file.

\section{Conclusion}
This template demonstrates basic \LaTeX{} features including mathematical formulas, tables, and lists. The bibliography system is now set up to work seamlessly with Bibby AI's citation management. Feel free to modify it for your needs!

\bibliographystyle{plain}
\bibliography{cite}

\end{document}\documentclass{article}
\usepackage[utf8]{inputenc}
\usepackage{booktabs}
\usepackage{amsmath}
\usepackage{natbib}

\title{My Document}
\author{Your Name}
\date{\today}

\begin{document}

\maketitle

\section{Introduction}
This is a sample document demonstrating various \LaTeX{} features including mathematical equations, tables, and formatting. You can cite papers using the citation search feature, and they will be automatically added to your \texttt{cite.bib} file.

\section{Mathematical Expressions}
Here are some example equations:

\subsection{Quadratic Formula}
The quadratic formula is given by:
\[x = \frac{-b \pm \sqrt{b^2 - 4ac}}{2a}\]

\subsection{Matrix Example}
A sample matrix:
\[
\begin{pmatrix}
a_{11} & a_{12} & a_{13} \\
a_{21} & a_{22} & a_{23} \\
a_{31} & a_{32} & a_{33}
\end{pmatrix}
\]

\section{Tables}
Below is a sample table showing data:

\begin{table}[h]
\centering
\begin{tabular}{@{}lrrr@{}}
\toprule
\textbf{Category} & \textbf{Value 1} & \textbf{Value 2} & \textbf{Total} \\
\midrule
Item A & 10.5 & 20.3 & 30.8 \\
Item B & 15.2 & 25.1 & 40.3 \\
Item C & 22.1 & 12.7 & 34.8 \\
\bottomrule
\end{tabular}
\caption{Sample Data Table}
\label{tab:sample}
\end{table}

\section{Lists}
\subsection{Bullet Points}
\begin{itemize}
  \item First item
  \item Second item
  \item Third item with sub-items:
    \begin{itemize}
      \item Sub-item 1
      \item Sub-item 2
    \end{itemize}
\end{itemize}

\subsection{Numbered List}
\begin{enumerate}
  \item First step
  \item Second step
  \item Third step
\end{enumerate}

\section{Citations and References}
When you use the citation search feature, references will be automatically managed. For example, you might cite a paper like this: \cite{example2024}. All citations will be automatically added to your \texttt{cite.bib} file.

\section{Conclusion}
This template demonstrates basic \LaTeX{} features including mathematical formulas, tables, and lists. The bibliography system is now set up to work seamlessly with Bibby AI's citation management. Feel free to modify it for your needs!

\bibliographystyle{plain}
\bibliography{cite}

\end{document}\documentclass{article}
\usepackage[utf8]{inputenc}
\usepackage{booktabs}
\usepackage{amsmath}
\usepackage{natbib}

\title{My Document}
\author{Your Name}
\date{\today}

\begin{document}

\maketitle

\section{Introduction}
This is a sample document demonstrating various \LaTeX{} features including mathematical equations, tables, and formatting. You can cite papers using the citation search feature, and they will be automatically added to your \texttt{cite.bib} file.

\section{Mathematical Expressions}
Here are some example equations:

\subsection{Quadratic Formula}
The quadratic formula is given by:
\[x = \frac{-b \pm \sqrt{b^2 - 4ac}}{2a}\]

\subsection{Matrix Example}
A sample matrix:
\[
\begin{pmatrix}
a_{11} & a_{12} & a_{13} \\
a_{21} & a_{22} & a_{23} \\
a_{31} & a_{32} & a_{33}
\end{pmatrix}
\]

\section{Tables}
Below is a sample table showing data:

\begin{table}[h]
\centering
\begin{tabular}{@{}lrrr@{}}
\toprule
\textbf{Category} & \textbf{Value 1} & \textbf{Value 2} & \textbf{Total} \\
\midrule
Item A & 10.5 & 20.3 & 30.8 \\
Item B & 15.2 & 25.1 & 40.3 \\
Item C & 22.1 & 12.7 & 34.8 \\
\bottomrule
\end{tabular}
\caption{Sample Data Table}
\label{tab:sample}
\end{table}

\section{Lists}
\subsection{Bullet Points}
\begin{itemize}
  \item First item
  \item Second item
  \item Third item with sub-items:
    \begin{itemize}
      \item Sub-item 1
      \item Sub-item 2
    \end{itemize}
\end{itemize}

\subsection{Numbered List}
\begin{enumerate}
  \item First step
  \item Second step
  \item Third step
\end{enumerate}

\section{Citations and References}
When you use the citation search feature, references will be automatically managed. For example, you might cite a paper like this: \cite{example2024}. All citations will be automatically added to your \texttt{cite.bib} file.

\section{Conclusion}
This template demonstrates basic \LaTeX{} features including mathematical formulas, tables, and lists. The bibliography system is now set up to work seamlessly with Bibby AI's citation management. Feel free to modify it for your needs!

\bibliographystyle{plain}
\bibliography{cite}

\end{document}\documentclass{article}
\usepackage[utf8]{inputenc}
\usepackage{booktabs}
\usepackage{amsmath}
\usepackage{natbib}

\title{My Document}
\author{Your Name}
\date{\today}

\begin{document}

\maketitle

\section{Introduction}
This is a sample document demonstrating various \LaTeX{} features including mathematical equations, tables, and formatting. You can cite papers using the citation search feature, and they will be automatically added to your \texttt{cite.bib} file.

\section{Mathematical Expressions}
Here are some example equations:

\subsection{Quadratic Formula}
The quadratic formula is given by:
\[x = \frac{-b \pm \sqrt{b^2 - 4ac}}{2a}\]

\subsection{Matrix Example}
A sample matrix:
\[
\begin{pmatrix}
a_{11} & a_{12} & a_{13} \\
a_{21} & a_{22} & a_{23} \\
a_{31} & a_{32} & a_{33}
\end{pmatrix}
\]

\section{Tables}
Below is a sample table showing data:

\begin{table}[h]
\centering
\begin{tabular}{@{}lrrr@{}}
\toprule
\textbf{Category} & \textbf{Value 1} & \textbf{Value 2} & \textbf{Total} \\
\midrule
Item A & 10.5 & 20.3 & 30.8 \\
Item B & 15.2 & 25.1 & 40.3 \\
Item C & 22.1 & 12.7 & 34.8 \\
\bottomrule
\end{tabular}
\caption{Sample Data Table}
\label{tab:sample}
\end{table}

\section{Lists}
\subsection{Bullet Points}
\begin{itemize}
  \item First item
  \item Second item
  \item Third item with sub-items:
    \begin{itemize}
      \item Sub-item 1
      \item Sub-item 2
    \end{itemize}
\end{itemize}

\subsection{Numbered List}
\begin{enumerate}
  \item First step
  \item Second step
  \item Third step
\end{enumerate}

\section{Citations and References}
When you use the citation search feature, references will be automatically managed. For example, you might cite a paper like this: \cite{example2024}. All citations will be automatically added to your \texttt{cite.bib} file.

\section{Conclusion}
This template demonstrates basic \LaTeX{} features including mathematical formulas, tables, and lists. The bibliography system is now set up to work seamlessly with Bibby AI's citation management. Feel free to modify it for your needs!

\bibliographystyle{plain}
\bibliography{cite}

\end{document}\documentclass{article}
\usepackage[utf8]{inputenc}
\usepackage{booktabs}
\usepackage{amsmath}
\usepackage{natbib}

\title{My Document}
\author{Your Name}
\date{\today}

\begin{document}

\maketitle

\section{Introduction}
This is a sample document demonstrating various \LaTeX{} features including mathematical equations, tables, and formatting. You can cite papers using the citation search feature, and they will be automatically added to your \texttt{cite.bib} file.

\section{Mathematical Expressions}
Here are some example equations:

\subsection{Quadratic Formula}
The quadratic formula is given by:
\[x = \frac{-b \pm \sqrt{b^2 - 4ac}}{2a}\]

\subsection{Matrix Example}
A sample matrix:
\[
\begin{pmatrix}
a_{11} & a_{12} & a_{13} \\
a_{21} & a_{22} & a_{23} \\
a_{31} & a_{32} & a_{33}
\end{pmatrix}
\]

\section{Tables}
Below is a sample table showing data:

\begin{table}[h]
\centering
\begin{tabular}{@{}lrrr@{}}
\toprule
\textbf{Category} & \textbf{Value 1} & \textbf{Value 2} & \textbf{Total} \\
\midrule
Item A & 10.5 & 20.3 & 30.8 \\
Item B & 15.2 & 25.1 & 40.3 \\
Item C & 22.1 & 12.7 & 34.8 \\
\bottomrule
\end{tabular}
\caption{Sample Data Table}
\label{tab:sample}
\end{table}

\section{Lists}
\subsection{Bullet Points}
\begin{itemize}
  \item First item
  \item Second item
  \item Third item with sub-items:
    \begin{itemize}
      \item Sub-item 1
      \item Sub-item 2
    \end{itemize}
\end{itemize}

\subsection{Numbered List}
\begin{enumerate}
  \item First step
  \item Second step
  \item Third step
\end{enumerate}

\section{Citations and References}
When you use the citation search feature, references will be automatically managed. For example, you might cite a paper like this: \cite{example2024}. All citations will be automatically added to your \texttt{cite.bib} file.

\section{Conclusion}
This template demonstrates basic \LaTeX{} features including mathematical formulas, tables, and lists. The bibliography system is now set up to work seamlessly with Bibby AI's citation management. Feel free to modify it for your needs!

\bibliographystyle{plain}
\bibliography{cite}

\end{document}\documentclass{article}
\usepackage[utf8]{inputenc}
\usepackage{booktabs}
\usepackage{amsmath}
\usepackage{natbib}

\title{My Document}
\author{Your Name}
\date{\today}

\begin{document}

\maketitle

\section{Introduction}
This is a sample document demonstrating various \LaTeX{} features including mathematical equations, tables, and formatting. You can cite papers using the citation search feature, and they will be automatically added to your \texttt{cite.bib} file.

\section{Mathematical Expressions}
Here are some example equations:

\subsection{Quadratic Formula}
The quadratic formula is given by:
\[x = \frac{-b \pm \sqrt{b^2 - 4ac}}{2a}\]

\subsection{Matrix Example}
A sample matrix:
\[
\begin{pmatrix}
a_{11} & a_{12} & a_{13} \\
a_{21} & a_{22} & a_{23} \\
a_{31} & a_{32} & a_{33}
\end{pmatrix}
\]

\section{Tables}
Below is a sample table showing data:

\begin{table}[h]
\centering
\begin{tabular}{@{}lrrr@{}}
\toprule
\textbf{Category} & \textbf{Value 1} & \textbf{Value 2} & \textbf{Total} \\
\midrule
Item A & 10.5 & 20.3 & 30.8 \\
Item B & 15.2 & 25.1 & 40.3 \\
Item C & 22.1 & 12.7 & 34.8 \\
\bottomrule
\end{tabular}
\caption{Sample Data Table}
\label{tab:sample}
\end{table}

\section{Lists}
\subsection{Bullet Points}
\begin{itemize}
  \item First item
  \item Second item
  \item Third item with sub-items:
    \begin{itemize}
      \item Sub-item 1
      \item Sub-item 2
    \end{itemize}
\end{itemize}

\subsection{Numbered List}
\begin{enumerate}
  \item First step
  \item Second step
  \item Third step
\end{enumerate}

\section{Citations and References}
When you use the citation search feature, references will be automatically managed. For example, you might cite a paper like this: \cite{example2024}. All citations will be automatically added to your \texttt{cite.bib} file.

\section{Conclusion}
This template demonstrates basic \LaTeX{} features including mathematical formulas, tables, and lists. The bibliography system is now set up to work seamlessly with Bibby AI's citation management. Feel free to modify it for your needs!

\bibliographystyle{plain}
\bibliography{cite}

\end{document}\documentclass{article}
\usepackage[utf8]{inputenc}
\usepackage{booktabs}
\usepackage{amsmath}
\usepackage{natbib}

\title{My Document}
\author{Your Name}
\date{\today}

\begin{document}

\maketitle

\section{Introduction}
This is a sample document demonstrating various \LaTeX{} features including mathematical equations, tables, and formatting. You can cite papers using the citation search feature, and they will be automatically added to your \texttt{cite.bib} file.

\section{Mathematical Expressions}
Here are some example equations:

\subsection{Quadratic Formula}
The quadratic formula is given by:
\[x = \frac{-b \pm \sqrt{b^2 - 4ac}}{2a}\]

\subsection{Matrix Example}
A sample matrix:
\[
\begin{pmatrix}
a_{11} & a_{12} & a_{13} \\
a_{21} & a_{22} & a_{23} \\
a_{31} & a_{32} & a_{33}
\end{pmatrix}
\]

\section{Tables}
Below is a sample table showing data:

\begin{table}[h]
\centering
\begin{tabular}{@{}lrrr@{}}
\toprule
\textbf{Category} & \textbf{Value 1} & \textbf{Value 2} & \textbf{Total} \\
\midrule
Item A & 10.5 & 20.3 & 30.8 \\
Item B & 15.2 & 25.1 & 40.3 \\
Item C & 22.1 & 12.7 & 34.8 \\
\bottomrule
\end{tabular}
\caption{Sample Data Table}
\label{tab:sample}
\end{table}

\section{Lists}
\subsection{Bullet Points}
\begin{itemize}
  \item First item
  \item Second item
  \item Third item with sub-items:
    \begin{itemize}
      \item Sub-item 1
      \item Sub-item 2
    \end{itemize}
\end{itemize}

\subsection{Numbered List}
\begin{enumerate}
  \item First step
  \item Second step
  \item Third step
\end{enumerate}

\section{Citations and References}
When you use the citation search feature, references will be automatically managed. For example, you might cite a paper like this: \cite{example2024}. All citations will be automatically added to your \texttt{cite.bib} file.

\section{Conclusion}
This template demonstrates basic \LaTeX{} features including mathematical formulas, tables, and lists. The bibliography system is now set up to work seamlessly with Bibby AI's citation management. Feel free to modify it for your needs!

\bibliographystyle{plain}
\bibliography{cite}

\end{document}\documentclass{article}
\usepackage[utf8]{inputenc}
\usepackage{booktabs}
\usepackage{amsmath}
\usepackage{natbib}

\title{My Document}
\author{Your Name}
\date{\today}

\begin{document}

\maketitle

\section{Introduction}
This is a sample document demonstrating various \LaTeX{} features including mathematical equations, tables, and formatting. You can cite papers using the citation search feature, and they will be automatically added to your \texttt{cite.bib} file.

\section{Mathematical Expressions}
Here are some example equations:

\subsection{Quadratic Formula}
The quadratic formula is given by:
\[x = \frac{-b \pm \sqrt{b^2 - 4ac}}{2a}\]

\subsection{Matrix Example}
A sample matrix:
\[
\begin{pmatrix}
a_{11} & a_{12} & a_{13} \\
a_{21} & a_{22} & a_{23} \\
a_{31} & a_{32} & a_{33}
\end{pmatrix}
\]

\section{Tables}
Below is a sample table showing data:

\begin{table}[h]
\centering
\begin{tabular}{@{}lrrr@{}}
\toprule
\textbf{Category} & \textbf{Value 1} & \textbf{Value 2} & \textbf{Total} \\
\midrule
Item A & 10.5 & 20.3 & 30.8 \\
Item B & 15.2 & 25.1 & 40.3 \\
Item C & 22.1 & 12.7 & 34.8 \\
\bottomrule
\end{tabular}
\caption{Sample Data Table}
\label{tab:sample}
\end{table}

\section{Lists}
\subsection{Bullet Points}
\begin{itemize}
  \item First item
  \item Second item
  \item Third item with sub-items:
    \begin{itemize}
      \item Sub-item 1
      \item Sub-item 2
    \end{itemize}
\end{itemize}

\subsection{Numbered List}
\begin{enumerate}
  \item First step
  \item Second step
  \item Third step
\end{enumerate}

\section{Citations and References}
When you use the citation search feature, references will be automatically managed. For example, you might cite a paper like this: \cite{example2024}. All citations will be automatically added to your \texttt{cite.bib} file.

\section{Conclusion}
This template demonstrates basic \LaTeX{} features including mathematical formulas, tables, and lists. The bibliography system is now set up to work seamlessly with Bibby AI's citation management. Feel free to modify it for your needs!

\bibliographystyle{plain}
\bibliography{cite}

\end{document}\documentclass{article}
\usepackage[utf8]{inputenc}
\usepackage{booktabs}
\usepackage{amsmath}
\usepackage{natbib}

\title{My Document}
\author{Your Name}
\date{\today}

\begin{document}

\maketitle

\section{Introduction}
This is a sample document demonstrating various \LaTeX{} features including mathematical equations, tables, and formatting. You can cite papers using the citation search feature, and they will be automatically added to your \texttt{cite.bib} file.

\section{Mathematical Expressions}
Here are some example equations:

\subsection{Quadratic Formula}
The quadratic formula is given by:
\[x = \frac{-b \pm \sqrt{b^2 - 4ac}}{2a}\]

\subsection{Matrix Example}
A sample matrix:
\[
\begin{pmatrix}
a_{11} & a_{12} & a_{13} \\
a_{21} & a_{22} & a_{23} \\
a_{31} & a_{32} & a_{33}
\end{pmatrix}
\]

\section{Tables}
Below is a sample table showing data:

\begin{table}[h]
\centering
\begin{tabular}{@{}lrrr@{}}
\toprule
\textbf{Category} & \textbf{Value 1} & \textbf{Value 2} & \textbf{Total} \\
\midrule
Item A & 10.5 & 20.3 & 30.8 \\
Item B & 15.2 & 25.1 & 40.3 \\
Item C & 22.1 & 12.7 & 34.8 \\
\bottomrule
\end{tabular}
\caption{Sample Data Table}
\label{tab:sample}
\end{table}

\section{Lists}
\subsection{Bullet Points}
\begin{itemize}
  \item First item
  \item Second item
  \item Third item with sub-items:
    \begin{itemize}
      \item Sub-item 1
      \item Sub-item 2
    \end{itemize}
\end{itemize}

\subsection{Numbered List}
\begin{enumerate}
  \item First step
  \item Second step
  \item Third step
\end{enumerate}

\section{Citations and References}
When you use the citation search feature, references will be automatically managed. For example, you might cite a paper like this: \cite{example2024}. All citations will be automatically added to your \texttt{cite.bib} file.

\section{Conclusion}
This template demonstrates basic \LaTeX{} features including mathematical formulas, tables, and lists. The bibliography system is now set up to work seamlessly with Bibby AI's citation management. Feel free to modify it for your needs!

\bibliographystyle{plain}
\bibliography{cite}

\end{document}\documentclass{article}
\usepackage[utf8]{inputenc}
\usepackage{booktabs}
\usepackage{amsmath}
\usepackage{natbib}

\title{My Document}
\author{Your Name}
\date{\today}

\begin{document}

\maketitle

\section{Introduction}
This is a sample document demonstrating various \LaTeX{} features including mathematical equations, tables, and formatting. You can cite papers using the citation search feature, and they will be automatically added to your \texttt{cite.bib} file.

\section{Mathematical Expressions}
Here are some example equations:

\subsection{Quadratic Formula}
The quadratic formula is given by:
\[x = \frac{-b \pm \sqrt{b^2 - 4ac}}{2a}\]

\subsection{Matrix Example}
A sample matrix:
\[
\begin{pmatrix}
a_{11} & a_{12} & a_{13} \\
a_{21} & a_{22} & a_{23} \\
a_{31} & a_{32} & a_{33}
\end{pmatrix}
\]

\section{Tables}
Below is a sample table showing data:

\begin{table}[h]
\centering
\begin{tabular}{@{}lrrr@{}}
\toprule
\textbf{Category} & \textbf{Value 1} & \textbf{Value 2} & \textbf{Total} \\
\midrule
Item A & 10.5 & 20.3 & 30.8 \\
Item B & 15.2 & 25.1 & 40.3 \\
Item C & 22.1 & 12.7 & 34.8 \\
\bottomrule
\end{tabular}
\caption{Sample Data Table}
\label{tab:sample}
\end{table}

\section{Lists}
\subsection{Bullet Points}
\begin{itemize}
  \item First item
  \item Second item
  \item Third item with sub-items:
    \begin{itemize}
      \item Sub-item 1
      \item Sub-item 2
    \end{itemize}
\end{itemize}

\subsection{Numbered List}
\begin{enumerate}
  \item First step
  \item Second step
  \item Third step
\end{enumerate}

\section{Citations and References}
When you use the citation search feature, references will be automatically managed. For example, you might cite a paper like this: \cite{example2024}. All citations will be automatically added to your \texttt{cite.bib} file.

\section{Conclusion}
This template demonstrates basic \LaTeX{} features including mathematical formulas, tables, and lists. The bibliography system is now set up to work seamlessly with Bibby AI's citation management. Feel free to modify it for your needs!

\bibliographystyle{plain}
\bibliography{cite}

\end{document}\documentclass{article}
\usepackage[utf8]{inputenc}
\usepackage{booktabs}
\usepackage{amsmath}
\usepackage{natbib}

\title{My Document}
\author{Your Name}
\date{\today}

\begin{document}

\maketitle

\section{Introduction}
This is a sample document demonstrating various \LaTeX{} features including mathematical equations, tables, and formatting. You can cite papers using the citation search feature, and they will be automatically added to your \texttt{cite.bib} file.

\section{Mathematical Expressions}
Here are some example equations:

\subsection{Quadratic Formula}
The quadratic formula is given by:
\[x = \frac{-b \pm \sqrt{b^2 - 4ac}}{2a}\]

\subsection{Matrix Example}
A sample matrix:
\[
\begin{pmatrix}
a_{11} & a_{12} & a_{13} \\
a_{21} & a_{22} & a_{23} \\
a_{31} & a_{32} & a_{33}
\end{pmatrix}
\]

\section{Tables}
Below is a sample table showing data:

\begin{table}[h]
\centering
\begin{tabular}{@{}lrrr@{}}
\toprule
\textbf{Category} & \textbf{Value 1} & \textbf{Value 2} & \textbf{Total} \\
\midrule
Item A & 10.5 & 20.3 & 30.8 \\
Item B & 15.2 & 25.1 & 40.3 \\
Item C & 22.1 & 12.7 & 34.8 \\
\bottomrule
\end{tabular}
\caption{Sample Data Table}
\label{tab:sample}
\end{table}

\section{Lists}
\subsection{Bullet Points}
\begin{itemize}
  \item First item
  \item Second item
  \item Third item with sub-items:
    \begin{itemize}
      \item Sub-item 1
      \item Sub-item 2
    \end{itemize}
\end{itemize}

\subsection{Numbered List}
\begin{enumerate}
  \item First step
  \item Second step
  \item Third step
\end{enumerate}

\section{Citations and References}
When you use the citation search feature, references will be automatically managed. For example, you might cite a paper like this: \cite{example2024}. All citations will be automatically added to your \texttt{cite.bib} file.

\section{Conclusion}
This template demonstrates basic \LaTeX{} features including mathematical formulas, tables, and lists. The bibliography system is now set up to work seamlessly with Bibby AI's citation management. Feel free to modify it for your needs!

\bibliographystyle{plain}
\bibliography{cite}

\end{document}\documentclass{article}
\usepackage[utf8]{inputenc}
\usepackage{booktabs}
\usepackage{amsmath}
\usepackage{natbib}

\title{My Document}
\author{Your Name}
\date{\today}

\begin{document}

\maketitle

\section{Introduction}
This is a sample document demonstrating various \LaTeX{} features including mathematical equations, tables, and formatting. You can cite papers using the citation search feature, and they will be automatically added to your \texttt{cite.bib} file.

\section{Mathematical Expressions}
Here are some example equations:

\subsection{Quadratic Formula}
The quadratic formula is given by:
\[x = \frac{-b \pm \sqrt{b^2 - 4ac}}{2a}\]

\subsection{Matrix Example}
A sample matrix:
\[
\begin{pmatrix}
a_{11} & a_{12} & a_{13} \\
a_{21} & a_{22} & a_{23} \\
a_{31} & a_{32} & a_{33}
\end{pmatrix}
\]

\section{Tables}
Below is a sample table showing data:

\begin{table}[h]
\centering
\begin{tabular}{@{}lrrr@{}}
\toprule
\textbf{Category} & \textbf{Value 1} & \textbf{Value 2} & \textbf{Total} \\
\midrule
Item A & 10.5 & 20.3 & 30.8 \\
Item B & 15.2 & 25.1 & 40.3 \\
Item C & 22.1 & 12.7 & 34.8 \\
\bottomrule
\end{tabular}
\caption{Sample Data Table}
\label{tab:sample}
\end{table}

\section{Lists}
\subsection{Bullet Points}
\begin{itemize}
  \item First item
  \item Second item
  \item Third item with sub-items:
    \begin{itemize}
      \item Sub-item 1
      \item Sub-item 2
    \end{itemize}
\end{itemize}

\subsection{Numbered List}
\begin{enumerate}
  \item First step
  \item Second step
  \item Third step
\end{enumerate}

\section{Citations and References}
When you use the citation search feature, references will be automatically managed. For example, you might cite a paper like this: \cite{example2024}. All citations will be automatically added to your \texttt{cite.bib} file.

\section{Conclusion}
This template demonstrates basic \LaTeX{} features including mathematical formulas, tables, and lists. The bibliography system is now set up to work seamlessly with Bibby AI's citation management. Feel free to modify it for your needs!

\bibliographystyle{plain}
\bibliography{cite}

\end{document}\documentclass{article}
\usepackage[utf8]{inputenc}
\usepackage{booktabs}
\usepackage{amsmath}
\usepackage{natbib}

\title{My Document}
\author{Your Name}
\date{\today}

\begin{document}

\maketitle

\section{Introduction}
This is a sample document demonstrating various \LaTeX{} features including mathematical equations, tables, and formatting. You can cite papers using the citation search feature, and they will be automatically added to your \texttt{cite.bib} file.

\section{Mathematical Expressions}
Here are some example equations:

\subsection{Quadratic Formula}
The quadratic formula is given by:
\[x = \frac{-b \pm \sqrt{b^2 - 4ac}}{2a}\]

\subsection{Matrix Example}
A sample matrix:
\[
\begin{pmatrix}
a_{11} & a_{12} & a_{13} \\
a_{21} & a_{22} & a_{23} \\
a_{31} & a_{32} & a_{33}
\end{pmatrix}
\]

\section{Tables}
Below is a sample table showing data:

\begin{table}[h]
\centering
\begin{tabular}{@{}lrrr@{}}
\toprule
\textbf{Category} & \textbf{Value 1} & \textbf{Value 2} & \textbf{Total} \\
\midrule
Item A & 10.5 & 20.3 & 30.8 \\
Item B & 15.2 & 25.1 & 40.3 \\
Item C & 22.1 & 12.7 & 34.8 \\
\bottomrule
\end{tabular}
\caption{Sample Data Table}
\label{tab:sample}
\end{table}

\section{Lists}
\subsection{Bullet Points}
\begin{itemize}
  \item First item
  \item Second item
  \item Third item with sub-items:
    \begin{itemize}
      \item Sub-item 1
      \item Sub-item 2
    \end{itemize}
\end{itemize}

\subsection{Numbered List}
\begin{enumerate}
  \item First step
  \item Second step
  \item Third step
\end{enumerate}

\section{Citations and References}
When you use the citation search feature, references will be automatically managed. For example, you might cite a paper like this: \cite{example2024}. All citations will be automatically added to your \texttt{cite.bib} file.

\section{Conclusion}
This template demonstrates basic \LaTeX{} features including mathematical formulas, tables, and lists. The bibliography system is now set up to work seamlessly with Bibby AI's citation management. Feel free to modify it for your needs!

\bibliographystyle{plain}
\bibliography{cite}

\end{document}\documentclass{article}
\usepackage[utf8]{inputenc}
\usepackage{booktabs}
\usepackage{amsmath}
\usepackage{natbib}

\title{My Document}
\author{Your Name}
\date{\today}

\begin{document}

\maketitle

\section{Introduction}
This is a sample document demonstrating various \LaTeX{} features including mathematical equations, tables, and formatting. You can cite papers using the citation search feature, and they will be automatically added to your \texttt{cite.bib} file.

\section{Mathematical Expressions}
Here are some example equations:

\subsection{Quadratic Formula}
The quadratic formula is given by:
\[x = \frac{-b \pm \sqrt{b^2 - 4ac}}{2a}\]

\subsection{Matrix Example}
A sample matrix:
\[
\begin{pmatrix}
a_{11} & a_{12} & a_{13} \\
a_{21} & a_{22} & a_{23} \\
a_{31} & a_{32} & a_{33}
\end{pmatrix}
\]

\section{Tables}
Below is a sample table showing data:

\begin{table}[h]
\centering
\begin{tabular}{@{}lrrr@{}}
\toprule
\textbf{Category} & \textbf{Value 1} & \textbf{Value 2} & \textbf{Total} \\
\midrule
Item A & 10.5 & 20.3 & 30.8 \\
Item B & 15.2 & 25.1 & 40.3 \\
Item C & 22.1 & 12.7 & 34.8 \\
\bottomrule
\end{tabular}
\caption{Sample Data Table}
\label{tab:sample}
\end{table}

\section{Lists}
\subsection{Bullet Points}
\begin{itemize}
  \item First item
  \item Second item
  \item Third item with sub-items:
    \begin{itemize}
      \item Sub-item 1
      \item Sub-item 2
    \end{itemize}
\end{itemize}

\subsection{Numbered List}
\begin{enumerate}
  \item First step
  \item Second step
  \item Third step
\end{enumerate}

\section{Citations and References}
When you use the citation search feature, references will be automatically managed. For example, you might cite a paper like this: \cite{example2024}. All citations will be automatically added to your \texttt{cite.bib} file.

\section{Conclusion}
This template demonstrates basic \LaTeX{} features including mathematical formulas, tables, and lists. The bibliography system is now set up to work seamlessly with Bibby AI's citation management. Feel free to modify it for your needs!

\bibliographystyle{plain}
\bibliography{cite}

\end{document}\documentclass{article}
\usepackage[utf8]{inputenc}
\usepackage{booktabs}
\usepackage{amsmath}
\usepackage{natbib}

\title{My Document}
\author{Your Name}
\date{\today}

\begin{document}

\maketitle

\section{Introduction}
This is a sample document demonstrating various \LaTeX{} features including mathematical equations, tables, and formatting. You can cite papers using the citation search feature, and they will be automatically added to your \texttt{cite.bib} file.

\section{Mathematical Expressions}
Here are some example equations:

\subsection{Quadratic Formula}
The quadratic formula is given by:
\[x = \frac{-b \pm \sqrt{b^2 - 4ac}}{2a}\]

\subsection{Matrix Example}
A sample matrix:
\[
\begin{pmatrix}
a_{11} & a_{12} & a_{13} \\
a_{21} & a_{22} & a_{23} \\
a_{31} & a_{32} & a_{33}
\end{pmatrix}
\]

\section{Tables}
Below is a sample table showing data:

\begin{table}[h]
\centering
\begin{tabular}{@{}lrrr@{}}
\toprule
\textbf{Category} & \textbf{Value 1} & \textbf{Value 2} & \textbf{Total} \\
\midrule
Item A & 10.5 & 20.3 & 30.8 \\
Item B & 15.2 & 25.1 & 40.3 \\
Item C & 22.1 & 12.7 & 34.8 \\
\bottomrule
\end{tabular}
\caption{Sample Data Table}
\label{tab:sample}
\end{table}

\section{Lists}
\subsection{Bullet Points}
\begin{itemize}
  \item First item
  \item Second item
  \item Third item with sub-items:
    \begin{itemize}
      \item Sub-item 1
      \item Sub-item 2
    \end{itemize}
\end{itemize}

\subsection{Numbered List}
\begin{enumerate}
  \item First step
  \item Second step
  \item Third step
\end{enumerate}

\section{Citations and References}
When you use the citation search feature, references will be automatically managed. For example, you might cite a paper like this: \cite{example2024}. All citations will be automatically added to your \texttt{cite.bib} file.

\section{Conclusion}
This template demonstrates basic \LaTeX{} features including mathematical formulas, tables, and lists. The bibliography system is now set up to work seamlessly with Bibby AI's citation management. Feel free to modify it for your needs!

\bibliographystyle{plain}
\bibliography{cite}

\end{document}\documentclass{article}
\usepackage[utf8]{inputenc}
\usepackage{booktabs}
\usepackage{amsmath}
\usepackage{natbib}

\title{My Document}
\author{Your Name}
\date{\today}

\begin{document}

\maketitle

\section{Introduction}
This is a sample document demonstrating various \LaTeX{} features including mathematical equations, tables, and formatting. You can cite papers using the citation search feature, and they will be automatically added to your \texttt{cite.bib} file.

\section{Mathematical Expressions}
Here are some example equations:

\subsection{Quadratic Formula}
The quadratic formula is given by:
\[x = \frac{-b \pm \sqrt{b^2 - 4ac}}{2a}\]

\subsection{Matrix Example}
A sample matrix:
\[
\begin{pmatrix}
a_{11} & a_{12} & a_{13} \\
a_{21} & a_{22} & a_{23} \\
a_{31} & a_{32} & a_{33}
\end{pmatrix}
\]

\section{Tables}
Below is a sample table showing data:

\begin{table}[h]
\centering
\begin{tabular}{@{}lrrr@{}}
\toprule
\textbf{Category} & \textbf{Value 1} & \textbf{Value 2} & \textbf{Total} \\
\midrule
Item A & 10.5 & 20.3 & 30.8 \\
Item B & 15.2 & 25.1 & 40.3 \\
Item C & 22.1 & 12.7 & 34.8 \\
\bottomrule
\end{tabular}
\caption{Sample Data Table}
\label{tab:sample}
\end{table}

\section{Lists}
\subsection{Bullet Points}
\begin{itemize}
  \item First item
  \item Second item
  \item Third item with sub-items:
    \begin{itemize}
      \item Sub-item 1
      \item Sub-item 2
    \end{itemize}
\end{itemize}

\subsection{Numbered List}
\begin{enumerate}
  \item First step
  \item Second step
  \item Third step
\end{enumerate}

\section{Citations and References}
When you use the citation search feature, references will be automatically managed. For example, you might cite a paper like this: \cite{example2024}. All citations will be automatically added to your \texttt{cite.bib} file.

\section{Conclusion}
This template demonstrates basic \LaTeX{} features including mathematical formulas, tables, and lists. The bibliography system is now set up to work seamlessly with Bibby AI's citation management. Feel free to modify it for your needs!

\bibliographystyle{plain}
\bibliography{cite}

\end{document}\documentclass{article}
\usepackage[utf8]{inputenc}
\usepackage{booktabs}
\usepackage{amsmath}
\usepackage{natbib}

\title{My Document}
\author{Your Name}
\date{\today}

\begin{document}

\maketitle

\section{Introduction}
This is a sample document demonstrating various \LaTeX{} features including mathematical equations, tables, and formatting. You can cite papers using the citation search feature, and they will be automatically added to your \texttt{cite.bib} file.

\section{Mathematical Expressions}
Here are some example equations:

\subsection{Quadratic Formula}
The quadratic formula is given by:
\[x = \frac{-b \pm \sqrt{b^2 - 4ac}}{2a}\]

\subsection{Matrix Example}
A sample matrix:
\[
\begin{pmatrix}
a_{11} & a_{12} & a_{13} \\
a_{21} & a_{22} & a_{23} \\
a_{31} & a_{32} & a_{33}
\end{pmatrix}
\]

\section{Tables}
Below is a sample table showing data:

\begin{table}[h]
\centering
\begin{tabular}{@{}lrrr@{}}
\toprule
\textbf{Category} & \textbf{Value 1} & \textbf{Value 2} & \textbf{Total} \\
\midrule
Item A & 10.5 & 20.3 & 30.8 \\
Item B & 15.2 & 25.1 & 40.3 \\
Item C & 22.1 & 12.7 & 34.8 \\
\bottomrule
\end{tabular}
\caption{Sample Data Table}
\label{tab:sample}
\end{table}

\section{Lists}
\subsection{Bullet Points}
\begin{itemize}
  \item First item
  \item Second item
  \item Third item with sub-items:
    \begin{itemize}
      \item Sub-item 1
      \item Sub-item 2
    \end{itemize}
\end{itemize}

\subsection{Numbered List}
\begin{enumerate}
  \item First step
  \item Second step
  \item Third step
\end{enumerate}

\section{Citations and References}
When you use the citation search feature, references will be automatically managed. For example, you might cite a paper like this: \cite{example2024}. All citations will be automatically added to your \texttt{cite.bib} file.

\section{Conclusion}
This template demonstrates basic \LaTeX{} features including mathematical formulas, tables, and lists. The bibliography system is now set up to work seamlessly with Bibby AI's citation management. Feel free to modify it for your needs!

\bibliographystyle{plain}
\bibliography{cite}

\end{document}\documentclass{article}
\usepackage[utf8]{inputenc}
\usepackage{booktabs}
\usepackage{amsmath}
\usepackage{natbib}

\title{My Document}
\author{Your Name}
\date{\today}

\begin{document}

\maketitle

\section{Introduction}
This is a sample document demonstrating various \LaTeX{} features including mathematical equations, tables, and formatting. You can cite papers using the citation search feature, and they will be automatically added to your \texttt{cite.bib} file.

\section{Mathematical Expressions}
Here are some example equations:

\subsection{Quadratic Formula}
The quadratic formula is given by:
\[x = \frac{-b \pm \sqrt{b^2 - 4ac}}{2a}\]

\subsection{Matrix Example}
A sample matrix:
\[
\begin{pmatrix}
a_{11} & a_{12} & a_{13} \\
a_{21} & a_{22} & a_{23} \\
a_{31} & a_{32} & a_{33}
\end{pmatrix}
\]

\section{Tables}
Below is a sample table showing data:

\begin{table}[h]
\centering
\begin{tabular}{@{}lrrr@{}}
\toprule
\textbf{Category} & \textbf{Value 1} & \textbf{Value 2} & \textbf{Total} \\
\midrule
Item A & 10.5 & 20.3 & 30.8 \\
Item B & 15.2 & 25.1 & 40.3 \\
Item C & 22.1 & 12.7 & 34.8 \\
\bottomrule
\end{tabular}
\caption{Sample Data Table}
\label{tab:sample}
\end{table}

\section{Lists}
\subsection{Bullet Points}
\begin{itemize}
  \item First item
  \item Second item
  \item Third item with sub-items:
    \begin{itemize}
      \item Sub-item 1
      \item Sub-item 2
    \end{itemize}
\end{itemize}

\subsection{Numbered List}
\begin{enumerate}
  \item First step
  \item Second step
  \item Third step
\end{enumerate}

\section{Citations and References}
When you use the citation search feature, references will be automatically managed. For example, you might cite a paper like this: \cite{example2024}. All citations will be automatically added to your \texttt{cite.bib} file.

\section{Conclusion}
This template demonstrates basic \LaTeX{} features including mathematical formulas, tables, and lists. The bibliography system is now set up to work seamlessly with Bibby AI's citation management. Feel free to modify it for your needs!

\bibliographystyle{plain}
\bibliography{cite}

\end{document}\documentclass{article}
\usepackage[utf8]{inputenc}
\usepackage{booktabs}
\usepackage{amsmath}
\usepackage{natbib}

\title{My Document}
\author{Your Name}
\date{\today}

\begin{document}

\maketitle

\section{Introduction}
This is a sample document demonstrating various \LaTeX{} features including mathematical equations, tables, and formatting. You can cite papers using the citation search feature, and they will be automatically added to your \texttt{cite.bib} file.

\section{Mathematical Expressions}
Here are some example equations:

\subsection{Quadratic Formula}
The quadratic formula is given by:
\[x = \frac{-b \pm \sqrt{b^2 - 4ac}}{2a}\]

\subsection{Matrix Example}
A sample matrix:
\[
\begin{pmatrix}
a_{11} & a_{12} & a_{13} \\
a_{21} & a_{22} & a_{23} \\
a_{31} & a_{32} & a_{33}
\end{pmatrix}
\]

\section{Tables}
Below is a sample table showing data:

\begin{table}[h]
\centering
\begin{tabular}{@{}lrrr@{}}
\toprule
\textbf{Category} & \textbf{Value 1} & \textbf{Value 2} & \textbf{Total} \\
\midrule
Item A & 10.5 & 20.3 & 30.8 \\
Item B & 15.2 & 25.1 & 40.3 \\
Item C & 22.1 & 12.7 & 34.8 \\
\bottomrule
\end{tabular}
\caption{Sample Data Table}
\label{tab:sample}
\end{table}

\section{Lists}
\subsection{Bullet Points}
\begin{itemize}
  \item First item
  \item Second item
  \item Third item with sub-items:
    \begin{itemize}
      \item Sub-item 1
      \item Sub-item 2
    \end{itemize}
\end{itemize}

\subsection{Numbered List}
\begin{enumerate}
  \item First step
  \item Second step
  \item Third step
\end{enumerate}

\section{Citations and References}
When you use the citation search feature, references will be automatically managed. For example, you might cite a paper like this: \cite{example2024}. All citations will be automatically added to your \texttt{cite.bib} file.

\section{Conclusion}
This template demonstrates basic \LaTeX{} features including mathematical formulas, tables, and lists. The bibliography system is now set up to work seamlessly with Bibby AI's citation management. Feel free to modify it for your needs!

\bibliographystyle{plain}
\bibliography{cite}

\end{document}\documentclass{article}
\usepackage[utf8]{inputenc}
\usepackage{booktabs}
\usepackage{amsmath}
\usepackage{natbib}

\title{My Document}
\author{Your Name}
\date{\today}

\begin{document}

\maketitle

\section{Introduction}
This is a sample document demonstrating various \LaTeX{} features including mathematical equations, tables, and formatting. You can cite papers using the citation search feature, and they will be automatically added to your \texttt{cite.bib} file.

\section{Mathematical Expressions}
Here are some example equations:

\subsection{Quadratic Formula}
The quadratic formula is given by:
\[x = \frac{-b \pm \sqrt{b^2 - 4ac}}{2a}\]

\subsection{Matrix Example}
A sample matrix:
\[
\begin{pmatrix}
a_{11} & a_{12} & a_{13} \\
a_{21} & a_{22} & a_{23} \\
a_{31} & a_{32} & a_{33}
\end{pmatrix}
\]

\section{Tables}
Below is a sample table showing data:

\begin{table}[h]
\centering
\begin{tabular}{@{}lrrr@{}}
\toprule
\textbf{Category} & \textbf{Value 1} & \textbf{Value 2} & \textbf{Total} \\
\midrule
Item A & 10.5 & 20.3 & 30.8 \\
Item B & 15.2 & 25.1 & 40.3 \\
Item C & 22.1 & 12.7 & 34.8 \\
\bottomrule
\end{tabular}
\caption{Sample Data Table}
\label{tab:sample}
\end{table}

\section{Lists}
\subsection{Bullet Points}
\begin{itemize}
  \item First item
  \item Second item
  \item Third item with sub-items:
    \begin{itemize}
      \item Sub-item 1
      \item Sub-item 2
    \end{itemize}
\end{itemize}

\subsection{Numbered List}
\begin{enumerate}
  \item First step
  \item Second step
  \item Third step
\end{enumerate}

\section{Citations and References}
When you use the citation search feature, references will be automatically managed. For example, you might cite a paper like this: \cite{example2024}. All citations will be automatically added to your \texttt{cite.bib} file.

\section{Conclusion}
This template demonstrates basic \LaTeX{} features including mathematical formulas, tables, and lists. The bibliography system is now set up to work seamlessly with Bibby AI's citation management. Feel free to modify it for your needs!

\bibliographystyle{plain}
\bibliography{cite}

\end{document}\documentclass{article}
\usepackage[utf8]{inputenc}
\usepackage{booktabs}
\usepackage{amsmath}
\usepackage{natbib}

\title{My Document}
\author{Your Name}
\date{\today}

\begin{document}

\maketitle

\section{Introduction}
This is a sample document demonstrating various \LaTeX{} features including mathematical equations, tables, and formatting. You can cite papers using the citation search feature, and they will be automatically added to your \texttt{cite.bib} file.

\section{Mathematical Expressions}
Here are some example equations:

\subsection{Quadratic Formula}
The quadratic formula is given by:
\[x = \frac{-b \pm \sqrt{b^2 - 4ac}}{2a}\]

\subsection{Matrix Example}
A sample matrix:
\[
\begin{pmatrix}
a_{11} & a_{12} & a_{13} \\
a_{21} & a_{22} & a_{23} \\
a_{31} & a_{32} & a_{33}
\end{pmatrix}
\]

\section{Tables}
Below is a sample table showing data:

\begin{table}[h]
\centering
\begin{tabular}{@{}lrrr@{}}
\toprule
\textbf{Category} & \textbf{Value 1} & \textbf{Value 2} & \textbf{Total} \\
\midrule
Item A & 10.5 & 20.3 & 30.8 \\
Item B & 15.2 & 25.1 & 40.3 \\
Item C & 22.1 & 12.7 & 34.8 \\
\bottomrule
\end{tabular}
\caption{Sample Data Table}
\label{tab:sample}
\end{table}

\section{Lists}
\subsection{Bullet Points}
\begin{itemize}
  \item First item
  \item Second item
  \item Third item with sub-items:
    \begin{itemize}
      \item Sub-item 1
      \item Sub-item 2
    \end{itemize}
\end{itemize}

\subsection{Numbered List}
\begin{enumerate}
  \item First step
  \item Second step
  \item Third step
\end{enumerate}

\section{Citations and References}
When you use the citation search feature, references will be automatically managed. For example, you might cite a paper like this: \cite{example2024}. All citations will be automatically added to your \texttt{cite.bib} file.

\section{Conclusion}
This template demonstrates basic \LaTeX{} features including mathematical formulas, tables, and lists. The bibliography system is now set up to work seamlessly with Bibby AI's citation management. Feel free to modify it for your needs!

\bibliographystyle{plain}
\bibliography{cite}

\end{document}\documentclass{article}
\usepackage[utf8]{inputenc}
\usepackage{booktabs}
\usepackage{amsmath}
\usepackage{natbib}

\title{My Document}
\author{Your Name}
\date{\today}

\begin{document}

\maketitle

\section{Introduction}
This is a sample document demonstrating various \LaTeX{} features including mathematical equations, tables, and formatting. You can cite papers using the citation search feature, and they will be automatically added to your \texttt{cite.bib} file.

\section{Mathematical Expressions}
Here are some example equations:

\subsection{Quadratic Formula}
The quadratic formula is given by:
\[x = \frac{-b \pm \sqrt{b^2 - 4ac}}{2a}\]

\subsection{Matrix Example}
A sample matrix:
\[
\begin{pmatrix}
a_{11} & a_{12} & a_{13} \\
a_{21} & a_{22} & a_{23} \\
a_{31} & a_{32} & a_{33}
\end{pmatrix}
\]

\section{Tables}
Below is a sample table showing data:

\begin{table}[h]
\centering
\begin{tabular}{@{}lrrr@{}}
\toprule
\textbf{Category} & \textbf{Value 1} & \textbf{Value 2} & \textbf{Total} \\
\midrule
Item A & 10.5 & 20.3 & 30.8 \\
Item B & 15.2 & 25.1 & 40.3 \\
Item C & 22.1 & 12.7 & 34.8 \\
\bottomrule
\end{tabular}
\caption{Sample Data Table}
\label{tab:sample}
\end{table}

\section{Lists}
\subsection{Bullet Points}
\begin{itemize}
  \item First item
  \item Second item
  \item Third item with sub-items:
    \begin{itemize}
      \item Sub-item 1
      \item Sub-item 2
    \end{itemize}
\end{itemize}

\subsection{Numbered List}
\begin{enumerate}
  \item First step
  \item Second step
  \item Third step
\end{enumerate}

\section{Citations and References}
When you use the citation search feature, references will be automatically managed. For example, you might cite a paper like this: \cite{example2024}. All citations will be automatically added to your \texttt{cite.bib} file.

\section{Conclusion}
This template demonstrates basic \LaTeX{} features including mathematical formulas, tables, and lists. The bibliography system is now set up to work seamlessly with Bibby AI's citation management. Feel free to modify it for your needs!

\bibliographystyle{plain}
\bibliography{cite}

\end{document}\documentclass{article}
\usepackage[utf8]{inputenc}
\usepackage{booktabs}
\usepackage{amsmath}
\usepackage{natbib}

\title{My Document}
\author{Your Name}
\date{\today}

\begin{document}

\maketitle

\section{Introduction}
This is a sample document demonstrating various \LaTeX{} features including mathematical equations, tables, and formatting. You can cite papers using the citation search feature, and they will be automatically added to your \texttt{cite.bib} file.

\section{Mathematical Expressions}
Here are some example equations:

\subsection{Quadratic Formula}
The quadratic formula is given by:
\[x = \frac{-b \pm \sqrt{b^2 - 4ac}}{2a}\]

\subsection{Matrix Example}
A sample matrix:
\[
\begin{pmatrix}
a_{11} & a_{12} & a_{13} \\
a_{21} & a_{22} & a_{23} \\
a_{31} & a_{32} & a_{33}
\end{pmatrix}
\]

\section{Tables}
Below is a sample table showing data:

\begin{table}[h]
\centering
\begin{tabular}{@{}lrrr@{}}
\toprule
\textbf{Category} & \textbf{Value 1} & \textbf{Value 2} & \textbf{Total} \\
\midrule
Item A & 10.5 & 20.3 & 30.8 \\
Item B & 15.2 & 25.1 & 40.3 \\
Item C & 22.1 & 12.7 & 34.8 \\
\bottomrule
\end{tabular}
\caption{Sample Data Table}
\label{tab:sample}
\end{table}

\section{Lists}
\subsection{Bullet Points}
\begin{itemize}
  \item First item
  \item Second item
  \item Third item with sub-items:
    \begin{itemize}
      \item Sub-item 1
      \item Sub-item 2
    \end{itemize}
\end{itemize}

\subsection{Numbered List}
\begin{enumerate}
  \item First step
  \item Second step
  \item Third step
\end{enumerate}

\section{Citations and References}
When you use the citation search feature, references will be automatically managed. For example, you might cite a paper like this: \cite{example2024}. All citations will be automatically added to your \texttt{cite.bib} file.

\section{Conclusion}
This template demonstrates basic \LaTeX{} features including mathematical formulas, tables, and lists. The bibliography system is now set up to work seamlessly with Bibby AI's citation management. Feel free to modify it for your needs!

\bibliographystyle{plain}
\bibliography{cite}

\end{document}\documentclass{article}
\usepackage[utf8]{inputenc}
\usepackage{booktabs}
\usepackage{amsmath}
\usepackage{natbib}

\title{My Document}
\author{Your Name}
\date{\today}

\begin{document}

\maketitle

\section{Introduction}
This is a sample document demonstrating various \LaTeX{} features including mathematical equations, tables, and formatting. You can cite papers using the citation search feature, and they will be automatically added to your \texttt{cite.bib} file.

\section{Mathematical Expressions}
Here are some example equations:

\subsection{Quadratic Formula}
The quadratic formula is given by:
\[x = \frac{-b \pm \sqrt{b^2 - 4ac}}{2a}\]

\subsection{Matrix Example}
A sample matrix:
\[
\begin{pmatrix}
a_{11} & a_{12} & a_{13} \\
a_{21} & a_{22} & a_{23} \\
a_{31} & a_{32} & a_{33}
\end{pmatrix}
\]

\section{Tables}
Below is a sample table showing data:

\begin{table}[h]
\centering
\begin{tabular}{@{}lrrr@{}}
\toprule
\textbf{Category} & \textbf{Value 1} & \textbf{Value 2} & \textbf{Total} \\
\midrule
Item A & 10.5 & 20.3 & 30.8 \\
Item B & 15.2 & 25.1 & 40.3 \\
Item C & 22.1 & 12.7 & 34.8 \\
\bottomrule
\end{tabular}
\caption{Sample Data Table}
\label{tab:sample}
\end{table}

\section{Lists}
\subsection{Bullet Points}
\begin{itemize}
  \item First item
  \item Second item
  \item Third item with sub-items:
    \begin{itemize}
      \item Sub-item 1
      \item Sub-item 2
    \end{itemize}
\end{itemize}

\subsection{Numbered List}
\begin{enumerate}
  \item First step
  \item Second step
  \item Third step
\end{enumerate}

\section{Citations and References}
When you use the citation search feature, references will be automatically managed. For example, you might cite a paper like this: \cite{example2024}. All citations will be automatically added to your \texttt{cite.bib} file.

\section{Conclusion}
This template demonstrates basic \LaTeX{} features including mathematical formulas, tables, and lists. The bibliography system is now set up to work seamlessly with Bibby AI's citation management. Feel free to modify it for your needs!

\bibliographystyle{plain}
\bibliography{cite}

\end{document}\documentclass{article}
\usepackage[utf8]{inputenc}
\usepackage{booktabs}
\usepackage{amsmath}
\usepackage{natbib}

\title{My Document}
\author{Your Name}
\date{\today}

\begin{document}

\maketitle

\section{Introduction}
This is a sample document demonstrating various \LaTeX{} features including mathematical equations, tables, and formatting. You can cite papers using the citation search feature, and they will be automatically added to your \texttt{cite.bib} file.

\section{Mathematical Expressions}
Here are some example equations:

\subsection{Quadratic Formula}
The quadratic formula is given by:
\[x = \frac{-b \pm \sqrt{b^2 - 4ac}}{2a}\]

\subsection{Matrix Example}
A sample matrix:
\[
\begin{pmatrix}
a_{11} & a_{12} & a_{13} \\
a_{21} & a_{22} & a_{23} \\
a_{31} & a_{32} & a_{33}
\end{pmatrix}
\]

\section{Tables}
Below is a sample table showing data:

\begin{table}[h]
\centering
\begin{tabular}{@{}lrrr@{}}
\toprule
\textbf{Category} & \textbf{Value 1} & \textbf{Value 2} & \textbf{Total} \\
\midrule
Item A & 10.5 & 20.3 & 30.8 \\
Item B & 15.2 & 25.1 & 40.3 \\
Item C & 22.1 & 12.7 & 34.8 \\
\bottomrule
\end{tabular}
\caption{Sample Data Table}
\label{tab:sample}
\end{table}

\section{Lists}
\subsection{Bullet Points}
\begin{itemize}
  \item First item
  \item Second item
  \item Third item with sub-items:
    \begin{itemize}
      \item Sub-item 1
      \item Sub-item 2
    \end{itemize}
\end{itemize}

\subsection{Numbered List}
\begin{enumerate}
  \item First step
  \item Second step
  \item Third step
\end{enumerate}

\section{Citations and References}
When you use the citation search feature, references will be automatically managed. For example, you might cite a paper like this: \cite{example2024}. All citations will be automatically added to your \texttt{cite.bib} file.

\section{Conclusion}
This template demonstrates basic \LaTeX{} features including mathematical formulas, tables, and lists. The bibliography system is now set up to work seamlessly with Bibby AI's citation management. Feel free to modify it for your needs!

\bibliographystyle{plain}
\bibliography{cite}

\end{document}\documentclass{article}
\usepackage[utf8]{inputenc}
\usepackage{booktabs}
\usepackage{amsmath}
\usepackage{natbib}

\title{My Document}
\author{Your Name}
\date{\today}

\begin{document}

\maketitle

\section{Introduction}
This is a sample document demonstrating various \LaTeX{} features including mathematical equations, tables, and formatting. You can cite papers using the citation search feature, and they will be automatically added to your \texttt{cite.bib} file.

\section{Mathematical Expressions}
Here are some example equations:

\subsection{Quadratic Formula}
The quadratic formula is given by:
\[x = \frac{-b \pm \sqrt{b^2 - 4ac}}{2a}\]

\subsection{Matrix Example}
A sample matrix:
\[
\begin{pmatrix}
a_{11} & a_{12} & a_{13} \\
a_{21} & a_{22} & a_{23} \\
a_{31} & a_{32} & a_{33}
\end{pmatrix}
\]

\section{Tables}
Below is a sample table showing data:

\begin{table}[h]
\centering
\begin{tabular}{@{}lrrr@{}}
\toprule
\textbf{Category} & \textbf{Value 1} & \textbf{Value 2} & \textbf{Total} \\
\midrule
Item A & 10.5 & 20.3 & 30.8 \\
Item B & 15.2 & 25.1 & 40.3 \\
Item C & 22.1 & 12.7 & 34.8 \\
\bottomrule
\end{tabular}
\caption{Sample Data Table}
\label{tab:sample}
\end{table}

\section{Lists}
\subsection{Bullet Points}
\begin{itemize}
  \item First item
  \item Second item
  \item Third item with sub-items:
    \begin{itemize}
      \item Sub-item 1
      \item Sub-item 2
    \end{itemize}
\end{itemize}

\subsection{Numbered List}
\begin{enumerate}
  \item First step
  \item Second step
  \item Third step
\end{enumerate}

\section{Citations and References}
When you use the citation search feature, references will be automatically managed. For example, you might cite a paper like this: \cite{example2024}. All citations will be automatically added to your \texttt{cite.bib} file.

\section{Conclusion}
This template demonstrates basic \LaTeX{} features including mathematical formulas, tables, and lists. The bibliography system is now set up to work seamlessly with Bibby AI's citation management. Feel free to modify it for your needs!

\bibliographystyle{plain}
\bibliography{cite}

\end{document}\documentclass{article}
\usepackage[utf8]{inputenc}
\usepackage{booktabs}
\usepackage{amsmath}
\usepackage{natbib}

\title{My Document}
\author{Your Name}
\date{\today}

\begin{document}

\maketitle

\section{Introduction}
This is a sample document demonstrating various \LaTeX{} features including mathematical equations, tables, and formatting. You can cite papers using the citation search feature, and they will be automatically added to your \texttt{cite.bib} file.

\section{Mathematical Expressions}
Here are some example equations:

\subsection{Quadratic Formula}
The quadratic formula is given by:
\[x = \frac{-b \pm \sqrt{b^2 - 4ac}}{2a}\]

\subsection{Matrix Example}
A sample matrix:
\[
\begin{pmatrix}
a_{11} & a_{12} & a_{13} \\
a_{21} & a_{22} & a_{23} \\
a_{31} & a_{32} & a_{33}
\end{pmatrix}
\]

\section{Tables}
Below is a sample table showing data:

\begin{table}[h]
\centering
\begin{tabular}{@{}lrrr@{}}
\toprule
\textbf{Category} & \textbf{Value 1} & \textbf{Value 2} & \textbf{Total} \\
\midrule
Item A & 10.5 & 20.3 & 30.8 \\
Item B & 15.2 & 25.1 & 40.3 \\
Item C & 22.1 & 12.7 & 34.8 \\
\bottomrule
\end{tabular}
\caption{Sample Data Table}
\label{tab:sample}
\end{table}

\section{Lists}
\subsection{Bullet Points}
\begin{itemize}
  \item First item
  \item Second item
  \item Third item with sub-items:
    \begin{itemize}
      \item Sub-item 1
      \item Sub-item 2
    \end{itemize}
\end{itemize}

\subsection{Numbered List}
\begin{enumerate}
  \item First step
  \item Second step
  \item Third step
\end{enumerate}

\section{Citations and References}
When you use the citation search feature, references will be automatically managed. For example, you might cite a paper like this: \cite{example2024}. All citations will be automatically added to your \texttt{cite.bib} file.

\section{Conclusion}
This template demonstrates basic \LaTeX{} features including mathematical formulas, tables, and lists. The bibliography system is now set up to work seamlessly with Bibby AI's citation management. Feel free to modify it for your needs!

\bibliographystyle{plain}
\bibliography{cite}

\end{document}\documentclass{article}
\usepackage[utf8]{inputenc}
\usepackage{booktabs}
\usepackage{amsmath}
\usepackage{natbib}

\title{My Document}
\author{Your Name}
\date{\today}

\begin{document}

\maketitle

\section{Introduction}
This is a sample document demonstrating various \LaTeX{} features including mathematical equations, tables, and formatting. You can cite papers using the citation search feature, and they will be automatically added to your \texttt{cite.bib} file.

\section{Mathematical Expressions}
Here are some example equations:

\subsection{Quadratic Formula}
The quadratic formula is given by:
\[x = \frac{-b \pm \sqrt{b^2 - 4ac}}{2a}\]

\subsection{Matrix Example}
A sample matrix:
\[
\begin{pmatrix}
a_{11} & a_{12} & a_{13} \\
a_{21} & a_{22} & a_{23} \\
a_{31} & a_{32} & a_{33}
\end{pmatrix}
\]

\section{Tables}
Below is a sample table showing data:

\begin{table}[h]
\centering
\begin{tabular}{@{}lrrr@{}}
\toprule
\textbf{Category} & \textbf{Value 1} & \textbf{Value 2} & \textbf{Total} \\
\midrule
Item A & 10.5 & 20.3 & 30.8 \\
Item B & 15.2 & 25.1 & 40.3 \\
Item C & 22.1 & 12.7 & 34.8 \\
\bottomrule
\end{tabular}
\caption{Sample Data Table}
\label{tab:sample}
\end{table}

\section{Lists}
\subsection{Bullet Points}
\begin{itemize}
  \item First item
  \item Second item
  \item Third item with sub-items:
    \begin{itemize}
      \item Sub-item 1
      \item Sub-item 2
    \end{itemize}
\end{itemize}

\subsection{Numbered List}
\begin{enumerate}
  \item First step
  \item Second step
  \item Third step
\end{enumerate}

\section{Citations and References}
When you use the citation search feature, references will be automatically managed. For example, you might cite a paper like this: \cite{example2024}. All citations will be automatically added to your \texttt{cite.bib} file.

\section{Conclusion}
This template demonstrates basic \LaTeX{} features including mathematical formulas, tables, and lists. The bibliography system is now set up to work seamlessly with Bibby AI's citation management. Feel free to modify it for your needs!

\bibliographystyle{plain}
\bibliography{cite}

\end{document}\documentclass{article}
\usepackage[utf8]{inputenc}
\usepackage{booktabs}
\usepackage{amsmath}
\usepackage{natbib}

\title{My Document}
\author{Your Name}
\date{\today}

\begin{document}

\maketitle

\section{Introduction}
This is a sample document demonstrating various \LaTeX{} features including mathematical equations, tables, and formatting. You can cite papers using the citation search feature, and they will be automatically added to your \texttt{cite.bib} file.

\section{Mathematical Expressions}
Here are some example equations:

\subsection{Quadratic Formula}
The quadratic formula is given by:
\[x = \frac{-b \pm \sqrt{b^2 - 4ac}}{2a}\]

\subsection{Matrix Example}
A sample matrix:
\[
\begin{pmatrix}
a_{11} & a_{12} & a_{13} \\
a_{21} & a_{22} & a_{23} \\
a_{31} & a_{32} & a_{33}
\end{pmatrix}
\]

\section{Tables}
Below is a sample table showing data:

\begin{table}[h]
\centering
\begin{tabular}{@{}lrrr@{}}
\toprule
\textbf{Category} & \textbf{Value 1} & \textbf{Value 2} & \textbf{Total} \\
\midrule
Item A & 10.5 & 20.3 & 30.8 \\
Item B & 15.2 & 25.1 & 40.3 \\
Item C & 22.1 & 12.7 & 34.8 \\
\bottomrule
\end{tabular}
\caption{Sample Data Table}
\label{tab:sample}
\end{table}

\section{Lists}
\subsection{Bullet Points}
\begin{itemize}
  \item First item
  \item Second item
  \item Third item with sub-items:
    \begin{itemize}
      \item Sub-item 1
      \item Sub-item 2
    \end{itemize}
\end{itemize}

\subsection{Numbered List}
\begin{enumerate}
  \item First step
  \item Second step
  \item Third step
\end{enumerate}

\section{Citations and References}
When you use the citation search feature, references will be automatically managed. For example, you might cite a paper like this: \cite{example2024}. All citations will be automatically added to your \texttt{cite.bib} file.

\section{Conclusion}
This template demonstrates basic \LaTeX{} features including mathematical formulas, tables, and lists. The bibliography system is now set up to work seamlessly with Bibby AI's citation management. Feel free to modify it for your needs!

\bibliographystyle{plain}
\bibliography{cite}

\end{document}\documentclass{article}
\usepackage[utf8]{inputenc}
\usepackage{booktabs}
\usepackage{amsmath}
\usepackage{natbib}

\title{My Document}
\author{Your Name}
\date{\today}

\begin{document}

\maketitle

\section{Introduction}
This is a sample document demonstrating various \LaTeX{} features including mathematical equations, tables, and formatting. You can cite papers using the citation search feature, and they will be automatically added to your \texttt{cite.bib} file.

\section{Mathematical Expressions}
Here are some example equations:

\subsection{Quadratic Formula}
The quadratic formula is given by:
\[x = \frac{-b \pm \sqrt{b^2 - 4ac}}{2a}\]

\subsection{Matrix Example}
A sample matrix:
\[
\begin{pmatrix}
a_{11} & a_{12} & a_{13} \\
a_{21} & a_{22} & a_{23} \\
a_{31} & a_{32} & a_{33}
\end{pmatrix}
\]

\section{Tables}
Below is a sample table showing data:

\begin{table}[h]
\centering
\begin{tabular}{@{}lrrr@{}}
\toprule
\textbf{Category} & \textbf{Value 1} & \textbf{Value 2} & \textbf{Total} \\
\midrule
Item A & 10.5 & 20.3 & 30.8 \\
Item B & 15.2 & 25.1 & 40.3 \\
Item C & 22.1 & 12.7 & 34.8 \\
\bottomrule
\end{tabular}
\caption{Sample Data Table}
\label{tab:sample}
\end{table}

\section{Lists}
\subsection{Bullet Points}
\begin{itemize}
  \item First item
  \item Second item
  \item Third item with sub-items:
    \begin{itemize}
      \item Sub-item 1
      \item Sub-item 2
    \end{itemize}
\end{itemize}

\subsection{Numbered List}
\begin{enumerate}
  \item First step
  \item Second step
  \item Third step
\end{enumerate}

\section{Citations and References}
When you use the citation search feature, references will be automatically managed. For example, you might cite a paper like this: \cite{example2024}. All citations will be automatically added to your \texttt{cite.bib} file.

\section{Conclusion}
This template demonstrates basic \LaTeX{} features including mathematical formulas, tables, and lists. The bibliography system is now set up to work seamlessly with Bibby AI's citation management. Feel free to modify it for your needs!

\bibliographystyle{plain}
\bibliography{cite}

\end{document}\documentclass{article}
\usepackage[utf8]{inputenc}
\usepackage{booktabs}
\usepackage{amsmath}
\usepackage{natbib}

\title{My Document}
\author{Your Name}
\date{\today}

\begin{document}

\maketitle

\section{Introduction}
This is a sample document demonstrating various \LaTeX{} features including mathematical equations, tables, and formatting. You can cite papers using the citation search feature, and they will be automatically added to your \texttt{cite.bib} file.

\section{Mathematical Expressions}
Here are some example equations:

\subsection{Quadratic Formula}
The quadratic formula is given by:
\[x = \frac{-b \pm \sqrt{b^2 - 4ac}}{2a}\]

\subsection{Matrix Example}
A sample matrix:
\[
\begin{pmatrix}
a_{11} & a_{12} & a_{13} \\
a_{21} & a_{22} & a_{23} \\
a_{31} & a_{32} & a_{33}
\end{pmatrix}
\]

\section{Tables}
Below is a sample table showing data:

\begin{table}[h]
\centering
\begin{tabular}{@{}lrrr@{}}
\toprule
\textbf{Category} & \textbf{Value 1} & \textbf{Value 2} & \textbf{Total} \\
\midrule
Item A & 10.5 & 20.3 & 30.8 \\
Item B & 15.2 & 25.1 & 40.3 \\
Item C & 22.1 & 12.7 & 34.8 \\
\bottomrule
\end{tabular}
\caption{Sample Data Table}
\label{tab:sample}
\end{table}

\section{Lists}
\subsection{Bullet Points}
\begin{itemize}
  \item First item
  \item Second item
  \item Third item with sub-items:
    \begin{itemize}
      \item Sub-item 1
      \item Sub-item 2
    \end{itemize}
\end{itemize}

\subsection{Numbered List}
\begin{enumerate}
  \item First step
  \item Second step
  \item Third step
\end{enumerate}

\section{Citations and References}
When you use the citation search feature, references will be automatically managed. For example, you might cite a paper like this: \cite{example2024}. All citations will be automatically added to your \texttt{cite.bib} file.

\section{Conclusion}
This template demonstrates basic \LaTeX{} features including mathematical formulas, tables, and lists. The bibliography system is now set up to work seamlessly with Bibby AI's citation management. Feel free to modify it for your needs!

\bibliographystyle{plain}
\bibliography{cite}

\end{document}\documentclass{article}
\usepackage[utf8]{inputenc}
\usepackage{booktabs}
\usepackage{amsmath}
\usepackage{natbib}

\title{My Document}
\author{Your Name}
\date{\today}

\begin{document}

\maketitle

\section{Introduction}
This is a sample document demonstrating various \LaTeX{} features including mathematical equations, tables, and formatting. You can cite papers using the citation search feature, and they will be automatically added to your \texttt{cite.bib} file.

\section{Mathematical Expressions}
Here are some example equations:

\subsection{Quadratic Formula}
The quadratic formula is given by:
\[x = \frac{-b \pm \sqrt{b^2 - 4ac}}{2a}\]

\subsection{Matrix Example}
A sample matrix:
\[
\begin{pmatrix}
a_{11} & a_{12} & a_{13} \\
a_{21} & a_{22} & a_{23} \\
a_{31} & a_{32} & a_{33}
\end{pmatrix}
\]

\section{Tables}
Below is a sample table showing data:

\begin{table}[h]
\centering
\begin{tabular}{@{}lrrr@{}}
\toprule
\textbf{Category} & \textbf{Value 1} & \textbf{Value 2} & \textbf{Total} \\
\midrule
Item A & 10.5 & 20.3 & 30.8 \\
Item B & 15.2 & 25.1 & 40.3 \\
Item C & 22.1 & 12.7 & 34.8 \\
\bottomrule
\end{tabular}
\caption{Sample Data Table}
\label{tab:sample}
\end{table}

\section{Lists}
\subsection{Bullet Points}
\begin{itemize}
  \item First item
  \item Second item
  \item Third item with sub-items:
    \begin{itemize}
      \item Sub-item 1
      \item Sub-item 2
    \end{itemize}
\end{itemize}

\subsection{Numbered List}
\begin{enumerate}
  \item First step
  \item Second step
  \item Third step
\end{enumerate}

\section{Citations and References}
When you use the citation search feature, references will be automatically managed. For example, you might cite a paper like this: \cite{example2024}. All citations will be automatically added to your \texttt{cite.bib} file.

\section{Conclusion}
This template demonstrates basic \LaTeX{} features including mathematical formulas, tables, and lists. The bibliography system is now set up to work seamlessly with Bibby AI's citation management. Feel free to modify it for your needs!

\bibliographystyle{plain}
\bibliography{cite}

\end{document}\documentclass{article}
\usepackage[utf8]{inputenc}
\usepackage{booktabs}
\usepackage{amsmath}
\usepackage{natbib}

\title{My Document}
\author{Your Name}
\date{\today}

\begin{document}

\maketitle

\section{Introduction}
This is a sample document demonstrating various \LaTeX{} features including mathematical equations, tables, and formatting. You can cite papers using the citation search feature, and they will be automatically added to your \texttt{cite.bib} file.

\section{Mathematical Expressions}
Here are some example equations:

\subsection{Quadratic Formula}
The quadratic formula is given by:
\[x = \frac{-b \pm \sqrt{b^2 - 4ac}}{2a}\]

\subsection{Matrix Example}
A sample matrix:
\[
\begin{pmatrix}
a_{11} & a_{12} & a_{13} \\
a_{21} & a_{22} & a_{23} \\
a_{31} & a_{32} & a_{33}
\end{pmatrix}
\]

\section{Tables}
Below is a sample table showing data:

\begin{table}[h]
\centering
\begin{tabular}{@{}lrrr@{}}
\toprule
\textbf{Category} & \textbf{Value 1} & \textbf{Value 2} & \textbf{Total} \\
\midrule
Item A & 10.5 & 20.3 & 30.8 \\
Item B & 15.2 & 25.1 & 40.3 \\
Item C & 22.1 & 12.7 & 34.8 \\
\bottomrule
\end{tabular}
\caption{Sample Data Table}
\label{tab:sample}
\end{table}

\section{Lists}
\subsection{Bullet Points}
\begin{itemize}
  \item First item
  \item Second item
  \item Third item with sub-items:
    \begin{itemize}
      \item Sub-item 1
      \item Sub-item 2
    \end{itemize}
\end{itemize}

\subsection{Numbered List}
\begin{enumerate}
  \item First step
  \item Second step
  \item Third step
\end{enumerate}

\section{Citations and References}
When you use the citation search feature, references will be automatically managed. For example, you might cite a paper like this: \cite{example2024}. All citations will be automatically added to your \texttt{cite.bib} file.

\section{Conclusion}
This template demonstrates basic \LaTeX{} features including mathematical formulas, tables, and lists. The bibliography system is now set up to work seamlessly with Bibby AI's citation management. Feel free to modify it for your needs!

\bibliographystyle{plain}
\bibliography{cite}

\end{document}\documentclass{article}
\usepackage[utf8]{inputenc}
\usepackage{booktabs}
\usepackage{amsmath}
\usepackage{natbib}

\title{My Document}
\author{Your Name}
\date{\today}

\begin{document}

\maketitle

\section{Introduction}
This is a sample document demonstrating various \LaTeX{} features including mathematical equations, tables, and formatting. You can cite papers using the citation search feature, and they will be automatically added to your \texttt{cite.bib} file.

\section{Mathematical Expressions}
Here are some example equations:

\subsection{Quadratic Formula}
The quadratic formula is given by:
\[x = \frac{-b \pm \sqrt{b^2 - 4ac}}{2a}\]

\subsection{Matrix Example}
A sample matrix:
\[
\begin{pmatrix}
a_{11} & a_{12} & a_{13} \\
a_{21} & a_{22} & a_{23} \\
a_{31} & a_{32} & a_{33}
\end{pmatrix}
\]

\section{Tables}
Below is a sample table showing data:

\begin{table}[h]
\centering
\begin{tabular}{@{}lrrr@{}}
\toprule
\textbf{Category} & \textbf{Value 1} & \textbf{Value 2} & \textbf{Total} \\
\midrule
Item A & 10.5 & 20.3 & 30.8 \\
Item B & 15.2 & 25.1 & 40.3 \\
Item C & 22.1 & 12.7 & 34.8 \\
\bottomrule
\end{tabular}
\caption{Sample Data Table}
\label{tab:sample}
\end{table}

\section{Lists}
\subsection{Bullet Points}
\begin{itemize}
  \item First item
  \item Second item
  \item Third item with sub-items:
    \begin{itemize}
      \item Sub-item 1
      \item Sub-item 2
    \end{itemize}
\end{itemize}

\subsection{Numbered List}
\begin{enumerate}
  \item First step
  \item Second step
  \item Third step
\end{enumerate}

\section{Citations and References}
When you use the citation search feature, references will be automatically managed. For example, you might cite a paper like this: \cite{example2024}. All citations will be automatically added to your \texttt{cite.bib} file.

\section{Conclusion}
This template demonstrates basic \LaTeX{} features including mathematical formulas, tables, and lists. The bibliography system is now set up to work seamlessly with Bibby AI's citation management. Feel free to modify it for your needs!

\bibliographystyle{plain}
\bibliography{cite}

\end{document}\documentclass{article}
\usepackage[utf8]{inputenc}
\usepackage{booktabs}
\usepackage{amsmath}
\usepackage{natbib}

\title{My Document}
\author{Your Name}
\date{\today}

\begin{document}

\maketitle

\section{Introduction}
This is a sample document demonstrating various \LaTeX{} features including mathematical equations, tables, and formatting. You can cite papers using the citation search feature, and they will be automatically added to your \texttt{cite.bib} file.

\section{Mathematical Expressions}
Here are some example equations:

\subsection{Quadratic Formula}
The quadratic formula is given by:
\[x = \frac{-b \pm \sqrt{b^2 - 4ac}}{2a}\]

\subsection{Matrix Example}
A sample matrix:
\[
\begin{pmatrix}
a_{11} & a_{12} & a_{13} \\
a_{21} & a_{22} & a_{23} \\
a_{31} & a_{32} & a_{33}
\end{pmatrix}
\]

\section{Tables}
Below is a sample table showing data:

\begin{table}[h]
\centering
\begin{tabular}{@{}lrrr@{}}
\toprule
\textbf{Category} & \textbf{Value 1} & \textbf{Value 2} & \textbf{Total} \\
\midrule
Item A & 10.5 & 20.3 & 30.8 \\
Item B & 15.2 & 25.1 & 40.3 \\
Item C & 22.1 & 12.7 & 34.8 \\
\bottomrule
\end{tabular}
\caption{Sample Data Table}
\label{tab:sample}
\end{table}

\section{Lists}
\subsection{Bullet Points}
\begin{itemize}
  \item First item
  \item Second item
  \item Third item with sub-items:
    \begin{itemize}
      \item Sub-item 1
      \item Sub-item 2
    \end{itemize}
\end{itemize}

\subsection{Numbered List}
\begin{enumerate}
  \item First step
  \item Second step
  \item Third step
\end{enumerate}

\section{Citations and References}
When you use the citation search feature, references will be automatically managed. For example, you might cite a paper like this: \cite{example2024}. All citations will be automatically added to your \texttt{cite.bib} file.

\section{Conclusion}
This template demonstrates basic \LaTeX{} features including mathematical formulas, tables, and lists. The bibliography system is now set up to work seamlessly with Bibby AI's citation management. Feel free to modify it for your needs!

\bibliographystyle{plain}
\bibliography{cite}

\end{document}\documentclass{article}
\usepackage[utf8]{inputenc}
\usepackage{booktabs}
\usepackage{amsmath}
\usepackage{natbib}

\title{My Document}
\author{Your Name}
\date{\today}

\begin{document}

\maketitle

\section{Introduction}
This is a sample document demonstrating various \LaTeX{} features including mathematical equations, tables, and formatting. You can cite papers using the citation search feature, and they will be automatically added to your \texttt{cite.bib} file.

\section{Mathematical Expressions}
Here are some example equations:

\subsection{Quadratic Formula}
The quadratic formula is given by:
\[x = \frac{-b \pm \sqrt{b^2 - 4ac}}{2a}\]

\subsection{Matrix Example}
A sample matrix:
\[
\begin{pmatrix}
a_{11} & a_{12} & a_{13} \\
a_{21} & a_{22} & a_{23} \\
a_{31} & a_{32} & a_{33}
\end{pmatrix}
\]

\section{Tables}
Below is a sample table showing data:

\begin{table}[h]
\centering
\begin{tabular}{@{}lrrr@{}}
\toprule
\textbf{Category} & \textbf{Value 1} & \textbf{Value 2} & \textbf{Total} \\
\midrule
Item A & 10.5 & 20.3 & 30.8 \\
Item B & 15.2 & 25.1 & 40.3 \\
Item C & 22.1 & 12.7 & 34.8 \\
\bottomrule
\end{tabular}
\caption{Sample Data Table}
\label{tab:sample}
\end{table}

\section{Lists}
\subsection{Bullet Points}
\begin{itemize}
  \item First item
  \item Second item
  \item Third item with sub-items:
    \begin{itemize}
      \item Sub-item 1
      \item Sub-item 2
    \end{itemize}
\end{itemize}

\subsection{Numbered List}
\begin{enumerate}
  \item First step
  \item Second step
  \item Third step
\end{enumerate}

\section{Citations and References}
When you use the citation search feature, references will be automatically managed. For example, you might cite a paper like this: \cite{example2024}. All citations will be automatically added to your \texttt{cite.bib} file.

\section{Conclusion}
This template demonstrates basic \LaTeX{} features including mathematical formulas, tables, and lists. The bibliography system is now set up to work seamlessly with Bibby AI's citation management. Feel free to modify it for your needs!

\bibliographystyle{plain}
\bibliography{cite}

\end{document}\documentclass{article}
\usepackage[utf8]{inputenc}
\usepackage{booktabs}
\usepackage{amsmath}
\usepackage{natbib}

\title{My Document}
\author{Your Name}
\date{\today}

\begin{document}

\maketitle

\section{Introduction}
This is a sample document demonstrating various \LaTeX{} features including mathematical equations, tables, and formatting. You can cite papers using the citation search feature, and they will be automatically added to your \texttt{cite.bib} file.

\section{Mathematical Expressions}
Here are some example equations:

\subsection{Quadratic Formula}
The quadratic formula is given by:
\[x = \frac{-b \pm \sqrt{b^2 - 4ac}}{2a}\]

\subsection{Matrix Example}
A sample matrix:
\[
\begin{pmatrix}
a_{11} & a_{12} & a_{13} \\
a_{21} & a_{22} & a_{23} \\
a_{31} & a_{32} & a_{33}
\end{pmatrix}
\]

\section{Tables}
Below is a sample table showing data:

\begin{table}[h]
\centering
\begin{tabular}{@{}lrrr@{}}
\toprule
\textbf{Category} & \textbf{Value 1} & \textbf{Value 2} & \textbf{Total} \\
\midrule
Item A & 10.5 & 20.3 & 30.8 \\
Item B & 15.2 & 25.1 & 40.3 \\
Item C & 22.1 & 12.7 & 34.8 \\
\bottomrule
\end{tabular}
\caption{Sample Data Table}
\label{tab:sample}
\end{table}

\section{Lists}
\subsection{Bullet Points}
\begin{itemize}
  \item First item
  \item Second item
  \item Third item with sub-items:
    \begin{itemize}
      \item Sub-item 1
      \item Sub-item 2
    \end{itemize}
\end{itemize}

\subsection{Numbered List}
\begin{enumerate}
  \item First step
  \item Second step
  \item Third step
\end{enumerate}

\section{Citations and References}
When you use the citation search feature, references will be automatically managed. For example, you might cite a paper like this: \cite{example2024}. All citations will be automatically added to your \texttt{cite.bib} file.

\section{Conclusion}
This template demonstrates basic \LaTeX{} features including mathematical formulas, tables, and lists. The bibliography system is now set up to work seamlessly with Bibby AI's citation management. Feel free to modify it for your needs!

\bibliographystyle{plain}
\bibliography{cite}

\end{document}\documentclass{article}
\usepackage[utf8]{inputenc}
\usepackage{booktabs}
\usepackage{amsmath}
\usepackage{natbib}

\title{My Document}
\author{Your Name}
\date{\today}

\begin{document}

\maketitle

\section{Introduction}
This is a sample document demonstrating various \LaTeX{} features including mathematical equations, tables, and formatting. You can cite papers using the citation search feature, and they will be automatically added to your \texttt{cite.bib} file.

\section{Mathematical Expressions}
Here are some example equations:

\subsection{Quadratic Formula}
The quadratic formula is given by:
\[x = \frac{-b \pm \sqrt{b^2 - 4ac}}{2a}\]

\subsection{Matrix Example}
A sample matrix:
\[
\begin{pmatrix}
a_{11} & a_{12} & a_{13} \\
a_{21} & a_{22} & a_{23} \\
a_{31} & a_{32} & a_{33}
\end{pmatrix}
\]

\section{Tables}
Below is a sample table showing data:

\begin{table}[h]
\centering
\begin{tabular}{@{}lrrr@{}}
\toprule
\textbf{Category} & \textbf{Value 1} & \textbf{Value 2} & \textbf{Total} \\
\midrule
Item A & 10.5 & 20.3 & 30.8 \\
Item B & 15.2 & 25.1 & 40.3 \\
Item C & 22.1 & 12.7 & 34.8 \\
\bottomrule
\end{tabular}
\caption{Sample Data Table}
\label{tab:sample}
\end{table}

\section{Lists}
\subsection{Bullet Points}
\begin{itemize}
  \item First item
  \item Second item
  \item Third item with sub-items:
    \begin{itemize}
      \item Sub-item 1
      \item Sub-item 2
    \end{itemize}
\end{itemize}

\subsection{Numbered List}
\begin{enumerate}
  \item First step
  \item Second step
  \item Third step
\end{enumerate}

\section{Citations and References}
When you use the citation search feature, references will be automatically managed. For example, you might cite a paper like this: \cite{example2024}. All citations will be automatically added to your \texttt{cite.bib} file.

\section{Conclusion}
This template demonstrates basic \LaTeX{} features including mathematical formulas, tables, and lists. The bibliography system is now set up to work seamlessly with Bibby AI's citation management. Feel free to modify it for your needs!

\bibliographystyle{plain}
\bibliography{cite}

\end{document}\documentclass{article}
\usepackage[utf8]{inputenc}
\usepackage{booktabs}
\usepackage{amsmath}
\usepackage{natbib}

\title{My Document}
\author{Your Name}
\date{\today}

\begin{document}

\maketitle

\section{Introduction}
This is a sample document demonstrating various \LaTeX{} features including mathematical equations, tables, and formatting. You can cite papers using the citation search feature, and they will be automatically added to your \texttt{cite.bib} file.

\section{Mathematical Expressions}
Here are some example equations:

\subsection{Quadratic Formula}
The quadratic formula is given by:
\[x = \frac{-b \pm \sqrt{b^2 - 4ac}}{2a}\]

\subsection{Matrix Example}
A sample matrix:
\[
\begin{pmatrix}
a_{11} & a_{12} & a_{13} \\
a_{21} & a_{22} & a_{23} \\
a_{31} & a_{32} & a_{33}
\end{pmatrix}
\]

\section{Tables}
Below is a sample table showing data:

\begin{table}[h]
\centering
\begin{tabular}{@{}lrrr@{}}
\toprule
\textbf{Category} & \textbf{Value 1} & \textbf{Value 2} & \textbf{Total} \\
\midrule
Item A & 10.5 & 20.3 & 30.8 \\
Item B & 15.2 & 25.1 & 40.3 \\
Item C & 22.1 & 12.7 & 34.8 \\
\bottomrule
\end{tabular}
\caption{Sample Data Table}
\label{tab:sample}
\end{table}

\section{Lists}
\subsection{Bullet Points}
\begin{itemize}
  \item First item
  \item Second item
  \item Third item with sub-items:
    \begin{itemize}
      \item Sub-item 1
      \item Sub-item 2
    \end{itemize}
\end{itemize}

\subsection{Numbered List}
\begin{enumerate}
  \item First step
  \item Second step
  \item Third step
\end{enumerate}

\section{Citations and References}
When you use the citation search feature, references will be automatically managed. For example, you might cite a paper like this: \cite{example2024}. All citations will be automatically added to your \texttt{cite.bib} file.

\section{Conclusion}
This template demonstrates basic \LaTeX{} features including mathematical formulas, tables, and lists. The bibliography system is now set up to work seamlessly with Bibby AI's citation management. Feel free to modify it for your needs!

\bibliographystyle{plain}
\bibliography{cite}

\end{document}\documentclass{article}
\usepackage[utf8]{inputenc}
\usepackage{booktabs}
\usepackage{amsmath}
\usepackage{natbib}

\title{My Document}
\author{Your Name}
\date{\today}

\begin{document}

\maketitle

\section{Introduction}
This is a sample document demonstrating various \LaTeX{} features including mathematical equations, tables, and formatting. You can cite papers using the citation search feature, and they will be automatically added to your \texttt{cite.bib} file.

\section{Mathematical Expressions}
Here are some example equations:

\subsection{Quadratic Formula}
The quadratic formula is given by:
\[x = \frac{-b \pm \sqrt{b^2 - 4ac}}{2a}\]

\subsection{Matrix Example}
A sample matrix:
\[
\begin{pmatrix}
a_{11} & a_{12} & a_{13} \\
a_{21} & a_{22} & a_{23} \\
a_{31} & a_{32} & a_{33}
\end{pmatrix}
\]

\section{Tables}
Below is a sample table showing data:

\begin{table}[h]
\centering
\begin{tabular}{@{}lrrr@{}}
\toprule
\textbf{Category} & \textbf{Value 1} & \textbf{Value 2} & \textbf{Total} \\
\midrule
Item A & 10.5 & 20.3 & 30.8 \\
Item B & 15.2 & 25.1 & 40.3 \\
Item C & 22.1 & 12.7 & 34.8 \\
\bottomrule
\end{tabular}
\caption{Sample Data Table}
\label{tab:sample}
\end{table}

\section{Lists}
\subsection{Bullet Points}
\begin{itemize}
  \item First item
  \item Second item
  \item Third item with sub-items:
    \begin{itemize}
      \item Sub-item 1
      \item Sub-item 2
    \end{itemize}
\end{itemize}

\subsection{Numbered List}
\begin{enumerate}
  \item First step
  \item Second step
  \item Third step
\end{enumerate}

\section{Citations and References}
When you use the citation search feature, references will be automatically managed. For example, you might cite a paper like this: \cite{example2024}. All citations will be automatically added to your \texttt{cite.bib} file.

\section{Conclusion}
This template demonstrates basic \LaTeX{} features including mathematical formulas, tables, and lists. The bibliography system is now set up to work seamlessly with Bibby AI's citation management. Feel free to modify it for your needs!

\bibliographystyle{plain}
\bibliography{cite}

\end{document}\documentclass{article}
\usepackage[utf8]{inputenc}
\usepackage{booktabs}
\usepackage{amsmath}
\usepackage{natbib}

\title{My Document}
\author{Your Name}
\date{\today}

\begin{document}

\maketitle

\section{Introduction}
This is a sample document demonstrating various \LaTeX{} features including mathematical equations, tables, and formatting. You can cite papers using the citation search feature, and they will be automatically added to your \texttt{cite.bib} file.

\section{Mathematical Expressions}
Here are some example equations:

\subsection{Quadratic Formula}
The quadratic formula is given by:
\[x = \frac{-b \pm \sqrt{b^2 - 4ac}}{2a}\]

\subsection{Matrix Example}
A sample matrix:
\[
\begin{pmatrix}
a_{11} & a_{12} & a_{13} \\
a_{21} & a_{22} & a_{23} \\
a_{31} & a_{32} & a_{33}
\end{pmatrix}
\]

\section{Tables}
Below is a sample table showing data:

\begin{table}[h]
\centering
\begin{tabular}{@{}lrrr@{}}
\toprule
\textbf{Category} & \textbf{Value 1} & \textbf{Value 2} & \textbf{Total} \\
\midrule
Item A & 10.5 & 20.3 & 30.8 \\
Item B & 15.2 & 25.1 & 40.3 \\
Item C & 22.1 & 12.7 & 34.8 \\
\bottomrule
\end{tabular}
\caption{Sample Data Table}
\label{tab:sample}
\end{table}

\section{Lists}
\subsection{Bullet Points}
\begin{itemize}
  \item First item
  \item Second item
  \item Third item with sub-items:
    \begin{itemize}
      \item Sub-item 1
      \item Sub-item 2
    \end{itemize}
\end{itemize}

\subsection{Numbered List}
\begin{enumerate}
  \item First step
  \item Second step
  \item Third step
\end{enumerate}

\section{Citations and References}
When you use the citation search feature, references will be automatically managed. For example, you might cite a paper like this: \cite{example2024}. All citations will be automatically added to your \texttt{cite.bib} file.

\section{Conclusion}
This template demonstrates basic \LaTeX{} features including mathematical formulas, tables, and lists. The bibliography system is now set up to work seamlessly with Bibby AI's citation management. Feel free to modify it for your needs!

\bibliographystyle{plain}
\bibliography{cite}

\end{document}\documentclass{article}
\usepackage[utf8]{inputenc}
\usepackage{booktabs}
\usepackage{amsmath}
\usepackage{natbib}

\title{My Document}
\author{Your Name}
\date{\today}

\begin{document}

\maketitle

\section{Introduction}
This is a sample document demonstrating various \LaTeX{} features including mathematical equations, tables, and formatting. You can cite papers using the citation search feature, and they will be automatically added to your \texttt{cite.bib} file.

\section{Mathematical Expressions}
Here are some example equations:

\subsection{Quadratic Formula}
The quadratic formula is given by:
\[x = \frac{-b \pm \sqrt{b^2 - 4ac}}{2a}\]

\subsection{Matrix Example}
A sample matrix:
\[
\begin{pmatrix}
a_{11} & a_{12} & a_{13} \\
a_{21} & a_{22} & a_{23} \\
a_{31} & a_{32} & a_{33}
\end{pmatrix}
\]

\section{Tables}
Below is a sample table showing data:

\begin{table}[h]
\centering
\begin{tabular}{@{}lrrr@{}}
\toprule
\textbf{Category} & \textbf{Value 1} & \textbf{Value 2} & \textbf{Total} \\
\midrule
Item A & 10.5 & 20.3 & 30.8 \\
Item B & 15.2 & 25.1 & 40.3 \\
Item C & 22.1 & 12.7 & 34.8 \\
\bottomrule
\end{tabular}
\caption{Sample Data Table}
\label{tab:sample}
\end{table}

\section{Lists}
\subsection{Bullet Points}
\begin{itemize}
  \item First item
  \item Second item
  \item Third item with sub-items:
    \begin{itemize}
      \item Sub-item 1
      \item Sub-item 2
    \end{itemize}
\end{itemize}

\subsection{Numbered List}
\begin{enumerate}
  \item First step
  \item Second step
  \item Third step
\end{enumerate}

\section{Citations and References}
When you use the citation search feature, references will be automatically managed. For example, you might cite a paper like this: \cite{example2024}. All citations will be automatically added to your \texttt{cite.bib} file.

\section{Conclusion}
This template demonstrates basic \LaTeX{} features including mathematical formulas, tables, and lists. The bibliography system is now set up to work seamlessly with Bibby AI's citation management. Feel free to modify it for your needs!

\bibliographystyle{plain}
\bibliography{cite}

\end{document}\documentclass{article}
\usepackage[utf8]{inputenc}
\usepackage{booktabs}
\usepackage{amsmath}
\usepackage{natbib}

\title{My Document}
\author{Your Name}
\date{\today}

\begin{document}

\maketitle

\section{Introduction}
This is a sample document demonstrating various \LaTeX{} features including mathematical equations, tables, and formatting. You can cite papers using the citation search feature, and they will be automatically added to your \texttt{cite.bib} file.

\section{Mathematical Expressions}
Here are some example equations:

\subsection{Quadratic Formula}
The quadratic formula is given by:
\[x = \frac{-b \pm \sqrt{b^2 - 4ac}}{2a}\]

\subsection{Matrix Example}
A sample matrix:
\[
\begin{pmatrix}
a_{11} & a_{12} & a_{13} \\
a_{21} & a_{22} & a_{23} \\
a_{31} & a_{32} & a_{33}
\end{pmatrix}
\]

\section{Tables}
Below is a sample table showing data:

\begin{table}[h]
\centering
\begin{tabular}{@{}lrrr@{}}
\toprule
\textbf{Category} & \textbf{Value 1} & \textbf{Value 2} & \textbf{Total} \\
\midrule
Item A & 10.5 & 20.3 & 30.8 \\
Item B & 15.2 & 25.1 & 40.3 \\
Item C & 22.1 & 12.7 & 34.8 \\
\bottomrule
\end{tabular}
\caption{Sample Data Table}
\label{tab:sample}
\end{table}

\section{Lists}
\subsection{Bullet Points}
\begin{itemize}
  \item First item
  \item Second item
  \item Third item with sub-items:
    \begin{itemize}
      \item Sub-item 1
      \item Sub-item 2
    \end{itemize}
\end{itemize}

\subsection{Numbered List}
\begin{enumerate}
  \item First step
  \item Second step
  \item Third step
\end{enumerate}

\section{Citations and References}
When you use the citation search feature, references will be automatically managed. For example, you might cite a paper like this: \cite{example2024}. All citations will be automatically added to your \texttt{cite.bib} file.

\section{Conclusion}
This template demonstrates basic \LaTeX{} features including mathematical formulas, tables, and lists. The bibliography system is now set up to work seamlessly with Bibby AI's citation management. Feel free to modify it for your needs!

\bibliographystyle{plain}
\bibliography{cite}

\end{document}\documentclass{article}
\usepackage[utf8]{inputenc}
\usepackage{booktabs}
\usepackage{amsmath}
\usepackage{natbib}

\title{My Document}
\author{Your Name}
\date{\today}

\begin{document}

\maketitle

\section{Introduction}
This is a sample document demonstrating various \LaTeX{} features including mathematical equations, tables, and formatting. You can cite papers using the citation search feature, and they will be automatically added to your \texttt{cite.bib} file.

\section{Mathematical Expressions}
Here are some example equations:

\subsection{Quadratic Formula}
The quadratic formula is given by:
\[x = \frac{-b \pm \sqrt{b^2 - 4ac}}{2a}\]

\subsection{Matrix Example}
A sample matrix:
\[
\begin{pmatrix}
a_{11} & a_{12} & a_{13} \\
a_{21} & a_{22} & a_{23} \\
a_{31} & a_{32} & a_{33}
\end{pmatrix}
\]

\section{Tables}
Below is a sample table showing data:

\begin{table}[h]
\centering
\begin{tabular}{@{}lrrr@{}}
\toprule
\textbf{Category} & \textbf{Value 1} & \textbf{Value 2} & \textbf{Total} \\
\midrule
Item A & 10.5 & 20.3 & 30.8 \\
Item B & 15.2 & 25.1 & 40.3 \\
Item C & 22.1 & 12.7 & 34.8 \\
\bottomrule
\end{tabular}
\caption{Sample Data Table}
\label{tab:sample}
\end{table}

\section{Lists}
\subsection{Bullet Points}
\begin{itemize}
  \item First item
  \item Second item
  \item Third item with sub-items:
    \begin{itemize}
      \item Sub-item 1
      \item Sub-item 2
    \end{itemize}
\end{itemize}

\subsection{Numbered List}
\begin{enumerate}
  \item First step
  \item Second step
  \item Third step
\end{enumerate}

\section{Citations and References}
When you use the citation search feature, references will be automatically managed. For example, you might cite a paper like this: \cite{example2024}. All citations will be automatically added to your \texttt{cite.bib} file.

\section{Conclusion}
This template demonstrates basic \LaTeX{} features including mathematical formulas, tables, and lists. The bibliography system is now set up to work seamlessly with Bibby AI's citation management. Feel free to modify it for your needs!

\bibliographystyle{plain}
\bibliography{cite}

\end{document}\documentclass{article}
\usepackage[utf8]{inputenc}
\usepackage{booktabs}
\usepackage{amsmath}
\usepackage{natbib}

\title{My Document}
\author{Your Name}
\date{\today}

\begin{document}

\maketitle

\section{Introduction}
This is a sample document demonstrating various \LaTeX{} features including mathematical equations, tables, and formatting. You can cite papers using the citation search feature, and they will be automatically added to your \texttt{cite.bib} file.

\section{Mathematical Expressions}
Here are some example equations:

\subsection{Quadratic Formula}
The quadratic formula is given by:
\[x = \frac{-b \pm \sqrt{b^2 - 4ac}}{2a}\]

\subsection{Matrix Example}
A sample matrix:
\[
\begin{pmatrix}
a_{11} & a_{12} & a_{13} \\
a_{21} & a_{22} & a_{23} \\
a_{31} & a_{32} & a_{33}
\end{pmatrix}
\]

\section{Tables}
Below is a sample table showing data:

\begin{table}[h]
\centering
\begin{tabular}{@{}lrrr@{}}
\toprule
\textbf{Category} & \textbf{Value 1} & \textbf{Value 2} & \textbf{Total} \\
\midrule
Item A & 10.5 & 20.3 & 30.8 \\
Item B & 15.2 & 25.1 & 40.3 \\
Item C & 22.1 & 12.7 & 34.8 \\
\bottomrule
\end{tabular}
\caption{Sample Data Table}
\label{tab:sample}
\end{table}

\section{Lists}
\subsection{Bullet Points}
\begin{itemize}
  \item First item
  \item Second item
  \item Third item with sub-items:
    \begin{itemize}
      \item Sub-item 1
      \item Sub-item 2
    \end{itemize}
\end{itemize}

\subsection{Numbered List}
\begin{enumerate}
  \item First step
  \item Second step
  \item Third step
\end{enumerate}

\section{Citations and References}
When you use the citation search feature, references will be automatically managed. For example, you might cite a paper like this: \cite{example2024}. All citations will be automatically added to your \texttt{cite.bib} file.

\section{Conclusion}
This template demonstrates basic \LaTeX{} features including mathematical formulas, tables, and lists. The bibliography system is now set up to work seamlessly with Bibby AI's citation management. Feel free to modify it for your needs!

\bibliographystyle{plain}
\bibliography{cite}

\end{document}\documentclass{article}
\usepackage[utf8]{inputenc}
\usepackage{booktabs}
\usepackage{amsmath}
\usepackage{natbib}

\title{My Document}
\author{Your Name}
\date{\today}

\begin{document}

\maketitle

\section{Introduction}
This is a sample document demonstrating various \LaTeX{} features including mathematical equations, tables, and formatting. You can cite papers using the citation search feature, and they will be automatically added to your \texttt{cite.bib} file.

\section{Mathematical Expressions}
Here are some example equations:

\subsection{Quadratic Formula}
The quadratic formula is given by:
\[x = \frac{-b \pm \sqrt{b^2 - 4ac}}{2a}\]

\subsection{Matrix Example}
A sample matrix:
\[
\begin{pmatrix}
a_{11} & a_{12} & a_{13} \\
a_{21} & a_{22} & a_{23} \\
a_{31} & a_{32} & a_{33}
\end{pmatrix}
\]

\section{Tables}
Below is a sample table showing data:

\begin{table}[h]
\centering
\begin{tabular}{@{}lrrr@{}}
\toprule
\textbf{Category} & \textbf{Value 1} & \textbf{Value 2} & \textbf{Total} \\
\midrule
Item A & 10.5 & 20.3 & 30.8 \\
Item B & 15.2 & 25.1 & 40.3 \\
Item C & 22.1 & 12.7 & 34.8 \\
\bottomrule
\end{tabular}
\caption{Sample Data Table}
\label{tab:sample}
\end{table}

\section{Lists}
\subsection{Bullet Points}
\begin{itemize}
  \item First item
  \item Second item
  \item Third item with sub-items:
    \begin{itemize}
      \item Sub-item 1
      \item Sub-item 2
    \end{itemize}
\end{itemize}

\subsection{Numbered List}
\begin{enumerate}
  \item First step
  \item Second step
  \item Third step
\end{enumerate}

\section{Citations and References}
When you use the citation search feature, references will be automatically managed. For example, you might cite a paper like this: \cite{example2024}. All citations will be automatically added to your \texttt{cite.bib} file.

\section{Conclusion}
This template demonstrates basic \LaTeX{} features including mathematical formulas, tables, and lists. The bibliography system is now set up to work seamlessly with Bibby AI's citation management. Feel free to modify it for your needs!

\bibliographystyle{plain}
\bibliography{cite}

\end{document}\documentclass{article}
\usepackage[utf8]{inputenc}
\usepackage{booktabs}
\usepackage{amsmath}
\usepackage{natbib}

\title{My Document}
\author{Your Name}
\date{\today}

\begin{document}

\maketitle

\section{Introduction}
This is a sample document demonstrating various \LaTeX{} features including mathematical equations, tables, and formatting. You can cite papers using the citation search feature, and they will be automatically added to your \texttt{cite.bib} file.

\section{Mathematical Expressions}
Here are some example equations:

\subsection{Quadratic Formula}
The quadratic formula is given by:
\[x = \frac{-b \pm \sqrt{b^2 - 4ac}}{2a}\]

\subsection{Matrix Example}
A sample matrix:
\[
\begin{pmatrix}
a_{11} & a_{12} & a_{13} \\
a_{21} & a_{22} & a_{23} \\
a_{31} & a_{32} & a_{33}
\end{pmatrix}
\]

\section{Tables}
Below is a sample table showing data:

\begin{table}[h]
\centering
\begin{tabular}{@{}lrrr@{}}
\toprule
\textbf{Category} & \textbf{Value 1} & \textbf{Value 2} & \textbf{Total} \\
\midrule
Item A & 10.5 & 20.3 & 30.8 \\
Item B & 15.2 & 25.1 & 40.3 \\
Item C & 22.1 & 12.7 & 34.8 \\
\bottomrule
\end{tabular}
\caption{Sample Data Table}
\label{tab:sample}
\end{table}

\section{Lists}
\subsection{Bullet Points}
\begin{itemize}
  \item First item
  \item Second item
  \item Third item with sub-items:
    \begin{itemize}
      \item Sub-item 1
      \item Sub-item 2
    \end{itemize}
\end{itemize}

\subsection{Numbered List}
\begin{enumerate}
  \item First step
  \item Second step
  \item Third step
\end{enumerate}

\section{Citations and References}
When you use the citation search feature, references will be automatically managed. For example, you might cite a paper like this: \cite{example2024}. All citations will be automatically added to your \texttt{cite.bib} file.

\section{Conclusion}
This template demonstrates basic \LaTeX{} features including mathematical formulas, tables, and lists. The bibliography system is now set up to work seamlessly with Bibby AI's citation management. Feel free to modify it for your needs!

\bibliographystyle{plain}
\bibliography{cite}

\end{document}\documentclass{article}
\usepackage[utf8]{inputenc}
\usepackage{booktabs}
\usepackage{amsmath}
\usepackage{natbib}

\title{My Document}
\author{Your Name}
\date{\today}

\begin{document}

\maketitle

\section{Introduction}
This is a sample document demonstrating various \LaTeX{} features including mathematical equations, tables, and formatting. You can cite papers using the citation search feature, and they will be automatically added to your \texttt{cite.bib} file.

\section{Mathematical Expressions}
Here are some example equations:

\subsection{Quadratic Formula}
The quadratic formula is given by:
\[x = \frac{-b \pm \sqrt{b^2 - 4ac}}{2a}\]

\subsection{Matrix Example}
A sample matrix:
\[
\begin{pmatrix}
a_{11} & a_{12} & a_{13} \\
a_{21} & a_{22} & a_{23} \\
a_{31} & a_{32} & a_{33}
\end{pmatrix}
\]

\section{Tables}
Below is a sample table showing data:

\begin{table}[h]
\centering
\begin{tabular}{@{}lrrr@{}}
\toprule
\textbf{Category} & \textbf{Value 1} & \textbf{Value 2} & \textbf{Total} \\
\midrule
Item A & 10.5 & 20.3 & 30.8 \\
Item B & 15.2 & 25.1 & 40.3 \\
Item C & 22.1 & 12.7 & 34.8 \\
\bottomrule
\end{tabular}
\caption{Sample Data Table}
\label{tab:sample}
\end{table}

\section{Lists}
\subsection{Bullet Points}
\begin{itemize}
  \item First item
  \item Second item
  \item Third item with sub-items:
    \begin{itemize}
      \item Sub-item 1
      \item Sub-item 2
    \end{itemize}
\end{itemize}

\subsection{Numbered List}
\begin{enumerate}
  \item First step
  \item Second step
  \item Third step
\end{enumerate}

\section{Citations and References}
When you use the citation search feature, references will be automatically managed. For example, you might cite a paper like this: \cite{example2024}. All citations will be automatically added to your \texttt{cite.bib} file.

\section{Conclusion}
This template demonstrates basic \LaTeX{} features including mathematical formulas, tables, and lists. The bibliography system is now set up to work seamlessly with Bibby AI's citation management. Feel free to modify it for your needs!

\bibliographystyle{plain}
\bibliography{cite}

\end{document}\documentclass{article}
\usepackage[utf8]{inputenc}
\usepackage{booktabs}
\usepackage{amsmath}
\usepackage{natbib}

\title{My Document}
\author{Your Name}
\date{\today}

\begin{document}

\maketitle

\section{Introduction}
This is a sample document demonstrating various \LaTeX{} features including mathematical equations, tables, and formatting. You can cite papers using the citation search feature, and they will be automatically added to your \texttt{cite.bib} file.

\section{Mathematical Expressions}
Here are some example equations:

\subsection{Quadratic Formula}
The quadratic formula is given by:
\[x = \frac{-b \pm \sqrt{b^2 - 4ac}}{2a}\]

\subsection{Matrix Example}
A sample matrix:
\[
\begin{pmatrix}
a_{11} & a_{12} & a_{13} \\
a_{21} & a_{22} & a_{23} \\
a_{31} & a_{32} & a_{33}
\end{pmatrix}
\]

\section{Tables}
Below is a sample table showing data:

\begin{table}[h]
\centering
\begin{tabular}{@{}lrrr@{}}
\toprule
\textbf{Category} & \textbf{Value 1} & \textbf{Value 2} & \textbf{Total} \\
\midrule
Item A & 10.5 & 20.3 & 30.8 \\
Item B & 15.2 & 25.1 & 40.3 \\
Item C & 22.1 & 12.7 & 34.8 \\
\bottomrule
\end{tabular}
\caption{Sample Data Table}
\label{tab:sample}
\end{table}

\section{Lists}
\subsection{Bullet Points}
\begin{itemize}
  \item First item
  \item Second item
  \item Third item with sub-items:
    \begin{itemize}
      \item Sub-item 1
      \item Sub-item 2
    \end{itemize}
\end{itemize}

\subsection{Numbered List}
\begin{enumerate}
  \item First step
  \item Second step
  \item Third step
\end{enumerate}

\section{Citations and References}
When you use the citation search feature, references will be automatically managed. For example, you might cite a paper like this: \cite{example2024}. All citations will be automatically added to your \texttt{cite.bib} file.

\section{Conclusion}
This template demonstrates basic \LaTeX{} features including mathematical formulas, tables, and lists. The bibliography system is now set up to work seamlessly with Bibby AI's citation management. Feel free to modify it for your needs!

\bibliographystyle{plain}
\bibliography{cite}

\end{document}\documentclass{article}
\usepackage[utf8]{inputenc}
\usepackage{booktabs}
\usepackage{amsmath}
\usepackage{natbib}

\title{My Document}
\author{Your Name}
\date{\today}

\begin{document}

\maketitle

\section{Introduction}
This is a sample document demonstrating various \LaTeX{} features including mathematical equations, tables, and formatting. You can cite papers using the citation search feature, and they will be automatically added to your \texttt{cite.bib} file.

\section{Mathematical Expressions}
Here are some example equations:

\subsection{Quadratic Formula}
The quadratic formula is given by:
\[x = \frac{-b \pm \sqrt{b^2 - 4ac}}{2a}\]

\subsection{Matrix Example}
A sample matrix:
\[
\begin{pmatrix}
a_{11} & a_{12} & a_{13} \\
a_{21} & a_{22} & a_{23} \\
a_{31} & a_{32} & a_{33}
\end{pmatrix}
\]

\section{Tables}
Below is a sample table showing data:

\begin{table}[h]
\centering
\begin{tabular}{@{}lrrr@{}}
\toprule
\textbf{Category} & \textbf{Value 1} & \textbf{Value 2} & \textbf{Total} \\
\midrule
Item A & 10.5 & 20.3 & 30.8 \\
Item B & 15.2 & 25.1 & 40.3 \\
Item C & 22.1 & 12.7 & 34.8 \\
\bottomrule
\end{tabular}
\caption{Sample Data Table}
\label{tab:sample}
\end{table}

\section{Lists}
\subsection{Bullet Points}
\begin{itemize}
  \item First item
  \item Second item
  \item Third item with sub-items:
    \begin{itemize}
      \item Sub-item 1
      \item Sub-item 2
    \end{itemize}
\end{itemize}

\subsection{Numbered List}
\begin{enumerate}
  \item First step
  \item Second step
  \item Third step
\end{enumerate}

\section{Citations and References}
When you use the citation search feature, references will be automatically managed. For example, you might cite a paper like this: \cite{example2024}. All citations will be automatically added to your \texttt{cite.bib} file.

\section{Conclusion}
This template demonstrates basic \LaTeX{} features including mathematical formulas, tables, and lists. The bibliography system is now set up to work seamlessly with Bibby AI's citation management. Feel free to modify it for your needs!

\bibliographystyle{plain}
\bibliography{cite}

\end{document}\documentclass{article}
\usepackage[utf8]{inputenc}
\usepackage{booktabs}
\usepackage{amsmath}
\usepackage{natbib}

\title{My Document}
\author{Your Name}
\date{\today}

\begin{document}

\maketitle

\section{Introduction}
This is a sample document demonstrating various \LaTeX{} features including mathematical equations, tables, and formatting. You can cite papers using the citation search feature, and they will be automatically added to your \texttt{cite.bib} file.

\section{Mathematical Expressions}
Here are some example equations:

\subsection{Quadratic Formula}
The quadratic formula is given by:
\[x = \frac{-b \pm \sqrt{b^2 - 4ac}}{2a}\]

\subsection{Matrix Example}
A sample matrix:
\[
\begin{pmatrix}
a_{11} & a_{12} & a_{13} \\
a_{21} & a_{22} & a_{23} \\
a_{31} & a_{32} & a_{33}
\end{pmatrix}
\]

\section{Tables}
Below is a sample table showing data:

\begin{table}[h]
\centering
\begin{tabular}{@{}lrrr@{}}
\toprule
\textbf{Category} & \textbf{Value 1} & \textbf{Value 2} & \textbf{Total} \\
\midrule
Item A & 10.5 & 20.3 & 30.8 \\
Item B & 15.2 & 25.1 & 40.3 \\
Item C & 22.1 & 12.7 & 34.8 \\
\bottomrule
\end{tabular}
\caption{Sample Data Table}
\label{tab:sample}
\end{table}

\section{Lists}
\subsection{Bullet Points}
\begin{itemize}
  \item First item
  \item Second item
  \item Third item with sub-items:
    \begin{itemize}
      \item Sub-item 1
      \item Sub-item 2
    \end{itemize}
\end{itemize}

\subsection{Numbered List}
\begin{enumerate}
  \item First step
  \item Second step
  \item Third step
\end{enumerate}

\section{Citations and References}
When you use the citation search feature, references will be automatically managed. For example, you might cite a paper like this: \cite{example2024}. All citations will be automatically added to your \texttt{cite.bib} file.

\section{Conclusion}
This template demonstrates basic \LaTeX{} features including mathematical formulas, tables, and lists. The bibliography system is now set up to work seamlessly with Bibby AI's citation management. Feel free to modify it for your needs!

\bibliographystyle{plain}
\bibliography{cite}

\end{document}\documentclass{article}
\usepackage[utf8]{inputenc}
\usepackage{booktabs}
\usepackage{amsmath}
\usepackage{natbib}

\title{My Document}
\author{Your Name}
\date{\today}

\begin{document}

\maketitle

\section{Introduction}
This is a sample document demonstrating various \LaTeX{} features including mathematical equations, tables, and formatting. You can cite papers using the citation search feature, and they will be automatically added to your \texttt{cite.bib} file.

\section{Mathematical Expressions}
Here are some example equations:

\subsection{Quadratic Formula}
The quadratic formula is given by:
\[x = \frac{-b \pm \sqrt{b^2 - 4ac}}{2a}\]

\subsection{Matrix Example}
A sample matrix:
\[
\begin{pmatrix}
a_{11} & a_{12} & a_{13} \\
a_{21} & a_{22} & a_{23} \\
a_{31} & a_{32} & a_{33}
\end{pmatrix}
\]

\section{Tables}
Below is a sample table showing data:

\begin{table}[h]
\centering
\begin{tabular}{@{}lrrr@{}}
\toprule
\textbf{Category} & \textbf{Value 1} & \textbf{Value 2} & \textbf{Total} \\
\midrule
Item A & 10.5 & 20.3 & 30.8 \\
Item B & 15.2 & 25.1 & 40.3 \\
Item C & 22.1 & 12.7 & 34.8 \\
\bottomrule
\end{tabular}
\caption{Sample Data Table}
\label{tab:sample}
\end{table}

\section{Lists}
\subsection{Bullet Points}
\begin{itemize}
  \item First item
  \item Second item
  \item Third item with sub-items:
    \begin{itemize}
      \item Sub-item 1
      \item Sub-item 2
    \end{itemize}
\end{itemize}

\subsection{Numbered List}
\begin{enumerate}
  \item First step
  \item Second step
  \item Third step
\end{enumerate}

\section{Citations and References}
When you use the citation search feature, references will be automatically managed. For example, you might cite a paper like this: \cite{example2024}. All citations will be automatically added to your \texttt{cite.bib} file.

\section{Conclusion}
This template demonstrates basic \LaTeX{} features including mathematical formulas, tables, and lists. The bibliography system is now set up to work seamlessly with Bibby AI's citation management. Feel free to modify it for your needs!

\bibliographystyle{plain}
\bibliography{cite}

\end{document}\documentclass{article}
\usepackage[utf8]{inputenc}
\usepackage{booktabs}
\usepackage{amsmath}
\usepackage{natbib}

\title{My Document}
\author{Your Name}
\date{\today}

\begin{document}

\maketitle

\section{Introduction}
This is a sample document demonstrating various \LaTeX{} features including mathematical equations, tables, and formatting. You can cite papers using the citation search feature, and they will be automatically added to your \texttt{cite.bib} file.

\section{Mathematical Expressions}
Here are some example equations:

\subsection{Quadratic Formula}
The quadratic formula is given by:
\[x = \frac{-b \pm \sqrt{b^2 - 4ac}}{2a}\]

\subsection{Matrix Example}
A sample matrix:
\[
\begin{pmatrix}
a_{11} & a_{12} & a_{13} \\
a_{21} & a_{22} & a_{23} \\
a_{31} & a_{32} & a_{33}
\end{pmatrix}
\]

\section{Tables}
Below is a sample table showing data:

\begin{table}[h]
\centering
\begin{tabular}{@{}lrrr@{}}
\toprule
\textbf{Category} & \textbf{Value 1} & \textbf{Value 2} & \textbf{Total} \\
\midrule
Item A & 10.5 & 20.3 & 30.8 \\
Item B & 15.2 & 25.1 & 40.3 \\
Item C & 22.1 & 12.7 & 34.8 \\
\bottomrule
\end{tabular}
\caption{Sample Data Table}
\label{tab:sample}
\end{table}

\section{Lists}
\subsection{Bullet Points}
\begin{itemize}
  \item First item
  \item Second item
  \item Third item with sub-items:
    \begin{itemize}
      \item Sub-item 1
      \item Sub-item 2
    \end{itemize}
\end{itemize}

\subsection{Numbered List}
\begin{enumerate}
  \item First step
  \item Second step
  \item Third step
\end{enumerate}

\section{Citations and References}
When you use the citation search feature, references will be automatically managed. For example, you might cite a paper like this: \cite{example2024}. All citations will be automatically added to your \texttt{cite.bib} file.

\section{Conclusion}
This template demonstrates basic \LaTeX{} features including mathematical formulas, tables, and lists. The bibliography system is now set up to work seamlessly with Bibby AI's citation management. Feel free to modify it for your needs!

\bibliographystyle{plain}
\bibliography{cite}

\end{document}\documentclass{article}
\usepackage[utf8]{inputenc}
\usepackage{booktabs}
\usepackage{amsmath}
\usepackage{natbib}

\title{My Document}
\author{Your Name}
\date{\today}

\begin{document}

\maketitle

\section{Introduction}
This is a sample document demonstrating various \LaTeX{} features including mathematical equations, tables, and formatting. You can cite papers using the citation search feature, and they will be automatically added to your \texttt{cite.bib} file.

\section{Mathematical Expressions}
Here are some example equations:

\subsection{Quadratic Formula}
The quadratic formula is given by:
\[x = \frac{-b \pm \sqrt{b^2 - 4ac}}{2a}\]

\subsection{Matrix Example}
A sample matrix:
\[
\begin{pmatrix}
a_{11} & a_{12} & a_{13} \\
a_{21} & a_{22} & a_{23} \\
a_{31} & a_{32} & a_{33}
\end{pmatrix}
\]

\section{Tables}
Below is a sample table showing data:

\begin{table}[h]
\centering
\begin{tabular}{@{}lrrr@{}}
\toprule
\textbf{Category} & \textbf{Value 1} & \textbf{Value 2} & \textbf{Total} \\
\midrule
Item A & 10.5 & 20.3 & 30.8 \\
Item B & 15.2 & 25.1 & 40.3 \\
Item C & 22.1 & 12.7 & 34.8 \\
\bottomrule
\end{tabular}
\caption{Sample Data Table}
\label{tab:sample}
\end{table}

\section{Lists}
\subsection{Bullet Points}
\begin{itemize}
  \item First item
  \item Second item
  \item Third item with sub-items:
    \begin{itemize}
      \item Sub-item 1
      \item Sub-item 2
    \end{itemize}
\end{itemize}

\subsection{Numbered List}
\begin{enumerate}
  \item First step
  \item Second step
  \item Third step
\end{enumerate}

\section{Citations and References}
When you use the citation search feature, references will be automatically managed. For example, you might cite a paper like this: \cite{example2024}. All citations will be automatically added to your \texttt{cite.bib} file.

\section{Conclusion}
This template demonstrates basic \LaTeX{} features including mathematical formulas, tables, and lists. The bibliography system is now set up to work seamlessly with Bibby AI's citation management. Feel free to modify it for your needs!

\bibliographystyle{plain}
\bibliography{cite}

\end{document}\documentclass{article}
\usepackage[utf8]{inputenc}
\usepackage{booktabs}
\usepackage{amsmath}
\usepackage{natbib}

\title{My Document}
\author{Your Name}
\date{\today}

\begin{document}

\maketitle

\section{Introduction}
This is a sample document demonstrating various \LaTeX{} features including mathematical equations, tables, and formatting. You can cite papers using the citation search feature, and they will be automatically added to your \texttt{cite.bib} file.

\section{Mathematical Expressions}
Here are some example equations:

\subsection{Quadratic Formula}
The quadratic formula is given by:
\[x = \frac{-b \pm \sqrt{b^2 - 4ac}}{2a}\]

\subsection{Matrix Example}
A sample matrix:
\[
\begin{pmatrix}
a_{11} & a_{12} & a_{13} \\
a_{21} & a_{22} & a_{23} \\
a_{31} & a_{32} & a_{33}
\end{pmatrix}
\]

\section{Tables}
Below is a sample table showing data:

\begin{table}[h]
\centering
\begin{tabular}{@{}lrrr@{}}
\toprule
\textbf{Category} & \textbf{Value 1} & \textbf{Value 2} & \textbf{Total} \\
\midrule
Item A & 10.5 & 20.3 & 30.8 \\
Item B & 15.2 & 25.1 & 40.3 \\
Item C & 22.1 & 12.7 & 34.8 \\
\bottomrule
\end{tabular}
\caption{Sample Data Table}
\label{tab:sample}
\end{table}

\section{Lists}
\subsection{Bullet Points}
\begin{itemize}
  \item First item
  \item Second item
  \item Third item with sub-items:
    \begin{itemize}
      \item Sub-item 1
      \item Sub-item 2
    \end{itemize}
\end{itemize}

\subsection{Numbered List}
\begin{enumerate}
  \item First step
  \item Second step
  \item Third step
\end{enumerate}

\section{Citations and References}
When you use the citation search feature, references will be automatically managed. For example, you might cite a paper like this: \cite{example2024}. All citations will be automatically added to your \texttt{cite.bib} file.

\section{Conclusion}
This template demonstrates basic \LaTeX{} features including mathematical formulas, tables, and lists. The bibliography system is now set up to work seamlessly with Bibby AI's citation management. Feel free to modify it for your needs!

\bibliographystyle{plain}
\bibliography{cite}

\end{document}\documentclass{article}
\usepackage[utf8]{inputenc}
\usepackage{booktabs}
\usepackage{amsmath}
\usepackage{natbib}

\title{My Document}
\author{Your Name}
\date{\today}

\begin{document}

\maketitle

\section{Introduction}
This is a sample document demonstrating various \LaTeX{} features including mathematical equations, tables, and formatting. You can cite papers using the citation search feature, and they will be automatically added to your \texttt{cite.bib} file.

\section{Mathematical Expressions}
Here are some example equations:

\subsection{Quadratic Formula}
The quadratic formula is given by:
\[x = \frac{-b \pm \sqrt{b^2 - 4ac}}{2a}\]

\subsection{Matrix Example}
A sample matrix:
\[
\begin{pmatrix}
a_{11} & a_{12} & a_{13} \\
a_{21} & a_{22} & a_{23} \\
a_{31} & a_{32} & a_{33}
\end{pmatrix}
\]

\section{Tables}
Below is a sample table showing data:

\begin{table}[h]
\centering
\begin{tabular}{@{}lrrr@{}}
\toprule
\textbf{Category} & \textbf{Value 1} & \textbf{Value 2} & \textbf{Total} \\
\midrule
Item A & 10.5 & 20.3 & 30.8 \\
Item B & 15.2 & 25.1 & 40.3 \\
Item C & 22.1 & 12.7 & 34.8 \\
\bottomrule
\end{tabular}
\caption{Sample Data Table}
\label{tab:sample}
\end{table}

\section{Lists}
\subsection{Bullet Points}
\begin{itemize}
  \item First item
  \item Second item
  \item Third item with sub-items:
    \begin{itemize}
      \item Sub-item 1
      \item Sub-item 2
    \end{itemize}
\end{itemize}

\subsection{Numbered List}
\begin{enumerate}
  \item First step
  \item Second step
  \item Third step
\end{enumerate}

\section{Citations and References}
When you use the citation search feature, references will be automatically managed. For example, you might cite a paper like this: \cite{example2024}. All citations will be automatically added to your \texttt{cite.bib} file.

\section{Conclusion}
This template demonstrates basic \LaTeX{} features including mathematical formulas, tables, and lists. The bibliography system is now set up to work seamlessly with Bibby AI's citation management. Feel free to modify it for your needs!

\bibliographystyle{plain}
\bibliography{cite}

\end{document}\documentclass{article}
\usepackage[utf8]{inputenc}
\usepackage{booktabs}
\usepackage{amsmath}
\usepackage{natbib}

\title{My Document}
\author{Your Name}
\date{\today}

\begin{document}

\maketitle

\section{Introduction}
This is a sample document demonstrating various \LaTeX{} features including mathematical equations, tables, and formatting. You can cite papers using the citation search feature, and they will be automatically added to your \texttt{cite.bib} file.

\section{Mathematical Expressions}
Here are some example equations:

\subsection{Quadratic Formula}
The quadratic formula is given by:
\[x = \frac{-b \pm \sqrt{b^2 - 4ac}}{2a}\]

\subsection{Matrix Example}
A sample matrix:
\[
\begin{pmatrix}
a_{11} & a_{12} & a_{13} \\
a_{21} & a_{22} & a_{23} \\
a_{31} & a_{32} & a_{33}
\end{pmatrix}
\]

\section{Tables}
Below is a sample table showing data:

\begin{table}[h]
\centering
\begin{tabular}{@{}lrrr@{}}
\toprule
\textbf{Category} & \textbf{Value 1} & \textbf{Value 2} & \textbf{Total} \\
\midrule
Item A & 10.5 & 20.3 & 30.8 \\
Item B & 15.2 & 25.1 & 40.3 \\
Item C & 22.1 & 12.7 & 34.8 \\
\bottomrule
\end{tabular}
\caption{Sample Data Table}
\label{tab:sample}
\end{table}

\section{Lists}
\subsection{Bullet Points}
\begin{itemize}
  \item First item
  \item Second item
  \item Third item with sub-items:
    \begin{itemize}
      \item Sub-item 1
      \item Sub-item 2
    \end{itemize}
\end{itemize}

\subsection{Numbered List}
\begin{enumerate}
  \item First step
  \item Second step
  \item Third step
\end{enumerate}

\section{Citations and References}
When you use the citation search feature, references will be automatically managed. For example, you might cite a paper like this: \cite{example2024}. All citations will be automatically added to your \texttt{cite.bib} file.

\section{Conclusion}
This template demonstrates basic \LaTeX{} features including mathematical formulas, tables, and lists. The bibliography system is now set up to work seamlessly with Bibby AI's citation management. Feel free to modify it for your needs!

\bibliographystyle{plain}
\bibliography{cite}

\end{document}\documentclass{article}
\usepackage[utf8]{inputenc}
\usepackage{booktabs}
\usepackage{amsmath}
\usepackage{natbib}

\title{My Document}
\author{Your Name}
\date{\today}

\begin{document}

\maketitle

\section{Introduction}
This is a sample document demonstrating various \LaTeX{} features including mathematical equations, tables, and formatting. You can cite papers using the citation search feature, and they will be automatically added to your \texttt{cite.bib} file.

\section{Mathematical Expressions}
Here are some example equations:

\subsection{Quadratic Formula}
The quadratic formula is given by:
\[x = \frac{-b \pm \sqrt{b^2 - 4ac}}{2a}\]

\subsection{Matrix Example}
A sample matrix:
\[
\begin{pmatrix}
a_{11} & a_{12} & a_{13} \\
a_{21} & a_{22} & a_{23} \\
a_{31} & a_{32} & a_{33}
\end{pmatrix}
\]

\section{Tables}
Below is a sample table showing data:

\begin{table}[h]
\centering
\begin{tabular}{@{}lrrr@{}}
\toprule
\textbf{Category} & \textbf{Value 1} & \textbf{Value 2} & \textbf{Total} \\
\midrule
Item A & 10.5 & 20.3 & 30.8 \\
Item B & 15.2 & 25.1 & 40.3 \\
Item C & 22.1 & 12.7 & 34.8 \\
\bottomrule
\end{tabular}
\caption{Sample Data Table}
\label{tab:sample}
\end{table}

\section{Lists}
\subsection{Bullet Points}
\begin{itemize}
  \item First item
  \item Second item
  \item Third item with sub-items:
    \begin{itemize}
      \item Sub-item 1
      \item Sub-item 2
    \end{itemize}
\end{itemize}

\subsection{Numbered List}
\begin{enumerate}
  \item First step
  \item Second step
  \item Third step
\end{enumerate}

\section{Citations and References}
When you use the citation search feature, references will be automatically managed. For example, you might cite a paper like this: \cite{example2024}. All citations will be automatically added to your \texttt{cite.bib} file.

\section{Conclusion}
This template demonstrates basic \LaTeX{} features including mathematical formulas, tables, and lists. The bibliography system is now set up to work seamlessly with Bibby AI's citation management. Feel free to modify it for your needs!

\bibliographystyle{plain}
\bibliography{cite}

\end{document}\documentclass{article}
\usepackage[utf8]{inputenc}
\usepackage{booktabs}
\usepackage{amsmath}
\usepackage{natbib}

\title{My Document}
\author{Your Name}
\date{\today}

\begin{document}

\maketitle

\section{Introduction}
This is a sample document demonstrating various \LaTeX{} features including mathematical equations, tables, and formatting. You can cite papers using the citation search feature, and they will be automatically added to your \texttt{cite.bib} file.

\section{Mathematical Expressions}
Here are some example equations:

\subsection{Quadratic Formula}
The quadratic formula is given by:
\[x = \frac{-b \pm \sqrt{b^2 - 4ac}}{2a}\]

\subsection{Matrix Example}
A sample matrix:
\[
\begin{pmatrix}
a_{11} & a_{12} & a_{13} \\
a_{21} & a_{22} & a_{23} \\
a_{31} & a_{32} & a_{33}
\end{pmatrix}
\]

\section{Tables}
Below is a sample table showing data:

\begin{table}[h]
\centering
\begin{tabular}{@{}lrrr@{}}
\toprule
\textbf{Category} & \textbf{Value 1} & \textbf{Value 2} & \textbf{Total} \\
\midrule
Item A & 10.5 & 20.3 & 30.8 \\
Item B & 15.2 & 25.1 & 40.3 \\
Item C & 22.1 & 12.7 & 34.8 \\
\bottomrule
\end{tabular}
\caption{Sample Data Table}
\label{tab:sample}
\end{table}

\section{Lists}
\subsection{Bullet Points}
\begin{itemize}
  \item First item
  \item Second item
  \item Third item with sub-items:
    \begin{itemize}
      \item Sub-item 1
      \item Sub-item 2
    \end{itemize}
\end{itemize}

\subsection{Numbered List}
\begin{enumerate}
  \item First step
  \item Second step
  \item Third step
\end{enumerate}

\section{Citations and References}
When you use the citation search feature, references will be automatically managed. For example, you might cite a paper like this: \cite{example2024}. All citations will be automatically added to your \texttt{cite.bib} file.

\section{Conclusion}
This template demonstrates basic \LaTeX{} features including mathematical formulas, tables, and lists. The bibliography system is now set up to work seamlessly with Bibby AI's citation management. Feel free to modify it for your needs!

\bibliographystyle{plain}
\bibliography{cite}

\end{document}\documentclass{article}
\usepackage[utf8]{inputenc}
\usepackage{booktabs}
\usepackage{amsmath}
\usepackage{natbib}

\title{My Document}
\author{Your Name}
\date{\today}

\begin{document}

\maketitle

\section{Introduction}
This is a sample document demonstrating various \LaTeX{} features including mathematical equations, tables, and formatting. You can cite papers using the citation search feature, and they will be automatically added to your \texttt{cite.bib} file.

\section{Mathematical Expressions}
Here are some example equations:

\subsection{Quadratic Formula}
The quadratic formula is given by:
\[x = \frac{-b \pm \sqrt{b^2 - 4ac}}{2a}\]

\subsection{Matrix Example}
A sample matrix:
\[
\begin{pmatrix}
a_{11} & a_{12} & a_{13} \\
a_{21} & a_{22} & a_{23} \\
a_{31} & a_{32} & a_{33}
\end{pmatrix}
\]

\section{Tables}
Below is a sample table showing data:

\begin{table}[h]
\centering
\begin{tabular}{@{}lrrr@{}}
\toprule
\textbf{Category} & \textbf{Value 1} & \textbf{Value 2} & \textbf{Total} \\
\midrule
Item A & 10.5 & 20.3 & 30.8 \\
Item B & 15.2 & 25.1 & 40.3 \\
Item C & 22.1 & 12.7 & 34.8 \\
\bottomrule
\end{tabular}
\caption{Sample Data Table}
\label{tab:sample}
\end{table}

\section{Lists}
\subsection{Bullet Points}
\begin{itemize}
  \item First item
  \item Second item
  \item Third item with sub-items:
    \begin{itemize}
      \item Sub-item 1
      \item Sub-item 2
    \end{itemize}
\end{itemize}

\subsection{Numbered List}
\begin{enumerate}
  \item First step
  \item Second step
  \item Third step
\end{enumerate}

\section{Citations and References}
When you use the citation search feature, references will be automatically managed. For example, you might cite a paper like this: \cite{example2024}. All citations will be automatically added to your \texttt{cite.bib} file.

\section{Conclusion}
This template demonstrates basic \LaTeX{} features including mathematical formulas, tables, and lists. The bibliography system is now set up to work seamlessly with Bibby AI's citation management. Feel free to modify it for your needs!

\bibliographystyle{plain}
\bibliography{cite}

\end{document}\documentclass{article}
\usepackage[utf8]{inputenc}
\usepackage{booktabs}
\usepackage{amsmath}
\usepackage{natbib}

\title{My Document}
\author{Your Name}
\date{\today}

\begin{document}

\maketitle

\section{Introduction}
This is a sample document demonstrating various \LaTeX{} features including mathematical equations, tables, and formatting. You can cite papers using the citation search feature, and they will be automatically added to your \texttt{cite.bib} file.

\section{Mathematical Expressions}
Here are some example equations:

\subsection{Quadratic Formula}
The quadratic formula is given by:
\[x = \frac{-b \pm \sqrt{b^2 - 4ac}}{2a}\]

\subsection{Matrix Example}
A sample matrix:
\[
\begin{pmatrix}
a_{11} & a_{12} & a_{13} \\
a_{21} & a_{22} & a_{23} \\
a_{31} & a_{32} & a_{33}
\end{pmatrix}
\]

\section{Tables}
Below is a sample table showing data:

\begin{table}[h]
\centering
\begin{tabular}{@{}lrrr@{}}
\toprule
\textbf{Category} & \textbf{Value 1} & \textbf{Value 2} & \textbf{Total} \\
\midrule
Item A & 10.5 & 20.3 & 30.8 \\
Item B & 15.2 & 25.1 & 40.3 \\
Item C & 22.1 & 12.7 & 34.8 \\
\bottomrule
\end{tabular}
\caption{Sample Data Table}
\label{tab:sample}
\end{table}

\section{Lists}
\subsection{Bullet Points}
\begin{itemize}
  \item First item
  \item Second item
  \item Third item with sub-items:
    \begin{itemize}
      \item Sub-item 1
      \item Sub-item 2
    \end{itemize}
\end{itemize}

\subsection{Numbered List}
\begin{enumerate}
  \item First step
  \item Second step
  \item Third step
\end{enumerate}

\section{Citations and References}
When you use the citation search feature, references will be automatically managed. For example, you might cite a paper like this: \cite{example2024}. All citations will be automatically added to your \texttt{cite.bib} file.

\section{Conclusion}
This template demonstrates basic \LaTeX{} features including mathematical formulas, tables, and lists. The bibliography system is now set up to work seamlessly with Bibby AI's citation management. Feel free to modify it for your needs!

\bibliographystyle{plain}
\bibliography{cite}

\end{document}\documentclass{article}
\usepackage[utf8]{inputenc}
\usepackage{booktabs}
\usepackage{amsmath}
\usepackage{natbib}

\title{My Document}
\author{Your Name}
\date{\today}

\begin{document}

\maketitle

\section{Introduction}
This is a sample document demonstrating various \LaTeX{} features including mathematical equations, tables, and formatting. You can cite papers using the citation search feature, and they will be automatically added to your \texttt{cite.bib} file.

\section{Mathematical Expressions}
Here are some example equations:

\subsection{Quadratic Formula}
The quadratic formula is given by:
\[x = \frac{-b \pm \sqrt{b^2 - 4ac}}{2a}\]

\subsection{Matrix Example}
A sample matrix:
\[
\begin{pmatrix}
a_{11} & a_{12} & a_{13} \\
a_{21} & a_{22} & a_{23} \\
a_{31} & a_{32} & a_{33}
\end{pmatrix}
\]

\section{Tables}
Below is a sample table showing data:

\begin{table}[h]
\centering
\begin{tabular}{@{}lrrr@{}}
\toprule
\textbf{Category} & \textbf{Value 1} & \textbf{Value 2} & \textbf{Total} \\
\midrule
Item A & 10.5 & 20.3 & 30.8 \\
Item B & 15.2 & 25.1 & 40.3 \\
Item C & 22.1 & 12.7 & 34.8 \\
\bottomrule
\end{tabular}
\caption{Sample Data Table}
\label{tab:sample}
\end{table}

\section{Lists}
\subsection{Bullet Points}
\begin{itemize}
  \item First item
  \item Second item
  \item Third item with sub-items:
    \begin{itemize}
      \item Sub-item 1
      \item Sub-item 2
    \end{itemize}
\end{itemize}

\subsection{Numbered List}
\begin{enumerate}
  \item First step
  \item Second step
  \item Third step
\end{enumerate}

\section{Citations and References}
When you use the citation search feature, references will be automatically managed. For example, you might cite a paper like this: \cite{example2024}. All citations will be automatically added to your \texttt{cite.bib} file.

\section{Conclusion}
This template demonstrates basic \LaTeX{} features including mathematical formulas, tables, and lists. The bibliography system is now set up to work seamlessly with Bibby AI's citation management. Feel free to modify it for your needs!

\bibliographystyle{plain}
\bibliography{cite}

\end{document}\documentclass{article}
\usepackage[utf8]{inputenc}
\usepackage{booktabs}
\usepackage{amsmath}
\usepackage{natbib}

\title{My Document}
\author{Your Name}
\date{\today}

\begin{document}

\maketitle

\section{Introduction}
This is a sample document demonstrating various \LaTeX{} features including mathematical equations, tables, and formatting. You can cite papers using the citation search feature, and they will be automatically added to your \texttt{cite.bib} file.

\section{Mathematical Expressions}
Here are some example equations:

\subsection{Quadratic Formula}
The quadratic formula is given by:
\[x = \frac{-b \pm \sqrt{b^2 - 4ac}}{2a}\]

\subsection{Matrix Example}
A sample matrix:
\[
\begin{pmatrix}
a_{11} & a_{12} & a_{13} \\
a_{21} & a_{22} & a_{23} \\
a_{31} & a_{32} & a_{33}
\end{pmatrix}
\]

\section{Tables}
Below is a sample table showing data:

\begin{table}[h]
\centering
\begin{tabular}{@{}lrrr@{}}
\toprule
\textbf{Category} & \textbf{Value 1} & \textbf{Value 2} & \textbf{Total} \\
\midrule
Item A & 10.5 & 20.3 & 30.8 \\
Item B & 15.2 & 25.1 & 40.3 \\
Item C & 22.1 & 12.7 & 34.8 \\
\bottomrule
\end{tabular}
\caption{Sample Data Table}
\label{tab:sample}
\end{table}

\section{Lists}
\subsection{Bullet Points}
\begin{itemize}
  \item First item
  \item Second item
  \item Third item with sub-items:
    \begin{itemize}
      \item Sub-item 1
      \item Sub-item 2
    \end{itemize}
\end{itemize}

\subsection{Numbered List}
\begin{enumerate}
  \item First step
  \item Second step
  \item Third step
\end{enumerate}

\section{Citations and References}
When you use the citation search feature, references will be automatically managed. For example, you might cite a paper like this: \cite{example2024}. All citations will be automatically added to your \texttt{cite.bib} file.

\section{Conclusion}
This template demonstrates basic \LaTeX{} features including mathematical formulas, tables, and lists. The bibliography system is now set up to work seamlessly with Bibby AI's citation management. Feel free to modify it for your needs!

\bibliographystyle{plain}
\bibliography{cite}

\end{document}\documentclass{article}
\usepackage[utf8]{inputenc}
\usepackage{booktabs}
\usepackage{amsmath}
\usepackage{natbib}

\title{My Document}
\author{Your Name}
\date{\today}

\begin{document}

\maketitle

\section{Introduction}
This is a sample document demonstrating various \LaTeX{} features including mathematical equations, tables, and formatting. You can cite papers using the citation search feature, and they will be automatically added to your \texttt{cite.bib} file.

\section{Mathematical Expressions}
Here are some example equations:

\subsection{Quadratic Formula}
The quadratic formula is given by:
\[x = \frac{-b \pm \sqrt{b^2 - 4ac}}{2a}\]

\subsection{Matrix Example}
A sample matrix:
\[
\begin{pmatrix}
a_{11} & a_{12} & a_{13} \\
a_{21} & a_{22} & a_{23} \\
a_{31} & a_{32} & a_{33}
\end{pmatrix}
\]

\section{Tables}
Below is a sample table showing data:

\begin{table}[h]
\centering
\begin{tabular}{@{}lrrr@{}}
\toprule
\textbf{Category} & \textbf{Value 1} & \textbf{Value 2} & \textbf{Total} \\
\midrule
Item A & 10.5 & 20.3 & 30.8 \\
Item B & 15.2 & 25.1 & 40.3 \\
Item C & 22.1 & 12.7 & 34.8 \\
\bottomrule
\end{tabular}
\caption{Sample Data Table}
\label{tab:sample}
\end{table}

\section{Lists}
\subsection{Bullet Points}
\begin{itemize}
  \item First item
  \item Second item
  \item Third item with sub-items:
    \begin{itemize}
      \item Sub-item 1
      \item Sub-item 2
    \end{itemize}
\end{itemize}

\subsection{Numbered List}
\begin{enumerate}
  \item First step
  \item Second step
  \item Third step
\end{enumerate}

\section{Citations and References}
When you use the citation search feature, references will be automatically managed. For example, you might cite a paper like this: \cite{example2024}. All citations will be automatically added to your \texttt{cite.bib} file.

\section{Conclusion}
This template demonstrates basic \LaTeX{} features including mathematical formulas, tables, and lists. The bibliography system is now set up to work seamlessly with Bibby AI's citation management. Feel free to modify it for your needs!

\bibliographystyle{plain}
\bibliography{cite}

\end{document}\documentclass{article}
\usepackage[utf8]{inputenc}
\usepackage{booktabs}
\usepackage{amsmath}
\usepackage{natbib}

\title{My Document}
\author{Your Name}
\date{\today}

\begin{document}

\maketitle

\section{Introduction}
This is a sample document demonstrating various \LaTeX{} features including mathematical equations, tables, and formatting. You can cite papers using the citation search feature, and they will be automatically added to your \texttt{cite.bib} file.

\section{Mathematical Expressions}
Here are some example equations:

\subsection{Quadratic Formula}
The quadratic formula is given by:
\[x = \frac{-b \pm \sqrt{b^2 - 4ac}}{2a}\]

\subsection{Matrix Example}
A sample matrix:
\[
\begin{pmatrix}
a_{11} & a_{12} & a_{13} \\
a_{21} & a_{22} & a_{23} \\
a_{31} & a_{32} & a_{33}
\end{pmatrix}
\]

\section{Tables}
Below is a sample table showing data:

\begin{table}[h]
\centering
\begin{tabular}{@{}lrrr@{}}
\toprule
\textbf{Category} & \textbf{Value 1} & \textbf{Value 2} & \textbf{Total} \\
\midrule
Item A & 10.5 & 20.3 & 30.8 \\
Item B & 15.2 & 25.1 & 40.3 \\
Item C & 22.1 & 12.7 & 34.8 \\
\bottomrule
\end{tabular}
\caption{Sample Data Table}
\label{tab:sample}
\end{table}

\section{Lists}
\subsection{Bullet Points}
\begin{itemize}
  \item First item
  \item Second item
  \item Third item with sub-items:
    \begin{itemize}
      \item Sub-item 1
      \item Sub-item 2
    \end{itemize}
\end{itemize}

\subsection{Numbered List}
\begin{enumerate}
  \item First step
  \item Second step
  \item Third step
\end{enumerate}

\section{Citations and References}
When you use the citation search feature, references will be automatically managed. For example, you might cite a paper like this: \cite{example2024}. All citations will be automatically added to your \texttt{cite.bib} file.

\section{Conclusion}
This template demonstrates basic \LaTeX{} features including mathematical formulas, tables, and lists. The bibliography system is now set up to work seamlessly with Bibby AI's citation management. Feel free to modify it for your needs!

\bibliographystyle{plain}
\bibliography{cite}

\end{document}\documentclass{article}
\usepackage[utf8]{inputenc}
\usepackage{booktabs}
\usepackage{amsmath}
\usepackage{natbib}

\title{My Document}
\author{Your Name}
\date{\today}

\begin{document}

\maketitle

\section{Introduction}
This is a sample document demonstrating various \LaTeX{} features including mathematical equations, tables, and formatting. You can cite papers using the citation search feature, and they will be automatically added to your \texttt{cite.bib} file.

\section{Mathematical Expressions}
Here are some example equations:

\subsection{Quadratic Formula}
The quadratic formula is given by:
\[x = \frac{-b \pm \sqrt{b^2 - 4ac}}{2a}\]

\subsection{Matrix Example}
A sample matrix:
\[
\begin{pmatrix}
a_{11} & a_{12} & a_{13} \\
a_{21} & a_{22} & a_{23} \\
a_{31} & a_{32} & a_{33}
\end{pmatrix}
\]

\section{Tables}
Below is a sample table showing data:

\begin{table}[h]
\centering
\begin{tabular}{@{}lrrr@{}}
\toprule
\textbf{Category} & \textbf{Value 1} & \textbf{Value 2} & \textbf{Total} \\
\midrule
Item A & 10.5 & 20.3 & 30.8 \\
Item B & 15.2 & 25.1 & 40.3 \\
Item C & 22.1 & 12.7 & 34.8 \\
\bottomrule
\end{tabular}
\caption{Sample Data Table}
\label{tab:sample}
\end{table}

\section{Lists}
\subsection{Bullet Points}
\begin{itemize}
  \item First item
  \item Second item
  \item Third item with sub-items:
    \begin{itemize}
      \item Sub-item 1
      \item Sub-item 2
    \end{itemize}
\end{itemize}

\subsection{Numbered List}
\begin{enumerate}
  \item First step
  \item Second step
  \item Third step
\end{enumerate}

\section{Citations and References}
When you use the citation search feature, references will be automatically managed. For example, you might cite a paper like this: \cite{example2024}. All citations will be automatically added to your \texttt{cite.bib} file.

\section{Conclusion}
This template demonstrates basic \LaTeX{} features including mathematical formulas, tables, and lists. The bibliography system is now set up to work seamlessly with Bibby AI's citation management. Feel free to modify it for your needs!

\bibliographystyle{plain}
\bibliography{cite}

\end{document}\documentclass{article}
\usepackage[utf8]{inputenc}
\usepackage{booktabs}
\usepackage{amsmath}
\usepackage{natbib}

\title{My Document}
\author{Your Name}
\date{\today}

\begin{document}

\maketitle

\section{Introduction}
This is a sample document demonstrating various \LaTeX{} features including mathematical equations, tables, and formatting. You can cite papers using the citation search feature, and they will be automatically added to your \texttt{cite.bib} file.

\section{Mathematical Expressions}
Here are some example equations:

\subsection{Quadratic Formula}
The quadratic formula is given by:
\[x = \frac{-b \pm \sqrt{b^2 - 4ac}}{2a}\]

\subsection{Matrix Example}
A sample matrix:
\[
\begin{pmatrix}
a_{11} & a_{12} & a_{13} \\
a_{21} & a_{22} & a_{23} \\
a_{31} & a_{32} & a_{33}
\end{pmatrix}
\]

\section{Tables}
Below is a sample table showing data:

\begin{table}[h]
\centering
\begin{tabular}{@{}lrrr@{}}
\toprule
\textbf{Category} & \textbf{Value 1} & \textbf{Value 2} & \textbf{Total} \\
\midrule
Item A & 10.5 & 20.3 & 30.8 \\
Item B & 15.2 & 25.1 & 40.3 \\
Item C & 22.1 & 12.7 & 34.8 \\
\bottomrule
\end{tabular}
\caption{Sample Data Table}
\label{tab:sample}
\end{table}

\section{Lists}
\subsection{Bullet Points}
\begin{itemize}
  \item First item
  \item Second item
  \item Third item with sub-items:
    \begin{itemize}
      \item Sub-item 1
      \item Sub-item 2
    \end{itemize}
\end{itemize}

\subsection{Numbered List}
\begin{enumerate}
  \item First step
  \item Second step
  \item Third step
\end{enumerate}

\section{Citations and References}
When you use the citation search feature, references will be automatically managed. For example, you might cite a paper like this: \cite{example2024}. All citations will be automatically added to your \texttt{cite.bib} file.

\section{Conclusion}
This template demonstrates basic \LaTeX{} features including mathematical formulas, tables, and lists. The bibliography system is now set up to work seamlessly with Bibby AI's citation management. Feel free to modify it for your needs!

\bibliographystyle{plain}
\bibliography{cite}

\end{document}\documentclass{article}
\usepackage[utf8]{inputenc}
\usepackage{booktabs}
\usepackage{amsmath}
\usepackage{natbib}

\title{My Document}
\author{Your Name}
\date{\today}

\begin{document}

\maketitle

\section{Introduction}
This is a sample document demonstrating various \LaTeX{} features including mathematical equations, tables, and formatting. You can cite papers using the citation search feature, and they will be automatically added to your \texttt{cite.bib} file.

\section{Mathematical Expressions}
Here are some example equations:

\subsection{Quadratic Formula}
The quadratic formula is given by:
\[x = \frac{-b \pm \sqrt{b^2 - 4ac}}{2a}\]

\subsection{Matrix Example}
A sample matrix:
\[
\begin{pmatrix}
a_{11} & a_{12} & a_{13} \\
a_{21} & a_{22} & a_{23} \\
a_{31} & a_{32} & a_{33}
\end{pmatrix}
\]

\section{Tables}
Below is a sample table showing data:

\begin{table}[h]
\centering
\begin{tabular}{@{}lrrr@{}}
\toprule
\textbf{Category} & \textbf{Value 1} & \textbf{Value 2} & \textbf{Total} \\
\midrule
Item A & 10.5 & 20.3 & 30.8 \\
Item B & 15.2 & 25.1 & 40.3 \\
Item C & 22.1 & 12.7 & 34.8 \\
\bottomrule
\end{tabular}
\caption{Sample Data Table}
\label{tab:sample}
\end{table}

\section{Lists}
\subsection{Bullet Points}
\begin{itemize}
  \item First item
  \item Second item
  \item Third item with sub-items:
    \begin{itemize}
      \item Sub-item 1
      \item Sub-item 2
    \end{itemize}
\end{itemize}

\subsection{Numbered List}
\begin{enumerate}
  \item First step
  \item Second step
  \item Third step
\end{enumerate}

\section{Citations and References}
When you use the citation search feature, references will be automatically managed. For example, you might cite a paper like this: \cite{example2024}. All citations will be automatically added to your \texttt{cite.bib} file.

\section{Conclusion}
This template demonstrates basic \LaTeX{} features including mathematical formulas, tables, and lists. The bibliography system is now set up to work seamlessly with Bibby AI's citation management. Feel free to modify it for your needs!

\bibliographystyle{plain}
\bibliography{cite}

\end{document}\documentclass{article}
\usepackage[utf8]{inputenc}
\usepackage{booktabs}
\usepackage{amsmath}
\usepackage{natbib}

\title{My Document}
\author{Your Name}
\date{\today}

\begin{document}

\maketitle

\section{Introduction}
This is a sample document demonstrating various \LaTeX{} features including mathematical equations, tables, and formatting. You can cite papers using the citation search feature, and they will be automatically added to your \texttt{cite.bib} file.

\section{Mathematical Expressions}
Here are some example equations:

\subsection{Quadratic Formula}
The quadratic formula is given by:
\[x = \frac{-b \pm \sqrt{b^2 - 4ac}}{2a}\]

\subsection{Matrix Example}
A sample matrix:
\[
\begin{pmatrix}
a_{11} & a_{12} & a_{13} \\
a_{21} & a_{22} & a_{23} \\
a_{31} & a_{32} & a_{33}
\end{pmatrix}
\]

\section{Tables}
Below is a sample table showing data:

\begin{table}[h]
\centering
\begin{tabular}{@{}lrrr@{}}
\toprule
\textbf{Category} & \textbf{Value 1} & \textbf{Value 2} & \textbf{Total} \\
\midrule
Item A & 10.5 & 20.3 & 30.8 \\
Item B & 15.2 & 25.1 & 40.3 \\
Item C & 22.1 & 12.7 & 34.8 \\
\bottomrule
\end{tabular}
\caption{Sample Data Table}
\label{tab:sample}
\end{table}

\section{Lists}
\subsection{Bullet Points}
\begin{itemize}
  \item First item
  \item Second item
  \item Third item with sub-items:
    \begin{itemize}
      \item Sub-item 1
      \item Sub-item 2
    \end{itemize}
\end{itemize}

\subsection{Numbered List}
\begin{enumerate}
  \item First step
  \item Second step
  \item Third step
\end{enumerate}

\section{Citations and References}
When you use the citation search feature, references will be automatically managed. For example, you might cite a paper like this: \cite{example2024}. All citations will be automatically added to your \texttt{cite.bib} file.

\section{Conclusion}
This template demonstrates basic \LaTeX{} features including mathematical formulas, tables, and lists. The bibliography system is now set up to work seamlessly with Bibby AI's citation management. Feel free to modify it for your needs!

\bibliographystyle{plain}
\bibliography{cite}

\end{document}\documentclass{article}
\usepackage[utf8]{inputenc}
\usepackage{booktabs}
\usepackage{amsmath}
\usepackage{natbib}

\title{My Document}
\author{Your Name}
\date{\today}

\begin{document}

\maketitle

\section{Introduction}
This is a sample document demonstrating various \LaTeX{} features including mathematical equations, tables, and formatting. You can cite papers using the citation search feature, and they will be automatically added to your \texttt{cite.bib} file.

\section{Mathematical Expressions}
Here are some example equations:

\subsection{Quadratic Formula}
The quadratic formula is given by:
\[x = \frac{-b \pm \sqrt{b^2 - 4ac}}{2a}\]

\subsection{Matrix Example}
A sample matrix:
\[
\begin{pmatrix}
a_{11} & a_{12} & a_{13} \\
a_{21} & a_{22} & a_{23} \\
a_{31} & a_{32} & a_{33}
\end{pmatrix}
\]

\section{Tables}
Below is a sample table showing data:

\begin{table}[h]
\centering
\begin{tabular}{@{}lrrr@{}}
\toprule
\textbf{Category} & \textbf{Value 1} & \textbf{Value 2} & \textbf{Total} \\
\midrule
Item A & 10.5 & 20.3 & 30.8 \\
Item B & 15.2 & 25.1 & 40.3 \\
Item C & 22.1 & 12.7 & 34.8 \\
\bottomrule
\end{tabular}
\caption{Sample Data Table}
\label{tab:sample}
\end{table}

\section{Lists}
\subsection{Bullet Points}
\begin{itemize}
  \item First item
  \item Second item
  \item Third item with sub-items:
    \begin{itemize}
      \item Sub-item 1
      \item Sub-item 2
    \end{itemize}
\end{itemize}

\subsection{Numbered List}
\begin{enumerate}
  \item First step
  \item Second step
  \item Third step
\end{enumerate}

\section{Citations and References}
When you use the citation search feature, references will be automatically managed. For example, you might cite a paper like this: \cite{example2024}. All citations will be automatically added to your \texttt{cite.bib} file.

\section{Conclusion}
This template demonstrates basic \LaTeX{} features including mathematical formulas, tables, and lists. The bibliography system is now set up to work seamlessly with Bibby AI's citation management. Feel free to modify it for your needs!

\bibliographystyle{plain}
\bibliography{cite}

\end{document}\documentclass{article}
\usepackage[utf8]{inputenc}
\usepackage{booktabs}
\usepackage{amsmath}
\usepackage{natbib}

\title{My Document}
\author{Your Name}
\date{\today}

\begin{document}

\maketitle

\section{Introduction}
This is a sample document demonstrating various \LaTeX{} features including mathematical equations, tables, and formatting. You can cite papers using the citation search feature, and they will be automatically added to your \texttt{cite.bib} file.

\section{Mathematical Expressions}
Here are some example equations:

\subsection{Quadratic Formula}
The quadratic formula is given by:
\[x = \frac{-b \pm \sqrt{b^2 - 4ac}}{2a}\]

\subsection{Matrix Example}
A sample matrix:
\[
\begin{pmatrix}
a_{11} & a_{12} & a_{13} \\
a_{21} & a_{22} & a_{23} \\
a_{31} & a_{32} & a_{33}
\end{pmatrix}
\]

\section{Tables}
Below is a sample table showing data:

\begin{table}[h]
\centering
\begin{tabular}{@{}lrrr@{}}
\toprule
\textbf{Category} & \textbf{Value 1} & \textbf{Value 2} & \textbf{Total} \\
\midrule
Item A & 10.5 & 20.3 & 30.8 \\
Item B & 15.2 & 25.1 & 40.3 \\
Item C & 22.1 & 12.7 & 34.8 \\
\bottomrule
\end{tabular}
\caption{Sample Data Table}
\label{tab:sample}
\end{table}

\section{Lists}
\subsection{Bullet Points}
\begin{itemize}
  \item First item
  \item Second item
  \item Third item with sub-items:
    \begin{itemize}
      \item Sub-item 1
      \item Sub-item 2
    \end{itemize}
\end{itemize}

\subsection{Numbered List}
\begin{enumerate}
  \item First step
  \item Second step
  \item Third step
\end{enumerate}

\section{Citations and References}
When you use the citation search feature, references will be automatically managed. For example, you might cite a paper like this: \cite{example2024}. All citations will be automatically added to your \texttt{cite.bib} file.

\section{Conclusion}
This template demonstrates basic \LaTeX{} features including mathematical formulas, tables, and lists. The bibliography system is now set up to work seamlessly with Bibby AI's citation management. Feel free to modify it for your needs!

\bibliographystyle{plain}
\bibliography{cite}

\end{document}\documentclass{article}
\usepackage[utf8]{inputenc}
\usepackage{booktabs}
\usepackage{amsmath}
\usepackage{natbib}

\title{My Document}
\author{Your Name}
\date{\today}

\begin{document}

\maketitle

\section{Introduction}
This is a sample document demonstrating various \LaTeX{} features including mathematical equations, tables, and formatting. You can cite papers using the citation search feature, and they will be automatically added to your \texttt{cite.bib} file.

\section{Mathematical Expressions}
Here are some example equations:

\subsection{Quadratic Formula}
The quadratic formula is given by:
\[x = \frac{-b \pm \sqrt{b^2 - 4ac}}{2a}\]

\subsection{Matrix Example}
A sample matrix:
\[
\begin{pmatrix}
a_{11} & a_{12} & a_{13} \\
a_{21} & a_{22} & a_{23} \\
a_{31} & a_{32} & a_{33}
\end{pmatrix}
\]

\section{Tables}
Below is a sample table showing data:

\begin{table}[h]
\centering
\begin{tabular}{@{}lrrr@{}}
\toprule
\textbf{Category} & \textbf{Value 1} & \textbf{Value 2} & \textbf{Total} \\
\midrule
Item A & 10.5 & 20.3 & 30.8 \\
Item B & 15.2 & 25.1 & 40.3 \\
Item C & 22.1 & 12.7 & 34.8 \\
\bottomrule
\end{tabular}
\caption{Sample Data Table}
\label{tab:sample}
\end{table}

\section{Lists}
\subsection{Bullet Points}
\begin{itemize}
  \item First item
  \item Second item
  \item Third item with sub-items:
    \begin{itemize}
      \item Sub-item 1
      \item Sub-item 2
    \end{itemize}
\end{itemize}

\subsection{Numbered List}
\begin{enumerate}
  \item First step
  \item Second step
  \item Third step
\end{enumerate}

\section{Citations and References}
When you use the citation search feature, references will be automatically managed. For example, you might cite a paper like this: \cite{example2024}. All citations will be automatically added to your \texttt{cite.bib} file.

\section{Conclusion}
This template demonstrates basic \LaTeX{} features including mathematical formulas, tables, and lists. The bibliography system is now set up to work seamlessly with Bibby AI's citation management. Feel free to modify it for your needs!

\bibliographystyle{plain}
\bibliography{cite}

\end{document}\documentclass{article}
\usepackage[utf8]{inputenc}
\usepackage{booktabs}
\usepackage{amsmath}
\usepackage{natbib}

\title{My Document}
\author{Your Name}
\date{\today}

\begin{document}

\maketitle

\section{Introduction}
This is a sample document demonstrating various \LaTeX{} features including mathematical equations, tables, and formatting. You can cite papers using the citation search feature, and they will be automatically added to your \texttt{cite.bib} file.

\section{Mathematical Expressions}
Here are some example equations:

\subsection{Quadratic Formula}
The quadratic formula is given by:
\[x = \frac{-b \pm \sqrt{b^2 - 4ac}}{2a}\]

\subsection{Matrix Example}
A sample matrix:
\[
\begin{pmatrix}
a_{11} & a_{12} & a_{13} \\
a_{21} & a_{22} & a_{23} \\
a_{31} & a_{32} & a_{33}
\end{pmatrix}
\]

\section{Tables}
Below is a sample table showing data:

\begin{table}[h]
\centering
\begin{tabular}{@{}lrrr@{}}
\toprule
\textbf{Category} & \textbf{Value 1} & \textbf{Value 2} & \textbf{Total} \\
\midrule
Item A & 10.5 & 20.3 & 30.8 \\
Item B & 15.2 & 25.1 & 40.3 \\
Item C & 22.1 & 12.7 & 34.8 \\
\bottomrule
\end{tabular}
\caption{Sample Data Table}
\label{tab:sample}
\end{table}

\section{Lists}
\subsection{Bullet Points}
\begin{itemize}
  \item First item
  \item Second item
  \item Third item with sub-items:
    \begin{itemize}
      \item Sub-item 1
      \item Sub-item 2
    \end{itemize}
\end{itemize}

\subsection{Numbered List}
\begin{enumerate}
  \item First step
  \item Second step
  \item Third step
\end{enumerate}

\section{Citations and References}
When you use the citation search feature, references will be automatically managed. For example, you might cite a paper like this: \cite{example2024}. All citations will be automatically added to your \texttt{cite.bib} file.

\section{Conclusion}
This template demonstrates basic \LaTeX{} features including mathematical formulas, tables, and lists. The bibliography system is now set up to work seamlessly with Bibby AI's citation management. Feel free to modify it for your needs!

\bibliographystyle{plain}
\bibliography{cite}

\end{document}\documentclass{article}
\usepackage[utf8]{inputenc}
\usepackage{booktabs}
\usepackage{amsmath}
\usepackage{natbib}

\title{My Document}
\author{Your Name}
\date{\today}

\begin{document}

\maketitle

\section{Introduction}
This is a sample document demonstrating various \LaTeX{} features including mathematical equations, tables, and formatting. You can cite papers using the citation search feature, and they will be automatically added to your \texttt{cite.bib} file.

\section{Mathematical Expressions}
Here are some example equations:

\subsection{Quadratic Formula}
The quadratic formula is given by:
\[x = \frac{-b \pm \sqrt{b^2 - 4ac}}{2a}\]

\subsection{Matrix Example}
A sample matrix:
\[
\begin{pmatrix}
a_{11} & a_{12} & a_{13} \\
a_{21} & a_{22} & a_{23} \\
a_{31} & a_{32} & a_{33}
\end{pmatrix}
\]

\section{Tables}
Below is a sample table showing data:

\begin{table}[h]
\centering
\begin{tabular}{@{}lrrr@{}}
\toprule
\textbf{Category} & \textbf{Value 1} & \textbf{Value 2} & \textbf{Total} \\
\midrule
Item A & 10.5 & 20.3 & 30.8 \\
Item B & 15.2 & 25.1 & 40.3 \\
Item C & 22.1 & 12.7 & 34.8 \\
\bottomrule
\end{tabular}
\caption{Sample Data Table}
\label{tab:sample}
\end{table}

\section{Lists}
\subsection{Bullet Points}
\begin{itemize}
  \item First item
  \item Second item
  \item Third item with sub-items:
    \begin{itemize}
      \item Sub-item 1
      \item Sub-item 2
    \end{itemize}
\end{itemize}

\subsection{Numbered List}
\begin{enumerate}
  \item First step
  \item Second step
  \item Third step
\end{enumerate}

\section{Citations and References}
When you use the citation search feature, references will be automatically managed. For example, you might cite a paper like this: \cite{example2024}. All citations will be automatically added to your \texttt{cite.bib} file.

\section{Conclusion}
This template demonstrates basic \LaTeX{} features including mathematical formulas, tables, and lists. The bibliography system is now set up to work seamlessly with Bibby AI's citation management. Feel free to modify it for your needs!

\bibliographystyle{plain}
\bibliography{cite}

\end{document}\documentclass{article}
\usepackage[utf8]{inputenc}
\usepackage{booktabs}
\usepackage{amsmath}
\usepackage{natbib}

\title{My Document}
\author{Your Name}
\date{\today}

\begin{document}

\maketitle

\section{Introduction}
This is a sample document demonstrating various \LaTeX{} features including mathematical equations, tables, and formatting. You can cite papers using the citation search feature, and they will be automatically added to your \texttt{cite.bib} file.

\section{Mathematical Expressions}
Here are some example equations:

\subsection{Quadratic Formula}
The quadratic formula is given by:
\[x = \frac{-b \pm \sqrt{b^2 - 4ac}}{2a}\]

\subsection{Matrix Example}
A sample matrix:
\[
\begin{pmatrix}
a_{11} & a_{12} & a_{13} \\
a_{21} & a_{22} & a_{23} \\
a_{31} & a_{32} & a_{33}
\end{pmatrix}
\]

\section{Tables}
Below is a sample table showing data:

\begin{table}[h]
\centering
\begin{tabular}{@{}lrrr@{}}
\toprule
\textbf{Category} & \textbf{Value 1} & \textbf{Value 2} & \textbf{Total} \\
\midrule
Item A & 10.5 & 20.3 & 30.8 \\
Item B & 15.2 & 25.1 & 40.3 \\
Item C & 22.1 & 12.7 & 34.8 \\
\bottomrule
\end{tabular}
\caption{Sample Data Table}
\label{tab:sample}
\end{table}

\section{Lists}
\subsection{Bullet Points}
\begin{itemize}
  \item First item
  \item Second item
  \item Third item with sub-items:
    \begin{itemize}
      \item Sub-item 1
      \item Sub-item 2
    \end{itemize}
\end{itemize}

\subsection{Numbered List}
\begin{enumerate}
  \item First step
  \item Second step
  \item Third step
\end{enumerate}

\section{Citations and References}
When you use the citation search feature, references will be automatically managed. For example, you might cite a paper like this: \cite{example2024}. All citations will be automatically added to your \texttt{cite.bib} file.

\section{Conclusion}
This template demonstrates basic \LaTeX{} features including mathematical formulas, tables, and lists. The bibliography system is now set up to work seamlessly with Bibby AI's citation management. Feel free to modify it for your needs!

\bibliographystyle{plain}
\bibliography{cite}

\end{document}\documentclass{article}
\usepackage[utf8]{inputenc}
\usepackage{booktabs}
\usepackage{amsmath}
\usepackage{natbib}

\title{My Document}
\author{Your Name}
\date{\today}

\begin{document}

\maketitle

\section{Introduction}
This is a sample document demonstrating various \LaTeX{} features including mathematical equations, tables, and formatting. You can cite papers using the citation search feature, and they will be automatically added to your \texttt{cite.bib} file.

\section{Mathematical Expressions}
Here are some example equations:

\subsection{Quadratic Formula}
The quadratic formula is given by:
\[x = \frac{-b \pm \sqrt{b^2 - 4ac}}{2a}\]

\subsection{Matrix Example}
A sample matrix:
\[
\begin{pmatrix}
a_{11} & a_{12} & a_{13} \\
a_{21} & a_{22} & a_{23} \\
a_{31} & a_{32} & a_{33}
\end{pmatrix}
\]

\section{Tables}
Below is a sample table showing data:

\begin{table}[h]
\centering
\begin{tabular}{@{}lrrr@{}}
\toprule
\textbf{Category} & \textbf{Value 1} & \textbf{Value 2} & \textbf{Total} \\
\midrule
Item A & 10.5 & 20.3 & 30.8 \\
Item B & 15.2 & 25.1 & 40.3 \\
Item C & 22.1 & 12.7 & 34.8 \\
\bottomrule
\end{tabular}
\caption{Sample Data Table}
\label{tab:sample}
\end{table}

\section{Lists}
\subsection{Bullet Points}
\begin{itemize}
  \item First item
  \item Second item
  \item Third item with sub-items:
    \begin{itemize}
      \item Sub-item 1
      \item Sub-item 2
    \end{itemize}
\end{itemize}

\subsection{Numbered List}
\begin{enumerate}
  \item First step
  \item Second step
  \item Third step
\end{enumerate}

\section{Citations and References}
When you use the citation search feature, references will be automatically managed. For example, you might cite a paper like this: \cite{example2024}. All citations will be automatically added to your \texttt{cite.bib} file.

\section{Conclusion}
This template demonstrates basic \LaTeX{} features including mathematical formulas, tables, and lists. The bibliography system is now set up to work seamlessly with Bibby AI's citation management. Feel free to modify it for your needs!

\bibliographystyle{plain}
\bibliography{cite}

\end{document}\documentclass{article}
\usepackage[utf8]{inputenc}
\usepackage{booktabs}
\usepackage{amsmath}
\usepackage{natbib}

\title{My Document}
\author{Your Name}
\date{\today}

\begin{document}

\maketitle

\section{Introduction}
This is a sample document demonstrating various \LaTeX{} features including mathematical equations, tables, and formatting. You can cite papers using the citation search feature, and they will be automatically added to your \texttt{cite.bib} file.

\section{Mathematical Expressions}
Here are some example equations:

\subsection{Quadratic Formula}
The quadratic formula is given by:
\[x = \frac{-b \pm \sqrt{b^2 - 4ac}}{2a}\]

\subsection{Matrix Example}
A sample matrix:
\[
\begin{pmatrix}
a_{11} & a_{12} & a_{13} \\
a_{21} & a_{22} & a_{23} \\
a_{31} & a_{32} & a_{33}
\end{pmatrix}
\]

\section{Tables}
Below is a sample table showing data:

\begin{table}[h]
\centering
\begin{tabular}{@{}lrrr@{}}
\toprule
\textbf{Category} & \textbf{Value 1} & \textbf{Value 2} & \textbf{Total} \\
\midrule
Item A & 10.5 & 20.3 & 30.8 \\
Item B & 15.2 & 25.1 & 40.3 \\
Item C & 22.1 & 12.7 & 34.8 \\
\bottomrule
\end{tabular}
\caption{Sample Data Table}
\label{tab:sample}
\end{table}

\section{Lists}
\subsection{Bullet Points}
\begin{itemize}
  \item First item
  \item Second item
  \item Third item with sub-items:
    \begin{itemize}
      \item Sub-item 1
      \item Sub-item 2
    \end{itemize}
\end{itemize}

\subsection{Numbered List}
\begin{enumerate}
  \item First step
  \item Second step
  \item Third step
\end{enumerate}

\section{Citations and References}
When you use the citation search feature, references will be automatically managed. For example, you might cite a paper like this: \cite{example2024}. All citations will be automatically added to your \texttt{cite.bib} file.

\section{Conclusion}
This template demonstrates basic \LaTeX{} features including mathematical formulas, tables, and lists. The bibliography system is now set up to work seamlessly with Bibby AI's citation management. Feel free to modify it for your needs!

\bibliographystyle{plain}
\bibliography{cite}

\end{document}\documentclass{article}
\usepackage[utf8]{inputenc}
\usepackage{booktabs}
\usepackage{amsmath}
\usepackage{natbib}

\title{My Document}
\author{Your Name}
\date{\today}

\begin{document}

\maketitle

\section{Introduction}
This is a sample document demonstrating various \LaTeX{} features including mathematical equations, tables, and formatting. You can cite papers using the citation search feature, and they will be automatically added to your \texttt{cite.bib} file.

\section{Mathematical Expressions}
Here are some example equations:

\subsection{Quadratic Formula}
The quadratic formula is given by:
\[x = \frac{-b \pm \sqrt{b^2 - 4ac}}{2a}\]

\subsection{Matrix Example}
A sample matrix:
\[
\begin{pmatrix}
a_{11} & a_{12} & a_{13} \\
a_{21} & a_{22} & a_{23} \\
a_{31} & a_{32} & a_{33}
\end{pmatrix}
\]

\section{Tables}
Below is a sample table showing data:

\begin{table}[h]
\centering
\begin{tabular}{@{}lrrr@{}}
\toprule
\textbf{Category} & \textbf{Value 1} & \textbf{Value 2} & \textbf{Total} \\
\midrule
Item A & 10.5 & 20.3 & 30.8 \\
Item B & 15.2 & 25.1 & 40.3 \\
Item C & 22.1 & 12.7 & 34.8 \\
\bottomrule
\end{tabular}
\caption{Sample Data Table}
\label{tab:sample}
\end{table}

\section{Lists}
\subsection{Bullet Points}
\begin{itemize}
  \item First item
  \item Second item
  \item Third item with sub-items:
    \begin{itemize}
      \item Sub-item 1
      \item Sub-item 2
    \end{itemize}
\end{itemize}

\subsection{Numbered List}
\begin{enumerate}
  \item First step
  \item Second step
  \item Third step
\end{enumerate}

\section{Citations and References}
When you use the citation search feature, references will be automatically managed. For example, you might cite a paper like this: \cite{example2024}. All citations will be automatically added to your \texttt{cite.bib} file.

\section{Conclusion}
This template demonstrates basic \LaTeX{} features including mathematical formulas, tables, and lists. The bibliography system is now set up to work seamlessly with Bibby AI's citation management. Feel free to modify it for your needs!

\bibliographystyle{plain}
\bibliography{cite}

\end{document}\documentclass{article}
\usepackage[utf8]{inputenc}
\usepackage{booktabs}
\usepackage{amsmath}
\usepackage{natbib}

\title{My Document}
\author{Your Name}
\date{\today}

\begin{document}

\maketitle

\section{Introduction}
This is a sample document demonstrating various \LaTeX{} features including mathematical equations, tables, and formatting. You can cite papers using the citation search feature, and they will be automatically added to your \texttt{cite.bib} file.

\section{Mathematical Expressions}
Here are some example equations:

\subsection{Quadratic Formula}
The quadratic formula is given by:
\[x = \frac{-b \pm \sqrt{b^2 - 4ac}}{2a}\]

\subsection{Matrix Example}
A sample matrix:
\[
\begin{pmatrix}
a_{11} & a_{12} & a_{13} \\
a_{21} & a_{22} & a_{23} \\
a_{31} & a_{32} & a_{33}
\end{pmatrix}
\]

\section{Tables}
Below is a sample table showing data:

\begin{table}[h]
\centering
\begin{tabular}{@{}lrrr@{}}
\toprule
\textbf{Category} & \textbf{Value 1} & \textbf{Value 2} & \textbf{Total} \\
\midrule
Item A & 10.5 & 20.3 & 30.8 \\
Item B & 15.2 & 25.1 & 40.3 \\
Item C & 22.1 & 12.7 & 34.8 \\
\bottomrule
\end{tabular}
\caption{Sample Data Table}
\label{tab:sample}
\end{table}

\section{Lists}
\subsection{Bullet Points}
\begin{itemize}
  \item First item
  \item Second item
  \item Third item with sub-items:
    \begin{itemize}
      \item Sub-item 1
      \item Sub-item 2
    \end{itemize}
\end{itemize}

\subsection{Numbered List}
\begin{enumerate}
  \item First step
  \item Second step
  \item Third step
\end{enumerate}

\section{Citations and References}
When you use the citation search feature, references will be automatically managed. For example, you might cite a paper like this: \cite{example2024}. All citations will be automatically added to your \texttt{cite.bib} file.

\section{Conclusion}
This template demonstrates basic \LaTeX{} features including mathematical formulas, tables, and lists. The bibliography system is now set up to work seamlessly with Bibby AI's citation management. Feel free to modify it for your needs!

\bibliographystyle{plain}
\bibliography{cite}

\end{document}\documentclass{article}
\usepackage[utf8]{inputenc}
\usepackage{booktabs}
\usepackage{amsmath}
\usepackage{natbib}

\title{My Document}
\author{Your Name}
\date{\today}

\begin{document}

\maketitle

\section{Introduction}
This is a sample document demonstrating various \LaTeX{} features including mathematical equations, tables, and formatting. You can cite papers using the citation search feature, and they will be automatically added to your \texttt{cite.bib} file.

\section{Mathematical Expressions}
Here are some example equations:

\subsection{Quadratic Formula}
The quadratic formula is given by:
\[x = \frac{-b \pm \sqrt{b^2 - 4ac}}{2a}\]

\subsection{Matrix Example}
A sample matrix:
\[
\begin{pmatrix}
a_{11} & a_{12} & a_{13} \\
a_{21} & a_{22} & a_{23} \\
a_{31} & a_{32} & a_{33}
\end{pmatrix}
\]

\section{Tables}
Below is a sample table showing data:

\begin{table}[h]
\centering
\begin{tabular}{@{}lrrr@{}}
\toprule
\textbf{Category} & \textbf{Value 1} & \textbf{Value 2} & \textbf{Total} \\
\midrule
Item A & 10.5 & 20.3 & 30.8 \\
Item B & 15.2 & 25.1 & 40.3 \\
Item C & 22.1 & 12.7 & 34.8 \\
\bottomrule
\end{tabular}
\caption{Sample Data Table}
\label{tab:sample}
\end{table}

\section{Lists}
\subsection{Bullet Points}
\begin{itemize}
  \item First item
  \item Second item
  \item Third item with sub-items:
    \begin{itemize}
      \item Sub-item 1
      \item Sub-item 2
    \end{itemize}
\end{itemize}

\subsection{Numbered List}
\begin{enumerate}
  \item First step
  \item Second step
  \item Third step
\end{enumerate}

\section{Citations and References}
When you use the citation search feature, references will be automatically managed. For example, you might cite a paper like this: \cite{example2024}. All citations will be automatically added to your \texttt{cite.bib} file.

\section{Conclusion}
This template demonstrates basic \LaTeX{} features including mathematical formulas, tables, and lists. The bibliography system is now set up to work seamlessly with Bibby AI's citation management. Feel free to modify it for your needs!

\bibliographystyle{plain}
\bibliography{cite}

\end{document}\documentclass{article}
\usepackage[utf8]{inputenc}
\usepackage{booktabs}
\usepackage{amsmath}
\usepackage{natbib}

\title{My Document}
\author{Your Name}
\date{\today}

\begin{document}

\maketitle

\section{Introduction}
This is a sample document demonstrating various \LaTeX{} features including mathematical equations, tables, and formatting. You can cite papers using the citation search feature, and they will be automatically added to your \texttt{cite.bib} file.

\section{Mathematical Expressions}
Here are some example equations:

\subsection{Quadratic Formula}
The quadratic formula is given by:
\[x = \frac{-b \pm \sqrt{b^2 - 4ac}}{2a}\]

\subsection{Matrix Example}
A sample matrix:
\[
\begin{pmatrix}
a_{11} & a_{12} & a_{13} \\
a_{21} & a_{22} & a_{23} \\
a_{31} & a_{32} & a_{33}
\end{pmatrix}
\]

\section{Tables}
Below is a sample table showing data:

\begin{table}[h]
\centering
\begin{tabular}{@{}lrrr@{}}
\toprule
\textbf{Category} & \textbf{Value 1} & \textbf{Value 2} & \textbf{Total} \\
\midrule
Item A & 10.5 & 20.3 & 30.8 \\
Item B & 15.2 & 25.1 & 40.3 \\
Item C & 22.1 & 12.7 & 34.8 \\
\bottomrule
\end{tabular}
\caption{Sample Data Table}
\label{tab:sample}
\end{table}

\section{Lists}
\subsection{Bullet Points}
\begin{itemize}
  \item First item
  \item Second item
  \item Third item with sub-items:
    \begin{itemize}
      \item Sub-item 1
      \item Sub-item 2
    \end{itemize}
\end{itemize}

\subsection{Numbered List}
\begin{enumerate}
  \item First step
  \item Second step
  \item Third step
\end{enumerate}

\section{Citations and References}
When you use the citation search feature, references will be automatically managed. For example, you might cite a paper like this: \cite{example2024}. All citations will be automatically added to your \texttt{cite.bib} file.

\section{Conclusion}
This template demonstrates basic \LaTeX{} features including mathematical formulas, tables, and lists. The bibliography system is now set up to work seamlessly with Bibby AI's citation management. Feel free to modify it for your needs!

\bibliographystyle{plain}
\bibliography{cite}

\end{document}\documentclass{article}
\usepackage[utf8]{inputenc}
\usepackage{booktabs}
\usepackage{amsmath}
\usepackage{natbib}

\title{My Document}
\author{Your Name}
\date{\today}

\begin{document}

\maketitle

\section{Introduction}
This is a sample document demonstrating various \LaTeX{} features including mathematical equations, tables, and formatting. You can cite papers using the citation search feature, and they will be automatically added to your \texttt{cite.bib} file.

\section{Mathematical Expressions}
Here are some example equations:

\subsection{Quadratic Formula}
The quadratic formula is given by:
\[x = \frac{-b \pm \sqrt{b^2 - 4ac}}{2a}\]

\subsection{Matrix Example}
A sample matrix:
\[
\begin{pmatrix}
a_{11} & a_{12} & a_{13} \\
a_{21} & a_{22} & a_{23} \\
a_{31} & a_{32} & a_{33}
\end{pmatrix}
\]

\section{Tables}
Below is a sample table showing data:

\begin{table}[h]
\centering
\begin{tabular}{@{}lrrr@{}}
\toprule
\textbf{Category} & \textbf{Value 1} & \textbf{Value 2} & \textbf{Total} \\
\midrule
Item A & 10.5 & 20.3 & 30.8 \\
Item B & 15.2 & 25.1 & 40.3 \\
Item C & 22.1 & 12.7 & 34.8 \\
\bottomrule
\end{tabular}
\caption{Sample Data Table}
\label{tab:sample}
\end{table}

\section{Lists}
\subsection{Bullet Points}
\begin{itemize}
  \item First item
  \item Second item
  \item Third item with sub-items:
    \begin{itemize}
      \item Sub-item 1
      \item Sub-item 2
    \end{itemize}
\end{itemize}

\subsection{Numbered List}
\begin{enumerate}
  \item First step
  \item Second step
  \item Third step
\end{enumerate}

\section{Citations and References}
When you use the citation search feature, references will be automatically managed. For example, you might cite a paper like this: \cite{example2024}. All citations will be automatically added to your \texttt{cite.bib} file.

\section{Conclusion}
This template demonstrates basic \LaTeX{} features including mathematical formulas, tables, and lists. The bibliography system is now set up to work seamlessly with Bibby AI's citation management. Feel free to modify it for your needs!

\bibliographystyle{plain}
\bibliography{cite}

\end{document}\documentclass{article}
\usepackage[utf8]{inputenc}
\usepackage{booktabs}
\usepackage{amsmath}
\usepackage{natbib}

\title{My Document}
\author{Your Name}
\date{\today}

\begin{document}

\maketitle

\section{Introduction}
This is a sample document demonstrating various \LaTeX{} features including mathematical equations, tables, and formatting. You can cite papers using the citation search feature, and they will be automatically added to your \texttt{cite.bib} file.

\section{Mathematical Expressions}
Here are some example equations:

\subsection{Quadratic Formula}
The quadratic formula is given by:
\[x = \frac{-b \pm \sqrt{b^2 - 4ac}}{2a}\]

\subsection{Matrix Example}
A sample matrix:
\[
\begin{pmatrix}
a_{11} & a_{12} & a_{13} \\
a_{21} & a_{22} & a_{23} \\
a_{31} & a_{32} & a_{33}
\end{pmatrix}
\]

\section{Tables}
Below is a sample table showing data:

\begin{table}[h]
\centering
\begin{tabular}{@{}lrrr@{}}
\toprule
\textbf{Category} & \textbf{Value 1} & \textbf{Value 2} & \textbf{Total} \\
\midrule
Item A & 10.5 & 20.3 & 30.8 \\
Item B & 15.2 & 25.1 & 40.3 \\
Item C & 22.1 & 12.7 & 34.8 \\
\bottomrule
\end{tabular}
\caption{Sample Data Table}
\label{tab:sample}
\end{table}

\section{Lists}
\subsection{Bullet Points}
\begin{itemize}
  \item First item
  \item Second item
  \item Third item with sub-items:
    \begin{itemize}
      \item Sub-item 1
      \item Sub-item 2
    \end{itemize}
\end{itemize}

\subsection{Numbered List}
\begin{enumerate}
  \item First step
  \item Second step
  \item Third step
\end{enumerate}

\section{Citations and References}
When you use the citation search feature, references will be automatically managed. For example, you might cite a paper like this: \cite{example2024}. All citations will be automatically added to your \texttt{cite.bib} file.

\section{Conclusion}
This template demonstrates basic \LaTeX{} features including mathematical formulas, tables, and lists. The bibliography system is now set up to work seamlessly with Bibby AI's citation management. Feel free to modify it for your needs!

\bibliographystyle{plain}
\bibliography{cite}

\end{document}\documentclass{article}
\usepackage[utf8]{inputenc}
\usepackage{booktabs}
\usepackage{amsmath}
\usepackage{natbib}

\title{My Document}
\author{Your Name}
\date{\today}

\begin{document}

\maketitle

\section{Introduction}
This is a sample document demonstrating various \LaTeX{} features including mathematical equations, tables, and formatting. You can cite papers using the citation search feature, and they will be automatically added to your \texttt{cite.bib} file.

\section{Mathematical Expressions}
Here are some example equations:

\subsection{Quadratic Formula}
The quadratic formula is given by:
\[x = \frac{-b \pm \sqrt{b^2 - 4ac}}{2a}\]

\subsection{Matrix Example}
A sample matrix:
\[
\begin{pmatrix}
a_{11} & a_{12} & a_{13} \\
a_{21} & a_{22} & a_{23} \\
a_{31} & a_{32} & a_{33}
\end{pmatrix}
\]

\section{Tables}
Below is a sample table showing data:

\begin{table}[h]
\centering
\begin{tabular}{@{}lrrr@{}}
\toprule
\textbf{Category} & \textbf{Value 1} & \textbf{Value 2} & \textbf{Total} \\
\midrule
Item A & 10.5 & 20.3 & 30.8 \\
Item B & 15.2 & 25.1 & 40.3 \\
Item C & 22.1 & 12.7 & 34.8 \\
\bottomrule
\end{tabular}
\caption{Sample Data Table}
\label{tab:sample}
\end{table}

\section{Lists}
\subsection{Bullet Points}
\begin{itemize}
  \item First item
  \item Second item
  \item Third item with sub-items:
    \begin{itemize}
      \item Sub-item 1
      \item Sub-item 2
    \end{itemize}
\end{itemize}

\subsection{Numbered List}
\begin{enumerate}
  \item First step
  \item Second step
  \item Third step
\end{enumerate}

\section{Citations and References}
When you use the citation search feature, references will be automatically managed. For example, you might cite a paper like this: \cite{example2024}. All citations will be automatically added to your \texttt{cite.bib} file.

\section{Conclusion}
This template demonstrates basic \LaTeX{} features including mathematical formulas, tables, and lists. The bibliography system is now set up to work seamlessly with Bibby AI's citation management. Feel free to modify it for your needs!

\bibliographystyle{plain}
\bibliography{cite}

\end{document}\documentclass{article}
\usepackage[utf8]{inputenc}
\usepackage{booktabs}
\usepackage{amsmath}
\usepackage{natbib}

\title{My Document}
\author{Your Name}
\date{\today}

\begin{document}

\maketitle

\section{Introduction}
This is a sample document demonstrating various \LaTeX{} features including mathematical equations, tables, and formatting. You can cite papers using the citation search feature, and they will be automatically added to your \texttt{cite.bib} file.

\section{Mathematical Expressions}
Here are some example equations:

\subsection{Quadratic Formula}
The quadratic formula is given by:
\[x = \frac{-b \pm \sqrt{b^2 - 4ac}}{2a}\]

\subsection{Matrix Example}
A sample matrix:
\[
\begin{pmatrix}
a_{11} & a_{12} & a_{13} \\
a_{21} & a_{22} & a_{23} \\
a_{31} & a_{32} & a_{33}
\end{pmatrix}
\]

\section{Tables}
Below is a sample table showing data:

\begin{table}[h]
\centering
\begin{tabular}{@{}lrrr@{}}
\toprule
\textbf{Category} & \textbf{Value 1} & \textbf{Value 2} & \textbf{Total} \\
\midrule
Item A & 10.5 & 20.3 & 30.8 \\
Item B & 15.2 & 25.1 & 40.3 \\
Item C & 22.1 & 12.7 & 34.8 \\
\bottomrule
\end{tabular}
\caption{Sample Data Table}
\label{tab:sample}
\end{table}

\section{Lists}
\subsection{Bullet Points}
\begin{itemize}
  \item First item
  \item Second item
  \item Third item with sub-items:
    \begin{itemize}
      \item Sub-item 1
      \item Sub-item 2
    \end{itemize}
\end{itemize}

\subsection{Numbered List}
\begin{enumerate}
  \item First step
  \item Second step
  \item Third step
\end{enumerate}

\section{Citations and References}
When you use the citation search feature, references will be automatically managed. For example, you might cite a paper like this: \cite{example2024}. All citations will be automatically added to your \texttt{cite.bib} file.

\section{Conclusion}
This template demonstrates basic \LaTeX{} features including mathematical formulas, tables, and lists. The bibliography system is now set up to work seamlessly with Bibby AI's citation management. Feel free to modify it for your needs!

\bibliographystyle{plain}
\bibliography{cite}

\end{document}\documentclass{article}
\usepackage[utf8]{inputenc}
\usepackage{booktabs}
\usepackage{amsmath}
\usepackage{natbib}

\title{My Document}
\author{Your Name}
\date{\today}

\begin{document}

\maketitle

\section{Introduction}
This is a sample document demonstrating various \LaTeX{} features including mathematical equations, tables, and formatting. You can cite papers using the citation search feature, and they will be automatically added to your \texttt{cite.bib} file.

\section{Mathematical Expressions}
Here are some example equations:

\subsection{Quadratic Formula}
The quadratic formula is given by:
\[x = \frac{-b \pm \sqrt{b^2 - 4ac}}{2a}\]

\subsection{Matrix Example}
A sample matrix:
\[
\begin{pmatrix}
a_{11} & a_{12} & a_{13} \\
a_{21} & a_{22} & a_{23} \\
a_{31} & a_{32} & a_{33}
\end{pmatrix}
\]

\section{Tables}
Below is a sample table showing data:

\begin{table}[h]
\centering
\begin{tabular}{@{}lrrr@{}}
\toprule
\textbf{Category} & \textbf{Value 1} & \textbf{Value 2} & \textbf{Total} \\
\midrule
Item A & 10.5 & 20.3 & 30.8 \\
Item B & 15.2 & 25.1 & 40.3 \\
Item C & 22.1 & 12.7 & 34.8 \\
\bottomrule
\end{tabular}
\caption{Sample Data Table}
\label{tab:sample}
\end{table}

\section{Lists}
\subsection{Bullet Points}
\begin{itemize}
  \item First item
  \item Second item
  \item Third item with sub-items:
    \begin{itemize}
      \item Sub-item 1
      \item Sub-item 2
    \end{itemize}
\end{itemize}

\subsection{Numbered List}
\begin{enumerate}
  \item First step
  \item Second step
  \item Third step
\end{enumerate}

\section{Citations and References}
When you use the citation search feature, references will be automatically managed. For example, you might cite a paper like this: \cite{example2024}. All citations will be automatically added to your \texttt{cite.bib} file.

\section{Conclusion}
This template demonstrates basic \LaTeX{} features including mathematical formulas, tables, and lists. The bibliography system is now set up to work seamlessly with Bibby AI's citation management. Feel free to modify it for your needs!

\bibliographystyle{plain}
\bibliography{cite}

\end{document}\documentclass{article}
\usepackage[utf8]{inputenc}
\usepackage{booktabs}
\usepackage{amsmath}
\usepackage{natbib}

\title{My Document}
\author{Your Name}
\date{\today}

\begin{document}

\maketitle

\section{Introduction}
This is a sample document demonstrating various \LaTeX{} features including mathematical equations, tables, and formatting. You can cite papers using the citation search feature, and they will be automatically added to your \texttt{cite.bib} file.

\section{Mathematical Expressions}
Here are some example equations:

\subsection{Quadratic Formula}
The quadratic formula is given by:
\[x = \frac{-b \pm \sqrt{b^2 - 4ac}}{2a}\]

\subsection{Matrix Example}
A sample matrix:
\[
\begin{pmatrix}
a_{11} & a_{12} & a_{13} \\
a_{21} & a_{22} & a_{23} \\
a_{31} & a_{32} & a_{33}
\end{pmatrix}
\]

\section{Tables}
Below is a sample table showing data:

\begin{table}[h]
\centering
\begin{tabular}{@{}lrrr@{}}
\toprule
\textbf{Category} & \textbf{Value 1} & \textbf{Value 2} & \textbf{Total} \\
\midrule
Item A & 10.5 & 20.3 & 30.8 \\
Item B & 15.2 & 25.1 & 40.3 \\
Item C & 22.1 & 12.7 & 34.8 \\
\bottomrule
\end{tabular}
\caption{Sample Data Table}
\label{tab:sample}
\end{table}

\section{Lists}
\subsection{Bullet Points}
\begin{itemize}
  \item First item
  \item Second item
  \item Third item with sub-items:
    \begin{itemize}
      \item Sub-item 1
      \item Sub-item 2
    \end{itemize}
\end{itemize}

\subsection{Numbered List}
\begin{enumerate}
  \item First step
  \item Second step
  \item Third step
\end{enumerate}

\section{Citations and References}
When you use the citation search feature, references will be automatically managed. For example, you might cite a paper like this: \cite{example2024}. All citations will be automatically added to your \texttt{cite.bib} file.

\section{Conclusion}
This template demonstrates basic \LaTeX{} features including mathematical formulas, tables, and lists. The bibliography system is now set up to work seamlessly with Bibby AI's citation management. Feel free to modify it for your needs!

\bibliographystyle{plain}
\bibliography{cite}

\end{document}\documentclass{article}
\usepackage[utf8]{inputenc}
\usepackage{booktabs}
\usepackage{amsmath}
\usepackage{natbib}

\title{My Document}
\author{Your Name}
\date{\today}

\begin{document}

\maketitle

\section{Introduction}
This is a sample document demonstrating various \LaTeX{} features including mathematical equations, tables, and formatting. You can cite papers using the citation search feature, and they will be automatically added to your \texttt{cite.bib} file.

\section{Mathematical Expressions}
Here are some example equations:

\subsection{Quadratic Formula}
The quadratic formula is given by:
\[x = \frac{-b \pm \sqrt{b^2 - 4ac}}{2a}\]

\subsection{Matrix Example}
A sample matrix:
\[
\begin{pmatrix}
a_{11} & a_{12} & a_{13} \\
a_{21} & a_{22} & a_{23} \\
a_{31} & a_{32} & a_{33}
\end{pmatrix}
\]

\section{Tables}
Below is a sample table showing data:

\begin{table}[h]
\centering
\begin{tabular}{@{}lrrr@{}}
\toprule
\textbf{Category} & \textbf{Value 1} & \textbf{Value 2} & \textbf{Total} \\
\midrule
Item A & 10.5 & 20.3 & 30.8 \\
Item B & 15.2 & 25.1 & 40.3 \\
Item C & 22.1 & 12.7 & 34.8 \\
\bottomrule
\end{tabular}
\caption{Sample Data Table}
\label{tab:sample}
\end{table}

\section{Lists}
\subsection{Bullet Points}
\begin{itemize}
  \item First item
  \item Second item
  \item Third item with sub-items:
    \begin{itemize}
      \item Sub-item 1
      \item Sub-item 2
    \end{itemize}
\end{itemize}

\subsection{Numbered List}
\begin{enumerate}
  \item First step
  \item Second step
  \item Third step
\end{enumerate}

\section{Citations and References}
When you use the citation search feature, references will be automatically managed. For example, you might cite a paper like this: \cite{example2024}. All citations will be automatically added to your \texttt{cite.bib} file.

\section{Conclusion}
This template demonstrates basic \LaTeX{} features including mathematical formulas, tables, and lists. The bibliography system is now set up to work seamlessly with Bibby AI's citation management. Feel free to modify it for your needs!

\bibliographystyle{plain}
\bibliography{cite}

\end{document}\documentclass{article}
\usepackage[utf8]{inputenc}
\usepackage{booktabs}
\usepackage{amsmath}
\usepackage{natbib}

\title{My Document}
\author{Your Name}
\date{\today}

\begin{document}

\maketitle

\section{Introduction}
This is a sample document demonstrating various \LaTeX{} features including mathematical equations, tables, and formatting. You can cite papers using the citation search feature, and they will be automatically added to your \texttt{cite.bib} file.

\section{Mathematical Expressions}
Here are some example equations:

\subsection{Quadratic Formula}
The quadratic formula is given by:
\[x = \frac{-b \pm \sqrt{b^2 - 4ac}}{2a}\]

\subsection{Matrix Example}
A sample matrix:
\[
\begin{pmatrix}
a_{11} & a_{12} & a_{13} \\
a_{21} & a_{22} & a_{23} \\
a_{31} & a_{32} & a_{33}
\end{pmatrix}
\]

\section{Tables}
Below is a sample table showing data:

\begin{table}[h]
\centering
\begin{tabular}{@{}lrrr@{}}
\toprule
\textbf{Category} & \textbf{Value 1} & \textbf{Value 2} & \textbf{Total} \\
\midrule
Item A & 10.5 & 20.3 & 30.8 \\
Item B & 15.2 & 25.1 & 40.3 \\
Item C & 22.1 & 12.7 & 34.8 \\
\bottomrule
\end{tabular}
\caption{Sample Data Table}
\label{tab:sample}
\end{table}

\section{Lists}
\subsection{Bullet Points}
\begin{itemize}
  \item First item
  \item Second item
  \item Third item with sub-items:
    \begin{itemize}
      \item Sub-item 1
      \item Sub-item 2
    \end{itemize}
\end{itemize}

\subsection{Numbered List}
\begin{enumerate}
  \item First step
  \item Second step
  \item Third step
\end{enumerate}

\section{Citations and References}
When you use the citation search feature, references will be automatically managed. For example, you might cite a paper like this: \cite{example2024}. All citations will be automatically added to your \texttt{cite.bib} file.

\section{Conclusion}
This template demonstrates basic \LaTeX{} features including mathematical formulas, tables, and lists. The bibliography system is now set up to work seamlessly with Bibby AI's citation management. Feel free to modify it for your needs!

\bibliographystyle{plain}
\bibliography{cite}

\end{document}\documentclass{article}
\usepackage[utf8]{inputenc}
\usepackage{booktabs}
\usepackage{amsmath}
\usepackage{natbib}

\title{My Document}
\author{Your Name}
\date{\today}

\begin{document}

\maketitle

\section{Introduction}
This is a sample document demonstrating various \LaTeX{} features including mathematical equations, tables, and formatting. You can cite papers using the citation search feature, and they will be automatically added to your \texttt{cite.bib} file.

\section{Mathematical Expressions}
Here are some example equations:

\subsection{Quadratic Formula}
The quadratic formula is given by:
\[x = \frac{-b \pm \sqrt{b^2 - 4ac}}{2a}\]

\subsection{Matrix Example}
A sample matrix:
\[
\begin{pmatrix}
a_{11} & a_{12} & a_{13} \\
a_{21} & a_{22} & a_{23} \\
a_{31} & a_{32} & a_{33}
\end{pmatrix}
\]

\section{Tables}
Below is a sample table showing data:

\begin{table}[h]
\centering
\begin{tabular}{@{}lrrr@{}}
\toprule
\textbf{Category} & \textbf{Value 1} & \textbf{Value 2} & \textbf{Total} \\
\midrule
Item A & 10.5 & 20.3 & 30.8 \\
Item B & 15.2 & 25.1 & 40.3 \\
Item C & 22.1 & 12.7 & 34.8 \\
\bottomrule
\end{tabular}
\caption{Sample Data Table}
\label{tab:sample}
\end{table}

\section{Lists}
\subsection{Bullet Points}
\begin{itemize}
  \item First item
  \item Second item
  \item Third item with sub-items:
    \begin{itemize}
      \item Sub-item 1
      \item Sub-item 2
    \end{itemize}
\end{itemize}

\subsection{Numbered List}
\begin{enumerate}
  \item First step
  \item Second step
  \item Third step
\end{enumerate}

\section{Citations and References}
When you use the citation search feature, references will be automatically managed. For example, you might cite a paper like this: \cite{example2024}. All citations will be automatically added to your \texttt{cite.bib} file.

\section{Conclusion}
This template demonstrates basic \LaTeX{} features including mathematical formulas, tables, and lists. The bibliography system is now set up to work seamlessly with Bibby AI's citation management. Feel free to modify it for your needs!

\bibliographystyle{plain}
\bibliography{cite}

\end{document}\documentclass{article}
\usepackage[utf8]{inputenc}
\usepackage{booktabs}
\usepackage{amsmath}
\usepackage{natbib}

\title{My Document}
\author{Your Name}
\date{\today}

\begin{document}

\maketitle

\section{Introduction}
This is a sample document demonstrating various \LaTeX{} features including mathematical equations, tables, and formatting. You can cite papers using the citation search feature, and they will be automatically added to your \texttt{cite.bib} file.

\section{Mathematical Expressions}
Here are some example equations:

\subsection{Quadratic Formula}
The quadratic formula is given by:
\[x = \frac{-b \pm \sqrt{b^2 - 4ac}}{2a}\]

\subsection{Matrix Example}
A sample matrix:
\[
\begin{pmatrix}
a_{11} & a_{12} & a_{13} \\
a_{21} & a_{22} & a_{23} \\
a_{31} & a_{32} & a_{33}
\end{pmatrix}
\]

\section{Tables}
Below is a sample table showing data:

\begin{table}[h]
\centering
\begin{tabular}{@{}lrrr@{}}
\toprule
\textbf{Category} & \textbf{Value 1} & \textbf{Value 2} & \textbf{Total} \\
\midrule
Item A & 10.5 & 20.3 & 30.8 \\
Item B & 15.2 & 25.1 & 40.3 \\
Item C & 22.1 & 12.7 & 34.8 \\
\bottomrule
\end{tabular}
\caption{Sample Data Table}
\label{tab:sample}
\end{table}

\section{Lists}
\subsection{Bullet Points}
\begin{itemize}
  \item First item
  \item Second item
  \item Third item with sub-items:
    \begin{itemize}
      \item Sub-item 1
      \item Sub-item 2
    \end{itemize}
\end{itemize}

\subsection{Numbered List}
\begin{enumerate}
  \item First step
  \item Second step
  \item Third step
\end{enumerate}

\section{Citations and References}
When you use the citation search feature, references will be automatically managed. For example, you might cite a paper like this: \cite{example2024}. All citations will be automatically added to your \texttt{cite.bib} file.

\section{Conclusion}
This template demonstrates basic \LaTeX{} features including mathematical formulas, tables, and lists. The bibliography system is now set up to work seamlessly with Bibby AI's citation management. Feel free to modify it for your needs!

\bibliographystyle{plain}
\bibliography{cite}

\end{document}\documentclass{article}
\usepackage[utf8]{inputenc}
\usepackage{booktabs}
\usepackage{amsmath}
\usepackage{natbib}

\title{My Document}
\author{Your Name}
\date{\today}

\begin{document}

\maketitle

\section{Introduction}
This is a sample document demonstrating various \LaTeX{} features including mathematical equations, tables, and formatting. You can cite papers using the citation search feature, and they will be automatically added to your \texttt{cite.bib} file.

\section{Mathematical Expressions}
Here are some example equations:

\subsection{Quadratic Formula}
The quadratic formula is given by:
\[x = \frac{-b \pm \sqrt{b^2 - 4ac}}{2a}\]

\subsection{Matrix Example}
A sample matrix:
\[
\begin{pmatrix}
a_{11} & a_{12} & a_{13} \\
a_{21} & a_{22} & a_{23} \\
a_{31} & a_{32} & a_{33}
\end{pmatrix}
\]

\section{Tables}
Below is a sample table showing data:

\begin{table}[h]
\centering
\begin{tabular}{@{}lrrr@{}}
\toprule
\textbf{Category} & \textbf{Value 1} & \textbf{Value 2} & \textbf{Total} \\
\midrule
Item A & 10.5 & 20.3 & 30.8 \\
Item B & 15.2 & 25.1 & 40.3 \\
Item C & 22.1 & 12.7 & 34.8 \\
\bottomrule
\end{tabular}
\caption{Sample Data Table}
\label{tab:sample}
\end{table}

\section{Lists}
\subsection{Bullet Points}
\begin{itemize}
  \item First item
  \item Second item
  \item Third item with sub-items:
    \begin{itemize}
      \item Sub-item 1
      \item Sub-item 2
    \end{itemize}
\end{itemize}

\subsection{Numbered List}
\begin{enumerate}
  \item First step
  \item Second step
  \item Third step
\end{enumerate}

\section{Citations and References}
When you use the citation search feature, references will be automatically managed. For example, you might cite a paper like this: \cite{example2024}. All citations will be automatically added to your \texttt{cite.bib} file.

\section{Conclusion}
This template demonstrates basic \LaTeX{} features including mathematical formulas, tables, and lists. The bibliography system is now set up to work seamlessly with Bibby AI's citation management. Feel free to modify it for your needs!

\bibliographystyle{plain}
\bibliography{cite}

\end{document}\documentclass{article}
\usepackage[utf8]{inputenc}
\usepackage{booktabs}
\usepackage{amsmath}
\usepackage{natbib}

\title{My Document}
\author{Your Name}
\date{\today}

\begin{document}

\maketitle

\section{Introduction}
This is a sample document demonstrating various \LaTeX{} features including mathematical equations, tables, and formatting. You can cite papers using the citation search feature, and they will be automatically added to your \texttt{cite.bib} file.

\section{Mathematical Expressions}
Here are some example equations:

\subsection{Quadratic Formula}
The quadratic formula is given by:
\[x = \frac{-b \pm \sqrt{b^2 - 4ac}}{2a}\]

\subsection{Matrix Example}
A sample matrix:
\[
\begin{pmatrix}
a_{11} & a_{12} & a_{13} \\
a_{21} & a_{22} & a_{23} \\
a_{31} & a_{32} & a_{33}
\end{pmatrix}
\]

\section{Tables}
Below is a sample table showing data:

\begin{table}[h]
\centering
\begin{tabular}{@{}lrrr@{}}
\toprule
\textbf{Category} & \textbf{Value 1} & \textbf{Value 2} & \textbf{Total} \\
\midrule
Item A & 10.5 & 20.3 & 30.8 \\
Item B & 15.2 & 25.1 & 40.3 \\
Item C & 22.1 & 12.7 & 34.8 \\
\bottomrule
\end{tabular}
\caption{Sample Data Table}
\label{tab:sample}
\end{table}

\section{Lists}
\subsection{Bullet Points}
\begin{itemize}
  \item First item
  \item Second item
  \item Third item with sub-items:
    \begin{itemize}
      \item Sub-item 1
      \item Sub-item 2
    \end{itemize}
\end{itemize}

\subsection{Numbered List}
\begin{enumerate}
  \item First step
  \item Second step
  \item Third step
\end{enumerate}

\section{Citations and References}
When you use the citation search feature, references will be automatically managed. For example, you might cite a paper like this: \cite{example2024}. All citations will be automatically added to your \texttt{cite.bib} file.

\section{Conclusion}
This template demonstrates basic \LaTeX{} features including mathematical formulas, tables, and lists. The bibliography system is now set up to work seamlessly with Bibby AI's citation management. Feel free to modify it for your needs!

\bibliographystyle{plain}
\bibliography{cite}

\end{document}\documentclass{article}
\usepackage[utf8]{inputenc}
\usepackage{booktabs}
\usepackage{amsmath}
\usepackage{natbib}

\title{My Document}
\author{Your Name}
\date{\today}

\begin{document}

\maketitle

\section{Introduction}
This is a sample document demonstrating various \LaTeX{} features including mathematical equations, tables, and formatting. You can cite papers using the citation search feature, and they will be automatically added to your \texttt{cite.bib} file.

\section{Mathematical Expressions}
Here are some example equations:

\subsection{Quadratic Formula}
The quadratic formula is given by:
\[x = \frac{-b \pm \sqrt{b^2 - 4ac}}{2a}\]

\subsection{Matrix Example}
A sample matrix:
\[
\begin{pmatrix}
a_{11} & a_{12} & a_{13} \\
a_{21} & a_{22} & a_{23} \\
a_{31} & a_{32} & a_{33}
\end{pmatrix}
\]

\section{Tables}
Below is a sample table showing data:

\begin{table}[h]
\centering
\begin{tabular}{@{}lrrr@{}}
\toprule
\textbf{Category} & \textbf{Value 1} & \textbf{Value 2} & \textbf{Total} \\
\midrule
Item A & 10.5 & 20.3 & 30.8 \\
Item B & 15.2 & 25.1 & 40.3 \\
Item C & 22.1 & 12.7 & 34.8 \\
\bottomrule
\end{tabular}
\caption{Sample Data Table}
\label{tab:sample}
\end{table}

\section{Lists}
\subsection{Bullet Points}
\begin{itemize}
  \item First item
  \item Second item
  \item Third item with sub-items:
    \begin{itemize}
      \item Sub-item 1
      \item Sub-item 2
    \end{itemize}
\end{itemize}

\subsection{Numbered List}
\begin{enumerate}
  \item First step
  \item Second step
  \item Third step
\end{enumerate}

\section{Citations and References}
When you use the citation search feature, references will be automatically managed. For example, you might cite a paper like this: \cite{example2024}. All citations will be automatically added to your \texttt{cite.bib} file.

\section{Conclusion}
This template demonstrates basic \LaTeX{} features including mathematical formulas, tables, and lists. The bibliography system is now set up to work seamlessly with Bibby AI's citation management. Feel free to modify it for your needs!

\bibliographystyle{plain}
\bibliography{cite}

\end{document}\documentclass{article}
\usepackage[utf8]{inputenc}
\usepackage{booktabs}
\usepackage{amsmath}
\usepackage{natbib}

\title{My Document}
\author{Your Name}
\date{\today}

\begin{document}

\maketitle

\section{Introduction}
This is a sample document demonstrating various \LaTeX{} features including mathematical equations, tables, and formatting. You can cite papers using the citation search feature, and they will be automatically added to your \texttt{cite.bib} file.

\section{Mathematical Expressions}
Here are some example equations:

\subsection{Quadratic Formula}
The quadratic formula is given by:
\[x = \frac{-b \pm \sqrt{b^2 - 4ac}}{2a}\]

\subsection{Matrix Example}
A sample matrix:
\[
\begin{pmatrix}
a_{11} & a_{12} & a_{13} \\
a_{21} & a_{22} & a_{23} \\
a_{31} & a_{32} & a_{33}
\end{pmatrix}
\]

\section{Tables}
Below is a sample table showing data:

\begin{table}[h]
\centering
\begin{tabular}{@{}lrrr@{}}
\toprule
\textbf{Category} & \textbf{Value 1} & \textbf{Value 2} & \textbf{Total} \\
\midrule
Item A & 10.5 & 20.3 & 30.8 \\
Item B & 15.2 & 25.1 & 40.3 \\
Item C & 22.1 & 12.7 & 34.8 \\
\bottomrule
\end{tabular}
\caption{Sample Data Table}
\label{tab:sample}
\end{table}

\section{Lists}
\subsection{Bullet Points}
\begin{itemize}
  \item First item
  \item Second item
  \item Third item with sub-items:
    \begin{itemize}
      \item Sub-item 1
      \item Sub-item 2
    \end{itemize}
\end{itemize}

\subsection{Numbered List}
\begin{enumerate}
  \item First step
  \item Second step
  \item Third step
\end{enumerate}

\section{Citations and References}
When you use the citation search feature, references will be automatically managed. For example, you might cite a paper like this: \cite{example2024}. All citations will be automatically added to your \texttt{cite.bib} file.

\section{Conclusion}
This template demonstrates basic \LaTeX{} features including mathematical formulas, tables, and lists. The bibliography system is now set up to work seamlessly with Bibby AI's citation management. Feel free to modify it for your needs!

\bibliographystyle{plain}
\bibliography{cite}

\end{document}\documentclass{article}
\usepackage[utf8]{inputenc}
\usepackage{booktabs}
\usepackage{amsmath}
\usepackage{natbib}

\title{My Document}
\author{Your Name}
\date{\today}

\begin{document}

\maketitle

\section{Introduction}
This is a sample document demonstrating various \LaTeX{} features including mathematical equations, tables, and formatting. You can cite papers using the citation search feature, and they will be automatically added to your \texttt{cite.bib} file.

\section{Mathematical Expressions}
Here are some example equations:

\subsection{Quadratic Formula}
The quadratic formula is given by:
\[x = \frac{-b \pm \sqrt{b^2 - 4ac}}{2a}\]

\subsection{Matrix Example}
A sample matrix:
\[
\begin{pmatrix}
a_{11} & a_{12} & a_{13} \\
a_{21} & a_{22} & a_{23} \\
a_{31} & a_{32} & a_{33}
\end{pmatrix}
\]

\section{Tables}
Below is a sample table showing data:

\begin{table}[h]
\centering
\begin{tabular}{@{}lrrr@{}}
\toprule
\textbf{Category} & \textbf{Value 1} & \textbf{Value 2} & \textbf{Total} \\
\midrule
Item A & 10.5 & 20.3 & 30.8 \\
Item B & 15.2 & 25.1 & 40.3 \\
Item C & 22.1 & 12.7 & 34.8 \\
\bottomrule
\end{tabular}
\caption{Sample Data Table}
\label{tab:sample}
\end{table}

\section{Lists}
\subsection{Bullet Points}
\begin{itemize}
  \item First item
  \item Second item
  \item Third item with sub-items:
    \begin{itemize}
      \item Sub-item 1
      \item Sub-item 2
    \end{itemize}
\end{itemize}

\subsection{Numbered List}
\begin{enumerate}
  \item First step
  \item Second step
  \item Third step
\end{enumerate}

\section{Citations and References}
When you use the citation search feature, references will be automatically managed. For example, you might cite a paper like this: \cite{example2024}. All citations will be automatically added to your \texttt{cite.bib} file.

\section{Conclusion}
This template demonstrates basic \LaTeX{} features including mathematical formulas, tables, and lists. The bibliography system is now set up to work seamlessly with Bibby AI's citation management. Feel free to modify it for your needs!

\bibliographystyle{plain}
\bibliography{cite}

\end{document}\documentclass{article}
\usepackage[utf8]{inputenc}
\usepackage{booktabs}
\usepackage{amsmath}
\usepackage{natbib}

\title{My Document}
\author{Your Name}
\date{\today}

\begin{document}

\maketitle

\section{Introduction}
This is a sample document demonstrating various \LaTeX{} features including mathematical equations, tables, and formatting. You can cite papers using the citation search feature, and they will be automatically added to your \texttt{cite.bib} file.

\section{Mathematical Expressions}
Here are some example equations:

\subsection{Quadratic Formula}
The quadratic formula is given by:
\[x = \frac{-b \pm \sqrt{b^2 - 4ac}}{2a}\]

\subsection{Matrix Example}
A sample matrix:
\[
\begin{pmatrix}
a_{11} & a_{12} & a_{13} \\
a_{21} & a_{22} & a_{23} \\
a_{31} & a_{32} & a_{33}
\end{pmatrix}
\]

\section{Tables}
Below is a sample table showing data:

\begin{table}[h]
\centering
\begin{tabular}{@{}lrrr@{}}
\toprule
\textbf{Category} & \textbf{Value 1} & \textbf{Value 2} & \textbf{Total} \\
\midrule
Item A & 10.5 & 20.3 & 30.8 \\
Item B & 15.2 & 25.1 & 40.3 \\
Item C & 22.1 & 12.7 & 34.8 \\
\bottomrule
\end{tabular}
\caption{Sample Data Table}
\label{tab:sample}
\end{table}

\section{Lists}
\subsection{Bullet Points}
\begin{itemize}
  \item First item
  \item Second item
  \item Third item with sub-items:
    \begin{itemize}
      \item Sub-item 1
      \item Sub-item 2
    \end{itemize}
\end{itemize}

\subsection{Numbered List}
\begin{enumerate}
  \item First step
  \item Second step
  \item Third step
\end{enumerate}

\section{Citations and References}
When you use the citation search feature, references will be automatically managed. For example, you might cite a paper like this: \cite{example2024}. All citations will be automatically added to your \texttt{cite.bib} file.

\section{Conclusion}
This template demonstrates basic \LaTeX{} features including mathematical formulas, tables, and lists. The bibliography system is now set up to work seamlessly with Bibby AI's citation management. Feel free to modify it for your needs!

\bibliographystyle{plain}
\bibliography{cite}

\end{document}\documentclass{article}
\usepackage[utf8]{inputenc}
\usepackage{booktabs}
\usepackage{amsmath}
\usepackage{natbib}

\title{My Document}
\author{Your Name}
\date{\today}

\begin{document}

\maketitle

\section{Introduction}
This is a sample document demonstrating various \LaTeX{} features including mathematical equations, tables, and formatting. You can cite papers using the citation search feature, and they will be automatically added to your \texttt{cite.bib} file.

\section{Mathematical Expressions}
Here are some example equations:

\subsection{Quadratic Formula}
The quadratic formula is given by:
\[x = \frac{-b \pm \sqrt{b^2 - 4ac}}{2a}\]

\subsection{Matrix Example}
A sample matrix:
\[
\begin{pmatrix}
a_{11} & a_{12} & a_{13} \\
a_{21} & a_{22} & a_{23} \\
a_{31} & a_{32} & a_{33}
\end{pmatrix}
\]

\section{Tables}
Below is a sample table showing data:

\begin{table}[h]
\centering
\begin{tabular}{@{}lrrr@{}}
\toprule
\textbf{Category} & \textbf{Value 1} & \textbf{Value 2} & \textbf{Total} \\
\midrule
Item A & 10.5 & 20.3 & 30.8 \\
Item B & 15.2 & 25.1 & 40.3 \\
Item C & 22.1 & 12.7 & 34.8 \\
\bottomrule
\end{tabular}
\caption{Sample Data Table}
\label{tab:sample}
\end{table}

\section{Lists}
\subsection{Bullet Points}
\begin{itemize}
  \item First item
  \item Second item
  \item Third item with sub-items:
    \begin{itemize}
      \item Sub-item 1
      \item Sub-item 2
    \end{itemize}
\end{itemize}

\subsection{Numbered List}
\begin{enumerate}
  \item First step
  \item Second step
  \item Third step
\end{enumerate}

\section{Citations and References}
When you use the citation search feature, references will be automatically managed. For example, you might cite a paper like this: \cite{example2024}. All citations will be automatically added to your \texttt{cite.bib} file.

\section{Conclusion}
This template demonstrates basic \LaTeX{} features including mathematical formulas, tables, and lists. The bibliography system is now set up to work seamlessly with Bibby AI's citation management. Feel free to modify it for your needs!

\bibliographystyle{plain}
\bibliography{cite}

\end{document}\documentclass{article}
\usepackage[utf8]{inputenc}
\usepackage{booktabs}
\usepackage{amsmath}
\usepackage{natbib}

\title{My Document}
\author{Your Name}
\date{\today}

\begin{document}

\maketitle

\section{Introduction}
This is a sample document demonstrating various \LaTeX{} features including mathematical equations, tables, and formatting. You can cite papers using the citation search feature, and they will be automatically added to your \texttt{cite.bib} file.

\section{Mathematical Expressions}
Here are some example equations:

\subsection{Quadratic Formula}
The quadratic formula is given by:
\[x = \frac{-b \pm \sqrt{b^2 - 4ac}}{2a}\]

\subsection{Matrix Example}
A sample matrix:
\[
\begin{pmatrix}
a_{11} & a_{12} & a_{13} \\
a_{21} & a_{22} & a_{23} \\
a_{31} & a_{32} & a_{33}
\end{pmatrix}
\]

\section{Tables}
Below is a sample table showing data:

\begin{table}[h]
\centering
\begin{tabular}{@{}lrrr@{}}
\toprule
\textbf{Category} & \textbf{Value 1} & \textbf{Value 2} & \textbf{Total} \\
\midrule
Item A & 10.5 & 20.3 & 30.8 \\
Item B & 15.2 & 25.1 & 40.3 \\
Item C & 22.1 & 12.7 & 34.8 \\
\bottomrule
\end{tabular}
\caption{Sample Data Table}
\label{tab:sample}
\end{table}

\section{Lists}
\subsection{Bullet Points}
\begin{itemize}
  \item First item
  \item Second item
  \item Third item with sub-items:
    \begin{itemize}
      \item Sub-item 1
      \item Sub-item 2
    \end{itemize}
\end{itemize}

\subsection{Numbered List}
\begin{enumerate}
  \item First step
  \item Second step
  \item Third step
\end{enumerate}

\section{Citations and References}
When you use the citation search feature, references will be automatically managed. For example, you might cite a paper like this: \cite{example2024}. All citations will be automatically added to your \texttt{cite.bib} file.

\section{Conclusion}
This template demonstrates basic \LaTeX{} features including mathematical formulas, tables, and lists. The bibliography system is now set up to work seamlessly with Bibby AI's citation management. Feel free to modify it for your needs!

\bibliographystyle{plain}
\bibliography{cite}

\end{document}\documentclass{article}
\usepackage[utf8]{inputenc}
\usepackage{booktabs}
\usepackage{amsmath}
\usepackage{natbib}

\title{My Document}
\author{Your Name}
\date{\today}

\begin{document}

\maketitle

\section{Introduction}
This is a sample document demonstrating various \LaTeX{} features including mathematical equations, tables, and formatting. You can cite papers using the citation search feature, and they will be automatically added to your \texttt{cite.bib} file.

\section{Mathematical Expressions}
Here are some example equations:

\subsection{Quadratic Formula}
The quadratic formula is given by:
\[x = \frac{-b \pm \sqrt{b^2 - 4ac}}{2a}\]

\subsection{Matrix Example}
A sample matrix:
\[
\begin{pmatrix}
a_{11} & a_{12} & a_{13} \\
a_{21} & a_{22} & a_{23} \\
a_{31} & a_{32} & a_{33}
\end{pmatrix}
\]

\section{Tables}
Below is a sample table showing data:

\begin{table}[h]
\centering
\begin{tabular}{@{}lrrr@{}}
\toprule
\textbf{Category} & \textbf{Value 1} & \textbf{Value 2} & \textbf{Total} \\
\midrule
Item A & 10.5 & 20.3 & 30.8 \\
Item B & 15.2 & 25.1 & 40.3 \\
Item C & 22.1 & 12.7 & 34.8 \\
\bottomrule
\end{tabular}
\caption{Sample Data Table}
\label{tab:sample}
\end{table}

\section{Lists}
\subsection{Bullet Points}
\begin{itemize}
  \item First item
  \item Second item
  \item Third item with sub-items:
    \begin{itemize}
      \item Sub-item 1
      \item Sub-item 2
    \end{itemize}
\end{itemize}

\subsection{Numbered List}
\begin{enumerate}
  \item First step
  \item Second step
  \item Third step
\end{enumerate}

\section{Citations and References}
When you use the citation search feature, references will be automatically managed. For example, you might cite a paper like this: \cite{example2024}. All citations will be automatically added to your \texttt{cite.bib} file.

\section{Conclusion}
This template demonstrates basic \LaTeX{} features including mathematical formulas, tables, and lists. The bibliography system is now set up to work seamlessly with Bibby AI's citation management. Feel free to modify it for your needs!

\bibliographystyle{plain}
\bibliography{cite}

\end{document}\documentclass{article}
\usepackage[utf8]{inputenc}
\usepackage{booktabs}
\usepackage{amsmath}
\usepackage{natbib}

\title{My Document}
\author{Your Name}
\date{\today}

\begin{document}

\maketitle

\section{Introduction}
This is a sample document demonstrating various \LaTeX{} features including mathematical equations, tables, and formatting. You can cite papers using the citation search feature, and they will be automatically added to your \texttt{cite.bib} file.

\section{Mathematical Expressions}
Here are some example equations:

\subsection{Quadratic Formula}
The quadratic formula is given by:
\[x = \frac{-b \pm \sqrt{b^2 - 4ac}}{2a}\]

\subsection{Matrix Example}
A sample matrix:
\[
\begin{pmatrix}
a_{11} & a_{12} & a_{13} \\
a_{21} & a_{22} & a_{23} \\
a_{31} & a_{32} & a_{33}
\end{pmatrix}
\]

\section{Tables}
Below is a sample table showing data:

\begin{table}[h]
\centering
\begin{tabular}{@{}lrrr@{}}
\toprule
\textbf{Category} & \textbf{Value 1} & \textbf{Value 2} & \textbf{Total} \\
\midrule
Item A & 10.5 & 20.3 & 30.8 \\
Item B & 15.2 & 25.1 & 40.3 \\
Item C & 22.1 & 12.7 & 34.8 \\
\bottomrule
\end{tabular}
\caption{Sample Data Table}
\label{tab:sample}
\end{table}

\section{Lists}
\subsection{Bullet Points}
\begin{itemize}
  \item First item
  \item Second item
  \item Third item with sub-items:
    \begin{itemize}
      \item Sub-item 1
      \item Sub-item 2
    \end{itemize}
\end{itemize}

\subsection{Numbered List}
\begin{enumerate}
  \item First step
  \item Second step
  \item Third step
\end{enumerate}

\section{Citations and References}
When you use the citation search feature, references will be automatically managed. For example, you might cite a paper like this: \cite{example2024}. All citations will be automatically added to your \texttt{cite.bib} file.

\section{Conclusion}
This template demonstrates basic \LaTeX{} features including mathematical formulas, tables, and lists. The bibliography system is now set up to work seamlessly with Bibby AI's citation management. Feel free to modify it for your needs!

\bibliographystyle{plain}
\bibliography{cite}

\end{document}\documentclass{article}
\usepackage[utf8]{inputenc}
\usepackage{booktabs}
\usepackage{amsmath}
\usepackage{natbib}

\title{My Document}
\author{Your Name}
\date{\today}

\begin{document}

\maketitle

\section{Introduction}
This is a sample document demonstrating various \LaTeX{} features including mathematical equations, tables, and formatting. You can cite papers using the citation search feature, and they will be automatically added to your \texttt{cite.bib} file.

\section{Mathematical Expressions}
Here are some example equations:

\subsection{Quadratic Formula}
The quadratic formula is given by:
\[x = \frac{-b \pm \sqrt{b^2 - 4ac}}{2a}\]

\subsection{Matrix Example}
A sample matrix:
\[
\begin{pmatrix}
a_{11} & a_{12} & a_{13} \\
a_{21} & a_{22} & a_{23} \\
a_{31} & a_{32} & a_{33}
\end{pmatrix}
\]

\section{Tables}
Below is a sample table showing data:

\begin{table}[h]
\centering
\begin{tabular}{@{}lrrr@{}}
\toprule
\textbf{Category} & \textbf{Value 1} & \textbf{Value 2} & \textbf{Total} \\
\midrule
Item A & 10.5 & 20.3 & 30.8 \\
Item B & 15.2 & 25.1 & 40.3 \\
Item C & 22.1 & 12.7 & 34.8 \\
\bottomrule
\end{tabular}
\caption{Sample Data Table}
\label{tab:sample}
\end{table}

\section{Lists}
\subsection{Bullet Points}
\begin{itemize}
  \item First item
  \item Second item
  \item Third item with sub-items:
    \begin{itemize}
      \item Sub-item 1
      \item Sub-item 2
    \end{itemize}
\end{itemize}

\subsection{Numbered List}
\begin{enumerate}
  \item First step
  \item Second step
  \item Third step
\end{enumerate}

\section{Citations and References}
When you use the citation search feature, references will be automatically managed. For example, you might cite a paper like this: \cite{example2024}. All citations will be automatically added to your \texttt{cite.bib} file.

\section{Conclusion}
This template demonstrates basic \LaTeX{} features including mathematical formulas, tables, and lists. The bibliography system is now set up to work seamlessly with Bibby AI's citation management. Feel free to modify it for your needs!

\bibliographystyle{plain}
\bibliography{cite}

\end{document}\documentclass{article}
\usepackage[utf8]{inputenc}
\usepackage{booktabs}
\usepackage{amsmath}
\usepackage{natbib}

\title{My Document}
\author{Your Name}
\date{\today}

\begin{document}

\maketitle

\section{Introduction}
This is a sample document demonstrating various \LaTeX{} features including mathematical equations, tables, and formatting. You can cite papers using the citation search feature, and they will be automatically added to your \texttt{cite.bib} file.

\section{Mathematical Expressions}
Here are some example equations:

\subsection{Quadratic Formula}
The quadratic formula is given by:
\[x = \frac{-b \pm \sqrt{b^2 - 4ac}}{2a}\]

\subsection{Matrix Example}
A sample matrix:
\[
\begin{pmatrix}
a_{11} & a_{12} & a_{13} \\
a_{21} & a_{22} & a_{23} \\
a_{31} & a_{32} & a_{33}
\end{pmatrix}
\]

\section{Tables}
Below is a sample table showing data:

\begin{table}[h]
\centering
\begin{tabular}{@{}lrrr@{}}
\toprule
\textbf{Category} & \textbf{Value 1} & \textbf{Value 2} & \textbf{Total} \\
\midrule
Item A & 10.5 & 20.3 & 30.8 \\
Item B & 15.2 & 25.1 & 40.3 \\
Item C & 22.1 & 12.7 & 34.8 \\
\bottomrule
\end{tabular}
\caption{Sample Data Table}
\label{tab:sample}
\end{table}

\section{Lists}
\subsection{Bullet Points}
\begin{itemize}
  \item First item
  \item Second item
  \item Third item with sub-items:
    \begin{itemize}
      \item Sub-item 1
      \item Sub-item 2
    \end{itemize}
\end{itemize}

\subsection{Numbered List}
\begin{enumerate}
  \item First step
  \item Second step
  \item Third step
\end{enumerate}

\section{Citations and References}
When you use the citation search feature, references will be automatically managed. For example, you might cite a paper like this: \cite{example2024}. All citations will be automatically added to your \texttt{cite.bib} file.

\section{Conclusion}
This template demonstrates basic \LaTeX{} features including mathematical formulas, tables, and lists. The bibliography system is now set up to work seamlessly with Bibby AI's citation management. Feel free to modify it for your needs!

\bibliographystyle{plain}
\bibliography{cite}

\end{document}\documentclass{article}
\usepackage[utf8]{inputenc}
\usepackage{booktabs}
\usepackage{amsmath}
\usepackage{natbib}

\title{My Document}
\author{Your Name}
\date{\today}

\begin{document}

\maketitle

\section{Introduction}
This is a sample document demonstrating various \LaTeX{} features including mathematical equations, tables, and formatting. You can cite papers using the citation search feature, and they will be automatically added to your \texttt{cite.bib} file.

\section{Mathematical Expressions}
Here are some example equations:

\subsection{Quadratic Formula}
The quadratic formula is given by:
\[x = \frac{-b \pm \sqrt{b^2 - 4ac}}{2a}\]

\subsection{Matrix Example}
A sample matrix:
\[
\begin{pmatrix}
a_{11} & a_{12} & a_{13} \\
a_{21} & a_{22} & a_{23} \\
a_{31} & a_{32} & a_{33}
\end{pmatrix}
\]

\section{Tables}
Below is a sample table showing data:

\begin{table}[h]
\centering
\begin{tabular}{@{}lrrr@{}}
\toprule
\textbf{Category} & \textbf{Value 1} & \textbf{Value 2} & \textbf{Total} \\
\midrule
Item A & 10.5 & 20.3 & 30.8 \\
Item B & 15.2 & 25.1 & 40.3 \\
Item C & 22.1 & 12.7 & 34.8 \\
\bottomrule
\end{tabular}
\caption{Sample Data Table}
\label{tab:sample}
\end{table}

\section{Lists}
\subsection{Bullet Points}
\begin{itemize}
  \item First item
  \item Second item
  \item Third item with sub-items:
    \begin{itemize}
      \item Sub-item 1
      \item Sub-item 2
    \end{itemize}
\end{itemize}

\subsection{Numbered List}
\begin{enumerate}
  \item First step
  \item Second step
  \item Third step
\end{enumerate}

\section{Citations and References}
When you use the citation search feature, references will be automatically managed. For example, you might cite a paper like this: \cite{example2024}. All citations will be automatically added to your \texttt{cite.bib} file.

\section{Conclusion}
This template demonstrates basic \LaTeX{} features including mathematical formulas, tables, and lists. The bibliography system is now set up to work seamlessly with Bibby AI's citation management. Feel free to modify it for your needs!

\bibliographystyle{plain}
\bibliography{cite}

\end{document}\documentclass{article}
\usepackage[utf8]{inputenc}
\usepackage{booktabs}
\usepackage{amsmath}
\usepackage{natbib}

\title{My Document}
\author{Your Name}
\date{\today}

\begin{document}

\maketitle

\section{Introduction}
This is a sample document demonstrating various \LaTeX{} features including mathematical equations, tables, and formatting. You can cite papers using the citation search feature, and they will be automatically added to your \texttt{cite.bib} file.

\section{Mathematical Expressions}
Here are some example equations:

\subsection{Quadratic Formula}
The quadratic formula is given by:
\[x = \frac{-b \pm \sqrt{b^2 - 4ac}}{2a}\]

\subsection{Matrix Example}
A sample matrix:
\[
\begin{pmatrix}
a_{11} & a_{12} & a_{13} \\
a_{21} & a_{22} & a_{23} \\
a_{31} & a_{32} & a_{33}
\end{pmatrix}
\]

\section{Tables}
Below is a sample table showing data:

\begin{table}[h]
\centering
\begin{tabular}{@{}lrrr@{}}
\toprule
\textbf{Category} & \textbf{Value 1} & \textbf{Value 2} & \textbf{Total} \\
\midrule
Item A & 10.5 & 20.3 & 30.8 \\
Item B & 15.2 & 25.1 & 40.3 \\
Item C & 22.1 & 12.7 & 34.8 \\
\bottomrule
\end{tabular}
\caption{Sample Data Table}
\label{tab:sample}
\end{table}

\section{Lists}
\subsection{Bullet Points}
\begin{itemize}
  \item First item
  \item Second item
  \item Third item with sub-items:
    \begin{itemize}
      \item Sub-item 1
      \item Sub-item 2
    \end{itemize}
\end{itemize}

\subsection{Numbered List}
\begin{enumerate}
  \item First step
  \item Second step
  \item Third step
\end{enumerate}

\section{Citations and References}
When you use the citation search feature, references will be automatically managed. For example, you might cite a paper like this: \cite{example2024}. All citations will be automatically added to your \texttt{cite.bib} file.

\section{Conclusion}
This template demonstrates basic \LaTeX{} features including mathematical formulas, tables, and lists. The bibliography system is now set up to work seamlessly with Bibby AI's citation management. Feel free to modify it for your needs!

\bibliographystyle{plain}
\bibliography{cite}

\end{document}\documentclass{article}
\usepackage[utf8]{inputenc}
\usepackage{booktabs}
\usepackage{amsmath}
\usepackage{natbib}

\title{My Document}
\author{Your Name}
\date{\today}

\begin{document}

\maketitle

\section{Introduction}
This is a sample document demonstrating various \LaTeX{} features including mathematical equations, tables, and formatting. You can cite papers using the citation search feature, and they will be automatically added to your \texttt{cite.bib} file.

\section{Mathematical Expressions}
Here are some example equations:

\subsection{Quadratic Formula}
The quadratic formula is given by:
\[x = \frac{-b \pm \sqrt{b^2 - 4ac}}{2a}\]

\subsection{Matrix Example}
A sample matrix:
\[
\begin{pmatrix}
a_{11} & a_{12} & a_{13} \\
a_{21} & a_{22} & a_{23} \\
a_{31} & a_{32} & a_{33}
\end{pmatrix}
\]

\section{Tables}
Below is a sample table showing data:

\begin{table}[h]
\centering
\begin{tabular}{@{}lrrr@{}}
\toprule
\textbf{Category} & \textbf{Value 1} & \textbf{Value 2} & \textbf{Total} \\
\midrule
Item A & 10.5 & 20.3 & 30.8 \\
Item B & 15.2 & 25.1 & 40.3 \\
Item C & 22.1 & 12.7 & 34.8 \\
\bottomrule
\end{tabular}
\caption{Sample Data Table}
\label{tab:sample}
\end{table}

\section{Lists}
\subsection{Bullet Points}
\begin{itemize}
  \item First item
  \item Second item
  \item Third item with sub-items:
    \begin{itemize}
      \item Sub-item 1
      \item Sub-item 2
    \end{itemize}
\end{itemize}

\subsection{Numbered List}
\begin{enumerate}
  \item First step
  \item Second step
  \item Third step
\end{enumerate}

\section{Citations and References}
When you use the citation search feature, references will be automatically managed. For example, you might cite a paper like this: \cite{example2024}. All citations will be automatically added to your \texttt{cite.bib} file.

\section{Conclusion}
This template demonstrates basic \LaTeX{} features including mathematical formulas, tables, and lists. The bibliography system is now set up to work seamlessly with Bibby AI's citation management. Feel free to modify it for your needs!

\bibliographystyle{plain}
\bibliography{cite}

\end{document}\documentclass{article}
\usepackage[utf8]{inputenc}
\usepackage{booktabs}
\usepackage{amsmath}
\usepackage{natbib}

\title{My Document}
\author{Your Name}
\date{\today}

\begin{document}

\maketitle

\section{Introduction}
This is a sample document demonstrating various \LaTeX{} features including mathematical equations, tables, and formatting. You can cite papers using the citation search feature, and they will be automatically added to your \texttt{cite.bib} file.

\section{Mathematical Expressions}
Here are some example equations:

\subsection{Quadratic Formula}
The quadratic formula is given by:
\[x = \frac{-b \pm \sqrt{b^2 - 4ac}}{2a}\]

\subsection{Matrix Example}
A sample matrix:
\[
\begin{pmatrix}
a_{11} & a_{12} & a_{13} \\
a_{21} & a_{22} & a_{23} \\
a_{31} & a_{32} & a_{33}
\end{pmatrix}
\]

\section{Tables}
Below is a sample table showing data:

\begin{table}[h]
\centering
\begin{tabular}{@{}lrrr@{}}
\toprule
\textbf{Category} & \textbf{Value 1} & \textbf{Value 2} & \textbf{Total} \\
\midrule
Item A & 10.5 & 20.3 & 30.8 \\
Item B & 15.2 & 25.1 & 40.3 \\
Item C & 22.1 & 12.7 & 34.8 \\
\bottomrule
\end{tabular}
\caption{Sample Data Table}
\label{tab:sample}
\end{table}

\section{Lists}
\subsection{Bullet Points}
\begin{itemize}
  \item First item
  \item Second item
  \item Third item with sub-items:
    \begin{itemize}
      \item Sub-item 1
      \item Sub-item 2
    \end{itemize}
\end{itemize}

\subsection{Numbered List}
\begin{enumerate}
  \item First step
  \item Second step
  \item Third step
\end{enumerate}

\section{Citations and References}
When you use the citation search feature, references will be automatically managed. For example, you might cite a paper like this: \cite{example2024}. All citations will be automatically added to your \texttt{cite.bib} file.

\section{Conclusion}
This template demonstrates basic \LaTeX{} features including mathematical formulas, tables, and lists. The bibliography system is now set up to work seamlessly with Bibby AI's citation management. Feel free to modify it for your needs!

\bibliographystyle{plain}
\bibliography{cite}

\end{document}\documentclass{article}
\usepackage[utf8]{inputenc}
\usepackage{booktabs}
\usepackage{amsmath}
\usepackage{natbib}

\title{My Document}
\author{Your Name}
\date{\today}

\begin{document}

\maketitle

\section{Introduction}
This is a sample document demonstrating various \LaTeX{} features including mathematical equations, tables, and formatting. You can cite papers using the citation search feature, and they will be automatically added to your \texttt{cite.bib} file.

\section{Mathematical Expressions}
Here are some example equations:

\subsection{Quadratic Formula}
The quadratic formula is given by:
\[x = \frac{-b \pm \sqrt{b^2 - 4ac}}{2a}\]

\subsection{Matrix Example}
A sample matrix:
\[
\begin{pmatrix}
a_{11} & a_{12} & a_{13} \\
a_{21} & a_{22} & a_{23} \\
a_{31} & a_{32} & a_{33}
\end{pmatrix}
\]

\section{Tables}
Below is a sample table showing data:

\begin{table}[h]
\centering
\begin{tabular}{@{}lrrr@{}}
\toprule
\textbf{Category} & \textbf{Value 1} & \textbf{Value 2} & \textbf{Total} \\
\midrule
Item A & 10.5 & 20.3 & 30.8 \\
Item B & 15.2 & 25.1 & 40.3 \\
Item C & 22.1 & 12.7 & 34.8 \\
\bottomrule
\end{tabular}
\caption{Sample Data Table}
\label{tab:sample}
\end{table}

\section{Lists}
\subsection{Bullet Points}
\begin{itemize}
  \item First item
  \item Second item
  \item Third item with sub-items:
    \begin{itemize}
      \item Sub-item 1
      \item Sub-item 2
    \end{itemize}
\end{itemize}

\subsection{Numbered List}
\begin{enumerate}
  \item First step
  \item Second step
  \item Third step
\end{enumerate}

\section{Citations and References}
When you use the citation search feature, references will be automatically managed. For example, you might cite a paper like this: \cite{example2024}. All citations will be automatically added to your \texttt{cite.bib} file.

\section{Conclusion}
This template demonstrates basic \LaTeX{} features including mathematical formulas, tables, and lists. The bibliography system is now set up to work seamlessly with Bibby AI's citation management. Feel free to modify it for your needs!

\bibliographystyle{plain}
\bibliography{cite}

\end{document}\documentclass{article}
\usepackage[utf8]{inputenc}
\usepackage{booktabs}
\usepackage{amsmath}
\usepackage{natbib}

\title{My Document}
\author{Your Name}
\date{\today}

\begin{document}

\maketitle

\section{Introduction}
This is a sample document demonstrating various \LaTeX{} features including mathematical equations, tables, and formatting. You can cite papers using the citation search feature, and they will be automatically added to your \texttt{cite.bib} file.

\section{Mathematical Expressions}
Here are some example equations:

\subsection{Quadratic Formula}
The quadratic formula is given by:
\[x = \frac{-b \pm \sqrt{b^2 - 4ac}}{2a}\]

\subsection{Matrix Example}
A sample matrix:
\[
\begin{pmatrix}
a_{11} & a_{12} & a_{13} \\
a_{21} & a_{22} & a_{23} \\
a_{31} & a_{32} & a_{33}
\end{pmatrix}
\]

\section{Tables}
Below is a sample table showing data:

\begin{table}[h]
\centering
\begin{tabular}{@{}lrrr@{}}
\toprule
\textbf{Category} & \textbf{Value 1} & \textbf{Value 2} & \textbf{Total} \\
\midrule
Item A & 10.5 & 20.3 & 30.8 \\
Item B & 15.2 & 25.1 & 40.3 \\
Item C & 22.1 & 12.7 & 34.8 \\
\bottomrule
\end{tabular}
\caption{Sample Data Table}
\label{tab:sample}
\end{table}

\section{Lists}
\subsection{Bullet Points}
\begin{itemize}
  \item First item
  \item Second item
  \item Third item with sub-items:
    \begin{itemize}
      \item Sub-item 1
      \item Sub-item 2
    \end{itemize}
\end{itemize}

\subsection{Numbered List}
\begin{enumerate}
  \item First step
  \item Second step
  \item Third step
\end{enumerate}

\section{Citations and References}
When you use the citation search feature, references will be automatically managed. For example, you might cite a paper like this: \cite{example2024}. All citations will be automatically added to your \texttt{cite.bib} file.

\section{Conclusion}
This template demonstrates basic \LaTeX{} features including mathematical formulas, tables, and lists. The bibliography system is now set up to work seamlessly with Bibby AI's citation management. Feel free to modify it for your needs!

\bibliographystyle{plain}
\bibliography{cite}

\end{document}\documentclass{article}
\usepackage[utf8]{inputenc}
\usepackage{booktabs}
\usepackage{amsmath}
\usepackage{natbib}

\title{My Document}
\author{Your Name}
\date{\today}

\begin{document}

\maketitle

\section{Introduction}
This is a sample document demonstrating various \LaTeX{} features including mathematical equations, tables, and formatting. You can cite papers using the citation search feature, and they will be automatically added to your \texttt{cite.bib} file.

\section{Mathematical Expressions}
Here are some example equations:

\subsection{Quadratic Formula}
The quadratic formula is given by:
\[x = \frac{-b \pm \sqrt{b^2 - 4ac}}{2a}\]

\subsection{Matrix Example}
A sample matrix:
\[
\begin{pmatrix}
a_{11} & a_{12} & a_{13} \\
a_{21} & a_{22} & a_{23} \\
a_{31} & a_{32} & a_{33}
\end{pmatrix}
\]

\section{Tables}
Below is a sample table showing data:

\begin{table}[h]
\centering
\begin{tabular}{@{}lrrr@{}}
\toprule
\textbf{Category} & \textbf{Value 1} & \textbf{Value 2} & \textbf{Total} \\
\midrule
Item A & 10.5 & 20.3 & 30.8 \\
Item B & 15.2 & 25.1 & 40.3 \\
Item C & 22.1 & 12.7 & 34.8 \\
\bottomrule
\end{tabular}
\caption{Sample Data Table}
\label{tab:sample}
\end{table}

\section{Lists}
\subsection{Bullet Points}
\begin{itemize}
  \item First item
  \item Second item
  \item Third item with sub-items:
    \begin{itemize}
      \item Sub-item 1
      \item Sub-item 2
    \end{itemize}
\end{itemize}

\subsection{Numbered List}
\begin{enumerate}
  \item First step
  \item Second step
  \item Third step
\end{enumerate}

\section{Citations and References}
When you use the citation search feature, references will be automatically managed. For example, you might cite a paper like this: \cite{example2024}. All citations will be automatically added to your \texttt{cite.bib} file.

\section{Conclusion}
This template demonstrates basic \LaTeX{} features including mathematical formulas, tables, and lists. The bibliography system is now set up to work seamlessly with Bibby AI's citation management. Feel free to modify it for your needs!

\bibliographystyle{plain}
\bibliography{cite}

\end{document}\documentclass{article}
\usepackage[utf8]{inputenc}
\usepackage{booktabs}
\usepackage{amsmath}
\usepackage{natbib}

\title{My Document}
\author{Your Name}
\date{\today}

\begin{document}

\maketitle

\section{Introduction}
This is a sample document demonstrating various \LaTeX{} features including mathematical equations, tables, and formatting. You can cite papers using the citation search feature, and they will be automatically added to your \texttt{cite.bib} file.

\section{Mathematical Expressions}
Here are some example equations:

\subsection{Quadratic Formula}
The quadratic formula is given by:
\[x = \frac{-b \pm \sqrt{b^2 - 4ac}}{2a}\]

\subsection{Matrix Example}
A sample matrix:
\[
\begin{pmatrix}
a_{11} & a_{12} & a_{13} \\
a_{21} & a_{22} & a_{23} \\
a_{31} & a_{32} & a_{33}
\end{pmatrix}
\]

\section{Tables}
Below is a sample table showing data:

\begin{table}[h]
\centering
\begin{tabular}{@{}lrrr@{}}
\toprule
\textbf{Category} & \textbf{Value 1} & \textbf{Value 2} & \textbf{Total} \\
\midrule
Item A & 10.5 & 20.3 & 30.8 \\
Item B & 15.2 & 25.1 & 40.3 \\
Item C & 22.1 & 12.7 & 34.8 \\
\bottomrule
\end{tabular}
\caption{Sample Data Table}
\label{tab:sample}
\end{table}

\section{Lists}
\subsection{Bullet Points}
\begin{itemize}
  \item First item
  \item Second item
  \item Third item with sub-items:
    \begin{itemize}
      \item Sub-item 1
      \item Sub-item 2
    \end{itemize}
\end{itemize}

\subsection{Numbered List}
\begin{enumerate}
  \item First step
  \item Second step
  \item Third step
\end{enumerate}

\section{Citations and References}
When you use the citation search feature, references will be automatically managed. For example, you might cite a paper like this: \cite{example2024}. All citations will be automatically added to your \texttt{cite.bib} file.

\section{Conclusion}
This template demonstrates basic \LaTeX{} features including mathematical formulas, tables, and lists. The bibliography system is now set up to work seamlessly with Bibby AI's citation management. Feel free to modify it for your needs!

\bibliographystyle{plain}
\bibliography{cite}

\end{document}\documentclass{article}
\usepackage[utf8]{inputenc}
\usepackage{booktabs}
\usepackage{amsmath}
\usepackage{natbib}

\title{My Document}
\author{Your Name}
\date{\today}

\begin{document}

\maketitle

\section{Introduction}
This is a sample document demonstrating various \LaTeX{} features including mathematical equations, tables, and formatting. You can cite papers using the citation search feature, and they will be automatically added to your \texttt{cite.bib} file.

\section{Mathematical Expressions}
Here are some example equations:

\subsection{Quadratic Formula}
The quadratic formula is given by:
\[x = \frac{-b \pm \sqrt{b^2 - 4ac}}{2a}\]

\subsection{Matrix Example}
A sample matrix:
\[
\begin{pmatrix}
a_{11} & a_{12} & a_{13} \\
a_{21} & a_{22} & a_{23} \\
a_{31} & a_{32} & a_{33}
\end{pmatrix}
\]

\section{Tables}
Below is a sample table showing data:

\begin{table}[h]
\centering
\begin{tabular}{@{}lrrr@{}}
\toprule
\textbf{Category} & \textbf{Value 1} & \textbf{Value 2} & \textbf{Total} \\
\midrule
Item A & 10.5 & 20.3 & 30.8 \\
Item B & 15.2 & 25.1 & 40.3 \\
Item C & 22.1 & 12.7 & 34.8 \\
\bottomrule
\end{tabular}
\caption{Sample Data Table}
\label{tab:sample}
\end{table}

\section{Lists}
\subsection{Bullet Points}
\begin{itemize}
  \item First item
  \item Second item
  \item Third item with sub-items:
    \begin{itemize}
      \item Sub-item 1
      \item Sub-item 2
    \end{itemize}
\end{itemize}

\subsection{Numbered List}
\begin{enumerate}
  \item First step
  \item Second step
  \item Third step
\end{enumerate}

\section{Citations and References}
When you use the citation search feature, references will be automatically managed. For example, you might cite a paper like this: \cite{example2024}. All citations will be automatically added to your \texttt{cite.bib} file.

\section{Conclusion}
This template demonstrates basic \LaTeX{} features including mathematical formulas, tables, and lists. The bibliography system is now set up to work seamlessly with Bibby AI's citation management. Feel free to modify it for your needs!

\bibliographystyle{plain}
\bibliography{cite}

\end{document}\documentclass{article}
\usepackage[utf8]{inputenc}
\usepackage{booktabs}
\usepackage{amsmath}
\usepackage{natbib}

\title{My Document}
\author{Your Name}
\date{\today}

\begin{document}

\maketitle

\section{Introduction}
This is a sample document demonstrating various \LaTeX{} features including mathematical equations, tables, and formatting. You can cite papers using the citation search feature, and they will be automatically added to your \texttt{cite.bib} file.

\section{Mathematical Expressions}
Here are some example equations:

\subsection{Quadratic Formula}
The quadratic formula is given by:
\[x = \frac{-b \pm \sqrt{b^2 - 4ac}}{2a}\]

\subsection{Matrix Example}
A sample matrix:
\[
\begin{pmatrix}
a_{11} & a_{12} & a_{13} \\
a_{21} & a_{22} & a_{23} \\
a_{31} & a_{32} & a_{33}
\end{pmatrix}
\]

\section{Tables}
Below is a sample table showing data:

\begin{table}[h]
\centering
\begin{tabular}{@{}lrrr@{}}
\toprule
\textbf{Category} & \textbf{Value 1} & \textbf{Value 2} & \textbf{Total} \\
\midrule
Item A & 10.5 & 20.3 & 30.8 \\
Item B & 15.2 & 25.1 & 40.3 \\
Item C & 22.1 & 12.7 & 34.8 \\
\bottomrule
\end{tabular}
\caption{Sample Data Table}
\label{tab:sample}
\end{table}

\section{Lists}
\subsection{Bullet Points}
\begin{itemize}
  \item First item
  \item Second item
  \item Third item with sub-items:
    \begin{itemize}
      \item Sub-item 1
      \item Sub-item 2
    \end{itemize}
\end{itemize}

\subsection{Numbered List}
\begin{enumerate}
  \item First step
  \item Second step
  \item Third step
\end{enumerate}

\section{Citations and References}
When you use the citation search feature, references will be automatically managed. For example, you might cite a paper like this: \cite{example2024}. All citations will be automatically added to your \texttt{cite.bib} file.

\section{Conclusion}
This template demonstrates basic \LaTeX{} features including mathematical formulas, tables, and lists. The bibliography system is now set up to work seamlessly with Bibby AI's citation management. Feel free to modify it for your needs!

\bibliographystyle{plain}
\bibliography{cite}

\end{document}\documentclass{article}
\usepackage[utf8]{inputenc}
\usepackage{booktabs}
\usepackage{amsmath}
\usepackage{natbib}

\title{My Document}
\author{Your Name}
\date{\today}

\begin{document}

\maketitle

\section{Introduction}
This is a sample document demonstrating various \LaTeX{} features including mathematical equations, tables, and formatting. You can cite papers using the citation search feature, and they will be automatically added to your \texttt{cite.bib} file.

\section{Mathematical Expressions}
Here are some example equations:

\subsection{Quadratic Formula}
The quadratic formula is given by:
\[x = \frac{-b \pm \sqrt{b^2 - 4ac}}{2a}\]

\subsection{Matrix Example}
A sample matrix:
\[
\begin{pmatrix}
a_{11} & a_{12} & a_{13} \\
a_{21} & a_{22} & a_{23} \\
a_{31} & a_{32} & a_{33}
\end{pmatrix}
\]

\section{Tables}
Below is a sample table showing data:

\begin{table}[h]
\centering
\begin{tabular}{@{}lrrr@{}}
\toprule
\textbf{Category} & \textbf{Value 1} & \textbf{Value 2} & \textbf{Total} \\
\midrule
Item A & 10.5 & 20.3 & 30.8 \\
Item B & 15.2 & 25.1 & 40.3 \\
Item C & 22.1 & 12.7 & 34.8 \\
\bottomrule
\end{tabular}
\caption{Sample Data Table}
\label{tab:sample}
\end{table}

\section{Lists}
\subsection{Bullet Points}
\begin{itemize}
  \item First item
  \item Second item
  \item Third item with sub-items:
    \begin{itemize}
      \item Sub-item 1
      \item Sub-item 2
    \end{itemize}
\end{itemize}

\subsection{Numbered List}
\begin{enumerate}
  \item First step
  \item Second step
  \item Third step
\end{enumerate}

\section{Citations and References}
When you use the citation search feature, references will be automatically managed. For example, you might cite a paper like this: \cite{example2024}. All citations will be automatically added to your \texttt{cite.bib} file.

\section{Conclusion}
This template demonstrates basic \LaTeX{} features including mathematical formulas, tables, and lists. The bibliography system is now set up to work seamlessly with Bibby AI's citation management. Feel free to modify it for your needs!

\bibliographystyle{plain}
\bibliography{cite}

\end{document}\documentclass{article}
\usepackage[utf8]{inputenc}
\usepackage{booktabs}
\usepackage{amsmath}
\usepackage{natbib}

\title{My Document}
\author{Your Name}
\date{\today}

\begin{document}

\maketitle

\section{Introduction}
This is a sample document demonstrating various \LaTeX{} features including mathematical equations, tables, and formatting. You can cite papers using the citation search feature, and they will be automatically added to your \texttt{cite.bib} file.

\section{Mathematical Expressions}
Here are some example equations:

\subsection{Quadratic Formula}
The quadratic formula is given by:
\[x = \frac{-b \pm \sqrt{b^2 - 4ac}}{2a}\]

\subsection{Matrix Example}
A sample matrix:
\[
\begin{pmatrix}
a_{11} & a_{12} & a_{13} \\
a_{21} & a_{22} & a_{23} \\
a_{31} & a_{32} & a_{33}
\end{pmatrix}
\]

\section{Tables}
Below is a sample table showing data:

\begin{table}[h]
\centering
\begin{tabular}{@{}lrrr@{}}
\toprule
\textbf{Category} & \textbf{Value 1} & \textbf{Value 2} & \textbf{Total} \\
\midrule
Item A & 10.5 & 20.3 & 30.8 \\
Item B & 15.2 & 25.1 & 40.3 \\
Item C & 22.1 & 12.7 & 34.8 \\
\bottomrule
\end{tabular}
\caption{Sample Data Table}
\label{tab:sample}
\end{table}

\section{Lists}
\subsection{Bullet Points}
\begin{itemize}
  \item First item
  \item Second item
  \item Third item with sub-items:
    \begin{itemize}
      \item Sub-item 1
      \item Sub-item 2
    \end{itemize}
\end{itemize}

\subsection{Numbered List}
\begin{enumerate}
  \item First step
  \item Second step
  \item Third step
\end{enumerate}

\section{Citations and References}
When you use the citation search feature, references will be automatically managed. For example, you might cite a paper like this: \cite{example2024}. All citations will be automatically added to your \texttt{cite.bib} file.

\section{Conclusion}
This template demonstrates basic \LaTeX{} features including mathematical formulas, tables, and lists. The bibliography system is now set up to work seamlessly with Bibby AI's citation management. Feel free to modify it for your needs!

\bibliographystyle{plain}
\bibliography{cite}

\end{document}\documentclass{article}
\usepackage[utf8]{inputenc}
\usepackage{booktabs}
\usepackage{amsmath}
\usepackage{natbib}

\title{My Document}
\author{Your Name}
\date{\today}

\begin{document}

\maketitle

\section{Introduction}
This is a sample document demonstrating various \LaTeX{} features including mathematical equations, tables, and formatting. You can cite papers using the citation search feature, and they will be automatically added to your \texttt{cite.bib} file.

\section{Mathematical Expressions}
Here are some example equations:

\subsection{Quadratic Formula}
The quadratic formula is given by:
\[x = \frac{-b \pm \sqrt{b^2 - 4ac}}{2a}\]

\subsection{Matrix Example}
A sample matrix:
\[
\begin{pmatrix}
a_{11} & a_{12} & a_{13} \\
a_{21} & a_{22} & a_{23} \\
a_{31} & a_{32} & a_{33}
\end{pmatrix}
\]

\section{Tables}
Below is a sample table showing data:

\begin{table}[h]
\centering
\begin{tabular}{@{}lrrr@{}}
\toprule
\textbf{Category} & \textbf{Value 1} & \textbf{Value 2} & \textbf{Total} \\
\midrule
Item A & 10.5 & 20.3 & 30.8 \\
Item B & 15.2 & 25.1 & 40.3 \\
Item C & 22.1 & 12.7 & 34.8 \\
\bottomrule
\end{tabular}
\caption{Sample Data Table}
\label{tab:sample}
\end{table}

\section{Lists}
\subsection{Bullet Points}
\begin{itemize}
  \item First item
  \item Second item
  \item Third item with sub-items:
    \begin{itemize}
      \item Sub-item 1
      \item Sub-item 2
    \end{itemize}
\end{itemize}

\subsection{Numbered List}
\begin{enumerate}
  \item First step
  \item Second step
  \item Third step
\end{enumerate}

\section{Citations and References}
When you use the citation search feature, references will be automatically managed. For example, you might cite a paper like this: \cite{example2024}. All citations will be automatically added to your \texttt{cite.bib} file.

\section{Conclusion}
This template demonstrates basic \LaTeX{} features including mathematical formulas, tables, and lists. The bibliography system is now set up to work seamlessly with Bibby AI's citation management. Feel free to modify it for your needs!

\bibliographystyle{plain}
\bibliography{cite}

\end{document}\documentclass{article}
\usepackage[utf8]{inputenc}
\usepackage{booktabs}
\usepackage{amsmath}
\usepackage{natbib}

\title{My Document}
\author{Your Name}
\date{\today}

\begin{document}

\maketitle

\section{Introduction}
This is a sample document demonstrating various \LaTeX{} features including mathematical equations, tables, and formatting. You can cite papers using the citation search feature, and they will be automatically added to your \texttt{cite.bib} file.

\section{Mathematical Expressions}
Here are some example equations:

\subsection{Quadratic Formula}
The quadratic formula is given by:
\[x = \frac{-b \pm \sqrt{b^2 - 4ac}}{2a}\]

\subsection{Matrix Example}
A sample matrix:
\[
\begin{pmatrix}
a_{11} & a_{12} & a_{13} \\
a_{21} & a_{22} & a_{23} \\
a_{31} & a_{32} & a_{33}
\end{pmatrix}
\]

\section{Tables}
Below is a sample table showing data:

\begin{table}[h]
\centering
\begin{tabular}{@{}lrrr@{}}
\toprule
\textbf{Category} & \textbf{Value 1} & \textbf{Value 2} & \textbf{Total} \\
\midrule
Item A & 10.5 & 20.3 & 30.8 \\
Item B & 15.2 & 25.1 & 40.3 \\
Item C & 22.1 & 12.7 & 34.8 \\
\bottomrule
\end{tabular}
\caption{Sample Data Table}
\label{tab:sample}
\end{table}

\section{Lists}
\subsection{Bullet Points}
\begin{itemize}
  \item First item
  \item Second item
  \item Third item with sub-items:
    \begin{itemize}
      \item Sub-item 1
      \item Sub-item 2
    \end{itemize}
\end{itemize}

\subsection{Numbered List}
\begin{enumerate}
  \item First step
  \item Second step
  \item Third step
\end{enumerate}

\section{Citations and References}
When you use the citation search feature, references will be automatically managed. For example, you might cite a paper like this: \cite{example2024}. All citations will be automatically added to your \texttt{cite.bib} file.

\section{Conclusion}
This template demonstrates basic \LaTeX{} features including mathematical formulas, tables, and lists. The bibliography system is now set up to work seamlessly with Bibby AI's citation management. Feel free to modify it for your needs!

\bibliographystyle{plain}
\bibliography{cite}

\end{document}\documentclass{article}
\usepackage[utf8]{inputenc}
\usepackage{booktabs}
\usepackage{amsmath}
\usepackage{natbib}

\title{My Document}
\author{Your Name}
\date{\today}

\begin{document}

\maketitle

\section{Introduction}
This is a sample document demonstrating various \LaTeX{} features including mathematical equations, tables, and formatting. You can cite papers using the citation search feature, and they will be automatically added to your \texttt{cite.bib} file.

\section{Mathematical Expressions}
Here are some example equations:

\subsection{Quadratic Formula}
The quadratic formula is given by:
\[x = \frac{-b \pm \sqrt{b^2 - 4ac}}{2a}\]

\subsection{Matrix Example}
A sample matrix:
\[
\begin{pmatrix}
a_{11} & a_{12} & a_{13} \\
a_{21} & a_{22} & a_{23} \\
a_{31} & a_{32} & a_{33}
\end{pmatrix}
\]

\section{Tables}
Below is a sample table showing data:

\begin{table}[h]
\centering
\begin{tabular}{@{}lrrr@{}}
\toprule
\textbf{Category} & \textbf{Value 1} & \textbf{Value 2} & \textbf{Total} \\
\midrule
Item A & 10.5 & 20.3 & 30.8 \\
Item B & 15.2 & 25.1 & 40.3 \\
Item C & 22.1 & 12.7 & 34.8 \\
\bottomrule
\end{tabular}
\caption{Sample Data Table}
\label{tab:sample}
\end{table}

\section{Lists}
\subsection{Bullet Points}
\begin{itemize}
  \item First item
  \item Second item
  \item Third item with sub-items:
    \begin{itemize}
      \item Sub-item 1
      \item Sub-item 2
    \end{itemize}
\end{itemize}

\subsection{Numbered List}
\begin{enumerate}
  \item First step
  \item Second step
  \item Third step
\end{enumerate}

\section{Citations and References}
When you use the citation search feature, references will be automatically managed. For example, you might cite a paper like this: \cite{example2024}. All citations will be automatically added to your \texttt{cite.bib} file.

\section{Conclusion}
This template demonstrates basic \LaTeX{} features including mathematical formulas, tables, and lists. The bibliography system is now set up to work seamlessly with Bibby AI's citation management. Feel free to modify it for your needs!

\bibliographystyle{plain}
\bibliography{cite}

\end{document}\documentclass{article}
\usepackage[utf8]{inputenc}
\usepackage{booktabs}
\usepackage{amsmath}
\usepackage{natbib}

\title{My Document}
\author{Your Name}
\date{\today}

\begin{document}

\maketitle

\section{Introduction}
This is a sample document demonstrating various \LaTeX{} features including mathematical equations, tables, and formatting. You can cite papers using the citation search feature, and they will be automatically added to your \texttt{cite.bib} file.

\section{Mathematical Expressions}
Here are some example equations:

\subsection{Quadratic Formula}
The quadratic formula is given by:
\[x = \frac{-b \pm \sqrt{b^2 - 4ac}}{2a}\]

\subsection{Matrix Example}
A sample matrix:
\[
\begin{pmatrix}
a_{11} & a_{12} & a_{13} \\
a_{21} & a_{22} & a_{23} \\
a_{31} & a_{32} & a_{33}
\end{pmatrix}
\]

\section{Tables}
Below is a sample table showing data:

\begin{table}[h]
\centering
\begin{tabular}{@{}lrrr@{}}
\toprule
\textbf{Category} & \textbf{Value 1} & \textbf{Value 2} & \textbf{Total} \\
\midrule
Item A & 10.5 & 20.3 & 30.8 \\
Item B & 15.2 & 25.1 & 40.3 \\
Item C & 22.1 & 12.7 & 34.8 \\
\bottomrule
\end{tabular}
\caption{Sample Data Table}
\label{tab:sample}
\end{table}

\section{Lists}
\subsection{Bullet Points}
\begin{itemize}
  \item First item
  \item Second item
  \item Third item with sub-items:
    \begin{itemize}
      \item Sub-item 1
      \item Sub-item 2
    \end{itemize}
\end{itemize}

\subsection{Numbered List}
\begin{enumerate}
  \item First step
  \item Second step
  \item Third step
\end{enumerate}

\section{Citations and References}
When you use the citation search feature, references will be automatically managed. For example, you might cite a paper like this: \cite{example2024}. All citations will be automatically added to your \texttt{cite.bib} file.

\section{Conclusion}
This template demonstrates basic \LaTeX{} features including mathematical formulas, tables, and lists. The bibliography system is now set up to work seamlessly with Bibby AI's citation management. Feel free to modify it for your needs!

\bibliographystyle{plain}
\bibliography{cite}

\end{document}\documentclass{article}
\usepackage[utf8]{inputenc}
\usepackage{booktabs}
\usepackage{amsmath}
\usepackage{natbib}

\title{My Document}
\author{Your Name}
\date{\today}

\begin{document}

\maketitle

\section{Introduction}
This is a sample document demonstrating various \LaTeX{} features including mathematical equations, tables, and formatting. You can cite papers using the citation search feature, and they will be automatically added to your \texttt{cite.bib} file.

\section{Mathematical Expressions}
Here are some example equations:

\subsection{Quadratic Formula}
The quadratic formula is given by:
\[x = \frac{-b \pm \sqrt{b^2 - 4ac}}{2a}\]

\subsection{Matrix Example}
A sample matrix:
\[
\begin{pmatrix}
a_{11} & a_{12} & a_{13} \\
a_{21} & a_{22} & a_{23} \\
a_{31} & a_{32} & a_{33}
\end{pmatrix}
\]

\section{Tables}
Below is a sample table showing data:

\begin{table}[h]
\centering
\begin{tabular}{@{}lrrr@{}}
\toprule
\textbf{Category} & \textbf{Value 1} & \textbf{Value 2} & \textbf{Total} \\
\midrule
Item A & 10.5 & 20.3 & 30.8 \\
Item B & 15.2 & 25.1 & 40.3 \\
Item C & 22.1 & 12.7 & 34.8 \\
\bottomrule
\end{tabular}
\caption{Sample Data Table}
\label{tab:sample}
\end{table}

\section{Lists}
\subsection{Bullet Points}
\begin{itemize}
  \item First item
  \item Second item
  \item Third item with sub-items:
    \begin{itemize}
      \item Sub-item 1
      \item Sub-item 2
    \end{itemize}
\end{itemize}

\subsection{Numbered List}
\begin{enumerate}
  \item First step
  \item Second step
  \item Third step
\end{enumerate}

\section{Citations and References}
When you use the citation search feature, references will be automatically managed. For example, you might cite a paper like this: \cite{example2024}. All citations will be automatically added to your \texttt{cite.bib} file.

\section{Conclusion}
This template demonstrates basic \LaTeX{} features including mathematical formulas, tables, and lists. The bibliography system is now set up to work seamlessly with Bibby AI's citation management. Feel free to modify it for your needs!

\bibliographystyle{plain}
\bibliography{cite}

\end{document}\documentclass{article}
\usepackage[utf8]{inputenc}
\usepackage{booktabs}
\usepackage{amsmath}
\usepackage{natbib}

\title{My Document}
\author{Your Name}
\date{\today}

\begin{document}

\maketitle

\section{Introduction}
This is a sample document demonstrating various \LaTeX{} features including mathematical equations, tables, and formatting. You can cite papers using the citation search feature, and they will be automatically added to your \texttt{cite.bib} file.

\section{Mathematical Expressions}
Here are some example equations:

\subsection{Quadratic Formula}
The quadratic formula is given by:
\[x = \frac{-b \pm \sqrt{b^2 - 4ac}}{2a}\]

\subsection{Matrix Example}
A sample matrix:
\[
\begin{pmatrix}
a_{11} & a_{12} & a_{13} \\
a_{21} & a_{22} & a_{23} \\
a_{31} & a_{32} & a_{33}
\end{pmatrix}
\]

\section{Tables}
Below is a sample table showing data:

\begin{table}[h]
\centering
\begin{tabular}{@{}lrrr@{}}
\toprule
\textbf{Category} & \textbf{Value 1} & \textbf{Value 2} & \textbf{Total} \\
\midrule
Item A & 10.5 & 20.3 & 30.8 \\
Item B & 15.2 & 25.1 & 40.3 \\
Item C & 22.1 & 12.7 & 34.8 \\
\bottomrule
\end{tabular}
\caption{Sample Data Table}
\label{tab:sample}
\end{table}

\section{Lists}
\subsection{Bullet Points}
\begin{itemize}
  \item First item
  \item Second item
  \item Third item with sub-items:
    \begin{itemize}
      \item Sub-item 1
      \item Sub-item 2
    \end{itemize}
\end{itemize}

\subsection{Numbered List}
\begin{enumerate}
  \item First step
  \item Second step
  \item Third step
\end{enumerate}

\section{Citations and References}
When you use the citation search feature, references will be automatically managed. For example, you might cite a paper like this: \cite{example2024}. All citations will be automatically added to your \texttt{cite.bib} file.

\section{Conclusion}
This template demonstrates basic \LaTeX{} features including mathematical formulas, tables, and lists. The bibliography system is now set up to work seamlessly with Bibby AI's citation management. Feel free to modify it for your needs!

\bibliographystyle{plain}
\bibliography{cite}

\end{document}\documentclass{article}
\usepackage[utf8]{inputenc}
\usepackage{booktabs}
\usepackage{amsmath}
\usepackage{natbib}

\title{My Document}
\author{Your Name}
\date{\today}

\begin{document}

\maketitle

\section{Introduction}
This is a sample document demonstrating various \LaTeX{} features including mathematical equations, tables, and formatting. You can cite papers using the citation search feature, and they will be automatically added to your \texttt{cite.bib} file.

\section{Mathematical Expressions}
Here are some example equations:

\subsection{Quadratic Formula}
The quadratic formula is given by:
\[x = \frac{-b \pm \sqrt{b^2 - 4ac}}{2a}\]

\subsection{Matrix Example}
A sample matrix:
\[
\begin{pmatrix}
a_{11} & a_{12} & a_{13} \\
a_{21} & a_{22} & a_{23} \\
a_{31} & a_{32} & a_{33}
\end{pmatrix}
\]

\section{Tables}
Below is a sample table showing data:

\begin{table}[h]
\centering
\begin{tabular}{@{}lrrr@{}}
\toprule
\textbf{Category} & \textbf{Value 1} & \textbf{Value 2} & \textbf{Total} \\
\midrule
Item A & 10.5 & 20.3 & 30.8 \\
Item B & 15.2 & 25.1 & 40.3 \\
Item C & 22.1 & 12.7 & 34.8 \\
\bottomrule
\end{tabular}
\caption{Sample Data Table}
\label{tab:sample}
\end{table}

\section{Lists}
\subsection{Bullet Points}
\begin{itemize}
  \item First item
  \item Second item
  \item Third item with sub-items:
    \begin{itemize}
      \item Sub-item 1
      \item Sub-item 2
    \end{itemize}
\end{itemize}

\subsection{Numbered List}
\begin{enumerate}
  \item First step
  \item Second step
  \item Third step
\end{enumerate}

\section{Citations and References}
When you use the citation search feature, references will be automatically managed. For example, you might cite a paper like this: \cite{example2024}. All citations will be automatically added to your \texttt{cite.bib} file.

\section{Conclusion}
This template demonstrates basic \LaTeX{} features including mathematical formulas, tables, and lists. The bibliography system is now set up to work seamlessly with Bibby AI's citation management. Feel free to modify it for your needs!

\bibliographystyle{plain}
\bibliography{cite}

\end{document}\documentclass{article}
\usepackage[utf8]{inputenc}
\usepackage{booktabs}
\usepackage{amsmath}
\usepackage{natbib}

\title{My Document}
\author{Your Name}
\date{\today}

\begin{document}

\maketitle

\section{Introduction}
This is a sample document demonstrating various \LaTeX{} features including mathematical equations, tables, and formatting. You can cite papers using the citation search feature, and they will be automatically added to your \texttt{cite.bib} file.

\section{Mathematical Expressions}
Here are some example equations:

\subsection{Quadratic Formula}
The quadratic formula is given by:
\[x = \frac{-b \pm \sqrt{b^2 - 4ac}}{2a}\]

\subsection{Matrix Example}
A sample matrix:
\[
\begin{pmatrix}
a_{11} & a_{12} & a_{13} \\
a_{21} & a_{22} & a_{23} \\
a_{31} & a_{32} & a_{33}
\end{pmatrix}
\]

\section{Tables}
Below is a sample table showing data:

\begin{table}[h]
\centering
\begin{tabular}{@{}lrrr@{}}
\toprule
\textbf{Category} & \textbf{Value 1} & \textbf{Value 2} & \textbf{Total} \\
\midrule
Item A & 10.5 & 20.3 & 30.8 \\
Item B & 15.2 & 25.1 & 40.3 \\
Item C & 22.1 & 12.7 & 34.8 \\
\bottomrule
\end{tabular}
\caption{Sample Data Table}
\label{tab:sample}
\end{table}

\section{Lists}
\subsection{Bullet Points}
\begin{itemize}
  \item First item
  \item Second item
  \item Third item with sub-items:
    \begin{itemize}
      \item Sub-item 1
      \item Sub-item 2
    \end{itemize}
\end{itemize}

\subsection{Numbered List}
\begin{enumerate}
  \item First step
  \item Second step
  \item Third step
\end{enumerate}

\section{Citations and References}
When you use the citation search feature, references will be automatically managed. For example, you might cite a paper like this: \cite{example2024}. All citations will be automatically added to your \texttt{cite.bib} file.

\section{Conclusion}
This template demonstrates basic \LaTeX{} features including mathematical formulas, tables, and lists. The bibliography system is now set up to work seamlessly with Bibby AI's citation management. Feel free to modify it for your needs!

\bibliographystyle{plain}
\bibliography{cite}

\end{document}\documentclass{article}
\usepackage[utf8]{inputenc}
\usepackage{booktabs}
\usepackage{amsmath}
\usepackage{natbib}

\title{My Document}
\author{Your Name}
\date{\today}

\begin{document}

\maketitle

\section{Introduction}
This is a sample document demonstrating various \LaTeX{} features including mathematical equations, tables, and formatting. You can cite papers using the citation search feature, and they will be automatically added to your \texttt{cite.bib} file.

\section{Mathematical Expressions}
Here are some example equations:

\subsection{Quadratic Formula}
The quadratic formula is given by:
\[x = \frac{-b \pm \sqrt{b^2 - 4ac}}{2a}\]

\subsection{Matrix Example}
A sample matrix:
\[
\begin{pmatrix}
a_{11} & a_{12} & a_{13} \\
a_{21} & a_{22} & a_{23} \\
a_{31} & a_{32} & a_{33}
\end{pmatrix}
\]

\section{Tables}
Below is a sample table showing data:

\begin{table}[h]
\centering
\begin{tabular}{@{}lrrr@{}}
\toprule
\textbf{Category} & \textbf{Value 1} & \textbf{Value 2} & \textbf{Total} \\
\midrule
Item A & 10.5 & 20.3 & 30.8 \\
Item B & 15.2 & 25.1 & 40.3 \\
Item C & 22.1 & 12.7 & 34.8 \\
\bottomrule
\end{tabular}
\caption{Sample Data Table}
\label{tab:sample}
\end{table}

\section{Lists}
\subsection{Bullet Points}
\begin{itemize}
  \item First item
  \item Second item
  \item Third item with sub-items:
    \begin{itemize}
      \item Sub-item 1
      \item Sub-item 2
    \end{itemize}
\end{itemize}

\subsection{Numbered List}
\begin{enumerate}
  \item First step
  \item Second step
  \item Third step
\end{enumerate}

\section{Citations and References}
When you use the citation search feature, references will be automatically managed. For example, you might cite a paper like this: \cite{example2024}. All citations will be automatically added to your \texttt{cite.bib} file.

\section{Conclusion}
This template demonstrates basic \LaTeX{} features including mathematical formulas, tables, and lists. The bibliography system is now set up to work seamlessly with Bibby AI's citation management. Feel free to modify it for your needs!

\bibliographystyle{plain}
\bibliography{cite}

\end{document}\documentclass{article}
\usepackage[utf8]{inputenc}
\usepackage{booktabs}
\usepackage{amsmath}
\usepackage{natbib}

\title{My Document}
\author{Your Name}
\date{\today}

\begin{document}

\maketitle

\section{Introduction}
This is a sample document demonstrating various \LaTeX{} features including mathematical equations, tables, and formatting. You can cite papers using the citation search feature, and they will be automatically added to your \texttt{cite.bib} file.

\section{Mathematical Expressions}
Here are some example equations:

\subsection{Quadratic Formula}
The quadratic formula is given by:
\[x = \frac{-b \pm \sqrt{b^2 - 4ac}}{2a}\]

\subsection{Matrix Example}
A sample matrix:
\[
\begin{pmatrix}
a_{11} & a_{12} & a_{13} \\
a_{21} & a_{22} & a_{23} \\
a_{31} & a_{32} & a_{33}
\end{pmatrix}
\]

\section{Tables}
Below is a sample table showing data:

\begin{table}[h]
\centering
\begin{tabular}{@{}lrrr@{}}
\toprule
\textbf{Category} & \textbf{Value 1} & \textbf{Value 2} & \textbf{Total} \\
\midrule
Item A & 10.5 & 20.3 & 30.8 \\
Item B & 15.2 & 25.1 & 40.3 \\
Item C & 22.1 & 12.7 & 34.8 \\
\bottomrule
\end{tabular}
\caption{Sample Data Table}
\label{tab:sample}
\end{table}

\section{Lists}
\subsection{Bullet Points}
\begin{itemize}
  \item First item
  \item Second item
  \item Third item with sub-items:
    \begin{itemize}
      \item Sub-item 1
      \item Sub-item 2
    \end{itemize}
\end{itemize}

\subsection{Numbered List}
\begin{enumerate}
  \item First step
  \item Second step
  \item Third step
\end{enumerate}

\section{Citations and References}
When you use the citation search feature, references will be automatically managed. For example, you might cite a paper like this: \cite{example2024}. All citations will be automatically added to your \texttt{cite.bib} file.

\section{Conclusion}
This template demonstrates basic \LaTeX{} features including mathematical formulas, tables, and lists. The bibliography system is now set up to work seamlessly with Bibby AI's citation management. Feel free to modify it for your needs!

\bibliographystyle{plain}
\bibliography{cite}

\end{document}\documentclass{article}
\usepackage[utf8]{inputenc}
\usepackage{booktabs}
\usepackage{amsmath}
\usepackage{natbib}

\title{My Document}
\author{Your Name}
\date{\today}

\begin{document}

\maketitle

\section{Introduction}
This is a sample document demonstrating various \LaTeX{} features including mathematical equations, tables, and formatting. You can cite papers using the citation search feature, and they will be automatically added to your \texttt{cite.bib} file.

\section{Mathematical Expressions}
Here are some example equations:

\subsection{Quadratic Formula}
The quadratic formula is given by:
\[x = \frac{-b \pm \sqrt{b^2 - 4ac}}{2a}\]

\subsection{Matrix Example}
A sample matrix:
\[
\begin{pmatrix}
a_{11} & a_{12} & a_{13} \\
a_{21} & a_{22} & a_{23} \\
a_{31} & a_{32} & a_{33}
\end{pmatrix}
\]

\section{Tables}
Below is a sample table showing data:

\begin{table}[h]
\centering
\begin{tabular}{@{}lrrr@{}}
\toprule
\textbf{Category} & \textbf{Value 1} & \textbf{Value 2} & \textbf{Total} \\
\midrule
Item A & 10.5 & 20.3 & 30.8 \\
Item B & 15.2 & 25.1 & 40.3 \\
Item C & 22.1 & 12.7 & 34.8 \\
\bottomrule
\end{tabular}
\caption{Sample Data Table}
\label{tab:sample}
\end{table}

\section{Lists}
\subsection{Bullet Points}
\begin{itemize}
  \item First item
  \item Second item
  \item Third item with sub-items:
    \begin{itemize}
      \item Sub-item 1
      \item Sub-item 2
    \end{itemize}
\end{itemize}

\subsection{Numbered List}
\begin{enumerate}
  \item First step
  \item Second step
  \item Third step
\end{enumerate}

\section{Citations and References}
When you use the citation search feature, references will be automatically managed. For example, you might cite a paper like this: \cite{example2024}. All citations will be automatically added to your \texttt{cite.bib} file.

\section{Conclusion}
This template demonstrates basic \LaTeX{} features including mathematical formulas, tables, and lists. The bibliography system is now set up to work seamlessly with Bibby AI's citation management. Feel free to modify it for your needs!

\bibliographystyle{plain}
\bibliography{cite}

\end{document}\documentclass{article}
\usepackage[utf8]{inputenc}
\usepackage{booktabs}
\usepackage{amsmath}
\usepackage{natbib}

\title{My Document}
\author{Your Name}
\date{\today}

\begin{document}

\maketitle

\section{Introduction}
This is a sample document demonstrating various \LaTeX{} features including mathematical equations, tables, and formatting. You can cite papers using the citation search feature, and they will be automatically added to your \texttt{cite.bib} file.

\section{Mathematical Expressions}
Here are some example equations:

\subsection{Quadratic Formula}
The quadratic formula is given by:
\[x = \frac{-b \pm \sqrt{b^2 - 4ac}}{2a}\]

\subsection{Matrix Example}
A sample matrix:
\[
\begin{pmatrix}
a_{11} & a_{12} & a_{13} \\
a_{21} & a_{22} & a_{23} \\
a_{31} & a_{32} & a_{33}
\end{pmatrix}
\]

\section{Tables}
Below is a sample table showing data:

\begin{table}[h]
\centering
\begin{tabular}{@{}lrrr@{}}
\toprule
\textbf{Category} & \textbf{Value 1} & \textbf{Value 2} & \textbf{Total} \\
\midrule
Item A & 10.5 & 20.3 & 30.8 \\
Item B & 15.2 & 25.1 & 40.3 \\
Item C & 22.1 & 12.7 & 34.8 \\
\bottomrule
\end{tabular}
\caption{Sample Data Table}
\label{tab:sample}
\end{table}

\section{Lists}
\subsection{Bullet Points}
\begin{itemize}
  \item First item
  \item Second item
  \item Third item with sub-items:
    \begin{itemize}
      \item Sub-item 1
      \item Sub-item 2
    \end{itemize}
\end{itemize}

\subsection{Numbered List}
\begin{enumerate}
  \item First step
  \item Second step
  \item Third step
\end{enumerate}

\section{Citations and References}
When you use the citation search feature, references will be automatically managed. For example, you might cite a paper like this: \cite{example2024}. All citations will be automatically added to your \texttt{cite.bib} file.

\section{Conclusion}
This template demonstrates basic \LaTeX{} features including mathematical formulas, tables, and lists. The bibliography system is now set up to work seamlessly with Bibby AI's citation management. Feel free to modify it for your needs!

\bibliographystyle{plain}
\bibliography{cite}

\end{document}\documentclass{article}
\usepackage[utf8]{inputenc}
\usepackage{booktabs}
\usepackage{amsmath}
\usepackage{natbib}

\title{My Document}
\author{Your Name}
\date{\today}

\begin{document}

\maketitle

\section{Introduction}
This is a sample document demonstrating various \LaTeX{} features including mathematical equations, tables, and formatting. You can cite papers using the citation search feature, and they will be automatically added to your \texttt{cite.bib} file.

\section{Mathematical Expressions}
Here are some example equations:

\subsection{Quadratic Formula}
The quadratic formula is given by:
\[x = \frac{-b \pm \sqrt{b^2 - 4ac}}{2a}\]

\subsection{Matrix Example}
A sample matrix:
\[
\begin{pmatrix}
a_{11} & a_{12} & a_{13} \\
a_{21} & a_{22} & a_{23} \\
a_{31} & a_{32} & a_{33}
\end{pmatrix}
\]

\section{Tables}
Below is a sample table showing data:

\begin{table}[h]
\centering
\begin{tabular}{@{}lrrr@{}}
\toprule
\textbf{Category} & \textbf{Value 1} & \textbf{Value 2} & \textbf{Total} \\
\midrule
Item A & 10.5 & 20.3 & 30.8 \\
Item B & 15.2 & 25.1 & 40.3 \\
Item C & 22.1 & 12.7 & 34.8 \\
\bottomrule
\end{tabular}
\caption{Sample Data Table}
\label{tab:sample}
\end{table}

\section{Lists}
\subsection{Bullet Points}
\begin{itemize}
  \item First item
  \item Second item
  \item Third item with sub-items:
    \begin{itemize}
      \item Sub-item 1
      \item Sub-item 2
    \end{itemize}
\end{itemize}

\subsection{Numbered List}
\begin{enumerate}
  \item First step
  \item Second step
  \item Third step
\end{enumerate}

\section{Citations and References}
When you use the citation search feature, references will be automatically managed. For example, you might cite a paper like this: \cite{example2024}. All citations will be automatically added to your \texttt{cite.bib} file.

\section{Conclusion}
This template demonstrates basic \LaTeX{} features including mathematical formulas, tables, and lists. The bibliography system is now set up to work seamlessly with Bibby AI's citation management. Feel free to modify it for your needs!

\bibliographystyle{plain}
\bibliography{cite}

\end{document}\documentclass{article}
\usepackage[utf8]{inputenc}
\usepackage{booktabs}
\usepackage{amsmath}
\usepackage{natbib}

\title{My Document}
\author{Your Name}
\date{\today}

\begin{document}

\maketitle

\section{Introduction}
This is a sample document demonstrating various \LaTeX{} features including mathematical equations, tables, and formatting. You can cite papers using the citation search feature, and they will be automatically added to your \texttt{cite.bib} file.

\section{Mathematical Expressions}
Here are some example equations:

\subsection{Quadratic Formula}
The quadratic formula is given by:
\[x = \frac{-b \pm \sqrt{b^2 - 4ac}}{2a}\]

\subsection{Matrix Example}
A sample matrix:
\[
\begin{pmatrix}
a_{11} & a_{12} & a_{13} \\
a_{21} & a_{22} & a_{23} \\
a_{31} & a_{32} & a_{33}
\end{pmatrix}
\]

\section{Tables}
Below is a sample table showing data:

\begin{table}[h]
\centering
\begin{tabular}{@{}lrrr@{}}
\toprule
\textbf{Category} & \textbf{Value 1} & \textbf{Value 2} & \textbf{Total} \\
\midrule
Item A & 10.5 & 20.3 & 30.8 \\
Item B & 15.2 & 25.1 & 40.3 \\
Item C & 22.1 & 12.7 & 34.8 \\
\bottomrule
\end{tabular}
\caption{Sample Data Table}
\label{tab:sample}
\end{table}

\section{Lists}
\subsection{Bullet Points}
\begin{itemize}
  \item First item
  \item Second item
  \item Third item with sub-items:
    \begin{itemize}
      \item Sub-item 1
      \item Sub-item 2
    \end{itemize}
\end{itemize}

\subsection{Numbered List}
\begin{enumerate}
  \item First step
  \item Second step
  \item Third step
\end{enumerate}

\section{Citations and References}
When you use the citation search feature, references will be automatically managed. For example, you might cite a paper like this: \cite{example2024}. All citations will be automatically added to your \texttt{cite.bib} file.

\section{Conclusion}
This template demonstrates basic \LaTeX{} features including mathematical formulas, tables, and lists. The bibliography system is now set up to work seamlessly with Bibby AI's citation management. Feel free to modify it for your needs!

\bibliographystyle{plain}
\bibliography{cite}

\end{document}\documentclass{article}
\usepackage[utf8]{inputenc}
\usepackage{booktabs}
\usepackage{amsmath}
\usepackage{natbib}

\title{My Document}
\author{Your Name}
\date{\today}

\begin{document}

\maketitle

\section{Introduction}
This is a sample document demonstrating various \LaTeX{} features including mathematical equations, tables, and formatting. You can cite papers using the citation search feature, and they will be automatically added to your \texttt{cite.bib} file.

\section{Mathematical Expressions}
Here are some example equations:

\subsection{Quadratic Formula}
The quadratic formula is given by:
\[x = \frac{-b \pm \sqrt{b^2 - 4ac}}{2a}\]

\subsection{Matrix Example}
A sample matrix:
\[
\begin{pmatrix}
a_{11} & a_{12} & a_{13} \\
a_{21} & a_{22} & a_{23} \\
a_{31} & a_{32} & a_{33}
\end{pmatrix}
\]

\section{Tables}
Below is a sample table showing data:

\begin{table}[h]
\centering
\begin{tabular}{@{}lrrr@{}}
\toprule
\textbf{Category} & \textbf{Value 1} & \textbf{Value 2} & \textbf{Total} \\
\midrule
Item A & 10.5 & 20.3 & 30.8 \\
Item B & 15.2 & 25.1 & 40.3 \\
Item C & 22.1 & 12.7 & 34.8 \\
\bottomrule
\end{tabular}
\caption{Sample Data Table}
\label{tab:sample}
\end{table}

\section{Lists}
\subsection{Bullet Points}
\begin{itemize}
  \item First item
  \item Second item
  \item Third item with sub-items:
    \begin{itemize}
      \item Sub-item 1
      \item Sub-item 2
    \end{itemize}
\end{itemize}

\subsection{Numbered List}
\begin{enumerate}
  \item First step
  \item Second step
  \item Third step
\end{enumerate}

\section{Citations and References}
When you use the citation search feature, references will be automatically managed. For example, you might cite a paper like this: \cite{example2024}. All citations will be automatically added to your \texttt{cite.bib} file.

\section{Conclusion}
This template demonstrates basic \LaTeX{} features including mathematical formulas, tables, and lists. The bibliography system is now set up to work seamlessly with Bibby AI's citation management. Feel free to modify it for your needs!

\bibliographystyle{plain}
\bibliography{cite}

\end{document}\documentclass{article}
\usepackage[utf8]{inputenc}
\usepackage{booktabs}
\usepackage{amsmath}
\usepackage{natbib}

\title{My Document}
\author{Your Name}
\date{\today}

\begin{document}

\maketitle

\section{Introduction}
This is a sample document demonstrating various \LaTeX{} features including mathematical equations, tables, and formatting. You can cite papers using the citation search feature, and they will be automatically added to your \texttt{cite.bib} file.

\section{Mathematical Expressions}
Here are some example equations:

\subsection{Quadratic Formula}
The quadratic formula is given by:
\[x = \frac{-b \pm \sqrt{b^2 - 4ac}}{2a}\]

\subsection{Matrix Example}
A sample matrix:
\[
\begin{pmatrix}
a_{11} & a_{12} & a_{13} \\
a_{21} & a_{22} & a_{23} \\
a_{31} & a_{32} & a_{33}
\end{pmatrix}
\]

\section{Tables}
Below is a sample table showing data:

\begin{table}[h]
\centering
\begin{tabular}{@{}lrrr@{}}
\toprule
\textbf{Category} & \textbf{Value 1} & \textbf{Value 2} & \textbf{Total} \\
\midrule
Item A & 10.5 & 20.3 & 30.8 \\
Item B & 15.2 & 25.1 & 40.3 \\
Item C & 22.1 & 12.7 & 34.8 \\
\bottomrule
\end{tabular}
\caption{Sample Data Table}
\label{tab:sample}
\end{table}

\section{Lists}
\subsection{Bullet Points}
\begin{itemize}
  \item First item
  \item Second item
  \item Third item with sub-items:
    \begin{itemize}
      \item Sub-item 1
      \item Sub-item 2
    \end{itemize}
\end{itemize}

\subsection{Numbered List}
\begin{enumerate}
  \item First step
  \item Second step
  \item Third step
\end{enumerate}

\section{Citations and References}
When you use the citation search feature, references will be automatically managed. For example, you might cite a paper like this: \cite{example2024}. All citations will be automatically added to your \texttt{cite.bib} file.

\section{Conclusion}
This template demonstrates basic \LaTeX{} features including mathematical formulas, tables, and lists. The bibliography system is now set up to work seamlessly with Bibby AI's citation management. Feel free to modify it for your needs!

\bibliographystyle{plain}
\bibliography{cite}

\end{document}\documentclass{article}
\usepackage[utf8]{inputenc}
\usepackage{booktabs}
\usepackage{amsmath}
\usepackage{natbib}

\title{My Document}
\author{Your Name}
\date{\today}

\begin{document}

\maketitle

\section{Introduction}
This is a sample document demonstrating various \LaTeX{} features including mathematical equations, tables, and formatting. You can cite papers using the citation search feature, and they will be automatically added to your \texttt{cite.bib} file.

\section{Mathematical Expressions}
Here are some example equations:

\subsection{Quadratic Formula}
The quadratic formula is given by:
\[x = \frac{-b \pm \sqrt{b^2 - 4ac}}{2a}\]

\subsection{Matrix Example}
A sample matrix:
\[
\begin{pmatrix}
a_{11} & a_{12} & a_{13} \\
a_{21} & a_{22} & a_{23} \\
a_{31} & a_{32} & a_{33}
\end{pmatrix}
\]

\section{Tables}
Below is a sample table showing data:

\begin{table}[h]
\centering
\begin{tabular}{@{}lrrr@{}}
\toprule
\textbf{Category} & \textbf{Value 1} & \textbf{Value 2} & \textbf{Total} \\
\midrule
Item A & 10.5 & 20.3 & 30.8 \\
Item B & 15.2 & 25.1 & 40.3 \\
Item C & 22.1 & 12.7 & 34.8 \\
\bottomrule
\end{tabular}
\caption{Sample Data Table}
\label{tab:sample}
\end{table}

\section{Lists}
\subsection{Bullet Points}
\begin{itemize}
  \item First item
  \item Second item
  \item Third item with sub-items:
    \begin{itemize}
      \item Sub-item 1
      \item Sub-item 2
    \end{itemize}
\end{itemize}

\subsection{Numbered List}
\begin{enumerate}
  \item First step
  \item Second step
  \item Third step
\end{enumerate}

\section{Citations and References}
When you use the citation search feature, references will be automatically managed. For example, you might cite a paper like this: \cite{example2024}. All citations will be automatically added to your \texttt{cite.bib} file.

\section{Conclusion}
This template demonstrates basic \LaTeX{} features including mathematical formulas, tables, and lists. The bibliography system is now set up to work seamlessly with Bibby AI's citation management. Feel free to modify it for your needs!

\bibliographystyle{plain}
\bibliography{cite}

\end{document}\documentclass{article}
\usepackage[utf8]{inputenc}
\usepackage{booktabs}
\usepackage{amsmath}
\usepackage{natbib}

\title{My Document}
\author{Your Name}
\date{\today}

\begin{document}

\maketitle

\section{Introduction}
This is a sample document demonstrating various \LaTeX{} features including mathematical equations, tables, and formatting. You can cite papers using the citation search feature, and they will be automatically added to your \texttt{cite.bib} file.

\section{Mathematical Expressions}
Here are some example equations:

\subsection{Quadratic Formula}
The quadratic formula is given by:
\[x = \frac{-b \pm \sqrt{b^2 - 4ac}}{2a}\]

\subsection{Matrix Example}
A sample matrix:
\[
\begin{pmatrix}
a_{11} & a_{12} & a_{13} \\
a_{21} & a_{22} & a_{23} \\
a_{31} & a_{32} & a_{33}
\end{pmatrix}
\]

\section{Tables}
Below is a sample table showing data:

\begin{table}[h]
\centering
\begin{tabular}{@{}lrrr@{}}
\toprule
\textbf{Category} & \textbf{Value 1} & \textbf{Value 2} & \textbf{Total} \\
\midrule
Item A & 10.5 & 20.3 & 30.8 \\
Item B & 15.2 & 25.1 & 40.3 \\
Item C & 22.1 & 12.7 & 34.8 \\
\bottomrule
\end{tabular}
\caption{Sample Data Table}
\label{tab:sample}
\end{table}

\section{Lists}
\subsection{Bullet Points}
\begin{itemize}
  \item First item
  \item Second item
  \item Third item with sub-items:
    \begin{itemize}
      \item Sub-item 1
      \item Sub-item 2
    \end{itemize}
\end{itemize}

\subsection{Numbered List}
\begin{enumerate}
  \item First step
  \item Second step
  \item Third step
\end{enumerate}

\section{Citations and References}
When you use the citation search feature, references will be automatically managed. For example, you might cite a paper like this: \cite{example2024}. All citations will be automatically added to your \texttt{cite.bib} file.

\section{Conclusion}
This template demonstrates basic \LaTeX{} features including mathematical formulas, tables, and lists. The bibliography system is now set up to work seamlessly with Bibby AI's citation management. Feel free to modify it for your needs!

\bibliographystyle{plain}
\bibliography{cite}

\end{document}\documentclass{article}
\usepackage[utf8]{inputenc}
\usepackage{booktabs}
\usepackage{amsmath}
\usepackage{natbib}

\title{My Document}
\author{Your Name}
\date{\today}

\begin{document}

\maketitle

\section{Introduction}
This is a sample document demonstrating various \LaTeX{} features including mathematical equations, tables, and formatting. You can cite papers using the citation search feature, and they will be automatically added to your \texttt{cite.bib} file.

\section{Mathematical Expressions}
Here are some example equations:

\subsection{Quadratic Formula}
The quadratic formula is given by:
\[x = \frac{-b \pm \sqrt{b^2 - 4ac}}{2a}\]

\subsection{Matrix Example}
A sample matrix:
\[
\begin{pmatrix}
a_{11} & a_{12} & a_{13} \\
a_{21} & a_{22} & a_{23} \\
a_{31} & a_{32} & a_{33}
\end{pmatrix}
\]

\section{Tables}
Below is a sample table showing data:

\begin{table}[h]
\centering
\begin{tabular}{@{}lrrr@{}}
\toprule
\textbf{Category} & \textbf{Value 1} & \textbf{Value 2} & \textbf{Total} \\
\midrule
Item A & 10.5 & 20.3 & 30.8 \\
Item B & 15.2 & 25.1 & 40.3 \\
Item C & 22.1 & 12.7 & 34.8 \\
\bottomrule
\end{tabular}
\caption{Sample Data Table}
\label{tab:sample}
\end{table}

\section{Lists}
\subsection{Bullet Points}
\begin{itemize}
  \item First item
  \item Second item
  \item Third item with sub-items:
    \begin{itemize}
      \item Sub-item 1
      \item Sub-item 2
    \end{itemize}
\end{itemize}

\subsection{Numbered List}
\begin{enumerate}
  \item First step
  \item Second step
  \item Third step
\end{enumerate}

\section{Citations and References}
When you use the citation search feature, references will be automatically managed. For example, you might cite a paper like this: \cite{example2024}. All citations will be automatically added to your \texttt{cite.bib} file.

\section{Conclusion}
This template demonstrates basic \LaTeX{} features including mathematical formulas, tables, and lists. The bibliography system is now set up to work seamlessly with Bibby AI's citation management. Feel free to modify it for your needs!

\bibliographystyle{plain}
\bibliography{cite}

\end{document}\documentclass{article}
\usepackage[utf8]{inputenc}
\usepackage{booktabs}
\usepackage{amsmath}
\usepackage{natbib}

\title{My Document}
\author{Your Name}
\date{\today}

\begin{document}

\maketitle

\section{Introduction}
This is a sample document demonstrating various \LaTeX{} features including mathematical equations, tables, and formatting. You can cite papers using the citation search feature, and they will be automatically added to your \texttt{cite.bib} file.

\section{Mathematical Expressions}
Here are some example equations:

\subsection{Quadratic Formula}
The quadratic formula is given by:
\[x = \frac{-b \pm \sqrt{b^2 - 4ac}}{2a}\]

\subsection{Matrix Example}
A sample matrix:
\[
\begin{pmatrix}
a_{11} & a_{12} & a_{13} \\
a_{21} & a_{22} & a_{23} \\
a_{31} & a_{32} & a_{33}
\end{pmatrix}
\]

\section{Tables}
Below is a sample table showing data:

\begin{table}[h]
\centering
\begin{tabular}{@{}lrrr@{}}
\toprule
\textbf{Category} & \textbf{Value 1} & \textbf{Value 2} & \textbf{Total} \\
\midrule
Item A & 10.5 & 20.3 & 30.8 \\
Item B & 15.2 & 25.1 & 40.3 \\
Item C & 22.1 & 12.7 & 34.8 \\
\bottomrule
\end{tabular}
\caption{Sample Data Table}
\label{tab:sample}
\end{table}

\section{Lists}
\subsection{Bullet Points}
\begin{itemize}
  \item First item
  \item Second item
  \item Third item with sub-items:
    \begin{itemize}
      \item Sub-item 1
      \item Sub-item 2
    \end{itemize}
\end{itemize}

\subsection{Numbered List}
\begin{enumerate}
  \item First step
  \item Second step
  \item Third step
\end{enumerate}

\section{Citations and References}
When you use the citation search feature, references will be automatically managed. For example, you might cite a paper like this: \cite{example2024}. All citations will be automatically added to your \texttt{cite.bib} file.

\section{Conclusion}
This template demonstrates basic \LaTeX{} features including mathematical formulas, tables, and lists. The bibliography system is now set up to work seamlessly with Bibby AI's citation management. Feel free to modify it for your needs!

\bibliographystyle{plain}
\bibliography{cite}

\end{document}\documentclass{article}
\usepackage[utf8]{inputenc}
\usepackage{booktabs}
\usepackage{amsmath}
\usepackage{natbib}

\title{My Document}
\author{Your Name}
\date{\today}

\begin{document}

\maketitle

\section{Introduction}
This is a sample document demonstrating various \LaTeX{} features including mathematical equations, tables, and formatting. You can cite papers using the citation search feature, and they will be automatically added to your \texttt{cite.bib} file.

\section{Mathematical Expressions}
Here are some example equations:

\subsection{Quadratic Formula}
The quadratic formula is given by:
\[x = \frac{-b \pm \sqrt{b^2 - 4ac}}{2a}\]

\subsection{Matrix Example}
A sample matrix:
\[
\begin{pmatrix}
a_{11} & a_{12} & a_{13} \\
a_{21} & a_{22} & a_{23} \\
a_{31} & a_{32} & a_{33}
\end{pmatrix}
\]

\section{Tables}
Below is a sample table showing data:

\begin{table}[h]
\centering
\begin{tabular}{@{}lrrr@{}}
\toprule
\textbf{Category} & \textbf{Value 1} & \textbf{Value 2} & \textbf{Total} \\
\midrule
Item A & 10.5 & 20.3 & 30.8 \\
Item B & 15.2 & 25.1 & 40.3 \\
Item C & 22.1 & 12.7 & 34.8 \\
\bottomrule
\end{tabular}
\caption{Sample Data Table}
\label{tab:sample}
\end{table}

\section{Lists}
\subsection{Bullet Points}
\begin{itemize}
  \item First item
  \item Second item
  \item Third item with sub-items:
    \begin{itemize}
      \item Sub-item 1
      \item Sub-item 2
    \end{itemize}
\end{itemize}

\subsection{Numbered List}
\begin{enumerate}
  \item First step
  \item Second step
  \item Third step
\end{enumerate}

\section{Citations and References}
When you use the citation search feature, references will be automatically managed. For example, you might cite a paper like this: \cite{example2024}. All citations will be automatically added to your \texttt{cite.bib} file.

\section{Conclusion}
This template demonstrates basic \LaTeX{} features including mathematical formulas, tables, and lists. The bibliography system is now set up to work seamlessly with Bibby AI's citation management. Feel free to modify it for your needs!

\bibliographystyle{plain}
\bibliography{cite}

\end{document}\documentclass{article}
\usepackage[utf8]{inputenc}
\usepackage{booktabs}
\usepackage{amsmath}
\usepackage{natbib}

\title{My Document}
\author{Your Name}
\date{\today}

\begin{document}

\maketitle

\section{Introduction}
This is a sample document demonstrating various \LaTeX{} features including mathematical equations, tables, and formatting. You can cite papers using the citation search feature, and they will be automatically added to your \texttt{cite.bib} file.

\section{Mathematical Expressions}
Here are some example equations:

\subsection{Quadratic Formula}
The quadratic formula is given by:
\[x = \frac{-b \pm \sqrt{b^2 - 4ac}}{2a}\]

\subsection{Matrix Example}
A sample matrix:
\[
\begin{pmatrix}
a_{11} & a_{12} & a_{13} \\
a_{21} & a_{22} & a_{23} \\
a_{31} & a_{32} & a_{33}
\end{pmatrix}
\]

\section{Tables}
Below is a sample table showing data:

\begin{table}[h]
\centering
\begin{tabular}{@{}lrrr@{}}
\toprule
\textbf{Category} & \textbf{Value 1} & \textbf{Value 2} & \textbf{Total} \\
\midrule
Item A & 10.5 & 20.3 & 30.8 \\
Item B & 15.2 & 25.1 & 40.3 \\
Item C & 22.1 & 12.7 & 34.8 \\
\bottomrule
\end{tabular}
\caption{Sample Data Table}
\label{tab:sample}
\end{table}

\section{Lists}
\subsection{Bullet Points}
\begin{itemize}
  \item First item
  \item Second item
  \item Third item with sub-items:
    \begin{itemize}
      \item Sub-item 1
      \item Sub-item 2
    \end{itemize}
\end{itemize}

\subsection{Numbered List}
\begin{enumerate}
  \item First step
  \item Second step
  \item Third step
\end{enumerate}

\section{Citations and References}
When you use the citation search feature, references will be automatically managed. For example, you might cite a paper like this: \cite{example2024}. All citations will be automatically added to your \texttt{cite.bib} file.

\section{Conclusion}
This template demonstrates basic \LaTeX{} features including mathematical formulas, tables, and lists. The bibliography system is now set up to work seamlessly with Bibby AI's citation management. Feel free to modify it for your needs!

\bibliographystyle{plain}
\bibliography{cite}

\end{document}\documentclass{article}
\usepackage[utf8]{inputenc}
\usepackage{booktabs}
\usepackage{amsmath}
\usepackage{natbib}

\title{My Document}
\author{Your Name}
\date{\today}

\begin{document}

\maketitle

\section{Introduction}
This is a sample document demonstrating various \LaTeX{} features including mathematical equations, tables, and formatting. You can cite papers using the citation search feature, and they will be automatically added to your \texttt{cite.bib} file.

\section{Mathematical Expressions}
Here are some example equations:

\subsection{Quadratic Formula}
The quadratic formula is given by:
\[x = \frac{-b \pm \sqrt{b^2 - 4ac}}{2a}\]

\subsection{Matrix Example}
A sample matrix:
\[
\begin{pmatrix}
a_{11} & a_{12} & a_{13} \\
a_{21} & a_{22} & a_{23} \\
a_{31} & a_{32} & a_{33}
\end{pmatrix}
\]

\section{Tables}
Below is a sample table showing data:

\begin{table}[h]
\centering
\begin{tabular}{@{}lrrr@{}}
\toprule
\textbf{Category} & \textbf{Value 1} & \textbf{Value 2} & \textbf{Total} \\
\midrule
Item A & 10.5 & 20.3 & 30.8 \\
Item B & 15.2 & 25.1 & 40.3 \\
Item C & 22.1 & 12.7 & 34.8 \\
\bottomrule
\end{tabular}
\caption{Sample Data Table}
\label{tab:sample}
\end{table}

\section{Lists}
\subsection{Bullet Points}
\begin{itemize}
  \item First item
  \item Second item
  \item Third item with sub-items:
    \begin{itemize}
      \item Sub-item 1
      \item Sub-item 2
    \end{itemize}
\end{itemize}

\subsection{Numbered List}
\begin{enumerate}
  \item First step
  \item Second step
  \item Third step
\end{enumerate}

\section{Citations and References}
When you use the citation search feature, references will be automatically managed. For example, you might cite a paper like this: \cite{example2024}. All citations will be automatically added to your \texttt{cite.bib} file.

\section{Conclusion}
This template demonstrates basic \LaTeX{} features including mathematical formulas, tables, and lists. The bibliography system is now set up to work seamlessly with Bibby AI's citation management. Feel free to modify it for your needs!

\bibliographystyle{plain}
\bibliography{cite}

\end{document}\documentclass{article}
\usepackage[utf8]{inputenc}
\usepackage{booktabs}
\usepackage{amsmath}
\usepackage{natbib}

\title{My Document}
\author{Your Name}
\date{\today}

\begin{document}

\maketitle

\section{Introduction}
This is a sample document demonstrating various \LaTeX{} features including mathematical equations, tables, and formatting. You can cite papers using the citation search feature, and they will be automatically added to your \texttt{cite.bib} file.

\section{Mathematical Expressions}
Here are some example equations:

\subsection{Quadratic Formula}
The quadratic formula is given by:
\[x = \frac{-b \pm \sqrt{b^2 - 4ac}}{2a}\]

\subsection{Matrix Example}
A sample matrix:
\[
\begin{pmatrix}
a_{11} & a_{12} & a_{13} \\
a_{21} & a_{22} & a_{23} \\
a_{31} & a_{32} & a_{33}
\end{pmatrix}
\]

\section{Tables}
Below is a sample table showing data:

\begin{table}[h]
\centering
\begin{tabular}{@{}lrrr@{}}
\toprule
\textbf{Category} & \textbf{Value 1} & \textbf{Value 2} & \textbf{Total} \\
\midrule
Item A & 10.5 & 20.3 & 30.8 \\
Item B & 15.2 & 25.1 & 40.3 \\
Item C & 22.1 & 12.7 & 34.8 \\
\bottomrule
\end{tabular}
\caption{Sample Data Table}
\label{tab:sample}
\end{table}

\section{Lists}
\subsection{Bullet Points}
\begin{itemize}
  \item First item
  \item Second item
  \item Third item with sub-items:
    \begin{itemize}
      \item Sub-item 1
      \item Sub-item 2
    \end{itemize}
\end{itemize}

\subsection{Numbered List}
\begin{enumerate}
  \item First step
  \item Second step
  \item Third step
\end{enumerate}

\section{Citations and References}
When you use the citation search feature, references will be automatically managed. For example, you might cite a paper like this: \cite{example2024}. All citations will be automatically added to your \texttt{cite.bib} file.

\section{Conclusion}
This template demonstrates basic \LaTeX{} features including mathematical formulas, tables, and lists. The bibliography system is now set up to work seamlessly with Bibby AI's citation management. Feel free to modify it for your needs!

\bibliographystyle{plain}
\bibliography{cite}

\end{document}\documentclass{article}
\usepackage[utf8]{inputenc}
\usepackage{booktabs}
\usepackage{amsmath}
\usepackage{natbib}

\title{My Document}
\author{Your Name}
\date{\today}

\begin{document}

\maketitle

\section{Introduction}
This is a sample document demonstrating various \LaTeX{} features including mathematical equations, tables, and formatting. You can cite papers using the citation search feature, and they will be automatically added to your \texttt{cite.bib} file.

\section{Mathematical Expressions}
Here are some example equations:

\subsection{Quadratic Formula}
The quadratic formula is given by:
\[x = \frac{-b \pm \sqrt{b^2 - 4ac}}{2a}\]

\subsection{Matrix Example}
A sample matrix:
\[
\begin{pmatrix}
a_{11} & a_{12} & a_{13} \\
a_{21} & a_{22} & a_{23} \\
a_{31} & a_{32} & a_{33}
\end{pmatrix}
\]

\section{Tables}
Below is a sample table showing data:

\begin{table}[h]
\centering
\begin{tabular}{@{}lrrr@{}}
\toprule
\textbf{Category} & \textbf{Value 1} & \textbf{Value 2} & \textbf{Total} \\
\midrule
Item A & 10.5 & 20.3 & 30.8 \\
Item B & 15.2 & 25.1 & 40.3 \\
Item C & 22.1 & 12.7 & 34.8 \\
\bottomrule
\end{tabular}
\caption{Sample Data Table}
\label{tab:sample}
\end{table}

\section{Lists}
\subsection{Bullet Points}
\begin{itemize}
  \item First item
  \item Second item
  \item Third item with sub-items:
    \begin{itemize}
      \item Sub-item 1
      \item Sub-item 2
    \end{itemize}
\end{itemize}

\subsection{Numbered List}
\begin{enumerate}
  \item First step
  \item Second step
  \item Third step
\end{enumerate}

\section{Citations and References}
When you use the citation search feature, references will be automatically managed. For example, you might cite a paper like this: \cite{example2024}. All citations will be automatically added to your \texttt{cite.bib} file.

\section{Conclusion}
This template demonstrates basic \LaTeX{} features including mathematical formulas, tables, and lists. The bibliography system is now set up to work seamlessly with Bibby AI's citation management. Feel free to modify it for your needs!

\bibliographystyle{plain}
\bibliography{cite}

\end{document}\documentclass{article}
\usepackage[utf8]{inputenc}
\usepackage{booktabs}
\usepackage{amsmath}
\usepackage{natbib}

\title{My Document}
\author{Your Name}
\date{\today}

\begin{document}

\maketitle

\section{Introduction}
This is a sample document demonstrating various \LaTeX{} features including mathematical equations, tables, and formatting. You can cite papers using the citation search feature, and they will be automatically added to your \texttt{cite.bib} file.

\section{Mathematical Expressions}
Here are some example equations:

\subsection{Quadratic Formula}
The quadratic formula is given by:
\[x = \frac{-b \pm \sqrt{b^2 - 4ac}}{2a}\]

\subsection{Matrix Example}
A sample matrix:
\[
\begin{pmatrix}
a_{11} & a_{12} & a_{13} \\
a_{21} & a_{22} & a_{23} \\
a_{31} & a_{32} & a_{33}
\end{pmatrix}
\]

\section{Tables}
Below is a sample table showing data:

\begin{table}[h]
\centering
\begin{tabular}{@{}lrrr@{}}
\toprule
\textbf{Category} & \textbf{Value 1} & \textbf{Value 2} & \textbf{Total} \\
\midrule
Item A & 10.5 & 20.3 & 30.8 \\
Item B & 15.2 & 25.1 & 40.3 \\
Item C & 22.1 & 12.7 & 34.8 \\
\bottomrule
\end{tabular}
\caption{Sample Data Table}
\label{tab:sample}
\end{table}

\section{Lists}
\subsection{Bullet Points}
\begin{itemize}
  \item First item
  \item Second item
  \item Third item with sub-items:
    \begin{itemize}
      \item Sub-item 1
      \item Sub-item 2
    \end{itemize}
\end{itemize}

\subsection{Numbered List}
\begin{enumerate}
  \item First step
  \item Second step
  \item Third step
\end{enumerate}

\section{Citations and References}
When you use the citation search feature, references will be automatically managed. For example, you might cite a paper like this: \cite{example2024}. All citations will be automatically added to your \texttt{cite.bib} file.

\section{Conclusion}
This template demonstrates basic \LaTeX{} features including mathematical formulas, tables, and lists. The bibliography system is now set up to work seamlessly with Bibby AI's citation management. Feel free to modify it for your needs!

\bibliographystyle{plain}
\bibliography{cite}

\end{document}\documentclass{article}
\usepackage[utf8]{inputenc}
\usepackage{booktabs}
\usepackage{amsmath}
\usepackage{natbib}

\title{My Document}
\author{Your Name}
\date{\today}

\begin{document}

\maketitle

\section{Introduction}
This is a sample document demonstrating various \LaTeX{} features including mathematical equations, tables, and formatting. You can cite papers using the citation search feature, and they will be automatically added to your \texttt{cite.bib} file.

\section{Mathematical Expressions}
Here are some example equations:

\subsection{Quadratic Formula}
The quadratic formula is given by:
\[x = \frac{-b \pm \sqrt{b^2 - 4ac}}{2a}\]

\subsection{Matrix Example}
A sample matrix:
\[
\begin{pmatrix}
a_{11} & a_{12} & a_{13} \\
a_{21} & a_{22} & a_{23} \\
a_{31} & a_{32} & a_{33}
\end{pmatrix}
\]

\section{Tables}
Below is a sample table showing data:

\begin{table}[h]
\centering
\begin{tabular}{@{}lrrr@{}}
\toprule
\textbf{Category} & \textbf{Value 1} & \textbf{Value 2} & \textbf{Total} \\
\midrule
Item A & 10.5 & 20.3 & 30.8 \\
Item B & 15.2 & 25.1 & 40.3 \\
Item C & 22.1 & 12.7 & 34.8 \\
\bottomrule
\end{tabular}
\caption{Sample Data Table}
\label{tab:sample}
\end{table}

\section{Lists}
\subsection{Bullet Points}
\begin{itemize}
  \item First item
  \item Second item
  \item Third item with sub-items:
    \begin{itemize}
      \item Sub-item 1
      \item Sub-item 2
    \end{itemize}
\end{itemize}

\subsection{Numbered List}
\begin{enumerate}
  \item First step
  \item Second step
  \item Third step
\end{enumerate}

\section{Citations and References}
When you use the citation search feature, references will be automatically managed. For example, you might cite a paper like this: \cite{example2024}. All citations will be automatically added to your \texttt{cite.bib} file.

\section{Conclusion}
This template demonstrates basic \LaTeX{} features including mathematical formulas, tables, and lists. The bibliography system is now set up to work seamlessly with Bibby AI's citation management. Feel free to modify it for your needs!

\bibliographystyle{plain}
\bibliography{cite}

\end{document}\documentclass{article}
\usepackage[utf8]{inputenc}
\usepackage{booktabs}
\usepackage{amsmath}
\usepackage{natbib}

\title{My Document}
\author{Your Name}
\date{\today}

\begin{document}

\maketitle

\section{Introduction}
This is a sample document demonstrating various \LaTeX{} features including mathematical equations, tables, and formatting. You can cite papers using the citation search feature, and they will be automatically added to your \texttt{cite.bib} file.

\section{Mathematical Expressions}
Here are some example equations:

\subsection{Quadratic Formula}
The quadratic formula is given by:
\[x = \frac{-b \pm \sqrt{b^2 - 4ac}}{2a}\]

\subsection{Matrix Example}
A sample matrix:
\[
\begin{pmatrix}
a_{11} & a_{12} & a_{13} \\
a_{21} & a_{22} & a_{23} \\
a_{31} & a_{32} & a_{33}
\end{pmatrix}
\]

\section{Tables}
Below is a sample table showing data:

\begin{table}[h]
\centering
\begin{tabular}{@{}lrrr@{}}
\toprule
\textbf{Category} & \textbf{Value 1} & \textbf{Value 2} & \textbf{Total} \\
\midrule
Item A & 10.5 & 20.3 & 30.8 \\
Item B & 15.2 & 25.1 & 40.3 \\
Item C & 22.1 & 12.7 & 34.8 \\
\bottomrule
\end{tabular}
\caption{Sample Data Table}
\label{tab:sample}
\end{table}

\section{Lists}
\subsection{Bullet Points}
\begin{itemize}
  \item First item
  \item Second item
  \item Third item with sub-items:
    \begin{itemize}
      \item Sub-item 1
      \item Sub-item 2
    \end{itemize}
\end{itemize}

\subsection{Numbered List}
\begin{enumerate}
  \item First step
  \item Second step
  \item Third step
\end{enumerate}

\section{Citations and References}
When you use the citation search feature, references will be automatically managed. For example, you might cite a paper like this: \cite{example2024}. All citations will be automatically added to your \texttt{cite.bib} file.

\section{Conclusion}
This template demonstrates basic \LaTeX{} features including mathematical formulas, tables, and lists. The bibliography system is now set up to work seamlessly with Bibby AI's citation management. Feel free to modify it for your needs!

\bibliographystyle{plain}
\bibliography{cite}

\end{document}\documentclass{article}
\usepackage[utf8]{inputenc}
\usepackage{booktabs}
\usepackage{amsmath}
\usepackage{natbib}

\title{My Document}
\author{Your Name}
\date{\today}

\begin{document}

\maketitle

\section{Introduction}
This is a sample document demonstrating various \LaTeX{} features including mathematical equations, tables, and formatting. You can cite papers using the citation search feature, and they will be automatically added to your \texttt{cite.bib} file.

\section{Mathematical Expressions}
Here are some example equations:

\subsection{Quadratic Formula}
The quadratic formula is given by:
\[x = \frac{-b \pm \sqrt{b^2 - 4ac}}{2a}\]

\subsection{Matrix Example}
A sample matrix:
\[
\begin{pmatrix}
a_{11} & a_{12} & a_{13} \\
a_{21} & a_{22} & a_{23} \\
a_{31} & a_{32} & a_{33}
\end{pmatrix}
\]

\section{Tables}
Below is a sample table showing data:

\begin{table}[h]
\centering
\begin{tabular}{@{}lrrr@{}}
\toprule
\textbf{Category} & \textbf{Value 1} & \textbf{Value 2} & \textbf{Total} \\
\midrule
Item A & 10.5 & 20.3 & 30.8 \\
Item B & 15.2 & 25.1 & 40.3 \\
Item C & 22.1 & 12.7 & 34.8 \\
\bottomrule
\end{tabular}
\caption{Sample Data Table}
\label{tab:sample}
\end{table}

\section{Lists}
\subsection{Bullet Points}
\begin{itemize}
  \item First item
  \item Second item
  \item Third item with sub-items:
    \begin{itemize}
      \item Sub-item 1
      \item Sub-item 2
    \end{itemize}
\end{itemize}

\subsection{Numbered List}
\begin{enumerate}
  \item First step
  \item Second step
  \item Third step
\end{enumerate}

\section{Citations and References}
When you use the citation search feature, references will be automatically managed. For example, you might cite a paper like this: \cite{example2024}. All citations will be automatically added to your \texttt{cite.bib} file.

\section{Conclusion}
This template demonstrates basic \LaTeX{} features including mathematical formulas, tables, and lists. The bibliography system is now set up to work seamlessly with Bibby AI's citation management. Feel free to modify it for your needs!

\bibliographystyle{plain}
\bibliography{cite}

\end{document}\documentclass{article}
\usepackage[utf8]{inputenc}
\usepackage{booktabs}
\usepackage{amsmath}
\usepackage{natbib}

\title{My Document}
\author{Your Name}
\date{\today}

\begin{document}

\maketitle

\section{Introduction}
This is a sample document demonstrating various \LaTeX{} features including mathematical equations, tables, and formatting. You can cite papers using the citation search feature, and they will be automatically added to your \texttt{cite.bib} file.

\section{Mathematical Expressions}
Here are some example equations:

\subsection{Quadratic Formula}
The quadratic formula is given by:
\[x = \frac{-b \pm \sqrt{b^2 - 4ac}}{2a}\]

\subsection{Matrix Example}
A sample matrix:
\[
\begin{pmatrix}
a_{11} & a_{12} & a_{13} \\
a_{21} & a_{22} & a_{23} \\
a_{31} & a_{32} & a_{33}
\end{pmatrix}
\]

\section{Tables}
Below is a sample table showing data:

\begin{table}[h]
\centering
\begin{tabular}{@{}lrrr@{}}
\toprule
\textbf{Category} & \textbf{Value 1} & \textbf{Value 2} & \textbf{Total} \\
\midrule
Item A & 10.5 & 20.3 & 30.8 \\
Item B & 15.2 & 25.1 & 40.3 \\
Item C & 22.1 & 12.7 & 34.8 \\
\bottomrule
\end{tabular}
\caption{Sample Data Table}
\label{tab:sample}
\end{table}

\section{Lists}
\subsection{Bullet Points}
\begin{itemize}
  \item First item
  \item Second item
  \item Third item with sub-items:
    \begin{itemize}
      \item Sub-item 1
      \item Sub-item 2
    \end{itemize}
\end{itemize}

\subsection{Numbered List}
\begin{enumerate}
  \item First step
  \item Second step
  \item Third step
\end{enumerate}

\section{Citations and References}
When you use the citation search feature, references will be automatically managed. For example, you might cite a paper like this: \cite{example2024}. All citations will be automatically added to your \texttt{cite.bib} file.

\section{Conclusion}
This template demonstrates basic \LaTeX{} features including mathematical formulas, tables, and lists. The bibliography system is now set up to work seamlessly with Bibby AI's citation management. Feel free to modify it for your needs!

\bibliographystyle{plain}
\bibliography{cite}

\end{document}\documentclass{article}
\usepackage[utf8]{inputenc}
\usepackage{booktabs}
\usepackage{amsmath}
\usepackage{natbib}

\title{My Document}
\author{Your Name}
\date{\today}

\begin{document}

\maketitle

\section{Introduction}
This is a sample document demonstrating various \LaTeX{} features including mathematical equations, tables, and formatting. You can cite papers using the citation search feature, and they will be automatically added to your \texttt{cite.bib} file.

\section{Mathematical Expressions}
Here are some example equations:

\subsection{Quadratic Formula}
The quadratic formula is given by:
\[x = \frac{-b \pm \sqrt{b^2 - 4ac}}{2a}\]

\subsection{Matrix Example}
A sample matrix:
\[
\begin{pmatrix}
a_{11} & a_{12} & a_{13} \\
a_{21} & a_{22} & a_{23} \\
a_{31} & a_{32} & a_{33}
\end{pmatrix}
\]

\section{Tables}
Below is a sample table showing data:

\begin{table}[h]
\centering
\begin{tabular}{@{}lrrr@{}}
\toprule
\textbf{Category} & \textbf{Value 1} & \textbf{Value 2} & \textbf{Total} \\
\midrule
Item A & 10.5 & 20.3 & 30.8 \\
Item B & 15.2 & 25.1 & 40.3 \\
Item C & 22.1 & 12.7 & 34.8 \\
\bottomrule
\end{tabular}
\caption{Sample Data Table}
\label{tab:sample}
\end{table}

\section{Lists}
\subsection{Bullet Points}
\begin{itemize}
  \item First item
  \item Second item
  \item Third item with sub-items:
    \begin{itemize}
      \item Sub-item 1
      \item Sub-item 2
    \end{itemize}
\end{itemize}

\subsection{Numbered List}
\begin{enumerate}
  \item First step
  \item Second step
  \item Third step
\end{enumerate}

\section{Citations and References}
When you use the citation search feature, references will be automatically managed. For example, you might cite a paper like this: \cite{example2024}. All citations will be automatically added to your \texttt{cite.bib} file.

\section{Conclusion}
This template demonstrates basic \LaTeX{} features including mathematical formulas, tables, and lists. The bibliography system is now set up to work seamlessly with Bibby AI's citation management. Feel free to modify it for your needs!

\bibliographystyle{plain}
\bibliography{cite}

\end{document}\documentclass{article}
\usepackage[utf8]{inputenc}
\usepackage{booktabs}
\usepackage{amsmath}
\usepackage{natbib}

\title{My Document}
\author{Your Name}
\date{\today}

\begin{document}

\maketitle

\section{Introduction}
This is a sample document demonstrating various \LaTeX{} features including mathematical equations, tables, and formatting. You can cite papers using the citation search feature, and they will be automatically added to your \texttt{cite.bib} file.

\section{Mathematical Expressions}
Here are some example equations:

\subsection{Quadratic Formula}
The quadratic formula is given by:
\[x = \frac{-b \pm \sqrt{b^2 - 4ac}}{2a}\]

\subsection{Matrix Example}
A sample matrix:
\[
\begin{pmatrix}
a_{11} & a_{12} & a_{13} \\
a_{21} & a_{22} & a_{23} \\
a_{31} & a_{32} & a_{33}
\end{pmatrix}
\]

\section{Tables}
Below is a sample table showing data:

\begin{table}[h]
\centering
\begin{tabular}{@{}lrrr@{}}
\toprule
\textbf{Category} & \textbf{Value 1} & \textbf{Value 2} & \textbf{Total} \\
\midrule
Item A & 10.5 & 20.3 & 30.8 \\
Item B & 15.2 & 25.1 & 40.3 \\
Item C & 22.1 & 12.7 & 34.8 \\
\bottomrule
\end{tabular}
\caption{Sample Data Table}
\label{tab:sample}
\end{table}

\section{Lists}
\subsection{Bullet Points}
\begin{itemize}
  \item First item
  \item Second item
  \item Third item with sub-items:
    \begin{itemize}
      \item Sub-item 1
      \item Sub-item 2
    \end{itemize}
\end{itemize}

\subsection{Numbered List}
\begin{enumerate}
  \item First step
  \item Second step
  \item Third step
\end{enumerate}

\section{Citations and References}
When you use the citation search feature, references will be automatically managed. For example, you might cite a paper like this: \cite{example2024}. All citations will be automatically added to your \texttt{cite.bib} file.

\section{Conclusion}
This template demonstrates basic \LaTeX{} features including mathematical formulas, tables, and lists. The bibliography system is now set up to work seamlessly with Bibby AI's citation management. Feel free to modify it for your needs!

\bibliographystyle{plain}
\bibliography{cite}

\end{document}\documentclass{article}
\usepackage[utf8]{inputenc}
\usepackage{booktabs}
\usepackage{amsmath}
\usepackage{natbib}

\title{My Document}
\author{Your Name}
\date{\today}

\begin{document}

\maketitle

\section{Introduction}
This is a sample document demonstrating various \LaTeX{} features including mathematical equations, tables, and formatting. You can cite papers using the citation search feature, and they will be automatically added to your \texttt{cite.bib} file.

\section{Mathematical Expressions}
Here are some example equations:

\subsection{Quadratic Formula}
The quadratic formula is given by:
\[x = \frac{-b \pm \sqrt{b^2 - 4ac}}{2a}\]

\subsection{Matrix Example}
A sample matrix:
\[
\begin{pmatrix}
a_{11} & a_{12} & a_{13} \\
a_{21} & a_{22} & a_{23} \\
a_{31} & a_{32} & a_{33}
\end{pmatrix}
\]

\section{Tables}
Below is a sample table showing data:

\begin{table}[h]
\centering
\begin{tabular}{@{}lrrr@{}}
\toprule
\textbf{Category} & \textbf{Value 1} & \textbf{Value 2} & \textbf{Total} \\
\midrule
Item A & 10.5 & 20.3 & 30.8 \\
Item B & 15.2 & 25.1 & 40.3 \\
Item C & 22.1 & 12.7 & 34.8 \\
\bottomrule
\end{tabular}
\caption{Sample Data Table}
\label{tab:sample}
\end{table}

\section{Lists}
\subsection{Bullet Points}
\begin{itemize}
  \item First item
  \item Second item
  \item Third item with sub-items:
    \begin{itemize}
      \item Sub-item 1
      \item Sub-item 2
    \end{itemize}
\end{itemize}

\subsection{Numbered List}
\begin{enumerate}
  \item First step
  \item Second step
  \item Third step
\end{enumerate}

\section{Citations and References}
When you use the citation search feature, references will be automatically managed. For example, you might cite a paper like this: \cite{example2024}. All citations will be automatically added to your \texttt{cite.bib} file.

\section{Conclusion}
This template demonstrates basic \LaTeX{} features including mathematical formulas, tables, and lists. The bibliography system is now set up to work seamlessly with Bibby AI's citation management. Feel free to modify it for your needs!

\bibliographystyle{plain}
\bibliography{cite}

\end{document}\documentclass{article}
\usepackage[utf8]{inputenc}
\usepackage{booktabs}
\usepackage{amsmath}
\usepackage{natbib}

\title{My Document}
\author{Your Name}
\date{\today}

\begin{document}

\maketitle

\section{Introduction}
This is a sample document demonstrating various \LaTeX{} features including mathematical equations, tables, and formatting. You can cite papers using the citation search feature, and they will be automatically added to your \texttt{cite.bib} file.

\section{Mathematical Expressions}
Here are some example equations:

\subsection{Quadratic Formula}
The quadratic formula is given by:
\[x = \frac{-b \pm \sqrt{b^2 - 4ac}}{2a}\]

\subsection{Matrix Example}
A sample matrix:
\[
\begin{pmatrix}
a_{11} & a_{12} & a_{13} \\
a_{21} & a_{22} & a_{23} \\
a_{31} & a_{32} & a_{33}
\end{pmatrix}
\]

\section{Tables}
Below is a sample table showing data:

\begin{table}[h]
\centering
\begin{tabular}{@{}lrrr@{}}
\toprule
\textbf{Category} & \textbf{Value 1} & \textbf{Value 2} & \textbf{Total} \\
\midrule
Item A & 10.5 & 20.3 & 30.8 \\
Item B & 15.2 & 25.1 & 40.3 \\
Item C & 22.1 & 12.7 & 34.8 \\
\bottomrule
\end{tabular}
\caption{Sample Data Table}
\label{tab:sample}
\end{table}

\section{Lists}
\subsection{Bullet Points}
\begin{itemize}
  \item First item
  \item Second item
  \item Third item with sub-items:
    \begin{itemize}
      \item Sub-item 1
      \item Sub-item 2
    \end{itemize}
\end{itemize}

\subsection{Numbered List}
\begin{enumerate}
  \item First step
  \item Second step
  \item Third step
\end{enumerate}

\section{Citations and References}
When you use the citation search feature, references will be automatically managed. For example, you might cite a paper like this: \cite{example2024}. All citations will be automatically added to your \texttt{cite.bib} file.

\section{Conclusion}
This template demonstrates basic \LaTeX{} features including mathematical formulas, tables, and lists. The bibliography system is now set up to work seamlessly with Bibby AI's citation management. Feel free to modify it for your needs!

\bibliographystyle{plain}
\bibliography{cite}

\end{document}\documentclass{article}
\usepackage[utf8]{inputenc}
\usepackage{booktabs}
\usepackage{amsmath}
\usepackage{natbib}

\title{My Document}
\author{Your Name}
\date{\today}

\begin{document}

\maketitle

\section{Introduction}
This is a sample document demonstrating various \LaTeX{} features including mathematical equations, tables, and formatting. You can cite papers using the citation search feature, and they will be automatically added to your \texttt{cite.bib} file.

\section{Mathematical Expressions}
Here are some example equations:

\subsection{Quadratic Formula}
The quadratic formula is given by:
\[x = \frac{-b \pm \sqrt{b^2 - 4ac}}{2a}\]

\subsection{Matrix Example}
A sample matrix:
\[
\begin{pmatrix}
a_{11} & a_{12} & a_{13} \\
a_{21} & a_{22} & a_{23} \\
a_{31} & a_{32} & a_{33}
\end{pmatrix}
\]

\section{Tables}
Below is a sample table showing data:

\begin{table}[h]
\centering
\begin{tabular}{@{}lrrr@{}}
\toprule
\textbf{Category} & \textbf{Value 1} & \textbf{Value 2} & \textbf{Total} \\
\midrule
Item A & 10.5 & 20.3 & 30.8 \\
Item B & 15.2 & 25.1 & 40.3 \\
Item C & 22.1 & 12.7 & 34.8 \\
\bottomrule
\end{tabular}
\caption{Sample Data Table}
\label{tab:sample}
\end{table}

\section{Lists}
\subsection{Bullet Points}
\begin{itemize}
  \item First item
  \item Second item
  \item Third item with sub-items:
    \begin{itemize}
      \item Sub-item 1
      \item Sub-item 2
    \end{itemize}
\end{itemize}

\subsection{Numbered List}
\begin{enumerate}
  \item First step
  \item Second step
  \item Third step
\end{enumerate}

\section{Citations and References}
When you use the citation search feature, references will be automatically managed. For example, you might cite a paper like this: \cite{example2024}. All citations will be automatically added to your \texttt{cite.bib} file.

\section{Conclusion}
This template demonstrates basic \LaTeX{} features including mathematical formulas, tables, and lists. The bibliography system is now set up to work seamlessly with Bibby AI's citation management. Feel free to modify it for your needs!

\bibliographystyle{plain}
\bibliography{cite}

\end{document}\documentclass{article}
\usepackage[utf8]{inputenc}
\usepackage{booktabs}
\usepackage{amsmath}
\usepackage{natbib}

\title{My Document}
\author{Your Name}
\date{\today}

\begin{document}

\maketitle

\section{Introduction}
This is a sample document demonstrating various \LaTeX{} features including mathematical equations, tables, and formatting. You can cite papers using the citation search feature, and they will be automatically added to your \texttt{cite.bib} file.

\section{Mathematical Expressions}
Here are some example equations:

\subsection{Quadratic Formula}
The quadratic formula is given by:
\[x = \frac{-b \pm \sqrt{b^2 - 4ac}}{2a}\]

\subsection{Matrix Example}
A sample matrix:
\[
\begin{pmatrix}
a_{11} & a_{12} & a_{13} \\
a_{21} & a_{22} & a_{23} \\
a_{31} & a_{32} & a_{33}
\end{pmatrix}
\]

\section{Tables}
Below is a sample table showing data:

\begin{table}[h]
\centering
\begin{tabular}{@{}lrrr@{}}
\toprule
\textbf{Category} & \textbf{Value 1} & \textbf{Value 2} & \textbf{Total} \\
\midrule
Item A & 10.5 & 20.3 & 30.8 \\
Item B & 15.2 & 25.1 & 40.3 \\
Item C & 22.1 & 12.7 & 34.8 \\
\bottomrule
\end{tabular}
\caption{Sample Data Table}
\label{tab:sample}
\end{table}

\section{Lists}
\subsection{Bullet Points}
\begin{itemize}
  \item First item
  \item Second item
  \item Third item with sub-items:
    \begin{itemize}
      \item Sub-item 1
      \item Sub-item 2
    \end{itemize}
\end{itemize}

\subsection{Numbered List}
\begin{enumerate}
  \item First step
  \item Second step
  \item Third step
\end{enumerate}

\section{Citations and References}
When you use the citation search feature, references will be automatically managed. For example, you might cite a paper like this: \cite{example2024}. All citations will be automatically added to your \texttt{cite.bib} file.

\section{Conclusion}
This template demonstrates basic \LaTeX{} features including mathematical formulas, tables, and lists. The bibliography system is now set up to work seamlessly with Bibby AI's citation management. Feel free to modify it for your needs!

\bibliographystyle{plain}
\bibliography{cite}

\end{document}\documentclass{article}
\usepackage[utf8]{inputenc}
\usepackage{booktabs}
\usepackage{amsmath}
\usepackage{natbib}

\title{My Document}
\author{Your Name}
\date{\today}

\begin{document}

\maketitle

\section{Introduction}
This is a sample document demonstrating various \LaTeX{} features including mathematical equations, tables, and formatting. You can cite papers using the citation search feature, and they will be automatically added to your \texttt{cite.bib} file.

\section{Mathematical Expressions}
Here are some example equations:

\subsection{Quadratic Formula}
The quadratic formula is given by:
\[x = \frac{-b \pm \sqrt{b^2 - 4ac}}{2a}\]

\subsection{Matrix Example}
A sample matrix:
\[
\begin{pmatrix}
a_{11} & a_{12} & a_{13} \\
a_{21} & a_{22} & a_{23} \\
a_{31} & a_{32} & a_{33}
\end{pmatrix}
\]

\section{Tables}
Below is a sample table showing data:

\begin{table}[h]
\centering
\begin{tabular}{@{}lrrr@{}}
\toprule
\textbf{Category} & \textbf{Value 1} & \textbf{Value 2} & \textbf{Total} \\
\midrule
Item A & 10.5 & 20.3 & 30.8 \\
Item B & 15.2 & 25.1 & 40.3 \\
Item C & 22.1 & 12.7 & 34.8 \\
\bottomrule
\end{tabular}
\caption{Sample Data Table}
\label{tab:sample}
\end{table}

\section{Lists}
\subsection{Bullet Points}
\begin{itemize}
  \item First item
  \item Second item
  \item Third item with sub-items:
    \begin{itemize}
      \item Sub-item 1
      \item Sub-item 2
    \end{itemize}
\end{itemize}

\subsection{Numbered List}
\begin{enumerate}
  \item First step
  \item Second step
  \item Third step
\end{enumerate}

\section{Citations and References}
When you use the citation search feature, references will be automatically managed. For example, you might cite a paper like this: \cite{example2024}. All citations will be automatically added to your \texttt{cite.bib} file.

\section{Conclusion}
This template demonstrates basic \LaTeX{} features including mathematical formulas, tables, and lists. The bibliography system is now set up to work seamlessly with Bibby AI's citation management. Feel free to modify it for your needs!

\bibliographystyle{plain}
\bibliography{cite}

\end{document}\documentclass{article}
\usepackage[utf8]{inputenc}
\usepackage{booktabs}
\usepackage{amsmath}
\usepackage{natbib}

\title{My Document}
\author{Your Name}
\date{\today}

\begin{document}

\maketitle

\section{Introduction}
This is a sample document demonstrating various \LaTeX{} features including mathematical equations, tables, and formatting. You can cite papers using the citation search feature, and they will be automatically added to your \texttt{cite.bib} file.

\section{Mathematical Expressions}
Here are some example equations:

\subsection{Quadratic Formula}
The quadratic formula is given by:
\[x = \frac{-b \pm \sqrt{b^2 - 4ac}}{2a}\]

\subsection{Matrix Example}
A sample matrix:
\[
\begin{pmatrix}
a_{11} & a_{12} & a_{13} \\
a_{21} & a_{22} & a_{23} \\
a_{31} & a_{32} & a_{33}
\end{pmatrix}
\]

\section{Tables}
Below is a sample table showing data:

\begin{table}[h]
\centering
\begin{tabular}{@{}lrrr@{}}
\toprule
\textbf{Category} & \textbf{Value 1} & \textbf{Value 2} & \textbf{Total} \\
\midrule
Item A & 10.5 & 20.3 & 30.8 \\
Item B & 15.2 & 25.1 & 40.3 \\
Item C & 22.1 & 12.7 & 34.8 \\
\bottomrule
\end{tabular}
\caption{Sample Data Table}
\label{tab:sample}
\end{table}

\section{Lists}
\subsection{Bullet Points}
\begin{itemize}
  \item First item
  \item Second item
  \item Third item with sub-items:
    \begin{itemize}
      \item Sub-item 1
      \item Sub-item 2
    \end{itemize}
\end{itemize}

\subsection{Numbered List}
\begin{enumerate}
  \item First step
  \item Second step
  \item Third step
\end{enumerate}

\section{Citations and References}
When you use the citation search feature, references will be automatically managed. For example, you might cite a paper like this: \cite{example2024}. All citations will be automatically added to your \texttt{cite.bib} file.

\section{Conclusion}
This template demonstrates basic \LaTeX{} features including mathematical formulas, tables, and lists. The bibliography system is now set up to work seamlessly with Bibby AI's citation management. Feel free to modify it for your needs!

\bibliographystyle{plain}
\bibliography{cite}

\end{document}\documentclass{article}
\usepackage[utf8]{inputenc}
\usepackage{booktabs}
\usepackage{amsmath}
\usepackage{natbib}

\title{My Document}
\author{Your Name}
\date{\today}

\begin{document}

\maketitle

\section{Introduction}
This is a sample document demonstrating various \LaTeX{} features including mathematical equations, tables, and formatting. You can cite papers using the citation search feature, and they will be automatically added to your \texttt{cite.bib} file.

\section{Mathematical Expressions}
Here are some example equations:

\subsection{Quadratic Formula}
The quadratic formula is given by:
\[x = \frac{-b \pm \sqrt{b^2 - 4ac}}{2a}\]

\subsection{Matrix Example}
A sample matrix:
\[
\begin{pmatrix}
a_{11} & a_{12} & a_{13} \\
a_{21} & a_{22} & a_{23} \\
a_{31} & a_{32} & a_{33}
\end{pmatrix}
\]

\section{Tables}
Below is a sample table showing data:

\begin{table}[h]
\centering
\begin{tabular}{@{}lrrr@{}}
\toprule
\textbf{Category} & \textbf{Value 1} & \textbf{Value 2} & \textbf{Total} \\
\midrule
Item A & 10.5 & 20.3 & 30.8 \\
Item B & 15.2 & 25.1 & 40.3 \\
Item C & 22.1 & 12.7 & 34.8 \\
\bottomrule
\end{tabular}
\caption{Sample Data Table}
\label{tab:sample}
\end{table}

\section{Lists}
\subsection{Bullet Points}
\begin{itemize}
  \item First item
  \item Second item
  \item Third item with sub-items:
    \begin{itemize}
      \item Sub-item 1
      \item Sub-item 2
    \end{itemize}
\end{itemize}

\subsection{Numbered List}
\begin{enumerate}
  \item First step
  \item Second step
  \item Third step
\end{enumerate}

\section{Citations and References}
When you use the citation search feature, references will be automatically managed. For example, you might cite a paper like this: \cite{example2024}. All citations will be automatically added to your \texttt{cite.bib} file.

\section{Conclusion}
This template demonstrates basic \LaTeX{} features including mathematical formulas, tables, and lists. The bibliography system is now set up to work seamlessly with Bibby AI's citation management. Feel free to modify it for your needs!

\bibliographystyle{plain}
\bibliography{cite}

\end{document}\documentclass{article}
\usepackage[utf8]{inputenc}
\usepackage{booktabs}
\usepackage{amsmath}
\usepackage{natbib}

\title{My Document}
\author{Your Name}
\date{\today}

\begin{document}

\maketitle

\section{Introduction}
This is a sample document demonstrating various \LaTeX{} features including mathematical equations, tables, and formatting. You can cite papers using the citation search feature, and they will be automatically added to your \texttt{cite.bib} file.

\section{Mathematical Expressions}
Here are some example equations:

\subsection{Quadratic Formula}
The quadratic formula is given by:
\[x = \frac{-b \pm \sqrt{b^2 - 4ac}}{2a}\]

\subsection{Matrix Example}
A sample matrix:
\[
\begin{pmatrix}
a_{11} & a_{12} & a_{13} \\
a_{21} & a_{22} & a_{23} \\
a_{31} & a_{32} & a_{33}
\end{pmatrix}
\]

\section{Tables}
Below is a sample table showing data:

\begin{table}[h]
\centering
\begin{tabular}{@{}lrrr@{}}
\toprule
\textbf{Category} & \textbf{Value 1} & \textbf{Value 2} & \textbf{Total} \\
\midrule
Item A & 10.5 & 20.3 & 30.8 \\
Item B & 15.2 & 25.1 & 40.3 \\
Item C & 22.1 & 12.7 & 34.8 \\
\bottomrule
\end{tabular}
\caption{Sample Data Table}
\label{tab:sample}
\end{table}

\section{Lists}
\subsection{Bullet Points}
\begin{itemize}
  \item First item
  \item Second item
  \item Third item with sub-items:
    \begin{itemize}
      \item Sub-item 1
      \item Sub-item 2
    \end{itemize}
\end{itemize}

\subsection{Numbered List}
\begin{enumerate}
  \item First step
  \item Second step
  \item Third step
\end{enumerate}

\section{Citations and References}
When you use the citation search feature, references will be automatically managed. For example, you might cite a paper like this: \cite{example2024}. All citations will be automatically added to your \texttt{cite.bib} file.

\section{Conclusion}
This template demonstrates basic \LaTeX{} features including mathematical formulas, tables, and lists. The bibliography system is now set up to work seamlessly with Bibby AI's citation management. Feel free to modify it for your needs!

\bibliographystyle{plain}
\bibliography{cite}

\end{document}\documentclass{article}
\usepackage[utf8]{inputenc}
\usepackage{booktabs}
\usepackage{amsmath}
\usepackage{natbib}

\title{My Document}
\author{Your Name}
\date{\today}

\begin{document}

\maketitle

\section{Introduction}
This is a sample document demonstrating various \LaTeX{} features including mathematical equations, tables, and formatting. You can cite papers using the citation search feature, and they will be automatically added to your \texttt{cite.bib} file.

\section{Mathematical Expressions}
Here are some example equations:

\subsection{Quadratic Formula}
The quadratic formula is given by:
\[x = \frac{-b \pm \sqrt{b^2 - 4ac}}{2a}\]

\subsection{Matrix Example}
A sample matrix:
\[
\begin{pmatrix}
a_{11} & a_{12} & a_{13} \\
a_{21} & a_{22} & a_{23} \\
a_{31} & a_{32} & a_{33}
\end{pmatrix}
\]

\section{Tables}
Below is a sample table showing data:

\begin{table}[h]
\centering
\begin{tabular}{@{}lrrr@{}}
\toprule
\textbf{Category} & \textbf{Value 1} & \textbf{Value 2} & \textbf{Total} \\
\midrule
Item A & 10.5 & 20.3 & 30.8 \\
Item B & 15.2 & 25.1 & 40.3 \\
Item C & 22.1 & 12.7 & 34.8 \\
\bottomrule
\end{tabular}
\caption{Sample Data Table}
\label{tab:sample}
\end{table}

\section{Lists}
\subsection{Bullet Points}
\begin{itemize}
  \item First item
  \item Second item
  \item Third item with sub-items:
    \begin{itemize}
      \item Sub-item 1
      \item Sub-item 2
    \end{itemize}
\end{itemize}

\subsection{Numbered List}
\begin{enumerate}
  \item First step
  \item Second step
  \item Third step
\end{enumerate}

\section{Citations and References}
When you use the citation search feature, references will be automatically managed. For example, you might cite a paper like this: \cite{example2024}. All citations will be automatically added to your \texttt{cite.bib} file.

\section{Conclusion}
This template demonstrates basic \LaTeX{} features including mathematical formulas, tables, and lists. The bibliography system is now set up to work seamlessly with Bibby AI's citation management. Feel free to modify it for your needs!

\bibliographystyle{plain}
\bibliography{cite}

\end{document}\documentclass{article}
\usepackage[utf8]{inputenc}
\usepackage{booktabs}
\usepackage{amsmath}
\usepackage{natbib}

\title{My Document}
\author{Your Name}
\date{\today}

\begin{document}

\maketitle

\section{Introduction}
This is a sample document demonstrating various \LaTeX{} features including mathematical equations, tables, and formatting. You can cite papers using the citation search feature, and they will be automatically added to your \texttt{cite.bib} file.

\section{Mathematical Expressions}
Here are some example equations:

\subsection{Quadratic Formula}
The quadratic formula is given by:
\[x = \frac{-b \pm \sqrt{b^2 - 4ac}}{2a}\]

\subsection{Matrix Example}
A sample matrix:
\[
\begin{pmatrix}
a_{11} & a_{12} & a_{13} \\
a_{21} & a_{22} & a_{23} \\
a_{31} & a_{32} & a_{33}
\end{pmatrix}
\]

\section{Tables}
Below is a sample table showing data:

\begin{table}[h]
\centering
\begin{tabular}{@{}lrrr@{}}
\toprule
\textbf{Category} & \textbf{Value 1} & \textbf{Value 2} & \textbf{Total} \\
\midrule
Item A & 10.5 & 20.3 & 30.8 \\
Item B & 15.2 & 25.1 & 40.3 \\
Item C & 22.1 & 12.7 & 34.8 \\
\bottomrule
\end{tabular}
\caption{Sample Data Table}
\label{tab:sample}
\end{table}

\section{Lists}
\subsection{Bullet Points}
\begin{itemize}
  \item First item
  \item Second item
  \item Third item with sub-items:
    \begin{itemize}
      \item Sub-item 1
      \item Sub-item 2
    \end{itemize}
\end{itemize}

\subsection{Numbered List}
\begin{enumerate}
  \item First step
  \item Second step
  \item Third step
\end{enumerate}

\section{Citations and References}
When you use the citation search feature, references will be automatically managed. For example, you might cite a paper like this: \cite{example2024}. All citations will be automatically added to your \texttt{cite.bib} file.

\section{Conclusion}
This template demonstrates basic \LaTeX{} features including mathematical formulas, tables, and lists. The bibliography system is now set up to work seamlessly with Bibby AI's citation management. Feel free to modify it for your needs!

\bibliographystyle{plain}
\bibliography{cite}

\end{document}\documentclass{article}
\usepackage[utf8]{inputenc}
\usepackage{booktabs}
\usepackage{amsmath}
\usepackage{natbib}

\title{My Document}
\author{Your Name}
\date{\today}

\begin{document}

\maketitle

\section{Introduction}
This is a sample document demonstrating various \LaTeX{} features including mathematical equations, tables, and formatting. You can cite papers using the citation search feature, and they will be automatically added to your \texttt{cite.bib} file.

\section{Mathematical Expressions}
Here are some example equations:

\subsection{Quadratic Formula}
The quadratic formula is given by:
\[x = \frac{-b \pm \sqrt{b^2 - 4ac}}{2a}\]

\subsection{Matrix Example}
A sample matrix:
\[
\begin{pmatrix}
a_{11} & a_{12} & a_{13} \\
a_{21} & a_{22} & a_{23} \\
a_{31} & a_{32} & a_{33}
\end{pmatrix}
\]

\section{Tables}
Below is a sample table showing data:

\begin{table}[h]
\centering
\begin{tabular}{@{}lrrr@{}}
\toprule
\textbf{Category} & \textbf{Value 1} & \textbf{Value 2} & \textbf{Total} \\
\midrule
Item A & 10.5 & 20.3 & 30.8 \\
Item B & 15.2 & 25.1 & 40.3 \\
Item C & 22.1 & 12.7 & 34.8 \\
\bottomrule
\end{tabular}
\caption{Sample Data Table}
\label{tab:sample}
\end{table}

\section{Lists}
\subsection{Bullet Points}
\begin{itemize}
  \item First item
  \item Second item
  \item Third item with sub-items:
    \begin{itemize}
      \item Sub-item 1
      \item Sub-item 2
    \end{itemize}
\end{itemize}

\subsection{Numbered List}
\begin{enumerate}
  \item First step
  \item Second step
  \item Third step
\end{enumerate}

\section{Citations and References}
When you use the citation search feature, references will be automatically managed. For example, you might cite a paper like this: \cite{example2024}. All citations will be automatically added to your \texttt{cite.bib} file.

\section{Conclusion}
This template demonstrates basic \LaTeX{} features including mathematical formulas, tables, and lists. The bibliography system is now set up to work seamlessly with Bibby AI's citation management. Feel free to modify it for your needs!

\bibliographystyle{plain}
\bibliography{cite}

\end{document}\documentclass{article}
\usepackage[utf8]{inputenc}
\usepackage{booktabs}
\usepackage{amsmath}
\usepackage{natbib}

\title{My Document}
\author{Your Name}
\date{\today}

\begin{document}

\maketitle

\section{Introduction}
This is a sample document demonstrating various \LaTeX{} features including mathematical equations, tables, and formatting. You can cite papers using the citation search feature, and they will be automatically added to your \texttt{cite.bib} file.

\section{Mathematical Expressions}
Here are some example equations:

\subsection{Quadratic Formula}
The quadratic formula is given by:
\[x = \frac{-b \pm \sqrt{b^2 - 4ac}}{2a}\]

\subsection{Matrix Example}
A sample matrix:
\[
\begin{pmatrix}
a_{11} & a_{12} & a_{13} \\
a_{21} & a_{22} & a_{23} \\
a_{31} & a_{32} & a_{33}
\end{pmatrix}
\]

\section{Tables}
Below is a sample table showing data:

\begin{table}[h]
\centering
\begin{tabular}{@{}lrrr@{}}
\toprule
\textbf{Category} & \textbf{Value 1} & \textbf{Value 2} & \textbf{Total} \\
\midrule
Item A & 10.5 & 20.3 & 30.8 \\
Item B & 15.2 & 25.1 & 40.3 \\
Item C & 22.1 & 12.7 & 34.8 \\
\bottomrule
\end{tabular}
\caption{Sample Data Table}
\label{tab:sample}
\end{table}

\section{Lists}
\subsection{Bullet Points}
\begin{itemize}
  \item First item
  \item Second item
  \item Third item with sub-items:
    \begin{itemize}
      \item Sub-item 1
      \item Sub-item 2
    \end{itemize}
\end{itemize}

\subsection{Numbered List}
\begin{enumerate}
  \item First step
  \item Second step
  \item Third step
\end{enumerate}

\section{Citations and References}
When you use the citation search feature, references will be automatically managed. For example, you might cite a paper like this: \cite{example2024}. All citations will be automatically added to your \texttt{cite.bib} file.

\section{Conclusion}
This template demonstrates basic \LaTeX{} features including mathematical formulas, tables, and lists. The bibliography system is now set up to work seamlessly with Bibby AI's citation management. Feel free to modify it for your needs!

\bibliographystyle{plain}
\bibliography{cite}

\end{document}\documentclass{article}
\usepackage[utf8]{inputenc}
\usepackage{booktabs}
\usepackage{amsmath}
\usepackage{natbib}

\title{My Document}
\author{Your Name}
\date{\today}

\begin{document}

\maketitle

\section{Introduction}
This is a sample document demonstrating various \LaTeX{} features including mathematical equations, tables, and formatting. You can cite papers using the citation search feature, and they will be automatically added to your \texttt{cite.bib} file.

\section{Mathematical Expressions}
Here are some example equations:

\subsection{Quadratic Formula}
The quadratic formula is given by:
\[x = \frac{-b \pm \sqrt{b^2 - 4ac}}{2a}\]

\subsection{Matrix Example}
A sample matrix:
\[
\begin{pmatrix}
a_{11} & a_{12} & a_{13} \\
a_{21} & a_{22} & a_{23} \\
a_{31} & a_{32} & a_{33}
\end{pmatrix}
\]

\section{Tables}
Below is a sample table showing data:

\begin{table}[h]
\centering
\begin{tabular}{@{}lrrr@{}}
\toprule
\textbf{Category} & \textbf{Value 1} & \textbf{Value 2} & \textbf{Total} \\
\midrule
Item A & 10.5 & 20.3 & 30.8 \\
Item B & 15.2 & 25.1 & 40.3 \\
Item C & 22.1 & 12.7 & 34.8 \\
\bottomrule
\end{tabular}
\caption{Sample Data Table}
\label{tab:sample}
\end{table}

\section{Lists}
\subsection{Bullet Points}
\begin{itemize}
  \item First item
  \item Second item
  \item Third item with sub-items:
    \begin{itemize}
      \item Sub-item 1
      \item Sub-item 2
    \end{itemize}
\end{itemize}

\subsection{Numbered List}
\begin{enumerate}
  \item First step
  \item Second step
  \item Third step
\end{enumerate}

\section{Citations and References}
When you use the citation search feature, references will be automatically managed. For example, you might cite a paper like this: \cite{example2024}. All citations will be automatically added to your \texttt{cite.bib} file.

\section{Conclusion}
This template demonstrates basic \LaTeX{} features including mathematical formulas, tables, and lists. The bibliography system is now set up to work seamlessly with Bibby AI's citation management. Feel free to modify it for your needs!

\bibliographystyle{plain}
\bibliography{cite}

\end{document}\documentclass{article}
\usepackage[utf8]{inputenc}
\usepackage{booktabs}
\usepackage{amsmath}
\usepackage{natbib}

\title{My Document}
\author{Your Name}
\date{\today}

\begin{document}

\maketitle

\section{Introduction}
This is a sample document demonstrating various \LaTeX{} features including mathematical equations, tables, and formatting. You can cite papers using the citation search feature, and they will be automatically added to your \texttt{cite.bib} file.

\section{Mathematical Expressions}
Here are some example equations:

\subsection{Quadratic Formula}
The quadratic formula is given by:
\[x = \frac{-b \pm \sqrt{b^2 - 4ac}}{2a}\]

\subsection{Matrix Example}
A sample matrix:
\[
\begin{pmatrix}
a_{11} & a_{12} & a_{13} \\
a_{21} & a_{22} & a_{23} \\
a_{31} & a_{32} & a_{33}
\end{pmatrix}
\]

\section{Tables}
Below is a sample table showing data:

\begin{table}[h]
\centering
\begin{tabular}{@{}lrrr@{}}
\toprule
\textbf{Category} & \textbf{Value 1} & \textbf{Value 2} & \textbf{Total} \\
\midrule
Item A & 10.5 & 20.3 & 30.8 \\
Item B & 15.2 & 25.1 & 40.3 \\
Item C & 22.1 & 12.7 & 34.8 \\
\bottomrule
\end{tabular}
\caption{Sample Data Table}
\label{tab:sample}
\end{table}

\section{Lists}
\subsection{Bullet Points}
\begin{itemize}
  \item First item
  \item Second item
  \item Third item with sub-items:
    \begin{itemize}
      \item Sub-item 1
      \item Sub-item 2
    \end{itemize}
\end{itemize}

\subsection{Numbered List}
\begin{enumerate}
  \item First step
  \item Second step
  \item Third step
\end{enumerate}

\section{Citations and References}
When you use the citation search feature, references will be automatically managed. For example, you might cite a paper like this: \cite{example2024}. All citations will be automatically added to your \texttt{cite.bib} file.

\section{Conclusion}
This template demonstrates basic \LaTeX{} features including mathematical formulas, tables, and lists. The bibliography system is now set up to work seamlessly with Bibby AI's citation management. Feel free to modify it for your needs!

\bibliographystyle{plain}
\bibliography{cite}

\end{document}\documentclass{article}
\usepackage[utf8]{inputenc}
\usepackage{booktabs}
\usepackage{amsmath}
\usepackage{natbib}

\title{My Document}
\author{Your Name}
\date{\today}

\begin{document}

\maketitle

\section{Introduction}
This is a sample document demonstrating various \LaTeX{} features including mathematical equations, tables, and formatting. You can cite papers using the citation search feature, and they will be automatically added to your \texttt{cite.bib} file.

\section{Mathematical Expressions}
Here are some example equations:

\subsection{Quadratic Formula}
The quadratic formula is given by:
\[x = \frac{-b \pm \sqrt{b^2 - 4ac}}{2a}\]

\subsection{Matrix Example}
A sample matrix:
\[
\begin{pmatrix}
a_{11} & a_{12} & a_{13} \\
a_{21} & a_{22} & a_{23} \\
a_{31} & a_{32} & a_{33}
\end{pmatrix}
\]

\section{Tables}
Below is a sample table showing data:

\begin{table}[h]
\centering
\begin{tabular}{@{}lrrr@{}}
\toprule
\textbf{Category} & \textbf{Value 1} & \textbf{Value 2} & \textbf{Total} \\
\midrule
Item A & 10.5 & 20.3 & 30.8 \\
Item B & 15.2 & 25.1 & 40.3 \\
Item C & 22.1 & 12.7 & 34.8 \\
\bottomrule
\end{tabular}
\caption{Sample Data Table}
\label{tab:sample}
\end{table}

\section{Lists}
\subsection{Bullet Points}
\begin{itemize}
  \item First item
  \item Second item
  \item Third item with sub-items:
    \begin{itemize}
      \item Sub-item 1
      \item Sub-item 2
    \end{itemize}
\end{itemize}

\subsection{Numbered List}
\begin{enumerate}
  \item First step
  \item Second step
  \item Third step
\end{enumerate}

\section{Citations and References}
When you use the citation search feature, references will be automatically managed. For example, you might cite a paper like this: \cite{example2024}. All citations will be automatically added to your \texttt{cite.bib} file.

\section{Conclusion}
This template demonstrates basic \LaTeX{} features including mathematical formulas, tables, and lists. The bibliography system is now set up to work seamlessly with Bibby AI's citation management. Feel free to modify it for your needs!

\bibliographystyle{plain}
\bibliography{cite}

\end{document}\documentclass{article}
\usepackage[utf8]{inputenc}
\usepackage{booktabs}
\usepackage{amsmath}
\usepackage{natbib}

\title{My Document}
\author{Your Name}
\date{\today}

\begin{document}

\maketitle

\section{Introduction}
This is a sample document demonstrating various \LaTeX{} features including mathematical equations, tables, and formatting. You can cite papers using the citation search feature, and they will be automatically added to your \texttt{cite.bib} file.

\section{Mathematical Expressions}
Here are some example equations:

\subsection{Quadratic Formula}
The quadratic formula is given by:
\[x = \frac{-b \pm \sqrt{b^2 - 4ac}}{2a}\]

\subsection{Matrix Example}
A sample matrix:
\[
\begin{pmatrix}
a_{11} & a_{12} & a_{13} \\
a_{21} & a_{22} & a_{23} \\
a_{31} & a_{32} & a_{33}
\end{pmatrix}
\]

\section{Tables}
Below is a sample table showing data:

\begin{table}[h]
\centering
\begin{tabular}{@{}lrrr@{}}
\toprule
\textbf{Category} & \textbf{Value 1} & \textbf{Value 2} & \textbf{Total} \\
\midrule
Item A & 10.5 & 20.3 & 30.8 \\
Item B & 15.2 & 25.1 & 40.3 \\
Item C & 22.1 & 12.7 & 34.8 \\
\bottomrule
\end{tabular}
\caption{Sample Data Table}
\label{tab:sample}
\end{table}

\section{Lists}
\subsection{Bullet Points}
\begin{itemize}
  \item First item
  \item Second item
  \item Third item with sub-items:
    \begin{itemize}
      \item Sub-item 1
      \item Sub-item 2
    \end{itemize}
\end{itemize}

\subsection{Numbered List}
\begin{enumerate}
  \item First step
  \item Second step
  \item Third step
\end{enumerate}

\section{Citations and References}
When you use the citation search feature, references will be automatically managed. For example, you might cite a paper like this: \cite{example2024}. All citations will be automatically added to your \texttt{cite.bib} file.

\section{Conclusion}
This template demonstrates basic \LaTeX{} features including mathematical formulas, tables, and lists. The bibliography system is now set up to work seamlessly with Bibby AI's citation management. Feel free to modify it for your needs!

\bibliographystyle{plain}
\bibliography{cite}

\end{document}\documentclass{article}
\usepackage[utf8]{inputenc}
\usepackage{booktabs}
\usepackage{amsmath}
\usepackage{natbib}

\title{My Document}
\author{Your Name}
\date{\today}

\begin{document}

\maketitle

\section{Introduction}
This is a sample document demonstrating various \LaTeX{} features including mathematical equations, tables, and formatting. You can cite papers using the citation search feature, and they will be automatically added to your \texttt{cite.bib} file.

\section{Mathematical Expressions}
Here are some example equations:

\subsection{Quadratic Formula}
The quadratic formula is given by:
\[x = \frac{-b \pm \sqrt{b^2 - 4ac}}{2a}\]

\subsection{Matrix Example}
A sample matrix:
\[
\begin{pmatrix}
a_{11} & a_{12} & a_{13} \\
a_{21} & a_{22} & a_{23} \\
a_{31} & a_{32} & a_{33}
\end{pmatrix}
\]

\section{Tables}
Below is a sample table showing data:

\begin{table}[h]
\centering
\begin{tabular}{@{}lrrr@{}}
\toprule
\textbf{Category} & \textbf{Value 1} & \textbf{Value 2} & \textbf{Total} \\
\midrule
Item A & 10.5 & 20.3 & 30.8 \\
Item B & 15.2 & 25.1 & 40.3 \\
Item C & 22.1 & 12.7 & 34.8 \\
\bottomrule
\end{tabular}
\caption{Sample Data Table}
\label{tab:sample}
\end{table}

\section{Lists}
\subsection{Bullet Points}
\begin{itemize}
  \item First item
  \item Second item
  \item Third item with sub-items:
    \begin{itemize}
      \item Sub-item 1
      \item Sub-item 2
    \end{itemize}
\end{itemize}

\subsection{Numbered List}
\begin{enumerate}
  \item First step
  \item Second step
  \item Third step
\end{enumerate}

\section{Citations and References}
When you use the citation search feature, references will be automatically managed. For example, you might cite a paper like this: \cite{example2024}. All citations will be automatically added to your \texttt{cite.bib} file.

\section{Conclusion}
This template demonstrates basic \LaTeX{} features including mathematical formulas, tables, and lists. The bibliography system is now set up to work seamlessly with Bibby AI's citation management. Feel free to modify it for your needs!

\bibliographystyle{plain}
\bibliography{cite}

\end{document}\documentclass{article}
\usepackage[utf8]{inputenc}
\usepackage{booktabs}
\usepackage{amsmath}
\usepackage{natbib}

\title{My Document}
\author{Your Name}
\date{\today}

\begin{document}

\maketitle

\section{Introduction}
This is a sample document demonstrating various \LaTeX{} features including mathematical equations, tables, and formatting. You can cite papers using the citation search feature, and they will be automatically added to your \texttt{cite.bib} file.

\section{Mathematical Expressions}
Here are some example equations:

\subsection{Quadratic Formula}
The quadratic formula is given by:
\[x = \frac{-b \pm \sqrt{b^2 - 4ac}}{2a}\]

\subsection{Matrix Example}
A sample matrix:
\[
\begin{pmatrix}
a_{11} & a_{12} & a_{13} \\
a_{21} & a_{22} & a_{23} \\
a_{31} & a_{32} & a_{33}
\end{pmatrix}
\]

\section{Tables}
Below is a sample table showing data:

\begin{table}[h]
\centering
\begin{tabular}{@{}lrrr@{}}
\toprule
\textbf{Category} & \textbf{Value 1} & \textbf{Value 2} & \textbf{Total} \\
\midrule
Item A & 10.5 & 20.3 & 30.8 \\
Item B & 15.2 & 25.1 & 40.3 \\
Item C & 22.1 & 12.7 & 34.8 \\
\bottomrule
\end{tabular}
\caption{Sample Data Table}
\label{tab:sample}
\end{table}

\section{Lists}
\subsection{Bullet Points}
\begin{itemize}
  \item First item
  \item Second item
  \item Third item with sub-items:
    \begin{itemize}
      \item Sub-item 1
      \item Sub-item 2
    \end{itemize}
\end{itemize}

\subsection{Numbered List}
\begin{enumerate}
  \item First step
  \item Second step
  \item Third step
\end{enumerate}

\section{Citations and References}
When you use the citation search feature, references will be automatically managed. For example, you might cite a paper like this: \cite{example2024}. All citations will be automatically added to your \texttt{cite.bib} file.

\section{Conclusion}
This template demonstrates basic \LaTeX{} features including mathematical formulas, tables, and lists. The bibliography system is now set up to work seamlessly with Bibby AI's citation management. Feel free to modify it for your needs!

\bibliographystyle{plain}
\bibliography{cite}

\end{document}\documentclass{article}
\usepackage[utf8]{inputenc}
\usepackage{booktabs}
\usepackage{amsmath}
\usepackage{natbib}

\title{My Document}
\author{Your Name}
\date{\today}

\begin{document}

\maketitle

\section{Introduction}
This is a sample document demonstrating various \LaTeX{} features including mathematical equations, tables, and formatting. You can cite papers using the citation search feature, and they will be automatically added to your \texttt{cite.bib} file.

\section{Mathematical Expressions}
Here are some example equations:

\subsection{Quadratic Formula}
The quadratic formula is given by:
\[x = \frac{-b \pm \sqrt{b^2 - 4ac}}{2a}\]

\subsection{Matrix Example}
A sample matrix:
\[
\begin{pmatrix}
a_{11} & a_{12} & a_{13} \\
a_{21} & a_{22} & a_{23} \\
a_{31} & a_{32} & a_{33}
\end{pmatrix}
\]

\section{Tables}
Below is a sample table showing data:

\begin{table}[h]
\centering
\begin{tabular}{@{}lrrr@{}}
\toprule
\textbf{Category} & \textbf{Value 1} & \textbf{Value 2} & \textbf{Total} \\
\midrule
Item A & 10.5 & 20.3 & 30.8 \\
Item B & 15.2 & 25.1 & 40.3 \\
Item C & 22.1 & 12.7 & 34.8 \\
\bottomrule
\end{tabular}
\caption{Sample Data Table}
\label{tab:sample}
\end{table}

\section{Lists}
\subsection{Bullet Points}
\begin{itemize}
  \item First item
  \item Second item
  \item Third item with sub-items:
    \begin{itemize}
      \item Sub-item 1
      \item Sub-item 2
    \end{itemize}
\end{itemize}

\subsection{Numbered List}
\begin{enumerate}
  \item First step
  \item Second step
  \item Third step
\end{enumerate}

\section{Citations and References}
When you use the citation search feature, references will be automatically managed. For example, you might cite a paper like this: \cite{example2024}. All citations will be automatically added to your \texttt{cite.bib} file.

\section{Conclusion}
This template demonstrates basic \LaTeX{} features including mathematical formulas, tables, and lists. The bibliography system is now set up to work seamlessly with Bibby AI's citation management. Feel free to modify it for your needs!

\bibliographystyle{plain}
\bibliography{cite}

\end{document}\documentclass{article}
\usepackage[utf8]{inputenc}
\usepackage{booktabs}
\usepackage{amsmath}
\usepackage{natbib}

\title{My Document}
\author{Your Name}
\date{\today}

\begin{document}

\maketitle

\section{Introduction}
This is a sample document demonstrating various \LaTeX{} features including mathematical equations, tables, and formatting. You can cite papers using the citation search feature, and they will be automatically added to your \texttt{cite.bib} file.

\section{Mathematical Expressions}
Here are some example equations:

\subsection{Quadratic Formula}
The quadratic formula is given by:
\[x = \frac{-b \pm \sqrt{b^2 - 4ac}}{2a}\]

\subsection{Matrix Example}
A sample matrix:
\[
\begin{pmatrix}
a_{11} & a_{12} & a_{13} \\
a_{21} & a_{22} & a_{23} \\
a_{31} & a_{32} & a_{33}
\end{pmatrix}
\]

\section{Tables}
Below is a sample table showing data:

\begin{table}[h]
\centering
\begin{tabular}{@{}lrrr@{}}
\toprule
\textbf{Category} & \textbf{Value 1} & \textbf{Value 2} & \textbf{Total} \\
\midrule
Item A & 10.5 & 20.3 & 30.8 \\
Item B & 15.2 & 25.1 & 40.3 \\
Item C & 22.1 & 12.7 & 34.8 \\
\bottomrule
\end{tabular}
\caption{Sample Data Table}
\label{tab:sample}
\end{table}

\section{Lists}
\subsection{Bullet Points}
\begin{itemize}
  \item First item
  \item Second item
  \item Third item with sub-items:
    \begin{itemize}
      \item Sub-item 1
      \item Sub-item 2
    \end{itemize}
\end{itemize}

\subsection{Numbered List}
\begin{enumerate}
  \item First step
  \item Second step
  \item Third step
\end{enumerate}

\section{Citations and References}
When you use the citation search feature, references will be automatically managed. For example, you might cite a paper like this: \cite{example2024}. All citations will be automatically added to your \texttt{cite.bib} file.

\section{Conclusion}
This template demonstrates basic \LaTeX{} features including mathematical formulas, tables, and lists. The bibliography system is now set up to work seamlessly with Bibby AI's citation management. Feel free to modify it for your needs!

\bibliographystyle{plain}
\bibliography{cite}

\end{document}\documentclass{article}
\usepackage[utf8]{inputenc}
\usepackage{booktabs}
\usepackage{amsmath}
\usepackage{natbib}

\title{My Document}
\author{Your Name}
\date{\today}

\begin{document}

\maketitle

\section{Introduction}
This is a sample document demonstrating various \LaTeX{} features including mathematical equations, tables, and formatting. You can cite papers using the citation search feature, and they will be automatically added to your \texttt{cite.bib} file.

\section{Mathematical Expressions}
Here are some example equations:

\subsection{Quadratic Formula}
The quadratic formula is given by:
\[x = \frac{-b \pm \sqrt{b^2 - 4ac}}{2a}\]

\subsection{Matrix Example}
A sample matrix:
\[
\begin{pmatrix}
a_{11} & a_{12} & a_{13} \\
a_{21} & a_{22} & a_{23} \\
a_{31} & a_{32} & a_{33}
\end{pmatrix}
\]

\section{Tables}
Below is a sample table showing data:

\begin{table}[h]
\centering
\begin{tabular}{@{}lrrr@{}}
\toprule
\textbf{Category} & \textbf{Value 1} & \textbf{Value 2} & \textbf{Total} \\
\midrule
Item A & 10.5 & 20.3 & 30.8 \\
Item B & 15.2 & 25.1 & 40.3 \\
Item C & 22.1 & 12.7 & 34.8 \\
\bottomrule
\end{tabular}
\caption{Sample Data Table}
\label{tab:sample}
\end{table}

\section{Lists}
\subsection{Bullet Points}
\begin{itemize}
  \item First item
  \item Second item
  \item Third item with sub-items:
    \begin{itemize}
      \item Sub-item 1
      \item Sub-item 2
    \end{itemize}
\end{itemize}

\subsection{Numbered List}
\begin{enumerate}
  \item First step
  \item Second step
  \item Third step
\end{enumerate}

\section{Citations and References}
When you use the citation search feature, references will be automatically managed. For example, you might cite a paper like this: \cite{example2024}. All citations will be automatically added to your \texttt{cite.bib} file.

\section{Conclusion}
This template demonstrates basic \LaTeX{} features including mathematical formulas, tables, and lists. The bibliography system is now set up to work seamlessly with Bibby AI's citation management. Feel free to modify it for your needs!

\bibliographystyle{plain}
\bibliography{cite}

\end{document}\documentclass{article}
\usepackage[utf8]{inputenc}
\usepackage{booktabs}
\usepackage{amsmath}
\usepackage{natbib}

\title{My Document}
\author{Your Name}
\date{\today}

\begin{document}

\maketitle

\section{Introduction}
This is a sample document demonstrating various \LaTeX{} features including mathematical equations, tables, and formatting. You can cite papers using the citation search feature, and they will be automatically added to your \texttt{cite.bib} file.

\section{Mathematical Expressions}
Here are some example equations:

\subsection{Quadratic Formula}
The quadratic formula is given by:
\[x = \frac{-b \pm \sqrt{b^2 - 4ac}}{2a}\]

\subsection{Matrix Example}
A sample matrix:
\[
\begin{pmatrix}
a_{11} & a_{12} & a_{13} \\
a_{21} & a_{22} & a_{23} \\
a_{31} & a_{32} & a_{33}
\end{pmatrix}
\]

\section{Tables}
Below is a sample table showing data:

\begin{table}[h]
\centering
\begin{tabular}{@{}lrrr@{}}
\toprule
\textbf{Category} & \textbf{Value 1} & \textbf{Value 2} & \textbf{Total} \\
\midrule
Item A & 10.5 & 20.3 & 30.8 \\
Item B & 15.2 & 25.1 & 40.3 \\
Item C & 22.1 & 12.7 & 34.8 \\
\bottomrule
\end{tabular}
\caption{Sample Data Table}
\label{tab:sample}
\end{table}

\section{Lists}
\subsection{Bullet Points}
\begin{itemize}
  \item First item
  \item Second item
  \item Third item with sub-items:
    \begin{itemize}
      \item Sub-item 1
      \item Sub-item 2
    \end{itemize}
\end{itemize}

\subsection{Numbered List}
\begin{enumerate}
  \item First step
  \item Second step
  \item Third step
\end{enumerate}

\section{Citations and References}
When you use the citation search feature, references will be automatically managed. For example, you might cite a paper like this: \cite{example2024}. All citations will be automatically added to your \texttt{cite.bib} file.

\section{Conclusion}
This template demonstrates basic \LaTeX{} features including mathematical formulas, tables, and lists. The bibliography system is now set up to work seamlessly with Bibby AI's citation management. Feel free to modify it for your needs!

\bibliographystyle{plain}
\bibliography{cite}

\end{document}\documentclass{article}
\usepackage[utf8]{inputenc}
\usepackage{booktabs}
\usepackage{amsmath}
\usepackage{natbib}

\title{My Document}
\author{Your Name}
\date{\today}

\begin{document}

\maketitle

\section{Introduction}
This is a sample document demonstrating various \LaTeX{} features including mathematical equations, tables, and formatting. You can cite papers using the citation search feature, and they will be automatically added to your \texttt{cite.bib} file.

\section{Mathematical Expressions}
Here are some example equations:

\subsection{Quadratic Formula}
The quadratic formula is given by:
\[x = \frac{-b \pm \sqrt{b^2 - 4ac}}{2a}\]

\subsection{Matrix Example}
A sample matrix:
\[
\begin{pmatrix}
a_{11} & a_{12} & a_{13} \\
a_{21} & a_{22} & a_{23} \\
a_{31} & a_{32} & a_{33}
\end{pmatrix}
\]

\section{Tables}
Below is a sample table showing data:

\begin{table}[h]
\centering
\begin{tabular}{@{}lrrr@{}}
\toprule
\textbf{Category} & \textbf{Value 1} & \textbf{Value 2} & \textbf{Total} \\
\midrule
Item A & 10.5 & 20.3 & 30.8 \\
Item B & 15.2 & 25.1 & 40.3 \\
Item C & 22.1 & 12.7 & 34.8 \\
\bottomrule
\end{tabular}
\caption{Sample Data Table}
\label{tab:sample}
\end{table}

\section{Lists}
\subsection{Bullet Points}
\begin{itemize}
  \item First item
  \item Second item
  \item Third item with sub-items:
    \begin{itemize}
      \item Sub-item 1
      \item Sub-item 2
    \end{itemize}
\end{itemize}

\subsection{Numbered List}
\begin{enumerate}
  \item First step
  \item Second step
  \item Third step
\end{enumerate}

\section{Citations and References}
When you use the citation search feature, references will be automatically managed. For example, you might cite a paper like this: \cite{example2024}. All citations will be automatically added to your \texttt{cite.bib} file.

\section{Conclusion}
This template demonstrates basic \LaTeX{} features including mathematical formulas, tables, and lists. The bibliography system is now set up to work seamlessly with Bibby AI's citation management. Feel free to modify it for your needs!

\bibliographystyle{plain}
\bibliography{cite}

\end{document}\documentclass{article}
\usepackage[utf8]{inputenc}
\usepackage{booktabs}
\usepackage{amsmath}
\usepackage{natbib}

\title{My Document}
\author{Your Name}
\date{\today}

\begin{document}

\maketitle

\section{Introduction}
This is a sample document demonstrating various \LaTeX{} features including mathematical equations, tables, and formatting. You can cite papers using the citation search feature, and they will be automatically added to your \texttt{cite.bib} file.

\section{Mathematical Expressions}
Here are some example equations:

\subsection{Quadratic Formula}
The quadratic formula is given by:
\[x = \frac{-b \pm \sqrt{b^2 - 4ac}}{2a}\]

\subsection{Matrix Example}
A sample matrix:
\[
\begin{pmatrix}
a_{11} & a_{12} & a_{13} \\
a_{21} & a_{22} & a_{23} \\
a_{31} & a_{32} & a_{33}
\end{pmatrix}
\]

\section{Tables}
Below is a sample table showing data:

\begin{table}[h]
\centering
\begin{tabular}{@{}lrrr@{}}
\toprule
\textbf{Category} & \textbf{Value 1} & \textbf{Value 2} & \textbf{Total} \\
\midrule
Item A & 10.5 & 20.3 & 30.8 \\
Item B & 15.2 & 25.1 & 40.3 \\
Item C & 22.1 & 12.7 & 34.8 \\
\bottomrule
\end{tabular}
\caption{Sample Data Table}
\label{tab:sample}
\end{table}

\section{Lists}
\subsection{Bullet Points}
\begin{itemize}
  \item First item
  \item Second item
  \item Third item with sub-items:
    \begin{itemize}
      \item Sub-item 1
      \item Sub-item 2
    \end{itemize}
\end{itemize}

\subsection{Numbered List}
\begin{enumerate}
  \item First step
  \item Second step
  \item Third step
\end{enumerate}

\section{Citations and References}
When you use the citation search feature, references will be automatically managed. For example, you might cite a paper like this: \cite{example2024}. All citations will be automatically added to your \texttt{cite.bib} file.

\section{Conclusion}
This template demonstrates basic \LaTeX{} features including mathematical formulas, tables, and lists. The bibliography system is now set up to work seamlessly with Bibby AI's citation management. Feel free to modify it for your needs!

\bibliographystyle{plain}
\bibliography{cite}

\end{document}\documentclass{article}
\usepackage[utf8]{inputenc}
\usepackage{booktabs}
\usepackage{amsmath}
\usepackage{natbib}

\title{My Document}
\author{Your Name}
\date{\today}

\begin{document}

\maketitle

\section{Introduction}
This is a sample document demonstrating various \LaTeX{} features including mathematical equations, tables, and formatting. You can cite papers using the citation search feature, and they will be automatically added to your \texttt{cite.bib} file.

\section{Mathematical Expressions}
Here are some example equations:

\subsection{Quadratic Formula}
The quadratic formula is given by:
\[x = \frac{-b \pm \sqrt{b^2 - 4ac}}{2a}\]

\subsection{Matrix Example}
A sample matrix:
\[
\begin{pmatrix}
a_{11} & a_{12} & a_{13} \\
a_{21} & a_{22} & a_{23} \\
a_{31} & a_{32} & a_{33}
\end{pmatrix}
\]

\section{Tables}
Below is a sample table showing data:

\begin{table}[h]
\centering
\begin{tabular}{@{}lrrr@{}}
\toprule
\textbf{Category} & \textbf{Value 1} & \textbf{Value 2} & \textbf{Total} \\
\midrule
Item A & 10.5 & 20.3 & 30.8 \\
Item B & 15.2 & 25.1 & 40.3 \\
Item C & 22.1 & 12.7 & 34.8 \\
\bottomrule
\end{tabular}
\caption{Sample Data Table}
\label{tab:sample}
\end{table}

\section{Lists}
\subsection{Bullet Points}
\begin{itemize}
  \item First item
  \item Second item
  \item Third item with sub-items:
    \begin{itemize}
      \item Sub-item 1
      \item Sub-item 2
    \end{itemize}
\end{itemize}

\subsection{Numbered List}
\begin{enumerate}
  \item First step
  \item Second step
  \item Third step
\end{enumerate}

\section{Citations and References}
When you use the citation search feature, references will be automatically managed. For example, you might cite a paper like this: \cite{example2024}. All citations will be automatically added to your \texttt{cite.bib} file.

\section{Conclusion}
This template demonstrates basic \LaTeX{} features including mathematical formulas, tables, and lists. The bibliography system is now set up to work seamlessly with Bibby AI's citation management. Feel free to modify it for your needs!

\bibliographystyle{plain}
\bibliography{cite}

\end{document}\documentclass{article}
\usepackage[utf8]{inputenc}
\usepackage{booktabs}
\usepackage{amsmath}
\usepackage{natbib}

\title{My Document}
\author{Your Name}
\date{\today}

\begin{document}

\maketitle

\section{Introduction}
This is a sample document demonstrating various \LaTeX{} features including mathematical equations, tables, and formatting. You can cite papers using the citation search feature, and they will be automatically added to your \texttt{cite.bib} file.

\section{Mathematical Expressions}
Here are some example equations:

\subsection{Quadratic Formula}
The quadratic formula is given by:
\[x = \frac{-b \pm \sqrt{b^2 - 4ac}}{2a}\]

\subsection{Matrix Example}
A sample matrix:
\[
\begin{pmatrix}
a_{11} & a_{12} & a_{13} \\
a_{21} & a_{22} & a_{23} \\
a_{31} & a_{32} & a_{33}
\end{pmatrix}
\]

\section{Tables}
Below is a sample table showing data:

\begin{table}[h]
\centering
\begin{tabular}{@{}lrrr@{}}
\toprule
\textbf{Category} & \textbf{Value 1} & \textbf{Value 2} & \textbf{Total} \\
\midrule
Item A & 10.5 & 20.3 & 30.8 \\
Item B & 15.2 & 25.1 & 40.3 \\
Item C & 22.1 & 12.7 & 34.8 \\
\bottomrule
\end{tabular}
\caption{Sample Data Table}
\label{tab:sample}
\end{table}

\section{Lists}
\subsection{Bullet Points}
\begin{itemize}
  \item First item
  \item Second item
  \item Third item with sub-items:
    \begin{itemize}
      \item Sub-item 1
      \item Sub-item 2
    \end{itemize}
\end{itemize}

\subsection{Numbered List}
\begin{enumerate}
  \item First step
  \item Second step
  \item Third step
\end{enumerate}

\section{Citations and References}
When you use the citation search feature, references will be automatically managed. For example, you might cite a paper like this: \cite{example2024}. All citations will be automatically added to your \texttt{cite.bib} file.

\section{Conclusion}
This template demonstrates basic \LaTeX{} features including mathematical formulas, tables, and lists. The bibliography system is now set up to work seamlessly with Bibby AI's citation management. Feel free to modify it for your needs!

\bibliographystyle{plain}
\bibliography{cite}

\end{document}\documentclass{article}
\usepackage[utf8]{inputenc}
\usepackage{booktabs}
\usepackage{amsmath}
\usepackage{natbib}

\title{My Document}
\author{Your Name}
\date{\today}

\begin{document}

\maketitle

\section{Introduction}
This is a sample document demonstrating various \LaTeX{} features including mathematical equations, tables, and formatting. You can cite papers using the citation search feature, and they will be automatically added to your \texttt{cite.bib} file.

\section{Mathematical Expressions}
Here are some example equations:

\subsection{Quadratic Formula}
The quadratic formula is given by:
\[x = \frac{-b \pm \sqrt{b^2 - 4ac}}{2a}\]

\subsection{Matrix Example}
A sample matrix:
\[
\begin{pmatrix}
a_{11} & a_{12} & a_{13} \\
a_{21} & a_{22} & a_{23} \\
a_{31} & a_{32} & a_{33}
\end{pmatrix}
\]

\section{Tables}
Below is a sample table showing data:

\begin{table}[h]
\centering
\begin{tabular}{@{}lrrr@{}}
\toprule
\textbf{Category} & \textbf{Value 1} & \textbf{Value 2} & \textbf{Total} \\
\midrule
Item A & 10.5 & 20.3 & 30.8 \\
Item B & 15.2 & 25.1 & 40.3 \\
Item C & 22.1 & 12.7 & 34.8 \\
\bottomrule
\end{tabular}
\caption{Sample Data Table}
\label{tab:sample}
\end{table}

\section{Lists}
\subsection{Bullet Points}
\begin{itemize}
  \item First item
  \item Second item
  \item Third item with sub-items:
    \begin{itemize}
      \item Sub-item 1
      \item Sub-item 2
    \end{itemize}
\end{itemize}

\subsection{Numbered List}
\begin{enumerate}
  \item First step
  \item Second step
  \item Third step
\end{enumerate}

\section{Citations and References}
When you use the citation search feature, references will be automatically managed. For example, you might cite a paper like this: \cite{example2024}. All citations will be automatically added to your \texttt{cite.bib} file.

\section{Conclusion}
This template demonstrates basic \LaTeX{} features including mathematical formulas, tables, and lists. The bibliography system is now set up to work seamlessly with Bibby AI's citation management. Feel free to modify it for your needs!

\bibliographystyle{plain}
\bibliography{cite}

\end{document}\documentclass{article}
\usepackage[utf8]{inputenc}
\usepackage{booktabs}
\usepackage{amsmath}
\usepackage{natbib}

\title{My Document}
\author{Your Name}
\date{\today}

\begin{document}

\maketitle

\section{Introduction}
This is a sample document demonstrating various \LaTeX{} features including mathematical equations, tables, and formatting. You can cite papers using the citation search feature, and they will be automatically added to your \texttt{cite.bib} file.

\section{Mathematical Expressions}
Here are some example equations:

\subsection{Quadratic Formula}
The quadratic formula is given by:
\[x = \frac{-b \pm \sqrt{b^2 - 4ac}}{2a}\]

\subsection{Matrix Example}
A sample matrix:
\[
\begin{pmatrix}
a_{11} & a_{12} & a_{13} \\
a_{21} & a_{22} & a_{23} \\
a_{31} & a_{32} & a_{33}
\end{pmatrix}
\]

\section{Tables}
Below is a sample table showing data:

\begin{table}[h]
\centering
\begin{tabular}{@{}lrrr@{}}
\toprule
\textbf{Category} & \textbf{Value 1} & \textbf{Value 2} & \textbf{Total} \\
\midrule
Item A & 10.5 & 20.3 & 30.8 \\
Item B & 15.2 & 25.1 & 40.3 \\
Item C & 22.1 & 12.7 & 34.8 \\
\bottomrule
\end{tabular}
\caption{Sample Data Table}
\label{tab:sample}
\end{table}

\section{Lists}
\subsection{Bullet Points}
\begin{itemize}
  \item First item
  \item Second item
  \item Third item with sub-items:
    \begin{itemize}
      \item Sub-item 1
      \item Sub-item 2
    \end{itemize}
\end{itemize}

\subsection{Numbered List}
\begin{enumerate}
  \item First step
  \item Second step
  \item Third step
\end{enumerate}

\section{Citations and References}
When you use the citation search feature, references will be automatically managed. For example, you might cite a paper like this: \cite{example2024}. All citations will be automatically added to your \texttt{cite.bib} file.

\section{Conclusion}
This template demonstrates basic \LaTeX{} features including mathematical formulas, tables, and lists. The bibliography system is now set up to work seamlessly with Bibby AI's citation management. Feel free to modify it for your needs!

\bibliographystyle{plain}
\bibliography{cite}

\end{document}\documentclass{article}
\usepackage[utf8]{inputenc}
\usepackage{booktabs}
\usepackage{amsmath}
\usepackage{natbib}

\title{My Document}
\author{Your Name}
\date{\today}

\begin{document}

\maketitle

\section{Introduction}
This is a sample document demonstrating various \LaTeX{} features including mathematical equations, tables, and formatting. You can cite papers using the citation search feature, and they will be automatically added to your \texttt{cite.bib} file.

\section{Mathematical Expressions}
Here are some example equations:

\subsection{Quadratic Formula}
The quadratic formula is given by:
\[x = \frac{-b \pm \sqrt{b^2 - 4ac}}{2a}\]

\subsection{Matrix Example}
A sample matrix:
\[
\begin{pmatrix}
a_{11} & a_{12} & a_{13} \\
a_{21} & a_{22} & a_{23} \\
a_{31} & a_{32} & a_{33}
\end{pmatrix}
\]

\section{Tables}
Below is a sample table showing data:

\begin{table}[h]
\centering
\begin{tabular}{@{}lrrr@{}}
\toprule
\textbf{Category} & \textbf{Value 1} & \textbf{Value 2} & \textbf{Total} \\
\midrule
Item A & 10.5 & 20.3 & 30.8 \\
Item B & 15.2 & 25.1 & 40.3 \\
Item C & 22.1 & 12.7 & 34.8 \\
\bottomrule
\end{tabular}
\caption{Sample Data Table}
\label{tab:sample}
\end{table}

\section{Lists}
\subsection{Bullet Points}
\begin{itemize}
  \item First item
  \item Second item
  \item Third item with sub-items:
    \begin{itemize}
      \item Sub-item 1
      \item Sub-item 2
    \end{itemize}
\end{itemize}

\subsection{Numbered List}
\begin{enumerate}
  \item First step
  \item Second step
  \item Third step
\end{enumerate}

\section{Citations and References}
When you use the citation search feature, references will be automatically managed. For example, you might cite a paper like this: \cite{example2024}. All citations will be automatically added to your \texttt{cite.bib} file.

\section{Conclusion}
This template demonstrates basic \LaTeX{} features including mathematical formulas, tables, and lists. The bibliography system is now set up to work seamlessly with Bibby AI's citation management. Feel free to modify it for your needs!

\bibliographystyle{plain}
\bibliography{cite}

\end{document}\documentclass{article}
\usepackage[utf8]{inputenc}
\usepackage{booktabs}
\usepackage{amsmath}
\usepackage{natbib}

\title{My Document}
\author{Your Name}
\date{\today}

\begin{document}

\maketitle

\section{Introduction}
This is a sample document demonstrating various \LaTeX{} features including mathematical equations, tables, and formatting. You can cite papers using the citation search feature, and they will be automatically added to your \texttt{cite.bib} file.

\section{Mathematical Expressions}
Here are some example equations:

\subsection{Quadratic Formula}
The quadratic formula is given by:
\[x = \frac{-b \pm \sqrt{b^2 - 4ac}}{2a}\]

\subsection{Matrix Example}
A sample matrix:
\[
\begin{pmatrix}
a_{11} & a_{12} & a_{13} \\
a_{21} & a_{22} & a_{23} \\
a_{31} & a_{32} & a_{33}
\end{pmatrix}
\]

\section{Tables}
Below is a sample table showing data:

\begin{table}[h]
\centering
\begin{tabular}{@{}lrrr@{}}
\toprule
\textbf{Category} & \textbf{Value 1} & \textbf{Value 2} & \textbf{Total} \\
\midrule
Item A & 10.5 & 20.3 & 30.8 \\
Item B & 15.2 & 25.1 & 40.3 \\
Item C & 22.1 & 12.7 & 34.8 \\
\bottomrule
\end{tabular}
\caption{Sample Data Table}
\label{tab:sample}
\end{table}

\section{Lists}
\subsection{Bullet Points}
\begin{itemize}
  \item First item
  \item Second item
  \item Third item with sub-items:
    \begin{itemize}
      \item Sub-item 1
      \item Sub-item 2
    \end{itemize}
\end{itemize}

\subsection{Numbered List}
\begin{enumerate}
  \item First step
  \item Second step
  \item Third step
\end{enumerate}

\section{Citations and References}
When you use the citation search feature, references will be automatically managed. For example, you might cite a paper like this: \cite{example2024}. All citations will be automatically added to your \texttt{cite.bib} file.

\section{Conclusion}
This template demonstrates basic \LaTeX{} features including mathematical formulas, tables, and lists. The bibliography system is now set up to work seamlessly with Bibby AI's citation management. Feel free to modify it for your needs!

\bibliographystyle{plain}
\bibliography{cite}

\end{document}\documentclass{article}
\usepackage[utf8]{inputenc}
\usepackage{booktabs}
\usepackage{amsmath}
\usepackage{natbib}

\title{My Document}
\author{Your Name}
\date{\today}

\begin{document}

\maketitle

\section{Introduction}
This is a sample document demonstrating various \LaTeX{} features including mathematical equations, tables, and formatting. You can cite papers using the citation search feature, and they will be automatically added to your \texttt{cite.bib} file.

\section{Mathematical Expressions}
Here are some example equations:

\subsection{Quadratic Formula}
The quadratic formula is given by:
\[x = \frac{-b \pm \sqrt{b^2 - 4ac}}{2a}\]

\subsection{Matrix Example}
A sample matrix:
\[
\begin{pmatrix}
a_{11} & a_{12} & a_{13} \\
a_{21} & a_{22} & a_{23} \\
a_{31} & a_{32} & a_{33}
\end{pmatrix}
\]

\section{Tables}
Below is a sample table showing data:

\begin{table}[h]
\centering
\begin{tabular}{@{}lrrr@{}}
\toprule
\textbf{Category} & \textbf{Value 1} & \textbf{Value 2} & \textbf{Total} \\
\midrule
Item A & 10.5 & 20.3 & 30.8 \\
Item B & 15.2 & 25.1 & 40.3 \\
Item C & 22.1 & 12.7 & 34.8 \\
\bottomrule
\end{tabular}
\caption{Sample Data Table}
\label{tab:sample}
\end{table}

\section{Lists}
\subsection{Bullet Points}
\begin{itemize}
  \item First item
  \item Second item
  \item Third item with sub-items:
    \begin{itemize}
      \item Sub-item 1
      \item Sub-item 2
    \end{itemize}
\end{itemize}

\subsection{Numbered List}
\begin{enumerate}
  \item First step
  \item Second step
  \item Third step
\end{enumerate}

\section{Citations and References}
When you use the citation search feature, references will be automatically managed. For example, you might cite a paper like this: \cite{example2024}. All citations will be automatically added to your \texttt{cite.bib} file.

\section{Conclusion}
This template demonstrates basic \LaTeX{} features including mathematical formulas, tables, and lists. The bibliography system is now set up to work seamlessly with Bibby AI's citation management. Feel free to modify it for your needs!

\bibliographystyle{plain}
\bibliography{cite}

\end{document}\documentclass{article}
\usepackage[utf8]{inputenc}
\usepackage{booktabs}
\usepackage{amsmath}
\usepackage{natbib}

\title{My Document}
\author{Your Name}
\date{\today}

\begin{document}

\maketitle

\section{Introduction}
This is a sample document demonstrating various \LaTeX{} features including mathematical equations, tables, and formatting. You can cite papers using the citation search feature, and they will be automatically added to your \texttt{cite.bib} file.

\section{Mathematical Expressions}
Here are some example equations:

\subsection{Quadratic Formula}
The quadratic formula is given by:
\[x = \frac{-b \pm \sqrt{b^2 - 4ac}}{2a}\]

\subsection{Matrix Example}
A sample matrix:
\[
\begin{pmatrix}
a_{11} & a_{12} & a_{13} \\
a_{21} & a_{22} & a_{23} \\
a_{31} & a_{32} & a_{33}
\end{pmatrix}
\]

\section{Tables}
Below is a sample table showing data:

\begin{table}[h]
\centering
\begin{tabular}{@{}lrrr@{}}
\toprule
\textbf{Category} & \textbf{Value 1} & \textbf{Value 2} & \textbf{Total} \\
\midrule
Item A & 10.5 & 20.3 & 30.8 \\
Item B & 15.2 & 25.1 & 40.3 \\
Item C & 22.1 & 12.7 & 34.8 \\
\bottomrule
\end{tabular}
\caption{Sample Data Table}
\label{tab:sample}
\end{table}

\section{Lists}
\subsection{Bullet Points}
\begin{itemize}
  \item First item
  \item Second item
  \item Third item with sub-items:
    \begin{itemize}
      \item Sub-item 1
      \item Sub-item 2
    \end{itemize}
\end{itemize}

\subsection{Numbered List}
\begin{enumerate}
  \item First step
  \item Second step
  \item Third step
\end{enumerate}

\section{Citations and References}
When you use the citation search feature, references will be automatically managed. For example, you might cite a paper like this: \cite{example2024}. All citations will be automatically added to your \texttt{cite.bib} file.

\section{Conclusion}
This template demonstrates basic \LaTeX{} features including mathematical formulas, tables, and lists. The bibliography system is now set up to work seamlessly with Bibby AI's citation management. Feel free to modify it for your needs!

\bibliographystyle{plain}
\bibliography{cite}

\end{document}\documentclass{article}
\usepackage[utf8]{inputenc}
\usepackage{booktabs}
\usepackage{amsmath}
\usepackage{natbib}

\title{My Document}
\author{Your Name}
\date{\today}

\begin{document}

\maketitle

\section{Introduction}
This is a sample document demonstrating various \LaTeX{} features including mathematical equations, tables, and formatting. You can cite papers using the citation search feature, and they will be automatically added to your \texttt{cite.bib} file.

\section{Mathematical Expressions}
Here are some example equations:

\subsection{Quadratic Formula}
The quadratic formula is given by:
\[x = \frac{-b \pm \sqrt{b^2 - 4ac}}{2a}\]

\subsection{Matrix Example}
A sample matrix:
\[
\begin{pmatrix}
a_{11} & a_{12} & a_{13} \\
a_{21} & a_{22} & a_{23} \\
a_{31} & a_{32} & a_{33}
\end{pmatrix}
\]

\section{Tables}
Below is a sample table showing data:

\begin{table}[h]
\centering
\begin{tabular}{@{}lrrr@{}}
\toprule
\textbf{Category} & \textbf{Value 1} & \textbf{Value 2} & \textbf{Total} \\
\midrule
Item A & 10.5 & 20.3 & 30.8 \\
Item B & 15.2 & 25.1 & 40.3 \\
Item C & 22.1 & 12.7 & 34.8 \\
\bottomrule
\end{tabular}
\caption{Sample Data Table}
\label{tab:sample}
\end{table}

\section{Lists}
\subsection{Bullet Points}
\begin{itemize}
  \item First item
  \item Second item
  \item Third item with sub-items:
    \begin{itemize}
      \item Sub-item 1
      \item Sub-item 2
    \end{itemize}
\end{itemize}

\subsection{Numbered List}
\begin{enumerate}
  \item First step
  \item Second step
  \item Third step
\end{enumerate}

\section{Citations and References}
When you use the citation search feature, references will be automatically managed. For example, you might cite a paper like this: \cite{example2024}. All citations will be automatically added to your \texttt{cite.bib} file.

\section{Conclusion}
This template demonstrates basic \LaTeX{} features including mathematical formulas, tables, and lists. The bibliography system is now set up to work seamlessly with Bibby AI's citation management. Feel free to modify it for your needs!

\bibliographystyle{plain}
\bibliography{cite}

\end{document}\documentclass{article}
\usepackage[utf8]{inputenc}
\usepackage{booktabs}
\usepackage{amsmath}
\usepackage{natbib}

\title{My Document}
\author{Your Name}
\date{\today}

\begin{document}

\maketitle

\section{Introduction}
This is a sample document demonstrating various \LaTeX{} features including mathematical equations, tables, and formatting. You can cite papers using the citation search feature, and they will be automatically added to your \texttt{cite.bib} file.

\section{Mathematical Expressions}
Here are some example equations:

\subsection{Quadratic Formula}
The quadratic formula is given by:
\[x = \frac{-b \pm \sqrt{b^2 - 4ac}}{2a}\]

\subsection{Matrix Example}
A sample matrix:
\[
\begin{pmatrix}
a_{11} & a_{12} & a_{13} \\
a_{21} & a_{22} & a_{23} \\
a_{31} & a_{32} & a_{33}
\end{pmatrix}
\]

\section{Tables}
Below is a sample table showing data:

\begin{table}[h]
\centering
\begin{tabular}{@{}lrrr@{}}
\toprule
\textbf{Category} & \textbf{Value 1} & \textbf{Value 2} & \textbf{Total} \\
\midrule
Item A & 10.5 & 20.3 & 30.8 \\
Item B & 15.2 & 25.1 & 40.3 \\
Item C & 22.1 & 12.7 & 34.8 \\
\bottomrule
\end{tabular}
\caption{Sample Data Table}
\label{tab:sample}
\end{table}

\section{Lists}
\subsection{Bullet Points}
\begin{itemize}
  \item First item
  \item Second item
  \item Third item with sub-items:
    \begin{itemize}
      \item Sub-item 1
      \item Sub-item 2
    \end{itemize}
\end{itemize}

\subsection{Numbered List}
\begin{enumerate}
  \item First step
  \item Second step
  \item Third step
\end{enumerate}

\section{Citations and References}
When you use the citation search feature, references will be automatically managed. For example, you might cite a paper like this: \cite{example2024}. All citations will be automatically added to your \texttt{cite.bib} file.

\section{Conclusion}
This template demonstrates basic \LaTeX{} features including mathematical formulas, tables, and lists. The bibliography system is now set up to work seamlessly with Bibby AI's citation management. Feel free to modify it for your needs!

\bibliographystyle{plain}
\bibliography{cite}

\end{document}\documentclass{article}
\usepackage[utf8]{inputenc}
\usepackage{booktabs}
\usepackage{amsmath}
\usepackage{natbib}

\title{My Document}
\author{Your Name}
\date{\today}

\begin{document}

\maketitle

\section{Introduction}
This is a sample document demonstrating various \LaTeX{} features including mathematical equations, tables, and formatting. You can cite papers using the citation search feature, and they will be automatically added to your \texttt{cite.bib} file.

\section{Mathematical Expressions}
Here are some example equations:

\subsection{Quadratic Formula}
The quadratic formula is given by:
\[x = \frac{-b \pm \sqrt{b^2 - 4ac}}{2a}\]

\subsection{Matrix Example}
A sample matrix:
\[
\begin{pmatrix}
a_{11} & a_{12} & a_{13} \\
a_{21} & a_{22} & a_{23} \\
a_{31} & a_{32} & a_{33}
\end{pmatrix}
\]

\section{Tables}
Below is a sample table showing data:

\begin{table}[h]
\centering
\begin{tabular}{@{}lrrr@{}}
\toprule
\textbf{Category} & \textbf{Value 1} & \textbf{Value 2} & \textbf{Total} \\
\midrule
Item A & 10.5 & 20.3 & 30.8 \\
Item B & 15.2 & 25.1 & 40.3 \\
Item C & 22.1 & 12.7 & 34.8 \\
\bottomrule
\end{tabular}
\caption{Sample Data Table}
\label{tab:sample}
\end{table}

\section{Lists}
\subsection{Bullet Points}
\begin{itemize}
  \item First item
  \item Second item
  \item Third item with sub-items:
    \begin{itemize}
      \item Sub-item 1
      \item Sub-item 2
    \end{itemize}
\end{itemize}

\subsection{Numbered List}
\begin{enumerate}
  \item First step
  \item Second step
  \item Third step
\end{enumerate}

\section{Citations and References}
When you use the citation search feature, references will be automatically managed. For example, you might cite a paper like this: \cite{example2024}. All citations will be automatically added to your \texttt{cite.bib} file.

\section{Conclusion}
This template demonstrates basic \LaTeX{} features including mathematical formulas, tables, and lists. The bibliography system is now set up to work seamlessly with Bibby AI's citation management. Feel free to modify it for your needs!

\bibliographystyle{plain}
\bibliography{cite}

\end{document}\documentclass{article}
\usepackage[utf8]{inputenc}
\usepackage{booktabs}
\usepackage{amsmath}
\usepackage{natbib}

\title{My Document}
\author{Your Name}
\date{\today}

\begin{document}

\maketitle

\section{Introduction}
This is a sample document demonstrating various \LaTeX{} features including mathematical equations, tables, and formatting. You can cite papers using the citation search feature, and they will be automatically added to your \texttt{cite.bib} file.

\section{Mathematical Expressions}
Here are some example equations:

\subsection{Quadratic Formula}
The quadratic formula is given by:
\[x = \frac{-b \pm \sqrt{b^2 - 4ac}}{2a}\]

\subsection{Matrix Example}
A sample matrix:
\[
\begin{pmatrix}
a_{11} & a_{12} & a_{13} \\
a_{21} & a_{22} & a_{23} \\
a_{31} & a_{32} & a_{33}
\end{pmatrix}
\]

\section{Tables}
Below is a sample table showing data:

\begin{table}[h]
\centering
\begin{tabular}{@{}lrrr@{}}
\toprule
\textbf{Category} & \textbf{Value 1} & \textbf{Value 2} & \textbf{Total} \\
\midrule
Item A & 10.5 & 20.3 & 30.8 \\
Item B & 15.2 & 25.1 & 40.3 \\
Item C & 22.1 & 12.7 & 34.8 \\
\bottomrule
\end{tabular}
\caption{Sample Data Table}
\label{tab:sample}
\end{table}

\section{Lists}
\subsection{Bullet Points}
\begin{itemize}
  \item First item
  \item Second item
  \item Third item with sub-items:
    \begin{itemize}
      \item Sub-item 1
      \item Sub-item 2
    \end{itemize}
\end{itemize}

\subsection{Numbered List}
\begin{enumerate}
  \item First step
  \item Second step
  \item Third step
\end{enumerate}

\section{Citations and References}
When you use the citation search feature, references will be automatically managed. For example, you might cite a paper like this: \cite{example2024}. All citations will be automatically added to your \texttt{cite.bib} file.

\section{Conclusion}
This template demonstrates basic \LaTeX{} features including mathematical formulas, tables, and lists. The bibliography system is now set up to work seamlessly with Bibby AI's citation management. Feel free to modify it for your needs!

\bibliographystyle{plain}
\bibliography{cite}

\end{document}\documentclass{article}
\usepackage[utf8]{inputenc}
\usepackage{booktabs}
\usepackage{amsmath}
\usepackage{natbib}

\title{My Document}
\author{Your Name}
\date{\today}

\begin{document}

\maketitle

\section{Introduction}
This is a sample document demonstrating various \LaTeX{} features including mathematical equations, tables, and formatting. You can cite papers using the citation search feature, and they will be automatically added to your \texttt{cite.bib} file.

\section{Mathematical Expressions}
Here are some example equations:

\subsection{Quadratic Formula}
The quadratic formula is given by:
\[x = \frac{-b \pm \sqrt{b^2 - 4ac}}{2a}\]

\subsection{Matrix Example}
A sample matrix:
\[
\begin{pmatrix}
a_{11} & a_{12} & a_{13} \\
a_{21} & a_{22} & a_{23} \\
a_{31} & a_{32} & a_{33}
\end{pmatrix}
\]

\section{Tables}
Below is a sample table showing data:

\begin{table}[h]
\centering
\begin{tabular}{@{}lrrr@{}}
\toprule
\textbf{Category} & \textbf{Value 1} & \textbf{Value 2} & \textbf{Total} \\
\midrule
Item A & 10.5 & 20.3 & 30.8 \\
Item B & 15.2 & 25.1 & 40.3 \\
Item C & 22.1 & 12.7 & 34.8 \\
\bottomrule
\end{tabular}
\caption{Sample Data Table}
\label{tab:sample}
\end{table}

\section{Lists}
\subsection{Bullet Points}
\begin{itemize}
  \item First item
  \item Second item
  \item Third item with sub-items:
    \begin{itemize}
      \item Sub-item 1
      \item Sub-item 2
    \end{itemize}
\end{itemize}

\subsection{Numbered List}
\begin{enumerate}
  \item First step
  \item Second step
  \item Third step
\end{enumerate}

\section{Citations and References}
When you use the citation search feature, references will be automatically managed. For example, you might cite a paper like this: \cite{example2024}. All citations will be automatically added to your \texttt{cite.bib} file.

\section{Conclusion}
This template demonstrates basic \LaTeX{} features including mathematical formulas, tables, and lists. The bibliography system is now set up to work seamlessly with Bibby AI's citation management. Feel free to modify it for your needs!

\bibliographystyle{plain}
\bibliography{cite}

\end{document}\documentclass{article}
\usepackage[utf8]{inputenc}
\usepackage{booktabs}
\usepackage{amsmath}
\usepackage{natbib}

\title{My Document}
\author{Your Name}
\date{\today}

\begin{document}

\maketitle

\section{Introduction}
This is a sample document demonstrating various \LaTeX{} features including mathematical equations, tables, and formatting. You can cite papers using the citation search feature, and they will be automatically added to your \texttt{cite.bib} file.

\section{Mathematical Expressions}
Here are some example equations:

\subsection{Quadratic Formula}
The quadratic formula is given by:
\[x = \frac{-b \pm \sqrt{b^2 - 4ac}}{2a}\]

\subsection{Matrix Example}
A sample matrix:
\[
\begin{pmatrix}
a_{11} & a_{12} & a_{13} \\
a_{21} & a_{22} & a_{23} \\
a_{31} & a_{32} & a_{33}
\end{pmatrix}
\]

\section{Tables}
Below is a sample table showing data:

\begin{table}[h]
\centering
\begin{tabular}{@{}lrrr@{}}
\toprule
\textbf{Category} & \textbf{Value 1} & \textbf{Value 2} & \textbf{Total} \\
\midrule
Item A & 10.5 & 20.3 & 30.8 \\
Item B & 15.2 & 25.1 & 40.3 \\
Item C & 22.1 & 12.7 & 34.8 \\
\bottomrule
\end{tabular}
\caption{Sample Data Table}
\label{tab:sample}
\end{table}

\section{Lists}
\subsection{Bullet Points}
\begin{itemize}
  \item First item
  \item Second item
  \item Third item with sub-items:
    \begin{itemize}
      \item Sub-item 1
      \item Sub-item 2
    \end{itemize}
\end{itemize}

\subsection{Numbered List}
\begin{enumerate}
  \item First step
  \item Second step
  \item Third step
\end{enumerate}

\section{Citations and References}
When you use the citation search feature, references will be automatically managed. For example, you might cite a paper like this: \cite{example2024}. All citations will be automatically added to your \texttt{cite.bib} file.

\section{Conclusion}
This template demonstrates basic \LaTeX{} features including mathematical formulas, tables, and lists. The bibliography system is now set up to work seamlessly with Bibby AI's citation management. Feel free to modify it for your needs!

\bibliographystyle{plain}
\bibliography{cite}

\end{document}\documentclass{article}
\usepackage[utf8]{inputenc}
\usepackage{booktabs}
\usepackage{amsmath}
\usepackage{natbib}

\title{My Document}
\author{Your Name}
\date{\today}

\begin{document}

\maketitle

\section{Introduction}
This is a sample document demonstrating various \LaTeX{} features including mathematical equations, tables, and formatting. You can cite papers using the citation search feature, and they will be automatically added to your \texttt{cite.bib} file.

\section{Mathematical Expressions}
Here are some example equations:

\subsection{Quadratic Formula}
The quadratic formula is given by:
\[x = \frac{-b \pm \sqrt{b^2 - 4ac}}{2a}\]

\subsection{Matrix Example}
A sample matrix:
\[
\begin{pmatrix}
a_{11} & a_{12} & a_{13} \\
a_{21} & a_{22} & a_{23} \\
a_{31} & a_{32} & a_{33}
\end{pmatrix}
\]

\section{Tables}
Below is a sample table showing data:

\begin{table}[h]
\centering
\begin{tabular}{@{}lrrr@{}}
\toprule
\textbf{Category} & \textbf{Value 1} & \textbf{Value 2} & \textbf{Total} \\
\midrule
Item A & 10.5 & 20.3 & 30.8 \\
Item B & 15.2 & 25.1 & 40.3 \\
Item C & 22.1 & 12.7 & 34.8 \\
\bottomrule
\end{tabular}
\caption{Sample Data Table}
\label{tab:sample}
\end{table}

\section{Lists}
\subsection{Bullet Points}
\begin{itemize}
  \item First item
  \item Second item
  \item Third item with sub-items:
    \begin{itemize}
      \item Sub-item 1
      \item Sub-item 2
    \end{itemize}
\end{itemize}

\subsection{Numbered List}
\begin{enumerate}
  \item First step
  \item Second step
  \item Third step
\end{enumerate}

\section{Citations and References}
When you use the citation search feature, references will be automatically managed. For example, you might cite a paper like this: \cite{example2024}. All citations will be automatically added to your \texttt{cite.bib} file.

\section{Conclusion}
This template demonstrates basic \LaTeX{} features including mathematical formulas, tables, and lists. The bibliography system is now set up to work seamlessly with Bibby AI's citation management. Feel free to modify it for your needs!

\bibliographystyle{plain}
\bibliography{cite}

\end{document}\documentclass{article}
\usepackage[utf8]{inputenc}
\usepackage{booktabs}
\usepackage{amsmath}
\usepackage{natbib}

\title{My Document}
\author{Your Name}
\date{\today}

\begin{document}

\maketitle

\section{Introduction}
This is a sample document demonstrating various \LaTeX{} features including mathematical equations, tables, and formatting. You can cite papers using the citation search feature, and they will be automatically added to your \texttt{cite.bib} file.

\section{Mathematical Expressions}
Here are some example equations:

\subsection{Quadratic Formula}
The quadratic formula is given by:
\[x = \frac{-b \pm \sqrt{b^2 - 4ac}}{2a}\]

\subsection{Matrix Example}
A sample matrix:
\[
\begin{pmatrix}
a_{11} & a_{12} & a_{13} \\
a_{21} & a_{22} & a_{23} \\
a_{31} & a_{32} & a_{33}
\end{pmatrix}
\]

\section{Tables}
Below is a sample table showing data:

\begin{table}[h]
\centering
\begin{tabular}{@{}lrrr@{}}
\toprule
\textbf{Category} & \textbf{Value 1} & \textbf{Value 2} & \textbf{Total} \\
\midrule
Item A & 10.5 & 20.3 & 30.8 \\
Item B & 15.2 & 25.1 & 40.3 \\
Item C & 22.1 & 12.7 & 34.8 \\
\bottomrule
\end{tabular}
\caption{Sample Data Table}
\label{tab:sample}
\end{table}

\section{Lists}
\subsection{Bullet Points}
\begin{itemize}
  \item First item
  \item Second item
  \item Third item with sub-items:
    \begin{itemize}
      \item Sub-item 1
      \item Sub-item 2
    \end{itemize}
\end{itemize}

\subsection{Numbered List}
\begin{enumerate}
  \item First step
  \item Second step
  \item Third step
\end{enumerate}

\section{Citations and References}
When you use the citation search feature, references will be automatically managed. For example, you might cite a paper like this: \cite{example2024}. All citations will be automatically added to your \texttt{cite.bib} file.

\section{Conclusion}
This template demonstrates basic \LaTeX{} features including mathematical formulas, tables, and lists. The bibliography system is now set up to work seamlessly with Bibby AI's citation management. Feel free to modify it for your needs!

\bibliographystyle{plain}
\bibliography{cite}

\end{document}\documentclass{article}
\usepackage[utf8]{inputenc}
\usepackage{booktabs}
\usepackage{amsmath}
\usepackage{natbib}

\title{My Document}
\author{Your Name}
\date{\today}

\begin{document}

\maketitle

\section{Introduction}
This is a sample document demonstrating various \LaTeX{} features including mathematical equations, tables, and formatting. You can cite papers using the citation search feature, and they will be automatically added to your \texttt{cite.bib} file.

\section{Mathematical Expressions}
Here are some example equations:

\subsection{Quadratic Formula}
The quadratic formula is given by:
\[x = \frac{-b \pm \sqrt{b^2 - 4ac}}{2a}\]

\subsection{Matrix Example}
A sample matrix:
\[
\begin{pmatrix}
a_{11} & a_{12} & a_{13} \\
a_{21} & a_{22} & a_{23} \\
a_{31} & a_{32} & a_{33}
\end{pmatrix}
\]

\section{Tables}
Below is a sample table showing data:

\begin{table}[h]
\centering
\begin{tabular}{@{}lrrr@{}}
\toprule
\textbf{Category} & \textbf{Value 1} & \textbf{Value 2} & \textbf{Total} \\
\midrule
Item A & 10.5 & 20.3 & 30.8 \\
Item B & 15.2 & 25.1 & 40.3 \\
Item C & 22.1 & 12.7 & 34.8 \\
\bottomrule
\end{tabular}
\caption{Sample Data Table}
\label{tab:sample}
\end{table}

\section{Lists}
\subsection{Bullet Points}
\begin{itemize}
  \item First item
  \item Second item
  \item Third item with sub-items:
    \begin{itemize}
      \item Sub-item 1
      \item Sub-item 2
    \end{itemize}
\end{itemize}

\subsection{Numbered List}
\begin{enumerate}
  \item First step
  \item Second step
  \item Third step
\end{enumerate}

\section{Citations and References}
When you use the citation search feature, references will be automatically managed. For example, you might cite a paper like this: \cite{example2024}. All citations will be automatically added to your \texttt{cite.bib} file.

\section{Conclusion}
This template demonstrates basic \LaTeX{} features including mathematical formulas, tables, and lists. The bibliography system is now set up to work seamlessly with Bibby AI's citation management. Feel free to modify it for your needs!

\bibliographystyle{plain}
\bibliography{cite}

\end{document}\documentclass{article}
\usepackage[utf8]{inputenc}
\usepackage{booktabs}
\usepackage{amsmath}
\usepackage{natbib}

\title{My Document}
\author{Your Name}
\date{\today}

\begin{document}

\maketitle

\section{Introduction}
This is a sample document demonstrating various \LaTeX{} features including mathematical equations, tables, and formatting. You can cite papers using the citation search feature, and they will be automatically added to your \texttt{cite.bib} file.

\section{Mathematical Expressions}
Here are some example equations:

\subsection{Quadratic Formula}
The quadratic formula is given by:
\[x = \frac{-b \pm \sqrt{b^2 - 4ac}}{2a}\]

\subsection{Matrix Example}
A sample matrix:
\[
\begin{pmatrix}
a_{11} & a_{12} & a_{13} \\
a_{21} & a_{22} & a_{23} \\
a_{31} & a_{32} & a_{33}
\end{pmatrix}
\]

\section{Tables}
Below is a sample table showing data:

\begin{table}[h]
\centering
\begin{tabular}{@{}lrrr@{}}
\toprule
\textbf{Category} & \textbf{Value 1} & \textbf{Value 2} & \textbf{Total} \\
\midrule
Item A & 10.5 & 20.3 & 30.8 \\
Item B & 15.2 & 25.1 & 40.3 \\
Item C & 22.1 & 12.7 & 34.8 \\
\bottomrule
\end{tabular}
\caption{Sample Data Table}
\label{tab:sample}
\end{table}

\section{Lists}
\subsection{Bullet Points}
\begin{itemize}
  \item First item
  \item Second item
  \item Third item with sub-items:
    \begin{itemize}
      \item Sub-item 1
      \item Sub-item 2
    \end{itemize}
\end{itemize}

\subsection{Numbered List}
\begin{enumerate}
  \item First step
  \item Second step
  \item Third step
\end{enumerate}

\section{Citations and References}
When you use the citation search feature, references will be automatically managed. For example, you might cite a paper like this: \cite{example2024}. All citations will be automatically added to your \texttt{cite.bib} file.

\section{Conclusion}
This template demonstrates basic \LaTeX{} features including mathematical formulas, tables, and lists. The bibliography system is now set up to work seamlessly with Bibby AI's citation management. Feel free to modify it for your needs!

\bibliographystyle{plain}
\bibliography{cite}

\end{document}\documentclass{article}
\usepackage[utf8]{inputenc}
\usepackage{booktabs}
\usepackage{amsmath}
\usepackage{natbib}

\title{My Document}
\author{Your Name}
\date{\today}

\begin{document}

\maketitle

\section{Introduction}
This is a sample document demonstrating various \LaTeX{} features including mathematical equations, tables, and formatting. You can cite papers using the citation search feature, and they will be automatically added to your \texttt{cite.bib} file.

\section{Mathematical Expressions}
Here are some example equations:

\subsection{Quadratic Formula}
The quadratic formula is given by:
\[x = \frac{-b \pm \sqrt{b^2 - 4ac}}{2a}\]

\subsection{Matrix Example}
A sample matrix:
\[
\begin{pmatrix}
a_{11} & a_{12} & a_{13} \\
a_{21} & a_{22} & a_{23} \\
a_{31} & a_{32} & a_{33}
\end{pmatrix}
\]

\section{Tables}
Below is a sample table showing data:

\begin{table}[h]
\centering
\begin{tabular}{@{}lrrr@{}}
\toprule
\textbf{Category} & \textbf{Value 1} & \textbf{Value 2} & \textbf{Total} \\
\midrule
Item A & 10.5 & 20.3 & 30.8 \\
Item B & 15.2 & 25.1 & 40.3 \\
Item C & 22.1 & 12.7 & 34.8 \\
\bottomrule
\end{tabular}
\caption{Sample Data Table}
\label{tab:sample}
\end{table}

\section{Lists}
\subsection{Bullet Points}
\begin{itemize}
  \item First item
  \item Second item
  \item Third item with sub-items:
    \begin{itemize}
      \item Sub-item 1
      \item Sub-item 2
    \end{itemize}
\end{itemize}

\subsection{Numbered List}
\begin{enumerate}
  \item First step
  \item Second step
  \item Third step
\end{enumerate}

\section{Citations and References}
When you use the citation search feature, references will be automatically managed. For example, you might cite a paper like this: \cite{example2024}. All citations will be automatically added to your \texttt{cite.bib} file.

\section{Conclusion}
This template demonstrates basic \LaTeX{} features including mathematical formulas, tables, and lists. The bibliography system is now set up to work seamlessly with Bibby AI's citation management. Feel free to modify it for your needs!

\bibliographystyle{plain}
\bibliography{cite}

\end{document}\documentclass{article}
\usepackage[utf8]{inputenc}
\usepackage{booktabs}
\usepackage{amsmath}
\usepackage{natbib}

\title{My Document}
\author{Your Name}
\date{\today}

\begin{document}

\maketitle

\section{Introduction}
This is a sample document demonstrating various \LaTeX{} features including mathematical equations, tables, and formatting. You can cite papers using the citation search feature, and they will be automatically added to your \texttt{cite.bib} file.

\section{Mathematical Expressions}
Here are some example equations:

\subsection{Quadratic Formula}
The quadratic formula is given by:
\[x = \frac{-b \pm \sqrt{b^2 - 4ac}}{2a}\]

\subsection{Matrix Example}
A sample matrix:
\[
\begin{pmatrix}
a_{11} & a_{12} & a_{13} \\
a_{21} & a_{22} & a_{23} \\
a_{31} & a_{32} & a_{33}
\end{pmatrix}
\]

\section{Tables}
Below is a sample table showing data:

\begin{table}[h]
\centering
\begin{tabular}{@{}lrrr@{}}
\toprule
\textbf{Category} & \textbf{Value 1} & \textbf{Value 2} & \textbf{Total} \\
\midrule
Item A & 10.5 & 20.3 & 30.8 \\
Item B & 15.2 & 25.1 & 40.3 \\
Item C & 22.1 & 12.7 & 34.8 \\
\bottomrule
\end{tabular}
\caption{Sample Data Table}
\label{tab:sample}
\end{table}

\section{Lists}
\subsection{Bullet Points}
\begin{itemize}
  \item First item
  \item Second item
  \item Third item with sub-items:
    \begin{itemize}
      \item Sub-item 1
      \item Sub-item 2
    \end{itemize}
\end{itemize}

\subsection{Numbered List}
\begin{enumerate}
  \item First step
  \item Second step
  \item Third step
\end{enumerate}

\section{Citations and References}
When you use the citation search feature, references will be automatically managed. For example, you might cite a paper like this: \cite{example2024}. All citations will be automatically added to your \texttt{cite.bib} file.

\section{Conclusion}
This template demonstrates basic \LaTeX{} features including mathematical formulas, tables, and lists. The bibliography system is now set up to work seamlessly with Bibby AI's citation management. Feel free to modify it for your needs!

\bibliographystyle{plain}
\bibliography{cite}

\end{document}\documentclass{article}
\usepackage[utf8]{inputenc}
\usepackage{booktabs}
\usepackage{amsmath}
\usepackage{natbib}

\title{My Document}
\author{Your Name}
\date{\today}

\begin{document}

\maketitle

\section{Introduction}
This is a sample document demonstrating various \LaTeX{} features including mathematical equations, tables, and formatting. You can cite papers using the citation search feature, and they will be automatically added to your \texttt{cite.bib} file.

\section{Mathematical Expressions}
Here are some example equations:

\subsection{Quadratic Formula}
The quadratic formula is given by:
\[x = \frac{-b \pm \sqrt{b^2 - 4ac}}{2a}\]

\subsection{Matrix Example}
A sample matrix:
\[
\begin{pmatrix}
a_{11} & a_{12} & a_{13} \\
a_{21} & a_{22} & a_{23} \\
a_{31} & a_{32} & a_{33}
\end{pmatrix}
\]

\section{Tables}
Below is a sample table showing data:

\begin{table}[h]
\centering
\begin{tabular}{@{}lrrr@{}}
\toprule
\textbf{Category} & \textbf{Value 1} & \textbf{Value 2} & \textbf{Total} \\
\midrule
Item A & 10.5 & 20.3 & 30.8 \\
Item B & 15.2 & 25.1 & 40.3 \\
Item C & 22.1 & 12.7 & 34.8 \\
\bottomrule
\end{tabular}
\caption{Sample Data Table}
\label{tab:sample}
\end{table}

\section{Lists}
\subsection{Bullet Points}
\begin{itemize}
  \item First item
  \item Second item
  \item Third item with sub-items:
    \begin{itemize}
      \item Sub-item 1
      \item Sub-item 2
    \end{itemize}
\end{itemize}

\subsection{Numbered List}
\begin{enumerate}
  \item First step
  \item Second step
  \item Third step
\end{enumerate}

\section{Citations and References}
When you use the citation search feature, references will be automatically managed. For example, you might cite a paper like this: \cite{example2024}. All citations will be automatically added to your \texttt{cite.bib} file.

\section{Conclusion}
This template demonstrates basic \LaTeX{} features including mathematical formulas, tables, and lists. The bibliography system is now set up to work seamlessly with Bibby AI's citation management. Feel free to modify it for your needs!

\bibliographystyle{plain}
\bibliography{cite}

\end{document}\documentclass{article}
\usepackage[utf8]{inputenc}
\usepackage{booktabs}
\usepackage{amsmath}
\usepackage{natbib}

\title{My Document}
\author{Your Name}
\date{\today}

\begin{document}

\maketitle

\section{Introduction}
This is a sample document demonstrating various \LaTeX{} features including mathematical equations, tables, and formatting. You can cite papers using the citation search feature, and they will be automatically added to your \texttt{cite.bib} file.

\section{Mathematical Expressions}
Here are some example equations:

\subsection{Quadratic Formula}
The quadratic formula is given by:
\[x = \frac{-b \pm \sqrt{b^2 - 4ac}}{2a}\]

\subsection{Matrix Example}
A sample matrix:
\[
\begin{pmatrix}
a_{11} & a_{12} & a_{13} \\
a_{21} & a_{22} & a_{23} \\
a_{31} & a_{32} & a_{33}
\end{pmatrix}
\]

\section{Tables}
Below is a sample table showing data:

\begin{table}[h]
\centering
\begin{tabular}{@{}lrrr@{}}
\toprule
\textbf{Category} & \textbf{Value 1} & \textbf{Value 2} & \textbf{Total} \\
\midrule
Item A & 10.5 & 20.3 & 30.8 \\
Item B & 15.2 & 25.1 & 40.3 \\
Item C & 22.1 & 12.7 & 34.8 \\
\bottomrule
\end{tabular}
\caption{Sample Data Table}
\label{tab:sample}
\end{table}

\section{Lists}
\subsection{Bullet Points}
\begin{itemize}
  \item First item
  \item Second item
  \item Third item with sub-items:
    \begin{itemize}
      \item Sub-item 1
      \item Sub-item 2
    \end{itemize}
\end{itemize}

\subsection{Numbered List}
\begin{enumerate}
  \item First step
  \item Second step
  \item Third step
\end{enumerate}

\section{Citations and References}
When you use the citation search feature, references will be automatically managed. For example, you might cite a paper like this: \cite{example2024}. All citations will be automatically added to your \texttt{cite.bib} file.

\section{Conclusion}
This template demonstrates basic \LaTeX{} features including mathematical formulas, tables, and lists. The bibliography system is now set up to work seamlessly with Bibby AI's citation management. Feel free to modify it for your needs!

\bibliographystyle{plain}
\bibliography{cite}

\end{document}\documentclass{article}
\usepackage[utf8]{inputenc}
\usepackage{booktabs}
\usepackage{amsmath}
\usepackage{natbib}

\title{My Document}
\author{Your Name}
\date{\today}

\begin{document}

\maketitle

\section{Introduction}
This is a sample document demonstrating various \LaTeX{} features including mathematical equations, tables, and formatting. You can cite papers using the citation search feature, and they will be automatically added to your \texttt{cite.bib} file.

\section{Mathematical Expressions}
Here are some example equations:

\subsection{Quadratic Formula}
The quadratic formula is given by:
\[x = \frac{-b \pm \sqrt{b^2 - 4ac}}{2a}\]

\subsection{Matrix Example}
A sample matrix:
\[
\begin{pmatrix}
a_{11} & a_{12} & a_{13} \\
a_{21} & a_{22} & a_{23} \\
a_{31} & a_{32} & a_{33}
\end{pmatrix}
\]

\section{Tables}
Below is a sample table showing data:

\begin{table}[h]
\centering
\begin{tabular}{@{}lrrr@{}}
\toprule
\textbf{Category} & \textbf{Value 1} & \textbf{Value 2} & \textbf{Total} \\
\midrule
Item A & 10.5 & 20.3 & 30.8 \\
Item B & 15.2 & 25.1 & 40.3 \\
Item C & 22.1 & 12.7 & 34.8 \\
\bottomrule
\end{tabular}
\caption{Sample Data Table}
\label{tab:sample}
\end{table}

\section{Lists}
\subsection{Bullet Points}
\begin{itemize}
  \item First item
  \item Second item
  \item Third item with sub-items:
    \begin{itemize}
      \item Sub-item 1
      \item Sub-item 2
    \end{itemize}
\end{itemize}

\subsection{Numbered List}
\begin{enumerate}
  \item First step
  \item Second step
  \item Third step
\end{enumerate}

\section{Citations and References}
When you use the citation search feature, references will be automatically managed. For example, you might cite a paper like this: \cite{example2024}. All citations will be automatically added to your \texttt{cite.bib} file.

\section{Conclusion}
This template demonstrates basic \LaTeX{} features including mathematical formulas, tables, and lists. The bibliography system is now set up to work seamlessly with Bibby AI's citation management. Feel free to modify it for your needs!

\bibliographystyle{plain}
\bibliography{cite}

\end{document}\documentclass{article}
\usepackage[utf8]{inputenc}
\usepackage{booktabs}
\usepackage{amsmath}
\usepackage{natbib}

\title{My Document}
\author{Your Name}
\date{\today}

\begin{document}

\maketitle

\section{Introduction}
This is a sample document demonstrating various \LaTeX{} features including mathematical equations, tables, and formatting. You can cite papers using the citation search feature, and they will be automatically added to your \texttt{cite.bib} file.

\section{Mathematical Expressions}
Here are some example equations:

\subsection{Quadratic Formula}
The quadratic formula is given by:
\[x = \frac{-b \pm \sqrt{b^2 - 4ac}}{2a}\]

\subsection{Matrix Example}
A sample matrix:
\[
\begin{pmatrix}
a_{11} & a_{12} & a_{13} \\
a_{21} & a_{22} & a_{23} \\
a_{31} & a_{32} & a_{33}
\end{pmatrix}
\]

\section{Tables}
Below is a sample table showing data:

\begin{table}[h]
\centering
\begin{tabular}{@{}lrrr@{}}
\toprule
\textbf{Category} & \textbf{Value 1} & \textbf{Value 2} & \textbf{Total} \\
\midrule
Item A & 10.5 & 20.3 & 30.8 \\
Item B & 15.2 & 25.1 & 40.3 \\
Item C & 22.1 & 12.7 & 34.8 \\
\bottomrule
\end{tabular}
\caption{Sample Data Table}
\label{tab:sample}
\end{table}

\section{Lists}
\subsection{Bullet Points}
\begin{itemize}
  \item First item
  \item Second item
  \item Third item with sub-items:
    \begin{itemize}
      \item Sub-item 1
      \item Sub-item 2
    \end{itemize}
\end{itemize}

\subsection{Numbered List}
\begin{enumerate}
  \item First step
  \item Second step
  \item Third step
\end{enumerate}

\section{Citations and References}
When you use the citation search feature, references will be automatically managed. For example, you might cite a paper like this: \cite{example2024}. All citations will be automatically added to your \texttt{cite.bib} file.

\section{Conclusion}
This template demonstrates basic \LaTeX{} features including mathematical formulas, tables, and lists. The bibliography system is now set up to work seamlessly with Bibby AI's citation management. Feel free to modify it for your needs!

\bibliographystyle{plain}
\bibliography{cite}

\end{document}\documentclass{article}
\usepackage[utf8]{inputenc}
\usepackage{booktabs}
\usepackage{amsmath}
\usepackage{natbib}

\title{My Document}
\author{Your Name}
\date{\today}

\begin{document}

\maketitle

\section{Introduction}
This is a sample document demonstrating various \LaTeX{} features including mathematical equations, tables, and formatting. You can cite papers using the citation search feature, and they will be automatically added to your \texttt{cite.bib} file.

\section{Mathematical Expressions}
Here are some example equations:

\subsection{Quadratic Formula}
The quadratic formula is given by:
\[x = \frac{-b \pm \sqrt{b^2 - 4ac}}{2a}\]

\subsection{Matrix Example}
A sample matrix:
\[
\begin{pmatrix}
a_{11} & a_{12} & a_{13} \\
a_{21} & a_{22} & a_{23} \\
a_{31} & a_{32} & a_{33}
\end{pmatrix}
\]

\section{Tables}
Below is a sample table showing data:

\begin{table}[h]
\centering
\begin{tabular}{@{}lrrr@{}}
\toprule
\textbf{Category} & \textbf{Value 1} & \textbf{Value 2} & \textbf{Total} \\
\midrule
Item A & 10.5 & 20.3 & 30.8 \\
Item B & 15.2 & 25.1 & 40.3 \\
Item C & 22.1 & 12.7 & 34.8 \\
\bottomrule
\end{tabular}
\caption{Sample Data Table}
\label{tab:sample}
\end{table}

\section{Lists}
\subsection{Bullet Points}
\begin{itemize}
  \item First item
  \item Second item
  \item Third item with sub-items:
    \begin{itemize}
      \item Sub-item 1
      \item Sub-item 2
    \end{itemize}
\end{itemize}

\subsection{Numbered List}
\begin{enumerate}
  \item First step
  \item Second step
  \item Third step
\end{enumerate}

\section{Citations and References}
When you use the citation search feature, references will be automatically managed. For example, you might cite a paper like this: \cite{example2024}. All citations will be automatically added to your \texttt{cite.bib} file.

\section{Conclusion}
This template demonstrates basic \LaTeX{} features including mathematical formulas, tables, and lists. The bibliography system is now set up to work seamlessly with Bibby AI's citation management. Feel free to modify it for your needs!

\bibliographystyle{plain}
\bibliography{cite}

\end{document}\documentclass{article}
\usepackage[utf8]{inputenc}
\usepackage{booktabs}
\usepackage{amsmath}
\usepackage{natbib}

\title{My Document}
\author{Your Name}
\date{\today}

\begin{document}

\maketitle

\section{Introduction}
This is a sample document demonstrating various \LaTeX{} features including mathematical equations, tables, and formatting. You can cite papers using the citation search feature, and they will be automatically added to your \texttt{cite.bib} file.

\section{Mathematical Expressions}
Here are some example equations:

\subsection{Quadratic Formula}
The quadratic formula is given by:
\[x = \frac{-b \pm \sqrt{b^2 - 4ac}}{2a}\]

\subsection{Matrix Example}
A sample matrix:
\[
\begin{pmatrix}
a_{11} & a_{12} & a_{13} \\
a_{21} & a_{22} & a_{23} \\
a_{31} & a_{32} & a_{33}
\end{pmatrix}
\]

\section{Tables}
Below is a sample table showing data:

\begin{table}[h]
\centering
\begin{tabular}{@{}lrrr@{}}
\toprule
\textbf{Category} & \textbf{Value 1} & \textbf{Value 2} & \textbf{Total} \\
\midrule
Item A & 10.5 & 20.3 & 30.8 \\
Item B & 15.2 & 25.1 & 40.3 \\
Item C & 22.1 & 12.7 & 34.8 \\
\bottomrule
\end{tabular}
\caption{Sample Data Table}
\label{tab:sample}
\end{table}

\section{Lists}
\subsection{Bullet Points}
\begin{itemize}
  \item First item
  \item Second item
  \item Third item with sub-items:
    \begin{itemize}
      \item Sub-item 1
      \item Sub-item 2
    \end{itemize}
\end{itemize}

\subsection{Numbered List}
\begin{enumerate}
  \item First step
  \item Second step
  \item Third step
\end{enumerate}

\section{Citations and References}
When you use the citation search feature, references will be automatically managed. For example, you might cite a paper like this: \cite{example2024}. All citations will be automatically added to your \texttt{cite.bib} file.

\section{Conclusion}
This template demonstrates basic \LaTeX{} features including mathematical formulas, tables, and lists. The bibliography system is now set up to work seamlessly with Bibby AI's citation management. Feel free to modify it for your needs!

\bibliographystyle{plain}
\bibliography{cite}

\end{document}\documentclass{article}
\usepackage[utf8]{inputenc}
\usepackage{booktabs}
\usepackage{amsmath}
\usepackage{natbib}

\title{My Document}
\author{Your Name}
\date{\today}

\begin{document}

\maketitle

\section{Introduction}
This is a sample document demonstrating various \LaTeX{} features including mathematical equations, tables, and formatting. You can cite papers using the citation search feature, and they will be automatically added to your \texttt{cite.bib} file.

\section{Mathematical Expressions}
Here are some example equations:

\subsection{Quadratic Formula}
The quadratic formula is given by:
\[x = \frac{-b \pm \sqrt{b^2 - 4ac}}{2a}\]

\subsection{Matrix Example}
A sample matrix:
\[
\begin{pmatrix}
a_{11} & a_{12} & a_{13} \\
a_{21} & a_{22} & a_{23} \\
a_{31} & a_{32} & a_{33}
\end{pmatrix}
\]

\section{Tables}
Below is a sample table showing data:

\begin{table}[h]
\centering
\begin{tabular}{@{}lrrr@{}}
\toprule
\textbf{Category} & \textbf{Value 1} & \textbf{Value 2} & \textbf{Total} \\
\midrule
Item A & 10.5 & 20.3 & 30.8 \\
Item B & 15.2 & 25.1 & 40.3 \\
Item C & 22.1 & 12.7 & 34.8 \\
\bottomrule
\end{tabular}
\caption{Sample Data Table}
\label{tab:sample}
\end{table}

\section{Lists}
\subsection{Bullet Points}
\begin{itemize}
  \item First item
  \item Second item
  \item Third item with sub-items:
    \begin{itemize}
      \item Sub-item 1
      \item Sub-item 2
    \end{itemize}
\end{itemize}

\subsection{Numbered List}
\begin{enumerate}
  \item First step
  \item Second step
  \item Third step
\end{enumerate}

\section{Citations and References}
When you use the citation search feature, references will be automatically managed. For example, you might cite a paper like this: \cite{example2024}. All citations will be automatically added to your \texttt{cite.bib} file.

\section{Conclusion}
This template demonstrates basic \LaTeX{} features including mathematical formulas, tables, and lists. The bibliography system is now set up to work seamlessly with Bibby AI's citation management. Feel free to modify it for your needs!

\bibliographystyle{plain}
\bibliography{cite}

\end{document}\documentclass{article}
\usepackage[utf8]{inputenc}
\usepackage{booktabs}
\usepackage{amsmath}
\usepackage{natbib}

\title{My Document}
\author{Your Name}
\date{\today}

\begin{document}

\maketitle

\section{Introduction}
This is a sample document demonstrating various \LaTeX{} features including mathematical equations, tables, and formatting. You can cite papers using the citation search feature, and they will be automatically added to your \texttt{cite.bib} file.

\section{Mathematical Expressions}
Here are some example equations:

\subsection{Quadratic Formula}
The quadratic formula is given by:
\[x = \frac{-b \pm \sqrt{b^2 - 4ac}}{2a}\]

\subsection{Matrix Example}
A sample matrix:
\[
\begin{pmatrix}
a_{11} & a_{12} & a_{13} \\
a_{21} & a_{22} & a_{23} \\
a_{31} & a_{32} & a_{33}
\end{pmatrix}
\]

\section{Tables}
Below is a sample table showing data:

\begin{table}[h]
\centering
\begin{tabular}{@{}lrrr@{}}
\toprule
\textbf{Category} & \textbf{Value 1} & \textbf{Value 2} & \textbf{Total} \\
\midrule
Item A & 10.5 & 20.3 & 30.8 \\
Item B & 15.2 & 25.1 & 40.3 \\
Item C & 22.1 & 12.7 & 34.8 \\
\bottomrule
\end{tabular}
\caption{Sample Data Table}
\label{tab:sample}
\end{table}

\section{Lists}
\subsection{Bullet Points}
\begin{itemize}
  \item First item
  \item Second item
  \item Third item with sub-items:
    \begin{itemize}
      \item Sub-item 1
      \item Sub-item 2
    \end{itemize}
\end{itemize}

\subsection{Numbered List}
\begin{enumerate}
  \item First step
  \item Second step
  \item Third step
\end{enumerate}

\section{Citations and References}
When you use the citation search feature, references will be automatically managed. For example, you might cite a paper like this: \cite{example2024}. All citations will be automatically added to your \texttt{cite.bib} file.

\section{Conclusion}
This template demonstrates basic \LaTeX{} features including mathematical formulas, tables, and lists. The bibliography system is now set up to work seamlessly with Bibby AI's citation management. Feel free to modify it for your needs!

\bibliographystyle{plain}
\bibliography{cite}

\end{document}\documentclass{article}
\usepackage[utf8]{inputenc}
\usepackage{booktabs}
\usepackage{amsmath}
\usepackage{natbib}

\title{My Document}
\author{Your Name}
\date{\today}

\begin{document}

\maketitle

\section{Introduction}
This is a sample document demonstrating various \LaTeX{} features including mathematical equations, tables, and formatting. You can cite papers using the citation search feature, and they will be automatically added to your \texttt{cite.bib} file.

\section{Mathematical Expressions}
Here are some example equations:

\subsection{Quadratic Formula}
The quadratic formula is given by:
\[x = \frac{-b \pm \sqrt{b^2 - 4ac}}{2a}\]

\subsection{Matrix Example}
A sample matrix:
\[
\begin{pmatrix}
a_{11} & a_{12} & a_{13} \\
a_{21} & a_{22} & a_{23} \\
a_{31} & a_{32} & a_{33}
\end{pmatrix}
\]

\section{Tables}
Below is a sample table showing data:

\begin{table}[h]
\centering
\begin{tabular}{@{}lrrr@{}}
\toprule
\textbf{Category} & \textbf{Value 1} & \textbf{Value 2} & \textbf{Total} \\
\midrule
Item A & 10.5 & 20.3 & 30.8 \\
Item B & 15.2 & 25.1 & 40.3 \\
Item C & 22.1 & 12.7 & 34.8 \\
\bottomrule
\end{tabular}
\caption{Sample Data Table}
\label{tab:sample}
\end{table}

\section{Lists}
\subsection{Bullet Points}
\begin{itemize}
  \item First item
  \item Second item
  \item Third item with sub-items:
    \begin{itemize}
      \item Sub-item 1
      \item Sub-item 2
    \end{itemize}
\end{itemize}

\subsection{Numbered List}
\begin{enumerate}
  \item First step
  \item Second step
  \item Third step
\end{enumerate}

\section{Citations and References}
When you use the citation search feature, references will be automatically managed. For example, you might cite a paper like this: \cite{example2024}. All citations will be automatically added to your \texttt{cite.bib} file.

\section{Conclusion}
This template demonstrates basic \LaTeX{} features including mathematical formulas, tables, and lists. The bibliography system is now set up to work seamlessly with Bibby AI's citation management. Feel free to modify it for your needs!

\bibliographystyle{plain}
\bibliography{cite}

\end{document}\documentclass{article}
\usepackage[utf8]{inputenc}
\usepackage{booktabs}
\usepackage{amsmath}
\usepackage{natbib}

\title{My Document}
\author{Your Name}
\date{\today}

\begin{document}

\maketitle

\section{Introduction}
This is a sample document demonstrating various \LaTeX{} features including mathematical equations, tables, and formatting. You can cite papers using the citation search feature, and they will be automatically added to your \texttt{cite.bib} file.

\section{Mathematical Expressions}
Here are some example equations:

\subsection{Quadratic Formula}
The quadratic formula is given by:
\[x = \frac{-b \pm \sqrt{b^2 - 4ac}}{2a}\]

\subsection{Matrix Example}
A sample matrix:
\[
\begin{pmatrix}
a_{11} & a_{12} & a_{13} \\
a_{21} & a_{22} & a_{23} \\
a_{31} & a_{32} & a_{33}
\end{pmatrix}
\]

\section{Tables}
Below is a sample table showing data:

\begin{table}[h]
\centering
\begin{tabular}{@{}lrrr@{}}
\toprule
\textbf{Category} & \textbf{Value 1} & \textbf{Value 2} & \textbf{Total} \\
\midrule
Item A & 10.5 & 20.3 & 30.8 \\
Item B & 15.2 & 25.1 & 40.3 \\
Item C & 22.1 & 12.7 & 34.8 \\
\bottomrule
\end{tabular}
\caption{Sample Data Table}
\label{tab:sample}
\end{table}

\section{Lists}
\subsection{Bullet Points}
\begin{itemize}
  \item First item
  \item Second item
  \item Third item with sub-items:
    \begin{itemize}
      \item Sub-item 1
      \item Sub-item 2
    \end{itemize}
\end{itemize}

\subsection{Numbered List}
\begin{enumerate}
  \item First step
  \item Second step
  \item Third step
\end{enumerate}

\section{Citations and References}
When you use the citation search feature, references will be automatically managed. For example, you might cite a paper like this: \cite{example2024}. All citations will be automatically added to your \texttt{cite.bib} file.

\section{Conclusion}
This template demonstrates basic \LaTeX{} features including mathematical formulas, tables, and lists. The bibliography system is now set up to work seamlessly with Bibby AI's citation management. Feel free to modify it for your needs!

\bibliographystyle{plain}
\bibliography{cite}

\end{document}\documentclass{article}
\usepackage[utf8]{inputenc}
\usepackage{booktabs}
\usepackage{amsmath}
\usepackage{natbib}

\title{My Document}
\author{Your Name}
\date{\today}

\begin{document}

\maketitle

\section{Introduction}
This is a sample document demonstrating various \LaTeX{} features including mathematical equations, tables, and formatting. You can cite papers using the citation search feature, and they will be automatically added to your \texttt{cite.bib} file.

\section{Mathematical Expressions}
Here are some example equations:

\subsection{Quadratic Formula}
The quadratic formula is given by:
\[x = \frac{-b \pm \sqrt{b^2 - 4ac}}{2a}\]

\subsection{Matrix Example}
A sample matrix:
\[
\begin{pmatrix}
a_{11} & a_{12} & a_{13} \\
a_{21} & a_{22} & a_{23} \\
a_{31} & a_{32} & a_{33}
\end{pmatrix}
\]

\section{Tables}
Below is a sample table showing data:

\begin{table}[h]
\centering
\begin{tabular}{@{}lrrr@{}}
\toprule
\textbf{Category} & \textbf{Value 1} & \textbf{Value 2} & \textbf{Total} \\
\midrule
Item A & 10.5 & 20.3 & 30.8 \\
Item B & 15.2 & 25.1 & 40.3 \\
Item C & 22.1 & 12.7 & 34.8 \\
\bottomrule
\end{tabular}
\caption{Sample Data Table}
\label{tab:sample}
\end{table}

\section{Lists}
\subsection{Bullet Points}
\begin{itemize}
  \item First item
  \item Second item
  \item Third item with sub-items:
    \begin{itemize}
      \item Sub-item 1
      \item Sub-item 2
    \end{itemize}
\end{itemize}

\subsection{Numbered List}
\begin{enumerate}
  \item First step
  \item Second step
  \item Third step
\end{enumerate}

\section{Citations and References}
When you use the citation search feature, references will be automatically managed. For example, you might cite a paper like this: \cite{example2024}. All citations will be automatically added to your \texttt{cite.bib} file.

\section{Conclusion}
This template demonstrates basic \LaTeX{} features including mathematical formulas, tables, and lists. The bibliography system is now set up to work seamlessly with Bibby AI's citation management. Feel free to modify it for your needs!

\bibliographystyle{plain}
\bibliography{cite}

\end{document}\documentclass{article}
\usepackage[utf8]{inputenc}
\usepackage{booktabs}
\usepackage{amsmath}
\usepackage{natbib}

\title{My Document}
\author{Your Name}
\date{\today}

\begin{document}

\maketitle

\section{Introduction}
This is a sample document demonstrating various \LaTeX{} features including mathematical equations, tables, and formatting. You can cite papers using the citation search feature, and they will be automatically added to your \texttt{cite.bib} file.

\section{Mathematical Expressions}
Here are some example equations:

\subsection{Quadratic Formula}
The quadratic formula is given by:
\[x = \frac{-b \pm \sqrt{b^2 - 4ac}}{2a}\]

\subsection{Matrix Example}
A sample matrix:
\[
\begin{pmatrix}
a_{11} & a_{12} & a_{13} \\
a_{21} & a_{22} & a_{23} \\
a_{31} & a_{32} & a_{33}
\end{pmatrix}
\]

\section{Tables}
Below is a sample table showing data:

\begin{table}[h]
\centering
\begin{tabular}{@{}lrrr@{}}
\toprule
\textbf{Category} & \textbf{Value 1} & \textbf{Value 2} & \textbf{Total} \\
\midrule
Item A & 10.5 & 20.3 & 30.8 \\
Item B & 15.2 & 25.1 & 40.3 \\
Item C & 22.1 & 12.7 & 34.8 \\
\bottomrule
\end{tabular}
\caption{Sample Data Table}
\label{tab:sample}
\end{table}

\section{Lists}
\subsection{Bullet Points}
\begin{itemize}
  \item First item
  \item Second item
  \item Third item with sub-items:
    \begin{itemize}
      \item Sub-item 1
      \item Sub-item 2
    \end{itemize}
\end{itemize}

\subsection{Numbered List}
\begin{enumerate}
  \item First step
  \item Second step
  \item Third step
\end{enumerate}

\section{Citations and References}
When you use the citation search feature, references will be automatically managed. For example, you might cite a paper like this: \cite{example2024}. All citations will be automatically added to your \texttt{cite.bib} file.

\section{Conclusion}
This template demonstrates basic \LaTeX{} features including mathematical formulas, tables, and lists. The bibliography system is now set up to work seamlessly with Bibby AI's citation management. Feel free to modify it for your needs!

\bibliographystyle{plain}
\bibliography{cite}

\end{document}\documentclass{article}
\usepackage[utf8]{inputenc}
\usepackage{booktabs}
\usepackage{amsmath}
\usepackage{natbib}

\title{My Document}
\author{Your Name}
\date{\today}

\begin{document}

\maketitle

\section{Introduction}
This is a sample document demonstrating various \LaTeX{} features including mathematical equations, tables, and formatting. You can cite papers using the citation search feature, and they will be automatically added to your \texttt{cite.bib} file.

\section{Mathematical Expressions}
Here are some example equations:

\subsection{Quadratic Formula}
The quadratic formula is given by:
\[x = \frac{-b \pm \sqrt{b^2 - 4ac}}{2a}\]

\subsection{Matrix Example}
A sample matrix:
\[
\begin{pmatrix}
a_{11} & a_{12} & a_{13} \\
a_{21} & a_{22} & a_{23} \\
a_{31} & a_{32} & a_{33}
\end{pmatrix}
\]

\section{Tables}
Below is a sample table showing data:

\begin{table}[h]
\centering
\begin{tabular}{@{}lrrr@{}}
\toprule
\textbf{Category} & \textbf{Value 1} & \textbf{Value 2} & \textbf{Total} \\
\midrule
Item A & 10.5 & 20.3 & 30.8 \\
Item B & 15.2 & 25.1 & 40.3 \\
Item C & 22.1 & 12.7 & 34.8 \\
\bottomrule
\end{tabular}
\caption{Sample Data Table}
\label{tab:sample}
\end{table}

\section{Lists}
\subsection{Bullet Points}
\begin{itemize}
  \item First item
  \item Second item
  \item Third item with sub-items:
    \begin{itemize}
      \item Sub-item 1
      \item Sub-item 2
    \end{itemize}
\end{itemize}

\subsection{Numbered List}
\begin{enumerate}
  \item First step
  \item Second step
  \item Third step
\end{enumerate}

\section{Citations and References}
When you use the citation search feature, references will be automatically managed. For example, you might cite a paper like this: \cite{example2024}. All citations will be automatically added to your \texttt{cite.bib} file.

\section{Conclusion}
This template demonstrates basic \LaTeX{} features including mathematical formulas, tables, and lists. The bibliography system is now set up to work seamlessly with Bibby AI's citation management. Feel free to modify it for your needs!

\bibliographystyle{plain}
\bibliography{cite}

\end{document}\documentclass{article}
\usepackage[utf8]{inputenc}
\usepackage{booktabs}
\usepackage{amsmath}
\usepackage{natbib}

\title{My Document}
\author{Your Name}
\date{\today}

\begin{document}

\maketitle

\section{Introduction}
This is a sample document demonstrating various \LaTeX{} features including mathematical equations, tables, and formatting. You can cite papers using the citation search feature, and they will be automatically added to your \texttt{cite.bib} file.

\section{Mathematical Expressions}
Here are some example equations:

\subsection{Quadratic Formula}
The quadratic formula is given by:
\[x = \frac{-b \pm \sqrt{b^2 - 4ac}}{2a}\]

\subsection{Matrix Example}
A sample matrix:
\[
\begin{pmatrix}
a_{11} & a_{12} & a_{13} \\
a_{21} & a_{22} & a_{23} \\
a_{31} & a_{32} & a_{33}
\end{pmatrix}
\]

\section{Tables}
Below is a sample table showing data:

\begin{table}[h]
\centering
\begin{tabular}{@{}lrrr@{}}
\toprule
\textbf{Category} & \textbf{Value 1} & \textbf{Value 2} & \textbf{Total} \\
\midrule
Item A & 10.5 & 20.3 & 30.8 \\
Item B & 15.2 & 25.1 & 40.3 \\
Item C & 22.1 & 12.7 & 34.8 \\
\bottomrule
\end{tabular}
\caption{Sample Data Table}
\label{tab:sample}
\end{table}

\section{Lists}
\subsection{Bullet Points}
\begin{itemize}
  \item First item
  \item Second item
  \item Third item with sub-items:
    \begin{itemize}
      \item Sub-item 1
      \item Sub-item 2
    \end{itemize}
\end{itemize}

\subsection{Numbered List}
\begin{enumerate}
  \item First step
  \item Second step
  \item Third step
\end{enumerate}

\section{Citations and References}
When you use the citation search feature, references will be automatically managed. For example, you might cite a paper like this: \cite{example2024}. All citations will be automatically added to your \texttt{cite.bib} file.

\section{Conclusion}
This template demonstrates basic \LaTeX{} features including mathematical formulas, tables, and lists. The bibliography system is now set up to work seamlessly with Bibby AI's citation management. Feel free to modify it for your needs!

\bibliographystyle{plain}
\bibliography{cite}

\end{document}\documentclass{article}
\usepackage[utf8]{inputenc}
\usepackage{booktabs}
\usepackage{amsmath}
\usepackage{natbib}

\title{My Document}
\author{Your Name}
\date{\today}

\begin{document}

\maketitle

\section{Introduction}
This is a sample document demonstrating various \LaTeX{} features including mathematical equations, tables, and formatting. You can cite papers using the citation search feature, and they will be automatically added to your \texttt{cite.bib} file.

\section{Mathematical Expressions}
Here are some example equations:

\subsection{Quadratic Formula}
The quadratic formula is given by:
\[x = \frac{-b \pm \sqrt{b^2 - 4ac}}{2a}\]

\subsection{Matrix Example}
A sample matrix:
\[
\begin{pmatrix}
a_{11} & a_{12} & a_{13} \\
a_{21} & a_{22} & a_{23} \\
a_{31} & a_{32} & a_{33}
\end{pmatrix}
\]

\section{Tables}
Below is a sample table showing data:

\begin{table}[h]
\centering
\begin{tabular}{@{}lrrr@{}}
\toprule
\textbf{Category} & \textbf{Value 1} & \textbf{Value 2} & \textbf{Total} \\
\midrule
Item A & 10.5 & 20.3 & 30.8 \\
Item B & 15.2 & 25.1 & 40.3 \\
Item C & 22.1 & 12.7 & 34.8 \\
\bottomrule
\end{tabular}
\caption{Sample Data Table}
\label{tab:sample}
\end{table}

\section{Lists}
\subsection{Bullet Points}
\begin{itemize}
  \item First item
  \item Second item
  \item Third item with sub-items:
    \begin{itemize}
      \item Sub-item 1
      \item Sub-item 2
    \end{itemize}
\end{itemize}

\subsection{Numbered List}
\begin{enumerate}
  \item First step
  \item Second step
  \item Third step
\end{enumerate}

\section{Citations and References}
When you use the citation search feature, references will be automatically managed. For example, you might cite a paper like this: \cite{example2024}. All citations will be automatically added to your \texttt{cite.bib} file.

\section{Conclusion}
This template demonstrates basic \LaTeX{} features including mathematical formulas, tables, and lists. The bibliography system is now set up to work seamlessly with Bibby AI's citation management. Feel free to modify it for your needs!

\bibliographystyle{plain}
\bibliography{cite}

\end{document}\documentclass{article}
\usepackage[utf8]{inputenc}
\usepackage{booktabs}
\usepackage{amsmath}
\usepackage{natbib}

\title{My Document}
\author{Your Name}
\date{\today}

\begin{document}

\maketitle

\section{Introduction}
This is a sample document demonstrating various \LaTeX{} features including mathematical equations, tables, and formatting. You can cite papers using the citation search feature, and they will be automatically added to your \texttt{cite.bib} file.

\section{Mathematical Expressions}
Here are some example equations:

\subsection{Quadratic Formula}
The quadratic formula is given by:
\[x = \frac{-b \pm \sqrt{b^2 - 4ac}}{2a}\]

\subsection{Matrix Example}
A sample matrix:
\[
\begin{pmatrix}
a_{11} & a_{12} & a_{13} \\
a_{21} & a_{22} & a_{23} \\
a_{31} & a_{32} & a_{33}
\end{pmatrix}
\]

\section{Tables}
Below is a sample table showing data:

\begin{table}[h]
\centering
\begin{tabular}{@{}lrrr@{}}
\toprule
\textbf{Category} & \textbf{Value 1} & \textbf{Value 2} & \textbf{Total} \\
\midrule
Item A & 10.5 & 20.3 & 30.8 \\
Item B & 15.2 & 25.1 & 40.3 \\
Item C & 22.1 & 12.7 & 34.8 \\
\bottomrule
\end{tabular}
\caption{Sample Data Table}
\label{tab:sample}
\end{table}

\section{Lists}
\subsection{Bullet Points}
\begin{itemize}
  \item First item
  \item Second item
  \item Third item with sub-items:
    \begin{itemize}
      \item Sub-item 1
      \item Sub-item 2
    \end{itemize}
\end{itemize}

\subsection{Numbered List}
\begin{enumerate}
  \item First step
  \item Second step
  \item Third step
\end{enumerate}

\section{Citations and References}
When you use the citation search feature, references will be automatically managed. For example, you might cite a paper like this: \cite{example2024}. All citations will be automatically added to your \texttt{cite.bib} file.

\section{Conclusion}
This template demonstrates basic \LaTeX{} features including mathematical formulas, tables, and lists. The bibliography system is now set up to work seamlessly with Bibby AI's citation management. Feel free to modify it for your needs!

\bibliographystyle{plain}
\bibliography{cite}

\end{document}\documentclass{article}
\usepackage[utf8]{inputenc}
\usepackage{booktabs}
\usepackage{amsmath}
\usepackage{natbib}

\title{My Document}
\author{Your Name}
\date{\today}

\begin{document}

\maketitle

\section{Introduction}
This is a sample document demonstrating various \LaTeX{} features including mathematical equations, tables, and formatting. You can cite papers using the citation search feature, and they will be automatically added to your \texttt{cite.bib} file.

\section{Mathematical Expressions}
Here are some example equations:

\subsection{Quadratic Formula}
The quadratic formula is given by:
\[x = \frac{-b \pm \sqrt{b^2 - 4ac}}{2a}\]

\subsection{Matrix Example}
A sample matrix:
\[
\begin{pmatrix}
a_{11} & a_{12} & a_{13} \\
a_{21} & a_{22} & a_{23} \\
a_{31} & a_{32} & a_{33}
\end{pmatrix}
\]

\section{Tables}
Below is a sample table showing data:

\begin{table}[h]
\centering
\begin{tabular}{@{}lrrr@{}}
\toprule
\textbf{Category} & \textbf{Value 1} & \textbf{Value 2} & \textbf{Total} \\
\midrule
Item A & 10.5 & 20.3 & 30.8 \\
Item B & 15.2 & 25.1 & 40.3 \\
Item C & 22.1 & 12.7 & 34.8 \\
\bottomrule
\end{tabular}
\caption{Sample Data Table}
\label{tab:sample}
\end{table}

\section{Lists}
\subsection{Bullet Points}
\begin{itemize}
  \item First item
  \item Second item
  \item Third item with sub-items:
    \begin{itemize}
      \item Sub-item 1
      \item Sub-item 2
    \end{itemize}
\end{itemize}

\subsection{Numbered List}
\begin{enumerate}
  \item First step
  \item Second step
  \item Third step
\end{enumerate}

\section{Citations and References}
When you use the citation search feature, references will be automatically managed. For example, you might cite a paper like this: \cite{example2024}. All citations will be automatically added to your \texttt{cite.bib} file.

\section{Conclusion}
This template demonstrates basic \LaTeX{} features including mathematical formulas, tables, and lists. The bibliography system is now set up to work seamlessly with Bibby AI's citation management. Feel free to modify it for your needs!

\bibliographystyle{plain}
\bibliography{cite}

\end{document}\documentclass{article}
\usepackage[utf8]{inputenc}
\usepackage{booktabs}
\usepackage{amsmath}
\usepackage{natbib}

\title{My Document}
\author{Your Name}
\date{\today}

\begin{document}

\maketitle

\section{Introduction}
This is a sample document demonstrating various \LaTeX{} features including mathematical equations, tables, and formatting. You can cite papers using the citation search feature, and they will be automatically added to your \texttt{cite.bib} file.

\section{Mathematical Expressions}
Here are some example equations:

\subsection{Quadratic Formula}
The quadratic formula is given by:
\[x = \frac{-b \pm \sqrt{b^2 - 4ac}}{2a}\]

\subsection{Matrix Example}
A sample matrix:
\[
\begin{pmatrix}
a_{11} & a_{12} & a_{13} \\
a_{21} & a_{22} & a_{23} \\
a_{31} & a_{32} & a_{33}
\end{pmatrix}
\]

\section{Tables}
Below is a sample table showing data:

\begin{table}[h]
\centering
\begin{tabular}{@{}lrrr@{}}
\toprule
\textbf{Category} & \textbf{Value 1} & \textbf{Value 2} & \textbf{Total} \\
\midrule
Item A & 10.5 & 20.3 & 30.8 \\
Item B & 15.2 & 25.1 & 40.3 \\
Item C & 22.1 & 12.7 & 34.8 \\
\bottomrule
\end{tabular}
\caption{Sample Data Table}
\label{tab:sample}
\end{table}

\section{Lists}
\subsection{Bullet Points}
\begin{itemize}
  \item First item
  \item Second item
  \item Third item with sub-items:
    \begin{itemize}
      \item Sub-item 1
      \item Sub-item 2
    \end{itemize}
\end{itemize}

\subsection{Numbered List}
\begin{enumerate}
  \item First step
  \item Second step
  \item Third step
\end{enumerate}

\section{Citations and References}
When you use the citation search feature, references will be automatically managed. For example, you might cite a paper like this: \cite{example2024}. All citations will be automatically added to your \texttt{cite.bib} file.

\section{Conclusion}
This template demonstrates basic \LaTeX{} features including mathematical formulas, tables, and lists. The bibliography system is now set up to work seamlessly with Bibby AI's citation management. Feel free to modify it for your needs!

\bibliographystyle{plain}
\bibliography{cite}

\end{document}\documentclass{article}
\usepackage[utf8]{inputenc}
\usepackage{booktabs}
\usepackage{amsmath}
\usepackage{natbib}

\title{My Document}
\author{Your Name}
\date{\today}

\begin{document}

\maketitle

\section{Introduction}
This is a sample document demonstrating various \LaTeX{} features including mathematical equations, tables, and formatting. You can cite papers using the citation search feature, and they will be automatically added to your \texttt{cite.bib} file.

\section{Mathematical Expressions}
Here are some example equations:

\subsection{Quadratic Formula}
The quadratic formula is given by:
\[x = \frac{-b \pm \sqrt{b^2 - 4ac}}{2a}\]

\subsection{Matrix Example}
A sample matrix:
\[
\begin{pmatrix}
a_{11} & a_{12} & a_{13} \\
a_{21} & a_{22} & a_{23} \\
a_{31} & a_{32} & a_{33}
\end{pmatrix}
\]

\section{Tables}
Below is a sample table showing data:

\begin{table}[h]
\centering
\begin{tabular}{@{}lrrr@{}}
\toprule
\textbf{Category} & \textbf{Value 1} & \textbf{Value 2} & \textbf{Total} \\
\midrule
Item A & 10.5 & 20.3 & 30.8 \\
Item B & 15.2 & 25.1 & 40.3 \\
Item C & 22.1 & 12.7 & 34.8 \\
\bottomrule
\end{tabular}
\caption{Sample Data Table}
\label{tab:sample}
\end{table}

\section{Lists}
\subsection{Bullet Points}
\begin{itemize}
  \item First item
  \item Second item
  \item Third item with sub-items:
    \begin{itemize}
      \item Sub-item 1
      \item Sub-item 2
    \end{itemize}
\end{itemize}

\subsection{Numbered List}
\begin{enumerate}
  \item First step
  \item Second step
  \item Third step
\end{enumerate}

\section{Citations and References}
When you use the citation search feature, references will be automatically managed. For example, you might cite a paper like this: \cite{example2024}. All citations will be automatically added to your \texttt{cite.bib} file.

\section{Conclusion}
This template demonstrates basic \LaTeX{} features including mathematical formulas, tables, and lists. The bibliography system is now set up to work seamlessly with Bibby AI's citation management. Feel free to modify it for your needs!

\bibliographystyle{plain}
\bibliography{cite}

\end{document}\documentclass{article}
\usepackage[utf8]{inputenc}
\usepackage{booktabs}
\usepackage{amsmath}
\usepackage{natbib}

\title{My Document}
\author{Your Name}
\date{\today}

\begin{document}

\maketitle

\section{Introduction}
This is a sample document demonstrating various \LaTeX{} features including mathematical equations, tables, and formatting. You can cite papers using the citation search feature, and they will be automatically added to your \texttt{cite.bib} file.

\section{Mathematical Expressions}
Here are some example equations:

\subsection{Quadratic Formula}
The quadratic formula is given by:
\[x = \frac{-b \pm \sqrt{b^2 - 4ac}}{2a}\]

\subsection{Matrix Example}
A sample matrix:
\[
\begin{pmatrix}
a_{11} & a_{12} & a_{13} \\
a_{21} & a_{22} & a_{23} \\
a_{31} & a_{32} & a_{33}
\end{pmatrix}
\]

\section{Tables}
Below is a sample table showing data:

\begin{table}[h]
\centering
\begin{tabular}{@{}lrrr@{}}
\toprule
\textbf{Category} & \textbf{Value 1} & \textbf{Value 2} & \textbf{Total} \\
\midrule
Item A & 10.5 & 20.3 & 30.8 \\
Item B & 15.2 & 25.1 & 40.3 \\
Item C & 22.1 & 12.7 & 34.8 \\
\bottomrule
\end{tabular}
\caption{Sample Data Table}
\label{tab:sample}
\end{table}

\section{Lists}
\subsection{Bullet Points}
\begin{itemize}
  \item First item
  \item Second item
  \item Third item with sub-items:
    \begin{itemize}
      \item Sub-item 1
      \item Sub-item 2
    \end{itemize}
\end{itemize}

\subsection{Numbered List}
\begin{enumerate}
  \item First step
  \item Second step
  \item Third step
\end{enumerate}

\section{Citations and References}
When you use the citation search feature, references will be automatically managed. For example, you might cite a paper like this: \cite{example2024}. All citations will be automatically added to your \texttt{cite.bib} file.

\section{Conclusion}
This template demonstrates basic \LaTeX{} features including mathematical formulas, tables, and lists. The bibliography system is now set up to work seamlessly with Bibby AI's citation management. Feel free to modify it for your needs!

\bibliographystyle{plain}
\bibliography{cite}

\end{document}\documentclass{article}
\usepackage[utf8]{inputenc}
\usepackage{booktabs}
\usepackage{amsmath}
\usepackage{natbib}

\title{My Document}
\author{Your Name}
\date{\today}

\begin{document}

\maketitle

\section{Introduction}
This is a sample document demonstrating various \LaTeX{} features including mathematical equations, tables, and formatting. You can cite papers using the citation search feature, and they will be automatically added to your \texttt{cite.bib} file.

\section{Mathematical Expressions}
Here are some example equations:

\subsection{Quadratic Formula}
The quadratic formula is given by:
\[x = \frac{-b \pm \sqrt{b^2 - 4ac}}{2a}\]

\subsection{Matrix Example}
A sample matrix:
\[
\begin{pmatrix}
a_{11} & a_{12} & a_{13} \\
a_{21} & a_{22} & a_{23} \\
a_{31} & a_{32} & a_{33}
\end{pmatrix}
\]

\section{Tables}
Below is a sample table showing data:

\begin{table}[h]
\centering
\begin{tabular}{@{}lrrr@{}}
\toprule
\textbf{Category} & \textbf{Value 1} & \textbf{Value 2} & \textbf{Total} \\
\midrule
Item A & 10.5 & 20.3 & 30.8 \\
Item B & 15.2 & 25.1 & 40.3 \\
Item C & 22.1 & 12.7 & 34.8 \\
\bottomrule
\end{tabular}
\caption{Sample Data Table}
\label{tab:sample}
\end{table}

\section{Lists}
\subsection{Bullet Points}
\begin{itemize}
  \item First item
  \item Second item
  \item Third item with sub-items:
    \begin{itemize}
      \item Sub-item 1
      \item Sub-item 2
    \end{itemize}
\end{itemize}

\subsection{Numbered List}
\begin{enumerate}
  \item First step
  \item Second step
  \item Third step
\end{enumerate}

\section{Citations and References}
When you use the citation search feature, references will be automatically managed. For example, you might cite a paper like this: \cite{example2024}. All citations will be automatically added to your \texttt{cite.bib} file.

\section{Conclusion}
This template demonstrates basic \LaTeX{} features including mathematical formulas, tables, and lists. The bibliography system is now set up to work seamlessly with Bibby AI's citation management. Feel free to modify it for your needs!

\bibliographystyle{plain}
\bibliography{cite}

\end{document}\documentclass{article}
\usepackage[utf8]{inputenc}
\usepackage{booktabs}
\usepackage{amsmath}
\usepackage{natbib}

\title{My Document}
\author{Your Name}
\date{\today}

\begin{document}

\maketitle

\section{Introduction}
This is a sample document demonstrating various \LaTeX{} features including mathematical equations, tables, and formatting. You can cite papers using the citation search feature, and they will be automatically added to your \texttt{cite.bib} file.

\section{Mathematical Expressions}
Here are some example equations:

\subsection{Quadratic Formula}
The quadratic formula is given by:
\[x = \frac{-b \pm \sqrt{b^2 - 4ac}}{2a}\]

\subsection{Matrix Example}
A sample matrix:
\[
\begin{pmatrix}
a_{11} & a_{12} & a_{13} \\
a_{21} & a_{22} & a_{23} \\
a_{31} & a_{32} & a_{33}
\end{pmatrix}
\]

\section{Tables}
Below is a sample table showing data:

\begin{table}[h]
\centering
\begin{tabular}{@{}lrrr@{}}
\toprule
\textbf{Category} & \textbf{Value 1} & \textbf{Value 2} & \textbf{Total} \\
\midrule
Item A & 10.5 & 20.3 & 30.8 \\
Item B & 15.2 & 25.1 & 40.3 \\
Item C & 22.1 & 12.7 & 34.8 \\
\bottomrule
\end{tabular}
\caption{Sample Data Table}
\label{tab:sample}
\end{table}

\section{Lists}
\subsection{Bullet Points}
\begin{itemize}
  \item First item
  \item Second item
  \item Third item with sub-items:
    \begin{itemize}
      \item Sub-item 1
      \item Sub-item 2
    \end{itemize}
\end{itemize}

\subsection{Numbered List}
\begin{enumerate}
  \item First step
  \item Second step
  \item Third step
\end{enumerate}

\section{Citations and References}
When you use the citation search feature, references will be automatically managed. For example, you might cite a paper like this: \cite{example2024}. All citations will be automatically added to your \texttt{cite.bib} file.

\section{Conclusion}
This template demonstrates basic \LaTeX{} features including mathematical formulas, tables, and lists. The bibliography system is now set up to work seamlessly with Bibby AI's citation management. Feel free to modify it for your needs!

\bibliographystyle{plain}
\bibliography{cite}

\end{document}\documentclass{article}
\usepackage[utf8]{inputenc}
\usepackage{booktabs}
\usepackage{amsmath}
\usepackage{natbib}

\title{My Document}
\author{Your Name}
\date{\today}

\begin{document}

\maketitle

\section{Introduction}
This is a sample document demonstrating various \LaTeX{} features including mathematical equations, tables, and formatting. You can cite papers using the citation search feature, and they will be automatically added to your \texttt{cite.bib} file.

\section{Mathematical Expressions}
Here are some example equations:

\subsection{Quadratic Formula}
The quadratic formula is given by:
\[x = \frac{-b \pm \sqrt{b^2 - 4ac}}{2a}\]

\subsection{Matrix Example}
A sample matrix:
\[
\begin{pmatrix}
a_{11} & a_{12} & a_{13} \\
a_{21} & a_{22} & a_{23} \\
a_{31} & a_{32} & a_{33}
\end{pmatrix}
\]

\section{Tables}
Below is a sample table showing data:

\begin{table}[h]
\centering
\begin{tabular}{@{}lrrr@{}}
\toprule
\textbf{Category} & \textbf{Value 1} & \textbf{Value 2} & \textbf{Total} \\
\midrule
Item A & 10.5 & 20.3 & 30.8 \\
Item B & 15.2 & 25.1 & 40.3 \\
Item C & 22.1 & 12.7 & 34.8 \\
\bottomrule
\end{tabular}
\caption{Sample Data Table}
\label{tab:sample}
\end{table}

\section{Lists}
\subsection{Bullet Points}
\begin{itemize}
  \item First item
  \item Second item
  \item Third item with sub-items:
    \begin{itemize}
      \item Sub-item 1
      \item Sub-item 2
    \end{itemize}
\end{itemize}

\subsection{Numbered List}
\begin{enumerate}
  \item First step
  \item Second step
  \item Third step
\end{enumerate}

\section{Citations and References}
When you use the citation search feature, references will be automatically managed. For example, you might cite a paper like this: \cite{example2024}. All citations will be automatically added to your \texttt{cite.bib} file.

\section{Conclusion}
This template demonstrates basic \LaTeX{} features including mathematical formulas, tables, and lists. The bibliography system is now set up to work seamlessly with Bibby AI's citation management. Feel free to modify it for your needs!

\bibliographystyle{plain}
\bibliography{cite}

\end{document}\documentclass{article}
\usepackage[utf8]{inputenc}
\usepackage{booktabs}
\usepackage{amsmath}
\usepackage{natbib}

\title{My Document}
\author{Your Name}
\date{\today}

\begin{document}

\maketitle

\section{Introduction}
This is a sample document demonstrating various \LaTeX{} features including mathematical equations, tables, and formatting. You can cite papers using the citation search feature, and they will be automatically added to your \texttt{cite.bib} file.

\section{Mathematical Expressions}
Here are some example equations:

\subsection{Quadratic Formula}
The quadratic formula is given by:
\[x = \frac{-b \pm \sqrt{b^2 - 4ac}}{2a}\]

\subsection{Matrix Example}
A sample matrix:
\[
\begin{pmatrix}
a_{11} & a_{12} & a_{13} \\
a_{21} & a_{22} & a_{23} \\
a_{31} & a_{32} & a_{33}
\end{pmatrix}
\]

\section{Tables}
Below is a sample table showing data:

\begin{table}[h]
\centering
\begin{tabular}{@{}lrrr@{}}
\toprule
\textbf{Category} & \textbf{Value 1} & \textbf{Value 2} & \textbf{Total} \\
\midrule
Item A & 10.5 & 20.3 & 30.8 \\
Item B & 15.2 & 25.1 & 40.3 \\
Item C & 22.1 & 12.7 & 34.8 \\
\bottomrule
\end{tabular}
\caption{Sample Data Table}
\label{tab:sample}
\end{table}

\section{Lists}
\subsection{Bullet Points}
\begin{itemize}
  \item First item
  \item Second item
  \item Third item with sub-items:
    \begin{itemize}
      \item Sub-item 1
      \item Sub-item 2
    \end{itemize}
\end{itemize}

\subsection{Numbered List}
\begin{enumerate}
  \item First step
  \item Second step
  \item Third step
\end{enumerate}

\section{Citations and References}
When you use the citation search feature, references will be automatically managed. For example, you might cite a paper like this: \cite{example2024}. All citations will be automatically added to your \texttt{cite.bib} file.

\section{Conclusion}
This template demonstrates basic \LaTeX{} features including mathematical formulas, tables, and lists. The bibliography system is now set up to work seamlessly with Bibby AI's citation management. Feel free to modify it for your needs!

\bibliographystyle{plain}
\bibliography{cite}

\end{document}\documentclass{article}
\usepackage[utf8]{inputenc}
\usepackage{booktabs}
\usepackage{amsmath}
\usepackage{natbib}

\title{My Document}
\author{Your Name}
\date{\today}

\begin{document}

\maketitle

\section{Introduction}
This is a sample document demonstrating various \LaTeX{} features including mathematical equations, tables, and formatting. You can cite papers using the citation search feature, and they will be automatically added to your \texttt{cite.bib} file.

\section{Mathematical Expressions}
Here are some example equations:

\subsection{Quadratic Formula}
The quadratic formula is given by:
\[x = \frac{-b \pm \sqrt{b^2 - 4ac}}{2a}\]

\subsection{Matrix Example}
A sample matrix:
\[
\begin{pmatrix}
a_{11} & a_{12} & a_{13} \\
a_{21} & a_{22} & a_{23} \\
a_{31} & a_{32} & a_{33}
\end{pmatrix}
\]

\section{Tables}
Below is a sample table showing data:

\begin{table}[h]
\centering
\begin{tabular}{@{}lrrr@{}}
\toprule
\textbf{Category} & \textbf{Value 1} & \textbf{Value 2} & \textbf{Total} \\
\midrule
Item A & 10.5 & 20.3 & 30.8 \\
Item B & 15.2 & 25.1 & 40.3 \\
Item C & 22.1 & 12.7 & 34.8 \\
\bottomrule
\end{tabular}
\caption{Sample Data Table}
\label{tab:sample}
\end{table}

\section{Lists}
\subsection{Bullet Points}
\begin{itemize}
  \item First item
  \item Second item
  \item Third item with sub-items:
    \begin{itemize}
      \item Sub-item 1
      \item Sub-item 2
    \end{itemize}
\end{itemize}

\subsection{Numbered List}
\begin{enumerate}
  \item First step
  \item Second step
  \item Third step
\end{enumerate}

\section{Citations and References}
When you use the citation search feature, references will be automatically managed. For example, you might cite a paper like this: \cite{example2024}. All citations will be automatically added to your \texttt{cite.bib} file.

\section{Conclusion}
This template demonstrates basic \LaTeX{} features including mathematical formulas, tables, and lists. The bibliography system is now set up to work seamlessly with Bibby AI's citation management. Feel free to modify it for your needs!

\bibliographystyle{plain}
\bibliography{cite}

\end{document}\documentclass{article}
\usepackage[utf8]{inputenc}
\usepackage{booktabs}
\usepackage{amsmath}
\usepackage{natbib}

\title{My Document}
\author{Your Name}
\date{\today}

\begin{document}

\maketitle

\section{Introduction}
This is a sample document demonstrating various \LaTeX{} features including mathematical equations, tables, and formatting. You can cite papers using the citation search feature, and they will be automatically added to your \texttt{cite.bib} file.

\section{Mathematical Expressions}
Here are some example equations:

\subsection{Quadratic Formula}
The quadratic formula is given by:
\[x = \frac{-b \pm \sqrt{b^2 - 4ac}}{2a}\]

\subsection{Matrix Example}
A sample matrix:
\[
\begin{pmatrix}
a_{11} & a_{12} & a_{13} \\
a_{21} & a_{22} & a_{23} \\
a_{31} & a_{32} & a_{33}
\end{pmatrix}
\]

\section{Tables}
Below is a sample table showing data:

\begin{table}[h]
\centering
\begin{tabular}{@{}lrrr@{}}
\toprule
\textbf{Category} & \textbf{Value 1} & \textbf{Value 2} & \textbf{Total} \\
\midrule
Item A & 10.5 & 20.3 & 30.8 \\
Item B & 15.2 & 25.1 & 40.3 \\
Item C & 22.1 & 12.7 & 34.8 \\
\bottomrule
\end{tabular}
\caption{Sample Data Table}
\label{tab:sample}
\end{table}

\section{Lists}
\subsection{Bullet Points}
\begin{itemize}
  \item First item
  \item Second item
  \item Third item with sub-items:
    \begin{itemize}
      \item Sub-item 1
      \item Sub-item 2
    \end{itemize}
\end{itemize}

\subsection{Numbered List}
\begin{enumerate}
  \item First step
  \item Second step
  \item Third step
\end{enumerate}

\section{Citations and References}
When you use the citation search feature, references will be automatically managed. For example, you might cite a paper like this: \cite{example2024}. All citations will be automatically added to your \texttt{cite.bib} file.

\section{Conclusion}
This template demonstrates basic \LaTeX{} features including mathematical formulas, tables, and lists. The bibliography system is now set up to work seamlessly with Bibby AI's citation management. Feel free to modify it for your needs!

\bibliographystyle{plain}
\bibliography{cite}

\end{document}\documentclass{article}
\usepackage[utf8]{inputenc}
\usepackage{booktabs}
\usepackage{amsmath}
\usepackage{natbib}

\title{My Document}
\author{Your Name}
\date{\today}

\begin{document}

\maketitle

\section{Introduction}
This is a sample document demonstrating various \LaTeX{} features including mathematical equations, tables, and formatting. You can cite papers using the citation search feature, and they will be automatically added to your \texttt{cite.bib} file.

\section{Mathematical Expressions}
Here are some example equations:

\subsection{Quadratic Formula}
The quadratic formula is given by:
\[x = \frac{-b \pm \sqrt{b^2 - 4ac}}{2a}\]

\subsection{Matrix Example}
A sample matrix:
\[
\begin{pmatrix}
a_{11} & a_{12} & a_{13} \\
a_{21} & a_{22} & a_{23} \\
a_{31} & a_{32} & a_{33}
\end{pmatrix}
\]

\section{Tables}
Below is a sample table showing data:

\begin{table}[h]
\centering
\begin{tabular}{@{}lrrr@{}}
\toprule
\textbf{Category} & \textbf{Value 1} & \textbf{Value 2} & \textbf{Total} \\
\midrule
Item A & 10.5 & 20.3 & 30.8 \\
Item B & 15.2 & 25.1 & 40.3 \\
Item C & 22.1 & 12.7 & 34.8 \\
\bottomrule
\end{tabular}
\caption{Sample Data Table}
\label{tab:sample}
\end{table}

\section{Lists}
\subsection{Bullet Points}
\begin{itemize}
  \item First item
  \item Second item
  \item Third item with sub-items:
    \begin{itemize}
      \item Sub-item 1
      \item Sub-item 2
    \end{itemize}
\end{itemize}

\subsection{Numbered List}
\begin{enumerate}
  \item First step
  \item Second step
  \item Third step
\end{enumerate}

\section{Citations and References}
When you use the citation search feature, references will be automatically managed. For example, you might cite a paper like this: \cite{example2024}. All citations will be automatically added to your \texttt{cite.bib} file.

\section{Conclusion}
This template demonstrates basic \LaTeX{} features including mathematical formulas, tables, and lists. The bibliography system is now set up to work seamlessly with Bibby AI's citation management. Feel free to modify it for your needs!

\bibliographystyle{plain}
\bibliography{cite}

\end{document}\documentclass{article}
\usepackage[utf8]{inputenc}
\usepackage{booktabs}
\usepackage{amsmath}
\usepackage{natbib}

\title{My Document}
\author{Your Name}
\date{\today}

\begin{document}

\maketitle

\section{Introduction}
This is a sample document demonstrating various \LaTeX{} features including mathematical equations, tables, and formatting. You can cite papers using the citation search feature, and they will be automatically added to your \texttt{cite.bib} file.

\section{Mathematical Expressions}
Here are some example equations:

\subsection{Quadratic Formula}
The quadratic formula is given by:
\[x = \frac{-b \pm \sqrt{b^2 - 4ac}}{2a}\]

\subsection{Matrix Example}
A sample matrix:
\[
\begin{pmatrix}
a_{11} & a_{12} & a_{13} \\
a_{21} & a_{22} & a_{23} \\
a_{31} & a_{32} & a_{33}
\end{pmatrix}
\]

\section{Tables}
Below is a sample table showing data:

\begin{table}[h]
\centering
\begin{tabular}{@{}lrrr@{}}
\toprule
\textbf{Category} & \textbf{Value 1} & \textbf{Value 2} & \textbf{Total} \\
\midrule
Item A & 10.5 & 20.3 & 30.8 \\
Item B & 15.2 & 25.1 & 40.3 \\
Item C & 22.1 & 12.7 & 34.8 \\
\bottomrule
\end{tabular}
\caption{Sample Data Table}
\label{tab:sample}
\end{table}

\section{Lists}
\subsection{Bullet Points}
\begin{itemize}
  \item First item
  \item Second item
  \item Third item with sub-items:
    \begin{itemize}
      \item Sub-item 1
      \item Sub-item 2
    \end{itemize}
\end{itemize}

\subsection{Numbered List}
\begin{enumerate}
  \item First step
  \item Second step
  \item Third step
\end{enumerate}

\section{Citations and References}
When you use the citation search feature, references will be automatically managed. For example, you might cite a paper like this: \cite{example2024}. All citations will be automatically added to your \texttt{cite.bib} file.

\section{Conclusion}
This template demonstrates basic \LaTeX{} features including mathematical formulas, tables, and lists. The bibliography system is now set up to work seamlessly with Bibby AI's citation management. Feel free to modify it for your needs!

\bibliographystyle{plain}
\bibliography{cite}

\end{document}\documentclass{article}
\usepackage[utf8]{inputenc}
\usepackage{booktabs}
\usepackage{amsmath}
\usepackage{natbib}

\title{My Document}
\author{Your Name}
\date{\today}

\begin{document}

\maketitle

\section{Introduction}
This is a sample document demonstrating various \LaTeX{} features including mathematical equations, tables, and formatting. You can cite papers using the citation search feature, and they will be automatically added to your \texttt{cite.bib} file.

\section{Mathematical Expressions}
Here are some example equations:

\subsection{Quadratic Formula}
The quadratic formula is given by:
\[x = \frac{-b \pm \sqrt{b^2 - 4ac}}{2a}\]

\subsection{Matrix Example}
A sample matrix:
\[
\begin{pmatrix}
a_{11} & a_{12} & a_{13} \\
a_{21} & a_{22} & a_{23} \\
a_{31} & a_{32} & a_{33}
\end{pmatrix}
\]

\section{Tables}
Below is a sample table showing data:

\begin{table}[h]
\centering
\begin{tabular}{@{}lrrr@{}}
\toprule
\textbf{Category} & \textbf{Value 1} & \textbf{Value 2} & \textbf{Total} \\
\midrule
Item A & 10.5 & 20.3 & 30.8 \\
Item B & 15.2 & 25.1 & 40.3 \\
Item C & 22.1 & 12.7 & 34.8 \\
\bottomrule
\end{tabular}
\caption{Sample Data Table}
\label{tab:sample}
\end{table}

\section{Lists}
\subsection{Bullet Points}
\begin{itemize}
  \item First item
  \item Second item
  \item Third item with sub-items:
    \begin{itemize}
      \item Sub-item 1
      \item Sub-item 2
    \end{itemize}
\end{itemize}

\subsection{Numbered List}
\begin{enumerate}
  \item First step
  \item Second step
  \item Third step
\end{enumerate}

\section{Citations and References}
When you use the citation search feature, references will be automatically managed. For example, you might cite a paper like this: \cite{example2024}. All citations will be automatically added to your \texttt{cite.bib} file.

\section{Conclusion}
This template demonstrates basic \LaTeX{} features including mathematical formulas, tables, and lists. The bibliography system is now set up to work seamlessly with Bibby AI's citation management. Feel free to modify it for your needs!

\bibliographystyle{plain}
\bibliography{cite}

\end{document}\documentclass{article}
\usepackage[utf8]{inputenc}
\usepackage{booktabs}
\usepackage{amsmath}
\usepackage{natbib}

\title{My Document}
\author{Your Name}
\date{\today}

\begin{document}

\maketitle

\section{Introduction}
This is a sample document demonstrating various \LaTeX{} features including mathematical equations, tables, and formatting. You can cite papers using the citation search feature, and they will be automatically added to your \texttt{cite.bib} file.

\section{Mathematical Expressions}
Here are some example equations:

\subsection{Quadratic Formula}
The quadratic formula is given by:
\[x = \frac{-b \pm \sqrt{b^2 - 4ac}}{2a}\]

\subsection{Matrix Example}
A sample matrix:
\[
\begin{pmatrix}
a_{11} & a_{12} & a_{13} \\
a_{21} & a_{22} & a_{23} \\
a_{31} & a_{32} & a_{33}
\end{pmatrix}
\]

\section{Tables}
Below is a sample table showing data:

\begin{table}[h]
\centering
\begin{tabular}{@{}lrrr@{}}
\toprule
\textbf{Category} & \textbf{Value 1} & \textbf{Value 2} & \textbf{Total} \\
\midrule
Item A & 10.5 & 20.3 & 30.8 \\
Item B & 15.2 & 25.1 & 40.3 \\
Item C & 22.1 & 12.7 & 34.8 \\
\bottomrule
\end{tabular}
\caption{Sample Data Table}
\label{tab:sample}
\end{table}

\section{Lists}
\subsection{Bullet Points}
\begin{itemize}
  \item First item
  \item Second item
  \item Third item with sub-items:
    \begin{itemize}
      \item Sub-item 1
      \item Sub-item 2
    \end{itemize}
\end{itemize}

\subsection{Numbered List}
\begin{enumerate}
  \item First step
  \item Second step
  \item Third step
\end{enumerate}

\section{Citations and References}
When you use the citation search feature, references will be automatically managed. For example, you might cite a paper like this: \cite{example2024}. All citations will be automatically added to your \texttt{cite.bib} file.

\section{Conclusion}
This template demonstrates basic \LaTeX{} features including mathematical formulas, tables, and lists. The bibliography system is now set up to work seamlessly with Bibby AI's citation management. Feel free to modify it for your needs!

\bibliographystyle{plain}
\bibliography{cite}

\end{document}\documentclass{article}
\usepackage[utf8]{inputenc}
\usepackage{booktabs}
\usepackage{amsmath}
\usepackage{natbib}

\title{My Document}
\author{Your Name}
\date{\today}

\begin{document}

\maketitle

\section{Introduction}
This is a sample document demonstrating various \LaTeX{} features including mathematical equations, tables, and formatting. You can cite papers using the citation search feature, and they will be automatically added to your \texttt{cite.bib} file.

\section{Mathematical Expressions}
Here are some example equations:

\subsection{Quadratic Formula}
The quadratic formula is given by:
\[x = \frac{-b \pm \sqrt{b^2 - 4ac}}{2a}\]

\subsection{Matrix Example}
A sample matrix:
\[
\begin{pmatrix}
a_{11} & a_{12} & a_{13} \\
a_{21} & a_{22} & a_{23} \\
a_{31} & a_{32} & a_{33}
\end{pmatrix}
\]

\section{Tables}
Below is a sample table showing data:

\begin{table}[h]
\centering
\begin{tabular}{@{}lrrr@{}}
\toprule
\textbf{Category} & \textbf{Value 1} & \textbf{Value 2} & \textbf{Total} \\
\midrule
Item A & 10.5 & 20.3 & 30.8 \\
Item B & 15.2 & 25.1 & 40.3 \\
Item C & 22.1 & 12.7 & 34.8 \\
\bottomrule
\end{tabular}
\caption{Sample Data Table}
\label{tab:sample}
\end{table}

\section{Lists}
\subsection{Bullet Points}
\begin{itemize}
  \item First item
  \item Second item
  \item Third item with sub-items:
    \begin{itemize}
      \item Sub-item 1
      \item Sub-item 2
    \end{itemize}
\end{itemize}

\subsection{Numbered List}
\begin{enumerate}
  \item First step
  \item Second step
  \item Third step
\end{enumerate}

\section{Citations and References}
When you use the citation search feature, references will be automatically managed. For example, you might cite a paper like this: \cite{example2024}. All citations will be automatically added to your \texttt{cite.bib} file.

\section{Conclusion}
This template demonstrates basic \LaTeX{} features including mathematical formulas, tables, and lists. The bibliography system is now set up to work seamlessly with Bibby AI's citation management. Feel free to modify it for your needs!

\bibliographystyle{plain}
\bibliography{cite}

\end{document}\documentclass{article}
\usepackage[utf8]{inputenc}
\usepackage{booktabs}
\usepackage{amsmath}
\usepackage{natbib}

\title{My Document}
\author{Your Name}
\date{\today}

\begin{document}

\maketitle

\section{Introduction}
This is a sample document demonstrating various \LaTeX{} features including mathematical equations, tables, and formatting. You can cite papers using the citation search feature, and they will be automatically added to your \texttt{cite.bib} file.

\section{Mathematical Expressions}
Here are some example equations:

\subsection{Quadratic Formula}
The quadratic formula is given by:
\[x = \frac{-b \pm \sqrt{b^2 - 4ac}}{2a}\]

\subsection{Matrix Example}
A sample matrix:
\[
\begin{pmatrix}
a_{11} & a_{12} & a_{13} \\
a_{21} & a_{22} & a_{23} \\
a_{31} & a_{32} & a_{33}
\end{pmatrix}
\]

\section{Tables}
Below is a sample table showing data:

\begin{table}[h]
\centering
\begin{tabular}{@{}lrrr@{}}
\toprule
\textbf{Category} & \textbf{Value 1} & \textbf{Value 2} & \textbf{Total} \\
\midrule
Item A & 10.5 & 20.3 & 30.8 \\
Item B & 15.2 & 25.1 & 40.3 \\
Item C & 22.1 & 12.7 & 34.8 \\
\bottomrule
\end{tabular}
\caption{Sample Data Table}
\label{tab:sample}
\end{table}

\section{Lists}
\subsection{Bullet Points}
\begin{itemize}
  \item First item
  \item Second item
  \item Third item with sub-items:
    \begin{itemize}
      \item Sub-item 1
      \item Sub-item 2
    \end{itemize}
\end{itemize}

\subsection{Numbered List}
\begin{enumerate}
  \item First step
  \item Second step
  \item Third step
\end{enumerate}

\section{Citations and References}
When you use the citation search feature, references will be automatically managed. For example, you might cite a paper like this: \cite{example2024}. All citations will be automatically added to your \texttt{cite.bib} file.

\section{Conclusion}
This template demonstrates basic \LaTeX{} features including mathematical formulas, tables, and lists. The bibliography system is now set up to work seamlessly with Bibby AI's citation management. Feel free to modify it for your needs!

\bibliographystyle{plain}
\bibliography{cite}

\end{document}\documentclass{article}
\usepackage[utf8]{inputenc}
\usepackage{booktabs}
\usepackage{amsmath}
\usepackage{natbib}

\title{My Document}
\author{Your Name}
\date{\today}

\begin{document}

\maketitle

\section{Introduction}
This is a sample document demonstrating various \LaTeX{} features including mathematical equations, tables, and formatting. You can cite papers using the citation search feature, and they will be automatically added to your \texttt{cite.bib} file.

\section{Mathematical Expressions}
Here are some example equations:

\subsection{Quadratic Formula}
The quadratic formula is given by:
\[x = \frac{-b \pm \sqrt{b^2 - 4ac}}{2a}\]

\subsection{Matrix Example}
A sample matrix:
\[
\begin{pmatrix}
a_{11} & a_{12} & a_{13} \\
a_{21} & a_{22} & a_{23} \\
a_{31} & a_{32} & a_{33}
\end{pmatrix}
\]

\section{Tables}
Below is a sample table showing data:

\begin{table}[h]
\centering
\begin{tabular}{@{}lrrr@{}}
\toprule
\textbf{Category} & \textbf{Value 1} & \textbf{Value 2} & \textbf{Total} \\
\midrule
Item A & 10.5 & 20.3 & 30.8 \\
Item B & 15.2 & 25.1 & 40.3 \\
Item C & 22.1 & 12.7 & 34.8 \\
\bottomrule
\end{tabular}
\caption{Sample Data Table}
\label{tab:sample}
\end{table}

\section{Lists}
\subsection{Bullet Points}
\begin{itemize}
  \item First item
  \item Second item
  \item Third item with sub-items:
    \begin{itemize}
      \item Sub-item 1
      \item Sub-item 2
    \end{itemize}
\end{itemize}

\subsection{Numbered List}
\begin{enumerate}
  \item First step
  \item Second step
  \item Third step
\end{enumerate}

\section{Citations and References}
When you use the citation search feature, references will be automatically managed. For example, you might cite a paper like this: \cite{example2024}. All citations will be automatically added to your \texttt{cite.bib} file.

\section{Conclusion}
This template demonstrates basic \LaTeX{} features including mathematical formulas, tables, and lists. The bibliography system is now set up to work seamlessly with Bibby AI's citation management. Feel free to modify it for your needs!

\bibliographystyle{plain}
\bibliography{cite}

\end{document}\documentclass{article}
\usepackage[utf8]{inputenc}
\usepackage{booktabs}
\usepackage{amsmath}
\usepackage{natbib}

\title{My Document}
\author{Your Name}
\date{\today}

\begin{document}

\maketitle

\section{Introduction}
This is a sample document demonstrating various \LaTeX{} features including mathematical equations, tables, and formatting. You can cite papers using the citation search feature, and they will be automatically added to your \texttt{cite.bib} file.

\section{Mathematical Expressions}
Here are some example equations:

\subsection{Quadratic Formula}
The quadratic formula is given by:
\[x = \frac{-b \pm \sqrt{b^2 - 4ac}}{2a}\]

\subsection{Matrix Example}
A sample matrix:
\[
\begin{pmatrix}
a_{11} & a_{12} & a_{13} \\
a_{21} & a_{22} & a_{23} \\
a_{31} & a_{32} & a_{33}
\end{pmatrix}
\]

\section{Tables}
Below is a sample table showing data:

\begin{table}[h]
\centering
\begin{tabular}{@{}lrrr@{}}
\toprule
\textbf{Category} & \textbf{Value 1} & \textbf{Value 2} & \textbf{Total} \\
\midrule
Item A & 10.5 & 20.3 & 30.8 \\
Item B & 15.2 & 25.1 & 40.3 \\
Item C & 22.1 & 12.7 & 34.8 \\
\bottomrule
\end{tabular}
\caption{Sample Data Table}
\label{tab:sample}
\end{table}

\section{Lists}
\subsection{Bullet Points}
\begin{itemize}
  \item First item
  \item Second item
  \item Third item with sub-items:
    \begin{itemize}
      \item Sub-item 1
      \item Sub-item 2
    \end{itemize}
\end{itemize}

\subsection{Numbered List}
\begin{enumerate}
  \item First step
  \item Second step
  \item Third step
\end{enumerate}

\section{Citations and References}
When you use the citation search feature, references will be automatically managed. For example, you might cite a paper like this: \cite{example2024}. All citations will be automatically added to your \texttt{cite.bib} file.

\section{Conclusion}
This template demonstrates basic \LaTeX{} features including mathematical formulas, tables, and lists. The bibliography system is now set up to work seamlessly with Bibby AI's citation management. Feel free to modify it for your needs!

\bibliographystyle{plain}
\bibliography{cite}

\end{document}\documentclass{article}
\usepackage[utf8]{inputenc}
\usepackage{booktabs}
\usepackage{amsmath}
\usepackage{natbib}

\title{My Document}
\author{Your Name}
\date{\today}

\begin{document}

\maketitle

\section{Introduction}
This is a sample document demonstrating various \LaTeX{} features including mathematical equations, tables, and formatting. You can cite papers using the citation search feature, and they will be automatically added to your \texttt{cite.bib} file.

\section{Mathematical Expressions}
Here are some example equations:

\subsection{Quadratic Formula}
The quadratic formula is given by:
\[x = \frac{-b \pm \sqrt{b^2 - 4ac}}{2a}\]

\subsection{Matrix Example}
A sample matrix:
\[
\begin{pmatrix}
a_{11} & a_{12} & a_{13} \\
a_{21} & a_{22} & a_{23} \\
a_{31} & a_{32} & a_{33}
\end{pmatrix}
\]

\section{Tables}
Below is a sample table showing data:

\begin{table}[h]
\centering
\begin{tabular}{@{}lrrr@{}}
\toprule
\textbf{Category} & \textbf{Value 1} & \textbf{Value 2} & \textbf{Total} \\
\midrule
Item A & 10.5 & 20.3 & 30.8 \\
Item B & 15.2 & 25.1 & 40.3 \\
Item C & 22.1 & 12.7 & 34.8 \\
\bottomrule
\end{tabular}
\caption{Sample Data Table}
\label{tab:sample}
\end{table}

\section{Lists}
\subsection{Bullet Points}
\begin{itemize}
  \item First item
  \item Second item
  \item Third item with sub-items:
    \begin{itemize}
      \item Sub-item 1
      \item Sub-item 2
    \end{itemize}
\end{itemize}

\subsection{Numbered List}
\begin{enumerate}
  \item First step
  \item Second step
  \item Third step
\end{enumerate}

\section{Citations and References}
When you use the citation search feature, references will be automatically managed. For example, you might cite a paper like this: \cite{example2024}. All citations will be automatically added to your \texttt{cite.bib} file.

\section{Conclusion}
This template demonstrates basic \LaTeX{} features including mathematical formulas, tables, and lists. The bibliography system is now set up to work seamlessly with Bibby AI's citation management. Feel free to modify it for your needs!

\bibliographystyle{plain}
\bibliography{cite}

\end{document}\documentclass{article}
\usepackage[utf8]{inputenc}
\usepackage{booktabs}
\usepackage{amsmath}
\usepackage{natbib}

\title{My Document}
\author{Your Name}
\date{\today}

\begin{document}

\maketitle

\section{Introduction}
This is a sample document demonstrating various \LaTeX{} features including mathematical equations, tables, and formatting. You can cite papers using the citation search feature, and they will be automatically added to your \texttt{cite.bib} file.

\section{Mathematical Expressions}
Here are some example equations:

\subsection{Quadratic Formula}
The quadratic formula is given by:
\[x = \frac{-b \pm \sqrt{b^2 - 4ac}}{2a}\]

\subsection{Matrix Example}
A sample matrix:
\[
\begin{pmatrix}
a_{11} & a_{12} & a_{13} \\
a_{21} & a_{22} & a_{23} \\
a_{31} & a_{32} & a_{33}
\end{pmatrix}
\]

\section{Tables}
Below is a sample table showing data:

\begin{table}[h]
\centering
\begin{tabular}{@{}lrrr@{}}
\toprule
\textbf{Category} & \textbf{Value 1} & \textbf{Value 2} & \textbf{Total} \\
\midrule
Item A & 10.5 & 20.3 & 30.8 \\
Item B & 15.2 & 25.1 & 40.3 \\
Item C & 22.1 & 12.7 & 34.8 \\
\bottomrule
\end{tabular}
\caption{Sample Data Table}
\label{tab:sample}
\end{table}

\section{Lists}
\subsection{Bullet Points}
\begin{itemize}
  \item First item
  \item Second item
  \item Third item with sub-items:
    \begin{itemize}
      \item Sub-item 1
      \item Sub-item 2
    \end{itemize}
\end{itemize}

\subsection{Numbered List}
\begin{enumerate}
  \item First step
  \item Second step
  \item Third step
\end{enumerate}

\section{Citations and References}
When you use the citation search feature, references will be automatically managed. For example, you might cite a paper like this: \cite{example2024}. All citations will be automatically added to your \texttt{cite.bib} file.

\section{Conclusion}
This template demonstrates basic \LaTeX{} features including mathematical formulas, tables, and lists. The bibliography system is now set up to work seamlessly with Bibby AI's citation management. Feel free to modify it for your needs!

\bibliographystyle{plain}
\bibliography{cite}

\end{document}\documentclass{article}
\usepackage[utf8]{inputenc}
\usepackage{booktabs}
\usepackage{amsmath}
\usepackage{natbib}

\title{My Document}
\author{Your Name}
\date{\today}

\begin{document}

\maketitle

\section{Introduction}
This is a sample document demonstrating various \LaTeX{} features including mathematical equations, tables, and formatting. You can cite papers using the citation search feature, and they will be automatically added to your \texttt{cite.bib} file.

\section{Mathematical Expressions}
Here are some example equations:

\subsection{Quadratic Formula}
The quadratic formula is given by:
\[x = \frac{-b \pm \sqrt{b^2 - 4ac}}{2a}\]

\subsection{Matrix Example}
A sample matrix:
\[
\begin{pmatrix}
a_{11} & a_{12} & a_{13} \\
a_{21} & a_{22} & a_{23} \\
a_{31} & a_{32} & a_{33}
\end{pmatrix}
\]

\section{Tables}
Below is a sample table showing data:

\begin{table}[h]
\centering
\begin{tabular}{@{}lrrr@{}}
\toprule
\textbf{Category} & \textbf{Value 1} & \textbf{Value 2} & \textbf{Total} \\
\midrule
Item A & 10.5 & 20.3 & 30.8 \\
Item B & 15.2 & 25.1 & 40.3 \\
Item C & 22.1 & 12.7 & 34.8 \\
\bottomrule
\end{tabular}
\caption{Sample Data Table}
\label{tab:sample}
\end{table}

\section{Lists}
\subsection{Bullet Points}
\begin{itemize}
  \item First item
  \item Second item
  \item Third item with sub-items:
    \begin{itemize}
      \item Sub-item 1
      \item Sub-item 2
    \end{itemize}
\end{itemize}

\subsection{Numbered List}
\begin{enumerate}
  \item First step
  \item Second step
  \item Third step
\end{enumerate}

\section{Citations and References}
When you use the citation search feature, references will be automatically managed. For example, you might cite a paper like this: \cite{example2024}. All citations will be automatically added to your \texttt{cite.bib} file.

\section{Conclusion}
This template demonstrates basic \LaTeX{} features including mathematical formulas, tables, and lists. The bibliography system is now set up to work seamlessly with Bibby AI's citation management. Feel free to modify it for your needs!

\bibliographystyle{plain}
\bibliography{cite}

\end{document}\documentclass{article}
\usepackage[utf8]{inputenc}
\usepackage{booktabs}
\usepackage{amsmath}
\usepackage{natbib}

\title{My Document}
\author{Your Name}
\date{\today}

\begin{document}

\maketitle

\section{Introduction}
This is a sample document demonstrating various \LaTeX{} features including mathematical equations, tables, and formatting. You can cite papers using the citation search feature, and they will be automatically added to your \texttt{cite.bib} file.

\section{Mathematical Expressions}
Here are some example equations:

\subsection{Quadratic Formula}
The quadratic formula is given by:
\[x = \frac{-b \pm \sqrt{b^2 - 4ac}}{2a}\]

\subsection{Matrix Example}
A sample matrix:
\[
\begin{pmatrix}
a_{11} & a_{12} & a_{13} \\
a_{21} & a_{22} & a_{23} \\
a_{31} & a_{32} & a_{33}
\end{pmatrix}
\]

\section{Tables}
Below is a sample table showing data:

\begin{table}[h]
\centering
\begin{tabular}{@{}lrrr@{}}
\toprule
\textbf{Category} & \textbf{Value 1} & \textbf{Value 2} & \textbf{Total} \\
\midrule
Item A & 10.5 & 20.3 & 30.8 \\
Item B & 15.2 & 25.1 & 40.3 \\
Item C & 22.1 & 12.7 & 34.8 \\
\bottomrule
\end{tabular}
\caption{Sample Data Table}
\label{tab:sample}
\end{table}

\section{Lists}
\subsection{Bullet Points}
\begin{itemize}
  \item First item
  \item Second item
  \item Third item with sub-items:
    \begin{itemize}
      \item Sub-item 1
      \item Sub-item 2
    \end{itemize}
\end{itemize}

\subsection{Numbered List}
\begin{enumerate}
  \item First step
  \item Second step
  \item Third step
\end{enumerate}

\section{Citations and References}
When you use the citation search feature, references will be automatically managed. For example, you might cite a paper like this: \cite{example2024}. All citations will be automatically added to your \texttt{cite.bib} file.

\section{Conclusion}
This template demonstrates basic \LaTeX{} features including mathematical formulas, tables, and lists. The bibliography system is now set up to work seamlessly with Bibby AI's citation management. Feel free to modify it for your needs!

\bibliographystyle{plain}
\bibliography{cite}

\end{document}\documentclass{article}
\usepackage[utf8]{inputenc}
\usepackage{booktabs}
\usepackage{amsmath}
\usepackage{natbib}

\title{My Document}
\author{Your Name}
\date{\today}

\begin{document}

\maketitle

\section{Introduction}
This is a sample document demonstrating various \LaTeX{} features including mathematical equations, tables, and formatting. You can cite papers using the citation search feature, and they will be automatically added to your \texttt{cite.bib} file.

\section{Mathematical Expressions}
Here are some example equations:

\subsection{Quadratic Formula}
The quadratic formula is given by:
\[x = \frac{-b \pm \sqrt{b^2 - 4ac}}{2a}\]

\subsection{Matrix Example}
A sample matrix:
\[
\begin{pmatrix}
a_{11} & a_{12} & a_{13} \\
a_{21} & a_{22} & a_{23} \\
a_{31} & a_{32} & a_{33}
\end{pmatrix}
\]

\section{Tables}
Below is a sample table showing data:

\begin{table}[h]
\centering
\begin{tabular}{@{}lrrr@{}}
\toprule
\textbf{Category} & \textbf{Value 1} & \textbf{Value 2} & \textbf{Total} \\
\midrule
Item A & 10.5 & 20.3 & 30.8 \\
Item B & 15.2 & 25.1 & 40.3 \\
Item C & 22.1 & 12.7 & 34.8 \\
\bottomrule
\end{tabular}
\caption{Sample Data Table}
\label{tab:sample}
\end{table}

\section{Lists}
\subsection{Bullet Points}
\begin{itemize}
  \item First item
  \item Second item
  \item Third item with sub-items:
    \begin{itemize}
      \item Sub-item 1
      \item Sub-item 2
    \end{itemize}
\end{itemize}

\subsection{Numbered List}
\begin{enumerate}
  \item First step
  \item Second step
  \item Third step
\end{enumerate}

\section{Citations and References}
When you use the citation search feature, references will be automatically managed. For example, you might cite a paper like this: \cite{example2024}. All citations will be automatically added to your \texttt{cite.bib} file.

\section{Conclusion}
This template demonstrates basic \LaTeX{} features including mathematical formulas, tables, and lists. The bibliography system is now set up to work seamlessly with Bibby AI's citation management. Feel free to modify it for your needs!

\bibliographystyle{plain}
\bibliography{cite}

\end{document}\documentclass{article}
\usepackage[utf8]{inputenc}
\usepackage{booktabs}
\usepackage{amsmath}
\usepackage{natbib}

\title{My Document}
\author{Your Name}
\date{\today}

\begin{document}

\maketitle

\section{Introduction}
This is a sample document demonstrating various \LaTeX{} features including mathematical equations, tables, and formatting. You can cite papers using the citation search feature, and they will be automatically added to your \texttt{cite.bib} file.

\section{Mathematical Expressions}
Here are some example equations:

\subsection{Quadratic Formula}
The quadratic formula is given by:
\[x = \frac{-b \pm \sqrt{b^2 - 4ac}}{2a}\]

\subsection{Matrix Example}
A sample matrix:
\[
\begin{pmatrix}
a_{11} & a_{12} & a_{13} \\
a_{21} & a_{22} & a_{23} \\
a_{31} & a_{32} & a_{33}
\end{pmatrix}
\]

\section{Tables}
Below is a sample table showing data:

\begin{table}[h]
\centering
\begin{tabular}{@{}lrrr@{}}
\toprule
\textbf{Category} & \textbf{Value 1} & \textbf{Value 2} & \textbf{Total} \\
\midrule
Item A & 10.5 & 20.3 & 30.8 \\
Item B & 15.2 & 25.1 & 40.3 \\
Item C & 22.1 & 12.7 & 34.8 \\
\bottomrule
\end{tabular}
\caption{Sample Data Table}
\label{tab:sample}
\end{table}

\section{Lists}
\subsection{Bullet Points}
\begin{itemize}
  \item First item
  \item Second item
  \item Third item with sub-items:
    \begin{itemize}
      \item Sub-item 1
      \item Sub-item 2
    \end{itemize}
\end{itemize}

\subsection{Numbered List}
\begin{enumerate}
  \item First step
  \item Second step
  \item Third step
\end{enumerate}

\section{Citations and References}
When you use the citation search feature, references will be automatically managed. For example, you might cite a paper like this: \cite{example2024}. All citations will be automatically added to your \texttt{cite.bib} file.

\section{Conclusion}
This template demonstrates basic \LaTeX{} features including mathematical formulas, tables, and lists. The bibliography system is now set up to work seamlessly with Bibby AI's citation management. Feel free to modify it for your needs!

\bibliographystyle{plain}
\bibliography{cite}

\end{document}\documentclass{article}
\usepackage[utf8]{inputenc}
\usepackage{booktabs}
\usepackage{amsmath}
\usepackage{natbib}

\title{My Document}
\author{Your Name}
\date{\today}

\begin{document}

\maketitle

\section{Introduction}
This is a sample document demonstrating various \LaTeX{} features including mathematical equations, tables, and formatting. You can cite papers using the citation search feature, and they will be automatically added to your \texttt{cite.bib} file.

\section{Mathematical Expressions}
Here are some example equations:

\subsection{Quadratic Formula}
The quadratic formula is given by:
\[x = \frac{-b \pm \sqrt{b^2 - 4ac}}{2a}\]

\subsection{Matrix Example}
A sample matrix:
\[
\begin{pmatrix}
a_{11} & a_{12} & a_{13} \\
a_{21} & a_{22} & a_{23} \\
a_{31} & a_{32} & a_{33}
\end{pmatrix}
\]

\section{Tables}
Below is a sample table showing data:

\begin{table}[h]
\centering
\begin{tabular}{@{}lrrr@{}}
\toprule
\textbf{Category} & \textbf{Value 1} & \textbf{Value 2} & \textbf{Total} \\
\midrule
Item A & 10.5 & 20.3 & 30.8 \\
Item B & 15.2 & 25.1 & 40.3 \\
Item C & 22.1 & 12.7 & 34.8 \\
\bottomrule
\end{tabular}
\caption{Sample Data Table}
\label{tab:sample}
\end{table}

\section{Lists}
\subsection{Bullet Points}
\begin{itemize}
  \item First item
  \item Second item
  \item Third item with sub-items:
    \begin{itemize}
      \item Sub-item 1
      \item Sub-item 2
    \end{itemize}
\end{itemize}

\subsection{Numbered List}
\begin{enumerate}
  \item First step
  \item Second step
  \item Third step
\end{enumerate}

\section{Citations and References}
When you use the citation search feature, references will be automatically managed. For example, you might cite a paper like this: \cite{example2024}. All citations will be automatically added to your \texttt{cite.bib} file.

\section{Conclusion}
This template demonstrates basic \LaTeX{} features including mathematical formulas, tables, and lists. The bibliography system is now set up to work seamlessly with Bibby AI's citation management. Feel free to modify it for your needs!

\bibliographystyle{plain}
\bibliography{cite}

\end{document}\documentclass{article}
\usepackage[utf8]{inputenc}
\usepackage{booktabs}
\usepackage{amsmath}
\usepackage{natbib}

\title{My Document}
\author{Your Name}
\date{\today}

\begin{document}

\maketitle

\section{Introduction}
This is a sample document demonstrating various \LaTeX{} features including mathematical equations, tables, and formatting. You can cite papers using the citation search feature, and they will be automatically added to your \texttt{cite.bib} file.

\section{Mathematical Expressions}
Here are some example equations:

\subsection{Quadratic Formula}
The quadratic formula is given by:
\[x = \frac{-b \pm \sqrt{b^2 - 4ac}}{2a}\]

\subsection{Matrix Example}
A sample matrix:
\[
\begin{pmatrix}
a_{11} & a_{12} & a_{13} \\
a_{21} & a_{22} & a_{23} \\
a_{31} & a_{32} & a_{33}
\end{pmatrix}
\]

\section{Tables}
Below is a sample table showing data:

\begin{table}[h]
\centering
\begin{tabular}{@{}lrrr@{}}
\toprule
\textbf{Category} & \textbf{Value 1} & \textbf{Value 2} & \textbf{Total} \\
\midrule
Item A & 10.5 & 20.3 & 30.8 \\
Item B & 15.2 & 25.1 & 40.3 \\
Item C & 22.1 & 12.7 & 34.8 \\
\bottomrule
\end{tabular}
\caption{Sample Data Table}
\label{tab:sample}
\end{table}

\section{Lists}
\subsection{Bullet Points}
\begin{itemize}
  \item First item
  \item Second item
  \item Third item with sub-items:
    \begin{itemize}
      \item Sub-item 1
      \item Sub-item 2
    \end{itemize}
\end{itemize}

\subsection{Numbered List}
\begin{enumerate}
  \item First step
  \item Second step
  \item Third step
\end{enumerate}

\section{Citations and References}
When you use the citation search feature, references will be automatically managed. For example, you might cite a paper like this: \cite{example2024}. All citations will be automatically added to your \texttt{cite.bib} file.

\section{Conclusion}
This template demonstrates basic \LaTeX{} features including mathematical formulas, tables, and lists. The bibliography system is now set up to work seamlessly with Bibby AI's citation management. Feel free to modify it for your needs!

\bibliographystyle{plain}
\bibliography{cite}

\end{document}\documentclass{article}
\usepackage[utf8]{inputenc}
\usepackage{booktabs}
\usepackage{amsmath}
\usepackage{natbib}

\title{My Document}
\author{Your Name}
\date{\today}

\begin{document}

\maketitle

\section{Introduction}
This is a sample document demonstrating various \LaTeX{} features including mathematical equations, tables, and formatting. You can cite papers using the citation search feature, and they will be automatically added to your \texttt{cite.bib} file.

\section{Mathematical Expressions}
Here are some example equations:

\subsection{Quadratic Formula}
The quadratic formula is given by:
\[x = \frac{-b \pm \sqrt{b^2 - 4ac}}{2a}\]

\subsection{Matrix Example}
A sample matrix:
\[
\begin{pmatrix}
a_{11} & a_{12} & a_{13} \\
a_{21} & a_{22} & a_{23} \\
a_{31} & a_{32} & a_{33}
\end{pmatrix}
\]

\section{Tables}
Below is a sample table showing data:

\begin{table}[h]
\centering
\begin{tabular}{@{}lrrr@{}}
\toprule
\textbf{Category} & \textbf{Value 1} & \textbf{Value 2} & \textbf{Total} \\
\midrule
Item A & 10.5 & 20.3 & 30.8 \\
Item B & 15.2 & 25.1 & 40.3 \\
Item C & 22.1 & 12.7 & 34.8 \\
\bottomrule
\end{tabular}
\caption{Sample Data Table}
\label{tab:sample}
\end{table}

\section{Lists}
\subsection{Bullet Points}
\begin{itemize}
  \item First item
  \item Second item
  \item Third item with sub-items:
    \begin{itemize}
      \item Sub-item 1
      \item Sub-item 2
    \end{itemize}
\end{itemize}

\subsection{Numbered List}
\begin{enumerate}
  \item First step
  \item Second step
  \item Third step
\end{enumerate}

\section{Citations and References}
When you use the citation search feature, references will be automatically managed. For example, you might cite a paper like this: \cite{example2024}. All citations will be automatically added to your \texttt{cite.bib} file.

\section{Conclusion}
This template demonstrates basic \LaTeX{} features including mathematical formulas, tables, and lists. The bibliography system is now set up to work seamlessly with Bibby AI's citation management. Feel free to modify it for your needs!

\bibliographystyle{plain}
\bibliography{cite}

\end{document}\documentclass{article}
\usepackage[utf8]{inputenc}
\usepackage{booktabs}
\usepackage{amsmath}
\usepackage{natbib}

\title{My Document}
\author{Your Name}
\date{\today}

\begin{document}

\maketitle

\section{Introduction}
This is a sample document demonstrating various \LaTeX{} features including mathematical equations, tables, and formatting. You can cite papers using the citation search feature, and they will be automatically added to your \texttt{cite.bib} file.

\section{Mathematical Expressions}
Here are some example equations:

\subsection{Quadratic Formula}
The quadratic formula is given by:
\[x = \frac{-b \pm \sqrt{b^2 - 4ac}}{2a}\]

\subsection{Matrix Example}
A sample matrix:
\[
\begin{pmatrix}
a_{11} & a_{12} & a_{13} \\
a_{21} & a_{22} & a_{23} \\
a_{31} & a_{32} & a_{33}
\end{pmatrix}
\]

\section{Tables}
Below is a sample table showing data:

\begin{table}[h]
\centering
\begin{tabular}{@{}lrrr@{}}
\toprule
\textbf{Category} & \textbf{Value 1} & \textbf{Value 2} & \textbf{Total} \\
\midrule
Item A & 10.5 & 20.3 & 30.8 \\
Item B & 15.2 & 25.1 & 40.3 \\
Item C & 22.1 & 12.7 & 34.8 \\
\bottomrule
\end{tabular}
\caption{Sample Data Table}
\label{tab:sample}
\end{table}

\section{Lists}
\subsection{Bullet Points}
\begin{itemize}
  \item First item
  \item Second item
  \item Third item with sub-items:
    \begin{itemize}
      \item Sub-item 1
      \item Sub-item 2
    \end{itemize}
\end{itemize}

\subsection{Numbered List}
\begin{enumerate}
  \item First step
  \item Second step
  \item Third step
\end{enumerate}

\section{Citations and References}
When you use the citation search feature, references will be automatically managed. For example, you might cite a paper like this: \cite{example2024}. All citations will be automatically added to your \texttt{cite.bib} file.

\section{Conclusion}
This template demonstrates basic \LaTeX{} features including mathematical formulas, tables, and lists. The bibliography system is now set up to work seamlessly with Bibby AI's citation management. Feel free to modify it for your needs!

\bibliographystyle{plain}
\bibliography{cite}

\end{document}\documentclass{article}
\usepackage[utf8]{inputenc}
\usepackage{booktabs}
\usepackage{amsmath}
\usepackage{natbib}

\title{My Document}
\author{Your Name}
\date{\today}

\begin{document}

\maketitle

\section{Introduction}
This is a sample document demonstrating various \LaTeX{} features including mathematical equations, tables, and formatting. You can cite papers using the citation search feature, and they will be automatically added to your \texttt{cite.bib} file.

\section{Mathematical Expressions}
Here are some example equations:

\subsection{Quadratic Formula}
The quadratic formula is given by:
\[x = \frac{-b \pm \sqrt{b^2 - 4ac}}{2a}\]

\subsection{Matrix Example}
A sample matrix:
\[
\begin{pmatrix}
a_{11} & a_{12} & a_{13} \\
a_{21} & a_{22} & a_{23} \\
a_{31} & a_{32} & a_{33}
\end{pmatrix}
\]

\section{Tables}
Below is a sample table showing data:

\begin{table}[h]
\centering
\begin{tabular}{@{}lrrr@{}}
\toprule
\textbf{Category} & \textbf{Value 1} & \textbf{Value 2} & \textbf{Total} \\
\midrule
Item A & 10.5 & 20.3 & 30.8 \\
Item B & 15.2 & 25.1 & 40.3 \\
Item C & 22.1 & 12.7 & 34.8 \\
\bottomrule
\end{tabular}
\caption{Sample Data Table}
\label{tab:sample}
\end{table}

\section{Lists}
\subsection{Bullet Points}
\begin{itemize}
  \item First item
  \item Second item
  \item Third item with sub-items:
    \begin{itemize}
      \item Sub-item 1
      \item Sub-item 2
    \end{itemize}
\end{itemize}

\subsection{Numbered List}
\begin{enumerate}
  \item First step
  \item Second step
  \item Third step
\end{enumerate}

\section{Citations and References}
When you use the citation search feature, references will be automatically managed. For example, you might cite a paper like this: \cite{example2024}. All citations will be automatically added to your \texttt{cite.bib} file.

\section{Conclusion}
This template demonstrates basic \LaTeX{} features including mathematical formulas, tables, and lists. The bibliography system is now set up to work seamlessly with Bibby AI's citation management. Feel free to modify it for your needs!

\bibliographystyle{plain}
\bibliography{cite}

\end{document}\documentclass{article}
\usepackage[utf8]{inputenc}
\usepackage{booktabs}
\usepackage{amsmath}
\usepackage{natbib}

\title{My Document}
\author{Your Name}
\date{\today}

\begin{document}

\maketitle

\section{Introduction}
This is a sample document demonstrating various \LaTeX{} features including mathematical equations, tables, and formatting. You can cite papers using the citation search feature, and they will be automatically added to your \texttt{cite.bib} file.

\section{Mathematical Expressions}
Here are some example equations:

\subsection{Quadratic Formula}
The quadratic formula is given by:
\[x = \frac{-b \pm \sqrt{b^2 - 4ac}}{2a}\]

\subsection{Matrix Example}
A sample matrix:
\[
\begin{pmatrix}
a_{11} & a_{12} & a_{13} \\
a_{21} & a_{22} & a_{23} \\
a_{31} & a_{32} & a_{33}
\end{pmatrix}
\]

\section{Tables}
Below is a sample table showing data:

\begin{table}[h]
\centering
\begin{tabular}{@{}lrrr@{}}
\toprule
\textbf{Category} & \textbf{Value 1} & \textbf{Value 2} & \textbf{Total} \\
\midrule
Item A & 10.5 & 20.3 & 30.8 \\
Item B & 15.2 & 25.1 & 40.3 \\
Item C & 22.1 & 12.7 & 34.8 \\
\bottomrule
\end{tabular}
\caption{Sample Data Table}
\label{tab:sample}
\end{table}

\section{Lists}
\subsection{Bullet Points}
\begin{itemize}
  \item First item
  \item Second item
  \item Third item with sub-items:
    \begin{itemize}
      \item Sub-item 1
      \item Sub-item 2
    \end{itemize}
\end{itemize}

\subsection{Numbered List}
\begin{enumerate}
  \item First step
  \item Second step
  \item Third step
\end{enumerate}

\section{Citations and References}
When you use the citation search feature, references will be automatically managed. For example, you might cite a paper like this: \cite{example2024}. All citations will be automatically added to your \texttt{cite.bib} file.

\section{Conclusion}
This template demonstrates basic \LaTeX{} features including mathematical formulas, tables, and lists. The bibliography system is now set up to work seamlessly with Bibby AI's citation management. Feel free to modify it for your needs!

\bibliographystyle{plain}
\bibliography{cite}

\end{document}\documentclass{article}
\usepackage[utf8]{inputenc}
\usepackage{booktabs}
\usepackage{amsmath}
\usepackage{natbib}

\title{My Document}
\author{Your Name}
\date{\today}

\begin{document}

\maketitle

\section{Introduction}
This is a sample document demonstrating various \LaTeX{} features including mathematical equations, tables, and formatting. You can cite papers using the citation search feature, and they will be automatically added to your \texttt{cite.bib} file.

\section{Mathematical Expressions}
Here are some example equations:

\subsection{Quadratic Formula}
The quadratic formula is given by:
\[x = \frac{-b \pm \sqrt{b^2 - 4ac}}{2a}\]

\subsection{Matrix Example}
A sample matrix:
\[
\begin{pmatrix}
a_{11} & a_{12} & a_{13} \\
a_{21} & a_{22} & a_{23} \\
a_{31} & a_{32} & a_{33}
\end{pmatrix}
\]

\section{Tables}
Below is a sample table showing data:

\begin{table}[h]
\centering
\begin{tabular}{@{}lrrr@{}}
\toprule
\textbf{Category} & \textbf{Value 1} & \textbf{Value 2} & \textbf{Total} \\
\midrule
Item A & 10.5 & 20.3 & 30.8 \\
Item B & 15.2 & 25.1 & 40.3 \\
Item C & 22.1 & 12.7 & 34.8 \\
\bottomrule
\end{tabular}
\caption{Sample Data Table}
\label{tab:sample}
\end{table}

\section{Lists}
\subsection{Bullet Points}
\begin{itemize}
  \item First item
  \item Second item
  \item Third item with sub-items:
    \begin{itemize}
      \item Sub-item 1
      \item Sub-item 2
    \end{itemize}
\end{itemize}

\subsection{Numbered List}
\begin{enumerate}
  \item First step
  \item Second step
  \item Third step
\end{enumerate}

\section{Citations and References}
When you use the citation search feature, references will be automatically managed. For example, you might cite a paper like this: \cite{example2024}. All citations will be automatically added to your \texttt{cite.bib} file.

\section{Conclusion}
This template demonstrates basic \LaTeX{} features including mathematical formulas, tables, and lists. The bibliography system is now set up to work seamlessly with Bibby AI's citation management. Feel free to modify it for your needs!

\bibliographystyle{plain}
\bibliography{cite}

\end{document}\documentclass{article}
\usepackage[utf8]{inputenc}
\usepackage{booktabs}
\usepackage{amsmath}
\usepackage{natbib}

\title{My Document}
\author{Your Name}
\date{\today}

\begin{document}

\maketitle

\section{Introduction}
This is a sample document demonstrating various \LaTeX{} features including mathematical equations, tables, and formatting. You can cite papers using the citation search feature, and they will be automatically added to your \texttt{cite.bib} file.

\section{Mathematical Expressions}
Here are some example equations:

\subsection{Quadratic Formula}
The quadratic formula is given by:
\[x = \frac{-b \pm \sqrt{b^2 - 4ac}}{2a}\]

\subsection{Matrix Example}
A sample matrix:
\[
\begin{pmatrix}
a_{11} & a_{12} & a_{13} \\
a_{21} & a_{22} & a_{23} \\
a_{31} & a_{32} & a_{33}
\end{pmatrix}
\]

\section{Tables}
Below is a sample table showing data:

\begin{table}[h]
\centering
\begin{tabular}{@{}lrrr@{}}
\toprule
\textbf{Category} & \textbf{Value 1} & \textbf{Value 2} & \textbf{Total} \\
\midrule
Item A & 10.5 & 20.3 & 30.8 \\
Item B & 15.2 & 25.1 & 40.3 \\
Item C & 22.1 & 12.7 & 34.8 \\
\bottomrule
\end{tabular}
\caption{Sample Data Table}
\label{tab:sample}
\end{table}

\section{Lists}
\subsection{Bullet Points}
\begin{itemize}
  \item First item
  \item Second item
  \item Third item with sub-items:
    \begin{itemize}
      \item Sub-item 1
      \item Sub-item 2
    \end{itemize}
\end{itemize}

\subsection{Numbered List}
\begin{enumerate}
  \item First step
  \item Second step
  \item Third step
\end{enumerate}

\section{Citations and References}
When you use the citation search feature, references will be automatically managed. For example, you might cite a paper like this: \cite{example2024}. All citations will be automatically added to your \texttt{cite.bib} file.

\section{Conclusion}
This template demonstrates basic \LaTeX{} features including mathematical formulas, tables, and lists. The bibliography system is now set up to work seamlessly with Bibby AI's citation management. Feel free to modify it for your needs!

\bibliographystyle{plain}
\bibliography{cite}

\end{document}\documentclass{article}
\usepackage[utf8]{inputenc}
\usepackage{booktabs}
\usepackage{amsmath}
\usepackage{natbib}

\title{My Document}
\author{Your Name}
\date{\today}

\begin{document}

\maketitle

\section{Introduction}
This is a sample document demonstrating various \LaTeX{} features including mathematical equations, tables, and formatting. You can cite papers using the citation search feature, and they will be automatically added to your \texttt{cite.bib} file.

\section{Mathematical Expressions}
Here are some example equations:

\subsection{Quadratic Formula}
The quadratic formula is given by:
\[x = \frac{-b \pm \sqrt{b^2 - 4ac}}{2a}\]

\subsection{Matrix Example}
A sample matrix:
\[
\begin{pmatrix}
a_{11} & a_{12} & a_{13} \\
a_{21} & a_{22} & a_{23} \\
a_{31} & a_{32} & a_{33}
\end{pmatrix}
\]

\section{Tables}
Below is a sample table showing data:

\begin{table}[h]
\centering
\begin{tabular}{@{}lrrr@{}}
\toprule
\textbf{Category} & \textbf{Value 1} & \textbf{Value 2} & \textbf{Total} \\
\midrule
Item A & 10.5 & 20.3 & 30.8 \\
Item B & 15.2 & 25.1 & 40.3 \\
Item C & 22.1 & 12.7 & 34.8 \\
\bottomrule
\end{tabular}
\caption{Sample Data Table}
\label{tab:sample}
\end{table}

\section{Lists}
\subsection{Bullet Points}
\begin{itemize}
  \item First item
  \item Second item
  \item Third item with sub-items:
    \begin{itemize}
      \item Sub-item 1
      \item Sub-item 2
    \end{itemize}
\end{itemize}

\subsection{Numbered List}
\begin{enumerate}
  \item First step
  \item Second step
  \item Third step
\end{enumerate}

\section{Citations and References}
When you use the citation search feature, references will be automatically managed. For example, you might cite a paper like this: \cite{example2024}. All citations will be automatically added to your \texttt{cite.bib} file.

\section{Conclusion}
This template demonstrates basic \LaTeX{} features including mathematical formulas, tables, and lists. The bibliography system is now set up to work seamlessly with Bibby AI's citation management. Feel free to modify it for your needs!

\bibliographystyle{plain}
\bibliography{cite}

\end{document}\documentclass{article}
\usepackage[utf8]{inputenc}
\usepackage{booktabs}
\usepackage{amsmath}
\usepackage{natbib}

\title{My Document}
\author{Your Name}
\date{\today}

\begin{document}

\maketitle

\section{Introduction}
This is a sample document demonstrating various \LaTeX{} features including mathematical equations, tables, and formatting. You can cite papers using the citation search feature, and they will be automatically added to your \texttt{cite.bib} file.

\section{Mathematical Expressions}
Here are some example equations:

\subsection{Quadratic Formula}
The quadratic formula is given by:
\[x = \frac{-b \pm \sqrt{b^2 - 4ac}}{2a}\]

\subsection{Matrix Example}
A sample matrix:
\[
\begin{pmatrix}
a_{11} & a_{12} & a_{13} \\
a_{21} & a_{22} & a_{23} \\
a_{31} & a_{32} & a_{33}
\end{pmatrix}
\]

\section{Tables}
Below is a sample table showing data:

\begin{table}[h]
\centering
\begin{tabular}{@{}lrrr@{}}
\toprule
\textbf{Category} & \textbf{Value 1} & \textbf{Value 2} & \textbf{Total} \\
\midrule
Item A & 10.5 & 20.3 & 30.8 \\
Item B & 15.2 & 25.1 & 40.3 \\
Item C & 22.1 & 12.7 & 34.8 \\
\bottomrule
\end{tabular}
\caption{Sample Data Table}
\label{tab:sample}
\end{table}

\section{Lists}
\subsection{Bullet Points}
\begin{itemize}
  \item First item
  \item Second item
  \item Third item with sub-items:
    \begin{itemize}
      \item Sub-item 1
      \item Sub-item 2
    \end{itemize}
\end{itemize}

\subsection{Numbered List}
\begin{enumerate}
  \item First step
  \item Second step
  \item Third step
\end{enumerate}

\section{Citations and References}
When you use the citation search feature, references will be automatically managed. For example, you might cite a paper like this: \cite{example2024}. All citations will be automatically added to your \texttt{cite.bib} file.

\section{Conclusion}
This template demonstrates basic \LaTeX{} features including mathematical formulas, tables, and lists. The bibliography system is now set up to work seamlessly with Bibby AI's citation management. Feel free to modify it for your needs!

\bibliographystyle{plain}
\bibliography{cite}

\end{document}\documentclass{article}
\usepackage[utf8]{inputenc}
\usepackage{booktabs}
\usepackage{amsmath}
\usepackage{natbib}

\title{My Document}
\author{Your Name}
\date{\today}

\begin{document}

\maketitle

\section{Introduction}
This is a sample document demonstrating various \LaTeX{} features including mathematical equations, tables, and formatting. You can cite papers using the citation search feature, and they will be automatically added to your \texttt{cite.bib} file.

\section{Mathematical Expressions}
Here are some example equations:

\subsection{Quadratic Formula}
The quadratic formula is given by:
\[x = \frac{-b \pm \sqrt{b^2 - 4ac}}{2a}\]

\subsection{Matrix Example}
A sample matrix:
\[
\begin{pmatrix}
a_{11} & a_{12} & a_{13} \\
a_{21} & a_{22} & a_{23} \\
a_{31} & a_{32} & a_{33}
\end{pmatrix}
\]

\section{Tables}
Below is a sample table showing data:

\begin{table}[h]
\centering
\begin{tabular}{@{}lrrr@{}}
\toprule
\textbf{Category} & \textbf{Value 1} & \textbf{Value 2} & \textbf{Total} \\
\midrule
Item A & 10.5 & 20.3 & 30.8 \\
Item B & 15.2 & 25.1 & 40.3 \\
Item C & 22.1 & 12.7 & 34.8 \\
\bottomrule
\end{tabular}
\caption{Sample Data Table}
\label{tab:sample}
\end{table}

\section{Lists}
\subsection{Bullet Points}
\begin{itemize}
  \item First item
  \item Second item
  \item Third item with sub-items:
    \begin{itemize}
      \item Sub-item 1
      \item Sub-item 2
    \end{itemize}
\end{itemize}

\subsection{Numbered List}
\begin{enumerate}
  \item First step
  \item Second step
  \item Third step
\end{enumerate}

\section{Citations and References}
When you use the citation search feature, references will be automatically managed. For example, you might cite a paper like this: \cite{example2024}. All citations will be automatically added to your \texttt{cite.bib} file.

\section{Conclusion}
This template demonstrates basic \LaTeX{} features including mathematical formulas, tables, and lists. The bibliography system is now set up to work seamlessly with Bibby AI's citation management. Feel free to modify it for your needs!

\bibliographystyle{plain}
\bibliography{cite}

\end{document}\documentclass{article}
\usepackage[utf8]{inputenc}
\usepackage{booktabs}
\usepackage{amsmath}
\usepackage{natbib}

\title{My Document}
\author{Your Name}
\date{\today}

\begin{document}

\maketitle

\section{Introduction}
This is a sample document demonstrating various \LaTeX{} features including mathematical equations, tables, and formatting. You can cite papers using the citation search feature, and they will be automatically added to your \texttt{cite.bib} file.

\section{Mathematical Expressions}
Here are some example equations:

\subsection{Quadratic Formula}
The quadratic formula is given by:
\[x = \frac{-b \pm \sqrt{b^2 - 4ac}}{2a}\]

\subsection{Matrix Example}
A sample matrix:
\[
\begin{pmatrix}
a_{11} & a_{12} & a_{13} \\
a_{21} & a_{22} & a_{23} \\
a_{31} & a_{32} & a_{33}
\end{pmatrix}
\]

\section{Tables}
Below is a sample table showing data:

\begin{table}[h]
\centering
\begin{tabular}{@{}lrrr@{}}
\toprule
\textbf{Category} & \textbf{Value 1} & \textbf{Value 2} & \textbf{Total} \\
\midrule
Item A & 10.5 & 20.3 & 30.8 \\
Item B & 15.2 & 25.1 & 40.3 \\
Item C & 22.1 & 12.7 & 34.8 \\
\bottomrule
\end{tabular}
\caption{Sample Data Table}
\label{tab:sample}
\end{table}

\section{Lists}
\subsection{Bullet Points}
\begin{itemize}
  \item First item
  \item Second item
  \item Third item with sub-items:
    \begin{itemize}
      \item Sub-item 1
      \item Sub-item 2
    \end{itemize}
\end{itemize}

\subsection{Numbered List}
\begin{enumerate}
  \item First step
  \item Second step
  \item Third step
\end{enumerate}

\section{Citations and References}
When you use the citation search feature, references will be automatically managed. For example, you might cite a paper like this: \cite{example2024}. All citations will be automatically added to your \texttt{cite.bib} file.

\section{Conclusion}
This template demonstrates basic \LaTeX{} features including mathematical formulas, tables, and lists. The bibliography system is now set up to work seamlessly with Bibby AI's citation management. Feel free to modify it for your needs!

\bibliographystyle{plain}
\bibliography{cite}

\end{document}\documentclass{article}
\usepackage[utf8]{inputenc}
\usepackage{booktabs}
\usepackage{amsmath}
\usepackage{natbib}

\title{My Document}
\author{Your Name}
\date{\today}

\begin{document}

\maketitle

\section{Introduction}
This is a sample document demonstrating various \LaTeX{} features including mathematical equations, tables, and formatting. You can cite papers using the citation search feature, and they will be automatically added to your \texttt{cite.bib} file.

\section{Mathematical Expressions}
Here are some example equations:

\subsection{Quadratic Formula}
The quadratic formula is given by:
\[x = \frac{-b \pm \sqrt{b^2 - 4ac}}{2a}\]

\subsection{Matrix Example}
A sample matrix:
\[
\begin{pmatrix}
a_{11} & a_{12} & a_{13} \\
a_{21} & a_{22} & a_{23} \\
a_{31} & a_{32} & a_{33}
\end{pmatrix}
\]

\section{Tables}
Below is a sample table showing data:

\begin{table}[h]
\centering
\begin{tabular}{@{}lrrr@{}}
\toprule
\textbf{Category} & \textbf{Value 1} & \textbf{Value 2} & \textbf{Total} \\
\midrule
Item A & 10.5 & 20.3 & 30.8 \\
Item B & 15.2 & 25.1 & 40.3 \\
Item C & 22.1 & 12.7 & 34.8 \\
\bottomrule
\end{tabular}
\caption{Sample Data Table}
\label{tab:sample}
\end{table}

\section{Lists}
\subsection{Bullet Points}
\begin{itemize}
  \item First item
  \item Second item
  \item Third item with sub-items:
    \begin{itemize}
      \item Sub-item 1
      \item Sub-item 2
    \end{itemize}
\end{itemize}

\subsection{Numbered List}
\begin{enumerate}
  \item First step
  \item Second step
  \item Third step
\end{enumerate}

\section{Citations and References}
When you use the citation search feature, references will be automatically managed. For example, you might cite a paper like this: \cite{example2024}. All citations will be automatically added to your \texttt{cite.bib} file.

\section{Conclusion}
This template demonstrates basic \LaTeX{} features including mathematical formulas, tables, and lists. The bibliography system is now set up to work seamlessly with Bibby AI's citation management. Feel free to modify it for your needs!

\bibliographystyle{plain}
\bibliography{cite}

\end{document}\documentclass{article}
\usepackage[utf8]{inputenc}
\usepackage{booktabs}
\usepackage{amsmath}
\usepackage{natbib}

\title{My Document}
\author{Your Name}
\date{\today}

\begin{document}

\maketitle

\section{Introduction}
This is a sample document demonstrating various \LaTeX{} features including mathematical equations, tables, and formatting. You can cite papers using the citation search feature, and they will be automatically added to your \texttt{cite.bib} file.

\section{Mathematical Expressions}
Here are some example equations:

\subsection{Quadratic Formula}
The quadratic formula is given by:
\[x = \frac{-b \pm \sqrt{b^2 - 4ac}}{2a}\]

\subsection{Matrix Example}
A sample matrix:
\[
\begin{pmatrix}
a_{11} & a_{12} & a_{13} \\
a_{21} & a_{22} & a_{23} \\
a_{31} & a_{32} & a_{33}
\end{pmatrix}
\]

\section{Tables}
Below is a sample table showing data:

\begin{table}[h]
\centering
\begin{tabular}{@{}lrrr@{}}
\toprule
\textbf{Category} & \textbf{Value 1} & \textbf{Value 2} & \textbf{Total} \\
\midrule
Item A & 10.5 & 20.3 & 30.8 \\
Item B & 15.2 & 25.1 & 40.3 \\
Item C & 22.1 & 12.7 & 34.8 \\
\bottomrule
\end{tabular}
\caption{Sample Data Table}
\label{tab:sample}
\end{table}

\section{Lists}
\subsection{Bullet Points}
\begin{itemize}
  \item First item
  \item Second item
  \item Third item with sub-items:
    \begin{itemize}
      \item Sub-item 1
      \item Sub-item 2
    \end{itemize}
\end{itemize}

\subsection{Numbered List}
\begin{enumerate}
  \item First step
  \item Second step
  \item Third step
\end{enumerate}

\section{Citations and References}
When you use the citation search feature, references will be automatically managed. For example, you might cite a paper like this: \cite{example2024}. All citations will be automatically added to your \texttt{cite.bib} file.

\section{Conclusion}
This template demonstrates basic \LaTeX{} features including mathematical formulas, tables, and lists. The bibliography system is now set up to work seamlessly with Bibby AI's citation management. Feel free to modify it for your needs!

\bibliographystyle{plain}
\bibliography{cite}

\end{document}\documentclass{article}
\usepackage[utf8]{inputenc}
\usepackage{booktabs}
\usepackage{amsmath}
\usepackage{natbib}

\title{My Document}
\author{Your Name}
\date{\today}

\begin{document}

\maketitle

\section{Introduction}
This is a sample document demonstrating various \LaTeX{} features including mathematical equations, tables, and formatting. You can cite papers using the citation search feature, and they will be automatically added to your \texttt{cite.bib} file.

\section{Mathematical Expressions}
Here are some example equations:

\subsection{Quadratic Formula}
The quadratic formula is given by:
\[x = \frac{-b \pm \sqrt{b^2 - 4ac}}{2a}\]

\subsection{Matrix Example}
A sample matrix:
\[
\begin{pmatrix}
a_{11} & a_{12} & a_{13} \\
a_{21} & a_{22} & a_{23} \\
a_{31} & a_{32} & a_{33}
\end{pmatrix}
\]

\section{Tables}
Below is a sample table showing data:

\begin{table}[h]
\centering
\begin{tabular}{@{}lrrr@{}}
\toprule
\textbf{Category} & \textbf{Value 1} & \textbf{Value 2} & \textbf{Total} \\
\midrule
Item A & 10.5 & 20.3 & 30.8 \\
Item B & 15.2 & 25.1 & 40.3 \\
Item C & 22.1 & 12.7 & 34.8 \\
\bottomrule
\end{tabular}
\caption{Sample Data Table}
\label{tab:sample}
\end{table}

\section{Lists}
\subsection{Bullet Points}
\begin{itemize}
  \item First item
  \item Second item
  \item Third item with sub-items:
    \begin{itemize}
      \item Sub-item 1
      \item Sub-item 2
    \end{itemize}
\end{itemize}

\subsection{Numbered List}
\begin{enumerate}
  \item First step
  \item Second step
  \item Third step
\end{enumerate}

\section{Citations and References}
When you use the citation search feature, references will be automatically managed. For example, you might cite a paper like this: \cite{example2024}. All citations will be automatically added to your \texttt{cite.bib} file.

\section{Conclusion}
This template demonstrates basic \LaTeX{} features including mathematical formulas, tables, and lists. The bibliography system is now set up to work seamlessly with Bibby AI's citation management. Feel free to modify it for your needs!

\bibliographystyle{plain}
\bibliography{cite}

\end{document}\documentclass{article}
\usepackage[utf8]{inputenc}
\usepackage{booktabs}
\usepackage{amsmath}
\usepackage{natbib}

\title{My Document}
\author{Your Name}
\date{\today}

\begin{document}

\maketitle

\section{Introduction}
This is a sample document demonstrating various \LaTeX{} features including mathematical equations, tables, and formatting. You can cite papers using the citation search feature, and they will be automatically added to your \texttt{cite.bib} file.

\section{Mathematical Expressions}
Here are some example equations:

\subsection{Quadratic Formula}
The quadratic formula is given by:
\[x = \frac{-b \pm \sqrt{b^2 - 4ac}}{2a}\]

\subsection{Matrix Example}
A sample matrix:
\[
\begin{pmatrix}
a_{11} & a_{12} & a_{13} \\
a_{21} & a_{22} & a_{23} \\
a_{31} & a_{32} & a_{33}
\end{pmatrix}
\]

\section{Tables}
Below is a sample table showing data:

\begin{table}[h]
\centering
\begin{tabular}{@{}lrrr@{}}
\toprule
\textbf{Category} & \textbf{Value 1} & \textbf{Value 2} & \textbf{Total} \\
\midrule
Item A & 10.5 & 20.3 & 30.8 \\
Item B & 15.2 & 25.1 & 40.3 \\
Item C & 22.1 & 12.7 & 34.8 \\
\bottomrule
\end{tabular}
\caption{Sample Data Table}
\label{tab:sample}
\end{table}

\section{Lists}
\subsection{Bullet Points}
\begin{itemize}
  \item First item
  \item Second item
  \item Third item with sub-items:
    \begin{itemize}
      \item Sub-item 1
      \item Sub-item 2
    \end{itemize}
\end{itemize}

\subsection{Numbered List}
\begin{enumerate}
  \item First step
  \item Second step
  \item Third step
\end{enumerate}

\section{Citations and References}
When you use the citation search feature, references will be automatically managed. For example, you might cite a paper like this: \cite{example2024}. All citations will be automatically added to your \texttt{cite.bib} file.

\section{Conclusion}
This template demonstrates basic \LaTeX{} features including mathematical formulas, tables, and lists. The bibliography system is now set up to work seamlessly with Bibby AI's citation management. Feel free to modify it for your needs!

\bibliographystyle{plain}
\bibliography{cite}

\end{document}\documentclass{article}
\usepackage[utf8]{inputenc}
\usepackage{booktabs}
\usepackage{amsmath}
\usepackage{natbib}

\title{My Document}
\author{Your Name}
\date{\today}

\begin{document}

\maketitle

\section{Introduction}
This is a sample document demonstrating various \LaTeX{} features including mathematical equations, tables, and formatting. You can cite papers using the citation search feature, and they will be automatically added to your \texttt{cite.bib} file.

\section{Mathematical Expressions}
Here are some example equations:

\subsection{Quadratic Formula}
The quadratic formula is given by:
\[x = \frac{-b \pm \sqrt{b^2 - 4ac}}{2a}\]

\subsection{Matrix Example}
A sample matrix:
\[
\begin{pmatrix}
a_{11} & a_{12} & a_{13} \\
a_{21} & a_{22} & a_{23} \\
a_{31} & a_{32} & a_{33}
\end{pmatrix}
\]

\section{Tables}
Below is a sample table showing data:

\begin{table}[h]
\centering
\begin{tabular}{@{}lrrr@{}}
\toprule
\textbf{Category} & \textbf{Value 1} & \textbf{Value 2} & \textbf{Total} \\
\midrule
Item A & 10.5 & 20.3 & 30.8 \\
Item B & 15.2 & 25.1 & 40.3 \\
Item C & 22.1 & 12.7 & 34.8 \\
\bottomrule
\end{tabular}
\caption{Sample Data Table}
\label{tab:sample}
\end{table}

\section{Lists}
\subsection{Bullet Points}
\begin{itemize}
  \item First item
  \item Second item
  \item Third item with sub-items:
    \begin{itemize}
      \item Sub-item 1
      \item Sub-item 2
    \end{itemize}
\end{itemize}

\subsection{Numbered List}
\begin{enumerate}
  \item First step
  \item Second step
  \item Third step
\end{enumerate}

\section{Citations and References}
When you use the citation search feature, references will be automatically managed. For example, you might cite a paper like this: \cite{example2024}. All citations will be automatically added to your \texttt{cite.bib} file.

\section{Conclusion}
This template demonstrates basic \LaTeX{} features including mathematical formulas, tables, and lists. The bibliography system is now set up to work seamlessly with Bibby AI's citation management. Feel free to modify it for your needs!

\bibliographystyle{plain}
\bibliography{cite}

\end{document}\documentclass{article}
\usepackage[utf8]{inputenc}
\usepackage{booktabs}
\usepackage{amsmath}
\usepackage{natbib}

\title{My Document}
\author{Your Name}
\date{\today}

\begin{document}

\maketitle

\section{Introduction}
This is a sample document demonstrating various \LaTeX{} features including mathematical equations, tables, and formatting. You can cite papers using the citation search feature, and they will be automatically added to your \texttt{cite.bib} file.

\section{Mathematical Expressions}
Here are some example equations:

\subsection{Quadratic Formula}
The quadratic formula is given by:
\[x = \frac{-b \pm \sqrt{b^2 - 4ac}}{2a}\]

\subsection{Matrix Example}
A sample matrix:
\[
\begin{pmatrix}
a_{11} & a_{12} & a_{13} \\
a_{21} & a_{22} & a_{23} \\
a_{31} & a_{32} & a_{33}
\end{pmatrix}
\]

\section{Tables}
Below is a sample table showing data:

\begin{table}[h]
\centering
\begin{tabular}{@{}lrrr@{}}
\toprule
\textbf{Category} & \textbf{Value 1} & \textbf{Value 2} & \textbf{Total} \\
\midrule
Item A & 10.5 & 20.3 & 30.8 \\
Item B & 15.2 & 25.1 & 40.3 \\
Item C & 22.1 & 12.7 & 34.8 \\
\bottomrule
\end{tabular}
\caption{Sample Data Table}
\label{tab:sample}
\end{table}

\section{Lists}
\subsection{Bullet Points}
\begin{itemize}
  \item First item
  \item Second item
  \item Third item with sub-items:
    \begin{itemize}
      \item Sub-item 1
      \item Sub-item 2
    \end{itemize}
\end{itemize}

\subsection{Numbered List}
\begin{enumerate}
  \item First step
  \item Second step
  \item Third step
\end{enumerate}

\section{Citations and References}
When you use the citation search feature, references will be automatically managed. For example, you might cite a paper like this: \cite{example2024}. All citations will be automatically added to your \texttt{cite.bib} file.

\section{Conclusion}
This template demonstrates basic \LaTeX{} features including mathematical formulas, tables, and lists. The bibliography system is now set up to work seamlessly with Bibby AI's citation management. Feel free to modify it for your needs!

\bibliographystyle{plain}
\bibliography{cite}

\end{document}\documentclass{article}
\usepackage[utf8]{inputenc}
\usepackage{booktabs}
\usepackage{amsmath}
\usepackage{natbib}

\title{My Document}
\author{Your Name}
\date{\today}

\begin{document}

\maketitle

\section{Introduction}
This is a sample document demonstrating various \LaTeX{} features including mathematical equations, tables, and formatting. You can cite papers using the citation search feature, and they will be automatically added to your \texttt{cite.bib} file.

\section{Mathematical Expressions}
Here are some example equations:

\subsection{Quadratic Formula}
The quadratic formula is given by:
\[x = \frac{-b \pm \sqrt{b^2 - 4ac}}{2a}\]

\subsection{Matrix Example}
A sample matrix:
\[
\begin{pmatrix}
a_{11} & a_{12} & a_{13} \\
a_{21} & a_{22} & a_{23} \\
a_{31} & a_{32} & a_{33}
\end{pmatrix}
\]

\section{Tables}
Below is a sample table showing data:

\begin{table}[h]
\centering
\begin{tabular}{@{}lrrr@{}}
\toprule
\textbf{Category} & \textbf{Value 1} & \textbf{Value 2} & \textbf{Total} \\
\midrule
Item A & 10.5 & 20.3 & 30.8 \\
Item B & 15.2 & 25.1 & 40.3 \\
Item C & 22.1 & 12.7 & 34.8 \\
\bottomrule
\end{tabular}
\caption{Sample Data Table}
\label{tab:sample}
\end{table}

\section{Lists}
\subsection{Bullet Points}
\begin{itemize}
  \item First item
  \item Second item
  \item Third item with sub-items:
    \begin{itemize}
      \item Sub-item 1
      \item Sub-item 2
    \end{itemize}
\end{itemize}

\subsection{Numbered List}
\begin{enumerate}
  \item First step
  \item Second step
  \item Third step
\end{enumerate}

\section{Citations and References}
When you use the citation search feature, references will be automatically managed. For example, you might cite a paper like this: \cite{example2024}. All citations will be automatically added to your \texttt{cite.bib} file.

\section{Conclusion}
This template demonstrates basic \LaTeX{} features including mathematical formulas, tables, and lists. The bibliography system is now set up to work seamlessly with Bibby AI's citation management. Feel free to modify it for your needs!

\bibliographystyle{plain}
\bibliography{cite}

\end{document}\documentclass{article}
\usepackage[utf8]{inputenc}
\usepackage{booktabs}
\usepackage{amsmath}
\usepackage{natbib}

\title{My Document}
\author{Your Name}
\date{\today}

\begin{document}

\maketitle

\section{Introduction}
This is a sample document demonstrating various \LaTeX{} features including mathematical equations, tables, and formatting. You can cite papers using the citation search feature, and they will be automatically added to your \texttt{cite.bib} file.

\section{Mathematical Expressions}
Here are some example equations:

\subsection{Quadratic Formula}
The quadratic formula is given by:
\[x = \frac{-b \pm \sqrt{b^2 - 4ac}}{2a}\]

\subsection{Matrix Example}
A sample matrix:
\[
\begin{pmatrix}
a_{11} & a_{12} & a_{13} \\
a_{21} & a_{22} & a_{23} \\
a_{31} & a_{32} & a_{33}
\end{pmatrix}
\]

\section{Tables}
Below is a sample table showing data:

\begin{table}[h]
\centering
\begin{tabular}{@{}lrrr@{}}
\toprule
\textbf{Category} & \textbf{Value 1} & \textbf{Value 2} & \textbf{Total} \\
\midrule
Item A & 10.5 & 20.3 & 30.8 \\
Item B & 15.2 & 25.1 & 40.3 \\
Item C & 22.1 & 12.7 & 34.8 \\
\bottomrule
\end{tabular}
\caption{Sample Data Table}
\label{tab:sample}
\end{table}

\section{Lists}
\subsection{Bullet Points}
\begin{itemize}
  \item First item
  \item Second item
  \item Third item with sub-items:
    \begin{itemize}
      \item Sub-item 1
      \item Sub-item 2
    \end{itemize}
\end{itemize}

\subsection{Numbered List}
\begin{enumerate}
  \item First step
  \item Second step
  \item Third step
\end{enumerate}

\section{Citations and References}
When you use the citation search feature, references will be automatically managed. For example, you might cite a paper like this: \cite{example2024}. All citations will be automatically added to your \texttt{cite.bib} file.

\section{Conclusion}
This template demonstrates basic \LaTeX{} features including mathematical formulas, tables, and lists. The bibliography system is now set up to work seamlessly with Bibby AI's citation management. Feel free to modify it for your needs!

\bibliographystyle{plain}
\bibliography{cite}

\end{document}\documentclass{article}
\usepackage[utf8]{inputenc}
\usepackage{booktabs}
\usepackage{amsmath}
\usepackage{natbib}

\title{My Document}
\author{Your Name}
\date{\today}

\begin{document}

\maketitle

\section{Introduction}
This is a sample document demonstrating various \LaTeX{} features including mathematical equations, tables, and formatting. You can cite papers using the citation search feature, and they will be automatically added to your \texttt{cite.bib} file.

\section{Mathematical Expressions}
Here are some example equations:

\subsection{Quadratic Formula}
The quadratic formula is given by:
\[x = \frac{-b \pm \sqrt{b^2 - 4ac}}{2a}\]

\subsection{Matrix Example}
A sample matrix:
\[
\begin{pmatrix}
a_{11} & a_{12} & a_{13} \\
a_{21} & a_{22} & a_{23} \\
a_{31} & a_{32} & a_{33}
\end{pmatrix}
\]

\section{Tables}
Below is a sample table showing data:

\begin{table}[h]
\centering
\begin{tabular}{@{}lrrr@{}}
\toprule
\textbf{Category} & \textbf{Value 1} & \textbf{Value 2} & \textbf{Total} \\
\midrule
Item A & 10.5 & 20.3 & 30.8 \\
Item B & 15.2 & 25.1 & 40.3 \\
Item C & 22.1 & 12.7 & 34.8 \\
\bottomrule
\end{tabular}
\caption{Sample Data Table}
\label{tab:sample}
\end{table}

\section{Lists}
\subsection{Bullet Points}
\begin{itemize}
  \item First item
  \item Second item
  \item Third item with sub-items:
    \begin{itemize}
      \item Sub-item 1
      \item Sub-item 2
    \end{itemize}
\end{itemize}

\subsection{Numbered List}
\begin{enumerate}
  \item First step
  \item Second step
  \item Third step
\end{enumerate}

\section{Citations and References}
When you use the citation search feature, references will be automatically managed. For example, you might cite a paper like this: \cite{example2024}. All citations will be automatically added to your \texttt{cite.bib} file.

\section{Conclusion}
This template demonstrates basic \LaTeX{} features including mathematical formulas, tables, and lists. The bibliography system is now set up to work seamlessly with Bibby AI's citation management. Feel free to modify it for your needs!

\bibliographystyle{plain}
\bibliography{cite}

\end{document}\documentclass{article}
\usepackage[utf8]{inputenc}
\usepackage{booktabs}
\usepackage{amsmath}
\usepackage{natbib}

\title{My Document}
\author{Your Name}
\date{\today}

\begin{document}

\maketitle

\section{Introduction}
This is a sample document demonstrating various \LaTeX{} features including mathematical equations, tables, and formatting. You can cite papers using the citation search feature, and they will be automatically added to your \texttt{cite.bib} file.

\section{Mathematical Expressions}
Here are some example equations:

\subsection{Quadratic Formula}
The quadratic formula is given by:
\[x = \frac{-b \pm \sqrt{b^2 - 4ac}}{2a}\]

\subsection{Matrix Example}
A sample matrix:
\[
\begin{pmatrix}
a_{11} & a_{12} & a_{13} \\
a_{21} & a_{22} & a_{23} \\
a_{31} & a_{32} & a_{33}
\end{pmatrix}
\]

\section{Tables}
Below is a sample table showing data:

\begin{table}[h]
\centering
\begin{tabular}{@{}lrrr@{}}
\toprule
\textbf{Category} & \textbf{Value 1} & \textbf{Value 2} & \textbf{Total} \\
\midrule
Item A & 10.5 & 20.3 & 30.8 \\
Item B & 15.2 & 25.1 & 40.3 \\
Item C & 22.1 & 12.7 & 34.8 \\
\bottomrule
\end{tabular}
\caption{Sample Data Table}
\label{tab:sample}
\end{table}

\section{Lists}
\subsection{Bullet Points}
\begin{itemize}
  \item First item
  \item Second item
  \item Third item with sub-items:
    \begin{itemize}
      \item Sub-item 1
      \item Sub-item 2
    \end{itemize}
\end{itemize}

\subsection{Numbered List}
\begin{enumerate}
  \item First step
  \item Second step
  \item Third step
\end{enumerate}

\section{Citations and References}
When you use the citation search feature, references will be automatically managed. For example, you might cite a paper like this: \cite{example2024}. All citations will be automatically added to your \texttt{cite.bib} file.

\section{Conclusion}
This template demonstrates basic \LaTeX{} features including mathematical formulas, tables, and lists. The bibliography system is now set up to work seamlessly with Bibby AI's citation management. Feel free to modify it for your needs!

\bibliographystyle{plain}
\bibliography{cite}

\end{document}\documentclass{article}
\usepackage[utf8]{inputenc}
\usepackage{booktabs}
\usepackage{amsmath}
\usepackage{natbib}

\title{My Document}
\author{Your Name}
\date{\today}

\begin{document}

\maketitle

\section{Introduction}
This is a sample document demonstrating various \LaTeX{} features including mathematical equations, tables, and formatting. You can cite papers using the citation search feature, and they will be automatically added to your \texttt{cite.bib} file.

\section{Mathematical Expressions}
Here are some example equations:

\subsection{Quadratic Formula}
The quadratic formula is given by:
\[x = \frac{-b \pm \sqrt{b^2 - 4ac}}{2a}\]

\subsection{Matrix Example}
A sample matrix:
\[
\begin{pmatrix}
a_{11} & a_{12} & a_{13} \\
a_{21} & a_{22} & a_{23} \\
a_{31} & a_{32} & a_{33}
\end{pmatrix}
\]

\section{Tables}
Below is a sample table showing data:

\begin{table}[h]
\centering
\begin{tabular}{@{}lrrr@{}}
\toprule
\textbf{Category} & \textbf{Value 1} & \textbf{Value 2} & \textbf{Total} \\
\midrule
Item A & 10.5 & 20.3 & 30.8 \\
Item B & 15.2 & 25.1 & 40.3 \\
Item C & 22.1 & 12.7 & 34.8 \\
\bottomrule
\end{tabular}
\caption{Sample Data Table}
\label{tab:sample}
\end{table}

\section{Lists}
\subsection{Bullet Points}
\begin{itemize}
  \item First item
  \item Second item
  \item Third item with sub-items:
    \begin{itemize}
      \item Sub-item 1
      \item Sub-item 2
    \end{itemize}
\end{itemize}

\subsection{Numbered List}
\begin{enumerate}
  \item First step
  \item Second step
  \item Third step
\end{enumerate}

\section{Citations and References}
When you use the citation search feature, references will be automatically managed. For example, you might cite a paper like this: \cite{example2024}. All citations will be automatically added to your \texttt{cite.bib} file.

\section{Conclusion}
This template demonstrates basic \LaTeX{} features including mathematical formulas, tables, and lists. The bibliography system is now set up to work seamlessly with Bibby AI's citation management. Feel free to modify it for your needs!

\bibliographystyle{plain}
\bibliography{cite}

\end{document}\documentclass{article}
\usepackage[utf8]{inputenc}
\usepackage{booktabs}
\usepackage{amsmath}
\usepackage{natbib}

\title{My Document}
\author{Your Name}
\date{\today}

\begin{document}

\maketitle

\section{Introduction}
This is a sample document demonstrating various \LaTeX{} features including mathematical equations, tables, and formatting. You can cite papers using the citation search feature, and they will be automatically added to your \texttt{cite.bib} file.

\section{Mathematical Expressions}
Here are some example equations:

\subsection{Quadratic Formula}
The quadratic formula is given by:
\[x = \frac{-b \pm \sqrt{b^2 - 4ac}}{2a}\]

\subsection{Matrix Example}
A sample matrix:
\[
\begin{pmatrix}
a_{11} & a_{12} & a_{13} \\
a_{21} & a_{22} & a_{23} \\
a_{31} & a_{32} & a_{33}
\end{pmatrix}
\]

\section{Tables}
Below is a sample table showing data:

\begin{table}[h]
\centering
\begin{tabular}{@{}lrrr@{}}
\toprule
\textbf{Category} & \textbf{Value 1} & \textbf{Value 2} & \textbf{Total} \\
\midrule
Item A & 10.5 & 20.3 & 30.8 \\
Item B & 15.2 & 25.1 & 40.3 \\
Item C & 22.1 & 12.7 & 34.8 \\
\bottomrule
\end{tabular}
\caption{Sample Data Table}
\label{tab:sample}
\end{table}

\section{Lists}
\subsection{Bullet Points}
\begin{itemize}
  \item First item
  \item Second item
  \item Third item with sub-items:
    \begin{itemize}
      \item Sub-item 1
      \item Sub-item 2
    \end{itemize}
\end{itemize}

\subsection{Numbered List}
\begin{enumerate}
  \item First step
  \item Second step
  \item Third step
\end{enumerate}

\section{Citations and References}
When you use the citation search feature, references will be automatically managed. For example, you might cite a paper like this: \cite{example2024}. All citations will be automatically added to your \texttt{cite.bib} file.

\section{Conclusion}
This template demonstrates basic \LaTeX{} features including mathematical formulas, tables, and lists. The bibliography system is now set up to work seamlessly with Bibby AI's citation management. Feel free to modify it for your needs!

\bibliographystyle{plain}
\bibliography{cite}

\end{document}\documentclass{article}
\usepackage[utf8]{inputenc}
\usepackage{booktabs}
\usepackage{amsmath}
\usepackage{natbib}

\title{My Document}
\author{Your Name}
\date{\today}

\begin{document}

\maketitle

\section{Introduction}
This is a sample document demonstrating various \LaTeX{} features including mathematical equations, tables, and formatting. You can cite papers using the citation search feature, and they will be automatically added to your \texttt{cite.bib} file.

\section{Mathematical Expressions}
Here are some example equations:

\subsection{Quadratic Formula}
The quadratic formula is given by:
\[x = \frac{-b \pm \sqrt{b^2 - 4ac}}{2a}\]

\subsection{Matrix Example}
A sample matrix:
\[
\begin{pmatrix}
a_{11} & a_{12} & a_{13} \\
a_{21} & a_{22} & a_{23} \\
a_{31} & a_{32} & a_{33}
\end{pmatrix}
\]

\section{Tables}
Below is a sample table showing data:

\begin{table}[h]
\centering
\begin{tabular}{@{}lrrr@{}}
\toprule
\textbf{Category} & \textbf{Value 1} & \textbf{Value 2} & \textbf{Total} \\
\midrule
Item A & 10.5 & 20.3 & 30.8 \\
Item B & 15.2 & 25.1 & 40.3 \\
Item C & 22.1 & 12.7 & 34.8 \\
\bottomrule
\end{tabular}
\caption{Sample Data Table}
\label{tab:sample}
\end{table}

\section{Lists}
\subsection{Bullet Points}
\begin{itemize}
  \item First item
  \item Second item
  \item Third item with sub-items:
    \begin{itemize}
      \item Sub-item 1
      \item Sub-item 2
    \end{itemize}
\end{itemize}

\subsection{Numbered List}
\begin{enumerate}
  \item First step
  \item Second step
  \item Third step
\end{enumerate}

\section{Citations and References}
When you use the citation search feature, references will be automatically managed. For example, you might cite a paper like this: \cite{example2024}. All citations will be automatically added to your \texttt{cite.bib} file.

\section{Conclusion}
This template demonstrates basic \LaTeX{} features including mathematical formulas, tables, and lists. The bibliography system is now set up to work seamlessly with Bibby AI's citation management. Feel free to modify it for your needs!

\bibliographystyle{plain}
\bibliography{cite}

\end{document}\documentclass{article}
\usepackage[utf8]{inputenc}
\usepackage{booktabs}
\usepackage{amsmath}
\usepackage{natbib}

\title{My Document}
\author{Your Name}
\date{\today}

\begin{document}

\maketitle

\section{Introduction}
This is a sample document demonstrating various \LaTeX{} features including mathematical equations, tables, and formatting. You can cite papers using the citation search feature, and they will be automatically added to your \texttt{cite.bib} file.

\section{Mathematical Expressions}
Here are some example equations:

\subsection{Quadratic Formula}
The quadratic formula is given by:
\[x = \frac{-b \pm \sqrt{b^2 - 4ac}}{2a}\]

\subsection{Matrix Example}
A sample matrix:
\[
\begin{pmatrix}
a_{11} & a_{12} & a_{13} \\
a_{21} & a_{22} & a_{23} \\
a_{31} & a_{32} & a_{33}
\end{pmatrix}
\]

\section{Tables}
Below is a sample table showing data:

\begin{table}[h]
\centering
\begin{tabular}{@{}lrrr@{}}
\toprule
\textbf{Category} & \textbf{Value 1} & \textbf{Value 2} & \textbf{Total} \\
\midrule
Item A & 10.5 & 20.3 & 30.8 \\
Item B & 15.2 & 25.1 & 40.3 \\
Item C & 22.1 & 12.7 & 34.8 \\
\bottomrule
\end{tabular}
\caption{Sample Data Table}
\label{tab:sample}
\end{table}

\section{Lists}
\subsection{Bullet Points}
\begin{itemize}
  \item First item
  \item Second item
  \item Third item with sub-items:
    \begin{itemize}
      \item Sub-item 1
      \item Sub-item 2
    \end{itemize}
\end{itemize}

\subsection{Numbered List}
\begin{enumerate}
  \item First step
  \item Second step
  \item Third step
\end{enumerate}

\section{Citations and References}
When you use the citation search feature, references will be automatically managed. For example, you might cite a paper like this: \cite{example2024}. All citations will be automatically added to your \texttt{cite.bib} file.

\section{Conclusion}
This template demonstrates basic \LaTeX{} features including mathematical formulas, tables, and lists. The bibliography system is now set up to work seamlessly with Bibby AI's citation management. Feel free to modify it for your needs!

\bibliographystyle{plain}
\bibliography{cite}

\end{document}\documentclass{article}
\usepackage[utf8]{inputenc}
\usepackage{booktabs}
\usepackage{amsmath}
\usepackage{natbib}

\title{My Document}
\author{Your Name}
\date{\today}

\begin{document}

\maketitle

\section{Introduction}
This is a sample document demonstrating various \LaTeX{} features including mathematical equations, tables, and formatting. You can cite papers using the citation search feature, and they will be automatically added to your \texttt{cite.bib} file.

\section{Mathematical Expressions}
Here are some example equations:

\subsection{Quadratic Formula}
The quadratic formula is given by:
\[x = \frac{-b \pm \sqrt{b^2 - 4ac}}{2a}\]

\subsection{Matrix Example}
A sample matrix:
\[
\begin{pmatrix}
a_{11} & a_{12} & a_{13} \\
a_{21} & a_{22} & a_{23} \\
a_{31} & a_{32} & a_{33}
\end{pmatrix}
\]

\section{Tables}
Below is a sample table showing data:

\begin{table}[h]
\centering
\begin{tabular}{@{}lrrr@{}}
\toprule
\textbf{Category} & \textbf{Value 1} & \textbf{Value 2} & \textbf{Total} \\
\midrule
Item A & 10.5 & 20.3 & 30.8 \\
Item B & 15.2 & 25.1 & 40.3 \\
Item C & 22.1 & 12.7 & 34.8 \\
\bottomrule
\end{tabular}
\caption{Sample Data Table}
\label{tab:sample}
\end{table}

\section{Lists}
\subsection{Bullet Points}
\begin{itemize}
  \item First item
  \item Second item
  \item Third item with sub-items:
    \begin{itemize}
      \item Sub-item 1
      \item Sub-item 2
    \end{itemize}
\end{itemize}

\subsection{Numbered List}
\begin{enumerate}
  \item First step
  \item Second step
  \item Third step
\end{enumerate}

\section{Citations and References}
When you use the citation search feature, references will be automatically managed. For example, you might cite a paper like this: \cite{example2024}. All citations will be automatically added to your \texttt{cite.bib} file.

\section{Conclusion}
This template demonstrates basic \LaTeX{} features including mathematical formulas, tables, and lists. The bibliography system is now set up to work seamlessly with Bibby AI's citation management. Feel free to modify it for your needs!

\bibliographystyle{plain}
\bibliography{cite}

\end{document}\documentclass{article}
\usepackage[utf8]{inputenc}
\usepackage{booktabs}
\usepackage{amsmath}
\usepackage{natbib}

\title{My Document}
\author{Your Name}
\date{\today}

\begin{document}

\maketitle

\section{Introduction}
This is a sample document demonstrating various \LaTeX{} features including mathematical equations, tables, and formatting. You can cite papers using the citation search feature, and they will be automatically added to your \texttt{cite.bib} file.

\section{Mathematical Expressions}
Here are some example equations:

\subsection{Quadratic Formula}
The quadratic formula is given by:
\[x = \frac{-b \pm \sqrt{b^2 - 4ac}}{2a}\]

\subsection{Matrix Example}
A sample matrix:
\[
\begin{pmatrix}
a_{11} & a_{12} & a_{13} \\
a_{21} & a_{22} & a_{23} \\
a_{31} & a_{32} & a_{33}
\end{pmatrix}
\]

\section{Tables}
Below is a sample table showing data:

\begin{table}[h]
\centering
\begin{tabular}{@{}lrrr@{}}
\toprule
\textbf{Category} & \textbf{Value 1} & \textbf{Value 2} & \textbf{Total} \\
\midrule
Item A & 10.5 & 20.3 & 30.8 \\
Item B & 15.2 & 25.1 & 40.3 \\
Item C & 22.1 & 12.7 & 34.8 \\
\bottomrule
\end{tabular}
\caption{Sample Data Table}
\label{tab:sample}
\end{table}

\section{Lists}
\subsection{Bullet Points}
\begin{itemize}
  \item First item
  \item Second item
  \item Third item with sub-items:
    \begin{itemize}
      \item Sub-item 1
      \item Sub-item 2
    \end{itemize}
\end{itemize}

\subsection{Numbered List}
\begin{enumerate}
  \item First step
  \item Second step
  \item Third step
\end{enumerate}

\section{Citations and References}
When you use the citation search feature, references will be automatically managed. For example, you might cite a paper like this: \cite{example2024}. All citations will be automatically added to your \texttt{cite.bib} file.

\section{Conclusion}
This template demonstrates basic \LaTeX{} features including mathematical formulas, tables, and lists. The bibliography system is now set up to work seamlessly with Bibby AI's citation management. Feel free to modify it for your needs!

\bibliographystyle{plain}
\bibliography{cite}

\end{document}\documentclass{article}
\usepackage[utf8]{inputenc}
\usepackage{booktabs}
\usepackage{amsmath}
\usepackage{natbib}

\title{My Document}
\author{Your Name}
\date{\today}

\begin{document}

\maketitle

\section{Introduction}
This is a sample document demonstrating various \LaTeX{} features including mathematical equations, tables, and formatting. You can cite papers using the citation search feature, and they will be automatically added to your \texttt{cite.bib} file.

\section{Mathematical Expressions}
Here are some example equations:

\subsection{Quadratic Formula}
The quadratic formula is given by:
\[x = \frac{-b \pm \sqrt{b^2 - 4ac}}{2a}\]

\subsection{Matrix Example}
A sample matrix:
\[
\begin{pmatrix}
a_{11} & a_{12} & a_{13} \\
a_{21} & a_{22} & a_{23} \\
a_{31} & a_{32} & a_{33}
\end{pmatrix}
\]

\section{Tables}
Below is a sample table showing data:

\begin{table}[h]
\centering
\begin{tabular}{@{}lrrr@{}}
\toprule
\textbf{Category} & \textbf{Value 1} & \textbf{Value 2} & \textbf{Total} \\
\midrule
Item A & 10.5 & 20.3 & 30.8 \\
Item B & 15.2 & 25.1 & 40.3 \\
Item C & 22.1 & 12.7 & 34.8 \\
\bottomrule
\end{tabular}
\caption{Sample Data Table}
\label{tab:sample}
\end{table}

\section{Lists}
\subsection{Bullet Points}
\begin{itemize}
  \item First item
  \item Second item
  \item Third item with sub-items:
    \begin{itemize}
      \item Sub-item 1
      \item Sub-item 2
    \end{itemize}
\end{itemize}

\subsection{Numbered List}
\begin{enumerate}
  \item First step
  \item Second step
  \item Third step
\end{enumerate}

\section{Citations and References}
When you use the citation search feature, references will be automatically managed. For example, you might cite a paper like this: \cite{example2024}. All citations will be automatically added to your \texttt{cite.bib} file.

\section{Conclusion}
This template demonstrates basic \LaTeX{} features including mathematical formulas, tables, and lists. The bibliography system is now set up to work seamlessly with Bibby AI's citation management. Feel free to modify it for your needs!

\bibliographystyle{plain}
\bibliography{cite}

\end{document}\documentclass{article}
\usepackage[utf8]{inputenc}
\usepackage{booktabs}
\usepackage{amsmath}
\usepackage{natbib}

\title{My Document}
\author{Your Name}
\date{\today}

\begin{document}

\maketitle

\section{Introduction}
This is a sample document demonstrating various \LaTeX{} features including mathematical equations, tables, and formatting. You can cite papers using the citation search feature, and they will be automatically added to your \texttt{cite.bib} file.

\section{Mathematical Expressions}
Here are some example equations:

\subsection{Quadratic Formula}
The quadratic formula is given by:
\[x = \frac{-b \pm \sqrt{b^2 - 4ac}}{2a}\]

\subsection{Matrix Example}
A sample matrix:
\[
\begin{pmatrix}
a_{11} & a_{12} & a_{13} \\
a_{21} & a_{22} & a_{23} \\
a_{31} & a_{32} & a_{33}
\end{pmatrix}
\]

\section{Tables}
Below is a sample table showing data:

\begin{table}[h]
\centering
\begin{tabular}{@{}lrrr@{}}
\toprule
\textbf{Category} & \textbf{Value 1} & \textbf{Value 2} & \textbf{Total} \\
\midrule
Item A & 10.5 & 20.3 & 30.8 \\
Item B & 15.2 & 25.1 & 40.3 \\
Item C & 22.1 & 12.7 & 34.8 \\
\bottomrule
\end{tabular}
\caption{Sample Data Table}
\label{tab:sample}
\end{table}

\section{Lists}
\subsection{Bullet Points}
\begin{itemize}
  \item First item
  \item Second item
  \item Third item with sub-items:
    \begin{itemize}
      \item Sub-item 1
      \item Sub-item 2
    \end{itemize}
\end{itemize}

\subsection{Numbered List}
\begin{enumerate}
  \item First step
  \item Second step
  \item Third step
\end{enumerate}

\section{Citations and References}
When you use the citation search feature, references will be automatically managed. For example, you might cite a paper like this: \cite{example2024}. All citations will be automatically added to your \texttt{cite.bib} file.

\section{Conclusion}
This template demonstrates basic \LaTeX{} features including mathematical formulas, tables, and lists. The bibliography system is now set up to work seamlessly with Bibby AI's citation management. Feel free to modify it for your needs!

\bibliographystyle{plain}
\bibliography{cite}

\end{document}\documentclass{article}
\usepackage[utf8]{inputenc}
\usepackage{booktabs}
\usepackage{amsmath}
\usepackage{natbib}

\title{My Document}
\author{Your Name}
\date{\today}

\begin{document}

\maketitle

\section{Introduction}
This is a sample document demonstrating various \LaTeX{} features including mathematical equations, tables, and formatting. You can cite papers using the citation search feature, and they will be automatically added to your \texttt{cite.bib} file.

\section{Mathematical Expressions}
Here are some example equations:

\subsection{Quadratic Formula}
The quadratic formula is given by:
\[x = \frac{-b \pm \sqrt{b^2 - 4ac}}{2a}\]

\subsection{Matrix Example}
A sample matrix:
\[
\begin{pmatrix}
a_{11} & a_{12} & a_{13} \\
a_{21} & a_{22} & a_{23} \\
a_{31} & a_{32} & a_{33}
\end{pmatrix}
\]

\section{Tables}
Below is a sample table showing data:

\begin{table}[h]
\centering
\begin{tabular}{@{}lrrr@{}}
\toprule
\textbf{Category} & \textbf{Value 1} & \textbf{Value 2} & \textbf{Total} \\
\midrule
Item A & 10.5 & 20.3 & 30.8 \\
Item B & 15.2 & 25.1 & 40.3 \\
Item C & 22.1 & 12.7 & 34.8 \\
\bottomrule
\end{tabular}
\caption{Sample Data Table}
\label{tab:sample}
\end{table}

\section{Lists}
\subsection{Bullet Points}
\begin{itemize}
  \item First item
  \item Second item
  \item Third item with sub-items:
    \begin{itemize}
      \item Sub-item 1
      \item Sub-item 2
    \end{itemize}
\end{itemize}

\subsection{Numbered List}
\begin{enumerate}
  \item First step
  \item Second step
  \item Third step
\end{enumerate}

\section{Citations and References}
When you use the citation search feature, references will be automatically managed. For example, you might cite a paper like this: \cite{example2024}. All citations will be automatically added to your \texttt{cite.bib} file.

\section{Conclusion}
This template demonstrates basic \LaTeX{} features including mathematical formulas, tables, and lists. The bibliography system is now set up to work seamlessly with Bibby AI's citation management. Feel free to modify it for your needs!

\bibliographystyle{plain}
\bibliography{cite}

\end{document}\documentclass{article}
\usepackage[utf8]{inputenc}
\usepackage{booktabs}
\usepackage{amsmath}
\usepackage{natbib}

\title{My Document}
\author{Your Name}
\date{\today}

\begin{document}

\maketitle

\section{Introduction}
This is a sample document demonstrating various \LaTeX{} features including mathematical equations, tables, and formatting. You can cite papers using the citation search feature, and they will be automatically added to your \texttt{cite.bib} file.

\section{Mathematical Expressions}
Here are some example equations:

\subsection{Quadratic Formula}
The quadratic formula is given by:
\[x = \frac{-b \pm \sqrt{b^2 - 4ac}}{2a}\]

\subsection{Matrix Example}
A sample matrix:
\[
\begin{pmatrix}
a_{11} & a_{12} & a_{13} \\
a_{21} & a_{22} & a_{23} \\
a_{31} & a_{32} & a_{33}
\end{pmatrix}
\]

\section{Tables}
Below is a sample table showing data:

\begin{table}[h]
\centering
\begin{tabular}{@{}lrrr@{}}
\toprule
\textbf{Category} & \textbf{Value 1} & \textbf{Value 2} & \textbf{Total} \\
\midrule
Item A & 10.5 & 20.3 & 30.8 \\
Item B & 15.2 & 25.1 & 40.3 \\
Item C & 22.1 & 12.7 & 34.8 \\
\bottomrule
\end{tabular}
\caption{Sample Data Table}
\label{tab:sample}
\end{table}

\section{Lists}
\subsection{Bullet Points}
\begin{itemize}
  \item First item
  \item Second item
  \item Third item with sub-items:
    \begin{itemize}
      \item Sub-item 1
      \item Sub-item 2
    \end{itemize}
\end{itemize}

\subsection{Numbered List}
\begin{enumerate}
  \item First step
  \item Second step
  \item Third step
\end{enumerate}

\section{Citations and References}
When you use the citation search feature, references will be automatically managed. For example, you might cite a paper like this: \cite{example2024}. All citations will be automatically added to your \texttt{cite.bib} file.

\section{Conclusion}
This template demonstrates basic \LaTeX{} features including mathematical formulas, tables, and lists. The bibliography system is now set up to work seamlessly with Bibby AI's citation management. Feel free to modify it for your needs!

\bibliographystyle{plain}
\bibliography{cite}

\end{document}\documentclass{article}
\usepackage[utf8]{inputenc}
\usepackage{booktabs}
\usepackage{amsmath}
\usepackage{natbib}

\title{My Document}
\author{Your Name}
\date{\today}

\begin{document}

\maketitle

\section{Introduction}
This is a sample document demonstrating various \LaTeX{} features including mathematical equations, tables, and formatting. You can cite papers using the citation search feature, and they will be automatically added to your \texttt{cite.bib} file.

\section{Mathematical Expressions}
Here are some example equations:

\subsection{Quadratic Formula}
The quadratic formula is given by:
\[x = \frac{-b \pm \sqrt{b^2 - 4ac}}{2a}\]

\subsection{Matrix Example}
A sample matrix:
\[
\begin{pmatrix}
a_{11} & a_{12} & a_{13} \\
a_{21} & a_{22} & a_{23} \\
a_{31} & a_{32} & a_{33}
\end{pmatrix}
\]

\section{Tables}
Below is a sample table showing data:

\begin{table}[h]
\centering
\begin{tabular}{@{}lrrr@{}}
\toprule
\textbf{Category} & \textbf{Value 1} & \textbf{Value 2} & \textbf{Total} \\
\midrule
Item A & 10.5 & 20.3 & 30.8 \\
Item B & 15.2 & 25.1 & 40.3 \\
Item C & 22.1 & 12.7 & 34.8 \\
\bottomrule
\end{tabular}
\caption{Sample Data Table}
\label{tab:sample}
\end{table}

\section{Lists}
\subsection{Bullet Points}
\begin{itemize}
  \item First item
  \item Second item
  \item Third item with sub-items:
    \begin{itemize}
      \item Sub-item 1
      \item Sub-item 2
    \end{itemize}
\end{itemize}

\subsection{Numbered List}
\begin{enumerate}
  \item First step
  \item Second step
  \item Third step
\end{enumerate}

\section{Citations and References}
When you use the citation search feature, references will be automatically managed. For example, you might cite a paper like this: \cite{example2024}. All citations will be automatically added to your \texttt{cite.bib} file.

\section{Conclusion}
This template demonstrates basic \LaTeX{} features including mathematical formulas, tables, and lists. The bibliography system is now set up to work seamlessly with Bibby AI's citation management. Feel free to modify it for your needs!

\bibliographystyle{plain}
\bibliography{cite}

\end{document}\documentclass{article}
\usepackage[utf8]{inputenc}
\usepackage{booktabs}
\usepackage{amsmath}
\usepackage{natbib}

\title{My Document}
\author{Your Name}
\date{\today}

\begin{document}

\maketitle

\section{Introduction}
This is a sample document demonstrating various \LaTeX{} features including mathematical equations, tables, and formatting. You can cite papers using the citation search feature, and they will be automatically added to your \texttt{cite.bib} file.

\section{Mathematical Expressions}
Here are some example equations:

\subsection{Quadratic Formula}
The quadratic formula is given by:
\[x = \frac{-b \pm \sqrt{b^2 - 4ac}}{2a}\]

\subsection{Matrix Example}
A sample matrix:
\[
\begin{pmatrix}
a_{11} & a_{12} & a_{13} \\
a_{21} & a_{22} & a_{23} \\
a_{31} & a_{32} & a_{33}
\end{pmatrix}
\]

\section{Tables}
Below is a sample table showing data:

\begin{table}[h]
\centering
\begin{tabular}{@{}lrrr@{}}
\toprule
\textbf{Category} & \textbf{Value 1} & \textbf{Value 2} & \textbf{Total} \\
\midrule
Item A & 10.5 & 20.3 & 30.8 \\
Item B & 15.2 & 25.1 & 40.3 \\
Item C & 22.1 & 12.7 & 34.8 \\
\bottomrule
\end{tabular}
\caption{Sample Data Table}
\label{tab:sample}
\end{table}

\section{Lists}
\subsection{Bullet Points}
\begin{itemize}
  \item First item
  \item Second item
  \item Third item with sub-items:
    \begin{itemize}
      \item Sub-item 1
      \item Sub-item 2
    \end{itemize}
\end{itemize}

\subsection{Numbered List}
\begin{enumerate}
  \item First step
  \item Second step
  \item Third step
\end{enumerate}

\section{Citations and References}
When you use the citation search feature, references will be automatically managed. For example, you might cite a paper like this: \cite{example2024}. All citations will be automatically added to your \texttt{cite.bib} file.

\section{Conclusion}
This template demonstrates basic \LaTeX{} features including mathematical formulas, tables, and lists. The bibliography system is now set up to work seamlessly with Bibby AI's citation management. Feel free to modify it for your needs!

\bibliographystyle{plain}
\bibliography{cite}

\end{document}\documentclass{article}
\usepackage[utf8]{inputenc}
\usepackage{booktabs}
\usepackage{amsmath}
\usepackage{natbib}

\title{My Document}
\author{Your Name}
\date{\today}

\begin{document}

\maketitle

\section{Introduction}
This is a sample document demonstrating various \LaTeX{} features including mathematical equations, tables, and formatting. You can cite papers using the citation search feature, and they will be automatically added to your \texttt{cite.bib} file.

\section{Mathematical Expressions}
Here are some example equations:

\subsection{Quadratic Formula}
The quadratic formula is given by:
\[x = \frac{-b \pm \sqrt{b^2 - 4ac}}{2a}\]

\subsection{Matrix Example}
A sample matrix:
\[
\begin{pmatrix}
a_{11} & a_{12} & a_{13} \\
a_{21} & a_{22} & a_{23} \\
a_{31} & a_{32} & a_{33}
\end{pmatrix}
\]

\section{Tables}
Below is a sample table showing data:

\begin{table}[h]
\centering
\begin{tabular}{@{}lrrr@{}}
\toprule
\textbf{Category} & \textbf{Value 1} & \textbf{Value 2} & \textbf{Total} \\
\midrule
Item A & 10.5 & 20.3 & 30.8 \\
Item B & 15.2 & 25.1 & 40.3 \\
Item C & 22.1 & 12.7 & 34.8 \\
\bottomrule
\end{tabular}
\caption{Sample Data Table}
\label{tab:sample}
\end{table}

\section{Lists}
\subsection{Bullet Points}
\begin{itemize}
  \item First item
  \item Second item
  \item Third item with sub-items:
    \begin{itemize}
      \item Sub-item 1
      \item Sub-item 2
    \end{itemize}
\end{itemize}

\subsection{Numbered List}
\begin{enumerate}
  \item First step
  \item Second step
  \item Third step
\end{enumerate}

\section{Citations and References}
When you use the citation search feature, references will be automatically managed. For example, you might cite a paper like this: \cite{example2024}. All citations will be automatically added to your \texttt{cite.bib} file.

\section{Conclusion}
This template demonstrates basic \LaTeX{} features including mathematical formulas, tables, and lists. The bibliography system is now set up to work seamlessly with Bibby AI's citation management. Feel free to modify it for your needs!

\bibliographystyle{plain}
\bibliography{cite}

\end{document}\documentclass{article}
\usepackage[utf8]{inputenc}
\usepackage{booktabs}
\usepackage{amsmath}
\usepackage{natbib}

\title{My Document}
\author{Your Name}
\date{\today}

\begin{document}

\maketitle

\section{Introduction}
This is a sample document demonstrating various \LaTeX{} features including mathematical equations, tables, and formatting. You can cite papers using the citation search feature, and they will be automatically added to your \texttt{cite.bib} file.

\section{Mathematical Expressions}
Here are some example equations:

\subsection{Quadratic Formula}
The quadratic formula is given by:
\[x = \frac{-b \pm \sqrt{b^2 - 4ac}}{2a}\]

\subsection{Matrix Example}
A sample matrix:
\[
\begin{pmatrix}
a_{11} & a_{12} & a_{13} \\
a_{21} & a_{22} & a_{23} \\
a_{31} & a_{32} & a_{33}
\end{pmatrix}
\]

\section{Tables}
Below is a sample table showing data:

\begin{table}[h]
\centering
\begin{tabular}{@{}lrrr@{}}
\toprule
\textbf{Category} & \textbf{Value 1} & \textbf{Value 2} & \textbf{Total} \\
\midrule
Item A & 10.5 & 20.3 & 30.8 \\
Item B & 15.2 & 25.1 & 40.3 \\
Item C & 22.1 & 12.7 & 34.8 \\
\bottomrule
\end{tabular}
\caption{Sample Data Table}
\label{tab:sample}
\end{table}

\section{Lists}
\subsection{Bullet Points}
\begin{itemize}
  \item First item
  \item Second item
  \item Third item with sub-items:
    \begin{itemize}
      \item Sub-item 1
      \item Sub-item 2
    \end{itemize}
\end{itemize}

\subsection{Numbered List}
\begin{enumerate}
  \item First step
  \item Second step
  \item Third step
\end{enumerate}

\section{Citations and References}
When you use the citation search feature, references will be automatically managed. For example, you might cite a paper like this: \cite{example2024}. All citations will be automatically added to your \texttt{cite.bib} file.

\section{Conclusion}
This template demonstrates basic \LaTeX{} features including mathematical formulas, tables, and lists. The bibliography system is now set up to work seamlessly with Bibby AI's citation management. Feel free to modify it for your needs!

\bibliographystyle{plain}
\bibliography{cite}

\end{document}\documentclass{article}
\usepackage[utf8]{inputenc}
\usepackage{booktabs}
\usepackage{amsmath}
\usepackage{natbib}

\title{My Document}
\author{Your Name}
\date{\today}

\begin{document}

\maketitle

\section{Introduction}
This is a sample document demonstrating various \LaTeX{} features including mathematical equations, tables, and formatting. You can cite papers using the citation search feature, and they will be automatically added to your \texttt{cite.bib} file.

\section{Mathematical Expressions}
Here are some example equations:

\subsection{Quadratic Formula}
The quadratic formula is given by:
\[x = \frac{-b \pm \sqrt{b^2 - 4ac}}{2a}\]

\subsection{Matrix Example}
A sample matrix:
\[
\begin{pmatrix}
a_{11} & a_{12} & a_{13} \\
a_{21} & a_{22} & a_{23} \\
a_{31} & a_{32} & a_{33}
\end{pmatrix}
\]

\section{Tables}
Below is a sample table showing data:

\begin{table}[h]
\centering
\begin{tabular}{@{}lrrr@{}}
\toprule
\textbf{Category} & \textbf{Value 1} & \textbf{Value 2} & \textbf{Total} \\
\midrule
Item A & 10.5 & 20.3 & 30.8 \\
Item B & 15.2 & 25.1 & 40.3 \\
Item C & 22.1 & 12.7 & 34.8 \\
\bottomrule
\end{tabular}
\caption{Sample Data Table}
\label{tab:sample}
\end{table}

\section{Lists}
\subsection{Bullet Points}
\begin{itemize}
  \item First item
  \item Second item
  \item Third item with sub-items:
    \begin{itemize}
      \item Sub-item 1
      \item Sub-item 2
    \end{itemize}
\end{itemize}

\subsection{Numbered List}
\begin{enumerate}
  \item First step
  \item Second step
  \item Third step
\end{enumerate}

\section{Citations and References}
When you use the citation search feature, references will be automatically managed. For example, you might cite a paper like this: \cite{example2024}. All citations will be automatically added to your \texttt{cite.bib} file.

\section{Conclusion}
This template demonstrates basic \LaTeX{} features including mathematical formulas, tables, and lists. The bibliography system is now set up to work seamlessly with Bibby AI's citation management. Feel free to modify it for your needs!

\bibliographystyle{plain}
\bibliography{cite}

\end{document}\documentclass{article}
\usepackage[utf8]{inputenc}
\usepackage{booktabs}
\usepackage{amsmath}
\usepackage{natbib}

\title{My Document}
\author{Your Name}
\date{\today}

\begin{document}

\maketitle

\section{Introduction}
This is a sample document demonstrating various \LaTeX{} features including mathematical equations, tables, and formatting. You can cite papers using the citation search feature, and they will be automatically added to your \texttt{cite.bib} file.

\section{Mathematical Expressions}
Here are some example equations:

\subsection{Quadratic Formula}
The quadratic formula is given by:
\[x = \frac{-b \pm \sqrt{b^2 - 4ac}}{2a}\]

\subsection{Matrix Example}
A sample matrix:
\[
\begin{pmatrix}
a_{11} & a_{12} & a_{13} \\
a_{21} & a_{22} & a_{23} \\
a_{31} & a_{32} & a_{33}
\end{pmatrix}
\]

\section{Tables}
Below is a sample table showing data:

\begin{table}[h]
\centering
\begin{tabular}{@{}lrrr@{}}
\toprule
\textbf{Category} & \textbf{Value 1} & \textbf{Value 2} & \textbf{Total} \\
\midrule
Item A & 10.5 & 20.3 & 30.8 \\
Item B & 15.2 & 25.1 & 40.3 \\
Item C & 22.1 & 12.7 & 34.8 \\
\bottomrule
\end{tabular}
\caption{Sample Data Table}
\label{tab:sample}
\end{table}

\section{Lists}
\subsection{Bullet Points}
\begin{itemize}
  \item First item
  \item Second item
  \item Third item with sub-items:
    \begin{itemize}
      \item Sub-item 1
      \item Sub-item 2
    \end{itemize}
\end{itemize}

\subsection{Numbered List}
\begin{enumerate}
  \item First step
  \item Second step
  \item Third step
\end{enumerate}

\section{Citations and References}
When you use the citation search feature, references will be automatically managed. For example, you might cite a paper like this: \cite{example2024}. All citations will be automatically added to your \texttt{cite.bib} file.

\section{Conclusion}
This template demonstrates basic \LaTeX{} features including mathematical formulas, tables, and lists. The bibliography system is now set up to work seamlessly with Bibby AI's citation management. Feel free to modify it for your needs!

\bibliographystyle{plain}
\bibliography{cite}

\end{document}\documentclass{article}
\usepackage[utf8]{inputenc}
\usepackage{booktabs}
\usepackage{amsmath}
\usepackage{natbib}

\title{My Document}
\author{Your Name}
\date{\today}

\begin{document}

\maketitle

\section{Introduction}
This is a sample document demonstrating various \LaTeX{} features including mathematical equations, tables, and formatting. You can cite papers using the citation search feature, and they will be automatically added to your \texttt{cite.bib} file.

\section{Mathematical Expressions}
Here are some example equations:

\subsection{Quadratic Formula}
The quadratic formula is given by:
\[x = \frac{-b \pm \sqrt{b^2 - 4ac}}{2a}\]

\subsection{Matrix Example}
A sample matrix:
\[
\begin{pmatrix}
a_{11} & a_{12} & a_{13} \\
a_{21} & a_{22} & a_{23} \\
a_{31} & a_{32} & a_{33}
\end{pmatrix}
\]

\section{Tables}
Below is a sample table showing data:

\begin{table}[h]
\centering
\begin{tabular}{@{}lrrr@{}}
\toprule
\textbf{Category} & \textbf{Value 1} & \textbf{Value 2} & \textbf{Total} \\
\midrule
Item A & 10.5 & 20.3 & 30.8 \\
Item B & 15.2 & 25.1 & 40.3 \\
Item C & 22.1 & 12.7 & 34.8 \\
\bottomrule
\end{tabular}
\caption{Sample Data Table}
\label{tab:sample}
\end{table}

\section{Lists}
\subsection{Bullet Points}
\begin{itemize}
  \item First item
  \item Second item
  \item Third item with sub-items:
    \begin{itemize}
      \item Sub-item 1
      \item Sub-item 2
    \end{itemize}
\end{itemize}

\subsection{Numbered List}
\begin{enumerate}
  \item First step
  \item Second step
  \item Third step
\end{enumerate}

\section{Citations and References}
When you use the citation search feature, references will be automatically managed. For example, you might cite a paper like this: \cite{example2024}. All citations will be automatically added to your \texttt{cite.bib} file.

\section{Conclusion}
This template demonstrates basic \LaTeX{} features including mathematical formulas, tables, and lists. The bibliography system is now set up to work seamlessly with Bibby AI's citation management. Feel free to modify it for your needs!

\bibliographystyle{plain}
\bibliography{cite}

\end{document}\documentclass{article}
\usepackage[utf8]{inputenc}
\usepackage{booktabs}
\usepackage{amsmath}
\usepackage{natbib}

\title{My Document}
\author{Your Name}
\date{\today}

\begin{document}

\maketitle

\section{Introduction}
This is a sample document demonstrating various \LaTeX{} features including mathematical equations, tables, and formatting. You can cite papers using the citation search feature, and they will be automatically added to your \texttt{cite.bib} file.

\section{Mathematical Expressions}
Here are some example equations:

\subsection{Quadratic Formula}
The quadratic formula is given by:
\[x = \frac{-b \pm \sqrt{b^2 - 4ac}}{2a}\]

\subsection{Matrix Example}
A sample matrix:
\[
\begin{pmatrix}
a_{11} & a_{12} & a_{13} \\
a_{21} & a_{22} & a_{23} \\
a_{31} & a_{32} & a_{33}
\end{pmatrix}
\]

\section{Tables}
Below is a sample table showing data:

\begin{table}[h]
\centering
\begin{tabular}{@{}lrrr@{}}
\toprule
\textbf{Category} & \textbf{Value 1} & \textbf{Value 2} & \textbf{Total} \\
\midrule
Item A & 10.5 & 20.3 & 30.8 \\
Item B & 15.2 & 25.1 & 40.3 \\
Item C & 22.1 & 12.7 & 34.8 \\
\bottomrule
\end{tabular}
\caption{Sample Data Table}
\label{tab:sample}
\end{table}

\section{Lists}
\subsection{Bullet Points}
\begin{itemize}
  \item First item
  \item Second item
  \item Third item with sub-items:
    \begin{itemize}
      \item Sub-item 1
      \item Sub-item 2
    \end{itemize}
\end{itemize}

\subsection{Numbered List}
\begin{enumerate}
  \item First step
  \item Second step
  \item Third step
\end{enumerate}

\section{Citations and References}
When you use the citation search feature, references will be automatically managed. For example, you might cite a paper like this: \cite{example2024}. All citations will be automatically added to your \texttt{cite.bib} file.

\section{Conclusion}
This template demonstrates basic \LaTeX{} features including mathematical formulas, tables, and lists. The bibliography system is now set up to work seamlessly with Bibby AI's citation management. Feel free to modify it for your needs!

\bibliographystyle{plain}
\bibliography{cite}

\end{document}\documentclass{article}
\usepackage[utf8]{inputenc}
\usepackage{booktabs}
\usepackage{amsmath}
\usepackage{natbib}

\title{My Document}
\author{Your Name}
\date{\today}

\begin{document}

\maketitle

\section{Introduction}
This is a sample document demonstrating various \LaTeX{} features including mathematical equations, tables, and formatting. You can cite papers using the citation search feature, and they will be automatically added to your \texttt{cite.bib} file.

\section{Mathematical Expressions}
Here are some example equations:

\subsection{Quadratic Formula}
The quadratic formula is given by:
\[x = \frac{-b \pm \sqrt{b^2 - 4ac}}{2a}\]

\subsection{Matrix Example}
A sample matrix:
\[
\begin{pmatrix}
a_{11} & a_{12} & a_{13} \\
a_{21} & a_{22} & a_{23} \\
a_{31} & a_{32} & a_{33}
\end{pmatrix}
\]

\section{Tables}
Below is a sample table showing data:

\begin{table}[h]
\centering
\begin{tabular}{@{}lrrr@{}}
\toprule
\textbf{Category} & \textbf{Value 1} & \textbf{Value 2} & \textbf{Total} \\
\midrule
Item A & 10.5 & 20.3 & 30.8 \\
Item B & 15.2 & 25.1 & 40.3 \\
Item C & 22.1 & 12.7 & 34.8 \\
\bottomrule
\end{tabular}
\caption{Sample Data Table}
\label{tab:sample}
\end{table}

\section{Lists}
\subsection{Bullet Points}
\begin{itemize}
  \item First item
  \item Second item
  \item Third item with sub-items:
    \begin{itemize}
      \item Sub-item 1
      \item Sub-item 2
    \end{itemize}
\end{itemize}

\subsection{Numbered List}
\begin{enumerate}
  \item First step
  \item Second step
  \item Third step
\end{enumerate}

\section{Citations and References}
When you use the citation search feature, references will be automatically managed. For example, you might cite a paper like this: \cite{example2024}. All citations will be automatically added to your \texttt{cite.bib} file.

\section{Conclusion}
This template demonstrates basic \LaTeX{} features including mathematical formulas, tables, and lists. The bibliography system is now set up to work seamlessly with Bibby AI's citation management. Feel free to modify it for your needs!

\bibliographystyle{plain}
\bibliography{cite}

\end{document}\documentclass{article}
\usepackage[utf8]{inputenc}
\usepackage{booktabs}
\usepackage{amsmath}
\usepackage{natbib}

\title{My Document}
\author{Your Name}
\date{\today}

\begin{document}

\maketitle

\section{Introduction}
This is a sample document demonstrating various \LaTeX{} features including mathematical equations, tables, and formatting. You can cite papers using the citation search feature, and they will be automatically added to your \texttt{cite.bib} file.

\section{Mathematical Expressions}
Here are some example equations:

\subsection{Quadratic Formula}
The quadratic formula is given by:
\[x = \frac{-b \pm \sqrt{b^2 - 4ac}}{2a}\]

\subsection{Matrix Example}
A sample matrix:
\[
\begin{pmatrix}
a_{11} & a_{12} & a_{13} \\
a_{21} & a_{22} & a_{23} \\
a_{31} & a_{32} & a_{33}
\end{pmatrix}
\]

\section{Tables}
Below is a sample table showing data:

\begin{table}[h]
\centering
\begin{tabular}{@{}lrrr@{}}
\toprule
\textbf{Category} & \textbf{Value 1} & \textbf{Value 2} & \textbf{Total} \\
\midrule
Item A & 10.5 & 20.3 & 30.8 \\
Item B & 15.2 & 25.1 & 40.3 \\
Item C & 22.1 & 12.7 & 34.8 \\
\bottomrule
\end{tabular}
\caption{Sample Data Table}
\label{tab:sample}
\end{table}

\section{Lists}
\subsection{Bullet Points}
\begin{itemize}
  \item First item
  \item Second item
  \item Third item with sub-items:
    \begin{itemize}
      \item Sub-item 1
      \item Sub-item 2
    \end{itemize}
\end{itemize}

\subsection{Numbered List}
\begin{enumerate}
  \item First step
  \item Second step
  \item Third step
\end{enumerate}

\section{Citations and References}
When you use the citation search feature, references will be automatically managed. For example, you might cite a paper like this: \cite{example2024}. All citations will be automatically added to your \texttt{cite.bib} file.

\section{Conclusion}
This template demonstrates basic \LaTeX{} features including mathematical formulas, tables, and lists. The bibliography system is now set up to work seamlessly with Bibby AI's citation management. Feel free to modify it for your needs!

\bibliographystyle{plain}
\bibliography{cite}

\end{document}\documentclass{article}
\usepackage[utf8]{inputenc}
\usepackage{booktabs}
\usepackage{amsmath}
\usepackage{natbib}

\title{My Document}
\author{Your Name}
\date{\today}

\begin{document}

\maketitle

\section{Introduction}
This is a sample document demonstrating various \LaTeX{} features including mathematical equations, tables, and formatting. You can cite papers using the citation search feature, and they will be automatically added to your \texttt{cite.bib} file.

\section{Mathematical Expressions}
Here are some example equations:

\subsection{Quadratic Formula}
The quadratic formula is given by:
\[x = \frac{-b \pm \sqrt{b^2 - 4ac}}{2a}\]

\subsection{Matrix Example}
A sample matrix:
\[
\begin{pmatrix}
a_{11} & a_{12} & a_{13} \\
a_{21} & a_{22} & a_{23} \\
a_{31} & a_{32} & a_{33}
\end{pmatrix}
\]

\section{Tables}
Below is a sample table showing data:

\begin{table}[h]
\centering
\begin{tabular}{@{}lrrr@{}}
\toprule
\textbf{Category} & \textbf{Value 1} & \textbf{Value 2} & \textbf{Total} \\
\midrule
Item A & 10.5 & 20.3 & 30.8 \\
Item B & 15.2 & 25.1 & 40.3 \\
Item C & 22.1 & 12.7 & 34.8 \\
\bottomrule
\end{tabular}
\caption{Sample Data Table}
\label{tab:sample}
\end{table}

\section{Lists}
\subsection{Bullet Points}
\begin{itemize}
  \item First item
  \item Second item
  \item Third item with sub-items:
    \begin{itemize}
      \item Sub-item 1
      \item Sub-item 2
    \end{itemize}
\end{itemize}

\subsection{Numbered List}
\begin{enumerate}
  \item First step
  \item Second step
  \item Third step
\end{enumerate}

\section{Citations and References}
When you use the citation search feature, references will be automatically managed. For example, you might cite a paper like this: \cite{example2024}. All citations will be automatically added to your \texttt{cite.bib} file.

\section{Conclusion}
This template demonstrates basic \LaTeX{} features including mathematical formulas, tables, and lists. The bibliography system is now set up to work seamlessly with Bibby AI's citation management. Feel free to modify it for your needs!

\bibliographystyle{plain}
\bibliography{cite}

\end{document}\documentclass{article}
\usepackage[utf8]{inputenc}
\usepackage{booktabs}
\usepackage{amsmath}
\usepackage{natbib}

\title{My Document}
\author{Your Name}
\date{\today}

\begin{document}

\maketitle

\section{Introduction}
This is a sample document demonstrating various \LaTeX{} features including mathematical equations, tables, and formatting. You can cite papers using the citation search feature, and they will be automatically added to your \texttt{cite.bib} file.

\section{Mathematical Expressions}
Here are some example equations:

\subsection{Quadratic Formula}
The quadratic formula is given by:
\[x = \frac{-b \pm \sqrt{b^2 - 4ac}}{2a}\]

\subsection{Matrix Example}
A sample matrix:
\[
\begin{pmatrix}
a_{11} & a_{12} & a_{13} \\
a_{21} & a_{22} & a_{23} \\
a_{31} & a_{32} & a_{33}
\end{pmatrix}
\]

\section{Tables}
Below is a sample table showing data:

\begin{table}[h]
\centering
\begin{tabular}{@{}lrrr@{}}
\toprule
\textbf{Category} & \textbf{Value 1} & \textbf{Value 2} & \textbf{Total} \\
\midrule
Item A & 10.5 & 20.3 & 30.8 \\
Item B & 15.2 & 25.1 & 40.3 \\
Item C & 22.1 & 12.7 & 34.8 \\
\bottomrule
\end{tabular}
\caption{Sample Data Table}
\label{tab:sample}
\end{table}

\section{Lists}
\subsection{Bullet Points}
\begin{itemize}
  \item First item
  \item Second item
  \item Third item with sub-items:
    \begin{itemize}
      \item Sub-item 1
      \item Sub-item 2
    \end{itemize}
\end{itemize}

\subsection{Numbered List}
\begin{enumerate}
  \item First step
  \item Second step
  \item Third step
\end{enumerate}

\section{Citations and References}
When you use the citation search feature, references will be automatically managed. For example, you might cite a paper like this: \cite{example2024}. All citations will be automatically added to your \texttt{cite.bib} file.

\section{Conclusion}
This template demonstrates basic \LaTeX{} features including mathematical formulas, tables, and lists. The bibliography system is now set up to work seamlessly with Bibby AI's citation management. Feel free to modify it for your needs!

\bibliographystyle{plain}
\bibliography{cite}

\end{document}\documentclass{article}
\usepackage[utf8]{inputenc}
\usepackage{booktabs}
\usepackage{amsmath}
\usepackage{natbib}

\title{My Document}
\author{Your Name}
\date{\today}

\begin{document}

\maketitle

\section{Introduction}
This is a sample document demonstrating various \LaTeX{} features including mathematical equations, tables, and formatting. You can cite papers using the citation search feature, and they will be automatically added to your \texttt{cite.bib} file.

\section{Mathematical Expressions}
Here are some example equations:

\subsection{Quadratic Formula}
The quadratic formula is given by:
\[x = \frac{-b \pm \sqrt{b^2 - 4ac}}{2a}\]

\subsection{Matrix Example}
A sample matrix:
\[
\begin{pmatrix}
a_{11} & a_{12} & a_{13} \\
a_{21} & a_{22} & a_{23} \\
a_{31} & a_{32} & a_{33}
\end{pmatrix}
\]

\section{Tables}
Below is a sample table showing data:

\begin{table}[h]
\centering
\begin{tabular}{@{}lrrr@{}}
\toprule
\textbf{Category} & \textbf{Value 1} & \textbf{Value 2} & \textbf{Total} \\
\midrule
Item A & 10.5 & 20.3 & 30.8 \\
Item B & 15.2 & 25.1 & 40.3 \\
Item C & 22.1 & 12.7 & 34.8 \\
\bottomrule
\end{tabular}
\caption{Sample Data Table}
\label{tab:sample}
\end{table}

\section{Lists}
\subsection{Bullet Points}
\begin{itemize}
  \item First item
  \item Second item
  \item Third item with sub-items:
    \begin{itemize}
      \item Sub-item 1
      \item Sub-item 2
    \end{itemize}
\end{itemize}

\subsection{Numbered List}
\begin{enumerate}
  \item First step
  \item Second step
  \item Third step
\end{enumerate}

\section{Citations and References}
When you use the citation search feature, references will be automatically managed. For example, you might cite a paper like this: \cite{example2024}. All citations will be automatically added to your \texttt{cite.bib} file.

\section{Conclusion}
This template demonstrates basic \LaTeX{} features including mathematical formulas, tables, and lists. The bibliography system is now set up to work seamlessly with Bibby AI's citation management. Feel free to modify it for your needs!

\bibliographystyle{plain}
\bibliography{cite}

\end{document}\documentclass{article}
\usepackage[utf8]{inputenc}
\usepackage{booktabs}
\usepackage{amsmath}
\usepackage{natbib}

\title{My Document}
\author{Your Name}
\date{\today}

\begin{document}

\maketitle

\section{Introduction}
This is a sample document demonstrating various \LaTeX{} features including mathematical equations, tables, and formatting. You can cite papers using the citation search feature, and they will be automatically added to your \texttt{cite.bib} file.

\section{Mathematical Expressions}
Here are some example equations:

\subsection{Quadratic Formula}
The quadratic formula is given by:
\[x = \frac{-b \pm \sqrt{b^2 - 4ac}}{2a}\]

\subsection{Matrix Example}
A sample matrix:
\[
\begin{pmatrix}
a_{11} & a_{12} & a_{13} \\
a_{21} & a_{22} & a_{23} \\
a_{31} & a_{32} & a_{33}
\end{pmatrix}
\]

\section{Tables}
Below is a sample table showing data:

\begin{table}[h]
\centering
\begin{tabular}{@{}lrrr@{}}
\toprule
\textbf{Category} & \textbf{Value 1} & \textbf{Value 2} & \textbf{Total} \\
\midrule
Item A & 10.5 & 20.3 & 30.8 \\
Item B & 15.2 & 25.1 & 40.3 \\
Item C & 22.1 & 12.7 & 34.8 \\
\bottomrule
\end{tabular}
\caption{Sample Data Table}
\label{tab:sample}
\end{table}

\section{Lists}
\subsection{Bullet Points}
\begin{itemize}
  \item First item
  \item Second item
  \item Third item with sub-items:
    \begin{itemize}
      \item Sub-item 1
      \item Sub-item 2
    \end{itemize}
\end{itemize}

\subsection{Numbered List}
\begin{enumerate}
  \item First step
  \item Second step
  \item Third step
\end{enumerate}

\section{Citations and References}
When you use the citation search feature, references will be automatically managed. For example, you might cite a paper like this: \cite{example2024}. All citations will be automatically added to your \texttt{cite.bib} file.

\section{Conclusion}
This template demonstrates basic \LaTeX{} features including mathematical formulas, tables, and lists. The bibliography system is now set up to work seamlessly with Bibby AI's citation management. Feel free to modify it for your needs!

\bibliographystyle{plain}
\bibliography{cite}

\end{document}\documentclass{article}
\usepackage[utf8]{inputenc}
\usepackage{booktabs}
\usepackage{amsmath}
\usepackage{natbib}

\title{My Document}
\author{Your Name}
\date{\today}

\begin{document}

\maketitle

\section{Introduction}
This is a sample document demonstrating various \LaTeX{} features including mathematical equations, tables, and formatting. You can cite papers using the citation search feature, and they will be automatically added to your \texttt{cite.bib} file.

\section{Mathematical Expressions}
Here are some example equations:

\subsection{Quadratic Formula}
The quadratic formula is given by:
\[x = \frac{-b \pm \sqrt{b^2 - 4ac}}{2a}\]

\subsection{Matrix Example}
A sample matrix:
\[
\begin{pmatrix}
a_{11} & a_{12} & a_{13} \\
a_{21} & a_{22} & a_{23} \\
a_{31} & a_{32} & a_{33}
\end{pmatrix}
\]

\section{Tables}
Below is a sample table showing data:

\begin{table}[h]
\centering
\begin{tabular}{@{}lrrr@{}}
\toprule
\textbf{Category} & \textbf{Value 1} & \textbf{Value 2} & \textbf{Total} \\
\midrule
Item A & 10.5 & 20.3 & 30.8 \\
Item B & 15.2 & 25.1 & 40.3 \\
Item C & 22.1 & 12.7 & 34.8 \\
\bottomrule
\end{tabular}
\caption{Sample Data Table}
\label{tab:sample}
\end{table}

\section{Lists}
\subsection{Bullet Points}
\begin{itemize}
  \item First item
  \item Second item
  \item Third item with sub-items:
    \begin{itemize}
      \item Sub-item 1
      \item Sub-item 2
    \end{itemize}
\end{itemize}

\subsection{Numbered List}
\begin{enumerate}
  \item First step
  \item Second step
  \item Third step
\end{enumerate}

\section{Citations and References}
When you use the citation search feature, references will be automatically managed. For example, you might cite a paper like this: \cite{example2024}. All citations will be automatically added to your \texttt{cite.bib} file.

\section{Conclusion}
This template demonstrates basic \LaTeX{} features including mathematical formulas, tables, and lists. The bibliography system is now set up to work seamlessly with Bibby AI's citation management. Feel free to modify it for your needs!

\bibliographystyle{plain}
\bibliography{cite}

\end{document}\documentclass{article}
\usepackage[utf8]{inputenc}
\usepackage{booktabs}
\usepackage{amsmath}
\usepackage{natbib}

\title{My Document}
\author{Your Name}
\date{\today}

\begin{document}

\maketitle

\section{Introduction}
This is a sample document demonstrating various \LaTeX{} features including mathematical equations, tables, and formatting. You can cite papers using the citation search feature, and they will be automatically added to your \texttt{cite.bib} file.

\section{Mathematical Expressions}
Here are some example equations:

\subsection{Quadratic Formula}
The quadratic formula is given by:
\[x = \frac{-b \pm \sqrt{b^2 - 4ac}}{2a}\]

\subsection{Matrix Example}
A sample matrix:
\[
\begin{pmatrix}
a_{11} & a_{12} & a_{13} \\
a_{21} & a_{22} & a_{23} \\
a_{31} & a_{32} & a_{33}
\end{pmatrix}
\]

\section{Tables}
Below is a sample table showing data:

\begin{table}[h]
\centering
\begin{tabular}{@{}lrrr@{}}
\toprule
\textbf{Category} & \textbf{Value 1} & \textbf{Value 2} & \textbf{Total} \\
\midrule
Item A & 10.5 & 20.3 & 30.8 \\
Item B & 15.2 & 25.1 & 40.3 \\
Item C & 22.1 & 12.7 & 34.8 \\
\bottomrule
\end{tabular}
\caption{Sample Data Table}
\label{tab:sample}
\end{table}

\section{Lists}
\subsection{Bullet Points}
\begin{itemize}
  \item First item
  \item Second item
  \item Third item with sub-items:
    \begin{itemize}
      \item Sub-item 1
      \item Sub-item 2
    \end{itemize}
\end{itemize}

\subsection{Numbered List}
\begin{enumerate}
  \item First step
  \item Second step
  \item Third step
\end{enumerate}

\section{Citations and References}
When you use the citation search feature, references will be automatically managed. For example, you might cite a paper like this: \cite{example2024}. All citations will be automatically added to your \texttt{cite.bib} file.

\section{Conclusion}
This template demonstrates basic \LaTeX{} features including mathematical formulas, tables, and lists. The bibliography system is now set up to work seamlessly with Bibby AI's citation management. Feel free to modify it for your needs!

\bibliographystyle{plain}
\bibliography{cite}

\end{document}\documentclass{article}
\usepackage[utf8]{inputenc}
\usepackage{booktabs}
\usepackage{amsmath}
\usepackage{natbib}

\title{My Document}
\author{Your Name}
\date{\today}

\begin{document}

\maketitle

\section{Introduction}
This is a sample document demonstrating various \LaTeX{} features including mathematical equations, tables, and formatting. You can cite papers using the citation search feature, and they will be automatically added to your \texttt{cite.bib} file.

\section{Mathematical Expressions}
Here are some example equations:

\subsection{Quadratic Formula}
The quadratic formula is given by:
\[x = \frac{-b \pm \sqrt{b^2 - 4ac}}{2a}\]

\subsection{Matrix Example}
A sample matrix:
\[
\begin{pmatrix}
a_{11} & a_{12} & a_{13} \\
a_{21} & a_{22} & a_{23} \\
a_{31} & a_{32} & a_{33}
\end{pmatrix}
\]

\section{Tables}
Below is a sample table showing data:

\begin{table}[h]
\centering
\begin{tabular}{@{}lrrr@{}}
\toprule
\textbf{Category} & \textbf{Value 1} & \textbf{Value 2} & \textbf{Total} \\
\midrule
Item A & 10.5 & 20.3 & 30.8 \\
Item B & 15.2 & 25.1 & 40.3 \\
Item C & 22.1 & 12.7 & 34.8 \\
\bottomrule
\end{tabular}
\caption{Sample Data Table}
\label{tab:sample}
\end{table}

\section{Lists}
\subsection{Bullet Points}
\begin{itemize}
  \item First item
  \item Second item
  \item Third item with sub-items:
    \begin{itemize}
      \item Sub-item 1
      \item Sub-item 2
    \end{itemize}
\end{itemize}

\subsection{Numbered List}
\begin{enumerate}
  \item First step
  \item Second step
  \item Third step
\end{enumerate}

\section{Citations and References}
When you use the citation search feature, references will be automatically managed. For example, you might cite a paper like this: \cite{example2024}. All citations will be automatically added to your \texttt{cite.bib} file.

\section{Conclusion}
This template demonstrates basic \LaTeX{} features including mathematical formulas, tables, and lists. The bibliography system is now set up to work seamlessly with Bibby AI's citation management. Feel free to modify it for your needs!

\bibliographystyle{plain}
\bibliography{cite}

\end{document}\documentclass{article}
\usepackage[utf8]{inputenc}
\usepackage{booktabs}
\usepackage{amsmath}
\usepackage{natbib}

\title{My Document}
\author{Your Name}
\date{\today}

\begin{document}

\maketitle

\section{Introduction}
This is a sample document demonstrating various \LaTeX{} features including mathematical equations, tables, and formatting. You can cite papers using the citation search feature, and they will be automatically added to your \texttt{cite.bib} file.

\section{Mathematical Expressions}
Here are some example equations:

\subsection{Quadratic Formula}
The quadratic formula is given by:
\[x = \frac{-b \pm \sqrt{b^2 - 4ac}}{2a}\]

\subsection{Matrix Example}
A sample matrix:
\[
\begin{pmatrix}
a_{11} & a_{12} & a_{13} \\
a_{21} & a_{22} & a_{23} \\
a_{31} & a_{32} & a_{33}
\end{pmatrix}
\]

\section{Tables}
Below is a sample table showing data:

\begin{table}[h]
\centering
\begin{tabular}{@{}lrrr@{}}
\toprule
\textbf{Category} & \textbf{Value 1} & \textbf{Value 2} & \textbf{Total} \\
\midrule
Item A & 10.5 & 20.3 & 30.8 \\
Item B & 15.2 & 25.1 & 40.3 \\
Item C & 22.1 & 12.7 & 34.8 \\
\bottomrule
\end{tabular}
\caption{Sample Data Table}
\label{tab:sample}
\end{table}

\section{Lists}
\subsection{Bullet Points}
\begin{itemize}
  \item First item
  \item Second item
  \item Third item with sub-items:
    \begin{itemize}
      \item Sub-item 1
      \item Sub-item 2
    \end{itemize}
\end{itemize}

\subsection{Numbered List}
\begin{enumerate}
  \item First step
  \item Second step
  \item Third step
\end{enumerate}

\section{Citations and References}
When you use the citation search feature, references will be automatically managed. For example, you might cite a paper like this: \cite{example2024}. All citations will be automatically added to your \texttt{cite.bib} file.

\section{Conclusion}
This template demonstrates basic \LaTeX{} features including mathematical formulas, tables, and lists. The bibliography system is now set up to work seamlessly with Bibby AI's citation management. Feel free to modify it for your needs!

\bibliographystyle{plain}
\bibliography{cite}

\end{document}\documentclass{article}
\usepackage[utf8]{inputenc}
\usepackage{booktabs}
\usepackage{amsmath}
\usepackage{natbib}

\title{My Document}
\author{Your Name}
\date{\today}

\begin{document}

\maketitle

\section{Introduction}
This is a sample document demonstrating various \LaTeX{} features including mathematical equations, tables, and formatting. You can cite papers using the citation search feature, and they will be automatically added to your \texttt{cite.bib} file.

\section{Mathematical Expressions}
Here are some example equations:

\subsection{Quadratic Formula}
The quadratic formula is given by:
\[x = \frac{-b \pm \sqrt{b^2 - 4ac}}{2a}\]

\subsection{Matrix Example}
A sample matrix:
\[
\begin{pmatrix}
a_{11} & a_{12} & a_{13} \\
a_{21} & a_{22} & a_{23} \\
a_{31} & a_{32} & a_{33}
\end{pmatrix}
\]

\section{Tables}
Below is a sample table showing data:

\begin{table}[h]
\centering
\begin{tabular}{@{}lrrr@{}}
\toprule
\textbf{Category} & \textbf{Value 1} & \textbf{Value 2} & \textbf{Total} \\
\midrule
Item A & 10.5 & 20.3 & 30.8 \\
Item B & 15.2 & 25.1 & 40.3 \\
Item C & 22.1 & 12.7 & 34.8 \\
\bottomrule
\end{tabular}
\caption{Sample Data Table}
\label{tab:sample}
\end{table}

\section{Lists}
\subsection{Bullet Points}
\begin{itemize}
  \item First item
  \item Second item
  \item Third item with sub-items:
    \begin{itemize}
      \item Sub-item 1
      \item Sub-item 2
    \end{itemize}
\end{itemize}

\subsection{Numbered List}
\begin{enumerate}
  \item First step
  \item Second step
  \item Third step
\end{enumerate}

\section{Citations and References}
When you use the citation search feature, references will be automatically managed. For example, you might cite a paper like this: \cite{example2024}. All citations will be automatically added to your \texttt{cite.bib} file.

\section{Conclusion}
This template demonstrates basic \LaTeX{} features including mathematical formulas, tables, and lists. The bibliography system is now set up to work seamlessly with Bibby AI's citation management. Feel free to modify it for your needs!

\bibliographystyle{plain}
\bibliography{cite}

\end{document}\documentclass{article}
\usepackage[utf8]{inputenc}
\usepackage{booktabs}
\usepackage{amsmath}
\usepackage{natbib}

\title{My Document}
\author{Your Name}
\date{\today}

\begin{document}

\maketitle

\section{Introduction}
This is a sample document demonstrating various \LaTeX{} features including mathematical equations, tables, and formatting. You can cite papers using the citation search feature, and they will be automatically added to your \texttt{cite.bib} file.

\section{Mathematical Expressions}
Here are some example equations:

\subsection{Quadratic Formula}
The quadratic formula is given by:
\[x = \frac{-b \pm \sqrt{b^2 - 4ac}}{2a}\]

\subsection{Matrix Example}
A sample matrix:
\[
\begin{pmatrix}
a_{11} & a_{12} & a_{13} \\
a_{21} & a_{22} & a_{23} \\
a_{31} & a_{32} & a_{33}
\end{pmatrix}
\]

\section{Tables}
Below is a sample table showing data:

\begin{table}[h]
\centering
\begin{tabular}{@{}lrrr@{}}
\toprule
\textbf{Category} & \textbf{Value 1} & \textbf{Value 2} & \textbf{Total} \\
\midrule
Item A & 10.5 & 20.3 & 30.8 \\
Item B & 15.2 & 25.1 & 40.3 \\
Item C & 22.1 & 12.7 & 34.8 \\
\bottomrule
\end{tabular}
\caption{Sample Data Table}
\label{tab:sample}
\end{table}

\section{Lists}
\subsection{Bullet Points}
\begin{itemize}
  \item First item
  \item Second item
  \item Third item with sub-items:
    \begin{itemize}
      \item Sub-item 1
      \item Sub-item 2
    \end{itemize}
\end{itemize}

\subsection{Numbered List}
\begin{enumerate}
  \item First step
  \item Second step
  \item Third step
\end{enumerate}

\section{Citations and References}
When you use the citation search feature, references will be automatically managed. For example, you might cite a paper like this: \cite{example2024}. All citations will be automatically added to your \texttt{cite.bib} file.

\section{Conclusion}
This template demonstrates basic \LaTeX{} features including mathematical formulas, tables, and lists. The bibliography system is now set up to work seamlessly with Bibby AI's citation management. Feel free to modify it for your needs!

\bibliographystyle{plain}
\bibliography{cite}

\end{document}\documentclass{article}
\usepackage[utf8]{inputenc}
\usepackage{booktabs}
\usepackage{amsmath}
\usepackage{natbib}

\title{My Document}
\author{Your Name}
\date{\today}

\begin{document}

\maketitle

\section{Introduction}
This is a sample document demonstrating various \LaTeX{} features including mathematical equations, tables, and formatting. You can cite papers using the citation search feature, and they will be automatically added to your \texttt{cite.bib} file.

\section{Mathematical Expressions}
Here are some example equations:

\subsection{Quadratic Formula}
The quadratic formula is given by:
\[x = \frac{-b \pm \sqrt{b^2 - 4ac}}{2a}\]

\subsection{Matrix Example}
A sample matrix:
\[
\begin{pmatrix}
a_{11} & a_{12} & a_{13} \\
a_{21} & a_{22} & a_{23} \\
a_{31} & a_{32} & a_{33}
\end{pmatrix}
\]

\section{Tables}
Below is a sample table showing data:

\begin{table}[h]
\centering
\begin{tabular}{@{}lrrr@{}}
\toprule
\textbf{Category} & \textbf{Value 1} & \textbf{Value 2} & \textbf{Total} \\
\midrule
Item A & 10.5 & 20.3 & 30.8 \\
Item B & 15.2 & 25.1 & 40.3 \\
Item C & 22.1 & 12.7 & 34.8 \\
\bottomrule
\end{tabular}
\caption{Sample Data Table}
\label{tab:sample}
\end{table}

\section{Lists}
\subsection{Bullet Points}
\begin{itemize}
  \item First item
  \item Second item
  \item Third item with sub-items:
    \begin{itemize}
      \item Sub-item 1
      \item Sub-item 2
    \end{itemize}
\end{itemize}

\subsection{Numbered List}
\begin{enumerate}
  \item First step
  \item Second step
  \item Third step
\end{enumerate}

\section{Citations and References}
When you use the citation search feature, references will be automatically managed. For example, you might cite a paper like this: \cite{example2024}. All citations will be automatically added to your \texttt{cite.bib} file.

\section{Conclusion}
This template demonstrates basic \LaTeX{} features including mathematical formulas, tables, and lists. The bibliography system is now set up to work seamlessly with Bibby AI's citation management. Feel free to modify it for your needs!

\bibliographystyle{plain}
\bibliography{cite}

\end{document}\documentclass{article}
\usepackage[utf8]{inputenc}
\usepackage{booktabs}
\usepackage{amsmath}
\usepackage{natbib}

\title{My Document}
\author{Your Name}
\date{\today}

\begin{document}

\maketitle

\section{Introduction}
This is a sample document demonstrating various \LaTeX{} features including mathematical equations, tables, and formatting. You can cite papers using the citation search feature, and they will be automatically added to your \texttt{cite.bib} file.

\section{Mathematical Expressions}
Here are some example equations:

\subsection{Quadratic Formula}
The quadratic formula is given by:
\[x = \frac{-b \pm \sqrt{b^2 - 4ac}}{2a}\]

\subsection{Matrix Example}
A sample matrix:
\[
\begin{pmatrix}
a_{11} & a_{12} & a_{13} \\
a_{21} & a_{22} & a_{23} \\
a_{31} & a_{32} & a_{33}
\end{pmatrix}
\]

\section{Tables}
Below is a sample table showing data:

\begin{table}[h]
\centering
\begin{tabular}{@{}lrrr@{}}
\toprule
\textbf{Category} & \textbf{Value 1} & \textbf{Value 2} & \textbf{Total} \\
\midrule
Item A & 10.5 & 20.3 & 30.8 \\
Item B & 15.2 & 25.1 & 40.3 \\
Item C & 22.1 & 12.7 & 34.8 \\
\bottomrule
\end{tabular}
\caption{Sample Data Table}
\label{tab:sample}
\end{table}

\section{Lists}
\subsection{Bullet Points}
\begin{itemize}
  \item First item
  \item Second item
  \item Third item with sub-items:
    \begin{itemize}
      \item Sub-item 1
      \item Sub-item 2
    \end{itemize}
\end{itemize}

\subsection{Numbered List}
\begin{enumerate}
  \item First step
  \item Second step
  \item Third step
\end{enumerate}

\section{Citations and References}
When you use the citation search feature, references will be automatically managed. For example, you might cite a paper like this: \cite{example2024}. All citations will be automatically added to your \texttt{cite.bib} file.

\section{Conclusion}
This template demonstrates basic \LaTeX{} features including mathematical formulas, tables, and lists. The bibliography system is now set up to work seamlessly with Bibby AI's citation management. Feel free to modify it for your needs!

\bibliographystyle{plain}
\bibliography{cite}

\end{document}\documentclass{article}
\usepackage[utf8]{inputenc}
\usepackage{booktabs}
\usepackage{amsmath}
\usepackage{natbib}

\title{My Document}
\author{Your Name}
\date{\today}

\begin{document}

\maketitle

\section{Introduction}
This is a sample document demonstrating various \LaTeX{} features including mathematical equations, tables, and formatting. You can cite papers using the citation search feature, and they will be automatically added to your \texttt{cite.bib} file.

\section{Mathematical Expressions}
Here are some example equations:

\subsection{Quadratic Formula}
The quadratic formula is given by:
\[x = \frac{-b \pm \sqrt{b^2 - 4ac}}{2a}\]

\subsection{Matrix Example}
A sample matrix:
\[
\begin{pmatrix}
a_{11} & a_{12} & a_{13} \\
a_{21} & a_{22} & a_{23} \\
a_{31} & a_{32} & a_{33}
\end{pmatrix}
\]

\section{Tables}
Below is a sample table showing data:

\begin{table}[h]
\centering
\begin{tabular}{@{}lrrr@{}}
\toprule
\textbf{Category} & \textbf{Value 1} & \textbf{Value 2} & \textbf{Total} \\
\midrule
Item A & 10.5 & 20.3 & 30.8 \\
Item B & 15.2 & 25.1 & 40.3 \\
Item C & 22.1 & 12.7 & 34.8 \\
\bottomrule
\end{tabular}
\caption{Sample Data Table}
\label{tab:sample}
\end{table}

\section{Lists}
\subsection{Bullet Points}
\begin{itemize}
  \item First item
  \item Second item
  \item Third item with sub-items:
    \begin{itemize}
      \item Sub-item 1
      \item Sub-item 2
    \end{itemize}
\end{itemize}

\subsection{Numbered List}
\begin{enumerate}
  \item First step
  \item Second step
  \item Third step
\end{enumerate}

\section{Citations and References}
When you use the citation search feature, references will be automatically managed. For example, you might cite a paper like this: \cite{example2024}. All citations will be automatically added to your \texttt{cite.bib} file.

\section{Conclusion}
This template demonstrates basic \LaTeX{} features including mathematical formulas, tables, and lists. The bibliography system is now set up to work seamlessly with Bibby AI's citation management. Feel free to modify it for your needs!

\bibliographystyle{plain}
\bibliography{cite}

\end{document}\documentclass{article}
\usepackage[utf8]{inputenc}
\usepackage{booktabs}
\usepackage{amsmath}
\usepackage{natbib}

\title{My Document}
\author{Your Name}
\date{\today}

\begin{document}

\maketitle

\section{Introduction}
This is a sample document demonstrating various \LaTeX{} features including mathematical equations, tables, and formatting. You can cite papers using the citation search feature, and they will be automatically added to your \texttt{cite.bib} file.

\section{Mathematical Expressions}
Here are some example equations:

\subsection{Quadratic Formula}
The quadratic formula is given by:
\[x = \frac{-b \pm \sqrt{b^2 - 4ac}}{2a}\]

\subsection{Matrix Example}
A sample matrix:
\[
\begin{pmatrix}
a_{11} & a_{12} & a_{13} \\
a_{21} & a_{22} & a_{23} \\
a_{31} & a_{32} & a_{33}
\end{pmatrix}
\]

\section{Tables}
Below is a sample table showing data:

\begin{table}[h]
\centering
\begin{tabular}{@{}lrrr@{}}
\toprule
\textbf{Category} & \textbf{Value 1} & \textbf{Value 2} & \textbf{Total} \\
\midrule
Item A & 10.5 & 20.3 & 30.8 \\
Item B & 15.2 & 25.1 & 40.3 \\
Item C & 22.1 & 12.7 & 34.8 \\
\bottomrule
\end{tabular}
\caption{Sample Data Table}
\label{tab:sample}
\end{table}

\section{Lists}
\subsection{Bullet Points}
\begin{itemize}
  \item First item
  \item Second item
  \item Third item with sub-items:
    \begin{itemize}
      \item Sub-item 1
      \item Sub-item 2
    \end{itemize}
\end{itemize}

\subsection{Numbered List}
\begin{enumerate}
  \item First step
  \item Second step
  \item Third step
\end{enumerate}

\section{Citations and References}
When you use the citation search feature, references will be automatically managed. For example, you might cite a paper like this: \cite{example2024}. All citations will be automatically added to your \texttt{cite.bib} file.

\section{Conclusion}
This template demonstrates basic \LaTeX{} features including mathematical formulas, tables, and lists. The bibliography system is now set up to work seamlessly with Bibby AI's citation management. Feel free to modify it for your needs!

\bibliographystyle{plain}
\bibliography{cite}

\end{document}\documentclass{article}
\usepackage[utf8]{inputenc}
\usepackage{booktabs}
\usepackage{amsmath}
\usepackage{natbib}

\title{My Document}
\author{Your Name}
\date{\today}

\begin{document}

\maketitle

\section{Introduction}
This is a sample document demonstrating various \LaTeX{} features including mathematical equations, tables, and formatting. You can cite papers using the citation search feature, and they will be automatically added to your \texttt{cite.bib} file.

\section{Mathematical Expressions}
Here are some example equations:

\subsection{Quadratic Formula}
The quadratic formula is given by:
\[x = \frac{-b \pm \sqrt{b^2 - 4ac}}{2a}\]

\subsection{Matrix Example}
A sample matrix:
\[
\begin{pmatrix}
a_{11} & a_{12} & a_{13} \\
a_{21} & a_{22} & a_{23} \\
a_{31} & a_{32} & a_{33}
\end{pmatrix}
\]

\section{Tables}
Below is a sample table showing data:

\begin{table}[h]
\centering
\begin{tabular}{@{}lrrr@{}}
\toprule
\textbf{Category} & \textbf{Value 1} & \textbf{Value 2} & \textbf{Total} \\
\midrule
Item A & 10.5 & 20.3 & 30.8 \\
Item B & 15.2 & 25.1 & 40.3 \\
Item C & 22.1 & 12.7 & 34.8 \\
\bottomrule
\end{tabular}
\caption{Sample Data Table}
\label{tab:sample}
\end{table}

\section{Lists}
\subsection{Bullet Points}
\begin{itemize}
  \item First item
  \item Second item
  \item Third item with sub-items:
    \begin{itemize}
      \item Sub-item 1
      \item Sub-item 2
    \end{itemize}
\end{itemize}

\subsection{Numbered List}
\begin{enumerate}
  \item First step
  \item Second step
  \item Third step
\end{enumerate}

\section{Citations and References}
When you use the citation search feature, references will be automatically managed. For example, you might cite a paper like this: \cite{example2024}. All citations will be automatically added to your \texttt{cite.bib} file.

\section{Conclusion}
This template demonstrates basic \LaTeX{} features including mathematical formulas, tables, and lists. The bibliography system is now set up to work seamlessly with Bibby AI's citation management. Feel free to modify it for your needs!

\bibliographystyle{plain}
\bibliography{cite}

\end{document}\documentclass{article}
\usepackage[utf8]{inputenc}
\usepackage{booktabs}
\usepackage{amsmath}
\usepackage{natbib}

\title{My Document}
\author{Your Name}
\date{\today}

\begin{document}

\maketitle

\section{Introduction}
This is a sample document demonstrating various \LaTeX{} features including mathematical equations, tables, and formatting. You can cite papers using the citation search feature, and they will be automatically added to your \texttt{cite.bib} file.

\section{Mathematical Expressions}
Here are some example equations:

\subsection{Quadratic Formula}
The quadratic formula is given by:
\[x = \frac{-b \pm \sqrt{b^2 - 4ac}}{2a}\]

\subsection{Matrix Example}
A sample matrix:
\[
\begin{pmatrix}
a_{11} & a_{12} & a_{13} \\
a_{21} & a_{22} & a_{23} \\
a_{31} & a_{32} & a_{33}
\end{pmatrix}
\]

\section{Tables}
Below is a sample table showing data:

\begin{table}[h]
\centering
\begin{tabular}{@{}lrrr@{}}
\toprule
\textbf{Category} & \textbf{Value 1} & \textbf{Value 2} & \textbf{Total} \\
\midrule
Item A & 10.5 & 20.3 & 30.8 \\
Item B & 15.2 & 25.1 & 40.3 \\
Item C & 22.1 & 12.7 & 34.8 \\
\bottomrule
\end{tabular}
\caption{Sample Data Table}
\label{tab:sample}
\end{table}

\section{Lists}
\subsection{Bullet Points}
\begin{itemize}
  \item First item
  \item Second item
  \item Third item with sub-items:
    \begin{itemize}
      \item Sub-item 1
      \item Sub-item 2
    \end{itemize}
\end{itemize}

\subsection{Numbered List}
\begin{enumerate}
  \item First step
  \item Second step
  \item Third step
\end{enumerate}

\section{Citations and References}
When you use the citation search feature, references will be automatically managed. For example, you might cite a paper like this: \cite{example2024}. All citations will be automatically added to your \texttt{cite.bib} file.

\section{Conclusion}
This template demonstrates basic \LaTeX{} features including mathematical formulas, tables, and lists. The bibliography system is now set up to work seamlessly with Bibby AI's citation management. Feel free to modify it for your needs!

\bibliographystyle{plain}
\bibliography{cite}

\end{document}\documentclass{article}
\usepackage[utf8]{inputenc}
\usepackage{booktabs}
\usepackage{amsmath}
\usepackage{natbib}

\title{My Document}
\author{Your Name}
\date{\today}

\begin{document}

\maketitle

\section{Introduction}
This is a sample document demonstrating various \LaTeX{} features including mathematical equations, tables, and formatting. You can cite papers using the citation search feature, and they will be automatically added to your \texttt{cite.bib} file.

\section{Mathematical Expressions}
Here are some example equations:

\subsection{Quadratic Formula}
The quadratic formula is given by:
\[x = \frac{-b \pm \sqrt{b^2 - 4ac}}{2a}\]

\subsection{Matrix Example}
A sample matrix:
\[
\begin{pmatrix}
a_{11} & a_{12} & a_{13} \\
a_{21} & a_{22} & a_{23} \\
a_{31} & a_{32} & a_{33}
\end{pmatrix}
\]

\section{Tables}
Below is a sample table showing data:

\begin{table}[h]
\centering
\begin{tabular}{@{}lrrr@{}}
\toprule
\textbf{Category} & \textbf{Value 1} & \textbf{Value 2} & \textbf{Total} \\
\midrule
Item A & 10.5 & 20.3 & 30.8 \\
Item B & 15.2 & 25.1 & 40.3 \\
Item C & 22.1 & 12.7 & 34.8 \\
\bottomrule
\end{tabular}
\caption{Sample Data Table}
\label{tab:sample}
\end{table}

\section{Lists}
\subsection{Bullet Points}
\begin{itemize}
  \item First item
  \item Second item
  \item Third item with sub-items:
    \begin{itemize}
      \item Sub-item 1
      \item Sub-item 2
    \end{itemize}
\end{itemize}

\subsection{Numbered List}
\begin{enumerate}
  \item First step
  \item Second step
  \item Third step
\end{enumerate}

\section{Citations and References}
When you use the citation search feature, references will be automatically managed. For example, you might cite a paper like this: \cite{example2024}. All citations will be automatically added to your \texttt{cite.bib} file.

\section{Conclusion}
This template demonstrates basic \LaTeX{} features including mathematical formulas, tables, and lists. The bibliography system is now set up to work seamlessly with Bibby AI's citation management. Feel free to modify it for your needs!

\bibliographystyle{plain}
\bibliography{cite}

\end{document}\documentclass{article}
\usepackage[utf8]{inputenc}
\usepackage{booktabs}
\usepackage{amsmath}
\usepackage{natbib}

\title{My Document}
\author{Your Name}
\date{\today}

\begin{document}

\maketitle

\section{Introduction}
This is a sample document demonstrating various \LaTeX{} features including mathematical equations, tables, and formatting. You can cite papers using the citation search feature, and they will be automatically added to your \texttt{cite.bib} file.

\section{Mathematical Expressions}
Here are some example equations:

\subsection{Quadratic Formula}
The quadratic formula is given by:
\[x = \frac{-b \pm \sqrt{b^2 - 4ac}}{2a}\]

\subsection{Matrix Example}
A sample matrix:
\[
\begin{pmatrix}
a_{11} & a_{12} & a_{13} \\
a_{21} & a_{22} & a_{23} \\
a_{31} & a_{32} & a_{33}
\end{pmatrix}
\]

\section{Tables}
Below is a sample table showing data:

\begin{table}[h]
\centering
\begin{tabular}{@{}lrrr@{}}
\toprule
\textbf{Category} & \textbf{Value 1} & \textbf{Value 2} & \textbf{Total} \\
\midrule
Item A & 10.5 & 20.3 & 30.8 \\
Item B & 15.2 & 25.1 & 40.3 \\
Item C & 22.1 & 12.7 & 34.8 \\
\bottomrule
\end{tabular}
\caption{Sample Data Table}
\label{tab:sample}
\end{table}

\section{Lists}
\subsection{Bullet Points}
\begin{itemize}
  \item First item
  \item Second item
  \item Third item with sub-items:
    \begin{itemize}
      \item Sub-item 1
      \item Sub-item 2
    \end{itemize}
\end{itemize}

\subsection{Numbered List}
\begin{enumerate}
  \item First step
  \item Second step
  \item Third step
\end{enumerate}

\section{Citations and References}
When you use the citation search feature, references will be automatically managed. For example, you might cite a paper like this: \cite{example2024}. All citations will be automatically added to your \texttt{cite.bib} file.

\section{Conclusion}
This template demonstrates basic \LaTeX{} features including mathematical formulas, tables, and lists. The bibliography system is now set up to work seamlessly with Bibby AI's citation management. Feel free to modify it for your needs!

\bibliographystyle{plain}
\bibliography{cite}

\end{document}\documentclass{article}
\usepackage[utf8]{inputenc}
\usepackage{booktabs}
\usepackage{amsmath}
\usepackage{natbib}

\title{My Document}
\author{Your Name}
\date{\today}

\begin{document}

\maketitle

\section{Introduction}
This is a sample document demonstrating various \LaTeX{} features including mathematical equations, tables, and formatting. You can cite papers using the citation search feature, and they will be automatically added to your \texttt{cite.bib} file.

\section{Mathematical Expressions}
Here are some example equations:

\subsection{Quadratic Formula}
The quadratic formula is given by:
\[x = \frac{-b \pm \sqrt{b^2 - 4ac}}{2a}\]

\subsection{Matrix Example}
A sample matrix:
\[
\begin{pmatrix}
a_{11} & a_{12} & a_{13} \\
a_{21} & a_{22} & a_{23} \\
a_{31} & a_{32} & a_{33}
\end{pmatrix}
\]

\section{Tables}
Below is a sample table showing data:

\begin{table}[h]
\centering
\begin{tabular}{@{}lrrr@{}}
\toprule
\textbf{Category} & \textbf{Value 1} & \textbf{Value 2} & \textbf{Total} \\
\midrule
Item A & 10.5 & 20.3 & 30.8 \\
Item B & 15.2 & 25.1 & 40.3 \\
Item C & 22.1 & 12.7 & 34.8 \\
\bottomrule
\end{tabular}
\caption{Sample Data Table}
\label{tab:sample}
\end{table}

\section{Lists}
\subsection{Bullet Points}
\begin{itemize}
  \item First item
  \item Second item
  \item Third item with sub-items:
    \begin{itemize}
      \item Sub-item 1
      \item Sub-item 2
    \end{itemize}
\end{itemize}

\subsection{Numbered List}
\begin{enumerate}
  \item First step
  \item Second step
  \item Third step
\end{enumerate}

\section{Citations and References}
When you use the citation search feature, references will be automatically managed. For example, you might cite a paper like this: \cite{example2024}. All citations will be automatically added to your \texttt{cite.bib} file.

\section{Conclusion}
This template demonstrates basic \LaTeX{} features including mathematical formulas, tables, and lists. The bibliography system is now set up to work seamlessly with Bibby AI's citation management. Feel free to modify it for your needs!

\bibliographystyle{plain}
\bibliography{cite}

\end{document}\documentclass{article}
\usepackage[utf8]{inputenc}
\usepackage{booktabs}
\usepackage{amsmath}
\usepackage{natbib}

\title{My Document}
\author{Your Name}
\date{\today}

\begin{document}

\maketitle

\section{Introduction}
This is a sample document demonstrating various \LaTeX{} features including mathematical equations, tables, and formatting. You can cite papers using the citation search feature, and they will be automatically added to your \texttt{cite.bib} file.

\section{Mathematical Expressions}
Here are some example equations:

\subsection{Quadratic Formula}
The quadratic formula is given by:
\[x = \frac{-b \pm \sqrt{b^2 - 4ac}}{2a}\]

\subsection{Matrix Example}
A sample matrix:
\[
\begin{pmatrix}
a_{11} & a_{12} & a_{13} \\
a_{21} & a_{22} & a_{23} \\
a_{31} & a_{32} & a_{33}
\end{pmatrix}
\]

\section{Tables}
Below is a sample table showing data:

\begin{table}[h]
\centering
\begin{tabular}{@{}lrrr@{}}
\toprule
\textbf{Category} & \textbf{Value 1} & \textbf{Value 2} & \textbf{Total} \\
\midrule
Item A & 10.5 & 20.3 & 30.8 \\
Item B & 15.2 & 25.1 & 40.3 \\
Item C & 22.1 & 12.7 & 34.8 \\
\bottomrule
\end{tabular}
\caption{Sample Data Table}
\label{tab:sample}
\end{table}

\section{Lists}
\subsection{Bullet Points}
\begin{itemize}
  \item First item
  \item Second item
  \item Third item with sub-items:
    \begin{itemize}
      \item Sub-item 1
      \item Sub-item 2
    \end{itemize}
\end{itemize}

\subsection{Numbered List}
\begin{enumerate}
  \item First step
  \item Second step
  \item Third step
\end{enumerate}

\section{Citations and References}
When you use the citation search feature, references will be automatically managed. For example, you might cite a paper like this: \cite{example2024}. All citations will be automatically added to your \texttt{cite.bib} file.

\section{Conclusion}
This template demonstrates basic \LaTeX{} features including mathematical formulas, tables, and lists. The bibliography system is now set up to work seamlessly with Bibby AI's citation management. Feel free to modify it for your needs!

\bibliographystyle{plain}
\bibliography{cite}

\end{document}\documentclass{article}
\usepackage[utf8]{inputenc}
\usepackage{booktabs}
\usepackage{amsmath}
\usepackage{natbib}

\title{My Document}
\author{Your Name}
\date{\today}

\begin{document}

\maketitle

\section{Introduction}
This is a sample document demonstrating various \LaTeX{} features including mathematical equations, tables, and formatting. You can cite papers using the citation search feature, and they will be automatically added to your \texttt{cite.bib} file.

\section{Mathematical Expressions}
Here are some example equations:

\subsection{Quadratic Formula}
The quadratic formula is given by:
\[x = \frac{-b \pm \sqrt{b^2 - 4ac}}{2a}\]

\subsection{Matrix Example}
A sample matrix:
\[
\begin{pmatrix}
a_{11} & a_{12} & a_{13} \\
a_{21} & a_{22} & a_{23} \\
a_{31} & a_{32} & a_{33}
\end{pmatrix}
\]

\section{Tables}
Below is a sample table showing data:

\begin{table}[h]
\centering
\begin{tabular}{@{}lrrr@{}}
\toprule
\textbf{Category} & \textbf{Value 1} & \textbf{Value 2} & \textbf{Total} \\
\midrule
Item A & 10.5 & 20.3 & 30.8 \\
Item B & 15.2 & 25.1 & 40.3 \\
Item C & 22.1 & 12.7 & 34.8 \\
\bottomrule
\end{tabular}
\caption{Sample Data Table}
\label{tab:sample}
\end{table}

\section{Lists}
\subsection{Bullet Points}
\begin{itemize}
  \item First item
  \item Second item
  \item Third item with sub-items:
    \begin{itemize}
      \item Sub-item 1
      \item Sub-item 2
    \end{itemize}
\end{itemize}

\subsection{Numbered List}
\begin{enumerate}
  \item First step
  \item Second step
  \item Third step
\end{enumerate}

\section{Citations and References}
When you use the citation search feature, references will be automatically managed. For example, you might cite a paper like this: \cite{example2024}. All citations will be automatically added to your \texttt{cite.bib} file.

\section{Conclusion}
This template demonstrates basic \LaTeX{} features including mathematical formulas, tables, and lists. The bibliography system is now set up to work seamlessly with Bibby AI's citation management. Feel free to modify it for your needs!

\bibliographystyle{plain}
\bibliography{cite}

\end{document}\documentclass{article}
\usepackage[utf8]{inputenc}
\usepackage{booktabs}
\usepackage{amsmath}
\usepackage{natbib}

\title{My Document}
\author{Your Name}
\date{\today}

\begin{document}

\maketitle

\section{Introduction}
This is a sample document demonstrating various \LaTeX{} features including mathematical equations, tables, and formatting. You can cite papers using the citation search feature, and they will be automatically added to your \texttt{cite.bib} file.

\section{Mathematical Expressions}
Here are some example equations:

\subsection{Quadratic Formula}
The quadratic formula is given by:
\[x = \frac{-b \pm \sqrt{b^2 - 4ac}}{2a}\]

\subsection{Matrix Example}
A sample matrix:
\[
\begin{pmatrix}
a_{11} & a_{12} & a_{13} \\
a_{21} & a_{22} & a_{23} \\
a_{31} & a_{32} & a_{33}
\end{pmatrix}
\]

\section{Tables}
Below is a sample table showing data:

\begin{table}[h]
\centering
\begin{tabular}{@{}lrrr@{}}
\toprule
\textbf{Category} & \textbf{Value 1} & \textbf{Value 2} & \textbf{Total} \\
\midrule
Item A & 10.5 & 20.3 & 30.8 \\
Item B & 15.2 & 25.1 & 40.3 \\
Item C & 22.1 & 12.7 & 34.8 \\
\bottomrule
\end{tabular}
\caption{Sample Data Table}
\label{tab:sample}
\end{table}

\section{Lists}
\subsection{Bullet Points}
\begin{itemize}
  \item First item
  \item Second item
  \item Third item with sub-items:
    \begin{itemize}
      \item Sub-item 1
      \item Sub-item 2
    \end{itemize}
\end{itemize}

\subsection{Numbered List}
\begin{enumerate}
  \item First step
  \item Second step
  \item Third step
\end{enumerate}

\section{Citations and References}
When you use the citation search feature, references will be automatically managed. For example, you might cite a paper like this: \cite{example2024}. All citations will be automatically added to your \texttt{cite.bib} file.

\section{Conclusion}
This template demonstrates basic \LaTeX{} features including mathematical formulas, tables, and lists. The bibliography system is now set up to work seamlessly with Bibby AI's citation management. Feel free to modify it for your needs!

\bibliographystyle{plain}
\bibliography{cite}

\end{document}\documentclass{article}
\usepackage[utf8]{inputenc}
\usepackage{booktabs}
\usepackage{amsmath}
\usepackage{natbib}

\title{My Document}
\author{Your Name}
\date{\today}

\begin{document}

\maketitle

\section{Introduction}
This is a sample document demonstrating various \LaTeX{} features including mathematical equations, tables, and formatting. You can cite papers using the citation search feature, and they will be automatically added to your \texttt{cite.bib} file.

\section{Mathematical Expressions}
Here are some example equations:

\subsection{Quadratic Formula}
The quadratic formula is given by:
\[x = \frac{-b \pm \sqrt{b^2 - 4ac}}{2a}\]

\subsection{Matrix Example}
A sample matrix:
\[
\begin{pmatrix}
a_{11} & a_{12} & a_{13} \\
a_{21} & a_{22} & a_{23} \\
a_{31} & a_{32} & a_{33}
\end{pmatrix}
\]

\section{Tables}
Below is a sample table showing data:

\begin{table}[h]
\centering
\begin{tabular}{@{}lrrr@{}}
\toprule
\textbf{Category} & \textbf{Value 1} & \textbf{Value 2} & \textbf{Total} \\
\midrule
Item A & 10.5 & 20.3 & 30.8 \\
Item B & 15.2 & 25.1 & 40.3 \\
Item C & 22.1 & 12.7 & 34.8 \\
\bottomrule
\end{tabular}
\caption{Sample Data Table}
\label{tab:sample}
\end{table}

\section{Lists}
\subsection{Bullet Points}
\begin{itemize}
  \item First item
  \item Second item
  \item Third item with sub-items:
    \begin{itemize}
      \item Sub-item 1
      \item Sub-item 2
    \end{itemize}
\end{itemize}

\subsection{Numbered List}
\begin{enumerate}
  \item First step
  \item Second step
  \item Third step
\end{enumerate}

\section{Citations and References}
When you use the citation search feature, references will be automatically managed. For example, you might cite a paper like this: \cite{example2024}. All citations will be automatically added to your \texttt{cite.bib} file.

\section{Conclusion}
This template demonstrates basic \LaTeX{} features including mathematical formulas, tables, and lists. The bibliography system is now set up to work seamlessly with Bibby AI's citation management. Feel free to modify it for your needs!

\bibliographystyle{plain}
\bibliography{cite}

\end{document}\documentclass{article}
\usepackage[utf8]{inputenc}
\usepackage{booktabs}
\usepackage{amsmath}
\usepackage{natbib}

\title{My Document}
\author{Your Name}
\date{\today}

\begin{document}

\maketitle

\section{Introduction}
This is a sample document demonstrating various \LaTeX{} features including mathematical equations, tables, and formatting. You can cite papers using the citation search feature, and they will be automatically added to your \texttt{cite.bib} file.

\section{Mathematical Expressions}
Here are some example equations:

\subsection{Quadratic Formula}
The quadratic formula is given by:
\[x = \frac{-b \pm \sqrt{b^2 - 4ac}}{2a}\]

\subsection{Matrix Example}
A sample matrix:
\[
\begin{pmatrix}
a_{11} & a_{12} & a_{13} \\
a_{21} & a_{22} & a_{23} \\
a_{31} & a_{32} & a_{33}
\end{pmatrix}
\]

\section{Tables}
Below is a sample table showing data:

\begin{table}[h]
\centering
\begin{tabular}{@{}lrrr@{}}
\toprule
\textbf{Category} & \textbf{Value 1} & \textbf{Value 2} & \textbf{Total} \\
\midrule
Item A & 10.5 & 20.3 & 30.8 \\
Item B & 15.2 & 25.1 & 40.3 \\
Item C & 22.1 & 12.7 & 34.8 \\
\bottomrule
\end{tabular}
\caption{Sample Data Table}
\label{tab:sample}
\end{table}

\section{Lists}
\subsection{Bullet Points}
\begin{itemize}
  \item First item
  \item Second item
  \item Third item with sub-items:
    \begin{itemize}
      \item Sub-item 1
      \item Sub-item 2
    \end{itemize}
\end{itemize}

\subsection{Numbered List}
\begin{enumerate}
  \item First step
  \item Second step
  \item Third step
\end{enumerate}

\section{Citations and References}
When you use the citation search feature, references will be automatically managed. For example, you might cite a paper like this: \cite{example2024}. All citations will be automatically added to your \texttt{cite.bib} file.

\section{Conclusion}
This template demonstrates basic \LaTeX{} features including mathematical formulas, tables, and lists. The bibliography system is now set up to work seamlessly with Bibby AI's citation management. Feel free to modify it for your needs!

\bibliographystyle{plain}
\bibliography{cite}

\end{document}\documentclass{article}
\usepackage[utf8]{inputenc}
\usepackage{booktabs}
\usepackage{amsmath}
\usepackage{natbib}

\title{My Document}
\author{Your Name}
\date{\today}

\begin{document}

\maketitle

\section{Introduction}
This is a sample document demonstrating various \LaTeX{} features including mathematical equations, tables, and formatting. You can cite papers using the citation search feature, and they will be automatically added to your \texttt{cite.bib} file.

\section{Mathematical Expressions}
Here are some example equations:

\subsection{Quadratic Formula}
The quadratic formula is given by:
\[x = \frac{-b \pm \sqrt{b^2 - 4ac}}{2a}\]

\subsection{Matrix Example}
A sample matrix:
\[
\begin{pmatrix}
a_{11} & a_{12} & a_{13} \\
a_{21} & a_{22} & a_{23} \\
a_{31} & a_{32} & a_{33}
\end{pmatrix}
\]

\section{Tables}
Below is a sample table showing data:

\begin{table}[h]
\centering
\begin{tabular}{@{}lrrr@{}}
\toprule
\textbf{Category} & \textbf{Value 1} & \textbf{Value 2} & \textbf{Total} \\
\midrule
Item A & 10.5 & 20.3 & 30.8 \\
Item B & 15.2 & 25.1 & 40.3 \\
Item C & 22.1 & 12.7 & 34.8 \\
\bottomrule
\end{tabular}
\caption{Sample Data Table}
\label{tab:sample}
\end{table}

\section{Lists}
\subsection{Bullet Points}
\begin{itemize}
  \item First item
  \item Second item
  \item Third item with sub-items:
    \begin{itemize}
      \item Sub-item 1
      \item Sub-item 2
    \end{itemize}
\end{itemize}

\subsection{Numbered List}
\begin{enumerate}
  \item First step
  \item Second step
  \item Third step
\end{enumerate}

\section{Citations and References}
When you use the citation search feature, references will be automatically managed. For example, you might cite a paper like this: \cite{example2024}. All citations will be automatically added to your \texttt{cite.bib} file.

\section{Conclusion}
This template demonstrates basic \LaTeX{} features including mathematical formulas, tables, and lists. The bibliography system is now set up to work seamlessly with Bibby AI's citation management. Feel free to modify it for your needs!

\bibliographystyle{plain}
\bibliography{cite}

\end{document}\documentclass{article}
\usepackage[utf8]{inputenc}
\usepackage{booktabs}
\usepackage{amsmath}
\usepackage{natbib}

\title{My Document}
\author{Your Name}
\date{\today}

\begin{document}

\maketitle

\section{Introduction}
This is a sample document demonstrating various \LaTeX{} features including mathematical equations, tables, and formatting. You can cite papers using the citation search feature, and they will be automatically added to your \texttt{cite.bib} file.

\section{Mathematical Expressions}
Here are some example equations:

\subsection{Quadratic Formula}
The quadratic formula is given by:
\[x = \frac{-b \pm \sqrt{b^2 - 4ac}}{2a}\]

\subsection{Matrix Example}
A sample matrix:
\[
\begin{pmatrix}
a_{11} & a_{12} & a_{13} \\
a_{21} & a_{22} & a_{23} \\
a_{31} & a_{32} & a_{33}
\end{pmatrix}
\]

\section{Tables}
Below is a sample table showing data:

\begin{table}[h]
\centering
\begin{tabular}{@{}lrrr@{}}
\toprule
\textbf{Category} & \textbf{Value 1} & \textbf{Value 2} & \textbf{Total} \\
\midrule
Item A & 10.5 & 20.3 & 30.8 \\
Item B & 15.2 & 25.1 & 40.3 \\
Item C & 22.1 & 12.7 & 34.8 \\
\bottomrule
\end{tabular}
\caption{Sample Data Table}
\label{tab:sample}
\end{table}

\section{Lists}
\subsection{Bullet Points}
\begin{itemize}
  \item First item
  \item Second item
  \item Third item with sub-items:
    \begin{itemize}
      \item Sub-item 1
      \item Sub-item 2
    \end{itemize}
\end{itemize}

\subsection{Numbered List}
\begin{enumerate}
  \item First step
  \item Second step
  \item Third step
\end{enumerate}

\section{Citations and References}
When you use the citation search feature, references will be automatically managed. For example, you might cite a paper like this: \cite{example2024}. All citations will be automatically added to your \texttt{cite.bib} file.

\section{Conclusion}
This template demonstrates basic \LaTeX{} features including mathematical formulas, tables, and lists. The bibliography system is now set up to work seamlessly with Bibby AI's citation management. Feel free to modify it for your needs!

\bibliographystyle{plain}
\bibliography{cite}

\end{document}\documentclass{article}
\usepackage[utf8]{inputenc}
\usepackage{booktabs}
\usepackage{amsmath}
\usepackage{natbib}

\title{My Document}
\author{Your Name}
\date{\today}

\begin{document}

\maketitle

\section{Introduction}
This is a sample document demonstrating various \LaTeX{} features including mathematical equations, tables, and formatting. You can cite papers using the citation search feature, and they will be automatically added to your \texttt{cite.bib} file.

\section{Mathematical Expressions}
Here are some example equations:

\subsection{Quadratic Formula}
The quadratic formula is given by:
\[x = \frac{-b \pm \sqrt{b^2 - 4ac}}{2a}\]

\subsection{Matrix Example}
A sample matrix:
\[
\begin{pmatrix}
a_{11} & a_{12} & a_{13} \\
a_{21} & a_{22} & a_{23} \\
a_{31} & a_{32} & a_{33}
\end{pmatrix}
\]

\section{Tables}
Below is a sample table showing data:

\begin{table}[h]
\centering
\begin{tabular}{@{}lrrr@{}}
\toprule
\textbf{Category} & \textbf{Value 1} & \textbf{Value 2} & \textbf{Total} \\
\midrule
Item A & 10.5 & 20.3 & 30.8 \\
Item B & 15.2 & 25.1 & 40.3 \\
Item C & 22.1 & 12.7 & 34.8 \\
\bottomrule
\end{tabular}
\caption{Sample Data Table}
\label{tab:sample}
\end{table}

\section{Lists}
\subsection{Bullet Points}
\begin{itemize}
  \item First item
  \item Second item
  \item Third item with sub-items:
    \begin{itemize}
      \item Sub-item 1
      \item Sub-item 2
    \end{itemize}
\end{itemize}

\subsection{Numbered List}
\begin{enumerate}
  \item First step
  \item Second step
  \item Third step
\end{enumerate}

\section{Citations and References}
When you use the citation search feature, references will be automatically managed. For example, you might cite a paper like this: \cite{example2024}. All citations will be automatically added to your \texttt{cite.bib} file.

\section{Conclusion}
This template demonstrates basic \LaTeX{} features including mathematical formulas, tables, and lists. The bibliography system is now set up to work seamlessly with Bibby AI's citation management. Feel free to modify it for your needs!

\bibliographystyle{plain}
\bibliography{cite}

\end{document}\documentclass{article}
\usepackage[utf8]{inputenc}
\usepackage{booktabs}
\usepackage{amsmath}
\usepackage{natbib}

\title{My Document}
\author{Your Name}
\date{\today}

\begin{document}

\maketitle

\section{Introduction}
This is a sample document demonstrating various \LaTeX{} features including mathematical equations, tables, and formatting. You can cite papers using the citation search feature, and they will be automatically added to your \texttt{cite.bib} file.

\section{Mathematical Expressions}
Here are some example equations:

\subsection{Quadratic Formula}
The quadratic formula is given by:
\[x = \frac{-b \pm \sqrt{b^2 - 4ac}}{2a}\]

\subsection{Matrix Example}
A sample matrix:
\[
\begin{pmatrix}
a_{11} & a_{12} & a_{13} \\
a_{21} & a_{22} & a_{23} \\
a_{31} & a_{32} & a_{33}
\end{pmatrix}
\]

\section{Tables}
Below is a sample table showing data:

\begin{table}[h]
\centering
\begin{tabular}{@{}lrrr@{}}
\toprule
\textbf{Category} & \textbf{Value 1} & \textbf{Value 2} & \textbf{Total} \\
\midrule
Item A & 10.5 & 20.3 & 30.8 \\
Item B & 15.2 & 25.1 & 40.3 \\
Item C & 22.1 & 12.7 & 34.8 \\
\bottomrule
\end{tabular}
\caption{Sample Data Table}
\label{tab:sample}
\end{table}

\section{Lists}
\subsection{Bullet Points}
\begin{itemize}
  \item First item
  \item Second item
  \item Third item with sub-items:
    \begin{itemize}
      \item Sub-item 1
      \item Sub-item 2
    \end{itemize}
\end{itemize}

\subsection{Numbered List}
\begin{enumerate}
  \item First step
  \item Second step
  \item Third step
\end{enumerate}

\section{Citations and References}
When you use the citation search feature, references will be automatically managed. For example, you might cite a paper like this: \cite{example2024}. All citations will be automatically added to your \texttt{cite.bib} file.

\section{Conclusion}
This template demonstrates basic \LaTeX{} features including mathematical formulas, tables, and lists. The bibliography system is now set up to work seamlessly with Bibby AI's citation management. Feel free to modify it for your needs!

\bibliographystyle{plain}
\bibliography{cite}

\end{document}\documentclass{article}
\usepackage[utf8]{inputenc}
\usepackage{booktabs}
\usepackage{amsmath}
\usepackage{natbib}

\title{My Document}
\author{Your Name}
\date{\today}

\begin{document}

\maketitle

\section{Introduction}
This is a sample document demonstrating various \LaTeX{} features including mathematical equations, tables, and formatting. You can cite papers using the citation search feature, and they will be automatically added to your \texttt{cite.bib} file.

\section{Mathematical Expressions}
Here are some example equations:

\subsection{Quadratic Formula}
The quadratic formula is given by:
\[x = \frac{-b \pm \sqrt{b^2 - 4ac}}{2a}\]

\subsection{Matrix Example}
A sample matrix:
\[
\begin{pmatrix}
a_{11} & a_{12} & a_{13} \\
a_{21} & a_{22} & a_{23} \\
a_{31} & a_{32} & a_{33}
\end{pmatrix}
\]

\section{Tables}
Below is a sample table showing data:

\begin{table}[h]
\centering
\begin{tabular}{@{}lrrr@{}}
\toprule
\textbf{Category} & \textbf{Value 1} & \textbf{Value 2} & \textbf{Total} \\
\midrule
Item A & 10.5 & 20.3 & 30.8 \\
Item B & 15.2 & 25.1 & 40.3 \\
Item C & 22.1 & 12.7 & 34.8 \\
\bottomrule
\end{tabular}
\caption{Sample Data Table}
\label{tab:sample}
\end{table}

\section{Lists}
\subsection{Bullet Points}
\begin{itemize}
  \item First item
  \item Second item
  \item Third item with sub-items:
    \begin{itemize}
      \item Sub-item 1
      \item Sub-item 2
    \end{itemize}
\end{itemize}

\subsection{Numbered List}
\begin{enumerate}
  \item First step
  \item Second step
  \item Third step
\end{enumerate}

\section{Citations and References}
When you use the citation search feature, references will be automatically managed. For example, you might cite a paper like this: \cite{example2024}. All citations will be automatically added to your \texttt{cite.bib} file.

\section{Conclusion}
This template demonstrates basic \LaTeX{} features including mathematical formulas, tables, and lists. The bibliography system is now set up to work seamlessly with Bibby AI's citation management. Feel free to modify it for your needs!

\bibliographystyle{plain}
\bibliography{cite}

\end{document}\documentclass{article}
\usepackage[utf8]{inputenc}
\usepackage{booktabs}
\usepackage{amsmath}
\usepackage{natbib}

\title{My Document}
\author{Your Name}
\date{\today}

\begin{document}

\maketitle

\section{Introduction}
This is a sample document demonstrating various \LaTeX{} features including mathematical equations, tables, and formatting. You can cite papers using the citation search feature, and they will be automatically added to your \texttt{cite.bib} file.

\section{Mathematical Expressions}
Here are some example equations:

\subsection{Quadratic Formula}
The quadratic formula is given by:
\[x = \frac{-b \pm \sqrt{b^2 - 4ac}}{2a}\]

\subsection{Matrix Example}
A sample matrix:
\[
\begin{pmatrix}
a_{11} & a_{12} & a_{13} \\
a_{21} & a_{22} & a_{23} \\
a_{31} & a_{32} & a_{33}
\end{pmatrix}
\]

\section{Tables}
Below is a sample table showing data:

\begin{table}[h]
\centering
\begin{tabular}{@{}lrrr@{}}
\toprule
\textbf{Category} & \textbf{Value 1} & \textbf{Value 2} & \textbf{Total} \\
\midrule
Item A & 10.5 & 20.3 & 30.8 \\
Item B & 15.2 & 25.1 & 40.3 \\
Item C & 22.1 & 12.7 & 34.8 \\
\bottomrule
\end{tabular}
\caption{Sample Data Table}
\label{tab:sample}
\end{table}

\section{Lists}
\subsection{Bullet Points}
\begin{itemize}
  \item First item
  \item Second item
  \item Third item with sub-items:
    \begin{itemize}
      \item Sub-item 1
      \item Sub-item 2
    \end{itemize}
\end{itemize}

\subsection{Numbered List}
\begin{enumerate}
  \item First step
  \item Second step
  \item Third step
\end{enumerate}

\section{Citations and References}
When you use the citation search feature, references will be automatically managed. For example, you might cite a paper like this: \cite{example2024}. All citations will be automatically added to your \texttt{cite.bib} file.

\section{Conclusion}
This template demonstrates basic \LaTeX{} features including mathematical formulas, tables, and lists. The bibliography system is now set up to work seamlessly with Bibby AI's citation management. Feel free to modify it for your needs!

\bibliographystyle{plain}
\bibliography{cite}

\end{document}\documentclass{article}
\usepackage[utf8]{inputenc}
\usepackage{booktabs}
\usepackage{amsmath}
\usepackage{natbib}

\title{My Document}
\author{Your Name}
\date{\today}

\begin{document}

\maketitle

\section{Introduction}
This is a sample document demonstrating various \LaTeX{} features including mathematical equations, tables, and formatting. You can cite papers using the citation search feature, and they will be automatically added to your \texttt{cite.bib} file.

\section{Mathematical Expressions}
Here are some example equations:

\subsection{Quadratic Formula}
The quadratic formula is given by:
\[x = \frac{-b \pm \sqrt{b^2 - 4ac}}{2a}\]

\subsection{Matrix Example}
A sample matrix:
\[
\begin{pmatrix}
a_{11} & a_{12} & a_{13} \\
a_{21} & a_{22} & a_{23} \\
a_{31} & a_{32} & a_{33}
\end{pmatrix}
\]

\section{Tables}
Below is a sample table showing data:

\begin{table}[h]
\centering
\begin{tabular}{@{}lrrr@{}}
\toprule
\textbf{Category} & \textbf{Value 1} & \textbf{Value 2} & \textbf{Total} \\
\midrule
Item A & 10.5 & 20.3 & 30.8 \\
Item B & 15.2 & 25.1 & 40.3 \\
Item C & 22.1 & 12.7 & 34.8 \\
\bottomrule
\end{tabular}
\caption{Sample Data Table}
\label{tab:sample}
\end{table}

\section{Lists}
\subsection{Bullet Points}
\begin{itemize}
  \item First item
  \item Second item
  \item Third item with sub-items:
    \begin{itemize}
      \item Sub-item 1
      \item Sub-item 2
    \end{itemize}
\end{itemize}

\subsection{Numbered List}
\begin{enumerate}
  \item First step
  \item Second step
  \item Third step
\end{enumerate}

\section{Citations and References}
When you use the citation search feature, references will be automatically managed. For example, you might cite a paper like this: \cite{example2024}. All citations will be automatically added to your \texttt{cite.bib} file.

\section{Conclusion}
This template demonstrates basic \LaTeX{} features including mathematical formulas, tables, and lists. The bibliography system is now set up to work seamlessly with Bibby AI's citation management. Feel free to modify it for your needs!

\bibliographystyle{plain}
\bibliography{cite}

\end{document}\documentclass{article}
\usepackage[utf8]{inputenc}
\usepackage{booktabs}
\usepackage{amsmath}
\usepackage{natbib}

\title{My Document}
\author{Your Name}
\date{\today}

\begin{document}

\maketitle

\section{Introduction}
This is a sample document demonstrating various \LaTeX{} features including mathematical equations, tables, and formatting. You can cite papers using the citation search feature, and they will be automatically added to your \texttt{cite.bib} file.

\section{Mathematical Expressions}
Here are some example equations:

\subsection{Quadratic Formula}
The quadratic formula is given by:
\[x = \frac{-b \pm \sqrt{b^2 - 4ac}}{2a}\]

\subsection{Matrix Example}
A sample matrix:
\[
\begin{pmatrix}
a_{11} & a_{12} & a_{13} \\
a_{21} & a_{22} & a_{23} \\
a_{31} & a_{32} & a_{33}
\end{pmatrix}
\]

\section{Tables}
Below is a sample table showing data:

\begin{table}[h]
\centering
\begin{tabular}{@{}lrrr@{}}
\toprule
\textbf{Category} & \textbf{Value 1} & \textbf{Value 2} & \textbf{Total} \\
\midrule
Item A & 10.5 & 20.3 & 30.8 \\
Item B & 15.2 & 25.1 & 40.3 \\
Item C & 22.1 & 12.7 & 34.8 \\
\bottomrule
\end{tabular}
\caption{Sample Data Table}
\label{tab:sample}
\end{table}

\section{Lists}
\subsection{Bullet Points}
\begin{itemize}
  \item First item
  \item Second item
  \item Third item with sub-items:
    \begin{itemize}
      \item Sub-item 1
      \item Sub-item 2
    \end{itemize}
\end{itemize}

\subsection{Numbered List}
\begin{enumerate}
  \item First step
  \item Second step
  \item Third step
\end{enumerate}

\section{Citations and References}
When you use the citation search feature, references will be automatically managed. For example, you might cite a paper like this: \cite{example2024}. All citations will be automatically added to your \texttt{cite.bib} file.

\section{Conclusion}
This template demonstrates basic \LaTeX{} features including mathematical formulas, tables, and lists. The bibliography system is now set up to work seamlessly with Bibby AI's citation management. Feel free to modify it for your needs!

\bibliographystyle{plain}
\bibliography{cite}

\end{document}\documentclass{article}
\usepackage[utf8]{inputenc}
\usepackage{booktabs}
\usepackage{amsmath}
\usepackage{natbib}

\title{My Document}
\author{Your Name}
\date{\today}

\begin{document}

\maketitle

\section{Introduction}
This is a sample document demonstrating various \LaTeX{} features including mathematical equations, tables, and formatting. You can cite papers using the citation search feature, and they will be automatically added to your \texttt{cite.bib} file.

\section{Mathematical Expressions}
Here are some example equations:

\subsection{Quadratic Formula}
The quadratic formula is given by:
\[x = \frac{-b \pm \sqrt{b^2 - 4ac}}{2a}\]

\subsection{Matrix Example}
A sample matrix:
\[
\begin{pmatrix}
a_{11} & a_{12} & a_{13} \\
a_{21} & a_{22} & a_{23} \\
a_{31} & a_{32} & a_{33}
\end{pmatrix}
\]

\section{Tables}
Below is a sample table showing data:

\begin{table}[h]
\centering
\begin{tabular}{@{}lrrr@{}}
\toprule
\textbf{Category} & \textbf{Value 1} & \textbf{Value 2} & \textbf{Total} \\
\midrule
Item A & 10.5 & 20.3 & 30.8 \\
Item B & 15.2 & 25.1 & 40.3 \\
Item C & 22.1 & 12.7 & 34.8 \\
\bottomrule
\end{tabular}
\caption{Sample Data Table}
\label{tab:sample}
\end{table}

\section{Lists}
\subsection{Bullet Points}
\begin{itemize}
  \item First item
  \item Second item
  \item Third item with sub-items:
    \begin{itemize}
      \item Sub-item 1
      \item Sub-item 2
    \end{itemize}
\end{itemize}

\subsection{Numbered List}
\begin{enumerate}
  \item First step
  \item Second step
  \item Third step
\end{enumerate}

\section{Citations and References}
When you use the citation search feature, references will be automatically managed. For example, you might cite a paper like this: \cite{example2024}. All citations will be automatically added to your \texttt{cite.bib} file.

\section{Conclusion}
This template demonstrates basic \LaTeX{} features including mathematical formulas, tables, and lists. The bibliography system is now set up to work seamlessly with Bibby AI's citation management. Feel free to modify it for your needs!

\bibliographystyle{plain}
\bibliography{cite}

\end{document}\documentclass{article}
\usepackage[utf8]{inputenc}
\usepackage{booktabs}
\usepackage{amsmath}
\usepackage{natbib}

\title{My Document}
\author{Your Name}
\date{\today}

\begin{document}

\maketitle

\section{Introduction}
This is a sample document demonstrating various \LaTeX{} features including mathematical equations, tables, and formatting. You can cite papers using the citation search feature, and they will be automatically added to your \texttt{cite.bib} file.

\section{Mathematical Expressions}
Here are some example equations:

\subsection{Quadratic Formula}
The quadratic formula is given by:
\[x = \frac{-b \pm \sqrt{b^2 - 4ac}}{2a}\]

\subsection{Matrix Example}
A sample matrix:
\[
\begin{pmatrix}
a_{11} & a_{12} & a_{13} \\
a_{21} & a_{22} & a_{23} \\
a_{31} & a_{32} & a_{33}
\end{pmatrix}
\]

\section{Tables}
Below is a sample table showing data:

\begin{table}[h]
\centering
\begin{tabular}{@{}lrrr@{}}
\toprule
\textbf{Category} & \textbf{Value 1} & \textbf{Value 2} & \textbf{Total} \\
\midrule
Item A & 10.5 & 20.3 & 30.8 \\
Item B & 15.2 & 25.1 & 40.3 \\
Item C & 22.1 & 12.7 & 34.8 \\
\bottomrule
\end{tabular}
\caption{Sample Data Table}
\label{tab:sample}
\end{table}

\section{Lists}
\subsection{Bullet Points}
\begin{itemize}
  \item First item
  \item Second item
  \item Third item with sub-items:
    \begin{itemize}
      \item Sub-item 1
      \item Sub-item 2
    \end{itemize}
\end{itemize}

\subsection{Numbered List}
\begin{enumerate}
  \item First step
  \item Second step
  \item Third step
\end{enumerate}

\section{Citations and References}
When you use the citation search feature, references will be automatically managed. For example, you might cite a paper like this: \cite{example2024}. All citations will be automatically added to your \texttt{cite.bib} file.

\section{Conclusion}
This template demonstrates basic \LaTeX{} features including mathematical formulas, tables, and lists. The bibliography system is now set up to work seamlessly with Bibby AI's citation management. Feel free to modify it for your needs!

\bibliographystyle{plain}
\bibliography{cite}

\end{document}\documentclass{article}
\usepackage[utf8]{inputenc}
\usepackage{booktabs}
\usepackage{amsmath}
\usepackage{natbib}

\title{My Document}
\author{Your Name}
\date{\today}

\begin{document}

\maketitle

\section{Introduction}
This is a sample document demonstrating various \LaTeX{} features including mathematical equations, tables, and formatting. You can cite papers using the citation search feature, and they will be automatically added to your \texttt{cite.bib} file.

\section{Mathematical Expressions}
Here are some example equations:

\subsection{Quadratic Formula}
The quadratic formula is given by:
\[x = \frac{-b \pm \sqrt{b^2 - 4ac}}{2a}\]

\subsection{Matrix Example}
A sample matrix:
\[
\begin{pmatrix}
a_{11} & a_{12} & a_{13} \\
a_{21} & a_{22} & a_{23} \\
a_{31} & a_{32} & a_{33}
\end{pmatrix}
\]

\section{Tables}
Below is a sample table showing data:

\begin{table}[h]
\centering
\begin{tabular}{@{}lrrr@{}}
\toprule
\textbf{Category} & \textbf{Value 1} & \textbf{Value 2} & \textbf{Total} \\
\midrule
Item A & 10.5 & 20.3 & 30.8 \\
Item B & 15.2 & 25.1 & 40.3 \\
Item C & 22.1 & 12.7 & 34.8 \\
\bottomrule
\end{tabular}
\caption{Sample Data Table}
\label{tab:sample}
\end{table}

\section{Lists}
\subsection{Bullet Points}
\begin{itemize}
  \item First item
  \item Second item
  \item Third item with sub-items:
    \begin{itemize}
      \item Sub-item 1
      \item Sub-item 2
    \end{itemize}
\end{itemize}

\subsection{Numbered List}
\begin{enumerate}
  \item First step
  \item Second step
  \item Third step
\end{enumerate}

\section{Citations and References}
When you use the citation search feature, references will be automatically managed. For example, you might cite a paper like this: \cite{example2024}. All citations will be automatically added to your \texttt{cite.bib} file.

\section{Conclusion}
This template demonstrates basic \LaTeX{} features including mathematical formulas, tables, and lists. The bibliography system is now set up to work seamlessly with Bibby AI's citation management. Feel free to modify it for your needs!

\bibliographystyle{plain}
\bibliography{cite}

\end{document}\documentclass{article}
\usepackage[utf8]{inputenc}
\usepackage{booktabs}
\usepackage{amsmath}
\usepackage{natbib}

\title{My Document}
\author{Your Name}
\date{\today}

\begin{document}

\maketitle

\section{Introduction}
This is a sample document demonstrating various \LaTeX{} features including mathematical equations, tables, and formatting. You can cite papers using the citation search feature, and they will be automatically added to your \texttt{cite.bib} file.

\section{Mathematical Expressions}
Here are some example equations:

\subsection{Quadratic Formula}
The quadratic formula is given by:
\[x = \frac{-b \pm \sqrt{b^2 - 4ac}}{2a}\]

\subsection{Matrix Example}
A sample matrix:
\[
\begin{pmatrix}
a_{11} & a_{12} & a_{13} \\
a_{21} & a_{22} & a_{23} \\
a_{31} & a_{32} & a_{33}
\end{pmatrix}
\]

\section{Tables}
Below is a sample table showing data:

\begin{table}[h]
\centering
\begin{tabular}{@{}lrrr@{}}
\toprule
\textbf{Category} & \textbf{Value 1} & \textbf{Value 2} & \textbf{Total} \\
\midrule
Item A & 10.5 & 20.3 & 30.8 \\
Item B & 15.2 & 25.1 & 40.3 \\
Item C & 22.1 & 12.7 & 34.8 \\
\bottomrule
\end{tabular}
\caption{Sample Data Table}
\label{tab:sample}
\end{table}

\section{Lists}
\subsection{Bullet Points}
\begin{itemize}
  \item First item
  \item Second item
  \item Third item with sub-items:
    \begin{itemize}
      \item Sub-item 1
      \item Sub-item 2
    \end{itemize}
\end{itemize}

\subsection{Numbered List}
\begin{enumerate}
  \item First step
  \item Second step
  \item Third step
\end{enumerate}

\section{Citations and References}
When you use the citation search feature, references will be automatically managed. For example, you might cite a paper like this: \cite{example2024}. All citations will be automatically added to your \texttt{cite.bib} file.

\section{Conclusion}
This template demonstrates basic \LaTeX{} features including mathematical formulas, tables, and lists. The bibliography system is now set up to work seamlessly with Bibby AI's citation management. Feel free to modify it for your needs!

\bibliographystyle{plain}
\bibliography{cite}

\end{document}\documentclass{article}
\usepackage[utf8]{inputenc}
\usepackage{booktabs}
\usepackage{amsmath}
\usepackage{natbib}

\title{My Document}
\author{Your Name}
\date{\today}

\begin{document}

\maketitle

\section{Introduction}
This is a sample document demonstrating various \LaTeX{} features including mathematical equations, tables, and formatting. You can cite papers using the citation search feature, and they will be automatically added to your \texttt{cite.bib} file.

\section{Mathematical Expressions}
Here are some example equations:

\subsection{Quadratic Formula}
The quadratic formula is given by:
\[x = \frac{-b \pm \sqrt{b^2 - 4ac}}{2a}\]

\subsection{Matrix Example}
A sample matrix:
\[
\begin{pmatrix}
a_{11} & a_{12} & a_{13} \\
a_{21} & a_{22} & a_{23} \\
a_{31} & a_{32} & a_{33}
\end{pmatrix}
\]

\section{Tables}
Below is a sample table showing data:

\begin{table}[h]
\centering
\begin{tabular}{@{}lrrr@{}}
\toprule
\textbf{Category} & \textbf{Value 1} & \textbf{Value 2} & \textbf{Total} \\
\midrule
Item A & 10.5 & 20.3 & 30.8 \\
Item B & 15.2 & 25.1 & 40.3 \\
Item C & 22.1 & 12.7 & 34.8 \\
\bottomrule
\end{tabular}
\caption{Sample Data Table}
\label{tab:sample}
\end{table}

\section{Lists}
\subsection{Bullet Points}
\begin{itemize}
  \item First item
  \item Second item
  \item Third item with sub-items:
    \begin{itemize}
      \item Sub-item 1
      \item Sub-item 2
    \end{itemize}
\end{itemize}

\subsection{Numbered List}
\begin{enumerate}
  \item First step
  \item Second step
  \item Third step
\end{enumerate}

\section{Citations and References}
When you use the citation search feature, references will be automatically managed. For example, you might cite a paper like this: \cite{example2024}. All citations will be automatically added to your \texttt{cite.bib} file.

\section{Conclusion}
This template demonstrates basic \LaTeX{} features including mathematical formulas, tables, and lists. The bibliography system is now set up to work seamlessly with Bibby AI's citation management. Feel free to modify it for your needs!

\bibliographystyle{plain}
\bibliography{cite}

\end{document}\documentclass{article}
\usepackage[utf8]{inputenc}
\usepackage{booktabs}
\usepackage{amsmath}
\usepackage{natbib}

\title{My Document}
\author{Your Name}
\date{\today}

\begin{document}

\maketitle

\section{Introduction}
This is a sample document demonstrating various \LaTeX{} features including mathematical equations, tables, and formatting. You can cite papers using the citation search feature, and they will be automatically added to your \texttt{cite.bib} file.

\section{Mathematical Expressions}
Here are some example equations:

\subsection{Quadratic Formula}
The quadratic formula is given by:
\[x = \frac{-b \pm \sqrt{b^2 - 4ac}}{2a}\]

\subsection{Matrix Example}
A sample matrix:
\[
\begin{pmatrix}
a_{11} & a_{12} & a_{13} \\
a_{21} & a_{22} & a_{23} \\
a_{31} & a_{32} & a_{33}
\end{pmatrix}
\]

\section{Tables}
Below is a sample table showing data:

\begin{table}[h]
\centering
\begin{tabular}{@{}lrrr@{}}
\toprule
\textbf{Category} & \textbf{Value 1} & \textbf{Value 2} & \textbf{Total} \\
\midrule
Item A & 10.5 & 20.3 & 30.8 \\
Item B & 15.2 & 25.1 & 40.3 \\
Item C & 22.1 & 12.7 & 34.8 \\
\bottomrule
\end{tabular}
\caption{Sample Data Table}
\label{tab:sample}
\end{table}

\section{Lists}
\subsection{Bullet Points}
\begin{itemize}
  \item First item
  \item Second item
  \item Third item with sub-items:
    \begin{itemize}
      \item Sub-item 1
      \item Sub-item 2
    \end{itemize}
\end{itemize}

\subsection{Numbered List}
\begin{enumerate}
  \item First step
  \item Second step
  \item Third step
\end{enumerate}

\section{Citations and References}
When you use the citation search feature, references will be automatically managed. For example, you might cite a paper like this: \cite{example2024}. All citations will be automatically added to your \texttt{cite.bib} file.

\section{Conclusion}
This template demonstrates basic \LaTeX{} features including mathematical formulas, tables, and lists. The bibliography system is now set up to work seamlessly with Bibby AI's citation management. Feel free to modify it for your needs!

\bibliographystyle{plain}
\bibliography{cite}

\end{document}\documentclass{article}
\usepackage[utf8]{inputenc}
\usepackage{booktabs}
\usepackage{amsmath}
\usepackage{natbib}

\title{My Document}
\author{Your Name}
\date{\today}

\begin{document}

\maketitle

\section{Introduction}
This is a sample document demonstrating various \LaTeX{} features including mathematical equations, tables, and formatting. You can cite papers using the citation search feature, and they will be automatically added to your \texttt{cite.bib} file.

\section{Mathematical Expressions}
Here are some example equations:

\subsection{Quadratic Formula}
The quadratic formula is given by:
\[x = \frac{-b \pm \sqrt{b^2 - 4ac}}{2a}\]

\subsection{Matrix Example}
A sample matrix:
\[
\begin{pmatrix}
a_{11} & a_{12} & a_{13} \\
a_{21} & a_{22} & a_{23} \\
a_{31} & a_{32} & a_{33}
\end{pmatrix}
\]

\section{Tables}
Below is a sample table showing data:

\begin{table}[h]
\centering
\begin{tabular}{@{}lrrr@{}}
\toprule
\textbf{Category} & \textbf{Value 1} & \textbf{Value 2} & \textbf{Total} \\
\midrule
Item A & 10.5 & 20.3 & 30.8 \\
Item B & 15.2 & 25.1 & 40.3 \\
Item C & 22.1 & 12.7 & 34.8 \\
\bottomrule
\end{tabular}
\caption{Sample Data Table}
\label{tab:sample}
\end{table}

\section{Lists}
\subsection{Bullet Points}
\begin{itemize}
  \item First item
  \item Second item
  \item Third item with sub-items:
    \begin{itemize}
      \item Sub-item 1
      \item Sub-item 2
    \end{itemize}
\end{itemize}

\subsection{Numbered List}
\begin{enumerate}
  \item First step
  \item Second step
  \item Third step
\end{enumerate}

\section{Citations and References}
When you use the citation search feature, references will be automatically managed. For example, you might cite a paper like this: \cite{example2024}. All citations will be automatically added to your \texttt{cite.bib} file.

\section{Conclusion}
This template demonstrates basic \LaTeX{} features including mathematical formulas, tables, and lists. The bibliography system is now set up to work seamlessly with Bibby AI's citation management. Feel free to modify it for your needs!

\bibliographystyle{plain}
\bibliography{cite}

\end{document}\documentclass{article}
\usepackage[utf8]{inputenc}
\usepackage{booktabs}
\usepackage{amsmath}
\usepackage{natbib}

\title{My Document}
\author{Your Name}
\date{\today}

\begin{document}

\maketitle

\section{Introduction}
This is a sample document demonstrating various \LaTeX{} features including mathematical equations, tables, and formatting. You can cite papers using the citation search feature, and they will be automatically added to your \texttt{cite.bib} file.

\section{Mathematical Expressions}
Here are some example equations:

\subsection{Quadratic Formula}
The quadratic formula is given by:
\[x = \frac{-b \pm \sqrt{b^2 - 4ac}}{2a}\]

\subsection{Matrix Example}
A sample matrix:
\[
\begin{pmatrix}
a_{11} & a_{12} & a_{13} \\
a_{21} & a_{22} & a_{23} \\
a_{31} & a_{32} & a_{33}
\end{pmatrix}
\]

\section{Tables}
Below is a sample table showing data:

\begin{table}[h]
\centering
\begin{tabular}{@{}lrrr@{}}
\toprule
\textbf{Category} & \textbf{Value 1} & \textbf{Value 2} & \textbf{Total} \\
\midrule
Item A & 10.5 & 20.3 & 30.8 \\
Item B & 15.2 & 25.1 & 40.3 \\
Item C & 22.1 & 12.7 & 34.8 \\
\bottomrule
\end{tabular}
\caption{Sample Data Table}
\label{tab:sample}
\end{table}

\section{Lists}
\subsection{Bullet Points}
\begin{itemize}
  \item First item
  \item Second item
  \item Third item with sub-items:
    \begin{itemize}
      \item Sub-item 1
      \item Sub-item 2
    \end{itemize}
\end{itemize}

\subsection{Numbered List}
\begin{enumerate}
  \item First step
  \item Second step
  \item Third step
\end{enumerate}

\section{Citations and References}
When you use the citation search feature, references will be automatically managed. For example, you might cite a paper like this: \cite{example2024}. All citations will be automatically added to your \texttt{cite.bib} file.

\section{Conclusion}
This template demonstrates basic \LaTeX{} features including mathematical formulas, tables, and lists. The bibliography system is now set up to work seamlessly with Bibby AI's citation management. Feel free to modify it for your needs!

\bibliographystyle{plain}
\bibliography{cite}

\end{document}\documentclass{article}
\usepackage[utf8]{inputenc}
\usepackage{booktabs}
\usepackage{amsmath}
\usepackage{natbib}

\title{My Document}
\author{Your Name}
\date{\today}

\begin{document}

\maketitle

\section{Introduction}
This is a sample document demonstrating various \LaTeX{} features including mathematical equations, tables, and formatting. You can cite papers using the citation search feature, and they will be automatically added to your \texttt{cite.bib} file.

\section{Mathematical Expressions}
Here are some example equations:

\subsection{Quadratic Formula}
The quadratic formula is given by:
\[x = \frac{-b \pm \sqrt{b^2 - 4ac}}{2a}\]

\subsection{Matrix Example}
A sample matrix:
\[
\begin{pmatrix}
a_{11} & a_{12} & a_{13} \\
a_{21} & a_{22} & a_{23} \\
a_{31} & a_{32} & a_{33}
\end{pmatrix}
\]

\section{Tables}
Below is a sample table showing data:

\begin{table}[h]
\centering
\begin{tabular}{@{}lrrr@{}}
\toprule
\textbf{Category} & \textbf{Value 1} & \textbf{Value 2} & \textbf{Total} \\
\midrule
Item A & 10.5 & 20.3 & 30.8 \\
Item B & 15.2 & 25.1 & 40.3 \\
Item C & 22.1 & 12.7 & 34.8 \\
\bottomrule
\end{tabular}
\caption{Sample Data Table}
\label{tab:sample}
\end{table}

\section{Lists}
\subsection{Bullet Points}
\begin{itemize}
  \item First item
  \item Second item
  \item Third item with sub-items:
    \begin{itemize}
      \item Sub-item 1
      \item Sub-item 2
    \end{itemize}
\end{itemize}

\subsection{Numbered List}
\begin{enumerate}
  \item First step
  \item Second step
  \item Third step
\end{enumerate}

\section{Citations and References}
When you use the citation search feature, references will be automatically managed. For example, you might cite a paper like this: \cite{example2024}. All citations will be automatically added to your \texttt{cite.bib} file.

\section{Conclusion}
This template demonstrates basic \LaTeX{} features including mathematical formulas, tables, and lists. The bibliography system is now set up to work seamlessly with Bibby AI's citation management. Feel free to modify it for your needs!

\bibliographystyle{plain}
\bibliography{cite}

\end{document}\documentclass{article}
\usepackage[utf8]{inputenc}
\usepackage{booktabs}
\usepackage{amsmath}
\usepackage{natbib}

\title{My Document}
\author{Your Name}
\date{\today}

\begin{document}

\maketitle

\section{Introduction}
This is a sample document demonstrating various \LaTeX{} features including mathematical equations, tables, and formatting. You can cite papers using the citation search feature, and they will be automatically added to your \texttt{cite.bib} file.

\section{Mathematical Expressions}
Here are some example equations:

\subsection{Quadratic Formula}
The quadratic formula is given by:
\[x = \frac{-b \pm \sqrt{b^2 - 4ac}}{2a}\]

\subsection{Matrix Example}
A sample matrix:
\[
\begin{pmatrix}
a_{11} & a_{12} & a_{13} \\
a_{21} & a_{22} & a_{23} \\
a_{31} & a_{32} & a_{33}
\end{pmatrix}
\]

\section{Tables}
Below is a sample table showing data:

\begin{table}[h]
\centering
\begin{tabular}{@{}lrrr@{}}
\toprule
\textbf{Category} & \textbf{Value 1} & \textbf{Value 2} & \textbf{Total} \\
\midrule
Item A & 10.5 & 20.3 & 30.8 \\
Item B & 15.2 & 25.1 & 40.3 \\
Item C & 22.1 & 12.7 & 34.8 \\
\bottomrule
\end{tabular}
\caption{Sample Data Table}
\label{tab:sample}
\end{table}

\section{Lists}
\subsection{Bullet Points}
\begin{itemize}
  \item First item
  \item Second item
  \item Third item with sub-items:
    \begin{itemize}
      \item Sub-item 1
      \item Sub-item 2
    \end{itemize}
\end{itemize}

\subsection{Numbered List}
\begin{enumerate}
  \item First step
  \item Second step
  \item Third step
\end{enumerate}

\section{Citations and References}
When you use the citation search feature, references will be automatically managed. For example, you might cite a paper like this: \cite{example2024}. All citations will be automatically added to your \texttt{cite.bib} file.

\section{Conclusion}
This template demonstrates basic \LaTeX{} features including mathematical formulas, tables, and lists. The bibliography system is now set up to work seamlessly with Bibby AI's citation management. Feel free to modify it for your needs!

\bibliographystyle{plain}
\bibliography{cite}

\end{document}\documentclass{article}
\usepackage[utf8]{inputenc}
\usepackage{booktabs}
\usepackage{amsmath}
\usepackage{natbib}

\title{My Document}
\author{Your Name}
\date{\today}

\begin{document}

\maketitle

\section{Introduction}
This is a sample document demonstrating various \LaTeX{} features including mathematical equations, tables, and formatting. You can cite papers using the citation search feature, and they will be automatically added to your \texttt{cite.bib} file.

\section{Mathematical Expressions}
Here are some example equations:

\subsection{Quadratic Formula}
The quadratic formula is given by:
\[x = \frac{-b \pm \sqrt{b^2 - 4ac}}{2a}\]

\subsection{Matrix Example}
A sample matrix:
\[
\begin{pmatrix}
a_{11} & a_{12} & a_{13} \\
a_{21} & a_{22} & a_{23} \\
a_{31} & a_{32} & a_{33}
\end{pmatrix}
\]

\section{Tables}
Below is a sample table showing data:

\begin{table}[h]
\centering
\begin{tabular}{@{}lrrr@{}}
\toprule
\textbf{Category} & \textbf{Value 1} & \textbf{Value 2} & \textbf{Total} \\
\midrule
Item A & 10.5 & 20.3 & 30.8 \\
Item B & 15.2 & 25.1 & 40.3 \\
Item C & 22.1 & 12.7 & 34.8 \\
\bottomrule
\end{tabular}
\caption{Sample Data Table}
\label{tab:sample}
\end{table}

\section{Lists}
\subsection{Bullet Points}
\begin{itemize}
  \item First item
  \item Second item
  \item Third item with sub-items:
    \begin{itemize}
      \item Sub-item 1
      \item Sub-item 2
    \end{itemize}
\end{itemize}

\subsection{Numbered List}
\begin{enumerate}
  \item First step
  \item Second step
  \item Third step
\end{enumerate}

\section{Citations and References}
When you use the citation search feature, references will be automatically managed. For example, you might cite a paper like this: \cite{example2024}. All citations will be automatically added to your \texttt{cite.bib} file.

\section{Conclusion}
This template demonstrates basic \LaTeX{} features including mathematical formulas, tables, and lists. The bibliography system is now set up to work seamlessly with Bibby AI's citation management. Feel free to modify it for your needs!

\bibliographystyle{plain}
\bibliography{cite}

\end{document}\documentclass{article}
\usepackage[utf8]{inputenc}
\usepackage{booktabs}
\usepackage{amsmath}
\usepackage{natbib}

\title{My Document}
\author{Your Name}
\date{\today}

\begin{document}

\maketitle

\section{Introduction}
This is a sample document demonstrating various \LaTeX{} features including mathematical equations, tables, and formatting. You can cite papers using the citation search feature, and they will be automatically added to your \texttt{cite.bib} file.

\section{Mathematical Expressions}
Here are some example equations:

\subsection{Quadratic Formula}
The quadratic formula is given by:
\[x = \frac{-b \pm \sqrt{b^2 - 4ac}}{2a}\]

\subsection{Matrix Example}
A sample matrix:
\[
\begin{pmatrix}
a_{11} & a_{12} & a_{13} \\
a_{21} & a_{22} & a_{23} \\
a_{31} & a_{32} & a_{33}
\end{pmatrix}
\]

\section{Tables}
Below is a sample table showing data:

\begin{table}[h]
\centering
\begin{tabular}{@{}lrrr@{}}
\toprule
\textbf{Category} & \textbf{Value 1} & \textbf{Value 2} & \textbf{Total} \\
\midrule
Item A & 10.5 & 20.3 & 30.8 \\
Item B & 15.2 & 25.1 & 40.3 \\
Item C & 22.1 & 12.7 & 34.8 \\
\bottomrule
\end{tabular}
\caption{Sample Data Table}
\label{tab:sample}
\end{table}

\section{Lists}
\subsection{Bullet Points}
\begin{itemize}
  \item First item
  \item Second item
  \item Third item with sub-items:
    \begin{itemize}
      \item Sub-item 1
      \item Sub-item 2
    \end{itemize}
\end{itemize}

\subsection{Numbered List}
\begin{enumerate}
  \item First step
  \item Second step
  \item Third step
\end{enumerate}

\section{Citations and References}
When you use the citation search feature, references will be automatically managed. For example, you might cite a paper like this: \cite{example2024}. All citations will be automatically added to your \texttt{cite.bib} file.

\section{Conclusion}
This template demonstrates basic \LaTeX{} features including mathematical formulas, tables, and lists. The bibliography system is now set up to work seamlessly with Bibby AI's citation management. Feel free to modify it for your needs!

\bibliographystyle{plain}
\bibliography{cite}

\end{document}\documentclass{article}
\usepackage[utf8]{inputenc}
\usepackage{booktabs}
\usepackage{amsmath}
\usepackage{natbib}

\title{My Document}
\author{Your Name}
\date{\today}

\begin{document}

\maketitle

\section{Introduction}
This is a sample document demonstrating various \LaTeX{} features including mathematical equations, tables, and formatting. You can cite papers using the citation search feature, and they will be automatically added to your \texttt{cite.bib} file.

\section{Mathematical Expressions}
Here are some example equations:

\subsection{Quadratic Formula}
The quadratic formula is given by:
\[x = \frac{-b \pm \sqrt{b^2 - 4ac}}{2a}\]

\subsection{Matrix Example}
A sample matrix:
\[
\begin{pmatrix}
a_{11} & a_{12} & a_{13} \\
a_{21} & a_{22} & a_{23} \\
a_{31} & a_{32} & a_{33}
\end{pmatrix}
\]

\section{Tables}
Below is a sample table showing data:

\begin{table}[h]
\centering
\begin{tabular}{@{}lrrr@{}}
\toprule
\textbf{Category} & \textbf{Value 1} & \textbf{Value 2} & \textbf{Total} \\
\midrule
Item A & 10.5 & 20.3 & 30.8 \\
Item B & 15.2 & 25.1 & 40.3 \\
Item C & 22.1 & 12.7 & 34.8 \\
\bottomrule
\end{tabular}
\caption{Sample Data Table}
\label{tab:sample}
\end{table}

\section{Lists}
\subsection{Bullet Points}
\begin{itemize}
  \item First item
  \item Second item
  \item Third item with sub-items:
    \begin{itemize}
      \item Sub-item 1
      \item Sub-item 2
    \end{itemize}
\end{itemize}

\subsection{Numbered List}
\begin{enumerate}
  \item First step
  \item Second step
  \item Third step
\end{enumerate}

\section{Citations and References}
When you use the citation search feature, references will be automatically managed. For example, you might cite a paper like this: \cite{example2024}. All citations will be automatically added to your \texttt{cite.bib} file.

\section{Conclusion}
This template demonstrates basic \LaTeX{} features including mathematical formulas, tables, and lists. The bibliography system is now set up to work seamlessly with Bibby AI's citation management. Feel free to modify it for your needs!

\bibliographystyle{plain}
\bibliography{cite}

\end{document}\documentclass{article}
\usepackage[utf8]{inputenc}
\usepackage{booktabs}
\usepackage{amsmath}
\usepackage{natbib}

\title{My Document}
\author{Your Name}
\date{\today}

\begin{document}

\maketitle

\section{Introduction}
This is a sample document demonstrating various \LaTeX{} features including mathematical equations, tables, and formatting. You can cite papers using the citation search feature, and they will be automatically added to your \texttt{cite.bib} file.

\section{Mathematical Expressions}
Here are some example equations:

\subsection{Quadratic Formula}
The quadratic formula is given by:
\[x = \frac{-b \pm \sqrt{b^2 - 4ac}}{2a}\]

\subsection{Matrix Example}
A sample matrix:
\[
\begin{pmatrix}
a_{11} & a_{12} & a_{13} \\
a_{21} & a_{22} & a_{23} \\
a_{31} & a_{32} & a_{33}
\end{pmatrix}
\]

\section{Tables}
Below is a sample table showing data:

\begin{table}[h]
\centering
\begin{tabular}{@{}lrrr@{}}
\toprule
\textbf{Category} & \textbf{Value 1} & \textbf{Value 2} & \textbf{Total} \\
\midrule
Item A & 10.5 & 20.3 & 30.8 \\
Item B & 15.2 & 25.1 & 40.3 \\
Item C & 22.1 & 12.7 & 34.8 \\
\bottomrule
\end{tabular}
\caption{Sample Data Table}
\label{tab:sample}
\end{table}

\section{Lists}
\subsection{Bullet Points}
\begin{itemize}
  \item First item
  \item Second item
  \item Third item with sub-items:
    \begin{itemize}
      \item Sub-item 1
      \item Sub-item 2
    \end{itemize}
\end{itemize}

\subsection{Numbered List}
\begin{enumerate}
  \item First step
  \item Second step
  \item Third step
\end{enumerate}

\section{Citations and References}
When you use the citation search feature, references will be automatically managed. For example, you might cite a paper like this: \cite{example2024}. All citations will be automatically added to your \texttt{cite.bib} file.

\section{Conclusion}
This template demonstrates basic \LaTeX{} features including mathematical formulas, tables, and lists. The bibliography system is now set up to work seamlessly with Bibby AI's citation management. Feel free to modify it for your needs!

\bibliographystyle{plain}
\bibliography{cite}

\end{document}\documentclass{article}
\usepackage[utf8]{inputenc}
\usepackage{booktabs}
\usepackage{amsmath}
\usepackage{natbib}

\title{My Document}
\author{Your Name}
\date{\today}

\begin{document}

\maketitle

\section{Introduction}
This is a sample document demonstrating various \LaTeX{} features including mathematical equations, tables, and formatting. You can cite papers using the citation search feature, and they will be automatically added to your \texttt{cite.bib} file.

\section{Mathematical Expressions}
Here are some example equations:

\subsection{Quadratic Formula}
The quadratic formula is given by:
\[x = \frac{-b \pm \sqrt{b^2 - 4ac}}{2a}\]

\subsection{Matrix Example}
A sample matrix:
\[
\begin{pmatrix}
a_{11} & a_{12} & a_{13} \\
a_{21} & a_{22} & a_{23} \\
a_{31} & a_{32} & a_{33}
\end{pmatrix}
\]

\section{Tables}
Below is a sample table showing data:

\begin{table}[h]
\centering
\begin{tabular}{@{}lrrr@{}}
\toprule
\textbf{Category} & \textbf{Value 1} & \textbf{Value 2} & \textbf{Total} \\
\midrule
Item A & 10.5 & 20.3 & 30.8 \\
Item B & 15.2 & 25.1 & 40.3 \\
Item C & 22.1 & 12.7 & 34.8 \\
\bottomrule
\end{tabular}
\caption{Sample Data Table}
\label{tab:sample}
\end{table}

\section{Lists}
\subsection{Bullet Points}
\begin{itemize}
  \item First item
  \item Second item
  \item Third item with sub-items:
    \begin{itemize}
      \item Sub-item 1
      \item Sub-item 2
    \end{itemize}
\end{itemize}

\subsection{Numbered List}
\begin{enumerate}
  \item First step
  \item Second step
  \item Third step
\end{enumerate}

\section{Citations and References}
When you use the citation search feature, references will be automatically managed. For example, you might cite a paper like this: \cite{example2024}. All citations will be automatically added to your \texttt{cite.bib} file.

\section{Conclusion}
This template demonstrates basic \LaTeX{} features including mathematical formulas, tables, and lists. The bibliography system is now set up to work seamlessly with Bibby AI's citation management. Feel free to modify it for your needs!

\bibliographystyle{plain}
\bibliography{cite}

\end{document}\documentclass{article}
\usepackage[utf8]{inputenc}
\usepackage{booktabs}
\usepackage{amsmath}
\usepackage{natbib}

\title{My Document}
\author{Your Name}
\date{\today}

\begin{document}

\maketitle

\section{Introduction}
This is a sample document demonstrating various \LaTeX{} features including mathematical equations, tables, and formatting. You can cite papers using the citation search feature, and they will be automatically added to your \texttt{cite.bib} file.

\section{Mathematical Expressions}
Here are some example equations:

\subsection{Quadratic Formula}
The quadratic formula is given by:
\[x = \frac{-b \pm \sqrt{b^2 - 4ac}}{2a}\]

\subsection{Matrix Example}
A sample matrix:
\[
\begin{pmatrix}
a_{11} & a_{12} & a_{13} \\
a_{21} & a_{22} & a_{23} \\
a_{31} & a_{32} & a_{33}
\end{pmatrix}
\]

\section{Tables}
Below is a sample table showing data:

\begin{table}[h]
\centering
\begin{tabular}{@{}lrrr@{}}
\toprule
\textbf{Category} & \textbf{Value 1} & \textbf{Value 2} & \textbf{Total} \\
\midrule
Item A & 10.5 & 20.3 & 30.8 \\
Item B & 15.2 & 25.1 & 40.3 \\
Item C & 22.1 & 12.7 & 34.8 \\
\bottomrule
\end{tabular}
\caption{Sample Data Table}
\label{tab:sample}
\end{table}

\section{Lists}
\subsection{Bullet Points}
\begin{itemize}
  \item First item
  \item Second item
  \item Third item with sub-items:
    \begin{itemize}
      \item Sub-item 1
      \item Sub-item 2
    \end{itemize}
\end{itemize}

\subsection{Numbered List}
\begin{enumerate}
  \item First step
  \item Second step
  \item Third step
\end{enumerate}

\section{Citations and References}
When you use the citation search feature, references will be automatically managed. For example, you might cite a paper like this: \cite{example2024}. All citations will be automatically added to your \texttt{cite.bib} file.

\section{Conclusion}
This template demonstrates basic \LaTeX{} features including mathematical formulas, tables, and lists. The bibliography system is now set up to work seamlessly with Bibby AI's citation management. Feel free to modify it for your needs!

\bibliographystyle{plain}
\bibliography{cite}

\end{document}\documentclass{article}
\usepackage[utf8]{inputenc}
\usepackage{booktabs}
\usepackage{amsmath}
\usepackage{natbib}

\title{My Document}
\author{Your Name}
\date{\today}

\begin{document}

\maketitle

\section{Introduction}
This is a sample document demonstrating various \LaTeX{} features including mathematical equations, tables, and formatting. You can cite papers using the citation search feature, and they will be automatically added to your \texttt{cite.bib} file.

\section{Mathematical Expressions}
Here are some example equations:

\subsection{Quadratic Formula}
The quadratic formula is given by:
\[x = \frac{-b \pm \sqrt{b^2 - 4ac}}{2a}\]

\subsection{Matrix Example}
A sample matrix:
\[
\begin{pmatrix}
a_{11} & a_{12} & a_{13} \\
a_{21} & a_{22} & a_{23} \\
a_{31} & a_{32} & a_{33}
\end{pmatrix}
\]

\section{Tables}
Below is a sample table showing data:

\begin{table}[h]
\centering
\begin{tabular}{@{}lrrr@{}}
\toprule
\textbf{Category} & \textbf{Value 1} & \textbf{Value 2} & \textbf{Total} \\
\midrule
Item A & 10.5 & 20.3 & 30.8 \\
Item B & 15.2 & 25.1 & 40.3 \\
Item C & 22.1 & 12.7 & 34.8 \\
\bottomrule
\end{tabular}
\caption{Sample Data Table}
\label{tab:sample}
\end{table}

\section{Lists}
\subsection{Bullet Points}
\begin{itemize}
  \item First item
  \item Second item
  \item Third item with sub-items:
    \begin{itemize}
      \item Sub-item 1
      \item Sub-item 2
    \end{itemize}
\end{itemize}

\subsection{Numbered List}
\begin{enumerate}
  \item First step
  \item Second step
  \item Third step
\end{enumerate}

\section{Citations and References}
When you use the citation search feature, references will be automatically managed. For example, you might cite a paper like this: \cite{example2024}. All citations will be automatically added to your \texttt{cite.bib} file.

\section{Conclusion}
This template demonstrates basic \LaTeX{} features including mathematical formulas, tables, and lists. The bibliography system is now set up to work seamlessly with Bibby AI's citation management. Feel free to modify it for your needs!

\bibliographystyle{plain}
\bibliography{cite}

\end{document}\documentclass{article}
\usepackage[utf8]{inputenc}
\usepackage{booktabs}
\usepackage{amsmath}
\usepackage{natbib}

\title{My Document}
\author{Your Name}
\date{\today}

\begin{document}

\maketitle

\section{Introduction}
This is a sample document demonstrating various \LaTeX{} features including mathematical equations, tables, and formatting. You can cite papers using the citation search feature, and they will be automatically added to your \texttt{cite.bib} file.

\section{Mathematical Expressions}
Here are some example equations:

\subsection{Quadratic Formula}
The quadratic formula is given by:
\[x = \frac{-b \pm \sqrt{b^2 - 4ac}}{2a}\]

\subsection{Matrix Example}
A sample matrix:
\[
\begin{pmatrix}
a_{11} & a_{12} & a_{13} \\
a_{21} & a_{22} & a_{23} \\
a_{31} & a_{32} & a_{33}
\end{pmatrix}
\]

\section{Tables}
Below is a sample table showing data:

\begin{table}[h]
\centering
\begin{tabular}{@{}lrrr@{}}
\toprule
\textbf{Category} & \textbf{Value 1} & \textbf{Value 2} & \textbf{Total} \\
\midrule
Item A & 10.5 & 20.3 & 30.8 \\
Item B & 15.2 & 25.1 & 40.3 \\
Item C & 22.1 & 12.7 & 34.8 \\
\bottomrule
\end{tabular}
\caption{Sample Data Table}
\label{tab:sample}
\end{table}

\section{Lists}
\subsection{Bullet Points}
\begin{itemize}
  \item First item
  \item Second item
  \item Third item with sub-items:
    \begin{itemize}
      \item Sub-item 1
      \item Sub-item 2
    \end{itemize}
\end{itemize}

\subsection{Numbered List}
\begin{enumerate}
  \item First step
  \item Second step
  \item Third step
\end{enumerate}

\section{Citations and References}
When you use the citation search feature, references will be automatically managed. For example, you might cite a paper like this: \cite{example2024}. All citations will be automatically added to your \texttt{cite.bib} file.

\section{Conclusion}
This template demonstrates basic \LaTeX{} features including mathematical formulas, tables, and lists. The bibliography system is now set up to work seamlessly with Bibby AI's citation management. Feel free to modify it for your needs!

\bibliographystyle{plain}
\bibliography{cite}

\end{document}\documentclass{article}
\usepackage[utf8]{inputenc}
\usepackage{booktabs}
\usepackage{amsmath}
\usepackage{natbib}

\title{My Document}
\author{Your Name}
\date{\today}

\begin{document}

\maketitle

\section{Introduction}
This is a sample document demonstrating various \LaTeX{} features including mathematical equations, tables, and formatting. You can cite papers using the citation search feature, and they will be automatically added to your \texttt{cite.bib} file.

\section{Mathematical Expressions}
Here are some example equations:

\subsection{Quadratic Formula}
The quadratic formula is given by:
\[x = \frac{-b \pm \sqrt{b^2 - 4ac}}{2a}\]

\subsection{Matrix Example}
A sample matrix:
\[
\begin{pmatrix}
a_{11} & a_{12} & a_{13} \\
a_{21} & a_{22} & a_{23} \\
a_{31} & a_{32} & a_{33}
\end{pmatrix}
\]

\section{Tables}
Below is a sample table showing data:

\begin{table}[h]
\centering
\begin{tabular}{@{}lrrr@{}}
\toprule
\textbf{Category} & \textbf{Value 1} & \textbf{Value 2} & \textbf{Total} \\
\midrule
Item A & 10.5 & 20.3 & 30.8 \\
Item B & 15.2 & 25.1 & 40.3 \\
Item C & 22.1 & 12.7 & 34.8 \\
\bottomrule
\end{tabular}
\caption{Sample Data Table}
\label{tab:sample}
\end{table}

\section{Lists}
\subsection{Bullet Points}
\begin{itemize}
  \item First item
  \item Second item
  \item Third item with sub-items:
    \begin{itemize}
      \item Sub-item 1
      \item Sub-item 2
    \end{itemize}
\end{itemize}

\subsection{Numbered List}
\begin{enumerate}
  \item First step
  \item Second step
  \item Third step
\end{enumerate}

\section{Citations and References}
When you use the citation search feature, references will be automatically managed. For example, you might cite a paper like this: \cite{example2024}. All citations will be automatically added to your \texttt{cite.bib} file.

\section{Conclusion}
This template demonstrates basic \LaTeX{} features including mathematical formulas, tables, and lists. The bibliography system is now set up to work seamlessly with Bibby AI's citation management. Feel free to modify it for your needs!

\bibliographystyle{plain}
\bibliography{cite}

\end{document}\documentclass{article}
\usepackage[utf8]{inputenc}
\usepackage{booktabs}
\usepackage{amsmath}
\usepackage{natbib}

\title{My Document}
\author{Your Name}
\date{\today}

\begin{document}

\maketitle

\section{Introduction}
This is a sample document demonstrating various \LaTeX{} features including mathematical equations, tables, and formatting. You can cite papers using the citation search feature, and they will be automatically added to your \texttt{cite.bib} file.

\section{Mathematical Expressions}
Here are some example equations:

\subsection{Quadratic Formula}
The quadratic formula is given by:
\[x = \frac{-b \pm \sqrt{b^2 - 4ac}}{2a}\]

\subsection{Matrix Example}
A sample matrix:
\[
\begin{pmatrix}
a_{11} & a_{12} & a_{13} \\
a_{21} & a_{22} & a_{23} \\
a_{31} & a_{32} & a_{33}
\end{pmatrix}
\]

\section{Tables}
Below is a sample table showing data:

\begin{table}[h]
\centering
\begin{tabular}{@{}lrrr@{}}
\toprule
\textbf{Category} & \textbf{Value 1} & \textbf{Value 2} & \textbf{Total} \\
\midrule
Item A & 10.5 & 20.3 & 30.8 \\
Item B & 15.2 & 25.1 & 40.3 \\
Item C & 22.1 & 12.7 & 34.8 \\
\bottomrule
\end{tabular}
\caption{Sample Data Table}
\label{tab:sample}
\end{table}

\section{Lists}
\subsection{Bullet Points}
\begin{itemize}
  \item First item
  \item Second item
  \item Third item with sub-items:
    \begin{itemize}
      \item Sub-item 1
      \item Sub-item 2
    \end{itemize}
\end{itemize}

\subsection{Numbered List}
\begin{enumerate}
  \item First step
  \item Second step
  \item Third step
\end{enumerate}

\section{Citations and References}
When you use the citation search feature, references will be automatically managed. For example, you might cite a paper like this: \cite{example2024}. All citations will be automatically added to your \texttt{cite.bib} file.

\section{Conclusion}
This template demonstrates basic \LaTeX{} features including mathematical formulas, tables, and lists. The bibliography system is now set up to work seamlessly with Bibby AI's citation management. Feel free to modify it for your needs!

\bibliographystyle{plain}
\bibliography{cite}

\end{document}\documentclass{article}
\usepackage[utf8]{inputenc}
\usepackage{booktabs}
\usepackage{amsmath}
\usepackage{natbib}

\title{My Document}
\author{Your Name}
\date{\today}

\begin{document}

\maketitle

\section{Introduction}
This is a sample document demonstrating various \LaTeX{} features including mathematical equations, tables, and formatting. You can cite papers using the citation search feature, and they will be automatically added to your \texttt{cite.bib} file.

\section{Mathematical Expressions}
Here are some example equations:

\subsection{Quadratic Formula}
The quadratic formula is given by:
\[x = \frac{-b \pm \sqrt{b^2 - 4ac}}{2a}\]

\subsection{Matrix Example}
A sample matrix:
\[
\begin{pmatrix}
a_{11} & a_{12} & a_{13} \\
a_{21} & a_{22} & a_{23} \\
a_{31} & a_{32} & a_{33}
\end{pmatrix}
\]

\section{Tables}
Below is a sample table showing data:

\begin{table}[h]
\centering
\begin{tabular}{@{}lrrr@{}}
\toprule
\textbf{Category} & \textbf{Value 1} & \textbf{Value 2} & \textbf{Total} \\
\midrule
Item A & 10.5 & 20.3 & 30.8 \\
Item B & 15.2 & 25.1 & 40.3 \\
Item C & 22.1 & 12.7 & 34.8 \\
\bottomrule
\end{tabular}
\caption{Sample Data Table}
\label{tab:sample}
\end{table}

\section{Lists}
\subsection{Bullet Points}
\begin{itemize}
  \item First item
  \item Second item
  \item Third item with sub-items:
    \begin{itemize}
      \item Sub-item 1
      \item Sub-item 2
    \end{itemize}
\end{itemize}

\subsection{Numbered List}
\begin{enumerate}
  \item First step
  \item Second step
  \item Third step
\end{enumerate}

\section{Citations and References}
When you use the citation search feature, references will be automatically managed. For example, you might cite a paper like this: \cite{example2024}. All citations will be automatically added to your \texttt{cite.bib} file.

\section{Conclusion}
This template demonstrates basic \LaTeX{} features including mathematical formulas, tables, and lists. The bibliography system is now set up to work seamlessly with Bibby AI's citation management. Feel free to modify it for your needs!

\bibliographystyle{plain}
\bibliography{cite}

\end{document}\documentclass{article}
\usepackage[utf8]{inputenc}
\usepackage{booktabs}
\usepackage{amsmath}
\usepackage{natbib}

\title{My Document}
\author{Your Name}
\date{\today}

\begin{document}

\maketitle

\section{Introduction}
This is a sample document demonstrating various \LaTeX{} features including mathematical equations, tables, and formatting. You can cite papers using the citation search feature, and they will be automatically added to your \texttt{cite.bib} file.

\section{Mathematical Expressions}
Here are some example equations:

\subsection{Quadratic Formula}
The quadratic formula is given by:
\[x = \frac{-b \pm \sqrt{b^2 - 4ac}}{2a}\]

\subsection{Matrix Example}
A sample matrix:
\[
\begin{pmatrix}
a_{11} & a_{12} & a_{13} \\
a_{21} & a_{22} & a_{23} \\
a_{31} & a_{32} & a_{33}
\end{pmatrix}
\]

\section{Tables}
Below is a sample table showing data:

\begin{table}[h]
\centering
\begin{tabular}{@{}lrrr@{}}
\toprule
\textbf{Category} & \textbf{Value 1} & \textbf{Value 2} & \textbf{Total} \\
\midrule
Item A & 10.5 & 20.3 & 30.8 \\
Item B & 15.2 & 25.1 & 40.3 \\
Item C & 22.1 & 12.7 & 34.8 \\
\bottomrule
\end{tabular}
\caption{Sample Data Table}
\label{tab:sample}
\end{table}

\section{Lists}
\subsection{Bullet Points}
\begin{itemize}
  \item First item
  \item Second item
  \item Third item with sub-items:
    \begin{itemize}
      \item Sub-item 1
      \item Sub-item 2
    \end{itemize}
\end{itemize}

\subsection{Numbered List}
\begin{enumerate}
  \item First step
  \item Second step
  \item Third step
\end{enumerate}

\section{Citations and References}
When you use the citation search feature, references will be automatically managed. For example, you might cite a paper like this: \cite{example2024}. All citations will be automatically added to your \texttt{cite.bib} file.

\section{Conclusion}
This template demonstrates basic \LaTeX{} features including mathematical formulas, tables, and lists. The bibliography system is now set up to work seamlessly with Bibby AI's citation management. Feel free to modify it for your needs!

\bibliographystyle{plain}
\bibliography{cite}

\end{document}\documentclass{article}
\usepackage[utf8]{inputenc}
\usepackage{booktabs}
\usepackage{amsmath}
\usepackage{natbib}

\title{My Document}
\author{Your Name}
\date{\today}

\begin{document}

\maketitle

\section{Introduction}
This is a sample document demonstrating various \LaTeX{} features including mathematical equations, tables, and formatting. You can cite papers using the citation search feature, and they will be automatically added to your \texttt{cite.bib} file.

\section{Mathematical Expressions}
Here are some example equations:

\subsection{Quadratic Formula}
The quadratic formula is given by:
\[x = \frac{-b \pm \sqrt{b^2 - 4ac}}{2a}\]

\subsection{Matrix Example}
A sample matrix:
\[
\begin{pmatrix}
a_{11} & a_{12} & a_{13} \\
a_{21} & a_{22} & a_{23} \\
a_{31} & a_{32} & a_{33}
\end{pmatrix}
\]

\section{Tables}
Below is a sample table showing data:

\begin{table}[h]
\centering
\begin{tabular}{@{}lrrr@{}}
\toprule
\textbf{Category} & \textbf{Value 1} & \textbf{Value 2} & \textbf{Total} \\
\midrule
Item A & 10.5 & 20.3 & 30.8 \\
Item B & 15.2 & 25.1 & 40.3 \\
Item C & 22.1 & 12.7 & 34.8 \\
\bottomrule
\end{tabular}
\caption{Sample Data Table}
\label{tab:sample}
\end{table}

\section{Lists}
\subsection{Bullet Points}
\begin{itemize}
  \item First item
  \item Second item
  \item Third item with sub-items:
    \begin{itemize}
      \item Sub-item 1
      \item Sub-item 2
    \end{itemize}
\end{itemize}

\subsection{Numbered List}
\begin{enumerate}
  \item First step
  \item Second step
  \item Third step
\end{enumerate}

\section{Citations and References}
When you use the citation search feature, references will be automatically managed. For example, you might cite a paper like this: \cite{example2024}. All citations will be automatically added to your \texttt{cite.bib} file.

\section{Conclusion}
This template demonstrates basic \LaTeX{} features including mathematical formulas, tables, and lists. The bibliography system is now set up to work seamlessly with Bibby AI's citation management. Feel free to modify it for your needs!

\bibliographystyle{plain}
\bibliography{cite}

\end{document}\documentclass{article}
\usepackage[utf8]{inputenc}
\usepackage{booktabs}
\usepackage{amsmath}
\usepackage{natbib}

\title{My Document}
\author{Your Name}
\date{\today}

\begin{document}

\maketitle

\section{Introduction}
This is a sample document demonstrating various \LaTeX{} features including mathematical equations, tables, and formatting. You can cite papers using the citation search feature, and they will be automatically added to your \texttt{cite.bib} file.

\section{Mathematical Expressions}
Here are some example equations:

\subsection{Quadratic Formula}
The quadratic formula is given by:
\[x = \frac{-b \pm \sqrt{b^2 - 4ac}}{2a}\]

\subsection{Matrix Example}
A sample matrix:
\[
\begin{pmatrix}
a_{11} & a_{12} & a_{13} \\
a_{21} & a_{22} & a_{23} \\
a_{31} & a_{32} & a_{33}
\end{pmatrix}
\]

\section{Tables}
Below is a sample table showing data:

\begin{table}[h]
\centering
\begin{tabular}{@{}lrrr@{}}
\toprule
\textbf{Category} & \textbf{Value 1} & \textbf{Value 2} & \textbf{Total} \\
\midrule
Item A & 10.5 & 20.3 & 30.8 \\
Item B & 15.2 & 25.1 & 40.3 \\
Item C & 22.1 & 12.7 & 34.8 \\
\bottomrule
\end{tabular}
\caption{Sample Data Table}
\label{tab:sample}
\end{table}

\section{Lists}
\subsection{Bullet Points}
\begin{itemize}
  \item First item
  \item Second item
  \item Third item with sub-items:
    \begin{itemize}
      \item Sub-item 1
      \item Sub-item 2
    \end{itemize}
\end{itemize}

\subsection{Numbered List}
\begin{enumerate}
  \item First step
  \item Second step
  \item Third step
\end{enumerate}

\section{Citations and References}
When you use the citation search feature, references will be automatically managed. For example, you might cite a paper like this: \cite{example2024}. All citations will be automatically added to your \texttt{cite.bib} file.

\section{Conclusion}
This template demonstrates basic \LaTeX{} features including mathematical formulas, tables, and lists. The bibliography system is now set up to work seamlessly with Bibby AI's citation management. Feel free to modify it for your needs!

\bibliographystyle{plain}
\bibliography{cite}

\end{document}\documentclass{article}
\usepackage[utf8]{inputenc}
\usepackage{booktabs}
\usepackage{amsmath}
\usepackage{natbib}

\title{My Document}
\author{Your Name}
\date{\today}

\begin{document}

\maketitle

\section{Introduction}
This is a sample document demonstrating various \LaTeX{} features including mathematical equations, tables, and formatting. You can cite papers using the citation search feature, and they will be automatically added to your \texttt{cite.bib} file.

\section{Mathematical Expressions}
Here are some example equations:

\subsection{Quadratic Formula}
The quadratic formula is given by:
\[x = \frac{-b \pm \sqrt{b^2 - 4ac}}{2a}\]

\subsection{Matrix Example}
A sample matrix:
\[
\begin{pmatrix}
a_{11} & a_{12} & a_{13} \\
a_{21} & a_{22} & a_{23} \\
a_{31} & a_{32} & a_{33}
\end{pmatrix}
\]

\section{Tables}
Below is a sample table showing data:

\begin{table}[h]
\centering
\begin{tabular}{@{}lrrr@{}}
\toprule
\textbf{Category} & \textbf{Value 1} & \textbf{Value 2} & \textbf{Total} \\
\midrule
Item A & 10.5 & 20.3 & 30.8 \\
Item B & 15.2 & 25.1 & 40.3 \\
Item C & 22.1 & 12.7 & 34.8 \\
\bottomrule
\end{tabular}
\caption{Sample Data Table}
\label{tab:sample}
\end{table}

\section{Lists}
\subsection{Bullet Points}
\begin{itemize}
  \item First item
  \item Second item
  \item Third item with sub-items:
    \begin{itemize}
      \item Sub-item 1
      \item Sub-item 2
    \end{itemize}
\end{itemize}

\subsection{Numbered List}
\begin{enumerate}
  \item First step
  \item Second step
  \item Third step
\end{enumerate}

\section{Citations and References}
When you use the citation search feature, references will be automatically managed. For example, you might cite a paper like this: \cite{example2024}. All citations will be automatically added to your \texttt{cite.bib} file.

\section{Conclusion}
This template demonstrates basic \LaTeX{} features including mathematical formulas, tables, and lists. The bibliography system is now set up to work seamlessly with Bibby AI's citation management. Feel free to modify it for your needs!

\bibliographystyle{plain}
\bibliography{cite}

\end{document}\documentclass{article}
\usepackage[utf8]{inputenc}
\usepackage{booktabs}
\usepackage{amsmath}
\usepackage{natbib}

\title{My Document}
\author{Your Name}
\date{\today}

\begin{document}

\maketitle

\section{Introduction}
This is a sample document demonstrating various \LaTeX{} features including mathematical equations, tables, and formatting. You can cite papers using the citation search feature, and they will be automatically added to your \texttt{cite.bib} file.

\section{Mathematical Expressions}
Here are some example equations:

\subsection{Quadratic Formula}
The quadratic formula is given by:
\[x = \frac{-b \pm \sqrt{b^2 - 4ac}}{2a}\]

\subsection{Matrix Example}
A sample matrix:
\[
\begin{pmatrix}
a_{11} & a_{12} & a_{13} \\
a_{21} & a_{22} & a_{23} \\
a_{31} & a_{32} & a_{33}
\end{pmatrix}
\]

\section{Tables}
Below is a sample table showing data:

\begin{table}[h]
\centering
\begin{tabular}{@{}lrrr@{}}
\toprule
\textbf{Category} & \textbf{Value 1} & \textbf{Value 2} & \textbf{Total} \\
\midrule
Item A & 10.5 & 20.3 & 30.8 \\
Item B & 15.2 & 25.1 & 40.3 \\
Item C & 22.1 & 12.7 & 34.8 \\
\bottomrule
\end{tabular}
\caption{Sample Data Table}
\label{tab:sample}
\end{table}

\section{Lists}
\subsection{Bullet Points}
\begin{itemize}
  \item First item
  \item Second item
  \item Third item with sub-items:
    \begin{itemize}
      \item Sub-item 1
      \item Sub-item 2
    \end{itemize}
\end{itemize}

\subsection{Numbered List}
\begin{enumerate}
  \item First step
  \item Second step
  \item Third step
\end{enumerate}

\section{Citations and References}
When you use the citation search feature, references will be automatically managed. For example, you might cite a paper like this: \cite{example2024}. All citations will be automatically added to your \texttt{cite.bib} file.

\section{Conclusion}
This template demonstrates basic \LaTeX{} features including mathematical formulas, tables, and lists. The bibliography system is now set up to work seamlessly with Bibby AI's citation management. Feel free to modify it for your needs!

\bibliographystyle{plain}
\bibliography{cite}

\end{document}\documentclass{article}
\usepackage[utf8]{inputenc}
\usepackage{booktabs}
\usepackage{amsmath}
\usepackage{natbib}

\title{My Document}
\author{Your Name}
\date{\today}

\begin{document}

\maketitle

\section{Introduction}
This is a sample document demonstrating various \LaTeX{} features including mathematical equations, tables, and formatting. You can cite papers using the citation search feature, and they will be automatically added to your \texttt{cite.bib} file.

\section{Mathematical Expressions}
Here are some example equations:

\subsection{Quadratic Formula}
The quadratic formula is given by:
\[x = \frac{-b \pm \sqrt{b^2 - 4ac}}{2a}\]

\subsection{Matrix Example}
A sample matrix:
\[
\begin{pmatrix}
a_{11} & a_{12} & a_{13} \\
a_{21} & a_{22} & a_{23} \\
a_{31} & a_{32} & a_{33}
\end{pmatrix}
\]

\section{Tables}
Below is a sample table showing data:

\begin{table}[h]
\centering
\begin{tabular}{@{}lrrr@{}}
\toprule
\textbf{Category} & \textbf{Value 1} & \textbf{Value 2} & \textbf{Total} \\
\midrule
Item A & 10.5 & 20.3 & 30.8 \\
Item B & 15.2 & 25.1 & 40.3 \\
Item C & 22.1 & 12.7 & 34.8 \\
\bottomrule
\end{tabular}
\caption{Sample Data Table}
\label{tab:sample}
\end{table}

\section{Lists}
\subsection{Bullet Points}
\begin{itemize}
  \item First item
  \item Second item
  \item Third item with sub-items:
    \begin{itemize}
      \item Sub-item 1
      \item Sub-item 2
    \end{itemize}
\end{itemize}

\subsection{Numbered List}
\begin{enumerate}
  \item First step
  \item Second step
  \item Third step
\end{enumerate}

\section{Citations and References}
When you use the citation search feature, references will be automatically managed. For example, you might cite a paper like this: \cite{example2024}. All citations will be automatically added to your \texttt{cite.bib} file.

\section{Conclusion}
This template demonstrates basic \LaTeX{} features including mathematical formulas, tables, and lists. The bibliography system is now set up to work seamlessly with Bibby AI's citation management. Feel free to modify it for your needs!

\bibliographystyle{plain}
\bibliography{cite}

\end{document}\documentclass{article}
\usepackage[utf8]{inputenc}
\usepackage{booktabs}
\usepackage{amsmath}
\usepackage{natbib}

\title{My Document}
\author{Your Name}
\date{\today}

\begin{document}

\maketitle

\section{Introduction}
This is a sample document demonstrating various \LaTeX{} features including mathematical equations, tables, and formatting. You can cite papers using the citation search feature, and they will be automatically added to your \texttt{cite.bib} file.

\section{Mathematical Expressions}
Here are some example equations:

\subsection{Quadratic Formula}
The quadratic formula is given by:
\[x = \frac{-b \pm \sqrt{b^2 - 4ac}}{2a}\]

\subsection{Matrix Example}
A sample matrix:
\[
\begin{pmatrix}
a_{11} & a_{12} & a_{13} \\
a_{21} & a_{22} & a_{23} \\
a_{31} & a_{32} & a_{33}
\end{pmatrix}
\]

\section{Tables}
Below is a sample table showing data:

\begin{table}[h]
\centering
\begin{tabular}{@{}lrrr@{}}
\toprule
\textbf{Category} & \textbf{Value 1} & \textbf{Value 2} & \textbf{Total} \\
\midrule
Item A & 10.5 & 20.3 & 30.8 \\
Item B & 15.2 & 25.1 & 40.3 \\
Item C & 22.1 & 12.7 & 34.8 \\
\bottomrule
\end{tabular}
\caption{Sample Data Table}
\label{tab:sample}
\end{table}

\section{Lists}
\subsection{Bullet Points}
\begin{itemize}
  \item First item
  \item Second item
  \item Third item with sub-items:
    \begin{itemize}
      \item Sub-item 1
      \item Sub-item 2
    \end{itemize}
\end{itemize}

\subsection{Numbered List}
\begin{enumerate}
  \item First step
  \item Second step
  \item Third step
\end{enumerate}

\section{Citations and References}
When you use the citation search feature, references will be automatically managed. For example, you might cite a paper like this: \cite{example2024}. All citations will be automatically added to your \texttt{cite.bib} file.

\section{Conclusion}
This template demonstrates basic \LaTeX{} features including mathematical formulas, tables, and lists. The bibliography system is now set up to work seamlessly with Bibby AI's citation management. Feel free to modify it for your needs!

\bibliographystyle{plain}
\bibliography{cite}

\end{document}\documentclass{article}
\usepackage[utf8]{inputenc}
\usepackage{booktabs}
\usepackage{amsmath}
\usepackage{natbib}

\title{My Document}
\author{Your Name}
\date{\today}

\begin{document}

\maketitle

\section{Introduction}
This is a sample document demonstrating various \LaTeX{} features including mathematical equations, tables, and formatting. You can cite papers using the citation search feature, and they will be automatically added to your \texttt{cite.bib} file.

\section{Mathematical Expressions}
Here are some example equations:

\subsection{Quadratic Formula}
The quadratic formula is given by:
\[x = \frac{-b \pm \sqrt{b^2 - 4ac}}{2a}\]

\subsection{Matrix Example}
A sample matrix:
\[
\begin{pmatrix}
a_{11} & a_{12} & a_{13} \\
a_{21} & a_{22} & a_{23} \\
a_{31} & a_{32} & a_{33}
\end{pmatrix}
\]

\section{Tables}
Below is a sample table showing data:

\begin{table}[h]
\centering
\begin{tabular}{@{}lrrr@{}}
\toprule
\textbf{Category} & \textbf{Value 1} & \textbf{Value 2} & \textbf{Total} \\
\midrule
Item A & 10.5 & 20.3 & 30.8 \\
Item B & 15.2 & 25.1 & 40.3 \\
Item C & 22.1 & 12.7 & 34.8 \\
\bottomrule
\end{tabular}
\caption{Sample Data Table}
\label{tab:sample}
\end{table}

\section{Lists}
\subsection{Bullet Points}
\begin{itemize}
  \item First item
  \item Second item
  \item Third item with sub-items:
    \begin{itemize}
      \item Sub-item 1
      \item Sub-item 2
    \end{itemize}
\end{itemize}

\subsection{Numbered List}
\begin{enumerate}
  \item First step
  \item Second step
  \item Third step
\end{enumerate}

\section{Citations and References}
When you use the citation search feature, references will be automatically managed. For example, you might cite a paper like this: \cite{example2024}. All citations will be automatically added to your \texttt{cite.bib} file.

\section{Conclusion}
This template demonstrates basic \LaTeX{} features including mathematical formulas, tables, and lists. The bibliography system is now set up to work seamlessly with Bibby AI's citation management. Feel free to modify it for your needs!

\bibliographystyle{plain}
\bibliography{cite}

\end{document}\documentclass{article}
\usepackage[utf8]{inputenc}
\usepackage{booktabs}
\usepackage{amsmath}
\usepackage{natbib}

\title{My Document}
\author{Your Name}
\date{\today}

\begin{document}

\maketitle

\section{Introduction}
This is a sample document demonstrating various \LaTeX{} features including mathematical equations, tables, and formatting. You can cite papers using the citation search feature, and they will be automatically added to your \texttt{cite.bib} file.

\section{Mathematical Expressions}
Here are some example equations:

\subsection{Quadratic Formula}
The quadratic formula is given by:
\[x = \frac{-b \pm \sqrt{b^2 - 4ac}}{2a}\]

\subsection{Matrix Example}
A sample matrix:
\[
\begin{pmatrix}
a_{11} & a_{12} & a_{13} \\
a_{21} & a_{22} & a_{23} \\
a_{31} & a_{32} & a_{33}
\end{pmatrix}
\]

\section{Tables}
Below is a sample table showing data:

\begin{table}[h]
\centering
\begin{tabular}{@{}lrrr@{}}
\toprule
\textbf{Category} & \textbf{Value 1} & \textbf{Value 2} & \textbf{Total} \\
\midrule
Item A & 10.5 & 20.3 & 30.8 \\
Item B & 15.2 & 25.1 & 40.3 \\
Item C & 22.1 & 12.7 & 34.8 \\
\bottomrule
\end{tabular}
\caption{Sample Data Table}
\label{tab:sample}
\end{table}

\section{Lists}
\subsection{Bullet Points}
\begin{itemize}
  \item First item
  \item Second item
  \item Third item with sub-items:
    \begin{itemize}
      \item Sub-item 1
      \item Sub-item 2
    \end{itemize}
\end{itemize}

\subsection{Numbered List}
\begin{enumerate}
  \item First step
  \item Second step
  \item Third step
\end{enumerate}

\section{Citations and References}
When you use the citation search feature, references will be automatically managed. For example, you might cite a paper like this: \cite{example2024}. All citations will be automatically added to your \texttt{cite.bib} file.

\section{Conclusion}
This template demonstrates basic \LaTeX{} features including mathematical formulas, tables, and lists. The bibliography system is now set up to work seamlessly with Bibby AI's citation management. Feel free to modify it for your needs!

\bibliographystyle{plain}
\bibliography{cite}

\end{document}\documentclass{article}
\usepackage[utf8]{inputenc}
\usepackage{booktabs}
\usepackage{amsmath}
\usepackage{natbib}

\title{My Document}
\author{Your Name}
\date{\today}

\begin{document}

\maketitle

\section{Introduction}
This is a sample document demonstrating various \LaTeX{} features including mathematical equations, tables, and formatting. You can cite papers using the citation search feature, and they will be automatically added to your \texttt{cite.bib} file.

\section{Mathematical Expressions}
Here are some example equations:

\subsection{Quadratic Formula}
The quadratic formula is given by:
\[x = \frac{-b \pm \sqrt{b^2 - 4ac}}{2a}\]

\subsection{Matrix Example}
A sample matrix:
\[
\begin{pmatrix}
a_{11} & a_{12} & a_{13} \\
a_{21} & a_{22} & a_{23} \\
a_{31} & a_{32} & a_{33}
\end{pmatrix}
\]

\section{Tables}
Below is a sample table showing data:

\begin{table}[h]
\centering
\begin{tabular}{@{}lrrr@{}}
\toprule
\textbf{Category} & \textbf{Value 1} & \textbf{Value 2} & \textbf{Total} \\
\midrule
Item A & 10.5 & 20.3 & 30.8 \\
Item B & 15.2 & 25.1 & 40.3 \\
Item C & 22.1 & 12.7 & 34.8 \\
\bottomrule
\end{tabular}
\caption{Sample Data Table}
\label{tab:sample}
\end{table}

\section{Lists}
\subsection{Bullet Points}
\begin{itemize}
  \item First item
  \item Second item
  \item Third item with sub-items:
    \begin{itemize}
      \item Sub-item 1
      \item Sub-item 2
    \end{itemize}
\end{itemize}

\subsection{Numbered List}
\begin{enumerate}
  \item First step
  \item Second step
  \item Third step
\end{enumerate}

\section{Citations and References}
When you use the citation search feature, references will be automatically managed. For example, you might cite a paper like this: \cite{example2024}. All citations will be automatically added to your \texttt{cite.bib} file.

\section{Conclusion}
This template demonstrates basic \LaTeX{} features including mathematical formulas, tables, and lists. The bibliography system is now set up to work seamlessly with Bibby AI's citation management. Feel free to modify it for your needs!

\bibliographystyle{plain}
\bibliography{cite}

\end{document}\documentclass{article}
\usepackage[utf8]{inputenc}
\usepackage{booktabs}
\usepackage{amsmath}
\usepackage{natbib}

\title{My Document}
\author{Your Name}
\date{\today}

\begin{document}

\maketitle

\section{Introduction}
This is a sample document demonstrating various \LaTeX{} features including mathematical equations, tables, and formatting. You can cite papers using the citation search feature, and they will be automatically added to your \texttt{cite.bib} file.

\section{Mathematical Expressions}
Here are some example equations:

\subsection{Quadratic Formula}
The quadratic formula is given by:
\[x = \frac{-b \pm \sqrt{b^2 - 4ac}}{2a}\]

\subsection{Matrix Example}
A sample matrix:
\[
\begin{pmatrix}
a_{11} & a_{12} & a_{13} \\
a_{21} & a_{22} & a_{23} \\
a_{31} & a_{32} & a_{33}
\end{pmatrix}
\]

\section{Tables}
Below is a sample table showing data:

\begin{table}[h]
\centering
\begin{tabular}{@{}lrrr@{}}
\toprule
\textbf{Category} & \textbf{Value 1} & \textbf{Value 2} & \textbf{Total} \\
\midrule
Item A & 10.5 & 20.3 & 30.8 \\
Item B & 15.2 & 25.1 & 40.3 \\
Item C & 22.1 & 12.7 & 34.8 \\
\bottomrule
\end{tabular}
\caption{Sample Data Table}
\label{tab:sample}
\end{table}

\section{Lists}
\subsection{Bullet Points}
\begin{itemize}
  \item First item
  \item Second item
  \item Third item with sub-items:
    \begin{itemize}
      \item Sub-item 1
      \item Sub-item 2
    \end{itemize}
\end{itemize}

\subsection{Numbered List}
\begin{enumerate}
  \item First step
  \item Second step
  \item Third step
\end{enumerate}

\section{Citations and References}
When you use the citation search feature, references will be automatically managed. For example, you might cite a paper like this: \cite{example2024}. All citations will be automatically added to your \texttt{cite.bib} file.

\section{Conclusion}
This template demonstrates basic \LaTeX{} features including mathematical formulas, tables, and lists. The bibliography system is now set up to work seamlessly with Bibby AI's citation management. Feel free to modify it for your needs!

\bibliographystyle{plain}
\bibliography{cite}

\end{document}\documentclass{article}
\usepackage[utf8]{inputenc}
\usepackage{booktabs}
\usepackage{amsmath}
\usepackage{natbib}

\title{My Document}
\author{Your Name}
\date{\today}

\begin{document}

\maketitle

\section{Introduction}
This is a sample document demonstrating various \LaTeX{} features including mathematical equations, tables, and formatting. You can cite papers using the citation search feature, and they will be automatically added to your \texttt{cite.bib} file.

\section{Mathematical Expressions}
Here are some example equations:

\subsection{Quadratic Formula}
The quadratic formula is given by:
\[x = \frac{-b \pm \sqrt{b^2 - 4ac}}{2a}\]

\subsection{Matrix Example}
A sample matrix:
\[
\begin{pmatrix}
a_{11} & a_{12} & a_{13} \\
a_{21} & a_{22} & a_{23} \\
a_{31} & a_{32} & a_{33}
\end{pmatrix}
\]

\section{Tables}
Below is a sample table showing data:

\begin{table}[h]
\centering
\begin{tabular}{@{}lrrr@{}}
\toprule
\textbf{Category} & \textbf{Value 1} & \textbf{Value 2} & \textbf{Total} \\
\midrule
Item A & 10.5 & 20.3 & 30.8 \\
Item B & 15.2 & 25.1 & 40.3 \\
Item C & 22.1 & 12.7 & 34.8 \\
\bottomrule
\end{tabular}
\caption{Sample Data Table}
\label{tab:sample}
\end{table}

\section{Lists}
\subsection{Bullet Points}
\begin{itemize}
  \item First item
  \item Second item
  \item Third item with sub-items:
    \begin{itemize}
      \item Sub-item 1
      \item Sub-item 2
    \end{itemize}
\end{itemize}

\subsection{Numbered List}
\begin{enumerate}
  \item First step
  \item Second step
  \item Third step
\end{enumerate}

\section{Citations and References}
When you use the citation search feature, references will be automatically managed. For example, you might cite a paper like this: \cite{example2024}. All citations will be automatically added to your \texttt{cite.bib} file.

\section{Conclusion}
This template demonstrates basic \LaTeX{} features including mathematical formulas, tables, and lists. The bibliography system is now set up to work seamlessly with Bibby AI's citation management. Feel free to modify it for your needs!

\bibliographystyle{plain}
\bibliography{cite}

\end{document}\documentclass{article}
\usepackage[utf8]{inputenc}
\usepackage{booktabs}
\usepackage{amsmath}
\usepackage{natbib}

\title{My Document}
\author{Your Name}
\date{\today}

\begin{document}

\maketitle

\section{Introduction}
This is a sample document demonstrating various \LaTeX{} features including mathematical equations, tables, and formatting. You can cite papers using the citation search feature, and they will be automatically added to your \texttt{cite.bib} file.

\section{Mathematical Expressions}
Here are some example equations:

\subsection{Quadratic Formula}
The quadratic formula is given by:
\[x = \frac{-b \pm \sqrt{b^2 - 4ac}}{2a}\]

\subsection{Matrix Example}
A sample matrix:
\[
\begin{pmatrix}
a_{11} & a_{12} & a_{13} \\
a_{21} & a_{22} & a_{23} \\
a_{31} & a_{32} & a_{33}
\end{pmatrix}
\]

\section{Tables}
Below is a sample table showing data:

\begin{table}[h]
\centering
\begin{tabular}{@{}lrrr@{}}
\toprule
\textbf{Category} & \textbf{Value 1} & \textbf{Value 2} & \textbf{Total} \\
\midrule
Item A & 10.5 & 20.3 & 30.8 \\
Item B & 15.2 & 25.1 & 40.3 \\
Item C & 22.1 & 12.7 & 34.8 \\
\bottomrule
\end{tabular}
\caption{Sample Data Table}
\label{tab:sample}
\end{table}

\section{Lists}
\subsection{Bullet Points}
\begin{itemize}
  \item First item
  \item Second item
  \item Third item with sub-items:
    \begin{itemize}
      \item Sub-item 1
      \item Sub-item 2
    \end{itemize}
\end{itemize}

\subsection{Numbered List}
\begin{enumerate}
  \item First step
  \item Second step
  \item Third step
\end{enumerate}

\section{Citations and References}
When you use the citation search feature, references will be automatically managed. For example, you might cite a paper like this: \cite{example2024}. All citations will be automatically added to your \texttt{cite.bib} file.

\section{Conclusion}
This template demonstrates basic \LaTeX{} features including mathematical formulas, tables, and lists. The bibliography system is now set up to work seamlessly with Bibby AI's citation management. Feel free to modify it for your needs!

\bibliographystyle{plain}
\bibliography{cite}

\end{document}\documentclass{article}
\usepackage[utf8]{inputenc}
\usepackage{booktabs}
\usepackage{amsmath}
\usepackage{natbib}

\title{My Document}
\author{Your Name}
\date{\today}

\begin{document}

\maketitle

\section{Introduction}
This is a sample document demonstrating various \LaTeX{} features including mathematical equations, tables, and formatting. You can cite papers using the citation search feature, and they will be automatically added to your \texttt{cite.bib} file.

\section{Mathematical Expressions}
Here are some example equations:

\subsection{Quadratic Formula}
The quadratic formula is given by:
\[x = \frac{-b \pm \sqrt{b^2 - 4ac}}{2a}\]

\subsection{Matrix Example}
A sample matrix:
\[
\begin{pmatrix}
a_{11} & a_{12} & a_{13} \\
a_{21} & a_{22} & a_{23} \\
a_{31} & a_{32} & a_{33}
\end{pmatrix}
\]

\section{Tables}
Below is a sample table showing data:

\begin{table}[h]
\centering
\begin{tabular}{@{}lrrr@{}}
\toprule
\textbf{Category} & \textbf{Value 1} & \textbf{Value 2} & \textbf{Total} \\
\midrule
Item A & 10.5 & 20.3 & 30.8 \\
Item B & 15.2 & 25.1 & 40.3 \\
Item C & 22.1 & 12.7 & 34.8 \\
\bottomrule
\end{tabular}
\caption{Sample Data Table}
\label{tab:sample}
\end{table}

\section{Lists}
\subsection{Bullet Points}
\begin{itemize}
  \item First item
  \item Second item
  \item Third item with sub-items:
    \begin{itemize}
      \item Sub-item 1
      \item Sub-item 2
    \end{itemize}
\end{itemize}

\subsection{Numbered List}
\begin{enumerate}
  \item First step
  \item Second step
  \item Third step
\end{enumerate}

\section{Citations and References}
When you use the citation search feature, references will be automatically managed. For example, you might cite a paper like this: \cite{example2024}. All citations will be automatically added to your \texttt{cite.bib} file.

\section{Conclusion}
This template demonstrates basic \LaTeX{} features including mathematical formulas, tables, and lists. The bibliography system is now set up to work seamlessly with Bibby AI's citation management. Feel free to modify it for your needs!

\bibliographystyle{plain}
\bibliography{cite}

\end{document}\documentclass{article}
\usepackage[utf8]{inputenc}
\usepackage{booktabs}
\usepackage{amsmath}
\usepackage{natbib}

\title{My Document}
\author{Your Name}
\date{\today}

\begin{document}

\maketitle

\section{Introduction}
This is a sample document demonstrating various \LaTeX{} features including mathematical equations, tables, and formatting. You can cite papers using the citation search feature, and they will be automatically added to your \texttt{cite.bib} file.

\section{Mathematical Expressions}
Here are some example equations:

\subsection{Quadratic Formula}
The quadratic formula is given by:
\[x = \frac{-b \pm \sqrt{b^2 - 4ac}}{2a}\]

\subsection{Matrix Example}
A sample matrix:
\[
\begin{pmatrix}
a_{11} & a_{12} & a_{13} \\
a_{21} & a_{22} & a_{23} \\
a_{31} & a_{32} & a_{33}
\end{pmatrix}
\]

\section{Tables}
Below is a sample table showing data:

\begin{table}[h]
\centering
\begin{tabular}{@{}lrrr@{}}
\toprule
\textbf{Category} & \textbf{Value 1} & \textbf{Value 2} & \textbf{Total} \\
\midrule
Item A & 10.5 & 20.3 & 30.8 \\
Item B & 15.2 & 25.1 & 40.3 \\
Item C & 22.1 & 12.7 & 34.8 \\
\bottomrule
\end{tabular}
\caption{Sample Data Table}
\label{tab:sample}
\end{table}

\section{Lists}
\subsection{Bullet Points}
\begin{itemize}
  \item First item
  \item Second item
  \item Third item with sub-items:
    \begin{itemize}
      \item Sub-item 1
      \item Sub-item 2
    \end{itemize}
\end{itemize}

\subsection{Numbered List}
\begin{enumerate}
  \item First step
  \item Second step
  \item Third step
\end{enumerate}

\section{Citations and References}
When you use the citation search feature, references will be automatically managed. For example, you might cite a paper like this: \cite{example2024}. All citations will be automatically added to your \texttt{cite.bib} file.

\section{Conclusion}
This template demonstrates basic \LaTeX{} features including mathematical formulas, tables, and lists. The bibliography system is now set up to work seamlessly with Bibby AI's citation management. Feel free to modify it for your needs!

\bibliographystyle{plain}
\bibliography{cite}

\end{document}\documentclass{article}
\usepackage[utf8]{inputenc}
\usepackage{booktabs}
\usepackage{amsmath}
\usepackage{natbib}

\title{My Document}
\author{Your Name}
\date{\today}

\begin{document}

\maketitle

\section{Introduction}
This is a sample document demonstrating various \LaTeX{} features including mathematical equations, tables, and formatting. You can cite papers using the citation search feature, and they will be automatically added to your \texttt{cite.bib} file.

\section{Mathematical Expressions}
Here are some example equations:

\subsection{Quadratic Formula}
The quadratic formula is given by:
\[x = \frac{-b \pm \sqrt{b^2 - 4ac}}{2a}\]

\subsection{Matrix Example}
A sample matrix:
\[
\begin{pmatrix}
a_{11} & a_{12} & a_{13} \\
a_{21} & a_{22} & a_{23} \\
a_{31} & a_{32} & a_{33}
\end{pmatrix}
\]

\section{Tables}
Below is a sample table showing data:

\begin{table}[h]
\centering
\begin{tabular}{@{}lrrr@{}}
\toprule
\textbf{Category} & \textbf{Value 1} & \textbf{Value 2} & \textbf{Total} \\
\midrule
Item A & 10.5 & 20.3 & 30.8 \\
Item B & 15.2 & 25.1 & 40.3 \\
Item C & 22.1 & 12.7 & 34.8 \\
\bottomrule
\end{tabular}
\caption{Sample Data Table}
\label{tab:sample}
\end{table}

\section{Lists}
\subsection{Bullet Points}
\begin{itemize}
  \item First item
  \item Second item
  \item Third item with sub-items:
    \begin{itemize}
      \item Sub-item 1
      \item Sub-item 2
    \end{itemize}
\end{itemize}

\subsection{Numbered List}
\begin{enumerate}
  \item First step
  \item Second step
  \item Third step
\end{enumerate}

\section{Citations and References}
When you use the citation search feature, references will be automatically managed. For example, you might cite a paper like this: \cite{example2024}. All citations will be automatically added to your \texttt{cite.bib} file.

\section{Conclusion}
This template demonstrates basic \LaTeX{} features including mathematical formulas, tables, and lists. The bibliography system is now set up to work seamlessly with Bibby AI's citation management. Feel free to modify it for your needs!

\bibliographystyle{plain}
\bibliography{cite}

\end{document}\documentclass{article}
\usepackage[utf8]{inputenc}
\usepackage{booktabs}
\usepackage{amsmath}
\usepackage{natbib}

\title{My Document}
\author{Your Name}
\date{\today}

\begin{document}

\maketitle

\section{Introduction}
This is a sample document demonstrating various \LaTeX{} features including mathematical equations, tables, and formatting. You can cite papers using the citation search feature, and they will be automatically added to your \texttt{cite.bib} file.

\section{Mathematical Expressions}
Here are some example equations:

\subsection{Quadratic Formula}
The quadratic formula is given by:
\[x = \frac{-b \pm \sqrt{b^2 - 4ac}}{2a}\]

\subsection{Matrix Example}
A sample matrix:
\[
\begin{pmatrix}
a_{11} & a_{12} & a_{13} \\
a_{21} & a_{22} & a_{23} \\
a_{31} & a_{32} & a_{33}
\end{pmatrix}
\]

\section{Tables}
Below is a sample table showing data:

\begin{table}[h]
\centering
\begin{tabular}{@{}lrrr@{}}
\toprule
\textbf{Category} & \textbf{Value 1} & \textbf{Value 2} & \textbf{Total} \\
\midrule
Item A & 10.5 & 20.3 & 30.8 \\
Item B & 15.2 & 25.1 & 40.3 \\
Item C & 22.1 & 12.7 & 34.8 \\
\bottomrule
\end{tabular}
\caption{Sample Data Table}
\label{tab:sample}
\end{table}

\section{Lists}
\subsection{Bullet Points}
\begin{itemize}
  \item First item
  \item Second item
  \item Third item with sub-items:
    \begin{itemize}
      \item Sub-item 1
      \item Sub-item 2
    \end{itemize}
\end{itemize}

\subsection{Numbered List}
\begin{enumerate}
  \item First step
  \item Second step
  \item Third step
\end{enumerate}

\section{Citations and References}
When you use the citation search feature, references will be automatically managed. For example, you might cite a paper like this: \cite{example2024}. All citations will be automatically added to your \texttt{cite.bib} file.

\section{Conclusion}
This template demonstrates basic \LaTeX{} features including mathematical formulas, tables, and lists. The bibliography system is now set up to work seamlessly with Bibby AI's citation management. Feel free to modify it for your needs!

\bibliographystyle{plain}
\bibliography{cite}

\end{document}\documentclass{article}
\usepackage[utf8]{inputenc}
\usepackage{booktabs}
\usepackage{amsmath}
\usepackage{natbib}

\title{My Document}
\author{Your Name}
\date{\today}

\begin{document}

\maketitle

\section{Introduction}
This is a sample document demonstrating various \LaTeX{} features including mathematical equations, tables, and formatting. You can cite papers using the citation search feature, and they will be automatically added to your \texttt{cite.bib} file.

\section{Mathematical Expressions}
Here are some example equations:

\subsection{Quadratic Formula}
The quadratic formula is given by:
\[x = \frac{-b \pm \sqrt{b^2 - 4ac}}{2a}\]

\subsection{Matrix Example}
A sample matrix:
\[
\begin{pmatrix}
a_{11} & a_{12} & a_{13} \\
a_{21} & a_{22} & a_{23} \\
a_{31} & a_{32} & a_{33}
\end{pmatrix}
\]

\section{Tables}
Below is a sample table showing data:

\begin{table}[h]
\centering
\begin{tabular}{@{}lrrr@{}}
\toprule
\textbf{Category} & \textbf{Value 1} & \textbf{Value 2} & \textbf{Total} \\
\midrule
Item A & 10.5 & 20.3 & 30.8 \\
Item B & 15.2 & 25.1 & 40.3 \\
Item C & 22.1 & 12.7 & 34.8 \\
\bottomrule
\end{tabular}
\caption{Sample Data Table}
\label{tab:sample}
\end{table}

\section{Lists}
\subsection{Bullet Points}
\begin{itemize}
  \item First item
  \item Second item
  \item Third item with sub-items:
    \begin{itemize}
      \item Sub-item 1
      \item Sub-item 2
    \end{itemize}
\end{itemize}

\subsection{Numbered List}
\begin{enumerate}
  \item First step
  \item Second step
  \item Third step
\end{enumerate}

\section{Citations and References}
When you use the citation search feature, references will be automatically managed. For example, you might cite a paper like this: \cite{example2024}. All citations will be automatically added to your \texttt{cite.bib} file.

\section{Conclusion}
This template demonstrates basic \LaTeX{} features including mathematical formulas, tables, and lists. The bibliography system is now set up to work seamlessly with Bibby AI's citation management. Feel free to modify it for your needs!

\bibliographystyle{plain}
\bibliography{cite}

\end{document}\documentclass{article}
\usepackage[utf8]{inputenc}
\usepackage{booktabs}
\usepackage{amsmath}
\usepackage{natbib}

\title{My Document}
\author{Your Name}
\date{\today}

\begin{document}

\maketitle

\section{Introduction}
This is a sample document demonstrating various \LaTeX{} features including mathematical equations, tables, and formatting. You can cite papers using the citation search feature, and they will be automatically added to your \texttt{cite.bib} file.

\section{Mathematical Expressions}
Here are some example equations:

\subsection{Quadratic Formula}
The quadratic formula is given by:
\[x = \frac{-b \pm \sqrt{b^2 - 4ac}}{2a}\]

\subsection{Matrix Example}
A sample matrix:
\[
\begin{pmatrix}
a_{11} & a_{12} & a_{13} \\
a_{21} & a_{22} & a_{23} \\
a_{31} & a_{32} & a_{33}
\end{pmatrix}
\]

\section{Tables}
Below is a sample table showing data:

\begin{table}[h]
\centering
\begin{tabular}{@{}lrrr@{}}
\toprule
\textbf{Category} & \textbf{Value 1} & \textbf{Value 2} & \textbf{Total} \\
\midrule
Item A & 10.5 & 20.3 & 30.8 \\
Item B & 15.2 & 25.1 & 40.3 \\
Item C & 22.1 & 12.7 & 34.8 \\
\bottomrule
\end{tabular}
\caption{Sample Data Table}
\label{tab:sample}
\end{table}

\section{Lists}
\subsection{Bullet Points}
\begin{itemize}
  \item First item
  \item Second item
  \item Third item with sub-items:
    \begin{itemize}
      \item Sub-item 1
      \item Sub-item 2
    \end{itemize}
\end{itemize}

\subsection{Numbered List}
\begin{enumerate}
  \item First step
  \item Second step
  \item Third step
\end{enumerate}

\section{Citations and References}
When you use the citation search feature, references will be automatically managed. For example, you might cite a paper like this: \cite{example2024}. All citations will be automatically added to your \texttt{cite.bib} file.

\section{Conclusion}
This template demonstrates basic \LaTeX{} features including mathematical formulas, tables, and lists. The bibliography system is now set up to work seamlessly with Bibby AI's citation management. Feel free to modify it for your needs!

\bibliographystyle{plain}
\bibliography{cite}

\end{document}\documentclass{article}
\usepackage[utf8]{inputenc}
\usepackage{booktabs}
\usepackage{amsmath}
\usepackage{natbib}

\title{My Document}
\author{Your Name}
\date{\today}

\begin{document}

\maketitle

\section{Introduction}
This is a sample document demonstrating various \LaTeX{} features including mathematical equations, tables, and formatting. You can cite papers using the citation search feature, and they will be automatically added to your \texttt{cite.bib} file.

\section{Mathematical Expressions}
Here are some example equations:

\subsection{Quadratic Formula}
The quadratic formula is given by:
\[x = \frac{-b \pm \sqrt{b^2 - 4ac}}{2a}\]

\subsection{Matrix Example}
A sample matrix:
\[
\begin{pmatrix}
a_{11} & a_{12} & a_{13} \\
a_{21} & a_{22} & a_{23} \\
a_{31} & a_{32} & a_{33}
\end{pmatrix}
\]

\section{Tables}
Below is a sample table showing data:

\begin{table}[h]
\centering
\begin{tabular}{@{}lrrr@{}}
\toprule
\textbf{Category} & \textbf{Value 1} & \textbf{Value 2} & \textbf{Total} \\
\midrule
Item A & 10.5 & 20.3 & 30.8 \\
Item B & 15.2 & 25.1 & 40.3 \\
Item C & 22.1 & 12.7 & 34.8 \\
\bottomrule
\end{tabular}
\caption{Sample Data Table}
\label{tab:sample}
\end{table}

\section{Lists}
\subsection{Bullet Points}
\begin{itemize}
  \item First item
  \item Second item
  \item Third item with sub-items:
    \begin{itemize}
      \item Sub-item 1
      \item Sub-item 2
    \end{itemize}
\end{itemize}

\subsection{Numbered List}
\begin{enumerate}
  \item First step
  \item Second step
  \item Third step
\end{enumerate}

\section{Citations and References}
When you use the citation search feature, references will be automatically managed. For example, you might cite a paper like this: \cite{example2024}. All citations will be automatically added to your \texttt{cite.bib} file.

\section{Conclusion}
This template demonstrates basic \LaTeX{} features including mathematical formulas, tables, and lists. The bibliography system is now set up to work seamlessly with Bibby AI's citation management. Feel free to modify it for your needs!

\bibliographystyle{plain}
\bibliography{cite}

\end{document}\documentclass{article}
\usepackage[utf8]{inputenc}
\usepackage{booktabs}
\usepackage{amsmath}
\usepackage{natbib}

\title{My Document}
\author{Your Name}
\date{\today}

\begin{document}

\maketitle

\section{Introduction}
This is a sample document demonstrating various \LaTeX{} features including mathematical equations, tables, and formatting. You can cite papers using the citation search feature, and they will be automatically added to your \texttt{cite.bib} file.

\section{Mathematical Expressions}
Here are some example equations:

\subsection{Quadratic Formula}
The quadratic formula is given by:
\[x = \frac{-b \pm \sqrt{b^2 - 4ac}}{2a}\]

\subsection{Matrix Example}
A sample matrix:
\[
\begin{pmatrix}
a_{11} & a_{12} & a_{13} \\
a_{21} & a_{22} & a_{23} \\
a_{31} & a_{32} & a_{33}
\end{pmatrix}
\]

\section{Tables}
Below is a sample table showing data:

\begin{table}[h]
\centering
\begin{tabular}{@{}lrrr@{}}
\toprule
\textbf{Category} & \textbf{Value 1} & \textbf{Value 2} & \textbf{Total} \\
\midrule
Item A & 10.5 & 20.3 & 30.8 \\
Item B & 15.2 & 25.1 & 40.3 \\
Item C & 22.1 & 12.7 & 34.8 \\
\bottomrule
\end{tabular}
\caption{Sample Data Table}
\label{tab:sample}
\end{table}

\section{Lists}
\subsection{Bullet Points}
\begin{itemize}
  \item First item
  \item Second item
  \item Third item with sub-items:
    \begin{itemize}
      \item Sub-item 1
      \item Sub-item 2
    \end{itemize}
\end{itemize}

\subsection{Numbered List}
\begin{enumerate}
  \item First step
  \item Second step
  \item Third step
\end{enumerate}

\section{Citations and References}
When you use the citation search feature, references will be automatically managed. For example, you might cite a paper like this: \cite{example2024}. All citations will be automatically added to your \texttt{cite.bib} file.

\section{Conclusion}
This template demonstrates basic \LaTeX{} features including mathematical formulas, tables, and lists. The bibliography system is now set up to work seamlessly with Bibby AI's citation management. Feel free to modify it for your needs!

\bibliographystyle{plain}
\bibliography{cite}

\end{document}\documentclass{article}
\usepackage[utf8]{inputenc}
\usepackage{booktabs}
\usepackage{amsmath}
\usepackage{natbib}

\title{My Document}
\author{Your Name}
\date{\today}

\begin{document}

\maketitle

\section{Introduction}
This is a sample document demonstrating various \LaTeX{} features including mathematical equations, tables, and formatting. You can cite papers using the citation search feature, and they will be automatically added to your \texttt{cite.bib} file.

\section{Mathematical Expressions}
Here are some example equations:

\subsection{Quadratic Formula}
The quadratic formula is given by:
\[x = \frac{-b \pm \sqrt{b^2 - 4ac}}{2a}\]

\subsection{Matrix Example}
A sample matrix:
\[
\begin{pmatrix}
a_{11} & a_{12} & a_{13} \\
a_{21} & a_{22} & a_{23} \\
a_{31} & a_{32} & a_{33}
\end{pmatrix}
\]

\section{Tables}
Below is a sample table showing data:

\begin{table}[h]
\centering
\begin{tabular}{@{}lrrr@{}}
\toprule
\textbf{Category} & \textbf{Value 1} & \textbf{Value 2} & \textbf{Total} \\
\midrule
Item A & 10.5 & 20.3 & 30.8 \\
Item B & 15.2 & 25.1 & 40.3 \\
Item C & 22.1 & 12.7 & 34.8 \\
\bottomrule
\end{tabular}
\caption{Sample Data Table}
\label{tab:sample}
\end{table}

\section{Lists}
\subsection{Bullet Points}
\begin{itemize}
  \item First item
  \item Second item
  \item Third item with sub-items:
    \begin{itemize}
      \item Sub-item 1
      \item Sub-item 2
    \end{itemize}
\end{itemize}

\subsection{Numbered List}
\begin{enumerate}
  \item First step
  \item Second step
  \item Third step
\end{enumerate}

\section{Citations and References}
When you use the citation search feature, references will be automatically managed. For example, you might cite a paper like this: \cite{example2024}. All citations will be automatically added to your \texttt{cite.bib} file.

\section{Conclusion}
This template demonstrates basic \LaTeX{} features including mathematical formulas, tables, and lists. The bibliography system is now set up to work seamlessly with Bibby AI's citation management. Feel free to modify it for your needs!

\bibliographystyle{plain}
\bibliography{cite}

\end{document}\documentclass{article}
\usepackage[utf8]{inputenc}
\usepackage{booktabs}
\usepackage{amsmath}
\usepackage{natbib}

\title{My Document}
\author{Your Name}
\date{\today}

\begin{document}

\maketitle

\section{Introduction}
This is a sample document demonstrating various \LaTeX{} features including mathematical equations, tables, and formatting. You can cite papers using the citation search feature, and they will be automatically added to your \texttt{cite.bib} file.

\section{Mathematical Expressions}
Here are some example equations:

\subsection{Quadratic Formula}
The quadratic formula is given by:
\[x = \frac{-b \pm \sqrt{b^2 - 4ac}}{2a}\]

\subsection{Matrix Example}
A sample matrix:
\[
\begin{pmatrix}
a_{11} & a_{12} & a_{13} \\
a_{21} & a_{22} & a_{23} \\
a_{31} & a_{32} & a_{33}
\end{pmatrix}
\]

\section{Tables}
Below is a sample table showing data:

\begin{table}[h]
\centering
\begin{tabular}{@{}lrrr@{}}
\toprule
\textbf{Category} & \textbf{Value 1} & \textbf{Value 2} & \textbf{Total} \\
\midrule
Item A & 10.5 & 20.3 & 30.8 \\
Item B & 15.2 & 25.1 & 40.3 \\
Item C & 22.1 & 12.7 & 34.8 \\
\bottomrule
\end{tabular}
\caption{Sample Data Table}
\label{tab:sample}
\end{table}

\section{Lists}
\subsection{Bullet Points}
\begin{itemize}
  \item First item
  \item Second item
  \item Third item with sub-items:
    \begin{itemize}
      \item Sub-item 1
      \item Sub-item 2
    \end{itemize}
\end{itemize}

\subsection{Numbered List}
\begin{enumerate}
  \item First step
  \item Second step
  \item Third step
\end{enumerate}

\section{Citations and References}
When you use the citation search feature, references will be automatically managed. For example, you might cite a paper like this: \cite{example2024}. All citations will be automatically added to your \texttt{cite.bib} file.

\section{Conclusion}
This template demonstrates basic \LaTeX{} features including mathematical formulas, tables, and lists. The bibliography system is now set up to work seamlessly with Bibby AI's citation management. Feel free to modify it for your needs!

\bibliographystyle{plain}
\bibliography{cite}

\end{document}\documentclass{article}
\usepackage[utf8]{inputenc}
\usepackage{booktabs}
\usepackage{amsmath}
\usepackage{natbib}

\title{My Document}
\author{Your Name}
\date{\today}

\begin{document}

\maketitle

\section{Introduction}
This is a sample document demonstrating various \LaTeX{} features including mathematical equations, tables, and formatting. You can cite papers using the citation search feature, and they will be automatically added to your \texttt{cite.bib} file.

\section{Mathematical Expressions}
Here are some example equations:

\subsection{Quadratic Formula}
The quadratic formula is given by:
\[x = \frac{-b \pm \sqrt{b^2 - 4ac}}{2a}\]

\subsection{Matrix Example}
A sample matrix:
\[
\begin{pmatrix}
a_{11} & a_{12} & a_{13} \\
a_{21} & a_{22} & a_{23} \\
a_{31} & a_{32} & a_{33}
\end{pmatrix}
\]

\section{Tables}
Below is a sample table showing data:

\begin{table}[h]
\centering
\begin{tabular}{@{}lrrr@{}}
\toprule
\textbf{Category} & \textbf{Value 1} & \textbf{Value 2} & \textbf{Total} \\
\midrule
Item A & 10.5 & 20.3 & 30.8 \\
Item B & 15.2 & 25.1 & 40.3 \\
Item C & 22.1 & 12.7 & 34.8 \\
\bottomrule
\end{tabular}
\caption{Sample Data Table}
\label{tab:sample}
\end{table}

\section{Lists}
\subsection{Bullet Points}
\begin{itemize}
  \item First item
  \item Second item
  \item Third item with sub-items:
    \begin{itemize}
      \item Sub-item 1
      \item Sub-item 2
    \end{itemize}
\end{itemize}

\subsection{Numbered List}
\begin{enumerate}
  \item First step
  \item Second step
  \item Third step
\end{enumerate}

\section{Citations and References}
When you use the citation search feature, references will be automatically managed. For example, you might cite a paper like this: \cite{example2024}. All citations will be automatically added to your \texttt{cite.bib} file.

\section{Conclusion}
This template demonstrates basic \LaTeX{} features including mathematical formulas, tables, and lists. The bibliography system is now set up to work seamlessly with Bibby AI's citation management. Feel free to modify it for your needs!

\bibliographystyle{plain}
\bibliography{cite}

\end{document}\documentclass{article}
\usepackage[utf8]{inputenc}
\usepackage{booktabs}
\usepackage{amsmath}
\usepackage{natbib}

\title{My Document}
\author{Your Name}
\date{\today}

\begin{document}

\maketitle

\section{Introduction}
This is a sample document demonstrating various \LaTeX{} features including mathematical equations, tables, and formatting. You can cite papers using the citation search feature, and they will be automatically added to your \texttt{cite.bib} file.

\section{Mathematical Expressions}
Here are some example equations:

\subsection{Quadratic Formula}
The quadratic formula is given by:
\[x = \frac{-b \pm \sqrt{b^2 - 4ac}}{2a}\]

\subsection{Matrix Example}
A sample matrix:
\[
\begin{pmatrix}
a_{11} & a_{12} & a_{13} \\
a_{21} & a_{22} & a_{23} \\
a_{31} & a_{32} & a_{33}
\end{pmatrix}
\]

\section{Tables}
Below is a sample table showing data:

\begin{table}[h]
\centering
\begin{tabular}{@{}lrrr@{}}
\toprule
\textbf{Category} & \textbf{Value 1} & \textbf{Value 2} & \textbf{Total} \\
\midrule
Item A & 10.5 & 20.3 & 30.8 \\
Item B & 15.2 & 25.1 & 40.3 \\
Item C & 22.1 & 12.7 & 34.8 \\
\bottomrule
\end{tabular}
\caption{Sample Data Table}
\label{tab:sample}
\end{table}

\section{Lists}
\subsection{Bullet Points}
\begin{itemize}
  \item First item
  \item Second item
  \item Third item with sub-items:
    \begin{itemize}
      \item Sub-item 1
      \item Sub-item 2
    \end{itemize}
\end{itemize}

\subsection{Numbered List}
\begin{enumerate}
  \item First step
  \item Second step
  \item Third step
\end{enumerate}

\section{Citations and References}
When you use the citation search feature, references will be automatically managed. For example, you might cite a paper like this: \cite{example2024}. All citations will be automatically added to your \texttt{cite.bib} file.

\section{Conclusion}
This template demonstrates basic \LaTeX{} features including mathematical formulas, tables, and lists. The bibliography system is now set up to work seamlessly with Bibby AI's citation management. Feel free to modify it for your needs!

\bibliographystyle{plain}
\bibliography{cite}

\end{document}\documentclass{article}
\usepackage[utf8]{inputenc}
\usepackage{booktabs}
\usepackage{amsmath}
\usepackage{natbib}

\title{My Document}
\author{Your Name}
\date{\today}

\begin{document}

\maketitle

\section{Introduction}
This is a sample document demonstrating various \LaTeX{} features including mathematical equations, tables, and formatting. You can cite papers using the citation search feature, and they will be automatically added to your \texttt{cite.bib} file.

\section{Mathematical Expressions}
Here are some example equations:

\subsection{Quadratic Formula}
The quadratic formula is given by:
\[x = \frac{-b \pm \sqrt{b^2 - 4ac}}{2a}\]

\subsection{Matrix Example}
A sample matrix:
\[
\begin{pmatrix}
a_{11} & a_{12} & a_{13} \\
a_{21} & a_{22} & a_{23} \\
a_{31} & a_{32} & a_{33}
\end{pmatrix}
\]

\section{Tables}
Below is a sample table showing data:

\begin{table}[h]
\centering
\begin{tabular}{@{}lrrr@{}}
\toprule
\textbf{Category} & \textbf{Value 1} & \textbf{Value 2} & \textbf{Total} \\
\midrule
Item A & 10.5 & 20.3 & 30.8 \\
Item B & 15.2 & 25.1 & 40.3 \\
Item C & 22.1 & 12.7 & 34.8 \\
\bottomrule
\end{tabular}
\caption{Sample Data Table}
\label{tab:sample}
\end{table}

\section{Lists}
\subsection{Bullet Points}
\begin{itemize}
  \item First item
  \item Second item
  \item Third item with sub-items:
    \begin{itemize}
      \item Sub-item 1
      \item Sub-item 2
    \end{itemize}
\end{itemize}

\subsection{Numbered List}
\begin{enumerate}
  \item First step
  \item Second step
  \item Third step
\end{enumerate}

\section{Citations and References}
When you use the citation search feature, references will be automatically managed. For example, you might cite a paper like this: \cite{example2024}. All citations will be automatically added to your \texttt{cite.bib} file.

\section{Conclusion}
This template demonstrates basic \LaTeX{} features including mathematical formulas, tables, and lists. The bibliography system is now set up to work seamlessly with Bibby AI's citation management. Feel free to modify it for your needs!

\bibliographystyle{plain}
\bibliography{cite}

\end{document}\documentclass{article}
\usepackage[utf8]{inputenc}
\usepackage{booktabs}
\usepackage{amsmath}
\usepackage{natbib}

\title{My Document}
\author{Your Name}
\date{\today}

\begin{document}

\maketitle

\section{Introduction}
This is a sample document demonstrating various \LaTeX{} features including mathematical equations, tables, and formatting. You can cite papers using the citation search feature, and they will be automatically added to your \texttt{cite.bib} file.

\section{Mathematical Expressions}
Here are some example equations:

\subsection{Quadratic Formula}
The quadratic formula is given by:
\[x = \frac{-b \pm \sqrt{b^2 - 4ac}}{2a}\]

\subsection{Matrix Example}
A sample matrix:
\[
\begin{pmatrix}
a_{11} & a_{12} & a_{13} \\
a_{21} & a_{22} & a_{23} \\
a_{31} & a_{32} & a_{33}
\end{pmatrix}
\]

\section{Tables}
Below is a sample table showing data:

\begin{table}[h]
\centering
\begin{tabular}{@{}lrrr@{}}
\toprule
\textbf{Category} & \textbf{Value 1} & \textbf{Value 2} & \textbf{Total} \\
\midrule
Item A & 10.5 & 20.3 & 30.8 \\
Item B & 15.2 & 25.1 & 40.3 \\
Item C & 22.1 & 12.7 & 34.8 \\
\bottomrule
\end{tabular}
\caption{Sample Data Table}
\label{tab:sample}
\end{table}

\section{Lists}
\subsection{Bullet Points}
\begin{itemize}
  \item First item
  \item Second item
  \item Third item with sub-items:
    \begin{itemize}
      \item Sub-item 1
      \item Sub-item 2
    \end{itemize}
\end{itemize}

\subsection{Numbered List}
\begin{enumerate}
  \item First step
  \item Second step
  \item Third step
\end{enumerate}

\section{Citations and References}
When you use the citation search feature, references will be automatically managed. For example, you might cite a paper like this: \cite{example2024}. All citations will be automatically added to your \texttt{cite.bib} file.

\section{Conclusion}
This template demonstrates basic \LaTeX{} features including mathematical formulas, tables, and lists. The bibliography system is now set up to work seamlessly with Bibby AI's citation management. Feel free to modify it for your needs!

\bibliographystyle{plain}
\bibliography{cite}

\end{document}\documentclass{article}
\usepackage[utf8]{inputenc}
\usepackage{booktabs}
\usepackage{amsmath}
\usepackage{natbib}

\title{My Document}
\author{Your Name}
\date{\today}

\begin{document}

\maketitle

\section{Introduction}
This is a sample document demonstrating various \LaTeX{} features including mathematical equations, tables, and formatting. You can cite papers using the citation search feature, and they will be automatically added to your \texttt{cite.bib} file.

\section{Mathematical Expressions}
Here are some example equations:

\subsection{Quadratic Formula}
The quadratic formula is given by:
\[x = \frac{-b \pm \sqrt{b^2 - 4ac}}{2a}\]

\subsection{Matrix Example}
A sample matrix:
\[
\begin{pmatrix}
a_{11} & a_{12} & a_{13} \\
a_{21} & a_{22} & a_{23} \\
a_{31} & a_{32} & a_{33}
\end{pmatrix}
\]

\section{Tables}
Below is a sample table showing data:

\begin{table}[h]
\centering
\begin{tabular}{@{}lrrr@{}}
\toprule
\textbf{Category} & \textbf{Value 1} & \textbf{Value 2} & \textbf{Total} \\
\midrule
Item A & 10.5 & 20.3 & 30.8 \\
Item B & 15.2 & 25.1 & 40.3 \\
Item C & 22.1 & 12.7 & 34.8 \\
\bottomrule
\end{tabular}
\caption{Sample Data Table}
\label{tab:sample}
\end{table}

\section{Lists}
\subsection{Bullet Points}
\begin{itemize}
  \item First item
  \item Second item
  \item Third item with sub-items:
    \begin{itemize}
      \item Sub-item 1
      \item Sub-item 2
    \end{itemize}
\end{itemize}

\subsection{Numbered List}
\begin{enumerate}
  \item First step
  \item Second step
  \item Third step
\end{enumerate}

\section{Citations and References}
When you use the citation search feature, references will be automatically managed. For example, you might cite a paper like this: \cite{example2024}. All citations will be automatically added to your \texttt{cite.bib} file.

\section{Conclusion}
This template demonstrates basic \LaTeX{} features including mathematical formulas, tables, and lists. The bibliography system is now set up to work seamlessly with Bibby AI's citation management. Feel free to modify it for your needs!

\bibliographystyle{plain}
\bibliography{cite}

\end{document}\documentclass{article}
\usepackage[utf8]{inputenc}
\usepackage{booktabs}
\usepackage{amsmath}
\usepackage{natbib}

\title{My Document}
\author{Your Name}
\date{\today}

\begin{document}

\maketitle

\section{Introduction}
This is a sample document demonstrating various \LaTeX{} features including mathematical equations, tables, and formatting. You can cite papers using the citation search feature, and they will be automatically added to your \texttt{cite.bib} file.

\section{Mathematical Expressions}
Here are some example equations:

\subsection{Quadratic Formula}
The quadratic formula is given by:
\[x = \frac{-b \pm \sqrt{b^2 - 4ac}}{2a}\]

\subsection{Matrix Example}
A sample matrix:
\[
\begin{pmatrix}
a_{11} & a_{12} & a_{13} \\
a_{21} & a_{22} & a_{23} \\
a_{31} & a_{32} & a_{33}
\end{pmatrix}
\]

\section{Tables}
Below is a sample table showing data:

\begin{table}[h]
\centering
\begin{tabular}{@{}lrrr@{}}
\toprule
\textbf{Category} & \textbf{Value 1} & \textbf{Value 2} & \textbf{Total} \\
\midrule
Item A & 10.5 & 20.3 & 30.8 \\
Item B & 15.2 & 25.1 & 40.3 \\
Item C & 22.1 & 12.7 & 34.8 \\
\bottomrule
\end{tabular}
\caption{Sample Data Table}
\label{tab:sample}
\end{table}

\section{Lists}
\subsection{Bullet Points}
\begin{itemize}
  \item First item
  \item Second item
  \item Third item with sub-items:
    \begin{itemize}
      \item Sub-item 1
      \item Sub-item 2
    \end{itemize}
\end{itemize}

\subsection{Numbered List}
\begin{enumerate}
  \item First step
  \item Second step
  \item Third step
\end{enumerate}

\section{Citations and References}
When you use the citation search feature, references will be automatically managed. For example, you might cite a paper like this: \cite{example2024}. All citations will be automatically added to your \texttt{cite.bib} file.

\section{Conclusion}
This template demonstrates basic \LaTeX{} features including mathematical formulas, tables, and lists. The bibliography system is now set up to work seamlessly with Bibby AI's citation management. Feel free to modify it for your needs!

\bibliographystyle{plain}
\bibliography{cite}

\end{document}\documentclass{article}
\usepackage[utf8]{inputenc}
\usepackage{booktabs}
\usepackage{amsmath}
\usepackage{natbib}

\title{My Document}
\author{Your Name}
\date{\today}

\begin{document}

\maketitle

\section{Introduction}
This is a sample document demonstrating various \LaTeX{} features including mathematical equations, tables, and formatting. You can cite papers using the citation search feature, and they will be automatically added to your \texttt{cite.bib} file.

\section{Mathematical Expressions}
Here are some example equations:

\subsection{Quadratic Formula}
The quadratic formula is given by:
\[x = \frac{-b \pm \sqrt{b^2 - 4ac}}{2a}\]

\subsection{Matrix Example}
A sample matrix:
\[
\begin{pmatrix}
a_{11} & a_{12} & a_{13} \\
a_{21} & a_{22} & a_{23} \\
a_{31} & a_{32} & a_{33}
\end{pmatrix}
\]

\section{Tables}
Below is a sample table showing data:

\begin{table}[h]
\centering
\begin{tabular}{@{}lrrr@{}}
\toprule
\textbf{Category} & \textbf{Value 1} & \textbf{Value 2} & \textbf{Total} \\
\midrule
Item A & 10.5 & 20.3 & 30.8 \\
Item B & 15.2 & 25.1 & 40.3 \\
Item C & 22.1 & 12.7 & 34.8 \\
\bottomrule
\end{tabular}
\caption{Sample Data Table}
\label{tab:sample}
\end{table}

\section{Lists}
\subsection{Bullet Points}
\begin{itemize}
  \item First item
  \item Second item
  \item Third item with sub-items:
    \begin{itemize}
      \item Sub-item 1
      \item Sub-item 2
    \end{itemize}
\end{itemize}

\subsection{Numbered List}
\begin{enumerate}
  \item First step
  \item Second step
  \item Third step
\end{enumerate}

\section{Citations and References}
When you use the citation search feature, references will be automatically managed. For example, you might cite a paper like this: \cite{example2024}. All citations will be automatically added to your \texttt{cite.bib} file.

\section{Conclusion}
This template demonstrates basic \LaTeX{} features including mathematical formulas, tables, and lists. The bibliography system is now set up to work seamlessly with Bibby AI's citation management. Feel free to modify it for your needs!

\bibliographystyle{plain}
\bibliography{cite}

\end{document}\documentclass{article}
\usepackage[utf8]{inputenc}
\usepackage{booktabs}
\usepackage{amsmath}
\usepackage{natbib}

\title{My Document}
\author{Your Name}
\date{\today}

\begin{document}

\maketitle

\section{Introduction}
This is a sample document demonstrating various \LaTeX{} features including mathematical equations, tables, and formatting. You can cite papers using the citation search feature, and they will be automatically added to your \texttt{cite.bib} file.

\section{Mathematical Expressions}
Here are some example equations:

\subsection{Quadratic Formula}
The quadratic formula is given by:
\[x = \frac{-b \pm \sqrt{b^2 - 4ac}}{2a}\]

\subsection{Matrix Example}
A sample matrix:
\[
\begin{pmatrix}
a_{11} & a_{12} & a_{13} \\
a_{21} & a_{22} & a_{23} \\
a_{31} & a_{32} & a_{33}
\end{pmatrix}
\]

\section{Tables}
Below is a sample table showing data:

\begin{table}[h]
\centering
\begin{tabular}{@{}lrrr@{}}
\toprule
\textbf{Category} & \textbf{Value 1} & \textbf{Value 2} & \textbf{Total} \\
\midrule
Item A & 10.5 & 20.3 & 30.8 \\
Item B & 15.2 & 25.1 & 40.3 \\
Item C & 22.1 & 12.7 & 34.8 \\
\bottomrule
\end{tabular}
\caption{Sample Data Table}
\label{tab:sample}
\end{table}

\section{Lists}
\subsection{Bullet Points}
\begin{itemize}
  \item First item
  \item Second item
  \item Third item with sub-items:
    \begin{itemize}
      \item Sub-item 1
      \item Sub-item 2
    \end{itemize}
\end{itemize}

\subsection{Numbered List}
\begin{enumerate}
  \item First step
  \item Second step
  \item Third step
\end{enumerate}

\section{Citations and References}
When you use the citation search feature, references will be automatically managed. For example, you might cite a paper like this: \cite{example2024}. All citations will be automatically added to your \texttt{cite.bib} file.

\section{Conclusion}
This template demonstrates basic \LaTeX{} features including mathematical formulas, tables, and lists. The bibliography system is now set up to work seamlessly with Bibby AI's citation management. Feel free to modify it for your needs!

\bibliographystyle{plain}
\bibliography{cite}

\end{document}\documentclass{article}
\usepackage[utf8]{inputenc}
\usepackage{booktabs}
\usepackage{amsmath}
\usepackage{natbib}

\title{My Document}
\author{Your Name}
\date{\today}

\begin{document}

\maketitle

\section{Introduction}
This is a sample document demonstrating various \LaTeX{} features including mathematical equations, tables, and formatting. You can cite papers using the citation search feature, and they will be automatically added to your \texttt{cite.bib} file.

\section{Mathematical Expressions}
Here are some example equations:

\subsection{Quadratic Formula}
The quadratic formula is given by:
\[x = \frac{-b \pm \sqrt{b^2 - 4ac}}{2a}\]

\subsection{Matrix Example}
A sample matrix:
\[
\begin{pmatrix}
a_{11} & a_{12} & a_{13} \\
a_{21} & a_{22} & a_{23} \\
a_{31} & a_{32} & a_{33}
\end{pmatrix}
\]

\section{Tables}
Below is a sample table showing data:

\begin{table}[h]
\centering
\begin{tabular}{@{}lrrr@{}}
\toprule
\textbf{Category} & \textbf{Value 1} & \textbf{Value 2} & \textbf{Total} \\
\midrule
Item A & 10.5 & 20.3 & 30.8 \\
Item B & 15.2 & 25.1 & 40.3 \\
Item C & 22.1 & 12.7 & 34.8 \\
\bottomrule
\end{tabular}
\caption{Sample Data Table}
\label{tab:sample}
\end{table}

\section{Lists}
\subsection{Bullet Points}
\begin{itemize}
  \item First item
  \item Second item
  \item Third item with sub-items:
    \begin{itemize}
      \item Sub-item 1
      \item Sub-item 2
    \end{itemize}
\end{itemize}

\subsection{Numbered List}
\begin{enumerate}
  \item First step
  \item Second step
  \item Third step
\end{enumerate}

\section{Citations and References}
When you use the citation search feature, references will be automatically managed. For example, you might cite a paper like this: \cite{example2024}. All citations will be automatically added to your \texttt{cite.bib} file.

\section{Conclusion}
This template demonstrates basic \LaTeX{} features including mathematical formulas, tables, and lists. The bibliography system is now set up to work seamlessly with Bibby AI's citation management. Feel free to modify it for your needs!

\bibliographystyle{plain}
\bibliography{cite}

\end{document}\documentclass{article}
\usepackage[utf8]{inputenc}
\usepackage{booktabs}
\usepackage{amsmath}
\usepackage{natbib}

\title{My Document}
\author{Your Name}
\date{\today}

\begin{document}

\maketitle

\section{Introduction}
This is a sample document demonstrating various \LaTeX{} features including mathematical equations, tables, and formatting. You can cite papers using the citation search feature, and they will be automatically added to your \texttt{cite.bib} file.

\section{Mathematical Expressions}
Here are some example equations:

\subsection{Quadratic Formula}
The quadratic formula is given by:
\[x = \frac{-b \pm \sqrt{b^2 - 4ac}}{2a}\]

\subsection{Matrix Example}
A sample matrix:
\[
\begin{pmatrix}
a_{11} & a_{12} & a_{13} \\
a_{21} & a_{22} & a_{23} \\
a_{31} & a_{32} & a_{33}
\end{pmatrix}
\]

\section{Tables}
Below is a sample table showing data:

\begin{table}[h]
\centering
\begin{tabular}{@{}lrrr@{}}
\toprule
\textbf{Category} & \textbf{Value 1} & \textbf{Value 2} & \textbf{Total} \\
\midrule
Item A & 10.5 & 20.3 & 30.8 \\
Item B & 15.2 & 25.1 & 40.3 \\
Item C & 22.1 & 12.7 & 34.8 \\
\bottomrule
\end{tabular}
\caption{Sample Data Table}
\label{tab:sample}
\end{table}

\section{Lists}
\subsection{Bullet Points}
\begin{itemize}
  \item First item
  \item Second item
  \item Third item with sub-items:
    \begin{itemize}
      \item Sub-item 1
      \item Sub-item 2
    \end{itemize}
\end{itemize}

\subsection{Numbered List}
\begin{enumerate}
  \item First step
  \item Second step
  \item Third step
\end{enumerate}

\section{Citations and References}
When you use the citation search feature, references will be automatically managed. For example, you might cite a paper like this: \cite{example2024}. All citations will be automatically added to your \texttt{cite.bib} file.

\section{Conclusion}
This template demonstrates basic \LaTeX{} features including mathematical formulas, tables, and lists. The bibliography system is now set up to work seamlessly with Bibby AI's citation management. Feel free to modify it for your needs!

\bibliographystyle{plain}
\bibliography{cite}

\end{document}\documentclass{article}
\usepackage[utf8]{inputenc}
\usepackage{booktabs}
\usepackage{amsmath}
\usepackage{natbib}

\title{My Document}
\author{Your Name}
\date{\today}

\begin{document}

\maketitle

\section{Introduction}
This is a sample document demonstrating various \LaTeX{} features including mathematical equations, tables, and formatting. You can cite papers using the citation search feature, and they will be automatically added to your \texttt{cite.bib} file.

\section{Mathematical Expressions}
Here are some example equations:

\subsection{Quadratic Formula}
The quadratic formula is given by:
\[x = \frac{-b \pm \sqrt{b^2 - 4ac}}{2a}\]

\subsection{Matrix Example}
A sample matrix:
\[
\begin{pmatrix}
a_{11} & a_{12} & a_{13} \\
a_{21} & a_{22} & a_{23} \\
a_{31} & a_{32} & a_{33}
\end{pmatrix}
\]

\section{Tables}
Below is a sample table showing data:

\begin{table}[h]
\centering
\begin{tabular}{@{}lrrr@{}}
\toprule
\textbf{Category} & \textbf{Value 1} & \textbf{Value 2} & \textbf{Total} \\
\midrule
Item A & 10.5 & 20.3 & 30.8 \\
Item B & 15.2 & 25.1 & 40.3 \\
Item C & 22.1 & 12.7 & 34.8 \\
\bottomrule
\end{tabular}
\caption{Sample Data Table}
\label{tab:sample}
\end{table}

\section{Lists}
\subsection{Bullet Points}
\begin{itemize}
  \item First item
  \item Second item
  \item Third item with sub-items:
    \begin{itemize}
      \item Sub-item 1
      \item Sub-item 2
    \end{itemize}
\end{itemize}

\subsection{Numbered List}
\begin{enumerate}
  \item First step
  \item Second step
  \item Third step
\end{enumerate}

\section{Citations and References}
When you use the citation search feature, references will be automatically managed. For example, you might cite a paper like this: \cite{example2024}. All citations will be automatically added to your \texttt{cite.bib} file.

\section{Conclusion}
This template demonstrates basic \LaTeX{} features including mathematical formulas, tables, and lists. The bibliography system is now set up to work seamlessly with Bibby AI's citation management. Feel free to modify it for your needs!

\bibliographystyle{plain}
\bibliography{cite}

\end{document}\documentclass{article}
\usepackage[utf8]{inputenc}
\usepackage{booktabs}
\usepackage{amsmath}
\usepackage{natbib}

\title{My Document}
\author{Your Name}
\date{\today}

\begin{document}

\maketitle

\section{Introduction}
This is a sample document demonstrating various \LaTeX{} features including mathematical equations, tables, and formatting. You can cite papers using the citation search feature, and they will be automatically added to your \texttt{cite.bib} file.

\section{Mathematical Expressions}
Here are some example equations:

\subsection{Quadratic Formula}
The quadratic formula is given by:
\[x = \frac{-b \pm \sqrt{b^2 - 4ac}}{2a}\]

\subsection{Matrix Example}
A sample matrix:
\[
\begin{pmatrix}
a_{11} & a_{12} & a_{13} \\
a_{21} & a_{22} & a_{23} \\
a_{31} & a_{32} & a_{33}
\end{pmatrix}
\]

\section{Tables}
Below is a sample table showing data:

\begin{table}[h]
\centering
\begin{tabular}{@{}lrrr@{}}
\toprule
\textbf{Category} & \textbf{Value 1} & \textbf{Value 2} & \textbf{Total} \\
\midrule
Item A & 10.5 & 20.3 & 30.8 \\
Item B & 15.2 & 25.1 & 40.3 \\
Item C & 22.1 & 12.7 & 34.8 \\
\bottomrule
\end{tabular}
\caption{Sample Data Table}
\label{tab:sample}
\end{table}

\section{Lists}
\subsection{Bullet Points}
\begin{itemize}
  \item First item
  \item Second item
  \item Third item with sub-items:
    \begin{itemize}
      \item Sub-item 1
      \item Sub-item 2
    \end{itemize}
\end{itemize}

\subsection{Numbered List}
\begin{enumerate}
  \item First step
  \item Second step
  \item Third step
\end{enumerate}

\section{Citations and References}
When you use the citation search feature, references will be automatically managed. For example, you might cite a paper like this: \cite{example2024}. All citations will be automatically added to your \texttt{cite.bib} file.

\section{Conclusion}
This template demonstrates basic \LaTeX{} features including mathematical formulas, tables, and lists. The bibliography system is now set up to work seamlessly with Bibby AI's citation management. Feel free to modify it for your needs!

\bibliographystyle{plain}
\bibliography{cite}

\end{document}\documentclass{article}
\usepackage[utf8]{inputenc}
\usepackage{booktabs}
\usepackage{amsmath}
\usepackage{natbib}

\title{My Document}
\author{Your Name}
\date{\today}

\begin{document}

\maketitle

\section{Introduction}
This is a sample document demonstrating various \LaTeX{} features including mathematical equations, tables, and formatting. You can cite papers using the citation search feature, and they will be automatically added to your \texttt{cite.bib} file.

\section{Mathematical Expressions}
Here are some example equations:

\subsection{Quadratic Formula}
The quadratic formula is given by:
\[x = \frac{-b \pm \sqrt{b^2 - 4ac}}{2a}\]

\subsection{Matrix Example}
A sample matrix:
\[
\begin{pmatrix}
a_{11} & a_{12} & a_{13} \\
a_{21} & a_{22} & a_{23} \\
a_{31} & a_{32} & a_{33}
\end{pmatrix}
\]

\section{Tables}
Below is a sample table showing data:

\begin{table}[h]
\centering
\begin{tabular}{@{}lrrr@{}}
\toprule
\textbf{Category} & \textbf{Value 1} & \textbf{Value 2} & \textbf{Total} \\
\midrule
Item A & 10.5 & 20.3 & 30.8 \\
Item B & 15.2 & 25.1 & 40.3 \\
Item C & 22.1 & 12.7 & 34.8 \\
\bottomrule
\end{tabular}
\caption{Sample Data Table}
\label{tab:sample}
\end{table}

\section{Lists}
\subsection{Bullet Points}
\begin{itemize}
  \item First item
  \item Second item
  \item Third item with sub-items:
    \begin{itemize}
      \item Sub-item 1
      \item Sub-item 2
    \end{itemize}
\end{itemize}

\subsection{Numbered List}
\begin{enumerate}
  \item First step
  \item Second step
  \item Third step
\end{enumerate}

\section{Citations and References}
When you use the citation search feature, references will be automatically managed. For example, you might cite a paper like this: \cite{example2024}. All citations will be automatically added to your \texttt{cite.bib} file.

\section{Conclusion}
This template demonstrates basic \LaTeX{} features including mathematical formulas, tables, and lists. The bibliography system is now set up to work seamlessly with Bibby AI's citation management. Feel free to modify it for your needs!

\bibliographystyle{plain}
\bibliography{cite}

\end{document}\documentclass{article}
\usepackage[utf8]{inputenc}
\usepackage{booktabs}
\usepackage{amsmath}
\usepackage{natbib}

\title{My Document}
\author{Your Name}
\date{\today}

\begin{document}

\maketitle

\section{Introduction}
This is a sample document demonstrating various \LaTeX{} features including mathematical equations, tables, and formatting. You can cite papers using the citation search feature, and they will be automatically added to your \texttt{cite.bib} file.

\section{Mathematical Expressions}
Here are some example equations:

\subsection{Quadratic Formula}
The quadratic formula is given by:
\[x = \frac{-b \pm \sqrt{b^2 - 4ac}}{2a}\]

\subsection{Matrix Example}
A sample matrix:
\[
\begin{pmatrix}
a_{11} & a_{12} & a_{13} \\
a_{21} & a_{22} & a_{23} \\
a_{31} & a_{32} & a_{33}
\end{pmatrix}
\]

\section{Tables}
Below is a sample table showing data:

\begin{table}[h]
\centering
\begin{tabular}{@{}lrrr@{}}
\toprule
\textbf{Category} & \textbf{Value 1} & \textbf{Value 2} & \textbf{Total} \\
\midrule
Item A & 10.5 & 20.3 & 30.8 \\
Item B & 15.2 & 25.1 & 40.3 \\
Item C & 22.1 & 12.7 & 34.8 \\
\bottomrule
\end{tabular}
\caption{Sample Data Table}
\label{tab:sample}
\end{table}

\section{Lists}
\subsection{Bullet Points}
\begin{itemize}
  \item First item
  \item Second item
  \item Third item with sub-items:
    \begin{itemize}
      \item Sub-item 1
      \item Sub-item 2
    \end{itemize}
\end{itemize}

\subsection{Numbered List}
\begin{enumerate}
  \item First step
  \item Second step
  \item Third step
\end{enumerate}

\section{Citations and References}
When you use the citation search feature, references will be automatically managed. For example, you might cite a paper like this: \cite{example2024}. All citations will be automatically added to your \texttt{cite.bib} file.

\section{Conclusion}
This template demonstrates basic \LaTeX{} features including mathematical formulas, tables, and lists. The bibliography system is now set up to work seamlessly with Bibby AI's citation management. Feel free to modify it for your needs!

\bibliographystyle{plain}
\bibliography{cite}

\end{document}\documentclass{article}
\usepackage[utf8]{inputenc}
\usepackage{booktabs}
\usepackage{amsmath}
\usepackage{natbib}

\title{My Document}
\author{Your Name}
\date{\today}

\begin{document}

\maketitle

\section{Introduction}
This is a sample document demonstrating various \LaTeX{} features including mathematical equations, tables, and formatting. You can cite papers using the citation search feature, and they will be automatically added to your \texttt{cite.bib} file.

\section{Mathematical Expressions}
Here are some example equations:

\subsection{Quadratic Formula}
The quadratic formula is given by:
\[x = \frac{-b \pm \sqrt{b^2 - 4ac}}{2a}\]

\subsection{Matrix Example}
A sample matrix:
\[
\begin{pmatrix}
a_{11} & a_{12} & a_{13} \\
a_{21} & a_{22} & a_{23} \\
a_{31} & a_{32} & a_{33}
\end{pmatrix}
\]

\section{Tables}
Below is a sample table showing data:

\begin{table}[h]
\centering
\begin{tabular}{@{}lrrr@{}}
\toprule
\textbf{Category} & \textbf{Value 1} & \textbf{Value 2} & \textbf{Total} \\
\midrule
Item A & 10.5 & 20.3 & 30.8 \\
Item B & 15.2 & 25.1 & 40.3 \\
Item C & 22.1 & 12.7 & 34.8 \\
\bottomrule
\end{tabular}
\caption{Sample Data Table}
\label{tab:sample}
\end{table}

\section{Lists}
\subsection{Bullet Points}
\begin{itemize}
  \item First item
  \item Second item
  \item Third item with sub-items:
    \begin{itemize}
      \item Sub-item 1
      \item Sub-item 2
    \end{itemize}
\end{itemize}

\subsection{Numbered List}
\begin{enumerate}
  \item First step
  \item Second step
  \item Third step
\end{enumerate}

\section{Citations and References}
When you use the citation search feature, references will be automatically managed. For example, you might cite a paper like this: \cite{example2024}. All citations will be automatically added to your \texttt{cite.bib} file.

\section{Conclusion}
This template demonstrates basic \LaTeX{} features including mathematical formulas, tables, and lists. The bibliography system is now set up to work seamlessly with Bibby AI's citation management. Feel free to modify it for your needs!

\bibliographystyle{plain}
\bibliography{cite}

\end{document}\documentclass{article}
\usepackage[utf8]{inputenc}
\usepackage{booktabs}
\usepackage{amsmath}
\usepackage{natbib}

\title{My Document}
\author{Your Name}
\date{\today}

\begin{document}

\maketitle

\section{Introduction}
This is a sample document demonstrating various \LaTeX{} features including mathematical equations, tables, and formatting. You can cite papers using the citation search feature, and they will be automatically added to your \texttt{cite.bib} file.

\section{Mathematical Expressions}
Here are some example equations:

\subsection{Quadratic Formula}
The quadratic formula is given by:
\[x = \frac{-b \pm \sqrt{b^2 - 4ac}}{2a}\]

\subsection{Matrix Example}
A sample matrix:
\[
\begin{pmatrix}
a_{11} & a_{12} & a_{13} \\
a_{21} & a_{22} & a_{23} \\
a_{31} & a_{32} & a_{33}
\end{pmatrix}
\]

\section{Tables}
Below is a sample table showing data:

\begin{table}[h]
\centering
\begin{tabular}{@{}lrrr@{}}
\toprule
\textbf{Category} & \textbf{Value 1} & \textbf{Value 2} & \textbf{Total} \\
\midrule
Item A & 10.5 & 20.3 & 30.8 \\
Item B & 15.2 & 25.1 & 40.3 \\
Item C & 22.1 & 12.7 & 34.8 \\
\bottomrule
\end{tabular}
\caption{Sample Data Table}
\label{tab:sample}
\end{table}

\section{Lists}
\subsection{Bullet Points}
\begin{itemize}
  \item First item
  \item Second item
  \item Third item with sub-items:
    \begin{itemize}
      \item Sub-item 1
      \item Sub-item 2
    \end{itemize}
\end{itemize}

\subsection{Numbered List}
\begin{enumerate}
  \item First step
  \item Second step
  \item Third step
\end{enumerate}

\section{Citations and References}
When you use the citation search feature, references will be automatically managed. For example, you might cite a paper like this: \cite{example2024}. All citations will be automatically added to your \texttt{cite.bib} file.

\section{Conclusion}
This template demonstrates basic \LaTeX{} features including mathematical formulas, tables, and lists. The bibliography system is now set up to work seamlessly with Bibby AI's citation management. Feel free to modify it for your needs!

\bibliographystyle{plain}
\bibliography{cite}

\end{document}\documentclass{article}
\usepackage[utf8]{inputenc}
\usepackage{booktabs}
\usepackage{amsmath}
\usepackage{natbib}

\title{My Document}
\author{Your Name}
\date{\today}

\begin{document}

\maketitle

\section{Introduction}
This is a sample document demonstrating various \LaTeX{} features including mathematical equations, tables, and formatting. You can cite papers using the citation search feature, and they will be automatically added to your \texttt{cite.bib} file.

\section{Mathematical Expressions}
Here are some example equations:

\subsection{Quadratic Formula}
The quadratic formula is given by:
\[x = \frac{-b \pm \sqrt{b^2 - 4ac}}{2a}\]

\subsection{Matrix Example}
A sample matrix:
\[
\begin{pmatrix}
a_{11} & a_{12} & a_{13} \\
a_{21} & a_{22} & a_{23} \\
a_{31} & a_{32} & a_{33}
\end{pmatrix}
\]

\section{Tables}
Below is a sample table showing data:

\begin{table}[h]
\centering
\begin{tabular}{@{}lrrr@{}}
\toprule
\textbf{Category} & \textbf{Value 1} & \textbf{Value 2} & \textbf{Total} \\
\midrule
Item A & 10.5 & 20.3 & 30.8 \\
Item B & 15.2 & 25.1 & 40.3 \\
Item C & 22.1 & 12.7 & 34.8 \\
\bottomrule
\end{tabular}
\caption{Sample Data Table}
\label{tab:sample}
\end{table}

\section{Lists}
\subsection{Bullet Points}
\begin{itemize}
  \item First item
  \item Second item
  \item Third item with sub-items:
    \begin{itemize}
      \item Sub-item 1
      \item Sub-item 2
    \end{itemize}
\end{itemize}

\subsection{Numbered List}
\begin{enumerate}
  \item First step
  \item Second step
  \item Third step
\end{enumerate}

\section{Citations and References}
When you use the citation search feature, references will be automatically managed. For example, you might cite a paper like this: \cite{example2024}. All citations will be automatically added to your \texttt{cite.bib} file.

\section{Conclusion}
This template demonstrates basic \LaTeX{} features including mathematical formulas, tables, and lists. The bibliography system is now set up to work seamlessly with Bibby AI's citation management. Feel free to modify it for your needs!

\bibliographystyle{plain}
\bibliography{cite}

\end{document}\documentclass{article}
\usepackage[utf8]{inputenc}
\usepackage{booktabs}
\usepackage{amsmath}
\usepackage{natbib}

\title{My Document}
\author{Your Name}
\date{\today}

\begin{document}

\maketitle

\section{Introduction}
This is a sample document demonstrating various \LaTeX{} features including mathematical equations, tables, and formatting. You can cite papers using the citation search feature, and they will be automatically added to your \texttt{cite.bib} file.

\section{Mathematical Expressions}
Here are some example equations:

\subsection{Quadratic Formula}
The quadratic formula is given by:
\[x = \frac{-b \pm \sqrt{b^2 - 4ac}}{2a}\]

\subsection{Matrix Example}
A sample matrix:
\[
\begin{pmatrix}
a_{11} & a_{12} & a_{13} \\
a_{21} & a_{22} & a_{23} \\
a_{31} & a_{32} & a_{33}
\end{pmatrix}
\]

\section{Tables}
Below is a sample table showing data:

\begin{table}[h]
\centering
\begin{tabular}{@{}lrrr@{}}
\toprule
\textbf{Category} & \textbf{Value 1} & \textbf{Value 2} & \textbf{Total} \\
\midrule
Item A & 10.5 & 20.3 & 30.8 \\
Item B & 15.2 & 25.1 & 40.3 \\
Item C & 22.1 & 12.7 & 34.8 \\
\bottomrule
\end{tabular}
\caption{Sample Data Table}
\label{tab:sample}
\end{table}

\section{Lists}
\subsection{Bullet Points}
\begin{itemize}
  \item First item
  \item Second item
  \item Third item with sub-items:
    \begin{itemize}
      \item Sub-item 1
      \item Sub-item 2
    \end{itemize}
\end{itemize}

\subsection{Numbered List}
\begin{enumerate}
  \item First step
  \item Second step
  \item Third step
\end{enumerate}

\section{Citations and References}
When you use the citation search feature, references will be automatically managed. For example, you might cite a paper like this: \cite{example2024}. All citations will be automatically added to your \texttt{cite.bib} file.

\section{Conclusion}
This template demonstrates basic \LaTeX{} features including mathematical formulas, tables, and lists. The bibliography system is now set up to work seamlessly with Bibby AI's citation management. Feel free to modify it for your needs!

\bibliographystyle{plain}
\bibliography{cite}

\end{document}\documentclass{article}
\usepackage[utf8]{inputenc}
\usepackage{booktabs}
\usepackage{amsmath}
\usepackage{natbib}

\title{My Document}
\author{Your Name}
\date{\today}

\begin{document}

\maketitle

\section{Introduction}
This is a sample document demonstrating various \LaTeX{} features including mathematical equations, tables, and formatting. You can cite papers using the citation search feature, and they will be automatically added to your \texttt{cite.bib} file.

\section{Mathematical Expressions}
Here are some example equations:

\subsection{Quadratic Formula}
The quadratic formula is given by:
\[x = \frac{-b \pm \sqrt{b^2 - 4ac}}{2a}\]

\subsection{Matrix Example}
A sample matrix:
\[
\begin{pmatrix}
a_{11} & a_{12} & a_{13} \\
a_{21} & a_{22} & a_{23} \\
a_{31} & a_{32} & a_{33}
\end{pmatrix}
\]

\section{Tables}
Below is a sample table showing data:

\begin{table}[h]
\centering
\begin{tabular}{@{}lrrr@{}}
\toprule
\textbf{Category} & \textbf{Value 1} & \textbf{Value 2} & \textbf{Total} \\
\midrule
Item A & 10.5 & 20.3 & 30.8 \\
Item B & 15.2 & 25.1 & 40.3 \\
Item C & 22.1 & 12.7 & 34.8 \\
\bottomrule
\end{tabular}
\caption{Sample Data Table}
\label{tab:sample}
\end{table}

\section{Lists}
\subsection{Bullet Points}
\begin{itemize}
  \item First item
  \item Second item
  \item Third item with sub-items:
    \begin{itemize}
      \item Sub-item 1
      \item Sub-item 2
    \end{itemize}
\end{itemize}

\subsection{Numbered List}
\begin{enumerate}
  \item First step
  \item Second step
  \item Third step
\end{enumerate}

\section{Citations and References}
When you use the citation search feature, references will be automatically managed. For example, you might cite a paper like this: \cite{example2024}. All citations will be automatically added to your \texttt{cite.bib} file.

\section{Conclusion}
This template demonstrates basic \LaTeX{} features including mathematical formulas, tables, and lists. The bibliography system is now set up to work seamlessly with Bibby AI's citation management. Feel free to modify it for your needs!

\bibliographystyle{plain}
\bibliography{cite}

\end{document}\documentclass{article}
\usepackage[utf8]{inputenc}
\usepackage{booktabs}
\usepackage{amsmath}
\usepackage{natbib}

\title{My Document}
\author{Your Name}
\date{\today}

\begin{document}

\maketitle

\section{Introduction}
This is a sample document demonstrating various \LaTeX{} features including mathematical equations, tables, and formatting. You can cite papers using the citation search feature, and they will be automatically added to your \texttt{cite.bib} file.

\section{Mathematical Expressions}
Here are some example equations:

\subsection{Quadratic Formula}
The quadratic formula is given by:
\[x = \frac{-b \pm \sqrt{b^2 - 4ac}}{2a}\]

\subsection{Matrix Example}
A sample matrix:
\[
\begin{pmatrix}
a_{11} & a_{12} & a_{13} \\
a_{21} & a_{22} & a_{23} \\
a_{31} & a_{32} & a_{33}
\end{pmatrix}
\]

\section{Tables}
Below is a sample table showing data:

\begin{table}[h]
\centering
\begin{tabular}{@{}lrrr@{}}
\toprule
\textbf{Category} & \textbf{Value 1} & \textbf{Value 2} & \textbf{Total} \\
\midrule
Item A & 10.5 & 20.3 & 30.8 \\
Item B & 15.2 & 25.1 & 40.3 \\
Item C & 22.1 & 12.7 & 34.8 \\
\bottomrule
\end{tabular}
\caption{Sample Data Table}
\label{tab:sample}
\end{table}

\section{Lists}
\subsection{Bullet Points}
\begin{itemize}
  \item First item
  \item Second item
  \item Third item with sub-items:
    \begin{itemize}
      \item Sub-item 1
      \item Sub-item 2
    \end{itemize}
\end{itemize}

\subsection{Numbered List}
\begin{enumerate}
  \item First step
  \item Second step
  \item Third step
\end{enumerate}

\section{Citations and References}
When you use the citation search feature, references will be automatically managed. For example, you might cite a paper like this: \cite{example2024}. All citations will be automatically added to your \texttt{cite.bib} file.

\section{Conclusion}
This template demonstrates basic \LaTeX{} features including mathematical formulas, tables, and lists. The bibliography system is now set up to work seamlessly with Bibby AI's citation management. Feel free to modify it for your needs!

\bibliographystyle{plain}
\bibliography{cite}

\end{document}\documentclass{article}
\usepackage[utf8]{inputenc}
\usepackage{booktabs}
\usepackage{amsmath}
\usepackage{natbib}

\title{My Document}
\author{Your Name}
\date{\today}

\begin{document}

\maketitle

\section{Introduction}
This is a sample document demonstrating various \LaTeX{} features including mathematical equations, tables, and formatting. You can cite papers using the citation search feature, and they will be automatically added to your \texttt{cite.bib} file.

\section{Mathematical Expressions}
Here are some example equations:

\subsection{Quadratic Formula}
The quadratic formula is given by:
\[x = \frac{-b \pm \sqrt{b^2 - 4ac}}{2a}\]

\subsection{Matrix Example}
A sample matrix:
\[
\begin{pmatrix}
a_{11} & a_{12} & a_{13} \\
a_{21} & a_{22} & a_{23} \\
a_{31} & a_{32} & a_{33}
\end{pmatrix}
\]

\section{Tables}
Below is a sample table showing data:

\begin{table}[h]
\centering
\begin{tabular}{@{}lrrr@{}}
\toprule
\textbf{Category} & \textbf{Value 1} & \textbf{Value 2} & \textbf{Total} \\
\midrule
Item A & 10.5 & 20.3 & 30.8 \\
Item B & 15.2 & 25.1 & 40.3 \\
Item C & 22.1 & 12.7 & 34.8 \\
\bottomrule
\end{tabular}
\caption{Sample Data Table}
\label{tab:sample}
\end{table}

\section{Lists}
\subsection{Bullet Points}
\begin{itemize}
  \item First item
  \item Second item
  \item Third item with sub-items:
    \begin{itemize}
      \item Sub-item 1
      \item Sub-item 2
    \end{itemize}
\end{itemize}

\subsection{Numbered List}
\begin{enumerate}
  \item First step
  \item Second step
  \item Third step
\end{enumerate}

\section{Citations and References}
When you use the citation search feature, references will be automatically managed. For example, you might cite a paper like this: \cite{example2024}. All citations will be automatically added to your \texttt{cite.bib} file.

\section{Conclusion}
This template demonstrates basic \LaTeX{} features including mathematical formulas, tables, and lists. The bibliography system is now set up to work seamlessly with Bibby AI's citation management. Feel free to modify it for your needs!

\bibliographystyle{plain}
\bibliography{cite}

\end{document}\documentclass{article}
\usepackage[utf8]{inputenc}
\usepackage{booktabs}
\usepackage{amsmath}
\usepackage{natbib}

\title{My Document}
\author{Your Name}
\date{\today}

\begin{document}

\maketitle

\section{Introduction}
This is a sample document demonstrating various \LaTeX{} features including mathematical equations, tables, and formatting. You can cite papers using the citation search feature, and they will be automatically added to your \texttt{cite.bib} file.

\section{Mathematical Expressions}
Here are some example equations:

\subsection{Quadratic Formula}
The quadratic formula is given by:
\[x = \frac{-b \pm \sqrt{b^2 - 4ac}}{2a}\]

\subsection{Matrix Example}
A sample matrix:
\[
\begin{pmatrix}
a_{11} & a_{12} & a_{13} \\
a_{21} & a_{22} & a_{23} \\
a_{31} & a_{32} & a_{33}
\end{pmatrix}
\]

\section{Tables}
Below is a sample table showing data:

\begin{table}[h]
\centering
\begin{tabular}{@{}lrrr@{}}
\toprule
\textbf{Category} & \textbf{Value 1} & \textbf{Value 2} & \textbf{Total} \\
\midrule
Item A & 10.5 & 20.3 & 30.8 \\
Item B & 15.2 & 25.1 & 40.3 \\
Item C & 22.1 & 12.7 & 34.8 \\
\bottomrule
\end{tabular}
\caption{Sample Data Table}
\label{tab:sample}
\end{table}

\section{Lists}
\subsection{Bullet Points}
\begin{itemize}
  \item First item
  \item Second item
  \item Third item with sub-items:
    \begin{itemize}
      \item Sub-item 1
      \item Sub-item 2
    \end{itemize}
\end{itemize}

\subsection{Numbered List}
\begin{enumerate}
  \item First step
  \item Second step
  \item Third step
\end{enumerate}

\section{Citations and References}
When you use the citation search feature, references will be automatically managed. For example, you might cite a paper like this: \cite{example2024}. All citations will be automatically added to your \texttt{cite.bib} file.

\section{Conclusion}
This template demonstrates basic \LaTeX{} features including mathematical formulas, tables, and lists. The bibliography system is now set up to work seamlessly with Bibby AI's citation management. Feel free to modify it for your needs!

\bibliographystyle{plain}
\bibliography{cite}

\end{document}\documentclass{article}
\usepackage[utf8]{inputenc}
\usepackage{booktabs}
\usepackage{amsmath}
\usepackage{natbib}

\title{My Document}
\author{Your Name}
\date{\today}

\begin{document}

\maketitle

\section{Introduction}
This is a sample document demonstrating various \LaTeX{} features including mathematical equations, tables, and formatting. You can cite papers using the citation search feature, and they will be automatically added to your \texttt{cite.bib} file.

\section{Mathematical Expressions}
Here are some example equations:

\subsection{Quadratic Formula}
The quadratic formula is given by:
\[x = \frac{-b \pm \sqrt{b^2 - 4ac}}{2a}\]

\subsection{Matrix Example}
A sample matrix:
\[
\begin{pmatrix}
a_{11} & a_{12} & a_{13} \\
a_{21} & a_{22} & a_{23} \\
a_{31} & a_{32} & a_{33}
\end{pmatrix}
\]

\section{Tables}
Below is a sample table showing data:

\begin{table}[h]
\centering
\begin{tabular}{@{}lrrr@{}}
\toprule
\textbf{Category} & \textbf{Value 1} & \textbf{Value 2} & \textbf{Total} \\
\midrule
Item A & 10.5 & 20.3 & 30.8 \\
Item B & 15.2 & 25.1 & 40.3 \\
Item C & 22.1 & 12.7 & 34.8 \\
\bottomrule
\end{tabular}
\caption{Sample Data Table}
\label{tab:sample}
\end{table}

\section{Lists}
\subsection{Bullet Points}
\begin{itemize}
  \item First item
  \item Second item
  \item Third item with sub-items:
    \begin{itemize}
      \item Sub-item 1
      \item Sub-item 2
    \end{itemize}
\end{itemize}

\subsection{Numbered List}
\begin{enumerate}
  \item First step
  \item Second step
  \item Third step
\end{enumerate}

\section{Citations and References}
When you use the citation search feature, references will be automatically managed. For example, you might cite a paper like this: \cite{example2024}. All citations will be automatically added to your \texttt{cite.bib} file.

\section{Conclusion}
This template demonstrates basic \LaTeX{} features including mathematical formulas, tables, and lists. The bibliography system is now set up to work seamlessly with Bibby AI's citation management. Feel free to modify it for your needs!

\bibliographystyle{plain}
\bibliography{cite}

\end{document}\documentclass{article}
\usepackage[utf8]{inputenc}
\usepackage{booktabs}
\usepackage{amsmath}
\usepackage{natbib}

\title{My Document}
\author{Your Name}
\date{\today}

\begin{document}

\maketitle

\section{Introduction}
This is a sample document demonstrating various \LaTeX{} features including mathematical equations, tables, and formatting. You can cite papers using the citation search feature, and they will be automatically added to your \texttt{cite.bib} file.

\section{Mathematical Expressions}
Here are some example equations:

\subsection{Quadratic Formula}
The quadratic formula is given by:
\[x = \frac{-b \pm \sqrt{b^2 - 4ac}}{2a}\]

\subsection{Matrix Example}
A sample matrix:
\[
\begin{pmatrix}
a_{11} & a_{12} & a_{13} \\
a_{21} & a_{22} & a_{23} \\
a_{31} & a_{32} & a_{33}
\end{pmatrix}
\]

\section{Tables}
Below is a sample table showing data:

\begin{table}[h]
\centering
\begin{tabular}{@{}lrrr@{}}
\toprule
\textbf{Category} & \textbf{Value 1} & \textbf{Value 2} & \textbf{Total} \\
\midrule
Item A & 10.5 & 20.3 & 30.8 \\
Item B & 15.2 & 25.1 & 40.3 \\
Item C & 22.1 & 12.7 & 34.8 \\
\bottomrule
\end{tabular}
\caption{Sample Data Table}
\label{tab:sample}
\end{table}

\section{Lists}
\subsection{Bullet Points}
\begin{itemize}
  \item First item
  \item Second item
  \item Third item with sub-items:
    \begin{itemize}
      \item Sub-item 1
      \item Sub-item 2
    \end{itemize}
\end{itemize}

\subsection{Numbered List}
\begin{enumerate}
  \item First step
  \item Second step
  \item Third step
\end{enumerate}

\section{Citations and References}
When you use the citation search feature, references will be automatically managed. For example, you might cite a paper like this: \cite{example2024}. All citations will be automatically added to your \texttt{cite.bib} file.

\section{Conclusion}
This template demonstrates basic \LaTeX{} features including mathematical formulas, tables, and lists. The bibliography system is now set up to work seamlessly with Bibby AI's citation management. Feel free to modify it for your needs!

\bibliographystyle{plain}
\bibliography{cite}

\end{document}\documentclass{article}
\usepackage[utf8]{inputenc}
\usepackage{booktabs}
\usepackage{amsmath}
\usepackage{natbib}

\title{My Document}
\author{Your Name}
\date{\today}

\begin{document}

\maketitle

\section{Introduction}
This is a sample document demonstrating various \LaTeX{} features including mathematical equations, tables, and formatting. You can cite papers using the citation search feature, and they will be automatically added to your \texttt{cite.bib} file.

\section{Mathematical Expressions}
Here are some example equations:

\subsection{Quadratic Formula}
The quadratic formula is given by:
\[x = \frac{-b \pm \sqrt{b^2 - 4ac}}{2a}\]

\subsection{Matrix Example}
A sample matrix:
\[
\begin{pmatrix}
a_{11} & a_{12} & a_{13} \\
a_{21} & a_{22} & a_{23} \\
a_{31} & a_{32} & a_{33}
\end{pmatrix}
\]

\section{Tables}
Below is a sample table showing data:

\begin{table}[h]
\centering
\begin{tabular}{@{}lrrr@{}}
\toprule
\textbf{Category} & \textbf{Value 1} & \textbf{Value 2} & \textbf{Total} \\
\midrule
Item A & 10.5 & 20.3 & 30.8 \\
Item B & 15.2 & 25.1 & 40.3 \\
Item C & 22.1 & 12.7 & 34.8 \\
\bottomrule
\end{tabular}
\caption{Sample Data Table}
\label{tab:sample}
\end{table}

\section{Lists}
\subsection{Bullet Points}
\begin{itemize}
  \item First item
  \item Second item
  \item Third item with sub-items:
    \begin{itemize}
      \item Sub-item 1
      \item Sub-item 2
    \end{itemize}
\end{itemize}

\subsection{Numbered List}
\begin{enumerate}
  \item First step
  \item Second step
  \item Third step
\end{enumerate}

\section{Citations and References}
When you use the citation search feature, references will be automatically managed. For example, you might cite a paper like this: \cite{example2024}. All citations will be automatically added to your \texttt{cite.bib} file.

\section{Conclusion}
This template demonstrates basic \LaTeX{} features including mathematical formulas, tables, and lists. The bibliography system is now set up to work seamlessly with Bibby AI's citation management. Feel free to modify it for your needs!

\bibliographystyle{plain}
\bibliography{cite}

\end{document}\documentclass{article}
\usepackage[utf8]{inputenc}
\usepackage{booktabs}
\usepackage{amsmath}
\usepackage{natbib}

\title{My Document}
\author{Your Name}
\date{\today}

\begin{document}

\maketitle

\section{Introduction}
This is a sample document demonstrating various \LaTeX{} features including mathematical equations, tables, and formatting. You can cite papers using the citation search feature, and they will be automatically added to your \texttt{cite.bib} file.

\section{Mathematical Expressions}
Here are some example equations:

\subsection{Quadratic Formula}
The quadratic formula is given by:
\[x = \frac{-b \pm \sqrt{b^2 - 4ac}}{2a}\]

\subsection{Matrix Example}
A sample matrix:
\[
\begin{pmatrix}
a_{11} & a_{12} & a_{13} \\
a_{21} & a_{22} & a_{23} \\
a_{31} & a_{32} & a_{33}
\end{pmatrix}
\]

\section{Tables}
Below is a sample table showing data:

\begin{table}[h]
\centering
\begin{tabular}{@{}lrrr@{}}
\toprule
\textbf{Category} & \textbf{Value 1} & \textbf{Value 2} & \textbf{Total} \\
\midrule
Item A & 10.5 & 20.3 & 30.8 \\
Item B & 15.2 & 25.1 & 40.3 \\
Item C & 22.1 & 12.7 & 34.8 \\
\bottomrule
\end{tabular}
\caption{Sample Data Table}
\label{tab:sample}
\end{table}

\section{Lists}
\subsection{Bullet Points}
\begin{itemize}
  \item First item
  \item Second item
  \item Third item with sub-items:
    \begin{itemize}
      \item Sub-item 1
      \item Sub-item 2
    \end{itemize}
\end{itemize}

\subsection{Numbered List}
\begin{enumerate}
  \item First step
  \item Second step
  \item Third step
\end{enumerate}

\section{Citations and References}
When you use the citation search feature, references will be automatically managed. For example, you might cite a paper like this: \cite{example2024}. All citations will be automatically added to your \texttt{cite.bib} file.

\section{Conclusion}
This template demonstrates basic \LaTeX{} features including mathematical formulas, tables, and lists. The bibliography system is now set up to work seamlessly with Bibby AI's citation management. Feel free to modify it for your needs!

\bibliographystyle{plain}
\bibliography{cite}

\end{document}\documentclass{article}
\usepackage[utf8]{inputenc}
\usepackage{booktabs}
\usepackage{amsmath}
\usepackage{natbib}

\title{My Document}
\author{Your Name}
\date{\today}

\begin{document}

\maketitle

\section{Introduction}
This is a sample document demonstrating various \LaTeX{} features including mathematical equations, tables, and formatting. You can cite papers using the citation search feature, and they will be automatically added to your \texttt{cite.bib} file.

\section{Mathematical Expressions}
Here are some example equations:

\subsection{Quadratic Formula}
The quadratic formula is given by:
\[x = \frac{-b \pm \sqrt{b^2 - 4ac}}{2a}\]

\subsection{Matrix Example}
A sample matrix:
\[
\begin{pmatrix}
a_{11} & a_{12} & a_{13} \\
a_{21} & a_{22} & a_{23} \\
a_{31} & a_{32} & a_{33}
\end{pmatrix}
\]

\section{Tables}
Below is a sample table showing data:

\begin{table}[h]
\centering
\begin{tabular}{@{}lrrr@{}}
\toprule
\textbf{Category} & \textbf{Value 1} & \textbf{Value 2} & \textbf{Total} \\
\midrule
Item A & 10.5 & 20.3 & 30.8 \\
Item B & 15.2 & 25.1 & 40.3 \\
Item C & 22.1 & 12.7 & 34.8 \\
\bottomrule
\end{tabular}
\caption{Sample Data Table}
\label{tab:sample}
\end{table}

\section{Lists}
\subsection{Bullet Points}
\begin{itemize}
  \item First item
  \item Second item
  \item Third item with sub-items:
    \begin{itemize}
      \item Sub-item 1
      \item Sub-item 2
    \end{itemize}
\end{itemize}

\subsection{Numbered List}
\begin{enumerate}
  \item First step
  \item Second step
  \item Third step
\end{enumerate}

\section{Citations and References}
When you use the citation search feature, references will be automatically managed. For example, you might cite a paper like this: \cite{example2024}. All citations will be automatically added to your \texttt{cite.bib} file.

\section{Conclusion}
This template demonstrates basic \LaTeX{} features including mathematical formulas, tables, and lists. The bibliography system is now set up to work seamlessly with Bibby AI's citation management. Feel free to modify it for your needs!

\bibliographystyle{plain}
\bibliography{cite}

\end{document}\documentclass{article}
\usepackage[utf8]{inputenc}
\usepackage{booktabs}
\usepackage{amsmath}
\usepackage{natbib}

\title{My Document}
\author{Your Name}
\date{\today}

\begin{document}

\maketitle

\section{Introduction}
This is a sample document demonstrating various \LaTeX{} features including mathematical equations, tables, and formatting. You can cite papers using the citation search feature, and they will be automatically added to your \texttt{cite.bib} file.

\section{Mathematical Expressions}
Here are some example equations:

\subsection{Quadratic Formula}
The quadratic formula is given by:
\[x = \frac{-b \pm \sqrt{b^2 - 4ac}}{2a}\]

\subsection{Matrix Example}
A sample matrix:
\[
\begin{pmatrix}
a_{11} & a_{12} & a_{13} \\
a_{21} & a_{22} & a_{23} \\
a_{31} & a_{32} & a_{33}
\end{pmatrix}
\]

\section{Tables}
Below is a sample table showing data:

\begin{table}[h]
\centering
\begin{tabular}{@{}lrrr@{}}
\toprule
\textbf{Category} & \textbf{Value 1} & \textbf{Value 2} & \textbf{Total} \\
\midrule
Item A & 10.5 & 20.3 & 30.8 \\
Item B & 15.2 & 25.1 & 40.3 \\
Item C & 22.1 & 12.7 & 34.8 \\
\bottomrule
\end{tabular}
\caption{Sample Data Table}
\label{tab:sample}
\end{table}

\section{Lists}
\subsection{Bullet Points}
\begin{itemize}
  \item First item
  \item Second item
  \item Third item with sub-items:
    \begin{itemize}
      \item Sub-item 1
      \item Sub-item 2
    \end{itemize}
\end{itemize}

\subsection{Numbered List}
\begin{enumerate}
  \item First step
  \item Second step
  \item Third step
\end{enumerate}

\section{Citations and References}
When you use the citation search feature, references will be automatically managed. For example, you might cite a paper like this: \cite{example2024}. All citations will be automatically added to your \texttt{cite.bib} file.

\section{Conclusion}
This template demonstrates basic \LaTeX{} features including mathematical formulas, tables, and lists. The bibliography system is now set up to work seamlessly with Bibby AI's citation management. Feel free to modify it for your needs!

\bibliographystyle{plain}
\bibliography{cite}

\end{document}\documentclass{article}
\usepackage[utf8]{inputenc}
\usepackage{booktabs}
\usepackage{amsmath}
\usepackage{natbib}

\title{My Document}
\author{Your Name}
\date{\today}

\begin{document}

\maketitle

\section{Introduction}
This is a sample document demonstrating various \LaTeX{} features including mathematical equations, tables, and formatting. You can cite papers using the citation search feature, and they will be automatically added to your \texttt{cite.bib} file.

\section{Mathematical Expressions}
Here are some example equations:

\subsection{Quadratic Formula}
The quadratic formula is given by:
\[x = \frac{-b \pm \sqrt{b^2 - 4ac}}{2a}\]

\subsection{Matrix Example}
A sample matrix:
\[
\begin{pmatrix}
a_{11} & a_{12} & a_{13} \\
a_{21} & a_{22} & a_{23} \\
a_{31} & a_{32} & a_{33}
\end{pmatrix}
\]

\section{Tables}
Below is a sample table showing data:

\begin{table}[h]
\centering
\begin{tabular}{@{}lrrr@{}}
\toprule
\textbf{Category} & \textbf{Value 1} & \textbf{Value 2} & \textbf{Total} \\
\midrule
Item A & 10.5 & 20.3 & 30.8 \\
Item B & 15.2 & 25.1 & 40.3 \\
Item C & 22.1 & 12.7 & 34.8 \\
\bottomrule
\end{tabular}
\caption{Sample Data Table}
\label{tab:sample}
\end{table}

\section{Lists}
\subsection{Bullet Points}
\begin{itemize}
  \item First item
  \item Second item
  \item Third item with sub-items:
    \begin{itemize}
      \item Sub-item 1
      \item Sub-item 2
    \end{itemize}
\end{itemize}

\subsection{Numbered List}
\begin{enumerate}
  \item First step
  \item Second step
  \item Third step
\end{enumerate}

\section{Citations and References}
When you use the citation search feature, references will be automatically managed. For example, you might cite a paper like this: \cite{example2024}. All citations will be automatically added to your \texttt{cite.bib} file.

\section{Conclusion}
This template demonstrates basic \LaTeX{} features including mathematical formulas, tables, and lists. The bibliography system is now set up to work seamlessly with Bibby AI's citation management. Feel free to modify it for your needs!

\bibliographystyle{plain}
\bibliography{cite}

\end{document}\documentclass{article}
\usepackage[utf8]{inputenc}
\usepackage{booktabs}
\usepackage{amsmath}
\usepackage{natbib}

\title{My Document}
\author{Your Name}
\date{\today}

\begin{document}

\maketitle

\section{Introduction}
This is a sample document demonstrating various \LaTeX{} features including mathematical equations, tables, and formatting. You can cite papers using the citation search feature, and they will be automatically added to your \texttt{cite.bib} file.

\section{Mathematical Expressions}
Here are some example equations:

\subsection{Quadratic Formula}
The quadratic formula is given by:
\[x = \frac{-b \pm \sqrt{b^2 - 4ac}}{2a}\]

\subsection{Matrix Example}
A sample matrix:
\[
\begin{pmatrix}
a_{11} & a_{12} & a_{13} \\
a_{21} & a_{22} & a_{23} \\
a_{31} & a_{32} & a_{33}
\end{pmatrix}
\]

\section{Tables}
Below is a sample table showing data:

\begin{table}[h]
\centering
\begin{tabular}{@{}lrrr@{}}
\toprule
\textbf{Category} & \textbf{Value 1} & \textbf{Value 2} & \textbf{Total} \\
\midrule
Item A & 10.5 & 20.3 & 30.8 \\
Item B & 15.2 & 25.1 & 40.3 \\
Item C & 22.1 & 12.7 & 34.8 \\
\bottomrule
\end{tabular}
\caption{Sample Data Table}
\label{tab:sample}
\end{table}

\section{Lists}
\subsection{Bullet Points}
\begin{itemize}
  \item First item
  \item Second item
  \item Third item with sub-items:
    \begin{itemize}
      \item Sub-item 1
      \item Sub-item 2
    \end{itemize}
\end{itemize}

\subsection{Numbered List}
\begin{enumerate}
  \item First step
  \item Second step
  \item Third step
\end{enumerate}

\section{Citations and References}
When you use the citation search feature, references will be automatically managed. For example, you might cite a paper like this: \cite{example2024}. All citations will be automatically added to your \texttt{cite.bib} file.

\section{Conclusion}
This template demonstrates basic \LaTeX{} features including mathematical formulas, tables, and lists. The bibliography system is now set up to work seamlessly with Bibby AI's citation management. Feel free to modify it for your needs!

\bibliographystyle{plain}
\bibliography{cite}

\end{document}\documentclass{article}
\usepackage[utf8]{inputenc}
\usepackage{booktabs}
\usepackage{amsmath}
\usepackage{natbib}

\title{My Document}
\author{Your Name}
\date{\today}

\begin{document}

\maketitle

\section{Introduction}
This is a sample document demonstrating various \LaTeX{} features including mathematical equations, tables, and formatting. You can cite papers using the citation search feature, and they will be automatically added to your \texttt{cite.bib} file.

\section{Mathematical Expressions}
Here are some example equations:

\subsection{Quadratic Formula}
The quadratic formula is given by:
\[x = \frac{-b \pm \sqrt{b^2 - 4ac}}{2a}\]

\subsection{Matrix Example}
A sample matrix:
\[
\begin{pmatrix}
a_{11} & a_{12} & a_{13} \\
a_{21} & a_{22} & a_{23} \\
a_{31} & a_{32} & a_{33}
\end{pmatrix}
\]

\section{Tables}
Below is a sample table showing data:

\begin{table}[h]
\centering
\begin{tabular}{@{}lrrr@{}}
\toprule
\textbf{Category} & \textbf{Value 1} & \textbf{Value 2} & \textbf{Total} \\
\midrule
Item A & 10.5 & 20.3 & 30.8 \\
Item B & 15.2 & 25.1 & 40.3 \\
Item C & 22.1 & 12.7 & 34.8 \\
\bottomrule
\end{tabular}
\caption{Sample Data Table}
\label{tab:sample}
\end{table}

\section{Lists}
\subsection{Bullet Points}
\begin{itemize}
  \item First item
  \item Second item
  \item Third item with sub-items:
    \begin{itemize}
      \item Sub-item 1
      \item Sub-item 2
    \end{itemize}
\end{itemize}

\subsection{Numbered List}
\begin{enumerate}
  \item First step
  \item Second step
  \item Third step
\end{enumerate}

\section{Citations and References}
When you use the citation search feature, references will be automatically managed. For example, you might cite a paper like this: \cite{example2024}. All citations will be automatically added to your \texttt{cite.bib} file.

\section{Conclusion}
This template demonstrates basic \LaTeX{} features including mathematical formulas, tables, and lists. The bibliography system is now set up to work seamlessly with Bibby AI's citation management. Feel free to modify it for your needs!

\bibliographystyle{plain}
\bibliography{cite}

\end{document}\documentclass{article}
\usepackage[utf8]{inputenc}
\usepackage{booktabs}
\usepackage{amsmath}
\usepackage{natbib}

\title{My Document}
\author{Your Name}
\date{\today}

\begin{document}

\maketitle

\section{Introduction}
This is a sample document demonstrating various \LaTeX{} features including mathematical equations, tables, and formatting. You can cite papers using the citation search feature, and they will be automatically added to your \texttt{cite.bib} file.

\section{Mathematical Expressions}
Here are some example equations:

\subsection{Quadratic Formula}
The quadratic formula is given by:
\[x = \frac{-b \pm \sqrt{b^2 - 4ac}}{2a}\]

\subsection{Matrix Example}
A sample matrix:
\[
\begin{pmatrix}
a_{11} & a_{12} & a_{13} \\
a_{21} & a_{22} & a_{23} \\
a_{31} & a_{32} & a_{33}
\end{pmatrix}
\]

\section{Tables}
Below is a sample table showing data:

\begin{table}[h]
\centering
\begin{tabular}{@{}lrrr@{}}
\toprule
\textbf{Category} & \textbf{Value 1} & \textbf{Value 2} & \textbf{Total} \\
\midrule
Item A & 10.5 & 20.3 & 30.8 \\
Item B & 15.2 & 25.1 & 40.3 \\
Item C & 22.1 & 12.7 & 34.8 \\
\bottomrule
\end{tabular}
\caption{Sample Data Table}
\label{tab:sample}
\end{table}

\section{Lists}
\subsection{Bullet Points}
\begin{itemize}
  \item First item
  \item Second item
  \item Third item with sub-items:
    \begin{itemize}
      \item Sub-item 1
      \item Sub-item 2
    \end{itemize}
\end{itemize}

\subsection{Numbered List}
\begin{enumerate}
  \item First step
  \item Second step
  \item Third step
\end{enumerate}

\section{Citations and References}
When you use the citation search feature, references will be automatically managed. For example, you might cite a paper like this: \cite{example2024}. All citations will be automatically added to your \texttt{cite.bib} file.

\section{Conclusion}
This template demonstrates basic \LaTeX{} features including mathematical formulas, tables, and lists. The bibliography system is now set up to work seamlessly with Bibby AI's citation management. Feel free to modify it for your needs!

\bibliographystyle{plain}
\bibliography{cite}

\end{document}\documentclass{article}
\usepackage[utf8]{inputenc}
\usepackage{booktabs}
\usepackage{amsmath}
\usepackage{natbib}

\title{My Document}
\author{Your Name}
\date{\today}

\begin{document}

\maketitle

\section{Introduction}
This is a sample document demonstrating various \LaTeX{} features including mathematical equations, tables, and formatting. You can cite papers using the citation search feature, and they will be automatically added to your \texttt{cite.bib} file.

\section{Mathematical Expressions}
Here are some example equations:

\subsection{Quadratic Formula}
The quadratic formula is given by:
\[x = \frac{-b \pm \sqrt{b^2 - 4ac}}{2a}\]

\subsection{Matrix Example}
A sample matrix:
\[
\begin{pmatrix}
a_{11} & a_{12} & a_{13} \\
a_{21} & a_{22} & a_{23} \\
a_{31} & a_{32} & a_{33}
\end{pmatrix}
\]

\section{Tables}
Below is a sample table showing data:

\begin{table}[h]
\centering
\begin{tabular}{@{}lrrr@{}}
\toprule
\textbf{Category} & \textbf{Value 1} & \textbf{Value 2} & \textbf{Total} \\
\midrule
Item A & 10.5 & 20.3 & 30.8 \\
Item B & 15.2 & 25.1 & 40.3 \\
Item C & 22.1 & 12.7 & 34.8 \\
\bottomrule
\end{tabular}
\caption{Sample Data Table}
\label{tab:sample}
\end{table}

\section{Lists}
\subsection{Bullet Points}
\begin{itemize}
  \item First item
  \item Second item
  \item Third item with sub-items:
    \begin{itemize}
      \item Sub-item 1
      \item Sub-item 2
    \end{itemize}
\end{itemize}

\subsection{Numbered List}
\begin{enumerate}
  \item First step
  \item Second step
  \item Third step
\end{enumerate}

\section{Citations and References}
When you use the citation search feature, references will be automatically managed. For example, you might cite a paper like this: \cite{example2024}. All citations will be automatically added to your \texttt{cite.bib} file.

\section{Conclusion}
This template demonstrates basic \LaTeX{} features including mathematical formulas, tables, and lists. The bibliography system is now set up to work seamlessly with Bibby AI's citation management. Feel free to modify it for your needs!

\bibliographystyle{plain}
\bibliography{cite}

\end{document}\documentclass{article}
\usepackage[utf8]{inputenc}
\usepackage{booktabs}
\usepackage{amsmath}
\usepackage{natbib}

\title{My Document}
\author{Your Name}
\date{\today}

\begin{document}

\maketitle

\section{Introduction}
This is a sample document demonstrating various \LaTeX{} features including mathematical equations, tables, and formatting. You can cite papers using the citation search feature, and they will be automatically added to your \texttt{cite.bib} file.

\section{Mathematical Expressions}
Here are some example equations:

\subsection{Quadratic Formula}
The quadratic formula is given by:
\[x = \frac{-b \pm \sqrt{b^2 - 4ac}}{2a}\]

\subsection{Matrix Example}
A sample matrix:
\[
\begin{pmatrix}
a_{11} & a_{12} & a_{13} \\
a_{21} & a_{22} & a_{23} \\
a_{31} & a_{32} & a_{33}
\end{pmatrix}
\]

\section{Tables}
Below is a sample table showing data:

\begin{table}[h]
\centering
\begin{tabular}{@{}lrrr@{}}
\toprule
\textbf{Category} & \textbf{Value 1} & \textbf{Value 2} & \textbf{Total} \\
\midrule
Item A & 10.5 & 20.3 & 30.8 \\
Item B & 15.2 & 25.1 & 40.3 \\
Item C & 22.1 & 12.7 & 34.8 \\
\bottomrule
\end{tabular}
\caption{Sample Data Table}
\label{tab:sample}
\end{table}

\section{Lists}
\subsection{Bullet Points}
\begin{itemize}
  \item First item
  \item Second item
  \item Third item with sub-items:
    \begin{itemize}
      \item Sub-item 1
      \item Sub-item 2
    \end{itemize}
\end{itemize}

\subsection{Numbered List}
\begin{enumerate}
  \item First step
  \item Second step
  \item Third step
\end{enumerate}

\section{Citations and References}
When you use the citation search feature, references will be automatically managed. For example, you might cite a paper like this: \cite{example2024}. All citations will be automatically added to your \texttt{cite.bib} file.

\section{Conclusion}
This template demonstrates basic \LaTeX{} features including mathematical formulas, tables, and lists. The bibliography system is now set up to work seamlessly with Bibby AI's citation management. Feel free to modify it for your needs!

\bibliographystyle{plain}
\bibliography{cite}

\end{document}\documentclass{article}
\usepackage[utf8]{inputenc}
\usepackage{booktabs}
\usepackage{amsmath}
\usepackage{natbib}

\title{My Document}
\author{Your Name}
\date{\today}

\begin{document}

\maketitle

\section{Introduction}
This is a sample document demonstrating various \LaTeX{} features including mathematical equations, tables, and formatting. You can cite papers using the citation search feature, and they will be automatically added to your \texttt{cite.bib} file.

\section{Mathematical Expressions}
Here are some example equations:

\subsection{Quadratic Formula}
The quadratic formula is given by:
\[x = \frac{-b \pm \sqrt{b^2 - 4ac}}{2a}\]

\subsection{Matrix Example}
A sample matrix:
\[
\begin{pmatrix}
a_{11} & a_{12} & a_{13} \\
a_{21} & a_{22} & a_{23} \\
a_{31} & a_{32} & a_{33}
\end{pmatrix}
\]

\section{Tables}
Below is a sample table showing data:

\begin{table}[h]
\centering
\begin{tabular}{@{}lrrr@{}}
\toprule
\textbf{Category} & \textbf{Value 1} & \textbf{Value 2} & \textbf{Total} \\
\midrule
Item A & 10.5 & 20.3 & 30.8 \\
Item B & 15.2 & 25.1 & 40.3 \\
Item C & 22.1 & 12.7 & 34.8 \\
\bottomrule
\end{tabular}
\caption{Sample Data Table}
\label{tab:sample}
\end{table}

\section{Lists}
\subsection{Bullet Points}
\begin{itemize}
  \item First item
  \item Second item
  \item Third item with sub-items:
    \begin{itemize}
      \item Sub-item 1
      \item Sub-item 2
    \end{itemize}
\end{itemize}

\subsection{Numbered List}
\begin{enumerate}
  \item First step
  \item Second step
  \item Third step
\end{enumerate}

\section{Citations and References}
When you use the citation search feature, references will be automatically managed. For example, you might cite a paper like this: \cite{example2024}. All citations will be automatically added to your \texttt{cite.bib} file.

\section{Conclusion}
This template demonstrates basic \LaTeX{} features including mathematical formulas, tables, and lists. The bibliography system is now set up to work seamlessly with Bibby AI's citation management. Feel free to modify it for your needs!

\bibliographystyle{plain}
\bibliography{cite}

\end{document}\documentclass{article}
\usepackage[utf8]{inputenc}
\usepackage{booktabs}
\usepackage{amsmath}
\usepackage{natbib}

\title{My Document}
\author{Your Name}
\date{\today}

\begin{document}

\maketitle

\section{Introduction}
This is a sample document demonstrating various \LaTeX{} features including mathematical equations, tables, and formatting. You can cite papers using the citation search feature, and they will be automatically added to your \texttt{cite.bib} file.

\section{Mathematical Expressions}
Here are some example equations:

\subsection{Quadratic Formula}
The quadratic formula is given by:
\[x = \frac{-b \pm \sqrt{b^2 - 4ac}}{2a}\]

\subsection{Matrix Example}
A sample matrix:
\[
\begin{pmatrix}
a_{11} & a_{12} & a_{13} \\
a_{21} & a_{22} & a_{23} \\
a_{31} & a_{32} & a_{33}
\end{pmatrix}
\]

\section{Tables}
Below is a sample table showing data:

\begin{table}[h]
\centering
\begin{tabular}{@{}lrrr@{}}
\toprule
\textbf{Category} & \textbf{Value 1} & \textbf{Value 2} & \textbf{Total} \\
\midrule
Item A & 10.5 & 20.3 & 30.8 \\
Item B & 15.2 & 25.1 & 40.3 \\
Item C & 22.1 & 12.7 & 34.8 \\
\bottomrule
\end{tabular}
\caption{Sample Data Table}
\label{tab:sample}
\end{table}

\section{Lists}
\subsection{Bullet Points}
\begin{itemize}
  \item First item
  \item Second item
  \item Third item with sub-items:
    \begin{itemize}
      \item Sub-item 1
      \item Sub-item 2
    \end{itemize}
\end{itemize}

\subsection{Numbered List}
\begin{enumerate}
  \item First step
  \item Second step
  \item Third step
\end{enumerate}

\section{Citations and References}
When you use the citation search feature, references will be automatically managed. For example, you might cite a paper like this: \cite{example2024}. All citations will be automatically added to your \texttt{cite.bib} file.

\section{Conclusion}
This template demonstrates basic \LaTeX{} features including mathematical formulas, tables, and lists. The bibliography system is now set up to work seamlessly with Bibby AI's citation management. Feel free to modify it for your needs!

\bibliographystyle{plain}
\bibliography{cite}

\end{document}\documentclass{article}
\usepackage[utf8]{inputenc}
\usepackage{booktabs}
\usepackage{amsmath}
\usepackage{natbib}

\title{My Document}
\author{Your Name}
\date{\today}

\begin{document}

\maketitle

\section{Introduction}
This is a sample document demonstrating various \LaTeX{} features including mathematical equations, tables, and formatting. You can cite papers using the citation search feature, and they will be automatically added to your \texttt{cite.bib} file.

\section{Mathematical Expressions}
Here are some example equations:

\subsection{Quadratic Formula}
The quadratic formula is given by:
\[x = \frac{-b \pm \sqrt{b^2 - 4ac}}{2a}\]

\subsection{Matrix Example}
A sample matrix:
\[
\begin{pmatrix}
a_{11} & a_{12} & a_{13} \\
a_{21} & a_{22} & a_{23} \\
a_{31} & a_{32} & a_{33}
\end{pmatrix}
\]

\section{Tables}
Below is a sample table showing data:

\begin{table}[h]
\centering
\begin{tabular}{@{}lrrr@{}}
\toprule
\textbf{Category} & \textbf{Value 1} & \textbf{Value 2} & \textbf{Total} \\
\midrule
Item A & 10.5 & 20.3 & 30.8 \\
Item B & 15.2 & 25.1 & 40.3 \\
Item C & 22.1 & 12.7 & 34.8 \\
\bottomrule
\end{tabular}
\caption{Sample Data Table}
\label{tab:sample}
\end{table}

\section{Lists}
\subsection{Bullet Points}
\begin{itemize}
  \item First item
  \item Second item
  \item Third item with sub-items:
    \begin{itemize}
      \item Sub-item 1
      \item Sub-item 2
    \end{itemize}
\end{itemize}

\subsection{Numbered List}
\begin{enumerate}
  \item First step
  \item Second step
  \item Third step
\end{enumerate}

\section{Citations and References}
When you use the citation search feature, references will be automatically managed. For example, you might cite a paper like this: \cite{example2024}. All citations will be automatically added to your \texttt{cite.bib} file.

\section{Conclusion}
This template demonstrates basic \LaTeX{} features including mathematical formulas, tables, and lists. The bibliography system is now set up to work seamlessly with Bibby AI's citation management. Feel free to modify it for your needs!

\bibliographystyle{plain}
\bibliography{cite}

\end{document}\documentclass{article}
\usepackage[utf8]{inputenc}
\usepackage{booktabs}
\usepackage{amsmath}
\usepackage{natbib}

\title{My Document}
\author{Your Name}
\date{\today}

\begin{document}

\maketitle

\section{Introduction}
This is a sample document demonstrating various \LaTeX{} features including mathematical equations, tables, and formatting. You can cite papers using the citation search feature, and they will be automatically added to your \texttt{cite.bib} file.

\section{Mathematical Expressions}
Here are some example equations:

\subsection{Quadratic Formula}
The quadratic formula is given by:
\[x = \frac{-b \pm \sqrt{b^2 - 4ac}}{2a}\]

\subsection{Matrix Example}
A sample matrix:
\[
\begin{pmatrix}
a_{11} & a_{12} & a_{13} \\
a_{21} & a_{22} & a_{23} \\
a_{31} & a_{32} & a_{33}
\end{pmatrix}
\]

\section{Tables}
Below is a sample table showing data:

\begin{table}[h]
\centering
\begin{tabular}{@{}lrrr@{}}
\toprule
\textbf{Category} & \textbf{Value 1} & \textbf{Value 2} & \textbf{Total} \\
\midrule
Item A & 10.5 & 20.3 & 30.8 \\
Item B & 15.2 & 25.1 & 40.3 \\
Item C & 22.1 & 12.7 & 34.8 \\
\bottomrule
\end{tabular}
\caption{Sample Data Table}
\label{tab:sample}
\end{table}

\section{Lists}
\subsection{Bullet Points}
\begin{itemize}
  \item First item
  \item Second item
  \item Third item with sub-items:
    \begin{itemize}
      \item Sub-item 1
      \item Sub-item 2
    \end{itemize}
\end{itemize}

\subsection{Numbered List}
\begin{enumerate}
  \item First step
  \item Second step
  \item Third step
\end{enumerate}

\section{Citations and References}
When you use the citation search feature, references will be automatically managed. For example, you might cite a paper like this: \cite{example2024}. All citations will be automatically added to your \texttt{cite.bib} file.

\section{Conclusion}
This template demonstrates basic \LaTeX{} features including mathematical formulas, tables, and lists. The bibliography system is now set up to work seamlessly with Bibby AI's citation management. Feel free to modify it for your needs!

\bibliographystyle{plain}
\bibliography{cite}

\end{document}\documentclass{article}
\usepackage[utf8]{inputenc}
\usepackage{booktabs}
\usepackage{amsmath}
\usepackage{natbib}

\title{My Document}
\author{Your Name}
\date{\today}

\begin{document}

\maketitle

\section{Introduction}
This is a sample document demonstrating various \LaTeX{} features including mathematical equations, tables, and formatting. You can cite papers using the citation search feature, and they will be automatically added to your \texttt{cite.bib} file.

\section{Mathematical Expressions}
Here are some example equations:

\subsection{Quadratic Formula}
The quadratic formula is given by:
\[x = \frac{-b \pm \sqrt{b^2 - 4ac}}{2a}\]

\subsection{Matrix Example}
A sample matrix:
\[
\begin{pmatrix}
a_{11} & a_{12} & a_{13} \\
a_{21} & a_{22} & a_{23} \\
a_{31} & a_{32} & a_{33}
\end{pmatrix}
\]

\section{Tables}
Below is a sample table showing data:

\begin{table}[h]
\centering
\begin{tabular}{@{}lrrr@{}}
\toprule
\textbf{Category} & \textbf{Value 1} & \textbf{Value 2} & \textbf{Total} \\
\midrule
Item A & 10.5 & 20.3 & 30.8 \\
Item B & 15.2 & 25.1 & 40.3 \\
Item C & 22.1 & 12.7 & 34.8 \\
\bottomrule
\end{tabular}
\caption{Sample Data Table}
\label{tab:sample}
\end{table}

\section{Lists}
\subsection{Bullet Points}
\begin{itemize}
  \item First item
  \item Second item
  \item Third item with sub-items:
    \begin{itemize}
      \item Sub-item 1
      \item Sub-item 2
    \end{itemize}
\end{itemize}

\subsection{Numbered List}
\begin{enumerate}
  \item First step
  \item Second step
  \item Third step
\end{enumerate}

\section{Citations and References}
When you use the citation search feature, references will be automatically managed. For example, you might cite a paper like this: \cite{example2024}. All citations will be automatically added to your \texttt{cite.bib} file.

\section{Conclusion}
This template demonstrates basic \LaTeX{} features including mathematical formulas, tables, and lists. The bibliography system is now set up to work seamlessly with Bibby AI's citation management. Feel free to modify it for your needs!

\bibliographystyle{plain}
\bibliography{cite}

\end{document}\documentclass{article}
\usepackage[utf8]{inputenc}
\usepackage{booktabs}
\usepackage{amsmath}
\usepackage{natbib}

\title{My Document}
\author{Your Name}
\date{\today}

\begin{document}

\maketitle

\section{Introduction}
This is a sample document demonstrating various \LaTeX{} features including mathematical equations, tables, and formatting. You can cite papers using the citation search feature, and they will be automatically added to your \texttt{cite.bib} file.

\section{Mathematical Expressions}
Here are some example equations:

\subsection{Quadratic Formula}
The quadratic formula is given by:
\[x = \frac{-b \pm \sqrt{b^2 - 4ac}}{2a}\]

\subsection{Matrix Example}
A sample matrix:
\[
\begin{pmatrix}
a_{11} & a_{12} & a_{13} \\
a_{21} & a_{22} & a_{23} \\
a_{31} & a_{32} & a_{33}
\end{pmatrix}
\]

\section{Tables}
Below is a sample table showing data:

\begin{table}[h]
\centering
\begin{tabular}{@{}lrrr@{}}
\toprule
\textbf{Category} & \textbf{Value 1} & \textbf{Value 2} & \textbf{Total} \\
\midrule
Item A & 10.5 & 20.3 & 30.8 \\
Item B & 15.2 & 25.1 & 40.3 \\
Item C & 22.1 & 12.7 & 34.8 \\
\bottomrule
\end{tabular}
\caption{Sample Data Table}
\label{tab:sample}
\end{table}

\section{Lists}
\subsection{Bullet Points}
\begin{itemize}
  \item First item
  \item Second item
  \item Third item with sub-items:
    \begin{itemize}
      \item Sub-item 1
      \item Sub-item 2
    \end{itemize}
\end{itemize}

\subsection{Numbered List}
\begin{enumerate}
  \item First step
  \item Second step
  \item Third step
\end{enumerate}

\section{Citations and References}
When you use the citation search feature, references will be automatically managed. For example, you might cite a paper like this: \cite{example2024}. All citations will be automatically added to your \texttt{cite.bib} file.

\section{Conclusion}
This template demonstrates basic \LaTeX{} features including mathematical formulas, tables, and lists. The bibliography system is now set up to work seamlessly with Bibby AI's citation management. Feel free to modify it for your needs!

\bibliographystyle{plain}
\bibliography{cite}

\end{document}\documentclass{article}
\usepackage[utf8]{inputenc}
\usepackage{booktabs}
\usepackage{amsmath}
\usepackage{natbib}

\title{My Document}
\author{Your Name}
\date{\today}

\begin{document}

\maketitle

\section{Introduction}
This is a sample document demonstrating various \LaTeX{} features including mathematical equations, tables, and formatting. You can cite papers using the citation search feature, and they will be automatically added to your \texttt{cite.bib} file.

\section{Mathematical Expressions}
Here are some example equations:

\subsection{Quadratic Formula}
The quadratic formula is given by:
\[x = \frac{-b \pm \sqrt{b^2 - 4ac}}{2a}\]

\subsection{Matrix Example}
A sample matrix:
\[
\begin{pmatrix}
a_{11} & a_{12} & a_{13} \\
a_{21} & a_{22} & a_{23} \\
a_{31} & a_{32} & a_{33}
\end{pmatrix}
\]

\section{Tables}
Below is a sample table showing data:

\begin{table}[h]
\centering
\begin{tabular}{@{}lrrr@{}}
\toprule
\textbf{Category} & \textbf{Value 1} & \textbf{Value 2} & \textbf{Total} \\
\midrule
Item A & 10.5 & 20.3 & 30.8 \\
Item B & 15.2 & 25.1 & 40.3 \\
Item C & 22.1 & 12.7 & 34.8 \\
\bottomrule
\end{tabular}
\caption{Sample Data Table}
\label{tab:sample}
\end{table}

\section{Lists}
\subsection{Bullet Points}
\begin{itemize}
  \item First item
  \item Second item
  \item Third item with sub-items:
    \begin{itemize}
      \item Sub-item 1
      \item Sub-item 2
    \end{itemize}
\end{itemize}

\subsection{Numbered List}
\begin{enumerate}
  \item First step
  \item Second step
  \item Third step
\end{enumerate}

\section{Citations and References}
When you use the citation search feature, references will be automatically managed. For example, you might cite a paper like this: \cite{example2024}. All citations will be automatically added to your \texttt{cite.bib} file.

\section{Conclusion}
This template demonstrates basic \LaTeX{} features including mathematical formulas, tables, and lists. The bibliography system is now set up to work seamlessly with Bibby AI's citation management. Feel free to modify it for your needs!

\bibliographystyle{plain}
\bibliography{cite}

\end{document}\documentclass{article}
\usepackage[utf8]{inputenc}
\usepackage{booktabs}
\usepackage{amsmath}
\usepackage{natbib}

\title{My Document}
\author{Your Name}
\date{\today}

\begin{document}

\maketitle

\section{Introduction}
This is a sample document demonstrating various \LaTeX{} features including mathematical equations, tables, and formatting. You can cite papers using the citation search feature, and they will be automatically added to your \texttt{cite.bib} file.

\section{Mathematical Expressions}
Here are some example equations:

\subsection{Quadratic Formula}
The quadratic formula is given by:
\[x = \frac{-b \pm \sqrt{b^2 - 4ac}}{2a}\]

\subsection{Matrix Example}
A sample matrix:
\[
\begin{pmatrix}
a_{11} & a_{12} & a_{13} \\
a_{21} & a_{22} & a_{23} \\
a_{31} & a_{32} & a_{33}
\end{pmatrix}
\]

\section{Tables}
Below is a sample table showing data:

\begin{table}[h]
\centering
\begin{tabular}{@{}lrrr@{}}
\toprule
\textbf{Category} & \textbf{Value 1} & \textbf{Value 2} & \textbf{Total} \\
\midrule
Item A & 10.5 & 20.3 & 30.8 \\
Item B & 15.2 & 25.1 & 40.3 \\
Item C & 22.1 & 12.7 & 34.8 \\
\bottomrule
\end{tabular}
\caption{Sample Data Table}
\label{tab:sample}
\end{table}

\section{Lists}
\subsection{Bullet Points}
\begin{itemize}
  \item First item
  \item Second item
  \item Third item with sub-items:
    \begin{itemize}
      \item Sub-item 1
      \item Sub-item 2
    \end{itemize}
\end{itemize}

\subsection{Numbered List}
\begin{enumerate}
  \item First step
  \item Second step
  \item Third step
\end{enumerate}

\section{Citations and References}
When you use the citation search feature, references will be automatically managed. For example, you might cite a paper like this: \cite{example2024}. All citations will be automatically added to your \texttt{cite.bib} file.

\section{Conclusion}
This template demonstrates basic \LaTeX{} features including mathematical formulas, tables, and lists. The bibliography system is now set up to work seamlessly with Bibby AI's citation management. Feel free to modify it for your needs!

\bibliographystyle{plain}
\bibliography{cite}

\end{document}\documentclass{article}
\usepackage[utf8]{inputenc}
\usepackage{booktabs}
\usepackage{amsmath}
\usepackage{natbib}

\title{My Document}
\author{Your Name}
\date{\today}

\begin{document}

\maketitle

\section{Introduction}
This is a sample document demonstrating various \LaTeX{} features including mathematical equations, tables, and formatting. You can cite papers using the citation search feature, and they will be automatically added to your \texttt{cite.bib} file.

\section{Mathematical Expressions}
Here are some example equations:

\subsection{Quadratic Formula}
The quadratic formula is given by:
\[x = \frac{-b \pm \sqrt{b^2 - 4ac}}{2a}\]

\subsection{Matrix Example}
A sample matrix:
\[
\begin{pmatrix}
a_{11} & a_{12} & a_{13} \\
a_{21} & a_{22} & a_{23} \\
a_{31} & a_{32} & a_{33}
\end{pmatrix}
\]

\section{Tables}
Below is a sample table showing data:

\begin{table}[h]
\centering
\begin{tabular}{@{}lrrr@{}}
\toprule
\textbf{Category} & \textbf{Value 1} & \textbf{Value 2} & \textbf{Total} \\
\midrule
Item A & 10.5 & 20.3 & 30.8 \\
Item B & 15.2 & 25.1 & 40.3 \\
Item C & 22.1 & 12.7 & 34.8 \\
\bottomrule
\end{tabular}
\caption{Sample Data Table}
\label{tab:sample}
\end{table}

\section{Lists}
\subsection{Bullet Points}
\begin{itemize}
  \item First item
  \item Second item
  \item Third item with sub-items:
    \begin{itemize}
      \item Sub-item 1
      \item Sub-item 2
    \end{itemize}
\end{itemize}

\subsection{Numbered List}
\begin{enumerate}
  \item First step
  \item Second step
  \item Third step
\end{enumerate}

\section{Citations and References}
When you use the citation search feature, references will be automatically managed. For example, you might cite a paper like this: \cite{example2024}. All citations will be automatically added to your \texttt{cite.bib} file.

\section{Conclusion}
This template demonstrates basic \LaTeX{} features including mathematical formulas, tables, and lists. The bibliography system is now set up to work seamlessly with Bibby AI's citation management. Feel free to modify it for your needs!

\bibliographystyle{plain}
\bibliography{cite}

\end{document}\documentclass{article}
\usepackage[utf8]{inputenc}
\usepackage{booktabs}
\usepackage{amsmath}
\usepackage{natbib}

\title{My Document}
\author{Your Name}
\date{\today}

\begin{document}

\maketitle

\section{Introduction}
This is a sample document demonstrating various \LaTeX{} features including mathematical equations, tables, and formatting. You can cite papers using the citation search feature, and they will be automatically added to your \texttt{cite.bib} file.

\section{Mathematical Expressions}
Here are some example equations:

\subsection{Quadratic Formula}
The quadratic formula is given by:
\[x = \frac{-b \pm \sqrt{b^2 - 4ac}}{2a}\]

\subsection{Matrix Example}
A sample matrix:
\[
\begin{pmatrix}
a_{11} & a_{12} & a_{13} \\
a_{21} & a_{22} & a_{23} \\
a_{31} & a_{32} & a_{33}
\end{pmatrix}
\]

\section{Tables}
Below is a sample table showing data:

\begin{table}[h]
\centering
\begin{tabular}{@{}lrrr@{}}
\toprule
\textbf{Category} & \textbf{Value 1} & \textbf{Value 2} & \textbf{Total} \\
\midrule
Item A & 10.5 & 20.3 & 30.8 \\
Item B & 15.2 & 25.1 & 40.3 \\
Item C & 22.1 & 12.7 & 34.8 \\
\bottomrule
\end{tabular}
\caption{Sample Data Table}
\label{tab:sample}
\end{table}

\section{Lists}
\subsection{Bullet Points}
\begin{itemize}
  \item First item
  \item Second item
  \item Third item with sub-items:
    \begin{itemize}
      \item Sub-item 1
      \item Sub-item 2
    \end{itemize}
\end{itemize}

\subsection{Numbered List}
\begin{enumerate}
  \item First step
  \item Second step
  \item Third step
\end{enumerate}

\section{Citations and References}
When you use the citation search feature, references will be automatically managed. For example, you might cite a paper like this: \cite{example2024}. All citations will be automatically added to your \texttt{cite.bib} file.

\section{Conclusion}
This template demonstrates basic \LaTeX{} features including mathematical formulas, tables, and lists. The bibliography system is now set up to work seamlessly with Bibby AI's citation management. Feel free to modify it for your needs!

\bibliographystyle{plain}
\bibliography{cite}

\end{document}\documentclass{article}
\usepackage[utf8]{inputenc}
\usepackage{booktabs}
\usepackage{amsmath}
\usepackage{natbib}

\title{My Document}
\author{Your Name}
\date{\today}

\begin{document}

\maketitle

\section{Introduction}
This is a sample document demonstrating various \LaTeX{} features including mathematical equations, tables, and formatting. You can cite papers using the citation search feature, and they will be automatically added to your \texttt{cite.bib} file.

\section{Mathematical Expressions}
Here are some example equations:

\subsection{Quadratic Formula}
The quadratic formula is given by:
\[x = \frac{-b \pm \sqrt{b^2 - 4ac}}{2a}\]

\subsection{Matrix Example}
A sample matrix:
\[
\begin{pmatrix}
a_{11} & a_{12} & a_{13} \\
a_{21} & a_{22} & a_{23} \\
a_{31} & a_{32} & a_{33}
\end{pmatrix}
\]

\section{Tables}
Below is a sample table showing data:

\begin{table}[h]
\centering
\begin{tabular}{@{}lrrr@{}}
\toprule
\textbf{Category} & \textbf{Value 1} & \textbf{Value 2} & \textbf{Total} \\
\midrule
Item A & 10.5 & 20.3 & 30.8 \\
Item B & 15.2 & 25.1 & 40.3 \\
Item C & 22.1 & 12.7 & 34.8 \\
\bottomrule
\end{tabular}
\caption{Sample Data Table}
\label{tab:sample}
\end{table}

\section{Lists}
\subsection{Bullet Points}
\begin{itemize}
  \item First item
  \item Second item
  \item Third item with sub-items:
    \begin{itemize}
      \item Sub-item 1
      \item Sub-item 2
    \end{itemize}
\end{itemize}

\subsection{Numbered List}
\begin{enumerate}
  \item First step
  \item Second step
  \item Third step
\end{enumerate}

\section{Citations and References}
When you use the citation search feature, references will be automatically managed. For example, you might cite a paper like this: \cite{example2024}. All citations will be automatically added to your \texttt{cite.bib} file.

\section{Conclusion}
This template demonstrates basic \LaTeX{} features including mathematical formulas, tables, and lists. The bibliography system is now set up to work seamlessly with Bibby AI's citation management. Feel free to modify it for your needs!

\bibliographystyle{plain}
\bibliography{cite}

\end{document}\documentclass{article}
\usepackage[utf8]{inputenc}
\usepackage{booktabs}
\usepackage{amsmath}
\usepackage{natbib}

\title{My Document}
\author{Your Name}
\date{\today}

\begin{document}

\maketitle

\section{Introduction}
This is a sample document demonstrating various \LaTeX{} features including mathematical equations, tables, and formatting. You can cite papers using the citation search feature, and they will be automatically added to your \texttt{cite.bib} file.

\section{Mathematical Expressions}
Here are some example equations:

\subsection{Quadratic Formula}
The quadratic formula is given by:
\[x = \frac{-b \pm \sqrt{b^2 - 4ac}}{2a}\]

\subsection{Matrix Example}
A sample matrix:
\[
\begin{pmatrix}
a_{11} & a_{12} & a_{13} \\
a_{21} & a_{22} & a_{23} \\
a_{31} & a_{32} & a_{33}
\end{pmatrix}
\]

\section{Tables}
Below is a sample table showing data:

\begin{table}[h]
\centering
\begin{tabular}{@{}lrrr@{}}
\toprule
\textbf{Category} & \textbf{Value 1} & \textbf{Value 2} & \textbf{Total} \\
\midrule
Item A & 10.5 & 20.3 & 30.8 \\
Item B & 15.2 & 25.1 & 40.3 \\
Item C & 22.1 & 12.7 & 34.8 \\
\bottomrule
\end{tabular}
\caption{Sample Data Table}
\label{tab:sample}
\end{table}

\section{Lists}
\subsection{Bullet Points}
\begin{itemize}
  \item First item
  \item Second item
  \item Third item with sub-items:
    \begin{itemize}
      \item Sub-item 1
      \item Sub-item 2
    \end{itemize}
\end{itemize}

\subsection{Numbered List}
\begin{enumerate}
  \item First step
  \item Second step
  \item Third step
\end{enumerate}

\section{Citations and References}
When you use the citation search feature, references will be automatically managed. For example, you might cite a paper like this: \cite{example2024}. All citations will be automatically added to your \texttt{cite.bib} file.

\section{Conclusion}
This template demonstrates basic \LaTeX{} features including mathematical formulas, tables, and lists. The bibliography system is now set up to work seamlessly with Bibby AI's citation management. Feel free to modify it for your needs!

\bibliographystyle{plain}
\bibliography{cite}

\end{document}\documentclass{article}
\usepackage[utf8]{inputenc}
\usepackage{booktabs}
\usepackage{amsmath}
\usepackage{natbib}

\title{My Document}
\author{Your Name}
\date{\today}

\begin{document}

\maketitle

\section{Introduction}
This is a sample document demonstrating various \LaTeX{} features including mathematical equations, tables, and formatting. You can cite papers using the citation search feature, and they will be automatically added to your \texttt{cite.bib} file.

\section{Mathematical Expressions}
Here are some example equations:

\subsection{Quadratic Formula}
The quadratic formula is given by:
\[x = \frac{-b \pm \sqrt{b^2 - 4ac}}{2a}\]

\subsection{Matrix Example}
A sample matrix:
\[
\begin{pmatrix}
a_{11} & a_{12} & a_{13} \\
a_{21} & a_{22} & a_{23} \\
a_{31} & a_{32} & a_{33}
\end{pmatrix}
\]

\section{Tables}
Below is a sample table showing data:

\begin{table}[h]
\centering
\begin{tabular}{@{}lrrr@{}}
\toprule
\textbf{Category} & \textbf{Value 1} & \textbf{Value 2} & \textbf{Total} \\
\midrule
Item A & 10.5 & 20.3 & 30.8 \\
Item B & 15.2 & 25.1 & 40.3 \\
Item C & 22.1 & 12.7 & 34.8 \\
\bottomrule
\end{tabular}
\caption{Sample Data Table}
\label{tab:sample}
\end{table}

\section{Lists}
\subsection{Bullet Points}
\begin{itemize}
  \item First item
  \item Second item
  \item Third item with sub-items:
    \begin{itemize}
      \item Sub-item 1
      \item Sub-item 2
    \end{itemize}
\end{itemize}

\subsection{Numbered List}
\begin{enumerate}
  \item First step
  \item Second step
  \item Third step
\end{enumerate}

\section{Citations and References}
When you use the citation search feature, references will be automatically managed. For example, you might cite a paper like this: \cite{example2024}. All citations will be automatically added to your \texttt{cite.bib} file.

\section{Conclusion}
This template demonstrates basic \LaTeX{} features including mathematical formulas, tables, and lists. The bibliography system is now set up to work seamlessly with Bibby AI's citation management. Feel free to modify it for your needs!

\bibliographystyle{plain}
\bibliography{cite}

\end{document}\documentclass{article}
\usepackage[utf8]{inputenc}
\usepackage{booktabs}
\usepackage{amsmath}
\usepackage{natbib}

\title{My Document}
\author{Your Name}
\date{\today}

\begin{document}

\maketitle

\section{Introduction}
This is a sample document demonstrating various \LaTeX{} features including mathematical equations, tables, and formatting. You can cite papers using the citation search feature, and they will be automatically added to your \texttt{cite.bib} file.

\section{Mathematical Expressions}
Here are some example equations:

\subsection{Quadratic Formula}
The quadratic formula is given by:
\[x = \frac{-b \pm \sqrt{b^2 - 4ac}}{2a}\]

\subsection{Matrix Example}
A sample matrix:
\[
\begin{pmatrix}
a_{11} & a_{12} & a_{13} \\
a_{21} & a_{22} & a_{23} \\
a_{31} & a_{32} & a_{33}
\end{pmatrix}
\]

\section{Tables}
Below is a sample table showing data:

\begin{table}[h]
\centering
\begin{tabular}{@{}lrrr@{}}
\toprule
\textbf{Category} & \textbf{Value 1} & \textbf{Value 2} & \textbf{Total} \\
\midrule
Item A & 10.5 & 20.3 & 30.8 \\
Item B & 15.2 & 25.1 & 40.3 \\
Item C & 22.1 & 12.7 & 34.8 \\
\bottomrule
\end{tabular}
\caption{Sample Data Table}
\label{tab:sample}
\end{table}

\section{Lists}
\subsection{Bullet Points}
\begin{itemize}
  \item First item
  \item Second item
  \item Third item with sub-items:
    \begin{itemize}
      \item Sub-item 1
      \item Sub-item 2
    \end{itemize}
\end{itemize}

\subsection{Numbered List}
\begin{enumerate}
  \item First step
  \item Second step
  \item Third step
\end{enumerate}

\section{Citations and References}
When you use the citation search feature, references will be automatically managed. For example, you might cite a paper like this: \cite{example2024}. All citations will be automatically added to your \texttt{cite.bib} file.

\section{Conclusion}
This template demonstrates basic \LaTeX{} features including mathematical formulas, tables, and lists. The bibliography system is now set up to work seamlessly with Bibby AI's citation management. Feel free to modify it for your needs!

\bibliographystyle{plain}
\bibliography{cite}

\end{document}\documentclass{article}
\usepackage[utf8]{inputenc}
\usepackage{booktabs}
\usepackage{amsmath}
\usepackage{natbib}

\title{My Document}
\author{Your Name}
\date{\today}

\begin{document}

\maketitle

\section{Introduction}
This is a sample document demonstrating various \LaTeX{} features including mathematical equations, tables, and formatting. You can cite papers using the citation search feature, and they will be automatically added to your \texttt{cite.bib} file.

\section{Mathematical Expressions}
Here are some example equations:

\subsection{Quadratic Formula}
The quadratic formula is given by:
\[x = \frac{-b \pm \sqrt{b^2 - 4ac}}{2a}\]

\subsection{Matrix Example}
A sample matrix:
\[
\begin{pmatrix}
a_{11} & a_{12} & a_{13} \\
a_{21} & a_{22} & a_{23} \\
a_{31} & a_{32} & a_{33}
\end{pmatrix}
\]

\section{Tables}
Below is a sample table showing data:

\begin{table}[h]
\centering
\begin{tabular}{@{}lrrr@{}}
\toprule
\textbf{Category} & \textbf{Value 1} & \textbf{Value 2} & \textbf{Total} \\
\midrule
Item A & 10.5 & 20.3 & 30.8 \\
Item B & 15.2 & 25.1 & 40.3 \\
Item C & 22.1 & 12.7 & 34.8 \\
\bottomrule
\end{tabular}
\caption{Sample Data Table}
\label{tab:sample}
\end{table}

\section{Lists}
\subsection{Bullet Points}
\begin{itemize}
  \item First item
  \item Second item
  \item Third item with sub-items:
    \begin{itemize}
      \item Sub-item 1
      \item Sub-item 2
    \end{itemize}
\end{itemize}

\subsection{Numbered List}
\begin{enumerate}
  \item First step
  \item Second step
  \item Third step
\end{enumerate}

\section{Citations and References}
When you use the citation search feature, references will be automatically managed. For example, you might cite a paper like this: \cite{example2024}. All citations will be automatically added to your \texttt{cite.bib} file.

\section{Conclusion}
This template demonstrates basic \LaTeX{} features including mathematical formulas, tables, and lists. The bibliography system is now set up to work seamlessly with Bibby AI's citation management. Feel free to modify it for your needs!

\bibliographystyle{plain}
\bibliography{cite}

\end{document}\documentclass{article}
\usepackage[utf8]{inputenc}
\usepackage{booktabs}
\usepackage{amsmath}
\usepackage{natbib}

\title{My Document}
\author{Your Name}
\date{\today}

\begin{document}

\maketitle

\section{Introduction}
This is a sample document demonstrating various \LaTeX{} features including mathematical equations, tables, and formatting. You can cite papers using the citation search feature, and they will be automatically added to your \texttt{cite.bib} file.

\section{Mathematical Expressions}
Here are some example equations:

\subsection{Quadratic Formula}
The quadratic formula is given by:
\[x = \frac{-b \pm \sqrt{b^2 - 4ac}}{2a}\]

\subsection{Matrix Example}
A sample matrix:
\[
\begin{pmatrix}
a_{11} & a_{12} & a_{13} \\
a_{21} & a_{22} & a_{23} \\
a_{31} & a_{32} & a_{33}
\end{pmatrix}
\]

\section{Tables}
Below is a sample table showing data:

\begin{table}[h]
\centering
\begin{tabular}{@{}lrrr@{}}
\toprule
\textbf{Category} & \textbf{Value 1} & \textbf{Value 2} & \textbf{Total} \\
\midrule
Item A & 10.5 & 20.3 & 30.8 \\
Item B & 15.2 & 25.1 & 40.3 \\
Item C & 22.1 & 12.7 & 34.8 \\
\bottomrule
\end{tabular}
\caption{Sample Data Table}
\label{tab:sample}
\end{table}

\section{Lists}
\subsection{Bullet Points}
\begin{itemize}
  \item First item
  \item Second item
  \item Third item with sub-items:
    \begin{itemize}
      \item Sub-item 1
      \item Sub-item 2
    \end{itemize}
\end{itemize}

\subsection{Numbered List}
\begin{enumerate}
  \item First step
  \item Second step
  \item Third step
\end{enumerate}

\section{Citations and References}
When you use the citation search feature, references will be automatically managed. For example, you might cite a paper like this: \cite{example2024}. All citations will be automatically added to your \texttt{cite.bib} file.

\section{Conclusion}
This template demonstrates basic \LaTeX{} features including mathematical formulas, tables, and lists. The bibliography system is now set up to work seamlessly with Bibby AI's citation management. Feel free to modify it for your needs!

\bibliographystyle{plain}
\bibliography{cite}

\end{document}\documentclass{article}
\usepackage[utf8]{inputenc}
\usepackage{booktabs}
\usepackage{amsmath}
\usepackage{natbib}

\title{My Document}
\author{Your Name}
\date{\today}

\begin{document}

\maketitle

\section{Introduction}
This is a sample document demonstrating various \LaTeX{} features including mathematical equations, tables, and formatting. You can cite papers using the citation search feature, and they will be automatically added to your \texttt{cite.bib} file.

\section{Mathematical Expressions}
Here are some example equations:

\subsection{Quadratic Formula}
The quadratic formula is given by:
\[x = \frac{-b \pm \sqrt{b^2 - 4ac}}{2a}\]

\subsection{Matrix Example}
A sample matrix:
\[
\begin{pmatrix}
a_{11} & a_{12} & a_{13} \\
a_{21} & a_{22} & a_{23} \\
a_{31} & a_{32} & a_{33}
\end{pmatrix}
\]

\section{Tables}
Below is a sample table showing data:

\begin{table}[h]
\centering
\begin{tabular}{@{}lrrr@{}}
\toprule
\textbf{Category} & \textbf{Value 1} & \textbf{Value 2} & \textbf{Total} \\
\midrule
Item A & 10.5 & 20.3 & 30.8 \\
Item B & 15.2 & 25.1 & 40.3 \\
Item C & 22.1 & 12.7 & 34.8 \\
\bottomrule
\end{tabular}
\caption{Sample Data Table}
\label{tab:sample}
\end{table}

\section{Lists}
\subsection{Bullet Points}
\begin{itemize}
  \item First item
  \item Second item
  \item Third item with sub-items:
    \begin{itemize}
      \item Sub-item 1
      \item Sub-item 2
    \end{itemize}
\end{itemize}

\subsection{Numbered List}
\begin{enumerate}
  \item First step
  \item Second step
  \item Third step
\end{enumerate}

\section{Citations and References}
When you use the citation search feature, references will be automatically managed. For example, you might cite a paper like this: \cite{example2024}. All citations will be automatically added to your \texttt{cite.bib} file.

\section{Conclusion}
This template demonstrates basic \LaTeX{} features including mathematical formulas, tables, and lists. The bibliography system is now set up to work seamlessly with Bibby AI's citation management. Feel free to modify it for your needs!

\bibliographystyle{plain}
\bibliography{cite}

\end{document}\documentclass{article}
\usepackage[utf8]{inputenc}
\usepackage{booktabs}
\usepackage{amsmath}
\usepackage{natbib}

\title{My Document}
\author{Your Name}
\date{\today}

\begin{document}

\maketitle

\section{Introduction}
This is a sample document demonstrating various \LaTeX{} features including mathematical equations, tables, and formatting. You can cite papers using the citation search feature, and they will be automatically added to your \texttt{cite.bib} file.

\section{Mathematical Expressions}
Here are some example equations:

\subsection{Quadratic Formula}
The quadratic formula is given by:
\[x = \frac{-b \pm \sqrt{b^2 - 4ac}}{2a}\]

\subsection{Matrix Example}
A sample matrix:
\[
\begin{pmatrix}
a_{11} & a_{12} & a_{13} \\
a_{21} & a_{22} & a_{23} \\
a_{31} & a_{32} & a_{33}
\end{pmatrix}
\]

\section{Tables}
Below is a sample table showing data:

\begin{table}[h]
\centering
\begin{tabular}{@{}lrrr@{}}
\toprule
\textbf{Category} & \textbf{Value 1} & \textbf{Value 2} & \textbf{Total} \\
\midrule
Item A & 10.5 & 20.3 & 30.8 \\
Item B & 15.2 & 25.1 & 40.3 \\
Item C & 22.1 & 12.7 & 34.8 \\
\bottomrule
\end{tabular}
\caption{Sample Data Table}
\label{tab:sample}
\end{table}

\section{Lists}
\subsection{Bullet Points}
\begin{itemize}
  \item First item
  \item Second item
  \item Third item with sub-items:
    \begin{itemize}
      \item Sub-item 1
      \item Sub-item 2
    \end{itemize}
\end{itemize}

\subsection{Numbered List}
\begin{enumerate}
  \item First step
  \item Second step
  \item Third step
\end{enumerate}

\section{Citations and References}
When you use the citation search feature, references will be automatically managed. For example, you might cite a paper like this: \cite{example2024}. All citations will be automatically added to your \texttt{cite.bib} file.

\section{Conclusion}
This template demonstrates basic \LaTeX{} features including mathematical formulas, tables, and lists. The bibliography system is now set up to work seamlessly with Bibby AI's citation management. Feel free to modify it for your needs!

\bibliographystyle{plain}
\bibliography{cite}

\end{document}\documentclass{article}
\usepackage[utf8]{inputenc}
\usepackage{booktabs}
\usepackage{amsmath}
\usepackage{natbib}

\title{My Document}
\author{Your Name}
\date{\today}

\begin{document}

\maketitle

\section{Introduction}
This is a sample document demonstrating various \LaTeX{} features including mathematical equations, tables, and formatting. You can cite papers using the citation search feature, and they will be automatically added to your \texttt{cite.bib} file.

\section{Mathematical Expressions}
Here are some example equations:

\subsection{Quadratic Formula}
The quadratic formula is given by:
\[x = \frac{-b \pm \sqrt{b^2 - 4ac}}{2a}\]

\subsection{Matrix Example}
A sample matrix:
\[
\begin{pmatrix}
a_{11} & a_{12} & a_{13} \\
a_{21} & a_{22} & a_{23} \\
a_{31} & a_{32} & a_{33}
\end{pmatrix}
\]

\section{Tables}
Below is a sample table showing data:

\begin{table}[h]
\centering
\begin{tabular}{@{}lrrr@{}}
\toprule
\textbf{Category} & \textbf{Value 1} & \textbf{Value 2} & \textbf{Total} \\
\midrule
Item A & 10.5 & 20.3 & 30.8 \\
Item B & 15.2 & 25.1 & 40.3 \\
Item C & 22.1 & 12.7 & 34.8 \\
\bottomrule
\end{tabular}
\caption{Sample Data Table}
\label{tab:sample}
\end{table}

\section{Lists}
\subsection{Bullet Points}
\begin{itemize}
  \item First item
  \item Second item
  \item Third item with sub-items:
    \begin{itemize}
      \item Sub-item 1
      \item Sub-item 2
    \end{itemize}
\end{itemize}

\subsection{Numbered List}
\begin{enumerate}
  \item First step
  \item Second step
  \item Third step
\end{enumerate}

\section{Citations and References}
When you use the citation search feature, references will be automatically managed. For example, you might cite a paper like this: \cite{example2024}. All citations will be automatically added to your \texttt{cite.bib} file.

\section{Conclusion}
This template demonstrates basic \LaTeX{} features including mathematical formulas, tables, and lists. The bibliography system is now set up to work seamlessly with Bibby AI's citation management. Feel free to modify it for your needs!

\bibliographystyle{plain}
\bibliography{cite}

\end{document}\documentclass{article}
\usepackage[utf8]{inputenc}
\usepackage{booktabs}
\usepackage{amsmath}
\usepackage{natbib}

\title{My Document}
\author{Your Name}
\date{\today}

\begin{document}

\maketitle

\section{Introduction}
This is a sample document demonstrating various \LaTeX{} features including mathematical equations, tables, and formatting. You can cite papers using the citation search feature, and they will be automatically added to your \texttt{cite.bib} file.

\section{Mathematical Expressions}
Here are some example equations:

\subsection{Quadratic Formula}
The quadratic formula is given by:
\[x = \frac{-b \pm \sqrt{b^2 - 4ac}}{2a}\]

\subsection{Matrix Example}
A sample matrix:
\[
\begin{pmatrix}
a_{11} & a_{12} & a_{13} \\
a_{21} & a_{22} & a_{23} \\
a_{31} & a_{32} & a_{33}
\end{pmatrix}
\]

\section{Tables}
Below is a sample table showing data:

\begin{table}[h]
\centering
\begin{tabular}{@{}lrrr@{}}
\toprule
\textbf{Category} & \textbf{Value 1} & \textbf{Value 2} & \textbf{Total} \\
\midrule
Item A & 10.5 & 20.3 & 30.8 \\
Item B & 15.2 & 25.1 & 40.3 \\
Item C & 22.1 & 12.7 & 34.8 \\
\bottomrule
\end{tabular}
\caption{Sample Data Table}
\label{tab:sample}
\end{table}

\section{Lists}
\subsection{Bullet Points}
\begin{itemize}
  \item First item
  \item Second item
  \item Third item with sub-items:
    \begin{itemize}
      \item Sub-item 1
      \item Sub-item 2
    \end{itemize}
\end{itemize}

\subsection{Numbered List}
\begin{enumerate}
  \item First step
  \item Second step
  \item Third step
\end{enumerate}

\section{Citations and References}
When you use the citation search feature, references will be automatically managed. For example, you might cite a paper like this: \cite{example2024}. All citations will be automatically added to your \texttt{cite.bib} file.

\section{Conclusion}
This template demonstrates basic \LaTeX{} features including mathematical formulas, tables, and lists. The bibliography system is now set up to work seamlessly with Bibby AI's citation management. Feel free to modify it for your needs!

\bibliographystyle{plain}
\bibliography{cite}

\end{document}\documentclass{article}
\usepackage[utf8]{inputenc}
\usepackage{booktabs}
\usepackage{amsmath}
\usepackage{natbib}

\title{My Document}
\author{Your Name}
\date{\today}

\begin{document}

\maketitle

\section{Introduction}
This is a sample document demonstrating various \LaTeX{} features including mathematical equations, tables, and formatting. You can cite papers using the citation search feature, and they will be automatically added to your \texttt{cite.bib} file.

\section{Mathematical Expressions}
Here are some example equations:

\subsection{Quadratic Formula}
The quadratic formula is given by:
\[x = \frac{-b \pm \sqrt{b^2 - 4ac}}{2a}\]

\subsection{Matrix Example}
A sample matrix:
\[
\begin{pmatrix}
a_{11} & a_{12} & a_{13} \\
a_{21} & a_{22} & a_{23} \\
a_{31} & a_{32} & a_{33}
\end{pmatrix}
\]

\section{Tables}
Below is a sample table showing data:

\begin{table}[h]
\centering
\begin{tabular}{@{}lrrr@{}}
\toprule
\textbf{Category} & \textbf{Value 1} & \textbf{Value 2} & \textbf{Total} \\
\midrule
Item A & 10.5 & 20.3 & 30.8 \\
Item B & 15.2 & 25.1 & 40.3 \\
Item C & 22.1 & 12.7 & 34.8 \\
\bottomrule
\end{tabular}
\caption{Sample Data Table}
\label{tab:sample}
\end{table}

\section{Lists}
\subsection{Bullet Points}
\begin{itemize}
  \item First item
  \item Second item
  \item Third item with sub-items:
    \begin{itemize}
      \item Sub-item 1
      \item Sub-item 2
    \end{itemize}
\end{itemize}

\subsection{Numbered List}
\begin{enumerate}
  \item First step
  \item Second step
  \item Third step
\end{enumerate}

\section{Citations and References}
When you use the citation search feature, references will be automatically managed. For example, you might cite a paper like this: \cite{example2024}. All citations will be automatically added to your \texttt{cite.bib} file.

\section{Conclusion}
This template demonstrates basic \LaTeX{} features including mathematical formulas, tables, and lists. The bibliography system is now set up to work seamlessly with Bibby AI's citation management. Feel free to modify it for your needs!

\bibliographystyle{plain}
\bibliography{cite}

\end{document}\documentclass{article}
\usepackage[utf8]{inputenc}
\usepackage{booktabs}
\usepackage{amsmath}
\usepackage{natbib}

\title{My Document}
\author{Your Name}
\date{\today}

\begin{document}

\maketitle

\section{Introduction}
This is a sample document demonstrating various \LaTeX{} features including mathematical equations, tables, and formatting. You can cite papers using the citation search feature, and they will be automatically added to your \texttt{cite.bib} file.

\section{Mathematical Expressions}
Here are some example equations:

\subsection{Quadratic Formula}
The quadratic formula is given by:
\[x = \frac{-b \pm \sqrt{b^2 - 4ac}}{2a}\]

\subsection{Matrix Example}
A sample matrix:
\[
\begin{pmatrix}
a_{11} & a_{12} & a_{13} \\
a_{21} & a_{22} & a_{23} \\
a_{31} & a_{32} & a_{33}
\end{pmatrix}
\]

\section{Tables}
Below is a sample table showing data:

\begin{table}[h]
\centering
\begin{tabular}{@{}lrrr@{}}
\toprule
\textbf{Category} & \textbf{Value 1} & \textbf{Value 2} & \textbf{Total} \\
\midrule
Item A & 10.5 & 20.3 & 30.8 \\
Item B & 15.2 & 25.1 & 40.3 \\
Item C & 22.1 & 12.7 & 34.8 \\
\bottomrule
\end{tabular}
\caption{Sample Data Table}
\label{tab:sample}
\end{table}

\section{Lists}
\subsection{Bullet Points}
\begin{itemize}
  \item First item
  \item Second item
  \item Third item with sub-items:
    \begin{itemize}
      \item Sub-item 1
      \item Sub-item 2
    \end{itemize}
\end{itemize}

\subsection{Numbered List}
\begin{enumerate}
  \item First step
  \item Second step
  \item Third step
\end{enumerate}

\section{Citations and References}
When you use the citation search feature, references will be automatically managed. For example, you might cite a paper like this: \cite{example2024}. All citations will be automatically added to your \texttt{cite.bib} file.

\section{Conclusion}
This template demonstrates basic \LaTeX{} features including mathematical formulas, tables, and lists. The bibliography system is now set up to work seamlessly with Bibby AI's citation management. Feel free to modify it for your needs!

\bibliographystyle{plain}
\bibliography{cite}

\end{document}\documentclass{article}
\usepackage[utf8]{inputenc}
\usepackage{booktabs}
\usepackage{amsmath}
\usepackage{natbib}

\title{My Document}
\author{Your Name}
\date{\today}

\begin{document}

\maketitle

\section{Introduction}
This is a sample document demonstrating various \LaTeX{} features including mathematical equations, tables, and formatting. You can cite papers using the citation search feature, and they will be automatically added to your \texttt{cite.bib} file.

\section{Mathematical Expressions}
Here are some example equations:

\subsection{Quadratic Formula}
The quadratic formula is given by:
\[x = \frac{-b \pm \sqrt{b^2 - 4ac}}{2a}\]

\subsection{Matrix Example}
A sample matrix:
\[
\begin{pmatrix}
a_{11} & a_{12} & a_{13} \\
a_{21} & a_{22} & a_{23} \\
a_{31} & a_{32} & a_{33}
\end{pmatrix}
\]

\section{Tables}
Below is a sample table showing data:

\begin{table}[h]
\centering
\begin{tabular}{@{}lrrr@{}}
\toprule
\textbf{Category} & \textbf{Value 1} & \textbf{Value 2} & \textbf{Total} \\
\midrule
Item A & 10.5 & 20.3 & 30.8 \\
Item B & 15.2 & 25.1 & 40.3 \\
Item C & 22.1 & 12.7 & 34.8 \\
\bottomrule
\end{tabular}
\caption{Sample Data Table}
\label{tab:sample}
\end{table}

\section{Lists}
\subsection{Bullet Points}
\begin{itemize}
  \item First item
  \item Second item
  \item Third item with sub-items:
    \begin{itemize}
      \item Sub-item 1
      \item Sub-item 2
    \end{itemize}
\end{itemize}

\subsection{Numbered List}
\begin{enumerate}
  \item First step
  \item Second step
  \item Third step
\end{enumerate}

\section{Citations and References}
When you use the citation search feature, references will be automatically managed. For example, you might cite a paper like this: \cite{example2024}. All citations will be automatically added to your \texttt{cite.bib} file.

\section{Conclusion}
This template demonstrates basic \LaTeX{} features including mathematical formulas, tables, and lists. The bibliography system is now set up to work seamlessly with Bibby AI's citation management. Feel free to modify it for your needs!

\bibliographystyle{plain}
\bibliography{cite}

\end{document}\documentclass{article}
\usepackage[utf8]{inputenc}
\usepackage{booktabs}
\usepackage{amsmath}
\usepackage{natbib}

\title{My Document}
\author{Your Name}
\date{\today}

\begin{document}

\maketitle

\section{Introduction}
This is a sample document demonstrating various \LaTeX{} features including mathematical equations, tables, and formatting. You can cite papers using the citation search feature, and they will be automatically added to your \texttt{cite.bib} file.

\section{Mathematical Expressions}
Here are some example equations:

\subsection{Quadratic Formula}
The quadratic formula is given by:
\[x = \frac{-b \pm \sqrt{b^2 - 4ac}}{2a}\]

\subsection{Matrix Example}
A sample matrix:
\[
\begin{pmatrix}
a_{11} & a_{12} & a_{13} \\
a_{21} & a_{22} & a_{23} \\
a_{31} & a_{32} & a_{33}
\end{pmatrix}
\]

\section{Tables}
Below is a sample table showing data:

\begin{table}[h]
\centering
\begin{tabular}{@{}lrrr@{}}
\toprule
\textbf{Category} & \textbf{Value 1} & \textbf{Value 2} & \textbf{Total} \\
\midrule
Item A & 10.5 & 20.3 & 30.8 \\
Item B & 15.2 & 25.1 & 40.3 \\
Item C & 22.1 & 12.7 & 34.8 \\
\bottomrule
\end{tabular}
\caption{Sample Data Table}
\label{tab:sample}
\end{table}

\section{Lists}
\subsection{Bullet Points}
\begin{itemize}
  \item First item
  \item Second item
  \item Third item with sub-items:
    \begin{itemize}
      \item Sub-item 1
      \item Sub-item 2
    \end{itemize}
\end{itemize}

\subsection{Numbered List}
\begin{enumerate}
  \item First step
  \item Second step
  \item Third step
\end{enumerate}

\section{Citations and References}
When you use the citation search feature, references will be automatically managed. For example, you might cite a paper like this: \cite{example2024}. All citations will be automatically added to your \texttt{cite.bib} file.

\section{Conclusion}
This template demonstrates basic \LaTeX{} features including mathematical formulas, tables, and lists. The bibliography system is now set up to work seamlessly with Bibby AI's citation management. Feel free to modify it for your needs!

\bibliographystyle{plain}
\bibliography{cite}

\end{document}\documentclass{article}
\usepackage[utf8]{inputenc}
\usepackage{booktabs}
\usepackage{amsmath}
\usepackage{natbib}

\title{My Document}
\author{Your Name}
\date{\today}

\begin{document}

\maketitle

\section{Introduction}
This is a sample document demonstrating various \LaTeX{} features including mathematical equations, tables, and formatting. You can cite papers using the citation search feature, and they will be automatically added to your \texttt{cite.bib} file.

\section{Mathematical Expressions}
Here are some example equations:

\subsection{Quadratic Formula}
The quadratic formula is given by:
\[x = \frac{-b \pm \sqrt{b^2 - 4ac}}{2a}\]

\subsection{Matrix Example}
A sample matrix:
\[
\begin{pmatrix}
a_{11} & a_{12} & a_{13} \\
a_{21} & a_{22} & a_{23} \\
a_{31} & a_{32} & a_{33}
\end{pmatrix}
\]

\section{Tables}
Below is a sample table showing data:

\begin{table}[h]
\centering
\begin{tabular}{@{}lrrr@{}}
\toprule
\textbf{Category} & \textbf{Value 1} & \textbf{Value 2} & \textbf{Total} \\
\midrule
Item A & 10.5 & 20.3 & 30.8 \\
Item B & 15.2 & 25.1 & 40.3 \\
Item C & 22.1 & 12.7 & 34.8 \\
\bottomrule
\end{tabular}
\caption{Sample Data Table}
\label{tab:sample}
\end{table}

\section{Lists}
\subsection{Bullet Points}
\begin{itemize}
  \item First item
  \item Second item
  \item Third item with sub-items:
    \begin{itemize}
      \item Sub-item 1
      \item Sub-item 2
    \end{itemize}
\end{itemize}

\subsection{Numbered List}
\begin{enumerate}
  \item First step
  \item Second step
  \item Third step
\end{enumerate}

\section{Citations and References}
When you use the citation search feature, references will be automatically managed. For example, you might cite a paper like this: \cite{example2024}. All citations will be automatically added to your \texttt{cite.bib} file.

\section{Conclusion}
This template demonstrates basic \LaTeX{} features including mathematical formulas, tables, and lists. The bibliography system is now set up to work seamlessly with Bibby AI's citation management. Feel free to modify it for your needs!

\bibliographystyle{plain}
\bibliography{cite}

\end{document}\documentclass{article}
\usepackage[utf8]{inputenc}
\usepackage{booktabs}
\usepackage{amsmath}
\usepackage{natbib}

\title{My Document}
\author{Your Name}
\date{\today}

\begin{document}

\maketitle

\section{Introduction}
This is a sample document demonstrating various \LaTeX{} features including mathematical equations, tables, and formatting. You can cite papers using the citation search feature, and they will be automatically added to your \texttt{cite.bib} file.

\section{Mathematical Expressions}
Here are some example equations:

\subsection{Quadratic Formula}
The quadratic formula is given by:
\[x = \frac{-b \pm \sqrt{b^2 - 4ac}}{2a}\]

\subsection{Matrix Example}
A sample matrix:
\[
\begin{pmatrix}
a_{11} & a_{12} & a_{13} \\
a_{21} & a_{22} & a_{23} \\
a_{31} & a_{32} & a_{33}
\end{pmatrix}
\]

\section{Tables}
Below is a sample table showing data:

\begin{table}[h]
\centering
\begin{tabular}{@{}lrrr@{}}
\toprule
\textbf{Category} & \textbf{Value 1} & \textbf{Value 2} & \textbf{Total} \\
\midrule
Item A & 10.5 & 20.3 & 30.8 \\
Item B & 15.2 & 25.1 & 40.3 \\
Item C & 22.1 & 12.7 & 34.8 \\
\bottomrule
\end{tabular}
\caption{Sample Data Table}
\label{tab:sample}
\end{table}

\section{Lists}
\subsection{Bullet Points}
\begin{itemize}
  \item First item
  \item Second item
  \item Third item with sub-items:
    \begin{itemize}
      \item Sub-item 1
      \item Sub-item 2
    \end{itemize}
\end{itemize}

\subsection{Numbered List}
\begin{enumerate}
  \item First step
  \item Second step
  \item Third step
\end{enumerate}

\section{Citations and References}
When you use the citation search feature, references will be automatically managed. For example, you might cite a paper like this: \cite{example2024}. All citations will be automatically added to your \texttt{cite.bib} file.

\section{Conclusion}
This template demonstrates basic \LaTeX{} features including mathematical formulas, tables, and lists. The bibliography system is now set up to work seamlessly with Bibby AI's citation management. Feel free to modify it for your needs!

\bibliographystyle{plain}
\bibliography{cite}

\end{document}\documentclass{article}
\usepackage[utf8]{inputenc}
\usepackage{booktabs}
\usepackage{amsmath}
\usepackage{natbib}

\title{My Document}
\author{Your Name}
\date{\today}

\begin{document}

\maketitle

\section{Introduction}
This is a sample document demonstrating various \LaTeX{} features including mathematical equations, tables, and formatting. You can cite papers using the citation search feature, and they will be automatically added to your \texttt{cite.bib} file.

\section{Mathematical Expressions}
Here are some example equations:

\subsection{Quadratic Formula}
The quadratic formula is given by:
\[x = \frac{-b \pm \sqrt{b^2 - 4ac}}{2a}\]

\subsection{Matrix Example}
A sample matrix:
\[
\begin{pmatrix}
a_{11} & a_{12} & a_{13} \\
a_{21} & a_{22} & a_{23} \\
a_{31} & a_{32} & a_{33}
\end{pmatrix}
\]

\section{Tables}
Below is a sample table showing data:

\begin{table}[h]
\centering
\begin{tabular}{@{}lrrr@{}}
\toprule
\textbf{Category} & \textbf{Value 1} & \textbf{Value 2} & \textbf{Total} \\
\midrule
Item A & 10.5 & 20.3 & 30.8 \\
Item B & 15.2 & 25.1 & 40.3 \\
Item C & 22.1 & 12.7 & 34.8 \\
\bottomrule
\end{tabular}
\caption{Sample Data Table}
\label{tab:sample}
\end{table}

\section{Lists}
\subsection{Bullet Points}
\begin{itemize}
  \item First item
  \item Second item
  \item Third item with sub-items:
    \begin{itemize}
      \item Sub-item 1
      \item Sub-item 2
    \end{itemize}
\end{itemize}

\subsection{Numbered List}
\begin{enumerate}
  \item First step
  \item Second step
  \item Third step
\end{enumerate}

\section{Citations and References}
When you use the citation search feature, references will be automatically managed. For example, you might cite a paper like this: \cite{example2024}. All citations will be automatically added to your \texttt{cite.bib} file.

\section{Conclusion}
This template demonstrates basic \LaTeX{} features including mathematical formulas, tables, and lists. The bibliography system is now set up to work seamlessly with Bibby AI's citation management. Feel free to modify it for your needs!

\bibliographystyle{plain}
\bibliography{cite}

\end{document}\documentclass{article}
\usepackage[utf8]{inputenc}
\usepackage{booktabs}
\usepackage{amsmath}
\usepackage{natbib}

\title{My Document}
\author{Your Name}
\date{\today}

\begin{document}

\maketitle

\section{Introduction}
This is a sample document demonstrating various \LaTeX{} features including mathematical equations, tables, and formatting. You can cite papers using the citation search feature, and they will be automatically added to your \texttt{cite.bib} file.

\section{Mathematical Expressions}
Here are some example equations:

\subsection{Quadratic Formula}
The quadratic formula is given by:
\[x = \frac{-b \pm \sqrt{b^2 - 4ac}}{2a}\]

\subsection{Matrix Example}
A sample matrix:
\[
\begin{pmatrix}
a_{11} & a_{12} & a_{13} \\
a_{21} & a_{22} & a_{23} \\
a_{31} & a_{32} & a_{33}
\end{pmatrix}
\]

\section{Tables}
Below is a sample table showing data:

\begin{table}[h]
\centering
\begin{tabular}{@{}lrrr@{}}
\toprule
\textbf{Category} & \textbf{Value 1} & \textbf{Value 2} & \textbf{Total} \\
\midrule
Item A & 10.5 & 20.3 & 30.8 \\
Item B & 15.2 & 25.1 & 40.3 \\
Item C & 22.1 & 12.7 & 34.8 \\
\bottomrule
\end{tabular}
\caption{Sample Data Table}
\label{tab:sample}
\end{table}

\section{Lists}
\subsection{Bullet Points}
\begin{itemize}
  \item First item
  \item Second item
  \item Third item with sub-items:
    \begin{itemize}
      \item Sub-item 1
      \item Sub-item 2
    \end{itemize}
\end{itemize}

\subsection{Numbered List}
\begin{enumerate}
  \item First step
  \item Second step
  \item Third step
\end{enumerate}

\section{Citations and References}
When you use the citation search feature, references will be automatically managed. For example, you might cite a paper like this: \cite{example2024}. All citations will be automatically added to your \texttt{cite.bib} file.

\section{Conclusion}
This template demonstrates basic \LaTeX{} features including mathematical formulas, tables, and lists. The bibliography system is now set up to work seamlessly with Bibby AI's citation management. Feel free to modify it for your needs!

\bibliographystyle{plain}
\bibliography{cite}

\end{document}\documentclass{article}
\usepackage[utf8]{inputenc}
\usepackage{booktabs}
\usepackage{amsmath}
\usepackage{natbib}

\title{My Document}
\author{Your Name}
\date{\today}

\begin{document}

\maketitle

\section{Introduction}
This is a sample document demonstrating various \LaTeX{} features including mathematical equations, tables, and formatting. You can cite papers using the citation search feature, and they will be automatically added to your \texttt{cite.bib} file.

\section{Mathematical Expressions}
Here are some example equations:

\subsection{Quadratic Formula}
The quadratic formula is given by:
\[x = \frac{-b \pm \sqrt{b^2 - 4ac}}{2a}\]

\subsection{Matrix Example}
A sample matrix:
\[
\begin{pmatrix}
a_{11} & a_{12} & a_{13} \\
a_{21} & a_{22} & a_{23} \\
a_{31} & a_{32} & a_{33}
\end{pmatrix}
\]

\section{Tables}
Below is a sample table showing data:

\begin{table}[h]
\centering
\begin{tabular}{@{}lrrr@{}}
\toprule
\textbf{Category} & \textbf{Value 1} & \textbf{Value 2} & \textbf{Total} \\
\midrule
Item A & 10.5 & 20.3 & 30.8 \\
Item B & 15.2 & 25.1 & 40.3 \\
Item C & 22.1 & 12.7 & 34.8 \\
\bottomrule
\end{tabular}
\caption{Sample Data Table}
\label{tab:sample}
\end{table}

\section{Lists}
\subsection{Bullet Points}
\begin{itemize}
  \item First item
  \item Second item
  \item Third item with sub-items:
    \begin{itemize}
      \item Sub-item 1
      \item Sub-item 2
    \end{itemize}
\end{itemize}

\subsection{Numbered List}
\begin{enumerate}
  \item First step
  \item Second step
  \item Third step
\end{enumerate}

\section{Citations and References}
When you use the citation search feature, references will be automatically managed. For example, you might cite a paper like this: \cite{example2024}. All citations will be automatically added to your \texttt{cite.bib} file.

\section{Conclusion}
This template demonstrates basic \LaTeX{} features including mathematical formulas, tables, and lists. The bibliography system is now set up to work seamlessly with Bibby AI's citation management. Feel free to modify it for your needs!

\bibliographystyle{plain}
\bibliography{cite}

\end{document}\documentclass{article}
\usepackage[utf8]{inputenc}
\usepackage{booktabs}
\usepackage{amsmath}
\usepackage{natbib}

\title{My Document}
\author{Your Name}
\date{\today}

\begin{document}

\maketitle

\section{Introduction}
This is a sample document demonstrating various \LaTeX{} features including mathematical equations, tables, and formatting. You can cite papers using the citation search feature, and they will be automatically added to your \texttt{cite.bib} file.

\section{Mathematical Expressions}
Here are some example equations:

\subsection{Quadratic Formula}
The quadratic formula is given by:
\[x = \frac{-b \pm \sqrt{b^2 - 4ac}}{2a}\]

\subsection{Matrix Example}
A sample matrix:
\[
\begin{pmatrix}
a_{11} & a_{12} & a_{13} \\
a_{21} & a_{22} & a_{23} \\
a_{31} & a_{32} & a_{33}
\end{pmatrix}
\]

\section{Tables}
Below is a sample table showing data:

\begin{table}[h]
\centering
\begin{tabular}{@{}lrrr@{}}
\toprule
\textbf{Category} & \textbf{Value 1} & \textbf{Value 2} & \textbf{Total} \\
\midrule
Item A & 10.5 & 20.3 & 30.8 \\
Item B & 15.2 & 25.1 & 40.3 \\
Item C & 22.1 & 12.7 & 34.8 \\
\bottomrule
\end{tabular}
\caption{Sample Data Table}
\label{tab:sample}
\end{table}

\section{Lists}
\subsection{Bullet Points}
\begin{itemize}
  \item First item
  \item Second item
  \item Third item with sub-items:
    \begin{itemize}
      \item Sub-item 1
      \item Sub-item 2
    \end{itemize}
\end{itemize}

\subsection{Numbered List}
\begin{enumerate}
  \item First step
  \item Second step
  \item Third step
\end{enumerate}

\section{Citations and References}
When you use the citation search feature, references will be automatically managed. For example, you might cite a paper like this: \cite{example2024}. All citations will be automatically added to your \texttt{cite.bib} file.

\section{Conclusion}
This template demonstrates basic \LaTeX{} features including mathematical formulas, tables, and lists. The bibliography system is now set up to work seamlessly with Bibby AI's citation management. Feel free to modify it for your needs!

\bibliographystyle{plain}
\bibliography{cite}

\end{document}\documentclass{article}
\usepackage[utf8]{inputenc}
\usepackage{booktabs}
\usepackage{amsmath}
\usepackage{natbib}

\title{My Document}
\author{Your Name}
\date{\today}

\begin{document}

\maketitle

\section{Introduction}
This is a sample document demonstrating various \LaTeX{} features including mathematical equations, tables, and formatting. You can cite papers using the citation search feature, and they will be automatically added to your \texttt{cite.bib} file.

\section{Mathematical Expressions}
Here are some example equations:

\subsection{Quadratic Formula}
The quadratic formula is given by:
\[x = \frac{-b \pm \sqrt{b^2 - 4ac}}{2a}\]

\subsection{Matrix Example}
A sample matrix:
\[
\begin{pmatrix}
a_{11} & a_{12} & a_{13} \\
a_{21} & a_{22} & a_{23} \\
a_{31} & a_{32} & a_{33}
\end{pmatrix}
\]

\section{Tables}
Below is a sample table showing data:

\begin{table}[h]
\centering
\begin{tabular}{@{}lrrr@{}}
\toprule
\textbf{Category} & \textbf{Value 1} & \textbf{Value 2} & \textbf{Total} \\
\midrule
Item A & 10.5 & 20.3 & 30.8 \\
Item B & 15.2 & 25.1 & 40.3 \\
Item C & 22.1 & 12.7 & 34.8 \\
\bottomrule
\end{tabular}
\caption{Sample Data Table}
\label{tab:sample}
\end{table}

\section{Lists}
\subsection{Bullet Points}
\begin{itemize}
  \item First item
  \item Second item
  \item Third item with sub-items:
    \begin{itemize}
      \item Sub-item 1
      \item Sub-item 2
    \end{itemize}
\end{itemize}

\subsection{Numbered List}
\begin{enumerate}
  \item First step
  \item Second step
  \item Third step
\end{enumerate}

\section{Citations and References}
When you use the citation search feature, references will be automatically managed. For example, you might cite a paper like this: \cite{example2024}. All citations will be automatically added to your \texttt{cite.bib} file.

\section{Conclusion}
This template demonstrates basic \LaTeX{} features including mathematical formulas, tables, and lists. The bibliography system is now set up to work seamlessly with Bibby AI's citation management. Feel free to modify it for your needs!

\bibliographystyle{plain}
\bibliography{cite}

\end{document}\documentclass{article}
\usepackage[utf8]{inputenc}
\usepackage{booktabs}
\usepackage{amsmath}
\usepackage{natbib}

\title{My Document}
\author{Your Name}
\date{\today}

\begin{document}

\maketitle

\section{Introduction}
This is a sample document demonstrating various \LaTeX{} features including mathematical equations, tables, and formatting. You can cite papers using the citation search feature, and they will be automatically added to your \texttt{cite.bib} file.

\section{Mathematical Expressions}
Here are some example equations:

\subsection{Quadratic Formula}
The quadratic formula is given by:
\[x = \frac{-b \pm \sqrt{b^2 - 4ac}}{2a}\]

\subsection{Matrix Example}
A sample matrix:
\[
\begin{pmatrix}
a_{11} & a_{12} & a_{13} \\
a_{21} & a_{22} & a_{23} \\
a_{31} & a_{32} & a_{33}
\end{pmatrix}
\]

\section{Tables}
Below is a sample table showing data:

\begin{table}[h]
\centering
\begin{tabular}{@{}lrrr@{}}
\toprule
\textbf{Category} & \textbf{Value 1} & \textbf{Value 2} & \textbf{Total} \\
\midrule
Item A & 10.5 & 20.3 & 30.8 \\
Item B & 15.2 & 25.1 & 40.3 \\
Item C & 22.1 & 12.7 & 34.8 \\
\bottomrule
\end{tabular}
\caption{Sample Data Table}
\label{tab:sample}
\end{table}

\section{Lists}
\subsection{Bullet Points}
\begin{itemize}
  \item First item
  \item Second item
  \item Third item with sub-items:
    \begin{itemize}
      \item Sub-item 1
      \item Sub-item 2
    \end{itemize}
\end{itemize}

\subsection{Numbered List}
\begin{enumerate}
  \item First step
  \item Second step
  \item Third step
\end{enumerate}

\section{Citations and References}
When you use the citation search feature, references will be automatically managed. For example, you might cite a paper like this: \cite{example2024}. All citations will be automatically added to your \texttt{cite.bib} file.

\section{Conclusion}
This template demonstrates basic \LaTeX{} features including mathematical formulas, tables, and lists. The bibliography system is now set up to work seamlessly with Bibby AI's citation management. Feel free to modify it for your needs!

\bibliographystyle{plain}
\bibliography{cite}

\end{document}\documentclass{article}
\usepackage[utf8]{inputenc}
\usepackage{booktabs}
\usepackage{amsmath}
\usepackage{natbib}

\title{My Document}
\author{Your Name}
\date{\today}

\begin{document}

\maketitle

\section{Introduction}
This is a sample document demonstrating various \LaTeX{} features including mathematical equations, tables, and formatting. You can cite papers using the citation search feature, and they will be automatically added to your \texttt{cite.bib} file.

\section{Mathematical Expressions}
Here are some example equations:

\subsection{Quadratic Formula}
The quadratic formula is given by:
\[x = \frac{-b \pm \sqrt{b^2 - 4ac}}{2a}\]

\subsection{Matrix Example}
A sample matrix:
\[
\begin{pmatrix}
a_{11} & a_{12} & a_{13} \\
a_{21} & a_{22} & a_{23} \\
a_{31} & a_{32} & a_{33}
\end{pmatrix}
\]

\section{Tables}
Below is a sample table showing data:

\begin{table}[h]
\centering
\begin{tabular}{@{}lrrr@{}}
\toprule
\textbf{Category} & \textbf{Value 1} & \textbf{Value 2} & \textbf{Total} \\
\midrule
Item A & 10.5 & 20.3 & 30.8 \\
Item B & 15.2 & 25.1 & 40.3 \\
Item C & 22.1 & 12.7 & 34.8 \\
\bottomrule
\end{tabular}
\caption{Sample Data Table}
\label{tab:sample}
\end{table}

\section{Lists}
\subsection{Bullet Points}
\begin{itemize}
  \item First item
  \item Second item
  \item Third item with sub-items:
    \begin{itemize}
      \item Sub-item 1
      \item Sub-item 2
    \end{itemize}
\end{itemize}

\subsection{Numbered List}
\begin{enumerate}
  \item First step
  \item Second step
  \item Third step
\end{enumerate}

\section{Citations and References}
When you use the citation search feature, references will be automatically managed. For example, you might cite a paper like this: \cite{example2024}. All citations will be automatically added to your \texttt{cite.bib} file.

\section{Conclusion}
This template demonstrates basic \LaTeX{} features including mathematical formulas, tables, and lists. The bibliography system is now set up to work seamlessly with Bibby AI's citation management. Feel free to modify it for your needs!

\bibliographystyle{plain}
\bibliography{cite}

\end{document}\documentclass{article}
\usepackage[utf8]{inputenc}
\usepackage{booktabs}
\usepackage{amsmath}
\usepackage{natbib}

\title{My Document}
\author{Your Name}
\date{\today}

\begin{document}

\maketitle

\section{Introduction}
This is a sample document demonstrating various \LaTeX{} features including mathematical equations, tables, and formatting. You can cite papers using the citation search feature, and they will be automatically added to your \texttt{cite.bib} file.

\section{Mathematical Expressions}
Here are some example equations:

\subsection{Quadratic Formula}
The quadratic formula is given by:
\[x = \frac{-b \pm \sqrt{b^2 - 4ac}}{2a}\]

\subsection{Matrix Example}
A sample matrix:
\[
\begin{pmatrix}
a_{11} & a_{12} & a_{13} \\
a_{21} & a_{22} & a_{23} \\
a_{31} & a_{32} & a_{33}
\end{pmatrix}
\]

\section{Tables}
Below is a sample table showing data:

\begin{table}[h]
\centering
\begin{tabular}{@{}lrrr@{}}
\toprule
\textbf{Category} & \textbf{Value 1} & \textbf{Value 2} & \textbf{Total} \\
\midrule
Item A & 10.5 & 20.3 & 30.8 \\
Item B & 15.2 & 25.1 & 40.3 \\
Item C & 22.1 & 12.7 & 34.8 \\
\bottomrule
\end{tabular}
\caption{Sample Data Table}
\label{tab:sample}
\end{table}

\section{Lists}
\subsection{Bullet Points}
\begin{itemize}
  \item First item
  \item Second item
  \item Third item with sub-items:
    \begin{itemize}
      \item Sub-item 1
      \item Sub-item 2
    \end{itemize}
\end{itemize}

\subsection{Numbered List}
\begin{enumerate}
  \item First step
  \item Second step
  \item Third step
\end{enumerate}

\section{Citations and References}
When you use the citation search feature, references will be automatically managed. For example, you might cite a paper like this: \cite{example2024}. All citations will be automatically added to your \texttt{cite.bib} file.

\section{Conclusion}
This template demonstrates basic \LaTeX{} features including mathematical formulas, tables, and lists. The bibliography system is now set up to work seamlessly with Bibby AI's citation management. Feel free to modify it for your needs!

\bibliographystyle{plain}
\bibliography{cite}

\end{document}\documentclass{article}
\usepackage[utf8]{inputenc}
\usepackage{booktabs}
\usepackage{amsmath}
\usepackage{natbib}

\title{My Document}
\author{Your Name}
\date{\today}

\begin{document}

\maketitle

\section{Introduction}
This is a sample document demonstrating various \LaTeX{} features including mathematical equations, tables, and formatting. You can cite papers using the citation search feature, and they will be automatically added to your \texttt{cite.bib} file.

\section{Mathematical Expressions}
Here are some example equations:

\subsection{Quadratic Formula}
The quadratic formula is given by:
\[x = \frac{-b \pm \sqrt{b^2 - 4ac}}{2a}\]

\subsection{Matrix Example}
A sample matrix:
\[
\begin{pmatrix}
a_{11} & a_{12} & a_{13} \\
a_{21} & a_{22} & a_{23} \\
a_{31} & a_{32} & a_{33}
\end{pmatrix}
\]

\section{Tables}
Below is a sample table showing data:

\begin{table}[h]
\centering
\begin{tabular}{@{}lrrr@{}}
\toprule
\textbf{Category} & \textbf{Value 1} & \textbf{Value 2} & \textbf{Total} \\
\midrule
Item A & 10.5 & 20.3 & 30.8 \\
Item B & 15.2 & 25.1 & 40.3 \\
Item C & 22.1 & 12.7 & 34.8 \\
\bottomrule
\end{tabular}
\caption{Sample Data Table}
\label{tab:sample}
\end{table}

\section{Lists}
\subsection{Bullet Points}
\begin{itemize}
  \item First item
  \item Second item
  \item Third item with sub-items:
    \begin{itemize}
      \item Sub-item 1
      \item Sub-item 2
    \end{itemize}
\end{itemize}

\subsection{Numbered List}
\begin{enumerate}
  \item First step
  \item Second step
  \item Third step
\end{enumerate}

\section{Citations and References}
When you use the citation search feature, references will be automatically managed. For example, you might cite a paper like this: \cite{example2024}. All citations will be automatically added to your \texttt{cite.bib} file.

\section{Conclusion}
This template demonstrates basic \LaTeX{} features including mathematical formulas, tables, and lists. The bibliography system is now set up to work seamlessly with Bibby AI's citation management. Feel free to modify it for your needs!

\bibliographystyle{plain}
\bibliography{cite}

\end{document}\documentclass{article}
\usepackage[utf8]{inputenc}
\usepackage{booktabs}
\usepackage{amsmath}
\usepackage{natbib}

\title{My Document}
\author{Your Name}
\date{\today}

\begin{document}

\maketitle

\section{Introduction}
This is a sample document demonstrating various \LaTeX{} features including mathematical equations, tables, and formatting. You can cite papers using the citation search feature, and they will be automatically added to your \texttt{cite.bib} file.

\section{Mathematical Expressions}
Here are some example equations:

\subsection{Quadratic Formula}
The quadratic formula is given by:
\[x = \frac{-b \pm \sqrt{b^2 - 4ac}}{2a}\]

\subsection{Matrix Example}
A sample matrix:
\[
\begin{pmatrix}
a_{11} & a_{12} & a_{13} \\
a_{21} & a_{22} & a_{23} \\
a_{31} & a_{32} & a_{33}
\end{pmatrix}
\]

\section{Tables}
Below is a sample table showing data:

\begin{table}[h]
\centering
\begin{tabular}{@{}lrrr@{}}
\toprule
\textbf{Category} & \textbf{Value 1} & \textbf{Value 2} & \textbf{Total} \\
\midrule
Item A & 10.5 & 20.3 & 30.8 \\
Item B & 15.2 & 25.1 & 40.3 \\
Item C & 22.1 & 12.7 & 34.8 \\
\bottomrule
\end{tabular}
\caption{Sample Data Table}
\label{tab:sample}
\end{table}

\section{Lists}
\subsection{Bullet Points}
\begin{itemize}
  \item First item
  \item Second item
  \item Third item with sub-items:
    \begin{itemize}
      \item Sub-item 1
      \item Sub-item 2
    \end{itemize}
\end{itemize}

\subsection{Numbered List}
\begin{enumerate}
  \item First step
  \item Second step
  \item Third step
\end{enumerate}

\section{Citations and References}
When you use the citation search feature, references will be automatically managed. For example, you might cite a paper like this: \cite{example2024}. All citations will be automatically added to your \texttt{cite.bib} file.

\section{Conclusion}
This template demonstrates basic \LaTeX{} features including mathematical formulas, tables, and lists. The bibliography system is now set up to work seamlessly with Bibby AI's citation management. Feel free to modify it for your needs!

\bibliographystyle{plain}
\bibliography{cite}

\end{document}\documentclass{article}
\usepackage[utf8]{inputenc}
\usepackage{booktabs}
\usepackage{amsmath}
\usepackage{natbib}

\title{My Document}
\author{Your Name}
\date{\today}

\begin{document}

\maketitle

\section{Introduction}
This is a sample document demonstrating various \LaTeX{} features including mathematical equations, tables, and formatting. You can cite papers using the citation search feature, and they will be automatically added to your \texttt{cite.bib} file.

\section{Mathematical Expressions}
Here are some example equations:

\subsection{Quadratic Formula}
The quadratic formula is given by:
\[x = \frac{-b \pm \sqrt{b^2 - 4ac}}{2a}\]

\subsection{Matrix Example}
A sample matrix:
\[
\begin{pmatrix}
a_{11} & a_{12} & a_{13} \\
a_{21} & a_{22} & a_{23} \\
a_{31} & a_{32} & a_{33}
\end{pmatrix}
\]

\section{Tables}
Below is a sample table showing data:

\begin{table}[h]
\centering
\begin{tabular}{@{}lrrr@{}}
\toprule
\textbf{Category} & \textbf{Value 1} & \textbf{Value 2} & \textbf{Total} \\
\midrule
Item A & 10.5 & 20.3 & 30.8 \\
Item B & 15.2 & 25.1 & 40.3 \\
Item C & 22.1 & 12.7 & 34.8 \\
\bottomrule
\end{tabular}
\caption{Sample Data Table}
\label{tab:sample}
\end{table}

\section{Lists}
\subsection{Bullet Points}
\begin{itemize}
  \item First item
  \item Second item
  \item Third item with sub-items:
    \begin{itemize}
      \item Sub-item 1
      \item Sub-item 2
    \end{itemize}
\end{itemize}

\subsection{Numbered List}
\begin{enumerate}
  \item First step
  \item Second step
  \item Third step
\end{enumerate}

\section{Citations and References}
When you use the citation search feature, references will be automatically managed. For example, you might cite a paper like this: \cite{example2024}. All citations will be automatically added to your \texttt{cite.bib} file.

\section{Conclusion}
This template demonstrates basic \LaTeX{} features including mathematical formulas, tables, and lists. The bibliography system is now set up to work seamlessly with Bibby AI's citation management. Feel free to modify it for your needs!

\bibliographystyle{plain}
\bibliography{cite}

\end{document}\documentclass{article}
\usepackage[utf8]{inputenc}
\usepackage{booktabs}
\usepackage{amsmath}
\usepackage{natbib}

\title{My Document}
\author{Your Name}
\date{\today}

\begin{document}

\maketitle

\section{Introduction}
This is a sample document demonstrating various \LaTeX{} features including mathematical equations, tables, and formatting. You can cite papers using the citation search feature, and they will be automatically added to your \texttt{cite.bib} file.

\section{Mathematical Expressions}
Here are some example equations:

\subsection{Quadratic Formula}
The quadratic formula is given by:
\[x = \frac{-b \pm \sqrt{b^2 - 4ac}}{2a}\]

\subsection{Matrix Example}
A sample matrix:
\[
\begin{pmatrix}
a_{11} & a_{12} & a_{13} \\
a_{21} & a_{22} & a_{23} \\
a_{31} & a_{32} & a_{33}
\end{pmatrix}
\]

\section{Tables}
Below is a sample table showing data:

\begin{table}[h]
\centering
\begin{tabular}{@{}lrrr@{}}
\toprule
\textbf{Category} & \textbf{Value 1} & \textbf{Value 2} & \textbf{Total} \\
\midrule
Item A & 10.5 & 20.3 & 30.8 \\
Item B & 15.2 & 25.1 & 40.3 \\
Item C & 22.1 & 12.7 & 34.8 \\
\bottomrule
\end{tabular}
\caption{Sample Data Table}
\label{tab:sample}
\end{table}

\section{Lists}
\subsection{Bullet Points}
\begin{itemize}
  \item First item
  \item Second item
  \item Third item with sub-items:
    \begin{itemize}
      \item Sub-item 1
      \item Sub-item 2
    \end{itemize}
\end{itemize}

\subsection{Numbered List}
\begin{enumerate}
  \item First step
  \item Second step
  \item Third step
\end{enumerate}

\section{Citations and References}
When you use the citation search feature, references will be automatically managed. For example, you might cite a paper like this: \cite{example2024}. All citations will be automatically added to your \texttt{cite.bib} file.

\section{Conclusion}
This template demonstrates basic \LaTeX{} features including mathematical formulas, tables, and lists. The bibliography system is now set up to work seamlessly with Bibby AI's citation management. Feel free to modify it for your needs!

\bibliographystyle{plain}
\bibliography{cite}

\end{document}\documentclass{article}
\usepackage[utf8]{inputenc}
\usepackage{booktabs}
\usepackage{amsmath}
\usepackage{natbib}

\title{My Document}
\author{Your Name}
\date{\today}

\begin{document}

\maketitle

\section{Introduction}
This is a sample document demonstrating various \LaTeX{} features including mathematical equations, tables, and formatting. You can cite papers using the citation search feature, and they will be automatically added to your \texttt{cite.bib} file.

\section{Mathematical Expressions}
Here are some example equations:

\subsection{Quadratic Formula}
The quadratic formula is given by:
\[x = \frac{-b \pm \sqrt{b^2 - 4ac}}{2a}\]

\subsection{Matrix Example}
A sample matrix:
\[
\begin{pmatrix}
a_{11} & a_{12} & a_{13} \\
a_{21} & a_{22} & a_{23} \\
a_{31} & a_{32} & a_{33}
\end{pmatrix}
\]

\section{Tables}
Below is a sample table showing data:

\begin{table}[h]
\centering
\begin{tabular}{@{}lrrr@{}}
\toprule
\textbf{Category} & \textbf{Value 1} & \textbf{Value 2} & \textbf{Total} \\
\midrule
Item A & 10.5 & 20.3 & 30.8 \\
Item B & 15.2 & 25.1 & 40.3 \\
Item C & 22.1 & 12.7 & 34.8 \\
\bottomrule
\end{tabular}
\caption{Sample Data Table}
\label{tab:sample}
\end{table}

\section{Lists}
\subsection{Bullet Points}
\begin{itemize}
  \item First item
  \item Second item
  \item Third item with sub-items:
    \begin{itemize}
      \item Sub-item 1
      \item Sub-item 2
    \end{itemize}
\end{itemize}

\subsection{Numbered List}
\begin{enumerate}
  \item First step
  \item Second step
  \item Third step
\end{enumerate}

\section{Citations and References}
When you use the citation search feature, references will be automatically managed. For example, you might cite a paper like this: \cite{example2024}. All citations will be automatically added to your \texttt{cite.bib} file.

\section{Conclusion}
This template demonstrates basic \LaTeX{} features including mathematical formulas, tables, and lists. The bibliography system is now set up to work seamlessly with Bibby AI's citation management. Feel free to modify it for your needs!

\bibliographystyle{plain}
\bibliography{cite}

\end{document}\documentclass{article}
\usepackage[utf8]{inputenc}
\usepackage{booktabs}
\usepackage{amsmath}
\usepackage{natbib}

\title{My Document}
\author{Your Name}
\date{\today}

\begin{document}

\maketitle

\section{Introduction}
This is a sample document demonstrating various \LaTeX{} features including mathematical equations, tables, and formatting. You can cite papers using the citation search feature, and they will be automatically added to your \texttt{cite.bib} file.

\section{Mathematical Expressions}
Here are some example equations:

\subsection{Quadratic Formula}
The quadratic formula is given by:
\[x = \frac{-b \pm \sqrt{b^2 - 4ac}}{2a}\]

\subsection{Matrix Example}
A sample matrix:
\[
\begin{pmatrix}
a_{11} & a_{12} & a_{13} \\
a_{21} & a_{22} & a_{23} \\
a_{31} & a_{32} & a_{33}
\end{pmatrix}
\]

\section{Tables}
Below is a sample table showing data:

\begin{table}[h]
\centering
\begin{tabular}{@{}lrrr@{}}
\toprule
\textbf{Category} & \textbf{Value 1} & \textbf{Value 2} & \textbf{Total} \\
\midrule
Item A & 10.5 & 20.3 & 30.8 \\
Item B & 15.2 & 25.1 & 40.3 \\
Item C & 22.1 & 12.7 & 34.8 \\
\bottomrule
\end{tabular}
\caption{Sample Data Table}
\label{tab:sample}
\end{table}

\section{Lists}
\subsection{Bullet Points}
\begin{itemize}
  \item First item
  \item Second item
  \item Third item with sub-items:
    \begin{itemize}
      \item Sub-item 1
      \item Sub-item 2
    \end{itemize}
\end{itemize}

\subsection{Numbered List}
\begin{enumerate}
  \item First step
  \item Second step
  \item Third step
\end{enumerate}

\section{Citations and References}
When you use the citation search feature, references will be automatically managed. For example, you might cite a paper like this: \cite{example2024}. All citations will be automatically added to your \texttt{cite.bib} file.

\section{Conclusion}
This template demonstrates basic \LaTeX{} features including mathematical formulas, tables, and lists. The bibliography system is now set up to work seamlessly with Bibby AI's citation management. Feel free to modify it for your needs!

\bibliographystyle{plain}
\bibliography{cite}

\end{document}\documentclass{article}
\usepackage[utf8]{inputenc}
\usepackage{booktabs}
\usepackage{amsmath}
\usepackage{natbib}

\title{My Document}
\author{Your Name}
\date{\today}

\begin{document}

\maketitle

\section{Introduction}
This is a sample document demonstrating various \LaTeX{} features including mathematical equations, tables, and formatting. You can cite papers using the citation search feature, and they will be automatically added to your \texttt{cite.bib} file.

\section{Mathematical Expressions}
Here are some example equations:

\subsection{Quadratic Formula}
The quadratic formula is given by:
\[x = \frac{-b \pm \sqrt{b^2 - 4ac}}{2a}\]

\subsection{Matrix Example}
A sample matrix:
\[
\begin{pmatrix}
a_{11} & a_{12} & a_{13} \\
a_{21} & a_{22} & a_{23} \\
a_{31} & a_{32} & a_{33}
\end{pmatrix}
\]

\section{Tables}
Below is a sample table showing data:

\begin{table}[h]
\centering
\begin{tabular}{@{}lrrr@{}}
\toprule
\textbf{Category} & \textbf{Value 1} & \textbf{Value 2} & \textbf{Total} \\
\midrule
Item A & 10.5 & 20.3 & 30.8 \\
Item B & 15.2 & 25.1 & 40.3 \\
Item C & 22.1 & 12.7 & 34.8 \\
\bottomrule
\end{tabular}
\caption{Sample Data Table}
\label{tab:sample}
\end{table}

\section{Lists}
\subsection{Bullet Points}
\begin{itemize}
  \item First item
  \item Second item
  \item Third item with sub-items:
    \begin{itemize}
      \item Sub-item 1
      \item Sub-item 2
    \end{itemize}
\end{itemize}

\subsection{Numbered List}
\begin{enumerate}
  \item First step
  \item Second step
  \item Third step
\end{enumerate}

\section{Citations and References}
When you use the citation search feature, references will be automatically managed. For example, you might cite a paper like this: \cite{example2024}. All citations will be automatically added to your \texttt{cite.bib} file.

\section{Conclusion}
This template demonstrates basic \LaTeX{} features including mathematical formulas, tables, and lists. The bibliography system is now set up to work seamlessly with Bibby AI's citation management. Feel free to modify it for your needs!

\bibliographystyle{plain}
\bibliography{cite}

\end{document}\documentclass{article}
\usepackage[utf8]{inputenc}
\usepackage{booktabs}
\usepackage{amsmath}
\usepackage{natbib}

\title{My Document}
\author{Your Name}
\date{\today}

\begin{document}

\maketitle

\section{Introduction}
This is a sample document demonstrating various \LaTeX{} features including mathematical equations, tables, and formatting. You can cite papers using the citation search feature, and they will be automatically added to your \texttt{cite.bib} file.

\section{Mathematical Expressions}
Here are some example equations:

\subsection{Quadratic Formula}
The quadratic formula is given by:
\[x = \frac{-b \pm \sqrt{b^2 - 4ac}}{2a}\]

\subsection{Matrix Example}
A sample matrix:
\[
\begin{pmatrix}
a_{11} & a_{12} & a_{13} \\
a_{21} & a_{22} & a_{23} \\
a_{31} & a_{32} & a_{33}
\end{pmatrix}
\]

\section{Tables}
Below is a sample table showing data:

\begin{table}[h]
\centering
\begin{tabular}{@{}lrrr@{}}
\toprule
\textbf{Category} & \textbf{Value 1} & \textbf{Value 2} & \textbf{Total} \\
\midrule
Item A & 10.5 & 20.3 & 30.8 \\
Item B & 15.2 & 25.1 & 40.3 \\
Item C & 22.1 & 12.7 & 34.8 \\
\bottomrule
\end{tabular}
\caption{Sample Data Table}
\label{tab:sample}
\end{table}

\section{Lists}
\subsection{Bullet Points}
\begin{itemize}
  \item First item
  \item Second item
  \item Third item with sub-items:
    \begin{itemize}
      \item Sub-item 1
      \item Sub-item 2
    \end{itemize}
\end{itemize}

\subsection{Numbered List}
\begin{enumerate}
  \item First step
  \item Second step
  \item Third step
\end{enumerate}

\section{Citations and References}
When you use the citation search feature, references will be automatically managed. For example, you might cite a paper like this: \cite{example2024}. All citations will be automatically added to your \texttt{cite.bib} file.

\section{Conclusion}
This template demonstrates basic \LaTeX{} features including mathematical formulas, tables, and lists. The bibliography system is now set up to work seamlessly with Bibby AI's citation management. Feel free to modify it for your needs!

\bibliographystyle{plain}
\bibliography{cite}

\end{document}\documentclass{article}
\usepackage[utf8]{inputenc}
\usepackage{booktabs}
\usepackage{amsmath}
\usepackage{natbib}

\title{My Document}
\author{Your Name}
\date{\today}

\begin{document}

\maketitle

\section{Introduction}
This is a sample document demonstrating various \LaTeX{} features including mathematical equations, tables, and formatting. You can cite papers using the citation search feature, and they will be automatically added to your \texttt{cite.bib} file.

\section{Mathematical Expressions}
Here are some example equations:

\subsection{Quadratic Formula}
The quadratic formula is given by:
\[x = \frac{-b \pm \sqrt{b^2 - 4ac}}{2a}\]

\subsection{Matrix Example}
A sample matrix:
\[
\begin{pmatrix}
a_{11} & a_{12} & a_{13} \\
a_{21} & a_{22} & a_{23} \\
a_{31} & a_{32} & a_{33}
\end{pmatrix}
\]

\section{Tables}
Below is a sample table showing data:

\begin{table}[h]
\centering
\begin{tabular}{@{}lrrr@{}}
\toprule
\textbf{Category} & \textbf{Value 1} & \textbf{Value 2} & \textbf{Total} \\
\midrule
Item A & 10.5 & 20.3 & 30.8 \\
Item B & 15.2 & 25.1 & 40.3 \\
Item C & 22.1 & 12.7 & 34.8 \\
\bottomrule
\end{tabular}
\caption{Sample Data Table}
\label{tab:sample}
\end{table}

\section{Lists}
\subsection{Bullet Points}
\begin{itemize}
  \item First item
  \item Second item
  \item Third item with sub-items:
    \begin{itemize}
      \item Sub-item 1
      \item Sub-item 2
    \end{itemize}
\end{itemize}

\subsection{Numbered List}
\begin{enumerate}
  \item First step
  \item Second step
  \item Third step
\end{enumerate}

\section{Citations and References}
When you use the citation search feature, references will be automatically managed. For example, you might cite a paper like this: \cite{example2024}. All citations will be automatically added to your \texttt{cite.bib} file.

\section{Conclusion}
This template demonstrates basic \LaTeX{} features including mathematical formulas, tables, and lists. The bibliography system is now set up to work seamlessly with Bibby AI's citation management. Feel free to modify it for your needs!

\bibliographystyle{plain}
\bibliography{cite}

\end{document}\documentclass{article}
\usepackage[utf8]{inputenc}
\usepackage{booktabs}
\usepackage{amsmath}
\usepackage{natbib}

\title{My Document}
\author{Your Name}
\date{\today}

\begin{document}

\maketitle

\section{Introduction}
This is a sample document demonstrating various \LaTeX{} features including mathematical equations, tables, and formatting. You can cite papers using the citation search feature, and they will be automatically added to your \texttt{cite.bib} file.

\section{Mathematical Expressions}
Here are some example equations:

\subsection{Quadratic Formula}
The quadratic formula is given by:
\[x = \frac{-b \pm \sqrt{b^2 - 4ac}}{2a}\]

\subsection{Matrix Example}
A sample matrix:
\[
\begin{pmatrix}
a_{11} & a_{12} & a_{13} \\
a_{21} & a_{22} & a_{23} \\
a_{31} & a_{32} & a_{33}
\end{pmatrix}
\]

\section{Tables}
Below is a sample table showing data:

\begin{table}[h]
\centering
\begin{tabular}{@{}lrrr@{}}
\toprule
\textbf{Category} & \textbf{Value 1} & \textbf{Value 2} & \textbf{Total} \\
\midrule
Item A & 10.5 & 20.3 & 30.8 \\
Item B & 15.2 & 25.1 & 40.3 \\
Item C & 22.1 & 12.7 & 34.8 \\
\bottomrule
\end{tabular}
\caption{Sample Data Table}
\label{tab:sample}
\end{table}

\section{Lists}
\subsection{Bullet Points}
\begin{itemize}
  \item First item
  \item Second item
  \item Third item with sub-items:
    \begin{itemize}
      \item Sub-item 1
      \item Sub-item 2
    \end{itemize}
\end{itemize}

\subsection{Numbered List}
\begin{enumerate}
  \item First step
  \item Second step
  \item Third step
\end{enumerate}

\section{Citations and References}
When you use the citation search feature, references will be automatically managed. For example, you might cite a paper like this: \cite{example2024}. All citations will be automatically added to your \texttt{cite.bib} file.

\section{Conclusion}
This template demonstrates basic \LaTeX{} features including mathematical formulas, tables, and lists. The bibliography system is now set up to work seamlessly with Bibby AI's citation management. Feel free to modify it for your needs!

\bibliographystyle{plain}
\bibliography{cite}

\end{document}\documentclass{article}
\usepackage[utf8]{inputenc}
\usepackage{booktabs}
\usepackage{amsmath}
\usepackage{natbib}

\title{My Document}
\author{Your Name}
\date{\today}

\begin{document}

\maketitle

\section{Introduction}
This is a sample document demonstrating various \LaTeX{} features including mathematical equations, tables, and formatting. You can cite papers using the citation search feature, and they will be automatically added to your \texttt{cite.bib} file.

\section{Mathematical Expressions}
Here are some example equations:

\subsection{Quadratic Formula}
The quadratic formula is given by:
\[x = \frac{-b \pm \sqrt{b^2 - 4ac}}{2a}\]

\subsection{Matrix Example}
A sample matrix:
\[
\begin{pmatrix}
a_{11} & a_{12} & a_{13} \\
a_{21} & a_{22} & a_{23} \\
a_{31} & a_{32} & a_{33}
\end{pmatrix}
\]

\section{Tables}
Below is a sample table showing data:

\begin{table}[h]
\centering
\begin{tabular}{@{}lrrr@{}}
\toprule
\textbf{Category} & \textbf{Value 1} & \textbf{Value 2} & \textbf{Total} \\
\midrule
Item A & 10.5 & 20.3 & 30.8 \\
Item B & 15.2 & 25.1 & 40.3 \\
Item C & 22.1 & 12.7 & 34.8 \\
\bottomrule
\end{tabular}
\caption{Sample Data Table}
\label{tab:sample}
\end{table}

\section{Lists}
\subsection{Bullet Points}
\begin{itemize}
  \item First item
  \item Second item
  \item Third item with sub-items:
    \begin{itemize}
      \item Sub-item 1
      \item Sub-item 2
    \end{itemize}
\end{itemize}

\subsection{Numbered List}
\begin{enumerate}
  \item First step
  \item Second step
  \item Third step
\end{enumerate}

\section{Citations and References}
When you use the citation search feature, references will be automatically managed. For example, you might cite a paper like this: \cite{example2024}. All citations will be automatically added to your \texttt{cite.bib} file.

\section{Conclusion}
This template demonstrates basic \LaTeX{} features including mathematical formulas, tables, and lists. The bibliography system is now set up to work seamlessly with Bibby AI's citation management. Feel free to modify it for your needs!

\bibliographystyle{plain}
\bibliography{cite}

\end{document}\documentclass{article}
\usepackage[utf8]{inputenc}
\usepackage{booktabs}
\usepackage{amsmath}
\usepackage{natbib}

\title{My Document}
\author{Your Name}
\date{\today}

\begin{document}

\maketitle

\section{Introduction}
This is a sample document demonstrating various \LaTeX{} features including mathematical equations, tables, and formatting. You can cite papers using the citation search feature, and they will be automatically added to your \texttt{cite.bib} file.

\section{Mathematical Expressions}
Here are some example equations:

\subsection{Quadratic Formula}
The quadratic formula is given by:
\[x = \frac{-b \pm \sqrt{b^2 - 4ac}}{2a}\]

\subsection{Matrix Example}
A sample matrix:
\[
\begin{pmatrix}
a_{11} & a_{12} & a_{13} \\
a_{21} & a_{22} & a_{23} \\
a_{31} & a_{32} & a_{33}
\end{pmatrix}
\]

\section{Tables}
Below is a sample table showing data:

\begin{table}[h]
\centering
\begin{tabular}{@{}lrrr@{}}
\toprule
\textbf{Category} & \textbf{Value 1} & \textbf{Value 2} & \textbf{Total} \\
\midrule
Item A & 10.5 & 20.3 & 30.8 \\
Item B & 15.2 & 25.1 & 40.3 \\
Item C & 22.1 & 12.7 & 34.8 \\
\bottomrule
\end{tabular}
\caption{Sample Data Table}
\label{tab:sample}
\end{table}

\section{Lists}
\subsection{Bullet Points}
\begin{itemize}
  \item First item
  \item Second item
  \item Third item with sub-items:
    \begin{itemize}
      \item Sub-item 1
      \item Sub-item 2
    \end{itemize}
\end{itemize}

\subsection{Numbered List}
\begin{enumerate}
  \item First step
  \item Second step
  \item Third step
\end{enumerate}

\section{Citations and References}
When you use the citation search feature, references will be automatically managed. For example, you might cite a paper like this: \cite{example2024}. All citations will be automatically added to your \texttt{cite.bib} file.

\section{Conclusion}
This template demonstrates basic \LaTeX{} features including mathematical formulas, tables, and lists. The bibliography system is now set up to work seamlessly with Bibby AI's citation management. Feel free to modify it for your needs!

\bibliographystyle{plain}
\bibliography{cite}

\end{document}\documentclass{article}
\usepackage[utf8]{inputenc}
\usepackage{booktabs}
\usepackage{amsmath}
\usepackage{natbib}

\title{My Document}
\author{Your Name}
\date{\today}

\begin{document}

\maketitle

\section{Introduction}
This is a sample document demonstrating various \LaTeX{} features including mathematical equations, tables, and formatting. You can cite papers using the citation search feature, and they will be automatically added to your \texttt{cite.bib} file.

\section{Mathematical Expressions}
Here are some example equations:

\subsection{Quadratic Formula}
The quadratic formula is given by:
\[x = \frac{-b \pm \sqrt{b^2 - 4ac}}{2a}\]

\subsection{Matrix Example}
A sample matrix:
\[
\begin{pmatrix}
a_{11} & a_{12} & a_{13} \\
a_{21} & a_{22} & a_{23} \\
a_{31} & a_{32} & a_{33}
\end{pmatrix}
\]

\section{Tables}
Below is a sample table showing data:

\begin{table}[h]
\centering
\begin{tabular}{@{}lrrr@{}}
\toprule
\textbf{Category} & \textbf{Value 1} & \textbf{Value 2} & \textbf{Total} \\
\midrule
Item A & 10.5 & 20.3 & 30.8 \\
Item B & 15.2 & 25.1 & 40.3 \\
Item C & 22.1 & 12.7 & 34.8 \\
\bottomrule
\end{tabular}
\caption{Sample Data Table}
\label{tab:sample}
\end{table}

\section{Lists}
\subsection{Bullet Points}
\begin{itemize}
  \item First item
  \item Second item
  \item Third item with sub-items:
    \begin{itemize}
      \item Sub-item 1
      \item Sub-item 2
    \end{itemize}
\end{itemize}

\subsection{Numbered List}
\begin{enumerate}
  \item First step
  \item Second step
  \item Third step
\end{enumerate}

\section{Citations and References}
When you use the citation search feature, references will be automatically managed. For example, you might cite a paper like this: \cite{example2024}. All citations will be automatically added to your \texttt{cite.bib} file.

\section{Conclusion}
This template demonstrates basic \LaTeX{} features including mathematical formulas, tables, and lists. The bibliography system is now set up to work seamlessly with Bibby AI's citation management. Feel free to modify it for your needs!

\bibliographystyle{plain}
\bibliography{cite}

\end{document}\documentclass{article}
\usepackage[utf8]{inputenc}
\usepackage{booktabs}
\usepackage{amsmath}
\usepackage{natbib}

\title{My Document}
\author{Your Name}
\date{\today}

\begin{document}

\maketitle

\section{Introduction}
This is a sample document demonstrating various \LaTeX{} features including mathematical equations, tables, and formatting. You can cite papers using the citation search feature, and they will be automatically added to your \texttt{cite.bib} file.

\section{Mathematical Expressions}
Here are some example equations:

\subsection{Quadratic Formula}
The quadratic formula is given by:
\[x = \frac{-b \pm \sqrt{b^2 - 4ac}}{2a}\]

\subsection{Matrix Example}
A sample matrix:
\[
\begin{pmatrix}
a_{11} & a_{12} & a_{13} \\
a_{21} & a_{22} & a_{23} \\
a_{31} & a_{32} & a_{33}
\end{pmatrix}
\]

\section{Tables}
Below is a sample table showing data:

\begin{table}[h]
\centering
\begin{tabular}{@{}lrrr@{}}
\toprule
\textbf{Category} & \textbf{Value 1} & \textbf{Value 2} & \textbf{Total} \\
\midrule
Item A & 10.5 & 20.3 & 30.8 \\
Item B & 15.2 & 25.1 & 40.3 \\
Item C & 22.1 & 12.7 & 34.8 \\
\bottomrule
\end{tabular}
\caption{Sample Data Table}
\label{tab:sample}
\end{table}

\section{Lists}
\subsection{Bullet Points}
\begin{itemize}
  \item First item
  \item Second item
  \item Third item with sub-items:
    \begin{itemize}
      \item Sub-item 1
      \item Sub-item 2
    \end{itemize}
\end{itemize}

\subsection{Numbered List}
\begin{enumerate}
  \item First step
  \item Second step
  \item Third step
\end{enumerate}

\section{Citations and References}
When you use the citation search feature, references will be automatically managed. For example, you might cite a paper like this: \cite{example2024}. All citations will be automatically added to your \texttt{cite.bib} file.

\section{Conclusion}
This template demonstrates basic \LaTeX{} features including mathematical formulas, tables, and lists. The bibliography system is now set up to work seamlessly with Bibby AI's citation management. Feel free to modify it for your needs!

\bibliographystyle{plain}
\bibliography{cite}

\end{document}\documentclass{article}
\usepackage[utf8]{inputenc}
\usepackage{booktabs}
\usepackage{amsmath}
\usepackage{natbib}

\title{My Document}
\author{Your Name}
\date{\today}

\begin{document}

\maketitle

\section{Introduction}
This is a sample document demonstrating various \LaTeX{} features including mathematical equations, tables, and formatting. You can cite papers using the citation search feature, and they will be automatically added to your \texttt{cite.bib} file.

\section{Mathematical Expressions}
Here are some example equations:

\subsection{Quadratic Formula}
The quadratic formula is given by:
\[x = \frac{-b \pm \sqrt{b^2 - 4ac}}{2a}\]

\subsection{Matrix Example}
A sample matrix:
\[
\begin{pmatrix}
a_{11} & a_{12} & a_{13} \\
a_{21} & a_{22} & a_{23} \\
a_{31} & a_{32} & a_{33}
\end{pmatrix}
\]

\section{Tables}
Below is a sample table showing data:

\begin{table}[h]
\centering
\begin{tabular}{@{}lrrr@{}}
\toprule
\textbf{Category} & \textbf{Value 1} & \textbf{Value 2} & \textbf{Total} \\
\midrule
Item A & 10.5 & 20.3 & 30.8 \\
Item B & 15.2 & 25.1 & 40.3 \\
Item C & 22.1 & 12.7 & 34.8 \\
\bottomrule
\end{tabular}
\caption{Sample Data Table}
\label{tab:sample}
\end{table}

\section{Lists}
\subsection{Bullet Points}
\begin{itemize}
  \item First item
  \item Second item
  \item Third item with sub-items:
    \begin{itemize}
      \item Sub-item 1
      \item Sub-item 2
    \end{itemize}
\end{itemize}

\subsection{Numbered List}
\begin{enumerate}
  \item First step
  \item Second step
  \item Third step
\end{enumerate}

\section{Citations and References}
When you use the citation search feature, references will be automatically managed. For example, you might cite a paper like this: \cite{example2024}. All citations will be automatically added to your \texttt{cite.bib} file.

\section{Conclusion}
This template demonstrates basic \LaTeX{} features including mathematical formulas, tables, and lists. The bibliography system is now set up to work seamlessly with Bibby AI's citation management. Feel free to modify it for your needs!

\bibliographystyle{plain}
\bibliography{cite}

\end{document}\documentclass{article}
\usepackage[utf8]{inputenc}
\usepackage{booktabs}
\usepackage{amsmath}
\usepackage{natbib}

\title{My Document}
\author{Your Name}
\date{\today}

\begin{document}

\maketitle

\section{Introduction}
This is a sample document demonstrating various \LaTeX{} features including mathematical equations, tables, and formatting. You can cite papers using the citation search feature, and they will be automatically added to your \texttt{cite.bib} file.

\section{Mathematical Expressions}
Here are some example equations:

\subsection{Quadratic Formula}
The quadratic formula is given by:
\[x = \frac{-b \pm \sqrt{b^2 - 4ac}}{2a}\]

\subsection{Matrix Example}
A sample matrix:
\[
\begin{pmatrix}
a_{11} & a_{12} & a_{13} \\
a_{21} & a_{22} & a_{23} \\
a_{31} & a_{32} & a_{33}
\end{pmatrix}
\]

\section{Tables}
Below is a sample table showing data:

\begin{table}[h]
\centering
\begin{tabular}{@{}lrrr@{}}
\toprule
\textbf{Category} & \textbf{Value 1} & \textbf{Value 2} & \textbf{Total} \\
\midrule
Item A & 10.5 & 20.3 & 30.8 \\
Item B & 15.2 & 25.1 & 40.3 \\
Item C & 22.1 & 12.7 & 34.8 \\
\bottomrule
\end{tabular}
\caption{Sample Data Table}
\label{tab:sample}
\end{table}

\section{Lists}
\subsection{Bullet Points}
\begin{itemize}
  \item First item
  \item Second item
  \item Third item with sub-items:
    \begin{itemize}
      \item Sub-item 1
      \item Sub-item 2
    \end{itemize}
\end{itemize}

\subsection{Numbered List}
\begin{enumerate}
  \item First step
  \item Second step
  \item Third step
\end{enumerate}

\section{Citations and References}
When you use the citation search feature, references will be automatically managed. For example, you might cite a paper like this: \cite{example2024}. All citations will be automatically added to your \texttt{cite.bib} file.

\section{Conclusion}
This template demonstrates basic \LaTeX{} features including mathematical formulas, tables, and lists. The bibliography system is now set up to work seamlessly with Bibby AI's citation management. Feel free to modify it for your needs!

\bibliographystyle{plain}
\bibliography{cite}

\end{document}\documentclass{article}
\usepackage[utf8]{inputenc}
\usepackage{booktabs}
\usepackage{amsmath}
\usepackage{natbib}

\title{My Document}
\author{Your Name}
\date{\today}

\begin{document}

\maketitle

\section{Introduction}
This is a sample document demonstrating various \LaTeX{} features including mathematical equations, tables, and formatting. You can cite papers using the citation search feature, and they will be automatically added to your \texttt{cite.bib} file.

\section{Mathematical Expressions}
Here are some example equations:

\subsection{Quadratic Formula}
The quadratic formula is given by:
\[x = \frac{-b \pm \sqrt{b^2 - 4ac}}{2a}\]

\subsection{Matrix Example}
A sample matrix:
\[
\begin{pmatrix}
a_{11} & a_{12} & a_{13} \\
a_{21} & a_{22} & a_{23} \\
a_{31} & a_{32} & a_{33}
\end{pmatrix}
\]

\section{Tables}
Below is a sample table showing data:

\begin{table}[h]
\centering
\begin{tabular}{@{}lrrr@{}}
\toprule
\textbf{Category} & \textbf{Value 1} & \textbf{Value 2} & \textbf{Total} \\
\midrule
Item A & 10.5 & 20.3 & 30.8 \\
Item B & 15.2 & 25.1 & 40.3 \\
Item C & 22.1 & 12.7 & 34.8 \\
\bottomrule
\end{tabular}
\caption{Sample Data Table}
\label{tab:sample}
\end{table}

\section{Lists}
\subsection{Bullet Points}
\begin{itemize}
  \item First item
  \item Second item
  \item Third item with sub-items:
    \begin{itemize}
      \item Sub-item 1
      \item Sub-item 2
    \end{itemize}
\end{itemize}

\subsection{Numbered List}
\begin{enumerate}
  \item First step
  \item Second step
  \item Third step
\end{enumerate}

\section{Citations and References}
When you use the citation search feature, references will be automatically managed. For example, you might cite a paper like this: \cite{example2024}. All citations will be automatically added to your \texttt{cite.bib} file.

\section{Conclusion}
This template demonstrates basic \LaTeX{} features including mathematical formulas, tables, and lists. The bibliography system is now set up to work seamlessly with Bibby AI's citation management. Feel free to modify it for your needs!

\bibliographystyle{plain}
\bibliography{cite}

\end{document}\documentclass{article}
\usepackage[utf8]{inputenc}
\usepackage{booktabs}
\usepackage{amsmath}
\usepackage{natbib}

\title{My Document}
\author{Your Name}
\date{\today}

\begin{document}

\maketitle

\section{Introduction}
This is a sample document demonstrating various \LaTeX{} features including mathematical equations, tables, and formatting. You can cite papers using the citation search feature, and they will be automatically added to your \texttt{cite.bib} file.

\section{Mathematical Expressions}
Here are some example equations:

\subsection{Quadratic Formula}
The quadratic formula is given by:
\[x = \frac{-b \pm \sqrt{b^2 - 4ac}}{2a}\]

\subsection{Matrix Example}
A sample matrix:
\[
\begin{pmatrix}
a_{11} & a_{12} & a_{13} \\
a_{21} & a_{22} & a_{23} \\
a_{31} & a_{32} & a_{33}
\end{pmatrix}
\]

\section{Tables}
Below is a sample table showing data:

\begin{table}[h]
\centering
\begin{tabular}{@{}lrrr@{}}
\toprule
\textbf{Category} & \textbf{Value 1} & \textbf{Value 2} & \textbf{Total} \\
\midrule
Item A & 10.5 & 20.3 & 30.8 \\
Item B & 15.2 & 25.1 & 40.3 \\
Item C & 22.1 & 12.7 & 34.8 \\
\bottomrule
\end{tabular}
\caption{Sample Data Table}
\label{tab:sample}
\end{table}

\section{Lists}
\subsection{Bullet Points}
\begin{itemize}
  \item First item
  \item Second item
  \item Third item with sub-items:
    \begin{itemize}
      \item Sub-item 1
      \item Sub-item 2
    \end{itemize}
\end{itemize}

\subsection{Numbered List}
\begin{enumerate}
  \item First step
  \item Second step
  \item Third step
\end{enumerate}

\section{Citations and References}
When you use the citation search feature, references will be automatically managed. For example, you might cite a paper like this: \cite{example2024}. All citations will be automatically added to your \texttt{cite.bib} file.

\section{Conclusion}
This template demonstrates basic \LaTeX{} features including mathematical formulas, tables, and lists. The bibliography system is now set up to work seamlessly with Bibby AI's citation management. Feel free to modify it for your needs!

\bibliographystyle{plain}
\bibliography{cite}

\end{document}\documentclass{article}
\usepackage[utf8]{inputenc}
\usepackage{booktabs}
\usepackage{amsmath}
\usepackage{natbib}

\title{My Document}
\author{Your Name}
\date{\today}

\begin{document}

\maketitle

\section{Introduction}
This is a sample document demonstrating various \LaTeX{} features including mathematical equations, tables, and formatting. You can cite papers using the citation search feature, and they will be automatically added to your \texttt{cite.bib} file.

\section{Mathematical Expressions}
Here are some example equations:

\subsection{Quadratic Formula}
The quadratic formula is given by:
\[x = \frac{-b \pm \sqrt{b^2 - 4ac}}{2a}\]

\subsection{Matrix Example}
A sample matrix:
\[
\begin{pmatrix}
a_{11} & a_{12} & a_{13} \\
a_{21} & a_{22} & a_{23} \\
a_{31} & a_{32} & a_{33}
\end{pmatrix}
\]

\section{Tables}
Below is a sample table showing data:

\begin{table}[h]
\centering
\begin{tabular}{@{}lrrr@{}}
\toprule
\textbf{Category} & \textbf{Value 1} & \textbf{Value 2} & \textbf{Total} \\
\midrule
Item A & 10.5 & 20.3 & 30.8 \\
Item B & 15.2 & 25.1 & 40.3 \\
Item C & 22.1 & 12.7 & 34.8 \\
\bottomrule
\end{tabular}
\caption{Sample Data Table}
\label{tab:sample}
\end{table}

\section{Lists}
\subsection{Bullet Points}
\begin{itemize}
  \item First item
  \item Second item
  \item Third item with sub-items:
    \begin{itemize}
      \item Sub-item 1
      \item Sub-item 2
    \end{itemize}
\end{itemize}

\subsection{Numbered List}
\begin{enumerate}
  \item First step
  \item Second step
  \item Third step
\end{enumerate}

\section{Citations and References}
When you use the citation search feature, references will be automatically managed. For example, you might cite a paper like this: \cite{example2024}. All citations will be automatically added to your \texttt{cite.bib} file.

\section{Conclusion}
This template demonstrates basic \LaTeX{} features including mathematical formulas, tables, and lists. The bibliography system is now set up to work seamlessly with Bibby AI's citation management. Feel free to modify it for your needs!

\bibliographystyle{plain}
\bibliography{cite}

\end{document}\documentclass{article}
\usepackage[utf8]{inputenc}
\usepackage{booktabs}
\usepackage{amsmath}
\usepackage{natbib}

\title{My Document}
\author{Your Name}
\date{\today}

\begin{document}

\maketitle

\section{Introduction}
This is a sample document demonstrating various \LaTeX{} features including mathematical equations, tables, and formatting. You can cite papers using the citation search feature, and they will be automatically added to your \texttt{cite.bib} file.

\section{Mathematical Expressions}
Here are some example equations:

\subsection{Quadratic Formula}
The quadratic formula is given by:
\[x = \frac{-b \pm \sqrt{b^2 - 4ac}}{2a}\]

\subsection{Matrix Example}
A sample matrix:
\[
\begin{pmatrix}
a_{11} & a_{12} & a_{13} \\
a_{21} & a_{22} & a_{23} \\
a_{31} & a_{32} & a_{33}
\end{pmatrix}
\]

\section{Tables}
Below is a sample table showing data:

\begin{table}[h]
\centering
\begin{tabular}{@{}lrrr@{}}
\toprule
\textbf{Category} & \textbf{Value 1} & \textbf{Value 2} & \textbf{Total} \\
\midrule
Item A & 10.5 & 20.3 & 30.8 \\
Item B & 15.2 & 25.1 & 40.3 \\
Item C & 22.1 & 12.7 & 34.8 \\
\bottomrule
\end{tabular}
\caption{Sample Data Table}
\label{tab:sample}
\end{table}

\section{Lists}
\subsection{Bullet Points}
\begin{itemize}
  \item First item
  \item Second item
  \item Third item with sub-items:
    \begin{itemize}
      \item Sub-item 1
      \item Sub-item 2
    \end{itemize}
\end{itemize}

\subsection{Numbered List}
\begin{enumerate}
  \item First step
  \item Second step
  \item Third step
\end{enumerate}

\section{Citations and References}
When you use the citation search feature, references will be automatically managed. For example, you might cite a paper like this: \cite{example2024}. All citations will be automatically added to your \texttt{cite.bib} file.

\section{Conclusion}
This template demonstrates basic \LaTeX{} features including mathematical formulas, tables, and lists. The bibliography system is now set up to work seamlessly with Bibby AI's citation management. Feel free to modify it for your needs!

\bibliographystyle{plain}
\bibliography{cite}

\end{document}\documentclass{article}
\usepackage[utf8]{inputenc}
\usepackage{booktabs}
\usepackage{amsmath}
\usepackage{natbib}

\title{My Document}
\author{Your Name}
\date{\today}

\begin{document}

\maketitle

\section{Introduction}
This is a sample document demonstrating various \LaTeX{} features including mathematical equations, tables, and formatting. You can cite papers using the citation search feature, and they will be automatically added to your \texttt{cite.bib} file.

\section{Mathematical Expressions}
Here are some example equations:

\subsection{Quadratic Formula}
The quadratic formula is given by:
\[x = \frac{-b \pm \sqrt{b^2 - 4ac}}{2a}\]

\subsection{Matrix Example}
A sample matrix:
\[
\begin{pmatrix}
a_{11} & a_{12} & a_{13} \\
a_{21} & a_{22} & a_{23} \\
a_{31} & a_{32} & a_{33}
\end{pmatrix}
\]

\section{Tables}
Below is a sample table showing data:

\begin{table}[h]
\centering
\begin{tabular}{@{}lrrr@{}}
\toprule
\textbf{Category} & \textbf{Value 1} & \textbf{Value 2} & \textbf{Total} \\
\midrule
Item A & 10.5 & 20.3 & 30.8 \\
Item B & 15.2 & 25.1 & 40.3 \\
Item C & 22.1 & 12.7 & 34.8 \\
\bottomrule
\end{tabular}
\caption{Sample Data Table}
\label{tab:sample}
\end{table}

\section{Lists}
\subsection{Bullet Points}
\begin{itemize}
  \item First item
  \item Second item
  \item Third item with sub-items:
    \begin{itemize}
      \item Sub-item 1
      \item Sub-item 2
    \end{itemize}
\end{itemize}

\subsection{Numbered List}
\begin{enumerate}
  \item First step
  \item Second step
  \item Third step
\end{enumerate}

\section{Citations and References}
When you use the citation search feature, references will be automatically managed. For example, you might cite a paper like this: \cite{example2024}. All citations will be automatically added to your \texttt{cite.bib} file.

\section{Conclusion}
This template demonstrates basic \LaTeX{} features including mathematical formulas, tables, and lists. The bibliography system is now set up to work seamlessly with Bibby AI's citation management. Feel free to modify it for your needs!

\bibliographystyle{plain}
\bibliography{cite}

\end{document}\documentclass{article}
\usepackage[utf8]{inputenc}
\usepackage{booktabs}
\usepackage{amsmath}
\usepackage{natbib}

\title{My Document}
\author{Your Name}
\date{\today}

\begin{document}

\maketitle

\section{Introduction}
This is a sample document demonstrating various \LaTeX{} features including mathematical equations, tables, and formatting. You can cite papers using the citation search feature, and they will be automatically added to your \texttt{cite.bib} file.

\section{Mathematical Expressions}
Here are some example equations:

\subsection{Quadratic Formula}
The quadratic formula is given by:
\[x = \frac{-b \pm \sqrt{b^2 - 4ac}}{2a}\]

\subsection{Matrix Example}
A sample matrix:
\[
\begin{pmatrix}
a_{11} & a_{12} & a_{13} \\
a_{21} & a_{22} & a_{23} \\
a_{31} & a_{32} & a_{33}
\end{pmatrix}
\]

\section{Tables}
Below is a sample table showing data:

\begin{table}[h]
\centering
\begin{tabular}{@{}lrrr@{}}
\toprule
\textbf{Category} & \textbf{Value 1} & \textbf{Value 2} & \textbf{Total} \\
\midrule
Item A & 10.5 & 20.3 & 30.8 \\
Item B & 15.2 & 25.1 & 40.3 \\
Item C & 22.1 & 12.7 & 34.8 \\
\bottomrule
\end{tabular}
\caption{Sample Data Table}
\label{tab:sample}
\end{table}

\section{Lists}
\subsection{Bullet Points}
\begin{itemize}
  \item First item
  \item Second item
  \item Third item with sub-items:
    \begin{itemize}
      \item Sub-item 1
      \item Sub-item 2
    \end{itemize}
\end{itemize}

\subsection{Numbered List}
\begin{enumerate}
  \item First step
  \item Second step
  \item Third step
\end{enumerate}

\section{Citations and References}
When you use the citation search feature, references will be automatically managed. For example, you might cite a paper like this: \cite{example2024}. All citations will be automatically added to your \texttt{cite.bib} file.

\section{Conclusion}
This template demonstrates basic \LaTeX{} features including mathematical formulas, tables, and lists. The bibliography system is now set up to work seamlessly with Bibby AI's citation management. Feel free to modify it for your needs!

\bibliographystyle{plain}
\bibliography{cite}

\end{document}\documentclass{article}
\usepackage[utf8]{inputenc}
\usepackage{booktabs}
\usepackage{amsmath}
\usepackage{natbib}

\title{My Document}
\author{Your Name}
\date{\today}

\begin{document}

\maketitle

\section{Introduction}
This is a sample document demonstrating various \LaTeX{} features including mathematical equations, tables, and formatting. You can cite papers using the citation search feature, and they will be automatically added to your \texttt{cite.bib} file.

\section{Mathematical Expressions}
Here are some example equations:

\subsection{Quadratic Formula}
The quadratic formula is given by:
\[x = \frac{-b \pm \sqrt{b^2 - 4ac}}{2a}\]

\subsection{Matrix Example}
A sample matrix:
\[
\begin{pmatrix}
a_{11} & a_{12} & a_{13} \\
a_{21} & a_{22} & a_{23} \\
a_{31} & a_{32} & a_{33}
\end{pmatrix}
\]

\section{Tables}
Below is a sample table showing data:

\begin{table}[h]
\centering
\begin{tabular}{@{}lrrr@{}}
\toprule
\textbf{Category} & \textbf{Value 1} & \textbf{Value 2} & \textbf{Total} \\
\midrule
Item A & 10.5 & 20.3 & 30.8 \\
Item B & 15.2 & 25.1 & 40.3 \\
Item C & 22.1 & 12.7 & 34.8 \\
\bottomrule
\end{tabular}
\caption{Sample Data Table}
\label{tab:sample}
\end{table}

\section{Lists}
\subsection{Bullet Points}
\begin{itemize}
  \item First item
  \item Second item
  \item Third item with sub-items:
    \begin{itemize}
      \item Sub-item 1
      \item Sub-item 2
    \end{itemize}
\end{itemize}

\subsection{Numbered List}
\begin{enumerate}
  \item First step
  \item Second step
  \item Third step
\end{enumerate}

\section{Citations and References}
When you use the citation search feature, references will be automatically managed. For example, you might cite a paper like this: \cite{example2024}. All citations will be automatically added to your \texttt{cite.bib} file.

\section{Conclusion}
This template demonstrates basic \LaTeX{} features including mathematical formulas, tables, and lists. The bibliography system is now set up to work seamlessly with Bibby AI's citation management. Feel free to modify it for your needs!

\bibliographystyle{plain}
\bibliography{cite}

\end{document}\documentclass{article}
\usepackage[utf8]{inputenc}
\usepackage{booktabs}
\usepackage{amsmath}
\usepackage{natbib}

\title{My Document}
\author{Your Name}
\date{\today}

\begin{document}

\maketitle

\section{Introduction}
This is a sample document demonstrating various \LaTeX{} features including mathematical equations, tables, and formatting. You can cite papers using the citation search feature, and they will be automatically added to your \texttt{cite.bib} file.

\section{Mathematical Expressions}
Here are some example equations:

\subsection{Quadratic Formula}
The quadratic formula is given by:
\[x = \frac{-b \pm \sqrt{b^2 - 4ac}}{2a}\]

\subsection{Matrix Example}
A sample matrix:
\[
\begin{pmatrix}
a_{11} & a_{12} & a_{13} \\
a_{21} & a_{22} & a_{23} \\
a_{31} & a_{32} & a_{33}
\end{pmatrix}
\]

\section{Tables}
Below is a sample table showing data:

\begin{table}[h]
\centering
\begin{tabular}{@{}lrrr@{}}
\toprule
\textbf{Category} & \textbf{Value 1} & \textbf{Value 2} & \textbf{Total} \\
\midrule
Item A & 10.5 & 20.3 & 30.8 \\
Item B & 15.2 & 25.1 & 40.3 \\
Item C & 22.1 & 12.7 & 34.8 \\
\bottomrule
\end{tabular}
\caption{Sample Data Table}
\label{tab:sample}
\end{table}

\section{Lists}
\subsection{Bullet Points}
\begin{itemize}
  \item First item
  \item Second item
  \item Third item with sub-items:
    \begin{itemize}
      \item Sub-item 1
      \item Sub-item 2
    \end{itemize}
\end{itemize}

\subsection{Numbered List}
\begin{enumerate}
  \item First step
  \item Second step
  \item Third step
\end{enumerate}

\section{Citations and References}
When you use the citation search feature, references will be automatically managed. For example, you might cite a paper like this: \cite{example2024}. All citations will be automatically added to your \texttt{cite.bib} file.

\section{Conclusion}
This template demonstrates basic \LaTeX{} features including mathematical formulas, tables, and lists. The bibliography system is now set up to work seamlessly with Bibby AI's citation management. Feel free to modify it for your needs!

\bibliographystyle{plain}
\bibliography{cite}

\end{document}\documentclass{article}
\usepackage[utf8]{inputenc}
\usepackage{booktabs}
\usepackage{amsmath}
\usepackage{natbib}

\title{My Document}
\author{Your Name}
\date{\today}

\begin{document}

\maketitle

\section{Introduction}
This is a sample document demonstrating various \LaTeX{} features including mathematical equations, tables, and formatting. You can cite papers using the citation search feature, and they will be automatically added to your \texttt{cite.bib} file.

\section{Mathematical Expressions}
Here are some example equations:

\subsection{Quadratic Formula}
The quadratic formula is given by:
\[x = \frac{-b \pm \sqrt{b^2 - 4ac}}{2a}\]

\subsection{Matrix Example}
A sample matrix:
\[
\begin{pmatrix}
a_{11} & a_{12} & a_{13} \\
a_{21} & a_{22} & a_{23} \\
a_{31} & a_{32} & a_{33}
\end{pmatrix}
\]

\section{Tables}
Below is a sample table showing data:

\begin{table}[h]
\centering
\begin{tabular}{@{}lrrr@{}}
\toprule
\textbf{Category} & \textbf{Value 1} & \textbf{Value 2} & \textbf{Total} \\
\midrule
Item A & 10.5 & 20.3 & 30.8 \\
Item B & 15.2 & 25.1 & 40.3 \\
Item C & 22.1 & 12.7 & 34.8 \\
\bottomrule
\end{tabular}
\caption{Sample Data Table}
\label{tab:sample}
\end{table}

\section{Lists}
\subsection{Bullet Points}
\begin{itemize}
  \item First item
  \item Second item
  \item Third item with sub-items:
    \begin{itemize}
      \item Sub-item 1
      \item Sub-item 2
    \end{itemize}
\end{itemize}

\subsection{Numbered List}
\begin{enumerate}
  \item First step
  \item Second step
  \item Third step
\end{enumerate}

\section{Citations and References}
When you use the citation search feature, references will be automatically managed. For example, you might cite a paper like this: \cite{example2024}. All citations will be automatically added to your \texttt{cite.bib} file.

\section{Conclusion}
This template demonstrates basic \LaTeX{} features including mathematical formulas, tables, and lists. The bibliography system is now set up to work seamlessly with Bibby AI's citation management. Feel free to modify it for your needs!

\bibliographystyle{plain}
\bibliography{cite}

\end{document}\documentclass{article}
\usepackage[utf8]{inputenc}
\usepackage{booktabs}
\usepackage{amsmath}
\usepackage{natbib}

\title{My Document}
\author{Your Name}
\date{\today}

\begin{document}

\maketitle

\section{Introduction}
This is a sample document demonstrating various \LaTeX{} features including mathematical equations, tables, and formatting. You can cite papers using the citation search feature, and they will be automatically added to your \texttt{cite.bib} file.

\section{Mathematical Expressions}
Here are some example equations:

\subsection{Quadratic Formula}
The quadratic formula is given by:
\[x = \frac{-b \pm \sqrt{b^2 - 4ac}}{2a}\]

\subsection{Matrix Example}
A sample matrix:
\[
\begin{pmatrix}
a_{11} & a_{12} & a_{13} \\
a_{21} & a_{22} & a_{23} \\
a_{31} & a_{32} & a_{33}
\end{pmatrix}
\]

\section{Tables}
Below is a sample table showing data:

\begin{table}[h]
\centering
\begin{tabular}{@{}lrrr@{}}
\toprule
\textbf{Category} & \textbf{Value 1} & \textbf{Value 2} & \textbf{Total} \\
\midrule
Item A & 10.5 & 20.3 & 30.8 \\
Item B & 15.2 & 25.1 & 40.3 \\
Item C & 22.1 & 12.7 & 34.8 \\
\bottomrule
\end{tabular}
\caption{Sample Data Table}
\label{tab:sample}
\end{table}

\section{Lists}
\subsection{Bullet Points}
\begin{itemize}
  \item First item
  \item Second item
  \item Third item with sub-items:
    \begin{itemize}
      \item Sub-item 1
      \item Sub-item 2
    \end{itemize}
\end{itemize}

\subsection{Numbered List}
\begin{enumerate}
  \item First step
  \item Second step
  \item Third step
\end{enumerate}

\section{Citations and References}
When you use the citation search feature, references will be automatically managed. For example, you might cite a paper like this: \cite{example2024}. All citations will be automatically added to your \texttt{cite.bib} file.

\section{Conclusion}
This template demonstrates basic \LaTeX{} features including mathematical formulas, tables, and lists. The bibliography system is now set up to work seamlessly with Bibby AI's citation management. Feel free to modify it for your needs!

\bibliographystyle{plain}
\bibliography{cite}

\end{document}\documentclass{article}
\usepackage[utf8]{inputenc}
\usepackage{booktabs}
\usepackage{amsmath}
\usepackage{natbib}

\title{My Document}
\author{Your Name}
\date{\today}

\begin{document}

\maketitle

\section{Introduction}
This is a sample document demonstrating various \LaTeX{} features including mathematical equations, tables, and formatting. You can cite papers using the citation search feature, and they will be automatically added to your \texttt{cite.bib} file.

\section{Mathematical Expressions}
Here are some example equations:

\subsection{Quadratic Formula}
The quadratic formula is given by:
\[x = \frac{-b \pm \sqrt{b^2 - 4ac}}{2a}\]

\subsection{Matrix Example}
A sample matrix:
\[
\begin{pmatrix}
a_{11} & a_{12} & a_{13} \\
a_{21} & a_{22} & a_{23} \\
a_{31} & a_{32} & a_{33}
\end{pmatrix}
\]

\section{Tables}
Below is a sample table showing data:

\begin{table}[h]
\centering
\begin{tabular}{@{}lrrr@{}}
\toprule
\textbf{Category} & \textbf{Value 1} & \textbf{Value 2} & \textbf{Total} \\
\midrule
Item A & 10.5 & 20.3 & 30.8 \\
Item B & 15.2 & 25.1 & 40.3 \\
Item C & 22.1 & 12.7 & 34.8 \\
\bottomrule
\end{tabular}
\caption{Sample Data Table}
\label{tab:sample}
\end{table}

\section{Lists}
\subsection{Bullet Points}
\begin{itemize}
  \item First item
  \item Second item
  \item Third item with sub-items:
    \begin{itemize}
      \item Sub-item 1
      \item Sub-item 2
    \end{itemize}
\end{itemize}

\subsection{Numbered List}
\begin{enumerate}
  \item First step
  \item Second step
  \item Third step
\end{enumerate}

\section{Citations and References}
When you use the citation search feature, references will be automatically managed. For example, you might cite a paper like this: \cite{example2024}. All citations will be automatically added to your \texttt{cite.bib} file.

\section{Conclusion}
This template demonstrates basic \LaTeX{} features including mathematical formulas, tables, and lists. The bibliography system is now set up to work seamlessly with Bibby AI's citation management. Feel free to modify it for your needs!

\bibliographystyle{plain}
\bibliography{cite}

\end{document}\documentclass{article}
\usepackage[utf8]{inputenc}
\usepackage{booktabs}
\usepackage{amsmath}
\usepackage{natbib}

\title{My Document}
\author{Your Name}
\date{\today}

\begin{document}

\maketitle

\section{Introduction}
This is a sample document demonstrating various \LaTeX{} features including mathematical equations, tables, and formatting. You can cite papers using the citation search feature, and they will be automatically added to your \texttt{cite.bib} file.

\section{Mathematical Expressions}
Here are some example equations:

\subsection{Quadratic Formula}
The quadratic formula is given by:
\[x = \frac{-b \pm \sqrt{b^2 - 4ac}}{2a}\]

\subsection{Matrix Example}
A sample matrix:
\[
\begin{pmatrix}
a_{11} & a_{12} & a_{13} \\
a_{21} & a_{22} & a_{23} \\
a_{31} & a_{32} & a_{33}
\end{pmatrix}
\]

\section{Tables}
Below is a sample table showing data:

\begin{table}[h]
\centering
\begin{tabular}{@{}lrrr@{}}
\toprule
\textbf{Category} & \textbf{Value 1} & \textbf{Value 2} & \textbf{Total} \\
\midrule
Item A & 10.5 & 20.3 & 30.8 \\
Item B & 15.2 & 25.1 & 40.3 \\
Item C & 22.1 & 12.7 & 34.8 \\
\bottomrule
\end{tabular}
\caption{Sample Data Table}
\label{tab:sample}
\end{table}

\section{Lists}
\subsection{Bullet Points}
\begin{itemize}
  \item First item
  \item Second item
  \item Third item with sub-items:
    \begin{itemize}
      \item Sub-item 1
      \item Sub-item 2
    \end{itemize}
\end{itemize}

\subsection{Numbered List}
\begin{enumerate}
  \item First step
  \item Second step
  \item Third step
\end{enumerate}

\section{Citations and References}
When you use the citation search feature, references will be automatically managed. For example, you might cite a paper like this: \cite{example2024}. All citations will be automatically added to your \texttt{cite.bib} file.

\section{Conclusion}
This template demonstrates basic \LaTeX{} features including mathematical formulas, tables, and lists. The bibliography system is now set up to work seamlessly with Bibby AI's citation management. Feel free to modify it for your needs!

\bibliographystyle{plain}
\bibliography{cite}

\end{document}\documentclass{article}
\usepackage[utf8]{inputenc}
\usepackage{booktabs}
\usepackage{amsmath}
\usepackage{natbib}

\title{My Document}
\author{Your Name}
\date{\today}

\begin{document}

\maketitle

\section{Introduction}
This is a sample document demonstrating various \LaTeX{} features including mathematical equations, tables, and formatting. You can cite papers using the citation search feature, and they will be automatically added to your \texttt{cite.bib} file.

\section{Mathematical Expressions}
Here are some example equations:

\subsection{Quadratic Formula}
The quadratic formula is given by:
\[x = \frac{-b \pm \sqrt{b^2 - 4ac}}{2a}\]

\subsection{Matrix Example}
A sample matrix:
\[
\begin{pmatrix}
a_{11} & a_{12} & a_{13} \\
a_{21} & a_{22} & a_{23} \\
a_{31} & a_{32} & a_{33}
\end{pmatrix}
\]

\section{Tables}
Below is a sample table showing data:

\begin{table}[h]
\centering
\begin{tabular}{@{}lrrr@{}}
\toprule
\textbf{Category} & \textbf{Value 1} & \textbf{Value 2} & \textbf{Total} \\
\midrule
Item A & 10.5 & 20.3 & 30.8 \\
Item B & 15.2 & 25.1 & 40.3 \\
Item C & 22.1 & 12.7 & 34.8 \\
\bottomrule
\end{tabular}
\caption{Sample Data Table}
\label{tab:sample}
\end{table}

\section{Lists}
\subsection{Bullet Points}
\begin{itemize}
  \item First item
  \item Second item
  \item Third item with sub-items:
    \begin{itemize}
      \item Sub-item 1
      \item Sub-item 2
    \end{itemize}
\end{itemize}

\subsection{Numbered List}
\begin{enumerate}
  \item First step
  \item Second step
  \item Third step
\end{enumerate}

\section{Citations and References}
When you use the citation search feature, references will be automatically managed. For example, you might cite a paper like this: \cite{example2024}. All citations will be automatically added to your \texttt{cite.bib} file.

\section{Conclusion}
This template demonstrates basic \LaTeX{} features including mathematical formulas, tables, and lists. The bibliography system is now set up to work seamlessly with Bibby AI's citation management. Feel free to modify it for your needs!

\bibliographystyle{plain}
\bibliography{cite}

\end{document}\documentclass{article}
\usepackage[utf8]{inputenc}
\usepackage{booktabs}
\usepackage{amsmath}
\usepackage{natbib}

\title{My Document}
\author{Your Name}
\date{\today}

\begin{document}

\maketitle

\section{Introduction}
This is a sample document demonstrating various \LaTeX{} features including mathematical equations, tables, and formatting. You can cite papers using the citation search feature, and they will be automatically added to your \texttt{cite.bib} file.

\section{Mathematical Expressions}
Here are some example equations:

\subsection{Quadratic Formula}
The quadratic formula is given by:
\[x = \frac{-b \pm \sqrt{b^2 - 4ac}}{2a}\]

\subsection{Matrix Example}
A sample matrix:
\[
\begin{pmatrix}
a_{11} & a_{12} & a_{13} \\
a_{21} & a_{22} & a_{23} \\
a_{31} & a_{32} & a_{33}
\end{pmatrix}
\]

\section{Tables}
Below is a sample table showing data:

\begin{table}[h]
\centering
\begin{tabular}{@{}lrrr@{}}
\toprule
\textbf{Category} & \textbf{Value 1} & \textbf{Value 2} & \textbf{Total} \\
\midrule
Item A & 10.5 & 20.3 & 30.8 \\
Item B & 15.2 & 25.1 & 40.3 \\
Item C & 22.1 & 12.7 & 34.8 \\
\bottomrule
\end{tabular}
\caption{Sample Data Table}
\label{tab:sample}
\end{table}

\section{Lists}
\subsection{Bullet Points}
\begin{itemize}
  \item First item
  \item Second item
  \item Third item with sub-items:
    \begin{itemize}
      \item Sub-item 1
      \item Sub-item 2
    \end{itemize}
\end{itemize}

\subsection{Numbered List}
\begin{enumerate}
  \item First step
  \item Second step
  \item Third step
\end{enumerate}

\section{Citations and References}
When you use the citation search feature, references will be automatically managed. For example, you might cite a paper like this: \cite{example2024}. All citations will be automatically added to your \texttt{cite.bib} file.

\section{Conclusion}
This template demonstrates basic \LaTeX{} features including mathematical formulas, tables, and lists. The bibliography system is now set up to work seamlessly with Bibby AI's citation management. Feel free to modify it for your needs!

\bibliographystyle{plain}
\bibliography{cite}

\end{document}\documentclass{article}
\usepackage[utf8]{inputenc}
\usepackage{booktabs}
\usepackage{amsmath}
\usepackage{natbib}

\title{My Document}
\author{Your Name}
\date{\today}

\begin{document}

\maketitle

\section{Introduction}
This is a sample document demonstrating various \LaTeX{} features including mathematical equations, tables, and formatting. You can cite papers using the citation search feature, and they will be automatically added to your \texttt{cite.bib} file.

\section{Mathematical Expressions}
Here are some example equations:

\subsection{Quadratic Formula}
The quadratic formula is given by:
\[x = \frac{-b \pm \sqrt{b^2 - 4ac}}{2a}\]

\subsection{Matrix Example}
A sample matrix:
\[
\begin{pmatrix}
a_{11} & a_{12} & a_{13} \\
a_{21} & a_{22} & a_{23} \\
a_{31} & a_{32} & a_{33}
\end{pmatrix}
\]

\section{Tables}
Below is a sample table showing data:

\begin{table}[h]
\centering
\begin{tabular}{@{}lrrr@{}}
\toprule
\textbf{Category} & \textbf{Value 1} & \textbf{Value 2} & \textbf{Total} \\
\midrule
Item A & 10.5 & 20.3 & 30.8 \\
Item B & 15.2 & 25.1 & 40.3 \\
Item C & 22.1 & 12.7 & 34.8 \\
\bottomrule
\end{tabular}
\caption{Sample Data Table}
\label{tab:sample}
\end{table}

\section{Lists}
\subsection{Bullet Points}
\begin{itemize}
  \item First item
  \item Second item
  \item Third item with sub-items:
    \begin{itemize}
      \item Sub-item 1
      \item Sub-item 2
    \end{itemize}
\end{itemize}

\subsection{Numbered List}
\begin{enumerate}
  \item First step
  \item Second step
  \item Third step
\end{enumerate}

\section{Citations and References}
When you use the citation search feature, references will be automatically managed. For example, you might cite a paper like this: \cite{example2024}. All citations will be automatically added to your \texttt{cite.bib} file.

\section{Conclusion}
This template demonstrates basic \LaTeX{} features including mathematical formulas, tables, and lists. The bibliography system is now set up to work seamlessly with Bibby AI's citation management. Feel free to modify it for your needs!

\bibliographystyle{plain}
\bibliography{cite}

\end{document}\documentclass{article}
\usepackage[utf8]{inputenc}
\usepackage{booktabs}
\usepackage{amsmath}
\usepackage{natbib}

\title{My Document}
\author{Your Name}
\date{\today}

\begin{document}

\maketitle

\section{Introduction}
This is a sample document demonstrating various \LaTeX{} features including mathematical equations, tables, and formatting. You can cite papers using the citation search feature, and they will be automatically added to your \texttt{cite.bib} file.

\section{Mathematical Expressions}
Here are some example equations:

\subsection{Quadratic Formula}
The quadratic formula is given by:
\[x = \frac{-b \pm \sqrt{b^2 - 4ac}}{2a}\]

\subsection{Matrix Example}
A sample matrix:
\[
\begin{pmatrix}
a_{11} & a_{12} & a_{13} \\
a_{21} & a_{22} & a_{23} \\
a_{31} & a_{32} & a_{33}
\end{pmatrix}
\]

\section{Tables}
Below is a sample table showing data:

\begin{table}[h]
\centering
\begin{tabular}{@{}lrrr@{}}
\toprule
\textbf{Category} & \textbf{Value 1} & \textbf{Value 2} & \textbf{Total} \\
\midrule
Item A & 10.5 & 20.3 & 30.8 \\
Item B & 15.2 & 25.1 & 40.3 \\
Item C & 22.1 & 12.7 & 34.8 \\
\bottomrule
\end{tabular}
\caption{Sample Data Table}
\label{tab:sample}
\end{table}

\section{Lists}
\subsection{Bullet Points}
\begin{itemize}
  \item First item
  \item Second item
  \item Third item with sub-items:
    \begin{itemize}
      \item Sub-item 1
      \item Sub-item 2
    \end{itemize}
\end{itemize}

\subsection{Numbered List}
\begin{enumerate}
  \item First step
  \item Second step
  \item Third step
\end{enumerate}

\section{Citations and References}
When you use the citation search feature, references will be automatically managed. For example, you might cite a paper like this: \cite{example2024}. All citations will be automatically added to your \texttt{cite.bib} file.

\section{Conclusion}
This template demonstrates basic \LaTeX{} features including mathematical formulas, tables, and lists. The bibliography system is now set up to work seamlessly with Bibby AI's citation management. Feel free to modify it for your needs!

\bibliographystyle{plain}
\bibliography{cite}

\end{document}\documentclass{article}
\usepackage[utf8]{inputenc}
\usepackage{booktabs}
\usepackage{amsmath}
\usepackage{natbib}

\title{My Document}
\author{Your Name}
\date{\today}

\begin{document}

\maketitle

\section{Introduction}
This is a sample document demonstrating various \LaTeX{} features including mathematical equations, tables, and formatting. You can cite papers using the citation search feature, and they will be automatically added to your \texttt{cite.bib} file.

\section{Mathematical Expressions}
Here are some example equations:

\subsection{Quadratic Formula}
The quadratic formula is given by:
\[x = \frac{-b \pm \sqrt{b^2 - 4ac}}{2a}\]

\subsection{Matrix Example}
A sample matrix:
\[
\begin{pmatrix}
a_{11} & a_{12} & a_{13} \\
a_{21} & a_{22} & a_{23} \\
a_{31} & a_{32} & a_{33}
\end{pmatrix}
\]

\section{Tables}
Below is a sample table showing data:

\begin{table}[h]
\centering
\begin{tabular}{@{}lrrr@{}}
\toprule
\textbf{Category} & \textbf{Value 1} & \textbf{Value 2} & \textbf{Total} \\
\midrule
Item A & 10.5 & 20.3 & 30.8 \\
Item B & 15.2 & 25.1 & 40.3 \\
Item C & 22.1 & 12.7 & 34.8 \\
\bottomrule
\end{tabular}
\caption{Sample Data Table}
\label{tab:sample}
\end{table}

\section{Lists}
\subsection{Bullet Points}
\begin{itemize}
  \item First item
  \item Second item
  \item Third item with sub-items:
    \begin{itemize}
      \item Sub-item 1
      \item Sub-item 2
    \end{itemize}
\end{itemize}

\subsection{Numbered List}
\begin{enumerate}
  \item First step
  \item Second step
  \item Third step
\end{enumerate}

\section{Citations and References}
When you use the citation search feature, references will be automatically managed. For example, you might cite a paper like this: \cite{example2024}. All citations will be automatically added to your \texttt{cite.bib} file.

\section{Conclusion}
This template demonstrates basic \LaTeX{} features including mathematical formulas, tables, and lists. The bibliography system is now set up to work seamlessly with Bibby AI's citation management. Feel free to modify it for your needs!

\bibliographystyle{plain}
\bibliography{cite}

\end{document}\documentclass{article}
\usepackage[utf8]{inputenc}
\usepackage{booktabs}
\usepackage{amsmath}
\usepackage{natbib}

\title{My Document}
\author{Your Name}
\date{\today}

\begin{document}

\maketitle

\section{Introduction}
This is a sample document demonstrating various \LaTeX{} features including mathematical equations, tables, and formatting. You can cite papers using the citation search feature, and they will be automatically added to your \texttt{cite.bib} file.

\section{Mathematical Expressions}
Here are some example equations:

\subsection{Quadratic Formula}
The quadratic formula is given by:
\[x = \frac{-b \pm \sqrt{b^2 - 4ac}}{2a}\]

\subsection{Matrix Example}
A sample matrix:
\[
\begin{pmatrix}
a_{11} & a_{12} & a_{13} \\
a_{21} & a_{22} & a_{23} \\
a_{31} & a_{32} & a_{33}
\end{pmatrix}
\]

\section{Tables}
Below is a sample table showing data:

\begin{table}[h]
\centering
\begin{tabular}{@{}lrrr@{}}
\toprule
\textbf{Category} & \textbf{Value 1} & \textbf{Value 2} & \textbf{Total} \\
\midrule
Item A & 10.5 & 20.3 & 30.8 \\
Item B & 15.2 & 25.1 & 40.3 \\
Item C & 22.1 & 12.7 & 34.8 \\
\bottomrule
\end{tabular}
\caption{Sample Data Table}
\label{tab:sample}
\end{table}

\section{Lists}
\subsection{Bullet Points}
\begin{itemize}
  \item First item
  \item Second item
  \item Third item with sub-items:
    \begin{itemize}
      \item Sub-item 1
      \item Sub-item 2
    \end{itemize}
\end{itemize}

\subsection{Numbered List}
\begin{enumerate}
  \item First step
  \item Second step
  \item Third step
\end{enumerate}

\section{Citations and References}
When you use the citation search feature, references will be automatically managed. For example, you might cite a paper like this: \cite{example2024}. All citations will be automatically added to your \texttt{cite.bib} file.

\section{Conclusion}
This template demonstrates basic \LaTeX{} features including mathematical formulas, tables, and lists. The bibliography system is now set up to work seamlessly with Bibby AI's citation management. Feel free to modify it for your needs!

\bibliographystyle{plain}
\bibliography{cite}

\end{document}\documentclass{article}
\usepackage[utf8]{inputenc}
\usepackage{booktabs}
\usepackage{amsmath}
\usepackage{natbib}

\title{My Document}
\author{Your Name}
\date{\today}

\begin{document}

\maketitle

\section{Introduction}
This is a sample document demonstrating various \LaTeX{} features including mathematical equations, tables, and formatting. You can cite papers using the citation search feature, and they will be automatically added to your \texttt{cite.bib} file.

\section{Mathematical Expressions}
Here are some example equations:

\subsection{Quadratic Formula}
The quadratic formula is given by:
\[x = \frac{-b \pm \sqrt{b^2 - 4ac}}{2a}\]

\subsection{Matrix Example}
A sample matrix:
\[
\begin{pmatrix}
a_{11} & a_{12} & a_{13} \\
a_{21} & a_{22} & a_{23} \\
a_{31} & a_{32} & a_{33}
\end{pmatrix}
\]

\section{Tables}
Below is a sample table showing data:

\begin{table}[h]
\centering
\begin{tabular}{@{}lrrr@{}}
\toprule
\textbf{Category} & \textbf{Value 1} & \textbf{Value 2} & \textbf{Total} \\
\midrule
Item A & 10.5 & 20.3 & 30.8 \\
Item B & 15.2 & 25.1 & 40.3 \\
Item C & 22.1 & 12.7 & 34.8 \\
\bottomrule
\end{tabular}
\caption{Sample Data Table}
\label{tab:sample}
\end{table}

\section{Lists}
\subsection{Bullet Points}
\begin{itemize}
  \item First item
  \item Second item
  \item Third item with sub-items:
    \begin{itemize}
      \item Sub-item 1
      \item Sub-item 2
    \end{itemize}
\end{itemize}

\subsection{Numbered List}
\begin{enumerate}
  \item First step
  \item Second step
  \item Third step
\end{enumerate}

\section{Citations and References}
When you use the citation search feature, references will be automatically managed. For example, you might cite a paper like this: \cite{example2024}. All citations will be automatically added to your \texttt{cite.bib} file.

\section{Conclusion}
This template demonstrates basic \LaTeX{} features including mathematical formulas, tables, and lists. The bibliography system is now set up to work seamlessly with Bibby AI's citation management. Feel free to modify it for your needs!

\bibliographystyle{plain}
\bibliography{cite}

\end{document}\documentclass{article}
\usepackage[utf8]{inputenc}
\usepackage{booktabs}
\usepackage{amsmath}
\usepackage{natbib}

\title{My Document}
\author{Your Name}
\date{\today}

\begin{document}

\maketitle

\section{Introduction}
This is a sample document demonstrating various \LaTeX{} features including mathematical equations, tables, and formatting. You can cite papers using the citation search feature, and they will be automatically added to your \texttt{cite.bib} file.

\section{Mathematical Expressions}
Here are some example equations:

\subsection{Quadratic Formula}
The quadratic formula is given by:
\[x = \frac{-b \pm \sqrt{b^2 - 4ac}}{2a}\]

\subsection{Matrix Example}
A sample matrix:
\[
\begin{pmatrix}
a_{11} & a_{12} & a_{13} \\
a_{21} & a_{22} & a_{23} \\
a_{31} & a_{32} & a_{33}
\end{pmatrix}
\]

\section{Tables}
Below is a sample table showing data:

\begin{table}[h]
\centering
\begin{tabular}{@{}lrrr@{}}
\toprule
\textbf{Category} & \textbf{Value 1} & \textbf{Value 2} & \textbf{Total} \\
\midrule
Item A & 10.5 & 20.3 & 30.8 \\
Item B & 15.2 & 25.1 & 40.3 \\
Item C & 22.1 & 12.7 & 34.8 \\
\bottomrule
\end{tabular}
\caption{Sample Data Table}
\label{tab:sample}
\end{table}

\section{Lists}
\subsection{Bullet Points}
\begin{itemize}
  \item First item
  \item Second item
  \item Third item with sub-items:
    \begin{itemize}
      \item Sub-item 1
      \item Sub-item 2
    \end{itemize}
\end{itemize}

\subsection{Numbered List}
\begin{enumerate}
  \item First step
  \item Second step
  \item Third step
\end{enumerate}

\section{Citations and References}
When you use the citation search feature, references will be automatically managed. For example, you might cite a paper like this: \cite{example2024}. All citations will be automatically added to your \texttt{cite.bib} file.

\section{Conclusion}
This template demonstrates basic \LaTeX{} features including mathematical formulas, tables, and lists. The bibliography system is now set up to work seamlessly with Bibby AI's citation management. Feel free to modify it for your needs!

\bibliographystyle{plain}
\bibliography{cite}

\end{document}\documentclass{article}
\usepackage[utf8]{inputenc}
\usepackage{booktabs}
\usepackage{amsmath}
\usepackage{natbib}

\title{My Document}
\author{Your Name}
\date{\today}

\begin{document}

\maketitle

\section{Introduction}
This is a sample document demonstrating various \LaTeX{} features including mathematical equations, tables, and formatting. You can cite papers using the citation search feature, and they will be automatically added to your \texttt{cite.bib} file.

\section{Mathematical Expressions}
Here are some example equations:

\subsection{Quadratic Formula}
The quadratic formula is given by:
\[x = \frac{-b \pm \sqrt{b^2 - 4ac}}{2a}\]

\subsection{Matrix Example}
A sample matrix:
\[
\begin{pmatrix}
a_{11} & a_{12} & a_{13} \\
a_{21} & a_{22} & a_{23} \\
a_{31} & a_{32} & a_{33}
\end{pmatrix}
\]

\section{Tables}
Below is a sample table showing data:

\begin{table}[h]
\centering
\begin{tabular}{@{}lrrr@{}}
\toprule
\textbf{Category} & \textbf{Value 1} & \textbf{Value 2} & \textbf{Total} \\
\midrule
Item A & 10.5 & 20.3 & 30.8 \\
Item B & 15.2 & 25.1 & 40.3 \\
Item C & 22.1 & 12.7 & 34.8 \\
\bottomrule
\end{tabular}
\caption{Sample Data Table}
\label{tab:sample}
\end{table}

\section{Lists}
\subsection{Bullet Points}
\begin{itemize}
  \item First item
  \item Second item
  \item Third item with sub-items:
    \begin{itemize}
      \item Sub-item 1
      \item Sub-item 2
    \end{itemize}
\end{itemize}

\subsection{Numbered List}
\begin{enumerate}
  \item First step
  \item Second step
  \item Third step
\end{enumerate}

\section{Citations and References}
When you use the citation search feature, references will be automatically managed. For example, you might cite a paper like this: \cite{example2024}. All citations will be automatically added to your \texttt{cite.bib} file.

\section{Conclusion}
This template demonstrates basic \LaTeX{} features including mathematical formulas, tables, and lists. The bibliography system is now set up to work seamlessly with Bibby AI's citation management. Feel free to modify it for your needs!

\bibliographystyle{plain}
\bibliography{cite}

\end{document}\documentclass{article}
\usepackage[utf8]{inputenc}
\usepackage{booktabs}
\usepackage{amsmath}
\usepackage{natbib}

\title{My Document}
\author{Your Name}
\date{\today}

\begin{document}

\maketitle

\section{Introduction}
This is a sample document demonstrating various \LaTeX{} features including mathematical equations, tables, and formatting. You can cite papers using the citation search feature, and they will be automatically added to your \texttt{cite.bib} file.

\section{Mathematical Expressions}
Here are some example equations:

\subsection{Quadratic Formula}
The quadratic formula is given by:
\[x = \frac{-b \pm \sqrt{b^2 - 4ac}}{2a}\]

\subsection{Matrix Example}
A sample matrix:
\[
\begin{pmatrix}
a_{11} & a_{12} & a_{13} \\
a_{21} & a_{22} & a_{23} \\
a_{31} & a_{32} & a_{33}
\end{pmatrix}
\]

\section{Tables}
Below is a sample table showing data:

\begin{table}[h]
\centering
\begin{tabular}{@{}lrrr@{}}
\toprule
\textbf{Category} & \textbf{Value 1} & \textbf{Value 2} & \textbf{Total} \\
\midrule
Item A & 10.5 & 20.3 & 30.8 \\
Item B & 15.2 & 25.1 & 40.3 \\
Item C & 22.1 & 12.7 & 34.8 \\
\bottomrule
\end{tabular}
\caption{Sample Data Table}
\label{tab:sample}
\end{table}

\section{Lists}
\subsection{Bullet Points}
\begin{itemize}
  \item First item
  \item Second item
  \item Third item with sub-items:
    \begin{itemize}
      \item Sub-item 1
      \item Sub-item 2
    \end{itemize}
\end{itemize}

\subsection{Numbered List}
\begin{enumerate}
  \item First step
  \item Second step
  \item Third step
\end{enumerate}

\section{Citations and References}
When you use the citation search feature, references will be automatically managed. For example, you might cite a paper like this: \cite{example2024}. All citations will be automatically added to your \texttt{cite.bib} file.

\section{Conclusion}
This template demonstrates basic \LaTeX{} features including mathematical formulas, tables, and lists. The bibliography system is now set up to work seamlessly with Bibby AI's citation management. Feel free to modify it for your needs!

\bibliographystyle{plain}
\bibliography{cite}

\end{document}\documentclass{article}
\usepackage[utf8]{inputenc}
\usepackage{booktabs}
\usepackage{amsmath}
\usepackage{natbib}

\title{My Document}
\author{Your Name}
\date{\today}

\begin{document}

\maketitle

\section{Introduction}
This is a sample document demonstrating various \LaTeX{} features including mathematical equations, tables, and formatting. You can cite papers using the citation search feature, and they will be automatically added to your \texttt{cite.bib} file.

\section{Mathematical Expressions}
Here are some example equations:

\subsection{Quadratic Formula}
The quadratic formula is given by:
\[x = \frac{-b \pm \sqrt{b^2 - 4ac}}{2a}\]

\subsection{Matrix Example}
A sample matrix:
\[
\begin{pmatrix}
a_{11} & a_{12} & a_{13} \\
a_{21} & a_{22} & a_{23} \\
a_{31} & a_{32} & a_{33}
\end{pmatrix}
\]

\section{Tables}
Below is a sample table showing data:

\begin{table}[h]
\centering
\begin{tabular}{@{}lrrr@{}}
\toprule
\textbf{Category} & \textbf{Value 1} & \textbf{Value 2} & \textbf{Total} \\
\midrule
Item A & 10.5 & 20.3 & 30.8 \\
Item B & 15.2 & 25.1 & 40.3 \\
Item C & 22.1 & 12.7 & 34.8 \\
\bottomrule
\end{tabular}
\caption{Sample Data Table}
\label{tab:sample}
\end{table}

\section{Lists}
\subsection{Bullet Points}
\begin{itemize}
  \item First item
  \item Second item
  \item Third item with sub-items:
    \begin{itemize}
      \item Sub-item 1
      \item Sub-item 2
    \end{itemize}
\end{itemize}

\subsection{Numbered List}
\begin{enumerate}
  \item First step
  \item Second step
  \item Third step
\end{enumerate}

\section{Citations and References}
When you use the citation search feature, references will be automatically managed. For example, you might cite a paper like this: \cite{example2024}. All citations will be automatically added to your \texttt{cite.bib} file.

\section{Conclusion}
This template demonstrates basic \LaTeX{} features including mathematical formulas, tables, and lists. The bibliography system is now set up to work seamlessly with Bibby AI's citation management. Feel free to modify it for your needs!

\bibliographystyle{plain}
\bibliography{cite}

\end{document}\documentclass{article}
\usepackage[utf8]{inputenc}
\usepackage{booktabs}
\usepackage{amsmath}
\usepackage{natbib}

\title{My Document}
\author{Your Name}
\date{\today}

\begin{document}

\maketitle

\section{Introduction}
This is a sample document demonstrating various \LaTeX{} features including mathematical equations, tables, and formatting. You can cite papers using the citation search feature, and they will be automatically added to your \texttt{cite.bib} file.

\section{Mathematical Expressions}
Here are some example equations:

\subsection{Quadratic Formula}
The quadratic formula is given by:
\[x = \frac{-b \pm \sqrt{b^2 - 4ac}}{2a}\]

\subsection{Matrix Example}
A sample matrix:
\[
\begin{pmatrix}
a_{11} & a_{12} & a_{13} \\
a_{21} & a_{22} & a_{23} \\
a_{31} & a_{32} & a_{33}
\end{pmatrix}
\]

\section{Tables}
Below is a sample table showing data:

\begin{table}[h]
\centering
\begin{tabular}{@{}lrrr@{}}
\toprule
\textbf{Category} & \textbf{Value 1} & \textbf{Value 2} & \textbf{Total} \\
\midrule
Item A & 10.5 & 20.3 & 30.8 \\
Item B & 15.2 & 25.1 & 40.3 \\
Item C & 22.1 & 12.7 & 34.8 \\
\bottomrule
\end{tabular}
\caption{Sample Data Table}
\label{tab:sample}
\end{table}

\section{Lists}
\subsection{Bullet Points}
\begin{itemize}
  \item First item
  \item Second item
  \item Third item with sub-items:
    \begin{itemize}
      \item Sub-item 1
      \item Sub-item 2
    \end{itemize}
\end{itemize}

\subsection{Numbered List}
\begin{enumerate}
  \item First step
  \item Second step
  \item Third step
\end{enumerate}

\section{Citations and References}
When you use the citation search feature, references will be automatically managed. For example, you might cite a paper like this: \cite{example2024}. All citations will be automatically added to your \texttt{cite.bib} file.

\section{Conclusion}
This template demonstrates basic \LaTeX{} features including mathematical formulas, tables, and lists. The bibliography system is now set up to work seamlessly with Bibby AI's citation management. Feel free to modify it for your needs!

\bibliographystyle{plain}
\bibliography{cite}

\end{document}\documentclass{article}
\usepackage[utf8]{inputenc}
\usepackage{booktabs}
\usepackage{amsmath}
\usepackage{natbib}

\title{My Document}
\author{Your Name}
\date{\today}

\begin{document}

\maketitle

\section{Introduction}
This is a sample document demonstrating various \LaTeX{} features including mathematical equations, tables, and formatting. You can cite papers using the citation search feature, and they will be automatically added to your \texttt{cite.bib} file.

\section{Mathematical Expressions}
Here are some example equations:

\subsection{Quadratic Formula}
The quadratic formula is given by:
\[x = \frac{-b \pm \sqrt{b^2 - 4ac}}{2a}\]

\subsection{Matrix Example}
A sample matrix:
\[
\begin{pmatrix}
a_{11} & a_{12} & a_{13} \\
a_{21} & a_{22} & a_{23} \\
a_{31} & a_{32} & a_{33}
\end{pmatrix}
\]

\section{Tables}
Below is a sample table showing data:

\begin{table}[h]
\centering
\begin{tabular}{@{}lrrr@{}}
\toprule
\textbf{Category} & \textbf{Value 1} & \textbf{Value 2} & \textbf{Total} \\
\midrule
Item A & 10.5 & 20.3 & 30.8 \\
Item B & 15.2 & 25.1 & 40.3 \\
Item C & 22.1 & 12.7 & 34.8 \\
\bottomrule
\end{tabular}
\caption{Sample Data Table}
\label{tab:sample}
\end{table}

\section{Lists}
\subsection{Bullet Points}
\begin{itemize}
  \item First item
  \item Second item
  \item Third item with sub-items:
    \begin{itemize}
      \item Sub-item 1
      \item Sub-item 2
    \end{itemize}
\end{itemize}

\subsection{Numbered List}
\begin{enumerate}
  \item First step
  \item Second step
  \item Third step
\end{enumerate}

\section{Citations and References}
When you use the citation search feature, references will be automatically managed. For example, you might cite a paper like this: \cite{example2024}. All citations will be automatically added to your \texttt{cite.bib} file.

\section{Conclusion}
This template demonstrates basic \LaTeX{} features including mathematical formulas, tables, and lists. The bibliography system is now set up to work seamlessly with Bibby AI's citation management. Feel free to modify it for your needs!

\bibliographystyle{plain}
\bibliography{cite}

\end{document}\documentclass{article}
\usepackage[utf8]{inputenc}
\usepackage{booktabs}
\usepackage{amsmath}
\usepackage{natbib}

\title{My Document}
\author{Your Name}
\date{\today}

\begin{document}

\maketitle

\section{Introduction}
This is a sample document demonstrating various \LaTeX{} features including mathematical equations, tables, and formatting. You can cite papers using the citation search feature, and they will be automatically added to your \texttt{cite.bib} file.

\section{Mathematical Expressions}
Here are some example equations:

\subsection{Quadratic Formula}
The quadratic formula is given by:
\[x = \frac{-b \pm \sqrt{b^2 - 4ac}}{2a}\]

\subsection{Matrix Example}
A sample matrix:
\[
\begin{pmatrix}
a_{11} & a_{12} & a_{13} \\
a_{21} & a_{22} & a_{23} \\
a_{31} & a_{32} & a_{33}
\end{pmatrix}
\]

\section{Tables}
Below is a sample table showing data:

\begin{table}[h]
\centering
\begin{tabular}{@{}lrrr@{}}
\toprule
\textbf{Category} & \textbf{Value 1} & \textbf{Value 2} & \textbf{Total} \\
\midrule
Item A & 10.5 & 20.3 & 30.8 \\
Item B & 15.2 & 25.1 & 40.3 \\
Item C & 22.1 & 12.7 & 34.8 \\
\bottomrule
\end{tabular}
\caption{Sample Data Table}
\label{tab:sample}
\end{table}

\section{Lists}
\subsection{Bullet Points}
\begin{itemize}
  \item First item
  \item Second item
  \item Third item with sub-items:
    \begin{itemize}
      \item Sub-item 1
      \item Sub-item 2
    \end{itemize}
\end{itemize}

\subsection{Numbered List}
\begin{enumerate}
  \item First step
  \item Second step
  \item Third step
\end{enumerate}

\section{Citations and References}
When you use the citation search feature, references will be automatically managed. For example, you might cite a paper like this: \cite{example2024}. All citations will be automatically added to your \texttt{cite.bib} file.

\section{Conclusion}
This template demonstrates basic \LaTeX{} features including mathematical formulas, tables, and lists. The bibliography system is now set up to work seamlessly with Bibby AI's citation management. Feel free to modify it for your needs!

\bibliographystyle{plain}
\bibliography{cite}

\end{document}\documentclass{article}
\usepackage[utf8]{inputenc}
\usepackage{booktabs}
\usepackage{amsmath}
\usepackage{natbib}

\title{My Document}
\author{Your Name}
\date{\today}

\begin{document}

\maketitle

\section{Introduction}
This is a sample document demonstrating various \LaTeX{} features including mathematical equations, tables, and formatting. You can cite papers using the citation search feature, and they will be automatically added to your \texttt{cite.bib} file.

\section{Mathematical Expressions}
Here are some example equations:

\subsection{Quadratic Formula}
The quadratic formula is given by:
\[x = \frac{-b \pm \sqrt{b^2 - 4ac}}{2a}\]

\subsection{Matrix Example}
A sample matrix:
\[
\begin{pmatrix}
a_{11} & a_{12} & a_{13} \\
a_{21} & a_{22} & a_{23} \\
a_{31} & a_{32} & a_{33}
\end{pmatrix}
\]

\section{Tables}
Below is a sample table showing data:

\begin{table}[h]
\centering
\begin{tabular}{@{}lrrr@{}}
\toprule
\textbf{Category} & \textbf{Value 1} & \textbf{Value 2} & \textbf{Total} \\
\midrule
Item A & 10.5 & 20.3 & 30.8 \\
Item B & 15.2 & 25.1 & 40.3 \\
Item C & 22.1 & 12.7 & 34.8 \\
\bottomrule
\end{tabular}
\caption{Sample Data Table}
\label{tab:sample}
\end{table}

\section{Lists}
\subsection{Bullet Points}
\begin{itemize}
  \item First item
  \item Second item
  \item Third item with sub-items:
    \begin{itemize}
      \item Sub-item 1
      \item Sub-item 2
    \end{itemize}
\end{itemize}

\subsection{Numbered List}
\begin{enumerate}
  \item First step
  \item Second step
  \item Third step
\end{enumerate}

\section{Citations and References}
When you use the citation search feature, references will be automatically managed. For example, you might cite a paper like this: \cite{example2024}. All citations will be automatically added to your \texttt{cite.bib} file.

\section{Conclusion}
This template demonstrates basic \LaTeX{} features including mathematical formulas, tables, and lists. The bibliography system is now set up to work seamlessly with Bibby AI's citation management. Feel free to modify it for your needs!

\bibliographystyle{plain}
\bibliography{cite}

\end{document}\documentclass{article}
\usepackage[utf8]{inputenc}
\usepackage{booktabs}
\usepackage{amsmath}
\usepackage{natbib}

\title{My Document}
\author{Your Name}
\date{\today}

\begin{document}

\maketitle

\section{Introduction}
This is a sample document demonstrating various \LaTeX{} features including mathematical equations, tables, and formatting. You can cite papers using the citation search feature, and they will be automatically added to your \texttt{cite.bib} file.

\section{Mathematical Expressions}
Here are some example equations:

\subsection{Quadratic Formula}
The quadratic formula is given by:
\[x = \frac{-b \pm \sqrt{b^2 - 4ac}}{2a}\]

\subsection{Matrix Example}
A sample matrix:
\[
\begin{pmatrix}
a_{11} & a_{12} & a_{13} \\
a_{21} & a_{22} & a_{23} \\
a_{31} & a_{32} & a_{33}
\end{pmatrix}
\]

\section{Tables}
Below is a sample table showing data:

\begin{table}[h]
\centering
\begin{tabular}{@{}lrrr@{}}
\toprule
\textbf{Category} & \textbf{Value 1} & \textbf{Value 2} & \textbf{Total} \\
\midrule
Item A & 10.5 & 20.3 & 30.8 \\
Item B & 15.2 & 25.1 & 40.3 \\
Item C & 22.1 & 12.7 & 34.8 \\
\bottomrule
\end{tabular}
\caption{Sample Data Table}
\label{tab:sample}
\end{table}

\section{Lists}
\subsection{Bullet Points}
\begin{itemize}
  \item First item
  \item Second item
  \item Third item with sub-items:
    \begin{itemize}
      \item Sub-item 1
      \item Sub-item 2
    \end{itemize}
\end{itemize}

\subsection{Numbered List}
\begin{enumerate}
  \item First step
  \item Second step
  \item Third step
\end{enumerate}

\section{Citations and References}
When you use the citation search feature, references will be automatically managed. For example, you might cite a paper like this: \cite{example2024}. All citations will be automatically added to your \texttt{cite.bib} file.

\section{Conclusion}
This template demonstrates basic \LaTeX{} features including mathematical formulas, tables, and lists. The bibliography system is now set up to work seamlessly with Bibby AI's citation management. Feel free to modify it for your needs!

\bibliographystyle{plain}
\bibliography{cite}

\end{document}\documentclass{article}
\usepackage[utf8]{inputenc}
\usepackage{booktabs}
\usepackage{amsmath}
\usepackage{natbib}

\title{My Document}
\author{Your Name}
\date{\today}

\begin{document}

\maketitle

\section{Introduction}
This is a sample document demonstrating various \LaTeX{} features including mathematical equations, tables, and formatting. You can cite papers using the citation search feature, and they will be automatically added to your \texttt{cite.bib} file.

\section{Mathematical Expressions}
Here are some example equations:

\subsection{Quadratic Formula}
The quadratic formula is given by:
\[x = \frac{-b \pm \sqrt{b^2 - 4ac}}{2a}\]

\subsection{Matrix Example}
A sample matrix:
\[
\begin{pmatrix}
a_{11} & a_{12} & a_{13} \\
a_{21} & a_{22} & a_{23} \\
a_{31} & a_{32} & a_{33}
\end{pmatrix}
\]

\section{Tables}
Below is a sample table showing data:

\begin{table}[h]
\centering
\begin{tabular}{@{}lrrr@{}}
\toprule
\textbf{Category} & \textbf{Value 1} & \textbf{Value 2} & \textbf{Total} \\
\midrule
Item A & 10.5 & 20.3 & 30.8 \\
Item B & 15.2 & 25.1 & 40.3 \\
Item C & 22.1 & 12.7 & 34.8 \\
\bottomrule
\end{tabular}
\caption{Sample Data Table}
\label{tab:sample}
\end{table}

\section{Lists}
\subsection{Bullet Points}
\begin{itemize}
  \item First item
  \item Second item
  \item Third item with sub-items:
    \begin{itemize}
      \item Sub-item 1
      \item Sub-item 2
    \end{itemize}
\end{itemize}

\subsection{Numbered List}
\begin{enumerate}
  \item First step
  \item Second step
  \item Third step
\end{enumerate}

\section{Citations and References}
When you use the citation search feature, references will be automatically managed. For example, you might cite a paper like this: \cite{example2024}. All citations will be automatically added to your \texttt{cite.bib} file.

\section{Conclusion}
This template demonstrates basic \LaTeX{} features including mathematical formulas, tables, and lists. The bibliography system is now set up to work seamlessly with Bibby AI's citation management. Feel free to modify it for your needs!

\bibliographystyle{plain}
\bibliography{cite}

\end{document}\documentclass{article}
\usepackage[utf8]{inputenc}
\usepackage{booktabs}
\usepackage{amsmath}
\usepackage{natbib}

\title{My Document}
\author{Your Name}
\date{\today}

\begin{document}

\maketitle

\section{Introduction}
This is a sample document demonstrating various \LaTeX{} features including mathematical equations, tables, and formatting. You can cite papers using the citation search feature, and they will be automatically added to your \texttt{cite.bib} file.

\section{Mathematical Expressions}
Here are some example equations:

\subsection{Quadratic Formula}
The quadratic formula is given by:
\[x = \frac{-b \pm \sqrt{b^2 - 4ac}}{2a}\]

\subsection{Matrix Example}
A sample matrix:
\[
\begin{pmatrix}
a_{11} & a_{12} & a_{13} \\
a_{21} & a_{22} & a_{23} \\
a_{31} & a_{32} & a_{33}
\end{pmatrix}
\]

\section{Tables}
Below is a sample table showing data:

\begin{table}[h]
\centering
\begin{tabular}{@{}lrrr@{}}
\toprule
\textbf{Category} & \textbf{Value 1} & \textbf{Value 2} & \textbf{Total} \\
\midrule
Item A & 10.5 & 20.3 & 30.8 \\
Item B & 15.2 & 25.1 & 40.3 \\
Item C & 22.1 & 12.7 & 34.8 \\
\bottomrule
\end{tabular}
\caption{Sample Data Table}
\label{tab:sample}
\end{table}

\section{Lists}
\subsection{Bullet Points}
\begin{itemize}
  \item First item
  \item Second item
  \item Third item with sub-items:
    \begin{itemize}
      \item Sub-item 1
      \item Sub-item 2
    \end{itemize}
\end{itemize}

\subsection{Numbered List}
\begin{enumerate}
  \item First step
  \item Second step
  \item Third step
\end{enumerate}

\section{Citations and References}
When you use the citation search feature, references will be automatically managed. For example, you might cite a paper like this: \cite{example2024}. All citations will be automatically added to your \texttt{cite.bib} file.

\section{Conclusion}
This template demonstrates basic \LaTeX{} features including mathematical formulas, tables, and lists. The bibliography system is now set up to work seamlessly with Bibby AI's citation management. Feel free to modify it for your needs!

\bibliographystyle{plain}
\bibliography{cite}

\end{document}\documentclass{article}
\usepackage[utf8]{inputenc}
\usepackage{booktabs}
\usepackage{amsmath}
\usepackage{natbib}

\title{My Document}
\author{Your Name}
\date{\today}

\begin{document}

\maketitle

\section{Introduction}
This is a sample document demonstrating various \LaTeX{} features including mathematical equations, tables, and formatting. You can cite papers using the citation search feature, and they will be automatically added to your \texttt{cite.bib} file.

\section{Mathematical Expressions}
Here are some example equations:

\subsection{Quadratic Formula}
The quadratic formula is given by:
\[x = \frac{-b \pm \sqrt{b^2 - 4ac}}{2a}\]

\subsection{Matrix Example}
A sample matrix:
\[
\begin{pmatrix}
a_{11} & a_{12} & a_{13} \\
a_{21} & a_{22} & a_{23} \\
a_{31} & a_{32} & a_{33}
\end{pmatrix}
\]

\section{Tables}
Below is a sample table showing data:

\begin{table}[h]
\centering
\begin{tabular}{@{}lrrr@{}}
\toprule
\textbf{Category} & \textbf{Value 1} & \textbf{Value 2} & \textbf{Total} \\
\midrule
Item A & 10.5 & 20.3 & 30.8 \\
Item B & 15.2 & 25.1 & 40.3 \\
Item C & 22.1 & 12.7 & 34.8 \\
\bottomrule
\end{tabular}
\caption{Sample Data Table}
\label{tab:sample}
\end{table}

\section{Lists}
\subsection{Bullet Points}
\begin{itemize}
  \item First item
  \item Second item
  \item Third item with sub-items:
    \begin{itemize}
      \item Sub-item 1
      \item Sub-item 2
    \end{itemize}
\end{itemize}

\subsection{Numbered List}
\begin{enumerate}
  \item First step
  \item Second step
  \item Third step
\end{enumerate}

\section{Citations and References}
When you use the citation search feature, references will be automatically managed. For example, you might cite a paper like this: \cite{example2024}. All citations will be automatically added to your \texttt{cite.bib} file.

\section{Conclusion}
This template demonstrates basic \LaTeX{} features including mathematical formulas, tables, and lists. The bibliography system is now set up to work seamlessly with Bibby AI's citation management. Feel free to modify it for your needs!

\bibliographystyle{plain}
\bibliography{cite}

\end{document}\documentclass{article}
\usepackage[utf8]{inputenc}
\usepackage{booktabs}
\usepackage{amsmath}
\usepackage{natbib}

\title{My Document}
\author{Your Name}
\date{\today}

\begin{document}

\maketitle

\section{Introduction}
This is a sample document demonstrating various \LaTeX{} features including mathematical equations, tables, and formatting. You can cite papers using the citation search feature, and they will be automatically added to your \texttt{cite.bib} file.

\section{Mathematical Expressions}
Here are some example equations:

\subsection{Quadratic Formula}
The quadratic formula is given by:
\[x = \frac{-b \pm \sqrt{b^2 - 4ac}}{2a}\]

\subsection{Matrix Example}
A sample matrix:
\[
\begin{pmatrix}
a_{11} & a_{12} & a_{13} \\
a_{21} & a_{22} & a_{23} \\
a_{31} & a_{32} & a_{33}
\end{pmatrix}
\]

\section{Tables}
Below is a sample table showing data:

\begin{table}[h]
\centering
\begin{tabular}{@{}lrrr@{}}
\toprule
\textbf{Category} & \textbf{Value 1} & \textbf{Value 2} & \textbf{Total} \\
\midrule
Item A & 10.5 & 20.3 & 30.8 \\
Item B & 15.2 & 25.1 & 40.3 \\
Item C & 22.1 & 12.7 & 34.8 \\
\bottomrule
\end{tabular}
\caption{Sample Data Table}
\label{tab:sample}
\end{table}

\section{Lists}
\subsection{Bullet Points}
\begin{itemize}
  \item First item
  \item Second item
  \item Third item with sub-items:
    \begin{itemize}
      \item Sub-item 1
      \item Sub-item 2
    \end{itemize}
\end{itemize}

\subsection{Numbered List}
\begin{enumerate}
  \item First step
  \item Second step
  \item Third step
\end{enumerate}

\section{Citations and References}
When you use the citation search feature, references will be automatically managed. For example, you might cite a paper like this: \cite{example2024}. All citations will be automatically added to your \texttt{cite.bib} file.

\section{Conclusion}
This template demonstrates basic \LaTeX{} features including mathematical formulas, tables, and lists. The bibliography system is now set up to work seamlessly with Bibby AI's citation management. Feel free to modify it for your needs!

\bibliographystyle{plain}
\bibliography{cite}

\end{document}\documentclass{article}
\usepackage[utf8]{inputenc}
\usepackage{booktabs}
\usepackage{amsmath}
\usepackage{natbib}

\title{My Document}
\author{Your Name}
\date{\today}

\begin{document}

\maketitle

\section{Introduction}
This is a sample document demonstrating various \LaTeX{} features including mathematical equations, tables, and formatting. You can cite papers using the citation search feature, and they will be automatically added to your \texttt{cite.bib} file.

\section{Mathematical Expressions}
Here are some example equations:

\subsection{Quadratic Formula}
The quadratic formula is given by:
\[x = \frac{-b \pm \sqrt{b^2 - 4ac}}{2a}\]

\subsection{Matrix Example}
A sample matrix:
\[
\begin{pmatrix}
a_{11} & a_{12} & a_{13} \\
a_{21} & a_{22} & a_{23} \\
a_{31} & a_{32} & a_{33}
\end{pmatrix}
\]

\section{Tables}
Below is a sample table showing data:

\begin{table}[h]
\centering
\begin{tabular}{@{}lrrr@{}}
\toprule
\textbf{Category} & \textbf{Value 1} & \textbf{Value 2} & \textbf{Total} \\
\midrule
Item A & 10.5 & 20.3 & 30.8 \\
Item B & 15.2 & 25.1 & 40.3 \\
Item C & 22.1 & 12.7 & 34.8 \\
\bottomrule
\end{tabular}
\caption{Sample Data Table}
\label{tab:sample}
\end{table}

\section{Lists}
\subsection{Bullet Points}
\begin{itemize}
  \item First item
  \item Second item
  \item Third item with sub-items:
    \begin{itemize}
      \item Sub-item 1
      \item Sub-item 2
    \end{itemize}
\end{itemize}

\subsection{Numbered List}
\begin{enumerate}
  \item First step
  \item Second step
  \item Third step
\end{enumerate}

\section{Citations and References}
When you use the citation search feature, references will be automatically managed. For example, you might cite a paper like this: \cite{example2024}. All citations will be automatically added to your \texttt{cite.bib} file.

\section{Conclusion}
This template demonstrates basic \LaTeX{} features including mathematical formulas, tables, and lists. The bibliography system is now set up to work seamlessly with Bibby AI's citation management. Feel free to modify it for your needs!

\bibliographystyle{plain}
\bibliography{cite}

\end{document}\documentclass{article}
\usepackage[utf8]{inputenc}
\usepackage{booktabs}
\usepackage{amsmath}
\usepackage{natbib}

\title{My Document}
\author{Your Name}
\date{\today}

\begin{document}

\maketitle

\section{Introduction}
This is a sample document demonstrating various \LaTeX{} features including mathematical equations, tables, and formatting. You can cite papers using the citation search feature, and they will be automatically added to your \texttt{cite.bib} file.

\section{Mathematical Expressions}
Here are some example equations:

\subsection{Quadratic Formula}
The quadratic formula is given by:
\[x = \frac{-b \pm \sqrt{b^2 - 4ac}}{2a}\]

\subsection{Matrix Example}
A sample matrix:
\[
\begin{pmatrix}
a_{11} & a_{12} & a_{13} \\
a_{21} & a_{22} & a_{23} \\
a_{31} & a_{32} & a_{33}
\end{pmatrix}
\]

\section{Tables}
Below is a sample table showing data:

\begin{table}[h]
\centering
\begin{tabular}{@{}lrrr@{}}
\toprule
\textbf{Category} & \textbf{Value 1} & \textbf{Value 2} & \textbf{Total} \\
\midrule
Item A & 10.5 & 20.3 & 30.8 \\
Item B & 15.2 & 25.1 & 40.3 \\
Item C & 22.1 & 12.7 & 34.8 \\
\bottomrule
\end{tabular}
\caption{Sample Data Table}
\label{tab:sample}
\end{table}

\section{Lists}
\subsection{Bullet Points}
\begin{itemize}
  \item First item
  \item Second item
  \item Third item with sub-items:
    \begin{itemize}
      \item Sub-item 1
      \item Sub-item 2
    \end{itemize}
\end{itemize}

\subsection{Numbered List}
\begin{enumerate}
  \item First step
  \item Second step
  \item Third step
\end{enumerate}

\section{Citations and References}
When you use the citation search feature, references will be automatically managed. For example, you might cite a paper like this: \cite{example2024}. All citations will be automatically added to your \texttt{cite.bib} file.

\section{Conclusion}
This template demonstrates basic \LaTeX{} features including mathematical formulas, tables, and lists. The bibliography system is now set up to work seamlessly with Bibby AI's citation management. Feel free to modify it for your needs!

\bibliographystyle{plain}
\bibliography{cite}

\end{document}\documentclass{article}
\usepackage[utf8]{inputenc}
\usepackage{booktabs}
\usepackage{amsmath}
\usepackage{natbib}

\title{My Document}
\author{Your Name}
\date{\today}

\begin{document}

\maketitle

\section{Introduction}
This is a sample document demonstrating various \LaTeX{} features including mathematical equations, tables, and formatting. You can cite papers using the citation search feature, and they will be automatically added to your \texttt{cite.bib} file.

\section{Mathematical Expressions}
Here are some example equations:

\subsection{Quadratic Formula}
The quadratic formula is given by:
\[x = \frac{-b \pm \sqrt{b^2 - 4ac}}{2a}\]

\subsection{Matrix Example}
A sample matrix:
\[
\begin{pmatrix}
a_{11} & a_{12} & a_{13} \\
a_{21} & a_{22} & a_{23} \\
a_{31} & a_{32} & a_{33}
\end{pmatrix}
\]

\section{Tables}
Below is a sample table showing data:

\begin{table}[h]
\centering
\begin{tabular}{@{}lrrr@{}}
\toprule
\textbf{Category} & \textbf{Value 1} & \textbf{Value 2} & \textbf{Total} \\
\midrule
Item A & 10.5 & 20.3 & 30.8 \\
Item B & 15.2 & 25.1 & 40.3 \\
Item C & 22.1 & 12.7 & 34.8 \\
\bottomrule
\end{tabular}
\caption{Sample Data Table}
\label{tab:sample}
\end{table}

\section{Lists}
\subsection{Bullet Points}
\begin{itemize}
  \item First item
  \item Second item
  \item Third item with sub-items:
    \begin{itemize}
      \item Sub-item 1
      \item Sub-item 2
    \end{itemize}
\end{itemize}

\subsection{Numbered List}
\begin{enumerate}
  \item First step
  \item Second step
  \item Third step
\end{enumerate}

\section{Citations and References}
When you use the citation search feature, references will be automatically managed. For example, you might cite a paper like this: \cite{example2024}. All citations will be automatically added to your \texttt{cite.bib} file.

\section{Conclusion}
This template demonstrates basic \LaTeX{} features including mathematical formulas, tables, and lists. The bibliography system is now set up to work seamlessly with Bibby AI's citation management. Feel free to modify it for your needs!

\bibliographystyle{plain}
\bibliography{cite}

\end{document}\documentclass{article}
\usepackage[utf8]{inputenc}
\usepackage{booktabs}
\usepackage{amsmath}
\usepackage{natbib}

\title{My Document}
\author{Your Name}
\date{\today}

\begin{document}

\maketitle

\section{Introduction}
This is a sample document demonstrating various \LaTeX{} features including mathematical equations, tables, and formatting. You can cite papers using the citation search feature, and they will be automatically added to your \texttt{cite.bib} file.

\section{Mathematical Expressions}
Here are some example equations:

\subsection{Quadratic Formula}
The quadratic formula is given by:
\[x = \frac{-b \pm \sqrt{b^2 - 4ac}}{2a}\]

\subsection{Matrix Example}
A sample matrix:
\[
\begin{pmatrix}
a_{11} & a_{12} & a_{13} \\
a_{21} & a_{22} & a_{23} \\
a_{31} & a_{32} & a_{33}
\end{pmatrix}
\]

\section{Tables}
Below is a sample table showing data:

\begin{table}[h]
\centering
\begin{tabular}{@{}lrrr@{}}
\toprule
\textbf{Category} & \textbf{Value 1} & \textbf{Value 2} & \textbf{Total} \\
\midrule
Item A & 10.5 & 20.3 & 30.8 \\
Item B & 15.2 & 25.1 & 40.3 \\
Item C & 22.1 & 12.7 & 34.8 \\
\bottomrule
\end{tabular}
\caption{Sample Data Table}
\label{tab:sample}
\end{table}

\section{Lists}
\subsection{Bullet Points}
\begin{itemize}
  \item First item
  \item Second item
  \item Third item with sub-items:
    \begin{itemize}
      \item Sub-item 1
      \item Sub-item 2
    \end{itemize}
\end{itemize}

\subsection{Numbered List}
\begin{enumerate}
  \item First step
  \item Second step
  \item Third step
\end{enumerate}

\section{Citations and References}
When you use the citation search feature, references will be automatically managed. For example, you might cite a paper like this: \cite{example2024}. All citations will be automatically added to your \texttt{cite.bib} file.

\section{Conclusion}
This template demonstrates basic \LaTeX{} features including mathematical formulas, tables, and lists. The bibliography system is now set up to work seamlessly with Bibby AI's citation management. Feel free to modify it for your needs!

\bibliographystyle{plain}
\bibliography{cite}

\end{document}\documentclass{article}
\usepackage[utf8]{inputenc}
\usepackage{booktabs}
\usepackage{amsmath}
\usepackage{natbib}

\title{My Document}
\author{Your Name}
\date{\today}

\begin{document}

\maketitle

\section{Introduction}
This is a sample document demonstrating various \LaTeX{} features including mathematical equations, tables, and formatting. You can cite papers using the citation search feature, and they will be automatically added to your \texttt{cite.bib} file.

\section{Mathematical Expressions}
Here are some example equations:

\subsection{Quadratic Formula}
The quadratic formula is given by:
\[x = \frac{-b \pm \sqrt{b^2 - 4ac}}{2a}\]

\subsection{Matrix Example}
A sample matrix:
\[
\begin{pmatrix}
a_{11} & a_{12} & a_{13} \\
a_{21} & a_{22} & a_{23} \\
a_{31} & a_{32} & a_{33}
\end{pmatrix}
\]

\section{Tables}
Below is a sample table showing data:

\begin{table}[h]
\centering
\begin{tabular}{@{}lrrr@{}}
\toprule
\textbf{Category} & \textbf{Value 1} & \textbf{Value 2} & \textbf{Total} \\
\midrule
Item A & 10.5 & 20.3 & 30.8 \\
Item B & 15.2 & 25.1 & 40.3 \\
Item C & 22.1 & 12.7 & 34.8 \\
\bottomrule
\end{tabular}
\caption{Sample Data Table}
\label{tab:sample}
\end{table}

\section{Lists}
\subsection{Bullet Points}
\begin{itemize}
  \item First item
  \item Second item
  \item Third item with sub-items:
    \begin{itemize}
      \item Sub-item 1
      \item Sub-item 2
    \end{itemize}
\end{itemize}

\subsection{Numbered List}
\begin{enumerate}
  \item First step
  \item Second step
  \item Third step
\end{enumerate}

\section{Citations and References}
When you use the citation search feature, references will be automatically managed. For example, you might cite a paper like this: \cite{example2024}. All citations will be automatically added to your \texttt{cite.bib} file.

\section{Conclusion}
This template demonstrates basic \LaTeX{} features including mathematical formulas, tables, and lists. The bibliography system is now set up to work seamlessly with Bibby AI's citation management. Feel free to modify it for your needs!

\bibliographystyle{plain}
\bibliography{cite}

\end{document}\documentclass{article}
\usepackage[utf8]{inputenc}
\usepackage{booktabs}
\usepackage{amsmath}
\usepackage{natbib}

\title{My Document}
\author{Your Name}
\date{\today}

\begin{document}

\maketitle

\section{Introduction}
This is a sample document demonstrating various \LaTeX{} features including mathematical equations, tables, and formatting. You can cite papers using the citation search feature, and they will be automatically added to your \texttt{cite.bib} file.

\section{Mathematical Expressions}
Here are some example equations:

\subsection{Quadratic Formula}
The quadratic formula is given by:
\[x = \frac{-b \pm \sqrt{b^2 - 4ac}}{2a}\]

\subsection{Matrix Example}
A sample matrix:
\[
\begin{pmatrix}
a_{11} & a_{12} & a_{13} \\
a_{21} & a_{22} & a_{23} \\
a_{31} & a_{32} & a_{33}
\end{pmatrix}
\]

\section{Tables}
Below is a sample table showing data:

\begin{table}[h]
\centering
\begin{tabular}{@{}lrrr@{}}
\toprule
\textbf{Category} & \textbf{Value 1} & \textbf{Value 2} & \textbf{Total} \\
\midrule
Item A & 10.5 & 20.3 & 30.8 \\
Item B & 15.2 & 25.1 & 40.3 \\
Item C & 22.1 & 12.7 & 34.8 \\
\bottomrule
\end{tabular}
\caption{Sample Data Table}
\label{tab:sample}
\end{table}

\section{Lists}
\subsection{Bullet Points}
\begin{itemize}
  \item First item
  \item Second item
  \item Third item with sub-items:
    \begin{itemize}
      \item Sub-item 1
      \item Sub-item 2
    \end{itemize}
\end{itemize}

\subsection{Numbered List}
\begin{enumerate}
  \item First step
  \item Second step
  \item Third step
\end{enumerate}

\section{Citations and References}
When you use the citation search feature, references will be automatically managed. For example, you might cite a paper like this: \cite{example2024}. All citations will be automatically added to your \texttt{cite.bib} file.

\section{Conclusion}
This template demonstrates basic \LaTeX{} features including mathematical formulas, tables, and lists. The bibliography system is now set up to work seamlessly with Bibby AI's citation management. Feel free to modify it for your needs!

\bibliographystyle{plain}
\bibliography{cite}

\end{document}\documentclass{article}
\usepackage[utf8]{inputenc}
\usepackage{booktabs}
\usepackage{amsmath}
\usepackage{natbib}

\title{My Document}
\author{Your Name}
\date{\today}

\begin{document}

\maketitle

\section{Introduction}
This is a sample document demonstrating various \LaTeX{} features including mathematical equations, tables, and formatting. You can cite papers using the citation search feature, and they will be automatically added to your \texttt{cite.bib} file.

\section{Mathematical Expressions}
Here are some example equations:

\subsection{Quadratic Formula}
The quadratic formula is given by:
\[x = \frac{-b \pm \sqrt{b^2 - 4ac}}{2a}\]

\subsection{Matrix Example}
A sample matrix:
\[
\begin{pmatrix}
a_{11} & a_{12} & a_{13} \\
a_{21} & a_{22} & a_{23} \\
a_{31} & a_{32} & a_{33}
\end{pmatrix}
\]

\section{Tables}
Below is a sample table showing data:

\begin{table}[h]
\centering
\begin{tabular}{@{}lrrr@{}}
\toprule
\textbf{Category} & \textbf{Value 1} & \textbf{Value 2} & \textbf{Total} \\
\midrule
Item A & 10.5 & 20.3 & 30.8 \\
Item B & 15.2 & 25.1 & 40.3 \\
Item C & 22.1 & 12.7 & 34.8 \\
\bottomrule
\end{tabular}
\caption{Sample Data Table}
\label{tab:sample}
\end{table}

\section{Lists}
\subsection{Bullet Points}
\begin{itemize}
  \item First item
  \item Second item
  \item Third item with sub-items:
    \begin{itemize}
      \item Sub-item 1
      \item Sub-item 2
    \end{itemize}
\end{itemize}

\subsection{Numbered List}
\begin{enumerate}
  \item First step
  \item Second step
  \item Third step
\end{enumerate}

\section{Citations and References}
When you use the citation search feature, references will be automatically managed. For example, you might cite a paper like this: \cite{example2024}. All citations will be automatically added to your \texttt{cite.bib} file.

\section{Conclusion}
This template demonstrates basic \LaTeX{} features including mathematical formulas, tables, and lists. The bibliography system is now set up to work seamlessly with Bibby AI's citation management. Feel free to modify it for your needs!

\bibliographystyle{plain}
\bibliography{cite}

\end{document}\documentclass{article}
\usepackage[utf8]{inputenc}
\usepackage{booktabs}
\usepackage{amsmath}
\usepackage{natbib}

\title{My Document}
\author{Your Name}
\date{\today}

\begin{document}

\maketitle

\section{Introduction}
This is a sample document demonstrating various \LaTeX{} features including mathematical equations, tables, and formatting. You can cite papers using the citation search feature, and they will be automatically added to your \texttt{cite.bib} file.

\section{Mathematical Expressions}
Here are some example equations:

\subsection{Quadratic Formula}
The quadratic formula is given by:
\[x = \frac{-b \pm \sqrt{b^2 - 4ac}}{2a}\]

\subsection{Matrix Example}
A sample matrix:
\[
\begin{pmatrix}
a_{11} & a_{12} & a_{13} \\
a_{21} & a_{22} & a_{23} \\
a_{31} & a_{32} & a_{33}
\end{pmatrix}
\]

\section{Tables}
Below is a sample table showing data:

\begin{table}[h]
\centering
\begin{tabular}{@{}lrrr@{}}
\toprule
\textbf{Category} & \textbf{Value 1} & \textbf{Value 2} & \textbf{Total} \\
\midrule
Item A & 10.5 & 20.3 & 30.8 \\
Item B & 15.2 & 25.1 & 40.3 \\
Item C & 22.1 & 12.7 & 34.8 \\
\bottomrule
\end{tabular}
\caption{Sample Data Table}
\label{tab:sample}
\end{table}

\section{Lists}
\subsection{Bullet Points}
\begin{itemize}
  \item First item
  \item Second item
  \item Third item with sub-items:
    \begin{itemize}
      \item Sub-item 1
      \item Sub-item 2
    \end{itemize}
\end{itemize}

\subsection{Numbered List}
\begin{enumerate}
  \item First step
  \item Second step
  \item Third step
\end{enumerate}

\section{Citations and References}
When you use the citation search feature, references will be automatically managed. For example, you might cite a paper like this: \cite{example2024}. All citations will be automatically added to your \texttt{cite.bib} file.

\section{Conclusion}
This template demonstrates basic \LaTeX{} features including mathematical formulas, tables, and lists. The bibliography system is now set up to work seamlessly with Bibby AI's citation management. Feel free to modify it for your needs!

\bibliographystyle{plain}
\bibliography{cite}

\end{document}\documentclass{article}
\usepackage[utf8]{inputenc}
\usepackage{booktabs}
\usepackage{amsmath}
\usepackage{natbib}

\title{My Document}
\author{Your Name}
\date{\today}

\begin{document}

\maketitle

\section{Introduction}
This is a sample document demonstrating various \LaTeX{} features including mathematical equations, tables, and formatting. You can cite papers using the citation search feature, and they will be automatically added to your \texttt{cite.bib} file.

\section{Mathematical Expressions}
Here are some example equations:

\subsection{Quadratic Formula}
The quadratic formula is given by:
\[x = \frac{-b \pm \sqrt{b^2 - 4ac}}{2a}\]

\subsection{Matrix Example}
A sample matrix:
\[
\begin{pmatrix}
a_{11} & a_{12} & a_{13} \\
a_{21} & a_{22} & a_{23} \\
a_{31} & a_{32} & a_{33}
\end{pmatrix}
\]

\section{Tables}
Below is a sample table showing data:

\begin{table}[h]
\centering
\begin{tabular}{@{}lrrr@{}}
\toprule
\textbf{Category} & \textbf{Value 1} & \textbf{Value 2} & \textbf{Total} \\
\midrule
Item A & 10.5 & 20.3 & 30.8 \\
Item B & 15.2 & 25.1 & 40.3 \\
Item C & 22.1 & 12.7 & 34.8 \\
\bottomrule
\end{tabular}
\caption{Sample Data Table}
\label{tab:sample}
\end{table}

\section{Lists}
\subsection{Bullet Points}
\begin{itemize}
  \item First item
  \item Second item
  \item Third item with sub-items:
    \begin{itemize}
      \item Sub-item 1
      \item Sub-item 2
    \end{itemize}
\end{itemize}

\subsection{Numbered List}
\begin{enumerate}
  \item First step
  \item Second step
  \item Third step
\end{enumerate}

\section{Citations and References}
When you use the citation search feature, references will be automatically managed. For example, you might cite a paper like this: \cite{example2024}. All citations will be automatically added to your \texttt{cite.bib} file.

\section{Conclusion}
This template demonstrates basic \LaTeX{} features including mathematical formulas, tables, and lists. The bibliography system is now set up to work seamlessly with Bibby AI's citation management. Feel free to modify it for your needs!

\bibliographystyle{plain}
\bibliography{cite}

\end{document}\documentclass{article}
\usepackage[utf8]{inputenc}
\usepackage{booktabs}
\usepackage{amsmath}
\usepackage{natbib}

\title{My Document}
\author{Your Name}
\date{\today}

\begin{document}

\maketitle

\section{Introduction}
This is a sample document demonstrating various \LaTeX{} features including mathematical equations, tables, and formatting. You can cite papers using the citation search feature, and they will be automatically added to your \texttt{cite.bib} file.

\section{Mathematical Expressions}
Here are some example equations:

\subsection{Quadratic Formula}
The quadratic formula is given by:
\[x = \frac{-b \pm \sqrt{b^2 - 4ac}}{2a}\]

\subsection{Matrix Example}
A sample matrix:
\[
\begin{pmatrix}
a_{11} & a_{12} & a_{13} \\
a_{21} & a_{22} & a_{23} \\
a_{31} & a_{32} & a_{33}
\end{pmatrix}
\]

\section{Tables}
Below is a sample table showing data:

\begin{table}[h]
\centering
\begin{tabular}{@{}lrrr@{}}
\toprule
\textbf{Category} & \textbf{Value 1} & \textbf{Value 2} & \textbf{Total} \\
\midrule
Item A & 10.5 & 20.3 & 30.8 \\
Item B & 15.2 & 25.1 & 40.3 \\
Item C & 22.1 & 12.7 & 34.8 \\
\bottomrule
\end{tabular}
\caption{Sample Data Table}
\label{tab:sample}
\end{table}

\section{Lists}
\subsection{Bullet Points}
\begin{itemize}
  \item First item
  \item Second item
  \item Third item with sub-items:
    \begin{itemize}
      \item Sub-item 1
      \item Sub-item 2
    \end{itemize}
\end{itemize}

\subsection{Numbered List}
\begin{enumerate}
  \item First step
  \item Second step
  \item Third step
\end{enumerate}

\section{Citations and References}
When you use the citation search feature, references will be automatically managed. For example, you might cite a paper like this: \cite{example2024}. All citations will be automatically added to your \texttt{cite.bib} file.

\section{Conclusion}
This template demonstrates basic \LaTeX{} features including mathematical formulas, tables, and lists. The bibliography system is now set up to work seamlessly with Bibby AI's citation management. Feel free to modify it for your needs!

\bibliographystyle{plain}
\bibliography{cite}

\end{document}\documentclass{article}
\usepackage[utf8]{inputenc}
\usepackage{booktabs}
\usepackage{amsmath}
\usepackage{natbib}

\title{My Document}
\author{Your Name}
\date{\today}

\begin{document}

\maketitle

\section{Introduction}
This is a sample document demonstrating various \LaTeX{} features including mathematical equations, tables, and formatting. You can cite papers using the citation search feature, and they will be automatically added to your \texttt{cite.bib} file.

\section{Mathematical Expressions}
Here are some example equations:

\subsection{Quadratic Formula}
The quadratic formula is given by:
\[x = \frac{-b \pm \sqrt{b^2 - 4ac}}{2a}\]

\subsection{Matrix Example}
A sample matrix:
\[
\begin{pmatrix}
a_{11} & a_{12} & a_{13} \\
a_{21} & a_{22} & a_{23} \\
a_{31} & a_{32} & a_{33}
\end{pmatrix}
\]

\section{Tables}
Below is a sample table showing data:

\begin{table}[h]
\centering
\begin{tabular}{@{}lrrr@{}}
\toprule
\textbf{Category} & \textbf{Value 1} & \textbf{Value 2} & \textbf{Total} \\
\midrule
Item A & 10.5 & 20.3 & 30.8 \\
Item B & 15.2 & 25.1 & 40.3 \\
Item C & 22.1 & 12.7 & 34.8 \\
\bottomrule
\end{tabular}
\caption{Sample Data Table}
\label{tab:sample}
\end{table}

\section{Lists}
\subsection{Bullet Points}
\begin{itemize}
  \item First item
  \item Second item
  \item Third item with sub-items:
    \begin{itemize}
      \item Sub-item 1
      \item Sub-item 2
    \end{itemize}
\end{itemize}

\subsection{Numbered List}
\begin{enumerate}
  \item First step
  \item Second step
  \item Third step
\end{enumerate}

\section{Citations and References}
When you use the citation search feature, references will be automatically managed. For example, you might cite a paper like this: \cite{example2024}. All citations will be automatically added to your \texttt{cite.bib} file.

\section{Conclusion}
This template demonstrates basic \LaTeX{} features including mathematical formulas, tables, and lists. The bibliography system is now set up to work seamlessly with Bibby AI's citation management. Feel free to modify it for your needs!

\bibliographystyle{plain}
\bibliography{cite}

\end{document}\documentclass{article}
\usepackage[utf8]{inputenc}
\usepackage{booktabs}
\usepackage{amsmath}
\usepackage{natbib}

\title{My Document}
\author{Your Name}
\date{\today}

\begin{document}

\maketitle

\section{Introduction}
This is a sample document demonstrating various \LaTeX{} features including mathematical equations, tables, and formatting. You can cite papers using the citation search feature, and they will be automatically added to your \texttt{cite.bib} file.

\section{Mathematical Expressions}
Here are some example equations:

\subsection{Quadratic Formula}
The quadratic formula is given by:
\[x = \frac{-b \pm \sqrt{b^2 - 4ac}}{2a}\]

\subsection{Matrix Example}
A sample matrix:
\[
\begin{pmatrix}
a_{11} & a_{12} & a_{13} \\
a_{21} & a_{22} & a_{23} \\
a_{31} & a_{32} & a_{33}
\end{pmatrix}
\]

\section{Tables}
Below is a sample table showing data:

\begin{table}[h]
\centering
\begin{tabular}{@{}lrrr@{}}
\toprule
\textbf{Category} & \textbf{Value 1} & \textbf{Value 2} & \textbf{Total} \\
\midrule
Item A & 10.5 & 20.3 & 30.8 \\
Item B & 15.2 & 25.1 & 40.3 \\
Item C & 22.1 & 12.7 & 34.8 \\
\bottomrule
\end{tabular}
\caption{Sample Data Table}
\label{tab:sample}
\end{table}

\section{Lists}
\subsection{Bullet Points}
\begin{itemize}
  \item First item
  \item Second item
  \item Third item with sub-items:
    \begin{itemize}
      \item Sub-item 1
      \item Sub-item 2
    \end{itemize}
\end{itemize}

\subsection{Numbered List}
\begin{enumerate}
  \item First step
  \item Second step
  \item Third step
\end{enumerate}

\section{Citations and References}
When you use the citation search feature, references will be automatically managed. For example, you might cite a paper like this: \cite{example2024}. All citations will be automatically added to your \texttt{cite.bib} file.

\section{Conclusion}
This template demonstrates basic \LaTeX{} features including mathematical formulas, tables, and lists. The bibliography system is now set up to work seamlessly with Bibby AI's citation management. Feel free to modify it for your needs!

\bibliographystyle{plain}
\bibliography{cite}

\end{document}\documentclass{article}
\usepackage[utf8]{inputenc}
\usepackage{booktabs}
\usepackage{amsmath}
\usepackage{natbib}

\title{My Document}
\author{Your Name}
\date{\today}

\begin{document}

\maketitle

\section{Introduction}
This is a sample document demonstrating various \LaTeX{} features including mathematical equations, tables, and formatting. You can cite papers using the citation search feature, and they will be automatically added to your \texttt{cite.bib} file.

\section{Mathematical Expressions}
Here are some example equations:

\subsection{Quadratic Formula}
The quadratic formula is given by:
\[x = \frac{-b \pm \sqrt{b^2 - 4ac}}{2a}\]

\subsection{Matrix Example}
A sample matrix:
\[
\begin{pmatrix}
a_{11} & a_{12} & a_{13} \\
a_{21} & a_{22} & a_{23} \\
a_{31} & a_{32} & a_{33}
\end{pmatrix}
\]

\section{Tables}
Below is a sample table showing data:

\begin{table}[h]
\centering
\begin{tabular}{@{}lrrr@{}}
\toprule
\textbf{Category} & \textbf{Value 1} & \textbf{Value 2} & \textbf{Total} \\
\midrule
Item A & 10.5 & 20.3 & 30.8 \\
Item B & 15.2 & 25.1 & 40.3 \\
Item C & 22.1 & 12.7 & 34.8 \\
\bottomrule
\end{tabular}
\caption{Sample Data Table}
\label{tab:sample}
\end{table}

\section{Lists}
\subsection{Bullet Points}
\begin{itemize}
  \item First item
  \item Second item
  \item Third item with sub-items:
    \begin{itemize}
      \item Sub-item 1
      \item Sub-item 2
    \end{itemize}
\end{itemize}

\subsection{Numbered List}
\begin{enumerate}
  \item First step
  \item Second step
  \item Third step
\end{enumerate}

\section{Citations and References}
When you use the citation search feature, references will be automatically managed. For example, you might cite a paper like this: \cite{example2024}. All citations will be automatically added to your \texttt{cite.bib} file.

\section{Conclusion}
This template demonstrates basic \LaTeX{} features including mathematical formulas, tables, and lists. The bibliography system is now set up to work seamlessly with Bibby AI's citation management. Feel free to modify it for your needs!

\bibliographystyle{plain}
\bibliography{cite}

\end{document}\documentclass{article}
\usepackage[utf8]{inputenc}
\usepackage{booktabs}
\usepackage{amsmath}
\usepackage{natbib}

\title{My Document}
\author{Your Name}
\date{\today}

\begin{document}

\maketitle

\section{Introduction}
This is a sample document demonstrating various \LaTeX{} features including mathematical equations, tables, and formatting. You can cite papers using the citation search feature, and they will be automatically added to your \texttt{cite.bib} file.

\section{Mathematical Expressions}
Here are some example equations:

\subsection{Quadratic Formula}
The quadratic formula is given by:
\[x = \frac{-b \pm \sqrt{b^2 - 4ac}}{2a}\]

\subsection{Matrix Example}
A sample matrix:
\[
\begin{pmatrix}
a_{11} & a_{12} & a_{13} \\
a_{21} & a_{22} & a_{23} \\
a_{31} & a_{32} & a_{33}
\end{pmatrix}
\]

\section{Tables}
Below is a sample table showing data:

\begin{table}[h]
\centering
\begin{tabular}{@{}lrrr@{}}
\toprule
\textbf{Category} & \textbf{Value 1} & \textbf{Value 2} & \textbf{Total} \\
\midrule
Item A & 10.5 & 20.3 & 30.8 \\
Item B & 15.2 & 25.1 & 40.3 \\
Item C & 22.1 & 12.7 & 34.8 \\
\bottomrule
\end{tabular}
\caption{Sample Data Table}
\label{tab:sample}
\end{table}

\section{Lists}
\subsection{Bullet Points}
\begin{itemize}
  \item First item
  \item Second item
  \item Third item with sub-items:
    \begin{itemize}
      \item Sub-item 1
      \item Sub-item 2
    \end{itemize}
\end{itemize}

\subsection{Numbered List}
\begin{enumerate}
  \item First step
  \item Second step
  \item Third step
\end{enumerate}

\section{Citations and References}
When you use the citation search feature, references will be automatically managed. For example, you might cite a paper like this: \cite{example2024}. All citations will be automatically added to your \texttt{cite.bib} file.

\section{Conclusion}
This template demonstrates basic \LaTeX{} features including mathematical formulas, tables, and lists. The bibliography system is now set up to work seamlessly with Bibby AI's citation management. Feel free to modify it for your needs!

\bibliographystyle{plain}
\bibliography{cite}

\end{document}\documentclass{article}
\usepackage[utf8]{inputenc}
\usepackage{booktabs}
\usepackage{amsmath}
\usepackage{natbib}

\title{My Document}
\author{Your Name}
\date{\today}

\begin{document}

\maketitle

\section{Introduction}
This is a sample document demonstrating various \LaTeX{} features including mathematical equations, tables, and formatting. You can cite papers using the citation search feature, and they will be automatically added to your \texttt{cite.bib} file.

\section{Mathematical Expressions}
Here are some example equations:

\subsection{Quadratic Formula}
The quadratic formula is given by:
\[x = \frac{-b \pm \sqrt{b^2 - 4ac}}{2a}\]

\subsection{Matrix Example}
A sample matrix:
\[
\begin{pmatrix}
a_{11} & a_{12} & a_{13} \\
a_{21} & a_{22} & a_{23} \\
a_{31} & a_{32} & a_{33}
\end{pmatrix}
\]

\section{Tables}
Below is a sample table showing data:

\begin{table}[h]
\centering
\begin{tabular}{@{}lrrr@{}}
\toprule
\textbf{Category} & \textbf{Value 1} & \textbf{Value 2} & \textbf{Total} \\
\midrule
Item A & 10.5 & 20.3 & 30.8 \\
Item B & 15.2 & 25.1 & 40.3 \\
Item C & 22.1 & 12.7 & 34.8 \\
\bottomrule
\end{tabular}
\caption{Sample Data Table}
\label{tab:sample}
\end{table}

\section{Lists}
\subsection{Bullet Points}
\begin{itemize}
  \item First item
  \item Second item
  \item Third item with sub-items:
    \begin{itemize}
      \item Sub-item 1
      \item Sub-item 2
    \end{itemize}
\end{itemize}

\subsection{Numbered List}
\begin{enumerate}
  \item First step
  \item Second step
  \item Third step
\end{enumerate}

\section{Citations and References}
When you use the citation search feature, references will be automatically managed. For example, you might cite a paper like this: \cite{example2024}. All citations will be automatically added to your \texttt{cite.bib} file.

\section{Conclusion}
This template demonstrates basic \LaTeX{} features including mathematical formulas, tables, and lists. The bibliography system is now set up to work seamlessly with Bibby AI's citation management. Feel free to modify it for your needs!

\bibliographystyle{plain}
\bibliography{cite}

\end{document}\documentclass{article}
\usepackage[utf8]{inputenc}
\usepackage{booktabs}
\usepackage{amsmath}
\usepackage{natbib}

\title{My Document}
\author{Your Name}
\date{\today}

\begin{document}

\maketitle

\section{Introduction}
This is a sample document demonstrating various \LaTeX{} features including mathematical equations, tables, and formatting. You can cite papers using the citation search feature, and they will be automatically added to your \texttt{cite.bib} file.

\section{Mathematical Expressions}
Here are some example equations:

\subsection{Quadratic Formula}
The quadratic formula is given by:
\[x = \frac{-b \pm \sqrt{b^2 - 4ac}}{2a}\]

\subsection{Matrix Example}
A sample matrix:
\[
\begin{pmatrix}
a_{11} & a_{12} & a_{13} \\
a_{21} & a_{22} & a_{23} \\
a_{31} & a_{32} & a_{33}
\end{pmatrix}
\]

\section{Tables}
Below is a sample table showing data:

\begin{table}[h]
\centering
\begin{tabular}{@{}lrrr@{}}
\toprule
\textbf{Category} & \textbf{Value 1} & \textbf{Value 2} & \textbf{Total} \\
\midrule
Item A & 10.5 & 20.3 & 30.8 \\
Item B & 15.2 & 25.1 & 40.3 \\
Item C & 22.1 & 12.7 & 34.8 \\
\bottomrule
\end{tabular}
\caption{Sample Data Table}
\label{tab:sample}
\end{table}

\section{Lists}
\subsection{Bullet Points}
\begin{itemize}
  \item First item
  \item Second item
  \item Third item with sub-items:
    \begin{itemize}
      \item Sub-item 1
      \item Sub-item 2
    \end{itemize}
\end{itemize}

\subsection{Numbered List}
\begin{enumerate}
  \item First step
  \item Second step
  \item Third step
\end{enumerate}

\section{Citations and References}
When you use the citation search feature, references will be automatically managed. For example, you might cite a paper like this: \cite{example2024}. All citations will be automatically added to your \texttt{cite.bib} file.

\section{Conclusion}
This template demonstrates basic \LaTeX{} features including mathematical formulas, tables, and lists. The bibliography system is now set up to work seamlessly with Bibby AI's citation management. Feel free to modify it for your needs!

\bibliographystyle{plain}
\bibliography{cite}

\end{document}\documentclass{article}
\usepackage[utf8]{inputenc}
\usepackage{booktabs}
\usepackage{amsmath}
\usepackage{natbib}

\title{My Document}
\author{Your Name}
\date{\today}

\begin{document}

\maketitle

\section{Introduction}
This is a sample document demonstrating various \LaTeX{} features including mathematical equations, tables, and formatting. You can cite papers using the citation search feature, and they will be automatically added to your \texttt{cite.bib} file.

\section{Mathematical Expressions}
Here are some example equations:

\subsection{Quadratic Formula}
The quadratic formula is given by:
\[x = \frac{-b \pm \sqrt{b^2 - 4ac}}{2a}\]

\subsection{Matrix Example}
A sample matrix:
\[
\begin{pmatrix}
a_{11} & a_{12} & a_{13} \\
a_{21} & a_{22} & a_{23} \\
a_{31} & a_{32} & a_{33}
\end{pmatrix}
\]

\section{Tables}
Below is a sample table showing data:

\begin{table}[h]
\centering
\begin{tabular}{@{}lrrr@{}}
\toprule
\textbf{Category} & \textbf{Value 1} & \textbf{Value 2} & \textbf{Total} \\
\midrule
Item A & 10.5 & 20.3 & 30.8 \\
Item B & 15.2 & 25.1 & 40.3 \\
Item C & 22.1 & 12.7 & 34.8 \\
\bottomrule
\end{tabular}
\caption{Sample Data Table}
\label{tab:sample}
\end{table}

\section{Lists}
\subsection{Bullet Points}
\begin{itemize}
  \item First item
  \item Second item
  \item Third item with sub-items:
    \begin{itemize}
      \item Sub-item 1
      \item Sub-item 2
    \end{itemize}
\end{itemize}

\subsection{Numbered List}
\begin{enumerate}
  \item First step
  \item Second step
  \item Third step
\end{enumerate}

\section{Citations and References}
When you use the citation search feature, references will be automatically managed. For example, you might cite a paper like this: \cite{example2024}. All citations will be automatically added to your \texttt{cite.bib} file.

\section{Conclusion}
This template demonstrates basic \LaTeX{} features including mathematical formulas, tables, and lists. The bibliography system is now set up to work seamlessly with Bibby AI's citation management. Feel free to modify it for your needs!

\bibliographystyle{plain}
\bibliography{cite}

\end{document}\documentclass{article}
\usepackage[utf8]{inputenc}
\usepackage{booktabs}
\usepackage{amsmath}
\usepackage{natbib}

\title{My Document}
\author{Your Name}
\date{\today}

\begin{document}

\maketitle

\section{Introduction}
This is a sample document demonstrating various \LaTeX{} features including mathematical equations, tables, and formatting. You can cite papers using the citation search feature, and they will be automatically added to your \texttt{cite.bib} file.

\section{Mathematical Expressions}
Here are some example equations:

\subsection{Quadratic Formula}
The quadratic formula is given by:
\[x = \frac{-b \pm \sqrt{b^2 - 4ac}}{2a}\]

\subsection{Matrix Example}
A sample matrix:
\[
\begin{pmatrix}
a_{11} & a_{12} & a_{13} \\
a_{21} & a_{22} & a_{23} \\
a_{31} & a_{32} & a_{33}
\end{pmatrix}
\]

\section{Tables}
Below is a sample table showing data:

\begin{table}[h]
\centering
\begin{tabular}{@{}lrrr@{}}
\toprule
\textbf{Category} & \textbf{Value 1} & \textbf{Value 2} & \textbf{Total} \\
\midrule
Item A & 10.5 & 20.3 & 30.8 \\
Item B & 15.2 & 25.1 & 40.3 \\
Item C & 22.1 & 12.7 & 34.8 \\
\bottomrule
\end{tabular}
\caption{Sample Data Table}
\label{tab:sample}
\end{table}

\section{Lists}
\subsection{Bullet Points}
\begin{itemize}
  \item First item
  \item Second item
  \item Third item with sub-items:
    \begin{itemize}
      \item Sub-item 1
      \item Sub-item 2
    \end{itemize}
\end{itemize}

\subsection{Numbered List}
\begin{enumerate}
  \item First step
  \item Second step
  \item Third step
\end{enumerate}

\section{Citations and References}
When you use the citation search feature, references will be automatically managed. For example, you might cite a paper like this: \cite{example2024}. All citations will be automatically added to your \texttt{cite.bib} file.

\section{Conclusion}
This template demonstrates basic \LaTeX{} features including mathematical formulas, tables, and lists. The bibliography system is now set up to work seamlessly with Bibby AI's citation management. Feel free to modify it for your needs!

\bibliographystyle{plain}
\bibliography{cite}

\end{document}\documentclass{article}
\usepackage[utf8]{inputenc}
\usepackage{booktabs}
\usepackage{amsmath}
\usepackage{natbib}

\title{My Document}
\author{Your Name}
\date{\today}

\begin{document}

\maketitle

\section{Introduction}
This is a sample document demonstrating various \LaTeX{} features including mathematical equations, tables, and formatting. You can cite papers using the citation search feature, and they will be automatically added to your \texttt{cite.bib} file.

\section{Mathematical Expressions}
Here are some example equations:

\subsection{Quadratic Formula}
The quadratic formula is given by:
\[x = \frac{-b \pm \sqrt{b^2 - 4ac}}{2a}\]

\subsection{Matrix Example}
A sample matrix:
\[
\begin{pmatrix}
a_{11} & a_{12} & a_{13} \\
a_{21} & a_{22} & a_{23} \\
a_{31} & a_{32} & a_{33}
\end{pmatrix}
\]

\section{Tables}
Below is a sample table showing data:

\begin{table}[h]
\centering
\begin{tabular}{@{}lrrr@{}}
\toprule
\textbf{Category} & \textbf{Value 1} & \textbf{Value 2} & \textbf{Total} \\
\midrule
Item A & 10.5 & 20.3 & 30.8 \\
Item B & 15.2 & 25.1 & 40.3 \\
Item C & 22.1 & 12.7 & 34.8 \\
\bottomrule
\end{tabular}
\caption{Sample Data Table}
\label{tab:sample}
\end{table}

\section{Lists}
\subsection{Bullet Points}
\begin{itemize}
  \item First item
  \item Second item
  \item Third item with sub-items:
    \begin{itemize}
      \item Sub-item 1
      \item Sub-item 2
    \end{itemize}
\end{itemize}

\subsection{Numbered List}
\begin{enumerate}
  \item First step
  \item Second step
  \item Third step
\end{enumerate}

\section{Citations and References}
When you use the citation search feature, references will be automatically managed. For example, you might cite a paper like this: \cite{example2024}. All citations will be automatically added to your \texttt{cite.bib} file.

\section{Conclusion}
This template demonstrates basic \LaTeX{} features including mathematical formulas, tables, and lists. The bibliography system is now set up to work seamlessly with Bibby AI's citation management. Feel free to modify it for your needs!

\bibliographystyle{plain}
\bibliography{cite}

\end{document}\documentclass{article}
\usepackage[utf8]{inputenc}
\usepackage{booktabs}
\usepackage{amsmath}
\usepackage{natbib}

\title{My Document}
\author{Your Name}
\date{\today}

\begin{document}

\maketitle

\section{Introduction}
This is a sample document demonstrating various \LaTeX{} features including mathematical equations, tables, and formatting. You can cite papers using the citation search feature, and they will be automatically added to your \texttt{cite.bib} file.

\section{Mathematical Expressions}
Here are some example equations:

\subsection{Quadratic Formula}
The quadratic formula is given by:
\[x = \frac{-b \pm \sqrt{b^2 - 4ac}}{2a}\]

\subsection{Matrix Example}
A sample matrix:
\[
\begin{pmatrix}
a_{11} & a_{12} & a_{13} \\
a_{21} & a_{22} & a_{23} \\
a_{31} & a_{32} & a_{33}
\end{pmatrix}
\]

\section{Tables}
Below is a sample table showing data:

\begin{table}[h]
\centering
\begin{tabular}{@{}lrrr@{}}
\toprule
\textbf{Category} & \textbf{Value 1} & \textbf{Value 2} & \textbf{Total} \\
\midrule
Item A & 10.5 & 20.3 & 30.8 \\
Item B & 15.2 & 25.1 & 40.3 \\
Item C & 22.1 & 12.7 & 34.8 \\
\bottomrule
\end{tabular}
\caption{Sample Data Table}
\label{tab:sample}
\end{table}

\section{Lists}
\subsection{Bullet Points}
\begin{itemize}
  \item First item
  \item Second item
  \item Third item with sub-items:
    \begin{itemize}
      \item Sub-item 1
      \item Sub-item 2
    \end{itemize}
\end{itemize}

\subsection{Numbered List}
\begin{enumerate}
  \item First step
  \item Second step
  \item Third step
\end{enumerate}

\section{Citations and References}
When you use the citation search feature, references will be automatically managed. For example, you might cite a paper like this: \cite{example2024}. All citations will be automatically added to your \texttt{cite.bib} file.

\section{Conclusion}
This template demonstrates basic \LaTeX{} features including mathematical formulas, tables, and lists. The bibliography system is now set up to work seamlessly with Bibby AI's citation management. Feel free to modify it for your needs!

\bibliographystyle{plain}
\bibliography{cite}

\end{document}\documentclass{article}
\usepackage[utf8]{inputenc}
\usepackage{booktabs}
\usepackage{amsmath}
\usepackage{natbib}

\title{My Document}
\author{Your Name}
\date{\today}

\begin{document}

\maketitle

\section{Introduction}
This is a sample document demonstrating various \LaTeX{} features including mathematical equations, tables, and formatting. You can cite papers using the citation search feature, and they will be automatically added to your \texttt{cite.bib} file.

\section{Mathematical Expressions}
Here are some example equations:

\subsection{Quadratic Formula}
The quadratic formula is given by:
\[x = \frac{-b \pm \sqrt{b^2 - 4ac}}{2a}\]

\subsection{Matrix Example}
A sample matrix:
\[
\begin{pmatrix}
a_{11} & a_{12} & a_{13} \\
a_{21} & a_{22} & a_{23} \\
a_{31} & a_{32} & a_{33}
\end{pmatrix}
\]

\section{Tables}
Below is a sample table showing data:

\begin{table}[h]
\centering
\begin{tabular}{@{}lrrr@{}}
\toprule
\textbf{Category} & \textbf{Value 1} & \textbf{Value 2} & \textbf{Total} \\
\midrule
Item A & 10.5 & 20.3 & 30.8 \\
Item B & 15.2 & 25.1 & 40.3 \\
Item C & 22.1 & 12.7 & 34.8 \\
\bottomrule
\end{tabular}
\caption{Sample Data Table}
\label{tab:sample}
\end{table}

\section{Lists}
\subsection{Bullet Points}
\begin{itemize}
  \item First item
  \item Second item
  \item Third item with sub-items:
    \begin{itemize}
      \item Sub-item 1
      \item Sub-item 2
    \end{itemize}
\end{itemize}

\subsection{Numbered List}
\begin{enumerate}
  \item First step
  \item Second step
  \item Third step
\end{enumerate}

\section{Citations and References}
When you use the citation search feature, references will be automatically managed. For example, you might cite a paper like this: \cite{example2024}. All citations will be automatically added to your \texttt{cite.bib} file.

\section{Conclusion}
This template demonstrates basic \LaTeX{} features including mathematical formulas, tables, and lists. The bibliography system is now set up to work seamlessly with Bibby AI's citation management. Feel free to modify it for your needs!

\bibliographystyle{plain}
\bibliography{cite}

\end{document}\documentclass{article}
\usepackage[utf8]{inputenc}
\usepackage{booktabs}
\usepackage{amsmath}
\usepackage{natbib}

\title{My Document}
\author{Your Name}
\date{\today}

\begin{document}

\maketitle

\section{Introduction}
This is a sample document demonstrating various \LaTeX{} features including mathematical equations, tables, and formatting. You can cite papers using the citation search feature, and they will be automatically added to your \texttt{cite.bib} file.

\section{Mathematical Expressions}
Here are some example equations:

\subsection{Quadratic Formula}
The quadratic formula is given by:
\[x = \frac{-b \pm \sqrt{b^2 - 4ac}}{2a}\]

\subsection{Matrix Example}
A sample matrix:
\[
\begin{pmatrix}
a_{11} & a_{12} & a_{13} \\
a_{21} & a_{22} & a_{23} \\
a_{31} & a_{32} & a_{33}
\end{pmatrix}
\]

\section{Tables}
Below is a sample table showing data:

\begin{table}[h]
\centering
\begin{tabular}{@{}lrrr@{}}
\toprule
\textbf{Category} & \textbf{Value 1} & \textbf{Value 2} & \textbf{Total} \\
\midrule
Item A & 10.5 & 20.3 & 30.8 \\
Item B & 15.2 & 25.1 & 40.3 \\
Item C & 22.1 & 12.7 & 34.8 \\
\bottomrule
\end{tabular}
\caption{Sample Data Table}
\label{tab:sample}
\end{table}

\section{Lists}
\subsection{Bullet Points}
\begin{itemize}
  \item First item
  \item Second item
  \item Third item with sub-items:
    \begin{itemize}
      \item Sub-item 1
      \item Sub-item 2
    \end{itemize}
\end{itemize}

\subsection{Numbered List}
\begin{enumerate}
  \item First step
  \item Second step
  \item Third step
\end{enumerate}

\section{Citations and References}
When you use the citation search feature, references will be automatically managed. For example, you might cite a paper like this: \cite{example2024}. All citations will be automatically added to your \texttt{cite.bib} file.

\section{Conclusion}
This template demonstrates basic \LaTeX{} features including mathematical formulas, tables, and lists. The bibliography system is now set up to work seamlessly with Bibby AI's citation management. Feel free to modify it for your needs!

\bibliographystyle{plain}
\bibliography{cite}

\end{document}\documentclass{article}
\usepackage[utf8]{inputenc}
\usepackage{booktabs}
\usepackage{amsmath}
\usepackage{natbib}

\title{My Document}
\author{Your Name}
\date{\today}

\begin{document}

\maketitle

\section{Introduction}
This is a sample document demonstrating various \LaTeX{} features including mathematical equations, tables, and formatting. You can cite papers using the citation search feature, and they will be automatically added to your \texttt{cite.bib} file.

\section{Mathematical Expressions}
Here are some example equations:

\subsection{Quadratic Formula}
The quadratic formula is given by:
\[x = \frac{-b \pm \sqrt{b^2 - 4ac}}{2a}\]

\subsection{Matrix Example}
A sample matrix:
\[
\begin{pmatrix}
a_{11} & a_{12} & a_{13} \\
a_{21} & a_{22} & a_{23} \\
a_{31} & a_{32} & a_{33}
\end{pmatrix}
\]

\section{Tables}
Below is a sample table showing data:

\begin{table}[h]
\centering
\begin{tabular}{@{}lrrr@{}}
\toprule
\textbf{Category} & \textbf{Value 1} & \textbf{Value 2} & \textbf{Total} \\
\midrule
Item A & 10.5 & 20.3 & 30.8 \\
Item B & 15.2 & 25.1 & 40.3 \\
Item C & 22.1 & 12.7 & 34.8 \\
\bottomrule
\end{tabular}
\caption{Sample Data Table}
\label{tab:sample}
\end{table}

\section{Lists}
\subsection{Bullet Points}
\begin{itemize}
  \item First item
  \item Second item
  \item Third item with sub-items:
    \begin{itemize}
      \item Sub-item 1
      \item Sub-item 2
    \end{itemize}
\end{itemize}

\subsection{Numbered List}
\begin{enumerate}
  \item First step
  \item Second step
  \item Third step
\end{enumerate}

\section{Citations and References}
When you use the citation search feature, references will be automatically managed. For example, you might cite a paper like this: \cite{example2024}. All citations will be automatically added to your \texttt{cite.bib} file.

\section{Conclusion}
This template demonstrates basic \LaTeX{} features including mathematical formulas, tables, and lists. The bibliography system is now set up to work seamlessly with Bibby AI's citation management. Feel free to modify it for your needs!

\bibliographystyle{plain}
\bibliography{cite}

\end{document}\documentclass{article}
\usepackage[utf8]{inputenc}
\usepackage{booktabs}
\usepackage{amsmath}
\usepackage{natbib}

\title{My Document}
\author{Your Name}
\date{\today}

\begin{document}

\maketitle

\section{Introduction}
This is a sample document demonstrating various \LaTeX{} features including mathematical equations, tables, and formatting. You can cite papers using the citation search feature, and they will be automatically added to your \texttt{cite.bib} file.

\section{Mathematical Expressions}
Here are some example equations:

\subsection{Quadratic Formula}
The quadratic formula is given by:
\[x = \frac{-b \pm \sqrt{b^2 - 4ac}}{2a}\]

\subsection{Matrix Example}
A sample matrix:
\[
\begin{pmatrix}
a_{11} & a_{12} & a_{13} \\
a_{21} & a_{22} & a_{23} \\
a_{31} & a_{32} & a_{33}
\end{pmatrix}
\]

\section{Tables}
Below is a sample table showing data:

\begin{table}[h]
\centering
\begin{tabular}{@{}lrrr@{}}
\toprule
\textbf{Category} & \textbf{Value 1} & \textbf{Value 2} & \textbf{Total} \\
\midrule
Item A & 10.5 & 20.3 & 30.8 \\
Item B & 15.2 & 25.1 & 40.3 \\
Item C & 22.1 & 12.7 & 34.8 \\
\bottomrule
\end{tabular}
\caption{Sample Data Table}
\label{tab:sample}
\end{table}

\section{Lists}
\subsection{Bullet Points}
\begin{itemize}
  \item First item
  \item Second item
  \item Third item with sub-items:
    \begin{itemize}
      \item Sub-item 1
      \item Sub-item 2
    \end{itemize}
\end{itemize}

\subsection{Numbered List}
\begin{enumerate}
  \item First step
  \item Second step
  \item Third step
\end{enumerate}

\section{Citations and References}
When you use the citation search feature, references will be automatically managed. For example, you might cite a paper like this: \cite{example2024}. All citations will be automatically added to your \texttt{cite.bib} file.

\section{Conclusion}
This template demonstrates basic \LaTeX{} features including mathematical formulas, tables, and lists. The bibliography system is now set up to work seamlessly with Bibby AI's citation management. Feel free to modify it for your needs!

\bibliographystyle{plain}
\bibliography{cite}

\end{document}\documentclass{article}
\usepackage[utf8]{inputenc}
\usepackage{booktabs}
\usepackage{amsmath}
\usepackage{natbib}

\title{My Document}
\author{Your Name}
\date{\today}

\begin{document}

\maketitle

\section{Introduction}
This is a sample document demonstrating various \LaTeX{} features including mathematical equations, tables, and formatting. You can cite papers using the citation search feature, and they will be automatically added to your \texttt{cite.bib} file.

\section{Mathematical Expressions}
Here are some example equations:

\subsection{Quadratic Formula}
The quadratic formula is given by:
\[x = \frac{-b \pm \sqrt{b^2 - 4ac}}{2a}\]

\subsection{Matrix Example}
A sample matrix:
\[
\begin{pmatrix}
a_{11} & a_{12} & a_{13} \\
a_{21} & a_{22} & a_{23} \\
a_{31} & a_{32} & a_{33}
\end{pmatrix}
\]

\section{Tables}
Below is a sample table showing data:

\begin{table}[h]
\centering
\begin{tabular}{@{}lrrr@{}}
\toprule
\textbf{Category} & \textbf{Value 1} & \textbf{Value 2} & \textbf{Total} \\
\midrule
Item A & 10.5 & 20.3 & 30.8 \\
Item B & 15.2 & 25.1 & 40.3 \\
Item C & 22.1 & 12.7 & 34.8 \\
\bottomrule
\end{tabular}
\caption{Sample Data Table}
\label{tab:sample}
\end{table}

\section{Lists}
\subsection{Bullet Points}
\begin{itemize}
  \item First item
  \item Second item
  \item Third item with sub-items:
    \begin{itemize}
      \item Sub-item 1
      \item Sub-item 2
    \end{itemize}
\end{itemize}

\subsection{Numbered List}
\begin{enumerate}
  \item First step
  \item Second step
  \item Third step
\end{enumerate}

\section{Citations and References}
When you use the citation search feature, references will be automatically managed. For example, you might cite a paper like this: \cite{example2024}. All citations will be automatically added to your \texttt{cite.bib} file.

\section{Conclusion}
This template demonstrates basic \LaTeX{} features including mathematical formulas, tables, and lists. The bibliography system is now set up to work seamlessly with Bibby AI's citation management. Feel free to modify it for your needs!

\bibliographystyle{plain}
\bibliography{cite}

\end{document}\documentclass{article}
\usepackage[utf8]{inputenc}
\usepackage{booktabs}
\usepackage{amsmath}
\usepackage{natbib}

\title{My Document}
\author{Your Name}
\date{\today}

\begin{document}

\maketitle

\section{Introduction}
This is a sample document demonstrating various \LaTeX{} features including mathematical equations, tables, and formatting. You can cite papers using the citation search feature, and they will be automatically added to your \texttt{cite.bib} file.

\section{Mathematical Expressions}
Here are some example equations:

\subsection{Quadratic Formula}
The quadratic formula is given by:
\[x = \frac{-b \pm \sqrt{b^2 - 4ac}}{2a}\]

\subsection{Matrix Example}
A sample matrix:
\[
\begin{pmatrix}
a_{11} & a_{12} & a_{13} \\
a_{21} & a_{22} & a_{23} \\
a_{31} & a_{32} & a_{33}
\end{pmatrix}
\]

\section{Tables}
Below is a sample table showing data:

\begin{table}[h]
\centering
\begin{tabular}{@{}lrrr@{}}
\toprule
\textbf{Category} & \textbf{Value 1} & \textbf{Value 2} & \textbf{Total} \\
\midrule
Item A & 10.5 & 20.3 & 30.8 \\
Item B & 15.2 & 25.1 & 40.3 \\
Item C & 22.1 & 12.7 & 34.8 \\
\bottomrule
\end{tabular}
\caption{Sample Data Table}
\label{tab:sample}
\end{table}

\section{Lists}
\subsection{Bullet Points}
\begin{itemize}
  \item First item
  \item Second item
  \item Third item with sub-items:
    \begin{itemize}
      \item Sub-item 1
      \item Sub-item 2
    \end{itemize}
\end{itemize}

\subsection{Numbered List}
\begin{enumerate}
  \item First step
  \item Second step
  \item Third step
\end{enumerate}

\section{Citations and References}
When you use the citation search feature, references will be automatically managed. For example, you might cite a paper like this: \cite{example2024}. All citations will be automatically added to your \texttt{cite.bib} file.

\section{Conclusion}
This template demonstrates basic \LaTeX{} features including mathematical formulas, tables, and lists. The bibliography system is now set up to work seamlessly with Bibby AI's citation management. Feel free to modify it for your needs!

\bibliographystyle{plain}
\bibliography{cite}

\end{document}\documentclass{article}
\usepackage[utf8]{inputenc}
\usepackage{booktabs}
\usepackage{amsmath}
\usepackage{natbib}

\title{My Document}
\author{Your Name}
\date{\today}

\begin{document}

\maketitle

\section{Introduction}
This is a sample document demonstrating various \LaTeX{} features including mathematical equations, tables, and formatting. You can cite papers using the citation search feature, and they will be automatically added to your \texttt{cite.bib} file.

\section{Mathematical Expressions}
Here are some example equations:

\subsection{Quadratic Formula}
The quadratic formula is given by:
\[x = \frac{-b \pm \sqrt{b^2 - 4ac}}{2a}\]

\subsection{Matrix Example}
A sample matrix:
\[
\begin{pmatrix}
a_{11} & a_{12} & a_{13} \\
a_{21} & a_{22} & a_{23} \\
a_{31} & a_{32} & a_{33}
\end{pmatrix}
\]

\section{Tables}
Below is a sample table showing data:

\begin{table}[h]
\centering
\begin{tabular}{@{}lrrr@{}}
\toprule
\textbf{Category} & \textbf{Value 1} & \textbf{Value 2} & \textbf{Total} \\
\midrule
Item A & 10.5 & 20.3 & 30.8 \\
Item B & 15.2 & 25.1 & 40.3 \\
Item C & 22.1 & 12.7 & 34.8 \\
\bottomrule
\end{tabular}
\caption{Sample Data Table}
\label{tab:sample}
\end{table}

\section{Lists}
\subsection{Bullet Points}
\begin{itemize}
  \item First item
  \item Second item
  \item Third item with sub-items:
    \begin{itemize}
      \item Sub-item 1
      \item Sub-item 2
    \end{itemize}
\end{itemize}

\subsection{Numbered List}
\begin{enumerate}
  \item First step
  \item Second step
  \item Third step
\end{enumerate}

\section{Citations and References}
When you use the citation search feature, references will be automatically managed. For example, you might cite a paper like this: \cite{example2024}. All citations will be automatically added to your \texttt{cite.bib} file.

\section{Conclusion}
This template demonstrates basic \LaTeX{} features including mathematical formulas, tables, and lists. The bibliography system is now set up to work seamlessly with Bibby AI's citation management. Feel free to modify it for your needs!

\bibliographystyle{plain}
\bibliography{cite}

\end{document}\documentclass{article}
\usepackage[utf8]{inputenc}
\usepackage{booktabs}
\usepackage{amsmath}
\usepackage{natbib}

\title{My Document}
\author{Your Name}
\date{\today}

\begin{document}

\maketitle

\section{Introduction}
This is a sample document demonstrating various \LaTeX{} features including mathematical equations, tables, and formatting. You can cite papers using the citation search feature, and they will be automatically added to your \texttt{cite.bib} file.

\section{Mathematical Expressions}
Here are some example equations:

\subsection{Quadratic Formula}
The quadratic formula is given by:
\[x = \frac{-b \pm \sqrt{b^2 - 4ac}}{2a}\]

\subsection{Matrix Example}
A sample matrix:
\[
\begin{pmatrix}
a_{11} & a_{12} & a_{13} \\
a_{21} & a_{22} & a_{23} \\
a_{31} & a_{32} & a_{33}
\end{pmatrix}
\]

\section{Tables}
Below is a sample table showing data:

\begin{table}[h]
\centering
\begin{tabular}{@{}lrrr@{}}
\toprule
\textbf{Category} & \textbf{Value 1} & \textbf{Value 2} & \textbf{Total} \\
\midrule
Item A & 10.5 & 20.3 & 30.8 \\
Item B & 15.2 & 25.1 & 40.3 \\
Item C & 22.1 & 12.7 & 34.8 \\
\bottomrule
\end{tabular}
\caption{Sample Data Table}
\label{tab:sample}
\end{table}

\section{Lists}
\subsection{Bullet Points}
\begin{itemize}
  \item First item
  \item Second item
  \item Third item with sub-items:
    \begin{itemize}
      \item Sub-item 1
      \item Sub-item 2
    \end{itemize}
\end{itemize}

\subsection{Numbered List}
\begin{enumerate}
  \item First step
  \item Second step
  \item Third step
\end{enumerate}

\section{Citations and References}
When you use the citation search feature, references will be automatically managed. For example, you might cite a paper like this: \cite{example2024}. All citations will be automatically added to your \texttt{cite.bib} file.

\section{Conclusion}
This template demonstrates basic \LaTeX{} features including mathematical formulas, tables, and lists. The bibliography system is now set up to work seamlessly with Bibby AI's citation management. Feel free to modify it for your needs!

\bibliographystyle{plain}
\bibliography{cite}

\end{document}\documentclass{article}
\usepackage[utf8]{inputenc}
\usepackage{booktabs}
\usepackage{amsmath}
\usepackage{natbib}

\title{My Document}
\author{Your Name}
\date{\today}

\begin{document}

\maketitle

\section{Introduction}
This is a sample document demonstrating various \LaTeX{} features including mathematical equations, tables, and formatting. You can cite papers using the citation search feature, and they will be automatically added to your \texttt{cite.bib} file.

\section{Mathematical Expressions}
Here are some example equations:

\subsection{Quadratic Formula}
The quadratic formula is given by:
\[x = \frac{-b \pm \sqrt{b^2 - 4ac}}{2a}\]

\subsection{Matrix Example}
A sample matrix:
\[
\begin{pmatrix}
a_{11} & a_{12} & a_{13} \\
a_{21} & a_{22} & a_{23} \\
a_{31} & a_{32} & a_{33}
\end{pmatrix}
\]

\section{Tables}
Below is a sample table showing data:

\begin{table}[h]
\centering
\begin{tabular}{@{}lrrr@{}}
\toprule
\textbf{Category} & \textbf{Value 1} & \textbf{Value 2} & \textbf{Total} \\
\midrule
Item A & 10.5 & 20.3 & 30.8 \\
Item B & 15.2 & 25.1 & 40.3 \\
Item C & 22.1 & 12.7 & 34.8 \\
\bottomrule
\end{tabular}
\caption{Sample Data Table}
\label{tab:sample}
\end{table}

\section{Lists}
\subsection{Bullet Points}
\begin{itemize}
  \item First item
  \item Second item
  \item Third item with sub-items:
    \begin{itemize}
      \item Sub-item 1
      \item Sub-item 2
    \end{itemize}
\end{itemize}

\subsection{Numbered List}
\begin{enumerate}
  \item First step
  \item Second step
  \item Third step
\end{enumerate}

\section{Citations and References}
When you use the citation search feature, references will be automatically managed. For example, you might cite a paper like this: \cite{example2024}. All citations will be automatically added to your \texttt{cite.bib} file.

\section{Conclusion}
This template demonstrates basic \LaTeX{} features including mathematical formulas, tables, and lists. The bibliography system is now set up to work seamlessly with Bibby AI's citation management. Feel free to modify it for your needs!

\bibliographystyle{plain}
\bibliography{cite}

\end{document}\documentclass{article}
\usepackage[utf8]{inputenc}
\usepackage{booktabs}
\usepackage{amsmath}
\usepackage{natbib}

\title{My Document}
\author{Your Name}
\date{\today}

\begin{document}

\maketitle

\section{Introduction}
This is a sample document demonstrating various \LaTeX{} features including mathematical equations, tables, and formatting. You can cite papers using the citation search feature, and they will be automatically added to your \texttt{cite.bib} file.

\section{Mathematical Expressions}
Here are some example equations:

\subsection{Quadratic Formula}
The quadratic formula is given by:
\[x = \frac{-b \pm \sqrt{b^2 - 4ac}}{2a}\]

\subsection{Matrix Example}
A sample matrix:
\[
\begin{pmatrix}
a_{11} & a_{12} & a_{13} \\
a_{21} & a_{22} & a_{23} \\
a_{31} & a_{32} & a_{33}
\end{pmatrix}
\]

\section{Tables}
Below is a sample table showing data:

\begin{table}[h]
\centering
\begin{tabular}{@{}lrrr@{}}
\toprule
\textbf{Category} & \textbf{Value 1} & \textbf{Value 2} & \textbf{Total} \\
\midrule
Item A & 10.5 & 20.3 & 30.8 \\
Item B & 15.2 & 25.1 & 40.3 \\
Item C & 22.1 & 12.7 & 34.8 \\
\bottomrule
\end{tabular}
\caption{Sample Data Table}
\label{tab:sample}
\end{table}

\section{Lists}
\subsection{Bullet Points}
\begin{itemize}
  \item First item
  \item Second item
  \item Third item with sub-items:
    \begin{itemize}
      \item Sub-item 1
      \item Sub-item 2
    \end{itemize}
\end{itemize}

\subsection{Numbered List}
\begin{enumerate}
  \item First step
  \item Second step
  \item Third step
\end{enumerate}

\section{Citations and References}
When you use the citation search feature, references will be automatically managed. For example, you might cite a paper like this: \cite{example2024}. All citations will be automatically added to your \texttt{cite.bib} file.

\section{Conclusion}
This template demonstrates basic \LaTeX{} features including mathematical formulas, tables, and lists. The bibliography system is now set up to work seamlessly with Bibby AI's citation management. Feel free to modify it for your needs!

\bibliographystyle{plain}
\bibliography{cite}

\end{document}\documentclass{article}
\usepackage[utf8]{inputenc}
\usepackage{booktabs}
\usepackage{amsmath}
\usepackage{natbib}

\title{My Document}
\author{Your Name}
\date{\today}

\begin{document}

\maketitle

\section{Introduction}
This is a sample document demonstrating various \LaTeX{} features including mathematical equations, tables, and formatting. You can cite papers using the citation search feature, and they will be automatically added to your \texttt{cite.bib} file.

\section{Mathematical Expressions}
Here are some example equations:

\subsection{Quadratic Formula}
The quadratic formula is given by:
\[x = \frac{-b \pm \sqrt{b^2 - 4ac}}{2a}\]

\subsection{Matrix Example}
A sample matrix:
\[
\begin{pmatrix}
a_{11} & a_{12} & a_{13} \\
a_{21} & a_{22} & a_{23} \\
a_{31} & a_{32} & a_{33}
\end{pmatrix}
\]

\section{Tables}
Below is a sample table showing data:

\begin{table}[h]
\centering
\begin{tabular}{@{}lrrr@{}}
\toprule
\textbf{Category} & \textbf{Value 1} & \textbf{Value 2} & \textbf{Total} \\
\midrule
Item A & 10.5 & 20.3 & 30.8 \\
Item B & 15.2 & 25.1 & 40.3 \\
Item C & 22.1 & 12.7 & 34.8 \\
\bottomrule
\end{tabular}
\caption{Sample Data Table}
\label{tab:sample}
\end{table}

\section{Lists}
\subsection{Bullet Points}
\begin{itemize}
  \item First item
  \item Second item
  \item Third item with sub-items:
    \begin{itemize}
      \item Sub-item 1
      \item Sub-item 2
    \end{itemize}
\end{itemize}

\subsection{Numbered List}
\begin{enumerate}
  \item First step
  \item Second step
  \item Third step
\end{enumerate}

\section{Citations and References}
When you use the citation search feature, references will be automatically managed. For example, you might cite a paper like this: \cite{example2024}. All citations will be automatically added to your \texttt{cite.bib} file.

\section{Conclusion}
This template demonstrates basic \LaTeX{} features including mathematical formulas, tables, and lists. The bibliography system is now set up to work seamlessly with Bibby AI's citation management. Feel free to modify it for your needs!

\bibliographystyle{plain}
\bibliography{cite}

\end{document}\documentclass{article}
\usepackage[utf8]{inputenc}
\usepackage{booktabs}
\usepackage{amsmath}
\usepackage{natbib}

\title{My Document}
\author{Your Name}
\date{\today}

\begin{document}

\maketitle

\section{Introduction}
This is a sample document demonstrating various \LaTeX{} features including mathematical equations, tables, and formatting. You can cite papers using the citation search feature, and they will be automatically added to your \texttt{cite.bib} file.

\section{Mathematical Expressions}
Here are some example equations:

\subsection{Quadratic Formula}
The quadratic formula is given by:
\[x = \frac{-b \pm \sqrt{b^2 - 4ac}}{2a}\]

\subsection{Matrix Example}
A sample matrix:
\[
\begin{pmatrix}
a_{11} & a_{12} & a_{13} \\
a_{21} & a_{22} & a_{23} \\
a_{31} & a_{32} & a_{33}
\end{pmatrix}
\]

\section{Tables}
Below is a sample table showing data:

\begin{table}[h]
\centering
\begin{tabular}{@{}lrrr@{}}
\toprule
\textbf{Category} & \textbf{Value 1} & \textbf{Value 2} & \textbf{Total} \\
\midrule
Item A & 10.5 & 20.3 & 30.8 \\
Item B & 15.2 & 25.1 & 40.3 \\
Item C & 22.1 & 12.7 & 34.8 \\
\bottomrule
\end{tabular}
\caption{Sample Data Table}
\label{tab:sample}
\end{table}

\section{Lists}
\subsection{Bullet Points}
\begin{itemize}
  \item First item
  \item Second item
  \item Third item with sub-items:
    \begin{itemize}
      \item Sub-item 1
      \item Sub-item 2
    \end{itemize}
\end{itemize}

\subsection{Numbered List}
\begin{enumerate}
  \item First step
  \item Second step
  \item Third step
\end{enumerate}

\section{Citations and References}
When you use the citation search feature, references will be automatically managed. For example, you might cite a paper like this: \cite{example2024}. All citations will be automatically added to your \texttt{cite.bib} file.

\section{Conclusion}
This template demonstrates basic \LaTeX{} features including mathematical formulas, tables, and lists. The bibliography system is now set up to work seamlessly with Bibby AI's citation management. Feel free to modify it for your needs!

\bibliographystyle{plain}
\bibliography{cite}

\end{document}\documentclass{article}
\usepackage[utf8]{inputenc}
\usepackage{booktabs}
\usepackage{amsmath}
\usepackage{natbib}

\title{My Document}
\author{Your Name}
\date{\today}

\begin{document}

\maketitle

\section{Introduction}
This is a sample document demonstrating various \LaTeX{} features including mathematical equations, tables, and formatting. You can cite papers using the citation search feature, and they will be automatically added to your \texttt{cite.bib} file.

\section{Mathematical Expressions}
Here are some example equations:

\subsection{Quadratic Formula}
The quadratic formula is given by:
\[x = \frac{-b \pm \sqrt{b^2 - 4ac}}{2a}\]

\subsection{Matrix Example}
A sample matrix:
\[
\begin{pmatrix}
a_{11} & a_{12} & a_{13} \\
a_{21} & a_{22} & a_{23} \\
a_{31} & a_{32} & a_{33}
\end{pmatrix}
\]

\section{Tables}
Below is a sample table showing data:

\begin{table}[h]
\centering
\begin{tabular}{@{}lrrr@{}}
\toprule
\textbf{Category} & \textbf{Value 1} & \textbf{Value 2} & \textbf{Total} \\
\midrule
Item A & 10.5 & 20.3 & 30.8 \\
Item B & 15.2 & 25.1 & 40.3 \\
Item C & 22.1 & 12.7 & 34.8 \\
\bottomrule
\end{tabular}
\caption{Sample Data Table}
\label{tab:sample}
\end{table}

\section{Lists}
\subsection{Bullet Points}
\begin{itemize}
  \item First item
  \item Second item
  \item Third item with sub-items:
    \begin{itemize}
      \item Sub-item 1
      \item Sub-item 2
    \end{itemize}
\end{itemize}

\subsection{Numbered List}
\begin{enumerate}
  \item First step
  \item Second step
  \item Third step
\end{enumerate}

\section{Citations and References}
When you use the citation search feature, references will be automatically managed. For example, you might cite a paper like this: \cite{example2024}. All citations will be automatically added to your \texttt{cite.bib} file.

\section{Conclusion}
This template demonstrates basic \LaTeX{} features including mathematical formulas, tables, and lists. The bibliography system is now set up to work seamlessly with Bibby AI's citation management. Feel free to modify it for your needs!

\bibliographystyle{plain}
\bibliography{cite}

\end{document}\documentclass{article}
\usepackage[utf8]{inputenc}
\usepackage{booktabs}
\usepackage{amsmath}
\usepackage{natbib}

\title{My Document}
\author{Your Name}
\date{\today}

\begin{document}

\maketitle

\section{Introduction}
This is a sample document demonstrating various \LaTeX{} features including mathematical equations, tables, and formatting. You can cite papers using the citation search feature, and they will be automatically added to your \texttt{cite.bib} file.

\section{Mathematical Expressions}
Here are some example equations:

\subsection{Quadratic Formula}
The quadratic formula is given by:
\[x = \frac{-b \pm \sqrt{b^2 - 4ac}}{2a}\]

\subsection{Matrix Example}
A sample matrix:
\[
\begin{pmatrix}
a_{11} & a_{12} & a_{13} \\
a_{21} & a_{22} & a_{23} \\
a_{31} & a_{32} & a_{33}
\end{pmatrix}
\]

\section{Tables}
Below is a sample table showing data:

\begin{table}[h]
\centering
\begin{tabular}{@{}lrrr@{}}
\toprule
\textbf{Category} & \textbf{Value 1} & \textbf{Value 2} & \textbf{Total} \\
\midrule
Item A & 10.5 & 20.3 & 30.8 \\
Item B & 15.2 & 25.1 & 40.3 \\
Item C & 22.1 & 12.7 & 34.8 \\
\bottomrule
\end{tabular}
\caption{Sample Data Table}
\label{tab:sample}
\end{table}

\section{Lists}
\subsection{Bullet Points}
\begin{itemize}
  \item First item
  \item Second item
  \item Third item with sub-items:
    \begin{itemize}
      \item Sub-item 1
      \item Sub-item 2
    \end{itemize}
\end{itemize}

\subsection{Numbered List}
\begin{enumerate}
  \item First step
  \item Second step
  \item Third step
\end{enumerate}

\section{Citations and References}
When you use the citation search feature, references will be automatically managed. For example, you might cite a paper like this: \cite{example2024}. All citations will be automatically added to your \texttt{cite.bib} file.

\section{Conclusion}
This template demonstrates basic \LaTeX{} features including mathematical formulas, tables, and lists. The bibliography system is now set up to work seamlessly with Bibby AI's citation management. Feel free to modify it for your needs!

\bibliographystyle{plain}
\bibliography{cite}

\end{document}\documentclass{article}
\usepackage[utf8]{inputenc}
\usepackage{booktabs}
\usepackage{amsmath}
\usepackage{natbib}

\title{My Document}
\author{Your Name}
\date{\today}

\begin{document}

\maketitle

\section{Introduction}
This is a sample document demonstrating various \LaTeX{} features including mathematical equations, tables, and formatting. You can cite papers using the citation search feature, and they will be automatically added to your \texttt{cite.bib} file.

\section{Mathematical Expressions}
Here are some example equations:

\subsection{Quadratic Formula}
The quadratic formula is given by:
\[x = \frac{-b \pm \sqrt{b^2 - 4ac}}{2a}\]

\subsection{Matrix Example}
A sample matrix:
\[
\begin{pmatrix}
a_{11} & a_{12} & a_{13} \\
a_{21} & a_{22} & a_{23} \\
a_{31} & a_{32} & a_{33}
\end{pmatrix}
\]

\section{Tables}
Below is a sample table showing data:

\begin{table}[h]
\centering
\begin{tabular}{@{}lrrr@{}}
\toprule
\textbf{Category} & \textbf{Value 1} & \textbf{Value 2} & \textbf{Total} \\
\midrule
Item A & 10.5 & 20.3 & 30.8 \\
Item B & 15.2 & 25.1 & 40.3 \\
Item C & 22.1 & 12.7 & 34.8 \\
\bottomrule
\end{tabular}
\caption{Sample Data Table}
\label{tab:sample}
\end{table}

\section{Lists}
\subsection{Bullet Points}
\begin{itemize}
  \item First item
  \item Second item
  \item Third item with sub-items:
    \begin{itemize}
      \item Sub-item 1
      \item Sub-item 2
    \end{itemize}
\end{itemize}

\subsection{Numbered List}
\begin{enumerate}
  \item First step
  \item Second step
  \item Third step
\end{enumerate}

\section{Citations and References}
When you use the citation search feature, references will be automatically managed. For example, you might cite a paper like this: \cite{example2024}. All citations will be automatically added to your \texttt{cite.bib} file.

\section{Conclusion}
This template demonstrates basic \LaTeX{} features including mathematical formulas, tables, and lists. The bibliography system is now set up to work seamlessly with Bibby AI's citation management. Feel free to modify it for your needs!

\bibliographystyle{plain}
\bibliography{cite}

\end{document}\documentclass{article}
\usepackage[utf8]{inputenc}
\usepackage{booktabs}
\usepackage{amsmath}
\usepackage{natbib}

\title{My Document}
\author{Your Name}
\date{\today}

\begin{document}

\maketitle

\section{Introduction}
This is a sample document demonstrating various \LaTeX{} features including mathematical equations, tables, and formatting. You can cite papers using the citation search feature, and they will be automatically added to your \texttt{cite.bib} file.

\section{Mathematical Expressions}
Here are some example equations:

\subsection{Quadratic Formula}
The quadratic formula is given by:
\[x = \frac{-b \pm \sqrt{b^2 - 4ac}}{2a}\]

\subsection{Matrix Example}
A sample matrix:
\[
\begin{pmatrix}
a_{11} & a_{12} & a_{13} \\
a_{21} & a_{22} & a_{23} \\
a_{31} & a_{32} & a_{33}
\end{pmatrix}
\]

\section{Tables}
Below is a sample table showing data:

\begin{table}[h]
\centering
\begin{tabular}{@{}lrrr@{}}
\toprule
\textbf{Category} & \textbf{Value 1} & \textbf{Value 2} & \textbf{Total} \\
\midrule
Item A & 10.5 & 20.3 & 30.8 \\
Item B & 15.2 & 25.1 & 40.3 \\
Item C & 22.1 & 12.7 & 34.8 \\
\bottomrule
\end{tabular}
\caption{Sample Data Table}
\label{tab:sample}
\end{table}

\section{Lists}
\subsection{Bullet Points}
\begin{itemize}
  \item First item
  \item Second item
  \item Third item with sub-items:
    \begin{itemize}
      \item Sub-item 1
      \item Sub-item 2
    \end{itemize}
\end{itemize}

\subsection{Numbered List}
\begin{enumerate}
  \item First step
  \item Second step
  \item Third step
\end{enumerate}

\section{Citations and References}
When you use the citation search feature, references will be automatically managed. For example, you might cite a paper like this: \cite{example2024}. All citations will be automatically added to your \texttt{cite.bib} file.

\section{Conclusion}
This template demonstrates basic \LaTeX{} features including mathematical formulas, tables, and lists. The bibliography system is now set up to work seamlessly with Bibby AI's citation management. Feel free to modify it for your needs!

\bibliographystyle{plain}
\bibliography{cite}

\end{document}\documentclass{article}
\usepackage[utf8]{inputenc}
\usepackage{booktabs}
\usepackage{amsmath}
\usepackage{natbib}

\title{My Document}
\author{Your Name}
\date{\today}

\begin{document}

\maketitle

\section{Introduction}
This is a sample document demonstrating various \LaTeX{} features including mathematical equations, tables, and formatting. You can cite papers using the citation search feature, and they will be automatically added to your \texttt{cite.bib} file.

\section{Mathematical Expressions}
Here are some example equations:

\subsection{Quadratic Formula}
The quadratic formula is given by:
\[x = \frac{-b \pm \sqrt{b^2 - 4ac}}{2a}\]

\subsection{Matrix Example}
A sample matrix:
\[
\begin{pmatrix}
a_{11} & a_{12} & a_{13} \\
a_{21} & a_{22} & a_{23} \\
a_{31} & a_{32} & a_{33}
\end{pmatrix}
\]

\section{Tables}
Below is a sample table showing data:

\begin{table}[h]
\centering
\begin{tabular}{@{}lrrr@{}}
\toprule
\textbf{Category} & \textbf{Value 1} & \textbf{Value 2} & \textbf{Total} \\
\midrule
Item A & 10.5 & 20.3 & 30.8 \\
Item B & 15.2 & 25.1 & 40.3 \\
Item C & 22.1 & 12.7 & 34.8 \\
\bottomrule
\end{tabular}
\caption{Sample Data Table}
\label{tab:sample}
\end{table}

\section{Lists}
\subsection{Bullet Points}
\begin{itemize}
  \item First item
  \item Second item
  \item Third item with sub-items:
    \begin{itemize}
      \item Sub-item 1
      \item Sub-item 2
    \end{itemize}
\end{itemize}

\subsection{Numbered List}
\begin{enumerate}
  \item First step
  \item Second step
  \item Third step
\end{enumerate}

\section{Citations and References}
When you use the citation search feature, references will be automatically managed. For example, you might cite a paper like this: \cite{example2024}. All citations will be automatically added to your \texttt{cite.bib} file.

\section{Conclusion}
This template demonstrates basic \LaTeX{} features including mathematical formulas, tables, and lists. The bibliography system is now set up to work seamlessly with Bibby AI's citation management. Feel free to modify it for your needs!

\bibliographystyle{plain}
\bibliography{cite}

\end{document}\documentclass{article}
\usepackage[utf8]{inputenc}
\usepackage{booktabs}
\usepackage{amsmath}
\usepackage{natbib}

\title{My Document}
\author{Your Name}
\date{\today}

\begin{document}

\maketitle

\section{Introduction}
This is a sample document demonstrating various \LaTeX{} features including mathematical equations, tables, and formatting. You can cite papers using the citation search feature, and they will be automatically added to your \texttt{cite.bib} file.

\section{Mathematical Expressions}
Here are some example equations:

\subsection{Quadratic Formula}
The quadratic formula is given by:
\[x = \frac{-b \pm \sqrt{b^2 - 4ac}}{2a}\]

\subsection{Matrix Example}
A sample matrix:
\[
\begin{pmatrix}
a_{11} & a_{12} & a_{13} \\
a_{21} & a_{22} & a_{23} \\
a_{31} & a_{32} & a_{33}
\end{pmatrix}
\]

\section{Tables}
Below is a sample table showing data:

\begin{table}[h]
\centering
\begin{tabular}{@{}lrrr@{}}
\toprule
\textbf{Category} & \textbf{Value 1} & \textbf{Value 2} & \textbf{Total} \\
\midrule
Item A & 10.5 & 20.3 & 30.8 \\
Item B & 15.2 & 25.1 & 40.3 \\
Item C & 22.1 & 12.7 & 34.8 \\
\bottomrule
\end{tabular}
\caption{Sample Data Table}
\label{tab:sample}
\end{table}

\section{Lists}
\subsection{Bullet Points}
\begin{itemize}
  \item First item
  \item Second item
  \item Third item with sub-items:
    \begin{itemize}
      \item Sub-item 1
      \item Sub-item 2
    \end{itemize}
\end{itemize}

\subsection{Numbered List}
\begin{enumerate}
  \item First step
  \item Second step
  \item Third step
\end{enumerate}

\section{Citations and References}
When you use the citation search feature, references will be automatically managed. For example, you might cite a paper like this: \cite{example2024}. All citations will be automatically added to your \texttt{cite.bib} file.

\section{Conclusion}
This template demonstrates basic \LaTeX{} features including mathematical formulas, tables, and lists. The bibliography system is now set up to work seamlessly with Bibby AI's citation management. Feel free to modify it for your needs!

\bibliographystyle{plain}
\bibliography{cite}

\end{document}\documentclass{article}
\usepackage[utf8]{inputenc}
\usepackage{booktabs}
\usepackage{amsmath}
\usepackage{natbib}

\title{My Document}
\author{Your Name}
\date{\today}

\begin{document}

\maketitle

\section{Introduction}
This is a sample document demonstrating various \LaTeX{} features including mathematical equations, tables, and formatting. You can cite papers using the citation search feature, and they will be automatically added to your \texttt{cite.bib} file.

\section{Mathematical Expressions}
Here are some example equations:

\subsection{Quadratic Formula}
The quadratic formula is given by:
\[x = \frac{-b \pm \sqrt{b^2 - 4ac}}{2a}\]

\subsection{Matrix Example}
A sample matrix:
\[
\begin{pmatrix}
a_{11} & a_{12} & a_{13} \\
a_{21} & a_{22} & a_{23} \\
a_{31} & a_{32} & a_{33}
\end{pmatrix}
\]

\section{Tables}
Below is a sample table showing data:

\begin{table}[h]
\centering
\begin{tabular}{@{}lrrr@{}}
\toprule
\textbf{Category} & \textbf{Value 1} & \textbf{Value 2} & \textbf{Total} \\
\midrule
Item A & 10.5 & 20.3 & 30.8 \\
Item B & 15.2 & 25.1 & 40.3 \\
Item C & 22.1 & 12.7 & 34.8 \\
\bottomrule
\end{tabular}
\caption{Sample Data Table}
\label{tab:sample}
\end{table}

\section{Lists}
\subsection{Bullet Points}
\begin{itemize}
  \item First item
  \item Second item
  \item Third item with sub-items:
    \begin{itemize}
      \item Sub-item 1
      \item Sub-item 2
    \end{itemize}
\end{itemize}

\subsection{Numbered List}
\begin{enumerate}
  \item First step
  \item Second step
  \item Third step
\end{enumerate}

\section{Citations and References}
When you use the citation search feature, references will be automatically managed. For example, you might cite a paper like this: \cite{example2024}. All citations will be automatically added to your \texttt{cite.bib} file.

\section{Conclusion}
This template demonstrates basic \LaTeX{} features including mathematical formulas, tables, and lists. The bibliography system is now set up to work seamlessly with Bibby AI's citation management. Feel free to modify it for your needs!

\bibliographystyle{plain}
\bibliography{cite}

\end{document}\documentclass{article}
\usepackage[utf8]{inputenc}
\usepackage{booktabs}
\usepackage{amsmath}
\usepackage{natbib}

\title{My Document}
\author{Your Name}
\date{\today}

\begin{document}

\maketitle

\section{Introduction}
This is a sample document demonstrating various \LaTeX{} features including mathematical equations, tables, and formatting. You can cite papers using the citation search feature, and they will be automatically added to your \texttt{cite.bib} file.

\section{Mathematical Expressions}
Here are some example equations:

\subsection{Quadratic Formula}
The quadratic formula is given by:
\[x = \frac{-b \pm \sqrt{b^2 - 4ac}}{2a}\]

\subsection{Matrix Example}
A sample matrix:
\[
\begin{pmatrix}
a_{11} & a_{12} & a_{13} \\
a_{21} & a_{22} & a_{23} \\
a_{31} & a_{32} & a_{33}
\end{pmatrix}
\]

\section{Tables}
Below is a sample table showing data:

\begin{table}[h]
\centering
\begin{tabular}{@{}lrrr@{}}
\toprule
\textbf{Category} & \textbf{Value 1} & \textbf{Value 2} & \textbf{Total} \\
\midrule
Item A & 10.5 & 20.3 & 30.8 \\
Item B & 15.2 & 25.1 & 40.3 \\
Item C & 22.1 & 12.7 & 34.8 \\
\bottomrule
\end{tabular}
\caption{Sample Data Table}
\label{tab:sample}
\end{table}

\section{Lists}
\subsection{Bullet Points}
\begin{itemize}
  \item First item
  \item Second item
  \item Third item with sub-items:
    \begin{itemize}
      \item Sub-item 1
      \item Sub-item 2
    \end{itemize}
\end{itemize}

\subsection{Numbered List}
\begin{enumerate}
  \item First step
  \item Second step
  \item Third step
\end{enumerate}

\section{Citations and References}
When you use the citation search feature, references will be automatically managed. For example, you might cite a paper like this: \cite{example2024}. All citations will be automatically added to your \texttt{cite.bib} file.

\section{Conclusion}
This template demonstrates basic \LaTeX{} features including mathematical formulas, tables, and lists. The bibliography system is now set up to work seamlessly with Bibby AI's citation management. Feel free to modify it for your needs!

\bibliographystyle{plain}
\bibliography{cite}

\end{document}\documentclass{article}
\usepackage[utf8]{inputenc}
\usepackage{booktabs}
\usepackage{amsmath}
\usepackage{natbib}

\title{My Document}
\author{Your Name}
\date{\today}

\begin{document}

\maketitle

\section{Introduction}
This is a sample document demonstrating various \LaTeX{} features including mathematical equations, tables, and formatting. You can cite papers using the citation search feature, and they will be automatically added to your \texttt{cite.bib} file.

\section{Mathematical Expressions}
Here are some example equations:

\subsection{Quadratic Formula}
The quadratic formula is given by:
\[x = \frac{-b \pm \sqrt{b^2 - 4ac}}{2a}\]

\subsection{Matrix Example}
A sample matrix:
\[
\begin{pmatrix}
a_{11} & a_{12} & a_{13} \\
a_{21} & a_{22} & a_{23} \\
a_{31} & a_{32} & a_{33}
\end{pmatrix}
\]

\section{Tables}
Below is a sample table showing data:

\begin{table}[h]
\centering
\begin{tabular}{@{}lrrr@{}}
\toprule
\textbf{Category} & \textbf{Value 1} & \textbf{Value 2} & \textbf{Total} \\
\midrule
Item A & 10.5 & 20.3 & 30.8 \\
Item B & 15.2 & 25.1 & 40.3 \\
Item C & 22.1 & 12.7 & 34.8 \\
\bottomrule
\end{tabular}
\caption{Sample Data Table}
\label{tab:sample}
\end{table}

\section{Lists}
\subsection{Bullet Points}
\begin{itemize}
  \item First item
  \item Second item
  \item Third item with sub-items:
    \begin{itemize}
      \item Sub-item 1
      \item Sub-item 2
    \end{itemize}
\end{itemize}

\subsection{Numbered List}
\begin{enumerate}
  \item First step
  \item Second step
  \item Third step
\end{enumerate}

\section{Citations and References}
When you use the citation search feature, references will be automatically managed. For example, you might cite a paper like this: \cite{example2024}. All citations will be automatically added to your \texttt{cite.bib} file.

\section{Conclusion}
This template demonstrates basic \LaTeX{} features including mathematical formulas, tables, and lists. The bibliography system is now set up to work seamlessly with Bibby AI's citation management. Feel free to modify it for your needs!

\bibliographystyle{plain}
\bibliography{cite}

\end{document}\documentclass{article}
\usepackage[utf8]{inputenc}
\usepackage{booktabs}
\usepackage{amsmath}
\usepackage{natbib}

\title{My Document}
\author{Your Name}
\date{\today}

\begin{document}

\maketitle

\section{Introduction}
This is a sample document demonstrating various \LaTeX{} features including mathematical equations, tables, and formatting. You can cite papers using the citation search feature, and they will be automatically added to your \texttt{cite.bib} file.

\section{Mathematical Expressions}
Here are some example equations:

\subsection{Quadratic Formula}
The quadratic formula is given by:
\[x = \frac{-b \pm \sqrt{b^2 - 4ac}}{2a}\]

\subsection{Matrix Example}
A sample matrix:
\[
\begin{pmatrix}
a_{11} & a_{12} & a_{13} \\
a_{21} & a_{22} & a_{23} \\
a_{31} & a_{32} & a_{33}
\end{pmatrix}
\]

\section{Tables}
Below is a sample table showing data:

\begin{table}[h]
\centering
\begin{tabular}{@{}lrrr@{}}
\toprule
\textbf{Category} & \textbf{Value 1} & \textbf{Value 2} & \textbf{Total} \\
\midrule
Item A & 10.5 & 20.3 & 30.8 \\
Item B & 15.2 & 25.1 & 40.3 \\
Item C & 22.1 & 12.7 & 34.8 \\
\bottomrule
\end{tabular}
\caption{Sample Data Table}
\label{tab:sample}
\end{table}

\section{Lists}
\subsection{Bullet Points}
\begin{itemize}
  \item First item
  \item Second item
  \item Third item with sub-items:
    \begin{itemize}
      \item Sub-item 1
      \item Sub-item 2
    \end{itemize}
\end{itemize}

\subsection{Numbered List}
\begin{enumerate}
  \item First step
  \item Second step
  \item Third step
\end{enumerate}

\section{Citations and References}
When you use the citation search feature, references will be automatically managed. For example, you might cite a paper like this: \cite{example2024}. All citations will be automatically added to your \texttt{cite.bib} file.

\section{Conclusion}
This template demonstrates basic \LaTeX{} features including mathematical formulas, tables, and lists. The bibliography system is now set up to work seamlessly with Bibby AI's citation management. Feel free to modify it for your needs!

\bibliographystyle{plain}
\bibliography{cite}

\end{document}\documentclass{article}
\usepackage[utf8]{inputenc}
\usepackage{booktabs}
\usepackage{amsmath}
\usepackage{natbib}

\title{My Document}
\author{Your Name}
\date{\today}

\begin{document}

\maketitle

\section{Introduction}
This is a sample document demonstrating various \LaTeX{} features including mathematical equations, tables, and formatting. You can cite papers using the citation search feature, and they will be automatically added to your \texttt{cite.bib} file.

\section{Mathematical Expressions}
Here are some example equations:

\subsection{Quadratic Formula}
The quadratic formula is given by:
\[x = \frac{-b \pm \sqrt{b^2 - 4ac}}{2a}\]

\subsection{Matrix Example}
A sample matrix:
\[
\begin{pmatrix}
a_{11} & a_{12} & a_{13} \\
a_{21} & a_{22} & a_{23} \\
a_{31} & a_{32} & a_{33}
\end{pmatrix}
\]

\section{Tables}
Below is a sample table showing data:

\begin{table}[h]
\centering
\begin{tabular}{@{}lrrr@{}}
\toprule
\textbf{Category} & \textbf{Value 1} & \textbf{Value 2} & \textbf{Total} \\
\midrule
Item A & 10.5 & 20.3 & 30.8 \\
Item B & 15.2 & 25.1 & 40.3 \\
Item C & 22.1 & 12.7 & 34.8 \\
\bottomrule
\end{tabular}
\caption{Sample Data Table}
\label{tab:sample}
\end{table}

\section{Lists}
\subsection{Bullet Points}
\begin{itemize}
  \item First item
  \item Second item
  \item Third item with sub-items:
    \begin{itemize}
      \item Sub-item 1
      \item Sub-item 2
    \end{itemize}
\end{itemize}

\subsection{Numbered List}
\begin{enumerate}
  \item First step
  \item Second step
  \item Third step
\end{enumerate}

\section{Citations and References}
When you use the citation search feature, references will be automatically managed. For example, you might cite a paper like this: \cite{example2024}. All citations will be automatically added to your \texttt{cite.bib} file.

\section{Conclusion}
This template demonstrates basic \LaTeX{} features including mathematical formulas, tables, and lists. The bibliography system is now set up to work seamlessly with Bibby AI's citation management. Feel free to modify it for your needs!

\bibliographystyle{plain}
\bibliography{cite}

\end{document}\documentclass{article}
\usepackage[utf8]{inputenc}
\usepackage{booktabs}
\usepackage{amsmath}
\usepackage{natbib}

\title{My Document}
\author{Your Name}
\date{\today}

\begin{document}

\maketitle

\section{Introduction}
This is a sample document demonstrating various \LaTeX{} features including mathematical equations, tables, and formatting. You can cite papers using the citation search feature, and they will be automatically added to your \texttt{cite.bib} file.

\section{Mathematical Expressions}
Here are some example equations:

\subsection{Quadratic Formula}
The quadratic formula is given by:
\[x = \frac{-b \pm \sqrt{b^2 - 4ac}}{2a}\]

\subsection{Matrix Example}
A sample matrix:
\[
\begin{pmatrix}
a_{11} & a_{12} & a_{13} \\
a_{21} & a_{22} & a_{23} \\
a_{31} & a_{32} & a_{33}
\end{pmatrix}
\]

\section{Tables}
Below is a sample table showing data:

\begin{table}[h]
\centering
\begin{tabular}{@{}lrrr@{}}
\toprule
\textbf{Category} & \textbf{Value 1} & \textbf{Value 2} & \textbf{Total} \\
\midrule
Item A & 10.5 & 20.3 & 30.8 \\
Item B & 15.2 & 25.1 & 40.3 \\
Item C & 22.1 & 12.7 & 34.8 \\
\bottomrule
\end{tabular}
\caption{Sample Data Table}
\label{tab:sample}
\end{table}

\section{Lists}
\subsection{Bullet Points}
\begin{itemize}
  \item First item
  \item Second item
  \item Third item with sub-items:
    \begin{itemize}
      \item Sub-item 1
      \item Sub-item 2
    \end{itemize}
\end{itemize}

\subsection{Numbered List}
\begin{enumerate}
  \item First step
  \item Second step
  \item Third step
\end{enumerate}

\section{Citations and References}
When you use the citation search feature, references will be automatically managed. For example, you might cite a paper like this: \cite{example2024}. All citations will be automatically added to your \texttt{cite.bib} file.

\section{Conclusion}
This template demonstrates basic \LaTeX{} features including mathematical formulas, tables, and lists. The bibliography system is now set up to work seamlessly with Bibby AI's citation management. Feel free to modify it for your needs!

\bibliographystyle{plain}
\bibliography{cite}

\end{document}\documentclass{article}
\usepackage[utf8]{inputenc}
\usepackage{booktabs}
\usepackage{amsmath}
\usepackage{natbib}

\title{My Document}
\author{Your Name}
\date{\today}

\begin{document}

\maketitle

\section{Introduction}
This is a sample document demonstrating various \LaTeX{} features including mathematical equations, tables, and formatting. You can cite papers using the citation search feature, and they will be automatically added to your \texttt{cite.bib} file.

\section{Mathematical Expressions}
Here are some example equations:

\subsection{Quadratic Formula}
The quadratic formula is given by:
\[x = \frac{-b \pm \sqrt{b^2 - 4ac}}{2a}\]

\subsection{Matrix Example}
A sample matrix:
\[
\begin{pmatrix}
a_{11} & a_{12} & a_{13} \\
a_{21} & a_{22} & a_{23} \\
a_{31} & a_{32} & a_{33}
\end{pmatrix}
\]

\section{Tables}
Below is a sample table showing data:

\begin{table}[h]
\centering
\begin{tabular}{@{}lrrr@{}}
\toprule
\textbf{Category} & \textbf{Value 1} & \textbf{Value 2} & \textbf{Total} \\
\midrule
Item A & 10.5 & 20.3 & 30.8 \\
Item B & 15.2 & 25.1 & 40.3 \\
Item C & 22.1 & 12.7 & 34.8 \\
\bottomrule
\end{tabular}
\caption{Sample Data Table}
\label{tab:sample}
\end{table}

\section{Lists}
\subsection{Bullet Points}
\begin{itemize}
  \item First item
  \item Second item
  \item Third item with sub-items:
    \begin{itemize}
      \item Sub-item 1
      \item Sub-item 2
    \end{itemize}
\end{itemize}

\subsection{Numbered List}
\begin{enumerate}
  \item First step
  \item Second step
  \item Third step
\end{enumerate}

\section{Citations and References}
When you use the citation search feature, references will be automatically managed. For example, you might cite a paper like this: \cite{example2024}. All citations will be automatically added to your \texttt{cite.bib} file.

\section{Conclusion}
This template demonstrates basic \LaTeX{} features including mathematical formulas, tables, and lists. The bibliography system is now set up to work seamlessly with Bibby AI's citation management. Feel free to modify it for your needs!

\bibliographystyle{plain}
\bibliography{cite}

\end{document}\documentclass{article}
\usepackage[utf8]{inputenc}
\usepackage{booktabs}
\usepackage{amsmath}
\usepackage{natbib}

\title{My Document}
\author{Your Name}
\date{\today}

\begin{document}

\maketitle

\section{Introduction}
This is a sample document demonstrating various \LaTeX{} features including mathematical equations, tables, and formatting. You can cite papers using the citation search feature, and they will be automatically added to your \texttt{cite.bib} file.

\section{Mathematical Expressions}
Here are some example equations:

\subsection{Quadratic Formula}
The quadratic formula is given by:
\[x = \frac{-b \pm \sqrt{b^2 - 4ac}}{2a}\]

\subsection{Matrix Example}
A sample matrix:
\[
\begin{pmatrix}
a_{11} & a_{12} & a_{13} \\
a_{21} & a_{22} & a_{23} \\
a_{31} & a_{32} & a_{33}
\end{pmatrix}
\]

\section{Tables}
Below is a sample table showing data:

\begin{table}[h]
\centering
\begin{tabular}{@{}lrrr@{}}
\toprule
\textbf{Category} & \textbf{Value 1} & \textbf{Value 2} & \textbf{Total} \\
\midrule
Item A & 10.5 & 20.3 & 30.8 \\
Item B & 15.2 & 25.1 & 40.3 \\
Item C & 22.1 & 12.7 & 34.8 \\
\bottomrule
\end{tabular}
\caption{Sample Data Table}
\label{tab:sample}
\end{table}

\section{Lists}
\subsection{Bullet Points}
\begin{itemize}
  \item First item
  \item Second item
  \item Third item with sub-items:
    \begin{itemize}
      \item Sub-item 1
      \item Sub-item 2
    \end{itemize}
\end{itemize}

\subsection{Numbered List}
\begin{enumerate}
  \item First step
  \item Second step
  \item Third step
\end{enumerate}

\section{Citations and References}
When you use the citation search feature, references will be automatically managed. For example, you might cite a paper like this: \cite{example2024}. All citations will be automatically added to your \texttt{cite.bib} file.

\section{Conclusion}
This template demonstrates basic \LaTeX{} features including mathematical formulas, tables, and lists. The bibliography system is now set up to work seamlessly with Bibby AI's citation management. Feel free to modify it for your needs!

\bibliographystyle{plain}
\bibliography{cite}

\end{document}\documentclass{article}
\usepackage[utf8]{inputenc}
\usepackage{booktabs}
\usepackage{amsmath}
\usepackage{natbib}

\title{My Document}
\author{Your Name}
\date{\today}

\begin{document}

\maketitle

\section{Introduction}
This is a sample document demonstrating various \LaTeX{} features including mathematical equations, tables, and formatting. You can cite papers using the citation search feature, and they will be automatically added to your \texttt{cite.bib} file.

\section{Mathematical Expressions}
Here are some example equations:

\subsection{Quadratic Formula}
The quadratic formula is given by:
\[x = \frac{-b \pm \sqrt{b^2 - 4ac}}{2a}\]

\subsection{Matrix Example}
A sample matrix:
\[
\begin{pmatrix}
a_{11} & a_{12} & a_{13} \\
a_{21} & a_{22} & a_{23} \\
a_{31} & a_{32} & a_{33}
\end{pmatrix}
\]

\section{Tables}
Below is a sample table showing data:

\begin{table}[h]
\centering
\begin{tabular}{@{}lrrr@{}}
\toprule
\textbf{Category} & \textbf{Value 1} & \textbf{Value 2} & \textbf{Total} \\
\midrule
Item A & 10.5 & 20.3 & 30.8 \\
Item B & 15.2 & 25.1 & 40.3 \\
Item C & 22.1 & 12.7 & 34.8 \\
\bottomrule
\end{tabular}
\caption{Sample Data Table}
\label{tab:sample}
\end{table}

\section{Lists}
\subsection{Bullet Points}
\begin{itemize}
  \item First item
  \item Second item
  \item Third item with sub-items:
    \begin{itemize}
      \item Sub-item 1
      \item Sub-item 2
    \end{itemize}
\end{itemize}

\subsection{Numbered List}
\begin{enumerate}
  \item First step
  \item Second step
  \item Third step
\end{enumerate}

\section{Citations and References}
When you use the citation search feature, references will be automatically managed. For example, you might cite a paper like this: \cite{example2024}. All citations will be automatically added to your \texttt{cite.bib} file.

\section{Conclusion}
This template demonstrates basic \LaTeX{} features including mathematical formulas, tables, and lists. The bibliography system is now set up to work seamlessly with Bibby AI's citation management. Feel free to modify it for your needs!

\bibliographystyle{plain}
\bibliography{cite}

\end{document}\documentclass{article}
\usepackage[utf8]{inputenc}
\usepackage{booktabs}
\usepackage{amsmath}
\usepackage{natbib}

\title{My Document}
\author{Your Name}
\date{\today}

\begin{document}

\maketitle

\section{Introduction}
This is a sample document demonstrating various \LaTeX{} features including mathematical equations, tables, and formatting. You can cite papers using the citation search feature, and they will be automatically added to your \texttt{cite.bib} file.

\section{Mathematical Expressions}
Here are some example equations:

\subsection{Quadratic Formula}
The quadratic formula is given by:
\[x = \frac{-b \pm \sqrt{b^2 - 4ac}}{2a}\]

\subsection{Matrix Example}
A sample matrix:
\[
\begin{pmatrix}
a_{11} & a_{12} & a_{13} \\
a_{21} & a_{22} & a_{23} \\
a_{31} & a_{32} & a_{33}
\end{pmatrix}
\]

\section{Tables}
Below is a sample table showing data:

\begin{table}[h]
\centering
\begin{tabular}{@{}lrrr@{}}
\toprule
\textbf{Category} & \textbf{Value 1} & \textbf{Value 2} & \textbf{Total} \\
\midrule
Item A & 10.5 & 20.3 & 30.8 \\
Item B & 15.2 & 25.1 & 40.3 \\
Item C & 22.1 & 12.7 & 34.8 \\
\bottomrule
\end{tabular}
\caption{Sample Data Table}
\label{tab:sample}
\end{table}

\section{Lists}
\subsection{Bullet Points}
\begin{itemize}
  \item First item
  \item Second item
  \item Third item with sub-items:
    \begin{itemize}
      \item Sub-item 1
      \item Sub-item 2
    \end{itemize}
\end{itemize}

\subsection{Numbered List}
\begin{enumerate}
  \item First step
  \item Second step
  \item Third step
\end{enumerate}

\section{Citations and References}
When you use the citation search feature, references will be automatically managed. For example, you might cite a paper like this: \cite{example2024}. All citations will be automatically added to your \texttt{cite.bib} file.

\section{Conclusion}
This template demonstrates basic \LaTeX{} features including mathematical formulas, tables, and lists. The bibliography system is now set up to work seamlessly with Bibby AI's citation management. Feel free to modify it for your needs!

\bibliographystyle{plain}
\bibliography{cite}

\end{document}\documentclass{article}
\usepackage[utf8]{inputenc}
\usepackage{booktabs}
\usepackage{amsmath}
\usepackage{natbib}

\title{My Document}
\author{Your Name}
\date{\today}

\begin{document}

\maketitle

\section{Introduction}
This is a sample document demonstrating various \LaTeX{} features including mathematical equations, tables, and formatting. You can cite papers using the citation search feature, and they will be automatically added to your \texttt{cite.bib} file.

\section{Mathematical Expressions}
Here are some example equations:

\subsection{Quadratic Formula}
The quadratic formula is given by:
\[x = \frac{-b \pm \sqrt{b^2 - 4ac}}{2a}\]

\subsection{Matrix Example}
A sample matrix:
\[
\begin{pmatrix}
a_{11} & a_{12} & a_{13} \\
a_{21} & a_{22} & a_{23} \\
a_{31} & a_{32} & a_{33}
\end{pmatrix}
\]

\section{Tables}
Below is a sample table showing data:

\begin{table}[h]
\centering
\begin{tabular}{@{}lrrr@{}}
\toprule
\textbf{Category} & \textbf{Value 1} & \textbf{Value 2} & \textbf{Total} \\
\midrule
Item A & 10.5 & 20.3 & 30.8 \\
Item B & 15.2 & 25.1 & 40.3 \\
Item C & 22.1 & 12.7 & 34.8 \\
\bottomrule
\end{tabular}
\caption{Sample Data Table}
\label{tab:sample}
\end{table}

\section{Lists}
\subsection{Bullet Points}
\begin{itemize}
  \item First item
  \item Second item
  \item Third item with sub-items:
    \begin{itemize}
      \item Sub-item 1
      \item Sub-item 2
    \end{itemize}
\end{itemize}

\subsection{Numbered List}
\begin{enumerate}
  \item First step
  \item Second step
  \item Third step
\end{enumerate}

\section{Citations and References}
When you use the citation search feature, references will be automatically managed. For example, you might cite a paper like this: \cite{example2024}. All citations will be automatically added to your \texttt{cite.bib} file.

\section{Conclusion}
This template demonstrates basic \LaTeX{} features including mathematical formulas, tables, and lists. The bibliography system is now set up to work seamlessly with Bibby AI's citation management. Feel free to modify it for your needs!

\bibliographystyle{plain}
\bibliography{cite}

\end{document}\documentclass{article}
\usepackage[utf8]{inputenc}
\usepackage{booktabs}
\usepackage{amsmath}
\usepackage{natbib}

\title{My Document}
\author{Your Name}
\date{\today}

\begin{document}

\maketitle

\section{Introduction}
This is a sample document demonstrating various \LaTeX{} features including mathematical equations, tables, and formatting. You can cite papers using the citation search feature, and they will be automatically added to your \texttt{cite.bib} file.

\section{Mathematical Expressions}
Here are some example equations:

\subsection{Quadratic Formula}
The quadratic formula is given by:
\[x = \frac{-b \pm \sqrt{b^2 - 4ac}}{2a}\]

\subsection{Matrix Example}
A sample matrix:
\[
\begin{pmatrix}
a_{11} & a_{12} & a_{13} \\
a_{21} & a_{22} & a_{23} \\
a_{31} & a_{32} & a_{33}
\end{pmatrix}
\]

\section{Tables}
Below is a sample table showing data:

\begin{table}[h]
\centering
\begin{tabular}{@{}lrrr@{}}
\toprule
\textbf{Category} & \textbf{Value 1} & \textbf{Value 2} & \textbf{Total} \\
\midrule
Item A & 10.5 & 20.3 & 30.8 \\
Item B & 15.2 & 25.1 & 40.3 \\
Item C & 22.1 & 12.7 & 34.8 \\
\bottomrule
\end{tabular}
\caption{Sample Data Table}
\label{tab:sample}
\end{table}

\section{Lists}
\subsection{Bullet Points}
\begin{itemize}
  \item First item
  \item Second item
  \item Third item with sub-items:
    \begin{itemize}
      \item Sub-item 1
      \item Sub-item 2
    \end{itemize}
\end{itemize}

\subsection{Numbered List}
\begin{enumerate}
  \item First step
  \item Second step
  \item Third step
\end{enumerate}

\section{Citations and References}
When you use the citation search feature, references will be automatically managed. For example, you might cite a paper like this: \cite{example2024}. All citations will be automatically added to your \texttt{cite.bib} file.

\section{Conclusion}
This template demonstrates basic \LaTeX{} features including mathematical formulas, tables, and lists. The bibliography system is now set up to work seamlessly with Bibby AI's citation management. Feel free to modify it for your needs!

\bibliographystyle{plain}
\bibliography{cite}

\end{document}\documentclass{article}
\usepackage[utf8]{inputenc}
\usepackage{booktabs}
\usepackage{amsmath}
\usepackage{natbib}

\title{My Document}
\author{Your Name}
\date{\today}

\begin{document}

\maketitle

\section{Introduction}
This is a sample document demonstrating various \LaTeX{} features including mathematical equations, tables, and formatting. You can cite papers using the citation search feature, and they will be automatically added to your \texttt{cite.bib} file.

\section{Mathematical Expressions}
Here are some example equations:

\subsection{Quadratic Formula}
The quadratic formula is given by:
\[x = \frac{-b \pm \sqrt{b^2 - 4ac}}{2a}\]

\subsection{Matrix Example}
A sample matrix:
\[
\begin{pmatrix}
a_{11} & a_{12} & a_{13} \\
a_{21} & a_{22} & a_{23} \\
a_{31} & a_{32} & a_{33}
\end{pmatrix}
\]

\section{Tables}
Below is a sample table showing data:

\begin{table}[h]
\centering
\begin{tabular}{@{}lrrr@{}}
\toprule
\textbf{Category} & \textbf{Value 1} & \textbf{Value 2} & \textbf{Total} \\
\midrule
Item A & 10.5 & 20.3 & 30.8 \\
Item B & 15.2 & 25.1 & 40.3 \\
Item C & 22.1 & 12.7 & 34.8 \\
\bottomrule
\end{tabular}
\caption{Sample Data Table}
\label{tab:sample}
\end{table}

\section{Lists}
\subsection{Bullet Points}
\begin{itemize}
  \item First item
  \item Second item
  \item Third item with sub-items:
    \begin{itemize}
      \item Sub-item 1
      \item Sub-item 2
    \end{itemize}
\end{itemize}

\subsection{Numbered List}
\begin{enumerate}
  \item First step
  \item Second step
  \item Third step
\end{enumerate}

\section{Citations and References}
When you use the citation search feature, references will be automatically managed. For example, you might cite a paper like this: \cite{example2024}. All citations will be automatically added to your \texttt{cite.bib} file.

\section{Conclusion}
This template demonstrates basic \LaTeX{} features including mathematical formulas, tables, and lists. The bibliography system is now set up to work seamlessly with Bibby AI's citation management. Feel free to modify it for your needs!

\bibliographystyle{plain}
\bibliography{cite}

\end{document}\documentclass{article}
\usepackage[utf8]{inputenc}
\usepackage{booktabs}
\usepackage{amsmath}
\usepackage{natbib}

\title{My Document}
\author{Your Name}
\date{\today}

\begin{document}

\maketitle

\section{Introduction}
This is a sample document demonstrating various \LaTeX{} features including mathematical equations, tables, and formatting. You can cite papers using the citation search feature, and they will be automatically added to your \texttt{cite.bib} file.

\section{Mathematical Expressions}
Here are some example equations:

\subsection{Quadratic Formula}
The quadratic formula is given by:
\[x = \frac{-b \pm \sqrt{b^2 - 4ac}}{2a}\]

\subsection{Matrix Example}
A sample matrix:
\[
\begin{pmatrix}
a_{11} & a_{12} & a_{13} \\
a_{21} & a_{22} & a_{23} \\
a_{31} & a_{32} & a_{33}
\end{pmatrix}
\]

\section{Tables}
Below is a sample table showing data:

\begin{table}[h]
\centering
\begin{tabular}{@{}lrrr@{}}
\toprule
\textbf{Category} & \textbf{Value 1} & \textbf{Value 2} & \textbf{Total} \\
\midrule
Item A & 10.5 & 20.3 & 30.8 \\
Item B & 15.2 & 25.1 & 40.3 \\
Item C & 22.1 & 12.7 & 34.8 \\
\bottomrule
\end{tabular}
\caption{Sample Data Table}
\label{tab:sample}
\end{table}

\section{Lists}
\subsection{Bullet Points}
\begin{itemize}
  \item First item
  \item Second item
  \item Third item with sub-items:
    \begin{itemize}
      \item Sub-item 1
      \item Sub-item 2
    \end{itemize}
\end{itemize}

\subsection{Numbered List}
\begin{enumerate}
  \item First step
  \item Second step
  \item Third step
\end{enumerate}

\section{Citations and References}
When you use the citation search feature, references will be automatically managed. For example, you might cite a paper like this: \cite{example2024}. All citations will be automatically added to your \texttt{cite.bib} file.

\section{Conclusion}
This template demonstrates basic \LaTeX{} features including mathematical formulas, tables, and lists. The bibliography system is now set up to work seamlessly with Bibby AI's citation management. Feel free to modify it for your needs!

\bibliographystyle{plain}
\bibliography{cite}

\end{document}\documentclass{article}
\usepackage[utf8]{inputenc}
\usepackage{booktabs}
\usepackage{amsmath}
\usepackage{natbib}

\title{My Document}
\author{Your Name}
\date{\today}

\begin{document}

\maketitle

\section{Introduction}
This is a sample document demonstrating various \LaTeX{} features including mathematical equations, tables, and formatting. You can cite papers using the citation search feature, and they will be automatically added to your \texttt{cite.bib} file.

\section{Mathematical Expressions}
Here are some example equations:

\subsection{Quadratic Formula}
The quadratic formula is given by:
\[x = \frac{-b \pm \sqrt{b^2 - 4ac}}{2a}\]

\subsection{Matrix Example}
A sample matrix:
\[
\begin{pmatrix}
a_{11} & a_{12} & a_{13} \\
a_{21} & a_{22} & a_{23} \\
a_{31} & a_{32} & a_{33}
\end{pmatrix}
\]

\section{Tables}
Below is a sample table showing data:

\begin{table}[h]
\centering
\begin{tabular}{@{}lrrr@{}}
\toprule
\textbf{Category} & \textbf{Value 1} & \textbf{Value 2} & \textbf{Total} \\
\midrule
Item A & 10.5 & 20.3 & 30.8 \\
Item B & 15.2 & 25.1 & 40.3 \\
Item C & 22.1 & 12.7 & 34.8 \\
\bottomrule
\end{tabular}
\caption{Sample Data Table}
\label{tab:sample}
\end{table}

\section{Lists}
\subsection{Bullet Points}
\begin{itemize}
  \item First item
  \item Second item
  \item Third item with sub-items:
    \begin{itemize}
      \item Sub-item 1
      \item Sub-item 2
    \end{itemize}
\end{itemize}

\subsection{Numbered List}
\begin{enumerate}
  \item First step
  \item Second step
  \item Third step
\end{enumerate}

\section{Citations and References}
When you use the citation search feature, references will be automatically managed. For example, you might cite a paper like this: \cite{example2024}. All citations will be automatically added to your \texttt{cite.bib} file.

\section{Conclusion}
This template demonstrates basic \LaTeX{} features including mathematical formulas, tables, and lists. The bibliography system is now set up to work seamlessly with Bibby AI's citation management. Feel free to modify it for your needs!

\bibliographystyle{plain}
\bibliography{cite}

\end{document}\documentclass{article}
\usepackage[utf8]{inputenc}
\usepackage{booktabs}
\usepackage{amsmath}
\usepackage{natbib}

\title{My Document}
\author{Your Name}
\date{\today}

\begin{document}

\maketitle

\section{Introduction}
This is a sample document demonstrating various \LaTeX{} features including mathematical equations, tables, and formatting. You can cite papers using the citation search feature, and they will be automatically added to your \texttt{cite.bib} file.

\section{Mathematical Expressions}
Here are some example equations:

\subsection{Quadratic Formula}
The quadratic formula is given by:
\[x = \frac{-b \pm \sqrt{b^2 - 4ac}}{2a}\]

\subsection{Matrix Example}
A sample matrix:
\[
\begin{pmatrix}
a_{11} & a_{12} & a_{13} \\
a_{21} & a_{22} & a_{23} \\
a_{31} & a_{32} & a_{33}
\end{pmatrix}
\]

\section{Tables}
Below is a sample table showing data:

\begin{table}[h]
\centering
\begin{tabular}{@{}lrrr@{}}
\toprule
\textbf{Category} & \textbf{Value 1} & \textbf{Value 2} & \textbf{Total} \\
\midrule
Item A & 10.5 & 20.3 & 30.8 \\
Item B & 15.2 & 25.1 & 40.3 \\
Item C & 22.1 & 12.7 & 34.8 \\
\bottomrule
\end{tabular}
\caption{Sample Data Table}
\label{tab:sample}
\end{table}

\section{Lists}
\subsection{Bullet Points}
\begin{itemize}
  \item First item
  \item Second item
  \item Third item with sub-items:
    \begin{itemize}
      \item Sub-item 1
      \item Sub-item 2
    \end{itemize}
\end{itemize}

\subsection{Numbered List}
\begin{enumerate}
  \item First step
  \item Second step
  \item Third step
\end{enumerate}

\section{Citations and References}
When you use the citation search feature, references will be automatically managed. For example, you might cite a paper like this: \cite{example2024}. All citations will be automatically added to your \texttt{cite.bib} file.

\section{Conclusion}
This template demonstrates basic \LaTeX{} features including mathematical formulas, tables, and lists. The bibliography system is now set up to work seamlessly with Bibby AI's citation management. Feel free to modify it for your needs!

\bibliographystyle{plain}
\bibliography{cite}

\end{document}\documentclass{article}
\usepackage[utf8]{inputenc}
\usepackage{booktabs}
\usepackage{amsmath}
\usepackage{natbib}

\title{My Document}
\author{Your Name}
\date{\today}

\begin{document}

\maketitle

\section{Introduction}
This is a sample document demonstrating various \LaTeX{} features including mathematical equations, tables, and formatting. You can cite papers using the citation search feature, and they will be automatically added to your \texttt{cite.bib} file.

\section{Mathematical Expressions}
Here are some example equations:

\subsection{Quadratic Formula}
The quadratic formula is given by:
\[x = \frac{-b \pm \sqrt{b^2 - 4ac}}{2a}\]

\subsection{Matrix Example}
A sample matrix:
\[
\begin{pmatrix}
a_{11} & a_{12} & a_{13} \\
a_{21} & a_{22} & a_{23} \\
a_{31} & a_{32} & a_{33}
\end{pmatrix}
\]

\section{Tables}
Below is a sample table showing data:

\begin{table}[h]
\centering
\begin{tabular}{@{}lrrr@{}}
\toprule
\textbf{Category} & \textbf{Value 1} & \textbf{Value 2} & \textbf{Total} \\
\midrule
Item A & 10.5 & 20.3 & 30.8 \\
Item B & 15.2 & 25.1 & 40.3 \\
Item C & 22.1 & 12.7 & 34.8 \\
\bottomrule
\end{tabular}
\caption{Sample Data Table}
\label{tab:sample}
\end{table}

\section{Lists}
\subsection{Bullet Points}
\begin{itemize}
  \item First item
  \item Second item
  \item Third item with sub-items:
    \begin{itemize}
      \item Sub-item 1
      \item Sub-item 2
    \end{itemize}
\end{itemize}

\subsection{Numbered List}
\begin{enumerate}
  \item First step
  \item Second step
  \item Third step
\end{enumerate}

\section{Citations and References}
When you use the citation search feature, references will be automatically managed. For example, you might cite a paper like this: \cite{example2024}. All citations will be automatically added to your \texttt{cite.bib} file.

\section{Conclusion}
This template demonstrates basic \LaTeX{} features including mathematical formulas, tables, and lists. The bibliography system is now set up to work seamlessly with Bibby AI's citation management. Feel free to modify it for your needs!

\bibliographystyle{plain}
\bibliography{cite}

\end{document}\documentclass{article}
\usepackage[utf8]{inputenc}
\usepackage{booktabs}
\usepackage{amsmath}
\usepackage{natbib}

\title{My Document}
\author{Your Name}
\date{\today}

\begin{document}

\maketitle

\section{Introduction}
This is a sample document demonstrating various \LaTeX{} features including mathematical equations, tables, and formatting. You can cite papers using the citation search feature, and they will be automatically added to your \texttt{cite.bib} file.

\section{Mathematical Expressions}
Here are some example equations:

\subsection{Quadratic Formula}
The quadratic formula is given by:
\[x = \frac{-b \pm \sqrt{b^2 - 4ac}}{2a}\]

\subsection{Matrix Example}
A sample matrix:
\[
\begin{pmatrix}
a_{11} & a_{12} & a_{13} \\
a_{21} & a_{22} & a_{23} \\
a_{31} & a_{32} & a_{33}
\end{pmatrix}
\]

\section{Tables}
Below is a sample table showing data:

\begin{table}[h]
\centering
\begin{tabular}{@{}lrrr@{}}
\toprule
\textbf{Category} & \textbf{Value 1} & \textbf{Value 2} & \textbf{Total} \\
\midrule
Item A & 10.5 & 20.3 & 30.8 \\
Item B & 15.2 & 25.1 & 40.3 \\
Item C & 22.1 & 12.7 & 34.8 \\
\bottomrule
\end{tabular}
\caption{Sample Data Table}
\label{tab:sample}
\end{table}

\section{Lists}
\subsection{Bullet Points}
\begin{itemize}
  \item First item
  \item Second item
  \item Third item with sub-items:
    \begin{itemize}
      \item Sub-item 1
      \item Sub-item 2
    \end{itemize}
\end{itemize}

\subsection{Numbered List}
\begin{enumerate}
  \item First step
  \item Second step
  \item Third step
\end{enumerate}

\section{Citations and References}
When you use the citation search feature, references will be automatically managed. For example, you might cite a paper like this: \cite{example2024}. All citations will be automatically added to your \texttt{cite.bib} file.

\section{Conclusion}
This template demonstrates basic \LaTeX{} features including mathematical formulas, tables, and lists. The bibliography system is now set up to work seamlessly with Bibby AI's citation management. Feel free to modify it for your needs!

\bibliographystyle{plain}
\bibliography{cite}

\end{document}\documentclass{article}
\usepackage[utf8]{inputenc}
\usepackage{booktabs}
\usepackage{amsmath}
\usepackage{natbib}

\title{My Document}
\author{Your Name}
\date{\today}

\begin{document}

\maketitle

\section{Introduction}
This is a sample document demonstrating various \LaTeX{} features including mathematical equations, tables, and formatting. You can cite papers using the citation search feature, and they will be automatically added to your \texttt{cite.bib} file.

\section{Mathematical Expressions}
Here are some example equations:

\subsection{Quadratic Formula}
The quadratic formula is given by:
\[x = \frac{-b \pm \sqrt{b^2 - 4ac}}{2a}\]

\subsection{Matrix Example}
A sample matrix:
\[
\begin{pmatrix}
a_{11} & a_{12} & a_{13} \\
a_{21} & a_{22} & a_{23} \\
a_{31} & a_{32} & a_{33}
\end{pmatrix}
\]

\section{Tables}
Below is a sample table showing data:

\begin{table}[h]
\centering
\begin{tabular}{@{}lrrr@{}}
\toprule
\textbf{Category} & \textbf{Value 1} & \textbf{Value 2} & \textbf{Total} \\
\midrule
Item A & 10.5 & 20.3 & 30.8 \\
Item B & 15.2 & 25.1 & 40.3 \\
Item C & 22.1 & 12.7 & 34.8 \\
\bottomrule
\end{tabular}
\caption{Sample Data Table}
\label{tab:sample}
\end{table}

\section{Lists}
\subsection{Bullet Points}
\begin{itemize}
  \item First item
  \item Second item
  \item Third item with sub-items:
    \begin{itemize}
      \item Sub-item 1
      \item Sub-item 2
    \end{itemize}
\end{itemize}

\subsection{Numbered List}
\begin{enumerate}
  \item First step
  \item Second step
  \item Third step
\end{enumerate}

\section{Citations and References}
When you use the citation search feature, references will be automatically managed. For example, you might cite a paper like this: \cite{example2024}. All citations will be automatically added to your \texttt{cite.bib} file.

\section{Conclusion}
This template demonstrates basic \LaTeX{} features including mathematical formulas, tables, and lists. The bibliography system is now set up to work seamlessly with Bibby AI's citation management. Feel free to modify it for your needs!

\bibliographystyle{plain}
\bibliography{cite}

\end{document}\documentclass{article}
\usepackage[utf8]{inputenc}
\usepackage{booktabs}
\usepackage{amsmath}
\usepackage{natbib}

\title{My Document}
\author{Your Name}
\date{\today}

\begin{document}

\maketitle

\section{Introduction}
This is a sample document demonstrating various \LaTeX{} features including mathematical equations, tables, and formatting. You can cite papers using the citation search feature, and they will be automatically added to your \texttt{cite.bib} file.

\section{Mathematical Expressions}
Here are some example equations:

\subsection{Quadratic Formula}
The quadratic formula is given by:
\[x = \frac{-b \pm \sqrt{b^2 - 4ac}}{2a}\]

\subsection{Matrix Example}
A sample matrix:
\[
\begin{pmatrix}
a_{11} & a_{12} & a_{13} \\
a_{21} & a_{22} & a_{23} \\
a_{31} & a_{32} & a_{33}
\end{pmatrix}
\]

\section{Tables}
Below is a sample table showing data:

\begin{table}[h]
\centering
\begin{tabular}{@{}lrrr@{}}
\toprule
\textbf{Category} & \textbf{Value 1} & \textbf{Value 2} & \textbf{Total} \\
\midrule
Item A & 10.5 & 20.3 & 30.8 \\
Item B & 15.2 & 25.1 & 40.3 \\
Item C & 22.1 & 12.7 & 34.8 \\
\bottomrule
\end{tabular}
\caption{Sample Data Table}
\label{tab:sample}
\end{table}

\section{Lists}
\subsection{Bullet Points}
\begin{itemize}
  \item First item
  \item Second item
  \item Third item with sub-items:
    \begin{itemize}
      \item Sub-item 1
      \item Sub-item 2
    \end{itemize}
\end{itemize}

\subsection{Numbered List}
\begin{enumerate}
  \item First step
  \item Second step
  \item Third step
\end{enumerate}

\section{Citations and References}
When you use the citation search feature, references will be automatically managed. For example, you might cite a paper like this: \cite{example2024}. All citations will be automatically added to your \texttt{cite.bib} file.

\section{Conclusion}
This template demonstrates basic \LaTeX{} features including mathematical formulas, tables, and lists. The bibliography system is now set up to work seamlessly with Bibby AI's citation management. Feel free to modify it for your needs!

\bibliographystyle{plain}
\bibliography{cite}

\end{document}\documentclass{article}
\usepackage[utf8]{inputenc}
\usepackage{booktabs}
\usepackage{amsmath}
\usepackage{natbib}

\title{My Document}
\author{Your Name}
\date{\today}

\begin{document}

\maketitle

\section{Introduction}
This is a sample document demonstrating various \LaTeX{} features including mathematical equations, tables, and formatting. You can cite papers using the citation search feature, and they will be automatically added to your \texttt{cite.bib} file.

\section{Mathematical Expressions}
Here are some example equations:

\subsection{Quadratic Formula}
The quadratic formula is given by:
\[x = \frac{-b \pm \sqrt{b^2 - 4ac}}{2a}\]

\subsection{Matrix Example}
A sample matrix:
\[
\begin{pmatrix}
a_{11} & a_{12} & a_{13} \\
a_{21} & a_{22} & a_{23} \\
a_{31} & a_{32} & a_{33}
\end{pmatrix}
\]

\section{Tables}
Below is a sample table showing data:

\begin{table}[h]
\centering
\begin{tabular}{@{}lrrr@{}}
\toprule
\textbf{Category} & \textbf{Value 1} & \textbf{Value 2} & \textbf{Total} \\
\midrule
Item A & 10.5 & 20.3 & 30.8 \\
Item B & 15.2 & 25.1 & 40.3 \\
Item C & 22.1 & 12.7 & 34.8 \\
\bottomrule
\end{tabular}
\caption{Sample Data Table}
\label{tab:sample}
\end{table}

\section{Lists}
\subsection{Bullet Points}
\begin{itemize}
  \item First item
  \item Second item
  \item Third item with sub-items:
    \begin{itemize}
      \item Sub-item 1
      \item Sub-item 2
    \end{itemize}
\end{itemize}

\subsection{Numbered List}
\begin{enumerate}
  \item First step
  \item Second step
  \item Third step
\end{enumerate}

\section{Citations and References}
When you use the citation search feature, references will be automatically managed. For example, you might cite a paper like this: \cite{example2024}. All citations will be automatically added to your \texttt{cite.bib} file.

\section{Conclusion}
This template demonstrates basic \LaTeX{} features including mathematical formulas, tables, and lists. The bibliography system is now set up to work seamlessly with Bibby AI's citation management. Feel free to modify it for your needs!

\bibliographystyle{plain}
\bibliography{cite}

\end{document}\documentclass{article}
\usepackage[utf8]{inputenc}
\usepackage{booktabs}
\usepackage{amsmath}
\usepackage{natbib}

\title{My Document}
\author{Your Name}
\date{\today}

\begin{document}

\maketitle

\section{Introduction}
This is a sample document demonstrating various \LaTeX{} features including mathematical equations, tables, and formatting. You can cite papers using the citation search feature, and they will be automatically added to your \texttt{cite.bib} file.

\section{Mathematical Expressions}
Here are some example equations:

\subsection{Quadratic Formula}
The quadratic formula is given by:
\[x = \frac{-b \pm \sqrt{b^2 - 4ac}}{2a}\]

\subsection{Matrix Example}
A sample matrix:
\[
\begin{pmatrix}
a_{11} & a_{12} & a_{13} \\
a_{21} & a_{22} & a_{23} \\
a_{31} & a_{32} & a_{33}
\end{pmatrix}
\]

\section{Tables}
Below is a sample table showing data:

\begin{table}[h]
\centering
\begin{tabular}{@{}lrrr@{}}
\toprule
\textbf{Category} & \textbf{Value 1} & \textbf{Value 2} & \textbf{Total} \\
\midrule
Item A & 10.5 & 20.3 & 30.8 \\
Item B & 15.2 & 25.1 & 40.3 \\
Item C & 22.1 & 12.7 & 34.8 \\
\bottomrule
\end{tabular}
\caption{Sample Data Table}
\label{tab:sample}
\end{table}

\section{Lists}
\subsection{Bullet Points}
\begin{itemize}
  \item First item
  \item Second item
  \item Third item with sub-items:
    \begin{itemize}
      \item Sub-item 1
      \item Sub-item 2
    \end{itemize}
\end{itemize}

\subsection{Numbered List}
\begin{enumerate}
  \item First step
  \item Second step
  \item Third step
\end{enumerate}

\section{Citations and References}
When you use the citation search feature, references will be automatically managed. For example, you might cite a paper like this: \cite{example2024}. All citations will be automatically added to your \texttt{cite.bib} file.

\section{Conclusion}
This template demonstrates basic \LaTeX{} features including mathematical formulas, tables, and lists. The bibliography system is now set up to work seamlessly with Bibby AI's citation management. Feel free to modify it for your needs!

\bibliographystyle{plain}
\bibliography{cite}

\end{document}\documentclass{article}
\usepackage[utf8]{inputenc}
\usepackage{booktabs}
\usepackage{amsmath}
\usepackage{natbib}

\title{My Document}
\author{Your Name}
\date{\today}

\begin{document}

\maketitle

\section{Introduction}
This is a sample document demonstrating various \LaTeX{} features including mathematical equations, tables, and formatting. You can cite papers using the citation search feature, and they will be automatically added to your \texttt{cite.bib} file.

\section{Mathematical Expressions}
Here are some example equations:

\subsection{Quadratic Formula}
The quadratic formula is given by:
\[x = \frac{-b \pm \sqrt{b^2 - 4ac}}{2a}\]

\subsection{Matrix Example}
A sample matrix:
\[
\begin{pmatrix}
a_{11} & a_{12} & a_{13} \\
a_{21} & a_{22} & a_{23} \\
a_{31} & a_{32} & a_{33}
\end{pmatrix}
\]

\section{Tables}
Below is a sample table showing data:

\begin{table}[h]
\centering
\begin{tabular}{@{}lrrr@{}}
\toprule
\textbf{Category} & \textbf{Value 1} & \textbf{Value 2} & \textbf{Total} \\
\midrule
Item A & 10.5 & 20.3 & 30.8 \\
Item B & 15.2 & 25.1 & 40.3 \\
Item C & 22.1 & 12.7 & 34.8 \\
\bottomrule
\end{tabular}
\caption{Sample Data Table}
\label{tab:sample}
\end{table}

\section{Lists}
\subsection{Bullet Points}
\begin{itemize}
  \item First item
  \item Second item
  \item Third item with sub-items:
    \begin{itemize}
      \item Sub-item 1
      \item Sub-item 2
    \end{itemize}
\end{itemize}

\subsection{Numbered List}
\begin{enumerate}
  \item First step
  \item Second step
  \item Third step
\end{enumerate}

\section{Citations and References}
When you use the citation search feature, references will be automatically managed. For example, you might cite a paper like this: \cite{example2024}. All citations will be automatically added to your \texttt{cite.bib} file.

\section{Conclusion}
This template demonstrates basic \LaTeX{} features including mathematical formulas, tables, and lists. The bibliography system is now set up to work seamlessly with Bibby AI's citation management. Feel free to modify it for your needs!

\bibliographystyle{plain}
\bibliography{cite}

\end{document}\documentclass{article}
\usepackage[utf8]{inputenc}
\usepackage{booktabs}
\usepackage{amsmath}
\usepackage{natbib}

\title{My Document}
\author{Your Name}
\date{\today}

\begin{document}

\maketitle

\section{Introduction}
This is a sample document demonstrating various \LaTeX{} features including mathematical equations, tables, and formatting. You can cite papers using the citation search feature, and they will be automatically added to your \texttt{cite.bib} file.

\section{Mathematical Expressions}
Here are some example equations:

\subsection{Quadratic Formula}
The quadratic formula is given by:
\[x = \frac{-b \pm \sqrt{b^2 - 4ac}}{2a}\]

\subsection{Matrix Example}
A sample matrix:
\[
\begin{pmatrix}
a_{11} & a_{12} & a_{13} \\
a_{21} & a_{22} & a_{23} \\
a_{31} & a_{32} & a_{33}
\end{pmatrix}
\]

\section{Tables}
Below is a sample table showing data:

\begin{table}[h]
\centering
\begin{tabular}{@{}lrrr@{}}
\toprule
\textbf{Category} & \textbf{Value 1} & \textbf{Value 2} & \textbf{Total} \\
\midrule
Item A & 10.5 & 20.3 & 30.8 \\
Item B & 15.2 & 25.1 & 40.3 \\
Item C & 22.1 & 12.7 & 34.8 \\
\bottomrule
\end{tabular}
\caption{Sample Data Table}
\label{tab:sample}
\end{table}

\section{Lists}
\subsection{Bullet Points}
\begin{itemize}
  \item First item
  \item Second item
  \item Third item with sub-items:
    \begin{itemize}
      \item Sub-item 1
      \item Sub-item 2
    \end{itemize}
\end{itemize}

\subsection{Numbered List}
\begin{enumerate}
  \item First step
  \item Second step
  \item Third step
\end{enumerate}

\section{Citations and References}
When you use the citation search feature, references will be automatically managed. For example, you might cite a paper like this: \cite{example2024}. All citations will be automatically added to your \texttt{cite.bib} file.

\section{Conclusion}
This template demonstrates basic \LaTeX{} features including mathematical formulas, tables, and lists. The bibliography system is now set up to work seamlessly with Bibby AI's citation management. Feel free to modify it for your needs!

\bibliographystyle{plain}
\bibliography{cite}

\end{document}\documentclass{article}
\usepackage[utf8]{inputenc}
\usepackage{booktabs}
\usepackage{amsmath}
\usepackage{natbib}

\title{My Document}
\author{Your Name}
\date{\today}

\begin{document}

\maketitle

\section{Introduction}
This is a sample document demonstrating various \LaTeX{} features including mathematical equations, tables, and formatting. You can cite papers using the citation search feature, and they will be automatically added to your \texttt{cite.bib} file.

\section{Mathematical Expressions}
Here are some example equations:

\subsection{Quadratic Formula}
The quadratic formula is given by:
\[x = \frac{-b \pm \sqrt{b^2 - 4ac}}{2a}\]

\subsection{Matrix Example}
A sample matrix:
\[
\begin{pmatrix}
a_{11} & a_{12} & a_{13} \\
a_{21} & a_{22} & a_{23} \\
a_{31} & a_{32} & a_{33}
\end{pmatrix}
\]

\section{Tables}
Below is a sample table showing data:

\begin{table}[h]
\centering
\begin{tabular}{@{}lrrr@{}}
\toprule
\textbf{Category} & \textbf{Value 1} & \textbf{Value 2} & \textbf{Total} \\
\midrule
Item A & 10.5 & 20.3 & 30.8 \\
Item B & 15.2 & 25.1 & 40.3 \\
Item C & 22.1 & 12.7 & 34.8 \\
\bottomrule
\end{tabular}
\caption{Sample Data Table}
\label{tab:sample}
\end{table}

\section{Lists}
\subsection{Bullet Points}
\begin{itemize}
  \item First item
  \item Second item
  \item Third item with sub-items:
    \begin{itemize}
      \item Sub-item 1
      \item Sub-item 2
    \end{itemize}
\end{itemize}

\subsection{Numbered List}
\begin{enumerate}
  \item First step
  \item Second step
  \item Third step
\end{enumerate}

\section{Citations and References}
When you use the citation search feature, references will be automatically managed. For example, you might cite a paper like this: \cite{example2024}. All citations will be automatically added to your \texttt{cite.bib} file.

\section{Conclusion}
This template demonstrates basic \LaTeX{} features including mathematical formulas, tables, and lists. The bibliography system is now set up to work seamlessly with Bibby AI's citation management. Feel free to modify it for your needs!

\bibliographystyle{plain}
\bibliography{cite}

\end{document}\documentclass{article}
\usepackage[utf8]{inputenc}
\usepackage{booktabs}
\usepackage{amsmath}
\usepackage{natbib}

\title{My Document}
\author{Your Name}
\date{\today}

\begin{document}

\maketitle

\section{Introduction}
This is a sample document demonstrating various \LaTeX{} features including mathematical equations, tables, and formatting. You can cite papers using the citation search feature, and they will be automatically added to your \texttt{cite.bib} file.

\section{Mathematical Expressions}
Here are some example equations:

\subsection{Quadratic Formula}
The quadratic formula is given by:
\[x = \frac{-b \pm \sqrt{b^2 - 4ac}}{2a}\]

\subsection{Matrix Example}
A sample matrix:
\[
\begin{pmatrix}
a_{11} & a_{12} & a_{13} \\
a_{21} & a_{22} & a_{23} \\
a_{31} & a_{32} & a_{33}
\end{pmatrix}
\]

\section{Tables}
Below is a sample table showing data:

\begin{table}[h]
\centering
\begin{tabular}{@{}lrrr@{}}
\toprule
\textbf{Category} & \textbf{Value 1} & \textbf{Value 2} & \textbf{Total} \\
\midrule
Item A & 10.5 & 20.3 & 30.8 \\
Item B & 15.2 & 25.1 & 40.3 \\
Item C & 22.1 & 12.7 & 34.8 \\
\bottomrule
\end{tabular}
\caption{Sample Data Table}
\label{tab:sample}
\end{table}

\section{Lists}
\subsection{Bullet Points}
\begin{itemize}
  \item First item
  \item Second item
  \item Third item with sub-items:
    \begin{itemize}
      \item Sub-item 1
      \item Sub-item 2
    \end{itemize}
\end{itemize}

\subsection{Numbered List}
\begin{enumerate}
  \item First step
  \item Second step
  \item Third step
\end{enumerate}

\section{Citations and References}
When you use the citation search feature, references will be automatically managed. For example, you might cite a paper like this: \cite{example2024}. All citations will be automatically added to your \texttt{cite.bib} file.

\section{Conclusion}
This template demonstrates basic \LaTeX{} features including mathematical formulas, tables, and lists. The bibliography system is now set up to work seamlessly with Bibby AI's citation management. Feel free to modify it for your needs!

\bibliographystyle{plain}
\bibliography{cite}

\end{document}\documentclass{article}
\usepackage[utf8]{inputenc}
\usepackage{booktabs}
\usepackage{amsmath}
\usepackage{natbib}

\title{My Document}
\author{Your Name}
\date{\today}

\begin{document}

\maketitle

\section{Introduction}
This is a sample document demonstrating various \LaTeX{} features including mathematical equations, tables, and formatting. You can cite papers using the citation search feature, and they will be automatically added to your \texttt{cite.bib} file.

\section{Mathematical Expressions}
Here are some example equations:

\subsection{Quadratic Formula}
The quadratic formula is given by:
\[x = \frac{-b \pm \sqrt{b^2 - 4ac}}{2a}\]

\subsection{Matrix Example}
A sample matrix:
\[
\begin{pmatrix}
a_{11} & a_{12} & a_{13} \\
a_{21} & a_{22} & a_{23} \\
a_{31} & a_{32} & a_{33}
\end{pmatrix}
\]

\section{Tables}
Below is a sample table showing data:

\begin{table}[h]
\centering
\begin{tabular}{@{}lrrr@{}}
\toprule
\textbf{Category} & \textbf{Value 1} & \textbf{Value 2} & \textbf{Total} \\
\midrule
Item A & 10.5 & 20.3 & 30.8 \\
Item B & 15.2 & 25.1 & 40.3 \\
Item C & 22.1 & 12.7 & 34.8 \\
\bottomrule
\end{tabular}
\caption{Sample Data Table}
\label{tab:sample}
\end{table}

\section{Lists}
\subsection{Bullet Points}
\begin{itemize}
  \item First item
  \item Second item
  \item Third item with sub-items:
    \begin{itemize}
      \item Sub-item 1
      \item Sub-item 2
    \end{itemize}
\end{itemize}

\subsection{Numbered List}
\begin{enumerate}
  \item First step
  \item Second step
  \item Third step
\end{enumerate}

\section{Citations and References}
When you use the citation search feature, references will be automatically managed. For example, you might cite a paper like this: \cite{example2024}. All citations will be automatically added to your \texttt{cite.bib} file.

\section{Conclusion}
This template demonstrates basic \LaTeX{} features including mathematical formulas, tables, and lists. The bibliography system is now set up to work seamlessly with Bibby AI's citation management. Feel free to modify it for your needs!

\bibliographystyle{plain}
\bibliography{cite}

\end{document}\documentclass{article}
\usepackage[utf8]{inputenc}
\usepackage{booktabs}
\usepackage{amsmath}
\usepackage{natbib}

\title{My Document}
\author{Your Name}
\date{\today}

\begin{document}

\maketitle

\section{Introduction}
This is a sample document demonstrating various \LaTeX{} features including mathematical equations, tables, and formatting. You can cite papers using the citation search feature, and they will be automatically added to your \texttt{cite.bib} file.

\section{Mathematical Expressions}
Here are some example equations:

\subsection{Quadratic Formula}
The quadratic formula is given by:
\[x = \frac{-b \pm \sqrt{b^2 - 4ac}}{2a}\]

\subsection{Matrix Example}
A sample matrix:
\[
\begin{pmatrix}
a_{11} & a_{12} & a_{13} \\
a_{21} & a_{22} & a_{23} \\
a_{31} & a_{32} & a_{33}
\end{pmatrix}
\]

\section{Tables}
Below is a sample table showing data:

\begin{table}[h]
\centering
\begin{tabular}{@{}lrrr@{}}
\toprule
\textbf{Category} & \textbf{Value 1} & \textbf{Value 2} & \textbf{Total} \\
\midrule
Item A & 10.5 & 20.3 & 30.8 \\
Item B & 15.2 & 25.1 & 40.3 \\
Item C & 22.1 & 12.7 & 34.8 \\
\bottomrule
\end{tabular}
\caption{Sample Data Table}
\label{tab:sample}
\end{table}

\section{Lists}
\subsection{Bullet Points}
\begin{itemize}
  \item First item
  \item Second item
  \item Third item with sub-items:
    \begin{itemize}
      \item Sub-item 1
      \item Sub-item 2
    \end{itemize}
\end{itemize}

\subsection{Numbered List}
\begin{enumerate}
  \item First step
  \item Second step
  \item Third step
\end{enumerate}

\section{Citations and References}
When you use the citation search feature, references will be automatically managed. For example, you might cite a paper like this: \cite{example2024}. All citations will be automatically added to your \texttt{cite.bib} file.

\section{Conclusion}
This template demonstrates basic \LaTeX{} features including mathematical formulas, tables, and lists. The bibliography system is now set up to work seamlessly with Bibby AI's citation management. Feel free to modify it for your needs!

\bibliographystyle{plain}
\bibliography{cite}

\end{document}\documentclass{article}
\usepackage[utf8]{inputenc}
\usepackage{booktabs}
\usepackage{amsmath}
\usepackage{natbib}

\title{My Document}
\author{Your Name}
\date{\today}

\begin{document}

\maketitle

\section{Introduction}
This is a sample document demonstrating various \LaTeX{} features including mathematical equations, tables, and formatting. You can cite papers using the citation search feature, and they will be automatically added to your \texttt{cite.bib} file.

\section{Mathematical Expressions}
Here are some example equations:

\subsection{Quadratic Formula}
The quadratic formula is given by:
\[x = \frac{-b \pm \sqrt{b^2 - 4ac}}{2a}\]

\subsection{Matrix Example}
A sample matrix:
\[
\begin{pmatrix}
a_{11} & a_{12} & a_{13} \\
a_{21} & a_{22} & a_{23} \\
a_{31} & a_{32} & a_{33}
\end{pmatrix}
\]

\section{Tables}
Below is a sample table showing data:

\begin{table}[h]
\centering
\begin{tabular}{@{}lrrr@{}}
\toprule
\textbf{Category} & \textbf{Value 1} & \textbf{Value 2} & \textbf{Total} \\
\midrule
Item A & 10.5 & 20.3 & 30.8 \\
Item B & 15.2 & 25.1 & 40.3 \\
Item C & 22.1 & 12.7 & 34.8 \\
\bottomrule
\end{tabular}
\caption{Sample Data Table}
\label{tab:sample}
\end{table}

\section{Lists}
\subsection{Bullet Points}
\begin{itemize}
  \item First item
  \item Second item
  \item Third item with sub-items:
    \begin{itemize}
      \item Sub-item 1
      \item Sub-item 2
    \end{itemize}
\end{itemize}

\subsection{Numbered List}
\begin{enumerate}
  \item First step
  \item Second step
  \item Third step
\end{enumerate}

\section{Citations and References}
When you use the citation search feature, references will be automatically managed. For example, you might cite a paper like this: \cite{example2024}. All citations will be automatically added to your \texttt{cite.bib} file.

\section{Conclusion}
This template demonstrates basic \LaTeX{} features including mathematical formulas, tables, and lists. The bibliography system is now set up to work seamlessly with Bibby AI's citation management. Feel free to modify it for your needs!

\bibliographystyle{plain}
\bibliography{cite}

\end{document}\documentclass{article}
\usepackage[utf8]{inputenc}
\usepackage{booktabs}
\usepackage{amsmath}
\usepackage{natbib}

\title{My Document}
\author{Your Name}
\date{\today}

\begin{document}

\maketitle

\section{Introduction}
This is a sample document demonstrating various \LaTeX{} features including mathematical equations, tables, and formatting. You can cite papers using the citation search feature, and they will be automatically added to your \texttt{cite.bib} file.

\section{Mathematical Expressions}
Here are some example equations:

\subsection{Quadratic Formula}
The quadratic formula is given by:
\[x = \frac{-b \pm \sqrt{b^2 - 4ac}}{2a}\]

\subsection{Matrix Example}
A sample matrix:
\[
\begin{pmatrix}
a_{11} & a_{12} & a_{13} \\
a_{21} & a_{22} & a_{23} \\
a_{31} & a_{32} & a_{33}
\end{pmatrix}
\]

\section{Tables}
Below is a sample table showing data:

\begin{table}[h]
\centering
\begin{tabular}{@{}lrrr@{}}
\toprule
\textbf{Category} & \textbf{Value 1} & \textbf{Value 2} & \textbf{Total} \\
\midrule
Item A & 10.5 & 20.3 & 30.8 \\
Item B & 15.2 & 25.1 & 40.3 \\
Item C & 22.1 & 12.7 & 34.8 \\
\bottomrule
\end{tabular}
\caption{Sample Data Table}
\label{tab:sample}
\end{table}

\section{Lists}
\subsection{Bullet Points}
\begin{itemize}
  \item First item
  \item Second item
  \item Third item with sub-items:
    \begin{itemize}
      \item Sub-item 1
      \item Sub-item 2
    \end{itemize}
\end{itemize}

\subsection{Numbered List}
\begin{enumerate}
  \item First step
  \item Second step
  \item Third step
\end{enumerate}

\section{Citations and References}
When you use the citation search feature, references will be automatically managed. For example, you might cite a paper like this: \cite{example2024}. All citations will be automatically added to your \texttt{cite.bib} file.

\section{Conclusion}
This template demonstrates basic \LaTeX{} features including mathematical formulas, tables, and lists. The bibliography system is now set up to work seamlessly with Bibby AI's citation management. Feel free to modify it for your needs!

\bibliographystyle{plain}
\bibliography{cite}

\end{document}\documentclass{article}
\usepackage[utf8]{inputenc}
\usepackage{booktabs}
\usepackage{amsmath}
\usepackage{natbib}

\title{My Document}
\author{Your Name}
\date{\today}

\begin{document}

\maketitle

\section{Introduction}
This is a sample document demonstrating various \LaTeX{} features including mathematical equations, tables, and formatting. You can cite papers using the citation search feature, and they will be automatically added to your \texttt{cite.bib} file.

\section{Mathematical Expressions}
Here are some example equations:

\subsection{Quadratic Formula}
The quadratic formula is given by:
\[x = \frac{-b \pm \sqrt{b^2 - 4ac}}{2a}\]

\subsection{Matrix Example}
A sample matrix:
\[
\begin{pmatrix}
a_{11} & a_{12} & a_{13} \\
a_{21} & a_{22} & a_{23} \\
a_{31} & a_{32} & a_{33}
\end{pmatrix}
\]

\section{Tables}
Below is a sample table showing data:

\begin{table}[h]
\centering
\begin{tabular}{@{}lrrr@{}}
\toprule
\textbf{Category} & \textbf{Value 1} & \textbf{Value 2} & \textbf{Total} \\
\midrule
Item A & 10.5 & 20.3 & 30.8 \\
Item B & 15.2 & 25.1 & 40.3 \\
Item C & 22.1 & 12.7 & 34.8 \\
\bottomrule
\end{tabular}
\caption{Sample Data Table}
\label{tab:sample}
\end{table}

\section{Lists}
\subsection{Bullet Points}
\begin{itemize}
  \item First item
  \item Second item
  \item Third item with sub-items:
    \begin{itemize}
      \item Sub-item 1
      \item Sub-item 2
    \end{itemize}
\end{itemize}

\subsection{Numbered List}
\begin{enumerate}
  \item First step
  \item Second step
  \item Third step
\end{enumerate}

\section{Citations and References}
When you use the citation search feature, references will be automatically managed. For example, you might cite a paper like this: \cite{example2024}. All citations will be automatically added to your \texttt{cite.bib} file.

\section{Conclusion}
This template demonstrates basic \LaTeX{} features including mathematical formulas, tables, and lists. The bibliography system is now set up to work seamlessly with Bibby AI's citation management. Feel free to modify it for your needs!

\bibliographystyle{plain}
\bibliography{cite}

\end{document}\documentclass{article}
\usepackage[utf8]{inputenc}
\usepackage{booktabs}
\usepackage{amsmath}
\usepackage{natbib}

\title{My Document}
\author{Your Name}
\date{\today}

\begin{document}

\maketitle

\section{Introduction}
This is a sample document demonstrating various \LaTeX{} features including mathematical equations, tables, and formatting. You can cite papers using the citation search feature, and they will be automatically added to your \texttt{cite.bib} file.

\section{Mathematical Expressions}
Here are some example equations:

\subsection{Quadratic Formula}
The quadratic formula is given by:
\[x = \frac{-b \pm \sqrt{b^2 - 4ac}}{2a}\]

\subsection{Matrix Example}
A sample matrix:
\[
\begin{pmatrix}
a_{11} & a_{12} & a_{13} \\
a_{21} & a_{22} & a_{23} \\
a_{31} & a_{32} & a_{33}
\end{pmatrix}
\]

\section{Tables}
Below is a sample table showing data:

\begin{table}[h]
\centering
\begin{tabular}{@{}lrrr@{}}
\toprule
\textbf{Category} & \textbf{Value 1} & \textbf{Value 2} & \textbf{Total} \\
\midrule
Item A & 10.5 & 20.3 & 30.8 \\
Item B & 15.2 & 25.1 & 40.3 \\
Item C & 22.1 & 12.7 & 34.8 \\
\bottomrule
\end{tabular}
\caption{Sample Data Table}
\label{tab:sample}
\end{table}

\section{Lists}
\subsection{Bullet Points}
\begin{itemize}
  \item First item
  \item Second item
  \item Third item with sub-items:
    \begin{itemize}
      \item Sub-item 1
      \item Sub-item 2
    \end{itemize}
\end{itemize}

\subsection{Numbered List}
\begin{enumerate}
  \item First step
  \item Second step
  \item Third step
\end{enumerate}

\section{Citations and References}
When you use the citation search feature, references will be automatically managed. For example, you might cite a paper like this: \cite{example2024}. All citations will be automatically added to your \texttt{cite.bib} file.

\section{Conclusion}
This template demonstrates basic \LaTeX{} features including mathematical formulas, tables, and lists. The bibliography system is now set up to work seamlessly with Bibby AI's citation management. Feel free to modify it for your needs!

\bibliographystyle{plain}
\bibliography{cite}

\end{document}\documentclass{article}
\usepackage[utf8]{inputenc}
\usepackage{booktabs}
\usepackage{amsmath}
\usepackage{natbib}

\title{My Document}
\author{Your Name}
\date{\today}

\begin{document}

\maketitle

\section{Introduction}
This is a sample document demonstrating various \LaTeX{} features including mathematical equations, tables, and formatting. You can cite papers using the citation search feature, and they will be automatically added to your \texttt{cite.bib} file.

\section{Mathematical Expressions}
Here are some example equations:

\subsection{Quadratic Formula}
The quadratic formula is given by:
\[x = \frac{-b \pm \sqrt{b^2 - 4ac}}{2a}\]

\subsection{Matrix Example}
A sample matrix:
\[
\begin{pmatrix}
a_{11} & a_{12} & a_{13} \\
a_{21} & a_{22} & a_{23} \\
a_{31} & a_{32} & a_{33}
\end{pmatrix}
\]

\section{Tables}
Below is a sample table showing data:

\begin{table}[h]
\centering
\begin{tabular}{@{}lrrr@{}}
\toprule
\textbf{Category} & \textbf{Value 1} & \textbf{Value 2} & \textbf{Total} \\
\midrule
Item A & 10.5 & 20.3 & 30.8 \\
Item B & 15.2 & 25.1 & 40.3 \\
Item C & 22.1 & 12.7 & 34.8 \\
\bottomrule
\end{tabular}
\caption{Sample Data Table}
\label{tab:sample}
\end{table}

\section{Lists}
\subsection{Bullet Points}
\begin{itemize}
  \item First item
  \item Second item
  \item Third item with sub-items:
    \begin{itemize}
      \item Sub-item 1
      \item Sub-item 2
    \end{itemize}
\end{itemize}

\subsection{Numbered List}
\begin{enumerate}
  \item First step
  \item Second step
  \item Third step
\end{enumerate}

\section{Citations and References}
When you use the citation search feature, references will be automatically managed. For example, you might cite a paper like this: \cite{example2024}. All citations will be automatically added to your \texttt{cite.bib} file.

\section{Conclusion}
This template demonstrates basic \LaTeX{} features including mathematical formulas, tables, and lists. The bibliography system is now set up to work seamlessly with Bibby AI's citation management. Feel free to modify it for your needs!

\bibliographystyle{plain}
\bibliography{cite}

\end{document}\documentclass{article}
\usepackage[utf8]{inputenc}
\usepackage{booktabs}
\usepackage{amsmath}
\usepackage{natbib}

\title{My Document}
\author{Your Name}
\date{\today}

\begin{document}

\maketitle

\section{Introduction}
This is a sample document demonstrating various \LaTeX{} features including mathematical equations, tables, and formatting. You can cite papers using the citation search feature, and they will be automatically added to your \texttt{cite.bib} file.

\section{Mathematical Expressions}
Here are some example equations:

\subsection{Quadratic Formula}
The quadratic formula is given by:
\[x = \frac{-b \pm \sqrt{b^2 - 4ac}}{2a}\]

\subsection{Matrix Example}
A sample matrix:
\[
\begin{pmatrix}
a_{11} & a_{12} & a_{13} \\
a_{21} & a_{22} & a_{23} \\
a_{31} & a_{32} & a_{33}
\end{pmatrix}
\]

\section{Tables}
Below is a sample table showing data:

\begin{table}[h]
\centering
\begin{tabular}{@{}lrrr@{}}
\toprule
\textbf{Category} & \textbf{Value 1} & \textbf{Value 2} & \textbf{Total} \\
\midrule
Item A & 10.5 & 20.3 & 30.8 \\
Item B & 15.2 & 25.1 & 40.3 \\
Item C & 22.1 & 12.7 & 34.8 \\
\bottomrule
\end{tabular}
\caption{Sample Data Table}
\label{tab:sample}
\end{table}

\section{Lists}
\subsection{Bullet Points}
\begin{itemize}
  \item First item
  \item Second item
  \item Third item with sub-items:
    \begin{itemize}
      \item Sub-item 1
      \item Sub-item 2
    \end{itemize}
\end{itemize}

\subsection{Numbered List}
\begin{enumerate}
  \item First step
  \item Second step
  \item Third step
\end{enumerate}

\section{Citations and References}
When you use the citation search feature, references will be automatically managed. For example, you might cite a paper like this: \cite{example2024}. All citations will be automatically added to your \texttt{cite.bib} file.

\section{Conclusion}
This template demonstrates basic \LaTeX{} features including mathematical formulas, tables, and lists. The bibliography system is now set up to work seamlessly with Bibby AI's citation management. Feel free to modify it for your needs!

\bibliographystyle{plain}
\bibliography{cite}

\end{document}\documentclass{article}
\usepackage[utf8]{inputenc}
\usepackage{booktabs}
\usepackage{amsmath}
\usepackage{natbib}

\title{My Document}
\author{Your Name}
\date{\today}

\begin{document}

\maketitle

\section{Introduction}
This is a sample document demonstrating various \LaTeX{} features including mathematical equations, tables, and formatting. You can cite papers using the citation search feature, and they will be automatically added to your \texttt{cite.bib} file.

\section{Mathematical Expressions}
Here are some example equations:

\subsection{Quadratic Formula}
The quadratic formula is given by:
\[x = \frac{-b \pm \sqrt{b^2 - 4ac}}{2a}\]

\subsection{Matrix Example}
A sample matrix:
\[
\begin{pmatrix}
a_{11} & a_{12} & a_{13} \\
a_{21} & a_{22} & a_{23} \\
a_{31} & a_{32} & a_{33}
\end{pmatrix}
\]

\section{Tables}
Below is a sample table showing data:

\begin{table}[h]
\centering
\begin{tabular}{@{}lrrr@{}}
\toprule
\textbf{Category} & \textbf{Value 1} & \textbf{Value 2} & \textbf{Total} \\
\midrule
Item A & 10.5 & 20.3 & 30.8 \\
Item B & 15.2 & 25.1 & 40.3 \\
Item C & 22.1 & 12.7 & 34.8 \\
\bottomrule
\end{tabular}
\caption{Sample Data Table}
\label{tab:sample}
\end{table}

\section{Lists}
\subsection{Bullet Points}
\begin{itemize}
  \item First item
  \item Second item
  \item Third item with sub-items:
    \begin{itemize}
      \item Sub-item 1
      \item Sub-item 2
    \end{itemize}
\end{itemize}

\subsection{Numbered List}
\begin{enumerate}
  \item First step
  \item Second step
  \item Third step
\end{enumerate}

\section{Citations and References}
When you use the citation search feature, references will be automatically managed. For example, you might cite a paper like this: \cite{example2024}. All citations will be automatically added to your \texttt{cite.bib} file.

\section{Conclusion}
This template demonstrates basic \LaTeX{} features including mathematical formulas, tables, and lists. The bibliography system is now set up to work seamlessly with Bibby AI's citation management. Feel free to modify it for your needs!

\bibliographystyle{plain}
\bibliography{cite}

\end{document}\documentclass{article}
\usepackage[utf8]{inputenc}
\usepackage{booktabs}
\usepackage{amsmath}
\usepackage{natbib}

\title{My Document}
\author{Your Name}
\date{\today}

\begin{document}

\maketitle

\section{Introduction}
This is a sample document demonstrating various \LaTeX{} features including mathematical equations, tables, and formatting. You can cite papers using the citation search feature, and they will be automatically added to your \texttt{cite.bib} file.

\section{Mathematical Expressions}
Here are some example equations:

\subsection{Quadratic Formula}
The quadratic formula is given by:
\[x = \frac{-b \pm \sqrt{b^2 - 4ac}}{2a}\]

\subsection{Matrix Example}
A sample matrix:
\[
\begin{pmatrix}
a_{11} & a_{12} & a_{13} \\
a_{21} & a_{22} & a_{23} \\
a_{31} & a_{32} & a_{33}
\end{pmatrix}
\]

\section{Tables}
Below is a sample table showing data:

\begin{table}[h]
\centering
\begin{tabular}{@{}lrrr@{}}
\toprule
\textbf{Category} & \textbf{Value 1} & \textbf{Value 2} & \textbf{Total} \\
\midrule
Item A & 10.5 & 20.3 & 30.8 \\
Item B & 15.2 & 25.1 & 40.3 \\
Item C & 22.1 & 12.7 & 34.8 \\
\bottomrule
\end{tabular}
\caption{Sample Data Table}
\label{tab:sample}
\end{table}

\section{Lists}
\subsection{Bullet Points}
\begin{itemize}
  \item First item
  \item Second item
  \item Third item with sub-items:
    \begin{itemize}
      \item Sub-item 1
      \item Sub-item 2
    \end{itemize}
\end{itemize}

\subsection{Numbered List}
\begin{enumerate}
  \item First step
  \item Second step
  \item Third step
\end{enumerate}

\section{Citations and References}
When you use the citation search feature, references will be automatically managed. For example, you might cite a paper like this: \cite{example2024}. All citations will be automatically added to your \texttt{cite.bib} file.

\section{Conclusion}
This template demonstrates basic \LaTeX{} features including mathematical formulas, tables, and lists. The bibliography system is now set up to work seamlessly with Bibby AI's citation management. Feel free to modify it for your needs!

\bibliographystyle{plain}
\bibliography{cite}

\end{document}\documentclass{article}
\usepackage[utf8]{inputenc}
\usepackage{booktabs}
\usepackage{amsmath}
\usepackage{natbib}

\title{My Document}
\author{Your Name}
\date{\today}

\begin{document}

\maketitle

\section{Introduction}
This is a sample document demonstrating various \LaTeX{} features including mathematical equations, tables, and formatting. You can cite papers using the citation search feature, and they will be automatically added to your \texttt{cite.bib} file.

\section{Mathematical Expressions}
Here are some example equations:

\subsection{Quadratic Formula}
The quadratic formula is given by:
\[x = \frac{-b \pm \sqrt{b^2 - 4ac}}{2a}\]

\subsection{Matrix Example}
A sample matrix:
\[
\begin{pmatrix}
a_{11} & a_{12} & a_{13} \\
a_{21} & a_{22} & a_{23} \\
a_{31} & a_{32} & a_{33}
\end{pmatrix}
\]

\section{Tables}
Below is a sample table showing data:

\begin{table}[h]
\centering
\begin{tabular}{@{}lrrr@{}}
\toprule
\textbf{Category} & \textbf{Value 1} & \textbf{Value 2} & \textbf{Total} \\
\midrule
Item A & 10.5 & 20.3 & 30.8 \\
Item B & 15.2 & 25.1 & 40.3 \\
Item C & 22.1 & 12.7 & 34.8 \\
\bottomrule
\end{tabular}
\caption{Sample Data Table}
\label{tab:sample}
\end{table}

\section{Lists}
\subsection{Bullet Points}
\begin{itemize}
  \item First item
  \item Second item
  \item Third item with sub-items:
    \begin{itemize}
      \item Sub-item 1
      \item Sub-item 2
    \end{itemize}
\end{itemize}

\subsection{Numbered List}
\begin{enumerate}
  \item First step
  \item Second step
  \item Third step
\end{enumerate}

\section{Citations and References}
When you use the citation search feature, references will be automatically managed. For example, you might cite a paper like this: \cite{example2024}. All citations will be automatically added to your \texttt{cite.bib} file.

\section{Conclusion}
This template demonstrates basic \LaTeX{} features including mathematical formulas, tables, and lists. The bibliography system is now set up to work seamlessly with Bibby AI's citation management. Feel free to modify it for your needs!

\bibliographystyle{plain}
\bibliography{cite}

\end{document}\documentclass{article}
\usepackage[utf8]{inputenc}
\usepackage{booktabs}
\usepackage{amsmath}
\usepackage{natbib}

\title{My Document}
\author{Your Name}
\date{\today}

\begin{document}

\maketitle

\section{Introduction}
This is a sample document demonstrating various \LaTeX{} features including mathematical equations, tables, and formatting. You can cite papers using the citation search feature, and they will be automatically added to your \texttt{cite.bib} file.

\section{Mathematical Expressions}
Here are some example equations:

\subsection{Quadratic Formula}
The quadratic formula is given by:
\[x = \frac{-b \pm \sqrt{b^2 - 4ac}}{2a}\]

\subsection{Matrix Example}
A sample matrix:
\[
\begin{pmatrix}
a_{11} & a_{12} & a_{13} \\
a_{21} & a_{22} & a_{23} \\
a_{31} & a_{32} & a_{33}
\end{pmatrix}
\]

\section{Tables}
Below is a sample table showing data:

\begin{table}[h]
\centering
\begin{tabular}{@{}lrrr@{}}
\toprule
\textbf{Category} & \textbf{Value 1} & \textbf{Value 2} & \textbf{Total} \\
\midrule
Item A & 10.5 & 20.3 & 30.8 \\
Item B & 15.2 & 25.1 & 40.3 \\
Item C & 22.1 & 12.7 & 34.8 \\
\bottomrule
\end{tabular}
\caption{Sample Data Table}
\label{tab:sample}
\end{table}

\section{Lists}
\subsection{Bullet Points}
\begin{itemize}
  \item First item
  \item Second item
  \item Third item with sub-items:
    \begin{itemize}
      \item Sub-item 1
      \item Sub-item 2
    \end{itemize}
\end{itemize}

\subsection{Numbered List}
\begin{enumerate}
  \item First step
  \item Second step
  \item Third step
\end{enumerate}

\section{Citations and References}
When you use the citation search feature, references will be automatically managed. For example, you might cite a paper like this: \cite{example2024}. All citations will be automatically added to your \texttt{cite.bib} file.

\section{Conclusion}
This template demonstrates basic \LaTeX{} features including mathematical formulas, tables, and lists. The bibliography system is now set up to work seamlessly with Bibby AI's citation management. Feel free to modify it for your needs!

\bibliographystyle{plain}
\bibliography{cite}

\end{document}\documentclass{article}
\usepackage[utf8]{inputenc}
\usepackage{booktabs}
\usepackage{amsmath}
\usepackage{natbib}

\title{My Document}
\author{Your Name}
\date{\today}

\begin{document}

\maketitle

\section{Introduction}
This is a sample document demonstrating various \LaTeX{} features including mathematical equations, tables, and formatting. You can cite papers using the citation search feature, and they will be automatically added to your \texttt{cite.bib} file.

\section{Mathematical Expressions}
Here are some example equations:

\subsection{Quadratic Formula}
The quadratic formula is given by:
\[x = \frac{-b \pm \sqrt{b^2 - 4ac}}{2a}\]

\subsection{Matrix Example}
A sample matrix:
\[
\begin{pmatrix}
a_{11} & a_{12} & a_{13} \\
a_{21} & a_{22} & a_{23} \\
a_{31} & a_{32} & a_{33}
\end{pmatrix}
\]

\section{Tables}
Below is a sample table showing data:

\begin{table}[h]
\centering
\begin{tabular}{@{}lrrr@{}}
\toprule
\textbf{Category} & \textbf{Value 1} & \textbf{Value 2} & \textbf{Total} \\
\midrule
Item A & 10.5 & 20.3 & 30.8 \\
Item B & 15.2 & 25.1 & 40.3 \\
Item C & 22.1 & 12.7 & 34.8 \\
\bottomrule
\end{tabular}
\caption{Sample Data Table}
\label{tab:sample}
\end{table}

\section{Lists}
\subsection{Bullet Points}
\begin{itemize}
  \item First item
  \item Second item
  \item Third item with sub-items:
    \begin{itemize}
      \item Sub-item 1
      \item Sub-item 2
    \end{itemize}
\end{itemize}

\subsection{Numbered List}
\begin{enumerate}
  \item First step
  \item Second step
  \item Third step
\end{enumerate}

\section{Citations and References}
When you use the citation search feature, references will be automatically managed. For example, you might cite a paper like this: \cite{example2024}. All citations will be automatically added to your \texttt{cite.bib} file.

\section{Conclusion}
This template demonstrates basic \LaTeX{} features including mathematical formulas, tables, and lists. The bibliography system is now set up to work seamlessly with Bibby AI's citation management. Feel free to modify it for your needs!

\bibliographystyle{plain}
\bibliography{cite}

\end{document}\documentclass{article}
\usepackage[utf8]{inputenc}
\usepackage{booktabs}
\usepackage{amsmath}
\usepackage{natbib}

\title{My Document}
\author{Your Name}
\date{\today}

\begin{document}

\maketitle

\section{Introduction}
This is a sample document demonstrating various \LaTeX{} features including mathematical equations, tables, and formatting. You can cite papers using the citation search feature, and they will be automatically added to your \texttt{cite.bib} file.

\section{Mathematical Expressions}
Here are some example equations:

\subsection{Quadratic Formula}
The quadratic formula is given by:
\[x = \frac{-b \pm \sqrt{b^2 - 4ac}}{2a}\]

\subsection{Matrix Example}
A sample matrix:
\[
\begin{pmatrix}
a_{11} & a_{12} & a_{13} \\
a_{21} & a_{22} & a_{23} \\
a_{31} & a_{32} & a_{33}
\end{pmatrix}
\]

\section{Tables}
Below is a sample table showing data:

\begin{table}[h]
\centering
\begin{tabular}{@{}lrrr@{}}
\toprule
\textbf{Category} & \textbf{Value 1} & \textbf{Value 2} & \textbf{Total} \\
\midrule
Item A & 10.5 & 20.3 & 30.8 \\
Item B & 15.2 & 25.1 & 40.3 \\
Item C & 22.1 & 12.7 & 34.8 \\
\bottomrule
\end{tabular}
\caption{Sample Data Table}
\label{tab:sample}
\end{table}

\section{Lists}
\subsection{Bullet Points}
\begin{itemize}
  \item First item
  \item Second item
  \item Third item with sub-items:
    \begin{itemize}
      \item Sub-item 1
      \item Sub-item 2
    \end{itemize}
\end{itemize}

\subsection{Numbered List}
\begin{enumerate}
  \item First step
  \item Second step
  \item Third step
\end{enumerate}

\section{Citations and References}
When you use the citation search feature, references will be automatically managed. For example, you might cite a paper like this: \cite{example2024}. All citations will be automatically added to your \texttt{cite.bib} file.

\section{Conclusion}
This template demonstrates basic \LaTeX{} features including mathematical formulas, tables, and lists. The bibliography system is now set up to work seamlessly with Bibby AI's citation management. Feel free to modify it for your needs!

\bibliographystyle{plain}
\bibliography{cite}

\end{document}\documentclass{article}
\usepackage[utf8]{inputenc}
\usepackage{booktabs}
\usepackage{amsmath}
\usepackage{natbib}

\title{My Document}
\author{Your Name}
\date{\today}

\begin{document}

\maketitle

\section{Introduction}
This is a sample document demonstrating various \LaTeX{} features including mathematical equations, tables, and formatting. You can cite papers using the citation search feature, and they will be automatically added to your \texttt{cite.bib} file.

\section{Mathematical Expressions}
Here are some example equations:

\subsection{Quadratic Formula}
The quadratic formula is given by:
\[x = \frac{-b \pm \sqrt{b^2 - 4ac}}{2a}\]

\subsection{Matrix Example}
A sample matrix:
\[
\begin{pmatrix}
a_{11} & a_{12} & a_{13} \\
a_{21} & a_{22} & a_{23} \\
a_{31} & a_{32} & a_{33}
\end{pmatrix}
\]

\section{Tables}
Below is a sample table showing data:

\begin{table}[h]
\centering
\begin{tabular}{@{}lrrr@{}}
\toprule
\textbf{Category} & \textbf{Value 1} & \textbf{Value 2} & \textbf{Total} \\
\midrule
Item A & 10.5 & 20.3 & 30.8 \\
Item B & 15.2 & 25.1 & 40.3 \\
Item C & 22.1 & 12.7 & 34.8 \\
\bottomrule
\end{tabular}
\caption{Sample Data Table}
\label{tab:sample}
\end{table}

\section{Lists}
\subsection{Bullet Points}
\begin{itemize}
  \item First item
  \item Second item
  \item Third item with sub-items:
    \begin{itemize}
      \item Sub-item 1
      \item Sub-item 2
    \end{itemize}
\end{itemize}

\subsection{Numbered List}
\begin{enumerate}
  \item First step
  \item Second step
  \item Third step
\end{enumerate}

\section{Citations and References}
When you use the citation search feature, references will be automatically managed. For example, you might cite a paper like this: \cite{example2024}. All citations will be automatically added to your \texttt{cite.bib} file.

\section{Conclusion}
This template demonstrates basic \LaTeX{} features including mathematical formulas, tables, and lists. The bibliography system is now set up to work seamlessly with Bibby AI's citation management. Feel free to modify it for your needs!

\bibliographystyle{plain}
\bibliography{cite}

\end{document}\documentclass{article}
\usepackage[utf8]{inputenc}
\usepackage{booktabs}
\usepackage{amsmath}
\usepackage{natbib}

\title{My Document}
\author{Your Name}
\date{\today}

\begin{document}

\maketitle

\section{Introduction}
This is a sample document demonstrating various \LaTeX{} features including mathematical equations, tables, and formatting. You can cite papers using the citation search feature, and they will be automatically added to your \texttt{cite.bib} file.

\section{Mathematical Expressions}
Here are some example equations:

\subsection{Quadratic Formula}
The quadratic formula is given by:
\[x = \frac{-b \pm \sqrt{b^2 - 4ac}}{2a}\]

\subsection{Matrix Example}
A sample matrix:
\[
\begin{pmatrix}
a_{11} & a_{12} & a_{13} \\
a_{21} & a_{22} & a_{23} \\
a_{31} & a_{32} & a_{33}
\end{pmatrix}
\]

\section{Tables}
Below is a sample table showing data:

\begin{table}[h]
\centering
\begin{tabular}{@{}lrrr@{}}
\toprule
\textbf{Category} & \textbf{Value 1} & \textbf{Value 2} & \textbf{Total} \\
\midrule
Item A & 10.5 & 20.3 & 30.8 \\
Item B & 15.2 & 25.1 & 40.3 \\
Item C & 22.1 & 12.7 & 34.8 \\
\bottomrule
\end{tabular}
\caption{Sample Data Table}
\label{tab:sample}
\end{table}

\section{Lists}
\subsection{Bullet Points}
\begin{itemize}
  \item First item
  \item Second item
  \item Third item with sub-items:
    \begin{itemize}
      \item Sub-item 1
      \item Sub-item 2
    \end{itemize}
\end{itemize}

\subsection{Numbered List}
\begin{enumerate}
  \item First step
  \item Second step
  \item Third step
\end{enumerate}

\section{Citations and References}
When you use the citation search feature, references will be automatically managed. For example, you might cite a paper like this: \cite{example2024}. All citations will be automatically added to your \texttt{cite.bib} file.

\section{Conclusion}
This template demonstrates basic \LaTeX{} features including mathematical formulas, tables, and lists. The bibliography system is now set up to work seamlessly with Bibby AI's citation management. Feel free to modify it for your needs!

\bibliographystyle{plain}
\bibliography{cite}

\end{document}\documentclass{article}
\usepackage[utf8]{inputenc}
\usepackage{booktabs}
\usepackage{amsmath}
\usepackage{natbib}

\title{My Document}
\author{Your Name}
\date{\today}

\begin{document}

\maketitle

\section{Introduction}
This is a sample document demonstrating various \LaTeX{} features including mathematical equations, tables, and formatting. You can cite papers using the citation search feature, and they will be automatically added to your \texttt{cite.bib} file.

\section{Mathematical Expressions}
Here are some example equations:

\subsection{Quadratic Formula}
The quadratic formula is given by:
\[x = \frac{-b \pm \sqrt{b^2 - 4ac}}{2a}\]

\subsection{Matrix Example}
A sample matrix:
\[
\begin{pmatrix}
a_{11} & a_{12} & a_{13} \\
a_{21} & a_{22} & a_{23} \\
a_{31} & a_{32} & a_{33}
\end{pmatrix}
\]

\section{Tables}
Below is a sample table showing data:

\begin{table}[h]
\centering
\begin{tabular}{@{}lrrr@{}}
\toprule
\textbf{Category} & \textbf{Value 1} & \textbf{Value 2} & \textbf{Total} \\
\midrule
Item A & 10.5 & 20.3 & 30.8 \\
Item B & 15.2 & 25.1 & 40.3 \\
Item C & 22.1 & 12.7 & 34.8 \\
\bottomrule
\end{tabular}
\caption{Sample Data Table}
\label{tab:sample}
\end{table}

\section{Lists}
\subsection{Bullet Points}
\begin{itemize}
  \item First item
  \item Second item
  \item Third item with sub-items:
    \begin{itemize}
      \item Sub-item 1
      \item Sub-item 2
    \end{itemize}
\end{itemize}

\subsection{Numbered List}
\begin{enumerate}
  \item First step
  \item Second step
  \item Third step
\end{enumerate}

\section{Citations and References}
When you use the citation search feature, references will be automatically managed. For example, you might cite a paper like this: \cite{example2024}. All citations will be automatically added to your \texttt{cite.bib} file.

\section{Conclusion}
This template demonstrates basic \LaTeX{} features including mathematical formulas, tables, and lists. The bibliography system is now set up to work seamlessly with Bibby AI's citation management. Feel free to modify it for your needs!

\bibliographystyle{plain}
\bibliography{cite}

\end{document}\documentclass{article}
\usepackage[utf8]{inputenc}
\usepackage{booktabs}
\usepackage{amsmath}
\usepackage{natbib}

\title{My Document}
\author{Your Name}
\date{\today}

\begin{document}

\maketitle

\section{Introduction}
This is a sample document demonstrating various \LaTeX{} features including mathematical equations, tables, and formatting. You can cite papers using the citation search feature, and they will be automatically added to your \texttt{cite.bib} file.

\section{Mathematical Expressions}
Here are some example equations:

\subsection{Quadratic Formula}
The quadratic formula is given by:
\[x = \frac{-b \pm \sqrt{b^2 - 4ac}}{2a}\]

\subsection{Matrix Example}
A sample matrix:
\[
\begin{pmatrix}
a_{11} & a_{12} & a_{13} \\
a_{21} & a_{22} & a_{23} \\
a_{31} & a_{32} & a_{33}
\end{pmatrix}
\]

\section{Tables}
Below is a sample table showing data:

\begin{table}[h]
\centering
\begin{tabular}{@{}lrrr@{}}
\toprule
\textbf{Category} & \textbf{Value 1} & \textbf{Value 2} & \textbf{Total} \\
\midrule
Item A & 10.5 & 20.3 & 30.8 \\
Item B & 15.2 & 25.1 & 40.3 \\
Item C & 22.1 & 12.7 & 34.8 \\
\bottomrule
\end{tabular}
\caption{Sample Data Table}
\label{tab:sample}
\end{table}

\section{Lists}
\subsection{Bullet Points}
\begin{itemize}
  \item First item
  \item Second item
  \item Third item with sub-items:
    \begin{itemize}
      \item Sub-item 1
      \item Sub-item 2
    \end{itemize}
\end{itemize}

\subsection{Numbered List}
\begin{enumerate}
  \item First step
  \item Second step
  \item Third step
\end{enumerate}

\section{Citations and References}
When you use the citation search feature, references will be automatically managed. For example, you might cite a paper like this: \cite{example2024}. All citations will be automatically added to your \texttt{cite.bib} file.

\section{Conclusion}
This template demonstrates basic \LaTeX{} features including mathematical formulas, tables, and lists. The bibliography system is now set up to work seamlessly with Bibby AI's citation management. Feel free to modify it for your needs!

\bibliographystyle{plain}
\bibliography{cite}

\end{document}\documentclass{article}
\usepackage[utf8]{inputenc}
\usepackage{booktabs}
\usepackage{amsmath}
\usepackage{natbib}

\title{My Document}
\author{Your Name}
\date{\today}

\begin{document}

\maketitle

\section{Introduction}
This is a sample document demonstrating various \LaTeX{} features including mathematical equations, tables, and formatting. You can cite papers using the citation search feature, and they will be automatically added to your \texttt{cite.bib} file.

\section{Mathematical Expressions}
Here are some example equations:

\subsection{Quadratic Formula}
The quadratic formula is given by:
\[x = \frac{-b \pm \sqrt{b^2 - 4ac}}{2a}\]

\subsection{Matrix Example}
A sample matrix:
\[
\begin{pmatrix}
a_{11} & a_{12} & a_{13} \\
a_{21} & a_{22} & a_{23} \\
a_{31} & a_{32} & a_{33}
\end{pmatrix}
\]

\section{Tables}
Below is a sample table showing data:

\begin{table}[h]
\centering
\begin{tabular}{@{}lrrr@{}}
\toprule
\textbf{Category} & \textbf{Value 1} & \textbf{Value 2} & \textbf{Total} \\
\midrule
Item A & 10.5 & 20.3 & 30.8 \\
Item B & 15.2 & 25.1 & 40.3 \\
Item C & 22.1 & 12.7 & 34.8 \\
\bottomrule
\end{tabular}
\caption{Sample Data Table}
\label{tab:sample}
\end{table}

\section{Lists}
\subsection{Bullet Points}
\begin{itemize}
  \item First item
  \item Second item
  \item Third item with sub-items:
    \begin{itemize}
      \item Sub-item 1
      \item Sub-item 2
    \end{itemize}
\end{itemize}

\subsection{Numbered List}
\begin{enumerate}
  \item First step
  \item Second step
  \item Third step
\end{enumerate}

\section{Citations and References}
When you use the citation search feature, references will be automatically managed. For example, you might cite a paper like this: \cite{example2024}. All citations will be automatically added to your \texttt{cite.bib} file.

\section{Conclusion}
This template demonstrates basic \LaTeX{} features including mathematical formulas, tables, and lists. The bibliography system is now set up to work seamlessly with Bibby AI's citation management. Feel free to modify it for your needs!

\bibliographystyle{plain}
\bibliography{cite}

\end{document}\documentclass{article}
\usepackage[utf8]{inputenc}
\usepackage{booktabs}
\usepackage{amsmath}
\usepackage{natbib}

\title{My Document}
\author{Your Name}
\date{\today}

\begin{document}

\maketitle

\section{Introduction}
This is a sample document demonstrating various \LaTeX{} features including mathematical equations, tables, and formatting. You can cite papers using the citation search feature, and they will be automatically added to your \texttt{cite.bib} file.

\section{Mathematical Expressions}
Here are some example equations:

\subsection{Quadratic Formula}
The quadratic formula is given by:
\[x = \frac{-b \pm \sqrt{b^2 - 4ac}}{2a}\]

\subsection{Matrix Example}
A sample matrix:
\[
\begin{pmatrix}
a_{11} & a_{12} & a_{13} \\
a_{21} & a_{22} & a_{23} \\
a_{31} & a_{32} & a_{33}
\end{pmatrix}
\]

\section{Tables}
Below is a sample table showing data:

\begin{table}[h]
\centering
\begin{tabular}{@{}lrrr@{}}
\toprule
\textbf{Category} & \textbf{Value 1} & \textbf{Value 2} & \textbf{Total} \\
\midrule
Item A & 10.5 & 20.3 & 30.8 \\
Item B & 15.2 & 25.1 & 40.3 \\
Item C & 22.1 & 12.7 & 34.8 \\
\bottomrule
\end{tabular}
\caption{Sample Data Table}
\label{tab:sample}
\end{table}

\section{Lists}
\subsection{Bullet Points}
\begin{itemize}
  \item First item
  \item Second item
  \item Third item with sub-items:
    \begin{itemize}
      \item Sub-item 1
      \item Sub-item 2
    \end{itemize}
\end{itemize}

\subsection{Numbered List}
\begin{enumerate}
  \item First step
  \item Second step
  \item Third step
\end{enumerate}

\section{Citations and References}
When you use the citation search feature, references will be automatically managed. For example, you might cite a paper like this: \cite{example2024}. All citations will be automatically added to your \texttt{cite.bib} file.

\section{Conclusion}
This template demonstrates basic \LaTeX{} features including mathematical formulas, tables, and lists. The bibliography system is now set up to work seamlessly with Bibby AI's citation management. Feel free to modify it for your needs!

\bibliographystyle{plain}
\bibliography{cite}

\end{document}\documentclass{article}
\usepackage[utf8]{inputenc}
\usepackage{booktabs}
\usepackage{amsmath}
\usepackage{natbib}

\title{My Document}
\author{Your Name}
\date{\today}

\begin{document}

\maketitle

\section{Introduction}
This is a sample document demonstrating various \LaTeX{} features including mathematical equations, tables, and formatting. You can cite papers using the citation search feature, and they will be automatically added to your \texttt{cite.bib} file.

\section{Mathematical Expressions}
Here are some example equations:

\subsection{Quadratic Formula}
The quadratic formula is given by:
\[x = \frac{-b \pm \sqrt{b^2 - 4ac}}{2a}\]

\subsection{Matrix Example}
A sample matrix:
\[
\begin{pmatrix}
a_{11} & a_{12} & a_{13} \\
a_{21} & a_{22} & a_{23} \\
a_{31} & a_{32} & a_{33}
\end{pmatrix}
\]

\section{Tables}
Below is a sample table showing data:

\begin{table}[h]
\centering
\begin{tabular}{@{}lrrr@{}}
\toprule
\textbf{Category} & \textbf{Value 1} & \textbf{Value 2} & \textbf{Total} \\
\midrule
Item A & 10.5 & 20.3 & 30.8 \\
Item B & 15.2 & 25.1 & 40.3 \\
Item C & 22.1 & 12.7 & 34.8 \\
\bottomrule
\end{tabular}
\caption{Sample Data Table}
\label{tab:sample}
\end{table}

\section{Lists}
\subsection{Bullet Points}
\begin{itemize}
  \item First item
  \item Second item
  \item Third item with sub-items:
    \begin{itemize}
      \item Sub-item 1
      \item Sub-item 2
    \end{itemize}
\end{itemize}

\subsection{Numbered List}
\begin{enumerate}
  \item First step
  \item Second step
  \item Third step
\end{enumerate}

\section{Citations and References}
When you use the citation search feature, references will be automatically managed. For example, you might cite a paper like this: \cite{example2024}. All citations will be automatically added to your \texttt{cite.bib} file.

\section{Conclusion}
This template demonstrates basic \LaTeX{} features including mathematical formulas, tables, and lists. The bibliography system is now set up to work seamlessly with Bibby AI's citation management. Feel free to modify it for your needs!

\bibliographystyle{plain}
\bibliography{cite}

\end{document}\documentclass{article}
\usepackage[utf8]{inputenc}
\usepackage{booktabs}
\usepackage{amsmath}
\usepackage{natbib}

\title{My Document}
\author{Your Name}
\date{\today}

\begin{document}

\maketitle

\section{Introduction}
This is a sample document demonstrating various \LaTeX{} features including mathematical equations, tables, and formatting. You can cite papers using the citation search feature, and they will be automatically added to your \texttt{cite.bib} file.

\section{Mathematical Expressions}
Here are some example equations:

\subsection{Quadratic Formula}
The quadratic formula is given by:
\[x = \frac{-b \pm \sqrt{b^2 - 4ac}}{2a}\]

\subsection{Matrix Example}
A sample matrix:
\[
\begin{pmatrix}
a_{11} & a_{12} & a_{13} \\
a_{21} & a_{22} & a_{23} \\
a_{31} & a_{32} & a_{33}
\end{pmatrix}
\]

\section{Tables}
Below is a sample table showing data:

\begin{table}[h]
\centering
\begin{tabular}{@{}lrrr@{}}
\toprule
\textbf{Category} & \textbf{Value 1} & \textbf{Value 2} & \textbf{Total} \\
\midrule
Item A & 10.5 & 20.3 & 30.8 \\
Item B & 15.2 & 25.1 & 40.3 \\
Item C & 22.1 & 12.7 & 34.8 \\
\bottomrule
\end{tabular}
\caption{Sample Data Table}
\label{tab:sample}
\end{table}

\section{Lists}
\subsection{Bullet Points}
\begin{itemize}
  \item First item
  \item Second item
  \item Third item with sub-items:
    \begin{itemize}
      \item Sub-item 1
      \item Sub-item 2
    \end{itemize}
\end{itemize}

\subsection{Numbered List}
\begin{enumerate}
  \item First step
  \item Second step
  \item Third step
\end{enumerate}

\section{Citations and References}
When you use the citation search feature, references will be automatically managed. For example, you might cite a paper like this: \cite{example2024}. All citations will be automatically added to your \texttt{cite.bib} file.

\section{Conclusion}
This template demonstrates basic \LaTeX{} features including mathematical formulas, tables, and lists. The bibliography system is now set up to work seamlessly with Bibby AI's citation management. Feel free to modify it for your needs!

\bibliographystyle{plain}
\bibliography{cite}

\end{document}\documentclass{article}
\usepackage[utf8]{inputenc}
\usepackage{booktabs}
\usepackage{amsmath}
\usepackage{natbib}

\title{My Document}
\author{Your Name}
\date{\today}

\begin{document}

\maketitle

\section{Introduction}
This is a sample document demonstrating various \LaTeX{} features including mathematical equations, tables, and formatting. You can cite papers using the citation search feature, and they will be automatically added to your \texttt{cite.bib} file.

\section{Mathematical Expressions}
Here are some example equations:

\subsection{Quadratic Formula}
The quadratic formula is given by:
\[x = \frac{-b \pm \sqrt{b^2 - 4ac}}{2a}\]

\subsection{Matrix Example}
A sample matrix:
\[
\begin{pmatrix}
a_{11} & a_{12} & a_{13} \\
a_{21} & a_{22} & a_{23} \\
a_{31} & a_{32} & a_{33}
\end{pmatrix}
\]

\section{Tables}
Below is a sample table showing data:

\begin{table}[h]
\centering
\begin{tabular}{@{}lrrr@{}}
\toprule
\textbf{Category} & \textbf{Value 1} & \textbf{Value 2} & \textbf{Total} \\
\midrule
Item A & 10.5 & 20.3 & 30.8 \\
Item B & 15.2 & 25.1 & 40.3 \\
Item C & 22.1 & 12.7 & 34.8 \\
\bottomrule
\end{tabular}
\caption{Sample Data Table}
\label{tab:sample}
\end{table}

\section{Lists}
\subsection{Bullet Points}
\begin{itemize}
  \item First item
  \item Second item
  \item Third item with sub-items:
    \begin{itemize}
      \item Sub-item 1
      \item Sub-item 2
    \end{itemize}
\end{itemize}

\subsection{Numbered List}
\begin{enumerate}
  \item First step
  \item Second step
  \item Third step
\end{enumerate}

\section{Citations and References}
When you use the citation search feature, references will be automatically managed. For example, you might cite a paper like this: \cite{example2024}. All citations will be automatically added to your \texttt{cite.bib} file.

\section{Conclusion}
This template demonstrates basic \LaTeX{} features including mathematical formulas, tables, and lists. The bibliography system is now set up to work seamlessly with Bibby AI's citation management. Feel free to modify it for your needs!

\bibliographystyle{plain}
\bibliography{cite}

\end{document}\documentclass{article}
\usepackage[utf8]{inputenc}
\usepackage{booktabs}
\usepackage{amsmath}
\usepackage{natbib}

\title{My Document}
\author{Your Name}
\date{\today}

\begin{document}

\maketitle

\section{Introduction}
This is a sample document demonstrating various \LaTeX{} features including mathematical equations, tables, and formatting. You can cite papers using the citation search feature, and they will be automatically added to your \texttt{cite.bib} file.

\section{Mathematical Expressions}
Here are some example equations:

\subsection{Quadratic Formula}
The quadratic formula is given by:
\[x = \frac{-b \pm \sqrt{b^2 - 4ac}}{2a}\]

\subsection{Matrix Example}
A sample matrix:
\[
\begin{pmatrix}
a_{11} & a_{12} & a_{13} \\
a_{21} & a_{22} & a_{23} \\
a_{31} & a_{32} & a_{33}
\end{pmatrix}
\]

\section{Tables}
Below is a sample table showing data:

\begin{table}[h]
\centering
\begin{tabular}{@{}lrrr@{}}
\toprule
\textbf{Category} & \textbf{Value 1} & \textbf{Value 2} & \textbf{Total} \\
\midrule
Item A & 10.5 & 20.3 & 30.8 \\
Item B & 15.2 & 25.1 & 40.3 \\
Item C & 22.1 & 12.7 & 34.8 \\
\bottomrule
\end{tabular}
\caption{Sample Data Table}
\label{tab:sample}
\end{table}

\section{Lists}
\subsection{Bullet Points}
\begin{itemize}
  \item First item
  \item Second item
  \item Third item with sub-items:
    \begin{itemize}
      \item Sub-item 1
      \item Sub-item 2
    \end{itemize}
\end{itemize}

\subsection{Numbered List}
\begin{enumerate}
  \item First step
  \item Second step
  \item Third step
\end{enumerate}

\section{Citations and References}
When you use the citation search feature, references will be automatically managed. For example, you might cite a paper like this: \cite{example2024}. All citations will be automatically added to your \texttt{cite.bib} file.

\section{Conclusion}
This template demonstrates basic \LaTeX{} features including mathematical formulas, tables, and lists. The bibliography system is now set up to work seamlessly with Bibby AI's citation management. Feel free to modify it for your needs!

\bibliographystyle{plain}
\bibliography{cite}

\end{document}\documentclass{article}
\usepackage[utf8]{inputenc}
\usepackage{booktabs}
\usepackage{amsmath}
\usepackage{natbib}

\title{My Document}
\author{Your Name}
\date{\today}

\begin{document}

\maketitle

\section{Introduction}
This is a sample document demonstrating various \LaTeX{} features including mathematical equations, tables, and formatting. You can cite papers using the citation search feature, and they will be automatically added to your \texttt{cite.bib} file.

\section{Mathematical Expressions}
Here are some example equations:

\subsection{Quadratic Formula}
The quadratic formula is given by:
\[x = \frac{-b \pm \sqrt{b^2 - 4ac}}{2a}\]

\subsection{Matrix Example}
A sample matrix:
\[
\begin{pmatrix}
a_{11} & a_{12} & a_{13} \\
a_{21} & a_{22} & a_{23} \\
a_{31} & a_{32} & a_{33}
\end{pmatrix}
\]

\section{Tables}
Below is a sample table showing data:

\begin{table}[h]
\centering
\begin{tabular}{@{}lrrr@{}}
\toprule
\textbf{Category} & \textbf{Value 1} & \textbf{Value 2} & \textbf{Total} \\
\midrule
Item A & 10.5 & 20.3 & 30.8 \\
Item B & 15.2 & 25.1 & 40.3 \\
Item C & 22.1 & 12.7 & 34.8 \\
\bottomrule
\end{tabular}
\caption{Sample Data Table}
\label{tab:sample}
\end{table}

\section{Lists}
\subsection{Bullet Points}
\begin{itemize}
  \item First item
  \item Second item
  \item Third item with sub-items:
    \begin{itemize}
      \item Sub-item 1
      \item Sub-item 2
    \end{itemize}
\end{itemize}

\subsection{Numbered List}
\begin{enumerate}
  \item First step
  \item Second step
  \item Third step
\end{enumerate}

\section{Citations and References}
When you use the citation search feature, references will be automatically managed. For example, you might cite a paper like this: \cite{example2024}. All citations will be automatically added to your \texttt{cite.bib} file.

\section{Conclusion}
This template demonstrates basic \LaTeX{} features including mathematical formulas, tables, and lists. The bibliography system is now set up to work seamlessly with Bibby AI's citation management. Feel free to modify it for your needs!

\bibliographystyle{plain}
\bibliography{cite}

\end{document}\documentclass{article}
\usepackage[utf8]{inputenc}
\usepackage{booktabs}
\usepackage{amsmath}
\usepackage{natbib}

\title{My Document}
\author{Your Name}
\date{\today}

\begin{document}

\maketitle

\section{Introduction}
This is a sample document demonstrating various \LaTeX{} features including mathematical equations, tables, and formatting. You can cite papers using the citation search feature, and they will be automatically added to your \texttt{cite.bib} file.

\section{Mathematical Expressions}
Here are some example equations:

\subsection{Quadratic Formula}
The quadratic formula is given by:
\[x = \frac{-b \pm \sqrt{b^2 - 4ac}}{2a}\]

\subsection{Matrix Example}
A sample matrix:
\[
\begin{pmatrix}
a_{11} & a_{12} & a_{13} \\
a_{21} & a_{22} & a_{23} \\
a_{31} & a_{32} & a_{33}
\end{pmatrix}
\]

\section{Tables}
Below is a sample table showing data:

\begin{table}[h]
\centering
\begin{tabular}{@{}lrrr@{}}
\toprule
\textbf{Category} & \textbf{Value 1} & \textbf{Value 2} & \textbf{Total} \\
\midrule
Item A & 10.5 & 20.3 & 30.8 \\
Item B & 15.2 & 25.1 & 40.3 \\
Item C & 22.1 & 12.7 & 34.8 \\
\bottomrule
\end{tabular}
\caption{Sample Data Table}
\label{tab:sample}
\end{table}

\section{Lists}
\subsection{Bullet Points}
\begin{itemize}
  \item First item
  \item Second item
  \item Third item with sub-items:
    \begin{itemize}
      \item Sub-item 1
      \item Sub-item 2
    \end{itemize}
\end{itemize}

\subsection{Numbered List}
\begin{enumerate}
  \item First step
  \item Second step
  \item Third step
\end{enumerate}

\section{Citations and References}
When you use the citation search feature, references will be automatically managed. For example, you might cite a paper like this: \cite{example2024}. All citations will be automatically added to your \texttt{cite.bib} file.

\section{Conclusion}
This template demonstrates basic \LaTeX{} features including mathematical formulas, tables, and lists. The bibliography system is now set up to work seamlessly with Bibby AI's citation management. Feel free to modify it for your needs!

\bibliographystyle{plain}
\bibliography{cite}

\end{document}\documentclass{article}
\usepackage[utf8]{inputenc}
\usepackage{booktabs}
\usepackage{amsmath}
\usepackage{natbib}

\title{My Document}
\author{Your Name}
\date{\today}

\begin{document}

\maketitle

\section{Introduction}
This is a sample document demonstrating various \LaTeX{} features including mathematical equations, tables, and formatting. You can cite papers using the citation search feature, and they will be automatically added to your \texttt{cite.bib} file.

\section{Mathematical Expressions}
Here are some example equations:

\subsection{Quadratic Formula}
The quadratic formula is given by:
\[x = \frac{-b \pm \sqrt{b^2 - 4ac}}{2a}\]

\subsection{Matrix Example}
A sample matrix:
\[
\begin{pmatrix}
a_{11} & a_{12} & a_{13} \\
a_{21} & a_{22} & a_{23} \\
a_{31} & a_{32} & a_{33}
\end{pmatrix}
\]

\section{Tables}
Below is a sample table showing data:

\begin{table}[h]
\centering
\begin{tabular}{@{}lrrr@{}}
\toprule
\textbf{Category} & \textbf{Value 1} & \textbf{Value 2} & \textbf{Total} \\
\midrule
Item A & 10.5 & 20.3 & 30.8 \\
Item B & 15.2 & 25.1 & 40.3 \\
Item C & 22.1 & 12.7 & 34.8 \\
\bottomrule
\end{tabular}
\caption{Sample Data Table}
\label{tab:sample}
\end{table}

\section{Lists}
\subsection{Bullet Points}
\begin{itemize}
  \item First item
  \item Second item
  \item Third item with sub-items:
    \begin{itemize}
      \item Sub-item 1
      \item Sub-item 2
    \end{itemize}
\end{itemize}

\subsection{Numbered List}
\begin{enumerate}
  \item First step
  \item Second step
  \item Third step
\end{enumerate}

\section{Citations and References}
When you use the citation search feature, references will be automatically managed. For example, you might cite a paper like this: \cite{example2024}. All citations will be automatically added to your \texttt{cite.bib} file.

\section{Conclusion}
This template demonstrates basic \LaTeX{} features including mathematical formulas, tables, and lists. The bibliography system is now set up to work seamlessly with Bibby AI's citation management. Feel free to modify it for your needs!

\bibliographystyle{plain}
\bibliography{cite}

\end{document}\documentclass{article}
\usepackage[utf8]{inputenc}
\usepackage{booktabs}
\usepackage{amsmath}
\usepackage{natbib}

\title{My Document}
\author{Your Name}
\date{\today}

\begin{document}

\maketitle

\section{Introduction}
This is a sample document demonstrating various \LaTeX{} features including mathematical equations, tables, and formatting. You can cite papers using the citation search feature, and they will be automatically added to your \texttt{cite.bib} file.

\section{Mathematical Expressions}
Here are some example equations:

\subsection{Quadratic Formula}
The quadratic formula is given by:
\[x = \frac{-b \pm \sqrt{b^2 - 4ac}}{2a}\]

\subsection{Matrix Example}
A sample matrix:
\[
\begin{pmatrix}
a_{11} & a_{12} & a_{13} \\
a_{21} & a_{22} & a_{23} \\
a_{31} & a_{32} & a_{33}
\end{pmatrix}
\]

\section{Tables}
Below is a sample table showing data:

\begin{table}[h]
\centering
\begin{tabular}{@{}lrrr@{}}
\toprule
\textbf{Category} & \textbf{Value 1} & \textbf{Value 2} & \textbf{Total} \\
\midrule
Item A & 10.5 & 20.3 & 30.8 \\
Item B & 15.2 & 25.1 & 40.3 \\
Item C & 22.1 & 12.7 & 34.8 \\
\bottomrule
\end{tabular}
\caption{Sample Data Table}
\label{tab:sample}
\end{table}

\section{Lists}
\subsection{Bullet Points}
\begin{itemize}
  \item First item
  \item Second item
  \item Third item with sub-items:
    \begin{itemize}
      \item Sub-item 1
      \item Sub-item 2
    \end{itemize}
\end{itemize}

\subsection{Numbered List}
\begin{enumerate}
  \item First step
  \item Second step
  \item Third step
\end{enumerate}

\section{Citations and References}
When you use the citation search feature, references will be automatically managed. For example, you might cite a paper like this: \cite{example2024}. All citations will be automatically added to your \texttt{cite.bib} file.

\section{Conclusion}
This template demonstrates basic \LaTeX{} features including mathematical formulas, tables, and lists. The bibliography system is now set up to work seamlessly with Bibby AI's citation management. Feel free to modify it for your needs!

\bibliographystyle{plain}
\bibliography{cite}

\end{document}\documentclass{article}
\usepackage[utf8]{inputenc}
\usepackage{booktabs}
\usepackage{amsmath}
\usepackage{natbib}

\title{My Document}
\author{Your Name}
\date{\today}

\begin{document}

\maketitle

\section{Introduction}
This is a sample document demonstrating various \LaTeX{} features including mathematical equations, tables, and formatting. You can cite papers using the citation search feature, and they will be automatically added to your \texttt{cite.bib} file.

\section{Mathematical Expressions}
Here are some example equations:

\subsection{Quadratic Formula}
The quadratic formula is given by:
\[x = \frac{-b \pm \sqrt{b^2 - 4ac}}{2a}\]

\subsection{Matrix Example}
A sample matrix:
\[
\begin{pmatrix}
a_{11} & a_{12} & a_{13} \\
a_{21} & a_{22} & a_{23} \\
a_{31} & a_{32} & a_{33}
\end{pmatrix}
\]

\section{Tables}
Below is a sample table showing data:

\begin{table}[h]
\centering
\begin{tabular}{@{}lrrr@{}}
\toprule
\textbf{Category} & \textbf{Value 1} & \textbf{Value 2} & \textbf{Total} \\
\midrule
Item A & 10.5 & 20.3 & 30.8 \\
Item B & 15.2 & 25.1 & 40.3 \\
Item C & 22.1 & 12.7 & 34.8 \\
\bottomrule
\end{tabular}
\caption{Sample Data Table}
\label{tab:sample}
\end{table}

\section{Lists}
\subsection{Bullet Points}
\begin{itemize}
  \item First item
  \item Second item
  \item Third item with sub-items:
    \begin{itemize}
      \item Sub-item 1
      \item Sub-item 2
    \end{itemize}
\end{itemize}

\subsection{Numbered List}
\begin{enumerate}
  \item First step
  \item Second step
  \item Third step
\end{enumerate}

\section{Citations and References}
When you use the citation search feature, references will be automatically managed. For example, you might cite a paper like this: \cite{example2024}. All citations will be automatically added to your \texttt{cite.bib} file.

\section{Conclusion}
This template demonstrates basic \LaTeX{} features including mathematical formulas, tables, and lists. The bibliography system is now set up to work seamlessly with Bibby AI's citation management. Feel free to modify it for your needs!

\bibliographystyle{plain}
\bibliography{cite}

\end{document}\documentclass{article}
\usepackage[utf8]{inputenc}
\usepackage{booktabs}
\usepackage{amsmath}
\usepackage{natbib}

\title{My Document}
\author{Your Name}
\date{\today}

\begin{document}

\maketitle

\section{Introduction}
This is a sample document demonstrating various \LaTeX{} features including mathematical equations, tables, and formatting. You can cite papers using the citation search feature, and they will be automatically added to your \texttt{cite.bib} file.

\section{Mathematical Expressions}
Here are some example equations:

\subsection{Quadratic Formula}
The quadratic formula is given by:
\[x = \frac{-b \pm \sqrt{b^2 - 4ac}}{2a}\]

\subsection{Matrix Example}
A sample matrix:
\[
\begin{pmatrix}
a_{11} & a_{12} & a_{13} \\
a_{21} & a_{22} & a_{23} \\
a_{31} & a_{32} & a_{33}
\end{pmatrix}
\]

\section{Tables}
Below is a sample table showing data:

\begin{table}[h]
\centering
\begin{tabular}{@{}lrrr@{}}
\toprule
\textbf{Category} & \textbf{Value 1} & \textbf{Value 2} & \textbf{Total} \\
\midrule
Item A & 10.5 & 20.3 & 30.8 \\
Item B & 15.2 & 25.1 & 40.3 \\
Item C & 22.1 & 12.7 & 34.8 \\
\bottomrule
\end{tabular}
\caption{Sample Data Table}
\label{tab:sample}
\end{table}

\section{Lists}
\subsection{Bullet Points}
\begin{itemize}
  \item First item
  \item Second item
  \item Third item with sub-items:
    \begin{itemize}
      \item Sub-item 1
      \item Sub-item 2
    \end{itemize}
\end{itemize}

\subsection{Numbered List}
\begin{enumerate}
  \item First step
  \item Second step
  \item Third step
\end{enumerate}

\section{Citations and References}
When you use the citation search feature, references will be automatically managed. For example, you might cite a paper like this: \cite{example2024}. All citations will be automatically added to your \texttt{cite.bib} file.

\section{Conclusion}
This template demonstrates basic \LaTeX{} features including mathematical formulas, tables, and lists. The bibliography system is now set up to work seamlessly with Bibby AI's citation management. Feel free to modify it for your needs!

\bibliographystyle{plain}
\bibliography{cite}

\end{document}\documentclass{article}
\usepackage[utf8]{inputenc}
\usepackage{booktabs}
\usepackage{amsmath}
\usepackage{natbib}

\title{My Document}
\author{Your Name}
\date{\today}

\begin{document}

\maketitle

\section{Introduction}
This is a sample document demonstrating various \LaTeX{} features including mathematical equations, tables, and formatting. You can cite papers using the citation search feature, and they will be automatically added to your \texttt{cite.bib} file.

\section{Mathematical Expressions}
Here are some example equations:

\subsection{Quadratic Formula}
The quadratic formula is given by:
\[x = \frac{-b \pm \sqrt{b^2 - 4ac}}{2a}\]

\subsection{Matrix Example}
A sample matrix:
\[
\begin{pmatrix}
a_{11} & a_{12} & a_{13} \\
a_{21} & a_{22} & a_{23} \\
a_{31} & a_{32} & a_{33}
\end{pmatrix}
\]

\section{Tables}
Below is a sample table showing data:

\begin{table}[h]
\centering
\begin{tabular}{@{}lrrr@{}}
\toprule
\textbf{Category} & \textbf{Value 1} & \textbf{Value 2} & \textbf{Total} \\
\midrule
Item A & 10.5 & 20.3 & 30.8 \\
Item B & 15.2 & 25.1 & 40.3 \\
Item C & 22.1 & 12.7 & 34.8 \\
\bottomrule
\end{tabular}
\caption{Sample Data Table}
\label{tab:sample}
\end{table}

\section{Lists}
\subsection{Bullet Points}
\begin{itemize}
  \item First item
  \item Second item
  \item Third item with sub-items:
    \begin{itemize}
      \item Sub-item 1
      \item Sub-item 2
    \end{itemize}
\end{itemize}

\subsection{Numbered List}
\begin{enumerate}
  \item First step
  \item Second step
  \item Third step
\end{enumerate}

\section{Citations and References}
When you use the citation search feature, references will be automatically managed. For example, you might cite a paper like this: \cite{example2024}. All citations will be automatically added to your \texttt{cite.bib} file.

\section{Conclusion}
This template demonstrates basic \LaTeX{} features including mathematical formulas, tables, and lists. The bibliography system is now set up to work seamlessly with Bibby AI's citation management. Feel free to modify it for your needs!

\bibliographystyle{plain}
\bibliography{cite}

\end{document}\documentclass{article}
\usepackage[utf8]{inputenc}
\usepackage{booktabs}
\usepackage{amsmath}
\usepackage{natbib}

\title{My Document}
\author{Your Name}
\date{\today}

\begin{document}

\maketitle

\section{Introduction}
This is a sample document demonstrating various \LaTeX{} features including mathematical equations, tables, and formatting. You can cite papers using the citation search feature, and they will be automatically added to your \texttt{cite.bib} file.

\section{Mathematical Expressions}
Here are some example equations:

\subsection{Quadratic Formula}
The quadratic formula is given by:
\[x = \frac{-b \pm \sqrt{b^2 - 4ac}}{2a}\]

\subsection{Matrix Example}
A sample matrix:
\[
\begin{pmatrix}
a_{11} & a_{12} & a_{13} \\
a_{21} & a_{22} & a_{23} \\
a_{31} & a_{32} & a_{33}
\end{pmatrix}
\]

\section{Tables}
Below is a sample table showing data:

\begin{table}[h]
\centering
\begin{tabular}{@{}lrrr@{}}
\toprule
\textbf{Category} & \textbf{Value 1} & \textbf{Value 2} & \textbf{Total} \\
\midrule
Item A & 10.5 & 20.3 & 30.8 \\
Item B & 15.2 & 25.1 & 40.3 \\
Item C & 22.1 & 12.7 & 34.8 \\
\bottomrule
\end{tabular}
\caption{Sample Data Table}
\label{tab:sample}
\end{table}

\section{Lists}
\subsection{Bullet Points}
\begin{itemize}
  \item First item
  \item Second item
  \item Third item with sub-items:
    \begin{itemize}
      \item Sub-item 1
      \item Sub-item 2
    \end{itemize}
\end{itemize}

\subsection{Numbered List}
\begin{enumerate}
  \item First step
  \item Second step
  \item Third step
\end{enumerate}

\section{Citations and References}
When you use the citation search feature, references will be automatically managed. For example, you might cite a paper like this: \cite{example2024}. All citations will be automatically added to your \texttt{cite.bib} file.

\section{Conclusion}
This template demonstrates basic \LaTeX{} features including mathematical formulas, tables, and lists. The bibliography system is now set up to work seamlessly with Bibby AI's citation management. Feel free to modify it for your needs!

\bibliographystyle{plain}
\bibliography{cite}

\end{document}\documentclass{article}
\usepackage[utf8]{inputenc}
\usepackage{booktabs}
\usepackage{amsmath}
\usepackage{natbib}

\title{My Document}
\author{Your Name}
\date{\today}

\begin{document}

\maketitle

\section{Introduction}
This is a sample document demonstrating various \LaTeX{} features including mathematical equations, tables, and formatting. You can cite papers using the citation search feature, and they will be automatically added to your \texttt{cite.bib} file.

\section{Mathematical Expressions}
Here are some example equations:

\subsection{Quadratic Formula}
The quadratic formula is given by:
\[x = \frac{-b \pm \sqrt{b^2 - 4ac}}{2a}\]

\subsection{Matrix Example}
A sample matrix:
\[
\begin{pmatrix}
a_{11} & a_{12} & a_{13} \\
a_{21} & a_{22} & a_{23} \\
a_{31} & a_{32} & a_{33}
\end{pmatrix}
\]

\section{Tables}
Below is a sample table showing data:

\begin{table}[h]
\centering
\begin{tabular}{@{}lrrr@{}}
\toprule
\textbf{Category} & \textbf{Value 1} & \textbf{Value 2} & \textbf{Total} \\
\midrule
Item A & 10.5 & 20.3 & 30.8 \\
Item B & 15.2 & 25.1 & 40.3 \\
Item C & 22.1 & 12.7 & 34.8 \\
\bottomrule
\end{tabular}
\caption{Sample Data Table}
\label{tab:sample}
\end{table}

\section{Lists}
\subsection{Bullet Points}
\begin{itemize}
  \item First item
  \item Second item
  \item Third item with sub-items:
    \begin{itemize}
      \item Sub-item 1
      \item Sub-item 2
    \end{itemize}
\end{itemize}

\subsection{Numbered List}
\begin{enumerate}
  \item First step
  \item Second step
  \item Third step
\end{enumerate}

\section{Citations and References}
When you use the citation search feature, references will be automatically managed. For example, you might cite a paper like this: \cite{example2024}. All citations will be automatically added to your \texttt{cite.bib} file.

\section{Conclusion}
This template demonstrates basic \LaTeX{} features including mathematical formulas, tables, and lists. The bibliography system is now set up to work seamlessly with Bibby AI's citation management. Feel free to modify it for your needs!

\bibliographystyle{plain}
\bibliography{cite}

\end{document}\documentclass{article}
\usepackage[utf8]{inputenc}
\usepackage{booktabs}
\usepackage{amsmath}
\usepackage{natbib}

\title{My Document}
\author{Your Name}
\date{\today}

\begin{document}

\maketitle

\section{Introduction}
This is a sample document demonstrating various \LaTeX{} features including mathematical equations, tables, and formatting. You can cite papers using the citation search feature, and they will be automatically added to your \texttt{cite.bib} file.

\section{Mathematical Expressions}
Here are some example equations:

\subsection{Quadratic Formula}
The quadratic formula is given by:
\[x = \frac{-b \pm \sqrt{b^2 - 4ac}}{2a}\]

\subsection{Matrix Example}
A sample matrix:
\[
\begin{pmatrix}
a_{11} & a_{12} & a_{13} \\
a_{21} & a_{22} & a_{23} \\
a_{31} & a_{32} & a_{33}
\end{pmatrix}
\]

\section{Tables}
Below is a sample table showing data:

\begin{table}[h]
\centering
\begin{tabular}{@{}lrrr@{}}
\toprule
\textbf{Category} & \textbf{Value 1} & \textbf{Value 2} & \textbf{Total} \\
\midrule
Item A & 10.5 & 20.3 & 30.8 \\
Item B & 15.2 & 25.1 & 40.3 \\
Item C & 22.1 & 12.7 & 34.8 \\
\bottomrule
\end{tabular}
\caption{Sample Data Table}
\label{tab:sample}
\end{table}

\section{Lists}
\subsection{Bullet Points}
\begin{itemize}
  \item First item
  \item Second item
  \item Third item with sub-items:
    \begin{itemize}
      \item Sub-item 1
      \item Sub-item 2
    \end{itemize}
\end{itemize}

\subsection{Numbered List}
\begin{enumerate}
  \item First step
  \item Second step
  \item Third step
\end{enumerate}

\section{Citations and References}
When you use the citation search feature, references will be automatically managed. For example, you might cite a paper like this: \cite{example2024}. All citations will be automatically added to your \texttt{cite.bib} file.

\section{Conclusion}
This template demonstrates basic \LaTeX{} features including mathematical formulas, tables, and lists. The bibliography system is now set up to work seamlessly with Bibby AI's citation management. Feel free to modify it for your needs!

\bibliographystyle{plain}
\bibliography{cite}

\end{document}\documentclass{article}
\usepackage[utf8]{inputenc}
\usepackage{booktabs}
\usepackage{amsmath}
\usepackage{natbib}

\title{My Document}
\author{Your Name}
\date{\today}

\begin{document}

\maketitle

\section{Introduction}
This is a sample document demonstrating various \LaTeX{} features including mathematical equations, tables, and formatting. You can cite papers using the citation search feature, and they will be automatically added to your \texttt{cite.bib} file.

\section{Mathematical Expressions}
Here are some example equations:

\subsection{Quadratic Formula}
The quadratic formula is given by:
\[x = \frac{-b \pm \sqrt{b^2 - 4ac}}{2a}\]

\subsection{Matrix Example}
A sample matrix:
\[
\begin{pmatrix}
a_{11} & a_{12} & a_{13} \\
a_{21} & a_{22} & a_{23} \\
a_{31} & a_{32} & a_{33}
\end{pmatrix}
\]

\section{Tables}
Below is a sample table showing data:

\begin{table}[h]
\centering
\begin{tabular}{@{}lrrr@{}}
\toprule
\textbf{Category} & \textbf{Value 1} & \textbf{Value 2} & \textbf{Total} \\
\midrule
Item A & 10.5 & 20.3 & 30.8 \\
Item B & 15.2 & 25.1 & 40.3 \\
Item C & 22.1 & 12.7 & 34.8 \\
\bottomrule
\end{tabular}
\caption{Sample Data Table}
\label{tab:sample}
\end{table}

\section{Lists}
\subsection{Bullet Points}
\begin{itemize}
  \item First item
  \item Second item
  \item Third item with sub-items:
    \begin{itemize}
      \item Sub-item 1
      \item Sub-item 2
    \end{itemize}
\end{itemize}

\subsection{Numbered List}
\begin{enumerate}
  \item First step
  \item Second step
  \item Third step
\end{enumerate}

\section{Citations and References}
When you use the citation search feature, references will be automatically managed. For example, you might cite a paper like this: \cite{example2024}. All citations will be automatically added to your \texttt{cite.bib} file.

\section{Conclusion}
This template demonstrates basic \LaTeX{} features including mathematical formulas, tables, and lists. The bibliography system is now set up to work seamlessly with Bibby AI's citation management. Feel free to modify it for your needs!

\bibliographystyle{plain}
\bibliography{cite}

\end{document}\documentclass{article}
\usepackage[utf8]{inputenc}
\usepackage{booktabs}
\usepackage{amsmath}
\usepackage{natbib}

\title{My Document}
\author{Your Name}
\date{\today}

\begin{document}

\maketitle

\section{Introduction}
This is a sample document demonstrating various \LaTeX{} features including mathematical equations, tables, and formatting. You can cite papers using the citation search feature, and they will be automatically added to your \texttt{cite.bib} file.

\section{Mathematical Expressions}
Here are some example equations:

\subsection{Quadratic Formula}
The quadratic formula is given by:
\[x = \frac{-b \pm \sqrt{b^2 - 4ac}}{2a}\]

\subsection{Matrix Example}
A sample matrix:
\[
\begin{pmatrix}
a_{11} & a_{12} & a_{13} \\
a_{21} & a_{22} & a_{23} \\
a_{31} & a_{32} & a_{33}
\end{pmatrix}
\]

\section{Tables}
Below is a sample table showing data:

\begin{table}[h]
\centering
\begin{tabular}{@{}lrrr@{}}
\toprule
\textbf{Category} & \textbf{Value 1} & \textbf{Value 2} & \textbf{Total} \\
\midrule
Item A & 10.5 & 20.3 & 30.8 \\
Item B & 15.2 & 25.1 & 40.3 \\
Item C & 22.1 & 12.7 & 34.8 \\
\bottomrule
\end{tabular}
\caption{Sample Data Table}
\label{tab:sample}
\end{table}

\section{Lists}
\subsection{Bullet Points}
\begin{itemize}
  \item First item
  \item Second item
  \item Third item with sub-items:
    \begin{itemize}
      \item Sub-item 1
      \item Sub-item 2
    \end{itemize}
\end{itemize}

\subsection{Numbered List}
\begin{enumerate}
  \item First step
  \item Second step
  \item Third step
\end{enumerate}

\section{Citations and References}
When you use the citation search feature, references will be automatically managed. For example, you might cite a paper like this: \cite{example2024}. All citations will be automatically added to your \texttt{cite.bib} file.

\section{Conclusion}
This template demonstrates basic \LaTeX{} features including mathematical formulas, tables, and lists. The bibliography system is now set up to work seamlessly with Bibby AI's citation management. Feel free to modify it for your needs!

\bibliographystyle{plain}
\bibliography{cite}

\end{document}\documentclass{article}
\usepackage[utf8]{inputenc}
\usepackage{booktabs}
\usepackage{amsmath}
\usepackage{natbib}

\title{My Document}
\author{Your Name}
\date{\today}

\begin{document}

\maketitle

\section{Introduction}
This is a sample document demonstrating various \LaTeX{} features including mathematical equations, tables, and formatting. You can cite papers using the citation search feature, and they will be automatically added to your \texttt{cite.bib} file.

\section{Mathematical Expressions}
Here are some example equations:

\subsection{Quadratic Formula}
The quadratic formula is given by:
\[x = \frac{-b \pm \sqrt{b^2 - 4ac}}{2a}\]

\subsection{Matrix Example}
A sample matrix:
\[
\begin{pmatrix}
a_{11} & a_{12} & a_{13} \\
a_{21} & a_{22} & a_{23} \\
a_{31} & a_{32} & a_{33}
\end{pmatrix}
\]

\section{Tables}
Below is a sample table showing data:

\begin{table}[h]
\centering
\begin{tabular}{@{}lrrr@{}}
\toprule
\textbf{Category} & \textbf{Value 1} & \textbf{Value 2} & \textbf{Total} \\
\midrule
Item A & 10.5 & 20.3 & 30.8 \\
Item B & 15.2 & 25.1 & 40.3 \\
Item C & 22.1 & 12.7 & 34.8 \\
\bottomrule
\end{tabular}
\caption{Sample Data Table}
\label{tab:sample}
\end{table}

\section{Lists}
\subsection{Bullet Points}
\begin{itemize}
  \item First item
  \item Second item
  \item Third item with sub-items:
    \begin{itemize}
      \item Sub-item 1
      \item Sub-item 2
    \end{itemize}
\end{itemize}

\subsection{Numbered List}
\begin{enumerate}
  \item First step
  \item Second step
  \item Third step
\end{enumerate}

\section{Citations and References}
When you use the citation search feature, references will be automatically managed. For example, you might cite a paper like this: \cite{example2024}. All citations will be automatically added to your \texttt{cite.bib} file.

\section{Conclusion}
This template demonstrates basic \LaTeX{} features including mathematical formulas, tables, and lists. The bibliography system is now set up to work seamlessly with Bibby AI's citation management. Feel free to modify it for your needs!

\bibliographystyle{plain}
\bibliography{cite}

\end{document}\documentclass{article}
\usepackage[utf8]{inputenc}
\usepackage{booktabs}
\usepackage{amsmath}
\usepackage{natbib}

\title{My Document}
\author{Your Name}
\date{\today}

\begin{document}

\maketitle

\section{Introduction}
This is a sample document demonstrating various \LaTeX{} features including mathematical equations, tables, and formatting. You can cite papers using the citation search feature, and they will be automatically added to your \texttt{cite.bib} file.

\section{Mathematical Expressions}
Here are some example equations:

\subsection{Quadratic Formula}
The quadratic formula is given by:
\[x = \frac{-b \pm \sqrt{b^2 - 4ac}}{2a}\]

\subsection{Matrix Example}
A sample matrix:
\[
\begin{pmatrix}
a_{11} & a_{12} & a_{13} \\
a_{21} & a_{22} & a_{23} \\
a_{31} & a_{32} & a_{33}
\end{pmatrix}
\]

\section{Tables}
Below is a sample table showing data:

\begin{table}[h]
\centering
\begin{tabular}{@{}lrrr@{}}
\toprule
\textbf{Category} & \textbf{Value 1} & \textbf{Value 2} & \textbf{Total} \\
\midrule
Item A & 10.5 & 20.3 & 30.8 \\
Item B & 15.2 & 25.1 & 40.3 \\
Item C & 22.1 & 12.7 & 34.8 \\
\bottomrule
\end{tabular}
\caption{Sample Data Table}
\label{tab:sample}
\end{table}

\section{Lists}
\subsection{Bullet Points}
\begin{itemize}
  \item First item
  \item Second item
  \item Third item with sub-items:
    \begin{itemize}
      \item Sub-item 1
      \item Sub-item 2
    \end{itemize}
\end{itemize}

\subsection{Numbered List}
\begin{enumerate}
  \item First step
  \item Second step
  \item Third step
\end{enumerate}

\section{Citations and References}
When you use the citation search feature, references will be automatically managed. For example, you might cite a paper like this: \cite{example2024}. All citations will be automatically added to your \texttt{cite.bib} file.

\section{Conclusion}
This template demonstrates basic \LaTeX{} features including mathematical formulas, tables, and lists. The bibliography system is now set up to work seamlessly with Bibby AI's citation management. Feel free to modify it for your needs!

\bibliographystyle{plain}
\bibliography{cite}

\end{document}\documentclass{article}
\usepackage[utf8]{inputenc}
\usepackage{booktabs}
\usepackage{amsmath}
\usepackage{natbib}

\title{My Document}
\author{Your Name}
\date{\today}

\begin{document}

\maketitle

\section{Introduction}
This is a sample document demonstrating various \LaTeX{} features including mathematical equations, tables, and formatting. You can cite papers using the citation search feature, and they will be automatically added to your \texttt{cite.bib} file.

\section{Mathematical Expressions}
Here are some example equations:

\subsection{Quadratic Formula}
The quadratic formula is given by:
\[x = \frac{-b \pm \sqrt{b^2 - 4ac}}{2a}\]

\subsection{Matrix Example}
A sample matrix:
\[
\begin{pmatrix}
a_{11} & a_{12} & a_{13} \\
a_{21} & a_{22} & a_{23} \\
a_{31} & a_{32} & a_{33}
\end{pmatrix}
\]

\section{Tables}
Below is a sample table showing data:

\begin{table}[h]
\centering
\begin{tabular}{@{}lrrr@{}}
\toprule
\textbf{Category} & \textbf{Value 1} & \textbf{Value 2} & \textbf{Total} \\
\midrule
Item A & 10.5 & 20.3 & 30.8 \\
Item B & 15.2 & 25.1 & 40.3 \\
Item C & 22.1 & 12.7 & 34.8 \\
\bottomrule
\end{tabular}
\caption{Sample Data Table}
\label{tab:sample}
\end{table}

\section{Lists}
\subsection{Bullet Points}
\begin{itemize}
  \item First item
  \item Second item
  \item Third item with sub-items:
    \begin{itemize}
      \item Sub-item 1
      \item Sub-item 2
    \end{itemize}
\end{itemize}

\subsection{Numbered List}
\begin{enumerate}
  \item First step
  \item Second step
  \item Third step
\end{enumerate}

\section{Citations and References}
When you use the citation search feature, references will be automatically managed. For example, you might cite a paper like this: \cite{example2024}. All citations will be automatically added to your \texttt{cite.bib} file.

\section{Conclusion}
This template demonstrates basic \LaTeX{} features including mathematical formulas, tables, and lists. The bibliography system is now set up to work seamlessly with Bibby AI's citation management. Feel free to modify it for your needs!

\bibliographystyle{plain}
\bibliography{cite}

\end{document}\documentclass{article}
\usepackage[utf8]{inputenc}
\usepackage{booktabs}
\usepackage{amsmath}
\usepackage{natbib}

\title{My Document}
\author{Your Name}
\date{\today}

\begin{document}

\maketitle

\section{Introduction}
This is a sample document demonstrating various \LaTeX{} features including mathematical equations, tables, and formatting. You can cite papers using the citation search feature, and they will be automatically added to your \texttt{cite.bib} file.

\section{Mathematical Expressions}
Here are some example equations:

\subsection{Quadratic Formula}
The quadratic formula is given by:
\[x = \frac{-b \pm \sqrt{b^2 - 4ac}}{2a}\]

\subsection{Matrix Example}
A sample matrix:
\[
\begin{pmatrix}
a_{11} & a_{12} & a_{13} \\
a_{21} & a_{22} & a_{23} \\
a_{31} & a_{32} & a_{33}
\end{pmatrix}
\]

\section{Tables}
Below is a sample table showing data:

\begin{table}[h]
\centering
\begin{tabular}{@{}lrrr@{}}
\toprule
\textbf{Category} & \textbf{Value 1} & \textbf{Value 2} & \textbf{Total} \\
\midrule
Item A & 10.5 & 20.3 & 30.8 \\
Item B & 15.2 & 25.1 & 40.3 \\
Item C & 22.1 & 12.7 & 34.8 \\
\bottomrule
\end{tabular}
\caption{Sample Data Table}
\label{tab:sample}
\end{table}

\section{Lists}
\subsection{Bullet Points}
\begin{itemize}
  \item First item
  \item Second item
  \item Third item with sub-items:
    \begin{itemize}
      \item Sub-item 1
      \item Sub-item 2
    \end{itemize}
\end{itemize}

\subsection{Numbered List}
\begin{enumerate}
  \item First step
  \item Second step
  \item Third step
\end{enumerate}

\section{Citations and References}
When you use the citation search feature, references will be automatically managed. For example, you might cite a paper like this: \cite{example2024}. All citations will be automatically added to your \texttt{cite.bib} file.

\section{Conclusion}
This template demonstrates basic \LaTeX{} features including mathematical formulas, tables, and lists. The bibliography system is now set up to work seamlessly with Bibby AI's citation management. Feel free to modify it for your needs!

\bibliographystyle{plain}
\bibliography{cite}

\end{document}\documentclass{article}
\usepackage[utf8]{inputenc}
\usepackage{booktabs}
\usepackage{amsmath}
\usepackage{natbib}

\title{My Document}
\author{Your Name}
\date{\today}

\begin{document}

\maketitle

\section{Introduction}
This is a sample document demonstrating various \LaTeX{} features including mathematical equations, tables, and formatting. You can cite papers using the citation search feature, and they will be automatically added to your \texttt{cite.bib} file.

\section{Mathematical Expressions}
Here are some example equations:

\subsection{Quadratic Formula}
The quadratic formula is given by:
\[x = \frac{-b \pm \sqrt{b^2 - 4ac}}{2a}\]

\subsection{Matrix Example}
A sample matrix:
\[
\begin{pmatrix}
a_{11} & a_{12} & a_{13} \\
a_{21} & a_{22} & a_{23} \\
a_{31} & a_{32} & a_{33}
\end{pmatrix}
\]

\section{Tables}
Below is a sample table showing data:

\begin{table}[h]
\centering
\begin{tabular}{@{}lrrr@{}}
\toprule
\textbf{Category} & \textbf{Value 1} & \textbf{Value 2} & \textbf{Total} \\
\midrule
Item A & 10.5 & 20.3 & 30.8 \\
Item B & 15.2 & 25.1 & 40.3 \\
Item C & 22.1 & 12.7 & 34.8 \\
\bottomrule
\end{tabular}
\caption{Sample Data Table}
\label{tab:sample}
\end{table}

\section{Lists}
\subsection{Bullet Points}
\begin{itemize}
  \item First item
  \item Second item
  \item Third item with sub-items:
    \begin{itemize}
      \item Sub-item 1
      \item Sub-item 2
    \end{itemize}
\end{itemize}

\subsection{Numbered List}
\begin{enumerate}
  \item First step
  \item Second step
  \item Third step
\end{enumerate}

\section{Citations and References}
When you use the citation search feature, references will be automatically managed. For example, you might cite a paper like this: \cite{example2024}. All citations will be automatically added to your \texttt{cite.bib} file.

\section{Conclusion}
This template demonstrates basic \LaTeX{} features including mathematical formulas, tables, and lists. The bibliography system is now set up to work seamlessly with Bibby AI's citation management. Feel free to modify it for your needs!

\bibliographystyle{plain}
\bibliography{cite}

\end{document}\documentclass{article}
\usepackage[utf8]{inputenc}
\usepackage{booktabs}
\usepackage{amsmath}
\usepackage{natbib}

\title{My Document}
\author{Your Name}
\date{\today}

\begin{document}

\maketitle

\section{Introduction}
This is a sample document demonstrating various \LaTeX{} features including mathematical equations, tables, and formatting. You can cite papers using the citation search feature, and they will be automatically added to your \texttt{cite.bib} file.

\section{Mathematical Expressions}
Here are some example equations:

\subsection{Quadratic Formula}
The quadratic formula is given by:
\[x = \frac{-b \pm \sqrt{b^2 - 4ac}}{2a}\]

\subsection{Matrix Example}
A sample matrix:
\[
\begin{pmatrix}
a_{11} & a_{12} & a_{13} \\
a_{21} & a_{22} & a_{23} \\
a_{31} & a_{32} & a_{33}
\end{pmatrix}
\]

\section{Tables}
Below is a sample table showing data:

\begin{table}[h]
\centering
\begin{tabular}{@{}lrrr@{}}
\toprule
\textbf{Category} & \textbf{Value 1} & \textbf{Value 2} & \textbf{Total} \\
\midrule
Item A & 10.5 & 20.3 & 30.8 \\
Item B & 15.2 & 25.1 & 40.3 \\
Item C & 22.1 & 12.7 & 34.8 \\
\bottomrule
\end{tabular}
\caption{Sample Data Table}
\label{tab:sample}
\end{table}

\section{Lists}
\subsection{Bullet Points}
\begin{itemize}
  \item First item
  \item Second item
  \item Third item with sub-items:
    \begin{itemize}
      \item Sub-item 1
      \item Sub-item 2
    \end{itemize}
\end{itemize}

\subsection{Numbered List}
\begin{enumerate}
  \item First step
  \item Second step
  \item Third step
\end{enumerate}

\section{Citations and References}
When you use the citation search feature, references will be automatically managed. For example, you might cite a paper like this: \cite{example2024}. All citations will be automatically added to your \texttt{cite.bib} file.

\section{Conclusion}
This template demonstrates basic \LaTeX{} features including mathematical formulas, tables, and lists. The bibliography system is now set up to work seamlessly with Bibby AI's citation management. Feel free to modify it for your needs!

\bibliographystyle{plain}
\bibliography{cite}

\end{document}\documentclass{article}
\usepackage[utf8]{inputenc}
\usepackage{booktabs}
\usepackage{amsmath}
\usepackage{natbib}

\title{My Document}
\author{Your Name}
\date{\today}

\begin{document}

\maketitle

\section{Introduction}
This is a sample document demonstrating various \LaTeX{} features including mathematical equations, tables, and formatting. You can cite papers using the citation search feature, and they will be automatically added to your \texttt{cite.bib} file.

\section{Mathematical Expressions}
Here are some example equations:

\subsection{Quadratic Formula}
The quadratic formula is given by:
\[x = \frac{-b \pm \sqrt{b^2 - 4ac}}{2a}\]

\subsection{Matrix Example}
A sample matrix:
\[
\begin{pmatrix}
a_{11} & a_{12} & a_{13} \\
a_{21} & a_{22} & a_{23} \\
a_{31} & a_{32} & a_{33}
\end{pmatrix}
\]

\section{Tables}
Below is a sample table showing data:

\begin{table}[h]
\centering
\begin{tabular}{@{}lrrr@{}}
\toprule
\textbf{Category} & \textbf{Value 1} & \textbf{Value 2} & \textbf{Total} \\
\midrule
Item A & 10.5 & 20.3 & 30.8 \\
Item B & 15.2 & 25.1 & 40.3 \\
Item C & 22.1 & 12.7 & 34.8 \\
\bottomrule
\end{tabular}
\caption{Sample Data Table}
\label{tab:sample}
\end{table}

\section{Lists}
\subsection{Bullet Points}
\begin{itemize}
  \item First item
  \item Second item
  \item Third item with sub-items:
    \begin{itemize}
      \item Sub-item 1
      \item Sub-item 2
    \end{itemize}
\end{itemize}

\subsection{Numbered List}
\begin{enumerate}
  \item First step
  \item Second step
  \item Third step
\end{enumerate}

\section{Citations and References}
When you use the citation search feature, references will be automatically managed. For example, you might cite a paper like this: \cite{example2024}. All citations will be automatically added to your \texttt{cite.bib} file.

\section{Conclusion}
This template demonstrates basic \LaTeX{} features including mathematical formulas, tables, and lists. The bibliography system is now set up to work seamlessly with Bibby AI's citation management. Feel free to modify it for your needs!

\bibliographystyle{plain}
\bibliography{cite}

\end{document}\documentclass{article}
\usepackage[utf8]{inputenc}
\usepackage{booktabs}
\usepackage{amsmath}
\usepackage{natbib}

\title{My Document}
\author{Your Name}
\date{\today}

\begin{document}

\maketitle

\section{Introduction}
This is a sample document demonstrating various \LaTeX{} features including mathematical equations, tables, and formatting. You can cite papers using the citation search feature, and they will be automatically added to your \texttt{cite.bib} file.

\section{Mathematical Expressions}
Here are some example equations:

\subsection{Quadratic Formula}
The quadratic formula is given by:
\[x = \frac{-b \pm \sqrt{b^2 - 4ac}}{2a}\]

\subsection{Matrix Example}
A sample matrix:
\[
\begin{pmatrix}
a_{11} & a_{12} & a_{13} \\
a_{21} & a_{22} & a_{23} \\
a_{31} & a_{32} & a_{33}
\end{pmatrix}
\]

\section{Tables}
Below is a sample table showing data:

\begin{table}[h]
\centering
\begin{tabular}{@{}lrrr@{}}
\toprule
\textbf{Category} & \textbf{Value 1} & \textbf{Value 2} & \textbf{Total} \\
\midrule
Item A & 10.5 & 20.3 & 30.8 \\
Item B & 15.2 & 25.1 & 40.3 \\
Item C & 22.1 & 12.7 & 34.8 \\
\bottomrule
\end{tabular}
\caption{Sample Data Table}
\label{tab:sample}
\end{table}

\section{Lists}
\subsection{Bullet Points}
\begin{itemize}
  \item First item
  \item Second item
  \item Third item with sub-items:
    \begin{itemize}
      \item Sub-item 1
      \item Sub-item 2
    \end{itemize}
\end{itemize}

\subsection{Numbered List}
\begin{enumerate}
  \item First step
  \item Second step
  \item Third step
\end{enumerate}

\section{Citations and References}
When you use the citation search feature, references will be automatically managed. For example, you might cite a paper like this: \cite{example2024}. All citations will be automatically added to your \texttt{cite.bib} file.

\section{Conclusion}
This template demonstrates basic \LaTeX{} features including mathematical formulas, tables, and lists. The bibliography system is now set up to work seamlessly with Bibby AI's citation management. Feel free to modify it for your needs!

\bibliographystyle{plain}
\bibliography{cite}

\end{document}\documentclass{article}
\usepackage[utf8]{inputenc}
\usepackage{booktabs}
\usepackage{amsmath}
\usepackage{natbib}

\title{My Document}
\author{Your Name}
\date{\today}

\begin{document}

\maketitle

\section{Introduction}
This is a sample document demonstrating various \LaTeX{} features including mathematical equations, tables, and formatting. You can cite papers using the citation search feature, and they will be automatically added to your \texttt{cite.bib} file.

\section{Mathematical Expressions}
Here are some example equations:

\subsection{Quadratic Formula}
The quadratic formula is given by:
\[x = \frac{-b \pm \sqrt{b^2 - 4ac}}{2a}\]

\subsection{Matrix Example}
A sample matrix:
\[
\begin{pmatrix}
a_{11} & a_{12} & a_{13} \\
a_{21} & a_{22} & a_{23} \\
a_{31} & a_{32} & a_{33}
\end{pmatrix}
\]

\section{Tables}
Below is a sample table showing data:

\begin{table}[h]
\centering
\begin{tabular}{@{}lrrr@{}}
\toprule
\textbf{Category} & \textbf{Value 1} & \textbf{Value 2} & \textbf{Total} \\
\midrule
Item A & 10.5 & 20.3 & 30.8 \\
Item B & 15.2 & 25.1 & 40.3 \\
Item C & 22.1 & 12.7 & 34.8 \\
\bottomrule
\end{tabular}
\caption{Sample Data Table}
\label{tab:sample}
\end{table}

\section{Lists}
\subsection{Bullet Points}
\begin{itemize}
  \item First item
  \item Second item
  \item Third item with sub-items:
    \begin{itemize}
      \item Sub-item 1
      \item Sub-item 2
    \end{itemize}
\end{itemize}

\subsection{Numbered List}
\begin{enumerate}
  \item First step
  \item Second step
  \item Third step
\end{enumerate}

\section{Citations and References}
When you use the citation search feature, references will be automatically managed. For example, you might cite a paper like this: \cite{example2024}. All citations will be automatically added to your \texttt{cite.bib} file.

\section{Conclusion}
This template demonstrates basic \LaTeX{} features including mathematical formulas, tables, and lists. The bibliography system is now set up to work seamlessly with Bibby AI's citation management. Feel free to modify it for your needs!

\bibliographystyle{plain}
\bibliography{cite}

\end{document}\documentclass{article}
\usepackage[utf8]{inputenc}
\usepackage{booktabs}
\usepackage{amsmath}
\usepackage{natbib}

\title{My Document}
\author{Your Name}
\date{\today}

\begin{document}

\maketitle

\section{Introduction}
This is a sample document demonstrating various \LaTeX{} features including mathematical equations, tables, and formatting. You can cite papers using the citation search feature, and they will be automatically added to your \texttt{cite.bib} file.

\section{Mathematical Expressions}
Here are some example equations:

\subsection{Quadratic Formula}
The quadratic formula is given by:
\[x = \frac{-b \pm \sqrt{b^2 - 4ac}}{2a}\]

\subsection{Matrix Example}
A sample matrix:
\[
\begin{pmatrix}
a_{11} & a_{12} & a_{13} \\
a_{21} & a_{22} & a_{23} \\
a_{31} & a_{32} & a_{33}
\end{pmatrix}
\]

\section{Tables}
Below is a sample table showing data:

\begin{table}[h]
\centering
\begin{tabular}{@{}lrrr@{}}
\toprule
\textbf{Category} & \textbf{Value 1} & \textbf{Value 2} & \textbf{Total} \\
\midrule
Item A & 10.5 & 20.3 & 30.8 \\
Item B & 15.2 & 25.1 & 40.3 \\
Item C & 22.1 & 12.7 & 34.8 \\
\bottomrule
\end{tabular}
\caption{Sample Data Table}
\label{tab:sample}
\end{table}

\section{Lists}
\subsection{Bullet Points}
\begin{itemize}
  \item First item
  \item Second item
  \item Third item with sub-items:
    \begin{itemize}
      \item Sub-item 1
      \item Sub-item 2
    \end{itemize}
\end{itemize}

\subsection{Numbered List}
\begin{enumerate}
  \item First step
  \item Second step
  \item Third step
\end{enumerate}

\section{Citations and References}
When you use the citation search feature, references will be automatically managed. For example, you might cite a paper like this: \cite{example2024}. All citations will be automatically added to your \texttt{cite.bib} file.

\section{Conclusion}
This template demonstrates basic \LaTeX{} features including mathematical formulas, tables, and lists. The bibliography system is now set up to work seamlessly with Bibby AI's citation management. Feel free to modify it for your needs!

\bibliographystyle{plain}
\bibliography{cite}

\end{document}\documentclass{article}
\usepackage[utf8]{inputenc}
\usepackage{booktabs}
\usepackage{amsmath}
\usepackage{natbib}

\title{My Document}
\author{Your Name}
\date{\today}

\begin{document}

\maketitle

\section{Introduction}
This is a sample document demonstrating various \LaTeX{} features including mathematical equations, tables, and formatting. You can cite papers using the citation search feature, and they will be automatically added to your \texttt{cite.bib} file.

\section{Mathematical Expressions}
Here are some example equations:

\subsection{Quadratic Formula}
The quadratic formula is given by:
\[x = \frac{-b \pm \sqrt{b^2 - 4ac}}{2a}\]

\subsection{Matrix Example}
A sample matrix:
\[
\begin{pmatrix}
a_{11} & a_{12} & a_{13} \\
a_{21} & a_{22} & a_{23} \\
a_{31} & a_{32} & a_{33}
\end{pmatrix}
\]

\section{Tables}
Below is a sample table showing data:

\begin{table}[h]
\centering
\begin{tabular}{@{}lrrr@{}}
\toprule
\textbf{Category} & \textbf{Value 1} & \textbf{Value 2} & \textbf{Total} \\
\midrule
Item A & 10.5 & 20.3 & 30.8 \\
Item B & 15.2 & 25.1 & 40.3 \\
Item C & 22.1 & 12.7 & 34.8 \\
\bottomrule
\end{tabular}
\caption{Sample Data Table}
\label{tab:sample}
\end{table}

\section{Lists}
\subsection{Bullet Points}
\begin{itemize}
  \item First item
  \item Second item
  \item Third item with sub-items:
    \begin{itemize}
      \item Sub-item 1
      \item Sub-item 2
    \end{itemize}
\end{itemize}

\subsection{Numbered List}
\begin{enumerate}
  \item First step
  \item Second step
  \item Third step
\end{enumerate}

\section{Citations and References}
When you use the citation search feature, references will be automatically managed. For example, you might cite a paper like this: \cite{example2024}. All citations will be automatically added to your \texttt{cite.bib} file.

\section{Conclusion}
This template demonstrates basic \LaTeX{} features including mathematical formulas, tables, and lists. The bibliography system is now set up to work seamlessly with Bibby AI's citation management. Feel free to modify it for your needs!

\bibliographystyle{plain}
\bibliography{cite}

\end{document}\documentclass{article}
\usepackage[utf8]{inputenc}
\usepackage{booktabs}
\usepackage{amsmath}
\usepackage{natbib}

\title{My Document}
\author{Your Name}
\date{\today}

\begin{document}

\maketitle

\section{Introduction}
This is a sample document demonstrating various \LaTeX{} features including mathematical equations, tables, and formatting. You can cite papers using the citation search feature, and they will be automatically added to your \texttt{cite.bib} file.

\section{Mathematical Expressions}
Here are some example equations:

\subsection{Quadratic Formula}
The quadratic formula is given by:
\[x = \frac{-b \pm \sqrt{b^2 - 4ac}}{2a}\]

\subsection{Matrix Example}
A sample matrix:
\[
\begin{pmatrix}
a_{11} & a_{12} & a_{13} \\
a_{21} & a_{22} & a_{23} \\
a_{31} & a_{32} & a_{33}
\end{pmatrix}
\]

\section{Tables}
Below is a sample table showing data:

\begin{table}[h]
\centering
\begin{tabular}{@{}lrrr@{}}
\toprule
\textbf{Category} & \textbf{Value 1} & \textbf{Value 2} & \textbf{Total} \\
\midrule
Item A & 10.5 & 20.3 & 30.8 \\
Item B & 15.2 & 25.1 & 40.3 \\
Item C & 22.1 & 12.7 & 34.8 \\
\bottomrule
\end{tabular}
\caption{Sample Data Table}
\label{tab:sample}
\end{table}

\section{Lists}
\subsection{Bullet Points}
\begin{itemize}
  \item First item
  \item Second item
  \item Third item with sub-items:
    \begin{itemize}
      \item Sub-item 1
      \item Sub-item 2
    \end{itemize}
\end{itemize}

\subsection{Numbered List}
\begin{enumerate}
  \item First step
  \item Second step
  \item Third step
\end{enumerate}

\section{Citations and References}
When you use the citation search feature, references will be automatically managed. For example, you might cite a paper like this: \cite{example2024}. All citations will be automatically added to your \texttt{cite.bib} file.

\section{Conclusion}
This template demonstrates basic \LaTeX{} features including mathematical formulas, tables, and lists. The bibliography system is now set up to work seamlessly with Bibby AI's citation management. Feel free to modify it for your needs!

\bibliographystyle{plain}
\bibliography{cite}

\end{document}\documentclass{article}
\usepackage[utf8]{inputenc}
\usepackage{booktabs}
\usepackage{amsmath}
\usepackage{natbib}

\title{My Document}
\author{Your Name}
\date{\today}

\begin{document}

\maketitle

\section{Introduction}
This is a sample document demonstrating various \LaTeX{} features including mathematical equations, tables, and formatting. You can cite papers using the citation search feature, and they will be automatically added to your \texttt{cite.bib} file.

\section{Mathematical Expressions}
Here are some example equations:

\subsection{Quadratic Formula}
The quadratic formula is given by:
\[x = \frac{-b \pm \sqrt{b^2 - 4ac}}{2a}\]

\subsection{Matrix Example}
A sample matrix:
\[
\begin{pmatrix}
a_{11} & a_{12} & a_{13} \\
a_{21} & a_{22} & a_{23} \\
a_{31} & a_{32} & a_{33}
\end{pmatrix}
\]

\section{Tables}
Below is a sample table showing data:

\begin{table}[h]
\centering
\begin{tabular}{@{}lrrr@{}}
\toprule
\textbf{Category} & \textbf{Value 1} & \textbf{Value 2} & \textbf{Total} \\
\midrule
Item A & 10.5 & 20.3 & 30.8 \\
Item B & 15.2 & 25.1 & 40.3 \\
Item C & 22.1 & 12.7 & 34.8 \\
\bottomrule
\end{tabular}
\caption{Sample Data Table}
\label{tab:sample}
\end{table}

\section{Lists}
\subsection{Bullet Points}
\begin{itemize}
  \item First item
  \item Second item
  \item Third item with sub-items:
    \begin{itemize}
      \item Sub-item 1
      \item Sub-item 2
    \end{itemize}
\end{itemize}

\subsection{Numbered List}
\begin{enumerate}
  \item First step
  \item Second step
  \item Third step
\end{enumerate}

\section{Citations and References}
When you use the citation search feature, references will be automatically managed. For example, you might cite a paper like this: \cite{example2024}. All citations will be automatically added to your \texttt{cite.bib} file.

\section{Conclusion}
This template demonstrates basic \LaTeX{} features including mathematical formulas, tables, and lists. The bibliography system is now set up to work seamlessly with Bibby AI's citation management. Feel free to modify it for your needs!

\bibliographystyle{plain}
\bibliography{cite}

\end{document}\documentclass{article}
\usepackage[utf8]{inputenc}
\usepackage{booktabs}
\usepackage{amsmath}
\usepackage{natbib}

\title{My Document}
\author{Your Name}
\date{\today}

\begin{document}

\maketitle

\section{Introduction}
This is a sample document demonstrating various \LaTeX{} features including mathematical equations, tables, and formatting. You can cite papers using the citation search feature, and they will be automatically added to your \texttt{cite.bib} file.

\section{Mathematical Expressions}
Here are some example equations:

\subsection{Quadratic Formula}
The quadratic formula is given by:
\[x = \frac{-b \pm \sqrt{b^2 - 4ac}}{2a}\]

\subsection{Matrix Example}
A sample matrix:
\[
\begin{pmatrix}
a_{11} & a_{12} & a_{13} \\
a_{21} & a_{22} & a_{23} \\
a_{31} & a_{32} & a_{33}
\end{pmatrix}
\]

\section{Tables}
Below is a sample table showing data:

\begin{table}[h]
\centering
\begin{tabular}{@{}lrrr@{}}
\toprule
\textbf{Category} & \textbf{Value 1} & \textbf{Value 2} & \textbf{Total} \\
\midrule
Item A & 10.5 & 20.3 & 30.8 \\
Item B & 15.2 & 25.1 & 40.3 \\
Item C & 22.1 & 12.7 & 34.8 \\
\bottomrule
\end{tabular}
\caption{Sample Data Table}
\label{tab:sample}
\end{table}

\section{Lists}
\subsection{Bullet Points}
\begin{itemize}
  \item First item
  \item Second item
  \item Third item with sub-items:
    \begin{itemize}
      \item Sub-item 1
      \item Sub-item 2
    \end{itemize}
\end{itemize}

\subsection{Numbered List}
\begin{enumerate}
  \item First step
  \item Second step
  \item Third step
\end{enumerate}

\section{Citations and References}
When you use the citation search feature, references will be automatically managed. For example, you might cite a paper like this: \cite{example2024}. All citations will be automatically added to your \texttt{cite.bib} file.

\section{Conclusion}
This template demonstrates basic \LaTeX{} features including mathematical formulas, tables, and lists. The bibliography system is now set up to work seamlessly with Bibby AI's citation management. Feel free to modify it for your needs!

\bibliographystyle{plain}
\bibliography{cite}

\end{document}\documentclass{article}
\usepackage[utf8]{inputenc}
\usepackage{booktabs}
\usepackage{amsmath}
\usepackage{natbib}

\title{My Document}
\author{Your Name}
\date{\today}

\begin{document}

\maketitle

\section{Introduction}
This is a sample document demonstrating various \LaTeX{} features including mathematical equations, tables, and formatting. You can cite papers using the citation search feature, and they will be automatically added to your \texttt{cite.bib} file.

\section{Mathematical Expressions}
Here are some example equations:

\subsection{Quadratic Formula}
The quadratic formula is given by:
\[x = \frac{-b \pm \sqrt{b^2 - 4ac}}{2a}\]

\subsection{Matrix Example}
A sample matrix:
\[
\begin{pmatrix}
a_{11} & a_{12} & a_{13} \\
a_{21} & a_{22} & a_{23} \\
a_{31} & a_{32} & a_{33}
\end{pmatrix}
\]

\section{Tables}
Below is a sample table showing data:

\begin{table}[h]
\centering
\begin{tabular}{@{}lrrr@{}}
\toprule
\textbf{Category} & \textbf{Value 1} & \textbf{Value 2} & \textbf{Total} \\
\midrule
Item A & 10.5 & 20.3 & 30.8 \\
Item B & 15.2 & 25.1 & 40.3 \\
Item C & 22.1 & 12.7 & 34.8 \\
\bottomrule
\end{tabular}
\caption{Sample Data Table}
\label{tab:sample}
\end{table}

\section{Lists}
\subsection{Bullet Points}
\begin{itemize}
  \item First item
  \item Second item
  \item Third item with sub-items:
    \begin{itemize}
      \item Sub-item 1
      \item Sub-item 2
    \end{itemize}
\end{itemize}

\subsection{Numbered List}
\begin{enumerate}
  \item First step
  \item Second step
  \item Third step
\end{enumerate}

\section{Citations and References}
When you use the citation search feature, references will be automatically managed. For example, you might cite a paper like this: \cite{example2024}. All citations will be automatically added to your \texttt{cite.bib} file.

\section{Conclusion}
This template demonstrates basic \LaTeX{} features including mathematical formulas, tables, and lists. The bibliography system is now set up to work seamlessly with Bibby AI's citation management. Feel free to modify it for your needs!

\bibliographystyle{plain}
\bibliography{cite}

\end{document}\documentclass{article}
\usepackage[utf8]{inputenc}
\usepackage{booktabs}
\usepackage{amsmath}
\usepackage{natbib}

\title{My Document}
\author{Your Name}
\date{\today}

\begin{document}

\maketitle

\section{Introduction}
This is a sample document demonstrating various \LaTeX{} features including mathematical equations, tables, and formatting. You can cite papers using the citation search feature, and they will be automatically added to your \texttt{cite.bib} file.

\section{Mathematical Expressions}
Here are some example equations:

\subsection{Quadratic Formula}
The quadratic formula is given by:
\[x = \frac{-b \pm \sqrt{b^2 - 4ac}}{2a}\]

\subsection{Matrix Example}
A sample matrix:
\[
\begin{pmatrix}
a_{11} & a_{12} & a_{13} \\
a_{21} & a_{22} & a_{23} \\
a_{31} & a_{32} & a_{33}
\end{pmatrix}
\]

\section{Tables}
Below is a sample table showing data:

\begin{table}[h]
\centering
\begin{tabular}{@{}lrrr@{}}
\toprule
\textbf{Category} & \textbf{Value 1} & \textbf{Value 2} & \textbf{Total} \\
\midrule
Item A & 10.5 & 20.3 & 30.8 \\
Item B & 15.2 & 25.1 & 40.3 \\
Item C & 22.1 & 12.7 & 34.8 \\
\bottomrule
\end{tabular}
\caption{Sample Data Table}
\label{tab:sample}
\end{table}

\section{Lists}
\subsection{Bullet Points}
\begin{itemize}
  \item First item
  \item Second item
  \item Third item with sub-items:
    \begin{itemize}
      \item Sub-item 1
      \item Sub-item 2
    \end{itemize}
\end{itemize}

\subsection{Numbered List}
\begin{enumerate}
  \item First step
  \item Second step
  \item Third step
\end{enumerate}

\section{Citations and References}
When you use the citation search feature, references will be automatically managed. For example, you might cite a paper like this: \cite{example2024}. All citations will be automatically added to your \texttt{cite.bib} file.

\section{Conclusion}
This template demonstrates basic \LaTeX{} features including mathematical formulas, tables, and lists. The bibliography system is now set up to work seamlessly with Bibby AI's citation management. Feel free to modify it for your needs!

\bibliographystyle{plain}
\bibliography{cite}

\end{document}\documentclass{article}
\usepackage[utf8]{inputenc}
\usepackage{booktabs}
\usepackage{amsmath}
\usepackage{natbib}

\title{My Document}
\author{Your Name}
\date{\today}

\begin{document}

\maketitle

\section{Introduction}
This is a sample document demonstrating various \LaTeX{} features including mathematical equations, tables, and formatting. You can cite papers using the citation search feature, and they will be automatically added to your \texttt{cite.bib} file.

\section{Mathematical Expressions}
Here are some example equations:

\subsection{Quadratic Formula}
The quadratic formula is given by:
\[x = \frac{-b \pm \sqrt{b^2 - 4ac}}{2a}\]

\subsection{Matrix Example}
A sample matrix:
\[
\begin{pmatrix}
a_{11} & a_{12} & a_{13} \\
a_{21} & a_{22} & a_{23} \\
a_{31} & a_{32} & a_{33}
\end{pmatrix}
\]

\section{Tables}
Below is a sample table showing data:

\begin{table}[h]
\centering
\begin{tabular}{@{}lrrr@{}}
\toprule
\textbf{Category} & \textbf{Value 1} & \textbf{Value 2} & \textbf{Total} \\
\midrule
Item A & 10.5 & 20.3 & 30.8 \\
Item B & 15.2 & 25.1 & 40.3 \\
Item C & 22.1 & 12.7 & 34.8 \\
\bottomrule
\end{tabular}
\caption{Sample Data Table}
\label{tab:sample}
\end{table}

\section{Lists}
\subsection{Bullet Points}
\begin{itemize}
  \item First item
  \item Second item
  \item Third item with sub-items:
    \begin{itemize}
      \item Sub-item 1
      \item Sub-item 2
    \end{itemize}
\end{itemize}

\subsection{Numbered List}
\begin{enumerate}
  \item First step
  \item Second step
  \item Third step
\end{enumerate}

\section{Citations and References}
When you use the citation search feature, references will be automatically managed. For example, you might cite a paper like this: \cite{example2024}. All citations will be automatically added to your \texttt{cite.bib} file.

\section{Conclusion}
This template demonstrates basic \LaTeX{} features including mathematical formulas, tables, and lists. The bibliography system is now set up to work seamlessly with Bibby AI's citation management. Feel free to modify it for your needs!

\bibliographystyle{plain}
\bibliography{cite}

\end{document}\documentclass{article}
\usepackage[utf8]{inputenc}
\usepackage{booktabs}
\usepackage{amsmath}
\usepackage{natbib}

\title{My Document}
\author{Your Name}
\date{\today}

\begin{document}

\maketitle

\section{Introduction}
This is a sample document demonstrating various \LaTeX{} features including mathematical equations, tables, and formatting. You can cite papers using the citation search feature, and they will be automatically added to your \texttt{cite.bib} file.

\section{Mathematical Expressions}
Here are some example equations:

\subsection{Quadratic Formula}
The quadratic formula is given by:
\[x = \frac{-b \pm \sqrt{b^2 - 4ac}}{2a}\]

\subsection{Matrix Example}
A sample matrix:
\[
\begin{pmatrix}
a_{11} & a_{12} & a_{13} \\
a_{21} & a_{22} & a_{23} \\
a_{31} & a_{32} & a_{33}
\end{pmatrix}
\]

\section{Tables}
Below is a sample table showing data:

\begin{table}[h]
\centering
\begin{tabular}{@{}lrrr@{}}
\toprule
\textbf{Category} & \textbf{Value 1} & \textbf{Value 2} & \textbf{Total} \\
\midrule
Item A & 10.5 & 20.3 & 30.8 \\
Item B & 15.2 & 25.1 & 40.3 \\
Item C & 22.1 & 12.7 & 34.8 \\
\bottomrule
\end{tabular}
\caption{Sample Data Table}
\label{tab:sample}
\end{table}

\section{Lists}
\subsection{Bullet Points}
\begin{itemize}
  \item First item
  \item Second item
  \item Third item with sub-items:
    \begin{itemize}
      \item Sub-item 1
      \item Sub-item 2
    \end{itemize}
\end{itemize}

\subsection{Numbered List}
\begin{enumerate}
  \item First step
  \item Second step
  \item Third step
\end{enumerate}

\section{Citations and References}
When you use the citation search feature, references will be automatically managed. For example, you might cite a paper like this: \cite{example2024}. All citations will be automatically added to your \texttt{cite.bib} file.

\section{Conclusion}
This template demonstrates basic \LaTeX{} features including mathematical formulas, tables, and lists. The bibliography system is now set up to work seamlessly with Bibby AI's citation management. Feel free to modify it for your needs!

\bibliographystyle{plain}
\bibliography{cite}

\end{document}\documentclass{article}
\usepackage[utf8]{inputenc}
\usepackage{booktabs}
\usepackage{amsmath}
\usepackage{natbib}

\title{My Document}
\author{Your Name}
\date{\today}

\begin{document}

\maketitle

\section{Introduction}
This is a sample document demonstrating various \LaTeX{} features including mathematical equations, tables, and formatting. You can cite papers using the citation search feature, and they will be automatically added to your \texttt{cite.bib} file.

\section{Mathematical Expressions}
Here are some example equations:

\subsection{Quadratic Formula}
The quadratic formula is given by:
\[x = \frac{-b \pm \sqrt{b^2 - 4ac}}{2a}\]

\subsection{Matrix Example}
A sample matrix:
\[
\begin{pmatrix}
a_{11} & a_{12} & a_{13} \\
a_{21} & a_{22} & a_{23} \\
a_{31} & a_{32} & a_{33}
\end{pmatrix}
\]

\section{Tables}
Below is a sample table showing data:

\begin{table}[h]
\centering
\begin{tabular}{@{}lrrr@{}}
\toprule
\textbf{Category} & \textbf{Value 1} & \textbf{Value 2} & \textbf{Total} \\
\midrule
Item A & 10.5 & 20.3 & 30.8 \\
Item B & 15.2 & 25.1 & 40.3 \\
Item C & 22.1 & 12.7 & 34.8 \\
\bottomrule
\end{tabular}
\caption{Sample Data Table}
\label{tab:sample}
\end{table}

\section{Lists}
\subsection{Bullet Points}
\begin{itemize}
  \item First item
  \item Second item
  \item Third item with sub-items:
    \begin{itemize}
      \item Sub-item 1
      \item Sub-item 2
    \end{itemize}
\end{itemize}

\subsection{Numbered List}
\begin{enumerate}
  \item First step
  \item Second step
  \item Third step
\end{enumerate}

\section{Citations and References}
When you use the citation search feature, references will be automatically managed. For example, you might cite a paper like this: \cite{example2024}. All citations will be automatically added to your \texttt{cite.bib} file.

\section{Conclusion}
This template demonstrates basic \LaTeX{} features including mathematical formulas, tables, and lists. The bibliography system is now set up to work seamlessly with Bibby AI's citation management. Feel free to modify it for your needs!

\bibliographystyle{plain}
\bibliography{cite}

\end{document}\documentclass{article}
\usepackage[utf8]{inputenc}
\usepackage{booktabs}
\usepackage{amsmath}
\usepackage{natbib}

\title{My Document}
\author{Your Name}
\date{\today}

\begin{document}

\maketitle

\section{Introduction}
This is a sample document demonstrating various \LaTeX{} features including mathematical equations, tables, and formatting. You can cite papers using the citation search feature, and they will be automatically added to your \texttt{cite.bib} file.

\section{Mathematical Expressions}
Here are some example equations:

\subsection{Quadratic Formula}
The quadratic formula is given by:
\[x = \frac{-b \pm \sqrt{b^2 - 4ac}}{2a}\]

\subsection{Matrix Example}
A sample matrix:
\[
\begin{pmatrix}
a_{11} & a_{12} & a_{13} \\
a_{21} & a_{22} & a_{23} \\
a_{31} & a_{32} & a_{33}
\end{pmatrix}
\]

\section{Tables}
Below is a sample table showing data:

\begin{table}[h]
\centering
\begin{tabular}{@{}lrrr@{}}
\toprule
\textbf{Category} & \textbf{Value 1} & \textbf{Value 2} & \textbf{Total} \\
\midrule
Item A & 10.5 & 20.3 & 30.8 \\
Item B & 15.2 & 25.1 & 40.3 \\
Item C & 22.1 & 12.7 & 34.8 \\
\bottomrule
\end{tabular}
\caption{Sample Data Table}
\label{tab:sample}
\end{table}

\section{Lists}
\subsection{Bullet Points}
\begin{itemize}
  \item First item
  \item Second item
  \item Third item with sub-items:
    \begin{itemize}
      \item Sub-item 1
      \item Sub-item 2
    \end{itemize}
\end{itemize}

\subsection{Numbered List}
\begin{enumerate}
  \item First step
  \item Second step
  \item Third step
\end{enumerate}

\section{Citations and References}
When you use the citation search feature, references will be automatically managed. For example, you might cite a paper like this: \cite{example2024}. All citations will be automatically added to your \texttt{cite.bib} file.

\section{Conclusion}
This template demonstrates basic \LaTeX{} features including mathematical formulas, tables, and lists. The bibliography system is now set up to work seamlessly with Bibby AI's citation management. Feel free to modify it for your needs!

\bibliographystyle{plain}
\bibliography{cite}

\end{document}\documentclass{article}
\usepackage[utf8]{inputenc}
\usepackage{booktabs}
\usepackage{amsmath}
\usepackage{natbib}

\title{My Document}
\author{Your Name}
\date{\today}

\begin{document}

\maketitle

\section{Introduction}
This is a sample document demonstrating various \LaTeX{} features including mathematical equations, tables, and formatting. You can cite papers using the citation search feature, and they will be automatically added to your \texttt{cite.bib} file.

\section{Mathematical Expressions}
Here are some example equations:

\subsection{Quadratic Formula}
The quadratic formula is given by:
\[x = \frac{-b \pm \sqrt{b^2 - 4ac}}{2a}\]

\subsection{Matrix Example}
A sample matrix:
\[
\begin{pmatrix}
a_{11} & a_{12} & a_{13} \\
a_{21} & a_{22} & a_{23} \\
a_{31} & a_{32} & a_{33}
\end{pmatrix}
\]

\section{Tables}
Below is a sample table showing data:

\begin{table}[h]
\centering
\begin{tabular}{@{}lrrr@{}}
\toprule
\textbf{Category} & \textbf{Value 1} & \textbf{Value 2} & \textbf{Total} \\
\midrule
Item A & 10.5 & 20.3 & 30.8 \\
Item B & 15.2 & 25.1 & 40.3 \\
Item C & 22.1 & 12.7 & 34.8 \\
\bottomrule
\end{tabular}
\caption{Sample Data Table}
\label{tab:sample}
\end{table}

\section{Lists}
\subsection{Bullet Points}
\begin{itemize}
  \item First item
  \item Second item
  \item Third item with sub-items:
    \begin{itemize}
      \item Sub-item 1
      \item Sub-item 2
    \end{itemize}
\end{itemize}

\subsection{Numbered List}
\begin{enumerate}
  \item First step
  \item Second step
  \item Third step
\end{enumerate}

\section{Citations and References}
When you use the citation search feature, references will be automatically managed. For example, you might cite a paper like this: \cite{example2024}. All citations will be automatically added to your \texttt{cite.bib} file.

\section{Conclusion}
This template demonstrates basic \LaTeX{} features including mathematical formulas, tables, and lists. The bibliography system is now set up to work seamlessly with Bibby AI's citation management. Feel free to modify it for your needs!

\bibliographystyle{plain}
\bibliography{cite}

\end{document}\documentclass{article}
\usepackage[utf8]{inputenc}
\usepackage{booktabs}
\usepackage{amsmath}
\usepackage{natbib}

\title{My Document}
\author{Your Name}
\date{\today}

\begin{document}

\maketitle

\section{Introduction}
This is a sample document demonstrating various \LaTeX{} features including mathematical equations, tables, and formatting. You can cite papers using the citation search feature, and they will be automatically added to your \texttt{cite.bib} file.

\section{Mathematical Expressions}
Here are some example equations:

\subsection{Quadratic Formula}
The quadratic formula is given by:
\[x = \frac{-b \pm \sqrt{b^2 - 4ac}}{2a}\]

\subsection{Matrix Example}
A sample matrix:
\[
\begin{pmatrix}
a_{11} & a_{12} & a_{13} \\
a_{21} & a_{22} & a_{23} \\
a_{31} & a_{32} & a_{33}
\end{pmatrix}
\]

\section{Tables}
Below is a sample table showing data:

\begin{table}[h]
\centering
\begin{tabular}{@{}lrrr@{}}
\toprule
\textbf{Category} & \textbf{Value 1} & \textbf{Value 2} & \textbf{Total} \\
\midrule
Item A & 10.5 & 20.3 & 30.8 \\
Item B & 15.2 & 25.1 & 40.3 \\
Item C & 22.1 & 12.7 & 34.8 \\
\bottomrule
\end{tabular}
\caption{Sample Data Table}
\label{tab:sample}
\end{table}

\section{Lists}
\subsection{Bullet Points}
\begin{itemize}
  \item First item
  \item Second item
  \item Third item with sub-items:
    \begin{itemize}
      \item Sub-item 1
      \item Sub-item 2
    \end{itemize}
\end{itemize}

\subsection{Numbered List}
\begin{enumerate}
  \item First step
  \item Second step
  \item Third step
\end{enumerate}

\section{Citations and References}
When you use the citation search feature, references will be automatically managed. For example, you might cite a paper like this: \cite{example2024}. All citations will be automatically added to your \texttt{cite.bib} file.

\section{Conclusion}
This template demonstrates basic \LaTeX{} features including mathematical formulas, tables, and lists. The bibliography system is now set up to work seamlessly with Bibby AI's citation management. Feel free to modify it for your needs!

\bibliographystyle{plain}
\bibliography{cite}

\end{document}\documentclass{article}
\usepackage[utf8]{inputenc}
\usepackage{booktabs}
\usepackage{amsmath}
\usepackage{natbib}

\title{My Document}
\author{Your Name}
\date{\today}

\begin{document}

\maketitle

\section{Introduction}
This is a sample document demonstrating various \LaTeX{} features including mathematical equations, tables, and formatting. You can cite papers using the citation search feature, and they will be automatically added to your \texttt{cite.bib} file.

\section{Mathematical Expressions}
Here are some example equations:

\subsection{Quadratic Formula}
The quadratic formula is given by:
\[x = \frac{-b \pm \sqrt{b^2 - 4ac}}{2a}\]

\subsection{Matrix Example}
A sample matrix:
\[
\begin{pmatrix}
a_{11} & a_{12} & a_{13} \\
a_{21} & a_{22} & a_{23} \\
a_{31} & a_{32} & a_{33}
\end{pmatrix}
\]

\section{Tables}
Below is a sample table showing data:

\begin{table}[h]
\centering
\begin{tabular}{@{}lrrr@{}}
\toprule
\textbf{Category} & \textbf{Value 1} & \textbf{Value 2} & \textbf{Total} \\
\midrule
Item A & 10.5 & 20.3 & 30.8 \\
Item B & 15.2 & 25.1 & 40.3 \\
Item C & 22.1 & 12.7 & 34.8 \\
\bottomrule
\end{tabular}
\caption{Sample Data Table}
\label{tab:sample}
\end{table}

\section{Lists}
\subsection{Bullet Points}
\begin{itemize}
  \item First item
  \item Second item
  \item Third item with sub-items:
    \begin{itemize}
      \item Sub-item 1
      \item Sub-item 2
    \end{itemize}
\end{itemize}

\subsection{Numbered List}
\begin{enumerate}
  \item First step
  \item Second step
  \item Third step
\end{enumerate}

\section{Citations and References}
When you use the citation search feature, references will be automatically managed. For example, you might cite a paper like this: \cite{example2024}. All citations will be automatically added to your \texttt{cite.bib} file.

\section{Conclusion}
This template demonstrates basic \LaTeX{} features including mathematical formulas, tables, and lists. The bibliography system is now set up to work seamlessly with Bibby AI's citation management. Feel free to modify it for your needs!

\bibliographystyle{plain}
\bibliography{cite}

\end{document}\documentclass{article}
\usepackage[utf8]{inputenc}
\usepackage{booktabs}
\usepackage{amsmath}
\usepackage{natbib}

\title{My Document}
\author{Your Name}
\date{\today}

\begin{document}

\maketitle

\section{Introduction}
This is a sample document demonstrating various \LaTeX{} features including mathematical equations, tables, and formatting. You can cite papers using the citation search feature, and they will be automatically added to your \texttt{cite.bib} file.

\section{Mathematical Expressions}
Here are some example equations:

\subsection{Quadratic Formula}
The quadratic formula is given by:
\[x = \frac{-b \pm \sqrt{b^2 - 4ac}}{2a}\]

\subsection{Matrix Example}
A sample matrix:
\[
\begin{pmatrix}
a_{11} & a_{12} & a_{13} \\
a_{21} & a_{22} & a_{23} \\
a_{31} & a_{32} & a_{33}
\end{pmatrix}
\]

\section{Tables}
Below is a sample table showing data:

\begin{table}[h]
\centering
\begin{tabular}{@{}lrrr@{}}
\toprule
\textbf{Category} & \textbf{Value 1} & \textbf{Value 2} & \textbf{Total} \\
\midrule
Item A & 10.5 & 20.3 & 30.8 \\
Item B & 15.2 & 25.1 & 40.3 \\
Item C & 22.1 & 12.7 & 34.8 \\
\bottomrule
\end{tabular}
\caption{Sample Data Table}
\label{tab:sample}
\end{table}

\section{Lists}
\subsection{Bullet Points}
\begin{itemize}
  \item First item
  \item Second item
  \item Third item with sub-items:
    \begin{itemize}
      \item Sub-item 1
      \item Sub-item 2
    \end{itemize}
\end{itemize}

\subsection{Numbered List}
\begin{enumerate}
  \item First step
  \item Second step
  \item Third step
\end{enumerate}

\section{Citations and References}
When you use the citation search feature, references will be automatically managed. For example, you might cite a paper like this: \cite{example2024}. All citations will be automatically added to your \texttt{cite.bib} file.

\section{Conclusion}
This template demonstrates basic \LaTeX{} features including mathematical formulas, tables, and lists. The bibliography system is now set up to work seamlessly with Bibby AI's citation management. Feel free to modify it for your needs!

\bibliographystyle{plain}
\bibliography{cite}

\end{document}\documentclass{article}
\usepackage[utf8]{inputenc}
\usepackage{booktabs}
\usepackage{amsmath}
\usepackage{natbib}

\title{My Document}
\author{Your Name}
\date{\today}

\begin{document}

\maketitle

\section{Introduction}
This is a sample document demonstrating various \LaTeX{} features including mathematical equations, tables, and formatting. You can cite papers using the citation search feature, and they will be automatically added to your \texttt{cite.bib} file.

\section{Mathematical Expressions}
Here are some example equations:

\subsection{Quadratic Formula}
The quadratic formula is given by:
\[x = \frac{-b \pm \sqrt{b^2 - 4ac}}{2a}\]

\subsection{Matrix Example}
A sample matrix:
\[
\begin{pmatrix}
a_{11} & a_{12} & a_{13} \\
a_{21} & a_{22} & a_{23} \\
a_{31} & a_{32} & a_{33}
\end{pmatrix}
\]

\section{Tables}
Below is a sample table showing data:

\begin{table}[h]
\centering
\begin{tabular}{@{}lrrr@{}}
\toprule
\textbf{Category} & \textbf{Value 1} & \textbf{Value 2} & \textbf{Total} \\
\midrule
Item A & 10.5 & 20.3 & 30.8 \\
Item B & 15.2 & 25.1 & 40.3 \\
Item C & 22.1 & 12.7 & 34.8 \\
\bottomrule
\end{tabular}
\caption{Sample Data Table}
\label{tab:sample}
\end{table}

\section{Lists}
\subsection{Bullet Points}
\begin{itemize}
  \item First item
  \item Second item
  \item Third item with sub-items:
    \begin{itemize}
      \item Sub-item 1
      \item Sub-item 2
    \end{itemize}
\end{itemize}

\subsection{Numbered List}
\begin{enumerate}
  \item First step
  \item Second step
  \item Third step
\end{enumerate}

\section{Citations and References}
When you use the citation search feature, references will be automatically managed. For example, you might cite a paper like this: \cite{example2024}. All citations will be automatically added to your \texttt{cite.bib} file.

\section{Conclusion}
This template demonstrates basic \LaTeX{} features including mathematical formulas, tables, and lists. The bibliography system is now set up to work seamlessly with Bibby AI's citation management. Feel free to modify it for your needs!

\bibliographystyle{plain}
\bibliography{cite}

\end{document}\documentclass{article}
\usepackage[utf8]{inputenc}
\usepackage{booktabs}
\usepackage{amsmath}
\usepackage{natbib}

\title{My Document}
\author{Your Name}
\date{\today}

\begin{document}

\maketitle

\section{Introduction}
This is a sample document demonstrating various \LaTeX{} features including mathematical equations, tables, and formatting. You can cite papers using the citation search feature, and they will be automatically added to your \texttt{cite.bib} file.

\section{Mathematical Expressions}
Here are some example equations:

\subsection{Quadratic Formula}
The quadratic formula is given by:
\[x = \frac{-b \pm \sqrt{b^2 - 4ac}}{2a}\]

\subsection{Matrix Example}
A sample matrix:
\[
\begin{pmatrix}
a_{11} & a_{12} & a_{13} \\
a_{21} & a_{22} & a_{23} \\
a_{31} & a_{32} & a_{33}
\end{pmatrix}
\]

\section{Tables}
Below is a sample table showing data:

\begin{table}[h]
\centering
\begin{tabular}{@{}lrrr@{}}
\toprule
\textbf{Category} & \textbf{Value 1} & \textbf{Value 2} & \textbf{Total} \\
\midrule
Item A & 10.5 & 20.3 & 30.8 \\
Item B & 15.2 & 25.1 & 40.3 \\
Item C & 22.1 & 12.7 & 34.8 \\
\bottomrule
\end{tabular}
\caption{Sample Data Table}
\label{tab:sample}
\end{table}

\section{Lists}
\subsection{Bullet Points}
\begin{itemize}
  \item First item
  \item Second item
  \item Third item with sub-items:
    \begin{itemize}
      \item Sub-item 1
      \item Sub-item 2
    \end{itemize}
\end{itemize}

\subsection{Numbered List}
\begin{enumerate}
  \item First step
  \item Second step
  \item Third step
\end{enumerate}

\section{Citations and References}
When you use the citation search feature, references will be automatically managed. For example, you might cite a paper like this: \cite{example2024}. All citations will be automatically added to your \texttt{cite.bib} file.

\section{Conclusion}
This template demonstrates basic \LaTeX{} features including mathematical formulas, tables, and lists. The bibliography system is now set up to work seamlessly with Bibby AI's citation management. Feel free to modify it for your needs!

\bibliographystyle{plain}
\bibliography{cite}

\end{document}\documentclass{article}
\usepackage[utf8]{inputenc}
\usepackage{booktabs}
\usepackage{amsmath}
\usepackage{natbib}

\title{My Document}
\author{Your Name}
\date{\today}

\begin{document}

\maketitle

\section{Introduction}
This is a sample document demonstrating various \LaTeX{} features including mathematical equations, tables, and formatting. You can cite papers using the citation search feature, and they will be automatically added to your \texttt{cite.bib} file.

\section{Mathematical Expressions}
Here are some example equations:

\subsection{Quadratic Formula}
The quadratic formula is given by:
\[x = \frac{-b \pm \sqrt{b^2 - 4ac}}{2a}\]

\subsection{Matrix Example}
A sample matrix:
\[
\begin{pmatrix}
a_{11} & a_{12} & a_{13} \\
a_{21} & a_{22} & a_{23} \\
a_{31} & a_{32} & a_{33}
\end{pmatrix}
\]

\section{Tables}
Below is a sample table showing data:

\begin{table}[h]
\centering
\begin{tabular}{@{}lrrr@{}}
\toprule
\textbf{Category} & \textbf{Value 1} & \textbf{Value 2} & \textbf{Total} \\
\midrule
Item A & 10.5 & 20.3 & 30.8 \\
Item B & 15.2 & 25.1 & 40.3 \\
Item C & 22.1 & 12.7 & 34.8 \\
\bottomrule
\end{tabular}
\caption{Sample Data Table}
\label{tab:sample}
\end{table}

\section{Lists}
\subsection{Bullet Points}
\begin{itemize}
  \item First item
  \item Second item
  \item Third item with sub-items:
    \begin{itemize}
      \item Sub-item 1
      \item Sub-item 2
    \end{itemize}
\end{itemize}

\subsection{Numbered List}
\begin{enumerate}
  \item First step
  \item Second step
  \item Third step
\end{enumerate}

\section{Citations and References}
When you use the citation search feature, references will be automatically managed. For example, you might cite a paper like this: \cite{example2024}. All citations will be automatically added to your \texttt{cite.bib} file.

\section{Conclusion}
This template demonstrates basic \LaTeX{} features including mathematical formulas, tables, and lists. The bibliography system is now set up to work seamlessly with Bibby AI's citation management. Feel free to modify it for your needs!

\bibliographystyle{plain}
\bibliography{cite}

\end{document}\documentclass{article}
\usepackage[utf8]{inputenc}
\usepackage{booktabs}
\usepackage{amsmath}
\usepackage{natbib}

\title{My Document}
\author{Your Name}
\date{\today}

\begin{document}

\maketitle

\section{Introduction}
This is a sample document demonstrating various \LaTeX{} features including mathematical equations, tables, and formatting. You can cite papers using the citation search feature, and they will be automatically added to your \texttt{cite.bib} file.

\section{Mathematical Expressions}
Here are some example equations:

\subsection{Quadratic Formula}
The quadratic formula is given by:
\[x = \frac{-b \pm \sqrt{b^2 - 4ac}}{2a}\]

\subsection{Matrix Example}
A sample matrix:
\[
\begin{pmatrix}
a_{11} & a_{12} & a_{13} \\
a_{21} & a_{22} & a_{23} \\
a_{31} & a_{32} & a_{33}
\end{pmatrix}
\]

\section{Tables}
Below is a sample table showing data:

\begin{table}[h]
\centering
\begin{tabular}{@{}lrrr@{}}
\toprule
\textbf{Category} & \textbf{Value 1} & \textbf{Value 2} & \textbf{Total} \\
\midrule
Item A & 10.5 & 20.3 & 30.8 \\
Item B & 15.2 & 25.1 & 40.3 \\
Item C & 22.1 & 12.7 & 34.8 \\
\bottomrule
\end{tabular}
\caption{Sample Data Table}
\label{tab:sample}
\end{table}

\section{Lists}
\subsection{Bullet Points}
\begin{itemize}
  \item First item
  \item Second item
  \item Third item with sub-items:
    \begin{itemize}
      \item Sub-item 1
      \item Sub-item 2
    \end{itemize}
\end{itemize}

\subsection{Numbered List}
\begin{enumerate}
  \item First step
  \item Second step
  \item Third step
\end{enumerate}

\section{Citations and References}
When you use the citation search feature, references will be automatically managed. For example, you might cite a paper like this: \cite{example2024}. All citations will be automatically added to your \texttt{cite.bib} file.

\section{Conclusion}
This template demonstrates basic \LaTeX{} features including mathematical formulas, tables, and lists. The bibliography system is now set up to work seamlessly with Bibby AI's citation management. Feel free to modify it for your needs!

\bibliographystyle{plain}
\bibliography{cite}

\end{document}\documentclass{article}
\usepackage[utf8]{inputenc}
\usepackage{booktabs}
\usepackage{amsmath}
\usepackage{natbib}

\title{My Document}
\author{Your Name}
\date{\today}

\begin{document}

\maketitle

\section{Introduction}
This is a sample document demonstrating various \LaTeX{} features including mathematical equations, tables, and formatting. You can cite papers using the citation search feature, and they will be automatically added to your \texttt{cite.bib} file.

\section{Mathematical Expressions}
Here are some example equations:

\subsection{Quadratic Formula}
The quadratic formula is given by:
\[x = \frac{-b \pm \sqrt{b^2 - 4ac}}{2a}\]

\subsection{Matrix Example}
A sample matrix:
\[
\begin{pmatrix}
a_{11} & a_{12} & a_{13} \\
a_{21} & a_{22} & a_{23} \\
a_{31} & a_{32} & a_{33}
\end{pmatrix}
\]

\section{Tables}
Below is a sample table showing data:

\begin{table}[h]
\centering
\begin{tabular}{@{}lrrr@{}}
\toprule
\textbf{Category} & \textbf{Value 1} & \textbf{Value 2} & \textbf{Total} \\
\midrule
Item A & 10.5 & 20.3 & 30.8 \\
Item B & 15.2 & 25.1 & 40.3 \\
Item C & 22.1 & 12.7 & 34.8 \\
\bottomrule
\end{tabular}
\caption{Sample Data Table}
\label{tab:sample}
\end{table}

\section{Lists}
\subsection{Bullet Points}
\begin{itemize}
  \item First item
  \item Second item
  \item Third item with sub-items:
    \begin{itemize}
      \item Sub-item 1
      \item Sub-item 2
    \end{itemize}
\end{itemize}

\subsection{Numbered List}
\begin{enumerate}
  \item First step
  \item Second step
  \item Third step
\end{enumerate}

\section{Citations and References}
When you use the citation search feature, references will be automatically managed. For example, you might cite a paper like this: \cite{example2024}. All citations will be automatically added to your \texttt{cite.bib} file.

\section{Conclusion}
This template demonstrates basic \LaTeX{} features including mathematical formulas, tables, and lists. The bibliography system is now set up to work seamlessly with Bibby AI's citation management. Feel free to modify it for your needs!

\bibliographystyle{plain}
\bibliography{cite}

\end{document}\documentclass{article}
\usepackage[utf8]{inputenc}
\usepackage{booktabs}
\usepackage{amsmath}
\usepackage{natbib}

\title{My Document}
\author{Your Name}
\date{\today}

\begin{document}

\maketitle

\section{Introduction}
This is a sample document demonstrating various \LaTeX{} features including mathematical equations, tables, and formatting. You can cite papers using the citation search feature, and they will be automatically added to your \texttt{cite.bib} file.

\section{Mathematical Expressions}
Here are some example equations:

\subsection{Quadratic Formula}
The quadratic formula is given by:
\[x = \frac{-b \pm \sqrt{b^2 - 4ac}}{2a}\]

\subsection{Matrix Example}
A sample matrix:
\[
\begin{pmatrix}
a_{11} & a_{12} & a_{13} \\
a_{21} & a_{22} & a_{23} \\
a_{31} & a_{32} & a_{33}
\end{pmatrix}
\]

\section{Tables}
Below is a sample table showing data:

\begin{table}[h]
\centering
\begin{tabular}{@{}lrrr@{}}
\toprule
\textbf{Category} & \textbf{Value 1} & \textbf{Value 2} & \textbf{Total} \\
\midrule
Item A & 10.5 & 20.3 & 30.8 \\
Item B & 15.2 & 25.1 & 40.3 \\
Item C & 22.1 & 12.7 & 34.8 \\
\bottomrule
\end{tabular}
\caption{Sample Data Table}
\label{tab:sample}
\end{table}

\section{Lists}
\subsection{Bullet Points}
\begin{itemize}
  \item First item
  \item Second item
  \item Third item with sub-items:
    \begin{itemize}
      \item Sub-item 1
      \item Sub-item 2
    \end{itemize}
\end{itemize}

\subsection{Numbered List}
\begin{enumerate}
  \item First step
  \item Second step
  \item Third step
\end{enumerate}

\section{Citations and References}
When you use the citation search feature, references will be automatically managed. For example, you might cite a paper like this: \cite{example2024}. All citations will be automatically added to your \texttt{cite.bib} file.

\section{Conclusion}
This template demonstrates basic \LaTeX{} features including mathematical formulas, tables, and lists. The bibliography system is now set up to work seamlessly with Bibby AI's citation management. Feel free to modify it for your needs!

\bibliographystyle{plain}
\bibliography{cite}

\end{document}\documentclass{article}
\usepackage[utf8]{inputenc}
\usepackage{booktabs}
\usepackage{amsmath}
\usepackage{natbib}

\title{My Document}
\author{Your Name}
\date{\today}

\begin{document}

\maketitle

\section{Introduction}
This is a sample document demonstrating various \LaTeX{} features including mathematical equations, tables, and formatting. You can cite papers using the citation search feature, and they will be automatically added to your \texttt{cite.bib} file.

\section{Mathematical Expressions}
Here are some example equations:

\subsection{Quadratic Formula}
The quadratic formula is given by:
\[x = \frac{-b \pm \sqrt{b^2 - 4ac}}{2a}\]

\subsection{Matrix Example}
A sample matrix:
\[
\begin{pmatrix}
a_{11} & a_{12} & a_{13} \\
a_{21} & a_{22} & a_{23} \\
a_{31} & a_{32} & a_{33}
\end{pmatrix}
\]

\section{Tables}
Below is a sample table showing data:

\begin{table}[h]
\centering
\begin{tabular}{@{}lrrr@{}}
\toprule
\textbf{Category} & \textbf{Value 1} & \textbf{Value 2} & \textbf{Total} \\
\midrule
Item A & 10.5 & 20.3 & 30.8 \\
Item B & 15.2 & 25.1 & 40.3 \\
Item C & 22.1 & 12.7 & 34.8 \\
\bottomrule
\end{tabular}
\caption{Sample Data Table}
\label{tab:sample}
\end{table}

\section{Lists}
\subsection{Bullet Points}
\begin{itemize}
  \item First item
  \item Second item
  \item Third item with sub-items:
    \begin{itemize}
      \item Sub-item 1
      \item Sub-item 2
    \end{itemize}
\end{itemize}

\subsection{Numbered List}
\begin{enumerate}
  \item First step
  \item Second step
  \item Third step
\end{enumerate}

\section{Citations and References}
When you use the citation search feature, references will be automatically managed. For example, you might cite a paper like this: \cite{example2024}. All citations will be automatically added to your \texttt{cite.bib} file.

\section{Conclusion}
This template demonstrates basic \LaTeX{} features including mathematical formulas, tables, and lists. The bibliography system is now set up to work seamlessly with Bibby AI's citation management. Feel free to modify it for your needs!

\bibliographystyle{plain}
\bibliography{cite}

\end{document}\documentclass{article}
\usepackage[utf8]{inputenc}
\usepackage{booktabs}
\usepackage{amsmath}
\usepackage{natbib}

\title{My Document}
\author{Your Name}
\date{\today}

\begin{document}

\maketitle

\section{Introduction}
This is a sample document demonstrating various \LaTeX{} features including mathematical equations, tables, and formatting. You can cite papers using the citation search feature, and they will be automatically added to your \texttt{cite.bib} file.

\section{Mathematical Expressions}
Here are some example equations:

\subsection{Quadratic Formula}
The quadratic formula is given by:
\[x = \frac{-b \pm \sqrt{b^2 - 4ac}}{2a}\]

\subsection{Matrix Example}
A sample matrix:
\[
\begin{pmatrix}
a_{11} & a_{12} & a_{13} \\
a_{21} & a_{22} & a_{23} \\
a_{31} & a_{32} & a_{33}
\end{pmatrix}
\]

\section{Tables}
Below is a sample table showing data:

\begin{table}[h]
\centering
\begin{tabular}{@{}lrrr@{}}
\toprule
\textbf{Category} & \textbf{Value 1} & \textbf{Value 2} & \textbf{Total} \\
\midrule
Item A & 10.5 & 20.3 & 30.8 \\
Item B & 15.2 & 25.1 & 40.3 \\
Item C & 22.1 & 12.7 & 34.8 \\
\bottomrule
\end{tabular}
\caption{Sample Data Table}
\label{tab:sample}
\end{table}

\section{Lists}
\subsection{Bullet Points}
\begin{itemize}
  \item First item
  \item Second item
  \item Third item with sub-items:
    \begin{itemize}
      \item Sub-item 1
      \item Sub-item 2
    \end{itemize}
\end{itemize}

\subsection{Numbered List}
\begin{enumerate}
  \item First step
  \item Second step
  \item Third step
\end{enumerate}

\section{Citations and References}
When you use the citation search feature, references will be automatically managed. For example, you might cite a paper like this: \cite{example2024}. All citations will be automatically added to your \texttt{cite.bib} file.

\section{Conclusion}
This template demonstrates basic \LaTeX{} features including mathematical formulas, tables, and lists. The bibliography system is now set up to work seamlessly with Bibby AI's citation management. Feel free to modify it for your needs!

\bibliographystyle{plain}
\bibliography{cite}

\end{document}\documentclass{article}
\usepackage[utf8]{inputenc}
\usepackage{booktabs}
\usepackage{amsmath}
\usepackage{natbib}

\title{My Document}
\author{Your Name}
\date{\today}

\begin{document}

\maketitle

\section{Introduction}
This is a sample document demonstrating various \LaTeX{} features including mathematical equations, tables, and formatting. You can cite papers using the citation search feature, and they will be automatically added to your \texttt{cite.bib} file.

\section{Mathematical Expressions}
Here are some example equations:

\subsection{Quadratic Formula}
The quadratic formula is given by:
\[x = \frac{-b \pm \sqrt{b^2 - 4ac}}{2a}\]

\subsection{Matrix Example}
A sample matrix:
\[
\begin{pmatrix}
a_{11} & a_{12} & a_{13} \\
a_{21} & a_{22} & a_{23} \\
a_{31} & a_{32} & a_{33}
\end{pmatrix}
\]

\section{Tables}
Below is a sample table showing data:

\begin{table}[h]
\centering
\begin{tabular}{@{}lrrr@{}}
\toprule
\textbf{Category} & \textbf{Value 1} & \textbf{Value 2} & \textbf{Total} \\
\midrule
Item A & 10.5 & 20.3 & 30.8 \\
Item B & 15.2 & 25.1 & 40.3 \\
Item C & 22.1 & 12.7 & 34.8 \\
\bottomrule
\end{tabular}
\caption{Sample Data Table}
\label{tab:sample}
\end{table}

\section{Lists}
\subsection{Bullet Points}
\begin{itemize}
  \item First item
  \item Second item
  \item Third item with sub-items:
    \begin{itemize}
      \item Sub-item 1
      \item Sub-item 2
    \end{itemize}
\end{itemize}

\subsection{Numbered List}
\begin{enumerate}
  \item First step
  \item Second step
  \item Third step
\end{enumerate}

\section{Citations and References}
When you use the citation search feature, references will be automatically managed. For example, you might cite a paper like this: \cite{example2024}. All citations will be automatically added to your \texttt{cite.bib} file.

\section{Conclusion}
This template demonstrates basic \LaTeX{} features including mathematical formulas, tables, and lists. The bibliography system is now set up to work seamlessly with Bibby AI's citation management. Feel free to modify it for your needs!

\bibliographystyle{plain}
\bibliography{cite}

\end{document}\documentclass{article}
\usepackage[utf8]{inputenc}
\usepackage{booktabs}
\usepackage{amsmath}
\usepackage{natbib}

\title{My Document}
\author{Your Name}
\date{\today}

\begin{document}

\maketitle

\section{Introduction}
This is a sample document demonstrating various \LaTeX{} features including mathematical equations, tables, and formatting. You can cite papers using the citation search feature, and they will be automatically added to your \texttt{cite.bib} file.

\section{Mathematical Expressions}
Here are some example equations:

\subsection{Quadratic Formula}
The quadratic formula is given by:
\[x = \frac{-b \pm \sqrt{b^2 - 4ac}}{2a}\]

\subsection{Matrix Example}
A sample matrix:
\[
\begin{pmatrix}
a_{11} & a_{12} & a_{13} \\
a_{21} & a_{22} & a_{23} \\
a_{31} & a_{32} & a_{33}
\end{pmatrix}
\]

\section{Tables}
Below is a sample table showing data:

\begin{table}[h]
\centering
\begin{tabular}{@{}lrrr@{}}
\toprule
\textbf{Category} & \textbf{Value 1} & \textbf{Value 2} & \textbf{Total} \\
\midrule
Item A & 10.5 & 20.3 & 30.8 \\
Item B & 15.2 & 25.1 & 40.3 \\
Item C & 22.1 & 12.7 & 34.8 \\
\bottomrule
\end{tabular}
\caption{Sample Data Table}
\label{tab:sample}
\end{table}

\section{Lists}
\subsection{Bullet Points}
\begin{itemize}
  \item First item
  \item Second item
  \item Third item with sub-items:
    \begin{itemize}
      \item Sub-item 1
      \item Sub-item 2
    \end{itemize}
\end{itemize}

\subsection{Numbered List}
\begin{enumerate}
  \item First step
  \item Second step
  \item Third step
\end{enumerate}

\section{Citations and References}
When you use the citation search feature, references will be automatically managed. For example, you might cite a paper like this: \cite{example2024}. All citations will be automatically added to your \texttt{cite.bib} file.

\section{Conclusion}
This template demonstrates basic \LaTeX{} features including mathematical formulas, tables, and lists. The bibliography system is now set up to work seamlessly with Bibby AI's citation management. Feel free to modify it for your needs!

\bibliographystyle{plain}
\bibliography{cite}

\end{document}\documentclass{article}
\usepackage[utf8]{inputenc}
\usepackage{booktabs}
\usepackage{amsmath}
\usepackage{natbib}

\title{My Document}
\author{Your Name}
\date{\today}

\begin{document}

\maketitle

\section{Introduction}
This is a sample document demonstrating various \LaTeX{} features including mathematical equations, tables, and formatting. You can cite papers using the citation search feature, and they will be automatically added to your \texttt{cite.bib} file.

\section{Mathematical Expressions}
Here are some example equations:

\subsection{Quadratic Formula}
The quadratic formula is given by:
\[x = \frac{-b \pm \sqrt{b^2 - 4ac}}{2a}\]

\subsection{Matrix Example}
A sample matrix:
\[
\begin{pmatrix}
a_{11} & a_{12} & a_{13} \\
a_{21} & a_{22} & a_{23} \\
a_{31} & a_{32} & a_{33}
\end{pmatrix}
\]

\section{Tables}
Below is a sample table showing data:

\begin{table}[h]
\centering
\begin{tabular}{@{}lrrr@{}}
\toprule
\textbf{Category} & \textbf{Value 1} & \textbf{Value 2} & \textbf{Total} \\
\midrule
Item A & 10.5 & 20.3 & 30.8 \\
Item B & 15.2 & 25.1 & 40.3 \\
Item C & 22.1 & 12.7 & 34.8 \\
\bottomrule
\end{tabular}
\caption{Sample Data Table}
\label{tab:sample}
\end{table}

\section{Lists}
\subsection{Bullet Points}
\begin{itemize}
  \item First item
  \item Second item
  \item Third item with sub-items:
    \begin{itemize}
      \item Sub-item 1
      \item Sub-item 2
    \end{itemize}
\end{itemize}

\subsection{Numbered List}
\begin{enumerate}
  \item First step
  \item Second step
  \item Third step
\end{enumerate}

\section{Citations and References}
When you use the citation search feature, references will be automatically managed. For example, you might cite a paper like this: \cite{example2024}. All citations will be automatically added to your \texttt{cite.bib} file.

\section{Conclusion}
This template demonstrates basic \LaTeX{} features including mathematical formulas, tables, and lists. The bibliography system is now set up to work seamlessly with Bibby AI's citation management. Feel free to modify it for your needs!

\bibliographystyle{plain}
\bibliography{cite}

\end{document}\documentclass{article}
\usepackage[utf8]{inputenc}
\usepackage{booktabs}
\usepackage{amsmath}
\usepackage{natbib}

\title{My Document}
\author{Your Name}
\date{\today}

\begin{document}

\maketitle

\section{Introduction}
This is a sample document demonstrating various \LaTeX{} features including mathematical equations, tables, and formatting. You can cite papers using the citation search feature, and they will be automatically added to your \texttt{cite.bib} file.

\section{Mathematical Expressions}
Here are some example equations:

\subsection{Quadratic Formula}
The quadratic formula is given by:
\[x = \frac{-b \pm \sqrt{b^2 - 4ac}}{2a}\]

\subsection{Matrix Example}
A sample matrix:
\[
\begin{pmatrix}
a_{11} & a_{12} & a_{13} \\
a_{21} & a_{22} & a_{23} \\
a_{31} & a_{32} & a_{33}
\end{pmatrix}
\]

\section{Tables}
Below is a sample table showing data:

\begin{table}[h]
\centering
\begin{tabular}{@{}lrrr@{}}
\toprule
\textbf{Category} & \textbf{Value 1} & \textbf{Value 2} & \textbf{Total} \\
\midrule
Item A & 10.5 & 20.3 & 30.8 \\
Item B & 15.2 & 25.1 & 40.3 \\
Item C & 22.1 & 12.7 & 34.8 \\
\bottomrule
\end{tabular}
\caption{Sample Data Table}
\label{tab:sample}
\end{table}

\section{Lists}
\subsection{Bullet Points}
\begin{itemize}
  \item First item
  \item Second item
  \item Third item with sub-items:
    \begin{itemize}
      \item Sub-item 1
      \item Sub-item 2
    \end{itemize}
\end{itemize}

\subsection{Numbered List}
\begin{enumerate}
  \item First step
  \item Second step
  \item Third step
\end{enumerate}

\section{Citations and References}
When you use the citation search feature, references will be automatically managed. For example, you might cite a paper like this: \cite{example2024}. All citations will be automatically added to your \texttt{cite.bib} file.

\section{Conclusion}
This template demonstrates basic \LaTeX{} features including mathematical formulas, tables, and lists. The bibliography system is now set up to work seamlessly with Bibby AI's citation management. Feel free to modify it for your needs!

\bibliographystyle{plain}
\bibliography{cite}

\end{document}\documentclass{article}
\usepackage[utf8]{inputenc}
\usepackage{booktabs}
\usepackage{amsmath}
\usepackage{natbib}

\title{My Document}
\author{Your Name}
\date{\today}

\begin{document}

\maketitle

\section{Introduction}
This is a sample document demonstrating various \LaTeX{} features including mathematical equations, tables, and formatting. You can cite papers using the citation search feature, and they will be automatically added to your \texttt{cite.bib} file.

\section{Mathematical Expressions}
Here are some example equations:

\subsection{Quadratic Formula}
The quadratic formula is given by:
\[x = \frac{-b \pm \sqrt{b^2 - 4ac}}{2a}\]

\subsection{Matrix Example}
A sample matrix:
\[
\begin{pmatrix}
a_{11} & a_{12} & a_{13} \\
a_{21} & a_{22} & a_{23} \\
a_{31} & a_{32} & a_{33}
\end{pmatrix}
\]

\section{Tables}
Below is a sample table showing data:

\begin{table}[h]
\centering
\begin{tabular}{@{}lrrr@{}}
\toprule
\textbf{Category} & \textbf{Value 1} & \textbf{Value 2} & \textbf{Total} \\
\midrule
Item A & 10.5 & 20.3 & 30.8 \\
Item B & 15.2 & 25.1 & 40.3 \\
Item C & 22.1 & 12.7 & 34.8 \\
\bottomrule
\end{tabular}
\caption{Sample Data Table}
\label{tab:sample}
\end{table}

\section{Lists}
\subsection{Bullet Points}
\begin{itemize}
  \item First item
  \item Second item
  \item Third item with sub-items:
    \begin{itemize}
      \item Sub-item 1
      \item Sub-item 2
    \end{itemize}
\end{itemize}

\subsection{Numbered List}
\begin{enumerate}
  \item First step
  \item Second step
  \item Third step
\end{enumerate}

\section{Citations and References}
When you use the citation search feature, references will be automatically managed. For example, you might cite a paper like this: \cite{example2024}. All citations will be automatically added to your \texttt{cite.bib} file.

\section{Conclusion}
This template demonstrates basic \LaTeX{} features including mathematical formulas, tables, and lists. The bibliography system is now set up to work seamlessly with Bibby AI's citation management. Feel free to modify it for your needs!

\bibliographystyle{plain}
\bibliography{cite}

\end{document}\documentclass{article}
\usepackage[utf8]{inputenc}
\usepackage{booktabs}
\usepackage{amsmath}
\usepackage{natbib}

\title{My Document}
\author{Your Name}
\date{\today}

\begin{document}

\maketitle

\section{Introduction}
This is a sample document demonstrating various \LaTeX{} features including mathematical equations, tables, and formatting. You can cite papers using the citation search feature, and they will be automatically added to your \texttt{cite.bib} file.

\section{Mathematical Expressions}
Here are some example equations:

\subsection{Quadratic Formula}
The quadratic formula is given by:
\[x = \frac{-b \pm \sqrt{b^2 - 4ac}}{2a}\]

\subsection{Matrix Example}
A sample matrix:
\[
\begin{pmatrix}
a_{11} & a_{12} & a_{13} \\
a_{21} & a_{22} & a_{23} \\
a_{31} & a_{32} & a_{33}
\end{pmatrix}
\]

\section{Tables}
Below is a sample table showing data:

\begin{table}[h]
\centering
\begin{tabular}{@{}lrrr@{}}
\toprule
\textbf{Category} & \textbf{Value 1} & \textbf{Value 2} & \textbf{Total} \\
\midrule
Item A & 10.5 & 20.3 & 30.8 \\
Item B & 15.2 & 25.1 & 40.3 \\
Item C & 22.1 & 12.7 & 34.8 \\
\bottomrule
\end{tabular}
\caption{Sample Data Table}
\label{tab:sample}
\end{table}

\section{Lists}
\subsection{Bullet Points}
\begin{itemize}
  \item First item
  \item Second item
  \item Third item with sub-items:
    \begin{itemize}
      \item Sub-item 1
      \item Sub-item 2
    \end{itemize}
\end{itemize}

\subsection{Numbered List}
\begin{enumerate}
  \item First step
  \item Second step
  \item Third step
\end{enumerate}

\section{Citations and References}
When you use the citation search feature, references will be automatically managed. For example, you might cite a paper like this: \cite{example2024}. All citations will be automatically added to your \texttt{cite.bib} file.

\section{Conclusion}
This template demonstrates basic \LaTeX{} features including mathematical formulas, tables, and lists. The bibliography system is now set up to work seamlessly with Bibby AI's citation management. Feel free to modify it for your needs!

\bibliographystyle{plain}
\bibliography{cite}

\end{document}\documentclass{article}
\usepackage[utf8]{inputenc}
\usepackage{booktabs}
\usepackage{amsmath}
\usepackage{natbib}

\title{My Document}
\author{Your Name}
\date{\today}

\begin{document}

\maketitle

\section{Introduction}
This is a sample document demonstrating various \LaTeX{} features including mathematical equations, tables, and formatting. You can cite papers using the citation search feature, and they will be automatically added to your \texttt{cite.bib} file.

\section{Mathematical Expressions}
Here are some example equations:

\subsection{Quadratic Formula}
The quadratic formula is given by:
\[x = \frac{-b \pm \sqrt{b^2 - 4ac}}{2a}\]

\subsection{Matrix Example}
A sample matrix:
\[
\begin{pmatrix}
a_{11} & a_{12} & a_{13} \\
a_{21} & a_{22} & a_{23} \\
a_{31} & a_{32} & a_{33}
\end{pmatrix}
\]

\section{Tables}
Below is a sample table showing data:

\begin{table}[h]
\centering
\begin{tabular}{@{}lrrr@{}}
\toprule
\textbf{Category} & \textbf{Value 1} & \textbf{Value 2} & \textbf{Total} \\
\midrule
Item A & 10.5 & 20.3 & 30.8 \\
Item B & 15.2 & 25.1 & 40.3 \\
Item C & 22.1 & 12.7 & 34.8 \\
\bottomrule
\end{tabular}
\caption{Sample Data Table}
\label{tab:sample}
\end{table}

\section{Lists}
\subsection{Bullet Points}
\begin{itemize}
  \item First item
  \item Second item
  \item Third item with sub-items:
    \begin{itemize}
      \item Sub-item 1
      \item Sub-item 2
    \end{itemize}
\end{itemize}

\subsection{Numbered List}
\begin{enumerate}
  \item First step
  \item Second step
  \item Third step
\end{enumerate}

\section{Citations and References}
When you use the citation search feature, references will be automatically managed. For example, you might cite a paper like this: \cite{example2024}. All citations will be automatically added to your \texttt{cite.bib} file.

\section{Conclusion}
This template demonstrates basic \LaTeX{} features including mathematical formulas, tables, and lists. The bibliography system is now set up to work seamlessly with Bibby AI's citation management. Feel free to modify it for your needs!

\bibliographystyle{plain}
\bibliography{cite}

\end{document}\documentclass{article}
\usepackage[utf8]{inputenc}
\usepackage{booktabs}
\usepackage{amsmath}
\usepackage{natbib}

\title{My Document}
\author{Your Name}
\date{\today}

\begin{document}

\maketitle

\section{Introduction}
This is a sample document demonstrating various \LaTeX{} features including mathematical equations, tables, and formatting. You can cite papers using the citation search feature, and they will be automatically added to your \texttt{cite.bib} file.

\section{Mathematical Expressions}
Here are some example equations:

\subsection{Quadratic Formula}
The quadratic formula is given by:
\[x = \frac{-b \pm \sqrt{b^2 - 4ac}}{2a}\]

\subsection{Matrix Example}
A sample matrix:
\[
\begin{pmatrix}
a_{11} & a_{12} & a_{13} \\
a_{21} & a_{22} & a_{23} \\
a_{31} & a_{32} & a_{33}
\end{pmatrix}
\]

\section{Tables}
Below is a sample table showing data:

\begin{table}[h]
\centering
\begin{tabular}{@{}lrrr@{}}
\toprule
\textbf{Category} & \textbf{Value 1} & \textbf{Value 2} & \textbf{Total} \\
\midrule
Item A & 10.5 & 20.3 & 30.8 \\
Item B & 15.2 & 25.1 & 40.3 \\
Item C & 22.1 & 12.7 & 34.8 \\
\bottomrule
\end{tabular}
\caption{Sample Data Table}
\label{tab:sample}
\end{table}

\section{Lists}
\subsection{Bullet Points}
\begin{itemize}
  \item First item
  \item Second item
  \item Third item with sub-items:
    \begin{itemize}
      \item Sub-item 1
      \item Sub-item 2
    \end{itemize}
\end{itemize}

\subsection{Numbered List}
\begin{enumerate}
  \item First step
  \item Second step
  \item Third step
\end{enumerate}

\section{Citations and References}
When you use the citation search feature, references will be automatically managed. For example, you might cite a paper like this: \cite{example2024}. All citations will be automatically added to your \texttt{cite.bib} file.

\section{Conclusion}
This template demonstrates basic \LaTeX{} features including mathematical formulas, tables, and lists. The bibliography system is now set up to work seamlessly with Bibby AI's citation management. Feel free to modify it for your needs!

\bibliographystyle{plain}
\bibliography{cite}

\end{document}\documentclass{article}
\usepackage[utf8]{inputenc}
\usepackage{booktabs}
\usepackage{amsmath}
\usepackage{natbib}

\title{My Document}
\author{Your Name}
\date{\today}

\begin{document}

\maketitle

\section{Introduction}
This is a sample document demonstrating various \LaTeX{} features including mathematical equations, tables, and formatting. You can cite papers using the citation search feature, and they will be automatically added to your \texttt{cite.bib} file.

\section{Mathematical Expressions}
Here are some example equations:

\subsection{Quadratic Formula}
The quadratic formula is given by:
\[x = \frac{-b \pm \sqrt{b^2 - 4ac}}{2a}\]

\subsection{Matrix Example}
A sample matrix:
\[
\begin{pmatrix}
a_{11} & a_{12} & a_{13} \\
a_{21} & a_{22} & a_{23} \\
a_{31} & a_{32} & a_{33}
\end{pmatrix}
\]

\section{Tables}
Below is a sample table showing data:

\begin{table}[h]
\centering
\begin{tabular}{@{}lrrr@{}}
\toprule
\textbf{Category} & \textbf{Value 1} & \textbf{Value 2} & \textbf{Total} \\
\midrule
Item A & 10.5 & 20.3 & 30.8 \\
Item B & 15.2 & 25.1 & 40.3 \\
Item C & 22.1 & 12.7 & 34.8 \\
\bottomrule
\end{tabular}
\caption{Sample Data Table}
\label{tab:sample}
\end{table}

\section{Lists}
\subsection{Bullet Points}
\begin{itemize}
  \item First item
  \item Second item
  \item Third item with sub-items:
    \begin{itemize}
      \item Sub-item 1
      \item Sub-item 2
    \end{itemize}
\end{itemize}

\subsection{Numbered List}
\begin{enumerate}
  \item First step
  \item Second step
  \item Third step
\end{enumerate}

\section{Citations and References}
When you use the citation search feature, references will be automatically managed. For example, you might cite a paper like this: \cite{example2024}. All citations will be automatically added to your \texttt{cite.bib} file.

\section{Conclusion}
This template demonstrates basic \LaTeX{} features including mathematical formulas, tables, and lists. The bibliography system is now set up to work seamlessly with Bibby AI's citation management. Feel free to modify it for your needs!

\bibliographystyle{plain}
\bibliography{cite}

\end{document}\documentclass{article}
\usepackage[utf8]{inputenc}
\usepackage{booktabs}
\usepackage{amsmath}
\usepackage{natbib}

\title{My Document}
\author{Your Name}
\date{\today}

\begin{document}

\maketitle

\section{Introduction}
This is a sample document demonstrating various \LaTeX{} features including mathematical equations, tables, and formatting. You can cite papers using the citation search feature, and they will be automatically added to your \texttt{cite.bib} file.

\section{Mathematical Expressions}
Here are some example equations:

\subsection{Quadratic Formula}
The quadratic formula is given by:
\[x = \frac{-b \pm \sqrt{b^2 - 4ac}}{2a}\]

\subsection{Matrix Example}
A sample matrix:
\[
\begin{pmatrix}
a_{11} & a_{12} & a_{13} \\
a_{21} & a_{22} & a_{23} \\
a_{31} & a_{32} & a_{33}
\end{pmatrix}
\]

\section{Tables}
Below is a sample table showing data:

\begin{table}[h]
\centering
\begin{tabular}{@{}lrrr@{}}
\toprule
\textbf{Category} & \textbf{Value 1} & \textbf{Value 2} & \textbf{Total} \\
\midrule
Item A & 10.5 & 20.3 & 30.8 \\
Item B & 15.2 & 25.1 & 40.3 \\
Item C & 22.1 & 12.7 & 34.8 \\
\bottomrule
\end{tabular}
\caption{Sample Data Table}
\label{tab:sample}
\end{table}

\section{Lists}
\subsection{Bullet Points}
\begin{itemize}
  \item First item
  \item Second item
  \item Third item with sub-items:
    \begin{itemize}
      \item Sub-item 1
      \item Sub-item 2
    \end{itemize}
\end{itemize}

\subsection{Numbered List}
\begin{enumerate}
  \item First step
  \item Second step
  \item Third step
\end{enumerate}

\section{Citations and References}
When you use the citation search feature, references will be automatically managed. For example, you might cite a paper like this: \cite{example2024}. All citations will be automatically added to your \texttt{cite.bib} file.

\section{Conclusion}
This template demonstrates basic \LaTeX{} features including mathematical formulas, tables, and lists. The bibliography system is now set up to work seamlessly with Bibby AI's citation management. Feel free to modify it for your needs!

\bibliographystyle{plain}
\bibliography{cite}

\end{document}\documentclass{article}
\usepackage[utf8]{inputenc}
\usepackage{booktabs}
\usepackage{amsmath}
\usepackage{natbib}

\title{My Document}
\author{Your Name}
\date{\today}

\begin{document}

\maketitle

\section{Introduction}
This is a sample document demonstrating various \LaTeX{} features including mathematical equations, tables, and formatting. You can cite papers using the citation search feature, and they will be automatically added to your \texttt{cite.bib} file.

\section{Mathematical Expressions}
Here are some example equations:

\subsection{Quadratic Formula}
The quadratic formula is given by:
\[x = \frac{-b \pm \sqrt{b^2 - 4ac}}{2a}\]

\subsection{Matrix Example}
A sample matrix:
\[
\begin{pmatrix}
a_{11} & a_{12} & a_{13} \\
a_{21} & a_{22} & a_{23} \\
a_{31} & a_{32} & a_{33}
\end{pmatrix}
\]

\section{Tables}
Below is a sample table showing data:

\begin{table}[h]
\centering
\begin{tabular}{@{}lrrr@{}}
\toprule
\textbf{Category} & \textbf{Value 1} & \textbf{Value 2} & \textbf{Total} \\
\midrule
Item A & 10.5 & 20.3 & 30.8 \\
Item B & 15.2 & 25.1 & 40.3 \\
Item C & 22.1 & 12.7 & 34.8 \\
\bottomrule
\end{tabular}
\caption{Sample Data Table}
\label{tab:sample}
\end{table}

\section{Lists}
\subsection{Bullet Points}
\begin{itemize}
  \item First item
  \item Second item
  \item Third item with sub-items:
    \begin{itemize}
      \item Sub-item 1
      \item Sub-item 2
    \end{itemize}
\end{itemize}

\subsection{Numbered List}
\begin{enumerate}
  \item First step
  \item Second step
  \item Third step
\end{enumerate}

\section{Citations and References}
When you use the citation search feature, references will be automatically managed. For example, you might cite a paper like this: \cite{example2024}. All citations will be automatically added to your \texttt{cite.bib} file.

\section{Conclusion}
This template demonstrates basic \LaTeX{} features including mathematical formulas, tables, and lists. The bibliography system is now set up to work seamlessly with Bibby AI's citation management. Feel free to modify it for your needs!

\bibliographystyle{plain}
\bibliography{cite}

\end{document}\documentclass{article}
\usepackage[utf8]{inputenc}
\usepackage{booktabs}
\usepackage{amsmath}
\usepackage{natbib}

\title{My Document}
\author{Your Name}
\date{\today}

\begin{document}

\maketitle

\section{Introduction}
This is a sample document demonstrating various \LaTeX{} features including mathematical equations, tables, and formatting. You can cite papers using the citation search feature, and they will be automatically added to your \texttt{cite.bib} file.

\section{Mathematical Expressions}
Here are some example equations:

\subsection{Quadratic Formula}
The quadratic formula is given by:
\[x = \frac{-b \pm \sqrt{b^2 - 4ac}}{2a}\]

\subsection{Matrix Example}
A sample matrix:
\[
\begin{pmatrix}
a_{11} & a_{12} & a_{13} \\
a_{21} & a_{22} & a_{23} \\
a_{31} & a_{32} & a_{33}
\end{pmatrix}
\]

\section{Tables}
Below is a sample table showing data:

\begin{table}[h]
\centering
\begin{tabular}{@{}lrrr@{}}
\toprule
\textbf{Category} & \textbf{Value 1} & \textbf{Value 2} & \textbf{Total} \\
\midrule
Item A & 10.5 & 20.3 & 30.8 \\
Item B & 15.2 & 25.1 & 40.3 \\
Item C & 22.1 & 12.7 & 34.8 \\
\bottomrule
\end{tabular}
\caption{Sample Data Table}
\label{tab:sample}
\end{table}

\section{Lists}
\subsection{Bullet Points}
\begin{itemize}
  \item First item
  \item Second item
  \item Third item with sub-items:
    \begin{itemize}
      \item Sub-item 1
      \item Sub-item 2
    \end{itemize}
\end{itemize}

\subsection{Numbered List}
\begin{enumerate}
  \item First step
  \item Second step
  \item Third step
\end{enumerate}

\section{Citations and References}
When you use the citation search feature, references will be automatically managed. For example, you might cite a paper like this: \cite{example2024}. All citations will be automatically added to your \texttt{cite.bib} file.

\section{Conclusion}
This template demonstrates basic \LaTeX{} features including mathematical formulas, tables, and lists. The bibliography system is now set up to work seamlessly with Bibby AI's citation management. Feel free to modify it for your needs!

\bibliographystyle{plain}
\bibliography{cite}

\end{document}\documentclass{article}
\usepackage[utf8]{inputenc}
\usepackage{booktabs}
\usepackage{amsmath}
\usepackage{natbib}

\title{My Document}
\author{Your Name}
\date{\today}

\begin{document}

\maketitle

\section{Introduction}
This is a sample document demonstrating various \LaTeX{} features including mathematical equations, tables, and formatting. You can cite papers using the citation search feature, and they will be automatically added to your \texttt{cite.bib} file.

\section{Mathematical Expressions}
Here are some example equations:

\subsection{Quadratic Formula}
The quadratic formula is given by:
\[x = \frac{-b \pm \sqrt{b^2 - 4ac}}{2a}\]

\subsection{Matrix Example}
A sample matrix:
\[
\begin{pmatrix}
a_{11} & a_{12} & a_{13} \\
a_{21} & a_{22} & a_{23} \\
a_{31} & a_{32} & a_{33}
\end{pmatrix}
\]

\section{Tables}
Below is a sample table showing data:

\begin{table}[h]
\centering
\begin{tabular}{@{}lrrr@{}}
\toprule
\textbf{Category} & \textbf{Value 1} & \textbf{Value 2} & \textbf{Total} \\
\midrule
Item A & 10.5 & 20.3 & 30.8 \\
Item B & 15.2 & 25.1 & 40.3 \\
Item C & 22.1 & 12.7 & 34.8 \\
\bottomrule
\end{tabular}
\caption{Sample Data Table}
\label{tab:sample}
\end{table}

\section{Lists}
\subsection{Bullet Points}
\begin{itemize}
  \item First item
  \item Second item
  \item Third item with sub-items:
    \begin{itemize}
      \item Sub-item 1
      \item Sub-item 2
    \end{itemize}
\end{itemize}

\subsection{Numbered List}
\begin{enumerate}
  \item First step
  \item Second step
  \item Third step
\end{enumerate}

\section{Citations and References}
When you use the citation search feature, references will be automatically managed. For example, you might cite a paper like this: \cite{example2024}. All citations will be automatically added to your \texttt{cite.bib} file.

\section{Conclusion}
This template demonstrates basic \LaTeX{} features including mathematical formulas, tables, and lists. The bibliography system is now set up to work seamlessly with Bibby AI's citation management. Feel free to modify it for your needs!

\bibliographystyle{plain}
\bibliography{cite}

\end{document}\documentclass{article}
\usepackage[utf8]{inputenc}
\usepackage{booktabs}
\usepackage{amsmath}
\usepackage{natbib}

\title{My Document}
\author{Your Name}
\date{\today}

\begin{document}

\maketitle

\section{Introduction}
This is a sample document demonstrating various \LaTeX{} features including mathematical equations, tables, and formatting. You can cite papers using the citation search feature, and they will be automatically added to your \texttt{cite.bib} file.

\section{Mathematical Expressions}
Here are some example equations:

\subsection{Quadratic Formula}
The quadratic formula is given by:
\[x = \frac{-b \pm \sqrt{b^2 - 4ac}}{2a}\]

\subsection{Matrix Example}
A sample matrix:
\[
\begin{pmatrix}
a_{11} & a_{12} & a_{13} \\
a_{21} & a_{22} & a_{23} \\
a_{31} & a_{32} & a_{33}
\end{pmatrix}
\]

\section{Tables}
Below is a sample table showing data:

\begin{table}[h]
\centering
\begin{tabular}{@{}lrrr@{}}
\toprule
\textbf{Category} & \textbf{Value 1} & \textbf{Value 2} & \textbf{Total} \\
\midrule
Item A & 10.5 & 20.3 & 30.8 \\
Item B & 15.2 & 25.1 & 40.3 \\
Item C & 22.1 & 12.7 & 34.8 \\
\bottomrule
\end{tabular}
\caption{Sample Data Table}
\label{tab:sample}
\end{table}

\section{Lists}
\subsection{Bullet Points}
\begin{itemize}
  \item First item
  \item Second item
  \item Third item with sub-items:
    \begin{itemize}
      \item Sub-item 1
      \item Sub-item 2
    \end{itemize}
\end{itemize}

\subsection{Numbered List}
\begin{enumerate}
  \item First step
  \item Second step
  \item Third step
\end{enumerate}

\section{Citations and References}
When you use the citation search feature, references will be automatically managed. For example, you might cite a paper like this: \cite{example2024}. All citations will be automatically added to your \texttt{cite.bib} file.

\section{Conclusion}
This template demonstrates basic \LaTeX{} features including mathematical formulas, tables, and lists. The bibliography system is now set up to work seamlessly with Bibby AI's citation management. Feel free to modify it for your needs!

\bibliographystyle{plain}
\bibliography{cite}

\end{document}\documentclass{article}
\usepackage[utf8]{inputenc}
\usepackage{booktabs}
\usepackage{amsmath}
\usepackage{natbib}

\title{My Document}
\author{Your Name}
\date{\today}

\begin{document}

\maketitle

\section{Introduction}
This is a sample document demonstrating various \LaTeX{} features including mathematical equations, tables, and formatting. You can cite papers using the citation search feature, and they will be automatically added to your \texttt{cite.bib} file.

\section{Mathematical Expressions}
Here are some example equations:

\subsection{Quadratic Formula}
The quadratic formula is given by:
\[x = \frac{-b \pm \sqrt{b^2 - 4ac}}{2a}\]

\subsection{Matrix Example}
A sample matrix:
\[
\begin{pmatrix}
a_{11} & a_{12} & a_{13} \\
a_{21} & a_{22} & a_{23} \\
a_{31} & a_{32} & a_{33}
\end{pmatrix}
\]

\section{Tables}
Below is a sample table showing data:

\begin{table}[h]
\centering
\begin{tabular}{@{}lrrr@{}}
\toprule
\textbf{Category} & \textbf{Value 1} & \textbf{Value 2} & \textbf{Total} \\
\midrule
Item A & 10.5 & 20.3 & 30.8 \\
Item B & 15.2 & 25.1 & 40.3 \\
Item C & 22.1 & 12.7 & 34.8 \\
\bottomrule
\end{tabular}
\caption{Sample Data Table}
\label{tab:sample}
\end{table}

\section{Lists}
\subsection{Bullet Points}
\begin{itemize}
  \item First item
  \item Second item
  \item Third item with sub-items:
    \begin{itemize}
      \item Sub-item 1
      \item Sub-item 2
    \end{itemize}
\end{itemize}

\subsection{Numbered List}
\begin{enumerate}
  \item First step
  \item Second step
  \item Third step
\end{enumerate}

\section{Citations and References}
When you use the citation search feature, references will be automatically managed. For example, you might cite a paper like this: \cite{example2024}. All citations will be automatically added to your \texttt{cite.bib} file.

\section{Conclusion}
This template demonstrates basic \LaTeX{} features including mathematical formulas, tables, and lists. The bibliography system is now set up to work seamlessly with Bibby AI's citation management. Feel free to modify it for your needs!

\bibliographystyle{plain}
\bibliography{cite}

\end{document}\documentclass{article}
\usepackage[utf8]{inputenc}
\usepackage{booktabs}
\usepackage{amsmath}
\usepackage{natbib}

\title{My Document}
\author{Your Name}
\date{\today}

\begin{document}

\maketitle

\section{Introduction}
This is a sample document demonstrating various \LaTeX{} features including mathematical equations, tables, and formatting. You can cite papers using the citation search feature, and they will be automatically added to your \texttt{cite.bib} file.

\section{Mathematical Expressions}
Here are some example equations:

\subsection{Quadratic Formula}
The quadratic formula is given by:
\[x = \frac{-b \pm \sqrt{b^2 - 4ac}}{2a}\]

\subsection{Matrix Example}
A sample matrix:
\[
\begin{pmatrix}
a_{11} & a_{12} & a_{13} \\
a_{21} & a_{22} & a_{23} \\
a_{31} & a_{32} & a_{33}
\end{pmatrix}
\]

\section{Tables}
Below is a sample table showing data:

\begin{table}[h]
\centering
\begin{tabular}{@{}lrrr@{}}
\toprule
\textbf{Category} & \textbf{Value 1} & \textbf{Value 2} & \textbf{Total} \\
\midrule
Item A & 10.5 & 20.3 & 30.8 \\
Item B & 15.2 & 25.1 & 40.3 \\
Item C & 22.1 & 12.7 & 34.8 \\
\bottomrule
\end{tabular}
\caption{Sample Data Table}
\label{tab:sample}
\end{table}

\section{Lists}
\subsection{Bullet Points}
\begin{itemize}
  \item First item
  \item Second item
  \item Third item with sub-items:
    \begin{itemize}
      \item Sub-item 1
      \item Sub-item 2
    \end{itemize}
\end{itemize}

\subsection{Numbered List}
\begin{enumerate}
  \item First step
  \item Second step
  \item Third step
\end{enumerate}

\section{Citations and References}
When you use the citation search feature, references will be automatically managed. For example, you might cite a paper like this: \cite{example2024}. All citations will be automatically added to your \texttt{cite.bib} file.

\section{Conclusion}
This template demonstrates basic \LaTeX{} features including mathematical formulas, tables, and lists. The bibliography system is now set up to work seamlessly with Bibby AI's citation management. Feel free to modify it for your needs!

\bibliographystyle{plain}
\bibliography{cite}

\end{document}\documentclass{article}
\usepackage[utf8]{inputenc}
\usepackage{booktabs}
\usepackage{amsmath}
\usepackage{natbib}

\title{My Document}
\author{Your Name}
\date{\today}

\begin{document}

\maketitle

\section{Introduction}
This is a sample document demonstrating various \LaTeX{} features including mathematical equations, tables, and formatting. You can cite papers using the citation search feature, and they will be automatically added to your \texttt{cite.bib} file.

\section{Mathematical Expressions}
Here are some example equations:

\subsection{Quadratic Formula}
The quadratic formula is given by:
\[x = \frac{-b \pm \sqrt{b^2 - 4ac}}{2a}\]

\subsection{Matrix Example}
A sample matrix:
\[
\begin{pmatrix}
a_{11} & a_{12} & a_{13} \\
a_{21} & a_{22} & a_{23} \\
a_{31} & a_{32} & a_{33}
\end{pmatrix}
\]

\section{Tables}
Below is a sample table showing data:

\begin{table}[h]
\centering
\begin{tabular}{@{}lrrr@{}}
\toprule
\textbf{Category} & \textbf{Value 1} & \textbf{Value 2} & \textbf{Total} \\
\midrule
Item A & 10.5 & 20.3 & 30.8 \\
Item B & 15.2 & 25.1 & 40.3 \\
Item C & 22.1 & 12.7 & 34.8 \\
\bottomrule
\end{tabular}
\caption{Sample Data Table}
\label{tab:sample}
\end{table}

\section{Lists}
\subsection{Bullet Points}
\begin{itemize}
  \item First item
  \item Second item
  \item Third item with sub-items:
    \begin{itemize}
      \item Sub-item 1
      \item Sub-item 2
    \end{itemize}
\end{itemize}

\subsection{Numbered List}
\begin{enumerate}
  \item First step
  \item Second step
  \item Third step
\end{enumerate}

\section{Citations and References}
When you use the citation search feature, references will be automatically managed. For example, you might cite a paper like this: \cite{example2024}. All citations will be automatically added to your \texttt{cite.bib} file.

\section{Conclusion}
This template demonstrates basic \LaTeX{} features including mathematical formulas, tables, and lists. The bibliography system is now set up to work seamlessly with Bibby AI's citation management. Feel free to modify it for your needs!

\bibliographystyle{plain}
\bibliography{cite}

\end{document}\documentclass{article}
\usepackage[utf8]{inputenc}
\usepackage{booktabs}
\usepackage{amsmath}
\usepackage{natbib}

\title{My Document}
\author{Your Name}
\date{\today}

\begin{document}

\maketitle

\section{Introduction}
This is a sample document demonstrating various \LaTeX{} features including mathematical equations, tables, and formatting. You can cite papers using the citation search feature, and they will be automatically added to your \texttt{cite.bib} file.

\section{Mathematical Expressions}
Here are some example equations:

\subsection{Quadratic Formula}
The quadratic formula is given by:
\[x = \frac{-b \pm \sqrt{b^2 - 4ac}}{2a}\]

\subsection{Matrix Example}
A sample matrix:
\[
\begin{pmatrix}
a_{11} & a_{12} & a_{13} \\
a_{21} & a_{22} & a_{23} \\
a_{31} & a_{32} & a_{33}
\end{pmatrix}
\]

\section{Tables}
Below is a sample table showing data:

\begin{table}[h]
\centering
\begin{tabular}{@{}lrrr@{}}
\toprule
\textbf{Category} & \textbf{Value 1} & \textbf{Value 2} & \textbf{Total} \\
\midrule
Item A & 10.5 & 20.3 & 30.8 \\
Item B & 15.2 & 25.1 & 40.3 \\
Item C & 22.1 & 12.7 & 34.8 \\
\bottomrule
\end{tabular}
\caption{Sample Data Table}
\label{tab:sample}
\end{table}

\section{Lists}
\subsection{Bullet Points}
\begin{itemize}
  \item First item
  \item Second item
  \item Third item with sub-items:
    \begin{itemize}
      \item Sub-item 1
      \item Sub-item 2
    \end{itemize}
\end{itemize}

\subsection{Numbered List}
\begin{enumerate}
  \item First step
  \item Second step
  \item Third step
\end{enumerate}

\section{Citations and References}
When you use the citation search feature, references will be automatically managed. For example, you might cite a paper like this: \cite{example2024}. All citations will be automatically added to your \texttt{cite.bib} file.

\section{Conclusion}
This template demonstrates basic \LaTeX{} features including mathematical formulas, tables, and lists. The bibliography system is now set up to work seamlessly with Bibby AI's citation management. Feel free to modify it for your needs!

\bibliographystyle{plain}
\bibliography{cite}

\end{document}\documentclass{article}
\usepackage[utf8]{inputenc}
\usepackage{booktabs}
\usepackage{amsmath}
\usepackage{natbib}

\title{My Document}
\author{Your Name}
\date{\today}

\begin{document}

\maketitle

\section{Introduction}
This is a sample document demonstrating various \LaTeX{} features including mathematical equations, tables, and formatting. You can cite papers using the citation search feature, and they will be automatically added to your \texttt{cite.bib} file.

\section{Mathematical Expressions}
Here are some example equations:

\subsection{Quadratic Formula}
The quadratic formula is given by:
\[x = \frac{-b \pm \sqrt{b^2 - 4ac}}{2a}\]

\subsection{Matrix Example}
A sample matrix:
\[
\begin{pmatrix}
a_{11} & a_{12} & a_{13} \\
a_{21} & a_{22} & a_{23} \\
a_{31} & a_{32} & a_{33}
\end{pmatrix}
\]

\section{Tables}
Below is a sample table showing data:

\begin{table}[h]
\centering
\begin{tabular}{@{}lrrr@{}}
\toprule
\textbf{Category} & \textbf{Value 1} & \textbf{Value 2} & \textbf{Total} \\
\midrule
Item A & 10.5 & 20.3 & 30.8 \\
Item B & 15.2 & 25.1 & 40.3 \\
Item C & 22.1 & 12.7 & 34.8 \\
\bottomrule
\end{tabular}
\caption{Sample Data Table}
\label{tab:sample}
\end{table}

\section{Lists}
\subsection{Bullet Points}
\begin{itemize}
  \item First item
  \item Second item
  \item Third item with sub-items:
    \begin{itemize}
      \item Sub-item 1
      \item Sub-item 2
    \end{itemize}
\end{itemize}

\subsection{Numbered List}
\begin{enumerate}
  \item First step
  \item Second step
  \item Third step
\end{enumerate}

\section{Citations and References}
When you use the citation search feature, references will be automatically managed. For example, you might cite a paper like this: \cite{example2024}. All citations will be automatically added to your \texttt{cite.bib} file.

\section{Conclusion}
This template demonstrates basic \LaTeX{} features including mathematical formulas, tables, and lists. The bibliography system is now set up to work seamlessly with Bibby AI's citation management. Feel free to modify it for your needs!

\bibliographystyle{plain}
\bibliography{cite}

\end{document}\documentclass{article}
\usepackage[utf8]{inputenc}
\usepackage{booktabs}
\usepackage{amsmath}
\usepackage{natbib}

\title{My Document}
\author{Your Name}
\date{\today}

\begin{document}

\maketitle

\section{Introduction}
This is a sample document demonstrating various \LaTeX{} features including mathematical equations, tables, and formatting. You can cite papers using the citation search feature, and they will be automatically added to your \texttt{cite.bib} file.

\section{Mathematical Expressions}
Here are some example equations:

\subsection{Quadratic Formula}
The quadratic formula is given by:
\[x = \frac{-b \pm \sqrt{b^2 - 4ac}}{2a}\]

\subsection{Matrix Example}
A sample matrix:
\[
\begin{pmatrix}
a_{11} & a_{12} & a_{13} \\
a_{21} & a_{22} & a_{23} \\
a_{31} & a_{32} & a_{33}
\end{pmatrix}
\]

\section{Tables}
Below is a sample table showing data:

\begin{table}[h]
\centering
\begin{tabular}{@{}lrrr@{}}
\toprule
\textbf{Category} & \textbf{Value 1} & \textbf{Value 2} & \textbf{Total} \\
\midrule
Item A & 10.5 & 20.3 & 30.8 \\
Item B & 15.2 & 25.1 & 40.3 \\
Item C & 22.1 & 12.7 & 34.8 \\
\bottomrule
\end{tabular}
\caption{Sample Data Table}
\label{tab:sample}
\end{table}

\section{Lists}
\subsection{Bullet Points}
\begin{itemize}
  \item First item
  \item Second item
  \item Third item with sub-items:
    \begin{itemize}
      \item Sub-item 1
      \item Sub-item 2
    \end{itemize}
\end{itemize}

\subsection{Numbered List}
\begin{enumerate}
  \item First step
  \item Second step
  \item Third step
\end{enumerate}

\section{Citations and References}
When you use the citation search feature, references will be automatically managed. For example, you might cite a paper like this: \cite{example2024}. All citations will be automatically added to your \texttt{cite.bib} file.

\section{Conclusion}
This template demonstrates basic \LaTeX{} features including mathematical formulas, tables, and lists. The bibliography system is now set up to work seamlessly with Bibby AI's citation management. Feel free to modify it for your needs!

\bibliographystyle{plain}
\bibliography{cite}

\end{document}\documentclass{article}
\usepackage[utf8]{inputenc}
\usepackage{booktabs}
\usepackage{amsmath}
\usepackage{natbib}

\title{My Document}
\author{Your Name}
\date{\today}

\begin{document}

\maketitle

\section{Introduction}
This is a sample document demonstrating various \LaTeX{} features including mathematical equations, tables, and formatting. You can cite papers using the citation search feature, and they will be automatically added to your \texttt{cite.bib} file.

\section{Mathematical Expressions}
Here are some example equations:

\subsection{Quadratic Formula}
The quadratic formula is given by:
\[x = \frac{-b \pm \sqrt{b^2 - 4ac}}{2a}\]

\subsection{Matrix Example}
A sample matrix:
\[
\begin{pmatrix}
a_{11} & a_{12} & a_{13} \\
a_{21} & a_{22} & a_{23} \\
a_{31} & a_{32} & a_{33}
\end{pmatrix}
\]

\section{Tables}
Below is a sample table showing data:

\begin{table}[h]
\centering
\begin{tabular}{@{}lrrr@{}}
\toprule
\textbf{Category} & \textbf{Value 1} & \textbf{Value 2} & \textbf{Total} \\
\midrule
Item A & 10.5 & 20.3 & 30.8 \\
Item B & 15.2 & 25.1 & 40.3 \\
Item C & 22.1 & 12.7 & 34.8 \\
\bottomrule
\end{tabular}
\caption{Sample Data Table}
\label{tab:sample}
\end{table}

\section{Lists}
\subsection{Bullet Points}
\begin{itemize}
  \item First item
  \item Second item
  \item Third item with sub-items:
    \begin{itemize}
      \item Sub-item 1
      \item Sub-item 2
    \end{itemize}
\end{itemize}

\subsection{Numbered List}
\begin{enumerate}
  \item First step
  \item Second step
  \item Third step
\end{enumerate}

\section{Citations and References}
When you use the citation search feature, references will be automatically managed. For example, you might cite a paper like this: \cite{example2024}. All citations will be automatically added to your \texttt{cite.bib} file.

\section{Conclusion}
This template demonstrates basic \LaTeX{} features including mathematical formulas, tables, and lists. The bibliography system is now set up to work seamlessly with Bibby AI's citation management. Feel free to modify it for your needs!

\bibliographystyle{plain}
\bibliography{cite}

\end{document}\documentclass{article}
\usepackage[utf8]{inputenc}
\usepackage{booktabs}
\usepackage{amsmath}
\usepackage{natbib}

\title{My Document}
\author{Your Name}
\date{\today}

\begin{document}

\maketitle

\section{Introduction}
This is a sample document demonstrating various \LaTeX{} features including mathematical equations, tables, and formatting. You can cite papers using the citation search feature, and they will be automatically added to your \texttt{cite.bib} file.

\section{Mathematical Expressions}
Here are some example equations:

\subsection{Quadratic Formula}
The quadratic formula is given by:
\[x = \frac{-b \pm \sqrt{b^2 - 4ac}}{2a}\]

\subsection{Matrix Example}
A sample matrix:
\[
\begin{pmatrix}
a_{11} & a_{12} & a_{13} \\
a_{21} & a_{22} & a_{23} \\
a_{31} & a_{32} & a_{33}
\end{pmatrix}
\]

\section{Tables}
Below is a sample table showing data:

\begin{table}[h]
\centering
\begin{tabular}{@{}lrrr@{}}
\toprule
\textbf{Category} & \textbf{Value 1} & \textbf{Value 2} & \textbf{Total} \\
\midrule
Item A & 10.5 & 20.3 & 30.8 \\
Item B & 15.2 & 25.1 & 40.3 \\
Item C & 22.1 & 12.7 & 34.8 \\
\bottomrule
\end{tabular}
\caption{Sample Data Table}
\label{tab:sample}
\end{table}

\section{Lists}
\subsection{Bullet Points}
\begin{itemize}
  \item First item
  \item Second item
  \item Third item with sub-items:
    \begin{itemize}
      \item Sub-item 1
      \item Sub-item 2
    \end{itemize}
\end{itemize}

\subsection{Numbered List}
\begin{enumerate}
  \item First step
  \item Second step
  \item Third step
\end{enumerate}

\section{Citations and References}
When you use the citation search feature, references will be automatically managed. For example, you might cite a paper like this: \cite{example2024}. All citations will be automatically added to your \texttt{cite.bib} file.

\section{Conclusion}
This template demonstrates basic \LaTeX{} features including mathematical formulas, tables, and lists. The bibliography system is now set up to work seamlessly with Bibby AI's citation management. Feel free to modify it for your needs!

\bibliographystyle{plain}
\bibliography{cite}

\end{document}\documentclass{article}
\usepackage[utf8]{inputenc}
\usepackage{booktabs}
\usepackage{amsmath}
\usepackage{natbib}

\title{My Document}
\author{Your Name}
\date{\today}

\begin{document}

\maketitle

\section{Introduction}
This is a sample document demonstrating various \LaTeX{} features including mathematical equations, tables, and formatting. You can cite papers using the citation search feature, and they will be automatically added to your \texttt{cite.bib} file.

\section{Mathematical Expressions}
Here are some example equations:

\subsection{Quadratic Formula}
The quadratic formula is given by:
\[x = \frac{-b \pm \sqrt{b^2 - 4ac}}{2a}\]

\subsection{Matrix Example}
A sample matrix:
\[
\begin{pmatrix}
a_{11} & a_{12} & a_{13} \\
a_{21} & a_{22} & a_{23} \\
a_{31} & a_{32} & a_{33}
\end{pmatrix}
\]

\section{Tables}
Below is a sample table showing data:

\begin{table}[h]
\centering
\begin{tabular}{@{}lrrr@{}}
\toprule
\textbf{Category} & \textbf{Value 1} & \textbf{Value 2} & \textbf{Total} \\
\midrule
Item A & 10.5 & 20.3 & 30.8 \\
Item B & 15.2 & 25.1 & 40.3 \\
Item C & 22.1 & 12.7 & 34.8 \\
\bottomrule
\end{tabular}
\caption{Sample Data Table}
\label{tab:sample}
\end{table}

\section{Lists}
\subsection{Bullet Points}
\begin{itemize}
  \item First item
  \item Second item
  \item Third item with sub-items:
    \begin{itemize}
      \item Sub-item 1
      \item Sub-item 2
    \end{itemize}
\end{itemize}

\subsection{Numbered List}
\begin{enumerate}
  \item First step
  \item Second step
  \item Third step
\end{enumerate}

\section{Citations and References}
When you use the citation search feature, references will be automatically managed. For example, you might cite a paper like this: \cite{example2024}. All citations will be automatically added to your \texttt{cite.bib} file.

\section{Conclusion}
This template demonstrates basic \LaTeX{} features including mathematical formulas, tables, and lists. The bibliography system is now set up to work seamlessly with Bibby AI's citation management. Feel free to modify it for your needs!

\bibliographystyle{plain}
\bibliography{cite}

\end{document}\documentclass{article}
\usepackage[utf8]{inputenc}
\usepackage{booktabs}
\usepackage{amsmath}
\usepackage{natbib}

\title{My Document}
\author{Your Name}
\date{\today}

\begin{document}

\maketitle

\section{Introduction}
This is a sample document demonstrating various \LaTeX{} features including mathematical equations, tables, and formatting. You can cite papers using the citation search feature, and they will be automatically added to your \texttt{cite.bib} file.

\section{Mathematical Expressions}
Here are some example equations:

\subsection{Quadratic Formula}
The quadratic formula is given by:
\[x = \frac{-b \pm \sqrt{b^2 - 4ac}}{2a}\]

\subsection{Matrix Example}
A sample matrix:
\[
\begin{pmatrix}
a_{11} & a_{12} & a_{13} \\
a_{21} & a_{22} & a_{23} \\
a_{31} & a_{32} & a_{33}
\end{pmatrix}
\]

\section{Tables}
Below is a sample table showing data:

\begin{table}[h]
\centering
\begin{tabular}{@{}lrrr@{}}
\toprule
\textbf{Category} & \textbf{Value 1} & \textbf{Value 2} & \textbf{Total} \\
\midrule
Item A & 10.5 & 20.3 & 30.8 \\
Item B & 15.2 & 25.1 & 40.3 \\
Item C & 22.1 & 12.7 & 34.8 \\
\bottomrule
\end{tabular}
\caption{Sample Data Table}
\label{tab:sample}
\end{table}

\section{Lists}
\subsection{Bullet Points}
\begin{itemize}
  \item First item
  \item Second item
  \item Third item with sub-items:
    \begin{itemize}
      \item Sub-item 1
      \item Sub-item 2
    \end{itemize}
\end{itemize}

\subsection{Numbered List}
\begin{enumerate}
  \item First step
  \item Second step
  \item Third step
\end{enumerate}

\section{Citations and References}
When you use the citation search feature, references will be automatically managed. For example, you might cite a paper like this: \cite{example2024}. All citations will be automatically added to your \texttt{cite.bib} file.

\section{Conclusion}
This template demonstrates basic \LaTeX{} features including mathematical formulas, tables, and lists. The bibliography system is now set up to work seamlessly with Bibby AI's citation management. Feel free to modify it for your needs!

\bibliographystyle{plain}
\bibliography{cite}

\end{document}\documentclass{article}
\usepackage[utf8]{inputenc}
\usepackage{booktabs}
\usepackage{amsmath}
\usepackage{natbib}

\title{My Document}
\author{Your Name}
\date{\today}

\begin{document}

\maketitle

\section{Introduction}
This is a sample document demonstrating various \LaTeX{} features including mathematical equations, tables, and formatting. You can cite papers using the citation search feature, and they will be automatically added to your \texttt{cite.bib} file.

\section{Mathematical Expressions}
Here are some example equations:

\subsection{Quadratic Formula}
The quadratic formula is given by:
\[x = \frac{-b \pm \sqrt{b^2 - 4ac}}{2a}\]

\subsection{Matrix Example}
A sample matrix:
\[
\begin{pmatrix}
a_{11} & a_{12} & a_{13} \\
a_{21} & a_{22} & a_{23} \\
a_{31} & a_{32} & a_{33}
\end{pmatrix}
\]

\section{Tables}
Below is a sample table showing data:

\begin{table}[h]
\centering
\begin{tabular}{@{}lrrr@{}}
\toprule
\textbf{Category} & \textbf{Value 1} & \textbf{Value 2} & \textbf{Total} \\
\midrule
Item A & 10.5 & 20.3 & 30.8 \\
Item B & 15.2 & 25.1 & 40.3 \\
Item C & 22.1 & 12.7 & 34.8 \\
\bottomrule
\end{tabular}
\caption{Sample Data Table}
\label{tab:sample}
\end{table}

\section{Lists}
\subsection{Bullet Points}
\begin{itemize}
  \item First item
  \item Second item
  \item Third item with sub-items:
    \begin{itemize}
      \item Sub-item 1
      \item Sub-item 2
    \end{itemize}
\end{itemize}

\subsection{Numbered List}
\begin{enumerate}
  \item First step
  \item Second step
  \item Third step
\end{enumerate}

\section{Citations and References}
When you use the citation search feature, references will be automatically managed. For example, you might cite a paper like this: \cite{example2024}. All citations will be automatically added to your \texttt{cite.bib} file.

\section{Conclusion}
This template demonstrates basic \LaTeX{} features including mathematical formulas, tables, and lists. The bibliography system is now set up to work seamlessly with Bibby AI's citation management. Feel free to modify it for your needs!

\bibliographystyle{plain}
\bibliography{cite}

\end{document}\documentclass{article}
\usepackage[utf8]{inputenc}
\usepackage{booktabs}
\usepackage{amsmath}
\usepackage{natbib}

\title{My Document}
\author{Your Name}
\date{\today}

\begin{document}

\maketitle

\section{Introduction}
This is a sample document demonstrating various \LaTeX{} features including mathematical equations, tables, and formatting. You can cite papers using the citation search feature, and they will be automatically added to your \texttt{cite.bib} file.

\section{Mathematical Expressions}
Here are some example equations:

\subsection{Quadratic Formula}
The quadratic formula is given by:
\[x = \frac{-b \pm \sqrt{b^2 - 4ac}}{2a}\]

\subsection{Matrix Example}
A sample matrix:
\[
\begin{pmatrix}
a_{11} & a_{12} & a_{13} \\
a_{21} & a_{22} & a_{23} \\
a_{31} & a_{32} & a_{33}
\end{pmatrix}
\]

\section{Tables}
Below is a sample table showing data:

\begin{table}[h]
\centering
\begin{tabular}{@{}lrrr@{}}
\toprule
\textbf{Category} & \textbf{Value 1} & \textbf{Value 2} & \textbf{Total} \\
\midrule
Item A & 10.5 & 20.3 & 30.8 \\
Item B & 15.2 & 25.1 & 40.3 \\
Item C & 22.1 & 12.7 & 34.8 \\
\bottomrule
\end{tabular}
\caption{Sample Data Table}
\label{tab:sample}
\end{table}

\section{Lists}
\subsection{Bullet Points}
\begin{itemize}
  \item First item
  \item Second item
  \item Third item with sub-items:
    \begin{itemize}
      \item Sub-item 1
      \item Sub-item 2
    \end{itemize}
\end{itemize}

\subsection{Numbered List}
\begin{enumerate}
  \item First step
  \item Second step
  \item Third step
\end{enumerate}

\section{Citations and References}
When you use the citation search feature, references will be automatically managed. For example, you might cite a paper like this: \cite{example2024}. All citations will be automatically added to your \texttt{cite.bib} file.

\section{Conclusion}
This template demonstrates basic \LaTeX{} features including mathematical formulas, tables, and lists. The bibliography system is now set up to work seamlessly with Bibby AI's citation management. Feel free to modify it for your needs!

\bibliographystyle{plain}
\bibliography{cite}

\end{document}\documentclass{article}
\usepackage[utf8]{inputenc}
\usepackage{booktabs}
\usepackage{amsmath}
\usepackage{natbib}

\title{My Document}
\author{Your Name}
\date{\today}

\begin{document}

\maketitle

\section{Introduction}
This is a sample document demonstrating various \LaTeX{} features including mathematical equations, tables, and formatting. You can cite papers using the citation search feature, and they will be automatically added to your \texttt{cite.bib} file.

\section{Mathematical Expressions}
Here are some example equations:

\subsection{Quadratic Formula}
The quadratic formula is given by:
\[x = \frac{-b \pm \sqrt{b^2 - 4ac}}{2a}\]

\subsection{Matrix Example}
A sample matrix:
\[
\begin{pmatrix}
a_{11} & a_{12} & a_{13} \\
a_{21} & a_{22} & a_{23} \\
a_{31} & a_{32} & a_{33}
\end{pmatrix}
\]

\section{Tables}
Below is a sample table showing data:

\begin{table}[h]
\centering
\begin{tabular}{@{}lrrr@{}}
\toprule
\textbf{Category} & \textbf{Value 1} & \textbf{Value 2} & \textbf{Total} \\
\midrule
Item A & 10.5 & 20.3 & 30.8 \\
Item B & 15.2 & 25.1 & 40.3 \\
Item C & 22.1 & 12.7 & 34.8 \\
\bottomrule
\end{tabular}
\caption{Sample Data Table}
\label{tab:sample}
\end{table}

\section{Lists}
\subsection{Bullet Points}
\begin{itemize}
  \item First item
  \item Second item
  \item Third item with sub-items:
    \begin{itemize}
      \item Sub-item 1
      \item Sub-item 2
    \end{itemize}
\end{itemize}

\subsection{Numbered List}
\begin{enumerate}
  \item First step
  \item Second step
  \item Third step
\end{enumerate}

\section{Citations and References}
When you use the citation search feature, references will be automatically managed. For example, you might cite a paper like this: \cite{example2024}. All citations will be automatically added to your \texttt{cite.bib} file.

\section{Conclusion}
This template demonstrates basic \LaTeX{} features including mathematical formulas, tables, and lists. The bibliography system is now set up to work seamlessly with Bibby AI's citation management. Feel free to modify it for your needs!

\bibliographystyle{plain}
\bibliography{cite}

\end{document}\documentclass{article}
\usepackage[utf8]{inputenc}
\usepackage{booktabs}
\usepackage{amsmath}
\usepackage{natbib}

\title{My Document}
\author{Your Name}
\date{\today}

\begin{document}

\maketitle

\section{Introduction}
This is a sample document demonstrating various \LaTeX{} features including mathematical equations, tables, and formatting. You can cite papers using the citation search feature, and they will be automatically added to your \texttt{cite.bib} file.

\section{Mathematical Expressions}
Here are some example equations:

\subsection{Quadratic Formula}
The quadratic formula is given by:
\[x = \frac{-b \pm \sqrt{b^2 - 4ac}}{2a}\]

\subsection{Matrix Example}
A sample matrix:
\[
\begin{pmatrix}
a_{11} & a_{12} & a_{13} \\
a_{21} & a_{22} & a_{23} \\
a_{31} & a_{32} & a_{33}
\end{pmatrix}
\]

\section{Tables}
Below is a sample table showing data:

\begin{table}[h]
\centering
\begin{tabular}{@{}lrrr@{}}
\toprule
\textbf{Category} & \textbf{Value 1} & \textbf{Value 2} & \textbf{Total} \\
\midrule
Item A & 10.5 & 20.3 & 30.8 \\
Item B & 15.2 & 25.1 & 40.3 \\
Item C & 22.1 & 12.7 & 34.8 \\
\bottomrule
\end{tabular}
\caption{Sample Data Table}
\label{tab:sample}
\end{table}

\section{Lists}
\subsection{Bullet Points}
\begin{itemize}
  \item First item
  \item Second item
  \item Third item with sub-items:
    \begin{itemize}
      \item Sub-item 1
      \item Sub-item 2
    \end{itemize}
\end{itemize}

\subsection{Numbered List}
\begin{enumerate}
  \item First step
  \item Second step
  \item Third step
\end{enumerate}

\section{Citations and References}
When you use the citation search feature, references will be automatically managed. For example, you might cite a paper like this: \cite{example2024}. All citations will be automatically added to your \texttt{cite.bib} file.

\section{Conclusion}
This template demonstrates basic \LaTeX{} features including mathematical formulas, tables, and lists. The bibliography system is now set up to work seamlessly with Bibby AI's citation management. Feel free to modify it for your needs!

\bibliographystyle{plain}
\bibliography{cite}

\end{document}\documentclass{article}
\usepackage[utf8]{inputenc}
\usepackage{booktabs}
\usepackage{amsmath}
\usepackage{natbib}

\title{My Document}
\author{Your Name}
\date{\today}

\begin{document}

\maketitle

\section{Introduction}
This is a sample document demonstrating various \LaTeX{} features including mathematical equations, tables, and formatting. You can cite papers using the citation search feature, and they will be automatically added to your \texttt{cite.bib} file.

\section{Mathematical Expressions}
Here are some example equations:

\subsection{Quadratic Formula}
The quadratic formula is given by:
\[x = \frac{-b \pm \sqrt{b^2 - 4ac}}{2a}\]

\subsection{Matrix Example}
A sample matrix:
\[
\begin{pmatrix}
a_{11} & a_{12} & a_{13} \\
a_{21} & a_{22} & a_{23} \\
a_{31} & a_{32} & a_{33}
\end{pmatrix}
\]

\section{Tables}
Below is a sample table showing data:

\begin{table}[h]
\centering
\begin{tabular}{@{}lrrr@{}}
\toprule
\textbf{Category} & \textbf{Value 1} & \textbf{Value 2} & \textbf{Total} \\
\midrule
Item A & 10.5 & 20.3 & 30.8 \\
Item B & 15.2 & 25.1 & 40.3 \\
Item C & 22.1 & 12.7 & 34.8 \\
\bottomrule
\end{tabular}
\caption{Sample Data Table}
\label{tab:sample}
\end{table}

\section{Lists}
\subsection{Bullet Points}
\begin{itemize}
  \item First item
  \item Second item
  \item Third item with sub-items:
    \begin{itemize}
      \item Sub-item 1
      \item Sub-item 2
    \end{itemize}
\end{itemize}

\subsection{Numbered List}
\begin{enumerate}
  \item First step
  \item Second step
  \item Third step
\end{enumerate}

\section{Citations and References}
When you use the citation search feature, references will be automatically managed. For example, you might cite a paper like this: \cite{example2024}. All citations will be automatically added to your \texttt{cite.bib} file.

\section{Conclusion}
This template demonstrates basic \LaTeX{} features including mathematical formulas, tables, and lists. The bibliography system is now set up to work seamlessly with Bibby AI's citation management. Feel free to modify it for your needs!

\bibliographystyle{plain}
\bibliography{cite}

\end{document}\documentclass{article}
\usepackage[utf8]{inputenc}
\usepackage{booktabs}
\usepackage{amsmath}
\usepackage{natbib}

\title{My Document}
\author{Your Name}
\date{\today}

\begin{document}

\maketitle

\section{Introduction}
This is a sample document demonstrating various \LaTeX{} features including mathematical equations, tables, and formatting. You can cite papers using the citation search feature, and they will be automatically added to your \texttt{cite.bib} file.

\section{Mathematical Expressions}
Here are some example equations:

\subsection{Quadratic Formula}
The quadratic formula is given by:
\[x = \frac{-b \pm \sqrt{b^2 - 4ac}}{2a}\]

\subsection{Matrix Example}
A sample matrix:
\[
\begin{pmatrix}
a_{11} & a_{12} & a_{13} \\
a_{21} & a_{22} & a_{23} \\
a_{31} & a_{32} & a_{33}
\end{pmatrix}
\]

\section{Tables}
Below is a sample table showing data:

\begin{table}[h]
\centering
\begin{tabular}{@{}lrrr@{}}
\toprule
\textbf{Category} & \textbf{Value 1} & \textbf{Value 2} & \textbf{Total} \\
\midrule
Item A & 10.5 & 20.3 & 30.8 \\
Item B & 15.2 & 25.1 & 40.3 \\
Item C & 22.1 & 12.7 & 34.8 \\
\bottomrule
\end{tabular}
\caption{Sample Data Table}
\label{tab:sample}
\end{table}

\section{Lists}
\subsection{Bullet Points}
\begin{itemize}
  \item First item
  \item Second item
  \item Third item with sub-items:
    \begin{itemize}
      \item Sub-item 1
      \item Sub-item 2
    \end{itemize}
\end{itemize}

\subsection{Numbered List}
\begin{enumerate}
  \item First step
  \item Second step
  \item Third step
\end{enumerate}

\section{Citations and References}
When you use the citation search feature, references will be automatically managed. For example, you might cite a paper like this: \cite{example2024}. All citations will be automatically added to your \texttt{cite.bib} file.

\section{Conclusion}
This template demonstrates basic \LaTeX{} features including mathematical formulas, tables, and lists. The bibliography system is now set up to work seamlessly with Bibby AI's citation management. Feel free to modify it for your needs!

\bibliographystyle{plain}
\bibliography{cite}

\end{document}\documentclass{article}
\usepackage[utf8]{inputenc}
\usepackage{booktabs}
\usepackage{amsmath}
\usepackage{natbib}

\title{My Document}
\author{Your Name}
\date{\today}

\begin{document}

\maketitle

\section{Introduction}
This is a sample document demonstrating various \LaTeX{} features including mathematical equations, tables, and formatting. You can cite papers using the citation search feature, and they will be automatically added to your \texttt{cite.bib} file.

\section{Mathematical Expressions}
Here are some example equations:

\subsection{Quadratic Formula}
The quadratic formula is given by:
\[x = \frac{-b \pm \sqrt{b^2 - 4ac}}{2a}\]

\subsection{Matrix Example}
A sample matrix:
\[
\begin{pmatrix}
a_{11} & a_{12} & a_{13} \\
a_{21} & a_{22} & a_{23} \\
a_{31} & a_{32} & a_{33}
\end{pmatrix}
\]

\section{Tables}
Below is a sample table showing data:

\begin{table}[h]
\centering
\begin{tabular}{@{}lrrr@{}}
\toprule
\textbf{Category} & \textbf{Value 1} & \textbf{Value 2} & \textbf{Total} \\
\midrule
Item A & 10.5 & 20.3 & 30.8 \\
Item B & 15.2 & 25.1 & 40.3 \\
Item C & 22.1 & 12.7 & 34.8 \\
\bottomrule
\end{tabular}
\caption{Sample Data Table}
\label{tab:sample}
\end{table}

\section{Lists}
\subsection{Bullet Points}
\begin{itemize}
  \item First item
  \item Second item
  \item Third item with sub-items:
    \begin{itemize}
      \item Sub-item 1
      \item Sub-item 2
    \end{itemize}
\end{itemize}

\subsection{Numbered List}
\begin{enumerate}
  \item First step
  \item Second step
  \item Third step
\end{enumerate}

\section{Citations and References}
When you use the citation search feature, references will be automatically managed. For example, you might cite a paper like this: \cite{example2024}. All citations will be automatically added to your \texttt{cite.bib} file.

\section{Conclusion}
This template demonstrates basic \LaTeX{} features including mathematical formulas, tables, and lists. The bibliography system is now set up to work seamlessly with Bibby AI's citation management. Feel free to modify it for your needs!

\bibliographystyle{plain}
\bibliography{cite}

\end{document}\documentclass{article}
\usepackage[utf8]{inputenc}
\usepackage{booktabs}
\usepackage{amsmath}
\usepackage{natbib}

\title{My Document}
\author{Your Name}
\date{\today}

\begin{document}

\maketitle

\section{Introduction}
This is a sample document demonstrating various \LaTeX{} features including mathematical equations, tables, and formatting. You can cite papers using the citation search feature, and they will be automatically added to your \texttt{cite.bib} file.

\section{Mathematical Expressions}
Here are some example equations:

\subsection{Quadratic Formula}
The quadratic formula is given by:
\[x = \frac{-b \pm \sqrt{b^2 - 4ac}}{2a}\]

\subsection{Matrix Example}
A sample matrix:
\[
\begin{pmatrix}
a_{11} & a_{12} & a_{13} \\
a_{21} & a_{22} & a_{23} \\
a_{31} & a_{32} & a_{33}
\end{pmatrix}
\]

\section{Tables}
Below is a sample table showing data:

\begin{table}[h]
\centering
\begin{tabular}{@{}lrrr@{}}
\toprule
\textbf{Category} & \textbf{Value 1} & \textbf{Value 2} & \textbf{Total} \\
\midrule
Item A & 10.5 & 20.3 & 30.8 \\
Item B & 15.2 & 25.1 & 40.3 \\
Item C & 22.1 & 12.7 & 34.8 \\
\bottomrule
\end{tabular}
\caption{Sample Data Table}
\label{tab:sample}
\end{table}

\section{Lists}
\subsection{Bullet Points}
\begin{itemize}
  \item First item
  \item Second item
  \item Third item with sub-items:
    \begin{itemize}
      \item Sub-item 1
      \item Sub-item 2
    \end{itemize}
\end{itemize}

\subsection{Numbered List}
\begin{enumerate}
  \item First step
  \item Second step
  \item Third step
\end{enumerate}

\section{Citations and References}
When you use the citation search feature, references will be automatically managed. For example, you might cite a paper like this: \cite{example2024}. All citations will be automatically added to your \texttt{cite.bib} file.

\section{Conclusion}
This template demonstrates basic \LaTeX{} features including mathematical formulas, tables, and lists. The bibliography system is now set up to work seamlessly with Bibby AI's citation management. Feel free to modify it for your needs!

\bibliographystyle{plain}
\bibliography{cite}

\end{document}\documentclass{article}
\usepackage[utf8]{inputenc}
\usepackage{booktabs}
\usepackage{amsmath}
\usepackage{natbib}

\title{My Document}
\author{Your Name}
\date{\today}

\begin{document}

\maketitle

\section{Introduction}
This is a sample document demonstrating various \LaTeX{} features including mathematical equations, tables, and formatting. You can cite papers using the citation search feature, and they will be automatically added to your \texttt{cite.bib} file.

\section{Mathematical Expressions}
Here are some example equations:

\subsection{Quadratic Formula}
The quadratic formula is given by:
\[x = \frac{-b \pm \sqrt{b^2 - 4ac}}{2a}\]

\subsection{Matrix Example}
A sample matrix:
\[
\begin{pmatrix}
a_{11} & a_{12} & a_{13} \\
a_{21} & a_{22} & a_{23} \\
a_{31} & a_{32} & a_{33}
\end{pmatrix}
\]

\section{Tables}
Below is a sample table showing data:

\begin{table}[h]
\centering
\begin{tabular}{@{}lrrr@{}}
\toprule
\textbf{Category} & \textbf{Value 1} & \textbf{Value 2} & \textbf{Total} \\
\midrule
Item A & 10.5 & 20.3 & 30.8 \\
Item B & 15.2 & 25.1 & 40.3 \\
Item C & 22.1 & 12.7 & 34.8 \\
\bottomrule
\end{tabular}
\caption{Sample Data Table}
\label{tab:sample}
\end{table}

\section{Lists}
\subsection{Bullet Points}
\begin{itemize}
  \item First item
  \item Second item
  \item Third item with sub-items:
    \begin{itemize}
      \item Sub-item 1
      \item Sub-item 2
    \end{itemize}
\end{itemize}

\subsection{Numbered List}
\begin{enumerate}
  \item First step
  \item Second step
  \item Third step
\end{enumerate}

\section{Citations and References}
When you use the citation search feature, references will be automatically managed. For example, you might cite a paper like this: \cite{example2024}. All citations will be automatically added to your \texttt{cite.bib} file.

\section{Conclusion}
This template demonstrates basic \LaTeX{} features including mathematical formulas, tables, and lists. The bibliography system is now set up to work seamlessly with Bibby AI's citation management. Feel free to modify it for your needs!

\bibliographystyle{plain}
\bibliography{cite}

\end{document}\documentclass{article}
\usepackage[utf8]{inputenc}
\usepackage{booktabs}
\usepackage{amsmath}
\usepackage{natbib}

\title{My Document}
\author{Your Name}
\date{\today}

\begin{document}

\maketitle

\section{Introduction}
This is a sample document demonstrating various \LaTeX{} features including mathematical equations, tables, and formatting. You can cite papers using the citation search feature, and they will be automatically added to your \texttt{cite.bib} file.

\section{Mathematical Expressions}
Here are some example equations:

\subsection{Quadratic Formula}
The quadratic formula is given by:
\[x = \frac{-b \pm \sqrt{b^2 - 4ac}}{2a}\]

\subsection{Matrix Example}
A sample matrix:
\[
\begin{pmatrix}
a_{11} & a_{12} & a_{13} \\
a_{21} & a_{22} & a_{23} \\
a_{31} & a_{32} & a_{33}
\end{pmatrix}
\]

\section{Tables}
Below is a sample table showing data:

\begin{table}[h]
\centering
\begin{tabular}{@{}lrrr@{}}
\toprule
\textbf{Category} & \textbf{Value 1} & \textbf{Value 2} & \textbf{Total} \\
\midrule
Item A & 10.5 & 20.3 & 30.8 \\
Item B & 15.2 & 25.1 & 40.3 \\
Item C & 22.1 & 12.7 & 34.8 \\
\bottomrule
\end{tabular}
\caption{Sample Data Table}
\label{tab:sample}
\end{table}

\section{Lists}
\subsection{Bullet Points}
\begin{itemize}
  \item First item
  \item Second item
  \item Third item with sub-items:
    \begin{itemize}
      \item Sub-item 1
      \item Sub-item 2
    \end{itemize}
\end{itemize}

\subsection{Numbered List}
\begin{enumerate}
  \item First step
  \item Second step
  \item Third step
\end{enumerate}

\section{Citations and References}
When you use the citation search feature, references will be automatically managed. For example, you might cite a paper like this: \cite{example2024}. All citations will be automatically added to your \texttt{cite.bib} file.

\section{Conclusion}
This template demonstrates basic \LaTeX{} features including mathematical formulas, tables, and lists. The bibliography system is now set up to work seamlessly with Bibby AI's citation management. Feel free to modify it for your needs!

\bibliographystyle{plain}
\bibliography{cite}

\end{document}\documentclass{article}
\usepackage[utf8]{inputenc}
\usepackage{booktabs}
\usepackage{amsmath}
\usepackage{natbib}

\title{My Document}
\author{Your Name}
\date{\today}

\begin{document}

\maketitle

\section{Introduction}
This is a sample document demonstrating various \LaTeX{} features including mathematical equations, tables, and formatting. You can cite papers using the citation search feature, and they will be automatically added to your \texttt{cite.bib} file.

\section{Mathematical Expressions}
Here are some example equations:

\subsection{Quadratic Formula}
The quadratic formula is given by:
\[x = \frac{-b \pm \sqrt{b^2 - 4ac}}{2a}\]

\subsection{Matrix Example}
A sample matrix:
\[
\begin{pmatrix}
a_{11} & a_{12} & a_{13} \\
a_{21} & a_{22} & a_{23} \\
a_{31} & a_{32} & a_{33}
\end{pmatrix}
\]

\section{Tables}
Below is a sample table showing data:

\begin{table}[h]
\centering
\begin{tabular}{@{}lrrr@{}}
\toprule
\textbf{Category} & \textbf{Value 1} & \textbf{Value 2} & \textbf{Total} \\
\midrule
Item A & 10.5 & 20.3 & 30.8 \\
Item B & 15.2 & 25.1 & 40.3 \\
Item C & 22.1 & 12.7 & 34.8 \\
\bottomrule
\end{tabular}
\caption{Sample Data Table}
\label{tab:sample}
\end{table}

\section{Lists}
\subsection{Bullet Points}
\begin{itemize}
  \item First item
  \item Second item
  \item Third item with sub-items:
    \begin{itemize}
      \item Sub-item 1
      \item Sub-item 2
    \end{itemize}
\end{itemize}

\subsection{Numbered List}
\begin{enumerate}
  \item First step
  \item Second step
  \item Third step
\end{enumerate}

\section{Citations and References}
When you use the citation search feature, references will be automatically managed. For example, you might cite a paper like this: \cite{example2024}. All citations will be automatically added to your \texttt{cite.bib} file.

\section{Conclusion}
This template demonstrates basic \LaTeX{} features including mathematical formulas, tables, and lists. The bibliography system is now set up to work seamlessly with Bibby AI's citation management. Feel free to modify it for your needs!

\bibliographystyle{plain}
\bibliography{cite}

\end{document}\documentclass{article}
\usepackage[utf8]{inputenc}
\usepackage{booktabs}
\usepackage{amsmath}
\usepackage{natbib}

\title{My Document}
\author{Your Name}
\date{\today}

\begin{document}

\maketitle

\section{Introduction}
This is a sample document demonstrating various \LaTeX{} features including mathematical equations, tables, and formatting. You can cite papers using the citation search feature, and they will be automatically added to your \texttt{cite.bib} file.

\section{Mathematical Expressions}
Here are some example equations:

\subsection{Quadratic Formula}
The quadratic formula is given by:
\[x = \frac{-b \pm \sqrt{b^2 - 4ac}}{2a}\]

\subsection{Matrix Example}
A sample matrix:
\[
\begin{pmatrix}
a_{11} & a_{12} & a_{13} \\
a_{21} & a_{22} & a_{23} \\
a_{31} & a_{32} & a_{33}
\end{pmatrix}
\]

\section{Tables}
Below is a sample table showing data:

\begin{table}[h]
\centering
\begin{tabular}{@{}lrrr@{}}
\toprule
\textbf{Category} & \textbf{Value 1} & \textbf{Value 2} & \textbf{Total} \\
\midrule
Item A & 10.5 & 20.3 & 30.8 \\
Item B & 15.2 & 25.1 & 40.3 \\
Item C & 22.1 & 12.7 & 34.8 \\
\bottomrule
\end{tabular}
\caption{Sample Data Table}
\label{tab:sample}
\end{table}

\section{Lists}
\subsection{Bullet Points}
\begin{itemize}
  \item First item
  \item Second item
  \item Third item with sub-items:
    \begin{itemize}
      \item Sub-item 1
      \item Sub-item 2
    \end{itemize}
\end{itemize}

\subsection{Numbered List}
\begin{enumerate}
  \item First step
  \item Second step
  \item Third step
\end{enumerate}

\section{Citations and References}
When you use the citation search feature, references will be automatically managed. For example, you might cite a paper like this: \cite{example2024}. All citations will be automatically added to your \texttt{cite.bib} file.

\section{Conclusion}
This template demonstrates basic \LaTeX{} features including mathematical formulas, tables, and lists. The bibliography system is now set up to work seamlessly with Bibby AI's citation management. Feel free to modify it for your needs!

\bibliographystyle{plain}
\bibliography{cite}

\end{document}\documentclass{article}
\usepackage[utf8]{inputenc}
\usepackage{booktabs}
\usepackage{amsmath}
\usepackage{natbib}

\title{My Document}
\author{Your Name}
\date{\today}

\begin{document}

\maketitle

\section{Introduction}
This is a sample document demonstrating various \LaTeX{} features including mathematical equations, tables, and formatting. You can cite papers using the citation search feature, and they will be automatically added to your \texttt{cite.bib} file.

\section{Mathematical Expressions}
Here are some example equations:

\subsection{Quadratic Formula}
The quadratic formula is given by:
\[x = \frac{-b \pm \sqrt{b^2 - 4ac}}{2a}\]

\subsection{Matrix Example}
A sample matrix:
\[
\begin{pmatrix}
a_{11} & a_{12} & a_{13} \\
a_{21} & a_{22} & a_{23} \\
a_{31} & a_{32} & a_{33}
\end{pmatrix}
\]

\section{Tables}
Below is a sample table showing data:

\begin{table}[h]
\centering
\begin{tabular}{@{}lrrr@{}}
\toprule
\textbf{Category} & \textbf{Value 1} & \textbf{Value 2} & \textbf{Total} \\
\midrule
Item A & 10.5 & 20.3 & 30.8 \\
Item B & 15.2 & 25.1 & 40.3 \\
Item C & 22.1 & 12.7 & 34.8 \\
\bottomrule
\end{tabular}
\caption{Sample Data Table}
\label{tab:sample}
\end{table}

\section{Lists}
\subsection{Bullet Points}
\begin{itemize}
  \item First item
  \item Second item
  \item Third item with sub-items:
    \begin{itemize}
      \item Sub-item 1
      \item Sub-item 2
    \end{itemize}
\end{itemize}

\subsection{Numbered List}
\begin{enumerate}
  \item First step
  \item Second step
  \item Third step
\end{enumerate}

\section{Citations and References}
When you use the citation search feature, references will be automatically managed. For example, you might cite a paper like this: \cite{example2024}. All citations will be automatically added to your \texttt{cite.bib} file.

\section{Conclusion}
This template demonstrates basic \LaTeX{} features including mathematical formulas, tables, and lists. The bibliography system is now set up to work seamlessly with Bibby AI's citation management. Feel free to modify it for your needs!

\bibliographystyle{plain}
\bibliography{cite}

\end{document}\documentclass{article}
\usepackage[utf8]{inputenc}
\usepackage{booktabs}
\usepackage{amsmath}
\usepackage{natbib}

\title{My Document}
\author{Your Name}
\date{\today}

\begin{document}

\maketitle

\section{Introduction}
This is a sample document demonstrating various \LaTeX{} features including mathematical equations, tables, and formatting. You can cite papers using the citation search feature, and they will be automatically added to your \texttt{cite.bib} file.

\section{Mathematical Expressions}
Here are some example equations:

\subsection{Quadratic Formula}
The quadratic formula is given by:
\[x = \frac{-b \pm \sqrt{b^2 - 4ac}}{2a}\]

\subsection{Matrix Example}
A sample matrix:
\[
\begin{pmatrix}
a_{11} & a_{12} & a_{13} \\
a_{21} & a_{22} & a_{23} \\
a_{31} & a_{32} & a_{33}
\end{pmatrix}
\]

\section{Tables}
Below is a sample table showing data:

\begin{table}[h]
\centering
\begin{tabular}{@{}lrrr@{}}
\toprule
\textbf{Category} & \textbf{Value 1} & \textbf{Value 2} & \textbf{Total} \\
\midrule
Item A & 10.5 & 20.3 & 30.8 \\
Item B & 15.2 & 25.1 & 40.3 \\
Item C & 22.1 & 12.7 & 34.8 \\
\bottomrule
\end{tabular}
\caption{Sample Data Table}
\label{tab:sample}
\end{table}

\section{Lists}
\subsection{Bullet Points}
\begin{itemize}
  \item First item
  \item Second item
  \item Third item with sub-items:
    \begin{itemize}
      \item Sub-item 1
      \item Sub-item 2
    \end{itemize}
\end{itemize}

\subsection{Numbered List}
\begin{enumerate}
  \item First step
  \item Second step
  \item Third step
\end{enumerate}

\section{Citations and References}
When you use the citation search feature, references will be automatically managed. For example, you might cite a paper like this: \cite{example2024}. All citations will be automatically added to your \texttt{cite.bib} file.

\section{Conclusion}
This template demonstrates basic \LaTeX{} features including mathematical formulas, tables, and lists. The bibliography system is now set up to work seamlessly with Bibby AI's citation management. Feel free to modify it for your needs!

\bibliographystyle{plain}
\bibliography{cite}

\end{document}\documentclass{article}
\usepackage[utf8]{inputenc}
\usepackage{booktabs}
\usepackage{amsmath}
\usepackage{natbib}

\title{My Document}
\author{Your Name}
\date{\today}

\begin{document}

\maketitle

\section{Introduction}
This is a sample document demonstrating various \LaTeX{} features including mathematical equations, tables, and formatting. You can cite papers using the citation search feature, and they will be automatically added to your \texttt{cite.bib} file.

\section{Mathematical Expressions}
Here are some example equations:

\subsection{Quadratic Formula}
The quadratic formula is given by:
\[x = \frac{-b \pm \sqrt{b^2 - 4ac}}{2a}\]

\subsection{Matrix Example}
A sample matrix:
\[
\begin{pmatrix}
a_{11} & a_{12} & a_{13} \\
a_{21} & a_{22} & a_{23} \\
a_{31} & a_{32} & a_{33}
\end{pmatrix}
\]

\section{Tables}
Below is a sample table showing data:

\begin{table}[h]
\centering
\begin{tabular}{@{}lrrr@{}}
\toprule
\textbf{Category} & \textbf{Value 1} & \textbf{Value 2} & \textbf{Total} \\
\midrule
Item A & 10.5 & 20.3 & 30.8 \\
Item B & 15.2 & 25.1 & 40.3 \\
Item C & 22.1 & 12.7 & 34.8 \\
\bottomrule
\end{tabular}
\caption{Sample Data Table}
\label{tab:sample}
\end{table}

\section{Lists}
\subsection{Bullet Points}
\begin{itemize}
  \item First item
  \item Second item
  \item Third item with sub-items:
    \begin{itemize}
      \item Sub-item 1
      \item Sub-item 2
    \end{itemize}
\end{itemize}

\subsection{Numbered List}
\begin{enumerate}
  \item First step
  \item Second step
  \item Third step
\end{enumerate}

\section{Citations and References}
When you use the citation search feature, references will be automatically managed. For example, you might cite a paper like this: \cite{example2024}. All citations will be automatically added to your \texttt{cite.bib} file.

\section{Conclusion}
This template demonstrates basic \LaTeX{} features including mathematical formulas, tables, and lists. The bibliography system is now set up to work seamlessly with Bibby AI's citation management. Feel free to modify it for your needs!

\bibliographystyle{plain}
\bibliography{cite}

\end{document}\documentclass{article}
\usepackage[utf8]{inputenc}
\usepackage{booktabs}
\usepackage{amsmath}
\usepackage{natbib}

\title{My Document}
\author{Your Name}
\date{\today}

\begin{document}

\maketitle

\section{Introduction}
This is a sample document demonstrating various \LaTeX{} features including mathematical equations, tables, and formatting. You can cite papers using the citation search feature, and they will be automatically added to your \texttt{cite.bib} file.

\section{Mathematical Expressions}
Here are some example equations:

\subsection{Quadratic Formula}
The quadratic formula is given by:
\[x = \frac{-b \pm \sqrt{b^2 - 4ac}}{2a}\]

\subsection{Matrix Example}
A sample matrix:
\[
\begin{pmatrix}
a_{11} & a_{12} & a_{13} \\
a_{21} & a_{22} & a_{23} \\
a_{31} & a_{32} & a_{33}
\end{pmatrix}
\]

\section{Tables}
Below is a sample table showing data:

\begin{table}[h]
\centering
\begin{tabular}{@{}lrrr@{}}
\toprule
\textbf{Category} & \textbf{Value 1} & \textbf{Value 2} & \textbf{Total} \\
\midrule
Item A & 10.5 & 20.3 & 30.8 \\
Item B & 15.2 & 25.1 & 40.3 \\
Item C & 22.1 & 12.7 & 34.8 \\
\bottomrule
\end{tabular}
\caption{Sample Data Table}
\label{tab:sample}
\end{table}

\section{Lists}
\subsection{Bullet Points}
\begin{itemize}
  \item First item
  \item Second item
  \item Third item with sub-items:
    \begin{itemize}
      \item Sub-item 1
      \item Sub-item 2
    \end{itemize}
\end{itemize}

\subsection{Numbered List}
\begin{enumerate}
  \item First step
  \item Second step
  \item Third step
\end{enumerate}

\section{Citations and References}
When you use the citation search feature, references will be automatically managed. For example, you might cite a paper like this: \cite{example2024}. All citations will be automatically added to your \texttt{cite.bib} file.

\section{Conclusion}
This template demonstrates basic \LaTeX{} features including mathematical formulas, tables, and lists. The bibliography system is now set up to work seamlessly with Bibby AI's citation management. Feel free to modify it for your needs!

\bibliographystyle{plain}
\bibliography{cite}

\end{document}\documentclass{article}
\usepackage[utf8]{inputenc}
\usepackage{booktabs}
\usepackage{amsmath}
\usepackage{natbib}

\title{My Document}
\author{Your Name}
\date{\today}

\begin{document}

\maketitle

\section{Introduction}
This is a sample document demonstrating various \LaTeX{} features including mathematical equations, tables, and formatting. You can cite papers using the citation search feature, and they will be automatically added to your \texttt{cite.bib} file.

\section{Mathematical Expressions}
Here are some example equations:

\subsection{Quadratic Formula}
The quadratic formula is given by:
\[x = \frac{-b \pm \sqrt{b^2 - 4ac}}{2a}\]

\subsection{Matrix Example}
A sample matrix:
\[
\begin{pmatrix}
a_{11} & a_{12} & a_{13} \\
a_{21} & a_{22} & a_{23} \\
a_{31} & a_{32} & a_{33}
\end{pmatrix}
\]

\section{Tables}
Below is a sample table showing data:

\begin{table}[h]
\centering
\begin{tabular}{@{}lrrr@{}}
\toprule
\textbf{Category} & \textbf{Value 1} & \textbf{Value 2} & \textbf{Total} \\
\midrule
Item A & 10.5 & 20.3 & 30.8 \\
Item B & 15.2 & 25.1 & 40.3 \\
Item C & 22.1 & 12.7 & 34.8 \\
\bottomrule
\end{tabular}
\caption{Sample Data Table}
\label{tab:sample}
\end{table}

\section{Lists}
\subsection{Bullet Points}
\begin{itemize}
  \item First item
  \item Second item
  \item Third item with sub-items:
    \begin{itemize}
      \item Sub-item 1
      \item Sub-item 2
    \end{itemize}
\end{itemize}

\subsection{Numbered List}
\begin{enumerate}
  \item First step
  \item Second step
  \item Third step
\end{enumerate}

\section{Citations and References}
When you use the citation search feature, references will be automatically managed. For example, you might cite a paper like this: \cite{example2024}. All citations will be automatically added to your \texttt{cite.bib} file.

\section{Conclusion}
This template demonstrates basic \LaTeX{} features including mathematical formulas, tables, and lists. The bibliography system is now set up to work seamlessly with Bibby AI's citation management. Feel free to modify it for your needs!

\bibliographystyle{plain}
\bibliography{cite}

\end{document}\documentclass{article}
\usepackage[utf8]{inputenc}
\usepackage{booktabs}
\usepackage{amsmath}
\usepackage{natbib}

\title{My Document}
\author{Your Name}
\date{\today}

\begin{document}

\maketitle

\section{Introduction}
This is a sample document demonstrating various \LaTeX{} features including mathematical equations, tables, and formatting. You can cite papers using the citation search feature, and they will be automatically added to your \texttt{cite.bib} file.

\section{Mathematical Expressions}
Here are some example equations:

\subsection{Quadratic Formula}
The quadratic formula is given by:
\[x = \frac{-b \pm \sqrt{b^2 - 4ac}}{2a}\]

\subsection{Matrix Example}
A sample matrix:
\[
\begin{pmatrix}
a_{11} & a_{12} & a_{13} \\
a_{21} & a_{22} & a_{23} \\
a_{31} & a_{32} & a_{33}
\end{pmatrix}
\]

\section{Tables}
Below is a sample table showing data:

\begin{table}[h]
\centering
\begin{tabular}{@{}lrrr@{}}
\toprule
\textbf{Category} & \textbf{Value 1} & \textbf{Value 2} & \textbf{Total} \\
\midrule
Item A & 10.5 & 20.3 & 30.8 \\
Item B & 15.2 & 25.1 & 40.3 \\
Item C & 22.1 & 12.7 & 34.8 \\
\bottomrule
\end{tabular}
\caption{Sample Data Table}
\label{tab:sample}
\end{table}

\section{Lists}
\subsection{Bullet Points}
\begin{itemize}
  \item First item
  \item Second item
  \item Third item with sub-items:
    \begin{itemize}
      \item Sub-item 1
      \item Sub-item 2
    \end{itemize}
\end{itemize}

\subsection{Numbered List}
\begin{enumerate}
  \item First step
  \item Second step
  \item Third step
\end{enumerate}

\section{Citations and References}
When you use the citation search feature, references will be automatically managed. For example, you might cite a paper like this: \cite{example2024}. All citations will be automatically added to your \texttt{cite.bib} file.

\section{Conclusion}
This template demonstrates basic \LaTeX{} features including mathematical formulas, tables, and lists. The bibliography system is now set up to work seamlessly with Bibby AI's citation management. Feel free to modify it for your needs!

\bibliographystyle{plain}
\bibliography{cite}

\end{document}\documentclass{article}
\usepackage[utf8]{inputenc}
\usepackage{booktabs}
\usepackage{amsmath}
\usepackage{natbib}

\title{My Document}
\author{Your Name}
\date{\today}

\begin{document}

\maketitle

\section{Introduction}
This is a sample document demonstrating various \LaTeX{} features including mathematical equations, tables, and formatting. You can cite papers using the citation search feature, and they will be automatically added to your \texttt{cite.bib} file.

\section{Mathematical Expressions}
Here are some example equations:

\subsection{Quadratic Formula}
The quadratic formula is given by:
\[x = \frac{-b \pm \sqrt{b^2 - 4ac}}{2a}\]

\subsection{Matrix Example}
A sample matrix:
\[
\begin{pmatrix}
a_{11} & a_{12} & a_{13} \\
a_{21} & a_{22} & a_{23} \\
a_{31} & a_{32} & a_{33}
\end{pmatrix}
\]

\section{Tables}
Below is a sample table showing data:

\begin{table}[h]
\centering
\begin{tabular}{@{}lrrr@{}}
\toprule
\textbf{Category} & \textbf{Value 1} & \textbf{Value 2} & \textbf{Total} \\
\midrule
Item A & 10.5 & 20.3 & 30.8 \\
Item B & 15.2 & 25.1 & 40.3 \\
Item C & 22.1 & 12.7 & 34.8 \\
\bottomrule
\end{tabular}
\caption{Sample Data Table}
\label{tab:sample}
\end{table}

\section{Lists}
\subsection{Bullet Points}
\begin{itemize}
  \item First item
  \item Second item
  \item Third item with sub-items:
    \begin{itemize}
      \item Sub-item 1
      \item Sub-item 2
    \end{itemize}
\end{itemize}

\subsection{Numbered List}
\begin{enumerate}
  \item First step
  \item Second step
  \item Third step
\end{enumerate}

\section{Citations and References}
When you use the citation search feature, references will be automatically managed. For example, you might cite a paper like this: \cite{example2024}. All citations will be automatically added to your \texttt{cite.bib} file.

\section{Conclusion}
This template demonstrates basic \LaTeX{} features including mathematical formulas, tables, and lists. The bibliography system is now set up to work seamlessly with Bibby AI's citation management. Feel free to modify it for your needs!

\bibliographystyle{plain}
\bibliography{cite}

\end{document}\documentclass{article}
\usepackage[utf8]{inputenc}
\usepackage{booktabs}
\usepackage{amsmath}
\usepackage{natbib}

\title{My Document}
\author{Your Name}
\date{\today}

\begin{document}

\maketitle

\section{Introduction}
This is a sample document demonstrating various \LaTeX{} features including mathematical equations, tables, and formatting. You can cite papers using the citation search feature, and they will be automatically added to your \texttt{cite.bib} file.

\section{Mathematical Expressions}
Here are some example equations:

\subsection{Quadratic Formula}
The quadratic formula is given by:
\[x = \frac{-b \pm \sqrt{b^2 - 4ac}}{2a}\]

\subsection{Matrix Example}
A sample matrix:
\[
\begin{pmatrix}
a_{11} & a_{12} & a_{13} \\
a_{21} & a_{22} & a_{23} \\
a_{31} & a_{32} & a_{33}
\end{pmatrix}
\]

\section{Tables}
Below is a sample table showing data:

\begin{table}[h]
\centering
\begin{tabular}{@{}lrrr@{}}
\toprule
\textbf{Category} & \textbf{Value 1} & \textbf{Value 2} & \textbf{Total} \\
\midrule
Item A & 10.5 & 20.3 & 30.8 \\
Item B & 15.2 & 25.1 & 40.3 \\
Item C & 22.1 & 12.7 & 34.8 \\
\bottomrule
\end{tabular}
\caption{Sample Data Table}
\label{tab:sample}
\end{table}

\section{Lists}
\subsection{Bullet Points}
\begin{itemize}
  \item First item
  \item Second item
  \item Third item with sub-items:
    \begin{itemize}
      \item Sub-item 1
      \item Sub-item 2
    \end{itemize}
\end{itemize}

\subsection{Numbered List}
\begin{enumerate}
  \item First step
  \item Second step
  \item Third step
\end{enumerate}

\section{Citations and References}
When you use the citation search feature, references will be automatically managed. For example, you might cite a paper like this: \cite{example2024}. All citations will be automatically added to your \texttt{cite.bib} file.

\section{Conclusion}
This template demonstrates basic \LaTeX{} features including mathematical formulas, tables, and lists. The bibliography system is now set up to work seamlessly with Bibby AI's citation management. Feel free to modify it for your needs!

\bibliographystyle{plain}
\bibliography{cite}

\end{document}\documentclass{article}
\usepackage[utf8]{inputenc}
\usepackage{booktabs}
\usepackage{amsmath}
\usepackage{natbib}

\title{My Document}
\author{Your Name}
\date{\today}

\begin{document}

\maketitle

\section{Introduction}
This is a sample document demonstrating various \LaTeX{} features including mathematical equations, tables, and formatting. You can cite papers using the citation search feature, and they will be automatically added to your \texttt{cite.bib} file.

\section{Mathematical Expressions}
Here are some example equations:

\subsection{Quadratic Formula}
The quadratic formula is given by:
\[x = \frac{-b \pm \sqrt{b^2 - 4ac}}{2a}\]

\subsection{Matrix Example}
A sample matrix:
\[
\begin{pmatrix}
a_{11} & a_{12} & a_{13} \\
a_{21} & a_{22} & a_{23} \\
a_{31} & a_{32} & a_{33}
\end{pmatrix}
\]

\section{Tables}
Below is a sample table showing data:

\begin{table}[h]
\centering
\begin{tabular}{@{}lrrr@{}}
\toprule
\textbf{Category} & \textbf{Value 1} & \textbf{Value 2} & \textbf{Total} \\
\midrule
Item A & 10.5 & 20.3 & 30.8 \\
Item B & 15.2 & 25.1 & 40.3 \\
Item C & 22.1 & 12.7 & 34.8 \\
\bottomrule
\end{tabular}
\caption{Sample Data Table}
\label{tab:sample}
\end{table}

\section{Lists}
\subsection{Bullet Points}
\begin{itemize}
  \item First item
  \item Second item
  \item Third item with sub-items:
    \begin{itemize}
      \item Sub-item 1
      \item Sub-item 2
    \end{itemize}
\end{itemize}

\subsection{Numbered List}
\begin{enumerate}
  \item First step
  \item Second step
  \item Third step
\end{enumerate}

\section{Citations and References}
When you use the citation search feature, references will be automatically managed. For example, you might cite a paper like this: \cite{example2024}. All citations will be automatically added to your \texttt{cite.bib} file.

\section{Conclusion}
This template demonstrates basic \LaTeX{} features including mathematical formulas, tables, and lists. The bibliography system is now set up to work seamlessly with Bibby AI's citation management. Feel free to modify it for your needs!

\bibliographystyle{plain}
\bibliography{cite}

\end{document}\documentclass{article}
\usepackage[utf8]{inputenc}
\usepackage{booktabs}
\usepackage{amsmath}
\usepackage{natbib}

\title{My Document}
\author{Your Name}
\date{\today}

\begin{document}

\maketitle

\section{Introduction}
This is a sample document demonstrating various \LaTeX{} features including mathematical equations, tables, and formatting. You can cite papers using the citation search feature, and they will be automatically added to your \texttt{cite.bib} file.

\section{Mathematical Expressions}
Here are some example equations:

\subsection{Quadratic Formula}
The quadratic formula is given by:
\[x = \frac{-b \pm \sqrt{b^2 - 4ac}}{2a}\]

\subsection{Matrix Example}
A sample matrix:
\[
\begin{pmatrix}
a_{11} & a_{12} & a_{13} \\
a_{21} & a_{22} & a_{23} \\
a_{31} & a_{32} & a_{33}
\end{pmatrix}
\]

\section{Tables}
Below is a sample table showing data:

\begin{table}[h]
\centering
\begin{tabular}{@{}lrrr@{}}
\toprule
\textbf{Category} & \textbf{Value 1} & \textbf{Value 2} & \textbf{Total} \\
\midrule
Item A & 10.5 & 20.3 & 30.8 \\
Item B & 15.2 & 25.1 & 40.3 \\
Item C & 22.1 & 12.7 & 34.8 \\
\bottomrule
\end{tabular}
\caption{Sample Data Table}
\label{tab:sample}
\end{table}

\section{Lists}
\subsection{Bullet Points}
\begin{itemize}
  \item First item
  \item Second item
  \item Third item with sub-items:
    \begin{itemize}
      \item Sub-item 1
      \item Sub-item 2
    \end{itemize}
\end{itemize}

\subsection{Numbered List}
\begin{enumerate}
  \item First step
  \item Second step
  \item Third step
\end{enumerate}

\section{Citations and References}
When you use the citation search feature, references will be automatically managed. For example, you might cite a paper like this: \cite{example2024}. All citations will be automatically added to your \texttt{cite.bib} file.

\section{Conclusion}
This template demonstrates basic \LaTeX{} features including mathematical formulas, tables, and lists. The bibliography system is now set up to work seamlessly with Bibby AI's citation management. Feel free to modify it for your needs!

\bibliographystyle{plain}
\bibliography{cite}

\end{document}\documentclass{article}
\usepackage[utf8]{inputenc}
\usepackage{booktabs}
\usepackage{amsmath}
\usepackage{natbib}

\title{My Document}
\author{Your Name}
\date{\today}

\begin{document}

\maketitle

\section{Introduction}
This is a sample document demonstrating various \LaTeX{} features including mathematical equations, tables, and formatting. You can cite papers using the citation search feature, and they will be automatically added to your \texttt{cite.bib} file.

\section{Mathematical Expressions}
Here are some example equations:

\subsection{Quadratic Formula}
The quadratic formula is given by:
\[x = \frac{-b \pm \sqrt{b^2 - 4ac}}{2a}\]

\subsection{Matrix Example}
A sample matrix:
\[
\begin{pmatrix}
a_{11} & a_{12} & a_{13} \\
a_{21} & a_{22} & a_{23} \\
a_{31} & a_{32} & a_{33}
\end{pmatrix}
\]

\section{Tables}
Below is a sample table showing data:

\begin{table}[h]
\centering
\begin{tabular}{@{}lrrr@{}}
\toprule
\textbf{Category} & \textbf{Value 1} & \textbf{Value 2} & \textbf{Total} \\
\midrule
Item A & 10.5 & 20.3 & 30.8 \\
Item B & 15.2 & 25.1 & 40.3 \\
Item C & 22.1 & 12.7 & 34.8 \\
\bottomrule
\end{tabular}
\caption{Sample Data Table}
\label{tab:sample}
\end{table}

\section{Lists}
\subsection{Bullet Points}
\begin{itemize}
  \item First item
  \item Second item
  \item Third item with sub-items:
    \begin{itemize}
      \item Sub-item 1
      \item Sub-item 2
    \end{itemize}
\end{itemize}

\subsection{Numbered List}
\begin{enumerate}
  \item First step
  \item Second step
  \item Third step
\end{enumerate}

\section{Citations and References}
When you use the citation search feature, references will be automatically managed. For example, you might cite a paper like this: \cite{example2024}. All citations will be automatically added to your \texttt{cite.bib} file.

\section{Conclusion}
This template demonstrates basic \LaTeX{} features including mathematical formulas, tables, and lists. The bibliography system is now set up to work seamlessly with Bibby AI's citation management. Feel free to modify it for your needs!

\bibliographystyle{plain}
\bibliography{cite}

\end{document}\documentclass{article}
\usepackage[utf8]{inputenc}
\usepackage{booktabs}
\usepackage{amsmath}
\usepackage{natbib}

\title{My Document}
\author{Your Name}
\date{\today}

\begin{document}

\maketitle

\section{Introduction}
This is a sample document demonstrating various \LaTeX{} features including mathematical equations, tables, and formatting. You can cite papers using the citation search feature, and they will be automatically added to your \texttt{cite.bib} file.

\section{Mathematical Expressions}
Here are some example equations:

\subsection{Quadratic Formula}
The quadratic formula is given by:
\[x = \frac{-b \pm \sqrt{b^2 - 4ac}}{2a}\]

\subsection{Matrix Example}
A sample matrix:
\[
\begin{pmatrix}
a_{11} & a_{12} & a_{13} \\
a_{21} & a_{22} & a_{23} \\
a_{31} & a_{32} & a_{33}
\end{pmatrix}
\]

\section{Tables}
Below is a sample table showing data:

\begin{table}[h]
\centering
\begin{tabular}{@{}lrrr@{}}
\toprule
\textbf{Category} & \textbf{Value 1} & \textbf{Value 2} & \textbf{Total} \\
\midrule
Item A & 10.5 & 20.3 & 30.8 \\
Item B & 15.2 & 25.1 & 40.3 \\
Item C & 22.1 & 12.7 & 34.8 \\
\bottomrule
\end{tabular}
\caption{Sample Data Table}
\label{tab:sample}
\end{table}

\section{Lists}
\subsection{Bullet Points}
\begin{itemize}
  \item First item
  \item Second item
  \item Third item with sub-items:
    \begin{itemize}
      \item Sub-item 1
      \item Sub-item 2
    \end{itemize}
\end{itemize}

\subsection{Numbered List}
\begin{enumerate}
  \item First step
  \item Second step
  \item Third step
\end{enumerate}

\section{Citations and References}
When you use the citation search feature, references will be automatically managed. For example, you might cite a paper like this: \cite{example2024}. All citations will be automatically added to your \texttt{cite.bib} file.

\section{Conclusion}
This template demonstrates basic \LaTeX{} features including mathematical formulas, tables, and lists. The bibliography system is now set up to work seamlessly with Bibby AI's citation management. Feel free to modify it for your needs!

\bibliographystyle{plain}
\bibliography{cite}

\end{document}\documentclass{article}
\usepackage[utf8]{inputenc}
\usepackage{booktabs}
\usepackage{amsmath}
\usepackage{natbib}

\title{My Document}
\author{Your Name}
\date{\today}

\begin{document}

\maketitle

\section{Introduction}
This is a sample document demonstrating various \LaTeX{} features including mathematical equations, tables, and formatting. You can cite papers using the citation search feature, and they will be automatically added to your \texttt{cite.bib} file.

\section{Mathematical Expressions}
Here are some example equations:

\subsection{Quadratic Formula}
The quadratic formula is given by:
\[x = \frac{-b \pm \sqrt{b^2 - 4ac}}{2a}\]

\subsection{Matrix Example}
A sample matrix:
\[
\begin{pmatrix}
a_{11} & a_{12} & a_{13} \\
a_{21} & a_{22} & a_{23} \\
a_{31} & a_{32} & a_{33}
\end{pmatrix}
\]

\section{Tables}
Below is a sample table showing data:

\begin{table}[h]
\centering
\begin{tabular}{@{}lrrr@{}}
\toprule
\textbf{Category} & \textbf{Value 1} & \textbf{Value 2} & \textbf{Total} \\
\midrule
Item A & 10.5 & 20.3 & 30.8 \\
Item B & 15.2 & 25.1 & 40.3 \\
Item C & 22.1 & 12.7 & 34.8 \\
\bottomrule
\end{tabular}
\caption{Sample Data Table}
\label{tab:sample}
\end{table}

\section{Lists}
\subsection{Bullet Points}
\begin{itemize}
  \item First item
  \item Second item
  \item Third item with sub-items:
    \begin{itemize}
      \item Sub-item 1
      \item Sub-item 2
    \end{itemize}
\end{itemize}

\subsection{Numbered List}
\begin{enumerate}
  \item First step
  \item Second step
  \item Third step
\end{enumerate}

\section{Citations and References}
When you use the citation search feature, references will be automatically managed. For example, you might cite a paper like this: \cite{example2024}. All citations will be automatically added to your \texttt{cite.bib} file.

\section{Conclusion}
This template demonstrates basic \LaTeX{} features including mathematical formulas, tables, and lists. The bibliography system is now set up to work seamlessly with Bibby AI's citation management. Feel free to modify it for your needs!

\bibliographystyle{plain}
\bibliography{cite}

\end{document}\documentclass{article}
\usepackage[utf8]{inputenc}
\usepackage{booktabs}
\usepackage{amsmath}
\usepackage{natbib}

\title{My Document}
\author{Your Name}
\date{\today}

\begin{document}

\maketitle

\section{Introduction}
This is a sample document demonstrating various \LaTeX{} features including mathematical equations, tables, and formatting. You can cite papers using the citation search feature, and they will be automatically added to your \texttt{cite.bib} file.

\section{Mathematical Expressions}
Here are some example equations:

\subsection{Quadratic Formula}
The quadratic formula is given by:
\[x = \frac{-b \pm \sqrt{b^2 - 4ac}}{2a}\]

\subsection{Matrix Example}
A sample matrix:
\[
\begin{pmatrix}
a_{11} & a_{12} & a_{13} \\
a_{21} & a_{22} & a_{23} \\
a_{31} & a_{32} & a_{33}
\end{pmatrix}
\]

\section{Tables}
Below is a sample table showing data:

\begin{table}[h]
\centering
\begin{tabular}{@{}lrrr@{}}
\toprule
\textbf{Category} & \textbf{Value 1} & \textbf{Value 2} & \textbf{Total} \\
\midrule
Item A & 10.5 & 20.3 & 30.8 \\
Item B & 15.2 & 25.1 & 40.3 \\
Item C & 22.1 & 12.7 & 34.8 \\
\bottomrule
\end{tabular}
\caption{Sample Data Table}
\label{tab:sample}
\end{table}

\section{Lists}
\subsection{Bullet Points}
\begin{itemize}
  \item First item
  \item Second item
  \item Third item with sub-items:
    \begin{itemize}
      \item Sub-item 1
      \item Sub-item 2
    \end{itemize}
\end{itemize}

\subsection{Numbered List}
\begin{enumerate}
  \item First step
  \item Second step
  \item Third step
\end{enumerate}

\section{Citations and References}
When you use the citation search feature, references will be automatically managed. For example, you might cite a paper like this: \cite{example2024}. All citations will be automatically added to your \texttt{cite.bib} file.

\section{Conclusion}
This template demonstrates basic \LaTeX{} features including mathematical formulas, tables, and lists. The bibliography system is now set up to work seamlessly with Bibby AI's citation management. Feel free to modify it for your needs!

\bibliographystyle{plain}
\bibliography{cite}

\end{document}\documentclass{article}
\usepackage[utf8]{inputenc}
\usepackage{booktabs}
\usepackage{amsmath}
\usepackage{natbib}

\title{My Document}
\author{Your Name}
\date{\today}

\begin{document}

\maketitle

\section{Introduction}
This is a sample document demonstrating various \LaTeX{} features including mathematical equations, tables, and formatting. You can cite papers using the citation search feature, and they will be automatically added to your \texttt{cite.bib} file.

\section{Mathematical Expressions}
Here are some example equations:

\subsection{Quadratic Formula}
The quadratic formula is given by:
\[x = \frac{-b \pm \sqrt{b^2 - 4ac}}{2a}\]

\subsection{Matrix Example}
A sample matrix:
\[
\begin{pmatrix}
a_{11} & a_{12} & a_{13} \\
a_{21} & a_{22} & a_{23} \\
a_{31} & a_{32} & a_{33}
\end{pmatrix}
\]

\section{Tables}
Below is a sample table showing data:

\begin{table}[h]
\centering
\begin{tabular}{@{}lrrr@{}}
\toprule
\textbf{Category} & \textbf{Value 1} & \textbf{Value 2} & \textbf{Total} \\
\midrule
Item A & 10.5 & 20.3 & 30.8 \\
Item B & 15.2 & 25.1 & 40.3 \\
Item C & 22.1 & 12.7 & 34.8 \\
\bottomrule
\end{tabular}
\caption{Sample Data Table}
\label{tab:sample}
\end{table}

\section{Lists}
\subsection{Bullet Points}
\begin{itemize}
  \item First item
  \item Second item
  \item Third item with sub-items:
    \begin{itemize}
      \item Sub-item 1
      \item Sub-item 2
    \end{itemize}
\end{itemize}

\subsection{Numbered List}
\begin{enumerate}
  \item First step
  \item Second step
  \item Third step
\end{enumerate}

\section{Citations and References}
When you use the citation search feature, references will be automatically managed. For example, you might cite a paper like this: \cite{example2024}. All citations will be automatically added to your \texttt{cite.bib} file.

\section{Conclusion}
This template demonstrates basic \LaTeX{} features including mathematical formulas, tables, and lists. The bibliography system is now set up to work seamlessly with Bibby AI's citation management. Feel free to modify it for your needs!

\bibliographystyle{plain}
\bibliography{cite}

\end{document}\documentclass{article}
\usepackage[utf8]{inputenc}
\usepackage{booktabs}
\usepackage{amsmath}
\usepackage{natbib}

\title{My Document}
\author{Your Name}
\date{\today}

\begin{document}

\maketitle

\section{Introduction}
This is a sample document demonstrating various \LaTeX{} features including mathematical equations, tables, and formatting. You can cite papers using the citation search feature, and they will be automatically added to your \texttt{cite.bib} file.

\section{Mathematical Expressions}
Here are some example equations:

\subsection{Quadratic Formula}
The quadratic formula is given by:
\[x = \frac{-b \pm \sqrt{b^2 - 4ac}}{2a}\]

\subsection{Matrix Example}
A sample matrix:
\[
\begin{pmatrix}
a_{11} & a_{12} & a_{13} \\
a_{21} & a_{22} & a_{23} \\
a_{31} & a_{32} & a_{33}
\end{pmatrix}
\]

\section{Tables}
Below is a sample table showing data:

\begin{table}[h]
\centering
\begin{tabular}{@{}lrrr@{}}
\toprule
\textbf{Category} & \textbf{Value 1} & \textbf{Value 2} & \textbf{Total} \\
\midrule
Item A & 10.5 & 20.3 & 30.8 \\
Item B & 15.2 & 25.1 & 40.3 \\
Item C & 22.1 & 12.7 & 34.8 \\
\bottomrule
\end{tabular}
\caption{Sample Data Table}
\label{tab:sample}
\end{table}

\section{Lists}
\subsection{Bullet Points}
\begin{itemize}
  \item First item
  \item Second item
  \item Third item with sub-items:
    \begin{itemize}
      \item Sub-item 1
      \item Sub-item 2
    \end{itemize}
\end{itemize}

\subsection{Numbered List}
\begin{enumerate}
  \item First step
  \item Second step
  \item Third step
\end{enumerate}

\section{Citations and References}
When you use the citation search feature, references will be automatically managed. For example, you might cite a paper like this: \cite{example2024}. All citations will be automatically added to your \texttt{cite.bib} file.

\section{Conclusion}
This template demonstrates basic \LaTeX{} features including mathematical formulas, tables, and lists. The bibliography system is now set up to work seamlessly with Bibby AI's citation management. Feel free to modify it for your needs!

\bibliographystyle{plain}
\bibliography{cite}

\end{document}\documentclass{article}
\usepackage[utf8]{inputenc}
\usepackage{booktabs}
\usepackage{amsmath}
\usepackage{natbib}

\title{My Document}
\author{Your Name}
\date{\today}

\begin{document}

\maketitle

\section{Introduction}
This is a sample document demonstrating various \LaTeX{} features including mathematical equations, tables, and formatting. You can cite papers using the citation search feature, and they will be automatically added to your \texttt{cite.bib} file.

\section{Mathematical Expressions}
Here are some example equations:

\subsection{Quadratic Formula}
The quadratic formula is given by:
\[x = \frac{-b \pm \sqrt{b^2 - 4ac}}{2a}\]

\subsection{Matrix Example}
A sample matrix:
\[
\begin{pmatrix}
a_{11} & a_{12} & a_{13} \\
a_{21} & a_{22} & a_{23} \\
a_{31} & a_{32} & a_{33}
\end{pmatrix}
\]

\section{Tables}
Below is a sample table showing data:

\begin{table}[h]
\centering
\begin{tabular}{@{}lrrr@{}}
\toprule
\textbf{Category} & \textbf{Value 1} & \textbf{Value 2} & \textbf{Total} \\
\midrule
Item A & 10.5 & 20.3 & 30.8 \\
Item B & 15.2 & 25.1 & 40.3 \\
Item C & 22.1 & 12.7 & 34.8 \\
\bottomrule
\end{tabular}
\caption{Sample Data Table}
\label{tab:sample}
\end{table}

\section{Lists}
\subsection{Bullet Points}
\begin{itemize}
  \item First item
  \item Second item
  \item Third item with sub-items:
    \begin{itemize}
      \item Sub-item 1
      \item Sub-item 2
    \end{itemize}
\end{itemize}

\subsection{Numbered List}
\begin{enumerate}
  \item First step
  \item Second step
  \item Third step
\end{enumerate}

\section{Citations and References}
When you use the citation search feature, references will be automatically managed. For example, you might cite a paper like this: \cite{example2024}. All citations will be automatically added to your \texttt{cite.bib} file.

\section{Conclusion}
This template demonstrates basic \LaTeX{} features including mathematical formulas, tables, and lists. The bibliography system is now set up to work seamlessly with Bibby AI's citation management. Feel free to modify it for your needs!

\bibliographystyle{plain}
\bibliography{cite}

\end{document}\documentclass{article}
\usepackage[utf8]{inputenc}
\usepackage{booktabs}
\usepackage{amsmath}
\usepackage{natbib}

\title{My Document}
\author{Your Name}
\date{\today}

\begin{document}

\maketitle

\section{Introduction}
This is a sample document demonstrating various \LaTeX{} features including mathematical equations, tables, and formatting. You can cite papers using the citation search feature, and they will be automatically added to your \texttt{cite.bib} file.

\section{Mathematical Expressions}
Here are some example equations:

\subsection{Quadratic Formula}
The quadratic formula is given by:
\[x = \frac{-b \pm \sqrt{b^2 - 4ac}}{2a}\]

\subsection{Matrix Example}
A sample matrix:
\[
\begin{pmatrix}
a_{11} & a_{12} & a_{13} \\
a_{21} & a_{22} & a_{23} \\
a_{31} & a_{32} & a_{33}
\end{pmatrix}
\]

\section{Tables}
Below is a sample table showing data:

\begin{table}[h]
\centering
\begin{tabular}{@{}lrrr@{}}
\toprule
\textbf{Category} & \textbf{Value 1} & \textbf{Value 2} & \textbf{Total} \\
\midrule
Item A & 10.5 & 20.3 & 30.8 \\
Item B & 15.2 & 25.1 & 40.3 \\
Item C & 22.1 & 12.7 & 34.8 \\
\bottomrule
\end{tabular}
\caption{Sample Data Table}
\label{tab:sample}
\end{table}

\section{Lists}
\subsection{Bullet Points}
\begin{itemize}
  \item First item
  \item Second item
  \item Third item with sub-items:
    \begin{itemize}
      \item Sub-item 1
      \item Sub-item 2
    \end{itemize}
\end{itemize}

\subsection{Numbered List}
\begin{enumerate}
  \item First step
  \item Second step
  \item Third step
\end{enumerate}

\section{Citations and References}
When you use the citation search feature, references will be automatically managed. For example, you might cite a paper like this: \cite{example2024}. All citations will be automatically added to your \texttt{cite.bib} file.

\section{Conclusion}
This template demonstrates basic \LaTeX{} features including mathematical formulas, tables, and lists. The bibliography system is now set up to work seamlessly with Bibby AI's citation management. Feel free to modify it for your needs!

\bibliographystyle{plain}
\bibliography{cite}

\end{document}\documentclass{article}
\usepackage[utf8]{inputenc}
\usepackage{booktabs}
\usepackage{amsmath}
\usepackage{natbib}

\title{My Document}
\author{Your Name}
\date{\today}

\begin{document}

\maketitle

\section{Introduction}
This is a sample document demonstrating various \LaTeX{} features including mathematical equations, tables, and formatting. You can cite papers using the citation search feature, and they will be automatically added to your \texttt{cite.bib} file.

\section{Mathematical Expressions}
Here are some example equations:

\subsection{Quadratic Formula}
The quadratic formula is given by:
\[x = \frac{-b \pm \sqrt{b^2 - 4ac}}{2a}\]

\subsection{Matrix Example}
A sample matrix:
\[
\begin{pmatrix}
a_{11} & a_{12} & a_{13} \\
a_{21} & a_{22} & a_{23} \\
a_{31} & a_{32} & a_{33}
\end{pmatrix}
\]

\section{Tables}
Below is a sample table showing data:

\begin{table}[h]
\centering
\begin{tabular}{@{}lrrr@{}}
\toprule
\textbf{Category} & \textbf{Value 1} & \textbf{Value 2} & \textbf{Total} \\
\midrule
Item A & 10.5 & 20.3 & 30.8 \\
Item B & 15.2 & 25.1 & 40.3 \\
Item C & 22.1 & 12.7 & 34.8 \\
\bottomrule
\end{tabular}
\caption{Sample Data Table}
\label{tab:sample}
\end{table}

\section{Lists}
\subsection{Bullet Points}
\begin{itemize}
  \item First item
  \item Second item
  \item Third item with sub-items:
    \begin{itemize}
      \item Sub-item 1
      \item Sub-item 2
    \end{itemize}
\end{itemize}

\subsection{Numbered List}
\begin{enumerate}
  \item First step
  \item Second step
  \item Third step
\end{enumerate}

\section{Citations and References}
When you use the citation search feature, references will be automatically managed. For example, you might cite a paper like this: \cite{example2024}. All citations will be automatically added to your \texttt{cite.bib} file.

\section{Conclusion}
This template demonstrates basic \LaTeX{} features including mathematical formulas, tables, and lists. The bibliography system is now set up to work seamlessly with Bibby AI's citation management. Feel free to modify it for your needs!

\bibliographystyle{plain}
\bibliography{cite}

\end{document}\documentclass{article}
\usepackage[utf8]{inputenc}
\usepackage{booktabs}
\usepackage{amsmath}
\usepackage{natbib}

\title{My Document}
\author{Your Name}
\date{\today}

\begin{document}

\maketitle

\section{Introduction}
This is a sample document demonstrating various \LaTeX{} features including mathematical equations, tables, and formatting. You can cite papers using the citation search feature, and they will be automatically added to your \texttt{cite.bib} file.

\section{Mathematical Expressions}
Here are some example equations:

\subsection{Quadratic Formula}
The quadratic formula is given by:
\[x = \frac{-b \pm \sqrt{b^2 - 4ac}}{2a}\]

\subsection{Matrix Example}
A sample matrix:
\[
\begin{pmatrix}
a_{11} & a_{12} & a_{13} \\
a_{21} & a_{22} & a_{23} \\
a_{31} & a_{32} & a_{33}
\end{pmatrix}
\]

\section{Tables}
Below is a sample table showing data:

\begin{table}[h]
\centering
\begin{tabular}{@{}lrrr@{}}
\toprule
\textbf{Category} & \textbf{Value 1} & \textbf{Value 2} & \textbf{Total} \\
\midrule
Item A & 10.5 & 20.3 & 30.8 \\
Item B & 15.2 & 25.1 & 40.3 \\
Item C & 22.1 & 12.7 & 34.8 \\
\bottomrule
\end{tabular}
\caption{Sample Data Table}
\label{tab:sample}
\end{table}

\section{Lists}
\subsection{Bullet Points}
\begin{itemize}
  \item First item
  \item Second item
  \item Third item with sub-items:
    \begin{itemize}
      \item Sub-item 1
      \item Sub-item 2
    \end{itemize}
\end{itemize}

\subsection{Numbered List}
\begin{enumerate}
  \item First step
  \item Second step
  \item Third step
\end{enumerate}

\section{Citations and References}
When you use the citation search feature, references will be automatically managed. For example, you might cite a paper like this: \cite{example2024}. All citations will be automatically added to your \texttt{cite.bib} file.

\section{Conclusion}
This template demonstrates basic \LaTeX{} features including mathematical formulas, tables, and lists. The bibliography system is now set up to work seamlessly with Bibby AI's citation management. Feel free to modify it for your needs!

\bibliographystyle{plain}
\bibliography{cite}

\end{document}\documentclass{article}
\usepackage[utf8]{inputenc}
\usepackage{booktabs}
\usepackage{amsmath}
\usepackage{natbib}

\title{My Document}
\author{Your Name}
\date{\today}

\begin{document}

\maketitle

\section{Introduction}
This is a sample document demonstrating various \LaTeX{} features including mathematical equations, tables, and formatting. You can cite papers using the citation search feature, and they will be automatically added to your \texttt{cite.bib} file.

\section{Mathematical Expressions}
Here are some example equations:

\subsection{Quadratic Formula}
The quadratic formula is given by:
\[x = \frac{-b \pm \sqrt{b^2 - 4ac}}{2a}\]

\subsection{Matrix Example}
A sample matrix:
\[
\begin{pmatrix}
a_{11} & a_{12} & a_{13} \\
a_{21} & a_{22} & a_{23} \\
a_{31} & a_{32} & a_{33}
\end{pmatrix}
\]

\section{Tables}
Below is a sample table showing data:

\begin{table}[h]
\centering
\begin{tabular}{@{}lrrr@{}}
\toprule
\textbf{Category} & \textbf{Value 1} & \textbf{Value 2} & \textbf{Total} \\
\midrule
Item A & 10.5 & 20.3 & 30.8 \\
Item B & 15.2 & 25.1 & 40.3 \\
Item C & 22.1 & 12.7 & 34.8 \\
\bottomrule
\end{tabular}
\caption{Sample Data Table}
\label{tab:sample}
\end{table}

\section{Lists}
\subsection{Bullet Points}
\begin{itemize}
  \item First item
  \item Second item
  \item Third item with sub-items:
    \begin{itemize}
      \item Sub-item 1
      \item Sub-item 2
    \end{itemize}
\end{itemize}

\subsection{Numbered List}
\begin{enumerate}
  \item First step
  \item Second step
  \item Third step
\end{enumerate}

\section{Citations and References}
When you use the citation search feature, references will be automatically managed. For example, you might cite a paper like this: \cite{example2024}. All citations will be automatically added to your \texttt{cite.bib} file.

\section{Conclusion}
This template demonstrates basic \LaTeX{} features including mathematical formulas, tables, and lists. The bibliography system is now set up to work seamlessly with Bibby AI's citation management. Feel free to modify it for your needs!

\bibliographystyle{plain}
\bibliography{cite}

\end{document}\documentclass{article}
\usepackage[utf8]{inputenc}
\usepackage{booktabs}
\usepackage{amsmath}
\usepackage{natbib}

\title{My Document}
\author{Your Name}
\date{\today}

\begin{document}

\maketitle

\section{Introduction}
This is a sample document demonstrating various \LaTeX{} features including mathematical equations, tables, and formatting. You can cite papers using the citation search feature, and they will be automatically added to your \texttt{cite.bib} file.

\section{Mathematical Expressions}
Here are some example equations:

\subsection{Quadratic Formula}
The quadratic formula is given by:
\[x = \frac{-b \pm \sqrt{b^2 - 4ac}}{2a}\]

\subsection{Matrix Example}
A sample matrix:
\[
\begin{pmatrix}
a_{11} & a_{12} & a_{13} \\
a_{21} & a_{22} & a_{23} \\
a_{31} & a_{32} & a_{33}
\end{pmatrix}
\]

\section{Tables}
Below is a sample table showing data:

\begin{table}[h]
\centering
\begin{tabular}{@{}lrrr@{}}
\toprule
\textbf{Category} & \textbf{Value 1} & \textbf{Value 2} & \textbf{Total} \\
\midrule
Item A & 10.5 & 20.3 & 30.8 \\
Item B & 15.2 & 25.1 & 40.3 \\
Item C & 22.1 & 12.7 & 34.8 \\
\bottomrule
\end{tabular}
\caption{Sample Data Table}
\label{tab:sample}
\end{table}

\section{Lists}
\subsection{Bullet Points}
\begin{itemize}
  \item First item
  \item Second item
  \item Third item with sub-items:
    \begin{itemize}
      \item Sub-item 1
      \item Sub-item 2
    \end{itemize}
\end{itemize}

\subsection{Numbered List}
\begin{enumerate}
  \item First step
  \item Second step
  \item Third step
\end{enumerate}

\section{Citations and References}
When you use the citation search feature, references will be automatically managed. For example, you might cite a paper like this: \cite{example2024}. All citations will be automatically added to your \texttt{cite.bib} file.

\section{Conclusion}
This template demonstrates basic \LaTeX{} features including mathematical formulas, tables, and lists. The bibliography system is now set up to work seamlessly with Bibby AI's citation management. Feel free to modify it for your needs!

\bibliographystyle{plain}
\bibliography{cite}

\end{document}\documentclass{article}
\usepackage[utf8]{inputenc}
\usepackage{booktabs}
\usepackage{amsmath}
\usepackage{natbib}

\title{My Document}
\author{Your Name}
\date{\today}

\begin{document}

\maketitle

\section{Introduction}
This is a sample document demonstrating various \LaTeX{} features including mathematical equations, tables, and formatting. You can cite papers using the citation search feature, and they will be automatically added to your \texttt{cite.bib} file.

\section{Mathematical Expressions}
Here are some example equations:

\subsection{Quadratic Formula}
The quadratic formula is given by:
\[x = \frac{-b \pm \sqrt{b^2 - 4ac}}{2a}\]

\subsection{Matrix Example}
A sample matrix:
\[
\begin{pmatrix}
a_{11} & a_{12} & a_{13} \\
a_{21} & a_{22} & a_{23} \\
a_{31} & a_{32} & a_{33}
\end{pmatrix}
\]

\section{Tables}
Below is a sample table showing data:

\begin{table}[h]
\centering
\begin{tabular}{@{}lrrr@{}}
\toprule
\textbf{Category} & \textbf{Value 1} & \textbf{Value 2} & \textbf{Total} \\
\midrule
Item A & 10.5 & 20.3 & 30.8 \\
Item B & 15.2 & 25.1 & 40.3 \\
Item C & 22.1 & 12.7 & 34.8 \\
\bottomrule
\end{tabular}
\caption{Sample Data Table}
\label{tab:sample}
\end{table}

\section{Lists}
\subsection{Bullet Points}
\begin{itemize}
  \item First item
  \item Second item
  \item Third item with sub-items:
    \begin{itemize}
      \item Sub-item 1
      \item Sub-item 2
    \end{itemize}
\end{itemize}

\subsection{Numbered List}
\begin{enumerate}
  \item First step
  \item Second step
  \item Third step
\end{enumerate}

\section{Citations and References}
When you use the citation search feature, references will be automatically managed. For example, you might cite a paper like this: \cite{example2024}. All citations will be automatically added to your \texttt{cite.bib} file.

\section{Conclusion}
This template demonstrates basic \LaTeX{} features including mathematical formulas, tables, and lists. The bibliography system is now set up to work seamlessly with Bibby AI's citation management. Feel free to modify it for your needs!

\bibliographystyle{plain}
\bibliography{cite}

\end{document}\documentclass{article}
\usepackage[utf8]{inputenc}
\usepackage{booktabs}
\usepackage{amsmath}
\usepackage{natbib}

\title{My Document}
\author{Your Name}
\date{\today}

\begin{document}

\maketitle

\section{Introduction}
This is a sample document demonstrating various \LaTeX{} features including mathematical equations, tables, and formatting. You can cite papers using the citation search feature, and they will be automatically added to your \texttt{cite.bib} file.

\section{Mathematical Expressions}
Here are some example equations:

\subsection{Quadratic Formula}
The quadratic formula is given by:
\[x = \frac{-b \pm \sqrt{b^2 - 4ac}}{2a}\]

\subsection{Matrix Example}
A sample matrix:
\[
\begin{pmatrix}
a_{11} & a_{12} & a_{13} \\
a_{21} & a_{22} & a_{23} \\
a_{31} & a_{32} & a_{33}
\end{pmatrix}
\]

\section{Tables}
Below is a sample table showing data:

\begin{table}[h]
\centering
\begin{tabular}{@{}lrrr@{}}
\toprule
\textbf{Category} & \textbf{Value 1} & \textbf{Value 2} & \textbf{Total} \\
\midrule
Item A & 10.5 & 20.3 & 30.8 \\
Item B & 15.2 & 25.1 & 40.3 \\
Item C & 22.1 & 12.7 & 34.8 \\
\bottomrule
\end{tabular}
\caption{Sample Data Table}
\label{tab:sample}
\end{table}

\section{Lists}
\subsection{Bullet Points}
\begin{itemize}
  \item First item
  \item Second item
  \item Third item with sub-items:
    \begin{itemize}
      \item Sub-item 1
      \item Sub-item 2
    \end{itemize}
\end{itemize}

\subsection{Numbered List}
\begin{enumerate}
  \item First step
  \item Second step
  \item Third step
\end{enumerate}

\section{Citations and References}
When you use the citation search feature, references will be automatically managed. For example, you might cite a paper like this: \cite{example2024}. All citations will be automatically added to your \texttt{cite.bib} file.

\section{Conclusion}
This template demonstrates basic \LaTeX{} features including mathematical formulas, tables, and lists. The bibliography system is now set up to work seamlessly with Bibby AI's citation management. Feel free to modify it for your needs!

\bibliographystyle{plain}
\bibliography{cite}

\end{document}\documentclass{article}
\usepackage[utf8]{inputenc}
\usepackage{booktabs}
\usepackage{amsmath}
\usepackage{natbib}

\title{My Document}
\author{Your Name}
\date{\today}

\begin{document}

\maketitle

\section{Introduction}
This is a sample document demonstrating various \LaTeX{} features including mathematical equations, tables, and formatting. You can cite papers using the citation search feature, and they will be automatically added to your \texttt{cite.bib} file.

\section{Mathematical Expressions}
Here are some example equations:

\subsection{Quadratic Formula}
The quadratic formula is given by:
\[x = \frac{-b \pm \sqrt{b^2 - 4ac}}{2a}\]

\subsection{Matrix Example}
A sample matrix:
\[
\begin{pmatrix}
a_{11} & a_{12} & a_{13} \\
a_{21} & a_{22} & a_{23} \\
a_{31} & a_{32} & a_{33}
\end{pmatrix}
\]

\section{Tables}
Below is a sample table showing data:

\begin{table}[h]
\centering
\begin{tabular}{@{}lrrr@{}}
\toprule
\textbf{Category} & \textbf{Value 1} & \textbf{Value 2} & \textbf{Total} \\
\midrule
Item A & 10.5 & 20.3 & 30.8 \\
Item B & 15.2 & 25.1 & 40.3 \\
Item C & 22.1 & 12.7 & 34.8 \\
\bottomrule
\end{tabular}
\caption{Sample Data Table}
\label{tab:sample}
\end{table}

\section{Lists}
\subsection{Bullet Points}
\begin{itemize}
  \item First item
  \item Second item
  \item Third item with sub-items:
    \begin{itemize}
      \item Sub-item 1
      \item Sub-item 2
    \end{itemize}
\end{itemize}

\subsection{Numbered List}
\begin{enumerate}
  \item First step
  \item Second step
  \item Third step
\end{enumerate}

\section{Citations and References}
When you use the citation search feature, references will be automatically managed. For example, you might cite a paper like this: \cite{example2024}. All citations will be automatically added to your \texttt{cite.bib} file.

\section{Conclusion}
This template demonstrates basic \LaTeX{} features including mathematical formulas, tables, and lists. The bibliography system is now set up to work seamlessly with Bibby AI's citation management. Feel free to modify it for your needs!

\bibliographystyle{plain}
\bibliography{cite}

\end{document}\documentclass{article}
\usepackage[utf8]{inputenc}
\usepackage{booktabs}
\usepackage{amsmath}
\usepackage{natbib}

\title{My Document}
\author{Your Name}
\date{\today}

\begin{document}

\maketitle

\section{Introduction}
This is a sample document demonstrating various \LaTeX{} features including mathematical equations, tables, and formatting. You can cite papers using the citation search feature, and they will be automatically added to your \texttt{cite.bib} file.

\section{Mathematical Expressions}
Here are some example equations:

\subsection{Quadratic Formula}
The quadratic formula is given by:
\[x = \frac{-b \pm \sqrt{b^2 - 4ac}}{2a}\]

\subsection{Matrix Example}
A sample matrix:
\[
\begin{pmatrix}
a_{11} & a_{12} & a_{13} \\
a_{21} & a_{22} & a_{23} \\
a_{31} & a_{32} & a_{33}
\end{pmatrix}
\]

\section{Tables}
Below is a sample table showing data:

\begin{table}[h]
\centering
\begin{tabular}{@{}lrrr@{}}
\toprule
\textbf{Category} & \textbf{Value 1} & \textbf{Value 2} & \textbf{Total} \\
\midrule
Item A & 10.5 & 20.3 & 30.8 \\
Item B & 15.2 & 25.1 & 40.3 \\
Item C & 22.1 & 12.7 & 34.8 \\
\bottomrule
\end{tabular}
\caption{Sample Data Table}
\label{tab:sample}
\end{table}

\section{Lists}
\subsection{Bullet Points}
\begin{itemize}
  \item First item
  \item Second item
  \item Third item with sub-items:
    \begin{itemize}
      \item Sub-item 1
      \item Sub-item 2
    \end{itemize}
\end{itemize}

\subsection{Numbered List}
\begin{enumerate}
  \item First step
  \item Second step
  \item Third step
\end{enumerate}

\section{Citations and References}
When you use the citation search feature, references will be automatically managed. For example, you might cite a paper like this: \cite{example2024}. All citations will be automatically added to your \texttt{cite.bib} file.

\section{Conclusion}
This template demonstrates basic \LaTeX{} features including mathematical formulas, tables, and lists. The bibliography system is now set up to work seamlessly with Bibby AI's citation management. Feel free to modify it for your needs!

\bibliographystyle{plain}
\bibliography{cite}

\end{document}\documentclass{article}
\usepackage[utf8]{inputenc}
\usepackage{booktabs}
\usepackage{amsmath}
\usepackage{natbib}

\title{My Document}
\author{Your Name}
\date{\today}

\begin{document}

\maketitle

\section{Introduction}
This is a sample document demonstrating various \LaTeX{} features including mathematical equations, tables, and formatting. You can cite papers using the citation search feature, and they will be automatically added to your \texttt{cite.bib} file.

\section{Mathematical Expressions}
Here are some example equations:

\subsection{Quadratic Formula}
The quadratic formula is given by:
\[x = \frac{-b \pm \sqrt{b^2 - 4ac}}{2a}\]

\subsection{Matrix Example}
A sample matrix:
\[
\begin{pmatrix}
a_{11} & a_{12} & a_{13} \\
a_{21} & a_{22} & a_{23} \\
a_{31} & a_{32} & a_{33}
\end{pmatrix}
\]

\section{Tables}
Below is a sample table showing data:

\begin{table}[h]
\centering
\begin{tabular}{@{}lrrr@{}}
\toprule
\textbf{Category} & \textbf{Value 1} & \textbf{Value 2} & \textbf{Total} \\
\midrule
Item A & 10.5 & 20.3 & 30.8 \\
Item B & 15.2 & 25.1 & 40.3 \\
Item C & 22.1 & 12.7 & 34.8 \\
\bottomrule
\end{tabular}
\caption{Sample Data Table}
\label{tab:sample}
\end{table}

\section{Lists}
\subsection{Bullet Points}
\begin{itemize}
  \item First item
  \item Second item
  \item Third item with sub-items:
    \begin{itemize}
      \item Sub-item 1
      \item Sub-item 2
    \end{itemize}
\end{itemize}

\subsection{Numbered List}
\begin{enumerate}
  \item First step
  \item Second step
  \item Third step
\end{enumerate}

\section{Citations and References}
When you use the citation search feature, references will be automatically managed. For example, you might cite a paper like this: \cite{example2024}. All citations will be automatically added to your \texttt{cite.bib} file.

\section{Conclusion}
This template demonstrates basic \LaTeX{} features including mathematical formulas, tables, and lists. The bibliography system is now set up to work seamlessly with Bibby AI's citation management. Feel free to modify it for your needs!

\bibliographystyle{plain}
\bibliography{cite}

\end{document}\documentclass{article}
\usepackage[utf8]{inputenc}
\usepackage{booktabs}
\usepackage{amsmath}
\usepackage{natbib}

\title{My Document}
\author{Your Name}
\date{\today}

\begin{document}

\maketitle

\section{Introduction}
This is a sample document demonstrating various \LaTeX{} features including mathematical equations, tables, and formatting. You can cite papers using the citation search feature, and they will be automatically added to your \texttt{cite.bib} file.

\section{Mathematical Expressions}
Here are some example equations:

\subsection{Quadratic Formula}
The quadratic formula is given by:
\[x = \frac{-b \pm \sqrt{b^2 - 4ac}}{2a}\]

\subsection{Matrix Example}
A sample matrix:
\[
\begin{pmatrix}
a_{11} & a_{12} & a_{13} \\
a_{21} & a_{22} & a_{23} \\
a_{31} & a_{32} & a_{33}
\end{pmatrix}
\]

\section{Tables}
Below is a sample table showing data:

\begin{table}[h]
\centering
\begin{tabular}{@{}lrrr@{}}
\toprule
\textbf{Category} & \textbf{Value 1} & \textbf{Value 2} & \textbf{Total} \\
\midrule
Item A & 10.5 & 20.3 & 30.8 \\
Item B & 15.2 & 25.1 & 40.3 \\
Item C & 22.1 & 12.7 & 34.8 \\
\bottomrule
\end{tabular}
\caption{Sample Data Table}
\label{tab:sample}
\end{table}

\section{Lists}
\subsection{Bullet Points}
\begin{itemize}
  \item First item
  \item Second item
  \item Third item with sub-items:
    \begin{itemize}
      \item Sub-item 1
      \item Sub-item 2
    \end{itemize}
\end{itemize}

\subsection{Numbered List}
\begin{enumerate}
  \item First step
  \item Second step
  \item Third step
\end{enumerate}

\section{Citations and References}
When you use the citation search feature, references will be automatically managed. For example, you might cite a paper like this: \cite{example2024}. All citations will be automatically added to your \texttt{cite.bib} file.

\section{Conclusion}
This template demonstrates basic \LaTeX{} features including mathematical formulas, tables, and lists. The bibliography system is now set up to work seamlessly with Bibby AI's citation management. Feel free to modify it for your needs!

\bibliographystyle{plain}
\bibliography{cite}

\end{document}\documentclass{article}
\usepackage[utf8]{inputenc}
\usepackage{booktabs}
\usepackage{amsmath}
\usepackage{natbib}

\title{My Document}
\author{Your Name}
\date{\today}

\begin{document}

\maketitle

\section{Introduction}
This is a sample document demonstrating various \LaTeX{} features including mathematical equations, tables, and formatting. You can cite papers using the citation search feature, and they will be automatically added to your \texttt{cite.bib} file.

\section{Mathematical Expressions}
Here are some example equations:

\subsection{Quadratic Formula}
The quadratic formula is given by:
\[x = \frac{-b \pm \sqrt{b^2 - 4ac}}{2a}\]

\subsection{Matrix Example}
A sample matrix:
\[
\begin{pmatrix}
a_{11} & a_{12} & a_{13} \\
a_{21} & a_{22} & a_{23} \\
a_{31} & a_{32} & a_{33}
\end{pmatrix}
\]

\section{Tables}
Below is a sample table showing data:

\begin{table}[h]
\centering
\begin{tabular}{@{}lrrr@{}}
\toprule
\textbf{Category} & \textbf{Value 1} & \textbf{Value 2} & \textbf{Total} \\
\midrule
Item A & 10.5 & 20.3 & 30.8 \\
Item B & 15.2 & 25.1 & 40.3 \\
Item C & 22.1 & 12.7 & 34.8 \\
\bottomrule
\end{tabular}
\caption{Sample Data Table}
\label{tab:sample}
\end{table}

\section{Lists}
\subsection{Bullet Points}
\begin{itemize}
  \item First item
  \item Second item
  \item Third item with sub-items:
    \begin{itemize}
      \item Sub-item 1
      \item Sub-item 2
    \end{itemize}
\end{itemize}

\subsection{Numbered List}
\begin{enumerate}
  \item First step
  \item Second step
  \item Third step
\end{enumerate}

\section{Citations and References}
When you use the citation search feature, references will be automatically managed. For example, you might cite a paper like this: \cite{example2024}. All citations will be automatically added to your \texttt{cite.bib} file.

\section{Conclusion}
This template demonstrates basic \LaTeX{} features including mathematical formulas, tables, and lists. The bibliography system is now set up to work seamlessly with Bibby AI's citation management. Feel free to modify it for your needs!

\bibliographystyle{plain}
\bibliography{cite}

\end{document}\documentclass{article}
\usepackage[utf8]{inputenc}
\usepackage{booktabs}
\usepackage{amsmath}
\usepackage{natbib}

\title{My Document}
\author{Your Name}
\date{\today}

\begin{document}

\maketitle

\section{Introduction}
This is a sample document demonstrating various \LaTeX{} features including mathematical equations, tables, and formatting. You can cite papers using the citation search feature, and they will be automatically added to your \texttt{cite.bib} file.

\section{Mathematical Expressions}
Here are some example equations:

\subsection{Quadratic Formula}
The quadratic formula is given by:
\[x = \frac{-b \pm \sqrt{b^2 - 4ac}}{2a}\]

\subsection{Matrix Example}
A sample matrix:
\[
\begin{pmatrix}
a_{11} & a_{12} & a_{13} \\
a_{21} & a_{22} & a_{23} \\
a_{31} & a_{32} & a_{33}
\end{pmatrix}
\]

\section{Tables}
Below is a sample table showing data:

\begin{table}[h]
\centering
\begin{tabular}{@{}lrrr@{}}
\toprule
\textbf{Category} & \textbf{Value 1} & \textbf{Value 2} & \textbf{Total} \\
\midrule
Item A & 10.5 & 20.3 & 30.8 \\
Item B & 15.2 & 25.1 & 40.3 \\
Item C & 22.1 & 12.7 & 34.8 \\
\bottomrule
\end{tabular}
\caption{Sample Data Table}
\label{tab:sample}
\end{table}

\section{Lists}
\subsection{Bullet Points}
\begin{itemize}
  \item First item
  \item Second item
  \item Third item with sub-items:
    \begin{itemize}
      \item Sub-item 1
      \item Sub-item 2
    \end{itemize}
\end{itemize}

\subsection{Numbered List}
\begin{enumerate}
  \item First step
  \item Second step
  \item Third step
\end{enumerate}

\section{Citations and References}
When you use the citation search feature, references will be automatically managed. For example, you might cite a paper like this: \cite{example2024}. All citations will be automatically added to your \texttt{cite.bib} file.

\section{Conclusion}
This template demonstrates basic \LaTeX{} features including mathematical formulas, tables, and lists. The bibliography system is now set up to work seamlessly with Bibby AI's citation management. Feel free to modify it for your needs!

\bibliographystyle{plain}
\bibliography{cite}

\end{document}\documentclass{article}
\usepackage[utf8]{inputenc}
\usepackage{booktabs}
\usepackage{amsmath}
\usepackage{natbib}

\title{My Document}
\author{Your Name}
\date{\today}

\begin{document}

\maketitle

\section{Introduction}
This is a sample document demonstrating various \LaTeX{} features including mathematical equations, tables, and formatting. You can cite papers using the citation search feature, and they will be automatically added to your \texttt{cite.bib} file.

\section{Mathematical Expressions}
Here are some example equations:

\subsection{Quadratic Formula}
The quadratic formula is given by:
\[x = \frac{-b \pm \sqrt{b^2 - 4ac}}{2a}\]

\subsection{Matrix Example}
A sample matrix:
\[
\begin{pmatrix}
a_{11} & a_{12} & a_{13} \\
a_{21} & a_{22} & a_{23} \\
a_{31} & a_{32} & a_{33}
\end{pmatrix}
\]

\section{Tables}
Below is a sample table showing data:

\begin{table}[h]
\centering
\begin{tabular}{@{}lrrr@{}}
\toprule
\textbf{Category} & \textbf{Value 1} & \textbf{Value 2} & \textbf{Total} \\
\midrule
Item A & 10.5 & 20.3 & 30.8 \\
Item B & 15.2 & 25.1 & 40.3 \\
Item C & 22.1 & 12.7 & 34.8 \\
\bottomrule
\end{tabular}
\caption{Sample Data Table}
\label{tab:sample}
\end{table}

\section{Lists}
\subsection{Bullet Points}
\begin{itemize}
  \item First item
  \item Second item
  \item Third item with sub-items:
    \begin{itemize}
      \item Sub-item 1
      \item Sub-item 2
    \end{itemize}
\end{itemize}

\subsection{Numbered List}
\begin{enumerate}
  \item First step
  \item Second step
  \item Third step
\end{enumerate}

\section{Citations and References}
When you use the citation search feature, references will be automatically managed. For example, you might cite a paper like this: \cite{example2024}. All citations will be automatically added to your \texttt{cite.bib} file.

\section{Conclusion}
This template demonstrates basic \LaTeX{} features including mathematical formulas, tables, and lists. The bibliography system is now set up to work seamlessly with Bibby AI's citation management. Feel free to modify it for your needs!

\bibliographystyle{plain}
\bibliography{cite}

\end{document}\documentclass{article}
\usepackage[utf8]{inputenc}
\usepackage{booktabs}
\usepackage{amsmath}
\usepackage{natbib}

\title{My Document}
\author{Your Name}
\date{\today}

\begin{document}

\maketitle

\section{Introduction}
This is a sample document demonstrating various \LaTeX{} features including mathematical equations, tables, and formatting. You can cite papers using the citation search feature, and they will be automatically added to your \texttt{cite.bib} file.

\section{Mathematical Expressions}
Here are some example equations:

\subsection{Quadratic Formula}
The quadratic formula is given by:
\[x = \frac{-b \pm \sqrt{b^2 - 4ac}}{2a}\]

\subsection{Matrix Example}
A sample matrix:
\[
\begin{pmatrix}
a_{11} & a_{12} & a_{13} \\
a_{21} & a_{22} & a_{23} \\
a_{31} & a_{32} & a_{33}
\end{pmatrix}
\]

\section{Tables}
Below is a sample table showing data:

\begin{table}[h]
\centering
\begin{tabular}{@{}lrrr@{}}
\toprule
\textbf{Category} & \textbf{Value 1} & \textbf{Value 2} & \textbf{Total} \\
\midrule
Item A & 10.5 & 20.3 & 30.8 \\
Item B & 15.2 & 25.1 & 40.3 \\
Item C & 22.1 & 12.7 & 34.8 \\
\bottomrule
\end{tabular}
\caption{Sample Data Table}
\label{tab:sample}
\end{table}

\section{Lists}
\subsection{Bullet Points}
\begin{itemize}
  \item First item
  \item Second item
  \item Third item with sub-items:
    \begin{itemize}
      \item Sub-item 1
      \item Sub-item 2
    \end{itemize}
\end{itemize}

\subsection{Numbered List}
\begin{enumerate}
  \item First step
  \item Second step
  \item Third step
\end{enumerate}

\section{Citations and References}
When you use the citation search feature, references will be automatically managed. For example, you might cite a paper like this: \cite{example2024}. All citations will be automatically added to your \texttt{cite.bib} file.

\section{Conclusion}
This template demonstrates basic \LaTeX{} features including mathematical formulas, tables, and lists. The bibliography system is now set up to work seamlessly with Bibby AI's citation management. Feel free to modify it for your needs!

\bibliographystyle{plain}
\bibliography{cite}

\end{document}\documentclass{article}
\usepackage[utf8]{inputenc}
\usepackage{booktabs}
\usepackage{amsmath}
\usepackage{natbib}

\title{My Document}
\author{Your Name}
\date{\today}

\begin{document}

\maketitle

\section{Introduction}
This is a sample document demonstrating various \LaTeX{} features including mathematical equations, tables, and formatting. You can cite papers using the citation search feature, and they will be automatically added to your \texttt{cite.bib} file.

\section{Mathematical Expressions}
Here are some example equations:

\subsection{Quadratic Formula}
The quadratic formula is given by:
\[x = \frac{-b \pm \sqrt{b^2 - 4ac}}{2a}\]

\subsection{Matrix Example}
A sample matrix:
\[
\begin{pmatrix}
a_{11} & a_{12} & a_{13} \\
a_{21} & a_{22} & a_{23} \\
a_{31} & a_{32} & a_{33}
\end{pmatrix}
\]

\section{Tables}
Below is a sample table showing data:

\begin{table}[h]
\centering
\begin{tabular}{@{}lrrr@{}}
\toprule
\textbf{Category} & \textbf{Value 1} & \textbf{Value 2} & \textbf{Total} \\
\midrule
Item A & 10.5 & 20.3 & 30.8 \\
Item B & 15.2 & 25.1 & 40.3 \\
Item C & 22.1 & 12.7 & 34.8 \\
\bottomrule
\end{tabular}
\caption{Sample Data Table}
\label{tab:sample}
\end{table}

\section{Lists}
\subsection{Bullet Points}
\begin{itemize}
  \item First item
  \item Second item
  \item Third item with sub-items:
    \begin{itemize}
      \item Sub-item 1
      \item Sub-item 2
    \end{itemize}
\end{itemize}

\subsection{Numbered List}
\begin{enumerate}
  \item First step
  \item Second step
  \item Third step
\end{enumerate}

\section{Citations and References}
When you use the citation search feature, references will be automatically managed. For example, you might cite a paper like this: \cite{example2024}. All citations will be automatically added to your \texttt{cite.bib} file.

\section{Conclusion}
This template demonstrates basic \LaTeX{} features including mathematical formulas, tables, and lists. The bibliography system is now set up to work seamlessly with Bibby AI's citation management. Feel free to modify it for your needs!

\bibliographystyle{plain}
\bibliography{cite}

\end{document}\documentclass{article}
\usepackage[utf8]{inputenc}
\usepackage{booktabs}
\usepackage{amsmath}
\usepackage{natbib}

\title{My Document}
\author{Your Name}
\date{\today}

\begin{document}

\maketitle

\section{Introduction}
This is a sample document demonstrating various \LaTeX{} features including mathematical equations, tables, and formatting. You can cite papers using the citation search feature, and they will be automatically added to your \texttt{cite.bib} file.

\section{Mathematical Expressions}
Here are some example equations:

\subsection{Quadratic Formula}
The quadratic formula is given by:
\[x = \frac{-b \pm \sqrt{b^2 - 4ac}}{2a}\]

\subsection{Matrix Example}
A sample matrix:
\[
\begin{pmatrix}
a_{11} & a_{12} & a_{13} \\
a_{21} & a_{22} & a_{23} \\
a_{31} & a_{32} & a_{33}
\end{pmatrix}
\]

\section{Tables}
Below is a sample table showing data:

\begin{table}[h]
\centering
\begin{tabular}{@{}lrrr@{}}
\toprule
\textbf{Category} & \textbf{Value 1} & \textbf{Value 2} & \textbf{Total} \\
\midrule
Item A & 10.5 & 20.3 & 30.8 \\
Item B & 15.2 & 25.1 & 40.3 \\
Item C & 22.1 & 12.7 & 34.8 \\
\bottomrule
\end{tabular}
\caption{Sample Data Table}
\label{tab:sample}
\end{table}

\section{Lists}
\subsection{Bullet Points}
\begin{itemize}
  \item First item
  \item Second item
  \item Third item with sub-items:
    \begin{itemize}
      \item Sub-item 1
      \item Sub-item 2
    \end{itemize}
\end{itemize}

\subsection{Numbered List}
\begin{enumerate}
  \item First step
  \item Second step
  \item Third step
\end{enumerate}

\section{Citations and References}
When you use the citation search feature, references will be automatically managed. For example, you might cite a paper like this: \cite{example2024}. All citations will be automatically added to your \texttt{cite.bib} file.

\section{Conclusion}
This template demonstrates basic \LaTeX{} features including mathematical formulas, tables, and lists. The bibliography system is now set up to work seamlessly with Bibby AI's citation management. Feel free to modify it for your needs!

\bibliographystyle{plain}
\bibliography{cite}

\end{document}\documentclass{article}
\usepackage[utf8]{inputenc}
\usepackage{booktabs}
\usepackage{amsmath}
\usepackage{natbib}

\title{My Document}
\author{Your Name}
\date{\today}

\begin{document}

\maketitle

\section{Introduction}
This is a sample document demonstrating various \LaTeX{} features including mathematical equations, tables, and formatting. You can cite papers using the citation search feature, and they will be automatically added to your \texttt{cite.bib} file.

\section{Mathematical Expressions}
Here are some example equations:

\subsection{Quadratic Formula}
The quadratic formula is given by:
\[x = \frac{-b \pm \sqrt{b^2 - 4ac}}{2a}\]

\subsection{Matrix Example}
A sample matrix:
\[
\begin{pmatrix}
a_{11} & a_{12} & a_{13} \\
a_{21} & a_{22} & a_{23} \\
a_{31} & a_{32} & a_{33}
\end{pmatrix}
\]

\section{Tables}
Below is a sample table showing data:

\begin{table}[h]
\centering
\begin{tabular}{@{}lrrr@{}}
\toprule
\textbf{Category} & \textbf{Value 1} & \textbf{Value 2} & \textbf{Total} \\
\midrule
Item A & 10.5 & 20.3 & 30.8 \\
Item B & 15.2 & 25.1 & 40.3 \\
Item C & 22.1 & 12.7 & 34.8 \\
\bottomrule
\end{tabular}
\caption{Sample Data Table}
\label{tab:sample}
\end{table}

\section{Lists}
\subsection{Bullet Points}
\begin{itemize}
  \item First item
  \item Second item
  \item Third item with sub-items:
    \begin{itemize}
      \item Sub-item 1
      \item Sub-item 2
    \end{itemize}
\end{itemize}

\subsection{Numbered List}
\begin{enumerate}
  \item First step
  \item Second step
  \item Third step
\end{enumerate}

\section{Citations and References}
When you use the citation search feature, references will be automatically managed. For example, you might cite a paper like this: \cite{example2024}. All citations will be automatically added to your \texttt{cite.bib} file.

\section{Conclusion}
This template demonstrates basic \LaTeX{} features including mathematical formulas, tables, and lists. The bibliography system is now set up to work seamlessly with Bibby AI's citation management. Feel free to modify it for your needs!

\bibliographystyle{plain}
\bibliography{cite}

\end{document}\documentclass{article}
\usepackage[utf8]{inputenc}
\usepackage{booktabs}
\usepackage{amsmath}
\usepackage{natbib}

\title{My Document}
\author{Your Name}
\date{\today}

\begin{document}

\maketitle

\section{Introduction}
This is a sample document demonstrating various \LaTeX{} features including mathematical equations, tables, and formatting. You can cite papers using the citation search feature, and they will be automatically added to your \texttt{cite.bib} file.

\section{Mathematical Expressions}
Here are some example equations:

\subsection{Quadratic Formula}
The quadratic formula is given by:
\[x = \frac{-b \pm \sqrt{b^2 - 4ac}}{2a}\]

\subsection{Matrix Example}
A sample matrix:
\[
\begin{pmatrix}
a_{11} & a_{12} & a_{13} \\
a_{21} & a_{22} & a_{23} \\
a_{31} & a_{32} & a_{33}
\end{pmatrix}
\]

\section{Tables}
Below is a sample table showing data:

\begin{table}[h]
\centering
\begin{tabular}{@{}lrrr@{}}
\toprule
\textbf{Category} & \textbf{Value 1} & \textbf{Value 2} & \textbf{Total} \\
\midrule
Item A & 10.5 & 20.3 & 30.8 \\
Item B & 15.2 & 25.1 & 40.3 \\
Item C & 22.1 & 12.7 & 34.8 \\
\bottomrule
\end{tabular}
\caption{Sample Data Table}
\label{tab:sample}
\end{table}

\section{Lists}
\subsection{Bullet Points}
\begin{itemize}
  \item First item
  \item Second item
  \item Third item with sub-items:
    \begin{itemize}
      \item Sub-item 1
      \item Sub-item 2
    \end{itemize}
\end{itemize}

\subsection{Numbered List}
\begin{enumerate}
  \item First step
  \item Second step
  \item Third step
\end{enumerate}

\section{Citations and References}
When you use the citation search feature, references will be automatically managed. For example, you might cite a paper like this: \cite{example2024}. All citations will be automatically added to your \texttt{cite.bib} file.

\section{Conclusion}
This template demonstrates basic \LaTeX{} features including mathematical formulas, tables, and lists. The bibliography system is now set up to work seamlessly with Bibby AI's citation management. Feel free to modify it for your needs!

\bibliographystyle{plain}
\bibliography{cite}

\end{document}\documentclass{article}
\usepackage[utf8]{inputenc}
\usepackage{booktabs}
\usepackage{amsmath}
\usepackage{natbib}

\title{My Document}
\author{Your Name}
\date{\today}

\begin{document}

\maketitle

\section{Introduction}
This is a sample document demonstrating various \LaTeX{} features including mathematical equations, tables, and formatting. You can cite papers using the citation search feature, and they will be automatically added to your \texttt{cite.bib} file.

\section{Mathematical Expressions}
Here are some example equations:

\subsection{Quadratic Formula}
The quadratic formula is given by:
\[x = \frac{-b \pm \sqrt{b^2 - 4ac}}{2a}\]

\subsection{Matrix Example}
A sample matrix:
\[
\begin{pmatrix}
a_{11} & a_{12} & a_{13} \\
a_{21} & a_{22} & a_{23} \\
a_{31} & a_{32} & a_{33}
\end{pmatrix}
\]

\section{Tables}
Below is a sample table showing data:

\begin{table}[h]
\centering
\begin{tabular}{@{}lrrr@{}}
\toprule
\textbf{Category} & \textbf{Value 1} & \textbf{Value 2} & \textbf{Total} \\
\midrule
Item A & 10.5 & 20.3 & 30.8 \\
Item B & 15.2 & 25.1 & 40.3 \\
Item C & 22.1 & 12.7 & 34.8 \\
\bottomrule
\end{tabular}
\caption{Sample Data Table}
\label{tab:sample}
\end{table}

\section{Lists}
\subsection{Bullet Points}
\begin{itemize}
  \item First item
  \item Second item
  \item Third item with sub-items:
    \begin{itemize}
      \item Sub-item 1
      \item Sub-item 2
    \end{itemize}
\end{itemize}

\subsection{Numbered List}
\begin{enumerate}
  \item First step
  \item Second step
  \item Third step
\end{enumerate}

\section{Citations and References}
When you use the citation search feature, references will be automatically managed. For example, you might cite a paper like this: \cite{example2024}. All citations will be automatically added to your \texttt{cite.bib} file.

\section{Conclusion}
This template demonstrates basic \LaTeX{} features including mathematical formulas, tables, and lists. The bibliography system is now set up to work seamlessly with Bibby AI's citation management. Feel free to modify it for your needs!

\bibliographystyle{plain}
\bibliography{cite}

\end{document}\documentclass{article}
\usepackage[utf8]{inputenc}
\usepackage{booktabs}
\usepackage{amsmath}
\usepackage{natbib}

\title{My Document}
\author{Your Name}
\date{\today}

\begin{document}

\maketitle

\section{Introduction}
This is a sample document demonstrating various \LaTeX{} features including mathematical equations, tables, and formatting. You can cite papers using the citation search feature, and they will be automatically added to your \texttt{cite.bib} file.

\section{Mathematical Expressions}
Here are some example equations:

\subsection{Quadratic Formula}
The quadratic formula is given by:
\[x = \frac{-b \pm \sqrt{b^2 - 4ac}}{2a}\]

\subsection{Matrix Example}
A sample matrix:
\[
\begin{pmatrix}
a_{11} & a_{12} & a_{13} \\
a_{21} & a_{22} & a_{23} \\
a_{31} & a_{32} & a_{33}
\end{pmatrix}
\]

\section{Tables}
Below is a sample table showing data:

\begin{table}[h]
\centering
\begin{tabular}{@{}lrrr@{}}
\toprule
\textbf{Category} & \textbf{Value 1} & \textbf{Value 2} & \textbf{Total} \\
\midrule
Item A & 10.5 & 20.3 & 30.8 \\
Item B & 15.2 & 25.1 & 40.3 \\
Item C & 22.1 & 12.7 & 34.8 \\
\bottomrule
\end{tabular}
\caption{Sample Data Table}
\label{tab:sample}
\end{table}

\section{Lists}
\subsection{Bullet Points}
\begin{itemize}
  \item First item
  \item Second item
  \item Third item with sub-items:
    \begin{itemize}
      \item Sub-item 1
      \item Sub-item 2
    \end{itemize}
\end{itemize}

\subsection{Numbered List}
\begin{enumerate}
  \item First step
  \item Second step
  \item Third step
\end{enumerate}

\section{Citations and References}
When you use the citation search feature, references will be automatically managed. For example, you might cite a paper like this: \cite{example2024}. All citations will be automatically added to your \texttt{cite.bib} file.

\section{Conclusion}
This template demonstrates basic \LaTeX{} features including mathematical formulas, tables, and lists. The bibliography system is now set up to work seamlessly with Bibby AI's citation management. Feel free to modify it for your needs!

\bibliographystyle{plain}
\bibliography{cite}

\end{document}\documentclass{article}
\usepackage[utf8]{inputenc}
\usepackage{booktabs}
\usepackage{amsmath}
\usepackage{natbib}

\title{My Document}
\author{Your Name}
\date{\today}

\begin{document}

\maketitle

\section{Introduction}
This is a sample document demonstrating various \LaTeX{} features including mathematical equations, tables, and formatting. You can cite papers using the citation search feature, and they will be automatically added to your \texttt{cite.bib} file.

\section{Mathematical Expressions}
Here are some example equations:

\subsection{Quadratic Formula}
The quadratic formula is given by:
\[x = \frac{-b \pm \sqrt{b^2 - 4ac}}{2a}\]

\subsection{Matrix Example}
A sample matrix:
\[
\begin{pmatrix}
a_{11} & a_{12} & a_{13} \\
a_{21} & a_{22} & a_{23} \\
a_{31} & a_{32} & a_{33}
\end{pmatrix}
\]

\section{Tables}
Below is a sample table showing data:

\begin{table}[h]
\centering
\begin{tabular}{@{}lrrr@{}}
\toprule
\textbf{Category} & \textbf{Value 1} & \textbf{Value 2} & \textbf{Total} \\
\midrule
Item A & 10.5 & 20.3 & 30.8 \\
Item B & 15.2 & 25.1 & 40.3 \\
Item C & 22.1 & 12.7 & 34.8 \\
\bottomrule
\end{tabular}
\caption{Sample Data Table}
\label{tab:sample}
\end{table}

\section{Lists}
\subsection{Bullet Points}
\begin{itemize}
  \item First item
  \item Second item
  \item Third item with sub-items:
    \begin{itemize}
      \item Sub-item 1
      \item Sub-item 2
    \end{itemize}
\end{itemize}

\subsection{Numbered List}
\begin{enumerate}
  \item First step
  \item Second step
  \item Third step
\end{enumerate}

\section{Citations and References}
When you use the citation search feature, references will be automatically managed. For example, you might cite a paper like this: \cite{example2024}. All citations will be automatically added to your \texttt{cite.bib} file.

\section{Conclusion}
This template demonstrates basic \LaTeX{} features including mathematical formulas, tables, and lists. The bibliography system is now set up to work seamlessly with Bibby AI's citation management. Feel free to modify it for your needs!

\bibliographystyle{plain}
\bibliography{cite}

\end{document}\documentclass{article}
\usepackage[utf8]{inputenc}
\usepackage{booktabs}
\usepackage{amsmath}
\usepackage{natbib}

\title{My Document}
\author{Your Name}
\date{\today}

\begin{document}

\maketitle

\section{Introduction}
This is a sample document demonstrating various \LaTeX{} features including mathematical equations, tables, and formatting. You can cite papers using the citation search feature, and they will be automatically added to your \texttt{cite.bib} file.

\section{Mathematical Expressions}
Here are some example equations:

\subsection{Quadratic Formula}
The quadratic formula is given by:
\[x = \frac{-b \pm \sqrt{b^2 - 4ac}}{2a}\]

\subsection{Matrix Example}
A sample matrix:
\[
\begin{pmatrix}
a_{11} & a_{12} & a_{13} \\
a_{21} & a_{22} & a_{23} \\
a_{31} & a_{32} & a_{33}
\end{pmatrix}
\]

\section{Tables}
Below is a sample table showing data:

\begin{table}[h]
\centering
\begin{tabular}{@{}lrrr@{}}
\toprule
\textbf{Category} & \textbf{Value 1} & \textbf{Value 2} & \textbf{Total} \\
\midrule
Item A & 10.5 & 20.3 & 30.8 \\
Item B & 15.2 & 25.1 & 40.3 \\
Item C & 22.1 & 12.7 & 34.8 \\
\bottomrule
\end{tabular}
\caption{Sample Data Table}
\label{tab:sample}
\end{table}

\section{Lists}
\subsection{Bullet Points}
\begin{itemize}
  \item First item
  \item Second item
  \item Third item with sub-items:
    \begin{itemize}
      \item Sub-item 1
      \item Sub-item 2
    \end{itemize}
\end{itemize}

\subsection{Numbered List}
\begin{enumerate}
  \item First step
  \item Second step
  \item Third step
\end{enumerate}

\section{Citations and References}
When you use the citation search feature, references will be automatically managed. For example, you might cite a paper like this: \cite{example2024}. All citations will be automatically added to your \texttt{cite.bib} file.

\section{Conclusion}
This template demonstrates basic \LaTeX{} features including mathematical formulas, tables, and lists. The bibliography system is now set up to work seamlessly with Bibby AI's citation management. Feel free to modify it for your needs!

\bibliographystyle{plain}
\bibliography{cite}

\end{document}\documentclass{article}
\usepackage[utf8]{inputenc}
\usepackage{booktabs}
\usepackage{amsmath}
\usepackage{natbib}

\title{My Document}
\author{Your Name}
\date{\today}

\begin{document}

\maketitle

\section{Introduction}
This is a sample document demonstrating various \LaTeX{} features including mathematical equations, tables, and formatting. You can cite papers using the citation search feature, and they will be automatically added to your \texttt{cite.bib} file.

\section{Mathematical Expressions}
Here are some example equations:

\subsection{Quadratic Formula}
The quadratic formula is given by:
\[x = \frac{-b \pm \sqrt{b^2 - 4ac}}{2a}\]

\subsection{Matrix Example}
A sample matrix:
\[
\begin{pmatrix}
a_{11} & a_{12} & a_{13} \\
a_{21} & a_{22} & a_{23} \\
a_{31} & a_{32} & a_{33}
\end{pmatrix}
\]

\section{Tables}
Below is a sample table showing data:

\begin{table}[h]
\centering
\begin{tabular}{@{}lrrr@{}}
\toprule
\textbf{Category} & \textbf{Value 1} & \textbf{Value 2} & \textbf{Total} \\
\midrule
Item A & 10.5 & 20.3 & 30.8 \\
Item B & 15.2 & 25.1 & 40.3 \\
Item C & 22.1 & 12.7 & 34.8 \\
\bottomrule
\end{tabular}
\caption{Sample Data Table}
\label{tab:sample}
\end{table}

\section{Lists}
\subsection{Bullet Points}
\begin{itemize}
  \item First item
  \item Second item
  \item Third item with sub-items:
    \begin{itemize}
      \item Sub-item 1
      \item Sub-item 2
    \end{itemize}
\end{itemize}

\subsection{Numbered List}
\begin{enumerate}
  \item First step
  \item Second step
  \item Third step
\end{enumerate}

\section{Citations and References}
When you use the citation search feature, references will be automatically managed. For example, you might cite a paper like this: \cite{example2024}. All citations will be automatically added to your \texttt{cite.bib} file.

\section{Conclusion}
This template demonstrates basic \LaTeX{} features including mathematical formulas, tables, and lists. The bibliography system is now set up to work seamlessly with Bibby AI's citation management. Feel free to modify it for your needs!

\bibliographystyle{plain}
\bibliography{cite}

\end{document}\documentclass{article}
\usepackage[utf8]{inputenc}
\usepackage{booktabs}
\usepackage{amsmath}
\usepackage{natbib}

\title{My Document}
\author{Your Name}
\date{\today}

\begin{document}

\maketitle

\section{Introduction}
This is a sample document demonstrating various \LaTeX{} features including mathematical equations, tables, and formatting. You can cite papers using the citation search feature, and they will be automatically added to your \texttt{cite.bib} file.

\section{Mathematical Expressions}
Here are some example equations:

\subsection{Quadratic Formula}
The quadratic formula is given by:
\[x = \frac{-b \pm \sqrt{b^2 - 4ac}}{2a}\]

\subsection{Matrix Example}
A sample matrix:
\[
\begin{pmatrix}
a_{11} & a_{12} & a_{13} \\
a_{21} & a_{22} & a_{23} \\
a_{31} & a_{32} & a_{33}
\end{pmatrix}
\]

\section{Tables}
Below is a sample table showing data:

\begin{table}[h]
\centering
\begin{tabular}{@{}lrrr@{}}
\toprule
\textbf{Category} & \textbf{Value 1} & \textbf{Value 2} & \textbf{Total} \\
\midrule
Item A & 10.5 & 20.3 & 30.8 \\
Item B & 15.2 & 25.1 & 40.3 \\
Item C & 22.1 & 12.7 & 34.8 \\
\bottomrule
\end{tabular}
\caption{Sample Data Table}
\label{tab:sample}
\end{table}

\section{Lists}
\subsection{Bullet Points}
\begin{itemize}
  \item First item
  \item Second item
  \item Third item with sub-items:
    \begin{itemize}
      \item Sub-item 1
      \item Sub-item 2
    \end{itemize}
\end{itemize}

\subsection{Numbered List}
\begin{enumerate}
  \item First step
  \item Second step
  \item Third step
\end{enumerate}

\section{Citations and References}
When you use the citation search feature, references will be automatically managed. For example, you might cite a paper like this: \cite{example2024}. All citations will be automatically added to your \texttt{cite.bib} file.

\section{Conclusion}
This template demonstrates basic \LaTeX{} features including mathematical formulas, tables, and lists. The bibliography system is now set up to work seamlessly with Bibby AI's citation management. Feel free to modify it for your needs!

\bibliographystyle{plain}
\bibliography{cite}

\end{document}\documentclass{article}
\usepackage[utf8]{inputenc}
\usepackage{booktabs}
\usepackage{amsmath}
\usepackage{natbib}

\title{My Document}
\author{Your Name}
\date{\today}

\begin{document}

\maketitle

\section{Introduction}
This is a sample document demonstrating various \LaTeX{} features including mathematical equations, tables, and formatting. You can cite papers using the citation search feature, and they will be automatically added to your \texttt{cite.bib} file.

\section{Mathematical Expressions}
Here are some example equations:

\subsection{Quadratic Formula}
The quadratic formula is given by:
\[x = \frac{-b \pm \sqrt{b^2 - 4ac}}{2a}\]

\subsection{Matrix Example}
A sample matrix:
\[
\begin{pmatrix}
a_{11} & a_{12} & a_{13} \\
a_{21} & a_{22} & a_{23} \\
a_{31} & a_{32} & a_{33}
\end{pmatrix}
\]

\section{Tables}
Below is a sample table showing data:

\begin{table}[h]
\centering
\begin{tabular}{@{}lrrr@{}}
\toprule
\textbf{Category} & \textbf{Value 1} & \textbf{Value 2} & \textbf{Total} \\
\midrule
Item A & 10.5 & 20.3 & 30.8 \\
Item B & 15.2 & 25.1 & 40.3 \\
Item C & 22.1 & 12.7 & 34.8 \\
\bottomrule
\end{tabular}
\caption{Sample Data Table}
\label{tab:sample}
\end{table}

\section{Lists}
\subsection{Bullet Points}
\begin{itemize}
  \item First item
  \item Second item
  \item Third item with sub-items:
    \begin{itemize}
      \item Sub-item 1
      \item Sub-item 2
    \end{itemize}
\end{itemize}

\subsection{Numbered List}
\begin{enumerate}
  \item First step
  \item Second step
  \item Third step
\end{enumerate}

\section{Citations and References}
When you use the citation search feature, references will be automatically managed. For example, you might cite a paper like this: \cite{example2024}. All citations will be automatically added to your \texttt{cite.bib} file.

\section{Conclusion}
This template demonstrates basic \LaTeX{} features including mathematical formulas, tables, and lists. The bibliography system is now set up to work seamlessly with Bibby AI's citation management. Feel free to modify it for your needs!

\bibliographystyle{plain}
\bibliography{cite}

\end{document}\documentclass{article}
\usepackage[utf8]{inputenc}
\usepackage{booktabs}
\usepackage{amsmath}
\usepackage{natbib}

\title{My Document}
\author{Your Name}
\date{\today}

\begin{document}

\maketitle

\section{Introduction}
This is a sample document demonstrating various \LaTeX{} features including mathematical equations, tables, and formatting. You can cite papers using the citation search feature, and they will be automatically added to your \texttt{cite.bib} file.

\section{Mathematical Expressions}
Here are some example equations:

\subsection{Quadratic Formula}
The quadratic formula is given by:
\[x = \frac{-b \pm \sqrt{b^2 - 4ac}}{2a}\]

\subsection{Matrix Example}
A sample matrix:
\[
\begin{pmatrix}
a_{11} & a_{12} & a_{13} \\
a_{21} & a_{22} & a_{23} \\
a_{31} & a_{32} & a_{33}
\end{pmatrix}
\]

\section{Tables}
Below is a sample table showing data:

\begin{table}[h]
\centering
\begin{tabular}{@{}lrrr@{}}
\toprule
\textbf{Category} & \textbf{Value 1} & \textbf{Value 2} & \textbf{Total} \\
\midrule
Item A & 10.5 & 20.3 & 30.8 \\
Item B & 15.2 & 25.1 & 40.3 \\
Item C & 22.1 & 12.7 & 34.8 \\
\bottomrule
\end{tabular}
\caption{Sample Data Table}
\label{tab:sample}
\end{table}

\section{Lists}
\subsection{Bullet Points}
\begin{itemize}
  \item First item
  \item Second item
  \item Third item with sub-items:
    \begin{itemize}
      \item Sub-item 1
      \item Sub-item 2
    \end{itemize}
\end{itemize}

\subsection{Numbered List}
\begin{enumerate}
  \item First step
  \item Second step
  \item Third step
\end{enumerate}

\section{Citations and References}
When you use the citation search feature, references will be automatically managed. For example, you might cite a paper like this: \cite{example2024}. All citations will be automatically added to your \texttt{cite.bib} file.

\section{Conclusion}
This template demonstrates basic \LaTeX{} features including mathematical formulas, tables, and lists. The bibliography system is now set up to work seamlessly with Bibby AI's citation management. Feel free to modify it for your needs!

\bibliographystyle{plain}
\bibliography{cite}

\end{document}\documentclass{article}
\usepackage[utf8]{inputenc}
\usepackage{booktabs}
\usepackage{amsmath}
\usepackage{natbib}

\title{My Document}
\author{Your Name}
\date{\today}

\begin{document}

\maketitle

\section{Introduction}
This is a sample document demonstrating various \LaTeX{} features including mathematical equations, tables, and formatting. You can cite papers using the citation search feature, and they will be automatically added to your \texttt{cite.bib} file.

\section{Mathematical Expressions}
Here are some example equations:

\subsection{Quadratic Formula}
The quadratic formula is given by:
\[x = \frac{-b \pm \sqrt{b^2 - 4ac}}{2a}\]

\subsection{Matrix Example}
A sample matrix:
\[
\begin{pmatrix}
a_{11} & a_{12} & a_{13} \\
a_{21} & a_{22} & a_{23} \\
a_{31} & a_{32} & a_{33}
\end{pmatrix}
\]

\section{Tables}
Below is a sample table showing data:

\begin{table}[h]
\centering
\begin{tabular}{@{}lrrr@{}}
\toprule
\textbf{Category} & \textbf{Value 1} & \textbf{Value 2} & \textbf{Total} \\
\midrule
Item A & 10.5 & 20.3 & 30.8 \\
Item B & 15.2 & 25.1 & 40.3 \\
Item C & 22.1 & 12.7 & 34.8 \\
\bottomrule
\end{tabular}
\caption{Sample Data Table}
\label{tab:sample}
\end{table}

\section{Lists}
\subsection{Bullet Points}
\begin{itemize}
  \item First item
  \item Second item
  \item Third item with sub-items:
    \begin{itemize}
      \item Sub-item 1
      \item Sub-item 2
    \end{itemize}
\end{itemize}

\subsection{Numbered List}
\begin{enumerate}
  \item First step
  \item Second step
  \item Third step
\end{enumerate}

\section{Citations and References}
When you use the citation search feature, references will be automatically managed. For example, you might cite a paper like this: \cite{example2024}. All citations will be automatically added to your \texttt{cite.bib} file.

\section{Conclusion}
This template demonstrates basic \LaTeX{} features including mathematical formulas, tables, and lists. The bibliography system is now set up to work seamlessly with Bibby AI's citation management. Feel free to modify it for your needs!

\bibliographystyle{plain}
\bibliography{cite}

\end{document}\documentclass{article}
\usepackage[utf8]{inputenc}
\usepackage{booktabs}
\usepackage{amsmath}
\usepackage{natbib}

\title{My Document}
\author{Your Name}
\date{\today}

\begin{document}

\maketitle

\section{Introduction}
This is a sample document demonstrating various \LaTeX{} features including mathematical equations, tables, and formatting. You can cite papers using the citation search feature, and they will be automatically added to your \texttt{cite.bib} file.

\section{Mathematical Expressions}
Here are some example equations:

\subsection{Quadratic Formula}
The quadratic formula is given by:
\[x = \frac{-b \pm \sqrt{b^2 - 4ac}}{2a}\]

\subsection{Matrix Example}
A sample matrix:
\[
\begin{pmatrix}
a_{11} & a_{12} & a_{13} \\
a_{21} & a_{22} & a_{23} \\
a_{31} & a_{32} & a_{33}
\end{pmatrix}
\]

\section{Tables}
Below is a sample table showing data:

\begin{table}[h]
\centering
\begin{tabular}{@{}lrrr@{}}
\toprule
\textbf{Category} & \textbf{Value 1} & \textbf{Value 2} & \textbf{Total} \\
\midrule
Item A & 10.5 & 20.3 & 30.8 \\
Item B & 15.2 & 25.1 & 40.3 \\
Item C & 22.1 & 12.7 & 34.8 \\
\bottomrule
\end{tabular}
\caption{Sample Data Table}
\label{tab:sample}
\end{table}

\section{Lists}
\subsection{Bullet Points}
\begin{itemize}
  \item First item
  \item Second item
  \item Third item with sub-items:
    \begin{itemize}
      \item Sub-item 1
      \item Sub-item 2
    \end{itemize}
\end{itemize}

\subsection{Numbered List}
\begin{enumerate}
  \item First step
  \item Second step
  \item Third step
\end{enumerate}

\section{Citations and References}
When you use the citation search feature, references will be automatically managed. For example, you might cite a paper like this: \cite{example2024}. All citations will be automatically added to your \texttt{cite.bib} file.

\section{Conclusion}
This template demonstrates basic \LaTeX{} features including mathematical formulas, tables, and lists. The bibliography system is now set up to work seamlessly with Bibby AI's citation management. Feel free to modify it for your needs!

\bibliographystyle{plain}
\bibliography{cite}

\end{document}\documentclass{article}
\usepackage[utf8]{inputenc}
\usepackage{booktabs}
\usepackage{amsmath}
\usepackage{natbib}

\title{My Document}
\author{Your Name}
\date{\today}

\begin{document}

\maketitle

\section{Introduction}
This is a sample document demonstrating various \LaTeX{} features including mathematical equations, tables, and formatting. You can cite papers using the citation search feature, and they will be automatically added to your \texttt{cite.bib} file.

\section{Mathematical Expressions}
Here are some example equations:

\subsection{Quadratic Formula}
The quadratic formula is given by:
\[x = \frac{-b \pm \sqrt{b^2 - 4ac}}{2a}\]

\subsection{Matrix Example}
A sample matrix:
\[
\begin{pmatrix}
a_{11} & a_{12} & a_{13} \\
a_{21} & a_{22} & a_{23} \\
a_{31} & a_{32} & a_{33}
\end{pmatrix}
\]

\section{Tables}
Below is a sample table showing data:

\begin{table}[h]
\centering
\begin{tabular}{@{}lrrr@{}}
\toprule
\textbf{Category} & \textbf{Value 1} & \textbf{Value 2} & \textbf{Total} \\
\midrule
Item A & 10.5 & 20.3 & 30.8 \\
Item B & 15.2 & 25.1 & 40.3 \\
Item C & 22.1 & 12.7 & 34.8 \\
\bottomrule
\end{tabular}
\caption{Sample Data Table}
\label{tab:sample}
\end{table}

\section{Lists}
\subsection{Bullet Points}
\begin{itemize}
  \item First item
  \item Second item
  \item Third item with sub-items:
    \begin{itemize}
      \item Sub-item 1
      \item Sub-item 2
    \end{itemize}
\end{itemize}

\subsection{Numbered List}
\begin{enumerate}
  \item First step
  \item Second step
  \item Third step
\end{enumerate}

\section{Citations and References}
When you use the citation search feature, references will be automatically managed. For example, you might cite a paper like this: \cite{example2024}. All citations will be automatically added to your \texttt{cite.bib} file.

\section{Conclusion}
This template demonstrates basic \LaTeX{} features including mathematical formulas, tables, and lists. The bibliography system is now set up to work seamlessly with Bibby AI's citation management. Feel free to modify it for your needs!

\bibliographystyle{plain}
\bibliography{cite}

\end{document}\documentclass{article}
\usepackage[utf8]{inputenc}
\usepackage{booktabs}
\usepackage{amsmath}
\usepackage{natbib}

\title{My Document}
\author{Your Name}
\date{\today}

\begin{document}

\maketitle

\section{Introduction}
This is a sample document demonstrating various \LaTeX{} features including mathematical equations, tables, and formatting. You can cite papers using the citation search feature, and they will be automatically added to your \texttt{cite.bib} file.

\section{Mathematical Expressions}
Here are some example equations:

\subsection{Quadratic Formula}
The quadratic formula is given by:
\[x = \frac{-b \pm \sqrt{b^2 - 4ac}}{2a}\]

\subsection{Matrix Example}
A sample matrix:
\[
\begin{pmatrix}
a_{11} & a_{12} & a_{13} \\
a_{21} & a_{22} & a_{23} \\
a_{31} & a_{32} & a_{33}
\end{pmatrix}
\]

\section{Tables}
Below is a sample table showing data:

\begin{table}[h]
\centering
\begin{tabular}{@{}lrrr@{}}
\toprule
\textbf{Category} & \textbf{Value 1} & \textbf{Value 2} & \textbf{Total} \\
\midrule
Item A & 10.5 & 20.3 & 30.8 \\
Item B & 15.2 & 25.1 & 40.3 \\
Item C & 22.1 & 12.7 & 34.8 \\
\bottomrule
\end{tabular}
\caption{Sample Data Table}
\label{tab:sample}
\end{table}

\section{Lists}
\subsection{Bullet Points}
\begin{itemize}
  \item First item
  \item Second item
  \item Third item with sub-items:
    \begin{itemize}
      \item Sub-item 1
      \item Sub-item 2
    \end{itemize}
\end{itemize}

\subsection{Numbered List}
\begin{enumerate}
  \item First step
  \item Second step
  \item Third step
\end{enumerate}

\section{Citations and References}
When you use the citation search feature, references will be automatically managed. For example, you might cite a paper like this: \cite{example2024}. All citations will be automatically added to your \texttt{cite.bib} file.

\section{Conclusion}
This template demonstrates basic \LaTeX{} features including mathematical formulas, tables, and lists. The bibliography system is now set up to work seamlessly with Bibby AI's citation management. Feel free to modify it for your needs!

\bibliographystyle{plain}
\bibliography{cite}

\end{document}\documentclass{article}
\usepackage[utf8]{inputenc}
\usepackage{booktabs}
\usepackage{amsmath}
\usepackage{natbib}

\title{My Document}
\author{Your Name}
\date{\today}

\begin{document}

\maketitle

\section{Introduction}
This is a sample document demonstrating various \LaTeX{} features including mathematical equations, tables, and formatting. You can cite papers using the citation search feature, and they will be automatically added to your \texttt{cite.bib} file.

\section{Mathematical Expressions}
Here are some example equations:

\subsection{Quadratic Formula}
The quadratic formula is given by:
\[x = \frac{-b \pm \sqrt{b^2 - 4ac}}{2a}\]

\subsection{Matrix Example}
A sample matrix:
\[
\begin{pmatrix}
a_{11} & a_{12} & a_{13} \\
a_{21} & a_{22} & a_{23} \\
a_{31} & a_{32} & a_{33}
\end{pmatrix}
\]

\section{Tables}
Below is a sample table showing data:

\begin{table}[h]
\centering
\begin{tabular}{@{}lrrr@{}}
\toprule
\textbf{Category} & \textbf{Value 1} & \textbf{Value 2} & \textbf{Total} \\
\midrule
Item A & 10.5 & 20.3 & 30.8 \\
Item B & 15.2 & 25.1 & 40.3 \\
Item C & 22.1 & 12.7 & 34.8 \\
\bottomrule
\end{tabular}
\caption{Sample Data Table}
\label{tab:sample}
\end{table}

\section{Lists}
\subsection{Bullet Points}
\begin{itemize}
  \item First item
  \item Second item
  \item Third item with sub-items:
    \begin{itemize}
      \item Sub-item 1
      \item Sub-item 2
    \end{itemize}
\end{itemize}

\subsection{Numbered List}
\begin{enumerate}
  \item First step
  \item Second step
  \item Third step
\end{enumerate}

\section{Citations and References}
When you use the citation search feature, references will be automatically managed. For example, you might cite a paper like this: \cite{example2024}. All citations will be automatically added to your \texttt{cite.bib} file.

\section{Conclusion}
This template demonstrates basic \LaTeX{} features including mathematical formulas, tables, and lists. The bibliography system is now set up to work seamlessly with Bibby AI's citation management. Feel free to modify it for your needs!

\bibliographystyle{plain}
\bibliography{cite}

\end{document}\documentclass{article}
\usepackage[utf8]{inputenc}
\usepackage{booktabs}
\usepackage{amsmath}
\usepackage{natbib}

\title{My Document}
\author{Your Name}
\date{\today}

\begin{document}

\maketitle

\section{Introduction}
This is a sample document demonstrating various \LaTeX{} features including mathematical equations, tables, and formatting. You can cite papers using the citation search feature, and they will be automatically added to your \texttt{cite.bib} file.

\section{Mathematical Expressions}
Here are some example equations:

\subsection{Quadratic Formula}
The quadratic formula is given by:
\[x = \frac{-b \pm \sqrt{b^2 - 4ac}}{2a}\]

\subsection{Matrix Example}
A sample matrix:
\[
\begin{pmatrix}
a_{11} & a_{12} & a_{13} \\
a_{21} & a_{22} & a_{23} \\
a_{31} & a_{32} & a_{33}
\end{pmatrix}
\]

\section{Tables}
Below is a sample table showing data:

\begin{table}[h]
\centering
\begin{tabular}{@{}lrrr@{}}
\toprule
\textbf{Category} & \textbf{Value 1} & \textbf{Value 2} & \textbf{Total} \\
\midrule
Item A & 10.5 & 20.3 & 30.8 \\
Item B & 15.2 & 25.1 & 40.3 \\
Item C & 22.1 & 12.7 & 34.8 \\
\bottomrule
\end{tabular}
\caption{Sample Data Table}
\label{tab:sample}
\end{table}

\section{Lists}
\subsection{Bullet Points}
\begin{itemize}
  \item First item
  \item Second item
  \item Third item with sub-items:
    \begin{itemize}
      \item Sub-item 1
      \item Sub-item 2
    \end{itemize}
\end{itemize}

\subsection{Numbered List}
\begin{enumerate}
  \item First step
  \item Second step
  \item Third step
\end{enumerate}

\section{Citations and References}
When you use the citation search feature, references will be automatically managed. For example, you might cite a paper like this: \cite{example2024}. All citations will be automatically added to your \texttt{cite.bib} file.

\section{Conclusion}
This template demonstrates basic \LaTeX{} features including mathematical formulas, tables, and lists. The bibliography system is now set up to work seamlessly with Bibby AI's citation management. Feel free to modify it for your needs!

\bibliographystyle{plain}
\bibliography{cite}

\end{document}\documentclass{article}
\usepackage[utf8]{inputenc}
\usepackage{booktabs}
\usepackage{amsmath}
\usepackage{natbib}

\title{My Document}
\author{Your Name}
\date{\today}

\begin{document}

\maketitle

\section{Introduction}
This is a sample document demonstrating various \LaTeX{} features including mathematical equations, tables, and formatting. You can cite papers using the citation search feature, and they will be automatically added to your \texttt{cite.bib} file.

\section{Mathematical Expressions}
Here are some example equations:

\subsection{Quadratic Formula}
The quadratic formula is given by:
\[x = \frac{-b \pm \sqrt{b^2 - 4ac}}{2a}\]

\subsection{Matrix Example}
A sample matrix:
\[
\begin{pmatrix}
a_{11} & a_{12} & a_{13} \\
a_{21} & a_{22} & a_{23} \\
a_{31} & a_{32} & a_{33}
\end{pmatrix}
\]

\section{Tables}
Below is a sample table showing data:

\begin{table}[h]
\centering
\begin{tabular}{@{}lrrr@{}}
\toprule
\textbf{Category} & \textbf{Value 1} & \textbf{Value 2} & \textbf{Total} \\
\midrule
Item A & 10.5 & 20.3 & 30.8 \\
Item B & 15.2 & 25.1 & 40.3 \\
Item C & 22.1 & 12.7 & 34.8 \\
\bottomrule
\end{tabular}
\caption{Sample Data Table}
\label{tab:sample}
\end{table}

\section{Lists}
\subsection{Bullet Points}
\begin{itemize}
  \item First item
  \item Second item
  \item Third item with sub-items:
    \begin{itemize}
      \item Sub-item 1
      \item Sub-item 2
    \end{itemize}
\end{itemize}

\subsection{Numbered List}
\begin{enumerate}
  \item First step
  \item Second step
  \item Third step
\end{enumerate}

\section{Citations and References}
When you use the citation search feature, references will be automatically managed. For example, you might cite a paper like this: \cite{example2024}. All citations will be automatically added to your \texttt{cite.bib} file.

\section{Conclusion}
This template demonstrates basic \LaTeX{} features including mathematical formulas, tables, and lists. The bibliography system is now set up to work seamlessly with Bibby AI's citation management. Feel free to modify it for your needs!

\bibliographystyle{plain}
\bibliography{cite}

\end{document}\documentclass{article}
\usepackage[utf8]{inputenc}
\usepackage{booktabs}
\usepackage{amsmath}
\usepackage{natbib}

\title{My Document}
\author{Your Name}
\date{\today}

\begin{document}

\maketitle

\section{Introduction}
This is a sample document demonstrating various \LaTeX{} features including mathematical equations, tables, and formatting. You can cite papers using the citation search feature, and they will be automatically added to your \texttt{cite.bib} file.

\section{Mathematical Expressions}
Here are some example equations:

\subsection{Quadratic Formula}
The quadratic formula is given by:
\[x = \frac{-b \pm \sqrt{b^2 - 4ac}}{2a}\]

\subsection{Matrix Example}
A sample matrix:
\[
\begin{pmatrix}
a_{11} & a_{12} & a_{13} \\
a_{21} & a_{22} & a_{23} \\
a_{31} & a_{32} & a_{33}
\end{pmatrix}
\]

\section{Tables}
Below is a sample table showing data:

\begin{table}[h]
\centering
\begin{tabular}{@{}lrrr@{}}
\toprule
\textbf{Category} & \textbf{Value 1} & \textbf{Value 2} & \textbf{Total} \\
\midrule
Item A & 10.5 & 20.3 & 30.8 \\
Item B & 15.2 & 25.1 & 40.3 \\
Item C & 22.1 & 12.7 & 34.8 \\
\bottomrule
\end{tabular}
\caption{Sample Data Table}
\label{tab:sample}
\end{table}

\section{Lists}
\subsection{Bullet Points}
\begin{itemize}
  \item First item
  \item Second item
  \item Third item with sub-items:
    \begin{itemize}
      \item Sub-item 1
      \item Sub-item 2
    \end{itemize}
\end{itemize}

\subsection{Numbered List}
\begin{enumerate}
  \item First step
  \item Second step
  \item Third step
\end{enumerate}

\section{Citations and References}
When you use the citation search feature, references will be automatically managed. For example, you might cite a paper like this: \cite{example2024}. All citations will be automatically added to your \texttt{cite.bib} file.

\section{Conclusion}
This template demonstrates basic \LaTeX{} features including mathematical formulas, tables, and lists. The bibliography system is now set up to work seamlessly with Bibby AI's citation management. Feel free to modify it for your needs!

\bibliographystyle{plain}
\bibliography{cite}

\end{document}\documentclass{article}
\usepackage[utf8]{inputenc}
\usepackage{booktabs}
\usepackage{amsmath}
\usepackage{natbib}

\title{My Document}
\author{Your Name}
\date{\today}

\begin{document}

\maketitle

\section{Introduction}
This is a sample document demonstrating various \LaTeX{} features including mathematical equations, tables, and formatting. You can cite papers using the citation search feature, and they will be automatically added to your \texttt{cite.bib} file.

\section{Mathematical Expressions}
Here are some example equations:

\subsection{Quadratic Formula}
The quadratic formula is given by:
\[x = \frac{-b \pm \sqrt{b^2 - 4ac}}{2a}\]

\subsection{Matrix Example}
A sample matrix:
\[
\begin{pmatrix}
a_{11} & a_{12} & a_{13} \\
a_{21} & a_{22} & a_{23} \\
a_{31} & a_{32} & a_{33}
\end{pmatrix}
\]

\section{Tables}
Below is a sample table showing data:

\begin{table}[h]
\centering
\begin{tabular}{@{}lrrr@{}}
\toprule
\textbf{Category} & \textbf{Value 1} & \textbf{Value 2} & \textbf{Total} \\
\midrule
Item A & 10.5 & 20.3 & 30.8 \\
Item B & 15.2 & 25.1 & 40.3 \\
Item C & 22.1 & 12.7 & 34.8 \\
\bottomrule
\end{tabular}
\caption{Sample Data Table}
\label{tab:sample}
\end{table}

\section{Lists}
\subsection{Bullet Points}
\begin{itemize}
  \item First item
  \item Second item
  \item Third item with sub-items:
    \begin{itemize}
      \item Sub-item 1
      \item Sub-item 2
    \end{itemize}
\end{itemize}

\subsection{Numbered List}
\begin{enumerate}
  \item First step
  \item Second step
  \item Third step
\end{enumerate}

\section{Citations and References}
When you use the citation search feature, references will be automatically managed. For example, you might cite a paper like this: \cite{example2024}. All citations will be automatically added to your \texttt{cite.bib} file.

\section{Conclusion}
This template demonstrates basic \LaTeX{} features including mathematical formulas, tables, and lists. The bibliography system is now set up to work seamlessly with Bibby AI's citation management. Feel free to modify it for your needs!

\bibliographystyle{plain}
\bibliography{cite}

\end{document}\documentclass{article}
\usepackage[utf8]{inputenc}
\usepackage{booktabs}
\usepackage{amsmath}
\usepackage{natbib}

\title{My Document}
\author{Your Name}
\date{\today}

\begin{document}

\maketitle

\section{Introduction}
This is a sample document demonstrating various \LaTeX{} features including mathematical equations, tables, and formatting. You can cite papers using the citation search feature, and they will be automatically added to your \texttt{cite.bib} file.

\section{Mathematical Expressions}
Here are some example equations:

\subsection{Quadratic Formula}
The quadratic formula is given by:
\[x = \frac{-b \pm \sqrt{b^2 - 4ac}}{2a}\]

\subsection{Matrix Example}
A sample matrix:
\[
\begin{pmatrix}
a_{11} & a_{12} & a_{13} \\
a_{21} & a_{22} & a_{23} \\
a_{31} & a_{32} & a_{33}
\end{pmatrix}
\]

\section{Tables}
Below is a sample table showing data:

\begin{table}[h]
\centering
\begin{tabular}{@{}lrrr@{}}
\toprule
\textbf{Category} & \textbf{Value 1} & \textbf{Value 2} & \textbf{Total} \\
\midrule
Item A & 10.5 & 20.3 & 30.8 \\
Item B & 15.2 & 25.1 & 40.3 \\
Item C & 22.1 & 12.7 & 34.8 \\
\bottomrule
\end{tabular}
\caption{Sample Data Table}
\label{tab:sample}
\end{table}

\section{Lists}
\subsection{Bullet Points}
\begin{itemize}
  \item First item
  \item Second item
  \item Third item with sub-items:
    \begin{itemize}
      \item Sub-item 1
      \item Sub-item 2
    \end{itemize}
\end{itemize}

\subsection{Numbered List}
\begin{enumerate}
  \item First step
  \item Second step
  \item Third step
\end{enumerate}

\section{Citations and References}
When you use the citation search feature, references will be automatically managed. For example, you might cite a paper like this: \cite{example2024}. All citations will be automatically added to your \texttt{cite.bib} file.

\section{Conclusion}
This template demonstrates basic \LaTeX{} features including mathematical formulas, tables, and lists. The bibliography system is now set up to work seamlessly with Bibby AI's citation management. Feel free to modify it for your needs!

\bibliographystyle{plain}
\bibliography{cite}

\end{document}\documentclass{article}
\usepackage[utf8]{inputenc}
\usepackage{booktabs}
\usepackage{amsmath}
\usepackage{natbib}

\title{My Document}
\author{Your Name}
\date{\today}

\begin{document}

\maketitle

\section{Introduction}
This is a sample document demonstrating various \LaTeX{} features including mathematical equations, tables, and formatting. You can cite papers using the citation search feature, and they will be automatically added to your \texttt{cite.bib} file.

\section{Mathematical Expressions}
Here are some example equations:

\subsection{Quadratic Formula}
The quadratic formula is given by:
\[x = \frac{-b \pm \sqrt{b^2 - 4ac}}{2a}\]

\subsection{Matrix Example}
A sample matrix:
\[
\begin{pmatrix}
a_{11} & a_{12} & a_{13} \\
a_{21} & a_{22} & a_{23} \\
a_{31} & a_{32} & a_{33}
\end{pmatrix}
\]

\section{Tables}
Below is a sample table showing data:

\begin{table}[h]
\centering
\begin{tabular}{@{}lrrr@{}}
\toprule
\textbf{Category} & \textbf{Value 1} & \textbf{Value 2} & \textbf{Total} \\
\midrule
Item A & 10.5 & 20.3 & 30.8 \\
Item B & 15.2 & 25.1 & 40.3 \\
Item C & 22.1 & 12.7 & 34.8 \\
\bottomrule
\end{tabular}
\caption{Sample Data Table}
\label{tab:sample}
\end{table}

\section{Lists}
\subsection{Bullet Points}
\begin{itemize}
  \item First item
  \item Second item
  \item Third item with sub-items:
    \begin{itemize}
      \item Sub-item 1
      \item Sub-item 2
    \end{itemize}
\end{itemize}

\subsection{Numbered List}
\begin{enumerate}
  \item First step
  \item Second step
  \item Third step
\end{enumerate}

\section{Citations and References}
When you use the citation search feature, references will be automatically managed. For example, you might cite a paper like this: \cite{example2024}. All citations will be automatically added to your \texttt{cite.bib} file.

\section{Conclusion}
This template demonstrates basic \LaTeX{} features including mathematical formulas, tables, and lists. The bibliography system is now set up to work seamlessly with Bibby AI's citation management. Feel free to modify it for your needs!

\bibliographystyle{plain}
\bibliography{cite}

\end{document}\documentclass{article}
\usepackage[utf8]{inputenc}
\usepackage{booktabs}
\usepackage{amsmath}
\usepackage{natbib}

\title{My Document}
\author{Your Name}
\date{\today}

\begin{document}

\maketitle

\section{Introduction}
This is a sample document demonstrating various \LaTeX{} features including mathematical equations, tables, and formatting. You can cite papers using the citation search feature, and they will be automatically added to your \texttt{cite.bib} file.

\section{Mathematical Expressions}
Here are some example equations:

\subsection{Quadratic Formula}
The quadratic formula is given by:
\[x = \frac{-b \pm \sqrt{b^2 - 4ac}}{2a}\]

\subsection{Matrix Example}
A sample matrix:
\[
\begin{pmatrix}
a_{11} & a_{12} & a_{13} \\
a_{21} & a_{22} & a_{23} \\
a_{31} & a_{32} & a_{33}
\end{pmatrix}
\]

\section{Tables}
Below is a sample table showing data:

\begin{table}[h]
\centering
\begin{tabular}{@{}lrrr@{}}
\toprule
\textbf{Category} & \textbf{Value 1} & \textbf{Value 2} & \textbf{Total} \\
\midrule
Item A & 10.5 & 20.3 & 30.8 \\
Item B & 15.2 & 25.1 & 40.3 \\
Item C & 22.1 & 12.7 & 34.8 \\
\bottomrule
\end{tabular}
\caption{Sample Data Table}
\label{tab:sample}
\end{table}

\section{Lists}
\subsection{Bullet Points}
\begin{itemize}
  \item First item
  \item Second item
  \item Third item with sub-items:
    \begin{itemize}
      \item Sub-item 1
      \item Sub-item 2
    \end{itemize}
\end{itemize}

\subsection{Numbered List}
\begin{enumerate}
  \item First step
  \item Second step
  \item Third step
\end{enumerate}

\section{Citations and References}
When you use the citation search feature, references will be automatically managed. For example, you might cite a paper like this: \cite{example2024}. All citations will be automatically added to your \texttt{cite.bib} file.

\section{Conclusion}
This template demonstrates basic \LaTeX{} features including mathematical formulas, tables, and lists. The bibliography system is now set up to work seamlessly with Bibby AI's citation management. Feel free to modify it for your needs!

\bibliographystyle{plain}
\bibliography{cite}

\end{document}\documentclass{article}
\usepackage[utf8]{inputenc}
\usepackage{booktabs}
\usepackage{amsmath}
\usepackage{natbib}

\title{My Document}
\author{Your Name}
\date{\today}

\begin{document}

\maketitle

\section{Introduction}
This is a sample document demonstrating various \LaTeX{} features including mathematical equations, tables, and formatting. You can cite papers using the citation search feature, and they will be automatically added to your \texttt{cite.bib} file.

\section{Mathematical Expressions}
Here are some example equations:

\subsection{Quadratic Formula}
The quadratic formula is given by:
\[x = \frac{-b \pm \sqrt{b^2 - 4ac}}{2a}\]

\subsection{Matrix Example}
A sample matrix:
\[
\begin{pmatrix}
a_{11} & a_{12} & a_{13} \\
a_{21} & a_{22} & a_{23} \\
a_{31} & a_{32} & a_{33}
\end{pmatrix}
\]

\section{Tables}
Below is a sample table showing data:

\begin{table}[h]
\centering
\begin{tabular}{@{}lrrr@{}}
\toprule
\textbf{Category} & \textbf{Value 1} & \textbf{Value 2} & \textbf{Total} \\
\midrule
Item A & 10.5 & 20.3 & 30.8 \\
Item B & 15.2 & 25.1 & 40.3 \\
Item C & 22.1 & 12.7 & 34.8 \\
\bottomrule
\end{tabular}
\caption{Sample Data Table}
\label{tab:sample}
\end{table}

\section{Lists}
\subsection{Bullet Points}
\begin{itemize}
  \item First item
  \item Second item
  \item Third item with sub-items:
    \begin{itemize}
      \item Sub-item 1
      \item Sub-item 2
    \end{itemize}
\end{itemize}

\subsection{Numbered List}
\begin{enumerate}
  \item First step
  \item Second step
  \item Third step
\end{enumerate}

\section{Citations and References}
When you use the citation search feature, references will be automatically managed. For example, you might cite a paper like this: \cite{example2024}. All citations will be automatically added to your \texttt{cite.bib} file.

\section{Conclusion}
This template demonstrates basic \LaTeX{} features including mathematical formulas, tables, and lists. The bibliography system is now set up to work seamlessly with Bibby AI's citation management. Feel free to modify it for your needs!

\bibliographystyle{plain}
\bibliography{cite}

\end{document}\documentclass{article}
\usepackage[utf8]{inputenc}
\usepackage{booktabs}
\usepackage{amsmath}
\usepackage{natbib}

\title{My Document}
\author{Your Name}
\date{\today}

\begin{document}

\maketitle

\section{Introduction}
This is a sample document demonstrating various \LaTeX{} features including mathematical equations, tables, and formatting. You can cite papers using the citation search feature, and they will be automatically added to your \texttt{cite.bib} file.

\section{Mathematical Expressions}
Here are some example equations:

\subsection{Quadratic Formula}
The quadratic formula is given by:
\[x = \frac{-b \pm \sqrt{b^2 - 4ac}}{2a}\]

\subsection{Matrix Example}
A sample matrix:
\[
\begin{pmatrix}
a_{11} & a_{12} & a_{13} \\
a_{21} & a_{22} & a_{23} \\
a_{31} & a_{32} & a_{33}
\end{pmatrix}
\]

\section{Tables}
Below is a sample table showing data:

\begin{table}[h]
\centering
\begin{tabular}{@{}lrrr@{}}
\toprule
\textbf{Category} & \textbf{Value 1} & \textbf{Value 2} & \textbf{Total} \\
\midrule
Item A & 10.5 & 20.3 & 30.8 \\
Item B & 15.2 & 25.1 & 40.3 \\
Item C & 22.1 & 12.7 & 34.8 \\
\bottomrule
\end{tabular}
\caption{Sample Data Table}
\label{tab:sample}
\end{table}

\section{Lists}
\subsection{Bullet Points}
\begin{itemize}
  \item First item
  \item Second item
  \item Third item with sub-items:
    \begin{itemize}
      \item Sub-item 1
      \item Sub-item 2
    \end{itemize}
\end{itemize}

\subsection{Numbered List}
\begin{enumerate}
  \item First step
  \item Second step
  \item Third step
\end{enumerate}

\section{Citations and References}
When you use the citation search feature, references will be automatically managed. For example, you might cite a paper like this: \cite{example2024}. All citations will be automatically added to your \texttt{cite.bib} file.

\section{Conclusion}
This template demonstrates basic \LaTeX{} features including mathematical formulas, tables, and lists. The bibliography system is now set up to work seamlessly with Bibby AI's citation management. Feel free to modify it for your needs!

\bibliographystyle{plain}
\bibliography{cite}

\end{document}\documentclass{article}
\usepackage[utf8]{inputenc}
\usepackage{booktabs}
\usepackage{amsmath}
\usepackage{natbib}

\title{My Document}
\author{Your Name}
\date{\today}

\begin{document}

\maketitle

\section{Introduction}
This is a sample document demonstrating various \LaTeX{} features including mathematical equations, tables, and formatting. You can cite papers using the citation search feature, and they will be automatically added to your \texttt{cite.bib} file.

\section{Mathematical Expressions}
Here are some example equations:

\subsection{Quadratic Formula}
The quadratic formula is given by:
\[x = \frac{-b \pm \sqrt{b^2 - 4ac}}{2a}\]

\subsection{Matrix Example}
A sample matrix:
\[
\begin{pmatrix}
a_{11} & a_{12} & a_{13} \\
a_{21} & a_{22} & a_{23} \\
a_{31} & a_{32} & a_{33}
\end{pmatrix}
\]

\section{Tables}
Below is a sample table showing data:

\begin{table}[h]
\centering
\begin{tabular}{@{}lrrr@{}}
\toprule
\textbf{Category} & \textbf{Value 1} & \textbf{Value 2} & \textbf{Total} \\
\midrule
Item A & 10.5 & 20.3 & 30.8 \\
Item B & 15.2 & 25.1 & 40.3 \\
Item C & 22.1 & 12.7 & 34.8 \\
\bottomrule
\end{tabular}
\caption{Sample Data Table}
\label{tab:sample}
\end{table}

\section{Lists}
\subsection{Bullet Points}
\begin{itemize}
  \item First item
  \item Second item
  \item Third item with sub-items:
    \begin{itemize}
      \item Sub-item 1
      \item Sub-item 2
    \end{itemize}
\end{itemize}

\subsection{Numbered List}
\begin{enumerate}
  \item First step
  \item Second step
  \item Third step
\end{enumerate}

\section{Citations and References}
When you use the citation search feature, references will be automatically managed. For example, you might cite a paper like this: \cite{example2024}. All citations will be automatically added to your \texttt{cite.bib} file.

\section{Conclusion}
This template demonstrates basic \LaTeX{} features including mathematical formulas, tables, and lists. The bibliography system is now set up to work seamlessly with Bibby AI's citation management. Feel free to modify it for your needs!

\bibliographystyle{plain}
\bibliography{cite}

\end{document}\documentclass{article}
\usepackage[utf8]{inputenc}
\usepackage{booktabs}
\usepackage{amsmath}
\usepackage{natbib}

\title{My Document}
\author{Your Name}
\date{\today}

\begin{document}

\maketitle

\section{Introduction}
This is a sample document demonstrating various \LaTeX{} features including mathematical equations, tables, and formatting. You can cite papers using the citation search feature, and they will be automatically added to your \texttt{cite.bib} file.

\section{Mathematical Expressions}
Here are some example equations:

\subsection{Quadratic Formula}
The quadratic formula is given by:
\[x = \frac{-b \pm \sqrt{b^2 - 4ac}}{2a}\]

\subsection{Matrix Example}
A sample matrix:
\[
\begin{pmatrix}
a_{11} & a_{12} & a_{13} \\
a_{21} & a_{22} & a_{23} \\
a_{31} & a_{32} & a_{33}
\end{pmatrix}
\]

\section{Tables}
Below is a sample table showing data:

\begin{table}[h]
\centering
\begin{tabular}{@{}lrrr@{}}
\toprule
\textbf{Category} & \textbf{Value 1} & \textbf{Value 2} & \textbf{Total} \\
\midrule
Item A & 10.5 & 20.3 & 30.8 \\
Item B & 15.2 & 25.1 & 40.3 \\
Item C & 22.1 & 12.7 & 34.8 \\
\bottomrule
\end{tabular}
\caption{Sample Data Table}
\label{tab:sample}
\end{table}

\section{Lists}
\subsection{Bullet Points}
\begin{itemize}
  \item First item
  \item Second item
  \item Third item with sub-items:
    \begin{itemize}
      \item Sub-item 1
      \item Sub-item 2
    \end{itemize}
\end{itemize}

\subsection{Numbered List}
\begin{enumerate}
  \item First step
  \item Second step
  \item Third step
\end{enumerate}

\section{Citations and References}
When you use the citation search feature, references will be automatically managed. For example, you might cite a paper like this: \cite{example2024}. All citations will be automatically added to your \texttt{cite.bib} file.

\section{Conclusion}
This template demonstrates basic \LaTeX{} features including mathematical formulas, tables, and lists. The bibliography system is now set up to work seamlessly with Bibby AI's citation management. Feel free to modify it for your needs!

\bibliographystyle{plain}
\bibliography{cite}

\end{document}\documentclass{article}
\usepackage[utf8]{inputenc}
\usepackage{booktabs}
\usepackage{amsmath}
\usepackage{natbib}

\title{My Document}
\author{Your Name}
\date{\today}

\begin{document}

\maketitle

\section{Introduction}
This is a sample document demonstrating various \LaTeX{} features including mathematical equations, tables, and formatting. You can cite papers using the citation search feature, and they will be automatically added to your \texttt{cite.bib} file.

\section{Mathematical Expressions}
Here are some example equations:

\subsection{Quadratic Formula}
The quadratic formula is given by:
\[x = \frac{-b \pm \sqrt{b^2 - 4ac}}{2a}\]

\subsection{Matrix Example}
A sample matrix:
\[
\begin{pmatrix}
a_{11} & a_{12} & a_{13} \\
a_{21} & a_{22} & a_{23} \\
a_{31} & a_{32} & a_{33}
\end{pmatrix}
\]

\section{Tables}
Below is a sample table showing data:

\begin{table}[h]
\centering
\begin{tabular}{@{}lrrr@{}}
\toprule
\textbf{Category} & \textbf{Value 1} & \textbf{Value 2} & \textbf{Total} \\
\midrule
Item A & 10.5 & 20.3 & 30.8 \\
Item B & 15.2 & 25.1 & 40.3 \\
Item C & 22.1 & 12.7 & 34.8 \\
\bottomrule
\end{tabular}
\caption{Sample Data Table}
\label{tab:sample}
\end{table}

\section{Lists}
\subsection{Bullet Points}
\begin{itemize}
  \item First item
  \item Second item
  \item Third item with sub-items:
    \begin{itemize}
      \item Sub-item 1
      \item Sub-item 2
    \end{itemize}
\end{itemize}

\subsection{Numbered List}
\begin{enumerate}
  \item First step
  \item Second step
  \item Third step
\end{enumerate}

\section{Citations and References}
When you use the citation search feature, references will be automatically managed. For example, you might cite a paper like this: \cite{example2024}. All citations will be automatically added to your \texttt{cite.bib} file.

\section{Conclusion}
This template demonstrates basic \LaTeX{} features including mathematical formulas, tables, and lists. The bibliography system is now set up to work seamlessly with Bibby AI's citation management. Feel free to modify it for your needs!

\bibliographystyle{plain}
\bibliography{cite}

\end{document}\documentclass{article}
\usepackage[utf8]{inputenc}
\usepackage{booktabs}
\usepackage{amsmath}
\usepackage{natbib}

\title{My Document}
\author{Your Name}
\date{\today}

\begin{document}

\maketitle

\section{Introduction}
This is a sample document demonstrating various \LaTeX{} features including mathematical equations, tables, and formatting. You can cite papers using the citation search feature, and they will be automatically added to your \texttt{cite.bib} file.

\section{Mathematical Expressions}
Here are some example equations:

\subsection{Quadratic Formula}
The quadratic formula is given by:
\[x = \frac{-b \pm \sqrt{b^2 - 4ac}}{2a}\]

\subsection{Matrix Example}
A sample matrix:
\[
\begin{pmatrix}
a_{11} & a_{12} & a_{13} \\
a_{21} & a_{22} & a_{23} \\
a_{31} & a_{32} & a_{33}
\end{pmatrix}
\]

\section{Tables}
Below is a sample table showing data:

\begin{table}[h]
\centering
\begin{tabular}{@{}lrrr@{}}
\toprule
\textbf{Category} & \textbf{Value 1} & \textbf{Value 2} & \textbf{Total} \\
\midrule
Item A & 10.5 & 20.3 & 30.8 \\
Item B & 15.2 & 25.1 & 40.3 \\
Item C & 22.1 & 12.7 & 34.8 \\
\bottomrule
\end{tabular}
\caption{Sample Data Table}
\label{tab:sample}
\end{table}

\section{Lists}
\subsection{Bullet Points}
\begin{itemize}
  \item First item
  \item Second item
  \item Third item with sub-items:
    \begin{itemize}
      \item Sub-item 1
      \item Sub-item 2
    \end{itemize}
\end{itemize}

\subsection{Numbered List}
\begin{enumerate}
  \item First step
  \item Second step
  \item Third step
\end{enumerate}

\section{Citations and References}
When you use the citation search feature, references will be automatically managed. For example, you might cite a paper like this: \cite{example2024}. All citations will be automatically added to your \texttt{cite.bib} file.

\section{Conclusion}
This template demonstrates basic \LaTeX{} features including mathematical formulas, tables, and lists. The bibliography system is now set up to work seamlessly with Bibby AI's citation management. Feel free to modify it for your needs!

\bibliographystyle{plain}
\bibliography{cite}

\end{document}\documentclass{article}
\usepackage[utf8]{inputenc}
\usepackage{booktabs}
\usepackage{amsmath}
\usepackage{natbib}

\title{My Document}
\author{Your Name}
\date{\today}

\begin{document}

\maketitle

\section{Introduction}
This is a sample document demonstrating various \LaTeX{} features including mathematical equations, tables, and formatting. You can cite papers using the citation search feature, and they will be automatically added to your \texttt{cite.bib} file.

\section{Mathematical Expressions}
Here are some example equations:

\subsection{Quadratic Formula}
The quadratic formula is given by:
\[x = \frac{-b \pm \sqrt{b^2 - 4ac}}{2a}\]

\subsection{Matrix Example}
A sample matrix:
\[
\begin{pmatrix}
a_{11} & a_{12} & a_{13} \\
a_{21} & a_{22} & a_{23} \\
a_{31} & a_{32} & a_{33}
\end{pmatrix}
\]

\section{Tables}
Below is a sample table showing data:

\begin{table}[h]
\centering
\begin{tabular}{@{}lrrr@{}}
\toprule
\textbf{Category} & \textbf{Value 1} & \textbf{Value 2} & \textbf{Total} \\
\midrule
Item A & 10.5 & 20.3 & 30.8 \\
Item B & 15.2 & 25.1 & 40.3 \\
Item C & 22.1 & 12.7 & 34.8 \\
\bottomrule
\end{tabular}
\caption{Sample Data Table}
\label{tab:sample}
\end{table}

\section{Lists}
\subsection{Bullet Points}
\begin{itemize}
  \item First item
  \item Second item
  \item Third item with sub-items:
    \begin{itemize}
      \item Sub-item 1
      \item Sub-item 2
    \end{itemize}
\end{itemize}

\subsection{Numbered List}
\begin{enumerate}
  \item First step
  \item Second step
  \item Third step
\end{enumerate}

\section{Citations and References}
When you use the citation search feature, references will be automatically managed. For example, you might cite a paper like this: \cite{example2024}. All citations will be automatically added to your \texttt{cite.bib} file.

\section{Conclusion}
This template demonstrates basic \LaTeX{} features including mathematical formulas, tables, and lists. The bibliography system is now set up to work seamlessly with Bibby AI's citation management. Feel free to modify it for your needs!

\bibliographystyle{plain}
\bibliography{cite}

\end{document}\documentclass{article}
\usepackage[utf8]{inputenc}
\usepackage{booktabs}
\usepackage{amsmath}
\usepackage{natbib}

\title{My Document}
\author{Your Name}
\date{\today}

\begin{document}

\maketitle

\section{Introduction}
This is a sample document demonstrating various \LaTeX{} features including mathematical equations, tables, and formatting. You can cite papers using the citation search feature, and they will be automatically added to your \texttt{cite.bib} file.

\section{Mathematical Expressions}
Here are some example equations:

\subsection{Quadratic Formula}
The quadratic formula is given by:
\[x = \frac{-b \pm \sqrt{b^2 - 4ac}}{2a}\]

\subsection{Matrix Example}
A sample matrix:
\[
\begin{pmatrix}
a_{11} & a_{12} & a_{13} \\
a_{21} & a_{22} & a_{23} \\
a_{31} & a_{32} & a_{33}
\end{pmatrix}
\]

\section{Tables}
Below is a sample table showing data:

\begin{table}[h]
\centering
\begin{tabular}{@{}lrrr@{}}
\toprule
\textbf{Category} & \textbf{Value 1} & \textbf{Value 2} & \textbf{Total} \\
\midrule
Item A & 10.5 & 20.3 & 30.8 \\
Item B & 15.2 & 25.1 & 40.3 \\
Item C & 22.1 & 12.7 & 34.8 \\
\bottomrule
\end{tabular}
\caption{Sample Data Table}
\label{tab:sample}
\end{table}

\section{Lists}
\subsection{Bullet Points}
\begin{itemize}
  \item First item
  \item Second item
  \item Third item with sub-items:
    \begin{itemize}
      \item Sub-item 1
      \item Sub-item 2
    \end{itemize}
\end{itemize}

\subsection{Numbered List}
\begin{enumerate}
  \item First step
  \item Second step
  \item Third step
\end{enumerate}

\section{Citations and References}
When you use the citation search feature, references will be automatically managed. For example, you might cite a paper like this: \cite{example2024}. All citations will be automatically added to your \texttt{cite.bib} file.

\section{Conclusion}
This template demonstrates basic \LaTeX{} features including mathematical formulas, tables, and lists. The bibliography system is now set up to work seamlessly with Bibby AI's citation management. Feel free to modify it for your needs!

\bibliographystyle{plain}
\bibliography{cite}

\end{document}\documentclass{article}
\usepackage[utf8]{inputenc}
\usepackage{booktabs}
\usepackage{amsmath}
\usepackage{natbib}

\title{My Document}
\author{Your Name}
\date{\today}

\begin{document}

\maketitle

\section{Introduction}
This is a sample document demonstrating various \LaTeX{} features including mathematical equations, tables, and formatting. You can cite papers using the citation search feature, and they will be automatically added to your \texttt{cite.bib} file.

\section{Mathematical Expressions}
Here are some example equations:

\subsection{Quadratic Formula}
The quadratic formula is given by:
\[x = \frac{-b \pm \sqrt{b^2 - 4ac}}{2a}\]

\subsection{Matrix Example}
A sample matrix:
\[
\begin{pmatrix}
a_{11} & a_{12} & a_{13} \\
a_{21} & a_{22} & a_{23} \\
a_{31} & a_{32} & a_{33}
\end{pmatrix}
\]

\section{Tables}
Below is a sample table showing data:

\begin{table}[h]
\centering
\begin{tabular}{@{}lrrr@{}}
\toprule
\textbf{Category} & \textbf{Value 1} & \textbf{Value 2} & \textbf{Total} \\
\midrule
Item A & 10.5 & 20.3 & 30.8 \\
Item B & 15.2 & 25.1 & 40.3 \\
Item C & 22.1 & 12.7 & 34.8 \\
\bottomrule
\end{tabular}
\caption{Sample Data Table}
\label{tab:sample}
\end{table}

\section{Lists}
\subsection{Bullet Points}
\begin{itemize}
  \item First item
  \item Second item
  \item Third item with sub-items:
    \begin{itemize}
      \item Sub-item 1
      \item Sub-item 2
    \end{itemize}
\end{itemize}

\subsection{Numbered List}
\begin{enumerate}
  \item First step
  \item Second step
  \item Third step
\end{enumerate}

\section{Citations and References}
When you use the citation search feature, references will be automatically managed. For example, you might cite a paper like this: \cite{example2024}. All citations will be automatically added to your \texttt{cite.bib} file.

\section{Conclusion}
This template demonstrates basic \LaTeX{} features including mathematical formulas, tables, and lists. The bibliography system is now set up to work seamlessly with Bibby AI's citation management. Feel free to modify it for your needs!

\bibliographystyle{plain}
\bibliography{cite}

\end{document}\documentclass{article}
\usepackage[utf8]{inputenc}
\usepackage{booktabs}
\usepackage{amsmath}
\usepackage{natbib}

\title{My Document}
\author{Your Name}
\date{\today}

\begin{document}

\maketitle

\section{Introduction}
This is a sample document demonstrating various \LaTeX{} features including mathematical equations, tables, and formatting. You can cite papers using the citation search feature, and they will be automatically added to your \texttt{cite.bib} file.

\section{Mathematical Expressions}
Here are some example equations:

\subsection{Quadratic Formula}
The quadratic formula is given by:
\[x = \frac{-b \pm \sqrt{b^2 - 4ac}}{2a}\]

\subsection{Matrix Example}
A sample matrix:
\[
\begin{pmatrix}
a_{11} & a_{12} & a_{13} \\
a_{21} & a_{22} & a_{23} \\
a_{31} & a_{32} & a_{33}
\end{pmatrix}
\]

\section{Tables}
Below is a sample table showing data:

\begin{table}[h]
\centering
\begin{tabular}{@{}lrrr@{}}
\toprule
\textbf{Category} & \textbf{Value 1} & \textbf{Value 2} & \textbf{Total} \\
\midrule
Item A & 10.5 & 20.3 & 30.8 \\
Item B & 15.2 & 25.1 & 40.3 \\
Item C & 22.1 & 12.7 & 34.8 \\
\bottomrule
\end{tabular}
\caption{Sample Data Table}
\label{tab:sample}
\end{table}

\section{Lists}
\subsection{Bullet Points}
\begin{itemize}
  \item First item
  \item Second item
  \item Third item with sub-items:
    \begin{itemize}
      \item Sub-item 1
      \item Sub-item 2
    \end{itemize}
\end{itemize}

\subsection{Numbered List}
\begin{enumerate}
  \item First step
  \item Second step
  \item Third step
\end{enumerate}

\section{Citations and References}
When you use the citation search feature, references will be automatically managed. For example, you might cite a paper like this: \cite{example2024}. All citations will be automatically added to your \texttt{cite.bib} file.

\section{Conclusion}
This template demonstrates basic \LaTeX{} features including mathematical formulas, tables, and lists. The bibliography system is now set up to work seamlessly with Bibby AI's citation management. Feel free to modify it for your needs!

\bibliographystyle{plain}
\bibliography{cite}

\end{document}\documentclass{article}
\usepackage[utf8]{inputenc}
\usepackage{booktabs}
\usepackage{amsmath}
\usepackage{natbib}

\title{My Document}
\author{Your Name}
\date{\today}

\begin{document}

\maketitle

\section{Introduction}
This is a sample document demonstrating various \LaTeX{} features including mathematical equations, tables, and formatting. You can cite papers using the citation search feature, and they will be automatically added to your \texttt{cite.bib} file.

\section{Mathematical Expressions}
Here are some example equations:

\subsection{Quadratic Formula}
The quadratic formula is given by:
\[x = \frac{-b \pm \sqrt{b^2 - 4ac}}{2a}\]

\subsection{Matrix Example}
A sample matrix:
\[
\begin{pmatrix}
a_{11} & a_{12} & a_{13} \\
a_{21} & a_{22} & a_{23} \\
a_{31} & a_{32} & a_{33}
\end{pmatrix}
\]

\section{Tables}
Below is a sample table showing data:

\begin{table}[h]
\centering
\begin{tabular}{@{}lrrr@{}}
\toprule
\textbf{Category} & \textbf{Value 1} & \textbf{Value 2} & \textbf{Total} \\
\midrule
Item A & 10.5 & 20.3 & 30.8 \\
Item B & 15.2 & 25.1 & 40.3 \\
Item C & 22.1 & 12.7 & 34.8 \\
\bottomrule
\end{tabular}
\caption{Sample Data Table}
\label{tab:sample}
\end{table}

\section{Lists}
\subsection{Bullet Points}
\begin{itemize}
  \item First item
  \item Second item
  \item Third item with sub-items:
    \begin{itemize}
      \item Sub-item 1
      \item Sub-item 2
    \end{itemize}
\end{itemize}

\subsection{Numbered List}
\begin{enumerate}
  \item First step
  \item Second step
  \item Third step
\end{enumerate}

\section{Citations and References}
When you use the citation search feature, references will be automatically managed. For example, you might cite a paper like this: \cite{example2024}. All citations will be automatically added to your \texttt{cite.bib} file.

\section{Conclusion}
This template demonstrates basic \LaTeX{} features including mathematical formulas, tables, and lists. The bibliography system is now set up to work seamlessly with Bibby AI's citation management. Feel free to modify it for your needs!

\bibliographystyle{plain}
\bibliography{cite}

\end{document}\documentclass{article}
\usepackage[utf8]{inputenc}
\usepackage{booktabs}
\usepackage{amsmath}
\usepackage{natbib}

\title{My Document}
\author{Your Name}
\date{\today}

\begin{document}

\maketitle

\section{Introduction}
This is a sample document demonstrating various \LaTeX{} features including mathematical equations, tables, and formatting. You can cite papers using the citation search feature, and they will be automatically added to your \texttt{cite.bib} file.

\section{Mathematical Expressions}
Here are some example equations:

\subsection{Quadratic Formula}
The quadratic formula is given by:
\[x = \frac{-b \pm \sqrt{b^2 - 4ac}}{2a}\]

\subsection{Matrix Example}
A sample matrix:
\[
\begin{pmatrix}
a_{11} & a_{12} & a_{13} \\
a_{21} & a_{22} & a_{23} \\
a_{31} & a_{32} & a_{33}
\end{pmatrix}
\]

\section{Tables}
Below is a sample table showing data:

\begin{table}[h]
\centering
\begin{tabular}{@{}lrrr@{}}
\toprule
\textbf{Category} & \textbf{Value 1} & \textbf{Value 2} & \textbf{Total} \\
\midrule
Item A & 10.5 & 20.3 & 30.8 \\
Item B & 15.2 & 25.1 & 40.3 \\
Item C & 22.1 & 12.7 & 34.8 \\
\bottomrule
\end{tabular}
\caption{Sample Data Table}
\label{tab:sample}
\end{table}

\section{Lists}
\subsection{Bullet Points}
\begin{itemize}
  \item First item
  \item Second item
  \item Third item with sub-items:
    \begin{itemize}
      \item Sub-item 1
      \item Sub-item 2
    \end{itemize}
\end{itemize}

\subsection{Numbered List}
\begin{enumerate}
  \item First step
  \item Second step
  \item Third step
\end{enumerate}

\section{Citations and References}
When you use the citation search feature, references will be automatically managed. For example, you might cite a paper like this: \cite{example2024}. All citations will be automatically added to your \texttt{cite.bib} file.

\section{Conclusion}
This template demonstrates basic \LaTeX{} features including mathematical formulas, tables, and lists. The bibliography system is now set up to work seamlessly with Bibby AI's citation management. Feel free to modify it for your needs!

\bibliographystyle{plain}
\bibliography{cite}

\end{document}\documentclass{article}
\usepackage[utf8]{inputenc}
\usepackage{booktabs}
\usepackage{amsmath}
\usepackage{natbib}

\title{My Document}
\author{Your Name}
\date{\today}

\begin{document}

\maketitle

\section{Introduction}
This is a sample document demonstrating various \LaTeX{} features including mathematical equations, tables, and formatting. You can cite papers using the citation search feature, and they will be automatically added to your \texttt{cite.bib} file.

\section{Mathematical Expressions}
Here are some example equations:

\subsection{Quadratic Formula}
The quadratic formula is given by:
\[x = \frac{-b \pm \sqrt{b^2 - 4ac}}{2a}\]

\subsection{Matrix Example}
A sample matrix:
\[
\begin{pmatrix}
a_{11} & a_{12} & a_{13} \\
a_{21} & a_{22} & a_{23} \\
a_{31} & a_{32} & a_{33}
\end{pmatrix}
\]

\section{Tables}
Below is a sample table showing data:

\begin{table}[h]
\centering
\begin{tabular}{@{}lrrr@{}}
\toprule
\textbf{Category} & \textbf{Value 1} & \textbf{Value 2} & \textbf{Total} \\
\midrule
Item A & 10.5 & 20.3 & 30.8 \\
Item B & 15.2 & 25.1 & 40.3 \\
Item C & 22.1 & 12.7 & 34.8 \\
\bottomrule
\end{tabular}
\caption{Sample Data Table}
\label{tab:sample}
\end{table}

\section{Lists}
\subsection{Bullet Points}
\begin{itemize}
  \item First item
  \item Second item
  \item Third item with sub-items:
    \begin{itemize}
      \item Sub-item 1
      \item Sub-item 2
    \end{itemize}
\end{itemize}

\subsection{Numbered List}
\begin{enumerate}
  \item First step
  \item Second step
  \item Third step
\end{enumerate}

\section{Citations and References}
When you use the citation search feature, references will be automatically managed. For example, you might cite a paper like this: \cite{example2024}. All citations will be automatically added to your \texttt{cite.bib} file.

\section{Conclusion}
This template demonstrates basic \LaTeX{} features including mathematical formulas, tables, and lists. The bibliography system is now set up to work seamlessly with Bibby AI's citation management. Feel free to modify it for your needs!

\bibliographystyle{plain}
\bibliography{cite}

\end{document}\documentclass{article}
\usepackage[utf8]{inputenc}
\usepackage{booktabs}
\usepackage{amsmath}
\usepackage{natbib}

\title{My Document}
\author{Your Name}
\date{\today}

\begin{document}

\maketitle

\section{Introduction}
This is a sample document demonstrating various \LaTeX{} features including mathematical equations, tables, and formatting. You can cite papers using the citation search feature, and they will be automatically added to your \texttt{cite.bib} file.

\section{Mathematical Expressions}
Here are some example equations:

\subsection{Quadratic Formula}
The quadratic formula is given by:
\[x = \frac{-b \pm \sqrt{b^2 - 4ac}}{2a}\]

\subsection{Matrix Example}
A sample matrix:
\[
\begin{pmatrix}
a_{11} & a_{12} & a_{13} \\
a_{21} & a_{22} & a_{23} \\
a_{31} & a_{32} & a_{33}
\end{pmatrix}
\]

\section{Tables}
Below is a sample table showing data:

\begin{table}[h]
\centering
\begin{tabular}{@{}lrrr@{}}
\toprule
\textbf{Category} & \textbf{Value 1} & \textbf{Value 2} & \textbf{Total} \\
\midrule
Item A & 10.5 & 20.3 & 30.8 \\
Item B & 15.2 & 25.1 & 40.3 \\
Item C & 22.1 & 12.7 & 34.8 \\
\bottomrule
\end{tabular}
\caption{Sample Data Table}
\label{tab:sample}
\end{table}

\section{Lists}
\subsection{Bullet Points}
\begin{itemize}
  \item First item
  \item Second item
  \item Third item with sub-items:
    \begin{itemize}
      \item Sub-item 1
      \item Sub-item 2
    \end{itemize}
\end{itemize}

\subsection{Numbered List}
\begin{enumerate}
  \item First step
  \item Second step
  \item Third step
\end{enumerate}

\section{Citations and References}
When you use the citation search feature, references will be automatically managed. For example, you might cite a paper like this: \cite{example2024}. All citations will be automatically added to your \texttt{cite.bib} file.

\section{Conclusion}
This template demonstrates basic \LaTeX{} features including mathematical formulas, tables, and lists. The bibliography system is now set up to work seamlessly with Bibby AI's citation management. Feel free to modify it for your needs!

\bibliographystyle{plain}
\bibliography{cite}

\end{document}\documentclass{article}
\usepackage[utf8]{inputenc}
\usepackage{booktabs}
\usepackage{amsmath}
\usepackage{natbib}

\title{My Document}
\author{Your Name}
\date{\today}

\begin{document}

\maketitle

\section{Introduction}
This is a sample document demonstrating various \LaTeX{} features including mathematical equations, tables, and formatting. You can cite papers using the citation search feature, and they will be automatically added to your \texttt{cite.bib} file.

\section{Mathematical Expressions}
Here are some example equations:

\subsection{Quadratic Formula}
The quadratic formula is given by:
\[x = \frac{-b \pm \sqrt{b^2 - 4ac}}{2a}\]

\subsection{Matrix Example}
A sample matrix:
\[
\begin{pmatrix}
a_{11} & a_{12} & a_{13} \\
a_{21} & a_{22} & a_{23} \\
a_{31} & a_{32} & a_{33}
\end{pmatrix}
\]

\section{Tables}
Below is a sample table showing data:

\begin{table}[h]
\centering
\begin{tabular}{@{}lrrr@{}}
\toprule
\textbf{Category} & \textbf{Value 1} & \textbf{Value 2} & \textbf{Total} \\
\midrule
Item A & 10.5 & 20.3 & 30.8 \\
Item B & 15.2 & 25.1 & 40.3 \\
Item C & 22.1 & 12.7 & 34.8 \\
\bottomrule
\end{tabular}
\caption{Sample Data Table}
\label{tab:sample}
\end{table}

\section{Lists}
\subsection{Bullet Points}
\begin{itemize}
  \item First item
  \item Second item
  \item Third item with sub-items:
    \begin{itemize}
      \item Sub-item 1
      \item Sub-item 2
    \end{itemize}
\end{itemize}

\subsection{Numbered List}
\begin{enumerate}
  \item First step
  \item Second step
  \item Third step
\end{enumerate}

\section{Citations and References}
When you use the citation search feature, references will be automatically managed. For example, you might cite a paper like this: \cite{example2024}. All citations will be automatically added to your \texttt{cite.bib} file.

\section{Conclusion}
This template demonstrates basic \LaTeX{} features including mathematical formulas, tables, and lists. The bibliography system is now set up to work seamlessly with Bibby AI's citation management. Feel free to modify it for your needs!

\bibliographystyle{plain}
\bibliography{cite}

\end{document}\documentclass{article}
\usepackage[utf8]{inputenc}
\usepackage{booktabs}
\usepackage{amsmath}
\usepackage{natbib}

\title{My Document}
\author{Your Name}
\date{\today}

\begin{document}

\maketitle

\section{Introduction}
This is a sample document demonstrating various \LaTeX{} features including mathematical equations, tables, and formatting. You can cite papers using the citation search feature, and they will be automatically added to your \texttt{cite.bib} file.

\section{Mathematical Expressions}
Here are some example equations:

\subsection{Quadratic Formula}
The quadratic formula is given by:
\[x = \frac{-b \pm \sqrt{b^2 - 4ac}}{2a}\]

\subsection{Matrix Example}
A sample matrix:
\[
\begin{pmatrix}
a_{11} & a_{12} & a_{13} \\
a_{21} & a_{22} & a_{23} \\
a_{31} & a_{32} & a_{33}
\end{pmatrix}
\]

\section{Tables}
Below is a sample table showing data:

\begin{table}[h]
\centering
\begin{tabular}{@{}lrrr@{}}
\toprule
\textbf{Category} & \textbf{Value 1} & \textbf{Value 2} & \textbf{Total} \\
\midrule
Item A & 10.5 & 20.3 & 30.8 \\
Item B & 15.2 & 25.1 & 40.3 \\
Item C & 22.1 & 12.7 & 34.8 \\
\bottomrule
\end{tabular}
\caption{Sample Data Table}
\label{tab:sample}
\end{table}

\section{Lists}
\subsection{Bullet Points}
\begin{itemize}
  \item First item
  \item Second item
  \item Third item with sub-items:
    \begin{itemize}
      \item Sub-item 1
      \item Sub-item 2
    \end{itemize}
\end{itemize}

\subsection{Numbered List}
\begin{enumerate}
  \item First step
  \item Second step
  \item Third step
\end{enumerate}

\section{Citations and References}
When you use the citation search feature, references will be automatically managed. For example, you might cite a paper like this: \cite{example2024}. All citations will be automatically added to your \texttt{cite.bib} file.

\section{Conclusion}
This template demonstrates basic \LaTeX{} features including mathematical formulas, tables, and lists. The bibliography system is now set up to work seamlessly with Bibby AI's citation management. Feel free to modify it for your needs!

\bibliographystyle{plain}
\bibliography{cite}

\end{document}\documentclass{article}
\usepackage[utf8]{inputenc}
\usepackage{booktabs}
\usepackage{amsmath}
\usepackage{natbib}

\title{My Document}
\author{Your Name}
\date{\today}

\begin{document}

\maketitle

\section{Introduction}
This is a sample document demonstrating various \LaTeX{} features including mathematical equations, tables, and formatting. You can cite papers using the citation search feature, and they will be automatically added to your \texttt{cite.bib} file.

\section{Mathematical Expressions}
Here are some example equations:

\subsection{Quadratic Formula}
The quadratic formula is given by:
\[x = \frac{-b \pm \sqrt{b^2 - 4ac}}{2a}\]

\subsection{Matrix Example}
A sample matrix:
\[
\begin{pmatrix}
a_{11} & a_{12} & a_{13} \\
a_{21} & a_{22} & a_{23} \\
a_{31} & a_{32} & a_{33}
\end{pmatrix}
\]

\section{Tables}
Below is a sample table showing data:

\begin{table}[h]
\centering
\begin{tabular}{@{}lrrr@{}}
\toprule
\textbf{Category} & \textbf{Value 1} & \textbf{Value 2} & \textbf{Total} \\
\midrule
Item A & 10.5 & 20.3 & 30.8 \\
Item B & 15.2 & 25.1 & 40.3 \\
Item C & 22.1 & 12.7 & 34.8 \\
\bottomrule
\end{tabular}
\caption{Sample Data Table}
\label{tab:sample}
\end{table}

\section{Lists}
\subsection{Bullet Points}
\begin{itemize}
  \item First item
  \item Second item
  \item Third item with sub-items:
    \begin{itemize}
      \item Sub-item 1
      \item Sub-item 2
    \end{itemize}
\end{itemize}

\subsection{Numbered List}
\begin{enumerate}
  \item First step
  \item Second step
  \item Third step
\end{enumerate}

\section{Citations and References}
When you use the citation search feature, references will be automatically managed. For example, you might cite a paper like this: \cite{example2024}. All citations will be automatically added to your \texttt{cite.bib} file.

\section{Conclusion}
This template demonstrates basic \LaTeX{} features including mathematical formulas, tables, and lists. The bibliography system is now set up to work seamlessly with Bibby AI's citation management. Feel free to modify it for your needs!

\bibliographystyle{plain}
\bibliography{cite}

\end{document}\documentclass{article}
\usepackage[utf8]{inputenc}
\usepackage{booktabs}
\usepackage{amsmath}
\usepackage{natbib}

\title{My Document}
\author{Your Name}
\date{\today}

\begin{document}

\maketitle

\section{Introduction}
This is a sample document demonstrating various \LaTeX{} features including mathematical equations, tables, and formatting. You can cite papers using the citation search feature, and they will be automatically added to your \texttt{cite.bib} file.

\section{Mathematical Expressions}
Here are some example equations:

\subsection{Quadratic Formula}
The quadratic formula is given by:
\[x = \frac{-b \pm \sqrt{b^2 - 4ac}}{2a}\]

\subsection{Matrix Example}
A sample matrix:
\[
\begin{pmatrix}
a_{11} & a_{12} & a_{13} \\
a_{21} & a_{22} & a_{23} \\
a_{31} & a_{32} & a_{33}
\end{pmatrix}
\]

\section{Tables}
Below is a sample table showing data:

\begin{table}[h]
\centering
\begin{tabular}{@{}lrrr@{}}
\toprule
\textbf{Category} & \textbf{Value 1} & \textbf{Value 2} & \textbf{Total} \\
\midrule
Item A & 10.5 & 20.3 & 30.8 \\
Item B & 15.2 & 25.1 & 40.3 \\
Item C & 22.1 & 12.7 & 34.8 \\
\bottomrule
\end{tabular}
\caption{Sample Data Table}
\label{tab:sample}
\end{table}

\section{Lists}
\subsection{Bullet Points}
\begin{itemize}
  \item First item
  \item Second item
  \item Third item with sub-items:
    \begin{itemize}
      \item Sub-item 1
      \item Sub-item 2
    \end{itemize}
\end{itemize}

\subsection{Numbered List}
\begin{enumerate}
  \item First step
  \item Second step
  \item Third step
\end{enumerate}

\section{Citations and References}
When you use the citation search feature, references will be automatically managed. For example, you might cite a paper like this: \cite{example2024}. All citations will be automatically added to your \texttt{cite.bib} file.

\section{Conclusion}
This template demonstrates basic \LaTeX{} features including mathematical formulas, tables, and lists. The bibliography system is now set up to work seamlessly with Bibby AI's citation management. Feel free to modify it for your needs!

\bibliographystyle{plain}
\bibliography{cite}

\end{document}\documentclass{article}
\usepackage[utf8]{inputenc}
\usepackage{booktabs}
\usepackage{amsmath}
\usepackage{natbib}

\title{My Document}
\author{Your Name}
\date{\today}

\begin{document}

\maketitle

\section{Introduction}
This is a sample document demonstrating various \LaTeX{} features including mathematical equations, tables, and formatting. You can cite papers using the citation search feature, and they will be automatically added to your \texttt{cite.bib} file.

\section{Mathematical Expressions}
Here are some example equations:

\subsection{Quadratic Formula}
The quadratic formula is given by:
\[x = \frac{-b \pm \sqrt{b^2 - 4ac}}{2a}\]

\subsection{Matrix Example}
A sample matrix:
\[
\begin{pmatrix}
a_{11} & a_{12} & a_{13} \\
a_{21} & a_{22} & a_{23} \\
a_{31} & a_{32} & a_{33}
\end{pmatrix}
\]

\section{Tables}
Below is a sample table showing data:

\begin{table}[h]
\centering
\begin{tabular}{@{}lrrr@{}}
\toprule
\textbf{Category} & \textbf{Value 1} & \textbf{Value 2} & \textbf{Total} \\
\midrule
Item A & 10.5 & 20.3 & 30.8 \\
Item B & 15.2 & 25.1 & 40.3 \\
Item C & 22.1 & 12.7 & 34.8 \\
\bottomrule
\end{tabular}
\caption{Sample Data Table}
\label{tab:sample}
\end{table}

\section{Lists}
\subsection{Bullet Points}
\begin{itemize}
  \item First item
  \item Second item
  \item Third item with sub-items:
    \begin{itemize}
      \item Sub-item 1
      \item Sub-item 2
    \end{itemize}
\end{itemize}

\subsection{Numbered List}
\begin{enumerate}
  \item First step
  \item Second step
  \item Third step
\end{enumerate}

\section{Citations and References}
When you use the citation search feature, references will be automatically managed. For example, you might cite a paper like this: \cite{example2024}. All citations will be automatically added to your \texttt{cite.bib} file.

\section{Conclusion}
This template demonstrates basic \LaTeX{} features including mathematical formulas, tables, and lists. The bibliography system is now set up to work seamlessly with Bibby AI's citation management. Feel free to modify it for your needs!

\bibliographystyle{plain}
\bibliography{cite}

\end{document}\documentclass{article}
\usepackage[utf8]{inputenc}
\usepackage{booktabs}
\usepackage{amsmath}
\usepackage{natbib}

\title{My Document}
\author{Your Name}
\date{\today}

\begin{document}

\maketitle

\section{Introduction}
This is a sample document demonstrating various \LaTeX{} features including mathematical equations, tables, and formatting. You can cite papers using the citation search feature, and they will be automatically added to your \texttt{cite.bib} file.

\section{Mathematical Expressions}
Here are some example equations:

\subsection{Quadratic Formula}
The quadratic formula is given by:
\[x = \frac{-b \pm \sqrt{b^2 - 4ac}}{2a}\]

\subsection{Matrix Example}
A sample matrix:
\[
\begin{pmatrix}
a_{11} & a_{12} & a_{13} \\
a_{21} & a_{22} & a_{23} \\
a_{31} & a_{32} & a_{33}
\end{pmatrix}
\]

\section{Tables}
Below is a sample table showing data:

\begin{table}[h]
\centering
\begin{tabular}{@{}lrrr@{}}
\toprule
\textbf{Category} & \textbf{Value 1} & \textbf{Value 2} & \textbf{Total} \\
\midrule
Item A & 10.5 & 20.3 & 30.8 \\
Item B & 15.2 & 25.1 & 40.3 \\
Item C & 22.1 & 12.7 & 34.8 \\
\bottomrule
\end{tabular}
\caption{Sample Data Table}
\label{tab:sample}
\end{table}

\section{Lists}
\subsection{Bullet Points}
\begin{itemize}
  \item First item
  \item Second item
  \item Third item with sub-items:
    \begin{itemize}
      \item Sub-item 1
      \item Sub-item 2
    \end{itemize}
\end{itemize}

\subsection{Numbered List}
\begin{enumerate}
  \item First step
  \item Second step
  \item Third step
\end{enumerate}

\section{Citations and References}
When you use the citation search feature, references will be automatically managed. For example, you might cite a paper like this: \cite{example2024}. All citations will be automatically added to your \texttt{cite.bib} file.

\section{Conclusion}
This template demonstrates basic \LaTeX{} features including mathematical formulas, tables, and lists. The bibliography system is now set up to work seamlessly with Bibby AI's citation management. Feel free to modify it for your needs!

\bibliographystyle{plain}
\bibliography{cite}

\end{document}\documentclass{article}
\usepackage[utf8]{inputenc}
\usepackage{booktabs}
\usepackage{amsmath}
\usepackage{natbib}

\title{My Document}
\author{Your Name}
\date{\today}

\begin{document}

\maketitle

\section{Introduction}
This is a sample document demonstrating various \LaTeX{} features including mathematical equations, tables, and formatting. You can cite papers using the citation search feature, and they will be automatically added to your \texttt{cite.bib} file.

\section{Mathematical Expressions}
Here are some example equations:

\subsection{Quadratic Formula}
The quadratic formula is given by:
\[x = \frac{-b \pm \sqrt{b^2 - 4ac}}{2a}\]

\subsection{Matrix Example}
A sample matrix:
\[
\begin{pmatrix}
a_{11} & a_{12} & a_{13} \\
a_{21} & a_{22} & a_{23} \\
a_{31} & a_{32} & a_{33}
\end{pmatrix}
\]

\section{Tables}
Below is a sample table showing data:

\begin{table}[h]
\centering
\begin{tabular}{@{}lrrr@{}}
\toprule
\textbf{Category} & \textbf{Value 1} & \textbf{Value 2} & \textbf{Total} \\
\midrule
Item A & 10.5 & 20.3 & 30.8 \\
Item B & 15.2 & 25.1 & 40.3 \\
Item C & 22.1 & 12.7 & 34.8 \\
\bottomrule
\end{tabular}
\caption{Sample Data Table}
\label{tab:sample}
\end{table}

\section{Lists}
\subsection{Bullet Points}
\begin{itemize}
  \item First item
  \item Second item
  \item Third item with sub-items:
    \begin{itemize}
      \item Sub-item 1
      \item Sub-item 2
    \end{itemize}
\end{itemize}

\subsection{Numbered List}
\begin{enumerate}
  \item First step
  \item Second step
  \item Third step
\end{enumerate}

\section{Citations and References}
When you use the citation search feature, references will be automatically managed. For example, you might cite a paper like this: \cite{example2024}. All citations will be automatically added to your \texttt{cite.bib} file.

\section{Conclusion}
This template demonstrates basic \LaTeX{} features including mathematical formulas, tables, and lists. The bibliography system is now set up to work seamlessly with Bibby AI's citation management. Feel free to modify it for your needs!

\bibliographystyle{plain}
\bibliography{cite}

\end{document}\documentclass{article}
\usepackage[utf8]{inputenc}
\usepackage{booktabs}
\usepackage{amsmath}
\usepackage{natbib}

\title{My Document}
\author{Your Name}
\date{\today}

\begin{document}

\maketitle

\section{Introduction}
This is a sample document demonstrating various \LaTeX{} features including mathematical equations, tables, and formatting. You can cite papers using the citation search feature, and they will be automatically added to your \texttt{cite.bib} file.

\section{Mathematical Expressions}
Here are some example equations:

\subsection{Quadratic Formula}
The quadratic formula is given by:
\[x = \frac{-b \pm \sqrt{b^2 - 4ac}}{2a}\]

\subsection{Matrix Example}
A sample matrix:
\[
\begin{pmatrix}
a_{11} & a_{12} & a_{13} \\
a_{21} & a_{22} & a_{23} \\
a_{31} & a_{32} & a_{33}
\end{pmatrix}
\]

\section{Tables}
Below is a sample table showing data:

\begin{table}[h]
\centering
\begin{tabular}{@{}lrrr@{}}
\toprule
\textbf{Category} & \textbf{Value 1} & \textbf{Value 2} & \textbf{Total} \\
\midrule
Item A & 10.5 & 20.3 & 30.8 \\
Item B & 15.2 & 25.1 & 40.3 \\
Item C & 22.1 & 12.7 & 34.8 \\
\bottomrule
\end{tabular}
\caption{Sample Data Table}
\label{tab:sample}
\end{table}

\section{Lists}
\subsection{Bullet Points}
\begin{itemize}
  \item First item
  \item Second item
  \item Third item with sub-items:
    \begin{itemize}
      \item Sub-item 1
      \item Sub-item 2
    \end{itemize}
\end{itemize}

\subsection{Numbered List}
\begin{enumerate}
  \item First step
  \item Second step
  \item Third step
\end{enumerate}

\section{Citations and References}
When you use the citation search feature, references will be automatically managed. For example, you might cite a paper like this: \cite{example2024}. All citations will be automatically added to your \texttt{cite.bib} file.

\section{Conclusion}
This template demonstrates basic \LaTeX{} features including mathematical formulas, tables, and lists. The bibliography system is now set up to work seamlessly with Bibby AI's citation management. Feel free to modify it for your needs!

\bibliographystyle{plain}
\bibliography{cite}

\end{document}\documentclass{article}
\usepackage[utf8]{inputenc}
\usepackage{booktabs}
\usepackage{amsmath}
\usepackage{natbib}

\title{My Document}
\author{Your Name}
\date{\today}

\begin{document}

\maketitle

\section{Introduction}
This is a sample document demonstrating various \LaTeX{} features including mathematical equations, tables, and formatting. You can cite papers using the citation search feature, and they will be automatically added to your \texttt{cite.bib} file.

\section{Mathematical Expressions}
Here are some example equations:

\subsection{Quadratic Formula}
The quadratic formula is given by:
\[x = \frac{-b \pm \sqrt{b^2 - 4ac}}{2a}\]

\subsection{Matrix Example}
A sample matrix:
\[
\begin{pmatrix}
a_{11} & a_{12} & a_{13} \\
a_{21} & a_{22} & a_{23} \\
a_{31} & a_{32} & a_{33}
\end{pmatrix}
\]

\section{Tables}
Below is a sample table showing data:

\begin{table}[h]
\centering
\begin{tabular}{@{}lrrr@{}}
\toprule
\textbf{Category} & \textbf{Value 1} & \textbf{Value 2} & \textbf{Total} \\
\midrule
Item A & 10.5 & 20.3 & 30.8 \\
Item B & 15.2 & 25.1 & 40.3 \\
Item C & 22.1 & 12.7 & 34.8 \\
\bottomrule
\end{tabular}
\caption{Sample Data Table}
\label{tab:sample}
\end{table}

\section{Lists}
\subsection{Bullet Points}
\begin{itemize}
  \item First item
  \item Second item
  \item Third item with sub-items:
    \begin{itemize}
      \item Sub-item 1
      \item Sub-item 2
    \end{itemize}
\end{itemize}

\subsection{Numbered List}
\begin{enumerate}
  \item First step
  \item Second step
  \item Third step
\end{enumerate}

\section{Citations and References}
When you use the citation search feature, references will be automatically managed. For example, you might cite a paper like this: \cite{example2024}. All citations will be automatically added to your \texttt{cite.bib} file.

\section{Conclusion}
This template demonstrates basic \LaTeX{} features including mathematical formulas, tables, and lists. The bibliography system is now set up to work seamlessly with Bibby AI's citation management. Feel free to modify it for your needs!

\bibliographystyle{plain}
\bibliography{cite}

\end{document}\documentclass{article}
\usepackage[utf8]{inputenc}
\usepackage{booktabs}
\usepackage{amsmath}
\usepackage{natbib}

\title{My Document}
\author{Your Name}
\date{\today}

\begin{document}

\maketitle

\section{Introduction}
This is a sample document demonstrating various \LaTeX{} features including mathematical equations, tables, and formatting. You can cite papers using the citation search feature, and they will be automatically added to your \texttt{cite.bib} file.

\section{Mathematical Expressions}
Here are some example equations:

\subsection{Quadratic Formula}
The quadratic formula is given by:
\[x = \frac{-b \pm \sqrt{b^2 - 4ac}}{2a}\]

\subsection{Matrix Example}
A sample matrix:
\[
\begin{pmatrix}
a_{11} & a_{12} & a_{13} \\
a_{21} & a_{22} & a_{23} \\
a_{31} & a_{32} & a_{33}
\end{pmatrix}
\]

\section{Tables}
Below is a sample table showing data:

\begin{table}[h]
\centering
\begin{tabular}{@{}lrrr@{}}
\toprule
\textbf{Category} & \textbf{Value 1} & \textbf{Value 2} & \textbf{Total} \\
\midrule
Item A & 10.5 & 20.3 & 30.8 \\
Item B & 15.2 & 25.1 & 40.3 \\
Item C & 22.1 & 12.7 & 34.8 \\
\bottomrule
\end{tabular}
\caption{Sample Data Table}
\label{tab:sample}
\end{table}

\section{Lists}
\subsection{Bullet Points}
\begin{itemize}
  \item First item
  \item Second item
  \item Third item with sub-items:
    \begin{itemize}
      \item Sub-item 1
      \item Sub-item 2
    \end{itemize}
\end{itemize}

\subsection{Numbered List}
\begin{enumerate}
  \item First step
  \item Second step
  \item Third step
\end{enumerate}

\section{Citations and References}
When you use the citation search feature, references will be automatically managed. For example, you might cite a paper like this: \cite{example2024}. All citations will be automatically added to your \texttt{cite.bib} file.

\section{Conclusion}
This template demonstrates basic \LaTeX{} features including mathematical formulas, tables, and lists. The bibliography system is now set up to work seamlessly with Bibby AI's citation management. Feel free to modify it for your needs!

\bibliographystyle{plain}
\bibliography{cite}

\end{document}\documentclass{article}
\usepackage[utf8]{inputenc}
\usepackage{booktabs}
\usepackage{amsmath}
\usepackage{natbib}

\title{My Document}
\author{Your Name}
\date{\today}

\begin{document}

\maketitle

\section{Introduction}
This is a sample document demonstrating various \LaTeX{} features including mathematical equations, tables, and formatting. You can cite papers using the citation search feature, and they will be automatically added to your \texttt{cite.bib} file.

\section{Mathematical Expressions}
Here are some example equations:

\subsection{Quadratic Formula}
The quadratic formula is given by:
\[x = \frac{-b \pm \sqrt{b^2 - 4ac}}{2a}\]

\subsection{Matrix Example}
A sample matrix:
\[
\begin{pmatrix}
a_{11} & a_{12} & a_{13} \\
a_{21} & a_{22} & a_{23} \\
a_{31} & a_{32} & a_{33}
\end{pmatrix}
\]

\section{Tables}
Below is a sample table showing data:

\begin{table}[h]
\centering
\begin{tabular}{@{}lrrr@{}}
\toprule
\textbf{Category} & \textbf{Value 1} & \textbf{Value 2} & \textbf{Total} \\
\midrule
Item A & 10.5 & 20.3 & 30.8 \\
Item B & 15.2 & 25.1 & 40.3 \\
Item C & 22.1 & 12.7 & 34.8 \\
\bottomrule
\end{tabular}
\caption{Sample Data Table}
\label{tab:sample}
\end{table}

\section{Lists}
\subsection{Bullet Points}
\begin{itemize}
  \item First item
  \item Second item
  \item Third item with sub-items:
    \begin{itemize}
      \item Sub-item 1
      \item Sub-item 2
    \end{itemize}
\end{itemize}

\subsection{Numbered List}
\begin{enumerate}
  \item First step
  \item Second step
  \item Third step
\end{enumerate}

\section{Citations and References}
When you use the citation search feature, references will be automatically managed. For example, you might cite a paper like this: \cite{example2024}. All citations will be automatically added to your \texttt{cite.bib} file.

\section{Conclusion}
This template demonstrates basic \LaTeX{} features including mathematical formulas, tables, and lists. The bibliography system is now set up to work seamlessly with Bibby AI's citation management. Feel free to modify it for your needs!

\bibliographystyle{plain}
\bibliography{cite}

\end{document}\documentclass{article}
\usepackage[utf8]{inputenc}
\usepackage{booktabs}
\usepackage{amsmath}
\usepackage{natbib}

\title{My Document}
\author{Your Name}
\date{\today}

\begin{document}

\maketitle

\section{Introduction}
This is a sample document demonstrating various \LaTeX{} features including mathematical equations, tables, and formatting. You can cite papers using the citation search feature, and they will be automatically added to your \texttt{cite.bib} file.

\section{Mathematical Expressions}
Here are some example equations:

\subsection{Quadratic Formula}
The quadratic formula is given by:
\[x = \frac{-b \pm \sqrt{b^2 - 4ac}}{2a}\]

\subsection{Matrix Example}
A sample matrix:
\[
\begin{pmatrix}
a_{11} & a_{12} & a_{13} \\
a_{21} & a_{22} & a_{23} \\
a_{31} & a_{32} & a_{33}
\end{pmatrix}
\]

\section{Tables}
Below is a sample table showing data:

\begin{table}[h]
\centering
\begin{tabular}{@{}lrrr@{}}
\toprule
\textbf{Category} & \textbf{Value 1} & \textbf{Value 2} & \textbf{Total} \\
\midrule
Item A & 10.5 & 20.3 & 30.8 \\
Item B & 15.2 & 25.1 & 40.3 \\
Item C & 22.1 & 12.7 & 34.8 \\
\bottomrule
\end{tabular}
\caption{Sample Data Table}
\label{tab:sample}
\end{table}

\section{Lists}
\subsection{Bullet Points}
\begin{itemize}
  \item First item
  \item Second item
  \item Third item with sub-items:
    \begin{itemize}
      \item Sub-item 1
      \item Sub-item 2
    \end{itemize}
\end{itemize}

\subsection{Numbered List}
\begin{enumerate}
  \item First step
  \item Second step
  \item Third step
\end{enumerate}

\section{Citations and References}
When you use the citation search feature, references will be automatically managed. For example, you might cite a paper like this: \cite{example2024}. All citations will be automatically added to your \texttt{cite.bib} file.

\section{Conclusion}
This template demonstrates basic \LaTeX{} features including mathematical formulas, tables, and lists. The bibliography system is now set up to work seamlessly with Bibby AI's citation management. Feel free to modify it for your needs!

\bibliographystyle{plain}
\bibliography{cite}

\end{document}\documentclass{article}
\usepackage[utf8]{inputenc}
\usepackage{booktabs}
\usepackage{amsmath}
\usepackage{natbib}

\title{My Document}
\author{Your Name}
\date{\today}

\begin{document}

\maketitle

\section{Introduction}
This is a sample document demonstrating various \LaTeX{} features including mathematical equations, tables, and formatting. You can cite papers using the citation search feature, and they will be automatically added to your \texttt{cite.bib} file.

\section{Mathematical Expressions}
Here are some example equations:

\subsection{Quadratic Formula}
The quadratic formula is given by:
\[x = \frac{-b \pm \sqrt{b^2 - 4ac}}{2a}\]

\subsection{Matrix Example}
A sample matrix:
\[
\begin{pmatrix}
a_{11} & a_{12} & a_{13} \\
a_{21} & a_{22} & a_{23} \\
a_{31} & a_{32} & a_{33}
\end{pmatrix}
\]

\section{Tables}
Below is a sample table showing data:

\begin{table}[h]
\centering
\begin{tabular}{@{}lrrr@{}}
\toprule
\textbf{Category} & \textbf{Value 1} & \textbf{Value 2} & \textbf{Total} \\
\midrule
Item A & 10.5 & 20.3 & 30.8 \\
Item B & 15.2 & 25.1 & 40.3 \\
Item C & 22.1 & 12.7 & 34.8 \\
\bottomrule
\end{tabular}
\caption{Sample Data Table}
\label{tab:sample}
\end{table}

\section{Lists}
\subsection{Bullet Points}
\begin{itemize}
  \item First item
  \item Second item
  \item Third item with sub-items:
    \begin{itemize}
      \item Sub-item 1
      \item Sub-item 2
    \end{itemize}
\end{itemize}

\subsection{Numbered List}
\begin{enumerate}
  \item First step
  \item Second step
  \item Third step
\end{enumerate}

\section{Citations and References}
When you use the citation search feature, references will be automatically managed. For example, you might cite a paper like this: \cite{example2024}. All citations will be automatically added to your \texttt{cite.bib} file.

\section{Conclusion}
This template demonstrates basic \LaTeX{} features including mathematical formulas, tables, and lists. The bibliography system is now set up to work seamlessly with Bibby AI's citation management. Feel free to modify it for your needs!

\bibliographystyle{plain}
\bibliography{cite}

\end{document}\documentclass{article}
\usepackage[utf8]{inputenc}
\usepackage{booktabs}
\usepackage{amsmath}
\usepackage{natbib}

\title{My Document}
\author{Your Name}
\date{\today}

\begin{document}

\maketitle

\section{Introduction}
This is a sample document demonstrating various \LaTeX{} features including mathematical equations, tables, and formatting. You can cite papers using the citation search feature, and they will be automatically added to your \texttt{cite.bib} file.

\section{Mathematical Expressions}
Here are some example equations:

\subsection{Quadratic Formula}
The quadratic formula is given by:
\[x = \frac{-b \pm \sqrt{b^2 - 4ac}}{2a}\]

\subsection{Matrix Example}
A sample matrix:
\[
\begin{pmatrix}
a_{11} & a_{12} & a_{13} \\
a_{21} & a_{22} & a_{23} \\
a_{31} & a_{32} & a_{33}
\end{pmatrix}
\]

\section{Tables}
Below is a sample table showing data:

\begin{table}[h]
\centering
\begin{tabular}{@{}lrrr@{}}
\toprule
\textbf{Category} & \textbf{Value 1} & \textbf{Value 2} & \textbf{Total} \\
\midrule
Item A & 10.5 & 20.3 & 30.8 \\
Item B & 15.2 & 25.1 & 40.3 \\
Item C & 22.1 & 12.7 & 34.8 \\
\bottomrule
\end{tabular}
\caption{Sample Data Table}
\label{tab:sample}
\end{table}

\section{Lists}
\subsection{Bullet Points}
\begin{itemize}
  \item First item
  \item Second item
  \item Third item with sub-items:
    \begin{itemize}
      \item Sub-item 1
      \item Sub-item 2
    \end{itemize}
\end{itemize}

\subsection{Numbered List}
\begin{enumerate}
  \item First step
  \item Second step
  \item Third step
\end{enumerate}

\section{Citations and References}
When you use the citation search feature, references will be automatically managed. For example, you might cite a paper like this: \cite{example2024}. All citations will be automatically added to your \texttt{cite.bib} file.

\section{Conclusion}
This template demonstrates basic \LaTeX{} features including mathematical formulas, tables, and lists. The bibliography system is now set up to work seamlessly with Bibby AI's citation management. Feel free to modify it for your needs!

\bibliographystyle{plain}
\bibliography{cite}

\end{document}\documentclass{article}
\usepackage[utf8]{inputenc}
\usepackage{booktabs}
\usepackage{amsmath}
\usepackage{natbib}

\title{My Document}
\author{Your Name}
\date{\today}

\begin{document}

\maketitle

\section{Introduction}
This is a sample document demonstrating various \LaTeX{} features including mathematical equations, tables, and formatting. You can cite papers using the citation search feature, and they will be automatically added to your \texttt{cite.bib} file.

\section{Mathematical Expressions}
Here are some example equations:

\subsection{Quadratic Formula}
The quadratic formula is given by:
\[x = \frac{-b \pm \sqrt{b^2 - 4ac}}{2a}\]

\subsection{Matrix Example}
A sample matrix:
\[
\begin{pmatrix}
a_{11} & a_{12} & a_{13} \\
a_{21} & a_{22} & a_{23} \\
a_{31} & a_{32} & a_{33}
\end{pmatrix}
\]

\section{Tables}
Below is a sample table showing data:

\begin{table}[h]
\centering
\begin{tabular}{@{}lrrr@{}}
\toprule
\textbf{Category} & \textbf{Value 1} & \textbf{Value 2} & \textbf{Total} \\
\midrule
Item A & 10.5 & 20.3 & 30.8 \\
Item B & 15.2 & 25.1 & 40.3 \\
Item C & 22.1 & 12.7 & 34.8 \\
\bottomrule
\end{tabular}
\caption{Sample Data Table}
\label{tab:sample}
\end{table}

\section{Lists}
\subsection{Bullet Points}
\begin{itemize}
  \item First item
  \item Second item
  \item Third item with sub-items:
    \begin{itemize}
      \item Sub-item 1
      \item Sub-item 2
    \end{itemize}
\end{itemize}

\subsection{Numbered List}
\begin{enumerate}
  \item First step
  \item Second step
  \item Third step
\end{enumerate}

\section{Citations and References}
When you use the citation search feature, references will be automatically managed. For example, you might cite a paper like this: \cite{example2024}. All citations will be automatically added to your \texttt{cite.bib} file.

\section{Conclusion}
This template demonstrates basic \LaTeX{} features including mathematical formulas, tables, and lists. The bibliography system is now set up to work seamlessly with Bibby AI's citation management. Feel free to modify it for your needs!

\bibliographystyle{plain}
\bibliography{cite}

\end{document}\documentclass{article}
\usepackage[utf8]{inputenc}
\usepackage{booktabs}
\usepackage{amsmath}
\usepackage{natbib}

\title{My Document}
\author{Your Name}
\date{\today}

\begin{document}

\maketitle

\section{Introduction}
This is a sample document demonstrating various \LaTeX{} features including mathematical equations, tables, and formatting. You can cite papers using the citation search feature, and they will be automatically added to your \texttt{cite.bib} file.

\section{Mathematical Expressions}
Here are some example equations:

\subsection{Quadratic Formula}
The quadratic formula is given by:
\[x = \frac{-b \pm \sqrt{b^2 - 4ac}}{2a}\]

\subsection{Matrix Example}
A sample matrix:
\[
\begin{pmatrix}
a_{11} & a_{12} & a_{13} \\
a_{21} & a_{22} & a_{23} \\
a_{31} & a_{32} & a_{33}
\end{pmatrix}
\]

\section{Tables}
Below is a sample table showing data:

\begin{table}[h]
\centering
\begin{tabular}{@{}lrrr@{}}
\toprule
\textbf{Category} & \textbf{Value 1} & \textbf{Value 2} & \textbf{Total} \\
\midrule
Item A & 10.5 & 20.3 & 30.8 \\
Item B & 15.2 & 25.1 & 40.3 \\
Item C & 22.1 & 12.7 & 34.8 \\
\bottomrule
\end{tabular}
\caption{Sample Data Table}
\label{tab:sample}
\end{table}

\section{Lists}
\subsection{Bullet Points}
\begin{itemize}
  \item First item
  \item Second item
  \item Third item with sub-items:
    \begin{itemize}
      \item Sub-item 1
      \item Sub-item 2
    \end{itemize}
\end{itemize}

\subsection{Numbered List}
\begin{enumerate}
  \item First step
  \item Second step
  \item Third step
\end{enumerate}

\section{Citations and References}
When you use the citation search feature, references will be automatically managed. For example, you might cite a paper like this: \cite{example2024}. All citations will be automatically added to your \texttt{cite.bib} file.

\section{Conclusion}
This template demonstrates basic \LaTeX{} features including mathematical formulas, tables, and lists. The bibliography system is now set up to work seamlessly with Bibby AI's citation management. Feel free to modify it for your needs!

\bibliographystyle{plain}
\bibliography{cite}

\end{document}\documentclass{article}
\usepackage[utf8]{inputenc}
\usepackage{booktabs}
\usepackage{amsmath}
\usepackage{natbib}

\title{My Document}
\author{Your Name}
\date{\today}

\begin{document}

\maketitle

\section{Introduction}
This is a sample document demonstrating various \LaTeX{} features including mathematical equations, tables, and formatting. You can cite papers using the citation search feature, and they will be automatically added to your \texttt{cite.bib} file.

\section{Mathematical Expressions}
Here are some example equations:

\subsection{Quadratic Formula}
The quadratic formula is given by:
\[x = \frac{-b \pm \sqrt{b^2 - 4ac}}{2a}\]

\subsection{Matrix Example}
A sample matrix:
\[
\begin{pmatrix}
a_{11} & a_{12} & a_{13} \\
a_{21} & a_{22} & a_{23} \\
a_{31} & a_{32} & a_{33}
\end{pmatrix}
\]

\section{Tables}
Below is a sample table showing data:

\begin{table}[h]
\centering
\begin{tabular}{@{}lrrr@{}}
\toprule
\textbf{Category} & \textbf{Value 1} & \textbf{Value 2} & \textbf{Total} \\
\midrule
Item A & 10.5 & 20.3 & 30.8 \\
Item B & 15.2 & 25.1 & 40.3 \\
Item C & 22.1 & 12.7 & 34.8 \\
\bottomrule
\end{tabular}
\caption{Sample Data Table}
\label{tab:sample}
\end{table}

\section{Lists}
\subsection{Bullet Points}
\begin{itemize}
  \item First item
  \item Second item
  \item Third item with sub-items:
    \begin{itemize}
      \item Sub-item 1
      \item Sub-item 2
    \end{itemize}
\end{itemize}

\subsection{Numbered List}
\begin{enumerate}
  \item First step
  \item Second step
  \item Third step
\end{enumerate}

\section{Citations and References}
When you use the citation search feature, references will be automatically managed. For example, you might cite a paper like this: \cite{example2024}. All citations will be automatically added to your \texttt{cite.bib} file.

\section{Conclusion}
This template demonstrates basic \LaTeX{} features including mathematical formulas, tables, and lists. The bibliography system is now set up to work seamlessly with Bibby AI's citation management. Feel free to modify it for your needs!

\bibliographystyle{plain}
\bibliography{cite}

\end{document}\documentclass{article}
\usepackage[utf8]{inputenc}
\usepackage{booktabs}
\usepackage{amsmath}
\usepackage{natbib}

\title{My Document}
\author{Your Name}
\date{\today}

\begin{document}

\maketitle

\section{Introduction}
This is a sample document demonstrating various \LaTeX{} features including mathematical equations, tables, and formatting. You can cite papers using the citation search feature, and they will be automatically added to your \texttt{cite.bib} file.

\section{Mathematical Expressions}
Here are some example equations:

\subsection{Quadratic Formula}
The quadratic formula is given by:
\[x = \frac{-b \pm \sqrt{b^2 - 4ac}}{2a}\]

\subsection{Matrix Example}
A sample matrix:
\[
\begin{pmatrix}
a_{11} & a_{12} & a_{13} \\
a_{21} & a_{22} & a_{23} \\
a_{31} & a_{32} & a_{33}
\end{pmatrix}
\]

\section{Tables}
Below is a sample table showing data:

\begin{table}[h]
\centering
\begin{tabular}{@{}lrrr@{}}
\toprule
\textbf{Category} & \textbf{Value 1} & \textbf{Value 2} & \textbf{Total} \\
\midrule
Item A & 10.5 & 20.3 & 30.8 \\
Item B & 15.2 & 25.1 & 40.3 \\
Item C & 22.1 & 12.7 & 34.8 \\
\bottomrule
\end{tabular}
\caption{Sample Data Table}
\label{tab:sample}
\end{table}

\section{Lists}
\subsection{Bullet Points}
\begin{itemize}
  \item First item
  \item Second item
  \item Third item with sub-items:
    \begin{itemize}
      \item Sub-item 1
      \item Sub-item 2
    \end{itemize}
\end{itemize}

\subsection{Numbered List}
\begin{enumerate}
  \item First step
  \item Second step
  \item Third step
\end{enumerate}

\section{Citations and References}
When you use the citation search feature, references will be automatically managed. For example, you might cite a paper like this: \cite{example2024}. All citations will be automatically added to your \texttt{cite.bib} file.

\section{Conclusion}
This template demonstrates basic \LaTeX{} features including mathematical formulas, tables, and lists. The bibliography system is now set up to work seamlessly with Bibby AI's citation management. Feel free to modify it for your needs!

\bibliographystyle{plain}
\bibliography{cite}

\end{document}\documentclass{article}
\usepackage[utf8]{inputenc}
\usepackage{booktabs}
\usepackage{amsmath}
\usepackage{natbib}

\title{My Document}
\author{Your Name}
\date{\today}

\begin{document}

\maketitle

\section{Introduction}
This is a sample document demonstrating various \LaTeX{} features including mathematical equations, tables, and formatting. You can cite papers using the citation search feature, and they will be automatically added to your \texttt{cite.bib} file.

\section{Mathematical Expressions}
Here are some example equations:

\subsection{Quadratic Formula}
The quadratic formula is given by:
\[x = \frac{-b \pm \sqrt{b^2 - 4ac}}{2a}\]

\subsection{Matrix Example}
A sample matrix:
\[
\begin{pmatrix}
a_{11} & a_{12} & a_{13} \\
a_{21} & a_{22} & a_{23} \\
a_{31} & a_{32} & a_{33}
\end{pmatrix}
\]

\section{Tables}
Below is a sample table showing data:

\begin{table}[h]
\centering
\begin{tabular}{@{}lrrr@{}}
\toprule
\textbf{Category} & \textbf{Value 1} & \textbf{Value 2} & \textbf{Total} \\
\midrule
Item A & 10.5 & 20.3 & 30.8 \\
Item B & 15.2 & 25.1 & 40.3 \\
Item C & 22.1 & 12.7 & 34.8 \\
\bottomrule
\end{tabular}
\caption{Sample Data Table}
\label{tab:sample}
\end{table}

\section{Lists}
\subsection{Bullet Points}
\begin{itemize}
  \item First item
  \item Second item
  \item Third item with sub-items:
    \begin{itemize}
      \item Sub-item 1
      \item Sub-item 2
    \end{itemize}
\end{itemize}

\subsection{Numbered List}
\begin{enumerate}
  \item First step
  \item Second step
  \item Third step
\end{enumerate}

\section{Citations and References}
When you use the citation search feature, references will be automatically managed. For example, you might cite a paper like this: \cite{example2024}. All citations will be automatically added to your \texttt{cite.bib} file.

\section{Conclusion}
This template demonstrates basic \LaTeX{} features including mathematical formulas, tables, and lists. The bibliography system is now set up to work seamlessly with Bibby AI's citation management. Feel free to modify it for your needs!

\bibliographystyle{plain}
\bibliography{cite}

\end{document}\documentclass{article}
\usepackage[utf8]{inputenc}
\usepackage{booktabs}
\usepackage{amsmath}
\usepackage{natbib}

\title{My Document}
\author{Your Name}
\date{\today}

\begin{document}

\maketitle

\section{Introduction}
This is a sample document demonstrating various \LaTeX{} features including mathematical equations, tables, and formatting. You can cite papers using the citation search feature, and they will be automatically added to your \texttt{cite.bib} file.

\section{Mathematical Expressions}
Here are some example equations:

\subsection{Quadratic Formula}
The quadratic formula is given by:
\[x = \frac{-b \pm \sqrt{b^2 - 4ac}}{2a}\]

\subsection{Matrix Example}
A sample matrix:
\[
\begin{pmatrix}
a_{11} & a_{12} & a_{13} \\
a_{21} & a_{22} & a_{23} \\
a_{31} & a_{32} & a_{33}
\end{pmatrix}
\]

\section{Tables}
Below is a sample table showing data:

\begin{table}[h]
\centering
\begin{tabular}{@{}lrrr@{}}
\toprule
\textbf{Category} & \textbf{Value 1} & \textbf{Value 2} & \textbf{Total} \\
\midrule
Item A & 10.5 & 20.3 & 30.8 \\
Item B & 15.2 & 25.1 & 40.3 \\
Item C & 22.1 & 12.7 & 34.8 \\
\bottomrule
\end{tabular}
\caption{Sample Data Table}
\label{tab:sample}
\end{table}

\section{Lists}
\subsection{Bullet Points}
\begin{itemize}
  \item First item
  \item Second item
  \item Third item with sub-items:
    \begin{itemize}
      \item Sub-item 1
      \item Sub-item 2
    \end{itemize}
\end{itemize}

\subsection{Numbered List}
\begin{enumerate}
  \item First step
  \item Second step
  \item Third step
\end{enumerate}

\section{Citations and References}
When you use the citation search feature, references will be automatically managed. For example, you might cite a paper like this: \cite{example2024}. All citations will be automatically added to your \texttt{cite.bib} file.

\section{Conclusion}
This template demonstrates basic \LaTeX{} features including mathematical formulas, tables, and lists. The bibliography system is now set up to work seamlessly with Bibby AI's citation management. Feel free to modify it for your needs!

\bibliographystyle{plain}
\bibliography{cite}

\end{document}\documentclass{article}
\usepackage[utf8]{inputenc}
\usepackage{booktabs}
\usepackage{amsmath}
\usepackage{natbib}

\title{My Document}
\author{Your Name}
\date{\today}

\begin{document}

\maketitle

\section{Introduction}
This is a sample document demonstrating various \LaTeX{} features including mathematical equations, tables, and formatting. You can cite papers using the citation search feature, and they will be automatically added to your \texttt{cite.bib} file.

\section{Mathematical Expressions}
Here are some example equations:

\subsection{Quadratic Formula}
The quadratic formula is given by:
\[x = \frac{-b \pm \sqrt{b^2 - 4ac}}{2a}\]

\subsection{Matrix Example}
A sample matrix:
\[
\begin{pmatrix}
a_{11} & a_{12} & a_{13} \\
a_{21} & a_{22} & a_{23} \\
a_{31} & a_{32} & a_{33}
\end{pmatrix}
\]

\section{Tables}
Below is a sample table showing data:

\begin{table}[h]
\centering
\begin{tabular}{@{}lrrr@{}}
\toprule
\textbf{Category} & \textbf{Value 1} & \textbf{Value 2} & \textbf{Total} \\
\midrule
Item A & 10.5 & 20.3 & 30.8 \\
Item B & 15.2 & 25.1 & 40.3 \\
Item C & 22.1 & 12.7 & 34.8 \\
\bottomrule
\end{tabular}
\caption{Sample Data Table}
\label{tab:sample}
\end{table}

\section{Lists}
\subsection{Bullet Points}
\begin{itemize}
  \item First item
  \item Second item
  \item Third item with sub-items:
    \begin{itemize}
      \item Sub-item 1
      \item Sub-item 2
    \end{itemize}
\end{itemize}

\subsection{Numbered List}
\begin{enumerate}
  \item First step
  \item Second step
  \item Third step
\end{enumerate}

\section{Citations and References}
When you use the citation search feature, references will be automatically managed. For example, you might cite a paper like this: \cite{example2024}. All citations will be automatically added to your \texttt{cite.bib} file.

\section{Conclusion}
This template demonstrates basic \LaTeX{} features including mathematical formulas, tables, and lists. The bibliography system is now set up to work seamlessly with Bibby AI's citation management. Feel free to modify it for your needs!

\bibliographystyle{plain}
\bibliography{cite}

\end{document}\documentclass{article}
\usepackage[utf8]{inputenc}
\usepackage{booktabs}
\usepackage{amsmath}
\usepackage{natbib}

\title{My Document}
\author{Your Name}
\date{\today}

\begin{document}

\maketitle

\section{Introduction}
This is a sample document demonstrating various \LaTeX{} features including mathematical equations, tables, and formatting. You can cite papers using the citation search feature, and they will be automatically added to your \texttt{cite.bib} file.

\section{Mathematical Expressions}
Here are some example equations:

\subsection{Quadratic Formula}
The quadratic formula is given by:
\[x = \frac{-b \pm \sqrt{b^2 - 4ac}}{2a}\]

\subsection{Matrix Example}
A sample matrix:
\[
\begin{pmatrix}
a_{11} & a_{12} & a_{13} \\
a_{21} & a_{22} & a_{23} \\
a_{31} & a_{32} & a_{33}
\end{pmatrix}
\]

\section{Tables}
Below is a sample table showing data:

\begin{table}[h]
\centering
\begin{tabular}{@{}lrrr@{}}
\toprule
\textbf{Category} & \textbf{Value 1} & \textbf{Value 2} & \textbf{Total} \\
\midrule
Item A & 10.5 & 20.3 & 30.8 \\
Item B & 15.2 & 25.1 & 40.3 \\
Item C & 22.1 & 12.7 & 34.8 \\
\bottomrule
\end{tabular}
\caption{Sample Data Table}
\label{tab:sample}
\end{table}

\section{Lists}
\subsection{Bullet Points}
\begin{itemize}
  \item First item
  \item Second item
  \item Third item with sub-items:
    \begin{itemize}
      \item Sub-item 1
      \item Sub-item 2
    \end{itemize}
\end{itemize}

\subsection{Numbered List}
\begin{enumerate}
  \item First step
  \item Second step
  \item Third step
\end{enumerate}

\section{Citations and References}
When you use the citation search feature, references will be automatically managed. For example, you might cite a paper like this: \cite{example2024}. All citations will be automatically added to your \texttt{cite.bib} file.

\section{Conclusion}
This template demonstrates basic \LaTeX{} features including mathematical formulas, tables, and lists. The bibliography system is now set up to work seamlessly with Bibby AI's citation management. Feel free to modify it for your needs!

\bibliographystyle{plain}
\bibliography{cite}

\end{document}\documentclass{article}
\usepackage[utf8]{inputenc}
\usepackage{booktabs}
\usepackage{amsmath}
\usepackage{natbib}

\title{My Document}
\author{Your Name}
\date{\today}

\begin{document}

\maketitle

\section{Introduction}
This is a sample document demonstrating various \LaTeX{} features including mathematical equations, tables, and formatting. You can cite papers using the citation search feature, and they will be automatically added to your \texttt{cite.bib} file.

\section{Mathematical Expressions}
Here are some example equations:

\subsection{Quadratic Formula}
The quadratic formula is given by:
\[x = \frac{-b \pm \sqrt{b^2 - 4ac}}{2a}\]

\subsection{Matrix Example}
A sample matrix:
\[
\begin{pmatrix}
a_{11} & a_{12} & a_{13} \\
a_{21} & a_{22} & a_{23} \\
a_{31} & a_{32} & a_{33}
\end{pmatrix}
\]

\section{Tables}
Below is a sample table showing data:

\begin{table}[h]
\centering
\begin{tabular}{@{}lrrr@{}}
\toprule
\textbf{Category} & \textbf{Value 1} & \textbf{Value 2} & \textbf{Total} \\
\midrule
Item A & 10.5 & 20.3 & 30.8 \\
Item B & 15.2 & 25.1 & 40.3 \\
Item C & 22.1 & 12.7 & 34.8 \\
\bottomrule
\end{tabular}
\caption{Sample Data Table}
\label{tab:sample}
\end{table}

\section{Lists}
\subsection{Bullet Points}
\begin{itemize}
  \item First item
  \item Second item
  \item Third item with sub-items:
    \begin{itemize}
      \item Sub-item 1
      \item Sub-item 2
    \end{itemize}
\end{itemize}

\subsection{Numbered List}
\begin{enumerate}
  \item First step
  \item Second step
  \item Third step
\end{enumerate}

\section{Citations and References}
When you use the citation search feature, references will be automatically managed. For example, you might cite a paper like this: \cite{example2024}. All citations will be automatically added to your \texttt{cite.bib} file.

\section{Conclusion}
This template demonstrates basic \LaTeX{} features including mathematical formulas, tables, and lists. The bibliography system is now set up to work seamlessly with Bibby AI's citation management. Feel free to modify it for your needs!

\bibliographystyle{plain}
\bibliography{cite}

\end{document}\documentclass{article}
\usepackage[utf8]{inputenc}
\usepackage{booktabs}
\usepackage{amsmath}
\usepackage{natbib}

\title{My Document}
\author{Your Name}
\date{\today}

\begin{document}

\maketitle

\section{Introduction}
This is a sample document demonstrating various \LaTeX{} features including mathematical equations, tables, and formatting. You can cite papers using the citation search feature, and they will be automatically added to your \texttt{cite.bib} file.

\section{Mathematical Expressions}
Here are some example equations:

\subsection{Quadratic Formula}
The quadratic formula is given by:
\[x = \frac{-b \pm \sqrt{b^2 - 4ac}}{2a}\]

\subsection{Matrix Example}
A sample matrix:
\[
\begin{pmatrix}
a_{11} & a_{12} & a_{13} \\
a_{21} & a_{22} & a_{23} \\
a_{31} & a_{32} & a_{33}
\end{pmatrix}
\]

\section{Tables}
Below is a sample table showing data:

\begin{table}[h]
\centering
\begin{tabular}{@{}lrrr@{}}
\toprule
\textbf{Category} & \textbf{Value 1} & \textbf{Value 2} & \textbf{Total} \\
\midrule
Item A & 10.5 & 20.3 & 30.8 \\
Item B & 15.2 & 25.1 & 40.3 \\
Item C & 22.1 & 12.7 & 34.8 \\
\bottomrule
\end{tabular}
\caption{Sample Data Table}
\label{tab:sample}
\end{table}

\section{Lists}
\subsection{Bullet Points}
\begin{itemize}
  \item First item
  \item Second item
  \item Third item with sub-items:
    \begin{itemize}
      \item Sub-item 1
      \item Sub-item 2
    \end{itemize}
\end{itemize}

\subsection{Numbered List}
\begin{enumerate}
  \item First step
  \item Second step
  \item Third step
\end{enumerate}

\section{Citations and References}
When you use the citation search feature, references will be automatically managed. For example, you might cite a paper like this: \cite{example2024}. All citations will be automatically added to your \texttt{cite.bib} file.

\section{Conclusion}
This template demonstrates basic \LaTeX{} features including mathematical formulas, tables, and lists. The bibliography system is now set up to work seamlessly with Bibby AI's citation management. Feel free to modify it for your needs!

\bibliographystyle{plain}
\bibliography{cite}

\end{document}\documentclass{article}
\usepackage[utf8]{inputenc}
\usepackage{booktabs}
\usepackage{amsmath}
\usepackage{natbib}

\title{My Document}
\author{Your Name}
\date{\today}

\begin{document}

\maketitle

\section{Introduction}
This is a sample document demonstrating various \LaTeX{} features including mathematical equations, tables, and formatting. You can cite papers using the citation search feature, and they will be automatically added to your \texttt{cite.bib} file.

\section{Mathematical Expressions}
Here are some example equations:

\subsection{Quadratic Formula}
The quadratic formula is given by:
\[x = \frac{-b \pm \sqrt{b^2 - 4ac}}{2a}\]

\subsection{Matrix Example}
A sample matrix:
\[
\begin{pmatrix}
a_{11} & a_{12} & a_{13} \\
a_{21} & a_{22} & a_{23} \\
a_{31} & a_{32} & a_{33}
\end{pmatrix}
\]

\section{Tables}
Below is a sample table showing data:

\begin{table}[h]
\centering
\begin{tabular}{@{}lrrr@{}}
\toprule
\textbf{Category} & \textbf{Value 1} & \textbf{Value 2} & \textbf{Total} \\
\midrule
Item A & 10.5 & 20.3 & 30.8 \\
Item B & 15.2 & 25.1 & 40.3 \\
Item C & 22.1 & 12.7 & 34.8 \\
\bottomrule
\end{tabular}
\caption{Sample Data Table}
\label{tab:sample}
\end{table}

\section{Lists}
\subsection{Bullet Points}
\begin{itemize}
  \item First item
  \item Second item
  \item Third item with sub-items:
    \begin{itemize}
      \item Sub-item 1
      \item Sub-item 2
    \end{itemize}
\end{itemize}

\subsection{Numbered List}
\begin{enumerate}
  \item First step
  \item Second step
  \item Third step
\end{enumerate}

\section{Citations and References}
When you use the citation search feature, references will be automatically managed. For example, you might cite a paper like this: \cite{example2024}. All citations will be automatically added to your \texttt{cite.bib} file.

\section{Conclusion}
This template demonstrates basic \LaTeX{} features including mathematical formulas, tables, and lists. The bibliography system is now set up to work seamlessly with Bibby AI's citation management. Feel free to modify it for your needs!

\bibliographystyle{plain}
\bibliography{cite}

\end{document}\documentclass{article}
\usepackage[utf8]{inputenc}
\usepackage{booktabs}
\usepackage{amsmath}
\usepackage{natbib}

\title{My Document}
\author{Your Name}
\date{\today}

\begin{document}

\maketitle

\section{Introduction}
This is a sample document demonstrating various \LaTeX{} features including mathematical equations, tables, and formatting. You can cite papers using the citation search feature, and they will be automatically added to your \texttt{cite.bib} file.

\section{Mathematical Expressions}
Here are some example equations:

\subsection{Quadratic Formula}
The quadratic formula is given by:
\[x = \frac{-b \pm \sqrt{b^2 - 4ac}}{2a}\]

\subsection{Matrix Example}
A sample matrix:
\[
\begin{pmatrix}
a_{11} & a_{12} & a_{13} \\
a_{21} & a_{22} & a_{23} \\
a_{31} & a_{32} & a_{33}
\end{pmatrix}
\]

\section{Tables}
Below is a sample table showing data:

\begin{table}[h]
\centering
\begin{tabular}{@{}lrrr@{}}
\toprule
\textbf{Category} & \textbf{Value 1} & \textbf{Value 2} & \textbf{Total} \\
\midrule
Item A & 10.5 & 20.3 & 30.8 \\
Item B & 15.2 & 25.1 & 40.3 \\
Item C & 22.1 & 12.7 & 34.8 \\
\bottomrule
\end{tabular}
\caption{Sample Data Table}
\label{tab:sample}
\end{table}

\section{Lists}
\subsection{Bullet Points}
\begin{itemize}
  \item First item
  \item Second item
  \item Third item with sub-items:
    \begin{itemize}
      \item Sub-item 1
      \item Sub-item 2
    \end{itemize}
\end{itemize}

\subsection{Numbered List}
\begin{enumerate}
  \item First step
  \item Second step
  \item Third step
\end{enumerate}

\section{Citations and References}
When you use the citation search feature, references will be automatically managed. For example, you might cite a paper like this: \cite{example2024}. All citations will be automatically added to your \texttt{cite.bib} file.

\section{Conclusion}
This template demonstrates basic \LaTeX{} features including mathematical formulas, tables, and lists. The bibliography system is now set up to work seamlessly with Bibby AI's citation management. Feel free to modify it for your needs!

\bibliographystyle{plain}
\bibliography{cite}

\end{document}\documentclass{article}
\usepackage[utf8]{inputenc}
\usepackage{booktabs}
\usepackage{amsmath}
\usepackage{natbib}

\title{My Document}
\author{Your Name}
\date{\today}

\begin{document}

\maketitle

\section{Introduction}
This is a sample document demonstrating various \LaTeX{} features including mathematical equations, tables, and formatting. You can cite papers using the citation search feature, and they will be automatically added to your \texttt{cite.bib} file.

\section{Mathematical Expressions}
Here are some example equations:

\subsection{Quadratic Formula}
The quadratic formula is given by:
\[x = \frac{-b \pm \sqrt{b^2 - 4ac}}{2a}\]

\subsection{Matrix Example}
A sample matrix:
\[
\begin{pmatrix}
a_{11} & a_{12} & a_{13} \\
a_{21} & a_{22} & a_{23} \\
a_{31} & a_{32} & a_{33}
\end{pmatrix}
\]

\section{Tables}
Below is a sample table showing data:

\begin{table}[h]
\centering
\begin{tabular}{@{}lrrr@{}}
\toprule
\textbf{Category} & \textbf{Value 1} & \textbf{Value 2} & \textbf{Total} \\
\midrule
Item A & 10.5 & 20.3 & 30.8 \\
Item B & 15.2 & 25.1 & 40.3 \\
Item C & 22.1 & 12.7 & 34.8 \\
\bottomrule
\end{tabular}
\caption{Sample Data Table}
\label{tab:sample}
\end{table}

\section{Lists}
\subsection{Bullet Points}
\begin{itemize}
  \item First item
  \item Second item
  \item Third item with sub-items:
    \begin{itemize}
      \item Sub-item 1
      \item Sub-item 2
    \end{itemize}
\end{itemize}

\subsection{Numbered List}
\begin{enumerate}
  \item First step
  \item Second step
  \item Third step
\end{enumerate}

\section{Citations and References}
When you use the citation search feature, references will be automatically managed. For example, you might cite a paper like this: \cite{example2024}. All citations will be automatically added to your \texttt{cite.bib} file.

\section{Conclusion}
This template demonstrates basic \LaTeX{} features including mathematical formulas, tables, and lists. The bibliography system is now set up to work seamlessly with Bibby AI's citation management. Feel free to modify it for your needs!

\bibliographystyle{plain}
\bibliography{cite}

\end{document}\documentclass{article}
\usepackage[utf8]{inputenc}
\usepackage{booktabs}
\usepackage{amsmath}
\usepackage{natbib}

\title{My Document}
\author{Your Name}
\date{\today}

\begin{document}

\maketitle

\section{Introduction}
This is a sample document demonstrating various \LaTeX{} features including mathematical equations, tables, and formatting. You can cite papers using the citation search feature, and they will be automatically added to your \texttt{cite.bib} file.

\section{Mathematical Expressions}
Here are some example equations:

\subsection{Quadratic Formula}
The quadratic formula is given by:
\[x = \frac{-b \pm \sqrt{b^2 - 4ac}}{2a}\]

\subsection{Matrix Example}
A sample matrix:
\[
\begin{pmatrix}
a_{11} & a_{12} & a_{13} \\
a_{21} & a_{22} & a_{23} \\
a_{31} & a_{32} & a_{33}
\end{pmatrix}
\]

\section{Tables}
Below is a sample table showing data:

\begin{table}[h]
\centering
\begin{tabular}{@{}lrrr@{}}
\toprule
\textbf{Category} & \textbf{Value 1} & \textbf{Value 2} & \textbf{Total} \\
\midrule
Item A & 10.5 & 20.3 & 30.8 \\
Item B & 15.2 & 25.1 & 40.3 \\
Item C & 22.1 & 12.7 & 34.8 \\
\bottomrule
\end{tabular}
\caption{Sample Data Table}
\label{tab:sample}
\end{table}

\section{Lists}
\subsection{Bullet Points}
\begin{itemize}
  \item First item
  \item Second item
  \item Third item with sub-items:
    \begin{itemize}
      \item Sub-item 1
      \item Sub-item 2
    \end{itemize}
\end{itemize}

\subsection{Numbered List}
\begin{enumerate}
  \item First step
  \item Second step
  \item Third step
\end{enumerate}

\section{Citations and References}
When you use the citation search feature, references will be automatically managed. For example, you might cite a paper like this: \cite{example2024}. All citations will be automatically added to your \texttt{cite.bib} file.

\section{Conclusion}
This template demonstrates basic \LaTeX{} features including mathematical formulas, tables, and lists. The bibliography system is now set up to work seamlessly with Bibby AI's citation management. Feel free to modify it for your needs!

\bibliographystyle{plain}
\bibliography{cite}

\end{document}\documentclass{article}
\usepackage[utf8]{inputenc}
\usepackage{booktabs}
\usepackage{amsmath}
\usepackage{natbib}

\title{My Document}
\author{Your Name}
\date{\today}

\begin{document}

\maketitle

\section{Introduction}
This is a sample document demonstrating various \LaTeX{} features including mathematical equations, tables, and formatting. You can cite papers using the citation search feature, and they will be automatically added to your \texttt{cite.bib} file.

\section{Mathematical Expressions}
Here are some example equations:

\subsection{Quadratic Formula}
The quadratic formula is given by:
\[x = \frac{-b \pm \sqrt{b^2 - 4ac}}{2a}\]

\subsection{Matrix Example}
A sample matrix:
\[
\begin{pmatrix}
a_{11} & a_{12} & a_{13} \\
a_{21} & a_{22} & a_{23} \\
a_{31} & a_{32} & a_{33}
\end{pmatrix}
\]

\section{Tables}
Below is a sample table showing data:

\begin{table}[h]
\centering
\begin{tabular}{@{}lrrr@{}}
\toprule
\textbf{Category} & \textbf{Value 1} & \textbf{Value 2} & \textbf{Total} \\
\midrule
Item A & 10.5 & 20.3 & 30.8 \\
Item B & 15.2 & 25.1 & 40.3 \\
Item C & 22.1 & 12.7 & 34.8 \\
\bottomrule
\end{tabular}
\caption{Sample Data Table}
\label{tab:sample}
\end{table}

\section{Lists}
\subsection{Bullet Points}
\begin{itemize}
  \item First item
  \item Second item
  \item Third item with sub-items:
    \begin{itemize}
      \item Sub-item 1
      \item Sub-item 2
    \end{itemize}
\end{itemize}

\subsection{Numbered List}
\begin{enumerate}
  \item First step
  \item Second step
  \item Third step
\end{enumerate}

\section{Citations and References}
When you use the citation search feature, references will be automatically managed. For example, you might cite a paper like this: \cite{example2024}. All citations will be automatically added to your \texttt{cite.bib} file.

\section{Conclusion}
This template demonstrates basic \LaTeX{} features including mathematical formulas, tables, and lists. The bibliography system is now set up to work seamlessly with Bibby AI's citation management. Feel free to modify it for your needs!

\bibliographystyle{plain}
\bibliography{cite}

\end{document}\documentclass{article}
\usepackage[utf8]{inputenc}
\usepackage{booktabs}
\usepackage{amsmath}
\usepackage{natbib}

\title{My Document}
\author{Your Name}
\date{\today}

\begin{document}

\maketitle

\section{Introduction}
This is a sample document demonstrating various \LaTeX{} features including mathematical equations, tables, and formatting. You can cite papers using the citation search feature, and they will be automatically added to your \texttt{cite.bib} file.

\section{Mathematical Expressions}
Here are some example equations:

\subsection{Quadratic Formula}
The quadratic formula is given by:
\[x = \frac{-b \pm \sqrt{b^2 - 4ac}}{2a}\]

\subsection{Matrix Example}
A sample matrix:
\[
\begin{pmatrix}
a_{11} & a_{12} & a_{13} \\
a_{21} & a_{22} & a_{23} \\
a_{31} & a_{32} & a_{33}
\end{pmatrix}
\]

\section{Tables}
Below is a sample table showing data:

\begin{table}[h]
\centering
\begin{tabular}{@{}lrrr@{}}
\toprule
\textbf{Category} & \textbf{Value 1} & \textbf{Value 2} & \textbf{Total} \\
\midrule
Item A & 10.5 & 20.3 & 30.8 \\
Item B & 15.2 & 25.1 & 40.3 \\
Item C & 22.1 & 12.7 & 34.8 \\
\bottomrule
\end{tabular}
\caption{Sample Data Table}
\label{tab:sample}
\end{table}

\section{Lists}
\subsection{Bullet Points}
\begin{itemize}
  \item First item
  \item Second item
  \item Third item with sub-items:
    \begin{itemize}
      \item Sub-item 1
      \item Sub-item 2
    \end{itemize}
\end{itemize}

\subsection{Numbered List}
\begin{enumerate}
  \item First step
  \item Second step
  \item Third step
\end{enumerate}

\section{Citations and References}
When you use the citation search feature, references will be automatically managed. For example, you might cite a paper like this: \cite{example2024}. All citations will be automatically added to your \texttt{cite.bib} file.

\section{Conclusion}
This template demonstrates basic \LaTeX{} features including mathematical formulas, tables, and lists. The bibliography system is now set up to work seamlessly with Bibby AI's citation management. Feel free to modify it for your needs!

\bibliographystyle{plain}
\bibliography{cite}

\end{document}\documentclass{article}
\usepackage[utf8]{inputenc}
\usepackage{booktabs}
\usepackage{amsmath}
\usepackage{natbib}

\title{My Document}
\author{Your Name}
\date{\today}

\begin{document}

\maketitle

\section{Introduction}
This is a sample document demonstrating various \LaTeX{} features including mathematical equations, tables, and formatting. You can cite papers using the citation search feature, and they will be automatically added to your \texttt{cite.bib} file.

\section{Mathematical Expressions}
Here are some example equations:

\subsection{Quadratic Formula}
The quadratic formula is given by:
\[x = \frac{-b \pm \sqrt{b^2 - 4ac}}{2a}\]

\subsection{Matrix Example}
A sample matrix:
\[
\begin{pmatrix}
a_{11} & a_{12} & a_{13} \\
a_{21} & a_{22} & a_{23} \\
a_{31} & a_{32} & a_{33}
\end{pmatrix}
\]

\section{Tables}
Below is a sample table showing data:

\begin{table}[h]
\centering
\begin{tabular}{@{}lrrr@{}}
\toprule
\textbf{Category} & \textbf{Value 1} & \textbf{Value 2} & \textbf{Total} \\
\midrule
Item A & 10.5 & 20.3 & 30.8 \\
Item B & 15.2 & 25.1 & 40.3 \\
Item C & 22.1 & 12.7 & 34.8 \\
\bottomrule
\end{tabular}
\caption{Sample Data Table}
\label{tab:sample}
\end{table}

\section{Lists}
\subsection{Bullet Points}
\begin{itemize}
  \item First item
  \item Second item
  \item Third item with sub-items:
    \begin{itemize}
      \item Sub-item 1
      \item Sub-item 2
    \end{itemize}
\end{itemize}

\subsection{Numbered List}
\begin{enumerate}
  \item First step
  \item Second step
  \item Third step
\end{enumerate}

\section{Citations and References}
When you use the citation search feature, references will be automatically managed. For example, you might cite a paper like this: \cite{example2024}. All citations will be automatically added to your \texttt{cite.bib} file.

\section{Conclusion}
This template demonstrates basic \LaTeX{} features including mathematical formulas, tables, and lists. The bibliography system is now set up to work seamlessly with Bibby AI's citation management. Feel free to modify it for your needs!

\bibliographystyle{plain}
\bibliography{cite}

\end{document}\documentclass{article}
\usepackage[utf8]{inputenc}
\usepackage{booktabs}
\usepackage{amsmath}
\usepackage{natbib}

\title{My Document}
\author{Your Name}
\date{\today}

\begin{document}

\maketitle

\section{Introduction}
This is a sample document demonstrating various \LaTeX{} features including mathematical equations, tables, and formatting. You can cite papers using the citation search feature, and they will be automatically added to your \texttt{cite.bib} file.

\section{Mathematical Expressions}
Here are some example equations:

\subsection{Quadratic Formula}
The quadratic formula is given by:
\[x = \frac{-b \pm \sqrt{b^2 - 4ac}}{2a}\]

\subsection{Matrix Example}
A sample matrix:
\[
\begin{pmatrix}
a_{11} & a_{12} & a_{13} \\
a_{21} & a_{22} & a_{23} \\
a_{31} & a_{32} & a_{33}
\end{pmatrix}
\]

\section{Tables}
Below is a sample table showing data:

\begin{table}[h]
\centering
\begin{tabular}{@{}lrrr@{}}
\toprule
\textbf{Category} & \textbf{Value 1} & \textbf{Value 2} & \textbf{Total} \\
\midrule
Item A & 10.5 & 20.3 & 30.8 \\
Item B & 15.2 & 25.1 & 40.3 \\
Item C & 22.1 & 12.7 & 34.8 \\
\bottomrule
\end{tabular}
\caption{Sample Data Table}
\label{tab:sample}
\end{table}

\section{Lists}
\subsection{Bullet Points}
\begin{itemize}
  \item First item
  \item Second item
  \item Third item with sub-items:
    \begin{itemize}
      \item Sub-item 1
      \item Sub-item 2
    \end{itemize}
\end{itemize}

\subsection{Numbered List}
\begin{enumerate}
  \item First step
  \item Second step
  \item Third step
\end{enumerate}

\section{Citations and References}
When you use the citation search feature, references will be automatically managed. For example, you might cite a paper like this: \cite{example2024}. All citations will be automatically added to your \texttt{cite.bib} file.

\section{Conclusion}
This template demonstrates basic \LaTeX{} features including mathematical formulas, tables, and lists. The bibliography system is now set up to work seamlessly with Bibby AI's citation management. Feel free to modify it for your needs!

\bibliographystyle{plain}
\bibliography{cite}

\end{document}\documentclass{article}
\usepackage[utf8]{inputenc}
\usepackage{booktabs}
\usepackage{amsmath}
\usepackage{natbib}

\title{My Document}
\author{Your Name}
\date{\today}

\begin{document}

\maketitle

\section{Introduction}
This is a sample document demonstrating various \LaTeX{} features including mathematical equations, tables, and formatting. You can cite papers using the citation search feature, and they will be automatically added to your \texttt{cite.bib} file.

\section{Mathematical Expressions}
Here are some example equations:

\subsection{Quadratic Formula}
The quadratic formula is given by:
\[x = \frac{-b \pm \sqrt{b^2 - 4ac}}{2a}\]

\subsection{Matrix Example}
A sample matrix:
\[
\begin{pmatrix}
a_{11} & a_{12} & a_{13} \\
a_{21} & a_{22} & a_{23} \\
a_{31} & a_{32} & a_{33}
\end{pmatrix}
\]

\section{Tables}
Below is a sample table showing data:

\begin{table}[h]
\centering
\begin{tabular}{@{}lrrr@{}}
\toprule
\textbf{Category} & \textbf{Value 1} & \textbf{Value 2} & \textbf{Total} \\
\midrule
Item A & 10.5 & 20.3 & 30.8 \\
Item B & 15.2 & 25.1 & 40.3 \\
Item C & 22.1 & 12.7 & 34.8 \\
\bottomrule
\end{tabular}
\caption{Sample Data Table}
\label{tab:sample}
\end{table}

\section{Lists}
\subsection{Bullet Points}
\begin{itemize}
  \item First item
  \item Second item
  \item Third item with sub-items:
    \begin{itemize}
      \item Sub-item 1
      \item Sub-item 2
    \end{itemize}
\end{itemize}

\subsection{Numbered List}
\begin{enumerate}
  \item First step
  \item Second step
  \item Third step
\end{enumerate}

\section{Citations and References}
When you use the citation search feature, references will be automatically managed. For example, you might cite a paper like this: \cite{example2024}. All citations will be automatically added to your \texttt{cite.bib} file.

\section{Conclusion}
This template demonstrates basic \LaTeX{} features including mathematical formulas, tables, and lists. The bibliography system is now set up to work seamlessly with Bibby AI's citation management. Feel free to modify it for your needs!

\bibliographystyle{plain}
\bibliography{cite}

\end{document}\documentclass{article}
\usepackage[utf8]{inputenc}
\usepackage{booktabs}
\usepackage{amsmath}
\usepackage{natbib}

\title{My Document}
\author{Your Name}
\date{\today}

\begin{document}

\maketitle

\section{Introduction}
This is a sample document demonstrating various \LaTeX{} features including mathematical equations, tables, and formatting. You can cite papers using the citation search feature, and they will be automatically added to your \texttt{cite.bib} file.

\section{Mathematical Expressions}
Here are some example equations:

\subsection{Quadratic Formula}
The quadratic formula is given by:
\[x = \frac{-b \pm \sqrt{b^2 - 4ac}}{2a}\]

\subsection{Matrix Example}
A sample matrix:
\[
\begin{pmatrix}
a_{11} & a_{12} & a_{13} \\
a_{21} & a_{22} & a_{23} \\
a_{31} & a_{32} & a_{33}
\end{pmatrix}
\]

\section{Tables}
Below is a sample table showing data:

\begin{table}[h]
\centering
\begin{tabular}{@{}lrrr@{}}
\toprule
\textbf{Category} & \textbf{Value 1} & \textbf{Value 2} & \textbf{Total} \\
\midrule
Item A & 10.5 & 20.3 & 30.8 \\
Item B & 15.2 & 25.1 & 40.3 \\
Item C & 22.1 & 12.7 & 34.8 \\
\bottomrule
\end{tabular}
\caption{Sample Data Table}
\label{tab:sample}
\end{table}

\section{Lists}
\subsection{Bullet Points}
\begin{itemize}
  \item First item
  \item Second item
  \item Third item with sub-items:
    \begin{itemize}
      \item Sub-item 1
      \item Sub-item 2
    \end{itemize}
\end{itemize}

\subsection{Numbered List}
\begin{enumerate}
  \item First step
  \item Second step
  \item Third step
\end{enumerate}

\section{Citations and References}
When you use the citation search feature, references will be automatically managed. For example, you might cite a paper like this: \cite{example2024}. All citations will be automatically added to your \texttt{cite.bib} file.

\section{Conclusion}
This template demonstrates basic \LaTeX{} features including mathematical formulas, tables, and lists. The bibliography system is now set up to work seamlessly with Bibby AI's citation management. Feel free to modify it for your needs!

\bibliographystyle{plain}
\bibliography{cite}

\end{document}\documentclass{article}
\usepackage[utf8]{inputenc}
\usepackage{booktabs}
\usepackage{amsmath}
\usepackage{natbib}

\title{My Document}
\author{Your Name}
\date{\today}

\begin{document}

\maketitle

\section{Introduction}
This is a sample document demonstrating various \LaTeX{} features including mathematical equations, tables, and formatting. You can cite papers using the citation search feature, and they will be automatically added to your \texttt{cite.bib} file.

\section{Mathematical Expressions}
Here are some example equations:

\subsection{Quadratic Formula}
The quadratic formula is given by:
\[x = \frac{-b \pm \sqrt{b^2 - 4ac}}{2a}\]

\subsection{Matrix Example}
A sample matrix:
\[
\begin{pmatrix}
a_{11} & a_{12} & a_{13} \\
a_{21} & a_{22} & a_{23} \\
a_{31} & a_{32} & a_{33}
\end{pmatrix}
\]

\section{Tables}
Below is a sample table showing data:

\begin{table}[h]
\centering
\begin{tabular}{@{}lrrr@{}}
\toprule
\textbf{Category} & \textbf{Value 1} & \textbf{Value 2} & \textbf{Total} \\
\midrule
Item A & 10.5 & 20.3 & 30.8 \\
Item B & 15.2 & 25.1 & 40.3 \\
Item C & 22.1 & 12.7 & 34.8 \\
\bottomrule
\end{tabular}
\caption{Sample Data Table}
\label{tab:sample}
\end{table}

\section{Lists}
\subsection{Bullet Points}
\begin{itemize}
  \item First item
  \item Second item
  \item Third item with sub-items:
    \begin{itemize}
      \item Sub-item 1
      \item Sub-item 2
    \end{itemize}
\end{itemize}

\subsection{Numbered List}
\begin{enumerate}
  \item First step
  \item Second step
  \item Third step
\end{enumerate}

\section{Citations and References}
When you use the citation search feature, references will be automatically managed. For example, you might cite a paper like this: \cite{example2024}. All citations will be automatically added to your \texttt{cite.bib} file.

\section{Conclusion}
This template demonstrates basic \LaTeX{} features including mathematical formulas, tables, and lists. The bibliography system is now set up to work seamlessly with Bibby AI's citation management. Feel free to modify it for your needs!

\bibliographystyle{plain}
\bibliography{cite}

\end{document}\documentclass{article}
\usepackage[utf8]{inputenc}
\usepackage{booktabs}
\usepackage{amsmath}
\usepackage{natbib}

\title{My Document}
\author{Your Name}
\date{\today}

\begin{document}

\maketitle

\section{Introduction}
This is a sample document demonstrating various \LaTeX{} features including mathematical equations, tables, and formatting. You can cite papers using the citation search feature, and they will be automatically added to your \texttt{cite.bib} file.

\section{Mathematical Expressions}
Here are some example equations:

\subsection{Quadratic Formula}
The quadratic formula is given by:
\[x = \frac{-b \pm \sqrt{b^2 - 4ac}}{2a}\]

\subsection{Matrix Example}
A sample matrix:
\[
\begin{pmatrix}
a_{11} & a_{12} & a_{13} \\
a_{21} & a_{22} & a_{23} \\
a_{31} & a_{32} & a_{33}
\end{pmatrix}
\]

\section{Tables}
Below is a sample table showing data:

\begin{table}[h]
\centering
\begin{tabular}{@{}lrrr@{}}
\toprule
\textbf{Category} & \textbf{Value 1} & \textbf{Value 2} & \textbf{Total} \\
\midrule
Item A & 10.5 & 20.3 & 30.8 \\
Item B & 15.2 & 25.1 & 40.3 \\
Item C & 22.1 & 12.7 & 34.8 \\
\bottomrule
\end{tabular}
\caption{Sample Data Table}
\label{tab:sample}
\end{table}

\section{Lists}
\subsection{Bullet Points}
\begin{itemize}
  \item First item
  \item Second item
  \item Third item with sub-items:
    \begin{itemize}
      \item Sub-item 1
      \item Sub-item 2
    \end{itemize}
\end{itemize}

\subsection{Numbered List}
\begin{enumerate}
  \item First step
  \item Second step
  \item Third step
\end{enumerate}

\section{Citations and References}
When you use the citation search feature, references will be automatically managed. For example, you might cite a paper like this: \cite{example2024}. All citations will be automatically added to your \texttt{cite.bib} file.

\section{Conclusion}
This template demonstrates basic \LaTeX{} features including mathematical formulas, tables, and lists. The bibliography system is now set up to work seamlessly with Bibby AI's citation management. Feel free to modify it for your needs!

\bibliographystyle{plain}
\bibliography{cite}

\end{document}\documentclass{article}
\usepackage[utf8]{inputenc}
\usepackage{booktabs}
\usepackage{amsmath}
\usepackage{natbib}

\title{My Document}
\author{Your Name}
\date{\today}

\begin{document}

\maketitle

\section{Introduction}
This is a sample document demonstrating various \LaTeX{} features including mathematical equations, tables, and formatting. You can cite papers using the citation search feature, and they will be automatically added to your \texttt{cite.bib} file.

\section{Mathematical Expressions}
Here are some example equations:

\subsection{Quadratic Formula}
The quadratic formula is given by:
\[x = \frac{-b \pm \sqrt{b^2 - 4ac}}{2a}\]

\subsection{Matrix Example}
A sample matrix:
\[
\begin{pmatrix}
a_{11} & a_{12} & a_{13} \\
a_{21} & a_{22} & a_{23} \\
a_{31} & a_{32} & a_{33}
\end{pmatrix}
\]

\section{Tables}
Below is a sample table showing data:

\begin{table}[h]
\centering
\begin{tabular}{@{}lrrr@{}}
\toprule
\textbf{Category} & \textbf{Value 1} & \textbf{Value 2} & \textbf{Total} \\
\midrule
Item A & 10.5 & 20.3 & 30.8 \\
Item B & 15.2 & 25.1 & 40.3 \\
Item C & 22.1 & 12.7 & 34.8 \\
\bottomrule
\end{tabular}
\caption{Sample Data Table}
\label{tab:sample}
\end{table}

\section{Lists}
\subsection{Bullet Points}
\begin{itemize}
  \item First item
  \item Second item
  \item Third item with sub-items:
    \begin{itemize}
      \item Sub-item 1
      \item Sub-item 2
    \end{itemize}
\end{itemize}

\subsection{Numbered List}
\begin{enumerate}
  \item First step
  \item Second step
  \item Third step
\end{enumerate}

\section{Citations and References}
When you use the citation search feature, references will be automatically managed. For example, you might cite a paper like this: \cite{example2024}. All citations will be automatically added to your \texttt{cite.bib} file.

\section{Conclusion}
This template demonstrates basic \LaTeX{} features including mathematical formulas, tables, and lists. The bibliography system is now set up to work seamlessly with Bibby AI's citation management. Feel free to modify it for your needs!

\bibliographystyle{plain}
\bibliography{cite}

\end{document}\documentclass{article}
\usepackage[utf8]{inputenc}
\usepackage{booktabs}
\usepackage{amsmath}
\usepackage{natbib}

\title{My Document}
\author{Your Name}
\date{\today}

\begin{document}

\maketitle

\section{Introduction}
This is a sample document demonstrating various \LaTeX{} features including mathematical equations, tables, and formatting. You can cite papers using the citation search feature, and they will be automatically added to your \texttt{cite.bib} file.

\section{Mathematical Expressions}
Here are some example equations:

\subsection{Quadratic Formula}
The quadratic formula is given by:
\[x = \frac{-b \pm \sqrt{b^2 - 4ac}}{2a}\]

\subsection{Matrix Example}
A sample matrix:
\[
\begin{pmatrix}
a_{11} & a_{12} & a_{13} \\
a_{21} & a_{22} & a_{23} \\
a_{31} & a_{32} & a_{33}
\end{pmatrix}
\]

\section{Tables}
Below is a sample table showing data:

\begin{table}[h]
\centering
\begin{tabular}{@{}lrrr@{}}
\toprule
\textbf{Category} & \textbf{Value 1} & \textbf{Value 2} & \textbf{Total} \\
\midrule
Item A & 10.5 & 20.3 & 30.8 \\
Item B & 15.2 & 25.1 & 40.3 \\
Item C & 22.1 & 12.7 & 34.8 \\
\bottomrule
\end{tabular}
\caption{Sample Data Table}
\label{tab:sample}
\end{table}

\section{Lists}
\subsection{Bullet Points}
\begin{itemize}
  \item First item
  \item Second item
  \item Third item with sub-items:
    \begin{itemize}
      \item Sub-item 1
      \item Sub-item 2
    \end{itemize}
\end{itemize}

\subsection{Numbered List}
\begin{enumerate}
  \item First step
  \item Second step
  \item Third step
\end{enumerate}

\section{Citations and References}
When you use the citation search feature, references will be automatically managed. For example, you might cite a paper like this: \cite{example2024}. All citations will be automatically added to your \texttt{cite.bib} file.

\section{Conclusion}
This template demonstrates basic \LaTeX{} features including mathematical formulas, tables, and lists. The bibliography system is now set up to work seamlessly with Bibby AI's citation management. Feel free to modify it for your needs!

\bibliographystyle{plain}
\bibliography{cite}

\end{document}\documentclass{article}
\usepackage[utf8]{inputenc}
\usepackage{booktabs}
\usepackage{amsmath}
\usepackage{natbib}

\title{My Document}
\author{Your Name}
\date{\today}

\begin{document}

\maketitle

\section{Introduction}
This is a sample document demonstrating various \LaTeX{} features including mathematical equations, tables, and formatting. You can cite papers using the citation search feature, and they will be automatically added to your \texttt{cite.bib} file.

\section{Mathematical Expressions}
Here are some example equations:

\subsection{Quadratic Formula}
The quadratic formula is given by:
\[x = \frac{-b \pm \sqrt{b^2 - 4ac}}{2a}\]

\subsection{Matrix Example}
A sample matrix:
\[
\begin{pmatrix}
a_{11} & a_{12} & a_{13} \\
a_{21} & a_{22} & a_{23} \\
a_{31} & a_{32} & a_{33}
\end{pmatrix}
\]

\section{Tables}
Below is a sample table showing data:

\begin{table}[h]
\centering
\begin{tabular}{@{}lrrr@{}}
\toprule
\textbf{Category} & \textbf{Value 1} & \textbf{Value 2} & \textbf{Total} \\
\midrule
Item A & 10.5 & 20.3 & 30.8 \\
Item B & 15.2 & 25.1 & 40.3 \\
Item C & 22.1 & 12.7 & 34.8 \\
\bottomrule
\end{tabular}
\caption{Sample Data Table}
\label{tab:sample}
\end{table}

\section{Lists}
\subsection{Bullet Points}
\begin{itemize}
  \item First item
  \item Second item
  \item Third item with sub-items:
    \begin{itemize}
      \item Sub-item 1
      \item Sub-item 2
    \end{itemize}
\end{itemize}

\subsection{Numbered List}
\begin{enumerate}
  \item First step
  \item Second step
  \item Third step
\end{enumerate}

\section{Citations and References}
When you use the citation search feature, references will be automatically managed. For example, you might cite a paper like this: \cite{example2024}. All citations will be automatically added to your \texttt{cite.bib} file.

\section{Conclusion}
This template demonstrates basic \LaTeX{} features including mathematical formulas, tables, and lists. The bibliography system is now set up to work seamlessly with Bibby AI's citation management. Feel free to modify it for your needs!

\bibliographystyle{plain}
\bibliography{cite}

\end{document}\documentclass{article}
\usepackage[utf8]{inputenc}
\usepackage{booktabs}
\usepackage{amsmath}
\usepackage{natbib}

\title{My Document}
\author{Your Name}
\date{\today}

\begin{document}

\maketitle

\section{Introduction}
This is a sample document demonstrating various \LaTeX{} features including mathematical equations, tables, and formatting. You can cite papers using the citation search feature, and they will be automatically added to your \texttt{cite.bib} file.

\section{Mathematical Expressions}
Here are some example equations:

\subsection{Quadratic Formula}
The quadratic formula is given by:
\[x = \frac{-b \pm \sqrt{b^2 - 4ac}}{2a}\]

\subsection{Matrix Example}
A sample matrix:
\[
\begin{pmatrix}
a_{11} & a_{12} & a_{13} \\
a_{21} & a_{22} & a_{23} \\
a_{31} & a_{32} & a_{33}
\end{pmatrix}
\]

\section{Tables}
Below is a sample table showing data:

\begin{table}[h]
\centering
\begin{tabular}{@{}lrrr@{}}
\toprule
\textbf{Category} & \textbf{Value 1} & \textbf{Value 2} & \textbf{Total} \\
\midrule
Item A & 10.5 & 20.3 & 30.8 \\
Item B & 15.2 & 25.1 & 40.3 \\
Item C & 22.1 & 12.7 & 34.8 \\
\bottomrule
\end{tabular}
\caption{Sample Data Table}
\label{tab:sample}
\end{table}

\section{Lists}
\subsection{Bullet Points}
\begin{itemize}
  \item First item
  \item Second item
  \item Third item with sub-items:
    \begin{itemize}
      \item Sub-item 1
      \item Sub-item 2
    \end{itemize}
\end{itemize}

\subsection{Numbered List}
\begin{enumerate}
  \item First step
  \item Second step
  \item Third step
\end{enumerate}

\section{Citations and References}
When you use the citation search feature, references will be automatically managed. For example, you might cite a paper like this: \cite{example2024}. All citations will be automatically added to your \texttt{cite.bib} file.

\section{Conclusion}
This template demonstrates basic \LaTeX{} features including mathematical formulas, tables, and lists. The bibliography system is now set up to work seamlessly with Bibby AI's citation management. Feel free to modify it for your needs!

\bibliographystyle{plain}
\bibliography{cite}

\end{document}\documentclass{article}
\usepackage[utf8]{inputenc}
\usepackage{booktabs}
\usepackage{amsmath}
\usepackage{natbib}

\title{My Document}
\author{Your Name}
\date{\today}

\begin{document}

\maketitle

\section{Introduction}
This is a sample document demonstrating various \LaTeX{} features including mathematical equations, tables, and formatting. You can cite papers using the citation search feature, and they will be automatically added to your \texttt{cite.bib} file.

\section{Mathematical Expressions}
Here are some example equations:

\subsection{Quadratic Formula}
The quadratic formula is given by:
\[x = \frac{-b \pm \sqrt{b^2 - 4ac}}{2a}\]

\subsection{Matrix Example}
A sample matrix:
\[
\begin{pmatrix}
a_{11} & a_{12} & a_{13} \\
a_{21} & a_{22} & a_{23} \\
a_{31} & a_{32} & a_{33}
\end{pmatrix}
\]

\section{Tables}
Below is a sample table showing data:

\begin{table}[h]
\centering
\begin{tabular}{@{}lrrr@{}}
\toprule
\textbf{Category} & \textbf{Value 1} & \textbf{Value 2} & \textbf{Total} \\
\midrule
Item A & 10.5 & 20.3 & 30.8 \\
Item B & 15.2 & 25.1 & 40.3 \\
Item C & 22.1 & 12.7 & 34.8 \\
\bottomrule
\end{tabular}
\caption{Sample Data Table}
\label{tab:sample}
\end{table}

\section{Lists}
\subsection{Bullet Points}
\begin{itemize}
  \item First item
  \item Second item
  \item Third item with sub-items:
    \begin{itemize}
      \item Sub-item 1
      \item Sub-item 2
    \end{itemize}
\end{itemize}

\subsection{Numbered List}
\begin{enumerate}
  \item First step
  \item Second step
  \item Third step
\end{enumerate}

\section{Citations and References}
When you use the citation search feature, references will be automatically managed. For example, you might cite a paper like this: \cite{example2024}. All citations will be automatically added to your \texttt{cite.bib} file.

\section{Conclusion}
This template demonstrates basic \LaTeX{} features including mathematical formulas, tables, and lists. The bibliography system is now set up to work seamlessly with Bibby AI's citation management. Feel free to modify it for your needs!

\bibliographystyle{plain}
\bibliography{cite}

\end{document}\documentclass{article}
\usepackage[utf8]{inputenc}
\usepackage{booktabs}
\usepackage{amsmath}
\usepackage{natbib}

\title{My Document}
\author{Your Name}
\date{\today}

\begin{document}

\maketitle

\section{Introduction}
This is a sample document demonstrating various \LaTeX{} features including mathematical equations, tables, and formatting. You can cite papers using the citation search feature, and they will be automatically added to your \texttt{cite.bib} file.

\section{Mathematical Expressions}
Here are some example equations:

\subsection{Quadratic Formula}
The quadratic formula is given by:
\[x = \frac{-b \pm \sqrt{b^2 - 4ac}}{2a}\]

\subsection{Matrix Example}
A sample matrix:
\[
\begin{pmatrix}
a_{11} & a_{12} & a_{13} \\
a_{21} & a_{22} & a_{23} \\
a_{31} & a_{32} & a_{33}
\end{pmatrix}
\]

\section{Tables}
Below is a sample table showing data:

\begin{table}[h]
\centering
\begin{tabular}{@{}lrrr@{}}
\toprule
\textbf{Category} & \textbf{Value 1} & \textbf{Value 2} & \textbf{Total} \\
\midrule
Item A & 10.5 & 20.3 & 30.8 \\
Item B & 15.2 & 25.1 & 40.3 \\
Item C & 22.1 & 12.7 & 34.8 \\
\bottomrule
\end{tabular}
\caption{Sample Data Table}
\label{tab:sample}
\end{table}

\section{Lists}
\subsection{Bullet Points}
\begin{itemize}
  \item First item
  \item Second item
  \item Third item with sub-items:
    \begin{itemize}
      \item Sub-item 1
      \item Sub-item 2
    \end{itemize}
\end{itemize}

\subsection{Numbered List}
\begin{enumerate}
  \item First step
  \item Second step
  \item Third step
\end{enumerate}

\section{Citations and References}
When you use the citation search feature, references will be automatically managed. For example, you might cite a paper like this: \cite{example2024}. All citations will be automatically added to your \texttt{cite.bib} file.

\section{Conclusion}
This template demonstrates basic \LaTeX{} features including mathematical formulas, tables, and lists. The bibliography system is now set up to work seamlessly with Bibby AI's citation management. Feel free to modify it for your needs!

\bibliographystyle{plain}
\bibliography{cite}

\end{document}\documentclass{article}
\usepackage[utf8]{inputenc}
\usepackage{booktabs}
\usepackage{amsmath}
\usepackage{natbib}

\title{My Document}
\author{Your Name}
\date{\today}

\begin{document}

\maketitle

\section{Introduction}
This is a sample document demonstrating various \LaTeX{} features including mathematical equations, tables, and formatting. You can cite papers using the citation search feature, and they will be automatically added to your \texttt{cite.bib} file.

\section{Mathematical Expressions}
Here are some example equations:

\subsection{Quadratic Formula}
The quadratic formula is given by:
\[x = \frac{-b \pm \sqrt{b^2 - 4ac}}{2a}\]

\subsection{Matrix Example}
A sample matrix:
\[
\begin{pmatrix}
a_{11} & a_{12} & a_{13} \\
a_{21} & a_{22} & a_{23} \\
a_{31} & a_{32} & a_{33}
\end{pmatrix}
\]

\section{Tables}
Below is a sample table showing data:

\begin{table}[h]
\centering
\begin{tabular}{@{}lrrr@{}}
\toprule
\textbf{Category} & \textbf{Value 1} & \textbf{Value 2} & \textbf{Total} \\
\midrule
Item A & 10.5 & 20.3 & 30.8 \\
Item B & 15.2 & 25.1 & 40.3 \\
Item C & 22.1 & 12.7 & 34.8 \\
\bottomrule
\end{tabular}
\caption{Sample Data Table}
\label{tab:sample}
\end{table}

\section{Lists}
\subsection{Bullet Points}
\begin{itemize}
  \item First item
  \item Second item
  \item Third item with sub-items:
    \begin{itemize}
      \item Sub-item 1
      \item Sub-item 2
    \end{itemize}
\end{itemize}

\subsection{Numbered List}
\begin{enumerate}
  \item First step
  \item Second step
  \item Third step
\end{enumerate}

\section{Citations and References}
When you use the citation search feature, references will be automatically managed. For example, you might cite a paper like this: \cite{example2024}. All citations will be automatically added to your \texttt{cite.bib} file.

\section{Conclusion}
This template demonstrates basic \LaTeX{} features including mathematical formulas, tables, and lists. The bibliography system is now set up to work seamlessly with Bibby AI's citation management. Feel free to modify it for your needs!

\bibliographystyle{plain}
\bibliography{cite}

\end{document}\documentclass{article}
\usepackage[utf8]{inputenc}
\usepackage{booktabs}
\usepackage{amsmath}
\usepackage{natbib}

\title{My Document}
\author{Your Name}
\date{\today}

\begin{document}

\maketitle

\section{Introduction}
This is a sample document demonstrating various \LaTeX{} features including mathematical equations, tables, and formatting. You can cite papers using the citation search feature, and they will be automatically added to your \texttt{cite.bib} file.

\section{Mathematical Expressions}
Here are some example equations:

\subsection{Quadratic Formula}
The quadratic formula is given by:
\[x = \frac{-b \pm \sqrt{b^2 - 4ac}}{2a}\]

\subsection{Matrix Example}
A sample matrix:
\[
\begin{pmatrix}
a_{11} & a_{12} & a_{13} \\
a_{21} & a_{22} & a_{23} \\
a_{31} & a_{32} & a_{33}
\end{pmatrix}
\]

\section{Tables}
Below is a sample table showing data:

\begin{table}[h]
\centering
\begin{tabular}{@{}lrrr@{}}
\toprule
\textbf{Category} & \textbf{Value 1} & \textbf{Value 2} & \textbf{Total} \\
\midrule
Item A & 10.5 & 20.3 & 30.8 \\
Item B & 15.2 & 25.1 & 40.3 \\
Item C & 22.1 & 12.7 & 34.8 \\
\bottomrule
\end{tabular}
\caption{Sample Data Table}
\label{tab:sample}
\end{table}

\section{Lists}
\subsection{Bullet Points}
\begin{itemize}
  \item First item
  \item Second item
  \item Third item with sub-items:
    \begin{itemize}
      \item Sub-item 1
      \item Sub-item 2
    \end{itemize}
\end{itemize}

\subsection{Numbered List}
\begin{enumerate}
  \item First step
  \item Second step
  \item Third step
\end{enumerate}

\section{Citations and References}
When you use the citation search feature, references will be automatically managed. For example, you might cite a paper like this: \cite{example2024}. All citations will be automatically added to your \texttt{cite.bib} file.

\section{Conclusion}
This template demonstrates basic \LaTeX{} features including mathematical formulas, tables, and lists. The bibliography system is now set up to work seamlessly with Bibby AI's citation management. Feel free to modify it for your needs!

\bibliographystyle{plain}
\bibliography{cite}

\end{document}\documentclass{article}
\usepackage[utf8]{inputenc}
\usepackage{booktabs}
\usepackage{amsmath}
\usepackage{natbib}

\title{My Document}
\author{Your Name}
\date{\today}

\begin{document}

\maketitle

\section{Introduction}
This is a sample document demonstrating various \LaTeX{} features including mathematical equations, tables, and formatting. You can cite papers using the citation search feature, and they will be automatically added to your \texttt{cite.bib} file.

\section{Mathematical Expressions}
Here are some example equations:

\subsection{Quadratic Formula}
The quadratic formula is given by:
\[x = \frac{-b \pm \sqrt{b^2 - 4ac}}{2a}\]

\subsection{Matrix Example}
A sample matrix:
\[
\begin{pmatrix}
a_{11} & a_{12} & a_{13} \\
a_{21} & a_{22} & a_{23} \\
a_{31} & a_{32} & a_{33}
\end{pmatrix}
\]

\section{Tables}
Below is a sample table showing data:

\begin{table}[h]
\centering
\begin{tabular}{@{}lrrr@{}}
\toprule
\textbf{Category} & \textbf{Value 1} & \textbf{Value 2} & \textbf{Total} \\
\midrule
Item A & 10.5 & 20.3 & 30.8 \\
Item B & 15.2 & 25.1 & 40.3 \\
Item C & 22.1 & 12.7 & 34.8 \\
\bottomrule
\end{tabular}
\caption{Sample Data Table}
\label{tab:sample}
\end{table}

\section{Lists}
\subsection{Bullet Points}
\begin{itemize}
  \item First item
  \item Second item
  \item Third item with sub-items:
    \begin{itemize}
      \item Sub-item 1
      \item Sub-item 2
    \end{itemize}
\end{itemize}

\subsection{Numbered List}
\begin{enumerate}
  \item First step
  \item Second step
  \item Third step
\end{enumerate}

\section{Citations and References}
When you use the citation search feature, references will be automatically managed. For example, you might cite a paper like this: \cite{example2024}. All citations will be automatically added to your \texttt{cite.bib} file.

\section{Conclusion}
This template demonstrates basic \LaTeX{} features including mathematical formulas, tables, and lists. The bibliography system is now set up to work seamlessly with Bibby AI's citation management. Feel free to modify it for your needs!

\bibliographystyle{plain}
\bibliography{cite}

\end{document}\documentclass{article}
\usepackage[utf8]{inputenc}
\usepackage{booktabs}
\usepackage{amsmath}
\usepackage{natbib}

\title{My Document}
\author{Your Name}
\date{\today}

\begin{document}

\maketitle

\section{Introduction}
This is a sample document demonstrating various \LaTeX{} features including mathematical equations, tables, and formatting. You can cite papers using the citation search feature, and they will be automatically added to your \texttt{cite.bib} file.

\section{Mathematical Expressions}
Here are some example equations:

\subsection{Quadratic Formula}
The quadratic formula is given by:
\[x = \frac{-b \pm \sqrt{b^2 - 4ac}}{2a}\]

\subsection{Matrix Example}
A sample matrix:
\[
\begin{pmatrix}
a_{11} & a_{12} & a_{13} \\
a_{21} & a_{22} & a_{23} \\
a_{31} & a_{32} & a_{33}
\end{pmatrix}
\]

\section{Tables}
Below is a sample table showing data:

\begin{table}[h]
\centering
\begin{tabular}{@{}lrrr@{}}
\toprule
\textbf{Category} & \textbf{Value 1} & \textbf{Value 2} & \textbf{Total} \\
\midrule
Item A & 10.5 & 20.3 & 30.8 \\
Item B & 15.2 & 25.1 & 40.3 \\
Item C & 22.1 & 12.7 & 34.8 \\
\bottomrule
\end{tabular}
\caption{Sample Data Table}
\label{tab:sample}
\end{table}

\section{Lists}
\subsection{Bullet Points}
\begin{itemize}
  \item First item
  \item Second item
  \item Third item with sub-items:
    \begin{itemize}
      \item Sub-item 1
      \item Sub-item 2
    \end{itemize}
\end{itemize}

\subsection{Numbered List}
\begin{enumerate}
  \item First step
  \item Second step
  \item Third step
\end{enumerate}

\section{Citations and References}
When you use the citation search feature, references will be automatically managed. For example, you might cite a paper like this: \cite{example2024}. All citations will be automatically added to your \texttt{cite.bib} file.

\section{Conclusion}
This template demonstrates basic \LaTeX{} features including mathematical formulas, tables, and lists. The bibliography system is now set up to work seamlessly with Bibby AI's citation management. Feel free to modify it for your needs!

\bibliographystyle{plain}
\bibliography{cite}

\end{document}\documentclass{article}
\usepackage[utf8]{inputenc}
\usepackage{booktabs}
\usepackage{amsmath}
\usepackage{natbib}

\title{My Document}
\author{Your Name}
\date{\today}

\begin{document}

\maketitle

\section{Introduction}
This is a sample document demonstrating various \LaTeX{} features including mathematical equations, tables, and formatting. You can cite papers using the citation search feature, and they will be automatically added to your \texttt{cite.bib} file.

\section{Mathematical Expressions}
Here are some example equations:

\subsection{Quadratic Formula}
The quadratic formula is given by:
\[x = \frac{-b \pm \sqrt{b^2 - 4ac}}{2a}\]

\subsection{Matrix Example}
A sample matrix:
\[
\begin{pmatrix}
a_{11} & a_{12} & a_{13} \\
a_{21} & a_{22} & a_{23} \\
a_{31} & a_{32} & a_{33}
\end{pmatrix}
\]

\section{Tables}
Below is a sample table showing data:

\begin{table}[h]
\centering
\begin{tabular}{@{}lrrr@{}}
\toprule
\textbf{Category} & \textbf{Value 1} & \textbf{Value 2} & \textbf{Total} \\
\midrule
Item A & 10.5 & 20.3 & 30.8 \\
Item B & 15.2 & 25.1 & 40.3 \\
Item C & 22.1 & 12.7 & 34.8 \\
\bottomrule
\end{tabular}
\caption{Sample Data Table}
\label{tab:sample}
\end{table}

\section{Lists}
\subsection{Bullet Points}
\begin{itemize}
  \item First item
  \item Second item
  \item Third item with sub-items:
    \begin{itemize}
      \item Sub-item 1
      \item Sub-item 2
    \end{itemize}
\end{itemize}

\subsection{Numbered List}
\begin{enumerate}
  \item First step
  \item Second step
  \item Third step
\end{enumerate}

\section{Citations and References}
When you use the citation search feature, references will be automatically managed. For example, you might cite a paper like this: \cite{example2024}. All citations will be automatically added to your \texttt{cite.bib} file.

\section{Conclusion}
This template demonstrates basic \LaTeX{} features including mathematical formulas, tables, and lists. The bibliography system is now set up to work seamlessly with Bibby AI's citation management. Feel free to modify it for your needs!

\bibliographystyle{plain}
\bibliography{cite}

\end{document}\documentclass{article}
\usepackage[utf8]{inputenc}
\usepackage{booktabs}
\usepackage{amsmath}
\usepackage{natbib}

\title{My Document}
\author{Your Name}
\date{\today}

\begin{document}

\maketitle

\section{Introduction}
This is a sample document demonstrating various \LaTeX{} features including mathematical equations, tables, and formatting. You can cite papers using the citation search feature, and they will be automatically added to your \texttt{cite.bib} file.

\section{Mathematical Expressions}
Here are some example equations:

\subsection{Quadratic Formula}
The quadratic formula is given by:
\[x = \frac{-b \pm \sqrt{b^2 - 4ac}}{2a}\]

\subsection{Matrix Example}
A sample matrix:
\[
\begin{pmatrix}
a_{11} & a_{12} & a_{13} \\
a_{21} & a_{22} & a_{23} \\
a_{31} & a_{32} & a_{33}
\end{pmatrix}
\]

\section{Tables}
Below is a sample table showing data:

\begin{table}[h]
\centering
\begin{tabular}{@{}lrrr@{}}
\toprule
\textbf{Category} & \textbf{Value 1} & \textbf{Value 2} & \textbf{Total} \\
\midrule
Item A & 10.5 & 20.3 & 30.8 \\
Item B & 15.2 & 25.1 & 40.3 \\
Item C & 22.1 & 12.7 & 34.8 \\
\bottomrule
\end{tabular}
\caption{Sample Data Table}
\label{tab:sample}
\end{table}

\section{Lists}
\subsection{Bullet Points}
\begin{itemize}
  \item First item
  \item Second item
  \item Third item with sub-items:
    \begin{itemize}
      \item Sub-item 1
      \item Sub-item 2
    \end{itemize}
\end{itemize}

\subsection{Numbered List}
\begin{enumerate}
  \item First step
  \item Second step
  \item Third step
\end{enumerate}

\section{Citations and References}
When you use the citation search feature, references will be automatically managed. For example, you might cite a paper like this: \cite{example2024}. All citations will be automatically added to your \texttt{cite.bib} file.

\section{Conclusion}
This template demonstrates basic \LaTeX{} features including mathematical formulas, tables, and lists. The bibliography system is now set up to work seamlessly with Bibby AI's citation management. Feel free to modify it for your needs!

\bibliographystyle{plain}
\bibliography{cite}

\end{document}\documentclass{article}
\usepackage[utf8]{inputenc}
\usepackage{booktabs}
\usepackage{amsmath}
\usepackage{natbib}

\title{My Document}
\author{Your Name}
\date{\today}

\begin{document}

\maketitle

\section{Introduction}
This is a sample document demonstrating various \LaTeX{} features including mathematical equations, tables, and formatting. You can cite papers using the citation search feature, and they will be automatically added to your \texttt{cite.bib} file.

\section{Mathematical Expressions}
Here are some example equations:

\subsection{Quadratic Formula}
The quadratic formula is given by:
\[x = \frac{-b \pm \sqrt{b^2 - 4ac}}{2a}\]

\subsection{Matrix Example}
A sample matrix:
\[
\begin{pmatrix}
a_{11} & a_{12} & a_{13} \\
a_{21} & a_{22} & a_{23} \\
a_{31} & a_{32} & a_{33}
\end{pmatrix}
\]

\section{Tables}
Below is a sample table showing data:

\begin{table}[h]
\centering
\begin{tabular}{@{}lrrr@{}}
\toprule
\textbf{Category} & \textbf{Value 1} & \textbf{Value 2} & \textbf{Total} \\
\midrule
Item A & 10.5 & 20.3 & 30.8 \\
Item B & 15.2 & 25.1 & 40.3 \\
Item C & 22.1 & 12.7 & 34.8 \\
\bottomrule
\end{tabular}
\caption{Sample Data Table}
\label{tab:sample}
\end{table}

\section{Lists}
\subsection{Bullet Points}
\begin{itemize}
  \item First item
  \item Second item
  \item Third item with sub-items:
    \begin{itemize}
      \item Sub-item 1
      \item Sub-item 2
    \end{itemize}
\end{itemize}

\subsection{Numbered List}
\begin{enumerate}
  \item First step
  \item Second step
  \item Third step
\end{enumerate}

\section{Citations and References}
When you use the citation search feature, references will be automatically managed. For example, you might cite a paper like this: \cite{example2024}. All citations will be automatically added to your \texttt{cite.bib} file.

\section{Conclusion}
This template demonstrates basic \LaTeX{} features including mathematical formulas, tables, and lists. The bibliography system is now set up to work seamlessly with Bibby AI's citation management. Feel free to modify it for your needs!

\bibliographystyle{plain}
\bibliography{cite}

\end{document}\documentclass{article}
\usepackage[utf8]{inputenc}
\usepackage{booktabs}
\usepackage{amsmath}
\usepackage{natbib}

\title{My Document}
\author{Your Name}
\date{\today}

\begin{document}

\maketitle

\section{Introduction}
This is a sample document demonstrating various \LaTeX{} features including mathematical equations, tables, and formatting. You can cite papers using the citation search feature, and they will be automatically added to your \texttt{cite.bib} file.

\section{Mathematical Expressions}
Here are some example equations:

\subsection{Quadratic Formula}
The quadratic formula is given by:
\[x = \frac{-b \pm \sqrt{b^2 - 4ac}}{2a}\]

\subsection{Matrix Example}
A sample matrix:
\[
\begin{pmatrix}
a_{11} & a_{12} & a_{13} \\
a_{21} & a_{22} & a_{23} \\
a_{31} & a_{32} & a_{33}
\end{pmatrix}
\]

\section{Tables}
Below is a sample table showing data:

\begin{table}[h]
\centering
\begin{tabular}{@{}lrrr@{}}
\toprule
\textbf{Category} & \textbf{Value 1} & \textbf{Value 2} & \textbf{Total} \\
\midrule
Item A & 10.5 & 20.3 & 30.8 \\
Item B & 15.2 & 25.1 & 40.3 \\
Item C & 22.1 & 12.7 & 34.8 \\
\bottomrule
\end{tabular}
\caption{Sample Data Table}
\label{tab:sample}
\end{table}

\section{Lists}
\subsection{Bullet Points}
\begin{itemize}
  \item First item
  \item Second item
  \item Third item with sub-items:
    \begin{itemize}
      \item Sub-item 1
      \item Sub-item 2
    \end{itemize}
\end{itemize}

\subsection{Numbered List}
\begin{enumerate}
  \item First step
  \item Second step
  \item Third step
\end{enumerate}

\section{Citations and References}
When you use the citation search feature, references will be automatically managed. For example, you might cite a paper like this: \cite{example2024}. All citations will be automatically added to your \texttt{cite.bib} file.

\section{Conclusion}
This template demonstrates basic \LaTeX{} features including mathematical formulas, tables, and lists. The bibliography system is now set up to work seamlessly with Bibby AI's citation management. Feel free to modify it for your needs!

\bibliographystyle{plain}
\bibliography{cite}

\end{document}\documentclass{article}
\usepackage[utf8]{inputenc}
\usepackage{booktabs}
\usepackage{amsmath}
\usepackage{natbib}

\title{My Document}
\author{Your Name}
\date{\today}

\begin{document}

\maketitle

\section{Introduction}
This is a sample document demonstrating various \LaTeX{} features including mathematical equations, tables, and formatting. You can cite papers using the citation search feature, and they will be automatically added to your \texttt{cite.bib} file.

\section{Mathematical Expressions}
Here are some example equations:

\subsection{Quadratic Formula}
The quadratic formula is given by:
\[x = \frac{-b \pm \sqrt{b^2 - 4ac}}{2a}\]

\subsection{Matrix Example}
A sample matrix:
\[
\begin{pmatrix}
a_{11} & a_{12} & a_{13} \\
a_{21} & a_{22} & a_{23} \\
a_{31} & a_{32} & a_{33}
\end{pmatrix}
\]

\section{Tables}
Below is a sample table showing data:

\begin{table}[h]
\centering
\begin{tabular}{@{}lrrr@{}}
\toprule
\textbf{Category} & \textbf{Value 1} & \textbf{Value 2} & \textbf{Total} \\
\midrule
Item A & 10.5 & 20.3 & 30.8 \\
Item B & 15.2 & 25.1 & 40.3 \\
Item C & 22.1 & 12.7 & 34.8 \\
\bottomrule
\end{tabular}
\caption{Sample Data Table}
\label{tab:sample}
\end{table}

\section{Lists}
\subsection{Bullet Points}
\begin{itemize}
  \item First item
  \item Second item
  \item Third item with sub-items:
    \begin{itemize}
      \item Sub-item 1
      \item Sub-item 2
    \end{itemize}
\end{itemize}

\subsection{Numbered List}
\begin{enumerate}
  \item First step
  \item Second step
  \item Third step
\end{enumerate}

\section{Citations and References}
When you use the citation search feature, references will be automatically managed. For example, you might cite a paper like this: \cite{example2024}. All citations will be automatically added to your \texttt{cite.bib} file.

\section{Conclusion}
This template demonstrates basic \LaTeX{} features including mathematical formulas, tables, and lists. The bibliography system is now set up to work seamlessly with Bibby AI's citation management. Feel free to modify it for your needs!

\bibliographystyle{plain}
\bibliography{cite}

\end{document}\documentclass{article}
\usepackage[utf8]{inputenc}
\usepackage{booktabs}
\usepackage{amsmath}
\usepackage{natbib}

\title{My Document}
\author{Your Name}
\date{\today}

\begin{document}

\maketitle

\section{Introduction}
This is a sample document demonstrating various \LaTeX{} features including mathematical equations, tables, and formatting. You can cite papers using the citation search feature, and they will be automatically added to your \texttt{cite.bib} file.

\section{Mathematical Expressions}
Here are some example equations:

\subsection{Quadratic Formula}
The quadratic formula is given by:
\[x = \frac{-b \pm \sqrt{b^2 - 4ac}}{2a}\]

\subsection{Matrix Example}
A sample matrix:
\[
\begin{pmatrix}
a_{11} & a_{12} & a_{13} \\
a_{21} & a_{22} & a_{23} \\
a_{31} & a_{32} & a_{33}
\end{pmatrix}
\]

\section{Tables}
Below is a sample table showing data:

\begin{table}[h]
\centering
\begin{tabular}{@{}lrrr@{}}
\toprule
\textbf{Category} & \textbf{Value 1} & \textbf{Value 2} & \textbf{Total} \\
\midrule
Item A & 10.5 & 20.3 & 30.8 \\
Item B & 15.2 & 25.1 & 40.3 \\
Item C & 22.1 & 12.7 & 34.8 \\
\bottomrule
\end{tabular}
\caption{Sample Data Table}
\label{tab:sample}
\end{table}

\section{Lists}
\subsection{Bullet Points}
\begin{itemize}
  \item First item
  \item Second item
  \item Third item with sub-items:
    \begin{itemize}
      \item Sub-item 1
      \item Sub-item 2
    \end{itemize}
\end{itemize}

\subsection{Numbered List}
\begin{enumerate}
  \item First step
  \item Second step
  \item Third step
\end{enumerate}

\section{Citations and References}
When you use the citation search feature, references will be automatically managed. For example, you might cite a paper like this: \cite{example2024}. All citations will be automatically added to your \texttt{cite.bib} file.

\section{Conclusion}
This template demonstrates basic \LaTeX{} features including mathematical formulas, tables, and lists. The bibliography system is now set up to work seamlessly with Bibby AI's citation management. Feel free to modify it for your needs!

\bibliographystyle{plain}
\bibliography{cite}

\end{document}\documentclass{article}
\usepackage[utf8]{inputenc}
\usepackage{booktabs}
\usepackage{amsmath}
\usepackage{natbib}

\title{My Document}
\author{Your Name}
\date{\today}

\begin{document}

\maketitle

\section{Introduction}
This is a sample document demonstrating various \LaTeX{} features including mathematical equations, tables, and formatting. You can cite papers using the citation search feature, and they will be automatically added to your \texttt{cite.bib} file.

\section{Mathematical Expressions}
Here are some example equations:

\subsection{Quadratic Formula}
The quadratic formula is given by:
\[x = \frac{-b \pm \sqrt{b^2 - 4ac}}{2a}\]

\subsection{Matrix Example}
A sample matrix:
\[
\begin{pmatrix}
a_{11} & a_{12} & a_{13} \\
a_{21} & a_{22} & a_{23} \\
a_{31} & a_{32} & a_{33}
\end{pmatrix}
\]

\section{Tables}
Below is a sample table showing data:

\begin{table}[h]
\centering
\begin{tabular}{@{}lrrr@{}}
\toprule
\textbf{Category} & \textbf{Value 1} & \textbf{Value 2} & \textbf{Total} \\
\midrule
Item A & 10.5 & 20.3 & 30.8 \\
Item B & 15.2 & 25.1 & 40.3 \\
Item C & 22.1 & 12.7 & 34.8 \\
\bottomrule
\end{tabular}
\caption{Sample Data Table}
\label{tab:sample}
\end{table}

\section{Lists}
\subsection{Bullet Points}
\begin{itemize}
  \item First item
  \item Second item
  \item Third item with sub-items:
    \begin{itemize}
      \item Sub-item 1
      \item Sub-item 2
    \end{itemize}
\end{itemize}

\subsection{Numbered List}
\begin{enumerate}
  \item First step
  \item Second step
  \item Third step
\end{enumerate}

\section{Citations and References}
When you use the citation search feature, references will be automatically managed. For example, you might cite a paper like this: \cite{example2024}. All citations will be automatically added to your \texttt{cite.bib} file.

\section{Conclusion}
This template demonstrates basic \LaTeX{} features including mathematical formulas, tables, and lists. The bibliography system is now set up to work seamlessly with Bibby AI's citation management. Feel free to modify it for your needs!

\bibliographystyle{plain}
\bibliography{cite}

\end{document}\documentclass{article}
\usepackage[utf8]{inputenc}
\usepackage{booktabs}
\usepackage{amsmath}
\usepackage{natbib}

\title{My Document}
\author{Your Name}
\date{\today}

\begin{document}

\maketitle

\section{Introduction}
This is a sample document demonstrating various \LaTeX{} features including mathematical equations, tables, and formatting. You can cite papers using the citation search feature, and they will be automatically added to your \texttt{cite.bib} file.

\section{Mathematical Expressions}
Here are some example equations:

\subsection{Quadratic Formula}
The quadratic formula is given by:
\[x = \frac{-b \pm \sqrt{b^2 - 4ac}}{2a}\]

\subsection{Matrix Example}
A sample matrix:
\[
\begin{pmatrix}
a_{11} & a_{12} & a_{13} \\
a_{21} & a_{22} & a_{23} \\
a_{31} & a_{32} & a_{33}
\end{pmatrix}
\]

\section{Tables}
Below is a sample table showing data:

\begin{table}[h]
\centering
\begin{tabular}{@{}lrrr@{}}
\toprule
\textbf{Category} & \textbf{Value 1} & \textbf{Value 2} & \textbf{Total} \\
\midrule
Item A & 10.5 & 20.3 & 30.8 \\
Item B & 15.2 & 25.1 & 40.3 \\
Item C & 22.1 & 12.7 & 34.8 \\
\bottomrule
\end{tabular}
\caption{Sample Data Table}
\label{tab:sample}
\end{table}

\section{Lists}
\subsection{Bullet Points}
\begin{itemize}
  \item First item
  \item Second item
  \item Third item with sub-items:
    \begin{itemize}
      \item Sub-item 1
      \item Sub-item 2
    \end{itemize}
\end{itemize}

\subsection{Numbered List}
\begin{enumerate}
  \item First step
  \item Second step
  \item Third step
\end{enumerate}

\section{Citations and References}
When you use the citation search feature, references will be automatically managed. For example, you might cite a paper like this: \cite{example2024}. All citations will be automatically added to your \texttt{cite.bib} file.

\section{Conclusion}
This template demonstrates basic \LaTeX{} features including mathematical formulas, tables, and lists. The bibliography system is now set up to work seamlessly with Bibby AI's citation management. Feel free to modify it for your needs!

\bibliographystyle{plain}
\bibliography{cite}

\end{document}\documentclass{article}
\usepackage[utf8]{inputenc}
\usepackage{booktabs}
\usepackage{amsmath}
\usepackage{natbib}

\title{My Document}
\author{Your Name}
\date{\today}

\begin{document}

\maketitle

\section{Introduction}
This is a sample document demonstrating various \LaTeX{} features including mathematical equations, tables, and formatting. You can cite papers using the citation search feature, and they will be automatically added to your \texttt{cite.bib} file.

\section{Mathematical Expressions}
Here are some example equations:

\subsection{Quadratic Formula}
The quadratic formula is given by:
\[x = \frac{-b \pm \sqrt{b^2 - 4ac}}{2a}\]

\subsection{Matrix Example}
A sample matrix:
\[
\begin{pmatrix}
a_{11} & a_{12} & a_{13} \\
a_{21} & a_{22} & a_{23} \\
a_{31} & a_{32} & a_{33}
\end{pmatrix}
\]

\section{Tables}
Below is a sample table showing data:

\begin{table}[h]
\centering
\begin{tabular}{@{}lrrr@{}}
\toprule
\textbf{Category} & \textbf{Value 1} & \textbf{Value 2} & \textbf{Total} \\
\midrule
Item A & 10.5 & 20.3 & 30.8 \\
Item B & 15.2 & 25.1 & 40.3 \\
Item C & 22.1 & 12.7 & 34.8 \\
\bottomrule
\end{tabular}
\caption{Sample Data Table}
\label{tab:sample}
\end{table}

\section{Lists}
\subsection{Bullet Points}
\begin{itemize}
  \item First item
  \item Second item
  \item Third item with sub-items:
    \begin{itemize}
      \item Sub-item 1
      \item Sub-item 2
    \end{itemize}
\end{itemize}

\subsection{Numbered List}
\begin{enumerate}
  \item First step
  \item Second step
  \item Third step
\end{enumerate}

\section{Citations and References}
When you use the citation search feature, references will be automatically managed. For example, you might cite a paper like this: \cite{example2024}. All citations will be automatically added to your \texttt{cite.bib} file.

\section{Conclusion}
This template demonstrates basic \LaTeX{} features including mathematical formulas, tables, and lists. The bibliography system is now set up to work seamlessly with Bibby AI's citation management. Feel free to modify it for your needs!

\bibliographystyle{plain}
\bibliography{cite}

\end{document}\documentclass{article}
\usepackage[utf8]{inputenc}
\usepackage{booktabs}
\usepackage{amsmath}
\usepackage{natbib}

\title{My Document}
\author{Your Name}
\date{\today}

\begin{document}

\maketitle

\section{Introduction}
This is a sample document demonstrating various \LaTeX{} features including mathematical equations, tables, and formatting. You can cite papers using the citation search feature, and they will be automatically added to your \texttt{cite.bib} file.

\section{Mathematical Expressions}
Here are some example equations:

\subsection{Quadratic Formula}
The quadratic formula is given by:
\[x = \frac{-b \pm \sqrt{b^2 - 4ac}}{2a}\]

\subsection{Matrix Example}
A sample matrix:
\[
\begin{pmatrix}
a_{11} & a_{12} & a_{13} \\
a_{21} & a_{22} & a_{23} \\
a_{31} & a_{32} & a_{33}
\end{pmatrix}
\]

\section{Tables}
Below is a sample table showing data:

\begin{table}[h]
\centering
\begin{tabular}{@{}lrrr@{}}
\toprule
\textbf{Category} & \textbf{Value 1} & \textbf{Value 2} & \textbf{Total} \\
\midrule
Item A & 10.5 & 20.3 & 30.8 \\
Item B & 15.2 & 25.1 & 40.3 \\
Item C & 22.1 & 12.7 & 34.8 \\
\bottomrule
\end{tabular}
\caption{Sample Data Table}
\label{tab:sample}
\end{table}

\section{Lists}
\subsection{Bullet Points}
\begin{itemize}
  \item First item
  \item Second item
  \item Third item with sub-items:
    \begin{itemize}
      \item Sub-item 1
      \item Sub-item 2
    \end{itemize}
\end{itemize}

\subsection{Numbered List}
\begin{enumerate}
  \item First step
  \item Second step
  \item Third step
\end{enumerate}

\section{Citations and References}
When you use the citation search feature, references will be automatically managed. For example, you might cite a paper like this: \cite{example2024}. All citations will be automatically added to your \texttt{cite.bib} file.

\section{Conclusion}
This template demonstrates basic \LaTeX{} features including mathematical formulas, tables, and lists. The bibliography system is now set up to work seamlessly with Bibby AI's citation management. Feel free to modify it for your needs!

\bibliographystyle{plain}
\bibliography{cite}

\end{document}\documentclass{article}
\usepackage[utf8]{inputenc}
\usepackage{booktabs}
\usepackage{amsmath}
\usepackage{natbib}

\title{My Document}
\author{Your Name}
\date{\today}

\begin{document}

\maketitle

\section{Introduction}
This is a sample document demonstrating various \LaTeX{} features including mathematical equations, tables, and formatting. You can cite papers using the citation search feature, and they will be automatically added to your \texttt{cite.bib} file.

\section{Mathematical Expressions}
Here are some example equations:

\subsection{Quadratic Formula}
The quadratic formula is given by:
\[x = \frac{-b \pm \sqrt{b^2 - 4ac}}{2a}\]

\subsection{Matrix Example}
A sample matrix:
\[
\begin{pmatrix}
a_{11} & a_{12} & a_{13} \\
a_{21} & a_{22} & a_{23} \\
a_{31} & a_{32} & a_{33}
\end{pmatrix}
\]

\section{Tables}
Below is a sample table showing data:

\begin{table}[h]
\centering
\begin{tabular}{@{}lrrr@{}}
\toprule
\textbf{Category} & \textbf{Value 1} & \textbf{Value 2} & \textbf{Total} \\
\midrule
Item A & 10.5 & 20.3 & 30.8 \\
Item B & 15.2 & 25.1 & 40.3 \\
Item C & 22.1 & 12.7 & 34.8 \\
\bottomrule
\end{tabular}
\caption{Sample Data Table}
\label{tab:sample}
\end{table}

\section{Lists}
\subsection{Bullet Points}
\begin{itemize}
  \item First item
  \item Second item
  \item Third item with sub-items:
    \begin{itemize}
      \item Sub-item 1
      \item Sub-item 2
    \end{itemize}
\end{itemize}

\subsection{Numbered List}
\begin{enumerate}
  \item First step
  \item Second step
  \item Third step
\end{enumerate}

\section{Citations and References}
When you use the citation search feature, references will be automatically managed. For example, you might cite a paper like this: \cite{example2024}. All citations will be automatically added to your \texttt{cite.bib} file.

\section{Conclusion}
This template demonstrates basic \LaTeX{} features including mathematical formulas, tables, and lists. The bibliography system is now set up to work seamlessly with Bibby AI's citation management. Feel free to modify it for your needs!

\bibliographystyle{plain}
\bibliography{cite}

\end{document}\documentclass{article}
\usepackage[utf8]{inputenc}
\usepackage{booktabs}
\usepackage{amsmath}
\usepackage{natbib}

\title{My Document}
\author{Your Name}
\date{\today}

\begin{document}

\maketitle

\section{Introduction}
This is a sample document demonstrating various \LaTeX{} features including mathematical equations, tables, and formatting. You can cite papers using the citation search feature, and they will be automatically added to your \texttt{cite.bib} file.

\section{Mathematical Expressions}
Here are some example equations:

\subsection{Quadratic Formula}
The quadratic formula is given by:
\[x = \frac{-b \pm \sqrt{b^2 - 4ac}}{2a}\]

\subsection{Matrix Example}
A sample matrix:
\[
\begin{pmatrix}
a_{11} & a_{12} & a_{13} \\
a_{21} & a_{22} & a_{23} \\
a_{31} & a_{32} & a_{33}
\end{pmatrix}
\]

\section{Tables}
Below is a sample table showing data:

\begin{table}[h]
\centering
\begin{tabular}{@{}lrrr@{}}
\toprule
\textbf{Category} & \textbf{Value 1} & \textbf{Value 2} & \textbf{Total} \\
\midrule
Item A & 10.5 & 20.3 & 30.8 \\
Item B & 15.2 & 25.1 & 40.3 \\
Item C & 22.1 & 12.7 & 34.8 \\
\bottomrule
\end{tabular}
\caption{Sample Data Table}
\label{tab:sample}
\end{table}

\section{Lists}
\subsection{Bullet Points}
\begin{itemize}
  \item First item
  \item Second item
  \item Third item with sub-items:
    \begin{itemize}
      \item Sub-item 1
      \item Sub-item 2
    \end{itemize}
\end{itemize}

\subsection{Numbered List}
\begin{enumerate}
  \item First step
  \item Second step
  \item Third step
\end{enumerate}

\section{Citations and References}
When you use the citation search feature, references will be automatically managed. For example, you might cite a paper like this: \cite{example2024}. All citations will be automatically added to your \texttt{cite.bib} file.

\section{Conclusion}
This template demonstrates basic \LaTeX{} features including mathematical formulas, tables, and lists. The bibliography system is now set up to work seamlessly with Bibby AI's citation management. Feel free to modify it for your needs!

\bibliographystyle{plain}
\bibliography{cite}

\end{document}\documentclass{article}
\usepackage[utf8]{inputenc}
\usepackage{booktabs}
\usepackage{amsmath}
\usepackage{natbib}

\title{My Document}
\author{Your Name}
\date{\today}

\begin{document}

\maketitle

\section{Introduction}
This is a sample document demonstrating various \LaTeX{} features including mathematical equations, tables, and formatting. You can cite papers using the citation search feature, and they will be automatically added to your \texttt{cite.bib} file.

\section{Mathematical Expressions}
Here are some example equations:

\subsection{Quadratic Formula}
The quadratic formula is given by:
\[x = \frac{-b \pm \sqrt{b^2 - 4ac}}{2a}\]

\subsection{Matrix Example}
A sample matrix:
\[
\begin{pmatrix}
a_{11} & a_{12} & a_{13} \\
a_{21} & a_{22} & a_{23} \\
a_{31} & a_{32} & a_{33}
\end{pmatrix}
\]

\section{Tables}
Below is a sample table showing data:

\begin{table}[h]
\centering
\begin{tabular}{@{}lrrr@{}}
\toprule
\textbf{Category} & \textbf{Value 1} & \textbf{Value 2} & \textbf{Total} \\
\midrule
Item A & 10.5 & 20.3 & 30.8 \\
Item B & 15.2 & 25.1 & 40.3 \\
Item C & 22.1 & 12.7 & 34.8 \\
\bottomrule
\end{tabular}
\caption{Sample Data Table}
\label{tab:sample}
\end{table}

\section{Lists}
\subsection{Bullet Points}
\begin{itemize}
  \item First item
  \item Second item
  \item Third item with sub-items:
    \begin{itemize}
      \item Sub-item 1
      \item Sub-item 2
    \end{itemize}
\end{itemize}

\subsection{Numbered List}
\begin{enumerate}
  \item First step
  \item Second step
  \item Third step
\end{enumerate}

\section{Citations and References}
When you use the citation search feature, references will be automatically managed. For example, you might cite a paper like this: \cite{example2024}. All citations will be automatically added to your \texttt{cite.bib} file.

\section{Conclusion}
This template demonstrates basic \LaTeX{} features including mathematical formulas, tables, and lists. The bibliography system is now set up to work seamlessly with Bibby AI's citation management. Feel free to modify it for your needs!

\bibliographystyle{plain}
\bibliography{cite}

\end{document}\documentclass{article}
\usepackage[utf8]{inputenc}
\usepackage{booktabs}
\usepackage{amsmath}
\usepackage{natbib}

\title{My Document}
\author{Your Name}
\date{\today}

\begin{document}

\maketitle

\section{Introduction}
This is a sample document demonstrating various \LaTeX{} features including mathematical equations, tables, and formatting. You can cite papers using the citation search feature, and they will be automatically added to your \texttt{cite.bib} file.

\section{Mathematical Expressions}
Here are some example equations:

\subsection{Quadratic Formula}
The quadratic formula is given by:
\[x = \frac{-b \pm \sqrt{b^2 - 4ac}}{2a}\]

\subsection{Matrix Example}
A sample matrix:
\[
\begin{pmatrix}
a_{11} & a_{12} & a_{13} \\
a_{21} & a_{22} & a_{23} \\
a_{31} & a_{32} & a_{33}
\end{pmatrix}
\]

\section{Tables}
Below is a sample table showing data:

\begin{table}[h]
\centering
\begin{tabular}{@{}lrrr@{}}
\toprule
\textbf{Category} & \textbf{Value 1} & \textbf{Value 2} & \textbf{Total} \\
\midrule
Item A & 10.5 & 20.3 & 30.8 \\
Item B & 15.2 & 25.1 & 40.3 \\
Item C & 22.1 & 12.7 & 34.8 \\
\bottomrule
\end{tabular}
\caption{Sample Data Table}
\label{tab:sample}
\end{table}

\section{Lists}
\subsection{Bullet Points}
\begin{itemize}
  \item First item
  \item Second item
  \item Third item with sub-items:
    \begin{itemize}
      \item Sub-item 1
      \item Sub-item 2
    \end{itemize}
\end{itemize}

\subsection{Numbered List}
\begin{enumerate}
  \item First step
  \item Second step
  \item Third step
\end{enumerate}

\section{Citations and References}
When you use the citation search feature, references will be automatically managed. For example, you might cite a paper like this: \cite{example2024}. All citations will be automatically added to your \texttt{cite.bib} file.

\section{Conclusion}
This template demonstrates basic \LaTeX{} features including mathematical formulas, tables, and lists. The bibliography system is now set up to work seamlessly with Bibby AI's citation management. Feel free to modify it for your needs!

\bibliographystyle{plain}
\bibliography{cite}

\end{document}\documentclass{article}
\usepackage[utf8]{inputenc}
\usepackage{booktabs}
\usepackage{amsmath}
\usepackage{natbib}

\title{My Document}
\author{Your Name}
\date{\today}

\begin{document}

\maketitle

\section{Introduction}
This is a sample document demonstrating various \LaTeX{} features including mathematical equations, tables, and formatting. You can cite papers using the citation search feature, and they will be automatically added to your \texttt{cite.bib} file.

\section{Mathematical Expressions}
Here are some example equations:

\subsection{Quadratic Formula}
The quadratic formula is given by:
\[x = \frac{-b \pm \sqrt{b^2 - 4ac}}{2a}\]

\subsection{Matrix Example}
A sample matrix:
\[
\begin{pmatrix}
a_{11} & a_{12} & a_{13} \\
a_{21} & a_{22} & a_{23} \\
a_{31} & a_{32} & a_{33}
\end{pmatrix}
\]

\section{Tables}
Below is a sample table showing data:

\begin{table}[h]
\centering
\begin{tabular}{@{}lrrr@{}}
\toprule
\textbf{Category} & \textbf{Value 1} & \textbf{Value 2} & \textbf{Total} \\
\midrule
Item A & 10.5 & 20.3 & 30.8 \\
Item B & 15.2 & 25.1 & 40.3 \\
Item C & 22.1 & 12.7 & 34.8 \\
\bottomrule
\end{tabular}
\caption{Sample Data Table}
\label{tab:sample}
\end{table}

\section{Lists}
\subsection{Bullet Points}
\begin{itemize}
  \item First item
  \item Second item
  \item Third item with sub-items:
    \begin{itemize}
      \item Sub-item 1
      \item Sub-item 2
    \end{itemize}
\end{itemize}

\subsection{Numbered List}
\begin{enumerate}
  \item First step
  \item Second step
  \item Third step
\end{enumerate}

\section{Citations and References}
When you use the citation search feature, references will be automatically managed. For example, you might cite a paper like this: \cite{example2024}. All citations will be automatically added to your \texttt{cite.bib} file.

\section{Conclusion}
This template demonstrates basic \LaTeX{} features including mathematical formulas, tables, and lists. The bibliography system is now set up to work seamlessly with Bibby AI's citation management. Feel free to modify it for your needs!

\bibliographystyle{plain}
\bibliography{cite}

\end{document}\documentclass{article}
\usepackage[utf8]{inputenc}
\usepackage{booktabs}
\usepackage{amsmath}
\usepackage{natbib}

\title{My Document}
\author{Your Name}
\date{\today}

\begin{document}

\maketitle

\section{Introduction}
This is a sample document demonstrating various \LaTeX{} features including mathematical equations, tables, and formatting. You can cite papers using the citation search feature, and they will be automatically added to your \texttt{cite.bib} file.

\section{Mathematical Expressions}
Here are some example equations:

\subsection{Quadratic Formula}
The quadratic formula is given by:
\[x = \frac{-b \pm \sqrt{b^2 - 4ac}}{2a}\]

\subsection{Matrix Example}
A sample matrix:
\[
\begin{pmatrix}
a_{11} & a_{12} & a_{13} \\
a_{21} & a_{22} & a_{23} \\
a_{31} & a_{32} & a_{33}
\end{pmatrix}
\]

\section{Tables}
Below is a sample table showing data:

\begin{table}[h]
\centering
\begin{tabular}{@{}lrrr@{}}
\toprule
\textbf{Category} & \textbf{Value 1} & \textbf{Value 2} & \textbf{Total} \\
\midrule
Item A & 10.5 & 20.3 & 30.8 \\
Item B & 15.2 & 25.1 & 40.3 \\
Item C & 22.1 & 12.7 & 34.8 \\
\bottomrule
\end{tabular}
\caption{Sample Data Table}
\label{tab:sample}
\end{table}

\section{Lists}
\subsection{Bullet Points}
\begin{itemize}
  \item First item
  \item Second item
  \item Third item with sub-items:
    \begin{itemize}
      \item Sub-item 1
      \item Sub-item 2
    \end{itemize}
\end{itemize}

\subsection{Numbered List}
\begin{enumerate}
  \item First step
  \item Second step
  \item Third step
\end{enumerate}

\section{Citations and References}
When you use the citation search feature, references will be automatically managed. For example, you might cite a paper like this: \cite{example2024}. All citations will be automatically added to your \texttt{cite.bib} file.

\section{Conclusion}
This template demonstrates basic \LaTeX{} features including mathematical formulas, tables, and lists. The bibliography system is now set up to work seamlessly with Bibby AI's citation management. Feel free to modify it for your needs!

\bibliographystyle{plain}
\bibliography{cite}

\end{document}\documentclass{article}
\usepackage[utf8]{inputenc}
\usepackage{booktabs}
\usepackage{amsmath}
\usepackage{natbib}

\title{My Document}
\author{Your Name}
\date{\today}

\begin{document}

\maketitle

\section{Introduction}
This is a sample document demonstrating various \LaTeX{} features including mathematical equations, tables, and formatting. You can cite papers using the citation search feature, and they will be automatically added to your \texttt{cite.bib} file.

\section{Mathematical Expressions}
Here are some example equations:

\subsection{Quadratic Formula}
The quadratic formula is given by:
\[x = \frac{-b \pm \sqrt{b^2 - 4ac}}{2a}\]

\subsection{Matrix Example}
A sample matrix:
\[
\begin{pmatrix}
a_{11} & a_{12} & a_{13} \\
a_{21} & a_{22} & a_{23} \\
a_{31} & a_{32} & a_{33}
\end{pmatrix}
\]

\section{Tables}
Below is a sample table showing data:

\begin{table}[h]
\centering
\begin{tabular}{@{}lrrr@{}}
\toprule
\textbf{Category} & \textbf{Value 1} & \textbf{Value 2} & \textbf{Total} \\
\midrule
Item A & 10.5 & 20.3 & 30.8 \\
Item B & 15.2 & 25.1 & 40.3 \\
Item C & 22.1 & 12.7 & 34.8 \\
\bottomrule
\end{tabular}
\caption{Sample Data Table}
\label{tab:sample}
\end{table}

\section{Lists}
\subsection{Bullet Points}
\begin{itemize}
  \item First item
  \item Second item
  \item Third item with sub-items:
    \begin{itemize}
      \item Sub-item 1
      \item Sub-item 2
    \end{itemize}
\end{itemize}

\subsection{Numbered List}
\begin{enumerate}
  \item First step
  \item Second step
  \item Third step
\end{enumerate}

\section{Citations and References}
When you use the citation search feature, references will be automatically managed. For example, you might cite a paper like this: \cite{example2024}. All citations will be automatically added to your \texttt{cite.bib} file.

\section{Conclusion}
This template demonstrates basic \LaTeX{} features including mathematical formulas, tables, and lists. The bibliography system is now set up to work seamlessly with Bibby AI's citation management. Feel free to modify it for your needs!

\bibliographystyle{plain}
\bibliography{cite}

\end{document}\documentclass{article}
\usepackage[utf8]{inputenc}
\usepackage{booktabs}
\usepackage{amsmath}
\usepackage{natbib}

\title{My Document}
\author{Your Name}
\date{\today}

\begin{document}

\maketitle

\section{Introduction}
This is a sample document demonstrating various \LaTeX{} features including mathematical equations, tables, and formatting. You can cite papers using the citation search feature, and they will be automatically added to your \texttt{cite.bib} file.

\section{Mathematical Expressions}
Here are some example equations:

\subsection{Quadratic Formula}
The quadratic formula is given by:
\[x = \frac{-b \pm \sqrt{b^2 - 4ac}}{2a}\]

\subsection{Matrix Example}
A sample matrix:
\[
\begin{pmatrix}
a_{11} & a_{12} & a_{13} \\
a_{21} & a_{22} & a_{23} \\
a_{31} & a_{32} & a_{33}
\end{pmatrix}
\]

\section{Tables}
Below is a sample table showing data:

\begin{table}[h]
\centering
\begin{tabular}{@{}lrrr@{}}
\toprule
\textbf{Category} & \textbf{Value 1} & \textbf{Value 2} & \textbf{Total} \\
\midrule
Item A & 10.5 & 20.3 & 30.8 \\
Item B & 15.2 & 25.1 & 40.3 \\
Item C & 22.1 & 12.7 & 34.8 \\
\bottomrule
\end{tabular}
\caption{Sample Data Table}
\label{tab:sample}
\end{table}

\section{Lists}
\subsection{Bullet Points}
\begin{itemize}
  \item First item
  \item Second item
  \item Third item with sub-items:
    \begin{itemize}
      \item Sub-item 1
      \item Sub-item 2
    \end{itemize}
\end{itemize}

\subsection{Numbered List}
\begin{enumerate}
  \item First step
  \item Second step
  \item Third step
\end{enumerate}

\section{Citations and References}
When you use the citation search feature, references will be automatically managed. For example, you might cite a paper like this: \cite{example2024}. All citations will be automatically added to your \texttt{cite.bib} file.

\section{Conclusion}
This template demonstrates basic \LaTeX{} features including mathematical formulas, tables, and lists. The bibliography system is now set up to work seamlessly with Bibby AI's citation management. Feel free to modify it for your needs!

\bibliographystyle{plain}
\bibliography{cite}

\end{document}\documentclass{article}
\usepackage[utf8]{inputenc}
\usepackage{booktabs}
\usepackage{amsmath}
\usepackage{natbib}

\title{My Document}
\author{Your Name}
\date{\today}

\begin{document}

\maketitle

\section{Introduction}
This is a sample document demonstrating various \LaTeX{} features including mathematical equations, tables, and formatting. You can cite papers using the citation search feature, and they will be automatically added to your \texttt{cite.bib} file.

\section{Mathematical Expressions}
Here are some example equations:

\subsection{Quadratic Formula}
The quadratic formula is given by:
\[x = \frac{-b \pm \sqrt{b^2 - 4ac}}{2a}\]

\subsection{Matrix Example}
A sample matrix:
\[
\begin{pmatrix}
a_{11} & a_{12} & a_{13} \\
a_{21} & a_{22} & a_{23} \\
a_{31} & a_{32} & a_{33}
\end{pmatrix}
\]

\section{Tables}
Below is a sample table showing data:

\begin{table}[h]
\centering
\begin{tabular}{@{}lrrr@{}}
\toprule
\textbf{Category} & \textbf{Value 1} & \textbf{Value 2} & \textbf{Total} \\
\midrule
Item A & 10.5 & 20.3 & 30.8 \\
Item B & 15.2 & 25.1 & 40.3 \\
Item C & 22.1 & 12.7 & 34.8 \\
\bottomrule
\end{tabular}
\caption{Sample Data Table}
\label{tab:sample}
\end{table}

\section{Lists}
\subsection{Bullet Points}
\begin{itemize}
  \item First item
  \item Second item
  \item Third item with sub-items:
    \begin{itemize}
      \item Sub-item 1
      \item Sub-item 2
    \end{itemize}
\end{itemize}

\subsection{Numbered List}
\begin{enumerate}
  \item First step
  \item Second step
  \item Third step
\end{enumerate}

\section{Citations and References}
When you use the citation search feature, references will be automatically managed. For example, you might cite a paper like this: \cite{example2024}. All citations will be automatically added to your \texttt{cite.bib} file.

\section{Conclusion}
This template demonstrates basic \LaTeX{} features including mathematical formulas, tables, and lists. The bibliography system is now set up to work seamlessly with Bibby AI's citation management. Feel free to modify it for your needs!

\bibliographystyle{plain}
\bibliography{cite}

\end{document}\documentclass{article}
\usepackage[utf8]{inputenc}
\usepackage{booktabs}
\usepackage{amsmath}
\usepackage{natbib}

\title{My Document}
\author{Your Name}
\date{\today}

\begin{document}

\maketitle

\section{Introduction}
This is a sample document demonstrating various \LaTeX{} features including mathematical equations, tables, and formatting. You can cite papers using the citation search feature, and they will be automatically added to your \texttt{cite.bib} file.

\section{Mathematical Expressions}
Here are some example equations:

\subsection{Quadratic Formula}
The quadratic formula is given by:
\[x = \frac{-b \pm \sqrt{b^2 - 4ac}}{2a}\]

\subsection{Matrix Example}
A sample matrix:
\[
\begin{pmatrix}
a_{11} & a_{12} & a_{13} \\
a_{21} & a_{22} & a_{23} \\
a_{31} & a_{32} & a_{33}
\end{pmatrix}
\]

\section{Tables}
Below is a sample table showing data:

\begin{table}[h]
\centering
\begin{tabular}{@{}lrrr@{}}
\toprule
\textbf{Category} & \textbf{Value 1} & \textbf{Value 2} & \textbf{Total} \\
\midrule
Item A & 10.5 & 20.3 & 30.8 \\
Item B & 15.2 & 25.1 & 40.3 \\
Item C & 22.1 & 12.7 & 34.8 \\
\bottomrule
\end{tabular}
\caption{Sample Data Table}
\label{tab:sample}
\end{table}

\section{Lists}
\subsection{Bullet Points}
\begin{itemize}
  \item First item
  \item Second item
  \item Third item with sub-items:
    \begin{itemize}
      \item Sub-item 1
      \item Sub-item 2
    \end{itemize}
\end{itemize}

\subsection{Numbered List}
\begin{enumerate}
  \item First step
  \item Second step
  \item Third step
\end{enumerate}

\section{Citations and References}
When you use the citation search feature, references will be automatically managed. For example, you might cite a paper like this: \cite{example2024}. All citations will be automatically added to your \texttt{cite.bib} file.

\section{Conclusion}
This template demonstrates basic \LaTeX{} features including mathematical formulas, tables, and lists. The bibliography system is now set up to work seamlessly with Bibby AI's citation management. Feel free to modify it for your needs!

\bibliographystyle{plain}
\bibliography{cite}

\end{document}\documentclass{article}
\usepackage[utf8]{inputenc}
\usepackage{booktabs}
\usepackage{amsmath}
\usepackage{natbib}

\title{My Document}
\author{Your Name}
\date{\today}

\begin{document}

\maketitle

\section{Introduction}
This is a sample document demonstrating various \LaTeX{} features including mathematical equations, tables, and formatting. You can cite papers using the citation search feature, and they will be automatically added to your \texttt{cite.bib} file.

\section{Mathematical Expressions}
Here are some example equations:

\subsection{Quadratic Formula}
The quadratic formula is given by:
\[x = \frac{-b \pm \sqrt{b^2 - 4ac}}{2a}\]

\subsection{Matrix Example}
A sample matrix:
\[
\begin{pmatrix}
a_{11} & a_{12} & a_{13} \\
a_{21} & a_{22} & a_{23} \\
a_{31} & a_{32} & a_{33}
\end{pmatrix}
\]

\section{Tables}
Below is a sample table showing data:

\begin{table}[h]
\centering
\begin{tabular}{@{}lrrr@{}}
\toprule
\textbf{Category} & \textbf{Value 1} & \textbf{Value 2} & \textbf{Total} \\
\midrule
Item A & 10.5 & 20.3 & 30.8 \\
Item B & 15.2 & 25.1 & 40.3 \\
Item C & 22.1 & 12.7 & 34.8 \\
\bottomrule
\end{tabular}
\caption{Sample Data Table}
\label{tab:sample}
\end{table}

\section{Lists}
\subsection{Bullet Points}
\begin{itemize}
  \item First item
  \item Second item
  \item Third item with sub-items:
    \begin{itemize}
      \item Sub-item 1
      \item Sub-item 2
    \end{itemize}
\end{itemize}

\subsection{Numbered List}
\begin{enumerate}
  \item First step
  \item Second step
  \item Third step
\end{enumerate}

\section{Citations and References}
When you use the citation search feature, references will be automatically managed. For example, you might cite a paper like this: \cite{example2024}. All citations will be automatically added to your \texttt{cite.bib} file.

\section{Conclusion}
This template demonstrates basic \LaTeX{} features including mathematical formulas, tables, and lists. The bibliography system is now set up to work seamlessly with Bibby AI's citation management. Feel free to modify it for your needs!

\bibliographystyle{plain}
\bibliography{cite}

\end{document}\documentclass{article}
\usepackage[utf8]{inputenc}
\usepackage{booktabs}
\usepackage{amsmath}
\usepackage{natbib}

\title{My Document}
\author{Your Name}
\date{\today}

\begin{document}

\maketitle

\section{Introduction}
This is a sample document demonstrating various \LaTeX{} features including mathematical equations, tables, and formatting. You can cite papers using the citation search feature, and they will be automatically added to your \texttt{cite.bib} file.

\section{Mathematical Expressions}
Here are some example equations:

\subsection{Quadratic Formula}
The quadratic formula is given by:
\[x = \frac{-b \pm \sqrt{b^2 - 4ac}}{2a}\]

\subsection{Matrix Example}
A sample matrix:
\[
\begin{pmatrix}
a_{11} & a_{12} & a_{13} \\
a_{21} & a_{22} & a_{23} \\
a_{31} & a_{32} & a_{33}
\end{pmatrix}
\]

\section{Tables}
Below is a sample table showing data:

\begin{table}[h]
\centering
\begin{tabular}{@{}lrrr@{}}
\toprule
\textbf{Category} & \textbf{Value 1} & \textbf{Value 2} & \textbf{Total} \\
\midrule
Item A & 10.5 & 20.3 & 30.8 \\
Item B & 15.2 & 25.1 & 40.3 \\
Item C & 22.1 & 12.7 & 34.8 \\
\bottomrule
\end{tabular}
\caption{Sample Data Table}
\label{tab:sample}
\end{table}

\section{Lists}
\subsection{Bullet Points}
\begin{itemize}
  \item First item
  \item Second item
  \item Third item with sub-items:
    \begin{itemize}
      \item Sub-item 1
      \item Sub-item 2
    \end{itemize}
\end{itemize}

\subsection{Numbered List}
\begin{enumerate}
  \item First step
  \item Second step
  \item Third step
\end{enumerate}

\section{Citations and References}
When you use the citation search feature, references will be automatically managed. For example, you might cite a paper like this: \cite{example2024}. All citations will be automatically added to your \texttt{cite.bib} file.

\section{Conclusion}
This template demonstrates basic \LaTeX{} features including mathematical formulas, tables, and lists. The bibliography system is now set up to work seamlessly with Bibby AI's citation management. Feel free to modify it for your needs!

\bibliographystyle{plain}
\bibliography{cite}

\end{document}\documentclass{article}
\usepackage[utf8]{inputenc}
\usepackage{booktabs}
\usepackage{amsmath}
\usepackage{natbib}

\title{My Document}
\author{Your Name}
\date{\today}

\begin{document}

\maketitle

\section{Introduction}
This is a sample document demonstrating various \LaTeX{} features including mathematical equations, tables, and formatting. You can cite papers using the citation search feature, and they will be automatically added to your \texttt{cite.bib} file.

\section{Mathematical Expressions}
Here are some example equations:

\subsection{Quadratic Formula}
The quadratic formula is given by:
\[x = \frac{-b \pm \sqrt{b^2 - 4ac}}{2a}\]

\subsection{Matrix Example}
A sample matrix:
\[
\begin{pmatrix}
a_{11} & a_{12} & a_{13} \\
a_{21} & a_{22} & a_{23} \\
a_{31} & a_{32} & a_{33}
\end{pmatrix}
\]

\section{Tables}
Below is a sample table showing data:

\begin{table}[h]
\centering
\begin{tabular}{@{}lrrr@{}}
\toprule
\textbf{Category} & \textbf{Value 1} & \textbf{Value 2} & \textbf{Total} \\
\midrule
Item A & 10.5 & 20.3 & 30.8 \\
Item B & 15.2 & 25.1 & 40.3 \\
Item C & 22.1 & 12.7 & 34.8 \\
\bottomrule
\end{tabular}
\caption{Sample Data Table}
\label{tab:sample}
\end{table}

\section{Lists}
\subsection{Bullet Points}
\begin{itemize}
  \item First item
  \item Second item
  \item Third item with sub-items:
    \begin{itemize}
      \item Sub-item 1
      \item Sub-item 2
    \end{itemize}
\end{itemize}

\subsection{Numbered List}
\begin{enumerate}
  \item First step
  \item Second step
  \item Third step
\end{enumerate}

\section{Citations and References}
When you use the citation search feature, references will be automatically managed. For example, you might cite a paper like this: \cite{example2024}. All citations will be automatically added to your \texttt{cite.bib} file.

\section{Conclusion}
This template demonstrates basic \LaTeX{} features including mathematical formulas, tables, and lists. The bibliography system is now set up to work seamlessly with Bibby AI's citation management. Feel free to modify it for your needs!

\bibliographystyle{plain}
\bibliography{cite}

\end{document}\documentclass{article}
\usepackage[utf8]{inputenc}
\usepackage{booktabs}
\usepackage{amsmath}
\usepackage{natbib}

\title{My Document}
\author{Your Name}
\date{\today}

\begin{document}

\maketitle

\section{Introduction}
This is a sample document demonstrating various \LaTeX{} features including mathematical equations, tables, and formatting. You can cite papers using the citation search feature, and they will be automatically added to your \texttt{cite.bib} file.

\section{Mathematical Expressions}
Here are some example equations:

\subsection{Quadratic Formula}
The quadratic formula is given by:
\[x = \frac{-b \pm \sqrt{b^2 - 4ac}}{2a}\]

\subsection{Matrix Example}
A sample matrix:
\[
\begin{pmatrix}
a_{11} & a_{12} & a_{13} \\
a_{21} & a_{22} & a_{23} \\
a_{31} & a_{32} & a_{33}
\end{pmatrix}
\]

\section{Tables}
Below is a sample table showing data:

\begin{table}[h]
\centering
\begin{tabular}{@{}lrrr@{}}
\toprule
\textbf{Category} & \textbf{Value 1} & \textbf{Value 2} & \textbf{Total} \\
\midrule
Item A & 10.5 & 20.3 & 30.8 \\
Item B & 15.2 & 25.1 & 40.3 \\
Item C & 22.1 & 12.7 & 34.8 \\
\bottomrule
\end{tabular}
\caption{Sample Data Table}
\label{tab:sample}
\end{table}

\section{Lists}
\subsection{Bullet Points}
\begin{itemize}
  \item First item
  \item Second item
  \item Third item with sub-items:
    \begin{itemize}
      \item Sub-item 1
      \item Sub-item 2
    \end{itemize}
\end{itemize}

\subsection{Numbered List}
\begin{enumerate}
  \item First step
  \item Second step
  \item Third step
\end{enumerate}

\section{Citations and References}
When you use the citation search feature, references will be automatically managed. For example, you might cite a paper like this: \cite{example2024}. All citations will be automatically added to your \texttt{cite.bib} file.

\section{Conclusion}
This template demonstrates basic \LaTeX{} features including mathematical formulas, tables, and lists. The bibliography system is now set up to work seamlessly with Bibby AI's citation management. Feel free to modify it for your needs!

\bibliographystyle{plain}
\bibliography{cite}

\end{document}\documentclass{article}
\usepackage[utf8]{inputenc}
\usepackage{booktabs}
\usepackage{amsmath}
\usepackage{natbib}

\title{My Document}
\author{Your Name}
\date{\today}

\begin{document}

\maketitle

\section{Introduction}
This is a sample document demonstrating various \LaTeX{} features including mathematical equations, tables, and formatting. You can cite papers using the citation search feature, and they will be automatically added to your \texttt{cite.bib} file.

\section{Mathematical Expressions}
Here are some example equations:

\subsection{Quadratic Formula}
The quadratic formula is given by:
\[x = \frac{-b \pm \sqrt{b^2 - 4ac}}{2a}\]

\subsection{Matrix Example}
A sample matrix:
\[
\begin{pmatrix}
a_{11} & a_{12} & a_{13} \\
a_{21} & a_{22} & a_{23} \\
a_{31} & a_{32} & a_{33}
\end{pmatrix}
\]

\section{Tables}
Below is a sample table showing data:

\begin{table}[h]
\centering
\begin{tabular}{@{}lrrr@{}}
\toprule
\textbf{Category} & \textbf{Value 1} & \textbf{Value 2} & \textbf{Total} \\
\midrule
Item A & 10.5 & 20.3 & 30.8 \\
Item B & 15.2 & 25.1 & 40.3 \\
Item C & 22.1 & 12.7 & 34.8 \\
\bottomrule
\end{tabular}
\caption{Sample Data Table}
\label{tab:sample}
\end{table}

\section{Lists}
\subsection{Bullet Points}
\begin{itemize}
  \item First item
  \item Second item
  \item Third item with sub-items:
    \begin{itemize}
      \item Sub-item 1
      \item Sub-item 2
    \end{itemize}
\end{itemize}

\subsection{Numbered List}
\begin{enumerate}
  \item First step
  \item Second step
  \item Third step
\end{enumerate}

\section{Citations and References}
When you use the citation search feature, references will be automatically managed. For example, you might cite a paper like this: \cite{example2024}. All citations will be automatically added to your \texttt{cite.bib} file.

\section{Conclusion}
This template demonstrates basic \LaTeX{} features including mathematical formulas, tables, and lists. The bibliography system is now set up to work seamlessly with Bibby AI's citation management. Feel free to modify it for your needs!

\bibliographystyle{plain}
\bibliography{cite}

\end{document}\documentclass{article}
\usepackage[utf8]{inputenc}
\usepackage{booktabs}
\usepackage{amsmath}
\usepackage{natbib}

\title{My Document}
\author{Your Name}
\date{\today}

\begin{document}

\maketitle

\section{Introduction}
This is a sample document demonstrating various \LaTeX{} features including mathematical equations, tables, and formatting. You can cite papers using the citation search feature, and they will be automatically added to your \texttt{cite.bib} file.

\section{Mathematical Expressions}
Here are some example equations:

\subsection{Quadratic Formula}
The quadratic formula is given by:
\[x = \frac{-b \pm \sqrt{b^2 - 4ac}}{2a}\]

\subsection{Matrix Example}
A sample matrix:
\[
\begin{pmatrix}
a_{11} & a_{12} & a_{13} \\
a_{21} & a_{22} & a_{23} \\
a_{31} & a_{32} & a_{33}
\end{pmatrix}
\]

\section{Tables}
Below is a sample table showing data:

\begin{table}[h]
\centering
\begin{tabular}{@{}lrrr@{}}
\toprule
\textbf{Category} & \textbf{Value 1} & \textbf{Value 2} & \textbf{Total} \\
\midrule
Item A & 10.5 & 20.3 & 30.8 \\
Item B & 15.2 & 25.1 & 40.3 \\
Item C & 22.1 & 12.7 & 34.8 \\
\bottomrule
\end{tabular}
\caption{Sample Data Table}
\label{tab:sample}
\end{table}

\section{Lists}
\subsection{Bullet Points}
\begin{itemize}
  \item First item
  \item Second item
  \item Third item with sub-items:
    \begin{itemize}
      \item Sub-item 1
      \item Sub-item 2
    \end{itemize}
\end{itemize}

\subsection{Numbered List}
\begin{enumerate}
  \item First step
  \item Second step
  \item Third step
\end{enumerate}

\section{Citations and References}
When you use the citation search feature, references will be automatically managed. For example, you might cite a paper like this: \cite{example2024}. All citations will be automatically added to your \texttt{cite.bib} file.

\section{Conclusion}
This template demonstrates basic \LaTeX{} features including mathematical formulas, tables, and lists. The bibliography system is now set up to work seamlessly with Bibby AI's citation management. Feel free to modify it for your needs!

\bibliographystyle{plain}
\bibliography{cite}

\end{document}\documentclass{article}
\usepackage[utf8]{inputenc}
\usepackage{booktabs}
\usepackage{amsmath}
\usepackage{natbib}

\title{My Document}
\author{Your Name}
\date{\today}

\begin{document}

\maketitle

\section{Introduction}
This is a sample document demonstrating various \LaTeX{} features including mathematical equations, tables, and formatting. You can cite papers using the citation search feature, and they will be automatically added to your \texttt{cite.bib} file.

\section{Mathematical Expressions}
Here are some example equations:

\subsection{Quadratic Formula}
The quadratic formula is given by:
\[x = \frac{-b \pm \sqrt{b^2 - 4ac}}{2a}\]

\subsection{Matrix Example}
A sample matrix:
\[
\begin{pmatrix}
a_{11} & a_{12} & a_{13} \\
a_{21} & a_{22} & a_{23} \\
a_{31} & a_{32} & a_{33}
\end{pmatrix}
\]

\section{Tables}
Below is a sample table showing data:

\begin{table}[h]
\centering
\begin{tabular}{@{}lrrr@{}}
\toprule
\textbf{Category} & \textbf{Value 1} & \textbf{Value 2} & \textbf{Total} \\
\midrule
Item A & 10.5 & 20.3 & 30.8 \\
Item B & 15.2 & 25.1 & 40.3 \\
Item C & 22.1 & 12.7 & 34.8 \\
\bottomrule
\end{tabular}
\caption{Sample Data Table}
\label{tab:sample}
\end{table}

\section{Lists}
\subsection{Bullet Points}
\begin{itemize}
  \item First item
  \item Second item
  \item Third item with sub-items:
    \begin{itemize}
      \item Sub-item 1
      \item Sub-item 2
    \end{itemize}
\end{itemize}

\subsection{Numbered List}
\begin{enumerate}
  \item First step
  \item Second step
  \item Third step
\end{enumerate}

\section{Citations and References}
When you use the citation search feature, references will be automatically managed. For example, you might cite a paper like this: \cite{example2024}. All citations will be automatically added to your \texttt{cite.bib} file.

\section{Conclusion}
This template demonstrates basic \LaTeX{} features including mathematical formulas, tables, and lists. The bibliography system is now set up to work seamlessly with Bibby AI's citation management. Feel free to modify it for your needs!

\bibliographystyle{plain}
\bibliography{cite}

\end{document}\documentclass{article}
\usepackage[utf8]{inputenc}
\usepackage{booktabs}
\usepackage{amsmath}
\usepackage{natbib}

\title{My Document}
\author{Your Name}
\date{\today}

\begin{document}

\maketitle

\section{Introduction}
This is a sample document demonstrating various \LaTeX{} features including mathematical equations, tables, and formatting. You can cite papers using the citation search feature, and they will be automatically added to your \texttt{cite.bib} file.

\section{Mathematical Expressions}
Here are some example equations:

\subsection{Quadratic Formula}
The quadratic formula is given by:
\[x = \frac{-b \pm \sqrt{b^2 - 4ac}}{2a}\]

\subsection{Matrix Example}
A sample matrix:
\[
\begin{pmatrix}
a_{11} & a_{12} & a_{13} \\
a_{21} & a_{22} & a_{23} \\
a_{31} & a_{32} & a_{33}
\end{pmatrix}
\]

\section{Tables}
Below is a sample table showing data:

\begin{table}[h]
\centering
\begin{tabular}{@{}lrrr@{}}
\toprule
\textbf{Category} & \textbf{Value 1} & \textbf{Value 2} & \textbf{Total} \\
\midrule
Item A & 10.5 & 20.3 & 30.8 \\
Item B & 15.2 & 25.1 & 40.3 \\
Item C & 22.1 & 12.7 & 34.8 \\
\bottomrule
\end{tabular}
\caption{Sample Data Table}
\label{tab:sample}
\end{table}

\section{Lists}
\subsection{Bullet Points}
\begin{itemize}
  \item First item
  \item Second item
  \item Third item with sub-items:
    \begin{itemize}
      \item Sub-item 1
      \item Sub-item 2
    \end{itemize}
\end{itemize}

\subsection{Numbered List}
\begin{enumerate}
  \item First step
  \item Second step
  \item Third step
\end{enumerate}

\section{Citations and References}
When you use the citation search feature, references will be automatically managed. For example, you might cite a paper like this: \cite{example2024}. All citations will be automatically added to your \texttt{cite.bib} file.

\section{Conclusion}
This template demonstrates basic \LaTeX{} features including mathematical formulas, tables, and lists. The bibliography system is now set up to work seamlessly with Bibby AI's citation management. Feel free to modify it for your needs!

\bibliographystyle{plain}
\bibliography{cite}

\end{document}\documentclass{article}
\usepackage[utf8]{inputenc}
\usepackage{booktabs}
\usepackage{amsmath}
\usepackage{natbib}

\title{My Document}
\author{Your Name}
\date{\today}

\begin{document}

\maketitle

\section{Introduction}
This is a sample document demonstrating various \LaTeX{} features including mathematical equations, tables, and formatting. You can cite papers using the citation search feature, and they will be automatically added to your \texttt{cite.bib} file.

\section{Mathematical Expressions}
Here are some example equations:

\subsection{Quadratic Formula}
The quadratic formula is given by:
\[x = \frac{-b \pm \sqrt{b^2 - 4ac}}{2a}\]

\subsection{Matrix Example}
A sample matrix:
\[
\begin{pmatrix}
a_{11} & a_{12} & a_{13} \\
a_{21} & a_{22} & a_{23} \\
a_{31} & a_{32} & a_{33}
\end{pmatrix}
\]

\section{Tables}
Below is a sample table showing data:

\begin{table}[h]
\centering
\begin{tabular}{@{}lrrr@{}}
\toprule
\textbf{Category} & \textbf{Value 1} & \textbf{Value 2} & \textbf{Total} \\
\midrule
Item A & 10.5 & 20.3 & 30.8 \\
Item B & 15.2 & 25.1 & 40.3 \\
Item C & 22.1 & 12.7 & 34.8 \\
\bottomrule
\end{tabular}
\caption{Sample Data Table}
\label{tab:sample}
\end{table}

\section{Lists}
\subsection{Bullet Points}
\begin{itemize}
  \item First item
  \item Second item
  \item Third item with sub-items:
    \begin{itemize}
      \item Sub-item 1
      \item Sub-item 2
    \end{itemize}
\end{itemize}

\subsection{Numbered List}
\begin{enumerate}
  \item First step
  \item Second step
  \item Third step
\end{enumerate}

\section{Citations and References}
When you use the citation search feature, references will be automatically managed. For example, you might cite a paper like this: \cite{example2024}. All citations will be automatically added to your \texttt{cite.bib} file.

\section{Conclusion}
This template demonstrates basic \LaTeX{} features including mathematical formulas, tables, and lists. The bibliography system is now set up to work seamlessly with Bibby AI's citation management. Feel free to modify it for your needs!

\bibliographystyle{plain}
\bibliography{cite}

\end{document}\documentclass{article}
\usepackage[utf8]{inputenc}
\usepackage{booktabs}
\usepackage{amsmath}
\usepackage{natbib}

\title{My Document}
\author{Your Name}
\date{\today}

\begin{document}

\maketitle

\section{Introduction}
This is a sample document demonstrating various \LaTeX{} features including mathematical equations, tables, and formatting. You can cite papers using the citation search feature, and they will be automatically added to your \texttt{cite.bib} file.

\section{Mathematical Expressions}
Here are some example equations:

\subsection{Quadratic Formula}
The quadratic formula is given by:
\[x = \frac{-b \pm \sqrt{b^2 - 4ac}}{2a}\]

\subsection{Matrix Example}
A sample matrix:
\[
\begin{pmatrix}
a_{11} & a_{12} & a_{13} \\
a_{21} & a_{22} & a_{23} \\
a_{31} & a_{32} & a_{33}
\end{pmatrix}
\]

\section{Tables}
Below is a sample table showing data:

\begin{table}[h]
\centering
\begin{tabular}{@{}lrrr@{}}
\toprule
\textbf{Category} & \textbf{Value 1} & \textbf{Value 2} & \textbf{Total} \\
\midrule
Item A & 10.5 & 20.3 & 30.8 \\
Item B & 15.2 & 25.1 & 40.3 \\
Item C & 22.1 & 12.7 & 34.8 \\
\bottomrule
\end{tabular}
\caption{Sample Data Table}
\label{tab:sample}
\end{table}

\section{Lists}
\subsection{Bullet Points}
\begin{itemize}
  \item First item
  \item Second item
  \item Third item with sub-items:
    \begin{itemize}
      \item Sub-item 1
      \item Sub-item 2
    \end{itemize}
\end{itemize}

\subsection{Numbered List}
\begin{enumerate}
  \item First step
  \item Second step
  \item Third step
\end{enumerate}

\section{Citations and References}
When you use the citation search feature, references will be automatically managed. For example, you might cite a paper like this: \cite{example2024}. All citations will be automatically added to your \texttt{cite.bib} file.

\section{Conclusion}
This template demonstrates basic \LaTeX{} features including mathematical formulas, tables, and lists. The bibliography system is now set up to work seamlessly with Bibby AI's citation management. Feel free to modify it for your needs!

\bibliographystyle{plain}
\bibliography{cite}

\end{document}\documentclass{article}
\usepackage[utf8]{inputenc}
\usepackage{booktabs}
\usepackage{amsmath}
\usepackage{natbib}

\title{My Document}
\author{Your Name}
\date{\today}

\begin{document}

\maketitle

\section{Introduction}
This is a sample document demonstrating various \LaTeX{} features including mathematical equations, tables, and formatting. You can cite papers using the citation search feature, and they will be automatically added to your \texttt{cite.bib} file.

\section{Mathematical Expressions}
Here are some example equations:

\subsection{Quadratic Formula}
The quadratic formula is given by:
\[x = \frac{-b \pm \sqrt{b^2 - 4ac}}{2a}\]

\subsection{Matrix Example}
A sample matrix:
\[
\begin{pmatrix}
a_{11} & a_{12} & a_{13} \\
a_{21} & a_{22} & a_{23} \\
a_{31} & a_{32} & a_{33}
\end{pmatrix}
\]

\section{Tables}
Below is a sample table showing data:

\begin{table}[h]
\centering
\begin{tabular}{@{}lrrr@{}}
\toprule
\textbf{Category} & \textbf{Value 1} & \textbf{Value 2} & \textbf{Total} \\
\midrule
Item A & 10.5 & 20.3 & 30.8 \\
Item B & 15.2 & 25.1 & 40.3 \\
Item C & 22.1 & 12.7 & 34.8 \\
\bottomrule
\end{tabular}
\caption{Sample Data Table}
\label{tab:sample}
\end{table}

\section{Lists}
\subsection{Bullet Points}
\begin{itemize}
  \item First item
  \item Second item
  \item Third item with sub-items:
    \begin{itemize}
      \item Sub-item 1
      \item Sub-item 2
    \end{itemize}
\end{itemize}

\subsection{Numbered List}
\begin{enumerate}
  \item First step
  \item Second step
  \item Third step
\end{enumerate}

\section{Citations and References}
When you use the citation search feature, references will be automatically managed. For example, you might cite a paper like this: \cite{example2024}. All citations will be automatically added to your \texttt{cite.bib} file.

\section{Conclusion}
This template demonstrates basic \LaTeX{} features including mathematical formulas, tables, and lists. The bibliography system is now set up to work seamlessly with Bibby AI's citation management. Feel free to modify it for your needs!

\bibliographystyle{plain}
\bibliography{cite}

\end{document}\documentclass{article}
\usepackage[utf8]{inputenc}
\usepackage{booktabs}
\usepackage{amsmath}
\usepackage{natbib}

\title{My Document}
\author{Your Name}
\date{\today}

\begin{document}

\maketitle

\section{Introduction}
This is a sample document demonstrating various \LaTeX{} features including mathematical equations, tables, and formatting. You can cite papers using the citation search feature, and they will be automatically added to your \texttt{cite.bib} file.

\section{Mathematical Expressions}
Here are some example equations:

\subsection{Quadratic Formula}
The quadratic formula is given by:
\[x = \frac{-b \pm \sqrt{b^2 - 4ac}}{2a}\]

\subsection{Matrix Example}
A sample matrix:
\[
\begin{pmatrix}
a_{11} & a_{12} & a_{13} \\
a_{21} & a_{22} & a_{23} \\
a_{31} & a_{32} & a_{33}
\end{pmatrix}
\]

\section{Tables}
Below is a sample table showing data:

\begin{table}[h]
\centering
\begin{tabular}{@{}lrrr@{}}
\toprule
\textbf{Category} & \textbf{Value 1} & \textbf{Value 2} & \textbf{Total} \\
\midrule
Item A & 10.5 & 20.3 & 30.8 \\
Item B & 15.2 & 25.1 & 40.3 \\
Item C & 22.1 & 12.7 & 34.8 \\
\bottomrule
\end{tabular}
\caption{Sample Data Table}
\label{tab:sample}
\end{table}

\section{Lists}
\subsection{Bullet Points}
\begin{itemize}
  \item First item
  \item Second item
  \item Third item with sub-items:
    \begin{itemize}
      \item Sub-item 1
      \item Sub-item 2
    \end{itemize}
\end{itemize}

\subsection{Numbered List}
\begin{enumerate}
  \item First step
  \item Second step
  \item Third step
\end{enumerate}

\section{Citations and References}
When you use the citation search feature, references will be automatically managed. For example, you might cite a paper like this: \cite{example2024}. All citations will be automatically added to your \texttt{cite.bib} file.

\section{Conclusion}
This template demonstrates basic \LaTeX{} features including mathematical formulas, tables, and lists. The bibliography system is now set up to work seamlessly with Bibby AI's citation management. Feel free to modify it for your needs!

\bibliographystyle{plain}
\bibliography{cite}

\end{document}\documentclass{article}
\usepackage[utf8]{inputenc}
\usepackage{booktabs}
\usepackage{amsmath}
\usepackage{natbib}

\title{My Document}
\author{Your Name}
\date{\today}

\begin{document}

\maketitle

\section{Introduction}
This is a sample document demonstrating various \LaTeX{} features including mathematical equations, tables, and formatting. You can cite papers using the citation search feature, and they will be automatically added to your \texttt{cite.bib} file.

\section{Mathematical Expressions}
Here are some example equations:

\subsection{Quadratic Formula}
The quadratic formula is given by:
\[x = \frac{-b \pm \sqrt{b^2 - 4ac}}{2a}\]

\subsection{Matrix Example}
A sample matrix:
\[
\begin{pmatrix}
a_{11} & a_{12} & a_{13} \\
a_{21} & a_{22} & a_{23} \\
a_{31} & a_{32} & a_{33}
\end{pmatrix}
\]

\section{Tables}
Below is a sample table showing data:

\begin{table}[h]
\centering
\begin{tabular}{@{}lrrr@{}}
\toprule
\textbf{Category} & \textbf{Value 1} & \textbf{Value 2} & \textbf{Total} \\
\midrule
Item A & 10.5 & 20.3 & 30.8 \\
Item B & 15.2 & 25.1 & 40.3 \\
Item C & 22.1 & 12.7 & 34.8 \\
\bottomrule
\end{tabular}
\caption{Sample Data Table}
\label{tab:sample}
\end{table}

\section{Lists}
\subsection{Bullet Points}
\begin{itemize}
  \item First item
  \item Second item
  \item Third item with sub-items:
    \begin{itemize}
      \item Sub-item 1
      \item Sub-item 2
    \end{itemize}
\end{itemize}

\subsection{Numbered List}
\begin{enumerate}
  \item First step
  \item Second step
  \item Third step
\end{enumerate}

\section{Citations and References}
When you use the citation search feature, references will be automatically managed. For example, you might cite a paper like this: \cite{example2024}. All citations will be automatically added to your \texttt{cite.bib} file.

\section{Conclusion}
This template demonstrates basic \LaTeX{} features including mathematical formulas, tables, and lists. The bibliography system is now set up to work seamlessly with Bibby AI's citation management. Feel free to modify it for your needs!

\bibliographystyle{plain}
\bibliography{cite}

\end{document}\documentclass{article}
\usepackage[utf8]{inputenc}
\usepackage{booktabs}
\usepackage{amsmath}
\usepackage{natbib}

\title{My Document}
\author{Your Name}
\date{\today}

\begin{document}

\maketitle

\section{Introduction}
This is a sample document demonstrating various \LaTeX{} features including mathematical equations, tables, and formatting. You can cite papers using the citation search feature, and they will be automatically added to your \texttt{cite.bib} file.

\section{Mathematical Expressions}
Here are some example equations:

\subsection{Quadratic Formula}
The quadratic formula is given by:
\[x = \frac{-b \pm \sqrt{b^2 - 4ac}}{2a}\]

\subsection{Matrix Example}
A sample matrix:
\[
\begin{pmatrix}
a_{11} & a_{12} & a_{13} \\
a_{21} & a_{22} & a_{23} \\
a_{31} & a_{32} & a_{33}
\end{pmatrix}
\]

\section{Tables}
Below is a sample table showing data:

\begin{table}[h]
\centering
\begin{tabular}{@{}lrrr@{}}
\toprule
\textbf{Category} & \textbf{Value 1} & \textbf{Value 2} & \textbf{Total} \\
\midrule
Item A & 10.5 & 20.3 & 30.8 \\
Item B & 15.2 & 25.1 & 40.3 \\
Item C & 22.1 & 12.7 & 34.8 \\
\bottomrule
\end{tabular}
\caption{Sample Data Table}
\label{tab:sample}
\end{table}

\section{Lists}
\subsection{Bullet Points}
\begin{itemize}
  \item First item
  \item Second item
  \item Third item with sub-items:
    \begin{itemize}
      \item Sub-item 1
      \item Sub-item 2
    \end{itemize}
\end{itemize}

\subsection{Numbered List}
\begin{enumerate}
  \item First step
  \item Second step
  \item Third step
\end{enumerate}

\section{Citations and References}
When you use the citation search feature, references will be automatically managed. For example, you might cite a paper like this: \cite{example2024}. All citations will be automatically added to your \texttt{cite.bib} file.

\section{Conclusion}
This template demonstrates basic \LaTeX{} features including mathematical formulas, tables, and lists. The bibliography system is now set up to work seamlessly with Bibby AI's citation management. Feel free to modify it for your needs!

\bibliographystyle{plain}
\bibliography{cite}

\end{document}\documentclass{article}
\usepackage[utf8]{inputenc}
\usepackage{booktabs}
\usepackage{amsmath}
\usepackage{natbib}

\title{My Document}
\author{Your Name}
\date{\today}

\begin{document}

\maketitle

\section{Introduction}
This is a sample document demonstrating various \LaTeX{} features including mathematical equations, tables, and formatting. You can cite papers using the citation search feature, and they will be automatically added to your \texttt{cite.bib} file.

\section{Mathematical Expressions}
Here are some example equations:

\subsection{Quadratic Formula}
The quadratic formula is given by:
\[x = \frac{-b \pm \sqrt{b^2 - 4ac}}{2a}\]

\subsection{Matrix Example}
A sample matrix:
\[
\begin{pmatrix}
a_{11} & a_{12} & a_{13} \\
a_{21} & a_{22} & a_{23} \\
a_{31} & a_{32} & a_{33}
\end{pmatrix}
\]

\section{Tables}
Below is a sample table showing data:

\begin{table}[h]
\centering
\begin{tabular}{@{}lrrr@{}}
\toprule
\textbf{Category} & \textbf{Value 1} & \textbf{Value 2} & \textbf{Total} \\
\midrule
Item A & 10.5 & 20.3 & 30.8 \\
Item B & 15.2 & 25.1 & 40.3 \\
Item C & 22.1 & 12.7 & 34.8 \\
\bottomrule
\end{tabular}
\caption{Sample Data Table}
\label{tab:sample}
\end{table}

\section{Lists}
\subsection{Bullet Points}
\begin{itemize}
  \item First item
  \item Second item
  \item Third item with sub-items:
    \begin{itemize}
      \item Sub-item 1
      \item Sub-item 2
    \end{itemize}
\end{itemize}

\subsection{Numbered List}
\begin{enumerate}
  \item First step
  \item Second step
  \item Third step
\end{enumerate}

\section{Citations and References}
When you use the citation search feature, references will be automatically managed. For example, you might cite a paper like this: \cite{example2024}. All citations will be automatically added to your \texttt{cite.bib} file.

\section{Conclusion}
This template demonstrates basic \LaTeX{} features including mathematical formulas, tables, and lists. The bibliography system is now set up to work seamlessly with Bibby AI's citation management. Feel free to modify it for your needs!

\bibliographystyle{plain}
\bibliography{cite}

\end{document}\documentclass{article}
\usepackage[utf8]{inputenc}
\usepackage{booktabs}
\usepackage{amsmath}
\usepackage{natbib}

\title{My Document}
\author{Your Name}
\date{\today}

\begin{document}

\maketitle

\section{Introduction}
This is a sample document demonstrating various \LaTeX{} features including mathematical equations, tables, and formatting. You can cite papers using the citation search feature, and they will be automatically added to your \texttt{cite.bib} file.

\section{Mathematical Expressions}
Here are some example equations:

\subsection{Quadratic Formula}
The quadratic formula is given by:
\[x = \frac{-b \pm \sqrt{b^2 - 4ac}}{2a}\]

\subsection{Matrix Example}
A sample matrix:
\[
\begin{pmatrix}
a_{11} & a_{12} & a_{13} \\
a_{21} & a_{22} & a_{23} \\
a_{31} & a_{32} & a_{33}
\end{pmatrix}
\]

\section{Tables}
Below is a sample table showing data:

\begin{table}[h]
\centering
\begin{tabular}{@{}lrrr@{}}
\toprule
\textbf{Category} & \textbf{Value 1} & \textbf{Value 2} & \textbf{Total} \\
\midrule
Item A & 10.5 & 20.3 & 30.8 \\
Item B & 15.2 & 25.1 & 40.3 \\
Item C & 22.1 & 12.7 & 34.8 \\
\bottomrule
\end{tabular}
\caption{Sample Data Table}
\label{tab:sample}
\end{table}

\section{Lists}
\subsection{Bullet Points}
\begin{itemize}
  \item First item
  \item Second item
  \item Third item with sub-items:
    \begin{itemize}
      \item Sub-item 1
      \item Sub-item 2
    \end{itemize}
\end{itemize}

\subsection{Numbered List}
\begin{enumerate}
  \item First step
  \item Second step
  \item Third step
\end{enumerate}

\section{Citations and References}
When you use the citation search feature, references will be automatically managed. For example, you might cite a paper like this: \cite{example2024}. All citations will be automatically added to your \texttt{cite.bib} file.

\section{Conclusion}
This template demonstrates basic \LaTeX{} features including mathematical formulas, tables, and lists. The bibliography system is now set up to work seamlessly with Bibby AI's citation management. Feel free to modify it for your needs!

\bibliographystyle{plain}
\bibliography{cite}

\end{document}\documentclass{article}
\usepackage[utf8]{inputenc}
\usepackage{booktabs}
\usepackage{amsmath}
\usepackage{natbib}

\title{My Document}
\author{Your Name}
\date{\today}

\begin{document}

\maketitle

\section{Introduction}
This is a sample document demonstrating various \LaTeX{} features including mathematical equations, tables, and formatting. You can cite papers using the citation search feature, and they will be automatically added to your \texttt{cite.bib} file.

\section{Mathematical Expressions}
Here are some example equations:

\subsection{Quadratic Formula}
The quadratic formula is given by:
\[x = \frac{-b \pm \sqrt{b^2 - 4ac}}{2a}\]

\subsection{Matrix Example}
A sample matrix:
\[
\begin{pmatrix}
a_{11} & a_{12} & a_{13} \\
a_{21} & a_{22} & a_{23} \\
a_{31} & a_{32} & a_{33}
\end{pmatrix}
\]

\section{Tables}
Below is a sample table showing data:

\begin{table}[h]
\centering
\begin{tabular}{@{}lrrr@{}}
\toprule
\textbf{Category} & \textbf{Value 1} & \textbf{Value 2} & \textbf{Total} \\
\midrule
Item A & 10.5 & 20.3 & 30.8 \\
Item B & 15.2 & 25.1 & 40.3 \\
Item C & 22.1 & 12.7 & 34.8 \\
\bottomrule
\end{tabular}
\caption{Sample Data Table}
\label{tab:sample}
\end{table}

\section{Lists}
\subsection{Bullet Points}
\begin{itemize}
  \item First item
  \item Second item
  \item Third item with sub-items:
    \begin{itemize}
      \item Sub-item 1
      \item Sub-item 2
    \end{itemize}
\end{itemize}

\subsection{Numbered List}
\begin{enumerate}
  \item First step
  \item Second step
  \item Third step
\end{enumerate}

\section{Citations and References}
When you use the citation search feature, references will be automatically managed. For example, you might cite a paper like this: \cite{example2024}. All citations will be automatically added to your \texttt{cite.bib} file.

\section{Conclusion}
This template demonstrates basic \LaTeX{} features including mathematical formulas, tables, and lists. The bibliography system is now set up to work seamlessly with Bibby AI's citation management. Feel free to modify it for your needs!

\bibliographystyle{plain}
\bibliography{cite}

\end{document}\documentclass{article}
\usepackage[utf8]{inputenc}
\usepackage{booktabs}
\usepackage{amsmath}
\usepackage{natbib}

\title{My Document}
\author{Your Name}
\date{\today}

\begin{document}

\maketitle

\section{Introduction}
This is a sample document demonstrating various \LaTeX{} features including mathematical equations, tables, and formatting. You can cite papers using the citation search feature, and they will be automatically added to your \texttt{cite.bib} file.

\section{Mathematical Expressions}
Here are some example equations:

\subsection{Quadratic Formula}
The quadratic formula is given by:
\[x = \frac{-b \pm \sqrt{b^2 - 4ac}}{2a}\]

\subsection{Matrix Example}
A sample matrix:
\[
\begin{pmatrix}
a_{11} & a_{12} & a_{13} \\
a_{21} & a_{22} & a_{23} \\
a_{31} & a_{32} & a_{33}
\end{pmatrix}
\]

\section{Tables}
Below is a sample table showing data:

\begin{table}[h]
\centering
\begin{tabular}{@{}lrrr@{}}
\toprule
\textbf{Category} & \textbf{Value 1} & \textbf{Value 2} & \textbf{Total} \\
\midrule
Item A & 10.5 & 20.3 & 30.8 \\
Item B & 15.2 & 25.1 & 40.3 \\
Item C & 22.1 & 12.7 & 34.8 \\
\bottomrule
\end{tabular}
\caption{Sample Data Table}
\label{tab:sample}
\end{table}

\section{Lists}
\subsection{Bullet Points}
\begin{itemize}
  \item First item
  \item Second item
  \item Third item with sub-items:
    \begin{itemize}
      \item Sub-item 1
      \item Sub-item 2
    \end{itemize}
\end{itemize}

\subsection{Numbered List}
\begin{enumerate}
  \item First step
  \item Second step
  \item Third step
\end{enumerate}

\section{Citations and References}
When you use the citation search feature, references will be automatically managed. For example, you might cite a paper like this: \cite{example2024}. All citations will be automatically added to your \texttt{cite.bib} file.

\section{Conclusion}
This template demonstrates basic \LaTeX{} features including mathematical formulas, tables, and lists. The bibliography system is now set up to work seamlessly with Bibby AI's citation management. Feel free to modify it for your needs!

\bibliographystyle{plain}
\bibliography{cite}

\end{document}\documentclass{article}
\usepackage[utf8]{inputenc}
\usepackage{booktabs}
\usepackage{amsmath}
\usepackage{natbib}

\title{My Document}
\author{Your Name}
\date{\today}

\begin{document}

\maketitle

\section{Introduction}
This is a sample document demonstrating various \LaTeX{} features including mathematical equations, tables, and formatting. You can cite papers using the citation search feature, and they will be automatically added to your \texttt{cite.bib} file.

\section{Mathematical Expressions}
Here are some example equations:

\subsection{Quadratic Formula}
The quadratic formula is given by:
\[x = \frac{-b \pm \sqrt{b^2 - 4ac}}{2a}\]

\subsection{Matrix Example}
A sample matrix:
\[
\begin{pmatrix}
a_{11} & a_{12} & a_{13} \\
a_{21} & a_{22} & a_{23} \\
a_{31} & a_{32} & a_{33}
\end{pmatrix}
\]

\section{Tables}
Below is a sample table showing data:

\begin{table}[h]
\centering
\begin{tabular}{@{}lrrr@{}}
\toprule
\textbf{Category} & \textbf{Value 1} & \textbf{Value 2} & \textbf{Total} \\
\midrule
Item A & 10.5 & 20.3 & 30.8 \\
Item B & 15.2 & 25.1 & 40.3 \\
Item C & 22.1 & 12.7 & 34.8 \\
\bottomrule
\end{tabular}
\caption{Sample Data Table}
\label{tab:sample}
\end{table}

\section{Lists}
\subsection{Bullet Points}
\begin{itemize}
  \item First item
  \item Second item
  \item Third item with sub-items:
    \begin{itemize}
      \item Sub-item 1
      \item Sub-item 2
    \end{itemize}
\end{itemize}

\subsection{Numbered List}
\begin{enumerate}
  \item First step
  \item Second step
  \item Third step
\end{enumerate}

\section{Citations and References}
When you use the citation search feature, references will be automatically managed. For example, you might cite a paper like this: \cite{example2024}. All citations will be automatically added to your \texttt{cite.bib} file.

\section{Conclusion}
This template demonstrates basic \LaTeX{} features including mathematical formulas, tables, and lists. The bibliography system is now set up to work seamlessly with Bibby AI's citation management. Feel free to modify it for your needs!

\bibliographystyle{plain}
\bibliography{cite}

\end{document}\documentclass{article}
\usepackage[utf8]{inputenc}
\usepackage{booktabs}
\usepackage{amsmath}
\usepackage{natbib}

\title{My Document}
\author{Your Name}
\date{\today}

\begin{document}

\maketitle

\section{Introduction}
This is a sample document demonstrating various \LaTeX{} features including mathematical equations, tables, and formatting. You can cite papers using the citation search feature, and they will be automatically added to your \texttt{cite.bib} file.

\section{Mathematical Expressions}
Here are some example equations:

\subsection{Quadratic Formula}
The quadratic formula is given by:
\[x = \frac{-b \pm \sqrt{b^2 - 4ac}}{2a}\]

\subsection{Matrix Example}
A sample matrix:
\[
\begin{pmatrix}
a_{11} & a_{12} & a_{13} \\
a_{21} & a_{22} & a_{23} \\
a_{31} & a_{32} & a_{33}
\end{pmatrix}
\]

\section{Tables}
Below is a sample table showing data:

\begin{table}[h]
\centering
\begin{tabular}{@{}lrrr@{}}
\toprule
\textbf{Category} & \textbf{Value 1} & \textbf{Value 2} & \textbf{Total} \\
\midrule
Item A & 10.5 & 20.3 & 30.8 \\
Item B & 15.2 & 25.1 & 40.3 \\
Item C & 22.1 & 12.7 & 34.8 \\
\bottomrule
\end{tabular}
\caption{Sample Data Table}
\label{tab:sample}
\end{table}

\section{Lists}
\subsection{Bullet Points}
\begin{itemize}
  \item First item
  \item Second item
  \item Third item with sub-items:
    \begin{itemize}
      \item Sub-item 1
      \item Sub-item 2
    \end{itemize}
\end{itemize}

\subsection{Numbered List}
\begin{enumerate}
  \item First step
  \item Second step
  \item Third step
\end{enumerate}

\section{Citations and References}
When you use the citation search feature, references will be automatically managed. For example, you might cite a paper like this: \cite{example2024}. All citations will be automatically added to your \texttt{cite.bib} file.

\section{Conclusion}
This template demonstrates basic \LaTeX{} features including mathematical formulas, tables, and lists. The bibliography system is now set up to work seamlessly with Bibby AI's citation management. Feel free to modify it for your needs!

\bibliographystyle{plain}
\bibliography{cite}

\end{document}\documentclass{article}
\usepackage[utf8]{inputenc}
\usepackage{booktabs}
\usepackage{amsmath}
\usepackage{natbib}

\title{My Document}
\author{Your Name}
\date{\today}

\begin{document}

\maketitle

\section{Introduction}
This is a sample document demonstrating various \LaTeX{} features including mathematical equations, tables, and formatting. You can cite papers using the citation search feature, and they will be automatically added to your \texttt{cite.bib} file.

\section{Mathematical Expressions}
Here are some example equations:

\subsection{Quadratic Formula}
The quadratic formula is given by:
\[x = \frac{-b \pm \sqrt{b^2 - 4ac}}{2a}\]

\subsection{Matrix Example}
A sample matrix:
\[
\begin{pmatrix}
a_{11} & a_{12} & a_{13} \\
a_{21} & a_{22} & a_{23} \\
a_{31} & a_{32} & a_{33}
\end{pmatrix}
\]

\section{Tables}
Below is a sample table showing data:

\begin{table}[h]
\centering
\begin{tabular}{@{}lrrr@{}}
\toprule
\textbf{Category} & \textbf{Value 1} & \textbf{Value 2} & \textbf{Total} \\
\midrule
Item A & 10.5 & 20.3 & 30.8 \\
Item B & 15.2 & 25.1 & 40.3 \\
Item C & 22.1 & 12.7 & 34.8 \\
\bottomrule
\end{tabular}
\caption{Sample Data Table}
\label{tab:sample}
\end{table}

\section{Lists}
\subsection{Bullet Points}
\begin{itemize}
  \item First item
  \item Second item
  \item Third item with sub-items:
    \begin{itemize}
      \item Sub-item 1
      \item Sub-item 2
    \end{itemize}
\end{itemize}

\subsection{Numbered List}
\begin{enumerate}
  \item First step
  \item Second step
  \item Third step
\end{enumerate}

\section{Citations and References}
When you use the citation search feature, references will be automatically managed. For example, you might cite a paper like this: \cite{example2024}. All citations will be automatically added to your \texttt{cite.bib} file.

\section{Conclusion}
This template demonstrates basic \LaTeX{} features including mathematical formulas, tables, and lists. The bibliography system is now set up to work seamlessly with Bibby AI's citation management. Feel free to modify it for your needs!

\bibliographystyle{plain}
\bibliography{cite}

\end{document}\documentclass{article}
\usepackage[utf8]{inputenc}
\usepackage{booktabs}
\usepackage{amsmath}
\usepackage{natbib}

\title{My Document}
\author{Your Name}
\date{\today}

\begin{document}

\maketitle

\section{Introduction}
This is a sample document demonstrating various \LaTeX{} features including mathematical equations, tables, and formatting. You can cite papers using the citation search feature, and they will be automatically added to your \texttt{cite.bib} file.

\section{Mathematical Expressions}
Here are some example equations:

\subsection{Quadratic Formula}
The quadratic formula is given by:
\[x = \frac{-b \pm \sqrt{b^2 - 4ac}}{2a}\]

\subsection{Matrix Example}
A sample matrix:
\[
\begin{pmatrix}
a_{11} & a_{12} & a_{13} \\
a_{21} & a_{22} & a_{23} \\
a_{31} & a_{32} & a_{33}
\end{pmatrix}
\]

\section{Tables}
Below is a sample table showing data:

\begin{table}[h]
\centering
\begin{tabular}{@{}lrrr@{}}
\toprule
\textbf{Category} & \textbf{Value 1} & \textbf{Value 2} & \textbf{Total} \\
\midrule
Item A & 10.5 & 20.3 & 30.8 \\
Item B & 15.2 & 25.1 & 40.3 \\
Item C & 22.1 & 12.7 & 34.8 \\
\bottomrule
\end{tabular}
\caption{Sample Data Table}
\label{tab:sample}
\end{table}

\section{Lists}
\subsection{Bullet Points}
\begin{itemize}
  \item First item
  \item Second item
  \item Third item with sub-items:
    \begin{itemize}
      \item Sub-item 1
      \item Sub-item 2
    \end{itemize}
\end{itemize}

\subsection{Numbered List}
\begin{enumerate}
  \item First step
  \item Second step
  \item Third step
\end{enumerate}

\section{Citations and References}
When you use the citation search feature, references will be automatically managed. For example, you might cite a paper like this: \cite{example2024}. All citations will be automatically added to your \texttt{cite.bib} file.

\section{Conclusion}
This template demonstrates basic \LaTeX{} features including mathematical formulas, tables, and lists. The bibliography system is now set up to work seamlessly with Bibby AI's citation management. Feel free to modify it for your needs!

\bibliographystyle{plain}
\bibliography{cite}

\end{document}\documentclass{article}
\usepackage[utf8]{inputenc}
\usepackage{booktabs}
\usepackage{amsmath}
\usepackage{natbib}

\title{My Document}
\author{Your Name}
\date{\today}

\begin{document}

\maketitle

\section{Introduction}
This is a sample document demonstrating various \LaTeX{} features including mathematical equations, tables, and formatting. You can cite papers using the citation search feature, and they will be automatically added to your \texttt{cite.bib} file.

\section{Mathematical Expressions}
Here are some example equations:

\subsection{Quadratic Formula}
The quadratic formula is given by:
\[x = \frac{-b \pm \sqrt{b^2 - 4ac}}{2a}\]

\subsection{Matrix Example}
A sample matrix:
\[
\begin{pmatrix}
a_{11} & a_{12} & a_{13} \\
a_{21} & a_{22} & a_{23} \\
a_{31} & a_{32} & a_{33}
\end{pmatrix}
\]

\section{Tables}
Below is a sample table showing data:

\begin{table}[h]
\centering
\begin{tabular}{@{}lrrr@{}}
\toprule
\textbf{Category} & \textbf{Value 1} & \textbf{Value 2} & \textbf{Total} \\
\midrule
Item A & 10.5 & 20.3 & 30.8 \\
Item B & 15.2 & 25.1 & 40.3 \\
Item C & 22.1 & 12.7 & 34.8 \\
\bottomrule
\end{tabular}
\caption{Sample Data Table}
\label{tab:sample}
\end{table}

\section{Lists}
\subsection{Bullet Points}
\begin{itemize}
  \item First item
  \item Second item
  \item Third item with sub-items:
    \begin{itemize}
      \item Sub-item 1
      \item Sub-item 2
    \end{itemize}
\end{itemize}

\subsection{Numbered List}
\begin{enumerate}
  \item First step
  \item Second step
  \item Third step
\end{enumerate}

\section{Citations and References}
When you use the citation search feature, references will be automatically managed. For example, you might cite a paper like this: \cite{example2024}. All citations will be automatically added to your \texttt{cite.bib} file.

\section{Conclusion}
This template demonstrates basic \LaTeX{} features including mathematical formulas, tables, and lists. The bibliography system is now set up to work seamlessly with Bibby AI's citation management. Feel free to modify it for your needs!

\bibliographystyle{plain}
\bibliography{cite}

\end{document}\documentclass{article}
\usepackage[utf8]{inputenc}
\usepackage{booktabs}
\usepackage{amsmath}
\usepackage{natbib}

\title{My Document}
\author{Your Name}
\date{\today}

\begin{document}

\maketitle

\section{Introduction}
This is a sample document demonstrating various \LaTeX{} features including mathematical equations, tables, and formatting. You can cite papers using the citation search feature, and they will be automatically added to your \texttt{cite.bib} file.

\section{Mathematical Expressions}
Here are some example equations:

\subsection{Quadratic Formula}
The quadratic formula is given by:
\[x = \frac{-b \pm \sqrt{b^2 - 4ac}}{2a}\]

\subsection{Matrix Example}
A sample matrix:
\[
\begin{pmatrix}
a_{11} & a_{12} & a_{13} \\
a_{21} & a_{22} & a_{23} \\
a_{31} & a_{32} & a_{33}
\end{pmatrix}
\]

\section{Tables}
Below is a sample table showing data:

\begin{table}[h]
\centering
\begin{tabular}{@{}lrrr@{}}
\toprule
\textbf{Category} & \textbf{Value 1} & \textbf{Value 2} & \textbf{Total} \\
\midrule
Item A & 10.5 & 20.3 & 30.8 \\
Item B & 15.2 & 25.1 & 40.3 \\
Item C & 22.1 & 12.7 & 34.8 \\
\bottomrule
\end{tabular}
\caption{Sample Data Table}
\label{tab:sample}
\end{table}

\section{Lists}
\subsection{Bullet Points}
\begin{itemize}
  \item First item
  \item Second item
  \item Third item with sub-items:
    \begin{itemize}
      \item Sub-item 1
      \item Sub-item 2
    \end{itemize}
\end{itemize}

\subsection{Numbered List}
\begin{enumerate}
  \item First step
  \item Second step
  \item Third step
\end{enumerate}

\section{Citations and References}
When you use the citation search feature, references will be automatically managed. For example, you might cite a paper like this: \cite{example2024}. All citations will be automatically added to your \texttt{cite.bib} file.

\section{Conclusion}
This template demonstrates basic \LaTeX{} features including mathematical formulas, tables, and lists. The bibliography system is now set up to work seamlessly with Bibby AI's citation management. Feel free to modify it for your needs!

\bibliographystyle{plain}
\bibliography{cite}

\end{document}\documentclass{article}
\usepackage[utf8]{inputenc}
\usepackage{booktabs}
\usepackage{amsmath}
\usepackage{natbib}

\title{My Document}
\author{Your Name}
\date{\today}

\begin{document}

\maketitle

\section{Introduction}
This is a sample document demonstrating various \LaTeX{} features including mathematical equations, tables, and formatting. You can cite papers using the citation search feature, and they will be automatically added to your \texttt{cite.bib} file.

\section{Mathematical Expressions}
Here are some example equations:

\subsection{Quadratic Formula}
The quadratic formula is given by:
\[x = \frac{-b \pm \sqrt{b^2 - 4ac}}{2a}\]

\subsection{Matrix Example}
A sample matrix:
\[
\begin{pmatrix}
a_{11} & a_{12} & a_{13} \\
a_{21} & a_{22} & a_{23} \\
a_{31} & a_{32} & a_{33}
\end{pmatrix}
\]

\section{Tables}
Below is a sample table showing data:

\begin{table}[h]
\centering
\begin{tabular}{@{}lrrr@{}}
\toprule
\textbf{Category} & \textbf{Value 1} & \textbf{Value 2} & \textbf{Total} \\
\midrule
Item A & 10.5 & 20.3 & 30.8 \\
Item B & 15.2 & 25.1 & 40.3 \\
Item C & 22.1 & 12.7 & 34.8 \\
\bottomrule
\end{tabular}
\caption{Sample Data Table}
\label{tab:sample}
\end{table}

\section{Lists}
\subsection{Bullet Points}
\begin{itemize}
  \item First item
  \item Second item
  \item Third item with sub-items:
    \begin{itemize}
      \item Sub-item 1
      \item Sub-item 2
    \end{itemize}
\end{itemize}

\subsection{Numbered List}
\begin{enumerate}
  \item First step
  \item Second step
  \item Third step
\end{enumerate}

\section{Citations and References}
When you use the citation search feature, references will be automatically managed. For example, you might cite a paper like this: \cite{example2024}. All citations will be automatically added to your \texttt{cite.bib} file.

\section{Conclusion}
This template demonstrates basic \LaTeX{} features including mathematical formulas, tables, and lists. The bibliography system is now set up to work seamlessly with Bibby AI's citation management. Feel free to modify it for your needs!

\bibliographystyle{plain}
\bibliography{cite}

\end{document}\documentclass{article}
\usepackage[utf8]{inputenc}
\usepackage{booktabs}
\usepackage{amsmath}
\usepackage{natbib}

\title{My Document}
\author{Your Name}
\date{\today}

\begin{document}

\maketitle

\section{Introduction}
This is a sample document demonstrating various \LaTeX{} features including mathematical equations, tables, and formatting. You can cite papers using the citation search feature, and they will be automatically added to your \texttt{cite.bib} file.

\section{Mathematical Expressions}
Here are some example equations:

\subsection{Quadratic Formula}
The quadratic formula is given by:
\[x = \frac{-b \pm \sqrt{b^2 - 4ac}}{2a}\]

\subsection{Matrix Example}
A sample matrix:
\[
\begin{pmatrix}
a_{11} & a_{12} & a_{13} \\
a_{21} & a_{22} & a_{23} \\
a_{31} & a_{32} & a_{33}
\end{pmatrix}
\]

\section{Tables}
Below is a sample table showing data:

\begin{table}[h]
\centering
\begin{tabular}{@{}lrrr@{}}
\toprule
\textbf{Category} & \textbf{Value 1} & \textbf{Value 2} & \textbf{Total} \\
\midrule
Item A & 10.5 & 20.3 & 30.8 \\
Item B & 15.2 & 25.1 & 40.3 \\
Item C & 22.1 & 12.7 & 34.8 \\
\bottomrule
\end{tabular}
\caption{Sample Data Table}
\label{tab:sample}
\end{table}

\section{Lists}
\subsection{Bullet Points}
\begin{itemize}
  \item First item
  \item Second item
  \item Third item with sub-items:
    \begin{itemize}
      \item Sub-item 1
      \item Sub-item 2
    \end{itemize}
\end{itemize}

\subsection{Numbered List}
\begin{enumerate}
  \item First step
  \item Second step
  \item Third step
\end{enumerate}

\section{Citations and References}
When you use the citation search feature, references will be automatically managed. For example, you might cite a paper like this: \cite{example2024}. All citations will be automatically added to your \texttt{cite.bib} file.

\section{Conclusion}
This template demonstrates basic \LaTeX{} features including mathematical formulas, tables, and lists. The bibliography system is now set up to work seamlessly with Bibby AI's citation management. Feel free to modify it for your needs!

\bibliographystyle{plain}
\bibliography{cite}

\end{document}\documentclass{article}
\usepackage[utf8]{inputenc}
\usepackage{booktabs}
\usepackage{amsmath}
\usepackage{natbib}

\title{My Document}
\author{Your Name}
\date{\today}

\begin{document}

\maketitle

\section{Introduction}
This is a sample document demonstrating various \LaTeX{} features including mathematical equations, tables, and formatting. You can cite papers using the citation search feature, and they will be automatically added to your \texttt{cite.bib} file.

\section{Mathematical Expressions}
Here are some example equations:

\subsection{Quadratic Formula}
The quadratic formula is given by:
\[x = \frac{-b \pm \sqrt{b^2 - 4ac}}{2a}\]

\subsection{Matrix Example}
A sample matrix:
\[
\begin{pmatrix}
a_{11} & a_{12} & a_{13} \\
a_{21} & a_{22} & a_{23} \\
a_{31} & a_{32} & a_{33}
\end{pmatrix}
\]

\section{Tables}
Below is a sample table showing data:

\begin{table}[h]
\centering
\begin{tabular}{@{}lrrr@{}}
\toprule
\textbf{Category} & \textbf{Value 1} & \textbf{Value 2} & \textbf{Total} \\
\midrule
Item A & 10.5 & 20.3 & 30.8 \\
Item B & 15.2 & 25.1 & 40.3 \\
Item C & 22.1 & 12.7 & 34.8 \\
\bottomrule
\end{tabular}
\caption{Sample Data Table}
\label{tab:sample}
\end{table}

\section{Lists}
\subsection{Bullet Points}
\begin{itemize}
  \item First item
  \item Second item
  \item Third item with sub-items:
    \begin{itemize}
      \item Sub-item 1
      \item Sub-item 2
    \end{itemize}
\end{itemize}

\subsection{Numbered List}
\begin{enumerate}
  \item First step
  \item Second step
  \item Third step
\end{enumerate}

\section{Citations and References}
When you use the citation search feature, references will be automatically managed. For example, you might cite a paper like this: \cite{example2024}. All citations will be automatically added to your \texttt{cite.bib} file.

\section{Conclusion}
This template demonstrates basic \LaTeX{} features including mathematical formulas, tables, and lists. The bibliography system is now set up to work seamlessly with Bibby AI's citation management. Feel free to modify it for your needs!

\bibliographystyle{plain}
\bibliography{cite}

\end{document}\documentclass{article}
\usepackage[utf8]{inputenc}
\usepackage{booktabs}
\usepackage{amsmath}
\usepackage{natbib}

\title{My Document}
\author{Your Name}
\date{\today}

\begin{document}

\maketitle

\section{Introduction}
This is a sample document demonstrating various \LaTeX{} features including mathematical equations, tables, and formatting. You can cite papers using the citation search feature, and they will be automatically added to your \texttt{cite.bib} file.

\section{Mathematical Expressions}
Here are some example equations:

\subsection{Quadratic Formula}
The quadratic formula is given by:
\[x = \frac{-b \pm \sqrt{b^2 - 4ac}}{2a}\]

\subsection{Matrix Example}
A sample matrix:
\[
\begin{pmatrix}
a_{11} & a_{12} & a_{13} \\
a_{21} & a_{22} & a_{23} \\
a_{31} & a_{32} & a_{33}
\end{pmatrix}
\]

\section{Tables}
Below is a sample table showing data:

\begin{table}[h]
\centering
\begin{tabular}{@{}lrrr@{}}
\toprule
\textbf{Category} & \textbf{Value 1} & \textbf{Value 2} & \textbf{Total} \\
\midrule
Item A & 10.5 & 20.3 & 30.8 \\
Item B & 15.2 & 25.1 & 40.3 \\
Item C & 22.1 & 12.7 & 34.8 \\
\bottomrule
\end{tabular}
\caption{Sample Data Table}
\label{tab:sample}
\end{table}

\section{Lists}
\subsection{Bullet Points}
\begin{itemize}
  \item First item
  \item Second item
  \item Third item with sub-items:
    \begin{itemize}
      \item Sub-item 1
      \item Sub-item 2
    \end{itemize}
\end{itemize}

\subsection{Numbered List}
\begin{enumerate}
  \item First step
  \item Second step
  \item Third step
\end{enumerate}

\section{Citations and References}
When you use the citation search feature, references will be automatically managed. For example, you might cite a paper like this: \cite{example2024}. All citations will be automatically added to your \texttt{cite.bib} file.

\section{Conclusion}
This template demonstrates basic \LaTeX{} features including mathematical formulas, tables, and lists. The bibliography system is now set up to work seamlessly with Bibby AI's citation management. Feel free to modify it for your needs!

\bibliographystyle{plain}
\bibliography{cite}

\end{document}\documentclass{article}
\usepackage[utf8]{inputenc}
\usepackage{booktabs}
\usepackage{amsmath}
\usepackage{natbib}

\title{My Document}
\author{Your Name}
\date{\today}

\begin{document}

\maketitle

\section{Introduction}
This is a sample document demonstrating various \LaTeX{} features including mathematical equations, tables, and formatting. You can cite papers using the citation search feature, and they will be automatically added to your \texttt{cite.bib} file.

\section{Mathematical Expressions}
Here are some example equations:

\subsection{Quadratic Formula}
The quadratic formula is given by:
\[x = \frac{-b \pm \sqrt{b^2 - 4ac}}{2a}\]

\subsection{Matrix Example}
A sample matrix:
\[
\begin{pmatrix}
a_{11} & a_{12} & a_{13} \\
a_{21} & a_{22} & a_{23} \\
a_{31} & a_{32} & a_{33}
\end{pmatrix}
\]

\section{Tables}
Below is a sample table showing data:

\begin{table}[h]
\centering
\begin{tabular}{@{}lrrr@{}}
\toprule
\textbf{Category} & \textbf{Value 1} & \textbf{Value 2} & \textbf{Total} \\
\midrule
Item A & 10.5 & 20.3 & 30.8 \\
Item B & 15.2 & 25.1 & 40.3 \\
Item C & 22.1 & 12.7 & 34.8 \\
\bottomrule
\end{tabular}
\caption{Sample Data Table}
\label{tab:sample}
\end{table}

\section{Lists}
\subsection{Bullet Points}
\begin{itemize}
  \item First item
  \item Second item
  \item Third item with sub-items:
    \begin{itemize}
      \item Sub-item 1
      \item Sub-item 2
    \end{itemize}
\end{itemize}

\subsection{Numbered List}
\begin{enumerate}
  \item First step
  \item Second step
  \item Third step
\end{enumerate}

\section{Citations and References}
When you use the citation search feature, references will be automatically managed. For example, you might cite a paper like this: \cite{example2024}. All citations will be automatically added to your \texttt{cite.bib} file.

\section{Conclusion}
This template demonstrates basic \LaTeX{} features including mathematical formulas, tables, and lists. The bibliography system is now set up to work seamlessly with Bibby AI's citation management. Feel free to modify it for your needs!

\bibliographystyle{plain}
\bibliography{cite}

\end{document}\documentclass{article}
\usepackage[utf8]{inputenc}
\usepackage{booktabs}
\usepackage{amsmath}
\usepackage{natbib}

\title{My Document}
\author{Your Name}
\date{\today}

\begin{document}

\maketitle

\section{Introduction}
This is a sample document demonstrating various \LaTeX{} features including mathematical equations, tables, and formatting. You can cite papers using the citation search feature, and they will be automatically added to your \texttt{cite.bib} file.

\section{Mathematical Expressions}
Here are some example equations:

\subsection{Quadratic Formula}
The quadratic formula is given by:
\[x = \frac{-b \pm \sqrt{b^2 - 4ac}}{2a}\]

\subsection{Matrix Example}
A sample matrix:
\[
\begin{pmatrix}
a_{11} & a_{12} & a_{13} \\
a_{21} & a_{22} & a_{23} \\
a_{31} & a_{32} & a_{33}
\end{pmatrix}
\]

\section{Tables}
Below is a sample table showing data:

\begin{table}[h]
\centering
\begin{tabular}{@{}lrrr@{}}
\toprule
\textbf{Category} & \textbf{Value 1} & \textbf{Value 2} & \textbf{Total} \\
\midrule
Item A & 10.5 & 20.3 & 30.8 \\
Item B & 15.2 & 25.1 & 40.3 \\
Item C & 22.1 & 12.7 & 34.8 \\
\bottomrule
\end{tabular}
\caption{Sample Data Table}
\label{tab:sample}
\end{table}

\section{Lists}
\subsection{Bullet Points}
\begin{itemize}
  \item First item
  \item Second item
  \item Third item with sub-items:
    \begin{itemize}
      \item Sub-item 1
      \item Sub-item 2
    \end{itemize}
\end{itemize}

\subsection{Numbered List}
\begin{enumerate}
  \item First step
  \item Second step
  \item Third step
\end{enumerate}

\section{Citations and References}
When you use the citation search feature, references will be automatically managed. For example, you might cite a paper like this: \cite{example2024}. All citations will be automatically added to your \texttt{cite.bib} file.

\section{Conclusion}
This template demonstrates basic \LaTeX{} features including mathematical formulas, tables, and lists. The bibliography system is now set up to work seamlessly with Bibby AI's citation management. Feel free to modify it for your needs!

\bibliographystyle{plain}
\bibliography{cite}

\end{document}\documentclass{article}
\usepackage[utf8]{inputenc}
\usepackage{booktabs}
\usepackage{amsmath}
\usepackage{natbib}

\title{My Document}
\author{Your Name}
\date{\today}

\begin{document}

\maketitle

\section{Introduction}
This is a sample document demonstrating various \LaTeX{} features including mathematical equations, tables, and formatting. You can cite papers using the citation search feature, and they will be automatically added to your \texttt{cite.bib} file.

\section{Mathematical Expressions}
Here are some example equations:

\subsection{Quadratic Formula}
The quadratic formula is given by:
\[x = \frac{-b \pm \sqrt{b^2 - 4ac}}{2a}\]

\subsection{Matrix Example}
A sample matrix:
\[
\begin{pmatrix}
a_{11} & a_{12} & a_{13} \\
a_{21} & a_{22} & a_{23} \\
a_{31} & a_{32} & a_{33}
\end{pmatrix}
\]

\section{Tables}
Below is a sample table showing data:

\begin{table}[h]
\centering
\begin{tabular}{@{}lrrr@{}}
\toprule
\textbf{Category} & \textbf{Value 1} & \textbf{Value 2} & \textbf{Total} \\
\midrule
Item A & 10.5 & 20.3 & 30.8 \\
Item B & 15.2 & 25.1 & 40.3 \\
Item C & 22.1 & 12.7 & 34.8 \\
\bottomrule
\end{tabular}
\caption{Sample Data Table}
\label{tab:sample}
\end{table}

\section{Lists}
\subsection{Bullet Points}
\begin{itemize}
  \item First item
  \item Second item
  \item Third item with sub-items:
    \begin{itemize}
      \item Sub-item 1
      \item Sub-item 2
    \end{itemize}
\end{itemize}

\subsection{Numbered List}
\begin{enumerate}
  \item First step
  \item Second step
  \item Third step
\end{enumerate}

\section{Citations and References}
When you use the citation search feature, references will be automatically managed. For example, you might cite a paper like this: \cite{example2024}. All citations will be automatically added to your \texttt{cite.bib} file.

\section{Conclusion}
This template demonstrates basic \LaTeX{} features including mathematical formulas, tables, and lists. The bibliography system is now set up to work seamlessly with Bibby AI's citation management. Feel free to modify it for your needs!

\bibliographystyle{plain}
\bibliography{cite}

\end{document}\documentclass{article}
\usepackage[utf8]{inputenc}
\usepackage{booktabs}
\usepackage{amsmath}
\usepackage{natbib}

\title{My Document}
\author{Your Name}
\date{\today}

\begin{document}

\maketitle

\section{Introduction}
This is a sample document demonstrating various \LaTeX{} features including mathematical equations, tables, and formatting. You can cite papers using the citation search feature, and they will be automatically added to your \texttt{cite.bib} file.

\section{Mathematical Expressions}
Here are some example equations:

\subsection{Quadratic Formula}
The quadratic formula is given by:
\[x = \frac{-b \pm \sqrt{b^2 - 4ac}}{2a}\]

\subsection{Matrix Example}
A sample matrix:
\[
\begin{pmatrix}
a_{11} & a_{12} & a_{13} \\
a_{21} & a_{22} & a_{23} \\
a_{31} & a_{32} & a_{33}
\end{pmatrix}
\]

\section{Tables}
Below is a sample table showing data:

\begin{table}[h]
\centering
\begin{tabular}{@{}lrrr@{}}
\toprule
\textbf{Category} & \textbf{Value 1} & \textbf{Value 2} & \textbf{Total} \\
\midrule
Item A & 10.5 & 20.3 & 30.8 \\
Item B & 15.2 & 25.1 & 40.3 \\
Item C & 22.1 & 12.7 & 34.8 \\
\bottomrule
\end{tabular}
\caption{Sample Data Table}
\label{tab:sample}
\end{table}

\section{Lists}
\subsection{Bullet Points}
\begin{itemize}
  \item First item
  \item Second item
  \item Third item with sub-items:
    \begin{itemize}
      \item Sub-item 1
      \item Sub-item 2
    \end{itemize}
\end{itemize}

\subsection{Numbered List}
\begin{enumerate}
  \item First step
  \item Second step
  \item Third step
\end{enumerate}

\section{Citations and References}
When you use the citation search feature, references will be automatically managed. For example, you might cite a paper like this: \cite{example2024}. All citations will be automatically added to your \texttt{cite.bib} file.

\section{Conclusion}
This template demonstrates basic \LaTeX{} features including mathematical formulas, tables, and lists. The bibliography system is now set up to work seamlessly with Bibby AI's citation management. Feel free to modify it for your needs!

\bibliographystyle{plain}
\bibliography{cite}

\end{document}\documentclass{article}
\usepackage[utf8]{inputenc}
\usepackage{booktabs}
\usepackage{amsmath}
\usepackage{natbib}

\title{My Document}
\author{Your Name}
\date{\today}

\begin{document}

\maketitle

\section{Introduction}
This is a sample document demonstrating various \LaTeX{} features including mathematical equations, tables, and formatting. You can cite papers using the citation search feature, and they will be automatically added to your \texttt{cite.bib} file.

\section{Mathematical Expressions}
Here are some example equations:

\subsection{Quadratic Formula}
The quadratic formula is given by:
\[x = \frac{-b \pm \sqrt{b^2 - 4ac}}{2a}\]

\subsection{Matrix Example}
A sample matrix:
\[
\begin{pmatrix}
a_{11} & a_{12} & a_{13} \\
a_{21} & a_{22} & a_{23} \\
a_{31} & a_{32} & a_{33}
\end{pmatrix}
\]

\section{Tables}
Below is a sample table showing data:

\begin{table}[h]
\centering
\begin{tabular}{@{}lrrr@{}}
\toprule
\textbf{Category} & \textbf{Value 1} & \textbf{Value 2} & \textbf{Total} \\
\midrule
Item A & 10.5 & 20.3 & 30.8 \\
Item B & 15.2 & 25.1 & 40.3 \\
Item C & 22.1 & 12.7 & 34.8 \\
\bottomrule
\end{tabular}
\caption{Sample Data Table}
\label{tab:sample}
\end{table}

\section{Lists}
\subsection{Bullet Points}
\begin{itemize}
  \item First item
  \item Second item
  \item Third item with sub-items:
    \begin{itemize}
      \item Sub-item 1
      \item Sub-item 2
    \end{itemize}
\end{itemize}

\subsection{Numbered List}
\begin{enumerate}
  \item First step
  \item Second step
  \item Third step
\end{enumerate}

\section{Citations and References}
When you use the citation search feature, references will be automatically managed. For example, you might cite a paper like this: \cite{example2024}. All citations will be automatically added to your \texttt{cite.bib} file.

\section{Conclusion}
This template demonstrates basic \LaTeX{} features including mathematical formulas, tables, and lists. The bibliography system is now set up to work seamlessly with Bibby AI's citation management. Feel free to modify it for your needs!

\bibliographystyle{plain}
\bibliography{cite}

\end{document}\documentclass{article}
\usepackage[utf8]{inputenc}
\usepackage{booktabs}
\usepackage{amsmath}
\usepackage{natbib}

\title{My Document}
\author{Your Name}
\date{\today}

\begin{document}

\maketitle

\section{Introduction}
This is a sample document demonstrating various \LaTeX{} features including mathematical equations, tables, and formatting. You can cite papers using the citation search feature, and they will be automatically added to your \texttt{cite.bib} file.

\section{Mathematical Expressions}
Here are some example equations:

\subsection{Quadratic Formula}
The quadratic formula is given by:
\[x = \frac{-b \pm \sqrt{b^2 - 4ac}}{2a}\]

\subsection{Matrix Example}
A sample matrix:
\[
\begin{pmatrix}
a_{11} & a_{12} & a_{13} \\
a_{21} & a_{22} & a_{23} \\
a_{31} & a_{32} & a_{33}
\end{pmatrix}
\]

\section{Tables}
Below is a sample table showing data:

\begin{table}[h]
\centering
\begin{tabular}{@{}lrrr@{}}
\toprule
\textbf{Category} & \textbf{Value 1} & \textbf{Value 2} & \textbf{Total} \\
\midrule
Item A & 10.5 & 20.3 & 30.8 \\
Item B & 15.2 & 25.1 & 40.3 \\
Item C & 22.1 & 12.7 & 34.8 \\
\bottomrule
\end{tabular}
\caption{Sample Data Table}
\label{tab:sample}
\end{table}

\section{Lists}
\subsection{Bullet Points}
\begin{itemize}
  \item First item
  \item Second item
  \item Third item with sub-items:
    \begin{itemize}
      \item Sub-item 1
      \item Sub-item 2
    \end{itemize}
\end{itemize}

\subsection{Numbered List}
\begin{enumerate}
  \item First step
  \item Second step
  \item Third step
\end{enumerate}

\section{Citations and References}
When you use the citation search feature, references will be automatically managed. For example, you might cite a paper like this: \cite{example2024}. All citations will be automatically added to your \texttt{cite.bib} file.

\section{Conclusion}
This template demonstrates basic \LaTeX{} features including mathematical formulas, tables, and lists. The bibliography system is now set up to work seamlessly with Bibby AI's citation management. Feel free to modify it for your needs!

\bibliographystyle{plain}
\bibliography{cite}

\end{document}\documentclass{article}
\usepackage[utf8]{inputenc}
\usepackage{booktabs}
\usepackage{amsmath}
\usepackage{natbib}

\title{My Document}
\author{Your Name}
\date{\today}

\begin{document}

\maketitle

\section{Introduction}
This is a sample document demonstrating various \LaTeX{} features including mathematical equations, tables, and formatting. You can cite papers using the citation search feature, and they will be automatically added to your \texttt{cite.bib} file.

\section{Mathematical Expressions}
Here are some example equations:

\subsection{Quadratic Formula}
The quadratic formula is given by:
\[x = \frac{-b \pm \sqrt{b^2 - 4ac}}{2a}\]

\subsection{Matrix Example}
A sample matrix:
\[
\begin{pmatrix}
a_{11} & a_{12} & a_{13} \\
a_{21} & a_{22} & a_{23} \\
a_{31} & a_{32} & a_{33}
\end{pmatrix}
\]

\section{Tables}
Below is a sample table showing data:

\begin{table}[h]
\centering
\begin{tabular}{@{}lrrr@{}}
\toprule
\textbf{Category} & \textbf{Value 1} & \textbf{Value 2} & \textbf{Total} \\
\midrule
Item A & 10.5 & 20.3 & 30.8 \\
Item B & 15.2 & 25.1 & 40.3 \\
Item C & 22.1 & 12.7 & 34.8 \\
\bottomrule
\end{tabular}
\caption{Sample Data Table}
\label{tab:sample}
\end{table}

\section{Lists}
\subsection{Bullet Points}
\begin{itemize}
  \item First item
  \item Second item
  \item Third item with sub-items:
    \begin{itemize}
      \item Sub-item 1
      \item Sub-item 2
    \end{itemize}
\end{itemize}

\subsection{Numbered List}
\begin{enumerate}
  \item First step
  \item Second step
  \item Third step
\end{enumerate}

\section{Citations and References}
When you use the citation search feature, references will be automatically managed. For example, you might cite a paper like this: \cite{example2024}. All citations will be automatically added to your \texttt{cite.bib} file.

\section{Conclusion}
This template demonstrates basic \LaTeX{} features including mathematical formulas, tables, and lists. The bibliography system is now set up to work seamlessly with Bibby AI's citation management. Feel free to modify it for your needs!

\bibliographystyle{plain}
\bibliography{cite}

\end{document}\documentclass{article}
\usepackage[utf8]{inputenc}
\usepackage{booktabs}
\usepackage{amsmath}
\usepackage{natbib}

\title{My Document}
\author{Your Name}
\date{\today}

\begin{document}

\maketitle

\section{Introduction}
This is a sample document demonstrating various \LaTeX{} features including mathematical equations, tables, and formatting. You can cite papers using the citation search feature, and they will be automatically added to your \texttt{cite.bib} file.

\section{Mathematical Expressions}
Here are some example equations:

\subsection{Quadratic Formula}
The quadratic formula is given by:
\[x = \frac{-b \pm \sqrt{b^2 - 4ac}}{2a}\]

\subsection{Matrix Example}
A sample matrix:
\[
\begin{pmatrix}
a_{11} & a_{12} & a_{13} \\
a_{21} & a_{22} & a_{23} \\
a_{31} & a_{32} & a_{33}
\end{pmatrix}
\]

\section{Tables}
Below is a sample table showing data:

\begin{table}[h]
\centering
\begin{tabular}{@{}lrrr@{}}
\toprule
\textbf{Category} & \textbf{Value 1} & \textbf{Value 2} & \textbf{Total} \\
\midrule
Item A & 10.5 & 20.3 & 30.8 \\
Item B & 15.2 & 25.1 & 40.3 \\
Item C & 22.1 & 12.7 & 34.8 \\
\bottomrule
\end{tabular}
\caption{Sample Data Table}
\label{tab:sample}
\end{table}

\section{Lists}
\subsection{Bullet Points}
\begin{itemize}
  \item First item
  \item Second item
  \item Third item with sub-items:
    \begin{itemize}
      \item Sub-item 1
      \item Sub-item 2
    \end{itemize}
\end{itemize}

\subsection{Numbered List}
\begin{enumerate}
  \item First step
  \item Second step
  \item Third step
\end{enumerate}

\section{Citations and References}
When you use the citation search feature, references will be automatically managed. For example, you might cite a paper like this: \cite{example2024}. All citations will be automatically added to your \texttt{cite.bib} file.

\section{Conclusion}
This template demonstrates basic \LaTeX{} features including mathematical formulas, tables, and lists. The bibliography system is now set up to work seamlessly with Bibby AI's citation management. Feel free to modify it for your needs!

\bibliographystyle{plain}
\bibliography{cite}

\end{document}\documentclass{article}
\usepackage[utf8]{inputenc}
\usepackage{booktabs}
\usepackage{amsmath}
\usepackage{natbib}

\title{My Document}
\author{Your Name}
\date{\today}

\begin{document}

\maketitle

\section{Introduction}
This is a sample document demonstrating various \LaTeX{} features including mathematical equations, tables, and formatting. You can cite papers using the citation search feature, and they will be automatically added to your \texttt{cite.bib} file.

\section{Mathematical Expressions}
Here are some example equations:

\subsection{Quadratic Formula}
The quadratic formula is given by:
\[x = \frac{-b \pm \sqrt{b^2 - 4ac}}{2a}\]

\subsection{Matrix Example}
A sample matrix:
\[
\begin{pmatrix}
a_{11} & a_{12} & a_{13} \\
a_{21} & a_{22} & a_{23} \\
a_{31} & a_{32} & a_{33}
\end{pmatrix}
\]

\section{Tables}
Below is a sample table showing data:

\begin{table}[h]
\centering
\begin{tabular}{@{}lrrr@{}}
\toprule
\textbf{Category} & \textbf{Value 1} & \textbf{Value 2} & \textbf{Total} \\
\midrule
Item A & 10.5 & 20.3 & 30.8 \\
Item B & 15.2 & 25.1 & 40.3 \\
Item C & 22.1 & 12.7 & 34.8 \\
\bottomrule
\end{tabular}
\caption{Sample Data Table}
\label{tab:sample}
\end{table}

\section{Lists}
\subsection{Bullet Points}
\begin{itemize}
  \item First item
  \item Second item
  \item Third item with sub-items:
    \begin{itemize}
      \item Sub-item 1
      \item Sub-item 2
    \end{itemize}
\end{itemize}

\subsection{Numbered List}
\begin{enumerate}
  \item First step
  \item Second step
  \item Third step
\end{enumerate}

\section{Citations and References}
When you use the citation search feature, references will be automatically managed. For example, you might cite a paper like this: \cite{example2024}. All citations will be automatically added to your \texttt{cite.bib} file.

\section{Conclusion}
This template demonstrates basic \LaTeX{} features including mathematical formulas, tables, and lists. The bibliography system is now set up to work seamlessly with Bibby AI's citation management. Feel free to modify it for your needs!

\bibliographystyle{plain}
\bibliography{cite}

\end{document}\documentclass{article}
\usepackage[utf8]{inputenc}
\usepackage{booktabs}
\usepackage{amsmath}
\usepackage{natbib}

\title{My Document}
\author{Your Name}
\date{\today}

\begin{document}

\maketitle

\section{Introduction}
This is a sample document demonstrating various \LaTeX{} features including mathematical equations, tables, and formatting. You can cite papers using the citation search feature, and they will be automatically added to your \texttt{cite.bib} file.

\section{Mathematical Expressions}
Here are some example equations:

\subsection{Quadratic Formula}
The quadratic formula is given by:
\[x = \frac{-b \pm \sqrt{b^2 - 4ac}}{2a}\]

\subsection{Matrix Example}
A sample matrix:
\[
\begin{pmatrix}
a_{11} & a_{12} & a_{13} \\
a_{21} & a_{22} & a_{23} \\
a_{31} & a_{32} & a_{33}
\end{pmatrix}
\]

\section{Tables}
Below is a sample table showing data:

\begin{table}[h]
\centering
\begin{tabular}{@{}lrrr@{}}
\toprule
\textbf{Category} & \textbf{Value 1} & \textbf{Value 2} & \textbf{Total} \\
\midrule
Item A & 10.5 & 20.3 & 30.8 \\
Item B & 15.2 & 25.1 & 40.3 \\
Item C & 22.1 & 12.7 & 34.8 \\
\bottomrule
\end{tabular}
\caption{Sample Data Table}
\label{tab:sample}
\end{table}

\section{Lists}
\subsection{Bullet Points}
\begin{itemize}
  \item First item
  \item Second item
  \item Third item with sub-items:
    \begin{itemize}
      \item Sub-item 1
      \item Sub-item 2
    \end{itemize}
\end{itemize}

\subsection{Numbered List}
\begin{enumerate}
  \item First step
  \item Second step
  \item Third step
\end{enumerate}

\section{Citations and References}
When you use the citation search feature, references will be automatically managed. For example, you might cite a paper like this: \cite{example2024}. All citations will be automatically added to your \texttt{cite.bib} file.

\section{Conclusion}
This template demonstrates basic \LaTeX{} features including mathematical formulas, tables, and lists. The bibliography system is now set up to work seamlessly with Bibby AI's citation management. Feel free to modify it for your needs!

\bibliographystyle{plain}
\bibliography{cite}

\end{document}\documentclass{article}
\usepackage[utf8]{inputenc}
\usepackage{booktabs}
\usepackage{amsmath}
\usepackage{natbib}

\title{My Document}
\author{Your Name}
\date{\today}

\begin{document}

\maketitle

\section{Introduction}
This is a sample document demonstrating various \LaTeX{} features including mathematical equations, tables, and formatting. You can cite papers using the citation search feature, and they will be automatically added to your \texttt{cite.bib} file.

\section{Mathematical Expressions}
Here are some example equations:

\subsection{Quadratic Formula}
The quadratic formula is given by:
\[x = \frac{-b \pm \sqrt{b^2 - 4ac}}{2a}\]

\subsection{Matrix Example}
A sample matrix:
\[
\begin{pmatrix}
a_{11} & a_{12} & a_{13} \\
a_{21} & a_{22} & a_{23} \\
a_{31} & a_{32} & a_{33}
\end{pmatrix}
\]

\section{Tables}
Below is a sample table showing data:

\begin{table}[h]
\centering
\begin{tabular}{@{}lrrr@{}}
\toprule
\textbf{Category} & \textbf{Value 1} & \textbf{Value 2} & \textbf{Total} \\
\midrule
Item A & 10.5 & 20.3 & 30.8 \\
Item B & 15.2 & 25.1 & 40.3 \\
Item C & 22.1 & 12.7 & 34.8 \\
\bottomrule
\end{tabular}
\caption{Sample Data Table}
\label{tab:sample}
\end{table}

\section{Lists}
\subsection{Bullet Points}
\begin{itemize}
  \item First item
  \item Second item
  \item Third item with sub-items:
    \begin{itemize}
      \item Sub-item 1
      \item Sub-item 2
    \end{itemize}
\end{itemize}

\subsection{Numbered List}
\begin{enumerate}
  \item First step
  \item Second step
  \item Third step
\end{enumerate}

\section{Citations and References}
When you use the citation search feature, references will be automatically managed. For example, you might cite a paper like this: \cite{example2024}. All citations will be automatically added to your \texttt{cite.bib} file.

\section{Conclusion}
This template demonstrates basic \LaTeX{} features including mathematical formulas, tables, and lists. The bibliography system is now set up to work seamlessly with Bibby AI's citation management. Feel free to modify it for your needs!

\bibliographystyle{plain}
\bibliography{cite}

\end{document}\documentclass{article}
\usepackage[utf8]{inputenc}
\usepackage{booktabs}
\usepackage{amsmath}
\usepackage{natbib}

\title{My Document}
\author{Your Name}
\date{\today}

\begin{document}

\maketitle

\section{Introduction}
This is a sample document demonstrating various \LaTeX{} features including mathematical equations, tables, and formatting. You can cite papers using the citation search feature, and they will be automatically added to your \texttt{cite.bib} file.

\section{Mathematical Expressions}
Here are some example equations:

\subsection{Quadratic Formula}
The quadratic formula is given by:
\[x = \frac{-b \pm \sqrt{b^2 - 4ac}}{2a}\]

\subsection{Matrix Example}
A sample matrix:
\[
\begin{pmatrix}
a_{11} & a_{12} & a_{13} \\
a_{21} & a_{22} & a_{23} \\
a_{31} & a_{32} & a_{33}
\end{pmatrix}
\]

\section{Tables}
Below is a sample table showing data:

\begin{table}[h]
\centering
\begin{tabular}{@{}lrrr@{}}
\toprule
\textbf{Category} & \textbf{Value 1} & \textbf{Value 2} & \textbf{Total} \\
\midrule
Item A & 10.5 & 20.3 & 30.8 \\
Item B & 15.2 & 25.1 & 40.3 \\
Item C & 22.1 & 12.7 & 34.8 \\
\bottomrule
\end{tabular}
\caption{Sample Data Table}
\label{tab:sample}
\end{table}

\section{Lists}
\subsection{Bullet Points}
\begin{itemize}
  \item First item
  \item Second item
  \item Third item with sub-items:
    \begin{itemize}
      \item Sub-item 1
      \item Sub-item 2
    \end{itemize}
\end{itemize}

\subsection{Numbered List}
\begin{enumerate}
  \item First step
  \item Second step
  \item Third step
\end{enumerate}

\section{Citations and References}
When you use the citation search feature, references will be automatically managed. For example, you might cite a paper like this: \cite{example2024}. All citations will be automatically added to your \texttt{cite.bib} file.

\section{Conclusion}
This template demonstrates basic \LaTeX{} features including mathematical formulas, tables, and lists. The bibliography system is now set up to work seamlessly with Bibby AI's citation management. Feel free to modify it for your needs!

\bibliographystyle{plain}
\bibliography{cite}

\end{document}\documentclass{article}
\usepackage[utf8]{inputenc}
\usepackage{booktabs}
\usepackage{amsmath}
\usepackage{natbib}

\title{My Document}
\author{Your Name}
\date{\today}

\begin{document}

\maketitle

\section{Introduction}
This is a sample document demonstrating various \LaTeX{} features including mathematical equations, tables, and formatting. You can cite papers using the citation search feature, and they will be automatically added to your \texttt{cite.bib} file.

\section{Mathematical Expressions}
Here are some example equations:

\subsection{Quadratic Formula}
The quadratic formula is given by:
\[x = \frac{-b \pm \sqrt{b^2 - 4ac}}{2a}\]

\subsection{Matrix Example}
A sample matrix:
\[
\begin{pmatrix}
a_{11} & a_{12} & a_{13} \\
a_{21} & a_{22} & a_{23} \\
a_{31} & a_{32} & a_{33}
\end{pmatrix}
\]

\section{Tables}
Below is a sample table showing data:

\begin{table}[h]
\centering
\begin{tabular}{@{}lrrr@{}}
\toprule
\textbf{Category} & \textbf{Value 1} & \textbf{Value 2} & \textbf{Total} \\
\midrule
Item A & 10.5 & 20.3 & 30.8 \\
Item B & 15.2 & 25.1 & 40.3 \\
Item C & 22.1 & 12.7 & 34.8 \\
\bottomrule
\end{tabular}
\caption{Sample Data Table}
\label{tab:sample}
\end{table}

\section{Lists}
\subsection{Bullet Points}
\begin{itemize}
  \item First item
  \item Second item
  \item Third item with sub-items:
    \begin{itemize}
      \item Sub-item 1
      \item Sub-item 2
    \end{itemize}
\end{itemize}

\subsection{Numbered List}
\begin{enumerate}
  \item First step
  \item Second step
  \item Third step
\end{enumerate}

\section{Citations and References}
When you use the citation search feature, references will be automatically managed. For example, you might cite a paper like this: \cite{example2024}. All citations will be automatically added to your \texttt{cite.bib} file.

\section{Conclusion}
This template demonstrates basic \LaTeX{} features including mathematical formulas, tables, and lists. The bibliography system is now set up to work seamlessly with Bibby AI's citation management. Feel free to modify it for your needs!

\bibliographystyle{plain}
\bibliography{cite}

\end{document}\documentclass{article}
\usepackage[utf8]{inputenc}
\usepackage{booktabs}
\usepackage{amsmath}
\usepackage{natbib}

\title{My Document}
\author{Your Name}
\date{\today}

\begin{document}

\maketitle

\section{Introduction}
This is a sample document demonstrating various \LaTeX{} features including mathematical equations, tables, and formatting. You can cite papers using the citation search feature, and they will be automatically added to your \texttt{cite.bib} file.

\section{Mathematical Expressions}
Here are some example equations:

\subsection{Quadratic Formula}
The quadratic formula is given by:
\[x = \frac{-b \pm \sqrt{b^2 - 4ac}}{2a}\]

\subsection{Matrix Example}
A sample matrix:
\[
\begin{pmatrix}
a_{11} & a_{12} & a_{13} \\
a_{21} & a_{22} & a_{23} \\
a_{31} & a_{32} & a_{33}
\end{pmatrix}
\]

\section{Tables}
Below is a sample table showing data:

\begin{table}[h]
\centering
\begin{tabular}{@{}lrrr@{}}
\toprule
\textbf{Category} & \textbf{Value 1} & \textbf{Value 2} & \textbf{Total} \\
\midrule
Item A & 10.5 & 20.3 & 30.8 \\
Item B & 15.2 & 25.1 & 40.3 \\
Item C & 22.1 & 12.7 & 34.8 \\
\bottomrule
\end{tabular}
\caption{Sample Data Table}
\label{tab:sample}
\end{table}

\section{Lists}
\subsection{Bullet Points}
\begin{itemize}
  \item First item
  \item Second item
  \item Third item with sub-items:
    \begin{itemize}
      \item Sub-item 1
      \item Sub-item 2
    \end{itemize}
\end{itemize}

\subsection{Numbered List}
\begin{enumerate}
  \item First step
  \item Second step
  \item Third step
\end{enumerate}

\section{Citations and References}
When you use the citation search feature, references will be automatically managed. For example, you might cite a paper like this: \cite{example2024}. All citations will be automatically added to your \texttt{cite.bib} file.

\section{Conclusion}
This template demonstrates basic \LaTeX{} features including mathematical formulas, tables, and lists. The bibliography system is now set up to work seamlessly with Bibby AI's citation management. Feel free to modify it for your needs!

\bibliographystyle{plain}
\bibliography{cite}

\end{document}\documentclass{article}
\usepackage[utf8]{inputenc}
\usepackage{booktabs}
\usepackage{amsmath}
\usepackage{natbib}

\title{My Document}
\author{Your Name}
\date{\today}

\begin{document}

\maketitle

\section{Introduction}
This is a sample document demonstrating various \LaTeX{} features including mathematical equations, tables, and formatting. You can cite papers using the citation search feature, and they will be automatically added to your \texttt{cite.bib} file.

\section{Mathematical Expressions}
Here are some example equations:

\subsection{Quadratic Formula}
The quadratic formula is given by:
\[x = \frac{-b \pm \sqrt{b^2 - 4ac}}{2a}\]

\subsection{Matrix Example}
A sample matrix:
\[
\begin{pmatrix}
a_{11} & a_{12} & a_{13} \\
a_{21} & a_{22} & a_{23} \\
a_{31} & a_{32} & a_{33}
\end{pmatrix}
\]

\section{Tables}
Below is a sample table showing data:

\begin{table}[h]
\centering
\begin{tabular}{@{}lrrr@{}}
\toprule
\textbf{Category} & \textbf{Value 1} & \textbf{Value 2} & \textbf{Total} \\
\midrule
Item A & 10.5 & 20.3 & 30.8 \\
Item B & 15.2 & 25.1 & 40.3 \\
Item C & 22.1 & 12.7 & 34.8 \\
\bottomrule
\end{tabular}
\caption{Sample Data Table}
\label{tab:sample}
\end{table}

\section{Lists}
\subsection{Bullet Points}
\begin{itemize}
  \item First item
  \item Second item
  \item Third item with sub-items:
    \begin{itemize}
      \item Sub-item 1
      \item Sub-item 2
    \end{itemize}
\end{itemize}

\subsection{Numbered List}
\begin{enumerate}
  \item First step
  \item Second step
  \item Third step
\end{enumerate}

\section{Citations and References}
When you use the citation search feature, references will be automatically managed. For example, you might cite a paper like this: \cite{example2024}. All citations will be automatically added to your \texttt{cite.bib} file.

\section{Conclusion}
This template demonstrates basic \LaTeX{} features including mathematical formulas, tables, and lists. The bibliography system is now set up to work seamlessly with Bibby AI's citation management. Feel free to modify it for your needs!

\bibliographystyle{plain}
\bibliography{cite}

\end{document}\documentclass{article}
\usepackage[utf8]{inputenc}
\usepackage{booktabs}
\usepackage{amsmath}
\usepackage{natbib}

\title{My Document}
\author{Your Name}
\date{\today}

\begin{document}

\maketitle

\section{Introduction}
This is a sample document demonstrating various \LaTeX{} features including mathematical equations, tables, and formatting. You can cite papers using the citation search feature, and they will be automatically added to your \texttt{cite.bib} file.

\section{Mathematical Expressions}
Here are some example equations:

\subsection{Quadratic Formula}
The quadratic formula is given by:
\[x = \frac{-b \pm \sqrt{b^2 - 4ac}}{2a}\]

\subsection{Matrix Example}
A sample matrix:
\[
\begin{pmatrix}
a_{11} & a_{12} & a_{13} \\
a_{21} & a_{22} & a_{23} \\
a_{31} & a_{32} & a_{33}
\end{pmatrix}
\]

\section{Tables}
Below is a sample table showing data:

\begin{table}[h]
\centering
\begin{tabular}{@{}lrrr@{}}
\toprule
\textbf{Category} & \textbf{Value 1} & \textbf{Value 2} & \textbf{Total} \\
\midrule
Item A & 10.5 & 20.3 & 30.8 \\
Item B & 15.2 & 25.1 & 40.3 \\
Item C & 22.1 & 12.7 & 34.8 \\
\bottomrule
\end{tabular}
\caption{Sample Data Table}
\label{tab:sample}
\end{table}

\section{Lists}
\subsection{Bullet Points}
\begin{itemize}
  \item First item
  \item Second item
  \item Third item with sub-items:
    \begin{itemize}
      \item Sub-item 1
      \item Sub-item 2
    \end{itemize}
\end{itemize}

\subsection{Numbered List}
\begin{enumerate}
  \item First step
  \item Second step
  \item Third step
\end{enumerate}

\section{Citations and References}
When you use the citation search feature, references will be automatically managed. For example, you might cite a paper like this: \cite{example2024}. All citations will be automatically added to your \texttt{cite.bib} file.

\section{Conclusion}
This template demonstrates basic \LaTeX{} features including mathematical formulas, tables, and lists. The bibliography system is now set up to work seamlessly with Bibby AI's citation management. Feel free to modify it for your needs!

\bibliographystyle{plain}
\bibliography{cite}

\end{document}\documentclass{article}
\usepackage[utf8]{inputenc}
\usepackage{booktabs}
\usepackage{amsmath}
\usepackage{natbib}

\title{My Document}
\author{Your Name}
\date{\today}

\begin{document}

\maketitle

\section{Introduction}
This is a sample document demonstrating various \LaTeX{} features including mathematical equations, tables, and formatting. You can cite papers using the citation search feature, and they will be automatically added to your \texttt{cite.bib} file.

\section{Mathematical Expressions}
Here are some example equations:

\subsection{Quadratic Formula}
The quadratic formula is given by:
\[x = \frac{-b \pm \sqrt{b^2 - 4ac}}{2a}\]

\subsection{Matrix Example}
A sample matrix:
\[
\begin{pmatrix}
a_{11} & a_{12} & a_{13} \\
a_{21} & a_{22} & a_{23} \\
a_{31} & a_{32} & a_{33}
\end{pmatrix}
\]

\section{Tables}
Below is a sample table showing data:

\begin{table}[h]
\centering
\begin{tabular}{@{}lrrr@{}}
\toprule
\textbf{Category} & \textbf{Value 1} & \textbf{Value 2} & \textbf{Total} \\
\midrule
Item A & 10.5 & 20.3 & 30.8 \\
Item B & 15.2 & 25.1 & 40.3 \\
Item C & 22.1 & 12.7 & 34.8 \\
\bottomrule
\end{tabular}
\caption{Sample Data Table}
\label{tab:sample}
\end{table}

\section{Lists}
\subsection{Bullet Points}
\begin{itemize}
  \item First item
  \item Second item
  \item Third item with sub-items:
    \begin{itemize}
      \item Sub-item 1
      \item Sub-item 2
    \end{itemize}
\end{itemize}

\subsection{Numbered List}
\begin{enumerate}
  \item First step
  \item Second step
  \item Third step
\end{enumerate}

\section{Citations and References}
When you use the citation search feature, references will be automatically managed. For example, you might cite a paper like this: \cite{example2024}. All citations will be automatically added to your \texttt{cite.bib} file.

\section{Conclusion}
This template demonstrates basic \LaTeX{} features including mathematical formulas, tables, and lists. The bibliography system is now set up to work seamlessly with Bibby AI's citation management. Feel free to modify it for your needs!

\bibliographystyle{plain}
\bibliography{cite}

\end{document}\documentclass{article}
\usepackage[utf8]{inputenc}
\usepackage{booktabs}
\usepackage{amsmath}
\usepackage{natbib}

\title{My Document}
\author{Your Name}
\date{\today}

\begin{document}

\maketitle

\section{Introduction}
This is a sample document demonstrating various \LaTeX{} features including mathematical equations, tables, and formatting. You can cite papers using the citation search feature, and they will be automatically added to your \texttt{cite.bib} file.

\section{Mathematical Expressions}
Here are some example equations:

\subsection{Quadratic Formula}
The quadratic formula is given by:
\[x = \frac{-b \pm \sqrt{b^2 - 4ac}}{2a}\]

\subsection{Matrix Example}
A sample matrix:
\[
\begin{pmatrix}
a_{11} & a_{12} & a_{13} \\
a_{21} & a_{22} & a_{23} \\
a_{31} & a_{32} & a_{33}
\end{pmatrix}
\]

\section{Tables}
Below is a sample table showing data:

\begin{table}[h]
\centering
\begin{tabular}{@{}lrrr@{}}
\toprule
\textbf{Category} & \textbf{Value 1} & \textbf{Value 2} & \textbf{Total} \\
\midrule
Item A & 10.5 & 20.3 & 30.8 \\
Item B & 15.2 & 25.1 & 40.3 \\
Item C & 22.1 & 12.7 & 34.8 \\
\bottomrule
\end{tabular}
\caption{Sample Data Table}
\label{tab:sample}
\end{table}

\section{Lists}
\subsection{Bullet Points}
\begin{itemize}
  \item First item
  \item Second item
  \item Third item with sub-items:
    \begin{itemize}
      \item Sub-item 1
      \item Sub-item 2
    \end{itemize}
\end{itemize}

\subsection{Numbered List}
\begin{enumerate}
  \item First step
  \item Second step
  \item Third step
\end{enumerate}

\section{Citations and References}
When you use the citation search feature, references will be automatically managed. For example, you might cite a paper like this: \cite{example2024}. All citations will be automatically added to your \texttt{cite.bib} file.

\section{Conclusion}
This template demonstrates basic \LaTeX{} features including mathematical formulas, tables, and lists. The bibliography system is now set up to work seamlessly with Bibby AI's citation management. Feel free to modify it for your needs!

\bibliographystyle{plain}
\bibliography{cite}

\end{document}\documentclass{article}
\usepackage[utf8]{inputenc}
\usepackage{booktabs}
\usepackage{amsmath}
\usepackage{natbib}

\title{My Document}
\author{Your Name}
\date{\today}

\begin{document}

\maketitle

\section{Introduction}
This is a sample document demonstrating various \LaTeX{} features including mathematical equations, tables, and formatting. You can cite papers using the citation search feature, and they will be automatically added to your \texttt{cite.bib} file.

\section{Mathematical Expressions}
Here are some example equations:

\subsection{Quadratic Formula}
The quadratic formula is given by:
\[x = \frac{-b \pm \sqrt{b^2 - 4ac}}{2a}\]

\subsection{Matrix Example}
A sample matrix:
\[
\begin{pmatrix}
a_{11} & a_{12} & a_{13} \\
a_{21} & a_{22} & a_{23} \\
a_{31} & a_{32} & a_{33}
\end{pmatrix}
\]

\section{Tables}
Below is a sample table showing data:

\begin{table}[h]
\centering
\begin{tabular}{@{}lrrr@{}}
\toprule
\textbf{Category} & \textbf{Value 1} & \textbf{Value 2} & \textbf{Total} \\
\midrule
Item A & 10.5 & 20.3 & 30.8 \\
Item B & 15.2 & 25.1 & 40.3 \\
Item C & 22.1 & 12.7 & 34.8 \\
\bottomrule
\end{tabular}
\caption{Sample Data Table}
\label{tab:sample}
\end{table}

\section{Lists}
\subsection{Bullet Points}
\begin{itemize}
  \item First item
  \item Second item
  \item Third item with sub-items:
    \begin{itemize}
      \item Sub-item 1
      \item Sub-item 2
    \end{itemize}
\end{itemize}

\subsection{Numbered List}
\begin{enumerate}
  \item First step
  \item Second step
  \item Third step
\end{enumerate}

\section{Citations and References}
When you use the citation search feature, references will be automatically managed. For example, you might cite a paper like this: \cite{example2024}. All citations will be automatically added to your \texttt{cite.bib} file.

\section{Conclusion}
This template demonstrates basic \LaTeX{} features including mathematical formulas, tables, and lists. The bibliography system is now set up to work seamlessly with Bibby AI's citation management. Feel free to modify it for your needs!

\bibliographystyle{plain}
\bibliography{cite}

\end{document}\documentclass{article}
\usepackage[utf8]{inputenc}
\usepackage{booktabs}
\usepackage{amsmath}
\usepackage{natbib}

\title{My Document}
\author{Your Name}
\date{\today}

\begin{document}

\maketitle

\section{Introduction}
This is a sample document demonstrating various \LaTeX{} features including mathematical equations, tables, and formatting. You can cite papers using the citation search feature, and they will be automatically added to your \texttt{cite.bib} file.

\section{Mathematical Expressions}
Here are some example equations:

\subsection{Quadratic Formula}
The quadratic formula is given by:
\[x = \frac{-b \pm \sqrt{b^2 - 4ac}}{2a}\]

\subsection{Matrix Example}
A sample matrix:
\[
\begin{pmatrix}
a_{11} & a_{12} & a_{13} \\
a_{21} & a_{22} & a_{23} \\
a_{31} & a_{32} & a_{33}
\end{pmatrix}
\]

\section{Tables}
Below is a sample table showing data:

\begin{table}[h]
\centering
\begin{tabular}{@{}lrrr@{}}
\toprule
\textbf{Category} & \textbf{Value 1} & \textbf{Value 2} & \textbf{Total} \\
\midrule
Item A & 10.5 & 20.3 & 30.8 \\
Item B & 15.2 & 25.1 & 40.3 \\
Item C & 22.1 & 12.7 & 34.8 \\
\bottomrule
\end{tabular}
\caption{Sample Data Table}
\label{tab:sample}
\end{table}

\section{Lists}
\subsection{Bullet Points}
\begin{itemize}
  \item First item
  \item Second item
  \item Third item with sub-items:
    \begin{itemize}
      \item Sub-item 1
      \item Sub-item 2
    \end{itemize}
\end{itemize}

\subsection{Numbered List}
\begin{enumerate}
  \item First step
  \item Second step
  \item Third step
\end{enumerate}

\section{Citations and References}
When you use the citation search feature, references will be automatically managed. For example, you might cite a paper like this: \cite{example2024}. All citations will be automatically added to your \texttt{cite.bib} file.

\section{Conclusion}
This template demonstrates basic \LaTeX{} features including mathematical formulas, tables, and lists. The bibliography system is now set up to work seamlessly with Bibby AI's citation management. Feel free to modify it for your needs!

\bibliographystyle{plain}
\bibliography{cite}

\end{document}\documentclass{article}
\usepackage[utf8]{inputenc}
\usepackage{booktabs}
\usepackage{amsmath}
\usepackage{natbib}

\title{My Document}
\author{Your Name}
\date{\today}

\begin{document}

\maketitle

\section{Introduction}
This is a sample document demonstrating various \LaTeX{} features including mathematical equations, tables, and formatting. You can cite papers using the citation search feature, and they will be automatically added to your \texttt{cite.bib} file.

\section{Mathematical Expressions}
Here are some example equations:

\subsection{Quadratic Formula}
The quadratic formula is given by:
\[x = \frac{-b \pm \sqrt{b^2 - 4ac}}{2a}\]

\subsection{Matrix Example}
A sample matrix:
\[
\begin{pmatrix}
a_{11} & a_{12} & a_{13} \\
a_{21} & a_{22} & a_{23} \\
a_{31} & a_{32} & a_{33}
\end{pmatrix}
\]

\section{Tables}
Below is a sample table showing data:

\begin{table}[h]
\centering
\begin{tabular}{@{}lrrr@{}}
\toprule
\textbf{Category} & \textbf{Value 1} & \textbf{Value 2} & \textbf{Total} \\
\midrule
Item A & 10.5 & 20.3 & 30.8 \\
Item B & 15.2 & 25.1 & 40.3 \\
Item C & 22.1 & 12.7 & 34.8 \\
\bottomrule
\end{tabular}
\caption{Sample Data Table}
\label{tab:sample}
\end{table}

\section{Lists}
\subsection{Bullet Points}
\begin{itemize}
  \item First item
  \item Second item
  \item Third item with sub-items:
    \begin{itemize}
      \item Sub-item 1
      \item Sub-item 2
    \end{itemize}
\end{itemize}

\subsection{Numbered List}
\begin{enumerate}
  \item First step
  \item Second step
  \item Third step
\end{enumerate}

\section{Citations and References}
When you use the citation search feature, references will be automatically managed. For example, you might cite a paper like this: \cite{example2024}. All citations will be automatically added to your \texttt{cite.bib} file.

\section{Conclusion}
This template demonstrates basic \LaTeX{} features including mathematical formulas, tables, and lists. The bibliography system is now set up to work seamlessly with Bibby AI's citation management. Feel free to modify it for your needs!

\bibliographystyle{plain}
\bibliography{cite}

\end{document}\documentclass{article}
\usepackage[utf8]{inputenc}
\usepackage{booktabs}
\usepackage{amsmath}
\usepackage{natbib}

\title{My Document}
\author{Your Name}
\date{\today}

\begin{document}

\maketitle

\section{Introduction}
This is a sample document demonstrating various \LaTeX{} features including mathematical equations, tables, and formatting. You can cite papers using the citation search feature, and they will be automatically added to your \texttt{cite.bib} file.

\section{Mathematical Expressions}
Here are some example equations:

\subsection{Quadratic Formula}
The quadratic formula is given by:
\[x = \frac{-b \pm \sqrt{b^2 - 4ac}}{2a}\]

\subsection{Matrix Example}
A sample matrix:
\[
\begin{pmatrix}
a_{11} & a_{12} & a_{13} \\
a_{21} & a_{22} & a_{23} \\
a_{31} & a_{32} & a_{33}
\end{pmatrix}
\]

\section{Tables}
Below is a sample table showing data:

\begin{table}[h]
\centering
\begin{tabular}{@{}lrrr@{}}
\toprule
\textbf{Category} & \textbf{Value 1} & \textbf{Value 2} & \textbf{Total} \\
\midrule
Item A & 10.5 & 20.3 & 30.8 \\
Item B & 15.2 & 25.1 & 40.3 \\
Item C & 22.1 & 12.7 & 34.8 \\
\bottomrule
\end{tabular}
\caption{Sample Data Table}
\label{tab:sample}
\end{table}

\section{Lists}
\subsection{Bullet Points}
\begin{itemize}
  \item First item
  \item Second item
  \item Third item with sub-items:
    \begin{itemize}
      \item Sub-item 1
      \item Sub-item 2
    \end{itemize}
\end{itemize}

\subsection{Numbered List}
\begin{enumerate}
  \item First step
  \item Second step
  \item Third step
\end{enumerate}

\section{Citations and References}
When you use the citation search feature, references will be automatically managed. For example, you might cite a paper like this: \cite{example2024}. All citations will be automatically added to your \texttt{cite.bib} file.

\section{Conclusion}
This template demonstrates basic \LaTeX{} features including mathematical formulas, tables, and lists. The bibliography system is now set up to work seamlessly with Bibby AI's citation management. Feel free to modify it for your needs!

\bibliographystyle{plain}
\bibliography{cite}

\end{document}\documentclass{article}
\usepackage[utf8]{inputenc}
\usepackage{booktabs}
\usepackage{amsmath}
\usepackage{natbib}

\title{My Document}
\author{Your Name}
\date{\today}

\begin{document}

\maketitle

\section{Introduction}
This is a sample document demonstrating various \LaTeX{} features including mathematical equations, tables, and formatting. You can cite papers using the citation search feature, and they will be automatically added to your \texttt{cite.bib} file.

\section{Mathematical Expressions}
Here are some example equations:

\subsection{Quadratic Formula}
The quadratic formula is given by:
\[x = \frac{-b \pm \sqrt{b^2 - 4ac}}{2a}\]

\subsection{Matrix Example}
A sample matrix:
\[
\begin{pmatrix}
a_{11} & a_{12} & a_{13} \\
a_{21} & a_{22} & a_{23} \\
a_{31} & a_{32} & a_{33}
\end{pmatrix}
\]

\section{Tables}
Below is a sample table showing data:

\begin{table}[h]
\centering
\begin{tabular}{@{}lrrr@{}}
\toprule
\textbf{Category} & \textbf{Value 1} & \textbf{Value 2} & \textbf{Total} \\
\midrule
Item A & 10.5 & 20.3 & 30.8 \\
Item B & 15.2 & 25.1 & 40.3 \\
Item C & 22.1 & 12.7 & 34.8 \\
\bottomrule
\end{tabular}
\caption{Sample Data Table}
\label{tab:sample}
\end{table}

\section{Lists}
\subsection{Bullet Points}
\begin{itemize}
  \item First item
  \item Second item
  \item Third item with sub-items:
    \begin{itemize}
      \item Sub-item 1
      \item Sub-item 2
    \end{itemize}
\end{itemize}

\subsection{Numbered List}
\begin{enumerate}
  \item First step
  \item Second step
  \item Third step
\end{enumerate}

\section{Citations and References}
When you use the citation search feature, references will be automatically managed. For example, you might cite a paper like this: \cite{example2024}. All citations will be automatically added to your \texttt{cite.bib} file.

\section{Conclusion}
This template demonstrates basic \LaTeX{} features including mathematical formulas, tables, and lists. The bibliography system is now set up to work seamlessly with Bibby AI's citation management. Feel free to modify it for your needs!

\bibliographystyle{plain}
\bibliography{cite}

\end{document}\documentclass{article}
\usepackage[utf8]{inputenc}
\usepackage{booktabs}
\usepackage{amsmath}
\usepackage{natbib}

\title{My Document}
\author{Your Name}
\date{\today}

\begin{document}

\maketitle

\section{Introduction}
This is a sample document demonstrating various \LaTeX{} features including mathematical equations, tables, and formatting. You can cite papers using the citation search feature, and they will be automatically added to your \texttt{cite.bib} file.

\section{Mathematical Expressions}
Here are some example equations:

\subsection{Quadratic Formula}
The quadratic formula is given by:
\[x = \frac{-b \pm \sqrt{b^2 - 4ac}}{2a}\]

\subsection{Matrix Example}
A sample matrix:
\[
\begin{pmatrix}
a_{11} & a_{12} & a_{13} \\
a_{21} & a_{22} & a_{23} \\
a_{31} & a_{32} & a_{33}
\end{pmatrix}
\]

\section{Tables}
Below is a sample table showing data:

\begin{table}[h]
\centering
\begin{tabular}{@{}lrrr@{}}
\toprule
\textbf{Category} & \textbf{Value 1} & \textbf{Value 2} & \textbf{Total} \\
\midrule
Item A & 10.5 & 20.3 & 30.8 \\
Item B & 15.2 & 25.1 & 40.3 \\
Item C & 22.1 & 12.7 & 34.8 \\
\bottomrule
\end{tabular}
\caption{Sample Data Table}
\label{tab:sample}
\end{table}

\section{Lists}
\subsection{Bullet Points}
\begin{itemize}
  \item First item
  \item Second item
  \item Third item with sub-items:
    \begin{itemize}
      \item Sub-item 1
      \item Sub-item 2
    \end{itemize}
\end{itemize}

\subsection{Numbered List}
\begin{enumerate}
  \item First step
  \item Second step
  \item Third step
\end{enumerate}

\section{Citations and References}
When you use the citation search feature, references will be automatically managed. For example, you might cite a paper like this: \cite{example2024}. All citations will be automatically added to your \texttt{cite.bib} file.

\section{Conclusion}
This template demonstrates basic \LaTeX{} features including mathematical formulas, tables, and lists. The bibliography system is now set up to work seamlessly with Bibby AI's citation management. Feel free to modify it for your needs!

\bibliographystyle{plain}
\bibliography{cite}

\end{document}\documentclass{article}
\usepackage[utf8]{inputenc}
\usepackage{booktabs}
\usepackage{amsmath}
\usepackage{natbib}

\title{My Document}
\author{Your Name}
\date{\today}

\begin{document}

\maketitle

\section{Introduction}
This is a sample document demonstrating various \LaTeX{} features including mathematical equations, tables, and formatting. You can cite papers using the citation search feature, and they will be automatically added to your \texttt{cite.bib} file.

\section{Mathematical Expressions}
Here are some example equations:

\subsection{Quadratic Formula}
The quadratic formula is given by:
\[x = \frac{-b \pm \sqrt{b^2 - 4ac}}{2a}\]

\subsection{Matrix Example}
A sample matrix:
\[
\begin{pmatrix}
a_{11} & a_{12} & a_{13} \\
a_{21} & a_{22} & a_{23} \\
a_{31} & a_{32} & a_{33}
\end{pmatrix}
\]

\section{Tables}
Below is a sample table showing data:

\begin{table}[h]
\centering
\begin{tabular}{@{}lrrr@{}}
\toprule
\textbf{Category} & \textbf{Value 1} & \textbf{Value 2} & \textbf{Total} \\
\midrule
Item A & 10.5 & 20.3 & 30.8 \\
Item B & 15.2 & 25.1 & 40.3 \\
Item C & 22.1 & 12.7 & 34.8 \\
\bottomrule
\end{tabular}
\caption{Sample Data Table}
\label{tab:sample}
\end{table}

\section{Lists}
\subsection{Bullet Points}
\begin{itemize}
  \item First item
  \item Second item
  \item Third item with sub-items:
    \begin{itemize}
      \item Sub-item 1
      \item Sub-item 2
    \end{itemize}
\end{itemize}

\subsection{Numbered List}
\begin{enumerate}
  \item First step
  \item Second step
  \item Third step
\end{enumerate}

\section{Citations and References}
When you use the citation search feature, references will be automatically managed. For example, you might cite a paper like this: \cite{example2024}. All citations will be automatically added to your \texttt{cite.bib} file.

\section{Conclusion}
This template demonstrates basic \LaTeX{} features including mathematical formulas, tables, and lists. The bibliography system is now set up to work seamlessly with Bibby AI's citation management. Feel free to modify it for your needs!

\bibliographystyle{plain}
\bibliography{cite}

\end{document}\documentclass{article}
\usepackage[utf8]{inputenc}
\usepackage{booktabs}
\usepackage{amsmath}
\usepackage{natbib}

\title{My Document}
\author{Your Name}
\date{\today}

\begin{document}

\maketitle

\section{Introduction}
This is a sample document demonstrating various \LaTeX{} features including mathematical equations, tables, and formatting. You can cite papers using the citation search feature, and they will be automatically added to your \texttt{cite.bib} file.

\section{Mathematical Expressions}
Here are some example equations:

\subsection{Quadratic Formula}
The quadratic formula is given by:
\[x = \frac{-b \pm \sqrt{b^2 - 4ac}}{2a}\]

\subsection{Matrix Example}
A sample matrix:
\[
\begin{pmatrix}
a_{11} & a_{12} & a_{13} \\
a_{21} & a_{22} & a_{23} \\
a_{31} & a_{32} & a_{33}
\end{pmatrix}
\]

\section{Tables}
Below is a sample table showing data:

\begin{table}[h]
\centering
\begin{tabular}{@{}lrrr@{}}
\toprule
\textbf{Category} & \textbf{Value 1} & \textbf{Value 2} & \textbf{Total} \\
\midrule
Item A & 10.5 & 20.3 & 30.8 \\
Item B & 15.2 & 25.1 & 40.3 \\
Item C & 22.1 & 12.7 & 34.8 \\
\bottomrule
\end{tabular}
\caption{Sample Data Table}
\label{tab:sample}
\end{table}

\section{Lists}
\subsection{Bullet Points}
\begin{itemize}
  \item First item
  \item Second item
  \item Third item with sub-items:
    \begin{itemize}
      \item Sub-item 1
      \item Sub-item 2
    \end{itemize}
\end{itemize}

\subsection{Numbered List}
\begin{enumerate}
  \item First step
  \item Second step
  \item Third step
\end{enumerate}

\section{Citations and References}
When you use the citation search feature, references will be automatically managed. For example, you might cite a paper like this: \cite{example2024}. All citations will be automatically added to your \texttt{cite.bib} file.

\section{Conclusion}
This template demonstrates basic \LaTeX{} features including mathematical formulas, tables, and lists. The bibliography system is now set up to work seamlessly with Bibby AI's citation management. Feel free to modify it for your needs!

\bibliographystyle{plain}
\bibliography{cite}

\end{document}\documentclass{article}
\usepackage[utf8]{inputenc}
\usepackage{booktabs}
\usepackage{amsmath}
\usepackage{natbib}

\title{My Document}
\author{Your Name}
\date{\today}

\begin{document}

\maketitle

\section{Introduction}
This is a sample document demonstrating various \LaTeX{} features including mathematical equations, tables, and formatting. You can cite papers using the citation search feature, and they will be automatically added to your \texttt{cite.bib} file.

\section{Mathematical Expressions}
Here are some example equations:

\subsection{Quadratic Formula}
The quadratic formula is given by:
\[x = \frac{-b \pm \sqrt{b^2 - 4ac}}{2a}\]

\subsection{Matrix Example}
A sample matrix:
\[
\begin{pmatrix}
a_{11} & a_{12} & a_{13} \\
a_{21} & a_{22} & a_{23} \\
a_{31} & a_{32} & a_{33}
\end{pmatrix}
\]

\section{Tables}
Below is a sample table showing data:

\begin{table}[h]
\centering
\begin{tabular}{@{}lrrr@{}}
\toprule
\textbf{Category} & \textbf{Value 1} & \textbf{Value 2} & \textbf{Total} \\
\midrule
Item A & 10.5 & 20.3 & 30.8 \\
Item B & 15.2 & 25.1 & 40.3 \\
Item C & 22.1 & 12.7 & 34.8 \\
\bottomrule
\end{tabular}
\caption{Sample Data Table}
\label{tab:sample}
\end{table}

\section{Lists}
\subsection{Bullet Points}
\begin{itemize}
  \item First item
  \item Second item
  \item Third item with sub-items:
    \begin{itemize}
      \item Sub-item 1
      \item Sub-item 2
    \end{itemize}
\end{itemize}

\subsection{Numbered List}
\begin{enumerate}
  \item First step
  \item Second step
  \item Third step
\end{enumerate}

\section{Citations and References}
When you use the citation search feature, references will be automatically managed. For example, you might cite a paper like this: \cite{example2024}. All citations will be automatically added to your \texttt{cite.bib} file.

\section{Conclusion}
This template demonstrates basic \LaTeX{} features including mathematical formulas, tables, and lists. The bibliography system is now set up to work seamlessly with Bibby AI's citation management. Feel free to modify it for your needs!

\bibliographystyle{plain}
\bibliography{cite}

\end{document}\documentclass{article}
\usepackage[utf8]{inputenc}
\usepackage{booktabs}
\usepackage{amsmath}
\usepackage{natbib}

\title{My Document}
\author{Your Name}
\date{\today}

\begin{document}

\maketitle

\section{Introduction}
This is a sample document demonstrating various \LaTeX{} features including mathematical equations, tables, and formatting. You can cite papers using the citation search feature, and they will be automatically added to your \texttt{cite.bib} file.

\section{Mathematical Expressions}
Here are some example equations:

\subsection{Quadratic Formula}
The quadratic formula is given by:
\[x = \frac{-b \pm \sqrt{b^2 - 4ac}}{2a}\]

\subsection{Matrix Example}
A sample matrix:
\[
\begin{pmatrix}
a_{11} & a_{12} & a_{13} \\
a_{21} & a_{22} & a_{23} \\
a_{31} & a_{32} & a_{33}
\end{pmatrix}
\]

\section{Tables}
Below is a sample table showing data:

\begin{table}[h]
\centering
\begin{tabular}{@{}lrrr@{}}
\toprule
\textbf{Category} & \textbf{Value 1} & \textbf{Value 2} & \textbf{Total} \\
\midrule
Item A & 10.5 & 20.3 & 30.8 \\
Item B & 15.2 & 25.1 & 40.3 \\
Item C & 22.1 & 12.7 & 34.8 \\
\bottomrule
\end{tabular}
\caption{Sample Data Table}
\label{tab:sample}
\end{table}

\section{Lists}
\subsection{Bullet Points}
\begin{itemize}
  \item First item
  \item Second item
  \item Third item with sub-items:
    \begin{itemize}
      \item Sub-item 1
      \item Sub-item 2
    \end{itemize}
\end{itemize}

\subsection{Numbered List}
\begin{enumerate}
  \item First step
  \item Second step
  \item Third step
\end{enumerate}

\section{Citations and References}
When you use the citation search feature, references will be automatically managed. For example, you might cite a paper like this: \cite{example2024}. All citations will be automatically added to your \texttt{cite.bib} file.

\section{Conclusion}
This template demonstrates basic \LaTeX{} features including mathematical formulas, tables, and lists. The bibliography system is now set up to work seamlessly with Bibby AI's citation management. Feel free to modify it for your needs!

\bibliographystyle{plain}
\bibliography{cite}

\end{document}\documentclass{article}
\usepackage[utf8]{inputenc}
\usepackage{booktabs}
\usepackage{amsmath}
\usepackage{natbib}

\title{My Document}
\author{Your Name}
\date{\today}

\begin{document}

\maketitle

\section{Introduction}
This is a sample document demonstrating various \LaTeX{} features including mathematical equations, tables, and formatting. You can cite papers using the citation search feature, and they will be automatically added to your \texttt{cite.bib} file.

\section{Mathematical Expressions}
Here are some example equations:

\subsection{Quadratic Formula}
The quadratic formula is given by:
\[x = \frac{-b \pm \sqrt{b^2 - 4ac}}{2a}\]

\subsection{Matrix Example}
A sample matrix:
\[
\begin{pmatrix}
a_{11} & a_{12} & a_{13} \\
a_{21} & a_{22} & a_{23} \\
a_{31} & a_{32} & a_{33}
\end{pmatrix}
\]

\section{Tables}
Below is a sample table showing data:

\begin{table}[h]
\centering
\begin{tabular}{@{}lrrr@{}}
\toprule
\textbf{Category} & \textbf{Value 1} & \textbf{Value 2} & \textbf{Total} \\
\midrule
Item A & 10.5 & 20.3 & 30.8 \\
Item B & 15.2 & 25.1 & 40.3 \\
Item C & 22.1 & 12.7 & 34.8 \\
\bottomrule
\end{tabular}
\caption{Sample Data Table}
\label{tab:sample}
\end{table}

\section{Lists}
\subsection{Bullet Points}
\begin{itemize}
  \item First item
  \item Second item
  \item Third item with sub-items:
    \begin{itemize}
      \item Sub-item 1
      \item Sub-item 2
    \end{itemize}
\end{itemize}

\subsection{Numbered List}
\begin{enumerate}
  \item First step
  \item Second step
  \item Third step
\end{enumerate}

\section{Citations and References}
When you use the citation search feature, references will be automatically managed. For example, you might cite a paper like this: \cite{example2024}. All citations will be automatically added to your \texttt{cite.bib} file.

\section{Conclusion}
This template demonstrates basic \LaTeX{} features including mathematical formulas, tables, and lists. The bibliography system is now set up to work seamlessly with Bibby AI's citation management. Feel free to modify it for your needs!

\bibliographystyle{plain}
\bibliography{cite}

\end{document}\documentclass{article}
\usepackage[utf8]{inputenc}
\usepackage{booktabs}
\usepackage{amsmath}
\usepackage{natbib}

\title{My Document}
\author{Your Name}
\date{\today}

\begin{document}

\maketitle

\section{Introduction}
This is a sample document demonstrating various \LaTeX{} features including mathematical equations, tables, and formatting. You can cite papers using the citation search feature, and they will be automatically added to your \texttt{cite.bib} file.

\section{Mathematical Expressions}
Here are some example equations:

\subsection{Quadratic Formula}
The quadratic formula is given by:
\[x = \frac{-b \pm \sqrt{b^2 - 4ac}}{2a}\]

\subsection{Matrix Example}
A sample matrix:
\[
\begin{pmatrix}
a_{11} & a_{12} & a_{13} \\
a_{21} & a_{22} & a_{23} \\
a_{31} & a_{32} & a_{33}
\end{pmatrix}
\]

\section{Tables}
Below is a sample table showing data:

\begin{table}[h]
\centering
\begin{tabular}{@{}lrrr@{}}
\toprule
\textbf{Category} & \textbf{Value 1} & \textbf{Value 2} & \textbf{Total} \\
\midrule
Item A & 10.5 & 20.3 & 30.8 \\
Item B & 15.2 & 25.1 & 40.3 \\
Item C & 22.1 & 12.7 & 34.8 \\
\bottomrule
\end{tabular}
\caption{Sample Data Table}
\label{tab:sample}
\end{table}

\section{Lists}
\subsection{Bullet Points}
\begin{itemize}
  \item First item
  \item Second item
  \item Third item with sub-items:
    \begin{itemize}
      \item Sub-item 1
      \item Sub-item 2
    \end{itemize}
\end{itemize}

\subsection{Numbered List}
\begin{enumerate}
  \item First step
  \item Second step
  \item Third step
\end{enumerate}

\section{Citations and References}
When you use the citation search feature, references will be automatically managed. For example, you might cite a paper like this: \cite{example2024}. All citations will be automatically added to your \texttt{cite.bib} file.

\section{Conclusion}
This template demonstrates basic \LaTeX{} features including mathematical formulas, tables, and lists. The bibliography system is now set up to work seamlessly with Bibby AI's citation management. Feel free to modify it for your needs!

\bibliographystyle{plain}
\bibliography{cite}

\end{document}\documentclass{article}
\usepackage[utf8]{inputenc}
\usepackage{booktabs}
\usepackage{amsmath}
\usepackage{natbib}

\title{My Document}
\author{Your Name}
\date{\today}

\begin{document}

\maketitle

\section{Introduction}
This is a sample document demonstrating various \LaTeX{} features including mathematical equations, tables, and formatting. You can cite papers using the citation search feature, and they will be automatically added to your \texttt{cite.bib} file.

\section{Mathematical Expressions}
Here are some example equations:

\subsection{Quadratic Formula}
The quadratic formula is given by:
\[x = \frac{-b \pm \sqrt{b^2 - 4ac}}{2a}\]

\subsection{Matrix Example}
A sample matrix:
\[
\begin{pmatrix}
a_{11} & a_{12} & a_{13} \\
a_{21} & a_{22} & a_{23} \\
a_{31} & a_{32} & a_{33}
\end{pmatrix}
\]

\section{Tables}
Below is a sample table showing data:

\begin{table}[h]
\centering
\begin{tabular}{@{}lrrr@{}}
\toprule
\textbf{Category} & \textbf{Value 1} & \textbf{Value 2} & \textbf{Total} \\
\midrule
Item A & 10.5 & 20.3 & 30.8 \\
Item B & 15.2 & 25.1 & 40.3 \\
Item C & 22.1 & 12.7 & 34.8 \\
\bottomrule
\end{tabular}
\caption{Sample Data Table}
\label{tab:sample}
\end{table}

\section{Lists}
\subsection{Bullet Points}
\begin{itemize}
  \item First item
  \item Second item
  \item Third item with sub-items:
    \begin{itemize}
      \item Sub-item 1
      \item Sub-item 2
    \end{itemize}
\end{itemize}

\subsection{Numbered List}
\begin{enumerate}
  \item First step
  \item Second step
  \item Third step
\end{enumerate}

\section{Citations and References}
When you use the citation search feature, references will be automatically managed. For example, you might cite a paper like this: \cite{example2024}. All citations will be automatically added to your \texttt{cite.bib} file.

\section{Conclusion}
This template demonstrates basic \LaTeX{} features including mathematical formulas, tables, and lists. The bibliography system is now set up to work seamlessly with Bibby AI's citation management. Feel free to modify it for your needs!

\bibliographystyle{plain}
\bibliography{cite}

\end{document}\documentclass{article}
\usepackage[utf8]{inputenc}
\usepackage{booktabs}
\usepackage{amsmath}
\usepackage{natbib}

\title{My Document}
\author{Your Name}
\date{\today}

\begin{document}

\maketitle

\section{Introduction}
This is a sample document demonstrating various \LaTeX{} features including mathematical equations, tables, and formatting. You can cite papers using the citation search feature, and they will be automatically added to your \texttt{cite.bib} file.

\section{Mathematical Expressions}
Here are some example equations:

\subsection{Quadratic Formula}
The quadratic formula is given by:
\[x = \frac{-b \pm \sqrt{b^2 - 4ac}}{2a}\]

\subsection{Matrix Example}
A sample matrix:
\[
\begin{pmatrix}
a_{11} & a_{12} & a_{13} \\
a_{21} & a_{22} & a_{23} \\
a_{31} & a_{32} & a_{33}
\end{pmatrix}
\]

\section{Tables}
Below is a sample table showing data:

\begin{table}[h]
\centering
\begin{tabular}{@{}lrrr@{}}
\toprule
\textbf{Category} & \textbf{Value 1} & \textbf{Value 2} & \textbf{Total} \\
\midrule
Item A & 10.5 & 20.3 & 30.8 \\
Item B & 15.2 & 25.1 & 40.3 \\
Item C & 22.1 & 12.7 & 34.8 \\
\bottomrule
\end{tabular}
\caption{Sample Data Table}
\label{tab:sample}
\end{table}

\section{Lists}
\subsection{Bullet Points}
\begin{itemize}
  \item First item
  \item Second item
  \item Third item with sub-items:
    \begin{itemize}
      \item Sub-item 1
      \item Sub-item 2
    \end{itemize}
\end{itemize}

\subsection{Numbered List}
\begin{enumerate}
  \item First step
  \item Second step
  \item Third step
\end{enumerate}

\section{Citations and References}
When you use the citation search feature, references will be automatically managed. For example, you might cite a paper like this: \cite{example2024}. All citations will be automatically added to your \texttt{cite.bib} file.

\section{Conclusion}
This template demonstrates basic \LaTeX{} features including mathematical formulas, tables, and lists. The bibliography system is now set up to work seamlessly with Bibby AI's citation management. Feel free to modify it for your needs!

\bibliographystyle{plain}
\bibliography{cite}

\end{document}\documentclass{article}
\usepackage[utf8]{inputenc}
\usepackage{booktabs}
\usepackage{amsmath}
\usepackage{natbib}

\title{My Document}
\author{Your Name}
\date{\today}

\begin{document}

\maketitle

\section{Introduction}
This is a sample document demonstrating various \LaTeX{} features including mathematical equations, tables, and formatting. You can cite papers using the citation search feature, and they will be automatically added to your \texttt{cite.bib} file.

\section{Mathematical Expressions}
Here are some example equations:

\subsection{Quadratic Formula}
The quadratic formula is given by:
\[x = \frac{-b \pm \sqrt{b^2 - 4ac}}{2a}\]

\subsection{Matrix Example}
A sample matrix:
\[
\begin{pmatrix}
a_{11} & a_{12} & a_{13} \\
a_{21} & a_{22} & a_{23} \\
a_{31} & a_{32} & a_{33}
\end{pmatrix}
\]

\section{Tables}
Below is a sample table showing data:

\begin{table}[h]
\centering
\begin{tabular}{@{}lrrr@{}}
\toprule
\textbf{Category} & \textbf{Value 1} & \textbf{Value 2} & \textbf{Total} \\
\midrule
Item A & 10.5 & 20.3 & 30.8 \\
Item B & 15.2 & 25.1 & 40.3 \\
Item C & 22.1 & 12.7 & 34.8 \\
\bottomrule
\end{tabular}
\caption{Sample Data Table}
\label{tab:sample}
\end{table}

\section{Lists}
\subsection{Bullet Points}
\begin{itemize}
  \item First item
  \item Second item
  \item Third item with sub-items:
    \begin{itemize}
      \item Sub-item 1
      \item Sub-item 2
    \end{itemize}
\end{itemize}

\subsection{Numbered List}
\begin{enumerate}
  \item First step
  \item Second step
  \item Third step
\end{enumerate}

\section{Citations and References}
When you use the citation search feature, references will be automatically managed. For example, you might cite a paper like this: \cite{example2024}. All citations will be automatically added to your \texttt{cite.bib} file.

\section{Conclusion}
This template demonstrates basic \LaTeX{} features including mathematical formulas, tables, and lists. The bibliography system is now set up to work seamlessly with Bibby AI's citation management. Feel free to modify it for your needs!

\bibliographystyle{plain}
\bibliography{cite}

\end{document}\documentclass{article}
\usepackage[utf8]{inputenc}
\usepackage{booktabs}
\usepackage{amsmath}
\usepackage{natbib}

\title{My Document}
\author{Your Name}
\date{\today}

\begin{document}

\maketitle

\section{Introduction}
This is a sample document demonstrating various \LaTeX{} features including mathematical equations, tables, and formatting. You can cite papers using the citation search feature, and they will be automatically added to your \texttt{cite.bib} file.

\section{Mathematical Expressions}
Here are some example equations:

\subsection{Quadratic Formula}
The quadratic formula is given by:
\[x = \frac{-b \pm \sqrt{b^2 - 4ac}}{2a}\]

\subsection{Matrix Example}
A sample matrix:
\[
\begin{pmatrix}
a_{11} & a_{12} & a_{13} \\
a_{21} & a_{22} & a_{23} \\
a_{31} & a_{32} & a_{33}
\end{pmatrix}
\]

\section{Tables}
Below is a sample table showing data:

\begin{table}[h]
\centering
\begin{tabular}{@{}lrrr@{}}
\toprule
\textbf{Category} & \textbf{Value 1} & \textbf{Value 2} & \textbf{Total} \\
\midrule
Item A & 10.5 & 20.3 & 30.8 \\
Item B & 15.2 & 25.1 & 40.3 \\
Item C & 22.1 & 12.7 & 34.8 \\
\bottomrule
\end{tabular}
\caption{Sample Data Table}
\label{tab:sample}
\end{table}

\section{Lists}
\subsection{Bullet Points}
\begin{itemize}
  \item First item
  \item Second item
  \item Third item with sub-items:
    \begin{itemize}
      \item Sub-item 1
      \item Sub-item 2
    \end{itemize}
\end{itemize}

\subsection{Numbered List}
\begin{enumerate}
  \item First step
  \item Second step
  \item Third step
\end{enumerate}

\section{Citations and References}
When you use the citation search feature, references will be automatically managed. For example, you might cite a paper like this: \cite{example2024}. All citations will be automatically added to your \texttt{cite.bib} file.

\section{Conclusion}
This template demonstrates basic \LaTeX{} features including mathematical formulas, tables, and lists. The bibliography system is now set up to work seamlessly with Bibby AI's citation management. Feel free to modify it for your needs!

\bibliographystyle{plain}
\bibliography{cite}

\end{document}\documentclass{article}
\usepackage[utf8]{inputenc}
\usepackage{booktabs}
\usepackage{amsmath}
\usepackage{natbib}

\title{My Document}
\author{Your Name}
\date{\today}

\begin{document}

\maketitle

\section{Introduction}
This is a sample document demonstrating various \LaTeX{} features including mathematical equations, tables, and formatting. You can cite papers using the citation search feature, and they will be automatically added to your \texttt{cite.bib} file.

\section{Mathematical Expressions}
Here are some example equations:

\subsection{Quadratic Formula}
The quadratic formula is given by:
\[x = \frac{-b \pm \sqrt{b^2 - 4ac}}{2a}\]

\subsection{Matrix Example}
A sample matrix:
\[
\begin{pmatrix}
a_{11} & a_{12} & a_{13} \\
a_{21} & a_{22} & a_{23} \\
a_{31} & a_{32} & a_{33}
\end{pmatrix}
\]

\section{Tables}
Below is a sample table showing data:

\begin{table}[h]
\centering
\begin{tabular}{@{}lrrr@{}}
\toprule
\textbf{Category} & \textbf{Value 1} & \textbf{Value 2} & \textbf{Total} \\
\midrule
Item A & 10.5 & 20.3 & 30.8 \\
Item B & 15.2 & 25.1 & 40.3 \\
Item C & 22.1 & 12.7 & 34.8 \\
\bottomrule
\end{tabular}
\caption{Sample Data Table}
\label{tab:sample}
\end{table}

\section{Lists}
\subsection{Bullet Points}
\begin{itemize}
  \item First item
  \item Second item
  \item Third item with sub-items:
    \begin{itemize}
      \item Sub-item 1
      \item Sub-item 2
    \end{itemize}
\end{itemize}

\subsection{Numbered List}
\begin{enumerate}
  \item First step
  \item Second step
  \item Third step
\end{enumerate}

\section{Citations and References}
When you use the citation search feature, references will be automatically managed. For example, you might cite a paper like this: \cite{example2024}. All citations will be automatically added to your \texttt{cite.bib} file.

\section{Conclusion}
This template demonstrates basic \LaTeX{} features including mathematical formulas, tables, and lists. The bibliography system is now set up to work seamlessly with Bibby AI's citation management. Feel free to modify it for your needs!

\bibliographystyle{plain}
\bibliography{cite}

\end{document}\documentclass{article}
\usepackage[utf8]{inputenc}
\usepackage{booktabs}
\usepackage{amsmath}
\usepackage{natbib}

\title{My Document}
\author{Your Name}
\date{\today}

\begin{document}

\maketitle

\section{Introduction}
This is a sample document demonstrating various \LaTeX{} features including mathematical equations, tables, and formatting. You can cite papers using the citation search feature, and they will be automatically added to your \texttt{cite.bib} file.

\section{Mathematical Expressions}
Here are some example equations:

\subsection{Quadratic Formula}
The quadratic formula is given by:
\[x = \frac{-b \pm \sqrt{b^2 - 4ac}}{2a}\]

\subsection{Matrix Example}
A sample matrix:
\[
\begin{pmatrix}
a_{11} & a_{12} & a_{13} \\
a_{21} & a_{22} & a_{23} \\
a_{31} & a_{32} & a_{33}
\end{pmatrix}
\]

\section{Tables}
Below is a sample table showing data:

\begin{table}[h]
\centering
\begin{tabular}{@{}lrrr@{}}
\toprule
\textbf{Category} & \textbf{Value 1} & \textbf{Value 2} & \textbf{Total} \\
\midrule
Item A & 10.5 & 20.3 & 30.8 \\
Item B & 15.2 & 25.1 & 40.3 \\
Item C & 22.1 & 12.7 & 34.8 \\
\bottomrule
\end{tabular}
\caption{Sample Data Table}
\label{tab:sample}
\end{table}

\section{Lists}
\subsection{Bullet Points}
\begin{itemize}
  \item First item
  \item Second item
  \item Third item with sub-items:
    \begin{itemize}
      \item Sub-item 1
      \item Sub-item 2
    \end{itemize}
\end{itemize}

\subsection{Numbered List}
\begin{enumerate}
  \item First step
  \item Second step
  \item Third step
\end{enumerate}

\section{Citations and References}
When you use the citation search feature, references will be automatically managed. For example, you might cite a paper like this: \cite{example2024}. All citations will be automatically added to your \texttt{cite.bib} file.

\section{Conclusion}
This template demonstrates basic \LaTeX{} features including mathematical formulas, tables, and lists. The bibliography system is now set up to work seamlessly with Bibby AI's citation management. Feel free to modify it for your needs!

\bibliographystyle{plain}
\bibliography{cite}

\end{document}\documentclass{article}
\usepackage[utf8]{inputenc}
\usepackage{booktabs}
\usepackage{amsmath}
\usepackage{natbib}

\title{My Document}
\author{Your Name}
\date{\today}

\begin{document}

\maketitle

\section{Introduction}
This is a sample document demonstrating various \LaTeX{} features including mathematical equations, tables, and formatting. You can cite papers using the citation search feature, and they will be automatically added to your \texttt{cite.bib} file.

\section{Mathematical Expressions}
Here are some example equations:

\subsection{Quadratic Formula}
The quadratic formula is given by:
\[x = \frac{-b \pm \sqrt{b^2 - 4ac}}{2a}\]

\subsection{Matrix Example}
A sample matrix:
\[
\begin{pmatrix}
a_{11} & a_{12} & a_{13} \\
a_{21} & a_{22} & a_{23} \\
a_{31} & a_{32} & a_{33}
\end{pmatrix}
\]

\section{Tables}
Below is a sample table showing data:

\begin{table}[h]
\centering
\begin{tabular}{@{}lrrr@{}}
\toprule
\textbf{Category} & \textbf{Value 1} & \textbf{Value 2} & \textbf{Total} \\
\midrule
Item A & 10.5 & 20.3 & 30.8 \\
Item B & 15.2 & 25.1 & 40.3 \\
Item C & 22.1 & 12.7 & 34.8 \\
\bottomrule
\end{tabular}
\caption{Sample Data Table}
\label{tab:sample}
\end{table}

\section{Lists}
\subsection{Bullet Points}
\begin{itemize}
  \item First item
  \item Second item
  \item Third item with sub-items:
    \begin{itemize}
      \item Sub-item 1
      \item Sub-item 2
    \end{itemize}
\end{itemize}

\subsection{Numbered List}
\begin{enumerate}
  \item First step
  \item Second step
  \item Third step
\end{enumerate}

\section{Citations and References}
When you use the citation search feature, references will be automatically managed. For example, you might cite a paper like this: \cite{example2024}. All citations will be automatically added to your \texttt{cite.bib} file.

\section{Conclusion}
This template demonstrates basic \LaTeX{} features including mathematical formulas, tables, and lists. The bibliography system is now set up to work seamlessly with Bibby AI's citation management. Feel free to modify it for your needs!

\bibliographystyle{plain}
\bibliography{cite}

\end{document}\documentclass{article}
\usepackage[utf8]{inputenc}
\usepackage{booktabs}
\usepackage{amsmath}
\usepackage{natbib}

\title{My Document}
\author{Your Name}
\date{\today}

\begin{document}

\maketitle

\section{Introduction}
This is a sample document demonstrating various \LaTeX{} features including mathematical equations, tables, and formatting. You can cite papers using the citation search feature, and they will be automatically added to your \texttt{cite.bib} file.

\section{Mathematical Expressions}
Here are some example equations:

\subsection{Quadratic Formula}
The quadratic formula is given by:
\[x = \frac{-b \pm \sqrt{b^2 - 4ac}}{2a}\]

\subsection{Matrix Example}
A sample matrix:
\[
\begin{pmatrix}
a_{11} & a_{12} & a_{13} \\
a_{21} & a_{22} & a_{23} \\
a_{31} & a_{32} & a_{33}
\end{pmatrix}
\]

\section{Tables}
Below is a sample table showing data:

\begin{table}[h]
\centering
\begin{tabular}{@{}lrrr@{}}
\toprule
\textbf{Category} & \textbf{Value 1} & \textbf{Value 2} & \textbf{Total} \\
\midrule
Item A & 10.5 & 20.3 & 30.8 \\
Item B & 15.2 & 25.1 & 40.3 \\
Item C & 22.1 & 12.7 & 34.8 \\
\bottomrule
\end{tabular}
\caption{Sample Data Table}
\label{tab:sample}
\end{table}

\section{Lists}
\subsection{Bullet Points}
\begin{itemize}
  \item First item
  \item Second item
  \item Third item with sub-items:
    \begin{itemize}
      \item Sub-item 1
      \item Sub-item 2
    \end{itemize}
\end{itemize}

\subsection{Numbered List}
\begin{enumerate}
  \item First step
  \item Second step
  \item Third step
\end{enumerate}

\section{Citations and References}
When you use the citation search feature, references will be automatically managed. For example, you might cite a paper like this: \cite{example2024}. All citations will be automatically added to your \texttt{cite.bib} file.

\section{Conclusion}
This template demonstrates basic \LaTeX{} features including mathematical formulas, tables, and lists. The bibliography system is now set up to work seamlessly with Bibby AI's citation management. Feel free to modify it for your needs!

\bibliographystyle{plain}
\bibliography{cite}

\end{document}\documentclass{article}
\usepackage[utf8]{inputenc}
\usepackage{booktabs}
\usepackage{amsmath}
\usepackage{natbib}

\title{My Document}
\author{Your Name}
\date{\today}

\begin{document}

\maketitle

\section{Introduction}
This is a sample document demonstrating various \LaTeX{} features including mathematical equations, tables, and formatting. You can cite papers using the citation search feature, and they will be automatically added to your \texttt{cite.bib} file.

\section{Mathematical Expressions}
Here are some example equations:

\subsection{Quadratic Formula}
The quadratic formula is given by:
\[x = \frac{-b \pm \sqrt{b^2 - 4ac}}{2a}\]

\subsection{Matrix Example}
A sample matrix:
\[
\begin{pmatrix}
a_{11} & a_{12} & a_{13} \\
a_{21} & a_{22} & a_{23} \\
a_{31} & a_{32} & a_{33}
\end{pmatrix}
\]

\section{Tables}
Below is a sample table showing data:

\begin{table}[h]
\centering
\begin{tabular}{@{}lrrr@{}}
\toprule
\textbf{Category} & \textbf{Value 1} & \textbf{Value 2} & \textbf{Total} \\
\midrule
Item A & 10.5 & 20.3 & 30.8 \\
Item B & 15.2 & 25.1 & 40.3 \\
Item C & 22.1 & 12.7 & 34.8 \\
\bottomrule
\end{tabular}
\caption{Sample Data Table}
\label{tab:sample}
\end{table}

\section{Lists}
\subsection{Bullet Points}
\begin{itemize}
  \item First item
  \item Second item
  \item Third item with sub-items:
    \begin{itemize}
      \item Sub-item 1
      \item Sub-item 2
    \end{itemize}
\end{itemize}

\subsection{Numbered List}
\begin{enumerate}
  \item First step
  \item Second step
  \item Third step
\end{enumerate}

\section{Citations and References}
When you use the citation search feature, references will be automatically managed. For example, you might cite a paper like this: \cite{example2024}. All citations will be automatically added to your \texttt{cite.bib} file.

\section{Conclusion}
This template demonstrates basic \LaTeX{} features including mathematical formulas, tables, and lists. The bibliography system is now set up to work seamlessly with Bibby AI's citation management. Feel free to modify it for your needs!

\bibliographystyle{plain}
\bibliography{cite}

\end{document}\documentclass{article}
\usepackage[utf8]{inputenc}
\usepackage{booktabs}
\usepackage{amsmath}
\usepackage{natbib}

\title{My Document}
\author{Your Name}
\date{\today}

\begin{document}

\maketitle

\section{Introduction}
This is a sample document demonstrating various \LaTeX{} features including mathematical equations, tables, and formatting. You can cite papers using the citation search feature, and they will be automatically added to your \texttt{cite.bib} file.

\section{Mathematical Expressions}
Here are some example equations:

\subsection{Quadratic Formula}
The quadratic formula is given by:
\[x = \frac{-b \pm \sqrt{b^2 - 4ac}}{2a}\]

\subsection{Matrix Example}
A sample matrix:
\[
\begin{pmatrix}
a_{11} & a_{12} & a_{13} \\
a_{21} & a_{22} & a_{23} \\
a_{31} & a_{32} & a_{33}
\end{pmatrix}
\]

\section{Tables}
Below is a sample table showing data:

\begin{table}[h]
\centering
\begin{tabular}{@{}lrrr@{}}
\toprule
\textbf{Category} & \textbf{Value 1} & \textbf{Value 2} & \textbf{Total} \\
\midrule
Item A & 10.5 & 20.3 & 30.8 \\
Item B & 15.2 & 25.1 & 40.3 \\
Item C & 22.1 & 12.7 & 34.8 \\
\bottomrule
\end{tabular}
\caption{Sample Data Table}
\label{tab:sample}
\end{table}

\section{Lists}
\subsection{Bullet Points}
\begin{itemize}
  \item First item
  \item Second item
  \item Third item with sub-items:
    \begin{itemize}
      \item Sub-item 1
      \item Sub-item 2
    \end{itemize}
\end{itemize}

\subsection{Numbered List}
\begin{enumerate}
  \item First step
  \item Second step
  \item Third step
\end{enumerate}

\section{Citations and References}
When you use the citation search feature, references will be automatically managed. For example, you might cite a paper like this: \cite{example2024}. All citations will be automatically added to your \texttt{cite.bib} file.

\section{Conclusion}
This template demonstrates basic \LaTeX{} features including mathematical formulas, tables, and lists. The bibliography system is now set up to work seamlessly with Bibby AI's citation management. Feel free to modify it for your needs!

\bibliographystyle{plain}
\bibliography{cite}

\end{document}\documentclass{article}
\usepackage[utf8]{inputenc}
\usepackage{booktabs}
\usepackage{amsmath}
\usepackage{natbib}

\title{My Document}
\author{Your Name}
\date{\today}

\begin{document}

\maketitle

\section{Introduction}
This is a sample document demonstrating various \LaTeX{} features including mathematical equations, tables, and formatting. You can cite papers using the citation search feature, and they will be automatically added to your \texttt{cite.bib} file.

\section{Mathematical Expressions}
Here are some example equations:

\subsection{Quadratic Formula}
The quadratic formula is given by:
\[x = \frac{-b \pm \sqrt{b^2 - 4ac}}{2a}\]

\subsection{Matrix Example}
A sample matrix:
\[
\begin{pmatrix}
a_{11} & a_{12} & a_{13} \\
a_{21} & a_{22} & a_{23} \\
a_{31} & a_{32} & a_{33}
\end{pmatrix}
\]

\section{Tables}
Below is a sample table showing data:

\begin{table}[h]
\centering
\begin{tabular}{@{}lrrr@{}}
\toprule
\textbf{Category} & \textbf{Value 1} & \textbf{Value 2} & \textbf{Total} \\
\midrule
Item A & 10.5 & 20.3 & 30.8 \\
Item B & 15.2 & 25.1 & 40.3 \\
Item C & 22.1 & 12.7 & 34.8 \\
\bottomrule
\end{tabular}
\caption{Sample Data Table}
\label{tab:sample}
\end{table}

\section{Lists}
\subsection{Bullet Points}
\begin{itemize}
  \item First item
  \item Second item
  \item Third item with sub-items:
    \begin{itemize}
      \item Sub-item 1
      \item Sub-item 2
    \end{itemize}
\end{itemize}

\subsection{Numbered List}
\begin{enumerate}
  \item First step
  \item Second step
  \item Third step
\end{enumerate}

\section{Citations and References}
When you use the citation search feature, references will be automatically managed. For example, you might cite a paper like this: \cite{example2024}. All citations will be automatically added to your \texttt{cite.bib} file.

\section{Conclusion}
This template demonstrates basic \LaTeX{} features including mathematical formulas, tables, and lists. The bibliography system is now set up to work seamlessly with Bibby AI's citation management. Feel free to modify it for your needs!

\bibliographystyle{plain}
\bibliography{cite}

\end{document}\documentclass{article}
\usepackage[utf8]{inputenc}
\usepackage{booktabs}
\usepackage{amsmath}
\usepackage{natbib}

\title{My Document}
\author{Your Name}
\date{\today}

\begin{document}

\maketitle

\section{Introduction}
This is a sample document demonstrating various \LaTeX{} features including mathematical equations, tables, and formatting. You can cite papers using the citation search feature, and they will be automatically added to your \texttt{cite.bib} file.

\section{Mathematical Expressions}
Here are some example equations:

\subsection{Quadratic Formula}
The quadratic formula is given by:
\[x = \frac{-b \pm \sqrt{b^2 - 4ac}}{2a}\]

\subsection{Matrix Example}
A sample matrix:
\[
\begin{pmatrix}
a_{11} & a_{12} & a_{13} \\
a_{21} & a_{22} & a_{23} \\
a_{31} & a_{32} & a_{33}
\end{pmatrix}
\]

\section{Tables}
Below is a sample table showing data:

\begin{table}[h]
\centering
\begin{tabular}{@{}lrrr@{}}
\toprule
\textbf{Category} & \textbf{Value 1} & \textbf{Value 2} & \textbf{Total} \\
\midrule
Item A & 10.5 & 20.3 & 30.8 \\
Item B & 15.2 & 25.1 & 40.3 \\
Item C & 22.1 & 12.7 & 34.8 \\
\bottomrule
\end{tabular}
\caption{Sample Data Table}
\label{tab:sample}
\end{table}

\section{Lists}
\subsection{Bullet Points}
\begin{itemize}
  \item First item
  \item Second item
  \item Third item with sub-items:
    \begin{itemize}
      \item Sub-item 1
      \item Sub-item 2
    \end{itemize}
\end{itemize}

\subsection{Numbered List}
\begin{enumerate}
  \item First step
  \item Second step
  \item Third step
\end{enumerate}

\section{Citations and References}
When you use the citation search feature, references will be automatically managed. For example, you might cite a paper like this: \cite{example2024}. All citations will be automatically added to your \texttt{cite.bib} file.

\section{Conclusion}
This template demonstrates basic \LaTeX{} features including mathematical formulas, tables, and lists. The bibliography system is now set up to work seamlessly with Bibby AI's citation management. Feel free to modify it for your needs!

\bibliographystyle{plain}
\bibliography{cite}

\end{document}\documentclass{article}
\usepackage[utf8]{inputenc}
\usepackage{booktabs}
\usepackage{amsmath}
\usepackage{natbib}

\title{My Document}
\author{Your Name}
\date{\today}

\begin{document}

\maketitle

\section{Introduction}
This is a sample document demonstrating various \LaTeX{} features including mathematical equations, tables, and formatting. You can cite papers using the citation search feature, and they will be automatically added to your \texttt{cite.bib} file.

\section{Mathematical Expressions}
Here are some example equations:

\subsection{Quadratic Formula}
The quadratic formula is given by:
\[x = \frac{-b \pm \sqrt{b^2 - 4ac}}{2a}\]

\subsection{Matrix Example}
A sample matrix:
\[
\begin{pmatrix}
a_{11} & a_{12} & a_{13} \\
a_{21} & a_{22} & a_{23} \\
a_{31} & a_{32} & a_{33}
\end{pmatrix}
\]

\section{Tables}
Below is a sample table showing data:

\begin{table}[h]
\centering
\begin{tabular}{@{}lrrr@{}}
\toprule
\textbf{Category} & \textbf{Value 1} & \textbf{Value 2} & \textbf{Total} \\
\midrule
Item A & 10.5 & 20.3 & 30.8 \\
Item B & 15.2 & 25.1 & 40.3 \\
Item C & 22.1 & 12.7 & 34.8 \\
\bottomrule
\end{tabular}
\caption{Sample Data Table}
\label{tab:sample}
\end{table}

\section{Lists}
\subsection{Bullet Points}
\begin{itemize}
  \item First item
  \item Second item
  \item Third item with sub-items:
    \begin{itemize}
      \item Sub-item 1
      \item Sub-item 2
    \end{itemize}
\end{itemize}

\subsection{Numbered List}
\begin{enumerate}
  \item First step
  \item Second step
  \item Third step
\end{enumerate}

\section{Citations and References}
When you use the citation search feature, references will be automatically managed. For example, you might cite a paper like this: \cite{example2024}. All citations will be automatically added to your \texttt{cite.bib} file.

\section{Conclusion}
This template demonstrates basic \LaTeX{} features including mathematical formulas, tables, and lists. The bibliography system is now set up to work seamlessly with Bibby AI's citation management. Feel free to modify it for your needs!

\bibliographystyle{plain}
\bibliography{cite}

\end{document}\documentclass{article}
\usepackage[utf8]{inputenc}
\usepackage{booktabs}
\usepackage{amsmath}
\usepackage{natbib}

\title{My Document}
\author{Your Name}
\date{\today}

\begin{document}

\maketitle

\section{Introduction}
This is a sample document demonstrating various \LaTeX{} features including mathematical equations, tables, and formatting. You can cite papers using the citation search feature, and they will be automatically added to your \texttt{cite.bib} file.

\section{Mathematical Expressions}
Here are some example equations:

\subsection{Quadratic Formula}
The quadratic formula is given by:
\[x = \frac{-b \pm \sqrt{b^2 - 4ac}}{2a}\]

\subsection{Matrix Example}
A sample matrix:
\[
\begin{pmatrix}
a_{11} & a_{12} & a_{13} \\
a_{21} & a_{22} & a_{23} \\
a_{31} & a_{32} & a_{33}
\end{pmatrix}
\]

\section{Tables}
Below is a sample table showing data:

\begin{table}[h]
\centering
\begin{tabular}{@{}lrrr@{}}
\toprule
\textbf{Category} & \textbf{Value 1} & \textbf{Value 2} & \textbf{Total} \\
\midrule
Item A & 10.5 & 20.3 & 30.8 \\
Item B & 15.2 & 25.1 & 40.3 \\
Item C & 22.1 & 12.7 & 34.8 \\
\bottomrule
\end{tabular}
\caption{Sample Data Table}
\label{tab:sample}
\end{table}

\section{Lists}
\subsection{Bullet Points}
\begin{itemize}
  \item First item
  \item Second item
  \item Third item with sub-items:
    \begin{itemize}
      \item Sub-item 1
      \item Sub-item 2
    \end{itemize}
\end{itemize}

\subsection{Numbered List}
\begin{enumerate}
  \item First step
  \item Second step
  \item Third step
\end{enumerate}

\section{Citations and References}
When you use the citation search feature, references will be automatically managed. For example, you might cite a paper like this: \cite{example2024}. All citations will be automatically added to your \texttt{cite.bib} file.

\section{Conclusion}
This template demonstrates basic \LaTeX{} features including mathematical formulas, tables, and lists. The bibliography system is now set up to work seamlessly with Bibby AI's citation management. Feel free to modify it for your needs!

\bibliographystyle{plain}
\bibliography{cite}

\end{document}\documentclass{article}
\usepackage[utf8]{inputenc}
\usepackage{booktabs}
\usepackage{amsmath}
\usepackage{natbib}

\title{My Document}
\author{Your Name}
\date{\today}

\begin{document}

\maketitle

\section{Introduction}
This is a sample document demonstrating various \LaTeX{} features including mathematical equations, tables, and formatting. You can cite papers using the citation search feature, and they will be automatically added to your \texttt{cite.bib} file.

\section{Mathematical Expressions}
Here are some example equations:

\subsection{Quadratic Formula}
The quadratic formula is given by:
\[x = \frac{-b \pm \sqrt{b^2 - 4ac}}{2a}\]

\subsection{Matrix Example}
A sample matrix:
\[
\begin{pmatrix}
a_{11} & a_{12} & a_{13} \\
a_{21} & a_{22} & a_{23} \\
a_{31} & a_{32} & a_{33}
\end{pmatrix}
\]

\section{Tables}
Below is a sample table showing data:

\begin{table}[h]
\centering
\begin{tabular}{@{}lrrr@{}}
\toprule
\textbf{Category} & \textbf{Value 1} & \textbf{Value 2} & \textbf{Total} \\
\midrule
Item A & 10.5 & 20.3 & 30.8 \\
Item B & 15.2 & 25.1 & 40.3 \\
Item C & 22.1 & 12.7 & 34.8 \\
\bottomrule
\end{tabular}
\caption{Sample Data Table}
\label{tab:sample}
\end{table}

\section{Lists}
\subsection{Bullet Points}
\begin{itemize}
  \item First item
  \item Second item
  \item Third item with sub-items:
    \begin{itemize}
      \item Sub-item 1
      \item Sub-item 2
    \end{itemize}
\end{itemize}

\subsection{Numbered List}
\begin{enumerate}
  \item First step
  \item Second step
  \item Third step
\end{enumerate}

\section{Citations and References}
When you use the citation search feature, references will be automatically managed. For example, you might cite a paper like this: \cite{example2024}. All citations will be automatically added to your \texttt{cite.bib} file.

\section{Conclusion}
This template demonstrates basic \LaTeX{} features including mathematical formulas, tables, and lists. The bibliography system is now set up to work seamlessly with Bibby AI's citation management. Feel free to modify it for your needs!

\bibliographystyle{plain}
\bibliography{cite}

\end{document}\documentclass{article}
\usepackage[utf8]{inputenc}
\usepackage{booktabs}
\usepackage{amsmath}
\usepackage{natbib}

\title{My Document}
\author{Your Name}
\date{\today}

\begin{document}

\maketitle

\section{Introduction}
This is a sample document demonstrating various \LaTeX{} features including mathematical equations, tables, and formatting. You can cite papers using the citation search feature, and they will be automatically added to your \texttt{cite.bib} file.

\section{Mathematical Expressions}
Here are some example equations:

\subsection{Quadratic Formula}
The quadratic formula is given by:
\[x = \frac{-b \pm \sqrt{b^2 - 4ac}}{2a}\]

\subsection{Matrix Example}
A sample matrix:
\[
\begin{pmatrix}
a_{11} & a_{12} & a_{13} \\
a_{21} & a_{22} & a_{23} \\
a_{31} & a_{32} & a_{33}
\end{pmatrix}
\]

\section{Tables}
Below is a sample table showing data:

\begin{table}[h]
\centering
\begin{tabular}{@{}lrrr@{}}
\toprule
\textbf{Category} & \textbf{Value 1} & \textbf{Value 2} & \textbf{Total} \\
\midrule
Item A & 10.5 & 20.3 & 30.8 \\
Item B & 15.2 & 25.1 & 40.3 \\
Item C & 22.1 & 12.7 & 34.8 \\
\bottomrule
\end{tabular}
\caption{Sample Data Table}
\label{tab:sample}
\end{table}

\section{Lists}
\subsection{Bullet Points}
\begin{itemize}
  \item First item
  \item Second item
  \item Third item with sub-items:
    \begin{itemize}
      \item Sub-item 1
      \item Sub-item 2
    \end{itemize}
\end{itemize}

\subsection{Numbered List}
\begin{enumerate}
  \item First step
  \item Second step
  \item Third step
\end{enumerate}

\section{Citations and References}
When you use the citation search feature, references will be automatically managed. For example, you might cite a paper like this: \cite{example2024}. All citations will be automatically added to your \texttt{cite.bib} file.

\section{Conclusion}
This template demonstrates basic \LaTeX{} features including mathematical formulas, tables, and lists. The bibliography system is now set up to work seamlessly with Bibby AI's citation management. Feel free to modify it for your needs!

\bibliographystyle{plain}
\bibliography{cite}

\end{document}\documentclass{article}
\usepackage[utf8]{inputenc}
\usepackage{booktabs}
\usepackage{amsmath}
\usepackage{natbib}

\title{My Document}
\author{Your Name}
\date{\today}

\begin{document}

\maketitle

\section{Introduction}
This is a sample document demonstrating various \LaTeX{} features including mathematical equations, tables, and formatting. You can cite papers using the citation search feature, and they will be automatically added to your \texttt{cite.bib} file.

\section{Mathematical Expressions}
Here are some example equations:

\subsection{Quadratic Formula}
The quadratic formula is given by:
\[x = \frac{-b \pm \sqrt{b^2 - 4ac}}{2a}\]

\subsection{Matrix Example}
A sample matrix:
\[
\begin{pmatrix}
a_{11} & a_{12} & a_{13} \\
a_{21} & a_{22} & a_{23} \\
a_{31} & a_{32} & a_{33}
\end{pmatrix}
\]

\section{Tables}
Below is a sample table showing data:

\begin{table}[h]
\centering
\begin{tabular}{@{}lrrr@{}}
\toprule
\textbf{Category} & \textbf{Value 1} & \textbf{Value 2} & \textbf{Total} \\
\midrule
Item A & 10.5 & 20.3 & 30.8 \\
Item B & 15.2 & 25.1 & 40.3 \\
Item C & 22.1 & 12.7 & 34.8 \\
\bottomrule
\end{tabular}
\caption{Sample Data Table}
\label{tab:sample}
\end{table}

\section{Lists}
\subsection{Bullet Points}
\begin{itemize}
  \item First item
  \item Second item
  \item Third item with sub-items:
    \begin{itemize}
      \item Sub-item 1
      \item Sub-item 2
    \end{itemize}
\end{itemize}

\subsection{Numbered List}
\begin{enumerate}
  \item First step
  \item Second step
  \item Third step
\end{enumerate}

\section{Citations and References}
When you use the citation search feature, references will be automatically managed. For example, you might cite a paper like this: \cite{example2024}. All citations will be automatically added to your \texttt{cite.bib} file.

\section{Conclusion}
This template demonstrates basic \LaTeX{} features including mathematical formulas, tables, and lists. The bibliography system is now set up to work seamlessly with Bibby AI's citation management. Feel free to modify it for your needs!

\bibliographystyle{plain}
\bibliography{cite}

\end{document}\documentclass{article}
\usepackage[utf8]{inputenc}
\usepackage{booktabs}
\usepackage{amsmath}
\usepackage{natbib}

\title{My Document}
\author{Your Name}
\date{\today}

\begin{document}

\maketitle

\section{Introduction}
This is a sample document demonstrating various \LaTeX{} features including mathematical equations, tables, and formatting. You can cite papers using the citation search feature, and they will be automatically added to your \texttt{cite.bib} file.

\section{Mathematical Expressions}
Here are some example equations:

\subsection{Quadratic Formula}
The quadratic formula is given by:
\[x = \frac{-b \pm \sqrt{b^2 - 4ac}}{2a}\]

\subsection{Matrix Example}
A sample matrix:
\[
\begin{pmatrix}
a_{11} & a_{12} & a_{13} \\
a_{21} & a_{22} & a_{23} \\
a_{31} & a_{32} & a_{33}
\end{pmatrix}
\]

\section{Tables}
Below is a sample table showing data:

\begin{table}[h]
\centering
\begin{tabular}{@{}lrrr@{}}
\toprule
\textbf{Category} & \textbf{Value 1} & \textbf{Value 2} & \textbf{Total} \\
\midrule
Item A & 10.5 & 20.3 & 30.8 \\
Item B & 15.2 & 25.1 & 40.3 \\
Item C & 22.1 & 12.7 & 34.8 \\
\bottomrule
\end{tabular}
\caption{Sample Data Table}
\label{tab:sample}
\end{table}

\section{Lists}
\subsection{Bullet Points}
\begin{itemize}
  \item First item
  \item Second item
  \item Third item with sub-items:
    \begin{itemize}
      \item Sub-item 1
      \item Sub-item 2
    \end{itemize}
\end{itemize}

\subsection{Numbered List}
\begin{enumerate}
  \item First step
  \item Second step
  \item Third step
\end{enumerate}

\section{Citations and References}
When you use the citation search feature, references will be automatically managed. For example, you might cite a paper like this: \cite{example2024}. All citations will be automatically added to your \texttt{cite.bib} file.

\section{Conclusion}
This template demonstrates basic \LaTeX{} features including mathematical formulas, tables, and lists. The bibliography system is now set up to work seamlessly with Bibby AI's citation management. Feel free to modify it for your needs!

\bibliographystyle{plain}
\bibliography{cite}

\end{document}\documentclass{article}
\usepackage[utf8]{inputenc}
\usepackage{booktabs}
\usepackage{amsmath}
\usepackage{natbib}

\title{My Document}
\author{Your Name}
\date{\today}

\begin{document}

\maketitle

\section{Introduction}
This is a sample document demonstrating various \LaTeX{} features including mathematical equations, tables, and formatting. You can cite papers using the citation search feature, and they will be automatically added to your \texttt{cite.bib} file.

\section{Mathematical Expressions}
Here are some example equations:

\subsection{Quadratic Formula}
The quadratic formula is given by:
\[x = \frac{-b \pm \sqrt{b^2 - 4ac}}{2a}\]

\subsection{Matrix Example}
A sample matrix:
\[
\begin{pmatrix}
a_{11} & a_{12} & a_{13} \\
a_{21} & a_{22} & a_{23} \\
a_{31} & a_{32} & a_{33}
\end{pmatrix}
\]

\section{Tables}
Below is a sample table showing data:

\begin{table}[h]
\centering
\begin{tabular}{@{}lrrr@{}}
\toprule
\textbf{Category} & \textbf{Value 1} & \textbf{Value 2} & \textbf{Total} \\
\midrule
Item A & 10.5 & 20.3 & 30.8 \\
Item B & 15.2 & 25.1 & 40.3 \\
Item C & 22.1 & 12.7 & 34.8 \\
\bottomrule
\end{tabular}
\caption{Sample Data Table}
\label{tab:sample}
\end{table}

\section{Lists}
\subsection{Bullet Points}
\begin{itemize}
  \item First item
  \item Second item
  \item Third item with sub-items:
    \begin{itemize}
      \item Sub-item 1
      \item Sub-item 2
    \end{itemize}
\end{itemize}

\subsection{Numbered List}
\begin{enumerate}
  \item First step
  \item Second step
  \item Third step
\end{enumerate}

\section{Citations and References}
When you use the citation search feature, references will be automatically managed. For example, you might cite a paper like this: \cite{example2024}. All citations will be automatically added to your \texttt{cite.bib} file.

\section{Conclusion}
This template demonstrates basic \LaTeX{} features including mathematical formulas, tables, and lists. The bibliography system is now set up to work seamlessly with Bibby AI's citation management. Feel free to modify it for your needs!

\bibliographystyle{plain}
\bibliography{cite}

\end{document}\documentclass{article}
\usepackage[utf8]{inputenc}
\usepackage{booktabs}
\usepackage{amsmath}
\usepackage{natbib}

\title{My Document}
\author{Your Name}
\date{\today}

\begin{document}

\maketitle

\section{Introduction}
This is a sample document demonstrating various \LaTeX{} features including mathematical equations, tables, and formatting. You can cite papers using the citation search feature, and they will be automatically added to your \texttt{cite.bib} file.

\section{Mathematical Expressions}
Here are some example equations:

\subsection{Quadratic Formula}
The quadratic formula is given by:
\[x = \frac{-b \pm \sqrt{b^2 - 4ac}}{2a}\]

\subsection{Matrix Example}
A sample matrix:
\[
\begin{pmatrix}
a_{11} & a_{12} & a_{13} \\
a_{21} & a_{22} & a_{23} \\
a_{31} & a_{32} & a_{33}
\end{pmatrix}
\]

\section{Tables}
Below is a sample table showing data:

\begin{table}[h]
\centering
\begin{tabular}{@{}lrrr@{}}
\toprule
\textbf{Category} & \textbf{Value 1} & \textbf{Value 2} & \textbf{Total} \\
\midrule
Item A & 10.5 & 20.3 & 30.8 \\
Item B & 15.2 & 25.1 & 40.3 \\
Item C & 22.1 & 12.7 & 34.8 \\
\bottomrule
\end{tabular}
\caption{Sample Data Table}
\label{tab:sample}
\end{table}

\section{Lists}
\subsection{Bullet Points}
\begin{itemize}
  \item First item
  \item Second item
  \item Third item with sub-items:
    \begin{itemize}
      \item Sub-item 1
      \item Sub-item 2
    \end{itemize}
\end{itemize}

\subsection{Numbered List}
\begin{enumerate}
  \item First step
  \item Second step
  \item Third step
\end{enumerate}

\section{Citations and References}
When you use the citation search feature, references will be automatically managed. For example, you might cite a paper like this: \cite{example2024}. All citations will be automatically added to your \texttt{cite.bib} file.

\section{Conclusion}
This template demonstrates basic \LaTeX{} features including mathematical formulas, tables, and lists. The bibliography system is now set up to work seamlessly with Bibby AI's citation management. Feel free to modify it for your needs!

\bibliographystyle{plain}
\bibliography{cite}

\end{document}\documentclass{article}
\usepackage[utf8]{inputenc}
\usepackage{booktabs}
\usepackage{amsmath}
\usepackage{natbib}

\title{My Document}
\author{Your Name}
\date{\today}

\begin{document}

\maketitle

\section{Introduction}
This is a sample document demonstrating various \LaTeX{} features including mathematical equations, tables, and formatting. You can cite papers using the citation search feature, and they will be automatically added to your \texttt{cite.bib} file.

\section{Mathematical Expressions}
Here are some example equations:

\subsection{Quadratic Formula}
The quadratic formula is given by:
\[x = \frac{-b \pm \sqrt{b^2 - 4ac}}{2a}\]

\subsection{Matrix Example}
A sample matrix:
\[
\begin{pmatrix}
a_{11} & a_{12} & a_{13} \\
a_{21} & a_{22} & a_{23} \\
a_{31} & a_{32} & a_{33}
\end{pmatrix}
\]

\section{Tables}
Below is a sample table showing data:

\begin{table}[h]
\centering
\begin{tabular}{@{}lrrr@{}}
\toprule
\textbf{Category} & \textbf{Value 1} & \textbf{Value 2} & \textbf{Total} \\
\midrule
Item A & 10.5 & 20.3 & 30.8 \\
Item B & 15.2 & 25.1 & 40.3 \\
Item C & 22.1 & 12.7 & 34.8 \\
\bottomrule
\end{tabular}
\caption{Sample Data Table}
\label{tab:sample}
\end{table}

\section{Lists}
\subsection{Bullet Points}
\begin{itemize}
  \item First item
  \item Second item
  \item Third item with sub-items:
    \begin{itemize}
      \item Sub-item 1
      \item Sub-item 2
    \end{itemize}
\end{itemize}

\subsection{Numbered List}
\begin{enumerate}
  \item First step
  \item Second step
  \item Third step
\end{enumerate}

\section{Citations and References}
When you use the citation search feature, references will be automatically managed. For example, you might cite a paper like this: \cite{example2024}. All citations will be automatically added to your \texttt{cite.bib} file.

\section{Conclusion}
This template demonstrates basic \LaTeX{} features including mathematical formulas, tables, and lists. The bibliography system is now set up to work seamlessly with Bibby AI's citation management. Feel free to modify it for your needs!

\bibliographystyle{plain}
\bibliography{cite}

\end{document}\documentclass{article}
\usepackage[utf8]{inputenc}
\usepackage{booktabs}
\usepackage{amsmath}
\usepackage{natbib}

\title{My Document}
\author{Your Name}
\date{\today}

\begin{document}

\maketitle

\section{Introduction}
This is a sample document demonstrating various \LaTeX{} features including mathematical equations, tables, and formatting. You can cite papers using the citation search feature, and they will be automatically added to your \texttt{cite.bib} file.

\section{Mathematical Expressions}
Here are some example equations:

\subsection{Quadratic Formula}
The quadratic formula is given by:
\[x = \frac{-b \pm \sqrt{b^2 - 4ac}}{2a}\]

\subsection{Matrix Example}
A sample matrix:
\[
\begin{pmatrix}
a_{11} & a_{12} & a_{13} \\
a_{21} & a_{22} & a_{23} \\
a_{31} & a_{32} & a_{33}
\end{pmatrix}
\]

\section{Tables}
Below is a sample table showing data:

\begin{table}[h]
\centering
\begin{tabular}{@{}lrrr@{}}
\toprule
\textbf{Category} & \textbf{Value 1} & \textbf{Value 2} & \textbf{Total} \\
\midrule
Item A & 10.5 & 20.3 & 30.8 \\
Item B & 15.2 & 25.1 & 40.3 \\
Item C & 22.1 & 12.7 & 34.8 \\
\bottomrule
\end{tabular}
\caption{Sample Data Table}
\label{tab:sample}
\end{table}

\section{Lists}
\subsection{Bullet Points}
\begin{itemize}
  \item First item
  \item Second item
  \item Third item with sub-items:
    \begin{itemize}
      \item Sub-item 1
      \item Sub-item 2
    \end{itemize}
\end{itemize}

\subsection{Numbered List}
\begin{enumerate}
  \item First step
  \item Second step
  \item Third step
\end{enumerate}

\section{Citations and References}
When you use the citation search feature, references will be automatically managed. For example, you might cite a paper like this: \cite{example2024}. All citations will be automatically added to your \texttt{cite.bib} file.

\section{Conclusion}
This template demonstrates basic \LaTeX{} features including mathematical formulas, tables, and lists. The bibliography system is now set up to work seamlessly with Bibby AI's citation management. Feel free to modify it for your needs!

\bibliographystyle{plain}
\bibliography{cite}

\end{document}\documentclass{article}
\usepackage[utf8]{inputenc}
\usepackage{booktabs}
\usepackage{amsmath}
\usepackage{natbib}

\title{My Document}
\author{Your Name}
\date{\today}

\begin{document}

\maketitle

\section{Introduction}
This is a sample document demonstrating various \LaTeX{} features including mathematical equations, tables, and formatting. You can cite papers using the citation search feature, and they will be automatically added to your \texttt{cite.bib} file.

\section{Mathematical Expressions}
Here are some example equations:

\subsection{Quadratic Formula}
The quadratic formula is given by:
\[x = \frac{-b \pm \sqrt{b^2 - 4ac}}{2a}\]

\subsection{Matrix Example}
A sample matrix:
\[
\begin{pmatrix}
a_{11} & a_{12} & a_{13} \\
a_{21} & a_{22} & a_{23} \\
a_{31} & a_{32} & a_{33}
\end{pmatrix}
\]

\section{Tables}
Below is a sample table showing data:

\begin{table}[h]
\centering
\begin{tabular}{@{}lrrr@{}}
\toprule
\textbf{Category} & \textbf{Value 1} & \textbf{Value 2} & \textbf{Total} \\
\midrule
Item A & 10.5 & 20.3 & 30.8 \\
Item B & 15.2 & 25.1 & 40.3 \\
Item C & 22.1 & 12.7 & 34.8 \\
\bottomrule
\end{tabular}
\caption{Sample Data Table}
\label{tab:sample}
\end{table}

\section{Lists}
\subsection{Bullet Points}
\begin{itemize}
  \item First item
  \item Second item
  \item Third item with sub-items:
    \begin{itemize}
      \item Sub-item 1
      \item Sub-item 2
    \end{itemize}
\end{itemize}

\subsection{Numbered List}
\begin{enumerate}
  \item First step
  \item Second step
  \item Third step
\end{enumerate}

\section{Citations and References}
When you use the citation search feature, references will be automatically managed. For example, you might cite a paper like this: \cite{example2024}. All citations will be automatically added to your \texttt{cite.bib} file.

\section{Conclusion}
This template demonstrates basic \LaTeX{} features including mathematical formulas, tables, and lists. The bibliography system is now set up to work seamlessly with Bibby AI's citation management. Feel free to modify it for your needs!

\bibliographystyle{plain}
\bibliography{cite}

\end{document}\documentclass{article}
\usepackage[utf8]{inputenc}
\usepackage{booktabs}
\usepackage{amsmath}
\usepackage{natbib}

\title{My Document}
\author{Your Name}
\date{\today}

\begin{document}

\maketitle

\section{Introduction}
This is a sample document demonstrating various \LaTeX{} features including mathematical equations, tables, and formatting. You can cite papers using the citation search feature, and they will be automatically added to your \texttt{cite.bib} file.

\section{Mathematical Expressions}
Here are some example equations:

\subsection{Quadratic Formula}
The quadratic formula is given by:
\[x = \frac{-b \pm \sqrt{b^2 - 4ac}}{2a}\]

\subsection{Matrix Example}
A sample matrix:
\[
\begin{pmatrix}
a_{11} & a_{12} & a_{13} \\
a_{21} & a_{22} & a_{23} \\
a_{31} & a_{32} & a_{33}
\end{pmatrix}
\]

\section{Tables}
Below is a sample table showing data:

\begin{table}[h]
\centering
\begin{tabular}{@{}lrrr@{}}
\toprule
\textbf{Category} & \textbf{Value 1} & \textbf{Value 2} & \textbf{Total} \\
\midrule
Item A & 10.5 & 20.3 & 30.8 \\
Item B & 15.2 & 25.1 & 40.3 \\
Item C & 22.1 & 12.7 & 34.8 \\
\bottomrule
\end{tabular}
\caption{Sample Data Table}
\label{tab:sample}
\end{table}

\section{Lists}
\subsection{Bullet Points}
\begin{itemize}
  \item First item
  \item Second item
  \item Third item with sub-items:
    \begin{itemize}
      \item Sub-item 1
      \item Sub-item 2
    \end{itemize}
\end{itemize}

\subsection{Numbered List}
\begin{enumerate}
  \item First step
  \item Second step
  \item Third step
\end{enumerate}

\section{Citations and References}
When you use the citation search feature, references will be automatically managed. For example, you might cite a paper like this: \cite{example2024}. All citations will be automatically added to your \texttt{cite.bib} file.

\section{Conclusion}
This template demonstrates basic \LaTeX{} features including mathematical formulas, tables, and lists. The bibliography system is now set up to work seamlessly with Bibby AI's citation management. Feel free to modify it for your needs!

\bibliographystyle{plain}
\bibliography{cite}

\end{document}\documentclass{article}
\usepackage[utf8]{inputenc}
\usepackage{booktabs}
\usepackage{amsmath}
\usepackage{natbib}

\title{My Document}
\author{Your Name}
\date{\today}

\begin{document}

\maketitle

\section{Introduction}
This is a sample document demonstrating various \LaTeX{} features including mathematical equations, tables, and formatting. You can cite papers using the citation search feature, and they will be automatically added to your \texttt{cite.bib} file.

\section{Mathematical Expressions}
Here are some example equations:

\subsection{Quadratic Formula}
The quadratic formula is given by:
\[x = \frac{-b \pm \sqrt{b^2 - 4ac}}{2a}\]

\subsection{Matrix Example}
A sample matrix:
\[
\begin{pmatrix}
a_{11} & a_{12} & a_{13} \\
a_{21} & a_{22} & a_{23} \\
a_{31} & a_{32} & a_{33}
\end{pmatrix}
\]

\section{Tables}
Below is a sample table showing data:

\begin{table}[h]
\centering
\begin{tabular}{@{}lrrr@{}}
\toprule
\textbf{Category} & \textbf{Value 1} & \textbf{Value 2} & \textbf{Total} \\
\midrule
Item A & 10.5 & 20.3 & 30.8 \\
Item B & 15.2 & 25.1 & 40.3 \\
Item C & 22.1 & 12.7 & 34.8 \\
\bottomrule
\end{tabular}
\caption{Sample Data Table}
\label{tab:sample}
\end{table}

\section{Lists}
\subsection{Bullet Points}
\begin{itemize}
  \item First item
  \item Second item
  \item Third item with sub-items:
    \begin{itemize}
      \item Sub-item 1
      \item Sub-item 2
    \end{itemize}
\end{itemize}

\subsection{Numbered List}
\begin{enumerate}
  \item First step
  \item Second step
  \item Third step
\end{enumerate}

\section{Citations and References}
When you use the citation search feature, references will be automatically managed. For example, you might cite a paper like this: \cite{example2024}. All citations will be automatically added to your \texttt{cite.bib} file.

\section{Conclusion}
This template demonstrates basic \LaTeX{} features including mathematical formulas, tables, and lists. The bibliography system is now set up to work seamlessly with Bibby AI's citation management. Feel free to modify it for your needs!

\bibliographystyle{plain}
\bibliography{cite}

\end{document}\documentclass{article}
\usepackage[utf8]{inputenc}
\usepackage{booktabs}
\usepackage{amsmath}
\usepackage{natbib}

\title{My Document}
\author{Your Name}
\date{\today}

\begin{document}

\maketitle

\section{Introduction}
This is a sample document demonstrating various \LaTeX{} features including mathematical equations, tables, and formatting. You can cite papers using the citation search feature, and they will be automatically added to your \texttt{cite.bib} file.

\section{Mathematical Expressions}
Here are some example equations:

\subsection{Quadratic Formula}
The quadratic formula is given by:
\[x = \frac{-b \pm \sqrt{b^2 - 4ac}}{2a}\]

\subsection{Matrix Example}
A sample matrix:
\[
\begin{pmatrix}
a_{11} & a_{12} & a_{13} \\
a_{21} & a_{22} & a_{23} \\
a_{31} & a_{32} & a_{33}
\end{pmatrix}
\]

\section{Tables}
Below is a sample table showing data:

\begin{table}[h]
\centering
\begin{tabular}{@{}lrrr@{}}
\toprule
\textbf{Category} & \textbf{Value 1} & \textbf{Value 2} & \textbf{Total} \\
\midrule
Item A & 10.5 & 20.3 & 30.8 \\
Item B & 15.2 & 25.1 & 40.3 \\
Item C & 22.1 & 12.7 & 34.8 \\
\bottomrule
\end{tabular}
\caption{Sample Data Table}
\label{tab:sample}
\end{table}

\section{Lists}
\subsection{Bullet Points}
\begin{itemize}
  \item First item
  \item Second item
  \item Third item with sub-items:
    \begin{itemize}
      \item Sub-item 1
      \item Sub-item 2
    \end{itemize}
\end{itemize}

\subsection{Numbered List}
\begin{enumerate}
  \item First step
  \item Second step
  \item Third step
\end{enumerate}

\section{Citations and References}
When you use the citation search feature, references will be automatically managed. For example, you might cite a paper like this: \cite{example2024}. All citations will be automatically added to your \texttt{cite.bib} file.

\section{Conclusion}
This template demonstrates basic \LaTeX{} features including mathematical formulas, tables, and lists. The bibliography system is now set up to work seamlessly with Bibby AI's citation management. Feel free to modify it for your needs!

\bibliographystyle{plain}
\bibliography{cite}

\end{document}\documentclass{article}
\usepackage[utf8]{inputenc}
\usepackage{booktabs}
\usepackage{amsmath}
\usepackage{natbib}

\title{My Document}
\author{Your Name}
\date{\today}

\begin{document}

\maketitle

\section{Introduction}
This is a sample document demonstrating various \LaTeX{} features including mathematical equations, tables, and formatting. You can cite papers using the citation search feature, and they will be automatically added to your \texttt{cite.bib} file.

\section{Mathematical Expressions}
Here are some example equations:

\subsection{Quadratic Formula}
The quadratic formula is given by:
\[x = \frac{-b \pm \sqrt{b^2 - 4ac}}{2a}\]

\subsection{Matrix Example}
A sample matrix:
\[
\begin{pmatrix}
a_{11} & a_{12} & a_{13} \\
a_{21} & a_{22} & a_{23} \\
a_{31} & a_{32} & a_{33}
\end{pmatrix}
\]

\section{Tables}
Below is a sample table showing data:

\begin{table}[h]
\centering
\begin{tabular}{@{}lrrr@{}}
\toprule
\textbf{Category} & \textbf{Value 1} & \textbf{Value 2} & \textbf{Total} \\
\midrule
Item A & 10.5 & 20.3 & 30.8 \\
Item B & 15.2 & 25.1 & 40.3 \\
Item C & 22.1 & 12.7 & 34.8 \\
\bottomrule
\end{tabular}
\caption{Sample Data Table}
\label{tab:sample}
\end{table}

\section{Lists}
\subsection{Bullet Points}
\begin{itemize}
  \item First item
  \item Second item
  \item Third item with sub-items:
    \begin{itemize}
      \item Sub-item 1
      \item Sub-item 2
    \end{itemize}
\end{itemize}

\subsection{Numbered List}
\begin{enumerate}
  \item First step
  \item Second step
  \item Third step
\end{enumerate}

\section{Citations and References}
When you use the citation search feature, references will be automatically managed. For example, you might cite a paper like this: \cite{example2024}. All citations will be automatically added to your \texttt{cite.bib} file.

\section{Conclusion}
This template demonstrates basic \LaTeX{} features including mathematical formulas, tables, and lists. The bibliography system is now set up to work seamlessly with Bibby AI's citation management. Feel free to modify it for your needs!

\bibliographystyle{plain}
\bibliography{cite}

\end{document}\documentclass{article}
\usepackage[utf8]{inputenc}
\usepackage{booktabs}
\usepackage{amsmath}
\usepackage{natbib}

\title{My Document}
\author{Your Name}
\date{\today}

\begin{document}

\maketitle

\section{Introduction}
This is a sample document demonstrating various \LaTeX{} features including mathematical equations, tables, and formatting. You can cite papers using the citation search feature, and they will be automatically added to your \texttt{cite.bib} file.

\section{Mathematical Expressions}
Here are some example equations:

\subsection{Quadratic Formula}
The quadratic formula is given by:
\[x = \frac{-b \pm \sqrt{b^2 - 4ac}}{2a}\]

\subsection{Matrix Example}
A sample matrix:
\[
\begin{pmatrix}
a_{11} & a_{12} & a_{13} \\
a_{21} & a_{22} & a_{23} \\
a_{31} & a_{32} & a_{33}
\end{pmatrix}
\]

\section{Tables}
Below is a sample table showing data:

\begin{table}[h]
\centering
\begin{tabular}{@{}lrrr@{}}
\toprule
\textbf{Category} & \textbf{Value 1} & \textbf{Value 2} & \textbf{Total} \\
\midrule
Item A & 10.5 & 20.3 & 30.8 \\
Item B & 15.2 & 25.1 & 40.3 \\
Item C & 22.1 & 12.7 & 34.8 \\
\bottomrule
\end{tabular}
\caption{Sample Data Table}
\label{tab:sample}
\end{table}

\section{Lists}
\subsection{Bullet Points}
\begin{itemize}
  \item First item
  \item Second item
  \item Third item with sub-items:
    \begin{itemize}
      \item Sub-item 1
      \item Sub-item 2
    \end{itemize}
\end{itemize}

\subsection{Numbered List}
\begin{enumerate}
  \item First step
  \item Second step
  \item Third step
\end{enumerate}

\section{Citations and References}
When you use the citation search feature, references will be automatically managed. For example, you might cite a paper like this: \cite{example2024}. All citations will be automatically added to your \texttt{cite.bib} file.

\section{Conclusion}
This template demonstrates basic \LaTeX{} features including mathematical formulas, tables, and lists. The bibliography system is now set up to work seamlessly with Bibby AI's citation management. Feel free to modify it for your needs!

\bibliographystyle{plain}
\bibliography{cite}

\end{document}\documentclass{article}
\usepackage[utf8]{inputenc}
\usepackage{booktabs}
\usepackage{amsmath}
\usepackage{natbib}

\title{My Document}
\author{Your Name}
\date{\today}

\begin{document}

\maketitle

\section{Introduction}
This is a sample document demonstrating various \LaTeX{} features including mathematical equations, tables, and formatting. You can cite papers using the citation search feature, and they will be automatically added to your \texttt{cite.bib} file.

\section{Mathematical Expressions}
Here are some example equations:

\subsection{Quadratic Formula}
The quadratic formula is given by:
\[x = \frac{-b \pm \sqrt{b^2 - 4ac}}{2a}\]

\subsection{Matrix Example}
A sample matrix:
\[
\begin{pmatrix}
a_{11} & a_{12} & a_{13} \\
a_{21} & a_{22} & a_{23} \\
a_{31} & a_{32} & a_{33}
\end{pmatrix}
\]

\section{Tables}
Below is a sample table showing data:

\begin{table}[h]
\centering
\begin{tabular}{@{}lrrr@{}}
\toprule
\textbf{Category} & \textbf{Value 1} & \textbf{Value 2} & \textbf{Total} \\
\midrule
Item A & 10.5 & 20.3 & 30.8 \\
Item B & 15.2 & 25.1 & 40.3 \\
Item C & 22.1 & 12.7 & 34.8 \\
\bottomrule
\end{tabular}
\caption{Sample Data Table}
\label{tab:sample}
\end{table}

\section{Lists}
\subsection{Bullet Points}
\begin{itemize}
  \item First item
  \item Second item
  \item Third item with sub-items:
    \begin{itemize}
      \item Sub-item 1
      \item Sub-item 2
    \end{itemize}
\end{itemize}

\subsection{Numbered List}
\begin{enumerate}
  \item First step
  \item Second step
  \item Third step
\end{enumerate}

\section{Citations and References}
When you use the citation search feature, references will be automatically managed. For example, you might cite a paper like this: \cite{example2024}. All citations will be automatically added to your \texttt{cite.bib} file.

\section{Conclusion}
This template demonstrates basic \LaTeX{} features including mathematical formulas, tables, and lists. The bibliography system is now set up to work seamlessly with Bibby AI's citation management. Feel free to modify it for your needs!

\bibliographystyle{plain}
\bibliography{cite}

\end{document}\documentclass{article}
\usepackage[utf8]{inputenc}
\usepackage{booktabs}
\usepackage{amsmath}
\usepackage{natbib}

\title{My Document}
\author{Your Name}
\date{\today}

\begin{document}

\maketitle

\section{Introduction}
This is a sample document demonstrating various \LaTeX{} features including mathematical equations, tables, and formatting. You can cite papers using the citation search feature, and they will be automatically added to your \texttt{cite.bib} file.

\section{Mathematical Expressions}
Here are some example equations:

\subsection{Quadratic Formula}
The quadratic formula is given by:
\[x = \frac{-b \pm \sqrt{b^2 - 4ac}}{2a}\]

\subsection{Matrix Example}
A sample matrix:
\[
\begin{pmatrix}
a_{11} & a_{12} & a_{13} \\
a_{21} & a_{22} & a_{23} \\
a_{31} & a_{32} & a_{33}
\end{pmatrix}
\]

\section{Tables}
Below is a sample table showing data:

\begin{table}[h]
\centering
\begin{tabular}{@{}lrrr@{}}
\toprule
\textbf{Category} & \textbf{Value 1} & \textbf{Value 2} & \textbf{Total} \\
\midrule
Item A & 10.5 & 20.3 & 30.8 \\
Item B & 15.2 & 25.1 & 40.3 \\
Item C & 22.1 & 12.7 & 34.8 \\
\bottomrule
\end{tabular}
\caption{Sample Data Table}
\label{tab:sample}
\end{table}

\section{Lists}
\subsection{Bullet Points}
\begin{itemize}
  \item First item
  \item Second item
  \item Third item with sub-items:
    \begin{itemize}
      \item Sub-item 1
      \item Sub-item 2
    \end{itemize}
\end{itemize}

\subsection{Numbered List}
\begin{enumerate}
  \item First step
  \item Second step
  \item Third step
\end{enumerate}

\section{Citations and References}
When you use the citation search feature, references will be automatically managed. For example, you might cite a paper like this: \cite{example2024}. All citations will be automatically added to your \texttt{cite.bib} file.

\section{Conclusion}
This template demonstrates basic \LaTeX{} features including mathematical formulas, tables, and lists. The bibliography system is now set up to work seamlessly with Bibby AI's citation management. Feel free to modify it for your needs!

\bibliographystyle{plain}
\bibliography{cite}

\end{document}\documentclass{article}
\usepackage[utf8]{inputenc}
\usepackage{booktabs}
\usepackage{amsmath}
\usepackage{natbib}

\title{My Document}
\author{Your Name}
\date{\today}

\begin{document}

\maketitle

\section{Introduction}
This is a sample document demonstrating various \LaTeX{} features including mathematical equations, tables, and formatting. You can cite papers using the citation search feature, and they will be automatically added to your \texttt{cite.bib} file.

\section{Mathematical Expressions}
Here are some example equations:

\subsection{Quadratic Formula}
The quadratic formula is given by:
\[x = \frac{-b \pm \sqrt{b^2 - 4ac}}{2a}\]

\subsection{Matrix Example}
A sample matrix:
\[
\begin{pmatrix}
a_{11} & a_{12} & a_{13} \\
a_{21} & a_{22} & a_{23} \\
a_{31} & a_{32} & a_{33}
\end{pmatrix}
\]

\section{Tables}
Below is a sample table showing data:

\begin{table}[h]
\centering
\begin{tabular}{@{}lrrr@{}}
\toprule
\textbf{Category} & \textbf{Value 1} & \textbf{Value 2} & \textbf{Total} \\
\midrule
Item A & 10.5 & 20.3 & 30.8 \\
Item B & 15.2 & 25.1 & 40.3 \\
Item C & 22.1 & 12.7 & 34.8 \\
\bottomrule
\end{tabular}
\caption{Sample Data Table}
\label{tab:sample}
\end{table}

\section{Lists}
\subsection{Bullet Points}
\begin{itemize}
  \item First item
  \item Second item
  \item Third item with sub-items:
    \begin{itemize}
      \item Sub-item 1
      \item Sub-item 2
    \end{itemize}
\end{itemize}

\subsection{Numbered List}
\begin{enumerate}
  \item First step
  \item Second step
  \item Third step
\end{enumerate}

\section{Citations and References}
When you use the citation search feature, references will be automatically managed. For example, you might cite a paper like this: \cite{example2024}. All citations will be automatically added to your \texttt{cite.bib} file.

\section{Conclusion}
This template demonstrates basic \LaTeX{} features including mathematical formulas, tables, and lists. The bibliography system is now set up to work seamlessly with Bibby AI's citation management. Feel free to modify it for your needs!

\bibliographystyle{plain}
\bibliography{cite}

\end{document}\documentclass{article}
\usepackage[utf8]{inputenc}
\usepackage{booktabs}
\usepackage{amsmath}
\usepackage{natbib}

\title{My Document}
\author{Your Name}
\date{\today}

\begin{document}

\maketitle

\section{Introduction}
This is a sample document demonstrating various \LaTeX{} features including mathematical equations, tables, and formatting. You can cite papers using the citation search feature, and they will be automatically added to your \texttt{cite.bib} file.

\section{Mathematical Expressions}
Here are some example equations:

\subsection{Quadratic Formula}
The quadratic formula is given by:
\[x = \frac{-b \pm \sqrt{b^2 - 4ac}}{2a}\]

\subsection{Matrix Example}
A sample matrix:
\[
\begin{pmatrix}
a_{11} & a_{12} & a_{13} \\
a_{21} & a_{22} & a_{23} \\
a_{31} & a_{32} & a_{33}
\end{pmatrix}
\]

\section{Tables}
Below is a sample table showing data:

\begin{table}[h]
\centering
\begin{tabular}{@{}lrrr@{}}
\toprule
\textbf{Category} & \textbf{Value 1} & \textbf{Value 2} & \textbf{Total} \\
\midrule
Item A & 10.5 & 20.3 & 30.8 \\
Item B & 15.2 & 25.1 & 40.3 \\
Item C & 22.1 & 12.7 & 34.8 \\
\bottomrule
\end{tabular}
\caption{Sample Data Table}
\label{tab:sample}
\end{table}

\section{Lists}
\subsection{Bullet Points}
\begin{itemize}
  \item First item
  \item Second item
  \item Third item with sub-items:
    \begin{itemize}
      \item Sub-item 1
      \item Sub-item 2
    \end{itemize}
\end{itemize}

\subsection{Numbered List}
\begin{enumerate}
  \item First step
  \item Second step
  \item Third step
\end{enumerate}

\section{Citations and References}
When you use the citation search feature, references will be automatically managed. For example, you might cite a paper like this: \cite{example2024}. All citations will be automatically added to your \texttt{cite.bib} file.

\section{Conclusion}
This template demonstrates basic \LaTeX{} features including mathematical formulas, tables, and lists. The bibliography system is now set up to work seamlessly with Bibby AI's citation management. Feel free to modify it for your needs!

\bibliographystyle{plain}
\bibliography{cite}

\end{document}\documentclass{article}
\usepackage[utf8]{inputenc}
\usepackage{booktabs}
\usepackage{amsmath}
\usepackage{natbib}

\title{My Document}
\author{Your Name}
\date{\today}

\begin{document}

\maketitle

\section{Introduction}
This is a sample document demonstrating various \LaTeX{} features including mathematical equations, tables, and formatting. You can cite papers using the citation search feature, and they will be automatically added to your \texttt{cite.bib} file.

\section{Mathematical Expressions}
Here are some example equations:

\subsection{Quadratic Formula}
The quadratic formula is given by:
\[x = \frac{-b \pm \sqrt{b^2 - 4ac}}{2a}\]

\subsection{Matrix Example}
A sample matrix:
\[
\begin{pmatrix}
a_{11} & a_{12} & a_{13} \\
a_{21} & a_{22} & a_{23} \\
a_{31} & a_{32} & a_{33}
\end{pmatrix}
\]

\section{Tables}
Below is a sample table showing data:

\begin{table}[h]
\centering
\begin{tabular}{@{}lrrr@{}}
\toprule
\textbf{Category} & \textbf{Value 1} & \textbf{Value 2} & \textbf{Total} \\
\midrule
Item A & 10.5 & 20.3 & 30.8 \\
Item B & 15.2 & 25.1 & 40.3 \\
Item C & 22.1 & 12.7 & 34.8 \\
\bottomrule
\end{tabular}
\caption{Sample Data Table}
\label{tab:sample}
\end{table}

\section{Lists}
\subsection{Bullet Points}
\begin{itemize}
  \item First item
  \item Second item
  \item Third item with sub-items:
    \begin{itemize}
      \item Sub-item 1
      \item Sub-item 2
    \end{itemize}
\end{itemize}

\subsection{Numbered List}
\begin{enumerate}
  \item First step
  \item Second step
  \item Third step
\end{enumerate}

\section{Citations and References}
When you use the citation search feature, references will be automatically managed. For example, you might cite a paper like this: \cite{example2024}. All citations will be automatically added to your \texttt{cite.bib} file.

\section{Conclusion}
This template demonstrates basic \LaTeX{} features including mathematical formulas, tables, and lists. The bibliography system is now set up to work seamlessly with Bibby AI's citation management. Feel free to modify it for your needs!

\bibliographystyle{plain}
\bibliography{cite}

\end{document}\documentclass{article}
\usepackage[utf8]{inputenc}
\usepackage{booktabs}
\usepackage{amsmath}
\usepackage{natbib}

\title{My Document}
\author{Your Name}
\date{\today}

\begin{document}

\maketitle

\section{Introduction}
This is a sample document demonstrating various \LaTeX{} features including mathematical equations, tables, and formatting. You can cite papers using the citation search feature, and they will be automatically added to your \texttt{cite.bib} file.

\section{Mathematical Expressions}
Here are some example equations:

\subsection{Quadratic Formula}
The quadratic formula is given by:
\[x = \frac{-b \pm \sqrt{b^2 - 4ac}}{2a}\]

\subsection{Matrix Example}
A sample matrix:
\[
\begin{pmatrix}
a_{11} & a_{12} & a_{13} \\
a_{21} & a_{22} & a_{23} \\
a_{31} & a_{32} & a_{33}
\end{pmatrix}
\]

\section{Tables}
Below is a sample table showing data:

\begin{table}[h]
\centering
\begin{tabular}{@{}lrrr@{}}
\toprule
\textbf{Category} & \textbf{Value 1} & \textbf{Value 2} & \textbf{Total} \\
\midrule
Item A & 10.5 & 20.3 & 30.8 \\
Item B & 15.2 & 25.1 & 40.3 \\
Item C & 22.1 & 12.7 & 34.8 \\
\bottomrule
\end{tabular}
\caption{Sample Data Table}
\label{tab:sample}
\end{table}

\section{Lists}
\subsection{Bullet Points}
\begin{itemize}
  \item First item
  \item Second item
  \item Third item with sub-items:
    \begin{itemize}
      \item Sub-item 1
      \item Sub-item 2
    \end{itemize}
\end{itemize}

\subsection{Numbered List}
\begin{enumerate}
  \item First step
  \item Second step
  \item Third step
\end{enumerate}

\section{Citations and References}
When you use the citation search feature, references will be automatically managed. For example, you might cite a paper like this: \cite{example2024}. All citations will be automatically added to your \texttt{cite.bib} file.

\section{Conclusion}
This template demonstrates basic \LaTeX{} features including mathematical formulas, tables, and lists. The bibliography system is now set up to work seamlessly with Bibby AI's citation management. Feel free to modify it for your needs!

\bibliographystyle{plain}
\bibliography{cite}

\end{document}\documentclass{article}
\usepackage[utf8]{inputenc}
\usepackage{booktabs}
\usepackage{amsmath}
\usepackage{natbib}

\title{My Document}
\author{Your Name}
\date{\today}

\begin{document}

\maketitle

\section{Introduction}
This is a sample document demonstrating various \LaTeX{} features including mathematical equations, tables, and formatting. You can cite papers using the citation search feature, and they will be automatically added to your \texttt{cite.bib} file.

\section{Mathematical Expressions}
Here are some example equations:

\subsection{Quadratic Formula}
The quadratic formula is given by:
\[x = \frac{-b \pm \sqrt{b^2 - 4ac}}{2a}\]

\subsection{Matrix Example}
A sample matrix:
\[
\begin{pmatrix}
a_{11} & a_{12} & a_{13} \\
a_{21} & a_{22} & a_{23} \\
a_{31} & a_{32} & a_{33}
\end{pmatrix}
\]

\section{Tables}
Below is a sample table showing data:

\begin{table}[h]
\centering
\begin{tabular}{@{}lrrr@{}}
\toprule
\textbf{Category} & \textbf{Value 1} & \textbf{Value 2} & \textbf{Total} \\
\midrule
Item A & 10.5 & 20.3 & 30.8 \\
Item B & 15.2 & 25.1 & 40.3 \\
Item C & 22.1 & 12.7 & 34.8 \\
\bottomrule
\end{tabular}
\caption{Sample Data Table}
\label{tab:sample}
\end{table}

\section{Lists}
\subsection{Bullet Points}
\begin{itemize}
  \item First item
  \item Second item
  \item Third item with sub-items:
    \begin{itemize}
      \item Sub-item 1
      \item Sub-item 2
    \end{itemize}
\end{itemize}

\subsection{Numbered List}
\begin{enumerate}
  \item First step
  \item Second step
  \item Third step
\end{enumerate}

\section{Citations and References}
When you use the citation search feature, references will be automatically managed. For example, you might cite a paper like this: \cite{example2024}. All citations will be automatically added to your \texttt{cite.bib} file.

\section{Conclusion}
This template demonstrates basic \LaTeX{} features including mathematical formulas, tables, and lists. The bibliography system is now set up to work seamlessly with Bibby AI's citation management. Feel free to modify it for your needs!

\bibliographystyle{plain}
\bibliography{cite}

\end{document}\documentclass{article}
\usepackage[utf8]{inputenc}
\usepackage{booktabs}
\usepackage{amsmath}
\usepackage{natbib}

\title{My Document}
\author{Your Name}
\date{\today}

\begin{document}

\maketitle

\section{Introduction}
This is a sample document demonstrating various \LaTeX{} features including mathematical equations, tables, and formatting. You can cite papers using the citation search feature, and they will be automatically added to your \texttt{cite.bib} file.

\section{Mathematical Expressions}
Here are some example equations:

\subsection{Quadratic Formula}
The quadratic formula is given by:
\[x = \frac{-b \pm \sqrt{b^2 - 4ac}}{2a}\]

\subsection{Matrix Example}
A sample matrix:
\[
\begin{pmatrix}
a_{11} & a_{12} & a_{13} \\
a_{21} & a_{22} & a_{23} \\
a_{31} & a_{32} & a_{33}
\end{pmatrix}
\]

\section{Tables}
Below is a sample table showing data:

\begin{table}[h]
\centering
\begin{tabular}{@{}lrrr@{}}
\toprule
\textbf{Category} & \textbf{Value 1} & \textbf{Value 2} & \textbf{Total} \\
\midrule
Item A & 10.5 & 20.3 & 30.8 \\
Item B & 15.2 & 25.1 & 40.3 \\
Item C & 22.1 & 12.7 & 34.8 \\
\bottomrule
\end{tabular}
\caption{Sample Data Table}
\label{tab:sample}
\end{table}

\section{Lists}
\subsection{Bullet Points}
\begin{itemize}
  \item First item
  \item Second item
  \item Third item with sub-items:
    \begin{itemize}
      \item Sub-item 1
      \item Sub-item 2
    \end{itemize}
\end{itemize}

\subsection{Numbered List}
\begin{enumerate}
  \item First step
  \item Second step
  \item Third step
\end{enumerate}

\section{Citations and References}
When you use the citation search feature, references will be automatically managed. For example, you might cite a paper like this: \cite{example2024}. All citations will be automatically added to your \texttt{cite.bib} file.

\section{Conclusion}
This template demonstrates basic \LaTeX{} features including mathematical formulas, tables, and lists. The bibliography system is now set up to work seamlessly with Bibby AI's citation management. Feel free to modify it for your needs!

\bibliographystyle{plain}
\bibliography{cite}

\end{document}\documentclass{article}
\usepackage[utf8]{inputenc}
\usepackage{booktabs}
\usepackage{amsmath}
\usepackage{natbib}

\title{My Document}
\author{Your Name}
\date{\today}

\begin{document}

\maketitle

\section{Introduction}
This is a sample document demonstrating various \LaTeX{} features including mathematical equations, tables, and formatting. You can cite papers using the citation search feature, and they will be automatically added to your \texttt{cite.bib} file.

\section{Mathematical Expressions}
Here are some example equations:

\subsection{Quadratic Formula}
The quadratic formula is given by:
\[x = \frac{-b \pm \sqrt{b^2 - 4ac}}{2a}\]

\subsection{Matrix Example}
A sample matrix:
\[
\begin{pmatrix}
a_{11} & a_{12} & a_{13} \\
a_{21} & a_{22} & a_{23} \\
a_{31} & a_{32} & a_{33}
\end{pmatrix}
\]

\section{Tables}
Below is a sample table showing data:

\begin{table}[h]
\centering
\begin{tabular}{@{}lrrr@{}}
\toprule
\textbf{Category} & \textbf{Value 1} & \textbf{Value 2} & \textbf{Total} \\
\midrule
Item A & 10.5 & 20.3 & 30.8 \\
Item B & 15.2 & 25.1 & 40.3 \\
Item C & 22.1 & 12.7 & 34.8 \\
\bottomrule
\end{tabular}
\caption{Sample Data Table}
\label{tab:sample}
\end{table}

\section{Lists}
\subsection{Bullet Points}
\begin{itemize}
  \item First item
  \item Second item
  \item Third item with sub-items:
    \begin{itemize}
      \item Sub-item 1
      \item Sub-item 2
    \end{itemize}
\end{itemize}

\subsection{Numbered List}
\begin{enumerate}
  \item First step
  \item Second step
  \item Third step
\end{enumerate}

\section{Citations and References}
When you use the citation search feature, references will be automatically managed. For example, you might cite a paper like this: \cite{example2024}. All citations will be automatically added to your \texttt{cite.bib} file.

\section{Conclusion}
This template demonstrates basic \LaTeX{} features including mathematical formulas, tables, and lists. The bibliography system is now set up to work seamlessly with Bibby AI's citation management. Feel free to modify it for your needs!

\bibliographystyle{plain}
\bibliography{cite}

\end{document}\documentclass{article}
\usepackage[utf8]{inputenc}
\usepackage{booktabs}
\usepackage{amsmath}
\usepackage{natbib}

\title{My Document}
\author{Your Name}
\date{\today}

\begin{document}

\maketitle

\section{Introduction}
This is a sample document demonstrating various \LaTeX{} features including mathematical equations, tables, and formatting. You can cite papers using the citation search feature, and they will be automatically added to your \texttt{cite.bib} file.

\section{Mathematical Expressions}
Here are some example equations:

\subsection{Quadratic Formula}
The quadratic formula is given by:
\[x = \frac{-b \pm \sqrt{b^2 - 4ac}}{2a}\]

\subsection{Matrix Example}
A sample matrix:
\[
\begin{pmatrix}
a_{11} & a_{12} & a_{13} \\
a_{21} & a_{22} & a_{23} \\
a_{31} & a_{32} & a_{33}
\end{pmatrix}
\]

\section{Tables}
Below is a sample table showing data:

\begin{table}[h]
\centering
\begin{tabular}{@{}lrrr@{}}
\toprule
\textbf{Category} & \textbf{Value 1} & \textbf{Value 2} & \textbf{Total} \\
\midrule
Item A & 10.5 & 20.3 & 30.8 \\
Item B & 15.2 & 25.1 & 40.3 \\
Item C & 22.1 & 12.7 & 34.8 \\
\bottomrule
\end{tabular}
\caption{Sample Data Table}
\label{tab:sample}
\end{table}

\section{Lists}
\subsection{Bullet Points}
\begin{itemize}
  \item First item
  \item Second item
  \item Third item with sub-items:
    \begin{itemize}
      \item Sub-item 1
      \item Sub-item 2
    \end{itemize}
\end{itemize}

\subsection{Numbered List}
\begin{enumerate}
  \item First step
  \item Second step
  \item Third step
\end{enumerate}

\section{Citations and References}
When you use the citation search feature, references will be automatically managed. For example, you might cite a paper like this: \cite{example2024}. All citations will be automatically added to your \texttt{cite.bib} file.

\section{Conclusion}
This template demonstrates basic \LaTeX{} features including mathematical formulas, tables, and lists. The bibliography system is now set up to work seamlessly with Bibby AI's citation management. Feel free to modify it for your needs!

\bibliographystyle{plain}
\bibliography{cite}

\end{document}\documentclass{article}
\usepackage[utf8]{inputenc}
\usepackage{booktabs}
\usepackage{amsmath}
\usepackage{natbib}

\title{My Document}
\author{Your Name}
\date{\today}

\begin{document}

\maketitle

\section{Introduction}
This is a sample document demonstrating various \LaTeX{} features including mathematical equations, tables, and formatting. You can cite papers using the citation search feature, and they will be automatically added to your \texttt{cite.bib} file.

\section{Mathematical Expressions}
Here are some example equations:

\subsection{Quadratic Formula}
The quadratic formula is given by:
\[x = \frac{-b \pm \sqrt{b^2 - 4ac}}{2a}\]

\subsection{Matrix Example}
A sample matrix:
\[
\begin{pmatrix}
a_{11} & a_{12} & a_{13} \\
a_{21} & a_{22} & a_{23} \\
a_{31} & a_{32} & a_{33}
\end{pmatrix}
\]

\section{Tables}
Below is a sample table showing data:

\begin{table}[h]
\centering
\begin{tabular}{@{}lrrr@{}}
\toprule
\textbf{Category} & \textbf{Value 1} & \textbf{Value 2} & \textbf{Total} \\
\midrule
Item A & 10.5 & 20.3 & 30.8 \\
Item B & 15.2 & 25.1 & 40.3 \\
Item C & 22.1 & 12.7 & 34.8 \\
\bottomrule
\end{tabular}
\caption{Sample Data Table}
\label{tab:sample}
\end{table}

\section{Lists}
\subsection{Bullet Points}
\begin{itemize}
  \item First item
  \item Second item
  \item Third item with sub-items:
    \begin{itemize}
      \item Sub-item 1
      \item Sub-item 2
    \end{itemize}
\end{itemize}

\subsection{Numbered List}
\begin{enumerate}
  \item First step
  \item Second step
  \item Third step
\end{enumerate}

\section{Citations and References}
When you use the citation search feature, references will be automatically managed. For example, you might cite a paper like this: \cite{example2024}. All citations will be automatically added to your \texttt{cite.bib} file.

\section{Conclusion}
This template demonstrates basic \LaTeX{} features including mathematical formulas, tables, and lists. The bibliography system is now set up to work seamlessly with Bibby AI's citation management. Feel free to modify it for your needs!

\bibliographystyle{plain}
\bibliography{cite}

\end{document}\documentclass{article}
\usepackage[utf8]{inputenc}
\usepackage{booktabs}
\usepackage{amsmath}
\usepackage{natbib}

\title{My Document}
\author{Your Name}
\date{\today}

\begin{document}

\maketitle

\section{Introduction}
This is a sample document demonstrating various \LaTeX{} features including mathematical equations, tables, and formatting. You can cite papers using the citation search feature, and they will be automatically added to your \texttt{cite.bib} file.

\section{Mathematical Expressions}
Here are some example equations:

\subsection{Quadratic Formula}
The quadratic formula is given by:
\[x = \frac{-b \pm \sqrt{b^2 - 4ac}}{2a}\]

\subsection{Matrix Example}
A sample matrix:
\[
\begin{pmatrix}
a_{11} & a_{12} & a_{13} \\
a_{21} & a_{22} & a_{23} \\
a_{31} & a_{32} & a_{33}
\end{pmatrix}
\]

\section{Tables}
Below is a sample table showing data:

\begin{table}[h]
\centering
\begin{tabular}{@{}lrrr@{}}
\toprule
\textbf{Category} & \textbf{Value 1} & \textbf{Value 2} & \textbf{Total} \\
\midrule
Item A & 10.5 & 20.3 & 30.8 \\
Item B & 15.2 & 25.1 & 40.3 \\
Item C & 22.1 & 12.7 & 34.8 \\
\bottomrule
\end{tabular}
\caption{Sample Data Table}
\label{tab:sample}
\end{table}

\section{Lists}
\subsection{Bullet Points}
\begin{itemize}
  \item First item
  \item Second item
  \item Third item with sub-items:
    \begin{itemize}
      \item Sub-item 1
      \item Sub-item 2
    \end{itemize}
\end{itemize}

\subsection{Numbered List}
\begin{enumerate}
  \item First step
  \item Second step
  \item Third step
\end{enumerate}

\section{Citations and References}
When you use the citation search feature, references will be automatically managed. For example, you might cite a paper like this: \cite{example2024}. All citations will be automatically added to your \texttt{cite.bib} file.

\section{Conclusion}
This template demonstrates basic \LaTeX{} features including mathematical formulas, tables, and lists. The bibliography system is now set up to work seamlessly with Bibby AI's citation management. Feel free to modify it for your needs!

\bibliographystyle{plain}
\bibliography{cite}

\end{document}\documentclass{article}
\usepackage[utf8]{inputenc}
\usepackage{booktabs}
\usepackage{amsmath}
\usepackage{natbib}

\title{My Document}
\author{Your Name}
\date{\today}

\begin{document}

\maketitle

\section{Introduction}
This is a sample document demonstrating various \LaTeX{} features including mathematical equations, tables, and formatting. You can cite papers using the citation search feature, and they will be automatically added to your \texttt{cite.bib} file.

\section{Mathematical Expressions}
Here are some example equations:

\subsection{Quadratic Formula}
The quadratic formula is given by:
\[x = \frac{-b \pm \sqrt{b^2 - 4ac}}{2a}\]

\subsection{Matrix Example}
A sample matrix:
\[
\begin{pmatrix}
a_{11} & a_{12} & a_{13} \\
a_{21} & a_{22} & a_{23} \\
a_{31} & a_{32} & a_{33}
\end{pmatrix}
\]

\section{Tables}
Below is a sample table showing data:

\begin{table}[h]
\centering
\begin{tabular}{@{}lrrr@{}}
\toprule
\textbf{Category} & \textbf{Value 1} & \textbf{Value 2} & \textbf{Total} \\
\midrule
Item A & 10.5 & 20.3 & 30.8 \\
Item B & 15.2 & 25.1 & 40.3 \\
Item C & 22.1 & 12.7 & 34.8 \\
\bottomrule
\end{tabular}
\caption{Sample Data Table}
\label{tab:sample}
\end{table}

\section{Lists}
\subsection{Bullet Points}
\begin{itemize}
  \item First item
  \item Second item
  \item Third item with sub-items:
    \begin{itemize}
      \item Sub-item 1
      \item Sub-item 2
    \end{itemize}
\end{itemize}

\subsection{Numbered List}
\begin{enumerate}
  \item First step
  \item Second step
  \item Third step
\end{enumerate}

\section{Citations and References}
When you use the citation search feature, references will be automatically managed. For example, you might cite a paper like this: \cite{example2024}. All citations will be automatically added to your \texttt{cite.bib} file.

\section{Conclusion}
This template demonstrates basic \LaTeX{} features including mathematical formulas, tables, and lists. The bibliography system is now set up to work seamlessly with Bibby AI's citation management. Feel free to modify it for your needs!

\bibliographystyle{plain}
\bibliography{cite}

\end{document}\documentclass{article}
\usepackage[utf8]{inputenc}
\usepackage{booktabs}
\usepackage{amsmath}
\usepackage{natbib}

\title{My Document}
\author{Your Name}
\date{\today}

\begin{document}

\maketitle

\section{Introduction}
This is a sample document demonstrating various \LaTeX{} features including mathematical equations, tables, and formatting. You can cite papers using the citation search feature, and they will be automatically added to your \texttt{cite.bib} file.

\section{Mathematical Expressions}
Here are some example equations:

\subsection{Quadratic Formula}
The quadratic formula is given by:
\[x = \frac{-b \pm \sqrt{b^2 - 4ac}}{2a}\]

\subsection{Matrix Example}
A sample matrix:
\[
\begin{pmatrix}
a_{11} & a_{12} & a_{13} \\
a_{21} & a_{22} & a_{23} \\
a_{31} & a_{32} & a_{33}
\end{pmatrix}
\]

\section{Tables}
Below is a sample table showing data:

\begin{table}[h]
\centering
\begin{tabular}{@{}lrrr@{}}
\toprule
\textbf{Category} & \textbf{Value 1} & \textbf{Value 2} & \textbf{Total} \\
\midrule
Item A & 10.5 & 20.3 & 30.8 \\
Item B & 15.2 & 25.1 & 40.3 \\
Item C & 22.1 & 12.7 & 34.8 \\
\bottomrule
\end{tabular}
\caption{Sample Data Table}
\label{tab:sample}
\end{table}

\section{Lists}
\subsection{Bullet Points}
\begin{itemize}
  \item First item
  \item Second item
  \item Third item with sub-items:
    \begin{itemize}
      \item Sub-item 1
      \item Sub-item 2
    \end{itemize}
\end{itemize}

\subsection{Numbered List}
\begin{enumerate}
  \item First step
  \item Second step
  \item Third step
\end{enumerate}

\section{Citations and References}
When you use the citation search feature, references will be automatically managed. For example, you might cite a paper like this: \cite{example2024}. All citations will be automatically added to your \texttt{cite.bib} file.

\section{Conclusion}
This template demonstrates basic \LaTeX{} features including mathematical formulas, tables, and lists. The bibliography system is now set up to work seamlessly with Bibby AI's citation management. Feel free to modify it for your needs!

\bibliographystyle{plain}
\bibliography{cite}

\end{document}\documentclass{article}
\usepackage[utf8]{inputenc}
\usepackage{booktabs}
\usepackage{amsmath}
\usepackage{natbib}

\title{My Document}
\author{Your Name}
\date{\today}

\begin{document}

\maketitle

\section{Introduction}
This is a sample document demonstrating various \LaTeX{} features including mathematical equations, tables, and formatting. You can cite papers using the citation search feature, and they will be automatically added to your \texttt{cite.bib} file.

\section{Mathematical Expressions}
Here are some example equations:

\subsection{Quadratic Formula}
The quadratic formula is given by:
\[x = \frac{-b \pm \sqrt{b^2 - 4ac}}{2a}\]

\subsection{Matrix Example}
A sample matrix:
\[
\begin{pmatrix}
a_{11} & a_{12} & a_{13} \\
a_{21} & a_{22} & a_{23} \\
a_{31} & a_{32} & a_{33}
\end{pmatrix}
\]

\section{Tables}
Below is a sample table showing data:

\begin{table}[h]
\centering
\begin{tabular}{@{}lrrr@{}}
\toprule
\textbf{Category} & \textbf{Value 1} & \textbf{Value 2} & \textbf{Total} \\
\midrule
Item A & 10.5 & 20.3 & 30.8 \\
Item B & 15.2 & 25.1 & 40.3 \\
Item C & 22.1 & 12.7 & 34.8 \\
\bottomrule
\end{tabular}
\caption{Sample Data Table}
\label{tab:sample}
\end{table}

\section{Lists}
\subsection{Bullet Points}
\begin{itemize}
  \item First item
  \item Second item
  \item Third item with sub-items:
    \begin{itemize}
      \item Sub-item 1
      \item Sub-item 2
    \end{itemize}
\end{itemize}

\subsection{Numbered List}
\begin{enumerate}
  \item First step
  \item Second step
  \item Third step
\end{enumerate}

\section{Citations and References}
When you use the citation search feature, references will be automatically managed. For example, you might cite a paper like this: \cite{example2024}. All citations will be automatically added to your \texttt{cite.bib} file.

\section{Conclusion}
This template demonstrates basic \LaTeX{} features including mathematical formulas, tables, and lists. The bibliography system is now set up to work seamlessly with Bibby AI's citation management. Feel free to modify it for your needs!

\bibliographystyle{plain}
\bibliography{cite}

\end{document}\documentclass{article}
\usepackage[utf8]{inputenc}
\usepackage{booktabs}
\usepackage{amsmath}
\usepackage{natbib}

\title{My Document}
\author{Your Name}
\date{\today}

\begin{document}

\maketitle

\section{Introduction}
This is a sample document demonstrating various \LaTeX{} features including mathematical equations, tables, and formatting. You can cite papers using the citation search feature, and they will be automatically added to your \texttt{cite.bib} file.

\section{Mathematical Expressions}
Here are some example equations:

\subsection{Quadratic Formula}
The quadratic formula is given by:
\[x = \frac{-b \pm \sqrt{b^2 - 4ac}}{2a}\]

\subsection{Matrix Example}
A sample matrix:
\[
\begin{pmatrix}
a_{11} & a_{12} & a_{13} \\
a_{21} & a_{22} & a_{23} \\
a_{31} & a_{32} & a_{33}
\end{pmatrix}
\]

\section{Tables}
Below is a sample table showing data:

\begin{table}[h]
\centering
\begin{tabular}{@{}lrrr@{}}
\toprule
\textbf{Category} & \textbf{Value 1} & \textbf{Value 2} & \textbf{Total} \\
\midrule
Item A & 10.5 & 20.3 & 30.8 \\
Item B & 15.2 & 25.1 & 40.3 \\
Item C & 22.1 & 12.7 & 34.8 \\
\bottomrule
\end{tabular}
\caption{Sample Data Table}
\label{tab:sample}
\end{table}

\section{Lists}
\subsection{Bullet Points}
\begin{itemize}
  \item First item
  \item Second item
  \item Third item with sub-items:
    \begin{itemize}
      \item Sub-item 1
      \item Sub-item 2
    \end{itemize}
\end{itemize}

\subsection{Numbered List}
\begin{enumerate}
  \item First step
  \item Second step
  \item Third step
\end{enumerate}

\section{Citations and References}
When you use the citation search feature, references will be automatically managed. For example, you might cite a paper like this: \cite{example2024}. All citations will be automatically added to your \texttt{cite.bib} file.

\section{Conclusion}
This template demonstrates basic \LaTeX{} features including mathematical formulas, tables, and lists. The bibliography system is now set up to work seamlessly with Bibby AI's citation management. Feel free to modify it for your needs!

\bibliographystyle{plain}
\bibliography{cite}

\end{document}\documentclass{article}
\usepackage[utf8]{inputenc}
\usepackage{booktabs}
\usepackage{amsmath}
\usepackage{natbib}

\title{My Document}
\author{Your Name}
\date{\today}

\begin{document}

\maketitle

\section{Introduction}
This is a sample document demonstrating various \LaTeX{} features including mathematical equations, tables, and formatting. You can cite papers using the citation search feature, and they will be automatically added to your \texttt{cite.bib} file.

\section{Mathematical Expressions}
Here are some example equations:

\subsection{Quadratic Formula}
The quadratic formula is given by:
\[x = \frac{-b \pm \sqrt{b^2 - 4ac}}{2a}\]

\subsection{Matrix Example}
A sample matrix:
\[
\begin{pmatrix}
a_{11} & a_{12} & a_{13} \\
a_{21} & a_{22} & a_{23} \\
a_{31} & a_{32} & a_{33}
\end{pmatrix}
\]

\section{Tables}
Below is a sample table showing data:

\begin{table}[h]
\centering
\begin{tabular}{@{}lrrr@{}}
\toprule
\textbf{Category} & \textbf{Value 1} & \textbf{Value 2} & \textbf{Total} \\
\midrule
Item A & 10.5 & 20.3 & 30.8 \\
Item B & 15.2 & 25.1 & 40.3 \\
Item C & 22.1 & 12.7 & 34.8 \\
\bottomrule
\end{tabular}
\caption{Sample Data Table}
\label{tab:sample}
\end{table}

\section{Lists}
\subsection{Bullet Points}
\begin{itemize}
  \item First item
  \item Second item
  \item Third item with sub-items:
    \begin{itemize}
      \item Sub-item 1
      \item Sub-item 2
    \end{itemize}
\end{itemize}

\subsection{Numbered List}
\begin{enumerate}
  \item First step
  \item Second step
  \item Third step
\end{enumerate}

\section{Citations and References}
When you use the citation search feature, references will be automatically managed. For example, you might cite a paper like this: \cite{example2024}. All citations will be automatically added to your \texttt{cite.bib} file.

\section{Conclusion}
This template demonstrates basic \LaTeX{} features including mathematical formulas, tables, and lists. The bibliography system is now set up to work seamlessly with Bibby AI's citation management. Feel free to modify it for your needs!

\bibliographystyle{plain}
\bibliography{cite}

\end{document}\documentclass{article}
\usepackage[utf8]{inputenc}
\usepackage{booktabs}
\usepackage{amsmath}
\usepackage{natbib}

\title{My Document}
\author{Your Name}
\date{\today}

\begin{document}

\maketitle

\section{Introduction}
This is a sample document demonstrating various \LaTeX{} features including mathematical equations, tables, and formatting. You can cite papers using the citation search feature, and they will be automatically added to your \texttt{cite.bib} file.

\section{Mathematical Expressions}
Here are some example equations:

\subsection{Quadratic Formula}
The quadratic formula is given by:
\[x = \frac{-b \pm \sqrt{b^2 - 4ac}}{2a}\]

\subsection{Matrix Example}
A sample matrix:
\[
\begin{pmatrix}
a_{11} & a_{12} & a_{13} \\
a_{21} & a_{22} & a_{23} \\
a_{31} & a_{32} & a_{33}
\end{pmatrix}
\]

\section{Tables}
Below is a sample table showing data:

\begin{table}[h]
\centering
\begin{tabular}{@{}lrrr@{}}
\toprule
\textbf{Category} & \textbf{Value 1} & \textbf{Value 2} & \textbf{Total} \\
\midrule
Item A & 10.5 & 20.3 & 30.8 \\
Item B & 15.2 & 25.1 & 40.3 \\
Item C & 22.1 & 12.7 & 34.8 \\
\bottomrule
\end{tabular}
\caption{Sample Data Table}
\label{tab:sample}
\end{table}

\section{Lists}
\subsection{Bullet Points}
\begin{itemize}
  \item First item
  \item Second item
  \item Third item with sub-items:
    \begin{itemize}
      \item Sub-item 1
      \item Sub-item 2
    \end{itemize}
\end{itemize}

\subsection{Numbered List}
\begin{enumerate}
  \item First step
  \item Second step
  \item Third step
\end{enumerate}

\section{Citations and References}
When you use the citation search feature, references will be automatically managed. For example, you might cite a paper like this: \cite{example2024}. All citations will be automatically added to your \texttt{cite.bib} file.

\section{Conclusion}
This template demonstrates basic \LaTeX{} features including mathematical formulas, tables, and lists. The bibliography system is now set up to work seamlessly with Bibby AI's citation management. Feel free to modify it for your needs!

\bibliographystyle{plain}
\bibliography{cite}

\end{document}\documentclass{article}
\usepackage[utf8]{inputenc}
\usepackage{booktabs}
\usepackage{amsmath}
\usepackage{natbib}

\title{My Document}
\author{Your Name}
\date{\today}

\begin{document}

\maketitle

\section{Introduction}
This is a sample document demonstrating various \LaTeX{} features including mathematical equations, tables, and formatting. You can cite papers using the citation search feature, and they will be automatically added to your \texttt{cite.bib} file.

\section{Mathematical Expressions}
Here are some example equations:

\subsection{Quadratic Formula}
The quadratic formula is given by:
\[x = \frac{-b \pm \sqrt{b^2 - 4ac}}{2a}\]

\subsection{Matrix Example}
A sample matrix:
\[
\begin{pmatrix}
a_{11} & a_{12} & a_{13} \\
a_{21} & a_{22} & a_{23} \\
a_{31} & a_{32} & a_{33}
\end{pmatrix}
\]

\section{Tables}
Below is a sample table showing data:

\begin{table}[h]
\centering
\begin{tabular}{@{}lrrr@{}}
\toprule
\textbf{Category} & \textbf{Value 1} & \textbf{Value 2} & \textbf{Total} \\
\midrule
Item A & 10.5 & 20.3 & 30.8 \\
Item B & 15.2 & 25.1 & 40.3 \\
Item C & 22.1 & 12.7 & 34.8 \\
\bottomrule
\end{tabular}
\caption{Sample Data Table}
\label{tab:sample}
\end{table}

\section{Lists}
\subsection{Bullet Points}
\begin{itemize}
  \item First item
  \item Second item
  \item Third item with sub-items:
    \begin{itemize}
      \item Sub-item 1
      \item Sub-item 2
    \end{itemize}
\end{itemize}

\subsection{Numbered List}
\begin{enumerate}
  \item First step
  \item Second step
  \item Third step
\end{enumerate}

\section{Citations and References}
When you use the citation search feature, references will be automatically managed. For example, you might cite a paper like this: \cite{example2024}. All citations will be automatically added to your \texttt{cite.bib} file.

\section{Conclusion}
This template demonstrates basic \LaTeX{} features including mathematical formulas, tables, and lists. The bibliography system is now set up to work seamlessly with Bibby AI's citation management. Feel free to modify it for your needs!

\bibliographystyle{plain}
\bibliography{cite}

\end{document}\documentclass{article}
\usepackage[utf8]{inputenc}
\usepackage{booktabs}
\usepackage{amsmath}
\usepackage{natbib}

\title{My Document}
\author{Your Name}
\date{\today}

\begin{document}

\maketitle

\section{Introduction}
This is a sample document demonstrating various \LaTeX{} features including mathematical equations, tables, and formatting. You can cite papers using the citation search feature, and they will be automatically added to your \texttt{cite.bib} file.

\section{Mathematical Expressions}
Here are some example equations:

\subsection{Quadratic Formula}
The quadratic formula is given by:
\[x = \frac{-b \pm \sqrt{b^2 - 4ac}}{2a}\]

\subsection{Matrix Example}
A sample matrix:
\[
\begin{pmatrix}
a_{11} & a_{12} & a_{13} \\
a_{21} & a_{22} & a_{23} \\
a_{31} & a_{32} & a_{33}
\end{pmatrix}
\]

\section{Tables}
Below is a sample table showing data:

\begin{table}[h]
\centering
\begin{tabular}{@{}lrrr@{}}
\toprule
\textbf{Category} & \textbf{Value 1} & \textbf{Value 2} & \textbf{Total} \\
\midrule
Item A & 10.5 & 20.3 & 30.8 \\
Item B & 15.2 & 25.1 & 40.3 \\
Item C & 22.1 & 12.7 & 34.8 \\
\bottomrule
\end{tabular}
\caption{Sample Data Table}
\label{tab:sample}
\end{table}

\section{Lists}
\subsection{Bullet Points}
\begin{itemize}
  \item First item
  \item Second item
  \item Third item with sub-items:
    \begin{itemize}
      \item Sub-item 1
      \item Sub-item 2
    \end{itemize}
\end{itemize}

\subsection{Numbered List}
\begin{enumerate}
  \item First step
  \item Second step
  \item Third step
\end{enumerate}

\section{Citations and References}
When you use the citation search feature, references will be automatically managed. For example, you might cite a paper like this: \cite{example2024}. All citations will be automatically added to your \texttt{cite.bib} file.

\section{Conclusion}
This template demonstrates basic \LaTeX{} features including mathematical formulas, tables, and lists. The bibliography system is now set up to work seamlessly with Bibby AI's citation management. Feel free to modify it for your needs!

\bibliographystyle{plain}
\bibliography{cite}

\end{document}\documentclass{article}
\usepackage[utf8]{inputenc}
\usepackage{booktabs}
\usepackage{amsmath}
\usepackage{natbib}

\title{My Document}
\author{Your Name}
\date{\today}

\begin{document}

\maketitle

\section{Introduction}
This is a sample document demonstrating various \LaTeX{} features including mathematical equations, tables, and formatting. You can cite papers using the citation search feature, and they will be automatically added to your \texttt{cite.bib} file.

\section{Mathematical Expressions}
Here are some example equations:

\subsection{Quadratic Formula}
The quadratic formula is given by:
\[x = \frac{-b \pm \sqrt{b^2 - 4ac}}{2a}\]

\subsection{Matrix Example}
A sample matrix:
\[
\begin{pmatrix}
a_{11} & a_{12} & a_{13} \\
a_{21} & a_{22} & a_{23} \\
a_{31} & a_{32} & a_{33}
\end{pmatrix}
\]

\section{Tables}
Below is a sample table showing data:

\begin{table}[h]
\centering
\begin{tabular}{@{}lrrr@{}}
\toprule
\textbf{Category} & \textbf{Value 1} & \textbf{Value 2} & \textbf{Total} \\
\midrule
Item A & 10.5 & 20.3 & 30.8 \\
Item B & 15.2 & 25.1 & 40.3 \\
Item C & 22.1 & 12.7 & 34.8 \\
\bottomrule
\end{tabular}
\caption{Sample Data Table}
\label{tab:sample}
\end{table}

\section{Lists}
\subsection{Bullet Points}
\begin{itemize}
  \item First item
  \item Second item
  \item Third item with sub-items:
    \begin{itemize}
      \item Sub-item 1
      \item Sub-item 2
    \end{itemize}
\end{itemize}

\subsection{Numbered List}
\begin{enumerate}
  \item First step
  \item Second step
  \item Third step
\end{enumerate}

\section{Citations and References}
When you use the citation search feature, references will be automatically managed. For example, you might cite a paper like this: \cite{example2024}. All citations will be automatically added to your \texttt{cite.bib} file.

\section{Conclusion}
This template demonstrates basic \LaTeX{} features including mathematical formulas, tables, and lists. The bibliography system is now set up to work seamlessly with Bibby AI's citation management. Feel free to modify it for your needs!

\bibliographystyle{plain}
\bibliography{cite}

\end{document}\documentclass{article}
\usepackage[utf8]{inputenc}
\usepackage{booktabs}
\usepackage{amsmath}
\usepackage{natbib}

\title{My Document}
\author{Your Name}
\date{\today}

\begin{document}

\maketitle

\section{Introduction}
This is a sample document demonstrating various \LaTeX{} features including mathematical equations, tables, and formatting. You can cite papers using the citation search feature, and they will be automatically added to your \texttt{cite.bib} file.

\section{Mathematical Expressions}
Here are some example equations:

\subsection{Quadratic Formula}
The quadratic formula is given by:
\[x = \frac{-b \pm \sqrt{b^2 - 4ac}}{2a}\]

\subsection{Matrix Example}
A sample matrix:
\[
\begin{pmatrix}
a_{11} & a_{12} & a_{13} \\
a_{21} & a_{22} & a_{23} \\
a_{31} & a_{32} & a_{33}
\end{pmatrix}
\]

\section{Tables}
Below is a sample table showing data:

\begin{table}[h]
\centering
\begin{tabular}{@{}lrrr@{}}
\toprule
\textbf{Category} & \textbf{Value 1} & \textbf{Value 2} & \textbf{Total} \\
\midrule
Item A & 10.5 & 20.3 & 30.8 \\
Item B & 15.2 & 25.1 & 40.3 \\
Item C & 22.1 & 12.7 & 34.8 \\
\bottomrule
\end{tabular}
\caption{Sample Data Table}
\label{tab:sample}
\end{table}

\section{Lists}
\subsection{Bullet Points}
\begin{itemize}
  \item First item
  \item Second item
  \item Third item with sub-items:
    \begin{itemize}
      \item Sub-item 1
      \item Sub-item 2
    \end{itemize}
\end{itemize}

\subsection{Numbered List}
\begin{enumerate}
  \item First step
  \item Second step
  \item Third step
\end{enumerate}

\section{Citations and References}
When you use the citation search feature, references will be automatically managed. For example, you might cite a paper like this: \cite{example2024}. All citations will be automatically added to your \texttt{cite.bib} file.

\section{Conclusion}
This template demonstrates basic \LaTeX{} features including mathematical formulas, tables, and lists. The bibliography system is now set up to work seamlessly with Bibby AI's citation management. Feel free to modify it for your needs!

\bibliographystyle{plain}
\bibliography{cite}

\end{document}\documentclass{article}
\usepackage[utf8]{inputenc}
\usepackage{booktabs}
\usepackage{amsmath}
\usepackage{natbib}

\title{My Document}
\author{Your Name}
\date{\today}

\begin{document}

\maketitle

\section{Introduction}
This is a sample document demonstrating various \LaTeX{} features including mathematical equations, tables, and formatting. You can cite papers using the citation search feature, and they will be automatically added to your \texttt{cite.bib} file.

\section{Mathematical Expressions}
Here are some example equations:

\subsection{Quadratic Formula}
The quadratic formula is given by:
\[x = \frac{-b \pm \sqrt{b^2 - 4ac}}{2a}\]

\subsection{Matrix Example}
A sample matrix:
\[
\begin{pmatrix}
a_{11} & a_{12} & a_{13} \\
a_{21} & a_{22} & a_{23} \\
a_{31} & a_{32} & a_{33}
\end{pmatrix}
\]

\section{Tables}
Below is a sample table showing data:

\begin{table}[h]
\centering
\begin{tabular}{@{}lrrr@{}}
\toprule
\textbf{Category} & \textbf{Value 1} & \textbf{Value 2} & \textbf{Total} \\
\midrule
Item A & 10.5 & 20.3 & 30.8 \\
Item B & 15.2 & 25.1 & 40.3 \\
Item C & 22.1 & 12.7 & 34.8 \\
\bottomrule
\end{tabular}
\caption{Sample Data Table}
\label{tab:sample}
\end{table}

\section{Lists}
\subsection{Bullet Points}
\begin{itemize}
  \item First item
  \item Second item
  \item Third item with sub-items:
    \begin{itemize}
      \item Sub-item 1
      \item Sub-item 2
    \end{itemize}
\end{itemize}

\subsection{Numbered List}
\begin{enumerate}
  \item First step
  \item Second step
  \item Third step
\end{enumerate}

\section{Citations and References}
When you use the citation search feature, references will be automatically managed. For example, you might cite a paper like this: \cite{example2024}. All citations will be automatically added to your \texttt{cite.bib} file.

\section{Conclusion}
This template demonstrates basic \LaTeX{} features including mathematical formulas, tables, and lists. The bibliography system is now set up to work seamlessly with Bibby AI's citation management. Feel free to modify it for your needs!

\bibliographystyle{plain}
\bibliography{cite}

\end{document}\documentclass{article}
\usepackage[utf8]{inputenc}
\usepackage{booktabs}
\usepackage{amsmath}
\usepackage{natbib}

\title{My Document}
\author{Your Name}
\date{\today}

\begin{document}

\maketitle

\section{Introduction}
This is a sample document demonstrating various \LaTeX{} features including mathematical equations, tables, and formatting. You can cite papers using the citation search feature, and they will be automatically added to your \texttt{cite.bib} file.

\section{Mathematical Expressions}
Here are some example equations:

\subsection{Quadratic Formula}
The quadratic formula is given by:
\[x = \frac{-b \pm \sqrt{b^2 - 4ac}}{2a}\]

\subsection{Matrix Example}
A sample matrix:
\[
\begin{pmatrix}
a_{11} & a_{12} & a_{13} \\
a_{21} & a_{22} & a_{23} \\
a_{31} & a_{32} & a_{33}
\end{pmatrix}
\]

\section{Tables}
Below is a sample table showing data:

\begin{table}[h]
\centering
\begin{tabular}{@{}lrrr@{}}
\toprule
\textbf{Category} & \textbf{Value 1} & \textbf{Value 2} & \textbf{Total} \\
\midrule
Item A & 10.5 & 20.3 & 30.8 \\
Item B & 15.2 & 25.1 & 40.3 \\
Item C & 22.1 & 12.7 & 34.8 \\
\bottomrule
\end{tabular}
\caption{Sample Data Table}
\label{tab:sample}
\end{table}

\section{Lists}
\subsection{Bullet Points}
\begin{itemize}
  \item First item
  \item Second item
  \item Third item with sub-items:
    \begin{itemize}
      \item Sub-item 1
      \item Sub-item 2
    \end{itemize}
\end{itemize}

\subsection{Numbered List}
\begin{enumerate}
  \item First step
  \item Second step
  \item Third step
\end{enumerate}

\section{Citations and References}
When you use the citation search feature, references will be automatically managed. For example, you might cite a paper like this: \cite{example2024}. All citations will be automatically added to your \texttt{cite.bib} file.

\section{Conclusion}
This template demonstrates basic \LaTeX{} features including mathematical formulas, tables, and lists. The bibliography system is now set up to work seamlessly with Bibby AI's citation management. Feel free to modify it for your needs!

\bibliographystyle{plain}
\bibliography{cite}

\end{document}\documentclass{article}
\usepackage[utf8]{inputenc}
\usepackage{booktabs}
\usepackage{amsmath}
\usepackage{natbib}

\title{My Document}
\author{Your Name}
\date{\today}

\begin{document}

\maketitle

\section{Introduction}
This is a sample document demonstrating various \LaTeX{} features including mathematical equations, tables, and formatting. You can cite papers using the citation search feature, and they will be automatically added to your \texttt{cite.bib} file.

\section{Mathematical Expressions}
Here are some example equations:

\subsection{Quadratic Formula}
The quadratic formula is given by:
\[x = \frac{-b \pm \sqrt{b^2 - 4ac}}{2a}\]

\subsection{Matrix Example}
A sample matrix:
\[
\begin{pmatrix}
a_{11} & a_{12} & a_{13} \\
a_{21} & a_{22} & a_{23} \\
a_{31} & a_{32} & a_{33}
\end{pmatrix}
\]

\section{Tables}
Below is a sample table showing data:

\begin{table}[h]
\centering
\begin{tabular}{@{}lrrr@{}}
\toprule
\textbf{Category} & \textbf{Value 1} & \textbf{Value 2} & \textbf{Total} \\
\midrule
Item A & 10.5 & 20.3 & 30.8 \\
Item B & 15.2 & 25.1 & 40.3 \\
Item C & 22.1 & 12.7 & 34.8 \\
\bottomrule
\end{tabular}
\caption{Sample Data Table}
\label{tab:sample}
\end{table}

\section{Lists}
\subsection{Bullet Points}
\begin{itemize}
  \item First item
  \item Second item
  \item Third item with sub-items:
    \begin{itemize}
      \item Sub-item 1
      \item Sub-item 2
    \end{itemize}
\end{itemize}

\subsection{Numbered List}
\begin{enumerate}
  \item First step
  \item Second step
  \item Third step
\end{enumerate}

\section{Citations and References}
When you use the citation search feature, references will be automatically managed. For example, you might cite a paper like this: \cite{example2024}. All citations will be automatically added to your \texttt{cite.bib} file.

\section{Conclusion}
This template demonstrates basic \LaTeX{} features including mathematical formulas, tables, and lists. The bibliography system is now set up to work seamlessly with Bibby AI's citation management. Feel free to modify it for your needs!

\bibliographystyle{plain}
\bibliography{cite}

\end{document}\documentclass{article}
\usepackage[utf8]{inputenc}
\usepackage{booktabs}
\usepackage{amsmath}
\usepackage{natbib}

\title{My Document}
\author{Your Name}
\date{\today}

\begin{document}

\maketitle

\section{Introduction}
This is a sample document demonstrating various \LaTeX{} features including mathematical equations, tables, and formatting. You can cite papers using the citation search feature, and they will be automatically added to your \texttt{cite.bib} file.

\section{Mathematical Expressions}
Here are some example equations:

\subsection{Quadratic Formula}
The quadratic formula is given by:
\[x = \frac{-b \pm \sqrt{b^2 - 4ac}}{2a}\]

\subsection{Matrix Example}
A sample matrix:
\[
\begin{pmatrix}
a_{11} & a_{12} & a_{13} \\
a_{21} & a_{22} & a_{23} \\
a_{31} & a_{32} & a_{33}
\end{pmatrix}
\]

\section{Tables}
Below is a sample table showing data:

\begin{table}[h]
\centering
\begin{tabular}{@{}lrrr@{}}
\toprule
\textbf{Category} & \textbf{Value 1} & \textbf{Value 2} & \textbf{Total} \\
\midrule
Item A & 10.5 & 20.3 & 30.8 \\
Item B & 15.2 & 25.1 & 40.3 \\
Item C & 22.1 & 12.7 & 34.8 \\
\bottomrule
\end{tabular}
\caption{Sample Data Table}
\label{tab:sample}
\end{table}

\section{Lists}
\subsection{Bullet Points}
\begin{itemize}
  \item First item
  \item Second item
  \item Third item with sub-items:
    \begin{itemize}
      \item Sub-item 1
      \item Sub-item 2
    \end{itemize}
\end{itemize}

\subsection{Numbered List}
\begin{enumerate}
  \item First step
  \item Second step
  \item Third step
\end{enumerate}

\section{Citations and References}
When you use the citation search feature, references will be automatically managed. For example, you might cite a paper like this: \cite{example2024}. All citations will be automatically added to your \texttt{cite.bib} file.

\section{Conclusion}
This template demonstrates basic \LaTeX{} features including mathematical formulas, tables, and lists. The bibliography system is now set up to work seamlessly with Bibby AI's citation management. Feel free to modify it for your needs!

\bibliographystyle{plain}
\bibliography{cite}

\end{document}\documentclass{article}
\usepackage[utf8]{inputenc}
\usepackage{booktabs}
\usepackage{amsmath}
\usepackage{natbib}

\title{My Document}
\author{Your Name}
\date{\today}

\begin{document}

\maketitle

\section{Introduction}
This is a sample document demonstrating various \LaTeX{} features including mathematical equations, tables, and formatting. You can cite papers using the citation search feature, and they will be automatically added to your \texttt{cite.bib} file.

\section{Mathematical Expressions}
Here are some example equations:

\subsection{Quadratic Formula}
The quadratic formula is given by:
\[x = \frac{-b \pm \sqrt{b^2 - 4ac}}{2a}\]

\subsection{Matrix Example}
A sample matrix:
\[
\begin{pmatrix}
a_{11} & a_{12} & a_{13} \\
a_{21} & a_{22} & a_{23} \\
a_{31} & a_{32} & a_{33}
\end{pmatrix}
\]

\section{Tables}
Below is a sample table showing data:

\begin{table}[h]
\centering
\begin{tabular}{@{}lrrr@{}}
\toprule
\textbf{Category} & \textbf{Value 1} & \textbf{Value 2} & \textbf{Total} \\
\midrule
Item A & 10.5 & 20.3 & 30.8 \\
Item B & 15.2 & 25.1 & 40.3 \\
Item C & 22.1 & 12.7 & 34.8 \\
\bottomrule
\end{tabular}
\caption{Sample Data Table}
\label{tab:sample}
\end{table}

\section{Lists}
\subsection{Bullet Points}
\begin{itemize}
  \item First item
  \item Second item
  \item Third item with sub-items:
    \begin{itemize}
      \item Sub-item 1
      \item Sub-item 2
    \end{itemize}
\end{itemize}

\subsection{Numbered List}
\begin{enumerate}
  \item First step
  \item Second step
  \item Third step
\end{enumerate}

\section{Citations and References}
When you use the citation search feature, references will be automatically managed. For example, you might cite a paper like this: \cite{example2024}. All citations will be automatically added to your \texttt{cite.bib} file.

\section{Conclusion}
This template demonstrates basic \LaTeX{} features including mathematical formulas, tables, and lists. The bibliography system is now set up to work seamlessly with Bibby AI's citation management. Feel free to modify it for your needs!

\bibliographystyle{plain}
\bibliography{cite}

\end{document}\documentclass{article}
\usepackage[utf8]{inputenc}
\usepackage{booktabs}
\usepackage{amsmath}
\usepackage{natbib}

\title{My Document}
\author{Your Name}
\date{\today}

\begin{document}

\maketitle

\section{Introduction}
This is a sample document demonstrating various \LaTeX{} features including mathematical equations, tables, and formatting. You can cite papers using the citation search feature, and they will be automatically added to your \texttt{cite.bib} file.

\section{Mathematical Expressions}
Here are some example equations:

\subsection{Quadratic Formula}
The quadratic formula is given by:
\[x = \frac{-b \pm \sqrt{b^2 - 4ac}}{2a}\]

\subsection{Matrix Example}
A sample matrix:
\[
\begin{pmatrix}
a_{11} & a_{12} & a_{13} \\
a_{21} & a_{22} & a_{23} \\
a_{31} & a_{32} & a_{33}
\end{pmatrix}
\]

\section{Tables}
Below is a sample table showing data:

\begin{table}[h]
\centering
\begin{tabular}{@{}lrrr@{}}
\toprule
\textbf{Category} & \textbf{Value 1} & \textbf{Value 2} & \textbf{Total} \\
\midrule
Item A & 10.5 & 20.3 & 30.8 \\
Item B & 15.2 & 25.1 & 40.3 \\
Item C & 22.1 & 12.7 & 34.8 \\
\bottomrule
\end{tabular}
\caption{Sample Data Table}
\label{tab:sample}
\end{table}

\section{Lists}
\subsection{Bullet Points}
\begin{itemize}
  \item First item
  \item Second item
  \item Third item with sub-items:
    \begin{itemize}
      \item Sub-item 1
      \item Sub-item 2
    \end{itemize}
\end{itemize}

\subsection{Numbered List}
\begin{enumerate}
  \item First step
  \item Second step
  \item Third step
\end{enumerate}

\section{Citations and References}
When you use the citation search feature, references will be automatically managed. For example, you might cite a paper like this: \cite{example2024}. All citations will be automatically added to your \texttt{cite.bib} file.

\section{Conclusion}
This template demonstrates basic \LaTeX{} features including mathematical formulas, tables, and lists. The bibliography system is now set up to work seamlessly with Bibby AI's citation management. Feel free to modify it for your needs!

\bibliographystyle{plain}
\bibliography{cite}

\end{document}\documentclass{article}
\usepackage[utf8]{inputenc}
\usepackage{booktabs}
\usepackage{amsmath}
\usepackage{natbib}

\title{My Document}
\author{Your Name}
\date{\today}

\begin{document}

\maketitle

\section{Introduction}
This is a sample document demonstrating various \LaTeX{} features including mathematical equations, tables, and formatting. You can cite papers using the citation search feature, and they will be automatically added to your \texttt{cite.bib} file.

\section{Mathematical Expressions}
Here are some example equations:

\subsection{Quadratic Formula}
The quadratic formula is given by:
\[x = \frac{-b \pm \sqrt{b^2 - 4ac}}{2a}\]

\subsection{Matrix Example}
A sample matrix:
\[
\begin{pmatrix}
a_{11} & a_{12} & a_{13} \\
a_{21} & a_{22} & a_{23} \\
a_{31} & a_{32} & a_{33}
\end{pmatrix}
\]

\section{Tables}
Below is a sample table showing data:

\begin{table}[h]
\centering
\begin{tabular}{@{}lrrr@{}}
\toprule
\textbf{Category} & \textbf{Value 1} & \textbf{Value 2} & \textbf{Total} \\
\midrule
Item A & 10.5 & 20.3 & 30.8 \\
Item B & 15.2 & 25.1 & 40.3 \\
Item C & 22.1 & 12.7 & 34.8 \\
\bottomrule
\end{tabular}
\caption{Sample Data Table}
\label{tab:sample}
\end{table}

\section{Lists}
\subsection{Bullet Points}
\begin{itemize}
  \item First item
  \item Second item
  \item Third item with sub-items:
    \begin{itemize}
      \item Sub-item 1
      \item Sub-item 2
    \end{itemize}
\end{itemize}

\subsection{Numbered List}
\begin{enumerate}
  \item First step
  \item Second step
  \item Third step
\end{enumerate}

\section{Citations and References}
When you use the citation search feature, references will be automatically managed. For example, you might cite a paper like this: \cite{example2024}. All citations will be automatically added to your \texttt{cite.bib} file.

\section{Conclusion}
This template demonstrates basic \LaTeX{} features including mathematical formulas, tables, and lists. The bibliography system is now set up to work seamlessly with Bibby AI's citation management. Feel free to modify it for your needs!

\bibliographystyle{plain}
\bibliography{cite}

\end{document}\documentclass{article}
\usepackage[utf8]{inputenc}
\usepackage{booktabs}
\usepackage{amsmath}
\usepackage{natbib}

\title{My Document}
\author{Your Name}
\date{\today}

\begin{document}

\maketitle

\section{Introduction}
This is a sample document demonstrating various \LaTeX{} features including mathematical equations, tables, and formatting. You can cite papers using the citation search feature, and they will be automatically added to your \texttt{cite.bib} file.

\section{Mathematical Expressions}
Here are some example equations:

\subsection{Quadratic Formula}
The quadratic formula is given by:
\[x = \frac{-b \pm \sqrt{b^2 - 4ac}}{2a}\]

\subsection{Matrix Example}
A sample matrix:
\[
\begin{pmatrix}
a_{11} & a_{12} & a_{13} \\
a_{21} & a_{22} & a_{23} \\
a_{31} & a_{32} & a_{33}
\end{pmatrix}
\]

\section{Tables}
Below is a sample table showing data:

\begin{table}[h]
\centering
\begin{tabular}{@{}lrrr@{}}
\toprule
\textbf{Category} & \textbf{Value 1} & \textbf{Value 2} & \textbf{Total} \\
\midrule
Item A & 10.5 & 20.3 & 30.8 \\
Item B & 15.2 & 25.1 & 40.3 \\
Item C & 22.1 & 12.7 & 34.8 \\
\bottomrule
\end{tabular}
\caption{Sample Data Table}
\label{tab:sample}
\end{table}

\section{Lists}
\subsection{Bullet Points}
\begin{itemize}
  \item First item
  \item Second item
  \item Third item with sub-items:
    \begin{itemize}
      \item Sub-item 1
      \item Sub-item 2
    \end{itemize}
\end{itemize}

\subsection{Numbered List}
\begin{enumerate}
  \item First step
  \item Second step
  \item Third step
\end{enumerate}

\section{Citations and References}
When you use the citation search feature, references will be automatically managed. For example, you might cite a paper like this: \cite{example2024}. All citations will be automatically added to your \texttt{cite.bib} file.

\section{Conclusion}
This template demonstrates basic \LaTeX{} features including mathematical formulas, tables, and lists. The bibliography system is now set up to work seamlessly with Bibby AI's citation management. Feel free to modify it for your needs!

\bibliographystyle{plain}
\bibliography{cite}

\end{document}\documentclass{article}
\usepackage[utf8]{inputenc}
\usepackage{booktabs}
\usepackage{amsmath}
\usepackage{natbib}

\title{My Document}
\author{Your Name}
\date{\today}

\begin{document}

\maketitle

\section{Introduction}
This is a sample document demonstrating various \LaTeX{} features including mathematical equations, tables, and formatting. You can cite papers using the citation search feature, and they will be automatically added to your \texttt{cite.bib} file.

\section{Mathematical Expressions}
Here are some example equations:

\subsection{Quadratic Formula}
The quadratic formula is given by:
\[x = \frac{-b \pm \sqrt{b^2 - 4ac}}{2a}\]

\subsection{Matrix Example}
A sample matrix:
\[
\begin{pmatrix}
a_{11} & a_{12} & a_{13} \\
a_{21} & a_{22} & a_{23} \\
a_{31} & a_{32} & a_{33}
\end{pmatrix}
\]

\section{Tables}
Below is a sample table showing data:

\begin{table}[h]
\centering
\begin{tabular}{@{}lrrr@{}}
\toprule
\textbf{Category} & \textbf{Value 1} & \textbf{Value 2} & \textbf{Total} \\
\midrule
Item A & 10.5 & 20.3 & 30.8 \\
Item B & 15.2 & 25.1 & 40.3 \\
Item C & 22.1 & 12.7 & 34.8 \\
\bottomrule
\end{tabular}
\caption{Sample Data Table}
\label{tab:sample}
\end{table}

\section{Lists}
\subsection{Bullet Points}
\begin{itemize}
  \item First item
  \item Second item
  \item Third item with sub-items:
    \begin{itemize}
      \item Sub-item 1
      \item Sub-item 2
    \end{itemize}
\end{itemize}

\subsection{Numbered List}
\begin{enumerate}
  \item First step
  \item Second step
  \item Third step
\end{enumerate}

\section{Citations and References}
When you use the citation search feature, references will be automatically managed. For example, you might cite a paper like this: \cite{example2024}. All citations will be automatically added to your \texttt{cite.bib} file.

\section{Conclusion}
This template demonstrates basic \LaTeX{} features including mathematical formulas, tables, and lists. The bibliography system is now set up to work seamlessly with Bibby AI's citation management. Feel free to modify it for your needs!

\bibliographystyle{plain}
\bibliography{cite}

\end{document}\documentclass{article}
\usepackage[utf8]{inputenc}
\usepackage{booktabs}
\usepackage{amsmath}
\usepackage{natbib}

\title{My Document}
\author{Your Name}
\date{\today}

\begin{document}

\maketitle

\section{Introduction}
This is a sample document demonstrating various \LaTeX{} features including mathematical equations, tables, and formatting. You can cite papers using the citation search feature, and they will be automatically added to your \texttt{cite.bib} file.

\section{Mathematical Expressions}
Here are some example equations:

\subsection{Quadratic Formula}
The quadratic formula is given by:
\[x = \frac{-b \pm \sqrt{b^2 - 4ac}}{2a}\]

\subsection{Matrix Example}
A sample matrix:
\[
\begin{pmatrix}
a_{11} & a_{12} & a_{13} \\
a_{21} & a_{22} & a_{23} \\
a_{31} & a_{32} & a_{33}
\end{pmatrix}
\]

\section{Tables}
Below is a sample table showing data:

\begin{table}[h]
\centering
\begin{tabular}{@{}lrrr@{}}
\toprule
\textbf{Category} & \textbf{Value 1} & \textbf{Value 2} & \textbf{Total} \\
\midrule
Item A & 10.5 & 20.3 & 30.8 \\
Item B & 15.2 & 25.1 & 40.3 \\
Item C & 22.1 & 12.7 & 34.8 \\
\bottomrule
\end{tabular}
\caption{Sample Data Table}
\label{tab:sample}
\end{table}

\section{Lists}
\subsection{Bullet Points}
\begin{itemize}
  \item First item
  \item Second item
  \item Third item with sub-items:
    \begin{itemize}
      \item Sub-item 1
      \item Sub-item 2
    \end{itemize}
\end{itemize}

\subsection{Numbered List}
\begin{enumerate}
  \item First step
  \item Second step
  \item Third step
\end{enumerate}

\section{Citations and References}
When you use the citation search feature, references will be automatically managed. For example, you might cite a paper like this: \cite{example2024}. All citations will be automatically added to your \texttt{cite.bib} file.

\section{Conclusion}
This template demonstrates basic \LaTeX{} features including mathematical formulas, tables, and lists. The bibliography system is now set up to work seamlessly with Bibby AI's citation management. Feel free to modify it for your needs!

\bibliographystyle{plain}
\bibliography{cite}

\end{document}\documentclass{article}
\usepackage[utf8]{inputenc}
\usepackage{booktabs}
\usepackage{amsmath}
\usepackage{natbib}

\title{My Document}
\author{Your Name}
\date{\today}

\begin{document}

\maketitle

\section{Introduction}
This is a sample document demonstrating various \LaTeX{} features including mathematical equations, tables, and formatting. You can cite papers using the citation search feature, and they will be automatically added to your \texttt{cite.bib} file.

\section{Mathematical Expressions}
Here are some example equations:

\subsection{Quadratic Formula}
The quadratic formula is given by:
\[x = \frac{-b \pm \sqrt{b^2 - 4ac}}{2a}\]

\subsection{Matrix Example}
A sample matrix:
\[
\begin{pmatrix}
a_{11} & a_{12} & a_{13} \\
a_{21} & a_{22} & a_{23} \\
a_{31} & a_{32} & a_{33}
\end{pmatrix}
\]

\section{Tables}
Below is a sample table showing data:

\begin{table}[h]
\centering
\begin{tabular}{@{}lrrr@{}}
\toprule
\textbf{Category} & \textbf{Value 1} & \textbf{Value 2} & \textbf{Total} \\
\midrule
Item A & 10.5 & 20.3 & 30.8 \\
Item B & 15.2 & 25.1 & 40.3 \\
Item C & 22.1 & 12.7 & 34.8 \\
\bottomrule
\end{tabular}
\caption{Sample Data Table}
\label{tab:sample}
\end{table}

\section{Lists}
\subsection{Bullet Points}
\begin{itemize}
  \item First item
  \item Second item
  \item Third item with sub-items:
    \begin{itemize}
      \item Sub-item 1
      \item Sub-item 2
    \end{itemize}
\end{itemize}

\subsection{Numbered List}
\begin{enumerate}
  \item First step
  \item Second step
  \item Third step
\end{enumerate}

\section{Citations and References}
When you use the citation search feature, references will be automatically managed. For example, you might cite a paper like this: \cite{example2024}. All citations will be automatically added to your \texttt{cite.bib} file.

\section{Conclusion}
This template demonstrates basic \LaTeX{} features including mathematical formulas, tables, and lists. The bibliography system is now set up to work seamlessly with Bibby AI's citation management. Feel free to modify it for your needs!

\bibliographystyle{plain}
\bibliography{cite}

\end{document}\documentclass{article}
\usepackage[utf8]{inputenc}
\usepackage{booktabs}
\usepackage{amsmath}
\usepackage{natbib}

\title{My Document}
\author{Your Name}
\date{\today}

\begin{document}

\maketitle

\section{Introduction}
This is a sample document demonstrating various \LaTeX{} features including mathematical equations, tables, and formatting. You can cite papers using the citation search feature, and they will be automatically added to your \texttt{cite.bib} file.

\section{Mathematical Expressions}
Here are some example equations:

\subsection{Quadratic Formula}
The quadratic formula is given by:
\[x = \frac{-b \pm \sqrt{b^2 - 4ac}}{2a}\]

\subsection{Matrix Example}
A sample matrix:
\[
\begin{pmatrix}
a_{11} & a_{12} & a_{13} \\
a_{21} & a_{22} & a_{23} \\
a_{31} & a_{32} & a_{33}
\end{pmatrix}
\]

\section{Tables}
Below is a sample table showing data:

\begin{table}[h]
\centering
\begin{tabular}{@{}lrrr@{}}
\toprule
\textbf{Category} & \textbf{Value 1} & \textbf{Value 2} & \textbf{Total} \\
\midrule
Item A & 10.5 & 20.3 & 30.8 \\
Item B & 15.2 & 25.1 & 40.3 \\
Item C & 22.1 & 12.7 & 34.8 \\
\bottomrule
\end{tabular}
\caption{Sample Data Table}
\label{tab:sample}
\end{table}

\section{Lists}
\subsection{Bullet Points}
\begin{itemize}
  \item First item
  \item Second item
  \item Third item with sub-items:
    \begin{itemize}
      \item Sub-item 1
      \item Sub-item 2
    \end{itemize}
\end{itemize}

\subsection{Numbered List}
\begin{enumerate}
  \item First step
  \item Second step
  \item Third step
\end{enumerate}

\section{Citations and References}
When you use the citation search feature, references will be automatically managed. For example, you might cite a paper like this: \cite{example2024}. All citations will be automatically added to your \texttt{cite.bib} file.

\section{Conclusion}
This template demonstrates basic \LaTeX{} features including mathematical formulas, tables, and lists. The bibliography system is now set up to work seamlessly with Bibby AI's citation management. Feel free to modify it for your needs!

\bibliographystyle{plain}
\bibliography{cite}

\end{document}\documentclass{article}
\usepackage[utf8]{inputenc}
\usepackage{booktabs}
\usepackage{amsmath}
\usepackage{natbib}

\title{My Document}
\author{Your Name}
\date{\today}

\begin{document}

\maketitle

\section{Introduction}
This is a sample document demonstrating various \LaTeX{} features including mathematical equations, tables, and formatting. You can cite papers using the citation search feature, and they will be automatically added to your \texttt{cite.bib} file.

\section{Mathematical Expressions}
Here are some example equations:

\subsection{Quadratic Formula}
The quadratic formula is given by:
\[x = \frac{-b \pm \sqrt{b^2 - 4ac}}{2a}\]

\subsection{Matrix Example}
A sample matrix:
\[
\begin{pmatrix}
a_{11} & a_{12} & a_{13} \\
a_{21} & a_{22} & a_{23} \\
a_{31} & a_{32} & a_{33}
\end{pmatrix}
\]

\section{Tables}
Below is a sample table showing data:

\begin{table}[h]
\centering
\begin{tabular}{@{}lrrr@{}}
\toprule
\textbf{Category} & \textbf{Value 1} & \textbf{Value 2} & \textbf{Total} \\
\midrule
Item A & 10.5 & 20.3 & 30.8 \\
Item B & 15.2 & 25.1 & 40.3 \\
Item C & 22.1 & 12.7 & 34.8 \\
\bottomrule
\end{tabular}
\caption{Sample Data Table}
\label{tab:sample}
\end{table}

\section{Lists}
\subsection{Bullet Points}
\begin{itemize}
  \item First item
  \item Second item
  \item Third item with sub-items:
    \begin{itemize}
      \item Sub-item 1
      \item Sub-item 2
    \end{itemize}
\end{itemize}

\subsection{Numbered List}
\begin{enumerate}
  \item First step
  \item Second step
  \item Third step
\end{enumerate}

\section{Citations and References}
When you use the citation search feature, references will be automatically managed. For example, you might cite a paper like this: \cite{example2024}. All citations will be automatically added to your \texttt{cite.bib} file.

\section{Conclusion}
This template demonstrates basic \LaTeX{} features including mathematical formulas, tables, and lists. The bibliography system is now set up to work seamlessly with Bibby AI's citation management. Feel free to modify it for your needs!

\bibliographystyle{plain}
\bibliography{cite}

\end{document}\documentclass{article}
\usepackage[utf8]{inputenc}
\usepackage{booktabs}
\usepackage{amsmath}
\usepackage{natbib}

\title{My Document}
\author{Your Name}
\date{\today}

\begin{document}

\maketitle

\section{Introduction}
This is a sample document demonstrating various \LaTeX{} features including mathematical equations, tables, and formatting. You can cite papers using the citation search feature, and they will be automatically added to your \texttt{cite.bib} file.

\section{Mathematical Expressions}
Here are some example equations:

\subsection{Quadratic Formula}
The quadratic formula is given by:
\[x = \frac{-b \pm \sqrt{b^2 - 4ac}}{2a}\]

\subsection{Matrix Example}
A sample matrix:
\[
\begin{pmatrix}
a_{11} & a_{12} & a_{13} \\
a_{21} & a_{22} & a_{23} \\
a_{31} & a_{32} & a_{33}
\end{pmatrix}
\]

\section{Tables}
Below is a sample table showing data:

\begin{table}[h]
\centering
\begin{tabular}{@{}lrrr@{}}
\toprule
\textbf{Category} & \textbf{Value 1} & \textbf{Value 2} & \textbf{Total} \\
\midrule
Item A & 10.5 & 20.3 & 30.8 \\
Item B & 15.2 & 25.1 & 40.3 \\
Item C & 22.1 & 12.7 & 34.8 \\
\bottomrule
\end{tabular}
\caption{Sample Data Table}
\label{tab:sample}
\end{table}

\section{Lists}
\subsection{Bullet Points}
\begin{itemize}
  \item First item
  \item Second item
  \item Third item with sub-items:
    \begin{itemize}
      \item Sub-item 1
      \item Sub-item 2
    \end{itemize}
\end{itemize}

\subsection{Numbered List}
\begin{enumerate}
  \item First step
  \item Second step
  \item Third step
\end{enumerate}

\section{Citations and References}
When you use the citation search feature, references will be automatically managed. For example, you might cite a paper like this: \cite{example2024}. All citations will be automatically added to your \texttt{cite.bib} file.

\section{Conclusion}
This template demonstrates basic \LaTeX{} features including mathematical formulas, tables, and lists. The bibliography system is now set up to work seamlessly with Bibby AI's citation management. Feel free to modify it for your needs!

\bibliographystyle{plain}
\bibliography{cite}

\end{document}\documentclass{article}
\usepackage[utf8]{inputenc}
\usepackage{booktabs}
\usepackage{amsmath}
\usepackage{natbib}

\title{My Document}
\author{Your Name}
\date{\today}

\begin{document}

\maketitle

\section{Introduction}
This is a sample document demonstrating various \LaTeX{} features including mathematical equations, tables, and formatting. You can cite papers using the citation search feature, and they will be automatically added to your \texttt{cite.bib} file.

\section{Mathematical Expressions}
Here are some example equations:

\subsection{Quadratic Formula}
The quadratic formula is given by:
\[x = \frac{-b \pm \sqrt{b^2 - 4ac}}{2a}\]

\subsection{Matrix Example}
A sample matrix:
\[
\begin{pmatrix}
a_{11} & a_{12} & a_{13} \\
a_{21} & a_{22} & a_{23} \\
a_{31} & a_{32} & a_{33}
\end{pmatrix}
\]

\section{Tables}
Below is a sample table showing data:

\begin{table}[h]
\centering
\begin{tabular}{@{}lrrr@{}}
\toprule
\textbf{Category} & \textbf{Value 1} & \textbf{Value 2} & \textbf{Total} \\
\midrule
Item A & 10.5 & 20.3 & 30.8 \\
Item B & 15.2 & 25.1 & 40.3 \\
Item C & 22.1 & 12.7 & 34.8 \\
\bottomrule
\end{tabular}
\caption{Sample Data Table}
\label{tab:sample}
\end{table}

\section{Lists}
\subsection{Bullet Points}
\begin{itemize}
  \item First item
  \item Second item
  \item Third item with sub-items:
    \begin{itemize}
      \item Sub-item 1
      \item Sub-item 2
    \end{itemize}
\end{itemize}

\subsection{Numbered List}
\begin{enumerate}
  \item First step
  \item Second step
  \item Third step
\end{enumerate}

\section{Citations and References}
When you use the citation search feature, references will be automatically managed. For example, you might cite a paper like this: \cite{example2024}. All citations will be automatically added to your \texttt{cite.bib} file.

\section{Conclusion}
This template demonstrates basic \LaTeX{} features including mathematical formulas, tables, and lists. The bibliography system is now set up to work seamlessly with Bibby AI's citation management. Feel free to modify it for your needs!

\bibliographystyle{plain}
\bibliography{cite}

\end{document}\documentclass{article}
\usepackage[utf8]{inputenc}
\usepackage{booktabs}
\usepackage{amsmath}
\usepackage{natbib}

\title{My Document}
\author{Your Name}
\date{\today}

\begin{document}

\maketitle

\section{Introduction}
This is a sample document demonstrating various \LaTeX{} features including mathematical equations, tables, and formatting. You can cite papers using the citation search feature, and they will be automatically added to your \texttt{cite.bib} file.

\section{Mathematical Expressions}
Here are some example equations:

\subsection{Quadratic Formula}
The quadratic formula is given by:
\[x = \frac{-b \pm \sqrt{b^2 - 4ac}}{2a}\]

\subsection{Matrix Example}
A sample matrix:
\[
\begin{pmatrix}
a_{11} & a_{12} & a_{13} \\
a_{21} & a_{22} & a_{23} \\
a_{31} & a_{32} & a_{33}
\end{pmatrix}
\]

\section{Tables}
Below is a sample table showing data:

\begin{table}[h]
\centering
\begin{tabular}{@{}lrrr@{}}
\toprule
\textbf{Category} & \textbf{Value 1} & \textbf{Value 2} & \textbf{Total} \\
\midrule
Item A & 10.5 & 20.3 & 30.8 \\
Item B & 15.2 & 25.1 & 40.3 \\
Item C & 22.1 & 12.7 & 34.8 \\
\bottomrule
\end{tabular}
\caption{Sample Data Table}
\label{tab:sample}
\end{table}

\section{Lists}
\subsection{Bullet Points}
\begin{itemize}
  \item First item
  \item Second item
  \item Third item with sub-items:
    \begin{itemize}
      \item Sub-item 1
      \item Sub-item 2
    \end{itemize}
\end{itemize}

\subsection{Numbered List}
\begin{enumerate}
  \item First step
  \item Second step
  \item Third step
\end{enumerate}

\section{Citations and References}
When you use the citation search feature, references will be automatically managed. For example, you might cite a paper like this: \cite{example2024}. All citations will be automatically added to your \texttt{cite.bib} file.

\section{Conclusion}
This template demonstrates basic \LaTeX{} features including mathematical formulas, tables, and lists. The bibliography system is now set up to work seamlessly with Bibby AI's citation management. Feel free to modify it for your needs!

\bibliographystyle{plain}
\bibliography{cite}

\end{document}\documentclass{article}
\usepackage[utf8]{inputenc}
\usepackage{booktabs}
\usepackage{amsmath}
\usepackage{natbib}

\title{My Document}
\author{Your Name}
\date{\today}

\begin{document}

\maketitle

\section{Introduction}
This is a sample document demonstrating various \LaTeX{} features including mathematical equations, tables, and formatting. You can cite papers using the citation search feature, and they will be automatically added to your \texttt{cite.bib} file.

\section{Mathematical Expressions}
Here are some example equations:

\subsection{Quadratic Formula}
The quadratic formula is given by:
\[x = \frac{-b \pm \sqrt{b^2 - 4ac}}{2a}\]

\subsection{Matrix Example}
A sample matrix:
\[
\begin{pmatrix}
a_{11} & a_{12} & a_{13} \\
a_{21} & a_{22} & a_{23} \\
a_{31} & a_{32} & a_{33}
\end{pmatrix}
\]

\section{Tables}
Below is a sample table showing data:

\begin{table}[h]
\centering
\begin{tabular}{@{}lrrr@{}}
\toprule
\textbf{Category} & \textbf{Value 1} & \textbf{Value 2} & \textbf{Total} \\
\midrule
Item A & 10.5 & 20.3 & 30.8 \\
Item B & 15.2 & 25.1 & 40.3 \\
Item C & 22.1 & 12.7 & 34.8 \\
\bottomrule
\end{tabular}
\caption{Sample Data Table}
\label{tab:sample}
\end{table}

\section{Lists}
\subsection{Bullet Points}
\begin{itemize}
  \item First item
  \item Second item
  \item Third item with sub-items:
    \begin{itemize}
      \item Sub-item 1
      \item Sub-item 2
    \end{itemize}
\end{itemize}

\subsection{Numbered List}
\begin{enumerate}
  \item First step
  \item Second step
  \item Third step
\end{enumerate}

\section{Citations and References}
When you use the citation search feature, references will be automatically managed. For example, you might cite a paper like this: \cite{example2024}. All citations will be automatically added to your \texttt{cite.bib} file.

\section{Conclusion}
This template demonstrates basic \LaTeX{} features including mathematical formulas, tables, and lists. The bibliography system is now set up to work seamlessly with Bibby AI's citation management. Feel free to modify it for your needs!

\bibliographystyle{plain}
\bibliography{cite}

\end{document}\documentclass{article}
\usepackage[utf8]{inputenc}
\usepackage{booktabs}
\usepackage{amsmath}
\usepackage{natbib}

\title{My Document}
\author{Your Name}
\date{\today}

\begin{document}

\maketitle

\section{Introduction}
This is a sample document demonstrating various \LaTeX{} features including mathematical equations, tables, and formatting. You can cite papers using the citation search feature, and they will be automatically added to your \texttt{cite.bib} file.

\section{Mathematical Expressions}
Here are some example equations:

\subsection{Quadratic Formula}
The quadratic formula is given by:
\[x = \frac{-b \pm \sqrt{b^2 - 4ac}}{2a}\]

\subsection{Matrix Example}
A sample matrix:
\[
\begin{pmatrix}
a_{11} & a_{12} & a_{13} \\
a_{21} & a_{22} & a_{23} \\
a_{31} & a_{32} & a_{33}
\end{pmatrix}
\]

\section{Tables}
Below is a sample table showing data:

\begin{table}[h]
\centering
\begin{tabular}{@{}lrrr@{}}
\toprule
\textbf{Category} & \textbf{Value 1} & \textbf{Value 2} & \textbf{Total} \\
\midrule
Item A & 10.5 & 20.3 & 30.8 \\
Item B & 15.2 & 25.1 & 40.3 \\
Item C & 22.1 & 12.7 & 34.8 \\
\bottomrule
\end{tabular}
\caption{Sample Data Table}
\label{tab:sample}
\end{table}

\section{Lists}
\subsection{Bullet Points}
\begin{itemize}
  \item First item
  \item Second item
  \item Third item with sub-items:
    \begin{itemize}
      \item Sub-item 1
      \item Sub-item 2
    \end{itemize}
\end{itemize}

\subsection{Numbered List}
\begin{enumerate}
  \item First step
  \item Second step
  \item Third step
\end{enumerate}

\section{Citations and References}
When you use the citation search feature, references will be automatically managed. For example, you might cite a paper like this: \cite{example2024}. All citations will be automatically added to your \texttt{cite.bib} file.

\section{Conclusion}
This template demonstrates basic \LaTeX{} features including mathematical formulas, tables, and lists. The bibliography system is now set up to work seamlessly with Bibby AI's citation management. Feel free to modify it for your needs!

\bibliographystyle{plain}
\bibliography{cite}

\end{document}\documentclass{article}
\usepackage[utf8]{inputenc}
\usepackage{booktabs}
\usepackage{amsmath}
\usepackage{natbib}

\title{My Document}
\author{Your Name}
\date{\today}

\begin{document}

\maketitle

\section{Introduction}
This is a sample document demonstrating various \LaTeX{} features including mathematical equations, tables, and formatting. You can cite papers using the citation search feature, and they will be automatically added to your \texttt{cite.bib} file.

\section{Mathematical Expressions}
Here are some example equations:

\subsection{Quadratic Formula}
The quadratic formula is given by:
\[x = \frac{-b \pm \sqrt{b^2 - 4ac}}{2a}\]

\subsection{Matrix Example}
A sample matrix:
\[
\begin{pmatrix}
a_{11} & a_{12} & a_{13} \\
a_{21} & a_{22} & a_{23} \\
a_{31} & a_{32} & a_{33}
\end{pmatrix}
\]

\section{Tables}
Below is a sample table showing data:

\begin{table}[h]
\centering
\begin{tabular}{@{}lrrr@{}}
\toprule
\textbf{Category} & \textbf{Value 1} & \textbf{Value 2} & \textbf{Total} \\
\midrule
Item A & 10.5 & 20.3 & 30.8 \\
Item B & 15.2 & 25.1 & 40.3 \\
Item C & 22.1 & 12.7 & 34.8 \\
\bottomrule
\end{tabular}
\caption{Sample Data Table}
\label{tab:sample}
\end{table}

\section{Lists}
\subsection{Bullet Points}
\begin{itemize}
  \item First item
  \item Second item
  \item Third item with sub-items:
    \begin{itemize}
      \item Sub-item 1
      \item Sub-item 2
    \end{itemize}
\end{itemize}

\subsection{Numbered List}
\begin{enumerate}
  \item First step
  \item Second step
  \item Third step
\end{enumerate}

\section{Citations and References}
When you use the citation search feature, references will be automatically managed. For example, you might cite a paper like this: \cite{example2024}. All citations will be automatically added to your \texttt{cite.bib} file.

\section{Conclusion}
This template demonstrates basic \LaTeX{} features including mathematical formulas, tables, and lists. The bibliography system is now set up to work seamlessly with Bibby AI's citation management. Feel free to modify it for your needs!

\bibliographystyle{plain}
\bibliography{cite}

\end{document}\documentclass{article}
\usepackage[utf8]{inputenc}
\usepackage{booktabs}
\usepackage{amsmath}
\usepackage{natbib}

\title{My Document}
\author{Your Name}
\date{\today}

\begin{document}

\maketitle

\section{Introduction}
This is a sample document demonstrating various \LaTeX{} features including mathematical equations, tables, and formatting. You can cite papers using the citation search feature, and they will be automatically added to your \texttt{cite.bib} file.

\section{Mathematical Expressions}
Here are some example equations:

\subsection{Quadratic Formula}
The quadratic formula is given by:
\[x = \frac{-b \pm \sqrt{b^2 - 4ac}}{2a}\]

\subsection{Matrix Example}
A sample matrix:
\[
\begin{pmatrix}
a_{11} & a_{12} & a_{13} \\
a_{21} & a_{22} & a_{23} \\
a_{31} & a_{32} & a_{33}
\end{pmatrix}
\]

\section{Tables}
Below is a sample table showing data:

\begin{table}[h]
\centering
\begin{tabular}{@{}lrrr@{}}
\toprule
\textbf{Category} & \textbf{Value 1} & \textbf{Value 2} & \textbf{Total} \\
\midrule
Item A & 10.5 & 20.3 & 30.8 \\
Item B & 15.2 & 25.1 & 40.3 \\
Item C & 22.1 & 12.7 & 34.8 \\
\bottomrule
\end{tabular}
\caption{Sample Data Table}
\label{tab:sample}
\end{table}

\section{Lists}
\subsection{Bullet Points}
\begin{itemize}
  \item First item
  \item Second item
  \item Third item with sub-items:
    \begin{itemize}
      \item Sub-item 1
      \item Sub-item 2
    \end{itemize}
\end{itemize}

\subsection{Numbered List}
\begin{enumerate}
  \item First step
  \item Second step
  \item Third step
\end{enumerate}

\section{Citations and References}
When you use the citation search feature, references will be automatically managed. For example, you might cite a paper like this: \cite{example2024}. All citations will be automatically added to your \texttt{cite.bib} file.

\section{Conclusion}
This template demonstrates basic \LaTeX{} features including mathematical formulas, tables, and lists. The bibliography system is now set up to work seamlessly with Bibby AI's citation management. Feel free to modify it for your needs!

\bibliographystyle{plain}
\bibliography{cite}

\end{document}\documentclass{article}
\usepackage[utf8]{inputenc}
\usepackage{booktabs}
\usepackage{amsmath}
\usepackage{natbib}

\title{My Document}
\author{Your Name}
\date{\today}

\begin{document}

\maketitle

\section{Introduction}
This is a sample document demonstrating various \LaTeX{} features including mathematical equations, tables, and formatting. You can cite papers using the citation search feature, and they will be automatically added to your \texttt{cite.bib} file.

\section{Mathematical Expressions}
Here are some example equations:

\subsection{Quadratic Formula}
The quadratic formula is given by:
\[x = \frac{-b \pm \sqrt{b^2 - 4ac}}{2a}\]

\subsection{Matrix Example}
A sample matrix:
\[
\begin{pmatrix}
a_{11} & a_{12} & a_{13} \\
a_{21} & a_{22} & a_{23} \\
a_{31} & a_{32} & a_{33}
\end{pmatrix}
\]

\section{Tables}
Below is a sample table showing data:

\begin{table}[h]
\centering
\begin{tabular}{@{}lrrr@{}}
\toprule
\textbf{Category} & \textbf{Value 1} & \textbf{Value 2} & \textbf{Total} \\
\midrule
Item A & 10.5 & 20.3 & 30.8 \\
Item B & 15.2 & 25.1 & 40.3 \\
Item C & 22.1 & 12.7 & 34.8 \\
\bottomrule
\end{tabular}
\caption{Sample Data Table}
\label{tab:sample}
\end{table}

\section{Lists}
\subsection{Bullet Points}
\begin{itemize}
  \item First item
  \item Second item
  \item Third item with sub-items:
    \begin{itemize}
      \item Sub-item 1
      \item Sub-item 2
    \end{itemize}
\end{itemize}

\subsection{Numbered List}
\begin{enumerate}
  \item First step
  \item Second step
  \item Third step
\end{enumerate}

\section{Citations and References}
When you use the citation search feature, references will be automatically managed. For example, you might cite a paper like this: \cite{example2024}. All citations will be automatically added to your \texttt{cite.bib} file.

\section{Conclusion}
This template demonstrates basic \LaTeX{} features including mathematical formulas, tables, and lists. The bibliography system is now set up to work seamlessly with Bibby AI's citation management. Feel free to modify it for your needs!

\bibliographystyle{plain}
\bibliography{cite}

\end{document}\documentclass{article}
\usepackage[utf8]{inputenc}
\usepackage{booktabs}
\usepackage{amsmath}
\usepackage{natbib}

\title{My Document}
\author{Your Name}
\date{\today}

\begin{document}

\maketitle

\section{Introduction}
This is a sample document demonstrating various \LaTeX{} features including mathematical equations, tables, and formatting. You can cite papers using the citation search feature, and they will be automatically added to your \texttt{cite.bib} file.

\section{Mathematical Expressions}
Here are some example equations:

\subsection{Quadratic Formula}
The quadratic formula is given by:
\[x = \frac{-b \pm \sqrt{b^2 - 4ac}}{2a}\]

\subsection{Matrix Example}
A sample matrix:
\[
\begin{pmatrix}
a_{11} & a_{12} & a_{13} \\
a_{21} & a_{22} & a_{23} \\
a_{31} & a_{32} & a_{33}
\end{pmatrix}
\]

\section{Tables}
Below is a sample table showing data:

\begin{table}[h]
\centering
\begin{tabular}{@{}lrrr@{}}
\toprule
\textbf{Category} & \textbf{Value 1} & \textbf{Value 2} & \textbf{Total} \\
\midrule
Item A & 10.5 & 20.3 & 30.8 \\
Item B & 15.2 & 25.1 & 40.3 \\
Item C & 22.1 & 12.7 & 34.8 \\
\bottomrule
\end{tabular}
\caption{Sample Data Table}
\label{tab:sample}
\end{table}

\section{Lists}
\subsection{Bullet Points}
\begin{itemize}
  \item First item
  \item Second item
  \item Third item with sub-items:
    \begin{itemize}
      \item Sub-item 1
      \item Sub-item 2
    \end{itemize}
\end{itemize}

\subsection{Numbered List}
\begin{enumerate}
  \item First step
  \item Second step
  \item Third step
\end{enumerate}

\section{Citations and References}
When you use the citation search feature, references will be automatically managed. For example, you might cite a paper like this: \cite{example2024}. All citations will be automatically added to your \texttt{cite.bib} file.

\section{Conclusion}
This template demonstrates basic \LaTeX{} features including mathematical formulas, tables, and lists. The bibliography system is now set up to work seamlessly with Bibby AI's citation management. Feel free to modify it for your needs!

\bibliographystyle{plain}
\bibliography{cite}

\end{document}\documentclass{article}
\usepackage[utf8]{inputenc}
\usepackage{booktabs}
\usepackage{amsmath}
\usepackage{natbib}

\title{My Document}
\author{Your Name}
\date{\today}

\begin{document}

\maketitle

\section{Introduction}
This is a sample document demonstrating various \LaTeX{} features including mathematical equations, tables, and formatting. You can cite papers using the citation search feature, and they will be automatically added to your \texttt{cite.bib} file.

\section{Mathematical Expressions}
Here are some example equations:

\subsection{Quadratic Formula}
The quadratic formula is given by:
\[x = \frac{-b \pm \sqrt{b^2 - 4ac}}{2a}\]

\subsection{Matrix Example}
A sample matrix:
\[
\begin{pmatrix}
a_{11} & a_{12} & a_{13} \\
a_{21} & a_{22} & a_{23} \\
a_{31} & a_{32} & a_{33}
\end{pmatrix}
\]

\section{Tables}
Below is a sample table showing data:

\begin{table}[h]
\centering
\begin{tabular}{@{}lrrr@{}}
\toprule
\textbf{Category} & \textbf{Value 1} & \textbf{Value 2} & \textbf{Total} \\
\midrule
Item A & 10.5 & 20.3 & 30.8 \\
Item B & 15.2 & 25.1 & 40.3 \\
Item C & 22.1 & 12.7 & 34.8 \\
\bottomrule
\end{tabular}
\caption{Sample Data Table}
\label{tab:sample}
\end{table}

\section{Lists}
\subsection{Bullet Points}
\begin{itemize}
  \item First item
  \item Second item
  \item Third item with sub-items:
    \begin{itemize}
      \item Sub-item 1
      \item Sub-item 2
    \end{itemize}
\end{itemize}

\subsection{Numbered List}
\begin{enumerate}
  \item First step
  \item Second step
  \item Third step
\end{enumerate}

\section{Citations and References}
When you use the citation search feature, references will be automatically managed. For example, you might cite a paper like this: \cite{example2024}. All citations will be automatically added to your \texttt{cite.bib} file.

\section{Conclusion}
This template demonstrates basic \LaTeX{} features including mathematical formulas, tables, and lists. The bibliography system is now set up to work seamlessly with Bibby AI's citation management. Feel free to modify it for your needs!

\bibliographystyle{plain}
\bibliography{cite}

\end{document}\documentclass{article}
\usepackage[utf8]{inputenc}
\usepackage{booktabs}
\usepackage{amsmath}
\usepackage{natbib}

\title{My Document}
\author{Your Name}
\date{\today}

\begin{document}

\maketitle

\section{Introduction}
This is a sample document demonstrating various \LaTeX{} features including mathematical equations, tables, and formatting. You can cite papers using the citation search feature, and they will be automatically added to your \texttt{cite.bib} file.

\section{Mathematical Expressions}
Here are some example equations:

\subsection{Quadratic Formula}
The quadratic formula is given by:
\[x = \frac{-b \pm \sqrt{b^2 - 4ac}}{2a}\]

\subsection{Matrix Example}
A sample matrix:
\[
\begin{pmatrix}
a_{11} & a_{12} & a_{13} \\
a_{21} & a_{22} & a_{23} \\
a_{31} & a_{32} & a_{33}
\end{pmatrix}
\]

\section{Tables}
Below is a sample table showing data:

\begin{table}[h]
\centering
\begin{tabular}{@{}lrrr@{}}
\toprule
\textbf{Category} & \textbf{Value 1} & \textbf{Value 2} & \textbf{Total} \\
\midrule
Item A & 10.5 & 20.3 & 30.8 \\
Item B & 15.2 & 25.1 & 40.3 \\
Item C & 22.1 & 12.7 & 34.8 \\
\bottomrule
\end{tabular}
\caption{Sample Data Table}
\label{tab:sample}
\end{table}

\section{Lists}
\subsection{Bullet Points}
\begin{itemize}
  \item First item
  \item Second item
  \item Third item with sub-items:
    \begin{itemize}
      \item Sub-item 1
      \item Sub-item 2
    \end{itemize}
\end{itemize}

\subsection{Numbered List}
\begin{enumerate}
  \item First step
  \item Second step
  \item Third step
\end{enumerate}

\section{Citations and References}
When you use the citation search feature, references will be automatically managed. For example, you might cite a paper like this: \cite{example2024}. All citations will be automatically added to your \texttt{cite.bib} file.

\section{Conclusion}
This template demonstrates basic \LaTeX{} features including mathematical formulas, tables, and lists. The bibliography system is now set up to work seamlessly with Bibby AI's citation management. Feel free to modify it for your needs!

\bibliographystyle{plain}
\bibliography{cite}

\end{document}\documentclass{article}
\usepackage[utf8]{inputenc}
\usepackage{booktabs}
\usepackage{amsmath}
\usepackage{natbib}

\title{My Document}
\author{Your Name}
\date{\today}

\begin{document}

\maketitle

\section{Introduction}
This is a sample document demonstrating various \LaTeX{} features including mathematical equations, tables, and formatting. You can cite papers using the citation search feature, and they will be automatically added to your \texttt{cite.bib} file.

\section{Mathematical Expressions}
Here are some example equations:

\subsection{Quadratic Formula}
The quadratic formula is given by:
\[x = \frac{-b \pm \sqrt{b^2 - 4ac}}{2a}\]

\subsection{Matrix Example}
A sample matrix:
\[
\begin{pmatrix}
a_{11} & a_{12} & a_{13} \\
a_{21} & a_{22} & a_{23} \\
a_{31} & a_{32} & a_{33}
\end{pmatrix}
\]

\section{Tables}
Below is a sample table showing data:

\begin{table}[h]
\centering
\begin{tabular}{@{}lrrr@{}}
\toprule
\textbf{Category} & \textbf{Value 1} & \textbf{Value 2} & \textbf{Total} \\
\midrule
Item A & 10.5 & 20.3 & 30.8 \\
Item B & 15.2 & 25.1 & 40.3 \\
Item C & 22.1 & 12.7 & 34.8 \\
\bottomrule
\end{tabular}
\caption{Sample Data Table}
\label{tab:sample}
\end{table}

\section{Lists}
\subsection{Bullet Points}
\begin{itemize}
  \item First item
  \item Second item
  \item Third item with sub-items:
    \begin{itemize}
      \item Sub-item 1
      \item Sub-item 2
    \end{itemize}
\end{itemize}

\subsection{Numbered List}
\begin{enumerate}
  \item First step
  \item Second step
  \item Third step
\end{enumerate}

\section{Citations and References}
When you use the citation search feature, references will be automatically managed. For example, you might cite a paper like this: \cite{example2024}. All citations will be automatically added to your \texttt{cite.bib} file.

\section{Conclusion}
This template demonstrates basic \LaTeX{} features including mathematical formulas, tables, and lists. The bibliography system is now set up to work seamlessly with Bibby AI's citation management. Feel free to modify it for your needs!

\bibliographystyle{plain}
\bibliography{cite}

\end{document}\documentclass{article}
\usepackage[utf8]{inputenc}
\usepackage{booktabs}
\usepackage{amsmath}
\usepackage{natbib}

\title{My Document}
\author{Your Name}
\date{\today}

\begin{document}

\maketitle

\section{Introduction}
This is a sample document demonstrating various \LaTeX{} features including mathematical equations, tables, and formatting. You can cite papers using the citation search feature, and they will be automatically added to your \texttt{cite.bib} file.

\section{Mathematical Expressions}
Here are some example equations:

\subsection{Quadratic Formula}
The quadratic formula is given by:
\[x = \frac{-b \pm \sqrt{b^2 - 4ac}}{2a}\]

\subsection{Matrix Example}
A sample matrix:
\[
\begin{pmatrix}
a_{11} & a_{12} & a_{13} \\
a_{21} & a_{22} & a_{23} \\
a_{31} & a_{32} & a_{33}
\end{pmatrix}
\]

\section{Tables}
Below is a sample table showing data:

\begin{table}[h]
\centering
\begin{tabular}{@{}lrrr@{}}
\toprule
\textbf{Category} & \textbf{Value 1} & \textbf{Value 2} & \textbf{Total} \\
\midrule
Item A & 10.5 & 20.3 & 30.8 \\
Item B & 15.2 & 25.1 & 40.3 \\
Item C & 22.1 & 12.7 & 34.8 \\
\bottomrule
\end{tabular}
\caption{Sample Data Table}
\label{tab:sample}
\end{table}

\section{Lists}
\subsection{Bullet Points}
\begin{itemize}
  \item First item
  \item Second item
  \item Third item with sub-items:
    \begin{itemize}
      \item Sub-item 1
      \item Sub-item 2
    \end{itemize}
\end{itemize}

\subsection{Numbered List}
\begin{enumerate}
  \item First step
  \item Second step
  \item Third step
\end{enumerate}

\section{Citations and References}
When you use the citation search feature, references will be automatically managed. For example, you might cite a paper like this: \cite{example2024}. All citations will be automatically added to your \texttt{cite.bib} file.

\section{Conclusion}
This template demonstrates basic \LaTeX{} features including mathematical formulas, tables, and lists. The bibliography system is now set up to work seamlessly with Bibby AI's citation management. Feel free to modify it for your needs!

\bibliographystyle{plain}
\bibliography{cite}

\end{document}\documentclass{article}
\usepackage[utf8]{inputenc}
\usepackage{booktabs}
\usepackage{amsmath}
\usepackage{natbib}

\title{My Document}
\author{Your Name}
\date{\today}

\begin{document}

\maketitle

\section{Introduction}
This is a sample document demonstrating various \LaTeX{} features including mathematical equations, tables, and formatting. You can cite papers using the citation search feature, and they will be automatically added to your \texttt{cite.bib} file.

\section{Mathematical Expressions}
Here are some example equations:

\subsection{Quadratic Formula}
The quadratic formula is given by:
\[x = \frac{-b \pm \sqrt{b^2 - 4ac}}{2a}\]

\subsection{Matrix Example}
A sample matrix:
\[
\begin{pmatrix}
a_{11} & a_{12} & a_{13} \\
a_{21} & a_{22} & a_{23} \\
a_{31} & a_{32} & a_{33}
\end{pmatrix}
\]

\section{Tables}
Below is a sample table showing data:

\begin{table}[h]
\centering
\begin{tabular}{@{}lrrr@{}}
\toprule
\textbf{Category} & \textbf{Value 1} & \textbf{Value 2} & \textbf{Total} \\
\midrule
Item A & 10.5 & 20.3 & 30.8 \\
Item B & 15.2 & 25.1 & 40.3 \\
Item C & 22.1 & 12.7 & 34.8 \\
\bottomrule
\end{tabular}
\caption{Sample Data Table}
\label{tab:sample}
\end{table}

\section{Lists}
\subsection{Bullet Points}
\begin{itemize}
  \item First item
  \item Second item
  \item Third item with sub-items:
    \begin{itemize}
      \item Sub-item 1
      \item Sub-item 2
    \end{itemize}
\end{itemize}

\subsection{Numbered List}
\begin{enumerate}
  \item First step
  \item Second step
  \item Third step
\end{enumerate}

\section{Citations and References}
When you use the citation search feature, references will be automatically managed. For example, you might cite a paper like this: \cite{example2024}. All citations will be automatically added to your \texttt{cite.bib} file.

\section{Conclusion}
This template demonstrates basic \LaTeX{} features including mathematical formulas, tables, and lists. The bibliography system is now set up to work seamlessly with Bibby AI's citation management. Feel free to modify it for your needs!

\bibliographystyle{plain}
\bibliography{cite}

\end{document}\documentclass{article}
\usepackage[utf8]{inputenc}
\usepackage{booktabs}
\usepackage{amsmath}
\usepackage{natbib}

\title{My Document}
\author{Your Name}
\date{\today}

\begin{document}

\maketitle

\section{Introduction}
This is a sample document demonstrating various \LaTeX{} features including mathematical equations, tables, and formatting. You can cite papers using the citation search feature, and they will be automatically added to your \texttt{cite.bib} file.

\section{Mathematical Expressions}
Here are some example equations:

\subsection{Quadratic Formula}
The quadratic formula is given by:
\[x = \frac{-b \pm \sqrt{b^2 - 4ac}}{2a}\]

\subsection{Matrix Example}
A sample matrix:
\[
\begin{pmatrix}
a_{11} & a_{12} & a_{13} \\
a_{21} & a_{22} & a_{23} \\
a_{31} & a_{32} & a_{33}
\end{pmatrix}
\]

\section{Tables}
Below is a sample table showing data:

\begin{table}[h]
\centering
\begin{tabular}{@{}lrrr@{}}
\toprule
\textbf{Category} & \textbf{Value 1} & \textbf{Value 2} & \textbf{Total} \\
\midrule
Item A & 10.5 & 20.3 & 30.8 \\
Item B & 15.2 & 25.1 & 40.3 \\
Item C & 22.1 & 12.7 & 34.8 \\
\bottomrule
\end{tabular}
\caption{Sample Data Table}
\label{tab:sample}
\end{table}

\section{Lists}
\subsection{Bullet Points}
\begin{itemize}
  \item First item
  \item Second item
  \item Third item with sub-items:
    \begin{itemize}
      \item Sub-item 1
      \item Sub-item 2
    \end{itemize}
\end{itemize}

\subsection{Numbered List}
\begin{enumerate}
  \item First step
  \item Second step
  \item Third step
\end{enumerate}

\section{Citations and References}
When you use the citation search feature, references will be automatically managed. For example, you might cite a paper like this: \cite{example2024}. All citations will be automatically added to your \texttt{cite.bib} file.

\section{Conclusion}
This template demonstrates basic \LaTeX{} features including mathematical formulas, tables, and lists. The bibliography system is now set up to work seamlessly with Bibby AI's citation management. Feel free to modify it for your needs!

\bibliographystyle{plain}
\bibliography{cite}

\end{document}\documentclass{article}
\usepackage[utf8]{inputenc}
\usepackage{booktabs}
\usepackage{amsmath}
\usepackage{natbib}

\title{My Document}
\author{Your Name}
\date{\today}

\begin{document}

\maketitle

\section{Introduction}
This is a sample document demonstrating various \LaTeX{} features including mathematical equations, tables, and formatting. You can cite papers using the citation search feature, and they will be automatically added to your \texttt{cite.bib} file.

\section{Mathematical Expressions}
Here are some example equations:

\subsection{Quadratic Formula}
The quadratic formula is given by:
\[x = \frac{-b \pm \sqrt{b^2 - 4ac}}{2a}\]

\subsection{Matrix Example}
A sample matrix:
\[
\begin{pmatrix}
a_{11} & a_{12} & a_{13} \\
a_{21} & a_{22} & a_{23} \\
a_{31} & a_{32} & a_{33}
\end{pmatrix}
\]

\section{Tables}
Below is a sample table showing data:

\begin{table}[h]
\centering
\begin{tabular}{@{}lrrr@{}}
\toprule
\textbf{Category} & \textbf{Value 1} & \textbf{Value 2} & \textbf{Total} \\
\midrule
Item A & 10.5 & 20.3 & 30.8 \\
Item B & 15.2 & 25.1 & 40.3 \\
Item C & 22.1 & 12.7 & 34.8 \\
\bottomrule
\end{tabular}
\caption{Sample Data Table}
\label{tab:sample}
\end{table}

\section{Lists}
\subsection{Bullet Points}
\begin{itemize}
  \item First item
  \item Second item
  \item Third item with sub-items:
    \begin{itemize}
      \item Sub-item 1
      \item Sub-item 2
    \end{itemize}
\end{itemize}

\subsection{Numbered List}
\begin{enumerate}
  \item First step
  \item Second step
  \item Third step
\end{enumerate}

\section{Citations and References}
When you use the citation search feature, references will be automatically managed. For example, you might cite a paper like this: \cite{example2024}. All citations will be automatically added to your \texttt{cite.bib} file.

\section{Conclusion}
This template demonstrates basic \LaTeX{} features including mathematical formulas, tables, and lists. The bibliography system is now set up to work seamlessly with Bibby AI's citation management. Feel free to modify it for your needs!

\bibliographystyle{plain}
\bibliography{cite}

\end{document}\documentclass{article}
\usepackage[utf8]{inputenc}
\usepackage{booktabs}
\usepackage{amsmath}
\usepackage{natbib}

\title{My Document}
\author{Your Name}
\date{\today}

\begin{document}

\maketitle

\section{Introduction}
This is a sample document demonstrating various \LaTeX{} features including mathematical equations, tables, and formatting. You can cite papers using the citation search feature, and they will be automatically added to your \texttt{cite.bib} file.

\section{Mathematical Expressions}
Here are some example equations:

\subsection{Quadratic Formula}
The quadratic formula is given by:
\[x = \frac{-b \pm \sqrt{b^2 - 4ac}}{2a}\]

\subsection{Matrix Example}
A sample matrix:
\[
\begin{pmatrix}
a_{11} & a_{12} & a_{13} \\
a_{21} & a_{22} & a_{23} \\
a_{31} & a_{32} & a_{33}
\end{pmatrix}
\]

\section{Tables}
Below is a sample table showing data:

\begin{table}[h]
\centering
\begin{tabular}{@{}lrrr@{}}
\toprule
\textbf{Category} & \textbf{Value 1} & \textbf{Value 2} & \textbf{Total} \\
\midrule
Item A & 10.5 & 20.3 & 30.8 \\
Item B & 15.2 & 25.1 & 40.3 \\
Item C & 22.1 & 12.7 & 34.8 \\
\bottomrule
\end{tabular}
\caption{Sample Data Table}
\label{tab:sample}
\end{table}

\section{Lists}
\subsection{Bullet Points}
\begin{itemize}
  \item First item
  \item Second item
  \item Third item with sub-items:
    \begin{itemize}
      \item Sub-item 1
      \item Sub-item 2
    \end{itemize}
\end{itemize}

\subsection{Numbered List}
\begin{enumerate}
  \item First step
  \item Second step
  \item Third step
\end{enumerate}

\section{Citations and References}
When you use the citation search feature, references will be automatically managed. For example, you might cite a paper like this: \cite{example2024}. All citations will be automatically added to your \texttt{cite.bib} file.

\section{Conclusion}
This template demonstrates basic \LaTeX{} features including mathematical formulas, tables, and lists. The bibliography system is now set up to work seamlessly with Bibby AI's citation management. Feel free to modify it for your needs!

\bibliographystyle{plain}
\bibliography{cite}

\end{document}\documentclass{article}
\usepackage[utf8]{inputenc}
\usepackage{booktabs}
\usepackage{amsmath}
\usepackage{natbib}

\title{My Document}
\author{Your Name}
\date{\today}

\begin{document}

\maketitle

\section{Introduction}
This is a sample document demonstrating various \LaTeX{} features including mathematical equations, tables, and formatting. You can cite papers using the citation search feature, and they will be automatically added to your \texttt{cite.bib} file.

\section{Mathematical Expressions}
Here are some example equations:

\subsection{Quadratic Formula}
The quadratic formula is given by:
\[x = \frac{-b \pm \sqrt{b^2 - 4ac}}{2a}\]

\subsection{Matrix Example}
A sample matrix:
\[
\begin{pmatrix}
a_{11} & a_{12} & a_{13} \\
a_{21} & a_{22} & a_{23} \\
a_{31} & a_{32} & a_{33}
\end{pmatrix}
\]

\section{Tables}
Below is a sample table showing data:

\begin{table}[h]
\centering
\begin{tabular}{@{}lrrr@{}}
\toprule
\textbf{Category} & \textbf{Value 1} & \textbf{Value 2} & \textbf{Total} \\
\midrule
Item A & 10.5 & 20.3 & 30.8 \\
Item B & 15.2 & 25.1 & 40.3 \\
Item C & 22.1 & 12.7 & 34.8 \\
\bottomrule
\end{tabular}
\caption{Sample Data Table}
\label{tab:sample}
\end{table}

\section{Lists}
\subsection{Bullet Points}
\begin{itemize}
  \item First item
  \item Second item
  \item Third item with sub-items:
    \begin{itemize}
      \item Sub-item 1
      \item Sub-item 2
    \end{itemize}
\end{itemize}

\subsection{Numbered List}
\begin{enumerate}
  \item First step
  \item Second step
  \item Third step
\end{enumerate}

\section{Citations and References}
When you use the citation search feature, references will be automatically managed. For example, you might cite a paper like this: \cite{example2024}. All citations will be automatically added to your \texttt{cite.bib} file.

\section{Conclusion}
This template demonstrates basic \LaTeX{} features including mathematical formulas, tables, and lists. The bibliography system is now set up to work seamlessly with Bibby AI's citation management. Feel free to modify it for your needs!

\bibliographystyle{plain}
\bibliography{cite}

\end{document}\documentclass{article}
\usepackage[utf8]{inputenc}
\usepackage{booktabs}
\usepackage{amsmath}
\usepackage{natbib}

\title{My Document}
\author{Your Name}
\date{\today}

\begin{document}

\maketitle

\section{Introduction}
This is a sample document demonstrating various \LaTeX{} features including mathematical equations, tables, and formatting. You can cite papers using the citation search feature, and they will be automatically added to your \texttt{cite.bib} file.

\section{Mathematical Expressions}
Here are some example equations:

\subsection{Quadratic Formula}
The quadratic formula is given by:
\[x = \frac{-b \pm \sqrt{b^2 - 4ac}}{2a}\]

\subsection{Matrix Example}
A sample matrix:
\[
\begin{pmatrix}
a_{11} & a_{12} & a_{13} \\
a_{21} & a_{22} & a_{23} \\
a_{31} & a_{32} & a_{33}
\end{pmatrix}
\]

\section{Tables}
Below is a sample table showing data:

\begin{table}[h]
\centering
\begin{tabular}{@{}lrrr@{}}
\toprule
\textbf{Category} & \textbf{Value 1} & \textbf{Value 2} & \textbf{Total} \\
\midrule
Item A & 10.5 & 20.3 & 30.8 \\
Item B & 15.2 & 25.1 & 40.3 \\
Item C & 22.1 & 12.7 & 34.8 \\
\bottomrule
\end{tabular}
\caption{Sample Data Table}
\label{tab:sample}
\end{table}

\section{Lists}
\subsection{Bullet Points}
\begin{itemize}
  \item First item
  \item Second item
  \item Third item with sub-items:
    \begin{itemize}
      \item Sub-item 1
      \item Sub-item 2
    \end{itemize}
\end{itemize}

\subsection{Numbered List}
\begin{enumerate}
  \item First step
  \item Second step
  \item Third step
\end{enumerate}

\section{Citations and References}
When you use the citation search feature, references will be automatically managed. For example, you might cite a paper like this: \cite{example2024}. All citations will be automatically added to your \texttt{cite.bib} file.

\section{Conclusion}
This template demonstrates basic \LaTeX{} features including mathematical formulas, tables, and lists. The bibliography system is now set up to work seamlessly with Bibby AI's citation management. Feel free to modify it for your needs!

\bibliographystyle{plain}
\bibliography{cite}

\end{document}\documentclass{article}
\usepackage[utf8]{inputenc}
\usepackage{booktabs}
\usepackage{amsmath}
\usepackage{natbib}

\title{My Document}
\author{Your Name}
\date{\today}

\begin{document}

\maketitle

\section{Introduction}
This is a sample document demonstrating various \LaTeX{} features including mathematical equations, tables, and formatting. You can cite papers using the citation search feature, and they will be automatically added to your \texttt{cite.bib} file.

\section{Mathematical Expressions}
Here are some example equations:

\subsection{Quadratic Formula}
The quadratic formula is given by:
\[x = \frac{-b \pm \sqrt{b^2 - 4ac}}{2a}\]

\subsection{Matrix Example}
A sample matrix:
\[
\begin{pmatrix}
a_{11} & a_{12} & a_{13} \\
a_{21} & a_{22} & a_{23} \\
a_{31} & a_{32} & a_{33}
\end{pmatrix}
\]

\section{Tables}
Below is a sample table showing data:

\begin{table}[h]
\centering
\begin{tabular}{@{}lrrr@{}}
\toprule
\textbf{Category} & \textbf{Value 1} & \textbf{Value 2} & \textbf{Total} \\
\midrule
Item A & 10.5 & 20.3 & 30.8 \\
Item B & 15.2 & 25.1 & 40.3 \\
Item C & 22.1 & 12.7 & 34.8 \\
\bottomrule
\end{tabular}
\caption{Sample Data Table}
\label{tab:sample}
\end{table}

\section{Lists}
\subsection{Bullet Points}
\begin{itemize}
  \item First item
  \item Second item
  \item Third item with sub-items:
    \begin{itemize}
      \item Sub-item 1
      \item Sub-item 2
    \end{itemize}
\end{itemize}

\subsection{Numbered List}
\begin{enumerate}
  \item First step
  \item Second step
  \item Third step
\end{enumerate}

\section{Citations and References}
When you use the citation search feature, references will be automatically managed. For example, you might cite a paper like this: \cite{example2024}. All citations will be automatically added to your \texttt{cite.bib} file.

\section{Conclusion}
This template demonstrates basic \LaTeX{} features including mathematical formulas, tables, and lists. The bibliography system is now set up to work seamlessly with Bibby AI's citation management. Feel free to modify it for your needs!

\bibliographystyle{plain}
\bibliography{cite}

\end{document}\documentclass{article}
\usepackage[utf8]{inputenc}
\usepackage{booktabs}
\usepackage{amsmath}
\usepackage{natbib}

\title{My Document}
\author{Your Name}
\date{\today}

\begin{document}

\maketitle

\section{Introduction}
This is a sample document demonstrating various \LaTeX{} features including mathematical equations, tables, and formatting. You can cite papers using the citation search feature, and they will be automatically added to your \texttt{cite.bib} file.

\section{Mathematical Expressions}
Here are some example equations:

\subsection{Quadratic Formula}
The quadratic formula is given by:
\[x = \frac{-b \pm \sqrt{b^2 - 4ac}}{2a}\]

\subsection{Matrix Example}
A sample matrix:
\[
\begin{pmatrix}
a_{11} & a_{12} & a_{13} \\
a_{21} & a_{22} & a_{23} \\
a_{31} & a_{32} & a_{33}
\end{pmatrix}
\]

\section{Tables}
Below is a sample table showing data:

\begin{table}[h]
\centering
\begin{tabular}{@{}lrrr@{}}
\toprule
\textbf{Category} & \textbf{Value 1} & \textbf{Value 2} & \textbf{Total} \\
\midrule
Item A & 10.5 & 20.3 & 30.8 \\
Item B & 15.2 & 25.1 & 40.3 \\
Item C & 22.1 & 12.7 & 34.8 \\
\bottomrule
\end{tabular}
\caption{Sample Data Table}
\label{tab:sample}
\end{table}

\section{Lists}
\subsection{Bullet Points}
\begin{itemize}
  \item First item
  \item Second item
  \item Third item with sub-items:
    \begin{itemize}
      \item Sub-item 1
      \item Sub-item 2
    \end{itemize}
\end{itemize}

\subsection{Numbered List}
\begin{enumerate}
  \item First step
  \item Second step
  \item Third step
\end{enumerate}

\section{Citations and References}
When you use the citation search feature, references will be automatically managed. For example, you might cite a paper like this: \cite{example2024}. All citations will be automatically added to your \texttt{cite.bib} file.

\section{Conclusion}
This template demonstrates basic \LaTeX{} features including mathematical formulas, tables, and lists. The bibliography system is now set up to work seamlessly with Bibby AI's citation management. Feel free to modify it for your needs!

\bibliographystyle{plain}
\bibliography{cite}

\end{document}\documentclass{article}
\usepackage[utf8]{inputenc}
\usepackage{booktabs}
\usepackage{amsmath}
\usepackage{natbib}

\title{My Document}
\author{Your Name}
\date{\today}

\begin{document}

\maketitle

\section{Introduction}
This is a sample document demonstrating various \LaTeX{} features including mathematical equations, tables, and formatting. You can cite papers using the citation search feature, and they will be automatically added to your \texttt{cite.bib} file.

\section{Mathematical Expressions}
Here are some example equations:

\subsection{Quadratic Formula}
The quadratic formula is given by:
\[x = \frac{-b \pm \sqrt{b^2 - 4ac}}{2a}\]

\subsection{Matrix Example}
A sample matrix:
\[
\begin{pmatrix}
a_{11} & a_{12} & a_{13} \\
a_{21} & a_{22} & a_{23} \\
a_{31} & a_{32} & a_{33}
\end{pmatrix}
\]

\section{Tables}
Below is a sample table showing data:

\begin{table}[h]
\centering
\begin{tabular}{@{}lrrr@{}}
\toprule
\textbf{Category} & \textbf{Value 1} & \textbf{Value 2} & \textbf{Total} \\
\midrule
Item A & 10.5 & 20.3 & 30.8 \\
Item B & 15.2 & 25.1 & 40.3 \\
Item C & 22.1 & 12.7 & 34.8 \\
\bottomrule
\end{tabular}
\caption{Sample Data Table}
\label{tab:sample}
\end{table}

\section{Lists}
\subsection{Bullet Points}
\begin{itemize}
  \item First item
  \item Second item
  \item Third item with sub-items:
    \begin{itemize}
      \item Sub-item 1
      \item Sub-item 2
    \end{itemize}
\end{itemize}

\subsection{Numbered List}
\begin{enumerate}
  \item First step
  \item Second step
  \item Third step
\end{enumerate}

\section{Citations and References}
When you use the citation search feature, references will be automatically managed. For example, you might cite a paper like this: \cite{example2024}. All citations will be automatically added to your \texttt{cite.bib} file.

\section{Conclusion}
This template demonstrates basic \LaTeX{} features including mathematical formulas, tables, and lists. The bibliography system is now set up to work seamlessly with Bibby AI's citation management. Feel free to modify it for your needs!

\bibliographystyle{plain}
\bibliography{cite}

\end{document}\documentclass{article}
\usepackage[utf8]{inputenc}
\usepackage{booktabs}
\usepackage{amsmath}
\usepackage{natbib}

\title{My Document}
\author{Your Name}
\date{\today}

\begin{document}

\maketitle

\section{Introduction}
This is a sample document demonstrating various \LaTeX{} features including mathematical equations, tables, and formatting. You can cite papers using the citation search feature, and they will be automatically added to your \texttt{cite.bib} file.

\section{Mathematical Expressions}
Here are some example equations:

\subsection{Quadratic Formula}
The quadratic formula is given by:
\[x = \frac{-b \pm \sqrt{b^2 - 4ac}}{2a}\]

\subsection{Matrix Example}
A sample matrix:
\[
\begin{pmatrix}
a_{11} & a_{12} & a_{13} \\
a_{21} & a_{22} & a_{23} \\
a_{31} & a_{32} & a_{33}
\end{pmatrix}
\]

\section{Tables}
Below is a sample table showing data:

\begin{table}[h]
\centering
\begin{tabular}{@{}lrrr@{}}
\toprule
\textbf{Category} & \textbf{Value 1} & \textbf{Value 2} & \textbf{Total} \\
\midrule
Item A & 10.5 & 20.3 & 30.8 \\
Item B & 15.2 & 25.1 & 40.3 \\
Item C & 22.1 & 12.7 & 34.8 \\
\bottomrule
\end{tabular}
\caption{Sample Data Table}
\label{tab:sample}
\end{table}

\section{Lists}
\subsection{Bullet Points}
\begin{itemize}
  \item First item
  \item Second item
  \item Third item with sub-items:
    \begin{itemize}
      \item Sub-item 1
      \item Sub-item 2
    \end{itemize}
\end{itemize}

\subsection{Numbered List}
\begin{enumerate}
  \item First step
  \item Second step
  \item Third step
\end{enumerate}

\section{Citations and References}
When you use the citation search feature, references will be automatically managed. For example, you might cite a paper like this: \cite{example2024}. All citations will be automatically added to your \texttt{cite.bib} file.

\section{Conclusion}
This template demonstrates basic \LaTeX{} features including mathematical formulas, tables, and lists. The bibliography system is now set up to work seamlessly with Bibby AI's citation management. Feel free to modify it for your needs!

\bibliographystyle{plain}
\bibliography{cite}

\end{document}\documentclass{article}
\usepackage[utf8]{inputenc}
\usepackage{booktabs}
\usepackage{amsmath}
\usepackage{natbib}

\title{My Document}
\author{Your Name}
\date{\today}

\begin{document}

\maketitle

\section{Introduction}
This is a sample document demonstrating various \LaTeX{} features including mathematical equations, tables, and formatting. You can cite papers using the citation search feature, and they will be automatically added to your \texttt{cite.bib} file.

\section{Mathematical Expressions}
Here are some example equations:

\subsection{Quadratic Formula}
The quadratic formula is given by:
\[x = \frac{-b \pm \sqrt{b^2 - 4ac}}{2a}\]

\subsection{Matrix Example}
A sample matrix:
\[
\begin{pmatrix}
a_{11} & a_{12} & a_{13} \\
a_{21} & a_{22} & a_{23} \\
a_{31} & a_{32} & a_{33}
\end{pmatrix}
\]

\section{Tables}
Below is a sample table showing data:

\begin{table}[h]
\centering
\begin{tabular}{@{}lrrr@{}}
\toprule
\textbf{Category} & \textbf{Value 1} & \textbf{Value 2} & \textbf{Total} \\
\midrule
Item A & 10.5 & 20.3 & 30.8 \\
Item B & 15.2 & 25.1 & 40.3 \\
Item C & 22.1 & 12.7 & 34.8 \\
\bottomrule
\end{tabular}
\caption{Sample Data Table}
\label{tab:sample}
\end{table}

\section{Lists}
\subsection{Bullet Points}
\begin{itemize}
  \item First item
  \item Second item
  \item Third item with sub-items:
    \begin{itemize}
      \item Sub-item 1
      \item Sub-item 2
    \end{itemize}
\end{itemize}

\subsection{Numbered List}
\begin{enumerate}
  \item First step
  \item Second step
  \item Third step
\end{enumerate}

\section{Citations and References}
When you use the citation search feature, references will be automatically managed. For example, you might cite a paper like this: \cite{example2024}. All citations will be automatically added to your \texttt{cite.bib} file.

\section{Conclusion}
This template demonstrates basic \LaTeX{} features including mathematical formulas, tables, and lists. The bibliography system is now set up to work seamlessly with Bibby AI's citation management. Feel free to modify it for your needs!

\bibliographystyle{plain}
\bibliography{cite}

\end{document}\documentclass{article}
\usepackage[utf8]{inputenc}
\usepackage{booktabs}
\usepackage{amsmath}
\usepackage{natbib}

\title{My Document}
\author{Your Name}
\date{\today}

\begin{document}

\maketitle

\section{Introduction}
This is a sample document demonstrating various \LaTeX{} features including mathematical equations, tables, and formatting. You can cite papers using the citation search feature, and they will be automatically added to your \texttt{cite.bib} file.

\section{Mathematical Expressions}
Here are some example equations:

\subsection{Quadratic Formula}
The quadratic formula is given by:
\[x = \frac{-b \pm \sqrt{b^2 - 4ac}}{2a}\]

\subsection{Matrix Example}
A sample matrix:
\[
\begin{pmatrix}
a_{11} & a_{12} & a_{13} \\
a_{21} & a_{22} & a_{23} \\
a_{31} & a_{32} & a_{33}
\end{pmatrix}
\]

\section{Tables}
Below is a sample table showing data:

\begin{table}[h]
\centering
\begin{tabular}{@{}lrrr@{}}
\toprule
\textbf{Category} & \textbf{Value 1} & \textbf{Value 2} & \textbf{Total} \\
\midrule
Item A & 10.5 & 20.3 & 30.8 \\
Item B & 15.2 & 25.1 & 40.3 \\
Item C & 22.1 & 12.7 & 34.8 \\
\bottomrule
\end{tabular}
\caption{Sample Data Table}
\label{tab:sample}
\end{table}

\section{Lists}
\subsection{Bullet Points}
\begin{itemize}
  \item First item
  \item Second item
  \item Third item with sub-items:
    \begin{itemize}
      \item Sub-item 1
      \item Sub-item 2
    \end{itemize}
\end{itemize}

\subsection{Numbered List}
\begin{enumerate}
  \item First step
  \item Second step
  \item Third step
\end{enumerate}

\section{Citations and References}
When you use the citation search feature, references will be automatically managed. For example, you might cite a paper like this: \cite{example2024}. All citations will be automatically added to your \texttt{cite.bib} file.

\section{Conclusion}
This template demonstrates basic \LaTeX{} features including mathematical formulas, tables, and lists. The bibliography system is now set up to work seamlessly with Bibby AI's citation management. Feel free to modify it for your needs!

\bibliographystyle{plain}
\bibliography{cite}

\end{document}\documentclass{article}
\usepackage[utf8]{inputenc}
\usepackage{booktabs}
\usepackage{amsmath}
\usepackage{natbib}

\title{My Document}
\author{Your Name}
\date{\today}

\begin{document}

\maketitle

\section{Introduction}
This is a sample document demonstrating various \LaTeX{} features including mathematical equations, tables, and formatting. You can cite papers using the citation search feature, and they will be automatically added to your \texttt{cite.bib} file.

\section{Mathematical Expressions}
Here are some example equations:

\subsection{Quadratic Formula}
The quadratic formula is given by:
\[x = \frac{-b \pm \sqrt{b^2 - 4ac}}{2a}\]

\subsection{Matrix Example}
A sample matrix:
\[
\begin{pmatrix}
a_{11} & a_{12} & a_{13} \\
a_{21} & a_{22} & a_{23} \\
a_{31} & a_{32} & a_{33}
\end{pmatrix}
\]

\section{Tables}
Below is a sample table showing data:

\begin{table}[h]
\centering
\begin{tabular}{@{}lrrr@{}}
\toprule
\textbf{Category} & \textbf{Value 1} & \textbf{Value 2} & \textbf{Total} \\
\midrule
Item A & 10.5 & 20.3 & 30.8 \\
Item B & 15.2 & 25.1 & 40.3 \\
Item C & 22.1 & 12.7 & 34.8 \\
\bottomrule
\end{tabular}
\caption{Sample Data Table}
\label{tab:sample}
\end{table}

\section{Lists}
\subsection{Bullet Points}
\begin{itemize}
  \item First item
  \item Second item
  \item Third item with sub-items:
    \begin{itemize}
      \item Sub-item 1
      \item Sub-item 2
    \end{itemize}
\end{itemize}

\subsection{Numbered List}
\begin{enumerate}
  \item First step
  \item Second step
  \item Third step
\end{enumerate}

\section{Citations and References}
When you use the citation search feature, references will be automatically managed. For example, you might cite a paper like this: \cite{example2024}. All citations will be automatically added to your \texttt{cite.bib} file.

\section{Conclusion}
This template demonstrates basic \LaTeX{} features including mathematical formulas, tables, and lists. The bibliography system is now set up to work seamlessly with Bibby AI's citation management. Feel free to modify it for your needs!

\bibliographystyle{plain}
\bibliography{cite}

\end{document}\documentclass{article}
\usepackage[utf8]{inputenc}
\usepackage{booktabs}
\usepackage{amsmath}
\usepackage{natbib}

\title{My Document}
\author{Your Name}
\date{\today}

\begin{document}

\maketitle

\section{Introduction}
This is a sample document demonstrating various \LaTeX{} features including mathematical equations, tables, and formatting. You can cite papers using the citation search feature, and they will be automatically added to your \texttt{cite.bib} file.

\section{Mathematical Expressions}
Here are some example equations:

\subsection{Quadratic Formula}
The quadratic formula is given by:
\[x = \frac{-b \pm \sqrt{b^2 - 4ac}}{2a}\]

\subsection{Matrix Example}
A sample matrix:
\[
\begin{pmatrix}
a_{11} & a_{12} & a_{13} \\
a_{21} & a_{22} & a_{23} \\
a_{31} & a_{32} & a_{33}
\end{pmatrix}
\]

\section{Tables}
Below is a sample table showing data:

\begin{table}[h]
\centering
\begin{tabular}{@{}lrrr@{}}
\toprule
\textbf{Category} & \textbf{Value 1} & \textbf{Value 2} & \textbf{Total} \\
\midrule
Item A & 10.5 & 20.3 & 30.8 \\
Item B & 15.2 & 25.1 & 40.3 \\
Item C & 22.1 & 12.7 & 34.8 \\
\bottomrule
\end{tabular}
\caption{Sample Data Table}
\label{tab:sample}
\end{table}

\section{Lists}
\subsection{Bullet Points}
\begin{itemize}
  \item First item
  \item Second item
  \item Third item with sub-items:
    \begin{itemize}
      \item Sub-item 1
      \item Sub-item 2
    \end{itemize}
\end{itemize}

\subsection{Numbered List}
\begin{enumerate}
  \item First step
  \item Second step
  \item Third step
\end{enumerate}

\section{Citations and References}
When you use the citation search feature, references will be automatically managed. For example, you might cite a paper like this: \cite{example2024}. All citations will be automatically added to your \texttt{cite.bib} file.

\section{Conclusion}
This template demonstrates basic \LaTeX{} features including mathematical formulas, tables, and lists. The bibliography system is now set up to work seamlessly with Bibby AI's citation management. Feel free to modify it for your needs!

\bibliographystyle{plain}
\bibliography{cite}

\end{document}\documentclass{article}
\usepackage[utf8]{inputenc}
\usepackage{booktabs}
\usepackage{amsmath}
\usepackage{natbib}

\title{My Document}
\author{Your Name}
\date{\today}

\begin{document}

\maketitle

\section{Introduction}
This is a sample document demonstrating various \LaTeX{} features including mathematical equations, tables, and formatting. You can cite papers using the citation search feature, and they will be automatically added to your \texttt{cite.bib} file.

\section{Mathematical Expressions}
Here are some example equations:

\subsection{Quadratic Formula}
The quadratic formula is given by:
\[x = \frac{-b \pm \sqrt{b^2 - 4ac}}{2a}\]

\subsection{Matrix Example}
A sample matrix:
\[
\begin{pmatrix}
a_{11} & a_{12} & a_{13} \\
a_{21} & a_{22} & a_{23} \\
a_{31} & a_{32} & a_{33}
\end{pmatrix}
\]

\section{Tables}
Below is a sample table showing data:

\begin{table}[h]
\centering
\begin{tabular}{@{}lrrr@{}}
\toprule
\textbf{Category} & \textbf{Value 1} & \textbf{Value 2} & \textbf{Total} \\
\midrule
Item A & 10.5 & 20.3 & 30.8 \\
Item B & 15.2 & 25.1 & 40.3 \\
Item C & 22.1 & 12.7 & 34.8 \\
\bottomrule
\end{tabular}
\caption{Sample Data Table}
\label{tab:sample}
\end{table}

\section{Lists}
\subsection{Bullet Points}
\begin{itemize}
  \item First item
  \item Second item
  \item Third item with sub-items:
    \begin{itemize}
      \item Sub-item 1
      \item Sub-item 2
    \end{itemize}
\end{itemize}

\subsection{Numbered List}
\begin{enumerate}
  \item First step
  \item Second step
  \item Third step
\end{enumerate}

\section{Citations and References}
When you use the citation search feature, references will be automatically managed. For example, you might cite a paper like this: \cite{example2024}. All citations will be automatically added to your \texttt{cite.bib} file.

\section{Conclusion}
This template demonstrates basic \LaTeX{} features including mathematical formulas, tables, and lists. The bibliography system is now set up to work seamlessly with Bibby AI's citation management. Feel free to modify it for your needs!

\bibliographystyle{plain}
\bibliography{cite}

\end{document}\documentclass{article}
\usepackage[utf8]{inputenc}
\usepackage{booktabs}
\usepackage{amsmath}
\usepackage{natbib}

\title{My Document}
\author{Your Name}
\date{\today}

\begin{document}

\maketitle

\section{Introduction}
This is a sample document demonstrating various \LaTeX{} features including mathematical equations, tables, and formatting. You can cite papers using the citation search feature, and they will be automatically added to your \texttt{cite.bib} file.

\section{Mathematical Expressions}
Here are some example equations:

\subsection{Quadratic Formula}
The quadratic formula is given by:
\[x = \frac{-b \pm \sqrt{b^2 - 4ac}}{2a}\]

\subsection{Matrix Example}
A sample matrix:
\[
\begin{pmatrix}
a_{11} & a_{12} & a_{13} \\
a_{21} & a_{22} & a_{23} \\
a_{31} & a_{32} & a_{33}
\end{pmatrix}
\]

\section{Tables}
Below is a sample table showing data:

\begin{table}[h]
\centering
\begin{tabular}{@{}lrrr@{}}
\toprule
\textbf{Category} & \textbf{Value 1} & \textbf{Value 2} & \textbf{Total} \\
\midrule
Item A & 10.5 & 20.3 & 30.8 \\
Item B & 15.2 & 25.1 & 40.3 \\
Item C & 22.1 & 12.7 & 34.8 \\
\bottomrule
\end{tabular}
\caption{Sample Data Table}
\label{tab:sample}
\end{table}

\section{Lists}
\subsection{Bullet Points}
\begin{itemize}
  \item First item
  \item Second item
  \item Third item with sub-items:
    \begin{itemize}
      \item Sub-item 1
      \item Sub-item 2
    \end{itemize}
\end{itemize}

\subsection{Numbered List}
\begin{enumerate}
  \item First step
  \item Second step
  \item Third step
\end{enumerate}

\section{Citations and References}
When you use the citation search feature, references will be automatically managed. For example, you might cite a paper like this: \cite{example2024}. All citations will be automatically added to your \texttt{cite.bib} file.

\section{Conclusion}
This template demonstrates basic \LaTeX{} features including mathematical formulas, tables, and lists. The bibliography system is now set up to work seamlessly with Bibby AI's citation management. Feel free to modify it for your needs!

\bibliographystyle{plain}
\bibliography{cite}

\end{document}\documentclass{article}
\usepackage[utf8]{inputenc}
\usepackage{booktabs}
\usepackage{amsmath}
\usepackage{natbib}

\title{My Document}
\author{Your Name}
\date{\today}

\begin{document}

\maketitle

\section{Introduction}
This is a sample document demonstrating various \LaTeX{} features including mathematical equations, tables, and formatting. You can cite papers using the citation search feature, and they will be automatically added to your \texttt{cite.bib} file.

\section{Mathematical Expressions}
Here are some example equations:

\subsection{Quadratic Formula}
The quadratic formula is given by:
\[x = \frac{-b \pm \sqrt{b^2 - 4ac}}{2a}\]

\subsection{Matrix Example}
A sample matrix:
\[
\begin{pmatrix}
a_{11} & a_{12} & a_{13} \\
a_{21} & a_{22} & a_{23} \\
a_{31} & a_{32} & a_{33}
\end{pmatrix}
\]

\section{Tables}
Below is a sample table showing data:

\begin{table}[h]
\centering
\begin{tabular}{@{}lrrr@{}}
\toprule
\textbf{Category} & \textbf{Value 1} & \textbf{Value 2} & \textbf{Total} \\
\midrule
Item A & 10.5 & 20.3 & 30.8 \\
Item B & 15.2 & 25.1 & 40.3 \\
Item C & 22.1 & 12.7 & 34.8 \\
\bottomrule
\end{tabular}
\caption{Sample Data Table}
\label{tab:sample}
\end{table}

\section{Lists}
\subsection{Bullet Points}
\begin{itemize}
  \item First item
  \item Second item
  \item Third item with sub-items:
    \begin{itemize}
      \item Sub-item 1
      \item Sub-item 2
    \end{itemize}
\end{itemize}

\subsection{Numbered List}
\begin{enumerate}
  \item First step
  \item Second step
  \item Third step
\end{enumerate}

\section{Citations and References}
When you use the citation search feature, references will be automatically managed. For example, you might cite a paper like this: \cite{example2024}. All citations will be automatically added to your \texttt{cite.bib} file.

\section{Conclusion}
This template demonstrates basic \LaTeX{} features including mathematical formulas, tables, and lists. The bibliography system is now set up to work seamlessly with Bibby AI's citation management. Feel free to modify it for your needs!

\bibliographystyle{plain}
\bibliography{cite}

\end{document}\documentclass{article}
\usepackage[utf8]{inputenc}
\usepackage{booktabs}
\usepackage{amsmath}
\usepackage{natbib}

\title{My Document}
\author{Your Name}
\date{\today}

\begin{document}

\maketitle

\section{Introduction}
This is a sample document demonstrating various \LaTeX{} features including mathematical equations, tables, and formatting. You can cite papers using the citation search feature, and they will be automatically added to your \texttt{cite.bib} file.

\section{Mathematical Expressions}
Here are some example equations:

\subsection{Quadratic Formula}
The quadratic formula is given by:
\[x = \frac{-b \pm \sqrt{b^2 - 4ac}}{2a}\]

\subsection{Matrix Example}
A sample matrix:
\[
\begin{pmatrix}
a_{11} & a_{12} & a_{13} \\
a_{21} & a_{22} & a_{23} \\
a_{31} & a_{32} & a_{33}
\end{pmatrix}
\]

\section{Tables}
Below is a sample table showing data:

\begin{table}[h]
\centering
\begin{tabular}{@{}lrrr@{}}
\toprule
\textbf{Category} & \textbf{Value 1} & \textbf{Value 2} & \textbf{Total} \\
\midrule
Item A & 10.5 & 20.3 & 30.8 \\
Item B & 15.2 & 25.1 & 40.3 \\
Item C & 22.1 & 12.7 & 34.8 \\
\bottomrule
\end{tabular}
\caption{Sample Data Table}
\label{tab:sample}
\end{table}

\section{Lists}
\subsection{Bullet Points}
\begin{itemize}
  \item First item
  \item Second item
  \item Third item with sub-items:
    \begin{itemize}
      \item Sub-item 1
      \item Sub-item 2
    \end{itemize}
\end{itemize}

\subsection{Numbered List}
\begin{enumerate}
  \item First step
  \item Second step
  \item Third step
\end{enumerate}

\section{Citations and References}
When you use the citation search feature, references will be automatically managed. For example, you might cite a paper like this: \cite{example2024}. All citations will be automatically added to your \texttt{cite.bib} file.

\section{Conclusion}
This template demonstrates basic \LaTeX{} features including mathematical formulas, tables, and lists. The bibliography system is now set up to work seamlessly with Bibby AI's citation management. Feel free to modify it for your needs!

\bibliographystyle{plain}
\bibliography{cite}

\end{document}\documentclass{article}
\usepackage[utf8]{inputenc}
\usepackage{booktabs}
\usepackage{amsmath}
\usepackage{natbib}

\title{My Document}
\author{Your Name}
\date{\today}

\begin{document}

\maketitle

\section{Introduction}
This is a sample document demonstrating various \LaTeX{} features including mathematical equations, tables, and formatting. You can cite papers using the citation search feature, and they will be automatically added to your \texttt{cite.bib} file.

\section{Mathematical Expressions}
Here are some example equations:

\subsection{Quadratic Formula}
The quadratic formula is given by:
\[x = \frac{-b \pm \sqrt{b^2 - 4ac}}{2a}\]

\subsection{Matrix Example}
A sample matrix:
\[
\begin{pmatrix}
a_{11} & a_{12} & a_{13} \\
a_{21} & a_{22} & a_{23} \\
a_{31} & a_{32} & a_{33}
\end{pmatrix}
\]

\section{Tables}
Below is a sample table showing data:

\begin{table}[h]
\centering
\begin{tabular}{@{}lrrr@{}}
\toprule
\textbf{Category} & \textbf{Value 1} & \textbf{Value 2} & \textbf{Total} \\
\midrule
Item A & 10.5 & 20.3 & 30.8 \\
Item B & 15.2 & 25.1 & 40.3 \\
Item C & 22.1 & 12.7 & 34.8 \\
\bottomrule
\end{tabular}
\caption{Sample Data Table}
\label{tab:sample}
\end{table}

\section{Lists}
\subsection{Bullet Points}
\begin{itemize}
  \item First item
  \item Second item
  \item Third item with sub-items:
    \begin{itemize}
      \item Sub-item 1
      \item Sub-item 2
    \end{itemize}
\end{itemize}

\subsection{Numbered List}
\begin{enumerate}
  \item First step
  \item Second step
  \item Third step
\end{enumerate}

\section{Citations and References}
When you use the citation search feature, references will be automatically managed. For example, you might cite a paper like this: \cite{example2024}. All citations will be automatically added to your \texttt{cite.bib} file.

\section{Conclusion}
This template demonstrates basic \LaTeX{} features including mathematical formulas, tables, and lists. The bibliography system is now set up to work seamlessly with Bibby AI's citation management. Feel free to modify it for your needs!

\bibliographystyle{plain}
\bibliography{cite}

\end{document}\documentclass{article}
\usepackage[utf8]{inputenc}
\usepackage{booktabs}
\usepackage{amsmath}
\usepackage{natbib}

\title{My Document}
\author{Your Name}
\date{\today}

\begin{document}

\maketitle

\section{Introduction}
This is a sample document demonstrating various \LaTeX{} features including mathematical equations, tables, and formatting. You can cite papers using the citation search feature, and they will be automatically added to your \texttt{cite.bib} file.

\section{Mathematical Expressions}
Here are some example equations:

\subsection{Quadratic Formula}
The quadratic formula is given by:
\[x = \frac{-b \pm \sqrt{b^2 - 4ac}}{2a}\]

\subsection{Matrix Example}
A sample matrix:
\[
\begin{pmatrix}
a_{11} & a_{12} & a_{13} \\
a_{21} & a_{22} & a_{23} \\
a_{31} & a_{32} & a_{33}
\end{pmatrix}
\]

\section{Tables}
Below is a sample table showing data:

\begin{table}[h]
\centering
\begin{tabular}{@{}lrrr@{}}
\toprule
\textbf{Category} & \textbf{Value 1} & \textbf{Value 2} & \textbf{Total} \\
\midrule
Item A & 10.5 & 20.3 & 30.8 \\
Item B & 15.2 & 25.1 & 40.3 \\
Item C & 22.1 & 12.7 & 34.8 \\
\bottomrule
\end{tabular}
\caption{Sample Data Table}
\label{tab:sample}
\end{table}

\section{Lists}
\subsection{Bullet Points}
\begin{itemize}
  \item First item
  \item Second item
  \item Third item with sub-items:
    \begin{itemize}
      \item Sub-item 1
      \item Sub-item 2
    \end{itemize}
\end{itemize}

\subsection{Numbered List}
\begin{enumerate}
  \item First step
  \item Second step
  \item Third step
\end{enumerate}

\section{Citations and References}
When you use the citation search feature, references will be automatically managed. For example, you might cite a paper like this: \cite{example2024}. All citations will be automatically added to your \texttt{cite.bib} file.

\section{Conclusion}
This template demonstrates basic \LaTeX{} features including mathematical formulas, tables, and lists. The bibliography system is now set up to work seamlessly with Bibby AI's citation management. Feel free to modify it for your needs!

\bibliographystyle{plain}
\bibliography{cite}

\end{document}\documentclass{article}
\usepackage[utf8]{inputenc}
\usepackage{booktabs}
\usepackage{amsmath}
\usepackage{natbib}

\title{My Document}
\author{Your Name}
\date{\today}

\begin{document}

\maketitle

\section{Introduction}
This is a sample document demonstrating various \LaTeX{} features including mathematical equations, tables, and formatting. You can cite papers using the citation search feature, and they will be automatically added to your \texttt{cite.bib} file.

\section{Mathematical Expressions}
Here are some example equations:

\subsection{Quadratic Formula}
The quadratic formula is given by:
\[x = \frac{-b \pm \sqrt{b^2 - 4ac}}{2a}\]

\subsection{Matrix Example}
A sample matrix:
\[
\begin{pmatrix}
a_{11} & a_{12} & a_{13} \\
a_{21} & a_{22} & a_{23} \\
a_{31} & a_{32} & a_{33}
\end{pmatrix}
\]

\section{Tables}
Below is a sample table showing data:

\begin{table}[h]
\centering
\begin{tabular}{@{}lrrr@{}}
\toprule
\textbf{Category} & \textbf{Value 1} & \textbf{Value 2} & \textbf{Total} \\
\midrule
Item A & 10.5 & 20.3 & 30.8 \\
Item B & 15.2 & 25.1 & 40.3 \\
Item C & 22.1 & 12.7 & 34.8 \\
\bottomrule
\end{tabular}
\caption{Sample Data Table}
\label{tab:sample}
\end{table}

\section{Lists}
\subsection{Bullet Points}
\begin{itemize}
  \item First item
  \item Second item
  \item Third item with sub-items:
    \begin{itemize}
      \item Sub-item 1
      \item Sub-item 2
    \end{itemize}
\end{itemize}

\subsection{Numbered List}
\begin{enumerate}
  \item First step
  \item Second step
  \item Third step
\end{enumerate}

\section{Citations and References}
When you use the citation search feature, references will be automatically managed. For example, you might cite a paper like this: \cite{example2024}. All citations will be automatically added to your \texttt{cite.bib} file.

\section{Conclusion}
This template demonstrates basic \LaTeX{} features including mathematical formulas, tables, and lists. The bibliography system is now set up to work seamlessly with Bibby AI's citation management. Feel free to modify it for your needs!

\bibliographystyle{plain}
\bibliography{cite}

\end{document}\documentclass{article}
\usepackage[utf8]{inputenc}
\usepackage{booktabs}
\usepackage{amsmath}
\usepackage{natbib}

\title{My Document}
\author{Your Name}
\date{\today}

\begin{document}

\maketitle

\section{Introduction}
This is a sample document demonstrating various \LaTeX{} features including mathematical equations, tables, and formatting. You can cite papers using the citation search feature, and they will be automatically added to your \texttt{cite.bib} file.

\section{Mathematical Expressions}
Here are some example equations:

\subsection{Quadratic Formula}
The quadratic formula is given by:
\[x = \frac{-b \pm \sqrt{b^2 - 4ac}}{2a}\]

\subsection{Matrix Example}
A sample matrix:
\[
\begin{pmatrix}
a_{11} & a_{12} & a_{13} \\
a_{21} & a_{22} & a_{23} \\
a_{31} & a_{32} & a_{33}
\end{pmatrix}
\]

\section{Tables}
Below is a sample table showing data:

\begin{table}[h]
\centering
\begin{tabular}{@{}lrrr@{}}
\toprule
\textbf{Category} & \textbf{Value 1} & \textbf{Value 2} & \textbf{Total} \\
\midrule
Item A & 10.5 & 20.3 & 30.8 \\
Item B & 15.2 & 25.1 & 40.3 \\
Item C & 22.1 & 12.7 & 34.8 \\
\bottomrule
\end{tabular}
\caption{Sample Data Table}
\label{tab:sample}
\end{table}

\section{Lists}
\subsection{Bullet Points}
\begin{itemize}
  \item First item
  \item Second item
  \item Third item with sub-items:
    \begin{itemize}
      \item Sub-item 1
      \item Sub-item 2
    \end{itemize}
\end{itemize}

\subsection{Numbered List}
\begin{enumerate}
  \item First step
  \item Second step
  \item Third step
\end{enumerate}

\section{Citations and References}
When you use the citation search feature, references will be automatically managed. For example, you might cite a paper like this: \cite{example2024}. All citations will be automatically added to your \texttt{cite.bib} file.

\section{Conclusion}
This template demonstrates basic \LaTeX{} features including mathematical formulas, tables, and lists. The bibliography system is now set up to work seamlessly with Bibby AI's citation management. Feel free to modify it for your needs!

\bibliographystyle{plain}
\bibliography{cite}

\end{document}\documentclass{article}
\usepackage[utf8]{inputenc}
\usepackage{booktabs}
\usepackage{amsmath}
\usepackage{natbib}

\title{My Document}
\author{Your Name}
\date{\today}

\begin{document}

\maketitle

\section{Introduction}
This is a sample document demonstrating various \LaTeX{} features including mathematical equations, tables, and formatting. You can cite papers using the citation search feature, and they will be automatically added to your \texttt{cite.bib} file.

\section{Mathematical Expressions}
Here are some example equations:

\subsection{Quadratic Formula}
The quadratic formula is given by:
\[x = \frac{-b \pm \sqrt{b^2 - 4ac}}{2a}\]

\subsection{Matrix Example}
A sample matrix:
\[
\begin{pmatrix}
a_{11} & a_{12} & a_{13} \\
a_{21} & a_{22} & a_{23} \\
a_{31} & a_{32} & a_{33}
\end{pmatrix}
\]

\section{Tables}
Below is a sample table showing data:

\begin{table}[h]
\centering
\begin{tabular}{@{}lrrr@{}}
\toprule
\textbf{Category} & \textbf{Value 1} & \textbf{Value 2} & \textbf{Total} \\
\midrule
Item A & 10.5 & 20.3 & 30.8 \\
Item B & 15.2 & 25.1 & 40.3 \\
Item C & 22.1 & 12.7 & 34.8 \\
\bottomrule
\end{tabular}
\caption{Sample Data Table}
\label{tab:sample}
\end{table}

\section{Lists}
\subsection{Bullet Points}
\begin{itemize}
  \item First item
  \item Second item
  \item Third item with sub-items:
    \begin{itemize}
      \item Sub-item 1
      \item Sub-item 2
    \end{itemize}
\end{itemize}

\subsection{Numbered List}
\begin{enumerate}
  \item First step
  \item Second step
  \item Third step
\end{enumerate}

\section{Citations and References}
When you use the citation search feature, references will be automatically managed. For example, you might cite a paper like this: \cite{example2024}. All citations will be automatically added to your \texttt{cite.bib} file.

\section{Conclusion}
This template demonstrates basic \LaTeX{} features including mathematical formulas, tables, and lists. The bibliography system is now set up to work seamlessly with Bibby AI's citation management. Feel free to modify it for your needs!

\bibliographystyle{plain}
\bibliography{cite}

\end{document}\documentclass{article}
\usepackage[utf8]{inputenc}
\usepackage{booktabs}
\usepackage{amsmath}
\usepackage{natbib}

\title{My Document}
\author{Your Name}
\date{\today}

\begin{document}

\maketitle

\section{Introduction}
This is a sample document demonstrating various \LaTeX{} features including mathematical equations, tables, and formatting. You can cite papers using the citation search feature, and they will be automatically added to your \texttt{cite.bib} file.

\section{Mathematical Expressions}
Here are some example equations:

\subsection{Quadratic Formula}
The quadratic formula is given by:
\[x = \frac{-b \pm \sqrt{b^2 - 4ac}}{2a}\]

\subsection{Matrix Example}
A sample matrix:
\[
\begin{pmatrix}
a_{11} & a_{12} & a_{13} \\
a_{21} & a_{22} & a_{23} \\
a_{31} & a_{32} & a_{33}
\end{pmatrix}
\]

\section{Tables}
Below is a sample table showing data:

\begin{table}[h]
\centering
\begin{tabular}{@{}lrrr@{}}
\toprule
\textbf{Category} & \textbf{Value 1} & \textbf{Value 2} & \textbf{Total} \\
\midrule
Item A & 10.5 & 20.3 & 30.8 \\
Item B & 15.2 & 25.1 & 40.3 \\
Item C & 22.1 & 12.7 & 34.8 \\
\bottomrule
\end{tabular}
\caption{Sample Data Table}
\label{tab:sample}
\end{table}

\section{Lists}
\subsection{Bullet Points}
\begin{itemize}
  \item First item
  \item Second item
  \item Third item with sub-items:
    \begin{itemize}
      \item Sub-item 1
      \item Sub-item 2
    \end{itemize}
\end{itemize}

\subsection{Numbered List}
\begin{enumerate}
  \item First step
  \item Second step
  \item Third step
\end{enumerate}

\section{Citations and References}
When you use the citation search feature, references will be automatically managed. For example, you might cite a paper like this: \cite{example2024}. All citations will be automatically added to your \texttt{cite.bib} file.

\section{Conclusion}
This template demonstrates basic \LaTeX{} features including mathematical formulas, tables, and lists. The bibliography system is now set up to work seamlessly with Bibby AI's citation management. Feel free to modify it for your needs!

\bibliographystyle{plain}
\bibliography{cite}

\end{document}\documentclass{article}
\usepackage[utf8]{inputenc}
\usepackage{booktabs}
\usepackage{amsmath}
\usepackage{natbib}

\title{My Document}
\author{Your Name}
\date{\today}

\begin{document}

\maketitle

\section{Introduction}
This is a sample document demonstrating various \LaTeX{} features including mathematical equations, tables, and formatting. You can cite papers using the citation search feature, and they will be automatically added to your \texttt{cite.bib} file.

\section{Mathematical Expressions}
Here are some example equations:

\subsection{Quadratic Formula}
The quadratic formula is given by:
\[x = \frac{-b \pm \sqrt{b^2 - 4ac}}{2a}\]

\subsection{Matrix Example}
A sample matrix:
\[
\begin{pmatrix}
a_{11} & a_{12} & a_{13} \\
a_{21} & a_{22} & a_{23} \\
a_{31} & a_{32} & a_{33}
\end{pmatrix}
\]

\section{Tables}
Below is a sample table showing data:

\begin{table}[h]
\centering
\begin{tabular}{@{}lrrr@{}}
\toprule
\textbf{Category} & \textbf{Value 1} & \textbf{Value 2} & \textbf{Total} \\
\midrule
Item A & 10.5 & 20.3 & 30.8 \\
Item B & 15.2 & 25.1 & 40.3 \\
Item C & 22.1 & 12.7 & 34.8 \\
\bottomrule
\end{tabular}
\caption{Sample Data Table}
\label{tab:sample}
\end{table}

\section{Lists}
\subsection{Bullet Points}
\begin{itemize}
  \item First item
  \item Second item
  \item Third item with sub-items:
    \begin{itemize}
      \item Sub-item 1
      \item Sub-item 2
    \end{itemize}
\end{itemize}

\subsection{Numbered List}
\begin{enumerate}
  \item First step
  \item Second step
  \item Third step
\end{enumerate}

\section{Citations and References}
When you use the citation search feature, references will be automatically managed. For example, you might cite a paper like this: \cite{example2024}. All citations will be automatically added to your \texttt{cite.bib} file.

\section{Conclusion}
This template demonstrates basic \LaTeX{} features including mathematical formulas, tables, and lists. The bibliography system is now set up to work seamlessly with Bibby AI's citation management. Feel free to modify it for your needs!

\bibliographystyle{plain}
\bibliography{cite}

\end{document}\documentclass{article}
\usepackage[utf8]{inputenc}
\usepackage{booktabs}
\usepackage{amsmath}
\usepackage{natbib}

\title{My Document}
\author{Your Name}
\date{\today}

\begin{document}

\maketitle

\section{Introduction}
This is a sample document demonstrating various \LaTeX{} features including mathematical equations, tables, and formatting. You can cite papers using the citation search feature, and they will be automatically added to your \texttt{cite.bib} file.

\section{Mathematical Expressions}
Here are some example equations:

\subsection{Quadratic Formula}
The quadratic formula is given by:
\[x = \frac{-b \pm \sqrt{b^2 - 4ac}}{2a}\]

\subsection{Matrix Example}
A sample matrix:
\[
\begin{pmatrix}
a_{11} & a_{12} & a_{13} \\
a_{21} & a_{22} & a_{23} \\
a_{31} & a_{32} & a_{33}
\end{pmatrix}
\]

\section{Tables}
Below is a sample table showing data:

\begin{table}[h]
\centering
\begin{tabular}{@{}lrrr@{}}
\toprule
\textbf{Category} & \textbf{Value 1} & \textbf{Value 2} & \textbf{Total} \\
\midrule
Item A & 10.5 & 20.3 & 30.8 \\
Item B & 15.2 & 25.1 & 40.3 \\
Item C & 22.1 & 12.7 & 34.8 \\
\bottomrule
\end{tabular}
\caption{Sample Data Table}
\label{tab:sample}
\end{table}

\section{Lists}
\subsection{Bullet Points}
\begin{itemize}
  \item First item
  \item Second item
  \item Third item with sub-items:
    \begin{itemize}
      \item Sub-item 1
      \item Sub-item 2
    \end{itemize}
\end{itemize}

\subsection{Numbered List}
\begin{enumerate}
  \item First step
  \item Second step
  \item Third step
\end{enumerate}

\section{Citations and References}
When you use the citation search feature, references will be automatically managed. For example, you might cite a paper like this: \cite{example2024}. All citations will be automatically added to your \texttt{cite.bib} file.

\section{Conclusion}
This template demonstrates basic \LaTeX{} features including mathematical formulas, tables, and lists. The bibliography system is now set up to work seamlessly with Bibby AI's citation management. Feel free to modify it for your needs!

\bibliographystyle{plain}
\bibliography{cite}

\end{document}\documentclass{article}
\usepackage[utf8]{inputenc}
\usepackage{booktabs}
\usepackage{amsmath}
\usepackage{natbib}

\title{My Document}
\author{Your Name}
\date{\today}

\begin{document}

\maketitle

\section{Introduction}
This is a sample document demonstrating various \LaTeX{} features including mathematical equations, tables, and formatting. You can cite papers using the citation search feature, and they will be automatically added to your \texttt{cite.bib} file.

\section{Mathematical Expressions}
Here are some example equations:

\subsection{Quadratic Formula}
The quadratic formula is given by:
\[x = \frac{-b \pm \sqrt{b^2 - 4ac}}{2a}\]

\subsection{Matrix Example}
A sample matrix:
\[
\begin{pmatrix}
a_{11} & a_{12} & a_{13} \\
a_{21} & a_{22} & a_{23} \\
a_{31} & a_{32} & a_{33}
\end{pmatrix}
\]

\section{Tables}
Below is a sample table showing data:

\begin{table}[h]
\centering
\begin{tabular}{@{}lrrr@{}}
\toprule
\textbf{Category} & \textbf{Value 1} & \textbf{Value 2} & \textbf{Total} \\
\midrule
Item A & 10.5 & 20.3 & 30.8 \\
Item B & 15.2 & 25.1 & 40.3 \\
Item C & 22.1 & 12.7 & 34.8 \\
\bottomrule
\end{tabular}
\caption{Sample Data Table}
\label{tab:sample}
\end{table}

\section{Lists}
\subsection{Bullet Points}
\begin{itemize}
  \item First item
  \item Second item
  \item Third item with sub-items:
    \begin{itemize}
      \item Sub-item 1
      \item Sub-item 2
    \end{itemize}
\end{itemize}

\subsection{Numbered List}
\begin{enumerate}
  \item First step
  \item Second step
  \item Third step
\end{enumerate}

\section{Citations and References}
When you use the citation search feature, references will be automatically managed. For example, you might cite a paper like this: \cite{example2024}. All citations will be automatically added to your \texttt{cite.bib} file.

\section{Conclusion}
This template demonstrates basic \LaTeX{} features including mathematical formulas, tables, and lists. The bibliography system is now set up to work seamlessly with Bibby AI's citation management. Feel free to modify it for your needs!

\bibliographystyle{plain}
\bibliography{cite}

\end{document}\documentclass{article}
\usepackage[utf8]{inputenc}
\usepackage{booktabs}
\usepackage{amsmath}
\usepackage{natbib}

\title{My Document}
\author{Your Name}
\date{\today}

\begin{document}

\maketitle

\section{Introduction}
This is a sample document demonstrating various \LaTeX{} features including mathematical equations, tables, and formatting. You can cite papers using the citation search feature, and they will be automatically added to your \texttt{cite.bib} file.

\section{Mathematical Expressions}
Here are some example equations:

\subsection{Quadratic Formula}
The quadratic formula is given by:
\[x = \frac{-b \pm \sqrt{b^2 - 4ac}}{2a}\]

\subsection{Matrix Example}
A sample matrix:
\[
\begin{pmatrix}
a_{11} & a_{12} & a_{13} \\
a_{21} & a_{22} & a_{23} \\
a_{31} & a_{32} & a_{33}
\end{pmatrix}
\]

\section{Tables}
Below is a sample table showing data:

\begin{table}[h]
\centering
\begin{tabular}{@{}lrrr@{}}
\toprule
\textbf{Category} & \textbf{Value 1} & \textbf{Value 2} & \textbf{Total} \\
\midrule
Item A & 10.5 & 20.3 & 30.8 \\
Item B & 15.2 & 25.1 & 40.3 \\
Item C & 22.1 & 12.7 & 34.8 \\
\bottomrule
\end{tabular}
\caption{Sample Data Table}
\label{tab:sample}
\end{table}

\section{Lists}
\subsection{Bullet Points}
\begin{itemize}
  \item First item
  \item Second item
  \item Third item with sub-items:
    \begin{itemize}
      \item Sub-item 1
      \item Sub-item 2
    \end{itemize}
\end{itemize}

\subsection{Numbered List}
\begin{enumerate}
  \item First step
  \item Second step
  \item Third step
\end{enumerate}

\section{Citations and References}
When you use the citation search feature, references will be automatically managed. For example, you might cite a paper like this: \cite{example2024}. All citations will be automatically added to your \texttt{cite.bib} file.

\section{Conclusion}
This template demonstrates basic \LaTeX{} features including mathematical formulas, tables, and lists. The bibliography system is now set up to work seamlessly with Bibby AI's citation management. Feel free to modify it for your needs!

\bibliographystyle{plain}
\bibliography{cite}

\end{document}\documentclass{article}
\usepackage[utf8]{inputenc}
\usepackage{booktabs}
\usepackage{amsmath}
\usepackage{natbib}

\title{My Document}
\author{Your Name}
\date{\today}

\begin{document}

\maketitle

\section{Introduction}
This is a sample document demonstrating various \LaTeX{} features including mathematical equations, tables, and formatting. You can cite papers using the citation search feature, and they will be automatically added to your \texttt{cite.bib} file.

\section{Mathematical Expressions}
Here are some example equations:

\subsection{Quadratic Formula}
The quadratic formula is given by:
\[x = \frac{-b \pm \sqrt{b^2 - 4ac}}{2a}\]

\subsection{Matrix Example}
A sample matrix:
\[
\begin{pmatrix}
a_{11} & a_{12} & a_{13} \\
a_{21} & a_{22} & a_{23} \\
a_{31} & a_{32} & a_{33}
\end{pmatrix}
\]

\section{Tables}
Below is a sample table showing data:

\begin{table}[h]
\centering
\begin{tabular}{@{}lrrr@{}}
\toprule
\textbf{Category} & \textbf{Value 1} & \textbf{Value 2} & \textbf{Total} \\
\midrule
Item A & 10.5 & 20.3 & 30.8 \\
Item B & 15.2 & 25.1 & 40.3 \\
Item C & 22.1 & 12.7 & 34.8 \\
\bottomrule
\end{tabular}
\caption{Sample Data Table}
\label{tab:sample}
\end{table}

\section{Lists}
\subsection{Bullet Points}
\begin{itemize}
  \item First item
  \item Second item
  \item Third item with sub-items:
    \begin{itemize}
      \item Sub-item 1
      \item Sub-item 2
    \end{itemize}
\end{itemize}

\subsection{Numbered List}
\begin{enumerate}
  \item First step
  \item Second step
  \item Third step
\end{enumerate}

\section{Citations and References}
When you use the citation search feature, references will be automatically managed. For example, you might cite a paper like this: \cite{example2024}. All citations will be automatically added to your \texttt{cite.bib} file.

\section{Conclusion}
This template demonstrates basic \LaTeX{} features including mathematical formulas, tables, and lists. The bibliography system is now set up to work seamlessly with Bibby AI's citation management. Feel free to modify it for your needs!

\bibliographystyle{plain}
\bibliography{cite}

\end{document}\documentclass{article}
\usepackage[utf8]{inputenc}
\usepackage{booktabs}
\usepackage{amsmath}
\usepackage{natbib}

\title{My Document}
\author{Your Name}
\date{\today}

\begin{document}

\maketitle

\section{Introduction}
This is a sample document demonstrating various \LaTeX{} features including mathematical equations, tables, and formatting. You can cite papers using the citation search feature, and they will be automatically added to your \texttt{cite.bib} file.

\section{Mathematical Expressions}
Here are some example equations:

\subsection{Quadratic Formula}
The quadratic formula is given by:
\[x = \frac{-b \pm \sqrt{b^2 - 4ac}}{2a}\]

\subsection{Matrix Example}
A sample matrix:
\[
\begin{pmatrix}
a_{11} & a_{12} & a_{13} \\
a_{21} & a_{22} & a_{23} \\
a_{31} & a_{32} & a_{33}
\end{pmatrix}
\]

\section{Tables}
Below is a sample table showing data:

\begin{table}[h]
\centering
\begin{tabular}{@{}lrrr@{}}
\toprule
\textbf{Category} & \textbf{Value 1} & \textbf{Value 2} & \textbf{Total} \\
\midrule
Item A & 10.5 & 20.3 & 30.8 \\
Item B & 15.2 & 25.1 & 40.3 \\
Item C & 22.1 & 12.7 & 34.8 \\
\bottomrule
\end{tabular}
\caption{Sample Data Table}
\label{tab:sample}
\end{table}

\section{Lists}
\subsection{Bullet Points}
\begin{itemize}
  \item First item
  \item Second item
  \item Third item with sub-items:
    \begin{itemize}
      \item Sub-item 1
      \item Sub-item 2
    \end{itemize}
\end{itemize}

\subsection{Numbered List}
\begin{enumerate}
  \item First step
  \item Second step
  \item Third step
\end{enumerate}

\section{Citations and References}
When you use the citation search feature, references will be automatically managed. For example, you might cite a paper like this: \cite{example2024}. All citations will be automatically added to your \texttt{cite.bib} file.

\section{Conclusion}
This template demonstrates basic \LaTeX{} features including mathematical formulas, tables, and lists. The bibliography system is now set up to work seamlessly with Bibby AI's citation management. Feel free to modify it for your needs!

\bibliographystyle{plain}
\bibliography{cite}

\end{document}\documentclass{article}
\usepackage[utf8]{inputenc}
\usepackage{booktabs}
\usepackage{amsmath}
\usepackage{natbib}

\title{My Document}
\author{Your Name}
\date{\today}

\begin{document}

\maketitle

\section{Introduction}
This is a sample document demonstrating various \LaTeX{} features including mathematical equations, tables, and formatting. You can cite papers using the citation search feature, and they will be automatically added to your \texttt{cite.bib} file.

\section{Mathematical Expressions}
Here are some example equations:

\subsection{Quadratic Formula}
The quadratic formula is given by:
\[x = \frac{-b \pm \sqrt{b^2 - 4ac}}{2a}\]

\subsection{Matrix Example}
A sample matrix:
\[
\begin{pmatrix}
a_{11} & a_{12} & a_{13} \\
a_{21} & a_{22} & a_{23} \\
a_{31} & a_{32} & a_{33}
\end{pmatrix}
\]

\section{Tables}
Below is a sample table showing data:

\begin{table}[h]
\centering
\begin{tabular}{@{}lrrr@{}}
\toprule
\textbf{Category} & \textbf{Value 1} & \textbf{Value 2} & \textbf{Total} \\
\midrule
Item A & 10.5 & 20.3 & 30.8 \\
Item B & 15.2 & 25.1 & 40.3 \\
Item C & 22.1 & 12.7 & 34.8 \\
\bottomrule
\end{tabular}
\caption{Sample Data Table}
\label{tab:sample}
\end{table}

\section{Lists}
\subsection{Bullet Points}
\begin{itemize}
  \item First item
  \item Second item
  \item Third item with sub-items:
    \begin{itemize}
      \item Sub-item 1
      \item Sub-item 2
    \end{itemize}
\end{itemize}

\subsection{Numbered List}
\begin{enumerate}
  \item First step
  \item Second step
  \item Third step
\end{enumerate}

\section{Citations and References}
When you use the citation search feature, references will be automatically managed. For example, you might cite a paper like this: \cite{example2024}. All citations will be automatically added to your \texttt{cite.bib} file.

\section{Conclusion}
This template demonstrates basic \LaTeX{} features including mathematical formulas, tables, and lists. The bibliography system is now set up to work seamlessly with Bibby AI's citation management. Feel free to modify it for your needs!

\bibliographystyle{plain}
\bibliography{cite}

\end{document}\documentclass{article}
\usepackage[utf8]{inputenc}
\usepackage{booktabs}
\usepackage{amsmath}
\usepackage{natbib}

\title{My Document}
\author{Your Name}
\date{\today}

\begin{document}

\maketitle

\section{Introduction}
This is a sample document demonstrating various \LaTeX{} features including mathematical equations, tables, and formatting. You can cite papers using the citation search feature, and they will be automatically added to your \texttt{cite.bib} file.

\section{Mathematical Expressions}
Here are some example equations:

\subsection{Quadratic Formula}
The quadratic formula is given by:
\[x = \frac{-b \pm \sqrt{b^2 - 4ac}}{2a}\]

\subsection{Matrix Example}
A sample matrix:
\[
\begin{pmatrix}
a_{11} & a_{12} & a_{13} \\
a_{21} & a_{22} & a_{23} \\
a_{31} & a_{32} & a_{33}
\end{pmatrix}
\]

\section{Tables}
Below is a sample table showing data:

\begin{table}[h]
\centering
\begin{tabular}{@{}lrrr@{}}
\toprule
\textbf{Category} & \textbf{Value 1} & \textbf{Value 2} & \textbf{Total} \\
\midrule
Item A & 10.5 & 20.3 & 30.8 \\
Item B & 15.2 & 25.1 & 40.3 \\
Item C & 22.1 & 12.7 & 34.8 \\
\bottomrule
\end{tabular}
\caption{Sample Data Table}
\label{tab:sample}
\end{table}

\section{Lists}
\subsection{Bullet Points}
\begin{itemize}
  \item First item
  \item Second item
  \item Third item with sub-items:
    \begin{itemize}
      \item Sub-item 1
      \item Sub-item 2
    \end{itemize}
\end{itemize}

\subsection{Numbered List}
\begin{enumerate}
  \item First step
  \item Second step
  \item Third step
\end{enumerate}

\section{Citations and References}
When you use the citation search feature, references will be automatically managed. For example, you might cite a paper like this: \cite{example2024}. All citations will be automatically added to your \texttt{cite.bib} file.

\section{Conclusion}
This template demonstrates basic \LaTeX{} features including mathematical formulas, tables, and lists. The bibliography system is now set up to work seamlessly with Bibby AI's citation management. Feel free to modify it for your needs!

\bibliographystyle{plain}
\bibliography{cite}

\end{document}\documentclass{article}
\usepackage[utf8]{inputenc}
\usepackage{booktabs}
\usepackage{amsmath}
\usepackage{natbib}

\title{My Document}
\author{Your Name}
\date{\today}

\begin{document}

\maketitle

\section{Introduction}
This is a sample document demonstrating various \LaTeX{} features including mathematical equations, tables, and formatting. You can cite papers using the citation search feature, and they will be automatically added to your \texttt{cite.bib} file.

\section{Mathematical Expressions}
Here are some example equations:

\subsection{Quadratic Formula}
The quadratic formula is given by:
\[x = \frac{-b \pm \sqrt{b^2 - 4ac}}{2a}\]

\subsection{Matrix Example}
A sample matrix:
\[
\begin{pmatrix}
a_{11} & a_{12} & a_{13} \\
a_{21} & a_{22} & a_{23} \\
a_{31} & a_{32} & a_{33}
\end{pmatrix}
\]

\section{Tables}
Below is a sample table showing data:

\begin{table}[h]
\centering
\begin{tabular}{@{}lrrr@{}}
\toprule
\textbf{Category} & \textbf{Value 1} & \textbf{Value 2} & \textbf{Total} \\
\midrule
Item A & 10.5 & 20.3 & 30.8 \\
Item B & 15.2 & 25.1 & 40.3 \\
Item C & 22.1 & 12.7 & 34.8 \\
\bottomrule
\end{tabular}
\caption{Sample Data Table}
\label{tab:sample}
\end{table}

\section{Lists}
\subsection{Bullet Points}
\begin{itemize}
  \item First item
  \item Second item
  \item Third item with sub-items:
    \begin{itemize}
      \item Sub-item 1
      \item Sub-item 2
    \end{itemize}
\end{itemize}

\subsection{Numbered List}
\begin{enumerate}
  \item First step
  \item Second step
  \item Third step
\end{enumerate}

\section{Citations and References}
When you use the citation search feature, references will be automatically managed. For example, you might cite a paper like this: \cite{example2024}. All citations will be automatically added to your \texttt{cite.bib} file.

\section{Conclusion}
This template demonstrates basic \LaTeX{} features including mathematical formulas, tables, and lists. The bibliography system is now set up to work seamlessly with Bibby AI's citation management. Feel free to modify it for your needs!

\bibliographystyle{plain}
\bibliography{cite}

\end{document}\documentclass{article}
\usepackage[utf8]{inputenc}
\usepackage{booktabs}
\usepackage{amsmath}
\usepackage{natbib}

\title{My Document}
\author{Your Name}
\date{\today}

\begin{document}

\maketitle

\section{Introduction}
This is a sample document demonstrating various \LaTeX{} features including mathematical equations, tables, and formatting. You can cite papers using the citation search feature, and they will be automatically added to your \texttt{cite.bib} file.

\section{Mathematical Expressions}
Here are some example equations:

\subsection{Quadratic Formula}
The quadratic formula is given by:
\[x = \frac{-b \pm \sqrt{b^2 - 4ac}}{2a}\]

\subsection{Matrix Example}
A sample matrix:
\[
\begin{pmatrix}
a_{11} & a_{12} & a_{13} \\
a_{21} & a_{22} & a_{23} \\
a_{31} & a_{32} & a_{33}
\end{pmatrix}
\]

\section{Tables}
Below is a sample table showing data:

\begin{table}[h]
\centering
\begin{tabular}{@{}lrrr@{}}
\toprule
\textbf{Category} & \textbf{Value 1} & \textbf{Value 2} & \textbf{Total} \\
\midrule
Item A & 10.5 & 20.3 & 30.8 \\
Item B & 15.2 & 25.1 & 40.3 \\
Item C & 22.1 & 12.7 & 34.8 \\
\bottomrule
\end{tabular}
\caption{Sample Data Table}
\label{tab:sample}
\end{table}

\section{Lists}
\subsection{Bullet Points}
\begin{itemize}
  \item First item
  \item Second item
  \item Third item with sub-items:
    \begin{itemize}
      \item Sub-item 1
      \item Sub-item 2
    \end{itemize}
\end{itemize}

\subsection{Numbered List}
\begin{enumerate}
  \item First step
  \item Second step
  \item Third step
\end{enumerate}

\section{Citations and References}
When you use the citation search feature, references will be automatically managed. For example, you might cite a paper like this: \cite{example2024}. All citations will be automatically added to your \texttt{cite.bib} file.

\section{Conclusion}
This template demonstrates basic \LaTeX{} features including mathematical formulas, tables, and lists. The bibliography system is now set up to work seamlessly with Bibby AI's citation management. Feel free to modify it for your needs!

\bibliographystyle{plain}
\bibliography{cite}

\end{document}\documentclass{article}
\usepackage[utf8]{inputenc}
\usepackage{booktabs}
\usepackage{amsmath}
\usepackage{natbib}

\title{My Document}
\author{Your Name}
\date{\today}

\begin{document}

\maketitle

\section{Introduction}
This is a sample document demonstrating various \LaTeX{} features including mathematical equations, tables, and formatting. You can cite papers using the citation search feature, and they will be automatically added to your \texttt{cite.bib} file.

\section{Mathematical Expressions}
Here are some example equations:

\subsection{Quadratic Formula}
The quadratic formula is given by:
\[x = \frac{-b \pm \sqrt{b^2 - 4ac}}{2a}\]

\subsection{Matrix Example}
A sample matrix:
\[
\begin{pmatrix}
a_{11} & a_{12} & a_{13} \\
a_{21} & a_{22} & a_{23} \\
a_{31} & a_{32} & a_{33}
\end{pmatrix}
\]

\section{Tables}
Below is a sample table showing data:

\begin{table}[h]
\centering
\begin{tabular}{@{}lrrr@{}}
\toprule
\textbf{Category} & \textbf{Value 1} & \textbf{Value 2} & \textbf{Total} \\
\midrule
Item A & 10.5 & 20.3 & 30.8 \\
Item B & 15.2 & 25.1 & 40.3 \\
Item C & 22.1 & 12.7 & 34.8 \\
\bottomrule
\end{tabular}
\caption{Sample Data Table}
\label{tab:sample}
\end{table}

\section{Lists}
\subsection{Bullet Points}
\begin{itemize}
  \item First item
  \item Second item
  \item Third item with sub-items:
    \begin{itemize}
      \item Sub-item 1
      \item Sub-item 2
    \end{itemize}
\end{itemize}

\subsection{Numbered List}
\begin{enumerate}
  \item First step
  \item Second step
  \item Third step
\end{enumerate}

\section{Citations and References}
When you use the citation search feature, references will be automatically managed. For example, you might cite a paper like this: \cite{example2024}. All citations will be automatically added to your \texttt{cite.bib} file.

\section{Conclusion}
This template demonstrates basic \LaTeX{} features including mathematical formulas, tables, and lists. The bibliography system is now set up to work seamlessly with Bibby AI's citation management. Feel free to modify it for your needs!

\bibliographystyle{plain}
\bibliography{cite}

\end{document}\documentclass{article}
\usepackage[utf8]{inputenc}
\usepackage{booktabs}
\usepackage{amsmath}
\usepackage{natbib}

\title{My Document}
\author{Your Name}
\date{\today}

\begin{document}

\maketitle

\section{Introduction}
This is a sample document demonstrating various \LaTeX{} features including mathematical equations, tables, and formatting. You can cite papers using the citation search feature, and they will be automatically added to your \texttt{cite.bib} file.

\section{Mathematical Expressions}
Here are some example equations:

\subsection{Quadratic Formula}
The quadratic formula is given by:
\[x = \frac{-b \pm \sqrt{b^2 - 4ac}}{2a}\]

\subsection{Matrix Example}
A sample matrix:
\[
\begin{pmatrix}
a_{11} & a_{12} & a_{13} \\
a_{21} & a_{22} & a_{23} \\
a_{31} & a_{32} & a_{33}
\end{pmatrix}
\]

\section{Tables}
Below is a sample table showing data:

\begin{table}[h]
\centering
\begin{tabular}{@{}lrrr@{}}
\toprule
\textbf{Category} & \textbf{Value 1} & \textbf{Value 2} & \textbf{Total} \\
\midrule
Item A & 10.5 & 20.3 & 30.8 \\
Item B & 15.2 & 25.1 & 40.3 \\
Item C & 22.1 & 12.7 & 34.8 \\
\bottomrule
\end{tabular}
\caption{Sample Data Table}
\label{tab:sample}
\end{table}

\section{Lists}
\subsection{Bullet Points}
\begin{itemize}
  \item First item
  \item Second item
  \item Third item with sub-items:
    \begin{itemize}
      \item Sub-item 1
      \item Sub-item 2
    \end{itemize}
\end{itemize}

\subsection{Numbered List}
\begin{enumerate}
  \item First step
  \item Second step
  \item Third step
\end{enumerate}

\section{Citations and References}
When you use the citation search feature, references will be automatically managed. For example, you might cite a paper like this: \cite{example2024}. All citations will be automatically added to your \texttt{cite.bib} file.

\section{Conclusion}
This template demonstrates basic \LaTeX{} features including mathematical formulas, tables, and lists. The bibliography system is now set up to work seamlessly with Bibby AI's citation management. Feel free to modify it for your needs!

\bibliographystyle{plain}
\bibliography{cite}

\end{document}\documentclass{article}
\usepackage[utf8]{inputenc}
\usepackage{booktabs}
\usepackage{amsmath}
\usepackage{natbib}

\title{My Document}
\author{Your Name}
\date{\today}

\begin{document}

\maketitle

\section{Introduction}
This is a sample document demonstrating various \LaTeX{} features including mathematical equations, tables, and formatting. You can cite papers using the citation search feature, and they will be automatically added to your \texttt{cite.bib} file.

\section{Mathematical Expressions}
Here are some example equations:

\subsection{Quadratic Formula}
The quadratic formula is given by:
\[x = \frac{-b \pm \sqrt{b^2 - 4ac}}{2a}\]

\subsection{Matrix Example}
A sample matrix:
\[
\begin{pmatrix}
a_{11} & a_{12} & a_{13} \\
a_{21} & a_{22} & a_{23} \\
a_{31} & a_{32} & a_{33}
\end{pmatrix}
\]

\section{Tables}
Below is a sample table showing data:

\begin{table}[h]
\centering
\begin{tabular}{@{}lrrr@{}}
\toprule
\textbf{Category} & \textbf{Value 1} & \textbf{Value 2} & \textbf{Total} \\
\midrule
Item A & 10.5 & 20.3 & 30.8 \\
Item B & 15.2 & 25.1 & 40.3 \\
Item C & 22.1 & 12.7 & 34.8 \\
\bottomrule
\end{tabular}
\caption{Sample Data Table}
\label{tab:sample}
\end{table}

\section{Lists}
\subsection{Bullet Points}
\begin{itemize}
  \item First item
  \item Second item
  \item Third item with sub-items:
    \begin{itemize}
      \item Sub-item 1
      \item Sub-item 2
    \end{itemize}
\end{itemize}

\subsection{Numbered List}
\begin{enumerate}
  \item First step
  \item Second step
  \item Third step
\end{enumerate}

\section{Citations and References}
When you use the citation search feature, references will be automatically managed. For example, you might cite a paper like this: \cite{example2024}. All citations will be automatically added to your \texttt{cite.bib} file.

\section{Conclusion}
This template demonstrates basic \LaTeX{} features including mathematical formulas, tables, and lists. The bibliography system is now set up to work seamlessly with Bibby AI's citation management. Feel free to modify it for your needs!

\bibliographystyle{plain}
\bibliography{cite}

\end{document}\documentclass{article}
\usepackage[utf8]{inputenc}
\usepackage{booktabs}
\usepackage{amsmath}
\usepackage{natbib}

\title{My Document}
\author{Your Name}
\date{\today}

\begin{document}

\maketitle

\section{Introduction}
This is a sample document demonstrating various \LaTeX{} features including mathematical equations, tables, and formatting. You can cite papers using the citation search feature, and they will be automatically added to your \texttt{cite.bib} file.

\section{Mathematical Expressions}
Here are some example equations:

\subsection{Quadratic Formula}
The quadratic formula is given by:
\[x = \frac{-b \pm \sqrt{b^2 - 4ac}}{2a}\]

\subsection{Matrix Example}
A sample matrix:
\[
\begin{pmatrix}
a_{11} & a_{12} & a_{13} \\
a_{21} & a_{22} & a_{23} \\
a_{31} & a_{32} & a_{33}
\end{pmatrix}
\]

\section{Tables}
Below is a sample table showing data:

\begin{table}[h]
\centering
\begin{tabular}{@{}lrrr@{}}
\toprule
\textbf{Category} & \textbf{Value 1} & \textbf{Value 2} & \textbf{Total} \\
\midrule
Item A & 10.5 & 20.3 & 30.8 \\
Item B & 15.2 & 25.1 & 40.3 \\
Item C & 22.1 & 12.7 & 34.8 \\
\bottomrule
\end{tabular}
\caption{Sample Data Table}
\label{tab:sample}
\end{table}

\section{Lists}
\subsection{Bullet Points}
\begin{itemize}
  \item First item
  \item Second item
  \item Third item with sub-items:
    \begin{itemize}
      \item Sub-item 1
      \item Sub-item 2
    \end{itemize}
\end{itemize}

\subsection{Numbered List}
\begin{enumerate}
  \item First step
  \item Second step
  \item Third step
\end{enumerate}

\section{Citations and References}
When you use the citation search feature, references will be automatically managed. For example, you might cite a paper like this: \cite{example2024}. All citations will be automatically added to your \texttt{cite.bib} file.

\section{Conclusion}
This template demonstrates basic \LaTeX{} features including mathematical formulas, tables, and lists. The bibliography system is now set up to work seamlessly with Bibby AI's citation management. Feel free to modify it for your needs!

\bibliographystyle{plain}
\bibliography{cite}

\end{document}\documentclass{article}
\usepackage[utf8]{inputenc}
\usepackage{booktabs}
\usepackage{amsmath}
\usepackage{natbib}

\title{My Document}
\author{Your Name}
\date{\today}

\begin{document}

\maketitle

\section{Introduction}
This is a sample document demonstrating various \LaTeX{} features including mathematical equations, tables, and formatting. You can cite papers using the citation search feature, and they will be automatically added to your \texttt{cite.bib} file.

\section{Mathematical Expressions}
Here are some example equations:

\subsection{Quadratic Formula}
The quadratic formula is given by:
\[x = \frac{-b \pm \sqrt{b^2 - 4ac}}{2a}\]

\subsection{Matrix Example}
A sample matrix:
\[
\begin{pmatrix}
a_{11} & a_{12} & a_{13} \\
a_{21} & a_{22} & a_{23} \\
a_{31} & a_{32} & a_{33}
\end{pmatrix}
\]

\section{Tables}
Below is a sample table showing data:

\begin{table}[h]
\centering
\begin{tabular}{@{}lrrr@{}}
\toprule
\textbf{Category} & \textbf{Value 1} & \textbf{Value 2} & \textbf{Total} \\
\midrule
Item A & 10.5 & 20.3 & 30.8 \\
Item B & 15.2 & 25.1 & 40.3 \\
Item C & 22.1 & 12.7 & 34.8 \\
\bottomrule
\end{tabular}
\caption{Sample Data Table}
\label{tab:sample}
\end{table}

\section{Lists}
\subsection{Bullet Points}
\begin{itemize}
  \item First item
  \item Second item
  \item Third item with sub-items:
    \begin{itemize}
      \item Sub-item 1
      \item Sub-item 2
    \end{itemize}
\end{itemize}

\subsection{Numbered List}
\begin{enumerate}
  \item First step
  \item Second step
  \item Third step
\end{enumerate}

\section{Citations and References}
When you use the citation search feature, references will be automatically managed. For example, you might cite a paper like this: \cite{example2024}. All citations will be automatically added to your \texttt{cite.bib} file.

\section{Conclusion}
This template demonstrates basic \LaTeX{} features including mathematical formulas, tables, and lists. The bibliography system is now set up to work seamlessly with Bibby AI's citation management. Feel free to modify it for your needs!

\bibliographystyle{plain}
\bibliography{cite}

\end{document}\documentclass{article}
\usepackage[utf8]{inputenc}
\usepackage{booktabs}
\usepackage{amsmath}
\usepackage{natbib}

\title{My Document}
\author{Your Name}
\date{\today}

\begin{document}

\maketitle

\section{Introduction}
This is a sample document demonstrating various \LaTeX{} features including mathematical equations, tables, and formatting. You can cite papers using the citation search feature, and they will be automatically added to your \texttt{cite.bib} file.

\section{Mathematical Expressions}
Here are some example equations:

\subsection{Quadratic Formula}
The quadratic formula is given by:
\[x = \frac{-b \pm \sqrt{b^2 - 4ac}}{2a}\]

\subsection{Matrix Example}
A sample matrix:
\[
\begin{pmatrix}
a_{11} & a_{12} & a_{13} \\
a_{21} & a_{22} & a_{23} \\
a_{31} & a_{32} & a_{33}
\end{pmatrix}
\]

\section{Tables}
Below is a sample table showing data:

\begin{table}[h]
\centering
\begin{tabular}{@{}lrrr@{}}
\toprule
\textbf{Category} & \textbf{Value 1} & \textbf{Value 2} & \textbf{Total} \\
\midrule
Item A & 10.5 & 20.3 & 30.8 \\
Item B & 15.2 & 25.1 & 40.3 \\
Item C & 22.1 & 12.7 & 34.8 \\
\bottomrule
\end{tabular}
\caption{Sample Data Table}
\label{tab:sample}
\end{table}

\section{Lists}
\subsection{Bullet Points}
\begin{itemize}
  \item First item
  \item Second item
  \item Third item with sub-items:
    \begin{itemize}
      \item Sub-item 1
      \item Sub-item 2
    \end{itemize}
\end{itemize}

\subsection{Numbered List}
\begin{enumerate}
  \item First step
  \item Second step
  \item Third step
\end{enumerate}

\section{Citations and References}
When you use the citation search feature, references will be automatically managed. For example, you might cite a paper like this: \cite{example2024}. All citations will be automatically added to your \texttt{cite.bib} file.

\section{Conclusion}
This template demonstrates basic \LaTeX{} features including mathematical formulas, tables, and lists. The bibliography system is now set up to work seamlessly with Bibby AI's citation management. Feel free to modify it for your needs!

\bibliographystyle{plain}
\bibliography{cite}

\end{document}\documentclass{article}
\usepackage[utf8]{inputenc}
\usepackage{booktabs}
\usepackage{amsmath}
\usepackage{natbib}

\title{My Document}
\author{Your Name}
\date{\today}

\begin{document}

\maketitle

\section{Introduction}
This is a sample document demonstrating various \LaTeX{} features including mathematical equations, tables, and formatting. You can cite papers using the citation search feature, and they will be automatically added to your \texttt{cite.bib} file.

\section{Mathematical Expressions}
Here are some example equations:

\subsection{Quadratic Formula}
The quadratic formula is given by:
\[x = \frac{-b \pm \sqrt{b^2 - 4ac}}{2a}\]

\subsection{Matrix Example}
A sample matrix:
\[
\begin{pmatrix}
a_{11} & a_{12} & a_{13} \\
a_{21} & a_{22} & a_{23} \\
a_{31} & a_{32} & a_{33}
\end{pmatrix}
\]

\section{Tables}
Below is a sample table showing data:

\begin{table}[h]
\centering
\begin{tabular}{@{}lrrr@{}}
\toprule
\textbf{Category} & \textbf{Value 1} & \textbf{Value 2} & \textbf{Total} \\
\midrule
Item A & 10.5 & 20.3 & 30.8 \\
Item B & 15.2 & 25.1 & 40.3 \\
Item C & 22.1 & 12.7 & 34.8 \\
\bottomrule
\end{tabular}
\caption{Sample Data Table}
\label{tab:sample}
\end{table}

\section{Lists}
\subsection{Bullet Points}
\begin{itemize}
  \item First item
  \item Second item
  \item Third item with sub-items:
    \begin{itemize}
      \item Sub-item 1
      \item Sub-item 2
    \end{itemize}
\end{itemize}

\subsection{Numbered List}
\begin{enumerate}
  \item First step
  \item Second step
  \item Third step
\end{enumerate}

\section{Citations and References}
When you use the citation search feature, references will be automatically managed. For example, you might cite a paper like this: \cite{example2024}. All citations will be automatically added to your \texttt{cite.bib} file.

\section{Conclusion}
This template demonstrates basic \LaTeX{} features including mathematical formulas, tables, and lists. The bibliography system is now set up to work seamlessly with Bibby AI's citation management. Feel free to modify it for your needs!

\bibliographystyle{plain}
\bibliography{cite}

\end{document}\documentclass{article}
\usepackage[utf8]{inputenc}
\usepackage{booktabs}
\usepackage{amsmath}
\usepackage{natbib}

\title{My Document}
\author{Your Name}
\date{\today}

\begin{document}

\maketitle

\section{Introduction}
This is a sample document demonstrating various \LaTeX{} features including mathematical equations, tables, and formatting. You can cite papers using the citation search feature, and they will be automatically added to your \texttt{cite.bib} file.

\section{Mathematical Expressions}
Here are some example equations:

\subsection{Quadratic Formula}
The quadratic formula is given by:
\[x = \frac{-b \pm \sqrt{b^2 - 4ac}}{2a}\]

\subsection{Matrix Example}
A sample matrix:
\[
\begin{pmatrix}
a_{11} & a_{12} & a_{13} \\
a_{21} & a_{22} & a_{23} \\
a_{31} & a_{32} & a_{33}
\end{pmatrix}
\]

\section{Tables}
Below is a sample table showing data:

\begin{table}[h]
\centering
\begin{tabular}{@{}lrrr@{}}
\toprule
\textbf{Category} & \textbf{Value 1} & \textbf{Value 2} & \textbf{Total} \\
\midrule
Item A & 10.5 & 20.3 & 30.8 \\
Item B & 15.2 & 25.1 & 40.3 \\
Item C & 22.1 & 12.7 & 34.8 \\
\bottomrule
\end{tabular}
\caption{Sample Data Table}
\label{tab:sample}
\end{table}

\section{Lists}
\subsection{Bullet Points}
\begin{itemize}
  \item First item
  \item Second item
  \item Third item with sub-items:
    \begin{itemize}
      \item Sub-item 1
      \item Sub-item 2
    \end{itemize}
\end{itemize}

\subsection{Numbered List}
\begin{enumerate}
  \item First step
  \item Second step
  \item Third step
\end{enumerate}

\section{Citations and References}
When you use the citation search feature, references will be automatically managed. For example, you might cite a paper like this: \cite{example2024}. All citations will be automatically added to your \texttt{cite.bib} file.

\section{Conclusion}
This template demonstrates basic \LaTeX{} features including mathematical formulas, tables, and lists. The bibliography system is now set up to work seamlessly with Bibby AI's citation management. Feel free to modify it for your needs!

\bibliographystyle{plain}
\bibliography{cite}

\end{document}\documentclass{article}
\usepackage[utf8]{inputenc}
\usepackage{booktabs}
\usepackage{amsmath}
\usepackage{natbib}

\title{My Document}
\author{Your Name}
\date{\today}

\begin{document}

\maketitle

\section{Introduction}
This is a sample document demonstrating various \LaTeX{} features including mathematical equations, tables, and formatting. You can cite papers using the citation search feature, and they will be automatically added to your \texttt{cite.bib} file.

\section{Mathematical Expressions}
Here are some example equations:

\subsection{Quadratic Formula}
The quadratic formula is given by:
\[x = \frac{-b \pm \sqrt{b^2 - 4ac}}{2a}\]

\subsection{Matrix Example}
A sample matrix:
\[
\begin{pmatrix}
a_{11} & a_{12} & a_{13} \\
a_{21} & a_{22} & a_{23} \\
a_{31} & a_{32} & a_{33}
\end{pmatrix}
\]

\section{Tables}
Below is a sample table showing data:

\begin{table}[h]
\centering
\begin{tabular}{@{}lrrr@{}}
\toprule
\textbf{Category} & \textbf{Value 1} & \textbf{Value 2} & \textbf{Total} \\
\midrule
Item A & 10.5 & 20.3 & 30.8 \\
Item B & 15.2 & 25.1 & 40.3 \\
Item C & 22.1 & 12.7 & 34.8 \\
\bottomrule
\end{tabular}
\caption{Sample Data Table}
\label{tab:sample}
\end{table}

\section{Lists}
\subsection{Bullet Points}
\begin{itemize}
  \item First item
  \item Second item
  \item Third item with sub-items:
    \begin{itemize}
      \item Sub-item 1
      \item Sub-item 2
    \end{itemize}
\end{itemize}

\subsection{Numbered List}
\begin{enumerate}
  \item First step
  \item Second step
  \item Third step
\end{enumerate}

\section{Citations and References}
When you use the citation search feature, references will be automatically managed. For example, you might cite a paper like this: \cite{example2024}. All citations will be automatically added to your \texttt{cite.bib} file.

\section{Conclusion}
This template demonstrates basic \LaTeX{} features including mathematical formulas, tables, and lists. The bibliography system is now set up to work seamlessly with Bibby AI's citation management. Feel free to modify it for your needs!

\bibliographystyle{plain}
\bibliography{cite}

\end{document}\documentclass{article}
\usepackage[utf8]{inputenc}
\usepackage{booktabs}
\usepackage{amsmath}
\usepackage{natbib}

\title{My Document}
\author{Your Name}
\date{\today}

\begin{document}

\maketitle

\section{Introduction}
This is a sample document demonstrating various \LaTeX{} features including mathematical equations, tables, and formatting. You can cite papers using the citation search feature, and they will be automatically added to your \texttt{cite.bib} file.

\section{Mathematical Expressions}
Here are some example equations:

\subsection{Quadratic Formula}
The quadratic formula is given by:
\[x = \frac{-b \pm \sqrt{b^2 - 4ac}}{2a}\]

\subsection{Matrix Example}
A sample matrix:
\[
\begin{pmatrix}
a_{11} & a_{12} & a_{13} \\
a_{21} & a_{22} & a_{23} \\
a_{31} & a_{32} & a_{33}
\end{pmatrix}
\]

\section{Tables}
Below is a sample table showing data:

\begin{table}[h]
\centering
\begin{tabular}{@{}lrrr@{}}
\toprule
\textbf{Category} & \textbf{Value 1} & \textbf{Value 2} & \textbf{Total} \\
\midrule
Item A & 10.5 & 20.3 & 30.8 \\
Item B & 15.2 & 25.1 & 40.3 \\
Item C & 22.1 & 12.7 & 34.8 \\
\bottomrule
\end{tabular}
\caption{Sample Data Table}
\label{tab:sample}
\end{table}

\section{Lists}
\subsection{Bullet Points}
\begin{itemize}
  \item First item
  \item Second item
  \item Third item with sub-items:
    \begin{itemize}
      \item Sub-item 1
      \item Sub-item 2
    \end{itemize}
\end{itemize}

\subsection{Numbered List}
\begin{enumerate}
  \item First step
  \item Second step
  \item Third step
\end{enumerate}

\section{Citations and References}
When you use the citation search feature, references will be automatically managed. For example, you might cite a paper like this: \cite{example2024}. All citations will be automatically added to your \texttt{cite.bib} file.

\section{Conclusion}
This template demonstrates basic \LaTeX{} features including mathematical formulas, tables, and lists. The bibliography system is now set up to work seamlessly with Bibby AI's citation management. Feel free to modify it for your needs!

\bibliographystyle{plain}
\bibliography{cite}

\end{document}\documentclass{article}
\usepackage[utf8]{inputenc}
\usepackage{booktabs}
\usepackage{amsmath}
\usepackage{natbib}

\title{My Document}
\author{Your Name}
\date{\today}

\begin{document}

\maketitle

\section{Introduction}
This is a sample document demonstrating various \LaTeX{} features including mathematical equations, tables, and formatting. You can cite papers using the citation search feature, and they will be automatically added to your \texttt{cite.bib} file.

\section{Mathematical Expressions}
Here are some example equations:

\subsection{Quadratic Formula}
The quadratic formula is given by:
\[x = \frac{-b \pm \sqrt{b^2 - 4ac}}{2a}\]

\subsection{Matrix Example}
A sample matrix:
\[
\begin{pmatrix}
a_{11} & a_{12} & a_{13} \\
a_{21} & a_{22} & a_{23} \\
a_{31} & a_{32} & a_{33}
\end{pmatrix}
\]

\section{Tables}
Below is a sample table showing data:

\begin{table}[h]
\centering
\begin{tabular}{@{}lrrr@{}}
\toprule
\textbf{Category} & \textbf{Value 1} & \textbf{Value 2} & \textbf{Total} \\
\midrule
Item A & 10.5 & 20.3 & 30.8 \\
Item B & 15.2 & 25.1 & 40.3 \\
Item C & 22.1 & 12.7 & 34.8 \\
\bottomrule
\end{tabular}
\caption{Sample Data Table}
\label{tab:sample}
\end{table}

\section{Lists}
\subsection{Bullet Points}
\begin{itemize}
  \item First item
  \item Second item
  \item Third item with sub-items:
    \begin{itemize}
      \item Sub-item 1
      \item Sub-item 2
    \end{itemize}
\end{itemize}

\subsection{Numbered List}
\begin{enumerate}
  \item First step
  \item Second step
  \item Third step
\end{enumerate}

\section{Citations and References}
When you use the citation search feature, references will be automatically managed. For example, you might cite a paper like this: \cite{example2024}. All citations will be automatically added to your \texttt{cite.bib} file.

\section{Conclusion}
This template demonstrates basic \LaTeX{} features including mathematical formulas, tables, and lists. The bibliography system is now set up to work seamlessly with Bibby AI's citation management. Feel free to modify it for your needs!

\bibliographystyle{plain}
\bibliography{cite}

\end{document}\documentclass{article}
\usepackage[utf8]{inputenc}
\usepackage{booktabs}
\usepackage{amsmath}
\usepackage{natbib}

\title{My Document}
\author{Your Name}
\date{\today}

\begin{document}

\maketitle

\section{Introduction}
This is a sample document demonstrating various \LaTeX{} features including mathematical equations, tables, and formatting. You can cite papers using the citation search feature, and they will be automatically added to your \texttt{cite.bib} file.

\section{Mathematical Expressions}
Here are some example equations:

\subsection{Quadratic Formula}
The quadratic formula is given by:
\[x = \frac{-b \pm \sqrt{b^2 - 4ac}}{2a}\]

\subsection{Matrix Example}
A sample matrix:
\[
\begin{pmatrix}
a_{11} & a_{12} & a_{13} \\
a_{21} & a_{22} & a_{23} \\
a_{31} & a_{32} & a_{33}
\end{pmatrix}
\]

\section{Tables}
Below is a sample table showing data:

\begin{table}[h]
\centering
\begin{tabular}{@{}lrrr@{}}
\toprule
\textbf{Category} & \textbf{Value 1} & \textbf{Value 2} & \textbf{Total} \\
\midrule
Item A & 10.5 & 20.3 & 30.8 \\
Item B & 15.2 & 25.1 & 40.3 \\
Item C & 22.1 & 12.7 & 34.8 \\
\bottomrule
\end{tabular}
\caption{Sample Data Table}
\label{tab:sample}
\end{table}

\section{Lists}
\subsection{Bullet Points}
\begin{itemize}
  \item First item
  \item Second item
  \item Third item with sub-items:
    \begin{itemize}
      \item Sub-item 1
      \item Sub-item 2
    \end{itemize}
\end{itemize}

\subsection{Numbered List}
\begin{enumerate}
  \item First step
  \item Second step
  \item Third step
\end{enumerate}

\section{Citations and References}
When you use the citation search feature, references will be automatically managed. For example, you might cite a paper like this: \cite{example2024}. All citations will be automatically added to your \texttt{cite.bib} file.

\section{Conclusion}
This template demonstrates basic \LaTeX{} features including mathematical formulas, tables, and lists. The bibliography system is now set up to work seamlessly with Bibby AI's citation management. Feel free to modify it for your needs!

\bibliographystyle{plain}
\bibliography{cite}

\end{document}\documentclass{article}
\usepackage[utf8]{inputenc}
\usepackage{booktabs}
\usepackage{amsmath}
\usepackage{natbib}

\title{My Document}
\author{Your Name}
\date{\today}

\begin{document}

\maketitle

\section{Introduction}
This is a sample document demonstrating various \LaTeX{} features including mathematical equations, tables, and formatting. You can cite papers using the citation search feature, and they will be automatically added to your \texttt{cite.bib} file.

\section{Mathematical Expressions}
Here are some example equations:

\subsection{Quadratic Formula}
The quadratic formula is given by:
\[x = \frac{-b \pm \sqrt{b^2 - 4ac}}{2a}\]

\subsection{Matrix Example}
A sample matrix:
\[
\begin{pmatrix}
a_{11} & a_{12} & a_{13} \\
a_{21} & a_{22} & a_{23} \\
a_{31} & a_{32} & a_{33}
\end{pmatrix}
\]

\section{Tables}
Below is a sample table showing data:

\begin{table}[h]
\centering
\begin{tabular}{@{}lrrr@{}}
\toprule
\textbf{Category} & \textbf{Value 1} & \textbf{Value 2} & \textbf{Total} \\
\midrule
Item A & 10.5 & 20.3 & 30.8 \\
Item B & 15.2 & 25.1 & 40.3 \\
Item C & 22.1 & 12.7 & 34.8 \\
\bottomrule
\end{tabular}
\caption{Sample Data Table}
\label{tab:sample}
\end{table}

\section{Lists}
\subsection{Bullet Points}
\begin{itemize}
  \item First item
  \item Second item
  \item Third item with sub-items:
    \begin{itemize}
      \item Sub-item 1
      \item Sub-item 2
    \end{itemize}
\end{itemize}

\subsection{Numbered List}
\begin{enumerate}
  \item First step
  \item Second step
  \item Third step
\end{enumerate}

\section{Citations and References}
When you use the citation search feature, references will be automatically managed. For example, you might cite a paper like this: \cite{example2024}. All citations will be automatically added to your \texttt{cite.bib} file.

\section{Conclusion}
This template demonstrates basic \LaTeX{} features including mathematical formulas, tables, and lists. The bibliography system is now set up to work seamlessly with Bibby AI's citation management. Feel free to modify it for your needs!

\bibliographystyle{plain}
\bibliography{cite}

\end{document}\documentclass{article}
\usepackage[utf8]{inputenc}
\usepackage{booktabs}
\usepackage{amsmath}
\usepackage{natbib}

\title{My Document}
\author{Your Name}
\date{\today}

\begin{document}

\maketitle

\section{Introduction}
This is a sample document demonstrating various \LaTeX{} features including mathematical equations, tables, and formatting. You can cite papers using the citation search feature, and they will be automatically added to your \texttt{cite.bib} file.

\section{Mathematical Expressions}
Here are some example equations:

\subsection{Quadratic Formula}
The quadratic formula is given by:
\[x = \frac{-b \pm \sqrt{b^2 - 4ac}}{2a}\]

\subsection{Matrix Example}
A sample matrix:
\[
\begin{pmatrix}
a_{11} & a_{12} & a_{13} \\
a_{21} & a_{22} & a_{23} \\
a_{31} & a_{32} & a_{33}
\end{pmatrix}
\]

\section{Tables}
Below is a sample table showing data:

\begin{table}[h]
\centering
\begin{tabular}{@{}lrrr@{}}
\toprule
\textbf{Category} & \textbf{Value 1} & \textbf{Value 2} & \textbf{Total} \\
\midrule
Item A & 10.5 & 20.3 & 30.8 \\
Item B & 15.2 & 25.1 & 40.3 \\
Item C & 22.1 & 12.7 & 34.8 \\
\bottomrule
\end{tabular}
\caption{Sample Data Table}
\label{tab:sample}
\end{table}

\section{Lists}
\subsection{Bullet Points}
\begin{itemize}
  \item First item
  \item Second item
  \item Third item with sub-items:
    \begin{itemize}
      \item Sub-item 1
      \item Sub-item 2
    \end{itemize}
\end{itemize}

\subsection{Numbered List}
\begin{enumerate}
  \item First step
  \item Second step
  \item Third step
\end{enumerate}

\section{Citations and References}
When you use the citation search feature, references will be automatically managed. For example, you might cite a paper like this: \cite{example2024}. All citations will be automatically added to your \texttt{cite.bib} file.

\section{Conclusion}
This template demonstrates basic \LaTeX{} features including mathematical formulas, tables, and lists. The bibliography system is now set up to work seamlessly with Bibby AI's citation management. Feel free to modify it for your needs!

\bibliographystyle{plain}
\bibliography{cite}

\end{document}\documentclass{article}
\usepackage[utf8]{inputenc}
\usepackage{booktabs}
\usepackage{amsmath}
\usepackage{natbib}

\title{My Document}
\author{Your Name}
\date{\today}

\begin{document}

\maketitle

\section{Introduction}
This is a sample document demonstrating various \LaTeX{} features including mathematical equations, tables, and formatting. You can cite papers using the citation search feature, and they will be automatically added to your \texttt{cite.bib} file.

\section{Mathematical Expressions}
Here are some example equations:

\subsection{Quadratic Formula}
The quadratic formula is given by:
\[x = \frac{-b \pm \sqrt{b^2 - 4ac}}{2a}\]

\subsection{Matrix Example}
A sample matrix:
\[
\begin{pmatrix}
a_{11} & a_{12} & a_{13} \\
a_{21} & a_{22} & a_{23} \\
a_{31} & a_{32} & a_{33}
\end{pmatrix}
\]

\section{Tables}
Below is a sample table showing data:

\begin{table}[h]
\centering
\begin{tabular}{@{}lrrr@{}}
\toprule
\textbf{Category} & \textbf{Value 1} & \textbf{Value 2} & \textbf{Total} \\
\midrule
Item A & 10.5 & 20.3 & 30.8 \\
Item B & 15.2 & 25.1 & 40.3 \\
Item C & 22.1 & 12.7 & 34.8 \\
\bottomrule
\end{tabular}
\caption{Sample Data Table}
\label{tab:sample}
\end{table}

\section{Lists}
\subsection{Bullet Points}
\begin{itemize}
  \item First item
  \item Second item
  \item Third item with sub-items:
    \begin{itemize}
      \item Sub-item 1
      \item Sub-item 2
    \end{itemize}
\end{itemize}

\subsection{Numbered List}
\begin{enumerate}
  \item First step
  \item Second step
  \item Third step
\end{enumerate}

\section{Citations and References}
When you use the citation search feature, references will be automatically managed. For example, you might cite a paper like this: \cite{example2024}. All citations will be automatically added to your \texttt{cite.bib} file.

\section{Conclusion}
This template demonstrates basic \LaTeX{} features including mathematical formulas, tables, and lists. The bibliography system is now set up to work seamlessly with Bibby AI's citation management. Feel free to modify it for your needs!

\bibliographystyle{plain}
\bibliography{cite}

\end{document}\documentclass{article}
\usepackage[utf8]{inputenc}
\usepackage{booktabs}
\usepackage{amsmath}
\usepackage{natbib}

\title{My Document}
\author{Your Name}
\date{\today}

\begin{document}

\maketitle

\section{Introduction}
This is a sample document demonstrating various \LaTeX{} features including mathematical equations, tables, and formatting. You can cite papers using the citation search feature, and they will be automatically added to your \texttt{cite.bib} file.

\section{Mathematical Expressions}
Here are some example equations:

\subsection{Quadratic Formula}
The quadratic formula is given by:
\[x = \frac{-b \pm \sqrt{b^2 - 4ac}}{2a}\]

\subsection{Matrix Example}
A sample matrix:
\[
\begin{pmatrix}
a_{11} & a_{12} & a_{13} \\
a_{21} & a_{22} & a_{23} \\
a_{31} & a_{32} & a_{33}
\end{pmatrix}
\]

\section{Tables}
Below is a sample table showing data:

\begin{table}[h]
\centering
\begin{tabular}{@{}lrrr@{}}
\toprule
\textbf{Category} & \textbf{Value 1} & \textbf{Value 2} & \textbf{Total} \\
\midrule
Item A & 10.5 & 20.3 & 30.8 \\
Item B & 15.2 & 25.1 & 40.3 \\
Item C & 22.1 & 12.7 & 34.8 \\
\bottomrule
\end{tabular}
\caption{Sample Data Table}
\label{tab:sample}
\end{table}

\section{Lists}
\subsection{Bullet Points}
\begin{itemize}
  \item First item
  \item Second item
  \item Third item with sub-items:
    \begin{itemize}
      \item Sub-item 1
      \item Sub-item 2
    \end{itemize}
\end{itemize}

\subsection{Numbered List}
\begin{enumerate}
  \item First step
  \item Second step
  \item Third step
\end{enumerate}

\section{Citations and References}
When you use the citation search feature, references will be automatically managed. For example, you might cite a paper like this: \cite{example2024}. All citations will be automatically added to your \texttt{cite.bib} file.

\section{Conclusion}
This template demonstrates basic \LaTeX{} features including mathematical formulas, tables, and lists. The bibliography system is now set up to work seamlessly with Bibby AI's citation management. Feel free to modify it for your needs!

\bibliographystyle{plain}
\bibliography{cite}

\end{document}\documentclass{article}
\usepackage[utf8]{inputenc}
\usepackage{booktabs}
\usepackage{amsmath}
\usepackage{natbib}

\title{My Document}
\author{Your Name}
\date{\today}

\begin{document}

\maketitle

\section{Introduction}
This is a sample document demonstrating various \LaTeX{} features including mathematical equations, tables, and formatting. You can cite papers using the citation search feature, and they will be automatically added to your \texttt{cite.bib} file.

\section{Mathematical Expressions}
Here are some example equations:

\subsection{Quadratic Formula}
The quadratic formula is given by:
\[x = \frac{-b \pm \sqrt{b^2 - 4ac}}{2a}\]

\subsection{Matrix Example}
A sample matrix:
\[
\begin{pmatrix}
a_{11} & a_{12} & a_{13} \\
a_{21} & a_{22} & a_{23} \\
a_{31} & a_{32} & a_{33}
\end{pmatrix}
\]

\section{Tables}
Below is a sample table showing data:

\begin{table}[h]
\centering
\begin{tabular}{@{}lrrr@{}}
\toprule
\textbf{Category} & \textbf{Value 1} & \textbf{Value 2} & \textbf{Total} \\
\midrule
Item A & 10.5 & 20.3 & 30.8 \\
Item B & 15.2 & 25.1 & 40.3 \\
Item C & 22.1 & 12.7 & 34.8 \\
\bottomrule
\end{tabular}
\caption{Sample Data Table}
\label{tab:sample}
\end{table}

\section{Lists}
\subsection{Bullet Points}
\begin{itemize}
  \item First item
  \item Second item
  \item Third item with sub-items:
    \begin{itemize}
      \item Sub-item 1
      \item Sub-item 2
    \end{itemize}
\end{itemize}

\subsection{Numbered List}
\begin{enumerate}
  \item First step
  \item Second step
  \item Third step
\end{enumerate}

\section{Citations and References}
When you use the citation search feature, references will be automatically managed. For example, you might cite a paper like this: \cite{example2024}. All citations will be automatically added to your \texttt{cite.bib} file.

\section{Conclusion}
This template demonstrates basic \LaTeX{} features including mathematical formulas, tables, and lists. The bibliography system is now set up to work seamlessly with Bibby AI's citation management. Feel free to modify it for your needs!

\bibliographystyle{plain}
\bibliography{cite}

\end{document}\documentclass{article}
\usepackage[utf8]{inputenc}
\usepackage{booktabs}
\usepackage{amsmath}
\usepackage{natbib}

\title{My Document}
\author{Your Name}
\date{\today}

\begin{document}

\maketitle

\section{Introduction}
This is a sample document demonstrating various \LaTeX{} features including mathematical equations, tables, and formatting. You can cite papers using the citation search feature, and they will be automatically added to your \texttt{cite.bib} file.

\section{Mathematical Expressions}
Here are some example equations:

\subsection{Quadratic Formula}
The quadratic formula is given by:
\[x = \frac{-b \pm \sqrt{b^2 - 4ac}}{2a}\]

\subsection{Matrix Example}
A sample matrix:
\[
\begin{pmatrix}
a_{11} & a_{12} & a_{13} \\
a_{21} & a_{22} & a_{23} \\
a_{31} & a_{32} & a_{33}
\end{pmatrix}
\]

\section{Tables}
Below is a sample table showing data:

\begin{table}[h]
\centering
\begin{tabular}{@{}lrrr@{}}
\toprule
\textbf{Category} & \textbf{Value 1} & \textbf{Value 2} & \textbf{Total} \\
\midrule
Item A & 10.5 & 20.3 & 30.8 \\
Item B & 15.2 & 25.1 & 40.3 \\
Item C & 22.1 & 12.7 & 34.8 \\
\bottomrule
\end{tabular}
\caption{Sample Data Table}
\label{tab:sample}
\end{table}

\section{Lists}
\subsection{Bullet Points}
\begin{itemize}
  \item First item
  \item Second item
  \item Third item with sub-items:
    \begin{itemize}
      \item Sub-item 1
      \item Sub-item 2
    \end{itemize}
\end{itemize}

\subsection{Numbered List}
\begin{enumerate}
  \item First step
  \item Second step
  \item Third step
\end{enumerate}

\section{Citations and References}
When you use the citation search feature, references will be automatically managed. For example, you might cite a paper like this: \cite{example2024}. All citations will be automatically added to your \texttt{cite.bib} file.

\section{Conclusion}
This template demonstrates basic \LaTeX{} features including mathematical formulas, tables, and lists. The bibliography system is now set up to work seamlessly with Bibby AI's citation management. Feel free to modify it for your needs!

\bibliographystyle{plain}
\bibliography{cite}

\end{document}\documentclass{article}
\usepackage[utf8]{inputenc}
\usepackage{booktabs}
\usepackage{amsmath}
\usepackage{natbib}

\title{My Document}
\author{Your Name}
\date{\today}

\begin{document}

\maketitle

\section{Introduction}
This is a sample document demonstrating various \LaTeX{} features including mathematical equations, tables, and formatting. You can cite papers using the citation search feature, and they will be automatically added to your \texttt{cite.bib} file.

\section{Mathematical Expressions}
Here are some example equations:

\subsection{Quadratic Formula}
The quadratic formula is given by:
\[x = \frac{-b \pm \sqrt{b^2 - 4ac}}{2a}\]

\subsection{Matrix Example}
A sample matrix:
\[
\begin{pmatrix}
a_{11} & a_{12} & a_{13} \\
a_{21} & a_{22} & a_{23} \\
a_{31} & a_{32} & a_{33}
\end{pmatrix}
\]

\section{Tables}
Below is a sample table showing data:

\begin{table}[h]
\centering
\begin{tabular}{@{}lrrr@{}}
\toprule
\textbf{Category} & \textbf{Value 1} & \textbf{Value 2} & \textbf{Total} \\
\midrule
Item A & 10.5 & 20.3 & 30.8 \\
Item B & 15.2 & 25.1 & 40.3 \\
Item C & 22.1 & 12.7 & 34.8 \\
\bottomrule
\end{tabular}
\caption{Sample Data Table}
\label{tab:sample}
\end{table}

\section{Lists}
\subsection{Bullet Points}
\begin{itemize}
  \item First item
  \item Second item
  \item Third item with sub-items:
    \begin{itemize}
      \item Sub-item 1
      \item Sub-item 2
    \end{itemize}
\end{itemize}

\subsection{Numbered List}
\begin{enumerate}
  \item First step
  \item Second step
  \item Third step
\end{enumerate}

\section{Citations and References}
When you use the citation search feature, references will be automatically managed. For example, you might cite a paper like this: \cite{example2024}. All citations will be automatically added to your \texttt{cite.bib} file.

\section{Conclusion}
This template demonstrates basic \LaTeX{} features including mathematical formulas, tables, and lists. The bibliography system is now set up to work seamlessly with Bibby AI's citation management. Feel free to modify it for your needs!

\bibliographystyle{plain}
\bibliography{cite}

\end{document}\documentclass{article}
\usepackage[utf8]{inputenc}
\usepackage{booktabs}
\usepackage{amsmath}
\usepackage{natbib}

\title{My Document}
\author{Your Name}
\date{\today}

\begin{document}

\maketitle

\section{Introduction}
This is a sample document demonstrating various \LaTeX{} features including mathematical equations, tables, and formatting. You can cite papers using the citation search feature, and they will be automatically added to your \texttt{cite.bib} file.

\section{Mathematical Expressions}
Here are some example equations:

\subsection{Quadratic Formula}
The quadratic formula is given by:
\[x = \frac{-b \pm \sqrt{b^2 - 4ac}}{2a}\]

\subsection{Matrix Example}
A sample matrix:
\[
\begin{pmatrix}
a_{11} & a_{12} & a_{13} \\
a_{21} & a_{22} & a_{23} \\
a_{31} & a_{32} & a_{33}
\end{pmatrix}
\]

\section{Tables}
Below is a sample table showing data:

\begin{table}[h]
\centering
\begin{tabular}{@{}lrrr@{}}
\toprule
\textbf{Category} & \textbf{Value 1} & \textbf{Value 2} & \textbf{Total} \\
\midrule
Item A & 10.5 & 20.3 & 30.8 \\
Item B & 15.2 & 25.1 & 40.3 \\
Item C & 22.1 & 12.7 & 34.8 \\
\bottomrule
\end{tabular}
\caption{Sample Data Table}
\label{tab:sample}
\end{table}

\section{Lists}
\subsection{Bullet Points}
\begin{itemize}
  \item First item
  \item Second item
  \item Third item with sub-items:
    \begin{itemize}
      \item Sub-item 1
      \item Sub-item 2
    \end{itemize}
\end{itemize}

\subsection{Numbered List}
\begin{enumerate}
  \item First step
  \item Second step
  \item Third step
\end{enumerate}

\section{Citations and References}
When you use the citation search feature, references will be automatically managed. For example, you might cite a paper like this: \cite{example2024}. All citations will be automatically added to your \texttt{cite.bib} file.

\section{Conclusion}
This template demonstrates basic \LaTeX{} features including mathematical formulas, tables, and lists. The bibliography system is now set up to work seamlessly with Bibby AI's citation management. Feel free to modify it for your needs!

\bibliographystyle{plain}
\bibliography{cite}

\end{document}\documentclass{article}
\usepackage[utf8]{inputenc}
\usepackage{booktabs}
\usepackage{amsmath}
\usepackage{natbib}

\title{My Document}
\author{Your Name}
\date{\today}

\begin{document}

\maketitle

\section{Introduction}
This is a sample document demonstrating various \LaTeX{} features including mathematical equations, tables, and formatting. You can cite papers using the citation search feature, and they will be automatically added to your \texttt{cite.bib} file.

\section{Mathematical Expressions}
Here are some example equations:

\subsection{Quadratic Formula}
The quadratic formula is given by:
\[x = \frac{-b \pm \sqrt{b^2 - 4ac}}{2a}\]

\subsection{Matrix Example}
A sample matrix:
\[
\begin{pmatrix}
a_{11} & a_{12} & a_{13} \\
a_{21} & a_{22} & a_{23} \\
a_{31} & a_{32} & a_{33}
\end{pmatrix}
\]

\section{Tables}
Below is a sample table showing data:

\begin{table}[h]
\centering
\begin{tabular}{@{}lrrr@{}}
\toprule
\textbf{Category} & \textbf{Value 1} & \textbf{Value 2} & \textbf{Total} \\
\midrule
Item A & 10.5 & 20.3 & 30.8 \\
Item B & 15.2 & 25.1 & 40.3 \\
Item C & 22.1 & 12.7 & 34.8 \\
\bottomrule
\end{tabular}
\caption{Sample Data Table}
\label{tab:sample}
\end{table}

\section{Lists}
\subsection{Bullet Points}
\begin{itemize}
  \item First item
  \item Second item
  \item Third item with sub-items:
    \begin{itemize}
      \item Sub-item 1
      \item Sub-item 2
    \end{itemize}
\end{itemize}

\subsection{Numbered List}
\begin{enumerate}
  \item First step
  \item Second step
  \item Third step
\end{enumerate}

\section{Citations and References}
When you use the citation search feature, references will be automatically managed. For example, you might cite a paper like this: \cite{example2024}. All citations will be automatically added to your \texttt{cite.bib} file.

\section{Conclusion}
This template demonstrates basic \LaTeX{} features including mathematical formulas, tables, and lists. The bibliography system is now set up to work seamlessly with Bibby AI's citation management. Feel free to modify it for your needs!

\bibliographystyle{plain}
\bibliography{cite}

\end{document}\documentclass{article}
\usepackage[utf8]{inputenc}
\usepackage{booktabs}
\usepackage{amsmath}
\usepackage{natbib}

\title{My Document}
\author{Your Name}
\date{\today}

\begin{document}

\maketitle

\section{Introduction}
This is a sample document demonstrating various \LaTeX{} features including mathematical equations, tables, and formatting. You can cite papers using the citation search feature, and they will be automatically added to your \texttt{cite.bib} file.

\section{Mathematical Expressions}
Here are some example equations:

\subsection{Quadratic Formula}
The quadratic formula is given by:
\[x = \frac{-b \pm \sqrt{b^2 - 4ac}}{2a}\]

\subsection{Matrix Example}
A sample matrix:
\[
\begin{pmatrix}
a_{11} & a_{12} & a_{13} \\
a_{21} & a_{22} & a_{23} \\
a_{31} & a_{32} & a_{33}
\end{pmatrix}
\]

\section{Tables}
Below is a sample table showing data:

\begin{table}[h]
\centering
\begin{tabular}{@{}lrrr@{}}
\toprule
\textbf{Category} & \textbf{Value 1} & \textbf{Value 2} & \textbf{Total} \\
\midrule
Item A & 10.5 & 20.3 & 30.8 \\
Item B & 15.2 & 25.1 & 40.3 \\
Item C & 22.1 & 12.7 & 34.8 \\
\bottomrule
\end{tabular}
\caption{Sample Data Table}
\label{tab:sample}
\end{table}

\section{Lists}
\subsection{Bullet Points}
\begin{itemize}
  \item First item
  \item Second item
  \item Third item with sub-items:
    \begin{itemize}
      \item Sub-item 1
      \item Sub-item 2
    \end{itemize}
\end{itemize}

\subsection{Numbered List}
\begin{enumerate}
  \item First step
  \item Second step
  \item Third step
\end{enumerate}

\section{Citations and References}
When you use the citation search feature, references will be automatically managed. For example, you might cite a paper like this: \cite{example2024}. All citations will be automatically added to your \texttt{cite.bib} file.

\section{Conclusion}
This template demonstrates basic \LaTeX{} features including mathematical formulas, tables, and lists. The bibliography system is now set up to work seamlessly with Bibby AI's citation management. Feel free to modify it for your needs!

\bibliographystyle{plain}
\bibliography{cite}

\end{document}\documentclass{article}
\usepackage[utf8]{inputenc}
\usepackage{booktabs}
\usepackage{amsmath}
\usepackage{natbib}

\title{My Document}
\author{Your Name}
\date{\today}

\begin{document}

\maketitle

\section{Introduction}
This is a sample document demonstrating various \LaTeX{} features including mathematical equations, tables, and formatting. You can cite papers using the citation search feature, and they will be automatically added to your \texttt{cite.bib} file.

\section{Mathematical Expressions}
Here are some example equations:

\subsection{Quadratic Formula}
The quadratic formula is given by:
\[x = \frac{-b \pm \sqrt{b^2 - 4ac}}{2a}\]

\subsection{Matrix Example}
A sample matrix:
\[
\begin{pmatrix}
a_{11} & a_{12} & a_{13} \\
a_{21} & a_{22} & a_{23} \\
a_{31} & a_{32} & a_{33}
\end{pmatrix}
\]

\section{Tables}
Below is a sample table showing data:

\begin{table}[h]
\centering
\begin{tabular}{@{}lrrr@{}}
\toprule
\textbf{Category} & \textbf{Value 1} & \textbf{Value 2} & \textbf{Total} \\
\midrule
Item A & 10.5 & 20.3 & 30.8 \\
Item B & 15.2 & 25.1 & 40.3 \\
Item C & 22.1 & 12.7 & 34.8 \\
\bottomrule
\end{tabular}
\caption{Sample Data Table}
\label{tab:sample}
\end{table}

\section{Lists}
\subsection{Bullet Points}
\begin{itemize}
  \item First item
  \item Second item
  \item Third item with sub-items:
    \begin{itemize}
      \item Sub-item 1
      \item Sub-item 2
    \end{itemize}
\end{itemize}

\subsection{Numbered List}
\begin{enumerate}
  \item First step
  \item Second step
  \item Third step
\end{enumerate}

\section{Citations and References}
When you use the citation search feature, references will be automatically managed. For example, you might cite a paper like this: \cite{example2024}. All citations will be automatically added to your \texttt{cite.bib} file.

\section{Conclusion}
This template demonstrates basic \LaTeX{} features including mathematical formulas, tables, and lists. The bibliography system is now set up to work seamlessly with Bibby AI's citation management. Feel free to modify it for your needs!

\bibliographystyle{plain}
\bibliography{cite}

\end{document}\documentclass{article}
\usepackage[utf8]{inputenc}
\usepackage{booktabs}
\usepackage{amsmath}
\usepackage{natbib}

\title{My Document}
\author{Your Name}
\date{\today}

\begin{document}

\maketitle

\section{Introduction}
This is a sample document demonstrating various \LaTeX{} features including mathematical equations, tables, and formatting. You can cite papers using the citation search feature, and they will be automatically added to your \texttt{cite.bib} file.

\section{Mathematical Expressions}
Here are some example equations:

\subsection{Quadratic Formula}
The quadratic formula is given by:
\[x = \frac{-b \pm \sqrt{b^2 - 4ac}}{2a}\]

\subsection{Matrix Example}
A sample matrix:
\[
\begin{pmatrix}
a_{11} & a_{12} & a_{13} \\
a_{21} & a_{22} & a_{23} \\
a_{31} & a_{32} & a_{33}
\end{pmatrix}
\]

\section{Tables}
Below is a sample table showing data:

\begin{table}[h]
\centering
\begin{tabular}{@{}lrrr@{}}
\toprule
\textbf{Category} & \textbf{Value 1} & \textbf{Value 2} & \textbf{Total} \\
\midrule
Item A & 10.5 & 20.3 & 30.8 \\
Item B & 15.2 & 25.1 & 40.3 \\
Item C & 22.1 & 12.7 & 34.8 \\
\bottomrule
\end{tabular}
\caption{Sample Data Table}
\label{tab:sample}
\end{table}

\section{Lists}
\subsection{Bullet Points}
\begin{itemize}
  \item First item
  \item Second item
  \item Third item with sub-items:
    \begin{itemize}
      \item Sub-item 1
      \item Sub-item 2
    \end{itemize}
\end{itemize}

\subsection{Numbered List}
\begin{enumerate}
  \item First step
  \item Second step
  \item Third step
\end{enumerate}

\section{Citations and References}
When you use the citation search feature, references will be automatically managed. For example, you might cite a paper like this: \cite{example2024}. All citations will be automatically added to your \texttt{cite.bib} file.

\section{Conclusion}
This template demonstrates basic \LaTeX{} features including mathematical formulas, tables, and lists. The bibliography system is now set up to work seamlessly with Bibby AI's citation management. Feel free to modify it for your needs!

\bibliographystyle{plain}
\bibliography{cite}

\end{document}\documentclass{article}
\usepackage[utf8]{inputenc}
\usepackage{booktabs}
\usepackage{amsmath}
\usepackage{natbib}

\title{My Document}
\author{Your Name}
\date{\today}

\begin{document}

\maketitle

\section{Introduction}
This is a sample document demonstrating various \LaTeX{} features including mathematical equations, tables, and formatting. You can cite papers using the citation search feature, and they will be automatically added to your \texttt{cite.bib} file.

\section{Mathematical Expressions}
Here are some example equations:

\subsection{Quadratic Formula}
The quadratic formula is given by:
\[x = \frac{-b \pm \sqrt{b^2 - 4ac}}{2a}\]

\subsection{Matrix Example}
A sample matrix:
\[
\begin{pmatrix}
a_{11} & a_{12} & a_{13} \\
a_{21} & a_{22} & a_{23} \\
a_{31} & a_{32} & a_{33}
\end{pmatrix}
\]

\section{Tables}
Below is a sample table showing data:

\begin{table}[h]
\centering
\begin{tabular}{@{}lrrr@{}}
\toprule
\textbf{Category} & \textbf{Value 1} & \textbf{Value 2} & \textbf{Total} \\
\midrule
Item A & 10.5 & 20.3 & 30.8 \\
Item B & 15.2 & 25.1 & 40.3 \\
Item C & 22.1 & 12.7 & 34.8 \\
\bottomrule
\end{tabular}
\caption{Sample Data Table}
\label{tab:sample}
\end{table}

\section{Lists}
\subsection{Bullet Points}
\begin{itemize}
  \item First item
  \item Second item
  \item Third item with sub-items:
    \begin{itemize}
      \item Sub-item 1
      \item Sub-item 2
    \end{itemize}
\end{itemize}

\subsection{Numbered List}
\begin{enumerate}
  \item First step
  \item Second step
  \item Third step
\end{enumerate}

\section{Citations and References}
When you use the citation search feature, references will be automatically managed. For example, you might cite a paper like this: \cite{example2024}. All citations will be automatically added to your \texttt{cite.bib} file.

\section{Conclusion}
This template demonstrates basic \LaTeX{} features including mathematical formulas, tables, and lists. The bibliography system is now set up to work seamlessly with Bibby AI's citation management. Feel free to modify it for your needs!

\bibliographystyle{plain}
\bibliography{cite}

\end{document}\documentclass{article}
\usepackage[utf8]{inputenc}
\usepackage{booktabs}
\usepackage{amsmath}
\usepackage{natbib}

\title{My Document}
\author{Your Name}
\date{\today}

\begin{document}

\maketitle

\section{Introduction}
This is a sample document demonstrating various \LaTeX{} features including mathematical equations, tables, and formatting. You can cite papers using the citation search feature, and they will be automatically added to your \texttt{cite.bib} file.

\section{Mathematical Expressions}
Here are some example equations:

\subsection{Quadratic Formula}
The quadratic formula is given by:
\[x = \frac{-b \pm \sqrt{b^2 - 4ac}}{2a}\]

\subsection{Matrix Example}
A sample matrix:
\[
\begin{pmatrix}
a_{11} & a_{12} & a_{13} \\
a_{21} & a_{22} & a_{23} \\
a_{31} & a_{32} & a_{33}
\end{pmatrix}
\]

\section{Tables}
Below is a sample table showing data:

\begin{table}[h]
\centering
\begin{tabular}{@{}lrrr@{}}
\toprule
\textbf{Category} & \textbf{Value 1} & \textbf{Value 2} & \textbf{Total} \\
\midrule
Item A & 10.5 & 20.3 & 30.8 \\
Item B & 15.2 & 25.1 & 40.3 \\
Item C & 22.1 & 12.7 & 34.8 \\
\bottomrule
\end{tabular}
\caption{Sample Data Table}
\label{tab:sample}
\end{table}

\section{Lists}
\subsection{Bullet Points}
\begin{itemize}
  \item First item
  \item Second item
  \item Third item with sub-items:
    \begin{itemize}
      \item Sub-item 1
      \item Sub-item 2
    \end{itemize}
\end{itemize}

\subsection{Numbered List}
\begin{enumerate}
  \item First step
  \item Second step
  \item Third step
\end{enumerate}

\section{Citations and References}
When you use the citation search feature, references will be automatically managed. For example, you might cite a paper like this: \cite{example2024}. All citations will be automatically added to your \texttt{cite.bib} file.

\section{Conclusion}
This template demonstrates basic \LaTeX{} features including mathematical formulas, tables, and lists. The bibliography system is now set up to work seamlessly with Bibby AI's citation management. Feel free to modify it for your needs!

\bibliographystyle{plain}
\bibliography{cite}

\end{document}\documentclass{article}
\usepackage[utf8]{inputenc}
\usepackage{booktabs}
\usepackage{amsmath}
\usepackage{natbib}

\title{My Document}
\author{Your Name}
\date{\today}

\begin{document}

\maketitle

\section{Introduction}
This is a sample document demonstrating various \LaTeX{} features including mathematical equations, tables, and formatting. You can cite papers using the citation search feature, and they will be automatically added to your \texttt{cite.bib} file.

\section{Mathematical Expressions}
Here are some example equations:

\subsection{Quadratic Formula}
The quadratic formula is given by:
\[x = \frac{-b \pm \sqrt{b^2 - 4ac}}{2a}\]

\subsection{Matrix Example}
A sample matrix:
\[
\begin{pmatrix}
a_{11} & a_{12} & a_{13} \\
a_{21} & a_{22} & a_{23} \\
a_{31} & a_{32} & a_{33}
\end{pmatrix}
\]

\section{Tables}
Below is a sample table showing data:

\begin{table}[h]
\centering
\begin{tabular}{@{}lrrr@{}}
\toprule
\textbf{Category} & \textbf{Value 1} & \textbf{Value 2} & \textbf{Total} \\
\midrule
Item A & 10.5 & 20.3 & 30.8 \\
Item B & 15.2 & 25.1 & 40.3 \\
Item C & 22.1 & 12.7 & 34.8 \\
\bottomrule
\end{tabular}
\caption{Sample Data Table}
\label{tab:sample}
\end{table}

\section{Lists}
\subsection{Bullet Points}
\begin{itemize}
  \item First item
  \item Second item
  \item Third item with sub-items:
    \begin{itemize}
      \item Sub-item 1
      \item Sub-item 2
    \end{itemize}
\end{itemize}

\subsection{Numbered List}
\begin{enumerate}
  \item First step
  \item Second step
  \item Third step
\end{enumerate}

\section{Citations and References}
When you use the citation search feature, references will be automatically managed. For example, you might cite a paper like this: \cite{example2024}. All citations will be automatically added to your \texttt{cite.bib} file.

\section{Conclusion}
This template demonstrates basic \LaTeX{} features including mathematical formulas, tables, and lists. The bibliography system is now set up to work seamlessly with Bibby AI's citation management. Feel free to modify it for your needs!

\bibliographystyle{plain}
\bibliography{cite}

\end{document}\documentclass{article}
\usepackage[utf8]{inputenc}
\usepackage{booktabs}
\usepackage{amsmath}
\usepackage{natbib}

\title{My Document}
\author{Your Name}
\date{\today}

\begin{document}

\maketitle

\section{Introduction}
This is a sample document demonstrating various \LaTeX{} features including mathematical equations, tables, and formatting. You can cite papers using the citation search feature, and they will be automatically added to your \texttt{cite.bib} file.

\section{Mathematical Expressions}
Here are some example equations:

\subsection{Quadratic Formula}
The quadratic formula is given by:
\[x = \frac{-b \pm \sqrt{b^2 - 4ac}}{2a}\]

\subsection{Matrix Example}
A sample matrix:
\[
\begin{pmatrix}
a_{11} & a_{12} & a_{13} \\
a_{21} & a_{22} & a_{23} \\
a_{31} & a_{32} & a_{33}
\end{pmatrix}
\]

\section{Tables}
Below is a sample table showing data:

\begin{table}[h]
\centering
\begin{tabular}{@{}lrrr@{}}
\toprule
\textbf{Category} & \textbf{Value 1} & \textbf{Value 2} & \textbf{Total} \\
\midrule
Item A & 10.5 & 20.3 & 30.8 \\
Item B & 15.2 & 25.1 & 40.3 \\
Item C & 22.1 & 12.7 & 34.8 \\
\bottomrule
\end{tabular}
\caption{Sample Data Table}
\label{tab:sample}
\end{table}

\section{Lists}
\subsection{Bullet Points}
\begin{itemize}
  \item First item
  \item Second item
  \item Third item with sub-items:
    \begin{itemize}
      \item Sub-item 1
      \item Sub-item 2
    \end{itemize}
\end{itemize}

\subsection{Numbered List}
\begin{enumerate}
  \item First step
  \item Second step
  \item Third step
\end{enumerate}

\section{Citations and References}
When you use the citation search feature, references will be automatically managed. For example, you might cite a paper like this: \cite{example2024}. All citations will be automatically added to your \texttt{cite.bib} file.

\section{Conclusion}
This template demonstrates basic \LaTeX{} features including mathematical formulas, tables, and lists. The bibliography system is now set up to work seamlessly with Bibby AI's citation management. Feel free to modify it for your needs!

\bibliographystyle{plain}
\bibliography{cite}

\end{document}\documentclass{article}
\usepackage[utf8]{inputenc}
\usepackage{booktabs}
\usepackage{amsmath}
\usepackage{natbib}

\title{My Document}
\author{Your Name}
\date{\today}

\begin{document}

\maketitle

\section{Introduction}
This is a sample document demonstrating various \LaTeX{} features including mathematical equations, tables, and formatting. You can cite papers using the citation search feature, and they will be automatically added to your \texttt{cite.bib} file.

\section{Mathematical Expressions}
Here are some example equations:

\subsection{Quadratic Formula}
The quadratic formula is given by:
\[x = \frac{-b \pm \sqrt{b^2 - 4ac}}{2a}\]

\subsection{Matrix Example}
A sample matrix:
\[
\begin{pmatrix}
a_{11} & a_{12} & a_{13} \\
a_{21} & a_{22} & a_{23} \\
a_{31} & a_{32} & a_{33}
\end{pmatrix}
\]

\section{Tables}
Below is a sample table showing data:

\begin{table}[h]
\centering
\begin{tabular}{@{}lrrr@{}}
\toprule
\textbf{Category} & \textbf{Value 1} & \textbf{Value 2} & \textbf{Total} \\
\midrule
Item A & 10.5 & 20.3 & 30.8 \\
Item B & 15.2 & 25.1 & 40.3 \\
Item C & 22.1 & 12.7 & 34.8 \\
\bottomrule
\end{tabular}
\caption{Sample Data Table}
\label{tab:sample}
\end{table}

\section{Lists}
\subsection{Bullet Points}
\begin{itemize}
  \item First item
  \item Second item
  \item Third item with sub-items:
    \begin{itemize}
      \item Sub-item 1
      \item Sub-item 2
    \end{itemize}
\end{itemize}

\subsection{Numbered List}
\begin{enumerate}
  \item First step
  \item Second step
  \item Third step
\end{enumerate}

\section{Citations and References}
When you use the citation search feature, references will be automatically managed. For example, you might cite a paper like this: \cite{example2024}. All citations will be automatically added to your \texttt{cite.bib} file.

\section{Conclusion}
This template demonstrates basic \LaTeX{} features including mathematical formulas, tables, and lists. The bibliography system is now set up to work seamlessly with Bibby AI's citation management. Feel free to modify it for your needs!

\bibliographystyle{plain}
\bibliography{cite}

\end{document}\documentclass{article}
\usepackage[utf8]{inputenc}
\usepackage{booktabs}
\usepackage{amsmath}
\usepackage{natbib}

\title{My Document}
\author{Your Name}
\date{\today}

\begin{document}

\maketitle

\section{Introduction}
This is a sample document demonstrating various \LaTeX{} features including mathematical equations, tables, and formatting. You can cite papers using the citation search feature, and they will be automatically added to your \texttt{cite.bib} file.

\section{Mathematical Expressions}
Here are some example equations:

\subsection{Quadratic Formula}
The quadratic formula is given by:
\[x = \frac{-b \pm \sqrt{b^2 - 4ac}}{2a}\]

\subsection{Matrix Example}
A sample matrix:
\[
\begin{pmatrix}
a_{11} & a_{12} & a_{13} \\
a_{21} & a_{22} & a_{23} \\
a_{31} & a_{32} & a_{33}
\end{pmatrix}
\]

\section{Tables}
Below is a sample table showing data:

\begin{table}[h]
\centering
\begin{tabular}{@{}lrrr@{}}
\toprule
\textbf{Category} & \textbf{Value 1} & \textbf{Value 2} & \textbf{Total} \\
\midrule
Item A & 10.5 & 20.3 & 30.8 \\
Item B & 15.2 & 25.1 & 40.3 \\
Item C & 22.1 & 12.7 & 34.8 \\
\bottomrule
\end{tabular}
\caption{Sample Data Table}
\label{tab:sample}
\end{table}

\section{Lists}
\subsection{Bullet Points}
\begin{itemize}
  \item First item
  \item Second item
  \item Third item with sub-items:
    \begin{itemize}
      \item Sub-item 1
      \item Sub-item 2
    \end{itemize}
\end{itemize}

\subsection{Numbered List}
\begin{enumerate}
  \item First step
  \item Second step
  \item Third step
\end{enumerate}

\section{Citations and References}
When you use the citation search feature, references will be automatically managed. For example, you might cite a paper like this: \cite{example2024}. All citations will be automatically added to your \texttt{cite.bib} file.

\section{Conclusion}
This template demonstrates basic \LaTeX{} features including mathematical formulas, tables, and lists. The bibliography system is now set up to work seamlessly with Bibby AI's citation management. Feel free to modify it for your needs!

\bibliographystyle{plain}
\bibliography{cite}

\end{document}\documentclass{article}
\usepackage[utf8]{inputenc}
\usepackage{booktabs}
\usepackage{amsmath}
\usepackage{natbib}

\title{My Document}
\author{Your Name}
\date{\today}

\begin{document}

\maketitle

\section{Introduction}
This is a sample document demonstrating various \LaTeX{} features including mathematical equations, tables, and formatting. You can cite papers using the citation search feature, and they will be automatically added to your \texttt{cite.bib} file.

\section{Mathematical Expressions}
Here are some example equations:

\subsection{Quadratic Formula}
The quadratic formula is given by:
\[x = \frac{-b \pm \sqrt{b^2 - 4ac}}{2a}\]

\subsection{Matrix Example}
A sample matrix:
\[
\begin{pmatrix}
a_{11} & a_{12} & a_{13} \\
a_{21} & a_{22} & a_{23} \\
a_{31} & a_{32} & a_{33}
\end{pmatrix}
\]

\section{Tables}
Below is a sample table showing data:

\begin{table}[h]
\centering
\begin{tabular}{@{}lrrr@{}}
\toprule
\textbf{Category} & \textbf{Value 1} & \textbf{Value 2} & \textbf{Total} \\
\midrule
Item A & 10.5 & 20.3 & 30.8 \\
Item B & 15.2 & 25.1 & 40.3 \\
Item C & 22.1 & 12.7 & 34.8 \\
\bottomrule
\end{tabular}
\caption{Sample Data Table}
\label{tab:sample}
\end{table}

\section{Lists}
\subsection{Bullet Points}
\begin{itemize}
  \item First item
  \item Second item
  \item Third item with sub-items:
    \begin{itemize}
      \item Sub-item 1
      \item Sub-item 2
    \end{itemize}
\end{itemize}

\subsection{Numbered List}
\begin{enumerate}
  \item First step
  \item Second step
  \item Third step
\end{enumerate}

\section{Citations and References}
When you use the citation search feature, references will be automatically managed. For example, you might cite a paper like this: \cite{example2024}. All citations will be automatically added to your \texttt{cite.bib} file.

\section{Conclusion}
This template demonstrates basic \LaTeX{} features including mathematical formulas, tables, and lists. The bibliography system is now set up to work seamlessly with Bibby AI's citation management. Feel free to modify it for your needs!

\bibliographystyle{plain}
\bibliography{cite}

\end{document}\documentclass{article}
\usepackage[utf8]{inputenc}
\usepackage{booktabs}
\usepackage{amsmath}
\usepackage{natbib}

\title{My Document}
\author{Your Name}
\date{\today}

\begin{document}

\maketitle

\section{Introduction}
This is a sample document demonstrating various \LaTeX{} features including mathematical equations, tables, and formatting. You can cite papers using the citation search feature, and they will be automatically added to your \texttt{cite.bib} file.

\section{Mathematical Expressions}
Here are some example equations:

\subsection{Quadratic Formula}
The quadratic formula is given by:
\[x = \frac{-b \pm \sqrt{b^2 - 4ac}}{2a}\]

\subsection{Matrix Example}
A sample matrix:
\[
\begin{pmatrix}
a_{11} & a_{12} & a_{13} \\
a_{21} & a_{22} & a_{23} \\
a_{31} & a_{32} & a_{33}
\end{pmatrix}
\]

\section{Tables}
Below is a sample table showing data:

\begin{table}[h]
\centering
\begin{tabular}{@{}lrrr@{}}
\toprule
\textbf{Category} & \textbf{Value 1} & \textbf{Value 2} & \textbf{Total} \\
\midrule
Item A & 10.5 & 20.3 & 30.8 \\
Item B & 15.2 & 25.1 & 40.3 \\
Item C & 22.1 & 12.7 & 34.8 \\
\bottomrule
\end{tabular}
\caption{Sample Data Table}
\label{tab:sample}
\end{table}

\section{Lists}
\subsection{Bullet Points}
\begin{itemize}
  \item First item
  \item Second item
  \item Third item with sub-items:
    \begin{itemize}
      \item Sub-item 1
      \item Sub-item 2
    \end{itemize}
\end{itemize}

\subsection{Numbered List}
\begin{enumerate}
  \item First step
  \item Second step
  \item Third step
\end{enumerate}

\section{Citations and References}
When you use the citation search feature, references will be automatically managed. For example, you might cite a paper like this: \cite{example2024}. All citations will be automatically added to your \texttt{cite.bib} file.

\section{Conclusion}
This template demonstrates basic \LaTeX{} features including mathematical formulas, tables, and lists. The bibliography system is now set up to work seamlessly with Bibby AI's citation management. Feel free to modify it for your needs!

\bibliographystyle{plain}
\bibliography{cite}

\end{document}\documentclass{article}
\usepackage[utf8]{inputenc}
\usepackage{booktabs}
\usepackage{amsmath}
\usepackage{natbib}

\title{My Document}
\author{Your Name}
\date{\today}

\begin{document}

\maketitle

\section{Introduction}
This is a sample document demonstrating various \LaTeX{} features including mathematical equations, tables, and formatting. You can cite papers using the citation search feature, and they will be automatically added to your \texttt{cite.bib} file.

\section{Mathematical Expressions}
Here are some example equations:

\subsection{Quadratic Formula}
The quadratic formula is given by:
\[x = \frac{-b \pm \sqrt{b^2 - 4ac}}{2a}\]

\subsection{Matrix Example}
A sample matrix:
\[
\begin{pmatrix}
a_{11} & a_{12} & a_{13} \\
a_{21} & a_{22} & a_{23} \\
a_{31} & a_{32} & a_{33}
\end{pmatrix}
\]

\section{Tables}
Below is a sample table showing data:

\begin{table}[h]
\centering
\begin{tabular}{@{}lrrr@{}}
\toprule
\textbf{Category} & \textbf{Value 1} & \textbf{Value 2} & \textbf{Total} \\
\midrule
Item A & 10.5 & 20.3 & 30.8 \\
Item B & 15.2 & 25.1 & 40.3 \\
Item C & 22.1 & 12.7 & 34.8 \\
\bottomrule
\end{tabular}
\caption{Sample Data Table}
\label{tab:sample}
\end{table}

\section{Lists}
\subsection{Bullet Points}
\begin{itemize}
  \item First item
  \item Second item
  \item Third item with sub-items:
    \begin{itemize}
      \item Sub-item 1
      \item Sub-item 2
    \end{itemize}
\end{itemize}

\subsection{Numbered List}
\begin{enumerate}
  \item First step
  \item Second step
  \item Third step
\end{enumerate}

\section{Citations and References}
When you use the citation search feature, references will be automatically managed. For example, you might cite a paper like this: \cite{example2024}. All citations will be automatically added to your \texttt{cite.bib} file.

\section{Conclusion}
This template demonstrates basic \LaTeX{} features including mathematical formulas, tables, and lists. The bibliography system is now set up to work seamlessly with Bibby AI's citation management. Feel free to modify it for your needs!

\bibliographystyle{plain}
\bibliography{cite}

\end{document}\documentclass{article}
\usepackage[utf8]{inputenc}
\usepackage{booktabs}
\usepackage{amsmath}
\usepackage{natbib}

\title{My Document}
\author{Your Name}
\date{\today}

\begin{document}

\maketitle

\section{Introduction}
This is a sample document demonstrating various \LaTeX{} features including mathematical equations, tables, and formatting. You can cite papers using the citation search feature, and they will be automatically added to your \texttt{cite.bib} file.

\section{Mathematical Expressions}
Here are some example equations:

\subsection{Quadratic Formula}
The quadratic formula is given by:
\[x = \frac{-b \pm \sqrt{b^2 - 4ac}}{2a}\]

\subsection{Matrix Example}
A sample matrix:
\[
\begin{pmatrix}
a_{11} & a_{12} & a_{13} \\
a_{21} & a_{22} & a_{23} \\
a_{31} & a_{32} & a_{33}
\end{pmatrix}
\]

\section{Tables}
Below is a sample table showing data:

\begin{table}[h]
\centering
\begin{tabular}{@{}lrrr@{}}
\toprule
\textbf{Category} & \textbf{Value 1} & \textbf{Value 2} & \textbf{Total} \\
\midrule
Item A & 10.5 & 20.3 & 30.8 \\
Item B & 15.2 & 25.1 & 40.3 \\
Item C & 22.1 & 12.7 & 34.8 \\
\bottomrule
\end{tabular}
\caption{Sample Data Table}
\label{tab:sample}
\end{table}

\section{Lists}
\subsection{Bullet Points}
\begin{itemize}
  \item First item
  \item Second item
  \item Third item with sub-items:
    \begin{itemize}
      \item Sub-item 1
      \item Sub-item 2
    \end{itemize}
\end{itemize}

\subsection{Numbered List}
\begin{enumerate}
  \item First step
  \item Second step
  \item Third step
\end{enumerate}

\section{Citations and References}
When you use the citation search feature, references will be automatically managed. For example, you might cite a paper like this: \cite{example2024}. All citations will be automatically added to your \texttt{cite.bib} file.

\section{Conclusion}
This template demonstrates basic \LaTeX{} features including mathematical formulas, tables, and lists. The bibliography system is now set up to work seamlessly with Bibby AI's citation management. Feel free to modify it for your needs!

\bibliographystyle{plain}
\bibliography{cite}

\end{document}\documentclass{article}
\usepackage[utf8]{inputenc}
\usepackage{booktabs}
\usepackage{amsmath}
\usepackage{natbib}

\title{My Document}
\author{Your Name}
\date{\today}

\begin{document}

\maketitle

\section{Introduction}
This is a sample document demonstrating various \LaTeX{} features including mathematical equations, tables, and formatting. You can cite papers using the citation search feature, and they will be automatically added to your \texttt{cite.bib} file.

\section{Mathematical Expressions}
Here are some example equations:

\subsection{Quadratic Formula}
The quadratic formula is given by:
\[x = \frac{-b \pm \sqrt{b^2 - 4ac}}{2a}\]

\subsection{Matrix Example}
A sample matrix:
\[
\begin{pmatrix}
a_{11} & a_{12} & a_{13} \\
a_{21} & a_{22} & a_{23} \\
a_{31} & a_{32} & a_{33}
\end{pmatrix}
\]

\section{Tables}
Below is a sample table showing data:

\begin{table}[h]
\centering
\begin{tabular}{@{}lrrr@{}}
\toprule
\textbf{Category} & \textbf{Value 1} & \textbf{Value 2} & \textbf{Total} \\
\midrule
Item A & 10.5 & 20.3 & 30.8 \\
Item B & 15.2 & 25.1 & 40.3 \\
Item C & 22.1 & 12.7 & 34.8 \\
\bottomrule
\end{tabular}
\caption{Sample Data Table}
\label{tab:sample}
\end{table}

\section{Lists}
\subsection{Bullet Points}
\begin{itemize}
  \item First item
  \item Second item
  \item Third item with sub-items:
    \begin{itemize}
      \item Sub-item 1
      \item Sub-item 2
    \end{itemize}
\end{itemize}

\subsection{Numbered List}
\begin{enumerate}
  \item First step
  \item Second step
  \item Third step
\end{enumerate}

\section{Citations and References}
When you use the citation search feature, references will be automatically managed. For example, you might cite a paper like this: \cite{example2024}. All citations will be automatically added to your \texttt{cite.bib} file.

\section{Conclusion}
This template demonstrates basic \LaTeX{} features including mathematical formulas, tables, and lists. The bibliography system is now set up to work seamlessly with Bibby AI's citation management. Feel free to modify it for your needs!

\bibliographystyle{plain}
\bibliography{cite}

\end{document}\documentclass{article}
\usepackage[utf8]{inputenc}
\usepackage{booktabs}
\usepackage{amsmath}
\usepackage{natbib}

\title{My Document}
\author{Your Name}
\date{\today}

\begin{document}

\maketitle

\section{Introduction}
This is a sample document demonstrating various \LaTeX{} features including mathematical equations, tables, and formatting. You can cite papers using the citation search feature, and they will be automatically added to your \texttt{cite.bib} file.

\section{Mathematical Expressions}
Here are some example equations:

\subsection{Quadratic Formula}
The quadratic formula is given by:
\[x = \frac{-b \pm \sqrt{b^2 - 4ac}}{2a}\]

\subsection{Matrix Example}
A sample matrix:
\[
\begin{pmatrix}
a_{11} & a_{12} & a_{13} \\
a_{21} & a_{22} & a_{23} \\
a_{31} & a_{32} & a_{33}
\end{pmatrix}
\]

\section{Tables}
Below is a sample table showing data:

\begin{table}[h]
\centering
\begin{tabular}{@{}lrrr@{}}
\toprule
\textbf{Category} & \textbf{Value 1} & \textbf{Value 2} & \textbf{Total} \\
\midrule
Item A & 10.5 & 20.3 & 30.8 \\
Item B & 15.2 & 25.1 & 40.3 \\
Item C & 22.1 & 12.7 & 34.8 \\
\bottomrule
\end{tabular}
\caption{Sample Data Table}
\label{tab:sample}
\end{table}

\section{Lists}
\subsection{Bullet Points}
\begin{itemize}
  \item First item
  \item Second item
  \item Third item with sub-items:
    \begin{itemize}
      \item Sub-item 1
      \item Sub-item 2
    \end{itemize}
\end{itemize}

\subsection{Numbered List}
\begin{enumerate}
  \item First step
  \item Second step
  \item Third step
\end{enumerate}

\section{Citations and References}
When you use the citation search feature, references will be automatically managed. For example, you might cite a paper like this: \cite{example2024}. All citations will be automatically added to your \texttt{cite.bib} file.

\section{Conclusion}
This template demonstrates basic \LaTeX{} features including mathematical formulas, tables, and lists. The bibliography system is now set up to work seamlessly with Bibby AI's citation management. Feel free to modify it for your needs!

\bibliographystyle{plain}
\bibliography{cite}

\end{document}\documentclass{article}
\usepackage[utf8]{inputenc}
\usepackage{booktabs}
\usepackage{amsmath}
\usepackage{natbib}

\title{My Document}
\author{Your Name}
\date{\today}

\begin{document}

\maketitle

\section{Introduction}
This is a sample document demonstrating various \LaTeX{} features including mathematical equations, tables, and formatting. You can cite papers using the citation search feature, and they will be automatically added to your \texttt{cite.bib} file.

\section{Mathematical Expressions}
Here are some example equations:

\subsection{Quadratic Formula}
The quadratic formula is given by:
\[x = \frac{-b \pm \sqrt{b^2 - 4ac}}{2a}\]

\subsection{Matrix Example}
A sample matrix:
\[
\begin{pmatrix}
a_{11} & a_{12} & a_{13} \\
a_{21} & a_{22} & a_{23} \\
a_{31} & a_{32} & a_{33}
\end{pmatrix}
\]

\section{Tables}
Below is a sample table showing data:

\begin{table}[h]
\centering
\begin{tabular}{@{}lrrr@{}}
\toprule
\textbf{Category} & \textbf{Value 1} & \textbf{Value 2} & \textbf{Total} \\
\midrule
Item A & 10.5 & 20.3 & 30.8 \\
Item B & 15.2 & 25.1 & 40.3 \\
Item C & 22.1 & 12.7 & 34.8 \\
\bottomrule
\end{tabular}
\caption{Sample Data Table}
\label{tab:sample}
\end{table}

\section{Lists}
\subsection{Bullet Points}
\begin{itemize}
  \item First item
  \item Second item
  \item Third item with sub-items:
    \begin{itemize}
      \item Sub-item 1
      \item Sub-item 2
    \end{itemize}
\end{itemize}

\subsection{Numbered List}
\begin{enumerate}
  \item First step
  \item Second step
  \item Third step
\end{enumerate}

\section{Citations and References}
When you use the citation search feature, references will be automatically managed. For example, you might cite a paper like this: \cite{example2024}. All citations will be automatically added to your \texttt{cite.bib} file.

\section{Conclusion}
This template demonstrates basic \LaTeX{} features including mathematical formulas, tables, and lists. The bibliography system is now set up to work seamlessly with Bibby AI's citation management. Feel free to modify it for your needs!

\bibliographystyle{plain}
\bibliography{cite}

\end{document}\documentclass{article}
\usepackage[utf8]{inputenc}
\usepackage{booktabs}
\usepackage{amsmath}
\usepackage{natbib}

\title{My Document}
\author{Your Name}
\date{\today}

\begin{document}

\maketitle

\section{Introduction}
This is a sample document demonstrating various \LaTeX{} features including mathematical equations, tables, and formatting. You can cite papers using the citation search feature, and they will be automatically added to your \texttt{cite.bib} file.

\section{Mathematical Expressions}
Here are some example equations:

\subsection{Quadratic Formula}
The quadratic formula is given by:
\[x = \frac{-b \pm \sqrt{b^2 - 4ac}}{2a}\]

\subsection{Matrix Example}
A sample matrix:
\[
\begin{pmatrix}
a_{11} & a_{12} & a_{13} \\
a_{21} & a_{22} & a_{23} \\
a_{31} & a_{32} & a_{33}
\end{pmatrix}
\]

\section{Tables}
Below is a sample table showing data:

\begin{table}[h]
\centering
\begin{tabular}{@{}lrrr@{}}
\toprule
\textbf{Category} & \textbf{Value 1} & \textbf{Value 2} & \textbf{Total} \\
\midrule
Item A & 10.5 & 20.3 & 30.8 \\
Item B & 15.2 & 25.1 & 40.3 \\
Item C & 22.1 & 12.7 & 34.8 \\
\bottomrule
\end{tabular}
\caption{Sample Data Table}
\label{tab:sample}
\end{table}

\section{Lists}
\subsection{Bullet Points}
\begin{itemize}
  \item First item
  \item Second item
  \item Third item with sub-items:
    \begin{itemize}
      \item Sub-item 1
      \item Sub-item 2
    \end{itemize}
\end{itemize}

\subsection{Numbered List}
\begin{enumerate}
  \item First step
  \item Second step
  \item Third step
\end{enumerate}

\section{Citations and References}
When you use the citation search feature, references will be automatically managed. For example, you might cite a paper like this: \cite{example2024}. All citations will be automatically added to your \texttt{cite.bib} file.

\section{Conclusion}
This template demonstrates basic \LaTeX{} features including mathematical formulas, tables, and lists. The bibliography system is now set up to work seamlessly with Bibby AI's citation management. Feel free to modify it for your needs!

\bibliographystyle{plain}
\bibliography{cite}

\end{document}\documentclass{article}
\usepackage[utf8]{inputenc}
\usepackage{booktabs}
\usepackage{amsmath}
\usepackage{natbib}

\title{My Document}
\author{Your Name}
\date{\today}

\begin{document}

\maketitle

\section{Introduction}
This is a sample document demonstrating various \LaTeX{} features including mathematical equations, tables, and formatting. You can cite papers using the citation search feature, and they will be automatically added to your \texttt{cite.bib} file.

\section{Mathematical Expressions}
Here are some example equations:

\subsection{Quadratic Formula}
The quadratic formula is given by:
\[x = \frac{-b \pm \sqrt{b^2 - 4ac}}{2a}\]

\subsection{Matrix Example}
A sample matrix:
\[
\begin{pmatrix}
a_{11} & a_{12} & a_{13} \\
a_{21} & a_{22} & a_{23} \\
a_{31} & a_{32} & a_{33}
\end{pmatrix}
\]

\section{Tables}
Below is a sample table showing data:

\begin{table}[h]
\centering
\begin{tabular}{@{}lrrr@{}}
\toprule
\textbf{Category} & \textbf{Value 1} & \textbf{Value 2} & \textbf{Total} \\
\midrule
Item A & 10.5 & 20.3 & 30.8 \\
Item B & 15.2 & 25.1 & 40.3 \\
Item C & 22.1 & 12.7 & 34.8 \\
\bottomrule
\end{tabular}
\caption{Sample Data Table}
\label{tab:sample}
\end{table}

\section{Lists}
\subsection{Bullet Points}
\begin{itemize}
  \item First item
  \item Second item
  \item Third item with sub-items:
    \begin{itemize}
      \item Sub-item 1
      \item Sub-item 2
    \end{itemize}
\end{itemize}

\subsection{Numbered List}
\begin{enumerate}
  \item First step
  \item Second step
  \item Third step
\end{enumerate}

\section{Citations and References}
When you use the citation search feature, references will be automatically managed. For example, you might cite a paper like this: \cite{example2024}. All citations will be automatically added to your \texttt{cite.bib} file.

\section{Conclusion}
This template demonstrates basic \LaTeX{} features including mathematical formulas, tables, and lists. The bibliography system is now set up to work seamlessly with Bibby AI's citation management. Feel free to modify it for your needs!

\bibliographystyle{plain}
\bibliography{cite}

\end{document}\documentclass{article}
\usepackage[utf8]{inputenc}
\usepackage{booktabs}
\usepackage{amsmath}
\usepackage{natbib}

\title{My Document}
\author{Your Name}
\date{\today}

\begin{document}

\maketitle

\section{Introduction}
This is a sample document demonstrating various \LaTeX{} features including mathematical equations, tables, and formatting. You can cite papers using the citation search feature, and they will be automatically added to your \texttt{cite.bib} file.

\section{Mathematical Expressions}
Here are some example equations:

\subsection{Quadratic Formula}
The quadratic formula is given by:
\[x = \frac{-b \pm \sqrt{b^2 - 4ac}}{2a}\]

\subsection{Matrix Example}
A sample matrix:
\[
\begin{pmatrix}
a_{11} & a_{12} & a_{13} \\
a_{21} & a_{22} & a_{23} \\
a_{31} & a_{32} & a_{33}
\end{pmatrix}
\]

\section{Tables}
Below is a sample table showing data:

\begin{table}[h]
\centering
\begin{tabular}{@{}lrrr@{}}
\toprule
\textbf{Category} & \textbf{Value 1} & \textbf{Value 2} & \textbf{Total} \\
\midrule
Item A & 10.5 & 20.3 & 30.8 \\
Item B & 15.2 & 25.1 & 40.3 \\
Item C & 22.1 & 12.7 & 34.8 \\
\bottomrule
\end{tabular}
\caption{Sample Data Table}
\label{tab:sample}
\end{table}

\section{Lists}
\subsection{Bullet Points}
\begin{itemize}
  \item First item
  \item Second item
  \item Third item with sub-items:
    \begin{itemize}
      \item Sub-item 1
      \item Sub-item 2
    \end{itemize}
\end{itemize}

\subsection{Numbered List}
\begin{enumerate}
  \item First step
  \item Second step
  \item Third step
\end{enumerate}

\section{Citations and References}
When you use the citation search feature, references will be automatically managed. For example, you might cite a paper like this: \cite{example2024}. All citations will be automatically added to your \texttt{cite.bib} file.

\section{Conclusion}
This template demonstrates basic \LaTeX{} features including mathematical formulas, tables, and lists. The bibliography system is now set up to work seamlessly with Bibby AI's citation management. Feel free to modify it for your needs!

\bibliographystyle{plain}
\bibliography{cite}

\end{document}\documentclass{article}
\usepackage[utf8]{inputenc}
\usepackage{booktabs}
\usepackage{amsmath}
\usepackage{natbib}

\title{My Document}
\author{Your Name}
\date{\today}

\begin{document}

\maketitle

\section{Introduction}
This is a sample document demonstrating various \LaTeX{} features including mathematical equations, tables, and formatting. You can cite papers using the citation search feature, and they will be automatically added to your \texttt{cite.bib} file.

\section{Mathematical Expressions}
Here are some example equations:

\subsection{Quadratic Formula}
The quadratic formula is given by:
\[x = \frac{-b \pm \sqrt{b^2 - 4ac}}{2a}\]

\subsection{Matrix Example}
A sample matrix:
\[
\begin{pmatrix}
a_{11} & a_{12} & a_{13} \\
a_{21} & a_{22} & a_{23} \\
a_{31} & a_{32} & a_{33}
\end{pmatrix}
\]

\section{Tables}
Below is a sample table showing data:

\begin{table}[h]
\centering
\begin{tabular}{@{}lrrr@{}}
\toprule
\textbf{Category} & \textbf{Value 1} & \textbf{Value 2} & \textbf{Total} \\
\midrule
Item A & 10.5 & 20.3 & 30.8 \\
Item B & 15.2 & 25.1 & 40.3 \\
Item C & 22.1 & 12.7 & 34.8 \\
\bottomrule
\end{tabular}
\caption{Sample Data Table}
\label{tab:sample}
\end{table}

\section{Lists}
\subsection{Bullet Points}
\begin{itemize}
  \item First item
  \item Second item
  \item Third item with sub-items:
    \begin{itemize}
      \item Sub-item 1
      \item Sub-item 2
    \end{itemize}
\end{itemize}

\subsection{Numbered List}
\begin{enumerate}
  \item First step
  \item Second step
  \item Third step
\end{enumerate}

\section{Citations and References}
When you use the citation search feature, references will be automatically managed. For example, you might cite a paper like this: \cite{example2024}. All citations will be automatically added to your \texttt{cite.bib} file.

\section{Conclusion}
This template demonstrates basic \LaTeX{} features including mathematical formulas, tables, and lists. The bibliography system is now set up to work seamlessly with Bibby AI's citation management. Feel free to modify it for your needs!

\bibliographystyle{plain}
\bibliography{cite}

\end{document}\documentclass{article}
\usepackage[utf8]{inputenc}
\usepackage{booktabs}
\usepackage{amsmath}
\usepackage{natbib}

\title{My Document}
\author{Your Name}
\date{\today}

\begin{document}

\maketitle

\section{Introduction}
This is a sample document demonstrating various \LaTeX{} features including mathematical equations, tables, and formatting. You can cite papers using the citation search feature, and they will be automatically added to your \texttt{cite.bib} file.

\section{Mathematical Expressions}
Here are some example equations:

\subsection{Quadratic Formula}
The quadratic formula is given by:
\[x = \frac{-b \pm \sqrt{b^2 - 4ac}}{2a}\]

\subsection{Matrix Example}
A sample matrix:
\[
\begin{pmatrix}
a_{11} & a_{12} & a_{13} \\
a_{21} & a_{22} & a_{23} \\
a_{31} & a_{32} & a_{33}
\end{pmatrix}
\]

\section{Tables}
Below is a sample table showing data:

\begin{table}[h]
\centering
\begin{tabular}{@{}lrrr@{}}
\toprule
\textbf{Category} & \textbf{Value 1} & \textbf{Value 2} & \textbf{Total} \\
\midrule
Item A & 10.5 & 20.3 & 30.8 \\
Item B & 15.2 & 25.1 & 40.3 \\
Item C & 22.1 & 12.7 & 34.8 \\
\bottomrule
\end{tabular}
\caption{Sample Data Table}
\label{tab:sample}
\end{table}

\section{Lists}
\subsection{Bullet Points}
\begin{itemize}
  \item First item
  \item Second item
  \item Third item with sub-items:
    \begin{itemize}
      \item Sub-item 1
      \item Sub-item 2
    \end{itemize}
\end{itemize}

\subsection{Numbered List}
\begin{enumerate}
  \item First step
  \item Second step
  \item Third step
\end{enumerate}

\section{Citations and References}
When you use the citation search feature, references will be automatically managed. For example, you might cite a paper like this: \cite{example2024}. All citations will be automatically added to your \texttt{cite.bib} file.

\section{Conclusion}
This template demonstrates basic \LaTeX{} features including mathematical formulas, tables, and lists. The bibliography system is now set up to work seamlessly with Bibby AI's citation management. Feel free to modify it for your needs!

\bibliographystyle{plain}
\bibliography{cite}

\end{document}\documentclass{article}
\usepackage[utf8]{inputenc}
\usepackage{booktabs}
\usepackage{amsmath}
\usepackage{natbib}

\title{My Document}
\author{Your Name}
\date{\today}

\begin{document}

\maketitle

\section{Introduction}
This is a sample document demonstrating various \LaTeX{} features including mathematical equations, tables, and formatting. You can cite papers using the citation search feature, and they will be automatically added to your \texttt{cite.bib} file.

\section{Mathematical Expressions}
Here are some example equations:

\subsection{Quadratic Formula}
The quadratic formula is given by:
\[x = \frac{-b \pm \sqrt{b^2 - 4ac}}{2a}\]

\subsection{Matrix Example}
A sample matrix:
\[
\begin{pmatrix}
a_{11} & a_{12} & a_{13} \\
a_{21} & a_{22} & a_{23} \\
a_{31} & a_{32} & a_{33}
\end{pmatrix}
\]

\section{Tables}
Below is a sample table showing data:

\begin{table}[h]
\centering
\begin{tabular}{@{}lrrr@{}}
\toprule
\textbf{Category} & \textbf{Value 1} & \textbf{Value 2} & \textbf{Total} \\
\midrule
Item A & 10.5 & 20.3 & 30.8 \\
Item B & 15.2 & 25.1 & 40.3 \\
Item C & 22.1 & 12.7 & 34.8 \\
\bottomrule
\end{tabular}
\caption{Sample Data Table}
\label{tab:sample}
\end{table}

\section{Lists}
\subsection{Bullet Points}
\begin{itemize}
  \item First item
  \item Second item
  \item Third item with sub-items:
    \begin{itemize}
      \item Sub-item 1
      \item Sub-item 2
    \end{itemize}
\end{itemize}

\subsection{Numbered List}
\begin{enumerate}
  \item First step
  \item Second step
  \item Third step
\end{enumerate}

\section{Citations and References}
When you use the citation search feature, references will be automatically managed. For example, you might cite a paper like this: \cite{example2024}. All citations will be automatically added to your \texttt{cite.bib} file.

\section{Conclusion}
This template demonstrates basic \LaTeX{} features including mathematical formulas, tables, and lists. The bibliography system is now set up to work seamlessly with Bibby AI's citation management. Feel free to modify it for your needs!

\bibliographystyle{plain}
\bibliography{cite}

\end{document}\documentclass{article}
\usepackage[utf8]{inputenc}
\usepackage{booktabs}
\usepackage{amsmath}
\usepackage{natbib}

\title{My Document}
\author{Your Name}
\date{\today}

\begin{document}

\maketitle

\section{Introduction}
This is a sample document demonstrating various \LaTeX{} features including mathematical equations, tables, and formatting. You can cite papers using the citation search feature, and they will be automatically added to your \texttt{cite.bib} file.

\section{Mathematical Expressions}
Here are some example equations:

\subsection{Quadratic Formula}
The quadratic formula is given by:
\[x = \frac{-b \pm \sqrt{b^2 - 4ac}}{2a}\]

\subsection{Matrix Example}
A sample matrix:
\[
\begin{pmatrix}
a_{11} & a_{12} & a_{13} \\
a_{21} & a_{22} & a_{23} \\
a_{31} & a_{32} & a_{33}
\end{pmatrix}
\]

\section{Tables}
Below is a sample table showing data:

\begin{table}[h]
\centering
\begin{tabular}{@{}lrrr@{}}
\toprule
\textbf{Category} & \textbf{Value 1} & \textbf{Value 2} & \textbf{Total} \\
\midrule
Item A & 10.5 & 20.3 & 30.8 \\
Item B & 15.2 & 25.1 & 40.3 \\
Item C & 22.1 & 12.7 & 34.8 \\
\bottomrule
\end{tabular}
\caption{Sample Data Table}
\label{tab:sample}
\end{table}

\section{Lists}
\subsection{Bullet Points}
\begin{itemize}
  \item First item
  \item Second item
  \item Third item with sub-items:
    \begin{itemize}
      \item Sub-item 1
      \item Sub-item 2
    \end{itemize}
\end{itemize}

\subsection{Numbered List}
\begin{enumerate}
  \item First step
  \item Second step
  \item Third step
\end{enumerate}

\section{Citations and References}
When you use the citation search feature, references will be automatically managed. For example, you might cite a paper like this: \cite{example2024}. All citations will be automatically added to your \texttt{cite.bib} file.

\section{Conclusion}
This template demonstrates basic \LaTeX{} features including mathematical formulas, tables, and lists. The bibliography system is now set up to work seamlessly with Bibby AI's citation management. Feel free to modify it for your needs!

\bibliographystyle{plain}
\bibliography{cite}

\end{document}\documentclass{article}
\usepackage[utf8]{inputenc}
\usepackage{booktabs}
\usepackage{amsmath}
\usepackage{natbib}

\title{My Document}
\author{Your Name}
\date{\today}

\begin{document}

\maketitle

\section{Introduction}
This is a sample document demonstrating various \LaTeX{} features including mathematical equations, tables, and formatting. You can cite papers using the citation search feature, and they will be automatically added to your \texttt{cite.bib} file.

\section{Mathematical Expressions}
Here are some example equations:

\subsection{Quadratic Formula}
The quadratic formula is given by:
\[x = \frac{-b \pm \sqrt{b^2 - 4ac}}{2a}\]

\subsection{Matrix Example}
A sample matrix:
\[
\begin{pmatrix}
a_{11} & a_{12} & a_{13} \\
a_{21} & a_{22} & a_{23} \\
a_{31} & a_{32} & a_{33}
\end{pmatrix}
\]

\section{Tables}
Below is a sample table showing data:

\begin{table}[h]
\centering
\begin{tabular}{@{}lrrr@{}}
\toprule
\textbf{Category} & \textbf{Value 1} & \textbf{Value 2} & \textbf{Total} \\
\midrule
Item A & 10.5 & 20.3 & 30.8 \\
Item B & 15.2 & 25.1 & 40.3 \\
Item C & 22.1 & 12.7 & 34.8 \\
\bottomrule
\end{tabular}
\caption{Sample Data Table}
\label{tab:sample}
\end{table}

\section{Lists}
\subsection{Bullet Points}
\begin{itemize}
  \item First item
  \item Second item
  \item Third item with sub-items:
    \begin{itemize}
      \item Sub-item 1
      \item Sub-item 2
    \end{itemize}
\end{itemize}

\subsection{Numbered List}
\begin{enumerate}
  \item First step
  \item Second step
  \item Third step
\end{enumerate}

\section{Citations and References}
When you use the citation search feature, references will be automatically managed. For example, you might cite a paper like this: \cite{example2024}. All citations will be automatically added to your \texttt{cite.bib} file.

\section{Conclusion}
This template demonstrates basic \LaTeX{} features including mathematical formulas, tables, and lists. The bibliography system is now set up to work seamlessly with Bibby AI's citation management. Feel free to modify it for your needs!

\bibliographystyle{plain}
\bibliography{cite}

\end{document}\documentclass{article}
\usepackage[utf8]{inputenc}
\usepackage{booktabs}
\usepackage{amsmath}
\usepackage{natbib}

\title{My Document}
\author{Your Name}
\date{\today}

\begin{document}

\maketitle

\section{Introduction}
This is a sample document demonstrating various \LaTeX{} features including mathematical equations, tables, and formatting. You can cite papers using the citation search feature, and they will be automatically added to your \texttt{cite.bib} file.

\section{Mathematical Expressions}
Here are some example equations:

\subsection{Quadratic Formula}
The quadratic formula is given by:
\[x = \frac{-b \pm \sqrt{b^2 - 4ac}}{2a}\]

\subsection{Matrix Example}
A sample matrix:
\[
\begin{pmatrix}
a_{11} & a_{12} & a_{13} \\
a_{21} & a_{22} & a_{23} \\
a_{31} & a_{32} & a_{33}
\end{pmatrix}
\]

\section{Tables}
Below is a sample table showing data:

\begin{table}[h]
\centering
\begin{tabular}{@{}lrrr@{}}
\toprule
\textbf{Category} & \textbf{Value 1} & \textbf{Value 2} & \textbf{Total} \\
\midrule
Item A & 10.5 & 20.3 & 30.8 \\
Item B & 15.2 & 25.1 & 40.3 \\
Item C & 22.1 & 12.7 & 34.8 \\
\bottomrule
\end{tabular}
\caption{Sample Data Table}
\label{tab:sample}
\end{table}

\section{Lists}
\subsection{Bullet Points}
\begin{itemize}
  \item First item
  \item Second item
  \item Third item with sub-items:
    \begin{itemize}
      \item Sub-item 1
      \item Sub-item 2
    \end{itemize}
\end{itemize}

\subsection{Numbered List}
\begin{enumerate}
  \item First step
  \item Second step
  \item Third step
\end{enumerate}

\section{Citations and References}
When you use the citation search feature, references will be automatically managed. For example, you might cite a paper like this: \cite{example2024}. All citations will be automatically added to your \texttt{cite.bib} file.

\section{Conclusion}
This template demonstrates basic \LaTeX{} features including mathematical formulas, tables, and lists. The bibliography system is now set up to work seamlessly with Bibby AI's citation management. Feel free to modify it for your needs!

\bibliographystyle{plain}
\bibliography{cite}

\end{document}\documentclass{article}
\usepackage[utf8]{inputenc}
\usepackage{booktabs}
\usepackage{amsmath}
\usepackage{natbib}

\title{My Document}
\author{Your Name}
\date{\today}

\begin{document}

\maketitle

\section{Introduction}
This is a sample document demonstrating various \LaTeX{} features including mathematical equations, tables, and formatting. You can cite papers using the citation search feature, and they will be automatically added to your \texttt{cite.bib} file.

\section{Mathematical Expressions}
Here are some example equations:

\subsection{Quadratic Formula}
The quadratic formula is given by:
\[x = \frac{-b \pm \sqrt{b^2 - 4ac}}{2a}\]

\subsection{Matrix Example}
A sample matrix:
\[
\begin{pmatrix}
a_{11} & a_{12} & a_{13} \\
a_{21} & a_{22} & a_{23} \\
a_{31} & a_{32} & a_{33}
\end{pmatrix}
\]

\section{Tables}
Below is a sample table showing data:

\begin{table}[h]
\centering
\begin{tabular}{@{}lrrr@{}}
\toprule
\textbf{Category} & \textbf{Value 1} & \textbf{Value 2} & \textbf{Total} \\
\midrule
Item A & 10.5 & 20.3 & 30.8 \\
Item B & 15.2 & 25.1 & 40.3 \\
Item C & 22.1 & 12.7 & 34.8 \\
\bottomrule
\end{tabular}
\caption{Sample Data Table}
\label{tab:sample}
\end{table}

\section{Lists}
\subsection{Bullet Points}
\begin{itemize}
  \item First item
  \item Second item
  \item Third item with sub-items:
    \begin{itemize}
      \item Sub-item 1
      \item Sub-item 2
    \end{itemize}
\end{itemize}

\subsection{Numbered List}
\begin{enumerate}
  \item First step
  \item Second step
  \item Third step
\end{enumerate}

\section{Citations and References}
When you use the citation search feature, references will be automatically managed. For example, you might cite a paper like this: \cite{example2024}. All citations will be automatically added to your \texttt{cite.bib} file.

\section{Conclusion}
This template demonstrates basic \LaTeX{} features including mathematical formulas, tables, and lists. The bibliography system is now set up to work seamlessly with Bibby AI's citation management. Feel free to modify it for your needs!

\bibliographystyle{plain}
\bibliography{cite}

\end{document}\documentclass{article}
\usepackage[utf8]{inputenc}
\usepackage{booktabs}
\usepackage{amsmath}
\usepackage{natbib}

\title{My Document}
\author{Your Name}
\date{\today}

\begin{document}

\maketitle

\section{Introduction}
This is a sample document demonstrating various \LaTeX{} features including mathematical equations, tables, and formatting. You can cite papers using the citation search feature, and they will be automatically added to your \texttt{cite.bib} file.

\section{Mathematical Expressions}
Here are some example equations:

\subsection{Quadratic Formula}
The quadratic formula is given by:
\[x = \frac{-b \pm \sqrt{b^2 - 4ac}}{2a}\]

\subsection{Matrix Example}
A sample matrix:
\[
\begin{pmatrix}
a_{11} & a_{12} & a_{13} \\
a_{21} & a_{22} & a_{23} \\
a_{31} & a_{32} & a_{33}
\end{pmatrix}
\]

\section{Tables}
Below is a sample table showing data:

\begin{table}[h]
\centering
\begin{tabular}{@{}lrrr@{}}
\toprule
\textbf{Category} & \textbf{Value 1} & \textbf{Value 2} & \textbf{Total} \\
\midrule
Item A & 10.5 & 20.3 & 30.8 \\
Item B & 15.2 & 25.1 & 40.3 \\
Item C & 22.1 & 12.7 & 34.8 \\
\bottomrule
\end{tabular}
\caption{Sample Data Table}
\label{tab:sample}
\end{table}

\section{Lists}
\subsection{Bullet Points}
\begin{itemize}
  \item First item
  \item Second item
  \item Third item with sub-items:
    \begin{itemize}
      \item Sub-item 1
      \item Sub-item 2
    \end{itemize}
\end{itemize}

\subsection{Numbered List}
\begin{enumerate}
  \item First step
  \item Second step
  \item Third step
\end{enumerate}

\section{Citations and References}
When you use the citation search feature, references will be automatically managed. For example, you might cite a paper like this: \cite{example2024}. All citations will be automatically added to your \texttt{cite.bib} file.

\section{Conclusion}
This template demonstrates basic \LaTeX{} features including mathematical formulas, tables, and lists. The bibliography system is now set up to work seamlessly with Bibby AI's citation management. Feel free to modify it for your needs!

\bibliographystyle{plain}
\bibliography{cite}

\end{document}\documentclass{article}
\usepackage[utf8]{inputenc}
\usepackage{booktabs}
\usepackage{amsmath}
\usepackage{natbib}

\title{My Document}
\author{Your Name}
\date{\today}

\begin{document}

\maketitle

\section{Introduction}
This is a sample document demonstrating various \LaTeX{} features including mathematical equations, tables, and formatting. You can cite papers using the citation search feature, and they will be automatically added to your \texttt{cite.bib} file.

\section{Mathematical Expressions}
Here are some example equations:

\subsection{Quadratic Formula}
The quadratic formula is given by:
\[x = \frac{-b \pm \sqrt{b^2 - 4ac}}{2a}\]

\subsection{Matrix Example}
A sample matrix:
\[
\begin{pmatrix}
a_{11} & a_{12} & a_{13} \\
a_{21} & a_{22} & a_{23} \\
a_{31} & a_{32} & a_{33}
\end{pmatrix}
\]

\section{Tables}
Below is a sample table showing data:

\begin{table}[h]
\centering
\begin{tabular}{@{}lrrr@{}}
\toprule
\textbf{Category} & \textbf{Value 1} & \textbf{Value 2} & \textbf{Total} \\
\midrule
Item A & 10.5 & 20.3 & 30.8 \\
Item B & 15.2 & 25.1 & 40.3 \\
Item C & 22.1 & 12.7 & 34.8 \\
\bottomrule
\end{tabular}
\caption{Sample Data Table}
\label{tab:sample}
\end{table}

\section{Lists}
\subsection{Bullet Points}
\begin{itemize}
  \item First item
  \item Second item
  \item Third item with sub-items:
    \begin{itemize}
      \item Sub-item 1
      \item Sub-item 2
    \end{itemize}
\end{itemize}

\subsection{Numbered List}
\begin{enumerate}
  \item First step
  \item Second step
  \item Third step
\end{enumerate}

\section{Citations and References}
When you use the citation search feature, references will be automatically managed. For example, you might cite a paper like this: \cite{example2024}. All citations will be automatically added to your \texttt{cite.bib} file.

\section{Conclusion}
This template demonstrates basic \LaTeX{} features including mathematical formulas, tables, and lists. The bibliography system is now set up to work seamlessly with Bibby AI's citation management. Feel free to modify it for your needs!

\bibliographystyle{plain}
\bibliography{cite}

\end{document}\documentclass{article}
\usepackage[utf8]{inputenc}
\usepackage{booktabs}
\usepackage{amsmath}
\usepackage{natbib}

\title{My Document}
\author{Your Name}
\date{\today}

\begin{document}

\maketitle

\section{Introduction}
This is a sample document demonstrating various \LaTeX{} features including mathematical equations, tables, and formatting. You can cite papers using the citation search feature, and they will be automatically added to your \texttt{cite.bib} file.

\section{Mathematical Expressions}
Here are some example equations:

\subsection{Quadratic Formula}
The quadratic formula is given by:
\[x = \frac{-b \pm \sqrt{b^2 - 4ac}}{2a}\]

\subsection{Matrix Example}
A sample matrix:
\[
\begin{pmatrix}
a_{11} & a_{12} & a_{13} \\
a_{21} & a_{22} & a_{23} \\
a_{31} & a_{32} & a_{33}
\end{pmatrix}
\]

\section{Tables}
Below is a sample table showing data:

\begin{table}[h]
\centering
\begin{tabular}{@{}lrrr@{}}
\toprule
\textbf{Category} & \textbf{Value 1} & \textbf{Value 2} & \textbf{Total} \\
\midrule
Item A & 10.5 & 20.3 & 30.8 \\
Item B & 15.2 & 25.1 & 40.3 \\
Item C & 22.1 & 12.7 & 34.8 \\
\bottomrule
\end{tabular}
\caption{Sample Data Table}
\label{tab:sample}
\end{table}

\section{Lists}
\subsection{Bullet Points}
\begin{itemize}
  \item First item
  \item Second item
  \item Third item with sub-items:
    \begin{itemize}
      \item Sub-item 1
      \item Sub-item 2
    \end{itemize}
\end{itemize}

\subsection{Numbered List}
\begin{enumerate}
  \item First step
  \item Second step
  \item Third step
\end{enumerate}

\section{Citations and References}
When you use the citation search feature, references will be automatically managed. For example, you might cite a paper like this: \cite{example2024}. All citations will be automatically added to your \texttt{cite.bib} file.

\section{Conclusion}
This template demonstrates basic \LaTeX{} features including mathematical formulas, tables, and lists. The bibliography system is now set up to work seamlessly with Bibby AI's citation management. Feel free to modify it for your needs!

\bibliographystyle{plain}
\bibliography{cite}

\end{document}\documentclass{article}
\usepackage[utf8]{inputenc}
\usepackage{booktabs}
\usepackage{amsmath}
\usepackage{natbib}

\title{My Document}
\author{Your Name}
\date{\today}

\begin{document}

\maketitle

\section{Introduction}
This is a sample document demonstrating various \LaTeX{} features including mathematical equations, tables, and formatting. You can cite papers using the citation search feature, and they will be automatically added to your \texttt{cite.bib} file.

\section{Mathematical Expressions}
Here are some example equations:

\subsection{Quadratic Formula}
The quadratic formula is given by:
\[x = \frac{-b \pm \sqrt{b^2 - 4ac}}{2a}\]

\subsection{Matrix Example}
A sample matrix:
\[
\begin{pmatrix}
a_{11} & a_{12} & a_{13} \\
a_{21} & a_{22} & a_{23} \\
a_{31} & a_{32} & a_{33}
\end{pmatrix}
\]

\section{Tables}
Below is a sample table showing data:

\begin{table}[h]
\centering
\begin{tabular}{@{}lrrr@{}}
\toprule
\textbf{Category} & \textbf{Value 1} & \textbf{Value 2} & \textbf{Total} \\
\midrule
Item A & 10.5 & 20.3 & 30.8 \\
Item B & 15.2 & 25.1 & 40.3 \\
Item C & 22.1 & 12.7 & 34.8 \\
\bottomrule
\end{tabular}
\caption{Sample Data Table}
\label{tab:sample}
\end{table}

\section{Lists}
\subsection{Bullet Points}
\begin{itemize}
  \item First item
  \item Second item
  \item Third item with sub-items:
    \begin{itemize}
      \item Sub-item 1
      \item Sub-item 2
    \end{itemize}
\end{itemize}

\subsection{Numbered List}
\begin{enumerate}
  \item First step
  \item Second step
  \item Third step
\end{enumerate}

\section{Citations and References}
When you use the citation search feature, references will be automatically managed. For example, you might cite a paper like this: \cite{example2024}. All citations will be automatically added to your \texttt{cite.bib} file.

\section{Conclusion}
This template demonstrates basic \LaTeX{} features including mathematical formulas, tables, and lists. The bibliography system is now set up to work seamlessly with Bibby AI's citation management. Feel free to modify it for your needs!

\bibliographystyle{plain}
\bibliography{cite}

\end{document}\documentclass{article}
\usepackage[utf8]{inputenc}
\usepackage{booktabs}
\usepackage{amsmath}
\usepackage{natbib}

\title{My Document}
\author{Your Name}
\date{\today}

\begin{document}

\maketitle

\section{Introduction}
This is a sample document demonstrating various \LaTeX{} features including mathematical equations, tables, and formatting. You can cite papers using the citation search feature, and they will be automatically added to your \texttt{cite.bib} file.

\section{Mathematical Expressions}
Here are some example equations:

\subsection{Quadratic Formula}
The quadratic formula is given by:
\[x = \frac{-b \pm \sqrt{b^2 - 4ac}}{2a}\]

\subsection{Matrix Example}
A sample matrix:
\[
\begin{pmatrix}
a_{11} & a_{12} & a_{13} \\
a_{21} & a_{22} & a_{23} \\
a_{31} & a_{32} & a_{33}
\end{pmatrix}
\]

\section{Tables}
Below is a sample table showing data:

\begin{table}[h]
\centering
\begin{tabular}{@{}lrrr@{}}
\toprule
\textbf{Category} & \textbf{Value 1} & \textbf{Value 2} & \textbf{Total} \\
\midrule
Item A & 10.5 & 20.3 & 30.8 \\
Item B & 15.2 & 25.1 & 40.3 \\
Item C & 22.1 & 12.7 & 34.8 \\
\bottomrule
\end{tabular}
\caption{Sample Data Table}
\label{tab:sample}
\end{table}

\section{Lists}
\subsection{Bullet Points}
\begin{itemize}
  \item First item
  \item Second item
  \item Third item with sub-items:
    \begin{itemize}
      \item Sub-item 1
      \item Sub-item 2
    \end{itemize}
\end{itemize}

\subsection{Numbered List}
\begin{enumerate}
  \item First step
  \item Second step
  \item Third step
\end{enumerate}

\section{Citations and References}
When you use the citation search feature, references will be automatically managed. For example, you might cite a paper like this: \cite{example2024}. All citations will be automatically added to your \texttt{cite.bib} file.

\section{Conclusion}
This template demonstrates basic \LaTeX{} features including mathematical formulas, tables, and lists. The bibliography system is now set up to work seamlessly with Bibby AI's citation management. Feel free to modify it for your needs!

\bibliographystyle{plain}
\bibliography{cite}

\end{document}\documentclass{article}
\usepackage[utf8]{inputenc}
\usepackage{booktabs}
\usepackage{amsmath}
\usepackage{natbib}

\title{My Document}
\author{Your Name}
\date{\today}

\begin{document}

\maketitle

\section{Introduction}
This is a sample document demonstrating various \LaTeX{} features including mathematical equations, tables, and formatting. You can cite papers using the citation search feature, and they will be automatically added to your \texttt{cite.bib} file.

\section{Mathematical Expressions}
Here are some example equations:

\subsection{Quadratic Formula}
The quadratic formula is given by:
\[x = \frac{-b \pm \sqrt{b^2 - 4ac}}{2a}\]

\subsection{Matrix Example}
A sample matrix:
\[
\begin{pmatrix}
a_{11} & a_{12} & a_{13} \\
a_{21} & a_{22} & a_{23} \\
a_{31} & a_{32} & a_{33}
\end{pmatrix}
\]

\section{Tables}
Below is a sample table showing data:

\begin{table}[h]
\centering
\begin{tabular}{@{}lrrr@{}}
\toprule
\textbf{Category} & \textbf{Value 1} & \textbf{Value 2} & \textbf{Total} \\
\midrule
Item A & 10.5 & 20.3 & 30.8 \\
Item B & 15.2 & 25.1 & 40.3 \\
Item C & 22.1 & 12.7 & 34.8 \\
\bottomrule
\end{tabular}
\caption{Sample Data Table}
\label{tab:sample}
\end{table}

\section{Lists}
\subsection{Bullet Points}
\begin{itemize}
  \item First item
  \item Second item
  \item Third item with sub-items:
    \begin{itemize}
      \item Sub-item 1
      \item Sub-item 2
    \end{itemize}
\end{itemize}

\subsection{Numbered List}
\begin{enumerate}
  \item First step
  \item Second step
  \item Third step
\end{enumerate}

\section{Citations and References}
When you use the citation search feature, references will be automatically managed. For example, you might cite a paper like this: \cite{example2024}. All citations will be automatically added to your \texttt{cite.bib} file.

\section{Conclusion}
This template demonstrates basic \LaTeX{} features including mathematical formulas, tables, and lists. The bibliography system is now set up to work seamlessly with Bibby AI's citation management. Feel free to modify it for your needs!

\bibliographystyle{plain}
\bibliography{cite}

\end{document}\documentclass{article}
\usepackage[utf8]{inputenc}
\usepackage{booktabs}
\usepackage{amsmath}
\usepackage{natbib}

\title{My Document}
\author{Your Name}
\date{\today}

\begin{document}

\maketitle

\section{Introduction}
This is a sample document demonstrating various \LaTeX{} features including mathematical equations, tables, and formatting. You can cite papers using the citation search feature, and they will be automatically added to your \texttt{cite.bib} file.

\section{Mathematical Expressions}
Here are some example equations:

\subsection{Quadratic Formula}
The quadratic formula is given by:
\[x = \frac{-b \pm \sqrt{b^2 - 4ac}}{2a}\]

\subsection{Matrix Example}
A sample matrix:
\[
\begin{pmatrix}
a_{11} & a_{12} & a_{13} \\
a_{21} & a_{22} & a_{23} \\
a_{31} & a_{32} & a_{33}
\end{pmatrix}
\]

\section{Tables}
Below is a sample table showing data:

\begin{table}[h]
\centering
\begin{tabular}{@{}lrrr@{}}
\toprule
\textbf{Category} & \textbf{Value 1} & \textbf{Value 2} & \textbf{Total} \\
\midrule
Item A & 10.5 & 20.3 & 30.8 \\
Item B & 15.2 & 25.1 & 40.3 \\
Item C & 22.1 & 12.7 & 34.8 \\
\bottomrule
\end{tabular}
\caption{Sample Data Table}
\label{tab:sample}
\end{table}

\section{Lists}
\subsection{Bullet Points}
\begin{itemize}
  \item First item
  \item Second item
  \item Third item with sub-items:
    \begin{itemize}
      \item Sub-item 1
      \item Sub-item 2
    \end{itemize}
\end{itemize}

\subsection{Numbered List}
\begin{enumerate}
  \item First step
  \item Second step
  \item Third step
\end{enumerate}

\section{Citations and References}
When you use the citation search feature, references will be automatically managed. For example, you might cite a paper like this: \cite{example2024}. All citations will be automatically added to your \texttt{cite.bib} file.

\section{Conclusion}
This template demonstrates basic \LaTeX{} features including mathematical formulas, tables, and lists. The bibliography system is now set up to work seamlessly with Bibby AI's citation management. Feel free to modify it for your needs!

\bibliographystyle{plain}
\bibliography{cite}

\end{document}\documentclass{article}
\usepackage[utf8]{inputenc}
\usepackage{booktabs}
\usepackage{amsmath}
\usepackage{natbib}

\title{My Document}
\author{Your Name}
\date{\today}

\begin{document}

\maketitle

\section{Introduction}
This is a sample document demonstrating various \LaTeX{} features including mathematical equations, tables, and formatting. You can cite papers using the citation search feature, and they will be automatically added to your \texttt{cite.bib} file.

\section{Mathematical Expressions}
Here are some example equations:

\subsection{Quadratic Formula}
The quadratic formula is given by:
\[x = \frac{-b \pm \sqrt{b^2 - 4ac}}{2a}\]

\subsection{Matrix Example}
A sample matrix:
\[
\begin{pmatrix}
a_{11} & a_{12} & a_{13} \\
a_{21} & a_{22} & a_{23} \\
a_{31} & a_{32} & a_{33}
\end{pmatrix}
\]

\section{Tables}
Below is a sample table showing data:

\begin{table}[h]
\centering
\begin{tabular}{@{}lrrr@{}}
\toprule
\textbf{Category} & \textbf{Value 1} & \textbf{Value 2} & \textbf{Total} \\
\midrule
Item A & 10.5 & 20.3 & 30.8 \\
Item B & 15.2 & 25.1 & 40.3 \\
Item C & 22.1 & 12.7 & 34.8 \\
\bottomrule
\end{tabular}
\caption{Sample Data Table}
\label{tab:sample}
\end{table}

\section{Lists}
\subsection{Bullet Points}
\begin{itemize}
  \item First item
  \item Second item
  \item Third item with sub-items:
    \begin{itemize}
      \item Sub-item 1
      \item Sub-item 2
    \end{itemize}
\end{itemize}

\subsection{Numbered List}
\begin{enumerate}
  \item First step
  \item Second step
  \item Third step
\end{enumerate}

\section{Citations and References}
When you use the citation search feature, references will be automatically managed. For example, you might cite a paper like this: \cite{example2024}. All citations will be automatically added to your \texttt{cite.bib} file.

\section{Conclusion}
This template demonstrates basic \LaTeX{} features including mathematical formulas, tables, and lists. The bibliography system is now set up to work seamlessly with Bibby AI's citation management. Feel free to modify it for your needs!

\bibliographystyle{plain}
\bibliography{cite}

\end{document}\documentclass{article}
\usepackage[utf8]{inputenc}
\usepackage{booktabs}
\usepackage{amsmath}
\usepackage{natbib}

\title{My Document}
\author{Your Name}
\date{\today}

\begin{document}

\maketitle

\section{Introduction}
This is a sample document demonstrating various \LaTeX{} features including mathematical equations, tables, and formatting. You can cite papers using the citation search feature, and they will be automatically added to your \texttt{cite.bib} file.

\section{Mathematical Expressions}
Here are some example equations:

\subsection{Quadratic Formula}
The quadratic formula is given by:
\[x = \frac{-b \pm \sqrt{b^2 - 4ac}}{2a}\]

\subsection{Matrix Example}
A sample matrix:
\[
\begin{pmatrix}
a_{11} & a_{12} & a_{13} \\
a_{21} & a_{22} & a_{23} \\
a_{31} & a_{32} & a_{33}
\end{pmatrix}
\]

\section{Tables}
Below is a sample table showing data:

\begin{table}[h]
\centering
\begin{tabular}{@{}lrrr@{}}
\toprule
\textbf{Category} & \textbf{Value 1} & \textbf{Value 2} & \textbf{Total} \\
\midrule
Item A & 10.5 & 20.3 & 30.8 \\
Item B & 15.2 & 25.1 & 40.3 \\
Item C & 22.1 & 12.7 & 34.8 \\
\bottomrule
\end{tabular}
\caption{Sample Data Table}
\label{tab:sample}
\end{table}

\section{Lists}
\subsection{Bullet Points}
\begin{itemize}
  \item First item
  \item Second item
  \item Third item with sub-items:
    \begin{itemize}
      \item Sub-item 1
      \item Sub-item 2
    \end{itemize}
\end{itemize}

\subsection{Numbered List}
\begin{enumerate}
  \item First step
  \item Second step
  \item Third step
\end{enumerate}

\section{Citations and References}
When you use the citation search feature, references will be automatically managed. For example, you might cite a paper like this: \cite{example2024}. All citations will be automatically added to your \texttt{cite.bib} file.

\section{Conclusion}
This template demonstrates basic \LaTeX{} features including mathematical formulas, tables, and lists. The bibliography system is now set up to work seamlessly with Bibby AI's citation management. Feel free to modify it for your needs!

\bibliographystyle{plain}
\bibliography{cite}

\end{document}\documentclass{article}
\usepackage[utf8]{inputenc}
\usepackage{booktabs}
\usepackage{amsmath}
\usepackage{natbib}

\title{My Document}
\author{Your Name}
\date{\today}

\begin{document}

\maketitle

\section{Introduction}
This is a sample document demonstrating various \LaTeX{} features including mathematical equations, tables, and formatting. You can cite papers using the citation search feature, and they will be automatically added to your \texttt{cite.bib} file.

\section{Mathematical Expressions}
Here are some example equations:

\subsection{Quadratic Formula}
The quadratic formula is given by:
\[x = \frac{-b \pm \sqrt{b^2 - 4ac}}{2a}\]

\subsection{Matrix Example}
A sample matrix:
\[
\begin{pmatrix}
a_{11} & a_{12} & a_{13} \\
a_{21} & a_{22} & a_{23} \\
a_{31} & a_{32} & a_{33}
\end{pmatrix}
\]

\section{Tables}
Below is a sample table showing data:

\begin{table}[h]
\centering
\begin{tabular}{@{}lrrr@{}}
\toprule
\textbf{Category} & \textbf{Value 1} & \textbf{Value 2} & \textbf{Total} \\
\midrule
Item A & 10.5 & 20.3 & 30.8 \\
Item B & 15.2 & 25.1 & 40.3 \\
Item C & 22.1 & 12.7 & 34.8 \\
\bottomrule
\end{tabular}
\caption{Sample Data Table}
\label{tab:sample}
\end{table}

\section{Lists}
\subsection{Bullet Points}
\begin{itemize}
  \item First item
  \item Second item
  \item Third item with sub-items:
    \begin{itemize}
      \item Sub-item 1
      \item Sub-item 2
    \end{itemize}
\end{itemize}

\subsection{Numbered List}
\begin{enumerate}
  \item First step
  \item Second step
  \item Third step
\end{enumerate}

\section{Citations and References}
When you use the citation search feature, references will be automatically managed. For example, you might cite a paper like this: \cite{example2024}. All citations will be automatically added to your \texttt{cite.bib} file.

\section{Conclusion}
This template demonstrates basic \LaTeX{} features including mathematical formulas, tables, and lists. The bibliography system is now set up to work seamlessly with Bibby AI's citation management. Feel free to modify it for your needs!

\bibliographystyle{plain}
\bibliography{cite}

\end{document}\documentclass{article}
\usepackage[utf8]{inputenc}
\usepackage{booktabs}
\usepackage{amsmath}
\usepackage{natbib}

\title{My Document}
\author{Your Name}
\date{\today}

\begin{document}

\maketitle

\section{Introduction}
This is a sample document demonstrating various \LaTeX{} features including mathematical equations, tables, and formatting. You can cite papers using the citation search feature, and they will be automatically added to your \texttt{cite.bib} file.

\section{Mathematical Expressions}
Here are some example equations:

\subsection{Quadratic Formula}
The quadratic formula is given by:
\[x = \frac{-b \pm \sqrt{b^2 - 4ac}}{2a}\]

\subsection{Matrix Example}
A sample matrix:
\[
\begin{pmatrix}
a_{11} & a_{12} & a_{13} \\
a_{21} & a_{22} & a_{23} \\
a_{31} & a_{32} & a_{33}
\end{pmatrix}
\]

\section{Tables}
Below is a sample table showing data:

\begin{table}[h]
\centering
\begin{tabular}{@{}lrrr@{}}
\toprule
\textbf{Category} & \textbf{Value 1} & \textbf{Value 2} & \textbf{Total} \\
\midrule
Item A & 10.5 & 20.3 & 30.8 \\
Item B & 15.2 & 25.1 & 40.3 \\
Item C & 22.1 & 12.7 & 34.8 \\
\bottomrule
\end{tabular}
\caption{Sample Data Table}
\label{tab:sample}
\end{table}

\section{Lists}
\subsection{Bullet Points}
\begin{itemize}
  \item First item
  \item Second item
  \item Third item with sub-items:
    \begin{itemize}
      \item Sub-item 1
      \item Sub-item 2
    \end{itemize}
\end{itemize}

\subsection{Numbered List}
\begin{enumerate}
  \item First step
  \item Second step
  \item Third step
\end{enumerate}

\section{Citations and References}
When you use the citation search feature, references will be automatically managed. For example, you might cite a paper like this: \cite{example2024}. All citations will be automatically added to your \texttt{cite.bib} file.

\section{Conclusion}
This template demonstrates basic \LaTeX{} features including mathematical formulas, tables, and lists. The bibliography system is now set up to work seamlessly with Bibby AI's citation management. Feel free to modify it for your needs!

\bibliographystyle{plain}
\bibliography{cite}

\end{document}\documentclass{article}
\usepackage[utf8]{inputenc}
\usepackage{booktabs}
\usepackage{amsmath}
\usepackage{natbib}

\title{My Document}
\author{Your Name}
\date{\today}

\begin{document}

\maketitle

\section{Introduction}
This is a sample document demonstrating various \LaTeX{} features including mathematical equations, tables, and formatting. You can cite papers using the citation search feature, and they will be automatically added to your \texttt{cite.bib} file.

\section{Mathematical Expressions}
Here are some example equations:

\subsection{Quadratic Formula}
The quadratic formula is given by:
\[x = \frac{-b \pm \sqrt{b^2 - 4ac}}{2a}\]

\subsection{Matrix Example}
A sample matrix:
\[
\begin{pmatrix}
a_{11} & a_{12} & a_{13} \\
a_{21} & a_{22} & a_{23} \\
a_{31} & a_{32} & a_{33}
\end{pmatrix}
\]

\section{Tables}
Below is a sample table showing data:

\begin{table}[h]
\centering
\begin{tabular}{@{}lrrr@{}}
\toprule
\textbf{Category} & \textbf{Value 1} & \textbf{Value 2} & \textbf{Total} \\
\midrule
Item A & 10.5 & 20.3 & 30.8 \\
Item B & 15.2 & 25.1 & 40.3 \\
Item C & 22.1 & 12.7 & 34.8 \\
\bottomrule
\end{tabular}
\caption{Sample Data Table}
\label{tab:sample}
\end{table}

\section{Lists}
\subsection{Bullet Points}
\begin{itemize}
  \item First item
  \item Second item
  \item Third item with sub-items:
    \begin{itemize}
      \item Sub-item 1
      \item Sub-item 2
    \end{itemize}
\end{itemize}

\subsection{Numbered List}
\begin{enumerate}
  \item First step
  \item Second step
  \item Third step
\end{enumerate}

\section{Citations and References}
When you use the citation search feature, references will be automatically managed. For example, you might cite a paper like this: \cite{example2024}. All citations will be automatically added to your \texttt{cite.bib} file.

\section{Conclusion}
This template demonstrates basic \LaTeX{} features including mathematical formulas, tables, and lists. The bibliography system is now set up to work seamlessly with Bibby AI's citation management. Feel free to modify it for your needs!

\bibliographystyle{plain}
\bibliography{cite}

\end{document}\documentclass{article}
\usepackage[utf8]{inputenc}
\usepackage{booktabs}
\usepackage{amsmath}
\usepackage{natbib}

\title{My Document}
\author{Your Name}
\date{\today}

\begin{document}

\maketitle

\section{Introduction}
This is a sample document demonstrating various \LaTeX{} features including mathematical equations, tables, and formatting. You can cite papers using the citation search feature, and they will be automatically added to your \texttt{cite.bib} file.

\section{Mathematical Expressions}
Here are some example equations:

\subsection{Quadratic Formula}
The quadratic formula is given by:
\[x = \frac{-b \pm \sqrt{b^2 - 4ac}}{2a}\]

\subsection{Matrix Example}
A sample matrix:
\[
\begin{pmatrix}
a_{11} & a_{12} & a_{13} \\
a_{21} & a_{22} & a_{23} \\
a_{31} & a_{32} & a_{33}
\end{pmatrix}
\]

\section{Tables}
Below is a sample table showing data:

\begin{table}[h]
\centering
\begin{tabular}{@{}lrrr@{}}
\toprule
\textbf{Category} & \textbf{Value 1} & \textbf{Value 2} & \textbf{Total} \\
\midrule
Item A & 10.5 & 20.3 & 30.8 \\
Item B & 15.2 & 25.1 & 40.3 \\
Item C & 22.1 & 12.7 & 34.8 \\
\bottomrule
\end{tabular}
\caption{Sample Data Table}
\label{tab:sample}
\end{table}

\section{Lists}
\subsection{Bullet Points}
\begin{itemize}
  \item First item
  \item Second item
  \item Third item with sub-items:
    \begin{itemize}
      \item Sub-item 1
      \item Sub-item 2
    \end{itemize}
\end{itemize}

\subsection{Numbered List}
\begin{enumerate}
  \item First step
  \item Second step
  \item Third step
\end{enumerate}

\section{Citations and References}
When you use the citation search feature, references will be automatically managed. For example, you might cite a paper like this: \cite{example2024}. All citations will be automatically added to your \texttt{cite.bib} file.

\section{Conclusion}
This template demonstrates basic \LaTeX{} features including mathematical formulas, tables, and lists. The bibliography system is now set up to work seamlessly with Bibby AI's citation management. Feel free to modify it for your needs!

\bibliographystyle{plain}
\bibliography{cite}

\end{document}\documentclass{article}
\usepackage[utf8]{inputenc}
\usepackage{booktabs}
\usepackage{amsmath}
\usepackage{natbib}

\title{My Document}
\author{Your Name}
\date{\today}

\begin{document}

\maketitle

\section{Introduction}
This is a sample document demonstrating various \LaTeX{} features including mathematical equations, tables, and formatting. You can cite papers using the citation search feature, and they will be automatically added to your \texttt{cite.bib} file.

\section{Mathematical Expressions}
Here are some example equations:

\subsection{Quadratic Formula}
The quadratic formula is given by:
\[x = \frac{-b \pm \sqrt{b^2 - 4ac}}{2a}\]

\subsection{Matrix Example}
A sample matrix:
\[
\begin{pmatrix}
a_{11} & a_{12} & a_{13} \\
a_{21} & a_{22} & a_{23} \\
a_{31} & a_{32} & a_{33}
\end{pmatrix}
\]

\section{Tables}
Below is a sample table showing data:

\begin{table}[h]
\centering
\begin{tabular}{@{}lrrr@{}}
\toprule
\textbf{Category} & \textbf{Value 1} & \textbf{Value 2} & \textbf{Total} \\
\midrule
Item A & 10.5 & 20.3 & 30.8 \\
Item B & 15.2 & 25.1 & 40.3 \\
Item C & 22.1 & 12.7 & 34.8 \\
\bottomrule
\end{tabular}
\caption{Sample Data Table}
\label{tab:sample}
\end{table}

\section{Lists}
\subsection{Bullet Points}
\begin{itemize}
  \item First item
  \item Second item
  \item Third item with sub-items:
    \begin{itemize}
      \item Sub-item 1
      \item Sub-item 2
    \end{itemize}
\end{itemize}

\subsection{Numbered List}
\begin{enumerate}
  \item First step
  \item Second step
  \item Third step
\end{enumerate}

\section{Citations and References}
When you use the citation search feature, references will be automatically managed. For example, you might cite a paper like this: \cite{example2024}. All citations will be automatically added to your \texttt{cite.bib} file.

\section{Conclusion}
This template demonstrates basic \LaTeX{} features including mathematical formulas, tables, and lists. The bibliography system is now set up to work seamlessly with Bibby AI's citation management. Feel free to modify it for your needs!

\bibliographystyle{plain}
\bibliography{cite}

\end{document}\documentclass{article}
\usepackage[utf8]{inputenc}
\usepackage{booktabs}
\usepackage{amsmath}
\usepackage{natbib}

\title{My Document}
\author{Your Name}
\date{\today}

\begin{document}

\maketitle

\section{Introduction}
This is a sample document demonstrating various \LaTeX{} features including mathematical equations, tables, and formatting. You can cite papers using the citation search feature, and they will be automatically added to your \texttt{cite.bib} file.

\section{Mathematical Expressions}
Here are some example equations:

\subsection{Quadratic Formula}
The quadratic formula is given by:
\[x = \frac{-b \pm \sqrt{b^2 - 4ac}}{2a}\]

\subsection{Matrix Example}
A sample matrix:
\[
\begin{pmatrix}
a_{11} & a_{12} & a_{13} \\
a_{21} & a_{22} & a_{23} \\
a_{31} & a_{32} & a_{33}
\end{pmatrix}
\]

\section{Tables}
Below is a sample table showing data:

\begin{table}[h]
\centering
\begin{tabular}{@{}lrrr@{}}
\toprule
\textbf{Category} & \textbf{Value 1} & \textbf{Value 2} & \textbf{Total} \\
\midrule
Item A & 10.5 & 20.3 & 30.8 \\
Item B & 15.2 & 25.1 & 40.3 \\
Item C & 22.1 & 12.7 & 34.8 \\
\bottomrule
\end{tabular}
\caption{Sample Data Table}
\label{tab:sample}
\end{table}

\section{Lists}
\subsection{Bullet Points}
\begin{itemize}
  \item First item
  \item Second item
  \item Third item with sub-items:
    \begin{itemize}
      \item Sub-item 1
      \item Sub-item 2
    \end{itemize}
\end{itemize}

\subsection{Numbered List}
\begin{enumerate}
  \item First step
  \item Second step
  \item Third step
\end{enumerate}

\section{Citations and References}
When you use the citation search feature, references will be automatically managed. For example, you might cite a paper like this: \cite{example2024}. All citations will be automatically added to your \texttt{cite.bib} file.

\section{Conclusion}
This template demonstrates basic \LaTeX{} features including mathematical formulas, tables, and lists. The bibliography system is now set up to work seamlessly with Bibby AI's citation management. Feel free to modify it for your needs!

\bibliographystyle{plain}
\bibliography{cite}

\end{document}\documentclass{article}
\usepackage[utf8]{inputenc}
\usepackage{booktabs}
\usepackage{amsmath}
\usepackage{natbib}

\title{My Document}
\author{Your Name}
\date{\today}

\begin{document}

\maketitle

\section{Introduction}
This is a sample document demonstrating various \LaTeX{} features including mathematical equations, tables, and formatting. You can cite papers using the citation search feature, and they will be automatically added to your \texttt{cite.bib} file.

\section{Mathematical Expressions}
Here are some example equations:

\subsection{Quadratic Formula}
The quadratic formula is given by:
\[x = \frac{-b \pm \sqrt{b^2 - 4ac}}{2a}\]

\subsection{Matrix Example}
A sample matrix:
\[
\begin{pmatrix}
a_{11} & a_{12} & a_{13} \\
a_{21} & a_{22} & a_{23} \\
a_{31} & a_{32} & a_{33}
\end{pmatrix}
\]

\section{Tables}
Below is a sample table showing data:

\begin{table}[h]
\centering
\begin{tabular}{@{}lrrr@{}}
\toprule
\textbf{Category} & \textbf{Value 1} & \textbf{Value 2} & \textbf{Total} \\
\midrule
Item A & 10.5 & 20.3 & 30.8 \\
Item B & 15.2 & 25.1 & 40.3 \\
Item C & 22.1 & 12.7 & 34.8 \\
\bottomrule
\end{tabular}
\caption{Sample Data Table}
\label{tab:sample}
\end{table}

\section{Lists}
\subsection{Bullet Points}
\begin{itemize}
  \item First item
  \item Second item
  \item Third item with sub-items:
    \begin{itemize}
      \item Sub-item 1
      \item Sub-item 2
    \end{itemize}
\end{itemize}

\subsection{Numbered List}
\begin{enumerate}
  \item First step
  \item Second step
  \item Third step
\end{enumerate}

\section{Citations and References}
When you use the citation search feature, references will be automatically managed. For example, you might cite a paper like this: \cite{example2024}. All citations will be automatically added to your \texttt{cite.bib} file.

\section{Conclusion}
This template demonstrates basic \LaTeX{} features including mathematical formulas, tables, and lists. The bibliography system is now set up to work seamlessly with Bibby AI's citation management. Feel free to modify it for your needs!

\bibliographystyle{plain}
\bibliography{cite}

\end{document}\documentclass{article}
\usepackage[utf8]{inputenc}
\usepackage{booktabs}
\usepackage{amsmath}
\usepackage{natbib}

\title{My Document}
\author{Your Name}
\date{\today}

\begin{document}

\maketitle

\section{Introduction}
This is a sample document demonstrating various \LaTeX{} features including mathematical equations, tables, and formatting. You can cite papers using the citation search feature, and they will be automatically added to your \texttt{cite.bib} file.

\section{Mathematical Expressions}
Here are some example equations:

\subsection{Quadratic Formula}
The quadratic formula is given by:
\[x = \frac{-b \pm \sqrt{b^2 - 4ac}}{2a}\]

\subsection{Matrix Example}
A sample matrix:
\[
\begin{pmatrix}
a_{11} & a_{12} & a_{13} \\
a_{21} & a_{22} & a_{23} \\
a_{31} & a_{32} & a_{33}
\end{pmatrix}
\]

\section{Tables}
Below is a sample table showing data:

\begin{table}[h]
\centering
\begin{tabular}{@{}lrrr@{}}
\toprule
\textbf{Category} & \textbf{Value 1} & \textbf{Value 2} & \textbf{Total} \\
\midrule
Item A & 10.5 & 20.3 & 30.8 \\
Item B & 15.2 & 25.1 & 40.3 \\
Item C & 22.1 & 12.7 & 34.8 \\
\bottomrule
\end{tabular}
\caption{Sample Data Table}
\label{tab:sample}
\end{table}

\section{Lists}
\subsection{Bullet Points}
\begin{itemize}
  \item First item
  \item Second item
  \item Third item with sub-items:
    \begin{itemize}
      \item Sub-item 1
      \item Sub-item 2
    \end{itemize}
\end{itemize}

\subsection{Numbered List}
\begin{enumerate}
  \item First step
  \item Second step
  \item Third step
\end{enumerate}

\section{Citations and References}
When you use the citation search feature, references will be automatically managed. For example, you might cite a paper like this: \cite{example2024}. All citations will be automatically added to your \texttt{cite.bib} file.

\section{Conclusion}
This template demonstrates basic \LaTeX{} features including mathematical formulas, tables, and lists. The bibliography system is now set up to work seamlessly with Bibby AI's citation management. Feel free to modify it for your needs!

\bibliographystyle{plain}
\bibliography{cite}

\end{document}\documentclass{article}
\usepackage[utf8]{inputenc}
\usepackage{booktabs}
\usepackage{amsmath}
\usepackage{natbib}

\title{My Document}
\author{Your Name}
\date{\today}

\begin{document}

\maketitle

\section{Introduction}
This is a sample document demonstrating various \LaTeX{} features including mathematical equations, tables, and formatting. You can cite papers using the citation search feature, and they will be automatically added to your \texttt{cite.bib} file.

\section{Mathematical Expressions}
Here are some example equations:

\subsection{Quadratic Formula}
The quadratic formula is given by:
\[x = \frac{-b \pm \sqrt{b^2 - 4ac}}{2a}\]

\subsection{Matrix Example}
A sample matrix:
\[
\begin{pmatrix}
a_{11} & a_{12} & a_{13} \\
a_{21} & a_{22} & a_{23} \\
a_{31} & a_{32} & a_{33}
\end{pmatrix}
\]

\section{Tables}
Below is a sample table showing data:

\begin{table}[h]
\centering
\begin{tabular}{@{}lrrr@{}}
\toprule
\textbf{Category} & \textbf{Value 1} & \textbf{Value 2} & \textbf{Total} \\
\midrule
Item A & 10.5 & 20.3 & 30.8 \\
Item B & 15.2 & 25.1 & 40.3 \\
Item C & 22.1 & 12.7 & 34.8 \\
\bottomrule
\end{tabular}
\caption{Sample Data Table}
\label{tab:sample}
\end{table}

\section{Lists}
\subsection{Bullet Points}
\begin{itemize}
  \item First item
  \item Second item
  \item Third item with sub-items:
    \begin{itemize}
      \item Sub-item 1
      \item Sub-item 2
    \end{itemize}
\end{itemize}

\subsection{Numbered List}
\begin{enumerate}
  \item First step
  \item Second step
  \item Third step
\end{enumerate}

\section{Citations and References}
When you use the citation search feature, references will be automatically managed. For example, you might cite a paper like this: \cite{example2024}. All citations will be automatically added to your \texttt{cite.bib} file.

\section{Conclusion}
This template demonstrates basic \LaTeX{} features including mathematical formulas, tables, and lists. The bibliography system is now set up to work seamlessly with Bibby AI's citation management. Feel free to modify it for your needs!

\bibliographystyle{plain}
\bibliography{cite}

\end{document}\documentclass{article}
\usepackage[utf8]{inputenc}
\usepackage{booktabs}
\usepackage{amsmath}
\usepackage{natbib}

\title{My Document}
\author{Your Name}
\date{\today}

\begin{document}

\maketitle

\section{Introduction}
This is a sample document demonstrating various \LaTeX{} features including mathematical equations, tables, and formatting. You can cite papers using the citation search feature, and they will be automatically added to your \texttt{cite.bib} file.

\section{Mathematical Expressions}
Here are some example equations:

\subsection{Quadratic Formula}
The quadratic formula is given by:
\[x = \frac{-b \pm \sqrt{b^2 - 4ac}}{2a}\]

\subsection{Matrix Example}
A sample matrix:
\[
\begin{pmatrix}
a_{11} & a_{12} & a_{13} \\
a_{21} & a_{22} & a_{23} \\
a_{31} & a_{32} & a_{33}
\end{pmatrix}
\]

\section{Tables}
Below is a sample table showing data:

\begin{table}[h]
\centering
\begin{tabular}{@{}lrrr@{}}
\toprule
\textbf{Category} & \textbf{Value 1} & \textbf{Value 2} & \textbf{Total} \\
\midrule
Item A & 10.5 & 20.3 & 30.8 \\
Item B & 15.2 & 25.1 & 40.3 \\
Item C & 22.1 & 12.7 & 34.8 \\
\bottomrule
\end{tabular}
\caption{Sample Data Table}
\label{tab:sample}
\end{table}

\section{Lists}
\subsection{Bullet Points}
\begin{itemize}
  \item First item
  \item Second item
  \item Third item with sub-items:
    \begin{itemize}
      \item Sub-item 1
      \item Sub-item 2
    \end{itemize}
\end{itemize}

\subsection{Numbered List}
\begin{enumerate}
  \item First step
  \item Second step
  \item Third step
\end{enumerate}

\section{Citations and References}
When you use the citation search feature, references will be automatically managed. For example, you might cite a paper like this: \cite{example2024}. All citations will be automatically added to your \texttt{cite.bib} file.

\section{Conclusion}
This template demonstrates basic \LaTeX{} features including mathematical formulas, tables, and lists. The bibliography system is now set up to work seamlessly with Bibby AI's citation management. Feel free to modify it for your needs!

\bibliographystyle{plain}
\bibliography{cite}

\end{document}\documentclass{article}
\usepackage[utf8]{inputenc}
\usepackage{booktabs}
\usepackage{amsmath}
\usepackage{natbib}

\title{My Document}
\author{Your Name}
\date{\today}

\begin{document}

\maketitle

\section{Introduction}
This is a sample document demonstrating various \LaTeX{} features including mathematical equations, tables, and formatting. You can cite papers using the citation search feature, and they will be automatically added to your \texttt{cite.bib} file.

\section{Mathematical Expressions}
Here are some example equations:

\subsection{Quadratic Formula}
The quadratic formula is given by:
\[x = \frac{-b \pm \sqrt{b^2 - 4ac}}{2a}\]

\subsection{Matrix Example}
A sample matrix:
\[
\begin{pmatrix}
a_{11} & a_{12} & a_{13} \\
a_{21} & a_{22} & a_{23} \\
a_{31} & a_{32} & a_{33}
\end{pmatrix}
\]

\section{Tables}
Below is a sample table showing data:

\begin{table}[h]
\centering
\begin{tabular}{@{}lrrr@{}}
\toprule
\textbf{Category} & \textbf{Value 1} & \textbf{Value 2} & \textbf{Total} \\
\midrule
Item A & 10.5 & 20.3 & 30.8 \\
Item B & 15.2 & 25.1 & 40.3 \\
Item C & 22.1 & 12.7 & 34.8 \\
\bottomrule
\end{tabular}
\caption{Sample Data Table}
\label{tab:sample}
\end{table}

\section{Lists}
\subsection{Bullet Points}
\begin{itemize}
  \item First item
  \item Second item
  \item Third item with sub-items:
    \begin{itemize}
      \item Sub-item 1
      \item Sub-item 2
    \end{itemize}
\end{itemize}

\subsection{Numbered List}
\begin{enumerate}
  \item First step
  \item Second step
  \item Third step
\end{enumerate}

\section{Citations and References}
When you use the citation search feature, references will be automatically managed. For example, you might cite a paper like this: \cite{example2024}. All citations will be automatically added to your \texttt{cite.bib} file.

\section{Conclusion}
This template demonstrates basic \LaTeX{} features including mathematical formulas, tables, and lists. The bibliography system is now set up to work seamlessly with Bibby AI's citation management. Feel free to modify it for your needs!

\bibliographystyle{plain}
\bibliography{cite}

\end{document}\documentclass{article}
\usepackage[utf8]{inputenc}
\usepackage{booktabs}
\usepackage{amsmath}
\usepackage{natbib}

\title{My Document}
\author{Your Name}
\date{\today}

\begin{document}

\maketitle

\section{Introduction}
This is a sample document demonstrating various \LaTeX{} features including mathematical equations, tables, and formatting. You can cite papers using the citation search feature, and they will be automatically added to your \texttt{cite.bib} file.

\section{Mathematical Expressions}
Here are some example equations:

\subsection{Quadratic Formula}
The quadratic formula is given by:
\[x = \frac{-b \pm \sqrt{b^2 - 4ac}}{2a}\]

\subsection{Matrix Example}
A sample matrix:
\[
\begin{pmatrix}
a_{11} & a_{12} & a_{13} \\
a_{21} & a_{22} & a_{23} \\
a_{31} & a_{32} & a_{33}
\end{pmatrix}
\]

\section{Tables}
Below is a sample table showing data:

\begin{table}[h]
\centering
\begin{tabular}{@{}lrrr@{}}
\toprule
\textbf{Category} & \textbf{Value 1} & \textbf{Value 2} & \textbf{Total} \\
\midrule
Item A & 10.5 & 20.3 & 30.8 \\
Item B & 15.2 & 25.1 & 40.3 \\
Item C & 22.1 & 12.7 & 34.8 \\
\bottomrule
\end{tabular}
\caption{Sample Data Table}
\label{tab:sample}
\end{table}

\section{Lists}
\subsection{Bullet Points}
\begin{itemize}
  \item First item
  \item Second item
  \item Third item with sub-items:
    \begin{itemize}
      \item Sub-item 1
      \item Sub-item 2
    \end{itemize}
\end{itemize}

\subsection{Numbered List}
\begin{enumerate}
  \item First step
  \item Second step
  \item Third step
\end{enumerate}

\section{Citations and References}
When you use the citation search feature, references will be automatically managed. For example, you might cite a paper like this: \cite{example2024}. All citations will be automatically added to your \texttt{cite.bib} file.

\section{Conclusion}
This template demonstrates basic \LaTeX{} features including mathematical formulas, tables, and lists. The bibliography system is now set up to work seamlessly with Bibby AI's citation management. Feel free to modify it for your needs!

\bibliographystyle{plain}
\bibliography{cite}

\end{document}\documentclass{article}
\usepackage[utf8]{inputenc}
\usepackage{booktabs}
\usepackage{amsmath}
\usepackage{natbib}

\title{My Document}
\author{Your Name}
\date{\today}

\begin{document}

\maketitle

\section{Introduction}
This is a sample document demonstrating various \LaTeX{} features including mathematical equations, tables, and formatting. You can cite papers using the citation search feature, and they will be automatically added to your \texttt{cite.bib} file.

\section{Mathematical Expressions}
Here are some example equations:

\subsection{Quadratic Formula}
The quadratic formula is given by:
\[x = \frac{-b \pm \sqrt{b^2 - 4ac}}{2a}\]

\subsection{Matrix Example}
A sample matrix:
\[
\begin{pmatrix}
a_{11} & a_{12} & a_{13} \\
a_{21} & a_{22} & a_{23} \\
a_{31} & a_{32} & a_{33}
\end{pmatrix}
\]

\section{Tables}
Below is a sample table showing data:

\begin{table}[h]
\centering
\begin{tabular}{@{}lrrr@{}}
\toprule
\textbf{Category} & \textbf{Value 1} & \textbf{Value 2} & \textbf{Total} \\
\midrule
Item A & 10.5 & 20.3 & 30.8 \\
Item B & 15.2 & 25.1 & 40.3 \\
Item C & 22.1 & 12.7 & 34.8 \\
\bottomrule
\end{tabular}
\caption{Sample Data Table}
\label{tab:sample}
\end{table}

\section{Lists}
\subsection{Bullet Points}
\begin{itemize}
  \item First item
  \item Second item
  \item Third item with sub-items:
    \begin{itemize}
      \item Sub-item 1
      \item Sub-item 2
    \end{itemize}
\end{itemize}

\subsection{Numbered List}
\begin{enumerate}
  \item First step
  \item Second step
  \item Third step
\end{enumerate}

\section{Citations and References}
When you use the citation search feature, references will be automatically managed. For example, you might cite a paper like this: \cite{example2024}. All citations will be automatically added to your \texttt{cite.bib} file.

\section{Conclusion}
This template demonstrates basic \LaTeX{} features including mathematical formulas, tables, and lists. The bibliography system is now set up to work seamlessly with Bibby AI's citation management. Feel free to modify it for your needs!

\bibliographystyle{plain}
\bibliography{cite}

\end{document}\documentclass{article}
\usepackage[utf8]{inputenc}
\usepackage{booktabs}
\usepackage{amsmath}
\usepackage{natbib}

\title{My Document}
\author{Your Name}
\date{\today}

\begin{document}

\maketitle

\section{Introduction}
This is a sample document demonstrating various \LaTeX{} features including mathematical equations, tables, and formatting. You can cite papers using the citation search feature, and they will be automatically added to your \texttt{cite.bib} file.

\section{Mathematical Expressions}
Here are some example equations:

\subsection{Quadratic Formula}
The quadratic formula is given by:
\[x = \frac{-b \pm \sqrt{b^2 - 4ac}}{2a}\]

\subsection{Matrix Example}
A sample matrix:
\[
\begin{pmatrix}
a_{11} & a_{12} & a_{13} \\
a_{21} & a_{22} & a_{23} \\
a_{31} & a_{32} & a_{33}
\end{pmatrix}
\]

\section{Tables}
Below is a sample table showing data:

\begin{table}[h]
\centering
\begin{tabular}{@{}lrrr@{}}
\toprule
\textbf{Category} & \textbf{Value 1} & \textbf{Value 2} & \textbf{Total} \\
\midrule
Item A & 10.5 & 20.3 & 30.8 \\
Item B & 15.2 & 25.1 & 40.3 \\
Item C & 22.1 & 12.7 & 34.8 \\
\bottomrule
\end{tabular}
\caption{Sample Data Table}
\label{tab:sample}
\end{table}

\section{Lists}
\subsection{Bullet Points}
\begin{itemize}
  \item First item
  \item Second item
  \item Third item with sub-items:
    \begin{itemize}
      \item Sub-item 1
      \item Sub-item 2
    \end{itemize}
\end{itemize}

\subsection{Numbered List}
\begin{enumerate}
  \item First step
  \item Second step
  \item Third step
\end{enumerate}

\section{Citations and References}
When you use the citation search feature, references will be automatically managed. For example, you might cite a paper like this: \cite{example2024}. All citations will be automatically added to your \texttt{cite.bib} file.

\section{Conclusion}
This template demonstrates basic \LaTeX{} features including mathematical formulas, tables, and lists. The bibliography system is now set up to work seamlessly with Bibby AI's citation management. Feel free to modify it for your needs!

\bibliographystyle{plain}
\bibliography{cite}

\end{document}\documentclass{article}
\usepackage[utf8]{inputenc}
\usepackage{booktabs}
\usepackage{amsmath}
\usepackage{natbib}

\title{My Document}
\author{Your Name}
\date{\today}

\begin{document}

\maketitle

\section{Introduction}
This is a sample document demonstrating various \LaTeX{} features including mathematical equations, tables, and formatting. You can cite papers using the citation search feature, and they will be automatically added to your \texttt{cite.bib} file.

\section{Mathematical Expressions}
Here are some example equations:

\subsection{Quadratic Formula}
The quadratic formula is given by:
\[x = \frac{-b \pm \sqrt{b^2 - 4ac}}{2a}\]

\subsection{Matrix Example}
A sample matrix:
\[
\begin{pmatrix}
a_{11} & a_{12} & a_{13} \\
a_{21} & a_{22} & a_{23} \\
a_{31} & a_{32} & a_{33}
\end{pmatrix}
\]

\section{Tables}
Below is a sample table showing data:

\begin{table}[h]
\centering
\begin{tabular}{@{}lrrr@{}}
\toprule
\textbf{Category} & \textbf{Value 1} & \textbf{Value 2} & \textbf{Total} \\
\midrule
Item A & 10.5 & 20.3 & 30.8 \\
Item B & 15.2 & 25.1 & 40.3 \\
Item C & 22.1 & 12.7 & 34.8 \\
\bottomrule
\end{tabular}
\caption{Sample Data Table}
\label{tab:sample}
\end{table}

\section{Lists}
\subsection{Bullet Points}
\begin{itemize}
  \item First item
  \item Second item
  \item Third item with sub-items:
    \begin{itemize}
      \item Sub-item 1
      \item Sub-item 2
    \end{itemize}
\end{itemize}

\subsection{Numbered List}
\begin{enumerate}
  \item First step
  \item Second step
  \item Third step
\end{enumerate}

\section{Citations and References}
When you use the citation search feature, references will be automatically managed. For example, you might cite a paper like this: \cite{example2024}. All citations will be automatically added to your \texttt{cite.bib} file.

\section{Conclusion}
This template demonstrates basic \LaTeX{} features including mathematical formulas, tables, and lists. The bibliography system is now set up to work seamlessly with Bibby AI's citation management. Feel free to modify it for your needs!

\bibliographystyle{plain}
\bibliography{cite}

\end{document}\documentclass{article}
\usepackage[utf8]{inputenc}
\usepackage{booktabs}
\usepackage{amsmath}
\usepackage{natbib}

\title{My Document}
\author{Your Name}
\date{\today}

\begin{document}

\maketitle

\section{Introduction}
This is a sample document demonstrating various \LaTeX{} features including mathematical equations, tables, and formatting. You can cite papers using the citation search feature, and they will be automatically added to your \texttt{cite.bib} file.

\section{Mathematical Expressions}
Here are some example equations:

\subsection{Quadratic Formula}
The quadratic formula is given by:
\[x = \frac{-b \pm \sqrt{b^2 - 4ac}}{2a}\]

\subsection{Matrix Example}
A sample matrix:
\[
\begin{pmatrix}
a_{11} & a_{12} & a_{13} \\
a_{21} & a_{22} & a_{23} \\
a_{31} & a_{32} & a_{33}
\end{pmatrix}
\]

\section{Tables}
Below is a sample table showing data:

\begin{table}[h]
\centering
\begin{tabular}{@{}lrrr@{}}
\toprule
\textbf{Category} & \textbf{Value 1} & \textbf{Value 2} & \textbf{Total} \\
\midrule
Item A & 10.5 & 20.3 & 30.8 \\
Item B & 15.2 & 25.1 & 40.3 \\
Item C & 22.1 & 12.7 & 34.8 \\
\bottomrule
\end{tabular}
\caption{Sample Data Table}
\label{tab:sample}
\end{table}

\section{Lists}
\subsection{Bullet Points}
\begin{itemize}
  \item First item
  \item Second item
  \item Third item with sub-items:
    \begin{itemize}
      \item Sub-item 1
      \item Sub-item 2
    \end{itemize}
\end{itemize}

\subsection{Numbered List}
\begin{enumerate}
  \item First step
  \item Second step
  \item Third step
\end{enumerate}

\section{Citations and References}
When you use the citation search feature, references will be automatically managed. For example, you might cite a paper like this: \cite{example2024}. All citations will be automatically added to your \texttt{cite.bib} file.

\section{Conclusion}
This template demonstrates basic \LaTeX{} features including mathematical formulas, tables, and lists. The bibliography system is now set up to work seamlessly with Bibby AI's citation management. Feel free to modify it for your needs!

\bibliographystyle{plain}
\bibliography{cite}

\end{document}\documentclass{article}
\usepackage[utf8]{inputenc}
\usepackage{booktabs}
\usepackage{amsmath}
\usepackage{natbib}

\title{My Document}
\author{Your Name}
\date{\today}

\begin{document}

\maketitle

\section{Introduction}
This is a sample document demonstrating various \LaTeX{} features including mathematical equations, tables, and formatting. You can cite papers using the citation search feature, and they will be automatically added to your \texttt{cite.bib} file.

\section{Mathematical Expressions}
Here are some example equations:

\subsection{Quadratic Formula}
The quadratic formula is given by:
\[x = \frac{-b \pm \sqrt{b^2 - 4ac}}{2a}\]

\subsection{Matrix Example}
A sample matrix:
\[
\begin{pmatrix}
a_{11} & a_{12} & a_{13} \\
a_{21} & a_{22} & a_{23} \\
a_{31} & a_{32} & a_{33}
\end{pmatrix}
\]

\section{Tables}
Below is a sample table showing data:

\begin{table}[h]
\centering
\begin{tabular}{@{}lrrr@{}}
\toprule
\textbf{Category} & \textbf{Value 1} & \textbf{Value 2} & \textbf{Total} \\
\midrule
Item A & 10.5 & 20.3 & 30.8 \\
Item B & 15.2 & 25.1 & 40.3 \\
Item C & 22.1 & 12.7 & 34.8 \\
\bottomrule
\end{tabular}
\caption{Sample Data Table}
\label{tab:sample}
\end{table}

\section{Lists}
\subsection{Bullet Points}
\begin{itemize}
  \item First item
  \item Second item
  \item Third item with sub-items:
    \begin{itemize}
      \item Sub-item 1
      \item Sub-item 2
    \end{itemize}
\end{itemize}

\subsection{Numbered List}
\begin{enumerate}
  \item First step
  \item Second step
  \item Third step
\end{enumerate}

\section{Citations and References}
When you use the citation search feature, references will be automatically managed. For example, you might cite a paper like this: \cite{example2024}. All citations will be automatically added to your \texttt{cite.bib} file.

\section{Conclusion}
This template demonstrates basic \LaTeX{} features including mathematical formulas, tables, and lists. The bibliography system is now set up to work seamlessly with Bibby AI's citation management. Feel free to modify it for your needs!

\bibliographystyle{plain}
\bibliography{cite}

\end{document}\documentclass{article}
\usepackage[utf8]{inputenc}
\usepackage{booktabs}
\usepackage{amsmath}
\usepackage{natbib}

\title{My Document}
\author{Your Name}
\date{\today}

\begin{document}

\maketitle

\section{Introduction}
This is a sample document demonstrating various \LaTeX{} features including mathematical equations, tables, and formatting. You can cite papers using the citation search feature, and they will be automatically added to your \texttt{cite.bib} file.

\section{Mathematical Expressions}
Here are some example equations:

\subsection{Quadratic Formula}
The quadratic formula is given by:
\[x = \frac{-b \pm \sqrt{b^2 - 4ac}}{2a}\]

\subsection{Matrix Example}
A sample matrix:
\[
\begin{pmatrix}
a_{11} & a_{12} & a_{13} \\
a_{21} & a_{22} & a_{23} \\
a_{31} & a_{32} & a_{33}
\end{pmatrix}
\]

\section{Tables}
Below is a sample table showing data:

\begin{table}[h]
\centering
\begin{tabular}{@{}lrrr@{}}
\toprule
\textbf{Category} & \textbf{Value 1} & \textbf{Value 2} & \textbf{Total} \\
\midrule
Item A & 10.5 & 20.3 & 30.8 \\
Item B & 15.2 & 25.1 & 40.3 \\
Item C & 22.1 & 12.7 & 34.8 \\
\bottomrule
\end{tabular}
\caption{Sample Data Table}
\label{tab:sample}
\end{table}

\section{Lists}
\subsection{Bullet Points}
\begin{itemize}
  \item First item
  \item Second item
  \item Third item with sub-items:
    \begin{itemize}
      \item Sub-item 1
      \item Sub-item 2
    \end{itemize}
\end{itemize}

\subsection{Numbered List}
\begin{enumerate}
  \item First step
  \item Second step
  \item Third step
\end{enumerate}

\section{Citations and References}
When you use the citation search feature, references will be automatically managed. For example, you might cite a paper like this: \cite{example2024}. All citations will be automatically added to your \texttt{cite.bib} file.

\section{Conclusion}
This template demonstrates basic \LaTeX{} features including mathematical formulas, tables, and lists. The bibliography system is now set up to work seamlessly with Bibby AI's citation management. Feel free to modify it for your needs!

\bibliographystyle{plain}
\bibliography{cite}

\end{document}\documentclass{article}
\usepackage[utf8]{inputenc}
\usepackage{booktabs}
\usepackage{amsmath}
\usepackage{natbib}

\title{My Document}
\author{Your Name}
\date{\today}

\begin{document}

\maketitle

\section{Introduction}
This is a sample document demonstrating various \LaTeX{} features including mathematical equations, tables, and formatting. You can cite papers using the citation search feature, and they will be automatically added to your \texttt{cite.bib} file.

\section{Mathematical Expressions}
Here are some example equations:

\subsection{Quadratic Formula}
The quadratic formula is given by:
\[x = \frac{-b \pm \sqrt{b^2 - 4ac}}{2a}\]

\subsection{Matrix Example}
A sample matrix:
\[
\begin{pmatrix}
a_{11} & a_{12} & a_{13} \\
a_{21} & a_{22} & a_{23} \\
a_{31} & a_{32} & a_{33}
\end{pmatrix}
\]

\section{Tables}
Below is a sample table showing data:

\begin{table}[h]
\centering
\begin{tabular}{@{}lrrr@{}}
\toprule
\textbf{Category} & \textbf{Value 1} & \textbf{Value 2} & \textbf{Total} \\
\midrule
Item A & 10.5 & 20.3 & 30.8 \\
Item B & 15.2 & 25.1 & 40.3 \\
Item C & 22.1 & 12.7 & 34.8 \\
\bottomrule
\end{tabular}
\caption{Sample Data Table}
\label{tab:sample}
\end{table}

\section{Lists}
\subsection{Bullet Points}
\begin{itemize}
  \item First item
  \item Second item
  \item Third item with sub-items:
    \begin{itemize}
      \item Sub-item 1
      \item Sub-item 2
    \end{itemize}
\end{itemize}

\subsection{Numbered List}
\begin{enumerate}
  \item First step
  \item Second step
  \item Third step
\end{enumerate}

\section{Citations and References}
When you use the citation search feature, references will be automatically managed. For example, you might cite a paper like this: \cite{example2024}. All citations will be automatically added to your \texttt{cite.bib} file.

\section{Conclusion}
This template demonstrates basic \LaTeX{} features including mathematical formulas, tables, and lists. The bibliography system is now set up to work seamlessly with Bibby AI's citation management. Feel free to modify it for your needs!

\bibliographystyle{plain}
\bibliography{cite}

\end{document}\documentclass{article}
\usepackage[utf8]{inputenc}
\usepackage{booktabs}
\usepackage{amsmath}
\usepackage{natbib}

\title{My Document}
\author{Your Name}
\date{\today}

\begin{document}

\maketitle

\section{Introduction}
This is a sample document demonstrating various \LaTeX{} features including mathematical equations, tables, and formatting. You can cite papers using the citation search feature, and they will be automatically added to your \texttt{cite.bib} file.

\section{Mathematical Expressions}
Here are some example equations:

\subsection{Quadratic Formula}
The quadratic formula is given by:
\[x = \frac{-b \pm \sqrt{b^2 - 4ac}}{2a}\]

\subsection{Matrix Example}
A sample matrix:
\[
\begin{pmatrix}
a_{11} & a_{12} & a_{13} \\
a_{21} & a_{22} & a_{23} \\
a_{31} & a_{32} & a_{33}
\end{pmatrix}
\]

\section{Tables}
Below is a sample table showing data:

\begin{table}[h]
\centering
\begin{tabular}{@{}lrrr@{}}
\toprule
\textbf{Category} & \textbf{Value 1} & \textbf{Value 2} & \textbf{Total} \\
\midrule
Item A & 10.5 & 20.3 & 30.8 \\
Item B & 15.2 & 25.1 & 40.3 \\
Item C & 22.1 & 12.7 & 34.8 \\
\bottomrule
\end{tabular}
\caption{Sample Data Table}
\label{tab:sample}
\end{table}

\section{Lists}
\subsection{Bullet Points}
\begin{itemize}
  \item First item
  \item Second item
  \item Third item with sub-items:
    \begin{itemize}
      \item Sub-item 1
      \item Sub-item 2
    \end{itemize}
\end{itemize}

\subsection{Numbered List}
\begin{enumerate}
  \item First step
  \item Second step
  \item Third step
\end{enumerate}

\section{Citations and References}
When you use the citation search feature, references will be automatically managed. For example, you might cite a paper like this: \cite{example2024}. All citations will be automatically added to your \texttt{cite.bib} file.

\section{Conclusion}
This template demonstrates basic \LaTeX{} features including mathematical formulas, tables, and lists. The bibliography system is now set up to work seamlessly with Bibby AI's citation management. Feel free to modify it for your needs!

\bibliographystyle{plain}
\bibliography{cite}

\end{document}\documentclass{article}
\usepackage[utf8]{inputenc}
\usepackage{booktabs}
\usepackage{amsmath}
\usepackage{natbib}

\title{My Document}
\author{Your Name}
\date{\today}

\begin{document}

\maketitle

\section{Introduction}
This is a sample document demonstrating various \LaTeX{} features including mathematical equations, tables, and formatting. You can cite papers using the citation search feature, and they will be automatically added to your \texttt{cite.bib} file.

\section{Mathematical Expressions}
Here are some example equations:

\subsection{Quadratic Formula}
The quadratic formula is given by:
\[x = \frac{-b \pm \sqrt{b^2 - 4ac}}{2a}\]

\subsection{Matrix Example}
A sample matrix:
\[
\begin{pmatrix}
a_{11} & a_{12} & a_{13} \\
a_{21} & a_{22} & a_{23} \\
a_{31} & a_{32} & a_{33}
\end{pmatrix}
\]

\section{Tables}
Below is a sample table showing data:

\begin{table}[h]
\centering
\begin{tabular}{@{}lrrr@{}}
\toprule
\textbf{Category} & \textbf{Value 1} & \textbf{Value 2} & \textbf{Total} \\
\midrule
Item A & 10.5 & 20.3 & 30.8 \\
Item B & 15.2 & 25.1 & 40.3 \\
Item C & 22.1 & 12.7 & 34.8 \\
\bottomrule
\end{tabular}
\caption{Sample Data Table}
\label{tab:sample}
\end{table}

\section{Lists}
\subsection{Bullet Points}
\begin{itemize}
  \item First item
  \item Second item
  \item Third item with sub-items:
    \begin{itemize}
      \item Sub-item 1
      \item Sub-item 2
    \end{itemize}
\end{itemize}

\subsection{Numbered List}
\begin{enumerate}
  \item First step
  \item Second step
  \item Third step
\end{enumerate}

\section{Citations and References}
When you use the citation search feature, references will be automatically managed. For example, you might cite a paper like this: \cite{example2024}. All citations will be automatically added to your \texttt{cite.bib} file.

\section{Conclusion}
This template demonstrates basic \LaTeX{} features including mathematical formulas, tables, and lists. The bibliography system is now set up to work seamlessly with Bibby AI's citation management. Feel free to modify it for your needs!

\bibliographystyle{plain}
\bibliography{cite}

\end{document}\documentclass{article}
\usepackage[utf8]{inputenc}
\usepackage{booktabs}
\usepackage{amsmath}
\usepackage{natbib}

\title{My Document}
\author{Your Name}
\date{\today}

\begin{document}

\maketitle

\section{Introduction}
This is a sample document demonstrating various \LaTeX{} features including mathematical equations, tables, and formatting. You can cite papers using the citation search feature, and they will be automatically added to your \texttt{cite.bib} file.

\section{Mathematical Expressions}
Here are some example equations:

\subsection{Quadratic Formula}
The quadratic formula is given by:
\[x = \frac{-b \pm \sqrt{b^2 - 4ac}}{2a}\]

\subsection{Matrix Example}
A sample matrix:
\[
\begin{pmatrix}
a_{11} & a_{12} & a_{13} \\
a_{21} & a_{22} & a_{23} \\
a_{31} & a_{32} & a_{33}
\end{pmatrix}
\]

\section{Tables}
Below is a sample table showing data:

\begin{table}[h]
\centering
\begin{tabular}{@{}lrrr@{}}
\toprule
\textbf{Category} & \textbf{Value 1} & \textbf{Value 2} & \textbf{Total} \\
\midrule
Item A & 10.5 & 20.3 & 30.8 \\
Item B & 15.2 & 25.1 & 40.3 \\
Item C & 22.1 & 12.7 & 34.8 \\
\bottomrule
\end{tabular}
\caption{Sample Data Table}
\label{tab:sample}
\end{table}

\section{Lists}
\subsection{Bullet Points}
\begin{itemize}
  \item First item
  \item Second item
  \item Third item with sub-items:
    \begin{itemize}
      \item Sub-item 1
      \item Sub-item 2
    \end{itemize}
\end{itemize}

\subsection{Numbered List}
\begin{enumerate}
  \item First step
  \item Second step
  \item Third step
\end{enumerate}

\section{Citations and References}
When you use the citation search feature, references will be automatically managed. For example, you might cite a paper like this: \cite{example2024}. All citations will be automatically added to your \texttt{cite.bib} file.

\section{Conclusion}
This template demonstrates basic \LaTeX{} features including mathematical formulas, tables, and lists. The bibliography system is now set up to work seamlessly with Bibby AI's citation management. Feel free to modify it for your needs!

\bibliographystyle{plain}
\bibliography{cite}

\end{document}\documentclass{article}
\usepackage[utf8]{inputenc}
\usepackage{booktabs}
\usepackage{amsmath}
\usepackage{natbib}

\title{My Document}
\author{Your Name}
\date{\today}

\begin{document}

\maketitle

\section{Introduction}
This is a sample document demonstrating various \LaTeX{} features including mathematical equations, tables, and formatting. You can cite papers using the citation search feature, and they will be automatically added to your \texttt{cite.bib} file.

\section{Mathematical Expressions}
Here are some example equations:

\subsection{Quadratic Formula}
The quadratic formula is given by:
\[x = \frac{-b \pm \sqrt{b^2 - 4ac}}{2a}\]

\subsection{Matrix Example}
A sample matrix:
\[
\begin{pmatrix}
a_{11} & a_{12} & a_{13} \\
a_{21} & a_{22} & a_{23} \\
a_{31} & a_{32} & a_{33}
\end{pmatrix}
\]

\section{Tables}
Below is a sample table showing data:

\begin{table}[h]
\centering
\begin{tabular}{@{}lrrr@{}}
\toprule
\textbf{Category} & \textbf{Value 1} & \textbf{Value 2} & \textbf{Total} \\
\midrule
Item A & 10.5 & 20.3 & 30.8 \\
Item B & 15.2 & 25.1 & 40.3 \\
Item C & 22.1 & 12.7 & 34.8 \\
\bottomrule
\end{tabular}
\caption{Sample Data Table}
\label{tab:sample}
\end{table}

\section{Lists}
\subsection{Bullet Points}
\begin{itemize}
  \item First item
  \item Second item
  \item Third item with sub-items:
    \begin{itemize}
      \item Sub-item 1
      \item Sub-item 2
    \end{itemize}
\end{itemize}

\subsection{Numbered List}
\begin{enumerate}
  \item First step
  \item Second step
  \item Third step
\end{enumerate}

\section{Citations and References}
When you use the citation search feature, references will be automatically managed. For example, you might cite a paper like this: \cite{example2024}. All citations will be automatically added to your \texttt{cite.bib} file.

\section{Conclusion}
This template demonstrates basic \LaTeX{} features including mathematical formulas, tables, and lists. The bibliography system is now set up to work seamlessly with Bibby AI's citation management. Feel free to modify it for your needs!

\bibliographystyle{plain}
\bibliography{cite}

\end{document}\documentclass{article}
\usepackage[utf8]{inputenc}
\usepackage{booktabs}
\usepackage{amsmath}
\usepackage{natbib}

\title{My Document}
\author{Your Name}
\date{\today}

\begin{document}

\maketitle

\section{Introduction}
This is a sample document demonstrating various \LaTeX{} features including mathematical equations, tables, and formatting. You can cite papers using the citation search feature, and they will be automatically added to your \texttt{cite.bib} file.

\section{Mathematical Expressions}
Here are some example equations:

\subsection{Quadratic Formula}
The quadratic formula is given by:
\[x = \frac{-b \pm \sqrt{b^2 - 4ac}}{2a}\]

\subsection{Matrix Example}
A sample matrix:
\[
\begin{pmatrix}
a_{11} & a_{12} & a_{13} \\
a_{21} & a_{22} & a_{23} \\
a_{31} & a_{32} & a_{33}
\end{pmatrix}
\]

\section{Tables}
Below is a sample table showing data:

\begin{table}[h]
\centering
\begin{tabular}{@{}lrrr@{}}
\toprule
\textbf{Category} & \textbf{Value 1} & \textbf{Value 2} & \textbf{Total} \\
\midrule
Item A & 10.5 & 20.3 & 30.8 \\
Item B & 15.2 & 25.1 & 40.3 \\
Item C & 22.1 & 12.7 & 34.8 \\
\bottomrule
\end{tabular}
\caption{Sample Data Table}
\label{tab:sample}
\end{table}

\section{Lists}
\subsection{Bullet Points}
\begin{itemize}
  \item First item
  \item Second item
  \item Third item with sub-items:
    \begin{itemize}
      \item Sub-item 1
      \item Sub-item 2
    \end{itemize}
\end{itemize}

\subsection{Numbered List}
\begin{enumerate}
  \item First step
  \item Second step
  \item Third step
\end{enumerate}

\section{Citations and References}
When you use the citation search feature, references will be automatically managed. For example, you might cite a paper like this: \cite{example2024}. All citations will be automatically added to your \texttt{cite.bib} file.

\section{Conclusion}
This template demonstrates basic \LaTeX{} features including mathematical formulas, tables, and lists. The bibliography system is now set up to work seamlessly with Bibby AI's citation management. Feel free to modify it for your needs!

\bibliographystyle{plain}
\bibliography{cite}

\end{document}\documentclass{article}
\usepackage[utf8]{inputenc}
\usepackage{booktabs}
\usepackage{amsmath}
\usepackage{natbib}

\title{My Document}
\author{Your Name}
\date{\today}

\begin{document}

\maketitle

\section{Introduction}
This is a sample document demonstrating various \LaTeX{} features including mathematical equations, tables, and formatting. You can cite papers using the citation search feature, and they will be automatically added to your \texttt{cite.bib} file.

\section{Mathematical Expressions}
Here are some example equations:

\subsection{Quadratic Formula}
The quadratic formula is given by:
\[x = \frac{-b \pm \sqrt{b^2 - 4ac}}{2a}\]

\subsection{Matrix Example}
A sample matrix:
\[
\begin{pmatrix}
a_{11} & a_{12} & a_{13} \\
a_{21} & a_{22} & a_{23} \\
a_{31} & a_{32} & a_{33}
\end{pmatrix}
\]

\section{Tables}
Below is a sample table showing data:

\begin{table}[h]
\centering
\begin{tabular}{@{}lrrr@{}}
\toprule
\textbf{Category} & \textbf{Value 1} & \textbf{Value 2} & \textbf{Total} \\
\midrule
Item A & 10.5 & 20.3 & 30.8 \\
Item B & 15.2 & 25.1 & 40.3 \\
Item C & 22.1 & 12.7 & 34.8 \\
\bottomrule
\end{tabular}
\caption{Sample Data Table}
\label{tab:sample}
\end{table}

\section{Lists}
\subsection{Bullet Points}
\begin{itemize}
  \item First item
  \item Second item
  \item Third item with sub-items:
    \begin{itemize}
      \item Sub-item 1
      \item Sub-item 2
    \end{itemize}
\end{itemize}

\subsection{Numbered List}
\begin{enumerate}
  \item First step
  \item Second step
  \item Third step
\end{enumerate}

\section{Citations and References}
When you use the citation search feature, references will be automatically managed. For example, you might cite a paper like this: \cite{example2024}. All citations will be automatically added to your \texttt{cite.bib} file.

\section{Conclusion}
This template demonstrates basic \LaTeX{} features including mathematical formulas, tables, and lists. The bibliography system is now set up to work seamlessly with Bibby AI's citation management. Feel free to modify it for your needs!

\bibliographystyle{plain}
\bibliography{cite}

\end{document}\documentclass{article}
\usepackage[utf8]{inputenc}
\usepackage{booktabs}
\usepackage{amsmath}
\usepackage{natbib}

\title{My Document}
\author{Your Name}
\date{\today}

\begin{document}

\maketitle

\section{Introduction}
This is a sample document demonstrating various \LaTeX{} features including mathematical equations, tables, and formatting. You can cite papers using the citation search feature, and they will be automatically added to your \texttt{cite.bib} file.

\section{Mathematical Expressions}
Here are some example equations:

\subsection{Quadratic Formula}
The quadratic formula is given by:
\[x = \frac{-b \pm \sqrt{b^2 - 4ac}}{2a}\]

\subsection{Matrix Example}
A sample matrix:
\[
\begin{pmatrix}
a_{11} & a_{12} & a_{13} \\
a_{21} & a_{22} & a_{23} \\
a_{31} & a_{32} & a_{33}
\end{pmatrix}
\]

\section{Tables}
Below is a sample table showing data:

\begin{table}[h]
\centering
\begin{tabular}{@{}lrrr@{}}
\toprule
\textbf{Category} & \textbf{Value 1} & \textbf{Value 2} & \textbf{Total} \\
\midrule
Item A & 10.5 & 20.3 & 30.8 \\
Item B & 15.2 & 25.1 & 40.3 \\
Item C & 22.1 & 12.7 & 34.8 \\
\bottomrule
\end{tabular}
\caption{Sample Data Table}
\label{tab:sample}
\end{table}

\section{Lists}
\subsection{Bullet Points}
\begin{itemize}
  \item First item
  \item Second item
  \item Third item with sub-items:
    \begin{itemize}
      \item Sub-item 1
      \item Sub-item 2
    \end{itemize}
\end{itemize}

\subsection{Numbered List}
\begin{enumerate}
  \item First step
  \item Second step
  \item Third step
\end{enumerate}

\section{Citations and References}
When you use the citation search feature, references will be automatically managed. For example, you might cite a paper like this: \cite{example2024}. All citations will be automatically added to your \texttt{cite.bib} file.

\section{Conclusion}
This template demonstrates basic \LaTeX{} features including mathematical formulas, tables, and lists. The bibliography system is now set up to work seamlessly with Bibby AI's citation management. Feel free to modify it for your needs!

\bibliographystyle{plain}
\bibliography{cite}

\end{document}\documentclass{article}
\usepackage[utf8]{inputenc}
\usepackage{booktabs}
\usepackage{amsmath}
\usepackage{natbib}

\title{My Document}
\author{Your Name}
\date{\today}

\begin{document}

\maketitle

\section{Introduction}
This is a sample document demonstrating various \LaTeX{} features including mathematical equations, tables, and formatting. You can cite papers using the citation search feature, and they will be automatically added to your \texttt{cite.bib} file.

\section{Mathematical Expressions}
Here are some example equations:

\subsection{Quadratic Formula}
The quadratic formula is given by:
\[x = \frac{-b \pm \sqrt{b^2 - 4ac}}{2a}\]

\subsection{Matrix Example}
A sample matrix:
\[
\begin{pmatrix}
a_{11} & a_{12} & a_{13} \\
a_{21} & a_{22} & a_{23} \\
a_{31} & a_{32} & a_{33}
\end{pmatrix}
\]

\section{Tables}
Below is a sample table showing data:

\begin{table}[h]
\centering
\begin{tabular}{@{}lrrr@{}}
\toprule
\textbf{Category} & \textbf{Value 1} & \textbf{Value 2} & \textbf{Total} \\
\midrule
Item A & 10.5 & 20.3 & 30.8 \\
Item B & 15.2 & 25.1 & 40.3 \\
Item C & 22.1 & 12.7 & 34.8 \\
\bottomrule
\end{tabular}
\caption{Sample Data Table}
\label{tab:sample}
\end{table}

\section{Lists}
\subsection{Bullet Points}
\begin{itemize}
  \item First item
  \item Second item
  \item Third item with sub-items:
    \begin{itemize}
      \item Sub-item 1
      \item Sub-item 2
    \end{itemize}
\end{itemize}

\subsection{Numbered List}
\begin{enumerate}
  \item First step
  \item Second step
  \item Third step
\end{enumerate}

\section{Citations and References}
When you use the citation search feature, references will be automatically managed. For example, you might cite a paper like this: \cite{example2024}. All citations will be automatically added to your \texttt{cite.bib} file.

\section{Conclusion}
This template demonstrates basic \LaTeX{} features including mathematical formulas, tables, and lists. The bibliography system is now set up to work seamlessly with Bibby AI's citation management. Feel free to modify it for your needs!

\bibliographystyle{plain}
\bibliography{cite}

\end{document}\documentclass{article}
\usepackage[utf8]{inputenc}
\usepackage{booktabs}
\usepackage{amsmath}
\usepackage{natbib}

\title{My Document}
\author{Your Name}
\date{\today}

\begin{document}

\maketitle

\section{Introduction}
This is a sample document demonstrating various \LaTeX{} features including mathematical equations, tables, and formatting. You can cite papers using the citation search feature, and they will be automatically added to your \texttt{cite.bib} file.

\section{Mathematical Expressions}
Here are some example equations:

\subsection{Quadratic Formula}
The quadratic formula is given by:
\[x = \frac{-b \pm \sqrt{b^2 - 4ac}}{2a}\]

\subsection{Matrix Example}
A sample matrix:
\[
\begin{pmatrix}
a_{11} & a_{12} & a_{13} \\
a_{21} & a_{22} & a_{23} \\
a_{31} & a_{32} & a_{33}
\end{pmatrix}
\]

\section{Tables}
Below is a sample table showing data:

\begin{table}[h]
\centering
\begin{tabular}{@{}lrrr@{}}
\toprule
\textbf{Category} & \textbf{Value 1} & \textbf{Value 2} & \textbf{Total} \\
\midrule
Item A & 10.5 & 20.3 & 30.8 \\
Item B & 15.2 & 25.1 & 40.3 \\
Item C & 22.1 & 12.7 & 34.8 \\
\bottomrule
\end{tabular}
\caption{Sample Data Table}
\label{tab:sample}
\end{table}

\section{Lists}
\subsection{Bullet Points}
\begin{itemize}
  \item First item
  \item Second item
  \item Third item with sub-items:
    \begin{itemize}
      \item Sub-item 1
      \item Sub-item 2
    \end{itemize}
\end{itemize}

\subsection{Numbered List}
\begin{enumerate}
  \item First step
  \item Second step
  \item Third step
\end{enumerate}

\section{Citations and References}
When you use the citation search feature, references will be automatically managed. For example, you might cite a paper like this: \cite{example2024}. All citations will be automatically added to your \texttt{cite.bib} file.

\section{Conclusion}
This template demonstrates basic \LaTeX{} features including mathematical formulas, tables, and lists. The bibliography system is now set up to work seamlessly with Bibby AI's citation management. Feel free to modify it for your needs!

\bibliographystyle{plain}
\bibliography{cite}

\end{document}\documentclass{article}
\usepackage[utf8]{inputenc}
\usepackage{booktabs}
\usepackage{amsmath}
\usepackage{natbib}

\title{My Document}
\author{Your Name}
\date{\today}

\begin{document}

\maketitle

\section{Introduction}
This is a sample document demonstrating various \LaTeX{} features including mathematical equations, tables, and formatting. You can cite papers using the citation search feature, and they will be automatically added to your \texttt{cite.bib} file.

\section{Mathematical Expressions}
Here are some example equations:

\subsection{Quadratic Formula}
The quadratic formula is given by:
\[x = \frac{-b \pm \sqrt{b^2 - 4ac}}{2a}\]

\subsection{Matrix Example}
A sample matrix:
\[
\begin{pmatrix}
a_{11} & a_{12} & a_{13} \\
a_{21} & a_{22} & a_{23} \\
a_{31} & a_{32} & a_{33}
\end{pmatrix}
\]

\section{Tables}
Below is a sample table showing data:

\begin{table}[h]
\centering
\begin{tabular}{@{}lrrr@{}}
\toprule
\textbf{Category} & \textbf{Value 1} & \textbf{Value 2} & \textbf{Total} \\
\midrule
Item A & 10.5 & 20.3 & 30.8 \\
Item B & 15.2 & 25.1 & 40.3 \\
Item C & 22.1 & 12.7 & 34.8 \\
\bottomrule
\end{tabular}
\caption{Sample Data Table}
\label{tab:sample}
\end{table}

\section{Lists}
\subsection{Bullet Points}
\begin{itemize}
  \item First item
  \item Second item
  \item Third item with sub-items:
    \begin{itemize}
      \item Sub-item 1
      \item Sub-item 2
    \end{itemize}
\end{itemize}

\subsection{Numbered List}
\begin{enumerate}
  \item First step
  \item Second step
  \item Third step
\end{enumerate}

\section{Citations and References}
When you use the citation search feature, references will be automatically managed. For example, you might cite a paper like this: \cite{example2024}. All citations will be automatically added to your \texttt{cite.bib} file.

\section{Conclusion}
This template demonstrates basic \LaTeX{} features including mathematical formulas, tables, and lists. The bibliography system is now set up to work seamlessly with Bibby AI's citation management. Feel free to modify it for your needs!

\bibliographystyle{plain}
\bibliography{cite}

\end{document}\documentclass{article}
\usepackage[utf8]{inputenc}
\usepackage{booktabs}
\usepackage{amsmath}
\usepackage{natbib}

\title{My Document}
\author{Your Name}
\date{\today}

\begin{document}

\maketitle

\section{Introduction}
This is a sample document demonstrating various \LaTeX{} features including mathematical equations, tables, and formatting. You can cite papers using the citation search feature, and they will be automatically added to your \texttt{cite.bib} file.

\section{Mathematical Expressions}
Here are some example equations:

\subsection{Quadratic Formula}
The quadratic formula is given by:
\[x = \frac{-b \pm \sqrt{b^2 - 4ac}}{2a}\]

\subsection{Matrix Example}
A sample matrix:
\[
\begin{pmatrix}
a_{11} & a_{12} & a_{13} \\
a_{21} & a_{22} & a_{23} \\
a_{31} & a_{32} & a_{33}
\end{pmatrix}
\]

\section{Tables}
Below is a sample table showing data:

\begin{table}[h]
\centering
\begin{tabular}{@{}lrrr@{}}
\toprule
\textbf{Category} & \textbf{Value 1} & \textbf{Value 2} & \textbf{Total} \\
\midrule
Item A & 10.5 & 20.3 & 30.8 \\
Item B & 15.2 & 25.1 & 40.3 \\
Item C & 22.1 & 12.7 & 34.8 \\
\bottomrule
\end{tabular}
\caption{Sample Data Table}
\label{tab:sample}
\end{table}

\section{Lists}
\subsection{Bullet Points}
\begin{itemize}
  \item First item
  \item Second item
  \item Third item with sub-items:
    \begin{itemize}
      \item Sub-item 1
      \item Sub-item 2
    \end{itemize}
\end{itemize}

\subsection{Numbered List}
\begin{enumerate}
  \item First step
  \item Second step
  \item Third step
\end{enumerate}

\section{Citations and References}
When you use the citation search feature, references will be automatically managed. For example, you might cite a paper like this: \cite{example2024}. All citations will be automatically added to your \texttt{cite.bib} file.

\section{Conclusion}
This template demonstrates basic \LaTeX{} features including mathematical formulas, tables, and lists. The bibliography system is now set up to work seamlessly with Bibby AI's citation management. Feel free to modify it for your needs!

\bibliographystyle{plain}
\bibliography{cite}

\end{document}\documentclass{article}
\usepackage[utf8]{inputenc}
\usepackage{booktabs}
\usepackage{amsmath}
\usepackage{natbib}

\title{My Document}
\author{Your Name}
\date{\today}

\begin{document}

\maketitle

\section{Introduction}
This is a sample document demonstrating various \LaTeX{} features including mathematical equations, tables, and formatting. You can cite papers using the citation search feature, and they will be automatically added to your \texttt{cite.bib} file.

\section{Mathematical Expressions}
Here are some example equations:

\subsection{Quadratic Formula}
The quadratic formula is given by:
\[x = \frac{-b \pm \sqrt{b^2 - 4ac}}{2a}\]

\subsection{Matrix Example}
A sample matrix:
\[
\begin{pmatrix}
a_{11} & a_{12} & a_{13} \\
a_{21} & a_{22} & a_{23} \\
a_{31} & a_{32} & a_{33}
\end{pmatrix}
\]

\section{Tables}
Below is a sample table showing data:

\begin{table}[h]
\centering
\begin{tabular}{@{}lrrr@{}}
\toprule
\textbf{Category} & \textbf{Value 1} & \textbf{Value 2} & \textbf{Total} \\
\midrule
Item A & 10.5 & 20.3 & 30.8 \\
Item B & 15.2 & 25.1 & 40.3 \\
Item C & 22.1 & 12.7 & 34.8 \\
\bottomrule
\end{tabular}
\caption{Sample Data Table}
\label{tab:sample}
\end{table}

\section{Lists}
\subsection{Bullet Points}
\begin{itemize}
  \item First item
  \item Second item
  \item Third item with sub-items:
    \begin{itemize}
      \item Sub-item 1
      \item Sub-item 2
    \end{itemize}
\end{itemize}

\subsection{Numbered List}
\begin{enumerate}
  \item First step
  \item Second step
  \item Third step
\end{enumerate}

\section{Citations and References}
When you use the citation search feature, references will be automatically managed. For example, you might cite a paper like this: \cite{example2024}. All citations will be automatically added to your \texttt{cite.bib} file.

\section{Conclusion}
This template demonstrates basic \LaTeX{} features including mathematical formulas, tables, and lists. The bibliography system is now set up to work seamlessly with Bibby AI's citation management. Feel free to modify it for your needs!

\bibliographystyle{plain}
\bibliography{cite}

\end{document}\documentclass{article}
\usepackage[utf8]{inputenc}
\usepackage{booktabs}
\usepackage{amsmath}
\usepackage{natbib}

\title{My Document}
\author{Your Name}
\date{\today}

\begin{document}

\maketitle

\section{Introduction}
This is a sample document demonstrating various \LaTeX{} features including mathematical equations, tables, and formatting. You can cite papers using the citation search feature, and they will be automatically added to your \texttt{cite.bib} file.

\section{Mathematical Expressions}
Here are some example equations:

\subsection{Quadratic Formula}
The quadratic formula is given by:
\[x = \frac{-b \pm \sqrt{b^2 - 4ac}}{2a}\]

\subsection{Matrix Example}
A sample matrix:
\[
\begin{pmatrix}
a_{11} & a_{12} & a_{13} \\
a_{21} & a_{22} & a_{23} \\
a_{31} & a_{32} & a_{33}
\end{pmatrix}
\]

\section{Tables}
Below is a sample table showing data:

\begin{table}[h]
\centering
\begin{tabular}{@{}lrrr@{}}
\toprule
\textbf{Category} & \textbf{Value 1} & \textbf{Value 2} & \textbf{Total} \\
\midrule
Item A & 10.5 & 20.3 & 30.8 \\
Item B & 15.2 & 25.1 & 40.3 \\
Item C & 22.1 & 12.7 & 34.8 \\
\bottomrule
\end{tabular}
\caption{Sample Data Table}
\label{tab:sample}
\end{table}

\section{Lists}
\subsection{Bullet Points}
\begin{itemize}
  \item First item
  \item Second item
  \item Third item with sub-items:
    \begin{itemize}
      \item Sub-item 1
      \item Sub-item 2
    \end{itemize}
\end{itemize}

\subsection{Numbered List}
\begin{enumerate}
  \item First step
  \item Second step
  \item Third step
\end{enumerate}

\section{Citations and References}
When you use the citation search feature, references will be automatically managed. For example, you might cite a paper like this: \cite{example2024}. All citations will be automatically added to your \texttt{cite.bib} file.

\section{Conclusion}
This template demonstrates basic \LaTeX{} features including mathematical formulas, tables, and lists. The bibliography system is now set up to work seamlessly with Bibby AI's citation management. Feel free to modify it for your needs!

\bibliographystyle{plain}
\bibliography{cite}

\end{document}\documentclass{article}
\usepackage[utf8]{inputenc}
\usepackage{booktabs}
\usepackage{amsmath}
\usepackage{natbib}

\title{My Document}
\author{Your Name}
\date{\today}

\begin{document}

\maketitle

\section{Introduction}
This is a sample document demonstrating various \LaTeX{} features including mathematical equations, tables, and formatting. You can cite papers using the citation search feature, and they will be automatically added to your \texttt{cite.bib} file.

\section{Mathematical Expressions}
Here are some example equations:

\subsection{Quadratic Formula}
The quadratic formula is given by:
\[x = \frac{-b \pm \sqrt{b^2 - 4ac}}{2a}\]

\subsection{Matrix Example}
A sample matrix:
\[
\begin{pmatrix}
a_{11} & a_{12} & a_{13} \\
a_{21} & a_{22} & a_{23} \\
a_{31} & a_{32} & a_{33}
\end{pmatrix}
\]

\section{Tables}
Below is a sample table showing data:

\begin{table}[h]
\centering
\begin{tabular}{@{}lrrr@{}}
\toprule
\textbf{Category} & \textbf{Value 1} & \textbf{Value 2} & \textbf{Total} \\
\midrule
Item A & 10.5 & 20.3 & 30.8 \\
Item B & 15.2 & 25.1 & 40.3 \\
Item C & 22.1 & 12.7 & 34.8 \\
\bottomrule
\end{tabular}
\caption{Sample Data Table}
\label{tab:sample}
\end{table}

\section{Lists}
\subsection{Bullet Points}
\begin{itemize}
  \item First item
  \item Second item
  \item Third item with sub-items:
    \begin{itemize}
      \item Sub-item 1
      \item Sub-item 2
    \end{itemize}
\end{itemize}

\subsection{Numbered List}
\begin{enumerate}
  \item First step
  \item Second step
  \item Third step
\end{enumerate}

\section{Citations and References}
When you use the citation search feature, references will be automatically managed. For example, you might cite a paper like this: \cite{example2024}. All citations will be automatically added to your \texttt{cite.bib} file.

\section{Conclusion}
This template demonstrates basic \LaTeX{} features including mathematical formulas, tables, and lists. The bibliography system is now set up to work seamlessly with Bibby AI's citation management. Feel free to modify it for your needs!

\bibliographystyle{plain}
\bibliography{cite}

\end{document}\documentclass{article}
\usepackage[utf8]{inputenc}
\usepackage{booktabs}
\usepackage{amsmath}
\usepackage{natbib}

\title{My Document}
\author{Your Name}
\date{\today}

\begin{document}

\maketitle

\section{Introduction}
This is a sample document demonstrating various \LaTeX{} features including mathematical equations, tables, and formatting. You can cite papers using the citation search feature, and they will be automatically added to your \texttt{cite.bib} file.

\section{Mathematical Expressions}
Here are some example equations:

\subsection{Quadratic Formula}
The quadratic formula is given by:
\[x = \frac{-b \pm \sqrt{b^2 - 4ac}}{2a}\]

\subsection{Matrix Example}
A sample matrix:
\[
\begin{pmatrix}
a_{11} & a_{12} & a_{13} \\
a_{21} & a_{22} & a_{23} \\
a_{31} & a_{32} & a_{33}
\end{pmatrix}
\]

\section{Tables}
Below is a sample table showing data:

\begin{table}[h]
\centering
\begin{tabular}{@{}lrrr@{}}
\toprule
\textbf{Category} & \textbf{Value 1} & \textbf{Value 2} & \textbf{Total} \\
\midrule
Item A & 10.5 & 20.3 & 30.8 \\
Item B & 15.2 & 25.1 & 40.3 \\
Item C & 22.1 & 12.7 & 34.8 \\
\bottomrule
\end{tabular}
\caption{Sample Data Table}
\label{tab:sample}
\end{table}

\section{Lists}
\subsection{Bullet Points}
\begin{itemize}
  \item First item
  \item Second item
  \item Third item with sub-items:
    \begin{itemize}
      \item Sub-item 1
      \item Sub-item 2
    \end{itemize}
\end{itemize}

\subsection{Numbered List}
\begin{enumerate}
  \item First step
  \item Second step
  \item Third step
\end{enumerate}

\section{Citations and References}
When you use the citation search feature, references will be automatically managed. For example, you might cite a paper like this: \cite{example2024}. All citations will be automatically added to your \texttt{cite.bib} file.

\section{Conclusion}
This template demonstrates basic \LaTeX{} features including mathematical formulas, tables, and lists. The bibliography system is now set up to work seamlessly with Bibby AI's citation management. Feel free to modify it for your needs!

\bibliographystyle{plain}
\bibliography{cite}

\end{document}\documentclass{article}
\usepackage[utf8]{inputenc}
\usepackage{booktabs}
\usepackage{amsmath}
\usepackage{natbib}

\title{My Document}
\author{Your Name}
\date{\today}

\begin{document}

\maketitle

\section{Introduction}
This is a sample document demonstrating various \LaTeX{} features including mathematical equations, tables, and formatting. You can cite papers using the citation search feature, and they will be automatically added to your \texttt{cite.bib} file.

\section{Mathematical Expressions}
Here are some example equations:

\subsection{Quadratic Formula}
The quadratic formula is given by:
\[x = \frac{-b \pm \sqrt{b^2 - 4ac}}{2a}\]

\subsection{Matrix Example}
A sample matrix:
\[
\begin{pmatrix}
a_{11} & a_{12} & a_{13} \\
a_{21} & a_{22} & a_{23} \\
a_{31} & a_{32} & a_{33}
\end{pmatrix}
\]

\section{Tables}
Below is a sample table showing data:

\begin{table}[h]
\centering
\begin{tabular}{@{}lrrr@{}}
\toprule
\textbf{Category} & \textbf{Value 1} & \textbf{Value 2} & \textbf{Total} \\
\midrule
Item A & 10.5 & 20.3 & 30.8 \\
Item B & 15.2 & 25.1 & 40.3 \\
Item C & 22.1 & 12.7 & 34.8 \\
\bottomrule
\end{tabular}
\caption{Sample Data Table}
\label{tab:sample}
\end{table}

\section{Lists}
\subsection{Bullet Points}
\begin{itemize}
  \item First item
  \item Second item
  \item Third item with sub-items:
    \begin{itemize}
      \item Sub-item 1
      \item Sub-item 2
    \end{itemize}
\end{itemize}

\subsection{Numbered List}
\begin{enumerate}
  \item First step
  \item Second step
  \item Third step
\end{enumerate}

\section{Citations and References}
When you use the citation search feature, references will be automatically managed. For example, you might cite a paper like this: \cite{example2024}. All citations will be automatically added to your \texttt{cite.bib} file.

\section{Conclusion}
This template demonstrates basic \LaTeX{} features including mathematical formulas, tables, and lists. The bibliography system is now set up to work seamlessly with Bibby AI's citation management. Feel free to modify it for your needs!

\bibliographystyle{plain}
\bibliography{cite}

\end{document}\documentclass{article}
\usepackage[utf8]{inputenc}
\usepackage{booktabs}
\usepackage{amsmath}
\usepackage{natbib}

\title{My Document}
\author{Your Name}
\date{\today}

\begin{document}

\maketitle

\section{Introduction}
This is a sample document demonstrating various \LaTeX{} features including mathematical equations, tables, and formatting. You can cite papers using the citation search feature, and they will be automatically added to your \texttt{cite.bib} file.

\section{Mathematical Expressions}
Here are some example equations:

\subsection{Quadratic Formula}
The quadratic formula is given by:
\[x = \frac{-b \pm \sqrt{b^2 - 4ac}}{2a}\]

\subsection{Matrix Example}
A sample matrix:
\[
\begin{pmatrix}
a_{11} & a_{12} & a_{13} \\
a_{21} & a_{22} & a_{23} \\
a_{31} & a_{32} & a_{33}
\end{pmatrix}
\]

\section{Tables}
Below is a sample table showing data:

\begin{table}[h]
\centering
\begin{tabular}{@{}lrrr@{}}
\toprule
\textbf{Category} & \textbf{Value 1} & \textbf{Value 2} & \textbf{Total} \\
\midrule
Item A & 10.5 & 20.3 & 30.8 \\
Item B & 15.2 & 25.1 & 40.3 \\
Item C & 22.1 & 12.7 & 34.8 \\
\bottomrule
\end{tabular}
\caption{Sample Data Table}
\label{tab:sample}
\end{table}

\section{Lists}
\subsection{Bullet Points}
\begin{itemize}
  \item First item
  \item Second item
  \item Third item with sub-items:
    \begin{itemize}
      \item Sub-item 1
      \item Sub-item 2
    \end{itemize}
\end{itemize}

\subsection{Numbered List}
\begin{enumerate}
  \item First step
  \item Second step
  \item Third step
\end{enumerate}

\section{Citations and References}
When you use the citation search feature, references will be automatically managed. For example, you might cite a paper like this: \cite{example2024}. All citations will be automatically added to your \texttt{cite.bib} file.

\section{Conclusion}
This template demonstrates basic \LaTeX{} features including mathematical formulas, tables, and lists. The bibliography system is now set up to work seamlessly with Bibby AI's citation management. Feel free to modify it for your needs!

\bibliographystyle{plain}
\bibliography{cite}

\end{document}\documentclass{article}
\usepackage[utf8]{inputenc}
\usepackage{booktabs}
\usepackage{amsmath}
\usepackage{natbib}

\title{My Document}
\author{Your Name}
\date{\today}

\begin{document}

\maketitle

\section{Introduction}
This is a sample document demonstrating various \LaTeX{} features including mathematical equations, tables, and formatting. You can cite papers using the citation search feature, and they will be automatically added to your \texttt{cite.bib} file.

\section{Mathematical Expressions}
Here are some example equations:

\subsection{Quadratic Formula}
The quadratic formula is given by:
\[x = \frac{-b \pm \sqrt{b^2 - 4ac}}{2a}\]

\subsection{Matrix Example}
A sample matrix:
\[
\begin{pmatrix}
a_{11} & a_{12} & a_{13} \\
a_{21} & a_{22} & a_{23} \\
a_{31} & a_{32} & a_{33}
\end{pmatrix}
\]

\section{Tables}
Below is a sample table showing data:

\begin{table}[h]
\centering
\begin{tabular}{@{}lrrr@{}}
\toprule
\textbf{Category} & \textbf{Value 1} & \textbf{Value 2} & \textbf{Total} \\
\midrule
Item A & 10.5 & 20.3 & 30.8 \\
Item B & 15.2 & 25.1 & 40.3 \\
Item C & 22.1 & 12.7 & 34.8 \\
\bottomrule
\end{tabular}
\caption{Sample Data Table}
\label{tab:sample}
\end{table}

\section{Lists}
\subsection{Bullet Points}
\begin{itemize}
  \item First item
  \item Second item
  \item Third item with sub-items:
    \begin{itemize}
      \item Sub-item 1
      \item Sub-item 2
    \end{itemize}
\end{itemize}

\subsection{Numbered List}
\begin{enumerate}
  \item First step
  \item Second step
  \item Third step
\end{enumerate}

\section{Citations and References}
When you use the citation search feature, references will be automatically managed. For example, you might cite a paper like this: \cite{example2024}. All citations will be automatically added to your \texttt{cite.bib} file.

\section{Conclusion}
This template demonstrates basic \LaTeX{} features including mathematical formulas, tables, and lists. The bibliography system is now set up to work seamlessly with Bibby AI's citation management. Feel free to modify it for your needs!

\bibliographystyle{plain}
\bibliography{cite}

\end{document}\documentclass{article}
\usepackage[utf8]{inputenc}
\usepackage{booktabs}
\usepackage{amsmath}
\usepackage{natbib}

\title{My Document}
\author{Your Name}
\date{\today}

\begin{document}

\maketitle

\section{Introduction}
This is a sample document demonstrating various \LaTeX{} features including mathematical equations, tables, and formatting. You can cite papers using the citation search feature, and they will be automatically added to your \texttt{cite.bib} file.

\section{Mathematical Expressions}
Here are some example equations:

\subsection{Quadratic Formula}
The quadratic formula is given by:
\[x = \frac{-b \pm \sqrt{b^2 - 4ac}}{2a}\]

\subsection{Matrix Example}
A sample matrix:
\[
\begin{pmatrix}
a_{11} & a_{12} & a_{13} \\
a_{21} & a_{22} & a_{23} \\
a_{31} & a_{32} & a_{33}
\end{pmatrix}
\]

\section{Tables}
Below is a sample table showing data:

\begin{table}[h]
\centering
\begin{tabular}{@{}lrrr@{}}
\toprule
\textbf{Category} & \textbf{Value 1} & \textbf{Value 2} & \textbf{Total} \\
\midrule
Item A & 10.5 & 20.3 & 30.8 \\
Item B & 15.2 & 25.1 & 40.3 \\
Item C & 22.1 & 12.7 & 34.8 \\
\bottomrule
\end{tabular}
\caption{Sample Data Table}
\label{tab:sample}
\end{table}

\section{Lists}
\subsection{Bullet Points}
\begin{itemize}
  \item First item
  \item Second item
  \item Third item with sub-items:
    \begin{itemize}
      \item Sub-item 1
      \item Sub-item 2
    \end{itemize}
\end{itemize}

\subsection{Numbered List}
\begin{enumerate}
  \item First step
  \item Second step
  \item Third step
\end{enumerate}

\section{Citations and References}
When you use the citation search feature, references will be automatically managed. For example, you might cite a paper like this: \cite{example2024}. All citations will be automatically added to your \texttt{cite.bib} file.

\section{Conclusion}
This template demonstrates basic \LaTeX{} features including mathematical formulas, tables, and lists. The bibliography system is now set up to work seamlessly with Bibby AI's citation management. Feel free to modify it for your needs!

\bibliographystyle{plain}
\bibliography{cite}

\end{document}\documentclass{article}
\usepackage[utf8]{inputenc}
\usepackage{booktabs}
\usepackage{amsmath}
\usepackage{natbib}

\title{My Document}
\author{Your Name}
\date{\today}

\begin{document}

\maketitle

\section{Introduction}
This is a sample document demonstrating various \LaTeX{} features including mathematical equations, tables, and formatting. You can cite papers using the citation search feature, and they will be automatically added to your \texttt{cite.bib} file.

\section{Mathematical Expressions}
Here are some example equations:

\subsection{Quadratic Formula}
The quadratic formula is given by:
\[x = \frac{-b \pm \sqrt{b^2 - 4ac}}{2a}\]

\subsection{Matrix Example}
A sample matrix:
\[
\begin{pmatrix}
a_{11} & a_{12} & a_{13} \\
a_{21} & a_{22} & a_{23} \\
a_{31} & a_{32} & a_{33}
\end{pmatrix}
\]

\section{Tables}
Below is a sample table showing data:

\begin{table}[h]
\centering
\begin{tabular}{@{}lrrr@{}}
\toprule
\textbf{Category} & \textbf{Value 1} & \textbf{Value 2} & \textbf{Total} \\
\midrule
Item A & 10.5 & 20.3 & 30.8 \\
Item B & 15.2 & 25.1 & 40.3 \\
Item C & 22.1 & 12.7 & 34.8 \\
\bottomrule
\end{tabular}
\caption{Sample Data Table}
\label{tab:sample}
\end{table}

\section{Lists}
\subsection{Bullet Points}
\begin{itemize}
  \item First item
  \item Second item
  \item Third item with sub-items:
    \begin{itemize}
      \item Sub-item 1
      \item Sub-item 2
    \end{itemize}
\end{itemize}

\subsection{Numbered List}
\begin{enumerate}
  \item First step
  \item Second step
  \item Third step
\end{enumerate}

\section{Citations and References}
When you use the citation search feature, references will be automatically managed. For example, you might cite a paper like this: \cite{example2024}. All citations will be automatically added to your \texttt{cite.bib} file.

\section{Conclusion}
This template demonstrates basic \LaTeX{} features including mathematical formulas, tables, and lists. The bibliography system is now set up to work seamlessly with Bibby AI's citation management. Feel free to modify it for your needs!

\bibliographystyle{plain}
\bibliography{cite}

\end{document}\documentclass{article}
\usepackage[utf8]{inputenc}
\usepackage{booktabs}
\usepackage{amsmath}
\usepackage{natbib}

\title{My Document}
\author{Your Name}
\date{\today}

\begin{document}

\maketitle

\section{Introduction}
This is a sample document demonstrating various \LaTeX{} features including mathematical equations, tables, and formatting. You can cite papers using the citation search feature, and they will be automatically added to your \texttt{cite.bib} file.

\section{Mathematical Expressions}
Here are some example equations:

\subsection{Quadratic Formula}
The quadratic formula is given by:
\[x = \frac{-b \pm \sqrt{b^2 - 4ac}}{2a}\]

\subsection{Matrix Example}
A sample matrix:
\[
\begin{pmatrix}
a_{11} & a_{12} & a_{13} \\
a_{21} & a_{22} & a_{23} \\
a_{31} & a_{32} & a_{33}
\end{pmatrix}
\]

\section{Tables}
Below is a sample table showing data:

\begin{table}[h]
\centering
\begin{tabular}{@{}lrrr@{}}
\toprule
\textbf{Category} & \textbf{Value 1} & \textbf{Value 2} & \textbf{Total} \\
\midrule
Item A & 10.5 & 20.3 & 30.8 \\
Item B & 15.2 & 25.1 & 40.3 \\
Item C & 22.1 & 12.7 & 34.8 \\
\bottomrule
\end{tabular}
\caption{Sample Data Table}
\label{tab:sample}
\end{table}

\section{Lists}
\subsection{Bullet Points}
\begin{itemize}
  \item First item
  \item Second item
  \item Third item with sub-items:
    \begin{itemize}
      \item Sub-item 1
      \item Sub-item 2
    \end{itemize}
\end{itemize}

\subsection{Numbered List}
\begin{enumerate}
  \item First step
  \item Second step
  \item Third step
\end{enumerate}

\section{Citations and References}
When you use the citation search feature, references will be automatically managed. For example, you might cite a paper like this: \cite{example2024}. All citations will be automatically added to your \texttt{cite.bib} file.

\section{Conclusion}
This template demonstrates basic \LaTeX{} features including mathematical formulas, tables, and lists. The bibliography system is now set up to work seamlessly with Bibby AI's citation management. Feel free to modify it for your needs!

\bibliographystyle{plain}
\bibliography{cite}

\end{document}\documentclass{article}
\usepackage[utf8]{inputenc}
\usepackage{booktabs}
\usepackage{amsmath}
\usepackage{natbib}

\title{My Document}
\author{Your Name}
\date{\today}

\begin{document}

\maketitle

\section{Introduction}
This is a sample document demonstrating various \LaTeX{} features including mathematical equations, tables, and formatting. You can cite papers using the citation search feature, and they will be automatically added to your \texttt{cite.bib} file.

\section{Mathematical Expressions}
Here are some example equations:

\subsection{Quadratic Formula}
The quadratic formula is given by:
\[x = \frac{-b \pm \sqrt{b^2 - 4ac}}{2a}\]

\subsection{Matrix Example}
A sample matrix:
\[
\begin{pmatrix}
a_{11} & a_{12} & a_{13} \\
a_{21} & a_{22} & a_{23} \\
a_{31} & a_{32} & a_{33}
\end{pmatrix}
\]

\section{Tables}
Below is a sample table showing data:

\begin{table}[h]
\centering
\begin{tabular}{@{}lrrr@{}}
\toprule
\textbf{Category} & \textbf{Value 1} & \textbf{Value 2} & \textbf{Total} \\
\midrule
Item A & 10.5 & 20.3 & 30.8 \\
Item B & 15.2 & 25.1 & 40.3 \\
Item C & 22.1 & 12.7 & 34.8 \\
\bottomrule
\end{tabular}
\caption{Sample Data Table}
\label{tab:sample}
\end{table}

\section{Lists}
\subsection{Bullet Points}
\begin{itemize}
  \item First item
  \item Second item
  \item Third item with sub-items:
    \begin{itemize}
      \item Sub-item 1
      \item Sub-item 2
    \end{itemize}
\end{itemize}

\subsection{Numbered List}
\begin{enumerate}
  \item First step
  \item Second step
  \item Third step
\end{enumerate}

\section{Citations and References}
When you use the citation search feature, references will be automatically managed. For example, you might cite a paper like this: \cite{example2024}. All citations will be automatically added to your \texttt{cite.bib} file.

\section{Conclusion}
This template demonstrates basic \LaTeX{} features including mathematical formulas, tables, and lists. The bibliography system is now set up to work seamlessly with Bibby AI's citation management. Feel free to modify it for your needs!

\bibliographystyle{plain}
\bibliography{cite}

\end{document}\documentclass{article}
\usepackage[utf8]{inputenc}
\usepackage{booktabs}
\usepackage{amsmath}
\usepackage{natbib}

\title{My Document}
\author{Your Name}
\date{\today}

\begin{document}

\maketitle

\section{Introduction}
This is a sample document demonstrating various \LaTeX{} features including mathematical equations, tables, and formatting. You can cite papers using the citation search feature, and they will be automatically added to your \texttt{cite.bib} file.

\section{Mathematical Expressions}
Here are some example equations:

\subsection{Quadratic Formula}
The quadratic formula is given by:
\[x = \frac{-b \pm \sqrt{b^2 - 4ac}}{2a}\]

\subsection{Matrix Example}
A sample matrix:
\[
\begin{pmatrix}
a_{11} & a_{12} & a_{13} \\
a_{21} & a_{22} & a_{23} \\
a_{31} & a_{32} & a_{33}
\end{pmatrix}
\]

\section{Tables}
Below is a sample table showing data:

\begin{table}[h]
\centering
\begin{tabular}{@{}lrrr@{}}
\toprule
\textbf{Category} & \textbf{Value 1} & \textbf{Value 2} & \textbf{Total} \\
\midrule
Item A & 10.5 & 20.3 & 30.8 \\
Item B & 15.2 & 25.1 & 40.3 \\
Item C & 22.1 & 12.7 & 34.8 \\
\bottomrule
\end{tabular}
\caption{Sample Data Table}
\label{tab:sample}
\end{table}

\section{Lists}
\subsection{Bullet Points}
\begin{itemize}
  \item First item
  \item Second item
  \item Third item with sub-items:
    \begin{itemize}
      \item Sub-item 1
      \item Sub-item 2
    \end{itemize}
\end{itemize}

\subsection{Numbered List}
\begin{enumerate}
  \item First step
  \item Second step
  \item Third step
\end{enumerate}

\section{Citations and References}
When you use the citation search feature, references will be automatically managed. For example, you might cite a paper like this: \cite{example2024}. All citations will be automatically added to your \texttt{cite.bib} file.

\section{Conclusion}
This template demonstrates basic \LaTeX{} features including mathematical formulas, tables, and lists. The bibliography system is now set up to work seamlessly with Bibby AI's citation management. Feel free to modify it for your needs!

\bibliographystyle{plain}
\bibliography{cite}

\end{document}\documentclass{article}
\usepackage[utf8]{inputenc}
\usepackage{booktabs}
\usepackage{amsmath}
\usepackage{natbib}

\title{My Document}
\author{Your Name}
\date{\today}

\begin{document}

\maketitle

\section{Introduction}
This is a sample document demonstrating various \LaTeX{} features including mathematical equations, tables, and formatting. You can cite papers using the citation search feature, and they will be automatically added to your \texttt{cite.bib} file.

\section{Mathematical Expressions}
Here are some example equations:

\subsection{Quadratic Formula}
The quadratic formula is given by:
\[x = \frac{-b \pm \sqrt{b^2 - 4ac}}{2a}\]

\subsection{Matrix Example}
A sample matrix:
\[
\begin{pmatrix}
a_{11} & a_{12} & a_{13} \\
a_{21} & a_{22} & a_{23} \\
a_{31} & a_{32} & a_{33}
\end{pmatrix}
\]

\section{Tables}
Below is a sample table showing data:

\begin{table}[h]
\centering
\begin{tabular}{@{}lrrr@{}}
\toprule
\textbf{Category} & \textbf{Value 1} & \textbf{Value 2} & \textbf{Total} \\
\midrule
Item A & 10.5 & 20.3 & 30.8 \\
Item B & 15.2 & 25.1 & 40.3 \\
Item C & 22.1 & 12.7 & 34.8 \\
\bottomrule
\end{tabular}
\caption{Sample Data Table}
\label{tab:sample}
\end{table}

\section{Lists}
\subsection{Bullet Points}
\begin{itemize}
  \item First item
  \item Second item
  \item Third item with sub-items:
    \begin{itemize}
      \item Sub-item 1
      \item Sub-item 2
    \end{itemize}
\end{itemize}

\subsection{Numbered List}
\begin{enumerate}
  \item First step
  \item Second step
  \item Third step
\end{enumerate}

\section{Citations and References}
When you use the citation search feature, references will be automatically managed. For example, you might cite a paper like this: \cite{example2024}. All citations will be automatically added to your \texttt{cite.bib} file.

\section{Conclusion}
This template demonstrates basic \LaTeX{} features including mathematical formulas, tables, and lists. The bibliography system is now set up to work seamlessly with Bibby AI's citation management. Feel free to modify it for your needs!

\bibliographystyle{plain}
\bibliography{cite}

\end{document}\documentclass{article}
\usepackage[utf8]{inputenc}
\usepackage{booktabs}
\usepackage{amsmath}
\usepackage{natbib}

\title{My Document}
\author{Your Name}
\date{\today}

\begin{document}

\maketitle

\section{Introduction}
This is a sample document demonstrating various \LaTeX{} features including mathematical equations, tables, and formatting. You can cite papers using the citation search feature, and they will be automatically added to your \texttt{cite.bib} file.

\section{Mathematical Expressions}
Here are some example equations:

\subsection{Quadratic Formula}
The quadratic formula is given by:
\[x = \frac{-b \pm \sqrt{b^2 - 4ac}}{2a}\]

\subsection{Matrix Example}
A sample matrix:
\[
\begin{pmatrix}
a_{11} & a_{12} & a_{13} \\
a_{21} & a_{22} & a_{23} \\
a_{31} & a_{32} & a_{33}
\end{pmatrix}
\]

\section{Tables}
Below is a sample table showing data:

\begin{table}[h]
\centering
\begin{tabular}{@{}lrrr@{}}
\toprule
\textbf{Category} & \textbf{Value 1} & \textbf{Value 2} & \textbf{Total} \\
\midrule
Item A & 10.5 & 20.3 & 30.8 \\
Item B & 15.2 & 25.1 & 40.3 \\
Item C & 22.1 & 12.7 & 34.8 \\
\bottomrule
\end{tabular}
\caption{Sample Data Table}
\label{tab:sample}
\end{table}

\section{Lists}
\subsection{Bullet Points}
\begin{itemize}
  \item First item
  \item Second item
  \item Third item with sub-items:
    \begin{itemize}
      \item Sub-item 1
      \item Sub-item 2
    \end{itemize}
\end{itemize}

\subsection{Numbered List}
\begin{enumerate}
  \item First step
  \item Second step
  \item Third step
\end{enumerate}

\section{Citations and References}
When you use the citation search feature, references will be automatically managed. For example, you might cite a paper like this: \cite{example2024}. All citations will be automatically added to your \texttt{cite.bib} file.

\section{Conclusion}
This template demonstrates basic \LaTeX{} features including mathematical formulas, tables, and lists. The bibliography system is now set up to work seamlessly with Bibby AI's citation management. Feel free to modify it for your needs!

\bibliographystyle{plain}
\bibliography{cite}

\end{document}\documentclass{article}
\usepackage[utf8]{inputenc}
\usepackage{booktabs}
\usepackage{amsmath}
\usepackage{natbib}

\title{My Document}
\author{Your Name}
\date{\today}

\begin{document}

\maketitle

\section{Introduction}
This is a sample document demonstrating various \LaTeX{} features including mathematical equations, tables, and formatting. You can cite papers using the citation search feature, and they will be automatically added to your \texttt{cite.bib} file.

\section{Mathematical Expressions}
Here are some example equations:

\subsection{Quadratic Formula}
The quadratic formula is given by:
\[x = \frac{-b \pm \sqrt{b^2 - 4ac}}{2a}\]

\subsection{Matrix Example}
A sample matrix:
\[
\begin{pmatrix}
a_{11} & a_{12} & a_{13} \\
a_{21} & a_{22} & a_{23} \\
a_{31} & a_{32} & a_{33}
\end{pmatrix}
\]

\section{Tables}
Below is a sample table showing data:

\begin{table}[h]
\centering
\begin{tabular}{@{}lrrr@{}}
\toprule
\textbf{Category} & \textbf{Value 1} & \textbf{Value 2} & \textbf{Total} \\
\midrule
Item A & 10.5 & 20.3 & 30.8 \\
Item B & 15.2 & 25.1 & 40.3 \\
Item C & 22.1 & 12.7 & 34.8 \\
\bottomrule
\end{tabular}
\caption{Sample Data Table}
\label{tab:sample}
\end{table}

\section{Lists}
\subsection{Bullet Points}
\begin{itemize}
  \item First item
  \item Second item
  \item Third item with sub-items:
    \begin{itemize}
      \item Sub-item 1
      \item Sub-item 2
    \end{itemize}
\end{itemize}

\subsection{Numbered List}
\begin{enumerate}
  \item First step
  \item Second step
  \item Third step
\end{enumerate}

\section{Citations and References}
When you use the citation search feature, references will be automatically managed. For example, you might cite a paper like this: \cite{example2024}. All citations will be automatically added to your \texttt{cite.bib} file.

\section{Conclusion}
This template demonstrates basic \LaTeX{} features including mathematical formulas, tables, and lists. The bibliography system is now set up to work seamlessly with Bibby AI's citation management. Feel free to modify it for your needs!

\bibliographystyle{plain}
\bibliography{cite}

\end{document}\documentclass{article}
\usepackage[utf8]{inputenc}
\usepackage{booktabs}
\usepackage{amsmath}
\usepackage{natbib}

\title{My Document}
\author{Your Name}
\date{\today}

\begin{document}

\maketitle

\section{Introduction}
This is a sample document demonstrating various \LaTeX{} features including mathematical equations, tables, and formatting. You can cite papers using the citation search feature, and they will be automatically added to your \texttt{cite.bib} file.

\section{Mathematical Expressions}
Here are some example equations:

\subsection{Quadratic Formula}
The quadratic formula is given by:
\[x = \frac{-b \pm \sqrt{b^2 - 4ac}}{2a}\]

\subsection{Matrix Example}
A sample matrix:
\[
\begin{pmatrix}
a_{11} & a_{12} & a_{13} \\
a_{21} & a_{22} & a_{23} \\
a_{31} & a_{32} & a_{33}
\end{pmatrix}
\]

\section{Tables}
Below is a sample table showing data:

\begin{table}[h]
\centering
\begin{tabular}{@{}lrrr@{}}
\toprule
\textbf{Category} & \textbf{Value 1} & \textbf{Value 2} & \textbf{Total} \\
\midrule
Item A & 10.5 & 20.3 & 30.8 \\
Item B & 15.2 & 25.1 & 40.3 \\
Item C & 22.1 & 12.7 & 34.8 \\
\bottomrule
\end{tabular}
\caption{Sample Data Table}
\label{tab:sample}
\end{table}

\section{Lists}
\subsection{Bullet Points}
\begin{itemize}
  \item First item
  \item Second item
  \item Third item with sub-items:
    \begin{itemize}
      \item Sub-item 1
      \item Sub-item 2
    \end{itemize}
\end{itemize}

\subsection{Numbered List}
\begin{enumerate}
  \item First step
  \item Second step
  \item Third step
\end{enumerate}

\section{Citations and References}
When you use the citation search feature, references will be automatically managed. For example, you might cite a paper like this: \cite{example2024}. All citations will be automatically added to your \texttt{cite.bib} file.

\section{Conclusion}
This template demonstrates basic \LaTeX{} features including mathematical formulas, tables, and lists. The bibliography system is now set up to work seamlessly with Bibby AI's citation management. Feel free to modify it for your needs!

\bibliographystyle{plain}
\bibliography{cite}

\end{document}\documentclass{article}
\usepackage[utf8]{inputenc}
\usepackage{booktabs}
\usepackage{amsmath}
\usepackage{natbib}

\title{My Document}
\author{Your Name}
\date{\today}

\begin{document}

\maketitle

\section{Introduction}
This is a sample document demonstrating various \LaTeX{} features including mathematical equations, tables, and formatting. You can cite papers using the citation search feature, and they will be automatically added to your \texttt{cite.bib} file.

\section{Mathematical Expressions}
Here are some example equations:

\subsection{Quadratic Formula}
The quadratic formula is given by:
\[x = \frac{-b \pm \sqrt{b^2 - 4ac}}{2a}\]

\subsection{Matrix Example}
A sample matrix:
\[
\begin{pmatrix}
a_{11} & a_{12} & a_{13} \\
a_{21} & a_{22} & a_{23} \\
a_{31} & a_{32} & a_{33}
\end{pmatrix}
\]

\section{Tables}
Below is a sample table showing data:

\begin{table}[h]
\centering
\begin{tabular}{@{}lrrr@{}}
\toprule
\textbf{Category} & \textbf{Value 1} & \textbf{Value 2} & \textbf{Total} \\
\midrule
Item A & 10.5 & 20.3 & 30.8 \\
Item B & 15.2 & 25.1 & 40.3 \\
Item C & 22.1 & 12.7 & 34.8 \\
\bottomrule
\end{tabular}
\caption{Sample Data Table}
\label{tab:sample}
\end{table}

\section{Lists}
\subsection{Bullet Points}
\begin{itemize}
  \item First item
  \item Second item
  \item Third item with sub-items:
    \begin{itemize}
      \item Sub-item 1
      \item Sub-item 2
    \end{itemize}
\end{itemize}

\subsection{Numbered List}
\begin{enumerate}
  \item First step
  \item Second step
  \item Third step
\end{enumerate}

\section{Citations and References}
When you use the citation search feature, references will be automatically managed. For example, you might cite a paper like this: \cite{example2024}. All citations will be automatically added to your \texttt{cite.bib} file.

\section{Conclusion}
This template demonstrates basic \LaTeX{} features including mathematical formulas, tables, and lists. The bibliography system is now set up to work seamlessly with Bibby AI's citation management. Feel free to modify it for your needs!

\bibliographystyle{plain}
\bibliography{cite}

\end{document}\documentclass{article}
\usepackage[utf8]{inputenc}
\usepackage{booktabs}
\usepackage{amsmath}
\usepackage{natbib}

\title{My Document}
\author{Your Name}
\date{\today}

\begin{document}

\maketitle

\section{Introduction}
This is a sample document demonstrating various \LaTeX{} features including mathematical equations, tables, and formatting. You can cite papers using the citation search feature, and they will be automatically added to your \texttt{cite.bib} file.

\section{Mathematical Expressions}
Here are some example equations:

\subsection{Quadratic Formula}
The quadratic formula is given by:
\[x = \frac{-b \pm \sqrt{b^2 - 4ac}}{2a}\]

\subsection{Matrix Example}
A sample matrix:
\[
\begin{pmatrix}
a_{11} & a_{12} & a_{13} \\
a_{21} & a_{22} & a_{23} \\
a_{31} & a_{32} & a_{33}
\end{pmatrix}
\]

\section{Tables}
Below is a sample table showing data:

\begin{table}[h]
\centering
\begin{tabular}{@{}lrrr@{}}
\toprule
\textbf{Category} & \textbf{Value 1} & \textbf{Value 2} & \textbf{Total} \\
\midrule
Item A & 10.5 & 20.3 & 30.8 \\
Item B & 15.2 & 25.1 & 40.3 \\
Item C & 22.1 & 12.7 & 34.8 \\
\bottomrule
\end{tabular}
\caption{Sample Data Table}
\label{tab:sample}
\end{table}

\section{Lists}
\subsection{Bullet Points}
\begin{itemize}
  \item First item
  \item Second item
  \item Third item with sub-items:
    \begin{itemize}
      \item Sub-item 1
      \item Sub-item 2
    \end{itemize}
\end{itemize}

\subsection{Numbered List}
\begin{enumerate}
  \item First step
  \item Second step
  \item Third step
\end{enumerate}

\section{Citations and References}
When you use the citation search feature, references will be automatically managed. For example, you might cite a paper like this: \cite{example2024}. All citations will be automatically added to your \texttt{cite.bib} file.

\section{Conclusion}
This template demonstrates basic \LaTeX{} features including mathematical formulas, tables, and lists. The bibliography system is now set up to work seamlessly with Bibby AI's citation management. Feel free to modify it for your needs!

\bibliographystyle{plain}
\bibliography{cite}

\end{document}\documentclass{article}
\usepackage[utf8]{inputenc}
\usepackage{booktabs}
\usepackage{amsmath}
\usepackage{natbib}

\title{My Document}
\author{Your Name}
\date{\today}

\begin{document}

\maketitle

\section{Introduction}
This is a sample document demonstrating various \LaTeX{} features including mathematical equations, tables, and formatting. You can cite papers using the citation search feature, and they will be automatically added to your \texttt{cite.bib} file.

\section{Mathematical Expressions}
Here are some example equations:

\subsection{Quadratic Formula}
The quadratic formula is given by:
\[x = \frac{-b \pm \sqrt{b^2 - 4ac}}{2a}\]

\subsection{Matrix Example}
A sample matrix:
\[
\begin{pmatrix}
a_{11} & a_{12} & a_{13} \\
a_{21} & a_{22} & a_{23} \\
a_{31} & a_{32} & a_{33}
\end{pmatrix}
\]

\section{Tables}
Below is a sample table showing data:

\begin{table}[h]
\centering
\begin{tabular}{@{}lrrr@{}}
\toprule
\textbf{Category} & \textbf{Value 1} & \textbf{Value 2} & \textbf{Total} \\
\midrule
Item A & 10.5 & 20.3 & 30.8 \\
Item B & 15.2 & 25.1 & 40.3 \\
Item C & 22.1 & 12.7 & 34.8 \\
\bottomrule
\end{tabular}
\caption{Sample Data Table}
\label{tab:sample}
\end{table}

\section{Lists}
\subsection{Bullet Points}
\begin{itemize}
  \item First item
  \item Second item
  \item Third item with sub-items:
    \begin{itemize}
      \item Sub-item 1
      \item Sub-item 2
    \end{itemize}
\end{itemize}

\subsection{Numbered List}
\begin{enumerate}
  \item First step
  \item Second step
  \item Third step
\end{enumerate}

\section{Citations and References}
When you use the citation search feature, references will be automatically managed. For example, you might cite a paper like this: \cite{example2024}. All citations will be automatically added to your \texttt{cite.bib} file.

\section{Conclusion}
This template demonstrates basic \LaTeX{} features including mathematical formulas, tables, and lists. The bibliography system is now set up to work seamlessly with Bibby AI's citation management. Feel free to modify it for your needs!

\bibliographystyle{plain}
\bibliography{cite}

\end{document}\documentclass{article}
\usepackage[utf8]{inputenc}
\usepackage{booktabs}
\usepackage{amsmath}
\usepackage{natbib}

\title{My Document}
\author{Your Name}
\date{\today}

\begin{document}

\maketitle

\section{Introduction}
This is a sample document demonstrating various \LaTeX{} features including mathematical equations, tables, and formatting. You can cite papers using the citation search feature, and they will be automatically added to your \texttt{cite.bib} file.

\section{Mathematical Expressions}
Here are some example equations:

\subsection{Quadratic Formula}
The quadratic formula is given by:
\[x = \frac{-b \pm \sqrt{b^2 - 4ac}}{2a}\]

\subsection{Matrix Example}
A sample matrix:
\[
\begin{pmatrix}
a_{11} & a_{12} & a_{13} \\
a_{21} & a_{22} & a_{23} \\
a_{31} & a_{32} & a_{33}
\end{pmatrix}
\]

\section{Tables}
Below is a sample table showing data:

\begin{table}[h]
\centering
\begin{tabular}{@{}lrrr@{}}
\toprule
\textbf{Category} & \textbf{Value 1} & \textbf{Value 2} & \textbf{Total} \\
\midrule
Item A & 10.5 & 20.3 & 30.8 \\
Item B & 15.2 & 25.1 & 40.3 \\
Item C & 22.1 & 12.7 & 34.8 \\
\bottomrule
\end{tabular}
\caption{Sample Data Table}
\label{tab:sample}
\end{table}

\section{Lists}
\subsection{Bullet Points}
\begin{itemize}
  \item First item
  \item Second item
  \item Third item with sub-items:
    \begin{itemize}
      \item Sub-item 1
      \item Sub-item 2
    \end{itemize}
\end{itemize}

\subsection{Numbered List}
\begin{enumerate}
  \item First step
  \item Second step
  \item Third step
\end{enumerate}

\section{Citations and References}
When you use the citation search feature, references will be automatically managed. For example, you might cite a paper like this: \cite{example2024}. All citations will be automatically added to your \texttt{cite.bib} file.

\section{Conclusion}
This template demonstrates basic \LaTeX{} features including mathematical formulas, tables, and lists. The bibliography system is now set up to work seamlessly with Bibby AI's citation management. Feel free to modify it for your needs!

\bibliographystyle{plain}
\bibliography{cite}

\end{document}\documentclass{article}
\usepackage[utf8]{inputenc}
\usepackage{booktabs}
\usepackage{amsmath}
\usepackage{natbib}

\title{My Document}
\author{Your Name}
\date{\today}

\begin{document}

\maketitle

\section{Introduction}
This is a sample document demonstrating various \LaTeX{} features including mathematical equations, tables, and formatting. You can cite papers using the citation search feature, and they will be automatically added to your \texttt{cite.bib} file.

\section{Mathematical Expressions}
Here are some example equations:

\subsection{Quadratic Formula}
The quadratic formula is given by:
\[x = \frac{-b \pm \sqrt{b^2 - 4ac}}{2a}\]

\subsection{Matrix Example}
A sample matrix:
\[
\begin{pmatrix}
a_{11} & a_{12} & a_{13} \\
a_{21} & a_{22} & a_{23} \\
a_{31} & a_{32} & a_{33}
\end{pmatrix}
\]

\section{Tables}
Below is a sample table showing data:

\begin{table}[h]
\centering
\begin{tabular}{@{}lrrr@{}}
\toprule
\textbf{Category} & \textbf{Value 1} & \textbf{Value 2} & \textbf{Total} \\
\midrule
Item A & 10.5 & 20.3 & 30.8 \\
Item B & 15.2 & 25.1 & 40.3 \\
Item C & 22.1 & 12.7 & 34.8 \\
\bottomrule
\end{tabular}
\caption{Sample Data Table}
\label{tab:sample}
\end{table}

\section{Lists}
\subsection{Bullet Points}
\begin{itemize}
  \item First item
  \item Second item
  \item Third item with sub-items:
    \begin{itemize}
      \item Sub-item 1
      \item Sub-item 2
    \end{itemize}
\end{itemize}

\subsection{Numbered List}
\begin{enumerate}
  \item First step
  \item Second step
  \item Third step
\end{enumerate}

\section{Citations and References}
When you use the citation search feature, references will be automatically managed. For example, you might cite a paper like this: \cite{example2024}. All citations will be automatically added to your \texttt{cite.bib} file.

\section{Conclusion}
This template demonstrates basic \LaTeX{} features including mathematical formulas, tables, and lists. The bibliography system is now set up to work seamlessly with Bibby AI's citation management. Feel free to modify it for your needs!

\bibliographystyle{plain}
\bibliography{cite}

\end{document}\documentclass{article}
\usepackage[utf8]{inputenc}
\usepackage{booktabs}
\usepackage{amsmath}
\usepackage{natbib}

\title{My Document}
\author{Your Name}
\date{\today}

\begin{document}

\maketitle

\section{Introduction}
This is a sample document demonstrating various \LaTeX{} features including mathematical equations, tables, and formatting. You can cite papers using the citation search feature, and they will be automatically added to your \texttt{cite.bib} file.

\section{Mathematical Expressions}
Here are some example equations:

\subsection{Quadratic Formula}
The quadratic formula is given by:
\[x = \frac{-b \pm \sqrt{b^2 - 4ac}}{2a}\]

\subsection{Matrix Example}
A sample matrix:
\[
\begin{pmatrix}
a_{11} & a_{12} & a_{13} \\
a_{21} & a_{22} & a_{23} \\
a_{31} & a_{32} & a_{33}
\end{pmatrix}
\]

\section{Tables}
Below is a sample table showing data:

\begin{table}[h]
\centering
\begin{tabular}{@{}lrrr@{}}
\toprule
\textbf{Category} & \textbf{Value 1} & \textbf{Value 2} & \textbf{Total} \\
\midrule
Item A & 10.5 & 20.3 & 30.8 \\
Item B & 15.2 & 25.1 & 40.3 \\
Item C & 22.1 & 12.7 & 34.8 \\
\bottomrule
\end{tabular}
\caption{Sample Data Table}
\label{tab:sample}
\end{table}

\section{Lists}
\subsection{Bullet Points}
\begin{itemize}
  \item First item
  \item Second item
  \item Third item with sub-items:
    \begin{itemize}
      \item Sub-item 1
      \item Sub-item 2
    \end{itemize}
\end{itemize}

\subsection{Numbered List}
\begin{enumerate}
  \item First step
  \item Second step
  \item Third step
\end{enumerate}

\section{Citations and References}
When you use the citation search feature, references will be automatically managed. For example, you might cite a paper like this: \cite{example2024}. All citations will be automatically added to your \texttt{cite.bib} file.

\section{Conclusion}
This template demonstrates basic \LaTeX{} features including mathematical formulas, tables, and lists. The bibliography system is now set up to work seamlessly with Bibby AI's citation management. Feel free to modify it for your needs!

\bibliographystyle{plain}
\bibliography{cite}

\end{document}\documentclass{article}
\usepackage[utf8]{inputenc}
\usepackage{booktabs}
\usepackage{amsmath}
\usepackage{natbib}

\title{My Document}
\author{Your Name}
\date{\today}

\begin{document}

\maketitle

\section{Introduction}
This is a sample document demonstrating various \LaTeX{} features including mathematical equations, tables, and formatting. You can cite papers using the citation search feature, and they will be automatically added to your \texttt{cite.bib} file.

\section{Mathematical Expressions}
Here are some example equations:

\subsection{Quadratic Formula}
The quadratic formula is given by:
\[x = \frac{-b \pm \sqrt{b^2 - 4ac}}{2a}\]

\subsection{Matrix Example}
A sample matrix:
\[
\begin{pmatrix}
a_{11} & a_{12} & a_{13} \\
a_{21} & a_{22} & a_{23} \\
a_{31} & a_{32} & a_{33}
\end{pmatrix}
\]

\section{Tables}
Below is a sample table showing data:

\begin{table}[h]
\centering
\begin{tabular}{@{}lrrr@{}}
\toprule
\textbf{Category} & \textbf{Value 1} & \textbf{Value 2} & \textbf{Total} \\
\midrule
Item A & 10.5 & 20.3 & 30.8 \\
Item B & 15.2 & 25.1 & 40.3 \\
Item C & 22.1 & 12.7 & 34.8 \\
\bottomrule
\end{tabular}
\caption{Sample Data Table}
\label{tab:sample}
\end{table}

\section{Lists}
\subsection{Bullet Points}
\begin{itemize}
  \item First item
  \item Second item
  \item Third item with sub-items:
    \begin{itemize}
      \item Sub-item 1
      \item Sub-item 2
    \end{itemize}
\end{itemize}

\subsection{Numbered List}
\begin{enumerate}
  \item First step
  \item Second step
  \item Third step
\end{enumerate}

\section{Citations and References}
When you use the citation search feature, references will be automatically managed. For example, you might cite a paper like this: \cite{example2024}. All citations will be automatically added to your \texttt{cite.bib} file.

\section{Conclusion}
This template demonstrates basic \LaTeX{} features including mathematical formulas, tables, and lists. The bibliography system is now set up to work seamlessly with Bibby AI's citation management. Feel free to modify it for your needs!

\bibliographystyle{plain}
\bibliography{cite}

\end{document}\documentclass{article}
\usepackage[utf8]{inputenc}
\usepackage{booktabs}
\usepackage{amsmath}
\usepackage{natbib}

\title{My Document}
\author{Your Name}
\date{\today}

\begin{document}

\maketitle

\section{Introduction}
This is a sample document demonstrating various \LaTeX{} features including mathematical equations, tables, and formatting. You can cite papers using the citation search feature, and they will be automatically added to your \texttt{cite.bib} file.

\section{Mathematical Expressions}
Here are some example equations:

\subsection{Quadratic Formula}
The quadratic formula is given by:
\[x = \frac{-b \pm \sqrt{b^2 - 4ac}}{2a}\]

\subsection{Matrix Example}
A sample matrix:
\[
\begin{pmatrix}
a_{11} & a_{12} & a_{13} \\
a_{21} & a_{22} & a_{23} \\
a_{31} & a_{32} & a_{33}
\end{pmatrix}
\]

\section{Tables}
Below is a sample table showing data:

\begin{table}[h]
\centering
\begin{tabular}{@{}lrrr@{}}
\toprule
\textbf{Category} & \textbf{Value 1} & \textbf{Value 2} & \textbf{Total} \\
\midrule
Item A & 10.5 & 20.3 & 30.8 \\
Item B & 15.2 & 25.1 & 40.3 \\
Item C & 22.1 & 12.7 & 34.8 \\
\bottomrule
\end{tabular}
\caption{Sample Data Table}
\label{tab:sample}
\end{table}

\section{Lists}
\subsection{Bullet Points}
\begin{itemize}
  \item First item
  \item Second item
  \item Third item with sub-items:
    \begin{itemize}
      \item Sub-item 1
      \item Sub-item 2
    \end{itemize}
\end{itemize}

\subsection{Numbered List}
\begin{enumerate}
  \item First step
  \item Second step
  \item Third step
\end{enumerate}

\section{Citations and References}
When you use the citation search feature, references will be automatically managed. For example, you might cite a paper like this: \cite{example2024}. All citations will be automatically added to your \texttt{cite.bib} file.

\section{Conclusion}
This template demonstrates basic \LaTeX{} features including mathematical formulas, tables, and lists. The bibliography system is now set up to work seamlessly with Bibby AI's citation management. Feel free to modify it for your needs!

\bibliographystyle{plain}
\bibliography{cite}

\end{document}\documentclass{article}
\usepackage[utf8]{inputenc}
\usepackage{booktabs}
\usepackage{amsmath}
\usepackage{natbib}

\title{My Document}
\author{Your Name}
\date{\today}

\begin{document}

\maketitle

\section{Introduction}
This is a sample document demonstrating various \LaTeX{} features including mathematical equations, tables, and formatting. You can cite papers using the citation search feature, and they will be automatically added to your \texttt{cite.bib} file.

\section{Mathematical Expressions}
Here are some example equations:

\subsection{Quadratic Formula}
The quadratic formula is given by:
\[x = \frac{-b \pm \sqrt{b^2 - 4ac}}{2a}\]

\subsection{Matrix Example}
A sample matrix:
\[
\begin{pmatrix}
a_{11} & a_{12} & a_{13} \\
a_{21} & a_{22} & a_{23} \\
a_{31} & a_{32} & a_{33}
\end{pmatrix}
\]

\section{Tables}
Below is a sample table showing data:

\begin{table}[h]
\centering
\begin{tabular}{@{}lrrr@{}}
\toprule
\textbf{Category} & \textbf{Value 1} & \textbf{Value 2} & \textbf{Total} \\
\midrule
Item A & 10.5 & 20.3 & 30.8 \\
Item B & 15.2 & 25.1 & 40.3 \\
Item C & 22.1 & 12.7 & 34.8 \\
\bottomrule
\end{tabular}
\caption{Sample Data Table}
\label{tab:sample}
\end{table}

\section{Lists}
\subsection{Bullet Points}
\begin{itemize}
  \item First item
  \item Second item
  \item Third item with sub-items:
    \begin{itemize}
      \item Sub-item 1
      \item Sub-item 2
    \end{itemize}
\end{itemize}

\subsection{Numbered List}
\begin{enumerate}
  \item First step
  \item Second step
  \item Third step
\end{enumerate}

\section{Citations and References}
When you use the citation search feature, references will be automatically managed. For example, you might cite a paper like this: \cite{example2024}. All citations will be automatically added to your \texttt{cite.bib} file.

\section{Conclusion}
This template demonstrates basic \LaTeX{} features including mathematical formulas, tables, and lists. The bibliography system is now set up to work seamlessly with Bibby AI's citation management. Feel free to modify it for your needs!

\bibliographystyle{plain}
\bibliography{cite}

\end{document}\documentclass{article}
\usepackage[utf8]{inputenc}
\usepackage{booktabs}
\usepackage{amsmath}
\usepackage{natbib}

\title{My Document}
\author{Your Name}
\date{\today}

\begin{document}

\maketitle

\section{Introduction}
This is a sample document demonstrating various \LaTeX{} features including mathematical equations, tables, and formatting. You can cite papers using the citation search feature, and they will be automatically added to your \texttt{cite.bib} file.

\section{Mathematical Expressions}
Here are some example equations:

\subsection{Quadratic Formula}
The quadratic formula is given by:
\[x = \frac{-b \pm \sqrt{b^2 - 4ac}}{2a}\]

\subsection{Matrix Example}
A sample matrix:
\[
\begin{pmatrix}
a_{11} & a_{12} & a_{13} \\
a_{21} & a_{22} & a_{23} \\
a_{31} & a_{32} & a_{33}
\end{pmatrix}
\]

\section{Tables}
Below is a sample table showing data:

\begin{table}[h]
\centering
\begin{tabular}{@{}lrrr@{}}
\toprule
\textbf{Category} & \textbf{Value 1} & \textbf{Value 2} & \textbf{Total} \\
\midrule
Item A & 10.5 & 20.3 & 30.8 \\
Item B & 15.2 & 25.1 & 40.3 \\
Item C & 22.1 & 12.7 & 34.8 \\
\bottomrule
\end{tabular}
\caption{Sample Data Table}
\label{tab:sample}
\end{table}

\section{Lists}
\subsection{Bullet Points}
\begin{itemize}
  \item First item
  \item Second item
  \item Third item with sub-items:
    \begin{itemize}
      \item Sub-item 1
      \item Sub-item 2
    \end{itemize}
\end{itemize}

\subsection{Numbered List}
\begin{enumerate}
  \item First step
  \item Second step
  \item Third step
\end{enumerate}

\section{Citations and References}
When you use the citation search feature, references will be automatically managed. For example, you might cite a paper like this: \cite{example2024}. All citations will be automatically added to your \texttt{cite.bib} file.

\section{Conclusion}
This template demonstrates basic \LaTeX{} features including mathematical formulas, tables, and lists. The bibliography system is now set up to work seamlessly with Bibby AI's citation management. Feel free to modify it for your needs!

\bibliographystyle{plain}
\bibliography{cite}

\end{document}\documentclass{article}
\usepackage[utf8]{inputenc}
\usepackage{booktabs}
\usepackage{amsmath}
\usepackage{natbib}

\title{My Document}
\author{Your Name}
\date{\today}

\begin{document}

\maketitle

\section{Introduction}
This is a sample document demonstrating various \LaTeX{} features including mathematical equations, tables, and formatting. You can cite papers using the citation search feature, and they will be automatically added to your \texttt{cite.bib} file.

\section{Mathematical Expressions}
Here are some example equations:

\subsection{Quadratic Formula}
The quadratic formula is given by:
\[x = \frac{-b \pm \sqrt{b^2 - 4ac}}{2a}\]

\subsection{Matrix Example}
A sample matrix:
\[
\begin{pmatrix}
a_{11} & a_{12} & a_{13} \\
a_{21} & a_{22} & a_{23} \\
a_{31} & a_{32} & a_{33}
\end{pmatrix}
\]

\section{Tables}
Below is a sample table showing data:

\begin{table}[h]
\centering
\begin{tabular}{@{}lrrr@{}}
\toprule
\textbf{Category} & \textbf{Value 1} & \textbf{Value 2} & \textbf{Total} \\
\midrule
Item A & 10.5 & 20.3 & 30.8 \\
Item B & 15.2 & 25.1 & 40.3 \\
Item C & 22.1 & 12.7 & 34.8 \\
\bottomrule
\end{tabular}
\caption{Sample Data Table}
\label{tab:sample}
\end{table}

\section{Lists}
\subsection{Bullet Points}
\begin{itemize}
  \item First item
  \item Second item
  \item Third item with sub-items:
    \begin{itemize}
      \item Sub-item 1
      \item Sub-item 2
    \end{itemize}
\end{itemize}

\subsection{Numbered List}
\begin{enumerate}
  \item First step
  \item Second step
  \item Third step
\end{enumerate}

\section{Citations and References}
When you use the citation search feature, references will be automatically managed. For example, you might cite a paper like this: \cite{example2024}. All citations will be automatically added to your \texttt{cite.bib} file.

\section{Conclusion}
This template demonstrates basic \LaTeX{} features including mathematical formulas, tables, and lists. The bibliography system is now set up to work seamlessly with Bibby AI's citation management. Feel free to modify it for your needs!

\bibliographystyle{plain}
\bibliography{cite}

\end{document}\documentclass{article}
\usepackage[utf8]{inputenc}
\usepackage{booktabs}
\usepackage{amsmath}
\usepackage{natbib}

\title{My Document}
\author{Your Name}
\date{\today}

\begin{document}

\maketitle

\section{Introduction}
This is a sample document demonstrating various \LaTeX{} features including mathematical equations, tables, and formatting. You can cite papers using the citation search feature, and they will be automatically added to your \texttt{cite.bib} file.

\section{Mathematical Expressions}
Here are some example equations:

\subsection{Quadratic Formula}
The quadratic formula is given by:
\[x = \frac{-b \pm \sqrt{b^2 - 4ac}}{2a}\]

\subsection{Matrix Example}
A sample matrix:
\[
\begin{pmatrix}
a_{11} & a_{12} & a_{13} \\
a_{21} & a_{22} & a_{23} \\
a_{31} & a_{32} & a_{33}
\end{pmatrix}
\]

\section{Tables}
Below is a sample table showing data:

\begin{table}[h]
\centering
\begin{tabular}{@{}lrrr@{}}
\toprule
\textbf{Category} & \textbf{Value 1} & \textbf{Value 2} & \textbf{Total} \\
\midrule
Item A & 10.5 & 20.3 & 30.8 \\
Item B & 15.2 & 25.1 & 40.3 \\
Item C & 22.1 & 12.7 & 34.8 \\
\bottomrule
\end{tabular}
\caption{Sample Data Table}
\label{tab:sample}
\end{table}

\section{Lists}
\subsection{Bullet Points}
\begin{itemize}
  \item First item
  \item Second item
  \item Third item with sub-items:
    \begin{itemize}
      \item Sub-item 1
      \item Sub-item 2
    \end{itemize}
\end{itemize}

\subsection{Numbered List}
\begin{enumerate}
  \item First step
  \item Second step
  \item Third step
\end{enumerate}

\section{Citations and References}
When you use the citation search feature, references will be automatically managed. For example, you might cite a paper like this: \cite{example2024}. All citations will be automatically added to your \texttt{cite.bib} file.

\section{Conclusion}
This template demonstrates basic \LaTeX{} features including mathematical formulas, tables, and lists. The bibliography system is now set up to work seamlessly with Bibby AI's citation management. Feel free to modify it for your needs!

\bibliographystyle{plain}
\bibliography{cite}

\end{document}\documentclass{article}
\usepackage[utf8]{inputenc}
\usepackage{booktabs}
\usepackage{amsmath}
\usepackage{natbib}

\title{My Document}
\author{Your Name}
\date{\today}

\begin{document}

\maketitle

\section{Introduction}
This is a sample document demonstrating various \LaTeX{} features including mathematical equations, tables, and formatting. You can cite papers using the citation search feature, and they will be automatically added to your \texttt{cite.bib} file.

\section{Mathematical Expressions}
Here are some example equations:

\subsection{Quadratic Formula}
The quadratic formula is given by:
\[x = \frac{-b \pm \sqrt{b^2 - 4ac}}{2a}\]

\subsection{Matrix Example}
A sample matrix:
\[
\begin{pmatrix}
a_{11} & a_{12} & a_{13} \\
a_{21} & a_{22} & a_{23} \\
a_{31} & a_{32} & a_{33}
\end{pmatrix}
\]

\section{Tables}
Below is a sample table showing data:

\begin{table}[h]
\centering
\begin{tabular}{@{}lrrr@{}}
\toprule
\textbf{Category} & \textbf{Value 1} & \textbf{Value 2} & \textbf{Total} \\
\midrule
Item A & 10.5 & 20.3 & 30.8 \\
Item B & 15.2 & 25.1 & 40.3 \\
Item C & 22.1 & 12.7 & 34.8 \\
\bottomrule
\end{tabular}
\caption{Sample Data Table}
\label{tab:sample}
\end{table}

\section{Lists}
\subsection{Bullet Points}
\begin{itemize}
  \item First item
  \item Second item
  \item Third item with sub-items:
    \begin{itemize}
      \item Sub-item 1
      \item Sub-item 2
    \end{itemize}
\end{itemize}

\subsection{Numbered List}
\begin{enumerate}
  \item First step
  \item Second step
  \item Third step
\end{enumerate}

\section{Citations and References}
When you use the citation search feature, references will be automatically managed. For example, you might cite a paper like this: \cite{example2024}. All citations will be automatically added to your \texttt{cite.bib} file.

\section{Conclusion}
This template demonstrates basic \LaTeX{} features including mathematical formulas, tables, and lists. The bibliography system is now set up to work seamlessly with Bibby AI's citation management. Feel free to modify it for your needs!

\bibliographystyle{plain}
\bibliography{cite}

\end{document}\documentclass{article}
\usepackage[utf8]{inputenc}
\usepackage{booktabs}
\usepackage{amsmath}
\usepackage{natbib}

\title{My Document}
\author{Your Name}
\date{\today}

\begin{document}

\maketitle

\section{Introduction}
This is a sample document demonstrating various \LaTeX{} features including mathematical equations, tables, and formatting. You can cite papers using the citation search feature, and they will be automatically added to your \texttt{cite.bib} file.

\section{Mathematical Expressions}
Here are some example equations:

\subsection{Quadratic Formula}
The quadratic formula is given by:
\[x = \frac{-b \pm \sqrt{b^2 - 4ac}}{2a}\]

\subsection{Matrix Example}
A sample matrix:
\[
\begin{pmatrix}
a_{11} & a_{12} & a_{13} \\
a_{21} & a_{22} & a_{23} \\
a_{31} & a_{32} & a_{33}
\end{pmatrix}
\]

\section{Tables}
Below is a sample table showing data:

\begin{table}[h]
\centering
\begin{tabular}{@{}lrrr@{}}
\toprule
\textbf{Category} & \textbf{Value 1} & \textbf{Value 2} & \textbf{Total} \\
\midrule
Item A & 10.5 & 20.3 & 30.8 \\
Item B & 15.2 & 25.1 & 40.3 \\
Item C & 22.1 & 12.7 & 34.8 \\
\bottomrule
\end{tabular}
\caption{Sample Data Table}
\label{tab:sample}
\end{table}

\section{Lists}
\subsection{Bullet Points}
\begin{itemize}
  \item First item
  \item Second item
  \item Third item with sub-items:
    \begin{itemize}
      \item Sub-item 1
      \item Sub-item 2
    \end{itemize}
\end{itemize}

\subsection{Numbered List}
\begin{enumerate}
  \item First step
  \item Second step
  \item Third step
\end{enumerate}

\section{Citations and References}
When you use the citation search feature, references will be automatically managed. For example, you might cite a paper like this: \cite{example2024}. All citations will be automatically added to your \texttt{cite.bib} file.

\section{Conclusion}
This template demonstrates basic \LaTeX{} features including mathematical formulas, tables, and lists. The bibliography system is now set up to work seamlessly with Bibby AI's citation management. Feel free to modify it for your needs!

\bibliographystyle{plain}
\bibliography{cite}

\end{document}\documentclass{article}
\usepackage[utf8]{inputenc}
\usepackage{booktabs}
\usepackage{amsmath}
\usepackage{natbib}

\title{My Document}
\author{Your Name}
\date{\today}

\begin{document}

\maketitle

\section{Introduction}
This is a sample document demonstrating various \LaTeX{} features including mathematical equations, tables, and formatting. You can cite papers using the citation search feature, and they will be automatically added to your \texttt{cite.bib} file.

\section{Mathematical Expressions}
Here are some example equations:

\subsection{Quadratic Formula}
The quadratic formula is given by:
\[x = \frac{-b \pm \sqrt{b^2 - 4ac}}{2a}\]

\subsection{Matrix Example}
A sample matrix:
\[
\begin{pmatrix}
a_{11} & a_{12} & a_{13} \\
a_{21} & a_{22} & a_{23} \\
a_{31} & a_{32} & a_{33}
\end{pmatrix}
\]

\section{Tables}
Below is a sample table showing data:

\begin{table}[h]
\centering
\begin{tabular}{@{}lrrr@{}}
\toprule
\textbf{Category} & \textbf{Value 1} & \textbf{Value 2} & \textbf{Total} \\
\midrule
Item A & 10.5 & 20.3 & 30.8 \\
Item B & 15.2 & 25.1 & 40.3 \\
Item C & 22.1 & 12.7 & 34.8 \\
\bottomrule
\end{tabular}
\caption{Sample Data Table}
\label{tab:sample}
\end{table}

\section{Lists}
\subsection{Bullet Points}
\begin{itemize}
  \item First item
  \item Second item
  \item Third item with sub-items:
    \begin{itemize}
      \item Sub-item 1
      \item Sub-item 2
    \end{itemize}
\end{itemize}

\subsection{Numbered List}
\begin{enumerate}
  \item First step
  \item Second step
  \item Third step
\end{enumerate}

\section{Citations and References}
When you use the citation search feature, references will be automatically managed. For example, you might cite a paper like this: \cite{example2024}. All citations will be automatically added to your \texttt{cite.bib} file.

\section{Conclusion}
This template demonstrates basic \LaTeX{} features including mathematical formulas, tables, and lists. The bibliography system is now set up to work seamlessly with Bibby AI's citation management. Feel free to modify it for your needs!

\bibliographystyle{plain}
\bibliography{cite}

\end{document}\documentclass{article}
\usepackage[utf8]{inputenc}
\usepackage{booktabs}
\usepackage{amsmath}
\usepackage{natbib}

\title{My Document}
\author{Your Name}
\date{\today}

\begin{document}

\maketitle

\section{Introduction}
This is a sample document demonstrating various \LaTeX{} features including mathematical equations, tables, and formatting. You can cite papers using the citation search feature, and they will be automatically added to your \texttt{cite.bib} file.

\section{Mathematical Expressions}
Here are some example equations:

\subsection{Quadratic Formula}
The quadratic formula is given by:
\[x = \frac{-b \pm \sqrt{b^2 - 4ac}}{2a}\]

\subsection{Matrix Example}
A sample matrix:
\[
\begin{pmatrix}
a_{11} & a_{12} & a_{13} \\
a_{21} & a_{22} & a_{23} \\
a_{31} & a_{32} & a_{33}
\end{pmatrix}
\]

\section{Tables}
Below is a sample table showing data:

\begin{table}[h]
\centering
\begin{tabular}{@{}lrrr@{}}
\toprule
\textbf{Category} & \textbf{Value 1} & \textbf{Value 2} & \textbf{Total} \\
\midrule
Item A & 10.5 & 20.3 & 30.8 \\
Item B & 15.2 & 25.1 & 40.3 \\
Item C & 22.1 & 12.7 & 34.8 \\
\bottomrule
\end{tabular}
\caption{Sample Data Table}
\label{tab:sample}
\end{table}

\section{Lists}
\subsection{Bullet Points}
\begin{itemize}
  \item First item
  \item Second item
  \item Third item with sub-items:
    \begin{itemize}
      \item Sub-item 1
      \item Sub-item 2
    \end{itemize}
\end{itemize}

\subsection{Numbered List}
\begin{enumerate}
  \item First step
  \item Second step
  \item Third step
\end{enumerate}

\section{Citations and References}
When you use the citation search feature, references will be automatically managed. For example, you might cite a paper like this: \cite{example2024}. All citations will be automatically added to your \texttt{cite.bib} file.

\section{Conclusion}
This template demonstrates basic \LaTeX{} features including mathematical formulas, tables, and lists. The bibliography system is now set up to work seamlessly with Bibby AI's citation management. Feel free to modify it for your needs!

\bibliographystyle{plain}
\bibliography{cite}

\end{document}\documentclass{article}
\usepackage[utf8]{inputenc}
\usepackage{booktabs}
\usepackage{amsmath}
\usepackage{natbib}

\title{My Document}
\author{Your Name}
\date{\today}

\begin{document}

\maketitle

\section{Introduction}
This is a sample document demonstrating various \LaTeX{} features including mathematical equations, tables, and formatting. You can cite papers using the citation search feature, and they will be automatically added to your \texttt{cite.bib} file.

\section{Mathematical Expressions}
Here are some example equations:

\subsection{Quadratic Formula}
The quadratic formula is given by:
\[x = \frac{-b \pm \sqrt{b^2 - 4ac}}{2a}\]

\subsection{Matrix Example}
A sample matrix:
\[
\begin{pmatrix}
a_{11} & a_{12} & a_{13} \\
a_{21} & a_{22} & a_{23} \\
a_{31} & a_{32} & a_{33}
\end{pmatrix}
\]

\section{Tables}
Below is a sample table showing data:

\begin{table}[h]
\centering
\begin{tabular}{@{}lrrr@{}}
\toprule
\textbf{Category} & \textbf{Value 1} & \textbf{Value 2} & \textbf{Total} \\
\midrule
Item A & 10.5 & 20.3 & 30.8 \\
Item B & 15.2 & 25.1 & 40.3 \\
Item C & 22.1 & 12.7 & 34.8 \\
\bottomrule
\end{tabular}
\caption{Sample Data Table}
\label{tab:sample}
\end{table}

\section{Lists}
\subsection{Bullet Points}
\begin{itemize}
  \item First item
  \item Second item
  \item Third item with sub-items:
    \begin{itemize}
      \item Sub-item 1
      \item Sub-item 2
    \end{itemize}
\end{itemize}

\subsection{Numbered List}
\begin{enumerate}
  \item First step
  \item Second step
  \item Third step
\end{enumerate}

\section{Citations and References}
When you use the citation search feature, references will be automatically managed. For example, you might cite a paper like this: \cite{example2024}. All citations will be automatically added to your \texttt{cite.bib} file.

\section{Conclusion}
This template demonstrates basic \LaTeX{} features including mathematical formulas, tables, and lists. The bibliography system is now set up to work seamlessly with Bibby AI's citation management. Feel free to modify it for your needs!

\bibliographystyle{plain}
\bibliography{cite}

\end{document}\documentclass{article}
\usepackage[utf8]{inputenc}
\usepackage{booktabs}
\usepackage{amsmath}
\usepackage{natbib}

\title{My Document}
\author{Your Name}
\date{\today}

\begin{document}

\maketitle

\section{Introduction}
This is a sample document demonstrating various \LaTeX{} features including mathematical equations, tables, and formatting. You can cite papers using the citation search feature, and they will be automatically added to your \texttt{cite.bib} file.

\section{Mathematical Expressions}
Here are some example equations:

\subsection{Quadratic Formula}
The quadratic formula is given by:
\[x = \frac{-b \pm \sqrt{b^2 - 4ac}}{2a}\]

\subsection{Matrix Example}
A sample matrix:
\[
\begin{pmatrix}
a_{11} & a_{12} & a_{13} \\
a_{21} & a_{22} & a_{23} \\
a_{31} & a_{32} & a_{33}
\end{pmatrix}
\]

\section{Tables}
Below is a sample table showing data:

\begin{table}[h]
\centering
\begin{tabular}{@{}lrrr@{}}
\toprule
\textbf{Category} & \textbf{Value 1} & \textbf{Value 2} & \textbf{Total} \\
\midrule
Item A & 10.5 & 20.3 & 30.8 \\
Item B & 15.2 & 25.1 & 40.3 \\
Item C & 22.1 & 12.7 & 34.8 \\
\bottomrule
\end{tabular}
\caption{Sample Data Table}
\label{tab:sample}
\end{table}

\section{Lists}
\subsection{Bullet Points}
\begin{itemize}
  \item First item
  \item Second item
  \item Third item with sub-items:
    \begin{itemize}
      \item Sub-item 1
      \item Sub-item 2
    \end{itemize}
\end{itemize}

\subsection{Numbered List}
\begin{enumerate}
  \item First step
  \item Second step
  \item Third step
\end{enumerate}

\section{Citations and References}
When you use the citation search feature, references will be automatically managed. For example, you might cite a paper like this: \cite{example2024}. All citations will be automatically added to your \texttt{cite.bib} file.

\section{Conclusion}
This template demonstrates basic \LaTeX{} features including mathematical formulas, tables, and lists. The bibliography system is now set up to work seamlessly with Bibby AI's citation management. Feel free to modify it for your needs!

\bibliographystyle{plain}
\bibliography{cite}

\end{document}\documentclass{article}
\usepackage[utf8]{inputenc}
\usepackage{booktabs}
\usepackage{amsmath}
\usepackage{natbib}

\title{My Document}
\author{Your Name}
\date{\today}

\begin{document}

\maketitle

\section{Introduction}
This is a sample document demonstrating various \LaTeX{} features including mathematical equations, tables, and formatting. You can cite papers using the citation search feature, and they will be automatically added to your \texttt{cite.bib} file.

\section{Mathematical Expressions}
Here are some example equations:

\subsection{Quadratic Formula}
The quadratic formula is given by:
\[x = \frac{-b \pm \sqrt{b^2 - 4ac}}{2a}\]

\subsection{Matrix Example}
A sample matrix:
\[
\begin{pmatrix}
a_{11} & a_{12} & a_{13} \\
a_{21} & a_{22} & a_{23} \\
a_{31} & a_{32} & a_{33}
\end{pmatrix}
\]

\section{Tables}
Below is a sample table showing data:

\begin{table}[h]
\centering
\begin{tabular}{@{}lrrr@{}}
\toprule
\textbf{Category} & \textbf{Value 1} & \textbf{Value 2} & \textbf{Total} \\
\midrule
Item A & 10.5 & 20.3 & 30.8 \\
Item B & 15.2 & 25.1 & 40.3 \\
Item C & 22.1 & 12.7 & 34.8 \\
\bottomrule
\end{tabular}
\caption{Sample Data Table}
\label{tab:sample}
\end{table}

\section{Lists}
\subsection{Bullet Points}
\begin{itemize}
  \item First item
  \item Second item
  \item Third item with sub-items:
    \begin{itemize}
      \item Sub-item 1
      \item Sub-item 2
    \end{itemize}
\end{itemize}

\subsection{Numbered List}
\begin{enumerate}
  \item First step
  \item Second step
  \item Third step
\end{enumerate}

\section{Citations and References}
When you use the citation search feature, references will be automatically managed. For example, you might cite a paper like this: \cite{example2024}. All citations will be automatically added to your \texttt{cite.bib} file.

\section{Conclusion}
This template demonstrates basic \LaTeX{} features including mathematical formulas, tables, and lists. The bibliography system is now set up to work seamlessly with Bibby AI's citation management. Feel free to modify it for your needs!

\bibliographystyle{plain}
\bibliography{cite}

\end{document}\documentclass{article}
\usepackage[utf8]{inputenc}
\usepackage{booktabs}
\usepackage{amsmath}
\usepackage{natbib}

\title{My Document}
\author{Your Name}
\date{\today}

\begin{document}

\maketitle

\section{Introduction}
This is a sample document demonstrating various \LaTeX{} features including mathematical equations, tables, and formatting. You can cite papers using the citation search feature, and they will be automatically added to your \texttt{cite.bib} file.

\section{Mathematical Expressions}
Here are some example equations:

\subsection{Quadratic Formula}
The quadratic formula is given by:
\[x = \frac{-b \pm \sqrt{b^2 - 4ac}}{2a}\]

\subsection{Matrix Example}
A sample matrix:
\[
\begin{pmatrix}
a_{11} & a_{12} & a_{13} \\
a_{21} & a_{22} & a_{23} \\
a_{31} & a_{32} & a_{33}
\end{pmatrix}
\]

\section{Tables}
Below is a sample table showing data:

\begin{table}[h]
\centering
\begin{tabular}{@{}lrrr@{}}
\toprule
\textbf{Category} & \textbf{Value 1} & \textbf{Value 2} & \textbf{Total} \\
\midrule
Item A & 10.5 & 20.3 & 30.8 \\
Item B & 15.2 & 25.1 & 40.3 \\
Item C & 22.1 & 12.7 & 34.8 \\
\bottomrule
\end{tabular}
\caption{Sample Data Table}
\label{tab:sample}
\end{table}

\section{Lists}
\subsection{Bullet Points}
\begin{itemize}
  \item First item
  \item Second item
  \item Third item with sub-items:
    \begin{itemize}
      \item Sub-item 1
      \item Sub-item 2
    \end{itemize}
\end{itemize}

\subsection{Numbered List}
\begin{enumerate}
  \item First step
  \item Second step
  \item Third step
\end{enumerate}

\section{Citations and References}
When you use the citation search feature, references will be automatically managed. For example, you might cite a paper like this: \cite{example2024}. All citations will be automatically added to your \texttt{cite.bib} file.

\section{Conclusion}
This template demonstrates basic \LaTeX{} features including mathematical formulas, tables, and lists. The bibliography system is now set up to work seamlessly with Bibby AI's citation management. Feel free to modify it for your needs!

\bibliographystyle{plain}
\bibliography{cite}

\end{document}\documentclass{article}
\usepackage[utf8]{inputenc}
\usepackage{booktabs}
\usepackage{amsmath}
\usepackage{natbib}

\title{My Document}
\author{Your Name}
\date{\today}

\begin{document}

\maketitle

\section{Introduction}
This is a sample document demonstrating various \LaTeX{} features including mathematical equations, tables, and formatting. You can cite papers using the citation search feature, and they will be automatically added to your \texttt{cite.bib} file.

\section{Mathematical Expressions}
Here are some example equations:

\subsection{Quadratic Formula}
The quadratic formula is given by:
\[x = \frac{-b \pm \sqrt{b^2 - 4ac}}{2a}\]

\subsection{Matrix Example}
A sample matrix:
\[
\begin{pmatrix}
a_{11} & a_{12} & a_{13} \\
a_{21} & a_{22} & a_{23} \\
a_{31} & a_{32} & a_{33}
\end{pmatrix}
\]

\section{Tables}
Below is a sample table showing data:

\begin{table}[h]
\centering
\begin{tabular}{@{}lrrr@{}}
\toprule
\textbf{Category} & \textbf{Value 1} & \textbf{Value 2} & \textbf{Total} \\
\midrule
Item A & 10.5 & 20.3 & 30.8 \\
Item B & 15.2 & 25.1 & 40.3 \\
Item C & 22.1 & 12.7 & 34.8 \\
\bottomrule
\end{tabular}
\caption{Sample Data Table}
\label{tab:sample}
\end{table}

\section{Lists}
\subsection{Bullet Points}
\begin{itemize}
  \item First item
  \item Second item
  \item Third item with sub-items:
    \begin{itemize}
      \item Sub-item 1
      \item Sub-item 2
    \end{itemize}
\end{itemize}

\subsection{Numbered List}
\begin{enumerate}
  \item First step
  \item Second step
  \item Third step
\end{enumerate}

\section{Citations and References}
When you use the citation search feature, references will be automatically managed. For example, you might cite a paper like this: \cite{example2024}. All citations will be automatically added to your \texttt{cite.bib} file.

\section{Conclusion}
This template demonstrates basic \LaTeX{} features including mathematical formulas, tables, and lists. The bibliography system is now set up to work seamlessly with Bibby AI's citation management. Feel free to modify it for your needs!

\bibliographystyle{plain}
\bibliography{cite}

\end{document}\documentclass{article}
\usepackage[utf8]{inputenc}
\usepackage{booktabs}
\usepackage{amsmath}
\usepackage{natbib}

\title{My Document}
\author{Your Name}
\date{\today}

\begin{document}

\maketitle

\section{Introduction}
This is a sample document demonstrating various \LaTeX{} features including mathematical equations, tables, and formatting. You can cite papers using the citation search feature, and they will be automatically added to your \texttt{cite.bib} file.

\section{Mathematical Expressions}
Here are some example equations:

\subsection{Quadratic Formula}
The quadratic formula is given by:
\[x = \frac{-b \pm \sqrt{b^2 - 4ac}}{2a}\]

\subsection{Matrix Example}
A sample matrix:
\[
\begin{pmatrix}
a_{11} & a_{12} & a_{13} \\
a_{21} & a_{22} & a_{23} \\
a_{31} & a_{32} & a_{33}
\end{pmatrix}
\]

\section{Tables}
Below is a sample table showing data:

\begin{table}[h]
\centering
\begin{tabular}{@{}lrrr@{}}
\toprule
\textbf{Category} & \textbf{Value 1} & \textbf{Value 2} & \textbf{Total} \\
\midrule
Item A & 10.5 & 20.3 & 30.8 \\
Item B & 15.2 & 25.1 & 40.3 \\
Item C & 22.1 & 12.7 & 34.8 \\
\bottomrule
\end{tabular}
\caption{Sample Data Table}
\label{tab:sample}
\end{table}

\section{Lists}
\subsection{Bullet Points}
\begin{itemize}
  \item First item
  \item Second item
  \item Third item with sub-items:
    \begin{itemize}
      \item Sub-item 1
      \item Sub-item 2
    \end{itemize}
\end{itemize}

\subsection{Numbered List}
\begin{enumerate}
  \item First step
  \item Second step
  \item Third step
\end{enumerate}

\section{Citations and References}
When you use the citation search feature, references will be automatically managed. For example, you might cite a paper like this: \cite{example2024}. All citations will be automatically added to your \texttt{cite.bib} file.

\section{Conclusion}
This template demonstrates basic \LaTeX{} features including mathematical formulas, tables, and lists. The bibliography system is now set up to work seamlessly with Bibby AI's citation management. Feel free to modify it for your needs!

\bibliographystyle{plain}
\bibliography{cite}

\end{document}\documentclass{article}
\usepackage[utf8]{inputenc}
\usepackage{booktabs}
\usepackage{amsmath}
\usepackage{natbib}

\title{My Document}
\author{Your Name}
\date{\today}

\begin{document}

\maketitle

\section{Introduction}
This is a sample document demonstrating various \LaTeX{} features including mathematical equations, tables, and formatting. You can cite papers using the citation search feature, and they will be automatically added to your \texttt{cite.bib} file.

\section{Mathematical Expressions}
Here are some example equations:

\subsection{Quadratic Formula}
The quadratic formula is given by:
\[x = \frac{-b \pm \sqrt{b^2 - 4ac}}{2a}\]

\subsection{Matrix Example}
A sample matrix:
\[
\begin{pmatrix}
a_{11} & a_{12} & a_{13} \\
a_{21} & a_{22} & a_{23} \\
a_{31} & a_{32} & a_{33}
\end{pmatrix}
\]

\section{Tables}
Below is a sample table showing data:

\begin{table}[h]
\centering
\begin{tabular}{@{}lrrr@{}}
\toprule
\textbf{Category} & \textbf{Value 1} & \textbf{Value 2} & \textbf{Total} \\
\midrule
Item A & 10.5 & 20.3 & 30.8 \\
Item B & 15.2 & 25.1 & 40.3 \\
Item C & 22.1 & 12.7 & 34.8 \\
\bottomrule
\end{tabular}
\caption{Sample Data Table}
\label{tab:sample}
\end{table}

\section{Lists}
\subsection{Bullet Points}
\begin{itemize}
  \item First item
  \item Second item
  \item Third item with sub-items:
    \begin{itemize}
      \item Sub-item 1
      \item Sub-item 2
    \end{itemize}
\end{itemize}

\subsection{Numbered List}
\begin{enumerate}
  \item First step
  \item Second step
  \item Third step
\end{enumerate}

\section{Citations and References}
When you use the citation search feature, references will be automatically managed. For example, you might cite a paper like this: \cite{example2024}. All citations will be automatically added to your \texttt{cite.bib} file.

\section{Conclusion}
This template demonstrates basic \LaTeX{} features including mathematical formulas, tables, and lists. The bibliography system is now set up to work seamlessly with Bibby AI's citation management. Feel free to modify it for your needs!

\bibliographystyle{plain}
\bibliography{cite}

\end{document}\documentclass{article}
\usepackage[utf8]{inputenc}
\usepackage{booktabs}
\usepackage{amsmath}
\usepackage{natbib}

\title{My Document}
\author{Your Name}
\date{\today}

\begin{document}

\maketitle

\section{Introduction}
This is a sample document demonstrating various \LaTeX{} features including mathematical equations, tables, and formatting. You can cite papers using the citation search feature, and they will be automatically added to your \texttt{cite.bib} file.

\section{Mathematical Expressions}
Here are some example equations:

\subsection{Quadratic Formula}
The quadratic formula is given by:
\[x = \frac{-b \pm \sqrt{b^2 - 4ac}}{2a}\]

\subsection{Matrix Example}
A sample matrix:
\[
\begin{pmatrix}
a_{11} & a_{12} & a_{13} \\
a_{21} & a_{22} & a_{23} \\
a_{31} & a_{32} & a_{33}
\end{pmatrix}
\]

\section{Tables}
Below is a sample table showing data:

\begin{table}[h]
\centering
\begin{tabular}{@{}lrrr@{}}
\toprule
\textbf{Category} & \textbf{Value 1} & \textbf{Value 2} & \textbf{Total} \\
\midrule
Item A & 10.5 & 20.3 & 30.8 \\
Item B & 15.2 & 25.1 & 40.3 \\
Item C & 22.1 & 12.7 & 34.8 \\
\bottomrule
\end{tabular}
\caption{Sample Data Table}
\label{tab:sample}
\end{table}

\section{Lists}
\subsection{Bullet Points}
\begin{itemize}
  \item First item
  \item Second item
  \item Third item with sub-items:
    \begin{itemize}
      \item Sub-item 1
      \item Sub-item 2
    \end{itemize}
\end{itemize}

\subsection{Numbered List}
\begin{enumerate}
  \item First step
  \item Second step
  \item Third step
\end{enumerate}

\section{Citations and References}
When you use the citation search feature, references will be automatically managed. For example, you might cite a paper like this: \cite{example2024}. All citations will be automatically added to your \texttt{cite.bib} file.

\section{Conclusion}
This template demonstrates basic \LaTeX{} features including mathematical formulas, tables, and lists. The bibliography system is now set up to work seamlessly with Bibby AI's citation management. Feel free to modify it for your needs!

\bibliographystyle{plain}
\bibliography{cite}

\end{document}\documentclass{article}
\usepackage[utf8]{inputenc}
\usepackage{booktabs}
\usepackage{amsmath}
\usepackage{natbib}

\title{My Document}
\author{Your Name}
\date{\today}

\begin{document}

\maketitle

\section{Introduction}
This is a sample document demonstrating various \LaTeX{} features including mathematical equations, tables, and formatting. You can cite papers using the citation search feature, and they will be automatically added to your \texttt{cite.bib} file.

\section{Mathematical Expressions}
Here are some example equations:

\subsection{Quadratic Formula}
The quadratic formula is given by:
\[x = \frac{-b \pm \sqrt{b^2 - 4ac}}{2a}\]

\subsection{Matrix Example}
A sample matrix:
\[
\begin{pmatrix}
a_{11} & a_{12} & a_{13} \\
a_{21} & a_{22} & a_{23} \\
a_{31} & a_{32} & a_{33}
\end{pmatrix}
\]

\section{Tables}
Below is a sample table showing data:

\begin{table}[h]
\centering
\begin{tabular}{@{}lrrr@{}}
\toprule
\textbf{Category} & \textbf{Value 1} & \textbf{Value 2} & \textbf{Total} \\
\midrule
Item A & 10.5 & 20.3 & 30.8 \\
Item B & 15.2 & 25.1 & 40.3 \\
Item C & 22.1 & 12.7 & 34.8 \\
\bottomrule
\end{tabular}
\caption{Sample Data Table}
\label{tab:sample}
\end{table}

\section{Lists}
\subsection{Bullet Points}
\begin{itemize}
  \item First item
  \item Second item
  \item Third item with sub-items:
    \begin{itemize}
      \item Sub-item 1
      \item Sub-item 2
    \end{itemize}
\end{itemize}

\subsection{Numbered List}
\begin{enumerate}
  \item First step
  \item Second step
  \item Third step
\end{enumerate}

\section{Citations and References}
When you use the citation search feature, references will be automatically managed. For example, you might cite a paper like this: \cite{example2024}. All citations will be automatically added to your \texttt{cite.bib} file.

\section{Conclusion}
This template demonstrates basic \LaTeX{} features including mathematical formulas, tables, and lists. The bibliography system is now set up to work seamlessly with Bibby AI's citation management. Feel free to modify it for your needs!

\bibliographystyle{plain}
\bibliography{cite}

\end{document}\documentclass{article}
\usepackage[utf8]{inputenc}
\usepackage{booktabs}
\usepackage{amsmath}
\usepackage{natbib}

\title{My Document}
\author{Your Name}
\date{\today}

\begin{document}

\maketitle

\section{Introduction}
This is a sample document demonstrating various \LaTeX{} features including mathematical equations, tables, and formatting. You can cite papers using the citation search feature, and they will be automatically added to your \texttt{cite.bib} file.

\section{Mathematical Expressions}
Here are some example equations:

\subsection{Quadratic Formula}
The quadratic formula is given by:
\[x = \frac{-b \pm \sqrt{b^2 - 4ac}}{2a}\]

\subsection{Matrix Example}
A sample matrix:
\[
\begin{pmatrix}
a_{11} & a_{12} & a_{13} \\
a_{21} & a_{22} & a_{23} \\
a_{31} & a_{32} & a_{33}
\end{pmatrix}
\]

\section{Tables}
Below is a sample table showing data:

\begin{table}[h]
\centering
\begin{tabular}{@{}lrrr@{}}
\toprule
\textbf{Category} & \textbf{Value 1} & \textbf{Value 2} & \textbf{Total} \\
\midrule
Item A & 10.5 & 20.3 & 30.8 \\
Item B & 15.2 & 25.1 & 40.3 \\
Item C & 22.1 & 12.7 & 34.8 \\
\bottomrule
\end{tabular}
\caption{Sample Data Table}
\label{tab:sample}
\end{table}

\section{Lists}
\subsection{Bullet Points}
\begin{itemize}
  \item First item
  \item Second item
  \item Third item with sub-items:
    \begin{itemize}
      \item Sub-item 1
      \item Sub-item 2
    \end{itemize}
\end{itemize}

\subsection{Numbered List}
\begin{enumerate}
  \item First step
  \item Second step
  \item Third step
\end{enumerate}

\section{Citations and References}
When you use the citation search feature, references will be automatically managed. For example, you might cite a paper like this: \cite{example2024}. All citations will be automatically added to your \texttt{cite.bib} file.

\section{Conclusion}
This template demonstrates basic \LaTeX{} features including mathematical formulas, tables, and lists. The bibliography system is now set up to work seamlessly with Bibby AI's citation management. Feel free to modify it for your needs!

\bibliographystyle{plain}
\bibliography{cite}

\end{document}\documentclass{article}
\usepackage[utf8]{inputenc}
\usepackage{booktabs}
\usepackage{amsmath}
\usepackage{natbib}

\title{My Document}
\author{Your Name}
\date{\today}

\begin{document}

\maketitle

\section{Introduction}
This is a sample document demonstrating various \LaTeX{} features including mathematical equations, tables, and formatting. You can cite papers using the citation search feature, and they will be automatically added to your \texttt{cite.bib} file.

\section{Mathematical Expressions}
Here are some example equations:

\subsection{Quadratic Formula}
The quadratic formula is given by:
\[x = \frac{-b \pm \sqrt{b^2 - 4ac}}{2a}\]

\subsection{Matrix Example}
A sample matrix:
\[
\begin{pmatrix}
a_{11} & a_{12} & a_{13} \\
a_{21} & a_{22} & a_{23} \\
a_{31} & a_{32} & a_{33}
\end{pmatrix}
\]

\section{Tables}
Below is a sample table showing data:

\begin{table}[h]
\centering
\begin{tabular}{@{}lrrr@{}}
\toprule
\textbf{Category} & \textbf{Value 1} & \textbf{Value 2} & \textbf{Total} \\
\midrule
Item A & 10.5 & 20.3 & 30.8 \\
Item B & 15.2 & 25.1 & 40.3 \\
Item C & 22.1 & 12.7 & 34.8 \\
\bottomrule
\end{tabular}
\caption{Sample Data Table}
\label{tab:sample}
\end{table}

\section{Lists}
\subsection{Bullet Points}
\begin{itemize}
  \item First item
  \item Second item
  \item Third item with sub-items:
    \begin{itemize}
      \item Sub-item 1
      \item Sub-item 2
    \end{itemize}
\end{itemize}

\subsection{Numbered List}
\begin{enumerate}
  \item First step
  \item Second step
  \item Third step
\end{enumerate}

\section{Citations and References}
When you use the citation search feature, references will be automatically managed. For example, you might cite a paper like this: \cite{example2024}. All citations will be automatically added to your \texttt{cite.bib} file.

\section{Conclusion}
This template demonstrates basic \LaTeX{} features including mathematical formulas, tables, and lists. The bibliography system is now set up to work seamlessly with Bibby AI's citation management. Feel free to modify it for your needs!

\bibliographystyle{plain}
\bibliography{cite}

\end{document}\documentclass{article}
\usepackage[utf8]{inputenc}
\usepackage{booktabs}
\usepackage{amsmath}
\usepackage{natbib}

\title{My Document}
\author{Your Name}
\date{\today}

\begin{document}

\maketitle

\section{Introduction}
This is a sample document demonstrating various \LaTeX{} features including mathematical equations, tables, and formatting. You can cite papers using the citation search feature, and they will be automatically added to your \texttt{cite.bib} file.

\section{Mathematical Expressions}
Here are some example equations:

\subsection{Quadratic Formula}
The quadratic formula is given by:
\[x = \frac{-b \pm \sqrt{b^2 - 4ac}}{2a}\]

\subsection{Matrix Example}
A sample matrix:
\[
\begin{pmatrix}
a_{11} & a_{12} & a_{13} \\
a_{21} & a_{22} & a_{23} \\
a_{31} & a_{32} & a_{33}
\end{pmatrix}
\]

\section{Tables}
Below is a sample table showing data:

\begin{table}[h]
\centering
\begin{tabular}{@{}lrrr@{}}
\toprule
\textbf{Category} & \textbf{Value 1} & \textbf{Value 2} & \textbf{Total} \\
\midrule
Item A & 10.5 & 20.3 & 30.8 \\
Item B & 15.2 & 25.1 & 40.3 \\
Item C & 22.1 & 12.7 & 34.8 \\
\bottomrule
\end{tabular}
\caption{Sample Data Table}
\label{tab:sample}
\end{table}

\section{Lists}
\subsection{Bullet Points}
\begin{itemize}
  \item First item
  \item Second item
  \item Third item with sub-items:
    \begin{itemize}
      \item Sub-item 1
      \item Sub-item 2
    \end{itemize}
\end{itemize}

\subsection{Numbered List}
\begin{enumerate}
  \item First step
  \item Second step
  \item Third step
\end{enumerate}

\section{Citations and References}
When you use the citation search feature, references will be automatically managed. For example, you might cite a paper like this: \cite{example2024}. All citations will be automatically added to your \texttt{cite.bib} file.

\section{Conclusion}
This template demonstrates basic \LaTeX{} features including mathematical formulas, tables, and lists. The bibliography system is now set up to work seamlessly with Bibby AI's citation management. Feel free to modify it for your needs!

\bibliographystyle{plain}
\bibliography{cite}

\end{document}\documentclass{article}
\usepackage[utf8]{inputenc}
\usepackage{booktabs}
\usepackage{amsmath}
\usepackage{natbib}

\title{My Document}
\author{Your Name}
\date{\today}

\begin{document}

\maketitle

\section{Introduction}
This is a sample document demonstrating various \LaTeX{} features including mathematical equations, tables, and formatting. You can cite papers using the citation search feature, and they will be automatically added to your \texttt{cite.bib} file.

\section{Mathematical Expressions}
Here are some example equations:

\subsection{Quadratic Formula}
The quadratic formula is given by:
\[x = \frac{-b \pm \sqrt{b^2 - 4ac}}{2a}\]

\subsection{Matrix Example}
A sample matrix:
\[
\begin{pmatrix}
a_{11} & a_{12} & a_{13} \\
a_{21} & a_{22} & a_{23} \\
a_{31} & a_{32} & a_{33}
\end{pmatrix}
\]

\section{Tables}
Below is a sample table showing data:

\begin{table}[h]
\centering
\begin{tabular}{@{}lrrr@{}}
\toprule
\textbf{Category} & \textbf{Value 1} & \textbf{Value 2} & \textbf{Total} \\
\midrule
Item A & 10.5 & 20.3 & 30.8 \\
Item B & 15.2 & 25.1 & 40.3 \\
Item C & 22.1 & 12.7 & 34.8 \\
\bottomrule
\end{tabular}
\caption{Sample Data Table}
\label{tab:sample}
\end{table}

\section{Lists}
\subsection{Bullet Points}
\begin{itemize}
  \item First item
  \item Second item
  \item Third item with sub-items:
    \begin{itemize}
      \item Sub-item 1
      \item Sub-item 2
    \end{itemize}
\end{itemize}

\subsection{Numbered List}
\begin{enumerate}
  \item First step
  \item Second step
  \item Third step
\end{enumerate}

\section{Citations and References}
When you use the citation search feature, references will be automatically managed. For example, you might cite a paper like this: \cite{example2024}. All citations will be automatically added to your \texttt{cite.bib} file.

\section{Conclusion}
This template demonstrates basic \LaTeX{} features including mathematical formulas, tables, and lists. The bibliography system is now set up to work seamlessly with Bibby AI's citation management. Feel free to modify it for your needs!

\bibliographystyle{plain}
\bibliography{cite}

\end{document}\documentclass{article}
\usepackage[utf8]{inputenc}
\usepackage{booktabs}
\usepackage{amsmath}
\usepackage{natbib}

\title{My Document}
\author{Your Name}
\date{\today}

\begin{document}

\maketitle

\section{Introduction}
This is a sample document demonstrating various \LaTeX{} features including mathematical equations, tables, and formatting. You can cite papers using the citation search feature, and they will be automatically added to your \texttt{cite.bib} file.

\section{Mathematical Expressions}
Here are some example equations:

\subsection{Quadratic Formula}
The quadratic formula is given by:
\[x = \frac{-b \pm \sqrt{b^2 - 4ac}}{2a}\]

\subsection{Matrix Example}
A sample matrix:
\[
\begin{pmatrix}
a_{11} & a_{12} & a_{13} \\
a_{21} & a_{22} & a_{23} \\
a_{31} & a_{32} & a_{33}
\end{pmatrix}
\]

\section{Tables}
Below is a sample table showing data:

\begin{table}[h]
\centering
\begin{tabular}{@{}lrrr@{}}
\toprule
\textbf{Category} & \textbf{Value 1} & \textbf{Value 2} & \textbf{Total} \\
\midrule
Item A & 10.5 & 20.3 & 30.8 \\
Item B & 15.2 & 25.1 & 40.3 \\
Item C & 22.1 & 12.7 & 34.8 \\
\bottomrule
\end{tabular}
\caption{Sample Data Table}
\label{tab:sample}
\end{table}

\section{Lists}
\subsection{Bullet Points}
\begin{itemize}
  \item First item
  \item Second item
  \item Third item with sub-items:
    \begin{itemize}
      \item Sub-item 1
      \item Sub-item 2
    \end{itemize}
\end{itemize}

\subsection{Numbered List}
\begin{enumerate}
  \item First step
  \item Second step
  \item Third step
\end{enumerate}

\section{Citations and References}
When you use the citation search feature, references will be automatically managed. For example, you might cite a paper like this: \cite{example2024}. All citations will be automatically added to your \texttt{cite.bib} file.

\section{Conclusion}
This template demonstrates basic \LaTeX{} features including mathematical formulas, tables, and lists. The bibliography system is now set up to work seamlessly with Bibby AI's citation management. Feel free to modify it for your needs!

\bibliographystyle{plain}
\bibliography{cite}

\end{document}\documentclass{article}
\usepackage[utf8]{inputenc}
\usepackage{booktabs}
\usepackage{amsmath}
\usepackage{natbib}

\title{My Document}
\author{Your Name}
\date{\today}

\begin{document}

\maketitle

\section{Introduction}
This is a sample document demonstrating various \LaTeX{} features including mathematical equations, tables, and formatting. You can cite papers using the citation search feature, and they will be automatically added to your \texttt{cite.bib} file.

\section{Mathematical Expressions}
Here are some example equations:

\subsection{Quadratic Formula}
The quadratic formula is given by:
\[x = \frac{-b \pm \sqrt{b^2 - 4ac}}{2a}\]

\subsection{Matrix Example}
A sample matrix:
\[
\begin{pmatrix}
a_{11} & a_{12} & a_{13} \\
a_{21} & a_{22} & a_{23} \\
a_{31} & a_{32} & a_{33}
\end{pmatrix}
\]

\section{Tables}
Below is a sample table showing data:

\begin{table}[h]
\centering
\begin{tabular}{@{}lrrr@{}}
\toprule
\textbf{Category} & \textbf{Value 1} & \textbf{Value 2} & \textbf{Total} \\
\midrule
Item A & 10.5 & 20.3 & 30.8 \\
Item B & 15.2 & 25.1 & 40.3 \\
Item C & 22.1 & 12.7 & 34.8 \\
\bottomrule
\end{tabular}
\caption{Sample Data Table}
\label{tab:sample}
\end{table}

\section{Lists}
\subsection{Bullet Points}
\begin{itemize}
  \item First item
  \item Second item
  \item Third item with sub-items:
    \begin{itemize}
      \item Sub-item 1
      \item Sub-item 2
    \end{itemize}
\end{itemize}

\subsection{Numbered List}
\begin{enumerate}
  \item First step
  \item Second step
  \item Third step
\end{enumerate}

\section{Citations and References}
When you use the citation search feature, references will be automatically managed. For example, you might cite a paper like this: \cite{example2024}. All citations will be automatically added to your \texttt{cite.bib} file.

\section{Conclusion}
This template demonstrates basic \LaTeX{} features including mathematical formulas, tables, and lists. The bibliography system is now set up to work seamlessly with Bibby AI's citation management. Feel free to modify it for your needs!

\bibliographystyle{plain}
\bibliography{cite}

\end{document}\documentclass{article}
\usepackage[utf8]{inputenc}
\usepackage{booktabs}
\usepackage{amsmath}
\usepackage{natbib}

\title{My Document}
\author{Your Name}
\date{\today}

\begin{document}

\maketitle

\section{Introduction}
This is a sample document demonstrating various \LaTeX{} features including mathematical equations, tables, and formatting. You can cite papers using the citation search feature, and they will be automatically added to your \texttt{cite.bib} file.

\section{Mathematical Expressions}
Here are some example equations:

\subsection{Quadratic Formula}
The quadratic formula is given by:
\[x = \frac{-b \pm \sqrt{b^2 - 4ac}}{2a}\]

\subsection{Matrix Example}
A sample matrix:
\[
\begin{pmatrix}
a_{11} & a_{12} & a_{13} \\
a_{21} & a_{22} & a_{23} \\
a_{31} & a_{32} & a_{33}
\end{pmatrix}
\]

\section{Tables}
Below is a sample table showing data:

\begin{table}[h]
\centering
\begin{tabular}{@{}lrrr@{}}
\toprule
\textbf{Category} & \textbf{Value 1} & \textbf{Value 2} & \textbf{Total} \\
\midrule
Item A & 10.5 & 20.3 & 30.8 \\
Item B & 15.2 & 25.1 & 40.3 \\
Item C & 22.1 & 12.7 & 34.8 \\
\bottomrule
\end{tabular}
\caption{Sample Data Table}
\label{tab:sample}
\end{table}

\section{Lists}
\subsection{Bullet Points}
\begin{itemize}
  \item First item
  \item Second item
  \item Third item with sub-items:
    \begin{itemize}
      \item Sub-item 1
      \item Sub-item 2
    \end{itemize}
\end{itemize}

\subsection{Numbered List}
\begin{enumerate}
  \item First step
  \item Second step
  \item Third step
\end{enumerate}

\section{Citations and References}
When you use the citation search feature, references will be automatically managed. For example, you might cite a paper like this: \cite{example2024}. All citations will be automatically added to your \texttt{cite.bib} file.

\section{Conclusion}
This template demonstrates basic \LaTeX{} features including mathematical formulas, tables, and lists. The bibliography system is now set up to work seamlessly with Bibby AI's citation management. Feel free to modify it for your needs!

\bibliographystyle{plain}
\bibliography{cite}

\end{document}\documentclass{article}
\usepackage[utf8]{inputenc}
\usepackage{booktabs}
\usepackage{amsmath}
\usepackage{natbib}

\title{My Document}
\author{Your Name}
\date{\today}

\begin{document}

\maketitle

\section{Introduction}
This is a sample document demonstrating various \LaTeX{} features including mathematical equations, tables, and formatting. You can cite papers using the citation search feature, and they will be automatically added to your \texttt{cite.bib} file.

\section{Mathematical Expressions}
Here are some example equations:

\subsection{Quadratic Formula}
The quadratic formula is given by:
\[x = \frac{-b \pm \sqrt{b^2 - 4ac}}{2a}\]

\subsection{Matrix Example}
A sample matrix:
\[
\begin{pmatrix}
a_{11} & a_{12} & a_{13} \\
a_{21} & a_{22} & a_{23} \\
a_{31} & a_{32} & a_{33}
\end{pmatrix}
\]

\section{Tables}
Below is a sample table showing data:

\begin{table}[h]
\centering
\begin{tabular}{@{}lrrr@{}}
\toprule
\textbf{Category} & \textbf{Value 1} & \textbf{Value 2} & \textbf{Total} \\
\midrule
Item A & 10.5 & 20.3 & 30.8 \\
Item B & 15.2 & 25.1 & 40.3 \\
Item C & 22.1 & 12.7 & 34.8 \\
\bottomrule
\end{tabular}
\caption{Sample Data Table}
\label{tab:sample}
\end{table}

\section{Lists}
\subsection{Bullet Points}
\begin{itemize}
  \item First item
  \item Second item
  \item Third item with sub-items:
    \begin{itemize}
      \item Sub-item 1
      \item Sub-item 2
    \end{itemize}
\end{itemize}

\subsection{Numbered List}
\begin{enumerate}
  \item First step
  \item Second step
  \item Third step
\end{enumerate}

\section{Citations and References}
When you use the citation search feature, references will be automatically managed. For example, you might cite a paper like this: \cite{example2024}. All citations will be automatically added to your \texttt{cite.bib} file.

\section{Conclusion}
This template demonstrates basic \LaTeX{} features including mathematical formulas, tables, and lists. The bibliography system is now set up to work seamlessly with Bibby AI's citation management. Feel free to modify it for your needs!

\bibliographystyle{plain}
\bibliography{cite}

\end{document}\documentclass{article}
\usepackage[utf8]{inputenc}
\usepackage{booktabs}
\usepackage{amsmath}
\usepackage{natbib}

\title{My Document}
\author{Your Name}
\date{\today}

\begin{document}

\maketitle

\section{Introduction}
This is a sample document demonstrating various \LaTeX{} features including mathematical equations, tables, and formatting. You can cite papers using the citation search feature, and they will be automatically added to your \texttt{cite.bib} file.

\section{Mathematical Expressions}
Here are some example equations:

\subsection{Quadratic Formula}
The quadratic formula is given by:
\[x = \frac{-b \pm \sqrt{b^2 - 4ac}}{2a}\]

\subsection{Matrix Example}
A sample matrix:
\[
\begin{pmatrix}
a_{11} & a_{12} & a_{13} \\
a_{21} & a_{22} & a_{23} \\
a_{31} & a_{32} & a_{33}
\end{pmatrix}
\]

\section{Tables}
Below is a sample table showing data:

\begin{table}[h]
\centering
\begin{tabular}{@{}lrrr@{}}
\toprule
\textbf{Category} & \textbf{Value 1} & \textbf{Value 2} & \textbf{Total} \\
\midrule
Item A & 10.5 & 20.3 & 30.8 \\
Item B & 15.2 & 25.1 & 40.3 \\
Item C & 22.1 & 12.7 & 34.8 \\
\bottomrule
\end{tabular}
\caption{Sample Data Table}
\label{tab:sample}
\end{table}

\section{Lists}
\subsection{Bullet Points}
\begin{itemize}
  \item First item
  \item Second item
  \item Third item with sub-items:
    \begin{itemize}
      \item Sub-item 1
      \item Sub-item 2
    \end{itemize}
\end{itemize}

\subsection{Numbered List}
\begin{enumerate}
  \item First step
  \item Second step
  \item Third step
\end{enumerate}

\section{Citations and References}
When you use the citation search feature, references will be automatically managed. For example, you might cite a paper like this: \cite{example2024}. All citations will be automatically added to your \texttt{cite.bib} file.

\section{Conclusion}
This template demonstrates basic \LaTeX{} features including mathematical formulas, tables, and lists. The bibliography system is now set up to work seamlessly with Bibby AI's citation management. Feel free to modify it for your needs!

\bibliographystyle{plain}
\bibliography{cite}

\end{document}\documentclass{article}
\usepackage[utf8]{inputenc}
\usepackage{booktabs}
\usepackage{amsmath}
\usepackage{natbib}

\title{My Document}
\author{Your Name}
\date{\today}

\begin{document}

\maketitle

\section{Introduction}
This is a sample document demonstrating various \LaTeX{} features including mathematical equations, tables, and formatting. You can cite papers using the citation search feature, and they will be automatically added to your \texttt{cite.bib} file.

\section{Mathematical Expressions}
Here are some example equations:

\subsection{Quadratic Formula}
The quadratic formula is given by:
\[x = \frac{-b \pm \sqrt{b^2 - 4ac}}{2a}\]

\subsection{Matrix Example}
A sample matrix:
\[
\begin{pmatrix}
a_{11} & a_{12} & a_{13} \\
a_{21} & a_{22} & a_{23} \\
a_{31} & a_{32} & a_{33}
\end{pmatrix}
\]

\section{Tables}
Below is a sample table showing data:

\begin{table}[h]
\centering
\begin{tabular}{@{}lrrr@{}}
\toprule
\textbf{Category} & \textbf{Value 1} & \textbf{Value 2} & \textbf{Total} \\
\midrule
Item A & 10.5 & 20.3 & 30.8 \\
Item B & 15.2 & 25.1 & 40.3 \\
Item C & 22.1 & 12.7 & 34.8 \\
\bottomrule
\end{tabular}
\caption{Sample Data Table}
\label{tab:sample}
\end{table}

\section{Lists}
\subsection{Bullet Points}
\begin{itemize}
  \item First item
  \item Second item
  \item Third item with sub-items:
    \begin{itemize}
      \item Sub-item 1
      \item Sub-item 2
    \end{itemize}
\end{itemize}

\subsection{Numbered List}
\begin{enumerate}
  \item First step
  \item Second step
  \item Third step
\end{enumerate}

\section{Citations and References}
When you use the citation search feature, references will be automatically managed. For example, you might cite a paper like this: \cite{example2024}. All citations will be automatically added to your \texttt{cite.bib} file.

\section{Conclusion}
This template demonstrates basic \LaTeX{} features including mathematical formulas, tables, and lists. The bibliography system is now set up to work seamlessly with Bibby AI's citation management. Feel free to modify it for your needs!

\bibliographystyle{plain}
\bibliography{cite}

\end{document}\documentclass{article}
\usepackage[utf8]{inputenc}
\usepackage{booktabs}
\usepackage{amsmath}
\usepackage{natbib}

\title{My Document}
\author{Your Name}
\date{\today}

\begin{document}

\maketitle

\section{Introduction}
This is a sample document demonstrating various \LaTeX{} features including mathematical equations, tables, and formatting. You can cite papers using the citation search feature, and they will be automatically added to your \texttt{cite.bib} file.

\section{Mathematical Expressions}
Here are some example equations:

\subsection{Quadratic Formula}
The quadratic formula is given by:
\[x = \frac{-b \pm \sqrt{b^2 - 4ac}}{2a}\]

\subsection{Matrix Example}
A sample matrix:
\[
\begin{pmatrix}
a_{11} & a_{12} & a_{13} \\
a_{21} & a_{22} & a_{23} \\
a_{31} & a_{32} & a_{33}
\end{pmatrix}
\]

\section{Tables}
Below is a sample table showing data:

\begin{table}[h]
\centering
\begin{tabular}{@{}lrrr@{}}
\toprule
\textbf{Category} & \textbf{Value 1} & \textbf{Value 2} & \textbf{Total} \\
\midrule
Item A & 10.5 & 20.3 & 30.8 \\
Item B & 15.2 & 25.1 & 40.3 \\
Item C & 22.1 & 12.7 & 34.8 \\
\bottomrule
\end{tabular}
\caption{Sample Data Table}
\label{tab:sample}
\end{table}

\section{Lists}
\subsection{Bullet Points}
\begin{itemize}
  \item First item
  \item Second item
  \item Third item with sub-items:
    \begin{itemize}
      \item Sub-item 1
      \item Sub-item 2
    \end{itemize}
\end{itemize}

\subsection{Numbered List}
\begin{enumerate}
  \item First step
  \item Second step
  \item Third step
\end{enumerate}

\section{Citations and References}
When you use the citation search feature, references will be automatically managed. For example, you might cite a paper like this: \cite{example2024}. All citations will be automatically added to your \texttt{cite.bib} file.

\section{Conclusion}
This template demonstrates basic \LaTeX{} features including mathematical formulas, tables, and lists. The bibliography system is now set up to work seamlessly with Bibby AI's citation management. Feel free to modify it for your needs!

\bibliographystyle{plain}
\bibliography{cite}

\end{document}\documentclass{article}
\usepackage[utf8]{inputenc}
\usepackage{booktabs}
\usepackage{amsmath}
\usepackage{natbib}

\title{My Document}
\author{Your Name}
\date{\today}

\begin{document}

\maketitle

\section{Introduction}
This is a sample document demonstrating various \LaTeX{} features including mathematical equations, tables, and formatting. You can cite papers using the citation search feature, and they will be automatically added to your \texttt{cite.bib} file.

\section{Mathematical Expressions}
Here are some example equations:

\subsection{Quadratic Formula}
The quadratic formula is given by:
\[x = \frac{-b \pm \sqrt{b^2 - 4ac}}{2a}\]

\subsection{Matrix Example}
A sample matrix:
\[
\begin{pmatrix}
a_{11} & a_{12} & a_{13} \\
a_{21} & a_{22} & a_{23} \\
a_{31} & a_{32} & a_{33}
\end{pmatrix}
\]

\section{Tables}
Below is a sample table showing data:

\begin{table}[h]
\centering
\begin{tabular}{@{}lrrr@{}}
\toprule
\textbf{Category} & \textbf{Value 1} & \textbf{Value 2} & \textbf{Total} \\
\midrule
Item A & 10.5 & 20.3 & 30.8 \\
Item B & 15.2 & 25.1 & 40.3 \\
Item C & 22.1 & 12.7 & 34.8 \\
\bottomrule
\end{tabular}
\caption{Sample Data Table}
\label{tab:sample}
\end{table}

\section{Lists}
\subsection{Bullet Points}
\begin{itemize}
  \item First item
  \item Second item
  \item Third item with sub-items:
    \begin{itemize}
      \item Sub-item 1
      \item Sub-item 2
    \end{itemize}
\end{itemize}

\subsection{Numbered List}
\begin{enumerate}
  \item First step
  \item Second step
  \item Third step
\end{enumerate}

\section{Citations and References}
When you use the citation search feature, references will be automatically managed. For example, you might cite a paper like this: \cite{example2024}. All citations will be automatically added to your \texttt{cite.bib} file.

\section{Conclusion}
This template demonstrates basic \LaTeX{} features including mathematical formulas, tables, and lists. The bibliography system is now set up to work seamlessly with Bibby AI's citation management. Feel free to modify it for your needs!

\bibliographystyle{plain}
\bibliography{cite}

\end{document}\documentclass{article}
\usepackage[utf8]{inputenc}
\usepackage{booktabs}
\usepackage{amsmath}
\usepackage{natbib}

\title{My Document}
\author{Your Name}
\date{\today}

\begin{document}

\maketitle

\section{Introduction}
This is a sample document demonstrating various \LaTeX{} features including mathematical equations, tables, and formatting. You can cite papers using the citation search feature, and they will be automatically added to your \texttt{cite.bib} file.

\section{Mathematical Expressions}
Here are some example equations:

\subsection{Quadratic Formula}
The quadratic formula is given by:
\[x = \frac{-b \pm \sqrt{b^2 - 4ac}}{2a}\]

\subsection{Matrix Example}
A sample matrix:
\[
\begin{pmatrix}
a_{11} & a_{12} & a_{13} \\
a_{21} & a_{22} & a_{23} \\
a_{31} & a_{32} & a_{33}
\end{pmatrix}
\]

\section{Tables}
Below is a sample table showing data:

\begin{table}[h]
\centering
\begin{tabular}{@{}lrrr@{}}
\toprule
\textbf{Category} & \textbf{Value 1} & \textbf{Value 2} & \textbf{Total} \\
\midrule
Item A & 10.5 & 20.3 & 30.8 \\
Item B & 15.2 & 25.1 & 40.3 \\
Item C & 22.1 & 12.7 & 34.8 \\
\bottomrule
\end{tabular}
\caption{Sample Data Table}
\label{tab:sample}
\end{table}

\section{Lists}
\subsection{Bullet Points}
\begin{itemize}
  \item First item
  \item Second item
  \item Third item with sub-items:
    \begin{itemize}
      \item Sub-item 1
      \item Sub-item 2
    \end{itemize}
\end{itemize}

\subsection{Numbered List}
\begin{enumerate}
  \item First step
  \item Second step
  \item Third step
\end{enumerate}

\section{Citations and References}
When you use the citation search feature, references will be automatically managed. For example, you might cite a paper like this: \cite{example2024}. All citations will be automatically added to your \texttt{cite.bib} file.

\section{Conclusion}
This template demonstrates basic \LaTeX{} features including mathematical formulas, tables, and lists. The bibliography system is now set up to work seamlessly with Bibby AI's citation management. Feel free to modify it for your needs!

\bibliographystyle{plain}
\bibliography{cite}

\end{document}\documentclass{article}
\usepackage[utf8]{inputenc}
\usepackage{booktabs}
\usepackage{amsmath}
\usepackage{natbib}

\title{My Document}
\author{Your Name}
\date{\today}

\begin{document}

\maketitle

\section{Introduction}
This is a sample document demonstrating various \LaTeX{} features including mathematical equations, tables, and formatting. You can cite papers using the citation search feature, and they will be automatically added to your \texttt{cite.bib} file.

\section{Mathematical Expressions}
Here are some example equations:

\subsection{Quadratic Formula}
The quadratic formula is given by:
\[x = \frac{-b \pm \sqrt{b^2 - 4ac}}{2a}\]

\subsection{Matrix Example}
A sample matrix:
\[
\begin{pmatrix}
a_{11} & a_{12} & a_{13} \\
a_{21} & a_{22} & a_{23} \\
a_{31} & a_{32} & a_{33}
\end{pmatrix}
\]

\section{Tables}
Below is a sample table showing data:

\begin{table}[h]
\centering
\begin{tabular}{@{}lrrr@{}}
\toprule
\textbf{Category} & \textbf{Value 1} & \textbf{Value 2} & \textbf{Total} \\
\midrule
Item A & 10.5 & 20.3 & 30.8 \\
Item B & 15.2 & 25.1 & 40.3 \\
Item C & 22.1 & 12.7 & 34.8 \\
\bottomrule
\end{tabular}
\caption{Sample Data Table}
\label{tab:sample}
\end{table}

\section{Lists}
\subsection{Bullet Points}
\begin{itemize}
  \item First item
  \item Second item
  \item Third item with sub-items:
    \begin{itemize}
      \item Sub-item 1
      \item Sub-item 2
    \end{itemize}
\end{itemize}

\subsection{Numbered List}
\begin{enumerate}
  \item First step
  \item Second step
  \item Third step
\end{enumerate}

\section{Citations and References}
When you use the citation search feature, references will be automatically managed. For example, you might cite a paper like this: \cite{example2024}. All citations will be automatically added to your \texttt{cite.bib} file.

\section{Conclusion}
This template demonstrates basic \LaTeX{} features including mathematical formulas, tables, and lists. The bibliography system is now set up to work seamlessly with Bibby AI's citation management. Feel free to modify it for your needs!

\bibliographystyle{plain}
\bibliography{cite}

\end{document}\documentclass{article}
\usepackage[utf8]{inputenc}
\usepackage{booktabs}
\usepackage{amsmath}
\usepackage{natbib}

\title{My Document}
\author{Your Name}
\date{\today}

\begin{document}

\maketitle

\section{Introduction}
This is a sample document demonstrating various \LaTeX{} features including mathematical equations, tables, and formatting. You can cite papers using the citation search feature, and they will be automatically added to your \texttt{cite.bib} file.

\section{Mathematical Expressions}
Here are some example equations:

\subsection{Quadratic Formula}
The quadratic formula is given by:
\[x = \frac{-b \pm \sqrt{b^2 - 4ac}}{2a}\]

\subsection{Matrix Example}
A sample matrix:
\[
\begin{pmatrix}
a_{11} & a_{12} & a_{13} \\
a_{21} & a_{22} & a_{23} \\
a_{31} & a_{32} & a_{33}
\end{pmatrix}
\]

\section{Tables}
Below is a sample table showing data:

\begin{table}[h]
\centering
\begin{tabular}{@{}lrrr@{}}
\toprule
\textbf{Category} & \textbf{Value 1} & \textbf{Value 2} & \textbf{Total} \\
\midrule
Item A & 10.5 & 20.3 & 30.8 \\
Item B & 15.2 & 25.1 & 40.3 \\
Item C & 22.1 & 12.7 & 34.8 \\
\bottomrule
\end{tabular}
\caption{Sample Data Table}
\label{tab:sample}
\end{table}

\section{Lists}
\subsection{Bullet Points}
\begin{itemize}
  \item First item
  \item Second item
  \item Third item with sub-items:
    \begin{itemize}
      \item Sub-item 1
      \item Sub-item 2
    \end{itemize}
\end{itemize}

\subsection{Numbered List}
\begin{enumerate}
  \item First step
  \item Second step
  \item Third step
\end{enumerate}

\section{Citations and References}
When you use the citation search feature, references will be automatically managed. For example, you might cite a paper like this: \cite{example2024}. All citations will be automatically added to your \texttt{cite.bib} file.

\section{Conclusion}
This template demonstrates basic \LaTeX{} features including mathematical formulas, tables, and lists. The bibliography system is now set up to work seamlessly with Bibby AI's citation management. Feel free to modify it for your needs!

\bibliographystyle{plain}
\bibliography{cite}

\end{document}\documentclass{article}
\usepackage[utf8]{inputenc}
\usepackage{booktabs}
\usepackage{amsmath}
\usepackage{natbib}

\title{My Document}
\author{Your Name}
\date{\today}

\begin{document}

\maketitle

\section{Introduction}
This is a sample document demonstrating various \LaTeX{} features including mathematical equations, tables, and formatting. You can cite papers using the citation search feature, and they will be automatically added to your \texttt{cite.bib} file.

\section{Mathematical Expressions}
Here are some example equations:

\subsection{Quadratic Formula}
The quadratic formula is given by:
\[x = \frac{-b \pm \sqrt{b^2 - 4ac}}{2a}\]

\subsection{Matrix Example}
A sample matrix:
\[
\begin{pmatrix}
a_{11} & a_{12} & a_{13} \\
a_{21} & a_{22} & a_{23} \\
a_{31} & a_{32} & a_{33}
\end{pmatrix}
\]

\section{Tables}
Below is a sample table showing data:

\begin{table}[h]
\centering
\begin{tabular}{@{}lrrr@{}}
\toprule
\textbf{Category} & \textbf{Value 1} & \textbf{Value 2} & \textbf{Total} \\
\midrule
Item A & 10.5 & 20.3 & 30.8 \\
Item B & 15.2 & 25.1 & 40.3 \\
Item C & 22.1 & 12.7 & 34.8 \\
\bottomrule
\end{tabular}
\caption{Sample Data Table}
\label{tab:sample}
\end{table}

\section{Lists}
\subsection{Bullet Points}
\begin{itemize}
  \item First item
  \item Second item
  \item Third item with sub-items:
    \begin{itemize}
      \item Sub-item 1
      \item Sub-item 2
    \end{itemize}
\end{itemize}

\subsection{Numbered List}
\begin{enumerate}
  \item First step
  \item Second step
  \item Third step
\end{enumerate}

\section{Citations and References}
When you use the citation search feature, references will be automatically managed. For example, you might cite a paper like this: \cite{example2024}. All citations will be automatically added to your \texttt{cite.bib} file.

\section{Conclusion}
This template demonstrates basic \LaTeX{} features including mathematical formulas, tables, and lists. The bibliography system is now set up to work seamlessly with Bibby AI's citation management. Feel free to modify it for your needs!

\bibliographystyle{plain}
\bibliography{cite}

\end{document}\documentclass{article}
\usepackage[utf8]{inputenc}
\usepackage{booktabs}
\usepackage{amsmath}
\usepackage{natbib}

\title{My Document}
\author{Your Name}
\date{\today}

\begin{document}

\maketitle

\section{Introduction}
This is a sample document demonstrating various \LaTeX{} features including mathematical equations, tables, and formatting. You can cite papers using the citation search feature, and they will be automatically added to your \texttt{cite.bib} file.

\section{Mathematical Expressions}
Here are some example equations:

\subsection{Quadratic Formula}
The quadratic formula is given by:
\[x = \frac{-b \pm \sqrt{b^2 - 4ac}}{2a}\]

\subsection{Matrix Example}
A sample matrix:
\[
\begin{pmatrix}
a_{11} & a_{12} & a_{13} \\
a_{21} & a_{22} & a_{23} \\
a_{31} & a_{32} & a_{33}
\end{pmatrix}
\]

\section{Tables}
Below is a sample table showing data:

\begin{table}[h]
\centering
\begin{tabular}{@{}lrrr@{}}
\toprule
\textbf{Category} & \textbf{Value 1} & \textbf{Value 2} & \textbf{Total} \\
\midrule
Item A & 10.5 & 20.3 & 30.8 \\
Item B & 15.2 & 25.1 & 40.3 \\
Item C & 22.1 & 12.7 & 34.8 \\
\bottomrule
\end{tabular}
\caption{Sample Data Table}
\label{tab:sample}
\end{table}

\section{Lists}
\subsection{Bullet Points}
\begin{itemize}
  \item First item
  \item Second item
  \item Third item with sub-items:
    \begin{itemize}
      \item Sub-item 1
      \item Sub-item 2
    \end{itemize}
\end{itemize}

\subsection{Numbered List}
\begin{enumerate}
  \item First step
  \item Second step
  \item Third step
\end{enumerate}

\section{Citations and References}
When you use the citation search feature, references will be automatically managed. For example, you might cite a paper like this: \cite{example2024}. All citations will be automatically added to your \texttt{cite.bib} file.

\section{Conclusion}
This template demonstrates basic \LaTeX{} features including mathematical formulas, tables, and lists. The bibliography system is now set up to work seamlessly with Bibby AI's citation management. Feel free to modify it for your needs!

\bibliographystyle{plain}
\bibliography{cite}

\end{document}\documentclass{article}
\usepackage[utf8]{inputenc}
\usepackage{booktabs}
\usepackage{amsmath}
\usepackage{natbib}

\title{My Document}
\author{Your Name}
\date{\today}

\begin{document}

\maketitle

\section{Introduction}
This is a sample document demonstrating various \LaTeX{} features including mathematical equations, tables, and formatting. You can cite papers using the citation search feature, and they will be automatically added to your \texttt{cite.bib} file.

\section{Mathematical Expressions}
Here are some example equations:

\subsection{Quadratic Formula}
The quadratic formula is given by:
\[x = \frac{-b \pm \sqrt{b^2 - 4ac}}{2a}\]

\subsection{Matrix Example}
A sample matrix:
\[
\begin{pmatrix}
a_{11} & a_{12} & a_{13} \\
a_{21} & a_{22} & a_{23} \\
a_{31} & a_{32} & a_{33}
\end{pmatrix}
\]

\section{Tables}
Below is a sample table showing data:

\begin{table}[h]
\centering
\begin{tabular}{@{}lrrr@{}}
\toprule
\textbf{Category} & \textbf{Value 1} & \textbf{Value 2} & \textbf{Total} \\
\midrule
Item A & 10.5 & 20.3 & 30.8 \\
Item B & 15.2 & 25.1 & 40.3 \\
Item C & 22.1 & 12.7 & 34.8 \\
\bottomrule
\end{tabular}
\caption{Sample Data Table}
\label{tab:sample}
\end{table}

\section{Lists}
\subsection{Bullet Points}
\begin{itemize}
  \item First item
  \item Second item
  \item Third item with sub-items:
    \begin{itemize}
      \item Sub-item 1
      \item Sub-item 2
    \end{itemize}
\end{itemize}

\subsection{Numbered List}
\begin{enumerate}
  \item First step
  \item Second step
  \item Third step
\end{enumerate}

\section{Citations and References}
When you use the citation search feature, references will be automatically managed. For example, you might cite a paper like this: \cite{example2024}. All citations will be automatically added to your \texttt{cite.bib} file.

\section{Conclusion}
This template demonstrates basic \LaTeX{} features including mathematical formulas, tables, and lists. The bibliography system is now set up to work seamlessly with Bibby AI's citation management. Feel free to modify it for your needs!

\bibliographystyle{plain}
\bibliography{cite}

\end{document}\documentclass{article}
\usepackage[utf8]{inputenc}
\usepackage{booktabs}
\usepackage{amsmath}
\usepackage{natbib}

\title{My Document}
\author{Your Name}
\date{\today}

\begin{document}

\maketitle

\section{Introduction}
This is a sample document demonstrating various \LaTeX{} features including mathematical equations, tables, and formatting. You can cite papers using the citation search feature, and they will be automatically added to your \texttt{cite.bib} file.

\section{Mathematical Expressions}
Here are some example equations:

\subsection{Quadratic Formula}
The quadratic formula is given by:
\[x = \frac{-b \pm \sqrt{b^2 - 4ac}}{2a}\]

\subsection{Matrix Example}
A sample matrix:
\[
\begin{pmatrix}
a_{11} & a_{12} & a_{13} \\
a_{21} & a_{22} & a_{23} \\
a_{31} & a_{32} & a_{33}
\end{pmatrix}
\]

\section{Tables}
Below is a sample table showing data:

\begin{table}[h]
\centering
\begin{tabular}{@{}lrrr@{}}
\toprule
\textbf{Category} & \textbf{Value 1} & \textbf{Value 2} & \textbf{Total} \\
\midrule
Item A & 10.5 & 20.3 & 30.8 \\
Item B & 15.2 & 25.1 & 40.3 \\
Item C & 22.1 & 12.7 & 34.8 \\
\bottomrule
\end{tabular}
\caption{Sample Data Table}
\label{tab:sample}
\end{table}

\section{Lists}
\subsection{Bullet Points}
\begin{itemize}
  \item First item
  \item Second item
  \item Third item with sub-items:
    \begin{itemize}
      \item Sub-item 1
      \item Sub-item 2
    \end{itemize}
\end{itemize}

\subsection{Numbered List}
\begin{enumerate}
  \item First step
  \item Second step
  \item Third step
\end{enumerate}

\section{Citations and References}
When you use the citation search feature, references will be automatically managed. For example, you might cite a paper like this: \cite{example2024}. All citations will be automatically added to your \texttt{cite.bib} file.

\section{Conclusion}
This template demonstrates basic \LaTeX{} features including mathematical formulas, tables, and lists. The bibliography system is now set up to work seamlessly with Bibby AI's citation management. Feel free to modify it for your needs!

\bibliographystyle{plain}
\bibliography{cite}

\end{document}\documentclass{article}
\usepackage[utf8]{inputenc}
\usepackage{booktabs}
\usepackage{amsmath}
\usepackage{natbib}

\title{My Document}
\author{Your Name}
\date{\today}

\begin{document}

\maketitle

\section{Introduction}
This is a sample document demonstrating various \LaTeX{} features including mathematical equations, tables, and formatting. You can cite papers using the citation search feature, and they will be automatically added to your \texttt{cite.bib} file.

\section{Mathematical Expressions}
Here are some example equations:

\subsection{Quadratic Formula}
The quadratic formula is given by:
\[x = \frac{-b \pm \sqrt{b^2 - 4ac}}{2a}\]

\subsection{Matrix Example}
A sample matrix:
\[
\begin{pmatrix}
a_{11} & a_{12} & a_{13} \\
a_{21} & a_{22} & a_{23} \\
a_{31} & a_{32} & a_{33}
\end{pmatrix}
\]

\section{Tables}
Below is a sample table showing data:

\begin{table}[h]
\centering
\begin{tabular}{@{}lrrr@{}}
\toprule
\textbf{Category} & \textbf{Value 1} & \textbf{Value 2} & \textbf{Total} \\
\midrule
Item A & 10.5 & 20.3 & 30.8 \\
Item B & 15.2 & 25.1 & 40.3 \\
Item C & 22.1 & 12.7 & 34.8 \\
\bottomrule
\end{tabular}
\caption{Sample Data Table}
\label{tab:sample}
\end{table}

\section{Lists}
\subsection{Bullet Points}
\begin{itemize}
  \item First item
  \item Second item
  \item Third item with sub-items:
    \begin{itemize}
      \item Sub-item 1
      \item Sub-item 2
    \end{itemize}
\end{itemize}

\subsection{Numbered List}
\begin{enumerate}
  \item First step
  \item Second step
  \item Third step
\end{enumerate}

\section{Citations and References}
When you use the citation search feature, references will be automatically managed. For example, you might cite a paper like this: \cite{example2024}. All citations will be automatically added to your \texttt{cite.bib} file.

\section{Conclusion}
This template demonstrates basic \LaTeX{} features including mathematical formulas, tables, and lists. The bibliography system is now set up to work seamlessly with Bibby AI's citation management. Feel free to modify it for your needs!

\bibliographystyle{plain}
\bibliography{cite}

\end{document}\documentclass{article}
\usepackage[utf8]{inputenc}
\usepackage{booktabs}
\usepackage{amsmath}
\usepackage{natbib}

\title{My Document}
\author{Your Name}
\date{\today}

\begin{document}

\maketitle

\section{Introduction}
This is a sample document demonstrating various \LaTeX{} features including mathematical equations, tables, and formatting. You can cite papers using the citation search feature, and they will be automatically added to your \texttt{cite.bib} file.

\section{Mathematical Expressions}
Here are some example equations:

\subsection{Quadratic Formula}
The quadratic formula is given by:
\[x = \frac{-b \pm \sqrt{b^2 - 4ac}}{2a}\]

\subsection{Matrix Example}
A sample matrix:
\[
\begin{pmatrix}
a_{11} & a_{12} & a_{13} \\
a_{21} & a_{22} & a_{23} \\
a_{31} & a_{32} & a_{33}
\end{pmatrix}
\]

\section{Tables}
Below is a sample table showing data:

\begin{table}[h]
\centering
\begin{tabular}{@{}lrrr@{}}
\toprule
\textbf{Category} & \textbf{Value 1} & \textbf{Value 2} & \textbf{Total} \\
\midrule
Item A & 10.5 & 20.3 & 30.8 \\
Item B & 15.2 & 25.1 & 40.3 \\
Item C & 22.1 & 12.7 & 34.8 \\
\bottomrule
\end{tabular}
\caption{Sample Data Table}
\label{tab:sample}
\end{table}

\section{Lists}
\subsection{Bullet Points}
\begin{itemize}
  \item First item
  \item Second item
  \item Third item with sub-items:
    \begin{itemize}
      \item Sub-item 1
      \item Sub-item 2
    \end{itemize}
\end{itemize}

\subsection{Numbered List}
\begin{enumerate}
  \item First step
  \item Second step
  \item Third step
\end{enumerate}

\section{Citations and References}
When you use the citation search feature, references will be automatically managed. For example, you might cite a paper like this: \cite{example2024}. All citations will be automatically added to your \texttt{cite.bib} file.

\section{Conclusion}
This template demonstrates basic \LaTeX{} features including mathematical formulas, tables, and lists. The bibliography system is now set up to work seamlessly with Bibby AI's citation management. Feel free to modify it for your needs!

\bibliographystyle{plain}
\bibliography{cite}

\end{document}\documentclass{article}
\usepackage[utf8]{inputenc}
\usepackage{booktabs}
\usepackage{amsmath}
\usepackage{natbib}

\title{My Document}
\author{Your Name}
\date{\today}

\begin{document}

\maketitle

\section{Introduction}
This is a sample document demonstrating various \LaTeX{} features including mathematical equations, tables, and formatting. You can cite papers using the citation search feature, and they will be automatically added to your \texttt{cite.bib} file.

\section{Mathematical Expressions}
Here are some example equations:

\subsection{Quadratic Formula}
The quadratic formula is given by:
\[x = \frac{-b \pm \sqrt{b^2 - 4ac}}{2a}\]

\subsection{Matrix Example}
A sample matrix:
\[
\begin{pmatrix}
a_{11} & a_{12} & a_{13} \\
a_{21} & a_{22} & a_{23} \\
a_{31} & a_{32} & a_{33}
\end{pmatrix}
\]

\section{Tables}
Below is a sample table showing data:

\begin{table}[h]
\centering
\begin{tabular}{@{}lrrr@{}}
\toprule
\textbf{Category} & \textbf{Value 1} & \textbf{Value 2} & \textbf{Total} \\
\midrule
Item A & 10.5 & 20.3 & 30.8 \\
Item B & 15.2 & 25.1 & 40.3 \\
Item C & 22.1 & 12.7 & 34.8 \\
\bottomrule
\end{tabular}
\caption{Sample Data Table}
\label{tab:sample}
\end{table}

\section{Lists}
\subsection{Bullet Points}
\begin{itemize}
  \item First item
  \item Second item
  \item Third item with sub-items:
    \begin{itemize}
      \item Sub-item 1
      \item Sub-item 2
    \end{itemize}
\end{itemize}

\subsection{Numbered List}
\begin{enumerate}
  \item First step
  \item Second step
  \item Third step
\end{enumerate}

\section{Citations and References}
When you use the citation search feature, references will be automatically managed. For example, you might cite a paper like this: \cite{example2024}. All citations will be automatically added to your \texttt{cite.bib} file.

\section{Conclusion}
This template demonstrates basic \LaTeX{} features including mathematical formulas, tables, and lists. The bibliography system is now set up to work seamlessly with Bibby AI's citation management. Feel free to modify it for your needs!

\bibliographystyle{plain}
\bibliography{cite}

\end{document}\documentclass{article}
\usepackage[utf8]{inputenc}
\usepackage{booktabs}
\usepackage{amsmath}
\usepackage{natbib}

\title{My Document}
\author{Your Name}
\date{\today}

\begin{document}

\maketitle

\section{Introduction}
This is a sample document demonstrating various \LaTeX{} features including mathematical equations, tables, and formatting. You can cite papers using the citation search feature, and they will be automatically added to your \texttt{cite.bib} file.

\section{Mathematical Expressions}
Here are some example equations:

\subsection{Quadratic Formula}
The quadratic formula is given by:
\[x = \frac{-b \pm \sqrt{b^2 - 4ac}}{2a}\]

\subsection{Matrix Example}
A sample matrix:
\[
\begin{pmatrix}
a_{11} & a_{12} & a_{13} \\
a_{21} & a_{22} & a_{23} \\
a_{31} & a_{32} & a_{33}
\end{pmatrix}
\]

\section{Tables}
Below is a sample table showing data:

\begin{table}[h]
\centering
\begin{tabular}{@{}lrrr@{}}
\toprule
\textbf{Category} & \textbf{Value 1} & \textbf{Value 2} & \textbf{Total} \\
\midrule
Item A & 10.5 & 20.3 & 30.8 \\
Item B & 15.2 & 25.1 & 40.3 \\
Item C & 22.1 & 12.7 & 34.8 \\
\bottomrule
\end{tabular}
\caption{Sample Data Table}
\label{tab:sample}
\end{table}

\section{Lists}
\subsection{Bullet Points}
\begin{itemize}
  \item First item
  \item Second item
  \item Third item with sub-items:
    \begin{itemize}
      \item Sub-item 1
      \item Sub-item 2
    \end{itemize}
\end{itemize}

\subsection{Numbered List}
\begin{enumerate}
  \item First step
  \item Second step
  \item Third step
\end{enumerate}

\section{Citations and References}
When you use the citation search feature, references will be automatically managed. For example, you might cite a paper like this: \cite{example2024}. All citations will be automatically added to your \texttt{cite.bib} file.

\section{Conclusion}
This template demonstrates basic \LaTeX{} features including mathematical formulas, tables, and lists. The bibliography system is now set up to work seamlessly with Bibby AI's citation management. Feel free to modify it for your needs!

\bibliographystyle{plain}
\bibliography{cite}

\end{document}\documentclass{article}
\usepackage[utf8]{inputenc}
\usepackage{booktabs}
\usepackage{amsmath}
\usepackage{natbib}

\title{My Document}
\author{Your Name}
\date{\today}

\begin{document}

\maketitle

\section{Introduction}
This is a sample document demonstrating various \LaTeX{} features including mathematical equations, tables, and formatting. You can cite papers using the citation search feature, and they will be automatically added to your \texttt{cite.bib} file.

\section{Mathematical Expressions}
Here are some example equations:

\subsection{Quadratic Formula}
The quadratic formula is given by:
\[x = \frac{-b \pm \sqrt{b^2 - 4ac}}{2a}\]

\subsection{Matrix Example}
A sample matrix:
\[
\begin{pmatrix}
a_{11} & a_{12} & a_{13} \\
a_{21} & a_{22} & a_{23} \\
a_{31} & a_{32} & a_{33}
\end{pmatrix}
\]

\section{Tables}
Below is a sample table showing data:

\begin{table}[h]
\centering
\begin{tabular}{@{}lrrr@{}}
\toprule
\textbf{Category} & \textbf{Value 1} & \textbf{Value 2} & \textbf{Total} \\
\midrule
Item A & 10.5 & 20.3 & 30.8 \\
Item B & 15.2 & 25.1 & 40.3 \\
Item C & 22.1 & 12.7 & 34.8 \\
\bottomrule
\end{tabular}
\caption{Sample Data Table}
\label{tab:sample}
\end{table}

\section{Lists}
\subsection{Bullet Points}
\begin{itemize}
  \item First item
  \item Second item
  \item Third item with sub-items:
    \begin{itemize}
      \item Sub-item 1
      \item Sub-item 2
    \end{itemize}
\end{itemize}

\subsection{Numbered List}
\begin{enumerate}
  \item First step
  \item Second step
  \item Third step
\end{enumerate}

\section{Citations and References}
When you use the citation search feature, references will be automatically managed. For example, you might cite a paper like this: \cite{example2024}. All citations will be automatically added to your \texttt{cite.bib} file.

\section{Conclusion}
This template demonstrates basic \LaTeX{} features including mathematical formulas, tables, and lists. The bibliography system is now set up to work seamlessly with Bibby AI's citation management. Feel free to modify it for your needs!

\bibliographystyle{plain}
\bibliography{cite}

\end{document}\documentclass{article}
\usepackage[utf8]{inputenc}
\usepackage{booktabs}
\usepackage{amsmath}
\usepackage{natbib}

\title{My Document}
\author{Your Name}
\date{\today}

\begin{document}

\maketitle

\section{Introduction}
This is a sample document demonstrating various \LaTeX{} features including mathematical equations, tables, and formatting. You can cite papers using the citation search feature, and they will be automatically added to your \texttt{cite.bib} file.

\section{Mathematical Expressions}
Here are some example equations:

\subsection{Quadratic Formula}
The quadratic formula is given by:
\[x = \frac{-b \pm \sqrt{b^2 - 4ac}}{2a}\]

\subsection{Matrix Example}
A sample matrix:
\[
\begin{pmatrix}
a_{11} & a_{12} & a_{13} \\
a_{21} & a_{22} & a_{23} \\
a_{31} & a_{32} & a_{33}
\end{pmatrix}
\]

\section{Tables}
Below is a sample table showing data:

\begin{table}[h]
\centering
\begin{tabular}{@{}lrrr@{}}
\toprule
\textbf{Category} & \textbf{Value 1} & \textbf{Value 2} & \textbf{Total} \\
\midrule
Item A & 10.5 & 20.3 & 30.8 \\
Item B & 15.2 & 25.1 & 40.3 \\
Item C & 22.1 & 12.7 & 34.8 \\
\bottomrule
\end{tabular}
\caption{Sample Data Table}
\label{tab:sample}
\end{table}

\section{Lists}
\subsection{Bullet Points}
\begin{itemize}
  \item First item
  \item Second item
  \item Third item with sub-items:
    \begin{itemize}
      \item Sub-item 1
      \item Sub-item 2
    \end{itemize}
\end{itemize}

\subsection{Numbered List}
\begin{enumerate}
  \item First step
  \item Second step
  \item Third step
\end{enumerate}

\section{Citations and References}
When you use the citation search feature, references will be automatically managed. For example, you might cite a paper like this: \cite{example2024}. All citations will be automatically added to your \texttt{cite.bib} file.

\section{Conclusion}
This template demonstrates basic \LaTeX{} features including mathematical formulas, tables, and lists. The bibliography system is now set up to work seamlessly with Bibby AI's citation management. Feel free to modify it for your needs!

\bibliographystyle{plain}
\bibliography{cite}

\end{document}\documentclass{article}
\usepackage[utf8]{inputenc}
\usepackage{booktabs}
\usepackage{amsmath}
\usepackage{natbib}

\title{My Document}
\author{Your Name}
\date{\today}

\begin{document}

\maketitle

\section{Introduction}
This is a sample document demonstrating various \LaTeX{} features including mathematical equations, tables, and formatting. You can cite papers using the citation search feature, and they will be automatically added to your \texttt{cite.bib} file.

\section{Mathematical Expressions}
Here are some example equations:

\subsection{Quadratic Formula}
The quadratic formula is given by:
\[x = \frac{-b \pm \sqrt{b^2 - 4ac}}{2a}\]

\subsection{Matrix Example}
A sample matrix:
\[
\begin{pmatrix}
a_{11} & a_{12} & a_{13} \\
a_{21} & a_{22} & a_{23} \\
a_{31} & a_{32} & a_{33}
\end{pmatrix}
\]

\section{Tables}
Below is a sample table showing data:

\begin{table}[h]
\centering
\begin{tabular}{@{}lrrr@{}}
\toprule
\textbf{Category} & \textbf{Value 1} & \textbf{Value 2} & \textbf{Total} \\
\midrule
Item A & 10.5 & 20.3 & 30.8 \\
Item B & 15.2 & 25.1 & 40.3 \\
Item C & 22.1 & 12.7 & 34.8 \\
\bottomrule
\end{tabular}
\caption{Sample Data Table}
\label{tab:sample}
\end{table}

\section{Lists}
\subsection{Bullet Points}
\begin{itemize}
  \item First item
  \item Second item
  \item Third item with sub-items:
    \begin{itemize}
      \item Sub-item 1
      \item Sub-item 2
    \end{itemize}
\end{itemize}

\subsection{Numbered List}
\begin{enumerate}
  \item First step
  \item Second step
  \item Third step
\end{enumerate}

\section{Citations and References}
When you use the citation search feature, references will be automatically managed. For example, you might cite a paper like this: \cite{example2024}. All citations will be automatically added to your \texttt{cite.bib} file.

\section{Conclusion}
This template demonstrates basic \LaTeX{} features including mathematical formulas, tables, and lists. The bibliography system is now set up to work seamlessly with Bibby AI's citation management. Feel free to modify it for your needs!

\bibliographystyle{plain}
\bibliography{cite}

\end{document}\documentclass{article}
\usepackage[utf8]{inputenc}
\usepackage{booktabs}
\usepackage{amsmath}
\usepackage{natbib}

\title{My Document}
\author{Your Name}
\date{\today}

\begin{document}

\maketitle

\section{Introduction}
This is a sample document demonstrating various \LaTeX{} features including mathematical equations, tables, and formatting. You can cite papers using the citation search feature, and they will be automatically added to your \texttt{cite.bib} file.

\section{Mathematical Expressions}
Here are some example equations:

\subsection{Quadratic Formula}
The quadratic formula is given by:
\[x = \frac{-b \pm \sqrt{b^2 - 4ac}}{2a}\]

\subsection{Matrix Example}
A sample matrix:
\[
\begin{pmatrix}
a_{11} & a_{12} & a_{13} \\
a_{21} & a_{22} & a_{23} \\
a_{31} & a_{32} & a_{33}
\end{pmatrix}
\]

\section{Tables}
Below is a sample table showing data:

\begin{table}[h]
\centering
\begin{tabular}{@{}lrrr@{}}
\toprule
\textbf{Category} & \textbf{Value 1} & \textbf{Value 2} & \textbf{Total} \\
\midrule
Item A & 10.5 & 20.3 & 30.8 \\
Item B & 15.2 & 25.1 & 40.3 \\
Item C & 22.1 & 12.7 & 34.8 \\
\bottomrule
\end{tabular}
\caption{Sample Data Table}
\label{tab:sample}
\end{table}

\section{Lists}
\subsection{Bullet Points}
\begin{itemize}
  \item First item
  \item Second item
  \item Third item with sub-items:
    \begin{itemize}
      \item Sub-item 1
      \item Sub-item 2
    \end{itemize}
\end{itemize}

\subsection{Numbered List}
\begin{enumerate}
  \item First step
  \item Second step
  \item Third step
\end{enumerate}

\section{Citations and References}
When you use the citation search feature, references will be automatically managed. For example, you might cite a paper like this: \cite{example2024}. All citations will be automatically added to your \texttt{cite.bib} file.

\section{Conclusion}
This template demonstrates basic \LaTeX{} features including mathematical formulas, tables, and lists. The bibliography system is now set up to work seamlessly with Bibby AI's citation management. Feel free to modify it for your needs!

\bibliographystyle{plain}
\bibliography{cite}

\end{document}\documentclass{article}
\usepackage[utf8]{inputenc}
\usepackage{booktabs}
\usepackage{amsmath}
\usepackage{natbib}

\title{My Document}
\author{Your Name}
\date{\today}

\begin{document}

\maketitle

\section{Introduction}
This is a sample document demonstrating various \LaTeX{} features including mathematical equations, tables, and formatting. You can cite papers using the citation search feature, and they will be automatically added to your \texttt{cite.bib} file.

\section{Mathematical Expressions}
Here are some example equations:

\subsection{Quadratic Formula}
The quadratic formula is given by:
\[x = \frac{-b \pm \sqrt{b^2 - 4ac}}{2a}\]

\subsection{Matrix Example}
A sample matrix:
\[
\begin{pmatrix}
a_{11} & a_{12} & a_{13} \\
a_{21} & a_{22} & a_{23} \\
a_{31} & a_{32} & a_{33}
\end{pmatrix}
\]

\section{Tables}
Below is a sample table showing data:

\begin{table}[h]
\centering
\begin{tabular}{@{}lrrr@{}}
\toprule
\textbf{Category} & \textbf{Value 1} & \textbf{Value 2} & \textbf{Total} \\
\midrule
Item A & 10.5 & 20.3 & 30.8 \\
Item B & 15.2 & 25.1 & 40.3 \\
Item C & 22.1 & 12.7 & 34.8 \\
\bottomrule
\end{tabular}
\caption{Sample Data Table}
\label{tab:sample}
\end{table}

\section{Lists}
\subsection{Bullet Points}
\begin{itemize}
  \item First item
  \item Second item
  \item Third item with sub-items:
    \begin{itemize}
      \item Sub-item 1
      \item Sub-item 2
    \end{itemize}
\end{itemize}

\subsection{Numbered List}
\begin{enumerate}
  \item First step
  \item Second step
  \item Third step
\end{enumerate}

\section{Citations and References}
When you use the citation search feature, references will be automatically managed. For example, you might cite a paper like this: \cite{example2024}. All citations will be automatically added to your \texttt{cite.bib} file.

\section{Conclusion}
This template demonstrates basic \LaTeX{} features including mathematical formulas, tables, and lists. The bibliography system is now set up to work seamlessly with Bibby AI's citation management. Feel free to modify it for your needs!

\bibliographystyle{plain}
\bibliography{cite}

\end{document}\documentclass{article}
\usepackage[utf8]{inputenc}
\usepackage{booktabs}
\usepackage{amsmath}
\usepackage{natbib}

\title{My Document}
\author{Your Name}
\date{\today}

\begin{document}

\maketitle

\section{Introduction}
This is a sample document demonstrating various \LaTeX{} features including mathematical equations, tables, and formatting. You can cite papers using the citation search feature, and they will be automatically added to your \texttt{cite.bib} file.

\section{Mathematical Expressions}
Here are some example equations:

\subsection{Quadratic Formula}
The quadratic formula is given by:
\[x = \frac{-b \pm \sqrt{b^2 - 4ac}}{2a}\]

\subsection{Matrix Example}
A sample matrix:
\[
\begin{pmatrix}
a_{11} & a_{12} & a_{13} \\
a_{21} & a_{22} & a_{23} \\
a_{31} & a_{32} & a_{33}
\end{pmatrix}
\]

\section{Tables}
Below is a sample table showing data:

\begin{table}[h]
\centering
\begin{tabular}{@{}lrrr@{}}
\toprule
\textbf{Category} & \textbf{Value 1} & \textbf{Value 2} & \textbf{Total} \\
\midrule
Item A & 10.5 & 20.3 & 30.8 \\
Item B & 15.2 & 25.1 & 40.3 \\
Item C & 22.1 & 12.7 & 34.8 \\
\bottomrule
\end{tabular}
\caption{Sample Data Table}
\label{tab:sample}
\end{table}

\section{Lists}
\subsection{Bullet Points}
\begin{itemize}
  \item First item
  \item Second item
  \item Third item with sub-items:
    \begin{itemize}
      \item Sub-item 1
      \item Sub-item 2
    \end{itemize}
\end{itemize}

\subsection{Numbered List}
\begin{enumerate}
  \item First step
  \item Second step
  \item Third step
\end{enumerate}

\section{Citations and References}
When you use the citation search feature, references will be automatically managed. For example, you might cite a paper like this: \cite{example2024}. All citations will be automatically added to your \texttt{cite.bib} file.

\section{Conclusion}
This template demonstrates basic \LaTeX{} features including mathematical formulas, tables, and lists. The bibliography system is now set up to work seamlessly with Bibby AI's citation management. Feel free to modify it for your needs!

\bibliographystyle{plain}
\bibliography{cite}

\end{document}\documentclass{article}
\usepackage[utf8]{inputenc}
\usepackage{booktabs}
\usepackage{amsmath}
\usepackage{natbib}

\title{My Document}
\author{Your Name}
\date{\today}

\begin{document}

\maketitle

\section{Introduction}
This is a sample document demonstrating various \LaTeX{} features including mathematical equations, tables, and formatting. You can cite papers using the citation search feature, and they will be automatically added to your \texttt{cite.bib} file.

\section{Mathematical Expressions}
Here are some example equations:

\subsection{Quadratic Formula}
The quadratic formula is given by:
\[x = \frac{-b \pm \sqrt{b^2 - 4ac}}{2a}\]

\subsection{Matrix Example}
A sample matrix:
\[
\begin{pmatrix}
a_{11} & a_{12} & a_{13} \\
a_{21} & a_{22} & a_{23} \\
a_{31} & a_{32} & a_{33}
\end{pmatrix}
\]

\section{Tables}
Below is a sample table showing data:

\begin{table}[h]
\centering
\begin{tabular}{@{}lrrr@{}}
\toprule
\textbf{Category} & \textbf{Value 1} & \textbf{Value 2} & \textbf{Total} \\
\midrule
Item A & 10.5 & 20.3 & 30.8 \\
Item B & 15.2 & 25.1 & 40.3 \\
Item C & 22.1 & 12.7 & 34.8 \\
\bottomrule
\end{tabular}
\caption{Sample Data Table}
\label{tab:sample}
\end{table}

\section{Lists}
\subsection{Bullet Points}
\begin{itemize}
  \item First item
  \item Second item
  \item Third item with sub-items:
    \begin{itemize}
      \item Sub-item 1
      \item Sub-item 2
    \end{itemize}
\end{itemize}

\subsection{Numbered List}
\begin{enumerate}
  \item First step
  \item Second step
  \item Third step
\end{enumerate}

\section{Citations and References}
When you use the citation search feature, references will be automatically managed. For example, you might cite a paper like this: \cite{example2024}. All citations will be automatically added to your \texttt{cite.bib} file.

\section{Conclusion}
This template demonstrates basic \LaTeX{} features including mathematical formulas, tables, and lists. The bibliography system is now set up to work seamlessly with Bibby AI's citation management. Feel free to modify it for your needs!

\bibliographystyle{plain}
\bibliography{cite}

\end{document}\documentclass{article}
\usepackage[utf8]{inputenc}
\usepackage{booktabs}
\usepackage{amsmath}
\usepackage{natbib}

\title{My Document}
\author{Your Name}
\date{\today}

\begin{document}

\maketitle

\section{Introduction}
This is a sample document demonstrating various \LaTeX{} features including mathematical equations, tables, and formatting. You can cite papers using the citation search feature, and they will be automatically added to your \texttt{cite.bib} file.

\section{Mathematical Expressions}
Here are some example equations:

\subsection{Quadratic Formula}
The quadratic formula is given by:
\[x = \frac{-b \pm \sqrt{b^2 - 4ac}}{2a}\]

\subsection{Matrix Example}
A sample matrix:
\[
\begin{pmatrix}
a_{11} & a_{12} & a_{13} \\
a_{21} & a_{22} & a_{23} \\
a_{31} & a_{32} & a_{33}
\end{pmatrix}
\]

\section{Tables}
Below is a sample table showing data:

\begin{table}[h]
\centering
\begin{tabular}{@{}lrrr@{}}
\toprule
\textbf{Category} & \textbf{Value 1} & \textbf{Value 2} & \textbf{Total} \\
\midrule
Item A & 10.5 & 20.3 & 30.8 \\
Item B & 15.2 & 25.1 & 40.3 \\
Item C & 22.1 & 12.7 & 34.8 \\
\bottomrule
\end{tabular}
\caption{Sample Data Table}
\label{tab:sample}
\end{table}

\section{Lists}
\subsection{Bullet Points}
\begin{itemize}
  \item First item
  \item Second item
  \item Third item with sub-items:
    \begin{itemize}
      \item Sub-item 1
      \item Sub-item 2
    \end{itemize}
\end{itemize}

\subsection{Numbered List}
\begin{enumerate}
  \item First step
  \item Second step
  \item Third step
\end{enumerate}

\section{Citations and References}
When you use the citation search feature, references will be automatically managed. For example, you might cite a paper like this: \cite{example2024}. All citations will be automatically added to your \texttt{cite.bib} file.

\section{Conclusion}
This template demonstrates basic \LaTeX{} features including mathematical formulas, tables, and lists. The bibliography system is now set up to work seamlessly with Bibby AI's citation management. Feel free to modify it for your needs!

\bibliographystyle{plain}
\bibliography{cite}

\end{document}\documentclass{article}
\usepackage[utf8]{inputenc}
\usepackage{booktabs}
\usepackage{amsmath}
\usepackage{natbib}

\title{My Document}
\author{Your Name}
\date{\today}

\begin{document}

\maketitle

\section{Introduction}
This is a sample document demonstrating various \LaTeX{} features including mathematical equations, tables, and formatting. You can cite papers using the citation search feature, and they will be automatically added to your \texttt{cite.bib} file.

\section{Mathematical Expressions}
Here are some example equations:

\subsection{Quadratic Formula}
The quadratic formula is given by:
\[x = \frac{-b \pm \sqrt{b^2 - 4ac}}{2a}\]

\subsection{Matrix Example}
A sample matrix:
\[
\begin{pmatrix}
a_{11} & a_{12} & a_{13} \\
a_{21} & a_{22} & a_{23} \\
a_{31} & a_{32} & a_{33}
\end{pmatrix}
\]

\section{Tables}
Below is a sample table showing data:

\begin{table}[h]
\centering
\begin{tabular}{@{}lrrr@{}}
\toprule
\textbf{Category} & \textbf{Value 1} & \textbf{Value 2} & \textbf{Total} \\
\midrule
Item A & 10.5 & 20.3 & 30.8 \\
Item B & 15.2 & 25.1 & 40.3 \\
Item C & 22.1 & 12.7 & 34.8 \\
\bottomrule
\end{tabular}
\caption{Sample Data Table}
\label{tab:sample}
\end{table}

\section{Lists}
\subsection{Bullet Points}
\begin{itemize}
  \item First item
  \item Second item
  \item Third item with sub-items:
    \begin{itemize}
      \item Sub-item 1
      \item Sub-item 2
    \end{itemize}
\end{itemize}

\subsection{Numbered List}
\begin{enumerate}
  \item First step
  \item Second step
  \item Third step
\end{enumerate}

\section{Citations and References}
When you use the citation search feature, references will be automatically managed. For example, you might cite a paper like this: \cite{example2024}. All citations will be automatically added to your \texttt{cite.bib} file.

\section{Conclusion}
This template demonstrates basic \LaTeX{} features including mathematical formulas, tables, and lists. The bibliography system is now set up to work seamlessly with Bibby AI's citation management. Feel free to modify it for your needs!

\bibliographystyle{plain}
\bibliography{cite}

\end{document}\documentclass{article}
\usepackage[utf8]{inputenc}
\usepackage{booktabs}
\usepackage{amsmath}
\usepackage{natbib}

\title{My Document}
\author{Your Name}
\date{\today}

\begin{document}

\maketitle

\section{Introduction}
This is a sample document demonstrating various \LaTeX{} features including mathematical equations, tables, and formatting. You can cite papers using the citation search feature, and they will be automatically added to your \texttt{cite.bib} file.

\section{Mathematical Expressions}
Here are some example equations:

\subsection{Quadratic Formula}
The quadratic formula is given by:
\[x = \frac{-b \pm \sqrt{b^2 - 4ac}}{2a}\]

\subsection{Matrix Example}
A sample matrix:
\[
\begin{pmatrix}
a_{11} & a_{12} & a_{13} \\
a_{21} & a_{22} & a_{23} \\
a_{31} & a_{32} & a_{33}
\end{pmatrix}
\]

\section{Tables}
Below is a sample table showing data:

\begin{table}[h]
\centering
\begin{tabular}{@{}lrrr@{}}
\toprule
\textbf{Category} & \textbf{Value 1} & \textbf{Value 2} & \textbf{Total} \\
\midrule
Item A & 10.5 & 20.3 & 30.8 \\
Item B & 15.2 & 25.1 & 40.3 \\
Item C & 22.1 & 12.7 & 34.8 \\
\bottomrule
\end{tabular}
\caption{Sample Data Table}
\label{tab:sample}
\end{table}

\section{Lists}
\subsection{Bullet Points}
\begin{itemize}
  \item First item
  \item Second item
  \item Third item with sub-items:
    \begin{itemize}
      \item Sub-item 1
      \item Sub-item 2
    \end{itemize}
\end{itemize}

\subsection{Numbered List}
\begin{enumerate}
  \item First step
  \item Second step
  \item Third step
\end{enumerate}

\section{Citations and References}
When you use the citation search feature, references will be automatically managed. For example, you might cite a paper like this: \cite{example2024}. All citations will be automatically added to your \texttt{cite.bib} file.

\section{Conclusion}
This template demonstrates basic \LaTeX{} features including mathematical formulas, tables, and lists. The bibliography system is now set up to work seamlessly with Bibby AI's citation management. Feel free to modify it for your needs!

\bibliographystyle{plain}
\bibliography{cite}

\end{document}\documentclass{article}
\usepackage[utf8]{inputenc}
\usepackage{booktabs}
\usepackage{amsmath}
\usepackage{natbib}

\title{My Document}
\author{Your Name}
\date{\today}

\begin{document}

\maketitle

\section{Introduction}
This is a sample document demonstrating various \LaTeX{} features including mathematical equations, tables, and formatting. You can cite papers using the citation search feature, and they will be automatically added to your \texttt{cite.bib} file.

\section{Mathematical Expressions}
Here are some example equations:

\subsection{Quadratic Formula}
The quadratic formula is given by:
\[x = \frac{-b \pm \sqrt{b^2 - 4ac}}{2a}\]

\subsection{Matrix Example}
A sample matrix:
\[
\begin{pmatrix}
a_{11} & a_{12} & a_{13} \\
a_{21} & a_{22} & a_{23} \\
a_{31} & a_{32} & a_{33}
\end{pmatrix}
\]

\section{Tables}
Below is a sample table showing data:

\begin{table}[h]
\centering
\begin{tabular}{@{}lrrr@{}}
\toprule
\textbf{Category} & \textbf{Value 1} & \textbf{Value 2} & \textbf{Total} \\
\midrule
Item A & 10.5 & 20.3 & 30.8 \\
Item B & 15.2 & 25.1 & 40.3 \\
Item C & 22.1 & 12.7 & 34.8 \\
\bottomrule
\end{tabular}
\caption{Sample Data Table}
\label{tab:sample}
\end{table}

\section{Lists}
\subsection{Bullet Points}
\begin{itemize}
  \item First item
  \item Second item
  \item Third item with sub-items:
    \begin{itemize}
      \item Sub-item 1
      \item Sub-item 2
    \end{itemize}
\end{itemize}

\subsection{Numbered List}
\begin{enumerate}
  \item First step
  \item Second step
  \item Third step
\end{enumerate}

\section{Citations and References}
When you use the citation search feature, references will be automatically managed. For example, you might cite a paper like this: \cite{example2024}. All citations will be automatically added to your \texttt{cite.bib} file.

\section{Conclusion}
This template demonstrates basic \LaTeX{} features including mathematical formulas, tables, and lists. The bibliography system is now set up to work seamlessly with Bibby AI's citation management. Feel free to modify it for your needs!

\bibliographystyle{plain}
\bibliography{cite}

\end{document}\documentclass{article}
\usepackage[utf8]{inputenc}
\usepackage{booktabs}
\usepackage{amsmath}
\usepackage{natbib}

\title{My Document}
\author{Your Name}
\date{\today}

\begin{document}

\maketitle

\section{Introduction}
This is a sample document demonstrating various \LaTeX{} features including mathematical equations, tables, and formatting. You can cite papers using the citation search feature, and they will be automatically added to your \texttt{cite.bib} file.

\section{Mathematical Expressions}
Here are some example equations:

\subsection{Quadratic Formula}
The quadratic formula is given by:
\[x = \frac{-b \pm \sqrt{b^2 - 4ac}}{2a}\]

\subsection{Matrix Example}
A sample matrix:
\[
\begin{pmatrix}
a_{11} & a_{12} & a_{13} \\
a_{21} & a_{22} & a_{23} \\
a_{31} & a_{32} & a_{33}
\end{pmatrix}
\]

\section{Tables}
Below is a sample table showing data:

\begin{table}[h]
\centering
\begin{tabular}{@{}lrrr@{}}
\toprule
\textbf{Category} & \textbf{Value 1} & \textbf{Value 2} & \textbf{Total} \\
\midrule
Item A & 10.5 & 20.3 & 30.8 \\
Item B & 15.2 & 25.1 & 40.3 \\
Item C & 22.1 & 12.7 & 34.8 \\
\bottomrule
\end{tabular}
\caption{Sample Data Table}
\label{tab:sample}
\end{table}

\section{Lists}
\subsection{Bullet Points}
\begin{itemize}
  \item First item
  \item Second item
  \item Third item with sub-items:
    \begin{itemize}
      \item Sub-item 1
      \item Sub-item 2
    \end{itemize}
\end{itemize}

\subsection{Numbered List}
\begin{enumerate}
  \item First step
  \item Second step
  \item Third step
\end{enumerate}

\section{Citations and References}
When you use the citation search feature, references will be automatically managed. For example, you might cite a paper like this: \cite{example2024}. All citations will be automatically added to your \texttt{cite.bib} file.

\section{Conclusion}
This template demonstrates basic \LaTeX{} features including mathematical formulas, tables, and lists. The bibliography system is now set up to work seamlessly with Bibby AI's citation management. Feel free to modify it for your needs!

\bibliographystyle{plain}
\bibliography{cite}

\end{document}\documentclass{article}
\usepackage[utf8]{inputenc}
\usepackage{booktabs}
\usepackage{amsmath}
\usepackage{natbib}

\title{My Document}
\author{Your Name}
\date{\today}

\begin{document}

\maketitle

\section{Introduction}
This is a sample document demonstrating various \LaTeX{} features including mathematical equations, tables, and formatting. You can cite papers using the citation search feature, and they will be automatically added to your \texttt{cite.bib} file.

\section{Mathematical Expressions}
Here are some example equations:

\subsection{Quadratic Formula}
The quadratic formula is given by:
\[x = \frac{-b \pm \sqrt{b^2 - 4ac}}{2a}\]

\subsection{Matrix Example}
A sample matrix:
\[
\begin{pmatrix}
a_{11} & a_{12} & a_{13} \\
a_{21} & a_{22} & a_{23} \\
a_{31} & a_{32} & a_{33}
\end{pmatrix}
\]

\section{Tables}
Below is a sample table showing data:

\begin{table}[h]
\centering
\begin{tabular}{@{}lrrr@{}}
\toprule
\textbf{Category} & \textbf{Value 1} & \textbf{Value 2} & \textbf{Total} \\
\midrule
Item A & 10.5 & 20.3 & 30.8 \\
Item B & 15.2 & 25.1 & 40.3 \\
Item C & 22.1 & 12.7 & 34.8 \\
\bottomrule
\end{tabular}
\caption{Sample Data Table}
\label{tab:sample}
\end{table}

\section{Lists}
\subsection{Bullet Points}
\begin{itemize}
  \item First item
  \item Second item
  \item Third item with sub-items:
    \begin{itemize}
      \item Sub-item 1
      \item Sub-item 2
    \end{itemize}
\end{itemize}

\subsection{Numbered List}
\begin{enumerate}
  \item First step
  \item Second step
  \item Third step
\end{enumerate}

\section{Citations and References}
When you use the citation search feature, references will be automatically managed. For example, you might cite a paper like this: \cite{example2024}. All citations will be automatically added to your \texttt{cite.bib} file.

\section{Conclusion}
This template demonstrates basic \LaTeX{} features including mathematical formulas, tables, and lists. The bibliography system is now set up to work seamlessly with Bibby AI's citation management. Feel free to modify it for your needs!

\bibliographystyle{plain}
\bibliography{cite}

\end{document}\documentclass{article}
\usepackage[utf8]{inputenc}
\usepackage{booktabs}
\usepackage{amsmath}
\usepackage{natbib}

\title{My Document}
\author{Your Name}
\date{\today}

\begin{document}

\maketitle

\section{Introduction}
This is a sample document demonstrating various \LaTeX{} features including mathematical equations, tables, and formatting. You can cite papers using the citation search feature, and they will be automatically added to your \texttt{cite.bib} file.

\section{Mathematical Expressions}
Here are some example equations:

\subsection{Quadratic Formula}
The quadratic formula is given by:
\[x = \frac{-b \pm \sqrt{b^2 - 4ac}}{2a}\]

\subsection{Matrix Example}
A sample matrix:
\[
\begin{pmatrix}
a_{11} & a_{12} & a_{13} \\
a_{21} & a_{22} & a_{23} \\
a_{31} & a_{32} & a_{33}
\end{pmatrix}
\]

\section{Tables}
Below is a sample table showing data:

\begin{table}[h]
\centering
\begin{tabular}{@{}lrrr@{}}
\toprule
\textbf{Category} & \textbf{Value 1} & \textbf{Value 2} & \textbf{Total} \\
\midrule
Item A & 10.5 & 20.3 & 30.8 \\
Item B & 15.2 & 25.1 & 40.3 \\
Item C & 22.1 & 12.7 & 34.8 \\
\bottomrule
\end{tabular}
\caption{Sample Data Table}
\label{tab:sample}
\end{table}

\section{Lists}
\subsection{Bullet Points}
\begin{itemize}
  \item First item
  \item Second item
  \item Third item with sub-items:
    \begin{itemize}
      \item Sub-item 1
      \item Sub-item 2
    \end{itemize}
\end{itemize}

\subsection{Numbered List}
\begin{enumerate}
  \item First step
  \item Second step
  \item Third step
\end{enumerate}

\section{Citations and References}
When you use the citation search feature, references will be automatically managed. For example, you might cite a paper like this: \cite{example2024}. All citations will be automatically added to your \texttt{cite.bib} file.

\section{Conclusion}
This template demonstrates basic \LaTeX{} features including mathematical formulas, tables, and lists. The bibliography system is now set up to work seamlessly with Bibby AI's citation management. Feel free to modify it for your needs!

\bibliographystyle{plain}
\bibliography{cite}

\end{document}\documentclass{article}
\usepackage[utf8]{inputenc}
\usepackage{booktabs}
\usepackage{amsmath}
\usepackage{natbib}

\title{My Document}
\author{Your Name}
\date{\today}

\begin{document}

\maketitle

\section{Introduction}
This is a sample document demonstrating various \LaTeX{} features including mathematical equations, tables, and formatting. You can cite papers using the citation search feature, and they will be automatically added to your \texttt{cite.bib} file.

\section{Mathematical Expressions}
Here are some example equations:

\subsection{Quadratic Formula}
The quadratic formula is given by:
\[x = \frac{-b \pm \sqrt{b^2 - 4ac}}{2a}\]

\subsection{Matrix Example}
A sample matrix:
\[
\begin{pmatrix}
a_{11} & a_{12} & a_{13} \\
a_{21} & a_{22} & a_{23} \\
a_{31} & a_{32} & a_{33}
\end{pmatrix}
\]

\section{Tables}
Below is a sample table showing data:

\begin{table}[h]
\centering
\begin{tabular}{@{}lrrr@{}}
\toprule
\textbf{Category} & \textbf{Value 1} & \textbf{Value 2} & \textbf{Total} \\
\midrule
Item A & 10.5 & 20.3 & 30.8 \\
Item B & 15.2 & 25.1 & 40.3 \\
Item C & 22.1 & 12.7 & 34.8 \\
\bottomrule
\end{tabular}
\caption{Sample Data Table}
\label{tab:sample}
\end{table}

\section{Lists}
\subsection{Bullet Points}
\begin{itemize}
  \item First item
  \item Second item
  \item Third item with sub-items:
    \begin{itemize}
      \item Sub-item 1
      \item Sub-item 2
    \end{itemize}
\end{itemize}

\subsection{Numbered List}
\begin{enumerate}
  \item First step
  \item Second step
  \item Third step
\end{enumerate}

\section{Citations and References}
When you use the citation search feature, references will be automatically managed. For example, you might cite a paper like this: \cite{example2024}. All citations will be automatically added to your \texttt{cite.bib} file.

\section{Conclusion}
This template demonstrates basic \LaTeX{} features including mathematical formulas, tables, and lists. The bibliography system is now set up to work seamlessly with Bibby AI's citation management. Feel free to modify it for your needs!

\bibliographystyle{plain}
\bibliography{cite}

\end{document}\documentclass{article}
\usepackage[utf8]{inputenc}
\usepackage{booktabs}
\usepackage{amsmath}
\usepackage{natbib}

\title{My Document}
\author{Your Name}
\date{\today}

\begin{document}

\maketitle

\section{Introduction}
This is a sample document demonstrating various \LaTeX{} features including mathematical equations, tables, and formatting. You can cite papers using the citation search feature, and they will be automatically added to your \texttt{cite.bib} file.

\section{Mathematical Expressions}
Here are some example equations:

\subsection{Quadratic Formula}
The quadratic formula is given by:
\[x = \frac{-b \pm \sqrt{b^2 - 4ac}}{2a}\]

\subsection{Matrix Example}
A sample matrix:
\[
\begin{pmatrix}
a_{11} & a_{12} & a_{13} \\
a_{21} & a_{22} & a_{23} \\
a_{31} & a_{32} & a_{33}
\end{pmatrix}
\]

\section{Tables}
Below is a sample table showing data:

\begin{table}[h]
\centering
\begin{tabular}{@{}lrrr@{}}
\toprule
\textbf{Category} & \textbf{Value 1} & \textbf{Value 2} & \textbf{Total} \\
\midrule
Item A & 10.5 & 20.3 & 30.8 \\
Item B & 15.2 & 25.1 & 40.3 \\
Item C & 22.1 & 12.7 & 34.8 \\
\bottomrule
\end{tabular}
\caption{Sample Data Table}
\label{tab:sample}
\end{table}

\section{Lists}
\subsection{Bullet Points}
\begin{itemize}
  \item First item
  \item Second item
  \item Third item with sub-items:
    \begin{itemize}
      \item Sub-item 1
      \item Sub-item 2
    \end{itemize}
\end{itemize}

\subsection{Numbered List}
\begin{enumerate}
  \item First step
  \item Second step
  \item Third step
\end{enumerate}

\section{Citations and References}
When you use the citation search feature, references will be automatically managed. For example, you might cite a paper like this: \cite{example2024}. All citations will be automatically added to your \texttt{cite.bib} file.

\section{Conclusion}
This template demonstrates basic \LaTeX{} features including mathematical formulas, tables, and lists. The bibliography system is now set up to work seamlessly with Bibby AI's citation management. Feel free to modify it for your needs!

\bibliographystyle{plain}
\bibliography{cite}

\end{document}\documentclass{article}
\usepackage[utf8]{inputenc}
\usepackage{booktabs}
\usepackage{amsmath}
\usepackage{natbib}

\title{My Document}
\author{Your Name}
\date{\today}

\begin{document}

\maketitle

\section{Introduction}
This is a sample document demonstrating various \LaTeX{} features including mathematical equations, tables, and formatting. You can cite papers using the citation search feature, and they will be automatically added to your \texttt{cite.bib} file.

\section{Mathematical Expressions}
Here are some example equations:

\subsection{Quadratic Formula}
The quadratic formula is given by:
\[x = \frac{-b \pm \sqrt{b^2 - 4ac}}{2a}\]

\subsection{Matrix Example}
A sample matrix:
\[
\begin{pmatrix}
a_{11} & a_{12} & a_{13} \\
a_{21} & a_{22} & a_{23} \\
a_{31} & a_{32} & a_{33}
\end{pmatrix}
\]

\section{Tables}
Below is a sample table showing data:

\begin{table}[h]
\centering
\begin{tabular}{@{}lrrr@{}}
\toprule
\textbf{Category} & \textbf{Value 1} & \textbf{Value 2} & \textbf{Total} \\
\midrule
Item A & 10.5 & 20.3 & 30.8 \\
Item B & 15.2 & 25.1 & 40.3 \\
Item C & 22.1 & 12.7 & 34.8 \\
\bottomrule
\end{tabular}
\caption{Sample Data Table}
\label{tab:sample}
\end{table}

\section{Lists}
\subsection{Bullet Points}
\begin{itemize}
  \item First item
  \item Second item
  \item Third item with sub-items:
    \begin{itemize}
      \item Sub-item 1
      \item Sub-item 2
    \end{itemize}
\end{itemize}

\subsection{Numbered List}
\begin{enumerate}
  \item First step
  \item Second step
  \item Third step
\end{enumerate}

\section{Citations and References}
When you use the citation search feature, references will be automatically managed. For example, you might cite a paper like this: \cite{example2024}. All citations will be automatically added to your \texttt{cite.bib} file.

\section{Conclusion}
This template demonstrates basic \LaTeX{} features including mathematical formulas, tables, and lists. The bibliography system is now set up to work seamlessly with Bibby AI's citation management. Feel free to modify it for your needs!

\bibliographystyle{plain}
\bibliography{cite}

\end{document}\documentclass{article}
\usepackage[utf8]{inputenc}
\usepackage{booktabs}
\usepackage{amsmath}
\usepackage{natbib}

\title{My Document}
\author{Your Name}
\date{\today}

\begin{document}

\maketitle

\section{Introduction}
This is a sample document demonstrating various \LaTeX{} features including mathematical equations, tables, and formatting. You can cite papers using the citation search feature, and they will be automatically added to your \texttt{cite.bib} file.

\section{Mathematical Expressions}
Here are some example equations:

\subsection{Quadratic Formula}
The quadratic formula is given by:
\[x = \frac{-b \pm \sqrt{b^2 - 4ac}}{2a}\]

\subsection{Matrix Example}
A sample matrix:
\[
\begin{pmatrix}
a_{11} & a_{12} & a_{13} \\
a_{21} & a_{22} & a_{23} \\
a_{31} & a_{32} & a_{33}
\end{pmatrix}
\]

\section{Tables}
Below is a sample table showing data:

\begin{table}[h]
\centering
\begin{tabular}{@{}lrrr@{}}
\toprule
\textbf{Category} & \textbf{Value 1} & \textbf{Value 2} & \textbf{Total} \\
\midrule
Item A & 10.5 & 20.3 & 30.8 \\
Item B & 15.2 & 25.1 & 40.3 \\
Item C & 22.1 & 12.7 & 34.8 \\
\bottomrule
\end{tabular}
\caption{Sample Data Table}
\label{tab:sample}
\end{table}

\section{Lists}
\subsection{Bullet Points}
\begin{itemize}
  \item First item
  \item Second item
  \item Third item with sub-items:
    \begin{itemize}
      \item Sub-item 1
      \item Sub-item 2
    \end{itemize}
\end{itemize}

\subsection{Numbered List}
\begin{enumerate}
  \item First step
  \item Second step
  \item Third step
\end{enumerate}

\section{Citations and References}
When you use the citation search feature, references will be automatically managed. For example, you might cite a paper like this: \cite{example2024}. All citations will be automatically added to your \texttt{cite.bib} file.

\section{Conclusion}
This template demonstrates basic \LaTeX{} features including mathematical formulas, tables, and lists. The bibliography system is now set up to work seamlessly with Bibby AI's citation management. Feel free to modify it for your needs!

\bibliographystyle{plain}
\bibliography{cite}

\end{document}\documentclass{article}
\usepackage[utf8]{inputenc}
\usepackage{booktabs}
\usepackage{amsmath}
\usepackage{natbib}

\title{My Document}
\author{Your Name}
\date{\today}

\begin{document}

\maketitle

\section{Introduction}
This is a sample document demonstrating various \LaTeX{} features including mathematical equations, tables, and formatting. You can cite papers using the citation search feature, and they will be automatically added to your \texttt{cite.bib} file.

\section{Mathematical Expressions}
Here are some example equations:

\subsection{Quadratic Formula}
The quadratic formula is given by:
\[x = \frac{-b \pm \sqrt{b^2 - 4ac}}{2a}\]

\subsection{Matrix Example}
A sample matrix:
\[
\begin{pmatrix}
a_{11} & a_{12} & a_{13} \\
a_{21} & a_{22} & a_{23} \\
a_{31} & a_{32} & a_{33}
\end{pmatrix}
\]

\section{Tables}
Below is a sample table showing data:

\begin{table}[h]
\centering
\begin{tabular}{@{}lrrr@{}}
\toprule
\textbf{Category} & \textbf{Value 1} & \textbf{Value 2} & \textbf{Total} \\
\midrule
Item A & 10.5 & 20.3 & 30.8 \\
Item B & 15.2 & 25.1 & 40.3 \\
Item C & 22.1 & 12.7 & 34.8 \\
\bottomrule
\end{tabular}
\caption{Sample Data Table}
\label{tab:sample}
\end{table}

\section{Lists}
\subsection{Bullet Points}
\begin{itemize}
  \item First item
  \item Second item
  \item Third item with sub-items:
    \begin{itemize}
      \item Sub-item 1
      \item Sub-item 2
    \end{itemize}
\end{itemize}

\subsection{Numbered List}
\begin{enumerate}
  \item First step
  \item Second step
  \item Third step
\end{enumerate}

\section{Citations and References}
When you use the citation search feature, references will be automatically managed. For example, you might cite a paper like this: \cite{example2024}. All citations will be automatically added to your \texttt{cite.bib} file.

\section{Conclusion}
This template demonstrates basic \LaTeX{} features including mathematical formulas, tables, and lists. The bibliography system is now set up to work seamlessly with Bibby AI's citation management. Feel free to modify it for your needs!

\bibliographystyle{plain}
\bibliography{cite}

\end{document}\documentclass{article}
\usepackage[utf8]{inputenc}
\usepackage{booktabs}
\usepackage{amsmath}
\usepackage{natbib}

\title{My Document}
\author{Your Name}
\date{\today}

\begin{document}

\maketitle

\section{Introduction}
This is a sample document demonstrating various \LaTeX{} features including mathematical equations, tables, and formatting. You can cite papers using the citation search feature, and they will be automatically added to your \texttt{cite.bib} file.

\section{Mathematical Expressions}
Here are some example equations:

\subsection{Quadratic Formula}
The quadratic formula is given by:
\[x = \frac{-b \pm \sqrt{b^2 - 4ac}}{2a}\]

\subsection{Matrix Example}
A sample matrix:
\[
\begin{pmatrix}
a_{11} & a_{12} & a_{13} \\
a_{21} & a_{22} & a_{23} \\
a_{31} & a_{32} & a_{33}
\end{pmatrix}
\]

\section{Tables}
Below is a sample table showing data:

\begin{table}[h]
\centering
\begin{tabular}{@{}lrrr@{}}
\toprule
\textbf{Category} & \textbf{Value 1} & \textbf{Value 2} & \textbf{Total} \\
\midrule
Item A & 10.5 & 20.3 & 30.8 \\
Item B & 15.2 & 25.1 & 40.3 \\
Item C & 22.1 & 12.7 & 34.8 \\
\bottomrule
\end{tabular}
\caption{Sample Data Table}
\label{tab:sample}
\end{table}

\section{Lists}
\subsection{Bullet Points}
\begin{itemize}
  \item First item
  \item Second item
  \item Third item with sub-items:
    \begin{itemize}
      \item Sub-item 1
      \item Sub-item 2
    \end{itemize}
\end{itemize}

\subsection{Numbered List}
\begin{enumerate}
  \item First step
  \item Second step
  \item Third step
\end{enumerate}

\section{Citations and References}
When you use the citation search feature, references will be automatically managed. For example, you might cite a paper like this: \cite{example2024}. All citations will be automatically added to your \texttt{cite.bib} file.

\section{Conclusion}
This template demonstrates basic \LaTeX{} features including mathematical formulas, tables, and lists. The bibliography system is now set up to work seamlessly with Bibby AI's citation management. Feel free to modify it for your needs!

\bibliographystyle{plain}
\bibliography{cite}

\end{document}\documentclass{article}
\usepackage[utf8]{inputenc}
\usepackage{booktabs}
\usepackage{amsmath}
\usepackage{natbib}

\title{My Document}
\author{Your Name}
\date{\today}

\begin{document}

\maketitle

\section{Introduction}
This is a sample document demonstrating various \LaTeX{} features including mathematical equations, tables, and formatting. You can cite papers using the citation search feature, and they will be automatically added to your \texttt{cite.bib} file.

\section{Mathematical Expressions}
Here are some example equations:

\subsection{Quadratic Formula}
The quadratic formula is given by:
\[x = \frac{-b \pm \sqrt{b^2 - 4ac}}{2a}\]

\subsection{Matrix Example}
A sample matrix:
\[
\begin{pmatrix}
a_{11} & a_{12} & a_{13} \\
a_{21} & a_{22} & a_{23} \\
a_{31} & a_{32} & a_{33}
\end{pmatrix}
\]

\section{Tables}
Below is a sample table showing data:

\begin{table}[h]
\centering
\begin{tabular}{@{}lrrr@{}}
\toprule
\textbf{Category} & \textbf{Value 1} & \textbf{Value 2} & \textbf{Total} \\
\midrule
Item A & 10.5 & 20.3 & 30.8 \\
Item B & 15.2 & 25.1 & 40.3 \\
Item C & 22.1 & 12.7 & 34.8 \\
\bottomrule
\end{tabular}
\caption{Sample Data Table}
\label{tab:sample}
\end{table}

\section{Lists}
\subsection{Bullet Points}
\begin{itemize}
  \item First item
  \item Second item
  \item Third item with sub-items:
    \begin{itemize}
      \item Sub-item 1
      \item Sub-item 2
    \end{itemize}
\end{itemize}

\subsection{Numbered List}
\begin{enumerate}
  \item First step
  \item Second step
  \item Third step
\end{enumerate}

\section{Citations and References}
When you use the citation search feature, references will be automatically managed. For example, you might cite a paper like this: \cite{example2024}. All citations will be automatically added to your \texttt{cite.bib} file.

\section{Conclusion}
This template demonstrates basic \LaTeX{} features including mathematical formulas, tables, and lists. The bibliography system is now set up to work seamlessly with Bibby AI's citation management. Feel free to modify it for your needs!

\bibliographystyle{plain}
\bibliography{cite}

\end{document}\documentclass{article}
\usepackage[utf8]{inputenc}
\usepackage{booktabs}
\usepackage{amsmath}
\usepackage{natbib}

\title{My Document}
\author{Your Name}
\date{\today}

\begin{document}

\maketitle

\section{Introduction}
This is a sample document demonstrating various \LaTeX{} features including mathematical equations, tables, and formatting. You can cite papers using the citation search feature, and they will be automatically added to your \texttt{cite.bib} file.

\section{Mathematical Expressions}
Here are some example equations:

\subsection{Quadratic Formula}
The quadratic formula is given by:
\[x = \frac{-b \pm \sqrt{b^2 - 4ac}}{2a}\]

\subsection{Matrix Example}
A sample matrix:
\[
\begin{pmatrix}
a_{11} & a_{12} & a_{13} \\
a_{21} & a_{22} & a_{23} \\
a_{31} & a_{32} & a_{33}
\end{pmatrix}
\]

\section{Tables}
Below is a sample table showing data:

\begin{table}[h]
\centering
\begin{tabular}{@{}lrrr@{}}
\toprule
\textbf{Category} & \textbf{Value 1} & \textbf{Value 2} & \textbf{Total} \\
\midrule
Item A & 10.5 & 20.3 & 30.8 \\
Item B & 15.2 & 25.1 & 40.3 \\
Item C & 22.1 & 12.7 & 34.8 \\
\bottomrule
\end{tabular}
\caption{Sample Data Table}
\label{tab:sample}
\end{table}

\section{Lists}
\subsection{Bullet Points}
\begin{itemize}
  \item First item
  \item Second item
  \item Third item with sub-items:
    \begin{itemize}
      \item Sub-item 1
      \item Sub-item 2
    \end{itemize}
\end{itemize}

\subsection{Numbered List}
\begin{enumerate}
  \item First step
  \item Second step
  \item Third step
\end{enumerate}

\section{Citations and References}
When you use the citation search feature, references will be automatically managed. For example, you might cite a paper like this: \cite{example2024}. All citations will be automatically added to your \texttt{cite.bib} file.

\section{Conclusion}
This template demonstrates basic \LaTeX{} features including mathematical formulas, tables, and lists. The bibliography system is now set up to work seamlessly with Bibby AI's citation management. Feel free to modify it for your needs!

\bibliographystyle{plain}
\bibliography{cite}

\end{document}\documentclass{article}
\usepackage[utf8]{inputenc}
\usepackage{booktabs}
\usepackage{amsmath}
\usepackage{natbib}

\title{My Document}
\author{Your Name}
\date{\today}

\begin{document}

\maketitle

\section{Introduction}
This is a sample document demonstrating various \LaTeX{} features including mathematical equations, tables, and formatting. You can cite papers using the citation search feature, and they will be automatically added to your \texttt{cite.bib} file.

\section{Mathematical Expressions}
Here are some example equations:

\subsection{Quadratic Formula}
The quadratic formula is given by:
\[x = \frac{-b \pm \sqrt{b^2 - 4ac}}{2a}\]

\subsection{Matrix Example}
A sample matrix:
\[
\begin{pmatrix}
a_{11} & a_{12} & a_{13} \\
a_{21} & a_{22} & a_{23} \\
a_{31} & a_{32} & a_{33}
\end{pmatrix}
\]

\section{Tables}
Below is a sample table showing data:

\begin{table}[h]
\centering
\begin{tabular}{@{}lrrr@{}}
\toprule
\textbf{Category} & \textbf{Value 1} & \textbf{Value 2} & \textbf{Total} \\
\midrule
Item A & 10.5 & 20.3 & 30.8 \\
Item B & 15.2 & 25.1 & 40.3 \\
Item C & 22.1 & 12.7 & 34.8 \\
\bottomrule
\end{tabular}
\caption{Sample Data Table}
\label{tab:sample}
\end{table}

\section{Lists}
\subsection{Bullet Points}
\begin{itemize}
  \item First item
  \item Second item
  \item Third item with sub-items:
    \begin{itemize}
      \item Sub-item 1
      \item Sub-item 2
    \end{itemize}
\end{itemize}

\subsection{Numbered List}
\begin{enumerate}
  \item First step
  \item Second step
  \item Third step
\end{enumerate}

\section{Citations and References}
When you use the citation search feature, references will be automatically managed. For example, you might cite a paper like this: \cite{example2024}. All citations will be automatically added to your \texttt{cite.bib} file.

\section{Conclusion}
This template demonstrates basic \LaTeX{} features including mathematical formulas, tables, and lists. The bibliography system is now set up to work seamlessly with Bibby AI's citation management. Feel free to modify it for your needs!

\bibliographystyle{plain}
\bibliography{cite}

\end{document}\documentclass{article}
\usepackage[utf8]{inputenc}
\usepackage{booktabs}
\usepackage{amsmath}
\usepackage{natbib}

\title{My Document}
\author{Your Name}
\date{\today}

\begin{document}

\maketitle

\section{Introduction}
This is a sample document demonstrating various \LaTeX{} features including mathematical equations, tables, and formatting. You can cite papers using the citation search feature, and they will be automatically added to your \texttt{cite.bib} file.

\section{Mathematical Expressions}
Here are some example equations:

\subsection{Quadratic Formula}
The quadratic formula is given by:
\[x = \frac{-b \pm \sqrt{b^2 - 4ac}}{2a}\]

\subsection{Matrix Example}
A sample matrix:
\[
\begin{pmatrix}
a_{11} & a_{12} & a_{13} \\
a_{21} & a_{22} & a_{23} \\
a_{31} & a_{32} & a_{33}
\end{pmatrix}
\]

\section{Tables}
Below is a sample table showing data:

\begin{table}[h]
\centering
\begin{tabular}{@{}lrrr@{}}
\toprule
\textbf{Category} & \textbf{Value 1} & \textbf{Value 2} & \textbf{Total} \\
\midrule
Item A & 10.5 & 20.3 & 30.8 \\
Item B & 15.2 & 25.1 & 40.3 \\
Item C & 22.1 & 12.7 & 34.8 \\
\bottomrule
\end{tabular}
\caption{Sample Data Table}
\label{tab:sample}
\end{table}

\section{Lists}
\subsection{Bullet Points}
\begin{itemize}
  \item First item
  \item Second item
  \item Third item with sub-items:
    \begin{itemize}
      \item Sub-item 1
      \item Sub-item 2
    \end{itemize}
\end{itemize}

\subsection{Numbered List}
\begin{enumerate}
  \item First step
  \item Second step
  \item Third step
\end{enumerate}

\section{Citations and References}
When you use the citation search feature, references will be automatically managed. For example, you might cite a paper like this: \cite{example2024}. All citations will be automatically added to your \texttt{cite.bib} file.

\section{Conclusion}
This template demonstrates basic \LaTeX{} features including mathematical formulas, tables, and lists. The bibliography system is now set up to work seamlessly with Bibby AI's citation management. Feel free to modify it for your needs!

\bibliographystyle{plain}
\bibliography{cite}

\end{document}\documentclass{article}
\usepackage[utf8]{inputenc}
\usepackage{booktabs}
\usepackage{amsmath}
\usepackage{natbib}

\title{My Document}
\author{Your Name}
\date{\today}

\begin{document}

\maketitle

\section{Introduction}
This is a sample document demonstrating various \LaTeX{} features including mathematical equations, tables, and formatting. You can cite papers using the citation search feature, and they will be automatically added to your \texttt{cite.bib} file.

\section{Mathematical Expressions}
Here are some example equations:

\subsection{Quadratic Formula}
The quadratic formula is given by:
\[x = \frac{-b \pm \sqrt{b^2 - 4ac}}{2a}\]

\subsection{Matrix Example}
A sample matrix:
\[
\begin{pmatrix}
a_{11} & a_{12} & a_{13} \\
a_{21} & a_{22} & a_{23} \\
a_{31} & a_{32} & a_{33}
\end{pmatrix}
\]

\section{Tables}
Below is a sample table showing data:

\begin{table}[h]
\centering
\begin{tabular}{@{}lrrr@{}}
\toprule
\textbf{Category} & \textbf{Value 1} & \textbf{Value 2} & \textbf{Total} \\
\midrule
Item A & 10.5 & 20.3 & 30.8 \\
Item B & 15.2 & 25.1 & 40.3 \\
Item C & 22.1 & 12.7 & 34.8 \\
\bottomrule
\end{tabular}
\caption{Sample Data Table}
\label{tab:sample}
\end{table}

\section{Lists}
\subsection{Bullet Points}
\begin{itemize}
  \item First item
  \item Second item
  \item Third item with sub-items:
    \begin{itemize}
      \item Sub-item 1
      \item Sub-item 2
    \end{itemize}
\end{itemize}

\subsection{Numbered List}
\begin{enumerate}
  \item First step
  \item Second step
  \item Third step
\end{enumerate}

\section{Citations and References}
When you use the citation search feature, references will be automatically managed. For example, you might cite a paper like this: \cite{example2024}. All citations will be automatically added to your \texttt{cite.bib} file.

\section{Conclusion}
This template demonstrates basic \LaTeX{} features including mathematical formulas, tables, and lists. The bibliography system is now set up to work seamlessly with Bibby AI's citation management. Feel free to modify it for your needs!

\bibliographystyle{plain}
\bibliography{cite}

\end{document}\documentclass{article}
\usepackage[utf8]{inputenc}
\usepackage{booktabs}
\usepackage{amsmath}
\usepackage{natbib}

\title{My Document}
\author{Your Name}
\date{\today}

\begin{document}

\maketitle

\section{Introduction}
This is a sample document demonstrating various \LaTeX{} features including mathematical equations, tables, and formatting. You can cite papers using the citation search feature, and they will be automatically added to your \texttt{cite.bib} file.

\section{Mathematical Expressions}
Here are some example equations:

\subsection{Quadratic Formula}
The quadratic formula is given by:
\[x = \frac{-b \pm \sqrt{b^2 - 4ac}}{2a}\]

\subsection{Matrix Example}
A sample matrix:
\[
\begin{pmatrix}
a_{11} & a_{12} & a_{13} \\
a_{21} & a_{22} & a_{23} \\
a_{31} & a_{32} & a_{33}
\end{pmatrix}
\]

\section{Tables}
Below is a sample table showing data:

\begin{table}[h]
\centering
\begin{tabular}{@{}lrrr@{}}
\toprule
\textbf{Category} & \textbf{Value 1} & \textbf{Value 2} & \textbf{Total} \\
\midrule
Item A & 10.5 & 20.3 & 30.8 \\
Item B & 15.2 & 25.1 & 40.3 \\
Item C & 22.1 & 12.7 & 34.8 \\
\bottomrule
\end{tabular}
\caption{Sample Data Table}
\label{tab:sample}
\end{table}

\section{Lists}
\subsection{Bullet Points}
\begin{itemize}
  \item First item
  \item Second item
  \item Third item with sub-items:
    \begin{itemize}
      \item Sub-item 1
      \item Sub-item 2
    \end{itemize}
\end{itemize}

\subsection{Numbered List}
\begin{enumerate}
  \item First step
  \item Second step
  \item Third step
\end{enumerate}

\section{Citations and References}
When you use the citation search feature, references will be automatically managed. For example, you might cite a paper like this: \cite{example2024}. All citations will be automatically added to your \texttt{cite.bib} file.

\section{Conclusion}
This template demonstrates basic \LaTeX{} features including mathematical formulas, tables, and lists. The bibliography system is now set up to work seamlessly with Bibby AI's citation management. Feel free to modify it for your needs!

\bibliographystyle{plain}
\bibliography{cite}

\end{document}\documentclass{article}
\usepackage[utf8]{inputenc}
\usepackage{booktabs}
\usepackage{amsmath}
\usepackage{natbib}

\title{My Document}
\author{Your Name}
\date{\today}

\begin{document}

\maketitle

\section{Introduction}
This is a sample document demonstrating various \LaTeX{} features including mathematical equations, tables, and formatting. You can cite papers using the citation search feature, and they will be automatically added to your \texttt{cite.bib} file.

\section{Mathematical Expressions}
Here are some example equations:

\subsection{Quadratic Formula}
The quadratic formula is given by:
\[x = \frac{-b \pm \sqrt{b^2 - 4ac}}{2a}\]

\subsection{Matrix Example}
A sample matrix:
\[
\begin{pmatrix}
a_{11} & a_{12} & a_{13} \\
a_{21} & a_{22} & a_{23} \\
a_{31} & a_{32} & a_{33}
\end{pmatrix}
\]

\section{Tables}
Below is a sample table showing data:

\begin{table}[h]
\centering
\begin{tabular}{@{}lrrr@{}}
\toprule
\textbf{Category} & \textbf{Value 1} & \textbf{Value 2} & \textbf{Total} \\
\midrule
Item A & 10.5 & 20.3 & 30.8 \\
Item B & 15.2 & 25.1 & 40.3 \\
Item C & 22.1 & 12.7 & 34.8 \\
\bottomrule
\end{tabular}
\caption{Sample Data Table}
\label{tab:sample}
\end{table}

\section{Lists}
\subsection{Bullet Points}
\begin{itemize}
  \item First item
  \item Second item
  \item Third item with sub-items:
    \begin{itemize}
      \item Sub-item 1
      \item Sub-item 2
    \end{itemize}
\end{itemize}

\subsection{Numbered List}
\begin{enumerate}
  \item First step
  \item Second step
  \item Third step
\end{enumerate}

\section{Citations and References}
When you use the citation search feature, references will be automatically managed. For example, you might cite a paper like this: \cite{example2024}. All citations will be automatically added to your \texttt{cite.bib} file.

\section{Conclusion}
This template demonstrates basic \LaTeX{} features including mathematical formulas, tables, and lists. The bibliography system is now set up to work seamlessly with Bibby AI's citation management. Feel free to modify it for your needs!

\bibliographystyle{plain}
\bibliography{cite}

\end{document}\documentclass{article}
\usepackage[utf8]{inputenc}
\usepackage{booktabs}
\usepackage{amsmath}
\usepackage{natbib}

\title{My Document}
\author{Your Name}
\date{\today}

\begin{document}

\maketitle

\section{Introduction}
This is a sample document demonstrating various \LaTeX{} features including mathematical equations, tables, and formatting. You can cite papers using the citation search feature, and they will be automatically added to your \texttt{cite.bib} file.

\section{Mathematical Expressions}
Here are some example equations:

\subsection{Quadratic Formula}
The quadratic formula is given by:
\[x = \frac{-b \pm \sqrt{b^2 - 4ac}}{2a}\]

\subsection{Matrix Example}
A sample matrix:
\[
\begin{pmatrix}
a_{11} & a_{12} & a_{13} \\
a_{21} & a_{22} & a_{23} \\
a_{31} & a_{32} & a_{33}
\end{pmatrix}
\]

\section{Tables}
Below is a sample table showing data:

\begin{table}[h]
\centering
\begin{tabular}{@{}lrrr@{}}
\toprule
\textbf{Category} & \textbf{Value 1} & \textbf{Value 2} & \textbf{Total} \\
\midrule
Item A & 10.5 & 20.3 & 30.8 \\
Item B & 15.2 & 25.1 & 40.3 \\
Item C & 22.1 & 12.7 & 34.8 \\
\bottomrule
\end{tabular}
\caption{Sample Data Table}
\label{tab:sample}
\end{table}

\section{Lists}
\subsection{Bullet Points}
\begin{itemize}
  \item First item
  \item Second item
  \item Third item with sub-items:
    \begin{itemize}
      \item Sub-item 1
      \item Sub-item 2
    \end{itemize}
\end{itemize}

\subsection{Numbered List}
\begin{enumerate}
  \item First step
  \item Second step
  \item Third step
\end{enumerate}

\section{Citations and References}
When you use the citation search feature, references will be automatically managed. For example, you might cite a paper like this: \cite{example2024}. All citations will be automatically added to your \texttt{cite.bib} file.

\section{Conclusion}
This template demonstrates basic \LaTeX{} features including mathematical formulas, tables, and lists. The bibliography system is now set up to work seamlessly with Bibby AI's citation management. Feel free to modify it for your needs!

\bibliographystyle{plain}
\bibliography{cite}

\end{document}\documentclass{article}
\usepackage[utf8]{inputenc}
\usepackage{booktabs}
\usepackage{amsmath}
\usepackage{natbib}

\title{My Document}
\author{Your Name}
\date{\today}

\begin{document}

\maketitle

\section{Introduction}
This is a sample document demonstrating various \LaTeX{} features including mathematical equations, tables, and formatting. You can cite papers using the citation search feature, and they will be automatically added to your \texttt{cite.bib} file.

\section{Mathematical Expressions}
Here are some example equations:

\subsection{Quadratic Formula}
The quadratic formula is given by:
\[x = \frac{-b \pm \sqrt{b^2 - 4ac}}{2a}\]

\subsection{Matrix Example}
A sample matrix:
\[
\begin{pmatrix}
a_{11} & a_{12} & a_{13} \\
a_{21} & a_{22} & a_{23} \\
a_{31} & a_{32} & a_{33}
\end{pmatrix}
\]

\section{Tables}
Below is a sample table showing data:

\begin{table}[h]
\centering
\begin{tabular}{@{}lrrr@{}}
\toprule
\textbf{Category} & \textbf{Value 1} & \textbf{Value 2} & \textbf{Total} \\
\midrule
Item A & 10.5 & 20.3 & 30.8 \\
Item B & 15.2 & 25.1 & 40.3 \\
Item C & 22.1 & 12.7 & 34.8 \\
\bottomrule
\end{tabular}
\caption{Sample Data Table}
\label{tab:sample}
\end{table}

\section{Lists}
\subsection{Bullet Points}
\begin{itemize}
  \item First item
  \item Second item
  \item Third item with sub-items:
    \begin{itemize}
      \item Sub-item 1
      \item Sub-item 2
    \end{itemize}
\end{itemize}

\subsection{Numbered List}
\begin{enumerate}
  \item First step
  \item Second step
  \item Third step
\end{enumerate}

\section{Citations and References}
When you use the citation search feature, references will be automatically managed. For example, you might cite a paper like this: \cite{example2024}. All citations will be automatically added to your \texttt{cite.bib} file.

\section{Conclusion}
This template demonstrates basic \LaTeX{} features including mathematical formulas, tables, and lists. The bibliography system is now set up to work seamlessly with Bibby AI's citation management. Feel free to modify it for your needs!

\bibliographystyle{plain}
\bibliography{cite}

\end{document}\documentclass{article}
\usepackage[utf8]{inputenc}
\usepackage{booktabs}
\usepackage{amsmath}
\usepackage{natbib}

\title{My Document}
\author{Your Name}
\date{\today}

\begin{document}

\maketitle

\section{Introduction}
This is a sample document demonstrating various \LaTeX{} features including mathematical equations, tables, and formatting. You can cite papers using the citation search feature, and they will be automatically added to your \texttt{cite.bib} file.

\section{Mathematical Expressions}
Here are some example equations:

\subsection{Quadratic Formula}
The quadratic formula is given by:
\[x = \frac{-b \pm \sqrt{b^2 - 4ac}}{2a}\]

\subsection{Matrix Example}
A sample matrix:
\[
\begin{pmatrix}
a_{11} & a_{12} & a_{13} \\
a_{21} & a_{22} & a_{23} \\
a_{31} & a_{32} & a_{33}
\end{pmatrix}
\]

\section{Tables}
Below is a sample table showing data:

\begin{table}[h]
\centering
\begin{tabular}{@{}lrrr@{}}
\toprule
\textbf{Category} & \textbf{Value 1} & \textbf{Value 2} & \textbf{Total} \\
\midrule
Item A & 10.5 & 20.3 & 30.8 \\
Item B & 15.2 & 25.1 & 40.3 \\
Item C & 22.1 & 12.7 & 34.8 \\
\bottomrule
\end{tabular}
\caption{Sample Data Table}
\label{tab:sample}
\end{table}

\section{Lists}
\subsection{Bullet Points}
\begin{itemize}
  \item First item
  \item Second item
  \item Third item with sub-items:
    \begin{itemize}
      \item Sub-item 1
      \item Sub-item 2
    \end{itemize}
\end{itemize}

\subsection{Numbered List}
\begin{enumerate}
  \item First step
  \item Second step
  \item Third step
\end{enumerate}

\section{Citations and References}
When you use the citation search feature, references will be automatically managed. For example, you might cite a paper like this: \cite{example2024}. All citations will be automatically added to your \texttt{cite.bib} file.

\section{Conclusion}
This template demonstrates basic \LaTeX{} features including mathematical formulas, tables, and lists. The bibliography system is now set up to work seamlessly with Bibby AI's citation management. Feel free to modify it for your needs!

\bibliographystyle{plain}
\bibliography{cite}

\end{document}\documentclass{article}
\usepackage[utf8]{inputenc}
\usepackage{booktabs}
\usepackage{amsmath}
\usepackage{natbib}

\title{My Document}
\author{Your Name}
\date{\today}

\begin{document}

\maketitle

\section{Introduction}
This is a sample document demonstrating various \LaTeX{} features including mathematical equations, tables, and formatting. You can cite papers using the citation search feature, and they will be automatically added to your \texttt{cite.bib} file.

\section{Mathematical Expressions}
Here are some example equations:

\subsection{Quadratic Formula}
The quadratic formula is given by:
\[x = \frac{-b \pm \sqrt{b^2 - 4ac}}{2a}\]

\subsection{Matrix Example}
A sample matrix:
\[
\begin{pmatrix}
a_{11} & a_{12} & a_{13} \\
a_{21} & a_{22} & a_{23} \\
a_{31} & a_{32} & a_{33}
\end{pmatrix}
\]

\section{Tables}
Below is a sample table showing data:

\begin{table}[h]
\centering
\begin{tabular}{@{}lrrr@{}}
\toprule
\textbf{Category} & \textbf{Value 1} & \textbf{Value 2} & \textbf{Total} \\
\midrule
Item A & 10.5 & 20.3 & 30.8 \\
Item B & 15.2 & 25.1 & 40.3 \\
Item C & 22.1 & 12.7 & 34.8 \\
\bottomrule
\end{tabular}
\caption{Sample Data Table}
\label{tab:sample}
\end{table}

\section{Lists}
\subsection{Bullet Points}
\begin{itemize}
  \item First item
  \item Second item
  \item Third item with sub-items:
    \begin{itemize}
      \item Sub-item 1
      \item Sub-item 2
    \end{itemize}
\end{itemize}

\subsection{Numbered List}
\begin{enumerate}
  \item First step
  \item Second step
  \item Third step
\end{enumerate}

\section{Citations and References}
When you use the citation search feature, references will be automatically managed. For example, you might cite a paper like this: \cite{example2024}. All citations will be automatically added to your \texttt{cite.bib} file.

\section{Conclusion}
This template demonstrates basic \LaTeX{} features including mathematical formulas, tables, and lists. The bibliography system is now set up to work seamlessly with Bibby AI's citation management. Feel free to modify it for your needs!

\bibliographystyle{plain}
\bibliography{cite}

\end{document}\documentclass{article}
\usepackage[utf8]{inputenc}
\usepackage{booktabs}
\usepackage{amsmath}
\usepackage{natbib}

\title{My Document}
\author{Your Name}
\date{\today}

\begin{document}

\maketitle

\section{Introduction}
This is a sample document demonstrating various \LaTeX{} features including mathematical equations, tables, and formatting. You can cite papers using the citation search feature, and they will be automatically added to your \texttt{cite.bib} file.

\section{Mathematical Expressions}
Here are some example equations:

\subsection{Quadratic Formula}
The quadratic formula is given by:
\[x = \frac{-b \pm \sqrt{b^2 - 4ac}}{2a}\]

\subsection{Matrix Example}
A sample matrix:
\[
\begin{pmatrix}
a_{11} & a_{12} & a_{13} \\
a_{21} & a_{22} & a_{23} \\
a_{31} & a_{32} & a_{33}
\end{pmatrix}
\]

\section{Tables}
Below is a sample table showing data:

\begin{table}[h]
\centering
\begin{tabular}{@{}lrrr@{}}
\toprule
\textbf{Category} & \textbf{Value 1} & \textbf{Value 2} & \textbf{Total} \\
\midrule
Item A & 10.5 & 20.3 & 30.8 \\
Item B & 15.2 & 25.1 & 40.3 \\
Item C & 22.1 & 12.7 & 34.8 \\
\bottomrule
\end{tabular}
\caption{Sample Data Table}
\label{tab:sample}
\end{table}

\section{Lists}
\subsection{Bullet Points}
\begin{itemize}
  \item First item
  \item Second item
  \item Third item with sub-items:
    \begin{itemize}
      \item Sub-item 1
      \item Sub-item 2
    \end{itemize}
\end{itemize}

\subsection{Numbered List}
\begin{enumerate}
  \item First step
  \item Second step
  \item Third step
\end{enumerate}

\section{Citations and References}
When you use the citation search feature, references will be automatically managed. For example, you might cite a paper like this: \cite{example2024}. All citations will be automatically added to your \texttt{cite.bib} file.

\section{Conclusion}
This template demonstrates basic \LaTeX{} features including mathematical formulas, tables, and lists. The bibliography system is now set up to work seamlessly with Bibby AI's citation management. Feel free to modify it for your needs!

\bibliographystyle{plain}
\bibliography{cite}

\end{document}\documentclass{article}
\usepackage[utf8]{inputenc}
\usepackage{booktabs}
\usepackage{amsmath}
\usepackage{natbib}

\title{My Document}
\author{Your Name}
\date{\today}

\begin{document}

\maketitle

\section{Introduction}
This is a sample document demonstrating various \LaTeX{} features including mathematical equations, tables, and formatting. You can cite papers using the citation search feature, and they will be automatically added to your \texttt{cite.bib} file.

\section{Mathematical Expressions}
Here are some example equations:

\subsection{Quadratic Formula}
The quadratic formula is given by:
\[x = \frac{-b \pm \sqrt{b^2 - 4ac}}{2a}\]

\subsection{Matrix Example}
A sample matrix:
\[
\begin{pmatrix}
a_{11} & a_{12} & a_{13} \\
a_{21} & a_{22} & a_{23} \\
a_{31} & a_{32} & a_{33}
\end{pmatrix}
\]

\section{Tables}
Below is a sample table showing data:

\begin{table}[h]
\centering
\begin{tabular}{@{}lrrr@{}}
\toprule
\textbf{Category} & \textbf{Value 1} & \textbf{Value 2} & \textbf{Total} \\
\midrule
Item A & 10.5 & 20.3 & 30.8 \\
Item B & 15.2 & 25.1 & 40.3 \\
Item C & 22.1 & 12.7 & 34.8 \\
\bottomrule
\end{tabular}
\caption{Sample Data Table}
\label{tab:sample}
\end{table}

\section{Lists}
\subsection{Bullet Points}
\begin{itemize}
  \item First item
  \item Second item
  \item Third item with sub-items:
    \begin{itemize}
      \item Sub-item 1
      \item Sub-item 2
    \end{itemize}
\end{itemize}

\subsection{Numbered List}
\begin{enumerate}
  \item First step
  \item Second step
  \item Third step
\end{enumerate}

\section{Citations and References}
When you use the citation search feature, references will be automatically managed. For example, you might cite a paper like this: \cite{example2024}. All citations will be automatically added to your \texttt{cite.bib} file.

\section{Conclusion}
This template demonstrates basic \LaTeX{} features including mathematical formulas, tables, and lists. The bibliography system is now set up to work seamlessly with Bibby AI's citation management. Feel free to modify it for your needs!

\bibliographystyle{plain}
\bibliography{cite}

\end{document}\documentclass{article}
\usepackage[utf8]{inputenc}
\usepackage{booktabs}
\usepackage{amsmath}
\usepackage{natbib}

\title{My Document}
\author{Your Name}
\date{\today}

\begin{document}

\maketitle

\section{Introduction}
This is a sample document demonstrating various \LaTeX{} features including mathematical equations, tables, and formatting. You can cite papers using the citation search feature, and they will be automatically added to your \texttt{cite.bib} file.

\section{Mathematical Expressions}
Here are some example equations:

\subsection{Quadratic Formula}
The quadratic formula is given by:
\[x = \frac{-b \pm \sqrt{b^2 - 4ac}}{2a}\]

\subsection{Matrix Example}
A sample matrix:
\[
\begin{pmatrix}
a_{11} & a_{12} & a_{13} \\
a_{21} & a_{22} & a_{23} \\
a_{31} & a_{32} & a_{33}
\end{pmatrix}
\]

\section{Tables}
Below is a sample table showing data:

\begin{table}[h]
\centering
\begin{tabular}{@{}lrrr@{}}
\toprule
\textbf{Category} & \textbf{Value 1} & \textbf{Value 2} & \textbf{Total} \\
\midrule
Item A & 10.5 & 20.3 & 30.8 \\
Item B & 15.2 & 25.1 & 40.3 \\
Item C & 22.1 & 12.7 & 34.8 \\
\bottomrule
\end{tabular}
\caption{Sample Data Table}
\label{tab:sample}
\end{table}

\section{Lists}
\subsection{Bullet Points}
\begin{itemize}
  \item First item
  \item Second item
  \item Third item with sub-items:
    \begin{itemize}
      \item Sub-item 1
      \item Sub-item 2
    \end{itemize}
\end{itemize}

\subsection{Numbered List}
\begin{enumerate}
  \item First step
  \item Second step
  \item Third step
\end{enumerate}

\section{Citations and References}
When you use the citation search feature, references will be automatically managed. For example, you might cite a paper like this: \cite{example2024}. All citations will be automatically added to your \texttt{cite.bib} file.

\section{Conclusion}
This template demonstrates basic \LaTeX{} features including mathematical formulas, tables, and lists. The bibliography system is now set up to work seamlessly with Bibby AI's citation management. Feel free to modify it for your needs!

\bibliographystyle{plain}
\bibliography{cite}

\end{document}\documentclass{article}
\usepackage[utf8]{inputenc}
\usepackage{booktabs}
\usepackage{amsmath}
\usepackage{natbib}

\title{My Document}
\author{Your Name}
\date{\today}

\begin{document}

\maketitle

\section{Introduction}
This is a sample document demonstrating various \LaTeX{} features including mathematical equations, tables, and formatting. You can cite papers using the citation search feature, and they will be automatically added to your \texttt{cite.bib} file.

\section{Mathematical Expressions}
Here are some example equations:

\subsection{Quadratic Formula}
The quadratic formula is given by:
\[x = \frac{-b \pm \sqrt{b^2 - 4ac}}{2a}\]

\subsection{Matrix Example}
A sample matrix:
\[
\begin{pmatrix}
a_{11} & a_{12} & a_{13} \\
a_{21} & a_{22} & a_{23} \\
a_{31} & a_{32} & a_{33}
\end{pmatrix}
\]

\section{Tables}
Below is a sample table showing data:

\begin{table}[h]
\centering
\begin{tabular}{@{}lrrr@{}}
\toprule
\textbf{Category} & \textbf{Value 1} & \textbf{Value 2} & \textbf{Total} \\
\midrule
Item A & 10.5 & 20.3 & 30.8 \\
Item B & 15.2 & 25.1 & 40.3 \\
Item C & 22.1 & 12.7 & 34.8 \\
\bottomrule
\end{tabular}
\caption{Sample Data Table}
\label{tab:sample}
\end{table}

\section{Lists}
\subsection{Bullet Points}
\begin{itemize}
  \item First item
  \item Second item
  \item Third item with sub-items:
    \begin{itemize}
      \item Sub-item 1
      \item Sub-item 2
    \end{itemize}
\end{itemize}

\subsection{Numbered List}
\begin{enumerate}
  \item First step
  \item Second step
  \item Third step
\end{enumerate}

\section{Citations and References}
When you use the citation search feature, references will be automatically managed. For example, you might cite a paper like this: \cite{example2024}. All citations will be automatically added to your \texttt{cite.bib} file.

\section{Conclusion}
This template demonstrates basic \LaTeX{} features including mathematical formulas, tables, and lists. The bibliography system is now set up to work seamlessly with Bibby AI's citation management. Feel free to modify it for your needs!

\bibliographystyle{plain}
\bibliography{cite}

\end{document}\documentclass{article}
\usepackage[utf8]{inputenc}
\usepackage{booktabs}
\usepackage{amsmath}
\usepackage{natbib}

\title{My Document}
\author{Your Name}
\date{\today}

\begin{document}

\maketitle

\section{Introduction}
This is a sample document demonstrating various \LaTeX{} features including mathematical equations, tables, and formatting. You can cite papers using the citation search feature, and they will be automatically added to your \texttt{cite.bib} file.

\section{Mathematical Expressions}
Here are some example equations:

\subsection{Quadratic Formula}
The quadratic formula is given by:
\[x = \frac{-b \pm \sqrt{b^2 - 4ac}}{2a}\]

\subsection{Matrix Example}
A sample matrix:
\[
\begin{pmatrix}
a_{11} & a_{12} & a_{13} \\
a_{21} & a_{22} & a_{23} \\
a_{31} & a_{32} & a_{33}
\end{pmatrix}
\]

\section{Tables}
Below is a sample table showing data:

\begin{table}[h]
\centering
\begin{tabular}{@{}lrrr@{}}
\toprule
\textbf{Category} & \textbf{Value 1} & \textbf{Value 2} & \textbf{Total} \\
\midrule
Item A & 10.5 & 20.3 & 30.8 \\
Item B & 15.2 & 25.1 & 40.3 \\
Item C & 22.1 & 12.7 & 34.8 \\
\bottomrule
\end{tabular}
\caption{Sample Data Table}
\label{tab:sample}
\end{table}

\section{Lists}
\subsection{Bullet Points}
\begin{itemize}
  \item First item
  \item Second item
  \item Third item with sub-items:
    \begin{itemize}
      \item Sub-item 1
      \item Sub-item 2
    \end{itemize}
\end{itemize}

\subsection{Numbered List}
\begin{enumerate}
  \item First step
  \item Second step
  \item Third step
\end{enumerate}

\section{Citations and References}
When you use the citation search feature, references will be automatically managed. For example, you might cite a paper like this: \cite{example2024}. All citations will be automatically added to your \texttt{cite.bib} file.

\section{Conclusion}
This template demonstrates basic \LaTeX{} features including mathematical formulas, tables, and lists. The bibliography system is now set up to work seamlessly with Bibby AI's citation management. Feel free to modify it for your needs!

\bibliographystyle{plain}
\bibliography{cite}

\end{document}\documentclass{article}
\usepackage[utf8]{inputenc}
\usepackage{booktabs}
\usepackage{amsmath}
\usepackage{natbib}

\title{My Document}
\author{Your Name}
\date{\today}

\begin{document}

\maketitle

\section{Introduction}
This is a sample document demonstrating various \LaTeX{} features including mathematical equations, tables, and formatting. You can cite papers using the citation search feature, and they will be automatically added to your \texttt{cite.bib} file.

\section{Mathematical Expressions}
Here are some example equations:

\subsection{Quadratic Formula}
The quadratic formula is given by:
\[x = \frac{-b \pm \sqrt{b^2 - 4ac}}{2a}\]

\subsection{Matrix Example}
A sample matrix:
\[
\begin{pmatrix}
a_{11} & a_{12} & a_{13} \\
a_{21} & a_{22} & a_{23} \\
a_{31} & a_{32} & a_{33}
\end{pmatrix}
\]

\section{Tables}
Below is a sample table showing data:

\begin{table}[h]
\centering
\begin{tabular}{@{}lrrr@{}}
\toprule
\textbf{Category} & \textbf{Value 1} & \textbf{Value 2} & \textbf{Total} \\
\midrule
Item A & 10.5 & 20.3 & 30.8 \\
Item B & 15.2 & 25.1 & 40.3 \\
Item C & 22.1 & 12.7 & 34.8 \\
\bottomrule
\end{tabular}
\caption{Sample Data Table}
\label{tab:sample}
\end{table}

\section{Lists}
\subsection{Bullet Points}
\begin{itemize}
  \item First item
  \item Second item
  \item Third item with sub-items:
    \begin{itemize}
      \item Sub-item 1
      \item Sub-item 2
    \end{itemize}
\end{itemize}

\subsection{Numbered List}
\begin{enumerate}
  \item First step
  \item Second step
  \item Third step
\end{enumerate}

\section{Citations and References}
When you use the citation search feature, references will be automatically managed. For example, you might cite a paper like this: \cite{example2024}. All citations will be automatically added to your \texttt{cite.bib} file.

\section{Conclusion}
This template demonstrates basic \LaTeX{} features including mathematical formulas, tables, and lists. The bibliography system is now set up to work seamlessly with Bibby AI's citation management. Feel free to modify it for your needs!

\bibliographystyle{plain}
\bibliography{cite}

\end{document}\documentclass{article}
\usepackage[utf8]{inputenc}
\usepackage{booktabs}
\usepackage{amsmath}
\usepackage{natbib}

\title{My Document}
\author{Your Name}
\date{\today}

\begin{document}

\maketitle

\section{Introduction}
This is a sample document demonstrating various \LaTeX{} features including mathematical equations, tables, and formatting. You can cite papers using the citation search feature, and they will be automatically added to your \texttt{cite.bib} file.

\section{Mathematical Expressions}
Here are some example equations:

\subsection{Quadratic Formula}
The quadratic formula is given by:
\[x = \frac{-b \pm \sqrt{b^2 - 4ac}}{2a}\]

\subsection{Matrix Example}
A sample matrix:
\[
\begin{pmatrix}
a_{11} & a_{12} & a_{13} \\
a_{21} & a_{22} & a_{23} \\
a_{31} & a_{32} & a_{33}
\end{pmatrix}
\]

\section{Tables}
Below is a sample table showing data:

\begin{table}[h]
\centering
\begin{tabular}{@{}lrrr@{}}
\toprule
\textbf{Category} & \textbf{Value 1} & \textbf{Value 2} & \textbf{Total} \\
\midrule
Item A & 10.5 & 20.3 & 30.8 \\
Item B & 15.2 & 25.1 & 40.3 \\
Item C & 22.1 & 12.7 & 34.8 \\
\bottomrule
\end{tabular}
\caption{Sample Data Table}
\label{tab:sample}
\end{table}

\section{Lists}
\subsection{Bullet Points}
\begin{itemize}
  \item First item
  \item Second item
  \item Third item with sub-items:
    \begin{itemize}
      \item Sub-item 1
      \item Sub-item 2
    \end{itemize}
\end{itemize}

\subsection{Numbered List}
\begin{enumerate}
  \item First step
  \item Second step
  \item Third step
\end{enumerate}

\section{Citations and References}
When you use the citation search feature, references will be automatically managed. For example, you might cite a paper like this: \cite{example2024}. All citations will be automatically added to your \texttt{cite.bib} file.

\section{Conclusion}
This template demonstrates basic \LaTeX{} features including mathematical formulas, tables, and lists. The bibliography system is now set up to work seamlessly with Bibby AI's citation management. Feel free to modify it for your needs!

\bibliographystyle{plain}
\bibliography{cite}

\end{document}\documentclass{article}
\usepackage[utf8]{inputenc}
\usepackage{booktabs}
\usepackage{amsmath}
\usepackage{natbib}

\title{My Document}
\author{Your Name}
\date{\today}

\begin{document}

\maketitle

\section{Introduction}
This is a sample document demonstrating various \LaTeX{} features including mathematical equations, tables, and formatting. You can cite papers using the citation search feature, and they will be automatically added to your \texttt{cite.bib} file.

\section{Mathematical Expressions}
Here are some example equations:

\subsection{Quadratic Formula}
The quadratic formula is given by:
\[x = \frac{-b \pm \sqrt{b^2 - 4ac}}{2a}\]

\subsection{Matrix Example}
A sample matrix:
\[
\begin{pmatrix}
a_{11} & a_{12} & a_{13} \\
a_{21} & a_{22} & a_{23} \\
a_{31} & a_{32} & a_{33}
\end{pmatrix}
\]

\section{Tables}
Below is a sample table showing data:

\begin{table}[h]
\centering
\begin{tabular}{@{}lrrr@{}}
\toprule
\textbf{Category} & \textbf{Value 1} & \textbf{Value 2} & \textbf{Total} \\
\midrule
Item A & 10.5 & 20.3 & 30.8 \\
Item B & 15.2 & 25.1 & 40.3 \\
Item C & 22.1 & 12.7 & 34.8 \\
\bottomrule
\end{tabular}
\caption{Sample Data Table}
\label{tab:sample}
\end{table}

\section{Lists}
\subsection{Bullet Points}
\begin{itemize}
  \item First item
  \item Second item
  \item Third item with sub-items:
    \begin{itemize}
      \item Sub-item 1
      \item Sub-item 2
    \end{itemize}
\end{itemize}

\subsection{Numbered List}
\begin{enumerate}
  \item First step
  \item Second step
  \item Third step
\end{enumerate}

\section{Citations and References}
When you use the citation search feature, references will be automatically managed. For example, you might cite a paper like this: \cite{example2024}. All citations will be automatically added to your \texttt{cite.bib} file.

\section{Conclusion}
This template demonstrates basic \LaTeX{} features including mathematical formulas, tables, and lists. The bibliography system is now set up to work seamlessly with Bibby AI's citation management. Feel free to modify it for your needs!

\bibliographystyle{plain}
\bibliography{cite}

\end{document}\documentclass{article}
\usepackage[utf8]{inputenc}
\usepackage{booktabs}
\usepackage{amsmath}
\usepackage{natbib}

\title{My Document}
\author{Your Name}
\date{\today}

\begin{document}

\maketitle

\section{Introduction}
This is a sample document demonstrating various \LaTeX{} features including mathematical equations, tables, and formatting. You can cite papers using the citation search feature, and they will be automatically added to your \texttt{cite.bib} file.

\section{Mathematical Expressions}
Here are some example equations:

\subsection{Quadratic Formula}
The quadratic formula is given by:
\[x = \frac{-b \pm \sqrt{b^2 - 4ac}}{2a}\]

\subsection{Matrix Example}
A sample matrix:
\[
\begin{pmatrix}
a_{11} & a_{12} & a_{13} \\
a_{21} & a_{22} & a_{23} \\
a_{31} & a_{32} & a_{33}
\end{pmatrix}
\]

\section{Tables}
Below is a sample table showing data:

\begin{table}[h]
\centering
\begin{tabular}{@{}lrrr@{}}
\toprule
\textbf{Category} & \textbf{Value 1} & \textbf{Value 2} & \textbf{Total} \\
\midrule
Item A & 10.5 & 20.3 & 30.8 \\
Item B & 15.2 & 25.1 & 40.3 \\
Item C & 22.1 & 12.7 & 34.8 \\
\bottomrule
\end{tabular}
\caption{Sample Data Table}
\label{tab:sample}
\end{table}

\section{Lists}
\subsection{Bullet Points}
\begin{itemize}
  \item First item
  \item Second item
  \item Third item with sub-items:
    \begin{itemize}
      \item Sub-item 1
      \item Sub-item 2
    \end{itemize}
\end{itemize}

\subsection{Numbered List}
\begin{enumerate}
  \item First step
  \item Second step
  \item Third step
\end{enumerate}

\section{Citations and References}
When you use the citation search feature, references will be automatically managed. For example, you might cite a paper like this: \cite{example2024}. All citations will be automatically added to your \texttt{cite.bib} file.

\section{Conclusion}
This template demonstrates basic \LaTeX{} features including mathematical formulas, tables, and lists. The bibliography system is now set up to work seamlessly with Bibby AI's citation management. Feel free to modify it for your needs!

\bibliographystyle{plain}
\bibliography{cite}

\end{document}\documentclass{article}
\usepackage[utf8]{inputenc}
\usepackage{booktabs}
\usepackage{amsmath}
\usepackage{natbib}

\title{My Document}
\author{Your Name}
\date{\today}

\begin{document}

\maketitle

\section{Introduction}
This is a sample document demonstrating various \LaTeX{} features including mathematical equations, tables, and formatting. You can cite papers using the citation search feature, and they will be automatically added to your \texttt{cite.bib} file.

\section{Mathematical Expressions}
Here are some example equations:

\subsection{Quadratic Formula}
The quadratic formula is given by:
\[x = \frac{-b \pm \sqrt{b^2 - 4ac}}{2a}\]

\subsection{Matrix Example}
A sample matrix:
\[
\begin{pmatrix}
a_{11} & a_{12} & a_{13} \\
a_{21} & a_{22} & a_{23} \\
a_{31} & a_{32} & a_{33}
\end{pmatrix}
\]

\section{Tables}
Below is a sample table showing data:

\begin{table}[h]
\centering
\begin{tabular}{@{}lrrr@{}}
\toprule
\textbf{Category} & \textbf{Value 1} & \textbf{Value 2} & \textbf{Total} \\
\midrule
Item A & 10.5 & 20.3 & 30.8 \\
Item B & 15.2 & 25.1 & 40.3 \\
Item C & 22.1 & 12.7 & 34.8 \\
\bottomrule
\end{tabular}
\caption{Sample Data Table}
\label{tab:sample}
\end{table}

\section{Lists}
\subsection{Bullet Points}
\begin{itemize}
  \item First item
  \item Second item
  \item Third item with sub-items:
    \begin{itemize}
      \item Sub-item 1
      \item Sub-item 2
    \end{itemize}
\end{itemize}

\subsection{Numbered List}
\begin{enumerate}
  \item First step
  \item Second step
  \item Third step
\end{enumerate}

\section{Citations and References}
When you use the citation search feature, references will be automatically managed. For example, you might cite a paper like this: \cite{example2024}. All citations will be automatically added to your \texttt{cite.bib} file.

\section{Conclusion}
This template demonstrates basic \LaTeX{} features including mathematical formulas, tables, and lists. The bibliography system is now set up to work seamlessly with Bibby AI's citation management. Feel free to modify it for your needs!

\bibliographystyle{plain}
\bibliography{cite}

\end{document}\documentclass{article}
\usepackage[utf8]{inputenc}
\usepackage{booktabs}
\usepackage{amsmath}
\usepackage{natbib}

\title{My Document}
\author{Your Name}
\date{\today}

\begin{document}

\maketitle

\section{Introduction}
This is a sample document demonstrating various \LaTeX{} features including mathematical equations, tables, and formatting. You can cite papers using the citation search feature, and they will be automatically added to your \texttt{cite.bib} file.

\section{Mathematical Expressions}
Here are some example equations:

\subsection{Quadratic Formula}
The quadratic formula is given by:
\[x = \frac{-b \pm \sqrt{b^2 - 4ac}}{2a}\]

\subsection{Matrix Example}
A sample matrix:
\[
\begin{pmatrix}
a_{11} & a_{12} & a_{13} \\
a_{21} & a_{22} & a_{23} \\
a_{31} & a_{32} & a_{33}
\end{pmatrix}
\]

\section{Tables}
Below is a sample table showing data:

\begin{table}[h]
\centering
\begin{tabular}{@{}lrrr@{}}
\toprule
\textbf{Category} & \textbf{Value 1} & \textbf{Value 2} & \textbf{Total} \\
\midrule
Item A & 10.5 & 20.3 & 30.8 \\
Item B & 15.2 & 25.1 & 40.3 \\
Item C & 22.1 & 12.7 & 34.8 \\
\bottomrule
\end{tabular}
\caption{Sample Data Table}
\label{tab:sample}
\end{table}

\section{Lists}
\subsection{Bullet Points}
\begin{itemize}
  \item First item
  \item Second item
  \item Third item with sub-items:
    \begin{itemize}
      \item Sub-item 1
      \item Sub-item 2
    \end{itemize}
\end{itemize}

\subsection{Numbered List}
\begin{enumerate}
  \item First step
  \item Second step
  \item Third step
\end{enumerate}

\section{Citations and References}
When you use the citation search feature, references will be automatically managed. For example, you might cite a paper like this: \cite{example2024}. All citations will be automatically added to your \texttt{cite.bib} file.

\section{Conclusion}
This template demonstrates basic \LaTeX{} features including mathematical formulas, tables, and lists. The bibliography system is now set up to work seamlessly with Bibby AI's citation management. Feel free to modify it for your needs!

\bibliographystyle{plain}
\bibliography{cite}

\end{document}\documentclass{article}
\usepackage[utf8]{inputenc}
\usepackage{booktabs}
\usepackage{amsmath}
\usepackage{natbib}

\title{My Document}
\author{Your Name}
\date{\today}

\begin{document}

\maketitle

\section{Introduction}
This is a sample document demonstrating various \LaTeX{} features including mathematical equations, tables, and formatting. You can cite papers using the citation search feature, and they will be automatically added to your \texttt{cite.bib} file.

\section{Mathematical Expressions}
Here are some example equations:

\subsection{Quadratic Formula}
The quadratic formula is given by:
\[x = \frac{-b \pm \sqrt{b^2 - 4ac}}{2a}\]

\subsection{Matrix Example}
A sample matrix:
\[
\begin{pmatrix}
a_{11} & a_{12} & a_{13} \\
a_{21} & a_{22} & a_{23} \\
a_{31} & a_{32} & a_{33}
\end{pmatrix}
\]

\section{Tables}
Below is a sample table showing data:

\begin{table}[h]
\centering
\begin{tabular}{@{}lrrr@{}}
\toprule
\textbf{Category} & \textbf{Value 1} & \textbf{Value 2} & \textbf{Total} \\
\midrule
Item A & 10.5 & 20.3 & 30.8 \\
Item B & 15.2 & 25.1 & 40.3 \\
Item C & 22.1 & 12.7 & 34.8 \\
\bottomrule
\end{tabular}
\caption{Sample Data Table}
\label{tab:sample}
\end{table}

\section{Lists}
\subsection{Bullet Points}
\begin{itemize}
  \item First item
  \item Second item
  \item Third item with sub-items:
    \begin{itemize}
      \item Sub-item 1
      \item Sub-item 2
    \end{itemize}
\end{itemize}

\subsection{Numbered List}
\begin{enumerate}
  \item First step
  \item Second step
  \item Third step
\end{enumerate}

\section{Citations and References}
When you use the citation search feature, references will be automatically managed. For example, you might cite a paper like this: \cite{example2024}. All citations will be automatically added to your \texttt{cite.bib} file.

\section{Conclusion}
This template demonstrates basic \LaTeX{} features including mathematical formulas, tables, and lists. The bibliography system is now set up to work seamlessly with Bibby AI's citation management. Feel free to modify it for your needs!

\bibliographystyle{plain}
\bibliography{cite}

\end{document}\documentclass{article}
\usepackage[utf8]{inputenc}
\usepackage{booktabs}
\usepackage{amsmath}
\usepackage{natbib}

\title{My Document}
\author{Your Name}
\date{\today}

\begin{document}

\maketitle

\section{Introduction}
This is a sample document demonstrating various \LaTeX{} features including mathematical equations, tables, and formatting. You can cite papers using the citation search feature, and they will be automatically added to your \texttt{cite.bib} file.

\section{Mathematical Expressions}
Here are some example equations:

\subsection{Quadratic Formula}
The quadratic formula is given by:
\[x = \frac{-b \pm \sqrt{b^2 - 4ac}}{2a}\]

\subsection{Matrix Example}
A sample matrix:
\[
\begin{pmatrix}
a_{11} & a_{12} & a_{13} \\
a_{21} & a_{22} & a_{23} \\
a_{31} & a_{32} & a_{33}
\end{pmatrix}
\]

\section{Tables}
Below is a sample table showing data:

\begin{table}[h]
\centering
\begin{tabular}{@{}lrrr@{}}
\toprule
\textbf{Category} & \textbf{Value 1} & \textbf{Value 2} & \textbf{Total} \\
\midrule
Item A & 10.5 & 20.3 & 30.8 \\
Item B & 15.2 & 25.1 & 40.3 \\
Item C & 22.1 & 12.7 & 34.8 \\
\bottomrule
\end{tabular}
\caption{Sample Data Table}
\label{tab:sample}
\end{table}

\section{Lists}
\subsection{Bullet Points}
\begin{itemize}
  \item First item
  \item Second item
  \item Third item with sub-items:
    \begin{itemize}
      \item Sub-item 1
      \item Sub-item 2
    \end{itemize}
\end{itemize}

\subsection{Numbered List}
\begin{enumerate}
  \item First step
  \item Second step
  \item Third step
\end{enumerate}

\section{Citations and References}
When you use the citation search feature, references will be automatically managed. For example, you might cite a paper like this: \cite{example2024}. All citations will be automatically added to your \texttt{cite.bib} file.

\section{Conclusion}
This template demonstrates basic \LaTeX{} features including mathematical formulas, tables, and lists. The bibliography system is now set up to work seamlessly with Bibby AI's citation management. Feel free to modify it for your needs!

\bibliographystyle{plain}
\bibliography{cite}

\end{document}\documentclass{article}
\usepackage[utf8]{inputenc}
\usepackage{booktabs}
\usepackage{amsmath}
\usepackage{natbib}

\title{My Document}
\author{Your Name}
\date{\today}

\begin{document}

\maketitle

\section{Introduction}
This is a sample document demonstrating various \LaTeX{} features including mathematical equations, tables, and formatting. You can cite papers using the citation search feature, and they will be automatically added to your \texttt{cite.bib} file.

\section{Mathematical Expressions}
Here are some example equations:

\subsection{Quadratic Formula}
The quadratic formula is given by:
\[x = \frac{-b \pm \sqrt{b^2 - 4ac}}{2a}\]

\subsection{Matrix Example}
A sample matrix:
\[
\begin{pmatrix}
a_{11} & a_{12} & a_{13} \\
a_{21} & a_{22} & a_{23} \\
a_{31} & a_{32} & a_{33}
\end{pmatrix}
\]

\section{Tables}
Below is a sample table showing data:

\begin{table}[h]
\centering
\begin{tabular}{@{}lrrr@{}}
\toprule
\textbf{Category} & \textbf{Value 1} & \textbf{Value 2} & \textbf{Total} \\
\midrule
Item A & 10.5 & 20.3 & 30.8 \\
Item B & 15.2 & 25.1 & 40.3 \\
Item C & 22.1 & 12.7 & 34.8 \\
\bottomrule
\end{tabular}
\caption{Sample Data Table}
\label{tab:sample}
\end{table}

\section{Lists}
\subsection{Bullet Points}
\begin{itemize}
  \item First item
  \item Second item
  \item Third item with sub-items:
    \begin{itemize}
      \item Sub-item 1
      \item Sub-item 2
    \end{itemize}
\end{itemize}

\subsection{Numbered List}
\begin{enumerate}
  \item First step
  \item Second step
  \item Third step
\end{enumerate}

\section{Citations and References}
When you use the citation search feature, references will be automatically managed. For example, you might cite a paper like this: \cite{example2024}. All citations will be automatically added to your \texttt{cite.bib} file.

\section{Conclusion}
This template demonstrates basic \LaTeX{} features including mathematical formulas, tables, and lists. The bibliography system is now set up to work seamlessly with Bibby AI's citation management. Feel free to modify it for your needs!

\bibliographystyle{plain}
\bibliography{cite}

\end{document}\documentclass{article}
\usepackage[utf8]{inputenc}
\usepackage{booktabs}
\usepackage{amsmath}
\usepackage{natbib}

\title{My Document}
\author{Your Name}
\date{\today}

\begin{document}

\maketitle

\section{Introduction}
This is a sample document demonstrating various \LaTeX{} features including mathematical equations, tables, and formatting. You can cite papers using the citation search feature, and they will be automatically added to your \texttt{cite.bib} file.

\section{Mathematical Expressions}
Here are some example equations:

\subsection{Quadratic Formula}
The quadratic formula is given by:
\[x = \frac{-b \pm \sqrt{b^2 - 4ac}}{2a}\]

\subsection{Matrix Example}
A sample matrix:
\[
\begin{pmatrix}
a_{11} & a_{12} & a_{13} \\
a_{21} & a_{22} & a_{23} \\
a_{31} & a_{32} & a_{33}
\end{pmatrix}
\]

\section{Tables}
Below is a sample table showing data:

\begin{table}[h]
\centering
\begin{tabular}{@{}lrrr@{}}
\toprule
\textbf{Category} & \textbf{Value 1} & \textbf{Value 2} & \textbf{Total} \\
\midrule
Item A & 10.5 & 20.3 & 30.8 \\
Item B & 15.2 & 25.1 & 40.3 \\
Item C & 22.1 & 12.7 & 34.8 \\
\bottomrule
\end{tabular}
\caption{Sample Data Table}
\label{tab:sample}
\end{table}

\section{Lists}
\subsection{Bullet Points}
\begin{itemize}
  \item First item
  \item Second item
  \item Third item with sub-items:
    \begin{itemize}
      \item Sub-item 1
      \item Sub-item 2
    \end{itemize}
\end{itemize}

\subsection{Numbered List}
\begin{enumerate}
  \item First step
  \item Second step
  \item Third step
\end{enumerate}

\section{Citations and References}
When you use the citation search feature, references will be automatically managed. For example, you might cite a paper like this: \cite{example2024}. All citations will be automatically added to your \texttt{cite.bib} file.

\section{Conclusion}
This template demonstrates basic \LaTeX{} features including mathematical formulas, tables, and lists. The bibliography system is now set up to work seamlessly with Bibby AI's citation management. Feel free to modify it for your needs!

\bibliographystyle{plain}
\bibliography{cite}

\end{document}\documentclass{article}
\usepackage[utf8]{inputenc}
\usepackage{booktabs}
\usepackage{amsmath}
\usepackage{natbib}

\title{My Document}
\author{Your Name}
\date{\today}

\begin{document}

\maketitle

\section{Introduction}
This is a sample document demonstrating various \LaTeX{} features including mathematical equations, tables, and formatting. You can cite papers using the citation search feature, and they will be automatically added to your \texttt{cite.bib} file.

\section{Mathematical Expressions}
Here are some example equations:

\subsection{Quadratic Formula}
The quadratic formula is given by:
\[x = \frac{-b \pm \sqrt{b^2 - 4ac}}{2a}\]

\subsection{Matrix Example}
A sample matrix:
\[
\begin{pmatrix}
a_{11} & a_{12} & a_{13} \\
a_{21} & a_{22} & a_{23} \\
a_{31} & a_{32} & a_{33}
\end{pmatrix}
\]

\section{Tables}
Below is a sample table showing data:

\begin{table}[h]
\centering
\begin{tabular}{@{}lrrr@{}}
\toprule
\textbf{Category} & \textbf{Value 1} & \textbf{Value 2} & \textbf{Total} \\
\midrule
Item A & 10.5 & 20.3 & 30.8 \\
Item B & 15.2 & 25.1 & 40.3 \\
Item C & 22.1 & 12.7 & 34.8 \\
\bottomrule
\end{tabular}
\caption{Sample Data Table}
\label{tab:sample}
\end{table}

\section{Lists}
\subsection{Bullet Points}
\begin{itemize}
  \item First item
  \item Second item
  \item Third item with sub-items:
    \begin{itemize}
      \item Sub-item 1
      \item Sub-item 2
    \end{itemize}
\end{itemize}

\subsection{Numbered List}
\begin{enumerate}
  \item First step
  \item Second step
  \item Third step
\end{enumerate}

\section{Citations and References}
When you use the citation search feature, references will be automatically managed. For example, you might cite a paper like this: \cite{example2024}. All citations will be automatically added to your \texttt{cite.bib} file.

\section{Conclusion}
This template demonstrates basic \LaTeX{} features including mathematical formulas, tables, and lists. The bibliography system is now set up to work seamlessly with Bibby AI's citation management. Feel free to modify it for your needs!

\bibliographystyle{plain}
\bibliography{cite}

\end{document}\documentclass{article}
\usepackage[utf8]{inputenc}
\usepackage{booktabs}
\usepackage{amsmath}
\usepackage{natbib}

\title{My Document}
\author{Your Name}
\date{\today}

\begin{document}

\maketitle

\section{Introduction}
This is a sample document demonstrating various \LaTeX{} features including mathematical equations, tables, and formatting. You can cite papers using the citation search feature, and they will be automatically added to your \texttt{cite.bib} file.

\section{Mathematical Expressions}
Here are some example equations:

\subsection{Quadratic Formula}
The quadratic formula is given by:
\[x = \frac{-b \pm \sqrt{b^2 - 4ac}}{2a}\]

\subsection{Matrix Example}
A sample matrix:
\[
\begin{pmatrix}
a_{11} & a_{12} & a_{13} \\
a_{21} & a_{22} & a_{23} \\
a_{31} & a_{32} & a_{33}
\end{pmatrix}
\]

\section{Tables}
Below is a sample table showing data:

\begin{table}[h]
\centering
\begin{tabular}{@{}lrrr@{}}
\toprule
\textbf{Category} & \textbf{Value 1} & \textbf{Value 2} & \textbf{Total} \\
\midrule
Item A & 10.5 & 20.3 & 30.8 \\
Item B & 15.2 & 25.1 & 40.3 \\
Item C & 22.1 & 12.7 & 34.8 \\
\bottomrule
\end{tabular}
\caption{Sample Data Table}
\label{tab:sample}
\end{table}

\section{Lists}
\subsection{Bullet Points}
\begin{itemize}
  \item First item
  \item Second item
  \item Third item with sub-items:
    \begin{itemize}
      \item Sub-item 1
      \item Sub-item 2
    \end{itemize}
\end{itemize}

\subsection{Numbered List}
\begin{enumerate}
  \item First step
  \item Second step
  \item Third step
\end{enumerate}

\section{Citations and References}
When you use the citation search feature, references will be automatically managed. For example, you might cite a paper like this: \cite{example2024}. All citations will be automatically added to your \texttt{cite.bib} file.

\section{Conclusion}
This template demonstrates basic \LaTeX{} features including mathematical formulas, tables, and lists. The bibliography system is now set up to work seamlessly with Bibby AI's citation management. Feel free to modify it for your needs!

\bibliographystyle{plain}
\bibliography{cite}

\end{document}\documentclass{article}
\usepackage[utf8]{inputenc}
\usepackage{booktabs}
\usepackage{amsmath}
\usepackage{natbib}

\title{My Document}
\author{Your Name}
\date{\today}

\begin{document}

\maketitle

\section{Introduction}
This is a sample document demonstrating various \LaTeX{} features including mathematical equations, tables, and formatting. You can cite papers using the citation search feature, and they will be automatically added to your \texttt{cite.bib} file.

\section{Mathematical Expressions}
Here are some example equations:

\subsection{Quadratic Formula}
The quadratic formula is given by:
\[x = \frac{-b \pm \sqrt{b^2 - 4ac}}{2a}\]

\subsection{Matrix Example}
A sample matrix:
\[
\begin{pmatrix}
a_{11} & a_{12} & a_{13} \\
a_{21} & a_{22} & a_{23} \\
a_{31} & a_{32} & a_{33}
\end{pmatrix}
\]

\section{Tables}
Below is a sample table showing data:

\begin{table}[h]
\centering
\begin{tabular}{@{}lrrr@{}}
\toprule
\textbf{Category} & \textbf{Value 1} & \textbf{Value 2} & \textbf{Total} \\
\midrule
Item A & 10.5 & 20.3 & 30.8 \\
Item B & 15.2 & 25.1 & 40.3 \\
Item C & 22.1 & 12.7 & 34.8 \\
\bottomrule
\end{tabular}
\caption{Sample Data Table}
\label{tab:sample}
\end{table}

\section{Lists}
\subsection{Bullet Points}
\begin{itemize}
  \item First item
  \item Second item
  \item Third item with sub-items:
    \begin{itemize}
      \item Sub-item 1
      \item Sub-item 2
    \end{itemize}
\end{itemize}

\subsection{Numbered List}
\begin{enumerate}
  \item First step
  \item Second step
  \item Third step
\end{enumerate}

\section{Citations and References}
When you use the citation search feature, references will be automatically managed. For example, you might cite a paper like this: \cite{example2024}. All citations will be automatically added to your \texttt{cite.bib} file.

\section{Conclusion}
This template demonstrates basic \LaTeX{} features including mathematical formulas, tables, and lists. The bibliography system is now set up to work seamlessly with Bibby AI's citation management. Feel free to modify it for your needs!

\bibliographystyle{plain}
\bibliography{cite}

\end{document}\documentclass{article}
\usepackage[utf8]{inputenc}
\usepackage{booktabs}
\usepackage{amsmath}
\usepackage{natbib}

\title{My Document}
\author{Your Name}
\date{\today}

\begin{document}

\maketitle

\section{Introduction}
This is a sample document demonstrating various \LaTeX{} features including mathematical equations, tables, and formatting. You can cite papers using the citation search feature, and they will be automatically added to your \texttt{cite.bib} file.

\section{Mathematical Expressions}
Here are some example equations:

\subsection{Quadratic Formula}
The quadratic formula is given by:
\[x = \frac{-b \pm \sqrt{b^2 - 4ac}}{2a}\]

\subsection{Matrix Example}
A sample matrix:
\[
\begin{pmatrix}
a_{11} & a_{12} & a_{13} \\
a_{21} & a_{22} & a_{23} \\
a_{31} & a_{32} & a_{33}
\end{pmatrix}
\]

\section{Tables}
Below is a sample table showing data:

\begin{table}[h]
\centering
\begin{tabular}{@{}lrrr@{}}
\toprule
\textbf{Category} & \textbf{Value 1} & \textbf{Value 2} & \textbf{Total} \\
\midrule
Item A & 10.5 & 20.3 & 30.8 \\
Item B & 15.2 & 25.1 & 40.3 \\
Item C & 22.1 & 12.7 & 34.8 \\
\bottomrule
\end{tabular}
\caption{Sample Data Table}
\label{tab:sample}
\end{table}

\section{Lists}
\subsection{Bullet Points}
\begin{itemize}
  \item First item
  \item Second item
  \item Third item with sub-items:
    \begin{itemize}
      \item Sub-item 1
      \item Sub-item 2
    \end{itemize}
\end{itemize}

\subsection{Numbered List}
\begin{enumerate}
  \item First step
  \item Second step
  \item Third step
\end{enumerate}

\section{Citations and References}
When you use the citation search feature, references will be automatically managed. For example, you might cite a paper like this: \cite{example2024}. All citations will be automatically added to your \texttt{cite.bib} file.

\section{Conclusion}
This template demonstrates basic \LaTeX{} features including mathematical formulas, tables, and lists. The bibliography system is now set up to work seamlessly with Bibby AI's citation management. Feel free to modify it for your needs!

\bibliographystyle{plain}
\bibliography{cite}

\end{document}\documentclass{article}
\usepackage[utf8]{inputenc}
\usepackage{booktabs}
\usepackage{amsmath}
\usepackage{natbib}

\title{My Document}
\author{Your Name}
\date{\today}

\begin{document}

\maketitle

\section{Introduction}
This is a sample document demonstrating various \LaTeX{} features including mathematical equations, tables, and formatting. You can cite papers using the citation search feature, and they will be automatically added to your \texttt{cite.bib} file.

\section{Mathematical Expressions}
Here are some example equations:

\subsection{Quadratic Formula}
The quadratic formula is given by:
\[x = \frac{-b \pm \sqrt{b^2 - 4ac}}{2a}\]

\subsection{Matrix Example}
A sample matrix:
\[
\begin{pmatrix}
a_{11} & a_{12} & a_{13} \\
a_{21} & a_{22} & a_{23} \\
a_{31} & a_{32} & a_{33}
\end{pmatrix}
\]

\section{Tables}
Below is a sample table showing data:

\begin{table}[h]
\centering
\begin{tabular}{@{}lrrr@{}}
\toprule
\textbf{Category} & \textbf{Value 1} & \textbf{Value 2} & \textbf{Total} \\
\midrule
Item A & 10.5 & 20.3 & 30.8 \\
Item B & 15.2 & 25.1 & 40.3 \\
Item C & 22.1 & 12.7 & 34.8 \\
\bottomrule
\end{tabular}
\caption{Sample Data Table}
\label{tab:sample}
\end{table}

\section{Lists}
\subsection{Bullet Points}
\begin{itemize}
  \item First item
  \item Second item
  \item Third item with sub-items:
    \begin{itemize}
      \item Sub-item 1
      \item Sub-item 2
    \end{itemize}
\end{itemize}

\subsection{Numbered List}
\begin{enumerate}
  \item First step
  \item Second step
  \item Third step
\end{enumerate}

\section{Citations and References}
When you use the citation search feature, references will be automatically managed. For example, you might cite a paper like this: \cite{example2024}. All citations will be automatically added to your \texttt{cite.bib} file.

\section{Conclusion}
This template demonstrates basic \LaTeX{} features including mathematical formulas, tables, and lists. The bibliography system is now set up to work seamlessly with Bibby AI's citation management. Feel free to modify it for your needs!

\bibliographystyle{plain}
\bibliography{cite}

\end{document}\documentclass{article}
\usepackage[utf8]{inputenc}
\usepackage{booktabs}
\usepackage{amsmath}
\usepackage{natbib}

\title{My Document}
\author{Your Name}
\date{\today}

\begin{document}

\maketitle

\section{Introduction}
This is a sample document demonstrating various \LaTeX{} features including mathematical equations, tables, and formatting. You can cite papers using the citation search feature, and they will be automatically added to your \texttt{cite.bib} file.

\section{Mathematical Expressions}
Here are some example equations:

\subsection{Quadratic Formula}
The quadratic formula is given by:
\[x = \frac{-b \pm \sqrt{b^2 - 4ac}}{2a}\]

\subsection{Matrix Example}
A sample matrix:
\[
\begin{pmatrix}
a_{11} & a_{12} & a_{13} \\
a_{21} & a_{22} & a_{23} \\
a_{31} & a_{32} & a_{33}
\end{pmatrix}
\]

\section{Tables}
Below is a sample table showing data:

\begin{table}[h]
\centering
\begin{tabular}{@{}lrrr@{}}
\toprule
\textbf{Category} & \textbf{Value 1} & \textbf{Value 2} & \textbf{Total} \\
\midrule
Item A & 10.5 & 20.3 & 30.8 \\
Item B & 15.2 & 25.1 & 40.3 \\
Item C & 22.1 & 12.7 & 34.8 \\
\bottomrule
\end{tabular}
\caption{Sample Data Table}
\label{tab:sample}
\end{table}

\section{Lists}
\subsection{Bullet Points}
\begin{itemize}
  \item First item
  \item Second item
  \item Third item with sub-items:
    \begin{itemize}
      \item Sub-item 1
      \item Sub-item 2
    \end{itemize}
\end{itemize}

\subsection{Numbered List}
\begin{enumerate}
  \item First step
  \item Second step
  \item Third step
\end{enumerate}

\section{Citations and References}
When you use the citation search feature, references will be automatically managed. For example, you might cite a paper like this: \cite{example2024}. All citations will be automatically added to your \texttt{cite.bib} file.

\section{Conclusion}
This template demonstrates basic \LaTeX{} features including mathematical formulas, tables, and lists. The bibliography system is now set up to work seamlessly with Bibby AI's citation management. Feel free to modify it for your needs!

\bibliographystyle{plain}
\bibliography{cite}

\end{document}\documentclass{article}
\usepackage[utf8]{inputenc}
\usepackage{booktabs}
\usepackage{amsmath}
\usepackage{natbib}

\title{My Document}
\author{Your Name}
\date{\today}

\begin{document}

\maketitle

\section{Introduction}
This is a sample document demonstrating various \LaTeX{} features including mathematical equations, tables, and formatting. You can cite papers using the citation search feature, and they will be automatically added to your \texttt{cite.bib} file.

\section{Mathematical Expressions}
Here are some example equations:

\subsection{Quadratic Formula}
The quadratic formula is given by:
\[x = \frac{-b \pm \sqrt{b^2 - 4ac}}{2a}\]

\subsection{Matrix Example}
A sample matrix:
\[
\begin{pmatrix}
a_{11} & a_{12} & a_{13} \\
a_{21} & a_{22} & a_{23} \\
a_{31} & a_{32} & a_{33}
\end{pmatrix}
\]

\section{Tables}
Below is a sample table showing data:

\begin{table}[h]
\centering
\begin{tabular}{@{}lrrr@{}}
\toprule
\textbf{Category} & \textbf{Value 1} & \textbf{Value 2} & \textbf{Total} \\
\midrule
Item A & 10.5 & 20.3 & 30.8 \\
Item B & 15.2 & 25.1 & 40.3 \\
Item C & 22.1 & 12.7 & 34.8 \\
\bottomrule
\end{tabular}
\caption{Sample Data Table}
\label{tab:sample}
\end{table}

\section{Lists}
\subsection{Bullet Points}
\begin{itemize}
  \item First item
  \item Second item
  \item Third item with sub-items:
    \begin{itemize}
      \item Sub-item 1
      \item Sub-item 2
    \end{itemize}
\end{itemize}

\subsection{Numbered List}
\begin{enumerate}
  \item First step
  \item Second step
  \item Third step
\end{enumerate}

\section{Citations and References}
When you use the citation search feature, references will be automatically managed. For example, you might cite a paper like this: \cite{example2024}. All citations will be automatically added to your \texttt{cite.bib} file.

\section{Conclusion}
This template demonstrates basic \LaTeX{} features including mathematical formulas, tables, and lists. The bibliography system is now set up to work seamlessly with Bibby AI's citation management. Feel free to modify it for your needs!

\bibliographystyle{plain}
\bibliography{cite}

\end{document}\documentclass{article}
\usepackage[utf8]{inputenc}
\usepackage{booktabs}
\usepackage{amsmath}
\usepackage{natbib}

\title{My Document}
\author{Your Name}
\date{\today}

\begin{document}

\maketitle

\section{Introduction}
This is a sample document demonstrating various \LaTeX{} features including mathematical equations, tables, and formatting. You can cite papers using the citation search feature, and they will be automatically added to your \texttt{cite.bib} file.

\section{Mathematical Expressions}
Here are some example equations:

\subsection{Quadratic Formula}
The quadratic formula is given by:
\[x = \frac{-b \pm \sqrt{b^2 - 4ac}}{2a}\]

\subsection{Matrix Example}
A sample matrix:
\[
\begin{pmatrix}
a_{11} & a_{12} & a_{13} \\
a_{21} & a_{22} & a_{23} \\
a_{31} & a_{32} & a_{33}
\end{pmatrix}
\]

\section{Tables}
Below is a sample table showing data:

\begin{table}[h]
\centering
\begin{tabular}{@{}lrrr@{}}
\toprule
\textbf{Category} & \textbf{Value 1} & \textbf{Value 2} & \textbf{Total} \\
\midrule
Item A & 10.5 & 20.3 & 30.8 \\
Item B & 15.2 & 25.1 & 40.3 \\
Item C & 22.1 & 12.7 & 34.8 \\
\bottomrule
\end{tabular}
\caption{Sample Data Table}
\label{tab:sample}
\end{table}

\section{Lists}
\subsection{Bullet Points}
\begin{itemize}
  \item First item
  \item Second item
  \item Third item with sub-items:
    \begin{itemize}
      \item Sub-item 1
      \item Sub-item 2
    \end{itemize}
\end{itemize}

\subsection{Numbered List}
\begin{enumerate}
  \item First step
  \item Second step
  \item Third step
\end{enumerate}

\section{Citations and References}
When you use the citation search feature, references will be automatically managed. For example, you might cite a paper like this: \cite{example2024}. All citations will be automatically added to your \texttt{cite.bib} file.

\section{Conclusion}
This template demonstrates basic \LaTeX{} features including mathematical formulas, tables, and lists. The bibliography system is now set up to work seamlessly with Bibby AI's citation management. Feel free to modify it for your needs!

\bibliographystyle{plain}
\bibliography{cite}

\end{document}\documentclass{article}
\usepackage[utf8]{inputenc}
\usepackage{booktabs}
\usepackage{amsmath}
\usepackage{natbib}

\title{My Document}
\author{Your Name}
\date{\today}

\begin{document}

\maketitle

\section{Introduction}
This is a sample document demonstrating various \LaTeX{} features including mathematical equations, tables, and formatting. You can cite papers using the citation search feature, and they will be automatically added to your \texttt{cite.bib} file.

\section{Mathematical Expressions}
Here are some example equations:

\subsection{Quadratic Formula}
The quadratic formula is given by:
\[x = \frac{-b \pm \sqrt{b^2 - 4ac}}{2a}\]

\subsection{Matrix Example}
A sample matrix:
\[
\begin{pmatrix}
a_{11} & a_{12} & a_{13} \\
a_{21} & a_{22} & a_{23} \\
a_{31} & a_{32} & a_{33}
\end{pmatrix}
\]

\section{Tables}
Below is a sample table showing data:

\begin{table}[h]
\centering
\begin{tabular}{@{}lrrr@{}}
\toprule
\textbf{Category} & \textbf{Value 1} & \textbf{Value 2} & \textbf{Total} \\
\midrule
Item A & 10.5 & 20.3 & 30.8 \\
Item B & 15.2 & 25.1 & 40.3 \\
Item C & 22.1 & 12.7 & 34.8 \\
\bottomrule
\end{tabular}
\caption{Sample Data Table}
\label{tab:sample}
\end{table}

\section{Lists}
\subsection{Bullet Points}
\begin{itemize}
  \item First item
  \item Second item
  \item Third item with sub-items:
    \begin{itemize}
      \item Sub-item 1
      \item Sub-item 2
    \end{itemize}
\end{itemize}

\subsection{Numbered List}
\begin{enumerate}
  \item First step
  \item Second step
  \item Third step
\end{enumerate}

\section{Citations and References}
When you use the citation search feature, references will be automatically managed. For example, you might cite a paper like this: \cite{example2024}. All citations will be automatically added to your \texttt{cite.bib} file.

\section{Conclusion}
This template demonstrates basic \LaTeX{} features including mathematical formulas, tables, and lists. The bibliography system is now set up to work seamlessly with Bibby AI's citation management. Feel free to modify it for your needs!

\bibliographystyle{plain}
\bibliography{cite}

\end{document}\documentclass{article}
\usepackage[utf8]{inputenc}
\usepackage{booktabs}
\usepackage{amsmath}
\usepackage{natbib}

\title{My Document}
\author{Your Name}
\date{\today}

\begin{document}

\maketitle

\section{Introduction}
This is a sample document demonstrating various \LaTeX{} features including mathematical equations, tables, and formatting. You can cite papers using the citation search feature, and they will be automatically added to your \texttt{cite.bib} file.

\section{Mathematical Expressions}
Here are some example equations:

\subsection{Quadratic Formula}
The quadratic formula is given by:
\[x = \frac{-b \pm \sqrt{b^2 - 4ac}}{2a}\]

\subsection{Matrix Example}
A sample matrix:
\[
\begin{pmatrix}
a_{11} & a_{12} & a_{13} \\
a_{21} & a_{22} & a_{23} \\
a_{31} & a_{32} & a_{33}
\end{pmatrix}
\]

\section{Tables}
Below is a sample table showing data:

\begin{table}[h]
\centering
\begin{tabular}{@{}lrrr@{}}
\toprule
\textbf{Category} & \textbf{Value 1} & \textbf{Value 2} & \textbf{Total} \\
\midrule
Item A & 10.5 & 20.3 & 30.8 \\
Item B & 15.2 & 25.1 & 40.3 \\
Item C & 22.1 & 12.7 & 34.8 \\
\bottomrule
\end{tabular}
\caption{Sample Data Table}
\label{tab:sample}
\end{table}

\section{Lists}
\subsection{Bullet Points}
\begin{itemize}
  \item First item
  \item Second item
  \item Third item with sub-items:
    \begin{itemize}
      \item Sub-item 1
      \item Sub-item 2
    \end{itemize}
\end{itemize}

\subsection{Numbered List}
\begin{enumerate}
  \item First step
  \item Second step
  \item Third step
\end{enumerate}

\section{Citations and References}
When you use the citation search feature, references will be automatically managed. For example, you might cite a paper like this: \cite{example2024}. All citations will be automatically added to your \texttt{cite.bib} file.

\section{Conclusion}
This template demonstrates basic \LaTeX{} features including mathematical formulas, tables, and lists. The bibliography system is now set up to work seamlessly with Bibby AI's citation management. Feel free to modify it for your needs!

\bibliographystyle{plain}
\bibliography{cite}

\end{document}\documentclass{article}
\usepackage[utf8]{inputenc}
\usepackage{booktabs}
\usepackage{amsmath}
\usepackage{natbib}

\title{My Document}
\author{Your Name}
\date{\today}

\begin{document}

\maketitle

\section{Introduction}
This is a sample document demonstrating various \LaTeX{} features including mathematical equations, tables, and formatting. You can cite papers using the citation search feature, and they will be automatically added to your \texttt{cite.bib} file.

\section{Mathematical Expressions}
Here are some example equations:

\subsection{Quadratic Formula}
The quadratic formula is given by:
\[x = \frac{-b \pm \sqrt{b^2 - 4ac}}{2a}\]

\subsection{Matrix Example}
A sample matrix:
\[
\begin{pmatrix}
a_{11} & a_{12} & a_{13} \\
a_{21} & a_{22} & a_{23} \\
a_{31} & a_{32} & a_{33}
\end{pmatrix}
\]

\section{Tables}
Below is a sample table showing data:

\begin{table}[h]
\centering
\begin{tabular}{@{}lrrr@{}}
\toprule
\textbf{Category} & \textbf{Value 1} & \textbf{Value 2} & \textbf{Total} \\
\midrule
Item A & 10.5 & 20.3 & 30.8 \\
Item B & 15.2 & 25.1 & 40.3 \\
Item C & 22.1 & 12.7 & 34.8 \\
\bottomrule
\end{tabular}
\caption{Sample Data Table}
\label{tab:sample}
\end{table}

\section{Lists}
\subsection{Bullet Points}
\begin{itemize}
  \item First item
  \item Second item
  \item Third item with sub-items:
    \begin{itemize}
      \item Sub-item 1
      \item Sub-item 2
    \end{itemize}
\end{itemize}

\subsection{Numbered List}
\begin{enumerate}
  \item First step
  \item Second step
  \item Third step
\end{enumerate}

\section{Citations and References}
When you use the citation search feature, references will be automatically managed. For example, you might cite a paper like this: \cite{example2024}. All citations will be automatically added to your \texttt{cite.bib} file.

\section{Conclusion}
This template demonstrates basic \LaTeX{} features including mathematical formulas, tables, and lists. The bibliography system is now set up to work seamlessly with Bibby AI's citation management. Feel free to modify it for your needs!

\bibliographystyle{plain}
\bibliography{cite}

\end{document}\documentclass{article}
\usepackage[utf8]{inputenc}
\usepackage{booktabs}
\usepackage{amsmath}
\usepackage{natbib}

\title{My Document}
\author{Your Name}
\date{\today}

\begin{document}

\maketitle

\section{Introduction}
This is a sample document demonstrating various \LaTeX{} features including mathematical equations, tables, and formatting. You can cite papers using the citation search feature, and they will be automatically added to your \texttt{cite.bib} file.

\section{Mathematical Expressions}
Here are some example equations:

\subsection{Quadratic Formula}
The quadratic formula is given by:
\[x = \frac{-b \pm \sqrt{b^2 - 4ac}}{2a}\]

\subsection{Matrix Example}
A sample matrix:
\[
\begin{pmatrix}
a_{11} & a_{12} & a_{13} \\
a_{21} & a_{22} & a_{23} \\
a_{31} & a_{32} & a_{33}
\end{pmatrix}
\]

\section{Tables}
Below is a sample table showing data:

\begin{table}[h]
\centering
\begin{tabular}{@{}lrrr@{}}
\toprule
\textbf{Category} & \textbf{Value 1} & \textbf{Value 2} & \textbf{Total} \\
\midrule
Item A & 10.5 & 20.3 & 30.8 \\
Item B & 15.2 & 25.1 & 40.3 \\
Item C & 22.1 & 12.7 & 34.8 \\
\bottomrule
\end{tabular}
\caption{Sample Data Table}
\label{tab:sample}
\end{table}

\section{Lists}
\subsection{Bullet Points}
\begin{itemize}
  \item First item
  \item Second item
  \item Third item with sub-items:
    \begin{itemize}
      \item Sub-item 1
      \item Sub-item 2
    \end{itemize}
\end{itemize}

\subsection{Numbered List}
\begin{enumerate}
  \item First step
  \item Second step
  \item Third step
\end{enumerate}

\section{Citations and References}
When you use the citation search feature, references will be automatically managed. For example, you might cite a paper like this: \cite{example2024}. All citations will be automatically added to your \texttt{cite.bib} file.

\section{Conclusion}
This template demonstrates basic \LaTeX{} features including mathematical formulas, tables, and lists. The bibliography system is now set up to work seamlessly with Bibby AI's citation management. Feel free to modify it for your needs!

\bibliographystyle{plain}
\bibliography{cite}

\end{document}\documentclass{article}
\usepackage[utf8]{inputenc}
\usepackage{booktabs}
\usepackage{amsmath}
\usepackage{natbib}

\title{My Document}
\author{Your Name}
\date{\today}

\begin{document}

\maketitle

\section{Introduction}
This is a sample document demonstrating various \LaTeX{} features including mathematical equations, tables, and formatting. You can cite papers using the citation search feature, and they will be automatically added to your \texttt{cite.bib} file.

\section{Mathematical Expressions}
Here are some example equations:

\subsection{Quadratic Formula}
The quadratic formula is given by:
\[x = \frac{-b \pm \sqrt{b^2 - 4ac}}{2a}\]

\subsection{Matrix Example}
A sample matrix:
\[
\begin{pmatrix}
a_{11} & a_{12} & a_{13} \\
a_{21} & a_{22} & a_{23} \\
a_{31} & a_{32} & a_{33}
\end{pmatrix}
\]

\section{Tables}
Below is a sample table showing data:

\begin{table}[h]
\centering
\begin{tabular}{@{}lrrr@{}}
\toprule
\textbf{Category} & \textbf{Value 1} & \textbf{Value 2} & \textbf{Total} \\
\midrule
Item A & 10.5 & 20.3 & 30.8 \\
Item B & 15.2 & 25.1 & 40.3 \\
Item C & 22.1 & 12.7 & 34.8 \\
\bottomrule
\end{tabular}
\caption{Sample Data Table}
\label{tab:sample}
\end{table}

\section{Lists}
\subsection{Bullet Points}
\begin{itemize}
  \item First item
  \item Second item
  \item Third item with sub-items:
    \begin{itemize}
      \item Sub-item 1
      \item Sub-item 2
    \end{itemize}
\end{itemize}

\subsection{Numbered List}
\begin{enumerate}
  \item First step
  \item Second step
  \item Third step
\end{enumerate}

\section{Citations and References}
When you use the citation search feature, references will be automatically managed. For example, you might cite a paper like this: \cite{example2024}. All citations will be automatically added to your \texttt{cite.bib} file.

\section{Conclusion}
This template demonstrates basic \LaTeX{} features including mathematical formulas, tables, and lists. The bibliography system is now set up to work seamlessly with Bibby AI's citation management. Feel free to modify it for your needs!

\bibliographystyle{plain}
\bibliography{cite}

\end{document}\documentclass{article}
\usepackage[utf8]{inputenc}
\usepackage{booktabs}
\usepackage{amsmath}
\usepackage{natbib}

\title{My Document}
\author{Your Name}
\date{\today}

\begin{document}

\maketitle

\section{Introduction}
This is a sample document demonstrating various \LaTeX{} features including mathematical equations, tables, and formatting. You can cite papers using the citation search feature, and they will be automatically added to your \texttt{cite.bib} file.

\section{Mathematical Expressions}
Here are some example equations:

\subsection{Quadratic Formula}
The quadratic formula is given by:
\[x = \frac{-b \pm \sqrt{b^2 - 4ac}}{2a}\]

\subsection{Matrix Example}
A sample matrix:
\[
\begin{pmatrix}
a_{11} & a_{12} & a_{13} \\
a_{21} & a_{22} & a_{23} \\
a_{31} & a_{32} & a_{33}
\end{pmatrix}
\]

\section{Tables}
Below is a sample table showing data:

\begin{table}[h]
\centering
\begin{tabular}{@{}lrrr@{}}
\toprule
\textbf{Category} & \textbf{Value 1} & \textbf{Value 2} & \textbf{Total} \\
\midrule
Item A & 10.5 & 20.3 & 30.8 \\
Item B & 15.2 & 25.1 & 40.3 \\
Item C & 22.1 & 12.7 & 34.8 \\
\bottomrule
\end{tabular}
\caption{Sample Data Table}
\label{tab:sample}
\end{table}

\section{Lists}
\subsection{Bullet Points}
\begin{itemize}
  \item First item
  \item Second item
  \item Third item with sub-items:
    \begin{itemize}
      \item Sub-item 1
      \item Sub-item 2
    \end{itemize}
\end{itemize}

\subsection{Numbered List}
\begin{enumerate}
  \item First step
  \item Second step
  \item Third step
\end{enumerate}

\section{Citations and References}
When you use the citation search feature, references will be automatically managed. For example, you might cite a paper like this: \cite{example2024}. All citations will be automatically added to your \texttt{cite.bib} file.

\section{Conclusion}
This template demonstrates basic \LaTeX{} features including mathematical formulas, tables, and lists. The bibliography system is now set up to work seamlessly with Bibby AI's citation management. Feel free to modify it for your needs!

\bibliographystyle{plain}
\bibliography{cite}

\end{document}\documentclass{article}
\usepackage[utf8]{inputenc}
\usepackage{booktabs}
\usepackage{amsmath}
\usepackage{natbib}

\title{My Document}
\author{Your Name}
\date{\today}

\begin{document}

\maketitle

\section{Introduction}
This is a sample document demonstrating various \LaTeX{} features including mathematical equations, tables, and formatting. You can cite papers using the citation search feature, and they will be automatically added to your \texttt{cite.bib} file.

\section{Mathematical Expressions}
Here are some example equations:

\subsection{Quadratic Formula}
The quadratic formula is given by:
\[x = \frac{-b \pm \sqrt{b^2 - 4ac}}{2a}\]

\subsection{Matrix Example}
A sample matrix:
\[
\begin{pmatrix}
a_{11} & a_{12} & a_{13} \\
a_{21} & a_{22} & a_{23} \\
a_{31} & a_{32} & a_{33}
\end{pmatrix}
\]

\section{Tables}
Below is a sample table showing data:

\begin{table}[h]
\centering
\begin{tabular}{@{}lrrr@{}}
\toprule
\textbf{Category} & \textbf{Value 1} & \textbf{Value 2} & \textbf{Total} \\
\midrule
Item A & 10.5 & 20.3 & 30.8 \\
Item B & 15.2 & 25.1 & 40.3 \\
Item C & 22.1 & 12.7 & 34.8 \\
\bottomrule
\end{tabular}
\caption{Sample Data Table}
\label{tab:sample}
\end{table}

\section{Lists}
\subsection{Bullet Points}
\begin{itemize}
  \item First item
  \item Second item
  \item Third item with sub-items:
    \begin{itemize}
      \item Sub-item 1
      \item Sub-item 2
    \end{itemize}
\end{itemize}

\subsection{Numbered List}
\begin{enumerate}
  \item First step
  \item Second step
  \item Third step
\end{enumerate}

\section{Citations and References}
When you use the citation search feature, references will be automatically managed. For example, you might cite a paper like this: \cite{example2024}. All citations will be automatically added to your \texttt{cite.bib} file.

\section{Conclusion}
This template demonstrates basic \LaTeX{} features including mathematical formulas, tables, and lists. The bibliography system is now set up to work seamlessly with Bibby AI's citation management. Feel free to modify it for your needs!

\bibliographystyle{plain}
\bibliography{cite}

\end{document}\documentclass{article}
\usepackage[utf8]{inputenc}
\usepackage{booktabs}
\usepackage{amsmath}
\usepackage{natbib}

\title{My Document}
\author{Your Name}
\date{\today}

\begin{document}

\maketitle

\section{Introduction}
This is a sample document demonstrating various \LaTeX{} features including mathematical equations, tables, and formatting. You can cite papers using the citation search feature, and they will be automatically added to your \texttt{cite.bib} file.

\section{Mathematical Expressions}
Here are some example equations:

\subsection{Quadratic Formula}
The quadratic formula is given by:
\[x = \frac{-b \pm \sqrt{b^2 - 4ac}}{2a}\]

\subsection{Matrix Example}
A sample matrix:
\[
\begin{pmatrix}
a_{11} & a_{12} & a_{13} \\
a_{21} & a_{22} & a_{23} \\
a_{31} & a_{32} & a_{33}
\end{pmatrix}
\]

\section{Tables}
Below is a sample table showing data:

\begin{table}[h]
\centering
\begin{tabular}{@{}lrrr@{}}
\toprule
\textbf{Category} & \textbf{Value 1} & \textbf{Value 2} & \textbf{Total} \\
\midrule
Item A & 10.5 & 20.3 & 30.8 \\
Item B & 15.2 & 25.1 & 40.3 \\
Item C & 22.1 & 12.7 & 34.8 \\
\bottomrule
\end{tabular}
\caption{Sample Data Table}
\label{tab:sample}
\end{table}

\section{Lists}
\subsection{Bullet Points}
\begin{itemize}
  \item First item
  \item Second item
  \item Third item with sub-items:
    \begin{itemize}
      \item Sub-item 1
      \item Sub-item 2
    \end{itemize}
\end{itemize}

\subsection{Numbered List}
\begin{enumerate}
  \item First step
  \item Second step
  \item Third step
\end{enumerate}

\section{Citations and References}
When you use the citation search feature, references will be automatically managed. For example, you might cite a paper like this: \cite{example2024}. All citations will be automatically added to your \texttt{cite.bib} file.

\section{Conclusion}
This template demonstrates basic \LaTeX{} features including mathematical formulas, tables, and lists. The bibliography system is now set up to work seamlessly with Bibby AI's citation management. Feel free to modify it for your needs!

\bibliographystyle{plain}
\bibliography{cite}

\end{document}\documentclass{article}
\usepackage[utf8]{inputenc}
\usepackage{booktabs}
\usepackage{amsmath}
\usepackage{natbib}

\title{My Document}
\author{Your Name}
\date{\today}

\begin{document}

\maketitle

\section{Introduction}
This is a sample document demonstrating various \LaTeX{} features including mathematical equations, tables, and formatting. You can cite papers using the citation search feature, and they will be automatically added to your \texttt{cite.bib} file.

\section{Mathematical Expressions}
Here are some example equations:

\subsection{Quadratic Formula}
The quadratic formula is given by:
\[x = \frac{-b \pm \sqrt{b^2 - 4ac}}{2a}\]

\subsection{Matrix Example}
A sample matrix:
\[
\begin{pmatrix}
a_{11} & a_{12} & a_{13} \\
a_{21} & a_{22} & a_{23} \\
a_{31} & a_{32} & a_{33}
\end{pmatrix}
\]

\section{Tables}
Below is a sample table showing data:

\begin{table}[h]
\centering
\begin{tabular}{@{}lrrr@{}}
\toprule
\textbf{Category} & \textbf{Value 1} & \textbf{Value 2} & \textbf{Total} \\
\midrule
Item A & 10.5 & 20.3 & 30.8 \\
Item B & 15.2 & 25.1 & 40.3 \\
Item C & 22.1 & 12.7 & 34.8 \\
\bottomrule
\end{tabular}
\caption{Sample Data Table}
\label{tab:sample}
\end{table}

\section{Lists}
\subsection{Bullet Points}
\begin{itemize}
  \item First item
  \item Second item
  \item Third item with sub-items:
    \begin{itemize}
      \item Sub-item 1
      \item Sub-item 2
    \end{itemize}
\end{itemize}

\subsection{Numbered List}
\begin{enumerate}
  \item First step
  \item Second step
  \item Third step
\end{enumerate}

\section{Citations and References}
When you use the citation search feature, references will be automatically managed. For example, you might cite a paper like this: \cite{example2024}. All citations will be automatically added to your \texttt{cite.bib} file.

\section{Conclusion}
This template demonstrates basic \LaTeX{} features including mathematical formulas, tables, and lists. The bibliography system is now set up to work seamlessly with Bibby AI's citation management. Feel free to modify it for your needs!

\bibliographystyle{plain}
\bibliography{cite}

\end{document}\documentclass{article}
\usepackage[utf8]{inputenc}
\usepackage{booktabs}
\usepackage{amsmath}
\usepackage{natbib}

\title{My Document}
\author{Your Name}
\date{\today}

\begin{document}

\maketitle

\section{Introduction}
This is a sample document demonstrating various \LaTeX{} features including mathematical equations, tables, and formatting. You can cite papers using the citation search feature, and they will be automatically added to your \texttt{cite.bib} file.

\section{Mathematical Expressions}
Here are some example equations:

\subsection{Quadratic Formula}
The quadratic formula is given by:
\[x = \frac{-b \pm \sqrt{b^2 - 4ac}}{2a}\]

\subsection{Matrix Example}
A sample matrix:
\[
\begin{pmatrix}
a_{11} & a_{12} & a_{13} \\
a_{21} & a_{22} & a_{23} \\
a_{31} & a_{32} & a_{33}
\end{pmatrix}
\]

\section{Tables}
Below is a sample table showing data:

\begin{table}[h]
\centering
\begin{tabular}{@{}lrrr@{}}
\toprule
\textbf{Category} & \textbf{Value 1} & \textbf{Value 2} & \textbf{Total} \\
\midrule
Item A & 10.5 & 20.3 & 30.8 \\
Item B & 15.2 & 25.1 & 40.3 \\
Item C & 22.1 & 12.7 & 34.8 \\
\bottomrule
\end{tabular}
\caption{Sample Data Table}
\label{tab:sample}
\end{table}

\section{Lists}
\subsection{Bullet Points}
\begin{itemize}
  \item First item
  \item Second item
  \item Third item with sub-items:
    \begin{itemize}
      \item Sub-item 1
      \item Sub-item 2
    \end{itemize}
\end{itemize}

\subsection{Numbered List}
\begin{enumerate}
  \item First step
  \item Second step
  \item Third step
\end{enumerate}

\section{Citations and References}
When you use the citation search feature, references will be automatically managed. For example, you might cite a paper like this: \cite{example2024}. All citations will be automatically added to your \texttt{cite.bib} file.

\section{Conclusion}
This template demonstrates basic \LaTeX{} features including mathematical formulas, tables, and lists. The bibliography system is now set up to work seamlessly with Bibby AI's citation management. Feel free to modify it for your needs!

\bibliographystyle{plain}
\bibliography{cite}

\end{document}\documentclass{article}
\usepackage[utf8]{inputenc}
\usepackage{booktabs}
\usepackage{amsmath}
\usepackage{natbib}

\title{My Document}
\author{Your Name}
\date{\today}

\begin{document}

\maketitle

\section{Introduction}
This is a sample document demonstrating various \LaTeX{} features including mathematical equations, tables, and formatting. You can cite papers using the citation search feature, and they will be automatically added to your \texttt{cite.bib} file.

\section{Mathematical Expressions}
Here are some example equations:

\subsection{Quadratic Formula}
The quadratic formula is given by:
\[x = \frac{-b \pm \sqrt{b^2 - 4ac}}{2a}\]

\subsection{Matrix Example}
A sample matrix:
\[
\begin{pmatrix}
a_{11} & a_{12} & a_{13} \\
a_{21} & a_{22} & a_{23} \\
a_{31} & a_{32} & a_{33}
\end{pmatrix}
\]

\section{Tables}
Below is a sample table showing data:

\begin{table}[h]
\centering
\begin{tabular}{@{}lrrr@{}}
\toprule
\textbf{Category} & \textbf{Value 1} & \textbf{Value 2} & \textbf{Total} \\
\midrule
Item A & 10.5 & 20.3 & 30.8 \\
Item B & 15.2 & 25.1 & 40.3 \\
Item C & 22.1 & 12.7 & 34.8 \\
\bottomrule
\end{tabular}
\caption{Sample Data Table}
\label{tab:sample}
\end{table}

\section{Lists}
\subsection{Bullet Points}
\begin{itemize}
  \item First item
  \item Second item
  \item Third item with sub-items:
    \begin{itemize}
      \item Sub-item 1
      \item Sub-item 2
    \end{itemize}
\end{itemize}

\subsection{Numbered List}
\begin{enumerate}
  \item First step
  \item Second step
  \item Third step
\end{enumerate}

\section{Citations and References}
When you use the citation search feature, references will be automatically managed. For example, you might cite a paper like this: \cite{example2024}. All citations will be automatically added to your \texttt{cite.bib} file.

\section{Conclusion}
This template demonstrates basic \LaTeX{} features including mathematical formulas, tables, and lists. The bibliography system is now set up to work seamlessly with Bibby AI's citation management. Feel free to modify it for your needs!

\bibliographystyle{plain}
\bibliography{cite}

\end{document}\documentclass{article}
\usepackage[utf8]{inputenc}
\usepackage{booktabs}
\usepackage{amsmath}
\usepackage{natbib}

\title{My Document}
\author{Your Name}
\date{\today}

\begin{document}

\maketitle

\section{Introduction}
This is a sample document demonstrating various \LaTeX{} features including mathematical equations, tables, and formatting. You can cite papers using the citation search feature, and they will be automatically added to your \texttt{cite.bib} file.

\section{Mathematical Expressions}
Here are some example equations:

\subsection{Quadratic Formula}
The quadratic formula is given by:
\[x = \frac{-b \pm \sqrt{b^2 - 4ac}}{2a}\]

\subsection{Matrix Example}
A sample matrix:
\[
\begin{pmatrix}
a_{11} & a_{12} & a_{13} \\
a_{21} & a_{22} & a_{23} \\
a_{31} & a_{32} & a_{33}
\end{pmatrix}
\]

\section{Tables}
Below is a sample table showing data:

\begin{table}[h]
\centering
\begin{tabular}{@{}lrrr@{}}
\toprule
\textbf{Category} & \textbf{Value 1} & \textbf{Value 2} & \textbf{Total} \\
\midrule
Item A & 10.5 & 20.3 & 30.8 \\
Item B & 15.2 & 25.1 & 40.3 \\
Item C & 22.1 & 12.7 & 34.8 \\
\bottomrule
\end{tabular}
\caption{Sample Data Table}
\label{tab:sample}
\end{table}

\section{Lists}
\subsection{Bullet Points}
\begin{itemize}
  \item First item
  \item Second item
  \item Third item with sub-items:
    \begin{itemize}
      \item Sub-item 1
      \item Sub-item 2
    \end{itemize}
\end{itemize}

\subsection{Numbered List}
\begin{enumerate}
  \item First step
  \item Second step
  \item Third step
\end{enumerate}

\section{Citations and References}
When you use the citation search feature, references will be automatically managed. For example, you might cite a paper like this: \cite{example2024}. All citations will be automatically added to your \texttt{cite.bib} file.

\section{Conclusion}
This template demonstrates basic \LaTeX{} features including mathematical formulas, tables, and lists. The bibliography system is now set up to work seamlessly with Bibby AI's citation management. Feel free to modify it for your needs!

\bibliographystyle{plain}
\bibliography{cite}

\end{document}\documentclass{article}
\usepackage[utf8]{inputenc}
\usepackage{booktabs}
\usepackage{amsmath}
\usepackage{natbib}

\title{My Document}
\author{Your Name}
\date{\today}

\begin{document}

\maketitle

\section{Introduction}
This is a sample document demonstrating various \LaTeX{} features including mathematical equations, tables, and formatting. You can cite papers using the citation search feature, and they will be automatically added to your \texttt{cite.bib} file.

\section{Mathematical Expressions}
Here are some example equations:

\subsection{Quadratic Formula}
The quadratic formula is given by:
\[x = \frac{-b \pm \sqrt{b^2 - 4ac}}{2a}\]

\subsection{Matrix Example}
A sample matrix:
\[
\begin{pmatrix}
a_{11} & a_{12} & a_{13} \\
a_{21} & a_{22} & a_{23} \\
a_{31} & a_{32} & a_{33}
\end{pmatrix}
\]

\section{Tables}
Below is a sample table showing data:

\begin{table}[h]
\centering
\begin{tabular}{@{}lrrr@{}}
\toprule
\textbf{Category} & \textbf{Value 1} & \textbf{Value 2} & \textbf{Total} \\
\midrule
Item A & 10.5 & 20.3 & 30.8 \\
Item B & 15.2 & 25.1 & 40.3 \\
Item C & 22.1 & 12.7 & 34.8 \\
\bottomrule
\end{tabular}
\caption{Sample Data Table}
\label{tab:sample}
\end{table}

\section{Lists}
\subsection{Bullet Points}
\begin{itemize}
  \item First item
  \item Second item
  \item Third item with sub-items:
    \begin{itemize}
      \item Sub-item 1
      \item Sub-item 2
    \end{itemize}
\end{itemize}

\subsection{Numbered List}
\begin{enumerate}
  \item First step
  \item Second step
  \item Third step
\end{enumerate}

\section{Citations and References}
When you use the citation search feature, references will be automatically managed. For example, you might cite a paper like this: \cite{example2024}. All citations will be automatically added to your \texttt{cite.bib} file.

\section{Conclusion}
This template demonstrates basic \LaTeX{} features including mathematical formulas, tables, and lists. The bibliography system is now set up to work seamlessly with Bibby AI's citation management. Feel free to modify it for your needs!

\bibliographystyle{plain}
\bibliography{cite}

\end{document}\documentclass{article}
\usepackage[utf8]{inputenc}
\usepackage{booktabs}
\usepackage{amsmath}
\usepackage{natbib}

\title{My Document}
\author{Your Name}
\date{\today}

\begin{document}

\maketitle

\section{Introduction}
This is a sample document demonstrating various \LaTeX{} features including mathematical equations, tables, and formatting. You can cite papers using the citation search feature, and they will be automatically added to your \texttt{cite.bib} file.

\section{Mathematical Expressions}
Here are some example equations:

\subsection{Quadratic Formula}
The quadratic formula is given by:
\[x = \frac{-b \pm \sqrt{b^2 - 4ac}}{2a}\]

\subsection{Matrix Example}
A sample matrix:
\[
\begin{pmatrix}
a_{11} & a_{12} & a_{13} \\
a_{21} & a_{22} & a_{23} \\
a_{31} & a_{32} & a_{33}
\end{pmatrix}
\]

\section{Tables}
Below is a sample table showing data:

\begin{table}[h]
\centering
\begin{tabular}{@{}lrrr@{}}
\toprule
\textbf{Category} & \textbf{Value 1} & \textbf{Value 2} & \textbf{Total} \\
\midrule
Item A & 10.5 & 20.3 & 30.8 \\
Item B & 15.2 & 25.1 & 40.3 \\
Item C & 22.1 & 12.7 & 34.8 \\
\bottomrule
\end{tabular}
\caption{Sample Data Table}
\label{tab:sample}
\end{table}

\section{Lists}
\subsection{Bullet Points}
\begin{itemize}
  \item First item
  \item Second item
  \item Third item with sub-items:
    \begin{itemize}
      \item Sub-item 1
      \item Sub-item 2
    \end{itemize}
\end{itemize}

\subsection{Numbered List}
\begin{enumerate}
  \item First step
  \item Second step
  \item Third step
\end{enumerate}

\section{Citations and References}
When you use the citation search feature, references will be automatically managed. For example, you might cite a paper like this: \cite{example2024}. All citations will be automatically added to your \texttt{cite.bib} file.

\section{Conclusion}
This template demonstrates basic \LaTeX{} features including mathematical formulas, tables, and lists. The bibliography system is now set up to work seamlessly with Bibby AI's citation management. Feel free to modify it for your needs!

\bibliographystyle{plain}
\bibliography{cite}

\end{document}\documentclass{article}
\usepackage[utf8]{inputenc}
\usepackage{booktabs}
\usepackage{amsmath}
\usepackage{natbib}

\title{My Document}
\author{Your Name}
\date{\today}

\begin{document}

\maketitle

\section{Introduction}
This is a sample document demonstrating various \LaTeX{} features including mathematical equations, tables, and formatting. You can cite papers using the citation search feature, and they will be automatically added to your \texttt{cite.bib} file.

\section{Mathematical Expressions}
Here are some example equations:

\subsection{Quadratic Formula}
The quadratic formula is given by:
\[x = \frac{-b \pm \sqrt{b^2 - 4ac}}{2a}\]

\subsection{Matrix Example}
A sample matrix:
\[
\begin{pmatrix}
a_{11} & a_{12} & a_{13} \\
a_{21} & a_{22} & a_{23} \\
a_{31} & a_{32} & a_{33}
\end{pmatrix}
\]

\section{Tables}
Below is a sample table showing data:

\begin{table}[h]
\centering
\begin{tabular}{@{}lrrr@{}}
\toprule
\textbf{Category} & \textbf{Value 1} & \textbf{Value 2} & \textbf{Total} \\
\midrule
Item A & 10.5 & 20.3 & 30.8 \\
Item B & 15.2 & 25.1 & 40.3 \\
Item C & 22.1 & 12.7 & 34.8 \\
\bottomrule
\end{tabular}
\caption{Sample Data Table}
\label{tab:sample}
\end{table}

\section{Lists}
\subsection{Bullet Points}
\begin{itemize}
  \item First item
  \item Second item
  \item Third item with sub-items:
    \begin{itemize}
      \item Sub-item 1
      \item Sub-item 2
    \end{itemize}
\end{itemize}

\subsection{Numbered List}
\begin{enumerate}
  \item First step
  \item Second step
  \item Third step
\end{enumerate}

\section{Citations and References}
When you use the citation search feature, references will be automatically managed. For example, you might cite a paper like this: \cite{example2024}. All citations will be automatically added to your \texttt{cite.bib} file.

\section{Conclusion}
This template demonstrates basic \LaTeX{} features including mathematical formulas, tables, and lists. The bibliography system is now set up to work seamlessly with Bibby AI's citation management. Feel free to modify it for your needs!

\bibliographystyle{plain}
\bibliography{cite}

\end{document}\documentclass{article}
\usepackage[utf8]{inputenc}
\usepackage{booktabs}
\usepackage{amsmath}
\usepackage{natbib}

\title{My Document}
\author{Your Name}
\date{\today}

\begin{document}

\maketitle

\section{Introduction}
This is a sample document demonstrating various \LaTeX{} features including mathematical equations, tables, and formatting. You can cite papers using the citation search feature, and they will be automatically added to your \texttt{cite.bib} file.

\section{Mathematical Expressions}
Here are some example equations:

\subsection{Quadratic Formula}
The quadratic formula is given by:
\[x = \frac{-b \pm \sqrt{b^2 - 4ac}}{2a}\]

\subsection{Matrix Example}
A sample matrix:
\[
\begin{pmatrix}
a_{11} & a_{12} & a_{13} \\
a_{21} & a_{22} & a_{23} \\
a_{31} & a_{32} & a_{33}
\end{pmatrix}
\]

\section{Tables}
Below is a sample table showing data:

\begin{table}[h]
\centering
\begin{tabular}{@{}lrrr@{}}
\toprule
\textbf{Category} & \textbf{Value 1} & \textbf{Value 2} & \textbf{Total} \\
\midrule
Item A & 10.5 & 20.3 & 30.8 \\
Item B & 15.2 & 25.1 & 40.3 \\
Item C & 22.1 & 12.7 & 34.8 \\
\bottomrule
\end{tabular}
\caption{Sample Data Table}
\label{tab:sample}
\end{table}

\section{Lists}
\subsection{Bullet Points}
\begin{itemize}
  \item First item
  \item Second item
  \item Third item with sub-items:
    \begin{itemize}
      \item Sub-item 1
      \item Sub-item 2
    \end{itemize}
\end{itemize}

\subsection{Numbered List}
\begin{enumerate}
  \item First step
  \item Second step
  \item Third step
\end{enumerate}

\section{Citations and References}
When you use the citation search feature, references will be automatically managed. For example, you might cite a paper like this: \cite{example2024}. All citations will be automatically added to your \texttt{cite.bib} file.

\section{Conclusion}
This template demonstrates basic \LaTeX{} features including mathematical formulas, tables, and lists. The bibliography system is now set up to work seamlessly with Bibby AI's citation management. Feel free to modify it for your needs!

\bibliographystyle{plain}
\bibliography{cite}

\end{document}\documentclass{article}
\usepackage[utf8]{inputenc}
\usepackage{booktabs}
\usepackage{amsmath}
\usepackage{natbib}

\title{My Document}
\author{Your Name}
\date{\today}

\begin{document}

\maketitle

\section{Introduction}
This is a sample document demonstrating various \LaTeX{} features including mathematical equations, tables, and formatting. You can cite papers using the citation search feature, and they will be automatically added to your \texttt{cite.bib} file.

\section{Mathematical Expressions}
Here are some example equations:

\subsection{Quadratic Formula}
The quadratic formula is given by:
\[x = \frac{-b \pm \sqrt{b^2 - 4ac}}{2a}\]

\subsection{Matrix Example}
A sample matrix:
\[
\begin{pmatrix}
a_{11} & a_{12} & a_{13} \\
a_{21} & a_{22} & a_{23} \\
a_{31} & a_{32} & a_{33}
\end{pmatrix}
\]

\section{Tables}
Below is a sample table showing data:

\begin{table}[h]
\centering
\begin{tabular}{@{}lrrr@{}}
\toprule
\textbf{Category} & \textbf{Value 1} & \textbf{Value 2} & \textbf{Total} \\
\midrule
Item A & 10.5 & 20.3 & 30.8 \\
Item B & 15.2 & 25.1 & 40.3 \\
Item C & 22.1 & 12.7 & 34.8 \\
\bottomrule
\end{tabular}
\caption{Sample Data Table}
\label{tab:sample}
\end{table}

\section{Lists}
\subsection{Bullet Points}
\begin{itemize}
  \item First item
  \item Second item
  \item Third item with sub-items:
    \begin{itemize}
      \item Sub-item 1
      \item Sub-item 2
    \end{itemize}
\end{itemize}

\subsection{Numbered List}
\begin{enumerate}
  \item First step
  \item Second step
  \item Third step
\end{enumerate}

\section{Citations and References}
When you use the citation search feature, references will be automatically managed. For example, you might cite a paper like this: \cite{example2024}. All citations will be automatically added to your \texttt{cite.bib} file.

\section{Conclusion}
This template demonstrates basic \LaTeX{} features including mathematical formulas, tables, and lists. The bibliography system is now set up to work seamlessly with Bibby AI's citation management. Feel free to modify it for your needs!

\bibliographystyle{plain}
\bibliography{cite}

\end{document}\documentclass{article}
\usepackage[utf8]{inputenc}
\usepackage{booktabs}
\usepackage{amsmath}
\usepackage{natbib}

\title{My Document}
\author{Your Name}
\date{\today}

\begin{document}

\maketitle

\section{Introduction}
This is a sample document demonstrating various \LaTeX{} features including mathematical equations, tables, and formatting. You can cite papers using the citation search feature, and they will be automatically added to your \texttt{cite.bib} file.

\section{Mathematical Expressions}
Here are some example equations:

\subsection{Quadratic Formula}
The quadratic formula is given by:
\[x = \frac{-b \pm \sqrt{b^2 - 4ac}}{2a}\]

\subsection{Matrix Example}
A sample matrix:
\[
\begin{pmatrix}
a_{11} & a_{12} & a_{13} \\
a_{21} & a_{22} & a_{23} \\
a_{31} & a_{32} & a_{33}
\end{pmatrix}
\]

\section{Tables}
Below is a sample table showing data:

\begin{table}[h]
\centering
\begin{tabular}{@{}lrrr@{}}
\toprule
\textbf{Category} & \textbf{Value 1} & \textbf{Value 2} & \textbf{Total} \\
\midrule
Item A & 10.5 & 20.3 & 30.8 \\
Item B & 15.2 & 25.1 & 40.3 \\
Item C & 22.1 & 12.7 & 34.8 \\
\bottomrule
\end{tabular}
\caption{Sample Data Table}
\label{tab:sample}
\end{table}

\section{Lists}
\subsection{Bullet Points}
\begin{itemize}
  \item First item
  \item Second item
  \item Third item with sub-items:
    \begin{itemize}
      \item Sub-item 1
      \item Sub-item 2
    \end{itemize}
\end{itemize}

\subsection{Numbered List}
\begin{enumerate}
  \item First step
  \item Second step
  \item Third step
\end{enumerate}

\section{Citations and References}
When you use the citation search feature, references will be automatically managed. For example, you might cite a paper like this: \cite{example2024}. All citations will be automatically added to your \texttt{cite.bib} file.

\section{Conclusion}
This template demonstrates basic \LaTeX{} features including mathematical formulas, tables, and lists. The bibliography system is now set up to work seamlessly with Bibby AI's citation management. Feel free to modify it for your needs!

\bibliographystyle{plain}
\bibliography{cite}

\end{document}\documentclass{article}
\usepackage[utf8]{inputenc}
\usepackage{booktabs}
\usepackage{amsmath}
\usepackage{natbib}

\title{My Document}
\author{Your Name}
\date{\today}

\begin{document}

\maketitle

\section{Introduction}
This is a sample document demonstrating various \LaTeX{} features including mathematical equations, tables, and formatting. You can cite papers using the citation search feature, and they will be automatically added to your \texttt{cite.bib} file.

\section{Mathematical Expressions}
Here are some example equations:

\subsection{Quadratic Formula}
The quadratic formula is given by:
\[x = \frac{-b \pm \sqrt{b^2 - 4ac}}{2a}\]

\subsection{Matrix Example}
A sample matrix:
\[
\begin{pmatrix}
a_{11} & a_{12} & a_{13} \\
a_{21} & a_{22} & a_{23} \\
a_{31} & a_{32} & a_{33}
\end{pmatrix}
\]

\section{Tables}
Below is a sample table showing data:

\begin{table}[h]
\centering
\begin{tabular}{@{}lrrr@{}}
\toprule
\textbf{Category} & \textbf{Value 1} & \textbf{Value 2} & \textbf{Total} \\
\midrule
Item A & 10.5 & 20.3 & 30.8 \\
Item B & 15.2 & 25.1 & 40.3 \\
Item C & 22.1 & 12.7 & 34.8 \\
\bottomrule
\end{tabular}
\caption{Sample Data Table}
\label{tab:sample}
\end{table}

\section{Lists}
\subsection{Bullet Points}
\begin{itemize}
  \item First item
  \item Second item
  \item Third item with sub-items:
    \begin{itemize}
      \item Sub-item 1
      \item Sub-item 2
    \end{itemize}
\end{itemize}

\subsection{Numbered List}
\begin{enumerate}
  \item First step
  \item Second step
  \item Third step
\end{enumerate}

\section{Citations and References}
When you use the citation search feature, references will be automatically managed. For example, you might cite a paper like this: \cite{example2024}. All citations will be automatically added to your \texttt{cite.bib} file.

\section{Conclusion}
This template demonstrates basic \LaTeX{} features including mathematical formulas, tables, and lists. The bibliography system is now set up to work seamlessly with Bibby AI's citation management. Feel free to modify it for your needs!

\bibliographystyle{plain}
\bibliography{cite}

\end{document}\documentclass{article}
\usepackage[utf8]{inputenc}
\usepackage{booktabs}
\usepackage{amsmath}
\usepackage{natbib}

\title{My Document}
\author{Your Name}
\date{\today}

\begin{document}

\maketitle

\section{Introduction}
This is a sample document demonstrating various \LaTeX{} features including mathematical equations, tables, and formatting. You can cite papers using the citation search feature, and they will be automatically added to your \texttt{cite.bib} file.

\section{Mathematical Expressions}
Here are some example equations:

\subsection{Quadratic Formula}
The quadratic formula is given by:
\[x = \frac{-b \pm \sqrt{b^2 - 4ac}}{2a}\]

\subsection{Matrix Example}
A sample matrix:
\[
\begin{pmatrix}
a_{11} & a_{12} & a_{13} \\
a_{21} & a_{22} & a_{23} \\
a_{31} & a_{32} & a_{33}
\end{pmatrix}
\]

\section{Tables}
Below is a sample table showing data:

\begin{table}[h]
\centering
\begin{tabular}{@{}lrrr@{}}
\toprule
\textbf{Category} & \textbf{Value 1} & \textbf{Value 2} & \textbf{Total} \\
\midrule
Item A & 10.5 & 20.3 & 30.8 \\
Item B & 15.2 & 25.1 & 40.3 \\
Item C & 22.1 & 12.7 & 34.8 \\
\bottomrule
\end{tabular}
\caption{Sample Data Table}
\label{tab:sample}
\end{table}

\section{Lists}
\subsection{Bullet Points}
\begin{itemize}
  \item First item
  \item Second item
  \item Third item with sub-items:
    \begin{itemize}
      \item Sub-item 1
      \item Sub-item 2
    \end{itemize}
\end{itemize}

\subsection{Numbered List}
\begin{enumerate}
  \item First step
  \item Second step
  \item Third step
\end{enumerate}

\section{Citations and References}
When you use the citation search feature, references will be automatically managed. For example, you might cite a paper like this: \cite{example2024}. All citations will be automatically added to your \texttt{cite.bib} file.

\section{Conclusion}
This template demonstrates basic \LaTeX{} features including mathematical formulas, tables, and lists. The bibliography system is now set up to work seamlessly with Bibby AI's citation management. Feel free to modify it for your needs!

\bibliographystyle{plain}
\bibliography{cite}

\end{document}\documentclass{article}
\usepackage[utf8]{inputenc}
\usepackage{booktabs}
\usepackage{amsmath}
\usepackage{natbib}

\title{My Document}
\author{Your Name}
\date{\today}

\begin{document}

\maketitle

\section{Introduction}
This is a sample document demonstrating various \LaTeX{} features including mathematical equations, tables, and formatting. You can cite papers using the citation search feature, and they will be automatically added to your \texttt{cite.bib} file.

\section{Mathematical Expressions}
Here are some example equations:

\subsection{Quadratic Formula}
The quadratic formula is given by:
\[x = \frac{-b \pm \sqrt{b^2 - 4ac}}{2a}\]

\subsection{Matrix Example}
A sample matrix:
\[
\begin{pmatrix}
a_{11} & a_{12} & a_{13} \\
a_{21} & a_{22} & a_{23} \\
a_{31} & a_{32} & a_{33}
\end{pmatrix}
\]

\section{Tables}
Below is a sample table showing data:

\begin{table}[h]
\centering
\begin{tabular}{@{}lrrr@{}}
\toprule
\textbf{Category} & \textbf{Value 1} & \textbf{Value 2} & \textbf{Total} \\
\midrule
Item A & 10.5 & 20.3 & 30.8 \\
Item B & 15.2 & 25.1 & 40.3 \\
Item C & 22.1 & 12.7 & 34.8 \\
\bottomrule
\end{tabular}
\caption{Sample Data Table}
\label{tab:sample}
\end{table}

\section{Lists}
\subsection{Bullet Points}
\begin{itemize}
  \item First item
  \item Second item
  \item Third item with sub-items:
    \begin{itemize}
      \item Sub-item 1
      \item Sub-item 2
    \end{itemize}
\end{itemize}

\subsection{Numbered List}
\begin{enumerate}
  \item First step
  \item Second step
  \item Third step
\end{enumerate}

\section{Citations and References}
When you use the citation search feature, references will be automatically managed. For example, you might cite a paper like this: \cite{example2024}. All citations will be automatically added to your \texttt{cite.bib} file.

\section{Conclusion}
This template demonstrates basic \LaTeX{} features including mathematical formulas, tables, and lists. The bibliography system is now set up to work seamlessly with Bibby AI's citation management. Feel free to modify it for your needs!

\bibliographystyle{plain}
\bibliography{cite}

\end{document}\documentclass{article}
\usepackage[utf8]{inputenc}
\usepackage{booktabs}
\usepackage{amsmath}
\usepackage{natbib}

\title{My Document}
\author{Your Name}
\date{\today}

\begin{document}

\maketitle

\section{Introduction}
This is a sample document demonstrating various \LaTeX{} features including mathematical equations, tables, and formatting. You can cite papers using the citation search feature, and they will be automatically added to your \texttt{cite.bib} file.

\section{Mathematical Expressions}
Here are some example equations:

\subsection{Quadratic Formula}
The quadratic formula is given by:
\[x = \frac{-b \pm \sqrt{b^2 - 4ac}}{2a}\]

\subsection{Matrix Example}
A sample matrix:
\[
\begin{pmatrix}
a_{11} & a_{12} & a_{13} \\
a_{21} & a_{22} & a_{23} \\
a_{31} & a_{32} & a_{33}
\end{pmatrix}
\]

\section{Tables}
Below is a sample table showing data:

\begin{table}[h]
\centering
\begin{tabular}{@{}lrrr@{}}
\toprule
\textbf{Category} & \textbf{Value 1} & \textbf{Value 2} & \textbf{Total} \\
\midrule
Item A & 10.5 & 20.3 & 30.8 \\
Item B & 15.2 & 25.1 & 40.3 \\
Item C & 22.1 & 12.7 & 34.8 \\
\bottomrule
\end{tabular}
\caption{Sample Data Table}
\label{tab:sample}
\end{table}

\section{Lists}
\subsection{Bullet Points}
\begin{itemize}
  \item First item
  \item Second item
  \item Third item with sub-items:
    \begin{itemize}
      \item Sub-item 1
      \item Sub-item 2
    \end{itemize}
\end{itemize}

\subsection{Numbered List}
\begin{enumerate}
  \item First step
  \item Second step
  \item Third step
\end{enumerate}

\section{Citations and References}
When you use the citation search feature, references will be automatically managed. For example, you might cite a paper like this: \cite{example2024}. All citations will be automatically added to your \texttt{cite.bib} file.

\section{Conclusion}
This template demonstrates basic \LaTeX{} features including mathematical formulas, tables, and lists. The bibliography system is now set up to work seamlessly with Bibby AI's citation management. Feel free to modify it for your needs!

\bibliographystyle{plain}
\bibliography{cite}

\end{document}\documentclass{article}
\usepackage[utf8]{inputenc}
\usepackage{booktabs}
\usepackage{amsmath}
\usepackage{natbib}

\title{My Document}
\author{Your Name}
\date{\today}

\begin{document}

\maketitle

\section{Introduction}
This is a sample document demonstrating various \LaTeX{} features including mathematical equations, tables, and formatting. You can cite papers using the citation search feature, and they will be automatically added to your \texttt{cite.bib} file.

\section{Mathematical Expressions}
Here are some example equations:

\subsection{Quadratic Formula}
The quadratic formula is given by:
\[x = \frac{-b \pm \sqrt{b^2 - 4ac}}{2a}\]

\subsection{Matrix Example}
A sample matrix:
\[
\begin{pmatrix}
a_{11} & a_{12} & a_{13} \\
a_{21} & a_{22} & a_{23} \\
a_{31} & a_{32} & a_{33}
\end{pmatrix}
\]

\section{Tables}
Below is a sample table showing data:

\begin{table}[h]
\centering
\begin{tabular}{@{}lrrr@{}}
\toprule
\textbf{Category} & \textbf{Value 1} & \textbf{Value 2} & \textbf{Total} \\
\midrule
Item A & 10.5 & 20.3 & 30.8 \\
Item B & 15.2 & 25.1 & 40.3 \\
Item C & 22.1 & 12.7 & 34.8 \\
\bottomrule
\end{tabular}
\caption{Sample Data Table}
\label{tab:sample}
\end{table}

\section{Lists}
\subsection{Bullet Points}
\begin{itemize}
  \item First item
  \item Second item
  \item Third item with sub-items:
    \begin{itemize}
      \item Sub-item 1
      \item Sub-item 2
    \end{itemize}
\end{itemize}

\subsection{Numbered List}
\begin{enumerate}
  \item First step
  \item Second step
  \item Third step
\end{enumerate}

\section{Citations and References}
When you use the citation search feature, references will be automatically managed. For example, you might cite a paper like this: \cite{example2024}. All citations will be automatically added to your \texttt{cite.bib} file.

\section{Conclusion}
This template demonstrates basic \LaTeX{} features including mathematical formulas, tables, and lists. The bibliography system is now set up to work seamlessly with Bibby AI's citation management. Feel free to modify it for your needs!

\bibliographystyle{plain}
\bibliography{cite}

\end{document}\documentclass{article}
\usepackage[utf8]{inputenc}
\usepackage{booktabs}
\usepackage{amsmath}
\usepackage{natbib}

\title{My Document}
\author{Your Name}
\date{\today}

\begin{document}

\maketitle

\section{Introduction}
This is a sample document demonstrating various \LaTeX{} features including mathematical equations, tables, and formatting. You can cite papers using the citation search feature, and they will be automatically added to your \texttt{cite.bib} file.

\section{Mathematical Expressions}
Here are some example equations:

\subsection{Quadratic Formula}
The quadratic formula is given by:
\[x = \frac{-b \pm \sqrt{b^2 - 4ac}}{2a}\]

\subsection{Matrix Example}
A sample matrix:
\[
\begin{pmatrix}
a_{11} & a_{12} & a_{13} \\
a_{21} & a_{22} & a_{23} \\
a_{31} & a_{32} & a_{33}
\end{pmatrix}
\]

\section{Tables}
Below is a sample table showing data:

\begin{table}[h]
\centering
\begin{tabular}{@{}lrrr@{}}
\toprule
\textbf{Category} & \textbf{Value 1} & \textbf{Value 2} & \textbf{Total} \\
\midrule
Item A & 10.5 & 20.3 & 30.8 \\
Item B & 15.2 & 25.1 & 40.3 \\
Item C & 22.1 & 12.7 & 34.8 \\
\bottomrule
\end{tabular}
\caption{Sample Data Table}
\label{tab:sample}
\end{table}

\section{Lists}
\subsection{Bullet Points}
\begin{itemize}
  \item First item
  \item Second item
  \item Third item with sub-items:
    \begin{itemize}
      \item Sub-item 1
      \item Sub-item 2
    \end{itemize}
\end{itemize}

\subsection{Numbered List}
\begin{enumerate}
  \item First step
  \item Second step
  \item Third step
\end{enumerate}

\section{Citations and References}
When you use the citation search feature, references will be automatically managed. For example, you might cite a paper like this: \cite{example2024}. All citations will be automatically added to your \texttt{cite.bib} file.

\section{Conclusion}
This template demonstrates basic \LaTeX{} features including mathematical formulas, tables, and lists. The bibliography system is now set up to work seamlessly with Bibby AI's citation management. Feel free to modify it for your needs!

\bibliographystyle{plain}
\bibliography{cite}

\end{document}\documentclass{article}
\usepackage[utf8]{inputenc}
\usepackage{booktabs}
\usepackage{amsmath}
\usepackage{natbib}

\title{My Document}
\author{Your Name}
\date{\today}

\begin{document}

\maketitle

\section{Introduction}
This is a sample document demonstrating various \LaTeX{} features including mathematical equations, tables, and formatting. You can cite papers using the citation search feature, and they will be automatically added to your \texttt{cite.bib} file.

\section{Mathematical Expressions}
Here are some example equations:

\subsection{Quadratic Formula}
The quadratic formula is given by:
\[x = \frac{-b \pm \sqrt{b^2 - 4ac}}{2a}\]

\subsection{Matrix Example}
A sample matrix:
\[
\begin{pmatrix}
a_{11} & a_{12} & a_{13} \\
a_{21} & a_{22} & a_{23} \\
a_{31} & a_{32} & a_{33}
\end{pmatrix}
\]

\section{Tables}
Below is a sample table showing data:

\begin{table}[h]
\centering
\begin{tabular}{@{}lrrr@{}}
\toprule
\textbf{Category} & \textbf{Value 1} & \textbf{Value 2} & \textbf{Total} \\
\midrule
Item A & 10.5 & 20.3 & 30.8 \\
Item B & 15.2 & 25.1 & 40.3 \\
Item C & 22.1 & 12.7 & 34.8 \\
\bottomrule
\end{tabular}
\caption{Sample Data Table}
\label{tab:sample}
\end{table}

\section{Lists}
\subsection{Bullet Points}
\begin{itemize}
  \item First item
  \item Second item
  \item Third item with sub-items:
    \begin{itemize}
      \item Sub-item 1
      \item Sub-item 2
    \end{itemize}
\end{itemize}

\subsection{Numbered List}
\begin{enumerate}
  \item First step
  \item Second step
  \item Third step
\end{enumerate}

\section{Citations and References}
When you use the citation search feature, references will be automatically managed. For example, you might cite a paper like this: \cite{example2024}. All citations will be automatically added to your \texttt{cite.bib} file.

\section{Conclusion}
This template demonstrates basic \LaTeX{} features including mathematical formulas, tables, and lists. The bibliography system is now set up to work seamlessly with Bibby AI's citation management. Feel free to modify it for your needs!

\bibliographystyle{plain}
\bibliography{cite}

\end{document}\documentclass{article}
\usepackage[utf8]{inputenc}
\usepackage{booktabs}
\usepackage{amsmath}
\usepackage{natbib}

\title{My Document}
\author{Your Name}
\date{\today}

\begin{document}

\maketitle

\section{Introduction}
This is a sample document demonstrating various \LaTeX{} features including mathematical equations, tables, and formatting. You can cite papers using the citation search feature, and they will be automatically added to your \texttt{cite.bib} file.

\section{Mathematical Expressions}
Here are some example equations:

\subsection{Quadratic Formula}
The quadratic formula is given by:
\[x = \frac{-b \pm \sqrt{b^2 - 4ac}}{2a}\]

\subsection{Matrix Example}
A sample matrix:
\[
\begin{pmatrix}
a_{11} & a_{12} & a_{13} \\
a_{21} & a_{22} & a_{23} \\
a_{31} & a_{32} & a_{33}
\end{pmatrix}
\]

\section{Tables}
Below is a sample table showing data:

\begin{table}[h]
\centering
\begin{tabular}{@{}lrrr@{}}
\toprule
\textbf{Category} & \textbf{Value 1} & \textbf{Value 2} & \textbf{Total} \\
\midrule
Item A & 10.5 & 20.3 & 30.8 \\
Item B & 15.2 & 25.1 & 40.3 \\
Item C & 22.1 & 12.7 & 34.8 \\
\bottomrule
\end{tabular}
\caption{Sample Data Table}
\label{tab:sample}
\end{table}

\section{Lists}
\subsection{Bullet Points}
\begin{itemize}
  \item First item
  \item Second item
  \item Third item with sub-items:
    \begin{itemize}
      \item Sub-item 1
      \item Sub-item 2
    \end{itemize}
\end{itemize}

\subsection{Numbered List}
\begin{enumerate}
  \item First step
  \item Second step
  \item Third step
\end{enumerate}

\section{Citations and References}
When you use the citation search feature, references will be automatically managed. For example, you might cite a paper like this: \cite{example2024}. All citations will be automatically added to your \texttt{cite.bib} file.

\section{Conclusion}
This template demonstrates basic \LaTeX{} features including mathematical formulas, tables, and lists. The bibliography system is now set up to work seamlessly with Bibby AI's citation management. Feel free to modify it for your needs!

\bibliographystyle{plain}
\bibliography{cite}

\end{document}\documentclass{article}
\usepackage[utf8]{inputenc}
\usepackage{booktabs}
\usepackage{amsmath}
\usepackage{natbib}

\title{My Document}
\author{Your Name}
\date{\today}

\begin{document}

\maketitle

\section{Introduction}
This is a sample document demonstrating various \LaTeX{} features including mathematical equations, tables, and formatting. You can cite papers using the citation search feature, and they will be automatically added to your \texttt{cite.bib} file.

\section{Mathematical Expressions}
Here are some example equations:

\subsection{Quadratic Formula}
The quadratic formula is given by:
\[x = \frac{-b \pm \sqrt{b^2 - 4ac}}{2a}\]

\subsection{Matrix Example}
A sample matrix:
\[
\begin{pmatrix}
a_{11} & a_{12} & a_{13} \\
a_{21} & a_{22} & a_{23} \\
a_{31} & a_{32} & a_{33}
\end{pmatrix}
\]

\section{Tables}
Below is a sample table showing data:

\begin{table}[h]
\centering
\begin{tabular}{@{}lrrr@{}}
\toprule
\textbf{Category} & \textbf{Value 1} & \textbf{Value 2} & \textbf{Total} \\
\midrule
Item A & 10.5 & 20.3 & 30.8 \\
Item B & 15.2 & 25.1 & 40.3 \\
Item C & 22.1 & 12.7 & 34.8 \\
\bottomrule
\end{tabular}
\caption{Sample Data Table}
\label{tab:sample}
\end{table}

\section{Lists}
\subsection{Bullet Points}
\begin{itemize}
  \item First item
  \item Second item
  \item Third item with sub-items:
    \begin{itemize}
      \item Sub-item 1
      \item Sub-item 2
    \end{itemize}
\end{itemize}

\subsection{Numbered List}
\begin{enumerate}
  \item First step
  \item Second step
  \item Third step
\end{enumerate}

\section{Citations and References}
When you use the citation search feature, references will be automatically managed. For example, you might cite a paper like this: \cite{example2024}. All citations will be automatically added to your \texttt{cite.bib} file.

\section{Conclusion}
This template demonstrates basic \LaTeX{} features including mathematical formulas, tables, and lists. The bibliography system is now set up to work seamlessly with Bibby AI's citation management. Feel free to modify it for your needs!

\bibliographystyle{plain}
\bibliography{cite}

\end{document}\documentclass{article}
\usepackage[utf8]{inputenc}
\usepackage{booktabs}
\usepackage{amsmath}
\usepackage{natbib}

\title{My Document}
\author{Your Name}
\date{\today}

\begin{document}

\maketitle

\section{Introduction}
This is a sample document demonstrating various \LaTeX{} features including mathematical equations, tables, and formatting. You can cite papers using the citation search feature, and they will be automatically added to your \texttt{cite.bib} file.

\section{Mathematical Expressions}
Here are some example equations:

\subsection{Quadratic Formula}
The quadratic formula is given by:
\[x = \frac{-b \pm \sqrt{b^2 - 4ac}}{2a}\]

\subsection{Matrix Example}
A sample matrix:
\[
\begin{pmatrix}
a_{11} & a_{12} & a_{13} \\
a_{21} & a_{22} & a_{23} \\
a_{31} & a_{32} & a_{33}
\end{pmatrix}
\]

\section{Tables}
Below is a sample table showing data:

\begin{table}[h]
\centering
\begin{tabular}{@{}lrrr@{}}
\toprule
\textbf{Category} & \textbf{Value 1} & \textbf{Value 2} & \textbf{Total} \\
\midrule
Item A & 10.5 & 20.3 & 30.8 \\
Item B & 15.2 & 25.1 & 40.3 \\
Item C & 22.1 & 12.7 & 34.8 \\
\bottomrule
\end{tabular}
\caption{Sample Data Table}
\label{tab:sample}
\end{table}

\section{Lists}
\subsection{Bullet Points}
\begin{itemize}
  \item First item
  \item Second item
  \item Third item with sub-items:
    \begin{itemize}
      \item Sub-item 1
      \item Sub-item 2
    \end{itemize}
\end{itemize}

\subsection{Numbered List}
\begin{enumerate}
  \item First step
  \item Second step
  \item Third step
\end{enumerate}

\section{Citations and References}
When you use the citation search feature, references will be automatically managed. For example, you might cite a paper like this: \cite{example2024}. All citations will be automatically added to your \texttt{cite.bib} file.

\section{Conclusion}
This template demonstrates basic \LaTeX{} features including mathematical formulas, tables, and lists. The bibliography system is now set up to work seamlessly with Bibby AI's citation management. Feel free to modify it for your needs!

\bibliographystyle{plain}
\bibliography{cite}

\end{document}\documentclass{article}
\usepackage[utf8]{inputenc}
\usepackage{booktabs}
\usepackage{amsmath}
\usepackage{natbib}

\title{My Document}
\author{Your Name}
\date{\today}

\begin{document}

\maketitle

\section{Introduction}
This is a sample document demonstrating various \LaTeX{} features including mathematical equations, tables, and formatting. You can cite papers using the citation search feature, and they will be automatically added to your \texttt{cite.bib} file.

\section{Mathematical Expressions}
Here are some example equations:

\subsection{Quadratic Formula}
The quadratic formula is given by:
\[x = \frac{-b \pm \sqrt{b^2 - 4ac}}{2a}\]

\subsection{Matrix Example}
A sample matrix:
\[
\begin{pmatrix}
a_{11} & a_{12} & a_{13} \\
a_{21} & a_{22} & a_{23} \\
a_{31} & a_{32} & a_{33}
\end{pmatrix}
\]

\section{Tables}
Below is a sample table showing data:

\begin{table}[h]
\centering
\begin{tabular}{@{}lrrr@{}}
\toprule
\textbf{Category} & \textbf{Value 1} & \textbf{Value 2} & \textbf{Total} \\
\midrule
Item A & 10.5 & 20.3 & 30.8 \\
Item B & 15.2 & 25.1 & 40.3 \\
Item C & 22.1 & 12.7 & 34.8 \\
\bottomrule
\end{tabular}
\caption{Sample Data Table}
\label{tab:sample}
\end{table}

\section{Lists}
\subsection{Bullet Points}
\begin{itemize}
  \item First item
  \item Second item
  \item Third item with sub-items:
    \begin{itemize}
      \item Sub-item 1
      \item Sub-item 2
    \end{itemize}
\end{itemize}

\subsection{Numbered List}
\begin{enumerate}
  \item First step
  \item Second step
  \item Third step
\end{enumerate}

\section{Citations and References}
When you use the citation search feature, references will be automatically managed. For example, you might cite a paper like this: \cite{example2024}. All citations will be automatically added to your \texttt{cite.bib} file.

\section{Conclusion}
This template demonstrates basic \LaTeX{} features including mathematical formulas, tables, and lists. The bibliography system is now set up to work seamlessly with Bibby AI's citation management. Feel free to modify it for your needs!

\bibliographystyle{plain}
\bibliography{cite}

\end{document}\documentclass{article}
\usepackage[utf8]{inputenc}
\usepackage{booktabs}
\usepackage{amsmath}
\usepackage{natbib}

\title{My Document}
\author{Your Name}
\date{\today}

\begin{document}

\maketitle

\section{Introduction}
This is a sample document demonstrating various \LaTeX{} features including mathematical equations, tables, and formatting. You can cite papers using the citation search feature, and they will be automatically added to your \texttt{cite.bib} file.

\section{Mathematical Expressions}
Here are some example equations:

\subsection{Quadratic Formula}
The quadratic formula is given by:
\[x = \frac{-b \pm \sqrt{b^2 - 4ac}}{2a}\]

\subsection{Matrix Example}
A sample matrix:
\[
\begin{pmatrix}
a_{11} & a_{12} & a_{13} \\
a_{21} & a_{22} & a_{23} \\
a_{31} & a_{32} & a_{33}
\end{pmatrix}
\]

\section{Tables}
Below is a sample table showing data:

\begin{table}[h]
\centering
\begin{tabular}{@{}lrrr@{}}
\toprule
\textbf{Category} & \textbf{Value 1} & \textbf{Value 2} & \textbf{Total} \\
\midrule
Item A & 10.5 & 20.3 & 30.8 \\
Item B & 15.2 & 25.1 & 40.3 \\
Item C & 22.1 & 12.7 & 34.8 \\
\bottomrule
\end{tabular}
\caption{Sample Data Table}
\label{tab:sample}
\end{table}

\section{Lists}
\subsection{Bullet Points}
\begin{itemize}
  \item First item
  \item Second item
  \item Third item with sub-items:
    \begin{itemize}
      \item Sub-item 1
      \item Sub-item 2
    \end{itemize}
\end{itemize}

\subsection{Numbered List}
\begin{enumerate}
  \item First step
  \item Second step
  \item Third step
\end{enumerate}

\section{Citations and References}
When you use the citation search feature, references will be automatically managed. For example, you might cite a paper like this: \cite{example2024}. All citations will be automatically added to your \texttt{cite.bib} file.

\section{Conclusion}
This template demonstrates basic \LaTeX{} features including mathematical formulas, tables, and lists. The bibliography system is now set up to work seamlessly with Bibby AI's citation management. Feel free to modify it for your needs!

\bibliographystyle{plain}
\bibliography{cite}

\end{document}\documentclass{article}
\usepackage[utf8]{inputenc}
\usepackage{booktabs}
\usepackage{amsmath}
\usepackage{natbib}

\title{My Document}
\author{Your Name}
\date{\today}

\begin{document}

\maketitle

\section{Introduction}
This is a sample document demonstrating various \LaTeX{} features including mathematical equations, tables, and formatting. You can cite papers using the citation search feature, and they will be automatically added to your \texttt{cite.bib} file.

\section{Mathematical Expressions}
Here are some example equations:

\subsection{Quadratic Formula}
The quadratic formula is given by:
\[x = \frac{-b \pm \sqrt{b^2 - 4ac}}{2a}\]

\subsection{Matrix Example}
A sample matrix:
\[
\begin{pmatrix}
a_{11} & a_{12} & a_{13} \\
a_{21} & a_{22} & a_{23} \\
a_{31} & a_{32} & a_{33}
\end{pmatrix}
\]

\section{Tables}
Below is a sample table showing data:

\begin{table}[h]
\centering
\begin{tabular}{@{}lrrr@{}}
\toprule
\textbf{Category} & \textbf{Value 1} & \textbf{Value 2} & \textbf{Total} \\
\midrule
Item A & 10.5 & 20.3 & 30.8 \\
Item B & 15.2 & 25.1 & 40.3 \\
Item C & 22.1 & 12.7 & 34.8 \\
\bottomrule
\end{tabular}
\caption{Sample Data Table}
\label{tab:sample}
\end{table}

\section{Lists}
\subsection{Bullet Points}
\begin{itemize}
  \item First item
  \item Second item
  \item Third item with sub-items:
    \begin{itemize}
      \item Sub-item 1
      \item Sub-item 2
    \end{itemize}
\end{itemize}

\subsection{Numbered List}
\begin{enumerate}
  \item First step
  \item Second step
  \item Third step
\end{enumerate}

\section{Citations and References}
When you use the citation search feature, references will be automatically managed. For example, you might cite a paper like this: \cite{example2024}. All citations will be automatically added to your \texttt{cite.bib} file.

\section{Conclusion}
This template demonstrates basic \LaTeX{} features including mathematical formulas, tables, and lists. The bibliography system is now set up to work seamlessly with Bibby AI's citation management. Feel free to modify it for your needs!

\bibliographystyle{plain}
\bibliography{cite}

\end{document}\documentclass{article}
\usepackage[utf8]{inputenc}
\usepackage{booktabs}
\usepackage{amsmath}
\usepackage{natbib}

\title{My Document}
\author{Your Name}
\date{\today}

\begin{document}

\maketitle

\section{Introduction}
This is a sample document demonstrating various \LaTeX{} features including mathematical equations, tables, and formatting. You can cite papers using the citation search feature, and they will be automatically added to your \texttt{cite.bib} file.

\section{Mathematical Expressions}
Here are some example equations:

\subsection{Quadratic Formula}
The quadratic formula is given by:
\[x = \frac{-b \pm \sqrt{b^2 - 4ac}}{2a}\]

\subsection{Matrix Example}
A sample matrix:
\[
\begin{pmatrix}
a_{11} & a_{12} & a_{13} \\
a_{21} & a_{22} & a_{23} \\
a_{31} & a_{32} & a_{33}
\end{pmatrix}
\]

\section{Tables}
Below is a sample table showing data:

\begin{table}[h]
\centering
\begin{tabular}{@{}lrrr@{}}
\toprule
\textbf{Category} & \textbf{Value 1} & \textbf{Value 2} & \textbf{Total} \\
\midrule
Item A & 10.5 & 20.3 & 30.8 \\
Item B & 15.2 & 25.1 & 40.3 \\
Item C & 22.1 & 12.7 & 34.8 \\
\bottomrule
\end{tabular}
\caption{Sample Data Table}
\label{tab:sample}
\end{table}

\section{Lists}
\subsection{Bullet Points}
\begin{itemize}
  \item First item
  \item Second item
  \item Third item with sub-items:
    \begin{itemize}
      \item Sub-item 1
      \item Sub-item 2
    \end{itemize}
\end{itemize}

\subsection{Numbered List}
\begin{enumerate}
  \item First step
  \item Second step
  \item Third step
\end{enumerate}

\section{Citations and References}
When you use the citation search feature, references will be automatically managed. For example, you might cite a paper like this: \cite{example2024}. All citations will be automatically added to your \texttt{cite.bib} file.

\section{Conclusion}
This template demonstrates basic \LaTeX{} features including mathematical formulas, tables, and lists. The bibliography system is now set up to work seamlessly with Bibby AI's citation management. Feel free to modify it for your needs!

\bibliographystyle{plain}
\bibliography{cite}

\end{document}\documentclass{article}
\usepackage[utf8]{inputenc}
\usepackage{booktabs}
\usepackage{amsmath}
\usepackage{natbib}

\title{My Document}
\author{Your Name}
\date{\today}

\begin{document}

\maketitle

\section{Introduction}
This is a sample document demonstrating various \LaTeX{} features including mathematical equations, tables, and formatting. You can cite papers using the citation search feature, and they will be automatically added to your \texttt{cite.bib} file.

\section{Mathematical Expressions}
Here are some example equations:

\subsection{Quadratic Formula}
The quadratic formula is given by:
\[x = \frac{-b \pm \sqrt{b^2 - 4ac}}{2a}\]

\subsection{Matrix Example}
A sample matrix:
\[
\begin{pmatrix}
a_{11} & a_{12} & a_{13} \\
a_{21} & a_{22} & a_{23} \\
a_{31} & a_{32} & a_{33}
\end{pmatrix}
\]

\section{Tables}
Below is a sample table showing data:

\begin{table}[h]
\centering
\begin{tabular}{@{}lrrr@{}}
\toprule
\textbf{Category} & \textbf{Value 1} & \textbf{Value 2} & \textbf{Total} \\
\midrule
Item A & 10.5 & 20.3 & 30.8 \\
Item B & 15.2 & 25.1 & 40.3 \\
Item C & 22.1 & 12.7 & 34.8 \\
\bottomrule
\end{tabular}
\caption{Sample Data Table}
\label{tab:sample}
\end{table}

\section{Lists}
\subsection{Bullet Points}
\begin{itemize}
  \item First item
  \item Second item
  \item Third item with sub-items:
    \begin{itemize}
      \item Sub-item 1
      \item Sub-item 2
    \end{itemize}
\end{itemize}

\subsection{Numbered List}
\begin{enumerate}
  \item First step
  \item Second step
  \item Third step
\end{enumerate}

\section{Citations and References}
When you use the citation search feature, references will be automatically managed. For example, you might cite a paper like this: \cite{example2024}. All citations will be automatically added to your \texttt{cite.bib} file.

\section{Conclusion}
This template demonstrates basic \LaTeX{} features including mathematical formulas, tables, and lists. The bibliography system is now set up to work seamlessly with Bibby AI's citation management. Feel free to modify it for your needs!

\bibliographystyle{plain}
\bibliography{cite}

\end{document}\documentclass{article}
\usepackage[utf8]{inputenc}
\usepackage{booktabs}
\usepackage{amsmath}
\usepackage{natbib}

\title{My Document}
\author{Your Name}
\date{\today}

\begin{document}

\maketitle

\section{Introduction}
This is a sample document demonstrating various \LaTeX{} features including mathematical equations, tables, and formatting. You can cite papers using the citation search feature, and they will be automatically added to your \texttt{cite.bib} file.

\section{Mathematical Expressions}
Here are some example equations:

\subsection{Quadratic Formula}
The quadratic formula is given by:
\[x = \frac{-b \pm \sqrt{b^2 - 4ac}}{2a}\]

\subsection{Matrix Example}
A sample matrix:
\[
\begin{pmatrix}
a_{11} & a_{12} & a_{13} \\
a_{21} & a_{22} & a_{23} \\
a_{31} & a_{32} & a_{33}
\end{pmatrix}
\]

\section{Tables}
Below is a sample table showing data:

\begin{table}[h]
\centering
\begin{tabular}{@{}lrrr@{}}
\toprule
\textbf{Category} & \textbf{Value 1} & \textbf{Value 2} & \textbf{Total} \\
\midrule
Item A & 10.5 & 20.3 & 30.8 \\
Item B & 15.2 & 25.1 & 40.3 \\
Item C & 22.1 & 12.7 & 34.8 \\
\bottomrule
\end{tabular}
\caption{Sample Data Table}
\label{tab:sample}
\end{table}

\section{Lists}
\subsection{Bullet Points}
\begin{itemize}
  \item First item
  \item Second item
  \item Third item with sub-items:
    \begin{itemize}
      \item Sub-item 1
      \item Sub-item 2
    \end{itemize}
\end{itemize}

\subsection{Numbered List}
\begin{enumerate}
  \item First step
  \item Second step
  \item Third step
\end{enumerate}

\section{Citations and References}
When you use the citation search feature, references will be automatically managed. For example, you might cite a paper like this: \cite{example2024}. All citations will be automatically added to your \texttt{cite.bib} file.

\section{Conclusion}
This template demonstrates basic \LaTeX{} features including mathematical formulas, tables, and lists. The bibliography system is now set up to work seamlessly with Bibby AI's citation management. Feel free to modify it for your needs!

\bibliographystyle{plain}
\bibliography{cite}

\end{document}\documentclass{article}
\usepackage[utf8]{inputenc}
\usepackage{booktabs}
\usepackage{amsmath}
\usepackage{natbib}

\title{My Document}
\author{Your Name}
\date{\today}

\begin{document}

\maketitle

\section{Introduction}
This is a sample document demonstrating various \LaTeX{} features including mathematical equations, tables, and formatting. You can cite papers using the citation search feature, and they will be automatically added to your \texttt{cite.bib} file.

\section{Mathematical Expressions}
Here are some example equations:

\subsection{Quadratic Formula}
The quadratic formula is given by:
\[x = \frac{-b \pm \sqrt{b^2 - 4ac}}{2a}\]

\subsection{Matrix Example}
A sample matrix:
\[
\begin{pmatrix}
a_{11} & a_{12} & a_{13} \\
a_{21} & a_{22} & a_{23} \\
a_{31} & a_{32} & a_{33}
\end{pmatrix}
\]

\section{Tables}
Below is a sample table showing data:

\begin{table}[h]
\centering
\begin{tabular}{@{}lrrr@{}}
\toprule
\textbf{Category} & \textbf{Value 1} & \textbf{Value 2} & \textbf{Total} \\
\midrule
Item A & 10.5 & 20.3 & 30.8 \\
Item B & 15.2 & 25.1 & 40.3 \\
Item C & 22.1 & 12.7 & 34.8 \\
\bottomrule
\end{tabular}
\caption{Sample Data Table}
\label{tab:sample}
\end{table}

\section{Lists}
\subsection{Bullet Points}
\begin{itemize}
  \item First item
  \item Second item
  \item Third item with sub-items:
    \begin{itemize}
      \item Sub-item 1
      \item Sub-item 2
    \end{itemize}
\end{itemize}

\subsection{Numbered List}
\begin{enumerate}
  \item First step
  \item Second step
  \item Third step
\end{enumerate}

\section{Citations and References}
When you use the citation search feature, references will be automatically managed. For example, you might cite a paper like this: \cite{example2024}. All citations will be automatically added to your \texttt{cite.bib} file.

\section{Conclusion}
This template demonstrates basic \LaTeX{} features including mathematical formulas, tables, and lists. The bibliography system is now set up to work seamlessly with Bibby AI's citation management. Feel free to modify it for your needs!

\bibliographystyle{plain}
\bibliography{cite}

\end{document}\documentclass{article}
\usepackage[utf8]{inputenc}
\usepackage{booktabs}
\usepackage{amsmath}
\usepackage{natbib}

\title{My Document}
\author{Your Name}
\date{\today}

\begin{document}

\maketitle

\section{Introduction}
This is a sample document demonstrating various \LaTeX{} features including mathematical equations, tables, and formatting. You can cite papers using the citation search feature, and they will be automatically added to your \texttt{cite.bib} file.

\section{Mathematical Expressions}
Here are some example equations:

\subsection{Quadratic Formula}
The quadratic formula is given by:
\[x = \frac{-b \pm \sqrt{b^2 - 4ac}}{2a}\]

\subsection{Matrix Example}
A sample matrix:
\[
\begin{pmatrix}
a_{11} & a_{12} & a_{13} \\
a_{21} & a_{22} & a_{23} \\
a_{31} & a_{32} & a_{33}
\end{pmatrix}
\]

\section{Tables}
Below is a sample table showing data:

\begin{table}[h]
\centering
\begin{tabular}{@{}lrrr@{}}
\toprule
\textbf{Category} & \textbf{Value 1} & \textbf{Value 2} & \textbf{Total} \\
\midrule
Item A & 10.5 & 20.3 & 30.8 \\
Item B & 15.2 & 25.1 & 40.3 \\
Item C & 22.1 & 12.7 & 34.8 \\
\bottomrule
\end{tabular}
\caption{Sample Data Table}
\label{tab:sample}
\end{table}

\section{Lists}
\subsection{Bullet Points}
\begin{itemize}
  \item First item
  \item Second item
  \item Third item with sub-items:
    \begin{itemize}
      \item Sub-item 1
      \item Sub-item 2
    \end{itemize}
\end{itemize}

\subsection{Numbered List}
\begin{enumerate}
  \item First step
  \item Second step
  \item Third step
\end{enumerate}

\section{Citations and References}
When you use the citation search feature, references will be automatically managed. For example, you might cite a paper like this: \cite{example2024}. All citations will be automatically added to your \texttt{cite.bib} file.

\section{Conclusion}
This template demonstrates basic \LaTeX{} features including mathematical formulas, tables, and lists. The bibliography system is now set up to work seamlessly with Bibby AI's citation management. Feel free to modify it for your needs!

\bibliographystyle{plain}
\bibliography{cite}

\end{document}\documentclass{article}
\usepackage[utf8]{inputenc}
\usepackage{booktabs}
\usepackage{amsmath}
\usepackage{natbib}

\title{My Document}
\author{Your Name}
\date{\today}

\begin{document}

\maketitle

\section{Introduction}
This is a sample document demonstrating various \LaTeX{} features including mathematical equations, tables, and formatting. You can cite papers using the citation search feature, and they will be automatically added to your \texttt{cite.bib} file.

\section{Mathematical Expressions}
Here are some example equations:

\subsection{Quadratic Formula}
The quadratic formula is given by:
\[x = \frac{-b \pm \sqrt{b^2 - 4ac}}{2a}\]

\subsection{Matrix Example}
A sample matrix:
\[
\begin{pmatrix}
a_{11} & a_{12} & a_{13} \\
a_{21} & a_{22} & a_{23} \\
a_{31} & a_{32} & a_{33}
\end{pmatrix}
\]

\section{Tables}
Below is a sample table showing data:

\begin{table}[h]
\centering
\begin{tabular}{@{}lrrr@{}}
\toprule
\textbf{Category} & \textbf{Value 1} & \textbf{Value 2} & \textbf{Total} \\
\midrule
Item A & 10.5 & 20.3 & 30.8 \\
Item B & 15.2 & 25.1 & 40.3 \\
Item C & 22.1 & 12.7 & 34.8 \\
\bottomrule
\end{tabular}
\caption{Sample Data Table}
\label{tab:sample}
\end{table}

\section{Lists}
\subsection{Bullet Points}
\begin{itemize}
  \item First item
  \item Second item
  \item Third item with sub-items:
    \begin{itemize}
      \item Sub-item 1
      \item Sub-item 2
    \end{itemize}
\end{itemize}

\subsection{Numbered List}
\begin{enumerate}
  \item First step
  \item Second step
  \item Third step
\end{enumerate}

\section{Citations and References}
When you use the citation search feature, references will be automatically managed. For example, you might cite a paper like this: \cite{example2024}. All citations will be automatically added to your \texttt{cite.bib} file.

\section{Conclusion}
This template demonstrates basic \LaTeX{} features including mathematical formulas, tables, and lists. The bibliography system is now set up to work seamlessly with Bibby AI's citation management. Feel free to modify it for your needs!

\bibliographystyle{plain}
\bibliography{cite}

\end{document}\documentclass{article}
\usepackage[utf8]{inputenc}
\usepackage{booktabs}
\usepackage{amsmath}
\usepackage{natbib}

\title{My Document}
\author{Your Name}
\date{\today}

\begin{document}

\maketitle

\section{Introduction}
This is a sample document demonstrating various \LaTeX{} features including mathematical equations, tables, and formatting. You can cite papers using the citation search feature, and they will be automatically added to your \texttt{cite.bib} file.

\section{Mathematical Expressions}
Here are some example equations:

\subsection{Quadratic Formula}
The quadratic formula is given by:
\[x = \frac{-b \pm \sqrt{b^2 - 4ac}}{2a}\]

\subsection{Matrix Example}
A sample matrix:
\[
\begin{pmatrix}
a_{11} & a_{12} & a_{13} \\
a_{21} & a_{22} & a_{23} \\
a_{31} & a_{32} & a_{33}
\end{pmatrix}
\]

\section{Tables}
Below is a sample table showing data:

\begin{table}[h]
\centering
\begin{tabular}{@{}lrrr@{}}
\toprule
\textbf{Category} & \textbf{Value 1} & \textbf{Value 2} & \textbf{Total} \\
\midrule
Item A & 10.5 & 20.3 & 30.8 \\
Item B & 15.2 & 25.1 & 40.3 \\
Item C & 22.1 & 12.7 & 34.8 \\
\bottomrule
\end{tabular}
\caption{Sample Data Table}
\label{tab:sample}
\end{table}

\section{Lists}
\subsection{Bullet Points}
\begin{itemize}
  \item First item
  \item Second item
  \item Third item with sub-items:
    \begin{itemize}
      \item Sub-item 1
      \item Sub-item 2
    \end{itemize}
\end{itemize}

\subsection{Numbered List}
\begin{enumerate}
  \item First step
  \item Second step
  \item Third step
\end{enumerate}

\section{Citations and References}
When you use the citation search feature, references will be automatically managed. For example, you might cite a paper like this: \cite{example2024}. All citations will be automatically added to your \texttt{cite.bib} file.

\section{Conclusion}
This template demonstrates basic \LaTeX{} features including mathematical formulas, tables, and lists. The bibliography system is now set up to work seamlessly with Bibby AI's citation management. Feel free to modify it for your needs!

\bibliographystyle{plain}
\bibliography{cite}

\end{document}\documentclass{article}
\usepackage[utf8]{inputenc}
\usepackage{booktabs}
\usepackage{amsmath}
\usepackage{natbib}

\title{My Document}
\author{Your Name}
\date{\today}

\begin{document}

\maketitle

\section{Introduction}
This is a sample document demonstrating various \LaTeX{} features including mathematical equations, tables, and formatting. You can cite papers using the citation search feature, and they will be automatically added to your \texttt{cite.bib} file.

\section{Mathematical Expressions}
Here are some example equations:

\subsection{Quadratic Formula}
The quadratic formula is given by:
\[x = \frac{-b \pm \sqrt{b^2 - 4ac}}{2a}\]

\subsection{Matrix Example}
A sample matrix:
\[
\begin{pmatrix}
a_{11} & a_{12} & a_{13} \\
a_{21} & a_{22} & a_{23} \\
a_{31} & a_{32} & a_{33}
\end{pmatrix}
\]

\section{Tables}
Below is a sample table showing data:

\begin{table}[h]
\centering
\begin{tabular}{@{}lrrr@{}}
\toprule
\textbf{Category} & \textbf{Value 1} & \textbf{Value 2} & \textbf{Total} \\
\midrule
Item A & 10.5 & 20.3 & 30.8 \\
Item B & 15.2 & 25.1 & 40.3 \\
Item C & 22.1 & 12.7 & 34.8 \\
\bottomrule
\end{tabular}
\caption{Sample Data Table}
\label{tab:sample}
\end{table}

\section{Lists}
\subsection{Bullet Points}
\begin{itemize}
  \item First item
  \item Second item
  \item Third item with sub-items:
    \begin{itemize}
      \item Sub-item 1
      \item Sub-item 2
    \end{itemize}
\end{itemize}

\subsection{Numbered List}
\begin{enumerate}
  \item First step
  \item Second step
  \item Third step
\end{enumerate}

\section{Citations and References}
When you use the citation search feature, references will be automatically managed. For example, you might cite a paper like this: \cite{example2024}. All citations will be automatically added to your \texttt{cite.bib} file.

\section{Conclusion}
This template demonstrates basic \LaTeX{} features including mathematical formulas, tables, and lists. The bibliography system is now set up to work seamlessly with Bibby AI's citation management. Feel free to modify it for your needs!

\bibliographystyle{plain}
\bibliography{cite}

\end{document}\documentclass{article}
\usepackage[utf8]{inputenc}
\usepackage{booktabs}
\usepackage{amsmath}
\usepackage{natbib}

\title{My Document}
\author{Your Name}
\date{\today}

\begin{document}

\maketitle

\section{Introduction}
This is a sample document demonstrating various \LaTeX{} features including mathematical equations, tables, and formatting. You can cite papers using the citation search feature, and they will be automatically added to your \texttt{cite.bib} file.

\section{Mathematical Expressions}
Here are some example equations:

\subsection{Quadratic Formula}
The quadratic formula is given by:
\[x = \frac{-b \pm \sqrt{b^2 - 4ac}}{2a}\]

\subsection{Matrix Example}
A sample matrix:
\[
\begin{pmatrix}
a_{11} & a_{12} & a_{13} \\
a_{21} & a_{22} & a_{23} \\
a_{31} & a_{32} & a_{33}
\end{pmatrix}
\]

\section{Tables}
Below is a sample table showing data:

\begin{table}[h]
\centering
\begin{tabular}{@{}lrrr@{}}
\toprule
\textbf{Category} & \textbf{Value 1} & \textbf{Value 2} & \textbf{Total} \\
\midrule
Item A & 10.5 & 20.3 & 30.8 \\
Item B & 15.2 & 25.1 & 40.3 \\
Item C & 22.1 & 12.7 & 34.8 \\
\bottomrule
\end{tabular}
\caption{Sample Data Table}
\label{tab:sample}
\end{table}

\section{Lists}
\subsection{Bullet Points}
\begin{itemize}
  \item First item
  \item Second item
  \item Third item with sub-items:
    \begin{itemize}
      \item Sub-item 1
      \item Sub-item 2
    \end{itemize}
\end{itemize}

\subsection{Numbered List}
\begin{enumerate}
  \item First step
  \item Second step
  \item Third step
\end{enumerate}

\section{Citations and References}
When you use the citation search feature, references will be automatically managed. For example, you might cite a paper like this: \cite{example2024}. All citations will be automatically added to your \texttt{cite.bib} file.

\section{Conclusion}
This template demonstrates basic \LaTeX{} features including mathematical formulas, tables, and lists. The bibliography system is now set up to work seamlessly with Bibby AI's citation management. Feel free to modify it for your needs!

\bibliographystyle{plain}
\bibliography{cite}

\end{document}\documentclass{article}
\usepackage[utf8]{inputenc}
\usepackage{booktabs}
\usepackage{amsmath}
\usepackage{natbib}

\title{My Document}
\author{Your Name}
\date{\today}

\begin{document}

\maketitle

\section{Introduction}
This is a sample document demonstrating various \LaTeX{} features including mathematical equations, tables, and formatting. You can cite papers using the citation search feature, and they will be automatically added to your \texttt{cite.bib} file.

\section{Mathematical Expressions}
Here are some example equations:

\subsection{Quadratic Formula}
The quadratic formula is given by:
\[x = \frac{-b \pm \sqrt{b^2 - 4ac}}{2a}\]

\subsection{Matrix Example}
A sample matrix:
\[
\begin{pmatrix}
a_{11} & a_{12} & a_{13} \\
a_{21} & a_{22} & a_{23} \\
a_{31} & a_{32} & a_{33}
\end{pmatrix}
\]

\section{Tables}
Below is a sample table showing data:

\begin{table}[h]
\centering
\begin{tabular}{@{}lrrr@{}}
\toprule
\textbf{Category} & \textbf{Value 1} & \textbf{Value 2} & \textbf{Total} \\
\midrule
Item A & 10.5 & 20.3 & 30.8 \\
Item B & 15.2 & 25.1 & 40.3 \\
Item C & 22.1 & 12.7 & 34.8 \\
\bottomrule
\end{tabular}
\caption{Sample Data Table}
\label{tab:sample}
\end{table}

\section{Lists}
\subsection{Bullet Points}
\begin{itemize}
  \item First item
  \item Second item
  \item Third item with sub-items:
    \begin{itemize}
      \item Sub-item 1
      \item Sub-item 2
    \end{itemize}
\end{itemize}

\subsection{Numbered List}
\begin{enumerate}
  \item First step
  \item Second step
  \item Third step
\end{enumerate}

\section{Citations and References}
When you use the citation search feature, references will be automatically managed. For example, you might cite a paper like this: \cite{example2024}. All citations will be automatically added to your \texttt{cite.bib} file.

\section{Conclusion}
This template demonstrates basic \LaTeX{} features including mathematical formulas, tables, and lists. The bibliography system is now set up to work seamlessly with Bibby AI's citation management. Feel free to modify it for your needs!

\bibliographystyle{plain}
\bibliography{cite}

\end{document}\documentclass{article}
\usepackage[utf8]{inputenc}
\usepackage{booktabs}
\usepackage{amsmath}
\usepackage{natbib}

\title{My Document}
\author{Your Name}
\date{\today}

\begin{document}

\maketitle

\section{Introduction}
This is a sample document demonstrating various \LaTeX{} features including mathematical equations, tables, and formatting. You can cite papers using the citation search feature, and they will be automatically added to your \texttt{cite.bib} file.

\section{Mathematical Expressions}
Here are some example equations:

\subsection{Quadratic Formula}
The quadratic formula is given by:
\[x = \frac{-b \pm \sqrt{b^2 - 4ac}}{2a}\]

\subsection{Matrix Example}
A sample matrix:
\[
\begin{pmatrix}
a_{11} & a_{12} & a_{13} \\
a_{21} & a_{22} & a_{23} \\
a_{31} & a_{32} & a_{33}
\end{pmatrix}
\]

\section{Tables}
Below is a sample table showing data:

\begin{table}[h]
\centering
\begin{tabular}{@{}lrrr@{}}
\toprule
\textbf{Category} & \textbf{Value 1} & \textbf{Value 2} & \textbf{Total} \\
\midrule
Item A & 10.5 & 20.3 & 30.8 \\
Item B & 15.2 & 25.1 & 40.3 \\
Item C & 22.1 & 12.7 & 34.8 \\
\bottomrule
\end{tabular}
\caption{Sample Data Table}
\label{tab:sample}
\end{table}

\section{Lists}
\subsection{Bullet Points}
\begin{itemize}
  \item First item
  \item Second item
  \item Third item with sub-items:
    \begin{itemize}
      \item Sub-item 1
      \item Sub-item 2
    \end{itemize}
\end{itemize}

\subsection{Numbered List}
\begin{enumerate}
  \item First step
  \item Second step
  \item Third step
\end{enumerate}

\section{Citations and References}
When you use the citation search feature, references will be automatically managed. For example, you might cite a paper like this: \cite{example2024}. All citations will be automatically added to your \texttt{cite.bib} file.

\section{Conclusion}
This template demonstrates basic \LaTeX{} features including mathematical formulas, tables, and lists. The bibliography system is now set up to work seamlessly with Bibby AI's citation management. Feel free to modify it for your needs!

\bibliographystyle{plain}
\bibliography{cite}

\end{document}\documentclass{article}
\usepackage[utf8]{inputenc}
\usepackage{booktabs}
\usepackage{amsmath}
\usepackage{natbib}

\title{My Document}
\author{Your Name}
\date{\today}

\begin{document}

\maketitle

\section{Introduction}
This is a sample document demonstrating various \LaTeX{} features including mathematical equations, tables, and formatting. You can cite papers using the citation search feature, and they will be automatically added to your \texttt{cite.bib} file.

\section{Mathematical Expressions}
Here are some example equations:

\subsection{Quadratic Formula}
The quadratic formula is given by:
\[x = \frac{-b \pm \sqrt{b^2 - 4ac}}{2a}\]

\subsection{Matrix Example}
A sample matrix:
\[
\begin{pmatrix}
a_{11} & a_{12} & a_{13} \\
a_{21} & a_{22} & a_{23} \\
a_{31} & a_{32} & a_{33}
\end{pmatrix}
\]

\section{Tables}
Below is a sample table showing data:

\begin{table}[h]
\centering
\begin{tabular}{@{}lrrr@{}}
\toprule
\textbf{Category} & \textbf{Value 1} & \textbf{Value 2} & \textbf{Total} \\
\midrule
Item A & 10.5 & 20.3 & 30.8 \\
Item B & 15.2 & 25.1 & 40.3 \\
Item C & 22.1 & 12.7 & 34.8 \\
\bottomrule
\end{tabular}
\caption{Sample Data Table}
\label{tab:sample}
\end{table}

\section{Lists}
\subsection{Bullet Points}
\begin{itemize}
  \item First item
  \item Second item
  \item Third item with sub-items:
    \begin{itemize}
      \item Sub-item 1
      \item Sub-item 2
    \end{itemize}
\end{itemize}

\subsection{Numbered List}
\begin{enumerate}
  \item First step
  \item Second step
  \item Third step
\end{enumerate}

\section{Citations and References}
When you use the citation search feature, references will be automatically managed. For example, you might cite a paper like this: \cite{example2024}. All citations will be automatically added to your \texttt{cite.bib} file.

\section{Conclusion}
This template demonstrates basic \LaTeX{} features including mathematical formulas, tables, and lists. The bibliography system is now set up to work seamlessly with Bibby AI's citation management. Feel free to modify it for your needs!

\bibliographystyle{plain}
\bibliography{cite}

\end{document}\documentclass{article}
\usepackage[utf8]{inputenc}
\usepackage{booktabs}
\usepackage{amsmath}
\usepackage{natbib}

\title{My Document}
\author{Your Name}
\date{\today}

\begin{document}

\maketitle

\section{Introduction}
This is a sample document demonstrating various \LaTeX{} features including mathematical equations, tables, and formatting. You can cite papers using the citation search feature, and they will be automatically added to your \texttt{cite.bib} file.

\section{Mathematical Expressions}
Here are some example equations:

\subsection{Quadratic Formula}
The quadratic formula is given by:
\[x = \frac{-b \pm \sqrt{b^2 - 4ac}}{2a}\]

\subsection{Matrix Example}
A sample matrix:
\[
\begin{pmatrix}
a_{11} & a_{12} & a_{13} \\
a_{21} & a_{22} & a_{23} \\
a_{31} & a_{32} & a_{33}
\end{pmatrix}
\]

\section{Tables}
Below is a sample table showing data:

\begin{table}[h]
\centering
\begin{tabular}{@{}lrrr@{}}
\toprule
\textbf{Category} & \textbf{Value 1} & \textbf{Value 2} & \textbf{Total} \\
\midrule
Item A & 10.5 & 20.3 & 30.8 \\
Item B & 15.2 & 25.1 & 40.3 \\
Item C & 22.1 & 12.7 & 34.8 \\
\bottomrule
\end{tabular}
\caption{Sample Data Table}
\label{tab:sample}
\end{table}

\section{Lists}
\subsection{Bullet Points}
\begin{itemize}
  \item First item
  \item Second item
  \item Third item with sub-items:
    \begin{itemize}
      \item Sub-item 1
      \item Sub-item 2
    \end{itemize}
\end{itemize}

\subsection{Numbered List}
\begin{enumerate}
  \item First step
  \item Second step
  \item Third step
\end{enumerate}

\section{Citations and References}
When you use the citation search feature, references will be automatically managed. For example, you might cite a paper like this: \cite{example2024}. All citations will be automatically added to your \texttt{cite.bib} file.

\section{Conclusion}
This template demonstrates basic \LaTeX{} features including mathematical formulas, tables, and lists. The bibliography system is now set up to work seamlessly with Bibby AI's citation management. Feel free to modify it for your needs!

\bibliographystyle{plain}
\bibliography{cite}

\end{document}\documentclass{article}
\usepackage[utf8]{inputenc}
\usepackage{booktabs}
\usepackage{amsmath}
\usepackage{natbib}

\title{My Document}
\author{Your Name}
\date{\today}

\begin{document}

\maketitle

\section{Introduction}
This is a sample document demonstrating various \LaTeX{} features including mathematical equations, tables, and formatting. You can cite papers using the citation search feature, and they will be automatically added to your \texttt{cite.bib} file.

\section{Mathematical Expressions}
Here are some example equations:

\subsection{Quadratic Formula}
The quadratic formula is given by:
\[x = \frac{-b \pm \sqrt{b^2 - 4ac}}{2a}\]

\subsection{Matrix Example}
A sample matrix:
\[
\begin{pmatrix}
a_{11} & a_{12} & a_{13} \\
a_{21} & a_{22} & a_{23} \\
a_{31} & a_{32} & a_{33}
\end{pmatrix}
\]

\section{Tables}
Below is a sample table showing data:

\begin{table}[h]
\centering
\begin{tabular}{@{}lrrr@{}}
\toprule
\textbf{Category} & \textbf{Value 1} & \textbf{Value 2} & \textbf{Total} \\
\midrule
Item A & 10.5 & 20.3 & 30.8 \\
Item B & 15.2 & 25.1 & 40.3 \\
Item C & 22.1 & 12.7 & 34.8 \\
\bottomrule
\end{tabular}
\caption{Sample Data Table}
\label{tab:sample}
\end{table}

\section{Lists}
\subsection{Bullet Points}
\begin{itemize}
  \item First item
  \item Second item
  \item Third item with sub-items:
    \begin{itemize}
      \item Sub-item 1
      \item Sub-item 2
    \end{itemize}
\end{itemize}

\subsection{Numbered List}
\begin{enumerate}
  \item First step
  \item Second step
  \item Third step
\end{enumerate}

\section{Citations and References}
When you use the citation search feature, references will be automatically managed. For example, you might cite a paper like this: \cite{example2024}. All citations will be automatically added to your \texttt{cite.bib} file.

\section{Conclusion}
This template demonstrates basic \LaTeX{} features including mathematical formulas, tables, and lists. The bibliography system is now set up to work seamlessly with Bibby AI's citation management. Feel free to modify it for your needs!

\bibliographystyle{plain}
\bibliography{cite}

\end{document}\documentclass{article}
\usepackage[utf8]{inputenc}
\usepackage{booktabs}
\usepackage{amsmath}
\usepackage{natbib}

\title{My Document}
\author{Your Name}
\date{\today}

\begin{document}

\maketitle

\section{Introduction}
This is a sample document demonstrating various \LaTeX{} features including mathematical equations, tables, and formatting. You can cite papers using the citation search feature, and they will be automatically added to your \texttt{cite.bib} file.

\section{Mathematical Expressions}
Here are some example equations:

\subsection{Quadratic Formula}
The quadratic formula is given by:
\[x = \frac{-b \pm \sqrt{b^2 - 4ac}}{2a}\]

\subsection{Matrix Example}
A sample matrix:
\[
\begin{pmatrix}
a_{11} & a_{12} & a_{13} \\
a_{21} & a_{22} & a_{23} \\
a_{31} & a_{32} & a_{33}
\end{pmatrix}
\]

\section{Tables}
Below is a sample table showing data:

\begin{table}[h]
\centering
\begin{tabular}{@{}lrrr@{}}
\toprule
\textbf{Category} & \textbf{Value 1} & \textbf{Value 2} & \textbf{Total} \\
\midrule
Item A & 10.5 & 20.3 & 30.8 \\
Item B & 15.2 & 25.1 & 40.3 \\
Item C & 22.1 & 12.7 & 34.8 \\
\bottomrule
\end{tabular}
\caption{Sample Data Table}
\label{tab:sample}
\end{table}

\section{Lists}
\subsection{Bullet Points}
\begin{itemize}
  \item First item
  \item Second item
  \item Third item with sub-items:
    \begin{itemize}
      \item Sub-item 1
      \item Sub-item 2
    \end{itemize}
\end{itemize}

\subsection{Numbered List}
\begin{enumerate}
  \item First step
  \item Second step
  \item Third step
\end{enumerate}

\section{Citations and References}
When you use the citation search feature, references will be automatically managed. For example, you might cite a paper like this: \cite{example2024}. All citations will be automatically added to your \texttt{cite.bib} file.

\section{Conclusion}
This template demonstrates basic \LaTeX{} features including mathematical formulas, tables, and lists. The bibliography system is now set up to work seamlessly with Bibby AI's citation management. Feel free to modify it for your needs!

\bibliographystyle{plain}
\bibliography{cite}

\end{document}\documentclass{article}
\usepackage[utf8]{inputenc}
\usepackage{booktabs}
\usepackage{amsmath}
\usepackage{natbib}

\title{My Document}
\author{Your Name}
\date{\today}

\begin{document}

\maketitle

\section{Introduction}
This is a sample document demonstrating various \LaTeX{} features including mathematical equations, tables, and formatting. You can cite papers using the citation search feature, and they will be automatically added to your \texttt{cite.bib} file.

\section{Mathematical Expressions}
Here are some example equations:

\subsection{Quadratic Formula}
The quadratic formula is given by:
\[x = \frac{-b \pm \sqrt{b^2 - 4ac}}{2a}\]

\subsection{Matrix Example}
A sample matrix:
\[
\begin{pmatrix}
a_{11} & a_{12} & a_{13} \\
a_{21} & a_{22} & a_{23} \\
a_{31} & a_{32} & a_{33}
\end{pmatrix}
\]

\section{Tables}
Below is a sample table showing data:

\begin{table}[h]
\centering
\begin{tabular}{@{}lrrr@{}}
\toprule
\textbf{Category} & \textbf{Value 1} & \textbf{Value 2} & \textbf{Total} \\
\midrule
Item A & 10.5 & 20.3 & 30.8 \\
Item B & 15.2 & 25.1 & 40.3 \\
Item C & 22.1 & 12.7 & 34.8 \\
\bottomrule
\end{tabular}
\caption{Sample Data Table}
\label{tab:sample}
\end{table}

\section{Lists}
\subsection{Bullet Points}
\begin{itemize}
  \item First item
  \item Second item
  \item Third item with sub-items:
    \begin{itemize}
      \item Sub-item 1
      \item Sub-item 2
    \end{itemize}
\end{itemize}

\subsection{Numbered List}
\begin{enumerate}
  \item First step
  \item Second step
  \item Third step
\end{enumerate}

\section{Citations and References}
When you use the citation search feature, references will be automatically managed. For example, you might cite a paper like this: \cite{example2024}. All citations will be automatically added to your \texttt{cite.bib} file.

\section{Conclusion}
This template demonstrates basic \LaTeX{} features including mathematical formulas, tables, and lists. The bibliography system is now set up to work seamlessly with Bibby AI's citation management. Feel free to modify it for your needs!

\bibliographystyle{plain}
\bibliography{cite}

\end{document}\documentclass{article}
\usepackage[utf8]{inputenc}
\usepackage{booktabs}
\usepackage{amsmath}
\usepackage{natbib}

\title{My Document}
\author{Your Name}
\date{\today}

\begin{document}

\maketitle

\section{Introduction}
This is a sample document demonstrating various \LaTeX{} features including mathematical equations, tables, and formatting. You can cite papers using the citation search feature, and they will be automatically added to your \texttt{cite.bib} file.

\section{Mathematical Expressions}
Here are some example equations:

\subsection{Quadratic Formula}
The quadratic formula is given by:
\[x = \frac{-b \pm \sqrt{b^2 - 4ac}}{2a}\]

\subsection{Matrix Example}
A sample matrix:
\[
\begin{pmatrix}
a_{11} & a_{12} & a_{13} \\
a_{21} & a_{22} & a_{23} \\
a_{31} & a_{32} & a_{33}
\end{pmatrix}
\]

\section{Tables}
Below is a sample table showing data:

\begin{table}[h]
\centering
\begin{tabular}{@{}lrrr@{}}
\toprule
\textbf{Category} & \textbf{Value 1} & \textbf{Value 2} & \textbf{Total} \\
\midrule
Item A & 10.5 & 20.3 & 30.8 \\
Item B & 15.2 & 25.1 & 40.3 \\
Item C & 22.1 & 12.7 & 34.8 \\
\bottomrule
\end{tabular}
\caption{Sample Data Table}
\label{tab:sample}
\end{table}

\section{Lists}
\subsection{Bullet Points}
\begin{itemize}
  \item First item
  \item Second item
  \item Third item with sub-items:
    \begin{itemize}
      \item Sub-item 1
      \item Sub-item 2
    \end{itemize}
\end{itemize}

\subsection{Numbered List}
\begin{enumerate}
  \item First step
  \item Second step
  \item Third step
\end{enumerate}

\section{Citations and References}
When you use the citation search feature, references will be automatically managed. For example, you might cite a paper like this: \cite{example2024}. All citations will be automatically added to your \texttt{cite.bib} file.

\section{Conclusion}
This template demonstrates basic \LaTeX{} features including mathematical formulas, tables, and lists. The bibliography system is now set up to work seamlessly with Bibby AI's citation management. Feel free to modify it for your needs!

\bibliographystyle{plain}
\bibliography{cite}

\end{document}\documentclass{article}
\usepackage[utf8]{inputenc}
\usepackage{booktabs}
\usepackage{amsmath}
\usepackage{natbib}

\title{My Document}
\author{Your Name}
\date{\today}

\begin{document}

\maketitle

\section{Introduction}
This is a sample document demonstrating various \LaTeX{} features including mathematical equations, tables, and formatting. You can cite papers using the citation search feature, and they will be automatically added to your \texttt{cite.bib} file.

\section{Mathematical Expressions}
Here are some example equations:

\subsection{Quadratic Formula}
The quadratic formula is given by:
\[x = \frac{-b \pm \sqrt{b^2 - 4ac}}{2a}\]

\subsection{Matrix Example}
A sample matrix:
\[
\begin{pmatrix}
a_{11} & a_{12} & a_{13} \\
a_{21} & a_{22} & a_{23} \\
a_{31} & a_{32} & a_{33}
\end{pmatrix}
\]

\section{Tables}
Below is a sample table showing data:

\begin{table}[h]
\centering
\begin{tabular}{@{}lrrr@{}}
\toprule
\textbf{Category} & \textbf{Value 1} & \textbf{Value 2} & \textbf{Total} \\
\midrule
Item A & 10.5 & 20.3 & 30.8 \\
Item B & 15.2 & 25.1 & 40.3 \\
Item C & 22.1 & 12.7 & 34.8 \\
\bottomrule
\end{tabular}
\caption{Sample Data Table}
\label{tab:sample}
\end{table}

\section{Lists}
\subsection{Bullet Points}
\begin{itemize}
  \item First item
  \item Second item
  \item Third item with sub-items:
    \begin{itemize}
      \item Sub-item 1
      \item Sub-item 2
    \end{itemize}
\end{itemize}

\subsection{Numbered List}
\begin{enumerate}
  \item First step
  \item Second step
  \item Third step
\end{enumerate}

\section{Citations and References}
When you use the citation search feature, references will be automatically managed. For example, you might cite a paper like this: \cite{example2024}. All citations will be automatically added to your \texttt{cite.bib} file.

\section{Conclusion}
This template demonstrates basic \LaTeX{} features including mathematical formulas, tables, and lists. The bibliography system is now set up to work seamlessly with Bibby AI's citation management. Feel free to modify it for your needs!

\bibliographystyle{plain}
\bibliography{cite}

\end{document}\documentclass{article}
\usepackage[utf8]{inputenc}
\usepackage{booktabs}
\usepackage{amsmath}
\usepackage{natbib}

\title{My Document}
\author{Your Name}
\date{\today}

\begin{document}

\maketitle

\section{Introduction}
This is a sample document demonstrating various \LaTeX{} features including mathematical equations, tables, and formatting. You can cite papers using the citation search feature, and they will be automatically added to your \texttt{cite.bib} file.

\section{Mathematical Expressions}
Here are some example equations:

\subsection{Quadratic Formula}
The quadratic formula is given by:
\[x = \frac{-b \pm \sqrt{b^2 - 4ac}}{2a}\]

\subsection{Matrix Example}
A sample matrix:
\[
\begin{pmatrix}
a_{11} & a_{12} & a_{13} \\
a_{21} & a_{22} & a_{23} \\
a_{31} & a_{32} & a_{33}
\end{pmatrix}
\]

\section{Tables}
Below is a sample table showing data:

\begin{table}[h]
\centering
\begin{tabular}{@{}lrrr@{}}
\toprule
\textbf{Category} & \textbf{Value 1} & \textbf{Value 2} & \textbf{Total} \\
\midrule
Item A & 10.5 & 20.3 & 30.8 \\
Item B & 15.2 & 25.1 & 40.3 \\
Item C & 22.1 & 12.7 & 34.8 \\
\bottomrule
\end{tabular}
\caption{Sample Data Table}
\label{tab:sample}
\end{table}

\section{Lists}
\subsection{Bullet Points}
\begin{itemize}
  \item First item
  \item Second item
  \item Third item with sub-items:
    \begin{itemize}
      \item Sub-item 1
      \item Sub-item 2
    \end{itemize}
\end{itemize}

\subsection{Numbered List}
\begin{enumerate}
  \item First step
  \item Second step
  \item Third step
\end{enumerate}

\section{Citations and References}
When you use the citation search feature, references will be automatically managed. For example, you might cite a paper like this: \cite{example2024}. All citations will be automatically added to your \texttt{cite.bib} file.

\section{Conclusion}
This template demonstrates basic \LaTeX{} features including mathematical formulas, tables, and lists. The bibliography system is now set up to work seamlessly with Bibby AI's citation management. Feel free to modify it for your needs!

\bibliographystyle{plain}
\bibliography{cite}

\end{document}\documentclass{article}
\usepackage[utf8]{inputenc}
\usepackage{booktabs}
\usepackage{amsmath}
\usepackage{natbib}

\title{My Document}
\author{Your Name}
\date{\today}

\begin{document}

\maketitle

\section{Introduction}
This is a sample document demonstrating various \LaTeX{} features including mathematical equations, tables, and formatting. You can cite papers using the citation search feature, and they will be automatically added to your \texttt{cite.bib} file.

\section{Mathematical Expressions}
Here are some example equations:

\subsection{Quadratic Formula}
The quadratic formula is given by:
\[x = \frac{-b \pm \sqrt{b^2 - 4ac}}{2a}\]

\subsection{Matrix Example}
A sample matrix:
\[
\begin{pmatrix}
a_{11} & a_{12} & a_{13} \\
a_{21} & a_{22} & a_{23} \\
a_{31} & a_{32} & a_{33}
\end{pmatrix}
\]

\section{Tables}
Below is a sample table showing data:

\begin{table}[h]
\centering
\begin{tabular}{@{}lrrr@{}}
\toprule
\textbf{Category} & \textbf{Value 1} & \textbf{Value 2} & \textbf{Total} \\
\midrule
Item A & 10.5 & 20.3 & 30.8 \\
Item B & 15.2 & 25.1 & 40.3 \\
Item C & 22.1 & 12.7 & 34.8 \\
\bottomrule
\end{tabular}
\caption{Sample Data Table}
\label{tab:sample}
\end{table}

\section{Lists}
\subsection{Bullet Points}
\begin{itemize}
  \item First item
  \item Second item
  \item Third item with sub-items:
    \begin{itemize}
      \item Sub-item 1
      \item Sub-item 2
    \end{itemize}
\end{itemize}

\subsection{Numbered List}
\begin{enumerate}
  \item First step
  \item Second step
  \item Third step
\end{enumerate}

\section{Citations and References}
When you use the citation search feature, references will be automatically managed. For example, you might cite a paper like this: \cite{example2024}. All citations will be automatically added to your \texttt{cite.bib} file.

\section{Conclusion}
This template demonstrates basic \LaTeX{} features including mathematical formulas, tables, and lists. The bibliography system is now set up to work seamlessly with Bibby AI's citation management. Feel free to modify it for your needs!

\bibliographystyle{plain}
\bibliography{cite}

\end{document}\documentclass{article}
\usepackage[utf8]{inputenc}
\usepackage{booktabs}
\usepackage{amsmath}
\usepackage{natbib}

\title{My Document}
\author{Your Name}
\date{\today}

\begin{document}

\maketitle

\section{Introduction}
This is a sample document demonstrating various \LaTeX{} features including mathematical equations, tables, and formatting. You can cite papers using the citation search feature, and they will be automatically added to your \texttt{cite.bib} file.

\section{Mathematical Expressions}
Here are some example equations:

\subsection{Quadratic Formula}
The quadratic formula is given by:
\[x = \frac{-b \pm \sqrt{b^2 - 4ac}}{2a}\]

\subsection{Matrix Example}
A sample matrix:
\[
\begin{pmatrix}
a_{11} & a_{12} & a_{13} \\
a_{21} & a_{22} & a_{23} \\
a_{31} & a_{32} & a_{33}
\end{pmatrix}
\]

\section{Tables}
Below is a sample table showing data:

\begin{table}[h]
\centering
\begin{tabular}{@{}lrrr@{}}
\toprule
\textbf{Category} & \textbf{Value 1} & \textbf{Value 2} & \textbf{Total} \\
\midrule
Item A & 10.5 & 20.3 & 30.8 \\
Item B & 15.2 & 25.1 & 40.3 \\
Item C & 22.1 & 12.7 & 34.8 \\
\bottomrule
\end{tabular}
\caption{Sample Data Table}
\label{tab:sample}
\end{table}

\section{Lists}
\subsection{Bullet Points}
\begin{itemize}
  \item First item
  \item Second item
  \item Third item with sub-items:
    \begin{itemize}
      \item Sub-item 1
      \item Sub-item 2
    \end{itemize}
\end{itemize}

\subsection{Numbered List}
\begin{enumerate}
  \item First step
  \item Second step
  \item Third step
\end{enumerate}

\section{Citations and References}
When you use the citation search feature, references will be automatically managed. For example, you might cite a paper like this: \cite{example2024}. All citations will be automatically added to your \texttt{cite.bib} file.

\section{Conclusion}
This template demonstrates basic \LaTeX{} features including mathematical formulas, tables, and lists. The bibliography system is now set up to work seamlessly with Bibby AI's citation management. Feel free to modify it for your needs!

\bibliographystyle{plain}
\bibliography{cite}

\end{document}\documentclass{article}
\usepackage[utf8]{inputenc}
\usepackage{booktabs}
\usepackage{amsmath}
\usepackage{natbib}

\title{My Document}
\author{Your Name}
\date{\today}

\begin{document}

\maketitle

\section{Introduction}
This is a sample document demonstrating various \LaTeX{} features including mathematical equations, tables, and formatting. You can cite papers using the citation search feature, and they will be automatically added to your \texttt{cite.bib} file.

\section{Mathematical Expressions}
Here are some example equations:

\subsection{Quadratic Formula}
The quadratic formula is given by:
\[x = \frac{-b \pm \sqrt{b^2 - 4ac}}{2a}\]

\subsection{Matrix Example}
A sample matrix:
\[
\begin{pmatrix}
a_{11} & a_{12} & a_{13} \\
a_{21} & a_{22} & a_{23} \\
a_{31} & a_{32} & a_{33}
\end{pmatrix}
\]

\section{Tables}
Below is a sample table showing data:

\begin{table}[h]
\centering
\begin{tabular}{@{}lrrr@{}}
\toprule
\textbf{Category} & \textbf{Value 1} & \textbf{Value 2} & \textbf{Total} \\
\midrule
Item A & 10.5 & 20.3 & 30.8 \\
Item B & 15.2 & 25.1 & 40.3 \\
Item C & 22.1 & 12.7 & 34.8 \\
\bottomrule
\end{tabular}
\caption{Sample Data Table}
\label{tab:sample}
\end{table}

\section{Lists}
\subsection{Bullet Points}
\begin{itemize}
  \item First item
  \item Second item
  \item Third item with sub-items:
    \begin{itemize}
      \item Sub-item 1
      \item Sub-item 2
    \end{itemize}
\end{itemize}

\subsection{Numbered List}
\begin{enumerate}
  \item First step
  \item Second step
  \item Third step
\end{enumerate}

\section{Citations and References}
When you use the citation search feature, references will be automatically managed. For example, you might cite a paper like this: \cite{example2024}. All citations will be automatically added to your \texttt{cite.bib} file.

\section{Conclusion}
This template demonstrates basic \LaTeX{} features including mathematical formulas, tables, and lists. The bibliography system is now set up to work seamlessly with Bibby AI's citation management. Feel free to modify it for your needs!

\bibliographystyle{plain}
\bibliography{cite}

\end{document}\documentclass{article}
\usepackage[utf8]{inputenc}
\usepackage{booktabs}
\usepackage{amsmath}
\usepackage{natbib}

\title{My Document}
\author{Your Name}
\date{\today}

\begin{document}

\maketitle

\section{Introduction}
This is a sample document demonstrating various \LaTeX{} features including mathematical equations, tables, and formatting. You can cite papers using the citation search feature, and they will be automatically added to your \texttt{cite.bib} file.

\section{Mathematical Expressions}
Here are some example equations:

\subsection{Quadratic Formula}
The quadratic formula is given by:
\[x = \frac{-b \pm \sqrt{b^2 - 4ac}}{2a}\]

\subsection{Matrix Example}
A sample matrix:
\[
\begin{pmatrix}
a_{11} & a_{12} & a_{13} \\
a_{21} & a_{22} & a_{23} \\
a_{31} & a_{32} & a_{33}
\end{pmatrix}
\]

\section{Tables}
Below is a sample table showing data:

\begin{table}[h]
\centering
\begin{tabular}{@{}lrrr@{}}
\toprule
\textbf{Category} & \textbf{Value 1} & \textbf{Value 2} & \textbf{Total} \\
\midrule
Item A & 10.5 & 20.3 & 30.8 \\
Item B & 15.2 & 25.1 & 40.3 \\
Item C & 22.1 & 12.7 & 34.8 \\
\bottomrule
\end{tabular}
\caption{Sample Data Table}
\label{tab:sample}
\end{table}

\section{Lists}
\subsection{Bullet Points}
\begin{itemize}
  \item First item
  \item Second item
  \item Third item with sub-items:
    \begin{itemize}
      \item Sub-item 1
      \item Sub-item 2
    \end{itemize}
\end{itemize}

\subsection{Numbered List}
\begin{enumerate}
  \item First step
  \item Second step
  \item Third step
\end{enumerate}

\section{Citations and References}
When you use the citation search feature, references will be automatically managed. For example, you might cite a paper like this: \cite{example2024}. All citations will be automatically added to your \texttt{cite.bib} file.

\section{Conclusion}
This template demonstrates basic \LaTeX{} features including mathematical formulas, tables, and lists. The bibliography system is now set up to work seamlessly with Bibby AI's citation management. Feel free to modify it for your needs!

\bibliographystyle{plain}
\bibliography{cite}

\end{document}\documentclass{article}
\usepackage[utf8]{inputenc}
\usepackage{booktabs}
\usepackage{amsmath}
\usepackage{natbib}

\title{My Document}
\author{Your Name}
\date{\today}

\begin{document}

\maketitle

\section{Introduction}
This is a sample document demonstrating various \LaTeX{} features including mathematical equations, tables, and formatting. You can cite papers using the citation search feature, and they will be automatically added to your \texttt{cite.bib} file.

\section{Mathematical Expressions}
Here are some example equations:

\subsection{Quadratic Formula}
The quadratic formula is given by:
\[x = \frac{-b \pm \sqrt{b^2 - 4ac}}{2a}\]

\subsection{Matrix Example}
A sample matrix:
\[
\begin{pmatrix}
a_{11} & a_{12} & a_{13} \\
a_{21} & a_{22} & a_{23} \\
a_{31} & a_{32} & a_{33}
\end{pmatrix}
\]

\section{Tables}
Below is a sample table showing data:

\begin{table}[h]
\centering
\begin{tabular}{@{}lrrr@{}}
\toprule
\textbf{Category} & \textbf{Value 1} & \textbf{Value 2} & \textbf{Total} \\
\midrule
Item A & 10.5 & 20.3 & 30.8 \\
Item B & 15.2 & 25.1 & 40.3 \\
Item C & 22.1 & 12.7 & 34.8 \\
\bottomrule
\end{tabular}
\caption{Sample Data Table}
\label{tab:sample}
\end{table}

\section{Lists}
\subsection{Bullet Points}
\begin{itemize}
  \item First item
  \item Second item
  \item Third item with sub-items:
    \begin{itemize}
      \item Sub-item 1
      \item Sub-item 2
    \end{itemize}
\end{itemize}

\subsection{Numbered List}
\begin{enumerate}
  \item First step
  \item Second step
  \item Third step
\end{enumerate}

\section{Citations and References}
When you use the citation search feature, references will be automatically managed. For example, you might cite a paper like this: \cite{example2024}. All citations will be automatically added to your \texttt{cite.bib} file.

\section{Conclusion}
This template demonstrates basic \LaTeX{} features including mathematical formulas, tables, and lists. The bibliography system is now set up to work seamlessly with Bibby AI's citation management. Feel free to modify it for your needs!

\bibliographystyle{plain}
\bibliography{cite}

\end{document}\documentclass{article}
\usepackage[utf8]{inputenc}
\usepackage{booktabs}
\usepackage{amsmath}
\usepackage{natbib}

\title{My Document}
\author{Your Name}
\date{\today}

\begin{document}

\maketitle

\section{Introduction}
This is a sample document demonstrating various \LaTeX{} features including mathematical equations, tables, and formatting. You can cite papers using the citation search feature, and they will be automatically added to your \texttt{cite.bib} file.

\section{Mathematical Expressions}
Here are some example equations:

\subsection{Quadratic Formula}
The quadratic formula is given by:
\[x = \frac{-b \pm \sqrt{b^2 - 4ac}}{2a}\]

\subsection{Matrix Example}
A sample matrix:
\[
\begin{pmatrix}
a_{11} & a_{12} & a_{13} \\
a_{21} & a_{22} & a_{23} \\
a_{31} & a_{32} & a_{33}
\end{pmatrix}
\]

\section{Tables}
Below is a sample table showing data:

\begin{table}[h]
\centering
\begin{tabular}{@{}lrrr@{}}
\toprule
\textbf{Category} & \textbf{Value 1} & \textbf{Value 2} & \textbf{Total} \\
\midrule
Item A & 10.5 & 20.3 & 30.8 \\
Item B & 15.2 & 25.1 & 40.3 \\
Item C & 22.1 & 12.7 & 34.8 \\
\bottomrule
\end{tabular}
\caption{Sample Data Table}
\label{tab:sample}
\end{table}

\section{Lists}
\subsection{Bullet Points}
\begin{itemize}
  \item First item
  \item Second item
  \item Third item with sub-items:
    \begin{itemize}
      \item Sub-item 1
      \item Sub-item 2
    \end{itemize}
\end{itemize}

\subsection{Numbered List}
\begin{enumerate}
  \item First step
  \item Second step
  \item Third step
\end{enumerate}

\section{Citations and References}
When you use the citation search feature, references will be automatically managed. For example, you might cite a paper like this: \cite{example2024}. All citations will be automatically added to your \texttt{cite.bib} file.

\section{Conclusion}
This template demonstrates basic \LaTeX{} features including mathematical formulas, tables, and lists. The bibliography system is now set up to work seamlessly with Bibby AI's citation management. Feel free to modify it for your needs!

\bibliographystyle{plain}
\bibliography{cite}

\end{document}\documentclass{article}
\usepackage[utf8]{inputenc}
\usepackage{booktabs}
\usepackage{amsmath}
\usepackage{natbib}

\title{My Document}
\author{Your Name}
\date{\today}

\begin{document}

\maketitle

\section{Introduction}
This is a sample document demonstrating various \LaTeX{} features including mathematical equations, tables, and formatting. You can cite papers using the citation search feature, and they will be automatically added to your \texttt{cite.bib} file.

\section{Mathematical Expressions}
Here are some example equations:

\subsection{Quadratic Formula}
The quadratic formula is given by:
\[x = \frac{-b \pm \sqrt{b^2 - 4ac}}{2a}\]

\subsection{Matrix Example}
A sample matrix:
\[
\begin{pmatrix}
a_{11} & a_{12} & a_{13} \\
a_{21} & a_{22} & a_{23} \\
a_{31} & a_{32} & a_{33}
\end{pmatrix}
\]

\section{Tables}
Below is a sample table showing data:

\begin{table}[h]
\centering
\begin{tabular}{@{}lrrr@{}}
\toprule
\textbf{Category} & \textbf{Value 1} & \textbf{Value 2} & \textbf{Total} \\
\midrule
Item A & 10.5 & 20.3 & 30.8 \\
Item B & 15.2 & 25.1 & 40.3 \\
Item C & 22.1 & 12.7 & 34.8 \\
\bottomrule
\end{tabular}
\caption{Sample Data Table}
\label{tab:sample}
\end{table}

\section{Lists}
\subsection{Bullet Points}
\begin{itemize}
  \item First item
  \item Second item
  \item Third item with sub-items:
    \begin{itemize}
      \item Sub-item 1
      \item Sub-item 2
    \end{itemize}
\end{itemize}

\subsection{Numbered List}
\begin{enumerate}
  \item First step
  \item Second step
  \item Third step
\end{enumerate}

\section{Citations and References}
When you use the citation search feature, references will be automatically managed. For example, you might cite a paper like this: \cite{example2024}. All citations will be automatically added to your \texttt{cite.bib} file.

\section{Conclusion}
This template demonstrates basic \LaTeX{} features including mathematical formulas, tables, and lists. The bibliography system is now set up to work seamlessly with Bibby AI's citation management. Feel free to modify it for your needs!

\bibliographystyle{plain}
\bibliography{cite}

\end{document}\documentclass{article}
\usepackage[utf8]{inputenc}
\usepackage{booktabs}
\usepackage{amsmath}
\usepackage{natbib}

\title{My Document}
\author{Your Name}
\date{\today}

\begin{document}

\maketitle

\section{Introduction}
This is a sample document demonstrating various \LaTeX{} features including mathematical equations, tables, and formatting. You can cite papers using the citation search feature, and they will be automatically added to your \texttt{cite.bib} file.

\section{Mathematical Expressions}
Here are some example equations:

\subsection{Quadratic Formula}
The quadratic formula is given by:
\[x = \frac{-b \pm \sqrt{b^2 - 4ac}}{2a}\]

\subsection{Matrix Example}
A sample matrix:
\[
\begin{pmatrix}
a_{11} & a_{12} & a_{13} \\
a_{21} & a_{22} & a_{23} \\
a_{31} & a_{32} & a_{33}
\end{pmatrix}
\]

\section{Tables}
Below is a sample table showing data:

\begin{table}[h]
\centering
\begin{tabular}{@{}lrrr@{}}
\toprule
\textbf{Category} & \textbf{Value 1} & \textbf{Value 2} & \textbf{Total} \\
\midrule
Item A & 10.5 & 20.3 & 30.8 \\
Item B & 15.2 & 25.1 & 40.3 \\
Item C & 22.1 & 12.7 & 34.8 \\
\bottomrule
\end{tabular}
\caption{Sample Data Table}
\label{tab:sample}
\end{table}

\section{Lists}
\subsection{Bullet Points}
\begin{itemize}
  \item First item
  \item Second item
  \item Third item with sub-items:
    \begin{itemize}
      \item Sub-item 1
      \item Sub-item 2
    \end{itemize}
\end{itemize}

\subsection{Numbered List}
\begin{enumerate}
  \item First step
  \item Second step
  \item Third step
\end{enumerate}

\section{Citations and References}
When you use the citation search feature, references will be automatically managed. For example, you might cite a paper like this: \cite{example2024}. All citations will be automatically added to your \texttt{cite.bib} file.

\section{Conclusion}
This template demonstrates basic \LaTeX{} features including mathematical formulas, tables, and lists. The bibliography system is now set up to work seamlessly with Bibby AI's citation management. Feel free to modify it for your needs!

\bibliographystyle{plain}
\bibliography{cite}

\end{document}\documentclass{article}
\usepackage[utf8]{inputenc}
\usepackage{booktabs}
\usepackage{amsmath}
\usepackage{natbib}

\title{My Document}
\author{Your Name}
\date{\today}

\begin{document}

\maketitle

\section{Introduction}
This is a sample document demonstrating various \LaTeX{} features including mathematical equations, tables, and formatting. You can cite papers using the citation search feature, and they will be automatically added to your \texttt{cite.bib} file.

\section{Mathematical Expressions}
Here are some example equations:

\subsection{Quadratic Formula}
The quadratic formula is given by:
\[x = \frac{-b \pm \sqrt{b^2 - 4ac}}{2a}\]

\subsection{Matrix Example}
A sample matrix:
\[
\begin{pmatrix}
a_{11} & a_{12} & a_{13} \\
a_{21} & a_{22} & a_{23} \\
a_{31} & a_{32} & a_{33}
\end{pmatrix}
\]

\section{Tables}
Below is a sample table showing data:

\begin{table}[h]
\centering
\begin{tabular}{@{}lrrr@{}}
\toprule
\textbf{Category} & \textbf{Value 1} & \textbf{Value 2} & \textbf{Total} \\
\midrule
Item A & 10.5 & 20.3 & 30.8 \\
Item B & 15.2 & 25.1 & 40.3 \\
Item C & 22.1 & 12.7 & 34.8 \\
\bottomrule
\end{tabular}
\caption{Sample Data Table}
\label{tab:sample}
\end{table}

\section{Lists}
\subsection{Bullet Points}
\begin{itemize}
  \item First item
  \item Second item
  \item Third item with sub-items:
    \begin{itemize}
      \item Sub-item 1
      \item Sub-item 2
    \end{itemize}
\end{itemize}

\subsection{Numbered List}
\begin{enumerate}
  \item First step
  \item Second step
  \item Third step
\end{enumerate}

\section{Citations and References}
When you use the citation search feature, references will be automatically managed. For example, you might cite a paper like this: \cite{example2024}. All citations will be automatically added to your \texttt{cite.bib} file.

\section{Conclusion}
This template demonstrates basic \LaTeX{} features including mathematical formulas, tables, and lists. The bibliography system is now set up to work seamlessly with Bibby AI's citation management. Feel free to modify it for your needs!

\bibliographystyle{plain}
\bibliography{cite}

\end{document}\documentclass{article}
\usepackage[utf8]{inputenc}
\usepackage{booktabs}
\usepackage{amsmath}
\usepackage{natbib}

\title{My Document}
\author{Your Name}
\date{\today}

\begin{document}

\maketitle

\section{Introduction}
This is a sample document demonstrating various \LaTeX{} features including mathematical equations, tables, and formatting. You can cite papers using the citation search feature, and they will be automatically added to your \texttt{cite.bib} file.

\section{Mathematical Expressions}
Here are some example equations:

\subsection{Quadratic Formula}
The quadratic formula is given by:
\[x = \frac{-b \pm \sqrt{b^2 - 4ac}}{2a}\]

\subsection{Matrix Example}
A sample matrix:
\[
\begin{pmatrix}
a_{11} & a_{12} & a_{13} \\
a_{21} & a_{22} & a_{23} \\
a_{31} & a_{32} & a_{33}
\end{pmatrix}
\]

\section{Tables}
Below is a sample table showing data:

\begin{table}[h]
\centering
\begin{tabular}{@{}lrrr@{}}
\toprule
\textbf{Category} & \textbf{Value 1} & \textbf{Value 2} & \textbf{Total} \\
\midrule
Item A & 10.5 & 20.3 & 30.8 \\
Item B & 15.2 & 25.1 & 40.3 \\
Item C & 22.1 & 12.7 & 34.8 \\
\bottomrule
\end{tabular}
\caption{Sample Data Table}
\label{tab:sample}
\end{table}

\section{Lists}
\subsection{Bullet Points}
\begin{itemize}
  \item First item
  \item Second item
  \item Third item with sub-items:
    \begin{itemize}
      \item Sub-item 1
      \item Sub-item 2
    \end{itemize}
\end{itemize}

\subsection{Numbered List}
\begin{enumerate}
  \item First step
  \item Second step
  \item Third step
\end{enumerate}

\section{Citations and References}
When you use the citation search feature, references will be automatically managed. For example, you might cite a paper like this: \cite{example2024}. All citations will be automatically added to your \texttt{cite.bib} file.

\section{Conclusion}
This template demonstrates basic \LaTeX{} features including mathematical formulas, tables, and lists. The bibliography system is now set up to work seamlessly with Bibby AI's citation management. Feel free to modify it for your needs!

\bibliographystyle{plain}
\bibliography{cite}

\end{document}\documentclass{article}
\usepackage[utf8]{inputenc}
\usepackage{booktabs}
\usepackage{amsmath}
\usepackage{natbib}

\title{My Document}
\author{Your Name}
\date{\today}

\begin{document}

\maketitle

\section{Introduction}
This is a sample document demonstrating various \LaTeX{} features including mathematical equations, tables, and formatting. You can cite papers using the citation search feature, and they will be automatically added to your \texttt{cite.bib} file.

\section{Mathematical Expressions}
Here are some example equations:

\subsection{Quadratic Formula}
The quadratic formula is given by:
\[x = \frac{-b \pm \sqrt{b^2 - 4ac}}{2a}\]

\subsection{Matrix Example}
A sample matrix:
\[
\begin{pmatrix}
a_{11} & a_{12} & a_{13} \\
a_{21} & a_{22} & a_{23} \\
a_{31} & a_{32} & a_{33}
\end{pmatrix}
\]

\section{Tables}
Below is a sample table showing data:

\begin{table}[h]
\centering
\begin{tabular}{@{}lrrr@{}}
\toprule
\textbf{Category} & \textbf{Value 1} & \textbf{Value 2} & \textbf{Total} \\
\midrule
Item A & 10.5 & 20.3 & 30.8 \\
Item B & 15.2 & 25.1 & 40.3 \\
Item C & 22.1 & 12.7 & 34.8 \\
\bottomrule
\end{tabular}
\caption{Sample Data Table}
\label{tab:sample}
\end{table}

\section{Lists}
\subsection{Bullet Points}
\begin{itemize}
  \item First item
  \item Second item
  \item Third item with sub-items:
    \begin{itemize}
      \item Sub-item 1
      \item Sub-item 2
    \end{itemize}
\end{itemize}

\subsection{Numbered List}
\begin{enumerate}
  \item First step
  \item Second step
  \item Third step
\end{enumerate}

\section{Citations and References}
When you use the citation search feature, references will be automatically managed. For example, you might cite a paper like this: \cite{example2024}. All citations will be automatically added to your \texttt{cite.bib} file.

\section{Conclusion}
This template demonstrates basic \LaTeX{} features including mathematical formulas, tables, and lists. The bibliography system is now set up to work seamlessly with Bibby AI's citation management. Feel free to modify it for your needs!

\bibliographystyle{plain}
\bibliography{cite}

\end{document}\documentclass{article}
\usepackage[utf8]{inputenc}
\usepackage{booktabs}
\usepackage{amsmath}
\usepackage{natbib}

\title{My Document}
\author{Your Name}
\date{\today}

\begin{document}

\maketitle

\section{Introduction}
This is a sample document demonstrating various \LaTeX{} features including mathematical equations, tables, and formatting. You can cite papers using the citation search feature, and they will be automatically added to your \texttt{cite.bib} file.

\section{Mathematical Expressions}
Here are some example equations:

\subsection{Quadratic Formula}
The quadratic formula is given by:
\[x = \frac{-b \pm \sqrt{b^2 - 4ac}}{2a}\]

\subsection{Matrix Example}
A sample matrix:
\[
\begin{pmatrix}
a_{11} & a_{12} & a_{13} \\
a_{21} & a_{22} & a_{23} \\
a_{31} & a_{32} & a_{33}
\end{pmatrix}
\]

\section{Tables}
Below is a sample table showing data:

\begin{table}[h]
\centering
\begin{tabular}{@{}lrrr@{}}
\toprule
\textbf{Category} & \textbf{Value 1} & \textbf{Value 2} & \textbf{Total} \\
\midrule
Item A & 10.5 & 20.3 & 30.8 \\
Item B & 15.2 & 25.1 & 40.3 \\
Item C & 22.1 & 12.7 & 34.8 \\
\bottomrule
\end{tabular}
\caption{Sample Data Table}
\label{tab:sample}
\end{table}

\section{Lists}
\subsection{Bullet Points}
\begin{itemize}
  \item First item
  \item Second item
  \item Third item with sub-items:
    \begin{itemize}
      \item Sub-item 1
      \item Sub-item 2
    \end{itemize}
\end{itemize}

\subsection{Numbered List}
\begin{enumerate}
  \item First step
  \item Second step
  \item Third step
\end{enumerate}

\section{Citations and References}
When you use the citation search feature, references will be automatically managed. For example, you might cite a paper like this: \cite{example2024}. All citations will be automatically added to your \texttt{cite.bib} file.

\section{Conclusion}
This template demonstrates basic \LaTeX{} features including mathematical formulas, tables, and lists. The bibliography system is now set up to work seamlessly with Bibby AI's citation management. Feel free to modify it for your needs!

\bibliographystyle{plain}
\bibliography{cite}

\end{document}\documentclass{article}
\usepackage[utf8]{inputenc}
\usepackage{booktabs}
\usepackage{amsmath}
\usepackage{natbib}

\title{My Document}
\author{Your Name}
\date{\today}

\begin{document}

\maketitle

\section{Introduction}
This is a sample document demonstrating various \LaTeX{} features including mathematical equations, tables, and formatting. You can cite papers using the citation search feature, and they will be automatically added to your \texttt{cite.bib} file.

\section{Mathematical Expressions}
Here are some example equations:

\subsection{Quadratic Formula}
The quadratic formula is given by:
\[x = \frac{-b \pm \sqrt{b^2 - 4ac}}{2a}\]

\subsection{Matrix Example}
A sample matrix:
\[
\begin{pmatrix}
a_{11} & a_{12} & a_{13} \\
a_{21} & a_{22} & a_{23} \\
a_{31} & a_{32} & a_{33}
\end{pmatrix}
\]

\section{Tables}
Below is a sample table showing data:

\begin{table}[h]
\centering
\begin{tabular}{@{}lrrr@{}}
\toprule
\textbf{Category} & \textbf{Value 1} & \textbf{Value 2} & \textbf{Total} \\
\midrule
Item A & 10.5 & 20.3 & 30.8 \\
Item B & 15.2 & 25.1 & 40.3 \\
Item C & 22.1 & 12.7 & 34.8 \\
\bottomrule
\end{tabular}
\caption{Sample Data Table}
\label{tab:sample}
\end{table}

\section{Lists}
\subsection{Bullet Points}
\begin{itemize}
  \item First item
  \item Second item
  \item Third item with sub-items:
    \begin{itemize}
      \item Sub-item 1
      \item Sub-item 2
    \end{itemize}
\end{itemize}

\subsection{Numbered List}
\begin{enumerate}
  \item First step
  \item Second step
  \item Third step
\end{enumerate}

\section{Citations and References}
When you use the citation search feature, references will be automatically managed. For example, you might cite a paper like this: \cite{example2024}. All citations will be automatically added to your \texttt{cite.bib} file.

\section{Conclusion}
This template demonstrates basic \LaTeX{} features including mathematical formulas, tables, and lists. The bibliography system is now set up to work seamlessly with Bibby AI's citation management. Feel free to modify it for your needs!

\bibliographystyle{plain}
\bibliography{cite}

\end{document}\documentclass{article}
\usepackage[utf8]{inputenc}
\usepackage{booktabs}
\usepackage{amsmath}
\usepackage{natbib}

\title{My Document}
\author{Your Name}
\date{\today}

\begin{document}

\maketitle

\section{Introduction}
This is a sample document demonstrating various \LaTeX{} features including mathematical equations, tables, and formatting. You can cite papers using the citation search feature, and they will be automatically added to your \texttt{cite.bib} file.

\section{Mathematical Expressions}
Here are some example equations:

\subsection{Quadratic Formula}
The quadratic formula is given by:
\[x = \frac{-b \pm \sqrt{b^2 - 4ac}}{2a}\]

\subsection{Matrix Example}
A sample matrix:
\[
\begin{pmatrix}
a_{11} & a_{12} & a_{13} \\
a_{21} & a_{22} & a_{23} \\
a_{31} & a_{32} & a_{33}
\end{pmatrix}
\]

\section{Tables}
Below is a sample table showing data:

\begin{table}[h]
\centering
\begin{tabular}{@{}lrrr@{}}
\toprule
\textbf{Category} & \textbf{Value 1} & \textbf{Value 2} & \textbf{Total} \\
\midrule
Item A & 10.5 & 20.3 & 30.8 \\
Item B & 15.2 & 25.1 & 40.3 \\
Item C & 22.1 & 12.7 & 34.8 \\
\bottomrule
\end{tabular}
\caption{Sample Data Table}
\label{tab:sample}
\end{table}

\section{Lists}
\subsection{Bullet Points}
\begin{itemize}
  \item First item
  \item Second item
  \item Third item with sub-items:
    \begin{itemize}
      \item Sub-item 1
      \item Sub-item 2
    \end{itemize}
\end{itemize}

\subsection{Numbered List}
\begin{enumerate}
  \item First step
  \item Second step
  \item Third step
\end{enumerate}

\section{Citations and References}
When you use the citation search feature, references will be automatically managed. For example, you might cite a paper like this: \cite{example2024}. All citations will be automatically added to your \texttt{cite.bib} file.

\section{Conclusion}
This template demonstrates basic \LaTeX{} features including mathematical formulas, tables, and lists. The bibliography system is now set up to work seamlessly with Bibby AI's citation management. Feel free to modify it for your needs!

\bibliographystyle{plain}
\bibliography{cite}

\end{document}\documentclass{article}
\usepackage[utf8]{inputenc}
\usepackage{booktabs}
\usepackage{amsmath}
\usepackage{natbib}

\title{My Document}
\author{Your Name}
\date{\today}

\begin{document}

\maketitle

\section{Introduction}
This is a sample document demonstrating various \LaTeX{} features including mathematical equations, tables, and formatting. You can cite papers using the citation search feature, and they will be automatically added to your \texttt{cite.bib} file.

\section{Mathematical Expressions}
Here are some example equations:

\subsection{Quadratic Formula}
The quadratic formula is given by:
\[x = \frac{-b \pm \sqrt{b^2 - 4ac}}{2a}\]

\subsection{Matrix Example}
A sample matrix:
\[
\begin{pmatrix}
a_{11} & a_{12} & a_{13} \\
a_{21} & a_{22} & a_{23} \\
a_{31} & a_{32} & a_{33}
\end{pmatrix}
\]

\section{Tables}
Below is a sample table showing data:

\begin{table}[h]
\centering
\begin{tabular}{@{}lrrr@{}}
\toprule
\textbf{Category} & \textbf{Value 1} & \textbf{Value 2} & \textbf{Total} \\
\midrule
Item A & 10.5 & 20.3 & 30.8 \\
Item B & 15.2 & 25.1 & 40.3 \\
Item C & 22.1 & 12.7 & 34.8 \\
\bottomrule
\end{tabular}
\caption{Sample Data Table}
\label{tab:sample}
\end{table}

\section{Lists}
\subsection{Bullet Points}
\begin{itemize}
  \item First item
  \item Second item
  \item Third item with sub-items:
    \begin{itemize}
      \item Sub-item 1
      \item Sub-item 2
    \end{itemize}
\end{itemize}

\subsection{Numbered List}
\begin{enumerate}
  \item First step
  \item Second step
  \item Third step
\end{enumerate}

\section{Citations and References}
When you use the citation search feature, references will be automatically managed. For example, you might cite a paper like this: \cite{example2024}. All citations will be automatically added to your \texttt{cite.bib} file.

\section{Conclusion}
This template demonstrates basic \LaTeX{} features including mathematical formulas, tables, and lists. The bibliography system is now set up to work seamlessly with Bibby AI's citation management. Feel free to modify it for your needs!

\bibliographystyle{plain}
\bibliography{cite}

\end{document}\documentclass{article}
\usepackage[utf8]{inputenc}
\usepackage{booktabs}
\usepackage{amsmath}
\usepackage{natbib}

\title{My Document}
\author{Your Name}
\date{\today}

\begin{document}

\maketitle

\section{Introduction}
This is a sample document demonstrating various \LaTeX{} features including mathematical equations, tables, and formatting. You can cite papers using the citation search feature, and they will be automatically added to your \texttt{cite.bib} file.

\section{Mathematical Expressions}
Here are some example equations:

\subsection{Quadratic Formula}
The quadratic formula is given by:
\[x = \frac{-b \pm \sqrt{b^2 - 4ac}}{2a}\]

\subsection{Matrix Example}
A sample matrix:
\[
\begin{pmatrix}
a_{11} & a_{12} & a_{13} \\
a_{21} & a_{22} & a_{23} \\
a_{31} & a_{32} & a_{33}
\end{pmatrix}
\]

\section{Tables}
Below is a sample table showing data:

\begin{table}[h]
\centering
\begin{tabular}{@{}lrrr@{}}
\toprule
\textbf{Category} & \textbf{Value 1} & \textbf{Value 2} & \textbf{Total} \\
\midrule
Item A & 10.5 & 20.3 & 30.8 \\
Item B & 15.2 & 25.1 & 40.3 \\
Item C & 22.1 & 12.7 & 34.8 \\
\bottomrule
\end{tabular}
\caption{Sample Data Table}
\label{tab:sample}
\end{table}

\section{Lists}
\subsection{Bullet Points}
\begin{itemize}
  \item First item
  \item Second item
  \item Third item with sub-items:
    \begin{itemize}
      \item Sub-item 1
      \item Sub-item 2
    \end{itemize}
\end{itemize}

\subsection{Numbered List}
\begin{enumerate}
  \item First step
  \item Second step
  \item Third step
\end{enumerate}

\section{Citations and References}
When you use the citation search feature, references will be automatically managed. For example, you might cite a paper like this: \cite{example2024}. All citations will be automatically added to your \texttt{cite.bib} file.

\section{Conclusion}
This template demonstrates basic \LaTeX{} features including mathematical formulas, tables, and lists. The bibliography system is now set up to work seamlessly with Bibby AI's citation management. Feel free to modify it for your needs!

\bibliographystyle{plain}
\bibliography{cite}

\end{document}\documentclass{article}
\usepackage[utf8]{inputenc}
\usepackage{booktabs}
\usepackage{amsmath}
\usepackage{natbib}

\title{My Document}
\author{Your Name}
\date{\today}

\begin{document}

\maketitle

\section{Introduction}
This is a sample document demonstrating various \LaTeX{} features including mathematical equations, tables, and formatting. You can cite papers using the citation search feature, and they will be automatically added to your \texttt{cite.bib} file.

\section{Mathematical Expressions}
Here are some example equations:

\subsection{Quadratic Formula}
The quadratic formula is given by:
\[x = \frac{-b \pm \sqrt{b^2 - 4ac}}{2a}\]

\subsection{Matrix Example}
A sample matrix:
\[
\begin{pmatrix}
a_{11} & a_{12} & a_{13} \\
a_{21} & a_{22} & a_{23} \\
a_{31} & a_{32} & a_{33}
\end{pmatrix}
\]

\section{Tables}
Below is a sample table showing data:

\begin{table}[h]
\centering
\begin{tabular}{@{}lrrr@{}}
\toprule
\textbf{Category} & \textbf{Value 1} & \textbf{Value 2} & \textbf{Total} \\
\midrule
Item A & 10.5 & 20.3 & 30.8 \\
Item B & 15.2 & 25.1 & 40.3 \\
Item C & 22.1 & 12.7 & 34.8 \\
\bottomrule
\end{tabular}
\caption{Sample Data Table}
\label{tab:sample}
\end{table}

\section{Lists}
\subsection{Bullet Points}
\begin{itemize}
  \item First item
  \item Second item
  \item Third item with sub-items:
    \begin{itemize}
      \item Sub-item 1
      \item Sub-item 2
    \end{itemize}
\end{itemize}

\subsection{Numbered List}
\begin{enumerate}
  \item First step
  \item Second step
  \item Third step
\end{enumerate}

\section{Citations and References}
When you use the citation search feature, references will be automatically managed. For example, you might cite a paper like this: \cite{example2024}. All citations will be automatically added to your \texttt{cite.bib} file.

\section{Conclusion}
This template demonstrates basic \LaTeX{} features including mathematical formulas, tables, and lists. The bibliography system is now set up to work seamlessly with Bibby AI's citation management. Feel free to modify it for your needs!

\bibliographystyle{plain}
\bibliography{cite}

\end{document}\documentclass{article}
\usepackage[utf8]{inputenc}
\usepackage{booktabs}
\usepackage{amsmath}
\usepackage{natbib}

\title{My Document}
\author{Your Name}
\date{\today}

\begin{document}

\maketitle

\section{Introduction}
This is a sample document demonstrating various \LaTeX{} features including mathematical equations, tables, and formatting. You can cite papers using the citation search feature, and they will be automatically added to your \texttt{cite.bib} file.

\section{Mathematical Expressions}
Here are some example equations:

\subsection{Quadratic Formula}
The quadratic formula is given by:
\[x = \frac{-b \pm \sqrt{b^2 - 4ac}}{2a}\]

\subsection{Matrix Example}
A sample matrix:
\[
\begin{pmatrix}
a_{11} & a_{12} & a_{13} \\
a_{21} & a_{22} & a_{23} \\
a_{31} & a_{32} & a_{33}
\end{pmatrix}
\]

\section{Tables}
Below is a sample table showing data:

\begin{table}[h]
\centering
\begin{tabular}{@{}lrrr@{}}
\toprule
\textbf{Category} & \textbf{Value 1} & \textbf{Value 2} & \textbf{Total} \\
\midrule
Item A & 10.5 & 20.3 & 30.8 \\
Item B & 15.2 & 25.1 & 40.3 \\
Item C & 22.1 & 12.7 & 34.8 \\
\bottomrule
\end{tabular}
\caption{Sample Data Table}
\label{tab:sample}
\end{table}

\section{Lists}
\subsection{Bullet Points}
\begin{itemize}
  \item First item
  \item Second item
  \item Third item with sub-items:
    \begin{itemize}
      \item Sub-item 1
      \item Sub-item 2
    \end{itemize}
\end{itemize}

\subsection{Numbered List}
\begin{enumerate}
  \item First step
  \item Second step
  \item Third step
\end{enumerate}

\section{Citations and References}
When you use the citation search feature, references will be automatically managed. For example, you might cite a paper like this: \cite{example2024}. All citations will be automatically added to your \texttt{cite.bib} file.

\section{Conclusion}
This template demonstrates basic \LaTeX{} features including mathematical formulas, tables, and lists. The bibliography system is now set up to work seamlessly with Bibby AI's citation management. Feel free to modify it for your needs!

\bibliographystyle{plain}
\bibliography{cite}

\end{document}\documentclass{article}
\usepackage[utf8]{inputenc}
\usepackage{booktabs}
\usepackage{amsmath}
\usepackage{natbib}

\title{My Document}
\author{Your Name}
\date{\today}

\begin{document}

\maketitle

\section{Introduction}
This is a sample document demonstrating various \LaTeX{} features including mathematical equations, tables, and formatting. You can cite papers using the citation search feature, and they will be automatically added to your \texttt{cite.bib} file.

\section{Mathematical Expressions}
Here are some example equations:

\subsection{Quadratic Formula}
The quadratic formula is given by:
\[x = \frac{-b \pm \sqrt{b^2 - 4ac}}{2a}\]

\subsection{Matrix Example}
A sample matrix:
\[
\begin{pmatrix}
a_{11} & a_{12} & a_{13} \\
a_{21} & a_{22} & a_{23} \\
a_{31} & a_{32} & a_{33}
\end{pmatrix}
\]

\section{Tables}
Below is a sample table showing data:

\begin{table}[h]
\centering
\begin{tabular}{@{}lrrr@{}}
\toprule
\textbf{Category} & \textbf{Value 1} & \textbf{Value 2} & \textbf{Total} \\
\midrule
Item A & 10.5 & 20.3 & 30.8 \\
Item B & 15.2 & 25.1 & 40.3 \\
Item C & 22.1 & 12.7 & 34.8 \\
\bottomrule
\end{tabular}
\caption{Sample Data Table}
\label{tab:sample}
\end{table}

\section{Lists}
\subsection{Bullet Points}
\begin{itemize}
  \item First item
  \item Second item
  \item Third item with sub-items:
    \begin{itemize}
      \item Sub-item 1
      \item Sub-item 2
    \end{itemize}
\end{itemize}

\subsection{Numbered List}
\begin{enumerate}
  \item First step
  \item Second step
  \item Third step
\end{enumerate}

\section{Citations and References}
When you use the citation search feature, references will be automatically managed. For example, you might cite a paper like this: \cite{example2024}. All citations will be automatically added to your \texttt{cite.bib} file.

\section{Conclusion}
This template demonstrates basic \LaTeX{} features including mathematical formulas, tables, and lists. The bibliography system is now set up to work seamlessly with Bibby AI's citation management. Feel free to modify it for your needs!

\bibliographystyle{plain}
\bibliography{cite}

\end{document}\documentclass{article}
\usepackage[utf8]{inputenc}
\usepackage{booktabs}
\usepackage{amsmath}
\usepackage{natbib}

\title{My Document}
\author{Your Name}
\date{\today}

\begin{document}

\maketitle

\section{Introduction}
This is a sample document demonstrating various \LaTeX{} features including mathematical equations, tables, and formatting. You can cite papers using the citation search feature, and they will be automatically added to your \texttt{cite.bib} file.

\section{Mathematical Expressions}
Here are some example equations:

\subsection{Quadratic Formula}
The quadratic formula is given by:
\[x = \frac{-b \pm \sqrt{b^2 - 4ac}}{2a}\]

\subsection{Matrix Example}
A sample matrix:
\[
\begin{pmatrix}
a_{11} & a_{12} & a_{13} \\
a_{21} & a_{22} & a_{23} \\
a_{31} & a_{32} & a_{33}
\end{pmatrix}
\]

\section{Tables}
Below is a sample table showing data:

\begin{table}[h]
\centering
\begin{tabular}{@{}lrrr@{}}
\toprule
\textbf{Category} & \textbf{Value 1} & \textbf{Value 2} & \textbf{Total} \\
\midrule
Item A & 10.5 & 20.3 & 30.8 \\
Item B & 15.2 & 25.1 & 40.3 \\
Item C & 22.1 & 12.7 & 34.8 \\
\bottomrule
\end{tabular}
\caption{Sample Data Table}
\label{tab:sample}
\end{table}

\section{Lists}
\subsection{Bullet Points}
\begin{itemize}
  \item First item
  \item Second item
  \item Third item with sub-items:
    \begin{itemize}
      \item Sub-item 1
      \item Sub-item 2
    \end{itemize}
\end{itemize}

\subsection{Numbered List}
\begin{enumerate}
  \item First step
  \item Second step
  \item Third step
\end{enumerate}

\section{Citations and References}
When you use the citation search feature, references will be automatically managed. For example, you might cite a paper like this: \cite{example2024}. All citations will be automatically added to your \texttt{cite.bib} file.

\section{Conclusion}
This template demonstrates basic \LaTeX{} features including mathematical formulas, tables, and lists. The bibliography system is now set up to work seamlessly with Bibby AI's citation management. Feel free to modify it for your needs!

\bibliographystyle{plain}
\bibliography{cite}

\end{document}\documentclass{article}
\usepackage[utf8]{inputenc}
\usepackage{booktabs}
\usepackage{amsmath}
\usepackage{natbib}

\title{My Document}
\author{Your Name}
\date{\today}

\begin{document}

\maketitle

\section{Introduction}
This is a sample document demonstrating various \LaTeX{} features including mathematical equations, tables, and formatting. You can cite papers using the citation search feature, and they will be automatically added to your \texttt{cite.bib} file.

\section{Mathematical Expressions}
Here are some example equations:

\subsection{Quadratic Formula}
The quadratic formula is given by:
\[x = \frac{-b \pm \sqrt{b^2 - 4ac}}{2a}\]

\subsection{Matrix Example}
A sample matrix:
\[
\begin{pmatrix}
a_{11} & a_{12} & a_{13} \\
a_{21} & a_{22} & a_{23} \\
a_{31} & a_{32} & a_{33}
\end{pmatrix}
\]

\section{Tables}
Below is a sample table showing data:

\begin{table}[h]
\centering
\begin{tabular}{@{}lrrr@{}}
\toprule
\textbf{Category} & \textbf{Value 1} & \textbf{Value 2} & \textbf{Total} \\
\midrule
Item A & 10.5 & 20.3 & 30.8 \\
Item B & 15.2 & 25.1 & 40.3 \\
Item C & 22.1 & 12.7 & 34.8 \\
\bottomrule
\end{tabular}
\caption{Sample Data Table}
\label{tab:sample}
\end{table}

\section{Lists}
\subsection{Bullet Points}
\begin{itemize}
  \item First item
  \item Second item
  \item Third item with sub-items:
    \begin{itemize}
      \item Sub-item 1
      \item Sub-item 2
    \end{itemize}
\end{itemize}

\subsection{Numbered List}
\begin{enumerate}
  \item First step
  \item Second step
  \item Third step
\end{enumerate}

\section{Citations and References}
When you use the citation search feature, references will be automatically managed. For example, you might cite a paper like this: \cite{example2024}. All citations will be automatically added to your \texttt{cite.bib} file.

\section{Conclusion}
This template demonstrates basic \LaTeX{} features including mathematical formulas, tables, and lists. The bibliography system is now set up to work seamlessly with Bibby AI's citation management. Feel free to modify it for your needs!

\bibliographystyle{plain}
\bibliography{cite}

\end{document}\documentclass{article}
\usepackage[utf8]{inputenc}
\usepackage{booktabs}
\usepackage{amsmath}
\usepackage{natbib}

\title{My Document}
\author{Your Name}
\date{\today}

\begin{document}

\maketitle

\section{Introduction}
This is a sample document demonstrating various \LaTeX{} features including mathematical equations, tables, and formatting. You can cite papers using the citation search feature, and they will be automatically added to your \texttt{cite.bib} file.

\section{Mathematical Expressions}
Here are some example equations:

\subsection{Quadratic Formula}
The quadratic formula is given by:
\[x = \frac{-b \pm \sqrt{b^2 - 4ac}}{2a}\]

\subsection{Matrix Example}
A sample matrix:
\[
\begin{pmatrix}
a_{11} & a_{12} & a_{13} \\
a_{21} & a_{22} & a_{23} \\
a_{31} & a_{32} & a_{33}
\end{pmatrix}
\]

\section{Tables}
Below is a sample table showing data:

\begin{table}[h]
\centering
\begin{tabular}{@{}lrrr@{}}
\toprule
\textbf{Category} & \textbf{Value 1} & \textbf{Value 2} & \textbf{Total} \\
\midrule
Item A & 10.5 & 20.3 & 30.8 \\
Item B & 15.2 & 25.1 & 40.3 \\
Item C & 22.1 & 12.7 & 34.8 \\
\bottomrule
\end{tabular}
\caption{Sample Data Table}
\label{tab:sample}
\end{table}

\section{Lists}
\subsection{Bullet Points}
\begin{itemize}
  \item First item
  \item Second item
  \item Third item with sub-items:
    \begin{itemize}
      \item Sub-item 1
      \item Sub-item 2
    \end{itemize}
\end{itemize}

\subsection{Numbered List}
\begin{enumerate}
  \item First step
  \item Second step
  \item Third step
\end{enumerate}

\section{Citations and References}
When you use the citation search feature, references will be automatically managed. For example, you might cite a paper like this: \cite{example2024}. All citations will be automatically added to your \texttt{cite.bib} file.

\section{Conclusion}
This template demonstrates basic \LaTeX{} features including mathematical formulas, tables, and lists. The bibliography system is now set up to work seamlessly with Bibby AI's citation management. Feel free to modify it for your needs!

\bibliographystyle{plain}
\bibliography{cite}

\end{document}\documentclass{article}
\usepackage[utf8]{inputenc}
\usepackage{booktabs}
\usepackage{amsmath}
\usepackage{natbib}

\title{My Document}
\author{Your Name}
\date{\today}

\begin{document}

\maketitle

\section{Introduction}
This is a sample document demonstrating various \LaTeX{} features including mathematical equations, tables, and formatting. You can cite papers using the citation search feature, and they will be automatically added to your \texttt{cite.bib} file.

\section{Mathematical Expressions}
Here are some example equations:

\subsection{Quadratic Formula}
The quadratic formula is given by:
\[x = \frac{-b \pm \sqrt{b^2 - 4ac}}{2a}\]

\subsection{Matrix Example}
A sample matrix:
\[
\begin{pmatrix}
a_{11} & a_{12} & a_{13} \\
a_{21} & a_{22} & a_{23} \\
a_{31} & a_{32} & a_{33}
\end{pmatrix}
\]

\section{Tables}
Below is a sample table showing data:

\begin{table}[h]
\centering
\begin{tabular}{@{}lrrr@{}}
\toprule
\textbf{Category} & \textbf{Value 1} & \textbf{Value 2} & \textbf{Total} \\
\midrule
Item A & 10.5 & 20.3 & 30.8 \\
Item B & 15.2 & 25.1 & 40.3 \\
Item C & 22.1 & 12.7 & 34.8 \\
\bottomrule
\end{tabular}
\caption{Sample Data Table}
\label{tab:sample}
\end{table}

\section{Lists}
\subsection{Bullet Points}
\begin{itemize}
  \item First item
  \item Second item
  \item Third item with sub-items:
    \begin{itemize}
      \item Sub-item 1
      \item Sub-item 2
    \end{itemize}
\end{itemize}

\subsection{Numbered List}
\begin{enumerate}
  \item First step
  \item Second step
  \item Third step
\end{enumerate}

\section{Citations and References}
When you use the citation search feature, references will be automatically managed. For example, you might cite a paper like this: \cite{example2024}. All citations will be automatically added to your \texttt{cite.bib} file.

\section{Conclusion}
This template demonstrates basic \LaTeX{} features including mathematical formulas, tables, and lists. The bibliography system is now set up to work seamlessly with Bibby AI's citation management. Feel free to modify it for your needs!

\bibliographystyle{plain}
\bibliography{cite}

\end{document}\documentclass{article}
\usepackage[utf8]{inputenc}
\usepackage{booktabs}
\usepackage{amsmath}
\usepackage{natbib}

\title{My Document}
\author{Your Name}
\date{\today}

\begin{document}

\maketitle

\section{Introduction}
This is a sample document demonstrating various \LaTeX{} features including mathematical equations, tables, and formatting. You can cite papers using the citation search feature, and they will be automatically added to your \texttt{cite.bib} file.

\section{Mathematical Expressions}
Here are some example equations:

\subsection{Quadratic Formula}
The quadratic formula is given by:
\[x = \frac{-b \pm \sqrt{b^2 - 4ac}}{2a}\]

\subsection{Matrix Example}
A sample matrix:
\[
\begin{pmatrix}
a_{11} & a_{12} & a_{13} \\
a_{21} & a_{22} & a_{23} \\
a_{31} & a_{32} & a_{33}
\end{pmatrix}
\]

\section{Tables}
Below is a sample table showing data:

\begin{table}[h]
\centering
\begin{tabular}{@{}lrrr@{}}
\toprule
\textbf{Category} & \textbf{Value 1} & \textbf{Value 2} & \textbf{Total} \\
\midrule
Item A & 10.5 & 20.3 & 30.8 \\
Item B & 15.2 & 25.1 & 40.3 \\
Item C & 22.1 & 12.7 & 34.8 \\
\bottomrule
\end{tabular}
\caption{Sample Data Table}
\label{tab:sample}
\end{table}

\section{Lists}
\subsection{Bullet Points}
\begin{itemize}
  \item First item
  \item Second item
  \item Third item with sub-items:
    \begin{itemize}
      \item Sub-item 1
      \item Sub-item 2
    \end{itemize}
\end{itemize}

\subsection{Numbered List}
\begin{enumerate}
  \item First step
  \item Second step
  \item Third step
\end{enumerate}

\section{Citations and References}
When you use the citation search feature, references will be automatically managed. For example, you might cite a paper like this: \cite{example2024}. All citations will be automatically added to your \texttt{cite.bib} file.

\section{Conclusion}
This template demonstrates basic \LaTeX{} features including mathematical formulas, tables, and lists. The bibliography system is now set up to work seamlessly with Bibby AI's citation management. Feel free to modify it for your needs!

\bibliographystyle{plain}
\bibliography{cite}

\end{document}\documentclass{article}
\usepackage[utf8]{inputenc}
\usepackage{booktabs}
\usepackage{amsmath}
\usepackage{natbib}

\title{My Document}
\author{Your Name}
\date{\today}

\begin{document}

\maketitle

\section{Introduction}
This is a sample document demonstrating various \LaTeX{} features including mathematical equations, tables, and formatting. You can cite papers using the citation search feature, and they will be automatically added to your \texttt{cite.bib} file.

\section{Mathematical Expressions}
Here are some example equations:

\subsection{Quadratic Formula}
The quadratic formula is given by:
\[x = \frac{-b \pm \sqrt{b^2 - 4ac}}{2a}\]

\subsection{Matrix Example}
A sample matrix:
\[
\begin{pmatrix}
a_{11} & a_{12} & a_{13} \\
a_{21} & a_{22} & a_{23} \\
a_{31} & a_{32} & a_{33}
\end{pmatrix}
\]

\section{Tables}
Below is a sample table showing data:

\begin{table}[h]
\centering
\begin{tabular}{@{}lrrr@{}}
\toprule
\textbf{Category} & \textbf{Value 1} & \textbf{Value 2} & \textbf{Total} \\
\midrule
Item A & 10.5 & 20.3 & 30.8 \\
Item B & 15.2 & 25.1 & 40.3 \\
Item C & 22.1 & 12.7 & 34.8 \\
\bottomrule
\end{tabular}
\caption{Sample Data Table}
\label{tab:sample}
\end{table}

\section{Lists}
\subsection{Bullet Points}
\begin{itemize}
  \item First item
  \item Second item
  \item Third item with sub-items:
    \begin{itemize}
      \item Sub-item 1
      \item Sub-item 2
    \end{itemize}
\end{itemize}

\subsection{Numbered List}
\begin{enumerate}
  \item First step
  \item Second step
  \item Third step
\end{enumerate}

\section{Citations and References}
When you use the citation search feature, references will be automatically managed. For example, you might cite a paper like this: \cite{example2024}. All citations will be automatically added to your \texttt{cite.bib} file.

\section{Conclusion}
This template demonstrates basic \LaTeX{} features including mathematical formulas, tables, and lists. The bibliography system is now set up to work seamlessly with Bibby AI's citation management. Feel free to modify it for your needs!

\bibliographystyle{plain}
\bibliography{cite}

\end{document}\documentclass{article}
\usepackage[utf8]{inputenc}
\usepackage{booktabs}
\usepackage{amsmath}
\usepackage{natbib}

\title{My Document}
\author{Your Name}
\date{\today}

\begin{document}

\maketitle

\section{Introduction}
This is a sample document demonstrating various \LaTeX{} features including mathematical equations, tables, and formatting. You can cite papers using the citation search feature, and they will be automatically added to your \texttt{cite.bib} file.

\section{Mathematical Expressions}
Here are some example equations:

\subsection{Quadratic Formula}
The quadratic formula is given by:
\[x = \frac{-b \pm \sqrt{b^2 - 4ac}}{2a}\]

\subsection{Matrix Example}
A sample matrix:
\[
\begin{pmatrix}
a_{11} & a_{12} & a_{13} \\
a_{21} & a_{22} & a_{23} \\
a_{31} & a_{32} & a_{33}
\end{pmatrix}
\]

\section{Tables}
Below is a sample table showing data:

\begin{table}[h]
\centering
\begin{tabular}{@{}lrrr@{}}
\toprule
\textbf{Category} & \textbf{Value 1} & \textbf{Value 2} & \textbf{Total} \\
\midrule
Item A & 10.5 & 20.3 & 30.8 \\
Item B & 15.2 & 25.1 & 40.3 \\
Item C & 22.1 & 12.7 & 34.8 \\
\bottomrule
\end{tabular}
\caption{Sample Data Table}
\label{tab:sample}
\end{table}

\section{Lists}
\subsection{Bullet Points}
\begin{itemize}
  \item First item
  \item Second item
  \item Third item with sub-items:
    \begin{itemize}
      \item Sub-item 1
      \item Sub-item 2
    \end{itemize}
\end{itemize}

\subsection{Numbered List}
\begin{enumerate}
  \item First step
  \item Second step
  \item Third step
\end{enumerate}

\section{Citations and References}
When you use the citation search feature, references will be automatically managed. For example, you might cite a paper like this: \cite{example2024}. All citations will be automatically added to your \texttt{cite.bib} file.

\section{Conclusion}
This template demonstrates basic \LaTeX{} features including mathematical formulas, tables, and lists. The bibliography system is now set up to work seamlessly with Bibby AI's citation management. Feel free to modify it for your needs!

\bibliographystyle{plain}
\bibliography{cite}

\end{document}\documentclass{article}
\usepackage[utf8]{inputenc}
\usepackage{booktabs}
\usepackage{amsmath}
\usepackage{natbib}

\title{My Document}
\author{Your Name}
\date{\today}

\begin{document}

\maketitle

\section{Introduction}
This is a sample document demonstrating various \LaTeX{} features including mathematical equations, tables, and formatting. You can cite papers using the citation search feature, and they will be automatically added to your \texttt{cite.bib} file.

\section{Mathematical Expressions}
Here are some example equations:

\subsection{Quadratic Formula}
The quadratic formula is given by:
\[x = \frac{-b \pm \sqrt{b^2 - 4ac}}{2a}\]

\subsection{Matrix Example}
A sample matrix:
\[
\begin{pmatrix}
a_{11} & a_{12} & a_{13} \\
a_{21} & a_{22} & a_{23} \\
a_{31} & a_{32} & a_{33}
\end{pmatrix}
\]

\section{Tables}
Below is a sample table showing data:

\begin{table}[h]
\centering
\begin{tabular}{@{}lrrr@{}}
\toprule
\textbf{Category} & \textbf{Value 1} & \textbf{Value 2} & \textbf{Total} \\
\midrule
Item A & 10.5 & 20.3 & 30.8 \\
Item B & 15.2 & 25.1 & 40.3 \\
Item C & 22.1 & 12.7 & 34.8 \\
\bottomrule
\end{tabular}
\caption{Sample Data Table}
\label{tab:sample}
\end{table}

\section{Lists}
\subsection{Bullet Points}
\begin{itemize}
  \item First item
  \item Second item
  \item Third item with sub-items:
    \begin{itemize}
      \item Sub-item 1
      \item Sub-item 2
    \end{itemize}
\end{itemize}

\subsection{Numbered List}
\begin{enumerate}
  \item First step
  \item Second step
  \item Third step
\end{enumerate}

\section{Citations and References}
When you use the citation search feature, references will be automatically managed. For example, you might cite a paper like this: \cite{example2024}. All citations will be automatically added to your \texttt{cite.bib} file.

\section{Conclusion}
This template demonstrates basic \LaTeX{} features including mathematical formulas, tables, and lists. The bibliography system is now set up to work seamlessly with Bibby AI's citation management. Feel free to modify it for your needs!

\bibliographystyle{plain}
\bibliography{cite}

\end{document}\documentclass{article}
\usepackage[utf8]{inputenc}
\usepackage{booktabs}
\usepackage{amsmath}
\usepackage{natbib}

\title{My Document}
\author{Your Name}
\date{\today}

\begin{document}

\maketitle

\section{Introduction}
This is a sample document demonstrating various \LaTeX{} features including mathematical equations, tables, and formatting. You can cite papers using the citation search feature, and they will be automatically added to your \texttt{cite.bib} file.

\section{Mathematical Expressions}
Here are some example equations:

\subsection{Quadratic Formula}
The quadratic formula is given by:
\[x = \frac{-b \pm \sqrt{b^2 - 4ac}}{2a}\]

\subsection{Matrix Example}
A sample matrix:
\[
\begin{pmatrix}
a_{11} & a_{12} & a_{13} \\
a_{21} & a_{22} & a_{23} \\
a_{31} & a_{32} & a_{33}
\end{pmatrix}
\]

\section{Tables}
Below is a sample table showing data:

\begin{table}[h]
\centering
\begin{tabular}{@{}lrrr@{}}
\toprule
\textbf{Category} & \textbf{Value 1} & \textbf{Value 2} & \textbf{Total} \\
\midrule
Item A & 10.5 & 20.3 & 30.8 \\
Item B & 15.2 & 25.1 & 40.3 \\
Item C & 22.1 & 12.7 & 34.8 \\
\bottomrule
\end{tabular}
\caption{Sample Data Table}
\label{tab:sample}
\end{table}

\section{Lists}
\subsection{Bullet Points}
\begin{itemize}
  \item First item
  \item Second item
  \item Third item with sub-items:
    \begin{itemize}
      \item Sub-item 1
      \item Sub-item 2
    \end{itemize}
\end{itemize}

\subsection{Numbered List}
\begin{enumerate}
  \item First step
  \item Second step
  \item Third step
\end{enumerate}

\section{Citations and References}
When you use the citation search feature, references will be automatically managed. For example, you might cite a paper like this: \cite{example2024}. All citations will be automatically added to your \texttt{cite.bib} file.

\section{Conclusion}
This template demonstrates basic \LaTeX{} features including mathematical formulas, tables, and lists. The bibliography system is now set up to work seamlessly with Bibby AI's citation management. Feel free to modify it for your needs!

\bibliographystyle{plain}
\bibliography{cite}

\end{document}\documentclass{article}
\usepackage[utf8]{inputenc}
\usepackage{booktabs}
\usepackage{amsmath}
\usepackage{natbib}

\title{My Document}
\author{Your Name}
\date{\today}

\begin{document}

\maketitle

\section{Introduction}
This is a sample document demonstrating various \LaTeX{} features including mathematical equations, tables, and formatting. You can cite papers using the citation search feature, and they will be automatically added to your \texttt{cite.bib} file.

\section{Mathematical Expressions}
Here are some example equations:

\subsection{Quadratic Formula}
The quadratic formula is given by:
\[x = \frac{-b \pm \sqrt{b^2 - 4ac}}{2a}\]

\subsection{Matrix Example}
A sample matrix:
\[
\begin{pmatrix}
a_{11} & a_{12} & a_{13} \\
a_{21} & a_{22} & a_{23} \\
a_{31} & a_{32} & a_{33}
\end{pmatrix}
\]

\section{Tables}
Below is a sample table showing data:

\begin{table}[h]
\centering
\begin{tabular}{@{}lrrr@{}}
\toprule
\textbf{Category} & \textbf{Value 1} & \textbf{Value 2} & \textbf{Total} \\
\midrule
Item A & 10.5 & 20.3 & 30.8 \\
Item B & 15.2 & 25.1 & 40.3 \\
Item C & 22.1 & 12.7 & 34.8 \\
\bottomrule
\end{tabular}
\caption{Sample Data Table}
\label{tab:sample}
\end{table}

\section{Lists}
\subsection{Bullet Points}
\begin{itemize}
  \item First item
  \item Second item
  \item Third item with sub-items:
    \begin{itemize}
      \item Sub-item 1
      \item Sub-item 2
    \end{itemize}
\end{itemize}

\subsection{Numbered List}
\begin{enumerate}
  \item First step
  \item Second step
  \item Third step
\end{enumerate}

\section{Citations and References}
When you use the citation search feature, references will be automatically managed. For example, you might cite a paper like this: \cite{example2024}. All citations will be automatically added to your \texttt{cite.bib} file.

\section{Conclusion}
This template demonstrates basic \LaTeX{} features including mathematical formulas, tables, and lists. The bibliography system is now set up to work seamlessly with Bibby AI's citation management. Feel free to modify it for your needs!

\bibliographystyle{plain}
\bibliography{cite}

\end{document}\documentclass{article}
\usepackage[utf8]{inputenc}
\usepackage{booktabs}
\usepackage{amsmath}
\usepackage{natbib}

\title{My Document}
\author{Your Name}
\date{\today}

\begin{document}

\maketitle

\section{Introduction}
This is a sample document demonstrating various \LaTeX{} features including mathematical equations, tables, and formatting. You can cite papers using the citation search feature, and they will be automatically added to your \texttt{cite.bib} file.

\section{Mathematical Expressions}
Here are some example equations:

\subsection{Quadratic Formula}
The quadratic formula is given by:
\[x = \frac{-b \pm \sqrt{b^2 - 4ac}}{2a}\]

\subsection{Matrix Example}
A sample matrix:
\[
\begin{pmatrix}
a_{11} & a_{12} & a_{13} \\
a_{21} & a_{22} & a_{23} \\
a_{31} & a_{32} & a_{33}
\end{pmatrix}
\]

\section{Tables}
Below is a sample table showing data:

\begin{table}[h]
\centering
\begin{tabular}{@{}lrrr@{}}
\toprule
\textbf{Category} & \textbf{Value 1} & \textbf{Value 2} & \textbf{Total} \\
\midrule
Item A & 10.5 & 20.3 & 30.8 \\
Item B & 15.2 & 25.1 & 40.3 \\
Item C & 22.1 & 12.7 & 34.8 \\
\bottomrule
\end{tabular}
\caption{Sample Data Table}
\label{tab:sample}
\end{table}

\section{Lists}
\subsection{Bullet Points}
\begin{itemize}
  \item First item
  \item Second item
  \item Third item with sub-items:
    \begin{itemize}
      \item Sub-item 1
      \item Sub-item 2
    \end{itemize}
\end{itemize}

\subsection{Numbered List}
\begin{enumerate}
  \item First step
  \item Second step
  \item Third step
\end{enumerate}

\section{Citations and References}
When you use the citation search feature, references will be automatically managed. For example, you might cite a paper like this: \cite{example2024}. All citations will be automatically added to your \texttt{cite.bib} file.

\section{Conclusion}
This template demonstrates basic \LaTeX{} features including mathematical formulas, tables, and lists. The bibliography system is now set up to work seamlessly with Bibby AI's citation management. Feel free to modify it for your needs!

\bibliographystyle{plain}
\bibliography{cite}

\end{document}\documentclass{article}
\usepackage[utf8]{inputenc}
\usepackage{booktabs}
\usepackage{amsmath}
\usepackage{natbib}

\title{My Document}
\author{Your Name}
\date{\today}

\begin{document}

\maketitle

\section{Introduction}
This is a sample document demonstrating various \LaTeX{} features including mathematical equations, tables, and formatting. You can cite papers using the citation search feature, and they will be automatically added to your \texttt{cite.bib} file.

\section{Mathematical Expressions}
Here are some example equations:

\subsection{Quadratic Formula}
The quadratic formula is given by:
\[x = \frac{-b \pm \sqrt{b^2 - 4ac}}{2a}\]

\subsection{Matrix Example}
A sample matrix:
\[
\begin{pmatrix}
a_{11} & a_{12} & a_{13} \\
a_{21} & a_{22} & a_{23} \\
a_{31} & a_{32} & a_{33}
\end{pmatrix}
\]

\section{Tables}
Below is a sample table showing data:

\begin{table}[h]
\centering
\begin{tabular}{@{}lrrr@{}}
\toprule
\textbf{Category} & \textbf{Value 1} & \textbf{Value 2} & \textbf{Total} \\
\midrule
Item A & 10.5 & 20.3 & 30.8 \\
Item B & 15.2 & 25.1 & 40.3 \\
Item C & 22.1 & 12.7 & 34.8 \\
\bottomrule
\end{tabular}
\caption{Sample Data Table}
\label{tab:sample}
\end{table}

\section{Lists}
\subsection{Bullet Points}
\begin{itemize}
  \item First item
  \item Second item
  \item Third item with sub-items:
    \begin{itemize}
      \item Sub-item 1
      \item Sub-item 2
    \end{itemize}
\end{itemize}

\subsection{Numbered List}
\begin{enumerate}
  \item First step
  \item Second step
  \item Third step
\end{enumerate}

\section{Citations and References}
When you use the citation search feature, references will be automatically managed. For example, you might cite a paper like this: \cite{example2024}. All citations will be automatically added to your \texttt{cite.bib} file.

\section{Conclusion}
This template demonstrates basic \LaTeX{} features including mathematical formulas, tables, and lists. The bibliography system is now set up to work seamlessly with Bibby AI's citation management. Feel free to modify it for your needs!

\bibliographystyle{plain}
\bibliography{cite}

\end{document}\documentclass{article}
\usepackage[utf8]{inputenc}
\usepackage{booktabs}
\usepackage{amsmath}
\usepackage{natbib}

\title{My Document}
\author{Your Name}
\date{\today}

\begin{document}

\maketitle

\section{Introduction}
This is a sample document demonstrating various \LaTeX{} features including mathematical equations, tables, and formatting. You can cite papers using the citation search feature, and they will be automatically added to your \texttt{cite.bib} file.

\section{Mathematical Expressions}
Here are some example equations:

\subsection{Quadratic Formula}
The quadratic formula is given by:
\[x = \frac{-b \pm \sqrt{b^2 - 4ac}}{2a}\]

\subsection{Matrix Example}
A sample matrix:
\[
\begin{pmatrix}
a_{11} & a_{12} & a_{13} \\
a_{21} & a_{22} & a_{23} \\
a_{31} & a_{32} & a_{33}
\end{pmatrix}
\]

\section{Tables}
Below is a sample table showing data:

\begin{table}[h]
\centering
\begin{tabular}{@{}lrrr@{}}
\toprule
\textbf{Category} & \textbf{Value 1} & \textbf{Value 2} & \textbf{Total} \\
\midrule
Item A & 10.5 & 20.3 & 30.8 \\
Item B & 15.2 & 25.1 & 40.3 \\
Item C & 22.1 & 12.7 & 34.8 \\
\bottomrule
\end{tabular}
\caption{Sample Data Table}
\label{tab:sample}
\end{table}

\section{Lists}
\subsection{Bullet Points}
\begin{itemize}
  \item First item
  \item Second item
  \item Third item with sub-items:
    \begin{itemize}
      \item Sub-item 1
      \item Sub-item 2
    \end{itemize}
\end{itemize}

\subsection{Numbered List}
\begin{enumerate}
  \item First step
  \item Second step
  \item Third step
\end{enumerate}

\section{Citations and References}
When you use the citation search feature, references will be automatically managed. For example, you might cite a paper like this: \cite{example2024}. All citations will be automatically added to your \texttt{cite.bib} file.

\section{Conclusion}
This template demonstrates basic \LaTeX{} features including mathematical formulas, tables, and lists. The bibliography system is now set up to work seamlessly with Bibby AI's citation management. Feel free to modify it for your needs!

\bibliographystyle{plain}
\bibliography{cite}

\end{document}\documentclass{article}
\usepackage[utf8]{inputenc}
\usepackage{booktabs}
\usepackage{amsmath}
\usepackage{natbib}

\title{My Document}
\author{Your Name}
\date{\today}

\begin{document}

\maketitle

\section{Introduction}
This is a sample document demonstrating various \LaTeX{} features including mathematical equations, tables, and formatting. You can cite papers using the citation search feature, and they will be automatically added to your \texttt{cite.bib} file.

\section{Mathematical Expressions}
Here are some example equations:

\subsection{Quadratic Formula}
The quadratic formula is given by:
\[x = \frac{-b \pm \sqrt{b^2 - 4ac}}{2a}\]

\subsection{Matrix Example}
A sample matrix:
\[
\begin{pmatrix}
a_{11} & a_{12} & a_{13} \\
a_{21} & a_{22} & a_{23} \\
a_{31} & a_{32} & a_{33}
\end{pmatrix}
\]

\section{Tables}
Below is a sample table showing data:

\begin{table}[h]
\centering
\begin{tabular}{@{}lrrr@{}}
\toprule
\textbf{Category} & \textbf{Value 1} & \textbf{Value 2} & \textbf{Total} \\
\midrule
Item A & 10.5 & 20.3 & 30.8 \\
Item B & 15.2 & 25.1 & 40.3 \\
Item C & 22.1 & 12.7 & 34.8 \\
\bottomrule
\end{tabular}
\caption{Sample Data Table}
\label{tab:sample}
\end{table}

\section{Lists}
\subsection{Bullet Points}
\begin{itemize}
  \item First item
  \item Second item
  \item Third item with sub-items:
    \begin{itemize}
      \item Sub-item 1
      \item Sub-item 2
    \end{itemize}
\end{itemize}

\subsection{Numbered List}
\begin{enumerate}
  \item First step
  \item Second step
  \item Third step
\end{enumerate}

\section{Citations and References}
When you use the citation search feature, references will be automatically managed. For example, you might cite a paper like this: \cite{example2024}. All citations will be automatically added to your \texttt{cite.bib} file.

\section{Conclusion}
This template demonstrates basic \LaTeX{} features including mathematical formulas, tables, and lists. The bibliography system is now set up to work seamlessly with Bibby AI's citation management. Feel free to modify it for your needs!

\bibliographystyle{plain}
\bibliography{cite}

\end{document}\documentclass{article}
\usepackage[utf8]{inputenc}
\usepackage{booktabs}
\usepackage{amsmath}
\usepackage{natbib}

\title{My Document}
\author{Your Name}
\date{\today}

\begin{document}

\maketitle

\section{Introduction}
This is a sample document demonstrating various \LaTeX{} features including mathematical equations, tables, and formatting. You can cite papers using the citation search feature, and they will be automatically added to your \texttt{cite.bib} file.

\section{Mathematical Expressions}
Here are some example equations:

\subsection{Quadratic Formula}
The quadratic formula is given by:
\[x = \frac{-b \pm \sqrt{b^2 - 4ac}}{2a}\]

\subsection{Matrix Example}
A sample matrix:
\[
\begin{pmatrix}
a_{11} & a_{12} & a_{13} \\
a_{21} & a_{22} & a_{23} \\
a_{31} & a_{32} & a_{33}
\end{pmatrix}
\]

\section{Tables}
Below is a sample table showing data:

\begin{table}[h]
\centering
\begin{tabular}{@{}lrrr@{}}
\toprule
\textbf{Category} & \textbf{Value 1} & \textbf{Value 2} & \textbf{Total} \\
\midrule
Item A & 10.5 & 20.3 & 30.8 \\
Item B & 15.2 & 25.1 & 40.3 \\
Item C & 22.1 & 12.7 & 34.8 \\
\bottomrule
\end{tabular}
\caption{Sample Data Table}
\label{tab:sample}
\end{table}

\section{Lists}
\subsection{Bullet Points}
\begin{itemize}
  \item First item
  \item Second item
  \item Third item with sub-items:
    \begin{itemize}
      \item Sub-item 1
      \item Sub-item 2
    \end{itemize}
\end{itemize}

\subsection{Numbered List}
\begin{enumerate}
  \item First step
  \item Second step
  \item Third step
\end{enumerate}

\section{Citations and References}
When you use the citation search feature, references will be automatically managed. For example, you might cite a paper like this: \cite{example2024}. All citations will be automatically added to your \texttt{cite.bib} file.

\section{Conclusion}
This template demonstrates basic \LaTeX{} features including mathematical formulas, tables, and lists. The bibliography system is now set up to work seamlessly with Bibby AI's citation management. Feel free to modify it for your needs!

\bibliographystyle{plain}
\bibliography{cite}

\end{document}\documentclass{article}
\usepackage[utf8]{inputenc}
\usepackage{booktabs}
\usepackage{amsmath}
\usepackage{natbib}

\title{My Document}
\author{Your Name}
\date{\today}

\begin{document}

\maketitle

\section{Introduction}
This is a sample document demonstrating various \LaTeX{} features including mathematical equations, tables, and formatting. You can cite papers using the citation search feature, and they will be automatically added to your \texttt{cite.bib} file.

\section{Mathematical Expressions}
Here are some example equations:

\subsection{Quadratic Formula}
The quadratic formula is given by:
\[x = \frac{-b \pm \sqrt{b^2 - 4ac}}{2a}\]

\subsection{Matrix Example}
A sample matrix:
\[
\begin{pmatrix}
a_{11} & a_{12} & a_{13} \\
a_{21} & a_{22} & a_{23} \\
a_{31} & a_{32} & a_{33}
\end{pmatrix}
\]

\section{Tables}
Below is a sample table showing data:

\begin{table}[h]
\centering
\begin{tabular}{@{}lrrr@{}}
\toprule
\textbf{Category} & \textbf{Value 1} & \textbf{Value 2} & \textbf{Total} \\
\midrule
Item A & 10.5 & 20.3 & 30.8 \\
Item B & 15.2 & 25.1 & 40.3 \\
Item C & 22.1 & 12.7 & 34.8 \\
\bottomrule
\end{tabular}
\caption{Sample Data Table}
\label{tab:sample}
\end{table}

\section{Lists}
\subsection{Bullet Points}
\begin{itemize}
  \item First item
  \item Second item
  \item Third item with sub-items:
    \begin{itemize}
      \item Sub-item 1
      \item Sub-item 2
    \end{itemize}
\end{itemize}

\subsection{Numbered List}
\begin{enumerate}
  \item First step
  \item Second step
  \item Third step
\end{enumerate}

\section{Citations and References}
When you use the citation search feature, references will be automatically managed. For example, you might cite a paper like this: \cite{example2024}. All citations will be automatically added to your \texttt{cite.bib} file.

\section{Conclusion}
This template demonstrates basic \LaTeX{} features including mathematical formulas, tables, and lists. The bibliography system is now set up to work seamlessly with Bibby AI's citation management. Feel free to modify it for your needs!

\bibliographystyle{plain}
\bibliography{cite}

\end{document}\documentclass{article}
\usepackage[utf8]{inputenc}
\usepackage{booktabs}
\usepackage{amsmath}
\usepackage{natbib}

\title{My Document}
\author{Your Name}
\date{\today}

\begin{document}

\maketitle

\section{Introduction}
This is a sample document demonstrating various \LaTeX{} features including mathematical equations, tables, and formatting. You can cite papers using the citation search feature, and they will be automatically added to your \texttt{cite.bib} file.

\section{Mathematical Expressions}
Here are some example equations:

\subsection{Quadratic Formula}
The quadratic formula is given by:
\[x = \frac{-b \pm \sqrt{b^2 - 4ac}}{2a}\]

\subsection{Matrix Example}
A sample matrix:
\[
\begin{pmatrix}
a_{11} & a_{12} & a_{13} \\
a_{21} & a_{22} & a_{23} \\
a_{31} & a_{32} & a_{33}
\end{pmatrix}
\]

\section{Tables}
Below is a sample table showing data:

\begin{table}[h]
\centering
\begin{tabular}{@{}lrrr@{}}
\toprule
\textbf{Category} & \textbf{Value 1} & \textbf{Value 2} & \textbf{Total} \\
\midrule
Item A & 10.5 & 20.3 & 30.8 \\
Item B & 15.2 & 25.1 & 40.3 \\
Item C & 22.1 & 12.7 & 34.8 \\
\bottomrule
\end{tabular}
\caption{Sample Data Table}
\label{tab:sample}
\end{table}

\section{Lists}
\subsection{Bullet Points}
\begin{itemize}
  \item First item
  \item Second item
  \item Third item with sub-items:
    \begin{itemize}
      \item Sub-item 1
      \item Sub-item 2
    \end{itemize}
\end{itemize}

\subsection{Numbered List}
\begin{enumerate}
  \item First step
  \item Second step
  \item Third step
\end{enumerate}

\section{Citations and References}
When you use the citation search feature, references will be automatically managed. For example, you might cite a paper like this: \cite{example2024}. All citations will be automatically added to your \texttt{cite.bib} file.

\section{Conclusion}
This template demonstrates basic \LaTeX{} features including mathematical formulas, tables, and lists. The bibliography system is now set up to work seamlessly with Bibby AI's citation management. Feel free to modify it for your needs!

\bibliographystyle{plain}
\bibliography{cite}

\end{document}\documentclass{article}
\usepackage[utf8]{inputenc}
\usepackage{booktabs}
\usepackage{amsmath}
\usepackage{natbib}

\title{My Document}
\author{Your Name}
\date{\today}

\begin{document}

\maketitle

\section{Introduction}
This is a sample document demonstrating various \LaTeX{} features including mathematical equations, tables, and formatting. You can cite papers using the citation search feature, and they will be automatically added to your \texttt{cite.bib} file.

\section{Mathematical Expressions}
Here are some example equations:

\subsection{Quadratic Formula}
The quadratic formula is given by:
\[x = \frac{-b \pm \sqrt{b^2 - 4ac}}{2a}\]

\subsection{Matrix Example}
A sample matrix:
\[
\begin{pmatrix}
a_{11} & a_{12} & a_{13} \\
a_{21} & a_{22} & a_{23} \\
a_{31} & a_{32} & a_{33}
\end{pmatrix}
\]

\section{Tables}
Below is a sample table showing data:

\begin{table}[h]
\centering
\begin{tabular}{@{}lrrr@{}}
\toprule
\textbf{Category} & \textbf{Value 1} & \textbf{Value 2} & \textbf{Total} \\
\midrule
Item A & 10.5 & 20.3 & 30.8 \\
Item B & 15.2 & 25.1 & 40.3 \\
Item C & 22.1 & 12.7 & 34.8 \\
\bottomrule
\end{tabular}
\caption{Sample Data Table}
\label{tab:sample}
\end{table}

\section{Lists}
\subsection{Bullet Points}
\begin{itemize}
  \item First item
  \item Second item
  \item Third item with sub-items:
    \begin{itemize}
      \item Sub-item 1
      \item Sub-item 2
    \end{itemize}
\end{itemize}

\subsection{Numbered List}
\begin{enumerate}
  \item First step
  \item Second step
  \item Third step
\end{enumerate}

\section{Citations and References}
When you use the citation search feature, references will be automatically managed. For example, you might cite a paper like this: \cite{example2024}. All citations will be automatically added to your \texttt{cite.bib} file.

\section{Conclusion}
This template demonstrates basic \LaTeX{} features including mathematical formulas, tables, and lists. The bibliography system is now set up to work seamlessly with Bibby AI's citation management. Feel free to modify it for your needs!

\bibliographystyle{plain}
\bibliography{cite}

\end{document}\documentclass{article}
\usepackage[utf8]{inputenc}
\usepackage{booktabs}
\usepackage{amsmath}
\usepackage{natbib}

\title{My Document}
\author{Your Name}
\date{\today}

\begin{document}

\maketitle

\section{Introduction}
This is a sample document demonstrating various \LaTeX{} features including mathematical equations, tables, and formatting. You can cite papers using the citation search feature, and they will be automatically added to your \texttt{cite.bib} file.

\section{Mathematical Expressions}
Here are some example equations:

\subsection{Quadratic Formula}
The quadratic formula is given by:
\[x = \frac{-b \pm \sqrt{b^2 - 4ac}}{2a}\]

\subsection{Matrix Example}
A sample matrix:
\[
\begin{pmatrix}
a_{11} & a_{12} & a_{13} \\
a_{21} & a_{22} & a_{23} \\
a_{31} & a_{32} & a_{33}
\end{pmatrix}
\]

\section{Tables}
Below is a sample table showing data:

\begin{table}[h]
\centering
\begin{tabular}{@{}lrrr@{}}
\toprule
\textbf{Category} & \textbf{Value 1} & \textbf{Value 2} & \textbf{Total} \\
\midrule
Item A & 10.5 & 20.3 & 30.8 \\
Item B & 15.2 & 25.1 & 40.3 \\
Item C & 22.1 & 12.7 & 34.8 \\
\bottomrule
\end{tabular}
\caption{Sample Data Table}
\label{tab:sample}
\end{table}

\section{Lists}
\subsection{Bullet Points}
\begin{itemize}
  \item First item
  \item Second item
  \item Third item with sub-items:
    \begin{itemize}
      \item Sub-item 1
      \item Sub-item 2
    \end{itemize}
\end{itemize}

\subsection{Numbered List}
\begin{enumerate}
  \item First step
  \item Second step
  \item Third step
\end{enumerate}

\section{Citations and References}
When you use the citation search feature, references will be automatically managed. For example, you might cite a paper like this: \cite{example2024}. All citations will be automatically added to your \texttt{cite.bib} file.

\section{Conclusion}
This template demonstrates basic \LaTeX{} features including mathematical formulas, tables, and lists. The bibliography system is now set up to work seamlessly with Bibby AI's citation management. Feel free to modify it for your needs!

\bibliographystyle{plain}
\bibliography{cite}

\end{document}\documentclass{article}
\usepackage[utf8]{inputenc}
\usepackage{booktabs}
\usepackage{amsmath}
\usepackage{natbib}

\title{My Document}
\author{Your Name}
\date{\today}

\begin{document}

\maketitle

\section{Introduction}
This is a sample document demonstrating various \LaTeX{} features including mathematical equations, tables, and formatting. You can cite papers using the citation search feature, and they will be automatically added to your \texttt{cite.bib} file.

\section{Mathematical Expressions}
Here are some example equations:

\subsection{Quadratic Formula}
The quadratic formula is given by:
\[x = \frac{-b \pm \sqrt{b^2 - 4ac}}{2a}\]

\subsection{Matrix Example}
A sample matrix:
\[
\begin{pmatrix}
a_{11} & a_{12} & a_{13} \\
a_{21} & a_{22} & a_{23} \\
a_{31} & a_{32} & a_{33}
\end{pmatrix}
\]

\section{Tables}
Below is a sample table showing data:

\begin{table}[h]
\centering
\begin{tabular}{@{}lrrr@{}}
\toprule
\textbf{Category} & \textbf{Value 1} & \textbf{Value 2} & \textbf{Total} \\
\midrule
Item A & 10.5 & 20.3 & 30.8 \\
Item B & 15.2 & 25.1 & 40.3 \\
Item C & 22.1 & 12.7 & 34.8 \\
\bottomrule
\end{tabular}
\caption{Sample Data Table}
\label{tab:sample}
\end{table}

\section{Lists}
\subsection{Bullet Points}
\begin{itemize}
  \item First item
  \item Second item
  \item Third item with sub-items:
    \begin{itemize}
      \item Sub-item 1
      \item Sub-item 2
    \end{itemize}
\end{itemize}

\subsection{Numbered List}
\begin{enumerate}
  \item First step
  \item Second step
  \item Third step
\end{enumerate}

\section{Citations and References}
When you use the citation search feature, references will be automatically managed. For example, you might cite a paper like this: \cite{example2024}. All citations will be automatically added to your \texttt{cite.bib} file.

\section{Conclusion}
This template demonstrates basic \LaTeX{} features including mathematical formulas, tables, and lists. The bibliography system is now set up to work seamlessly with Bibby AI's citation management. Feel free to modify it for your needs!

\bibliographystyle{plain}
\bibliography{cite}

\end{document}\documentclass{article}
\usepackage[utf8]{inputenc}
\usepackage{booktabs}
\usepackage{amsmath}
\usepackage{natbib}

\title{My Document}
\author{Your Name}
\date{\today}

\begin{document}

\maketitle

\section{Introduction}
This is a sample document demonstrating various \LaTeX{} features including mathematical equations, tables, and formatting. You can cite papers using the citation search feature, and they will be automatically added to your \texttt{cite.bib} file.

\section{Mathematical Expressions}
Here are some example equations:

\subsection{Quadratic Formula}
The quadratic formula is given by:
\[x = \frac{-b \pm \sqrt{b^2 - 4ac}}{2a}\]

\subsection{Matrix Example}
A sample matrix:
\[
\begin{pmatrix}
a_{11} & a_{12} & a_{13} \\
a_{21} & a_{22} & a_{23} \\
a_{31} & a_{32} & a_{33}
\end{pmatrix}
\]

\section{Tables}
Below is a sample table showing data:

\begin{table}[h]
\centering
\begin{tabular}{@{}lrrr@{}}
\toprule
\textbf{Category} & \textbf{Value 1} & \textbf{Value 2} & \textbf{Total} \\
\midrule
Item A & 10.5 & 20.3 & 30.8 \\
Item B & 15.2 & 25.1 & 40.3 \\
Item C & 22.1 & 12.7 & 34.8 \\
\bottomrule
\end{tabular}
\caption{Sample Data Table}
\label{tab:sample}
\end{table}

\section{Lists}
\subsection{Bullet Points}
\begin{itemize}
  \item First item
  \item Second item
  \item Third item with sub-items:
    \begin{itemize}
      \item Sub-item 1
      \item Sub-item 2
    \end{itemize}
\end{itemize}

\subsection{Numbered List}
\begin{enumerate}
  \item First step
  \item Second step
  \item Third step
\end{enumerate}

\section{Citations and References}
When you use the citation search feature, references will be automatically managed. For example, you might cite a paper like this: \cite{example2024}. All citations will be automatically added to your \texttt{cite.bib} file.

\section{Conclusion}
This template demonstrates basic \LaTeX{} features including mathematical formulas, tables, and lists. The bibliography system is now set up to work seamlessly with Bibby AI's citation management. Feel free to modify it for your needs!

\bibliographystyle{plain}
\bibliography{cite}

\end{document}\documentclass{article}
\usepackage[utf8]{inputenc}
\usepackage{booktabs}
\usepackage{amsmath}
\usepackage{natbib}

\title{My Document}
\author{Your Name}
\date{\today}

\begin{document}

\maketitle

\section{Introduction}
This is a sample document demonstrating various \LaTeX{} features including mathematical equations, tables, and formatting. You can cite papers using the citation search feature, and they will be automatically added to your \texttt{cite.bib} file.

\section{Mathematical Expressions}
Here are some example equations:

\subsection{Quadratic Formula}
The quadratic formula is given by:
\[x = \frac{-b \pm \sqrt{b^2 - 4ac}}{2a}\]

\subsection{Matrix Example}
A sample matrix:
\[
\begin{pmatrix}
a_{11} & a_{12} & a_{13} \\
a_{21} & a_{22} & a_{23} \\
a_{31} & a_{32} & a_{33}
\end{pmatrix}
\]

\section{Tables}
Below is a sample table showing data:

\begin{table}[h]
\centering
\begin{tabular}{@{}lrrr@{}}
\toprule
\textbf{Category} & \textbf{Value 1} & \textbf{Value 2} & \textbf{Total} \\
\midrule
Item A & 10.5 & 20.3 & 30.8 \\
Item B & 15.2 & 25.1 & 40.3 \\
Item C & 22.1 & 12.7 & 34.8 \\
\bottomrule
\end{tabular}
\caption{Sample Data Table}
\label{tab:sample}
\end{table}

\section{Lists}
\subsection{Bullet Points}
\begin{itemize}
  \item First item
  \item Second item
  \item Third item with sub-items:
    \begin{itemize}
      \item Sub-item 1
      \item Sub-item 2
    \end{itemize}
\end{itemize}

\subsection{Numbered List}
\begin{enumerate}
  \item First step
  \item Second step
  \item Third step
\end{enumerate}

\section{Citations and References}
When you use the citation search feature, references will be automatically managed. For example, you might cite a paper like this: \cite{example2024}. All citations will be automatically added to your \texttt{cite.bib} file.

\section{Conclusion}
This template demonstrates basic \LaTeX{} features including mathematical formulas, tables, and lists. The bibliography system is now set up to work seamlessly with Bibby AI's citation management. Feel free to modify it for your needs!

\bibliographystyle{plain}
\bibliography{cite}

\end{document}\documentclass{article}
\usepackage[utf8]{inputenc}
\usepackage{booktabs}
\usepackage{amsmath}
\usepackage{natbib}

\title{My Document}
\author{Your Name}
\date{\today}

\begin{document}

\maketitle

\section{Introduction}
This is a sample document demonstrating various \LaTeX{} features including mathematical equations, tables, and formatting. You can cite papers using the citation search feature, and they will be automatically added to your \texttt{cite.bib} file.

\section{Mathematical Expressions}
Here are some example equations:

\subsection{Quadratic Formula}
The quadratic formula is given by:
\[x = \frac{-b \pm \sqrt{b^2 - 4ac}}{2a}\]

\subsection{Matrix Example}
A sample matrix:
\[
\begin{pmatrix}
a_{11} & a_{12} & a_{13} \\
a_{21} & a_{22} & a_{23} \\
a_{31} & a_{32} & a_{33}
\end{pmatrix}
\]

\section{Tables}
Below is a sample table showing data:

\begin{table}[h]
\centering
\begin{tabular}{@{}lrrr@{}}
\toprule
\textbf{Category} & \textbf{Value 1} & \textbf{Value 2} & \textbf{Total} \\
\midrule
Item A & 10.5 & 20.3 & 30.8 \\
Item B & 15.2 & 25.1 & 40.3 \\
Item C & 22.1 & 12.7 & 34.8 \\
\bottomrule
\end{tabular}
\caption{Sample Data Table}
\label{tab:sample}
\end{table}

\section{Lists}
\subsection{Bullet Points}
\begin{itemize}
  \item First item
  \item Second item
  \item Third item with sub-items:
    \begin{itemize}
      \item Sub-item 1
      \item Sub-item 2
    \end{itemize}
\end{itemize}

\subsection{Numbered List}
\begin{enumerate}
  \item First step
  \item Second step
  \item Third step
\end{enumerate}

\section{Citations and References}
When you use the citation search feature, references will be automatically managed. For example, you might cite a paper like this: \cite{example2024}. All citations will be automatically added to your \texttt{cite.bib} file.

\section{Conclusion}
This template demonstrates basic \LaTeX{} features including mathematical formulas, tables, and lists. The bibliography system is now set up to work seamlessly with Bibby AI's citation management. Feel free to modify it for your needs!

\bibliographystyle{plain}
\bibliography{cite}

\end{document}\documentclass{article}
\usepackage[utf8]{inputenc}
\usepackage{booktabs}
\usepackage{amsmath}
\usepackage{natbib}

\title{My Document}
\author{Your Name}
\date{\today}

\begin{document}

\maketitle

\section{Introduction}
This is a sample document demonstrating various \LaTeX{} features including mathematical equations, tables, and formatting. You can cite papers using the citation search feature, and they will be automatically added to your \texttt{cite.bib} file.

\section{Mathematical Expressions}
Here are some example equations:

\subsection{Quadratic Formula}
The quadratic formula is given by:
\[x = \frac{-b \pm \sqrt{b^2 - 4ac}}{2a}\]

\subsection{Matrix Example}
A sample matrix:
\[
\begin{pmatrix}
a_{11} & a_{12} & a_{13} \\
a_{21} & a_{22} & a_{23} \\
a_{31} & a_{32} & a_{33}
\end{pmatrix}
\]

\section{Tables}
Below is a sample table showing data:

\begin{table}[h]
\centering
\begin{tabular}{@{}lrrr@{}}
\toprule
\textbf{Category} & \textbf{Value 1} & \textbf{Value 2} & \textbf{Total} \\
\midrule
Item A & 10.5 & 20.3 & 30.8 \\
Item B & 15.2 & 25.1 & 40.3 \\
Item C & 22.1 & 12.7 & 34.8 \\
\bottomrule
\end{tabular}
\caption{Sample Data Table}
\label{tab:sample}
\end{table}

\section{Lists}
\subsection{Bullet Points}
\begin{itemize}
  \item First item
  \item Second item
  \item Third item with sub-items:
    \begin{itemize}
      \item Sub-item 1
      \item Sub-item 2
    \end{itemize}
\end{itemize}

\subsection{Numbered List}
\begin{enumerate}
  \item First step
  \item Second step
  \item Third step
\end{enumerate}

\section{Citations and References}
When you use the citation search feature, references will be automatically managed. For example, you might cite a paper like this: \cite{example2024}. All citations will be automatically added to your \texttt{cite.bib} file.

\section{Conclusion}
This template demonstrates basic \LaTeX{} features including mathematical formulas, tables, and lists. The bibliography system is now set up to work seamlessly with Bibby AI's citation management. Feel free to modify it for your needs!

\bibliographystyle{plain}
\bibliography{cite}

\end{document}\documentclass{article}
\usepackage[utf8]{inputenc}
\usepackage{booktabs}
\usepackage{amsmath}
\usepackage{natbib}

\title{My Document}
\author{Your Name}
\date{\today}

\begin{document}

\maketitle

\section{Introduction}
This is a sample document demonstrating various \LaTeX{} features including mathematical equations, tables, and formatting. You can cite papers using the citation search feature, and they will be automatically added to your \texttt{cite.bib} file.

\section{Mathematical Expressions}
Here are some example equations:

\subsection{Quadratic Formula}
The quadratic formula is given by:
\[x = \frac{-b \pm \sqrt{b^2 - 4ac}}{2a}\]

\subsection{Matrix Example}
A sample matrix:
\[
\begin{pmatrix}
a_{11} & a_{12} & a_{13} \\
a_{21} & a_{22} & a_{23} \\
a_{31} & a_{32} & a_{33}
\end{pmatrix}
\]

\section{Tables}
Below is a sample table showing data:

\begin{table}[h]
\centering
\begin{tabular}{@{}lrrr@{}}
\toprule
\textbf{Category} & \textbf{Value 1} & \textbf{Value 2} & \textbf{Total} \\
\midrule
Item A & 10.5 & 20.3 & 30.8 \\
Item B & 15.2 & 25.1 & 40.3 \\
Item C & 22.1 & 12.7 & 34.8 \\
\bottomrule
\end{tabular}
\caption{Sample Data Table}
\label{tab:sample}
\end{table}

\section{Lists}
\subsection{Bullet Points}
\begin{itemize}
  \item First item
  \item Second item
  \item Third item with sub-items:
    \begin{itemize}
      \item Sub-item 1
      \item Sub-item 2
    \end{itemize}
\end{itemize}

\subsection{Numbered List}
\begin{enumerate}
  \item First step
  \item Second step
  \item Third step
\end{enumerate}

\section{Citations and References}
When you use the citation search feature, references will be automatically managed. For example, you might cite a paper like this: \cite{example2024}. All citations will be automatically added to your \texttt{cite.bib} file.

\section{Conclusion}
This template demonstrates basic \LaTeX{} features including mathematical formulas, tables, and lists. The bibliography system is now set up to work seamlessly with Bibby AI's citation management. Feel free to modify it for your needs!

\bibliographystyle{plain}
\bibliography{cite}

\end{document}\documentclass{article}
\usepackage[utf8]{inputenc}
\usepackage{booktabs}
\usepackage{amsmath}
\usepackage{natbib}

\title{My Document}
\author{Your Name}
\date{\today}

\begin{document}

\maketitle

\section{Introduction}
This is a sample document demonstrating various \LaTeX{} features including mathematical equations, tables, and formatting. You can cite papers using the citation search feature, and they will be automatically added to your \texttt{cite.bib} file.

\section{Mathematical Expressions}
Here are some example equations:

\subsection{Quadratic Formula}
The quadratic formula is given by:
\[x = \frac{-b \pm \sqrt{b^2 - 4ac}}{2a}\]

\subsection{Matrix Example}
A sample matrix:
\[
\begin{pmatrix}
a_{11} & a_{12} & a_{13} \\
a_{21} & a_{22} & a_{23} \\
a_{31} & a_{32} & a_{33}
\end{pmatrix}
\]

\section{Tables}
Below is a sample table showing data:

\begin{table}[h]
\centering
\begin{tabular}{@{}lrrr@{}}
\toprule
\textbf{Category} & \textbf{Value 1} & \textbf{Value 2} & \textbf{Total} \\
\midrule
Item A & 10.5 & 20.3 & 30.8 \\
Item B & 15.2 & 25.1 & 40.3 \\
Item C & 22.1 & 12.7 & 34.8 \\
\bottomrule
\end{tabular}
\caption{Sample Data Table}
\label{tab:sample}
\end{table}

\section{Lists}
\subsection{Bullet Points}
\begin{itemize}
  \item First item
  \item Second item
  \item Third item with sub-items:
    \begin{itemize}
      \item Sub-item 1
      \item Sub-item 2
    \end{itemize}
\end{itemize}

\subsection{Numbered List}
\begin{enumerate}
  \item First step
  \item Second step
  \item Third step
\end{enumerate}

\section{Citations and References}
When you use the citation search feature, references will be automatically managed. For example, you might cite a paper like this: \cite{example2024}. All citations will be automatically added to your \texttt{cite.bib} file.

\section{Conclusion}
This template demonstrates basic \LaTeX{} features including mathematical formulas, tables, and lists. The bibliography system is now set up to work seamlessly with Bibby AI's citation management. Feel free to modify it for your needs!

\bibliographystyle{plain}
\bibliography{cite}

\end{document}\documentclass{article}
\usepackage[utf8]{inputenc}
\usepackage{booktabs}
\usepackage{amsmath}
\usepackage{natbib}

\title{My Document}
\author{Your Name}
\date{\today}

\begin{document}

\maketitle

\section{Introduction}
This is a sample document demonstrating various \LaTeX{} features including mathematical equations, tables, and formatting. You can cite papers using the citation search feature, and they will be automatically added to your \texttt{cite.bib} file.

\section{Mathematical Expressions}
Here are some example equations:

\subsection{Quadratic Formula}
The quadratic formula is given by:
\[x = \frac{-b \pm \sqrt{b^2 - 4ac}}{2a}\]

\subsection{Matrix Example}
A sample matrix:
\[
\begin{pmatrix}
a_{11} & a_{12} & a_{13} \\
a_{21} & a_{22} & a_{23} \\
a_{31} & a_{32} & a_{33}
\end{pmatrix}
\]

\section{Tables}
Below is a sample table showing data:

\begin{table}[h]
\centering
\begin{tabular}{@{}lrrr@{}}
\toprule
\textbf{Category} & \textbf{Value 1} & \textbf{Value 2} & \textbf{Total} \\
\midrule
Item A & 10.5 & 20.3 & 30.8 \\
Item B & 15.2 & 25.1 & 40.3 \\
Item C & 22.1 & 12.7 & 34.8 \\
\bottomrule
\end{tabular}
\caption{Sample Data Table}
\label{tab:sample}
\end{table}

\section{Lists}
\subsection{Bullet Points}
\begin{itemize}
  \item First item
  \item Second item
  \item Third item with sub-items:
    \begin{itemize}
      \item Sub-item 1
      \item Sub-item 2
    \end{itemize}
\end{itemize}

\subsection{Numbered List}
\begin{enumerate}
  \item First step
  \item Second step
  \item Third step
\end{enumerate}

\section{Citations and References}
When you use the citation search feature, references will be automatically managed. For example, you might cite a paper like this: \cite{example2024}. All citations will be automatically added to your \texttt{cite.bib} file.

\section{Conclusion}
This template demonstrates basic \LaTeX{} features including mathematical formulas, tables, and lists. The bibliography system is now set up to work seamlessly with Bibby AI's citation management. Feel free to modify it for your needs!

\bibliographystyle{plain}
\bibliography{cite}

\end{document}\documentclass{article}
\usepackage[utf8]{inputenc}
\usepackage{booktabs}
\usepackage{amsmath}
\usepackage{natbib}

\title{My Document}
\author{Your Name}
\date{\today}

\begin{document}

\maketitle

\section{Introduction}
This is a sample document demonstrating various \LaTeX{} features including mathematical equations, tables, and formatting. You can cite papers using the citation search feature, and they will be automatically added to your \texttt{cite.bib} file.

\section{Mathematical Expressions}
Here are some example equations:

\subsection{Quadratic Formula}
The quadratic formula is given by:
\[x = \frac{-b \pm \sqrt{b^2 - 4ac}}{2a}\]

\subsection{Matrix Example}
A sample matrix:
\[
\begin{pmatrix}
a_{11} & a_{12} & a_{13} \\
a_{21} & a_{22} & a_{23} \\
a_{31} & a_{32} & a_{33}
\end{pmatrix}
\]

\section{Tables}
Below is a sample table showing data:

\begin{table}[h]
\centering
\begin{tabular}{@{}lrrr@{}}
\toprule
\textbf{Category} & \textbf{Value 1} & \textbf{Value 2} & \textbf{Total} \\
\midrule
Item A & 10.5 & 20.3 & 30.8 \\
Item B & 15.2 & 25.1 & 40.3 \\
Item C & 22.1 & 12.7 & 34.8 \\
\bottomrule
\end{tabular}
\caption{Sample Data Table}
\label{tab:sample}
\end{table}

\section{Lists}
\subsection{Bullet Points}
\begin{itemize}
  \item First item
  \item Second item
  \item Third item with sub-items:
    \begin{itemize}
      \item Sub-item 1
      \item Sub-item 2
    \end{itemize}
\end{itemize}

\subsection{Numbered List}
\begin{enumerate}
  \item First step
  \item Second step
  \item Third step
\end{enumerate}

\section{Citations and References}
When you use the citation search feature, references will be automatically managed. For example, you might cite a paper like this: \cite{example2024}. All citations will be automatically added to your \texttt{cite.bib} file.

\section{Conclusion}
This template demonstrates basic \LaTeX{} features including mathematical formulas, tables, and lists. The bibliography system is now set up to work seamlessly with Bibby AI's citation management. Feel free to modify it for your needs!

\bibliographystyle{plain}
\bibliography{cite}

\end{document}\documentclass{article}
\usepackage[utf8]{inputenc}
\usepackage{booktabs}
\usepackage{amsmath}
\usepackage{natbib}

\title{My Document}
\author{Your Name}
\date{\today}

\begin{document}

\maketitle

\section{Introduction}
This is a sample document demonstrating various \LaTeX{} features including mathematical equations, tables, and formatting. You can cite papers using the citation search feature, and they will be automatically added to your \texttt{cite.bib} file.

\section{Mathematical Expressions}
Here are some example equations:

\subsection{Quadratic Formula}
The quadratic formula is given by:
\[x = \frac{-b \pm \sqrt{b^2 - 4ac}}{2a}\]

\subsection{Matrix Example}
A sample matrix:
\[
\begin{pmatrix}
a_{11} & a_{12} & a_{13} \\
a_{21} & a_{22} & a_{23} \\
a_{31} & a_{32} & a_{33}
\end{pmatrix}
\]

\section{Tables}
Below is a sample table showing data:

\begin{table}[h]
\centering
\begin{tabular}{@{}lrrr@{}}
\toprule
\textbf{Category} & \textbf{Value 1} & \textbf{Value 2} & \textbf{Total} \\
\midrule
Item A & 10.5 & 20.3 & 30.8 \\
Item B & 15.2 & 25.1 & 40.3 \\
Item C & 22.1 & 12.7 & 34.8 \\
\bottomrule
\end{tabular}
\caption{Sample Data Table}
\label{tab:sample}
\end{table}

\section{Lists}
\subsection{Bullet Points}
\begin{itemize}
  \item First item
  \item Second item
  \item Third item with sub-items:
    \begin{itemize}
      \item Sub-item 1
      \item Sub-item 2
    \end{itemize}
\end{itemize}

\subsection{Numbered List}
\begin{enumerate}
  \item First step
  \item Second step
  \item Third step
\end{enumerate}

\section{Citations and References}
When you use the citation search feature, references will be automatically managed. For example, you might cite a paper like this: \cite{example2024}. All citations will be automatically added to your \texttt{cite.bib} file.

\section{Conclusion}
This template demonstrates basic \LaTeX{} features including mathematical formulas, tables, and lists. The bibliography system is now set up to work seamlessly with Bibby AI's citation management. Feel free to modify it for your needs!

\bibliographystyle{plain}
\bibliography{cite}

\end{document}\documentclass{article}
\usepackage[utf8]{inputenc}
\usepackage{booktabs}
\usepackage{amsmath}
\usepackage{natbib}

\title{My Document}
\author{Your Name}
\date{\today}

\begin{document}

\maketitle

\section{Introduction}
This is a sample document demonstrating various \LaTeX{} features including mathematical equations, tables, and formatting. You can cite papers using the citation search feature, and they will be automatically added to your \texttt{cite.bib} file.

\section{Mathematical Expressions}
Here are some example equations:

\subsection{Quadratic Formula}
The quadratic formula is given by:
\[x = \frac{-b \pm \sqrt{b^2 - 4ac}}{2a}\]

\subsection{Matrix Example}
A sample matrix:
\[
\begin{pmatrix}
a_{11} & a_{12} & a_{13} \\
a_{21} & a_{22} & a_{23} \\
a_{31} & a_{32} & a_{33}
\end{pmatrix}
\]

\section{Tables}
Below is a sample table showing data:

\begin{table}[h]
\centering
\begin{tabular}{@{}lrrr@{}}
\toprule
\textbf{Category} & \textbf{Value 1} & \textbf{Value 2} & \textbf{Total} \\
\midrule
Item A & 10.5 & 20.3 & 30.8 \\
Item B & 15.2 & 25.1 & 40.3 \\
Item C & 22.1 & 12.7 & 34.8 \\
\bottomrule
\end{tabular}
\caption{Sample Data Table}
\label{tab:sample}
\end{table}

\section{Lists}
\subsection{Bullet Points}
\begin{itemize}
  \item First item
  \item Second item
  \item Third item with sub-items:
    \begin{itemize}
      \item Sub-item 1
      \item Sub-item 2
    \end{itemize}
\end{itemize}

\subsection{Numbered List}
\begin{enumerate}
  \item First step
  \item Second step
  \item Third step
\end{enumerate}

\section{Citations and References}
When you use the citation search feature, references will be automatically managed. For example, you might cite a paper like this: \cite{example2024}. All citations will be automatically added to your \texttt{cite.bib} file.

\section{Conclusion}
This template demonstrates basic \LaTeX{} features including mathematical formulas, tables, and lists. The bibliography system is now set up to work seamlessly with Bibby AI's citation management. Feel free to modify it for your needs!

\bibliographystyle{plain}
\bibliography{cite}

\end{document}\documentclass{article}
\usepackage[utf8]{inputenc}
\usepackage{booktabs}
\usepackage{amsmath}
\usepackage{natbib}

\title{My Document}
\author{Your Name}
\date{\today}

\begin{document}

\maketitle

\section{Introduction}
This is a sample document demonstrating various \LaTeX{} features including mathematical equations, tables, and formatting. You can cite papers using the citation search feature, and they will be automatically added to your \texttt{cite.bib} file.

\section{Mathematical Expressions}
Here are some example equations:

\subsection{Quadratic Formula}
The quadratic formula is given by:
\[x = \frac{-b \pm \sqrt{b^2 - 4ac}}{2a}\]

\subsection{Matrix Example}
A sample matrix:
\[
\begin{pmatrix}
a_{11} & a_{12} & a_{13} \\
a_{21} & a_{22} & a_{23} \\
a_{31} & a_{32} & a_{33}
\end{pmatrix}
\]

\section{Tables}
Below is a sample table showing data:

\begin{table}[h]
\centering
\begin{tabular}{@{}lrrr@{}}
\toprule
\textbf{Category} & \textbf{Value 1} & \textbf{Value 2} & \textbf{Total} \\
\midrule
Item A & 10.5 & 20.3 & 30.8 \\
Item B & 15.2 & 25.1 & 40.3 \\
Item C & 22.1 & 12.7 & 34.8 \\
\bottomrule
\end{tabular}
\caption{Sample Data Table}
\label{tab:sample}
\end{table}

\section{Lists}
\subsection{Bullet Points}
\begin{itemize}
  \item First item
  \item Second item
  \item Third item with sub-items:
    \begin{itemize}
      \item Sub-item 1
      \item Sub-item 2
    \end{itemize}
\end{itemize}

\subsection{Numbered List}
\begin{enumerate}
  \item First step
  \item Second step
  \item Third step
\end{enumerate}

\section{Citations and References}
When you use the citation search feature, references will be automatically managed. For example, you might cite a paper like this: \cite{example2024}. All citations will be automatically added to your \texttt{cite.bib} file.

\section{Conclusion}
This template demonstrates basic \LaTeX{} features including mathematical formulas, tables, and lists. The bibliography system is now set up to work seamlessly with Bibby AI's citation management. Feel free to modify it for your needs!

\bibliographystyle{plain}
\bibliography{cite}

\end{document}\documentclass{article}
\usepackage[utf8]{inputenc}
\usepackage{booktabs}
\usepackage{amsmath}
\usepackage{natbib}

\title{My Document}
\author{Your Name}
\date{\today}

\begin{document}

\maketitle

\section{Introduction}
This is a sample document demonstrating various \LaTeX{} features including mathematical equations, tables, and formatting. You can cite papers using the citation search feature, and they will be automatically added to your \texttt{cite.bib} file.

\section{Mathematical Expressions}
Here are some example equations:

\subsection{Quadratic Formula}
The quadratic formula is given by:
\[x = \frac{-b \pm \sqrt{b^2 - 4ac}}{2a}\]

\subsection{Matrix Example}
A sample matrix:
\[
\begin{pmatrix}
a_{11} & a_{12} & a_{13} \\
a_{21} & a_{22} & a_{23} \\
a_{31} & a_{32} & a_{33}
\end{pmatrix}
\]

\section{Tables}
Below is a sample table showing data:

\begin{table}[h]
\centering
\begin{tabular}{@{}lrrr@{}}
\toprule
\textbf{Category} & \textbf{Value 1} & \textbf{Value 2} & \textbf{Total} \\
\midrule
Item A & 10.5 & 20.3 & 30.8 \\
Item B & 15.2 & 25.1 & 40.3 \\
Item C & 22.1 & 12.7 & 34.8 \\
\bottomrule
\end{tabular}
\caption{Sample Data Table}
\label{tab:sample}
\end{table}

\section{Lists}
\subsection{Bullet Points}
\begin{itemize}
  \item First item
  \item Second item
  \item Third item with sub-items:
    \begin{itemize}
      \item Sub-item 1
      \item Sub-item 2
    \end{itemize}
\end{itemize}

\subsection{Numbered List}
\begin{enumerate}
  \item First step
  \item Second step
  \item Third step
\end{enumerate}

\section{Citations and References}
When you use the citation search feature, references will be automatically managed. For example, you might cite a paper like this: \cite{example2024}. All citations will be automatically added to your \texttt{cite.bib} file.

\section{Conclusion}
This template demonstrates basic \LaTeX{} features including mathematical formulas, tables, and lists. The bibliography system is now set up to work seamlessly with Bibby AI's citation management. Feel free to modify it for your needs!

\bibliographystyle{plain}
\bibliography{cite}

\end{document}\documentclass{article}
\usepackage[utf8]{inputenc}
\usepackage{booktabs}
\usepackage{amsmath}
\usepackage{natbib}

\title{My Document}
\author{Your Name}
\date{\today}

\begin{document}

\maketitle

\section{Introduction}
This is a sample document demonstrating various \LaTeX{} features including mathematical equations, tables, and formatting. You can cite papers using the citation search feature, and they will be automatically added to your \texttt{cite.bib} file.

\section{Mathematical Expressions}
Here are some example equations:

\subsection{Quadratic Formula}
The quadratic formula is given by:
\[x = \frac{-b \pm \sqrt{b^2 - 4ac}}{2a}\]

\subsection{Matrix Example}
A sample matrix:
\[
\begin{pmatrix}
a_{11} & a_{12} & a_{13} \\
a_{21} & a_{22} & a_{23} \\
a_{31} & a_{32} & a_{33}
\end{pmatrix}
\]

\section{Tables}
Below is a sample table showing data:

\begin{table}[h]
\centering
\begin{tabular}{@{}lrrr@{}}
\toprule
\textbf{Category} & \textbf{Value 1} & \textbf{Value 2} & \textbf{Total} \\
\midrule
Item A & 10.5 & 20.3 & 30.8 \\
Item B & 15.2 & 25.1 & 40.3 \\
Item C & 22.1 & 12.7 & 34.8 \\
\bottomrule
\end{tabular}
\caption{Sample Data Table}
\label{tab:sample}
\end{table}

\section{Lists}
\subsection{Bullet Points}
\begin{itemize}
  \item First item
  \item Second item
  \item Third item with sub-items:
    \begin{itemize}
      \item Sub-item 1
      \item Sub-item 2
    \end{itemize}
\end{itemize}

\subsection{Numbered List}
\begin{enumerate}
  \item First step
  \item Second step
  \item Third step
\end{enumerate}

\section{Citations and References}
When you use the citation search feature, references will be automatically managed. For example, you might cite a paper like this: \cite{example2024}. All citations will be automatically added to your \texttt{cite.bib} file.

\section{Conclusion}
This template demonstrates basic \LaTeX{} features including mathematical formulas, tables, and lists. The bibliography system is now set up to work seamlessly with Bibby AI's citation management. Feel free to modify it for your needs!

\bibliographystyle{plain}
\bibliography{cite}

\end{document}\documentclass{article}
\usepackage[utf8]{inputenc}
\usepackage{booktabs}
\usepackage{amsmath}
\usepackage{natbib}

\title{My Document}
\author{Your Name}
\date{\today}

\begin{document}

\maketitle

\section{Introduction}
This is a sample document demonstrating various \LaTeX{} features including mathematical equations, tables, and formatting. You can cite papers using the citation search feature, and they will be automatically added to your \texttt{cite.bib} file.

\section{Mathematical Expressions}
Here are some example equations:

\subsection{Quadratic Formula}
The quadratic formula is given by:
\[x = \frac{-b \pm \sqrt{b^2 - 4ac}}{2a}\]

\subsection{Matrix Example}
A sample matrix:
\[
\begin{pmatrix}
a_{11} & a_{12} & a_{13} \\
a_{21} & a_{22} & a_{23} \\
a_{31} & a_{32} & a_{33}
\end{pmatrix}
\]

\section{Tables}
Below is a sample table showing data:

\begin{table}[h]
\centering
\begin{tabular}{@{}lrrr@{}}
\toprule
\textbf{Category} & \textbf{Value 1} & \textbf{Value 2} & \textbf{Total} \\
\midrule
Item A & 10.5 & 20.3 & 30.8 \\
Item B & 15.2 & 25.1 & 40.3 \\
Item C & 22.1 & 12.7 & 34.8 \\
\bottomrule
\end{tabular}
\caption{Sample Data Table}
\label{tab:sample}
\end{table}

\section{Lists}
\subsection{Bullet Points}
\begin{itemize}
  \item First item
  \item Second item
  \item Third item with sub-items:
    \begin{itemize}
      \item Sub-item 1
      \item Sub-item 2
    \end{itemize}
\end{itemize}

\subsection{Numbered List}
\begin{enumerate}
  \item First step
  \item Second step
  \item Third step
\end{enumerate}

\section{Citations and References}
When you use the citation search feature, references will be automatically managed. For example, you might cite a paper like this: \cite{example2024}. All citations will be automatically added to your \texttt{cite.bib} file.

\section{Conclusion}
This template demonstrates basic \LaTeX{} features including mathematical formulas, tables, and lists. The bibliography system is now set up to work seamlessly with Bibby AI's citation management. Feel free to modify it for your needs!

\bibliographystyle{plain}
\bibliography{cite}

\end{document}\documentclass{article}
\usepackage[utf8]{inputenc}
\usepackage{booktabs}
\usepackage{amsmath}
\usepackage{natbib}

\title{My Document}
\author{Your Name}
\date{\today}

\begin{document}

\maketitle

\section{Introduction}
This is a sample document demonstrating various \LaTeX{} features including mathematical equations, tables, and formatting. You can cite papers using the citation search feature, and they will be automatically added to your \texttt{cite.bib} file.

\section{Mathematical Expressions}
Here are some example equations:

\subsection{Quadratic Formula}
The quadratic formula is given by:
\[x = \frac{-b \pm \sqrt{b^2 - 4ac}}{2a}\]

\subsection{Matrix Example}
A sample matrix:
\[
\begin{pmatrix}
a_{11} & a_{12} & a_{13} \\
a_{21} & a_{22} & a_{23} \\
a_{31} & a_{32} & a_{33}
\end{pmatrix}
\]

\section{Tables}
Below is a sample table showing data:

\begin{table}[h]
\centering
\begin{tabular}{@{}lrrr@{}}
\toprule
\textbf{Category} & \textbf{Value 1} & \textbf{Value 2} & \textbf{Total} \\
\midrule
Item A & 10.5 & 20.3 & 30.8 \\
Item B & 15.2 & 25.1 & 40.3 \\
Item C & 22.1 & 12.7 & 34.8 \\
\bottomrule
\end{tabular}
\caption{Sample Data Table}
\label{tab:sample}
\end{table}

\section{Lists}
\subsection{Bullet Points}
\begin{itemize}
  \item First item
  \item Second item
  \item Third item with sub-items:
    \begin{itemize}
      \item Sub-item 1
      \item Sub-item 2
    \end{itemize}
\end{itemize}

\subsection{Numbered List}
\begin{enumerate}
  \item First step
  \item Second step
  \item Third step
\end{enumerate}

\section{Citations and References}
When you use the citation search feature, references will be automatically managed. For example, you might cite a paper like this: \cite{example2024}. All citations will be automatically added to your \texttt{cite.bib} file.

\section{Conclusion}
This template demonstrates basic \LaTeX{} features including mathematical formulas, tables, and lists. The bibliography system is now set up to work seamlessly with Bibby AI's citation management. Feel free to modify it for your needs!

\bibliographystyle{plain}
\bibliography{cite}

\end{document}\documentclass{article}
\usepackage[utf8]{inputenc}
\usepackage{booktabs}
\usepackage{amsmath}
\usepackage{natbib}

\title{My Document}
\author{Your Name}
\date{\today}

\begin{document}

\maketitle

\section{Introduction}
This is a sample document demonstrating various \LaTeX{} features including mathematical equations, tables, and formatting. You can cite papers using the citation search feature, and they will be automatically added to your \texttt{cite.bib} file.

\section{Mathematical Expressions}
Here are some example equations:

\subsection{Quadratic Formula}
The quadratic formula is given by:
\[x = \frac{-b \pm \sqrt{b^2 - 4ac}}{2a}\]

\subsection{Matrix Example}
A sample matrix:
\[
\begin{pmatrix}
a_{11} & a_{12} & a_{13} \\
a_{21} & a_{22} & a_{23} \\
a_{31} & a_{32} & a_{33}
\end{pmatrix}
\]

\section{Tables}
Below is a sample table showing data:

\begin{table}[h]
\centering
\begin{tabular}{@{}lrrr@{}}
\toprule
\textbf{Category} & \textbf{Value 1} & \textbf{Value 2} & \textbf{Total} \\
\midrule
Item A & 10.5 & 20.3 & 30.8 \\
Item B & 15.2 & 25.1 & 40.3 \\
Item C & 22.1 & 12.7 & 34.8 \\
\bottomrule
\end{tabular}
\caption{Sample Data Table}
\label{tab:sample}
\end{table}

\section{Lists}
\subsection{Bullet Points}
\begin{itemize}
  \item First item
  \item Second item
  \item Third item with sub-items:
    \begin{itemize}
      \item Sub-item 1
      \item Sub-item 2
    \end{itemize}
\end{itemize}

\subsection{Numbered List}
\begin{enumerate}
  \item First step
  \item Second step
  \item Third step
\end{enumerate}

\section{Citations and References}
When you use the citation search feature, references will be automatically managed. For example, you might cite a paper like this: \cite{example2024}. All citations will be automatically added to your \texttt{cite.bib} file.

\section{Conclusion}
This template demonstrates basic \LaTeX{} features including mathematical formulas, tables, and lists. The bibliography system is now set up to work seamlessly with Bibby AI's citation management. Feel free to modify it for your needs!

\bibliographystyle{plain}
\bibliography{cite}

\end{document}\documentclass{article}
\usepackage[utf8]{inputenc}
\usepackage{booktabs}
\usepackage{amsmath}
\usepackage{natbib}

\title{My Document}
\author{Your Name}
\date{\today}

\begin{document}

\maketitle

\section{Introduction}
This is a sample document demonstrating various \LaTeX{} features including mathematical equations, tables, and formatting. You can cite papers using the citation search feature, and they will be automatically added to your \texttt{cite.bib} file.

\section{Mathematical Expressions}
Here are some example equations:

\subsection{Quadratic Formula}
The quadratic formula is given by:
\[x = \frac{-b \pm \sqrt{b^2 - 4ac}}{2a}\]

\subsection{Matrix Example}
A sample matrix:
\[
\begin{pmatrix}
a_{11} & a_{12} & a_{13} \\
a_{21} & a_{22} & a_{23} \\
a_{31} & a_{32} & a_{33}
\end{pmatrix}
\]

\section{Tables}
Below is a sample table showing data:

\begin{table}[h]
\centering
\begin{tabular}{@{}lrrr@{}}
\toprule
\textbf{Category} & \textbf{Value 1} & \textbf{Value 2} & \textbf{Total} \\
\midrule
Item A & 10.5 & 20.3 & 30.8 \\
Item B & 15.2 & 25.1 & 40.3 \\
Item C & 22.1 & 12.7 & 34.8 \\
\bottomrule
\end{tabular}
\caption{Sample Data Table}
\label{tab:sample}
\end{table}

\section{Lists}
\subsection{Bullet Points}
\begin{itemize}
  \item First item
  \item Second item
  \item Third item with sub-items:
    \begin{itemize}
      \item Sub-item 1
      \item Sub-item 2
    \end{itemize}
\end{itemize}

\subsection{Numbered List}
\begin{enumerate}
  \item First step
  \item Second step
  \item Third step
\end{enumerate}

\section{Citations and References}
When you use the citation search feature, references will be automatically managed. For example, you might cite a paper like this: \cite{example2024}. All citations will be automatically added to your \texttt{cite.bib} file.

\section{Conclusion}
This template demonstrates basic \LaTeX{} features including mathematical formulas, tables, and lists. The bibliography system is now set up to work seamlessly with Bibby AI's citation management. Feel free to modify it for your needs!

\bibliographystyle{plain}
\bibliography{cite}

\end{document}\documentclass{article}
\usepackage[utf8]{inputenc}
\usepackage{booktabs}
\usepackage{amsmath}
\usepackage{natbib}

\title{My Document}
\author{Your Name}
\date{\today}

\begin{document}

\maketitle

\section{Introduction}
This is a sample document demonstrating various \LaTeX{} features including mathematical equations, tables, and formatting. You can cite papers using the citation search feature, and they will be automatically added to your \texttt{cite.bib} file.

\section{Mathematical Expressions}
Here are some example equations:

\subsection{Quadratic Formula}
The quadratic formula is given by:
\[x = \frac{-b \pm \sqrt{b^2 - 4ac}}{2a}\]

\subsection{Matrix Example}
A sample matrix:
\[
\begin{pmatrix}
a_{11} & a_{12} & a_{13} \\
a_{21} & a_{22} & a_{23} \\
a_{31} & a_{32} & a_{33}
\end{pmatrix}
\]

\section{Tables}
Below is a sample table showing data:

\begin{table}[h]
\centering
\begin{tabular}{@{}lrrr@{}}
\toprule
\textbf{Category} & \textbf{Value 1} & \textbf{Value 2} & \textbf{Total} \\
\midrule
Item A & 10.5 & 20.3 & 30.8 \\
Item B & 15.2 & 25.1 & 40.3 \\
Item C & 22.1 & 12.7 & 34.8 \\
\bottomrule
\end{tabular}
\caption{Sample Data Table}
\label{tab:sample}
\end{table}

\section{Lists}
\subsection{Bullet Points}
\begin{itemize}
  \item First item
  \item Second item
  \item Third item with sub-items:
    \begin{itemize}
      \item Sub-item 1
      \item Sub-item 2
    \end{itemize}
\end{itemize}

\subsection{Numbered List}
\begin{enumerate}
  \item First step
  \item Second step
  \item Third step
\end{enumerate}

\section{Citations and References}
When you use the citation search feature, references will be automatically managed. For example, you might cite a paper like this: \cite{example2024}. All citations will be automatically added to your \texttt{cite.bib} file.

\section{Conclusion}
This template demonstrates basic \LaTeX{} features including mathematical formulas, tables, and lists. The bibliography system is now set up to work seamlessly with Bibby AI's citation management. Feel free to modify it for your needs!

\bibliographystyle{plain}
\bibliography{cite}

\end{document}\documentclass{article}
\usepackage[utf8]{inputenc}
\usepackage{booktabs}
\usepackage{amsmath}
\usepackage{natbib}

\title{My Document}
\author{Your Name}
\date{\today}

\begin{document}

\maketitle

\section{Introduction}
This is a sample document demonstrating various \LaTeX{} features including mathematical equations, tables, and formatting. You can cite papers using the citation search feature, and they will be automatically added to your \texttt{cite.bib} file.

\section{Mathematical Expressions}
Here are some example equations:

\subsection{Quadratic Formula}
The quadratic formula is given by:
\[x = \frac{-b \pm \sqrt{b^2 - 4ac}}{2a}\]

\subsection{Matrix Example}
A sample matrix:
\[
\begin{pmatrix}
a_{11} & a_{12} & a_{13} \\
a_{21} & a_{22} & a_{23} \\
a_{31} & a_{32} & a_{33}
\end{pmatrix}
\]

\section{Tables}
Below is a sample table showing data:

\begin{table}[h]
\centering
\begin{tabular}{@{}lrrr@{}}
\toprule
\textbf{Category} & \textbf{Value 1} & \textbf{Value 2} & \textbf{Total} \\
\midrule
Item A & 10.5 & 20.3 & 30.8 \\
Item B & 15.2 & 25.1 & 40.3 \\
Item C & 22.1 & 12.7 & 34.8 \\
\bottomrule
\end{tabular}
\caption{Sample Data Table}
\label{tab:sample}
\end{table}

\section{Lists}
\subsection{Bullet Points}
\begin{itemize}
  \item First item
  \item Second item
  \item Third item with sub-items:
    \begin{itemize}
      \item Sub-item 1
      \item Sub-item 2
    \end{itemize}
\end{itemize}

\subsection{Numbered List}
\begin{enumerate}
  \item First step
  \item Second step
  \item Third step
\end{enumerate}

\section{Citations and References}
When you use the citation search feature, references will be automatically managed. For example, you might cite a paper like this: \cite{example2024}. All citations will be automatically added to your \texttt{cite.bib} file.

\section{Conclusion}
This template demonstrates basic \LaTeX{} features including mathematical formulas, tables, and lists. The bibliography system is now set up to work seamlessly with Bibby AI's citation management. Feel free to modify it for your needs!

\bibliographystyle{plain}
\bibliography{cite}

\end{document}\documentclass{article}
\usepackage[utf8]{inputenc}
\usepackage{booktabs}
\usepackage{amsmath}
\usepackage{natbib}

\title{My Document}
\author{Your Name}
\date{\today}

\begin{document}

\maketitle

\section{Introduction}
This is a sample document demonstrating various \LaTeX{} features including mathematical equations, tables, and formatting. You can cite papers using the citation search feature, and they will be automatically added to your \texttt{cite.bib} file.

\section{Mathematical Expressions}
Here are some example equations:

\subsection{Quadratic Formula}
The quadratic formula is given by:
\[x = \frac{-b \pm \sqrt{b^2 - 4ac}}{2a}\]

\subsection{Matrix Example}
A sample matrix:
\[
\begin{pmatrix}
a_{11} & a_{12} & a_{13} \\
a_{21} & a_{22} & a_{23} \\
a_{31} & a_{32} & a_{33}
\end{pmatrix}
\]

\section{Tables}
Below is a sample table showing data:

\begin{table}[h]
\centering
\begin{tabular}{@{}lrrr@{}}
\toprule
\textbf{Category} & \textbf{Value 1} & \textbf{Value 2} & \textbf{Total} \\
\midrule
Item A & 10.5 & 20.3 & 30.8 \\
Item B & 15.2 & 25.1 & 40.3 \\
Item C & 22.1 & 12.7 & 34.8 \\
\bottomrule
\end{tabular}
\caption{Sample Data Table}
\label{tab:sample}
\end{table}

\section{Lists}
\subsection{Bullet Points}
\begin{itemize}
  \item First item
  \item Second item
  \item Third item with sub-items:
    \begin{itemize}
      \item Sub-item 1
      \item Sub-item 2
    \end{itemize}
\end{itemize}

\subsection{Numbered List}
\begin{enumerate}
  \item First step
  \item Second step
  \item Third step
\end{enumerate}

\section{Citations and References}
When you use the citation search feature, references will be automatically managed. For example, you might cite a paper like this: \cite{example2024}. All citations will be automatically added to your \texttt{cite.bib} file.

\section{Conclusion}
This template demonstrates basic \LaTeX{} features including mathematical formulas, tables, and lists. The bibliography system is now set up to work seamlessly with Bibby AI's citation management. Feel free to modify it for your needs!

\bibliographystyle{plain}
\bibliography{cite}

\end{document}\documentclass{article}
\usepackage[utf8]{inputenc}
\usepackage{booktabs}
\usepackage{amsmath}
\usepackage{natbib}

\title{My Document}
\author{Your Name}
\date{\today}

\begin{document}

\maketitle

\section{Introduction}
This is a sample document demonstrating various \LaTeX{} features including mathematical equations, tables, and formatting. You can cite papers using the citation search feature, and they will be automatically added to your \texttt{cite.bib} file.

\section{Mathematical Expressions}
Here are some example equations:

\subsection{Quadratic Formula}
The quadratic formula is given by:
\[x = \frac{-b \pm \sqrt{b^2 - 4ac}}{2a}\]

\subsection{Matrix Example}
A sample matrix:
\[
\begin{pmatrix}
a_{11} & a_{12} & a_{13} \\
a_{21} & a_{22} & a_{23} \\
a_{31} & a_{32} & a_{33}
\end{pmatrix}
\]

\section{Tables}
Below is a sample table showing data:

\begin{table}[h]
\centering
\begin{tabular}{@{}lrrr@{}}
\toprule
\textbf{Category} & \textbf{Value 1} & \textbf{Value 2} & \textbf{Total} \\
\midrule
Item A & 10.5 & 20.3 & 30.8 \\
Item B & 15.2 & 25.1 & 40.3 \\
Item C & 22.1 & 12.7 & 34.8 \\
\bottomrule
\end{tabular}
\caption{Sample Data Table}
\label{tab:sample}
\end{table}

\section{Lists}
\subsection{Bullet Points}
\begin{itemize}
  \item First item
  \item Second item
  \item Third item with sub-items:
    \begin{itemize}
      \item Sub-item 1
      \item Sub-item 2
    \end{itemize}
\end{itemize}

\subsection{Numbered List}
\begin{enumerate}
  \item First step
  \item Second step
  \item Third step
\end{enumerate}

\section{Citations and References}
When you use the citation search feature, references will be automatically managed. For example, you might cite a paper like this: \cite{example2024}. All citations will be automatically added to your \texttt{cite.bib} file.

\section{Conclusion}
This template demonstrates basic \LaTeX{} features including mathematical formulas, tables, and lists. The bibliography system is now set up to work seamlessly with Bibby AI's citation management. Feel free to modify it for your needs!

\bibliographystyle{plain}
\bibliography{cite}

\end{document}\documentclass{article}
\usepackage[utf8]{inputenc}
\usepackage{booktabs}
\usepackage{amsmath}
\usepackage{natbib}

\title{My Document}
\author{Your Name}
\date{\today}

\begin{document}

\maketitle

\section{Introduction}
This is a sample document demonstrating various \LaTeX{} features including mathematical equations, tables, and formatting. You can cite papers using the citation search feature, and they will be automatically added to your \texttt{cite.bib} file.

\section{Mathematical Expressions}
Here are some example equations:

\subsection{Quadratic Formula}
The quadratic formula is given by:
\[x = \frac{-b \pm \sqrt{b^2 - 4ac}}{2a}\]

\subsection{Matrix Example}
A sample matrix:
\[
\begin{pmatrix}
a_{11} & a_{12} & a_{13} \\
a_{21} & a_{22} & a_{23} \\
a_{31} & a_{32} & a_{33}
\end{pmatrix}
\]

\section{Tables}
Below is a sample table showing data:

\begin{table}[h]
\centering
\begin{tabular}{@{}lrrr@{}}
\toprule
\textbf{Category} & \textbf{Value 1} & \textbf{Value 2} & \textbf{Total} \\
\midrule
Item A & 10.5 & 20.3 & 30.8 \\
Item B & 15.2 & 25.1 & 40.3 \\
Item C & 22.1 & 12.7 & 34.8 \\
\bottomrule
\end{tabular}
\caption{Sample Data Table}
\label{tab:sample}
\end{table}

\section{Lists}
\subsection{Bullet Points}
\begin{itemize}
  \item First item
  \item Second item
  \item Third item with sub-items:
    \begin{itemize}
      \item Sub-item 1
      \item Sub-item 2
    \end{itemize}
\end{itemize}

\subsection{Numbered List}
\begin{enumerate}
  \item First step
  \item Second step
  \item Third step
\end{enumerate}

\section{Citations and References}
When you use the citation search feature, references will be automatically managed. For example, you might cite a paper like this: \cite{example2024}. All citations will be automatically added to your \texttt{cite.bib} file.

\section{Conclusion}
This template demonstrates basic \LaTeX{} features including mathematical formulas, tables, and lists. The bibliography system is now set up to work seamlessly with Bibby AI's citation management. Feel free to modify it for your needs!

\bibliographystyle{plain}
\bibliography{cite}

\end{document}\documentclass{article}
\usepackage[utf8]{inputenc}
\usepackage{booktabs}
\usepackage{amsmath}
\usepackage{natbib}

\title{My Document}
\author{Your Name}
\date{\today}

\begin{document}

\maketitle

\section{Introduction}
This is a sample document demonstrating various \LaTeX{} features including mathematical equations, tables, and formatting. You can cite papers using the citation search feature, and they will be automatically added to your \texttt{cite.bib} file.

\section{Mathematical Expressions}
Here are some example equations:

\subsection{Quadratic Formula}
The quadratic formula is given by:
\[x = \frac{-b \pm \sqrt{b^2 - 4ac}}{2a}\]

\subsection{Matrix Example}
A sample matrix:
\[
\begin{pmatrix}
a_{11} & a_{12} & a_{13} \\
a_{21} & a_{22} & a_{23} \\
a_{31} & a_{32} & a_{33}
\end{pmatrix}
\]

\section{Tables}
Below is a sample table showing data:

\begin{table}[h]
\centering
\begin{tabular}{@{}lrrr@{}}
\toprule
\textbf{Category} & \textbf{Value 1} & \textbf{Value 2} & \textbf{Total} \\
\midrule
Item A & 10.5 & 20.3 & 30.8 \\
Item B & 15.2 & 25.1 & 40.3 \\
Item C & 22.1 & 12.7 & 34.8 \\
\bottomrule
\end{tabular}
\caption{Sample Data Table}
\label{tab:sample}
\end{table}

\section{Lists}
\subsection{Bullet Points}
\begin{itemize}
  \item First item
  \item Second item
  \item Third item with sub-items:
    \begin{itemize}
      \item Sub-item 1
      \item Sub-item 2
    \end{itemize}
\end{itemize}

\subsection{Numbered List}
\begin{enumerate}
  \item First step
  \item Second step
  \item Third step
\end{enumerate}

\section{Citations and References}
When you use the citation search feature, references will be automatically managed. For example, you might cite a paper like this: \cite{example2024}. All citations will be automatically added to your \texttt{cite.bib} file.

\section{Conclusion}
This template demonstrates basic \LaTeX{} features including mathematical formulas, tables, and lists. The bibliography system is now set up to work seamlessly with Bibby AI's citation management. Feel free to modify it for your needs!

\bibliographystyle{plain}
\bibliography{cite}

\end{document}\documentclass{article}
\usepackage[utf8]{inputenc}
\usepackage{booktabs}
\usepackage{amsmath}
\usepackage{natbib}

\title{My Document}
\author{Your Name}
\date{\today}

\begin{document}

\maketitle

\section{Introduction}
This is a sample document demonstrating various \LaTeX{} features including mathematical equations, tables, and formatting. You can cite papers using the citation search feature, and they will be automatically added to your \texttt{cite.bib} file.

\section{Mathematical Expressions}
Here are some example equations:

\subsection{Quadratic Formula}
The quadratic formula is given by:
\[x = \frac{-b \pm \sqrt{b^2 - 4ac}}{2a}\]

\subsection{Matrix Example}
A sample matrix:
\[
\begin{pmatrix}
a_{11} & a_{12} & a_{13} \\
a_{21} & a_{22} & a_{23} \\
a_{31} & a_{32} & a_{33}
\end{pmatrix}
\]

\section{Tables}
Below is a sample table showing data:

\begin{table}[h]
\centering
\begin{tabular}{@{}lrrr@{}}
\toprule
\textbf{Category} & \textbf{Value 1} & \textbf{Value 2} & \textbf{Total} \\
\midrule
Item A & 10.5 & 20.3 & 30.8 \\
Item B & 15.2 & 25.1 & 40.3 \\
Item C & 22.1 & 12.7 & 34.8 \\
\bottomrule
\end{tabular}
\caption{Sample Data Table}
\label{tab:sample}
\end{table}

\section{Lists}
\subsection{Bullet Points}
\begin{itemize}
  \item First item
  \item Second item
  \item Third item with sub-items:
    \begin{itemize}
      \item Sub-item 1
      \item Sub-item 2
    \end{itemize}
\end{itemize}

\subsection{Numbered List}
\begin{enumerate}
  \item First step
  \item Second step
  \item Third step
\end{enumerate}

\section{Citations and References}
When you use the citation search feature, references will be automatically managed. For example, you might cite a paper like this: \cite{example2024}. All citations will be automatically added to your \texttt{cite.bib} file.

\section{Conclusion}
This template demonstrates basic \LaTeX{} features including mathematical formulas, tables, and lists. The bibliography system is now set up to work seamlessly with Bibby AI's citation management. Feel free to modify it for your needs!

\bibliographystyle{plain}
\bibliography{cite}

\end{document}\documentclass{article}
\usepackage[utf8]{inputenc}
\usepackage{booktabs}
\usepackage{amsmath}
\usepackage{natbib}

\title{My Document}
\author{Your Name}
\date{\today}

\begin{document}

\maketitle

\section{Introduction}
This is a sample document demonstrating various \LaTeX{} features including mathematical equations, tables, and formatting. You can cite papers using the citation search feature, and they will be automatically added to your \texttt{cite.bib} file.

\section{Mathematical Expressions}
Here are some example equations:

\subsection{Quadratic Formula}
The quadratic formula is given by:
\[x = \frac{-b \pm \sqrt{b^2 - 4ac}}{2a}\]

\subsection{Matrix Example}
A sample matrix:
\[
\begin{pmatrix}
a_{11} & a_{12} & a_{13} \\
a_{21} & a_{22} & a_{23} \\
a_{31} & a_{32} & a_{33}
\end{pmatrix}
\]

\section{Tables}
Below is a sample table showing data:

\begin{table}[h]
\centering
\begin{tabular}{@{}lrrr@{}}
\toprule
\textbf{Category} & \textbf{Value 1} & \textbf{Value 2} & \textbf{Total} \\
\midrule
Item A & 10.5 & 20.3 & 30.8 \\
Item B & 15.2 & 25.1 & 40.3 \\
Item C & 22.1 & 12.7 & 34.8 \\
\bottomrule
\end{tabular}
\caption{Sample Data Table}
\label{tab:sample}
\end{table}

\section{Lists}
\subsection{Bullet Points}
\begin{itemize}
  \item First item
  \item Second item
  \item Third item with sub-items:
    \begin{itemize}
      \item Sub-item 1
      \item Sub-item 2
    \end{itemize}
\end{itemize}

\subsection{Numbered List}
\begin{enumerate}
  \item First step
  \item Second step
  \item Third step
\end{enumerate}

\section{Citations and References}
When you use the citation search feature, references will be automatically managed. For example, you might cite a paper like this: \cite{example2024}. All citations will be automatically added to your \texttt{cite.bib} file.

\section{Conclusion}
This template demonstrates basic \LaTeX{} features including mathematical formulas, tables, and lists. The bibliography system is now set up to work seamlessly with Bibby AI's citation management. Feel free to modify it for your needs!

\bibliographystyle{plain}
\bibliography{cite}

\end{document}\documentclass{article}
\usepackage[utf8]{inputenc}
\usepackage{booktabs}
\usepackage{amsmath}
\usepackage{natbib}

\title{My Document}
\author{Your Name}
\date{\today}

\begin{document}

\maketitle

\section{Introduction}
This is a sample document demonstrating various \LaTeX{} features including mathematical equations, tables, and formatting. You can cite papers using the citation search feature, and they will be automatically added to your \texttt{cite.bib} file.

\section{Mathematical Expressions}
Here are some example equations:

\subsection{Quadratic Formula}
The quadratic formula is given by:
\[x = \frac{-b \pm \sqrt{b^2 - 4ac}}{2a}\]

\subsection{Matrix Example}
A sample matrix:
\[
\begin{pmatrix}
a_{11} & a_{12} & a_{13} \\
a_{21} & a_{22} & a_{23} \\
a_{31} & a_{32} & a_{33}
\end{pmatrix}
\]

\section{Tables}
Below is a sample table showing data:

\begin{table}[h]
\centering
\begin{tabular}{@{}lrrr@{}}
\toprule
\textbf{Category} & \textbf{Value 1} & \textbf{Value 2} & \textbf{Total} \\
\midrule
Item A & 10.5 & 20.3 & 30.8 \\
Item B & 15.2 & 25.1 & 40.3 \\
Item C & 22.1 & 12.7 & 34.8 \\
\bottomrule
\end{tabular}
\caption{Sample Data Table}
\label{tab:sample}
\end{table}

\section{Lists}
\subsection{Bullet Points}
\begin{itemize}
  \item First item
  \item Second item
  \item Third item with sub-items:
    \begin{itemize}
      \item Sub-item 1
      \item Sub-item 2
    \end{itemize}
\end{itemize}

\subsection{Numbered List}
\begin{enumerate}
  \item First step
  \item Second step
  \item Third step
\end{enumerate}

\section{Citations and References}
When you use the citation search feature, references will be automatically managed. For example, you might cite a paper like this: \cite{example2024}. All citations will be automatically added to your \texttt{cite.bib} file.

\section{Conclusion}
This template demonstrates basic \LaTeX{} features including mathematical formulas, tables, and lists. The bibliography system is now set up to work seamlessly with Bibby AI's citation management. Feel free to modify it for your needs!

\bibliographystyle{plain}
\bibliography{cite}

\end{document}\documentclass{article}
\usepackage[utf8]{inputenc}
\usepackage{booktabs}
\usepackage{amsmath}
\usepackage{natbib}

\title{My Document}
\author{Your Name}
\date{\today}

\begin{document}

\maketitle

\section{Introduction}
This is a sample document demonstrating various \LaTeX{} features including mathematical equations, tables, and formatting. You can cite papers using the citation search feature, and they will be automatically added to your \texttt{cite.bib} file.

\section{Mathematical Expressions}
Here are some example equations:

\subsection{Quadratic Formula}
The quadratic formula is given by:
\[x = \frac{-b \pm \sqrt{b^2 - 4ac}}{2a}\]

\subsection{Matrix Example}
A sample matrix:
\[
\begin{pmatrix}
a_{11} & a_{12} & a_{13} \\
a_{21} & a_{22} & a_{23} \\
a_{31} & a_{32} & a_{33}
\end{pmatrix}
\]

\section{Tables}
Below is a sample table showing data:

\begin{table}[h]
\centering
\begin{tabular}{@{}lrrr@{}}
\toprule
\textbf{Category} & \textbf{Value 1} & \textbf{Value 2} & \textbf{Total} \\
\midrule
Item A & 10.5 & 20.3 & 30.8 \\
Item B & 15.2 & 25.1 & 40.3 \\
Item C & 22.1 & 12.7 & 34.8 \\
\bottomrule
\end{tabular}
\caption{Sample Data Table}
\label{tab:sample}
\end{table}

\section{Lists}
\subsection{Bullet Points}
\begin{itemize}
  \item First item
  \item Second item
  \item Third item with sub-items:
    \begin{itemize}
      \item Sub-item 1
      \item Sub-item 2
    \end{itemize}
\end{itemize}

\subsection{Numbered List}
\begin{enumerate}
  \item First step
  \item Second step
  \item Third step
\end{enumerate}

\section{Citations and References}
When you use the citation search feature, references will be automatically managed. For example, you might cite a paper like this: \cite{example2024}. All citations will be automatically added to your \texttt{cite.bib} file.

\section{Conclusion}
This template demonstrates basic \LaTeX{} features including mathematical formulas, tables, and lists. The bibliography system is now set up to work seamlessly with Bibby AI's citation management. Feel free to modify it for your needs!

\bibliographystyle{plain}
\bibliography{cite}

\end{document}\documentclass{article}
\usepackage[utf8]{inputenc}
\usepackage{booktabs}
\usepackage{amsmath}
\usepackage{natbib}

\title{My Document}
\author{Your Name}
\date{\today}

\begin{document}

\maketitle

\section{Introduction}
This is a sample document demonstrating various \LaTeX{} features including mathematical equations, tables, and formatting. You can cite papers using the citation search feature, and they will be automatically added to your \texttt{cite.bib} file.

\section{Mathematical Expressions}
Here are some example equations:

\subsection{Quadratic Formula}
The quadratic formula is given by:
\[x = \frac{-b \pm \sqrt{b^2 - 4ac}}{2a}\]

\subsection{Matrix Example}
A sample matrix:
\[
\begin{pmatrix}
a_{11} & a_{12} & a_{13} \\
a_{21} & a_{22} & a_{23} \\
a_{31} & a_{32} & a_{33}
\end{pmatrix}
\]

\section{Tables}
Below is a sample table showing data:

\begin{table}[h]
\centering
\begin{tabular}{@{}lrrr@{}}
\toprule
\textbf{Category} & \textbf{Value 1} & \textbf{Value 2} & \textbf{Total} \\
\midrule
Item A & 10.5 & 20.3 & 30.8 \\
Item B & 15.2 & 25.1 & 40.3 \\
Item C & 22.1 & 12.7 & 34.8 \\
\bottomrule
\end{tabular}
\caption{Sample Data Table}
\label{tab:sample}
\end{table}

\section{Lists}
\subsection{Bullet Points}
\begin{itemize}
  \item First item
  \item Second item
  \item Third item with sub-items:
    \begin{itemize}
      \item Sub-item 1
      \item Sub-item 2
    \end{itemize}
\end{itemize}

\subsection{Numbered List}
\begin{enumerate}
  \item First step
  \item Second step
  \item Third step
\end{enumerate}

\section{Citations and References}
When you use the citation search feature, references will be automatically managed. For example, you might cite a paper like this: \cite{example2024}. All citations will be automatically added to your \texttt{cite.bib} file.

\section{Conclusion}
This template demonstrates basic \LaTeX{} features including mathematical formulas, tables, and lists. The bibliography system is now set up to work seamlessly with Bibby AI's citation management. Feel free to modify it for your needs!

\bibliographystyle{plain}
\bibliography{cite}

\end{document}\documentclass{article}
\usepackage[utf8]{inputenc}
\usepackage{booktabs}
\usepackage{amsmath}
\usepackage{natbib}

\title{My Document}
\author{Your Name}
\date{\today}

\begin{document}

\maketitle

\section{Introduction}
This is a sample document demonstrating various \LaTeX{} features including mathematical equations, tables, and formatting. You can cite papers using the citation search feature, and they will be automatically added to your \texttt{cite.bib} file.

\section{Mathematical Expressions}
Here are some example equations:

\subsection{Quadratic Formula}
The quadratic formula is given by:
\[x = \frac{-b \pm \sqrt{b^2 - 4ac}}{2a}\]

\subsection{Matrix Example}
A sample matrix:
\[
\begin{pmatrix}
a_{11} & a_{12} & a_{13} \\
a_{21} & a_{22} & a_{23} \\
a_{31} & a_{32} & a_{33}
\end{pmatrix}
\]

\section{Tables}
Below is a sample table showing data:

\begin{table}[h]
\centering
\begin{tabular}{@{}lrrr@{}}
\toprule
\textbf{Category} & \textbf{Value 1} & \textbf{Value 2} & \textbf{Total} \\
\midrule
Item A & 10.5 & 20.3 & 30.8 \\
Item B & 15.2 & 25.1 & 40.3 \\
Item C & 22.1 & 12.7 & 34.8 \\
\bottomrule
\end{tabular}
\caption{Sample Data Table}
\label{tab:sample}
\end{table}

\section{Lists}
\subsection{Bullet Points}
\begin{itemize}
  \item First item
  \item Second item
  \item Third item with sub-items:
    \begin{itemize}
      \item Sub-item 1
      \item Sub-item 2
    \end{itemize}
\end{itemize}

\subsection{Numbered List}
\begin{enumerate}
  \item First step
  \item Second step
  \item Third step
\end{enumerate}

\section{Citations and References}
When you use the citation search feature, references will be automatically managed. For example, you might cite a paper like this: \cite{example2024}. All citations will be automatically added to your \texttt{cite.bib} file.

\section{Conclusion}
This template demonstrates basic \LaTeX{} features including mathematical formulas, tables, and lists. The bibliography system is now set up to work seamlessly with Bibby AI's citation management. Feel free to modify it for your needs!

\bibliographystyle{plain}
\bibliography{cite}

\end{document}\documentclass{article}
\usepackage[utf8]{inputenc}
\usepackage{booktabs}
\usepackage{amsmath}
\usepackage{natbib}

\title{My Document}
\author{Your Name}
\date{\today}

\begin{document}

\maketitle

\section{Introduction}
This is a sample document demonstrating various \LaTeX{} features including mathematical equations, tables, and formatting. You can cite papers using the citation search feature, and they will be automatically added to your \texttt{cite.bib} file.

\section{Mathematical Expressions}
Here are some example equations:

\subsection{Quadratic Formula}
The quadratic formula is given by:
\[x = \frac{-b \pm \sqrt{b^2 - 4ac}}{2a}\]

\subsection{Matrix Example}
A sample matrix:
\[
\begin{pmatrix}
a_{11} & a_{12} & a_{13} \\
a_{21} & a_{22} & a_{23} \\
a_{31} & a_{32} & a_{33}
\end{pmatrix}
\]

\section{Tables}
Below is a sample table showing data:

\begin{table}[h]
\centering
\begin{tabular}{@{}lrrr@{}}
\toprule
\textbf{Category} & \textbf{Value 1} & \textbf{Value 2} & \textbf{Total} \\
\midrule
Item A & 10.5 & 20.3 & 30.8 \\
Item B & 15.2 & 25.1 & 40.3 \\
Item C & 22.1 & 12.7 & 34.8 \\
\bottomrule
\end{tabular}
\caption{Sample Data Table}
\label{tab:sample}
\end{table}

\section{Lists}
\subsection{Bullet Points}
\begin{itemize}
  \item First item
  \item Second item
  \item Third item with sub-items:
    \begin{itemize}
      \item Sub-item 1
      \item Sub-item 2
    \end{itemize}
\end{itemize}

\subsection{Numbered List}
\begin{enumerate}
  \item First step
  \item Second step
  \item Third step
\end{enumerate}

\section{Citations and References}
When you use the citation search feature, references will be automatically managed. For example, you might cite a paper like this: \cite{example2024}. All citations will be automatically added to your \texttt{cite.bib} file.

\section{Conclusion}
This template demonstrates basic \LaTeX{} features including mathematical formulas, tables, and lists. The bibliography system is now set up to work seamlessly with Bibby AI's citation management. Feel free to modify it for your needs!

\bibliographystyle{plain}
\bibliography{cite}

\end{document}\documentclass{article}
\usepackage[utf8]{inputenc}
\usepackage{booktabs}
\usepackage{amsmath}
\usepackage{natbib}

\title{My Document}
\author{Your Name}
\date{\today}

\begin{document}

\maketitle

\section{Introduction}
This is a sample document demonstrating various \LaTeX{} features including mathematical equations, tables, and formatting. You can cite papers using the citation search feature, and they will be automatically added to your \texttt{cite.bib} file.

\section{Mathematical Expressions}
Here are some example equations:

\subsection{Quadratic Formula}
The quadratic formula is given by:
\[x = \frac{-b \pm \sqrt{b^2 - 4ac}}{2a}\]

\subsection{Matrix Example}
A sample matrix:
\[
\begin{pmatrix}
a_{11} & a_{12} & a_{13} \\
a_{21} & a_{22} & a_{23} \\
a_{31} & a_{32} & a_{33}
\end{pmatrix}
\]

\section{Tables}
Below is a sample table showing data:

\begin{table}[h]
\centering
\begin{tabular}{@{}lrrr@{}}
\toprule
\textbf{Category} & \textbf{Value 1} & \textbf{Value 2} & \textbf{Total} \\
\midrule
Item A & 10.5 & 20.3 & 30.8 \\
Item B & 15.2 & 25.1 & 40.3 \\
Item C & 22.1 & 12.7 & 34.8 \\
\bottomrule
\end{tabular}
\caption{Sample Data Table}
\label{tab:sample}
\end{table}

\section{Lists}
\subsection{Bullet Points}
\begin{itemize}
  \item First item
  \item Second item
  \item Third item with sub-items:
    \begin{itemize}
      \item Sub-item 1
      \item Sub-item 2
    \end{itemize}
\end{itemize}

\subsection{Numbered List}
\begin{enumerate}
  \item First step
  \item Second step
  \item Third step
\end{enumerate}

\section{Citations and References}
When you use the citation search feature, references will be automatically managed. For example, you might cite a paper like this: \cite{example2024}. All citations will be automatically added to your \texttt{cite.bib} file.

\section{Conclusion}
This template demonstrates basic \LaTeX{} features including mathematical formulas, tables, and lists. The bibliography system is now set up to work seamlessly with Bibby AI's citation management. Feel free to modify it for your needs!

\bibliographystyle{plain}
\bibliography{cite}

\end{document}\documentclass{article}
\usepackage[utf8]{inputenc}
\usepackage{booktabs}
\usepackage{amsmath}
\usepackage{natbib}

\title{My Document}
\author{Your Name}
\date{\today}

\begin{document}

\maketitle

\section{Introduction}
This is a sample document demonstrating various \LaTeX{} features including mathematical equations, tables, and formatting. You can cite papers using the citation search feature, and they will be automatically added to your \texttt{cite.bib} file.

\section{Mathematical Expressions}
Here are some example equations:

\subsection{Quadratic Formula}
The quadratic formula is given by:
\[x = \frac{-b \pm \sqrt{b^2 - 4ac}}{2a}\]

\subsection{Matrix Example}
A sample matrix:
\[
\begin{pmatrix}
a_{11} & a_{12} & a_{13} \\
a_{21} & a_{22} & a_{23} \\
a_{31} & a_{32} & a_{33}
\end{pmatrix}
\]

\section{Tables}
Below is a sample table showing data:

\begin{table}[h]
\centering
\begin{tabular}{@{}lrrr@{}}
\toprule
\textbf{Category} & \textbf{Value 1} & \textbf{Value 2} & \textbf{Total} \\
\midrule
Item A & 10.5 & 20.3 & 30.8 \\
Item B & 15.2 & 25.1 & 40.3 \\
Item C & 22.1 & 12.7 & 34.8 \\
\bottomrule
\end{tabular}
\caption{Sample Data Table}
\label{tab:sample}
\end{table}

\section{Lists}
\subsection{Bullet Points}
\begin{itemize}
  \item First item
  \item Second item
  \item Third item with sub-items:
    \begin{itemize}
      \item Sub-item 1
      \item Sub-item 2
    \end{itemize}
\end{itemize}

\subsection{Numbered List}
\begin{enumerate}
  \item First step
  \item Second step
  \item Third step
\end{enumerate}

\section{Citations and References}
When you use the citation search feature, references will be automatically managed. For example, you might cite a paper like this: \cite{example2024}. All citations will be automatically added to your \texttt{cite.bib} file.

\section{Conclusion}
This template demonstrates basic \LaTeX{} features including mathematical formulas, tables, and lists. The bibliography system is now set up to work seamlessly with Bibby AI's citation management. Feel free to modify it for your needs!

\bibliographystyle{plain}
\bibliography{cite}

\end{document}\documentclass{article}
\usepackage[utf8]{inputenc}
\usepackage{booktabs}
\usepackage{amsmath}
\usepackage{natbib}

\title{My Document}
\author{Your Name}
\date{\today}

\begin{document}

\maketitle

\section{Introduction}
This is a sample document demonstrating various \LaTeX{} features including mathematical equations, tables, and formatting. You can cite papers using the citation search feature, and they will be automatically added to your \texttt{cite.bib} file.

\section{Mathematical Expressions}
Here are some example equations:

\subsection{Quadratic Formula}
The quadratic formula is given by:
\[x = \frac{-b \pm \sqrt{b^2 - 4ac}}{2a}\]

\subsection{Matrix Example}
A sample matrix:
\[
\begin{pmatrix}
a_{11} & a_{12} & a_{13} \\
a_{21} & a_{22} & a_{23} \\
a_{31} & a_{32} & a_{33}
\end{pmatrix}
\]

\section{Tables}
Below is a sample table showing data:

\begin{table}[h]
\centering
\begin{tabular}{@{}lrrr@{}}
\toprule
\textbf{Category} & \textbf{Value 1} & \textbf{Value 2} & \textbf{Total} \\
\midrule
Item A & 10.5 & 20.3 & 30.8 \\
Item B & 15.2 & 25.1 & 40.3 \\
Item C & 22.1 & 12.7 & 34.8 \\
\bottomrule
\end{tabular}
\caption{Sample Data Table}
\label{tab:sample}
\end{table}

\section{Lists}
\subsection{Bullet Points}
\begin{itemize}
  \item First item
  \item Second item
  \item Third item with sub-items:
    \begin{itemize}
      \item Sub-item 1
      \item Sub-item 2
    \end{itemize}
\end{itemize}

\subsection{Numbered List}
\begin{enumerate}
  \item First step
  \item Second step
  \item Third step
\end{enumerate}

\section{Citations and References}
When you use the citation search feature, references will be automatically managed. For example, you might cite a paper like this: \cite{example2024}. All citations will be automatically added to your \texttt{cite.bib} file.

\section{Conclusion}
This template demonstrates basic \LaTeX{} features including mathematical formulas, tables, and lists. The bibliography system is now set up to work seamlessly with Bibby AI's citation management. Feel free to modify it for your needs!

\bibliographystyle{plain}
\bibliography{cite}

\end{document}\documentclass{article}
\usepackage[utf8]{inputenc}
\usepackage{booktabs}
\usepackage{amsmath}
\usepackage{natbib}

\title{My Document}
\author{Your Name}
\date{\today}

\begin{document}

\maketitle

\section{Introduction}
This is a sample document demonstrating various \LaTeX{} features including mathematical equations, tables, and formatting. You can cite papers using the citation search feature, and they will be automatically added to your \texttt{cite.bib} file.

\section{Mathematical Expressions}
Here are some example equations:

\subsection{Quadratic Formula}
The quadratic formula is given by:
\[x = \frac{-b \pm \sqrt{b^2 - 4ac}}{2a}\]

\subsection{Matrix Example}
A sample matrix:
\[
\begin{pmatrix}
a_{11} & a_{12} & a_{13} \\
a_{21} & a_{22} & a_{23} \\
a_{31} & a_{32} & a_{33}
\end{pmatrix}
\]

\section{Tables}
Below is a sample table showing data:

\begin{table}[h]
\centering
\begin{tabular}{@{}lrrr@{}}
\toprule
\textbf{Category} & \textbf{Value 1} & \textbf{Value 2} & \textbf{Total} \\
\midrule
Item A & 10.5 & 20.3 & 30.8 \\
Item B & 15.2 & 25.1 & 40.3 \\
Item C & 22.1 & 12.7 & 34.8 \\
\bottomrule
\end{tabular}
\caption{Sample Data Table}
\label{tab:sample}
\end{table}

\section{Lists}
\subsection{Bullet Points}
\begin{itemize}
  \item First item
  \item Second item
  \item Third item with sub-items:
    \begin{itemize}
      \item Sub-item 1
      \item Sub-item 2
    \end{itemize}
\end{itemize}

\subsection{Numbered List}
\begin{enumerate}
  \item First step
  \item Second step
  \item Third step
\end{enumerate}

\section{Citations and References}
When you use the citation search feature, references will be automatically managed. For example, you might cite a paper like this: \cite{example2024}. All citations will be automatically added to your \texttt{cite.bib} file.

\section{Conclusion}
This template demonstrates basic \LaTeX{} features including mathematical formulas, tables, and lists. The bibliography system is now set up to work seamlessly with Bibby AI's citation management. Feel free to modify it for your needs!

\bibliographystyle{plain}
\bibliography{cite}

\end{document}\documentclass{article}
\usepackage[utf8]{inputenc}
\usepackage{booktabs}
\usepackage{amsmath}
\usepackage{natbib}

\title{My Document}
\author{Your Name}
\date{\today}

\begin{document}

\maketitle

\section{Introduction}
This is a sample document demonstrating various \LaTeX{} features including mathematical equations, tables, and formatting. You can cite papers using the citation search feature, and they will be automatically added to your \texttt{cite.bib} file.

\section{Mathematical Expressions}
Here are some example equations:

\subsection{Quadratic Formula}
The quadratic formula is given by:
\[x = \frac{-b \pm \sqrt{b^2 - 4ac}}{2a}\]

\subsection{Matrix Example}
A sample matrix:
\[
\begin{pmatrix}
a_{11} & a_{12} & a_{13} \\
a_{21} & a_{22} & a_{23} \\
a_{31} & a_{32} & a_{33}
\end{pmatrix}
\]

\section{Tables}
Below is a sample table showing data:

\begin{table}[h]
\centering
\begin{tabular}{@{}lrrr@{}}
\toprule
\textbf{Category} & \textbf{Value 1} & \textbf{Value 2} & \textbf{Total} \\
\midrule
Item A & 10.5 & 20.3 & 30.8 \\
Item B & 15.2 & 25.1 & 40.3 \\
Item C & 22.1 & 12.7 & 34.8 \\
\bottomrule
\end{tabular}
\caption{Sample Data Table}
\label{tab:sample}
\end{table}

\section{Lists}
\subsection{Bullet Points}
\begin{itemize}
  \item First item
  \item Second item
  \item Third item with sub-items:
    \begin{itemize}
      \item Sub-item 1
      \item Sub-item 2
    \end{itemize}
\end{itemize}

\subsection{Numbered List}
\begin{enumerate}
  \item First step
  \item Second step
  \item Third step
\end{enumerate}

\section{Citations and References}
When you use the citation search feature, references will be automatically managed. For example, you might cite a paper like this: \cite{example2024}. All citations will be automatically added to your \texttt{cite.bib} file.

\section{Conclusion}
This template demonstrates basic \LaTeX{} features including mathematical formulas, tables, and lists. The bibliography system is now set up to work seamlessly with Bibby AI's citation management. Feel free to modify it for your needs!

\bibliographystyle{plain}
\bibliography{cite}

\end{document}\documentclass{article}
\usepackage[utf8]{inputenc}
\usepackage{booktabs}
\usepackage{amsmath}
\usepackage{natbib}

\title{My Document}
\author{Your Name}
\date{\today}

\begin{document}

\maketitle

\section{Introduction}
This is a sample document demonstrating various \LaTeX{} features including mathematical equations, tables, and formatting. You can cite papers using the citation search feature, and they will be automatically added to your \texttt{cite.bib} file.

\section{Mathematical Expressions}
Here are some example equations:

\subsection{Quadratic Formula}
The quadratic formula is given by:
\[x = \frac{-b \pm \sqrt{b^2 - 4ac}}{2a}\]

\subsection{Matrix Example}
A sample matrix:
\[
\begin{pmatrix}
a_{11} & a_{12} & a_{13} \\
a_{21} & a_{22} & a_{23} \\
a_{31} & a_{32} & a_{33}
\end{pmatrix}
\]

\section{Tables}
Below is a sample table showing data:

\begin{table}[h]
\centering
\begin{tabular}{@{}lrrr@{}}
\toprule
\textbf{Category} & \textbf{Value 1} & \textbf{Value 2} & \textbf{Total} \\
\midrule
Item A & 10.5 & 20.3 & 30.8 \\
Item B & 15.2 & 25.1 & 40.3 \\
Item C & 22.1 & 12.7 & 34.8 \\
\bottomrule
\end{tabular}
\caption{Sample Data Table}
\label{tab:sample}
\end{table}

\section{Lists}
\subsection{Bullet Points}
\begin{itemize}
  \item First item
  \item Second item
  \item Third item with sub-items:
    \begin{itemize}
      \item Sub-item 1
      \item Sub-item 2
    \end{itemize}
\end{itemize}

\subsection{Numbered List}
\begin{enumerate}
  \item First step
  \item Second step
  \item Third step
\end{enumerate}

\section{Citations and References}
When you use the citation search feature, references will be automatically managed. For example, you might cite a paper like this: \cite{example2024}. All citations will be automatically added to your \texttt{cite.bib} file.

\section{Conclusion}
This template demonstrates basic \LaTeX{} features including mathematical formulas, tables, and lists. The bibliography system is now set up to work seamlessly with Bibby AI's citation management. Feel free to modify it for your needs!

\bibliographystyle{plain}
\bibliography{cite}

\end{document}\documentclass{article}
\usepackage[utf8]{inputenc}
\usepackage{booktabs}
\usepackage{amsmath}
\usepackage{natbib}

\title{My Document}
\author{Your Name}
\date{\today}

\begin{document}

\maketitle

\section{Introduction}
This is a sample document demonstrating various \LaTeX{} features including mathematical equations, tables, and formatting. You can cite papers using the citation search feature, and they will be automatically added to your \texttt{cite.bib} file.

\section{Mathematical Expressions}
Here are some example equations:

\subsection{Quadratic Formula}
The quadratic formula is given by:
\[x = \frac{-b \pm \sqrt{b^2 - 4ac}}{2a}\]

\subsection{Matrix Example}
A sample matrix:
\[
\begin{pmatrix}
a_{11} & a_{12} & a_{13} \\
a_{21} & a_{22} & a_{23} \\
a_{31} & a_{32} & a_{33}
\end{pmatrix}
\]

\section{Tables}
Below is a sample table showing data:

\begin{table}[h]
\centering
\begin{tabular}{@{}lrrr@{}}
\toprule
\textbf{Category} & \textbf{Value 1} & \textbf{Value 2} & \textbf{Total} \\
\midrule
Item A & 10.5 & 20.3 & 30.8 \\
Item B & 15.2 & 25.1 & 40.3 \\
Item C & 22.1 & 12.7 & 34.8 \\
\bottomrule
\end{tabular}
\caption{Sample Data Table}
\label{tab:sample}
\end{table}

\section{Lists}
\subsection{Bullet Points}
\begin{itemize}
  \item First item
  \item Second item
  \item Third item with sub-items:
    \begin{itemize}
      \item Sub-item 1
      \item Sub-item 2
    \end{itemize}
\end{itemize}

\subsection{Numbered List}
\begin{enumerate}
  \item First step
  \item Second step
  \item Third step
\end{enumerate}

\section{Citations and References}
When you use the citation search feature, references will be automatically managed. For example, you might cite a paper like this: \cite{example2024}. All citations will be automatically added to your \texttt{cite.bib} file.

\section{Conclusion}
This template demonstrates basic \LaTeX{} features including mathematical formulas, tables, and lists. The bibliography system is now set up to work seamlessly with Bibby AI's citation management. Feel free to modify it for your needs!

\bibliographystyle{plain}
\bibliography{cite}

\end{document}\documentclass{article}
\usepackage[utf8]{inputenc}
\usepackage{booktabs}
\usepackage{amsmath}
\usepackage{natbib}

\title{My Document}
\author{Your Name}
\date{\today}

\begin{document}

\maketitle

\section{Introduction}
This is a sample document demonstrating various \LaTeX{} features including mathematical equations, tables, and formatting. You can cite papers using the citation search feature, and they will be automatically added to your \texttt{cite.bib} file.

\section{Mathematical Expressions}
Here are some example equations:

\subsection{Quadratic Formula}
The quadratic formula is given by:
\[x = \frac{-b \pm \sqrt{b^2 - 4ac}}{2a}\]

\subsection{Matrix Example}
A sample matrix:
\[
\begin{pmatrix}
a_{11} & a_{12} & a_{13} \\
a_{21} & a_{22} & a_{23} \\
a_{31} & a_{32} & a_{33}
\end{pmatrix}
\]

\section{Tables}
Below is a sample table showing data:

\begin{table}[h]
\centering
\begin{tabular}{@{}lrrr@{}}
\toprule
\textbf{Category} & \textbf{Value 1} & \textbf{Value 2} & \textbf{Total} \\
\midrule
Item A & 10.5 & 20.3 & 30.8 \\
Item B & 15.2 & 25.1 & 40.3 \\
Item C & 22.1 & 12.7 & 34.8 \\
\bottomrule
\end{tabular}
\caption{Sample Data Table}
\label{tab:sample}
\end{table}

\section{Lists}
\subsection{Bullet Points}
\begin{itemize}
  \item First item
  \item Second item
  \item Third item with sub-items:
    \begin{itemize}
      \item Sub-item 1
      \item Sub-item 2
    \end{itemize}
\end{itemize}

\subsection{Numbered List}
\begin{enumerate}
  \item First step
  \item Second step
  \item Third step
\end{enumerate}

\section{Citations and References}
When you use the citation search feature, references will be automatically managed. For example, you might cite a paper like this: \cite{example2024}. All citations will be automatically added to your \texttt{cite.bib} file.

\section{Conclusion}
This template demonstrates basic \LaTeX{} features including mathematical formulas, tables, and lists. The bibliography system is now set up to work seamlessly with Bibby AI's citation management. Feel free to modify it for your needs!

\bibliographystyle{plain}
\bibliography{cite}

\end{document}\documentclass{article}
\usepackage[utf8]{inputenc}
\usepackage{booktabs}
\usepackage{amsmath}
\usepackage{natbib}

\title{My Document}
\author{Your Name}
\date{\today}

\begin{document}

\maketitle

\section{Introduction}
This is a sample document demonstrating various \LaTeX{} features including mathematical equations, tables, and formatting. You can cite papers using the citation search feature, and they will be automatically added to your \texttt{cite.bib} file.

\section{Mathematical Expressions}
Here are some example equations:

\subsection{Quadratic Formula}
The quadratic formula is given by:
\[x = \frac{-b \pm \sqrt{b^2 - 4ac}}{2a}\]

\subsection{Matrix Example}
A sample matrix:
\[
\begin{pmatrix}
a_{11} & a_{12} & a_{13} \\
a_{21} & a_{22} & a_{23} \\
a_{31} & a_{32} & a_{33}
\end{pmatrix}
\]

\section{Tables}
Below is a sample table showing data:

\begin{table}[h]
\centering
\begin{tabular}{@{}lrrr@{}}
\toprule
\textbf{Category} & \textbf{Value 1} & \textbf{Value 2} & \textbf{Total} \\
\midrule
Item A & 10.5 & 20.3 & 30.8 \\
Item B & 15.2 & 25.1 & 40.3 \\
Item C & 22.1 & 12.7 & 34.8 \\
\bottomrule
\end{tabular}
\caption{Sample Data Table}
\label{tab:sample}
\end{table}

\section{Lists}
\subsection{Bullet Points}
\begin{itemize}
  \item First item
  \item Second item
  \item Third item with sub-items:
    \begin{itemize}
      \item Sub-item 1
      \item Sub-item 2
    \end{itemize}
\end{itemize}

\subsection{Numbered List}
\begin{enumerate}
  \item First step
  \item Second step
  \item Third step
\end{enumerate}

\section{Citations and References}
When you use the citation search feature, references will be automatically managed. For example, you might cite a paper like this: \cite{example2024}. All citations will be automatically added to your \texttt{cite.bib} file.

\section{Conclusion}
This template demonstrates basic \LaTeX{} features including mathematical formulas, tables, and lists. The bibliography system is now set up to work seamlessly with Bibby AI's citation management. Feel free to modify it for your needs!

\bibliographystyle{plain}
\bibliography{cite}

\end{document}\documentclass{article}
\usepackage[utf8]{inputenc}
\usepackage{booktabs}
\usepackage{amsmath}
\usepackage{natbib}

\title{My Document}
\author{Your Name}
\date{\today}

\begin{document}

\maketitle

\section{Introduction}
This is a sample document demonstrating various \LaTeX{} features including mathematical equations, tables, and formatting. You can cite papers using the citation search feature, and they will be automatically added to your \texttt{cite.bib} file.

\section{Mathematical Expressions}
Here are some example equations:

\subsection{Quadratic Formula}
The quadratic formula is given by:
\[x = \frac{-b \pm \sqrt{b^2 - 4ac}}{2a}\]

\subsection{Matrix Example}
A sample matrix:
\[
\begin{pmatrix}
a_{11} & a_{12} & a_{13} \\
a_{21} & a_{22} & a_{23} \\
a_{31} & a_{32} & a_{33}
\end{pmatrix}
\]

\section{Tables}
Below is a sample table showing data:

\begin{table}[h]
\centering
\begin{tabular}{@{}lrrr@{}}
\toprule
\textbf{Category} & \textbf{Value 1} & \textbf{Value 2} & \textbf{Total} \\
\midrule
Item A & 10.5 & 20.3 & 30.8 \\
Item B & 15.2 & 25.1 & 40.3 \\
Item C & 22.1 & 12.7 & 34.8 \\
\bottomrule
\end{tabular}
\caption{Sample Data Table}
\label{tab:sample}
\end{table}

\section{Lists}
\subsection{Bullet Points}
\begin{itemize}
  \item First item
  \item Second item
  \item Third item with sub-items:
    \begin{itemize}
      \item Sub-item 1
      \item Sub-item 2
    \end{itemize}
\end{itemize}

\subsection{Numbered List}
\begin{enumerate}
  \item First step
  \item Second step
  \item Third step
\end{enumerate}

\section{Citations and References}
When you use the citation search feature, references will be automatically managed. For example, you might cite a paper like this: \cite{example2024}. All citations will be automatically added to your \texttt{cite.bib} file.

\section{Conclusion}
This template demonstrates basic \LaTeX{} features including mathematical formulas, tables, and lists. The bibliography system is now set up to work seamlessly with Bibby AI's citation management. Feel free to modify it for your needs!

\bibliographystyle{plain}
\bibliography{cite}

\end{document}\documentclass{article}
\usepackage[utf8]{inputenc}
\usepackage{booktabs}
\usepackage{amsmath}
\usepackage{natbib}

\title{My Document}
\author{Your Name}
\date{\today}

\begin{document}

\maketitle

\section{Introduction}
This is a sample document demonstrating various \LaTeX{} features including mathematical equations, tables, and formatting. You can cite papers using the citation search feature, and they will be automatically added to your \texttt{cite.bib} file.

\section{Mathematical Expressions}
Here are some example equations:

\subsection{Quadratic Formula}
The quadratic formula is given by:
\[x = \frac{-b \pm \sqrt{b^2 - 4ac}}{2a}\]

\subsection{Matrix Example}
A sample matrix:
\[
\begin{pmatrix}
a_{11} & a_{12} & a_{13} \\
a_{21} & a_{22} & a_{23} \\
a_{31} & a_{32} & a_{33}
\end{pmatrix}
\]

\section{Tables}
Below is a sample table showing data:

\begin{table}[h]
\centering
\begin{tabular}{@{}lrrr@{}}
\toprule
\textbf{Category} & \textbf{Value 1} & \textbf{Value 2} & \textbf{Total} \\
\midrule
Item A & 10.5 & 20.3 & 30.8 \\
Item B & 15.2 & 25.1 & 40.3 \\
Item C & 22.1 & 12.7 & 34.8 \\
\bottomrule
\end{tabular}
\caption{Sample Data Table}
\label{tab:sample}
\end{table}

\section{Lists}
\subsection{Bullet Points}
\begin{itemize}
  \item First item
  \item Second item
  \item Third item with sub-items:
    \begin{itemize}
      \item Sub-item 1
      \item Sub-item 2
    \end{itemize}
\end{itemize}

\subsection{Numbered List}
\begin{enumerate}
  \item First step
  \item Second step
  \item Third step
\end{enumerate}

\section{Citations and References}
When you use the citation search feature, references will be automatically managed. For example, you might cite a paper like this: \cite{example2024}. All citations will be automatically added to your \texttt{cite.bib} file.

\section{Conclusion}
This template demonstrates basic \LaTeX{} features including mathematical formulas, tables, and lists. The bibliography system is now set up to work seamlessly with Bibby AI's citation management. Feel free to modify it for your needs!

\bibliographystyle{plain}
\bibliography{cite}

\end{document}\documentclass{article}
\usepackage[utf8]{inputenc}
\usepackage{booktabs}
\usepackage{amsmath}
\usepackage{natbib}

\title{My Document}
\author{Your Name}
\date{\today}

\begin{document}

\maketitle

\section{Introduction}
This is a sample document demonstrating various \LaTeX{} features including mathematical equations, tables, and formatting. You can cite papers using the citation search feature, and they will be automatically added to your \texttt{cite.bib} file.

\section{Mathematical Expressions}
Here are some example equations:

\subsection{Quadratic Formula}
The quadratic formula is given by:
\[x = \frac{-b \pm \sqrt{b^2 - 4ac}}{2a}\]

\subsection{Matrix Example}
A sample matrix:
\[
\begin{pmatrix}
a_{11} & a_{12} & a_{13} \\
a_{21} & a_{22} & a_{23} \\
a_{31} & a_{32} & a_{33}
\end{pmatrix}
\]

\section{Tables}
Below is a sample table showing data:

\begin{table}[h]
\centering
\begin{tabular}{@{}lrrr@{}}
\toprule
\textbf{Category} & \textbf{Value 1} & \textbf{Value 2} & \textbf{Total} \\
\midrule
Item A & 10.5 & 20.3 & 30.8 \\
Item B & 15.2 & 25.1 & 40.3 \\
Item C & 22.1 & 12.7 & 34.8 \\
\bottomrule
\end{tabular}
\caption{Sample Data Table}
\label{tab:sample}
\end{table}

\section{Lists}
\subsection{Bullet Points}
\begin{itemize}
  \item First item
  \item Second item
  \item Third item with sub-items:
    \begin{itemize}
      \item Sub-item 1
      \item Sub-item 2
    \end{itemize}
\end{itemize}

\subsection{Numbered List}
\begin{enumerate}
  \item First step
  \item Second step
  \item Third step
\end{enumerate}

\section{Citations and References}
When you use the citation search feature, references will be automatically managed. For example, you might cite a paper like this: \cite{example2024}. All citations will be automatically added to your \texttt{cite.bib} file.

\section{Conclusion}
This template demonstrates basic \LaTeX{} features including mathematical formulas, tables, and lists. The bibliography system is now set up to work seamlessly with Bibby AI's citation management. Feel free to modify it for your needs!

\bibliographystyle{plain}
\bibliography{cite}

\end{document}\documentclass{article}
\usepackage[utf8]{inputenc}
\usepackage{booktabs}
\usepackage{amsmath}
\usepackage{natbib}

\title{My Document}
\author{Your Name}
\date{\today}

\begin{document}

\maketitle

\section{Introduction}
This is a sample document demonstrating various \LaTeX{} features including mathematical equations, tables, and formatting. You can cite papers using the citation search feature, and they will be automatically added to your \texttt{cite.bib} file.

\section{Mathematical Expressions}
Here are some example equations:

\subsection{Quadratic Formula}
The quadratic formula is given by:
\[x = \frac{-b \pm \sqrt{b^2 - 4ac}}{2a}\]

\subsection{Matrix Example}
A sample matrix:
\[
\begin{pmatrix}
a_{11} & a_{12} & a_{13} \\
a_{21} & a_{22} & a_{23} \\
a_{31} & a_{32} & a_{33}
\end{pmatrix}
\]

\section{Tables}
Below is a sample table showing data:

\begin{table}[h]
\centering
\begin{tabular}{@{}lrrr@{}}
\toprule
\textbf{Category} & \textbf{Value 1} & \textbf{Value 2} & \textbf{Total} \\
\midrule
Item A & 10.5 & 20.3 & 30.8 \\
Item B & 15.2 & 25.1 & 40.3 \\
Item C & 22.1 & 12.7 & 34.8 \\
\bottomrule
\end{tabular}
\caption{Sample Data Table}
\label{tab:sample}
\end{table}

\section{Lists}
\subsection{Bullet Points}
\begin{itemize}
  \item First item
  \item Second item
  \item Third item with sub-items:
    \begin{itemize}
      \item Sub-item 1
      \item Sub-item 2
    \end{itemize}
\end{itemize}

\subsection{Numbered List}
\begin{enumerate}
  \item First step
  \item Second step
  \item Third step
\end{enumerate}

\section{Citations and References}
When you use the citation search feature, references will be automatically managed. For example, you might cite a paper like this: \cite{example2024}. All citations will be automatically added to your \texttt{cite.bib} file.

\section{Conclusion}
This template demonstrates basic \LaTeX{} features including mathematical formulas, tables, and lists. The bibliography system is now set up to work seamlessly with Bibby AI's citation management. Feel free to modify it for your needs!

\bibliographystyle{plain}
\bibliography{cite}

\end{document}\documentclass{article}
\usepackage[utf8]{inputenc}
\usepackage{booktabs}
\usepackage{amsmath}
\usepackage{natbib}

\title{My Document}
\author{Your Name}
\date{\today}

\begin{document}

\maketitle

\section{Introduction}
This is a sample document demonstrating various \LaTeX{} features including mathematical equations, tables, and formatting. You can cite papers using the citation search feature, and they will be automatically added to your \texttt{cite.bib} file.

\section{Mathematical Expressions}
Here are some example equations:

\subsection{Quadratic Formula}
The quadratic formula is given by:
\[x = \frac{-b \pm \sqrt{b^2 - 4ac}}{2a}\]

\subsection{Matrix Example}
A sample matrix:
\[
\begin{pmatrix}
a_{11} & a_{12} & a_{13} \\
a_{21} & a_{22} & a_{23} \\
a_{31} & a_{32} & a_{33}
\end{pmatrix}
\]

\section{Tables}
Below is a sample table showing data:

\begin{table}[h]
\centering
\begin{tabular}{@{}lrrr@{}}
\toprule
\textbf{Category} & \textbf{Value 1} & \textbf{Value 2} & \textbf{Total} \\
\midrule
Item A & 10.5 & 20.3 & 30.8 \\
Item B & 15.2 & 25.1 & 40.3 \\
Item C & 22.1 & 12.7 & 34.8 \\
\bottomrule
\end{tabular}
\caption{Sample Data Table}
\label{tab:sample}
\end{table}

\section{Lists}
\subsection{Bullet Points}
\begin{itemize}
  \item First item
  \item Second item
  \item Third item with sub-items:
    \begin{itemize}
      \item Sub-item 1
      \item Sub-item 2
    \end{itemize}
\end{itemize}

\subsection{Numbered List}
\begin{enumerate}
  \item First step
  \item Second step
  \item Third step
\end{enumerate}

\section{Citations and References}
When you use the citation search feature, references will be automatically managed. For example, you might cite a paper like this: \cite{example2024}. All citations will be automatically added to your \texttt{cite.bib} file.

\section{Conclusion}
This template demonstrates basic \LaTeX{} features including mathematical formulas, tables, and lists. The bibliography system is now set up to work seamlessly with Bibby AI's citation management. Feel free to modify it for your needs!

\bibliographystyle{plain}
\bibliography{cite}

\end{document}\documentclass{article}
\usepackage[utf8]{inputenc}
\usepackage{booktabs}
\usepackage{amsmath}
\usepackage{natbib}

\title{My Document}
\author{Your Name}
\date{\today}

\begin{document}

\maketitle

\section{Introduction}
This is a sample document demonstrating various \LaTeX{} features including mathematical equations, tables, and formatting. You can cite papers using the citation search feature, and they will be automatically added to your \texttt{cite.bib} file.

\section{Mathematical Expressions}
Here are some example equations:

\subsection{Quadratic Formula}
The quadratic formula is given by:
\[x = \frac{-b \pm \sqrt{b^2 - 4ac}}{2a}\]

\subsection{Matrix Example}
A sample matrix:
\[
\begin{pmatrix}
a_{11} & a_{12} & a_{13} \\
a_{21} & a_{22} & a_{23} \\
a_{31} & a_{32} & a_{33}
\end{pmatrix}
\]

\section{Tables}
Below is a sample table showing data:

\begin{table}[h]
\centering
\begin{tabular}{@{}lrrr@{}}
\toprule
\textbf{Category} & \textbf{Value 1} & \textbf{Value 2} & \textbf{Total} \\
\midrule
Item A & 10.5 & 20.3 & 30.8 \\
Item B & 15.2 & 25.1 & 40.3 \\
Item C & 22.1 & 12.7 & 34.8 \\
\bottomrule
\end{tabular}
\caption{Sample Data Table}
\label{tab:sample}
\end{table}

\section{Lists}
\subsection{Bullet Points}
\begin{itemize}
  \item First item
  \item Second item
  \item Third item with sub-items:
    \begin{itemize}
      \item Sub-item 1
      \item Sub-item 2
    \end{itemize}
\end{itemize}

\subsection{Numbered List}
\begin{enumerate}
  \item First step
  \item Second step
  \item Third step
\end{enumerate}

\section{Citations and References}
When you use the citation search feature, references will be automatically managed. For example, you might cite a paper like this: \cite{example2024}. All citations will be automatically added to your \texttt{cite.bib} file.

\section{Conclusion}
This template demonstrates basic \LaTeX{} features including mathematical formulas, tables, and lists. The bibliography system is now set up to work seamlessly with Bibby AI's citation management. Feel free to modify it for your needs!

\bibliographystyle{plain}
\bibliography{cite}

\end{document}\documentclass{article}
\usepackage[utf8]{inputenc}
\usepackage{booktabs}
\usepackage{amsmath}
\usepackage{natbib}

\title{My Document}
\author{Your Name}
\date{\today}

\begin{document}

\maketitle

\section{Introduction}
This is a sample document demonstrating various \LaTeX{} features including mathematical equations, tables, and formatting. You can cite papers using the citation search feature, and they will be automatically added to your \texttt{cite.bib} file.

\section{Mathematical Expressions}
Here are some example equations:

\subsection{Quadratic Formula}
The quadratic formula is given by:
\[x = \frac{-b \pm \sqrt{b^2 - 4ac}}{2a}\]

\subsection{Matrix Example}
A sample matrix:
\[
\begin{pmatrix}
a_{11} & a_{12} & a_{13} \\
a_{21} & a_{22} & a_{23} \\
a_{31} & a_{32} & a_{33}
\end{pmatrix}
\]

\section{Tables}
Below is a sample table showing data:

\begin{table}[h]
\centering
\begin{tabular}{@{}lrrr@{}}
\toprule
\textbf{Category} & \textbf{Value 1} & \textbf{Value 2} & \textbf{Total} \\
\midrule
Item A & 10.5 & 20.3 & 30.8 \\
Item B & 15.2 & 25.1 & 40.3 \\
Item C & 22.1 & 12.7 & 34.8 \\
\bottomrule
\end{tabular}
\caption{Sample Data Table}
\label{tab:sample}
\end{table}

\section{Lists}
\subsection{Bullet Points}
\begin{itemize}
  \item First item
  \item Second item
  \item Third item with sub-items:
    \begin{itemize}
      \item Sub-item 1
      \item Sub-item 2
    \end{itemize}
\end{itemize}

\subsection{Numbered List}
\begin{enumerate}
  \item First step
  \item Second step
  \item Third step
\end{enumerate}

\section{Citations and References}
When you use the citation search feature, references will be automatically managed. For example, you might cite a paper like this: \cite{example2024}. All citations will be automatically added to your \texttt{cite.bib} file.

\section{Conclusion}
This template demonstrates basic \LaTeX{} features including mathematical formulas, tables, and lists. The bibliography system is now set up to work seamlessly with Bibby AI's citation management. Feel free to modify it for your needs!

\bibliographystyle{plain}
\bibliography{cite}

\end{document}\documentclass{article}
\usepackage[utf8]{inputenc}
\usepackage{booktabs}
\usepackage{amsmath}
\usepackage{natbib}

\title{My Document}
\author{Your Name}
\date{\today}

\begin{document}

\maketitle

\section{Introduction}
This is a sample document demonstrating various \LaTeX{} features including mathematical equations, tables, and formatting. You can cite papers using the citation search feature, and they will be automatically added to your \texttt{cite.bib} file.

\section{Mathematical Expressions}
Here are some example equations:

\subsection{Quadratic Formula}
The quadratic formula is given by:
\[x = \frac{-b \pm \sqrt{b^2 - 4ac}}{2a}\]

\subsection{Matrix Example}
A sample matrix:
\[
\begin{pmatrix}
a_{11} & a_{12} & a_{13} \\
a_{21} & a_{22} & a_{23} \\
a_{31} & a_{32} & a_{33}
\end{pmatrix}
\]

\section{Tables}
Below is a sample table showing data:

\begin{table}[h]
\centering
\begin{tabular}{@{}lrrr@{}}
\toprule
\textbf{Category} & \textbf{Value 1} & \textbf{Value 2} & \textbf{Total} \\
\midrule
Item A & 10.5 & 20.3 & 30.8 \\
Item B & 15.2 & 25.1 & 40.3 \\
Item C & 22.1 & 12.7 & 34.8 \\
\bottomrule
\end{tabular}
\caption{Sample Data Table}
\label{tab:sample}
\end{table}

\section{Lists}
\subsection{Bullet Points}
\begin{itemize}
  \item First item
  \item Second item
  \item Third item with sub-items:
    \begin{itemize}
      \item Sub-item 1
      \item Sub-item 2
    \end{itemize}
\end{itemize}

\subsection{Numbered List}
\begin{enumerate}
  \item First step
  \item Second step
  \item Third step
\end{enumerate}

\section{Citations and References}
When you use the citation search feature, references will be automatically managed. For example, you might cite a paper like this: \cite{example2024}. All citations will be automatically added to your \texttt{cite.bib} file.

\section{Conclusion}
This template demonstrates basic \LaTeX{} features including mathematical formulas, tables, and lists. The bibliography system is now set up to work seamlessly with Bibby AI's citation management. Feel free to modify it for your needs!

\bibliographystyle{plain}
\bibliography{cite}

\end{document}\documentclass{article}
\usepackage[utf8]{inputenc}
\usepackage{booktabs}
\usepackage{amsmath}
\usepackage{natbib}

\title{My Document}
\author{Your Name}
\date{\today}

\begin{document}

\maketitle

\section{Introduction}
This is a sample document demonstrating various \LaTeX{} features including mathematical equations, tables, and formatting. You can cite papers using the citation search feature, and they will be automatically added to your \texttt{cite.bib} file.

\section{Mathematical Expressions}
Here are some example equations:

\subsection{Quadratic Formula}
The quadratic formula is given by:
\[x = \frac{-b \pm \sqrt{b^2 - 4ac}}{2a}\]

\subsection{Matrix Example}
A sample matrix:
\[
\begin{pmatrix}
a_{11} & a_{12} & a_{13} \\
a_{21} & a_{22} & a_{23} \\
a_{31} & a_{32} & a_{33}
\end{pmatrix}
\]

\section{Tables}
Below is a sample table showing data:

\begin{table}[h]
\centering
\begin{tabular}{@{}lrrr@{}}
\toprule
\textbf{Category} & \textbf{Value 1} & \textbf{Value 2} & \textbf{Total} \\
\midrule
Item A & 10.5 & 20.3 & 30.8 \\
Item B & 15.2 & 25.1 & 40.3 \\
Item C & 22.1 & 12.7 & 34.8 \\
\bottomrule
\end{tabular}
\caption{Sample Data Table}
\label{tab:sample}
\end{table}

\section{Lists}
\subsection{Bullet Points}
\begin{itemize}
  \item First item
  \item Second item
  \item Third item with sub-items:
    \begin{itemize}
      \item Sub-item 1
      \item Sub-item 2
    \end{itemize}
\end{itemize}

\subsection{Numbered List}
\begin{enumerate}
  \item First step
  \item Second step
  \item Third step
\end{enumerate}

\section{Citations and References}
When you use the citation search feature, references will be automatically managed. For example, you might cite a paper like this: \cite{example2024}. All citations will be automatically added to your \texttt{cite.bib} file.

\section{Conclusion}
This template demonstrates basic \LaTeX{} features including mathematical formulas, tables, and lists. The bibliography system is now set up to work seamlessly with Bibby AI's citation management. Feel free to modify it for your needs!

\bibliographystyle{plain}
\bibliography{cite}

\end{document}\documentclass{article}
\usepackage[utf8]{inputenc}
\usepackage{booktabs}
\usepackage{amsmath}
\usepackage{natbib}

\title{My Document}
\author{Your Name}
\date{\today}

\begin{document}

\maketitle

\section{Introduction}
This is a sample document demonstrating various \LaTeX{} features including mathematical equations, tables, and formatting. You can cite papers using the citation search feature, and they will be automatically added to your \texttt{cite.bib} file.

\section{Mathematical Expressions}
Here are some example equations:

\subsection{Quadratic Formula}
The quadratic formula is given by:
\[x = \frac{-b \pm \sqrt{b^2 - 4ac}}{2a}\]

\subsection{Matrix Example}
A sample matrix:
\[
\begin{pmatrix}
a_{11} & a_{12} & a_{13} \\
a_{21} & a_{22} & a_{23} \\
a_{31} & a_{32} & a_{33}
\end{pmatrix}
\]

\section{Tables}
Below is a sample table showing data:

\begin{table}[h]
\centering
\begin{tabular}{@{}lrrr@{}}
\toprule
\textbf{Category} & \textbf{Value 1} & \textbf{Value 2} & \textbf{Total} \\
\midrule
Item A & 10.5 & 20.3 & 30.8 \\
Item B & 15.2 & 25.1 & 40.3 \\
Item C & 22.1 & 12.7 & 34.8 \\
\bottomrule
\end{tabular}
\caption{Sample Data Table}
\label{tab:sample}
\end{table}

\section{Lists}
\subsection{Bullet Points}
\begin{itemize}
  \item First item
  \item Second item
  \item Third item with sub-items:
    \begin{itemize}
      \item Sub-item 1
      \item Sub-item 2
    \end{itemize}
\end{itemize}

\subsection{Numbered List}
\begin{enumerate}
  \item First step
  \item Second step
  \item Third step
\end{enumerate}

\section{Citations and References}
When you use the citation search feature, references will be automatically managed. For example, you might cite a paper like this: \cite{example2024}. All citations will be automatically added to your \texttt{cite.bib} file.

\section{Conclusion}
This template demonstrates basic \LaTeX{} features including mathematical formulas, tables, and lists. The bibliography system is now set up to work seamlessly with Bibby AI's citation management. Feel free to modify it for your needs!

\bibliographystyle{plain}
\bibliography{cite}

\end{document}\documentclass{article}
\usepackage[utf8]{inputenc}
\usepackage{booktabs}
\usepackage{amsmath}
\usepackage{natbib}

\title{My Document}
\author{Your Name}
\date{\today}

\begin{document}

\maketitle

\section{Introduction}
This is a sample document demonstrating various \LaTeX{} features including mathematical equations, tables, and formatting. You can cite papers using the citation search feature, and they will be automatically added to your \texttt{cite.bib} file.

\section{Mathematical Expressions}
Here are some example equations:

\subsection{Quadratic Formula}
The quadratic formula is given by:
\[x = \frac{-b \pm \sqrt{b^2 - 4ac}}{2a}\]

\subsection{Matrix Example}
A sample matrix:
\[
\begin{pmatrix}
a_{11} & a_{12} & a_{13} \\
a_{21} & a_{22} & a_{23} \\
a_{31} & a_{32} & a_{33}
\end{pmatrix}
\]

\section{Tables}
Below is a sample table showing data:

\begin{table}[h]
\centering
\begin{tabular}{@{}lrrr@{}}
\toprule
\textbf{Category} & \textbf{Value 1} & \textbf{Value 2} & \textbf{Total} \\
\midrule
Item A & 10.5 & 20.3 & 30.8 \\
Item B & 15.2 & 25.1 & 40.3 \\
Item C & 22.1 & 12.7 & 34.8 \\
\bottomrule
\end{tabular}
\caption{Sample Data Table}
\label{tab:sample}
\end{table}

\section{Lists}
\subsection{Bullet Points}
\begin{itemize}
  \item First item
  \item Second item
  \item Third item with sub-items:
    \begin{itemize}
      \item Sub-item 1
      \item Sub-item 2
    \end{itemize}
\end{itemize}

\subsection{Numbered List}
\begin{enumerate}
  \item First step
  \item Second step
  \item Third step
\end{enumerate}

\section{Citations and References}
When you use the citation search feature, references will be automatically managed. For example, you might cite a paper like this: \cite{example2024}. All citations will be automatically added to your \texttt{cite.bib} file.

\section{Conclusion}
This template demonstrates basic \LaTeX{} features including mathematical formulas, tables, and lists. The bibliography system is now set up to work seamlessly with Bibby AI's citation management. Feel free to modify it for your needs!

\bibliographystyle{plain}
\bibliography{cite}

\end{document}\documentclass{article}
\usepackage[utf8]{inputenc}
\usepackage{booktabs}
\usepackage{amsmath}
\usepackage{natbib}

\title{My Document}
\author{Your Name}
\date{\today}

\begin{document}

\maketitle

\section{Introduction}
This is a sample document demonstrating various \LaTeX{} features including mathematical equations, tables, and formatting. You can cite papers using the citation search feature, and they will be automatically added to your \texttt{cite.bib} file.

\section{Mathematical Expressions}
Here are some example equations:

\subsection{Quadratic Formula}
The quadratic formula is given by:
\[x = \frac{-b \pm \sqrt{b^2 - 4ac}}{2a}\]

\subsection{Matrix Example}
A sample matrix:
\[
\begin{pmatrix}
a_{11} & a_{12} & a_{13} \\
a_{21} & a_{22} & a_{23} \\
a_{31} & a_{32} & a_{33}
\end{pmatrix}
\]

\section{Tables}
Below is a sample table showing data:

\begin{table}[h]
\centering
\begin{tabular}{@{}lrrr@{}}
\toprule
\textbf{Category} & \textbf{Value 1} & \textbf{Value 2} & \textbf{Total} \\
\midrule
Item A & 10.5 & 20.3 & 30.8 \\
Item B & 15.2 & 25.1 & 40.3 \\
Item C & 22.1 & 12.7 & 34.8 \\
\bottomrule
\end{tabular}
\caption{Sample Data Table}
\label{tab:sample}
\end{table}

\section{Lists}
\subsection{Bullet Points}
\begin{itemize}
  \item First item
  \item Second item
  \item Third item with sub-items:
    \begin{itemize}
      \item Sub-item 1
      \item Sub-item 2
    \end{itemize}
\end{itemize}

\subsection{Numbered List}
\begin{enumerate}
  \item First step
  \item Second step
  \item Third step
\end{enumerate}

\section{Citations and References}
When you use the citation search feature, references will be automatically managed. For example, you might cite a paper like this: \cite{example2024}. All citations will be automatically added to your \texttt{cite.bib} file.

\section{Conclusion}
This template demonstrates basic \LaTeX{} features including mathematical formulas, tables, and lists. The bibliography system is now set up to work seamlessly with Bibby AI's citation management. Feel free to modify it for your needs!

\bibliographystyle{plain}
\bibliography{cite}

\end{document}\documentclass{article}
\usepackage[utf8]{inputenc}
\usepackage{booktabs}
\usepackage{amsmath}
\usepackage{natbib}

\title{My Document}
\author{Your Name}
\date{\today}

\begin{document}

\maketitle

\section{Introduction}
This is a sample document demonstrating various \LaTeX{} features including mathematical equations, tables, and formatting. You can cite papers using the citation search feature, and they will be automatically added to your \texttt{cite.bib} file.

\section{Mathematical Expressions}
Here are some example equations:

\subsection{Quadratic Formula}
The quadratic formula is given by:
\[x = \frac{-b \pm \sqrt{b^2 - 4ac}}{2a}\]

\subsection{Matrix Example}
A sample matrix:
\[
\begin{pmatrix}
a_{11} & a_{12} & a_{13} \\
a_{21} & a_{22} & a_{23} \\
a_{31} & a_{32} & a_{33}
\end{pmatrix}
\]

\section{Tables}
Below is a sample table showing data:

\begin{table}[h]
\centering
\begin{tabular}{@{}lrrr@{}}
\toprule
\textbf{Category} & \textbf{Value 1} & \textbf{Value 2} & \textbf{Total} \\
\midrule
Item A & 10.5 & 20.3 & 30.8 \\
Item B & 15.2 & 25.1 & 40.3 \\
Item C & 22.1 & 12.7 & 34.8 \\
\bottomrule
\end{tabular}
\caption{Sample Data Table}
\label{tab:sample}
\end{table}

\section{Lists}
\subsection{Bullet Points}
\begin{itemize}
  \item First item
  \item Second item
  \item Third item with sub-items:
    \begin{itemize}
      \item Sub-item 1
      \item Sub-item 2
    \end{itemize}
\end{itemize}

\subsection{Numbered List}
\begin{enumerate}
  \item First step
  \item Second step
  \item Third step
\end{enumerate}

\section{Citations and References}
When you use the citation search feature, references will be automatically managed. For example, you might cite a paper like this: \cite{example2024}. All citations will be automatically added to your \texttt{cite.bib} file.

\section{Conclusion}
This template demonstrates basic \LaTeX{} features including mathematical formulas, tables, and lists. The bibliography system is now set up to work seamlessly with Bibby AI's citation management. Feel free to modify it for your needs!

\bibliographystyle{plain}
\bibliography{cite}

\end{document}\documentclass{article}
\usepackage[utf8]{inputenc}
\usepackage{booktabs}
\usepackage{amsmath}
\usepackage{natbib}

\title{My Document}
\author{Your Name}
\date{\today}

\begin{document}

\maketitle

\section{Introduction}
This is a sample document demonstrating various \LaTeX{} features including mathematical equations, tables, and formatting. You can cite papers using the citation search feature, and they will be automatically added to your \texttt{cite.bib} file.

\section{Mathematical Expressions}
Here are some example equations:

\subsection{Quadratic Formula}
The quadratic formula is given by:
\[x = \frac{-b \pm \sqrt{b^2 - 4ac}}{2a}\]

\subsection{Matrix Example}
A sample matrix:
\[
\begin{pmatrix}
a_{11} & a_{12} & a_{13} \\
a_{21} & a_{22} & a_{23} \\
a_{31} & a_{32} & a_{33}
\end{pmatrix}
\]

\section{Tables}
Below is a sample table showing data:

\begin{table}[h]
\centering
\begin{tabular}{@{}lrrr@{}}
\toprule
\textbf{Category} & \textbf{Value 1} & \textbf{Value 2} & \textbf{Total} \\
\midrule
Item A & 10.5 & 20.3 & 30.8 \\
Item B & 15.2 & 25.1 & 40.3 \\
Item C & 22.1 & 12.7 & 34.8 \\
\bottomrule
\end{tabular}
\caption{Sample Data Table}
\label{tab:sample}
\end{table}

\section{Lists}
\subsection{Bullet Points}
\begin{itemize}
  \item First item
  \item Second item
  \item Third item with sub-items:
    \begin{itemize}
      \item Sub-item 1
      \item Sub-item 2
    \end{itemize}
\end{itemize}

\subsection{Numbered List}
\begin{enumerate}
  \item First step
  \item Second step
  \item Third step
\end{enumerate}

\section{Citations and References}
When you use the citation search feature, references will be automatically managed. For example, you might cite a paper like this: \cite{example2024}. All citations will be automatically added to your \texttt{cite.bib} file.

\section{Conclusion}
This template demonstrates basic \LaTeX{} features including mathematical formulas, tables, and lists. The bibliography system is now set up to work seamlessly with Bibby AI's citation management. Feel free to modify it for your needs!

\bibliographystyle{plain}
\bibliography{cite}

\end{document}\documentclass{article}
\usepackage[utf8]{inputenc}
\usepackage{booktabs}
\usepackage{amsmath}
\usepackage{natbib}

\title{My Document}
\author{Your Name}
\date{\today}

\begin{document}

\maketitle

\section{Introduction}
This is a sample document demonstrating various \LaTeX{} features including mathematical equations, tables, and formatting. You can cite papers using the citation search feature, and they will be automatically added to your \texttt{cite.bib} file.

\section{Mathematical Expressions}
Here are some example equations:

\subsection{Quadratic Formula}
The quadratic formula is given by:
\[x = \frac{-b \pm \sqrt{b^2 - 4ac}}{2a}\]

\subsection{Matrix Example}
A sample matrix:
\[
\begin{pmatrix}
a_{11} & a_{12} & a_{13} \\
a_{21} & a_{22} & a_{23} \\
a_{31} & a_{32} & a_{33}
\end{pmatrix}
\]

\section{Tables}
Below is a sample table showing data:

\begin{table}[h]
\centering
\begin{tabular}{@{}lrrr@{}}
\toprule
\textbf{Category} & \textbf{Value 1} & \textbf{Value 2} & \textbf{Total} \\
\midrule
Item A & 10.5 & 20.3 & 30.8 \\
Item B & 15.2 & 25.1 & 40.3 \\
Item C & 22.1 & 12.7 & 34.8 \\
\bottomrule
\end{tabular}
\caption{Sample Data Table}
\label{tab:sample}
\end{table}

\section{Lists}
\subsection{Bullet Points}
\begin{itemize}
  \item First item
  \item Second item
  \item Third item with sub-items:
    \begin{itemize}
      \item Sub-item 1
      \item Sub-item 2
    \end{itemize}
\end{itemize}

\subsection{Numbered List}
\begin{enumerate}
  \item First step
  \item Second step
  \item Third step
\end{enumerate}

\section{Citations and References}
When you use the citation search feature, references will be automatically managed. For example, you might cite a paper like this: \cite{example2024}. All citations will be automatically added to your \texttt{cite.bib} file.

\section{Conclusion}
This template demonstrates basic \LaTeX{} features including mathematical formulas, tables, and lists. The bibliography system is now set up to work seamlessly with Bibby AI's citation management. Feel free to modify it for your needs!

\bibliographystyle{plain}
\bibliography{cite}

\end{document}\documentclass{article}
\usepackage[utf8]{inputenc}
\usepackage{booktabs}
\usepackage{amsmath}
\usepackage{natbib}

\title{My Document}
\author{Your Name}
\date{\today}

\begin{document}

\maketitle

\section{Introduction}
This is a sample document demonstrating various \LaTeX{} features including mathematical equations, tables, and formatting. You can cite papers using the citation search feature, and they will be automatically added to your \texttt{cite.bib} file.

\section{Mathematical Expressions}
Here are some example equations:

\subsection{Quadratic Formula}
The quadratic formula is given by:
\[x = \frac{-b \pm \sqrt{b^2 - 4ac}}{2a}\]

\subsection{Matrix Example}
A sample matrix:
\[
\begin{pmatrix}
a_{11} & a_{12} & a_{13} \\
a_{21} & a_{22} & a_{23} \\
a_{31} & a_{32} & a_{33}
\end{pmatrix}
\]

\section{Tables}
Below is a sample table showing data:

\begin{table}[h]
\centering
\begin{tabular}{@{}lrrr@{}}
\toprule
\textbf{Category} & \textbf{Value 1} & \textbf{Value 2} & \textbf{Total} \\
\midrule
Item A & 10.5 & 20.3 & 30.8 \\
Item B & 15.2 & 25.1 & 40.3 \\
Item C & 22.1 & 12.7 & 34.8 \\
\bottomrule
\end{tabular}
\caption{Sample Data Table}
\label{tab:sample}
\end{table}

\section{Lists}
\subsection{Bullet Points}
\begin{itemize}
  \item First item
  \item Second item
  \item Third item with sub-items:
    \begin{itemize}
      \item Sub-item 1
      \item Sub-item 2
    \end{itemize}
\end{itemize}

\subsection{Numbered List}
\begin{enumerate}
  \item First step
  \item Second step
  \item Third step
\end{enumerate}

\section{Citations and References}
When you use the citation search feature, references will be automatically managed. For example, you might cite a paper like this: \cite{example2024}. All citations will be automatically added to your \texttt{cite.bib} file.

\section{Conclusion}
This template demonstrates basic \LaTeX{} features including mathematical formulas, tables, and lists. The bibliography system is now set up to work seamlessly with Bibby AI's citation management. Feel free to modify it for your needs!

\bibliographystyle{plain}
\bibliography{cite}

\end{document}\documentclass{article}
\usepackage[utf8]{inputenc}
\usepackage{booktabs}
\usepackage{amsmath}
\usepackage{natbib}

\title{My Document}
\author{Your Name}
\date{\today}

\begin{document}

\maketitle

\section{Introduction}
This is a sample document demonstrating various \LaTeX{} features including mathematical equations, tables, and formatting. You can cite papers using the citation search feature, and they will be automatically added to your \texttt{cite.bib} file.

\section{Mathematical Expressions}
Here are some example equations:

\subsection{Quadratic Formula}
The quadratic formula is given by:
\[x = \frac{-b \pm \sqrt{b^2 - 4ac}}{2a}\]

\subsection{Matrix Example}
A sample matrix:
\[
\begin{pmatrix}
a_{11} & a_{12} & a_{13} \\
a_{21} & a_{22} & a_{23} \\
a_{31} & a_{32} & a_{33}
\end{pmatrix}
\]

\section{Tables}
Below is a sample table showing data:

\begin{table}[h]
\centering
\begin{tabular}{@{}lrrr@{}}
\toprule
\textbf{Category} & \textbf{Value 1} & \textbf{Value 2} & \textbf{Total} \\
\midrule
Item A & 10.5 & 20.3 & 30.8 \\
Item B & 15.2 & 25.1 & 40.3 \\
Item C & 22.1 & 12.7 & 34.8 \\
\bottomrule
\end{tabular}
\caption{Sample Data Table}
\label{tab:sample}
\end{table}

\section{Lists}
\subsection{Bullet Points}
\begin{itemize}
  \item First item
  \item Second item
  \item Third item with sub-items:
    \begin{itemize}
      \item Sub-item 1
      \item Sub-item 2
    \end{itemize}
\end{itemize}

\subsection{Numbered List}
\begin{enumerate}
  \item First step
  \item Second step
  \item Third step
\end{enumerate}

\section{Citations and References}
When you use the citation search feature, references will be automatically managed. For example, you might cite a paper like this: \cite{example2024}. All citations will be automatically added to your \texttt{cite.bib} file.

\section{Conclusion}
This template demonstrates basic \LaTeX{} features including mathematical formulas, tables, and lists. The bibliography system is now set up to work seamlessly with Bibby AI's citation management. Feel free to modify it for your needs!

\bibliographystyle{plain}
\bibliography{cite}

\end{document}\documentclass{article}
\usepackage[utf8]{inputenc}
\usepackage{booktabs}
\usepackage{amsmath}
\usepackage{natbib}

\title{My Document}
\author{Your Name}
\date{\today}

\begin{document}

\maketitle

\section{Introduction}
This is a sample document demonstrating various \LaTeX{} features including mathematical equations, tables, and formatting. You can cite papers using the citation search feature, and they will be automatically added to your \texttt{cite.bib} file.

\section{Mathematical Expressions}
Here are some example equations:

\subsection{Quadratic Formula}
The quadratic formula is given by:
\[x = \frac{-b \pm \sqrt{b^2 - 4ac}}{2a}\]

\subsection{Matrix Example}
A sample matrix:
\[
\begin{pmatrix}
a_{11} & a_{12} & a_{13} \\
a_{21} & a_{22} & a_{23} \\
a_{31} & a_{32} & a_{33}
\end{pmatrix}
\]

\section{Tables}
Below is a sample table showing data:

\begin{table}[h]
\centering
\begin{tabular}{@{}lrrr@{}}
\toprule
\textbf{Category} & \textbf{Value 1} & \textbf{Value 2} & \textbf{Total} \\
\midrule
Item A & 10.5 & 20.3 & 30.8 \\
Item B & 15.2 & 25.1 & 40.3 \\
Item C & 22.1 & 12.7 & 34.8 \\
\bottomrule
\end{tabular}
\caption{Sample Data Table}
\label{tab:sample}
\end{table}

\section{Lists}
\subsection{Bullet Points}
\begin{itemize}
  \item First item
  \item Second item
  \item Third item with sub-items:
    \begin{itemize}
      \item Sub-item 1
      \item Sub-item 2
    \end{itemize}
\end{itemize}

\subsection{Numbered List}
\begin{enumerate}
  \item First step
  \item Second step
  \item Third step
\end{enumerate}

\section{Citations and References}
When you use the citation search feature, references will be automatically managed. For example, you might cite a paper like this: \cite{example2024}. All citations will be automatically added to your \texttt{cite.bib} file.

\section{Conclusion}
This template demonstrates basic \LaTeX{} features including mathematical formulas, tables, and lists. The bibliography system is now set up to work seamlessly with Bibby AI's citation management. Feel free to modify it for your needs!

\bibliographystyle{plain}
\bibliography{cite}

\end{document}\documentclass{article}
\usepackage[utf8]{inputenc}
\usepackage{booktabs}
\usepackage{amsmath}
\usepackage{natbib}

\title{My Document}
\author{Your Name}
\date{\today}

\begin{document}

\maketitle

\section{Introduction}
This is a sample document demonstrating various \LaTeX{} features including mathematical equations, tables, and formatting. You can cite papers using the citation search feature, and they will be automatically added to your \texttt{cite.bib} file.

\section{Mathematical Expressions}
Here are some example equations:

\subsection{Quadratic Formula}
The quadratic formula is given by:
\[x = \frac{-b \pm \sqrt{b^2 - 4ac}}{2a}\]

\subsection{Matrix Example}
A sample matrix:
\[
\begin{pmatrix}
a_{11} & a_{12} & a_{13} \\
a_{21} & a_{22} & a_{23} \\
a_{31} & a_{32} & a_{33}
\end{pmatrix}
\]

\section{Tables}
Below is a sample table showing data:

\begin{table}[h]
\centering
\begin{tabular}{@{}lrrr@{}}
\toprule
\textbf{Category} & \textbf{Value 1} & \textbf{Value 2} & \textbf{Total} \\
\midrule
Item A & 10.5 & 20.3 & 30.8 \\
Item B & 15.2 & 25.1 & 40.3 \\
Item C & 22.1 & 12.7 & 34.8 \\
\bottomrule
\end{tabular}
\caption{Sample Data Table}
\label{tab:sample}
\end{table}

\section{Lists}
\subsection{Bullet Points}
\begin{itemize}
  \item First item
  \item Second item
  \item Third item with sub-items:
    \begin{itemize}
      \item Sub-item 1
      \item Sub-item 2
    \end{itemize}
\end{itemize}

\subsection{Numbered List}
\begin{enumerate}
  \item First step
  \item Second step
  \item Third step
\end{enumerate}

\section{Citations and References}
When you use the citation search feature, references will be automatically managed. For example, you might cite a paper like this: \cite{example2024}. All citations will be automatically added to your \texttt{cite.bib} file.

\section{Conclusion}
This template demonstrates basic \LaTeX{} features including mathematical formulas, tables, and lists. The bibliography system is now set up to work seamlessly with Bibby AI's citation management. Feel free to modify it for your needs!

\bibliographystyle{plain}
\bibliography{cite}

\end{document}\documentclass{article}
\usepackage[utf8]{inputenc}
\usepackage{booktabs}
\usepackage{amsmath}
\usepackage{natbib}

\title{My Document}
\author{Your Name}
\date{\today}

\begin{document}

\maketitle

\section{Introduction}
This is a sample document demonstrating various \LaTeX{} features including mathematical equations, tables, and formatting. You can cite papers using the citation search feature, and they will be automatically added to your \texttt{cite.bib} file.

\section{Mathematical Expressions}
Here are some example equations:

\subsection{Quadratic Formula}
The quadratic formula is given by:
\[x = \frac{-b \pm \sqrt{b^2 - 4ac}}{2a}\]

\subsection{Matrix Example}
A sample matrix:
\[
\begin{pmatrix}
a_{11} & a_{12} & a_{13} \\
a_{21} & a_{22} & a_{23} \\
a_{31} & a_{32} & a_{33}
\end{pmatrix}
\]

\section{Tables}
Below is a sample table showing data:

\begin{table}[h]
\centering
\begin{tabular}{@{}lrrr@{}}
\toprule
\textbf{Category} & \textbf{Value 1} & \textbf{Value 2} & \textbf{Total} \\
\midrule
Item A & 10.5 & 20.3 & 30.8 \\
Item B & 15.2 & 25.1 & 40.3 \\
Item C & 22.1 & 12.7 & 34.8 \\
\bottomrule
\end{tabular}
\caption{Sample Data Table}
\label{tab:sample}
\end{table}

\section{Lists}
\subsection{Bullet Points}
\begin{itemize}
  \item First item
  \item Second item
  \item Third item with sub-items:
    \begin{itemize}
      \item Sub-item 1
      \item Sub-item 2
    \end{itemize}
\end{itemize}

\subsection{Numbered List}
\begin{enumerate}
  \item First step
  \item Second step
  \item Third step
\end{enumerate}

\section{Citations and References}
When you use the citation search feature, references will be automatically managed. For example, you might cite a paper like this: \cite{example2024}. All citations will be automatically added to your \texttt{cite.bib} file.

\section{Conclusion}
This template demonstrates basic \LaTeX{} features including mathematical formulas, tables, and lists. The bibliography system is now set up to work seamlessly with Bibby AI's citation management. Feel free to modify it for your needs!

\bibliographystyle{plain}
\bibliography{cite}

\end{document}\documentclass{article}
\usepackage[utf8]{inputenc}
\usepackage{booktabs}
\usepackage{amsmath}
\usepackage{natbib}

\title{My Document}
\author{Your Name}
\date{\today}

\begin{document}

\maketitle

\section{Introduction}
This is a sample document demonstrating various \LaTeX{} features including mathematical equations, tables, and formatting. You can cite papers using the citation search feature, and they will be automatically added to your \texttt{cite.bib} file.

\section{Mathematical Expressions}
Here are some example equations:

\subsection{Quadratic Formula}
The quadratic formula is given by:
\[x = \frac{-b \pm \sqrt{b^2 - 4ac}}{2a}\]

\subsection{Matrix Example}
A sample matrix:
\[
\begin{pmatrix}
a_{11} & a_{12} & a_{13} \\
a_{21} & a_{22} & a_{23} \\
a_{31} & a_{32} & a_{33}
\end{pmatrix}
\]

\section{Tables}
Below is a sample table showing data:

\begin{table}[h]
\centering
\begin{tabular}{@{}lrrr@{}}
\toprule
\textbf{Category} & \textbf{Value 1} & \textbf{Value 2} & \textbf{Total} \\
\midrule
Item A & 10.5 & 20.3 & 30.8 \\
Item B & 15.2 & 25.1 & 40.3 \\
Item C & 22.1 & 12.7 & 34.8 \\
\bottomrule
\end{tabular}
\caption{Sample Data Table}
\label{tab:sample}
\end{table}

\section{Lists}
\subsection{Bullet Points}
\begin{itemize}
  \item First item
  \item Second item
  \item Third item with sub-items:
    \begin{itemize}
      \item Sub-item 1
      \item Sub-item 2
    \end{itemize}
\end{itemize}

\subsection{Numbered List}
\begin{enumerate}
  \item First step
  \item Second step
  \item Third step
\end{enumerate}

\section{Citations and References}
When you use the citation search feature, references will be automatically managed. For example, you might cite a paper like this: \cite{example2024}. All citations will be automatically added to your \texttt{cite.bib} file.

\section{Conclusion}
This template demonstrates basic \LaTeX{} features including mathematical formulas, tables, and lists. The bibliography system is now set up to work seamlessly with Bibby AI's citation management. Feel free to modify it for your needs!

\bibliographystyle{plain}
\bibliography{cite}

\end{document}\documentclass{article}
\usepackage[utf8]{inputenc}
\usepackage{booktabs}
\usepackage{amsmath}
\usepackage{natbib}

\title{My Document}
\author{Your Name}
\date{\today}

\begin{document}

\maketitle

\section{Introduction}
This is a sample document demonstrating various \LaTeX{} features including mathematical equations, tables, and formatting. You can cite papers using the citation search feature, and they will be automatically added to your \texttt{cite.bib} file.

\section{Mathematical Expressions}
Here are some example equations:

\subsection{Quadratic Formula}
The quadratic formula is given by:
\[x = \frac{-b \pm \sqrt{b^2 - 4ac}}{2a}\]

\subsection{Matrix Example}
A sample matrix:
\[
\begin{pmatrix}
a_{11} & a_{12} & a_{13} \\
a_{21} & a_{22} & a_{23} \\
a_{31} & a_{32} & a_{33}
\end{pmatrix}
\]

\section{Tables}
Below is a sample table showing data:

\begin{table}[h]
\centering
\begin{tabular}{@{}lrrr@{}}
\toprule
\textbf{Category} & \textbf{Value 1} & \textbf{Value 2} & \textbf{Total} \\
\midrule
Item A & 10.5 & 20.3 & 30.8 \\
Item B & 15.2 & 25.1 & 40.3 \\
Item C & 22.1 & 12.7 & 34.8 \\
\bottomrule
\end{tabular}
\caption{Sample Data Table}
\label{tab:sample}
\end{table}

\section{Lists}
\subsection{Bullet Points}
\begin{itemize}
  \item First item
  \item Second item
  \item Third item with sub-items:
    \begin{itemize}
      \item Sub-item 1
      \item Sub-item 2
    \end{itemize}
\end{itemize}

\subsection{Numbered List}
\begin{enumerate}
  \item First step
  \item Second step
  \item Third step
\end{enumerate}

\section{Citations and References}
When you use the citation search feature, references will be automatically managed. For example, you might cite a paper like this: \cite{example2024}. All citations will be automatically added to your \texttt{cite.bib} file.

\section{Conclusion}
This template demonstrates basic \LaTeX{} features including mathematical formulas, tables, and lists. The bibliography system is now set up to work seamlessly with Bibby AI's citation management. Feel free to modify it for your needs!

\bibliographystyle{plain}
\bibliography{cite}

\end{document}\documentclass{article}
\usepackage[utf8]{inputenc}
\usepackage{booktabs}
\usepackage{amsmath}
\usepackage{natbib}

\title{My Document}
\author{Your Name}
\date{\today}

\begin{document}

\maketitle

\section{Introduction}
This is a sample document demonstrating various \LaTeX{} features including mathematical equations, tables, and formatting. You can cite papers using the citation search feature, and they will be automatically added to your \texttt{cite.bib} file.

\section{Mathematical Expressions}
Here are some example equations:

\subsection{Quadratic Formula}
The quadratic formula is given by:
\[x = \frac{-b \pm \sqrt{b^2 - 4ac}}{2a}\]

\subsection{Matrix Example}
A sample matrix:
\[
\begin{pmatrix}
a_{11} & a_{12} & a_{13} \\
a_{21} & a_{22} & a_{23} \\
a_{31} & a_{32} & a_{33}
\end{pmatrix}
\]

\section{Tables}
Below is a sample table showing data:

\begin{table}[h]
\centering
\begin{tabular}{@{}lrrr@{}}
\toprule
\textbf{Category} & \textbf{Value 1} & \textbf{Value 2} & \textbf{Total} \\
\midrule
Item A & 10.5 & 20.3 & 30.8 \\
Item B & 15.2 & 25.1 & 40.3 \\
Item C & 22.1 & 12.7 & 34.8 \\
\bottomrule
\end{tabular}
\caption{Sample Data Table}
\label{tab:sample}
\end{table}

\section{Lists}
\subsection{Bullet Points}
\begin{itemize}
  \item First item
  \item Second item
  \item Third item with sub-items:
    \begin{itemize}
      \item Sub-item 1
      \item Sub-item 2
    \end{itemize}
\end{itemize}

\subsection{Numbered List}
\begin{enumerate}
  \item First step
  \item Second step
  \item Third step
\end{enumerate}

\section{Citations and References}
When you use the citation search feature, references will be automatically managed. For example, you might cite a paper like this: \cite{example2024}. All citations will be automatically added to your \texttt{cite.bib} file.

\section{Conclusion}
This template demonstrates basic \LaTeX{} features including mathematical formulas, tables, and lists. The bibliography system is now set up to work seamlessly with Bibby AI's citation management. Feel free to modify it for your needs!

\bibliographystyle{plain}
\bibliography{cite}

\end{document}\documentclass{article}
\usepackage[utf8]{inputenc}
\usepackage{booktabs}
\usepackage{amsmath}
\usepackage{natbib}

\title{My Document}
\author{Your Name}
\date{\today}

\begin{document}

\maketitle

\section{Introduction}
This is a sample document demonstrating various \LaTeX{} features including mathematical equations, tables, and formatting. You can cite papers using the citation search feature, and they will be automatically added to your \texttt{cite.bib} file.

\section{Mathematical Expressions}
Here are some example equations:

\subsection{Quadratic Formula}
The quadratic formula is given by:
\[x = \frac{-b \pm \sqrt{b^2 - 4ac}}{2a}\]

\subsection{Matrix Example}
A sample matrix:
\[
\begin{pmatrix}
a_{11} & a_{12} & a_{13} \\
a_{21} & a_{22} & a_{23} \\
a_{31} & a_{32} & a_{33}
\end{pmatrix}
\]

\section{Tables}
Below is a sample table showing data:

\begin{table}[h]
\centering
\begin{tabular}{@{}lrrr@{}}
\toprule
\textbf{Category} & \textbf{Value 1} & \textbf{Value 2} & \textbf{Total} \\
\midrule
Item A & 10.5 & 20.3 & 30.8 \\
Item B & 15.2 & 25.1 & 40.3 \\
Item C & 22.1 & 12.7 & 34.8 \\
\bottomrule
\end{tabular}
\caption{Sample Data Table}
\label{tab:sample}
\end{table}

\section{Lists}
\subsection{Bullet Points}
\begin{itemize}
  \item First item
  \item Second item
  \item Third item with sub-items:
    \begin{itemize}
      \item Sub-item 1
      \item Sub-item 2
    \end{itemize}
\end{itemize}

\subsection{Numbered List}
\begin{enumerate}
  \item First step
  \item Second step
  \item Third step
\end{enumerate}

\section{Citations and References}
When you use the citation search feature, references will be automatically managed. For example, you might cite a paper like this: \cite{example2024}. All citations will be automatically added to your \texttt{cite.bib} file.

\section{Conclusion}
This template demonstrates basic \LaTeX{} features including mathematical formulas, tables, and lists. The bibliography system is now set up to work seamlessly with Bibby AI's citation management. Feel free to modify it for your needs!

\bibliographystyle{plain}
\bibliography{cite}

\end{document}\documentclass{article}
\usepackage[utf8]{inputenc}
\usepackage{booktabs}
\usepackage{amsmath}
\usepackage{natbib}

\title{My Document}
\author{Your Name}
\date{\today}

\begin{document}

\maketitle

\section{Introduction}
This is a sample document demonstrating various \LaTeX{} features including mathematical equations, tables, and formatting. You can cite papers using the citation search feature, and they will be automatically added to your \texttt{cite.bib} file.

\section{Mathematical Expressions}
Here are some example equations:

\subsection{Quadratic Formula}
The quadratic formula is given by:
\[x = \frac{-b \pm \sqrt{b^2 - 4ac}}{2a}\]

\subsection{Matrix Example}
A sample matrix:
\[
\begin{pmatrix}
a_{11} & a_{12} & a_{13} \\
a_{21} & a_{22} & a_{23} \\
a_{31} & a_{32} & a_{33}
\end{pmatrix}
\]

\section{Tables}
Below is a sample table showing data:

\begin{table}[h]
\centering
\begin{tabular}{@{}lrrr@{}}
\toprule
\textbf{Category} & \textbf{Value 1} & \textbf{Value 2} & \textbf{Total} \\
\midrule
Item A & 10.5 & 20.3 & 30.8 \\
Item B & 15.2 & 25.1 & 40.3 \\
Item C & 22.1 & 12.7 & 34.8 \\
\bottomrule
\end{tabular}
\caption{Sample Data Table}
\label{tab:sample}
\end{table}

\section{Lists}
\subsection{Bullet Points}
\begin{itemize}
  \item First item
  \item Second item
  \item Third item with sub-items:
    \begin{itemize}
      \item Sub-item 1
      \item Sub-item 2
    \end{itemize}
\end{itemize}

\subsection{Numbered List}
\begin{enumerate}
  \item First step
  \item Second step
  \item Third step
\end{enumerate}

\section{Citations and References}
When you use the citation search feature, references will be automatically managed. For example, you might cite a paper like this: \cite{example2024}. All citations will be automatically added to your \texttt{cite.bib} file.

\section{Conclusion}
This template demonstrates basic \LaTeX{} features including mathematical formulas, tables, and lists. The bibliography system is now set up to work seamlessly with Bibby AI's citation management. Feel free to modify it for your needs!

\bibliographystyle{plain}
\bibliography{cite}

\end{document}\documentclass{article}
\usepackage[utf8]{inputenc}
\usepackage{booktabs}
\usepackage{amsmath}
\usepackage{natbib}

\title{My Document}
\author{Your Name}
\date{\today}

\begin{document}

\maketitle

\section{Introduction}
This is a sample document demonstrating various \LaTeX{} features including mathematical equations, tables, and formatting. You can cite papers using the citation search feature, and they will be automatically added to your \texttt{cite.bib} file.

\section{Mathematical Expressions}
Here are some example equations:

\subsection{Quadratic Formula}
The quadratic formula is given by:
\[x = \frac{-b \pm \sqrt{b^2 - 4ac}}{2a}\]

\subsection{Matrix Example}
A sample matrix:
\[
\begin{pmatrix}
a_{11} & a_{12} & a_{13} \\
a_{21} & a_{22} & a_{23} \\
a_{31} & a_{32} & a_{33}
\end{pmatrix}
\]

\section{Tables}
Below is a sample table showing data:

\begin{table}[h]
\centering
\begin{tabular}{@{}lrrr@{}}
\toprule
\textbf{Category} & \textbf{Value 1} & \textbf{Value 2} & \textbf{Total} \\
\midrule
Item A & 10.5 & 20.3 & 30.8 \\
Item B & 15.2 & 25.1 & 40.3 \\
Item C & 22.1 & 12.7 & 34.8 \\
\bottomrule
\end{tabular}
\caption{Sample Data Table}
\label{tab:sample}
\end{table}

\section{Lists}
\subsection{Bullet Points}
\begin{itemize}
  \item First item
  \item Second item
  \item Third item with sub-items:
    \begin{itemize}
      \item Sub-item 1
      \item Sub-item 2
    \end{itemize}
\end{itemize}

\subsection{Numbered List}
\begin{enumerate}
  \item First step
  \item Second step
  \item Third step
\end{enumerate}

\section{Citations and References}
When you use the citation search feature, references will be automatically managed. For example, you might cite a paper like this: \cite{example2024}. All citations will be automatically added to your \texttt{cite.bib} file.

\section{Conclusion}
This template demonstrates basic \LaTeX{} features including mathematical formulas, tables, and lists. The bibliography system is now set up to work seamlessly with Bibby AI's citation management. Feel free to modify it for your needs!

\bibliographystyle{plain}
\bibliography{cite}

\end{document}\documentclass{article}
\usepackage[utf8]{inputenc}
\usepackage{booktabs}
\usepackage{amsmath}
\usepackage{natbib}

\title{My Document}
\author{Your Name}
\date{\today}

\begin{document}

\maketitle

\section{Introduction}
This is a sample document demonstrating various \LaTeX{} features including mathematical equations, tables, and formatting. You can cite papers using the citation search feature, and they will be automatically added to your \texttt{cite.bib} file.

\section{Mathematical Expressions}
Here are some example equations:

\subsection{Quadratic Formula}
The quadratic formula is given by:
\[x = \frac{-b \pm \sqrt{b^2 - 4ac}}{2a}\]

\subsection{Matrix Example}
A sample matrix:
\[
\begin{pmatrix}
a_{11} & a_{12} & a_{13} \\
a_{21} & a_{22} & a_{23} \\
a_{31} & a_{32} & a_{33}
\end{pmatrix}
\]

\section{Tables}
Below is a sample table showing data:

\begin{table}[h]
\centering
\begin{tabular}{@{}lrrr@{}}
\toprule
\textbf{Category} & \textbf{Value 1} & \textbf{Value 2} & \textbf{Total} \\
\midrule
Item A & 10.5 & 20.3 & 30.8 \\
Item B & 15.2 & 25.1 & 40.3 \\
Item C & 22.1 & 12.7 & 34.8 \\
\bottomrule
\end{tabular}
\caption{Sample Data Table}
\label{tab:sample}
\end{table}

\section{Lists}
\subsection{Bullet Points}
\begin{itemize}
  \item First item
  \item Second item
  \item Third item with sub-items:
    \begin{itemize}
      \item Sub-item 1
      \item Sub-item 2
    \end{itemize}
\end{itemize}

\subsection{Numbered List}
\begin{enumerate}
  \item First step
  \item Second step
  \item Third step
\end{enumerate}

\section{Citations and References}
When you use the citation search feature, references will be automatically managed. For example, you might cite a paper like this: \cite{example2024}. All citations will be automatically added to your \texttt{cite.bib} file.

\section{Conclusion}
This template demonstrates basic \LaTeX{} features including mathematical formulas, tables, and lists. The bibliography system is now set up to work seamlessly with Bibby AI's citation management. Feel free to modify it for your needs!

\bibliographystyle{plain}
\bibliography{cite}

\end{document}\documentclass{article}
\usepackage[utf8]{inputenc}
\usepackage{booktabs}
\usepackage{amsmath}
\usepackage{natbib}

\title{My Document}
\author{Your Name}
\date{\today}

\begin{document}

\maketitle

\section{Introduction}
This is a sample document demonstrating various \LaTeX{} features including mathematical equations, tables, and formatting. You can cite papers using the citation search feature, and they will be automatically added to your \texttt{cite.bib} file.

\section{Mathematical Expressions}
Here are some example equations:

\subsection{Quadratic Formula}
The quadratic formula is given by:
\[x = \frac{-b \pm \sqrt{b^2 - 4ac}}{2a}\]

\subsection{Matrix Example}
A sample matrix:
\[
\begin{pmatrix}
a_{11} & a_{12} & a_{13} \\
a_{21} & a_{22} & a_{23} \\
a_{31} & a_{32} & a_{33}
\end{pmatrix}
\]

\section{Tables}
Below is a sample table showing data:

\begin{table}[h]
\centering
\begin{tabular}{@{}lrrr@{}}
\toprule
\textbf{Category} & \textbf{Value 1} & \textbf{Value 2} & \textbf{Total} \\
\midrule
Item A & 10.5 & 20.3 & 30.8 \\
Item B & 15.2 & 25.1 & 40.3 \\
Item C & 22.1 & 12.7 & 34.8 \\
\bottomrule
\end{tabular}
\caption{Sample Data Table}
\label{tab:sample}
\end{table}

\section{Lists}
\subsection{Bullet Points}
\begin{itemize}
  \item First item
  \item Second item
  \item Third item with sub-items:
    \begin{itemize}
      \item Sub-item 1
      \item Sub-item 2
    \end{itemize}
\end{itemize}

\subsection{Numbered List}
\begin{enumerate}
  \item First step
  \item Second step
  \item Third step
\end{enumerate}

\section{Citations and References}
When you use the citation search feature, references will be automatically managed. For example, you might cite a paper like this: \cite{example2024}. All citations will be automatically added to your \texttt{cite.bib} file.

\section{Conclusion}
This template demonstrates basic \LaTeX{} features including mathematical formulas, tables, and lists. The bibliography system is now set up to work seamlessly with Bibby AI's citation management. Feel free to modify it for your needs!

\bibliographystyle{plain}
\bibliography{cite}

\end{document}\documentclass{article}
\usepackage[utf8]{inputenc}
\usepackage{booktabs}
\usepackage{amsmath}
\usepackage{natbib}

\title{My Document}
\author{Your Name}
\date{\today}

\begin{document}

\maketitle

\section{Introduction}
This is a sample document demonstrating various \LaTeX{} features including mathematical equations, tables, and formatting. You can cite papers using the citation search feature, and they will be automatically added to your \texttt{cite.bib} file.

\section{Mathematical Expressions}
Here are some example equations:

\subsection{Quadratic Formula}
The quadratic formula is given by:
\[x = \frac{-b \pm \sqrt{b^2 - 4ac}}{2a}\]

\subsection{Matrix Example}
A sample matrix:
\[
\begin{pmatrix}
a_{11} & a_{12} & a_{13} \\
a_{21} & a_{22} & a_{23} \\
a_{31} & a_{32} & a_{33}
\end{pmatrix}
\]

\section{Tables}
Below is a sample table showing data:

\begin{table}[h]
\centering
\begin{tabular}{@{}lrrr@{}}
\toprule
\textbf{Category} & \textbf{Value 1} & \textbf{Value 2} & \textbf{Total} \\
\midrule
Item A & 10.5 & 20.3 & 30.8 \\
Item B & 15.2 & 25.1 & 40.3 \\
Item C & 22.1 & 12.7 & 34.8 \\
\bottomrule
\end{tabular}
\caption{Sample Data Table}
\label{tab:sample}
\end{table}

\section{Lists}
\subsection{Bullet Points}
\begin{itemize}
  \item First item
  \item Second item
  \item Third item with sub-items:
    \begin{itemize}
      \item Sub-item 1
      \item Sub-item 2
    \end{itemize}
\end{itemize}

\subsection{Numbered List}
\begin{enumerate}
  \item First step
  \item Second step
  \item Third step
\end{enumerate}

\section{Citations and References}
When you use the citation search feature, references will be automatically managed. For example, you might cite a paper like this: \cite{example2024}. All citations will be automatically added to your \texttt{cite.bib} file.

\section{Conclusion}
This template demonstrates basic \LaTeX{} features including mathematical formulas, tables, and lists. The bibliography system is now set up to work seamlessly with Bibby AI's citation management. Feel free to modify it for your needs!

\bibliographystyle{plain}
\bibliography{cite}

\end{document}\documentclass{article}
\usepackage[utf8]{inputenc}
\usepackage{booktabs}
\usepackage{amsmath}
\usepackage{natbib}

\title{My Document}
\author{Your Name}
\date{\today}

\begin{document}

\maketitle

\section{Introduction}
This is a sample document demonstrating various \LaTeX{} features including mathematical equations, tables, and formatting. You can cite papers using the citation search feature, and they will be automatically added to your \texttt{cite.bib} file.

\section{Mathematical Expressions}
Here are some example equations:

\subsection{Quadratic Formula}
The quadratic formula is given by:
\[x = \frac{-b \pm \sqrt{b^2 - 4ac}}{2a}\]

\subsection{Matrix Example}
A sample matrix:
\[
\begin{pmatrix}
a_{11} & a_{12} & a_{13} \\
a_{21} & a_{22} & a_{23} \\
a_{31} & a_{32} & a_{33}
\end{pmatrix}
\]

\section{Tables}
Below is a sample table showing data:

\begin{table}[h]
\centering
\begin{tabular}{@{}lrrr@{}}
\toprule
\textbf{Category} & \textbf{Value 1} & \textbf{Value 2} & \textbf{Total} \\
\midrule
Item A & 10.5 & 20.3 & 30.8 \\
Item B & 15.2 & 25.1 & 40.3 \\
Item C & 22.1 & 12.7 & 34.8 \\
\bottomrule
\end{tabular}
\caption{Sample Data Table}
\label{tab:sample}
\end{table}

\section{Lists}
\subsection{Bullet Points}
\begin{itemize}
  \item First item
  \item Second item
  \item Third item with sub-items:
    \begin{itemize}
      \item Sub-item 1
      \item Sub-item 2
    \end{itemize}
\end{itemize}

\subsection{Numbered List}
\begin{enumerate}
  \item First step
  \item Second step
  \item Third step
\end{enumerate}

\section{Citations and References}
When you use the citation search feature, references will be automatically managed. For example, you might cite a paper like this: \cite{example2024}. All citations will be automatically added to your \texttt{cite.bib} file.

\section{Conclusion}
This template demonstrates basic \LaTeX{} features including mathematical formulas, tables, and lists. The bibliography system is now set up to work seamlessly with Bibby AI's citation management. Feel free to modify it for your needs!

\bibliographystyle{plain}
\bibliography{cite}

\end{document}\documentclass{article}
\usepackage[utf8]{inputenc}
\usepackage{booktabs}
\usepackage{amsmath}
\usepackage{natbib}

\title{My Document}
\author{Your Name}
\date{\today}

\begin{document}

\maketitle

\section{Introduction}
This is a sample document demonstrating various \LaTeX{} features including mathematical equations, tables, and formatting. You can cite papers using the citation search feature, and they will be automatically added to your \texttt{cite.bib} file.

\section{Mathematical Expressions}
Here are some example equations:

\subsection{Quadratic Formula}
The quadratic formula is given by:
\[x = \frac{-b \pm \sqrt{b^2 - 4ac}}{2a}\]

\subsection{Matrix Example}
A sample matrix:
\[
\begin{pmatrix}
a_{11} & a_{12} & a_{13} \\
a_{21} & a_{22} & a_{23} \\
a_{31} & a_{32} & a_{33}
\end{pmatrix}
\]

\section{Tables}
Below is a sample table showing data:

\begin{table}[h]
\centering
\begin{tabular}{@{}lrrr@{}}
\toprule
\textbf{Category} & \textbf{Value 1} & \textbf{Value 2} & \textbf{Total} \\
\midrule
Item A & 10.5 & 20.3 & 30.8 \\
Item B & 15.2 & 25.1 & 40.3 \\
Item C & 22.1 & 12.7 & 34.8 \\
\bottomrule
\end{tabular}
\caption{Sample Data Table}
\label{tab:sample}
\end{table}

\section{Lists}
\subsection{Bullet Points}
\begin{itemize}
  \item First item
  \item Second item
  \item Third item with sub-items:
    \begin{itemize}
      \item Sub-item 1
      \item Sub-item 2
    \end{itemize}
\end{itemize}

\subsection{Numbered List}
\begin{enumerate}
  \item First step
  \item Second step
  \item Third step
\end{enumerate}

\section{Citations and References}
When you use the citation search feature, references will be automatically managed. For example, you might cite a paper like this: \cite{example2024}. All citations will be automatically added to your \texttt{cite.bib} file.

\section{Conclusion}
This template demonstrates basic \LaTeX{} features including mathematical formulas, tables, and lists. The bibliography system is now set up to work seamlessly with Bibby AI's citation management. Feel free to modify it for your needs!

\bibliographystyle{plain}
\bibliography{cite}

\end{document}\documentclass{article}
\usepackage[utf8]{inputenc}
\usepackage{booktabs}
\usepackage{amsmath}
\usepackage{natbib}

\title{My Document}
\author{Your Name}
\date{\today}

\begin{document}

\maketitle

\section{Introduction}
This is a sample document demonstrating various \LaTeX{} features including mathematical equations, tables, and formatting. You can cite papers using the citation search feature, and they will be automatically added to your \texttt{cite.bib} file.

\section{Mathematical Expressions}
Here are some example equations:

\subsection{Quadratic Formula}
The quadratic formula is given by:
\[x = \frac{-b \pm \sqrt{b^2 - 4ac}}{2a}\]

\subsection{Matrix Example}
A sample matrix:
\[
\begin{pmatrix}
a_{11} & a_{12} & a_{13} \\
a_{21} & a_{22} & a_{23} \\
a_{31} & a_{32} & a_{33}
\end{pmatrix}
\]

\section{Tables}
Below is a sample table showing data:

\begin{table}[h]
\centering
\begin{tabular}{@{}lrrr@{}}
\toprule
\textbf{Category} & \textbf{Value 1} & \textbf{Value 2} & \textbf{Total} \\
\midrule
Item A & 10.5 & 20.3 & 30.8 \\
Item B & 15.2 & 25.1 & 40.3 \\
Item C & 22.1 & 12.7 & 34.8 \\
\bottomrule
\end{tabular}
\caption{Sample Data Table}
\label{tab:sample}
\end{table}

\section{Lists}
\subsection{Bullet Points}
\begin{itemize}
  \item First item
  \item Second item
  \item Third item with sub-items:
    \begin{itemize}
      \item Sub-item 1
      \item Sub-item 2
    \end{itemize}
\end{itemize}

\subsection{Numbered List}
\begin{enumerate}
  \item First step
  \item Second step
  \item Third step
\end{enumerate}

\section{Citations and References}
When you use the citation search feature, references will be automatically managed. For example, you might cite a paper like this: \cite{example2024}. All citations will be automatically added to your \texttt{cite.bib} file.

\section{Conclusion}
This template demonstrates basic \LaTeX{} features including mathematical formulas, tables, and lists. The bibliography system is now set up to work seamlessly with Bibby AI's citation management. Feel free to modify it for your needs!

\bibliographystyle{plain}
\bibliography{cite}

\end{document}\documentclass{article}
\usepackage[utf8]{inputenc}
\usepackage{booktabs}
\usepackage{amsmath}
\usepackage{natbib}

\title{My Document}
\author{Your Name}
\date{\today}

\begin{document}

\maketitle

\section{Introduction}
This is a sample document demonstrating various \LaTeX{} features including mathematical equations, tables, and formatting. You can cite papers using the citation search feature, and they will be automatically added to your \texttt{cite.bib} file.

\section{Mathematical Expressions}
Here are some example equations:

\subsection{Quadratic Formula}
The quadratic formula is given by:
\[x = \frac{-b \pm \sqrt{b^2 - 4ac}}{2a}\]

\subsection{Matrix Example}
A sample matrix:
\[
\begin{pmatrix}
a_{11} & a_{12} & a_{13} \\
a_{21} & a_{22} & a_{23} \\
a_{31} & a_{32} & a_{33}
\end{pmatrix}
\]

\section{Tables}
Below is a sample table showing data:

\begin{table}[h]
\centering
\begin{tabular}{@{}lrrr@{}}
\toprule
\textbf{Category} & \textbf{Value 1} & \textbf{Value 2} & \textbf{Total} \\
\midrule
Item A & 10.5 & 20.3 & 30.8 \\
Item B & 15.2 & 25.1 & 40.3 \\
Item C & 22.1 & 12.7 & 34.8 \\
\bottomrule
\end{tabular}
\caption{Sample Data Table}
\label{tab:sample}
\end{table}

\section{Lists}
\subsection{Bullet Points}
\begin{itemize}
  \item First item
  \item Second item
  \item Third item with sub-items:
    \begin{itemize}
      \item Sub-item 1
      \item Sub-item 2
    \end{itemize}
\end{itemize}

\subsection{Numbered List}
\begin{enumerate}
  \item First step
  \item Second step
  \item Third step
\end{enumerate}

\section{Citations and References}
When you use the citation search feature, references will be automatically managed. For example, you might cite a paper like this: \cite{example2024}. All citations will be automatically added to your \texttt{cite.bib} file.

\section{Conclusion}
This template demonstrates basic \LaTeX{} features including mathematical formulas, tables, and lists. The bibliography system is now set up to work seamlessly with Bibby AI's citation management. Feel free to modify it for your needs!

\bibliographystyle{plain}
\bibliography{cite}

\end{document}\documentclass{article}
\usepackage[utf8]{inputenc}
\usepackage{booktabs}
\usepackage{amsmath}
\usepackage{natbib}

\title{My Document}
\author{Your Name}
\date{\today}

\begin{document}

\maketitle

\section{Introduction}
This is a sample document demonstrating various \LaTeX{} features including mathematical equations, tables, and formatting. You can cite papers using the citation search feature, and they will be automatically added to your \texttt{cite.bib} file.

\section{Mathematical Expressions}
Here are some example equations:

\subsection{Quadratic Formula}
The quadratic formula is given by:
\[x = \frac{-b \pm \sqrt{b^2 - 4ac}}{2a}\]

\subsection{Matrix Example}
A sample matrix:
\[
\begin{pmatrix}
a_{11} & a_{12} & a_{13} \\
a_{21} & a_{22} & a_{23} \\
a_{31} & a_{32} & a_{33}
\end{pmatrix}
\]

\section{Tables}
Below is a sample table showing data:

\begin{table}[h]
\centering
\begin{tabular}{@{}lrrr@{}}
\toprule
\textbf{Category} & \textbf{Value 1} & \textbf{Value 2} & \textbf{Total} \\
\midrule
Item A & 10.5 & 20.3 & 30.8 \\
Item B & 15.2 & 25.1 & 40.3 \\
Item C & 22.1 & 12.7 & 34.8 \\
\bottomrule
\end{tabular}
\caption{Sample Data Table}
\label{tab:sample}
\end{table}

\section{Lists}
\subsection{Bullet Points}
\begin{itemize}
  \item First item
  \item Second item
  \item Third item with sub-items:
    \begin{itemize}
      \item Sub-item 1
      \item Sub-item 2
    \end{itemize}
\end{itemize}

\subsection{Numbered List}
\begin{enumerate}
  \item First step
  \item Second step
  \item Third step
\end{enumerate}

\section{Citations and References}
When you use the citation search feature, references will be automatically managed. For example, you might cite a paper like this: \cite{example2024}. All citations will be automatically added to your \texttt{cite.bib} file.

\section{Conclusion}
This template demonstrates basic \LaTeX{} features including mathematical formulas, tables, and lists. The bibliography system is now set up to work seamlessly with Bibby AI's citation management. Feel free to modify it for your needs!

\bibliographystyle{plain}
\bibliography{cite}

\end{document}\documentclass{article}
\usepackage[utf8]{inputenc}
\usepackage{booktabs}
\usepackage{amsmath}
\usepackage{natbib}

\title{My Document}
\author{Your Name}
\date{\today}

\begin{document}

\maketitle

\section{Introduction}
This is a sample document demonstrating various \LaTeX{} features including mathematical equations, tables, and formatting. You can cite papers using the citation search feature, and they will be automatically added to your \texttt{cite.bib} file.

\section{Mathematical Expressions}
Here are some example equations:

\subsection{Quadratic Formula}
The quadratic formula is given by:
\[x = \frac{-b \pm \sqrt{b^2 - 4ac}}{2a}\]

\subsection{Matrix Example}
A sample matrix:
\[
\begin{pmatrix}
a_{11} & a_{12} & a_{13} \\
a_{21} & a_{22} & a_{23} \\
a_{31} & a_{32} & a_{33}
\end{pmatrix}
\]

\section{Tables}
Below is a sample table showing data:

\begin{table}[h]
\centering
\begin{tabular}{@{}lrrr@{}}
\toprule
\textbf{Category} & \textbf{Value 1} & \textbf{Value 2} & \textbf{Total} \\
\midrule
Item A & 10.5 & 20.3 & 30.8 \\
Item B & 15.2 & 25.1 & 40.3 \\
Item C & 22.1 & 12.7 & 34.8 \\
\bottomrule
\end{tabular}
\caption{Sample Data Table}
\label{tab:sample}
\end{table}

\section{Lists}
\subsection{Bullet Points}
\begin{itemize}
  \item First item
  \item Second item
  \item Third item with sub-items:
    \begin{itemize}
      \item Sub-item 1
      \item Sub-item 2
    \end{itemize}
\end{itemize}

\subsection{Numbered List}
\begin{enumerate}
  \item First step
  \item Second step
  \item Third step
\end{enumerate}

\section{Citations and References}
When you use the citation search feature, references will be automatically managed. For example, you might cite a paper like this: \cite{example2024}. All citations will be automatically added to your \texttt{cite.bib} file.

\section{Conclusion}
This template demonstrates basic \LaTeX{} features including mathematical formulas, tables, and lists. The bibliography system is now set up to work seamlessly with Bibby AI's citation management. Feel free to modify it for your needs!

\bibliographystyle{plain}
\bibliography{cite}

\end{document}\documentclass{article}
\usepackage[utf8]{inputenc}
\usepackage{booktabs}
\usepackage{amsmath}
\usepackage{natbib}

\title{My Document}
\author{Your Name}
\date{\today}

\begin{document}

\maketitle

\section{Introduction}
This is a sample document demonstrating various \LaTeX{} features including mathematical equations, tables, and formatting. You can cite papers using the citation search feature, and they will be automatically added to your \texttt{cite.bib} file.

\section{Mathematical Expressions}
Here are some example equations:

\subsection{Quadratic Formula}
The quadratic formula is given by:
\[x = \frac{-b \pm \sqrt{b^2 - 4ac}}{2a}\]

\subsection{Matrix Example}
A sample matrix:
\[
\begin{pmatrix}
a_{11} & a_{12} & a_{13} \\
a_{21} & a_{22} & a_{23} \\
a_{31} & a_{32} & a_{33}
\end{pmatrix}
\]

\section{Tables}
Below is a sample table showing data:

\begin{table}[h]
\centering
\begin{tabular}{@{}lrrr@{}}
\toprule
\textbf{Category} & \textbf{Value 1} & \textbf{Value 2} & \textbf{Total} \\
\midrule
Item A & 10.5 & 20.3 & 30.8 \\
Item B & 15.2 & 25.1 & 40.3 \\
Item C & 22.1 & 12.7 & 34.8 \\
\bottomrule
\end{tabular}
\caption{Sample Data Table}
\label{tab:sample}
\end{table}

\section{Lists}
\subsection{Bullet Points}
\begin{itemize}
  \item First item
  \item Second item
  \item Third item with sub-items:
    \begin{itemize}
      \item Sub-item 1
      \item Sub-item 2
    \end{itemize}
\end{itemize}

\subsection{Numbered List}
\begin{enumerate}
  \item First step
  \item Second step
  \item Third step
\end{enumerate}

\section{Citations and References}
When you use the citation search feature, references will be automatically managed. For example, you might cite a paper like this: \cite{example2024}. All citations will be automatically added to your \texttt{cite.bib} file.

\section{Conclusion}
This template demonstrates basic \LaTeX{} features including mathematical formulas, tables, and lists. The bibliography system is now set up to work seamlessly with Bibby AI's citation management. Feel free to modify it for your needs!

\bibliographystyle{plain}
\bibliography{cite}

\end{document}\documentclass{article}
\usepackage[utf8]{inputenc}
\usepackage{booktabs}
\usepackage{amsmath}
\usepackage{natbib}

\title{My Document}
\author{Your Name}
\date{\today}

\begin{document}

\maketitle

\section{Introduction}
This is a sample document demonstrating various \LaTeX{} features including mathematical equations, tables, and formatting. You can cite papers using the citation search feature, and they will be automatically added to your \texttt{cite.bib} file.

\section{Mathematical Expressions}
Here are some example equations:

\subsection{Quadratic Formula}
The quadratic formula is given by:
\[x = \frac{-b \pm \sqrt{b^2 - 4ac}}{2a}\]

\subsection{Matrix Example}
A sample matrix:
\[
\begin{pmatrix}
a_{11} & a_{12} & a_{13} \\
a_{21} & a_{22} & a_{23} \\
a_{31} & a_{32} & a_{33}
\end{pmatrix}
\]

\section{Tables}
Below is a sample table showing data:

\begin{table}[h]
\centering
\begin{tabular}{@{}lrrr@{}}
\toprule
\textbf{Category} & \textbf{Value 1} & \textbf{Value 2} & \textbf{Total} \\
\midrule
Item A & 10.5 & 20.3 & 30.8 \\
Item B & 15.2 & 25.1 & 40.3 \\
Item C & 22.1 & 12.7 & 34.8 \\
\bottomrule
\end{tabular}
\caption{Sample Data Table}
\label{tab:sample}
\end{table}

\section{Lists}
\subsection{Bullet Points}
\begin{itemize}
  \item First item
  \item Second item
  \item Third item with sub-items:
    \begin{itemize}
      \item Sub-item 1
      \item Sub-item 2
    \end{itemize}
\end{itemize}

\subsection{Numbered List}
\begin{enumerate}
  \item First step
  \item Second step
  \item Third step
\end{enumerate}

\section{Citations and References}
When you use the citation search feature, references will be automatically managed. For example, you might cite a paper like this: \cite{example2024}. All citations will be automatically added to your \texttt{cite.bib} file.

\section{Conclusion}
This template demonstrates basic \LaTeX{} features including mathematical formulas, tables, and lists. The bibliography system is now set up to work seamlessly with Bibby AI's citation management. Feel free to modify it for your needs!

\bibliographystyle{plain}
\bibliography{cite}

\end{document}\documentclass{article}
\usepackage[utf8]{inputenc}
\usepackage{booktabs}
\usepackage{amsmath}
\usepackage{natbib}

\title{My Document}
\author{Your Name}
\date{\today}

\begin{document}

\maketitle

\section{Introduction}
This is a sample document demonstrating various \LaTeX{} features including mathematical equations, tables, and formatting. You can cite papers using the citation search feature, and they will be automatically added to your \texttt{cite.bib} file.

\section{Mathematical Expressions}
Here are some example equations:

\subsection{Quadratic Formula}
The quadratic formula is given by:
\[x = \frac{-b \pm \sqrt{b^2 - 4ac}}{2a}\]

\subsection{Matrix Example}
A sample matrix:
\[
\begin{pmatrix}
a_{11} & a_{12} & a_{13} \\
a_{21} & a_{22} & a_{23} \\
a_{31} & a_{32} & a_{33}
\end{pmatrix}
\]

\section{Tables}
Below is a sample table showing data:

\begin{table}[h]
\centering
\begin{tabular}{@{}lrrr@{}}
\toprule
\textbf{Category} & \textbf{Value 1} & \textbf{Value 2} & \textbf{Total} \\
\midrule
Item A & 10.5 & 20.3 & 30.8 \\
Item B & 15.2 & 25.1 & 40.3 \\
Item C & 22.1 & 12.7 & 34.8 \\
\bottomrule
\end{tabular}
\caption{Sample Data Table}
\label{tab:sample}
\end{table}

\section{Lists}
\subsection{Bullet Points}
\begin{itemize}
  \item First item
  \item Second item
  \item Third item with sub-items:
    \begin{itemize}
      \item Sub-item 1
      \item Sub-item 2
    \end{itemize}
\end{itemize}

\subsection{Numbered List}
\begin{enumerate}
  \item First step
  \item Second step
  \item Third step
\end{enumerate}

\section{Citations and References}
When you use the citation search feature, references will be automatically managed. For example, you might cite a paper like this: \cite{example2024}. All citations will be automatically added to your \texttt{cite.bib} file.

\section{Conclusion}
This template demonstrates basic \LaTeX{} features including mathematical formulas, tables, and lists. The bibliography system is now set up to work seamlessly with Bibby AI's citation management. Feel free to modify it for your needs!

\bibliographystyle{plain}
\bibliography{cite}

\end{document}\documentclass{article}
\usepackage[utf8]{inputenc}
\usepackage{booktabs}
\usepackage{amsmath}
\usepackage{natbib}

\title{My Document}
\author{Your Name}
\date{\today}

\begin{document}

\maketitle

\section{Introduction}
This is a sample document demonstrating various \LaTeX{} features including mathematical equations, tables, and formatting. You can cite papers using the citation search feature, and they will be automatically added to your \texttt{cite.bib} file.

\section{Mathematical Expressions}
Here are some example equations:

\subsection{Quadratic Formula}
The quadratic formula is given by:
\[x = \frac{-b \pm \sqrt{b^2 - 4ac}}{2a}\]

\subsection{Matrix Example}
A sample matrix:
\[
\begin{pmatrix}
a_{11} & a_{12} & a_{13} \\
a_{21} & a_{22} & a_{23} \\
a_{31} & a_{32} & a_{33}
\end{pmatrix}
\]

\section{Tables}
Below is a sample table showing data:

\begin{table}[h]
\centering
\begin{tabular}{@{}lrrr@{}}
\toprule
\textbf{Category} & \textbf{Value 1} & \textbf{Value 2} & \textbf{Total} \\
\midrule
Item A & 10.5 & 20.3 & 30.8 \\
Item B & 15.2 & 25.1 & 40.3 \\
Item C & 22.1 & 12.7 & 34.8 \\
\bottomrule
\end{tabular}
\caption{Sample Data Table}
\label{tab:sample}
\end{table}

\section{Lists}
\subsection{Bullet Points}
\begin{itemize}
  \item First item
  \item Second item
  \item Third item with sub-items:
    \begin{itemize}
      \item Sub-item 1
      \item Sub-item 2
    \end{itemize}
\end{itemize}

\subsection{Numbered List}
\begin{enumerate}
  \item First step
  \item Second step
  \item Third step
\end{enumerate}

\section{Citations and References}
When you use the citation search feature, references will be automatically managed. For example, you might cite a paper like this: \cite{example2024}. All citations will be automatically added to your \texttt{cite.bib} file.

\section{Conclusion}
This template demonstrates basic \LaTeX{} features including mathematical formulas, tables, and lists. The bibliography system is now set up to work seamlessly with Bibby AI's citation management. Feel free to modify it for your needs!

\bibliographystyle{plain}
\bibliography{cite}

\end{document}\documentclass{article}
\usepackage[utf8]{inputenc}
\usepackage{booktabs}
\usepackage{amsmath}
\usepackage{natbib}

\title{My Document}
\author{Your Name}
\date{\today}

\begin{document}

\maketitle

\section{Introduction}
This is a sample document demonstrating various \LaTeX{} features including mathematical equations, tables, and formatting. You can cite papers using the citation search feature, and they will be automatically added to your \texttt{cite.bib} file.

\section{Mathematical Expressions}
Here are some example equations:

\subsection{Quadratic Formula}
The quadratic formula is given by:
\[x = \frac{-b \pm \sqrt{b^2 - 4ac}}{2a}\]

\subsection{Matrix Example}
A sample matrix:
\[
\begin{pmatrix}
a_{11} & a_{12} & a_{13} \\
a_{21} & a_{22} & a_{23} \\
a_{31} & a_{32} & a_{33}
\end{pmatrix}
\]

\section{Tables}
Below is a sample table showing data:

\begin{table}[h]
\centering
\begin{tabular}{@{}lrrr@{}}
\toprule
\textbf{Category} & \textbf{Value 1} & \textbf{Value 2} & \textbf{Total} \\
\midrule
Item A & 10.5 & 20.3 & 30.8 \\
Item B & 15.2 & 25.1 & 40.3 \\
Item C & 22.1 & 12.7 & 34.8 \\
\bottomrule
\end{tabular}
\caption{Sample Data Table}
\label{tab:sample}
\end{table}

\section{Lists}
\subsection{Bullet Points}
\begin{itemize}
  \item First item
  \item Second item
  \item Third item with sub-items:
    \begin{itemize}
      \item Sub-item 1
      \item Sub-item 2
    \end{itemize}
\end{itemize}

\subsection{Numbered List}
\begin{enumerate}
  \item First step
  \item Second step
  \item Third step
\end{enumerate}

\section{Citations and References}
When you use the citation search feature, references will be automatically managed. For example, you might cite a paper like this: \cite{example2024}. All citations will be automatically added to your \texttt{cite.bib} file.

\section{Conclusion}
This template demonstrates basic \LaTeX{} features including mathematical formulas, tables, and lists. The bibliography system is now set up to work seamlessly with Bibby AI's citation management. Feel free to modify it for your needs!

\bibliographystyle{plain}
\bibliography{cite}

\end{document}\documentclass{article}
\usepackage[utf8]{inputenc}
\usepackage{booktabs}
\usepackage{amsmath}
\usepackage{natbib}

\title{My Document}
\author{Your Name}
\date{\today}

\begin{document}

\maketitle

\section{Introduction}
This is a sample document demonstrating various \LaTeX{} features including mathematical equations, tables, and formatting. You can cite papers using the citation search feature, and they will be automatically added to your \texttt{cite.bib} file.

\section{Mathematical Expressions}
Here are some example equations:

\subsection{Quadratic Formula}
The quadratic formula is given by:
\[x = \frac{-b \pm \sqrt{b^2 - 4ac}}{2a}\]

\subsection{Matrix Example}
A sample matrix:
\[
\begin{pmatrix}
a_{11} & a_{12} & a_{13} \\
a_{21} & a_{22} & a_{23} \\
a_{31} & a_{32} & a_{33}
\end{pmatrix}
\]

\section{Tables}
Below is a sample table showing data:

\begin{table}[h]
\centering
\begin{tabular}{@{}lrrr@{}}
\toprule
\textbf{Category} & \textbf{Value 1} & \textbf{Value 2} & \textbf{Total} \\
\midrule
Item A & 10.5 & 20.3 & 30.8 \\
Item B & 15.2 & 25.1 & 40.3 \\
Item C & 22.1 & 12.7 & 34.8 \\
\bottomrule
\end{tabular}
\caption{Sample Data Table}
\label{tab:sample}
\end{table}

\section{Lists}
\subsection{Bullet Points}
\begin{itemize}
  \item First item
  \item Second item
  \item Third item with sub-items:
    \begin{itemize}
      \item Sub-item 1
      \item Sub-item 2
    \end{itemize}
\end{itemize}

\subsection{Numbered List}
\begin{enumerate}
  \item First step
  \item Second step
  \item Third step
\end{enumerate}

\section{Citations and References}
When you use the citation search feature, references will be automatically managed. For example, you might cite a paper like this: \cite{example2024}. All citations will be automatically added to your \texttt{cite.bib} file.

\section{Conclusion}
This template demonstrates basic \LaTeX{} features including mathematical formulas, tables, and lists. The bibliography system is now set up to work seamlessly with Bibby AI's citation management. Feel free to modify it for your needs!

\bibliographystyle{plain}
\bibliography{cite}

\end{document}\documentclass{article}
\usepackage[utf8]{inputenc}
\usepackage{booktabs}
\usepackage{amsmath}
\usepackage{natbib}

\title{My Document}
\author{Your Name}
\date{\today}

\begin{document}

\maketitle

\section{Introduction}
This is a sample document demonstrating various \LaTeX{} features including mathematical equations, tables, and formatting. You can cite papers using the citation search feature, and they will be automatically added to your \texttt{cite.bib} file.

\section{Mathematical Expressions}
Here are some example equations:

\subsection{Quadratic Formula}
The quadratic formula is given by:
\[x = \frac{-b \pm \sqrt{b^2 - 4ac}}{2a}\]

\subsection{Matrix Example}
A sample matrix:
\[
\begin{pmatrix}
a_{11} & a_{12} & a_{13} \\
a_{21} & a_{22} & a_{23} \\
a_{31} & a_{32} & a_{33}
\end{pmatrix}
\]

\section{Tables}
Below is a sample table showing data:

\begin{table}[h]
\centering
\begin{tabular}{@{}lrrr@{}}
\toprule
\textbf{Category} & \textbf{Value 1} & \textbf{Value 2} & \textbf{Total} \\
\midrule
Item A & 10.5 & 20.3 & 30.8 \\
Item B & 15.2 & 25.1 & 40.3 \\
Item C & 22.1 & 12.7 & 34.8 \\
\bottomrule
\end{tabular}
\caption{Sample Data Table}
\label{tab:sample}
\end{table}

\section{Lists}
\subsection{Bullet Points}
\begin{itemize}
  \item First item
  \item Second item
  \item Third item with sub-items:
    \begin{itemize}
      \item Sub-item 1
      \item Sub-item 2
    \end{itemize}
\end{itemize}

\subsection{Numbered List}
\begin{enumerate}
  \item First step
  \item Second step
  \item Third step
\end{enumerate}

\section{Citations and References}
When you use the citation search feature, references will be automatically managed. For example, you might cite a paper like this: \cite{example2024}. All citations will be automatically added to your \texttt{cite.bib} file.

\section{Conclusion}
This template demonstrates basic \LaTeX{} features including mathematical formulas, tables, and lists. The bibliography system is now set up to work seamlessly with Bibby AI's citation management. Feel free to modify it for your needs!

\bibliographystyle{plain}
\bibliography{cite}

\end{document}\documentclass{article}
\usepackage[utf8]{inputenc}
\usepackage{booktabs}
\usepackage{amsmath}
\usepackage{natbib}

\title{My Document}
\author{Your Name}
\date{\today}

\begin{document}

\maketitle

\section{Introduction}
This is a sample document demonstrating various \LaTeX{} features including mathematical equations, tables, and formatting. You can cite papers using the citation search feature, and they will be automatically added to your \texttt{cite.bib} file.

\section{Mathematical Expressions}
Here are some example equations:

\subsection{Quadratic Formula}
The quadratic formula is given by:
\[x = \frac{-b \pm \sqrt{b^2 - 4ac}}{2a}\]

\subsection{Matrix Example}
A sample matrix:
\[
\begin{pmatrix}
a_{11} & a_{12} & a_{13} \\
a_{21} & a_{22} & a_{23} \\
a_{31} & a_{32} & a_{33}
\end{pmatrix}
\]

\section{Tables}
Below is a sample table showing data:

\begin{table}[h]
\centering
\begin{tabular}{@{}lrrr@{}}
\toprule
\textbf{Category} & \textbf{Value 1} & \textbf{Value 2} & \textbf{Total} \\
\midrule
Item A & 10.5 & 20.3 & 30.8 \\
Item B & 15.2 & 25.1 & 40.3 \\
Item C & 22.1 & 12.7 & 34.8 \\
\bottomrule
\end{tabular}
\caption{Sample Data Table}
\label{tab:sample}
\end{table}

\section{Lists}
\subsection{Bullet Points}
\begin{itemize}
  \item First item
  \item Second item
  \item Third item with sub-items:
    \begin{itemize}
      \item Sub-item 1
      \item Sub-item 2
    \end{itemize}
\end{itemize}

\subsection{Numbered List}
\begin{enumerate}
  \item First step
  \item Second step
  \item Third step
\end{enumerate}

\section{Citations and References}
When you use the citation search feature, references will be automatically managed. For example, you might cite a paper like this: \cite{example2024}. All citations will be automatically added to your \texttt{cite.bib} file.

\section{Conclusion}
This template demonstrates basic \LaTeX{} features including mathematical formulas, tables, and lists. The bibliography system is now set up to work seamlessly with Bibby AI's citation management. Feel free to modify it for your needs!

\bibliographystyle{plain}
\bibliography{cite}

\end{document}\documentclass{article}
\usepackage[utf8]{inputenc}
\usepackage{booktabs}
\usepackage{amsmath}
\usepackage{natbib}

\title{My Document}
\author{Your Name}
\date{\today}

\begin{document}

\maketitle

\section{Introduction}
This is a sample document demonstrating various \LaTeX{} features including mathematical equations, tables, and formatting. You can cite papers using the citation search feature, and they will be automatically added to your \texttt{cite.bib} file.

\section{Mathematical Expressions}
Here are some example equations:

\subsection{Quadratic Formula}
The quadratic formula is given by:
\[x = \frac{-b \pm \sqrt{b^2 - 4ac}}{2a}\]

\subsection{Matrix Example}
A sample matrix:
\[
\begin{pmatrix}
a_{11} & a_{12} & a_{13} \\
a_{21} & a_{22} & a_{23} \\
a_{31} & a_{32} & a_{33}
\end{pmatrix}
\]

\section{Tables}
Below is a sample table showing data:

\begin{table}[h]
\centering
\begin{tabular}{@{}lrrr@{}}
\toprule
\textbf{Category} & \textbf{Value 1} & \textbf{Value 2} & \textbf{Total} \\
\midrule
Item A & 10.5 & 20.3 & 30.8 \\
Item B & 15.2 & 25.1 & 40.3 \\
Item C & 22.1 & 12.7 & 34.8 \\
\bottomrule
\end{tabular}
\caption{Sample Data Table}
\label{tab:sample}
\end{table}

\section{Lists}
\subsection{Bullet Points}
\begin{itemize}
  \item First item
  \item Second item
  \item Third item with sub-items:
    \begin{itemize}
      \item Sub-item 1
      \item Sub-item 2
    \end{itemize}
\end{itemize}

\subsection{Numbered List}
\begin{enumerate}
  \item First step
  \item Second step
  \item Third step
\end{enumerate}

\section{Citations and References}
When you use the citation search feature, references will be automatically managed. For example, you might cite a paper like this: \cite{example2024}. All citations will be automatically added to your \texttt{cite.bib} file.

\section{Conclusion}
This template demonstrates basic \LaTeX{} features including mathematical formulas, tables, and lists. The bibliography system is now set up to work seamlessly with Bibby AI's citation management. Feel free to modify it for your needs!

\bibliographystyle{plain}
\bibliography{cite}

\end{document}\documentclass{article}
\usepackage[utf8]{inputenc}
\usepackage{booktabs}
\usepackage{amsmath}
\usepackage{natbib}

\title{My Document}
\author{Your Name}
\date{\today}

\begin{document}

\maketitle

\section{Introduction}
This is a sample document demonstrating various \LaTeX{} features including mathematical equations, tables, and formatting. You can cite papers using the citation search feature, and they will be automatically added to your \texttt{cite.bib} file.

\section{Mathematical Expressions}
Here are some example equations:

\subsection{Quadratic Formula}
The quadratic formula is given by:
\[x = \frac{-b \pm \sqrt{b^2 - 4ac}}{2a}\]

\subsection{Matrix Example}
A sample matrix:
\[
\begin{pmatrix}
a_{11} & a_{12} & a_{13} \\
a_{21} & a_{22} & a_{23} \\
a_{31} & a_{32} & a_{33}
\end{pmatrix}
\]

\section{Tables}
Below is a sample table showing data:

\begin{table}[h]
\centering
\begin{tabular}{@{}lrrr@{}}
\toprule
\textbf{Category} & \textbf{Value 1} & \textbf{Value 2} & \textbf{Total} \\
\midrule
Item A & 10.5 & 20.3 & 30.8 \\
Item B & 15.2 & 25.1 & 40.3 \\
Item C & 22.1 & 12.7 & 34.8 \\
\bottomrule
\end{tabular}
\caption{Sample Data Table}
\label{tab:sample}
\end{table}

\section{Lists}
\subsection{Bullet Points}
\begin{itemize}
  \item First item
  \item Second item
  \item Third item with sub-items:
    \begin{itemize}
      \item Sub-item 1
      \item Sub-item 2
    \end{itemize}
\end{itemize}

\subsection{Numbered List}
\begin{enumerate}
  \item First step
  \item Second step
  \item Third step
\end{enumerate}

\section{Citations and References}
When you use the citation search feature, references will be automatically managed. For example, you might cite a paper like this: \cite{example2024}. All citations will be automatically added to your \texttt{cite.bib} file.

\section{Conclusion}
This template demonstrates basic \LaTeX{} features including mathematical formulas, tables, and lists. The bibliography system is now set up to work seamlessly with Bibby AI's citation management. Feel free to modify it for your needs!

\bibliographystyle{plain}
\bibliography{cite}

\end{document}\documentclass{article}
\usepackage[utf8]{inputenc}
\usepackage{booktabs}
\usepackage{amsmath}
\usepackage{natbib}

\title{My Document}
\author{Your Name}
\date{\today}

\begin{document}

\maketitle

\section{Introduction}
This is a sample document demonstrating various \LaTeX{} features including mathematical equations, tables, and formatting. You can cite papers using the citation search feature, and they will be automatically added to your \texttt{cite.bib} file.

\section{Mathematical Expressions}
Here are some example equations:

\subsection{Quadratic Formula}
The quadratic formula is given by:
\[x = \frac{-b \pm \sqrt{b^2 - 4ac}}{2a}\]

\subsection{Matrix Example}
A sample matrix:
\[
\begin{pmatrix}
a_{11} & a_{12} & a_{13} \\
a_{21} & a_{22} & a_{23} \\
a_{31} & a_{32} & a_{33}
\end{pmatrix}
\]

\section{Tables}
Below is a sample table showing data:

\begin{table}[h]
\centering
\begin{tabular}{@{}lrrr@{}}
\toprule
\textbf{Category} & \textbf{Value 1} & \textbf{Value 2} & \textbf{Total} \\
\midrule
Item A & 10.5 & 20.3 & 30.8 \\
Item B & 15.2 & 25.1 & 40.3 \\
Item C & 22.1 & 12.7 & 34.8 \\
\bottomrule
\end{tabular}
\caption{Sample Data Table}
\label{tab:sample}
\end{table}

\section{Lists}
\subsection{Bullet Points}
\begin{itemize}
  \item First item
  \item Second item
  \item Third item with sub-items:
    \begin{itemize}
      \item Sub-item 1
      \item Sub-item 2
    \end{itemize}
\end{itemize}

\subsection{Numbered List}
\begin{enumerate}
  \item First step
  \item Second step
  \item Third step
\end{enumerate}

\section{Citations and References}
When you use the citation search feature, references will be automatically managed. For example, you might cite a paper like this: \cite{example2024}. All citations will be automatically added to your \texttt{cite.bib} file.

\section{Conclusion}
This template demonstrates basic \LaTeX{} features including mathematical formulas, tables, and lists. The bibliography system is now set up to work seamlessly with Bibby AI's citation management. Feel free to modify it for your needs!

\bibliographystyle{plain}
\bibliography{cite}

\end{document}\documentclass{article}
\usepackage[utf8]{inputenc}
\usepackage{booktabs}
\usepackage{amsmath}
\usepackage{natbib}

\title{My Document}
\author{Your Name}
\date{\today}

\begin{document}

\maketitle

\section{Introduction}
This is a sample document demonstrating various \LaTeX{} features including mathematical equations, tables, and formatting. You can cite papers using the citation search feature, and they will be automatically added to your \texttt{cite.bib} file.

\section{Mathematical Expressions}
Here are some example equations:

\subsection{Quadratic Formula}
The quadratic formula is given by:
\[x = \frac{-b \pm \sqrt{b^2 - 4ac}}{2a}\]

\subsection{Matrix Example}
A sample matrix:
\[
\begin{pmatrix}
a_{11} & a_{12} & a_{13} \\
a_{21} & a_{22} & a_{23} \\
a_{31} & a_{32} & a_{33}
\end{pmatrix}
\]

\section{Tables}
Below is a sample table showing data:

\begin{table}[h]
\centering
\begin{tabular}{@{}lrrr@{}}
\toprule
\textbf{Category} & \textbf{Value 1} & \textbf{Value 2} & \textbf{Total} \\
\midrule
Item A & 10.5 & 20.3 & 30.8 \\
Item B & 15.2 & 25.1 & 40.3 \\
Item C & 22.1 & 12.7 & 34.8 \\
\bottomrule
\end{tabular}
\caption{Sample Data Table}
\label{tab:sample}
\end{table}

\section{Lists}
\subsection{Bullet Points}
\begin{itemize}
  \item First item
  \item Second item
  \item Third item with sub-items:
    \begin{itemize}
      \item Sub-item 1
      \item Sub-item 2
    \end{itemize}
\end{itemize}

\subsection{Numbered List}
\begin{enumerate}
  \item First step
  \item Second step
  \item Third step
\end{enumerate}

\section{Citations and References}
When you use the citation search feature, references will be automatically managed. For example, you might cite a paper like this: \cite{example2024}. All citations will be automatically added to your \texttt{cite.bib} file.

\section{Conclusion}
This template demonstrates basic \LaTeX{} features including mathematical formulas, tables, and lists. The bibliography system is now set up to work seamlessly with Bibby AI's citation management. Feel free to modify it for your needs!

\bibliographystyle{plain}
\bibliography{cite}

\end{document}\documentclass{article}
\usepackage[utf8]{inputenc}
\usepackage{booktabs}
\usepackage{amsmath}
\usepackage{natbib}

\title{My Document}
\author{Your Name}
\date{\today}

\begin{document}

\maketitle

\section{Introduction}
This is a sample document demonstrating various \LaTeX{} features including mathematical equations, tables, and formatting. You can cite papers using the citation search feature, and they will be automatically added to your \texttt{cite.bib} file.

\section{Mathematical Expressions}
Here are some example equations:

\subsection{Quadratic Formula}
The quadratic formula is given by:
\[x = \frac{-b \pm \sqrt{b^2 - 4ac}}{2a}\]

\subsection{Matrix Example}
A sample matrix:
\[
\begin{pmatrix}
a_{11} & a_{12} & a_{13} \\
a_{21} & a_{22} & a_{23} \\
a_{31} & a_{32} & a_{33}
\end{pmatrix}
\]

\section{Tables}
Below is a sample table showing data:

\begin{table}[h]
\centering
\begin{tabular}{@{}lrrr@{}}
\toprule
\textbf{Category} & \textbf{Value 1} & \textbf{Value 2} & \textbf{Total} \\
\midrule
Item A & 10.5 & 20.3 & 30.8 \\
Item B & 15.2 & 25.1 & 40.3 \\
Item C & 22.1 & 12.7 & 34.8 \\
\bottomrule
\end{tabular}
\caption{Sample Data Table}
\label{tab:sample}
\end{table}

\section{Lists}
\subsection{Bullet Points}
\begin{itemize}
  \item First item
  \item Second item
  \item Third item with sub-items:
    \begin{itemize}
      \item Sub-item 1
      \item Sub-item 2
    \end{itemize}
\end{itemize}

\subsection{Numbered List}
\begin{enumerate}
  \item First step
  \item Second step
  \item Third step
\end{enumerate}

\section{Citations and References}
When you use the citation search feature, references will be automatically managed. For example, you might cite a paper like this: \cite{example2024}. All citations will be automatically added to your \texttt{cite.bib} file.

\section{Conclusion}
This template demonstrates basic \LaTeX{} features including mathematical formulas, tables, and lists. The bibliography system is now set up to work seamlessly with Bibby AI's citation management. Feel free to modify it for your needs!

\bibliographystyle{plain}
\bibliography{cite}

\end{document}\documentclass{article}
\usepackage[utf8]{inputenc}
\usepackage{booktabs}
\usepackage{amsmath}
\usepackage{natbib}

\title{My Document}
\author{Your Name}
\date{\today}

\begin{document}

\maketitle

\section{Introduction}
This is a sample document demonstrating various \LaTeX{} features including mathematical equations, tables, and formatting. You can cite papers using the citation search feature, and they will be automatically added to your \texttt{cite.bib} file.

\section{Mathematical Expressions}
Here are some example equations:

\subsection{Quadratic Formula}
The quadratic formula is given by:
\[x = \frac{-b \pm \sqrt{b^2 - 4ac}}{2a}\]

\subsection{Matrix Example}
A sample matrix:
\[
\begin{pmatrix}
a_{11} & a_{12} & a_{13} \\
a_{21} & a_{22} & a_{23} \\
a_{31} & a_{32} & a_{33}
\end{pmatrix}
\]

\section{Tables}
Below is a sample table showing data:

\begin{table}[h]
\centering
\begin{tabular}{@{}lrrr@{}}
\toprule
\textbf{Category} & \textbf{Value 1} & \textbf{Value 2} & \textbf{Total} \\
\midrule
Item A & 10.5 & 20.3 & 30.8 \\
Item B & 15.2 & 25.1 & 40.3 \\
Item C & 22.1 & 12.7 & 34.8 \\
\bottomrule
\end{tabular}
\caption{Sample Data Table}
\label{tab:sample}
\end{table}

\section{Lists}
\subsection{Bullet Points}
\begin{itemize}
  \item First item
  \item Second item
  \item Third item with sub-items:
    \begin{itemize}
      \item Sub-item 1
      \item Sub-item 2
    \end{itemize}
\end{itemize}

\subsection{Numbered List}
\begin{enumerate}
  \item First step
  \item Second step
  \item Third step
\end{enumerate}

\section{Citations and References}
When you use the citation search feature, references will be automatically managed. For example, you might cite a paper like this: \cite{example2024}. All citations will be automatically added to your \texttt{cite.bib} file.

\section{Conclusion}
This template demonstrates basic \LaTeX{} features including mathematical formulas, tables, and lists. The bibliography system is now set up to work seamlessly with Bibby AI's citation management. Feel free to modify it for your needs!

\bibliographystyle{plain}
\bibliography{cite}

\end{document}\documentclass{article}
\usepackage[utf8]{inputenc}
\usepackage{booktabs}
\usepackage{amsmath}
\usepackage{natbib}

\title{My Document}
\author{Your Name}
\date{\today}

\begin{document}

\maketitle

\section{Introduction}
This is a sample document demonstrating various \LaTeX{} features including mathematical equations, tables, and formatting. You can cite papers using the citation search feature, and they will be automatically added to your \texttt{cite.bib} file.

\section{Mathematical Expressions}
Here are some example equations:

\subsection{Quadratic Formula}
The quadratic formula is given by:
\[x = \frac{-b \pm \sqrt{b^2 - 4ac}}{2a}\]

\subsection{Matrix Example}
A sample matrix:
\[
\begin{pmatrix}
a_{11} & a_{12} & a_{13} \\
a_{21} & a_{22} & a_{23} \\
a_{31} & a_{32} & a_{33}
\end{pmatrix}
\]

\section{Tables}
Below is a sample table showing data:

\begin{table}[h]
\centering
\begin{tabular}{@{}lrrr@{}}
\toprule
\textbf{Category} & \textbf{Value 1} & \textbf{Value 2} & \textbf{Total} \\
\midrule
Item A & 10.5 & 20.3 & 30.8 \\
Item B & 15.2 & 25.1 & 40.3 \\
Item C & 22.1 & 12.7 & 34.8 \\
\bottomrule
\end{tabular}
\caption{Sample Data Table}
\label{tab:sample}
\end{table}

\section{Lists}
\subsection{Bullet Points}
\begin{itemize}
  \item First item
  \item Second item
  \item Third item with sub-items:
    \begin{itemize}
      \item Sub-item 1
      \item Sub-item 2
    \end{itemize}
\end{itemize}

\subsection{Numbered List}
\begin{enumerate}
  \item First step
  \item Second step
  \item Third step
\end{enumerate}

\section{Citations and References}
When you use the citation search feature, references will be automatically managed. For example, you might cite a paper like this: \cite{example2024}. All citations will be automatically added to your \texttt{cite.bib} file.

\section{Conclusion}
This template demonstrates basic \LaTeX{} features including mathematical formulas, tables, and lists. The bibliography system is now set up to work seamlessly with Bibby AI's citation management. Feel free to modify it for your needs!

\bibliographystyle{plain}
\bibliography{cite}

\end{document}\documentclass{article}
\usepackage[utf8]{inputenc}
\usepackage{booktabs}
\usepackage{amsmath}
\usepackage{natbib}

\title{My Document}
\author{Your Name}
\date{\today}

\begin{document}

\maketitle

\section{Introduction}
This is a sample document demonstrating various \LaTeX{} features including mathematical equations, tables, and formatting. You can cite papers using the citation search feature, and they will be automatically added to your \texttt{cite.bib} file.

\section{Mathematical Expressions}
Here are some example equations:

\subsection{Quadratic Formula}
The quadratic formula is given by:
\[x = \frac{-b \pm \sqrt{b^2 - 4ac}}{2a}\]

\subsection{Matrix Example}
A sample matrix:
\[
\begin{pmatrix}
a_{11} & a_{12} & a_{13} \\
a_{21} & a_{22} & a_{23} \\
a_{31} & a_{32} & a_{33}
\end{pmatrix}
\]

\section{Tables}
Below is a sample table showing data:

\begin{table}[h]
\centering
\begin{tabular}{@{}lrrr@{}}
\toprule
\textbf{Category} & \textbf{Value 1} & \textbf{Value 2} & \textbf{Total} \\
\midrule
Item A & 10.5 & 20.3 & 30.8 \\
Item B & 15.2 & 25.1 & 40.3 \\
Item C & 22.1 & 12.7 & 34.8 \\
\bottomrule
\end{tabular}
\caption{Sample Data Table}
\label{tab:sample}
\end{table}

\section{Lists}
\subsection{Bullet Points}
\begin{itemize}
  \item First item
  \item Second item
  \item Third item with sub-items:
    \begin{itemize}
      \item Sub-item 1
      \item Sub-item 2
    \end{itemize}
\end{itemize}

\subsection{Numbered List}
\begin{enumerate}
  \item First step
  \item Second step
  \item Third step
\end{enumerate}

\section{Citations and References}
When you use the citation search feature, references will be automatically managed. For example, you might cite a paper like this: \cite{example2024}. All citations will be automatically added to your \texttt{cite.bib} file.

\section{Conclusion}
This template demonstrates basic \LaTeX{} features including mathematical formulas, tables, and lists. The bibliography system is now set up to work seamlessly with Bibby AI's citation management. Feel free to modify it for your needs!

\bibliographystyle{plain}
\bibliography{cite}

\end{document}\documentclass{article}
\usepackage[utf8]{inputenc}
\usepackage{booktabs}
\usepackage{amsmath}
\usepackage{natbib}

\title{My Document}
\author{Your Name}
\date{\today}

\begin{document}

\maketitle

\section{Introduction}
This is a sample document demonstrating various \LaTeX{} features including mathematical equations, tables, and formatting. You can cite papers using the citation search feature, and they will be automatically added to your \texttt{cite.bib} file.

\section{Mathematical Expressions}
Here are some example equations:

\subsection{Quadratic Formula}
The quadratic formula is given by:
\[x = \frac{-b \pm \sqrt{b^2 - 4ac}}{2a}\]

\subsection{Matrix Example}
A sample matrix:
\[
\begin{pmatrix}
a_{11} & a_{12} & a_{13} \\
a_{21} & a_{22} & a_{23} \\
a_{31} & a_{32} & a_{33}
\end{pmatrix}
\]

\section{Tables}
Below is a sample table showing data:

\begin{table}[h]
\centering
\begin{tabular}{@{}lrrr@{}}
\toprule
\textbf{Category} & \textbf{Value 1} & \textbf{Value 2} & \textbf{Total} \\
\midrule
Item A & 10.5 & 20.3 & 30.8 \\
Item B & 15.2 & 25.1 & 40.3 \\
Item C & 22.1 & 12.7 & 34.8 \\
\bottomrule
\end{tabular}
\caption{Sample Data Table}
\label{tab:sample}
\end{table}

\section{Lists}
\subsection{Bullet Points}
\begin{itemize}
  \item First item
  \item Second item
  \item Third item with sub-items:
    \begin{itemize}
      \item Sub-item 1
      \item Sub-item 2
    \end{itemize}
\end{itemize}

\subsection{Numbered List}
\begin{enumerate}
  \item First step
  \item Second step
  \item Third step
\end{enumerate}

\section{Citations and References}
When you use the citation search feature, references will be automatically managed. For example, you might cite a paper like this: \cite{example2024}. All citations will be automatically added to your \texttt{cite.bib} file.

\section{Conclusion}
This template demonstrates basic \LaTeX{} features including mathematical formulas, tables, and lists. The bibliography system is now set up to work seamlessly with Bibby AI's citation management. Feel free to modify it for your needs!

\bibliographystyle{plain}
\bibliography{cite}

\end{document}\documentclass{article}
\usepackage[utf8]{inputenc}
\usepackage{booktabs}
\usepackage{amsmath}
\usepackage{natbib}

\title{My Document}
\author{Your Name}
\date{\today}

\begin{document}

\maketitle

\section{Introduction}
This is a sample document demonstrating various \LaTeX{} features including mathematical equations, tables, and formatting. You can cite papers using the citation search feature, and they will be automatically added to your \texttt{cite.bib} file.

\section{Mathematical Expressions}
Here are some example equations:

\subsection{Quadratic Formula}
The quadratic formula is given by:
\[x = \frac{-b \pm \sqrt{b^2 - 4ac}}{2a}\]

\subsection{Matrix Example}
A sample matrix:
\[
\begin{pmatrix}
a_{11} & a_{12} & a_{13} \\
a_{21} & a_{22} & a_{23} \\
a_{31} & a_{32} & a_{33}
\end{pmatrix}
\]

\section{Tables}
Below is a sample table showing data:

\begin{table}[h]
\centering
\begin{tabular}{@{}lrrr@{}}
\toprule
\textbf{Category} & \textbf{Value 1} & \textbf{Value 2} & \textbf{Total} \\
\midrule
Item A & 10.5 & 20.3 & 30.8 \\
Item B & 15.2 & 25.1 & 40.3 \\
Item C & 22.1 & 12.7 & 34.8 \\
\bottomrule
\end{tabular}
\caption{Sample Data Table}
\label{tab:sample}
\end{table}

\section{Lists}
\subsection{Bullet Points}
\begin{itemize}
  \item First item
  \item Second item
  \item Third item with sub-items:
    \begin{itemize}
      \item Sub-item 1
      \item Sub-item 2
    \end{itemize}
\end{itemize}

\subsection{Numbered List}
\begin{enumerate}
  \item First step
  \item Second step
  \item Third step
\end{enumerate}

\section{Citations and References}
When you use the citation search feature, references will be automatically managed. For example, you might cite a paper like this: \cite{example2024}. All citations will be automatically added to your \texttt{cite.bib} file.

\section{Conclusion}
This template demonstrates basic \LaTeX{} features including mathematical formulas, tables, and lists. The bibliography system is now set up to work seamlessly with Bibby AI's citation management. Feel free to modify it for your needs!

\bibliographystyle{plain}
\bibliography{cite}

\end{document}\documentclass{article}
\usepackage[utf8]{inputenc}
\usepackage{booktabs}
\usepackage{amsmath}
\usepackage{natbib}

\title{My Document}
\author{Your Name}
\date{\today}

\begin{document}

\maketitle

\section{Introduction}
This is a sample document demonstrating various \LaTeX{} features including mathematical equations, tables, and formatting. You can cite papers using the citation search feature, and they will be automatically added to your \texttt{cite.bib} file.

\section{Mathematical Expressions}
Here are some example equations:

\subsection{Quadratic Formula}
The quadratic formula is given by:
\[x = \frac{-b \pm \sqrt{b^2 - 4ac}}{2a}\]

\subsection{Matrix Example}
A sample matrix:
\[
\begin{pmatrix}
a_{11} & a_{12} & a_{13} \\
a_{21} & a_{22} & a_{23} \\
a_{31} & a_{32} & a_{33}
\end{pmatrix}
\]

\section{Tables}
Below is a sample table showing data:

\begin{table}[h]
\centering
\begin{tabular}{@{}lrrr@{}}
\toprule
\textbf{Category} & \textbf{Value 1} & \textbf{Value 2} & \textbf{Total} \\
\midrule
Item A & 10.5 & 20.3 & 30.8 \\
Item B & 15.2 & 25.1 & 40.3 \\
Item C & 22.1 & 12.7 & 34.8 \\
\bottomrule
\end{tabular}
\caption{Sample Data Table}
\label{tab:sample}
\end{table}

\section{Lists}
\subsection{Bullet Points}
\begin{itemize}
  \item First item
  \item Second item
  \item Third item with sub-items:
    \begin{itemize}
      \item Sub-item 1
      \item Sub-item 2
    \end{itemize}
\end{itemize}

\subsection{Numbered List}
\begin{enumerate}
  \item First step
  \item Second step
  \item Third step
\end{enumerate}

\section{Citations and References}
When you use the citation search feature, references will be automatically managed. For example, you might cite a paper like this: \cite{example2024}. All citations will be automatically added to your \texttt{cite.bib} file.

\section{Conclusion}
This template demonstrates basic \LaTeX{} features including mathematical formulas, tables, and lists. The bibliography system is now set up to work seamlessly with Bibby AI's citation management. Feel free to modify it for your needs!

\bibliographystyle{plain}
\bibliography{cite}

\end{document}\documentclass{article}
\usepackage[utf8]{inputenc}
\usepackage{booktabs}
\usepackage{amsmath}
\usepackage{natbib}

\title{My Document}
\author{Your Name}
\date{\today}

\begin{document}

\maketitle

\section{Introduction}
This is a sample document demonstrating various \LaTeX{} features including mathematical equations, tables, and formatting. You can cite papers using the citation search feature, and they will be automatically added to your \texttt{cite.bib} file.

\section{Mathematical Expressions}
Here are some example equations:

\subsection{Quadratic Formula}
The quadratic formula is given by:
\[x = \frac{-b \pm \sqrt{b^2 - 4ac}}{2a}\]

\subsection{Matrix Example}
A sample matrix:
\[
\begin{pmatrix}
a_{11} & a_{12} & a_{13} \\
a_{21} & a_{22} & a_{23} \\
a_{31} & a_{32} & a_{33}
\end{pmatrix}
\]

\section{Tables}
Below is a sample table showing data:

\begin{table}[h]
\centering
\begin{tabular}{@{}lrrr@{}}
\toprule
\textbf{Category} & \textbf{Value 1} & \textbf{Value 2} & \textbf{Total} \\
\midrule
Item A & 10.5 & 20.3 & 30.8 \\
Item B & 15.2 & 25.1 & 40.3 \\
Item C & 22.1 & 12.7 & 34.8 \\
\bottomrule
\end{tabular}
\caption{Sample Data Table}
\label{tab:sample}
\end{table}

\section{Lists}
\subsection{Bullet Points}
\begin{itemize}
  \item First item
  \item Second item
  \item Third item with sub-items:
    \begin{itemize}
      \item Sub-item 1
      \item Sub-item 2
    \end{itemize}
\end{itemize}

\subsection{Numbered List}
\begin{enumerate}
  \item First step
  \item Second step
  \item Third step
\end{enumerate}

\section{Citations and References}
When you use the citation search feature, references will be automatically managed. For example, you might cite a paper like this: \cite{example2024}. All citations will be automatically added to your \texttt{cite.bib} file.

\section{Conclusion}
This template demonstrates basic \LaTeX{} features including mathematical formulas, tables, and lists. The bibliography system is now set up to work seamlessly with Bibby AI's citation management. Feel free to modify it for your needs!

\bibliographystyle{plain}
\bibliography{cite}

\end{document}\documentclass{article}
\usepackage[utf8]{inputenc}
\usepackage{booktabs}
\usepackage{amsmath}
\usepackage{natbib}

\title{My Document}
\author{Your Name}
\date{\today}

\begin{document}

\maketitle

\section{Introduction}
This is a sample document demonstrating various \LaTeX{} features including mathematical equations, tables, and formatting. You can cite papers using the citation search feature, and they will be automatically added to your \texttt{cite.bib} file.

\section{Mathematical Expressions}
Here are some example equations:

\subsection{Quadratic Formula}
The quadratic formula is given by:
\[x = \frac{-b \pm \sqrt{b^2 - 4ac}}{2a}\]

\subsection{Matrix Example}
A sample matrix:
\[
\begin{pmatrix}
a_{11} & a_{12} & a_{13} \\
a_{21} & a_{22} & a_{23} \\
a_{31} & a_{32} & a_{33}
\end{pmatrix}
\]

\section{Tables}
Below is a sample table showing data:

\begin{table}[h]
\centering
\begin{tabular}{@{}lrrr@{}}
\toprule
\textbf{Category} & \textbf{Value 1} & \textbf{Value 2} & \textbf{Total} \\
\midrule
Item A & 10.5 & 20.3 & 30.8 \\
Item B & 15.2 & 25.1 & 40.3 \\
Item C & 22.1 & 12.7 & 34.8 \\
\bottomrule
\end{tabular}
\caption{Sample Data Table}
\label{tab:sample}
\end{table}

\section{Lists}
\subsection{Bullet Points}
\begin{itemize}
  \item First item
  \item Second item
  \item Third item with sub-items:
    \begin{itemize}
      \item Sub-item 1
      \item Sub-item 2
    \end{itemize}
\end{itemize}

\subsection{Numbered List}
\begin{enumerate}
  \item First step
  \item Second step
  \item Third step
\end{enumerate}

\section{Citations and References}
When you use the citation search feature, references will be automatically managed. For example, you might cite a paper like this: \cite{example2024}. All citations will be automatically added to your \texttt{cite.bib} file.

\section{Conclusion}
This template demonstrates basic \LaTeX{} features including mathematical formulas, tables, and lists. The bibliography system is now set up to work seamlessly with Bibby AI's citation management. Feel free to modify it for your needs!

\bibliographystyle{plain}
\bibliography{cite}

\end{document}\documentclass{article}
\usepackage[utf8]{inputenc}
\usepackage{booktabs}
\usepackage{amsmath}
\usepackage{natbib}

\title{My Document}
\author{Your Name}
\date{\today}

\begin{document}

\maketitle

\section{Introduction}
This is a sample document demonstrating various \LaTeX{} features including mathematical equations, tables, and formatting. You can cite papers using the citation search feature, and they will be automatically added to your \texttt{cite.bib} file.

\section{Mathematical Expressions}
Here are some example equations:

\subsection{Quadratic Formula}
The quadratic formula is given by:
\[x = \frac{-b \pm \sqrt{b^2 - 4ac}}{2a}\]

\subsection{Matrix Example}
A sample matrix:
\[
\begin{pmatrix}
a_{11} & a_{12} & a_{13} \\
a_{21} & a_{22} & a_{23} \\
a_{31} & a_{32} & a_{33}
\end{pmatrix}
\]

\section{Tables}
Below is a sample table showing data:

\begin{table}[h]
\centering
\begin{tabular}{@{}lrrr@{}}
\toprule
\textbf{Category} & \textbf{Value 1} & \textbf{Value 2} & \textbf{Total} \\
\midrule
Item A & 10.5 & 20.3 & 30.8 \\
Item B & 15.2 & 25.1 & 40.3 \\
Item C & 22.1 & 12.7 & 34.8 \\
\bottomrule
\end{tabular}
\caption{Sample Data Table}
\label{tab:sample}
\end{table}

\section{Lists}
\subsection{Bullet Points}
\begin{itemize}
  \item First item
  \item Second item
  \item Third item with sub-items:
    \begin{itemize}
      \item Sub-item 1
      \item Sub-item 2
    \end{itemize}
\end{itemize}

\subsection{Numbered List}
\begin{enumerate}
  \item First step
  \item Second step
  \item Third step
\end{enumerate}

\section{Citations and References}
When you use the citation search feature, references will be automatically managed. For example, you might cite a paper like this: \cite{example2024}. All citations will be automatically added to your \texttt{cite.bib} file.

\section{Conclusion}
This template demonstrates basic \LaTeX{} features including mathematical formulas, tables, and lists. The bibliography system is now set up to work seamlessly with Bibby AI's citation management. Feel free to modify it for your needs!

\bibliographystyle{plain}
\bibliography{cite}

\end{document}\documentclass{article}
\usepackage[utf8]{inputenc}
\usepackage{booktabs}
\usepackage{amsmath}
\usepackage{natbib}

\title{My Document}
\author{Your Name}
\date{\today}

\begin{document}

\maketitle

\section{Introduction}
This is a sample document demonstrating various \LaTeX{} features including mathematical equations, tables, and formatting. You can cite papers using the citation search feature, and they will be automatically added to your \texttt{cite.bib} file.

\section{Mathematical Expressions}
Here are some example equations:

\subsection{Quadratic Formula}
The quadratic formula is given by:
\[x = \frac{-b \pm \sqrt{b^2 - 4ac}}{2a}\]

\subsection{Matrix Example}
A sample matrix:
\[
\begin{pmatrix}
a_{11} & a_{12} & a_{13} \\
a_{21} & a_{22} & a_{23} \\
a_{31} & a_{32} & a_{33}
\end{pmatrix}
\]

\section{Tables}
Below is a sample table showing data:

\begin{table}[h]
\centering
\begin{tabular}{@{}lrrr@{}}
\toprule
\textbf{Category} & \textbf{Value 1} & \textbf{Value 2} & \textbf{Total} \\
\midrule
Item A & 10.5 & 20.3 & 30.8 \\
Item B & 15.2 & 25.1 & 40.3 \\
Item C & 22.1 & 12.7 & 34.8 \\
\bottomrule
\end{tabular}
\caption{Sample Data Table}
\label{tab:sample}
\end{table}

\section{Lists}
\subsection{Bullet Points}
\begin{itemize}
  \item First item
  \item Second item
  \item Third item with sub-items:
    \begin{itemize}
      \item Sub-item 1
      \item Sub-item 2
    \end{itemize}
\end{itemize}

\subsection{Numbered List}
\begin{enumerate}
  \item First step
  \item Second step
  \item Third step
\end{enumerate}

\section{Citations and References}
When you use the citation search feature, references will be automatically managed. For example, you might cite a paper like this: \cite{example2024}. All citations will be automatically added to your \texttt{cite.bib} file.

\section{Conclusion}
This template demonstrates basic \LaTeX{} features including mathematical formulas, tables, and lists. The bibliography system is now set up to work seamlessly with Bibby AI's citation management. Feel free to modify it for your needs!

\bibliographystyle{plain}
\bibliography{cite}

\end{document}\documentclass{article}
\usepackage[utf8]{inputenc}
\usepackage{booktabs}
\usepackage{amsmath}
\usepackage{natbib}

\title{My Document}
\author{Your Name}
\date{\today}

\begin{document}

\maketitle

\section{Introduction}
This is a sample document demonstrating various \LaTeX{} features including mathematical equations, tables, and formatting. You can cite papers using the citation search feature, and they will be automatically added to your \texttt{cite.bib} file.

\section{Mathematical Expressions}
Here are some example equations:

\subsection{Quadratic Formula}
The quadratic formula is given by:
\[x = \frac{-b \pm \sqrt{b^2 - 4ac}}{2a}\]

\subsection{Matrix Example}
A sample matrix:
\[
\begin{pmatrix}
a_{11} & a_{12} & a_{13} \\
a_{21} & a_{22} & a_{23} \\
a_{31} & a_{32} & a_{33}
\end{pmatrix}
\]

\section{Tables}
Below is a sample table showing data:

\begin{table}[h]
\centering
\begin{tabular}{@{}lrrr@{}}
\toprule
\textbf{Category} & \textbf{Value 1} & \textbf{Value 2} & \textbf{Total} \\
\midrule
Item A & 10.5 & 20.3 & 30.8 \\
Item B & 15.2 & 25.1 & 40.3 \\
Item C & 22.1 & 12.7 & 34.8 \\
\bottomrule
\end{tabular}
\caption{Sample Data Table}
\label{tab:sample}
\end{table}

\section{Lists}
\subsection{Bullet Points}
\begin{itemize}
  \item First item
  \item Second item
  \item Third item with sub-items:
    \begin{itemize}
      \item Sub-item 1
      \item Sub-item 2
    \end{itemize}
\end{itemize}

\subsection{Numbered List}
\begin{enumerate}
  \item First step
  \item Second step
  \item Third step
\end{enumerate}

\section{Citations and References}
When you use the citation search feature, references will be automatically managed. For example, you might cite a paper like this: \cite{example2024}. All citations will be automatically added to your \texttt{cite.bib} file.

\section{Conclusion}
This template demonstrates basic \LaTeX{} features including mathematical formulas, tables, and lists. The bibliography system is now set up to work seamlessly with Bibby AI's citation management. Feel free to modify it for your needs!

\bibliographystyle{plain}
\bibliography{cite}

\end{document}\documentclass{article}
\usepackage[utf8]{inputenc}
\usepackage{booktabs}
\usepackage{amsmath}
\usepackage{natbib}

\title{My Document}
\author{Your Name}
\date{\today}

\begin{document}

\maketitle

\section{Introduction}
This is a sample document demonstrating various \LaTeX{} features including mathematical equations, tables, and formatting. You can cite papers using the citation search feature, and they will be automatically added to your \texttt{cite.bib} file.

\section{Mathematical Expressions}
Here are some example equations:

\subsection{Quadratic Formula}
The quadratic formula is given by:
\[x = \frac{-b \pm \sqrt{b^2 - 4ac}}{2a}\]

\subsection{Matrix Example}
A sample matrix:
\[
\begin{pmatrix}
a_{11} & a_{12} & a_{13} \\
a_{21} & a_{22} & a_{23} \\
a_{31} & a_{32} & a_{33}
\end{pmatrix}
\]

\section{Tables}
Below is a sample table showing data:

\begin{table}[h]
\centering
\begin{tabular}{@{}lrrr@{}}
\toprule
\textbf{Category} & \textbf{Value 1} & \textbf{Value 2} & \textbf{Total} \\
\midrule
Item A & 10.5 & 20.3 & 30.8 \\
Item B & 15.2 & 25.1 & 40.3 \\
Item C & 22.1 & 12.7 & 34.8 \\
\bottomrule
\end{tabular}
\caption{Sample Data Table}
\label{tab:sample}
\end{table}

\section{Lists}
\subsection{Bullet Points}
\begin{itemize}
  \item First item
  \item Second item
  \item Third item with sub-items:
    \begin{itemize}
      \item Sub-item 1
      \item Sub-item 2
    \end{itemize}
\end{itemize}

\subsection{Numbered List}
\begin{enumerate}
  \item First step
  \item Second step
  \item Third step
\end{enumerate}

\section{Citations and References}
When you use the citation search feature, references will be automatically managed. For example, you might cite a paper like this: \cite{example2024}. All citations will be automatically added to your \texttt{cite.bib} file.

\section{Conclusion}
This template demonstrates basic \LaTeX{} features including mathematical formulas, tables, and lists. The bibliography system is now set up to work seamlessly with Bibby AI's citation management. Feel free to modify it for your needs!

\bibliographystyle{plain}
\bibliography{cite}

\end{document}\documentclass{article}
\usepackage[utf8]{inputenc}
\usepackage{booktabs}
\usepackage{amsmath}
\usepackage{natbib}

\title{My Document}
\author{Your Name}
\date{\today}

\begin{document}

\maketitle

\section{Introduction}
This is a sample document demonstrating various \LaTeX{} features including mathematical equations, tables, and formatting. You can cite papers using the citation search feature, and they will be automatically added to your \texttt{cite.bib} file.

\section{Mathematical Expressions}
Here are some example equations:

\subsection{Quadratic Formula}
The quadratic formula is given by:
\[x = \frac{-b \pm \sqrt{b^2 - 4ac}}{2a}\]

\subsection{Matrix Example}
A sample matrix:
\[
\begin{pmatrix}
a_{11} & a_{12} & a_{13} \\
a_{21} & a_{22} & a_{23} \\
a_{31} & a_{32} & a_{33}
\end{pmatrix}
\]

\section{Tables}
Below is a sample table showing data:

\begin{table}[h]
\centering
\begin{tabular}{@{}lrrr@{}}
\toprule
\textbf{Category} & \textbf{Value 1} & \textbf{Value 2} & \textbf{Total} \\
\midrule
Item A & 10.5 & 20.3 & 30.8 \\
Item B & 15.2 & 25.1 & 40.3 \\
Item C & 22.1 & 12.7 & 34.8 \\
\bottomrule
\end{tabular}
\caption{Sample Data Table}
\label{tab:sample}
\end{table}

\section{Lists}
\subsection{Bullet Points}
\begin{itemize}
  \item First item
  \item Second item
  \item Third item with sub-items:
    \begin{itemize}
      \item Sub-item 1
      \item Sub-item 2
    \end{itemize}
\end{itemize}

\subsection{Numbered List}
\begin{enumerate}
  \item First step
  \item Second step
  \item Third step
\end{enumerate}

\section{Citations and References}
When you use the citation search feature, references will be automatically managed. For example, you might cite a paper like this: \cite{example2024}. All citations will be automatically added to your \texttt{cite.bib} file.

\section{Conclusion}
This template demonstrates basic \LaTeX{} features including mathematical formulas, tables, and lists. The bibliography system is now set up to work seamlessly with Bibby AI's citation management. Feel free to modify it for your needs!

\bibliographystyle{plain}
\bibliography{cite}

\end{document}\documentclass{article}
\usepackage[utf8]{inputenc}
\usepackage{booktabs}
\usepackage{amsmath}
\usepackage{natbib}

\title{My Document}
\author{Your Name}
\date{\today}

\begin{document}

\maketitle

\section{Introduction}
This is a sample document demonstrating various \LaTeX{} features including mathematical equations, tables, and formatting. You can cite papers using the citation search feature, and they will be automatically added to your \texttt{cite.bib} file.

\section{Mathematical Expressions}
Here are some example equations:

\subsection{Quadratic Formula}
The quadratic formula is given by:
\[x = \frac{-b \pm \sqrt{b^2 - 4ac}}{2a}\]

\subsection{Matrix Example}
A sample matrix:
\[
\begin{pmatrix}
a_{11} & a_{12} & a_{13} \\
a_{21} & a_{22} & a_{23} \\
a_{31} & a_{32} & a_{33}
\end{pmatrix}
\]

\section{Tables}
Below is a sample table showing data:

\begin{table}[h]
\centering
\begin{tabular}{@{}lrrr@{}}
\toprule
\textbf{Category} & \textbf{Value 1} & \textbf{Value 2} & \textbf{Total} \\
\midrule
Item A & 10.5 & 20.3 & 30.8 \\
Item B & 15.2 & 25.1 & 40.3 \\
Item C & 22.1 & 12.7 & 34.8 \\
\bottomrule
\end{tabular}
\caption{Sample Data Table}
\label{tab:sample}
\end{table}

\section{Lists}
\subsection{Bullet Points}
\begin{itemize}
  \item First item
  \item Second item
  \item Third item with sub-items:
    \begin{itemize}
      \item Sub-item 1
      \item Sub-item 2
    \end{itemize}
\end{itemize}

\subsection{Numbered List}
\begin{enumerate}
  \item First step
  \item Second step
  \item Third step
\end{enumerate}

\section{Citations and References}
When you use the citation search feature, references will be automatically managed. For example, you might cite a paper like this: \cite{example2024}. All citations will be automatically added to your \texttt{cite.bib} file.

\section{Conclusion}
This template demonstrates basic \LaTeX{} features including mathematical formulas, tables, and lists. The bibliography system is now set up to work seamlessly with Bibby AI's citation management. Feel free to modify it for your needs!

\bibliographystyle{plain}
\bibliography{cite}

\end{document}\documentclass{article}
\usepackage[utf8]{inputenc}
\usepackage{booktabs}
\usepackage{amsmath}
\usepackage{natbib}

\title{My Document}
\author{Your Name}
\date{\today}

\begin{document}

\maketitle

\section{Introduction}
This is a sample document demonstrating various \LaTeX{} features including mathematical equations, tables, and formatting. You can cite papers using the citation search feature, and they will be automatically added to your \texttt{cite.bib} file.

\section{Mathematical Expressions}
Here are some example equations:

\subsection{Quadratic Formula}
The quadratic formula is given by:
\[x = \frac{-b \pm \sqrt{b^2 - 4ac}}{2a}\]

\subsection{Matrix Example}
A sample matrix:
\[
\begin{pmatrix}
a_{11} & a_{12} & a_{13} \\
a_{21} & a_{22} & a_{23} \\
a_{31} & a_{32} & a_{33}
\end{pmatrix}
\]

\section{Tables}
Below is a sample table showing data:

\begin{table}[h]
\centering
\begin{tabular}{@{}lrrr@{}}
\toprule
\textbf{Category} & \textbf{Value 1} & \textbf{Value 2} & \textbf{Total} \\
\midrule
Item A & 10.5 & 20.3 & 30.8 \\
Item B & 15.2 & 25.1 & 40.3 \\
Item C & 22.1 & 12.7 & 34.8 \\
\bottomrule
\end{tabular}
\caption{Sample Data Table}
\label{tab:sample}
\end{table}

\section{Lists}
\subsection{Bullet Points}
\begin{itemize}
  \item First item
  \item Second item
  \item Third item with sub-items:
    \begin{itemize}
      \item Sub-item 1
      \item Sub-item 2
    \end{itemize}
\end{itemize}

\subsection{Numbered List}
\begin{enumerate}
  \item First step
  \item Second step
  \item Third step
\end{enumerate}

\section{Citations and References}
When you use the citation search feature, references will be automatically managed. For example, you might cite a paper like this: \cite{example2024}. All citations will be automatically added to your \texttt{cite.bib} file.

\section{Conclusion}
This template demonstrates basic \LaTeX{} features including mathematical formulas, tables, and lists. The bibliography system is now set up to work seamlessly with Bibby AI's citation management. Feel free to modify it for your needs!

\bibliographystyle{plain}
\bibliography{cite}

\end{document}\documentclass{article}
\usepackage[utf8]{inputenc}
\usepackage{booktabs}
\usepackage{amsmath}
\usepackage{natbib}

\title{My Document}
\author{Your Name}
\date{\today}

\begin{document}

\maketitle

\section{Introduction}
This is a sample document demonstrating various \LaTeX{} features including mathematical equations, tables, and formatting. You can cite papers using the citation search feature, and they will be automatically added to your \texttt{cite.bib} file.

\section{Mathematical Expressions}
Here are some example equations:

\subsection{Quadratic Formula}
The quadratic formula is given by:
\[x = \frac{-b \pm \sqrt{b^2 - 4ac}}{2a}\]

\subsection{Matrix Example}
A sample matrix:
\[
\begin{pmatrix}
a_{11} & a_{12} & a_{13} \\
a_{21} & a_{22} & a_{23} \\
a_{31} & a_{32} & a_{33}
\end{pmatrix}
\]

\section{Tables}
Below is a sample table showing data:

\begin{table}[h]
\centering
\begin{tabular}{@{}lrrr@{}}
\toprule
\textbf{Category} & \textbf{Value 1} & \textbf{Value 2} & \textbf{Total} \\
\midrule
Item A & 10.5 & 20.3 & 30.8 \\
Item B & 15.2 & 25.1 & 40.3 \\
Item C & 22.1 & 12.7 & 34.8 \\
\bottomrule
\end{tabular}
\caption{Sample Data Table}
\label{tab:sample}
\end{table}

\section{Lists}
\subsection{Bullet Points}
\begin{itemize}
  \item First item
  \item Second item
  \item Third item with sub-items:
    \begin{itemize}
      \item Sub-item 1
      \item Sub-item 2
    \end{itemize}
\end{itemize}

\subsection{Numbered List}
\begin{enumerate}
  \item First step
  \item Second step
  \item Third step
\end{enumerate}

\section{Citations and References}
When you use the citation search feature, references will be automatically managed. For example, you might cite a paper like this: \cite{example2024}. All citations will be automatically added to your \texttt{cite.bib} file.

\section{Conclusion}
This template demonstrates basic \LaTeX{} features including mathematical formulas, tables, and lists. The bibliography system is now set up to work seamlessly with Bibby AI's citation management. Feel free to modify it for your needs!

\bibliographystyle{plain}
\bibliography{cite}

\end{document}\documentclass{article}
\usepackage[utf8]{inputenc}
\usepackage{booktabs}
\usepackage{amsmath}
\usepackage{natbib}

\title{My Document}
\author{Your Name}
\date{\today}

\begin{document}

\maketitle

\section{Introduction}
This is a sample document demonstrating various \LaTeX{} features including mathematical equations, tables, and formatting. You can cite papers using the citation search feature, and they will be automatically added to your \texttt{cite.bib} file.

\section{Mathematical Expressions}
Here are some example equations:

\subsection{Quadratic Formula}
The quadratic formula is given by:
\[x = \frac{-b \pm \sqrt{b^2 - 4ac}}{2a}\]

\subsection{Matrix Example}
A sample matrix:
\[
\begin{pmatrix}
a_{11} & a_{12} & a_{13} \\
a_{21} & a_{22} & a_{23} \\
a_{31} & a_{32} & a_{33}
\end{pmatrix}
\]

\section{Tables}
Below is a sample table showing data:

\begin{table}[h]
\centering
\begin{tabular}{@{}lrrr@{}}
\toprule
\textbf{Category} & \textbf{Value 1} & \textbf{Value 2} & \textbf{Total} \\
\midrule
Item A & 10.5 & 20.3 & 30.8 \\
Item B & 15.2 & 25.1 & 40.3 \\
Item C & 22.1 & 12.7 & 34.8 \\
\bottomrule
\end{tabular}
\caption{Sample Data Table}
\label{tab:sample}
\end{table}

\section{Lists}
\subsection{Bullet Points}
\begin{itemize}
  \item First item
  \item Second item
  \item Third item with sub-items:
    \begin{itemize}
      \item Sub-item 1
      \item Sub-item 2
    \end{itemize}
\end{itemize}

\subsection{Numbered List}
\begin{enumerate}
  \item First step
  \item Second step
  \item Third step
\end{enumerate}

\section{Citations and References}
When you use the citation search feature, references will be automatically managed. For example, you might cite a paper like this: \cite{example2024}. All citations will be automatically added to your \texttt{cite.bib} file.

\section{Conclusion}
This template demonstrates basic \LaTeX{} features including mathematical formulas, tables, and lists. The bibliography system is now set up to work seamlessly with Bibby AI's citation management. Feel free to modify it for your needs!

\bibliographystyle{plain}
\bibliography{cite}

\end{document}\documentclass{article}
\usepackage[utf8]{inputenc}
\usepackage{booktabs}
\usepackage{amsmath}
\usepackage{natbib}

\title{My Document}
\author{Your Name}
\date{\today}

\begin{document}

\maketitle

\section{Introduction}
This is a sample document demonstrating various \LaTeX{} features including mathematical equations, tables, and formatting. You can cite papers using the citation search feature, and they will be automatically added to your \texttt{cite.bib} file.

\section{Mathematical Expressions}
Here are some example equations:

\subsection{Quadratic Formula}
The quadratic formula is given by:
\[x = \frac{-b \pm \sqrt{b^2 - 4ac}}{2a}\]

\subsection{Matrix Example}
A sample matrix:
\[
\begin{pmatrix}
a_{11} & a_{12} & a_{13} \\
a_{21} & a_{22} & a_{23} \\
a_{31} & a_{32} & a_{33}
\end{pmatrix}
\]

\section{Tables}
Below is a sample table showing data:

\begin{table}[h]
\centering
\begin{tabular}{@{}lrrr@{}}
\toprule
\textbf{Category} & \textbf{Value 1} & \textbf{Value 2} & \textbf{Total} \\
\midrule
Item A & 10.5 & 20.3 & 30.8 \\
Item B & 15.2 & 25.1 & 40.3 \\
Item C & 22.1 & 12.7 & 34.8 \\
\bottomrule
\end{tabular}
\caption{Sample Data Table}
\label{tab:sample}
\end{table}

\section{Lists}
\subsection{Bullet Points}
\begin{itemize}
  \item First item
  \item Second item
  \item Third item with sub-items:
    \begin{itemize}
      \item Sub-item 1
      \item Sub-item 2
    \end{itemize}
\end{itemize}

\subsection{Numbered List}
\begin{enumerate}
  \item First step
  \item Second step
  \item Third step
\end{enumerate}

\section{Citations and References}
When you use the citation search feature, references will be automatically managed. For example, you might cite a paper like this: \cite{example2024}. All citations will be automatically added to your \texttt{cite.bib} file.

\section{Conclusion}
This template demonstrates basic \LaTeX{} features including mathematical formulas, tables, and lists. The bibliography system is now set up to work seamlessly with Bibby AI's citation management. Feel free to modify it for your needs!

\bibliographystyle{plain}
\bibliography{cite}

\end{document}\documentclass{article}
\usepackage[utf8]{inputenc}
\usepackage{booktabs}
\usepackage{amsmath}
\usepackage{natbib}

\title{My Document}
\author{Your Name}
\date{\today}

\begin{document}

\maketitle

\section{Introduction}
This is a sample document demonstrating various \LaTeX{} features including mathematical equations, tables, and formatting. You can cite papers using the citation search feature, and they will be automatically added to your \texttt{cite.bib} file.

\section{Mathematical Expressions}
Here are some example equations:

\subsection{Quadratic Formula}
The quadratic formula is given by:
\[x = \frac{-b \pm \sqrt{b^2 - 4ac}}{2a}\]

\subsection{Matrix Example}
A sample matrix:
\[
\begin{pmatrix}
a_{11} & a_{12} & a_{13} \\
a_{21} & a_{22} & a_{23} \\
a_{31} & a_{32} & a_{33}
\end{pmatrix}
\]

\section{Tables}
Below is a sample table showing data:

\begin{table}[h]
\centering
\begin{tabular}{@{}lrrr@{}}
\toprule
\textbf{Category} & \textbf{Value 1} & \textbf{Value 2} & \textbf{Total} \\
\midrule
Item A & 10.5 & 20.3 & 30.8 \\
Item B & 15.2 & 25.1 & 40.3 \\
Item C & 22.1 & 12.7 & 34.8 \\
\bottomrule
\end{tabular}
\caption{Sample Data Table}
\label{tab:sample}
\end{table}

\section{Lists}
\subsection{Bullet Points}
\begin{itemize}
  \item First item
  \item Second item
  \item Third item with sub-items:
    \begin{itemize}
      \item Sub-item 1
      \item Sub-item 2
    \end{itemize}
\end{itemize}

\subsection{Numbered List}
\begin{enumerate}
  \item First step
  \item Second step
  \item Third step
\end{enumerate}

\section{Citations and References}
When you use the citation search feature, references will be automatically managed. For example, you might cite a paper like this: \cite{example2024}. All citations will be automatically added to your \texttt{cite.bib} file.

\section{Conclusion}
This template demonstrates basic \LaTeX{} features including mathematical formulas, tables, and lists. The bibliography system is now set up to work seamlessly with Bibby AI's citation management. Feel free to modify it for your needs!

\bibliographystyle{plain}
\bibliography{cite}

\end{document}\documentclass{article}
\usepackage[utf8]{inputenc}
\usepackage{booktabs}
\usepackage{amsmath}
\usepackage{natbib}

\title{My Document}
\author{Your Name}
\date{\today}

\begin{document}

\maketitle

\section{Introduction}
This is a sample document demonstrating various \LaTeX{} features including mathematical equations, tables, and formatting. You can cite papers using the citation search feature, and they will be automatically added to your \texttt{cite.bib} file.

\section{Mathematical Expressions}
Here are some example equations:

\subsection{Quadratic Formula}
The quadratic formula is given by:
\[x = \frac{-b \pm \sqrt{b^2 - 4ac}}{2a}\]

\subsection{Matrix Example}
A sample matrix:
\[
\begin{pmatrix}
a_{11} & a_{12} & a_{13} \\
a_{21} & a_{22} & a_{23} \\
a_{31} & a_{32} & a_{33}
\end{pmatrix}
\]

\section{Tables}
Below is a sample table showing data:

\begin{table}[h]
\centering
\begin{tabular}{@{}lrrr@{}}
\toprule
\textbf{Category} & \textbf{Value 1} & \textbf{Value 2} & \textbf{Total} \\
\midrule
Item A & 10.5 & 20.3 & 30.8 \\
Item B & 15.2 & 25.1 & 40.3 \\
Item C & 22.1 & 12.7 & 34.8 \\
\bottomrule
\end{tabular}
\caption{Sample Data Table}
\label{tab:sample}
\end{table}

\section{Lists}
\subsection{Bullet Points}
\begin{itemize}
  \item First item
  \item Second item
  \item Third item with sub-items:
    \begin{itemize}
      \item Sub-item 1
      \item Sub-item 2
    \end{itemize}
\end{itemize}

\subsection{Numbered List}
\begin{enumerate}
  \item First step
  \item Second step
  \item Third step
\end{enumerate}

\section{Citations and References}
When you use the citation search feature, references will be automatically managed. For example, you might cite a paper like this: \cite{example2024}. All citations will be automatically added to your \texttt{cite.bib} file.

\section{Conclusion}
This template demonstrates basic \LaTeX{} features including mathematical formulas, tables, and lists. The bibliography system is now set up to work seamlessly with Bibby AI's citation management. Feel free to modify it for your needs!

\bibliographystyle{plain}
\bibliography{cite}

\end{document}\documentclass{article}
\usepackage[utf8]{inputenc}
\usepackage{booktabs}
\usepackage{amsmath}
\usepackage{natbib}

\title{My Document}
\author{Your Name}
\date{\today}

\begin{document}

\maketitle

\section{Introduction}
This is a sample document demonstrating various \LaTeX{} features including mathematical equations, tables, and formatting. You can cite papers using the citation search feature, and they will be automatically added to your \texttt{cite.bib} file.

\section{Mathematical Expressions}
Here are some example equations:

\subsection{Quadratic Formula}
The quadratic formula is given by:
\[x = \frac{-b \pm \sqrt{b^2 - 4ac}}{2a}\]

\subsection{Matrix Example}
A sample matrix:
\[
\begin{pmatrix}
a_{11} & a_{12} & a_{13} \\
a_{21} & a_{22} & a_{23} \\
a_{31} & a_{32} & a_{33}
\end{pmatrix}
\]

\section{Tables}
Below is a sample table showing data:

\begin{table}[h]
\centering
\begin{tabular}{@{}lrrr@{}}
\toprule
\textbf{Category} & \textbf{Value 1} & \textbf{Value 2} & \textbf{Total} \\
\midrule
Item A & 10.5 & 20.3 & 30.8 \\
Item B & 15.2 & 25.1 & 40.3 \\
Item C & 22.1 & 12.7 & 34.8 \\
\bottomrule
\end{tabular}
\caption{Sample Data Table}
\label{tab:sample}
\end{table}

\section{Lists}
\subsection{Bullet Points}
\begin{itemize}
  \item First item
  \item Second item
  \item Third item with sub-items:
    \begin{itemize}
      \item Sub-item 1
      \item Sub-item 2
    \end{itemize}
\end{itemize}

\subsection{Numbered List}
\begin{enumerate}
  \item First step
  \item Second step
  \item Third step
\end{enumerate}

\section{Citations and References}
When you use the citation search feature, references will be automatically managed. For example, you might cite a paper like this: \cite{example2024}. All citations will be automatically added to your \texttt{cite.bib} file.

\section{Conclusion}
This template demonstrates basic \LaTeX{} features including mathematical formulas, tables, and lists. The bibliography system is now set up to work seamlessly with Bibby AI's citation management. Feel free to modify it for your needs!

\bibliographystyle{plain}
\bibliography{cite}

\end{document}\documentclass{article}
\usepackage[utf8]{inputenc}
\usepackage{booktabs}
\usepackage{amsmath}
\usepackage{natbib}

\title{My Document}
\author{Your Name}
\date{\today}

\begin{document}

\maketitle

\section{Introduction}
This is a sample document demonstrating various \LaTeX{} features including mathematical equations, tables, and formatting. You can cite papers using the citation search feature, and they will be automatically added to your \texttt{cite.bib} file.

\section{Mathematical Expressions}
Here are some example equations:

\subsection{Quadratic Formula}
The quadratic formula is given by:
\[x = \frac{-b \pm \sqrt{b^2 - 4ac}}{2a}\]

\subsection{Matrix Example}
A sample matrix:
\[
\begin{pmatrix}
a_{11} & a_{12} & a_{13} \\
a_{21} & a_{22} & a_{23} \\
a_{31} & a_{32} & a_{33}
\end{pmatrix}
\]

\section{Tables}
Below is a sample table showing data:

\begin{table}[h]
\centering
\begin{tabular}{@{}lrrr@{}}
\toprule
\textbf{Category} & \textbf{Value 1} & \textbf{Value 2} & \textbf{Total} \\
\midrule
Item A & 10.5 & 20.3 & 30.8 \\
Item B & 15.2 & 25.1 & 40.3 \\
Item C & 22.1 & 12.7 & 34.8 \\
\bottomrule
\end{tabular}
\caption{Sample Data Table}
\label{tab:sample}
\end{table}

\section{Lists}
\subsection{Bullet Points}
\begin{itemize}
  \item First item
  \item Second item
  \item Third item with sub-items:
    \begin{itemize}
      \item Sub-item 1
      \item Sub-item 2
    \end{itemize}
\end{itemize}

\subsection{Numbered List}
\begin{enumerate}
  \item First step
  \item Second step
  \item Third step
\end{enumerate}

\section{Citations and References}
When you use the citation search feature, references will be automatically managed. For example, you might cite a paper like this: \cite{example2024}. All citations will be automatically added to your \texttt{cite.bib} file.

\section{Conclusion}
This template demonstrates basic \LaTeX{} features including mathematical formulas, tables, and lists. The bibliography system is now set up to work seamlessly with Bibby AI's citation management. Feel free to modify it for your needs!

\bibliographystyle{plain}
\bibliography{cite}

\end{document}\documentclass{article}
\usepackage[utf8]{inputenc}
\usepackage{booktabs}
\usepackage{amsmath}
\usepackage{natbib}

\title{My Document}
\author{Your Name}
\date{\today}

\begin{document}

\maketitle

\section{Introduction}
This is a sample document demonstrating various \LaTeX{} features including mathematical equations, tables, and formatting. You can cite papers using the citation search feature, and they will be automatically added to your \texttt{cite.bib} file.

\section{Mathematical Expressions}
Here are some example equations:

\subsection{Quadratic Formula}
The quadratic formula is given by:
\[x = \frac{-b \pm \sqrt{b^2 - 4ac}}{2a}\]

\subsection{Matrix Example}
A sample matrix:
\[
\begin{pmatrix}
a_{11} & a_{12} & a_{13} \\
a_{21} & a_{22} & a_{23} \\
a_{31} & a_{32} & a_{33}
\end{pmatrix}
\]

\section{Tables}
Below is a sample table showing data:

\begin{table}[h]
\centering
\begin{tabular}{@{}lrrr@{}}
\toprule
\textbf{Category} & \textbf{Value 1} & \textbf{Value 2} & \textbf{Total} \\
\midrule
Item A & 10.5 & 20.3 & 30.8 \\
Item B & 15.2 & 25.1 & 40.3 \\
Item C & 22.1 & 12.7 & 34.8 \\
\bottomrule
\end{tabular}
\caption{Sample Data Table}
\label{tab:sample}
\end{table}

\section{Lists}
\subsection{Bullet Points}
\begin{itemize}
  \item First item
  \item Second item
  \item Third item with sub-items:
    \begin{itemize}
      \item Sub-item 1
      \item Sub-item 2
    \end{itemize}
\end{itemize}

\subsection{Numbered List}
\begin{enumerate}
  \item First step
  \item Second step
  \item Third step
\end{enumerate}

\section{Citations and References}
When you use the citation search feature, references will be automatically managed. For example, you might cite a paper like this: \cite{example2024}. All citations will be automatically added to your \texttt{cite.bib} file.

\section{Conclusion}
This template demonstrates basic \LaTeX{} features including mathematical formulas, tables, and lists. The bibliography system is now set up to work seamlessly with Bibby AI's citation management. Feel free to modify it for your needs!

\bibliographystyle{plain}
\bibliography{cite}

\end{document}\documentclass{article}
\usepackage[utf8]{inputenc}
\usepackage{booktabs}
\usepackage{amsmath}
\usepackage{natbib}

\title{My Document}
\author{Your Name}
\date{\today}

\begin{document}

\maketitle

\section{Introduction}
This is a sample document demonstrating various \LaTeX{} features including mathematical equations, tables, and formatting. You can cite papers using the citation search feature, and they will be automatically added to your \texttt{cite.bib} file.

\section{Mathematical Expressions}
Here are some example equations:

\subsection{Quadratic Formula}
The quadratic formula is given by:
\[x = \frac{-b \pm \sqrt{b^2 - 4ac}}{2a}\]

\subsection{Matrix Example}
A sample matrix:
\[
\begin{pmatrix}
a_{11} & a_{12} & a_{13} \\
a_{21} & a_{22} & a_{23} \\
a_{31} & a_{32} & a_{33}
\end{pmatrix}
\]

\section{Tables}
Below is a sample table showing data:

\begin{table}[h]
\centering
\begin{tabular}{@{}lrrr@{}}
\toprule
\textbf{Category} & \textbf{Value 1} & \textbf{Value 2} & \textbf{Total} \\
\midrule
Item A & 10.5 & 20.3 & 30.8 \\
Item B & 15.2 & 25.1 & 40.3 \\
Item C & 22.1 & 12.7 & 34.8 \\
\bottomrule
\end{tabular}
\caption{Sample Data Table}
\label{tab:sample}
\end{table}

\section{Lists}
\subsection{Bullet Points}
\begin{itemize}
  \item First item
  \item Second item
  \item Third item with sub-items:
    \begin{itemize}
      \item Sub-item 1
      \item Sub-item 2
    \end{itemize}
\end{itemize}

\subsection{Numbered List}
\begin{enumerate}
  \item First step
  \item Second step
  \item Third step
\end{enumerate}

\section{Citations and References}
When you use the citation search feature, references will be automatically managed. For example, you might cite a paper like this: \cite{example2024}. All citations will be automatically added to your \texttt{cite.bib} file.

\section{Conclusion}
This template demonstrates basic \LaTeX{} features including mathematical formulas, tables, and lists. The bibliography system is now set up to work seamlessly with Bibby AI's citation management. Feel free to modify it for your needs!

\bibliographystyle{plain}
\bibliography{cite}

\end{document}\documentclass{article}
\usepackage[utf8]{inputenc}
\usepackage{booktabs}
\usepackage{amsmath}
\usepackage{natbib}

\title{My Document}
\author{Your Name}
\date{\today}

\begin{document}

\maketitle

\section{Introduction}
This is a sample document demonstrating various \LaTeX{} features including mathematical equations, tables, and formatting. You can cite papers using the citation search feature, and they will be automatically added to your \texttt{cite.bib} file.

\section{Mathematical Expressions}
Here are some example equations:

\subsection{Quadratic Formula}
The quadratic formula is given by:
\[x = \frac{-b \pm \sqrt{b^2 - 4ac}}{2a}\]

\subsection{Matrix Example}
A sample matrix:
\[
\begin{pmatrix}
a_{11} & a_{12} & a_{13} \\
a_{21} & a_{22} & a_{23} \\
a_{31} & a_{32} & a_{33}
\end{pmatrix}
\]

\section{Tables}
Below is a sample table showing data:

\begin{table}[h]
\centering
\begin{tabular}{@{}lrrr@{}}
\toprule
\textbf{Category} & \textbf{Value 1} & \textbf{Value 2} & \textbf{Total} \\
\midrule
Item A & 10.5 & 20.3 & 30.8 \\
Item B & 15.2 & 25.1 & 40.3 \\
Item C & 22.1 & 12.7 & 34.8 \\
\bottomrule
\end{tabular}
\caption{Sample Data Table}
\label{tab:sample}
\end{table}

\section{Lists}
\subsection{Bullet Points}
\begin{itemize}
  \item First item
  \item Second item
  \item Third item with sub-items:
    \begin{itemize}
      \item Sub-item 1
      \item Sub-item 2
    \end{itemize}
\end{itemize}

\subsection{Numbered List}
\begin{enumerate}
  \item First step
  \item Second step
  \item Third step
\end{enumerate}

\section{Citations and References}
When you use the citation search feature, references will be automatically managed. For example, you might cite a paper like this: \cite{example2024}. All citations will be automatically added to your \texttt{cite.bib} file.

\section{Conclusion}
This template demonstrates basic \LaTeX{} features including mathematical formulas, tables, and lists. The bibliography system is now set up to work seamlessly with Bibby AI's citation management. Feel free to modify it for your needs!

\bibliographystyle{plain}
\bibliography{cite}

\end{document}\documentclass{article}
\usepackage[utf8]{inputenc}
\usepackage{booktabs}
\usepackage{amsmath}
\usepackage{natbib}

\title{My Document}
\author{Your Name}
\date{\today}

\begin{document}

\maketitle

\section{Introduction}
This is a sample document demonstrating various \LaTeX{} features including mathematical equations, tables, and formatting. You can cite papers using the citation search feature, and they will be automatically added to your \texttt{cite.bib} file.

\section{Mathematical Expressions}
Here are some example equations:

\subsection{Quadratic Formula}
The quadratic formula is given by:
\[x = \frac{-b \pm \sqrt{b^2 - 4ac}}{2a}\]

\subsection{Matrix Example}
A sample matrix:
\[
\begin{pmatrix}
a_{11} & a_{12} & a_{13} \\
a_{21} & a_{22} & a_{23} \\
a_{31} & a_{32} & a_{33}
\end{pmatrix}
\]

\section{Tables}
Below is a sample table showing data:

\begin{table}[h]
\centering
\begin{tabular}{@{}lrrr@{}}
\toprule
\textbf{Category} & \textbf{Value 1} & \textbf{Value 2} & \textbf{Total} \\
\midrule
Item A & 10.5 & 20.3 & 30.8 \\
Item B & 15.2 & 25.1 & 40.3 \\
Item C & 22.1 & 12.7 & 34.8 \\
\bottomrule
\end{tabular}
\caption{Sample Data Table}
\label{tab:sample}
\end{table}

\section{Lists}
\subsection{Bullet Points}
\begin{itemize}
  \item First item
  \item Second item
  \item Third item with sub-items:
    \begin{itemize}
      \item Sub-item 1
      \item Sub-item 2
    \end{itemize}
\end{itemize}

\subsection{Numbered List}
\begin{enumerate}
  \item First step
  \item Second step
  \item Third step
\end{enumerate}

\section{Citations and References}
When you use the citation search feature, references will be automatically managed. For example, you might cite a paper like this: \cite{example2024}. All citations will be automatically added to your \texttt{cite.bib} file.

\section{Conclusion}
This template demonstrates basic \LaTeX{} features including mathematical formulas, tables, and lists. The bibliography system is now set up to work seamlessly with Bibby AI's citation management. Feel free to modify it for your needs!

\bibliographystyle{plain}
\bibliography{cite}

\end{document}\documentclass{article}
\usepackage[utf8]{inputenc}
\usepackage{booktabs}
\usepackage{amsmath}
\usepackage{natbib}

\title{My Document}
\author{Your Name}
\date{\today}

\begin{document}

\maketitle

\section{Introduction}
This is a sample document demonstrating various \LaTeX{} features including mathematical equations, tables, and formatting. You can cite papers using the citation search feature, and they will be automatically added to your \texttt{cite.bib} file.

\section{Mathematical Expressions}
Here are some example equations:

\subsection{Quadratic Formula}
The quadratic formula is given by:
\[x = \frac{-b \pm \sqrt{b^2 - 4ac}}{2a}\]

\subsection{Matrix Example}
A sample matrix:
\[
\begin{pmatrix}
a_{11} & a_{12} & a_{13} \\
a_{21} & a_{22} & a_{23} \\
a_{31} & a_{32} & a_{33}
\end{pmatrix}
\]

\section{Tables}
Below is a sample table showing data:

\begin{table}[h]
\centering
\begin{tabular}{@{}lrrr@{}}
\toprule
\textbf{Category} & \textbf{Value 1} & \textbf{Value 2} & \textbf{Total} \\
\midrule
Item A & 10.5 & 20.3 & 30.8 \\
Item B & 15.2 & 25.1 & 40.3 \\
Item C & 22.1 & 12.7 & 34.8 \\
\bottomrule
\end{tabular}
\caption{Sample Data Table}
\label{tab:sample}
\end{table}

\section{Lists}
\subsection{Bullet Points}
\begin{itemize}
  \item First item
  \item Second item
  \item Third item with sub-items:
    \begin{itemize}
      \item Sub-item 1
      \item Sub-item 2
    \end{itemize}
\end{itemize}

\subsection{Numbered List}
\begin{enumerate}
  \item First step
  \item Second step
  \item Third step
\end{enumerate}

\section{Citations and References}
When you use the citation search feature, references will be automatically managed. For example, you might cite a paper like this: \cite{example2024}. All citations will be automatically added to your \texttt{cite.bib} file.

\section{Conclusion}
This template demonstrates basic \LaTeX{} features including mathematical formulas, tables, and lists. The bibliography system is now set up to work seamlessly with Bibby AI's citation management. Feel free to modify it for your needs!

\bibliographystyle{plain}
\bibliography{cite}

\end{document}\documentclass{article}
\usepackage[utf8]{inputenc}
\usepackage{booktabs}
\usepackage{amsmath}
\usepackage{natbib}

\title{My Document}
\author{Your Name}
\date{\today}

\begin{document}

\maketitle

\section{Introduction}
This is a sample document demonstrating various \LaTeX{} features including mathematical equations, tables, and formatting. You can cite papers using the citation search feature, and they will be automatically added to your \texttt{cite.bib} file.

\section{Mathematical Expressions}
Here are some example equations:

\subsection{Quadratic Formula}
The quadratic formula is given by:
\[x = \frac{-b \pm \sqrt{b^2 - 4ac}}{2a}\]

\subsection{Matrix Example}
A sample matrix:
\[
\begin{pmatrix}
a_{11} & a_{12} & a_{13} \\
a_{21} & a_{22} & a_{23} \\
a_{31} & a_{32} & a_{33}
\end{pmatrix}
\]

\section{Tables}
Below is a sample table showing data:

\begin{table}[h]
\centering
\begin{tabular}{@{}lrrr@{}}
\toprule
\textbf{Category} & \textbf{Value 1} & \textbf{Value 2} & \textbf{Total} \\
\midrule
Item A & 10.5 & 20.3 & 30.8 \\
Item B & 15.2 & 25.1 & 40.3 \\
Item C & 22.1 & 12.7 & 34.8 \\
\bottomrule
\end{tabular}
\caption{Sample Data Table}
\label{tab:sample}
\end{table}

\section{Lists}
\subsection{Bullet Points}
\begin{itemize}
  \item First item
  \item Second item
  \item Third item with sub-items:
    \begin{itemize}
      \item Sub-item 1
      \item Sub-item 2
    \end{itemize}
\end{itemize}

\subsection{Numbered List}
\begin{enumerate}
  \item First step
  \item Second step
  \item Third step
\end{enumerate}

\section{Citations and References}
When you use the citation search feature, references will be automatically managed. For example, you might cite a paper like this: \cite{example2024}. All citations will be automatically added to your \texttt{cite.bib} file.

\section{Conclusion}
This template demonstrates basic \LaTeX{} features including mathematical formulas, tables, and lists. The bibliography system is now set up to work seamlessly with Bibby AI's citation management. Feel free to modify it for your needs!

\bibliographystyle{plain}
\bibliography{cite}

\end{document}\documentclass{article}
\usepackage[utf8]{inputenc}
\usepackage{booktabs}
\usepackage{amsmath}
\usepackage{natbib}

\title{My Document}
\author{Your Name}
\date{\today}

\begin{document}

\maketitle

\section{Introduction}
This is a sample document demonstrating various \LaTeX{} features including mathematical equations, tables, and formatting. You can cite papers using the citation search feature, and they will be automatically added to your \texttt{cite.bib} file.

\section{Mathematical Expressions}
Here are some example equations:

\subsection{Quadratic Formula}
The quadratic formula is given by:
\[x = \frac{-b \pm \sqrt{b^2 - 4ac}}{2a}\]

\subsection{Matrix Example}
A sample matrix:
\[
\begin{pmatrix}
a_{11} & a_{12} & a_{13} \\
a_{21} & a_{22} & a_{23} \\
a_{31} & a_{32} & a_{33}
\end{pmatrix}
\]

\section{Tables}
Below is a sample table showing data:

\begin{table}[h]
\centering
\begin{tabular}{@{}lrrr@{}}
\toprule
\textbf{Category} & \textbf{Value 1} & \textbf{Value 2} & \textbf{Total} \\
\midrule
Item A & 10.5 & 20.3 & 30.8 \\
Item B & 15.2 & 25.1 & 40.3 \\
Item C & 22.1 & 12.7 & 34.8 \\
\bottomrule
\end{tabular}
\caption{Sample Data Table}
\label{tab:sample}
\end{table}

\section{Lists}
\subsection{Bullet Points}
\begin{itemize}
  \item First item
  \item Second item
  \item Third item with sub-items:
    \begin{itemize}
      \item Sub-item 1
      \item Sub-item 2
    \end{itemize}
\end{itemize}

\subsection{Numbered List}
\begin{enumerate}
  \item First step
  \item Second step
  \item Third step
\end{enumerate}

\section{Citations and References}
When you use the citation search feature, references will be automatically managed. For example, you might cite a paper like this: \cite{example2024}. All citations will be automatically added to your \texttt{cite.bib} file.

\section{Conclusion}
This template demonstrates basic \LaTeX{} features including mathematical formulas, tables, and lists. The bibliography system is now set up to work seamlessly with Bibby AI's citation management. Feel free to modify it for your needs!

\bibliographystyle{plain}
\bibliography{cite}

\end{document}\documentclass{article}
\usepackage[utf8]{inputenc}
\usepackage{booktabs}
\usepackage{amsmath}
\usepackage{natbib}

\title{My Document}
\author{Your Name}
\date{\today}

\begin{document}

\maketitle

\section{Introduction}
This is a sample document demonstrating various \LaTeX{} features including mathematical equations, tables, and formatting. You can cite papers using the citation search feature, and they will be automatically added to your \texttt{cite.bib} file.

\section{Mathematical Expressions}
Here are some example equations:

\subsection{Quadratic Formula}
The quadratic formula is given by:
\[x = \frac{-b \pm \sqrt{b^2 - 4ac}}{2a}\]

\subsection{Matrix Example}
A sample matrix:
\[
\begin{pmatrix}
a_{11} & a_{12} & a_{13} \\
a_{21} & a_{22} & a_{23} \\
a_{31} & a_{32} & a_{33}
\end{pmatrix}
\]

\section{Tables}
Below is a sample table showing data:

\begin{table}[h]
\centering
\begin{tabular}{@{}lrrr@{}}
\toprule
\textbf{Category} & \textbf{Value 1} & \textbf{Value 2} & \textbf{Total} \\
\midrule
Item A & 10.5 & 20.3 & 30.8 \\
Item B & 15.2 & 25.1 & 40.3 \\
Item C & 22.1 & 12.7 & 34.8 \\
\bottomrule
\end{tabular}
\caption{Sample Data Table}
\label{tab:sample}
\end{table}

\section{Lists}
\subsection{Bullet Points}
\begin{itemize}
  \item First item
  \item Second item
  \item Third item with sub-items:
    \begin{itemize}
      \item Sub-item 1
      \item Sub-item 2
    \end{itemize}
\end{itemize}

\subsection{Numbered List}
\begin{enumerate}
  \item First step
  \item Second step
  \item Third step
\end{enumerate}

\section{Citations and References}
When you use the citation search feature, references will be automatically managed. For example, you might cite a paper like this: \cite{example2024}. All citations will be automatically added to your \texttt{cite.bib} file.

\section{Conclusion}
This template demonstrates basic \LaTeX{} features including mathematical formulas, tables, and lists. The bibliography system is now set up to work seamlessly with Bibby AI's citation management. Feel free to modify it for your needs!

\bibliographystyle{plain}
\bibliography{cite}

\end{document}\documentclass{article}
\usepackage[utf8]{inputenc}
\usepackage{booktabs}
\usepackage{amsmath}
\usepackage{natbib}

\title{My Document}
\author{Your Name}
\date{\today}

\begin{document}

\maketitle

\section{Introduction}
This is a sample document demonstrating various \LaTeX{} features including mathematical equations, tables, and formatting. You can cite papers using the citation search feature, and they will be automatically added to your \texttt{cite.bib} file.

\section{Mathematical Expressions}
Here are some example equations:

\subsection{Quadratic Formula}
The quadratic formula is given by:
\[x = \frac{-b \pm \sqrt{b^2 - 4ac}}{2a}\]

\subsection{Matrix Example}
A sample matrix:
\[
\begin{pmatrix}
a_{11} & a_{12} & a_{13} \\
a_{21} & a_{22} & a_{23} \\
a_{31} & a_{32} & a_{33}
\end{pmatrix}
\]

\section{Tables}
Below is a sample table showing data:

\begin{table}[h]
\centering
\begin{tabular}{@{}lrrr@{}}
\toprule
\textbf{Category} & \textbf{Value 1} & \textbf{Value 2} & \textbf{Total} \\
\midrule
Item A & 10.5 & 20.3 & 30.8 \\
Item B & 15.2 & 25.1 & 40.3 \\
Item C & 22.1 & 12.7 & 34.8 \\
\bottomrule
\end{tabular}
\caption{Sample Data Table}
\label{tab:sample}
\end{table}

\section{Lists}
\subsection{Bullet Points}
\begin{itemize}
  \item First item
  \item Second item
  \item Third item with sub-items:
    \begin{itemize}
      \item Sub-item 1
      \item Sub-item 2
    \end{itemize}
\end{itemize}

\subsection{Numbered List}
\begin{enumerate}
  \item First step
  \item Second step
  \item Third step
\end{enumerate}

\section{Citations and References}
When you use the citation search feature, references will be automatically managed. For example, you might cite a paper like this: \cite{example2024}. All citations will be automatically added to your \texttt{cite.bib} file.

\section{Conclusion}
This template demonstrates basic \LaTeX{} features including mathematical formulas, tables, and lists. The bibliography system is now set up to work seamlessly with Bibby AI's citation management. Feel free to modify it for your needs!

\bibliographystyle{plain}
\bibliography{cite}

\end{document}\documentclass{article}
\usepackage[utf8]{inputenc}
\usepackage{booktabs}
\usepackage{amsmath}
\usepackage{natbib}

\title{My Document}
\author{Your Name}
\date{\today}

\begin{document}

\maketitle

\section{Introduction}
This is a sample document demonstrating various \LaTeX{} features including mathematical equations, tables, and formatting. You can cite papers using the citation search feature, and they will be automatically added to your \texttt{cite.bib} file.

\section{Mathematical Expressions}
Here are some example equations:

\subsection{Quadratic Formula}
The quadratic formula is given by:
\[x = \frac{-b \pm \sqrt{b^2 - 4ac}}{2a}\]

\subsection{Matrix Example}
A sample matrix:
\[
\begin{pmatrix}
a_{11} & a_{12} & a_{13} \\
a_{21} & a_{22} & a_{23} \\
a_{31} & a_{32} & a_{33}
\end{pmatrix}
\]

\section{Tables}
Below is a sample table showing data:

\begin{table}[h]
\centering
\begin{tabular}{@{}lrrr@{}}
\toprule
\textbf{Category} & \textbf{Value 1} & \textbf{Value 2} & \textbf{Total} \\
\midrule
Item A & 10.5 & 20.3 & 30.8 \\
Item B & 15.2 & 25.1 & 40.3 \\
Item C & 22.1 & 12.7 & 34.8 \\
\bottomrule
\end{tabular}
\caption{Sample Data Table}
\label{tab:sample}
\end{table}

\section{Lists}
\subsection{Bullet Points}
\begin{itemize}
  \item First item
  \item Second item
  \item Third item with sub-items:
    \begin{itemize}
      \item Sub-item 1
      \item Sub-item 2
    \end{itemize}
\end{itemize}

\subsection{Numbered List}
\begin{enumerate}
  \item First step
  \item Second step
  \item Third step
\end{enumerate}

\section{Citations and References}
When you use the citation search feature, references will be automatically managed. For example, you might cite a paper like this: \cite{example2024}. All citations will be automatically added to your \texttt{cite.bib} file.

\section{Conclusion}
This template demonstrates basic \LaTeX{} features including mathematical formulas, tables, and lists. The bibliography system is now set up to work seamlessly with Bibby AI's citation management. Feel free to modify it for your needs!

\bibliographystyle{plain}
\bibliography{cite}

\end{document}\documentclass{article}
\usepackage[utf8]{inputenc}
\usepackage{booktabs}
\usepackage{amsmath}
\usepackage{natbib}

\title{My Document}
\author{Your Name}
\date{\today}

\begin{document}

\maketitle

\section{Introduction}
This is a sample document demonstrating various \LaTeX{} features including mathematical equations, tables, and formatting. You can cite papers using the citation search feature, and they will be automatically added to your \texttt{cite.bib} file.

\section{Mathematical Expressions}
Here are some example equations:

\subsection{Quadratic Formula}
The quadratic formula is given by:
\[x = \frac{-b \pm \sqrt{b^2 - 4ac}}{2a}\]

\subsection{Matrix Example}
A sample matrix:
\[
\begin{pmatrix}
a_{11} & a_{12} & a_{13} \\
a_{21} & a_{22} & a_{23} \\
a_{31} & a_{32} & a_{33}
\end{pmatrix}
\]

\section{Tables}
Below is a sample table showing data:

\begin{table}[h]
\centering
\begin{tabular}{@{}lrrr@{}}
\toprule
\textbf{Category} & \textbf{Value 1} & \textbf{Value 2} & \textbf{Total} \\
\midrule
Item A & 10.5 & 20.3 & 30.8 \\
Item B & 15.2 & 25.1 & 40.3 \\
Item C & 22.1 & 12.7 & 34.8 \\
\bottomrule
\end{tabular}
\caption{Sample Data Table}
\label{tab:sample}
\end{table}

\section{Lists}
\subsection{Bullet Points}
\begin{itemize}
  \item First item
  \item Second item
  \item Third item with sub-items:
    \begin{itemize}
      \item Sub-item 1
      \item Sub-item 2
    \end{itemize}
\end{itemize}

\subsection{Numbered List}
\begin{enumerate}
  \item First step
  \item Second step
  \item Third step
\end{enumerate}

\section{Citations and References}
When you use the citation search feature, references will be automatically managed. For example, you might cite a paper like this: \cite{example2024}. All citations will be automatically added to your \texttt{cite.bib} file.

\section{Conclusion}
This template demonstrates basic \LaTeX{} features including mathematical formulas, tables, and lists. The bibliography system is now set up to work seamlessly with Bibby AI's citation management. Feel free to modify it for your needs!

\bibliographystyle{plain}
\bibliography{cite}

\end{document}\documentclass{article}
\usepackage[utf8]{inputenc}
\usepackage{booktabs}
\usepackage{amsmath}
\usepackage{natbib}

\title{My Document}
\author{Your Name}
\date{\today}

\begin{document}

\maketitle

\section{Introduction}
This is a sample document demonstrating various \LaTeX{} features including mathematical equations, tables, and formatting. You can cite papers using the citation search feature, and they will be automatically added to your \texttt{cite.bib} file.

\section{Mathematical Expressions}
Here are some example equations:

\subsection{Quadratic Formula}
The quadratic formula is given by:
\[x = \frac{-b \pm \sqrt{b^2 - 4ac}}{2a}\]

\subsection{Matrix Example}
A sample matrix:
\[
\begin{pmatrix}
a_{11} & a_{12} & a_{13} \\
a_{21} & a_{22} & a_{23} \\
a_{31} & a_{32} & a_{33}
\end{pmatrix}
\]

\section{Tables}
Below is a sample table showing data:

\begin{table}[h]
\centering
\begin{tabular}{@{}lrrr@{}}
\toprule
\textbf{Category} & \textbf{Value 1} & \textbf{Value 2} & \textbf{Total} \\
\midrule
Item A & 10.5 & 20.3 & 30.8 \\
Item B & 15.2 & 25.1 & 40.3 \\
Item C & 22.1 & 12.7 & 34.8 \\
\bottomrule
\end{tabular}
\caption{Sample Data Table}
\label{tab:sample}
\end{table}

\section{Lists}
\subsection{Bullet Points}
\begin{itemize}
  \item First item
  \item Second item
  \item Third item with sub-items:
    \begin{itemize}
      \item Sub-item 1
      \item Sub-item 2
    \end{itemize}
\end{itemize}

\subsection{Numbered List}
\begin{enumerate}
  \item First step
  \item Second step
  \item Third step
\end{enumerate}

\section{Citations and References}
When you use the citation search feature, references will be automatically managed. For example, you might cite a paper like this: \cite{example2024}. All citations will be automatically added to your \texttt{cite.bib} file.

\section{Conclusion}
This template demonstrates basic \LaTeX{} features including mathematical formulas, tables, and lists. The bibliography system is now set up to work seamlessly with Bibby AI's citation management. Feel free to modify it for your needs!

\bibliographystyle{plain}
\bibliography{cite}

\end{document}\documentclass{article}
\usepackage[utf8]{inputenc}
\usepackage{booktabs}
\usepackage{amsmath}
\usepackage{natbib}

\title{My Document}
\author{Your Name}
\date{\today}

\begin{document}

\maketitle

\section{Introduction}
This is a sample document demonstrating various \LaTeX{} features including mathematical equations, tables, and formatting. You can cite papers using the citation search feature, and they will be automatically added to your \texttt{cite.bib} file.

\section{Mathematical Expressions}
Here are some example equations:

\subsection{Quadratic Formula}
The quadratic formula is given by:
\[x = \frac{-b \pm \sqrt{b^2 - 4ac}}{2a}\]

\subsection{Matrix Example}
A sample matrix:
\[
\begin{pmatrix}
a_{11} & a_{12} & a_{13} \\
a_{21} & a_{22} & a_{23} \\
a_{31} & a_{32} & a_{33}
\end{pmatrix}
\]

\section{Tables}
Below is a sample table showing data:

\begin{table}[h]
\centering
\begin{tabular}{@{}lrrr@{}}
\toprule
\textbf{Category} & \textbf{Value 1} & \textbf{Value 2} & \textbf{Total} \\
\midrule
Item A & 10.5 & 20.3 & 30.8 \\
Item B & 15.2 & 25.1 & 40.3 \\
Item C & 22.1 & 12.7 & 34.8 \\
\bottomrule
\end{tabular}
\caption{Sample Data Table}
\label{tab:sample}
\end{table}

\section{Lists}
\subsection{Bullet Points}
\begin{itemize}
  \item First item
  \item Second item
  \item Third item with sub-items:
    \begin{itemize}
      \item Sub-item 1
      \item Sub-item 2
    \end{itemize}
\end{itemize}

\subsection{Numbered List}
\begin{enumerate}
  \item First step
  \item Second step
  \item Third step
\end{enumerate}

\section{Citations and References}
When you use the citation search feature, references will be automatically managed. For example, you might cite a paper like this: \cite{example2024}. All citations will be automatically added to your \texttt{cite.bib} file.

\section{Conclusion}
This template demonstrates basic \LaTeX{} features including mathematical formulas, tables, and lists. The bibliography system is now set up to work seamlessly with Bibby AI's citation management. Feel free to modify it for your needs!

\bibliographystyle{plain}
\bibliography{cite}

\end{document}\documentclass{article}
\usepackage[utf8]{inputenc}
\usepackage{booktabs}
\usepackage{amsmath}
\usepackage{natbib}

\title{My Document}
\author{Your Name}
\date{\today}

\begin{document}

\maketitle

\section{Introduction}
This is a sample document demonstrating various \LaTeX{} features including mathematical equations, tables, and formatting. You can cite papers using the citation search feature, and they will be automatically added to your \texttt{cite.bib} file.

\section{Mathematical Expressions}
Here are some example equations:

\subsection{Quadratic Formula}
The quadratic formula is given by:
\[x = \frac{-b \pm \sqrt{b^2 - 4ac}}{2a}\]

\subsection{Matrix Example}
A sample matrix:
\[
\begin{pmatrix}
a_{11} & a_{12} & a_{13} \\
a_{21} & a_{22} & a_{23} \\
a_{31} & a_{32} & a_{33}
\end{pmatrix}
\]

\section{Tables}
Below is a sample table showing data:

\begin{table}[h]
\centering
\begin{tabular}{@{}lrrr@{}}
\toprule
\textbf{Category} & \textbf{Value 1} & \textbf{Value 2} & \textbf{Total} \\
\midrule
Item A & 10.5 & 20.3 & 30.8 \\
Item B & 15.2 & 25.1 & 40.3 \\
Item C & 22.1 & 12.7 & 34.8 \\
\bottomrule
\end{tabular}
\caption{Sample Data Table}
\label{tab:sample}
\end{table}

\section{Lists}
\subsection{Bullet Points}
\begin{itemize}
  \item First item
  \item Second item
  \item Third item with sub-items:
    \begin{itemize}
      \item Sub-item 1
      \item Sub-item 2
    \end{itemize}
\end{itemize}

\subsection{Numbered List}
\begin{enumerate}
  \item First step
  \item Second step
  \item Third step
\end{enumerate}

\section{Citations and References}
When you use the citation search feature, references will be automatically managed. For example, you might cite a paper like this: \cite{example2024}. All citations will be automatically added to your \texttt{cite.bib} file.

\section{Conclusion}
This template demonstrates basic \LaTeX{} features including mathematical formulas, tables, and lists. The bibliography system is now set up to work seamlessly with Bibby AI's citation management. Feel free to modify it for your needs!

\bibliographystyle{plain}
\bibliography{cite}

\end{document}\documentclass{article}
\usepackage[utf8]{inputenc}
\usepackage{booktabs}
\usepackage{amsmath}
\usepackage{natbib}

\title{My Document}
\author{Your Name}
\date{\today}

\begin{document}

\maketitle

\section{Introduction}
This is a sample document demonstrating various \LaTeX{} features including mathematical equations, tables, and formatting. You can cite papers using the citation search feature, and they will be automatically added to your \texttt{cite.bib} file.

\section{Mathematical Expressions}
Here are some example equations:

\subsection{Quadratic Formula}
The quadratic formula is given by:
\[x = \frac{-b \pm \sqrt{b^2 - 4ac}}{2a}\]

\subsection{Matrix Example}
A sample matrix:
\[
\begin{pmatrix}
a_{11} & a_{12} & a_{13} \\
a_{21} & a_{22} & a_{23} \\
a_{31} & a_{32} & a_{33}
\end{pmatrix}
\]

\section{Tables}
Below is a sample table showing data:

\begin{table}[h]
\centering
\begin{tabular}{@{}lrrr@{}}
\toprule
\textbf{Category} & \textbf{Value 1} & \textbf{Value 2} & \textbf{Total} \\
\midrule
Item A & 10.5 & 20.3 & 30.8 \\
Item B & 15.2 & 25.1 & 40.3 \\
Item C & 22.1 & 12.7 & 34.8 \\
\bottomrule
\end{tabular}
\caption{Sample Data Table}
\label{tab:sample}
\end{table}

\section{Lists}
\subsection{Bullet Points}
\begin{itemize}
  \item First item
  \item Second item
  \item Third item with sub-items:
    \begin{itemize}
      \item Sub-item 1
      \item Sub-item 2
    \end{itemize}
\end{itemize}

\subsection{Numbered List}
\begin{enumerate}
  \item First step
  \item Second step
  \item Third step
\end{enumerate}

\section{Citations and References}
When you use the citation search feature, references will be automatically managed. For example, you might cite a paper like this: \cite{example2024}. All citations will be automatically added to your \texttt{cite.bib} file.

\section{Conclusion}
This template demonstrates basic \LaTeX{} features including mathematical formulas, tables, and lists. The bibliography system is now set up to work seamlessly with Bibby AI's citation management. Feel free to modify it for your needs!

\bibliographystyle{plain}
\bibliography{cite}

\end{document}\documentclass{article}
\usepackage[utf8]{inputenc}
\usepackage{booktabs}
\usepackage{amsmath}
\usepackage{natbib}

\title{My Document}
\author{Your Name}
\date{\today}

\begin{document}

\maketitle

\section{Introduction}
This is a sample document demonstrating various \LaTeX{} features including mathematical equations, tables, and formatting. You can cite papers using the citation search feature, and they will be automatically added to your \texttt{cite.bib} file.

\section{Mathematical Expressions}
Here are some example equations:

\subsection{Quadratic Formula}
The quadratic formula is given by:
\[x = \frac{-b \pm \sqrt{b^2 - 4ac}}{2a}\]

\subsection{Matrix Example}
A sample matrix:
\[
\begin{pmatrix}
a_{11} & a_{12} & a_{13} \\
a_{21} & a_{22} & a_{23} \\
a_{31} & a_{32} & a_{33}
\end{pmatrix}
\]

\section{Tables}
Below is a sample table showing data:

\begin{table}[h]
\centering
\begin{tabular}{@{}lrrr@{}}
\toprule
\textbf{Category} & \textbf{Value 1} & \textbf{Value 2} & \textbf{Total} \\
\midrule
Item A & 10.5 & 20.3 & 30.8 \\
Item B & 15.2 & 25.1 & 40.3 \\
Item C & 22.1 & 12.7 & 34.8 \\
\bottomrule
\end{tabular}
\caption{Sample Data Table}
\label{tab:sample}
\end{table}

\section{Lists}
\subsection{Bullet Points}
\begin{itemize}
  \item First item
  \item Second item
  \item Third item with sub-items:
    \begin{itemize}
      \item Sub-item 1
      \item Sub-item 2
    \end{itemize}
\end{itemize}

\subsection{Numbered List}
\begin{enumerate}
  \item First step
  \item Second step
  \item Third step
\end{enumerate}

\section{Citations and References}
When you use the citation search feature, references will be automatically managed. For example, you might cite a paper like this: \cite{example2024}. All citations will be automatically added to your \texttt{cite.bib} file.

\section{Conclusion}
This template demonstrates basic \LaTeX{} features including mathematical formulas, tables, and lists. The bibliography system is now set up to work seamlessly with Bibby AI's citation management. Feel free to modify it for your needs!

\bibliographystyle{plain}
\bibliography{cite}

\end{document}\documentclass{article}
\usepackage[utf8]{inputenc}
\usepackage{booktabs}
\usepackage{amsmath}
\usepackage{natbib}

\title{My Document}
\author{Your Name}
\date{\today}

\begin{document}

\maketitle

\section{Introduction}
This is a sample document demonstrating various \LaTeX{} features including mathematical equations, tables, and formatting. You can cite papers using the citation search feature, and they will be automatically added to your \texttt{cite.bib} file.

\section{Mathematical Expressions}
Here are some example equations:

\subsection{Quadratic Formula}
The quadratic formula is given by:
\[x = \frac{-b \pm \sqrt{b^2 - 4ac}}{2a}\]

\subsection{Matrix Example}
A sample matrix:
\[
\begin{pmatrix}
a_{11} & a_{12} & a_{13} \\
a_{21} & a_{22} & a_{23} \\
a_{31} & a_{32} & a_{33}
\end{pmatrix}
\]

\section{Tables}
Below is a sample table showing data:

\begin{table}[h]
\centering
\begin{tabular}{@{}lrrr@{}}
\toprule
\textbf{Category} & \textbf{Value 1} & \textbf{Value 2} & \textbf{Total} \\
\midrule
Item A & 10.5 & 20.3 & 30.8 \\
Item B & 15.2 & 25.1 & 40.3 \\
Item C & 22.1 & 12.7 & 34.8 \\
\bottomrule
\end{tabular}
\caption{Sample Data Table}
\label{tab:sample}
\end{table}

\section{Lists}
\subsection{Bullet Points}
\begin{itemize}
  \item First item
  \item Second item
  \item Third item with sub-items:
    \begin{itemize}
      \item Sub-item 1
      \item Sub-item 2
    \end{itemize}
\end{itemize}

\subsection{Numbered List}
\begin{enumerate}
  \item First step
  \item Second step
  \item Third step
\end{enumerate}

\section{Citations and References}
When you use the citation search feature, references will be automatically managed. For example, you might cite a paper like this: \cite{example2024}. All citations will be automatically added to your \texttt{cite.bib} file.

\section{Conclusion}
This template demonstrates basic \LaTeX{} features including mathematical formulas, tables, and lists. The bibliography system is now set up to work seamlessly with Bibby AI's citation management. Feel free to modify it for your needs!

\bibliographystyle{plain}
\bibliography{cite}

\end{document}\documentclass{article}
\usepackage[utf8]{inputenc}
\usepackage{booktabs}
\usepackage{amsmath}
\usepackage{natbib}

\title{My Document}
\author{Your Name}
\date{\today}

\begin{document}

\maketitle

\section{Introduction}
This is a sample document demonstrating various \LaTeX{} features including mathematical equations, tables, and formatting. You can cite papers using the citation search feature, and they will be automatically added to your \texttt{cite.bib} file.

\section{Mathematical Expressions}
Here are some example equations:

\subsection{Quadratic Formula}
The quadratic formula is given by:
\[x = \frac{-b \pm \sqrt{b^2 - 4ac}}{2a}\]

\subsection{Matrix Example}
A sample matrix:
\[
\begin{pmatrix}
a_{11} & a_{12} & a_{13} \\
a_{21} & a_{22} & a_{23} \\
a_{31} & a_{32} & a_{33}
\end{pmatrix}
\]

\section{Tables}
Below is a sample table showing data:

\begin{table}[h]
\centering
\begin{tabular}{@{}lrrr@{}}
\toprule
\textbf{Category} & \textbf{Value 1} & \textbf{Value 2} & \textbf{Total} \\
\midrule
Item A & 10.5 & 20.3 & 30.8 \\
Item B & 15.2 & 25.1 & 40.3 \\
Item C & 22.1 & 12.7 & 34.8 \\
\bottomrule
\end{tabular}
\caption{Sample Data Table}
\label{tab:sample}
\end{table}

\section{Lists}
\subsection{Bullet Points}
\begin{itemize}
  \item First item
  \item Second item
  \item Third item with sub-items:
    \begin{itemize}
      \item Sub-item 1
      \item Sub-item 2
    \end{itemize}
\end{itemize}

\subsection{Numbered List}
\begin{enumerate}
  \item First step
  \item Second step
  \item Third step
\end{enumerate}

\section{Citations and References}
When you use the citation search feature, references will be automatically managed. For example, you might cite a paper like this: \cite{example2024}. All citations will be automatically added to your \texttt{cite.bib} file.

\section{Conclusion}
This template demonstrates basic \LaTeX{} features including mathematical formulas, tables, and lists. The bibliography system is now set up to work seamlessly with Bibby AI's citation management. Feel free to modify it for your needs!

\bibliographystyle{plain}
\bibliography{cite}

\end{document}\documentclass{article}
\usepackage[utf8]{inputenc}
\usepackage{booktabs}
\usepackage{amsmath}
\usepackage{natbib}

\title{My Document}
\author{Your Name}
\date{\today}

\begin{document}

\maketitle

\section{Introduction}
This is a sample document demonstrating various \LaTeX{} features including mathematical equations, tables, and formatting. You can cite papers using the citation search feature, and they will be automatically added to your \texttt{cite.bib} file.

\section{Mathematical Expressions}
Here are some example equations:

\subsection{Quadratic Formula}
The quadratic formula is given by:
\[x = \frac{-b \pm \sqrt{b^2 - 4ac}}{2a}\]

\subsection{Matrix Example}
A sample matrix:
\[
\begin{pmatrix}
a_{11} & a_{12} & a_{13} \\
a_{21} & a_{22} & a_{23} \\
a_{31} & a_{32} & a_{33}
\end{pmatrix}
\]

\section{Tables}
Below is a sample table showing data:

\begin{table}[h]
\centering
\begin{tabular}{@{}lrrr@{}}
\toprule
\textbf{Category} & \textbf{Value 1} & \textbf{Value 2} & \textbf{Total} \\
\midrule
Item A & 10.5 & 20.3 & 30.8 \\
Item B & 15.2 & 25.1 & 40.3 \\
Item C & 22.1 & 12.7 & 34.8 \\
\bottomrule
\end{tabular}
\caption{Sample Data Table}
\label{tab:sample}
\end{table}

\section{Lists}
\subsection{Bullet Points}
\begin{itemize}
  \item First item
  \item Second item
  \item Third item with sub-items:
    \begin{itemize}
      \item Sub-item 1
      \item Sub-item 2
    \end{itemize}
\end{itemize}

\subsection{Numbered List}
\begin{enumerate}
  \item First step
  \item Second step
  \item Third step
\end{enumerate}

\section{Citations and References}
When you use the citation search feature, references will be automatically managed. For example, you might cite a paper like this: \cite{example2024}. All citations will be automatically added to your \texttt{cite.bib} file.

\section{Conclusion}
This template demonstrates basic \LaTeX{} features including mathematical formulas, tables, and lists. The bibliography system is now set up to work seamlessly with Bibby AI's citation management. Feel free to modify it for your needs!

\bibliographystyle{plain}
\bibliography{cite}

\end{document}\documentclass{article}
\usepackage[utf8]{inputenc}
\usepackage{booktabs}
\usepackage{amsmath}
\usepackage{natbib}

\title{My Document}
\author{Your Name}
\date{\today}

\begin{document}

\maketitle

\section{Introduction}
This is a sample document demonstrating various \LaTeX{} features including mathematical equations, tables, and formatting. You can cite papers using the citation search feature, and they will be automatically added to your \texttt{cite.bib} file.

\section{Mathematical Expressions}
Here are some example equations:

\subsection{Quadratic Formula}
The quadratic formula is given by:
\[x = \frac{-b \pm \sqrt{b^2 - 4ac}}{2a}\]

\subsection{Matrix Example}
A sample matrix:
\[
\begin{pmatrix}
a_{11} & a_{12} & a_{13} \\
a_{21} & a_{22} & a_{23} \\
a_{31} & a_{32} & a_{33}
\end{pmatrix}
\]

\section{Tables}
Below is a sample table showing data:

\begin{table}[h]
\centering
\begin{tabular}{@{}lrrr@{}}
\toprule
\textbf{Category} & \textbf{Value 1} & \textbf{Value 2} & \textbf{Total} \\
\midrule
Item A & 10.5 & 20.3 & 30.8 \\
Item B & 15.2 & 25.1 & 40.3 \\
Item C & 22.1 & 12.7 & 34.8 \\
\bottomrule
\end{tabular}
\caption{Sample Data Table}
\label{tab:sample}
\end{table}

\section{Lists}
\subsection{Bullet Points}
\begin{itemize}
  \item First item
  \item Second item
  \item Third item with sub-items:
    \begin{itemize}
      \item Sub-item 1
      \item Sub-item 2
    \end{itemize}
\end{itemize}

\subsection{Numbered List}
\begin{enumerate}
  \item First step
  \item Second step
  \item Third step
\end{enumerate}

\section{Citations and References}
When you use the citation search feature, references will be automatically managed. For example, you might cite a paper like this: \cite{example2024}. All citations will be automatically added to your \texttt{cite.bib} file.

\section{Conclusion}
This template demonstrates basic \LaTeX{} features including mathematical formulas, tables, and lists. The bibliography system is now set up to work seamlessly with Bibby AI's citation management. Feel free to modify it for your needs!

\bibliographystyle{plain}
\bibliography{cite}

\end{document}\documentclass{article}
\usepackage[utf8]{inputenc}
\usepackage{booktabs}
\usepackage{amsmath}
\usepackage{natbib}

\title{My Document}
\author{Your Name}
\date{\today}

\begin{document}

\maketitle

\section{Introduction}
This is a sample document demonstrating various \LaTeX{} features including mathematical equations, tables, and formatting. You can cite papers using the citation search feature, and they will be automatically added to your \texttt{cite.bib} file.

\section{Mathematical Expressions}
Here are some example equations:

\subsection{Quadratic Formula}
The quadratic formula is given by:
\[x = \frac{-b \pm \sqrt{b^2 - 4ac}}{2a}\]

\subsection{Matrix Example}
A sample matrix:
\[
\begin{pmatrix}
a_{11} & a_{12} & a_{13} \\
a_{21} & a_{22} & a_{23} \\
a_{31} & a_{32} & a_{33}
\end{pmatrix}
\]

\section{Tables}
Below is a sample table showing data:

\begin{table}[h]
\centering
\begin{tabular}{@{}lrrr@{}}
\toprule
\textbf{Category} & \textbf{Value 1} & \textbf{Value 2} & \textbf{Total} \\
\midrule
Item A & 10.5 & 20.3 & 30.8 \\
Item B & 15.2 & 25.1 & 40.3 \\
Item C & 22.1 & 12.7 & 34.8 \\
\bottomrule
\end{tabular}
\caption{Sample Data Table}
\label{tab:sample}
\end{table}

\section{Lists}
\subsection{Bullet Points}
\begin{itemize}
  \item First item
  \item Second item
  \item Third item with sub-items:
    \begin{itemize}
      \item Sub-item 1
      \item Sub-item 2
    \end{itemize}
\end{itemize}

\subsection{Numbered List}
\begin{enumerate}
  \item First step
  \item Second step
  \item Third step
\end{enumerate}

\section{Citations and References}
When you use the citation search feature, references will be automatically managed. For example, you might cite a paper like this: \cite{example2024}. All citations will be automatically added to your \texttt{cite.bib} file.

\section{Conclusion}
This template demonstrates basic \LaTeX{} features including mathematical formulas, tables, and lists. The bibliography system is now set up to work seamlessly with Bibby AI's citation management. Feel free to modify it for your needs!

\bibliographystyle{plain}
\bibliography{cite}

\end{document}\documentclass{article}
\usepackage[utf8]{inputenc}
\usepackage{booktabs}
\usepackage{amsmath}
\usepackage{natbib}

\title{My Document}
\author{Your Name}
\date{\today}

\begin{document}

\maketitle

\section{Introduction}
This is a sample document demonstrating various \LaTeX{} features including mathematical equations, tables, and formatting. You can cite papers using the citation search feature, and they will be automatically added to your \texttt{cite.bib} file.

\section{Mathematical Expressions}
Here are some example equations:

\subsection{Quadratic Formula}
The quadratic formula is given by:
\[x = \frac{-b \pm \sqrt{b^2 - 4ac}}{2a}\]

\subsection{Matrix Example}
A sample matrix:
\[
\begin{pmatrix}
a_{11} & a_{12} & a_{13} \\
a_{21} & a_{22} & a_{23} \\
a_{31} & a_{32} & a_{33}
\end{pmatrix}
\]

\section{Tables}
Below is a sample table showing data:

\begin{table}[h]
\centering
\begin{tabular}{@{}lrrr@{}}
\toprule
\textbf{Category} & \textbf{Value 1} & \textbf{Value 2} & \textbf{Total} \\
\midrule
Item A & 10.5 & 20.3 & 30.8 \\
Item B & 15.2 & 25.1 & 40.3 \\
Item C & 22.1 & 12.7 & 34.8 \\
\bottomrule
\end{tabular}
\caption{Sample Data Table}
\label{tab:sample}
\end{table}

\section{Lists}
\subsection{Bullet Points}
\begin{itemize}
  \item First item
  \item Second item
  \item Third item with sub-items:
    \begin{itemize}
      \item Sub-item 1
      \item Sub-item 2
    \end{itemize}
\end{itemize}

\subsection{Numbered List}
\begin{enumerate}
  \item First step
  \item Second step
  \item Third step
\end{enumerate}

\section{Citations and References}
When you use the citation search feature, references will be automatically managed. For example, you might cite a paper like this: \cite{example2024}. All citations will be automatically added to your \texttt{cite.bib} file.

\section{Conclusion}
This template demonstrates basic \LaTeX{} features including mathematical formulas, tables, and lists. The bibliography system is now set up to work seamlessly with Bibby AI's citation management. Feel free to modify it for your needs!

\bibliographystyle{plain}
\bibliography{cite}

\end{document}\documentclass{article}
\usepackage[utf8]{inputenc}
\usepackage{booktabs}
\usepackage{amsmath}
\usepackage{natbib}

\title{My Document}
\author{Your Name}
\date{\today}

\begin{document}

\maketitle

\section{Introduction}
This is a sample document demonstrating various \LaTeX{} features including mathematical equations, tables, and formatting. You can cite papers using the citation search feature, and they will be automatically added to your \texttt{cite.bib} file.

\section{Mathematical Expressions}
Here are some example equations:

\subsection{Quadratic Formula}
The quadratic formula is given by:
\[x = \frac{-b \pm \sqrt{b^2 - 4ac}}{2a}\]

\subsection{Matrix Example}
A sample matrix:
\[
\begin{pmatrix}
a_{11} & a_{12} & a_{13} \\
a_{21} & a_{22} & a_{23} \\
a_{31} & a_{32} & a_{33}
\end{pmatrix}
\]

\section{Tables}
Below is a sample table showing data:

\begin{table}[h]
\centering
\begin{tabular}{@{}lrrr@{}}
\toprule
\textbf{Category} & \textbf{Value 1} & \textbf{Value 2} & \textbf{Total} \\
\midrule
Item A & 10.5 & 20.3 & 30.8 \\
Item B & 15.2 & 25.1 & 40.3 \\
Item C & 22.1 & 12.7 & 34.8 \\
\bottomrule
\end{tabular}
\caption{Sample Data Table}
\label{tab:sample}
\end{table}

\section{Lists}
\subsection{Bullet Points}
\begin{itemize}
  \item First item
  \item Second item
  \item Third item with sub-items:
    \begin{itemize}
      \item Sub-item 1
      \item Sub-item 2
    \end{itemize}
\end{itemize}

\subsection{Numbered List}
\begin{enumerate}
  \item First step
  \item Second step
  \item Third step
\end{enumerate}

\section{Citations and References}
When you use the citation search feature, references will be automatically managed. For example, you might cite a paper like this: \cite{example2024}. All citations will be automatically added to your \texttt{cite.bib} file.

\section{Conclusion}
This template demonstrates basic \LaTeX{} features including mathematical formulas, tables, and lists. The bibliography system is now set up to work seamlessly with Bibby AI's citation management. Feel free to modify it for your needs!

\bibliographystyle{plain}
\bibliography{cite}

\end{document}\documentclass{article}
\usepackage[utf8]{inputenc}
\usepackage{booktabs}
\usepackage{amsmath}
\usepackage{natbib}

\title{My Document}
\author{Your Name}
\date{\today}

\begin{document}

\maketitle

\section{Introduction}
This is a sample document demonstrating various \LaTeX{} features including mathematical equations, tables, and formatting. You can cite papers using the citation search feature, and they will be automatically added to your \texttt{cite.bib} file.

\section{Mathematical Expressions}
Here are some example equations:

\subsection{Quadratic Formula}
The quadratic formula is given by:
\[x = \frac{-b \pm \sqrt{b^2 - 4ac}}{2a}\]

\subsection{Matrix Example}
A sample matrix:
\[
\begin{pmatrix}
a_{11} & a_{12} & a_{13} \\
a_{21} & a_{22} & a_{23} \\
a_{31} & a_{32} & a_{33}
\end{pmatrix}
\]

\section{Tables}
Below is a sample table showing data:

\begin{table}[h]
\centering
\begin{tabular}{@{}lrrr@{}}
\toprule
\textbf{Category} & \textbf{Value 1} & \textbf{Value 2} & \textbf{Total} \\
\midrule
Item A & 10.5 & 20.3 & 30.8 \\
Item B & 15.2 & 25.1 & 40.3 \\
Item C & 22.1 & 12.7 & 34.8 \\
\bottomrule
\end{tabular}
\caption{Sample Data Table}
\label{tab:sample}
\end{table}

\section{Lists}
\subsection{Bullet Points}
\begin{itemize}
  \item First item
  \item Second item
  \item Third item with sub-items:
    \begin{itemize}
      \item Sub-item 1
      \item Sub-item 2
    \end{itemize}
\end{itemize}

\subsection{Numbered List}
\begin{enumerate}
  \item First step
  \item Second step
  \item Third step
\end{enumerate}

\section{Citations and References}
When you use the citation search feature, references will be automatically managed. For example, you might cite a paper like this: \cite{example2024}. All citations will be automatically added to your \texttt{cite.bib} file.

\section{Conclusion}
This template demonstrates basic \LaTeX{} features including mathematical formulas, tables, and lists. The bibliography system is now set up to work seamlessly with Bibby AI's citation management. Feel free to modify it for your needs!

\bibliographystyle{plain}
\bibliography{cite}

\end{document}\documentclass{article}
\usepackage[utf8]{inputenc}
\usepackage{booktabs}
\usepackage{amsmath}
\usepackage{natbib}

\title{My Document}
\author{Your Name}
\date{\today}

\begin{document}

\maketitle

\section{Introduction}
This is a sample document demonstrating various \LaTeX{} features including mathematical equations, tables, and formatting. You can cite papers using the citation search feature, and they will be automatically added to your \texttt{cite.bib} file.

\section{Mathematical Expressions}
Here are some example equations:

\subsection{Quadratic Formula}
The quadratic formula is given by:
\[x = \frac{-b \pm \sqrt{b^2 - 4ac}}{2a}\]

\subsection{Matrix Example}
A sample matrix:
\[
\begin{pmatrix}
a_{11} & a_{12} & a_{13} \\
a_{21} & a_{22} & a_{23} \\
a_{31} & a_{32} & a_{33}
\end{pmatrix}
\]

\section{Tables}
Below is a sample table showing data:

\begin{table}[h]
\centering
\begin{tabular}{@{}lrrr@{}}
\toprule
\textbf{Category} & \textbf{Value 1} & \textbf{Value 2} & \textbf{Total} \\
\midrule
Item A & 10.5 & 20.3 & 30.8 \\
Item B & 15.2 & 25.1 & 40.3 \\
Item C & 22.1 & 12.7 & 34.8 \\
\bottomrule
\end{tabular}
\caption{Sample Data Table}
\label{tab:sample}
\end{table}

\section{Lists}
\subsection{Bullet Points}
\begin{itemize}
  \item First item
  \item Second item
  \item Third item with sub-items:
    \begin{itemize}
      \item Sub-item 1
      \item Sub-item 2
    \end{itemize}
\end{itemize}

\subsection{Numbered List}
\begin{enumerate}
  \item First step
  \item Second step
  \item Third step
\end{enumerate}

\section{Citations and References}
When you use the citation search feature, references will be automatically managed. For example, you might cite a paper like this: \cite{example2024}. All citations will be automatically added to your \texttt{cite.bib} file.

\section{Conclusion}
This template demonstrates basic \LaTeX{} features including mathematical formulas, tables, and lists. The bibliography system is now set up to work seamlessly with Bibby AI's citation management. Feel free to modify it for your needs!

\bibliographystyle{plain}
\bibliography{cite}

\end{document}\documentclass{article}
\usepackage[utf8]{inputenc}
\usepackage{booktabs}
\usepackage{amsmath}
\usepackage{natbib}

\title{My Document}
\author{Your Name}
\date{\today}

\begin{document}

\maketitle

\section{Introduction}
This is a sample document demonstrating various \LaTeX{} features including mathematical equations, tables, and formatting. You can cite papers using the citation search feature, and they will be automatically added to your \texttt{cite.bib} file.

\section{Mathematical Expressions}
Here are some example equations:

\subsection{Quadratic Formula}
The quadratic formula is given by:
\[x = \frac{-b \pm \sqrt{b^2 - 4ac}}{2a}\]

\subsection{Matrix Example}
A sample matrix:
\[
\begin{pmatrix}
a_{11} & a_{12} & a_{13} \\
a_{21} & a_{22} & a_{23} \\
a_{31} & a_{32} & a_{33}
\end{pmatrix}
\]

\section{Tables}
Below is a sample table showing data:

\begin{table}[h]
\centering
\begin{tabular}{@{}lrrr@{}}
\toprule
\textbf{Category} & \textbf{Value 1} & \textbf{Value 2} & \textbf{Total} \\
\midrule
Item A & 10.5 & 20.3 & 30.8 \\
Item B & 15.2 & 25.1 & 40.3 \\
Item C & 22.1 & 12.7 & 34.8 \\
\bottomrule
\end{tabular}
\caption{Sample Data Table}
\label{tab:sample}
\end{table}

\section{Lists}
\subsection{Bullet Points}
\begin{itemize}
  \item First item
  \item Second item
  \item Third item with sub-items:
    \begin{itemize}
      \item Sub-item 1
      \item Sub-item 2
    \end{itemize}
\end{itemize}

\subsection{Numbered List}
\begin{enumerate}
  \item First step
  \item Second step
  \item Third step
\end{enumerate}

\section{Citations and References}
When you use the citation search feature, references will be automatically managed. For example, you might cite a paper like this: \cite{example2024}. All citations will be automatically added to your \texttt{cite.bib} file.

\section{Conclusion}
This template demonstrates basic \LaTeX{} features including mathematical formulas, tables, and lists. The bibliography system is now set up to work seamlessly with Bibby AI's citation management. Feel free to modify it for your needs!

\bibliographystyle{plain}
\bibliography{cite}

\end{document}\documentclass{article}
\usepackage[utf8]{inputenc}
\usepackage{booktabs}
\usepackage{amsmath}
\usepackage{natbib}

\title{My Document}
\author{Your Name}
\date{\today}

\begin{document}

\maketitle

\section{Introduction}
This is a sample document demonstrating various \LaTeX{} features including mathematical equations, tables, and formatting. You can cite papers using the citation search feature, and they will be automatically added to your \texttt{cite.bib} file.

\section{Mathematical Expressions}
Here are some example equations:

\subsection{Quadratic Formula}
The quadratic formula is given by:
\[x = \frac{-b \pm \sqrt{b^2 - 4ac}}{2a}\]

\subsection{Matrix Example}
A sample matrix:
\[
\begin{pmatrix}
a_{11} & a_{12} & a_{13} \\
a_{21} & a_{22} & a_{23} \\
a_{31} & a_{32} & a_{33}
\end{pmatrix}
\]

\section{Tables}
Below is a sample table showing data:

\begin{table}[h]
\centering
\begin{tabular}{@{}lrrr@{}}
\toprule
\textbf{Category} & \textbf{Value 1} & \textbf{Value 2} & \textbf{Total} \\
\midrule
Item A & 10.5 & 20.3 & 30.8 \\
Item B & 15.2 & 25.1 & 40.3 \\
Item C & 22.1 & 12.7 & 34.8 \\
\bottomrule
\end{tabular}
\caption{Sample Data Table}
\label{tab:sample}
\end{table}

\section{Lists}
\subsection{Bullet Points}
\begin{itemize}
  \item First item
  \item Second item
  \item Third item with sub-items:
    \begin{itemize}
      \item Sub-item 1
      \item Sub-item 2
    \end{itemize}
\end{itemize}

\subsection{Numbered List}
\begin{enumerate}
  \item First step
  \item Second step
  \item Third step
\end{enumerate}

\section{Citations and References}
When you use the citation search feature, references will be automatically managed. For example, you might cite a paper like this: \cite{example2024}. All citations will be automatically added to your \texttt{cite.bib} file.

\section{Conclusion}
This template demonstrates basic \LaTeX{} features including mathematical formulas, tables, and lists. The bibliography system is now set up to work seamlessly with Bibby AI's citation management. Feel free to modify it for your needs!

\bibliographystyle{plain}
\bibliography{cite}

\end{document}\documentclass{article}
\usepackage[utf8]{inputenc}
\usepackage{booktabs}
\usepackage{amsmath}
\usepackage{natbib}

\title{My Document}
\author{Your Name}
\date{\today}

\begin{document}

\maketitle

\section{Introduction}
This is a sample document demonstrating various \LaTeX{} features including mathematical equations, tables, and formatting. You can cite papers using the citation search feature, and they will be automatically added to your \texttt{cite.bib} file.

\section{Mathematical Expressions}
Here are some example equations:

\subsection{Quadratic Formula}
The quadratic formula is given by:
\[x = \frac{-b \pm \sqrt{b^2 - 4ac}}{2a}\]

\subsection{Matrix Example}
A sample matrix:
\[
\begin{pmatrix}
a_{11} & a_{12} & a_{13} \\
a_{21} & a_{22} & a_{23} \\
a_{31} & a_{32} & a_{33}
\end{pmatrix}
\]

\section{Tables}
Below is a sample table showing data:

\begin{table}[h]
\centering
\begin{tabular}{@{}lrrr@{}}
\toprule
\textbf{Category} & \textbf{Value 1} & \textbf{Value 2} & \textbf{Total} \\
\midrule
Item A & 10.5 & 20.3 & 30.8 \\
Item B & 15.2 & 25.1 & 40.3 \\
Item C & 22.1 & 12.7 & 34.8 \\
\bottomrule
\end{tabular}
\caption{Sample Data Table}
\label{tab:sample}
\end{table}

\section{Lists}
\subsection{Bullet Points}
\begin{itemize}
  \item First item
  \item Second item
  \item Third item with sub-items:
    \begin{itemize}
      \item Sub-item 1
      \item Sub-item 2
    \end{itemize}
\end{itemize}

\subsection{Numbered List}
\begin{enumerate}
  \item First step
  \item Second step
  \item Third step
\end{enumerate}

\section{Citations and References}
When you use the citation search feature, references will be automatically managed. For example, you might cite a paper like this: \cite{example2024}. All citations will be automatically added to your \texttt{cite.bib} file.

\section{Conclusion}
This template demonstrates basic \LaTeX{} features including mathematical formulas, tables, and lists. The bibliography system is now set up to work seamlessly with Bibby AI's citation management. Feel free to modify it for your needs!

\bibliographystyle{plain}
\bibliography{cite}

\end{document}\documentclass{article}
\usepackage[utf8]{inputenc}
\usepackage{booktabs}
\usepackage{amsmath}
\usepackage{natbib}

\title{My Document}
\author{Your Name}
\date{\today}

\begin{document}

\maketitle

\section{Introduction}
This is a sample document demonstrating various \LaTeX{} features including mathematical equations, tables, and formatting. You can cite papers using the citation search feature, and they will be automatically added to your \texttt{cite.bib} file.

\section{Mathematical Expressions}
Here are some example equations:

\subsection{Quadratic Formula}
The quadratic formula is given by:
\[x = \frac{-b \pm \sqrt{b^2 - 4ac}}{2a}\]

\subsection{Matrix Example}
A sample matrix:
\[
\begin{pmatrix}
a_{11} & a_{12} & a_{13} \\
a_{21} & a_{22} & a_{23} \\
a_{31} & a_{32} & a_{33}
\end{pmatrix}
\]

\section{Tables}
Below is a sample table showing data:

\begin{table}[h]
\centering
\begin{tabular}{@{}lrrr@{}}
\toprule
\textbf{Category} & \textbf{Value 1} & \textbf{Value 2} & \textbf{Total} \\
\midrule
Item A & 10.5 & 20.3 & 30.8 \\
Item B & 15.2 & 25.1 & 40.3 \\
Item C & 22.1 & 12.7 & 34.8 \\
\bottomrule
\end{tabular}
\caption{Sample Data Table}
\label{tab:sample}
\end{table}

\section{Lists}
\subsection{Bullet Points}
\begin{itemize}
  \item First item
  \item Second item
  \item Third item with sub-items:
    \begin{itemize}
      \item Sub-item 1
      \item Sub-item 2
    \end{itemize}
\end{itemize}

\subsection{Numbered List}
\begin{enumerate}
  \item First step
  \item Second step
  \item Third step
\end{enumerate}

\section{Citations and References}
When you use the citation search feature, references will be automatically managed. For example, you might cite a paper like this: \cite{example2024}. All citations will be automatically added to your \texttt{cite.bib} file.

\section{Conclusion}
This template demonstrates basic \LaTeX{} features including mathematical formulas, tables, and lists. The bibliography system is now set up to work seamlessly with Bibby AI's citation management. Feel free to modify it for your needs!

\bibliographystyle{plain}
\bibliography{cite}

\end{document}\documentclass{article}
\usepackage[utf8]{inputenc}
\usepackage{booktabs}
\usepackage{amsmath}
\usepackage{natbib}

\title{My Document}
\author{Your Name}
\date{\today}

\begin{document}

\maketitle

\section{Introduction}
This is a sample document demonstrating various \LaTeX{} features including mathematical equations, tables, and formatting. You can cite papers using the citation search feature, and they will be automatically added to your \texttt{cite.bib} file.

\section{Mathematical Expressions}
Here are some example equations:

\subsection{Quadratic Formula}
The quadratic formula is given by:
\[x = \frac{-b \pm \sqrt{b^2 - 4ac}}{2a}\]

\subsection{Matrix Example}
A sample matrix:
\[
\begin{pmatrix}
a_{11} & a_{12} & a_{13} \\
a_{21} & a_{22} & a_{23} \\
a_{31} & a_{32} & a_{33}
\end{pmatrix}
\]

\section{Tables}
Below is a sample table showing data:

\begin{table}[h]
\centering
\begin{tabular}{@{}lrrr@{}}
\toprule
\textbf{Category} & \textbf{Value 1} & \textbf{Value 2} & \textbf{Total} \\
\midrule
Item A & 10.5 & 20.3 & 30.8 \\
Item B & 15.2 & 25.1 & 40.3 \\
Item C & 22.1 & 12.7 & 34.8 \\
\bottomrule
\end{tabular}
\caption{Sample Data Table}
\label{tab:sample}
\end{table}

\section{Lists}
\subsection{Bullet Points}
\begin{itemize}
  \item First item
  \item Second item
  \item Third item with sub-items:
    \begin{itemize}
      \item Sub-item 1
      \item Sub-item 2
    \end{itemize}
\end{itemize}

\subsection{Numbered List}
\begin{enumerate}
  \item First step
  \item Second step
  \item Third step
\end{enumerate}

\section{Citations and References}
When you use the citation search feature, references will be automatically managed. For example, you might cite a paper like this: \cite{example2024}. All citations will be automatically added to your \texttt{cite.bib} file.

\section{Conclusion}
This template demonstrates basic \LaTeX{} features including mathematical formulas, tables, and lists. The bibliography system is now set up to work seamlessly with Bibby AI's citation management. Feel free to modify it for your needs!

\bibliographystyle{plain}
\bibliography{cite}

\end{document}\documentclass{article}
\usepackage[utf8]{inputenc}
\usepackage{booktabs}
\usepackage{amsmath}
\usepackage{natbib}

\title{My Document}
\author{Your Name}
\date{\today}

\begin{document}

\maketitle

\section{Introduction}
This is a sample document demonstrating various \LaTeX{} features including mathematical equations, tables, and formatting. You can cite papers using the citation search feature, and they will be automatically added to your \texttt{cite.bib} file.

\section{Mathematical Expressions}
Here are some example equations:

\subsection{Quadratic Formula}
The quadratic formula is given by:
\[x = \frac{-b \pm \sqrt{b^2 - 4ac}}{2a}\]

\subsection{Matrix Example}
A sample matrix:
\[
\begin{pmatrix}
a_{11} & a_{12} & a_{13} \\
a_{21} & a_{22} & a_{23} \\
a_{31} & a_{32} & a_{33}
\end{pmatrix}
\]

\section{Tables}
Below is a sample table showing data:

\begin{table}[h]
\centering
\begin{tabular}{@{}lrrr@{}}
\toprule
\textbf{Category} & \textbf{Value 1} & \textbf{Value 2} & \textbf{Total} \\
\midrule
Item A & 10.5 & 20.3 & 30.8 \\
Item B & 15.2 & 25.1 & 40.3 \\
Item C & 22.1 & 12.7 & 34.8 \\
\bottomrule
\end{tabular}
\caption{Sample Data Table}
\label{tab:sample}
\end{table}

\section{Lists}
\subsection{Bullet Points}
\begin{itemize}
  \item First item
  \item Second item
  \item Third item with sub-items:
    \begin{itemize}
      \item Sub-item 1
      \item Sub-item 2
    \end{itemize}
\end{itemize}

\subsection{Numbered List}
\begin{enumerate}
  \item First step
  \item Second step
  \item Third step
\end{enumerate}

\section{Citations and References}
When you use the citation search feature, references will be automatically managed. For example, you might cite a paper like this: \cite{example2024}. All citations will be automatically added to your \texttt{cite.bib} file.

\section{Conclusion}
This template demonstrates basic \LaTeX{} features including mathematical formulas, tables, and lists. The bibliography system is now set up to work seamlessly with Bibby AI's citation management. Feel free to modify it for your needs!

\bibliographystyle{plain}
\bibliography{cite}

\end{document}\documentclass{article}
\usepackage[utf8]{inputenc}
\usepackage{booktabs}
\usepackage{amsmath}
\usepackage{natbib}

\title{My Document}
\author{Your Name}
\date{\today}

\begin{document}

\maketitle

\section{Introduction}
This is a sample document demonstrating various \LaTeX{} features including mathematical equations, tables, and formatting. You can cite papers using the citation search feature, and they will be automatically added to your \texttt{cite.bib} file.

\section{Mathematical Expressions}
Here are some example equations:

\subsection{Quadratic Formula}
The quadratic formula is given by:
\[x = \frac{-b \pm \sqrt{b^2 - 4ac}}{2a}\]

\subsection{Matrix Example}
A sample matrix:
\[
\begin{pmatrix}
a_{11} & a_{12} & a_{13} \\
a_{21} & a_{22} & a_{23} \\
a_{31} & a_{32} & a_{33}
\end{pmatrix}
\]

\section{Tables}
Below is a sample table showing data:

\begin{table}[h]
\centering
\begin{tabular}{@{}lrrr@{}}
\toprule
\textbf{Category} & \textbf{Value 1} & \textbf{Value 2} & \textbf{Total} \\
\midrule
Item A & 10.5 & 20.3 & 30.8 \\
Item B & 15.2 & 25.1 & 40.3 \\
Item C & 22.1 & 12.7 & 34.8 \\
\bottomrule
\end{tabular}
\caption{Sample Data Table}
\label{tab:sample}
\end{table}

\section{Lists}
\subsection{Bullet Points}
\begin{itemize}
  \item First item
  \item Second item
  \item Third item with sub-items:
    \begin{itemize}
      \item Sub-item 1
      \item Sub-item 2
    \end{itemize}
\end{itemize}

\subsection{Numbered List}
\begin{enumerate}
  \item First step
  \item Second step
  \item Third step
\end{enumerate}

\section{Citations and References}
When you use the citation search feature, references will be automatically managed. For example, you might cite a paper like this: \cite{example2024}. All citations will be automatically added to your \texttt{cite.bib} file.

\section{Conclusion}
This template demonstrates basic \LaTeX{} features including mathematical formulas, tables, and lists. The bibliography system is now set up to work seamlessly with Bibby AI's citation management. Feel free to modify it for your needs!

\bibliographystyle{plain}
\bibliography{cite}

\end{document}\documentclass{article}
\usepackage[utf8]{inputenc}
\usepackage{booktabs}
\usepackage{amsmath}
\usepackage{natbib}

\title{My Document}
\author{Your Name}
\date{\today}

\begin{document}

\maketitle

\section{Introduction}
This is a sample document demonstrating various \LaTeX{} features including mathematical equations, tables, and formatting. You can cite papers using the citation search feature, and they will be automatically added to your \texttt{cite.bib} file.

\section{Mathematical Expressions}
Here are some example equations:

\subsection{Quadratic Formula}
The quadratic formula is given by:
\[x = \frac{-b \pm \sqrt{b^2 - 4ac}}{2a}\]

\subsection{Matrix Example}
A sample matrix:
\[
\begin{pmatrix}
a_{11} & a_{12} & a_{13} \\
a_{21} & a_{22} & a_{23} \\
a_{31} & a_{32} & a_{33}
\end{pmatrix}
\]

\section{Tables}
Below is a sample table showing data:

\begin{table}[h]
\centering
\begin{tabular}{@{}lrrr@{}}
\toprule
\textbf{Category} & \textbf{Value 1} & \textbf{Value 2} & \textbf{Total} \\
\midrule
Item A & 10.5 & 20.3 & 30.8 \\
Item B & 15.2 & 25.1 & 40.3 \\
Item C & 22.1 & 12.7 & 34.8 \\
\bottomrule
\end{tabular}
\caption{Sample Data Table}
\label{tab:sample}
\end{table}

\section{Lists}
\subsection{Bullet Points}
\begin{itemize}
  \item First item
  \item Second item
  \item Third item with sub-items:
    \begin{itemize}
      \item Sub-item 1
      \item Sub-item 2
    \end{itemize}
\end{itemize}

\subsection{Numbered List}
\begin{enumerate}
  \item First step
  \item Second step
  \item Third step
\end{enumerate}

\section{Citations and References}
When you use the citation search feature, references will be automatically managed. For example, you might cite a paper like this: \cite{example2024}. All citations will be automatically added to your \texttt{cite.bib} file.

\section{Conclusion}
This template demonstrates basic \LaTeX{} features including mathematical formulas, tables, and lists. The bibliography system is now set up to work seamlessly with Bibby AI's citation management. Feel free to modify it for your needs!

\bibliographystyle{plain}
\bibliography{cite}

\end{document}\documentclass{article}
\usepackage[utf8]{inputenc}
\usepackage{booktabs}
\usepackage{amsmath}
\usepackage{natbib}

\title{My Document}
\author{Your Name}
\date{\today}

\begin{document}

\maketitle

\section{Introduction}
This is a sample document demonstrating various \LaTeX{} features including mathematical equations, tables, and formatting. You can cite papers using the citation search feature, and they will be automatically added to your \texttt{cite.bib} file.

\section{Mathematical Expressions}
Here are some example equations:

\subsection{Quadratic Formula}
The quadratic formula is given by:
\[x = \frac{-b \pm \sqrt{b^2 - 4ac}}{2a}\]

\subsection{Matrix Example}
A sample matrix:
\[
\begin{pmatrix}
a_{11} & a_{12} & a_{13} \\
a_{21} & a_{22} & a_{23} \\
a_{31} & a_{32} & a_{33}
\end{pmatrix}
\]

\section{Tables}
Below is a sample table showing data:

\begin{table}[h]
\centering
\begin{tabular}{@{}lrrr@{}}
\toprule
\textbf{Category} & \textbf{Value 1} & \textbf{Value 2} & \textbf{Total} \\
\midrule
Item A & 10.5 & 20.3 & 30.8 \\
Item B & 15.2 & 25.1 & 40.3 \\
Item C & 22.1 & 12.7 & 34.8 \\
\bottomrule
\end{tabular}
\caption{Sample Data Table}
\label{tab:sample}
\end{table}

\section{Lists}
\subsection{Bullet Points}
\begin{itemize}
  \item First item
  \item Second item
  \item Third item with sub-items:
    \begin{itemize}
      \item Sub-item 1
      \item Sub-item 2
    \end{itemize}
\end{itemize}

\subsection{Numbered List}
\begin{enumerate}
  \item First step
  \item Second step
  \item Third step
\end{enumerate}

\section{Citations and References}
When you use the citation search feature, references will be automatically managed. For example, you might cite a paper like this: \cite{example2024}. All citations will be automatically added to your \texttt{cite.bib} file.

\section{Conclusion}
This template demonstrates basic \LaTeX{} features including mathematical formulas, tables, and lists. The bibliography system is now set up to work seamlessly with Bibby AI's citation management. Feel free to modify it for your needs!

\bibliographystyle{plain}
\bibliography{cite}

\end{document}\documentclass{article}
\usepackage[utf8]{inputenc}
\usepackage{booktabs}
\usepackage{amsmath}
\usepackage{natbib}

\title{My Document}
\author{Your Name}
\date{\today}

\begin{document}

\maketitle

\section{Introduction}
This is a sample document demonstrating various \LaTeX{} features including mathematical equations, tables, and formatting. You can cite papers using the citation search feature, and they will be automatically added to your \texttt{cite.bib} file.

\section{Mathematical Expressions}
Here are some example equations:

\subsection{Quadratic Formula}
The quadratic formula is given by:
\[x = \frac{-b \pm \sqrt{b^2 - 4ac}}{2a}\]

\subsection{Matrix Example}
A sample matrix:
\[
\begin{pmatrix}
a_{11} & a_{12} & a_{13} \\
a_{21} & a_{22} & a_{23} \\
a_{31} & a_{32} & a_{33}
\end{pmatrix}
\]

\section{Tables}
Below is a sample table showing data:

\begin{table}[h]
\centering
\begin{tabular}{@{}lrrr@{}}
\toprule
\textbf{Category} & \textbf{Value 1} & \textbf{Value 2} & \textbf{Total} \\
\midrule
Item A & 10.5 & 20.3 & 30.8 \\
Item B & 15.2 & 25.1 & 40.3 \\
Item C & 22.1 & 12.7 & 34.8 \\
\bottomrule
\end{tabular}
\caption{Sample Data Table}
\label{tab:sample}
\end{table}

\section{Lists}
\subsection{Bullet Points}
\begin{itemize}
  \item First item
  \item Second item
  \item Third item with sub-items:
    \begin{itemize}
      \item Sub-item 1
      \item Sub-item 2
    \end{itemize}
\end{itemize}

\subsection{Numbered List}
\begin{enumerate}
  \item First step
  \item Second step
  \item Third step
\end{enumerate}

\section{Citations and References}
When you use the citation search feature, references will be automatically managed. For example, you might cite a paper like this: \cite{example2024}. All citations will be automatically added to your \texttt{cite.bib} file.

\section{Conclusion}
This template demonstrates basic \LaTeX{} features including mathematical formulas, tables, and lists. The bibliography system is now set up to work seamlessly with Bibby AI's citation management. Feel free to modify it for your needs!

\bibliographystyle{plain}
\bibliography{cite}

\end{document}\documentclass{article}
\usepackage[utf8]{inputenc}
\usepackage{booktabs}
\usepackage{amsmath}
\usepackage{natbib}

\title{My Document}
\author{Your Name}
\date{\today}

\begin{document}

\maketitle

\section{Introduction}
This is a sample document demonstrating various \LaTeX{} features including mathematical equations, tables, and formatting. You can cite papers using the citation search feature, and they will be automatically added to your \texttt{cite.bib} file.

\section{Mathematical Expressions}
Here are some example equations:

\subsection{Quadratic Formula}
The quadratic formula is given by:
\[x = \frac{-b \pm \sqrt{b^2 - 4ac}}{2a}\]

\subsection{Matrix Example}
A sample matrix:
\[
\begin{pmatrix}
a_{11} & a_{12} & a_{13} \\
a_{21} & a_{22} & a_{23} \\
a_{31} & a_{32} & a_{33}
\end{pmatrix}
\]

\section{Tables}
Below is a sample table showing data:

\begin{table}[h]
\centering
\begin{tabular}{@{}lrrr@{}}
\toprule
\textbf{Category} & \textbf{Value 1} & \textbf{Value 2} & \textbf{Total} \\
\midrule
Item A & 10.5 & 20.3 & 30.8 \\
Item B & 15.2 & 25.1 & 40.3 \\
Item C & 22.1 & 12.7 & 34.8 \\
\bottomrule
\end{tabular}
\caption{Sample Data Table}
\label{tab:sample}
\end{table}

\section{Lists}
\subsection{Bullet Points}
\begin{itemize}
  \item First item
  \item Second item
  \item Third item with sub-items:
    \begin{itemize}
      \item Sub-item 1
      \item Sub-item 2
    \end{itemize}
\end{itemize}

\subsection{Numbered List}
\begin{enumerate}
  \item First step
  \item Second step
  \item Third step
\end{enumerate}

\section{Citations and References}
When you use the citation search feature, references will be automatically managed. For example, you might cite a paper like this: \cite{example2024}. All citations will be automatically added to your \texttt{cite.bib} file.

\section{Conclusion}
This template demonstrates basic \LaTeX{} features including mathematical formulas, tables, and lists. The bibliography system is now set up to work seamlessly with Bibby AI's citation management. Feel free to modify it for your needs!

\bibliographystyle{plain}
\bibliography{cite}

\end{document}\documentclass{article}
\usepackage[utf8]{inputenc}
\usepackage{booktabs}
\usepackage{amsmath}
\usepackage{natbib}

\title{My Document}
\author{Your Name}
\date{\today}

\begin{document}

\maketitle

\section{Introduction}
This is a sample document demonstrating various \LaTeX{} features including mathematical equations, tables, and formatting. You can cite papers using the citation search feature, and they will be automatically added to your \texttt{cite.bib} file.

\section{Mathematical Expressions}
Here are some example equations:

\subsection{Quadratic Formula}
The quadratic formula is given by:
\[x = \frac{-b \pm \sqrt{b^2 - 4ac}}{2a}\]

\subsection{Matrix Example}
A sample matrix:
\[
\begin{pmatrix}
a_{11} & a_{12} & a_{13} \\
a_{21} & a_{22} & a_{23} \\
a_{31} & a_{32} & a_{33}
\end{pmatrix}
\]

\section{Tables}
Below is a sample table showing data:

\begin{table}[h]
\centering
\begin{tabular}{@{}lrrr@{}}
\toprule
\textbf{Category} & \textbf{Value 1} & \textbf{Value 2} & \textbf{Total} \\
\midrule
Item A & 10.5 & 20.3 & 30.8 \\
Item B & 15.2 & 25.1 & 40.3 \\
Item C & 22.1 & 12.7 & 34.8 \\
\bottomrule
\end{tabular}
\caption{Sample Data Table}
\label{tab:sample}
\end{table}

\section{Lists}
\subsection{Bullet Points}
\begin{itemize}
  \item First item
  \item Second item
  \item Third item with sub-items:
    \begin{itemize}
      \item Sub-item 1
      \item Sub-item 2
    \end{itemize}
\end{itemize}

\subsection{Numbered List}
\begin{enumerate}
  \item First step
  \item Second step
  \item Third step
\end{enumerate}

\section{Citations and References}
When you use the citation search feature, references will be automatically managed. For example, you might cite a paper like this: \cite{example2024}. All citations will be automatically added to your \texttt{cite.bib} file.

\section{Conclusion}
This template demonstrates basic \LaTeX{} features including mathematical formulas, tables, and lists. The bibliography system is now set up to work seamlessly with Bibby AI's citation management. Feel free to modify it for your needs!

\bibliographystyle{plain}
\bibliography{cite}

\end{document}\documentclass{article}
\usepackage[utf8]{inputenc}
\usepackage{booktabs}
\usepackage{amsmath}
\usepackage{natbib}

\title{My Document}
\author{Your Name}
\date{\today}

\begin{document}

\maketitle

\section{Introduction}
This is a sample document demonstrating various \LaTeX{} features including mathematical equations, tables, and formatting. You can cite papers using the citation search feature, and they will be automatically added to your \texttt{cite.bib} file.

\section{Mathematical Expressions}
Here are some example equations:

\subsection{Quadratic Formula}
The quadratic formula is given by:
\[x = \frac{-b \pm \sqrt{b^2 - 4ac}}{2a}\]

\subsection{Matrix Example}
A sample matrix:
\[
\begin{pmatrix}
a_{11} & a_{12} & a_{13} \\
a_{21} & a_{22} & a_{23} \\
a_{31} & a_{32} & a_{33}
\end{pmatrix}
\]

\section{Tables}
Below is a sample table showing data:

\begin{table}[h]
\centering
\begin{tabular}{@{}lrrr@{}}
\toprule
\textbf{Category} & \textbf{Value 1} & \textbf{Value 2} & \textbf{Total} \\
\midrule
Item A & 10.5 & 20.3 & 30.8 \\
Item B & 15.2 & 25.1 & 40.3 \\
Item C & 22.1 & 12.7 & 34.8 \\
\bottomrule
\end{tabular}
\caption{Sample Data Table}
\label{tab:sample}
\end{table}

\section{Lists}
\subsection{Bullet Points}
\begin{itemize}
  \item First item
  \item Second item
  \item Third item with sub-items:
    \begin{itemize}
      \item Sub-item 1
      \item Sub-item 2
    \end{itemize}
\end{itemize}

\subsection{Numbered List}
\begin{enumerate}
  \item First step
  \item Second step
  \item Third step
\end{enumerate}

\section{Citations and References}
When you use the citation search feature, references will be automatically managed. For example, you might cite a paper like this: \cite{example2024}. All citations will be automatically added to your \texttt{cite.bib} file.

\section{Conclusion}
This template demonstrates basic \LaTeX{} features including mathematical formulas, tables, and lists. The bibliography system is now set up to work seamlessly with Bibby AI's citation management. Feel free to modify it for your needs!

\bibliographystyle{plain}
\bibliography{cite}

\end{document}\documentclass{article}
\usepackage[utf8]{inputenc}
\usepackage{booktabs}
\usepackage{amsmath}
\usepackage{natbib}

\title{My Document}
\author{Your Name}
\date{\today}

\begin{document}

\maketitle

\section{Introduction}
This is a sample document demonstrating various \LaTeX{} features including mathematical equations, tables, and formatting. You can cite papers using the citation search feature, and they will be automatically added to your \texttt{cite.bib} file.

\section{Mathematical Expressions}
Here are some example equations:

\subsection{Quadratic Formula}
The quadratic formula is given by:
\[x = \frac{-b \pm \sqrt{b^2 - 4ac}}{2a}\]

\subsection{Matrix Example}
A sample matrix:
\[
\begin{pmatrix}
a_{11} & a_{12} & a_{13} \\
a_{21} & a_{22} & a_{23} \\
a_{31} & a_{32} & a_{33}
\end{pmatrix}
\]

\section{Tables}
Below is a sample table showing data:

\begin{table}[h]
\centering
\begin{tabular}{@{}lrrr@{}}
\toprule
\textbf{Category} & \textbf{Value 1} & \textbf{Value 2} & \textbf{Total} \\
\midrule
Item A & 10.5 & 20.3 & 30.8 \\
Item B & 15.2 & 25.1 & 40.3 \\
Item C & 22.1 & 12.7 & 34.8 \\
\bottomrule
\end{tabular}
\caption{Sample Data Table}
\label{tab:sample}
\end{table}

\section{Lists}
\subsection{Bullet Points}
\begin{itemize}
  \item First item
  \item Second item
  \item Third item with sub-items:
    \begin{itemize}
      \item Sub-item 1
      \item Sub-item 2
    \end{itemize}
\end{itemize}

\subsection{Numbered List}
\begin{enumerate}
  \item First step
  \item Second step
  \item Third step
\end{enumerate}

\section{Citations and References}
When you use the citation search feature, references will be automatically managed. For example, you might cite a paper like this: \cite{example2024}. All citations will be automatically added to your \texttt{cite.bib} file.

\section{Conclusion}
This template demonstrates basic \LaTeX{} features including mathematical formulas, tables, and lists. The bibliography system is now set up to work seamlessly with Bibby AI's citation management. Feel free to modify it for your needs!

\bibliographystyle{plain}
\bibliography{cite}

\end{document}\documentclass{article}
\usepackage[utf8]{inputenc}
\usepackage{booktabs}
\usepackage{amsmath}
\usepackage{natbib}

\title{My Document}
\author{Your Name}
\date{\today}

\begin{document}

\maketitle

\section{Introduction}
This is a sample document demonstrating various \LaTeX{} features including mathematical equations, tables, and formatting. You can cite papers using the citation search feature, and they will be automatically added to your \texttt{cite.bib} file.

\section{Mathematical Expressions}
Here are some example equations:

\subsection{Quadratic Formula}
The quadratic formula is given by:
\[x = \frac{-b \pm \sqrt{b^2 - 4ac}}{2a}\]

\subsection{Matrix Example}
A sample matrix:
\[
\begin{pmatrix}
a_{11} & a_{12} & a_{13} \\
a_{21} & a_{22} & a_{23} \\
a_{31} & a_{32} & a_{33}
\end{pmatrix}
\]

\section{Tables}
Below is a sample table showing data:

\begin{table}[h]
\centering
\begin{tabular}{@{}lrrr@{}}
\toprule
\textbf{Category} & \textbf{Value 1} & \textbf{Value 2} & \textbf{Total} \\
\midrule
Item A & 10.5 & 20.3 & 30.8 \\
Item B & 15.2 & 25.1 & 40.3 \\
Item C & 22.1 & 12.7 & 34.8 \\
\bottomrule
\end{tabular}
\caption{Sample Data Table}
\label{tab:sample}
\end{table}

\section{Lists}
\subsection{Bullet Points}
\begin{itemize}
  \item First item
  \item Second item
  \item Third item with sub-items:
    \begin{itemize}
      \item Sub-item 1
      \item Sub-item 2
    \end{itemize}
\end{itemize}

\subsection{Numbered List}
\begin{enumerate}
  \item First step
  \item Second step
  \item Third step
\end{enumerate}

\section{Citations and References}
When you use the citation search feature, references will be automatically managed. For example, you might cite a paper like this: \cite{example2024}. All citations will be automatically added to your \texttt{cite.bib} file.

\section{Conclusion}
This template demonstrates basic \LaTeX{} features including mathematical formulas, tables, and lists. The bibliography system is now set up to work seamlessly with Bibby AI's citation management. Feel free to modify it for your needs!

\bibliographystyle{plain}
\bibliography{cite}

\end{document}\documentclass{article}
\usepackage[utf8]{inputenc}
\usepackage{booktabs}
\usepackage{amsmath}
\usepackage{natbib}

\title{My Document}
\author{Your Name}
\date{\today}

\begin{document}

\maketitle

\section{Introduction}
This is a sample document demonstrating various \LaTeX{} features including mathematical equations, tables, and formatting. You can cite papers using the citation search feature, and they will be automatically added to your \texttt{cite.bib} file.

\section{Mathematical Expressions}
Here are some example equations:

\subsection{Quadratic Formula}
The quadratic formula is given by:
\[x = \frac{-b \pm \sqrt{b^2 - 4ac}}{2a}\]

\subsection{Matrix Example}
A sample matrix:
\[
\begin{pmatrix}
a_{11} & a_{12} & a_{13} \\
a_{21} & a_{22} & a_{23} \\
a_{31} & a_{32} & a_{33}
\end{pmatrix}
\]

\section{Tables}
Below is a sample table showing data:

\begin{table}[h]
\centering
\begin{tabular}{@{}lrrr@{}}
\toprule
\textbf{Category} & \textbf{Value 1} & \textbf{Value 2} & \textbf{Total} \\
\midrule
Item A & 10.5 & 20.3 & 30.8 \\
Item B & 15.2 & 25.1 & 40.3 \\
Item C & 22.1 & 12.7 & 34.8 \\
\bottomrule
\end{tabular}
\caption{Sample Data Table}
\label{tab:sample}
\end{table}

\section{Lists}
\subsection{Bullet Points}
\begin{itemize}
  \item First item
  \item Second item
  \item Third item with sub-items:
    \begin{itemize}
      \item Sub-item 1
      \item Sub-item 2
    \end{itemize}
\end{itemize}

\subsection{Numbered List}
\begin{enumerate}
  \item First step
  \item Second step
  \item Third step
\end{enumerate}

\section{Citations and References}
When you use the citation search feature, references will be automatically managed. For example, you might cite a paper like this: \cite{example2024}. All citations will be automatically added to your \texttt{cite.bib} file.

\section{Conclusion}
This template demonstrates basic \LaTeX{} features including mathematical formulas, tables, and lists. The bibliography system is now set up to work seamlessly with Bibby AI's citation management. Feel free to modify it for your needs!

\bibliographystyle{plain}
\bibliography{cite}

\end{document}\documentclass{article}
\usepackage[utf8]{inputenc}
\usepackage{booktabs}
\usepackage{amsmath}
\usepackage{natbib}

\title{My Document}
\author{Your Name}
\date{\today}

\begin{document}

\maketitle

\section{Introduction}
This is a sample document demonstrating various \LaTeX{} features including mathematical equations, tables, and formatting. You can cite papers using the citation search feature, and they will be automatically added to your \texttt{cite.bib} file.

\section{Mathematical Expressions}
Here are some example equations:

\subsection{Quadratic Formula}
The quadratic formula is given by:
\[x = \frac{-b \pm \sqrt{b^2 - 4ac}}{2a}\]

\subsection{Matrix Example}
A sample matrix:
\[
\begin{pmatrix}
a_{11} & a_{12} & a_{13} \\
a_{21} & a_{22} & a_{23} \\
a_{31} & a_{32} & a_{33}
\end{pmatrix}
\]

\section{Tables}
Below is a sample table showing data:

\begin{table}[h]
\centering
\begin{tabular}{@{}lrrr@{}}
\toprule
\textbf{Category} & \textbf{Value 1} & \textbf{Value 2} & \textbf{Total} \\
\midrule
Item A & 10.5 & 20.3 & 30.8 \\
Item B & 15.2 & 25.1 & 40.3 \\
Item C & 22.1 & 12.7 & 34.8 \\
\bottomrule
\end{tabular}
\caption{Sample Data Table}
\label{tab:sample}
\end{table}

\section{Lists}
\subsection{Bullet Points}
\begin{itemize}
  \item First item
  \item Second item
  \item Third item with sub-items:
    \begin{itemize}
      \item Sub-item 1
      \item Sub-item 2
    \end{itemize}
\end{itemize}

\subsection{Numbered List}
\begin{enumerate}
  \item First step
  \item Second step
  \item Third step
\end{enumerate}

\section{Citations and References}
When you use the citation search feature, references will be automatically managed. For example, you might cite a paper like this: \cite{example2024}. All citations will be automatically added to your \texttt{cite.bib} file.

\section{Conclusion}
This template demonstrates basic \LaTeX{} features including mathematical formulas, tables, and lists. The bibliography system is now set up to work seamlessly with Bibby AI's citation management. Feel free to modify it for your needs!

\bibliographystyle{plain}
\bibliography{cite}

\end{document}\documentclass{article}
\usepackage[utf8]{inputenc}
\usepackage{booktabs}
\usepackage{amsmath}
\usepackage{natbib}

\title{My Document}
\author{Your Name}
\date{\today}

\begin{document}

\maketitle

\section{Introduction}
This is a sample document demonstrating various \LaTeX{} features including mathematical equations, tables, and formatting. You can cite papers using the citation search feature, and they will be automatically added to your \texttt{cite.bib} file.

\section{Mathematical Expressions}
Here are some example equations:

\subsection{Quadratic Formula}
The quadratic formula is given by:
\[x = \frac{-b \pm \sqrt{b^2 - 4ac}}{2a}\]

\subsection{Matrix Example}
A sample matrix:
\[
\begin{pmatrix}
a_{11} & a_{12} & a_{13} \\
a_{21} & a_{22} & a_{23} \\
a_{31} & a_{32} & a_{33}
\end{pmatrix}
\]

\section{Tables}
Below is a sample table showing data:

\begin{table}[h]
\centering
\begin{tabular}{@{}lrrr@{}}
\toprule
\textbf{Category} & \textbf{Value 1} & \textbf{Value 2} & \textbf{Total} \\
\midrule
Item A & 10.5 & 20.3 & 30.8 \\
Item B & 15.2 & 25.1 & 40.3 \\
Item C & 22.1 & 12.7 & 34.8 \\
\bottomrule
\end{tabular}
\caption{Sample Data Table}
\label{tab:sample}
\end{table}

\section{Lists}
\subsection{Bullet Points}
\begin{itemize}
  \item First item
  \item Second item
  \item Third item with sub-items:
    \begin{itemize}
      \item Sub-item 1
      \item Sub-item 2
    \end{itemize}
\end{itemize}

\subsection{Numbered List}
\begin{enumerate}
  \item First step
  \item Second step
  \item Third step
\end{enumerate}

\section{Citations and References}
When you use the citation search feature, references will be automatically managed. For example, you might cite a paper like this: \cite{example2024}. All citations will be automatically added to your \texttt{cite.bib} file.

\section{Conclusion}
This template demonstrates basic \LaTeX{} features including mathematical formulas, tables, and lists. The bibliography system is now set up to work seamlessly with Bibby AI's citation management. Feel free to modify it for your needs!

\bibliographystyle{plain}
\bibliography{cite}

\end{document}\documentclass{article}
\usepackage[utf8]{inputenc}
\usepackage{booktabs}
\usepackage{amsmath}
\usepackage{natbib}

\title{My Document}
\author{Your Name}
\date{\today}

\begin{document}

\maketitle

\section{Introduction}
This is a sample document demonstrating various \LaTeX{} features including mathematical equations, tables, and formatting. You can cite papers using the citation search feature, and they will be automatically added to your \texttt{cite.bib} file.

\section{Mathematical Expressions}
Here are some example equations:

\subsection{Quadratic Formula}
The quadratic formula is given by:
\[x = \frac{-b \pm \sqrt{b^2 - 4ac}}{2a}\]

\subsection{Matrix Example}
A sample matrix:
\[
\begin{pmatrix}
a_{11} & a_{12} & a_{13} \\
a_{21} & a_{22} & a_{23} \\
a_{31} & a_{32} & a_{33}
\end{pmatrix}
\]

\section{Tables}
Below is a sample table showing data:

\begin{table}[h]
\centering
\begin{tabular}{@{}lrrr@{}}
\toprule
\textbf{Category} & \textbf{Value 1} & \textbf{Value 2} & \textbf{Total} \\
\midrule
Item A & 10.5 & 20.3 & 30.8 \\
Item B & 15.2 & 25.1 & 40.3 \\
Item C & 22.1 & 12.7 & 34.8 \\
\bottomrule
\end{tabular}
\caption{Sample Data Table}
\label{tab:sample}
\end{table}

\section{Lists}
\subsection{Bullet Points}
\begin{itemize}
  \item First item
  \item Second item
  \item Third item with sub-items:
    \begin{itemize}
      \item Sub-item 1
      \item Sub-item 2
    \end{itemize}
\end{itemize}

\subsection{Numbered List}
\begin{enumerate}
  \item First step
  \item Second step
  \item Third step
\end{enumerate}

\section{Citations and References}
When you use the citation search feature, references will be automatically managed. For example, you might cite a paper like this: \cite{example2024}. All citations will be automatically added to your \texttt{cite.bib} file.

\section{Conclusion}
This template demonstrates basic \LaTeX{} features including mathematical formulas, tables, and lists. The bibliography system is now set up to work seamlessly with Bibby AI's citation management. Feel free to modify it for your needs!

\bibliographystyle{plain}
\bibliography{cite}

\end{document}\documentclass{article}
\usepackage[utf8]{inputenc}
\usepackage{booktabs}
\usepackage{amsmath}
\usepackage{natbib}

\title{My Document}
\author{Your Name}
\date{\today}

\begin{document}

\maketitle

\section{Introduction}
This is a sample document demonstrating various \LaTeX{} features including mathematical equations, tables, and formatting. You can cite papers using the citation search feature, and they will be automatically added to your \texttt{cite.bib} file.

\section{Mathematical Expressions}
Here are some example equations:

\subsection{Quadratic Formula}
The quadratic formula is given by:
\[x = \frac{-b \pm \sqrt{b^2 - 4ac}}{2a}\]

\subsection{Matrix Example}
A sample matrix:
\[
\begin{pmatrix}
a_{11} & a_{12} & a_{13} \\
a_{21} & a_{22} & a_{23} \\
a_{31} & a_{32} & a_{33}
\end{pmatrix}
\]

\section{Tables}
Below is a sample table showing data:

\begin{table}[h]
\centering
\begin{tabular}{@{}lrrr@{}}
\toprule
\textbf{Category} & \textbf{Value 1} & \textbf{Value 2} & \textbf{Total} \\
\midrule
Item A & 10.5 & 20.3 & 30.8 \\
Item B & 15.2 & 25.1 & 40.3 \\
Item C & 22.1 & 12.7 & 34.8 \\
\bottomrule
\end{tabular}
\caption{Sample Data Table}
\label{tab:sample}
\end{table}

\section{Lists}
\subsection{Bullet Points}
\begin{itemize}
  \item First item
  \item Second item
  \item Third item with sub-items:
    \begin{itemize}
      \item Sub-item 1
      \item Sub-item 2
    \end{itemize}
\end{itemize}

\subsection{Numbered List}
\begin{enumerate}
  \item First step
  \item Second step
  \item Third step
\end{enumerate}

\section{Citations and References}
When you use the citation search feature, references will be automatically managed. For example, you might cite a paper like this: \cite{example2024}. All citations will be automatically added to your \texttt{cite.bib} file.

\section{Conclusion}
This template demonstrates basic \LaTeX{} features including mathematical formulas, tables, and lists. The bibliography system is now set up to work seamlessly with Bibby AI's citation management. Feel free to modify it for your needs!

\bibliographystyle{plain}
\bibliography{cite}

\end{document}\documentclass{article}
\usepackage[utf8]{inputenc}
\usepackage{booktabs}
\usepackage{amsmath}
\usepackage{natbib}

\title{My Document}
\author{Your Name}
\date{\today}

\begin{document}

\maketitle

\section{Introduction}
This is a sample document demonstrating various \LaTeX{} features including mathematical equations, tables, and formatting. You can cite papers using the citation search feature, and they will be automatically added to your \texttt{cite.bib} file.

\section{Mathematical Expressions}
Here are some example equations:

\subsection{Quadratic Formula}
The quadratic formula is given by:
\[x = \frac{-b \pm \sqrt{b^2 - 4ac}}{2a}\]

\subsection{Matrix Example}
A sample matrix:
\[
\begin{pmatrix}
a_{11} & a_{12} & a_{13} \\
a_{21} & a_{22} & a_{23} \\
a_{31} & a_{32} & a_{33}
\end{pmatrix}
\]

\section{Tables}
Below is a sample table showing data:

\begin{table}[h]
\centering
\begin{tabular}{@{}lrrr@{}}
\toprule
\textbf{Category} & \textbf{Value 1} & \textbf{Value 2} & \textbf{Total} \\
\midrule
Item A & 10.5 & 20.3 & 30.8 \\
Item B & 15.2 & 25.1 & 40.3 \\
Item C & 22.1 & 12.7 & 34.8 \\
\bottomrule
\end{tabular}
\caption{Sample Data Table}
\label{tab:sample}
\end{table}

\section{Lists}
\subsection{Bullet Points}
\begin{itemize}
  \item First item
  \item Second item
  \item Third item with sub-items:
    \begin{itemize}
      \item Sub-item 1
      \item Sub-item 2
    \end{itemize}
\end{itemize}

\subsection{Numbered List}
\begin{enumerate}
  \item First step
  \item Second step
  \item Third step
\end{enumerate}

\section{Citations and References}
When you use the citation search feature, references will be automatically managed. For example, you might cite a paper like this: \cite{example2024}. All citations will be automatically added to your \texttt{cite.bib} file.

\section{Conclusion}
This template demonstrates basic \LaTeX{} features including mathematical formulas, tables, and lists. The bibliography system is now set up to work seamlessly with Bibby AI's citation management. Feel free to modify it for your needs!

\bibliographystyle{plain}
\bibliography{cite}

\end{document}\documentclass{article}
\usepackage[utf8]{inputenc}
\usepackage{booktabs}
\usepackage{amsmath}
\usepackage{natbib}

\title{My Document}
\author{Your Name}
\date{\today}

\begin{document}

\maketitle

\section{Introduction}
This is a sample document demonstrating various \LaTeX{} features including mathematical equations, tables, and formatting. You can cite papers using the citation search feature, and they will be automatically added to your \texttt{cite.bib} file.

\section{Mathematical Expressions}
Here are some example equations:

\subsection{Quadratic Formula}
The quadratic formula is given by:
\[x = \frac{-b \pm \sqrt{b^2 - 4ac}}{2a}\]

\subsection{Matrix Example}
A sample matrix:
\[
\begin{pmatrix}
a_{11} & a_{12} & a_{13} \\
a_{21} & a_{22} & a_{23} \\
a_{31} & a_{32} & a_{33}
\end{pmatrix}
\]

\section{Tables}
Below is a sample table showing data:

\begin{table}[h]
\centering
\begin{tabular}{@{}lrrr@{}}
\toprule
\textbf{Category} & \textbf{Value 1} & \textbf{Value 2} & \textbf{Total} \\
\midrule
Item A & 10.5 & 20.3 & 30.8 \\
Item B & 15.2 & 25.1 & 40.3 \\
Item C & 22.1 & 12.7 & 34.8 \\
\bottomrule
\end{tabular}
\caption{Sample Data Table}
\label{tab:sample}
\end{table}

\section{Lists}
\subsection{Bullet Points}
\begin{itemize}
  \item First item
  \item Second item
  \item Third item with sub-items:
    \begin{itemize}
      \item Sub-item 1
      \item Sub-item 2
    \end{itemize}
\end{itemize}

\subsection{Numbered List}
\begin{enumerate}
  \item First step
  \item Second step
  \item Third step
\end{enumerate}

\section{Citations and References}
When you use the citation search feature, references will be automatically managed. For example, you might cite a paper like this: \cite{example2024}. All citations will be automatically added to your \texttt{cite.bib} file.

\section{Conclusion}
This template demonstrates basic \LaTeX{} features including mathematical formulas, tables, and lists. The bibliography system is now set up to work seamlessly with Bibby AI's citation management. Feel free to modify it for your needs!

\bibliographystyle{plain}
\bibliography{cite}

\end{document}\documentclass{article}
\usepackage[utf8]{inputenc}
\usepackage{booktabs}
\usepackage{amsmath}
\usepackage{natbib}

\title{My Document}
\author{Your Name}
\date{\today}

\begin{document}

\maketitle

\section{Introduction}
This is a sample document demonstrating various \LaTeX{} features including mathematical equations, tables, and formatting. You can cite papers using the citation search feature, and they will be automatically added to your \texttt{cite.bib} file.

\section{Mathematical Expressions}
Here are some example equations:

\subsection{Quadratic Formula}
The quadratic formula is given by:
\[x = \frac{-b \pm \sqrt{b^2 - 4ac}}{2a}\]

\subsection{Matrix Example}
A sample matrix:
\[
\begin{pmatrix}
a_{11} & a_{12} & a_{13} \\
a_{21} & a_{22} & a_{23} \\
a_{31} & a_{32} & a_{33}
\end{pmatrix}
\]

\section{Tables}
Below is a sample table showing data:

\begin{table}[h]
\centering
\begin{tabular}{@{}lrrr@{}}
\toprule
\textbf{Category} & \textbf{Value 1} & \textbf{Value 2} & \textbf{Total} \\
\midrule
Item A & 10.5 & 20.3 & 30.8 \\
Item B & 15.2 & 25.1 & 40.3 \\
Item C & 22.1 & 12.7 & 34.8 \\
\bottomrule
\end{tabular}
\caption{Sample Data Table}
\label{tab:sample}
\end{table}

\section{Lists}
\subsection{Bullet Points}
\begin{itemize}
  \item First item
  \item Second item
  \item Third item with sub-items:
    \begin{itemize}
      \item Sub-item 1
      \item Sub-item 2
    \end{itemize}
\end{itemize}

\subsection{Numbered List}
\begin{enumerate}
  \item First step
  \item Second step
  \item Third step
\end{enumerate}

\section{Citations and References}
When you use the citation search feature, references will be automatically managed. For example, you might cite a paper like this: \cite{example2024}. All citations will be automatically added to your \texttt{cite.bib} file.

\section{Conclusion}
This template demonstrates basic \LaTeX{} features including mathematical formulas, tables, and lists. The bibliography system is now set up to work seamlessly with Bibby AI's citation management. Feel free to modify it for your needs!

\bibliographystyle{plain}
\bibliography{cite}

\end{document}\documentclass{article}
\usepackage[utf8]{inputenc}
\usepackage{booktabs}
\usepackage{amsmath}
\usepackage{natbib}

\title{My Document}
\author{Your Name}
\date{\today}

\begin{document}

\maketitle

\section{Introduction}
This is a sample document demonstrating various \LaTeX{} features including mathematical equations, tables, and formatting. You can cite papers using the citation search feature, and they will be automatically added to your \texttt{cite.bib} file.

\section{Mathematical Expressions}
Here are some example equations:

\subsection{Quadratic Formula}
The quadratic formula is given by:
\[x = \frac{-b \pm \sqrt{b^2 - 4ac}}{2a}\]

\subsection{Matrix Example}
A sample matrix:
\[
\begin{pmatrix}
a_{11} & a_{12} & a_{13} \\
a_{21} & a_{22} & a_{23} \\
a_{31} & a_{32} & a_{33}
\end{pmatrix}
\]

\section{Tables}
Below is a sample table showing data:

\begin{table}[h]
\centering
\begin{tabular}{@{}lrrr@{}}
\toprule
\textbf{Category} & \textbf{Value 1} & \textbf{Value 2} & \textbf{Total} \\
\midrule
Item A & 10.5 & 20.3 & 30.8 \\
Item B & 15.2 & 25.1 & 40.3 \\
Item C & 22.1 & 12.7 & 34.8 \\
\bottomrule
\end{tabular}
\caption{Sample Data Table}
\label{tab:sample}
\end{table}

\section{Lists}
\subsection{Bullet Points}
\begin{itemize}
  \item First item
  \item Second item
  \item Third item with sub-items:
    \begin{itemize}
      \item Sub-item 1
      \item Sub-item 2
    \end{itemize}
\end{itemize}

\subsection{Numbered List}
\begin{enumerate}
  \item First step
  \item Second step
  \item Third step
\end{enumerate}

\section{Citations and References}
When you use the citation search feature, references will be automatically managed. For example, you might cite a paper like this: \cite{example2024}. All citations will be automatically added to your \texttt{cite.bib} file.

\section{Conclusion}
This template demonstrates basic \LaTeX{} features including mathematical formulas, tables, and lists. The bibliography system is now set up to work seamlessly with Bibby AI's citation management. Feel free to modify it for your needs!

\bibliographystyle{plain}
\bibliography{cite}

\end{document}\documentclass{article}
\usepackage[utf8]{inputenc}
\usepackage{booktabs}
\usepackage{amsmath}
\usepackage{natbib}

\title{My Document}
\author{Your Name}
\date{\today}

\begin{document}

\maketitle

\section{Introduction}
This is a sample document demonstrating various \LaTeX{} features including mathematical equations, tables, and formatting. You can cite papers using the citation search feature, and they will be automatically added to your \texttt{cite.bib} file.

\section{Mathematical Expressions}
Here are some example equations:

\subsection{Quadratic Formula}
The quadratic formula is given by:
\[x = \frac{-b \pm \sqrt{b^2 - 4ac}}{2a}\]

\subsection{Matrix Example}
A sample matrix:
\[
\begin{pmatrix}
a_{11} & a_{12} & a_{13} \\
a_{21} & a_{22} & a_{23} \\
a_{31} & a_{32} & a_{33}
\end{pmatrix}
\]

\section{Tables}
Below is a sample table showing data:

\begin{table}[h]
\centering
\begin{tabular}{@{}lrrr@{}}
\toprule
\textbf{Category} & \textbf{Value 1} & \textbf{Value 2} & \textbf{Total} \\
\midrule
Item A & 10.5 & 20.3 & 30.8 \\
Item B & 15.2 & 25.1 & 40.3 \\
Item C & 22.1 & 12.7 & 34.8 \\
\bottomrule
\end{tabular}
\caption{Sample Data Table}
\label{tab:sample}
\end{table}

\section{Lists}
\subsection{Bullet Points}
\begin{itemize}
  \item First item
  \item Second item
  \item Third item with sub-items:
    \begin{itemize}
      \item Sub-item 1
      \item Sub-item 2
    \end{itemize}
\end{itemize}

\subsection{Numbered List}
\begin{enumerate}
  \item First step
  \item Second step
  \item Third step
\end{enumerate}

\section{Citations and References}
When you use the citation search feature, references will be automatically managed. For example, you might cite a paper like this: \cite{example2024}. All citations will be automatically added to your \texttt{cite.bib} file.

\section{Conclusion}
This template demonstrates basic \LaTeX{} features including mathematical formulas, tables, and lists. The bibliography system is now set up to work seamlessly with Bibby AI's citation management. Feel free to modify it for your needs!

\bibliographystyle{plain}
\bibliography{cite}

\end{document}\documentclass{article}
\usepackage[utf8]{inputenc}
\usepackage{booktabs}
\usepackage{amsmath}
\usepackage{natbib}

\title{My Document}
\author{Your Name}
\date{\today}

\begin{document}

\maketitle

\section{Introduction}
This is a sample document demonstrating various \LaTeX{} features including mathematical equations, tables, and formatting. You can cite papers using the citation search feature, and they will be automatically added to your \texttt{cite.bib} file.

\section{Mathematical Expressions}
Here are some example equations:

\subsection{Quadratic Formula}
The quadratic formula is given by:
\[x = \frac{-b \pm \sqrt{b^2 - 4ac}}{2a}\]

\subsection{Matrix Example}
A sample matrix:
\[
\begin{pmatrix}
a_{11} & a_{12} & a_{13} \\
a_{21} & a_{22} & a_{23} \\
a_{31} & a_{32} & a_{33}
\end{pmatrix}
\]

\section{Tables}
Below is a sample table showing data:

\begin{table}[h]
\centering
\begin{tabular}{@{}lrrr@{}}
\toprule
\textbf{Category} & \textbf{Value 1} & \textbf{Value 2} & \textbf{Total} \\
\midrule
Item A & 10.5 & 20.3 & 30.8 \\
Item B & 15.2 & 25.1 & 40.3 \\
Item C & 22.1 & 12.7 & 34.8 \\
\bottomrule
\end{tabular}
\caption{Sample Data Table}
\label{tab:sample}
\end{table}

\section{Lists}
\subsection{Bullet Points}
\begin{itemize}
  \item First item
  \item Second item
  \item Third item with sub-items:
    \begin{itemize}
      \item Sub-item 1
      \item Sub-item 2
    \end{itemize}
\end{itemize}

\subsection{Numbered List}
\begin{enumerate}
  \item First step
  \item Second step
  \item Third step
\end{enumerate}

\section{Citations and References}
When you use the citation search feature, references will be automatically managed. For example, you might cite a paper like this: \cite{example2024}. All citations will be automatically added to your \texttt{cite.bib} file.

\section{Conclusion}
This template demonstrates basic \LaTeX{} features including mathematical formulas, tables, and lists. The bibliography system is now set up to work seamlessly with Bibby AI's citation management. Feel free to modify it for your needs!

\bibliographystyle{plain}
\bibliography{cite}

\end{document}\documentclass{article}
\usepackage[utf8]{inputenc}
\usepackage{booktabs}
\usepackage{amsmath}
\usepackage{natbib}

\title{My Document}
\author{Your Name}
\date{\today}

\begin{document}

\maketitle

\section{Introduction}
This is a sample document demonstrating various \LaTeX{} features including mathematical equations, tables, and formatting. You can cite papers using the citation search feature, and they will be automatically added to your \texttt{cite.bib} file.

\section{Mathematical Expressions}
Here are some example equations:

\subsection{Quadratic Formula}
The quadratic formula is given by:
\[x = \frac{-b \pm \sqrt{b^2 - 4ac}}{2a}\]

\subsection{Matrix Example}
A sample matrix:
\[
\begin{pmatrix}
a_{11} & a_{12} & a_{13} \\
a_{21} & a_{22} & a_{23} \\
a_{31} & a_{32} & a_{33}
\end{pmatrix}
\]

\section{Tables}
Below is a sample table showing data:

\begin{table}[h]
\centering
\begin{tabular}{@{}lrrr@{}}
\toprule
\textbf{Category} & \textbf{Value 1} & \textbf{Value 2} & \textbf{Total} \\
\midrule
Item A & 10.5 & 20.3 & 30.8 \\
Item B & 15.2 & 25.1 & 40.3 \\
Item C & 22.1 & 12.7 & 34.8 \\
\bottomrule
\end{tabular}
\caption{Sample Data Table}
\label{tab:sample}
\end{table}

\section{Lists}
\subsection{Bullet Points}
\begin{itemize}
  \item First item
  \item Second item
  \item Third item with sub-items:
    \begin{itemize}
      \item Sub-item 1
      \item Sub-item 2
    \end{itemize}
\end{itemize}

\subsection{Numbered List}
\begin{enumerate}
  \item First step
  \item Second step
  \item Third step
\end{enumerate}

\section{Citations and References}
When you use the citation search feature, references will be automatically managed. For example, you might cite a paper like this: \cite{example2024}. All citations will be automatically added to your \texttt{cite.bib} file.

\section{Conclusion}
This template demonstrates basic \LaTeX{} features including mathematical formulas, tables, and lists. The bibliography system is now set up to work seamlessly with Bibby AI's citation management. Feel free to modify it for your needs!

\bibliographystyle{plain}
\bibliography{cite}

\end{document}\documentclass{article}
\usepackage[utf8]{inputenc}
\usepackage{booktabs}
\usepackage{amsmath}
\usepackage{natbib}

\title{My Document}
\author{Your Name}
\date{\today}

\begin{document}

\maketitle

\section{Introduction}
This is a sample document demonstrating various \LaTeX{} features including mathematical equations, tables, and formatting. You can cite papers using the citation search feature, and they will be automatically added to your \texttt{cite.bib} file.

\section{Mathematical Expressions}
Here are some example equations:

\subsection{Quadratic Formula}
The quadratic formula is given by:
\[x = \frac{-b \pm \sqrt{b^2 - 4ac}}{2a}\]

\subsection{Matrix Example}
A sample matrix:
\[
\begin{pmatrix}
a_{11} & a_{12} & a_{13} \\
a_{21} & a_{22} & a_{23} \\
a_{31} & a_{32} & a_{33}
\end{pmatrix}
\]

\section{Tables}
Below is a sample table showing data:

\begin{table}[h]
\centering
\begin{tabular}{@{}lrrr@{}}
\toprule
\textbf{Category} & \textbf{Value 1} & \textbf{Value 2} & \textbf{Total} \\
\midrule
Item A & 10.5 & 20.3 & 30.8 \\
Item B & 15.2 & 25.1 & 40.3 \\
Item C & 22.1 & 12.7 & 34.8 \\
\bottomrule
\end{tabular}
\caption{Sample Data Table}
\label{tab:sample}
\end{table}

\section{Lists}
\subsection{Bullet Points}
\begin{itemize}
  \item First item
  \item Second item
  \item Third item with sub-items:
    \begin{itemize}
      \item Sub-item 1
      \item Sub-item 2
    \end{itemize}
\end{itemize}

\subsection{Numbered List}
\begin{enumerate}
  \item First step
  \item Second step
  \item Third step
\end{enumerate}

\section{Citations and References}
When you use the citation search feature, references will be automatically managed. For example, you might cite a paper like this: \cite{example2024}. All citations will be automatically added to your \texttt{cite.bib} file.

\section{Conclusion}
This template demonstrates basic \LaTeX{} features including mathematical formulas, tables, and lists. The bibliography system is now set up to work seamlessly with Bibby AI's citation management. Feel free to modify it for your needs!

\bibliographystyle{plain}
\bibliography{cite}

\end{document}\documentclass{article}
\usepackage[utf8]{inputenc}
\usepackage{booktabs}
\usepackage{amsmath}
\usepackage{natbib}

\title{My Document}
\author{Your Name}
\date{\today}

\begin{document}

\maketitle

\section{Introduction}
This is a sample document demonstrating various \LaTeX{} features including mathematical equations, tables, and formatting. You can cite papers using the citation search feature, and they will be automatically added to your \texttt{cite.bib} file.

\section{Mathematical Expressions}
Here are some example equations:

\subsection{Quadratic Formula}
The quadratic formula is given by:
\[x = \frac{-b \pm \sqrt{b^2 - 4ac}}{2a}\]

\subsection{Matrix Example}
A sample matrix:
\[
\begin{pmatrix}
a_{11} & a_{12} & a_{13} \\
a_{21} & a_{22} & a_{23} \\
a_{31} & a_{32} & a_{33}
\end{pmatrix}
\]

\section{Tables}
Below is a sample table showing data:

\begin{table}[h]
\centering
\begin{tabular}{@{}lrrr@{}}
\toprule
\textbf{Category} & \textbf{Value 1} & \textbf{Value 2} & \textbf{Total} \\
\midrule
Item A & 10.5 & 20.3 & 30.8 \\
Item B & 15.2 & 25.1 & 40.3 \\
Item C & 22.1 & 12.7 & 34.8 \\
\bottomrule
\end{tabular}
\caption{Sample Data Table}
\label{tab:sample}
\end{table}

\section{Lists}
\subsection{Bullet Points}
\begin{itemize}
  \item First item
  \item Second item
  \item Third item with sub-items:
    \begin{itemize}
      \item Sub-item 1
      \item Sub-item 2
    \end{itemize}
\end{itemize}

\subsection{Numbered List}
\begin{enumerate}
  \item First step
  \item Second step
  \item Third step
\end{enumerate}

\section{Citations and References}
When you use the citation search feature, references will be automatically managed. For example, you might cite a paper like this: \cite{example2024}. All citations will be automatically added to your \texttt{cite.bib} file.

\section{Conclusion}
This template demonstrates basic \LaTeX{} features including mathematical formulas, tables, and lists. The bibliography system is now set up to work seamlessly with Bibby AI's citation management. Feel free to modify it for your needs!

\bibliographystyle{plain}
\bibliography{cite}

\end{document}\documentclass{article}
\usepackage[utf8]{inputenc}
\usepackage{booktabs}
\usepackage{amsmath}
\usepackage{natbib}

\title{My Document}
\author{Your Name}
\date{\today}

\begin{document}

\maketitle

\section{Introduction}
This is a sample document demonstrating various \LaTeX{} features including mathematical equations, tables, and formatting. You can cite papers using the citation search feature, and they will be automatically added to your \texttt{cite.bib} file.

\section{Mathematical Expressions}
Here are some example equations:

\subsection{Quadratic Formula}
The quadratic formula is given by:
\[x = \frac{-b \pm \sqrt{b^2 - 4ac}}{2a}\]

\subsection{Matrix Example}
A sample matrix:
\[
\begin{pmatrix}
a_{11} & a_{12} & a_{13} \\
a_{21} & a_{22} & a_{23} \\
a_{31} & a_{32} & a_{33}
\end{pmatrix}
\]

\section{Tables}
Below is a sample table showing data:

\begin{table}[h]
\centering
\begin{tabular}{@{}lrrr@{}}
\toprule
\textbf{Category} & \textbf{Value 1} & \textbf{Value 2} & \textbf{Total} \\
\midrule
Item A & 10.5 & 20.3 & 30.8 \\
Item B & 15.2 & 25.1 & 40.3 \\
Item C & 22.1 & 12.7 & 34.8 \\
\bottomrule
\end{tabular}
\caption{Sample Data Table}
\label{tab:sample}
\end{table}

\section{Lists}
\subsection{Bullet Points}
\begin{itemize}
  \item First item
  \item Second item
  \item Third item with sub-items:
    \begin{itemize}
      \item Sub-item 1
      \item Sub-item 2
    \end{itemize}
\end{itemize}

\subsection{Numbered List}
\begin{enumerate}
  \item First step
  \item Second step
  \item Third step
\end{enumerate}

\section{Citations and References}
When you use the citation search feature, references will be automatically managed. For example, you might cite a paper like this: \cite{example2024}. All citations will be automatically added to your \texttt{cite.bib} file.

\section{Conclusion}
This template demonstrates basic \LaTeX{} features including mathematical formulas, tables, and lists. The bibliography system is now set up to work seamlessly with Bibby AI's citation management. Feel free to modify it for your needs!

\bibliographystyle{plain}
\bibliography{cite}

\end{document}\documentclass{article}
\usepackage[utf8]{inputenc}
\usepackage{booktabs}
\usepackage{amsmath}
\usepackage{natbib}

\title{My Document}
\author{Your Name}
\date{\today}

\begin{document}

\maketitle

\section{Introduction}
This is a sample document demonstrating various \LaTeX{} features including mathematical equations, tables, and formatting. You can cite papers using the citation search feature, and they will be automatically added to your \texttt{cite.bib} file.

\section{Mathematical Expressions}
Here are some example equations:

\subsection{Quadratic Formula}
The quadratic formula is given by:
\[x = \frac{-b \pm \sqrt{b^2 - 4ac}}{2a}\]

\subsection{Matrix Example}
A sample matrix:
\[
\begin{pmatrix}
a_{11} & a_{12} & a_{13} \\
a_{21} & a_{22} & a_{23} \\
a_{31} & a_{32} & a_{33}
\end{pmatrix}
\]

\section{Tables}
Below is a sample table showing data:

\begin{table}[h]
\centering
\begin{tabular}{@{}lrrr@{}}
\toprule
\textbf{Category} & \textbf{Value 1} & \textbf{Value 2} & \textbf{Total} \\
\midrule
Item A & 10.5 & 20.3 & 30.8 \\
Item B & 15.2 & 25.1 & 40.3 \\
Item C & 22.1 & 12.7 & 34.8 \\
\bottomrule
\end{tabular}
\caption{Sample Data Table}
\label{tab:sample}
\end{table}

\section{Lists}
\subsection{Bullet Points}
\begin{itemize}
  \item First item
  \item Second item
  \item Third item with sub-items:
    \begin{itemize}
      \item Sub-item 1
      \item Sub-item 2
    \end{itemize}
\end{itemize}

\subsection{Numbered List}
\begin{enumerate}
  \item First step
  \item Second step
  \item Third step
\end{enumerate}

\section{Citations and References}
When you use the citation search feature, references will be automatically managed. For example, you might cite a paper like this: \cite{example2024}. All citations will be automatically added to your \texttt{cite.bib} file.

\section{Conclusion}
This template demonstrates basic \LaTeX{} features including mathematical formulas, tables, and lists. The bibliography system is now set up to work seamlessly with Bibby AI's citation management. Feel free to modify it for your needs!

\bibliographystyle{plain}
\bibliography{cite}

\end{document}\documentclass{article}
\usepackage[utf8]{inputenc}
\usepackage{booktabs}
\usepackage{amsmath}
\usepackage{natbib}

\title{My Document}
\author{Your Name}
\date{\today}

\begin{document}

\maketitle

\section{Introduction}
This is a sample document demonstrating various \LaTeX{} features including mathematical equations, tables, and formatting. You can cite papers using the citation search feature, and they will be automatically added to your \texttt{cite.bib} file.

\section{Mathematical Expressions}
Here are some example equations:

\subsection{Quadratic Formula}
The quadratic formula is given by:
\[x = \frac{-b \pm \sqrt{b^2 - 4ac}}{2a}\]

\subsection{Matrix Example}
A sample matrix:
\[
\begin{pmatrix}
a_{11} & a_{12} & a_{13} \\
a_{21} & a_{22} & a_{23} \\
a_{31} & a_{32} & a_{33}
\end{pmatrix}
\]

\section{Tables}
Below is a sample table showing data:

\begin{table}[h]
\centering
\begin{tabular}{@{}lrrr@{}}
\toprule
\textbf{Category} & \textbf{Value 1} & \textbf{Value 2} & \textbf{Total} \\
\midrule
Item A & 10.5 & 20.3 & 30.8 \\
Item B & 15.2 & 25.1 & 40.3 \\
Item C & 22.1 & 12.7 & 34.8 \\
\bottomrule
\end{tabular}
\caption{Sample Data Table}
\label{tab:sample}
\end{table}

\section{Lists}
\subsection{Bullet Points}
\begin{itemize}
  \item First item
  \item Second item
  \item Third item with sub-items:
    \begin{itemize}
      \item Sub-item 1
      \item Sub-item 2
    \end{itemize}
\end{itemize}

\subsection{Numbered List}
\begin{enumerate}
  \item First step
  \item Second step
  \item Third step
\end{enumerate}

\section{Citations and References}
When you use the citation search feature, references will be automatically managed. For example, you might cite a paper like this: \cite{example2024}. All citations will be automatically added to your \texttt{cite.bib} file.

\section{Conclusion}
This template demonstrates basic \LaTeX{} features including mathematical formulas, tables, and lists. The bibliography system is now set up to work seamlessly with Bibby AI's citation management. Feel free to modify it for your needs!

\bibliographystyle{plain}
\bibliography{cite}

\end{document}\documentclass{article}
\usepackage[utf8]{inputenc}
\usepackage{booktabs}
\usepackage{amsmath}
\usepackage{natbib}

\title{My Document}
\author{Your Name}
\date{\today}

\begin{document}

\maketitle

\section{Introduction}
This is a sample document demonstrating various \LaTeX{} features including mathematical equations, tables, and formatting. You can cite papers using the citation search feature, and they will be automatically added to your \texttt{cite.bib} file.

\section{Mathematical Expressions}
Here are some example equations:

\subsection{Quadratic Formula}
The quadratic formula is given by:
\[x = \frac{-b \pm \sqrt{b^2 - 4ac}}{2a}\]

\subsection{Matrix Example}
A sample matrix:
\[
\begin{pmatrix}
a_{11} & a_{12} & a_{13} \\
a_{21} & a_{22} & a_{23} \\
a_{31} & a_{32} & a_{33}
\end{pmatrix}
\]

\section{Tables}
Below is a sample table showing data:

\begin{table}[h]
\centering
\begin{tabular}{@{}lrrr@{}}
\toprule
\textbf{Category} & \textbf{Value 1} & \textbf{Value 2} & \textbf{Total} \\
\midrule
Item A & 10.5 & 20.3 & 30.8 \\
Item B & 15.2 & 25.1 & 40.3 \\
Item C & 22.1 & 12.7 & 34.8 \\
\bottomrule
\end{tabular}
\caption{Sample Data Table}
\label{tab:sample}
\end{table}

\section{Lists}
\subsection{Bullet Points}
\begin{itemize}
  \item First item
  \item Second item
  \item Third item with sub-items:
    \begin{itemize}
      \item Sub-item 1
      \item Sub-item 2
    \end{itemize}
\end{itemize}

\subsection{Numbered List}
\begin{enumerate}
  \item First step
  \item Second step
  \item Third step
\end{enumerate}

\section{Citations and References}
When you use the citation search feature, references will be automatically managed. For example, you might cite a paper like this: \cite{example2024}. All citations will be automatically added to your \texttt{cite.bib} file.

\section{Conclusion}
This template demonstrates basic \LaTeX{} features including mathematical formulas, tables, and lists. The bibliography system is now set up to work seamlessly with Bibby AI's citation management. Feel free to modify it for your needs!

\bibliographystyle{plain}
\bibliography{cite}

\end{document}\documentclass{article}
\usepackage[utf8]{inputenc}
\usepackage{booktabs}
\usepackage{amsmath}
\usepackage{natbib}

\title{My Document}
\author{Your Name}
\date{\today}

\begin{document}

\maketitle

\section{Introduction}
This is a sample document demonstrating various \LaTeX{} features including mathematical equations, tables, and formatting. You can cite papers using the citation search feature, and they will be automatically added to your \texttt{cite.bib} file.

\section{Mathematical Expressions}
Here are some example equations:

\subsection{Quadratic Formula}
The quadratic formula is given by:
\[x = \frac{-b \pm \sqrt{b^2 - 4ac}}{2a}\]

\subsection{Matrix Example}
A sample matrix:
\[
\begin{pmatrix}
a_{11} & a_{12} & a_{13} \\
a_{21} & a_{22} & a_{23} \\
a_{31} & a_{32} & a_{33}
\end{pmatrix}
\]

\section{Tables}
Below is a sample table showing data:

\begin{table}[h]
\centering
\begin{tabular}{@{}lrrr@{}}
\toprule
\textbf{Category} & \textbf{Value 1} & \textbf{Value 2} & \textbf{Total} \\
\midrule
Item A & 10.5 & 20.3 & 30.8 \\
Item B & 15.2 & 25.1 & 40.3 \\
Item C & 22.1 & 12.7 & 34.8 \\
\bottomrule
\end{tabular}
\caption{Sample Data Table}
\label{tab:sample}
\end{table}

\section{Lists}
\subsection{Bullet Points}
\begin{itemize}
  \item First item
  \item Second item
  \item Third item with sub-items:
    \begin{itemize}
      \item Sub-item 1
      \item Sub-item 2
    \end{itemize}
\end{itemize}

\subsection{Numbered List}
\begin{enumerate}
  \item First step
  \item Second step
  \item Third step
\end{enumerate}

\section{Citations and References}
When you use the citation search feature, references will be automatically managed. For example, you might cite a paper like this: \cite{example2024}. All citations will be automatically added to your \texttt{cite.bib} file.

\section{Conclusion}
This template demonstrates basic \LaTeX{} features including mathematical formulas, tables, and lists. The bibliography system is now set up to work seamlessly with Bibby AI's citation management. Feel free to modify it for your needs!

\bibliographystyle{plain}
\bibliography{cite}

\end{document}\documentclass{article}
\usepackage[utf8]{inputenc}
\usepackage{booktabs}
\usepackage{amsmath}
\usepackage{natbib}

\title{My Document}
\author{Your Name}
\date{\today}

\begin{document}

\maketitle

\section{Introduction}
This is a sample document demonstrating various \LaTeX{} features including mathematical equations, tables, and formatting. You can cite papers using the citation search feature, and they will be automatically added to your \texttt{cite.bib} file.

\section{Mathematical Expressions}
Here are some example equations:

\subsection{Quadratic Formula}
The quadratic formula is given by:
\[x = \frac{-b \pm \sqrt{b^2 - 4ac}}{2a}\]

\subsection{Matrix Example}
A sample matrix:
\[
\begin{pmatrix}
a_{11} & a_{12} & a_{13} \\
a_{21} & a_{22} & a_{23} \\
a_{31} & a_{32} & a_{33}
\end{pmatrix}
\]

\section{Tables}
Below is a sample table showing data:

\begin{table}[h]
\centering
\begin{tabular}{@{}lrrr@{}}
\toprule
\textbf{Category} & \textbf{Value 1} & \textbf{Value 2} & \textbf{Total} \\
\midrule
Item A & 10.5 & 20.3 & 30.8 \\
Item B & 15.2 & 25.1 & 40.3 \\
Item C & 22.1 & 12.7 & 34.8 \\
\bottomrule
\end{tabular}
\caption{Sample Data Table}
\label{tab:sample}
\end{table}

\section{Lists}
\subsection{Bullet Points}
\begin{itemize}
  \item First item
  \item Second item
  \item Third item with sub-items:
    \begin{itemize}
      \item Sub-item 1
      \item Sub-item 2
    \end{itemize}
\end{itemize}

\subsection{Numbered List}
\begin{enumerate}
  \item First step
  \item Second step
  \item Third step
\end{enumerate}

\section{Citations and References}
When you use the citation search feature, references will be automatically managed. For example, you might cite a paper like this: \cite{example2024}. All citations will be automatically added to your \texttt{cite.bib} file.

\section{Conclusion}
This template demonstrates basic \LaTeX{} features including mathematical formulas, tables, and lists. The bibliography system is now set up to work seamlessly with Bibby AI's citation management. Feel free to modify it for your needs!

\bibliographystyle{plain}
\bibliography{cite}

\end{document}\documentclass{article}
\usepackage[utf8]{inputenc}
\usepackage{booktabs}
\usepackage{amsmath}
\usepackage{natbib}

\title{My Document}
\author{Your Name}
\date{\today}

\begin{document}

\maketitle

\section{Introduction}
This is a sample document demonstrating various \LaTeX{} features including mathematical equations, tables, and formatting. You can cite papers using the citation search feature, and they will be automatically added to your \texttt{cite.bib} file.

\section{Mathematical Expressions}
Here are some example equations:

\subsection{Quadratic Formula}
The quadratic formula is given by:
\[x = \frac{-b \pm \sqrt{b^2 - 4ac}}{2a}\]

\subsection{Matrix Example}
A sample matrix:
\[
\begin{pmatrix}
a_{11} & a_{12} & a_{13} \\
a_{21} & a_{22} & a_{23} \\
a_{31} & a_{32} & a_{33}
\end{pmatrix}
\]

\section{Tables}
Below is a sample table showing data:

\begin{table}[h]
\centering
\begin{tabular}{@{}lrrr@{}}
\toprule
\textbf{Category} & \textbf{Value 1} & \textbf{Value 2} & \textbf{Total} \\
\midrule
Item A & 10.5 & 20.3 & 30.8 \\
Item B & 15.2 & 25.1 & 40.3 \\
Item C & 22.1 & 12.7 & 34.8 \\
\bottomrule
\end{tabular}
\caption{Sample Data Table}
\label{tab:sample}
\end{table}

\section{Lists}
\subsection{Bullet Points}
\begin{itemize}
  \item First item
  \item Second item
  \item Third item with sub-items:
    \begin{itemize}
      \item Sub-item 1
      \item Sub-item 2
    \end{itemize}
\end{itemize}

\subsection{Numbered List}
\begin{enumerate}
  \item First step
  \item Second step
  \item Third step
\end{enumerate}

\section{Citations and References}
When you use the citation search feature, references will be automatically managed. For example, you might cite a paper like this: \cite{example2024}. All citations will be automatically added to your \texttt{cite.bib} file.

\section{Conclusion}
This template demonstrates basic \LaTeX{} features including mathematical formulas, tables, and lists. The bibliography system is now set up to work seamlessly with Bibby AI's citation management. Feel free to modify it for your needs!

\bibliographystyle{plain}
\bibliography{cite}

\end{document}\documentclass{article}
\usepackage[utf8]{inputenc}
\usepackage{booktabs}
\usepackage{amsmath}
\usepackage{natbib}

\title{My Document}
\author{Your Name}
\date{\today}

\begin{document}

\maketitle

\section{Introduction}
This is a sample document demonstrating various \LaTeX{} features including mathematical equations, tables, and formatting. You can cite papers using the citation search feature, and they will be automatically added to your \texttt{cite.bib} file.

\section{Mathematical Expressions}
Here are some example equations:

\subsection{Quadratic Formula}
The quadratic formula is given by:
\[x = \frac{-b \pm \sqrt{b^2 - 4ac}}{2a}\]

\subsection{Matrix Example}
A sample matrix:
\[
\begin{pmatrix}
a_{11} & a_{12} & a_{13} \\
a_{21} & a_{22} & a_{23} \\
a_{31} & a_{32} & a_{33}
\end{pmatrix}
\]

\section{Tables}
Below is a sample table showing data:

\begin{table}[h]
\centering
\begin{tabular}{@{}lrrr@{}}
\toprule
\textbf{Category} & \textbf{Value 1} & \textbf{Value 2} & \textbf{Total} \\
\midrule
Item A & 10.5 & 20.3 & 30.8 \\
Item B & 15.2 & 25.1 & 40.3 \\
Item C & 22.1 & 12.7 & 34.8 \\
\bottomrule
\end{tabular}
\caption{Sample Data Table}
\label{tab:sample}
\end{table}

\section{Lists}
\subsection{Bullet Points}
\begin{itemize}
  \item First item
  \item Second item
  \item Third item with sub-items:
    \begin{itemize}
      \item Sub-item 1
      \item Sub-item 2
    \end{itemize}
\end{itemize}

\subsection{Numbered List}
\begin{enumerate}
  \item First step
  \item Second step
  \item Third step
\end{enumerate}

\section{Citations and References}
When you use the citation search feature, references will be automatically managed. For example, you might cite a paper like this: \cite{example2024}. All citations will be automatically added to your \texttt{cite.bib} file.

\section{Conclusion}
This template demonstrates basic \LaTeX{} features including mathematical formulas, tables, and lists. The bibliography system is now set up to work seamlessly with Bibby AI's citation management. Feel free to modify it for your needs!

\bibliographystyle{plain}
\bibliography{cite}

\end{document}\documentclass{article}
\usepackage[utf8]{inputenc}
\usepackage{booktabs}
\usepackage{amsmath}
\usepackage{natbib}

\title{My Document}
\author{Your Name}
\date{\today}

\begin{document}

\maketitle

\section{Introduction}
This is a sample document demonstrating various \LaTeX{} features including mathematical equations, tables, and formatting. You can cite papers using the citation search feature, and they will be automatically added to your \texttt{cite.bib} file.

\section{Mathematical Expressions}
Here are some example equations:

\subsection{Quadratic Formula}
The quadratic formula is given by:
\[x = \frac{-b \pm \sqrt{b^2 - 4ac}}{2a}\]

\subsection{Matrix Example}
A sample matrix:
\[
\begin{pmatrix}
a_{11} & a_{12} & a_{13} \\
a_{21} & a_{22} & a_{23} \\
a_{31} & a_{32} & a_{33}
\end{pmatrix}
\]

\section{Tables}
Below is a sample table showing data:

\begin{table}[h]
\centering
\begin{tabular}{@{}lrrr@{}}
\toprule
\textbf{Category} & \textbf{Value 1} & \textbf{Value 2} & \textbf{Total} \\
\midrule
Item A & 10.5 & 20.3 & 30.8 \\
Item B & 15.2 & 25.1 & 40.3 \\
Item C & 22.1 & 12.7 & 34.8 \\
\bottomrule
\end{tabular}
\caption{Sample Data Table}
\label{tab:sample}
\end{table}

\section{Lists}
\subsection{Bullet Points}
\begin{itemize}
  \item First item
  \item Second item
  \item Third item with sub-items:
    \begin{itemize}
      \item Sub-item 1
      \item Sub-item 2
    \end{itemize}
\end{itemize}

\subsection{Numbered List}
\begin{enumerate}
  \item First step
  \item Second step
  \item Third step
\end{enumerate}

\section{Citations and References}
When you use the citation search feature, references will be automatically managed. For example, you might cite a paper like this: \cite{example2024}. All citations will be automatically added to your \texttt{cite.bib} file.

\section{Conclusion}
This template demonstrates basic \LaTeX{} features including mathematical formulas, tables, and lists. The bibliography system is now set up to work seamlessly with Bibby AI's citation management. Feel free to modify it for your needs!

\bibliographystyle{plain}
\bibliography{cite}

\end{document}\documentclass{article}
\usepackage[utf8]{inputenc}
\usepackage{booktabs}
\usepackage{amsmath}
\usepackage{natbib}

\title{My Document}
\author{Your Name}
\date{\today}

\begin{document}

\maketitle

\section{Introduction}
This is a sample document demonstrating various \LaTeX{} features including mathematical equations, tables, and formatting. You can cite papers using the citation search feature, and they will be automatically added to your \texttt{cite.bib} file.

\section{Mathematical Expressions}
Here are some example equations:

\subsection{Quadratic Formula}
The quadratic formula is given by:
\[x = \frac{-b \pm \sqrt{b^2 - 4ac}}{2a}\]

\subsection{Matrix Example}
A sample matrix:
\[
\begin{pmatrix}
a_{11} & a_{12} & a_{13} \\
a_{21} & a_{22} & a_{23} \\
a_{31} & a_{32} & a_{33}
\end{pmatrix}
\]

\section{Tables}
Below is a sample table showing data:

\begin{table}[h]
\centering
\begin{tabular}{@{}lrrr@{}}
\toprule
\textbf{Category} & \textbf{Value 1} & \textbf{Value 2} & \textbf{Total} \\
\midrule
Item A & 10.5 & 20.3 & 30.8 \\
Item B & 15.2 & 25.1 & 40.3 \\
Item C & 22.1 & 12.7 & 34.8 \\
\bottomrule
\end{tabular}
\caption{Sample Data Table}
\label{tab:sample}
\end{table}

\section{Lists}
\subsection{Bullet Points}
\begin{itemize}
  \item First item
  \item Second item
  \item Third item with sub-items:
    \begin{itemize}
      \item Sub-item 1
      \item Sub-item 2
    \end{itemize}
\end{itemize}

\subsection{Numbered List}
\begin{enumerate}
  \item First step
  \item Second step
  \item Third step
\end{enumerate}

\section{Citations and References}
When you use the citation search feature, references will be automatically managed. For example, you might cite a paper like this: \cite{example2024}. All citations will be automatically added to your \texttt{cite.bib} file.

\section{Conclusion}
This template demonstrates basic \LaTeX{} features including mathematical formulas, tables, and lists. The bibliography system is now set up to work seamlessly with Bibby AI's citation management. Feel free to modify it for your needs!

\bibliographystyle{plain}
\bibliography{cite}

\end{document}\documentclass{article}
\usepackage[utf8]{inputenc}
\usepackage{booktabs}
\usepackage{amsmath}
\usepackage{natbib}

\title{My Document}
\author{Your Name}
\date{\today}

\begin{document}

\maketitle

\section{Introduction}
This is a sample document demonstrating various \LaTeX{} features including mathematical equations, tables, and formatting. You can cite papers using the citation search feature, and they will be automatically added to your \texttt{cite.bib} file.

\section{Mathematical Expressions}
Here are some example equations:

\subsection{Quadratic Formula}
The quadratic formula is given by:
\[x = \frac{-b \pm \sqrt{b^2 - 4ac}}{2a}\]

\subsection{Matrix Example}
A sample matrix:
\[
\begin{pmatrix}
a_{11} & a_{12} & a_{13} \\
a_{21} & a_{22} & a_{23} \\
a_{31} & a_{32} & a_{33}
\end{pmatrix}
\]

\section{Tables}
Below is a sample table showing data:

\begin{table}[h]
\centering
\begin{tabular}{@{}lrrr@{}}
\toprule
\textbf{Category} & \textbf{Value 1} & \textbf{Value 2} & \textbf{Total} \\
\midrule
Item A & 10.5 & 20.3 & 30.8 \\
Item B & 15.2 & 25.1 & 40.3 \\
Item C & 22.1 & 12.7 & 34.8 \\
\bottomrule
\end{tabular}
\caption{Sample Data Table}
\label{tab:sample}
\end{table}

\section{Lists}
\subsection{Bullet Points}
\begin{itemize}
  \item First item
  \item Second item
  \item Third item with sub-items:
    \begin{itemize}
      \item Sub-item 1
      \item Sub-item 2
    \end{itemize}
\end{itemize}

\subsection{Numbered List}
\begin{enumerate}
  \item First step
  \item Second step
  \item Third step
\end{enumerate}

\section{Citations and References}
When you use the citation search feature, references will be automatically managed. For example, you might cite a paper like this: \cite{example2024}. All citations will be automatically added to your \texttt{cite.bib} file.

\section{Conclusion}
This template demonstrates basic \LaTeX{} features including mathematical formulas, tables, and lists. The bibliography system is now set up to work seamlessly with Bibby AI's citation management. Feel free to modify it for your needs!

\bibliographystyle{plain}
\bibliography{cite}

\end{document}\documentclass{article}
\usepackage[utf8]{inputenc}
\usepackage{booktabs}
\usepackage{amsmath}
\usepackage{natbib}

\title{My Document}
\author{Your Name}
\date{\today}

\begin{document}

\maketitle

\section{Introduction}
This is a sample document demonstrating various \LaTeX{} features including mathematical equations, tables, and formatting. You can cite papers using the citation search feature, and they will be automatically added to your \texttt{cite.bib} file.

\section{Mathematical Expressions}
Here are some example equations:

\subsection{Quadratic Formula}
The quadratic formula is given by:
\[x = \frac{-b \pm \sqrt{b^2 - 4ac}}{2a}\]

\subsection{Matrix Example}
A sample matrix:
\[
\begin{pmatrix}
a_{11} & a_{12} & a_{13} \\
a_{21} & a_{22} & a_{23} \\
a_{31} & a_{32} & a_{33}
\end{pmatrix}
\]

\section{Tables}
Below is a sample table showing data:

\begin{table}[h]
\centering
\begin{tabular}{@{}lrrr@{}}
\toprule
\textbf{Category} & \textbf{Value 1} & \textbf{Value 2} & \textbf{Total} \\
\midrule
Item A & 10.5 & 20.3 & 30.8 \\
Item B & 15.2 & 25.1 & 40.3 \\
Item C & 22.1 & 12.7 & 34.8 \\
\bottomrule
\end{tabular}
\caption{Sample Data Table}
\label{tab:sample}
\end{table}

\section{Lists}
\subsection{Bullet Points}
\begin{itemize}
  \item First item
  \item Second item
  \item Third item with sub-items:
    \begin{itemize}
      \item Sub-item 1
      \item Sub-item 2
    \end{itemize}
\end{itemize}

\subsection{Numbered List}
\begin{enumerate}
  \item First step
  \item Second step
  \item Third step
\end{enumerate}

\section{Citations and References}
When you use the citation search feature, references will be automatically managed. For example, you might cite a paper like this: \cite{example2024}. All citations will be automatically added to your \texttt{cite.bib} file.

\section{Conclusion}
This template demonstrates basic \LaTeX{} features including mathematical formulas, tables, and lists. The bibliography system is now set up to work seamlessly with Bibby AI's citation management. Feel free to modify it for your needs!

\bibliographystyle{plain}
\bibliography{cite}

\end{document}\documentclass{article}
\usepackage[utf8]{inputenc}
\usepackage{booktabs}
\usepackage{amsmath}
\usepackage{natbib}

\title{My Document}
\author{Your Name}
\date{\today}

\begin{document}

\maketitle

\section{Introduction}
This is a sample document demonstrating various \LaTeX{} features including mathematical equations, tables, and formatting. You can cite papers using the citation search feature, and they will be automatically added to your \texttt{cite.bib} file.

\section{Mathematical Expressions}
Here are some example equations:

\subsection{Quadratic Formula}
The quadratic formula is given by:
\[x = \frac{-b \pm \sqrt{b^2 - 4ac}}{2a}\]

\subsection{Matrix Example}
A sample matrix:
\[
\begin{pmatrix}
a_{11} & a_{12} & a_{13} \\
a_{21} & a_{22} & a_{23} \\
a_{31} & a_{32} & a_{33}
\end{pmatrix}
\]

\section{Tables}
Below is a sample table showing data:

\begin{table}[h]
\centering
\begin{tabular}{@{}lrrr@{}}
\toprule
\textbf{Category} & \textbf{Value 1} & \textbf{Value 2} & \textbf{Total} \\
\midrule
Item A & 10.5 & 20.3 & 30.8 \\
Item B & 15.2 & 25.1 & 40.3 \\
Item C & 22.1 & 12.7 & 34.8 \\
\bottomrule
\end{tabular}
\caption{Sample Data Table}
\label{tab:sample}
\end{table}

\section{Lists}
\subsection{Bullet Points}
\begin{itemize}
  \item First item
  \item Second item
  \item Third item with sub-items:
    \begin{itemize}
      \item Sub-item 1
      \item Sub-item 2
    \end{itemize}
\end{itemize}

\subsection{Numbered List}
\begin{enumerate}
  \item First step
  \item Second step
  \item Third step
\end{enumerate}

\section{Citations and References}
When you use the citation search feature, references will be automatically managed. For example, you might cite a paper like this: \cite{example2024}. All citations will be automatically added to your \texttt{cite.bib} file.

\section{Conclusion}
This template demonstrates basic \LaTeX{} features including mathematical formulas, tables, and lists. The bibliography system is now set up to work seamlessly with Bibby AI's citation management. Feel free to modify it for your needs!

\bibliographystyle{plain}
\bibliography{cite}

\end{document}\documentclass{article}
\usepackage[utf8]{inputenc}
\usepackage{booktabs}
\usepackage{amsmath}
\usepackage{natbib}

\title{My Document}
\author{Your Name}
\date{\today}

\begin{document}

\maketitle

\section{Introduction}
This is a sample document demonstrating various \LaTeX{} features including mathematical equations, tables, and formatting. You can cite papers using the citation search feature, and they will be automatically added to your \texttt{cite.bib} file.

\section{Mathematical Expressions}
Here are some example equations:

\subsection{Quadratic Formula}
The quadratic formula is given by:
\[x = \frac{-b \pm \sqrt{b^2 - 4ac}}{2a}\]

\subsection{Matrix Example}
A sample matrix:
\[
\begin{pmatrix}
a_{11} & a_{12} & a_{13} \\
a_{21} & a_{22} & a_{23} \\
a_{31} & a_{32} & a_{33}
\end{pmatrix}
\]

\section{Tables}
Below is a sample table showing data:

\begin{table}[h]
\centering
\begin{tabular}{@{}lrrr@{}}
\toprule
\textbf{Category} & \textbf{Value 1} & \textbf{Value 2} & \textbf{Total} \\
\midrule
Item A & 10.5 & 20.3 & 30.8 \\
Item B & 15.2 & 25.1 & 40.3 \\
Item C & 22.1 & 12.7 & 34.8 \\
\bottomrule
\end{tabular}
\caption{Sample Data Table}
\label{tab:sample}
\end{table}

\section{Lists}
\subsection{Bullet Points}
\begin{itemize}
  \item First item
  \item Second item
  \item Third item with sub-items:
    \begin{itemize}
      \item Sub-item 1
      \item Sub-item 2
    \end{itemize}
\end{itemize}

\subsection{Numbered List}
\begin{enumerate}
  \item First step
  \item Second step
  \item Third step
\end{enumerate}

\section{Citations and References}
When you use the citation search feature, references will be automatically managed. For example, you might cite a paper like this: \cite{example2024}. All citations will be automatically added to your \texttt{cite.bib} file.

\section{Conclusion}
This template demonstrates basic \LaTeX{} features including mathematical formulas, tables, and lists. The bibliography system is now set up to work seamlessly with Bibby AI's citation management. Feel free to modify it for your needs!

\bibliographystyle{plain}
\bibliography{cite}

\end{document}\documentclass{article}
\usepackage[utf8]{inputenc}
\usepackage{booktabs}
\usepackage{amsmath}
\usepackage{natbib}

\title{My Document}
\author{Your Name}
\date{\today}

\begin{document}

\maketitle

\section{Introduction}
This is a sample document demonstrating various \LaTeX{} features including mathematical equations, tables, and formatting. You can cite papers using the citation search feature, and they will be automatically added to your \texttt{cite.bib} file.

\section{Mathematical Expressions}
Here are some example equations:

\subsection{Quadratic Formula}
The quadratic formula is given by:
\[x = \frac{-b \pm \sqrt{b^2 - 4ac}}{2a}\]

\subsection{Matrix Example}
A sample matrix:
\[
\begin{pmatrix}
a_{11} & a_{12} & a_{13} \\
a_{21} & a_{22} & a_{23} \\
a_{31} & a_{32} & a_{33}
\end{pmatrix}
\]

\section{Tables}
Below is a sample table showing data:

\begin{table}[h]
\centering
\begin{tabular}{@{}lrrr@{}}
\toprule
\textbf{Category} & \textbf{Value 1} & \textbf{Value 2} & \textbf{Total} \\
\midrule
Item A & 10.5 & 20.3 & 30.8 \\
Item B & 15.2 & 25.1 & 40.3 \\
Item C & 22.1 & 12.7 & 34.8 \\
\bottomrule
\end{tabular}
\caption{Sample Data Table}
\label{tab:sample}
\end{table}

\section{Lists}
\subsection{Bullet Points}
\begin{itemize}
  \item First item
  \item Second item
  \item Third item with sub-items:
    \begin{itemize}
      \item Sub-item 1
      \item Sub-item 2
    \end{itemize}
\end{itemize}

\subsection{Numbered List}
\begin{enumerate}
  \item First step
  \item Second step
  \item Third step
\end{enumerate}

\section{Citations and References}
When you use the citation search feature, references will be automatically managed. For example, you might cite a paper like this: \cite{example2024}. All citations will be automatically added to your \texttt{cite.bib} file.

\section{Conclusion}
This template demonstrates basic \LaTeX{} features including mathematical formulas, tables, and lists. The bibliography system is now set up to work seamlessly with Bibby AI's citation management. Feel free to modify it for your needs!

\bibliographystyle{plain}
\bibliography{cite}

\end{document}\documentclass{article}
\usepackage[utf8]{inputenc}
\usepackage{booktabs}
\usepackage{amsmath}
\usepackage{natbib}

\title{My Document}
\author{Your Name}
\date{\today}

\begin{document}

\maketitle

\section{Introduction}
This is a sample document demonstrating various \LaTeX{} features including mathematical equations, tables, and formatting. You can cite papers using the citation search feature, and they will be automatically added to your \texttt{cite.bib} file.

\section{Mathematical Expressions}
Here are some example equations:

\subsection{Quadratic Formula}
The quadratic formula is given by:
\[x = \frac{-b \pm \sqrt{b^2 - 4ac}}{2a}\]

\subsection{Matrix Example}
A sample matrix:
\[
\begin{pmatrix}
a_{11} & a_{12} & a_{13} \\
a_{21} & a_{22} & a_{23} \\
a_{31} & a_{32} & a_{33}
\end{pmatrix}
\]

\section{Tables}
Below is a sample table showing data:

\begin{table}[h]
\centering
\begin{tabular}{@{}lrrr@{}}
\toprule
\textbf{Category} & \textbf{Value 1} & \textbf{Value 2} & \textbf{Total} \\
\midrule
Item A & 10.5 & 20.3 & 30.8 \\
Item B & 15.2 & 25.1 & 40.3 \\
Item C & 22.1 & 12.7 & 34.8 \\
\bottomrule
\end{tabular}
\caption{Sample Data Table}
\label{tab:sample}
\end{table}

\section{Lists}
\subsection{Bullet Points}
\begin{itemize}
  \item First item
  \item Second item
  \item Third item with sub-items:
    \begin{itemize}
      \item Sub-item 1
      \item Sub-item 2
    \end{itemize}
\end{itemize}

\subsection{Numbered List}
\begin{enumerate}
  \item First step
  \item Second step
  \item Third step
\end{enumerate}

\section{Citations and References}
When you use the citation search feature, references will be automatically managed. For example, you might cite a paper like this: \cite{example2024}. All citations will be automatically added to your \texttt{cite.bib} file.

\section{Conclusion}
This template demonstrates basic \LaTeX{} features including mathematical formulas, tables, and lists. The bibliography system is now set up to work seamlessly with Bibby AI's citation management. Feel free to modify it for your needs!

\bibliographystyle{plain}
\bibliography{cite}

\end{document}\documentclass{article}
\usepackage[utf8]{inputenc}
\usepackage{booktabs}
\usepackage{amsmath}
\usepackage{natbib}

\title{My Document}
\author{Your Name}
\date{\today}

\begin{document}

\maketitle

\section{Introduction}
This is a sample document demonstrating various \LaTeX{} features including mathematical equations, tables, and formatting. You can cite papers using the citation search feature, and they will be automatically added to your \texttt{cite.bib} file.

\section{Mathematical Expressions}
Here are some example equations:

\subsection{Quadratic Formula}
The quadratic formula is given by:
\[x = \frac{-b \pm \sqrt{b^2 - 4ac}}{2a}\]

\subsection{Matrix Example}
A sample matrix:
\[
\begin{pmatrix}
a_{11} & a_{12} & a_{13} \\
a_{21} & a_{22} & a_{23} \\
a_{31} & a_{32} & a_{33}
\end{pmatrix}
\]

\section{Tables}
Below is a sample table showing data:

\begin{table}[h]
\centering
\begin{tabular}{@{}lrrr@{}}
\toprule
\textbf{Category} & \textbf{Value 1} & \textbf{Value 2} & \textbf{Total} \\
\midrule
Item A & 10.5 & 20.3 & 30.8 \\
Item B & 15.2 & 25.1 & 40.3 \\
Item C & 22.1 & 12.7 & 34.8 \\
\bottomrule
\end{tabular}
\caption{Sample Data Table}
\label{tab:sample}
\end{table}

\section{Lists}
\subsection{Bullet Points}
\begin{itemize}
  \item First item
  \item Second item
  \item Third item with sub-items:
    \begin{itemize}
      \item Sub-item 1
      \item Sub-item 2
    \end{itemize}
\end{itemize}

\subsection{Numbered List}
\begin{enumerate}
  \item First step
  \item Second step
  \item Third step
\end{enumerate}

\section{Citations and References}
When you use the citation search feature, references will be automatically managed. For example, you might cite a paper like this: \cite{example2024}. All citations will be automatically added to your \texttt{cite.bib} file.

\section{Conclusion}
This template demonstrates basic \LaTeX{} features including mathematical formulas, tables, and lists. The bibliography system is now set up to work seamlessly with Bibby AI's citation management. Feel free to modify it for your needs!

\bibliographystyle{plain}
\bibliography{cite}

\end{document}\documentclass{article}
\usepackage[utf8]{inputenc}
\usepackage{booktabs}
\usepackage{amsmath}
\usepackage{natbib}

\title{My Document}
\author{Your Name}
\date{\today}

\begin{document}

\maketitle

\section{Introduction}
This is a sample document demonstrating various \LaTeX{} features including mathematical equations, tables, and formatting. You can cite papers using the citation search feature, and they will be automatically added to your \texttt{cite.bib} file.

\section{Mathematical Expressions}
Here are some example equations:

\subsection{Quadratic Formula}
The quadratic formula is given by:
\[x = \frac{-b \pm \sqrt{b^2 - 4ac}}{2a}\]

\subsection{Matrix Example}
A sample matrix:
\[
\begin{pmatrix}
a_{11} & a_{12} & a_{13} \\
a_{21} & a_{22} & a_{23} \\
a_{31} & a_{32} & a_{33}
\end{pmatrix}
\]

\section{Tables}
Below is a sample table showing data:

\begin{table}[h]
\centering
\begin{tabular}{@{}lrrr@{}}
\toprule
\textbf{Category} & \textbf{Value 1} & \textbf{Value 2} & \textbf{Total} \\
\midrule
Item A & 10.5 & 20.3 & 30.8 \\
Item B & 15.2 & 25.1 & 40.3 \\
Item C & 22.1 & 12.7 & 34.8 \\
\bottomrule
\end{tabular}
\caption{Sample Data Table}
\label{tab:sample}
\end{table}

\section{Lists}
\subsection{Bullet Points}
\begin{itemize}
  \item First item
  \item Second item
  \item Third item with sub-items:
    \begin{itemize}
      \item Sub-item 1
      \item Sub-item 2
    \end{itemize}
\end{itemize}

\subsection{Numbered List}
\begin{enumerate}
  \item First step
  \item Second step
  \item Third step
\end{enumerate}

\section{Citations and References}
When you use the citation search feature, references will be automatically managed. For example, you might cite a paper like this: \cite{example2024}. All citations will be automatically added to your \texttt{cite.bib} file.

\section{Conclusion}
This template demonstrates basic \LaTeX{} features including mathematical formulas, tables, and lists. The bibliography system is now set up to work seamlessly with Bibby AI's citation management. Feel free to modify it for your needs!

\bibliographystyle{plain}
\bibliography{cite}

\end{document}\documentclass{article}
\usepackage[utf8]{inputenc}
\usepackage{booktabs}
\usepackage{amsmath}
\usepackage{natbib}

\title{My Document}
\author{Your Name}
\date{\today}

\begin{document}

\maketitle

\section{Introduction}
This is a sample document demonstrating various \LaTeX{} features including mathematical equations, tables, and formatting. You can cite papers using the citation search feature, and they will be automatically added to your \texttt{cite.bib} file.

\section{Mathematical Expressions}
Here are some example equations:

\subsection{Quadratic Formula}
The quadratic formula is given by:
\[x = \frac{-b \pm \sqrt{b^2 - 4ac}}{2a}\]

\subsection{Matrix Example}
A sample matrix:
\[
\begin{pmatrix}
a_{11} & a_{12} & a_{13} \\
a_{21} & a_{22} & a_{23} \\
a_{31} & a_{32} & a_{33}
\end{pmatrix}
\]

\section{Tables}
Below is a sample table showing data:

\begin{table}[h]
\centering
\begin{tabular}{@{}lrrr@{}}
\toprule
\textbf{Category} & \textbf{Value 1} & \textbf{Value 2} & \textbf{Total} \\
\midrule
Item A & 10.5 & 20.3 & 30.8 \\
Item B & 15.2 & 25.1 & 40.3 \\
Item C & 22.1 & 12.7 & 34.8 \\
\bottomrule
\end{tabular}
\caption{Sample Data Table}
\label{tab:sample}
\end{table}

\section{Lists}
\subsection{Bullet Points}
\begin{itemize}
  \item First item
  \item Second item
  \item Third item with sub-items:
    \begin{itemize}
      \item Sub-item 1
      \item Sub-item 2
    \end{itemize}
\end{itemize}

\subsection{Numbered List}
\begin{enumerate}
  \item First step
  \item Second step
  \item Third step
\end{enumerate}

\section{Citations and References}
When you use the citation search feature, references will be automatically managed. For example, you might cite a paper like this: \cite{example2024}. All citations will be automatically added to your \texttt{cite.bib} file.

\section{Conclusion}
This template demonstrates basic \LaTeX{} features including mathematical formulas, tables, and lists. The bibliography system is now set up to work seamlessly with Bibby AI's citation management. Feel free to modify it for your needs!

\bibliographystyle{plain}
\bibliography{cite}

\end{document}\documentclass{article}
\usepackage[utf8]{inputenc}
\usepackage{booktabs}
\usepackage{amsmath}
\usepackage{natbib}

\title{My Document}
\author{Your Name}
\date{\today}

\begin{document}

\maketitle

\section{Introduction}
This is a sample document demonstrating various \LaTeX{} features including mathematical equations, tables, and formatting. You can cite papers using the citation search feature, and they will be automatically added to your \texttt{cite.bib} file.

\section{Mathematical Expressions}
Here are some example equations:

\subsection{Quadratic Formula}
The quadratic formula is given by:
\[x = \frac{-b \pm \sqrt{b^2 - 4ac}}{2a}\]

\subsection{Matrix Example}
A sample matrix:
\[
\begin{pmatrix}
a_{11} & a_{12} & a_{13} \\
a_{21} & a_{22} & a_{23} \\
a_{31} & a_{32} & a_{33}
\end{pmatrix}
\]

\section{Tables}
Below is a sample table showing data:

\begin{table}[h]
\centering
\begin{tabular}{@{}lrrr@{}}
\toprule
\textbf{Category} & \textbf{Value 1} & \textbf{Value 2} & \textbf{Total} \\
\midrule
Item A & 10.5 & 20.3 & 30.8 \\
Item B & 15.2 & 25.1 & 40.3 \\
Item C & 22.1 & 12.7 & 34.8 \\
\bottomrule
\end{tabular}
\caption{Sample Data Table}
\label{tab:sample}
\end{table}

\section{Lists}
\subsection{Bullet Points}
\begin{itemize}
  \item First item
  \item Second item
  \item Third item with sub-items:
    \begin{itemize}
      \item Sub-item 1
      \item Sub-item 2
    \end{itemize}
\end{itemize}

\subsection{Numbered List}
\begin{enumerate}
  \item First step
  \item Second step
  \item Third step
\end{enumerate}

\section{Citations and References}
When you use the citation search feature, references will be automatically managed. For example, you might cite a paper like this: \cite{example2024}. All citations will be automatically added to your \texttt{cite.bib} file.

\section{Conclusion}
This template demonstrates basic \LaTeX{} features including mathematical formulas, tables, and lists. The bibliography system is now set up to work seamlessly with Bibby AI's citation management. Feel free to modify it for your needs!

\bibliographystyle{plain}
\bibliography{cite}

\end{document}\documentclass{article}
\usepackage[utf8]{inputenc}
\usepackage{booktabs}
\usepackage{amsmath}
\usepackage{natbib}

\title{My Document}
\author{Your Name}
\date{\today}

\begin{document}

\maketitle

\section{Introduction}
This is a sample document demonstrating various \LaTeX{} features including mathematical equations, tables, and formatting. You can cite papers using the citation search feature, and they will be automatically added to your \texttt{cite.bib} file.

\section{Mathematical Expressions}
Here are some example equations:

\subsection{Quadratic Formula}
The quadratic formula is given by:
\[x = \frac{-b \pm \sqrt{b^2 - 4ac}}{2a}\]

\subsection{Matrix Example}
A sample matrix:
\[
\begin{pmatrix}
a_{11} & a_{12} & a_{13} \\
a_{21} & a_{22} & a_{23} \\
a_{31} & a_{32} & a_{33}
\end{pmatrix}
\]

\section{Tables}
Below is a sample table showing data:

\begin{table}[h]
\centering
\begin{tabular}{@{}lrrr@{}}
\toprule
\textbf{Category} & \textbf{Value 1} & \textbf{Value 2} & \textbf{Total} \\
\midrule
Item A & 10.5 & 20.3 & 30.8 \\
Item B & 15.2 & 25.1 & 40.3 \\
Item C & 22.1 & 12.7 & 34.8 \\
\bottomrule
\end{tabular}
\caption{Sample Data Table}
\label{tab:sample}
\end{table}

\section{Lists}
\subsection{Bullet Points}
\begin{itemize}
  \item First item
  \item Second item
  \item Third item with sub-items:
    \begin{itemize}
      \item Sub-item 1
      \item Sub-item 2
    \end{itemize}
\end{itemize}

\subsection{Numbered List}
\begin{enumerate}
  \item First step
  \item Second step
  \item Third step
\end{enumerate}

\section{Citations and References}
When you use the citation search feature, references will be automatically managed. For example, you might cite a paper like this: \cite{example2024}. All citations will be automatically added to your \texttt{cite.bib} file.

\section{Conclusion}
This template demonstrates basic \LaTeX{} features including mathematical formulas, tables, and lists. The bibliography system is now set up to work seamlessly with Bibby AI's citation management. Feel free to modify it for your needs!

\bibliographystyle{plain}
\bibliography{cite}

\end{document}\documentclass{article}
\usepackage[utf8]{inputenc}
\usepackage{booktabs}
\usepackage{amsmath}
\usepackage{natbib}

\title{My Document}
\author{Your Name}
\date{\today}

\begin{document}

\maketitle

\section{Introduction}
This is a sample document demonstrating various \LaTeX{} features including mathematical equations, tables, and formatting. You can cite papers using the citation search feature, and they will be automatically added to your \texttt{cite.bib} file.

\section{Mathematical Expressions}
Here are some example equations:

\subsection{Quadratic Formula}
The quadratic formula is given by:
\[x = \frac{-b \pm \sqrt{b^2 - 4ac}}{2a}\]

\subsection{Matrix Example}
A sample matrix:
\[
\begin{pmatrix}
a_{11} & a_{12} & a_{13} \\
a_{21} & a_{22} & a_{23} \\
a_{31} & a_{32} & a_{33}
\end{pmatrix}
\]

\section{Tables}
Below is a sample table showing data:

\begin{table}[h]
\centering
\begin{tabular}{@{}lrrr@{}}
\toprule
\textbf{Category} & \textbf{Value 1} & \textbf{Value 2} & \textbf{Total} \\
\midrule
Item A & 10.5 & 20.3 & 30.8 \\
Item B & 15.2 & 25.1 & 40.3 \\
Item C & 22.1 & 12.7 & 34.8 \\
\bottomrule
\end{tabular}
\caption{Sample Data Table}
\label{tab:sample}
\end{table}

\section{Lists}
\subsection{Bullet Points}
\begin{itemize}
  \item First item
  \item Second item
  \item Third item with sub-items:
    \begin{itemize}
      \item Sub-item 1
      \item Sub-item 2
    \end{itemize}
\end{itemize}

\subsection{Numbered List}
\begin{enumerate}
  \item First step
  \item Second step
  \item Third step
\end{enumerate}

\section{Citations and References}
When you use the citation search feature, references will be automatically managed. For example, you might cite a paper like this: \cite{example2024}. All citations will be automatically added to your \texttt{cite.bib} file.

\section{Conclusion}
This template demonstrates basic \LaTeX{} features including mathematical formulas, tables, and lists. The bibliography system is now set up to work seamlessly with Bibby AI's citation management. Feel free to modify it for your needs!

\bibliographystyle{plain}
\bibliography{cite}

\end{document}\documentclass{article}
\usepackage[utf8]{inputenc}
\usepackage{booktabs}
\usepackage{amsmath}
\usepackage{natbib}

\title{My Document}
\author{Your Name}
\date{\today}

\begin{document}

\maketitle

\section{Introduction}
This is a sample document demonstrating various \LaTeX{} features including mathematical equations, tables, and formatting. You can cite papers using the citation search feature, and they will be automatically added to your \texttt{cite.bib} file.

\section{Mathematical Expressions}
Here are some example equations:

\subsection{Quadratic Formula}
The quadratic formula is given by:
\[x = \frac{-b \pm \sqrt{b^2 - 4ac}}{2a}\]

\subsection{Matrix Example}
A sample matrix:
\[
\begin{pmatrix}
a_{11} & a_{12} & a_{13} \\
a_{21} & a_{22} & a_{23} \\
a_{31} & a_{32} & a_{33}
\end{pmatrix}
\]

\section{Tables}
Below is a sample table showing data:

\begin{table}[h]
\centering
\begin{tabular}{@{}lrrr@{}}
\toprule
\textbf{Category} & \textbf{Value 1} & \textbf{Value 2} & \textbf{Total} \\
\midrule
Item A & 10.5 & 20.3 & 30.8 \\
Item B & 15.2 & 25.1 & 40.3 \\
Item C & 22.1 & 12.7 & 34.8 \\
\bottomrule
\end{tabular}
\caption{Sample Data Table}
\label{tab:sample}
\end{table}

\section{Lists}
\subsection{Bullet Points}
\begin{itemize}
  \item First item
  \item Second item
  \item Third item with sub-items:
    \begin{itemize}
      \item Sub-item 1
      \item Sub-item 2
    \end{itemize}
\end{itemize}

\subsection{Numbered List}
\begin{enumerate}
  \item First step
  \item Second step
  \item Third step
\end{enumerate}

\section{Citations and References}
When you use the citation search feature, references will be automatically managed. For example, you might cite a paper like this: \cite{example2024}. All citations will be automatically added to your \texttt{cite.bib} file.

\section{Conclusion}
This template demonstrates basic \LaTeX{} features including mathematical formulas, tables, and lists. The bibliography system is now set up to work seamlessly with Bibby AI's citation management. Feel free to modify it for your needs!

\bibliographystyle{plain}
\bibliography{cite}

\end{document}\documentclass{article}
\usepackage[utf8]{inputenc}
\usepackage{booktabs}
\usepackage{amsmath}
\usepackage{natbib}

\title{My Document}
\author{Your Name}
\date{\today}

\begin{document}

\maketitle

\section{Introduction}
This is a sample document demonstrating various \LaTeX{} features including mathematical equations, tables, and formatting. You can cite papers using the citation search feature, and they will be automatically added to your \texttt{cite.bib} file.

\section{Mathematical Expressions}
Here are some example equations:

\subsection{Quadratic Formula}
The quadratic formula is given by:
\[x = \frac{-b \pm \sqrt{b^2 - 4ac}}{2a}\]

\subsection{Matrix Example}
A sample matrix:
\[
\begin{pmatrix}
a_{11} & a_{12} & a_{13} \\
a_{21} & a_{22} & a_{23} \\
a_{31} & a_{32} & a_{33}
\end{pmatrix}
\]

\section{Tables}
Below is a sample table showing data:

\begin{table}[h]
\centering
\begin{tabular}{@{}lrrr@{}}
\toprule
\textbf{Category} & \textbf{Value 1} & \textbf{Value 2} & \textbf{Total} \\
\midrule
Item A & 10.5 & 20.3 & 30.8 \\
Item B & 15.2 & 25.1 & 40.3 \\
Item C & 22.1 & 12.7 & 34.8 \\
\bottomrule
\end{tabular}
\caption{Sample Data Table}
\label{tab:sample}
\end{table}

\section{Lists}
\subsection{Bullet Points}
\begin{itemize}
  \item First item
  \item Second item
  \item Third item with sub-items:
    \begin{itemize}
      \item Sub-item 1
      \item Sub-item 2
    \end{itemize}
\end{itemize}

\subsection{Numbered List}
\begin{enumerate}
  \item First step
  \item Second step
  \item Third step
\end{enumerate}

\section{Citations and References}
When you use the citation search feature, references will be automatically managed. For example, you might cite a paper like this: \cite{example2024}. All citations will be automatically added to your \texttt{cite.bib} file.

\section{Conclusion}
This template demonstrates basic \LaTeX{} features including mathematical formulas, tables, and lists. The bibliography system is now set up to work seamlessly with Bibby AI's citation management. Feel free to modify it for your needs!

\bibliographystyle{plain}
\bibliography{cite}

\end{document}\documentclass{article}
\usepackage[utf8]{inputenc}
\usepackage{booktabs}
\usepackage{amsmath}
\usepackage{natbib}

\title{My Document}
\author{Your Name}
\date{\today}

\begin{document}

\maketitle

\section{Introduction}
This is a sample document demonstrating various \LaTeX{} features including mathematical equations, tables, and formatting. You can cite papers using the citation search feature, and they will be automatically added to your \texttt{cite.bib} file.

\section{Mathematical Expressions}
Here are some example equations:

\subsection{Quadratic Formula}
The quadratic formula is given by:
\[x = \frac{-b \pm \sqrt{b^2 - 4ac}}{2a}\]

\subsection{Matrix Example}
A sample matrix:
\[
\begin{pmatrix}
a_{11} & a_{12} & a_{13} \\
a_{21} & a_{22} & a_{23} \\
a_{31} & a_{32} & a_{33}
\end{pmatrix}
\]

\section{Tables}
Below is a sample table showing data:

\begin{table}[h]
\centering
\begin{tabular}{@{}lrrr@{}}
\toprule
\textbf{Category} & \textbf{Value 1} & \textbf{Value 2} & \textbf{Total} \\
\midrule
Item A & 10.5 & 20.3 & 30.8 \\
Item B & 15.2 & 25.1 & 40.3 \\
Item C & 22.1 & 12.7 & 34.8 \\
\bottomrule
\end{tabular}
\caption{Sample Data Table}
\label{tab:sample}
\end{table}

\section{Lists}
\subsection{Bullet Points}
\begin{itemize}
  \item First item
  \item Second item
  \item Third item with sub-items:
    \begin{itemize}
      \item Sub-item 1
      \item Sub-item 2
    \end{itemize}
\end{itemize}

\subsection{Numbered List}
\begin{enumerate}
  \item First step
  \item Second step
  \item Third step
\end{enumerate}

\section{Citations and References}
When you use the citation search feature, references will be automatically managed. For example, you might cite a paper like this: \cite{example2024}. All citations will be automatically added to your \texttt{cite.bib} file.

\section{Conclusion}
This template demonstrates basic \LaTeX{} features including mathematical formulas, tables, and lists. The bibliography system is now set up to work seamlessly with Bibby AI's citation management. Feel free to modify it for your needs!

\bibliographystyle{plain}
\bibliography{cite}

\end{document}\documentclass{article}
\usepackage[utf8]{inputenc}
\usepackage{booktabs}
\usepackage{amsmath}
\usepackage{natbib}

\title{My Document}
\author{Your Name}
\date{\today}

\begin{document}

\maketitle

\section{Introduction}
This is a sample document demonstrating various \LaTeX{} features including mathematical equations, tables, and formatting. You can cite papers using the citation search feature, and they will be automatically added to your \texttt{cite.bib} file.

\section{Mathematical Expressions}
Here are some example equations:

\subsection{Quadratic Formula}
The quadratic formula is given by:
\[x = \frac{-b \pm \sqrt{b^2 - 4ac}}{2a}\]

\subsection{Matrix Example}
A sample matrix:
\[
\begin{pmatrix}
a_{11} & a_{12} & a_{13} \\
a_{21} & a_{22} & a_{23} \\
a_{31} & a_{32} & a_{33}
\end{pmatrix}
\]

\section{Tables}
Below is a sample table showing data:

\begin{table}[h]
\centering
\begin{tabular}{@{}lrrr@{}}
\toprule
\textbf{Category} & \textbf{Value 1} & \textbf{Value 2} & \textbf{Total} \\
\midrule
Item A & 10.5 & 20.3 & 30.8 \\
Item B & 15.2 & 25.1 & 40.3 \\
Item C & 22.1 & 12.7 & 34.8 \\
\bottomrule
\end{tabular}
\caption{Sample Data Table}
\label{tab:sample}
\end{table}

\section{Lists}
\subsection{Bullet Points}
\begin{itemize}
  \item First item
  \item Second item
  \item Third item with sub-items:
    \begin{itemize}
      \item Sub-item 1
      \item Sub-item 2
    \end{itemize}
\end{itemize}

\subsection{Numbered List}
\begin{enumerate}
  \item First step
  \item Second step
  \item Third step
\end{enumerate}

\section{Citations and References}
When you use the citation search feature, references will be automatically managed. For example, you might cite a paper like this: \cite{example2024}. All citations will be automatically added to your \texttt{cite.bib} file.

\section{Conclusion}
This template demonstrates basic \LaTeX{} features including mathematical formulas, tables, and lists. The bibliography system is now set up to work seamlessly with Bibby AI's citation management. Feel free to modify it for your needs!

\bibliographystyle{plain}
\bibliography{cite}

\end{document}\documentclass{article}
\usepackage[utf8]{inputenc}
\usepackage{booktabs}
\usepackage{amsmath}
\usepackage{natbib}

\title{My Document}
\author{Your Name}
\date{\today}

\begin{document}

\maketitle

\section{Introduction}
This is a sample document demonstrating various \LaTeX{} features including mathematical equations, tables, and formatting. You can cite papers using the citation search feature, and they will be automatically added to your \texttt{cite.bib} file.

\section{Mathematical Expressions}
Here are some example equations:

\subsection{Quadratic Formula}
The quadratic formula is given by:
\[x = \frac{-b \pm \sqrt{b^2 - 4ac}}{2a}\]

\subsection{Matrix Example}
A sample matrix:
\[
\begin{pmatrix}
a_{11} & a_{12} & a_{13} \\
a_{21} & a_{22} & a_{23} \\
a_{31} & a_{32} & a_{33}
\end{pmatrix}
\]

\section{Tables}
Below is a sample table showing data:

\begin{table}[h]
\centering
\begin{tabular}{@{}lrrr@{}}
\toprule
\textbf{Category} & \textbf{Value 1} & \textbf{Value 2} & \textbf{Total} \\
\midrule
Item A & 10.5 & 20.3 & 30.8 \\
Item B & 15.2 & 25.1 & 40.3 \\
Item C & 22.1 & 12.7 & 34.8 \\
\bottomrule
\end{tabular}
\caption{Sample Data Table}
\label{tab:sample}
\end{table}

\section{Lists}
\subsection{Bullet Points}
\begin{itemize}
  \item First item
  \item Second item
  \item Third item with sub-items:
    \begin{itemize}
      \item Sub-item 1
      \item Sub-item 2
    \end{itemize}
\end{itemize}

\subsection{Numbered List}
\begin{enumerate}
  \item First step
  \item Second step
  \item Third step
\end{enumerate}

\section{Citations and References}
When you use the citation search feature, references will be automatically managed. For example, you might cite a paper like this: \cite{example2024}. All citations will be automatically added to your \texttt{cite.bib} file.

\section{Conclusion}
This template demonstrates basic \LaTeX{} features including mathematical formulas, tables, and lists. The bibliography system is now set up to work seamlessly with Bibby AI's citation management. Feel free to modify it for your needs!

\bibliographystyle{plain}
\bibliography{cite}

\end{document}\documentclass{article}
\usepackage[utf8]{inputenc}
\usepackage{booktabs}
\usepackage{amsmath}
\usepackage{natbib}

\title{My Document}
\author{Your Name}
\date{\today}

\begin{document}

\maketitle

\section{Introduction}
This is a sample document demonstrating various \LaTeX{} features including mathematical equations, tables, and formatting. You can cite papers using the citation search feature, and they will be automatically added to your \texttt{cite.bib} file.

\section{Mathematical Expressions}
Here are some example equations:

\subsection{Quadratic Formula}
The quadratic formula is given by:
\[x = \frac{-b \pm \sqrt{b^2 - 4ac}}{2a}\]

\subsection{Matrix Example}
A sample matrix:
\[
\begin{pmatrix}
a_{11} & a_{12} & a_{13} \\
a_{21} & a_{22} & a_{23} \\
a_{31} & a_{32} & a_{33}
\end{pmatrix}
\]

\section{Tables}
Below is a sample table showing data:

\begin{table}[h]
\centering
\begin{tabular}{@{}lrrr@{}}
\toprule
\textbf{Category} & \textbf{Value 1} & \textbf{Value 2} & \textbf{Total} \\
\midrule
Item A & 10.5 & 20.3 & 30.8 \\
Item B & 15.2 & 25.1 & 40.3 \\
Item C & 22.1 & 12.7 & 34.8 \\
\bottomrule
\end{tabular}
\caption{Sample Data Table}
\label{tab:sample}
\end{table}

\section{Lists}
\subsection{Bullet Points}
\begin{itemize}
  \item First item
  \item Second item
  \item Third item with sub-items:
    \begin{itemize}
      \item Sub-item 1
      \item Sub-item 2
    \end{itemize}
\end{itemize}

\subsection{Numbered List}
\begin{enumerate}
  \item First step
  \item Second step
  \item Third step
\end{enumerate}

\section{Citations and References}
When you use the citation search feature, references will be automatically managed. For example, you might cite a paper like this: \cite{example2024}. All citations will be automatically added to your \texttt{cite.bib} file.

\section{Conclusion}
This template demonstrates basic \LaTeX{} features including mathematical formulas, tables, and lists. The bibliography system is now set up to work seamlessly with Bibby AI's citation management. Feel free to modify it for your needs!

\bibliographystyle{plain}
\bibliography{cite}

\end{document}\documentclass{article}
\usepackage[utf8]{inputenc}
\usepackage{booktabs}
\usepackage{amsmath}
\usepackage{natbib}

\title{My Document}
\author{Your Name}
\date{\today}

\begin{document}

\maketitle

\section{Introduction}
This is a sample document demonstrating various \LaTeX{} features including mathematical equations, tables, and formatting. You can cite papers using the citation search feature, and they will be automatically added to your \texttt{cite.bib} file.

\section{Mathematical Expressions}
Here are some example equations:

\subsection{Quadratic Formula}
The quadratic formula is given by:
\[x = \frac{-b \pm \sqrt{b^2 - 4ac}}{2a}\]

\subsection{Matrix Example}
A sample matrix:
\[
\begin{pmatrix}
a_{11} & a_{12} & a_{13} \\
a_{21} & a_{22} & a_{23} \\
a_{31} & a_{32} & a_{33}
\end{pmatrix}
\]

\section{Tables}
Below is a sample table showing data:

\begin{table}[h]
\centering
\begin{tabular}{@{}lrrr@{}}
\toprule
\textbf{Category} & \textbf{Value 1} & \textbf{Value 2} & \textbf{Total} \\
\midrule
Item A & 10.5 & 20.3 & 30.8 \\
Item B & 15.2 & 25.1 & 40.3 \\
Item C & 22.1 & 12.7 & 34.8 \\
\bottomrule
\end{tabular}
\caption{Sample Data Table}
\label{tab:sample}
\end{table}

\section{Lists}
\subsection{Bullet Points}
\begin{itemize}
  \item First item
  \item Second item
  \item Third item with sub-items:
    \begin{itemize}
      \item Sub-item 1
      \item Sub-item 2
    \end{itemize}
\end{itemize}

\subsection{Numbered List}
\begin{enumerate}
  \item First step
  \item Second step
  \item Third step
\end{enumerate}

\section{Citations and References}
When you use the citation search feature, references will be automatically managed. For example, you might cite a paper like this: \cite{example2024}. All citations will be automatically added to your \texttt{cite.bib} file.

\section{Conclusion}
This template demonstrates basic \LaTeX{} features including mathematical formulas, tables, and lists. The bibliography system is now set up to work seamlessly with Bibby AI's citation management. Feel free to modify it for your needs!

\bibliographystyle{plain}
\bibliography{cite}

\end{document}\documentclass{article}
\usepackage[utf8]{inputenc}
\usepackage{booktabs}
\usepackage{amsmath}
\usepackage{natbib}

\title{My Document}
\author{Your Name}
\date{\today}

\begin{document}

\maketitle

\section{Introduction}
This is a sample document demonstrating various \LaTeX{} features including mathematical equations, tables, and formatting. You can cite papers using the citation search feature, and they will be automatically added to your \texttt{cite.bib} file.

\section{Mathematical Expressions}
Here are some example equations:

\subsection{Quadratic Formula}
The quadratic formula is given by:
\[x = \frac{-b \pm \sqrt{b^2 - 4ac}}{2a}\]

\subsection{Matrix Example}
A sample matrix:
\[
\begin{pmatrix}
a_{11} & a_{12} & a_{13} \\
a_{21} & a_{22} & a_{23} \\
a_{31} & a_{32} & a_{33}
\end{pmatrix}
\]

\section{Tables}
Below is a sample table showing data:

\begin{table}[h]
\centering
\begin{tabular}{@{}lrrr@{}}
\toprule
\textbf{Category} & \textbf{Value 1} & \textbf{Value 2} & \textbf{Total} \\
\midrule
Item A & 10.5 & 20.3 & 30.8 \\
Item B & 15.2 & 25.1 & 40.3 \\
Item C & 22.1 & 12.7 & 34.8 \\
\bottomrule
\end{tabular}
\caption{Sample Data Table}
\label{tab:sample}
\end{table}

\section{Lists}
\subsection{Bullet Points}
\begin{itemize}
  \item First item
  \item Second item
  \item Third item with sub-items:
    \begin{itemize}
      \item Sub-item 1
      \item Sub-item 2
    \end{itemize}
\end{itemize}

\subsection{Numbered List}
\begin{enumerate}
  \item First step
  \item Second step
  \item Third step
\end{enumerate}

\section{Citations and References}
When you use the citation search feature, references will be automatically managed. For example, you might cite a paper like this: \cite{example2024}. All citations will be automatically added to your \texttt{cite.bib} file.

\section{Conclusion}
This template demonstrates basic \LaTeX{} features including mathematical formulas, tables, and lists. The bibliography system is now set up to work seamlessly with Bibby AI's citation management. Feel free to modify it for your needs!

\bibliographystyle{plain}
\bibliography{cite}

\end{document}\documentclass{article}
\usepackage[utf8]{inputenc}
\usepackage{booktabs}
\usepackage{amsmath}
\usepackage{natbib}

\title{My Document}
\author{Your Name}
\date{\today}

\begin{document}

\maketitle

\section{Introduction}
This is a sample document demonstrating various \LaTeX{} features including mathematical equations, tables, and formatting. You can cite papers using the citation search feature, and they will be automatically added to your \texttt{cite.bib} file.

\section{Mathematical Expressions}
Here are some example equations:

\subsection{Quadratic Formula}
The quadratic formula is given by:
\[x = \frac{-b \pm \sqrt{b^2 - 4ac}}{2a}\]

\subsection{Matrix Example}
A sample matrix:
\[
\begin{pmatrix}
a_{11} & a_{12} & a_{13} \\
a_{21} & a_{22} & a_{23} \\
a_{31} & a_{32} & a_{33}
\end{pmatrix}
\]

\section{Tables}
Below is a sample table showing data:

\begin{table}[h]
\centering
\begin{tabular}{@{}lrrr@{}}
\toprule
\textbf{Category} & \textbf{Value 1} & \textbf{Value 2} & \textbf{Total} \\
\midrule
Item A & 10.5 & 20.3 & 30.8 \\
Item B & 15.2 & 25.1 & 40.3 \\
Item C & 22.1 & 12.7 & 34.8 \\
\bottomrule
\end{tabular}
\caption{Sample Data Table}
\label{tab:sample}
\end{table}

\section{Lists}
\subsection{Bullet Points}
\begin{itemize}
  \item First item
  \item Second item
  \item Third item with sub-items:
    \begin{itemize}
      \item Sub-item 1
      \item Sub-item 2
    \end{itemize}
\end{itemize}

\subsection{Numbered List}
\begin{enumerate}
  \item First step
  \item Second step
  \item Third step
\end{enumerate}

\section{Citations and References}
When you use the citation search feature, references will be automatically managed. For example, you might cite a paper like this: \cite{example2024}. All citations will be automatically added to your \texttt{cite.bib} file.

\section{Conclusion}
This template demonstrates basic \LaTeX{} features including mathematical formulas, tables, and lists. The bibliography system is now set up to work seamlessly with Bibby AI's citation management. Feel free to modify it for your needs!

\bibliographystyle{plain}
\bibliography{cite}

\end{document}\documentclass{article}
\usepackage[utf8]{inputenc}
\usepackage{booktabs}
\usepackage{amsmath}
\usepackage{natbib}

\title{My Document}
\author{Your Name}
\date{\today}

\begin{document}

\maketitle

\section{Introduction}
This is a sample document demonstrating various \LaTeX{} features including mathematical equations, tables, and formatting. You can cite papers using the citation search feature, and they will be automatically added to your \texttt{cite.bib} file.

\section{Mathematical Expressions}
Here are some example equations:

\subsection{Quadratic Formula}
The quadratic formula is given by:
\[x = \frac{-b \pm \sqrt{b^2 - 4ac}}{2a}\]

\subsection{Matrix Example}
A sample matrix:
\[
\begin{pmatrix}
a_{11} & a_{12} & a_{13} \\
a_{21} & a_{22} & a_{23} \\
a_{31} & a_{32} & a_{33}
\end{pmatrix}
\]

\section{Tables}
Below is a sample table showing data:

\begin{table}[h]
\centering
\begin{tabular}{@{}lrrr@{}}
\toprule
\textbf{Category} & \textbf{Value 1} & \textbf{Value 2} & \textbf{Total} \\
\midrule
Item A & 10.5 & 20.3 & 30.8 \\
Item B & 15.2 & 25.1 & 40.3 \\
Item C & 22.1 & 12.7 & 34.8 \\
\bottomrule
\end{tabular}
\caption{Sample Data Table}
\label{tab:sample}
\end{table}

\section{Lists}
\subsection{Bullet Points}
\begin{itemize}
  \item First item
  \item Second item
  \item Third item with sub-items:
    \begin{itemize}
      \item Sub-item 1
      \item Sub-item 2
    \end{itemize}
\end{itemize}

\subsection{Numbered List}
\begin{enumerate}
  \item First step
  \item Second step
  \item Third step
\end{enumerate}

\section{Citations and References}
When you use the citation search feature, references will be automatically managed. For example, you might cite a paper like this: \cite{example2024}. All citations will be automatically added to your \texttt{cite.bib} file.

\section{Conclusion}
This template demonstrates basic \LaTeX{} features including mathematical formulas, tables, and lists. The bibliography system is now set up to work seamlessly with Bibby AI's citation management. Feel free to modify it for your needs!

\bibliographystyle{plain}
\bibliography{cite}

\end{document}\documentclass{article}
\usepackage[utf8]{inputenc}
\usepackage{booktabs}
\usepackage{amsmath}
\usepackage{natbib}

\title{My Document}
\author{Your Name}
\date{\today}

\begin{document}

\maketitle

\section{Introduction}
This is a sample document demonstrating various \LaTeX{} features including mathematical equations, tables, and formatting. You can cite papers using the citation search feature, and they will be automatically added to your \texttt{cite.bib} file.

\section{Mathematical Expressions}
Here are some example equations:

\subsection{Quadratic Formula}
The quadratic formula is given by:
\[x = \frac{-b \pm \sqrt{b^2 - 4ac}}{2a}\]

\subsection{Matrix Example}
A sample matrix:
\[
\begin{pmatrix}
a_{11} & a_{12} & a_{13} \\
a_{21} & a_{22} & a_{23} \\
a_{31} & a_{32} & a_{33}
\end{pmatrix}
\]

\section{Tables}
Below is a sample table showing data:

\begin{table}[h]
\centering
\begin{tabular}{@{}lrrr@{}}
\toprule
\textbf{Category} & \textbf{Value 1} & \textbf{Value 2} & \textbf{Total} \\
\midrule
Item A & 10.5 & 20.3 & 30.8 \\
Item B & 15.2 & 25.1 & 40.3 \\
Item C & 22.1 & 12.7 & 34.8 \\
\bottomrule
\end{tabular}
\caption{Sample Data Table}
\label{tab:sample}
\end{table}

\section{Lists}
\subsection{Bullet Points}
\begin{itemize}
  \item First item
  \item Second item
  \item Third item with sub-items:
    \begin{itemize}
      \item Sub-item 1
      \item Sub-item 2
    \end{itemize}
\end{itemize}

\subsection{Numbered List}
\begin{enumerate}
  \item First step
  \item Second step
  \item Third step
\end{enumerate}

\section{Citations and References}
When you use the citation search feature, references will be automatically managed. For example, you might cite a paper like this: \cite{example2024}. All citations will be automatically added to your \texttt{cite.bib} file.

\section{Conclusion}
This template demonstrates basic \LaTeX{} features including mathematical formulas, tables, and lists. The bibliography system is now set up to work seamlessly with Bibby AI's citation management. Feel free to modify it for your needs!

\bibliographystyle{plain}
\bibliography{cite}

\end{document}\documentclass{article}
\usepackage[utf8]{inputenc}
\usepackage{booktabs}
\usepackage{amsmath}
\usepackage{natbib}

\title{My Document}
\author{Your Name}
\date{\today}

\begin{document}

\maketitle

\section{Introduction}
This is a sample document demonstrating various \LaTeX{} features including mathematical equations, tables, and formatting. You can cite papers using the citation search feature, and they will be automatically added to your \texttt{cite.bib} file.

\section{Mathematical Expressions}
Here are some example equations:

\subsection{Quadratic Formula}
The quadratic formula is given by:
\[x = \frac{-b \pm \sqrt{b^2 - 4ac}}{2a}\]

\subsection{Matrix Example}
A sample matrix:
\[
\begin{pmatrix}
a_{11} & a_{12} & a_{13} \\
a_{21} & a_{22} & a_{23} \\
a_{31} & a_{32} & a_{33}
\end{pmatrix}
\]

\section{Tables}
Below is a sample table showing data:

\begin{table}[h]
\centering
\begin{tabular}{@{}lrrr@{}}
\toprule
\textbf{Category} & \textbf{Value 1} & \textbf{Value 2} & \textbf{Total} \\
\midrule
Item A & 10.5 & 20.3 & 30.8 \\
Item B & 15.2 & 25.1 & 40.3 \\
Item C & 22.1 & 12.7 & 34.8 \\
\bottomrule
\end{tabular}
\caption{Sample Data Table}
\label{tab:sample}
\end{table}

\section{Lists}
\subsection{Bullet Points}
\begin{itemize}
  \item First item
  \item Second item
  \item Third item with sub-items:
    \begin{itemize}
      \item Sub-item 1
      \item Sub-item 2
    \end{itemize}
\end{itemize}

\subsection{Numbered List}
\begin{enumerate}
  \item First step
  \item Second step
  \item Third step
\end{enumerate}

\section{Citations and References}
When you use the citation search feature, references will be automatically managed. For example, you might cite a paper like this: \cite{example2024}. All citations will be automatically added to your \texttt{cite.bib} file.

\section{Conclusion}
This template demonstrates basic \LaTeX{} features including mathematical formulas, tables, and lists. The bibliography system is now set up to work seamlessly with Bibby AI's citation management. Feel free to modify it for your needs!

\bibliographystyle{plain}
\bibliography{cite}

\end{document}\documentclass{article}
\usepackage[utf8]{inputenc}
\usepackage{booktabs}
\usepackage{amsmath}
\usepackage{natbib}

\title{My Document}
\author{Your Name}
\date{\today}

\begin{document}

\maketitle

\section{Introduction}
This is a sample document demonstrating various \LaTeX{} features including mathematical equations, tables, and formatting. You can cite papers using the citation search feature, and they will be automatically added to your \texttt{cite.bib} file.

\section{Mathematical Expressions}
Here are some example equations:

\subsection{Quadratic Formula}
The quadratic formula is given by:
\[x = \frac{-b \pm \sqrt{b^2 - 4ac}}{2a}\]

\subsection{Matrix Example}
A sample matrix:
\[
\begin{pmatrix}
a_{11} & a_{12} & a_{13} \\
a_{21} & a_{22} & a_{23} \\
a_{31} & a_{32} & a_{33}
\end{pmatrix}
\]

\section{Tables}
Below is a sample table showing data:

\begin{table}[h]
\centering
\begin{tabular}{@{}lrrr@{}}
\toprule
\textbf{Category} & \textbf{Value 1} & \textbf{Value 2} & \textbf{Total} \\
\midrule
Item A & 10.5 & 20.3 & 30.8 \\
Item B & 15.2 & 25.1 & 40.3 \\
Item C & 22.1 & 12.7 & 34.8 \\
\bottomrule
\end{tabular}
\caption{Sample Data Table}
\label{tab:sample}
\end{table}

\section{Lists}
\subsection{Bullet Points}
\begin{itemize}
  \item First item
  \item Second item
  \item Third item with sub-items:
    \begin{itemize}
      \item Sub-item 1
      \item Sub-item 2
    \end{itemize}
\end{itemize}

\subsection{Numbered List}
\begin{enumerate}
  \item First step
  \item Second step
  \item Third step
\end{enumerate}

\section{Citations and References}
When you use the citation search feature, references will be automatically managed. For example, you might cite a paper like this: \cite{example2024}. All citations will be automatically added to your \texttt{cite.bib} file.

\section{Conclusion}
This template demonstrates basic \LaTeX{} features including mathematical formulas, tables, and lists. The bibliography system is now set up to work seamlessly with Bibby AI's citation management. Feel free to modify it for your needs!

\bibliographystyle{plain}
\bibliography{cite}

\end{document}\documentclass{article}
\usepackage[utf8]{inputenc}
\usepackage{booktabs}
\usepackage{amsmath}
\usepackage{natbib}

\title{My Document}
\author{Your Name}
\date{\today}

\begin{document}

\maketitle

\section{Introduction}
This is a sample document demonstrating various \LaTeX{} features including mathematical equations, tables, and formatting. You can cite papers using the citation search feature, and they will be automatically added to your \texttt{cite.bib} file.

\section{Mathematical Expressions}
Here are some example equations:

\subsection{Quadratic Formula}
The quadratic formula is given by:
\[x = \frac{-b \pm \sqrt{b^2 - 4ac}}{2a}\]

\subsection{Matrix Example}
A sample matrix:
\[
\begin{pmatrix}
a_{11} & a_{12} & a_{13} \\
a_{21} & a_{22} & a_{23} \\
a_{31} & a_{32} & a_{33}
\end{pmatrix}
\]

\section{Tables}
Below is a sample table showing data:

\begin{table}[h]
\centering
\begin{tabular}{@{}lrrr@{}}
\toprule
\textbf{Category} & \textbf{Value 1} & \textbf{Value 2} & \textbf{Total} \\
\midrule
Item A & 10.5 & 20.3 & 30.8 \\
Item B & 15.2 & 25.1 & 40.3 \\
Item C & 22.1 & 12.7 & 34.8 \\
\bottomrule
\end{tabular}
\caption{Sample Data Table}
\label{tab:sample}
\end{table}

\section{Lists}
\subsection{Bullet Points}
\begin{itemize}
  \item First item
  \item Second item
  \item Third item with sub-items:
    \begin{itemize}
      \item Sub-item 1
      \item Sub-item 2
    \end{itemize}
\end{itemize}

\subsection{Numbered List}
\begin{enumerate}
  \item First step
  \item Second step
  \item Third step
\end{enumerate}

\section{Citations and References}
When you use the citation search feature, references will be automatically managed. For example, you might cite a paper like this: \cite{example2024}. All citations will be automatically added to your \texttt{cite.bib} file.

\section{Conclusion}
This template demonstrates basic \LaTeX{} features including mathematical formulas, tables, and lists. The bibliography system is now set up to work seamlessly with Bibby AI's citation management. Feel free to modify it for your needs!

\bibliographystyle{plain}
\bibliography{cite}

\end{document}\documentclass{article}
\usepackage[utf8]{inputenc}
\usepackage{booktabs}
\usepackage{amsmath}
\usepackage{natbib}

\title{My Document}
\author{Your Name}
\date{\today}

\begin{document}

\maketitle

\section{Introduction}
This is a sample document demonstrating various \LaTeX{} features including mathematical equations, tables, and formatting. You can cite papers using the citation search feature, and they will be automatically added to your \texttt{cite.bib} file.

\section{Mathematical Expressions}
Here are some example equations:

\subsection{Quadratic Formula}
The quadratic formula is given by:
\[x = \frac{-b \pm \sqrt{b^2 - 4ac}}{2a}\]

\subsection{Matrix Example}
A sample matrix:
\[
\begin{pmatrix}
a_{11} & a_{12} & a_{13} \\
a_{21} & a_{22} & a_{23} \\
a_{31} & a_{32} & a_{33}
\end{pmatrix}
\]

\section{Tables}
Below is a sample table showing data:

\begin{table}[h]
\centering
\begin{tabular}{@{}lrrr@{}}
\toprule
\textbf{Category} & \textbf{Value 1} & \textbf{Value 2} & \textbf{Total} \\
\midrule
Item A & 10.5 & 20.3 & 30.8 \\
Item B & 15.2 & 25.1 & 40.3 \\
Item C & 22.1 & 12.7 & 34.8 \\
\bottomrule
\end{tabular}
\caption{Sample Data Table}
\label{tab:sample}
\end{table}

\section{Lists}
\subsection{Bullet Points}
\begin{itemize}
  \item First item
  \item Second item
  \item Third item with sub-items:
    \begin{itemize}
      \item Sub-item 1
      \item Sub-item 2
    \end{itemize}
\end{itemize}

\subsection{Numbered List}
\begin{enumerate}
  \item First step
  \item Second step
  \item Third step
\end{enumerate}

\section{Citations and References}
When you use the citation search feature, references will be automatically managed. For example, you might cite a paper like this: \cite{example2024}. All citations will be automatically added to your \texttt{cite.bib} file.

\section{Conclusion}
This template demonstrates basic \LaTeX{} features including mathematical formulas, tables, and lists. The bibliography system is now set up to work seamlessly with Bibby AI's citation management. Feel free to modify it for your needs!

\bibliographystyle{plain}
\bibliography{cite}

\end{document}\documentclass{article}
\usepackage[utf8]{inputenc}
\usepackage{booktabs}
\usepackage{amsmath}
\usepackage{natbib}

\title{My Document}
\author{Your Name}
\date{\today}

\begin{document}

\maketitle

\section{Introduction}
This is a sample document demonstrating various \LaTeX{} features including mathematical equations, tables, and formatting. You can cite papers using the citation search feature, and they will be automatically added to your \texttt{cite.bib} file.

\section{Mathematical Expressions}
Here are some example equations:

\subsection{Quadratic Formula}
The quadratic formula is given by:
\[x = \frac{-b \pm \sqrt{b^2 - 4ac}}{2a}\]

\subsection{Matrix Example}
A sample matrix:
\[
\begin{pmatrix}
a_{11} & a_{12} & a_{13} \\
a_{21} & a_{22} & a_{23} \\
a_{31} & a_{32} & a_{33}
\end{pmatrix}
\]

\section{Tables}
Below is a sample table showing data:

\begin{table}[h]
\centering
\begin{tabular}{@{}lrrr@{}}
\toprule
\textbf{Category} & \textbf{Value 1} & \textbf{Value 2} & \textbf{Total} \\
\midrule
Item A & 10.5 & 20.3 & 30.8 \\
Item B & 15.2 & 25.1 & 40.3 \\
Item C & 22.1 & 12.7 & 34.8 \\
\bottomrule
\end{tabular}
\caption{Sample Data Table}
\label{tab:sample}
\end{table}

\section{Lists}
\subsection{Bullet Points}
\begin{itemize}
  \item First item
  \item Second item
  \item Third item with sub-items:
    \begin{itemize}
      \item Sub-item 1
      \item Sub-item 2
    \end{itemize}
\end{itemize}

\subsection{Numbered List}
\begin{enumerate}
  \item First step
  \item Second step
  \item Third step
\end{enumerate}

\section{Citations and References}
When you use the citation search feature, references will be automatically managed. For example, you might cite a paper like this: \cite{example2024}. All citations will be automatically added to your \texttt{cite.bib} file.

\section{Conclusion}
This template demonstrates basic \LaTeX{} features including mathematical formulas, tables, and lists. The bibliography system is now set up to work seamlessly with Bibby AI's citation management. Feel free to modify it for your needs!

\bibliographystyle{plain}
\bibliography{cite}

\end{document}\documentclass{article}
\usepackage[utf8]{inputenc}
\usepackage{booktabs}
\usepackage{amsmath}
\usepackage{natbib}

\title{My Document}
\author{Your Name}
\date{\today}

\begin{document}

\maketitle

\section{Introduction}
This is a sample document demonstrating various \LaTeX{} features including mathematical equations, tables, and formatting. You can cite papers using the citation search feature, and they will be automatically added to your \texttt{cite.bib} file.

\section{Mathematical Expressions}
Here are some example equations:

\subsection{Quadratic Formula}
The quadratic formula is given by:
\[x = \frac{-b \pm \sqrt{b^2 - 4ac}}{2a}\]

\subsection{Matrix Example}
A sample matrix:
\[
\begin{pmatrix}
a_{11} & a_{12} & a_{13} \\
a_{21} & a_{22} & a_{23} \\
a_{31} & a_{32} & a_{33}
\end{pmatrix}
\]

\section{Tables}
Below is a sample table showing data:

\begin{table}[h]
\centering
\begin{tabular}{@{}lrrr@{}}
\toprule
\textbf{Category} & \textbf{Value 1} & \textbf{Value 2} & \textbf{Total} \\
\midrule
Item A & 10.5 & 20.3 & 30.8 \\
Item B & 15.2 & 25.1 & 40.3 \\
Item C & 22.1 & 12.7 & 34.8 \\
\bottomrule
\end{tabular}
\caption{Sample Data Table}
\label{tab:sample}
\end{table}

\section{Lists}
\subsection{Bullet Points}
\begin{itemize}
  \item First item
  \item Second item
  \item Third item with sub-items:
    \begin{itemize}
      \item Sub-item 1
      \item Sub-item 2
    \end{itemize}
\end{itemize}

\subsection{Numbered List}
\begin{enumerate}
  \item First step
  \item Second step
  \item Third step
\end{enumerate}

\section{Citations and References}
When you use the citation search feature, references will be automatically managed. For example, you might cite a paper like this: \cite{example2024}. All citations will be automatically added to your \texttt{cite.bib} file.

\section{Conclusion}
This template demonstrates basic \LaTeX{} features including mathematical formulas, tables, and lists. The bibliography system is now set up to work seamlessly with Bibby AI's citation management. Feel free to modify it for your needs!

\bibliographystyle{plain}
\bibliography{cite}

\end{document}\documentclass{article}
\usepackage[utf8]{inputenc}
\usepackage{booktabs}
\usepackage{amsmath}
\usepackage{natbib}

\title{My Document}
\author{Your Name}
\date{\today}

\begin{document}

\maketitle

\section{Introduction}
This is a sample document demonstrating various \LaTeX{} features including mathematical equations, tables, and formatting. You can cite papers using the citation search feature, and they will be automatically added to your \texttt{cite.bib} file.

\section{Mathematical Expressions}
Here are some example equations:

\subsection{Quadratic Formula}
The quadratic formula is given by:
\[x = \frac{-b \pm \sqrt{b^2 - 4ac}}{2a}\]

\subsection{Matrix Example}
A sample matrix:
\[
\begin{pmatrix}
a_{11} & a_{12} & a_{13} \\
a_{21} & a_{22} & a_{23} \\
a_{31} & a_{32} & a_{33}
\end{pmatrix}
\]

\section{Tables}
Below is a sample table showing data:

\begin{table}[h]
\centering
\begin{tabular}{@{}lrrr@{}}
\toprule
\textbf{Category} & \textbf{Value 1} & \textbf{Value 2} & \textbf{Total} \\
\midrule
Item A & 10.5 & 20.3 & 30.8 \\
Item B & 15.2 & 25.1 & 40.3 \\
Item C & 22.1 & 12.7 & 34.8 \\
\bottomrule
\end{tabular}
\caption{Sample Data Table}
\label{tab:sample}
\end{table}

\section{Lists}
\subsection{Bullet Points}
\begin{itemize}
  \item First item
  \item Second item
  \item Third item with sub-items:
    \begin{itemize}
      \item Sub-item 1
      \item Sub-item 2
    \end{itemize}
\end{itemize}

\subsection{Numbered List}
\begin{enumerate}
  \item First step
  \item Second step
  \item Third step
\end{enumerate}

\section{Citations and References}
When you use the citation search feature, references will be automatically managed. For example, you might cite a paper like this: \cite{example2024}. All citations will be automatically added to your \texttt{cite.bib} file.

\section{Conclusion}
This template demonstrates basic \LaTeX{} features including mathematical formulas, tables, and lists. The bibliography system is now set up to work seamlessly with Bibby AI's citation management. Feel free to modify it for your needs!

\bibliographystyle{plain}
\bibliography{cite}

\end{document}\documentclass{article}
\usepackage[utf8]{inputenc}
\usepackage{booktabs}
\usepackage{amsmath}
\usepackage{natbib}

\title{My Document}
\author{Your Name}
\date{\today}

\begin{document}

\maketitle

\section{Introduction}
This is a sample document demonstrating various \LaTeX{} features including mathematical equations, tables, and formatting. You can cite papers using the citation search feature, and they will be automatically added to your \texttt{cite.bib} file.

\section{Mathematical Expressions}
Here are some example equations:

\subsection{Quadratic Formula}
The quadratic formula is given by:
\[x = \frac{-b \pm \sqrt{b^2 - 4ac}}{2a}\]

\subsection{Matrix Example}
A sample matrix:
\[
\begin{pmatrix}
a_{11} & a_{12} & a_{13} \\
a_{21} & a_{22} & a_{23} \\
a_{31} & a_{32} & a_{33}
\end{pmatrix}
\]

\section{Tables}
Below is a sample table showing data:

\begin{table}[h]
\centering
\begin{tabular}{@{}lrrr@{}}
\toprule
\textbf{Category} & \textbf{Value 1} & \textbf{Value 2} & \textbf{Total} \\
\midrule
Item A & 10.5 & 20.3 & 30.8 \\
Item B & 15.2 & 25.1 & 40.3 \\
Item C & 22.1 & 12.7 & 34.8 \\
\bottomrule
\end{tabular}
\caption{Sample Data Table}
\label{tab:sample}
\end{table}

\section{Lists}
\subsection{Bullet Points}
\begin{itemize}
  \item First item
  \item Second item
  \item Third item with sub-items:
    \begin{itemize}
      \item Sub-item 1
      \item Sub-item 2
    \end{itemize}
\end{itemize}

\subsection{Numbered List}
\begin{enumerate}
  \item First step
  \item Second step
  \item Third step
\end{enumerate}

\section{Citations and References}
When you use the citation search feature, references will be automatically managed. For example, you might cite a paper like this: \cite{example2024}. All citations will be automatically added to your \texttt{cite.bib} file.

\section{Conclusion}
This template demonstrates basic \LaTeX{} features including mathematical formulas, tables, and lists. The bibliography system is now set up to work seamlessly with Bibby AI's citation management. Feel free to modify it for your needs!

\bibliographystyle{plain}
\bibliography{cite}

\end{document}\documentclass{article}
\usepackage[utf8]{inputenc}
\usepackage{booktabs}
\usepackage{amsmath}
\usepackage{natbib}

\title{My Document}
\author{Your Name}
\date{\today}

\begin{document}

\maketitle

\section{Introduction}
This is a sample document demonstrating various \LaTeX{} features including mathematical equations, tables, and formatting. You can cite papers using the citation search feature, and they will be automatically added to your \texttt{cite.bib} file.

\section{Mathematical Expressions}
Here are some example equations:

\subsection{Quadratic Formula}
The quadratic formula is given by:
\[x = \frac{-b \pm \sqrt{b^2 - 4ac}}{2a}\]

\subsection{Matrix Example}
A sample matrix:
\[
\begin{pmatrix}
a_{11} & a_{12} & a_{13} \\
a_{21} & a_{22} & a_{23} \\
a_{31} & a_{32} & a_{33}
\end{pmatrix}
\]

\section{Tables}
Below is a sample table showing data:

\begin{table}[h]
\centering
\begin{tabular}{@{}lrrr@{}}
\toprule
\textbf{Category} & \textbf{Value 1} & \textbf{Value 2} & \textbf{Total} \\
\midrule
Item A & 10.5 & 20.3 & 30.8 \\
Item B & 15.2 & 25.1 & 40.3 \\
Item C & 22.1 & 12.7 & 34.8 \\
\bottomrule
\end{tabular}
\caption{Sample Data Table}
\label{tab:sample}
\end{table}

\section{Lists}
\subsection{Bullet Points}
\begin{itemize}
  \item First item
  \item Second item
  \item Third item with sub-items:
    \begin{itemize}
      \item Sub-item 1
      \item Sub-item 2
    \end{itemize}
\end{itemize}

\subsection{Numbered List}
\begin{enumerate}
  \item First step
  \item Second step
  \item Third step
\end{enumerate}

\section{Citations and References}
When you use the citation search feature, references will be automatically managed. For example, you might cite a paper like this: \cite{example2024}. All citations will be automatically added to your \texttt{cite.bib} file.

\section{Conclusion}
This template demonstrates basic \LaTeX{} features including mathematical formulas, tables, and lists. The bibliography system is now set up to work seamlessly with Bibby AI's citation management. Feel free to modify it for your needs!

\bibliographystyle{plain}
\bibliography{cite}

\end{document}\documentclass{article}
\usepackage[utf8]{inputenc}
\usepackage{booktabs}
\usepackage{amsmath}
\usepackage{natbib}

\title{My Document}
\author{Your Name}
\date{\today}

\begin{document}

\maketitle

\section{Introduction}
This is a sample document demonstrating various \LaTeX{} features including mathematical equations, tables, and formatting. You can cite papers using the citation search feature, and they will be automatically added to your \texttt{cite.bib} file.

\section{Mathematical Expressions}
Here are some example equations:

\subsection{Quadratic Formula}
The quadratic formula is given by:
\[x = \frac{-b \pm \sqrt{b^2 - 4ac}}{2a}\]

\subsection{Matrix Example}
A sample matrix:
\[
\begin{pmatrix}
a_{11} & a_{12} & a_{13} \\
a_{21} & a_{22} & a_{23} \\
a_{31} & a_{32} & a_{33}
\end{pmatrix}
\]

\section{Tables}
Below is a sample table showing data:

\begin{table}[h]
\centering
\begin{tabular}{@{}lrrr@{}}
\toprule
\textbf{Category} & \textbf{Value 1} & \textbf{Value 2} & \textbf{Total} \\
\midrule
Item A & 10.5 & 20.3 & 30.8 \\
Item B & 15.2 & 25.1 & 40.3 \\
Item C & 22.1 & 12.7 & 34.8 \\
\bottomrule
\end{tabular}
\caption{Sample Data Table}
\label{tab:sample}
\end{table}

\section{Lists}
\subsection{Bullet Points}
\begin{itemize}
  \item First item
  \item Second item
  \item Third item with sub-items:
    \begin{itemize}
      \item Sub-item 1
      \item Sub-item 2
    \end{itemize}
\end{itemize}

\subsection{Numbered List}
\begin{enumerate}
  \item First step
  \item Second step
  \item Third step
\end{enumerate}

\section{Citations and References}
When you use the citation search feature, references will be automatically managed. For example, you might cite a paper like this: \cite{example2024}. All citations will be automatically added to your \texttt{cite.bib} file.

\section{Conclusion}
This template demonstrates basic \LaTeX{} features including mathematical formulas, tables, and lists. The bibliography system is now set up to work seamlessly with Bibby AI's citation management. Feel free to modify it for your needs!

\bibliographystyle{plain}
\bibliography{cite}

\end{document}\documentclass{article}
\usepackage[utf8]{inputenc}
\usepackage{booktabs}
\usepackage{amsmath}
\usepackage{natbib}

\title{My Document}
\author{Your Name}
\date{\today}

\begin{document}

\maketitle

\section{Introduction}
This is a sample document demonstrating various \LaTeX{} features including mathematical equations, tables, and formatting. You can cite papers using the citation search feature, and they will be automatically added to your \texttt{cite.bib} file.

\section{Mathematical Expressions}
Here are some example equations:

\subsection{Quadratic Formula}
The quadratic formula is given by:
\[x = \frac{-b \pm \sqrt{b^2 - 4ac}}{2a}\]

\subsection{Matrix Example}
A sample matrix:
\[
\begin{pmatrix}
a_{11} & a_{12} & a_{13} \\
a_{21} & a_{22} & a_{23} \\
a_{31} & a_{32} & a_{33}
\end{pmatrix}
\]

\section{Tables}
Below is a sample table showing data:

\begin{table}[h]
\centering
\begin{tabular}{@{}lrrr@{}}
\toprule
\textbf{Category} & \textbf{Value 1} & \textbf{Value 2} & \textbf{Total} \\
\midrule
Item A & 10.5 & 20.3 & 30.8 \\
Item B & 15.2 & 25.1 & 40.3 \\
Item C & 22.1 & 12.7 & 34.8 \\
\bottomrule
\end{tabular}
\caption{Sample Data Table}
\label{tab:sample}
\end{table}

\section{Lists}
\subsection{Bullet Points}
\begin{itemize}
  \item First item
  \item Second item
  \item Third item with sub-items:
    \begin{itemize}
      \item Sub-item 1
      \item Sub-item 2
    \end{itemize}
\end{itemize}

\subsection{Numbered List}
\begin{enumerate}
  \item First step
  \item Second step
  \item Third step
\end{enumerate}

\section{Citations and References}
When you use the citation search feature, references will be automatically managed. For example, you might cite a paper like this: \cite{example2024}. All citations will be automatically added to your \texttt{cite.bib} file.

\section{Conclusion}
This template demonstrates basic \LaTeX{} features including mathematical formulas, tables, and lists. The bibliography system is now set up to work seamlessly with Bibby AI's citation management. Feel free to modify it for your needs!

\bibliographystyle{plain}
\bibliography{cite}

\end{document}\documentclass{article}
\usepackage[utf8]{inputenc}
\usepackage{booktabs}
\usepackage{amsmath}
\usepackage{natbib}

\title{My Document}
\author{Your Name}
\date{\today}

\begin{document}

\maketitle

\section{Introduction}
This is a sample document demonstrating various \LaTeX{} features including mathematical equations, tables, and formatting. You can cite papers using the citation search feature, and they will be automatically added to your \texttt{cite.bib} file.

\section{Mathematical Expressions}
Here are some example equations:

\subsection{Quadratic Formula}
The quadratic formula is given by:
\[x = \frac{-b \pm \sqrt{b^2 - 4ac}}{2a}\]

\subsection{Matrix Example}
A sample matrix:
\[
\begin{pmatrix}
a_{11} & a_{12} & a_{13} \\
a_{21} & a_{22} & a_{23} \\
a_{31} & a_{32} & a_{33}
\end{pmatrix}
\]

\section{Tables}
Below is a sample table showing data:

\begin{table}[h]
\centering
\begin{tabular}{@{}lrrr@{}}
\toprule
\textbf{Category} & \textbf{Value 1} & \textbf{Value 2} & \textbf{Total} \\
\midrule
Item A & 10.5 & 20.3 & 30.8 \\
Item B & 15.2 & 25.1 & 40.3 \\
Item C & 22.1 & 12.7 & 34.8 \\
\bottomrule
\end{tabular}
\caption{Sample Data Table}
\label{tab:sample}
\end{table}

\section{Lists}
\subsection{Bullet Points}
\begin{itemize}
  \item First item
  \item Second item
  \item Third item with sub-items:
    \begin{itemize}
      \item Sub-item 1
      \item Sub-item 2
    \end{itemize}
\end{itemize}

\subsection{Numbered List}
\begin{enumerate}
  \item First step
  \item Second step
  \item Third step
\end{enumerate}

\section{Citations and References}
When you use the citation search feature, references will be automatically managed. For example, you might cite a paper like this: \cite{example2024}. All citations will be automatically added to your \texttt{cite.bib} file.

\section{Conclusion}
This template demonstrates basic \LaTeX{} features including mathematical formulas, tables, and lists. The bibliography system is now set up to work seamlessly with Bibby AI's citation management. Feel free to modify it for your needs!

\bibliographystyle{plain}
\bibliography{cite}

\end{document}\documentclass{article}
\usepackage[utf8]{inputenc}
\usepackage{booktabs}
\usepackage{amsmath}
\usepackage{natbib}

\title{My Document}
\author{Your Name}
\date{\today}

\begin{document}

\maketitle

\section{Introduction}
This is a sample document demonstrating various \LaTeX{} features including mathematical equations, tables, and formatting. You can cite papers using the citation search feature, and they will be automatically added to your \texttt{cite.bib} file.

\section{Mathematical Expressions}
Here are some example equations:

\subsection{Quadratic Formula}
The quadratic formula is given by:
\[x = \frac{-b \pm \sqrt{b^2 - 4ac}}{2a}\]

\subsection{Matrix Example}
A sample matrix:
\[
\begin{pmatrix}
a_{11} & a_{12} & a_{13} \\
a_{21} & a_{22} & a_{23} \\
a_{31} & a_{32} & a_{33}
\end{pmatrix}
\]

\section{Tables}
Below is a sample table showing data:

\begin{table}[h]
\centering
\begin{tabular}{@{}lrrr@{}}
\toprule
\textbf{Category} & \textbf{Value 1} & \textbf{Value 2} & \textbf{Total} \\
\midrule
Item A & 10.5 & 20.3 & 30.8 \\
Item B & 15.2 & 25.1 & 40.3 \\
Item C & 22.1 & 12.7 & 34.8 \\
\bottomrule
\end{tabular}
\caption{Sample Data Table}
\label{tab:sample}
\end{table}

\section{Lists}
\subsection{Bullet Points}
\begin{itemize}
  \item First item
  \item Second item
  \item Third item with sub-items:
    \begin{itemize}
      \item Sub-item 1
      \item Sub-item 2
    \end{itemize}
\end{itemize}

\subsection{Numbered List}
\begin{enumerate}
  \item First step
  \item Second step
  \item Third step
\end{enumerate}

\section{Citations and References}
When you use the citation search feature, references will be automatically managed. For example, you might cite a paper like this: \cite{example2024}. All citations will be automatically added to your \texttt{cite.bib} file.

\section{Conclusion}
This template demonstrates basic \LaTeX{} features including mathematical formulas, tables, and lists. The bibliography system is now set up to work seamlessly with Bibby AI's citation management. Feel free to modify it for your needs!

\bibliographystyle{plain}
\bibliography{cite}

\end{document}